\documentclass{book}
%\usepackage[letterpaper, portrait, margin=1cm]{geometry}
%\usepackage[letterpaper, bindingoffset=0.2in, left=1in,right=1in,top=.5in,bottom=.5in,footskip=.25in,marginparwidth=5em]{geometry}
\usepackage[letterpaper, left=1in,right=1in,top=.5in,bottom=.5in,footskip=.25in,marginparwidth=1cm]{geometry}
% ---------------------
% mini-table-of-content
% ---------------------
\usepackage{minitoc}
\setcounter{minitocdepth}{1}
\setlength{\mtcindent}{24pt}
\setcounter{secnumdepth}{-2}
%\renewcommand{\mtcfont}{\small\rm}
%\renewcommand{\mtcSfont}{\small\bf}
%\usepackage{setspace}
%\usepackage{tocloft}
%\setlength\cftparskip{-1.2pt}
%\setlength\cftbeforesecskip{1.3pt}
%\setlength\cftaftertoctitleskip{2pt}
%\renewcommand{\cftsecafterpnum}{\hspace*{02.0em}}
%\renewcommand{\cftsubsecafterpnum}{\hspace*{02.0em}}

% ---------------------------
% Chinese Characters Packages
% ---------------------------
\usepackage{fontspec} 
\usepackage{xeCJK}
\setmainfont{Times}
\setCJKmainfont{BiauKai}
\newfontfamily\sblgoodhebrew{SBL BibLit}[Script=Hebrew,Contextuals=Alternate]
\newfontfamily\sblgoodgreek{SBL BibLit}[Script=Greek,Contextuals=Alternate]

\usepackage{ifpdf,cite,algorithmic,url,tikz}
\usepackage[cmex10]{amsmath}

% ---------------------------
% Hebrew Characters Packages
% ---------------------------
\usepackage{polyglossia}
\setmainfont{Times New Roman}

% -------
% General
% -------
\usepackage{multicol}
\usepackage{multirow}
\usepackage{color,colortbl}
\usepackage{xparse}
\usepackage{pbox}
\usepackage{stackengine}
\usepackage{titlesec}% http://ctan.org/pkg/titlesec
\usepackage{tabularx}
\usepackage{xltabular}
\usepackage{titlesec}
\usepackage{makecell}
\newcommand{\sectionbreak}{\clearpage}

\author{
  Editor, Michael Chan\\
  \texttt{michaelchan\_wahyan@yahoo.com.hk}
}
\usepackage{tocloft}

\usepackage{hyperref}
\hypersetup{
    colorlinks=true, % set true if you want colored links
    linktoc   =all , % set to all if you want both sections and subsections linked
    linkcolor =blue, % choose some color if you want links to stand out
}

% ----------
% Afterword
% ----------
\usepackage{marginnote}
\usepackage{sectsty}
\usepackage{ragged2e}
\usepackage{lineno}
\usepackage{xcolor}
\usepackage{paracol}

\begin{document}

\clearpage
%% temporary titles
% command to provide stretchy vertical space in proportion
\newcommand\nbvspace[1][3]{\vspace*{\stretch{#1}}}
% allow some slack to avoid under/overfull boxes
\newcommand\nbstretchyspace{\spaceskip0.5em plus 0.25em minus 0.25em}
% To improve spacing on titlepages
\newcommand{\nbtitlestretch}{\spaceskip0.6em}
\pagestyle{empty}
\begin{center}
\bfseries
\nbvspace[1]
\Huge
{%\nbtitlestretch
\Large
\textbf{宣道會洪恩堂 粵語講道逐字稿 2020-present \\
       Youtube Channel: Graceflow Church
       }}

\nbvspace[1]

{\large
Editor: Michael\\
\texttt{michaelchan\_wahyan@yahoo.com.hk}
}

\nbvspace[1]

{\large
Revision: \texttt{v1.3}\\
Last Update: \today
}


\vfill
\begin{tikzpicture}
    %remove comment for OT cover%\node (0,0) [opacity=0.03]{\includegraphics[width=15cm]{../bible_out/ot_frontcover.png}} ;
    %remove comment for NT cover%\node (0,0) [opacity=0.03]{\includegraphics[width=15cm]{../bible_out/christ_on_cross.png}} ;
    %remove comment for Bible cover%\node (0,0) [xshift=0.8cm, yshift=+2cm, opacity=0.03]{\includegraphics[width=10cm]{./christ_on_cross.png}} ;
    %remove comment for Bible cover%\node (0,0) [              yshift=-2cm, opacity=0.03]{\includegraphics[width=14cm]{./ot_frontcover.png}} ;
\end{tikzpicture}
\vfill

\end{center}

\newpage

\setcounter{tocdepth}{0}
\dominitoc
\begin{multicols}{3}
\addtocontents{toc}{\protect\hypertarget{toc}{}}
\tableofcontents
\end{multicols}

\large
%\twocolumn

% the color definition syntax is as follow:
% \definecolor{name}{system}{definition}
% example: a mono-channel color can be defined as
%          \definecolor{Gray}{gray}{0.9}
% example: an rgb-3-channel color can be defined as
%          \definecolor{LightCyan}{rgb}{0.88,1,1}
%          \definecolor{pink}{rgb}{0.68,0,0.68}

\definecolor{CUV1LightRed}{rgb}{1,0.75,0.75}     % for CUV1
\definecolor{LZZVLightGray}{rgb}{0.9,0.9,0.9}    % for LZZ
\definecolor{KJVVLightGreen}{rgb}{0.75,1,0.85}   % for KJV
\definecolor{CUV2LightYellow}{rgb}{1,1,0.75}     % for CUV2
\definecolor{CNVVLightBrown}{rgb}{1,0.85,0.7}    % for CNV
\definecolor{NRSVLightBlue}{rgb}{0.75,1,1}       % for NRSV
\definecolor{WENLLightPurple}{rgb}{0.95,0.85,0.9}% for WENL
\definecolor{TCV19PaleGreen}{rgb}{0.85,1,0.95}   % for TCV19
\definecolor{MSGVLightWhite}{rgb}{0.98,0.98,0.98}% for MSGV
\definecolor{NETSLightRed}{rgb}{1,0.75,0.75}     % for NETS
\definecolor{JPS1917LightYellow}{rgb}{1,1,0.75}  % for JPS1917
\definecolor{SBLGNTPaleRed}{rgb}{1,0.85,0.80}    % for SBLGNT

\section{目錄\small{(順時)}}
\label{sec:index_chronic}
{ \scriptsize


\begin{xltabular}{\textwidth}{|p{0.15\textwidth} p{0.6\textwidth}|p{0.07\textwidth} p{0.1\textwidth}|}
\hline
帖撒羅尼迦前書 3:3-5:9 & \hyperref[sec:iY8ed2cNTa4]{2020.08.02《患難是擺在面前的》 帖撒羅尼迦前書 3\_3,5\_9 郭鴻標牧師} & 2020-08-10 & \href{https://youtube.com/watch?v=iY8ed2cNTa4}{\texttt{iY8ed2cNTa4}} \\
哈該書 2:10-19 & \hyperref[sec:N89ZkrKQWlk]{2020.08.09《生命的回轉與盼望》哈該書 2\_10-19( 和修版)  老耀雄傳道} & 2020-08-20 & \href{https://youtube.com/watch?v=N89ZkrKQWlk}{\texttt{N89ZkrKQWlk}} \\
彼得前書 2:9 & \hyperref[sec:UM_spSYKywQ]{2020.08.16《不可以逃避》彼得前書 2\_9 陳志鵬傳道} & 2020-08-20 & \href{https://youtube.com/watch?v=UM-spSYKywQ}{\texttt{UM-spSYKywQ}} \\
哈該書 2:1-9 & \hyperref[sec:xKlgziHaUYU]{2020.08.23《困境中的勇氣》 哈該書 2\_1-9 陳麗蟬牧師} & 2020-09-05 & \href{https://youtube.com/watch?v=xKlgziHaUYU}{\texttt{xKlgziHaUYU}} \\
約伯記 1:1-5 & \hyperref[sec:YQ9vapt0XGE]{2020.08.30《闔府蒙福》 約伯記 1\_1-5 陳一華牧師} & 2020-09-06 & \href{https://youtube.com/watch?v=YQ9vapt0XGE}{\texttt{YQ9vapt0XGE}} \\
彼得後書 1:3-11 & \hyperref[sec:OSwpE_DOZH0]{2020.09.06《信徒長進的人生》 彼得後書 1\_3-11 謝靄薇牧師} & 2020-09-09 & \href{https://youtube.com/watch?v=OSwpE-DOZH0}{\texttt{OSwpE-DOZH0}} \\
士師記 6:11-18 & \hyperref[sec:RqJiqD22j9E]{2020.09.13《從軟弱到剛強》 士師記 6\_11-18 (和修版) 老耀雄傳道} & 2020-09-15 & \href{https://youtube.com/watch?v=RqJiqD22j9E}{\texttt{RqJiqD22j9E}} \\
約伯記 17:15 & \hyperref[sec:7bIUWEs7F1E]{2020.09.27《HOPE ̶l̶e̶s̶s̶》 約伯記 17\_15 朱子溢傳道} & 2020-09-30 & \href{https://youtube.com/watch?v=7bIUWEs7F1E}{\texttt{7bIUWEs7F1E}} \\
哥林多後書 12:7-10 & \hyperref[sec:6c0M0guWu_4]{2020.10.04《軟弱化剛強》 哥林多後書 12\_7-10 陳一華牧師} & 2020-10-06 & \href{https://youtube.com/watch?v=6c0M0guWu-4}{\texttt{6c0M0guWu-4}} \\
詩篇 126:1-6 & \hyperref[sec:6ptt2Cl0aVs]{2020.10.11《流淚撒種者的歡呼》 詩篇 126\_1-6 蔣文忠博士} & 2020-10-14 & \href{https://youtube.com/watch?v=6ptt2Cl0aVs}{\texttt{6ptt2Cl0aVs}} \\
哥林多前書 12:1-31 & \hyperref[sec:YoAst8SyrWc]{2020.10.18《凝聚每分光》 哥林多前書 12\_1-31 陳志鵬傳道} & 2020-10-22 & \href{https://youtube.com/watch?v=YoAst8SyrWc}{\texttt{YoAst8SyrWc}} \\
腓立比書 3:7-14-20-21 & \hyperref[sec:L_6ydn1BUoo]{2020.10.25 講道《我要得甚麼的獎賞》 腓立比書 3\_7-14,20-21( 和修版)  老耀雄傳道} & 2020-10-28 & \href{https://youtube.com/watch?v=L-6ydn1BUoo}{\texttt{L-6ydn1BUoo}} \\
猶大書  & \hyperref[sec:xDCCx3kJCvo]{2020.11.08《在至聖真道上造就自己》 猶大書 一 17-25 ( 和修版) 老耀雄傳道} & 2020-11-11 & \href{https://youtube.com/watch?v=xDCCx3kJCvo}{\texttt{xDCCx3kJCvo}} \\
路加福音 17:11-19 & \hyperref[sec:buLtIY201P0]{2020.11.15《施恩望報》 路加福音 17\_11-19 梁永善牧師} & 2020-11-17 & \href{https://youtube.com/watch?v=buLtIY201P0}{\texttt{buLtIY201P0}} \\
瑪拉基書 2:17-3:6 & \hyperref[sec:UB7GZHuROmQ]{2020.11.22 講道《信心的等候》 瑪拉基書 2\_17-3\_6 謝靄薇牧師} & 2020-11-26 & \href{https://youtube.com/watch?v=UB7GZHuROmQ}{\texttt{UB7GZHuROmQ}} \\
約翰福音 13:1-17 & \hyperref[sec:B0mTZIy27eI]{2020.11.29 講道《活出愛》 約翰福音 13\_1-17 曾敏姑娘} & 2020-12-03 & \href{https://youtube.com/watch?v=B0mTZIy27eI}{\texttt{B0mTZIy27eI}} \\
使徒行傳 8:4-8-26-36-38-40 & \hyperref[sec:pZ9tffANuvA]{2020.12.06《神的用人----腓利的生命素質》 使徒行傳 8\_4-8,26-36,38-40 陳朗欣傳道} & 2020-12-08 & \href{https://youtube.com/watch?v=pZ9tffANuvA}{\texttt{pZ9tffANuvA}} \\
帖撒羅尼迦前書 1:1-10 & \hyperref[sec:BO948tI9cPE]{2020.12.13《合神心意的教會》帖撒羅尼迦前書 1\_1-10( 和修版)  老耀雄傳道} & 2020-12-18 & \href{https://youtube.com/watch?v=BO948tI9cPE}{\texttt{BO948tI9cPE}} \\
路加福音 2:1-14 & \hyperref[sec:7YQ_zudEJ1I]{2020.12.20《因有一嬰孩為我們生》路加福音 2\_1-14( 和修版)  曾敏姑娘} & 2020-12-25 & \href{https://youtube.com/watch?v=7YQ-zudEJ1I}{\texttt{7YQ-zudEJ1I}} \\
詩篇 103:1 & \hyperref[sec:p_pFg5YcRiU]{2020.12.27《心的稱頌》詩篇 103\_1 - 5( 和修版)  陳麗蟬牧師} & 2020-12-27 & \href{https://youtube.com/watch?v=p-pFg5YcRiU}{\texttt{p-pFg5YcRiU}} \\
約翰一書 4:7-12 & \hyperref[sec:r3qTUjnU0eY]{2021.01.03《生命欠一個奇蹟》約翰一書 4\_7-12( 和修版)  老耀雄傳道} & 2021-01-04 & \href{https://youtube.com/watch?v=r3qTUjnU0eY}{\texttt{r3qTUjnU0eY}} \\
創世記 1:1-31 & \hyperref[sec:kzUx2k7DXdQ]{2021.01.10《起初,神創造天地》創世記 1\_1-31( 和修版)  曾敏姑娘} & 2021-01-11 & \href{https://youtube.com/watch?v=kzUx2k7DXdQ}{\texttt{kzUx2k7DXdQ}} \\
提摩太後書 2:1-7 & \hyperref[sec:hPE_EBTYs3o]{2021.01.17《傳下去》提摩太後書 2\_1-7( 和修版)  張美薇牧師} & 2021-01-18 & \href{https://youtube.com/watch?v=hPE_EBTYs3o}{\texttt{hPE\_EBTYs3o}} \\
出埃及記 3:1-15 & \hyperref[sec:sVUIck2trk4]{2021.01.24《再思上主的呼召》出埃及記 3\_1-15( 和修版)  曾裔貴牧師} & 2021-01-26 & \href{https://youtube.com/watch?v=sVUIck2trk4}{\texttt{sVUIck2trk4}} \\
耶利米書 1:1-10 & \hyperref[sec:tlLyIq5teaA]{2021.01.31《亂世中的召命》耶利米書 1\_1-10( 和修版)  楊穎兒導師} & 2021-02-02 & \href{https://youtube.com/watch?v=tlLyIq5teaA}{\texttt{tlLyIq5teaA}} \\
創世記 2:4-25 & \hyperref[sec:xeKLoIITcbQ]{2021.02.07《唯獨你是不可取替》創世記 2\_4-25( 和修版)  曾敏姑娘} & 2021-02-10 & \href{https://youtube.com/watch?v=xeKLoIITcbQ}{\texttt{xeKLoIITcbQ}} \\
約翰福音 6:1-15 & \hyperref[sec:nPhRUxT71P8]{2021.02.14《年年有餘》約翰福音 6\_1-15( 和修版)  老耀雄牧師} & 2021-02-16 & \href{https://youtube.com/watch?v=nPhRUxT71P8}{\texttt{nPhRUxT71P8}} \\
腓立比書 2:12-18 & \hyperref[sec:t7FE7L_V2U0]{2021.02.21《天國的拓展》腓立比書 2\_12-18( 和修版)  黎光偉牧師} & 2021-02-22 & \href{https://youtube.com/watch?v=t7FE7L-V2U0}{\texttt{t7FE7L-V2U0}} \\
約翰福音 13:34-35 & \hyperref[sec:PDRx0v0AuiQ]{2021.02.28《新的命令》約翰福音 13\_34-35( 和修版)  林佩心姑娘} & 2021-02-28 & \href{https://youtube.com/watch?v=PDRx0v0AuiQ}{\texttt{PDRx0v0AuiQ}} \\
撒母耳記上 16:1-13 & \hyperref[sec:JLt4gSCI6Y8]{2021.03.07《神怎樣看我們》撒母耳記上 16\_1-13( 和修版)  老耀雄牧師} & 2021-03-11 & \href{https://youtube.com/watch?v=JLt4gSCI6Y8}{\texttt{JLt4gSCI6Y8}} \\
約翰福音 1:1-28 & \hyperref[sec:2_wlsYE_36Q]{2021.03.14《教會存在於見證中》約翰福音 1\_1-28( 和修版)  陳韋安博士} & 2021-03-18 & \href{https://youtube.com/watch?v=2_wlsYE-36Q}{\texttt{2\_wlsYE-36Q}} \\
以賽亞書 6:1-7 & \hyperref[sec:oH_gH6L8AqQ]{2021.03.21《聖哉!聖哉!聖哉!》以賽亞書 6\_1-7( 和修版)  曾敏姑娘} & 2021-03-25 & \href{https://youtube.com/watch?v=oH_gH6L8AqQ}{\texttt{oH\_gH6L8AqQ}} \\
使徒行傳 10:34-35 & \hyperref[sec:7LBy_VBf8iw]{2021.03.28《誰是主宰》使徒行傳 10\_34-35( 和修版)  鄧泰康傳道} & 2021-03-29 & \href{https://youtube.com/watch?v=7LBy_VBf8iw}{\texttt{7LBy\_VBf8iw}} \\
約翰福音路加福音 23:32-49 & \hyperref[sec:rYjqMvMSdYU]{2021.04.04《我們所做的,我們知道嗎?》路加福音 23\_32-49;約翰福音19\_30( 和修版)  老耀雄牧師} & 2021-04-06 & \href{https://youtube.com/watch?v=rYjqMvMSdYU}{\texttt{rYjqMvMSdYU}} \\
約翰福音 6:7-9 & \hyperref[sec:q2mVvzFJiBo]{2021.04.11《一位平凡但重要的使徒--安德烈》約翰福音 6\_7-9( 和修版)  郭鴻標牧師} & 2021-04-11 & \href{https://youtube.com/watch?v=q2mVvzFJiBo}{\texttt{q2mVvzFJiBo}} \\
羅馬書 5:1-11 & \hyperref[sec:D7R35FEAwH0]{2021.04.18《我們快樂是因為……》羅馬書 5\_1-11( 和修版)  曾敏姑娘} & 2021-04-20 & \href{https://youtube.com/watch?v=D7R35FEAwH0}{\texttt{D7R35FEAwH0}} \\
路加福音 5:17-20-24-26 & \hyperref[sec:NkgRuGpNbys]{2021.04.25《信心的服侍》路加福音 5\_17-20,24-26( 和修版)  楊穎兒導師} & 2021-04-26 & \href{https://youtube.com/watch?v=NkgRuGpNbys}{\texttt{NkgRuGpNbys}} \\
約翰福音 5:1-15 & \hyperref[sec:FhlMWYYtvXU]{2021.05.02《你要痊癒嗎?》約翰福音 5\_1-15( 和修版)  曾敏姑娘} & 2021-05-03 & \href{https://youtube.com/watch?v=FhlMWYYtvXU}{\texttt{FhlMWYYtvXU}} \\
馬太福音 20:20-28 & \hyperref[sec:k0ge2og3_Pw]{2021.05.09《一位母親的要求》馬太福音 20\_20-28( 和修版)  老耀雄牧師} & 2021-05-10 & \href{https://youtube.com/watch?v=k0ge2og3_Pw}{\texttt{k0ge2og3\_Pw}} \\
利未記 25:1-24 & \hyperref[sec:yXf3anXpfOw]{2021.05.16《神悅納人的禧年》 利未記 25\_1-24 ( 和修版)  楊穎兒導師} & 2021-05-17 & \href{https://youtube.com/watch?v=yXf3anXpfOw}{\texttt{yXf3anXpfOw}} \\
帖撒羅尼迦前書 1:1-10 & \hyperref[sec:HojbxuWCThI]{2021.05.19《「逆/疫」是有信,有望,有愛》 帖撒羅尼迦前書 1\_1-10 老耀雄牧師} & 2021-05-19 & \href{https://youtube.com/watch?v=HojbxuWCThI}{\texttt{HojbxuWCThI}} \\
馬可福音 2:18-22 & \hyperref[sec:0Rk0UlP9K50]{2021.05.23《新皮袋》馬可福音 2\_18-22( 和修版)  吳景霖傳道} & 2021-05-25 & \href{https://youtube.com/watch?v=0Rk0UlP9K50}{\texttt{0Rk0UlP9K50}} \\
尼希米記 4:7-14 & \hyperref[sec:JhzI3zqDYAI]{2021.05.30《內外受敵,卻不被困住》 尼希米記 4\_7-14 ( 和修版)  陳一華牧師} & 2021-05-30 & \href{https://youtube.com/watch?v=JhzI3zqDYAI}{\texttt{JhzI3zqDYAI}} \\
路加福音 4:18-19 & \hyperref[sec:11ItqReN_00]{2021.06.06《釋放與捆綁》 路加福音 4\_18-19 ( 和修版)  雷子健牧師} & 2021-06-06 & \href{https://youtube.com/watch?v=11ItqReN_00}{\texttt{11ItqReN\_00}} \\
路加福音 2:36-38 & \hyperref[sec:hNtKC6DPyIM]{2021.06.13《使命人生》 路加福音 2\_36-38 ( 和修版)  老耀雄牧師} & 2021-06-15 & \href{https://youtube.com/watch?v=hNtKC6DPyIM}{\texttt{hNtKC6DPyIM}} \\
詩篇 139:13-18 & \hyperref[sec:T6XL5VlhAWY]{2021.06.20《無可取替的愛》詩篇 139\_13-18( 和修版)  曾敏姑娘} & 2021-06-21 & \href{https://youtube.com/watch?v=T6XL5VlhAWY}{\texttt{T6XL5VlhAWY}} \\
歌羅西書 1:16-32 & \hyperref[sec:M84h3HizR5w]{2021.06.27《不要自欺欺人》 歌羅西書 1\_16-32 ( 和修版)  鄧泰康傳道} & 2021-06-29 & \href{https://youtube.com/watch?v=M84h3HizR5w}{\texttt{M84h3HizR5w}} \\
列王記上 19:1-18 & \hyperref[sec:F4zrvJPl9fU]{2021.07.04《你在這裡做甚麼?》 列王記上 19\_1-18 ( 和修版)  老耀雄牧師} & 2021-07-04 & \href{https://youtube.com/watch?v=F4zrvJPl9fU}{\texttt{F4zrvJPl9fU}} \\
彼得前書 1:1-6-7-22-25 & \hyperref[sec:LNNbZTaQW6E]{2021.07.11《奇妙的道》 彼得前書 1\_1,6-7,22-25 (和修版) 陳一華牧師} & 2021-07-13 & \href{https://youtube.com/watch?v=LNNbZTaQW6E}{\texttt{LNNbZTaQW6E}} \\
腓立比書 2:25-30 & \hyperref[sec:5qNMTaHZqr4]{2021.07.18《隱藏的事奉》 腓立比書 2\_25-30 (和修版) 曾敏姑娘} & 2021-07-20 & \href{https://youtube.com/watch?v=5qNMTaHZqr4}{\texttt{5qNMTaHZqr4}} \\
馬太福音列王記下 5:1-14 & \hyperref[sec:eRsLDyLgpD0]{2021.07.25《活著為主》 列王記下 5\_1-14;馬太福音22\_36-40 (和修版) 高淑芬傳道} & 2021-07-27 & \href{https://youtube.com/watch?v=eRsLDyLgpD0}{\texttt{eRsLDyLgpD0}} \\
創世記 12:1-20 & \hyperref[sec:G34IvTirgOo]{2021.08.01《絕不離地的信仰》 創世記 12\_1-20 (和修版) 老耀雄牧師} & 2021-08-04 & \href{https://youtube.com/watch?v=G34IvTirgOo}{\texttt{G34IvTirgOo}} \\
以弗所書 4:11-16 & \hyperref[sec:_2_W0T2LR08]{2021.08.08《同心合意建立基督的身體》 以弗所書 4\_11-16 (和修版) 曾敏姑娘} & 2021-08-11 & \href{https://youtube.com/watch?v=_2_W0T2LR08}{\texttt{\_2\_W0T2LR08}} \\
列王記下彼得後書 3:9 & \hyperref[sec:mw0QSsEsPEU]{2021.08.15《宣教的異象》 彼得後書 3\_9;列王記下 5\_1-27 (和修版) 高淑芬傳道} & 2021-08-19 & \href{https://youtube.com/watch?v=mw0QSsEsPEU}{\texttt{mw0QSsEsPEU}} \\
使徒行傳 16:6-10 & \hyperref[sec:h6Nm0N8GuaU]{2021.08.22《使命必達》 使徒行傳 16\_6-10 (和修版) 謝靄薇牧師} & 2021-08-25 & \href{https://youtube.com/watch?v=h6Nm0N8GuaU}{\texttt{h6Nm0N8GuaU}} \\
詩篇 150:1-6 & \hyperref[sec:65pDMB41cZg]{2021.08.29《當來讚美神》 詩篇 150\_1-6 (和修版) 張美薇牧師} & 2021-09-01 & \href{https://youtube.com/watch?v=65pDMB41cZg}{\texttt{65pDMB41cZg}} \\
彼得前書 4:7-11 & \hyperref[sec:xUmmDjoz5Ds]{2021.09.05《在苦難中的信仰實踐》 彼得前書 4\_7-11 (和修版) 老耀雄牧師} & 2021-09-07 & \href{https://youtube.com/watch?v=xUmmDjoz5Ds}{\texttt{xUmmDjoz5Ds}} \\
哥林多後書 4:7-15 & \hyperref[sec:oCu8Aak_B2s]{2021.09.12《無懼患難的人生》 哥林多後書 4\_7-15 (和修版) 曾敏姑娘} & 2021-09-16 & \href{https://youtube.com/watch?v=oCu8Aak_B2s}{\texttt{oCu8Aak\_B2s}} \\
馬太福音 12:18-21 & \hyperref[sec:fusOcxn2XoI]{2021.09.19 講道《壓傷的蘆葦》 馬太福音 12\_18-21 (和修版) 王楚男傳道} & 2021-09-23 & \href{https://youtube.com/watch?v=fusOcxn2XoI}{\texttt{fusOcxn2XoI}} \\
路加福音 5:1-11 & \hyperref[sec:_ayvrCyPl7U]{2021.09.26 講道《與神同工》 路加福音 5\_1-11 (和修版) 楊明麗姑娘} & 2021-09-29 & \href{https://youtube.com/watch?v=_ayvrCyPl7U}{\texttt{\_ayvrCyPl7U}} \\
列王記上 2:1-9 & \hyperref[sec:BYuafMUFtfY]{2021.10.03 《活在神的應許中》 列王記上 2\_1-9 (和修版) 老耀雄牧師} & 2021-10-08 & \href{https://youtube.com/watch?v=BYuafMUFtfY}{\texttt{BYuafMUFtfY}} \\
哥林多前書 12:12-27 & \hyperref[sec:P_diqejYp0w]{2021.10.10 《蒙神喜悅的事奉》哥林多前書 12\_12-27 曾敏姑娘} & 2021-10-12 & \href{https://youtube.com/watch?v=P-diqejYp0w}{\texttt{P-diqejYp0w}} \\
希伯來書 12:1-3 & \hyperref[sec:9yhz2ddk1BQ]{2021.10.17 《得力秘訣》 希伯來書 12\_1-3 (和修版) 陳秀賢牧師} & 2021-10-18 & \href{https://youtube.com/watch?v=9yhz2ddk1BQ}{\texttt{9yhz2ddk1BQ}} \\
創世記 5:1-5 & \hyperref[sec:PriyvxgsE30]{2021.10.24 《耶和華眼中的亞當》 創世記 5\_1-5 (新譯本) 楊穎兒傳道} & 2021-10-27 & \href{https://youtube.com/watch?v=PriyvxgsE30}{\texttt{PriyvxgsE30}} \\
哥林多後書 4:1-18 & \hyperref[sec:L2pB3hPJTOs]{2021.10.31《不喪膽的事奉》 哥林多後書 4\_1-18 (和修版) 林誠信牧師} & 2021-11-01 & \href{https://youtube.com/watch?v=L2pB3hPJTOs}{\texttt{L2pB3hPJTOs}} \\
詩篇 23:1-6 & \hyperref[sec:zUN5t6YF3yU]{2021.11.07 《好牧者之歌》 詩篇 23\_1-6 (和修版) 曾敏姑娘} & 2021-11-08 & \href{https://youtube.com/watch?v=zUN5t6YF3yU}{\texttt{zUN5t6YF3yU}} \\
以斯拉記 6:13-22 & \hyperref[sec:TujTTaedR4Y]{2021.11.14《 神的恩典 》以斯拉記 6\_13-22 (和修版) 老耀雄牧師} & 2021-11-15 & \href{https://youtube.com/watch?v=TujTTaedR4Y}{\texttt{TujTTaedR4Y}} \\
路加福音 10:38-42 & \hyperref[sec:sbhfqkg2NGI]{2021.11.21 《接待上主之道》 路加福音 10\_38-42 (和修版) 趙崇明博士} & 2021-11-25 & \href{https://youtube.com/watch?v=sbhfqkg2NGI}{\texttt{sbhfqkg2NGI}} \\
撒母耳記下啟示錄 12:22-23 & \hyperref[sec:S_q0D_6qEuY]{2021.11.28 《 我信永生 》撒母耳記下 12\_22-23;啟示錄 21\_3-6 (和修版) 陳麗蟬牧師} & 2021-11-29 & \href{https://youtube.com/watch?v=S-q0D_6qEuY}{\texttt{S-q0D\_6qEuY}} \\
撒母耳記上 13:5-14 & \hyperref[sec:foX2kGxGGAY]{2021.12.05 《事奉者的眼界》 撒母耳記上 13\_5-14 (和修版) 老耀雄牧師} & 2021-12-08 & \href{https://youtube.com/watch?v=foX2kGxGGAY}{\texttt{foX2kGxGGAY}} \\
列王記上 22:14 & \hyperref[sec:ds7i5doUxn0]{2021.12.12《只為真道爭辯》 列王記上 22\_14 (和修版) 郭鴻標牧師} & 2021-12-14 & \href{https://youtube.com/watch?v=ds7i5doUxn0}{\texttt{ds7i5doUxn0}} \\
馬太福音 2:1-12 & \hyperref[sec:yEoqVyieILA]{2021.12.19《來敬拜我們的王》 馬太福音 2\_1-12 (和修版) 曾敏姑娘} & 2021-12-19 & \href{https://youtube.com/watch?v=yEoqVyieILA}{\texttt{yEoqVyieILA}} \\
路加福音 2:8-14 & \hyperref[sec:_NykIK_k97E]{2021.12.26《聖誕節的意義》 路加福音 2\_8-14 (和修版) 曾裔貴牧師} & 2021-12-31 & \href{https://youtube.com/watch?v=_NykIK-k97E}{\texttt{\_NykIK-k97E}} \\
歌羅西書 1:15-23 & \hyperref[sec:yWaJULCl_Mc]{2022.01.02《尊貴無比的基督》 歌羅西書 1\_15-23 (和修版) 老耀雄牧師} & 2022-01-05 & \href{https://youtube.com/watch?v=yWaJULCl-Mc}{\texttt{yWaJULCl-Mc}} \\
路加福音 7:36-50 & \hyperref[sec:VMnlnDfbr9Q]{2022.01.09 《以愛還愛》 路加福音 7\_36-50 (和修版) 曾敏姑娘} & 2022-01-09 & \href{https://youtube.com/watch?v=VMnlnDfbr9Q}{\texttt{VMnlnDfbr9Q}} \\
啟示錄 21:1-8 & \hyperref[sec:biU1pOyNb74]{2022.01.16 《神的帳幕在人間》 啟示錄 21\_1-8 (和修版) 蔣文忠博士} & 2022-01-16 & \href{https://youtube.com/watch?v=biU1pOyNb74}{\texttt{biU1pOyNb74}} \\
約翰福音 1:6-8-15 & \hyperref[sec:0IxLPtCeZ88]{2022.01.23《見證》 約翰福音 1\_6-8,15 (和修版) 曾裔貴牧師} & 2022-01-26 & \href{https://youtube.com/watch?v=0IxLPtCeZ88}{\texttt{0IxLPtCeZ88}} \\
約翰福音 4:3-26 & \hyperref[sec:lmTkPW75UmA]{2022.01.30 《耶穌領人歸順的奧妙之處》 約翰福音4\_3-26 (新漢語譯本) 楊穎兒傳道} & 2022-02-03 & \href{https://youtube.com/watch?v=lmTkPW75UmA}{\texttt{lmTkPW75UmA}} \\
歌羅西書 2:6-23 & \hyperref[sec:aSFgvj6BhUY]{2022.02.06 《尊貴無比的基督信仰》 歌羅西書 2\_6-23 (和修版) 老耀雄牧師} & 2022-02-06 & \href{https://youtube.com/watch?v=aSFgvj6BhUY}{\texttt{aSFgvj6BhUY}} \\
出埃及記 17:8-16 & \hyperref[sec:_xVz_wL7W7s]{2022.02.13 《耶和華尼西》 出埃及記 17\_8-16 (和修版) 曾敏傳道} & 2022-02-13 & \href{https://youtube.com/watch?v=-xVz-wL7W7s}{\texttt{-xVz-wL7W7s}} \\
詩篇 1:1-3 & \hyperref[sec:jmOEpWc4LtY]{2022.02.20 《有福的人》 詩篇 1\_1-3 (和修版) 張美薇牧師} & 2022-02-22 & \href{https://youtube.com/watch?v=jmOEpWc4LtY}{\texttt{jmOEpWc4LtY}} \\
詩篇 27:1-14 & \hyperref[sec:PKTsSbC_DIE]{2022.02.27 《基督光--燃亮幽暗心》 詩篇 27\_1-14 (和修版) 陳麗蟬牧師} & 2022-02-27 & \href{https://youtube.com/watch?v=PKTsSbC-DIE}{\texttt{PKTsSbC-DIE}} \\
歌羅西書 3:1-17 & \hyperref[sec:jiG98Pa_BV4]{2022.03.06 《尊貴無比的基督徒身份》 歌羅西書3\_1-17 (和修版) 老耀雄牧師} & 2022-03-08 & \href{https://youtube.com/watch?v=jiG98Pa_BV4}{\texttt{jiG98Pa\_BV4}} \\
約翰福音 14:27 & \hyperref[sec:RMhyNLglQvM]{2022.03.13 《願主賜你們平安》 約翰福音 14\_27 (和修版) 曾敏姑娘} & 2022-03-14 & \href{https://youtube.com/watch?v=RMhyNLglQvM}{\texttt{RMhyNLglQvM}} \\
哥林多前書 12:4-11 & \hyperref[sec:bW0AAWns91U]{2022.03.20 《聖靈隨己意分給各人的恩賜》 哥林多前書 12\_4-11 (和修版) 何智雄牧師} & 2022-03-20 & \href{https://youtube.com/watch?v=bW0AAWns91U}{\texttt{bW0AAWns91U}} \\
哥林多後書 5:17 & \hyperref[sec:atdLdwQVuag]{2022.03.27 《新造的人》 哥林多後書 5\_17 (和修版) 曾裔貴牧師} & 2022-03-27 & \href{https://youtube.com/watch?v=atdLdwQVuag}{\texttt{atdLdwQVuag}} \\
馬太福音 26:36-46 & \hyperref[sec:LGE2dbGZ_40]{2022.04.03 《願你的旨意成全》 馬太福音 26\_36-46 (和修版) 曾敏姑娘} & 2022-04-04 & \href{https://youtube.com/watch?v=LGE2dbGZ-40}{\texttt{LGE2dbGZ-40}} \\
創世記 1:26-28 & \hyperref[sec:ra6IMFq2wMc]{2022.04.10 《因為愛,所以創作》 創世記1\_26-28;2\_18-23;3\_7,21 (新普及譯本) 陳朗欣傳道} & 2022-04-10 & \href{https://youtube.com/watch?v=ra6IMFq2wMc}{\texttt{ra6IMFq2wMc}} \\
馬可福音 15:33-39 & \hyperref[sec:DcwFkXHzJa4]{2022.04.15 《為甚麼離棄我》 馬可福音 15\_33-39 (和修版) 老耀雄牧師} & 2022-04-18 & \href{https://youtube.com/watch?v=DcwFkXHzJa4}{\texttt{DcwFkXHzJa4}} \\
約翰福音 20:24-29 & \hyperref[sec:lmsk0116lHM]{2022.04.17 《信心的轉化》 約翰福音 20\_24-29 (和修版) 老耀雄牧師} & 2022-04-18 & \href{https://youtube.com/watch?v=lmsk0116lHM}{\texttt{lmsk0116lHM}} \\
彼得前書 1:1-3 & \hyperref[sec:O3rEhqjezKM]{2022.04.24 《主復活的盼望》 彼得前書1\_1-3 (和修版) 謝靄薇牧師} & 2022-04-25 & \href{https://youtube.com/watch?v=O3rEhqjezKM}{\texttt{O3rEhqjezKM}} \\
馬可福音 8:27-37 & \hyperref[sec:D5wSPpsbnpU]{2022.05.01 《你願意付甚麼代價 ?》 馬可福音 8\_27-37 (和修版) 老耀雄牧師} & 2022-05-06 & \href{https://youtube.com/watch?v=D5wSPpsbnpU}{\texttt{D5wSPpsbnpU}} \\
撒母耳記上 14:6-15 & \hyperref[sec:8C2XRS_C_0M]{2022.05.08 《轉敗為勝的秘訣》 撒母耳記上 14\_6-15 (和修版) 陳一華牧師} & 2022-05-09 & \href{https://youtube.com/watch?v=8C2XRS-C-0M}{\texttt{8C2XRS-C-0M}} \\
馬可福音 16:14-20 & \hyperref[sec:A3yUdy51iFc]{2022.05.15 《宣道人 . 宣教人》 馬可福音 16\_14-20 (和修版) 游志豪傳道} & 2022-05-17 & \href{https://youtube.com/watch?v=A3yUdy51iFc}{\texttt{A3yUdy51iFc}} \\
馬太福音 7:21-23 & \hyperref[sec:PjWDhON898w]{2022.05.22 《我從來不認識你們》 馬太福音 7\_21-23 (和修版) 曾敏姑娘} & 2022-05-25 & \href{https://youtube.com/watch?v=PjWDhON898w}{\texttt{PjWDhON898w}} \\
使徒行傳路加福音 24:44-49 & \hyperref[sec:5T528Sf9V9Y]{2022.05.29 《宣教任務( 一) :務要 傳 道》 路加福音24\_44-49;使徒行傳1\_6-11(新漢語譯本) 楊穎兒傳道} & 2022-05-31 & \href{https://youtube.com/watch?v=5T528Sf9V9Y}{\texttt{5T528Sf9V9Y}} \\
啟示錄 2:1-7 & \hyperref[sec:abVWbFOM8VQ]{2022.06.05 《重回信仰生命的起跑線上》 啟示錄 2\_1-7 (和修版) 老耀雄牧師} & 2022-06-07 & \href{https://youtube.com/watch?v=abVWbFOM8VQ}{\texttt{abVWbFOM8VQ}} \\
民數記 24:1-25:3 & \hyperref[sec:qiH_9wRKv74]{2022.06.12 《豐厚?缺乏?》 民數記 24\_1,25\_3 (和修版) 曾敏姑娘} & 2022-06-12 & \href{https://youtube.com/watch?v=qiH_9wRKv74}{\texttt{qiH\_9wRKv74}} \\
約書亞記創世記 5:21-24 & \hyperref[sec:L3GT0EvH0Vg]{2022.06.19 《我和我一家必定事奉耶和華》 創世記 5\_21-24;約書亞記 24\_15 (和修版) 高淑芬傳道} & 2022-06-22 & \href{https://youtube.com/watch?v=L3GT0EvH0Vg}{\texttt{L3GT0EvH0Vg}} \\
詩篇 103:1-5 & \hyperref[sec:PeueLB8OHRY]{2022.06.26 《拒絕忘記》 曾裔貴牧師 詩篇 103\_1-5 (和修版)} & 2022-07-01 & \href{https://youtube.com/watch?v=PeueLB8OHRY}{\texttt{PeueLB8OHRY}} \\
啟示錄 3:14-22 & \hyperref[sec:zIFZXKCsIdk]{2022.07.03 《自以為是的屬靈生命》 老耀雄牧師 啟示錄 3\_14-22 (和修版)} & 2022-07-05 & \href{https://youtube.com/watch?v=zIFZXKCsIdk}{\texttt{zIFZXKCsIdk}} \\
哥林多前書 13:13 & \hyperref[sec:EJ_GoCr8XKk]{2022.07.10 《愛的呼喚》 哥林多前書 13\_13 (和修版) 陳一華牧師} & 2022-07-14 & \href{https://youtube.com/watch?v=EJ_GoCr8XKk}{\texttt{EJ\_GoCr8XKk}} \\
路得記 1:1-22 & \hyperref[sec:2v_ng9IARWc]{2022.07.17 《 忠心:是禍抑是褔? 》路得記 1\_1-22 梁子勇牧師} & 2022-07-22 & \href{https://youtube.com/watch?v=2v_ng9IARWc}{\texttt{2v\_ng9IARWc}} \\
出埃及記 12:40-41 & \hyperref[sec:rVzIB96PGyI]{2022.07.24 《榮耀的軍隊》 出埃及記 12\_40-41 (和修版) 曾敏姑娘} & 2022-07-25 & \href{https://youtube.com/watch?v=rVzIB96PGyI}{\texttt{rVzIB96PGyI}} \\
路加福音 10:1-11-16-20 & \hyperref[sec:pNmpdKyDl3E]{2022.07.31 《宣教任務( 二) :熟讀這篇武林秘笈》 路加福音 10\_1-11,16-20 ( 漢語新譯本)  楊穎兒傳道} & 2022-08-04 & \href{https://youtube.com/watch?v=pNmpdKyDl3E}{\texttt{pNmpdKyDl3E}} \\
啟示錄 4:1-11-20:11-15 & \hyperref[sec:MN9at1kI2W8]{2022.08.07 《世界的終局》 啟示錄 4\_1-11,20\_11-15 (和修版) 老耀雄牧師} & 2022-08-09 & \href{https://youtube.com/watch?v=MN9at1kI2W8}{\texttt{MN9at1kI2W8}} \\
約翰福音 15:1-8 & \hyperref[sec:lBPL13HHJ8A]{2022.08.14 《活出結果子的生命》 約翰福音 15\_1-8 (和修版) 曾敏姑娘} & 2022-08-18 & \href{https://youtube.com/watch?v=lBPL13HHJ8A}{\texttt{lBPL13HHJ8A}} \\
加拉太書 5:13-26 & \hyperref[sec:J30380FlPik]{2022.08.21 《聖靈的果子》 加拉太書 5\_13-26 (和修版) 黎光偉牧師} & 2022-08-22 & \href{https://youtube.com/watch?v=J30380FlPik}{\texttt{J30380FlPik}} \\
提摩太前書 6:1-19 & \hyperref[sec:D4j0WaSCINU]{2022.08.28 《假道理與真財富》 提摩太前書 6\_1-19 (和修版) 曾裔貴牧師} & 2022-09-01 & \href{https://youtube.com/watch?v=D4j0WaSCINU}{\texttt{D4j0WaSCINU}} \\
約翰福音 9:1-7-13-17-30-41 & \hyperref[sec:qTaSFAQFmRY]{2022.09.04 《屬靈的看見》 約翰福音 9\_1-7,13-17,30-41 (和修版) 老耀雄牧師} & 2022-09-04 & \href{https://youtube.com/watch?v=qTaSFAQFmRY}{\texttt{qTaSFAQFmRY}} \\
馬太福音使徒行傳羅馬書 28:19-20 & \hyperref[sec:u7aPgc4mo0Y]{2022.09.18 《承擔宣教的使命》 馬太福音 28\_19-20;使徒行傳 1\_8;羅馬書 10\_14-15 (和修版) 余家寶傳道} & 2022-09-21 & \href{https://youtube.com/watch?v=u7aPgc4mo0Y}{\texttt{u7aPgc4mo0Y}} \\
阿摩司書 5:1-17 & \hyperref[sec:QekXUaNsm08]{2022.09.25 《尋求耶和華者必存活》 阿摩司書 5\_1-17 (和修版) 陳麗蟬牧師} & 2022-09-29 & \href{https://youtube.com/watch?v=QekXUaNsm08}{\texttt{QekXUaNsm08}} \\
馬可福音 5:21-43 & \hyperref[sec:bLMot_knMH4]{2022.10.02 《一個怎樣的信心》馬可福音 5\_21-43 (和修版) 老耀雄牧師} & 2022-10-03 & \href{https://youtube.com/watch?v=bLMot_knMH4}{\texttt{bLMot\_knMH4}} \\
馬太福音 25:31-46 & \hyperref[sec:KSHhb73QdyA]{2022.10.09 《作一隻蒙天父賜福的綿羊》 馬太福音 25\_31-46 (和修版) 曾敏姑娘} & 2022-10-15 & \href{https://youtube.com/watch?v=KSHhb73QdyA}{\texttt{KSHhb73QdyA}} \\
以弗所書 4:7 & \hyperref[sec:rTQz8Qod6Ek]{2022.10.16 《基督所賜的恩》 以弗所書 4\_7, 11-16節 (和修版) 張美薇牧師} & 2022-10-18 & \href{https://youtube.com/watch?v=rTQz8Qod6Ek}{\texttt{rTQz8Qod6Ek}} \\
提摩太前書 3:14-16 & \hyperref[sec:ngXOvQd6nwI]{2022.10.23 《屬靈的家》 提摩太前書 3\_14-16 (和修版) 謝靄薇牧師} & 2022-10-27 & \href{https://youtube.com/watch?v=ngXOvQd6nwI}{\texttt{ngXOvQd6nwI}} \\
出埃及記 13:17-22 & \hyperref[sec:KbMX2DPfTSA]{2022.10.30 《曠野中的更新》 出埃及記 13\_17-22 節 (和修版) 董家驊牧師} & 2022-10-30 & \href{https://youtube.com/watch?v=KbMX2DPfTSA}{\texttt{KbMX2DPfTSA}} \\
尼希米記 10:28-39 & \hyperref[sec:WZ8nJrtvmJs]{2022.11.06 《重新立約》 尼希米記 10\_28-39 (和修版) 老耀雄牧師} & 2022-11-13 & \href{https://youtube.com/watch?v=WZ8nJrtvmJs}{\texttt{WZ8nJrtvmJs}} \\
馬太福音 7:24-27 & \hyperref[sec:VN5NAedAhq4]{2022.11.13 《建立堅不可摧的生命根基》馬太福音7\_24-27(和修版) 曾敏傳道} & 2022-11-13 & \href{https://youtube.com/watch?v=VN5NAedAhq4}{\texttt{VN5NAedAhq4}} \\
但以理書 2:18-6:10-9:3 & \hyperref[sec:bh_7j2paksc]{2022.11.20 《讓但以理的生命勒著你和我》但以理書2\_18,6\_10,9\_3(和修版) 高淑芬傳道} & 2022-11-25 & \href{https://youtube.com/watch?v=bh-7j2paksc}{\texttt{bh-7j2paksc}} \\
希伯來書 10:19-25 & \hyperref[sec:Hx9S3W_wWOQ]{2022.11.27 《讓我們定睛十字架》希伯來書 10\_19-25(和修版) 曾裔貴牧師} & 2022-11-29 & \href{https://youtube.com/watch?v=Hx9S3W_wWOQ}{\texttt{Hx9S3W\_wWOQ}} \\
哥林多前書 3:10-15 & \hyperref[sec:qxsO_vpno54]{2022.12.04 《建立蒙主喜悅的生命工程》哥林多前書 3\_10-15(和修版) 曾敏姑娘} & 2022-12-08 & \href{https://youtube.com/watch?v=qxsO_vpno54}{\texttt{qxsO\_vpno54}} \\
約翰福音 14:16-17-25-26-15:26-27-16:5-15 & \hyperref[sec:OuveUNRd_9o]{2022.12.11 《認識聖靈》 約翰福音 14\_16-17,25-26,15\_26-27,16\_5-15 (和修版) 講員 : 陳秀賢牧師} & 2022-12-12 & \href{https://youtube.com/watch?v=OuveUNRd-9o}{\texttt{OuveUNRd-9o}} \\
使徒行傳 16:6-10 & \hyperref[sec:pCQk0pIrnYA]{2022.12.18 《順著聖靈而行》使徒行傳 16\_6-10 (和修版) 鄧泰康傳道} & 2022-12-23 & \href{https://youtube.com/watch?v=pCQk0pIrnYA}{\texttt{pCQk0pIrnYA}} \\
路加福音 19:1-10 & \hyperref[sec:IOHR3FgR_rc]{2022.12.25 《生命的呼喚》 路加福音 19\_1-10 (和修版) 講員 : 老耀雄牧師} & 2022-12-26 & \href{https://youtube.com/watch?v=IOHR3FgR-rc}{\texttt{IOHR3FgR-rc}} \\
約翰福音 1:1-13 & \hyperref[sec:jrHnpDp52JU]{2023.01.01 《擁抱那非一般的道》 約翰福音 1\_1-13 (和修版) 講員 : 曾敏傳道} & 2023-01-01 & \href{https://youtube.com/watch?v=jrHnpDp52JU}{\texttt{jrHnpDp52JU}} \\
哥林多前書 12:4-27 & \hyperref[sec:mxoNZ6Tp6YY]{2023.01.15 《多元中的合一》哥林多前書12\_4-27 (和修版) 趙崇明博士} & 2023-01-19 & \href{https://youtube.com/watch?v=mxoNZ6Tp6YY}{\texttt{mxoNZ6Tp6YY}} \\
約書亞記申命記 1:5 & \hyperref[sec:8tGomcGt7oY]{2023.01.22 《讓約書亞的忠心勒著你和我》 約書亞記 1\_5;申命記 1\_37-38;約書亞記 24\_15 (和修版) 講員 : 高淑芬傳道} & 2023-01-26 & \href{https://youtube.com/watch?v=8tGomcGt7oY}{\texttt{8tGomcGt7oY}} \\
哥林多後書 8:9 & \hyperref[sec:_P22pS_FVkg]{2023.01.29 《無比恩典》哥林多後書 8\_9 (和修版) 陳一華牧師} & 2023-02-10 & \href{https://youtube.com/watch?v=_P22pS_FVkg}{\texttt{\_P22pS\_FVkg}} \\
士師記 19:1-30 & \hyperref[sec:CnIFsRI_OWk]{2023.02.05 《那時,以色列中沒有王!》士師記 19\_1-30 (和修版) 曾敏傳道} & 2023-02-10 & \href{https://youtube.com/watch?v=CnIFsRI_OWk}{\texttt{CnIFsRI\_OWk}} \\
路加福音 10:25-37 & \hyperref[sec:QG8KdqrlJlA]{2023.02.12 《誰是我的鄰舍?我是誰的鄰舍?》 路加福音 10\_25-37 (和修版) 彭秀芹導師} & 2023-02-13 & \href{https://youtube.com/watch?v=QG8KdqrlJlA}{\texttt{QG8KdqrlJlA}} \\
提摩太後書 2:19-22 & \hyperref[sec:_JHmy8bvPEI]{2023.02.19 《成為聖潔的器皿》提摩太後書 2\_19-22 (和修版) 張美薇牧師} & 2023-02-21 & \href{https://youtube.com/watch?v=_JHmy8bvPEI}{\texttt{\_JHmy8bvPEI}} \\
約翰一書 4:17-19 & \hyperref[sec:zwEIaeiy0j4]{2023.02.26 《幸福永伴隨》 約翰一書 4\_17-19 (和修版) 講員 : 曾裔貴牧師} & 2023-03-19 & \href{https://youtube.com/watch?v=zwEIaeiy0j4}{\texttt{zwEIaeiy0j4}} \\
哈該書 1:1-15 & \hyperref[sec:gwCa31uE_PU]{2023.03.05 《起來,建造神的家》 哈該書 1\_1-15 (和修版) 講員 : 曾敏傳道} & 2023-03-20 & \href{https://youtube.com/watch?v=gwCa31uE_PU}{\texttt{gwCa31uE\_PU}} \\
哈該書 2:1-23 & \hyperref[sec:F6OnOYSIQDU]{2023.03.12 《從今日起,我必賜福.》 哈該書 2\_1-23 (和修版) 講員 : 曾敏傳道} & 2023-03-22 & \href{https://youtube.com/watch?v=F6OnOYSIQDU}{\texttt{F6OnOYSIQDU}} \\
路加福音 19:1-10 & \hyperref[sec:bMJz4f2t79c]{2023.03.19 《我要這樣地委身……》 路加福音 19\_1-10 (和修版) 講員 : 張天和牧師} & 2023-03-22 & \href{https://youtube.com/watch?v=bMJz4f2t79c}{\texttt{bMJz4f2t79c}} \\
加拉太書 3:21-29 & \hyperref[sec:nJ0r0oqtEjI]{2023.03.26 《律法之上》 加拉太書3\_21-29 (和修版) 陳韋安牧師} & 2023-03-27 & \href{https://youtube.com/watch?v=nJ0r0oqtEjI}{\texttt{nJ0r0oqtEjI}} \\
約翰福音 11:28-44 & \hyperref[sec:7_miM8zb42g]{2023.04.02 《跟你休戚與共的耶穌》 約翰福音11\_28-44( 新漢語譯本) 彭秀芹導師} & 2023-04-04 & \href{https://youtube.com/watch?v=7_miM8zb42g}{\texttt{7\_miM8zb42g}} \\
馬可福音 10:35-45 & \hyperref[sec:AV2ZegtPBEc]{2023.04.09 《耶穌的苦杯,我們能喝嗎?》馬可福音10\_35-45 (和修版) 曾敏傳道} & 2023-04-10 & \href{https://youtube.com/watch?v=AV2ZegtPBEc}{\texttt{AV2ZegtPBEc}} \\
路加福音 11:1-13 & \hyperref[sec:b_zpDLNhNUk]{2023.04.16 《宣教任務( 四) :一切從禱告開始》 路加福音 11\_1-13 ( 新漢語譯本) 楊穎兒傳道} & 2023-04-19 & \href{https://youtube.com/watch?v=b-zpDLNhNUk}{\texttt{b-zpDLNhNUk}} \\
使徒行傳 5:29-42 & \hyperref[sec:c61oKrE6RXs]{2023.04.23 《神聖的耶穌》使徒行傳 5\_29-42 (和修版) 曾裔貴牧師} & 2023-04-28 & \href{https://youtube.com/watch?v=c61oKrE6RXs}{\texttt{c61oKrE6RXs}} \\
馬太福音 25:1-13 & \hyperref[sec:_IeOxEq_8vc]{2023.04.30 《有危機意識,認清問題的核心,行動起來》馬太福音 25\_1-13 (和修版) 郭鴻標牧師} & 2023-05-01 & \href{https://youtube.com/watch?v=-IeOxEq_8vc}{\texttt{-IeOxEq\_8vc}} \\
以賽亞書 1:1-31 & \hyperref[sec:pT_5txip_oE]{2023.05.07 《重拾對罪的警惕 重新領受神愛的管教 重燃悔改的動力》 以賽亞書 1\_1-31 (和修版) 講員 : 高淑芬傳道} & 2023-05-10 & \href{https://youtube.com/watch?v=pT-5txip_oE}{\texttt{pT-5txip\_oE}} \\
馬太福音 15:21-28 & \hyperref[sec:5NIuKMKgBWE]{2023.05.14 《一位祝福孩子的母親》馬太福音15\_21-28 (和修版) 彭秀芹導師} & 2023-05-18 & \href{https://youtube.com/watch?v=5NIuKMKgBWE}{\texttt{5NIuKMKgBWE}} \\
猶大書 1:1-4-17-25 & \hyperref[sec:XkIka_Z9Y3w]{2023.05.21 《聖道- 我將你的話藏在心裡》猶大書 1\_1-4,17-25 (和修版) 游志豪傳道} & 2023-05-21 & \href{https://youtube.com/watch?v=XkIka_Z9Y3w}{\texttt{XkIka\_Z9Y3w}} \\
羅馬書 15:22 & \hyperref[sec:33y9hhVAb_s]{2023.05.28 《五餅二魚的團契》羅馬書 15\_22~29 (和修版) 陳麗蟬牧師} & 2023-06-01 & \href{https://youtube.com/watch?v=33y9hhVAb-s}{\texttt{33y9hhVAb-s}} \\
使徒行傳 9:1-31 & \hyperref[sec:MMKQi99_fYg]{2023.06.04 《一個全新的你》 使徒行傳 9\_1-31 (和修版) 曾敏傳道} & 2023-06-07 & \href{https://youtube.com/watch?v=MMKQi99_fYg}{\texttt{MMKQi99\_fYg}} \\
馬太福音 6:19-34 & \hyperref[sec:B5mDgas2LUk]{2023.06.11 《今日的恩典》 馬太福音 6\_19-34 (和修版) 講員 : 曾裔貴牧師} & 2023-06-11 & \href{https://youtube.com/watch?v=B5mDgas2LUk}{\texttt{B5mDgas2LUk}} \\
路得記 4:13-17 & \hyperref[sec:NXj4I_yTkE8]{2023.06.18 《犧牲的愛；愛的犧牲》 路得記 4\_13-17 (和修版) 講員 : 梁子勇牧師} & 2023-06-18 & \href{https://youtube.com/watch?v=NXj4I-yTkE8}{\texttt{NXj4I-yTkE8}} \\
路加福音 1:5-25 & \hyperref[sec:3dWTg4gpwpQ]{2023.06.25 《不相稱的禱告》 路加福音 1\_5-25;57-66 節 (和修版) 講員 : 陳韋安牧師} & 2023-06-25 & \href{https://youtube.com/watch?v=3dWTg4gpwpQ}{\texttt{3dWTg4gpwpQ}} \\
約翰福音 21:15-19 & \hyperref[sec:qziNPLIfqJE]{2023.07.02 《你愛我比這些更深嗎?》約翰福音 21\_15-19 (和修版) 彭秀芹導師} & 2023-07-03 & \href{https://youtube.com/watch?v=qziNPLIfqJE}{\texttt{qziNPLIfqJE}} \\
尼希米記 11:1-3 & \hyperref[sec:5F_Xzyjh7_k]{2023.07.09 《幸福的感覺從行動開始》尼希米記 11\_1-3 (和修版) 曾裔貴牧師} & 2023-07-11 & \href{https://youtube.com/watch?v=5F-Xzyjh7-k}{\texttt{5F-Xzyjh7-k}} \\
路加福音 8:4-15 & \hyperref[sec:pH3FS1HN2vk]{2023.07.16 《生命豐盛只因有......》 路加福音 8\_4-15 (和修版) 講員 : 曾敏傳道} & 2023-07-17 & \href{https://youtube.com/watch?v=pH3FS1HN2vk}{\texttt{pH3FS1HN2vk}} \\
歌羅西書 2:1-19 & \hyperref[sec:Jxn7bTnB5O8]{2023.07.23 《打穩根基》 歌羅西書2\_1-19 (和修版) 謝靄薇牧師} & 2023-07-25 & \href{https://youtube.com/watch?v=Jxn7bTnB5O8}{\texttt{Jxn7bTnB5O8}} \\
馬可福音 5:1-20 & \hyperref[sec:5NHMIEcWoqY]{2023.07.30 《多麼大的憐憫》 馬可福音 5\_1-20 (和修版) 講員 : 黃宏達傳道} & 2023-08-01 & \href{https://youtube.com/watch?v=5NHMIEcWoqY}{\texttt{5NHMIEcWoqY}} \\
彼得後書 1:3-11 & \hyperref[sec:cnxh3H5LvDs]{2023.08.06 《活出敬虔,豐富的生命》 彼得後書 1\_3-11 (和修版) 曾敏傳道} & 2023-08-07 & \href{https://youtube.com/watch?v=cnxh3H5LvDs}{\texttt{cnxh3H5LvDs}} \\
使徒行傳 17:16-31 & \hyperref[sec:vjbxBlVTElQ]{2023.08.13 《進入人群的福音》 使徒行傳 17\_16-31 (和修版) 講員 : 陳麗蟬牧師} & 2023-08-16 & \href{https://youtube.com/watch?v=vjbxBlVTElQ}{\texttt{vjbxBlVTElQ}} \\
以弗所書 4:21-32 & \hyperref[sec:cIHnc4_JX1s]{2023.08.20 《脫去舊人‧ 穿上新人》以弗所書 4\_21-32 (和修版) 張美薇牧師} & 2023-08-22 & \href{https://youtube.com/watch?v=cIHnc4-JX1s}{\texttt{cIHnc4-JX1s}} \\
馬太福音 28:19 & \hyperref[sec:TnVaeXBM1r0]{2023.08.27 《你們要去,使萬民作我的門徒》 馬太福音 28\_19 (和修版) 講員 : 高淑芬傳道} & 2023-09-01 & \href{https://youtube.com/watch?v=TnVaeXBM1r0}{\texttt{TnVaeXBM1r0}} \\
詩篇 1:1-6 & \hyperref[sec:iQ2CTi_2bNs]{2023.09.03 《晝夜思想神的話語》 詩篇 1\_1-6 (和修版) 講員 : 老耀雄牧師} & 2023-09-04 & \href{https://youtube.com/watch?v=iQ2CTi_2bNs}{\texttt{iQ2CTi\_2bNs}} \\
馬太福音 11:1-19 & \hyperref[sec:Wzp5u5hQ940]{2023.09.10 《不一樣的彌賽亞》 馬太福音11\_1-19 (和修版) 講員 : 謝靄薇牧師} & 2023-09-17 & \href{https://youtube.com/watch?v=Wzp5u5hQ940}{\texttt{Wzp5u5hQ940}} \\
哥林多後書 8:1-9 & \hyperref[sec:BzeYjKv9RX8]{2023.09.17 《活出一個樂捐慷慨的生命》 哥林多後書8\_1-9 (和修版) 曾敏傳道} & 2023-09-18 & \href{https://youtube.com/watch?v=BzeYjKv9RX8}{\texttt{BzeYjKv9RX8}} \\
列王記下 5:1-5 & \hyperref[sec:1E5Bibwk4hc]{2023.09.24 《神使用小人物》 列王記下5\_1-5 (和修版) 陳一華牧師} & 2023-09-25 & \href{https://youtube.com/watch?v=1E5Bibwk4hc}{\texttt{1E5Bibwk4hc}} \\
馬太福音 6:9-15 & \hyperref[sec:ILOp7icU0bo]{2023.10.01 《渴望與追求》 馬太福音6\_9-15 (和修版) 黎光偉牧師} & 2023-10-03 & \href{https://youtube.com/watch?v=ILOp7icU0bo}{\texttt{ILOp7icU0bo}} \\
約翰一書 4:19-21 & \hyperref[sec:514XvjAlzD8]{2023.10.15 《我們愛,因為神先愛我們》 約翰一書4\_19-21 (和修版) 曾敏傳道} & 2023-10-16 & \href{https://youtube.com/watch?v=514XvjAlzD8}{\texttt{514XvjAlzD8}} \\
馬太福音 20:29-34 & \hyperref[sec:EQ76QKt8XLA]{2023.10.22 《因 信 奇 癒》 馬太福音20\_29-34 (和修版) 陳一華牧師} & 2023-10-31 & \href{https://youtube.com/watch?v=EQ76QKt8XLA}{\texttt{EQ76QKt8XLA}} \\
約珥書 2:10-14-28-32 & \hyperref[sec:VjKjBj0FEho]{2023.11.05 《審判或復興》 約珥書2\_10-14,28-32 (和修版) 講員 : 老耀雄牧師} & 2023-11-10 & \href{https://youtube.com/watch?v=VjKjBj0FEho}{\texttt{VjKjBj0FEho}} \\
馬太福音 11:25-30 & \hyperref[sec:EtvLUaUHtqc]{2023.11.12 《你們要來 , 在主裡得安息》 馬太福音11\_25-30 (和修版) 講員 : 高淑芬傳道} & 2023-11-15 & \href{https://youtube.com/watch?v=EtvLUaUHtqc}{\texttt{EtvLUaUHtqc}} \\
腓利門書 1:1-4-20 & \hyperref[sec:xQAERVXii10]{2023.10.29 42周年堂慶聯合崇拜《友伴．前行》腓利門書 1\_1,4,20 周曉暉牧師} & 2023-11-18 & \href{https://youtube.com/watch?v=xQAERVXii10}{\texttt{xQAERVXii10}} \\
羅馬書 12:15 & \hyperref[sec:3lGUeO3H154]{2023.11.19 《愛裡同行》 羅馬書12\_15 (和修版) 講員 : 曾敏傳道} & 2023-11-20 & \href{https://youtube.com/watch?v=3lGUeO3H154}{\texttt{3lGUeO3H154}} \\
路加福音 15:22-32 & \hyperref[sec:3jZxMYDYvEI]{2023.11.26 《天父的擁抱》路加福音15\_22-32 (和修版) 陳韋安牧師} & 2023-11-27 & \href{https://youtube.com/watch?v=3jZxMYDYvEI}{\texttt{3jZxMYDYvEI}} \\
詩篇 126:5-6 & \hyperref[sec:fKZdgGDua7A]{2023.12.03 《誰是收割者?》 詩篇 126\_5-6 (和修版) 謝靄薇牧師} & 2023-12-07 & \href{https://youtube.com/watch?v=fKZdgGDua7A}{\texttt{fKZdgGDua7A}} \\
馬太福音 6:9-13 & \hyperref[sec:7y90EbIBR_w]{2023.12.10 《主禱文》 馬太福音6\_9-13 (和修版) 講員 : 陳秀賢牧師} & 2023-12-10 & \href{https://youtube.com/watch?v=7y90EbIBR_w}{\texttt{7y90EbIBR\_w}} \\
創世記 32:22-32 & \hyperref[sec:drs0Z9UDTE0]{2023.12.17 《真正的得勝》 創世記32\_22-32 (和修版) 楊小玲傳道} & 2023-12-20 & \href{https://youtube.com/watch?v=drs0Z9UDTE0}{\texttt{drs0Z9UDTE0}} \\
路加福音 15:11-32 & \hyperref[sec:ZLhqPiukKhM]{2023.12.24 《回家》 路加福音15\_11-32 (和修版) 曾敏傳道} & 2023-12-27 & \href{https://youtube.com/watch?v=ZLhqPiukKhM}{\texttt{ZLhqPiukKhM}} \\
詩篇 100:1-5 & \hyperref[sec:6_IYA7gmRpA]{2023.12.31 《謝主洪恩》 詩篇 100\_1-5 (和修版) 講員 : 高淑芬傳道} & 2023-12-31 & \href{https://youtube.com/watch?v=6_IYA7gmRpA}{\texttt{6\_IYA7gmRpA}} \\
撒母耳記上 15:17-23 & \hyperref[sec:r5CRKC0_eBw]{2024.01.07 《作主聽命的兒女》 撒母耳記上 15\_17-23 (和修版) 曾敏傳道} & 2024-01-11 & \href{https://youtube.com/watch?v=r5CRKC0_eBw}{\texttt{r5CRKC0\_eBw}} \\
馬太福音 7:1-6 & \hyperref[sec:wCnaB7fTUZk]{2024.01.14 《以愛結連》 馬太福音 7\_1-6 (和修版) 曾裔貴牧師} & 2024-01-16 & \href{https://youtube.com/watch?v=wCnaB7fTUZk}{\texttt{wCnaB7fTUZk}} \\
約書亞記 1:1-9 & \hyperref[sec:sZSvUc8jKA0]{2024.01.21 《使命的承傳》 約書亞記 1\_1-9 (和修版) 郭文池院長} & 2024-01-23 & \href{https://youtube.com/watch?v=sZSvUc8jKA0}{\texttt{sZSvUc8jKA0}} \\
路加福音 23:40 & \hyperref[sec:dSTGrtAycig]{2024.01.28 《談敬畏》 路加福音 23\_40 (和修版) 陳韋安牧師} & 2024-01-30 & \href{https://youtube.com/watch?v=dSTGrtAycig}{\texttt{dSTGrtAycig}} \\
出埃及記 33:12-17 & \hyperref[sec:HRSr7qaJnjw]{2024.02.04 《認識神的摩西》 出埃及記33\_12-17 (和修版) 樓恩德牧師} & 2024-02-04 & \href{https://youtube.com/watch?v=HRSr7qaJnjw}{\texttt{HRSr7qaJnjw}} \\
路加福音 12:33-40 & \hyperref[sec:IFtn6wInBFk]{2024.02.11 《龍年祈福》 路加福音12\_33-40 (和修版) 曾裔貴牧師} & 2024-02-13 & \href{https://youtube.com/watch?v=IFtn6wInBFk}{\texttt{IFtn6wInBFk}} \\
創世記 11:1-9 & \hyperref[sec:Yl0EevfFnJU]{2024.02.18 《宣教使命:從分散一族到召聚萬國》 創世記11\_1-9 (和修版) 楊穎兒宣教士} & 2024-02-20 & \href{https://youtube.com/watch?v=Yl0EevfFnJU}{\texttt{Yl0EevfFnJU}} \\
提摩太後書 1:1-7 & \hyperref[sec:lg5Gw9lhP5s]{2024.02.25 《最有影響力的栽培員》 提摩太後書1\_1-7 (和修版) 講員 : 曾敏姑娘} & 2024-02-28 & \href{https://youtube.com/watch?v=lg5Gw9lhP5s}{\texttt{lg5Gw9lhP5s}} \\
約翰福音 4:31-38 & \hyperref[sec:DSsURlo4WPg]{2024.03.03 《不一樣的看見》 約翰福音 4\_31-38 (和修版) 講員 : 梁友東牧師} & 2024-03-08 & \href{https://youtube.com/watch?v=DSsURlo4WPg}{\texttt{DSsURlo4WPg}} \\
腓立比書 1:1-11 & \hyperref[sec:sd0jGt0zjkQ]{2024.03.10 《健康的關係 ‧ 喜樂的泉源》 腓立比書 1\_1-11 (和修版) 張美薇牧師} & 2024-03-12 & \href{https://youtube.com/watch?v=sd0jGt0zjkQ}{\texttt{sd0jGt0zjkQ}} \\
路加福音 13:22-30 & \hyperref[sec:wuCk59M9Tmo]{2024.03.17 《當努力進窄門》 路加福音13\_22-30 (和修版) 講員 : 曾敏姑娘} & 2024-03-18 & \href{https://youtube.com/watch?v=wuCk59M9Tmo}{\texttt{wuCk59M9Tmo}} \\
羅馬書 8:31-34 & \hyperref[sec:hGBqBs0DEW4]{2024.03.31 《一切從復活開始》 羅馬書 8\_31-34 (和修版) 曾裔貴牧師} & 2024-04-01 & \href{https://youtube.com/watch?v=hGBqBs0DEW4}{\texttt{hGBqBs0DEW4}} \\
  & \hyperref[sec:Ux8UdMSUx6c]{2024.03.31 復活節佈道會《活在每一天》 李炳光牧師} & 2024-04-02 & \href{https://youtube.com/watch?v=Ux8UdMSUx6c}{\texttt{Ux8UdMSUx6c}} \\
雅各書 1:1-8 & \hyperref[sec:EYrQ_D9F5Wo]{2024.04.07 《作個成熟的基督徒》 雅各書1\_1-8 (和修版) 講員 : 曾敏姑娘} & 2024-04-08 & \href{https://youtube.com/watch?v=EYrQ-D9F5Wo}{\texttt{EYrQ-D9F5Wo}} \\
哥林多前書 4:1-6 & \hyperref[sec:i_n9uN87S9M]{2024.04.14 《克服身份危機》 哥林多前書4\_1-6 (和修版) 講員 : 曾裔貴牧師} & 2024-04-14 & \href{https://youtube.com/watch?v=i-n9uN87S9M}{\texttt{i-n9uN87S9M}} \\
歷代志上彼得前書 16:23-24 & \hyperref[sec:jSl7w85PRnU]{2024.04.21 《宣教由 「信」做起》 歷代志上16\_23-24;彼得前書1\_17節 (和修版) 講員 : 高淑芬姑娘} & 2024-04-23 & \href{https://youtube.com/watch?v=jSl7w85PRnU}{\texttt{jSl7w85PRnU}} \\
腓立比書 2:17-30 & \hyperref[sec:bW5i5_FWoTI]{2024.04.28 《同心事主的知己》 腓立比書2\_17-30 (和修版) 講員 : 馮耀榮牧師} & 2024-05-01 & \href{https://youtube.com/watch?v=bW5i5-FWoTI}{\texttt{bW5i5-FWoTI}} \\
雅各書 1:12-18 & \hyperref[sec:W1Okagljeug]{2024.05.05 《過得勝的生活》 雅各書1\_12-18 (和修版) 曾敏姑娘} & 2024-05-06 & \href{https://youtube.com/watch?v=W1Okagljeug}{\texttt{W1Okagljeug}} \\
箴言 31:1-9 & \hyperref[sec:Q94i4NHF3WY]{2024.05.12 《利慕伊勒王的真言》 箴言31\_1-9 (和修版) 講員 : 謝靄薇牧師} & 2024-05-13 & \href{https://youtube.com/watch?v=Q94i4NHF3WY}{\texttt{Q94i4NHF3WY}} \\
馬可福音 4:35-41 & \hyperref[sec:mgPJ9gnhbtQ]{2024.05.19 《有主在船中》 馬可福音4\_35-41 (和修版) 講員 : 陳一華牧師} & 2024-05-21 & \href{https://youtube.com/watch?v=mgPJ9gnhbtQ}{\texttt{mgPJ9gnhbtQ}} \\
詩篇 46:1-11 & \hyperref[sec:JJ0c8EbTgFo]{2024.05.26 《上帝是我們患難中隨時的幫助》 詩篇46\_1-11 (和修版) 鄧泰康傳道} & 2024-05-28 & \href{https://youtube.com/watch?v=JJ0c8EbTgFo}{\texttt{JJ0c8EbTgFo}} \\
希伯來書 11:1-3 & \hyperref[sec:QSB7W2xKaig]{2024.06.02 《真實的信心》 希伯來書11\_1-3 (和修版) 曾敏姑娘} & 2024-06-03 & \href{https://youtube.com/watch?v=QSB7W2xKaig}{\texttt{QSB7W2xKaig}} \\
詩篇 118:1-13 & \hyperref[sec:EkAEiOY56rU]{2024.06.09 《與神同行的人生》 詩篇118\_1-13 (和修版) 曾裔貴牧師} & 2024-06-13 & \href{https://youtube.com/watch?v=EkAEiOY56rU}{\texttt{EkAEiOY56rU}} \\
使徒行傳 6:1-7 & \hyperref[sec:WqOKT5MNC_k]{2024.06.16 《lor 上帝做藉口之:當有投訴時》 使徒行傳6\_1-7 (和修版) 張雲開老師} & 2024-06-17 & \href{https://youtube.com/watch?v=WqOKT5MNC_k}{\texttt{WqOKT5MNC\_k}} \\
創世記 18:16-33 & \hyperref[sec:RaRAzKcNgEg]{2024.06.23 《與神同行》創世記 18\_16-33 (和修版) 許淑芬牧師} & 2024-06-24 & \href{https://youtube.com/watch?v=RaRAzKcNgEg}{\texttt{RaRAzKcNgEg}} \\
歌羅西書 1:4-6-13-23 & \hyperref[sec:6OJe2UUeX7c]{2024.06.30 《頂級耶穌》 歌羅西書1\_4-6,13-23 (和修版) 曾立華牧師} & 2024-07-01 & \href{https://youtube.com/watch?v=6OJe2UUeX7c}{\texttt{6OJe2UUeX7c}} \\
哥林多後書 3:4-6 & \hyperref[sec:W9_hqGLMr2Q]{2024.07.07 《我要感謝你》 哥林多後書3\_4-6 (和修版) 曾裔貴牧師} & 2024-07-08 & \href{https://youtube.com/watch?v=W9_hqGLMr2Q}{\texttt{W9\_hqGLMr2Q}} \\
約翰福音希伯來書 13:15 & \hyperref[sec:CrvvxFK_CfQ]{2024.07.14 《敬拜為主.為主敬拜》 希伯來書13\_15;約翰福音4\_24 (和修版) 吉中鳴牧師} & 2024-07-15 & \href{https://youtube.com/watch?v=CrvvxFK-CfQ}{\texttt{CrvvxFK-CfQ}} \\
啟示錄詩篇 24:7-10-100:5 & \hyperref[sec:in2H2R6OX8g]{2024.07.21 《Church of Turkey and Smyrna》 詩篇 24\_7-10,100\_5;啟示錄 2\_8-11( NIV)  Pastor Ozan} & 2024-07-21 & \href{https://youtube.com/watch?v=in2H2R6OX8g}{\texttt{in2H2R6OX8g}} \\
馬太福音 9:35-38 & \hyperref[sec:1eBxYFSJlEI]{2024.08.04 《求主差遣》 馬太福音 9\_35-38 (和修版) 譚立晶姑娘} & 2024-08-06 & \href{https://youtube.com/watch?v=1eBxYFSJlEI}{\texttt{1eBxYFSJlEI}} \\
以弗所書 4:20-24 & \hyperref[sec:sujYDAwGYXA]{2024.08.11 《帶著耶穌》以弗所書4\_20-24 (和修版) 曾裔貴牧師} & 2024-08-12 & \href{https://youtube.com/watch?v=sujYDAwGYXA}{\texttt{sujYDAwGYXA}} \\
約翰福音 1:14-18 & \hyperref[sec:_zoyMoXK3Ek]{2024.08.18 《上帝的說話》約翰福音1\_14-18 (和修版) 陳韋安牧師} & 2024-08-19 & \href{https://youtube.com/watch?v=-zoyMoXK3Ek}{\texttt{-zoyMoXK3Ek}} \\
馬可福音 2:1-12 & \hyperref[sec:QmBSplE5ueA]{2024.09.01 《行出看得見的信心》 馬可福音2\_1-12 (和修版) 講員 : 陳燕雯傳道} & 2024-09-02 & \href{https://youtube.com/watch?v=QmBSplE5ueA}{\texttt{QmBSplE5ueA}} \\
傳道書 6:1-12 & \hyperref[sec:XLqWIi0boGI]{2024.08.25 《你想人生》 傳道書6\_1-12 (和修版) 講員 : 周曉暉牧師} & 2024-09-04 & \href{https://youtube.com/watch?v=XLqWIi0boGI}{\texttt{XLqWIi0boGI}} \\
以弗所書 6:18-20 & \hyperref[sec:6BPEnUSuXaU]{2024.09.08 《凡人的祈禱》 以弗所書6\_18-20 (和修版) 講員 : 曾裔貴牧師} & 2024-09-08 & \href{https://youtube.com/watch?v=6BPEnUSuXaU}{\texttt{6BPEnUSuXaU}} \\
約翰福音 12:20-26 & \hyperref[sec:w9qOKksHuCE]{2024.09.15 《作一粒落在地裡的麥子》約翰福音12\_20-26 (和修版) 曾敏姑娘} & 2024-09-16 & \href{https://youtube.com/watch?v=w9qOKksHuCE}{\texttt{w9qOKksHuCE}} \\
申命記 18:15-22 & \hyperref[sec:upkPX3gy4Ao]{2024.09.22 《時代先知的工作》 申命記18\_15-22 (和修版) 講員 : 黃永強牧師} & 2024-09-24 & \href{https://youtube.com/watch?v=upkPX3gy4Ao}{\texttt{upkPX3gy4Ao}} \\
尼希米記 8:1-12 & \hyperref[sec:dufzJl9BQpc]{2024.09.29 《建立讀經群體》 尼希米記8\_1-12 (和修版) 講員 : 蔣文忠傳道} & 2024-10-01 & \href{https://youtube.com/watch?v=dufzJl9BQpc}{\texttt{dufzJl9BQpc}} \\
使徒行傳 16:6-10 & \hyperref[sec:8mvS2zOfCvg]{2024.10.06 《來自遠方的呼喚,你們聽見嗎?》 使徒行傳16\_6-10 (和修版) 講員 : 楊穎兒姑娘} & 2024-10-08 & \href{https://youtube.com/watch?v=8mvS2zOfCvg}{\texttt{8mvS2zOfCvg}} \\
歌羅西書 4:7-15 & \hyperref[sec:jVaUl1bB0os]{2024.10.13 《來自保羅團隊的問安》 歌羅西書4\_7-15 (和修版) 講員 : 曾裔貴牧師} & 2024-10-14 & \href{https://youtube.com/watch?v=jVaUl1bB0os}{\texttt{jVaUl1bB0os}} \\
列王記上 17:8-16 & \hyperref[sec:dXTahKH7Xis]{2024.10.20 《順從主的話,上帝必供應》 列王記上17\_8-16 (和修版) 講員 : 譚立晶傳道} & 2024-10-21 & \href{https://youtube.com/watch?v=dXTahKH7Xis}{\texttt{dXTahKH7Xis}} \\
路加福音 8:1-3 & \hyperref[sec:xtW6V_2_PsE]{2024.11.03 《跟隨基督的腳蹤》 路加福音8\_1-3 (和修版) 講員 : 馮浩鎏醫生} & 2024-11-04 & \href{https://youtube.com/watch?v=xtW6V_2_PsE}{\texttt{xtW6V\_2\_PsE}} \\
歷代志下 1:7-12 & \hyperref[sec:lQqzMqR_pN8]{2024.11.10 《蒙神喜悅的禱告》 歷代志下1\_7-12 (和修版) 陳一華牧師} & 2024-11-11 & \href{https://youtube.com/watch?v=lQqzMqR-pN8}{\texttt{lQqzMqR-pN8}} \\
創世記 16:7-14 & \hyperref[sec:Y3F7I4ibIlg]{2024.11.17 《神與你同行》 創世記16\_7-14 (和修版) 講員 : 許淑芬牧師} & 2024-11-17 & \href{https://youtube.com/watch?v=Y3F7I4ibIlg}{\texttt{Y3F7I4ibIlg}} \\
提多書 1:10-2:1 & \hyperref[sec:lRt1SGTCNSE]{2024.11.24 《講真話的教會》 提多書1\_10-2\_1 (和修版) 曾裔貴牧師} & 2024-11-25 & \href{https://youtube.com/watch?v=lRt1SGTCNSE}{\texttt{lRt1SGTCNSE}} \\
提多書 2:11-15 & \hyperref[sec:KpxypRPehLQ]{2024.12.01 《你仲返教會?點解ge ?》 提多書2\_11-15 (和修版) 講員 : 曾裔貴牧師} & 2024-12-01 & \href{https://youtube.com/watch?v=KpxypRPehLQ}{\texttt{KpxypRPehLQ}} \\
哥林多後書 4:5-7 & \hyperref[sec:3EdQX7vZPlY]{2024.12.08 《光明之子》 哥林多後書4\_5-7 (和修版) 陳韋安牧師} & 2024-12-09 & \href{https://youtube.com/watch?v=3EdQX7vZPlY}{\texttt{3EdQX7vZPlY}} \\
創世記 39:2 & \hyperref[sec:rX2pkPqwylA]{2024.12.15 《活出神的同在》 創世記39\_2 (和修版) 盧思源姑娘} & 2024-12-17 & \href{https://youtube.com/watch?v=rX2pkPqwylA}{\texttt{rX2pkPqwylA}} \\
創世記 4:1-16 & \hyperref[sec:j8cE70aJucM]{2024.12.22 《愛子》 創世記4\_1-16 (和修版) 陳晟希傳道} & 2024-12-23 & \href{https://youtube.com/watch?v=j8cE70aJucM}{\texttt{j8cE70aJucM}} \\
使徒行傳 26:1-23 & \hyperref[sec:tR6kQF63JQE]{2024.12.29 《沒有違背那從天上來的異象》 使徒行傳26\_1-23 (和修版) 講員 : 李富成牧師} & 2024-12-29 & \href{https://youtube.com/watch?v=tR6kQF63JQE}{\texttt{tR6kQF63JQE}} \\
瑪拉基書 3:6-12 & \hyperref[sec:0jw5FB1CQ6s]{2025.01.05 《我們一起來試驗神》 瑪拉基書3\_6-12 (和修版) 講員 : 曾裔貴牧師} & 2025-01-05 & \href{https://youtube.com/watch?v=0jw5FB1CQ6s}{\texttt{0jw5FB1CQ6s}} \\
哥林多前書 10:23-11:1 & \hyperref[sec:gjk_0mIGHss]{2025.01.12 《忠於使命,效法基督》 哥林多前書10\_23-11\_1 (和修版) 講員 : 曾敏姑娘} & 2025-01-14 & \href{https://youtube.com/watch?v=gjk-0mIGHss}{\texttt{gjk-0mIGHss}} \\
約書亞記 1:12-18 & \hyperref[sec:15mYsSTo_Pc]{2025.01.19 《不做局外人》 約書亞記1\_12-18 (和修版) 講員 : 鄧泰康傳道} & 2025-01-19 & \href{https://youtube.com/watch?v=15mYsSTo-Pc}{\texttt{15mYsSTo-Pc}} \\
路加福音 6:20 & \hyperref[sec:8UN_bXXTUSo]{2025.01.26 《崇拜與群體德行陶塑》 路加福音6\_20 (和修版) 講員 : 何世傑牧師} & 2025-02-01 & \href{https://youtube.com/watch?v=8UN_bXXTUSo}{\texttt{8UN\_bXXTUSo}} \\
箴言 29:18 & \hyperref[sec:9D6qqGrY1os]{2025.02.02 《洪恩家齊躍動》 箴言29\_18 (和修版) 陳麗蟬牧師} & 2025-02-03 & \href{https://youtube.com/watch?v=9D6qqGrY1os}{\texttt{9D6qqGrY1os}} \\
約翰一書 5:18-21 & \hyperref[sec:Ho71_1KRGq8]{2025.02.09 《信仰要尋真》 約翰一書5\_18-21 (和修版) 曾裔貴牧師} & 2025-02-10 & \href{https://youtube.com/watch?v=Ho71-1KRGq8}{\texttt{Ho71-1KRGq8}} \\
腓立比書  & \hyperref[sec:ulbWKYhlm8s]{2025.02.16 《做一個內心強大的基督徒》 《腓立比書》4-10-15 (和修版) 李文耀牧師} & 2025-02-17 & \href{https://youtube.com/watch?v=ulbWKYhlm8s}{\texttt{ulbWKYhlm8s}} \\
哥林多後書 8:1-15 & \hyperref[sec:TScbhjJAAD8]{2025.02.23 《信徒愛心的實在》 哥林多後書8\_1-15 (和修版) 講員 : 陳淑娟牧師} & 2025-02-25 & \href{https://youtube.com/watch?v=TScbhjJAAD8}{\texttt{TScbhjJAAD8}} \\
使徒行傳 9:10-20 & \hyperref[sec:QcQ_nHz9bjU]{2025.03.02 《善待不可愛的人》 使徒行傳9\_10-20 (和修版) 講員 : 李富成牧師} & 2025-03-02 & \href{https://youtube.com/watch?v=QcQ_nHz9bjU}{\texttt{QcQ\_nHz9bjU}} \\
使徒行傳 8:4-8 & \hyperref[sec:oUBZwJOpFGQ]{2025.03.09 《聖靈引領的傳道》 使徒行傳8\_4-8 (和修版) 講員 : 曾裔貴牧師} & 2025-03-09 & \href{https://youtube.com/watch?v=oUBZwJOpFGQ}{\texttt{oUBZwJOpFGQ}} \\
路加福音 22:39-46 & \hyperref[sec:td1EIi6hWco]{2025.03.16 《願你旨意成就》路加福音22\_39-46 (和修版) 講員 : 曾敏姑娘} & 2025-03-16 & \href{https://youtube.com/watch?v=td1EIi6hWco}{\texttt{td1EIi6hWco}} \\
哥林多前書 15:55-58 & \hyperref[sec:wH4miLAekc8]{2025.03.23《絕不徒然的事奉》 哥林多前書15\_55-58 (和修版) 講員 : 林誠信牧師} & 2025-03-25 & \href{https://youtube.com/watch?v=wH4miLAekc8}{\texttt{wH4miLAekc8}} \\
哥林多後書 4:16-5:10 & \hyperref[sec:gznE_6kHeYc]{2025.04.06《活潑的盼望》 哥林多後書4\_16-5\_10 講員 : 曾敏姑娘} & 2025-04-06 & \href{https://youtube.com/watch?v=gznE-6kHeYc}{\texttt{gznE-6kHeYc}} \\
使徒行傳 1:6 & \hyperref[sec:AR6jAGDUTDM]{2025.04.13《直到地極》 使徒行傳1\_6―8 講員 : 李思敬牧師} & 2025-04-15 & \href{https://youtube.com/watch?v=AR6jAGDUTDM}{\texttt{AR6jAGDUTDM}} \\
羅馬書 8:28 & \hyperref[sec:95mRci2_UGA]{2025.04.20《哭泣的祝福》 羅馬書8\_28 講員 : 曾裔貴牧師} & 2025-04-23 & \href{https://youtube.com/watch?v=95mRci2-UGA}{\texttt{95mRci2-UGA}} \\
路加福音 1:39-45 & \hyperref[sec:o8E8B2LqAe4]{2025.05.04《載道,信道,行道》 路加福音1\_39-45 講員 : 蔣文忠傳道} & 2025-05-04 & \href{https://youtube.com/watch?v=o8E8B2LqAe4}{\texttt{o8E8B2LqAe4}} \\
路加福音 1:39-45 & \hyperref[sec:vW2xmW8knK8]{2025.05.04《載道,信道,行道》 路加福音1\_39-45 講員 : 蔣文忠傳道} & 2025-05-13 & \href{https://youtube.com/watch?v=vW2xmW8knK8}{\texttt{vW2xmW8knK8}} \\
希伯來書 13:7-13 & \hyperref[sec:KCtIY5_ZgcE]{2025.05.11《祂的故事,我的故事》 希伯來書 13\_7-13 講員 : 曾裔貴牧師} & 2025-05-13 & \href{https://youtube.com/watch?v=KCtIY5_ZgcE}{\texttt{KCtIY5\_ZgcE}} \\
路加福音 7:36-50 & \hyperref[sec:_TgXxd1W4v8]{2025.05.18《看見尊貴》 路加福音7\_36-50 講員 : 黃國維牧師} & 2025-05-20 & \href{https://youtube.com/watch?v=-TgXxd1W4v8}{\texttt{-TgXxd1W4v8}} \\
撒母耳記上 1:1-28 & \hyperref[sec:xtvbpi6Ud4M]{2025.05.25《借錢梗要還》 經文:撒母耳記上1\_1-28 講員 : 張雲開老師} & 2025-05-27 & \href{https://youtube.com/watch?v=xtvbpi6Ud4M}{\texttt{xtvbpi6Ud4M}} \\
約書亞記 14:6-15 & \hyperref[sec:HMy_7Yhf_bE]{2025.06.01 《憑信得著應許的勇者――迦勒》 經文:約書亞記14\_6-15 講員 : 鄧泰康傳道} & 2025-06-06 & \href{https://youtube.com/watch?v=HMy-7Yhf-bE}{\texttt{HMy-7Yhf-bE}} \\
路加福音 15:11-24 & \hyperref[sec:iSHaOgcIOTI]{2025.06.15 《筵席的再思》 經文:路加福音15\_11-24 講員 : 曾裔貴牧師} & 2025-06-17 & \href{https://youtube.com/watch?v=iSHaOgcIOTI}{\texttt{iSHaOgcIOTI}} \\
創世記 5:21-24 & \hyperref[sec:XbnJouzOuIw]{2025.06.22 《與神同行人生路》 經文:創世記5\_21-24 講員 : 許淑芬牧師} & 2025-06-24 & \href{https://youtube.com/watch?v=XbnJouzOuIw}{\texttt{XbnJouzOuIw}} \\
哥林多後書 4:16-18 & \hyperref[sec:jY3KiVb_LXo]{2025.06.29 《如果我們的信仰是一次變身》 經文:哥林多後書4\_16-18 講員︰陳韋安牧師} & 2025-07-08 & \href{https://youtube.com/watch?v=jY3KiVb_LXo}{\texttt{jY3KiVb\_LXo}} \\
猶大書 1:20-21 & \hyperref[sec:m1thLVxgZE8]{2025.07.06 《保守自己常在神的愛中》 經文:猶大書1\_20-21 楊明麗姑娘} & 2025-07-08 & \href{https://youtube.com/watch?v=m1thLVxgZE8}{\texttt{m1thLVxgZE8}} \\
哥林多前書 13:1-13 & \hyperref[sec:V1IfIE0yrWg]{2025.07.13 《求主賜異象》 經文:哥林多前書13\_1-13 曾裔貴牧師} & 2025-07-15 & \href{https://youtube.com/watch?v=V1IfIE0yrWg}{\texttt{V1IfIE0yrWg}} \\
希伯來書 5:11-6:1 & \hyperref[sec:tbFVD8f2sac]{2025 7 27主日崇拜講道:《竭力成為一個成熟的基督徒》 經文:希伯來書5\_11-6\_1 曾敏姑娘} & 2025-07-30 & \href{https://youtube.com/watch?v=tbFVD8f2sac}{\texttt{tbFVD8f2sac}} \\
使徒行傳 26:19-32 & \hyperref[sec:YkzDcCtXDZ8]{2025.08.03:《回應天上的呼召》 經文:使徒行傳26\_19-32 講員 : 杜文瀚牧師} & 2025-08-07 & \href{https://youtube.com/watch?v=YkzDcCtXDZ8}{\texttt{YkzDcCtXDZ8}} \\
加拉太書 5:13-16 & \hyperref[sec:6RiKDzQP6RM]{2025.08.10 《跟隨,跟隨》 經文:加拉太書5\_13-16 曾裔貴牧師} & 2025-08-13 & \href{https://youtube.com/watch?v=6RiKDzQP6RM}{\texttt{6RiKDzQP6RM}} \\
馬太福音 22:15-22 & \hyperref[sec:OFK7CEvBZU8]{2025.08.17《這是誰的?》 經文:馬太福音22\_15-22 周克禮牧師} & 2025-08-30 & \href{https://youtube.com/watch?v=OFK7CEvBZU8}{\texttt{OFK7CEvBZU8}} \\
啟示錄 2:1-7 & \hyperref[sec:Lov4FCvlm6Q]{2025.08.24 《永恆的嘉獎-我知你的心》 經文:啟示錄2\_1-7 黃永康牧師} & 2025-08-30 & \href{https://youtube.com/watch?v=Lov4FCvlm6Q}{\texttt{Lov4FCvlm6Q}} \\
希伯來書 4:12-16 & \hyperref[sec:mR6Dgq2WzwQ]{2025 8 31主日講道2025.08.31 《真理的揭露與恩典的扶持》 經文:希伯來書4\_12-16 譚立晶傳道} & 2025-09-06 & \href{https://youtube.com/watch?v=mR6Dgq2WzwQ}{\texttt{mR6Dgq2WzwQ}} \\
列王記下 5:1-7 & \hyperref[sec:lUBiAGHTBWo]{2025.09.07 《大時代中的福音見證》:列王記下5\_1-7 (和修版) 黃國樑傳道} & 2025-09-11 & \href{https://youtube.com/watch?v=lUBiAGHTBWo}{\texttt{lUBiAGHTBWo}} \\
馬太福音 25:31-46 & \hyperref[sec:pWTRHiVMfgo]{2025.09.14《綿羊定山羊》:馬太福音25\_31-46 (和修版) 講員 :盧思源姑娘} & 2025-09-20 & \href{https://youtube.com/watch?v=pWTRHiVMfgo}{\texttt{pWTRHiVMfgo}} \\
使徒行傳 6:1-7:60 & \hyperref[sec:pPe9_H3XXfI]{2025.09.21《作主忠心的僕人》使徒行傳6\_1-7\_60 曾敏姑娘} & 2025-09-23 & \href{https://youtube.com/watch?v=pPe9-H3XXfI}{\texttt{pPe9-H3XXfI}} \\
馬太福音 9:35 & \hyperref[sec:WCmVYVwec7g]{2025.09.28《耶穌的「三濟」事業:教導的「理濟」傳道的「道濟」,醫治的「利濟」》馬太福音9\_35 (和修版) 郭鴻標牧師} & 2025-10-01 & \href{https://youtube.com/watch?v=WCmVYVwec7g}{\texttt{WCmVYVwec7g}} \\
馬太福音 9:35 & \hyperref[sec:ZmW9HFGb5CU]{2025.09.28《耶穌的「三濟」事業:教導的「理濟」傳道的「道濟」,醫治的「利濟」》馬太福音9\_35 (和修版) 郭鴻標牧師} & 2025-10-01 & \href{https://youtube.com/watch?v=ZmW9HFGb5CU}{\texttt{ZmW9HFGb5CU}} \\
以賽亞書 54:1-10 & \hyperref[sec:j0UisY02_wA]{2025.10.05《神在你荒涼中說話》:以賽亞書54\_1-10(和修版) 曾裔貴牧師} & 2025-10-07 & \href{https://youtube.com/watch?v=j0UisY02_wA}{\texttt{j0UisY02\_wA}} \\
約翰福音 5:24 & \hyperref[sec:Bm5kMjeNAdQ]{2025 10 12《惶恐中的平安》:約翰福音5\_24 (和修版) 講員 : 樓恩德牧師} & 2025-10-15 & \href{https://youtube.com/watch?v=Bm5kMjeNAdQ}{\texttt{Bm5kMjeNAdQ}} \\
以賽亞書 43:18-19 & \hyperref[sec:RkNxGeCkTHA]{2025.10.19《信仰中不斷變新的東西》:以賽亞書43\_18-19(和修版) 陳韋安牧師} & 2025-11-03 & \href{https://youtube.com/watch?v=RkNxGeCkTHA}{\texttt{RkNxGeCkTHA}} \\
馬太福音 16:13-19 & \hyperref[sec:Uvg7nmv7Igo]{2025.11.02 《教會的基石》 經文:馬太福音16\_13-19 陳麗蟬牧師} & 2025-11-04 & \href{https://youtube.com/watch?v=Uvg7nmv7Igo}{\texttt{Uvg7nmv7Igo}} \\
傳道書 7:1-18 & \hyperref[sec:7BCvZD9NU64]{2025.11.9《義人的滅亡》:傳道書7\_1-18 (和修版) 講員 :張雲開老師} & 2025-11-12 & \href{https://youtube.com/watch?v=7BCvZD9NU64}{\texttt{7BCvZD9NU64}} \\
馬太福音 10:34-39 & \hyperref[sec:7vPzFHXDGtY]{2025.11.16《與神同行,超越困難》:馬太福音10\_34-39 (和修版) 講員 :許淑芬牧師} & 2025-11-19 & \href{https://youtube.com/watch?v=7vPzFHXDGtY}{\texttt{7vPzFHXDGtY}} \\
馬太福音 28:16-20 & \hyperref[sec:11YyQ2efYE8]{2025.11.23《尚待完成的使命》:馬太福音28\_16-20 (和修版) 講員 :張美薇牧師} & 2025-11-29 & \href{https://youtube.com/watch?v=11YyQ2efYE8}{\texttt{11YyQ2efYE8}} \\
馬太福音 5:43-48 & \hyperref[sec:L0es3lIs6g8]{2025.11.30《我們要完全如同天父》:馬太福音5\_43-48 (和修版)} & 2025-12-03 & \href{https://youtube.com/watch?v=L0es3lIs6g8}{\texttt{L0es3lIs6g8}} \\
腓立比書 2:21 & \hyperref[sec:WQHTSm8b4dk]{2025.12.07《求耶穌基督的事》:腓立比書2\_21 (和修版) 講員 : 謝靄薇牧師} & 2025-12-10 & \href{https://youtube.com/watch?v=WQHTSm8b4dk}{\texttt{WQHTSm8b4dk}} \\
馬太福音 6:25-34 & \hyperref[sec:moCCv_ijPQc]{2025.12.14《今日的恩典》:馬太福音6\_25-34(和修版) 曾裔貴牧師} & 2025-12-16 & \href{https://youtube.com/watch?v=moCCv_ijPQc}{\texttt{moCCv\_ijPQc}} \\
約翰福音 10:7-18 & \hyperref[sec:fk3jYby_MrQ]{2025.12.21《祂來了,是叫人得生命》約翰福音10\_7-18(和修版) 李炳光牧師} & 2025-12-23 & \href{https://youtube.com/watch?v=fk3jYby_MrQ}{\texttt{fk3jYby\_MrQ}} \\
哥林多後書 8:1-5 & \hyperref[sec:0LnkonZdJmk]{2025.12.28《我拿甚麼奉獻我的主》哥林多後書8\_1-5(和修版) 李富成牧師} & 2025-12-30 & \href{https://youtube.com/watch?v=0LnkonZdJmk}{\texttt{0LnkonZdJmk}} \\
創世記 12:1-9 & \hyperref[sec:VGwZW_K3or4]{2026.01.04《起來》:創世記12\_1-9 (和修版) 曾敏姑娘} & 2026-01-07 & \href{https://youtube.com/watch?v=VGwZW-K3or4}{\texttt{VGwZW-K3or4}} \\
歌羅西書 4:7-9 & \hyperref[sec:8tVCcvQzY68]{2026.01.11《當保羅被囚》:歌羅西書4\_7-9 (和修版) 曾裔貴牧師} & 2026-01-14 & \href{https://youtube.com/watch?v=8tVCcvQzY68}{\texttt{8tVCcvQzY68}} \\
路加福音 10:25-37 & \hyperref[sec:sojB_TREQFc]{2026.01.18《走近你不想走近的人》:路加福音10\_25-37 (和修版) 何世傑牧師} & 2026-01-20 & \href{https://youtube.com/watch?v=sojB-TREQFc}{\texttt{sojB-TREQFc}} \\
民數記 20:14-21 & \hyperref[sec:wcZQh6png8E]{2026.01.25《生命的關口》:民數記20\_14-21 (和修版) 陳韋安牧師} & 2026-01-26 & \href{https://youtube.com/watch?v=wcZQh6png8E}{\texttt{wcZQh6png8E}} \\
箴言 3:5-8 & \hyperref[sec:w_aZ5m39lOc]{2026.02.01《謙卑事奉---生命的轉向》:箴言3\_5-8(和修版) 楊明麗姑娘} & 2026-02-04 & \href{https://youtube.com/watch?v=w-aZ5m39lOc}{\texttt{w-aZ5m39lOc}} \\
哥林多前書 9:19-27 & \hyperref[sec:4qbR6KQJgy4]{2026.02.15《不再「人有我有」的事奉》:哥林多前書9\_19-27) 曾裔貴牧師} & 2026-02-19 & \href{https://youtube.com/watch?v=4qbR6KQJgy4}{\texttt{4qbR6KQJgy4}} \\
哥林多後書 12:7-11 & \hyperref[sec:POs_KWCBzOE]{2026.02.22《因祂用厚恩待我》 哥林多後書12\_7-11 (和修版) 李富成牧師} & 2026-02-25 & \href{https://youtube.com/watch?v=POs-KWCBzOE}{\texttt{POs-KWCBzOE}} \\
馬可福音路加福音 6:34 & \hyperref[sec:S89_Yf51Xes]{2026.03.08 馬可福音6\_34;路加福音15\_20 (和修版) 曾裔貴牧師} & 2026-03-11 & \href{https://youtube.com/watch?v=S89-Yf51Xes}{\texttt{S89-Yf51Xes}} \\
約翰福音 13:1-17 & \hyperref[sec:Jj4ZioVnCDo]{2026.03.15 《做騷的耶穌》:約翰福音13\_1-17 (和修版) 張雲開老師} & 2026-03-18 & \href{https://youtube.com/watch?v=Jj4ZioVnCDo}{\texttt{Jj4ZioVnCDo}} \\
哈該書 1:1-15 & \hyperref[sec:zmeyo_PNxjc]{2026.03.22 《你們要省察自己的行為》:哈該書1\_1-15 (和修版) 陳淑娟牧師} & 2026-03-25 & \href{https://youtube.com/watch?v=zmeyo_PNxjc}{\texttt{zmeyo\_PNxjc}} \\
約翰福音 12:23-28 & \hyperref[sec:_kJpwvPJ95U]{2026.03.29 《得榮耀的時候》:約翰福音12\_23-28 (和修版) 黃國維牧師} & 2026-04-02 & \href{https://youtube.com/watch?v=-kJpwvPJ95U}{\texttt{-kJpwvPJ95U}} \\
哥林多前書 15:14-19 & \hyperref[sec:ke2vMZ9uUT4]{2026.04.05 《復活,不僅是認同》:哥林多前書15\_14-19(和修版) 曾裔貴牧師} & 2026-04-07 & \href{https://youtube.com/watch?v=ke2vMZ9uUT4}{\texttt{ke2vMZ9uUT4}} \\
哈巴谷書 3:16-19 & \hyperref[sec:y558B5MlYc4]{2026.04.12 《黑夜中的信心之歌》:哈巴谷書3\_16-19(和修版) 曾敏傳道} & 2026-04-14 & \href{https://youtube.com/watch?v=y558B5MlYc4}{\texttt{y558B5MlYc4}} \\
歌羅西書 2:6-15 & \hyperref[sec:R7nxbZBAviM]{2026.04.19 《在基督裡生根建造滿豐盛》:歌羅西書2\_6-15(和修版) 講員 : 譚立晶傳道} & 2026-04-21 & \href{https://youtube.com/watch?v=R7nxbZBAviM}{\texttt{R7nxbZBAviM}} \\
希伯來書 12:14-28 & \hyperref[sec:WokQkPhl0sU]{2026.04.26 《存感恩與神同行》:希伯來書12\_14-28(和修版) 許淑芬牧師} & 2026-05-01 & \href{https://youtube.com/watch?v=WokQkPhl0sU}{\texttt{WokQkPhl0sU}} \\
羅馬書 12:1-3 & \hyperref[sec:DhmQL4xjGtQ]{2026.05.03 《合乎中道的自我認識》:羅馬書12\_1-3節(和修版) 曾裔貴牧師} & 2026-05-07 & \href{https://youtube.com/watch?v=DhmQL4xjGtQ}{\texttt{DhmQL4xjGtQ}} \\
彼得前書 3:13-22 & \hyperref[sec:WnFpUR4Lsgw]{2026.05.10 《聖道為軸的崇拜》:彼得前書3\_13-22(和修版) 譚靜芝博士} & 2026-05-14 & \href{https://youtube.com/watch?v=WnFpUR4Lsgw}{\texttt{WnFpUR4Lsgw}} \\
創世記 13:8-9 & \hyperref[sec:W5xhymK0FGw]{2026.05.17 《有情義的人蒙福》:創世記13\_8-9(和修版) 郭鴻標牧師} & 2026-05-20 & \href{https://youtube.com/watch?v=W5xhymK0FGw}{\texttt{W5xhymK0FGw}} \\
約翰福音 7:37-39 & \hyperref[sec:3yIpsNheSeo]{2026.05.24 《湧出活水的生命》:約翰福音7\_37-39(和修版) 陳麗蟬牧師} & 2026-05-27 & \href{https://youtube.com/watch?v=3yIpsNheSeo}{\texttt{3yIpsNheSeo}} \\
創世記 39:1-23 & \hyperref[sec:kpqG_RsMI0U]{2026.05.31 《有神與無神之別》:創世記39\_1-23(和修版) 杜文翰牧師} & 2026-06-05 & \href{https://youtube.com/watch?v=kpqG-RsMI0U}{\texttt{kpqG-RsMI0U}} \\
路得記 2:4-18 & \hyperref[sec:iob9mODZZqY]{2026.06.07 《跟從神!向有需要者施恩》:路得記2\_4-18(和修版) 翟凱欣傳道} & 2026-06-10 & \href{https://youtube.com/watch?v=iob9mODZZqY}{\texttt{iob9mODZZqY}} \\
馬太福音 19:27-30 & \hyperref[sec:csMea85HLY0]{2026.06.14 《恩典,超越計算》:馬太福音19\_27-30 曾裔貴牧師} & 2026-06-19 & \href{https://youtube.com/watch?v=csMea85HLY0}{\texttt{csMea85HLY0}} \\
路加福音 8:40-56 & \hyperref[sec:V7GlEPh1gqw]{2026.06.21 《父.愛》:路加福音8\_40-56(和修版) 曾敬宗牧師} & 2026-06-25 & \href{https://youtube.com/watch?v=V7GlEPh1gqw}{\texttt{V7GlEPh1gqw}} \\
\end{xltabular}
}
\newpage



\section{帖撒羅尼迦前書 3:3-5:9}
\label{sec:iY8ed2cNTa4}
\textbf{2020.08.02《患難是擺在面前的》 帖撒羅尼迦前書 3\_3,5\_9 郭鴻標牧師}
\newline
\newline
連結: \href{https://youtube.com/watch?v=iY8ed2cNTa4}{\texttt{https://youtube.com/watch?v=iY8ed2cNTa4}} ~~~~ 語音日期: 2020-08-10
\newline
\newline
\hyperref[sec:code]{\small{< < < PREV SERMON < < <}}
~
\hyperref[sec:index_chronic]{\small{[返順時目]}}
~
\hyperref[sec:index_scriptual]{\small{[返順卷目]}}
~
\hyperref[sec:N89ZkrKQWlk]{\small{> > > NEXT SERMON > > >}}
\newline
\newline
帖撒羅尼迦前書 3:3-5:9
\newline
\begin{longtable}{cl}
\hline
\hline
章節 & 經文 (和合本修訂版)\\
\hline
3:3 & \begin{tabularx}{0.7\textwidth}{X} 免得有人被這些患難動搖。因為你們自己知道，我們受患難原是命定的。 \end{tabularx} \\ \\ \relax
3:4 & \begin{tabularx}{0.7\textwidth}{X} 我們在你們那裡的時候，曾預先告訴你們，我們必受患難；你們知道，這果然發生了。 \end{tabularx} \\ \\ \relax
3:5 & \begin{tabularx}{0.7\textwidth}{X} 為此，既然我不能再忍，就差派人去，要知道你們的信心如何，恐怕那誘惑人的果真誘惑了你們，以致我們的勞苦歸於徒然。 \end{tabularx} \\ \\ \relax
3:6 & \begin{tabularx}{0.7\textwidth}{X} 但是，提摩太剛從你們那裡回來，將你們信心和愛心的好消息報給我們，又說你們常常記念我們，切切想見我們，如同我們想見你們一樣。 \end{tabularx} \\ \\ \relax
3:7 & \begin{tabularx}{0.7\textwidth}{X} 所以，弟兄們，我們在一切困苦患難中，因著你們的信心得到鼓勵。 \end{tabularx} \\ \\ \relax
3:8 & \begin{tabularx}{0.7\textwidth}{X} 如今你們若靠主站立得穩，我們就得生了。 \end{tabularx} \\ \\ \relax
3:9 & \begin{tabularx}{0.7\textwidth}{X} 我們在神面前，因著你們滿有喜樂。為這一切喜樂，我們能用怎樣的感謝為你們報答神呢？ \end{tabularx} \\ \\ \relax
3:10 & \begin{tabularx}{0.7\textwidth}{X} 我們晝夜切切祈求要見你們的面，來補足你們信心的不足。 \end{tabularx} \\ \\ \relax
3:11 & \begin{tabularx}{0.7\textwidth}{X} 願我們的父神自己和我們的主耶穌，為我們開路到你們那裡去。 \end{tabularx} \\ \\ \relax
3:12 & \begin{tabularx}{0.7\textwidth}{X} 又願主使你們彼此相愛的心，和愛眾人的心，都能增長，充足，如同我們愛你們一樣， \end{tabularx} \\ \\ \relax
3:13 & \begin{tabularx}{0.7\textwidth}{X} 好堅固你們的心，使你們在我們的主耶穌同他眾聖徒來臨的時候，在我們父神面前，成為聖潔，無可指責。阿們！ \end{tabularx} \\ \\ \relax
4:1 & \begin{tabularx}{0.7\textwidth}{X} 末了，弟兄們，我們靠著主耶穌求你們，勸你們，既然你們領受了我們的教導，知道該怎樣行事為人，討神的喜悅，其實你們也正這樣行，我勸你們要更加努力。 \end{tabularx} \\ \\ \relax
4:2 & \begin{tabularx}{0.7\textwidth}{X} 你們原知道，我們憑主耶穌傳給你們甚麼命令。 \end{tabularx} \\ \\ \relax
4:3 & \begin{tabularx}{0.7\textwidth}{X} 神的旨意就是要你們成為聖潔，遠避淫行； \end{tabularx} \\ \\ \relax
4:4 & \begin{tabularx}{0.7\textwidth}{X} 要你們各人知道怎樣用聖潔、尊貴控制自己的身體 ， \end{tabularx} \\ \\ \relax
4:5 & \begin{tabularx}{0.7\textwidth}{X} 不放縱私慾的邪情，像不認識神的外邦人。 \end{tabularx} \\ \\ \relax
4:6 & \begin{tabularx}{0.7\textwidth}{X} 不准有人在這事上越軌，佔他弟兄的便宜；因為這一類的事，主必報應，正如我預先對你們說過，又切切警告過你們的。 \end{tabularx} \\ \\ \relax
4:7 & \begin{tabularx}{0.7\textwidth}{X} 神召我們本不是要我們沾染污穢，而是要我們聖潔。 \end{tabularx} \\ \\ \relax
4:8 & \begin{tabularx}{0.7\textwidth}{X} 所以，那棄絕這教導的不是棄絕人，而是棄絕那把自己的聖靈賜給你們的神。 \end{tabularx} \\ \\ \relax
4:9 & \begin{tabularx}{0.7\textwidth}{X} 有關弟兄間的手足之情，不用人寫信給你們，因為你們自己蒙了神的教導要彼此相愛。 \end{tabularx} \\ \\ \relax
4:10 & \begin{tabularx}{0.7\textwidth}{X} 你們向全馬其頓的眾弟兄固然是這樣行，但我勸弟兄們要更加努力。 \end{tabularx} \\ \\ \relax
4:11 & \begin{tabularx}{0.7\textwidth}{X} 要立志過安靜的生活，管自己的事，親手做工，正如我們從前吩咐你們的， \end{tabularx} \\ \\ \relax
4:12 & \begin{tabularx}{0.7\textwidth}{X} 好使你們的行為能得外人的尊敬，同時也不依賴任何人。 \end{tabularx} \\ \\ \relax
4:13 & \begin{tabularx}{0.7\textwidth}{X} 弟兄們，至於已睡了的人，我們不願意你們不知道，恐怕你們憂傷，像那些沒有指望的人一樣。 \end{tabularx} \\ \\ \relax
4:14 & \begin{tabularx}{0.7\textwidth}{X} 既然我們信耶穌死了，復活了，那些已經在耶穌裡睡了的人，神也必將他們與耶穌一同帶來。 \end{tabularx} \\ \\ \relax
4:15 & \begin{tabularx}{0.7\textwidth}{X} 我們照主的話告訴你們一件事：我們這活著還存留到主來臨的人，絕不會在那已經睡了的人之先。 \end{tabularx} \\ \\ \relax
4:16 & \begin{tabularx}{0.7\textwidth}{X} 因為，召集令一發，天使長的呼聲一叫，神的號角一吹，主必親自從天降臨；那在基督裡死了的人必先復活， \end{tabularx} \\ \\ \relax
4:17 & \begin{tabularx}{0.7\textwidth}{X} 然後我們這些活著還存留的人必和他們一同被提到雲裡，在空中與主相會。這樣，我們就要和主永遠同在。 \end{tabularx} \\ \\ \relax
4:18 & \begin{tabularx}{0.7\textwidth}{X} 所以，你們當用這些話彼此勸勉。 \end{tabularx} \\ \\ \relax
5:1 & \begin{tabularx}{0.7\textwidth}{X} 弟兄們，關於那時候和日期，不用人寫信給你們， \end{tabularx} \\ \\ \relax
5:2 & \begin{tabularx}{0.7\textwidth}{X} 因為你們自己明明知道，主的日子來到會像賊在夜間突然來到一樣。 \end{tabularx} \\ \\ \relax
5:3 & \begin{tabularx}{0.7\textwidth}{X} 人正說平安穩定的時候，災禍忽然臨到他們，如同陣痛臨到懷胎的婦人一樣，他們絕逃脫不了。 \end{tabularx} \\ \\ \relax
5:4 & \begin{tabularx}{0.7\textwidth}{X} 弟兄們，你們並不在黑暗裡，那日子不會像賊一樣臨到你們。 \end{tabularx} \\ \\ \relax
5:5 & \begin{tabularx}{0.7\textwidth}{X} 你們都是光明之子，都是白晝之子；我們不屬黑夜，也不屬幽暗。 \end{tabularx} \\ \\ \relax
5:6 & \begin{tabularx}{0.7\textwidth}{X} 所以，我們不要沉睡，像別人一樣，總要警醒謹慎。 \end{tabularx} \\ \\ \relax
5:7 & \begin{tabularx}{0.7\textwidth}{X} 因為睡了的人是在夜間睡，醉了的人是在夜間醉。 \end{tabularx} \\ \\ \relax
5:8 & \begin{tabularx}{0.7\textwidth}{X} 但既然我們屬於白晝，就應當謹慎，把信和愛當作護心鏡遮胸，把得救的盼望當作頭盔戴上。 \end{tabularx} \\ \\ \relax
5:9 & \begin{tabularx}{0.7\textwidth}{X} 因為神不是預定我們受懲罰，而是預定我們藉著我們的主耶穌基督得救。 \end{tabularx} \\ \\
[1ex]
\hline
\hline
\end{longtable}
$^{1}$我們一起同心去祈禱.
我們親愛的天父上帝.
因為我們可以去敬拜你.
我們時刻的要思想.
我們的生命是一份禮物.
是一份恩賜.
我們實在要再去思想.
當我們發現到.
我們有了一份恩賜.
當我們發現到.
生活上的事情.
好像不是我們可以掌握的時候.
我們更加要知道.
我們只不過是很普通的人.
在我們以上有我們的上帝.
你是我們生命的主.
在這個充滿混亂的時候.
我們要回歸於你.
再一次在你的面前去謙卑.
讓我們去承認.
你是我們生命的主.
讓我們在敬拜的時候.
得到你自己聖經的教導.
讓你的聖靈住在我們的內心.
讓我們有力量去面對各種挑戰.
因為我們知道.
我們所盼望的.
不是在地上.
而是在永恆的國度裡面.
我們的祈禱奉耶穌基督的名.
阿們.
讓我和洪恩堂的弟兄姊妹.
來分享聖經的說話.
我心裡有一種很特別的感受.
雖然我看不到你們.
但是我感受到你們在一起崇拜.
最近也很感恩.
因為經歷了.
雖然外面的環境有很多變化.
但是神又讓我的情緒穩定.

$^{41}$可以每天去工作.
特別是看聖經.
這件事是我以前經歷不足的事.
不過這一兩年我經歷了很多.
今天我將我看聖經的心得.
和大家分享.
剛才陳傳道已經將那兩節經文讀了出來.
所以我不再重複了.
在這兩段經文裡.
有兩個詞是很特別的.
一個詞是預定.
第二個詞是明定.
好了,我會問.
究竟這兩段經文裡.
這兩個詞.
在希臘文的原文裡.
是不是同一個字?.
如果兩個是不同的字.
究竟原來的意思是什麼?.
透過追查經文的背景.
去發現一些道理.
我將這個發現定為一個講道的題目.
就是患難是擺在面前的.
我將講道分為三個部分.
第一個部分是神預定人是得救的.
第二個部分就是受苦.
就是擺在我們面前的.
第三個部分就是神的旨意.
是人的聖潔.
這三個部分都是從這兩段經文中歸納出來的.
首先我們看看第一部分.
神預定人得救.
在帖撒羅尼加前書第五章第九節.
因為神不是預定我們受刑.
乃是預定我們直著.
我們主耶穌基督得救.
在這段經文裡.
這個預定.
希利文.
Attack.

$^{81}$是神預定我們直著耶穌基督得救的.
這是很清楚的.
不過很多基督徒.
他們聽過我們基督教神學裡面.
有一種叫做預定論.
英文叫做Predestination.
不過未必很清楚.
知道基督教的預定論.
Predestination.
和希臘哲學裡面的決定論.
Determinism.
有什麼根本上的分別.
因為在希臘哲學裡面的決定論.
Determinism.
是有一種悲劇的意識.
就好像神明.
總是和人作對.
希臘神話中.
有一個故事.
就是西西弗斯.
Sisyphus.
他被懲罰.
他要將石頭推上山.
每一次.
當他推石頭到山頂的時候.
石頭又會滾回山上.
西西弗斯.
他就要重重複複地.
將石頭推上山.
又滾回山上.
又推上山.
這些動作是重重複複的.
這個故事就形容.
他的人生是一個悲劇的人生.
所以你看希臘的神話故事.
是有很多悲劇的英雄.
故事的主題其實是天意弄人.
很多很多年前.
有一套電視劇.
叫做亞信的故事.

$^{121}$主題曲第一句就說.
命運是對手.
永不低頭.
我認為這是亞洲版本的決定論.
Determinism.
命運是作弄人的.
值得我們去問的時候.
我們的人生.
是不是被命運作弄呢?.
如果宇宙中間有神.
這位神.
是不是又是天道英才呢?.
因為我們看電影經常有這些題材.
好了.
貼在羅尼加前書第五章第九節就說.
因為神不是預定我們受刑.
乃是預定我們藉著我們主耶穌基督得救.
神不是預定來懲罰我們的.
不是預定我們去受苦的.
乃是預定我們藉著耶穌基督來得救的.
請大家留心這節經文.
我們一起去思想.
當我們遇到挫折.
苦難的時候.
我們是不應該埋怨神的.
我們是赤身出於無胎.
亦赤身歸為塵土.
我們根本一無所有.
我們由出生到死亡之間.
我們每一個人都努力為自己增添很多很多.
我們稱為心愛的物.
我們覺得失去一些心愛的物.
是會很失望的.
當我們失去健康.
就更失望.
如果我們失去親人朋友.
亦會很失望.
因為有感情.
災難,患難,苦難,痛苦.
我們可以說.

$^{161}$一方面是客觀的事情.
有件事我們才會有這種感受.
另一方面.
你可以說.
是我們心裡有患難的感受.
所以有外面的.
亦有裡面的.
兩個層面.
在新冠狀病毒威脅下.
我們會感覺到災難,患難,苦難,痛苦.
我們生活出現了變化.
內心感到無力.
最近有一些牧者和弟兄姊妹跟我分享.
早一段時間.
他們滿懷希望.
因為以為疫情好轉.
教會漸漸可以恢復現場崇拜.
特別有些教會老人家比較多.
他們因為不習慣上網.
他們很渴望見面.
很渴望聊天.
所以很開心.
但不久後我們聽到傳聞.
教會要煞車.
因為疫情變化了.
對很多牧者來說.
他們都覺得打亂了陣腳.
因為已經想著如何恢復正常的教會生活.
但弟兄姊妹都是一樣.
等了這麼久.
現在又說要用Zoom.
都不知道何時才能見面.
那種感受落差.
很不開心.
很無奈.
現在都不知道還要等多久.
接下來前書五章九節說.
因為神不是預定我們受刑.
是預定我們藉著耶穌基督得救.
所以我勉勵大家.

$^{201}$記住.
神不是預定我們承受新冠狀病毒的痛苦.
是預定我們在基督耶穌裡面得救.
你可能說我如何應付.
我自己和大家一樣.
每天都會受外在的事影響感受.
情緒可以起起伏伏.
每個人都是一樣.
不過.
我們可以選擇.
就是去祈禱.
求神讓我們能夠去處理.
或者控制.
我們自己的情緒.
因為我們如果不開心.
我們覺得沒有動力.
我們其實是在消耗我們的時間.
又或者我們會令到身邊的人.
添一些煩惱.
倒不如我們更加祈禱.
求神讓我們善用每一分一秒.
去做可以做的事.
因為時光飛逝.
我經歷一段時間都覺得很混亂.
生活真的很不正常.
但是後來想想.
如果這樣下去不是辦法.
所以很認真的祈禱.
我又發覺神真的聽我們的祈禱.
結果可以自己恢復正常.
雖然外在的生活不正常.
我們現在的情況都常態化了.
但是我們自己調節到生活正常都可以.
可以不受外在的環境影響你那麼多.
每一天我都數著做了什麼.
多讀聖經.
將你看聖經的成果寫下來.
然後跟別人分享.
我見到很多弟兄姊妹都是這樣.
他們在WhatsApp群組裡.

$^{241}$分享很多有造就性的文章.
非常非常好.
近來因為早一段時間.
七月底.
全灣葵青區有他的培靈會.
我認識的傳道人都很忙.
八月有港九培靈研經大會.
我認識的弟兄.
他幫忙做翻譯.
還從廣東話翻譯到普通話.
聽聞大家很欣賞他.
我想大家可以用時間.
在這方面.
有人說.
從早看到晚.
由研經會到培靈會.
都是一個很豐盛的筵席.
第二部分我想說的是.
患難是擺在面前的.
我以下所說的.
是有經文的根據.
我會告訴你為什麼我會這樣解釋.
在《帖撒羅尼迦前書三章三節》說.
免得有人被豬搬的患難搖動.
因為你們自己知道.
我們受患難.
原是命定的.
是什麼意思呢?.
我們要看經文的背景.
《帖撒羅尼迦前書三章三節》說.
我們既不能再忍.
就願意獨自等在雅典.
打發我們的兄弟.
在基督福音上.
作神執事的提摩太前去.
去兼顧你們.
並在你們所信的道上.
勸慰你們.
為什麼保羅不能再忍呢?.
因為保羅.

$^{281}$他在離開《帖撒羅尼迦》之後.
有猶太人去迫害教會.
《帖撒羅尼迦》的弟兄姊妹.
他面對很多衝擊.
所以保羅就說.
去猜險提摩太.
去帖撒羅尼迦教會.
去兼顧這班弟兄姊妹.
保羅所指的患難是什麼呢?.
是主要猶太人.
去敵視那些.
傳耶穌基督福音的弟兄姊妹.
不是後來羅馬帝國.
去迫害信徒的行動.
其實《新約聖經》有不同的時期.
不同.
大大小小的迫害都有.
命定這個詞.
希利文是matter(有關).
《新約聖經》出現了24詞.
是說不同的情況.
值得討論的問題是.
究竟matter(有關)這個詞.
是應該譯作命定.
還是事情就是這樣發生呢?.
以下我舉一批經文給大家.
讓你了解一下.
究竟是什麼回事.
首先是路加福音第二章第34節.
西面為他們祝福.
又對孩子的母親瑪利亞說.
「這孩子彼立」.
是要叫以色列中許多人跌倒.
許多人興起.
又要作偽旁話柄.
叫許多人心裡的異念顯露出來.
你自己的心也都像刀刺頭一樣.
彼立.
「彼立」這個字可以譯作上帝命定.
但也可以譯作「被建立」.

$^{321}$剛才那裡不是「被立」.
那個字就是這裡了.
就是這個字.
另一段經文就是.
《馬太福音》第28章第6節.
「他不在這裡.
照他所說的已經復活了.
請門來看安放主的地方」.
「安放」是「Akaito」.
又可以譯作上帝命定的意思.
但也可以是很平實的.
「被放在那裡」的意思.
有兩個意思都可以.
另一段經文.
《路加福音》第3章第9節.
「現在虎子已經放在樹根上.
凡不見好果子的樹就撼下來.
凋在火裡」.
「放在」「Kaitai」.
也可以譯作上帝命定的意思.
但也可以很平實的.
「就是放在那裡」.
而《帖撒羅尼加前書》第3章第3節.
「免得有人被豬般患難搖動.
因為你們自己知道.
我們受患難原是命定的」.
在這裡我看了這麼多節經文.
再看《上文下理》.
我覺得這節經文.
其實可以譯作.
患難原是放在我們面前的.
可以這樣譯.
為什麼我這樣說呢?.
譯那一節很重要的經文.
就是希伯來書第12章第2節.
「仰望為我們信心創始成中的耶穌.
他因那擺在面前的喜樂.
就輕看羞辱.
忍受了十字架的苦難.
便坐在神寶座的右邊」.

$^{361}$「擺在前面」.
Proclimenes這個字的字根.
Prokaitai.
Pro的意思是「預先」.
Prokaitai可以譯作什麼呢?.
是「預先安排」.
帖撒羅尼迦前書第3章第4節說.
「我們在你們那裡的時候.
預先告訴你們.
我們必受患難.
以後果然應驗你們也知道」.
預先告訴是另一個字.
Prolegomen.
本來是預先告訴帖撒羅尼迦教會的弟兄姊妹知道.
患難是擺在你們面前的.
當大家遇到患難的時候.
你就不要動搖了.
經文的意思是.
你去傳耶穌基督的福音.
你必然會遇到患難.
當然我們可以推論.
所有人生的事都是神命定的.
這個一定對的.
不過帖撒羅尼迦前書第3章第3節的意思是.
當你帖撒羅尼迦教會的弟兄姊妹.
你去傳耶穌基督的福音的時候.
你應該有心理準備.
有猶太人會加給你們患難的.
我們都應該聽過.
人生不如意的事.
十常八九.
我們很清楚知道.
我們所活著的世界.
不是完美的.
早一段時間.
歐美的國家領袖.
都染上了新冠狀病毒.
你想一想.
一般人染病有什麼稀奇呢?.
歐美國家的醫療體系是很完善的.

$^{401}$但是都應付不了.
這麼大量的新冠狀病毒確診的病人.
我們認為可以應付的挑戰.
結果發現是怎樣呢?.
是遠遠超過人能夠承擔的能力範圍.
當我們覺得每一件事.
都是在我們控制之下的時候.
我們當然不會接受.
患難是放在面前.
什麼都能解決的.
現在我們應該很明白.
不是神預定我們面對新冠狀病毒的威脅.
但是新冠狀病毒的威脅.
就是這樣放在我們面前.
現在還在說.
世衛還在說.
可能這種情況會持續幾年.
這樣的時候你就會發覺.
我們的生活有很多變化.
當然有很多人都還不能夠內心安定.
其中一個原因是他們還沒有認識神.
在這段時間我自己的領受就是.
當世界越是混亂的時候.
教會就是做事的時候.
我這樣說其實上人不是那麼容易明白.
好像我們很黑心那樣.
不是.
因為耶穌來就是要將救恩給人.
為什麼?因為人在一個罪惡黑暗的生活下.
不過人是不知道的.
所以祂拯救我們.
讓我們脫離黑暗.
現在這個新冠狀病毒.
令我們體會到.
我們根本不可以控制這個世界.
我們並且很無能,很無助.
我們甚至不知道將來會怎樣.
如果我們要面對這樣的困境.
我們一定要知道.
誰掌管這個世界.

$^{441}$我們的把握在哪裡.
這樣說明顯地教會要做事.
教會要傳福音.
好了,你和我都是基督徒.
我們和世人一樣.
我們每天都面對這些事的時候.
如果我們和世人一樣.
好像沒有什麼盼望.
或者很煩惱.
或者很無助的時候.
你又想一想.
我們又能做些什麼呢?.
我們做不到什麼,教會又能做些什麼呢?.
所以是不是應該在這個時刻去想.
神祂想我們怎樣.
這是很重要的問題.
神祂很想我們有見證.
把握機會傳福音.
但是往往就是因為我們自己.
沒有在這個時刻回歸於神.
我們和其他人一樣.
哎呀,很煩啊.
其實人之常情.
我們不要怪責.
但是我們可以互相提醒.
在這個時刻我們再看清楚.
患難擺在我們面前.
我們應該回歸於神.
不是神要我們預定我們受苦.
這個很重要.
我希望大家互相勉勵.
第三部分就是神的旨意是人的聖潔.
你看看《貼沙羅尼加前書》第四章第三節.
就很清楚.
這叫經文行.
神的旨意就是要你們成為聖潔.
遠避淫行.
面對新冠狀病毒威脅的時候.
我們應該去問.
神對我們人生有什麼旨意呢?.

$^{481}$究竟我們有一個什麼樣的生活方式呢?.
現在大部分人.
應該是全部人在香港都要戴口罩.
使用潔手液.
我們都知道細菌病毒是很可怕的.
你看不到它.
那些是最可怕的.
細菌病毒你看不到它.
但是它真實存在.
還可以人傳人.
還可以令你的肺部受損.
現在又說可能腦部也受損.
甚至奪去生命.
雖然我們都不太清楚.
這些病毒是什麼原因出來.
是不是跟吃蝙蝠有關.
但是你想想.
人吃野生動物的目的是什麼呢?.
其實說真的我們很幸福.
我們可以吃很多種肉.
有雞肉,牛肉,豬肉,羊肉等等.
其實為什麼人一定要吃野味呢?.
是不是?.
新冠狀病毒的出現.
令我們醒覺.
人類口腹之肉.
可以帶來人類的災難.
現在我們每個人都戴著口罩.
說真的呼吸都要用很大氣力.
所以現在有人說.
我們要留一口氣亂逃.
不說這麼多無謂的事情.
是不是這樣?.
我們也會這樣想.
就是說我們平時說了很多都沒什麼用.
是不是這個意思?.
你可以想想.
我們平時說的話有多少價值呢?.
我們的日常生活有沒有規律呢?.
這個問題又來了.

$^{521}$現在我們越來越注重規律.
你又要飲食又要各種東西.
希望我們真的要醒覺.
我們也知道生活要有規律.
要乾淨.
但是你想想.
只是手忙腳亂.
每個人都知道.
你要把心也結在一起才行.
是不是?.
在彈性上班的方式下.
我們會發覺原來生活.
習以為常的模式.
是可以變的.
那裡工作上的影響.
會讓我們有很多感受.
我們會感覺到自己的工作.
並不是不可取代的.
因為你現在可以安排在家中工作.
又或者發覺自己.
在某一些範圍裡面.
你仍然是不可取代的.
但是總的來說.
這個危機就在這裡了.
究竟我有什麼必然的存在價值.
大家會看到.
很多行業是受影響的.
零售業.
或者飲食業等等.
我們每個人都努力求生存.
我看過一些報導.
有些自由藝術人.
現在不能演出.
他就沒有錢收入了.
有一個女士.
因為有小孩.
她要從事小巴的清潔工作.
還要晚上出工去開工.
我想我們可以想像到.
很多人在擔心生活上的事情.

$^{561}$我們希望神保守大家的生活.
我也聽過一些人說.
他做一些設計工作.
他不用上班是沒有工資的.
每個月只剩下四分之一的工資.
聽到其實很慘.
我們應該為弟兄姊妹祈禱.
求福.
不過要用帖撒羅尼迦前書.
四章三節去為他祈禱.
神的旨意就是要你們成為聖潔.
遠避淫行.
去勉勵弟兄姊妹.
在這個時候去問.
怎樣可以調教我們的生活呢?.
怎樣去調教我們的內心呢?.
怎樣去讓我們活出聖潔呢?.
有些牧者就說.
新冠狀病毒就是很明顯.
神對人的罪惡的審判.
這是要讓人回歸於神的機會.
是對的.
有牧者就說.
不要浪費了新冠狀病毒的機會.
對的.
我們不應該只是想逃避病毒這麼簡單.
我們要思想一些深層的事情.
我們不單要想到衛生上的潔淨.
同時要想到道德上的聖潔.
屬靈上的聖潔.
當我們看到香港.
早一段時間叫那些人從疫區回來的時候.
要自我隔離14天.
但是你發覺有些人不聽話.
他覺得在家裡很悶.
就自己走了出去.
你會發覺原來人的自律是一個問題.
縱使有規矩.
但是人又想辦法不守規矩這麼有趣.
自律背後就是一個功德心.

$^{601}$向社會負責的一個問題.
你發覺那些人應該要留在家裡.
卻是四處走的那些人.
其實沒有什麼功德心.
為什麼沒有什麼功德心呢?.
因為他沒有向神負責的意識.
沒有罪疚感.
他只是說我做錯了.
沒有人發現就沒有問題.
他覺得沒有人發現就沒有問題.
但是他不想到你做錯了就是做錯了.
什麼是問題呢?.
當自己被人發現做錯的時候.
很多人都是這麼說的.
被人說一句就不好.
這是什麼問題呢?.
是面子的問題.
不是說沒有人發現你就沒有問題.
被人發現說一句就不好意思.
那個錯其實在神的面前就是錯的.
這就是因為他沒有聖潔的觀念.
沒有神又沒有標準.
他自己就是標準.
在新冠狀病毒的情況下.
我們看到很多人性的東西.
我們應該要好好謹記.
貼沙羅里加前書四章三節.
那個教訓就是.
神的旨意是要你們成為聖潔.
遠避淫行.
這個真的很值得我們去思想.
我聽聞現在的人看醫生都少了.
因為我們每個人都戴了口罩.
傷風鼻塞好像也少了.
很希望我們靠著我們的神.
去過每一天.
不是無奈地去過每一天.
不是被動地去過每一天.
而是我們每一天都有充實的.
以及可以做到一些對神.

$^{641}$對家人.
對他人有價值的事情.
總結來說.
貼沙羅里加前書第五章第九節說.
因為神不是預定我們受刑.
乃是預定我們直著耶穌基督得救.
我們應該謹記的是什麼.
我們在神裡面.
神預定我們在直著耶穌基督得救.
我們的神給我們的恩典是很豐富的.
而貼沙羅里加前書第三章第三節說.
免得有人被蛛絆患難去搖動.
因為你們自己知道我們受患難.
原是命定的.
這節經文不是告訴我們無奈.
不是要去訴苦.
而是說一個現實的情況.
患難就是擺在我們面前.
我們這個世界是一個墮落了的世界.
我們就不要被蛛絆的患難搖動.
貼沙羅里加前書四章三節的教訓.
神的旨意是要你們成為聖傑.
遠避淫行.
神的旨意是要我們有聖傑的人生.
我們不單單想到衛生上的潔淨.
同時要想到道德上的聖傑.
屬靈的聖傑.
最後我再說兩句.
其實在這段時間.
很多很多人的內心情緒.
是承擔不了的.
所以那些減壓.
那些叫我們遇到難處.
找人聊天訴苦.
不要冤枉自己的訊息.
百分之百是對的.
我完全認同的.
但是我還是很想鼓勵大家.
我們不單是將人帶到人的面前.
我們要將人帶到神的面前.

$^{681}$因為我們要靠神的話語.
來承載應付我們內心裡.
很複雜的情緒.
那些情緒好像波濤洶湧.
我們找朋友聊天是很對的.
找家人聊天是很對的.
一定有作用的.
有人了解你是好的.
但是你想解決問題.
大家都知道最後是自己去面對的.
但是我憑什麼有力量呢?.
力量從何而來呢?.
我們應該要想.
是從造天地的耶和華而來.
所以我是很鼓勵大家.
記住這個時候.
福音的訊息是很重要的.
將人帶到神的面前.
讓神親自在人的心裡說話.
而神的話一句勝過我們千言萬語.
我們相信神的話有能力.
那不如就讓神的話成為眾人的幫助.
當然我們當然要安慰人.
我們要鼓勵人.
這肯定是對的.
這沒有錯的.
我們鼓勵人不要灰心.
減壓.
這些完全是對的.
但是你要記住.
這些話只不過是人的說話.
帶來一定的影響力.
但是沒有決定性.
我要說的是神的話.
神自己的話.
神的應許.
他才可以解答人心裡的謎團.
所以我很勉勵大家.
在這個時候.
什麼時候都要想.

$^{721}$回到神那裡.
讀聖經.
祈禱.
默想經文的意思.
我們一起祈禱.
我們親愛的天父上帝.
因為你的說話是我們生命的力量.
我們知道在一個很困難的時候.
我們當然是很需要彼此關懷.
互相勉勵說一些正面的話.
但是真真正正有能力的.
那些話語從哪裡來呢?.
是從你那裡來的.
我們懇求你.
賜下你的恩賢.
賜下你的說話.
成為我們生活的力量.
讓我們經歷到你的話語的能力.
經歷到你的聖靈的能力.
讓我們面對著外面.
像波浪澎湃的情況的時候.
我們的內心就像耶穌基督.
在船上平靜風浪的時候.
一樣的穩定.
我們是人.
我們不是神.
我們唯有在你的面前去求.
你說過你會將我們不能得.
又不配得的平安給我們.
讓我們面對.
擺在我們面前的患難的時候.
我們真的經歷到這一種的平安.
以至我們更加肯定的.
神你是真實的.
也讓我們紀念那些很艱苦的人.
很多現在的工作.
朝不保夕的人.
很多患了病的人.
或者有些人在醫院.
他的家人探不到的人.

$^{761}$我們紀念很多比我們更加慘的人.
讓他們體會到神的愛.
我們的祈禱奉耶穌基督的名.
阿們.
\newpage



\section{哈該書 2:10-19}
\label{sec:N89ZkrKQWlk}
\textbf{2020.08.09《生命的回轉與盼望》哈該書 2\_10-19( 和修版)  老耀雄傳道}
\newline
\newline
連結: \href{https://youtube.com/watch?v=N89ZkrKQWlk}{\texttt{https://youtube.com/watch?v=N89ZkrKQWlk}} ~~~~ 語音日期: 2020-08-20
\newline
\newline
\hyperref[sec:iY8ed2cNTa4]{\small{< < < PREV SERMON < < <}}
~
\hyperref[sec:index_chronic]{\small{[返順時目]}}
~
\hyperref[sec:index_scriptual]{\small{[返順卷目]}}
~
\hyperref[sec:UM_spSYKywQ]{\small{> > > NEXT SERMON > > >}}
\newline
\newline
哈該書 2:10-19
\newline
\begin{longtable}{cl}
\hline
\hline
章節 & 經文 (和合本修訂版)\\
\hline
2:10 & \begin{tabularx}{0.7\textwidth}{X} 大流士王第二年九月二十四日，耶和華的話臨到哈該先知，說： \end{tabularx} \\ \\ \relax
2:11 & \begin{tabularx}{0.7\textwidth}{X} 「萬軍之耶和華如此說，你要向祭司請教律法，說： \end{tabularx} \\ \\ \relax
2:12 & \begin{tabularx}{0.7\textwidth}{X} 『看哪，若有人用衣服的邊兜聖肉，這衣服的邊接觸了餅，或湯，或酒，或油，或別的食物，這些是否成為聖呢？』」祭司回答說：「不。」 \end{tabularx} \\ \\ \relax
2:13 & \begin{tabularx}{0.7\textwidth}{X} 哈該又說：「若有人因摸屍體染了不潔淨，然後接觸任何東西，這東西就變為不潔淨嗎？」祭司回答說：「必不潔淨。」 \end{tabularx} \\ \\ \relax
2:14 & \begin{tabularx}{0.7\textwidth}{X} 於是哈該說：「耶和華說，在我面前這民如此，這國也是如此；他們手裡的各樣工作都是如此；他們在那裡所獻的都不潔淨。」 \end{tabularx} \\ \\ \relax
2:15 & \begin{tabularx}{0.7\textwidth}{X} 「現在，你們心裡要想一想，從今日起，耶和華的殿還沒有一塊石頭放在石頭上的情況。 \end{tabularx} \\ \\ \relax
2:16 & \begin{tabularx}{0.7\textwidth}{X} 那時你們怎麼了？有人來到二十斗的穀堆那裡，卻只得了十斗；有人來到酒池那裡要取五十桶，卻只得了二十桶。 \end{tabularx} \\ \\ \relax
2:17 & \begin{tabularx}{0.7\textwidth}{X} 我以焚風、霉爛、冰雹攻擊你們，和你們手上的各樣工作，你們仍不歸向我。這是耶和華說的。 \end{tabularx} \\ \\ \relax
2:18 & \begin{tabularx}{0.7\textwidth}{X} 你們心裡要想一想，從今日起，就是從這九月二十四日起，從立耶和華殿根基的日子起，你們心裡想一想： \end{tabularx} \\ \\ \relax
2:19 & \begin{tabularx}{0.7\textwidth}{X} 倉裡還有穀種嗎？葡萄樹、無花果樹、石榴樹、橄欖樹雖沒有結果子，從今日起，我必賜福。」 \end{tabularx} \\ \\
[1ex]
\hline
\hline
\end{longtable}
$^{1}$洪恩嘉的弟兄姊妹平安.
在疫情之下我們相信一段時間.
在這樣的情況下做了崇拜.
有一次英國某一間報紙舉辦了一個獎金豐富的處境問答遊戲.
題目是我們一起參與一下也可以.
一個處境問答遊戲.
題目是有三個科學家承擔著一個熱氣球環遊世界.
但處境是熱氣球正在漏氣.
他們決定終止這個行程.
而且快快脆脆地掉到地面.
不過可惜負膀太重.
如果負膀不負重的話熱氣球會直接掉到地面.
所以一定要把所有物品都扔掉.
因為要減重.
但扔掉所有物品也好.
仍然是太重.
他們需要在三個人之中選出一個.
然後殘忍地把他扔出熱氣球.
犧牲一個救回兩個.
以減輕熱氣球的重量.
為了讓大家對三個科學家有認識.
所以公佈了這三個科學家的背景.
第一個是環保專家.
這個環保專家的研究可以令人類免受環境污染而面臨死亡.
大家知道現在這個世界環境污染很嚴重.
第二個是核能專家.
他有能力阻止全球性的核戰爭.
使地球不會遭受滅亡的厄運.
我們知道現在核子戰爭的機會仍然存在.
甚至核電廠如果洩漏了輻射或爆炸.
造成的災難也會很全面.
第三個是糧食專家.
這個糧食專家有一個能力.
他可以在不毛之地.
即是在沙漠地帶.
種植多種糧食.
使人類不會遭受飢餓的厄運.
參加處境遊戲的人.
他要回答一個問題.
你以甚麼理由將哪個科學家調出熱氣球呢?.

$^{41}$由於報紙宣傳廣泛.
獎金豐富.
很多人寄了答案.
有些是大學教授.
有些是退休學者.
提出方方面面種種自然科學理據.
證明調哪個科學家出熱氣球.
會對世界最有利.
不要一個環保專家會好一點.
不要一個核能專家會好一點.
不要一個糧食專家會好一點.
或者調轉頭.
保留哪些科學家對地球最有利呢?.
我們心目中有沒有答案?.
報刊最後選中的答案.
令很多人大跌眼鏡.
因為獎金由一個十歲的小朋友獲得.
他的答案很簡單.
將最肥身體最重的科學家調出熱氣球.
是令熱氣球減重的最佳方法.
消息公佈後.
大家認為有心作弄.
玩弄嗎?.
這麼簡單的答案.
為何還要宣傳去詢問呢?.
報刊後來出了一個聲明.
解釋這次遊戲的目的.
不是想玩弄你們.
是想了解到.
人在一個特殊的處境下.
能否直接看到問題的核心.
而沒有被其他東西.
或次要的資料.
側看.
即是說.
我們在一個特別的環境下.
能否直接看出問題.
而沒有其他東西阻礙我們.
雖然有些讀者不太同意.
但最後的結果.

$^{81}$就是不了了之.
其實很多人都不能正確了解問題的核心.
就像舊約.
以色列人在被擄歸回的處境下.
同樣有一個問題.
令他們百思不得其解.
這個問題是甚麼呢?.
就是這批以色列人.
仍然有敬拜.
仍然有獻祭.
但為何得不到神的祝福呢?.
為何生活得這麼困苦呢?.
他們不明白.
直至才知道哈該的出現.
他們終於知道各種原因.
今天我們就透過哈該書2章10至19節.
從以色列人兩次對信仰的回顧.
去檢視我們信仰的優先次序.
在講道之先我們一起有個祈禱.
「真主,當我們聆聽你話語之先.
我們需要你的寬恕和赦免.
因為我們縱然已經認識你.
但我們仍然生活在一個以自己為中心的信仰生活.
我們不要將主放在我們生命裡最高處.
求因主潔淨我們每一個人.
除去我們心裡的罪.
好讓我們能夠守潔心清地聆聽你的話語.
然後按你的心意去過我們的信仰生活.
我們祈禱是奉主耶穌基督的聖名祈求.
阿們.
我們首先去了解當時的以色列人.
處於一個什麼處境之下的生活.
在公元前586年.
巴比倫尼布格尼撒攻入耶路撒冷之後.
直接滅了南國.
因此整個以色列基本上已經滅亡了.
尼布格尼撒將以色列人擄去巴比倫.
百姓正式成為一個亡國奴.
直到公元前539年.
波斯王古列攻陷巴比倫.

$^{121}$當時估計有100萬猶太人.
散居在波斯不同的地方.
以色列人無論在巴比倫.
或之後的波斯這個地方.
其實已經住了差不多50多年.
很多以色列人已經在這裡落地生根.
50年有兒有女.
在古時來說50多歲的孫子也有了.
當波斯王古列准許猶太人返回祖國耶路撒冷.
重建聖殿的時候.
很多人根本就不想再返回耶路撒冷生活.
為什麼.
因為很多以色列人已經適應了新的生活.
不想離開波斯.
而且波斯的生活和文化水平.
比當時落後的耶路撒冷好很多.
住慣城市.
叫你去住鄉村.
也是一個挑戰.
根據以色列記第二章記載.
當時所羅巴伯只帶領著42,360人返回耶路撒冷.
你算一算.
100萬散居在波斯的以色列人.
只有42,360人願意返回.
其實5\%也不夠.
那你說這5\%是不是很愛神.
也應該是.
這42,360人.
離開了相對文化和質素高的波斯.
然後去到落後的耶路撒冷.
他們都是愛主.
他們歸回之後第一件事就是重建聖殿.
很可惜.
建了兩年就停工.
停了多久.
一停就停了16年.
停了十年後.
神興起先知哈該.
再一次考慮所羅巴伯和以色列人.
透過哈該.

$^{161}$向以色列人說了四個訊息.
要求以色列重新歸回到神的身邊.
今天港圖經文二章十節十九節.
就是哈該向以色列人說的第三個訊息.
這次訊息有一個特點.
特點是甚麼呢.
就是要求以色列人.
以兩個日子為界限.
或是一個標準.
去回顧過去16年的生活和展望將來.
明白了經文的背景後.
我們開始進入經文.
哈該對以色列人的第一次信仰回顧.
就是在6月24日之前.
6月24日之前.
以色列人是甚麼呢.
就是以色列人不將神放在守衛.
引致神的攻擊.
經文第一次提到回顧在二章十五節.
他說現在你們心裡要想一想.
新譯本亦做了反省.
原文是指放棄你們內心.
所以想一想是一種從內心去檢視的意思.
十五節下.
從今日起.
耶和華的殿環會有一塊石頭.
放在石頭上的情況.
這次經文是根據上文一章十五節所指.
百姓為耶和華的殿開始工作的日子.
經文所指的今天.
就是大利烏王第二年6月24日.
哈該要求這班42360個歸回的以色列民.
甚至當時居住在當地耶路撒冷的居民.
認真地去內心去檢視一下.
重建聖殿之前的日子.
究竟他們的信仰狀態是甚麼.
當時百姓對神的事毫不在意.
即是沒有將神放在他們生命的首位.
所羅巴伯在波斯王古列元年.
即公元538年歸回了耶路撒冷.

$^{201}$隔了一年就開始重建聖殿.
但到了大利烏王第二年的6月24日.
聖殿仍然未建成.
因為重建的工程停了16年.
十五節下.
耶和華的殿還沒有一塊石頭放在石頭上的情況.
百姓讓神的殿荒涼.
聖殿連最底部的根基都未建造好.
聖殿的冷清亦是回應一張九折.
因為我們的殿荒涼你們各人卻顧自己的房屋.
下改指出並不是百姓沒有能力去重建聖殿.
而是百姓只顧自己的房屋.
自己的生活多過想重建聖殿.
百姓沒有將神的事放在心裡.
百姓縱然只顧自己的事.
努力工作又如何呢?.
經文說百姓想得到二十斗菊.
但只得到十斗.
有修成但只是期望的一半.
葡萄是釀酒的主要材料.
他們想釀五十桶酒.
但只得二十桶.
修成只得四十個百分比.
一半都不到.
對一張九折說你們盼望多得.
但所得的卻少的一種回應.
以色列人回歸的目的.
本身是要重建聖殿.
但現實上他們不再理會聖殿的重建.
只專注於自己的生活.
縱然多努力也好.
但也得不到預期的效果.
以色列人不單對重建聖殿冷淡.
就連神憤怒的感覺也麻木.
神以焚風,煤爛,冰雹.
三種自然災害來懲罰百姓.
其實在摩西時代.
神也曾經用薄災去懲罰埃及人.
《申命記》亦提到.
如果以色列人行惡理希耶和說.

$^{241}$神會以焚風,煤爛來攻擊以色列.
所以如果以色列人能以史為鑒.
就能知道天災背後.
神正在等他們的迴轉.
大經文說.
你們仍不歸向我.
究竟是百姓心裡像發露般僵硬.
還是他們不知道,不明白呢?.
三節經文裡可以簡單說到一個事實.
就是人對神的嗜好不關心.
連神憤怒的攻擊也察覺不到.
甚至不迴轉.
只顧著自己.
只顧著自己的房屋.
只顧著自己的生活.
為何以色列文會淪落到這個地步呢?.
在二章十至十四節.
神透過哈該向百姓解釋了原因.
哈該向祭司發問了兩個聖結的條例.
都是律法的一部分.
第一個律法在十二節.
說到伊甘勝在聖肉.
然後伊甘再接觸到餅,湯,酒,油.
和其他食物.
這些被伊甘接觸過的東西.
這些東西可否當做聖物呢?.
祭司的回答是不算.
只有被聖肉直接接觸過伊甘才是聖物.
而透過伊甘而接觸到的一切物件.
不能視為聖物.
這個結證條例記載在利未記六章二十七至二十八節.
十三節.
說到人若接觸到屍體而污穢.
其實律例裡不單止接觸屍體.
連接觸到死人的骨頭.
或者走過的墳墓都一樣.
污穢即不潔.
接觸了這些東西之後.
接觸了屍體或者死人的骨頭.
或者墳墓之後再接觸其他物件.

$^{281}$這些物件算不算不潔呢?.
祭司的回覆是都成為不潔.
這些物件無需要直接接觸到屍體.
只需要間接透過人而接觸到.
它就已經成為一個不潔之物.
這個條例在利未記十一章二十八節.
以及文掃記十九章二十二節有記錄.
透過兩個律法的比較.
法性知道不潔的傳播力.
比潔淨的傳播力更大.
成聖的物件要直接與源頭有接觸才可以.
而不潔的物件只需要間接的接觸.
就已經成為不潔.
成聖在當事人一定會知道.
因為它是真實地接觸過那個聖物.
但是成為不潔呢?.
當事人可能會不知道.
因為不潔的物件的產生.
是無需要與不潔的來源有直接的關係.
只需要間接的接觸就已經被污染.
跟著哈該就講出了以色列人的核心問題.
十四節.
於是哈該說.
於是是一個轉接詞.
原文亦叫為者緣因.
即是為了上述的緣因.
加強了上文和下文的因果關係.
為了上述的緣因.
耶和華說.
在我面前這民如此.
這國也是如此.
這民是指以色列的所有人民.
這國當時已經沒有南國和北國了.
不過仍是指以色列整個民族.
也是如此.
根據上文對不潔之物的理解.
整個民族所有人民在耶和華面前.
都是不潔的.
由於人民本身已經是不潔.
因此他們手上所做的工.

$^{321}$都是不潔.
手下所播的種.
所栽種的樹.
所有的收成.
都是不潔.
就算百姓獻上.
在壇上獻祭.
所獻的祭物.
都是不潔.
百姓不潔的原因.
就是他們對神的事毫不重視.
不將神放在生命中的首位.
對聖殿的重建.
頭閒之心.
社群所犯的罪.
使全體百姓的生產力.
和他們的敬拜.
和建制.
全部變為不潔.
十六年以來.
他們生活在神的憤怒當中.
他們仍然不知道過終的原因.
直至此刻.
他們才如夢初醒.
我在舞會教初信舞蹈班的時候.
很多時候都會用這個例子作為比喻.
我們有一杯清水.
清水雖然經過100\%高溫消毒.
很值得我們去喝.
但只要我們在這杯水上.
滴一滴坑渠水進去.
一滴而已.
這杯100\%經過高溫消毒的水.
滴一滴坑渠水進去.
你還會不會喝?.
整杯清水就是因為一滴坑渠水而污染了.
而杯子本身是裝乾淨的水.
但因為水被污染.
杯子也變成了.
不是裝乾淨水的杯子.

$^{361}$而是裝一個污水的杯子.
罪就是那滴坑渠水.
我們每一天滴了多少坑渠水進我們的杯子裡.
在我們的生命裡.
我們每一天一滴.
每一天一滴.
那杯子的水.
是固定的.
每一天滴一滴.
污染會越來越嚴重.
被老歸回的以色列人所犯的罪就是.
不將神放在生命中的首位.
而一滴坑渠水滴了16年.
各位紅衣人家弟兄姊妹.
我們的生命是不是和以色列人一樣.
只是看到自己的需要.
而看不到神的需要.
我們是否也有這樣的生命.
而看不到神的需要.
我們對神的提醒有沒有敏感度.
神可能用不同的處境或不同的方法.
不同的人或事去提醒我們.
叫我們歸回祂身邊.
我們敏不敏銳神對我們的提醒.
因為我們聽而不見.
視而不聞.
我們是否只顧著解決目前的困境.
而看不到問題的核心.
盼望我們可以在這一刻.
我們好像黑丐提醒以色列人一樣.
我們再次檢視一下我們的生命.
重新檢視一下我們的生命.
我們有沒有將神放在我們生命中的首位.
還是神在我們生命中只是可有可無.
黑丐才知對以色列人要求第二個回想.
就在9月24日起.
回轉之後對神傳一個敬畏的心.
和充滿盼望.
英文第18節提到你們心裡要想一想.
從今天起就是從9月24日起.

$^{401}$黑丐這次要求以色列民從9月24日開始.
再一次用內心去檢視一下.
他們開始重建聖殿之後的改變.
9月24日是什麼日子?.
就是聖殿奠基的日子.
為聖殿建立根基的一個儀式.
不再是十五節那種沒有一塊石頭.
放在石頭上的情況的光景.
百姓在6月24日開始重建.
三個月之後就舉行了奠基的儀式.
這三個月看到百姓的心影回轉.
在一章十二節提到.
百姓也在耶和華面前傳敬畏的心.
百姓對神是心傳敬畏.
而不是對神投閒置散.
百姓願意心意回轉.
不再只是掛念自己的事.
而是更加紀念神家的事.
為了神有一個好的居所.
開始重建聖殿.
百姓傳敬畏的心.
促使他們加倍專注.
在聖殿重建的事工上.
聖殿最後在四年之後.
大利烏王第六年亞達月初三.
即是公元前515年3月12日建成.
百姓在聖殿建成後.
隨即舉行如意節的慶祝.
聖殿的建成對整個以色列民來說.
是有另一個重要的意義.
在《以色列記錄》將21節.
新譯本亦這樣說.
可以被老歸回的以色列人.
和所有脫離當地民族的污穢的人.
一同吃者楊糕.
尋求耶和華以色列的神.
即是百姓可以由不潔淨成為潔淨.
聖殿的建成使他們可以在.
波斯異族統治之下.
仍然可以為耶和華.

$^{441}$能夠活出分別為性的生活.
在異族統治之下.
仍然可以為耶和華.
活出分別為性的信仰生活.
在14節耶和華對以色列文化中.
也是不潔的指責.
即是不重建聖殿.
只顧自己手中的東西.
有憲制而沒有敬畏的心.
由於他們願意回轉.
在不久的將來.
能夠成為聖潔的子民.
這個意義對於當時的以色列百姓.
是很深遠和鼓舞.
深意的回轉.
除了看到未來潔淨身份的盼望.
也看到修成的盼望.
經文第十九節提到.
倉裡有穀種.
新列本認為倉裡有沒有種子.
過往在旱風,霉爛,冰雹的環境下.
影響了修成.
影響了生產.
修成自然就不好.
提問背後表示倉裡沒有儲糧.
即是沒有糧可種.
葡萄樹,無花果樹,石榴樹,橄欖樹都沒有結果字.
九月二十四日.
這個日子不是一個修割的季節.
葡萄樹,無花果樹,石榴樹的修成.
在八月份.
即是表示過了.
過了結果的時間.
橄欖樹就在十一月.
修成期未到.
在現實的情況下.
在眼見的情況下.
倉裡不但沒有種子.
樹上亦沒有果實.
一切都是不樂觀.

$^{481}$倉裡沒有東西可以種.
樹沒有果實.
可以讓你摘下來.
一切都是不確定,不樂觀.
不過盼望的根源在於最後那句.
從今天起.
我必賜福於你們.
當神將咒咒變為祝福.
表示神閱立百姓的迴轉.
所以眼前雖然見到倉裡沒有剩餘的種子.
眼裡見不到樹上有任何的修成.
不過盼望不是在於眼前所見到的景象.
盼望是在於我們見不到的將來.
因為神不再伸手攻擊.
神的賜福表示今後的日子必然修成豐盛.
百姓只需要安然等待.
等候神的工作便可.
黑鬼要求百姓對生命作一次內心的檢視.
就是要他們明白到迴轉的重要性.
以往百姓種種的不是.
迴轉了之後.
一切因著神的賜福而改變.
除了可以回復聖潔身份.
又可以地亦為他們效力.
日後不會再有要二十斗菊.
只得十斗菊的情況.
又不會要五十桶葡萄酒.
只得二十桶葡萄酒這種光景.
我每次在一些浸禮儀式裡.
我都聽到很多得救的見證.
每一個媳婦都講述她們生命裡的一些故事.
由不認識神.
或是敵視神.
到願意最後歸回神的大家庭.
見證分享林林總總.
但都有一個共通點.
共通點是甚麼呢?.
就是明白到生命只有回歸到神的身邊.
然後才可以體驗到神的賜福.
亦因著回歸到神的身邊.

$^{521}$然後才有各種不同的盼望.
各位紅衣家弟兄姊妹.
當以色列的百姓願意回轉的時候.
他以耶和華的事為首位.
百姓以重建聖殿來實踐對神的重視.
亦藉此得到神的祝福.
今日我們紅衣家不需要重建聖殿.
但我們需要重整我們信仰的焦點.
我們信仰的焦點是甚麼?.
我們所信的核心是甚麼?.
我們回教會的目的是甚麼?.
我們的信仰是否一份可以分別為性的生活?.
我們今日所擁抱的信仰.
是一個不可或缺的信仰.
還是一個可有可無的信仰?.
如何體驗或體現到我們對信仰的認真和投入?.
這一切都需要我們反覆思考和認真檢視.
今日講題是生命的回轉和盼望.
透過回顧.
以色列人明白到.
他們回歸耶路撒冷後.
沒有將神放在生命中的首位.
這種不潔的罪.
污染了他們一切手中所做的功.
甚至連獻的祭都被污染.
透過先知的提醒.
以色列人願意回轉.
並且實踐在重建聖殿一件事上.
以至他們縱然眼前看不到有甚麼修成.
有甚麼糧食.
但仍然滿懷信心盼望.
適合以色列民透過兩個日期回顧信仰.
我們可否像以色列人一樣.
同樣以一個日期作為我們的界線?.
以一個日期去回顧我們.
展望我們個人和教會信仰的實踐和使命?.
或者我大膽一點.
我們可否以疫情作為我們洪恩家的分界線?.
在教會方面.
疫情之前.

$^{561}$洪恩家有沒有使命的焦點?.
如果有.
有沒有具體的方案和實行的計劃?.
我們可以在疫情過後去無全清.
繼續發展下去.
如果沒有.
疫情過後.
我們能否確立洪恩家的使命?.
我們能否確立洪恩家的意象?.
意象和使命不是一個口號.
而是一個具體的計劃和時間表.
或者一個中長期的具體計劃.
或者一個可以評估進度的時間表.
這會否是我們洪恩家在疫情過後的盼望?.
眼裡看不到的.
盼望將來會有.
正如昔日的以色列人一樣.
他們看見的葡萄樹,無花果樹,石榴樹,橄欖樹.
全部都沒有果子.
但我們不只是看見眼前.
而是在神的祝福下.
在神的祝福洪恩家.
興起洪恩家之下.
我們可否有一個盼望.
看到日後的果實累累的將來?.
在個人方面.
信仰告訴我們.
信耶穌生命是豐盛的.
如果我們今天覺得生命是不豐盛的話.
我們會否反思一下為甚麼?.
是否我們心裡有一種隱藏.
未見的罪?我們察覺不到.
以色列人的罪.
讓他們手上一切的功.
都不被神所閱立.
我們有沒有一種隱藏.
在我們內心深處.
我們不知道的罪?.
無論是直接的,間接的,明知的,無意的.
這些罪都能導致我們不能享受神賜給我們的福.

$^{601}$神賜給我們的福是一個應許.
但如果因著我們的罪而享受不到.
是很可惜的.
經文說無論我們有多不潔.
只要我們願意回轉.
神就會閱立我們.
同樣的,約翰一書一章九節都說過.
我們若認自己的罪神是信實的.
必要赦免我們的罪.
洗淨我們一切的不義.
透過今天的信息.
願意我們洪順嘉每一個弟兄姊妹.
我們都能夠誠心在神面前.
坦然去承認我們有意,無意.
干犯了得罪神的罪.
在行為上我們能夠實踐對神的重視.
可以享受神賜給我們的豐盛人生.
我們一起去祈禱.
七日主我們知道你是聖潔的.
你絕對不容許.
你絕對不能夠接納一切不潔的事.
擺在你的面前.
因此我們任何不潔的思想言行.
都需要你的赦免.
求你讓我們有一個清潔的心.
求你讓我們將你擺在我們生命的首位.
讓我們尊主為大.
讓我們忠心地被你差遣.
讓我們信服你一切的大靈.
我們的祈禱是奉主耶穌基督的聖名祈求.
阿們.
願主祝福大家.
\newpage



\section{彼得前書 2:9}
\label{sec:UM_spSYKywQ}
\textbf{2020.08.16《不可以逃避》彼得前書 2\_9 陳志鵬傳道}
\newline
\newline
連結: \href{https://youtube.com/watch?v=UM-spSYKywQ}{\texttt{https://youtube.com/watch?v=UM-spSYKywQ}} ~~~~ 語音日期: 2020-08-20
\newline
\newline
\hyperref[sec:N89ZkrKQWlk]{\small{< < < PREV SERMON < < <}}
~
\hyperref[sec:index_chronic]{\small{[返順時目]}}
~
\hyperref[sec:index_scriptual]{\small{[返順卷目]}}
~
\hyperref[sec:xKlgziHaUYU]{\small{> > > NEXT SERMON > > >}}
\newline
\newline
彼得前書 2:9
\newline
\begin{longtable}{cl}
\hline
\hline
章節 & 經文 (和合本修訂版)\\
\hline
2:9 & \begin{tabularx}{0.7\textwidth}{X} 不過，你們是被揀選的一族，是君尊的祭司，是神聖的國度，是屬神的子民，要使你們宣揚那召你們出黑暗入奇妙光明者的美德。 \end{tabularx} \\ \\
[1ex]
\hline
\hline
\end{longtable}
$^{1}$弟兄姊妹早晨.
和雲鳳剛才的感受都是很接近的.
這首詩歌除了歌詞很能夠觸動到人之外.
這首詩歌也牽動到弟兄姊妹想起昔日的人和事.
其實有很多片段在腦海浮現.
當時大家很開心能夠在主的聖殿當中一起敬拜.
但今天就不能夠了.
因為我們仍然有限聚令.
我相信弟兄姊妹都習慣了崇拜的新常態.
就是在家中透過教會的直播去參與主日崇拜.
我盼望我們都一起學習.
我們學習有一個忍耐的心.
也操練我們常存盼望.
我們相信有朝一日我們能夠在教會再一次相聚.
大概五年前我曾經參加過一個活動.
是一個籌款暴行活動.
是奧比斯所主辦的一個暴行籌款活動.
這個暴行籌款活動和一般的暴行籌款活動有點不同.
不知道大家有沒有參加過一些暴行籌款活動.
一般的暴行籌款活動是在日間進行.
我們以前在紅銀堂也參與過.
我們去過青衣那裡有暴行籌款.
我們也曾經去過山頂的暴行籌款.
你不會在晚上去山頂去暴行籌款.
你也不會在晚上去青衣那裡去暴行籌款.
但是奧比斯的暴行籌款.
一反常態和平時的暴行籌款不同.
他就將暴行籌款的時間放在晚上.
他叫這個活動叫moonwalk萬俠行.
當時我服侍的教會大概是20歲到35歲的年輕人.
我也學習了耶穌基督道成肉身的服侍.
我進到他們當中.
參與了他們的團隊和他們一起去行.
一起去參與這個活動.
我們大概晚上九點鐘在長沙灣聚集.
這張照片就是我們在長沙灣聚集所拍的照片.
九點鐘在長沙灣聚集.
按照大會安排的路線.
我們就經過深水Po .
經過九龍塘.

$^{41}$然後就到大會指定的終點就是觀塘.
我們去到觀塘的時候大概是凌晨五點鐘.
奧比斯是一個關心失明人士的慈善團體.
他致力透過預防和治療失明.
去改變失明人士的生活.
正因為奧比斯有這個關心.
去關心失明人士.
所以在這個步行籌款當中.
他特意安排了一個環節.
這個環節就是參加者要蒙著眼睛.
會走一段路.
在這段路裡面會有很多障礙.
或者有人會噴水弄濕你.
最終目的就是這個機構期望.
參加者透過體驗失明人士的無助.
體驗失明人士的活在黑暗裡面的心情.
我相信在我們中間大部分的頂智妹.
都不是失明的.
大家的視力都是正常的.
但是其實我們都曾經活在失明的當中.
我們曾經活在黑暗的裡面.
我們被罪惡的權勢所合制住.
上帝差派耶穌基督去救贖我們.
將我們從黑暗的當中.
遷移到一個光明的地方.
就正如今天我們所看的經文.
《聖經》前書二章九節裡所說.
我上個月已經和大家說過.
在2020年的下半年.
我會和不同的經文一起思考.
聖經裡面關於教會的真理.
盼望我們眾人都能夠謙卑在神的面前.
一同去學習.
上次不知道大家還記不記得我分享過甚麼.
我和大家分享過教會的根基.
我用馬太科林第十六章十三至二十節的經文和大家分享.
當時彼得很清晰地認出耶穌基督的身份.
他說你是基督你是永生神的兒子.
然後耶穌基督怎樣和彼得說呢.
我還告訴你你是彼得.

$^{81}$我要將我的教會建造在這盆石上.
陰間的權並不能勝過他.
我說教會是建立在彼得對耶穌身份的認信之上.
說完教會的根基.
今天就想和大家說教會的使命.
教會存在在地上究竟有甚麼意義呢.
或者我們透過彼得前書二章九節.
一同去思考這個課題教會的使命.
由於只有一節經文.
可能你未必習慣一節經文.
只有一節經文.
一節經文不如我們一起讀一讀彼得前書二章九節.
預備起.
唯有你們是被揀選的屬類.
是有君尊的祭司.
是聖潔的國度.
是屬神的子民.
要叫你們傳揚那.
滲入奇妙光明者的美德.
在這裡一開始的時候彼得說.
唯有你們.
你們究竟是知道甚麼人呢.
你們是誰呢.
我相信大英姐妹都很聰明.
你們就是彼得前書的修信人.
這個說法沒錯.
不過我要再問.
彼得前書的修信人究竟是一班甚麼人呢.
彼得前書是彼得所寫的.
彼得在第一章第一節這樣說.
耶穌基督的使徒彼得.
寫信給那分散在本都街.
太加伯多加亞細亞彼得尼.
寄居的一帶的地方.
就是當時羅馬帝國東邊的地方.
當時彼得在羅馬寫信.
去給現今土耳其的一帶教會.
這一帶教會究竟是怎樣的教會呢.
其實當時這個地區大部分都是外邦人所住.
所以彼得現在要寫信給當地的教會.

$^{121}$是一班外邦人的教會.
換言之你們就是當時外邦教會.
在彼得前書的代名詞.
當彼得一說你們的時候.
就是說到那班外邦的教會.
第二節彼得前書二章九節.
很多時候都會成為我們個人在事奉上.
或傳福音上的一個鼓勵.
但是我請大家今天去思考這段經文的時候.
從教會的角度去看一看這段經文.
事實上彼得當時不是寫信給一個人.
彼得是寫信給整間教會.
聖靈向眾教會的說話.
凡有意聽的就應當聽.
唯有你們這個表達.
是凸顯了一件事.
如果在《新普及》的譯本是這樣說的.
但你們卻不是這樣.
表示到現在彼得所寫信給那班外邦教會.
和上文所描述的一班人.
是有明顯的不同.
那一班人究竟是一班什麼樣的人呢.
如果我們按照第七節和第八節所說的.
那一班是不信耶穌的人.
而彼得現在要寫信給一班信了耶穌的人.
彼得要寫信給一班教會.
教會的群體和這些不信的人明顯有分別.
究竟有什麼分別呢.
彼得前書二章九節.
是對教會的群體作出了四方面的形容.
這四方面的形容就凸顯了教會的獨特性.
信教主的群體和不信主的群體是有分別的.
這四個形容詞也凸顯了教會的本質.
彼得對教會作了哪四方面的形容呢.
我們來看看彼得前書二章九節.
彼得說:唯有你們是被揀選的族類.
是有君尊的祭司.
是聖潔的國度.
是屬神的子民.
彼得就用了這四方面去表達.

$^{161}$彼得說:教會就是被揀選的族類.
教會就是有君尊的祭司.
教會就是聖潔的國度.
教會就是屬神的子民.
時間關係我沒辦法逐個很詳細地.
跟大家去思考當中的意思.
但這四個經文的內容相當豐富.
我只能挑選其中一些重點跟大家去思考.
第一個就是被揀選的族類.
這個形容強調上帝在揀選的事情上.
採取一個主動還是被動.
採取一個主動.
教會是被上帝所選上的一個群體.
上帝揀選這個課題.
有兩方面值得我們去留心.
第一方面就是揀選是完全出於上帝的恩典.
既然上帝是做主動的人.
在揀選的事情上根本是毫無功勞可言.
唯獨是上帝的恩典.
上帝的恩典淋到你和我的中間.
我們才能夠進入天國裡面.
我們才能夠組成一間教會.
第二揀選是完全出於上帝的慈愛.
上帝揀選你和我.
不是基於我們的背景.
也不是基於我們有多高的學歷.
也不是基於我們的家財有多少.
上帝的揀選是出於祂的慈愛.
我想弟兄姊妹也有這樣的經歷.
就是你去超級市場或街市買菜或水果.
你也會去挑選.
你挑選的過程中會想什麼呢?.
這個炆雞豆皮.
這個蘋果看上去有點凹凸不平.
可能快要壞了就不要了.
你會挑選什麼?.
你要挑選最好的才會買.
但是上帝的挑選.
上帝的揀選不是用人的標準.
上帝的揀選是用祂的慈愛作為祂的標準.

$^{201}$去揀選你揀選我.
是出於上帝的愛.
上帝愛你.
彼得說教會除了是被揀選的族類.
教會也是君尊的祭司.
君尊祭司可以理解為什麼呢?.
就是屬於王尊貴的祭司群體.
王就是指導我們的大君王耶穌基督.
祭司是人和神之間的橋樑.
人藉著祭司和上帝建立關係.
人藉著祭司的獻祭.
去獻上他的讚美.
獻上一個認罪的禱告.
教會作為大君王耶穌基督尊貴的祭司群體.
教會就要將人帶到上帝的面前.
所以保羅在《歌羅西書》一章二十八節裡說.
我們傳人他,是用諸般的智慧.
勸誡各人,教導各人.
要把各人在基督裡完完全全地引到神的面前.
弟兄姊妹,教會是祭司的群體.
教會是要將人帶到上帝的面前.
或者我從另一個表達去說.
教會是要讓人看到基督.
弟兄姊妹,我們的教會.
有沒有讓人看到基督在我們中間?.
除了有君尊的祭司.
第三個描述就是聖潔的國度.
聖潔不是在說一些很高尚,很完美的道德標準.
聖潔也不是在說毫無指責.
聖潔是在說被分別出來的意思.
最近因為新冠肺炎的個案.
就飆升得很厲害.
很多專家說現在是第三輪的疫情.
因為疫情實在太嚴峻.
政府早前宣佈全民檢測.
遲些會再公佈一些詳細的內容.
檢測是屬於自願性的.
所以成效是如何呢?.
言之尚早.
我們看看到時有多少市民參與其中.

$^{241}$我們就知道成效是如何.
不過檢測的目的是什麼呢?.
檢測的目的是要將有感染的人.
從整群人當中抽出來.
然後將他隔離.
然後讓他接受治療.
這個就是分別出來的意思.
弟兄姊妹.
教會在萬民當中.
藉著耶穌基督的救贖被分別出來.
彼得最後的一個形容.
是屬神的子民.
屬神的子民.
呂鎮忠牧師翻譯聖經.
他是按照聖經的原文直譯出來的.
如果你留意呂鎮忠牧師的譯本.
這一句話.
呂牧師怎樣譯呢?.
他就說屬上帝的子民.
他怎樣譯呢?.
是子民做上帝產業的.
我覺得這個比起和諧本的翻譯.
更加能夠帶出彼得原來想帶出來的意思.
上帝是把人買熟過來.
成為他的產業.
上帝是百分之一百擁有教會這個群體.
上帝也是百分之一百在教會裡面掌權.
求主真的讓我們能夠明白.
聖經裡面所說.
我們屬神的子民.
既然屬神的子民的時候.
是他買熟回來的.
我們有沒有一份信心.
相信上帝是保守我們.
保守我們這個群體直到見他的面.
縱然現在在疫情的底下.
政局很動盪.
但我們必然建立他所建立的群體.
因為我們是他買熟回來的.
我們是尊貴的.

$^{281}$大概我已經把四個描述和大家分享過.
不過有一點很值得我們去留意.
大家有沒有覺得這四個描述好像似曾相識一樣.
其實彼得某程度上在引用聖經.
背誦舊約聖經.
因為這四個描述分別來自兩段舊約經文.
一段經文是以賽亞書43章20至21節.
一段經文是第一個形容被揀選的族類.
和第四個形容屬神的子民的出處.
第二段經文是出埃及記第十九章第六節.
你們要歸我作祭司的國度為聖潔的國民.
這是第二個形容君尊的祭司.
和第三個形容聖潔的國度的經文的出處.
出埃及記是在說什麼呢.
出埃及記是在說一群神的選民.
在埃及當中被法老奴役.
神在當中行了很大的神蹟.
將以色列民從埃及當中拯救出來.
帶他們去加拿大.
以賽亞書一段經文.
以賽亞書第四十章開始.
其實是在說有朝一日.
以色列人會被擄到巴比倫當中.
彼得很巧妙地透過這兩段經文.
彼得似乎在暗示新約的信仰群體.
同樣像舊約的神選民一樣.
是經歷過上帝的釋放.
其實教會是上帝從罪惡的權勢當中.
買熟回來的一個群體.
如果用彼得的話去說.
就是出黑暗 入奇妙光明.
昔日上帝在萬族當中.
去揀選以色列民.
以他們作為地上的代表.
今日上帝揀選你和我進入教會當中.
教會亦是上帝今時今日在地上的代表.
昔日上帝興起以色列民.
從萬族當中將他們分別出來.
要他們成為祭司的角度.
今日教會.

$^{321}$上帝興起教會.
是要我們將人帶到上帝面前.
成為耶穌基督祭司的群體.
弟兄姊妹.
教會是被揀選的族類.
是有君尊的祭司.
是聖潔的角度.
是屬神的子民.
這四個身份你有多認同?.
教會這四個身份.
有沒有令你感覺到尊貴?.
有沒有令你感覺到與別不同?.
彼得不單止將教會的本質.
很清楚地指出來.
他更加將上帝對教會的心意.
透過第九節的下半節去表達出來.
第九節的下半節說.
要叫你們宣揚那照你們出暗入奇妙光明者的美德.
宣揚這個字.
在原來的文字.
一個複合字.
由兩個字所組成.
一個字叫出離.
另一個字叫宣告.
所以中文聖經和合本.
將它翻譯為宣揚.
是非常好的翻譯.
亦很貼切.
我翻查過很多不同的譯本.
現代中文譯本.
新譯本.
呂鎮忠譯本.
新漢語譯本.
大部份的中文譯本都是宣揚.
不過.
讓我找到一個譯本.
它不是這樣的翻譯.
新普及譯本的翻譯是怎樣呢?.
新普及譯本的翻譯.
是讓我們更加明白宣揚的意思.

$^{361}$我再說一次.
頂主們.
新普及譯本的翻譯.
是讓我們對宣揚這個字.
有更加多的理解.
宣揚不是只讓人聽見.
宣揚亦是要讓人看見.
宣揚不是我們用口去表達.
宣揚是有行動去配合的.
究竟今時今日.
教會宣揚甚麼呢?.
彼得在那段經文裡.
教會是宣揚上帝的美德.
美德這個字原本是結出的成就.
通常是指道德方面的事情.
不過有些聖經學者指出.
美德這個字如果是用來形容上帝的話.
就不是在說道德的層面.
而是在說上帝的大能.
和上帝榮耀的大作為.
在彼得前書.
寫到二章九節的時候.
上帝榮耀的大作為是甚麼呢?.
上帝榮耀的大作為.
就是上帝將自己所揀選的教會.
從黑暗當中帶進光明.
頂子梅彼得指出.
教會的本質就是上帝的揀選.
是被分別出來的祭司群體.
也是聖祭的角度.
也是屬上帝的指紋.
上帝所揀選的教會.
要教會宣告上帝的大作為.
要讓上帝得榮耀.
這隻經文對今時今日.
我們教會群體有甚麼提醒呢?.
我們又如何回應.
這節經文對我們的教導呢?.
我給你看看一幅圖畫.
這就是新普及本的經文.

$^{401}$美國總統林肯曾經說過一句話.
You cannot escape the responsibility of tomorrow by erasing it today.
如果用中文翻譯.
就是你不能藉由今天的逃避.
去逃避明天的責任.
今天明天教會都有一個不能逃避的使命.
我再說一次.
教會今天和明天都有一個不能逃避的使命.
這個使命是甚麼呢?.
這個使命是.
教會要在黑暗的世代見證上帝的大作為.
也要讓人看見上帝的榮耀.
我們讀聖經.
我們很喜歡金句式地讀.
甚麼是金句式呢?.
就是讀到有一兩節的聖經.
很觸動到自己.
又或者對自己很有幫助.
又或者這兩節的聖經.
你感覺到很有意思.
於是我們就抽這兩節的聖經出來.
成為我們的金句.
我們可能會背誦了它.
或者我們會適時地運用.
不過這樣的讀經方法有好處.
也有不好處.
你將神的話背誦在心裡.
一定有好處.
作為你隨時的幫助.
但也有一些弊處.
因為你未必會很清楚.
究竟這段經文本身想表達的意思.
彼得前書二章九節就有類似的情況.
我們很多時候就用彼得前書二章九節.
作為我們個人事奉.
又或者個人傳福音上的提醒.
不過當我們將這段經文.
放回彼得前書整卷的聖經當中.
我們對這段經文會有更豐富的理解.
大概彼得前書分開四個部分.

$^{441}$第一個部分就是第一章第一節到第十二節.
這是一個問安感恩頌讚的部分.
然後第一章的第十三節到第二章的第十節.
在這裡彼得將教會的身份.
講得很清楚明白.
然後第二章的第十一節.
到第四章的第十一節.
這裡是說基督的群體在世上受苦的意義.
就是見證基督.
然後第四章十二節到第五章最後.
就是說基督的群體在世上有受苦的心志.
要行隱靠主.
原來彼得前書二章九節.
這一節的經文是彼得寫出來.
要預備當時活在小亞細亞外邦教會的信徒.
預備他們的心志.
預備他們受苦的心志.
預備他們去見證基督的心志.
其實彼得前書的修信人.
小亞細亞的教會.
他們生活在一個怎樣的境況裡面.
他們是生活在一個不信的世代.
他們是生活在一個不認識基督信仰的社會.
甚至是一個排斥敵視基督信仰的世界.
由於價值觀的不同.
基督群體在不信的世代生活.
會面對很多挑戰.
又或者會受到很多排斥.
如果我們細心去閱讀整卷的彼得前書.
你會發現當時的信徒.
他們正在面對一些艱難.
究竟是甚麼.
其實當中都很具體.
包括甚麼呢.
被人毀謗.
也會被人威嚇.
也會被人惡待.
難怪彼得在施信的開頭會這樣說.
如今在百般的試煉當中.
是暫時的優秀.

$^{481}$整卷的彼得前書.
彼得是要安慰信徒.
縱然你活在一個不信的群體當中.
縱然你活在一個敵視基督信仰的世界.
你仍然要堅忍.
恆忍靠主.
要活好美好的見證.
其實歷世歷代的教會.
都是這樣在黑暗的時代當中去掙扎.
弟兄姊妹.
我們同樣都活在一個彎曲薄茂的世代.
我們同樣都活在一個黑暗的城市.
在過去幾年.
我相信大家有很深的體會.
我們見到我們所居住的城市.
有很多的不同.
很多的轉變.
作為擁有公權力的政府.
它的施政往往都是令港大市民沮喪.
更加嚴重的就是政府的施政.
它未能盡它的責任去彰顯社會的公義.
在去年社會運動當中.
不同陣營都有不同的暴力的情況.
而這些暴力也是越來越嚴重.
在過去的十年.
全球很多的國家都推行一個量化寬鬆的政策.
在這個政策的底下.
資產價格的抬升.
而一些的資產就會慢慢落入一小撮人的手.
這個情況最終造成什麼呢?.
貧富懸殊是越來越嚴重.
根據統計處的報告.
香港的堅尼系數.
堅尼系數是貧富懸殊的一個指數.
當這個指數越來越接近1的時候.
表示貧富懸殊是非常之嚴重.
在2006年堅尼系數的數字是0.533.
香港十年過後2016年.
香港的堅尼系數是0.539.
是飆升得很厲害.

$^{521}$創了45年的新高.
香港大概有兩成的人口活在貧窮當中.
又或者被政府界定為貧窮的一個群體.
社會上貧富懸殊越來越嚴重.
背後在說什麼呢?.
背後在說人的貪婪.
如果人不是貪婪的時候.
貧富懸殊不會這麼嚴重.
弟兄姊妹.
我們是活在一個罪惡不斷膨脹的城市.
不知道你能否感受到.
活在一個罪惡滔天的年代.
作為上帝地上的代表.
作為福音的出口.
教會應該如何去面對呢?.
香港有三個基督教的宗派.
每一年的六月第三個星期的星期三.
都會聚在一起.
會有相交的時間.
會一起去探討一些牧養的問題.
這個都成為一個常規.
今年六月十七日.
就是三宗的聚會的那一天.
今年的三宗聚會的主題很有意思.
很有時代感.
主題是什麼呢?.
可能弟兄姊妹會知道.
主題就是「復興教會,時代照命」.
是不是很有時代感?.
讀起來都覺得很厲害.
其實教會最大和最根本的使命是什麼呢?.
教會最大和最根本的使命.
就是要見證上帝的大作為.
什麼叫上帝的大作為呢?.
上帝的大作為可以有很多討論.
你可以說是上帝的創造.
又或者是上帝的救贖.
但是如果我們放在彼得前書的內容裡.
上帝的大作為是在說耶穌基督的福音.
教會要見證上帝的大作為.

$^{561}$教會要傳揚耶穌基督的福音.
問一下弟兄姊妹一個問題.
如果今天你要出去傳福音.
你會選擇哪一節的經文.
去和你的福音對象分享呢?.
想到沒有?.
我猜的.
我純粹猜的.
我沒有實際的答案.
因為也見不到大家.
我猜大部分的弟兄姊妹.
都會用《有限福音》三章十六節去說.
那裡說得很清楚.
神愛世人甚至將他的獨生子賜給他們.
叫一切信他的不自滅亡反得永生.
我記得我第一次跟弟兄姊妹報道傳福音的時候.
我當時初信一個多月.
我就跟著她出去傳福音.
弟兄就引用這節經文.
去和一個坐在公園的叔叔傳福音.
那時候我就說很厲害.
一出口就知道這節經文在哪裡.
那時候我覺得很驚訝.
後來發覺原來很多人都會用這節經文去傳福音.
這節經文也是好的.
不過我有另一個選擇給大家.
大家可以去參詳一下.
去想一想.
這節經文是《馬太方》四章第十七節.
《馬太方》四章第十七節說的是什麼呢?.
是耶穌的說話.
他說「天國近了,你們應當悔改」.
天國是什麼意思?.
天國是上帝的權能.
藉著道成肉身的耶穌基督.
在地上去彰顯出來.
天國是在說上帝的權能.
當我們去傳福音的時候.
我們有沒有說到上帝的權能.
透過耶穌基督道成肉身的行動去彰顯出來呢?.

$^{601}$弟兄姊妹.
福音不是說我們個人和上帝的關係.
福音也不是說我們得救了.
我們上天堂.
我們靈魂得救了.
福音不只是說這件事.
福音是說上帝的權能.
藉著耶穌基督道成肉身.
來到這個世界當中去彰顯出來.
弟兄姊妹.
如果我們去傳福音.
我們沒有將這個內容.
和我們福音對象分享的時候.
我盼望在今天過後.
當我們再去傳福音的時候.
我們將福音一個完整的意思.
去和對方分享.
剛才我說教會是不能逃避一件事.
就是黑暗的世代.
要去見證上帝.
讓上帝得榮耀.
甚麼是上帝得榮耀呢?.
上帝得榮耀也有很多討論.
可能有些人說.
上帝得榮耀就是.
上帝屬性的彰顯.
但是如果我們放在彼得前書.
整卷書的內容裡面.
讓上帝得榮耀是指.
基督的群體在末世寄居的日子.
無論環境如何.
仍然持守信仰的核心.
和持守信仰的純正.
這就是讓上帝得榮耀.
在彼得前書的意思.
因為種種的原因.
內地在2015年開始.
逐漸收緊宗教的政策.
所以你會見到.
有很多拆十字架的情況.

$^{641}$你也會見到.
有很多傳道人牧師被政府拘留.
黃賢牧師也被政府拘留.
並且經過審訊被判監.
他現在在監牢當中.
在今年年初23日.
國務院對主理香港事務的官員.
也作出了一些調動.
被派去主理港澳事務辦公室的主任.
就是曾經在浙江省.
推動拆十字架的那一位官員.
這個調動會不會帶來.
香港宗教事務的一些轉變呢?.
答案只有兩個.
一是會.
一是不會.
當然我們不希望會有任何的轉變.
但是萬一真的有轉變的時候.
弟兄姊妹.
你預備好了嗎?.
你預備好為基督去見證.
去見證基督.
你預備好受苦的心智了嗎?.
我覺得彼得前書整卷的聖經.
很值得我們弟兄姊妹去誓約.
在當中能夠讓我們去學習得到.
在一個艱難的時代.
在一個黑暗的時代.
怎樣去為主去作辯證.
怎樣去讓上帝得榮耀.
這是彼得去寫彼得前書的心意.
最後我想和大家去唱一首詩歌.
這幅圖畫最後送給大家.
這幅漫畫是說上帝將我們從黑暗當中拯救出來.
上帝差遣你和我.
差遣教會.
進到黑暗的世代當中.
去為他做見證.
這首我想和大家唱一首詩歌.
是《興起發光》.

$^{681}$神的恩典看顧我.
罪裡釋放新的我.
從頭活過不計過錯.
完全的愛超過深海.
從今天起我決意成就你之意.
而我卻要把主的愛來囂觀這世界.
你偶然生活黑暗裡.
亦會有愁和掛慮.
因光定別光照你對你沒有捨棄.
神我要發光於這地上.
要將福音來傳頌.
神我要發光於這地上.
要將救恩宣告.
神我要發光於這地上.
要將真光來延亮.
神我要發光於這地上.
還贊助奇妙救恩.
求你興起一生向前闖.
主恩為無窮.
我請詩琴繼續彈.
《興起發光》.
是我們上一個月的善明詩.
當我們去唱頌這些善明詩的時候.
我們內心裡有沒有被.
上帝的救贖.
或者詩歌裡的訊息.
勒住了.
詩歌說.
我要發光於這地上.
要將福音傳頌.
要發光於這地上.
要將救恩宣告.
丙姐妹.
今天如果有一個.
未信的朋友去探訪你.
又或者你要去見一個.
未信主的親友.
你能不能夠讓他.
在你身上見到基督.
你能不能夠將他.

$^{721}$將一個完整的福音.
去和他分享.
求主幫助我們.
《興起發光》.
我們再唱這首詩歌一次.
神的恩典看顧我.
罪裡釋放新的我.
從頭活過不計過錯.
完全的愛超過深海.
從今天起.
我決意成就你之意.
而我更要化處力外.
來囂觀這世界.
你偶然生活黑暗裡.
亦會憂愁和掛慮.
因光定別光照你.
對你沒有捨棄.
神我要發光於這地上.
要將福音來傳頌.
神我要發光於這地上.
要將舊恩宣講.
神我要發光於這地上.
要將真光來獻亮.
神我要發光於這地上.
眾讚主奇妙救恩.
求你興起一身向前創.
主因為無窮.
以上言論不代表本台立場.
\newpage



\section{哈該書 2:1-9}
\label{sec:xKlgziHaUYU}
\textbf{2020.08.23《困境中的勇氣》 哈該書 2\_1-9 陳麗蟬牧師}
\newline
\newline
連結: \href{https://youtube.com/watch?v=xKlgziHaUYU}{\texttt{https://youtube.com/watch?v=xKlgziHaUYU}} ~~~~ 語音日期: 2020-09-05
\newline
\newline
\hyperref[sec:UM_spSYKywQ]{\small{< < < PREV SERMON < < <}}
~
\hyperref[sec:index_chronic]{\small{[返順時目]}}
~
\hyperref[sec:index_scriptual]{\small{[返順卷目]}}
~
\hyperref[sec:YQ9vapt0XGE]{\small{> > > NEXT SERMON > > >}}
\newline
\newline
哈該書 2:1-9
\newline
\begin{longtable}{cl}
\hline
\hline
章節 & 經文 (和合本修訂版)\\
\hline
2:1 & \begin{tabularx}{0.7\textwidth}{X} 七月二十一日，耶和華的話藉哈該先知傳講，說： \end{tabularx} \\ \\ \relax
2:2 & \begin{tabularx}{0.7\textwidth}{X} 「你要曉諭撒拉鐵的兒子猶大省長所羅巴伯、約撒答的兒子約書亞大祭司，和所有倖存的百姓，說： \end{tabularx} \\ \\ \relax
2:3 & \begin{tabularx}{0.7\textwidth}{X} 『你們中間存留的，有誰見過這殿從前的榮耀呢？現在你們看如何？在你們眼中豈不是如同無有嗎？ \end{tabularx} \\ \\ \relax
2:4 & \begin{tabularx}{0.7\textwidth}{X} 所羅巴伯啊，現在，你當剛強！這是耶和華說的。約撒答的兒子約書亞大祭司啊，你當剛強！這是耶和華說的。這地的百姓啊，你們都當剛強做工，因為我與你們同在。這是萬軍之耶和華說的。 \end{tabularx} \\ \\ \relax
2:5 & \begin{tabularx}{0.7\textwidth}{X} 這是照著你們出埃及時我與你們立約的話。我的靈仍要住在你們中間，你們不必懼怕。 \end{tabularx} \\ \\ \relax
2:6 & \begin{tabularx}{0.7\textwidth}{X} 萬軍之耶和華如此說：過些時候，我必再一次震動天地、滄海與乾地。 \end{tabularx} \\ \\ \relax
2:7 & \begin{tabularx}{0.7\textwidth}{X} 我必震動萬國，萬國的珍寶都必運來，我就使這殿充滿榮耀。這是萬軍之耶和華說的。 \end{tabularx} \\ \\ \relax
2:8 & \begin{tabularx}{0.7\textwidth}{X} 銀子是我的，金子也是我的。這是萬軍之耶和華說的。 \end{tabularx} \\ \\ \relax
2:9 & \begin{tabularx}{0.7\textwidth}{X} 這後來的殿的榮耀必大過先前的榮耀。這是萬軍之耶和華說的。在這地方我必賜平安。這是萬軍之耶和華說的。』」 \end{tabularx} \\ \\
[1ex]
\hline
\hline
\end{longtable}
$^{1}$各位姐妹平安.
今天很感恩可以來到洪恩家.
跟大家一起用上帝的話.
彼此勉勵.
在十九世紀的時候.
英國有一個很出名的政治家.
他叫格蘭斯敦.
他在國會裡面做了四任財政大臣.
也做了十二年四任首相.
到他年老的時候.
有人問他.
這麼成功.
可以留在國會裡面做這麼久.
秘訣在哪裡呢?.
他就用一個故事來回答訪問他的人.
他說有一條路從郊外一直通到倫敦.
很特別.
那些馬在這裡主要行這條路的拉車馬.
很容易死的.
不知道為什麼.
後來有人就走了.
去探究一下原因.
就發覺到這條路又平坦又通暢.
沒有什麼阻滯.
所以那些馬經不起任何一些崎嶇的路.
所以當他行慣了這條這麼平的路.
一旦有時要行一些比較崎嶇的路的時候.
他們就很容易死.
因為太過平坦.
以致到那些筋肉或者氣力.
根本就沒有機會來走去鍛鍊.
那位老宰首相就繼續說.
他說人生的路上也有很多艱難的事情.
但這些唯有遇到艱難的事情的時候.
才能夠叫人繼續奮勇向前.
能夠繼續堅持到底.
也都唯有去遇到這些逆境的時候.
我們才可以將上帝給我們每一個人的潛力.
可以盡量發揮出來.
原來這個老宰首相成功的秘訣.

$^{41}$就是很積極地去面對每一個的逆境.
真的,在人生的路上.
或者在事奉的路上.
都會遇到很多困難.
但如果一遇到困難的時候.
我們就退縮,我們就停下來.
什麼都不做的時候.
我們很容易就變成像一個死的活人一樣.
永遠都沒有進步的.
剛才讀的那段經文.
就是說以色列人回歸耶路撒冷去重建聖殿.
但當我們建了不久之後就停了下來.
所以耶和華就呼籲黑丐先知去鼓勵他.
我們就先看一下當中整件事件的背景.
在公元前606年的時候.
以色列人被擄到巴比鄰.
但是上帝又應許他們.
70年之後他們會回歸.
到了公元前536年的時候.
耶和華就搞動了波斯王的心.
他就准許了42,000的以色列人.
回歸耶路撒冷重建家園.
重建聖殿.
他們最初都很熱心.
回歸第二年的時候已經開始動工.
不過離開了這裡幾十年.
住在這裡慢慢的霸佔了這個地方的人.
就覺得他們阻礙了.
回來搶了他們的地方.
所以就想了很多辦法去阻止他們.
恐嚇他們.
就是這樣被人罵一罵.
嚇一嚇他們就停工了.
這次一停工就停了20年.
他們就忘記了為什麼他們要回來耶路撒冷.
他們只想著自己建好自己的房子.
希望住得好一點.
但是他們最終要回來重建聖殿的地方.
就擲物重新了.
仍然被那些野獸來居住.

$^{81}$他們是將敬拜上帝的地方置諸不理.
其實當他們這樣來看上帝的殿的時候.
他們就是不將上帝放在眼裡.
所以耶和華就興起下蓋.
去勸這班子民.
去想一下在過往這十九二十年裡面.
他們自己想著自己建房子.
他們都很努力耕種做事.
但是總是好像做什麼都沒有什麼成功.
又不是很能夠賺到錢.
就是因為他們沒有將上帝放在首位.
就好像耶穌基督所說.
沒有去先求上帝的國和上帝的義.
當我們以上帝為首的時候.
所有我們所需要的東西.
上帝都會給過我們的.
我們就看一下第一方面.
當時他們遇到什麼困難.
經文開始說.
當天是七月二十一日.
當然是猶太人他們是七月二十一日.
按照利未記來說.
七月是以色列人的新年.
二十一日應該是鑄棚節最後一天.
這個鑄棚節是以色列人慶祝收成.
當然當天他們沒有什麼收成.
我還沒有下這幅powerpoint.
可以上去嗎?.
好,謝謝.
當天他們都沒有豐收.
但他們仍然會慶祝.
除了慶祝農作物的收成之外.
他們也會感恩他們的祖先.
從埃及地被耶和華拯救了出來.
他們在這個曠野裡四十年鑄像棚的日子.
所以每逢這些節日.
他們都會有很多人集中在一起.
在這個時間.
耶和華呼籲哈該到他們的中間.
鼓勵省長索羅巴伯.

$^{121}$鼓勵大祭司約書.
他跟他們說.
你們已經停頓了二十年.
我之前跟你們說過.
你們又要開工了.
但開工之後只有一個月.
好像又有些事情停頓了.
還沒完全停頓.
他們要繼續的剛強.
繼續的來作工.
為什麼呢?.
因為當中有些人.
特別是一些老人家.
他之前見過所羅門建的聖殿.
我們剛才看到.
所羅門的聖殿是金碧輝煌的.
當時所羅門聖殿.
全部的東西都是.
他的父親大衛幫他預備材料.
所以所有的東西都是金做的.
而那些木材.
最好是從黎巴嫩運來的.
還有各個很出色的工匠.
都來到這裡幫忙.
整整建了這個聖殿七年.
但到現在.
所羅門的父親帶領他們回來建這個聖殿.
只是用一些木板造成.
你看金燈台也是用一整塊金來造成.
然後呢.
我去不到下一幅.
對了.
但現在只是一個爛地一片.
而且也估計了.
將來他們建的聖殿.
也不再像昔日所羅門建的聖殿那麼輝煌.
這裡只是一些木板.
還有將波斯王給他們的一些財寶.
和退還一些聖殿的器皿.
給他們回去.

$^{161}$所以他們就預計到.
他們的資源有限.
做事的人都不及以前所羅門的人那麼多.
和那麼有工藝.
再加上當時一些外族.
隨時會來煩擾他們.
說你來霸佔了我們的土地.
所以有些人可能會覺得.
建來做甚麼呢?.
以前所羅門的殿那麼漂亮那麼輝煌.
現在沒有那麼多金體.
又沒有那麼宏偉.
究竟建來做甚麼呢?.
那麼少的殿.
建來做甚麼呢?.
都是空空的.
可能他們說了一些負面的話.
所以令到那班人又想想.
做的東西都沒有那麼漂亮.
第二次做的東西會更好.
為甚麼現在會做得更差呢?.
所以他們就有些沮喪.
我記得在我之前教會事奉的時候.
有一次在這個報道聚會裡面.
我寫了一個短劇.
大約五分鐘左右.
作為我講道之前的一個引入.
我知道有些姐妹都很擅長演戲.
請她來幫我去做.
誰知道她看過劇本之後.
她看過了這個劇本之後.
她就給了我一些意見.
她說劇本又不夠深入.
令到我們的演員都不能夠發揮演技.
我以前在國方話劇團那裡.
劇本多深入.
令到很感人.
不單看的人都很感動.
做的人都很感動.
她一大堆東西潑了冷水來之後.

$^{201}$我覺得我真的沒有心再做下去了.
但我心裡想不是的.
我正式做了一個小時的國方話劇.
我只是五分鐘的引子.
想清楚的時候.
不可以放棄.
有時一些負面的說話.
或者負面的情緒.
很容易蠶食人的心靈.
很容易令到人感到失望.
感到灰心.
很想放棄.
好像在《箴言》裡面這樣說.
她說.
《箴言》二十四章裡面這樣說.
她說你在患難之日.
若膽怯.
你的力量就微小.
她說你在患難之日.
若膽怯.
你的力量就微小.
當我們明明要去做一些事的時候.
若果因為人的一些說話.
或者因為一些情景的時候.
我們就會退縮了.
所以上帝就呼籲哈該.
走去為這班人加力.
所以在人生的路上.
或者在事奉的路程裡面.
有時我們會灰心喪志.
有時我們會覺得沒有力量.
沒有衝勁.
甚至覺得那些事做來也沒有用.
做來做去.
說來說去.
人們也不容易相信耶穌.
做來做去幹什麼.
但是我們作為上帝的子民的時候.
我們一定要對上帝有信心.
他一定會帶領我們去走每一步的路.

$^{241}$接著我們看看第二方面.
第二方面我們又看看這個.
我想我不太懂得控制這個part point.
第二方面.
不如下面幫我控制行不行.
就是講勇氣的由來.
是這個.
勇氣.
勇氣那裡.
是啦,謝謝.
所以當先知出來做這個宣告的時候.
他怎麼說呢.
他說猶大的省長索羅巴伯.
整個帶領索羅巴伯.
他說現在你當剛強.
這是耶和華說的.
約撒達的兒子約書亞大祭司.
你當剛強.
這是耶和華說的.
這個是宗教界的領袖.
他說這裡的百姓.
你們都當剛強做工.
因為我與你們同在.
這是萬軍的耶和華說的.
這裡講到百姓的層面.
就是說整個一起回歸回去的.
無論你是領袖還是平民百姓都好.
他們都要剛強壯膽地做工.
不要再被周圍的人的一些說話.
或者一些情景.
來嚇著他們.
他們要有一個堅定的心智.
有一個強壯的身軀.
不可以灰心氣累.
因為為什麼.
上帝與他們同在.
這個上帝究竟是怎麼樣呢.
在單單第四節裡面.
已經講了三次.
耶和華與你同在.

$^{281}$耶和華這個神又是怎麼樣呢.
當他出埃及的時候.
在第三章就講到.
耶和華呼籲摩西帶領那班以色列人出埃及.
當時摩西問耶和華.
當時他還不知道他叫耶和華.
當人問我憑著什麼.
來帶他們出埃及的時候.
我怎麼說呢.
耶和華就跟他說.
你就告訴他們.
帶領他們的神是那位自有永有的神.
也就是說我們的上帝是色在,金在,以後永在的上帝.
在這麼長的時間.
他就將這班以色列人.
從埃及本來是做奴隸的地方拯救了出來.
並且去允許這班以色列人.
永遠保護他們.
還有他形容耶和華是什麼呢.
是萬軍的耶和華.
其實當剛才主席為我們讀出的時候.
我們看到這是萬軍的耶和華說的.
出現了很多次.
出現了六次.
這是萬軍的耶和華說的.
不單單代表我們的上帝有一萬隊軍隊.
如果敵人有一萬零一隊軍隊我們都輸.
他原本的文字意思就是我們的上帝.
是那位全能的上帝.
他這次重建這個聖殿.
可能沒錯.
是真的沒有昔日所羅門建這個聖殿的時候.
有很多的金銀珠寶.
有很多的工匠.
有很多的漂亮的木材.
也都預知了.
建出來的聖殿一定沒有過往那麼輝煌.
但不要緊.
在這次整個工程裡面.
萬軍的耶和華.

$^{321}$全能的上帝與他們同工.
這個全能的上帝與我們同工.
究竟感覺是怎樣的呢?.
每天上班.
我都要走到金鐘站去搭地鐵.
要走進金鐘站之前.
有一道很大的玻璃門.
很重的.
所以每一次我進去的時候.
我都要用九牛二虎之力大力地推進去.
有一天.
我走來走去的時候.
又是把馬準備好.
兩隻手一推的時候.
那天有些奇妙的事情發生.
當我想用力.
剛剛碰到那個手柄的時候.
我發覺不用再用力.
那道門自動推開了.
真是奇妙.
當我向上看的時候.
我看到有一個比我高很多的男生.
他在上面一推.
推開了那道玻璃門.
就進去了.
讓我明白到.
體會到甚麼叫做萬軍的耶和華同在.
與我們同工的時候.
那份的.
有一個字是爽.
爽的感覺.
就是自己都要用力.
站在那裡都要去推.
但一發現原來.
上高有人來加力的時候.
所以我們好像.
實質又不是用了很多力的時候.
這件事就成功了.
上帝的同在.
上帝的同工.

$^{361}$萬軍的耶和華.
全能的上帝.
與我們同在.
與我們同工.
就是我們一定要.
一定要自己親自作工.
然後上帝又在自己的上面.
去工作.
所以順著聖靈的風去工作的時候.
就很容易.
我們就去處理到.
好,然後就.
下一幅的powerpoint.
榮耀的角度.
剛才那幅,對不起.
榮耀的角度那裡.
第二章第九節.
「這後來的電的榮耀.
必大過先前的榮耀.
這是萬軍之耶和華說的.
在這地方我必似平安.
這是萬軍之耶和華說的」.
聖殿可以說是.
以色列人一個很感到驕傲的建築物.
也是以色列人信仰的中心.
所以黑蓋先知.
走來走去跟一群國姓說.
不是要求他們.
建一個像鎖郎門的聖殿那麼漂亮.
因為知道了.
無論在財力,物力,人力.
都不能夠跟鎖郎門時代相提並論.
祂只是在鼓勵這些以色列人.
不要因為現在的能力有限.
就自卑起來.
因為上帝不是看人.
或者建築物的外表.
是看人的內心.
去看人是否真的以上帝為首.
真正敬畏上帝的人.

$^{401}$一定會時時刻刻都經歷上帝同在的力量.
上帝在這裡告訴你.
將來的殿一定比先前的殿更加榮耀.
上帝就將千禧年這個聖殿.
和當日黑蓋先知.
他們那個時代所建的聖殿.
連在一起來說.
所以當時那群子民都不知道.
不過就已經給了他們一個盼望.
就是將來.
所以他們不要灰心.
將來必有一個比所羅門的殿更加榮耀的聖殿出現.
所以他們現在要做的.
就是專心地去工作.
因為上帝應許在將來.
就是千禧年的時候.
必定會有一個更加大的聖殿出現.
我們又看看啟示錄.
啟示錄怎樣去描述將來的殿是怎樣的.
啟示錄第21章.
「我沒有看見城內有殿.
因主,全能者神和羔羊就是城的殿.
那城內不辱日月光照.
因為有神的榮耀光照.
將來的聖殿就是上帝自己.
在將來的聖殿裡面.
光明是不需要有燈.
也不需要靠太陽光來照.
因為羔羊就是耶穌基督.
就是這個城.
就是這個聖殿的燈.
他們要藉著城的光行走.
地上的君王要把自己的榮耀帶給那城.
主神,全能者和羔羊為城的殿.
那個殿的輝煌.
因為祂自己就是上帝自己.
耶穌基督就是終極的聖殿.
所以在祂的裡面就有這個榮美.
有這個風色.
所以黑葵仙姐在說這句話的時候.

$^{441}$其實她也不知道自己在說什麼.
但原來上帝就在指示.
在將來終極的那個殿.
就是耶穌基督自己.
萬國都要來去敬拜祂.
這個就是我們堅持要去等待的.
縱然地上有困難,有逆境.
但我們要去堅心地去等待.
去有這個盼望.
今天就很想介紹一首詩歌.
叫做《唱一首天上的歌》.
唱一首天上的歌.
麻煩你最後一句話.
現在懂得唱的可以跟我一起唱.
下一個部分麻煩你.
聖母的歌 是一朵的花.
彎彎有情我的聖母.
聖母的歌 是一朵的花.
彎彎有情我的聖母.
我要唱啊 耶穌啊.
唱一首天上的歌.
多想你呼籲.
心裡的憂傷 全都散.
這首歌是一位叫做小敏的姊妹所寫的.
這個小敏是河南省的人.
不過她沒有怎麼讀過書.
她也從來沒有學過音樂.
上帝聖靈給她這方面的恩賜.
她寫了一千六百多首的短詩歌.
就收納在一本迦南詩選裡面.
每一首歌都非常感人.
這首是她其中的一首歌.
有一套電影在2013年的時候.
有一套電影叫做《1942》.
不知道有沒有人看過.
是馮小剛導演拍攝的.
他應該有獲獎的.
《1942》這個故事是講在河南省.
在1942年的時候遇到大饑荒.
那裡也接近大戰.

$^{481}$所以整個河南省的人民非常痛苦.
整個故事就寫這個背景.
當他拍攝完的時候.
馮小剛導演為了更加真切了解河南省的情況.
他就走到河南省.
試一下找有沒有老人家仍然在那裡.
就去講一下昔日的那些警方.
應該是七十幾八十幾歲的老人家.
真的被他找到兩個.
他就想著其實他經歷了那麼多.
經歷了饑荒,戰爭,內戰,共產黨.
一定會有很多事情.
但他見到這兩位老人家是很喜樂,很開心的.
他就問一下.
為什麼你大半輩子的人生都那麼困難.
為什麼你仍然那麼開心呢?.
那兩個老人家就告訴他.
我相信耶穌嘛.
導演就問信耶穌又怎麼樣?.
老人家就說信耶穌將來可以上天堂.
上天堂又怎麼樣?.
再餓過,再走難過,再經歷戰爭.
那兩位老人家就說.
在天堂那裡喝杯水你都不會餓死.
然後那兩位老人家就唱了小敏姐妹這一首的詩歌.
所以馮小剛就覺得很感人.
不過他覺得一首這麼輕鬆的詩歌.
對比整個故事就比較淒慘.
好像不太配合.
所以就找人將它修改了.
他就用一個暗淡一點的來寫.
我們試一下這次怎麼唱.
我會先唱我要唱哪一首歌.
就是尾四那四行.
用一個慢版.
然後下去唱下面那一版.
先唱這一版.
然後唱下面那一版.
準備.
Go.

$^{521}$我要唱哪一首歌.
唱一首天生的歌.
土上的烏雲 心底的憂傷.
穿越山谷.
唱一首.
一二三四.
生命啊 死了的歌.
萬物的寂寞的寂寞.
生命啊 死了的歌.
萬物的寂寞的寂寞.
我要唱哪一首歌.
唱一首天生的歌.
土上的烏雲 心底的憂傷.
穿越山谷.
生命的河 喜樂的河.
真是在聖經的描述裡面.
在將來的國度裡面.
真是有這個生命的河.
不過我們都知道.
這個聖靈已經在我們信徒的心靈裡面.
給我們有這個生命的活水.
以致這個生命的活水.
不斷在我們生命裡面湧流的時候.
我們就會有力量去堅持到底.
我們不懼怕任何的艱難.
我們有勇氣去面對.
我們再唱一次.
就這樣唱這裡.
好嗎?.
記住這句話.
《生命的河》.
生命的河 喜樂的河.
我們要將它散.
生命的河 喜樂的河.
我們要將它散.
我要唱哪一首歌.
唱一首天生的歌.
我們同心協力.
天父鎮教我們謝謝你.
因為你知道有萬軍的耶和華.

$^{561}$全能的上帝.
在我們生命當中的時候.
我們就有這個勇氣.
繼續去面對各樣的事情.
我們謝謝聖靈.
你真的在我們生命當中.
將這個活水江河.
在我們心中不斷的湧流.
以致在我們頭上的烏雲.
我們心中的憂傷.
都可以因著這個生命的活水.
湧流的時候.
我們可以去消除這一切.
我們謝謝你.
我們再一次謝謝我們的主聖靈.
求主你幫助我們.
以致我們時刻的去醒覺到.
在我們生命當中.
有聖靈成為我們自保的時候.
我們就無懼任何的環境.
祝福紅印堂的弟兄姊妹.
在現在這個疫情的底下.
主多謝你.
讓他們有機會將整個教會翻新.
他們在不同的流程裡面.
有不同的主題.
有不同的顏色.
讓在過程裡面.
主可能都會有些難處.
主就讓一班的同工.
一班的弟兄姊妹.
同心合意的去衝破.
以致到這一座是一座榮耀的殿.
榮耀你的殿.
是一個去為主你的福音.
去發出這個訊息的殿.
求主的工作.
多謝你瑞聽禱告.
是奉靠主耶穌基督.
得勝明治而求.

$^{601}$好累啊.
\newpage



\section{約伯記 1:1-5}
\label{sec:YQ9vapt0XGE}
\textbf{2020.08.30《闔府蒙福》 約伯記 1\_1-5 陳一華牧師}
\newline
\newline
連結: \href{https://youtube.com/watch?v=YQ9vapt0XGE}{\texttt{https://youtube.com/watch?v=YQ9vapt0XGE}} ~~~~ 語音日期: 2020-09-06
\newline
\newline
\hyperref[sec:xKlgziHaUYU]{\small{< < < PREV SERMON < < <}}
~
\hyperref[sec:index_chronic]{\small{[返順時目]}}
~
\hyperref[sec:index_scriptual]{\small{[返順卷目]}}
~
\hyperref[sec:OSwpE_DOZH0]{\small{> > > NEXT SERMON > > >}}
\newline
\newline
約伯記 1:1-5
\newline
\begin{longtable}{cl}
\hline
\hline
章節 & 經文 (和合本修訂版)\\
\hline
1:1 & \begin{tabularx}{0.7\textwidth}{X} 烏斯地有一個人名叫約伯。這人完全、正直、敬畏神、遠離惡事。 \end{tabularx} \\ \\ \relax
1:2 & \begin{tabularx}{0.7\textwidth}{X} 他生了七個兒子，三個女兒。 \end{tabularx} \\ \\ \relax
1:3 & \begin{tabularx}{0.7\textwidth}{X} 他的家產有七千隻羊，三千匹駱駝，五百對牛，五百匹母驢，並有許多僕婢。這人在東方人中為至大。 \end{tabularx} \\ \\ \relax
1:4 & \begin{tabularx}{0.7\textwidth}{X} 他的兒子按著日子各在自己家裡擺設宴席，派人去請他們的三個姊妹來，與他們一同吃喝。 \end{tabularx} \\ \\ \relax
1:5 & \begin{tabularx}{0.7\textwidth}{X} 宴席的日子過了，約伯派人去叫他們自潔。他清早起來，按著他們眾人的數目獻燔祭，因為他說：「恐怕我的兒子犯了罪，心中背棄神。」約伯常常這樣行。 \end{tabularx} \\ \\
[1ex]
\hline
\hline
\end{longtable}
$^{1}$大兄姐妹早晨.
很高興今天和大家一起敬拜.
雖然我們還是要靠網上來敬拜.
不過我們在網上.
我們同樣可以來專心.
來親近上帝.
今天很高興在崇拜當中.
和大家唱詩歌.
又祈禱.
還有一些詩歌也很激勵我們.
各位弟兄姊妹們.
為了今天早上敬拜上帝的機會.
獻上我們的感恩.
今天和大家思想的題目.
叫合輔蒙福.
喜歡這個題目嗎?.
我相信這個題目人見人愛.
平時我們說合輔統請.
我現在也是合輔統請.
統請什麼?.
統請全家蒙福.
如何能夠合輔蒙福呢?.
我們今天就選擇《約伯記》.
《約伯記》第一章一到五節.
在這裡我們可以看到.
一說到約伯.
我想大家都很熟悉這個名字.
但《約伯記》我不知道你有沒有看過.
可能你覺得舊約的聖經很深奧.
你就可能沒有打開.
舊約的經文去看.
但是你對約伯也很有印象.
原因是什麼呢?.
原因是你在自己生活不順利的時候.
或者遇到困難的時候.
遇到逆境的時候.
你就很快想起約伯.
為什麼呢?.
因為有同路人.
他曾經受過患難.

$^{41}$受過苦難比我們更加大.
所以你就會很快想起約伯.
希望有共鳴.
有同路人去支持你.
結果可能《約伯記》對你來說.
是半生熟也不一定.
人物你會熟悉.
但經文可能你會陌生.
不過無論如何也好.
今天早上我們看這段聖經的時候.
各位弟兄姊妹.
你要對約伯的印象有少許改變.
因為很多時候我們一想起約伯.
就覺得他很淒涼.
覺得他怎樣.
頭髮蓬鬆.
山瘡從頭長到腳.
一個人獨自坐在路邊.
神情落寞.
這個就是我們對約伯的印象.
如果你有這樣的印象.
今天早上我要幫你洗脫這個形象.
要少許的革新.
為什麼呢?.
因為你剛才的印象.
都是約伯遭遇患難之後的景象.
但在約伯未遭受患難之前.
他的情況卻不是這樣的.
他的情景是非常光彩的.
令人羨慕的.
何以見得呢?.
我們試試逐點來看.
首先我們看看約伯是生於哪個地方.
聖經說約伯生於烏斯地.
烏斯地是什麼地方呢?.
烏斯地是約旦河東的地方.
可以跳到約旦河東的地方.
有幅地圖.
約旦河東和約旦河西是很不同的.
約旦河西是以色列人的應許地.

$^{81}$約旦河東是阿拉伯人和以東人居住的地方.
烏斯地剛好在阿拉伯地和以東地之間.
叫烏斯地.
約伯是生於約旦河東的人.
即是說約伯不是希伯來人.
即不是猶太人.
在猶太人生活當中.
他是一個愛邦人也不一定.
不過很感恩.
他透過什麼渠道認識天父.
敬拜耶和華這位真神.
所以約伯當時在那個地方生活.
他的背景是很不錯的.
首先我們發覺.
原來約伯這個人個人名聲是非常好的.
為什麼?.
有十二個字去形容他.
不得了.
頂層門邊十二個字.
這十二個字就是約伯完全正直.
敬畏上帝.
遠離惡事.
這十二個字是不是很棒呢?.
如果有人用這十二個字去形容你.
不要說十二個字.
八個字都很感恩.
原來約伯這麼好.
完全正直.
敬畏上帝.
遠離惡事.
所以你見到他個人的聲望.
個人的名聲.
是非常榮耀上帝的.
這個人一說起他的名字.
就豎起手指.
這個人是頂呱呱.
不單是一個好人.
而是一個敬畏上帝的人.
這是第一.
他的名聲.

$^{121}$第二.
聖經記載他的財富.
如果你發覺.
剛才主席讀經文.
你留心數一數.
你會發覺不得了.
他的牛羊多到多少?.
成千這麼厲害.
駱駝也不少.
牛也五百對.
你有沒有留意到?.
各位頂智妹妹.
如果你逐一去買.
逐一去計算.
你會發覺他的家財滿貫.
你會發覺這位約伯.
他的財富真是富甲一方.
所以我們說他是烏斯地的首富.
是當之無愧.
大家認同嗎?.
你想想這麼多的牛羊駱駝.
在當時舊約的年代.
遠古的時候.
能夠在畜牧的日子當中.
擁有這麼多的牛羊.
實在是甚麼?.
實在是顯示出他富有的地方.
這就是他個人的財富.
不單是這樣.
聖經還描述.
他是一個名門望族.
為甚麼呢?.
因為這個家族裡.
竟然有多少位兒子.
約伯有七個兒子.
有三個女兒.
加起來有十個這麼多.
你說厲害嗎?.
如果在香港來說.
生十個真是不得了.

$^{161}$各位頂智妹妹.
有十兄弟姐妹.
真是很難得.
如果我們假選.
假設我們估計.
他的子女結婚.
生了子女.
生了孫子.
非常多.
正所謂兒孫滿堂.
所以你會發覺.
這個家庭,這個家族.
亦是非常壯大.
一個很有實力的家庭.
當時我們知道.
在遠古的時代.
如果你的子女人數多.
這亦是一個很重要的.
顯示出你的家底.
你的家庭的實力.
我們看到約伯.
如果你數一數剛才所說的.
你會發覺他真是甚麼?.
真是令人羨慕.
流口水.
但他這麼好的背景.
同時亦令魔鬼很妒忌.
認為他愛上帝.
不過因為上帝賜福給他.
他對上帝說.
如果你有一天.
將他所有的財富.
所有擁有的東西.
全部收回.
全部消除.
你看他還敬不敬畏你.
這就表現出約伯.
為何會遭受魔鬼的攻擊.
是因為魔鬼很妒忌他.
我們對約伯的背景描述.

$^{201}$暫時停一停.
接下來我們要看.
約伯之所以成功.
能夠有光彩的一面.
能夠有美好的見證.
不是單單他自身的因素.
當然他個人的因素很重要.
但另一個因素大家不要忽略.
就是剛才聖經描述.
原來他對家庭的重視.
亦是一個很重要的成功因素.
所記載的經文字句.
不是很多.
可能是幾十隻字.
但字裡行間.
你會找到.
看得見原來約伯.
能夠顯示出.
被神大大的賜福.
能夠在家庭當中.
傳家到蒙上帝的恩惠.
其實是有原因的.
為甚麼呢?.
因為約伯對家人的關懷,關愛.
做了四件事.
第一件事是甚麼呢?.
就是約伯對家人給予尊重.
對家人有尊重的心態.
這是很難得的.
為甚麼呢?.
你知道.
古舊日子.
遠古時代.
通常都是比較封建的社會.
大家留意.
尤其是在慕容,耕種,畜牧的年代.
做爸爸是很有權威的.
做爸爸是一家之主.
可以一言堂.
但約伯雖然有這樣的地位.

$^{241}$有這樣的權威.
但他沒有用盡.
特別當他知道子女長大了.
子女們都有自己的家庭.
他更加懂得尊重兒女.
給兒女們空間.
沒有管轄他們.
所以你會發覺.
當他們的兒女長大後.
如果約伯說.
你們不要常常都飲飲食食.
你們不要吃喝玩樂.
要減少,節約,懂得管理自己.
如果約伯說這些話.
都可以的.
不過他明白到.
他的兒女都是年輕.
在那個階段對吃喝玩樂.
都願意多些經驗.
所以約伯給予尊重.
第二個表現.
約伯給他們的家人兒女.
有個關心,關愛在.
何以見得呢?.
因為當他們的兒女.
每一次飲宴完畢後.
約伯都叮囑他們.
苦口婆心提醒他們.
你們要自潔.
提醒他們自潔.
不是單單沐浴潔淨自己身體那麼簡單.
原來當時的社會.
提到自潔還包括.
心靈的自省.
反省自己的思想,言語行為.
反思自己所作所為.
有沒有得罪人.
有沒有得罪上帝.
這是心靈的自省.
所以約伯提醒他們的兒女.

$^{281}$要自潔.
大家都知道.
很多時候出事.
都是在高興場合發生.
你認不認同.
我們中國人有一句說話.
叫做樂極生悲.
為什麼.
因為很多時候.
你看報紙.
通常有時吵架,打架.
都是在飲宴時發生.
有沒有留意到.
又或者被人抓.
被警察.
要去測試.
為什麼.
因為可能他走後駕駛都未定.
你會發覺到.
原來飲宴是高興的場合.
但亦是一個陷阱.
為什麼.
因為人一高興.
興奮過頭的時候.
有時會飲大了.
說話會說大了.
或者不應該說就說出來.
又說錯了.
又說大了.
又說的話無中生有.
所以很多時候.
很容易在說話之間得罪人.
或者得罪上帝.
甚至說話當上帝不存在.
所以各位頂子母們.
這就是我們人的弱點.
很容易在高興的場合.
很興奮的場合.
我們沒有警覺的心.
很容易跟別人一樣.

$^{321}$做了不應該做的事情.
所以我們更加明白.
為什麼約伯作為爸爸.
他真的很關心子女.
很想念他們.
亦很擔心他們離棄上帝.
為什麼.
因為他們經常吃喝玩樂.
飲飲厭厭的時候.
或者在過程當中.
說了不應該說的話.
做了不應該做的事情.
自然就會掉落犯罪陷阱.
得罪人又得罪神.
所以你見到約伯.
對子女很關注.
提醒他們.
這就是第二點.
我們見到約伯對家人的表現.
第三點.
我們見到約伯對家人的關心.
為家人守望.
因為每一次飲宴之後.
原來約伯不單提醒子女.
甚至約伯對自己的子女.
每一個都為他們獻犯罪.
你有沒有留意到.
個別獻犯罪.
不是一次過做完.
十個子女的犯罪.
而是每一個子女都為他們個別獻犯罪.
即是為他們個別的代禱.
為他們個別的懇求.
求上帝寬恕.
求上帝赦免.
如果他們有做錯的地方.
或者他們有未顯然的罪行.
未認的罪.
約伯都為他們代求.
為他們認罪.

$^{361}$可見這個爸爸的心態.
是一個慈父的心.
表露無遺.
今時今日我們做父母就很明白.
當中做父母就很明白.
有時說是子女結婚了.
但你都很牽掛他們.
他們有甚麼頭暈身痛.
你就很擔心.
電話不停打過去.
WhatsApp不停問候他們.
關心他們.
如果他們不是頭暈身痛.
原來媽媽爸爸要做手術.
你更加憂心掛慮.
不知手術危不危險.
不知手術過程順不順利.
當時你不單自己祈禱.
你很迫切地找身邊的弟兄姊妹.
你的親密戰友.
你的死黨.
都叫他們為你的子女做手術祈禱.
雖然父母見到子女長大.
但為人父母長憂是99.
各位弟兄姊妹們.
這就是父母的心腸.
所以約伯同樣也是.
他不單給家人關心.
亦為家人守望.
求主施恩憐憫保守他們.
使他們的腳不要走錯路.
這就是為父的心態.
所以約伯之所以成功.
是因為他以身作則.
表現做父母對兒女的守望.
守護他們.
這是一個非常難得的美好的榜樣.
第四點.
約伯對他的子女.
不單有關心有守望.

$^{401}$而且這種行動並非興之所至去做.
不是一時高興就為子女祈禱.
一時心血來潮就為子女代求.
不是的.
原來約伯這種愛家人的心.
為家人的守望.
他是持之以恆.
並非一時三刻.
並非偶爾做一下.
而是經常這樣做.
難怪他真的可以維繫整個家庭.
或者整個家族都走在上帝的心意當中.
大家明白這一點何等重要.
所以今天家庭興不興旺.
家庭蒙不蒙福.
其實作為父母.
作為家庭一份子.
每一個人都有責任.
尤其是做父母的.
更加有一個不能夠推卸的責任.
所以你發覺到.
如果我們要合乎蒙福的時候.
今天約伯已經示範了一個非常好的榜樣.
擺在我們眼前.
所以接下來.
我相信如果我問大家要舉手的話.
我猜全場都舉手了.
雖然今天不是報道會.
今天是一般的崇拜.
如果我問大家.
誰想家人蒙福.
我想大家都舉手.
誰想家人蒙恩.
大家都舉手.
這是我們學求的目標.
但是怎樣可以做到呢?.
牧師接下來建議有四個秘訣.
可以達成合乎蒙福.
第一個建議牧師給大家的是.
要為家人守望.

$^{441}$各位頂智慕們.
真的很重要.
無論你家人是否相信耶穌.
你都要為他們守望.
時時分享.
說話之間提醒.
因為家人很容易會做上帝不喜歡的事情.
或者說了上帝不喜悅的話.
甚至家人有時候因為個人的私慾.
掉進了陷阱.
這些情慾私慾的隱憂是非常之大的.
我們不要小看它.
所以我們作為家人.
就應該為家人守望.
看看哪個兄弟姐妹.
最近有些不對勁.
或者留意他們最近的生活.
有沒有問題出現.
有沒有進入試探.
這就是我們作為基督徒的家人.
我們可以做到的.
就是為家人多些守望.
其實有時候我們都會想.
我們很多時候都很愛惜家人.
我們為家人不要感染新冠病毒.
我們會提供很多防疫的措施.
用品.
經常都怕他們口罩不夠.
眼罩又不夠.
面罩又不夠.
又怕他們搓手液又不夠.
總之我們很關心他們在肉身上的防疫措施.
但是同時間我們會不會都應該去關心他們.
在屬靈的防疫措施呢?.
如果我們留意屬靈的防疫措施的時候.
我們就應該多些為他們祈禱.
為他們守望.
求上帝保守他們.
有盡我們的責任去提醒我們的家人.
來勸勉他們.

$^{481}$來鼓勵他們走在上帝的正路當中.
這是第一件事我們可以做的.
第二件事我們可以做的是什麼呢?.
就是我們可以為家人去代禱.
這是我們最基本的動作.
就是我們為他們祈禱.
有時候在家裡.
這麼多兄弟姐妹.
你會發覺最小的又是你.
排到最後的又是你.
身形最小的又是你.
說話小聲的又是你.
總之你好像在家裡沒有地位.
沒有人重視.
不過上帝是重視你的.
只要你多些為家人祈禱.
為他們提名的代禱.
個別的為他們祈禱.
求上帝保守你的家人.
求上帝令你家人的心軟化.
能夠接受耶穌做救主.
這就是你和我每一個人都可以做到的.
只要你這樣來祈禱的時候.
為你家人代禱的時候.
我相信上帝會聽你的祈禱.
在適當的時候機會.
上帝就工作令你的家人.
一個一個來歸順主耶穌.
各位短期節目的這些見證.
在你和我身邊都很多時候聽到.
發覺家人沒辦法信耶穌.
或者老人家很難信耶穌.
很硬頸的.
很難過,毫無禮貌的.
那你怎樣帶他們信耶穌呢?.
牧師也遇過這些人.
但是上帝給我有機會.
和他們做個人工作.
在床邊慰問他們.
在床邊和他們聊天.

$^{521}$慢慢地老人家肯信耶穌.
願意信耶穌.
又願意接受水禮.
其實你看到你的祈禱是不會白費的.
只要你多些為家人禱告.
上帝就可以行其事.
施展祂大能的作為.
記得千萬不要讓家人覺得.
上帝好像不存在.
今天人最嚴重,最大的問題是什麼呢?.
就是當上帝沒有到.
當祂不存在.
目中無神.
這是最可怕的事情.
各位短期節目的你認同嗎?.
如果當一個人對上帝目中無神的時候.
祂的破壞,祂做出來的行為.
真的可以令人恐怕,恐懼.
所以我們真的需要求上帝幫助我們.
我們叫我們家人.
目中有上帝.
都願意聽上帝的話.
以致我們全家都得到上帝的喜悅.
這是第二件事我們可以做的.
第三件事我們可以做的是.
為家人報道.
即是我們為家人傳福音.
無論弟兄姊妹平時帶父母,家人.
回教會,報道會.
這是一個很好的行動.
其實我們真的要多關心家人.
不單單是他們肉身上的需要.
關心家人肉身的需要.
我們可能都做得很多.
請他們吃飯,逛街,購物.
和他們去日本,台灣旅行.
這些你和我都做了很多.
為家人在肉體上的需要.
而做出的愛心關懷.
但是我們有沒有對家人靈魂的需要.

$^{561}$他們靈魂的得救.
來做一些功夫.
這是很重要的.
所以特別在你家裡.
還有未信主的家人.
弟兄姊妹們.
記得要為家人報道.
現在的科技真的很先進.
你有沒有留意到.
就算沒有了實體崇拜.
沒有了實體報道.
你都可以和家人傳福音.
至於你有部手機.
你一上YouTube.
弟兄姊妹們.
你找一些見證的視頻.
找一些見證的節目.
難不難找?.
不難找的.
很多的.
好像我們早一個星期五.
我們福音城會晚會.
都是原定是實體.
結果沒有實體.
但是那些弟兄姊妹.
在網上看YouTube.
介紹家人去看.
又有陳友的見證.
又有謝牧師的講道.
真的一個星期都不夠.
有八千多人上網去看.
弟兄姊妹們.
你有沒有留意到.
原來你今天利用這個科技.
是突破了很多空間的限制.
時間的限制.
你只要將你的手機放大在電視上.
給你父母看.
他們就可以看到.
今天傳福音.

$^{601}$信得給我們很多機會.
有一次父親節的時候.
我們有逆境同行.
愛是不保留的Online音樂會.
同樣透過YouTube.
透過網上的發放.
有四十萬人點擊去看.
我知道有不少弟兄姊妹的家庭.
都預了時間.
下午坐定.
吃個包.
喝杯茶.
開著電視.
和家人欣賞音樂會.
看當中見證的分享.
這些都是很好的傳福音的途徑.
特別現在疫情期間.
沒有實體聚會.
大家可以利用這些方法.
來傳福音.
讓我們突破了種種的障礙.
第四件事我們可以為家人做的.
就是為家人崇拜.
即是大家可以在家裡.
設立家庭崇拜.
就算你說.
牧師我家人不是全部都信耶穌.
只有兩名.
兩名都夠了.
耶穌說地上有兩個人同心祈禱.
我在他們當中.
你抓住你姐姐.
抓住你哥哥.
弟弟妹妹.
和他們一起祈禱.
簡單地說.
兩個人一起祈禱.
讀聖經.
敬拜上帝.
如果有三個更加好.

$^{641}$有四個就更加不錯.
全家信主就更加好.
好像我自己的家庭.
由我爸爸的年代開始.
由家庭崇拜.
到我這一代.
家庭崇拜.
到Mary Stephen那一代.
由家庭崇拜.
到孫子那一代.
都一起來.
家庭崇拜.
這是上帝給我們的恩典.
將一個很好的家庭崇拜模式.
一代傳一代下去.
其實家庭崇拜不是很困難.
叫家人一起.
大家手牽手也好.
不牽手也行.
總之一起唱詩歌.
一起讀經文.
一起聊天.
一起祈禱.
很快就半小時.
整個小時就完結.
所以家庭崇拜.
是培育我們的家人.
從小就開始認識上帝.
從小就開始得到靈明方面的餵養.
當然如果只有一個人.
一個人沒有辦法.
一個人唯有單獨祈禱.
但你努力.
希望上帝讓你明年家裡有兩個人.
後年有三個人.
更加好.
有一個家庭崇拜.
到我結束的時候.
頂指滿門.
家庭真的很重要.

$^{681}$為什麼?.
因為我們中國人有一句話.
叫做收齊置評.
收齊置評.
其中就是要齊家.
修心齊家.
治國平天下.
原來家是很重要的.
如果我想做一些大事.
但家庭不能和諧.
就很難做大事.
收齊置評.
大家很認同.
大家也知道.
如果我們社會.
需要健康的社會.
而這個社會更加需要.
有更加多幸福的家庭.
蒙福的家庭.
因為這個社會.
是由這些家庭作為核心去建立.
如果我們能夠更多和諧.
幸福的家庭.
這個社會的建立更加健康.
所以是非常重要.
當然聖經也很重視家庭的觀念.
所以在提摩太前書.
第三章第五節告訴我們.
若有人不能夠管理自己的家.
焉能管理上帝的家呢?.
這句話也讓我們看到.
原來聖經很重視每一個人的家庭.
作為我們事奉的基礎.
如果我們的家庭不能夠蒙福.
不能夠敬畏上帝.
和諧,幸福的時候.
我們在事奉上.
真的有一個很大的阻礙.
最後要跟大家談得近一點.
就是近來大家知道.

$^{721}$疫情令很多教會.
門也封了.
沒有了活動.
沒有了崇拜.
很多東西都沒有了.
不過你有沒有留意到.
家庭的門是門常開的.
家庭的門是永遠都打開的.
只要我們能夠好好將.
我們屬靈的追求.
靈性的培育.
能夠用心機.
將它建立在我們家庭裡面.
我相信我們的家庭.
一定是一個蒙福的家庭.
願意今早我們看這段經文.
透過約伯一個美好的典範.
叫我們更加建立一個很清晰的目標.
就是我們追求我們的家庭.
是一個合乎蒙福的家庭.
我們低頭祈禱.
慈愛天父多謝你.
透過今天的經文.
再一次提醒我們.
原來你給我們有溫暖的家.
快樂的家.
幸福的家.
都不是必然的事情.
是需要我們努力去耕耘.
努力去依靠上帝.
盡我們的本份.
正如約伯一樣.
能夠有樹立美好的榜樣.
令到他的兒女.
都認識上帝.
敬畏上帝.
天父同樣求你祝福.
我們紅恩堂的弟兄姊妹.
我們的家庭.
我們的兒女.

$^{761}$都能夠蒙福.
都能夠因著我們的見證.
認識上帝.
能夠敬畏上帝.
以致我們的家庭.
都是全家得著恩典的家庭.
多謝你.
祈禱奉耶穌的名求.
阿們.
\newpage



\section{彼得後書 1:3-11}
\label{sec:OSwpE_DOZH0}
\textbf{2020.09.06《信徒長進的人生》 彼得後書 1\_3-11 謝靄薇牧師}
\newline
\newline
連結: \href{https://youtube.com/watch?v=OSwpE-DOZH0}{\texttt{https://youtube.com/watch?v=OSwpE-DOZH0}} ~~~~ 語音日期: 2020-09-09
\newline
\newline
\hyperref[sec:YQ9vapt0XGE]{\small{< < < PREV SERMON < < <}}
~
\hyperref[sec:index_chronic]{\small{[返順時目]}}
~
\hyperref[sec:index_scriptual]{\small{[返順卷目]}}
~
\hyperref[sec:RqJiqD22j9E]{\small{> > > NEXT SERMON > > >}}
\newline
\newline
彼得後書 1:3-11
\newline
\begin{longtable}{cl}
\hline
\hline
章節 & 經文 (和合本修訂版)\\
\hline
1:3 & \begin{tabularx}{0.7\textwidth}{X} 神的神能已把一切關乎生命和虔敬的事賜給我們，因我們認識那用自己榮耀和美德召我們的神。 \end{tabularx} \\ \\ \relax
1:4 & \begin{tabularx}{0.7\textwidth}{X} 因此，他已把又寶貴又極大的應許賜給我們，使我們既脫離世上從情慾來的敗壞，就得分享神的本性。 \end{tabularx} \\ \\ \relax
1:5 & \begin{tabularx}{0.7\textwidth}{X} 正因這緣故，你們要分外地努力。有了信心，又要加上德行；有了德行，又要加上知識； \end{tabularx} \\ \\ \relax
1:6 & \begin{tabularx}{0.7\textwidth}{X} 有了知識，又要加上節制；有了節制，又要加上忍耐；有了忍耐，又要加上虔敬； \end{tabularx} \\ \\ \relax
1:7 & \begin{tabularx}{0.7\textwidth}{X} 有了虔敬，又要加上愛弟兄的心；有了愛弟兄的心，又要加上愛眾人的心。 \end{tabularx} \\ \\ \relax
1:8 & \begin{tabularx}{0.7\textwidth}{X} 你們有了這幾樣，再繼續增長，就必使你們在認識我們的主耶穌基督上，不至於懶散和不結果子了。 \end{tabularx} \\ \\ \relax
1:9 & \begin{tabularx}{0.7\textwidth}{X} 沒有這幾樣的人就是瞎眼，是短視，忘了他過去的罪已經得了潔淨。 \end{tabularx} \\ \\ \relax
1:10 & \begin{tabularx}{0.7\textwidth}{X} 所以，弟兄們，要更加努力，使你們的蒙召和被選堅定不移。你們實行這幾樣，就永不失腳。 \end{tabularx} \\ \\ \relax
1:11 & \begin{tabularx}{0.7\textwidth}{X} 這樣，必叫你們豐豐富富地得以進入我們主－救主耶穌基督永遠的國度。 \end{tabularx} \\ \\
[1ex]
\hline
\hline
\end{longtable}
$^{1}$(在座的各位).
各位弟兄姊妹平安.
我們在家中的.
雖然見不到你們.
但我們求神在眼中引領我們.
我們一起有一個祈禱.
我們恩主當我們來到你面前的時候.
雖然我們受到疫情有很多的限制.
我們弟兄姊妹他們在家中.
和我們好像有很大的距離.
但我們知道神無論你在哪個地方.
你都可以親近,熟悉你的人.
由我們在你面前守著心清的.
我們去等候你.
願你帶領我們的心.
讓孩子在你面前所說的.
心裡所想的.
都蒙你月亮.
奉主耶穌基督的名頭.
阿們.
今天我們看這段經文的時候.
首先我們先了解一下.
當時期彼得寫的後書的時候.
當時的情況我們了解一下.
原來那個時候.
正正是這班基督徒.
受到尼羅王很大的逼迫.
特別是很劇烈的時候.
教會裡面出現了假教師.
於是彼得就要在當中.
去揭發他們一些不好的動機.
同時要去抵擋這班假教師的異端邪說.
所以他要再次的.
在這本書裡面去強調.
耶穌基督的真理.
要高舉這本聖經的權威.
要讓他們知道信心的根基在哪裡.
還有耶穌基督再回來.
這些真理要讓他們明白.
所以我們看到雖然.

$^{41}$這班基督徒.
有些因為假教師的緣故.
可能他們離開了教會.
但是發現越大的逼迫.
基督徒的人數並沒有因此而消失.
反而是越來越多.
我們發現信徒在受逼迫的光景裡面.
他們怎樣可以經得起考驗呢?.
怎樣可以堅持信仰到底呢?.
其實很在乎信徒的生命.
今天我們已經相信了耶穌基督.
我們有了耶穌基督的生命.
不如我們從經文裡面.
去檢視一下我們的生命是怎樣的.
在彼得後書的一章一節.
講到耶穌基督的僕人和仕徒西門彼得.
寫信給那些因我們的神.
和救主耶穌基督的義.
與我們同德一樣寶貴信心的人.
今天你和我都一樣得到寶貴的信心.
所以我們第一方面要看信心.
信心我們大家都知道.
我們在不同的時間我們都信主了.
我們信主的時候.
神都應許我們.
好像羅馬書十章九節那裡.
他講到你要考慮.
宣認耶穌基督.
真的相信祂和死律復活.
就必得救.
因為是神給我們的應許.
但是真正的信心都有表現.
我不說深信就行了.
我又不用去祈禱.
又不用去上教會.
我們的外在的表現是怎樣.
我們現在每個人都有手機.
我們一按就可以看到經文.
可能你在聚會的時候.
你都會看手機.

$^{81}$但不知道你在看什麼.
不知道是不是在看經文.
但是在十幾年前.
我們信徒他們會買聖經.
會帶聖經回教會.
看自己的聖經.
但是有些信徒他就很害怕.
很害怕別人知道他是信徒.
所以他會包著聖經.
回到教會很神秘地打開聖經.
那就要看你的心態.
你的心態到底我們當神是什麼呢?.
為什麼不可以見到光呢?.
講這個外在的表現.
信心.
我們又看看.
就是《已不所需》的二章八節說.
你們得救是本會恩.
也因著信.
你知道這個信是神給你們的.
我們因著信.
耶穌基督給了你們.
不是因為自己是神所賜的.
也不是因為行為.
免得有人自誇.
所以我們知道.
得救是本會恩.
是完全是耶穌基督賜給我們的.
有了這個信心.
我們就有一個永恆的生命.
我們可以得到拯救.
就因著有這個信心.
我們可以進入信心的恩典裡面.
所以彼得在這裡.
很珍貴地提醒我們.
我們今天之所以有這麼寶貴的信心.
是因著耶穌.
為我們付出了一個很重的代價.
去洗淨我們的罪.
以致給我們一個新的生命.

$^{121}$我們就可以因著耶穌基督的義.
去承受神的義.
這是很重要的.
然後.
在剛才你們讀的第三節.
祂告訴我們.
神的神能以把一切.
關乎生命和乾淨的事賜給我們.
因我們認識那樣自己榮耀和美德.
召我們的神.
看到嗎?.
耶穌基督從世界裡面選召了你.
祂用榮耀和美德去召你.
祂把生命和乾淨的事賜給我們.
所以我們第二方面.
就要看看生命和乾淨的事是怎樣的.
生命就是指屬靈的生命.
我們未信主的時候.
我們很注重肉身.
但今天我們信了耶穌基督之後.
我們有改變.
我們有新的生命.
我們成為一個新造的人.
而乾淨的事是什麼呢?.
就是我們領受神恩典的人.
我們應該敬虔度日.
我們應該敬畏神.
用敬畏的態度去服侍我們的神.
要遵行祂的旨意.
是一個敬虔的人.
他們得到神的救恩.
就應該是這樣的.
所以我們看到.
我們的屬靈的生命.
是應該要成長的.
也是必須要成長的.
但成長不是必然的.
要做點事.
就像信主的時候.
上初信班.

$^{161}$你一定要留意何時可以上課.
你要去才可以.
不是報名.
上課之後.
教了你一些東西.
你要去思想一下.
你學到東西了.
這樣你才會有進步.
所以我們看到.
原來是要開始去做的.
不是坐著等著別人給你飯吃.
這樣你就不能成長了.
所以成長不是必然的.
你必須要有計劃的.
曾經在外國有一個老太太.
叫做羅斯.
她87歲.
她去讀大學.
那學期結束的時候.
學校就邀請她.
一個這麼年輕的學生.
就叫她演講.
她講了一句名句.
很轟動全世界.
她講什麼呢?.
她說不是因為年老.
而停止玩樂.
而是停止玩樂.
人就會變老了.
的確的.
你要青春常駐.
要經常快樂.
那你必須要讓自己笑口常開.
要保持幽默和風趣.
同時.
要看到有一個目標.
你要有夢想.
當你夢想沒有了的時候.
你就暗淡無光.
所以有目標你就會有行動.

$^{201}$有動力去做事.
所以變老和長大.
是很大的差別.
變老就是任何人都會變老的.
但是長大.
不是每個人都會長大.
你會看到長大的意思是什麼呢?.
你必須在不斷的蛻變當中.
去尋找成長的機會.
然後去加以利用.
你要活得無怨無悔.
那些上一年級人.
他從來都不會為他所做的事.
去後悔的.
但反而年輕的時候.
有些事想做,但沒有做.
就覺得很遺憾.
所以.
如果人心裡懷著一個悔恨的話.
他就很怕死亡.
所以我們第二方面.
我們來看看.
一個成長的信徒的生命是怎樣的呢?.
就是說我們熟悉的生命.
我們要長進.
首先我們看到的是.
原來我們要追求.
熟悉的生命要長進的態度.
就是要憤外的努力.
以前的版本就是.
憤外的殷勤.
要很努力的.
同時.
在我們讀的五至七節裡面.
就會看到.
當你願意去追求.
長進的時候.
你的熟悉的生命.
就會擁有以下的美德.
以下的美德是什麼呢?.

$^{241}$這些美德是一步一步來的.
首先.
就看到有信心.
原來開始.
有了信心.
這個信心他先提.
為什麼呢?.
因為信心是一個基礎.
沒有它是不行的.
我們有了耶穌基督的生命.
我們進入耶穌基督裡面的時候.
我們就跟耶穌基督的聖情有份.
當我們有了耶穌基督的聖情的時候.
在彼得前書一章七節裡面說到.
使你們的信心既被考驗.
就比那被火試煉.
仍然能壞的金紙更顯寶貴.
可以在耶穌基督顯現的時候.
得著稱讚榮耀尊貴.
當你去煉金紙的時候.
你會看到.
金紙.
外面賣得很貴.
要煉.
蒸金.
經過火煉.
高溫度裡面燒出來的.
在後書裡面.
前書說到.
金紙形容我們的信心.
要這樣的考驗.
所以可見這個考驗不是簡單的事.
但是世上的金.
是蒸金.
人們很看重的.
都會過去的.
而信心經過考驗的時候.
這個形容到耶穌基督回來的時候.
你可以得到稱讚.
得到榮耀.

$^{281}$得到尊貴.
是多麼寶貴的事.
所以信心很重要.
有了信心又要加上德行.
德行就是我們信了主之後.
我們不是只信的.
因為雅各書教導我們.
有信心沒有行為是死的.
所以你要外在有表現的.
所以人們看到你.
這個基督徒.
這個品質的標牌.
真是信得過的.
這個德行.
所以你應該要有一個好的品德.
去表現出來.
因為你已經有耶穌基督的生命.
你外在.
正如聖經說的.
要發出清香的味道.
所以德行上.
都不是臭名的.
要發出好的品德.
然後.
他又加上什麼呢?.
加上知識.
知識不是叫人自高自大嗎?.
唯有愛心可以造就人.
為什麼這個道理加上知識呢?.
可見知識不是只講頭腦上的知識.
而是怎樣去找尋認識耶穌基督的知識.
大家都有機會去探人.
當你去探人的時候.
主人就會客氣叫你請坐.
我相信你多數都會坐下.
你不會慢慢去檢查別人的椅子.
除非真的有問題.
否則通常我們就會坐下.
然後去聊天.
可能你坐下後會覺得很舒服.

$^{321}$其實我們在信耶穌的過程中.
都是我們用單純的信心.
我們信了主之後.
然後在以後的日子中.
這個過程中.
我們去經歷.
經歷神是這麼實在的.
我們經常聽到老人家.
他們去講見證.
他們可能沒什麼機會讀書.
更加不會去讀神學.
但是為什麼他們對神的信心.
是打風也打不掉呢?.
就是因為他們在艱難的過程中.
在順境的過程中.
看到神怎樣一步一步地帶他們經過.
所以他們在經歷的時候.
他們就很實在.
所以這裡說的知識.
就是因為你認識神.
而有一個知識是很寶貴的.
有了知識又要節制.
節制就是提醒我們什麼呢?.
有時候我們欲望無窮.
有些東西想學.
很多東西要學.
但是我們要節制.
我們不要只學字.
每樣東西都不精.
有一個小朋友.
他很想爸爸陪他.
爸爸經常給他一些錢獎勵他.
於是他存了很多錢.
有一天他就問爸爸.
爸爸你一天賺多少錢呢?.
他說我這裡很多錢.
不如我給你.
你陪我玩.
其實你看到這個小朋友很嚮往什麼呢?.
就是爸爸陪他.

$^{361}$其實做大人的.
去上班是很平常的事.
但是如果我們沒有節制的話.
我們不懂得去計劃.
如何去分配我們的時間.
就是你永遠都沒有時間陪他.
你看人家統計裡面.
有多少時間留在家陪家人.
或者有多少時間.
和他們的小朋友一起聊天.
其實少之又少.
如果我們沒有節制.
不好好地去分配我們的時間.
永遠都沒有時間.
忍耐.
特別是在香港.
大家都很急.
塞車就有很多埋怨.
但是你見到拍拖的人又不同.
等女生等多久都無所謂.
很有忍耐.
見到之後就不知道.
你見到就是因為忍耐真的很難.
在忍耐上都是熟齡的果子.
我們在事奉的路上.
是需要忍耐的.
我們不能一步登天.
不能一信主就要懂很多東西.
所以信主的人.
我們不會那麼快去學習事奉的事.
要慢慢學習.
就像做導師.
他平時都要看聖經.
好好鑽研.
經過培訓.
差不多一天了.
他才可以做一個稱職的導師.
在熟齡的路上.
都是要這樣.
要慢慢地去浸泡.

$^{401}$我們要忍耐等候.
我們就會成長成熟.
就好像一個運動員.
他要經過相當的時間.
怎樣去訓練他.
去考核.
或者比賽.
他才有今天的成績.
所以真的要忍耐的.
不會馬上就看到效果的.
然後要加上虔敬.
虔敬就是很在乎.
你處事的態度.
你說你是一個信神的人.
你的處事是怎樣呢?.
譬如你是一個老闆.
你對你的員工會不會很苛刻呢?.
或者你這些人比較好談的話.
你就樣樣都可以.
那個人是看不順眼的.
你就刁難他.
但是你當神是什麼呢?.
如果有神的話.
如果有公義的話.
那你為什麼要用這樣的態度去對待人呢?.
虔敬.
如果你對神是一個很虔敬的人.
那你對人呢.
應該是很恭敬的.
不應該對這些人很好的.
對那些你看不順眼的.
就很偏心.
你對他不好.
這個不是神要我們做的態度.
就好像大衛.
他的兒子亞撒龍.
造反的時候.
碰到一個師妹.
師妹看到他就罵他.
連屬下都看不慣.

$^{441}$為什麼罵皇帝呢?.
而大衛說.
不要緊.
可能神叫他罵我的.
可能神因為他罵我.
會施恩給我.
如果是你又會怎樣呢?.
反過來跟他對罵.
但是跟他對罵的時候.
其實就很失儀態.
又失見證.
因為我們是屬神的兒女.
我們既然要虔敬的話.
我們對人呢.
我們的態度會不會是失了見證.
而沒辦法在神的面前.
我們都是很差的態度.
被人看到.
我還敢信主.
所以亞國書教我們.
就是要快快的聽.
慢慢的說.
慢慢的動腦.
當你想清楚再說.
想清楚再發脾氣的話.
相信你就會減低很多.
不標的犯錯.
有了虔敬.
又要愛弟兄的心.
我們的心.
就是提醒我們.
在弟兄姊妹當中.
我們有些很熟的.
有些可能很生疏的.
但我們都要在別人面前.
是一樣的.
一視同仁.
用耶穌基督的愛去愛他.
因為聖經說.
你們若有彼此相愛的心.

$^{481}$眾人因此就認出.
你們是我的門徒.
所以我們是一樣的.
無論我們跟他熟不熟.
我們都應該關心他.
用耶穌基督的愛去愛他.
愛眾人的心.
不單止是我們主內的弟兄姊妹.
你也有外面的朋友.
有些人回到教會之後.
全部都是教會的朋友.
這樣是不行的.
我們也應該有外面的.
微信主的朋友.
不然我們怎麼傳福音.
神愛世人.
我們也應該學耶穌基督.
我們也要愛外面的.
微信的朋友.
讓他們認識神.
以致我們求神給我們.
諸般的智慧.
將他們帶到神的面前.
讓他們能夠歸順主.
第八節就告訴我們.
我們有了這麼多.
我們還要做什麼?.
繼續增長.
我們不要自滿.
我們要繼續增長.
所以我們第三方面.
我們看看.
當我們自滿了.
覺得可以了.
我又有幾樣.
這麼辛苦了.
熬到現在了.
我休息一下也行吧.
原來這裡說到.
停止進步是有些危機的.

$^{521}$懶散.
和不潔果子.
可以閒懶了.
你看看.
你好像種東西.
你種東西.
你也要陽光雨水.
你要隨從.
施肥.
很多事情要做的.
是不是?.
不會坐著就有果子吃的.
可能那些是前人種的.
所以說到就是.
我們懶懶也不行的.
你要修成.
你要去努力的.
要做這麼多事情.
是為了什麼?.
要潔果子.
所以如果我們不長進的話.
這個危機就會出現了.
懶散.
聖經說.
我們不努力的話.
在前的在後.
在後的在前.
那我們就退步了.
我們不要說.
我上了教育很多年了.
我信主你不知道在哪裡.
但不是這樣比的.
如果我們不努力.
不長進的話.
在後的就在前了.
同時更加恐怖的.
聖經說到.
我們不長進的話.
就是那個又惡又懶的僕人.
連你所有的都要給別人.

$^{561}$而你變成了在什麼光景?.
就是哀哭切齒.
在耶穌的門外.
同時沒有份.
是不是很慘?.
所以我們真的要努力.
不要退後.
第二個危機.
就是黑暗.
短視.
黑暗就是說.
不是說你真的瞎了.
在《約翰一書二章十一節》.
說到唯獨那恨弟兄的.
是在黑暗裡.
也在黑暗裡行走.
不知往哪裡去.
因為黑暗使他的眼睛黑了.
因為說到.
是屬靈方面你瞎了.
你做了錯事.
而短視呢.
就是你只能看到眼前的.
眼前有機會賺大錢.
別人這樣做.
我跟著去.
一窩蜂去.
或者你賭一把.
將你整副身家放進去.
為了什麼呢?.
其實有時候我們很短視.
或者為了升職.
或者想著一些手段.
怎樣令自己得到名利.
其實當我們有這樣的想法時.
其實無形之中.
我們已經出賣了兒女的名份.
因為我們成為神的兒女.
我們應該要遵行神的旨意.
但是我們做了這麼多事.

$^{601}$但這不是信徒應有的本份.
就好像一隻小鳥.
進了一間屋裡.
其實這些情況.
我們在錦繡大堂經常都看到.
有一隻小鳥飛了進去.
兒子很怕他撞到窗戶.
很怕他受傷.
撞來撞去都走不出去.
其實小鳥可能說.
我明明看到外面的.
為什麼飛不出去呢?.
他不知道.
玻璃是一個無形的障礙.
其實同樣的.
在我們的人生裡.
我們稍微有些貪念.
或者我們眷戀名利的時候.
這些都是無形的障礙.
難阻了耶穌基督進入我們的心門.
難阻了祂將福氣給我們.
而我們只顧著眼前.
這些障礙難阻了我們.
不可以向前.
這是很可惜的.
我們要去注重永恆的.
你看看你今天.
雖然你可能很風光.
或者很有名氣.
但是都會過去的.
有一天我們兩腳一伸.
什麼都沒有了.
但是我們去到神面前.
我們有什麼交給祂呢?.
我們在這個地上的時候.
我們要去追求永恆的東西.
不要這麼短視的.
只顧著別人有的.
我們也要有.
而浪費我們的光陰.

$^{641}$第三個更加可惜的.
就是我們忘了過去的罪.
已經得到潔淨.
第一次我們想到.
我們是被耶穌重駕買贖回來的.
祂救了我們.
但我們竟然忘恩負義.
我們的眼睛注重罪惡.
我們忘記了我們已經蒙潔淨.
我們應該過一個分別為勝的生活.
但我們一點都沒有.
耶穌基督給了我們生命的樣式.
我們沒有過新的生命的樣子.
這是很可惜的.
不能活出一個生命.
所以難怪彼得.
他知道他日子無多.
所以他不斷提醒我們.
雖然你們已經知道這件事.
你們回到教會很多年了.
你們什麼都知道了.
但我還是要提醒你們.
你們已有的真道上得到堅固.
現在很穩定.
我還是要常常提醒你們這些事.
他在老人家不斷提醒我們.
因為他自己走過這些路.
所以他不想我們走錯路.
所以在十三節他提醒我們.
弟兄們要更加努力.
要我們努力.
使你們的夢照和被選堅定不移.
今日神選照你.
給你有這麼異類的名份.
給你有這個二人的地位.
我們不要退後.
甚至失去.
這就很可惜.
因為神在萬人當中選擇你.
所以我們有了這幾樣.

$^{681}$我們就永不失腳.
他叫我們不斷地去成長.
要進步.
成長的表現是什麼呢?.
就是結果節.
要努力.
要殷勤去工作.
當我們靠聖靈去幫助我們.
保守我們的心.
我們就可以站立得穩.
盼望我們真的可以得到神.
讓我們穩步向前.
以致能夠榮譽我們的主.
能夠進入神的榮譽國度.
我們一起祈禱.
你這麼恩待我們.
你選照我們.
最後求你幫助我們.
透過經文提醒我們.
不要卻步.
不要懶惰.
讓我們能夠穩步向前.
繼續的.
一步一步的.
去認見你的面.
都能夠被你稱讚.
我們是由你由忠心.
由良善的僕人.
求你恩待我們.
加力給我們.
讓我們在情境不穩定的時候.
我們仍然可以為你忠心到底.
我們這樣祈禱交托.
奉主耶穌基督明求.
阿們.
\newpage



\section{士師記 6:11-18}
\label{sec:RqJiqD22j9E}
\textbf{2020.09.13《從軟弱到剛強》 士師記 6\_11-18 (和修版) 老耀雄傳道}
\newline
\newline
連結: \href{https://youtube.com/watch?v=RqJiqD22j9E}{\texttt{https://youtube.com/watch?v=RqJiqD22j9E}} ~~~~ 語音日期: 2020-09-15
\newline
\newline
\hyperref[sec:OSwpE_DOZH0]{\small{< < < PREV SERMON < < <}}
~
\hyperref[sec:index_chronic]{\small{[返順時目]}}
~
\hyperref[sec:index_scriptual]{\small{[返順卷目]}}
~
\hyperref[sec:7bIUWEs7F1E]{\small{> > > NEXT SERMON > > >}}
\newline
\newline
士師記 6:11-18
\newline
\begin{longtable}{cl}
\hline
\hline
章節 & 經文 (和合本修訂版)\\
\hline
6:11 & \begin{tabularx}{0.7\textwidth}{X} 耶和華的使者到了俄弗拉，坐在亞比以謝族約阿施的橡樹下。約阿施的兒子基甸正在醡酒池那裡打麥子，為了躲避米甸人。 \end{tabularx} \\ \\ \relax
6:12 & \begin{tabularx}{0.7\textwidth}{X} 耶和華的使者向基甸顯現，對他說：「大能的勇士啊，耶和華與你同在！」 \end{tabularx} \\ \\ \relax
6:13 & \begin{tabularx}{0.7\textwidth}{X} 基甸對他說：「主啊，請容許我說，耶和華若與我們同在，我們怎麼會遭遇這一切事呢？我們的列祖告訴我們：『耶和華領我們從埃及上來』，他那奇妙的作為在哪裡呢？現在耶和華卻丟棄了我們，把我們交在米甸人的手掌中。」 \end{tabularx} \\ \\ \relax
6:14 & \begin{tabularx}{0.7\textwidth}{X} 耶和華轉向基甸，說：「去，靠著你這能力拯救以色列脫離米甸人的手掌。我豈不是已經差遣了你嗎？」 \end{tabularx} \\ \\ \relax
6:15 & \begin{tabularx}{0.7\textwidth}{X} 基甸對他說：「主啊，請容許我說，我怎能拯救以色列呢？看哪，我這一支在瑪拿西支派中是最貧寒的，我在我父家又是最微小的。」 \end{tabularx} \\ \\ \relax
6:16 & \begin{tabularx}{0.7\textwidth}{X} 耶和華對他說：「我與你同在，你就必擊敗米甸，如擊打一個人。」 \end{tabularx} \\ \\ \relax
6:17 & \begin{tabularx}{0.7\textwidth}{X} 基甸對他說：「我若在你眼前蒙恩，求你給我一個證據，證明是你在跟我說話。 \end{tabularx} \\ \\ \relax
6:18 & \begin{tabularx}{0.7\textwidth}{X} 求你不要離開這裡，等我回來，將供物帶來，供在你面前。」他說：「我必等你回來。」 \end{tabularx} \\ \\
[1ex]
\hline
\hline
\end{longtable}
$^{1}$洪影家,弟兄姊妹平安.
有很多人都說過.
如果要評價一個人的勇氣.
是需要參照他在面對困難的時候.
他所作出的決定和行為.
然後才能夠準確地作出一個正確的評價.
因為一般人在遇到困難和危險的時候.
他們都會有一個很好的思考.
大多數人都會選擇逃避.
不太想去面對.
唯有真正有勇氣的人.
才會選擇迎難而上.
並且作出解決困難的部署和計劃.
有人問.
這種勇氣是否與生俱來.
如果不是與生俱來的話.
有沒有方法可以培養呢.
或許坊間有很多不同的方法可以提供.
但作為信徒的我們.
最佳的方法就是從信仰裡提升我們勇氣的能力.
今天我很想透過《攝施記錄》.
將十一至十八節.
機電的一段生命故事.
和大家一起反思.
一個從軟弱到剛強的秘訣.
我們一起去祈禱吧.
因為無論我們的能力是多麼軟弱.
你都願意使用我們.
一切計劃的成功不是因為我們的能力.
而是因為主你的大能.
你透過我們微小的能力.
成就你偉大的計劃.
願榮耀歸給我們至高至大的主.
我們的祈禱是奉靠主耶穌基督的聖名祈求.
開始的時候.
我們首先和大家講講《攝施記》的背景.
《攝施記》一開始就講到.
約書亞死了之後.
以色列人就去求問耶和華.
他說在我們中間.

$^{41}$誰會是首先攻打迦南人呢?.
與他們爭戰.
原來《攝施記》是一個戰爭的故事.
約書亞在生的時候.
他未能帶領以色列民完全地去征服迦南.
因此當約書亞死了之後.
十二支派仍然要在所屬地域.
與迦南當地的部族去打仗.
《攝施記》結束的時候.
經文是這樣說.
那時以色列中沒有王.
各人照自己眼中看為對的去做.
以色列民在《攝施記》開始的時候.
有事的時候都會求問耶和華.
顯示以色列民以耶和華為首.
對耶和華有一份敬畏的心.
到結束的時候.
以色列民就各人照自己眼中看為對的去做.
顯示耶和華在以色列民心中已經沒有地位.
以色列民過的是一種.
以自我為中心的生活形態.
我們知道第一個.
《攝施記》沒有太多缺點.
但到最後一個《攝施記》參選的時候.
你就會說.
什麼這樣的人都可以做攝施嗎?.
真是一孩不如一孩.
由以耶和華為先.
去到當耶和華沒有刀.
我們就看到《攝施記》所記載的.
是一群屬靈生命一直倒退的人物和故事.
在《攝施記》裡面.
縱然以色列人是多麼不堪.
多麼遠離耶和華.
但每當以色列民被欺壓.
向耶和華呼求的時候.
耶和華就興起攝施去幫助以色列人去打敗敵人.
不過一旦太平盛世.
以色列人再次得罪耶和華.
以色列人的屬靈生命.

$^{81}$就好像走兩步就退三步.
最後退到一個地步.
就好像經文結束的時候所說.
他們眼中已經沒有耶和華.
心中所想的就只有自己.
基建是其中一個設施.
它沒有愛特列那麼完美.
但又不至於去到參選那麼不堪.
不過不失.
中中挺挺.
就好像我們這樣.
我們不是最好的那個.
我深信我們也不是最差的那個.
正正如此.
基建的生命更值得我們細心地審視.
看看有什麼值得反省和學習的地方.
基建的生命經歷很豐富.
《攝施記》用了三章的經文.
第六章到第八章.
講述它由夢照講到死亡.
所以一篇講道.
應該講不完整個基建的故事.
今天我只是選擇了.
講述它一個屬靈生命的片段.
希望下次有機會可以講其他的部分.
我們開始進入經文.
基建生命中第一個反省.
就是信仰的偏差.
會令人產生不滿和埋怨.
第六章一開始.
就講到以色列人被米顛人欺壓了七年.
每當以色列人殺種之後.
米顛人就來攻打他們.
毀壞他們所有的農作物.
並且沒有留下任何食物.
連牛,羊,驢都沒有留下.
可以這麼說.
以色列人的確生活得很苦.
所以以色列人再次去呼求耶和華.
耶和華如何回應以色列呢?.

$^{121}$在第六章第十二節和基建講.
大能的勇士耶和華與你同在.
你看基建如何回應.
基建在十三節就這樣回應.
主啊!請容許我說.
耶和華若與我們同在.
我們怎麼會遭遇這一切事呢?.
我們的列祖告訴我們.
耶和華領我們從埃及上來.
他那奇妙的作為在哪裡呢?.
現在耶和華卻丟棄了我們.
把我們交在米顛人的手掌中.
基建認為.
如果耶和華真的與我們同在.
以色列就不會像現在這樣的光景.
而且基建也重提埃及的歷史.
縱然基建承認以往耶和華與以色列人的關係.
也承認以往耶和華拯救過以色列.
不過他今天質疑.
他質疑今天是不是這樣.
所以他說.
他那奇妙的作為在哪裡呢?.
他不相信.
正正因為基建認為.
耶和華現在不是與以色列同在.
所以當耶和華與我們同在.
他就作出了一個信仰的回應.
他說:耶和華丟棄了我們.
把以色列交在米顛人的手掌中.
這段對話反映了基建對現狀的不滿.
因為以色列人已經受了米顛人七年的欺壓.
以至基建要躲在榨酒池裡打麥子.
他不敢在空曠的地方這樣做.
因為他很怕被米顛人看到.
搶了他的收成.
基建充滿了埋怨.
他認為如果真是保守以色列人的話.
為什麼任由以色列被米顛人欺壓了七年.
而不施行拯救呢?.
這是基建心裡的那種埋怨.

$^{161}$不過諷刺的是.
在六章八至十節.
先知也曾經譴責以色列人.
同樣先知也提到埃及這段歷史.
但先知最後的總結就跟基廉不一樣.
先知怎樣說呢?.
「以色列的神如此說:.
我曾領你們從埃及上來.
從慧奴之家出來.
救你們脫離埃及人的手.
脫離一切欺壓你們的人之手.
我從你們面前趕出他們.
把他們的地賜給你們.
然後說:我是耶和華你們的神.
你們住在阿摩尼人的地.
不可敬畏他們的神明.
但你們卻不聽從我的話」.
同樣的一段埃及拯救的歷史.
先知說經歷過耶和華的拯救.
以色列人仍然不聽從耶和華的話.
你去敬拜別神.
而基廉呢?.
基廉只專注於耶和華不與以色列人同在.
耶和華不施行拯救的結果.
基廉完全沒有理會到最終的原因.
就是以色列人沒有聽從耶和華的命令.
這個重要的因素.
同一段相同的歷史背景.
為何基廉看出來的東西.
與先知所說的有這麼大的偏差呢?.
這種信仰的偏差.
很值得我們深切反省.
原來這種偏差.
是基廉從結果去看的信仰.
基廉認為.
如果耶和華真的與以色列同在.
就不會有壞事發生.
所以現在以色列人被米甸人欺壓.
這就表示耶和華沒有與以色列人同在.
這是基廉的想法.

$^{201}$這也是基廉信仰上的偏差.
基廉需要修正.
以色列遭遇到的.
無論是好還是壞.
耶和華都可以與以色列人同在.
耶和華是否與以色列人同在.
關鍵並不是以色列的好與壞.
縱然他們不聽神的命令.
敬拜別神.
耶和華是否與以色列人同在.
關鍵在於耶和華是否願意.
而不是以色列的行為.
只要耶和華願意.
耶和華隨時隨地都可以與以色列人同在.
另一個信仰偏差就是.
基廉認為就算以色列人犯罪.
犯了多大的罪.
也好.
耶和華也有責任即時拯救他們.
因此既然耶和華七年都沒有出手拯救.
就表示耶和華沒有與以色列人同在.
基廉也需要修正這種信仰偏差.
耶和華拯救以色列並不是出於責任.
是出於耶和華的憐憫.
是一種恩典.
當以色列人不聽從神的命令.
去敬拜別神的時候.
是以色列人主動破壞了.
他與耶和華的關係.
是以色列人破壞了.
他與神所立的約.
耶和華沒有一個必然的責任.
向以色列人拯救.
因此一切對以色列民的拯救.
都是基於耶和華的憐憫和恩典.
無論在舍斯基.
還是在舊約的其他經卷.
耶和華每一次對以色列人作出拯救.
都不是出於責任.
而是出於耶和華的恩典和憐憫.

$^{241}$其實基廉一開始與耶和華的對話裡.
已經顯示出一種認知上的偏差.
可能大家沒有留意到.
耶和華一開始的時候.
在第十二節.
大能的勇士.
耶和華與你同在.
基廉在第十三節的回應是.
主啊 請容許我說.
耶和華若與我們同在.
耶和華只是說與你基廉同在.
而基廉就認為.
耶和華是與我們以色列同在.
耶和華指的是個人.
基廉指的是民族.
基廉認為耶和華的同在.
必然是與以色列整個群體.
但是耶和華的同在可以選擇個人.
而不是選擇群體.
當耶和華向基廉說與你同在.
就讓我們明白.
耶和華不一定與他的眾子民同在.
但是耶和華確實是與基廉同在.
即使當時整體的以色列子民.
都不聽從神的命令.
但是耶和華可以選擇任何一個子民.
去建立一段特殊的關係.
我嘗試用一個例子去說明.
主耶穌是教會的頭.
主耶穌必定與教會同在.
但是如果教會所走的路.
違背了主耶穌的心意.
而且越走越遠.
教會不再遵守主耶穌的吩咐.
主耶穌的確可以在一段時間內.
不與整體教會同在.
但是主耶穌仍然可以在這間半逆的教會裡.
去揀選他的信徒同在.
並且協助這間教會.
走回主耶穌的心意當中.

$^{281}$因此弟兄姊妹.
如果你今天對教會很失望.
你對教會很心灰意冷.
你認為教會偏離了神的心意.
甚至教會離經半道.
請你不要離開教會.
請你繼續為教會守望.
因為主耶穌可能會呼召你.
藉著你帶領教會走回主耶穌的心意當中.
我們要深信.
教會無論面對順境還是逆境.
主耶穌仍然有他的主權.
主耶穌可以選擇與整個教會群體同在.
又同時可以與個別的信徒同在.
洪英嘉的弟兄姊妹.
我們對信仰有沒有偏差.
當我們正在患難中的時候.
我們有沒有埋怨神.
為什麼不出手幫自己.
我們有沒有反省過自己.
可能有機會自己一直都在得罪神.
但自己卻從來都不知道.
因此從來都沒有向神認罪悔改.
我們有沒有另一種信仰偏差.
就是當我們生活得好的時候.
我們就認為是神的保守.
我們生活得不好的時候.
就是認為神沒有保守我們.
神的保守和看顧.
不是按照我們的物質生活.
覺得好與不好作為一個衡量標準.
而是按照我們與主耶穌的關係.
我們的屬靈生命是否成長來衡量.
弟兄姊妹.
這種信仰偏差.
會導致我們屬靈生命不會成長.
而生命不成長.
就導致我們不再依靠神.
我們不再依靠神.
就會導致我們不敢對神委身.

$^{321}$最後就像舌斯基最後的結論.
跌入了一個屬靈生命倒退的情況.
甚願洪英嘉的弟兄姊妹.
不要跌入這個信仰偏差.
導致我們對神產生埋怨和不滿.
基甸生命的第二個反省就是.
在軟弱裡面.
我們可以接受神的呼召.
如果現在形容基甸的屬靈剛警.
我會這樣說.
基甸對信仰極度的偏差和埋怨.
它與耶和華的關係.
就像潛水艇一樣.
潛入了深海.
很長時間也浮不上來.
不過縱然如此.
耶和華對基甸的憐憫.
也是恩典滿滿.
耶和華對基甸的負面回應.
他如何處理呢?.
第十四節.
去!靠著你這能力拯救以色列脫離米顛人的手掌.
我豈不是已經差遣了你嗎?.
原來耶和華在呼召基甸.
好!既然基甸這麼為以色列人著想.
現在以色列正受米顛人的壓迫.
就靠你現在的能力.
去拯救你的民族脫離米顛人的欺責吧.
我們留意一下.
耶和華這句說話.
他沒有特別加添基甸的能力.
他只是說靠著你這能力.
拯救以色列脫離米顛人的手掌.
耶和華想幫助基甸修正他以往信仰的偏差.
基甸的偏差就是以結果去看待耶和華.
耶和華想基甸專注在耶和華身上.
就是差遣者.
耶和華才是這一場對米顛人戰爭的勝利關鍵.
而被差遣的基甸的能力.
一點都不重要.

$^{361}$耶和華對基甸呼召的重點.
就是要基甸不要再以人為中心.
因為人總是將自己的能力和缺點作出一個比較.
從而認為自己是否有這個資格.
去擔任一些重要的角色呢?.
因此耶和華沒有跟基甸說.
加級你的能力向前行.
他沒有說加級你任何能力.
他只說去靠著你現在這個能力去拯救以色列.
縱然耶和華對基甸一片苦心.
恩典滿滿.
但這一刻基甸仍然在軟弱之中.
很軟弱.
基甸在十五節如何回應呢?.
主啊請容許我說.
我怎能拯救以色列呢?.
看啊我在馬來西亞支派中是最貧悍的.
我在我父家是最微小的.
基甸認為自己不配.
沒有能力.
沒有資格.
我哪夠班呢?.
基甸的回應我們聽起來一點也不陌生.
因為我們一定記得.
在耶和華呼召摩西的時候.
摩西在《出阿吉記》三章十一節的時候.
同樣說過類似的話.
我是什麼人?.
竟能去見法老.
把以色列人從埃及領出來呢?.
摩西同樣認為自己沒有能力沒有資格.
耶和華當時如何回應摩西呢?.
神說我必與你同在.
這句與你同在其實就是.
耶和華一見基甸的時候已經跟他說了.
不過當時基甸是什麼心情呢?.
火遮眼.
以解錯誤.
因此耶和華在十六節再說一次.
加強基甸的信心.

$^{401}$我與你同在.
你就必擊敗米顛如擊打一個人.
基甸兩次聽到耶和華跟他說.
我與你同在.
他開始建立一點信心.
他要求一個證據.
這個證據就是.
耶和華要等他.
等他帶回一些祭物.
親自向耶和華獻祭.
基甸的獻祭甚具象徵意義.
為什麼呢?.
因為獻祭代表基甸願意在神面前認罪悔改.
承認他對神的偏差.
承認他對神的看法錯誤.
也肯定了以往先知責備以色列的原因.
沒有聽耶和華的話.
因此他作了一個具體的認罪行動.
向神獻祭.
基甸要求耶和華提供證據.
其實不止一次.
在往後的日子.
他不斷要求耶和華提供證據.
加強他的信心.
如果再看下去.
包括了兩次楊無的神蹟.
一次要求楊無濕身.
另一次要求楊無乾身.
神對基甸很有恩典.
每一次都滿足他.
神不斷與基甸同在.
也讓基甸的信心一路一路增加.
最後基甸帶領以色列人打了兩場很重要的戰事.
我們一定會很有印象.
以三百個以色列士兵.
去打贏十三萬五千個米顛人.
三百對十三萬五千.
正常來說這種叫以卵擊石.
但是神與基甸同在.
縱然只有三百.

$^{441}$有耶和華的幫手.
十三萬五千人.
完完全全被擊倒.
因此基甸能夠完成耶和華對他的呼召.
就是將以色列人拯救出來.
基甸在心靈狀況最軟弱的時候.
卻可以接受耶和華的呼召.
這的確成為了一個很好的反省.
只要神願意.
神可以使用任何一個人.
只要神願意.
哪怕那個人的屬靈光景如此不濟.
被召的那個人.
不需要輕看自己.
不需要怕自己誤會.
不需要覺得自己能力不足.
被召的那個人.
也不需要諸多藉口去推搪.
因為最後你都推搪不到.
被召者只需要清清楚楚地知道.
這個呼召是從神而來的.
被召者只需要清楚地知道.
神與他同在就足夠了.
基甸被召這個生命片段.
讓我們清楚知道.
使命能不能達成.
不是取決於我們的身份或能力.
使命能不能完成.
完全取決於神的大能.
以及神與我們的關係.
紅映嘉弟兄姊妹.
今天的講題是從軟弱到剛強.
基甸這段生命的經歷.
給了我們一個很寶貴的教訓.
基甸的軟弱是因為信仰的偏差.
不過他願意認罪悔改.
而基甸能夠剛強起來.
帶領著以色列人打敗米天人.
關鍵的是基甸能夠認清一件事.
呼召者是誰.

$^{481}$呼召者才是能夠令基甸剛強起來的關鍵點.
呼召者既是萬君的耶和華.
他有大能力.
基甸就無需擔心自己的能力.
因為他的能力是從耶和華而來.
呼召者既是掌管萬有的主.
他有極大的主權.
基甸只需要信服主耶穌的大靈就足夠了.
呼召者既是說有就有的創造者.
他的話語充滿智慧.
基甸只需要謙卑聆聽就足夠了.
呼召者既然願意與基甸同行.
那麼基甸又何懼之有呢.
因此從基甸這一段生命的歷程.
讓我們學習到一件事.
我們縱然在軟弱當中.
我們仍然因著我們所相信的.
可以變得剛強.
昔日耶和華從埃及地將以色列民拯救出來.
成為以色列民族的拯救主.
昔日主耶穌亦在各國他山.
在十字架上用寶血將我們從罪惡中拯救出來.
成為我們這群愛邦信徒的救主.
我們能夠剛強的原因不在於我們.
而在於我們所相信的主耶穌.
只要主耶穌願意.
只要主耶穌與我們每一個同在.
只要主耶穌向我們發出呼召.
我們縱然身處在軟弱當中.
我們仍然可以變得堅強領受召命.
承擔主所交託的召命.
弟兄姊妹.
在疫情底下.
我們每一個都懸弱.
我們每一個都擔心.
我們每一個都憂慮.
我們不知道這場疫情何時完結.
我們都不知道自己會否成為感染者.
其實我們生活在一個惶恐不安當中.
我們與基甸一樣.

$^{521}$基甸要在酒渣池裡面打墨紙.
因為他心裡面恐懼.
他怕米顛人發現後.
搶走他的修成.
基甸生活在一個惶恐不安當中.
我們今天也是.
但是基甸願意重新認識信仰.
改變了他對信仰的偏差.
他向耶穌說認罪,獻祭.
並且再次重新領受從神而來的呼召.
他就不再害怕.
基甸憑著與神同在的信心和能力.
最終完成神給他的使命.
今天只要我們能夠認定.
呼召我們的是誰.
並且深信主耶穌與我們同在.
我們就能夠安然產生一種信賴,倚靠的心.
主與我們同在的平安.
使我們信仰的生命可以從軟弱化為剛強.
並且能夠承擔使命,被神使用.
同樣的弟兄姊妹.
我們願不願意生命有這種轉變.
我們願不願意被神所使用.
我們是否想每天都經歷一種與主同在的恩典呢?.
我們願不願意.
我們願不願意被神所使用.
一種與主同在的恩典呢?.
我深信.
我們每一位基督的跟隨者.
都會從心底裡回應.
我願意.
差遣我吧.
我們有一個祈禱.
祈求與你同在的今天.
這一刻.
求主你兼顧我們每一個的心.
讓你的話語栽種在我們的心裡.
並且讓聖靈親自的囂觀.
按著你的心意.
一步一步地讓我們的屬靈生命能夠茁壯地成長.

$^{561}$並且甘心樂意地被你使用.
終一生去高舉主耶穌基督的聖名.
我們的祈禱是奉主耶穌基督的聖名祈求.
阿們.
願至祝福大家.
\newpage



\section{約伯記 17:15}
\label{sec:7bIUWEs7F1E}
\textbf{2020.09.27《HOPE ̶l̶e̶s̶s̶》 約伯記 17\_15 朱子溢傳道}
\newline
\newline
連結: \href{https://youtube.com/watch?v=7bIUWEs7F1E}{\texttt{https://youtube.com/watch?v=7bIUWEs7F1E}} ~~~~ 語音日期: 2020-09-30
\newline
\newline
\hyperref[sec:RqJiqD22j9E]{\small{< < < PREV SERMON < < <}}
~
\hyperref[sec:index_chronic]{\small{[返順時目]}}
~
\hyperref[sec:index_scriptual]{\small{[返順卷目]}}
~
\hyperref[sec:6c0M0guWu_4]{\small{> > > NEXT SERMON > > >}}
\newline
\newline
約伯記 17:15
\newline
\begin{longtable}{cl}
\hline
\hline
章節 & 經文 (和合本修訂版)\\
\hline
17:15 & \begin{tabularx}{0.7\textwidth}{X} 這樣，我的盼望在哪裡呢？我所盼望的，誰能看見呢？ \end{tabularx} \\ \\
[1ex]
\hline
\hline
\end{longtable}
$^{1}$各位大英姐妹平安.
今日來到洪恩家的時候.
很多熟悉的面孔.
有一份戀意.
當這麼多年過了.
彼此很認識的時候.
這種關係是很難得的.
今天教會邀請我去講這個講題的時候.
給了一個主題.
是生與死.
確實是的.
面對生死的時候.
在我的個人經歷裡.
或者在信仰的反省裡.
或者在服事裡的實踐裡.
都看到很多時候.
我們面對死亡或者不計生.
老病死的過程裡.
在一個階段的時候.
很容易萌生一個疑問.
究竟為什麼會發生在我身上.
為什麼會發生在我們這個時代身上.
為什麼會發生在現在呢.
無論是什麼時間發生.
你都會覺得很早.
來得很突然.
所以都發現很多的.
面對死亡的時候.
或者生死之間的時候.
總是會有一種.
好像找不到那種希望.
好像缺乏了那份盼望.
好像有一絲絲的絕望.
萌生在意念裡.
所以今天我的主題.
確實是希望和絕望.
因為希望和絕望.
確實是一線之間.
弟兄姊妹.
我們今天都問很多的問題.

$^{41}$無論我們戴著口罩.
無論我們這一年所經歷的.
無論我們家庭的時候.
面對不同的挑戰.
不同的困難.
你都可能問.
為什麼會發生在現在這一刻.
為什麼我要在這個時候去面對.
當我去看到聖經裡面.
究竟是不是現在這個時代.
是最艱難的時代呢.
我們看到很多的.
聖經裡面的人物.
他都面對不同大小的艱難.
有先知面對他那個世代的人.
那個盟瑟.
那種失望.
那種失意.
那個環境讓他們那種灰心.
又或者在我們一個很重要的詩人.
如果你有聖經.
你可以打開詩篇的69篇.
詩篇69篇.
有一位詩人.
大衛他面對一個困難的時候.
他說他的狀態.
就好像深陷於淤泥裡面.
沒有立腳之地.
他去到深水之中.
那個瓶子比他的身體慢.
他因為呼求去告困乏.
他連喉嚨都叫得乾了.
他的狀態就是一個很沉溺的狀態.
好像是在說.
當我們面對很多的受煩.
困難.
挑戰的時候.
我們真的好像窒息一樣.
被那個環境讓我們.
面對的困難真的很大.

$^{81}$可能是病患.
可能是我們環境的處境.
我們覺得呼吸不來.
大衛就是面對這樣的情景.
當我們說無辜恨我的.
是沒有原因恨我的.
竟然比我的頭髮還要多.
本來我們清清道地的.
可以拍拍肩膀的.
可以傍著的.
竟然今天當我不認識的.
當我是陌路人.
我已經用盡我的辦法.
用麻布去當衣裳.
但是其他人怎樣呢.
取笑他.
很多時候我們.
有時候我們面對不同的艱難挑戰的時候.
我們用我們的信仰原則.
我們堅持自己的信仰原則.
可能在你和家人的相處裡面.
可能在職場裡面.
可能照顧孩子.
教導他們的過程裡面.
你會發現.
這樣做不是最好的嗎.
但是在我們的經驗裡面.
可能大家也會同意.
或者有相同的經歷.
就是被人取笑.
甚至質疑我們這個信仰.
或者這個決定.
其實你是不是真的.
你憑什麼.
每個人都這樣做的.
用不著這樣吧.
很多時候我們都面對這種艱難.
是吧.
縱使大衛我們認識他.
是一個很有信心的人.

$^{121}$和神很親近.
但是他的處境下來的時候.
他也是人.
也同樣是面對這種艱難.
甚至乎他說.
你以為我聽不到嗎.
你根本坐在城門口討論我.
用走圖.
用歌詞 歌曲.
去取笑他.
以他為靶 笑病.
各位.
這種人生的經歷.
走到這樣的時候.
或許也會令人萌生一種絕望的態度.
同樣.
今天主席為我們讀出的一段經文裡面.
約伯我們很認識.
很正直 很二人.
但是.
他的日子.
他說自己這樣去.
他好像向著一個空氣.
或者一個不知道向著什麼去說.
他說我的日子已經過去了.
謀略 我心的願望都斷絕了.
因為他背後經歷的就是.
家散人亡.
就是一夜之間 一日之間.
一切都沒有了.
我們說有時候死亡.
或者我們去想像死亡.
因為我們沒有人經歷過死亡.
今天在座的全部都是有活著的.
所以當面對死亡的時候.
我們說其實是一種終極的失去.
當然我們知道我們將來有盼望.
但是我們都知道.
在地上的生命是要完結.
所以這個終極的失去.

$^{161}$讓我們想到其實我們.
步向生命的終結的時候.
我們不停經歷著.
擁有但是慢慢失去.
生命是一樣的.
財富是一樣的.
或者你的健康.
或者你和人的關係.
或者你就是.
基本上我們都是在學習.
或者經歷著不停地失去.
當你孩子的時候成長.
他開始學習擁有.
但是他也學習經歷失去.
年輕人最多就是.
失戀是他們最大的問題.
擁有但是又要失去.
所以那個時候很多人都是.
可能有一個這樣的.
就是尋想就是說.
那我的盼望還可以在哪裡.
我所盼望誰可以看見.
如果我們說我們的.
信仰的答案就是面對絕望.
就沒有答案的話.
那我們真的不能夠再.
繼續說我們可以怎樣.
但是信仰就是在基督信仰裡面.
耶穌來就是告訴我們.
祂是帶著這份活著的盼望給我們.
所以今天我很想和大家分享.
就是其實當我們去面對.
一個艱難的處境.
甚至乎可能是一個.
面對終有一些不同的情況.
失去的時候.
我們還可以怎樣去.
堅守著這份盼望呢.
在信仰裡面從四個向度.
四個角度或者四個點子.

$^{201}$去到反省.
我們今天我們的處境裡面.
信仰可以成為我們什麼的盼望呢.
怎樣我們可以在絕望裡面.
好像這個小小的chicky位.
畫一畫就成為我們的盼望呢.
而怎樣可以在現在這個世代裡面.
我們在一個絕望.
很容易跌入絕望的環境裡面.
我們可以帶出這個盼望的訊息.
我們可以活出這個活潑的見證.
我們一起有個禱告開始.
主系今天我們來到你的面前.
領受你對我們的說話.
求聖靈今天在我們當中.
去教導我們.
分娩我們.
提醒我們.
還有鼓勵我們.
求天父去與我們同在.
今天我們要看見.
你親自在我們當中.
我們又要看見你的話語.
成為我們的幫助.
我們這樣禱告一望.
奉賴耶穌基督得勝靈自求.
阿們.
第一點呢.
我很想和大家分享.
是我們的天父.
當我們面對.
其實如果你沒有擁有過.
你就不用面對失去.
但是我們很多時候擁有的時候.
面對失去是很艱難.
我們在我們的成長過程裡面.
我們可以慢慢去學習.
但是有一個.
有些情況是.
我們突然失去了原本擁有的.

$^{241}$我們就好像很慌亂.
因為失去了倚靠.
但是信仰裡面.
其中一樣就是.
耶穌很清楚地對門徒說.
你要去禱告的時候.
要怎樣禱告呢.
就是要說.
父啊.
願人都尊你的名為聖.
願你的國降臨.
願你的旨意成就在地.
我們知道.
主禱文教導我們的.
就是在天上的父.
有聖經學者就說.
整個主禱文裡面.
三願三求裡面.
最重要的一點.
反而不是三個願三個求.
而是第一個的稱謂.
當我們說.
我們向神祈禱的時候.
我們不是單單向做我們的位置.
不是單單向創造天地的位置.
更加我們要向天上的父去禱告.
而在聖經裡面.
很清楚.
其實也有原文說.
連天上的都可以生去.
父是直接的.
因為我們的主.
就是我們的爸爸.
當我們最直接叫的時候.
爸爸啊.
就在說我們和神的關係.
是充滿著這種愛的關係.
充滿著接納的關係.
充滿著會有包容.
會有教導.

$^{281}$可以在當中得到保護.
有倚靠.
這種父與兒女的關係.
是相當直接的.
也成為我們面對任何的艱難.
我們都有天上的爸爸.
來幫著我們.
好像箭肉一樣.
我們掉下去.
兜著我們.
在我的事奉裡面.
我經歷過一次.
很深刻的一個經歷.
當我想到如果面對不同的艱難.
或者處境裡面.
失卻了父母的愛.
究竟是怎樣的呢?.
我在醫院服務的.
也感受到醫院發生的事.
就好像社會的小縮影.
社會發生了什麼事.
在醫院裡面就好像都呈現了.
有一段時間我在聯合醫院.
負責青少年和兒童的病房.
大概三年前.
特別是一開學.
我不知道大家有沒有印象.
有些新聞.
不停地報導.
其實我們一年.
很多人都面對很多的艱難.
所以其實一年有五六百.
甚至乎去到七八百人.
選擇用一些輕生的方法.
去面對自己的生命.
或者想要解決的問題.
但是三年前特別.
多了一些報導.
就是很年輕的少年人.
他們就好像每一天都有一宗.

$^{321}$那段時間醫院就會.
接收到很多救得回來的病人.
或者一些少年人可以關心他們.
很快其實半年左右.
好像過去了的消息.
但是在兩年多前.
大概是年底的時候.
我不知道大家有沒有印象.
有什麼那時候很轟動的新聞.
關於一些家庭.
就是在屯門.
很近我們.
屯門有一個小朋友.
被人虐兒離開了世界.
很震撼我們.
覺得這個世界為什麼會這樣.
這個社會發生什麼事.
那段時間因為大家很關注.
所以那些保安.
或者怎麼樣.
一見到撞雨了就報警了.
就送他進醫院.
那段時間在醫院.
就有一個我們有這樣的叫法.
叫做金魚缸.
知道什麼叫金魚缸嗎.
就是那些魚不停游來游去.
你走不開的.
為什麼叫金魚缸.
因為平時不會這樣走來走去.
那段時間很多小朋友.
他們不是很重病.
二來就是家長不在身邊.
因為沒有監護人.
我們一般在醫院裡.
其實醫護就是醫治.
照顧的都是監護人.
但是那段時間因為他們是監護人.
更加不可以陪著他們.
就是要分開他們.

$^{361}$有一個情景就是發生在那一次的經驗.
就是我去探望一個病床邊的小朋友的時候.
當然試探望他的家長.
在病床上已經有一個十幾歲的少女.
談了兩句的時候.
我就很不好意思地問那個媽媽.
我就跟她說.
其實你躺在你的膝蓋的那個.
是誰啊.
為什麼我這樣問呢.
因為在我的畫面上就看到.
大概兩歲多.
比我小兒子還大一點.
兩歲多的一個小女孩.
躺在大腿上躺得很舒服.
那家長怎麼回答我呢.
為什麼我這樣問呢.
因為很少有兩個姐妹一起在同一個病房.
但是這個畫面好像很特別.
那家長說其實他也不知道是誰.
只不過是隔壁床跟他的女兒玩了一會.
他覺得很累就躺在他的大腿上.
他也不知道怎麼辦.
最後我們就抱了他回床.
讓他睡得舒服一點.
我不知道你有沒有這一刻.
剛才這一刻我說完之後.
你會不會想起那個畫面.
其實那時候有一兩秒.
我覺得那個環境很溫暖.
一個小妹妹躺得很舒服.
其實你看到也覺得很心甜.
但是這個只是一兩秒.
過了一兩秒之後.
我的心就很緊張.
為什麼這樣說呢.
我覺得一個小女孩兩歲多.
面對一個陌生的人.
他不理會應不應該可以這麼舒服地躺在他身上.
我心裡就想.

$^{401}$這個小朋友應該很大機會.
就是因為他成長裡缺乏了父母的愛.
當一個陌生的人.
稍微給他一點關心關愛.
他已經可以完全依賴在他身上.
其實代表這個小朋友很缺乏安全感.
各位.
很想跟大家分享的是.
原來一個人在他成長裡.
如果缺乏了父母的愛或是一個關係的時候.
其實是不容易的.
但是我們的信仰就告訴他.
無論任何的艱難都好.
我們的主就是我們天上的爸爸.
他從來不會離開我們.
我們人是很有限的.
今天我自己作為爸爸.
我都知道我所謂的愛都很有限.
我都以為自己.
當我大女兒出世抱著他的時候.
我覺得我這一生人.
可以為他奉獻一切.
就是犧牲所有.
但是當他成長的過程裡.
我都知道自己因為自己的限制.
我不能夠真的和我的承諾完全兌現.
但是我們的主不是.
所以今天我都教我的孩子.
他說最重要的是.
有天上的爸爸的保護帶領.
面前有什麼困難都好.
我們就可以容易去面對.
我們有些教會.
經常討論到青少年的服侍或是模樣.
有些教會是這樣說的.
就是代父式的模樣.
因為這個世代裡.
我們知道其實家庭的關係都不容易經營.
很多時候因為不同的緣故.
年青人在成長裡好像缺乏了父母.

$^{441}$所以教會的回應就是.
既然我們的主是天上的爸爸.
我們今天都可以成為代父.
貼身去回應年青人的.
不是訴求而是生命成長的需要.
以至在當中他們明白什麼是天父的愛.
各位.
我們在我們的生命的階段裡面.
當我們面對很大的艱難的時候.
在醫院裡面見到很多臨終的病人的時候.
很多時候他們都好像回到那些福圖畫那樣.
孩子一般.
有那種單純的時候.
就可以找到尋.
第二點是「O」.
「O」是活生生的意思.
有機的意思.
因為耶穌曾經說過.
我們的神是什麼神呢?.
是亞伯拉罕,以撒和雅各的神.
祂從來不是死人的神.
是活人的神.
各位姐妹.
我們的主是那位昨日,今天.
直到永遠都不變的主.
所謂不變是說.
不是祂不動.
不是祂不做事.
不是祂不保護,不保守.
不是今天顯神蹟.
而是祂不變是昔日都顯神蹟.
我們相信祂今天仍然參與在我們的生命裡面.
昔日很多人去到.
去到呼求神的時候.
神回應他.
那些就好像雲彩的見證.
告訴我們今天我們去呼求.
我們仍然在不同的艱難裡面.
神都會回應我們.
有一次我的女兒.

$^{481}$第二個女兒.
大概四五歲的時候.
她想起一個玩具.
在櫃頂的箱子裡.
她就用盡辦法想拿到櫃頂的箱子.
她無論是搬椅子.
嘗試爬一下.
她都沒有辦法拿到那個玩具.
有時候小朋友是這樣的.
她想起那個東西就很想拿.
很想即時拿.
她就哭了.
開始不滿.
不滿完她就開始覺得很絕望.
很無望.
根本拿不到.
她就來爸爸的面前跟我說.
我拿不到啊,爸爸.
我就嘗試去引導她想想方法.
其實你用盡方法.
你還有一個方法還沒出.
我真的想不到.
沒有方法了.
我就跟她說.
你可以找人幫你的.
她就靈機一觸.
是喔,爸爸你幫我拿行不行.
我真的不用說就拿了那個箱子下來.
她就玩得很開心.
各位,其實呢.
作為爸爸的.
我又要給她嘗試.
但其實我已經預計了幫她拿.
不過在過程裡.
很想等她去問我.
因為我也很想嘗試完之後.
去找到一些辦法.
這種等.
我又不是真的在等.
在看著她.

$^{521}$正如在路加福音十五章裡.
失很多的比喻.
其中一個就是失.
失兒子的比喻.
失楊的比喻.
失錢的比喻.
失兒子的比喻裡面說.
那個慈父就在等待浪子回頭.
這種等待在粵文裡面是一種語法.
是在不停地等.
不是一個被動的在等.
一個不停地等.
為甚麼這樣說.
因為那種等就好像.
我剛才泊車過來要一點時間.
我相信老全都在等我.
他不是真的在等.
他在裝裝看我來不來.
所以當慈父在等.
就是裝裝看.
看看他在不在.
才可以當浪子來的時候.
他的小兒子來的時候.
他第一時間給那些僕人.
其他人第一時間出去的原因.
就是因為他不停地在等.
他沒有停過他在等.
我們的主就是.
在我們生命的不同困難的時候.
他能夠被我們.
去倚靠他.
有一個宣教士.
去到一些土著裡面.
去翻譯聖經.
但是他們土著的土話裡面.
沒有了一個字.
就是倚靠.
他不知道怎麼翻譯.
他就叫了那個族長過來.
你試試靠在那張木椅上.

$^{561}$他就坐在那裡.
他坐在那裡的時候.
宣教士就跟他說.
你提高你的腳.
你的腳不要碰地.
你就自然整個躺在那張椅上.
他就問.
這是什麼話來的.
那個族長就.
就這樣說了一些說話.
這個宣教士就用這個.
說話來去.
去翻譯成倚靠.
就是當我們的生命就放在這張椅上.
就好像我們.
去倚靠神.
我們將生命完全給了神.
我們相信祂.
會承托我們生命的重量.
所以這個就是倚靠.
各位.
這個倚靠背後是需要我們.
這份信心的.
需要我們知道神仍然在.
需要知道.
他昔日是做神蹟.
今天就像這幅圖畫.
都是可以.
今天都會顯出神蹟.
第三個.
點子是.
我們生命的目的.
我們生命的目標.
當有人去問.
耶穌的時候.
這個人為什麼會失明.
肯定是不是他犯罪.
或者他的父母犯罪.
耶穌回答.
不是這個人犯罪.

$^{601}$不是他的父母.
而是要在他的生命上.
顯出神的作為.
很多人面對很多的.
生命的挑戰困難病患的時候.
都問為什麼.
改變.
或者我們不想遇到的艱難的時候.
我們都會問為什麼會發生在.
我們生命上的.
不是愛我們的嗎.
不是給我們最好的嗎.
為什麼要面對逆境困難.
想的嗎?不想的,沒有人想的.
今天我們都不會說.
擁抱一些的挑戰和困難.
我們都在這裡禱告神.
如果今天可以.
除去所有的.
不是除去所有的口罩.
如果社會的狀態讓我們真的可以.
除下口罩走出去就好了.
都不是很高的要求.
因為沒有人想面對.
一些這樣的.
面對一些困難.
我不想的.
同樣.
我想很多人.
面對不同的.
生命的階段裡面.
都要去面對這種.
等待著的.
困難過去.
我都不少時間.
帶不同的大學生.
去醫院做服務.
有一次在中大.
去教這個program.
通常都完了一個program.

$^{641}$我就會跟他們說.
名額有限,要給機會給其他人.
你可以推薦其他人.
下次就讓給其他人.
給機會其他人.
他們通常明白就不會再報名.
中大有一個學生女孩子.
20歲左右.
去報名.
第一次我帶著她.
去做服務的時候.
她都給我一些印象.
但不是太深.
每一次她的分享都有眼淚.
或者很感性.
但都不是很了解她的情況.
但第二次她報名的時候.
那我就知道了.
因為她要說多一點自己報名.
她就是.
看到.
第二次的服務.
去聯合醫院.
兒童病房.
她就說我真的要去.
因為她原來.
她第一個學期.
去做服務的原因.
她就告訴我.
就因為她小時候住過聯合醫院.
在她小五的時候.
她和她媽媽.
一起送去聯合醫院.
分開隔離.
原因她就住在淘大.
她就經歷了.
那時候的沙士.
她說她小五.
她說她不知道媽媽怎樣.
她也不知道自己會怎樣.

$^{681}$她就經歷了幾個月.
在醫院裡面的.
捱或者很絕望.
那時候都感恩.
事實上真的很感恩.
就是她除了護士.
有院務去探她.
以致她那時候.
認識了神.
在教會成長.
為何她會加入服務呢.
就算她和她媽媽.
最後出院了.
沒事了,可以治療了.
但是因為這個陰影.
在一個十歲的小女孩裡面.
她生死未卜.
就被隔離了.
那個陰影令她.
十幾年來.
她又沒有睡得好.
經過醫院都在顫抖.
她說她生命真的不行.
到她二十歲的時候.
不能讓這個陰影纏繞十年以上.
她已經問了很多問題.
為何神要這樣.
為何我要面對這些.
為何我仍然在.
好像在沉溺的精神狀態裡面.
她看到服務.
她就說.
我試試用另一個身份試試.
所以原來她第一個學期.
是很挑戰她的.
她就覺得我真的要.
試試吧.
誰知她成功.
她可以進醫院服務.
以至她第二次服務.

$^{721}$她知道是服務.
聯合醫院的小朋友.
她就想起自己的經歷.
並且完了第二次服務之後.
她立志要去.
第二個學位.
去讀醫護.
其實是很大的改變.
我為這個年青人很感恩.
我為她的生命的改變很感恩.
但同時間.
我稍為想像她.
過往的十年.
我今天看到她的改變.
但那十年期間.
我相信她沒有經歷.
很多人的改變.
其實很不容易.
3650天.
每一天.
都這麼艱難.
但當她經歷了.
這十年.
當她知道原來.
她可以說.
我明白為何我要遇到這個病.
以致她改變了.
她知道神給她的旨意.
就是要服侍.
有病的人.
以她的生命經歷去服侍.
這群人.
我知道原來最難的時候.
是等待的功課.
但當經過回頭望的時候.
她的信心.
就因為被火屋煉.
以致.
她能夠再面對.
任何艱難都好.

$^{761}$她的生命更加寶貴.
所以.
等待是不容易.
但我們知道.
有時回頭望.
我們才明白今天的處境.
帶給我們生命的甚麼目的.
能夠在將來的時候.
我們回望.
自己聽,告訴神.
是的.
你的旨意.
比我們所想.
所求.
我們的最後一點.
我們的主.
是意馬內地的主.
因為.
祂與我們同在.
在約書亞記.
當約書亞接了一個.
超越她能力的棒的時候.
她說她不行.
很艱難.
可能在工作裡.
在事奉裡.
在家庭裡.
我們有時也會面對.
一些責任是我們承擔不起.
但是.
神應許.
幾次告訴她.
我不是吩咐過你嗎.
你應該剛強壯膽.
不需要懼怕,不需要驚慌.
因為你無論往哪裡去.
耶和華的神必與你同在.
因為我們的主.
就是那位與我們同在的主.
祂是意馬內地.

$^{801}$與我們同在的主.
各位.
我們或多或少.
也在經歷一個.
無形的教育.
不知道大家知不知道.
我有時想這個教育.
怎樣稱呼它呢?就叫孤兒教育.
其實香港很少孤兒.
大部分.
都是可能.
現在的孤兒院都轉成兒童院.
因為家庭有困難.
所以就送進去.
為什麼有一個很強的.
孤兒教育在我們的.
潛移文化裡.
已經是盛傳.
可以這樣說.
有一次,我不知道你有沒有遭遇.
就是我在茶餐廳.
吃東西的時候.
聽到隔壁桌子.
一個媽媽對著她的小孩,六歲左右.
小學生.
好吃了你.
可能下次沒得吃.
媽媽可能隨時就走.
很難說.
好吃了.
有沒有聽過?可能自己也說過.
跟自己的小孩.
好做功課.
有得做就好做.
要珍惜.
你想想深一層.
一個小學生六歲.
他為什麼要想像媽媽突然.
發生交通意外呢?.
他突然有危險.

$^{841}$是不是有這個可能?.
我明白的,背後就是想他獨立.
我們說不要讓他依賴很多.
成為港孩怎麼行.
其實背後.
我更加相信.
可能是文化.
史言.
我阿公.
在我成長的過程裡.
他陪了我很多.
他經常都跟我說.
當他十多二十歲的時候在香港.
他經歷什麼呢?.
他經歷的時候.
日本.
入侵香港.
那時候是民安隊.
負責搬運物資.
他說他在這裡.
就好像後面鐘的距離.
那個炸彈就炸下來.
他所有的朋友.
全部身亡.
那一刻令他很震撼.
所以.
他很明白珍惜生命.
和生命的重要.
他經常叮囑我們.
有書讀就好讀.
其實.
我們那一代.
上一代也好.
其實就是經歷苦難.
香港就是一個逃難下來的城市.
其實就是很多的苦難.
以致那一代人.
他認為他的經歷就是這樣.
就要跟下一代去說.
所以我們不是經歷了一兩代.

$^{881}$這一種.
隨時要靠自己.
這個.
文化的.
意識已經.
潛藏在我們的生命裡.
意識或想像裡.
但這個是不是聖經所說.
或我們的信仰所說.
雖然不是.
所以我們在我們的.
生命裡面.
我們好像做得不好.
做得不好的時候神會.
凋棄我們.
但剛才聖經很清楚.
告訴我們.
要與你同在.
我們是與你同在.
因為神就是與我們同在的主.
他不會因為我們怎樣.
而去離棄我們.
這幅圖畫.
就是.
耶穌抱著羊.
牽著小朋友.
有人這樣說的.
不知道有沒有聽過.
中東人養羊.
一群羊走在前面.
總有一些比較頑皮的.
離開了那一條群.
你知道的.
羊仔一離群.
很容易被.
其他猛獸.
擔了他,當然不可以讓他離群.
那怎麼辦.
很簡單.
你養寵物.

$^{921}$你都會懲罰他.
當然.
懲罰之前.
指引一下他.
牧羊人有橡皮筋.
橡皮筋.
橡皮筋就是有勾那邊.
就是勾一勾他的腳.
指引他走回去.
他就走回去.
但總有一些不長進的.
勾他兩下.
他都離群.
那就用另外那邊.
植不粒那邊就是橡皮筋.
打都不用想.
他不是寵物.
他只是羊.
用資產.
但是牧羊人.
打得來,中東人不只是打痛他.
打斷他的腳.
是這樣的嗎.
打斷他的腳.
為什麼要打斷他.
不是打斷他,是丟他在野地.
他也不是像.
這些小羊那樣單手抱.
他就像皮草.
兩隻腳.
擔在肩膀上.
牧羊人擔在.
身體上.
是重的,但是在過程裡.
因為腳服完.
要接駁骨頭.
要時間.
那段時間,這隻羊.
聞到牧羊人的氣味.
以致當他康復.

$^{961}$再走回前面的時候.
他不會離群的原因.
因為他認得牧羊人的味道.
所以他才要用這個方法.
去叫那隻羊.
去面對這麼大的困難.
這麼痛楚.
各位,我們的神就是這樣.
他是與我們同在,但是他不會.
單單讓我們保護.
我們完全不面對困難.
但是他願意承擔.
我們生命的重量.
所以我們就可以.
像這張椅子那樣.
依靠我們的主.
各位,當我們.
所謂面對這些.
艱難的時候,盼望.
成為我們很重要.
讓我們可以堅忍.
在這些困難的裡面.
所以《約伯》裡面其中一個特質.
就是好像.
貼在《加錢書》.
保羅老師所寫的.
信望愛裡面.
望那個堅忍.
堅忍才是.
約伯面對他的處境.
一個最大的生命特質.
而且堅忍.
和背後.
就是要有一份盼望我們的主.
所以我們才.
存在那份忍耐.
有神學家說.
在教會裡面我們宗養.
宣揚就是信.
望和愛,但是.

$^{1001}$最少講卻是.
盼望,我們經常講愛.
經常講信.
但是盼望在我們的生命裡面.
在我們的處境裡面才是.
最能夠回應.
我們現在的生命的狀態.
昔日教會歷史的.
殉道者他們面對很大的.
艱難,他們在鬥獸場裡面.
他們要死了,但是.
為何他們仍然可以手牽手.
他們說遙望這個盼望.
堅忍到底.
或許我們今天的時代會面對.
不同的困難,我們都要有這種殉道者.
這種精神,他們不是.
因為他們人的能力.
而是他們能夠知道.
將來這份盼望.
就是神已經給過我們.
我們直接釘在十字架裡面.
告訴他,面前任何困難都好.
他就是我們的父.
他可以給我們倚靠.
與我們同在,昔日今日都.
不變的主,所以我們.
能夠去燃點.
這個希望的燭光,縱使.
黑暗死亡要來.
但是我們這個生命的燭光.
能夠照在人的面前.
不是放在斗底下.
是可以放在燈台上.
成為這個見證.
今日我們的教會.
大兄姐妹,面對困難.
不能少說.
但是如果我們能夠藉著這個困難.
鍛鍊他們的信心.

$^{1041}$我們就能夠在其他人面前告訴他們.
這就是美好的見證.
因為我們生命裡面,被人看到.
我們背後的那位牧者.
他的幫助,他的指引,他的能力.
我最後.
算了,不要緊.
有幅畫面.
不過.
我按不上去.
那麼.
小小的一個結束.
小的片段,就是有一次.
在嶺南的後山裡面.
嶺南大學裡面.
最多就是貓貓.
因為最高峰的時候.
有100隻野貓.
有些義工.
經常去餵牠們.
很疼牠們.
很多都是貓奴,嶺南的學生.
有一次我看到一個女孩.
拿著相機.
拍山後.
我問她.
是否拍貓.
她說不是,拍甚麼.
她說拍.
毛蟲.
毛蟲?.
我問她為何拍毛蟲.
原來她很喜歡毛蟲.
為何喜歡毛蟲.
她說喜歡毛蟲的原因是.
毛蟲結繭後成為蝴蝶.
她覺得過程.
很有意思.
所以她特別喜歡毛蟲.
我不能評論她的喜好.

$^{1081}$但背後都告訴我.
原來很多人.
心裡都在等待.
一些盼望.
所以我們嶺南有一班年青人.
年提.
為年提代領.
所謂.
在現在的艱難裡.
社會裡面.
最艱難的就是在校園裡面.
好像被人看不見.
但又面對很多困難.
又沒有人幫助.
他們很想去服侍他們.
所以又是待又是服侍.
這個主題.
我都會在想.
原來我們生命裡面.
身邊的人都很多人在等待人領.
所謂等待就是.
等待.
人去領他.
成為.
神的兒女.
所以我想我們身邊的人.
我們都面對很多艱難.
好像顧不上自己.
但我們亦都深信.
當我們經歷困難.
我們能夠幫一些.
等待人領領.
以致他們能夠.
在信仰裡面.
經歷到真實.
我們一起有個祈禱.
七恩主我們多謝你.
因為我們在困難裡面.
我們不會落在絕望裡面.
我們在死刃.

$^{1121}$的幽谷的時候.
你都與我們同在.
你是我們的牧者.
我們可以有盼望.
是因為你的像綱.
成為我們的安慰和指引.
教導.
今日我們大英姐妹.
在這個時代裡面.
或多或少面對不同處境的艱難.
困難.
甚至乎經歷一些失去.
但我們永遠都不會失去.
這個活潑的盼望.
就是從耶穌基督.
你親自去給過我們的時候.
就求你去.
充滿我們.
給我們有信心.
去倚靠你.
去走過面前.
所有的挑戰.
以致我們能夠去到.
好像魚鷹.
展翅重騰.
求你去.
祝福我們.
幫助我們.
帶領我們.
讓洪錦嘉的每個大英姐妹.
牧者在你面前.
去領受你更多的.
祝福和恩典.
以致我們成為.
運彩的見證人.
去到這個社區.
這個地方.
這個時代很多的家庭裡面.
成為這個.
美好的見證.

$^{1161}$成為這個時代的.
崩台.
祝福我們.
我們在你們面前的禱告.
讓望是奉賴耶穌基督.
得勝靈自救.
\newpage



\section{哥林多後書 12:7-10}
\label{sec:6c0M0guWu_4}
\textbf{2020.10.04《軟弱化剛強》 哥林多後書 12\_7-10 陳一華牧師}
\newline
\newline
連結: \href{https://youtube.com/watch?v=6c0M0guWu-4}{\texttt{https://youtube.com/watch?v=6c0M0guWu-4}} ~~~~ 語音日期: 2020-10-06
\newline
\newline
\hyperref[sec:7bIUWEs7F1E]{\small{< < < PREV SERMON < < <}}
~
\hyperref[sec:index_chronic]{\small{[返順時目]}}
~
\hyperref[sec:index_scriptual]{\small{[返順卷目]}}
~
\hyperref[sec:6ptt2Cl0aVs]{\small{> > > NEXT SERMON > > >}}
\newline
\newline
哥林多後書 12:7-10
\newline
\begin{longtable}{cl}
\hline
\hline
章節 & 經文 (和合本修訂版)\\
\hline
12:7 & \begin{tabularx}{0.7\textwidth}{X} 又恐怕我因所得的啟示太高深，就過於高抬自己，所以有一根刺加在我身上，就是撒但的差役來折磨我，免得我過於高抬自己。 \end{tabularx} \\ \\ \relax
12:8 & \begin{tabularx}{0.7\textwidth}{X} 為了這事，我曾三次求主使這根刺離開我。 \end{tabularx} \\ \\ \relax
12:9 & \begin{tabularx}{0.7\textwidth}{X} 他對我說：「我的恩典是夠你用的，因為我的能力是在人的軟弱上顯得完全。」所以，我更喜歡誇耀自己的軟弱，好使基督的能力覆庇我。 \end{tabularx} \\ \\ \relax
12:10 & \begin{tabularx}{0.7\textwidth}{X} 為基督的緣故，我以軟弱、凌辱、艱難、迫害、困苦為可喜樂的事；因為我甚麼時候軟弱，甚麼時候就剛強了。 \end{tabularx} \\ \\
[1ex]
\hline
\hline
\end{longtable}
$^{1}$各位大兄姐妹,早安!.
今早我們回到神日家,開心嗎?.
很感恩,因為我們好像沒見面很久了.
今天終於可以實體敬拜.
好像拍掌不夠厚.
不如我們一起說Hallelujah讚美主.
預備,一二三.
Hallelujah讚美主.
我們等了很久了.
終於今早有機會能夠大家聚集敬拜天父.
所以牧師因為最近環境的緣故.
想起一個題目,叫做軟弱化剛強.
多謝陳全道剛才幫我們讀出這段經文.
我相信不少大兄姐妹都會有印象.
如果你靈修讀過的話.
我想很多人都不喜歡軟弱這兩個字.
因為軟弱這兩個字一聽到.
印象就是軟婆婆,無精打采.
站著也站不住,走也走不動.
坐也坐不住,常常躺在床上.
這就叫做軟弱.
所以不太多人喜歡軟弱這兩個字.
大家認同嗎?.
但有時軟弱是自己拿回來的.
有沒有留意到?.
有些人不喜歡戴口罩.
不喜歡社交距離.
然後就中招了.
我想大家都知道最近誰中招了.
所以有些軟弱是自己拿回來的.
但有些軟弱是被動的.
即是你不想要,不想得到.
總是無緣無故發生在你身上.
大家相不相信有這樣的事情?.
有時你不喜歡軟弱,不喜歡逆境.
但到我們選擇嗎?.
不到我們選擇.
有時我們很被動就面對這些情況.
正如保羅.
保羅遇到的軟弱其實很多.

$^{41}$因為他傳福音,建立教會過程中.
都不是一帆風順.
我相信大家都知道.
但這些軟弱是事公上帶給他的軟弱.
但他在自己身上也有軟弱.
為甚麼呢?.
因為他告訴大家.
他說在我身上有一根甚麼?.
刺.
這根刺厲害了.
看完節目.
說到這根刺.
你要記得保羅這樣描述.
是一個如意識或圖像化.
讓你多一點明白.
所以有一根刺.
讀到這裡不要想這根刺有多長.
有多粗.
有多尖.
刺下去有多痛.
不是這樣去想.
為甚麼呢?.
因為是如意識說.
原來保羅說我身上有一根刺.
實證是想告訴大家.
原來他身上有一個軟弱.
這個軟弱就像一根刺.
甚麼是軟弱呢?.
根據《解經家》的解釋.
保羅這根刺軟弱.
是指他患了眼疾.
為甚麼保羅會無緣無故患了眼疾呢?.
實證不是無緣無故.
是有原因的.
為甚麼呢?.
原來大家知道保羅為主的緣故.
傳福音.
建立教會.
但他經常被人攻擊.
受迫不.

$^{81}$又被人圍著他.
拉他去坐牢.
有沒有這些印象?.
特別有些人不肯接受福音的時候.
覺得福音令他們的生意不得興隆.
然後就抓保羅.
有沒有這樣的印象?.
所以保羅為主的緣故.
他坐牢的次數都很多.
一個人坐牢.
就跟眼疾有關係.
因為大家都知道.
監房的環境肯定不乾淨.
也不會很衛生.
再加上古時候坐牢的地方.
如果有一盞油燈.
都算很難得.
我想可能連油燈都沒有.
你明白嗎?.
只看天的日子.
天晴一點就看到光.
天暗就看不到.
如果保羅又好不幸.
人坐牢.
他坐牢.
他坐在內監.
就更加淪盡.
你明白我的意思嗎?.
就當作重犯般.
在這樣的時候.
更加沒有光線可以看到.
不過保羅有一個很明顯的習慣.
他喜歡閱讀.
很喜歡看書.
雖然當時沒有書可以看.
但有些洋卷.
你聽過洋卷嗎?.
字刻在洋卷上.
一卷卷的.
所以保羅對那些來探監的人說.

$^{121}$你們來探監的時候.
不用帶葡萄色.
不用帶利賓納.
都不用帶這些.
來探我的時候.
最重要帶什麼?.
我見外衣.
可能有點冷.
帶外衣和皮卷.
因為他很喜歡閱讀.
你知道光線不夠.
弄一弄眼睛.
是否受影響.
就受影響了.
保羅不單喜歡閱讀.
原來他也很喜歡寫書.
所以有好幾卷福音書.
稱為監獄書信.
大家有沒有印象?.
因為這個緣故.
又喜歡讀書.
又喜歡寫書.
但光線不夠.
久而久之.
他患上眼疾.
眼疾有多慘.
有多辛苦.
我想未患過.
真的不知道.
如果患過眼疾的人.
就明白眼疾是很辛苦.
很痛苦.
坐也不是.
走也不是.
想做什麼都做不到.
為什麼?.
眼睛張不開.
大家明白我的意思嗎?.
或者我簡單說說.
可能你會明白多一點.

$^{161}$大家有沒有試過眼瞼?.
眼瞼試過嗎?.
即是眼乾.
我現在明白.
女士們出街.
手袋多數有什麼呢?.
白花油.
棍籠油.
男士出街.
可能要帶眼藥水.
有沒有人帶眼藥水在身邊?.
有一次我這樣經歷.
開車在市區.
開回元朗.
誰不知走到丁馬大橋.
眼睛突然刺痛.
刺痛得不得了.
我知道這是眼瞼.
沒有眼淚水.
但你又開車在丁馬大橋.
你怎麼辦呢?.
停也停不到.
迫著什麼呢?.
一隻手拿著tie 盤.
其實都危險駕駛.
另一隻手不斷.
希望有眼淚水出來.
舒服一下.
我開始明白.
少少眼瞼都令你很辛苦.
何況是眼窒.
大家明白嗎?.
所以對保羅來說.
他肯求主.
求主將病拿走.
他說有三次肯求主.
有沒有留意到?.
三次.
我相信保羅為這個病.
不是單單祈禱三次這麼少.

$^{201}$你試試自己不舒服.
發燒感冒傷風.
病了也不止祈三次.
你明白我的意思嗎?.
他說祈三次的意思是.
我有三次特別的祈禱.
這三次特別的祈禱.
可能我猜是禁食祈禱.
或者通宵祈禱.
或者不眠不休.
很懇切的祈禱.
好像求到主.
你真的要幫我.
不幫我我不行了.
我死定了.
你明白嗎?.
好像這樣的情況.
但三次很懇切.
很迫切的祈禱.
都不見效.
為什麼?.
根刺還在嗎?.
還繼續在.
真是很令人費解.
為什麼上帝不聽保羅的祈禱.
不將根刺拔走.
讓他可以暢暢快快.
有氣有力的事奉上帝.
這樣就好了.
但是.
打回原狀.
根刺仍然都在.
所以我相信.
保羅當時的心情應該很低落.
為什麼呢?.
又有眼疾.
又不能醫好.
又影響他的生活.
上帝又不聽祈禱.
為什麼上帝這樣對待我呢?.

$^{241}$你明白我的意思嗎?.
很多時候.
人的情緒就出來了.
上帝我這麼好對你.
我又傳福音.
又聽你的話.
又建立教會.
為什麼你會這樣對待我呢?.
各位弟兄姊妹們.
這種情況很容易令我們.
怨天.
又人.
很容易令我們粉身碎骨.
為什麼呢?.
因為我們很不如意.
因為我們很冤枉.
我們不是做壞人.
也不是做壞事.
為什麼上帝這樣對待我們呢?.
各位弟兄姊妹們.
有這樣的想法是很正常的.
這種情緒就出來了.
不過.
奇在剛才陳全道讀到.
幾年間.
保羅說.
我為基督的緣故.
就以這些吉濫的不靈慾羞恥.
當為可喜樂的.
有沒有讀到這四個字?.
可喜樂的.
嘩!真是.
很不開心.
明明是很負面的.
很不開心的.
說得不好聽.
罵少半句都偷笑.
你明白我的意思嗎?.
為什麼上帝這樣對待我.
現在保羅竟然說.

$^{281}$將這些軟弱的吉濫.
當成什麼呢?.
可喜樂的.
我相信弟兄姊妹們都很想明白.
什麼原因.
令到保羅可以戰勝這些軟弱呢?.
我今早要跟大家說.
起碼有兩個原因.
是令到保羅對軟弱的體會.
與別不同.
第一.
他察覺到.
原來這個軟弱.
這根刺.
上帝還存在在我身上.
是因為要幫我去認識自己.
大家明白這個意思嗎?.
原來人很多時候.
都不是那麼認識自己.
看人很多時候.
很喜歡比較.
我高過你.
我大隻過你.
我靚過你.
我有錢過你.
人很多時候都這樣比較.
但人很多時候看不到自己的問題.
這是一個最大的問題.
除非什麼?.
有軟弱出現的時候.
他就看到自己的問題.
所以我在醫院探病.
做院務的時候.
我見到有不少的病人.
出院前後.
人有很大分別.
你相不相信?.
有些人沒有病之前.
又可以行,又可以吃,又可以睡.
又可以打.

$^{321}$什麼都可以的時候.
他真的不可味.
眼睛關在哪裡?.
關在這裡.
懂不懂看人?.
不懂看人.
不過當他.
不知為何.
突然之間手腳冰冷.
出汗.
個人發燒.
個人就淚滴.
這樣的時候.
被人送進醫院.
睡在病床上.
那時候.
就真的什麼都不行.
大家明白嗎?.
在這種情況下.
我發覺不少病人.
開始看自己.
看得比較客觀一點.
你相不相信?.
看自己.
看得比較合乎中度一點.
因為他不覺得自己是天上有地下無.
他會發覺自己.
原來我也有軟弱.
所以我見過不少病人.
大病一場之後.
不知為何.
兩夫妻的關係好轉.
你有沒有見過這種情形?.
大病一場之後.
出院之後.
那些兄弟姐妹.
大家的關係.
好像拉近了一點.
你相不相信?.
原來是真的.

$^{361}$一個人病了之後.
出院之後.
開始檢視自己.
開始檢討自己.
我為什麼和人的關係這麼差?.
我為什麼經常插嘴.
又不能輸.
你明白我的意思嗎?.
有些人是不能輸的.
我見過這些人.
他姓吳叫輸得.
就是要贏你.
但是病完之後.
他發現自己的軟弱.
他知道.
我不過是血肉之軀.
何德何能呢?.
大家就發覺到.
原來上帝透過軟弱.
其中一個目的.
是要幫我們認識自己.
各位頂智媒們.
在疫情當中.
今次的新冠病毒.
你發覺自己很多事情.
想做都做不到.
你想去旅行.
去不去到?.
你想坐飛機.
能不能飛?.
什麼都不行.
頂智媒們.
差點連出門都不行.
是吧?.
好像香港政府.
還讓我們出門.
如果限什麼令.
禁什麼令.
連門口都出不了.
我兒子在美國.

$^{401}$試過有一段時間.
連門口都出不了.
各位頂智媒們.
你發覺人.
是否真的在軟弱當中.
看到原來我不過是一個.
很普通的人.
是吧?.
所以保羅明白.
為什麼上帝將那根刺.
仍然放在他身上.
第二個原因.
為什麼上帝不拔走那根刺呢?.
我相信因為.
上帝透過那根刺.
要讓保羅放下自己.
明白嗎?.
放下自己.
放下自己什麼?.
放下自己的執著.
是吧?.
放下自己的強項.
人是很有趣的.
人人都很喜歡抓住一些東西.
明白嗎?.
有什麼就抓什麼.
有些人很貪心.
尤其是錢.
抓得緊緊的.
想著有錢人就很有安全感.
人人都抓住一些自己的強項.
自己的優點.
自己的執著.
所以保羅告訴別人.
我是純種的希伯來人.
猶太人.
為什麼?.
因為我是希伯來人所生的希伯來人.
我第八天就受國禮.
厲害吧?.

$^{441}$我從小就明白律法.
律例典章.
厲害吧?.
我受教於加瑪列門下.
在麻省哈佛畢業.
很厲害吧?.
我很熱心迫不及待教會.
厲不厲害?.
很厲害.
後來我悔改信主之後.
大馬色路上主向我顯現.
我迴轉悔改.
我信奉基督.
從此之後.
我熱心傳福音.
厲害吧?.
而且我得到的啟示.
是零碩多.
比其他使徒更多.
所以不覺得保羅寫的經書特別多.
你有沒有留意到?.
他說曾經有一種屬靈的經歷.
就是被上帝提到三層天.
聽到奧秘的聲音.
厲害吧?.
厲害到爆炸.
超厲害.
如果保羅要驕傲.
容易嗎?.
他要展示高度.
能不能展示?.
可以展示很多高度.
但是上帝現在告訴他.
一斤刺.
就可以砌低.
你相不相信?.
一種病.
就可以將你的強項.
化為烏有.
大家相不相信?.

$^{481}$我在醫院見到不少病人.
有年長的.
有中年的.
有年輕的.
有學業有成的.
事業有成的.
錢很多的.
底子沒有.
以為可以很多東西誇口.
但是一個病.
令他身不由己.
你相不相信?.
一種病.
令他所擁有的.
都化為烏有.
有沒有得誇口?.
沒有得誇口.
所以現在保羅明白.
為什麼那一斤刺.
仍然在他身上.
因為他很容易.
心花怒放.
輕飄飄.
為什麼?.
因為他什麼都比別人厲害.
大家明白嗎?.
但現在一種的硬質.
令保羅什麼都無可誇口.
所以因為這個緣故.
保羅就開始發覺到.
我以前倚賴的東西.
現在真是什麼?.
是一無所有.
難怪他對林多的教會說.
我在你們當中.
立了一個志向.
什麼志向呢?.
就是告訴你們.
我在你們當中.
我並不知道什麼.

$^{521}$只知道主耶穌和他的十字架.
大家對這本聖經有印象嗎?.
其實保羅懂不懂?.
我想問大家.
他懂很多東西.
很有知識.
很有學問.
但是他告訴哥林多教會.
在你們當中.
我並不知道什麼.
只知道一件事.
就是主耶穌和他的十字架.
頂止門門.
原來這一斤刺.
是幫助我們放下了自己的執著.
自己的強項.
自己的優勢.
當我們這樣做的時候.
很快地.
我們會把我們的視線眼光.
轉移去仰望什麼?.
既然人什麼都不靠.
唯有靠什麼?.
唯有靠上帝的恩典.
所以主耶穌對保羅說.
我的恩典是什麼?.
救你用.
因為我的能力.
在人的軟弱上顯得完全.
所以保羅說.
我更加喜歡誇自己的軟弱.
好叫基督的能力來到什麼?.
奉備我.
各位頂止門門.
讀到這裡的時候.
我們現在開始明白.
為什麼保羅可以誇軟弱為什麼?.
為剛卻.
因為他一放下了自己的時候.
他就轉移仰望上帝的恩典.

$^{561}$上帝的恩典是什麼?.
上帝的恩典就等同上帝的能力.
所以當他仰望上帝的恩典的時候.
上帝的能力就賦予在他身上.
結果他就覺得自己不再軟弱.
所以最後那句怎麼說?.
我什麼時候軟弱.
什麼時候剛強了.
頂止門門.
在當中做父母的.
做爺爺奶奶.
阿公阿婆那些.
你有個孫七八歲.
或者六七歲的時候.
你猜猜你最想他學的東西是什麼?.
頂止門門.
或者做父母的.
你的子女六七歲.
七八歲的時候.
你最想他學什麼?.
可能有人想他學彈琴.
跳芭蕾舞.
柔道.
空手道.
強身.
很多東西都想小孩學.
跆拳道.
頂止門門.
我猜有一件事你一定很想他學.
為什麼呢?.
這件事他學了之後.
這個小孩從此之後.
你不需要太擔心他.
因為他懂得救自己的命.
我講得這麼白.
你知不知道要學什麼?.
當然要學聖經.
聖經上帝會拯救他.
保護他一生.
知不知道學游泳.

$^{601}$你的孫子六七歲.
七八歲的時候.
你想不想他學游泳?.
想啊.
立即參加游泳班.
初級到中級.
什麼級都不要緊.
有個教練教他.
為什麼?.
你希望有一天他學會游泳.
怕不怕淹死?.
沉船都不用怕.
為什麼?.
他起碼熬到幾分鐘.
熬到半小時.
他懂得游泳.
懂得救自己.
頂止門門.
這個小孩何時真正學會游泳呢?.
教練由第一課開始.
給他一塊浮板.
叫他游.
踢腳.
搖手.
上了幾課之後.
開始不錯了.
由初級到中級.
教練跟他說.
有一天.
小孩今早游泳.
你不要用浮板.
拿一塊浮板.
你就這樣游到我面前.
我就站在你面前.
你游過來.
你放心游過來.
如果小孩聽說有信心.
向著教練飛撲過去.
你猜他學不學會游泳?.
學會了.

$^{641}$那一刻他學會游泳.
他掌握到秘訣.
但如果他永遠拿著浮板.
他學不學會游泳?.
他永遠都學不會游泳.
因為他所有的依賴就在浮板上.
直至有一天.
他放下浮板.
憑信心.
游去教練那裡.
他就成功了.
頂止門門.
人生有很多東西是需要放下的.
如果你不放下的話.
你會很辛苦.
你會時時都陷在軟弱當中.
但當你願意放下的時候.
保羅現在明白.
原來我放下了自己的強項優勢.
我的執著的時候.
我轉眼仰望上帝的恩典.
上帝的能力.
看顧我.
保護我.
所以我可以告訴大家.
我什麼時候軟弱.
什麼時候就可以剛強.
所以今天這篇訊息.
求主幫助我們.
讓我們在疫情還沒有過去的時候.
我們都有很多軟弱.
但不要埋怨這些軟弱.
軟弱當中.
上帝有他的心意.
有他的計劃.
求主幫助我們.
學習保羅.
他的榜樣.
化軟弱為剛強.
我們祈禱.

$^{681}$天父上帝多謝你.
透過經文給我們的教導.
是十分寶貴.
讓我們看到的.
不只是自己的強項優勢了不起的地方.
讓我們看到的是我們的不足.
我們的軟弱.
我們的瑕疵.
好讓我們更加謙卑在你的面前.
我們事奉上帝.
我們與人相處.
都能夠存在一個謙卑的態度.
以致我們能夠.
可以在任何的地方.
來見證榮耀你.
化軟弱為剛強.
我們求心祈禱.
奉耶穌的名頭.
阿們.
謝謝大家.
\newpage



\section{詩篇 126:1-6}
\label{sec:6ptt2Cl0aVs}
\textbf{2020.10.11《流淚撒種者的歡呼》 詩篇 126\_1-6 蔣文忠博士}
\newline
\newline
連結: \href{https://youtube.com/watch?v=6ptt2Cl0aVs}{\texttt{https://youtube.com/watch?v=6ptt2Cl0aVs}} ~~~~ 語音日期: 2020-10-14
\newline
\newline
\hyperref[sec:6c0M0guWu_4]{\small{< < < PREV SERMON < < <}}
~
\hyperref[sec:index_chronic]{\small{[返順時目]}}
~
\hyperref[sec:index_scriptual]{\small{[返順卷目]}}
~
\hyperref[sec:YoAst8SyrWc]{\small{> > > NEXT SERMON > > >}}
\newline
\newline
詩篇 126:1-6
\newline
\begin{longtable}{cl}
\hline
\hline
章節 & 經文 (和合本修訂版)\\
\hline
126:1 & \begin{tabularx}{0.7\textwidth}{X} 當耶和華使錫安被擄的人歸回的時候，我們好像做夢的人。 \end{tabularx} \\ \\ \relax
126:2 & \begin{tabularx}{0.7\textwidth}{X} 那時，我們滿口喜笑、滿舌歡呼；那時，列國中就有人說：「耶和華為他們行了大事！」 \end{tabularx} \\ \\ \relax
126:3 & \begin{tabularx}{0.7\textwidth}{X} 耶和華果然為我們行了大事，我們就歡喜。 \end{tabularx} \\ \\ \relax
126:4 & \begin{tabularx}{0.7\textwidth}{X} 耶和華啊，求你使我們這些被擄的人歸回，好像尼革夫的河水復流。 \end{tabularx} \\ \\ \relax
126:5 & \begin{tabularx}{0.7\textwidth}{X} 流淚撒種的，必歡呼收割！ \end{tabularx} \\ \\ \relax
126:6 & \begin{tabularx}{0.7\textwidth}{X} 那帶種流淚出去的，必歡呼地帶禾捆回來！ \end{tabularx} \\ \\
[1ex]
\hline
\hline
\end{longtable}
$^{1}$(字幕提供:陳志全).
(陳志全) 頂智妹很感恩.
能夠和大家一齊見面.
在疫症讓崇拜進行的時候.
我們能夠聚集.
亦有頂智妹在網上.
一齊和我們在家中敬拜.
這個我們相信是一個.
今天經常用的字眼叫做新常態.
不知道大家有沒有聽過.
其實也是一個很講究一種.
似乎在一個世紀疫症之後.
對於一個不只是教會的崇拜.
其實是整個人類的生活.
全球的人類生活的一種影響.
究竟我們能不能夠在當中去適應呢?.
又或者其實我們當要轉變.
或者被轉變的時候.
其實我們的生命又能夠怎樣在.
這種我們不能夠把握.
不能夠在用人的力量去控制的.
一個大的挑戰之下.
我們作為上帝的子民.
我們作為上帝呼召的群體.
作為教會.
我們又能夠怎樣去在信仰當中實踐.
又或者是.
老傳道邀請我講的一個盼望的主題.
就是怎樣能夠持守盼望.
繼續向前行呢?.
特別是剛才陳傳道跟我們禱告.
在香港當中也有很多的挑戰.
在政治上有很多的壓力.
究竟我們作為教會的群體.
怎樣把握我們有的自由.
應該有的自由.
信仰的自由.
宗教的自由.
繼續去為主在我們這個社區當中.
在我們身處的地方當中作見證呢?.

$^{41}$我相信也是我們在一個新常態環境之下.
對於我們教會的群體.
一個很大的挑戰.
求主憐憫幫助.
讓我們今天聚集.
我們更加用一種珍惜的心懷去面對.
也都是更加去.
躺開我們的心靈.
去問上帝.
今天我們要.
能夠怎樣為祂作供呢?.
在一種今時今日的狀態當中.
我們一群人聚集的時候.
好像有一種不是必然的感覺.
怎樣能夠繼續.
祂為主去.
實踐祂的使命呢?.
去聆聽祂的話.
去回應祂的呼召呢?.
下幕今天我們藉著.
詩篇一百二十六篇.
這個詩篇.
一首上行之詩.
一首都是猶太人他們經歷很大的苦難.
經歷很大的創傷.
歸回他們的地方的時候唱的詩歌.
成為我們一個.
屬靈上面的一個提醒.
一個能夠給我們一些屬靈上面的資源.
去讓我們的生命當中.
怎樣去回應他們的主.
也認識更深我們的主.
是怎樣去保守祂的子民.
去向前走下去.
我們也一起有個禱告.
然後去聆聽今天這首詩篇.
上帝的話.
天上的多謝你讓我們今天.
我們能夠聚集.
也是隔了一段時間.

$^{81}$我們又在一個可以放寬的形勢之下.
可以聚集.
但是我們知道.
我們仍然面對很嚴峻的疫症.
全球的影響.
也是一段時間才能夠去恢復.
但是在想的時候.
其實這種新的常態之下.
恢復之後是一種怎樣的形態呢?.
都不是人能夠可料得到.
我們求主憐憫.
讓我們面對這些接種而來的挑戰.
這些改變.
我們心靈上的盛載.
能夠在主的話當中.
在主的道當中.
我們紮根.
我們穩妥在你的話之中.
去迎面這些新的挑戰.
求你幫助我們今天.
一起聚集.
一起聆聽詩篇.
一百二十六篇的時候.
我們都能夠真的謙卑我們的心.
開放我們的心.
隨著聖靈的引導.
去讓我們的心能夠聽得明白你的話.
藉著詩人向你的禱告.
向你的呼求.
向你的呈獻.
成為我們今天同樣.
紅恩同樣的弟兄姊妹.
一起向你所去禱告的詩歌.
求你親自帶領以下的時間.
同心祈禱.
奉主耶穌得勝.
明治已強.
阿們.
告訴大家的題目就是.
流淚殺種者的歡呼.

$^{121}$可能大家剛才讀的時候都很熟悉.
是不是?.
第五節.
流淚殺種的必歡呼收割.
這是我們常常傳福音都會用的金句.
但我猜我們大家也未必常常有機會.
是整首詩篇這樣讀.
今天我都和大家比較詳細地讀這首詩篇.
(字幕組:我已經不想再看了).
可以幫我按一下下一張.
謝謝.
不好意思.
我先和大家介紹上行之詩.
這個詩篇.
由120篇至134篇.
裡面有十五首詩歌.
我大致上很簡單的和大家引子而已.
詩篇裡面.
很多人就說其實是一個.
上帝的屬於他的子民.
這群朝聖者.
去登耶和華的山的時候.
去預備自己的心靈去唱詩歌.
但是更多的聖經學者都會體會到.
其實詩篇的成書.
在被擄的時期.
之後大家可能都聽過被擄的時期.
是巴比倫將整個猶大亡國.
滅了整個國家.
然後就將猶撒冷的城拆掉.
將聖殿燒毀.
將他們裡面的很多器具都搬到加比倫.
不單止這樣.
可不可以回到剛才那一張.
不好意思.
我不是很控制得到.
再上一張.
那麼.
他裡面就是.
就是子民裡面很多人呢.

$^{161}$都去被擄的巴比倫.
一種異地的生活.
一種離鄉別井的生活.
但是回去的時候是七十年之後.
所以很矛盾是不是?.
你移民七十年.
他現在整個家庭都回去了.
是不是?.
你的子女你的孫子.
基本上都說了外語了.
你都不想走.
人生有時候都是去到最後的階段.
在當中安養晚年.
但是上帝呼召回子民.
全部人.
就是連根拔起回到故鄉.
是不容易的.
心情相當複雜.
很矛盾.
很多的暢商的記憶都回來了.
很多的舊事.
其實帶給來那種.
對於生命的影響.
其實那種回憶都會回來.
但是被擄者相信上帝的心意.
回去歸回是有它的意思.
所以他們就會唱上行之詩.
表示這群子民歸回他們的國土.
同時間我提醒大家.
就是他讀詩篇的時候.
他有一些特點.
其中一種就是.
他會有一種去處境化的特點.
什麼意思呢?.
就是說他沒有很刻意地.
將一首詩篇原先的時候.
那個歷史處境告訴你.
他不想你推敲.
是什麼時候發生的事.
又或者我們不需要去想.

$^{201}$就讓我們進入私人的狀態.
就可以了.
我們不需要去限制著那種歷史.
就是我們今天活在香港.
對於活在千多年前的.
猶太被擄回歸的狀態的時候.
他不需要我們去考究真實的事情.
隨著聖靈的感動.
就代入在我們今天的光景.
進入詩篇裡面.
私人的狀態.
感受到私人的處境.
跟我們好像有一種相關性.
就可以了.
這種就是一種.
我們不是在讀歷史書.
我們不是在讀五經.
他有一種很清晰的歷史背景.
在詩篇裡面有時候是不會有的.
他就說要讓我們全程投入.
到他的詩句裡面就可以了.
所以裡面也有很多的圖像.
讓讀者進入私人和主的關係.
這首詩篇裡面.
六節裡面.
簡單裡面也有他的圖像.
麻煩下一張.
我簡介這首詩篇的大綱.
一至三節.
剛才讀的時候.
大家都能夠聽到.
其實他有一個.
每當.
第一節.
當姚華帶我們回去的時候.
想當年.
回想起來.
是一個已經發生了的事情.
所以就說是一個益述.
一個過程.

$^{241}$益述什麼呢?.
益述一些意想不到的作為.
想都沒想過會發生的作為.
那圖像就是.
讓好像做夢的人.
有重複的字眼.
就是姚華行左大事.
第二節是在說外邦人.
有人說.
在他們當中.
姚華行左大事.
而第三節.
就是這群子民子說.
姚華果然為我們行左大事.
一會兒我可以再詳細和大家分享這一點.
到了第二段的時候.
就是一個禱告.
很多人都說.
很多醒學者都說.
詩篇每一首都有禱告.
如果真的用一個禱告的字眼的時候.
真正的禱告在這篇詩篇裡面.
就是第四節.
求你使我們被擄人歸回.
這個就是在說他們當下.
他們回想了過去.
立於當下.
他們的困難還沒解決過.
記住.
回想過去之後.
不是代表問題解決了.
他仍然活在當下的處境.
他們可能是被擄之中.
你知道被擄也有很複雜的過程.
有幾批人.
在不同時間回到當下.
都是一段很長的時間.
我們猜不到.
也不知道是哪一段時間.
但是他們還沒走.

$^{281}$他們就讓上帝.
求上帝帶領他們回到當下的時候.
就好像南地的河水湧流.
這裡也是在說上帝將要成就的作為.
然後重複的字眼就是.
流淚殺腫的.
必歡呼收割.
其實也有不同的.
在第六節裡面.
他還說了帶著什麼.
帶著種子出去.
帶著禾菌回來.
收成.
背負著.
大概不只是說他背著背囊.
就是拿了什麼出去這麼簡單.
帶了什麼回來這麼簡單.
是在說他們一種心靈上的壓力.
心靈上背負著一種什麼心情出去.
如今背負著一種什麼心情回來的狀態.
如果大家看這兩個圖像.
好像做夢的人是在說.
這群以色列人.
如何經歷上帝.
但是到了第二段的圖像.
好像南地的河水浮流.
就是在說上帝將要成就的作為.
這個南地的圖像.
是說什麼呢?.
我先解釋這個.
南地是一個在以色列南端的.
一個接壤埃及的地方.
其實也是一些荒漠之地.
按一些考古的人去了解.
其實這個荒漠之地長期都是很乾旱的.
但是每年裡面就有一些季節性的暴雨.
使得那些乾旱的河床湧流著河水.
瞬間就湧流.
其實就是在說上帝的作為.
是有一種突如其來.

$^{321}$卻是帶來了一種很豐盛.
很豐滿.
甚至滿溢的一種這樣的狀況.
這個就是在說南地的河水湧流.
我們逐段去看.
麻煩你下一下.
我給大家在螢幕上.
是和修本.
大家應該是剛才讀的那本和修本.
那麼大家再聽一聽和修本的版本.
當耶和華洗石安.
庇祿的人歸回時候.
我們好像做夢的人.
那時我們滿口喜笑.
滿嘴歡呼.
那時列國中就有人說.
為他們行了大事.
果然為我們行了大事.
我們就歡喜.
所以強調多兩個那時.
當時發生什麼事.
我就問.
是一種怎樣的歡喜呢.
是不是我們等了很久.
沒有聚餐.
然後出去吃一餐.
自助餐.
這個時候不要去.
大家千萬不要再.
這時候大家忍一忍手.
但我們很想去旅行.
很想去吃一餐.
那種是我們香港人那種.
開心的一種.
一群朋友一群家人.
可以這樣聚集.
享受的時間.
但是.
這裡說的那種怎樣的開心呢.
怎樣的喜樂呢.

$^{361}$我也說漏了一樣東西.
其實開心喜樂的字.
整首詩歌都充滿著.
其實基本上.
第一段有兩次.
第二段也有兩次.
這一段你看到第二節.
滿口喜笑滿舌歡呼.
第二個就是第三節.
我們就歡喜.
好像很開心那樣.
但那種開心不是純粹一種.
得到一種物質的滿足.
一種觀感上.
觀能上的享受.
那種比較表層的開心而已.
剛才說了背景.
他們.
被擄.
一段很長的日子.
你想像一下.
如果我們有一個大病.
或者是.
好像這次的疫症.
一段長的日子.
什麼時候會.
可以恢復呢.
什麼時候會沒事呢.
什麼時候會醫得好呢.
有時候甚至等到我們自己都.
發呆的.
不知道大家有沒有經歷過一些狀況.
可能未必發生在自己身上.
或者發生在自己家人裡面.
或者發生在社會裡面.
但是.
好像突然間.
法官宣告無罪釋放.
那種的.
解放.

$^{401}$那種意想不到的開心.
就是這裡.
在說的.
這份的歡喜.
這份歡呼的狀態.
為什麼這麼說呢.
你看到第三節.
第二節尾.
在列國中有人這麼說.
大家留意列國的人.
不是善男淳女.
也不是以色列的朋友.
往往是以色列的敵人.
往往都在這裡.
嘲弄以色列.
猶大.
當他們被擄的時候.
他會說他們的神.
為什麼不幫你打仗呢.
為什麼不保護你們呢.
要讓你們落到這樣的田地呢.
其實以色列歷代.
都是一種他們.
列強爭奪的兵家之地.
所以不會是一個.
對他們帶著.
善良的心.
或者帶著同情的心.
去對待以色列.
連列國都說.
聽到這樣的消息.
原來以色列的人.
以色列的上帝.
帶回.
以色列人.
歸回他們的國土.
留心第三節.
你現在看二三節的關係.
大概意思就是.
他竟然.

$^{441}$這群以色列人.
是要聽到.
外邦人.
傳來的消息.
他才會開心.
所以果然.
中文的字譯得不錯.
原來真的這麼做.
我們不知道.
我們.
今天.
我們每一個人.
你的心靈狀態.
你的情緒狀態.
是不是很沉鬱呢.
對我來說也很沉鬱.
其實有時候也睡得不好.
也感恩昨晚睡得很好.
可以有精神.
有時候主日也擔心.
如果睡得不好.
要趕渡.
上台睡不著.
或者發台瘟.
說不出要說的話.
這些是傳達人的壓力.
有時候你問牧師傳達人.
每天狀態好不好.
可能是整天的事.
我想大家的狀態.
會覺得自己沉鬱了很久.
其實從去年六月開始.
我已經開始很沉重.
到今天.
也不覺得那種沉重的心情.
有一種開懷的感覺.
但是我們的指望在哪裡呢.
我們的盼望在哪裡呢.
我們會不會好像獨到私人.
他憶想起上帝.

$^{481}$其實在我們曾經的生命當中.
有一些意想不到的作為.
我們把持著這份上帝的作為.
這份就是他的屬性.
這份就是上帝本身.
是一個這樣看待人的上帝.
他在一個最危難的時候.
最意想不到的時候.
私行去拯救.
讓我們得到一份偶然.
意想不到的驚喜.
這份就是一種這樣的喜樂.
是這首詩篇頭三節所流露出來的.
也跟大家分享.
其實在這個疫症當中.
也要多謝謝牧師.
我爸爸在八月底的時候.
就受洗歸入教會.
其實他一直都很抗拒的.
我媽媽在上年度就去世了.
她在床邊也受洗歸住.
所以其實也感恩.
兩位老人家在人生最後階段.
都能夠在主面前.
去接受這份救恩.
去認耶穌是主.
是一種心裡面很大的安慰.
想不到.
意想不到.
以他們整個大半生的經歷.
我覺得機會很微.
很難的.
我真的硬著頭皮闖進去.
跟他說福音.
其實只是一種堅持去說.
到我媽媽去世的時候.
我爸爸就開始有一點鬆動.
那種對於信仰.
但是到他發病五月的時候.
中風發病的時候.

$^{521}$他都說你不要說這些.
第一次跟牧師一起去探望他的時候.
他立刻說你們不要說這些.
走吧.
立刻有很大的抗拒.
然後疫症基本上沒有機會探望到他.
所以有時候.
都是一些很偶然的機會.
見一面.
說一句說為你禱告.
到了八月底的時候.
有機會帶他去覆診的時候.
我就邀請牧師.
就是鄭牧師和陳牧師.
一起去住的安老院當中.
有機會跟他分享.
很感恩.
他竟然在那個時候就很認定.
主是他生命的主.
這種改變.
這種突如其來.
意想不到的作為.
都不是我們作為後輩兒女.
可以想得到.
就算我口才多好.
多熟悉聖經.
都不一定.
我想你們傳福音的人都明白.
福音不是一個人的口才或者準備.
就傳得到的.
信心.
接受不是一個人.
能夠預備對方做得到.
但這是上帝臨到的作為.
上帝讓我們家人.
去認定他是那位.
似人意想不到的上帝的作為.
按他的時間.
有他的大事要他去成就.
弟兄姊妹.

$^{561}$我們所信的就是一位.
這樣去成就大事的上帝.
求主讓我們去憶想.
這個憶想也很重要.
我們的信仰.
是一個歷史.
一直傳遞到今天.
你個人的歷史一樣.
教會歷史.
甚至說從初期教會到今天.
我們的教會.
整個人類的歷史.
上帝的作為.
都是始終如一.
信實不變.
去憐憫信他的人.
講到第二段的時候.
就是「確信上帝將要成就」的作為.
悉述過去.
他都要立於當下.
然後展望將來.
前路對他們來說都是很迷茫.
前路都是很暗淡.
前路是很多未知之數.
要他們去投身.
要他們去嘗試.
要他們去進入.
所以這班被擄的人.
他們求耶和華.
我讀一次第四節六節給大家聽.
耶和華求你使我們.
這些被擄人歸回.
好象南地的河水伏流.
流淚殺腫的必歡呼收割.
那帶腫流淚出氣的.
必歡呼帶禍菌歸回.
剛才我解釋了.
這個「好象南地的河水」.
就求上帝你好像以往一樣.
去對待我們.

$^{601}$剛罕的河床.
你可以瞬間去充滿著湧流的河水.
他們在求什麼呢?.
這個禱告很有意思.
第四節說.
求你使我們這些被擄的人歸回.
我比較喜歡呂鎮忠的譯法.
恢復我們的故業.
這裡有兩個不同.
跟我們和合本和修本的譯法.
第一個就是.
和合本和修本是說人的.
被擄的人.
但是呂鎮忠的譯法就譯成了一個.
的情景.
故業.
繁榮昌盛的過去的日子.
那種大家一起生活的.
那種社會的生活的狀態.
是說那種整個以色列的.
當日的生活.
當然有人在當中.
去構成這個生活.
但是重點有一點不同.
人就是歸回.
回到那個地方.
但是在說整個光景的時候.
他就用恢復.
restore.
恢復我們的處境.
所以和大家一起想.
這個恢復.
剛才也說了.
我們當處於一個.
今天的香港處境的時候.
不知道大家有沒有想過.
上帝求你恢復我們昔日的面貌.
有沒有感其過肚.
可能很多人都說其實是恢復不了的.
你不會回到以往的日子.

$^{641}$那一定是的.
大家不要奢望著.
可以回到兩年前的.
可能香港的狀態還沒有回到.
所以我們要有信心.
繼續面對面前的改變.
就算我們人的狀態回不了.
很簡單.
你病過.
跌倒.
你康復的時候.
你的骨頭也回不了你之前.
大家的額頭你明白.
大家知道.
心裡明白.
那個是一直這樣衰退的.
我們要接受這個就是人的身體.
我們要因然接受.
這個是我們身體的狀態.
不會100\%恢復回以往.
但是恢復仍然有意思.
恢復上帝主權之下的生活.
恢復上帝引領著這群子民.
以他為中心的敬拜.
彼此守望.
彼此分享.
那種生活.
即使外來環境轉變.
我們仍然可以祈禱恢復我們.
這群上帝子民.
那種向人和人之間.
向這個社區那種的.
服侍的生活.
我們可以這樣祈禱.
求主恢復我們這個故業.
大同心就在第五第六節.
提到我講到的主題.
流淚殺統者.
這首詩來到這一步.
很特別.

$^{681}$我覺得很特別.
回到一個人的狀態.
是不是?.
他講完上帝的作為很偉大.
我們可以讚嘆敬拜.
但是他留在人.
這是對人的表達.
對人的描述.
是一個比喻.
一個圖畫.
剛才講了一個圖畫.
除了南地復原之後.
還有流淚殺種.
歡呼收割.
這個圖畫.
他不是真的在講他們全部是奸眾.
大家明白.
他是在用殺種者的身份.
殺種者面對的挑戰.
或者那種的修成.
來去描述這班.
回到去的人的生命狀態.
大家明白嗎?.
都是在描述我們.
所以我們.
我還有時間和大家一起去投入.
下去流淚殺種者.
這個意境裡面.
但我們可能要問.
其實我們.
我不知道.
可能來到洪水橋.
就多一點接觸.
可能還是用奸眾的人來修成.
大部分.
我住在九龍.
住在城市的人.
基本上.
當然大家其實都沒有分別.
香港來說.

$^{721}$我這麼說而已.
大家不是用一種.
就是.
農奸的生活.
來修成.
但是如果你想象.
你認識一些人.
你明白的.
如果真的.
以前我媽媽的親戚.
就在大陸裡面.
真的做奸田.
就真的這樣去過.
用奸田來過活的生活.
大家明白就是.
很多因素.
不是一個商業社會裡面.
想象得到的.
天氣.
你的外來環境.
你有沒有蟲.
住你.
咬你.
你有沒有一些好的.
就是下雨.
然後就日曬.
就是有太陽的.
那種照管等等.
很多的因素不是你控制的.
你不必然有飯吃.
其實.
它更加有一種.
就是你有修成的時候.
帶著一種感恩的心.
去面對.
跟我們.
習慣了.
一定有糧出的.
一定一定會有.
基本上是.

$^{761}$就是必然有一種.
就好像付出的有回報.
那種心態是有點不同的.
因為更多更多因素.
去影響一個耕種者.
他要面對.
那個大環境的挑戰.
外來環境的挑戰.
我再用一個圖畫.
跟大家去思想一下.
這個是梵高的作品.
1888年他畫的時候.
當時很多這些殺種者.
用《聖經》意象來畫的圖畫.
或者叫做《夕陽下的殺種者》.
這幅畫裡面.
或者當時的畫裡面.
跟畫的有些共通點的.
就是將那個殺種的人.
畫得比較很有線條的.
還有很有力量的.
他就很堅定.
這樣去將種子殺在地裡面.
不好意思.
可能螢幕上不太清楚.
他一隻手拿開袋子.
一隻手就將種子殺在地裡面.
一個很簡單.
很有力量.
但是一個很有尊嚴.
他要做的工作就是做這件事.
他不需要做很多花巧的東西.
也不需要做很多的東西.
他是一個殺種者.
就做足殺種者要做的功夫.
就足夠.
這幅畫令我很深刻.
就是因為他頭頂上有一個很大的太陽.
你知道梵谷畫這些星星.
太陽都畫得很出色.

$^{801}$好像真的在你面前流動著.
散發熱力.
一種這樣的狀態.
這個太陽在他頭頂上.
彷彿像上帝的照管一樣.
保守他的工作一樣.
其實這幅畫的左下角.
他一隻手伸下去.
將種子殺在地上的時候.
其實他有一些.
受著太陽光的光照.
那些種子都好像黃金閃閃發光一樣.
這些黃金流在地上.
表示著.
殺種者的功夫.
是得著更高更超越的上帝.
這個陽光.
就是保守著他.
照顧著他.
讓他的種子.
能夠得到好的保守.
落到好的土地當中.
能夠有收成.
鄧子梅詩篇裡面.
用殺種來比喻一些被擄人.
還要流淚殺種.
流淚是一種哀傷的狀態.
是一種辛酸的狀態.
當你讀到流淚殺種.
是表示一個出去.
去到異邦之地的人的時候.
但是上帝告訴你.
你一定會歡呼收割.
你今天就歸回你的土地.
是一個歡呼收割的日子.
很有意思.
隔了七十年.
詩篇裡面將到.
被擄和歸回.
是放為一個一整體的事情來看.

$^{841}$上帝整件事上.
是有保守.
拆開來看.
我不詳細講了.
當時被擄的時候.
很多爭議性.
有些人覺得留下.
有些人覺得要離開.
才是上帝的心意.
所以有不同的先知.
有不同的討論.
但是到了詩篇.
作者的時候.
過了幾十年了.
能夠將整件.
這麼沉重.
這麼意想不到的事情.
在歷史當中發生的事情.
能夠作一個整合.
作一個整理.
原來我們出去的時候.
是流淚殺眾出去.
當今天歸回的時候.
是帶著歡呼聲.
帶著和菌.
帶著修成.
這樣回去到他的國土當中.
我想講一個見證.
最近大家知道疫症的事.
很多教會有不同的工作.
教會在油麻地區.
傳統我也熟悉.
所以他做的事情令我很深刻.
令我帶著一個.
曾經也去當中.
帶著大學生.
曾經去參與過.
他一直有一個叫做.
由淺入深的一個計劃.
在油麻地區.

$^{881}$最初一兩年前的時候.
還沒有疫症.
其實他只是做一些.
很基本的工作.
怕怕糖水.
或者他們本身不是很多資金.
到了一些傍晚時間.
就去一些麵包店買一些低價麵包.
送給一些需要食物的街坊.
給他們去溫飽.
後來就開始覺得.
希望能夠給他們一些更好的食物.
就在教會當中去煮食.
但是教會又不准用明火.
所以用一些電磁爐.
很簡單的.
只是做一些簡單的飯盒.
但是也有一些暖飯.
去送給街坊去吃.
到了疫症的時候.
需求很大.
開始我去探望他的時候.
他那一次一整就整了七十個.
他跟我說.
以往只是做三十個飯盒.
今天竟然這麼神奇.
做了七十個飯盒.
而且剛剛送給街坊.
和一些少數族裔.
在油麻地區遊蕩的少數族裔.
去分享.
後來他們面對一些.
大家知道爆發疫症.
大廈有人確診.
於是教會就很小心.
行政,安全.
就不准聚會.
一向就是逢星期三晚去聚會.
那怎麼辦呢?.
很多人有需要.

$^{921}$其實更大需要.
但是竟然教會因為疫症之下.
限制了.
怎麼將一份拍案帶給他們呢?.
後來他想到.
不如就收集了街坊的電話.
結果就找義工逐個街坊.
帶電話給他們.
反而更好.
因為跟他們就更多溝通.
以前可能只是送飯的時候.
你也知道說不了很多句.
祝福你.
或者說一兩句.
他就拿飯走.
不是一個很能夠對話溝通的時間.
但是竟然.
當用電話溝通的時候.
就一對一地了解他們的狀態.
聽到更多辛酸又流眼淚的故事.
和了解到他們的經濟實況.
以至他們有限的資源.
可以真的盡量去因應他們的需要.
給他們應該有的資源.
後來開始限聚令.
那時候三月底.
就開始.
大家知道什麼叫密難民.
密記就不能夠通宵營業.
於是原來很多人住在那裡.
不知道元朗或者附近有沒有.
於是他們很多探望的朋友.
他們突然間.
另一個需要了.
就是沒有房子住.
這位傳媒人就.
他說很難.
從來沒有做過這些事情.
怎麼能夠將探望帶給他們呢?.
後來很神奇.

$^{961}$他們就聯絡了很多不同的人士.
在那個地區當中.
有些是良心的業主.
有些就是做社福界的一些弟兄姊妹.
朋友.
甚至微信的人.
於是聯繫了的時候.
就盡量幫這些密難民.
可以安排一些短期的住宿的機會.
他說有一個故事.
分享的時候也令我很觸動.
他說.
有一位朋友上去的時候.
拿著一個遙控器.
就是好像這樣.
按電視的時候.
他說跟這位傳媒人說.
很久都沒有試過.
用遙控器來選擇.
我看的電台.
其實住在這裡挺好的.
激勵了他.
他覺得需要找工作.
需要讓他的生活.
能夠提升到他的水準.
我問他.
怎樣能夠可以.
每一個危機.
似乎你都轉為一個契機.
似乎你都可以帶來一個.
更高層次的服事.
你怎樣能夠做得到?.
他回答的時候.
也令我很深刻.
也是我今天.
跟大家鼓勵的話.
他說其實我也是.
行步見步.
不過他覺得一個上帝.
一個值得去信靠.

$^{1001}$他會帶給人不怕夢的上帝.
只要他走出那一步.
就好像一個殺種者.
他走出那一步.
雖然面對著很辛酸.
流著眼淚的狀態.
但相信上帝會使用.
上帝會將這些種子.
化為祝福到人的一個實際行動.
所以每次他帶來一個.
見到困難的時候.
心情都很沉重.
但他仍然走出去.
仍然去做他能夠.
覺得做得到的事.
結果就見到.
一幅一幅.
就是受祝福.
有喜樂.
歡笑的圖案.
到最後.
一個總結的.
今天的意義.
所以我想說的是.
人的作為.
和上帝作為的關係.
大家不要輕看我們的作為.
如果你是一個殺種者.
我們有自己的崗位.
有我們自己應該做的那一份.
然後相信上帝.
就會使用我們所做的事.
化為一個我們意想不到.
或者我們計劃不到的祝福.
去帶給別人.
很大的一個影響.
不要輕看我們所做的事.
甚至很皮毛.
很低微的事.
你只要堅持.

$^{1041}$好像殺種者那樣心腸.
繼續走出去做的時候.
上帝會看得見.
上帝也會使用.
所以信心的確據.
就是基於.
好像詩人一樣.
他想起上帝.
是會拯救人.
念敏人的上帝.
所以今天.
他也相信.
上帝繼續在他的光景之中.
繼續去念敏人.
去祝福人.
我們面對今天的處境.
我們怎麼自處呢?.
我們的盼望在乎哪裡呢?.
是不是在乎那位.
出人意外的作為的上帝呢?.
我們的盼望.
其實是基於那一位.
他真是實實在在.
與人同在.
而且是他曾經應許過.
一定會實現的那位上帝.
我們可能也會想到很多作為.
是很實在.
無可指責.
我們的盼望就基於.
這位上帝.
我想讀一段.
這本書叫《公義在望》.
我也帶了一些來的.
一會兒大家也可以去選購.
其中的.
《盼望在公義的神》.
的最後幾句.
作為講道的結束.
他說.

$^{1081}$當我們真心相信聖經.
和耶穌所見證的上帝.
乃是一位公義的上帝的時候.
我們就成為這黑暗世界裡的光.
這盼望的確是屬於我們的.
我們只要知取便可.
只要我們願意相信.
就可以在這個疲乏的世界裡.
為神作偉大的見證.
我們不要輕看我們.
在上帝的手中.
在上帝的角度當中.
為每一個卑微的人所作的.
都可以成為他偉大的見證.
我們一起低頭禱告.
讓我們今天藉著詩篇一百二十六篇.
我們能夠體會你.
是一位怎樣的上帝.
你的作為出人意外.
你此下的平安喜樂.
都是出人意外.
我們就學習.
作流隊殺種者的形象.
走進我們的社群.
我們的社區當中.
將你的盼望的福音.
帶到人的中間.
如果堅定我們的信心.
讓我們持守著.
這份對你的盼望.
盼望在乎你.
你這位信實.
私人作為的上帝.
求你帶領著.
紅印加裡面所有的弟兄姊妹.
為生命.
都藉著這首詩篇.
得著激勵.
得著盼望.
同心禱告,奉主名求.

$^{1121}$盼望.
我先回到這裡,不好意思.
因為我剛才最後.
因為我帶了一些書和卡給大家介紹.
這就是一個聆聽運動.
我們在這兩年裡.
做了一個動作.
就是去做一些卡.
就好像一些玩卡遊戲的卡.
上年做了一套叫做.
道聽途說卡.
是用詩篇.
不是,是用八幅的經文.
來去用圖像.
本來我們主要都是對著年輕人做的.
就是用一些圖像.
讓他們去分享生命的故事.
然後圖像.
每一個圖畫後面.
就是一個八幅的經文.
如果他們有興趣再看回圖像背後的信息.
就跟他們分享八幅的經文.
今年做了一套叫做情感詩篇.
不好意思.
剛才那張.
黃色那個.
的聆聽卡.
這個就是特別針對我們覺得今天年輕一代.
都很多.
剛才大家說的那種心情的.
就是各樣的複雜,沉重,幽悶等等.
我們就用了詩篇.
來去作為一個卡的信息.
我們第一套就用八幅.
第二套就用詩篇.
都是用圖畫.
作為一個橋樑.
跟他們分享.
然後分享完之後.
就用經文引導他們認識福音.

$^{1161}$如果大家有興趣.
可以一會兒去買.
有簡單的說明.
也有QR Code.
你可以再下載那些詳細的說明.
在卡的套裝裡面有的.
大家不用擔心.
有介紹一下怎樣使用.
我們覺得怎樣用都是最好的方法.
剛才我說的就是簡單.
用圖畫說故事.
然後就用經文.
跟他們引導.
去認識上帝的話.
另外也有一些書籍.
剛才除了這本《公義在望》.
希望可以附和一下主題.
如果大家覺得.
《盼望》這本書是很適切.
就是今天的香港的狀態.
作者曾經去過盧旺達.
九零五代的時候大屠殺.
之後那種帶給人盼望和復活的體驗.
去寫下這本書.
可能你覺得在香港今天的處境.
可能也有些對應性.
這本書是《公義在望》.
怎樣在世界中心中存著盼望.
如果大家覺得值得.
可以一會兒在外面圓了崇拜.
大家可以去購買.
另外有幾本就是.
《迎向政治的呼召》.
是中神一班老師.
從聖經和神學當中去介紹.
怎樣去教導他們.
在今天的社會當中去作見證.
也有一本叫《走向廉潔與共融》.
是講群體復和的關係.
王貴博士也是從聖經和神學當中去分享.

$^{1201}$怎樣建立群體.
中間那本就是.
我的一本著作.
這本講讀經.
我相信讀經大家都認識.
也是鼓勵大家在這段時間當中.
去好好地去裝備.
是很重要的.
我們如果沒有能夠.
更多自己時間的時候.
其實讀《屬靈書籍》.
或《聖經》.
是作為我們自己準備好自己.
預備新的挑戰.
新的事情來臨的時候.
一種我們作為信徒的生命.
向上帝負責任.
以及去回應的一個機會.
所以盼望大家.
鼓勵大家.
如果覺得適合的時候.
都可以選購.
作為你自己在一些.
《屬靈》生命上裝備的材料.
好,那我今天分享到這裡.
\newpage



\section{哥林多前書 12:1-31}
\label{sec:YoAst8SyrWc}
\textbf{2020.10.18《凝聚每分光》 哥林多前書 12\_1-31 陳志鵬傳道}
\newline
\newline
連結: \href{https://youtube.com/watch?v=YoAst8SyrWc}{\texttt{https://youtube.com/watch?v=YoAst8SyrWc}} ~~~~ 語音日期: 2020-10-22
\newline
\newline
\hyperref[sec:6ptt2Cl0aVs]{\small{< < < PREV SERMON < < <}}
~
\hyperref[sec:index_chronic]{\small{[返順時目]}}
~
\hyperref[sec:index_scriptual]{\small{[返順卷目]}}
~
\hyperref[sec:L_6ydn1BUoo]{\small{> > > NEXT SERMON > > >}}
\newline
\newline
哥林多前書 12:1-31
\newline
\begin{longtable}{cl}
\hline
\hline
章節 & 經文 (和合本修訂版)\\
\hline
12:1 & \begin{tabularx}{0.7\textwidth}{X} 弟兄們，關於屬靈的恩賜，我不願意你們不明白。 \end{tabularx} \\ \\ \relax
12:2 & \begin{tabularx}{0.7\textwidth}{X} 你們知道，你們作外邦人的時候，隨事被引誘，受了迷惑去拜不會出聲的偶像。 \end{tabularx} \\ \\ \relax
12:3 & \begin{tabularx}{0.7\textwidth}{X} 所以，我要你們知道，被神的靈感動的，沒有人會說「耶穌該受詛咒」；若不是被聖靈感動的，也沒有人能說「耶穌是主」。 \end{tabularx} \\ \\ \relax
12:4 & \begin{tabularx}{0.7\textwidth}{X} 恩賜有許多種，卻是同一位聖靈所賜。 \end{tabularx} \\ \\ \relax
12:5 & \begin{tabularx}{0.7\textwidth}{X} 事奉有許多種，卻是事奉同一位主。 \end{tabularx} \\ \\ \relax
12:6 & \begin{tabularx}{0.7\textwidth}{X} 工作有許多種，卻是同一位神在萬人中運行萬事。 \end{tabularx} \\ \\ \relax
12:7 & \begin{tabularx}{0.7\textwidth}{X} 聖靈彰顯在各人身上，是要使人得益處。 \end{tabularx} \\ \\ \relax
12:8 & \begin{tabularx}{0.7\textwidth}{X} 有人藉著聖靈領受智慧的言語；有人也靠著同一位聖靈領受知識的言語； \end{tabularx} \\ \\ \relax
12:9 & \begin{tabularx}{0.7\textwidth}{X} 又有人由同一位聖靈領受信心；還有人由同一位聖靈領受醫病的恩賜； \end{tabularx} \\ \\ \relax
12:10 & \begin{tabularx}{0.7\textwidth}{X} 又有人能行異能，又有人能作先知，又有人能辨別諸靈，又有人能說方言，又有人能翻方言。 \end{tabularx} \\ \\ \relax
12:11 & \begin{tabularx}{0.7\textwidth}{X} 這一切都是由惟一的、同一位聖靈所運行，隨著自己的旨意分給各人的。 \end{tabularx} \\ \\ \relax
12:12 & \begin{tabularx}{0.7\textwidth}{X} 就如身體是一個，卻有許多肢體，身體的肢體雖多，仍是一個身體；基督也是這樣。 \end{tabularx} \\ \\ \relax
12:13 & \begin{tabularx}{0.7\textwidth}{X} 我們無論是猶太人是希臘人，是為奴的是自主的，都從一位聖靈受洗成了一個身體，並且共享這位聖靈。 \end{tabularx} \\ \\ \relax
12:14 & \begin{tabularx}{0.7\textwidth}{X} 身體原不只是一個肢體，而是許多肢體。 \end{tabularx} \\ \\ \relax
12:15 & \begin{tabularx}{0.7\textwidth}{X} 假如腳說：「我不是手，所以不屬於身體」，它不能因此就不屬於身體。 \end{tabularx} \\ \\ \relax
12:16 & \begin{tabularx}{0.7\textwidth}{X} 假如耳朵說：「我不是眼睛，所以不屬於身體」，它也不能因此就不屬於身體。 \end{tabularx} \\ \\ \relax
12:17 & \begin{tabularx}{0.7\textwidth}{X} 假如全身是眼睛，聽覺在哪裡呢？假如全身是耳朵，嗅覺在哪裡呢？ \end{tabularx} \\ \\ \relax
12:18 & \begin{tabularx}{0.7\textwidth}{X} 但現在神隨自己的意思把肢體一一安置在身體上了。 \end{tabularx} \\ \\ \relax
12:19 & \begin{tabularx}{0.7\textwidth}{X} 假如全都是一個肢體，身體在哪裡呢？ \end{tabularx} \\ \\ \relax
12:20 & \begin{tabularx}{0.7\textwidth}{X} 但現在肢體雖多，身體還是一個。 \end{tabularx} \\ \\ \relax
12:21 & \begin{tabularx}{0.7\textwidth}{X} 眼睛不能對手說：「我用不著你。」頭也不能對腳說：「我用不著你。」 \end{tabularx} \\ \\ \relax
12:22 & \begin{tabularx}{0.7\textwidth}{X} 不但如此，身上的肢體，人以為軟弱的，更是不可缺少的； \end{tabularx} \\ \\ \relax
12:23 & \begin{tabularx}{0.7\textwidth}{X} 身上的肢體，我們認為不體面的，越發給它加上體面；我們不雅觀的，越發裝飾得雅觀。 \end{tabularx} \\ \\ \relax
12:24 & \begin{tabularx}{0.7\textwidth}{X} 我們雅觀的肢體自然用不著裝飾；但神配搭這身子，把加倍的體面給那有缺欠的肢體， \end{tabularx} \\ \\ \relax
12:25 & \begin{tabularx}{0.7\textwidth}{X} 免得身體不協調，總要肢體彼此照顧。 \end{tabularx} \\ \\ \relax
12:26 & \begin{tabularx}{0.7\textwidth}{X} 假如一個肢體受苦，所有的肢體就一同受苦；假如一個肢體得光榮，所有的肢體就一同快樂。 \end{tabularx} \\ \\ \relax
12:27 & \begin{tabularx}{0.7\textwidth}{X} 你們是基督的身體，並且各自都是肢體。 \end{tabularx} \\ \\ \relax
12:28 & \begin{tabularx}{0.7\textwidth}{X} 神在教會所設立的：第一是使徒；第二是先知；第三是教師；其次是行異能的；再次是醫病的恩賜，幫助人的，治理事的，說方言的。 \end{tabularx} \\ \\ \relax
12:29 & \begin{tabularx}{0.7\textwidth}{X} 難道個個都是使徒嗎？難道個個都是先知嗎？難道個個都是教師嗎？難道個個都是行異能的嗎？ \end{tabularx} \\ \\ \relax
12:30 & \begin{tabularx}{0.7\textwidth}{X} 難道個個都是有醫病的恩賜嗎？難道個個都是說方言的嗎？難道個個都是翻方言的嗎？ \end{tabularx} \\ \\ \relax
12:31 & \begin{tabularx}{0.7\textwidth}{X} 你們要追求那更大的恩賜。我現今把最妙的道指示你們。 \end{tabularx} \\ \\
[1ex]
\hline
\hline
\end{longtable}
$^{1}$很高興能夠再次在港台.
跟大家分享聖經.
教會恢復現場崇拜已經有三個主日.
很欣賞弟兄姊妹的投入.
其實看到很多弟兄姊妹急不及待.
都要來到教會現場參與崇拜.
很欣賞大家.
不過我們也體諒一些弟兄姊妹.
因為種種原因.
他們未能夠回來.
我們也期望在未來的日子.
如果在政府更加放寬的情況下.
我們有更多的弟兄姊妹能夠回來.
在教會當中.
在禮堂當中.
我們一同去敬拜.
求主幫助我們.
讓弟兄姊妹有通達的道路.
如果說到.
結出年青華人的音樂家.
就是彈琴的鋼琴家.
你會想起誰呢?.
朗朗.
朗朗?.
朗朗是別人的兒子.
是朗朗.
可能我年紀比較大.
我想起李雲迪.
大家認不認識李雲迪?.
原來我翻查過資料.
李雲迪和朗朗差不多年紀.
甚至是同年紀.
但不知為何.
我總覺得李雲迪好像很早出道.
問大家一個問題.
麻煩powerpoint的弟兄姊妹.
這個問題就是.
一個標準的鋼琴.
究竟有多少個黑白鍵呢?.
給你四個選擇.

$^{41}$A.
99.
B.
88.
C.
77.
D.
66.
C.
多少呢?.
C.
根據一些人說.
通常選擇題.
A就一定不是.
最後的也未必是.
中間的兩個就很接近.
有些做multiple choice的朋友.
就會有這樣的經驗.
選擇A的舉手.
有弟兄姊妹.
選擇B的舉手.
有幾位弟兄姊妹.
選擇C的舉手.
C的舉手比較多.
難道做multiple choice.
C是最容易中的嗎?.
選擇D的舉手.
有哪位弟兄姊妹?.
也有些弟兄姊妹.
不同的數目也有弟兄姊妹去猜.
究竟有多少呢?.
答案就是B的.
有88個琴鍵.
在一個標準鋼琴裡.
黑白鍵加起來有88個.
黑鍵有36個.
白鍵有52個.
根據有限的認知.
一部鋼琴的琴鍵.
會有一條至三條的弦線在當中.

$^{81}$弦線有些粗些.
粗些的時候.
奏出來的聲音會比較沉.
長些,粗些的話.
奏出來的聲音會比較沉.
會比較低音些.
但短些,幼些的.
奏出來的聲音會比較高音些.
其實每一個琴鍵.
都是調校到不同的音階.
不同的音階.
每一個琴鍵都有不同的音階.
但如果每一個琴鍵.
都調校到相同的音階.
你猜效果會是怎樣呢?.
看一條短片給大家.
這條短片會告訴大家.
《小鵬》.
大家喜歡哪一個琴鍵的音調呢?.
我相信是第一個吧.
不可能喜歡一個.
「嘟嘟嘟嘟嘟嘟嘟嘟嘟嘟」的音調的琴.
第一個琴鍵就是剛才所說的.
每一個鍵都有不同的音階.
第二個琴鍵就是調校到相同的音階.
所以出來的效果.
沒有第一個琴鍵那麼理想.
這個短片是Android公司的宣傳片.
我留意到最後有一個口號.
口號是.
「Be together, not the same」.
這個口號我覺得很有意思.
正正是回應今天我想和大家分享的經文.
就是哥林多前書第十二章的訊息.
在過去三堂的訊息.
我和大家分享了聖經裡面關於教會的真理.
在七月份的時候.
我和大家分享過教會的根基.
我們說彼得和耶穌說.
你是基督你是永生神的兒子.

$^{121}$教會的根基就是建立在彼得的認信之上.
然後在八月份.
我和大家分享到教會的使命.
我們用彼得前書的二章九節.
去講教會的使命.
教會的使命是什麼呢?.
教會的使命就是在黑暗的世代當中.
作明亮的燈台.
在上一個月.
我和大家講教會發展的策略.
我就用馬太福音的第28章第16至20節和大家分享.
我們說教會發展的策略.
又或者這樣說.
教會發展的基礎策略.
就是在門徒訓練當中.
教會的發展有很多策略.
但基礎就在門徒訓練當中.
今天我想和大家分享教會另一個真理.
這個真理就是有關教會的成長.
我會選用哥林多前書第12章經文和大家分享.
哥林多前書是保羅寫給哥林多教會的一封書信.
哥林多教會是保羅在第二次宣教旅程親自建立的一間教會.
當時大概是主後的50年至52年之間.
保羅寫信給哥林多教會是在他第三次宣教旅程當中.
當時是主後的53至56年.
保羅寫哥林多前書的時候.
他身處的地方就是在爾忽所這個地方.
保羅為什麼要寫哥林多前書呢?.
有多個因素促使保羅寫了一本長達16章的聖經給哥林多教會.
剛剛提到保羅在第三次宣教旅程中停留在爾忽所.
當時在爾忽所保羅做宣教工作的一位人士來探望他.
這個人叫做格羅氏.
格羅氏的家人到爾忽所探望保羅的時候.
把哥林多教會裡面的一些實況告訴保羅.
他告訴保羅在哥林多教會裡面有分黨,結黨,有紛爭.
於是保羅就馬上寫了哥林多前書的第一章到第四章.
保羅準備找人將這封信帶到哥林多.
但在帶之前又有一個人來探望保羅.
這個人是誰呢?.
這個是哥林多教會所猜出來的一個使者.

$^{161}$原來哥林多教會的人有些問題想問保羅.
於是他就寫了一封信給保羅.
叫一個人帶去給保羅.
這個人帶了一封信給保羅之後.
也跟保羅說了哥林多教會裡面的一些問題.
於是保羅就馬上聽到這些問題之後.
就馬上寫了哥林多前書的第五章和第六章.
接著打後的經文.
哥林多前書的第七章到第十五章.
就是保羅回覆哥林多教會寄來的這封書信.
在當中保羅處理了很多的問題.
所以今天我和大家去看的哥林多前書的第十二章.
就是保羅回信的其中一個部分.
究竟保羅在一張聖經要去處理哥林多教會什麼問題呢?.
我們留意經文的第一節.
在程序表的經文裡面的第一節我們留意.
他們說弟兄們論到屬靈的恩賜.
我不願意你們不明白.
原來保羅要回應哥林多教會有關屬靈恩賜的問題.
保羅回應這個問題不是單單在第十二章裡面.
保羅回應屬靈恩賜的問題.
由第十二章一直說到第十四章.
就是說保羅用了三張聖經.
去回應哥林多教會信問有關屬靈恩賜的問題.
什麼是恩賜呢?.
恩賜是什麼呢?.
在希臘文當中有兩個字.
在中文聖經華本裡面同樣是翻譯成恩賜.
其中一個字是強調屬靈恩賜的本質和來源.
屬靈恩賜不是一種天賦的才能.
屬靈恩賜是聖靈用超自然的方式賜給每一個基督徒.
正如在第十一節那裡說.
這一切都是這位聖靈所運行.
隨己而分給各人.
另外一個希臘字翻譯成為恩賜的.
它強調屬靈恩賜是上帝恩典的次語.
就是說恩賜就是上帝給每一個基督徒的禮物.
究竟上帝把恩賜給每一個基督徒的目的是什麼呢?.
我們留心第七節的經文.
第七節說.

$^{201}$聖靈顯在各人身上是叫人得益處.
原來上帝把恩賜賜給每一個基督徒.
不是要他拿了恩賜之後自我感覺良好.
也不是讓基督徒拿了恩賜之後.
去炫耀自己有多麼高超.
有多麼屬靈.
上帝把恩賜交給教會.
交給每一個基督徒.
是有個期望.
期望他們能夠去建立教會.
用上帝所給他們的恩賜去建立教會.
這個就是上帝賜予恩賜的目的.
簡單來說恩賜是什麼呢?.
恩賜不是與生俱來的才能和能力.
是人接受了耶穌基督之後.
上帝藉著聖靈賜給每一個基督徒.
目的就是要建立教會.
而每一個基督徒所接受的恩賜是不同的.
每一個基督徒所接受的恩賜都是很獨特的.
正如第四節那裡這樣說.
「恩賜原有分別,聖靈卻是一位」.
弟兄姊妹,你有什麼恩賜?.
又或者你知不知道你自己有什麼恩賜?.
可能你會說,我不是很清楚.
其實坦白說,我問過很多弟兄姊妹.
你知不知道你自己有什麼恩賜?.
弟兄姊妹往往的回覆就是.
「我不是很清楚」.
我想很多時候弟兄姊妹未必不是不知道.
不過我們弟兄姊妹習慣了比較謙虛.
不敢覺得如果說出來.
「我有這方面的恩賜」.
就好像很「駕世」.
所以弟兄姊妹寧願去內斂一點.
「我都不是很清楚」.
但其實我們是需要知道.
上帝給了我們什麼恩賜.
因為如果我們不知道上帝給了我們什麼恩賜.
我們怎樣建立教會?.
我們要建立教會.

$^{241}$我們要知道上帝給了我們什麼恩賜.
以致我能夠在教會裡.
去用出來,去建立整間教會.
所以每一個的弟兄姊妹.
你必須要知道自己有什麼恩賜.
上帝給了你什麼恩賜.
有一篇文章記載.
有一個人去問葛培里牧師.
葛培里牧師早幾年去世.
這個人去問葛培里牧師.
「我怎樣才知道自己的恩賜?」.
葛培里牧師回覆了四點.
第一點他說.
你要認定上帝至少給了你一個恩賜.
而你要相信.
你要知道這些恩賜是要運用出來的.
為了上帝的榮耀.
你的恩賜是要運用出來的.
第一點就是說.
你認定上帝至少給了你一種的恩賜.
第二點他說.
要相信經過一個仔細的禱告.
你才能夠知道上帝給了你什麼恩賜.
又或者你要經過一個很詳細的禱告.
你才能夠發掘到自己的恩賜.
我們要懇求上帝去引導我們.
讓我們知道他給了我們什麼恩賜.
第三點就是.
我們要清楚聖經裡面所說的恩賜是一件什麼事情.
任何的恩賜都不可以離開真理的框架.
就是說上帝將恩賜給過我們.
是在真理裡面的.
我們不可以在真理以外的地方去尋找恩賜.
恩賜必須要在真理的框架裡面.
第四點.
最後郭牧師回覆這個人所說的.
你要透過你的事奉去發掘你的恩賜.
又或者你去向你的屬靈的長輩.
或者教會的牧者去查詢你究竟有什麼恩賜.
其實作為教會的傳道同工.

$^{281}$我們其實有責任去幫助大家.
去明白上帝給了你什麼恩賜.
在《耳窟所書》第四章七到十一節裡說.
為要成全聖徒各盡其職.
建立基督的身體.
這是每一個傳道同工都需要去做的.
就是要裝備大家.
讓大家知道自己有什麼恩賜.
以至到教會能夠去建立.
所以弟兄姊妹.
如果有傳道同工邀請你去事奉的時候.
你的回應應該是怎樣.
好 我就去.
既然你叫我去做 我就去做.
因為你在透過你事奉的過程當中.
你就知道自己的恩賜.
教會成長的關鍵就在於.
我們有沒有運用上帝給我們的恩賜.
如果沒有上帝的帶領.
我相信沒有一個人能夠來到教會當中.
就是說你來到紅恩堂裡面.
有上帝的帶領.
如果沒有上帝的心意.
也沒有人離開這間教會.
弟兄姊妹.
你在紅恩堂當中聚會.
上帝將你帶領來到這裡.
上帝也將恩賜交給你.
你就應該在教會裡面.
將上帝給你的恩賜發揮出來.
剛才跟大家說.
怎樣去發掘恩賜.
如果你真的不知道自己有甚麼恩賜的時候.
你就按剛才我所分享的四方面去發掘.
說到恩賜的發掘.
我也有一些體會.
我不知道大家弟兄姊妹覺得.
我唱詩歌的情況是怎樣.
我不是自我感覺良好.
不過我自己覺得.

$^{321}$唱詩歌也算是勉強合格.
不是很難聽的那種.
但是我可以告訴大家.
在我還沒有在神學院畢業.
在神學院的四年裡.
從來沒有同學邀請我帶敬拜.
我只是經常坐在下面很渴望.
我也很想學習一下帶敬拜.
有沒有這樣的機會.
沒有.
神學院裡面有很多奇人異士.
他們很厲害.
未輪到你.
我只是坐在那裡.
加上我覺得自己唱歌也不太好.
到了教會的事奉.
我去到一間年輕人居多的教會.
他們很多弟兄姊妹.
音樂的恩賜是很厲害的.
上帝給了他們很多音樂的恩賜.
有些弟兄姊妹自學彈琴.
有些弟兄姊妹自學彈結他.
很多弟兄姊妹唱歌很好聽.
我在當中帶敬拜的機會不是很多.
不過因為你是傳道同工.
所以教會也給你機會帶敬拜.
所以我在當中學習敬拜.
在過程當中.
我發現自己開始覺得.
自己唱詩歌方面有上帝而來的恩賜.
來到紅印堂事奉.
我又發覺上帝不斷將音樂.
或者唱詩歌的恩賜交給我.
以致到今天我唱詩歌.
比起之前那間教會的事奉.
是比較理想一點的.
不過不說不知道.
有一樣東西很特別.
大家有時也會去唱卡拉OK.
有吧.

$^{361}$我發覺當我去唱卡拉OK的時候.
我是唱不到的.
我發覺唱得很難聽.
為什麼平時不是這樣.
平時唱詩歌是可以的.
唱詩歌OK就代表唱歌OK.
唱卡拉OK應該是OK的.
但原來不是.
當我去唱卡拉OK的時候.
我發覺好像舌頭打結.
又不聽又走音.
高音又上不到.
幸好在房間裡.
否則在大庭廣眾就很難聽.
在我的年代.
唱卡拉OK有些是在大堂唱的.
不知道大家知不知道.
如果是這樣就真的很難聽.
但回到教會唱詩歌.
又OK的.
又不覺得難聽.
自此之後我體會到.
當你願意去事奉的時候.
上帝會將相應的恩賜給你.
如果在你事奉的過程當中.
你發覺好像一般般.
旁人亦有一些意見給你的時候.
可能你會清楚地知道.
原來在這方面的恩賜.
我未必是那麼強.
你就要向另一方面去發展.
求主幫助我們.
教會要成長.
我們就要運用上帝給我們的恩賜.
去建立現在的教會.
保羅在講述屬靈恩賜的時間.
在十二節到二十六節.
用了一個比喻去講述屬靈恩賜.
這個比喻我相信大家都很熟悉.
就是身體的比喻.

$^{401}$或者我們在下一張影片裡面看看.
這個身子的比喻.
在這個身子的比喻裡面.
大概可以分開兩段.
第一段就是第十二節到第十三節.
這裡講到一個身子眾多的肢體.
然後就是第十四節到二十六節.
就是肢體彼此相顧.
第十二到第十三節.
保羅就說.
就如身子是一個.
卻有許多的肢體.
而且肢體衰多仍是一個身子.
基督也是這樣.
之後在第十三節.
保羅用了當時哥林多教會裡面.
這個群體的成員.
去講述身子只有一個.
卻有許多的肢體.
根據一些資料.
保羅時代哥林多教會裡面.
有些甚麼的弟兄姊妹呢.
原來都有不同階層的弟兄姊妹在當中.
在哥林多教會裡面.
有一些叫做自由人.
有一些叫做平民百姓.
也有一些比較富有的.
簡單來說.
在哥林多教會裡面.
真的甚麼人都有.
不過這句說話要小心一點說.
教會甚麼人都有.
看你用甚麼語氣說.
我記得在我舞會事奉的時候.
有一位傳道童工報告的時候.
不知為何突然間彈了一句出來.
可能他的語氣比較輕佻.
教會真的甚麼人都有.
帶出來的效果比較負面.
但這句說話是沒有錯的.

$^{441}$教會真的甚麼人都有.
在哥林多教會裡面.
我們都可以視猶太人,希利人為奴的,自主的.
即是說在教會裡面.
有很多種不同的弟兄姊妹.
不過縱然有不同的弟兄姊妹.
保羅怎樣說.
都從一位聖靈受洗.
成了一個身體.
飲於一位聖靈.
甚麼是飲於一位聖靈呢.
可能在中文的翻譯當中.
我們比較難去明白甚麼是飲於一位聖靈.
但如果我從另一個字去表達.
大家可能就會明白.
飲這個字有囂觀的意思.
即是說我們被一個聖靈囂觀.
有很多別的宗派.
解釋這節聖經的時候.
從一位聖靈受洗.
就是我們受洗的時候.
然後後面的飲於一位聖靈.
就是聖靈的囂觀,聖靈的充滿.
所以因著聖靈的充滿.
就有聖靈充滿的外顯出來.
但是保羅在這裡所說的.
從一位聖靈受洗.
飲於一位聖靈.
其實是在說同一件事.
即是說我們雖然有不同的人種.
我們有不同階層的人士.
組成這間教會.
但是我們都是從一位聖靈受洗.
建立同一個身體的底下.
成為一個身子.
保羅就用了身子這個比喻.
去解釋屬靈恩賜.
即是說恩賜有很多.
但是這麼多的恩賜.
都是從一位聖靈而來.

$^{481}$接著保羅在第十四到二十六節.
就說到肢體與肢體之間的關係.
究竟是怎樣.
不過很有趣.
保羅在第十五到第十六節.
是這樣說的.
切若腳說:我不是手,所以不屬乎身子.
他不能因此就不屬乎身子.
切若耳說:我不是眼,所以不屬乎身子.
他不能因此就不屬乎身子.
有趣的地方就是.
為何腳會說:我不是手,所以我不是身子.
為何耳朵會說:我不是眼,所以不屬乎身子.
其實在身體裡.
不同的器官.
不同的肢體.
都有不同的功用.
我在神學院裡有一個同學.
去做學生報道的時候.
經常會與同學玩一個小把戲.
他怎樣說呢?.
他說:你試試掩著耳朵.
一邊,不是兩邊,掩著一邊耳朵.
你試試掩著一邊耳朵.
然後我在這裡說話.
你能否分辨到聲音從何而來.
似乎分辨不到.
但如果你放開了兩隻耳朵.
你能夠分辨到聲音從何而來.
他玩另一個把戲.
就是要單眼.
看看大家能否做到.
你拿出兩隻手指.
單眼.
單眼之後,看看可否.
不行.
但如果是兩隻眼.
你就可以做到這個動作.
即是手指對手指.
單眼就很容易猜錯了.

$^{521}$這裡說甚麼呢?.
我們身體每一個器官.
每一個肢體.
都有它獨特的用處.
你試試跟隔壁的人說.
如果今天要你.
拿走你身體其中一個器官.
你要拿走誰呢?.
看看隔壁的人會否回答你.
隔壁的人可能會說.
你是不是有些神經失常呢?.
我相信沒有一個人.
願意將自己身體任何一個器官.
去撤除了.
因為在我們身體上.
每一個肢體.
每一個器官.
都有它的用處.
就算很細微的器官.
或者肢體.
都很有用處.
有人說.
當你沒有腳趾尾.
你走路是很有困難的.
站著也很有困難.
信不信?.
不知道?.
你今晚試試.
把腳趾尾撤掉.
看看你的情況是怎樣.
但不要說我教你的.
因為是你自己做的.
身體上每一個器官.
每一個肢體.
都是很獨特.
都有它的用處.
保羅用了身子這個比喻.
去講述屬靈恩賜.
屬靈恩賜有不同.
上帝把.

$^{561}$藉著聖靈.
把恩賜給每一個弟兄姊妹.
每一個弟兄姊妹.
所領受的恩賜都不同.
但是.
他們就在同一個教會裡面.
把上帝給他們的恩賜.
去發揮出來.
這個就是保羅.
去講解屬靈恩賜的時候.
所用的比喻.
不過有趣的地方是.
保羅去講.
用身子的比喻.
去講屬靈恩賜的時候.
帶出了另一個的真理.
這個真理就是教會的合一.
教會的合一.
保羅說.
在整個的身子裡面.
有很多不同的肢體.
而肢體與肢體之間.
是不可以排斥的.
在第25節裡面怎樣說.
免得身上分門別類.
總要肢體彼此相顧.
可能中文聖經的翻譯.
做分門別類.
是比較斯文了一點.
原來的意思.
分門別類是講撕裂.
是講撕裂.
就是說如果一個的恩賜.
或者身體上的肢體.
是彼此排斥的時候.
就會造成撕裂.
所以保羅說.
肢體與肢體之間.
總要彼此相顧.
免得彼此撕裂.

$^{601}$我想撕裂這個形容詞.
在今時今日.
大家都會相當有很深的體會.
在教會裡面.
受著過去的政治的風暴.
或者社會的氣氛.
教會裡面.
也有一種撕裂的現象出現.
於是乎有一些人提出.
分析牧養.
不知道大家知不知道甚麼是分析牧養.
就是說不同政見的人.
分開一堆是政見的.
牧養他.
另一堆是另一個政見的.
牧養他.
基本上我不太認同這個做法.
不過我明白有它的處境性.
為甚麼我不認同它的做法呢.
在哥林多前書第十二章裡說.
身子只有一個.
卻有許多的肢體.
肢體與肢體之間.
剛才我說過.
是不可以排斥的.
為甚麼不同政見的立場.
不可以在同一間教會裡面出現呢.
我到現在都不太明白這件事.
在我牧養的團契裡面.
其實有不同的頂子妹.
有他自己的意見.
但是他們相處上.
都是非常和諧.
都很有聖徒的.
各種彼此相顧的表現出現.
頂子妹.
不要因著社會的一些氣氛.
又或者因著社會的一些政見.
而將教會撕裂.
一樣是上帝所不起願.

$^{641}$在保羅說的屬靈恩賜的時候.
他用生子去做一個比喻.
然後帶出教會合一的真理.
究竟教會的合一是在說甚麼呢.
這張圖很特別.
大家有沒有玩過砌圖.
砌圖有不同的細塊.
可能你說有些形狀是一樣的.
是沒錯的.
它的形狀是一樣的.
但是顏色是不同的.
或者它裡面的圖案是不同的.
但是當不同的圖案.
按照著設計圖去合在一起的時候.
就能夠呈現出一幅很漂亮的圖畫.
我的小朋友也有玩砌圖的習慣.
有個親戚送了一幅一千塊的砌圖.
他和太太一起去砌圖.
可能他很喜歡卡通人物多拉蒙.
所以用了很短的時間.
他們兩個砌了這幅砌圖.
砌圖是由不同的小塊去組成.
每個小塊都有不同的圖案.
但是砌成一幅很美麗的圖畫.
這就是教會的合一.
教會的合一就是.
如果用英文去說就是Unity in Diversity.
就是說在合一裡面有它的多元性.
教會有它的多元性.
教會有不同的弟兄姊妹.
不同階層的弟兄姊妹.
不同恩賜的弟兄姊妹.
在同一間教會.
但是弟兄姊妹與弟兄姊妹之間.
彼此相顧合一在一間教會當中.
這才是上帝心意當中的教會.
教會的合一除了Unity in Diversity之外.
有另一句說話我也想給大家聽聽.
這句說話其實我翻查了很久.
我也不知道究竟誰是第一個去說這番話.

$^{681}$在西方教會很流行這句說話.
它說.
我用中文說.
在機要的事上要合一.
在非機要的事上要有自由.
在一切的事情上要有愛心.
我想這裡大概是說.
教會要合一就需要有愛心.
保羅用哥林多前書第十二章去說.
屬靈恩賜.
他用身子做一個比喻.
身子裡面的肢體不能夠彼此排斥.
恩賜在教會裡面的運用.
必須要在基督的愛裡面運用.
有一節的聖經.
我相信兩節的聖經大家都很熟悉.
在二弗所書的第四章第二到第三節.
詩歌也有唱.
凡事謙虛溫柔忍耐.
用愛心互相寬容.
用和平彼此聯絡.
竭力保守聖靈所賜.
合而為一的心.
這個就是保羅在二弗所書裡面說.
合一已經賜給我們.
我們要做的就是.
竭力保守聖靈所賜的合而為一的心.
我們很知道這節經文說什麼.
凡事謙虛溫柔忍耐.
用愛心互相寬容.
用和平彼此聯絡.
竭力保守聖靈所賜合而為一的心.
有些弟兄姊妹背都背得出來.
我們很清楚知道.
要合一就是凡事謙虛溫柔忍耐.
但什麼破壞教會的合一呢.
如果你將這條題目問一些小學生或中學生.
他們很快可以回答你.
知道是怎樣嗎.
在前面加個筆字就行了.

$^{721}$凡事不謙虛不溫柔不忍耐.
不用愛心互相寬容.
不用和平彼此聯絡.
就做不到合一的心.
沒錯.
只要我們不謙虛不溫柔不忍耐.
不用愛心互相寬容.
不用和平彼此聯絡.
我們就破壞了上帝所賜合而為一的心.
曾經擔任第52屆港九培靈議經大會的.
其中一位港議員薛若光先生.
他在他的講道裡曾經這樣說過.
他說什麼事情會妨礙教會的合一呢.
就是當我們不高舉主.
只高舉人的時候.
合一就會受到破壞.
保羅,翟比,哥林多教會分黨就是其中一個例子.
當人去高舉我是屬於誰的.
我是屬於誰的.
我是屬於阿波羅的.
我是屬於保羅的.
我是屬於基發的時候.
教會的合一就會破壞.
因為有人高舉了自己.
不高舉主耶穌.
在教會裡面.
凡是有人只是高舉自己.
而不高舉主耶穌.
就是正在破壞教會的合一.
所以弟兄姊妹.
在教會裡面.
千萬不要高舉自己.
你要高舉的就是主耶穌.
有什麼不是領受回來的.
有什麼是需要可誇的呢.
我們只誇主耶穌基督並他的十字架.
求主幫助我們.
狄根斯曾經說過一句話很有意思.
他怎樣說呢.
教會無非是一個聚人.

$^{761}$帶領著一群的聚人.
什麼意思呢.
就是說我們做帶領的.
又或者信徒領袖做帶領的.
其實沒有什麼可誇的.
我們跟弟兄姊妹一模一樣.
我們都是領受同一個的聖力.
我們都是同受同一個恩典.
所以在教會裡面.
求主幫助我們.
當我們說要去合一.
我們真的需要高舉主耶穌基督.
求主幫助我們.
我們有個祈禱.
親愛的主好多謝你.
多謝你透過保羅.
去回覆哥林多教會.
關於屬靈恩賜的問題.
帶出一個教會的真理.
就是教會必須要合一.
因為教會是在基督的身子裡面.
我們互為之體.
我們要彼此相顧.
免得彼此的撕裂.
免得分門別類.
求主幫助我們.
當我們要在疫情過後.
繼續向前發展的時候.
求主幫助我們.
凡事謙虛 溫柔 忍耐.
用愛心互相寬容.
用和平彼此聯絡.
竭力保守你所賜下合二為一的心.
多謝主你垂聽我們在你面前的禱告.
祈禱奉耶穌基督的聖名.
阿們.
\newpage



\section{腓立比書 3:7-14-20-21}
\label{sec:L_6ydn1BUoo}
\textbf{2020.10.25 講道《我要得甚麼的獎賞》 腓立比書 3\_7-14,20-21( 和修版)  老耀雄傳道}
\newline
\newline
連結: \href{https://youtube.com/watch?v=L-6ydn1BUoo}{\texttt{https://youtube.com/watch?v=L-6ydn1BUoo}} ~~~~ 語音日期: 2020-10-28
\newline
\newline
\hyperref[sec:YoAst8SyrWc]{\small{< < < PREV SERMON < < <}}
~
\hyperref[sec:index_chronic]{\small{[返順時目]}}
~
\hyperref[sec:index_scriptual]{\small{[返順卷目]}}
~
\hyperref[sec:xDCCx3kJCvo]{\small{> > > NEXT SERMON > > >}}
\newline
\newline
腓立比書 3:7-14-20-21
\newline
\begin{longtable}{cl}
\hline
\hline
章節 & 經文 (和合本修訂版)\\
\hline
3:7 & \begin{tabularx}{0.7\textwidth}{X} 只是我先前以為對我是有益的，我現在因基督的緣故而當作是有損的。 \end{tabularx} \\ \\ \relax
3:8 & \begin{tabularx}{0.7\textwidth}{X} 不但如此，我已把萬事當作是有損的，因我以認識我主基督耶穌為至寶。我為他已經丟棄萬事，看作糞土，為要贏得基督， \end{tabularx} \\ \\ \relax
3:9 & \begin{tabularx}{0.7\textwidth}{X} 並且得以在他裡面，不是有自己因律法而得的義，而是有信基督的義，就是基於信，從神而來的義， \end{tabularx} \\ \\ \relax
3:10 & \begin{tabularx}{0.7\textwidth}{X} 使我認識基督，知道他復活的大能，並且知道和他一同受苦，效法他的死， \end{tabularx} \\ \\ \relax
3:11 & \begin{tabularx}{0.7\textwidth}{X} 或許我也得以從死人中復活。 \end{tabularx} \\ \\ \relax
3:12 & \begin{tabularx}{0.7\textwidth}{X} 這不是說我已經得著了，已經完全了；而是竭力追求，或許可以得著基督耶穌所要我得著的。 \end{tabularx} \\ \\ \relax
3:13 & \begin{tabularx}{0.7\textwidth}{X} 弟兄們，我不是以為自己已經得著了；我只有一件事，就是忘記背後，努力面前的， \end{tabularx} \\ \\ \relax
3:14 & \begin{tabularx}{0.7\textwidth}{X} 向著標竿直跑，要得神在基督耶穌裡從上面召我來得的獎賞。 \end{tabularx} \\ \\ \relax
3:15 & \begin{tabularx}{0.7\textwidth}{X} 所以，我們中間凡是成熟的人，總要存這樣的心；若在甚麼事上存別樣的心，神也會把這些事指示你們。 \end{tabularx} \\ \\ \relax
3:16 & \begin{tabularx}{0.7\textwidth}{X} 然而，我們達到甚麼地步，就當照這個地步行。 \end{tabularx} \\ \\ \relax
3:17 & \begin{tabularx}{0.7\textwidth}{X} 弟兄們，你們要一同效法我，也當留意看那些效法我們榜樣的人。 \end{tabularx} \\ \\ \relax
3:18 & \begin{tabularx}{0.7\textwidth}{X} 因為，我屢次告訴你們，現在又流淚告訴你們：許多人行事是基督十字架的仇敵。 \end{tabularx} \\ \\ \relax
3:19 & \begin{tabularx}{0.7\textwidth}{X} 他們的結局就是滅亡。他們的神明是自己的肚腹；他們以自己的羞辱為光榮，專以地上的事為念。 \end{tabularx} \\ \\ \relax
3:20 & \begin{tabularx}{0.7\textwidth}{X} 我們卻是天上的國民，並且等候救主，就是主耶穌基督從天上降臨。 \end{tabularx} \\ \\ \relax
3:21 & \begin{tabularx}{0.7\textwidth}{X} 他要按著那能使萬有歸服自己的大能，把我們這卑賤的身體改變形狀，和他自己榮耀的身體相似。 \end{tabularx} \\ \\
[1ex]
\hline
\hline
\end{longtable}
$^{1}$洪影家弟兄姊妹平安.
在錦宣家服侍的時候.
我也有一個習慣.
很多時候每一年.
我都和部委會去做退休.
當然做退休一定有一種檢視的作用.
去看看我們那個部.
在過往一年.
究竟做得好不好.
固然這種是很事工化.
但是退休的下半部.
我們就會檢討.
究竟我們的事奉.
究竟事奉官是不是正確.
事奉固然要講果孝.
但是事奉本身的事奉人員的心態.
更加重要.
我們要看看我們和人相處的態度.
會不會有需要反省一下.
會不會和人家的關係.
這個事奉是親密了.
還是疏離了.
是什麼影響到我們事奉的態度.
這個都是很多時候.
在我以往在錦宣服侍的時候.
那個退休會的一個主題.
就是怎樣能夠可以.
讓我們那個事奉.
能夠帶到一個屬靈生命的成長.
在今年年頭.
聖靈對我一個很大的提醒.
作為牧者.
我們當然很希望.
我們的會友.
是可以有屬靈生命的成長.
很想他們每一天都有靈修的習慣.
但是怎樣可以讓那些會友.
能夠有靈修呢.
聖靈對我一個很大的提醒就是.
如果你不以身作則.

$^{41}$你又怎能夠讓你的會友.
都可以這樣做呢.
所以我在年頭三月的時候.
我就和我當時.
錦宣團契的會友說.
不如我將我每天的靈修.
寫成一種札記.
然後轉發給他們.
雖然對於他們.
這些不是親自的領受.
都是一些易受的材料.
但是我覺得都是一種.
屬靈的反省.
透過一個生命的反思.
或許對他們信仰有幫助.
而最重要最重要的是.
在那個靈修.
或信仰的歷程裡面.
我願意和他們一起走.
讓他們知道.
每天靈修既然這麼重要的話.
牧者和他們一起走.
他們是不孤單的.
因為我會和他們一起走.
這條信仰的路.
而我過往的靈修札記.
現在也都上載了.
在教會的網頁裡面.
盼望大家都可以和我一起.
在一個屬靈的信仰路上.
一起走.
而這種信仰經歷的回顧.
展望或者實踐.
亦都是基督宗教的傳統.
史陶保羅就是我們的先導者.
他在《彼勒比書》.
向當代的信徒說.
忘記背後努力面前.
向著標竿直跑.
這一種就是一種回顧過去.

$^{81}$檢視現在.
展望將來的信仰的重整.
今天我很想透過.
《彼勒比書》這三章.
七至十四節和二十至二十一節的經文.
重溫一下.
保羅在這個信仰整理之下.
的一段心路歷程.
我們一起去祈禱.
提醒我們有聖徒相通的恩典.
透過主的恩典.
讓先賢的生命.
和我們成為一個流通的管子.
保羅所經歷的.
成為我們的提醒和勸勉.
甚至每一個聽到的.
都能夠領受這種生命的傳承.
叫我們的生命.
能夠透過保羅的見證.
讓我們不斷地成長.
最終被你使用.
榮耀主名.
我們的祈禱.
是奉主耶穌基督得勝明智祈求.
阿們.
開始之前.
我們先了解一下.
肥勒比書的神學立場.
肥勒比書有一個.
很重要的神學主題.
三一神.
基本的屬性.
完全在基督裡面.
彰顯出來了.
而肥勒比書二章一至十一節.
是描述基督論最清楚的一段.
基督救贖了我們.
要使我們有份於他的樣式.
與基督樣式有份.
也構成了肥勒比書裡面.

$^{121}$的一些分題.
譬如受苦.
譬如喜樂.
合一.
努力得到賞賞.
面對困難的反應.
都是全書的分題的一部分.
因此當我們去閱讀.
肥勒比書的時候.
就可以看到保羅的信仰觀.
他對喜樂.
生命.
和成敗得失的觀點.
是完完全全.
聯繫於主耶穌基督身上的.
基督.
就是保羅整個生命的核心價值.
也是保羅要求當代信徒.
應該要效法的功課.
保羅在三章.
十三至十四節.
對自己的信仰說了一個總結.
是涉及有關自己的過去.
現在和未來.
也是一種回顧.
以及對將來的展望.
他對於過去怎麼說呢?.
他說應該要忘記背後.
對於現在.
他就要努力面前.
對於將來.
就是向著標竿直跑.
而這種信仰的總結.
的終極目標.
是要得到從神而來的獎賞.
所以今天講到的重點.
不是講基督論.
不是講教會論.
也不是講受苦喜樂.
或者合一.

$^{161}$而是單單是講.
得獎賞.
所以講題是.
我要得什麼獎賞.
既然保羅很明確地說.
他很想要得到.
在基督耶穌裡.
從上面召我來得的獎賞.
那我們就從保羅這句話.
去了解內裡的屬靈意義和凡是.
我們開始進入經文.
保羅在第十四節這樣說.
他想要得到神的賞賜.
那就是得到沒有呢?.
十三節講到明.
他說我不是以為自己已經得著了.
所以答案是應該還沒有得到.
雖然保羅說他還沒有得到獎賞.
但是他為這份獎賞付出了什麼呢?.
七至八節這樣說.
只是我先前以為對我是有益的.
我現在因基督的緣故而當作是有損的.
不但如此.
我以把萬事當作是有損的.
因我已認識我主基督耶穌為至寶.
我為他已經刁棄萬事.
看作奮土為要贏得基督.
保羅在這裡提出了他的信仰心態.
就是以前覺得對他是有益的事.
現在覺得不但沒有益.
反而是有害.
以前保羅認為有益的.
是指什麼呢?.
就是指第五至六節.
保羅所說的兩件事.
第一件事就是他的身份.
保羅的父母給了他一個猶太人的身份.
種族上保羅是屬於以色列民.
而且保羅也受了割禮.
因此保羅從身份上.

$^{201}$成為了屬臣子民的一份子.
而保羅也曾以這個身份為榮.
亦感到自豪.
第二件事就是指保羅的行為.
保羅按照法利賽人的方式.
去進行摩西的律法.
並且全然投入其中.
所以保羅認為.
按照摩西律法上的要求來說.
他是無可指責的.
無論是屬臣子民的身份.
還是遵守摩西的行為.
保羅都認為自己已經擁有了.
因此保羅很有信心地.
他覺得終有一天.
這個身份和這個成就.
會有助於他在臣面前稱為義人.
以往保羅認為.
臣所喜悅和接納的有益條件.
現在保羅完完全全地否定了這個想法.
而且覺得是有損的.
有損即是不好.
為何覺得不好呢?.
因為保羅發現這兩個條件.
都不是臣真正閱納的因素.
臣並不單單閱納以色列人.
臣亦會閱納其他種族.
讓他們成為臣的子民.
而人的行為不可能達到完全.
因此不可能以行為去取悅臣.
而保羅過往卻緊緊地擁抱著.
不捨得放棄.
那又損害了甚麼呢?.
就是這種自以為可以得到義的態度.
亦是因為這種自以為擁抱了真理的態度.
妨礙了保羅去認識主耶穌基督.
以及要贏得基督.
保羅說忘記背後的真正意思.
就是要捨棄以往這種自以為擁有了真理.
或者自以為已經得到的信仰價值觀.

$^{241}$保羅在忘記背後之後.
他就說努力面前.
他現在面前要努力的究竟是甚麼呢?.
保羅說現在要為了贏了基督.
就將萬事都丟棄甚至視為糞土.
保羅認為世界上沒有任何一件事.
重要得過認識主耶穌基督.
以及要贏得基督.
對於認識主耶穌基督我們比較容易理解.
主耶穌基督是舊約所預言的彌賽亞.
主耶穌基督是生命的主.
也是拯救的主.
主耶穌基督是三一的真神.
主耶穌基督有神人異性的本質.
似乎關乎認識的層面都是知識性的.
那甚麼叫贏得基督呢?.
新漢語譯本以及和合本都譯為.
得著基督.
詩高譯本譯為贏得基督.
但基督怎能得著.
怎能賺得.
或怎能贏得呢?.
希臘原文的意思是.
使他屬於自己.
如果要基督屬於自己.
唯一的方法就是.
我們和基督的生命連繫.
就好像你中有我.
我中有你.
生命的連結是屬於屬靈上的體驗.
因此保羅更生了的信仰.
不單止在知識性上要認識耶穌.
還要在生命的經歷上.
要體驗和主耶穌生命的結連.
保羅為了達到這個目的.
他可以將一切都拋棄.
因此雖然說.
他仍然未得到這份獎賞.
但他已經為此而付出了代價.
這個代價就是放棄以往的信仰觀.

$^{281}$並且為了認識和以及和基督生命的結連.
他可以放棄一切.
為的是要得到基督.
洪英嘉弟兄姊妹.
當保羅說.
凋棄萬事漢作糞土.
為要贏得基督.
這是否也是我們的信仰態度.
如果贏得基督是指和基督生命的結連.
那麼我們和基督的關係.
是否只是單單停留在知性層面上.
而沒有屬靈的關係呢?.
昔日保羅更新了他的信仰觀念.
認為所有妨礙他認識和結連基督的身份和行為.
都是有害無益的.
今天當我們視為有益.
對我們有益的信仰核心和態度.
其實我們有沒有反省過?.
我們的事奉態度.
或者我們的屬靈生命.
會否成為了一種.
阻礙了我們屬靈生命成長的絆腳石呢?.
這個很值得我們每一個弟兄姊妹.
作一個深層次的反省.
保羅說向著彪竿直跑.
彪竿背後的想像.
又從何而來呢?.
十四節.
他說:要得神在基督耶穌裡.
從上面召我來得的獎賞.
原文的主句子的結構是.
要得神的賞賜.
而在耶穌基督裡從上面召我來.
這是一個形容詞子句.
用來形容獎賞的性質.
這個形容詞子句有兩個字.
我們值得留意.
第一個是「召」.
這個字可以是動詞也可以是名詞.
有呼召或者邀請的意思.

$^{321}$第二個要留意的是「上面」這個字.
這是一個形容詞.
是形容呼召這個名詞.
意思是在基督耶穌裡.
從上面對我的呼召.
新漢語將上面譯為天上.
即是從天上的呼召.
和磨修版上的意思大致相同.
意思是在基督耶穌裡.
從天上對我的呼召.
如果我們再看呂鎮忠或詩高逆本.
他將上高這個形容詞.
譯為向上.
而且還將這個形容詞的位置調了.
放在獎賞之前.
用來做獎賞.
而不是形容呼召.
意思是在基督耶穌裡.
呼召我向上去得的賞賜.
說了這麼多.
其實我們可以從上面.
重點是什麼呢?.
重點是和修版和新漢語的焦點.
是呼召從上高而來.
呂鎮忠和詩高的焦點是.
獎賞是向上高得到.
究竟兩種的譯法.
和差異的焦點有什麼不同?.
和修版強調的.
呼召從天上而來.
得獎賞.
並沒有說.
將來才可以得到.
呂鎮忠強調獎賞是從往上取得的.
反映獎賞未必是現在才可以得到.
獎賞是將來的事.
不過無論將形容詞譯為上邊.
天上或向上也好.
這個獎賞也不屬於地上.
是屬於耶穌基督.

$^{361}$分別是.
這是從天上而來的呼召.
還是得到從天上而來的獎賞.
所以如果我們的問題是.
獎賞從哪裡來的?.
我們可以說獎賞從天上而來.
但如果我們將問題說.
怎樣才可以得到天上的賞賜?.
我們就要說.
一定要達到基督耶穌從天上的呼召.
才可以得到.
因此我們就必要知道.
這個從天上而來的呼召.
內容是甚麼呢?.
第十節就說了.
這個呼召的內容就是使我們認識基督.
知道祂復活的大能.
並且知道和祂一同受苦.
效法祂的死.
這段經文不單包括了知識性的層面.
還有實踐的層面.
認識基督.
知道祂復活的大能.
知道和祂一同受苦.
這種認識和知道.
都是屬於知性的層面.
這種認識和第八節說.
認識主耶穌基督為至寶都是一樣的.
只是在第十節.
祂有再具體的內容去指出.
認識甚麼.
就是基督受苦的事奉生命.
和祂的生和死.
耶穌基督事奉所受的苦.
以及最後在十字架上的死亡.
正正是回應天父從上而來的呼召.
耶穌基督三天之後的復活.
亦都是獲得天父從天上而來的賞賞.
因為天父將主耶穌從卑微處升上至高.
讓主耶穌坐在天父的右邊.

$^{401}$如果贏得基督.
就是和主耶穌生命的結連.
效法祂的死.
就是過一個以基督為中心的信仰生命.
因為基督的死.
就是踐行一個信服天父的信仰生命.
主耶穌放下主權.
按照天父的旨意而活出這個信仰.
這就是耶穌基督.
從天上對每一個信徒的呼召.
這個呼召就是過一個以基督為中心的信仰生活.
而唯有達到這個呼召.
才可以得到這份賞賜.
同性戀家庭姐妹.
我們是否願意為基督過一個受苦的信仰生活呢.
一個效法基督的死的信仰生活.
一個放下主權.
按照主耶穌的旨意.
以基督為中心的信仰生活.
對我們來說是否一個遙不可及的信仰要求.
或者換另一個說法.
我們可以為主耶穌受哪一種苦.
我有一次回國內事奉.
講完一個主題講座之後.
有一個Q and A的環節.
有一個國內的同工問了我一個問題.
他說他和太太都是傳道人.
有一個女兒.
她看到國內的經濟開始蓬勃了.
好起來了.
有些會友的生活質素也改善了.
可以提升兒女的生活質素.
例如可以給錢上功課輔導班.
改善成績.
或者可以參加一些國外的交流活動.
提升個人的履歷.
有利進入一些心儀的學校.
有些比較富裕的還可以送兒女出國.
他說他覺得很掙扎.
究竟應否放棄一個傳道人的位份.

$^{441}$這樣可以令到家庭和女兒的生活改善.
他的Q and A問題是.
他覺得對女兒很有虧欠.
他覺得我做傳道人自己受苦不要緊.
但要女兒和我一起受苦.
會否覺得對不起女兒.
我當時的回應是.
如果這個生活質素比別人差.
就是為主受苦的代價.
你願意付出嗎.
為主而受苦.
是否一點回報也沒有.
最後我提醒了這位國內的傳道人.
這種對女兒有虧欠的心態.
會否是一種扭曲了的信仰價值觀.
女兒能夠過一個豐盛物質生活的童年.
是否真的比陪伴父母過一個對信仰毀身的童年好.
女兒將來的屬靈生命成長.
是否比物質生活更重要.
因此我和他分享了戴德生的一生.
戴德生一生奉獻給中國.
固然成為一個信仰的佳話.
但神恩待他的不是宣教的成就.
不單單是宣教成就.
而是神不斷祝福和使用戴德生的後人.
因為戴德生五代都成為宣教士.
被主使用.
戴德生的兒子戴傳人 戴傳兒.
他的孫子戴永勉 曾孫 戴紹增 袁孫 戴繼宗.
每一個都是被神使用的忠心僕人.
五代人都成全和實踐宣教士的使命.
如果我們的子女將來能奉獻給神使用.
這是我們的福氣.
子女在我們身邊成長.
見證信仰的真實 見證神的大能.
是否比物質享受更豐富.
為主受苦 是否只是受苦而沒有回報.
當然不是.
因此保羅為我們進一步解釋獎賞的內容.
獎賞的內容是甚麼呢.

$^{481}$縱然保羅說要得到天上的獎賞.
是否表示獎賞只能在將來才得到呢.
而現在是一無所有呢.
雖然我之前說過和修版和呂振中的易法有差異.
但我選擇和修版的易法.
即是呼召主耶穌從天上而來.
但賞賜未必要在將來才得到.
保羅在第九節回應了這個說法.
第九節說「不是有自己因律法而得的義.
而是有信基督的義.
就是基於信從神而來的義」.
今生所得到的是有信基督的義.
就是基於信從神而來的義.
中文譯本是.
「我現在有的義是因信基督而有的」.
我現在有的義.
即是保羅無需在將來才能擁有.
現在已經擁有了.
究竟保羅現在所擁有的義是甚麼呢.
義是指公正,公義,正直.
從信仰的角度.
是神要求人應當有的特性.
不過這個特性被罪污染了.
人無法完完全全彰顯出來.
最多只是一種殘破的顯現.
這種殘破的顯現.
虧損了神的榮耀.
對於猶太人他們認為.
可以滿足律法的要求而獲得這種義.
滿足了神所頒布的律法要求.
就可以成為義人.
但沒有人做得到.
保羅以為自己可以做得到.
但最後發覺這個方法是失敗的.
因此保羅說他的義不是從律法而來.
而是由相信基督而來.
即是這種義不是由自身的行為得到.
而是一種愛家的義.
以標準和角度去審判我們是否一個義人.
保羅說他現在所擁有的義.

$^{521}$即是神對保羅的一種信仰判斷.
是神認為保羅是一個甚麼人.
這個信仰判斷從來不是保羅自己判斷自己.
更不是當時的世界對保羅的判斷.
當保羅以相信基督能夠在基督面前.
被基督視他為義.
這就是保羅能夠即時得到的獎賞.
這個獎賞隨之為保羅帶來一個新的身份.
就是天國的子民.
神的兒女.
君尊的子子.
這也是今天所有人.
包括你和我.
當我們承認主耶穌的身份.
當我們承認主耶穌是我們生命的主.
拯救的主.
我們便即時可以獲得的獎賞.
另一種保羅要得到賞賜的內容.
就是他未能夠得到的.
未能夠即時得到的.
這也符合呂鎮忠的本內容.
在第21節.
按著那能使萬有歸服自己的大能.
使身體改變形狀和他自己榮耀的身體相似.
使萬有歸服自己的大能.
即是神完完全全掌權和作王的狀態.
這種狀態是指到將來某一天.
當萬有都歸於神的那一刻.
就是指末日.
新天新地的來臨.
卑賤的身體改變形狀.
是指信徒將來身體的復活.
今天我們的身體會摧殘.
會老死.
甚至化為無尤.
但復活之後的身體是永恆的.
他能夠承受永恆的生命.
一個永恆的身體承托著一個永恆的生命.
這個永恆的身體亦是一個榮美的身體.
而這個榮美的身體.

$^{561}$是與神的身體很相似.
將來所擁有的賞賜.
就是與神有份.
有著神的榮美.
這個份現在也有.
但不完全.
將來新天新地就完全了.
這個永恆的賞賜.
就是戰勝了人類最大的宿敵.
死亡.
人在地球的歷史上.
從來未能戰勝過死亡.
因為人現在的身體會摧殘.
會扭壞.
死亡隨著身體的殘破和扭壞.
而奪取了我們的生命.
當我們復活之後.
我們的新身體就不扭壞了.
我們就不再受死亡的威脅.
保羅所要得到的賞賜.
既有即時的.
即是天國的子民.
神的兒女.
和山川的祭司.
又有永恆的.
就是不扭壞的身體.
與主榮耀的身體相似.
因著這一份今生.
以及將來永恆的賞賜.
所以保羅能夠很決斷地說.
我可以凋棄萬事.
看作糞土.
為要贏得基督.
因為贏得基督.
就即時有一個新的身份.
因為贏得基督.
就能夠與基督有份.
將來就能夠得到不扭的身體.
作為賞賞.
同學們.

$^{601}$我今天講題是.
我要得甚麼的賞賞.
究竟我們今天作為信徒.
我們所關心的.
是我們會得到甚麼的賞賞呢.
是一份即時可以獲得的身份嗎.
如果是.
我們有沒有珍惜到自己這個身份.
我們有沒有踐行這個身份所要承擔的責任.
作為神的兒女.
我們有沒有推動天國的價值.
作為神的子民.
我們有沒有關心過神家的需要.
作為一個君尊的祭司.
我們有沒有為別人獻上禱告的祭.
如果我們關心的是一個永恆的生命.
我們現在所做的.
究竟有多少與永恆的價值是有關的.
傳福音.
靈魂得救.
是不是我們經常會牽掛的事呢.
甚至我們每一個屬神的子民.
都能夠肩負屬靈的新身份.
帶給我們的責任和使命.
當我們身處逆境.
面對為主而需要受苦的時候.
我們有一個永恆的盼望.
可以承托著我們的信仰.
保羅昔日想要得到的賞賜.
都是今天我們信徒所想要得到的賞賜.
我們一起去祈禱.
因為你的恩典不斷.
並且毫無保留地施予給我們.
你賜下救恩.
讓我們生命改變.
糾正了我們的價值觀.
讓我們行在主的喜悅的道路之上.
能夠遠離死亡之路.
求主你繼續向我們施恩.
繼續帶領我們忠心的事奉你.

$^{641}$按著你的心意而行.
願主你大大使用.
榮耀全歸於你.
我們的祈禱是奉主耶穌基督的聖名祈求.
願至祝福大家.
\newpage



\section{猶大書 }
\label{sec:xDCCx3kJCvo}
\textbf{2020.11.08《在至聖真道上造就自己》 猶大書 一 17-25 ( 和修版) 老耀雄傳道}
\newline
\newline
連結: \href{https://youtube.com/watch?v=xDCCx3kJCvo}{\texttt{https://youtube.com/watch?v=xDCCx3kJCvo}} ~~~~ 語音日期: 2020-11-11
\newline
\newline
\hyperref[sec:L_6ydn1BUoo]{\small{< < < PREV SERMON < < <}}
~
\hyperref[sec:index_chronic]{\small{[返順時目]}}
~
\hyperref[sec:index_scriptual]{\small{[返順卷目]}}
~
\hyperref[sec:buLtIY201P0]{\small{> > > NEXT SERMON > > >}}
\newline
\newline
$^{1}$梁國雄議員.
主席.
其他事項.
\newpage



\section{路加福音 17:11-19}
\label{sec:buLtIY201P0}
\textbf{2020.11.15《施恩望報》 路加福音 17\_11-19 梁永善牧師}
\newline
\newline
連結: \href{https://youtube.com/watch?v=buLtIY201P0}{\texttt{https://youtube.com/watch?v=buLtIY201P0}} ~~~~ 語音日期: 2020-11-17
\newline
\newline
\hyperref[sec:xDCCx3kJCvo]{\small{< < < PREV SERMON < < <}}
~
\hyperref[sec:index_chronic]{\small{[返順時目]}}
~
\hyperref[sec:index_scriptual]{\small{[返順卷目]}}
~
\hyperref[sec:UB7GZHuROmQ]{\small{> > > NEXT SERMON > > >}}
\newline
\newline
路加福音 17:11-19
\newline
\begin{longtable}{cl}
\hline
\hline
章節 & 經文 (和合本修訂版)\\
\hline
17:11 & \begin{tabularx}{0.7\textwidth}{X} 耶穌往耶路撒冷去，經過撒瑪利亞和加利利中間的地區。 \end{tabularx} \\ \\ \relax
17:12 & \begin{tabularx}{0.7\textwidth}{X} 他進入一個村子，有十個痲瘋病人迎面而來，遠遠地站著， \end{tabularx} \\ \\ \relax
17:13 & \begin{tabularx}{0.7\textwidth}{X} 高聲說：「耶穌，老師啊，可憐我們吧！」 \end{tabularx} \\ \\ \relax
17:14 & \begin{tabularx}{0.7\textwidth}{X} 耶穌看見，就對他們說：「你們去，把身體給祭司檢查。」他們正去的時候就潔淨了。 \end{tabularx} \\ \\ \relax
17:15 & \begin{tabularx}{0.7\textwidth}{X} 其中有一個見自己已經好了，就回來大聲歸榮耀給神， \end{tabularx} \\ \\ \relax
17:16 & \begin{tabularx}{0.7\textwidth}{X} 又俯伏在耶穌腳前感謝他。這人是撒瑪利亞人。 \end{tabularx} \\ \\ \relax
17:17 & \begin{tabularx}{0.7\textwidth}{X} 耶穌回答說：「潔淨了的不是十個人嗎？那九個在哪裡呢？ \end{tabularx} \\ \\ \relax
17:18 & \begin{tabularx}{0.7\textwidth}{X} 除了這外族人，再沒有別人回來歸榮耀給神嗎？」 \end{tabularx} \\ \\ \relax
17:19 & \begin{tabularx}{0.7\textwidth}{X} 於是他對那人說：「起來，走吧，你的信救了你！」 \end{tabularx} \\ \\
[1ex]
\hline
\hline
\end{longtable}
$^{1}$可能有些人不太知道,我自己本身也是在仙道會長大的,我13歲就回教會了,就回了仙道會..
那時候,聽到的吧?我們應該是錦繡堂的分堂,錦繡是北仙的分堂,.
也就是說,跟我有關係,因為我是回北宣的,所以那個關係特別密切..
更加難得能夠在這麼大的地方跟大家分享,這是一個祝福..
今天跟大家分享的題目就叫做「思因忘部」..
「思因忘部」,我們先想一想,其實是有問題的..
為什麼呢?上一代,或者中文都教導我們,.
思因莫忘部,忘部莫思因..
這樣的話,我想你就不要去做這些好事了..
很簡單的舉一個例子,如果我們來大教去幫助人,.
有回報的話,這個不是幫助,這個只是投資..
這段經文大家都很熟悉的,很簡單的..
那耶穌就問,咦,我醫好了那個,不是有十個嗎?.
只有一個回來,那其他的九個在哪裡呢?.
想不想?想.那耶穌這個想法,對不對呢?.
是不是思因忘部呢?有時候我們很熟悉的經文,.
反而我們沒有很詳細的去探究,所以這個我們需要反省的..
很簡單的舉一個例子,原來我們知道耶穌來到這個地方,.
不是無意之中撞到他們的,絕對不是偶然遇上的驚喜..
耶穌是刻意的來尋找他..
如果你留意,現在耶穌其實要去耶路撒冷,是預備受死..
但是聖經說道,他經過撒瑪利亞和加利利,.
我這樣看就覺得很普通,你看看這個地圖,最底的南部就是耶路撒冷,.
中間的是撒瑪利亞,北邊才是加利利..
耶穌去耶路撒冷的時候,是不需要經過撒瑪利亞,.
上回加利利,再下去,即是混路走..
耶穌是刻意上去加利利,尋找這十個人,不是偶然碰到的..
不要忘記,當我們還沒有認識神的時候,.
神已經為我們預備了救恩,每一個做父母的都會知道,.
我們未祈求以先,神就已經為我們來預備..
很簡單舉個例,當細蚊仔未出世的時候,.
父母已經預備了細蚊仔的衣服,床等東西,.
誰知道能不能生,不會的,先預備..
當他出世的時候,新一代已經開始為子女買教育保險,.
供七年八年之後,將來他讀大學的時候,.
就可以不用那麼辛苦,誰知道能不能讀成書,.
誰知道能不能長大,但你都會為他預備,.
近期很難受,但我會接受,甚至鼓勵,.
很多年青夫婦離開香港,去外地,我教會了很多人走,.
最快走的是12月4日,有一對夫婦又去英國了,.

$^{41}$但這些人走的,十個有九個家庭都是為了子女,.
但他們讀幼稚園,讀12年級的細蚊仔,.
是不會要求父母,你為我將來的幸福著想離開,.
不會的,但父母才去預備,聖經很清楚,.
原來當我們還罪人的時候,神已經拯救我們,.
但不單拯救,祂會祝福我們,我們一起來讀《馬太福音》第六章第八節,.
我們一起來讀出,所以很寶貴,.
當我們還不知道自己有什麼需要的時候,神已經知道,.
要我們預備,民間的信仰是人來,但去尋找神,.
但今天是神尋找我們,第二方面我們要看的時候,.
罪是蠶食人,麻風病,麻風病在猶太人心目中是罪人,.
正如一個人患了愛滋病,你會想,你需要什麼?.
你搞什麼令你患了愛滋病?原來可能是受感染,.
輸血的時候,但在聖經裡面,大家知道一些較為出名又犯錯的人,.
他都是有麻風被刑罰的,摩西的姐姐米利安背後來批評摩西,.
結果她患了麻風,有一個仙姿的門徒叫基哈西,.
仙姿來醫好了亞蘭王的元帥乃曼,她本來有麻風,.
但後來終於在約旦河洗七次,就乾淨了,.
那時候她送禮物給仙姿,仙姿很厲害,從來沒有直接見過她,.
最後奔赴門徒,我不要,要她走,那個門徒叫基哈西,.
後來她貪心追上去,說我們老師都回心轉意,.
給些禮物,結果怎樣呢?當她回去的時候,仙姿就說你貪心,.
所以怎樣?患麻風,所以他們犯罪的時候很容易患麻風病,.
大家知不知道乃曼的元帥叫什麼?叫乃媽,不簡單,.
當她離開離大人的時候,就患麻風病,所以麻風病被視為罪,.
而當時患麻風病的人,有病的人,有人關心,.
有人離大人照顧你,你會差很遠,我看到我們的程序表,.
有些弟兄姊妹,他們的健康出問題,我們都會登出來為他們禱告,.
甚至我們會傳訊息給他們,是一個安慰,.
但當日的麻風病者被視為罪人,所以他們不單身體受到傷害,.
他們心靈都會被傷害,還有就是罪是要令到人和神,.
人和人之間產生隔膜,很簡單舉個例子,.
聖經裡面記載,原來患麻風病的人就要被趕離家,.
不可以住在家裡,因為怕會傳染,年輕一代不太知道,.
五六十年代之前香港都有麻風病,.
香港的麻風病人當時送去知不知道什麼地方?.
喜寧洲,所以大家都很有那個知識,送去喜寧洲,.
是一個孤島,你不能走,原來聖經裡面說到,.
他要被趕走,他不可以出街,如果出街的時候,.
迎面有人來的時候,搖晃晃,大聲說,我麻風的,.

$^{81}$讓人們避開他,不然人們就用石頭扔他,.
所以你看到的時候,一個麻風病者其實是很傷,.
但我們更加體會到罪和麻風是很相似,.
一個患麻風病的人一開始的時候,.
不會爛到出面,可能在身體某一部份,.
外面的人看不到,但你沒有好好醫治,.
沒有好好處理的時候,麻風就越來越大,.
你麻風爛到出面,當一個人犯罪的時候,.
如果他成功,反而害死他,.
他不會因為這個緣故而停止犯罪,.
他是會越犯越大,越踩越深,.
所以罪是慢慢的感染,而罪無奈,.
我想做父母,我們做牧者都很深體會一件事,.
負面的力量永遠是大過正面的力量,.
正面的力量往往是不及負面的力量影響都大,.
十二個探子走去探路,十個說打不贏,.
有兩個說可以,結果是怎樣?.
那些人就埋怨,是受那十個人影響,.
很簡單舉一個例子,如果有十個橙,.
其中有一個橙發霉,你沒有拿走那個發霉的橙,.
我想問那九個好的橙,可否感染到這個發霉的橙不發霉?.
可不可以?是不可以的,.
但你不拿走那個發霉的橙,.
那九個沒有池的橙就會怎樣?.
池又會淡好,所以有時很無奈,.
所以父母有時會說,好東西又不見你學得那麼好,.
計數你又計得那麼差,計多少計不成,.
答排幾個回來,一計就計到,.
這個就是很大的悲哀,.
還有就是罪是令到人和人人和神之間有教學,.
就正如這個新冠狀病毒那樣,.
如果你中了招的時候,你就要怎樣?被隔離,.
我有個會友是教書的,在我們那間叫冥恩小學教書,.
誰知他就行山,原來有一個人中了招,.
他被視為緊密接觸者,現在要隔離十四天,.
不過都感恩,已經檢測了兩次,都是陰性反應,.
如果陽性就大件事了,因為他有回來崇拜的,.
求主保守,他要被隔離,所以你留意到,.
如果一個人犯罪的時候,帶來的就是傷痛,.
很簡單,一個丈夫拋妻棄女,令結新歡,.

$^{121}$後來還要迎接第三者入門,要逼自己的女兒叫媽媽,.
那個女兒說,叫媽媽?都不是她生的,.
結果甚至要自殺,不肯來叫媽媽,爸爸又逼她,.
一個家庭如果夫婦之間相愛互相尊重,.
他的兒女都在一個溫暖的情況下長大,是很sweet的,.
但問題當一有罪的時候,就令到這些美景破碎,.
還有就是罪是不會分等級的,.
不只是窮的人,沒有學識的人,才會犯罪,.
就好像這個病毒一樣,不只是老人院的老人家才會有事,.
你看英女皇的兒子又中,那個經常不梳頭的英國首相又中,.
連Donald Trump都中的,足球明星又中,.
斯朗又中,那些大隻佬明星又中,湯漢斯又中,.
你明不明白罪是污染每一個人,.
如果一個越有權力的人犯的錯,是影響到整個社會,.
越有地位的人犯的錯,影響性是越大,.
這個是很可悲的,所以罪是沒有分等級的,.
但我經常在消極裡面都找到積極的,.
在這裡我看到一個某個程度是好處,.
大家看不看到什麼好處,這十個人住在一起,.
見不見一個叫撒瑪利亞人,撒瑪利亞人就是混血兒,.
雜種,我們說是真是雜種,沒有什麼特別意思,.
大家知道猶太人是很憎外邦人,更加憎撒瑪利亞人,.
為什麼呢?撒瑪利亞人是猶太人混血兒,.
所以記得他們上去加利利的時候,.
他們都要繞圈走的,不肯進去撒瑪利亞這個地方,.
我舉一個例子,當我這條攤子是否得了,.
當我是YSL的老闆,大家知道YSL又出胎又出手袋,.
我當然很憎恨那些對手,那些什麼Cardia,LV,.
我很憎恨他們,但我最憎恨的是什麼?.
羅湖城賣YSL那些,因為他說是YSL,.
但那些質料很差,你明不明白我的意思?.
你抄我的東西,冒我的東西,是不是更加憎恨?.
因為大家知道撒瑪利亞人要抄他們的律法,.
但又說有祭司,但祭司又不是利未人,.
他們又有聖殿的敬拜,就抄一些減一些,所以特別憎恨他們,.
所以他們之間的關係是很差的,.
但這次這十個人竟然又住在一起,沒事的,.
我又抄開問問大家,知不知道YSL全名是什麼?.
姓羅,不是,是You Stupid Lady,.
那你都可以來帶我去買我的東西,所以就不要買YSL了,.

$^{161}$當他們在一起是沒有問題的,舉兩個例子大家就會明白,.
豪宅獨立屋是關門閉戶的,大家就比較,.
他現在日本車,我現在德國車,大家德國車就鬥牌子,.
我的模特兒好一點,相同牌子就鬥車牌,.
他的9413,我的1328,我勁過他,.
你留意到有錢的就關門閉戶,.
但我不知道大家有沒有經歷過或者見過那些安置區,.
現在很少,安置區一定開門的,互相幫助,.
喂,幫我看著我的兒子,我現在先去洗澡,.
公共浴室,行,你回來幫我看著我的狗,.
然後輪到我去,大家是互相幫助的,讀過書的都會知道,.
考第一和第二的那個表面好像很好,.
你很勁,92分多我2分,你真是厲害,.
你有沒有留意考第2和第3的那個表面,.
真是朋友的,是沒有問題的,你很勁,22分多我2分,.
我只有20分,你真是厲害,他真的覺得他很厲害,.
多我2分,你這樣都多我2分,.
即是考第2和第3的那個表面,真是朋友,.
所以大家那10個人可以住在一起的,你有沒有留意,.
大家鬥,我比你爛,很勁,沒有,所以有時候都較為特別,.
好了,耶穌主動去尋找人,但是不要忘記,.
人是要有回應的,神已經為我們預備了救恩,.
但都要我們相信他,我們一起來帶他去讀路加班第17章,.
12,13兩節,我們一起來讀出,.
(教會教學).
如果你留意到的時候,耶穌去到的時候,.
但這些人聽到耶穌要來到村的時候,馬上走出來,.
但大家留意,一會兒我們會對照,他們的距離是怎樣,.
他們是遠遠的站住,他不敢靠近耶穌身邊的人,.
但他很迫切的大聲呼叫,夫子,可憐我,.
你記住這10個人未必一定是窮人,可能是有錢人,.
他穩盡名醫,但都醫不到,這一刻他知道耶穌來了,.
他聽過耶穌一些的事,他期望,或者他相信,.
這個可能是他最後的機會,所以他把握住這個機會,.
他要大聲很迫切的帶他去呼叫,.
我想我們基督徒都是有這樣的經歷,其實是很無奈,.
當你在平順警方的時候,你都有祈禱,.
但不會很緊張,甚至可有可無,.
但當你在困境當中,你的禱告都迫切些,.
跟神的關係都好像緊密些,.

$^{201}$原來有時候我們都是臨急抱主腳,.
所以都求神幫助我們反省,為何他要遠遠站住?.
因為他知道自己是一個麻風病者,可能會傳染,.
但他更加知道自己是一個卑污的人,他不敢靠近主那裡,.
今天有時候我們忽略了一個很正確的思維,.
我們以為神是可親可愛,這是沒有錯的,.
神是滿有恩典憐憫的神,但我們往往輕忽了,.
神是尊貴榮耀滿有能力的神,.
在現在這個世代裡面,很多基督徒失去了敬畏神的心,.
以致他們很輕率生活,我們要反省,.
所以鼓勵在鏡頭前的弟兄姊妹,.
如果你沒有什麼特別的事,鼓勵你實體來崇拜,.
為什麼呢?始終在家裡的敬拜是差一點的,.
雖然理性知道神在你當中,.
但有時候WhatsApp,Recheck,你又要走去看了,.
我甚至見到很多教會都有這個問題,.
很多教會都是在實踐聖經的原則,.
在崇拜的時候,在前的是在後,在後的是在前,.
很早來的那些坐在後面,.
遲到的那些沒有位置坐,迫著要坐前面,.
你們有沒有留意?你們都有一點點類似我的情況,.
當年我建立冥恩堂的時候,買了第一個物業,.
是地下來的,但有些叫主力牆,.
就像你們的牆,不過你們這個牆小一點,.
我們那個牆大一點,相配你們這個牆,.
所以我們都像你那樣掛了個電視機在這裡,.
讓在牆後面的人可以來來去去看到,.
但原來那些位置是最熱門的位置,.
大家最早回去就說那些位置,.
特別還要選那個講完,.
這個講完,他講到一定有教訓的,.
於是那些人就會坐在這個牆後面,.
就是當我來,他一味認同你所講的,.
時運高,你看不到我,我看不到你,.
但是我反省,有些牧者真的講到的時候,.
是有催眠作用的,不得不承認這件事,.
但不要忘記,原來我們敬拜我們的神,.
神在我們當中,我們怎麼可以輕率,.
怎麼可以隨便,有時候很慘,.
我們做牧者就是這樣,.

$^{241}$剛才我都問有沒有開冷氣,原來有,.
我是一個很怕熱的人,.
我當年的教堂,馬來西亞的那一頁,.
有一個台,很矮的,我們不可以舉手,.
現場幾個人射燈射著,是很近的射著,.
但是沒有冷氣,台下就有冷氣,.
一邊射著一邊冒汗,我還記得有位姐妹,.
不知道是不是太急趕回來,忘記穿外套,.
穿了一件背心,和她男朋友坐在一起,.
坐著的時候,很冷,她男朋友應該讀物理,.
有些科學頭腦,知道摩擦是產生熱能的,.
很冷,現在是盲鼻按摩,和他擦背,.
後面那個姨姨眼都亮了,眼火都爆了,.
崇拜著,你為什麼要擦背?.
所以你留意到,剛才我來的時候,.
衣服都歪了,沒有領帶,為什麼呢?.
因為我是怕熱的人,.
但是我敬拜的時候,我一定會打領帶,.
把衣服放下去,因為這是對神的尊敬,.
所以求神來幫助我們,.
我們有沒有來,但是好好的,有這個意識,.
神是可畏,我不可以輕信..
這次耶穌如何來醫治她呢?.
我們一起讀路加福音第十七章第十四節,.
我們一起讀出..
記住神的目的一定是幫助人,.
但是方法可以不同的,.
如果你看路加福音第五章的時候,.
耶穌曾經醫治一個麻風病者,.
耶穌如何醫治他呢?.
走過去,摸他的臉,.
跟他說,我要叫你痊癒,.
這個人就得到醫治了,.
為什麼耶穌要摸他呢?.
耶穌叫一句話,.
拉撒路復活,死了死日都起來,.
為什麼要摸他呢?.
就是耶穌不但要醫治他身體的病,.
更加要醫治他心靈的傷痛,因為被人拒絕,.
所以耶穌摸他,.

$^{281}$這件經文給我都很有幫助,.
有時候我回家的時候,.
這段時間壓力很大,.
洗手了嗎?.
還沒有,還沒有,還沒有,.
還不去洗手,.
放心,放心,.
耶穌摸完那個麻風病者都沒有洗手的,.
都沒事的,.
神會憐憫保護的,.
這裡都救了我,.
你明白我的意思嗎?.
好了,但是這次呢,.
是十個,.
我相信耶穌為事摸這麼多,.
摸十個很煩,.
記住耶穌基督他做的事情,.
目的一定是建立生命,.
但是方法可以不同,.
耶穌跟他說,.
你去,.
被祭司察看,.
但是聖經裡面很清楚說,.
他們去的時候就得到潔淨,.
為什麼要去被祭司察看呢?.
原來有律法,.
有一個規定,.
我們一起來帶大家去看,.
你還沒記第十四章的第二節,.
一起來讀出,.
(大麻風病的自己,.
自己去看自己的,.
要帶回去見祭司,.
其餘的經文你自己看下去,.
原來做什麼呢?.
一個麻風病可能被醫好了,.
一樣要做的,.
第一,.
拿剃刀出來,.
剃一下毛,.

$^{321}$剃毛完褲,.
剃完毛,.
完整的皮膚,.
要給他看,.
所以叫剃毛完褲,.
好了,.
然後祭司就會來看,.
逐一逐一檢查,.
看看你還有沒有麻風,.
沒有,.
那麼祭司就會說,.
你好了,.
於是你要獻祭,.
然後宣告,.
這個人好了,.
可以回到自己的家,.
回到自己的社群裡面居住了,.
所以是由祭司去判斷,.
但是如果這十個人說,.
不要玩了,.
現在在爛,.
在臭,.
叫我去見祭司,.
去到也是這樣的時候,.
很弱的,.
可能得罪祭司,.
那麼我以後還要寫回頭的,.
但是他們要去試,.
信心不是頭腦的認知,.
信心是要有行動去回應,.
今天我們不是只聽到,.
我們要去行動,.
但是聖經很清楚說,.
他們去的時候,.
怎樣?.
就得到潔淨,.
所以很寶貴,.
我們看到一件事就是,.
醫生可能會治好那個人的病,.
但是卻不能夠解決人的罪,.

$^{361}$還有,.
律法或者法律只會定這個人有罪或者無罪,.
甚至可能枉判一些人,.
但是律法並不能夠令一個人成疑,.
從來沒有一個人坐完牢之後,.
就可以覺悟前非,.
甚至他的仇恨感是更加大的,.
如果過分的判斷,.
所以律法改變不了人,.
強權壓力並不能夠得到真正令人信服,.
但是不要忘記,.
只有聖潔的主,.
可以讓我們知道自己是一個罪人,.
但是他也賦予了這個力量,.
叫我們不再去犯罪,.
叫我們生命得到成熟,.
我們可以被稱為義,.
很寶貴,.
好了,來到重點了,.
就是施恩忘報,.
這十個人都康復了,.
耶穌就問,.
我治了不是十個嗎?.
那九個去了哪裡?.
正如我一開始說,.
那耶穌的思維對不對?.
我舉兩個例子,.
看看大家怎麼看,.
我不知道你有沒有這個經歷,.
你坐電梯的時候,.
電梯門快關了,.
你看到有人衝過來的時候,.
你按什麼按鈕?.
你真是好,.
我不是的,.
我是按關門按鈕,.
快點,.
等你進不來,.
這個是我的,.
我會做的事,.

$^{401}$我頗頑皮的,.
現在信了耶穌之後,.
我當然不會按這個按鈕,.
雖然有時候心裡也想按,.
但是不太協調,.
我還是會按開門按鈕,.
但是那個人一衝進來的時候,.
我沒有舉重,.
還啤啤你,.
你有什麼感受?.
我想你不是為了他而說謝謝,.
但是他這樣的態度,.
是令你難過的,.
又或者,.
大家有沒有看過人,.
在禮拜堂行婚禮?.
我們做牧者,.
看的東西多很多,.
多過你們,.
所以我們付出,.
其實是比你們好一點,.
我們會多一點祝福,.
我們坐在台上看到很多東西,.
譬如新郎新娘,.
對著聖壇,.
有什麼小動作,.
你看不到,.
你只是看背部,.
是不是?.
好了,.
婚禮結束了,.
新人不是衝出去嗎?.
接著主席或司儀說,.
好,我們再請,.
什麼三位太太回來,.
耶!.
然後就回來,.
然後拿著咪高峰,.
說謝謝,.
是不是?.

$^{441}$你這麼多年,.
養育之恩,.
栽培我成材,.
放心,.
不用你養了,.
現在老婆養的了,.
OK?.
接著就說,.
我們很想表達我們的謝意,.
我們有一點點東西想送給你,.
接著那些兄弟,.
拿一束花出來,.
一對新人就走過去,.
送給那對新人的父母,.
你們看不到的東西,.
我看到的,.
傳統的男人,.
男兒有淚不輕彈,.
但是看到十對有七對,.
那些爸爸在這個時候,.
看到兒子或者媳婦,.
拿一些花過來的時候,.
不斷流眼淚,.
為什麼他會流眼淚呢?.
有兩個可能性,.
第一就是,.
養了你三十年,.
得一束花就慘了,.
蝕水,.
比買螞蟻還慘,.
真可憐,.
還有,.
雖然從九月過了,.
清明還沒到,.
就送一束花給我,.
哭了一句,.
是不是?.
但是第二個可能性是什麼?.
他很感動,.
因為兒子和女兒,.

$^{481}$當著這麼多人的面,.
去多謝他,.
我想問,.
每一個做父母,.
難道他養兒女,.
養二十多三十年,.
目的是為了這句話嗎?.
如果是這樣的時候,.
虧也虧死你了,.
是不是?.
所以你記住,.
他不是為了這句,.
但當他這樣做的時候,.
是很令他感動,.
這個神對我們的期望,.
詩人在神的靈感動下,.
就寫下了,.
鼓勵我們歌頌讚美我們的神,.
我們一起讀詩篇79篇第13節,.
和105篇第1節,.
我們一起讀出,.
多謝人,.
我很想大家聽完這篇道之後,.
都至少做一件事,.
這是我一直以來都會做的,.
我也鼓勵很多弟兄姊妹做這件事,.
無論我坐巴士,.
坐小巴,.
我兌完卡之後,.
我一定會說多謝,.
或者勞煩,.
有些司機很愕然,.
有些沒有反應,.
有些很開心,.
但我不是在乎他的反應,.
我也要多謝,.
不用給錢嗎?.
已經給了,.
但我仍然要多謝,.
特別是我們對一些基層的人,.

$^{521}$去快餐店,.
我會自己盡量的,.
拿一些杯子和碟子去,.
你去外國,.
外國是自己拿去的,.
如果有人來收的時候,.
我會說多謝,.
不要看輕這句話,.
當疫情很嚴重的時候,.
澳洲的超市一樣被人搶購一空,.
很多人上完課,.
沒多久又要再上課,.
於是自己一些網民發起的運動是什麼呢?.
就是我們要去表示我們的謝意,.
於是很多人就找這些貼紙寫一些多謝的話,.
貼在那些貨架上,.
你知不知道這兩句話可能是鼓勵到一些人,.
譬如我有一次,.
無論是因為我坐火車,.
廁所裡面有些人來做清潔,.
我會說多謝,.
甚至有一次我記得我去The One,.
去那個洗手間,.
洗完手有個伯伯在,.
伯伯辛苦你了,多謝你,.
伯伯說,.
我收工的,.
我都不是,多謝你,.
觸動到伯伯,.
你真是好人,.
我開有限公司,.
好人有限,.
有什麼時候,.
不是啊,.
昨天有個人罵我,.
說我不用做,地上濕的,.
什麼事,.
不停地罵我,.
這些人就不應該,.
你有沒有留意到那些伯伯呢?.

$^{561}$因為我見到很多伯伯,.
不來不來男廁都有女人,.
我不知道女廁有沒有男人清潔,.
不是很公平,.
男廁有女人清潔,.
但女廁沒有男人清潔,.
我們永遠都受歧視,.
我們男人是弱者,.
女廁商場那些,.
是要拿鎖匙開門的,.
我們男廁不用拿鎖匙就開門,.
我們都受保護,.
你有沒有留意,.
做這些人多數是上了年紀,.
很真實地說,.
如果他經濟好,.
生活好,.
他應該享萬福,.
為什麼還要在廁所,.
你帶他做這些事,.
所以這些人我們不可以輕視,.
我們要去說多謝,.
你說一句多謝是不死的,.
但你知不知道那些人可能很觸動,.
這是我自己的原則,.
我的愛就是這樣流露出來,.
所以今天我們繼續要說什麼呢?.
就是「思人忘步」,.
記住這個思是什麼呢?.
是思想的思,.
我的女兒的名字叫思恩,.
我改名思恩,.
我只是學問回報,.
是真的,.
我們再讀這一段經文,.
你留意到這個撒瑪利亞人,.
就是去多謝耶穌,.
這個撒瑪利亞人可能在律法的事上,.
對耶穌的了解,.
或者對律法的認識不及猶太人深,.

$^{601}$但他在敬拜,.
在感恩的事上,.
他是超越了那些猶太人,.
還有你不要想著耶穌站在村子那裡,.
你帶他去等他回來多謝,.
不是的,耶穌已經走了,.
正在去耶路撒冷,.
但耶穌是較為出名的,.
所以較為容易知道耶穌的行程,.
他是去追著去找耶穌,.
多謝耶穌,.
最近國內有一件事令我很觸動,.
我都很欣賞,.
相信他們的父母一定對小孩子教導非常好,.
不知道大家知不知道這個很小的新聞,.
有一個學生去搭巴士,.
但後來上車發現原來他沒有帶錢包,.
這個都是我行常會做的事情,.
但我會帶電話,機不可失,.
但司機非常好,.
司機不是不收錢,.
司機自己拿了一元給小孩子,.
叫他入錢,.
但想不到後來這個小孩子要還一元給司機,.
要等這輛車過,.
結果在較為冷的天氣下,.
他站了兩個小時,.
車都很疏,.
這樣他終於找到司機,.
上去還一元給司機,.
和多謝司機,.
我相信如果我們中國每一個的子民,.
都有這個小孩子這樣的性格的話,.
我們才有真正的強國,.
這種是很高質素,.
是等兩個小時去多謝,.
今天我們有沒有這個多謝的心?.
有時沒有,.
還有,.
記不記得我剛才讀的那段經文,.

$^{641}$沒有打出來,.
這個人是怎樣來多謝耶穌?.
是俯伏在耶穌的腳前,.
聽到一個對比嗎?.
在麻風病的時候,.
軟軟站著,高聲呼叫,.
但當他結症之後,.
他是在耶穌的腳前,.
來多謝耶穌,.
剛才我已經說了,.
罪是破壞人與人之間的感情,.
這個年輕人販毒,.
毒品很多,那些小丸子,.
拘留了很久,.
後來終於審了,.
法庭判他有罪,.
然後法官問有沒有話要說?.
他不是求減刑,.
他說我的小孩從出生到現在,.
我還未抱過他,.
可不可以讓我抱抱他?.
法官不准,.
但他是容讓他的同居女友,.
抱著小孩,.
他在法人的前面,.
他伸手摸這個小孩,.
摸完之後,.
他一邊摸一邊哭,.
法官說夠了,坐牢吧,.
去多謝法官讓他摸這個小孩,.
什麼讓他和小丸子,.
他的親生兒子,.
產生這個隔膜?.
原來因為他犯了罪,.
而不是法官,.
但在這個病毒下,.
我們特別深感受,.
人與人之間很多的長,.
我們一起來看二弗所書第二章,.
十四十五兩節,.

$^{681}$我們一起來讀書,.
(二弗所書第二章,十四十五兩節).
原來我們人與人之間有很多無形的長在,.
現在吃東西有些板隔著,.
但有很多無形的長,.
人的罪令我們互相挑戰對方,.
鄙視對方,.
但信仰不只是為了上天堂,.
是叫我們生命改變,.
我們學習到什麼叫接納,.
什麼叫體諒,.
什麼叫寬恕,.
於是這面牆可以拆去,.
很寶貴,.
所以為什麼這個人可以俯伏在耶穌的腳前,.
來向祂道謝,.
耶穌問了九個去了哪裡,.
有時候我們會不會犯了這個錯,.
我們急需的時候不斷呼求神,.
但我們又不懂得如何回來道謝神,.
我那時候建立了教會大埔,.
星期三晚我們傳統的祈禱會,.
現在已經減少了,.
一個月一次,.
因為沒有什麼人回來,.
偶然有一次來,.
我們教會很復興,.
很熱切,.
星期三晚的祈禱會,.
有九十多人跪在地上向神禱告,.
如果看到這間教會就很興旺,.
但是為什麼會這麼多人呢,.
因為第二天是會考方訪,.
當年我們很多人來讀中五,.
明天是審判日,.
所以求神憐憫,求神引導,.
求神安慰,.
真的超過九十人,.
但是一個星期之後,.
我就和耶穌一樣在問,.

$^{721}$那個星期的祈禱會,.
只有十二三個人,.
有幾個人是傳道人,.
有兩個人是執事,.
真的會有五隻手指都已經鎖上了,.
我也在問這個問題,.
整個星期起不是有九十多人回來參加祈禱會嗎,.
為什麼今天只有這麼少呢,.
那七十多人去了哪裡呢,.
因為我一樣,.
我們利用神多於愛我們的神,.
這是我們的悲哀,.
但是不要忘記是因上加因,.
當一個人來去呼求,.
來去多謝耶穌的時候,.
耶穌怎麼說呢,.
起來,走吧,你的信就了你,.
那九個人有沒有經歷過這件事呢,.
是沒有的,.
他們身體得到醫治,.
但是他們的罪得不到赦免,.
但是這個人回來多謝主的時候,.
得到主的赦免,.
有時候我們很容易埋怨,.
一樣不好,那樣不好,.
但是不要忘記原來我們已經比其他人好很多了,.
你想一下這些小朋友,.
這些在國內的人,.
他們怎麼讀書,他們沒有書桌,.
他們沒有桌燈,他們skylight,.
拿著一個螺就晃著,.
如果你看有線的中國新聞組,.
我記得有一次我很感受心,.
就因為疫情緣故,.
在廣州有個人在街市那裡,.
賣什麼鵝頸,鴨舌,.
就好像這個港台,原來下面這個是空的,.
原來他的女兒一年級就在這個洞裡面,.
去讀書,.
因為不可以留她在家裡,.

$^{761}$在街市這樣生活,.
我們已經幸福很多了,.
很多時候我們認為神已經欠了我們,.
有時候我們都會覺得我們的父母欠我們,.
很悲哀的,.
我有兩個孫子,.
今年五歲,.
我兒子也很厲害,.
請了五個工人去看這兩個孫子,.
很厲害,.
有四個菲傭,一個印傭去看他,.
很厲害,那四個菲傭一個叫爺傭,.
一個叫媽傭,不是包,.
有兩個叫外傭,.
爺去做工人,媽去做工人,.
叫爺傭,媽傭,.
外公外婆就叫外傭,.
這四個叫菲傭,是菲的菲,.
不是工人,.
但在做工人的事,.
工人有錢拿,我們要拿錢去,.
我每次去的時候都會買東西給他,.
糟了,這個月是必然,.
有一次我去探望我孫子的時候,.
他們很開心,.
有人叫我,我很開心,.
然後問,.
雞蛋仔呢?.
因為他很喜歡吃雞蛋仔,.
這次沒有買,.
很快眼淚汪汪,.
為什麼不買給我?.
我心想,.
法例規定我每次來都要買雞蛋仔,.
我不是欠你,.
他罵我,.
多過份,.
他叫我死纏不休地去買,.
這些是我們的悲哀,.
你明白我的意思嗎?.

$^{801}$但你記住,.
我們的神沒有欠我們,.
只有我們欠我們的神,.
今天神才是我們的神,.
神才是我們的債權人,.
不是我們是神的債權人,.
反而神是破產者,.
因為神付出的永遠收不回對等的東西,.
我們怎樣努力都不能回報神在我們身上這麼莫大的恩典,.
不要忘記我們叫紅恩堂,.
神恩典真的好像洪水一樣湧流在我們身上,.
我們怎樣去回報?.
如果我們連可以做的事情都不去回報就很悲哀,.
所以今天我們要去記得,.
有時你不知道你做的事情有甚麼回報,.
這裡不詳細說,.
原來我當年去幫助一個家庭,.
我已經忘記了,.
原來我還給了他們二千元,.
但他們只是在教會很短時間,.
後來有一次,.
我去渥大陸一間很出名的中學,.
尼達去港島的時候,.
有個小女孩走來跟我說,.
牧師,認不認得我?.
我有句標準答案,.
很面善,但我忘記了,.
後來跟我說,.
當年你是小女孩,.
她還說到當年我們家庭最困難的時候,.
你給我們二千元,.
雖然你記得但沒有還給我,.
她說她現在已經出來做社工,.
因為是受我影響,.
所以她說到甚麼呢?.
我本來可以帶機構去工作,.
但現在我在深水Po ,.
專服務那些露宿者和天台,.
或者住在tong 房的人,.
我沒有想過,.

$^{841}$其實我到現在都記不起這個人,.
因為我對會友認識很深,.
是很短時間,.
但你想不到原來很多年之後有這個那刻,.
我今早出來的時候,.
我看傳道書,.
尾二那一章,.
就是你要殺你的種,.
不管環境好不好,.
可能有一天會收割,.
這是我們要去做的,.
他求神幫助我們,.
他求我們去做,.
這篇道對我們有甚麼實際回應呢?.
至少有五方面我們可以做到的,.
第一,.
就是我們要好好的歌頌敬拜我們的神,.
為甚麼我們要重視敬拜的生活呢?.
就是我們透過詩歌,.
是我們去歌頌我們的神,.
我們一起來讀希伯來書,.
第十三章的第十五節,.
我們一起來讀出,.
(何況與天科下者其數,.
上上以中下下為計之分身,.
且有心之立,成了主,.
而不知人,.
就在於天的過之).
所以崇拜不是我們要領受甚麼,.
這是學武的心,.
但我們一個星期寫一次回來,.
我們唱詩的時間很短,.
我們沒有歌頌我們的神,.
你想想,.
如果你聽到你的子女在別人面前說,.
我老爸煮菜很好吃,.
頂ye ,.
像一級,.
我媽媽身形超越肥媽,.
不單單這樣,.

$^{881}$她煮菜也一樣厲害過肥媽,.
我想你聽完之後會怎樣?.
會很開心,.
還有,.
她以後會繼續吃很多好菜,.
所以今天我們有沒有歌頌讚美我們的神?.
剛才我說,.
不信的人可能比我們更勁虔,.
這是我們帶隊去伊朗時拍的照片,.
這間清真寺很漂亮,.
我們進去時,.
那些人不斷地大聲大叫,.
我也有點尷尬,.
我們帶隊去拍照,.
但你有沒有留意到,.
他們是很勁虔地去祈禱,.
不會理會你們,.
後面黑色那堆不是垃圾,.
那是媽媽,.
他們兒子來的時候,.
他攪著媽媽,.
媽媽也很勁虔地去禱告,.
我們要檢視我們,.
究竟我們的敬拜生活是怎樣,.
這是感恩的第一步,.
第二步,.
我們有沒有好好地去見證我們的神?.
有兩方面,.
我們一起來讀書篇,.
第十八篇是第四十九節,.
三十五篇是第十八節,.
我們一起來讀書,.
媽媽,.
我愛你們的精靈,.
我愛大會中的主席,.
我愛大會中的教會,.
這裡很清楚,.
我要在微信的人當中,.
分享我自己的經歷,.
見證著,.

$^{921}$第二,我在弟兄姊妹當中,.
我也是一樣的,.
你要分享神在我身上的恩典,.
不是一定要立刻將它帶到信耶穌,.
但你要分享這個恩典,.
所以我很憎恨,.
我不是客套,.
我是用很感恩的心,.
去多謝大家,.
讓我有機會來講道,.
這是我的祝福,.
神有很多恩典,.
如果我現在帶它到信耶穌,.
會較為複雜,.
但也會透過攝師來載我,.
他真的會載我,.
這是我的祝福,.
所以你要記住,.
你要有一個敏感的心,.
因為會再講,.
你要多謝天父,.
還有就是要獻上金心的祭,.
什麼叫金心祭?.
不是律法規定的,.
是金心情願,.
什麼叫律法規定,.
什麼叫金心情願?.
很簡單舉個例子,.
我很興奮,.
我女兒今早跟我說,.
爹地,.
昨天你又這麼辛苦,.
今天你又不斷地講道,.
今晚請你吃東西,.
你喜歡吃泰國還是日本的東西?.
很享受,.
很珍惜自己可以給錢,.
但今晚他要請我吃飯,.
我不是真的沒錢吃飯,.
但這個邀請,.

$^{961}$是很觸動我的心,.
於是他就得到更多,.
我會疼他,.
我不會馬上給他錢,.
但很肯定,.
他會得到的祝福一定是多過他所付上,.
他不會說,.
我付了那些不是家用,.
是查錢,.
因為我女兒出來工作時,.
我頭三個月不用給家用,.
但他很好,.
他後來第二個月已經付了,.
然後他就說,.
爸爸,給家用,.
講話要清楚一點,.
我們解聖經都要解得很清楚,.
這個不是家用,.
你給我的那些,.
是連管理費都不夠給的,.
他是雙倍給我的,.
但我又要體諒他不是賺很多錢,.
我這些是查錢,.
你明不明白我的意思,.
他不會說請我已經給了你,.
我何須請你吃飯,.
法例都沒有規定他要請我吃飯,.
但這是甘心,.
我是很感動,.
我們的神都是很容易哄的,.
你看看哈該書,.
神說他們上山取木料,.
我就高興,.
不是要建好一間聖殿,.
你就已經很感動,.
所以我們沒有做這件事,.
我們一起讀《基伯來書》13章第16節,.
原來你去幫助有需要的人,.
這也是獻祭給神,.
好像在菲律賓的教會,.

$^{1001}$送禮物給保羅,.
保羅說我月立,.
就好像你們獻祭給神一樣,.
你幫助一些有需要的人,.
就做在耶穌身上,.
最後就是我們要樂於事奉,.
《基伯來書》12章第28節,.
我們一起讀,.
這裡很清楚的,.
尼達去說,.
當我們經歷到神的恩典,.
我們要感恩,.
感恩是一個愛,.
我們要愛事奉神,.
不可以輕率,不可以隨便,.
我到今天一個中文字都不會打,.
但你不要只聽我來講,.
講中文字很簡單,.
我用手寫,.
有些人跟我打,.
我會出一整篇講章,.
Fullscript稍後也會介紹一本書,.
都是我的講章,.
這時候我是比其他人辛苦一點,.
因為我不懂得練腦,.
但我不可以輕率,.
因為我不是只對弟兄姊妹,.
我要面對我的神,.
人生裡沒有辛苦這個字,.
有另一個字是我自己改,.
叫新樂,.
新樂,但快樂,.
因為是祝福,.
我下午會去另一間教會,.
他們五十一週年,.
昨晚我跟他們講了培靈會,.
他們叫我今天下午再講另一堂,.
我說可以,.
是祝福,.
所以求神幫助我們,.

$^{1041}$最後一點就是做敏感的人,.
敏銳於人和神在你身上,.
然後去感恩,.
為什麼我可以這麼多年,.
我不是一個好牧者,.
但我不會懶,.
我很享受時分,.
因為很感受到神對我的愛,.
很多很多祝福,.
人家送給我的卡,.
我都會掃描,.
在電腦,.
很簡單舉一個例子,.
有一包漁夫之寶,.
有一次我回教會時,.
我有點不舒服,.
我就問同工們有沒有後躺,.
他們說沒有,.
我就坐在第一行預備聚會,.
坐下來後,.
沒有人拍我背,.
小聲跟我說,.
牧師,我有包後躺,.
吃過的,給你,.
我說你吃過的,.
怎麼給我,.
很害怕,.
原來他吃了兩顆,.
他聽到我問有沒有後躺,.
他就把後躺給我,.
這個我很疼的回憶,.
當年讀中學時,.
我還記得我買過雲吞麵請他吃,.
那時很害羞,.
我等他走,.
好像我偷了東西,.
但他是一位好姐妹,.
在威爾斯做護士,.
當年沙士第68個中招的人就是他,.
但都很愛住,.

$^{1081}$他送一包後躺給我,.
不是那包後躺,.
而是背後那份熱愛,.
我帶隊去旅行,.
一年我只是第一次,.
在杜拜轉機,.
每個人都去玩,.
吃早餐,.
但我不會浪費這些東西,.
我就躲在麥當勞,.
一邊充電,.
一邊看書,.
有個團友很好,.
看到我沒有吃早餐,.
就去星巴克買杯咖啡,.
做早餐給我,.
有少許尷尬,.
我在麥當勞喝星巴克的咖啡,.
是是是,.
但我都會記下,.
你記住,.
當你敏銳於人和神在你身上的眾作,.
你整個人是不同的,.
神幫助我們,.
在以後的日子裡,.
當你坐小巴,.
坐地鐵,.
看到司機向你道謝,.
看到一些人其實很辛苦,.
他們來和你收碟,.
清潔桌面,.
說聲謝謝,.
更加還有你的家人,.
幸好你會去做這些事,.
很難說,.
WhatsApp給他,.
你想想如果深圳的發展,.
只是說深圳有這麼大的發展,.
但沒有提到香港,.
當年剛改革開放,.

$^{1121}$是香港給深圳和國內很多援助,.
絕口不提,.
我想我們都有一些不太舒服,.
但我們都不是為了要你道謝,.
第一句多謝是應該的,.
所以求神幫助我們,.
對我們的模樣,.
很願意這篇訊息成為大家的幫助,.
我低頭,我吐槽,.
天文教你自己用你的話語成為我們生命的幫助和激勵,.
這篇我們很熟悉的經文,.
讓我們能夠看到當中的淑玲的教導,.
也讓我們不是單聽道,.
讓我們行道,.
多謝你所賦予的恩典,.
讓我們真的來到好地的回報,.
雖然我們所擺上的微不足道,.
但讓我們真的竭力而為,.
以致我們能夠永遙主你,.
多謝你,.
聽下這篇禱告,.
\newpage



\section{瑪拉基書 2:17-3:6}
\label{sec:UB7GZHuROmQ}
\textbf{2020.11.22 講道《信心的等候》 瑪拉基書 2\_17-3\_6 謝靄薇牧師}
\newline
\newline
連結: \href{https://youtube.com/watch?v=UB7GZHuROmQ}{\texttt{https://youtube.com/watch?v=UB7GZHuROmQ}} ~~~~ 語音日期: 2020-11-26
\newline
\newline
\hyperref[sec:buLtIY201P0]{\small{< < < PREV SERMON < < <}}
~
\hyperref[sec:index_chronic]{\small{[返順時目]}}
~
\hyperref[sec:index_scriptual]{\small{[返順卷目]}}
~
\hyperref[sec:B0mTZIy27eI]{\small{> > > NEXT SERMON > > >}}
\newline
\newline
瑪拉基書 2:17-3:6
\newline
\begin{longtable}{cl}
\hline
\hline
章節 & 經文 (和合本修訂版)\\
\hline
2:17 & \begin{tabularx}{0.7\textwidth}{X} 你們用言語使耶和華厭煩，卻說：「我們在何事上使他厭煩呢？」因為你們說：「凡行惡的，耶和華看為善，並且喜愛他們；」又說：「公平的神在哪裡呢？」 \end{tabularx} \\ \\ \relax
3:1 & \begin{tabularx}{0.7\textwidth}{X} 萬軍之耶和華說：「看哪，我要差遣我的使者在我前面預備道路。你們所尋求的主必忽然來到他的殿；立約的使者，就是你們所仰慕的，看哪，快要來到。」 \end{tabularx} \\ \\ \relax
3:2 & \begin{tabularx}{0.7\textwidth}{X} 他來的日子，誰能當得起呢？他顯現的時候，誰能立得住呢？因為他如煉金匠的火，如漂洗者的鹼。 \end{tabularx} \\ \\ \relax
3:3 & \begin{tabularx}{0.7\textwidth}{X} 他必坐下如煉淨銀子的人，必潔淨利未人，熬煉他們像金銀一樣；他們就憑公義獻供物給耶和華。 \end{tabularx} \\ \\ \relax
3:4 & \begin{tabularx}{0.7\textwidth}{X} 那時，猶大和耶路撒冷所獻的供物必蒙耶和華悅納，彷彿古時之日、上古之年。 \end{tabularx} \\ \\ \relax
3:5 & \begin{tabularx}{0.7\textwidth}{X} 萬軍之耶和華說：「我必臨近你們，施行審判。我必速速作見證，警戒那些行邪術的、犯姦淫的、起假誓的、剝削雇工工錢的、欺壓孤兒寡婦的、屈枉寄居者的和不敬畏我的人。」 \end{tabularx} \\ \\ \relax
3:6 & \begin{tabularx}{0.7\textwidth}{X} 「我－耶和華是不改變的；所以，雅各的子孫啊，你們不致滅亡。 \end{tabularx} \\ \\
[1ex]
\hline
\hline
\end{longtable}
$^{1}$各位弟兄姊妹,平安!.
在我們看神話之前,我們一起祈禱.
每一天,每一次我們來到你面前,我們真是有何等的感恩.
因為我們不是必然可以相見.
但我們能聚集在一起的時候,我們能夠來到你面前.
我們能夠聽到你自己的話語,與主在聖靈眼當中感動我們的心.
若然我們有些不蒙你喜悅的,求主挪開.
若我們能夠專心一意的,能夠朝見你.
以致我們能夠聽到你的吩咐,我們能夠按著你的心意而行.
有讓孩子心中的意念,口中的言語,蒙你如立.
奉主耶穌基督命求,阿們.
在這個世界上,誰不自私呢?.
但是人為了自己的私己,而去歪曲了真理.
其實這些事情,在馬拉基先知的時代已經存在.
神已經厭煩他們的所作所為.
他們以為神是監牧.
以為這樣胡作非為是神不容許的.
其實神是要懲罰犯各種罪的假信徒.
我們首先看看聖經經文.
開始你們剛才讀過的經文.
我們看看.
你們用言語使耶和華厭煩.
卻說我們在何事上使他厭煩呢?.
因為你們說凡行惡的,而說看為善,並且喜愛他們.
又說公平的神在哪裡呢?.
當以色列人提出這些有關公義的質詢的時候.
好像我們今天被冤枉,受到冤屈,或者被欺壓的時候.
我們也會說,公理何在?.
你會覺得這些問題,神不是公義的嗎?.
他在哪裡呢?.
或者你看到,惡人得勢.
好人好像受盡很多冤屈.
好像很遭殃.
你會發出一個聲音,公義的神在哪裡呢?.
我們看看以色列人,他們所表達的態度.
必須從他們被捕過後,受到一些艱苦的生活.
從他們的背景我們了解一下.
相信大家都很有印象.
就好像以色列人,他們出埃及.
他們在抗議上沒有水喝.

$^{41}$他們開始埋怨神和摩西王.
他們說寧願留在埃及,做奴隸.
我都可以坐在玉波旁邊,有餅吃.
甚至他們說,誰給我們肉吃呢?.
記得我在埃及的時候,都不用花錢去魚吃.
有黃瓜,西瓜,韭菜,蔥,蒜.
現在來到這裡,什麼都沒有.
就只有麻辣.
我留在埃及好過,死因抗議.
他們在埋怨.
他們見到人犯罪,覺得神沒有擊打他們.
他們以為神認可這些罪.
他們見惡人興旺,以為神喜悅惡人.
他們由於嫉妒別人的成功.
將所犯的罪和神的忍耐混淆了.
由於他們的想法,他們彼此議論.
甚至他們去埋怨神.
他們沒有將這些埋怨帶到神的面前.
反而互相的議論.
你要想,互相的議論會出現什麼結果呢?.
本來這些人是屬於神的.
如果大家和和氣氣,一鼓作氣,多有激勵.
但是現在很多議論.
埋怨神的時候,變成大家彼此的洩氣.
曾經有一個玄町.
他很滿意為主人打理花園.
但他巡視花園的時候.
發現有一對他悉心照顧的花被人摘了.
他很生氣.
我轉頭來發現原來聽人說.
這朵花是主人摘下的.
其實這個時候他真的無話可說.
因為花園是主人的.
主人有權去摘,去除,去保留任何一朵花.
同樣這個世界是屬於神的.
神都可以支配一切的.
就算他讓人籌下一切.
但是成事在神.
神有很多計劃在他的手裡.
我們看到全部計劃都在神的手裡.

$^{81}$我們要學習什麼呢?.
就是將自己交託給神.
神為了這群百姓所發的言語去憂慮.
所以他透過馬拉基先知告訴他們.
他必須施行審判.
沒有一個邪惡的人可以逃脫審判.
拜偶像,淫亂,對神和自己的配偶不忠心的.
都要受審判.
他要差事者為他的來臨去做準備.
神是不改變的.
他是信實的.
但是這群以色列人變了.
人是會變的.
他們變了怎樣呢?.
他們離開了神.
但是他要再次和神的關係更新的話.
他們唯有迴轉.
他們迴轉歸向神.
所以萬君之言而說.
「汗衙我要差矣,我的事者在我前面,預備道路」.
他要預告.
他不僅要在主要君王來之前.
他要在他預備道路之前.
他還要在他預備道路之外.
他還要宣告.
告訴他要來的主和納入的使者.
君王要出巡的時候.
一定要有人先去準備.
記得我小時候是住在台灣的.
我們會日月潭.
在日月潭有很多觀光客去的地方.
但是發現那天有地方封了.
不讓人過去.
當然他怕閒雜人.
原來那天是當時的總統在那裡.
所以分得很清楚.
有些人是閒雜人,不可以過去的.
你會看到原來這個重要的人物.
要出現之前.
一定很多人做了很多的事情.

$^{121}$他們以前可能他道路.
各方面都要去鋪平車轍.
令到車子過去的時候.
是很放心的.
而且很安全的.
你會看到既然一個重要君王要出現的時候.
那些人都要這樣去做.
來臨之前.
這個先鋒他做的事情.
在我們聖經上.
我們可以在施洗若寒身上.
看到這個實現了.
看到施洗若寒.
他在出現的時候他做什麼呢?.
他叫人悔改,他預備人的心.
他宣告悔改的洗禮.
洗罪得赦.
他告訴人家日期滿了.
整個國近了.
你們要悔改,要信福音.
之後耶穌就來了.
另一位立約的使者.
就是人所仰望的耶穌基督.
你們所尋求的主.
必忽然來到他的殿.
立約的使者就是你們所仰望的漢.
快要來到.
主耶穌就要來了.
很近了.
當我們看立約的使者.
其實就是講到耶穌基督.
他來這個世界.
他是要拯救世人.
他不是要定世人的罪.
主確實來了.
他來到聖殿.
他成為了新約的宗保.
他取代了這個殿.
和其中的憲制.
記不記得這個經文.

$^{161}$有人說這個殿是要四十六年.
才可以建好的.
但耶穌說要三天重建.
其實他的意思是什麼呢?.
其實他指的是自己的身體.
就是這個殿.
所以當耶穌基督死在十字架斷氣的時候.
這個聖殿的萬寺.
聖所和至聖所中間有一個萬寺.
這個萬寺.
由上到下裂為兩半.
地震動.
盆石崩裂.
墳墓也開了.
有許多已碎的聖徒的身體也復活了.
我們看到一個很奇怪的事.
就是這個萬寺裂開了.
就是講到耶穌基督.
這個萬寺是代表他的身體.
耶穌復活之後.
他從墳墓裡出來.
進了聖城.
向許多人顯現.
這個灘門是講到那些復活了的人.
那些被我們離世的信徒.
伯夫長和一起看守的人.
看到這些事情.
由地震.
又看到剛才我們讀的經文.
所講的這些事情.
他很害怕.
但是他說他真的是神的兒子.
耶穌基督為我們付了極重的代價.
以弗所書2章14節說到.
因為他自己是我們的和平.
使雙方合而為一.
雙方是講到猶太人和外邦人.
本來是不合的.
不可以相容的.
猶太人是鄙視外邦人的.

$^{201}$因為他們可以合而為一.
他拆毀了中間的牆.
是耶穌基督.
他拆毀了牆.
是由他的救恩完成的.
而他自己的身體.
也終止了這個怨仇.
所以費吊那記載了法上的規條.
說要使兩方藉著自己.
造成一個新人促成了和平.
兩方就是講到猶太人和外邦人.
他見到耶穌基督完成了救恩.
是為了開了一條又新又活的路.
進到神的面前.
不用再靠祭物.
不用再靠祭司.
唯有靠耶穌基督.
他見到耶穌基督去成就一件事.
促成了和平.
記載十字架上消滅了怨仇.
就藉著十字架使雙方歸為一體.
與上帝和好.
猶太人和外邦人本來不合的.
本來我們都是遠離開神的.
但今天透過耶穌基督.
令到我們不但歸為一體.
我們還與神有和好的關係.
這是靠著耶穌基督帶給我們的.
所以引申到並且來傳和平的福音.
給你們遠處的人.
也傳和平給那些近處的人.
所以今天無論在世界各地.
或是我們在近處.
我們都有這個福氣.
是神很大的恩典給了我們.
因為我們雙方是藉著耶穌基督.
在同一位聖靈裡得以進到父的面前.
這是耶穌基督帶給我們的恩典.
我們有沒有看到呢?.
「只有你們不再是外人或客裡.

$^{241}$是與聖徒同國.
是與上帝加裡的人」.
今天去在座每一個人.
無論是不同的背景.
但當我們來到神裡.
原來我們在神裡是家裡.
家裡與客裡是不同的.
以前有機會去旅行.
現在被迫要留在這裡.
但你會覺得你做客裡.
就算你在外面過一天.
你也覺得自己的家.
自己的床還是比在外面舒服.
雖然自己的窩是一般般.
但你也覺得比在外面舒服.
所以今天我們在神的家裡.
是很美好的.
這是促成耶穌基督.
祂為我們所付的代價.
帶給我們.
而我們去看經文的時候.
我們有什麼感受呢?.
我們有沒有想到.
原來我們是神家裡的人的時候.
我們有沒有想到.
這完全是神的恩典.
祂給我們很大的恩典.
我們不再是客裡.
但耶穌基督來這個世界的目的.
是做什麼呢?.
祂不是要殺你們.
不是要毀滅.
祂是要清潔,要廉政.
祂仍然要將福氣給屬祂的百姓.
所以你們所尋求的主.
必不然來叫祂的殿.
納入的使者.
就是你們所仰慕的汗臘.
快要來到.
主耶穌快要回來了.

$^{281}$在我們剛才讀經文的第二節.
說到祂來的日子.
誰能當得起呢?.
在祂顯現的時候.
誰能當得起呢?.
其實這個經文.
是著重在耶穌基督第二次回來.
祂第一次回來的時候.
沒有人是當不起的.
是吧?.
當不起.
祂來是做萬人的救贖主.
祂第一次來.
現在第二次.
這個經文是著重在第二次.
因為祂如煉金像的火.
如飄逝者的鹼.
你見到祂必坐下.
如煉證銀子的人.
必結證利未人.
而他們像金銀一樣.
他們就憑公義獻公物給耶華.
我們發現這些都還沒有出現.
到現在還沒有見到.
蒙耶華月立的公物.
彷彿古時之日.
上古之年.
這些經文還沒有出現.
他們所獻的祭.
蒙神月立.
但現在還沒有出現.
反而在耶穌基督升天之後.
大概四十年的時候.
羅馬君就攻取了耶路撒冷.
然後這班猶太人就要跟他們搏命.
他們就把聖殿的圍籠燒毀了.
燒毀之後.
這班猶太人就繼續跟羅馬君打.
當羅馬君很氣憤的時候.
就把聖殿拆毀.

$^{321}$毀滅淨盡.
完全沒有了.
這個預言實現了.
所以.
既然聖殿都拆毀了.
獻祭也沒有了.
這件事也沒有了.
所以還沒有看到.
是神月立他們所獻上的.
第一次來.
我們還沒有見到這個審判.
各種強暴.
欺壓的事不斷增加.
不法的事也越來越多.
甚至連耶穌基督自己.
也要被釘在十字架上.
但第二次回來.
他之前藉著不同的災難.
讓以色列人受到毆斂.
然後他就跟天降臨.
要施行審判.
要結淨他的百姓.
使得他們夢神月立.
是必然的事.
他來的日子.
誰能當得起呢?.
他顯現的時候.
誰能立得住呢?.
才認為沒有一個人.
可以通過我們的神.
實施這個嚴重的考驗.
不過對於煉金的人.
他的目的是什麼呢?.
煉金人的目的不是要毀滅.
其實他要煉淨.
剛才我們看到.
因為他如煉金匠的火.
如漂洗者的鹼.
當主耶穌降臨的時候.
他要像火熬煉金子一樣.

$^{361}$去除他信徒的渣滓.
消去不潔的物.
好像除去罪惡.
除去壞的東西.
雜質.
臭的,污穢的.
他要除去.
這個鹼就好像那些.
以前人用肥皂.
去洗那些骯髒的東西.
現在有很多不同的牌子.
你將污穢的東西漂白.
對於那些做壞事的.
行惡的,破壞神律法的.
就要面對彼主去除滅.
必坐下如煉淨銀子的人.
必潔淨利未人.
而煉他們像金人一樣.
所以那時候煉銀.
會將火升得很旺.
然後將銀融在裡面.
就可以除去那些雜質.
其實當中特別針對利未人.
這批聖職人員.
他們在神眼中.
本來很寶貴.
但是他們被玷污了.
大家很熟悉的.
就好像.
以利的兩個兒子.
他竟然是做什麼?.
做祭司.
但是聖經怎麼說呢?.
聖經形容他們兩個兒子是無賴.
是不認識耶和華.
不認識耶和華.
但是他們是做祭司.
還有人家去獻祭的時候.
那些肉.
那些僕人是用三個齒的叉.

$^{401}$去叉那些肉的.
特別是那些人說.
我要燒了脂肪.
之後就給你了.
但是他們不願意.
他們說我要生的.
我不要熟的.
去搶那些肉.
聖經是什麼體統呢?.
但是聖經明明讓我們看到.
燒脂肪是表示獻給神的.
獻火祭.
輕香的火祭是獻給神的.
但是他們竟然去搶.
像什麼呢?.
做祭司.
我們看到這個年輕人.
他們所行所為.
在神眼中看到相當嚴重.
神都是警告以來.
他說你尊重你的兒子.
是多於神.
我們看到.
他說將百姓獻上美好的祭物.
養肥自己.
神要截斷他的防備.
和祖宗家的防備.
所以後代是沒有老人的.
中年已經死了.
因為他們做得太過份.
特別是父親老邁的時候.
都聽到別人說.
他們兩個兒子怎樣呢?.
做了種種惡事之餘.
他還和在會幕門前.
事後的婦人同寢.
真的很離譜.
但他做盡惡事.
神怎會不理會呢?.
你會看到.

$^{441}$因為他如煉金像的火.
神很憤怒,他要煉盡.
所以需要經過.
熬煉金人的師父.
去熬煉他們.
經過熬煉之後.
這群聖者學會了甚麼呢?.
學會了公義.
他們會按神的律法去行.
以正確的方式.
來向神獻祭.
他們不再憑外表.
做樣子去獻祭.
任隨意的獻祭.
他們可以憑公義.
去獻公物給神.
這就像經文所說的.
那時猶大和耶路撒冷所獻的公物.
必忘耶和華月立.
彷彿古時之日.
上古之年.
當聖職人員不隨便亂來的時候.
按正規去獻上祭物的時候.
神自然就會月立.
他們所獻上的公物.
就像古時摩西阿布拉罕.
他們的時代所獻的祭一樣.
所以經文第五節就說到.
萬君子耶華說.
我必臨近你們施行審判.
我必速速作見證.
警戒那些行邪術的.
犯奸淫的.
起假勢的.
剝削故宮.
宮前的.
欺壓孤兒寡婦的.
屈往寄居者的.
和不敬畏我的人.
我們不要漠視神的警告.

$^{481}$在1938年9月21日.
大約兩點半的時候.
在美國東岸.
就上了一個極度的強風.
這個巨風在襲擊地方.
強到連阿拉斯加.
和佛蒙特州的地震儀.
都很清楚地記錄了.
在這個巨風來之前.
有一天.
有一個人住在長島.
他就買了一個氣象測量儀器.
讓他帶回家.
氣象測量儀器的指示針.
就指著有巨風.
或者龍捲風.
但這個人就不太相信.
他就把這個儀器.
扔來扔去.
又把它扔在角落.
搖來搖去.
最後又是那個答案.
又是巨風和龍捲風.
他很生氣.
於是他就找回店主去理論.
然後他回家的時候.
他的房子已經沒有了.
因為他不相信那個儀器.
所說的指示.
所以那間房子被他捲走了.
他自己因為不相信的緣故.
都差點沒命.
警告是很重要的.
原來神都有警告法老王.
當你看到那麼多個災.
他不知道嗎?.
但他就當神沒有到.
很多個警告.
所以神很嚴厲的.
去懲罰法老和埃及人.

$^{521}$其實都成為我們的借鏡.
我們要信服神.
把我們生命的主權.
交給我們的神.
我們的學鳳.
降服在他的手裡.
我爺爺說:是不改變的.
所以雅各的子孫.
你們不自然滅亡.
這裡說到以色列的傳族.
或者他們的傳國仍然有傳流.
因為神的慈愛沒有改變.
所以雅各的子孫.
他們仍然可以保存.
按屬靈來說.
是個別的以色列人.
凡是有信心.
按著亞伯拉罕的信心.
這個宗旨去行的人.
他們仍然可以保存.
是真的以色列人.
雅各真的子孫.
他們就會有得救.
沒有滅亡.
但是都有一些個別的.
沒有信心的以色列人.
他們在馬來西亞先知.
去講這些預言之前.
和之後已經有一些.
雅各的子孫是因為不信而滅亡.
他們原本這些怨聲.
其實可以轉化成祈禱.
其實你們如果和人有過節的話.
最好找回那個人.
如果沒有辦法.
其實你可以將怨聲放在祈禱裡.
神可樂以聽人的祈禱.
祂願意在我們一切患難當中.
去安慰我們.
祂的安慰是無可限量的.

$^{561}$我們將我們的生命.
是交託給神.
因為主權是由祂去管理的.
祂說什麼呢?.
凡勞苦擔重擔的人.
可以到我這裡來.
因為耶穌的額是容易的.
是輕散的.
所以我們一定要學教託.
在英國有一個有錢人家.
有一次他們去度假的時候.
他的兒子竟然掉在泳池裡.
差點淹死.
花王的兒子救了他.
這個主人很感激.
這個爸爸媽媽.
就問花王.
我們有什麼可以報答你們呢?.
他說我的兒子想讀大學.
他想做醫生.
於是這個父母.
毫無猶豫的.
就答應了供他讀大學.
過了幾年之後.
這個邱吉爾就在德克蘭開會.
開會的時候他感染了肺炎.
英國就下令要找最好的醫生去治他.
這個醫生去找人治是誰呢?.
這個最好的醫生就是治好他.
他就是法明的.
佩林明醫生.
這個醫生也是花王的兒子.
他救了邱吉爾兩次.
你說人生是不是很奇妙呢?.
看到神有祂安排在每個人身上.
祂有祂的計劃.
我們只要盡我們所能.
擺盡自己給神.
任神去用.
神會在我們身上.

$^{601}$或者在我們的社會上.
或者在我們的國家裡.
祂有成就祂的大事.
所以我們只可以向神說.
我的生命在你的手裡.
求你引領我們.
有什麼後果我們不要介懷.
因為萬事都互相效力.
叫愛神的人得益處.
我們需要用信心去接受.
去依靠神.
我們一起去祈禱吧.
愛你的主 愛你的神.
讓我們有一群弟兄姊妹.
他們仍然願意愛你.
願意聽你的話.
幫助他們在難處裡.
仍然願意專心等候你.
仰望你.
又願意去服侍你.
幫助我們有信心去仰望.
依靠你.
因為我們不知道前路是怎樣.
但是主你是最清楚的.
你知道我們的未來.
你知道我們下一刻是怎樣.
求你幫助我們.
能夠將這個重擔.
卸給你.
以致我們有平安.
我們能夠勇敢地去走每一步.
我們知道每一步.
你都已經計劃了.
求你幫助我們不偏不倚的.
我們跟著你.
就能夠去到目的地.
求你在當中使用我們.
建立我們的信心.
我們這樣祈禱交托.
奉主耶穌基督的名教.

$^{641}$阿們.
多謝謝牧師的分享.
謝牧師說到.
神有祂的計劃.
我們深信神一定有祂的旨意.
就算在洪恩家也是.
神對洪恩家有祂的心意.
我們能否找到.
神在洪恩家的心意.
然後像謝牧師所說.
獻上我們自己.
還有謝牧師有一句.
講道裡有一個很好的提醒.
我們能否將我們的怨性.
化為一個禱告.
我相信我們人生有.
我們的生命還未完結.
\newpage



\section{約翰福音 13:1-17}
\label{sec:B0mTZIy27eI}
\textbf{2020.11.29 講道《活出愛》 約翰福音 13\_1-17 曾敏姑娘}
\newline
\newline
連結: \href{https://youtube.com/watch?v=B0mTZIy27eI}{\texttt{https://youtube.com/watch?v=B0mTZIy27eI}} ~~~~ 語音日期: 2020-12-03
\newline
\newline
\hyperref[sec:UB7GZHuROmQ]{\small{< < < PREV SERMON < < <}}
~
\hyperref[sec:index_chronic]{\small{[返順時目]}}
~
\hyperref[sec:index_scriptual]{\small{[返順卷目]}}
~
\hyperref[sec:pZ9tffANuvA]{\small{> > > NEXT SERMON > > >}}
\newline
\newline
約翰福音 13:1-17
\newline
\begin{longtable}{cl}
\hline
\hline
章節 & 經文 (和合本修訂版)\\
\hline
13:1 & \begin{tabularx}{0.7\textwidth}{X} 逾越節以前，耶穌知道自己離世歸父的時候到了。他一向愛世間屬自己的人，就愛他們到底。 \end{tabularx} \\ \\ \relax
13:2 & \begin{tabularx}{0.7\textwidth}{X} 晚餐的時候，魔鬼已把出賣耶穌的意思放在加略人西門的兒子猶大心裡。 \end{tabularx} \\ \\ \relax
13:3 & \begin{tabularx}{0.7\textwidth}{X} 耶穌知道父已把萬有交在他手裡，且知道自己是從神出來的，又要回到神那裡去， \end{tabularx} \\ \\ \relax
13:4 & \begin{tabularx}{0.7\textwidth}{X} 就離席站起來，脫了衣服，拿一條手巾束腰， \end{tabularx} \\ \\ \relax
13:5 & \begin{tabularx}{0.7\textwidth}{X} 隨後把水倒在盆裡，開始洗門徒的腳，並用束腰的手巾擦乾。 \end{tabularx} \\ \\ \relax
13:6 & \begin{tabularx}{0.7\textwidth}{X} 到了西門‧彼得跟前，彼得對他說：「主啊，你洗我的腳嗎？」 \end{tabularx} \\ \\ \relax
13:7 & \begin{tabularx}{0.7\textwidth}{X} 耶穌回答他說：「我所做的，你現在不知道，但以後會明白。」 \end{tabularx} \\ \\ \relax
13:8 & \begin{tabularx}{0.7\textwidth}{X} 彼得對他說：「你絕對不可以洗我的腳！」耶穌回答他：「我若不洗你，你就與我無份了。」 \end{tabularx} \\ \\ \relax
13:9 & \begin{tabularx}{0.7\textwidth}{X} 西門‧彼得對他說：「主啊，不僅是我的腳，連手和頭也要洗！」 \end{tabularx} \\ \\ \relax
13:10 & \begin{tabularx}{0.7\textwidth}{X} 耶穌對他說：「凡洗過澡的人不需要再洗，只要把腳一洗，全身就乾淨了。你們是乾淨的，然而不都是乾淨的。」 \end{tabularx} \\ \\ \relax
13:11 & \begin{tabularx}{0.7\textwidth}{X} 耶穌已知道要出賣他的是誰，因此說「你們不都是乾淨的」。 \end{tabularx} \\ \\ \relax
13:12 & \begin{tabularx}{0.7\textwidth}{X} 耶穌洗完了他們的腳，就穿上衣服，又坐下，對他們說：「我為你們所做的，你們明白嗎？ \end{tabularx} \\ \\ \relax
13:13 & \begin{tabularx}{0.7\textwidth}{X} 你們稱呼我老師，稱呼我主，你們說的不錯，我本來就是。 \end{tabularx} \\ \\ \relax
13:14 & \begin{tabularx}{0.7\textwidth}{X} 我是你們的主，你們的老師，尚且洗你們的腳，你們也應當彼此洗腳。 \end{tabularx} \\ \\ \relax
13:15 & \begin{tabularx}{0.7\textwidth}{X} 我給你們作了榜樣，為要你們照著我為你們所做的去做。 \end{tabularx} \\ \\ \relax
13:16 & \begin{tabularx}{0.7\textwidth}{X} 我實實在在地告訴你們，僕人不大於主人；奉差的人也不大於差他的人。 \end{tabularx} \\ \\ \relax
13:17 & \begin{tabularx}{0.7\textwidth}{X} 你們既知道這些事，若是去實行就有福了。 \end{tabularx} \\ \\
[1ex]
\hline
\hline
\end{longtable}
$^{1}$弟兄姊妹平安.
我記得上一次在這裡講道的時候.
是七月尾的時間.
眨眼間時間過得很快.
又到了十一月尾了.
在過程中自己也有很多的經歷.
我自己一直去預備的時候.
覺得神是首先要透過今天這段經文對我說話.
我將上帝對我所說的一一的和大家分享.
我們一同去禱告.
讓我們仍然有生命氣息來到你的面前.
很願意今天你就伺候你的話語.
成為我們腳前的燈路上的光.
願意聖靈充滿我們每一個的心.
不單止讓我們明白.
而是能夠讓我們坐然起行去活出來.
我們仰望你的憐憫.
奉耶穌基督的名字祈禱.
阿們.
有一句很重要的話.
我們看聖經時.
很容易被這句話打動.
就是「神就是愛」.
教會是神的家.
所以教會會充滿神的愛.
弟兄姊妹在洪恩家生活的年日也不一樣.
我現在看下去也不是每個人都認識.
或許我自己是其中一個比較新的.
在這裡生活的一個成員.
不知道是否有些人在洪恩家生活了一,兩個星期.
有些人就很長.
長到在洪恩堂創堂時.
當時是感受堂,直堂時.
大家就已經過來.
是我們的直堂勇士.
如果是這樣的話.
真的是紅內了.
無論如何.
不知道大家對洪恩家的感覺是怎樣.
覺得我們這裡是一個怎樣的群體.

$^{41}$我們是否一個充滿了愛的群體呢?.
如果要一個群體充滿了愛.
我們真的很需要.
很需要每一個弟兄姊妹.
在這裡有所付出.
今天就讓我們從耶穌基督的身上.
學習這個愛的功課.
大家如果有機會的話.
不知道可不可以手上都有聖經呢?.
我們今天一起翻開.
《約翰福音》十三章一到十七節.
記載了耶穌為門徒洗腳的事件.
整個事件是一個愛的宣言.
表明了耶穌基督完全愛那些屬於祂的人.
第一件值得我們問的是.
究竟耶穌基督愛的動機是甚麼?.
為何在這一刻.
祂要表達祂的那份愛.
背後的動機是甚麼?.
十三章一節這樣說.
「逾越節以前,耶穌知道自己離世歸附的時候到了,.
祂一向愛世間屬自己的人,.
就愛他們到底..
耶穌知道自己即將離世,.
回到天父的身邊..
當祂回到天父的身邊之後,.
門徒將會在地面上面對很大很多的挑戰,.
例如會面對逼迫,被捉拿,.
還有很多很多的困難..
在這個時候,.
耶穌很想給門徒一個鼓勵和安慰,.
很想告訴他們,.
祂是徹底地愛他們的..
就算有一天,祂不在身邊,.
祂要門徒記得祂那份愛,.
就可以得到那份力量..
困難的時候,.
門徒真的很需要耶穌基督那份愛,.
有了這份愛,.
他們在信仰上就不會那麼容易被動搖了..

$^{81}$其實在這一刻,.
當我們仔細去看的時候,.
我們不要抽離背景去看,.
你看回背景的時候,.
這一刻最需要鼓勵和愛的那一位是誰?.
是耶穌自己..
因為耶穌知道自己即將被門徒出賣,.
經歷不公義的審判,.
被鞭打,.
然後被釘死在十字架上..
祂將要為了背負人類罪惡的緣故,.
而經歷一個巨大的痛苦..
這一段路,.
是要祂自己一個人去經歷,.
去面對..
所以祂這一刻,.
是最需要安慰和鼓勵的..
但是在這一刻,.
我們看到的,.
耶穌基督不是不停地求,.
「你記住,.
我將要受苦了,.
我將要犧牲了,.
你記得安慰我,.
記得陪我一起,.
你們不要忘記,.
我很孤單的.」.
祂不是這樣說,.
祂也沒有求這些安慰,.
祂反而在這樣的形勢之下,.
盡力去幫助那些相信和跟從的人,.
堅固他們的信心,.
讓他們知道祂那份愛..
你可以看到,.
祂那份愛是多麼深,.
多麼大..
耶穌基督在這裡,.
要透過洗腳的行動,.
讓門徒感受祂的愛,.
這不是唯一的動機..

$^{121}$耶穌透過祂為門徒洗腳這件事,.
也要門徒去學習彼此洗腳..
這是經文裡很重要的地方..
我們一起看看約翰福音13章,.
13章我們去到後一點,.
12節打後,.
13章12節這樣說,.
「耶穌洗完了他們的腳,.
就穿上衣服,.
又坐下對他們說,.
我為你們所做的,你們明白嗎?.
你們稱呼我老師,.
稱呼我主,.
你們說的不錯,.
我本來就是..
我是你們的主,你們的老師,.
尚且洗你們的腳,.
你們也應當彼此洗腳..
我給你們作了榜樣,.
為要你們照著我為你們所做的去做..
我實實在在地告訴你們,.
僕人不大於主人,.
奉差的人也不大於差他的人.」.
耶穌基督是這群門徒的夫子,.
他們的老師,.
老師怎樣做,.
他希望學生也怎樣做..
耶穌是他們的主人,.
主人怎樣做,.
他也希望僕人怎樣跟著做..
耶穌自己為門徒洗腳,.
他做了一個榜樣,.
他顯示他自己徹底的那份愛,.
他做完這個榜樣之後,.
他也希望透過這個行動,.
要讓門徒學彼此相愛,.
彼此服侍,.
彼此洗腳..
要洗腳是為了什麼?.
不只是去除污穢,.

$^{161}$是學愛對方,.
服侍對方..
愛的動機,.
耶穌仍然很想透過這個行動,.
讓門徒得到鼓勵,得到肯定..
他離開之後,.
門徒會很孤單,.
面對很多挑戰..
剛才說了,.
他希望門徒記住他那份愛,.
是的,.
但是在地面上,.
耶穌還要做多一件事,.
他希望門徒記住彼此的愛..
在身邊的這些人,.
就是可以時時刻刻陪伴著他們,.
去面對這些經歷的人..
這些人可以給予對方愛,.
鼓勵和支持..
當信仰面對挑戰的時候,.
正正因為身邊有這些人的鼓勵,.
他們不會失腳..
所以這個愛的動機,.
也是耶穌要為這些門徒建立一個安全網,.
很重要..
今天這裡讓我自己有兩個思考..
第一件事,.
耶穌基督愛我們這麼徹底,.
自己經歷這麼大的困難,.
將會面對這麼大的挑戰,.
祂想的,.
都是祂所愛的門徒..
今天,.
我和你,.
會不會愛耶穌多一點?.
我剛才在下面唱詩歌的時候,.
我很感動,.
也很難過..
我在想,.
這麼久以來,.

$^{201}$經常說祈禱,.
經常為自己祈禱,.
祈禱這個,祈禱那個,.
有不同的事工,.
教會有不同的發展需要,.
祈禱,祈禱,祈禱..
為弟兄姊妹患病,.
不同的家庭需要,.
等等等等..
一路祈禱,祈禱,祈禱..
我想我有多渴望,.
停下來問耶穌,.
耶穌,.
你想我怎樣做,.
你才會很喜悅..
耶穌,你想我怎樣哄你開心..
耶穌,你想我怎樣過這些時間,.
單單為你,.
不是為了不同的事情,.
求啊,求啊..
耶穌,我可以為什麼,.
是你真的很喜歡,.
很開心,.
你也很感受到我的那份愛..
我真的覺得很慚愧,.
我不是說那些代求,.
自己的祈禱不重要,.
而是我在想,.
我關心的是什麼,.
我祈禱的焦點是什麼,.
有多少時候真正掛念耶穌,.
單單的愛祂,.
我自己也很慚愧..
在這裡我也挑戰弟兄姊妹去想一件事,.
我們很渴望耶穌基督的愛,.
祂的陪伴,.
我們有沒有想過耶穌的孤單,.
祂也很渴望我們愛的回應呢?.
我挑戰我們大家也是,.
一起單單的,單純的愛耶穌,.

$^{241}$問耶穌,.
耶穌,我們怎樣做,.
怎樣才能讓你感受到那份愛,.
是你喜悅的..
就好像我們跟愛人說話,.
我們所愛的另一半說話,.
怎樣才能讓配偶真正感到很溫暖,.
去愛呢?.
不是在他面前求啊求啊..
這是第一點,.
第二點我自己思考的就是,.
耶穌很想我們彼此之間去愛對方,.
服侍對方,.
我們這個服侍有做多少呢?.
我來到洪恩家的時間並不長,.
大概是兩個月左右,.
在這兩個月裡,.
也有很多很奇妙的經歷,.
我來不久,.
就立刻有弟兄姊妹約我吃飯,.
聊天,.
其實是很溫暖的,.
我也試過最後離開教會,.
晚上的時間,.
有弟兄姊妹陪我一起鎖門,.
這是為了彼此的陪伴,.
鼓勵,.
也有試過我自己去做壁報的時候,.
有弟兄姊妹主動請英幫忙,.
試過剛來辦公室的時候,.
有兩位同工很熱心地遞上這些清潔物品,.
讓我好好地清潔我的地方,.
其實那小小的動作是很溫暖的,.
也有試過和弟兄姊妹出去吃飯,.
弟兄姊妹照顧我,.
很細心地在吃飯的時候,.
還有,.
今天我去港島的時候,.
背後真的有很多的愛,.
有很多弟兄姊妹為我去禱告,.

$^{281}$我真的很感受到彼此相愛,.
彼此服侍的那份溫暖,.
真的很寶貴,.
但是在這裡,.
我在挑戰,.
我們是否只停在這個層面,.
可不可以再做多一點,.
不單止我們彼此認識的,.
相熟的,.
不認識的,.
不相熟的,.
大家看上去是在座的,.
是否每一位弟兄姊妹你都已經認識,.
你很熟悉他們了,.
可不可以將這份愛發出去,.
讓在座的每一個弟兄姊妹,.
或者今天在看直播的,.
都感受到這份愛,.
彼此服侍的,.
這一刻,.
耶穌透過這個群體,.
讓我們感受這份愛,.
在孤單裡感受到有支持,.
這不是只有一個人要去做,.
是整個群體一起去做,.
這是第二點我想到的,.
接下來的挑戰更大,.
剛才說愛的動機是什麼,.
耶穌很想門徒得著鼓勵,.
得著肯定,.
透過他自己的愛,.
透過門徒彼此的愛,.
第二件事,.
我們要問,.
愛的對象是什麼人,.
是否只是愛那些我們認識的,.
我們熟悉的,.
我們喜歡的,.
我們覺得很好的人呢?.
耶穌很愛那十二個門徒,.

$^{321}$但是十二個門徒並不是全然可愛,.
土人喜歡的,.
我舉個例子,.
彼得個性很急躁,.
衝動,.
後來三次不認主,.
雅各和約翰性情剛烈,.
曾經很希望將來耶穌坐在榮耀的寶座上的時候,.
他們可以坐在耶穌的左右,.
門徒爭造大,.
很喜歡權力,.
多瑪信心很軟弱,.
很多疑,.
加略人猶大,.
貪戀錢財,.
最後出賣了耶穌,.
耶穌被捉拿的那一夜,.
那群和耶穌生活了幾年,.
一同吃喝,.
休息,.
一同服侍的門徒是怎樣的?.
被捉拿的那一夜,.
門徒是耶穌被捉拿的那一夜,.
門徒為了保存自己的性命,.
四散,.
沒有人站在主的身邊,.
被你帶著一群這樣的門徒,.
你愛不愛得下呢?.
耶穌認識門徒的過去,.
他的現在和將來,.
他完全知道他們是怎樣的人,.
包括猶大在內,.
那個不是忽然眼前感到很驚嚇,.
原來猶大會出賣我,.
不是,.
他知道,.
是知道,.
但他選擇去愛他們每一個,.
他堅持愛他們每一個,.
包括猶大在內,.

$^{361}$今天我和你又如何?.
可不可以挑戰自己?.
不要只愛那些我們很熟悉的,.
喜歡的,.
我們覺得很好的,.
電影姐妹,.
而是愛那些自己,.
甚至乎,.
總之總是相處上不是那麼容易,.
總是有些困難的,.
不是那麼喜歡的,.
挑戰自己去愛那些,.
總是覺得他差一點的人,.
我們自己也有很多問題,.
如果上帝也按照我們這個問題,.
去選擇愛不愛我們,.
我和你也毫無盼望,.
但神就是這樣選擇愛我們,.
堅持到底,.
我們可不可以像神一樣,.
像耶穌一樣,.
堅持愛人,.
不是因為對方是怎樣的人,.
這才是困難,.
我深信在未來的日子,.
我們洪恩家會面對很大的挑戰,.
在裡面,.
就需要我和你起來,.
付出愛,.
活出愛,.
挑戰自己,.
愛得不一樣,.
耶穌怎樣愛,.
我們怎樣愛,.
第三件事很重要,.
要問的是,.
剛才說愛的對象是甚麼人,.
第三件事要問,.
愛的程度可以有多深,.
13章1節這樣說,.

$^{401}$他一向愛世間屬自己的人,.
就愛他們到底,.
到底的意思是指,.
很完全極致的意思,.
耶穌基督徹底的愛門徒,.
對門徒的愛從來沒有減少過,.
由他開始揀選他們,.
帶領他們,.
到他要面對生死考驗的日子,.
他都是這樣堅持愛他們,.
沒有變過,.
這份愛是很困難的,.
你看我們人,.
要這樣堅持愛是很困難的,.
我們人總是看很多外在的條件,.
看很多事決定去不去愛一個人,.
我們人也會變化,.
因為我們自己就是這樣充滿變化的一個人,.
可能開始這一刻很愛,.
愛下去,愛不下去,我不想愛,.
之後就沒有了那份愛,.
或者由頭到尾都不想愛,.
因為一些原因,.
但耶穌基督是很厲害,.
一直愛,堅持愛,.
所以門徒很有盼望,.
所以今天我和你才有這個盼望,.
這一份這麼深的極致的完全的愛,.
包括了什麼元素呢?.
不簡單,.
第一個元素是謙卑的付出,.
謙卑的付出,.
付出的意思是要有實際的行動,.
不只是口說我很愛你,.
你是透過行動告訴別人你很愛對方,.
耶穌基督的行動在這段經文很自然,.
就是為門徒洗腳,.
這個洗腳是什麼一回事呢?.
我們就想像多一點,.
其實大家經常看世界名畫,.

$^{441}$最後的晚餐的名畫,.
吃飯,.
很多時候看到的圖畫,.
就是一張長的橫向的桌子,.
然後耶穌坐在中間,.
然後門徒在左邊坐著,.
有一張椅子,.
但實際中東的情況,.
在耶穌的年代,.
不是這樣坐的,.
他坐著,.
你可以想像,.
門徒圍著一張很矮小的桌子,.
很矮小,.
不高的,.
圍著,.
怎麼坐呢?.
不是拿著椅子坐直,.
而是有一個斜靠的駁墊,.
有一個駁墊在下面,.
斜靠向左,.
頭枕在左手的手臂上,.
然後腳怎麼放呢?.
向外伸,.
向外對著桌子來說,.
向外伸,.
亦很自然,.
因為腳是用來行路的,.
那時候你沒有想像,.
那些鞋一定是很多名牌的鞋,.
很漂亮的鞋,.
那時候的鞋非常簡單,.
行路時,.
弄得整個腳都是塵,.
很骯髒,.
所以腳向外伸是正常的,.
否則,.
那些骯髒骯髒的腳,.
對著桌子很影響食慾,.
在過程中,.

$^{481}$大家都是這樣靠著坐著,.
腳很骯髒,.
我經常提及,.
在這個時候,.
你想像一下耶穌,.
站起來,.
離開那張薄薄的墊,.
然後脫下上衣,.
拿一條手巾綁著腰,.
然後倒水在盆裡,.
就去幫門徒洗腳,.
逐個洗,.
洗完腳,.
又用他綁著腰的手巾抹乾,.
整個過程,.
他都不是被迫的,.
他是甘心樂意,.
主動想去洗的,.
經文留了很多空間,.
讓我們去想洗腳這件事,.
其實經文沒有詳細講,.
但有些空間去想,.
例如,.
耶穌是否只是倒了一盆水,.
倒了一盆水,.
那盆水就刺激了,.
應該是誰第一個洗腳,.
那個最有好處,.
是乾淨的,.
我一定要做第一個,.
之後的拿同一盆水,.
全部都拿著黑黑的水,.
我相信不是,.
我相信耶穌基督會換水,.
這個門徒洗完,.
洗完之後,.
水都黑了,.
幫他抹乾,.
倒了水,.
再拿水,.

$^{521}$再幫第二個,.
如是者,.
至少十二次,.
這個動作不簡單,.
逐個去做,.
每次做這個動作,.
重複的,.
再一次提醒門徒,.
我很愛你,.
耶穌很愛門徒,.
這個行動很難,.
還有,.
耶穌為這個門徒洗腳,.
旁邊的門徒都看到,.
一邊看著,.
這個動作一邊重複,.
重複地進行,.
耶穌不斷地說,.
我愛你,.
愛你,.
愛你們每一個,.
透過行動,.
祂要說這件事,.
這個行動不是一般的行動,.
是一個很謙卑的行動,.
為什麼?.
洗腳在那時候被視為很低微,.
很卑賤,.
通常留給社會地位很低的人,.
在猶太社會裡,.
什麼叫社會地位很低?.
聽著,.
就是外邦的奴隸,.
婦女,.
孩子,.
學生,.
這些是當時群體被視為社會地位很低的,.
才做的,.
地位崇高的人,.
不會為地位低的人洗腳,.

$^{561}$同輩的人,.
也不會彼此洗腳,.
猶太的奴隸,.
也不會這樣洗腳,.
這樣是很卑微,.
簡直是有損身份,.
留給那些奴隸,.
那些低一點的人就好了,.
現在,.
你難以想像,.
耶穌基督,.
身為夫子,.
為門徒洗腳,.
更甚的是,.
十三章三節這樣說,.
耶穌知道,.
父已把萬有交在他手裡,.
且知道自己是從神而來的,.
又要回到神那裡去,.
耶穌知道自己有多尊貴的身份,.
祂從神來的,.
祂回到神那裡,.
父將萬有交在祂手裡,.
祂多麼尊貴,多麼崇高,.
在這個時候,.
怎麼會和洗腳扯上關係呢?.
各位萬有的主,.
我絕對不會迷一個這麼卑微的蟻民,.
如果從天空看下來,.
我們真的很小,.
不知道能否看到,.
還要為人洗腳,.
你想像不到,.
但祂即使知道自己是這樣的身份,.
祂沒有拒絕這麼卑微的服侍,.
祂可以降低自己,.
可以降低自己,.
你可見到祂的愛心,.
其實我覺得在過程中,.
固然付諸行動不容易,.

$^{601}$更難的是謙卑,.
只要我們在社會上,.
有一定的位置地位,.
人生中有一定的歷練,.
有一定的人生經驗,.
在家中有一定的位置,.
你叫自己謙卑下來,.
是很難的,.
說就容易,.
但耶穌就想我們學,.
你要活出這份愛,.
你就是謙卑的付出,.
放下你的頭,.
放下你的身段,.
這樣才可以去愛,.
我自己很感動,.
以前認識一位姊妹,.
非常專業,.
在精神輔導的領域,.
在訓練的層面,.
她經常可以去不同的公司,.
做很多訓練,.
她是一個領袖,.
那些公司不是小型的,.
普通的公司,.
是很專業的群體,.
她真的在裡面培訓領袖,.
她是領袖中的領袖,.
一直做這些事,.
但回到教會,.
她願意放下身段,.
做什麼事奉?.
你通常見到這些,.
很多經驗,.
可能她的專業資格幫到我們,.
一定要做領袖,.
一定要做一些我們覺得很顯眼的,.
很重要的位置,.
不是,.
她是特意選一些很卑微的位置,.

$^{641}$籌備飯吃,.
買東西吃,.
不用提我的名字,.
不需要說我有多偉大,.
不要特意提我公司,.
我可以的,我很開心,.
籌備飯吃,.
你想像一下,.
回到教會,.
有個弟兄,有個姐妹,.
見到地方污糟,.
自己主動,.
不等教會的清潔,.
自己拿來掃一下,.
搬一下桌,.
見到聖經都爛了,.
主動拿來包,.
你想像一下,.
這個弟兄姊妹是從哪裡來的?.
想像不到某間公司的CEO,.
回來教會掃地,.
又損身份,.
不是的,.
耶穌基督都為門徒洗腳,.
為何不行?.
為何回來還要說我是誰,.
我做甚麼是很重要的?.
回來這裡,愛的家,.
放下這個身段,.
付出愛,.
因為耶穌是這樣走的,.
第二個很重要的元素是犧牲社稷,.
這是愛的最高境界,.
在《約翰福音》十三章八節,.
彼得當時對耶穌這樣說,.
他說,你絕對不可以洗我的腳,.
耶穌回答他,.
我若不洗你,你就與我無份了,.
你看看,.
覺得有點奇怪,.

$^{681}$你在說甚麼?.
為何洗腳,.
跟我有甚麼份?.
原來洗腳的意義,.
耶穌在說甚麼?.
是指向他將要在十字架上,.
被釘十字架,.
犧牲,流出寶血,.
使人罪得赦免這件事,.
因為耶穌基督所做的救贖工作,.
門徒的罪,.
才得以潔淨,.
他們才可以與耶穌基督結連,.
他們在耶穌基督裡有份,.
若非因為耶穌基督的救贖,.
門徒不可以與基督有真正的結連,.
不可以在末世的福分裡,.
他們有份,.
原來洗腳的意義很深,.
所以耶穌透過整個動作,.
這番說話都是提醒門徒,.
愛到甚麼地步?.
不是只洗你的腳,.
不是因為我很謙卑,.
而是因為我愛的動作,.
最極致是我為你犧牲,.
我為了你,.
把我的性命都不要,.
讓你有真正的福氣,.
犧牲,捨己,.
我想起了一幅圖畫,.
我想起一幅圖畫,.
在地震裡,.
這是真實的事件,.
我不記得是哪個區,.
哪個國家的地震,.
是一個大災難,.
在地震後,.
房屋大量倒塌,.
其中有一幢樓,.

$^{721}$有一個小小的地方,.
其實整幢樓都碎了,.
塌了,.
有一個位置,.
救援人員去挖掘時,.
一直挖掘,.
挖掘了塵,磚,泥,.
離世後,.
發現有一對夫婦,.
那對夫婦是抱在一起的動作,.
夫婦二人已經離世,.
但他們緊緊的抱在一起,.
身上鋪滿泥沙的石頭,.
壓在上面,.
應該在地震時,.
就是這樣壓死了他們,.
於是救援人員需要處理遺體,.
處理屍體,.
搬起他們時,.
竟然發覺下面有一顆小生命,.
未曾離世,.
於是他們拯救這個小生命,.
大家才看到,.
原來這對父母,.
為了讓小朋友可以活下去,.
他們選擇抱著對方,.
保護這個小朋友,.
犧牲自己的生命,.
這真是最偉大的,.
在我們人類世界看到,.
犧牲的愛,.
犧牲自己,.
成全他們所愛,.
我們看到這些,.
總覺得很感動,.
耶穌基督更令我們感動,.
犧牲自己,.
讓萬人得蒙神的愛,.
去捨去自己生命,.
這不簡單,.

$^{761}$尤其是剛才大家看到,.
我們愛的對象是甚麼人,.
是否一定是很可愛的,.
你喜歡的,你覺得好的人,.
不是,.
耶穌也為那些,.
我們覺得不可愛的,.
我們不喜歡的,.
甚至我們覺得他不是好人,.
犧牲,.
其實耶穌會愛這些人,.
祂自己做了最大最大的示範,.
今天,.
祂也呼召我和你,.
挑戰自己,.
付出犧牲的愛,.
如果大家有機會,.
再慢慢看回耶穌的門徒的經歷,.
我沒有記錯,.
除了猶大自己是自殺的,.
除了約翰,.
其他大部分都是殉道的,.
耶穌放棄,.
捨棄自己生命,.
去拯救世人,.
離開罪惡,.
耶穌的門徒,.
也走上犧牲自己生命,.
為福音,.
為教會緣故,.
不是停在那裡,.
殉道時繼續下去,.
有很多宣教士,.
走上了捨棄犧牲的路,.
為了愛人,.
愛的最高境界是捨棄犧牲,.
今天所挑戰的是什麼?.
這裡哪裡有捨棄犧牲呢?.
有人逼迫我們嗎?.
有人拿槍打我們嗎?.

$^{801}$沒有,.
只是挑戰一件事,.
放下自己,.
這裡很難寫出犧牲,.
去到一個地步,.
拿出生命,.
但可以放下自己,.
放下自己,.
為了愛別人,.
付出最大的代價,.
這裡我留空給大家想,.
什麼叫放下自己,.
為愛別人付出最大的代價,.
我只希望,.
在接下來的時間,.
我們看到洪恩嘉,.
在這裡我們活出愛,.
我們愛上帝,.
我們彼此相愛,.
我們願意謙卑付出,.
我們願意捨棄犧牲,.
讓這個地方成為一個.
真的充滿愛的群體,.
真的在洪水橋成為明亮的燈台,.
在這裡我們一同禱告,.
天父上帝,.
藉著耶穌基督洗腳的事件,.
讓我們看到,.
愛的功課,.
真的充滿了挑戰,.
我們說就很容易,.
做出來就很困難,.
但求天父幫助我們,.
給我們力量,.
祝福洪恩嘉,.
成為一個大暖爐,.
在洪水橋,.
溫暖多人的生命,.
為了讓上帝你得榮耀,.
為了讓上帝你的心意得到滿足,.

$^{841}$奉耶穌基督的名字祈禱,.
阿們..
\newpage



\section{使徒行傳 8:4-8-26-36-38-40}
\label{sec:pZ9tffANuvA}
\textbf{2020.12.06《神的用人----腓利的生命素質》 使徒行傳 8\_4-8,26-36,38-40 陳朗欣傳道}
\newline
\newline
連結: \href{https://youtube.com/watch?v=pZ9tffANuvA}{\texttt{https://youtube.com/watch?v=pZ9tffANuvA}} ~~~~ 語音日期: 2020-12-08
\newline
\newline
\hyperref[sec:B0mTZIy27eI]{\small{< < < PREV SERMON < < <}}
~
\hyperref[sec:index_chronic]{\small{[返順時目]}}
~
\hyperref[sec:index_scriptual]{\small{[返順卷目]}}
~
\hyperref[sec:BO948tI9cPE]{\small{> > > NEXT SERMON > > >}}
\newline
\newline
使徒行傳 8:4-8-26-36-38-40
\newline
\begin{longtable}{cl}
\hline
\hline
章節 & 經文 (和合本修訂版)\\
\hline
8:4 & \begin{tabularx}{0.7\textwidth}{X} 那些分散的人往各地去傳福音的道。 \end{tabularx} \\ \\ \relax
8:5 & \begin{tabularx}{0.7\textwidth}{X} 腓利下撒瑪利亞城去，向當地人宣講基督。 \end{tabularx} \\ \\ \relax
8:6 & \begin{tabularx}{0.7\textwidth}{X} 眾人都聚精會神，同心合意地聽腓利所說的話，一邊聽他的話，一邊看他所行的神蹟。 \end{tabularx} \\ \\ \relax
8:7 & \begin{tabularx}{0.7\textwidth}{X} 因為有許多人被污靈附著，那些污靈大聲呼叫，從他們身上出來；還有許多癱瘓的、瘸腿的都得了醫治。 \end{tabularx} \\ \\ \relax
8:8 & \begin{tabularx}{0.7\textwidth}{X} 那城裡，有極大的喜樂。 \end{tabularx} \\ \\ \relax
8:9 & \begin{tabularx}{0.7\textwidth}{X} 有一個人名叫西門，向來在那城裡行邪術，自命為大人物，使撒瑪利亞的居民驚奇。 \end{tabularx} \\ \\ \relax
8:10 & \begin{tabularx}{0.7\textwidth}{X} 所有的人，從小到大都聽從他，說：「這個人就是神的能力，那稱為大能者的。」 \end{tabularx} \\ \\ \relax
8:11 & \begin{tabularx}{0.7\textwidth}{X} 他們聽從他，因他很久以來用邪術使他們驚奇。 \end{tabularx} \\ \\ \relax
8:12 & \begin{tabularx}{0.7\textwidth}{X} 當他們信了腓利所傳神國的福音和耶穌基督的名，連男帶女都受了洗。 \end{tabularx} \\ \\ \relax
8:13 & \begin{tabularx}{0.7\textwidth}{X} 西門自己也信了；既受了洗，就常與腓利在一處，看見他所行的神蹟和大異能，就覺得很驚奇。 \end{tabularx} \\ \\ \relax
8:14 & \begin{tabularx}{0.7\textwidth}{X} 在耶路撒冷的使徒聽見撒瑪利亞人領受了神的道，就打發彼得和約翰到他們那裡去。 \end{tabularx} \\ \\ \relax
8:15 & \begin{tabularx}{0.7\textwidth}{X} 兩個人下去，就為他們禱告，要讓他們領受聖靈， \end{tabularx} \\ \\ \relax
8:16 & \begin{tabularx}{0.7\textwidth}{X} 因為聖靈還沒有降在他們任何一個人身上，他們只奉主耶穌的名受了洗。 \end{tabularx} \\ \\ \relax
8:17 & \begin{tabularx}{0.7\textwidth}{X} 於是使徒按手在他們頭上，他們就領受了聖靈。 \end{tabularx} \\ \\ \relax
8:18 & \begin{tabularx}{0.7\textwidth}{X} 西門看見使徒一按手，就有聖靈賜下，就拿錢給使徒， \end{tabularx} \\ \\ \relax
8:19 & \begin{tabularx}{0.7\textwidth}{X} 說：「請把這權柄也給我，使我手按著誰，誰就可以領受聖靈。」 \end{tabularx} \\ \\ \relax
8:20 & \begin{tabularx}{0.7\textwidth}{X} 彼得對他說：「你的銀子和你一同滅亡吧！因為你想神的恩賜是可以用錢買的。 \end{tabularx} \\ \\ \relax
8:21 & \begin{tabularx}{0.7\textwidth}{X} 你在這道上無份無關；因為你在神面前心懷不正。 \end{tabularx} \\ \\ \relax
8:22 & \begin{tabularx}{0.7\textwidth}{X} 你要為你這樣的惡而悔改，祈求主，或者你心裡的意念可得赦免。 \end{tabularx} \\ \\ \relax
8:23 & \begin{tabularx}{0.7\textwidth}{X} 我看出你正在苦膽之中，被不義捆綁著。」 \end{tabularx} \\ \\ \relax
8:24 & \begin{tabularx}{0.7\textwidth}{X} 西門回答說：「請你們為我求主，使你們所說的，沒有一樣臨到我身上。」 \end{tabularx} \\ \\ \relax
8:25 & \begin{tabularx}{0.7\textwidth}{X} 使徒既作了見證，並且宣講了主的道，就回耶路撒冷去，一路在撒瑪利亞好些村莊傳揚福音。 \end{tabularx} \\ \\ \relax
8:26 & \begin{tabularx}{0.7\textwidth}{X} 有主的一個使者對腓利說：「起來！向南走，往那從耶路撒冷下迦薩的路上去。」那路是曠野。 \end{tabularx} \\ \\ \relax
8:27 & \begin{tabularx}{0.7\textwidth}{X} 腓利就起身去了。不料，有一個埃塞俄比亞人，是個有大權的太監，在埃塞俄比亞女王甘大基的手下總管銀庫，他上耶路撒冷去禮拜。 \end{tabularx} \\ \\ \relax
8:28 & \begin{tabularx}{0.7\textwidth}{X} 回程中，他坐在車上，正念著以賽亞先知的書， \end{tabularx} \\ \\ \relax
8:29 & \begin{tabularx}{0.7\textwidth}{X} 聖靈對腓利說：「你去！靠近那車走。」 \end{tabularx} \\ \\ \relax
8:30 & \begin{tabularx}{0.7\textwidth}{X} 腓利就跑到太監那裡，聽見他正在念以賽亞先知的書，就說：「你明白你所念的嗎？」 \end{tabularx} \\ \\ \relax
8:31 & \begin{tabularx}{0.7\textwidth}{X} 他說：「沒有人指教我，怎能明白呢？」於是他請腓利上車，與他同坐。 \end{tabularx} \\ \\ \relax
8:32 & \begin{tabularx}{0.7\textwidth}{X} 他所念的那段經文是這樣：「他像羊被牽去宰殺，又像羔羊在剪毛的人手下無聲，他也是這樣不開口。 \end{tabularx} \\ \\ \relax
8:33 & \begin{tabularx}{0.7\textwidth}{X} 他卑微的時候，得不到公義的審判，誰能述說他的身世？因為他的生命從地上被奪去。」 \end{tabularx} \\ \\ \relax
8:34 & \begin{tabularx}{0.7\textwidth}{X} 太監回答腓利說：「請問，先知說這話是指誰，是指自己，還是指別人呢？」 \end{tabularx} \\ \\ \relax
8:35 & \begin{tabularx}{0.7\textwidth}{X} 腓利就開口，從這段經文開始，對他傳講耶穌的福音。 \end{tabularx} \\ \\ \relax
8:36 & \begin{tabularx}{0.7\textwidth}{X} 二人正沿路往前走，到了有水的地方，太監說：「看哪！這裡有水，有甚麼能阻止我受洗呢？」 \end{tabularx} \\ \\ \relax
8:37 & \begin{tabularx}{0.7\textwidth}{X} 腓利說：『你若是一心相信，就可以。』他回答：『我信耶穌基督是神的兒子。』 \end{tabularx} \\ \\ \relax
8:38 & \begin{tabularx}{0.7\textwidth}{X} 於是他吩咐把車停下來，腓利和太監二人一同下到水裡，腓利就給他施洗。 \end{tabularx} \\ \\ \relax
8:39 & \begin{tabularx}{0.7\textwidth}{X} 他們從水裡上來，主的靈把腓利提了去，太監再也看不見他了，就歡歡喜喜地上路。 \end{tabularx} \\ \\ \relax
8:40 & \begin{tabularx}{0.7\textwidth}{X} 後來有人在亞鎖都遇見腓利；他走遍那地方，在各城宣揚福音，一直到凱撒利亞。 \end{tabularx} \\ \\
[1ex]
\hline
\hline
\end{longtable}
$^{1}$謝謝主席.
還有很多謝洪恩棠的邀請.
讓我有機會和大家一起敬拜.
大家好,螢光幕前的弟兄姊妹.
自我介紹,我叫陳朗恩.
大家可以叫我英文名Alison.
我是一個會畫畫的宣教人.
和一個會傳道的藝術家.
這些是我以前的作品.
這是送給我母會宣道會沙覺堂的一幅十字架作品.
有八呎高.
花了兩個月的時間去做.
這是第一次有人去買我的畫作.
是一個在台灣的酒店的畫作.
陳很幽默.
當我探索在哪裡服侍他的過程.
去過不同的地方.
中東,北非,非洲,東南亞.
很不經意地和當地人畫畫的時候.
不知不覺做了藝術宣教.
我以為做宣教士就等於要丟下畫筆.
但神教導我.
不是,而是將我的專業和信仰結合.
神教帶我到香港猜拳師公聯會做動員工作.
動員即是經常說話.
呼之欲人去參與猜拳.
但參與猜拳並不只是宣教士的一個角色.
是很多種.
動員是其中一種.
即使做宣教士也不只是專職宣教士.
有機會再和大家講解.
宣教和猜拳都是源自「mission」這個字的翻譯.
拉丁文的意思是上帝的使命.
萊特這位宣教前輩有這樣說.
宣教的使命是源自上帝的心.
祂的心傳到我們的心.
宣教使命就是一位全球的上帝.
祂全球的子民在全球的工作.
我這樣理解就像媽媽想兒子回來喝湯.
兒子遍佈世界各地.

$^{41}$其實是希望世界各地的人認識神的心.
祂對他們的愛.
「洛桑使命」「洛桑信約」也有這樣說.
普世的福音傳揚是要求宗教會.
教會全體將整傳的福音帶給全世界.
這裡有三個「傳」.
教會全體是全球的教會.
還有整體的像基督.
整傳的福音不單單是呼召人決自信耶穌.
其實是整本聖經.
所有界別議題的福音內涵和信息.
全世界所有的受眾物和全球所有的議題.
包括全球化,貧窮,被忽略的群體等等.
這是真正的宣教.
我們說自己是福音派教會.
「肥利」有一個美名.
就是全福音的「肥利」.
我們在祂的生命中看看.
一位初期教會神所用的領袖.
七位執事之一.
祂有什麼生命值得我們學習.
故事的開始是說耶路撒冷有迫不臨遁.
有巨大的衝擊臨到耶路撒冷教會.
所以當時除了使徒之外.
所有信徒都四散世界各地.
不是世界各地 是猶太各地.
包括在耶路撒冷即是紅色那一點.
去到猶太地黃色的地方.
和撒瑪利亞上面粉紅色的地方.
「肥利」是其中一個主角.
他去到撒瑪利亞城.
在撒瑪利亞城的時候.
有一天主跟他說.
「你去南走」.
「你去耶路撒冷 去加薩那條路」.
我想問你.
如果有一天神猜天使跟你說話.
「詩詩 你吃完午餐」.
不是 祂也沒有說.
「你去南走 向黃金海岸」.

$^{81}$「去九龍的方向那條路」.
你會怎樣呢?.
也要打點一下.
「黃金海岸很近」就去吧.
但如果再遠一點.
可能要走幾天的路程.
但無論如何.
作者路奸很有趣.
他形容「肥利」在二十七節.
就動身出發.
就好像媽媽吩咐子女.
「拿個碗來」就拿個碗.
我相信「肥利」也要打點.
但他是立即的.
到他聽到聖靈的提醒.
「你去吧 靠近馬車旁邊走」.
三十節說他跑過去.
這是我想說的.
「肥利」第一個生命素質.
就是他的聽功.
他不單單聽到.
而是行動迅速 果斷 無遲延.
如果神對我們的吩咐.
你聽到的時候.
我們會怎樣呢?.
願不願意聽到然後去走呢?.
第二個生命素質是.
「肥利」的講功.
他的講話.
當「肥利」到達馬車.
就聽到裡面的人很大聲.
讀一段文章.
我們知道他正在讀《二賽亞書》.
但「肥利」當時就讀《舊約》.
讀《二賽亞先知書》.
然後他辨別到讀到哪一段.
如果「肥利」神那麼遠.
差天使叫你去到那裡.
再叫你靠近馬車.
當然不是聽人讀聖經.

$^{121}$「肥利」就行動了.
他就明知故問.
「你讀的你明不明?」.
對方邀請「肥利」上車.
還問他.
「先知說的是自己還是其他人?」.
「肥利」就在那裡.
35字.
「肥利」就在那段經文開始.
講到耶穌基督的福音.
「肥利」沒有說.
「這一段我不是那麼熟悉」.
表面上就說「不如我和你說第二段」.
他不是.
他是由那裡說起.
別人讀哪一段.
就在那裡說起.
講到耶穌基督的福音.
還講洗禮.
我記得以前《掌門人》有個遊戲.
是猜歌的.
播放音樂的前段.
明星在那裡聽.
一聽前段就鬥快猜到這首歌是哪一首流行曲.
他能夠聽到音樂的前段.
已經猜到哪一首歌是因為他很熟悉.
「肥利」一樣.
聽到對方讀哪一段.
他就知道在說哪一段.
然後就在那裡又通到.
講到耶穌基督的福音.
這就是我想說的「講功」.
有熟悉聖經.
隨時隨地都可以分享.
其實呢.
我們飛一飛.
其實在很多時候.
我們都是這樣去講聖經的.
我認識一位宣教士.
黃聲鋒牧師和葛志西師母.

$^{161}$他們去到馬達加斯加.
當他們經常接待短宣隊.
黃師母其實是由早到晚都在廚房工作.
你可能說讀完神學去做宣教士.
結果竟然是做主婦.
但是黃師母沒有因為這樣而放下事奉.
其實廚房就變成了她的主教學的教室.
每一個進廚房的弟兄姊妹.
都會聽到她講她生命的故事.
她怎樣理解神的話 神的使命.
其實就是在上課.
有個美名就是「廚房神學」.
有一位蒙古宣教士也是這樣說.
當蒙古國不准公開傳福音.
他作為宣教士.
其實能夠分享福音的時候.
就是零散地.
無論任何時候.
例如他們辦學校 搞幼稚園.
辦有聲有色 有表演的畢業禮.
完結了 校長說完話.
宣教士真的說了一點點.
其實他就是在說訊息.
或者在餐廳的時候.
又在說訊息.
是零散的.
說完你都不覺的.
其實就是講道.
其實藝術和宣教都是一樣.
當有任何的場景.
無論是公眾的場景.
或者是一個私人的場館 展覽廳.
只要有藝術活動.
藝術家就可以藉著藝術作品.
或者藝術活動和人交流.
由交流作品到交流生命.
作品本身可以有道在裡面.
但是由對話開始.
我們由交流作品.
到交流到我和你的生命.

$^{201}$我們在對話的時候.
可以知道對方的生命.
對某些事情的看法.
交流信仰.
我有機會跟一個國際宣教團隊.
叫Tako Team.
他們去埃及短訓的時候.
我有一個成員是埃及的魔術師.
有成員是來自加拿大的空中雜技人.
和澳洲的音樂人.
還有我和烏克蘭的姐妹一起去畫畫.
我們沿著尼羅河流域.
表演了十一場的表演.
結果接觸了二千人.
有六百四十人願意認識更多的福音.
九十人是跟隨著.
一百二十人填回應表.
願意做大使命基督徒.
阿爾巴尼亞這個東歐國家的領袖.
曾經邀請Tako Team.
他們去表演一系烈怒天音樂會.
結果是公開呼召.
這就是他們音樂會的內容.
和他們公開呼召的情節.
結果有百多二百個穆斯林.
決志信耶穌.
當藝術表演團隊離開後.
當地教會領袖很努力跟進.
結果幾年後.
有事成為二十間教會.
你看到在音樂表演.
這個Rock Music.
即是滾石樂隊音樂會的時候.
你看到這一幕.
是音樂會表演的其中一個情節.
但我們看到.
其實是在說耶穌釘十字架.
藝術是可以再度的.
其實是講.
講功.

$^{241}$作為藝術家.
我們利用藝術作品去講道.
如果你是作為主婦.
你可否用你的食物.
你的廚藝去講道.
有藝術家.
有宣教士在宣教工場.
透過福音飯局.
當地是法國文化.
在非洲的地方.
用北京田鴨去招聚朋友.
然後搞著搞著.
搞到當地的穆斯林朋友不吃豬肉.
就算去那餐飯煮豬肉.
他們也照樣來.
如果你是醫療人員.
你可否用醫療去講道.
有弟兄姊妹.
他們去以色列服侍當地大屠殺的倖存者.
也是公公婆婆.
當他們用推拿.
拔罐.
舒緩他們的痛症.
令他們心花怒放的時候.
就講耶穌.
然後有人決志.
如果你是有物業的.
你可否用它來講道.
有弟兄姊妹.
在疫情前.
他們搞一個民宿.
免費的民宿.
不賺錢的.
去線上宣傳招募.
歡迎周遊列國的.
以色列年青人來香港住.
當他們入完軍訓.
當兵前後.
他們一定喜歡去外國旅遊一年.
他們見到免費.

$^{281}$當然來住.
結果.
跟你住.
就約他們吃飯.
帶他們遊香港.
過程中就講道.
交流生命信仰.
其實是一個講功.
我們如果看回肥利的場景.
其實更加看到神很憐憫外國的官.
我們回到前面的powerpoint.
其實不知道大家知不知道.
馬車裡的人是甚麼人.
我們知道他很虔誠.
他由他所在的地方.
即現今蘇丹的位置.
走二千多公里去到耶路撒冷朝聖.
花一程的時間.
都要二十日或四十日坐馬車.
你想想他請大假.
花費一大筆旅費.
為了去朝聖.
你猜猜這位外國有權有勢的太監.
即是江大姬女皇手下的財政大臣.
他能否去到聖殿.
我們看看聖殿的結構.
你見到黃色的範圍是外邦人院.
在廣場旁邊有門廊.
當年的文士耶穌都一樣.
在門廊教導人.
這個黃色廣場的位置.
我們知道當年擺了很多攤檔.
換牛換羊和換稅.
聖殿的銀錢很嘈吵.
其實像街市一樣.
到淺橙色的地方.
女院只能夠以色列的女性男性都可以進入.
但到深橙色的地方.
只能夠以色列的男性才可以進入.
到白色的地方只有祭司才可以進入.

$^{321}$原來在黃色廣場到淺橙色的地方.
有一級台階.
那裡有雙語的標示.
冒險闖入的外邦人.
一切死傷皆與人無尤.
這位愛國的大臣.
他的血統是非洲裔.
他不是猶太人.
他頂多去到廣場.
即是很嘈吵的街市.
還有他作為太監.
正如《申命記》所說.
身體曾經有過手術.
生殖器受損.
不能夠進入主的聚會.
一個有權有勢的太監.
即是愛國的財政大臣.
花這麼多代價.
竟然被拒絕.
好像是摸門釘.
你想想當他回程.
他讀到《以賽亞書》.
讀到這段經文.
他像楊高一樣被抓去宰殺.
受到不公正的對待.
他都不開口.
誰提到他的後代.
因為他的性命已經從地上被奪去.
咦 為什麼《以賽亞書》說的那個人.
這麼像我.
又是沒有後代.
又是被拒絕.
他覺得有人明白他的處境.
所以他很想知道.
《以賽亞書》說的那個人是誰.
是說他自己還是其他人.
當肥利有聽功 講功.
神用他去講聖經.
去回應他心靈的問題.
經文的前後都有提及到罪.

$^{361}$我相信太監一邊聽肥利的講解.
一邊認識原來先知以賽亞.
這裡所預言的耶穌.
就是那個尼賽雅.
亦都是承擔我們的罪的義僕.
我願意將我的罪孽交給他.
以至當後來見到有水的地方.
太監又明白洗禮的意義.
我們在這裡洗禮何妨.
這是肥利的心靈.
亦都是神使用他.
還有可能我們說.
這是愛國大神.
坐馬車有司機.
當然要好好招呼他.
但當肥利去到撒瑪利亞城.
他一樣熱切地與當地人分享信仰.
肥利有看功.
有神的眼光.
撒瑪利亞人是甚麼民族.
與猶太人沒有聯絡.
是敵對的.
被猶太人看不起.
我們去每一個地方.
總有當地人有民族優越感.
例如當地菲律賓人.
非洲裔在當地做工人.
總是低一層.
我發現世界各地的民族.
都有不合乎神的心意特質.
但肥利沒有這樣看他們.
即是在肥利眼中.
無論愛國大官.
肥利撒瑪利亞人.
甚至他作為猶太人自己.
都是同等的.
不同的民族在神面前都很尊貴.
如果我們有肥利的眼光.
每一個民族都按神的形象去做.
我相信我們看元彬.

$^{401}$這個韓國男神.
和我們巴基斯坦人的朋友.
經常搞導賞團的阿文.
是一樣的帥.
我們看我們的行政長官.
或其他官員.
都和我們尊敬的記者一樣.
還有我們看外國美女.
和我們經常見到的工人.
印尼姐姐.
都是一樣漂亮.
還有就算看英女皇.
和我們香港尼泊爾女孩.
一樣尊貴.
她們都是神體重寶貴的人.
當我有機會.
不知不覺.
做藝術宣教.
和別人畫畫.
純粹是消閒.
但我發現.
神一直吸引著我的眼光.
我有機會對著.
馬達加斯加的小女孩.
中東的敘利亞小朋友.
東馬的伊班族小女孩.
又或者是非洲的黑裔.
或者是白人.
或者是華人.
當我有機會.
無論他是穆斯林也好.
是基督徒也好.
我對著他們每一個人.
至少五分鐘.
對著一個人至少五分鐘.
讓我看到.
原來祂很漂亮.
祂很獨特.
每一個民族都是這麼尊貴.
這麼寶貴.

$^{441}$我覺得神吸引著我的眼光.
越來越明白祂的心意.
這是神對藝術宣教.
或者福音工作者的心靈的吸引.
我相信當你一邊煮菜.
一邊體會不同顏色的食物.
都已經營養均衡.
這簡直是神的創造.
在每一個行動.
神都會對我們有心靈的塑造.
如果有機會給你一包糖.
派給面前的小朋友.
我相信你會平均分配.
如果那些糖是代表教會的資源.
是代表我們的時間.
是代表我們的金錢.
而那些小朋友現在分佈世界各地.
這個地圖只顯示歐亞非.
有些小朋友在紅色的國家.
即是極度迫不及待的基督教國家.
有些小朋友在紫色的國家.
是非常迫不及待的基督教國家.
有些小朋友在綠色的標誌.
他們是難民.
有些小朋友在橙色的位置.
是被人口販賣.
被迫勞動或性剝削的受害者.
有些小朋友在間條的地方.
代表他們未認識神.
他們的群族沒有一個基督徒.
整個群族低於2\%是基督徒.
如果神要我們分糖的時候.
我們可以怎樣去回應呢?.
我們可以用《士徒人傳》一章八節.
去形容當聖靈降臨在你們身上.
你們就會得著能力.
在耶路撒冷,猶太傳地,撒瑪利亞.
直到天涯海角作為我的見證人.
傳講關於我的事.
當耶穌這樣說.

$^{481}$其實耶路撒冷和猶太傳地.
是有「同埋」這個字的.
即是你要去到哪裡.
耶路撒冷,猶太傳地,撒瑪利亞.
和天涯海角.
是不會少了一樣東西.
「同埋」是同步進行的.
多管齊下.
肥地另一樣的提供.
除了看到人平等尊貴.
他還看到最終領人去到神面前.
當他們洗完禮.
神帶他離開.
很有趣.
聖經形容.
神的靈將肥地帶走.
當時太監看不到肥地.
但經文形容.
他歡呼喜氣地繼續前行.
即是肥地成功將他帶到神面前.
即使他不再在一起.
但對象仍然願意繼續.
歡呼氣地與神前行.
如果我們也有這個心.
即是領人去到神面前.
或當我們所依靠的牧者離開我們時.
我們看不到他.
我們也一樣靠著神.
歡呼氣地前行.
第四宮就是酒宮.
剛才主席讀英文讀到四十節.
此時肥地發覺自己去了很遠的北部城鎮亞瑟都.
其實很有趣.
首先肥地是難民.
他在耶路撒冷住得好好的.
做教會執事.
結果要逃難去到撒米尼亞城.
即是神用得著任何處境.
我們成為難民也可以是傳福音.
當神差天使叫他去耶路撒冷.

$^{521}$去加薩的路就是馬車的位置.
就在那裡遇見了太監.
當他洗完禮.
神就將他帶到亞瑟都.
即是上面那個很遠的北部城鎮.
但結果肥地是怎樣呢.
他去到凱撒尼亞.
即是上面第五紫色的位置.
和沿途的城鎮傳福音.
當他已經去了一個很遠的北部城鎮.
但肥地願意再行遠一點.
再北一點.
肥地有這個心.
如果神吩咐我們去服侍一個人.
又或者一個族群.
我們願不願意履行.
甚至行遠一點 行多一步.
最後一個生命的素質是肥地的土供.
其實他去到撒瑪利亞的時候.
不是依治港鬼的.
除了港福耶穌是基督之外.
其實為甚麼可以能夠依治港鬼呢.
他之前是為他們祈禱.
其實是每一個的祈禱就行了一個神蹟.
祈禱其實是能夠行神蹟.
在1949年新中國成立的時候.
我們知道當年的領導人覺得.
基督教是外國勢力.
所以趕走所有中國境內的外國人.
對他們來說是外國人.
但事實上是世界各地的宣教士.
結果有幾十年普世的教會.
都不知道中國教會的情況.
他們只能夠祈禱.
為中國的教會代禱.
希望他們能夠站立穩.
在共產政權一個很高度壓迫的統治下.
中國教會信徒的信心堅固.
後來很感恩.
到80年代零零星星有些消息.

$^{561}$原來中國教會不但沒有滅絕.
結果開花遍地是急速的增長.
到今天我們都知道.
中國教會其中一個特質是祈禱很厲害.
他們很喜歡祈禱.
在疫情期間我有一個朋友.
他有參與國內的祈禱會.
早上六時有幾百人在網上.
有一個線上的青年宣教大會.
叫做Mission China.
在呼召的時候.
聊天室是水蛇春般長.
不停有「我願意」,「我願意」.
不停很久.
這不單是發生在現在.
我們知道中國有很多宣教士興起.
我們亦認識一些中國的青年人.
我相信這是世界各地的教會.
為中國教會祈禱的果子.
在二十年代.
我們知道全球都有一次經濟大衰退.
當年美國股市倒下.
全球都陷入經濟大衰退.
當年中國有一個反基督教的運動.
無論是文宣各樣的事.
政府都做事.
當時的文人都做文宣反對基督教.
但在同一個年代.
在美國中國內地會.
他們呼召二百個新宣教士到中國傳福音.
你想想經濟衰退.
在洛杉磯,紐約有很多貧民失業.
中國亦有迫不得基督教的浪潮.
但宣教工人呼召新宣教士到中國傳福音.
要錢,要人.
結果經過兩年的禱告.
和呼喚全球教會為這件事祈禱.
結果兩年後有二百零三位去到中國.
去做福音工作.
今日我們都可以透過祈禱行神蹟.

$^{601}$我們知道除了宣教士.
有很多種角色.
動員是一個.
接待是一個.
差遣是一個.
做宣教都有很多種形式.
但有一樣東西我們每一個人都可以參與.
是祈禱和奉獻.
在這裡介紹兩個資源.
有一個是雙語.
有一個是英文.
都有中文.
第一個叫闖開的門.
Open Doors Hong Kong.
如果搜索它的程式.
你可以Open Doors Hong Kong.
他們每一年都有一個更新的全球首望名單.
就是搜集了全球頭五十個國家.
聚迫的基督教會.
當地是有基督徒.
但他們很辛苦.
因為他們面對著各種威脅.
無論是自由還是生命.
例如我們知道第一位的北韓.
在新冠疫情下.
他們資源缺乏.
他們只知道很多人死了.
好像一個幽靈病.
死了都不知道.
在闖開的門的網站.
最新故事.
我們知道西非尼日利亞.
雖然在疫情下街市都停擺.
但極端分子沒有停下他們的工作去殺害基督徒.
他們仍然巡邏去殺害基督徒.
結果有一個牧師被殺.
在普通每日巡邏的村子.
看看有甚麼異樣的時候.
他的太太.
還有兩個女兒.

$^{641}$還有他肚子裡的嬰兒.
他們很辛苦.
本身他們的文化是.
夫家的家人會照顧遺孀.
但因為他們是基督徒.
他們得不到照顧.
很不容易等到政府有救濟物資.
結果不輪到他們.
因為他們是基督徒.
但他們很感恩.
透過開門的童工.
被救濟包圍.
渡過難關.
結果他們繼續去讚美神.
我們可以用禱告和金錢.
紀念其他地方的弟兄姊妹.
這個叫做普世宣教手冊.
在網站或程式.
每幾天會更新一個國家的資料.
包括他們的人民地理情況.
還有祈禱蒙英運.
和為他們繼續祈禱的事項.
前幾天是英國.
今天我們在《肥利的生命》裡.
看到一個不一樣的猜拳人.
一個舊領袖.
但他被神所用.
在本地跨文化去做不同的工作.
因為他有聽的素質.
敏銳 信服聖靈的聲音.
他有講的素質.
他熟悉聖經.
隨時隨地可以分享.
還有看到聖經怎樣回應人的需要.
他有看的眼光.
他把人的尊貴平等.
還有終極目標是將人帶到神面前.
還有他願意為神走多一里路.
還有他願意祈禱.
願意神都祝福紅印堂的弟兄姊妹.

$^{681}$我們一起學習.
效法肥利.
有一個神的用人.
願神大大使用你們.
謝謝.
\newpage



\section{帖撒羅尼迦前書 1:1-10}
\label{sec:BO948tI9cPE}
\textbf{2020.12.13《合神心意的教會》帖撒羅尼迦前書 1\_1-10( 和修版)  老耀雄傳道}
\newline
\newline
連結: \href{https://youtube.com/watch?v=BO948tI9cPE}{\texttt{https://youtube.com/watch?v=BO948tI9cPE}} ~~~~ 語音日期: 2020-12-18
\newline
\newline
\hyperref[sec:pZ9tffANuvA]{\small{< < < PREV SERMON < < <}}
~
\hyperref[sec:index_chronic]{\small{[返順時目]}}
~
\hyperref[sec:index_scriptual]{\small{[返順卷目]}}
~
\hyperref[sec:7YQ_zudEJ1I]{\small{> > > NEXT SERMON > > >}}
\newline
\newline
帖撒羅尼迦前書 1:1-10
\newline
\begin{longtable}{cl}
\hline
\hline
章節 & 經文 (和合本修訂版)\\
\hline
1:1 & \begin{tabularx}{0.7\textwidth}{X} 保羅、西拉、提摩太寫信給帖撒羅尼迦在父神和主耶穌基督裡的教會。願恩惠、平安歸給你們！ \end{tabularx} \\ \\ \relax
1:2 & \begin{tabularx}{0.7\textwidth}{X} 我們為你們眾人常常感謝神，禱告的時候提到你們， \end{tabularx} \\ \\ \relax
1:3 & \begin{tabularx}{0.7\textwidth}{X} 在我們的父神面前，不住地記念你們因信心所做的工作，因愛心所受的勞苦，因盼望我們主耶穌基督所存的堅忍。 \end{tabularx} \\ \\ \relax
1:4 & \begin{tabularx}{0.7\textwidth}{X} 神所愛的弟兄啊，我知道你們是蒙揀選的； \end{tabularx} \\ \\ \relax
1:5 & \begin{tabularx}{0.7\textwidth}{X} 因為我們的福音傳到你們那裡，不僅在言語，也在能力，也在聖靈和充足的確信。你們知道，我們在你們那裡，為你們的緣故是怎樣為人。 \end{tabularx} \\ \\ \relax
1:6 & \begin{tabularx}{0.7\textwidth}{X} 你們成為效法我們，更效法主的人，因聖靈所激發的喜樂，在大患難中領受了真道， \end{tabularx} \\ \\ \relax
1:7 & \begin{tabularx}{0.7\textwidth}{X} 從此你們作了馬其頓和亞該亞所有信主的人的榜樣。 \end{tabularx} \\ \\ \relax
1:8 & \begin{tabularx}{0.7\textwidth}{X} 因為主的道已經從你們那裡傳播出去，你們向神的信心不只在馬其頓和亞該亞，就是在各處也都傳開了，所以不用我們說甚麼話。 \end{tabularx} \\ \\ \relax
1:9 & \begin{tabularx}{0.7\textwidth}{X} 因為他們自己已經傳講我們是怎樣進到你們那裡，你們是怎樣離棄偶像，歸向神來服侍那又真又活的神， \end{tabularx} \\ \\ \relax
1:10 & \begin{tabularx}{0.7\textwidth}{X} 等候他兒子從天降臨，就是神使他從死人中復活的那位救我們脫離將來憤怒的耶穌。 \end{tabularx} \\ \\
[1ex]
\hline
\hline
\end{longtable}
$^{1}$洪恩嘉的弟兄姊妹平安.
如果有人問你教會是一個甚麼地方.
大部份人都會這樣說.
教會是一個敬拜神的地方.
如果我們再深入一點去問.
教會存在在世界上有甚麼目的.
我想每一個人都有不同的答案.
不過最後都會指向同一個方向.
教會存在在世界的目的就是要見證神和榮耀神.
一個能夠見證神的教會.
必然是一間合神心意的教會.
今日的講題就是合神心意的教會.
我們試試從帖撒羅尼加前書一至十節.
透過保羅對帖撒羅尼加的教會.
的教導讓我們學習三個實踐信仰的關鍵.
我們一起有個祈禱.
「真是恩主,求你臨到我們當中開我們的心.
讓我們將你的話語栽種在心裡.
成為承托著我們信仰生命的基礎.
好叫我們信仰的根基穩妥.
不會隨風飄去,隨水飄流.
成為一個沒有根的追求者」.
我們的祈禱是奉教主耶穌基督的聖名祈求.
阿們.
開始的時候我要交代一下帖撒羅尼加前書的背景.
帖撒羅尼加教會是在保羅第二次宣教旅程的時候建立的.
當時和他一起的同工還有西拉和提摩泰.
帖撒羅尼加前書的背景是在保羅第二次宣教旅程的時候建立的.
帖撒羅尼加前書的成書日期很早.
估計在耶穌被釘十字架之後.
二十年左右就已經成書了.
根據《使徒行傳》十七章一至九節的記載.
我們知道在保羅在帖撒羅尼加城的一個信徒.
耶穌的家裡面.
最少有連續三個安息日的時間進行福音的工作.
結果有些猶太人或者是希臘人信主當中也包括了不少婦女.
但是很快就引來一些不信的猶太人迫害.
引起暴動和騷亂.
保羅和西拉差一點就被他們拘捕了.
所以當地的信徒就安排他們立刻離開帖撒羅尼加城.

$^{41}$但是這批不信的猶太人仍然不放過保羅.
就算保羅離開了帖撒羅尼加城.
去了比利亞.
他們仍然跟到那裡.
還聳動那些人去抓保羅.
最後保羅就只留下西拉和提摩泰.
自己一個人走到雅典避開他們.
所以帖撒羅尼加城的教會是由一班初信主的信徒所組成.
他們缺乏了牧者一個靠近身的模樣.
而且面對很多的挑戰.
情況實在令人擔心.
所以保羅在雅典和西拉提摩泰會合之後.
就立刻打發了提摩泰去帖撒羅尼加.
去了解教會當時的情況.
後來提摩泰在哥林多和保羅會合了.
並且帶來帖撒羅尼加城的一些消息.
保羅是在這樣的處境下去寫帖撒羅尼加前書.
保羅寫教會的書信裡面.
一般開始的時候都有一段感恩或者祝福的圖文.
而帖撒羅尼加前書裡面就有三段這一類的感恩祈禱.
今天講到的經文就是保羅第一段的感恩圖文.
我們開始進入經文.
一間合神心的教會.
應該可以見到信望愛的生命見證.
在一章二節裡面保羅透露.
在第三節就詳細講明.
在第三節裡面有一個動詞叫做紀念.
在保羅未收到提摩泰向他回報的時候.
他心裡時時想念這批信徒.
他們究竟有沒有跌倒呢?.
他們究竟有沒有軟弱呢?.
因此當保羅聽到帖撒羅尼加城的信徒.
展現到一種信望愛的見證的時候.
其實在保羅所寫的經卷裡面.
有很多地方都在講信望愛.
我舉個例子.
在哥林多前書十三章我們很熟悉的.
講到信望愛裡面.
最大的就是愛.
在羅馬書保羅所講的盼望.

$^{81}$應該要歡歡喜喜.
在歌羅西書他講到信心的對象.
是主耶穌基督.
在迦拉太書他講到信心.
是要連著人愛才有果效.
雖然不同經卷都有講到信望愛這個主題.
但每一次的重點都不同.
而保羅對帖城的信徒所講的信望愛.
有他獨特之處.
就是這種信望愛.
是有對應的實際行為.
我們在第三節看到.
他說因信心所做的工作.
這裡講的信心是和工作有關的.
是什麼工作呢?.
我們看看第九節.
他說你們是怎樣離棄偶像歸向神.
來服侍那又真又活的神.
你們是指什麼人呢?.
估計是希臘人.
而不是在猶太會堂信主的猶太人.
這群外邦信徒信了主之後.
就離開了原先的信仰.
願意毀身去服侍神.
保羅為了一群初信主的信徒感恩.
因為他們接受了主耶穌基督之後.
以信心去毀身服侍.
同樣在第三節.
因愛心所受的勞苦.
勞苦的意思是指身體上的勞苦工作.
有辛勞至身體疲乏的含義.
這裡所講的愛心.
而不同第三章十二節.
彼此相愛的心愛眾人的心.
以及第九章.
以及第四章九節手足之情.
是同一個字根.
所以愛心是傳書.
是前後呼應.
貼撒萊尼加的信徒.

$^{121}$以愛心去幫助有需要的人.
甚至到了身體很疲乏的地步.
這種透過愛心去服侍.
也是對應一章九節裡.
所講到要服侍那又真又活的神.
的一種實踐.
怎樣服侍那又真又活的神呢?.
就是以愛心去關心.
幫助他身邊的每一個人.
保羅為了他們的一種愛心.
第三節最後講到.
因盼望我們主耶穌基督所傳的忍耐.
這裡所指的盼望.
和一章二章三章四章五章.
所講的盼望是一樣的.
保羅所指的盼望.
就是對主再來的盼望.
而在這個盼望裡.
信徒就需要有忍耐的行為和心態.
這個盼望的動力.
使到貼撒羅尼加信徒的生命裡.
產生一種忍耐.
在困難的情況裡.
不氣餒.
不放棄.
堅毅地前進.
這個就是貼撒羅尼加信徒所體現.
因盼望主再來所傳的堅忍的行為.
各位在紅磚街的弟兄姊妹.
保羅在圖文裡.
為什麼呢?.
因為信徒能夠以信心去服侍人.
以愛心去愛身邊的人.
以主再來的盼望.
去忍耐面前的困難.
如果教會裡.
我們找不到信望愛的生命見證.
我們怎能說自己是一間合神心意的教會呢?.
貼撒羅尼加的信徒.
在信望愛的信仰實踐裡.

$^{161}$對我們有什麼啟發呢?.
今年我們教會.
做了一些翻新工程.
教會在地下.
到禮堂之間的樓梯.
有三幅海報.
主題就是信望愛.
而三幅海報裡.
都有不同的副題.
分別就是.
信靠基督.
生命有力.
常存盼望.
無懼風浪.
愛神愛人.
直到永恆.
今天的社會.
面對很嚴重的撕裂.
我們說不知道該怎麼做.
我們能不能夠對不同政見的人.
多一份愛心.
甚至對一些我們看不順眼的人.
我們能夠心存歡舒.
當我們面對前面的路.
感到灰心失望的時候.
能不能夠有主在來.
主在來的盼望.
而產生一種忍耐和堅持呢?.
甚至信望愛的屬靈特質.
常存在我們洪恩家的信仰實踐裡.
第二個合神心意的教會.
能夠彰顯一個毀身的生命.
第四節.
神所愛弟兄.
我知道你們是揀選的.
夢揀選.
揀選這個字是一個受格名詞.
即是完全被動.
神做主動.
是神揀選貼上羅尼加的信徒.

$^{201}$作為教會的一員.
貼上羅尼加教會.
如何成為神所愛.
而又夢揀選的呢?.
第五節一開始就用.
「因為」這個字.
來帶出.
為何他們能夠被神所愛的原因.
第一個原因.
是因為保羅將福音.
帶到貼上羅尼加.
經文說我們的福音傳到你們那裡.
如果沒有保羅將福音.
傳入貼上羅尼加.
貼上羅尼加的信徒.
又如何信主呢?.
又如何建立教會呢?.
所以傳福音很重要.
我們知道保羅有三次宣教旅程.
而貼上羅尼加教會.
是保羅第二次宣教旅程當中.
在菲律賓建立了教會之後.
然後再來到貼上羅尼加城.
然後才建立的.
保羅在哥林多後書.
曾經說過.
他向外邦人傳福音的經歷.
他怎樣說呢?.
他說「我被他們多受勞苦.
多下監牢.
受鞭打是過重的.
誣死是屢次有的.
被猶太人鞭打了五次.
每次四十砍去一下.
被棍打了三次.
被石頭打了一次.
遇著船壞三次.
三次一晝一夜在深海裡.
有旅行遠路遭江河的危險.
盜賊的危險.

$^{241}$同族的危險.
外邦人的危險.
城裡的危險.
曠野的危險.
海中的危險.
假弟兄的危險.
受勞碌受困苦.
多次不得睡.
又饑又渴.
多次不得食.
受寒冷赤身怒體」.
一般人聽到這個經歷.
一定打退堂鼓.
這麼辛苦為了什麼呢?.
沒有誰願意為傳福音而毀身.
如果要去到這麼盡的話.
所以保羅夢召毀身向外邦人傳福音.
在整個過程裡.
保羅付出了很大的代價.
但是正正由於保羅願意付這個代價.
貼撒羅尼加的信徒.
才有機會聽到福音.
然後繼而信主.
貼撒羅尼加能夠被神所愛的第二個原因.
是因為福音本身的大能.
是跟保羅從聖靈而來的信心.
同樣的第五節.
「不僅在言語也在能力也在聖靈和充足的確信」.
新漢語也很準確.
「不僅是藉著言語也是藉著能力.
藉著聖靈和堅定的信念」.
堅定的信念.
不是說保羅對福音的信心.
而是指保羅對聖靈能力的信心.
是確信神對聽道者的請求.
所以這一節的意思就是.
福音的傳遞不單單是透過傳福音者所說的言語.
也是透過福音本身的能力.
而且也是傳福音者帶著從聖靈而來的信心.
確信聽道者是蒙神請求.

$^{281}$保羅來到貼撒羅尼加傳福音.
是帶著聖靈而來的信心.
貼撒羅尼加教會能夠備臣所愛的第三個原因.
是保羅帶著一個無愧的生命來作見證.
在第五節的下半節.
他說「你們知道我們在你們那裡.
為你們的緣故是一個怎樣的人」.
經文本身沒有說到怎若為人的內容.
不過從經文的上下文我們可以找到答案.
在第二章十節保羅說.
「我們向你們信主的人.
是何等的聖潔,公義,無可指責.
有你們作見證也有神作見證」.
所以我們知道保羅在貼撒羅尼加城的行為.
是聖潔的,公義的,無可指責的.
而這些行為也是貼撒羅尼加信徒可以見證的.
可以這樣說.
這是無愧的生命見證.
貼撒羅尼加教會之所以備臣所愛.
是因為保羅從無計較的付出.
滿有聖靈而來的信心.
並且實踐著一個無愧的生命見證.
將福音帶到貼撒羅尼加.
才令到貼撒羅尼加教會有美好的結果.
保羅,這種忠於召命的生命.
是否我們很羨慕呢?.
弟子們,教會是否能夠備臣所使用.
教會是否合臣心意.
其實每個信徒都有責任.
如果每一個信徒都不願意毀身的話.
教會是沒有可能會復興的.
亦不可能讓臣大大使用.
我們可能沒有保羅這麼捱得苦.
亦沒有他這麼堅強的心智.
但是我們有沒有一個願意毀身的心呢?.
我們相信一切都是從心出發的.
而臣有閱納每一個願意毀身的心.
臣會使用每一個願意毀身的信徒和教會.
只要我們願意.
臣會就大大使用.

$^{321}$第三個合臣心意的教會.
就是以效法主為使命.
在第六節.
並且在大能之中蒙了聖靈所賜的喜樂.
領受真道就效法我們.
也效法了主.
新漢語將經文的次序改了一點.
你們效法了我們.
又效法了主.
在大患難中.
懷著聖靈所賜的喜樂.
領受真道.
這個譯法強調了.
效法了我們又效法了主這個信息.
他們如何效法保羅和主耶穌呢?.
就是他們在大患難中.
帶著聖靈的喜樂.
接受了真道.
領受真道就是接受了上帝的道.
即是福音.
貼撒羅來加信徒在兩種狀態之下.
去接受福音.
第一種接受福音的狀態是在大患難中.
這裡所指的大能未必是指.
基督徒在信仰上受到羅馬政府的大迫害.
那種嚴重.
事實上羅馬皇帝尼祿皇帝.
規模迫害基督徒的時間.
大概在公元64年左右.
而保羅在貼撒羅來加傳福音的時間.
大概在公元53年左右.
距離真正的大迫害相差十年.
我們知道保羅在貼撒羅來加傳福音的時候.
有猶太人和希臘人信主.
包括婦女.
這些希臘的新信徒.
是轉信一個陌生.
而且備受社會懷疑的宗教.
如果他們拒絕參與以往的.
那些祭祀或者.

$^{361}$敬拜偶像的活動的時候.
可能就已經得罪了他們那些非基督徒的朋友.
另外好像剛剛信主的爺孫一樣.
信主不久就被拉到官府.
要向地方官保證保羅不再在生事.
也是一種大患難的狀況.
同樣一些有傳統猶太教背景的信徒.
他們的家族成員會認為.
不再維持祖先的傳統.
就等於對家庭不關心.
不負責任.
令到猶太的信徒左右做人難.
所以大能的意思.
可能是指信徒在日常生活中所遇到的衝突.
他們被排擠.
甚至被取笑.
第二種接受福音的狀態就是.
蒙了聖靈所賜的喜樂之後.
這段經文除了可以理解為.
聖靈所帶來的喜樂之外.
也可以解釋為.
聖靈與信徒同在成為喜樂之源.
所以第六節可以理解為.
貼薩萊加信徒成為了使徒和主的效法者.
是在於他們在患難中帶著聖靈的喜樂.
接受了福音.
而第七至第八節他就這樣說.
從此你們作了馬其頓和阿蓋亞所有信主的人的榜樣.
因為主的道已經從你們那裡傳播出去.
你們向神的信心不只在馬其頓和阿蓋亞.
就是在各處也都傳開了.
所以不用我們說什麼話.
這兩節是對貼薩萊加信徒.
在第六節的表現的一種肯定和稱讚.
當時他們身邊的母牧者.
保羅,西拉,提摩泰都不在他們身邊.
但是貼薩萊加信徒能夠成為.
馬其頓和阿蓋亞信主的人的榜樣.
是因為他們在生命上願意效法主.
貼薩萊加信徒因著效法主的生命能夠活出信仰.

$^{401}$在苦難中仍然充滿著喜樂.
接受福音並且成為別人的榜樣.
我們經常說.
生命影響生命.
貼薩萊加信徒為馬其頓和阿蓋亞的人.
展示了另一種生命力.
這種生命力讓人驚歎.
所以保羅說.
不用我們說什麼.
信徒本身的生命比傳道者的言語.
更能展現出信仰的真實.
貼薩萊加初熟的果子.
雖然沒有仕途.
沒有傳道人的身份.
但他們同樣是一種福音的使者.
因為他們不單單是自己改變了.
也透過生命的見證.
改變了很多人的生命.
洪英嘉弟兄姊妹.
三個合神心意的教會的關鍵.
可以歸納為.
能夠體現信望愛的生命.
能夠帶著照明願意毀生的生命.
能夠以效法主在大患難中.
仍有喜樂並且成為別人榜樣的生命.
當信徒能夠展現.
任何一種生命的特質.
其實都可以讓個人或教會被神所使用.
當然我們知道.
每一種的屬靈生命.
都需要我們持續不斷的屬靈的操練.
然後才能夠結出果實.
關鍵是我們願不願意花時間.
去做這些靈修操練呢?.
本來講到這裡就完結了.
但我想就貼薩萊加的背景再分享多一點.
在《仕途行傳》的記載裡.
沒有提及保羅在貼薩萊加城實際逗留了多久.
最準確的資料.
只記錄了連續三個安息日.

$^{441}$進行了傳方的工作.
之後就提到一班不信的猶太人.
召集了市中一些壞份子.
在城中引起騷動.
他們衝入耶穌家中.
要搜捕保羅.
所以保羅要即時離開貼薩萊加.
根據一些學者的研究.
估計保羅在貼薩萊加牧養信徒的日子.
是不會超過半年.
即是說.
貼薩萊加城的信徒.
由開始信耶穌.
信主.
再接受保羅的培訓.
其實最多只有半年時間左右.
接著.
保羅就離開了.
西拉不在.
提摩泰也不在.
在這樣的環境下.
貼薩萊加的信徒.
竟然可以有上述生命表現的特質.
你會不會有些好奇.
所以再追問下去.
保羅在這半年裡.
對他們怎樣培訓.
培訓了甚麼.
以致到.
牧者離開了.
仍然可以成為馬其頓.
阿蓋亞.
信徒的榜樣.
在二章十二至十三節.
保羅這樣說.
「要叫你們行事對得起.
那召你們進他國.
得他榮耀的神.
因為你們聽見我們所傳的.
神的道就領受了.

$^{481}$不是以人的道.
乃是以神的道.
這個道實在在是神的.
並且運行在你們信主的心裡」.
保羅想說甚麼呢.
保羅想說.
行事為人.
和你的召命有關.
這個召命.
是由神所發出的.
在二弗所說的四章.
保羅剛才說過.
既然蒙召.
行事為人.
就當與蒙召的恩相稱.
如果我們看整個新約.
在彼得前書二章更加說到.
你們蒙召原是為此.
因基督也為你們受過苦.
給你們留下榜樣.
叫你們追隨他的腳蹤行.
如果耶穌基督已經成為.
貼實羅來加信徒的榜樣.
而信徒有想要效法主耶穌的話.
我們就能夠.
活出一種在苦難裡有喜樂.
在召命裡體現信望愛的行為.
這是不是一個很好的解釋呢.
如果這個信息就是我們要找的答案.
只要信徒肯以效法主耶穌的心態.
去回應神對他發出的召命.
這樣每一個信徒.
都可以活出.
上述三種生命的特質.
弟兄姊妹.
我們實在很需要問自己.
今天我們為誰而活呢.
我們為自己而活.
大部分都是.
為我們的理想而活.

$^{521}$大部分人都是.
因為我們為主而活.
這是需要神對我們有額外的恩典.
最後我想引用.
龔永軒博士在聖詔出聖徒的一段說話.
結束我們今天的講道.
他說教會有四個身份.
其中一個身份就是聖潔的角度.
教會被稱為聖潔的角度.
並不是因為信徒聚得赦免而成為聖潔.
不是.
而是教會是被神分別出來的意思.
教會能夠被神分別出來.
就是成為一個聖潔的角度.
教會和世界的群體應該是不同的.
教會都會犯罪.
教會的弟兄姊妹都會犯罪.
教會所犯的罪和危險程度.
不會比其他世上的群體少.
但是神的聖潔就說明了.
神永遠都不會放棄教會.
神不會放棄將教會分別為聖.
也不會被教會因罪而消滅.
而且神永遠忠於祂對自己對教會的應許.
祂愛教會為教會社會.
教會是一個以人為單位的群體.
有人的地方必然有罪惡以及不潔.
所以教會被稱為聖潔.
不是因為有人有聖潔的本質.
而是主耶穌基督不斷將教會分別出來.
分別為聖.
將教會分別出來成為聖潔.
作為祂的一個子民.
因為主耶穌基督是一個聖者.
今天的洪恩嘉.
神仍然會不斷地將我們分別出來.
我們願不願意為主分別為聖.
叫我們的生活是為主而活的.
為主去結出一個福音的果子.
而我們願不願意毀身去回應.

$^{561}$神對我們的呼召.
我們一起去祈禱.
因為我們知道福音本是神的大能.
當我們將福音講傳的時候.
不是因為我們有什麼能力.
更不是因為我們的生命是什麼聖潔.
而是我們帶出福音的大能.
當我們願意被你使用的時候.
求主讓我們清楚明白.
不是因為我們的能力.
只是主的恩典牽著我們.
讓我們每一個願意為主服侍的人.
我們都有從主而來的那種恩典.
有從主而來的那種能力.
最重要的是主要我們同在的那種平安和喜樂.
讓我們整個生命能夠被你使用.
這是恩典.
所以主我們求你讓.
我們同一家的每一個弟兄姊妹都願意毀身去事奉你.
讓我們每一個人都願意被你使用.
以至這間教會讓人見到分別為勝的特質.
看到這間教會與這個世界不一樣.
與這個世界的價值觀不一樣.
因為這間教會所行的所展現的.
是主耶穌那種順望愛那種真善美的精神.
願主繼續祝福每一個紅人家.
我們祈禱奉教主耶穌的名 阿們.
\newpage



\section{路加福音 2:1-14}
\label{sec:7YQ_zudEJ1I}
\textbf{2020.12.20《因有一嬰孩為我們生》路加福音 2\_1-14( 和修版)  曾敏姑娘}
\newline
\newline
連結: \href{https://youtube.com/watch?v=7YQ-zudEJ1I}{\texttt{https://youtube.com/watch?v=7YQ-zudEJ1I}} ~~~~ 語音日期: 2020-12-25
\newline
\newline
\hyperref[sec:BO948tI9cPE]{\small{< < < PREV SERMON < < <}}
~
\hyperref[sec:index_chronic]{\small{[返順時目]}}
~
\hyperref[sec:index_scriptual]{\small{[返順卷目]}}
~
\hyperref[sec:p_pFg5YcRiU]{\small{> > > NEXT SERMON > > >}}
\newline
\newline
路加福音 2:1-14
\newline
\begin{longtable}{cl}
\hline
\hline
章節 & 經文 (和合本修訂版)\\
\hline
2:1 & \begin{tabularx}{0.7\textwidth}{X} 在那些日子，凱撒奧古斯都降旨，叫全國人民都登記戶籍。 \end{tabularx} \\ \\ \relax
2:2 & \begin{tabularx}{0.7\textwidth}{X} 這第一次登記戶籍是在居里扭作敘利亞總督的時候行的。 \end{tabularx} \\ \\ \relax
2:3 & \begin{tabularx}{0.7\textwidth}{X} 眾人各歸各城，辦理登記。 \end{tabularx} \\ \\ \relax
2:4 & \begin{tabularx}{0.7\textwidth}{X} 約瑟也從加利利的拿撒勒城上猶太去，到了大衛的城名叫伯利恆，因為他是大衛家族的人， \end{tabularx} \\ \\ \relax
2:5 & \begin{tabularx}{0.7\textwidth}{X} 要和他所聘之妻馬利亞一同登記戶籍。那時馬利亞已經懷孕。 \end{tabularx} \\ \\ \relax
2:6 & \begin{tabularx}{0.7\textwidth}{X} 他們在那裡的時候，馬利亞的產期到了， \end{tabularx} \\ \\ \relax
2:7 & \begin{tabularx}{0.7\textwidth}{X} 就生了頭胎的兒子，用布包起來，放在馬槽裡，因為客店裡沒有地方。 \end{tabularx} \\ \\ \relax
2:8 & \begin{tabularx}{0.7\textwidth}{X} 在伯利恆的野外有牧羊人，夜間值班看守羊群。 \end{tabularx} \\ \\ \relax
2:9 & \begin{tabularx}{0.7\textwidth}{X} 有主的一個使者站在他們旁邊，主的榮光四面照著他們，牧羊人就很懼怕。 \end{tabularx} \\ \\ \relax
2:10 & \begin{tabularx}{0.7\textwidth}{X} 那天使對他們說：「不要懼怕！看哪！因為我報給你們大喜的信息，是關乎萬民的： \end{tabularx} \\ \\ \relax
2:11 & \begin{tabularx}{0.7\textwidth}{X} 因今天在大衛的城裡，為你們生了救主，就是主基督。 \end{tabularx} \\ \\ \relax
2:12 & \begin{tabularx}{0.7\textwidth}{X} 你們要看見一個嬰孩，包著布，臥在馬槽裡，那就是給你們的記號。」 \end{tabularx} \\ \\ \relax
2:13 & \begin{tabularx}{0.7\textwidth}{X} 忽然，有一大隊天兵同那天使讚美神說： \end{tabularx} \\ \\ \relax
2:14 & \begin{tabularx}{0.7\textwidth}{X} 「在至高之處榮耀歸與神！在地上平安歸與他所喜悅的人！」 \end{tabularx} \\ \\
[1ex]
\hline
\hline
\end{longtable}
$^{1}$請用餐具和食物的盤子.
放在桌上.
請用餐具和食物的盤子.
放在桌上.
請用餐具和食物的盤子.
放在桌上.
請用餐具和食物的盤子.
放在桌上.
請用餐具和食物的盤子.
放在桌上.
請用餐具和食物的盤子.
放在桌上.
請用餐具和食物的盤子.
放在桌上.
請用餐具和食物的盤子.
放在桌上.
請用餐具和食物的盤子.
放在桌上.
請用餐具和食物的盤子.
放在桌上.
請用餐具和食物的盤子.
放在桌上.
請用餐具和食物的盤子.
放在桌上.
請用餐具和食物的盤子.
放在桌上.
請用餐具和食物的盤子.
放在桌上.
請用餐具和食物的盤子.
放在桌上.
請用餐具和食物的盤子.
放在桌上.
請用餐具和食物的盤子.
放在桌上.
請用餐具和食物的盤子.
放在桌上.
請用餐具和食物的盤子.
放在桌上.
請用餐具和食物的盤子.
放在桌上.

$^{41}$請用餐具和食物的盤子.
放在桌上.
請用餐具和食物的盤子.
放在桌上.
請用餐具和食物的盤子.
放在桌上.
請用餐具和食物的盤子.
放在桌上.
請用餐具和食物的盤子.
放在桌上.
請用餐具和食物的盤子.
放在桌上.
請用餐具和食物的盤子.
放在桌上.
請用餐具和食物的盤子.
放在桌上.
請用餐具和食物的盤子.
放在桌上.
請用餐具和食物的盤子.
放在桌上.
請用餐具和食物的盤子.
放在桌上.
請用餐具和食物的盤子.
放在桌上.
請用餐具和食物的盤子.
放在桌上.
請用餐具和食物的盤子.
放在桌上.
請用餐具和食物的盤子.
放在桌上.
請用餐具和食物的盤子.
放在桌上.
請用餐具和食物的盤子.
放在桌上.
請用餐具和食物的盤子.
放在桌上.
請用餐具和食物的盤子.
放在桌上.
請用餐具和食物的盤子.
放在桌上.

$^{81}$請用餐具和食物的盤子.
放在桌上.
請用餐具和食物的盤子.
放在桌上.
請用餐具和食物的盤子.
放在桌上.
請用餐具和食物的盤子.
放在桌上.
請用餐具和食物的盤子.
放在桌上.
請用餐具和食物的盤子.
放在桌上.
請用餐具和食物的盤子.
放在桌上.
請用餐具和食物的盤子.
放在桌上.
請用餐具和食物的盤子.
放在桌上.
請用餐具和食物的盤子.
放在桌上.
請用餐具和食物的盤子.
放在桌上.
請用餐具和食物的盤子.
放在桌上.
請用餐具和食物的盤子.
放在桌上.
請用餐具和食物的盤子.
放在桌上.
請用餐具和食物的盤子.
放在桌上.
請用餐具和食物的盤子.
放在桌上.
請用餐具和食物的盤子.
放在桌上.
請用餐具和食物的盤子.
放在桌上.
請用餐具和食物的盤子.
放在桌上.
請用餐具和食物的盤子.
放在桌上.

$^{121}$請用餐具和食物的盤子.
放在桌上.
請用餐具和食物的盤子.
放在桌上.
請用餐具和食物的盤子.
放在桌上.
請用餐具和食物的盤子.
放在桌上.
請用餐具和食物的盤子.
放在桌上.
請用餐具和食物的盤子.
放在桌上.
請用餐具和食物的盤子.
放在桌上.
請用餐具和食物的盤子.
放在桌上.
請用餐具和食物的盤子.
放在桌上.
請用餐具和食物的盤子.
放在桌上.
請用餐具和食物的盤子.
放在桌上.
請用餐具和食物的盤子.
放在桌上.
請用餐具和食物的盤子.
放在桌上.
請用餐具和食物的盤子.
放在桌上.
請用餐具和食物的盤子.
放在桌上.
請用餐具和食物的盤子.
放在桌上.
請用餐具和食物的盤子.
放在桌上.
請用餐具和食物的盤子.
放在桌上.
請用餐具和食物的盤子.
放在桌上.
請用餐具和食物的盤子.
放在桌上.

$^{161}$請用餐具和食物的盤子.
放在桌上.
請用餐具和食物的盤子.
放在桌上.
請用餐具和食物的盤子.
放在桌上.
請用餐具和食物的盤子.
放在桌上.
請用餐具和食物的盤子.
放在桌上.
請用餐具和食物的盤子.
放在桌上.
請用餐具和食物的盤子.
放在桌上.
請用餐具和食物的盤子.
放在桌上.
請用餐具和食物的盤子.
放在桌上.
請用餐具和食物的盤子.
放在桌上.
請用餐具和食物的盤子.
放在桌上.
請用餐具和食物的盤子.
放在桌上.
請用餐具和食物的盤子.
放在桌上.
請用餐具和食物的盤子.
放在桌上.
請用餐具和食物的盤子.
放在桌上.
請用餐具和食物的盤子.
放在桌上.
請用餐具和食物的盤子.
放在桌上.
請用餐具和食物的盤子.
放在桌上.
請用餐具和食物的盤子.
放在桌上.
請用餐具和食物的盤子.
放在桌上.

$^{201}$請用餐具和食物的盤子.
放在桌上.
請用餐具和食物的盤子.
放在桌上.
請用餐具和食物的盤子.
放在桌上.
請用餐具和食物的盤子.
放在桌上.
請用餐具和食物的盤子.
放在桌上.
請用餐具和食物的盤子.
放在桌上.
請用餐具和食物的盤子.
放在桌上.
請用餐具和食物的盤子.
放在桌上.
請用餐具和食物的盤子.
放在桌上.
請用餐具和食物的盤子.
放在桌上.
請用餐具和食物的盤子.
放在桌上.
請用餐具和食物的盤子.
放在桌上.
請用餐具和食物的盤子.
放在桌上.
請用餐具和食物的盤子.
放在桌上.
請用餐具和食物的盤子.
放在桌上.
請用餐具和食物的盤子.
放在桌上.
請用餐具和食物的盤子.
放在桌上.
請用餐具和食物的盤子.
放在桌上.
請用餐具和食物的盤子.
放在桌上.
請用餐具和食物的盤子.
放在桌上.

$^{241}$請用餐具和食物的盤子.
放在桌上.
請用餐具和食物的盤子.
放在桌上.
請用餐具和食物的盤子.
放在桌上.
請用餐具和食物的盤子.
放在桌上.
請用餐具和食物的盤子.
放在桌上.
請用餐具和食物的盤子.
放在桌上.
請用餐具和食物的盤子.
放在桌上.
請用餐具和食物的盤子.
放在桌上.
請用餐具和食物的盤子.
放在桌上.
請用餐具和食物的盤子.
放在桌上.
請用餐具和食物的盤子.
放在桌上.
請用餐具和食物的盤子.
放在桌上.
請用餐具和食物的盤子.
放在桌上.
請用餐具和食物的盤子.
放在桌上.
請用餐具和食物的盤子.
放在桌上.
請用餐具和食物的盤子.
放在桌上.
請用餐具和食物的盤子.
放在桌上.
請用餐具和食物的盤子.
放在桌上.
請用餐具和食物的盤子.
放在桌上.
請用餐具和食物的盤子.
放在桌上.

$^{281}$請用餐具和食物的盤子.
放在桌上.
請用餐具和食物的盤子.
放在桌上.
請用餐具和食物的盤子.
放在桌上.
請用餐具和食物的盤子.
放在桌上.
請用餐具和食物的盤子.
放在桌上.
請用餐具和食物的盤子.
放在桌上.
請用餐具和食物的盤子.
放在桌上.
請用餐具和食物的盤子.
放在桌上.
請用餐具和食物的盤子.
放在桌上.
請用餐具和食物的盤子.
放在桌上.
請用餐具和食物的盤子.
放在桌上.
請用餐具和食物的盤子.
放在桌上.
請用餐具和食物的盤子.
放在桌上.
請用餐具和食物的盤子.
放在桌上.
請用餐具和食物的盤子.
放在桌上.
請用餐具和食物的盤子.
放在桌上.
請用餐具和食物的盤子.
放在桌上.
請用餐具和食物的盤子.
放在桌上.
請用餐具和食物的盤子.
放在桌上.
請用餐具和食物的盤子.
放在桌上.

$^{321}$請用餐具和食物的盤子.
放在桌上.
請用餐具和食物的盤子.
放在桌上.
請用餐具和食物的盤子.
放在桌上.
請用餐具和食物的盤子.
放在桌上.
請用餐具和食物的盤子.
放在桌上.
請用餐具和食物的盤子.
放在桌上.
請用餐具和食物的盤子.
放在桌上.
請用餐具和食物的盤子.
放在桌上.
請用餐具和食物的盤子.
放在桌上.
請用餐具和食物的盤子.
放在桌上.
請用餐具和食物的盤子.
放在桌上.
請用餐具和食物的盤子.
放在桌上.
請用餐具和食物的盤子.
放在桌上.
請用餐具和食物的盤子.
放在桌上.
請用餐具和食物的盤子.
放在桌上.
請用餐具和食物的盤子.
放在桌上.
請用餐具和食物的盤子.
放在桌上.
請用餐具和食物的盤子.
放在桌上.
請用餐具和食物的盤子.
放在桌上.
請用餐具和食物的盤子.
放在桌上.

$^{361}$請用餐具和食物的盤子.
放在桌上.
請用餐具和食物的盤子.
放在桌上.
請用餐具和食物的盤子.
放在桌上.
請用餐具和食物的盤子.
放在桌上.
請用餐具和食物的盤子.
放在桌上.
請用餐具和食物的盤子.
放在桌上.
請用餐具和食物的盤子.
放在桌上.
請用餐具和食物的盤子.
放在桌上.
請用餐具和食物的盤子.
放在桌上.
請用餐具和食物的盤子.
放在桌上.
請用餐具和食物的盤子.
放在桌上.
請用餐具和食物的盤子.
放在桌上.
請用餐具和食物的盤子.
放在桌上.
請用餐具和食物的盤子.
放在桌上.
請用餐具和食物的盤子.
放在桌上.
請用餐具和食物的盤子.
放在桌上.
請用餐具和食物的盤子.
放在桌上.
請用餐具和食物的盤子.
放在桌上.
請用餐具和食物的盤子.
放在桌上.
請用餐具和食物的盤子.
放在桌上.

$^{401}$請用餐具和食物的盤子.
放在桌上.
請用餐具和食物的盤子.
放在桌上.
請用餐具和食物的盤子.
放在桌上.
請用餐具和食物的盤子.
放在桌上.
請用餐具和食物的盤子.
放在桌上.
請用餐具和食物的盤子.
放在桌上.
請用餐具和食物的盤子.
放在桌上.
請用餐具和食物的盤子.
放在桌上.
請用餐具和食物的盤子.
放在桌上.
請用餐具和食物的盤子.
放在桌上.
請用餐具和食物的盤子.
放在桌上.
請用餐具和食物的盤子.
放在桌上.
請用餐具和食物的盤子.
放在桌上.
請用餐具和食物的盤子.
放在桌上.
請用餐具和食物的盤子.
放在桌上.
請用餐具和食物的盤子.
放在桌上.
請用餐具和食物的盤子.
放在桌上.
請用餐具和食物的盤子.
放在桌上.
請用餐具和食物的盤子.
放在桌上.
請用餐具和食物的盤子.
放在桌上.

$^{441}$請用餐具和食物的盤子.
放在桌上.
請用餐具和食物的盤子.
放在桌上.
請用餐具和食物的盤子.
放在桌上.
請用餐具和食物的盤子.
放在桌上.
請用餐具和食物的盤子.
放在桌上.
請用餐具和食物的盤子.
放在桌上.
請用餐具和食物的盤子.
放在桌上.
請用餐具和食物的盤子.
放在桌上.
請用餐具和食物的盤子.
放在桌上.
請用餐具和食物的盤子.
放在桌上.
請用餐具和食物的盤子.
放在桌上.
請用餐具和食物的盤子.
放在桌上.
請用餐具和食物的盤子.
放在桌上.
請用餐具和食物的盤子.
放在桌上.
請用餐具和食物的盤子.
放在桌上.
請用餐具和食物的盤子.
放在桌上.
請用餐具和食物的盤子.
放在桌上.
請用餐具和食物的盤子.
放在桌上.
請用餐具和食物的盤子.
放在桌上.
請用餐具和食物的盤子.
放在桌上.

$^{481}$請用餐具和食物的盤子.
放在桌上.
請用餐具和食物的盤子.
放在桌上.
請用餐具和食物的盤子.
放在桌上.
請用餐具和食物的盤子.
放在桌上.
請用餐具和食物的盤子.
放在桌上.
請用餐具和食物的盤子.
放在桌上.
請用餐具和食物的盤子.
放在桌上.
請用餐具和食物的盤子.
放在桌上.
請用餐具和食物的盤子.
放在桌上.
請用餐具和食物的盤子.
放在桌上.
請用餐具和食物的盤子.
放在桌上.
請用餐具和食物的盤子.
放在桌上.
請用餐具和食物的盤子.
放在桌上.
請用餐具和食物的盤子.
放在桌上.
請用餐具和食物的盤子.
放在桌上.
請用餐具和食物的盤子.
放在桌上.
請用餐具和食物的盤子.
放在桌上.
請用餐具和食物的盤子.
放在桌上.
請用餐具和食物的盤子.
放在桌上.
請用餐具和食物的盤子.
放在桌上.

$^{521}$請用餐具和食物的盤子.
放在桌上.
請用餐具和食物的盤子.
放在桌上.
請用餐具和食物的盤子.
放在桌上.
請用餐具和食物的盤子.
放在桌上.
請用餐具和食物的盤子.
放在桌上.
請用餐具和食物的盤子.
放在桌上.
請用餐具和食物的盤子.
放在桌上.
請用餐具和食物的盤子.
放在桌上.
請用餐具和食物的盤子.
放在桌上.
請用餐具和食物的盤子.
放在桌上.
請用餐具和食物的盤子.
放在桌上.
請用餐具和食物的盤子.
放在桌上.
請用餐具和食物的盤子.
放在桌上.
請用餐具和食物的盤子.
放在桌上.
請用餐具和食物的盤子.
放在桌上.
請用餐具和食物的盤子.
放在桌上.
請用餐具和食物的盤子.
放在桌上.
請用餐具和食物的盤子.
放在桌上.
請用餐具和食物的盤子.
放在桌上.
請用餐具和食物的盤子.
放在桌上.

$^{561}$請用餐具和食物的盤子.
放在桌上.
請用餐具和食物的盤子.
放在桌上.
請用餐具和食物的盤子.
放在桌上.
請用餐具和食物的盤子.
放在桌上.
請用餐具和食物的盤子.
放在桌上.
請用餐具和食物的盤子.
放在桌上.
請用餐具和食物的盤子.
放在桌上.
請用餐具和食物的盤子.
放在桌上.
請用餐具和食物的盤子.
放在桌上.
請用餐具和食物的盤子.
放在桌上.
請用餐具和食物的盤子.
放在桌上.
請用餐具和食物的盤子.
放在桌上.
請用餐具和食物的盤子.
放在桌上.
請用餐具和食物的盤子.
放在桌上.
請用餐具和食物的盤子.
放在桌上.
請用餐具和食物的盤子.
放在桌上.
請用餐具和食物的盤子.
放在桌上.
請用餐具和食物的盤子.
放在桌上.
請用餐具和食物的盤子.
放在桌上.
請用餐具和食物的盤子.
放在桌上.

$^{601}$請用餐具和食物的盤子.
放在桌上.
請用餐具和食物的盤子.
放在桌上.
\newpage



\section{詩篇 103:1}
\label{sec:p_pFg5YcRiU}
\textbf{2020.12.27《心的稱頌》詩篇 103\_1 - 5( 和修版)  陳麗蟬牧師}
\newline
\newline
連結: \href{https://youtube.com/watch?v=p-pFg5YcRiU}{\texttt{https://youtube.com/watch?v=p-pFg5YcRiU}} ~~~~ 語音日期: 2020-12-27
\newline
\newline
\hyperref[sec:7YQ_zudEJ1I]{\small{< < < PREV SERMON < < <}}
~
\hyperref[sec:index_chronic]{\small{[返順時目]}}
~
\hyperref[sec:index_scriptual]{\small{[返順卷目]}}
~
\hyperref[sec:r3qTUjnU0eY]{\small{> > > NEXT SERMON > > >}}
\newline
\newline
詩篇 103:1
\newline
\begin{longtable}{cl}
\hline
\hline
章節 & 經文 (和合本修訂版)\\
\hline
103:1 & \begin{tabularx}{0.7\textwidth}{X} 我的心哪，你要稱頌耶和華！凡在我裡面的，都要稱頌他的聖名！ \end{tabularx} \\ \\
[1ex]
\hline
\hline
\end{longtable}
$^{1}$請自我介紹.
《2020年國慶日》主題曲 張碩晟議員.
各位姐妹早安.
今天是2020年最後一次的主日崇拜.
我竟然有幸來到紅恩堂.
跟大家一齊歌頌讚美神.
用上帝的話來給予鼓勵.
在中國有一個很出名的商人.
他叫做張碩敏.
不過可能你和我都不太認識.
我也是靠上網才能認識到這個人.
他將中國的商品帶到歐洲的市場.
原來在70年代的時候.
在歐洲,美國各方面都沒有中國的商品.
在當中的市場.
是他第一個帶進去的.
他做事非常積極.
他說每一樣簡單的事.
只要你用心去做好的時候.
它可能會成為一個不簡單的事.
一件很平凡的事.
只要我們用心去做好的時候.
這件事可以變得不平凡.
所以他鼓勵所有的員工.
做事的時候只要用心去做的話.
得出來的果效是很不同的.
有三個工人在砌磚.
有人問他們三個.
你們在做什麼呢?.
第一個工人就好戲不好戲.
你看不到的嗎?.
我在砌磚.
這個工人只是在做事.
第二個就回答他.
我每一個小時可以拿到160元的工.
好一點.
因為要看錢份上.
第三個工人就很開心.
很輕鬆地說.
我現在在建造一間世界上最偉大的樓宇.

$^{41}$我想這個第三個工人.
他有一個遠大的目標.
他心裡面有一份積極的時候.
他更加用心去做.
他出來的質量和效果是很不同的.
如果我們像第二個工人一樣.
為了錢去看的時候.
他都會盡責的.
不過可能那個心態就是.
誰叫你窮頂硬上呢?.
不行不行.
我們要有一個遠大的目標.
我們用心去做好每一樣上帝給我們的事情.
你去看這個詩篇.
第103篇第一召第五節的時候.
我們知道他一點點的背景.
很清楚知道他是大衛的作品.
因為其中一句他說.
他用美麗使你所願的得以知足.
以至你如英返老還童.
在第五節他這樣說.
所以人們覺得可能大衛已經是練腦了.
但是他因為堅心去跟著神的時候.
他覺得何時他去投靠神的時候.
他的力量就來了.
所以人們就猜了.
這個是在大衛練腦的時候寫成的.
並且在他的兒子叛亂之後.
也在他因為驕傲犯了罪之後.
上帝懲罰了之後.
他就寫成了這一篇感恩的詩篇.
我們就看一下.
在第一方面我們看一下心的稱頌.
我們這次就只看心那方面.
第一第二節大衛這樣說.
我的心你要稱頌耶和華.
凡在我裡面的都要稱頌他的名.
第二節也是同樣這樣說.
我的心你要稱頌耶和華.
所以這兩節聖經裡面.

$^{81}$不斷地說我的心你要稱頌.
出現了三次稱頌這個字.
在聖經裡面凡出現三次的字.
就說明他是很確實的.
是不會變的.
就如我們看.
在《熱賽書》.
當熱賽先知見到上帝的時候.
他就說.
聖哉聖哉聖哉萬軍的耶和華.
他就很絕對肯定.
我們的上帝是一位聖潔的上帝.
在這裡出現了三次.
就是說大衛很清楚.
也永遠都不會變的.
他的心要去稱頌上帝.
這篇詩篇有一個特別的.
就是他用這個心.
他邀請他自己的心.
來去對上帝的感恩.
心是指什麼呢?.
他原本的文字是說到.
是指到生命.
整個生命.
整個人.
是裡面包括了我們裡面的感情.
那個意志.
和我們整個人裡面的知識所在乎.
其實這個希伯來人和我們中國人都差不多一樣.
我們什麼時候督實小孩子.
給一點心機去做事.
不要在玩來玩去.
不要心在做事.
手在玩玩具.
眼又看電視.
不行.
當我們說用這個心的時候.
就指到整個人.
整個人.
都要來去稱頌我們的上帝.

$^{121}$好了.
聖經.
耶穌基督都說過.
我們要用心靈和誠實去敬拜上帝.
當他和井旁的婦人.
在討論究竟要在基里森山那裡來敬拜.
還是在耶路撒冷來敬拜的時候.
他就說了.
最重要就是用心靈和誠實來去敬拜祂.
就是說.
要用我們整個生命.
按著真理.
用整個生命來去讚美.
敬拜我們的上帝.
整個人.
當我們用心.
我們整個人去做每一樣事情.
出來的效果.
很不一樣的.
是不是?.
我記得.
有時候我放假.
一放假的時候.
我會和我家人去爬山做運動.
有時候看到那些人.
站在一邊.
在那裡做早操.
我們就走過去做.
我就看著師父.
師父本來就帶著那些人說.
一二三四.
說著說著.
他就改了不同的字.
一二三出來.
為什麼說著說著就不同了呢?.
我後來看看周圍環境.
每個人天天去習慣的.
就很用心去做運動.
我們這些懶惰豬.
久不久才走一次的.

$^{161}$我們就怎樣呢?.
一二三四.
可能師父不值得.
他就一二三.
你可以說出力.
你不夠膽就不出力了.
我就明白了.
有時候有些事情.
當我們出力.
盡力去做的時候.
得出來的果效是不同的.
當我們真真正正.
出力用心去做運動的時候.
那個果效就會出來.
但當我們孖孖呼呼的時候.
做的都是這樣.
那我們怎樣呢?.
我們要用心去稱頌我們的上帝.
用心去敬拜我們的上帝.
就是當我們星期日.
早上要崇拜了.
我們星期六早一點睡覺.
星期日早一點出門.
然後很安靜.
很欣賞在紅恩堂裡面.
有些姐妹.
或者平時都見到有些姐妹.
很早就會來到這裡.
他們安靜心靈.
他們來禱告.
看到嶺昌,詩琴那些.
一路很早就回來了.
他們就怎樣呢?.
準備他們要做的事情.
準備他們的心.
他們來去敬拜.
要來去稱頌.
好了.
第二方面.
都是心的稱頌.

$^{201}$但這次我們看「稱頌」這個字.
都是心要稱頌耶和華.
犯在我裡面的.
也要來去稱頌他的聖名.
「稱頌」這個字.
在英文「稱頌他的聖名」.
在英文裡面怎樣翻譯呢?.
就是Bless his holy name.
祝福他的名.
我們覺得.
我們哪裡有資格.
去祝福上帝的名啊.
是不是?.
我們沒有這樣的資格.
但是呢.
「祝福」這個字.
或者「稱頌」這個字.
在舊約或者新約裡面.
它的文字.
都是這樣.
就是將我們每一個人.
我們有我們自己成功的地方.
或者我們有我們的財產.
有我們的資源.
總之呢.
就將我自己所擁有的東西.
全數最好的最好.
都交出來給上帝.
將我們最好最好所有的東西.
都交給上帝.
就是說.
要稱頌上帝.
稱頌他的聖名.
就說將最好的東西.
最好的禮物.
總之所有的東西.
都交回給上帝.
但是.
上帝他自己.
本身已經是.

$^{241}$充滿了榮耀.
充滿了尊貴.
還有什麼東西.
是很好的給過他呢?.
《南斯篇》五十篇都這樣提醒我們.
臨終的白獸是我的.
千山的生畜都是我的.
山上高的飛鳥.
我都知道.
田野的走獸.
也都是屬於我的.
就算我肚餓.
我也都不需要告訴你.
因為世界和其中所充滿的.
都是屬於我的.
所以呢.
什麼都是上帝的.
我們有什麼東西可以給我們的上帝呢?.
香港在近十多年.
都經歷很多次的疫情.
我記得印象最深刻的一次.
就是這個禽流感.
那時候整個香港很害怕.
因為所有的雞都殺光了.
全個香港沒有雞吃過.
我記得.
我很清楚很印象很深刻.
有一次坐小巴的時候.
後面就坐了一些.
可能拜慣神的一些女士.
她就說.
糟了.
全香港的雞都殺光了.
沒有雞拜神.
那些人就說.
怕拜神的時候.
宰一隻鴨子.
不知道合不合她的口味.
吃了這麼多年雞.
那就真是慘了.

$^{281}$那另外旁邊的那位太太就說.
是的.
不過我才知道.
我封封大禮事給她.
任由她自己喜歡買什麼.
當我們這樣想.
我想拜偶像的人.
就會很擔心那些東西不合她的口味.
很聰明的人就封封大禮事給她.
她更聰明的就把大禮事放在那裡.
拜完神之後就放回自己的袋子裡.
但是我們的上帝.
祂說這些東西全部都是屬於她的.
她根本不需要.
在詩篇這個十五篇就提醒我們.
「他凡以感恩獻祭的.
就是榮耀我.
那按正路而行的.
我必使他得著神的救恩」.
原來上帝要我們最好的禮物.
就是將我們的心.
將我們整個人.
都呈獻給上帝.
將我們的意志.
我們的感情.
我們的情緒.
都交給上帝.
縮小範圍一點.
在敬拜的時候.
我們怎樣將我們整個心.
整個人交給上帝呢?.
我們有一個決心.
就是在崇拜的裡面.
無論我之前在崇拜的時候.
發生了什麼事情.
我都要決心.
專心一致的去讚美.
去歌頌.
去聆聽.
我記得在我做會友的時候.

$^{321}$有位導師.
他因為去越南難民營那裡.
去做工作.
就染了一些菌.
很快就去世了.
但我很記得.
有一幕在兩三個星期之後.
我們團體就出外.
街頭派單張.
他太太帶著肚子.
堅持和我們一起.
出到街頭去暴動.
我記得當天我看到的時候.
很感動.
讓我看到一件事.
就是什麼是整個生命.
交給上帝.
他不會因著自己的丈夫.
離開了世界.
他就不參與那次的暴動.
他堅持繼續的.
去有這個福音的行動.
所以當我們在崇拜的時候.
我都要學習這種態度.
無論之前發生什麼事情.
我都用一個最好的精神狀態.
最好的心境.
我們來去敬拜.
在唱詩的時候.
我究竟是從心裡對上帝.
來發出讚美.
還是我只是在唸口往.
我們在做面部的運動呢.
或者當我們在聽道的時候.
有沒有專心的去聆聽.
從上帝而來的說話呢.
其實為了紅印堂很感恩.
因為大家很關心.
看著這個紅印堂.
來進行崇拜.

$^{361}$但很多教會都鼓勵.
弟兄姊妹的手機.
用手機看程序表.
用手機看經文.
其實是一個很大的引誘.
當一邊看著手機的時候.
究竟看著歌詞.
看著經文呢.
還是看著自己的Facebook.
看著自己的WhatsApp呢.
有時候我們教會.
在我們錦繡堂當值的時候.
我們坐在最後一排的.
我們就看到前面.
起碼近我們一兩行的弟兄姊妹.
她們在回覆WhatsApp.
她們在做什麼.
當然當我去觀察人家.
在做什麼的時候.
其實也代表我不太專心崇拜.
求神赦免.
但是我覺得是一個.
我們要去學習一件事.
無論其他人怎麼樣.
我們自己去忠心的.
整個人的.
將我們自己最好的.
交給上帝.
第三方面.
就是心的感恩.
我的心.
你要稱頌耶和華.
耶和華他創造萬有.
他是一位非常豐富.
也是極之卓越的上帝.
沒錯.
他就是什麼都有.
不過整個的正義.
就是說我們要將我們的心智.
將我們最好的精神狀態.

$^{401}$將我們整個的生命.
來去交給上帝.
我們不需要去和其他人比較.
哪個人好一點.
哪個人沒有那麼好.
我又怎麼樣.
不需要的.
上帝所看的.
只是看我們自己獨立的人.
有沒有做到最好.
來去呈獻給他.
所以在這裡說.
我們要去稱頌我們的耶和華.
那麼上帝究竟有什麼那麼好的呢.
大衛就繼續說.
他赦免你一切的罪孽.
醫治你一切的疾病.
他也是救贖我們的命.
脫離死亡.
大衛他自己去經歷到.
當他犯了姦淫.
和畢氏巴犯了姦淫之後.
他不單犯了姦淫罪.
他同樣犯了殺人罪.
因為他怕事情被釋穿.
所以就將女子的丈夫.
調到最危險的地方.
讓人去殺死.
他同樣犯了殺人罪.
另外他自己驕傲.
當他整個國家慢慢穩定了.
他又想數一下自己.
究竟壯丁有多少人.
以致將來可以怎麼打仗.
他犯了這個驕傲的罪.
他不聽耶和華的勸告.
他自己去數一下.
究竟我的軍力有多厲害.
結果上帝又懲罰他.
以致發生了瘟疫.

$^{441}$在一天之內.
傳過的地方死了七萬人.
所以他經歷到自己犯了罪.
而帶來了很多人的死亡.
但當他知道犯下了彌天大罪之後.
他向上帝認罪悔改.
上帝就釋放他的心靈.
所以當他犯了姦淫罪.
又殺了別人的丈夫之後.
上帝就懲罰他.
叫他因姦淫而生的兒子.
很快就去世了.
但他很難過.
有一段時間他知道兒子已經不行了.
回到天家之後.
他很快重新得力.
所以他經歷過.
上帝耶和華赦免他一切的罪孽.
醫治他一切的疾病.
也救贖他的命脫離死亡.
雖然他一直在講自己過往的生命經歷.
但他也在講一件事.
就是彌賽亞救贖我們每一個人脫離死亡的權勢.
凡是信靠耶穌基督的.
都可以和耶穌基督一樣有復活的生命.
這是大衛自己的經歷.
我們也看到.
在現今這個世代裡面.
我們看到有很多人犯了彌天大罪.
但上帝也讓他們的心靈得到釋放.
當然他們在地上高.
他們同樣都要受法律的制裁.
有時我去以林書室買東西.
買書買禮物.
有一次我去到的時候.
我看到一個人很面熟.
是誰呢?.
很有禮貌.
不過樣子也有一點兇.
但後來想了想.

$^{481}$是的 我記得了.
他們那裡有些基督教的機構.
他們都會找一些坐牢的.
甚至是終身監禁.
或者提早釋放的人在那裡工作.
我認得那位弟兄.
以前看新聞.
看到他犯了什麼罪的時候.
覺得很討厭.
但我覺得在監獄信了主之後.
在書室面前.
我看到的是一個很斯文.
很有禮貌的弟兄.
他在那裡介紹一些貨品.
我覺得這就是我們的上帝.
去赦免我們一切的罪孽.
當然他也受到制裁.
他要坐牢.
但當他完成了刑罰之後.
他出來又是一個新的生命.
除了這樣之外.
他還說以仁愛慈悲為我們的冠冕.
他在這裡說到.
上帝給我們每一個人.
有不同的專長.
我經常說很欣賞老傳道的行程.
一拍下去.
我就走到一邊去了.
我們看到這裡.
我們看到有彈琴的恩賜.
有領唱的恩賜.
有做電腦做音響的恩賜.
但當上帝給我們運用這些恩賜的時候.
我們發覺感覺是很不同的.
如果我們弟兄姊妹.
我們有機會將我們從上帝而得到的那份專長.
那份恩賜.
我們用出來事奉的時候.
我們生命會覺得有那份很滿足.
很幸福.

$^{521}$我覺得被上帝運用那份感覺.
整個人都會醒一些.
所以他說以仁愛慈悲為我們的冠冕.
他用美物使你所願的得而滿足.
以致你如英返老還童.
說到在上帝的裡面.
看到我們每一個人的身心靈都能夠得到力量.
在每一方面他都會給過我們.
以致我們的生命會不斷的更新.
所以弟兄姊妹不用怕自己老了.
老了一年過了2020年.
不要緊.
只要我們的生命在耶穌基督裡面不斷的更新的時候.
每一個人都可以如英返老還童.
在1972年的時候.
新加坡的觀光局局長就交了一份報告給李光耀總理.
報告就說.
他說我們新加坡這個地方甚麼都沒有.
他說人家埃及都有一個金字塔.
中國又有萬里長城.
日本又有一個富士山.
美國的夏威夷又有陽光與海灘.
所以他們可以好好發展這個旅遊業.
不過我們新加坡甚麼都沒有.
完全的那些明星古蹟甚麼都沒有.
一年四季那個日頭就曬到不得了.
所以你說我們靠旅遊業來打理這個經濟的時候.
實在沒有這樣的可能.
李光耀總理看到這份報告的時候非常生氣.
我想像一下會不會把這份報告扔在地上.
他就在這份報告裡面寫了一行字.
他說上帝給了我們每天都有燦爛的陽光.
已經非常的足夠了.
你還想要甚麼呢?.
然後他們就要重新整理那份報告.
重新計劃所有的事.
最後我們就利用上帝給新加坡一年四季都有燦爛的陽光.
他們開始種花養草.
以致到沒多久幾年之後.
新加坡就成為了世界上很出名的花園城市.

$^{561}$新加坡給我們的感覺很多花花草草.
燦爛的陽光也都很乾淨.
我想李光耀總理這樣就抓到上帝給他們這個地方.
他們最重要的是甚麼的時候.
他懂得怎樣去運用.
所以我們懂事的.
他把上帝給我們恩賜來運用這件事是很重要的.
一個不懂得感恩的人.
就算上帝給他很多恩賜也好.
給他很多機會也好.
他就經常覺得.
我這行政沒有我老傳道那麼好.
真的沒甚麼用.
我做主席也沒有我Steven做得那麼好.
他經常都在一個很憂鬱的情況裡面.
所以我們要知道一件事.
這個感恩的生命.
不在於我自己生命裡面擁有多少的恩賜.
擁有多少的才幹.
或者擁有多少的物質.
而是一種的態度.
就是我們把握著上帝給我的那份的恩賜.
把握著那件事.
我就好好地去運用.
有人就說.
感恩是幸福的約事.
我們懂得去感恩的時候.
拿著這把感恩的鎖匙.
我們開了那扇門.
在裡面充滿了各樣的幸福.
香港人就習慣在年尾的時候.
就會數算上帝給我們的恩典.
現在距離2021年還有四天左右.
我們還有四天的時間.
可以數算一下過往一年的裡面.
上帝給了我們甚麼.
我細細地逐樣逐樣地數出來.
不要只是猛數著口罩,隔板,限聚令.
又不可以怎樣.
我相信在每一件的事情.

$^{601}$我們用上帝的眼光.
我們一定可以數到上帝很多的恩典.
以致我們藉著不斷的感恩.
我們可以能夠打開2021年.
一個幸福的一年.
我們同心吐構.
天父上帝我們多謝你.
我們要一生的稱頌你.
我們要整個生命來歌頌你.
你實在是在我們每一個人的生命裡面.
給我們很多很多的恩賜.
給我們有很多很多的經歷.
全部這些經歷我們都知道是從你而來.
我們知道我們所擁有的一切.
我們的好處不在你以外.
我們多謝你.
求你賜福給我們每一位的兄弟姐妹.
我們的同工.
我們懂得事事數算主理的恩典.
我們不再埋怨其他人很多.
其他人又怎樣.
我們只是從我們的生命裡面.
我們看到上帝你大能的作為.
求主幫助.
讓我們整個人的生命.
我們整個心時時刻刻的稱頌你.
讚美你.
我們願意你得到至高的榮耀和讚美.
多謝你誓聽禱告.
去奉靠主耶穌基督得勝的名字而求.
Amen.
\newpage



\section{約翰一書 4:7-12}
\label{sec:r3qTUjnU0eY}
\textbf{2021.01.03《生命欠一個奇蹟》約翰一書 4\_7-12( 和修版)  老耀雄傳道}
\newline
\newline
連結: \href{https://youtube.com/watch?v=r3qTUjnU0eY}{\texttt{https://youtube.com/watch?v=r3qTUjnU0eY}} ~~~~ 語音日期: 2021-01-04
\newline
\newline
\hyperref[sec:p_pFg5YcRiU]{\small{< < < PREV SERMON < < <}}
~
\hyperref[sec:index_chronic]{\small{[返順時目]}}
~
\hyperref[sec:index_scriptual]{\small{[返順卷目]}}
~
\hyperref[sec:kzUx2k7DXdQ]{\small{> > > NEXT SERMON > > >}}
\newline
\newline
約翰一書 4:7-12
\newline
\begin{longtable}{cl}
\hline
\hline
章節 & 經文 (和合本修訂版)\\
\hline
4:7 & \begin{tabularx}{0.7\textwidth}{X} 親愛的，我們要彼此相愛，因為愛是從神來的。凡有愛的都是由神而生，並且認識神。 \end{tabularx} \\ \\ \relax
4:8 & \begin{tabularx}{0.7\textwidth}{X} 沒有愛的就不認識神，因為神就是愛。 \end{tabularx} \\ \\ \relax
4:9 & \begin{tabularx}{0.7\textwidth}{X} 神差他獨一的兒子到世上來，使我們藉著他得生命；由此，神對我們的愛就顯明了。 \end{tabularx} \\ \\ \relax
4:10 & \begin{tabularx}{0.7\textwidth}{X} 不是我們愛神，而是神愛我們，差他的兒子為我們的罪作了贖罪祭；這就是愛。 \end{tabularx} \\ \\ \relax
4:11 & \begin{tabularx}{0.7\textwidth}{X} 親愛的，既然神這樣愛我們，我們也要彼此相愛。 \end{tabularx} \\ \\ \relax
4:12 & \begin{tabularx}{0.7\textwidth}{X} 從來沒有人見過神，我們若彼此相愛，神就住在我們裡面，他的愛在我們裡面得以完滿了。 \end{tabularx} \\ \\
[1ex]
\hline
\hline
\end{longtable}
$^{1}$洪永嘉的弟妹平安.
中國人有句說話是這麼說的.
他說一年之計在於春.
一日之計在於晨.
一生之計在於勤.
這番說話造就了中國人一個勤奮的性格.
以及在年頭為一年裡做好計劃.
做好準備.
如果說一年的計劃是這麼重要的話.
我在想用什麼心態去面對.
以及執行這個計劃就更加重要.
在2021年的第一個崇拜.
我想和大家分享一個生命的奇蹟.
在幾年之前已經有一個這樣的報導.
有一對夫婦在油麻地開設一個燒臘店.
店鋪開了二十多年了.
他們兩夫婦當時見到小市民生活越來越艱辛.
所以為了街坊能夠有一餐飽飯吃.
他們明知每賣出一個飯盒就虧了一個半.
仍然堅持賣11元一盒燒肉飯連泥湯.
堅持了三年都不加價.
令到街坊真的能夠有平價飯吃.
不要小看他們.
燒臘店每天賣出六百盒燒肉飯.
四百盒雙拼飯.
一天就有一千個人受惠.
老闆說每天要補貼一千五千元.
是的,大家沒有聽錯.
每天老闆要補貼一千五千元.
每月補貼四萬五千元.
是不是覺得很奇怪?.
在香港地,虧本生意也有人做嗎?.
原來這對夫婦計算過.
他們將批發燒肉所賺到的.
用來補貼這邊虧的.
所以剛剛打和,不是虧很多.
他的仗義行為影響了他的徒弟.
他在深水Po 賣燒肉飯.
每盒賣十元.
效法他施父幫助貧窮人.

$^{41}$我們對這種行為.
我們都會說老闆有企業良心.
他沒有賺到盡.
做生意也顧著街坊.
真是人間有情.
沒錯,的確是人間有情.
不過如果我們說人間有情.
我們能不能說人間有愛呢?.
燒肉舖老闆對街坊的心.
是一份情還是一份愛呢?.
我們相信無論是情也好.
或者是愛也好.
都要老闆對他者.
即是其他人有所行動.
而他者亦要對燒肉舖老闆的心意.
作出回應.
整個過程雙方是需要有互動的.
究竟聖經裡面有沒有類似的典範.
讓我們學習呢?.
今天我想透過約翰一書四章七至十二節.
去認識愛的三個不同的層面.
我們一起去祈禱.
讓我們在新的一年開始.
首先去聆聽你的話語.
求你讓你的道去鞏固我們.
建立我們的生命.
最後被你使用.
我們去祈禱.
奉靠主耶穌聖名祈求.
阿們.
各位主媒.
剛才主席都要求我們一起讀這六節經文.
有沒有留意到每一節的經文裡面.
都有「愛」這個字.
所不同的是這個「愛」字.
是作為句子的主詞.
或者授詞.
或者動詞.
原來這六節經文裡面的「愛」.
原文都是同一個字根.

$^{81}$所以我們就知道.
這六節經文的信息是相關的.
而且一定和愛有關.
我們說愛的第一個層面就是.
愛是神本質的一部分.
約翰一開始就說了.
他在信仰上的一個認知.
這個認知在第七節和第八節.
我們就看到了.
他說凡有愛的都是由神而生.
並且認識神.
沒有愛的就不認識神.
究竟約翰這個認知是根據什麼來建立的呢?.
我們在中文譯本就看不出.
那個第一個認識的字的時態.
是現在的時態.
是表示有愛心的人.
現在已經在認識神.
第二句沒有愛的人就不認識神.
經文提及的認識.
是一個過去的時態.
就表示了不單單是現在不認識.
更加是從過去到現在.
都從來未認識神.
為什麼約翰認為這些沒有愛心的人.
是從以往到現在都不認識神呢?.
約翰提供了一個答案的部分.
他說愛是從神而來的.
由神而生.
愛是從神而來.
表示了愛是有源頭的.
這個源頭就是神了.
如果愛心要透過神而來.
那沒有愛心的人.
就沒有可能認識到愛的根源了.
唯有那些有愛心的人.
才能夠認識到愛的根源就是神.
那什麼叫凡有愛的都是由神而生呢?.
當然「生」這個字.
不是李啟維那種由母親生兒子的「生」.

$^{121}$這個「生」字既可以解作「生育」.
也可以解作「導致」或「引導」.
即是愛的產生是由神所導致引導的.
沒有了神的引導.
根本產生不了愛.
既然愛的動力是由神所導致的.
沒有愛心的人.
那又如何認識神呢?.
唯有被神導引而產生愛的人.
才是真正認識神.
約翰在第八節尾.
也說出了他對這個認知的終極答案.
他說因為神就是愛.
神就是愛.
包括了愛是從神而來.
也包括了愛是由神而生.
神就是愛.
說明了愛是神本質的一部份.
也是神的一個重要的本質.
神如果不是愛.
就不是約翰所認識的神.
正正因為神就是愛.
所以約翰能夠在他的信仰生活上體驗到.
有愛心的都認識神.
沒有愛心的就不認識神.
如果神就是愛.
那麼神去愛一個他者的時候.
他會不會像這個燒肉店的老闆.
也計算過代價呢?.
燒肉店的老闆.
雖然每個月虧蝕四萬五千元.
但是他計算過這個代價.
然後才去做.
如果神有沒有計算過代價.
去為我們付出呢?.
以前聽過一個故事.
說到一隻蠍子很想過河.
他就請求青蛙背他過河.
青蛙說不行.
如果在過河期間.

$^{161}$在我背上針我一下.
那我就白白枉死.
蠍子就說.
如果我叛徒刺死你.
那我就跟你一樣掉進河裡.
我一樣會死的.
我不會這麼蠢的.
青蛙於是就相信了他.
就背他過河.
誰知道河中間青蛙背部一痛.
就知道一定是被蠍子刺到.
青蛙很不服氣.
在死的一剎那他就問蠍子.
你為什麼要刺我?.
難道你也不怕死嗎?.
蠍子就回答.
他說我也不想死.
不過這是我的本性.
本性使然.
我也控制不了.
故事是有延續的.
多數人都會理解.
去解讀這個故事.
蠍子的天性就是這樣.
就算沒了命也改不了.
真的死性不改.
其實人也差不多.
死性不改.
不過當然有些人用另一個角度去說這個故事.
就是用青蛙的角度去解讀這個故事.
他說青蛙是沒有義務去背負蠍子過河.
青蛙只要拒絕就可以了.
他不需要冒任何風險.
但是青蛙為什麼願意這麼做呢?.
是因為青蛙都認為蠍子.
都會愛惜自己的生命.
經過青蛙自己的評估.
他覺得蠍子的風險很低.
所以願意幫忙.
如果從這個角度去看.

$^{201}$每個人都會先計算代價.
然後才幫助別人.
神當然不是那隻蠍子.
因為蠍子控制不了他的本性.
但是神會不會像青蛙一樣.
當他幫助別人的時候.
都要計算一下要付出什麼代價呢?.
一看就告訴我們.
神就是愛.
這個帶有愛的本質的神.
當他去愛一個他者.
就是其他人的時候.
他的價值的取向和人類的標準.
是會有所不同.
所以愛的第二個層面就是最寶貴和捨己的愛.
神對他者做了什麼呢?.
在第九節和第十節.
他講了兩件事.
第一是神猜他獨一的兒子來到這個世界上.
第二就是猜他的兒子為我們的罪作了贖罪制.
第九節所提到的獨一的兒子.
並不是單純解作唯一的兒子.
兒子是一個形容為獨一無二極之珍貴的意思.
是神最珍貴最獨一無二的兒子.
就是耶穌.
對於三一的神.
耶穌就等同於神祂自己.
也可以說是神最寶貴.
神對他者的愛的表達.
是將最寶貴的耶穌猜到這個世界上來.
第十節所提到的是.
耶穌為罪作了贖罪制.
贖罪制的意思就是指.
耶穌為我們犧牲贖了我們的罪.
使我們免了罪的責任.
我們知道這個贖罪制其實就是.
耶穌為我們走上十字架受死.
是一條捨己犧牲的路.
神對他者的愛.
以兩種具體的行動表達了出來.

$^{241}$第一是將最好的最珍貴的.
施予了給他者.
第二就是為他者犧牲.
這個犧牲可以連性命都不要.
完全為他者的緣故.
而願意無條件的犧牲.
神為人做的這兩件事.
其實都是為了一個目的.
這個目的不是為了神自己.
而是為了對方他者.
就是第九節.
使我們直著他得生命.
得生命如果單純理解為得到生命.
我覺得不能夠溢出這個字的精髓.
這個字可以解釋為.
可以活出某種行為的方式.
充滿了活力賦予生命力的意思.
神將最好的給我們.
為我們犧牲生命.
為的是要我們能夠過一個充滿活力.
有動力精彩的生活方式.
當我們生命能夠以充滿活力.
有動力精彩的生活去展現的時候.
第九節就說.
神對我們的愛就顯明了.
我們生命的改變.
能夠顯明神對我們的愛.
當神對我們付出最好的時候.
為我們犧牲的時候.
第十節就說.
這個就是愛了.
神對我們的行動能夠顯明.
這個就是愛.
如果大家不忘記.
當年汶川大地震有一個故事.
今天還有人在流傳.
地震過後.
救援隊進行黃金72小時的探殺.
即搜索工作.
務求盡快找到生還者.

$^{281}$當時有個報道這樣說.
記載了一個獲救的故事.
搜索隊在一個瓦礫下面.
聽到有嬰孩的呼聲.
於是隊伍就立即向下發掘.
發掘到的時候.
已經看到一個女性的屍體.
似乎女死者已經死了超過一天.
這個女死者被發現的時候.
身體的形狀是屈膝跪下.
而且手抱著一個嬰孩.
原來這個姿勢就是要保護著嬰孩.
免了被下塌下來的物件所傷害.
而事實亦都證明如此.
嬰孩獲救.
接受檢查.
除了營養不良之外.
身體上沒有骨折.
沒有外傷.
事後證明.
嬰孩是個女死者.
是母子的關係.
這個嬰孩獲救的故事.
展現了一種.
媽媽或母親為了保護孩子的犧牲行為.
母親將最好的給了孩子.
為了保護孩子.
用身體去擋住下塌的石頭.
為了孩子能夠獲救的機會大一些.
犧牲自己的生命都在所不惜.
母親對孩子的愛.
目的是想孩子能夠繼續生存下去.
弟兄姊妹.
地上的母親尚且可以如此對待她的兒子.
更何況我們天上的父呢.
天父本性就是充滿愛.
當漢子說神在愛的時候.
神沒有計算過.
祂不會不洩才做.
祂亦不像青蛙那樣.

$^{321}$計算過風險評估.
然後才去做.
天父是將最好的給了我們.
哪怕是為我們犧牲.
目的是想我們的生命能夠更加豐盛.
今日我們的生命改變了.
是因為與神的愛有關.
究竟今日我們的生命.
與神的愛可以有甚麼改變呢?.
弟兄姊妹.
今日你的生命有甚麼改變呢?.
你可能說.
我已經藉著耶穌與神和好.
你可能說.
我已經確認是神的兒女.
還不夠嗎?.
我們現在可以戰勝死亡.
與永恆有份了.
我們可以藉著神的能力.
戰勝魔鬼了.
沒錯,是呀.
弟兄姊妹.
我們的身份是改變了.
但我們的生命質素有沒有改變呢?.
我們的生命的改變不能夠離開神的愛.
當我們能夠將最好的給人的時候.
愛人的心才顯明.
當我們能夠向人展現一個犧牲的愛的時候.
才能夠與神的愛相稱.
在此我懇求天父去饒恕我們.
饒恕我們每一個.
因為我們的愛與你相比.
實在相差很遠.
約翰開始時跟我們說了神的愛的真理.
然後解釋了神愛人的兩種表達方式和目的.
現在約翰要跟我們說明.
我們要怎樣去回應神的愛.
愛的第三個層面就是彼此相愛的超越性.
約翰在第十一節說.
既然神這樣愛我們.

$^{361}$我們也要彼此相愛.
大家覺得很奇怪嗎?.
為何約翰不是說既然神這樣愛我們.
我們也應該這樣去愛神呢?.
我們的標準不是你對我好.
我對你好嗎?.
難道約翰寫錯了嗎?.
但約翰在六節經文中.
他分別說了三次彼此相愛.
在第七節說我們要彼此相愛.
在第十一節說我們也要彼此相愛.
在第十二節說我們若彼此相愛.
我們要彼此相愛的「要」字是呼求的語氣.
約翰呼求信徒.
你們要這樣做.
第十一節「也要」是命令詞.
是一種命令的語氣.
因為神既然這樣愛我們.
我們也必須彼此相愛.
這個回應是一種命令.
亦是理所當然的責任.
但我們不明白.
我們仍然要問.
為何回應對象不是神呢?.
神愛我們.
我們應該愛神.
但約翰卻教我們.
你們要彼此相愛.
愛的對象是你旁邊那個.
為何?.
弟兄姊妹我想大家想想.
為何愛的回應對象不是神.
而是身邊的人.
神愛我們.
我們應該愛神.
約翰說你們要彼此相愛.
神愛我們.
但你們要彼此相愛.
答案在第十二節.
「我們若彼此相愛」是條件句子.

$^{401}$反映條件發生時.
便有接下來的結果.
結果解釋為何要彼此相愛.
約翰說「我們若彼此相愛」.
神就住在我們裡.
這句說話說明.
如果人與人之間能彼此相愛.
這種彼此相愛的關係.
有一種超越性.
因為神不單單參與其中.
神更加住在彼此相愛的人裡.
是人與神最親密的狀態.
是人與神關係的最高峰.
即人與人之間的彼此相愛.
神人的關係.
可以連結在一起.
弟兄姊妹你覺不覺得很奇怪.
人與人之間的相愛關係.
本來只是個人與個人的關係.
但因著神的參與.
可以成為彼此與神的一種聯繫.
這種愛的聯繫就是一種永恆的價值.
你想想.
從一個熟悉的情操.
轉變為一個熟悉永恆的意義.
如果我說這是一個奇蹟的話.
不知道你會不會同意.
本來只是個人與個人的關係.
但因著神的參與.
我們有了一個永恆的意義.
這個彼此相愛不是單純個人與個人.
而是三個人.
神與彼此之間.
不單單是這樣.
日漢再將這個奇蹟推上高峰.
他接著說.
他的愛在我們裡得以圓滿了.
人與人之間的彼此相愛.
竟然關係到神的愛可以達至完全圓滿.
100分.

$^{441}$神的愛在人世間有很多彰顯的處境.
當然我們知道有很多的規矩.
但在彼此相愛的行動裡.
神的愛竟然可以有一個圓滿的效果.
100分.
人類的一種彼此相愛的行為.
竟然關係到神的愛的本質.
可以達至圓滿.
各位姐妹.
如果我們在實踐彼此相愛的時候.
神就住在我們裡面.
如果這已經是奇蹟的話.
當我們實踐彼此相愛的時候.
神不單住在我們裡面.
並且神的愛得以完全充滿完美.
如果我說這是奇蹟之中的奇蹟.
你又會否同意?.
如果彼此相愛能夠與神的關係結連.
彼此相愛能夠將神的愛推到最圓滿的境界.
我相信這個回應的意義.
已經超越了第七節的呼求.
也超越了第十一節的責任和命令.
這個超越了呼求.
又超越了責任和命令的回應.
是人與神的一種親密的聯繫.
是最深入的關係.
是一種永恆的關係.
當我們用彼此相愛作為回應神的同時.
亦是神的愛透過我們彼此相愛的行動.
而達至最好,最大,最美的彰顯.
原來約翰想告訴我們.
回應神對我們愛最好的方法.
就是彼此相愛.
當我們說神很愛我們.
我們知道神對我們的愛是很好的.
我們要回應.
如何去回應?.
約翰告訴我們.
最好的回應方法.
彼此相愛.

$^{481}$想和大家分享一個見證.
有一個未信主的姐妹.
她是家庭主婦.
照顧著一個小孩.
有一天她發現丈夫有婚外情.
她當然要和丈夫攤牌.
她的丈夫跟姐妹說.
我們已經沒有感情了.
男方還主動提出離婚.
姐妹當然很傷心.
很想挽回.
但覺得丈夫心意已決.
也無法挽回.
好,來就來吧.
爭取到小孩的撫養權.
另外要求丈夫付適當的贍養費.
香港很正常.
誰知丈夫將公司便宜賣給第三者.
即是沒有贍養費了.
或者贍養費變得很少.
姐妹無奈接受所有離婚條件.
她想一切重新開始.
只要兒子還在身邊.
怎樣辛苦都值得.
接下來她要面對日後的經濟問題.
她要找工作.
但工作不夠請工人.
照顧不了小孩.
小孩還在讀書.
上班半天.
下午午餐如何?.
媽媽上班後.
下午午餐如何解決?.
這是日常的問題.
難道每天都要叫外賣?.
麥當勞?.
姐妹這個處境被另一間教會的信徒知道了.
即是輾轉.
她願意為小孩煮午飯.
亦願意晚上為姐妹買菜.

$^{521}$讓姐妹晚上回來.
不用去街市.
就有菜煮飯.
姐妹覺得這個教徒很有愛心.
很想多謝他.
見面的時候除了多謝他.
還問他為何願意幫忙.
他不認識小孩.
那個信徒說.
我回教會這麼久.
學會了一件事.
就是耶穌很愛我.
我現在這樣做.
也是學耶穌.
當然最後兩個人都做了好朋友.
最後姐妹都加入了.
受浸教會.
不過是不同的教會.
但兩個不同的.
這個故事告訴我們.
雖然不屬於同一間教會.
但卻可以實踐彼此相愛.
各位姐妹.
今天講題是生命欠一個奇蹟.
如果彼此相愛.
能夠造就一個人神關係的奇蹟.
如果彼此相愛.
能夠造就神的愛.
充滿在人間的奇蹟.
我想問我們有什麼理由.
不去實踐彼此相愛呢?.
約翰在六節經文告訴我們.
神的本質就是愛.
神有愛的本質.
神的每一個行動都帶有愛.
神差派最寶貴的耶穌.
為我們的罪走上十字架受死.
目的是要我們的生命.
可以活得更豐盛.
當我們已經是一個新造的人.

$^{561}$當我們已經是神的兒女.
當我們為神對我們的愛.
作出彼此相愛的回應的時候.
我們與神的關係會變得更親密.
而且達至神的愛.
在人間更加圓滿.
各位姐妹.
主席說.
我們今天是2021年.
第一個的主日崇拜.
2021年的年頭.
當我們要策劃今年的計劃事工的時候.
我們能否將彼此相愛.
加入在我們的計劃裡面呢?.
其實在11月底.
教會本來也想和大家進行一個分享會.
去講2021年我們教會有什麼事工.
未到2021年之前.
其實教會與很多NGO合作過.
做了很多扶貧的活動.
我們與陸在元朗派二手玩具.
我們與洪孩儀合作派冬至的送包.
我們基督教聯會有很多禮物給我們.
我們送給街坊.
我們與教官合作.
我們有很多福袋.
送給我們洪水橋的街坊.
2021年我們很想.
仍然繼續與很多NGO合作.
舞會的合作.
然後我們在洪水橋做扶貧的工作.
透過扶貧接觸我們的街坊.
透過扶貧能否做到主耶穌的彼此相愛.
縱然我們不認識祂.
但他們都是神所愛的人.
我們洪恩家就在洪水橋.
2021年能否實踐彼此相愛的聖經原則.
當我們實踐彼此相愛的時候.
神的愛就達至最好,最美,最大的效果.
這是教會的事工發展.

$^{601}$個人又如何呢?.
個人如何能夠實踐彼此相愛呢?.
我很想挑戰大家.
彼此相愛一定是由禱告和關心開始.
沒有禱告的托著.
你的愛是激發不起來的.
愛是從神而來的.
你與神的關係是怎樣.
你的愛就是怎樣.
你與神的關係越疏離.
你就越沒有愛心.
能否參與我們教會的祈禱會呢?.
與我們一起同心合意的禱告.
為了我們的會友.
為了我們的香港.
為了我們的教會獻上我們的禱告祭.
神的愛就在你的生命裡能夠彰顯.
關心一下我們身邊的人.
給他一個電話.
WhatsApp他.
如果你說我身邊的人已經關心了.
回來教會.
與我們一起參與扶貧活動.
接觸一些你不認識的街坊.
去派米.
去派口罩.
去玩玩具.
這些都是我們在2021年會做的.
期望弟兄姊妹的生命能夠展現神的愛.
在我們每一個彼此弟兄姊妹的心裡.
我們一起祈禱.
你愛我們去到一個地步.
就是為我們死.
我們要怎樣回應呢?.
你要求我們彼此相愛.
因為你就在我們中間.
並且得以完全.
求主讓我們在今年能夠切實去實踐.
我們對別人的關愛.
以回應你對我們的恩典.

$^{641}$我們的祈禱是奉主耶穌基督得勝明智祈求.
阿們.
願至祝福大家.
\newpage



\section{創世記 1:1-31}
\label{sec:kzUx2k7DXdQ}
\textbf{2021.01.10《起初,神創造天地》創世記 1\_1-31( 和修版)  曾敏姑娘}
\newline
\newline
連結: \href{https://youtube.com/watch?v=kzUx2k7DXdQ}{\texttt{https://youtube.com/watch?v=kzUx2k7DXdQ}} ~~~~ 語音日期: 2021-01-11
\newline
\newline
\hyperref[sec:r3qTUjnU0eY]{\small{< < < PREV SERMON < < <}}
~
\hyperref[sec:index_chronic]{\small{[返順時目]}}
~
\hyperref[sec:index_scriptual]{\small{[返順卷目]}}
~
\hyperref[sec:hPE_EBTYs3o]{\small{> > > NEXT SERMON > > >}}
\newline
\newline
創世記 1:1-31
\newline
\begin{longtable}{cl}
\hline
\hline
章節 & 經文 (和合本修訂版)\\
\hline
1:1 & \begin{tabularx}{0.7\textwidth}{X} 起初，神創造天地。 \end{tabularx} \\ \\ \relax
1:2 & \begin{tabularx}{0.7\textwidth}{X} 地是空虛混沌，深淵上面一片黑暗；神的靈運行在水面上。 \end{tabularx} \\ \\ \relax
1:3 & \begin{tabularx}{0.7\textwidth}{X} 神說：「要有光」，就有了光。 \end{tabularx} \\ \\ \relax
1:4 & \begin{tabularx}{0.7\textwidth}{X} 神看光是好的，於是神就把光和暗分開。 \end{tabularx} \\ \\ \relax
1:5 & \begin{tabularx}{0.7\textwidth}{X} 神稱光為「晝」，稱暗為「夜」。有晚上，有早晨，這是第一日。 \end{tabularx} \\ \\ \relax
1:6 & \begin{tabularx}{0.7\textwidth}{X} 神說：「眾水之間要有穹蒼，把水和水分開。」 \end{tabularx} \\ \\ \relax
1:7 & \begin{tabularx}{0.7\textwidth}{X} 神就造了穹蒼，把穹蒼以下的水和穹蒼以上的水分開。事就這樣成了。 \end{tabularx} \\ \\ \relax
1:8 & \begin{tabularx}{0.7\textwidth}{X} 神稱穹蒼為「天」。有晚上，有早晨，這是第二日。 \end{tabularx} \\ \\ \relax
1:9 & \begin{tabularx}{0.7\textwidth}{X} 神說：「天下面的水要聚集在一處，使乾地露出來。」事就這樣成了。 \end{tabularx} \\ \\ \relax
1:10 & \begin{tabularx}{0.7\textwidth}{X} 神稱乾地為「地」，稱聚集在一起的水為「海」。神看為好的。 \end{tabularx} \\ \\ \relax
1:11 & \begin{tabularx}{0.7\textwidth}{X} 神說：「地要長出植物，就是含種子的五穀菜蔬，和會結果子、果子裡有種子的樹，在地上各從其類。」事就這樣成了。 \end{tabularx} \\ \\ \relax
1:12 & \begin{tabularx}{0.7\textwidth}{X} 於是地長出了植物：含種子的五穀菜蔬，各從其類；會結果子、果子裡有種子的樹，各從其類。神看為好的。 \end{tabularx} \\ \\ \relax
1:13 & \begin{tabularx}{0.7\textwidth}{X} 有晚上，有早晨，這是第三日。 \end{tabularx} \\ \\ \relax
1:14 & \begin{tabularx}{0.7\textwidth}{X} 神說：「天上要有光體來分晝夜，讓它們作記號，定季節、日子、年份， \end{tabularx} \\ \\ \relax
1:15 & \begin{tabularx}{0.7\textwidth}{X} 它們要在天空發光，照在地上。」事就這樣成了。 \end{tabularx} \\ \\ \relax
1:16 & \begin{tabularx}{0.7\textwidth}{X} 於是神造了兩個大光體，大的管晝，小的管夜，又造了星辰。 \end{tabularx} \\ \\ \relax
1:17 & \begin{tabularx}{0.7\textwidth}{X} 神把這些光體擺列在天空，照在地上， \end{tabularx} \\ \\ \relax
1:18 & \begin{tabularx}{0.7\textwidth}{X} 管理晝夜，分別光暗。神看為好的。 \end{tabularx} \\ \\ \relax
1:19 & \begin{tabularx}{0.7\textwidth}{X} 有晚上，有早晨，這是第四日。 \end{tabularx} \\ \\ \relax
1:20 & \begin{tabularx}{0.7\textwidth}{X} 神說：「水要滋生眾多有生命之物；要有鳥飛在地面以上，天空之中。」 \end{tabularx} \\ \\ \relax
1:21 & \begin{tabularx}{0.7\textwidth}{X} 神就創造了大魚和在水裡滋生的各樣活動的生物，各從其類，以及各樣有翅膀的鳥，各從其類。神看為好的。 \end{tabularx} \\ \\ \relax
1:22 & \begin{tabularx}{0.7\textwidth}{X} 神就賜福給這一切，說：「要繁殖增多，充滿在海的水裡；飛鳥也要在地上增多。」 \end{tabularx} \\ \\ \relax
1:23 & \begin{tabularx}{0.7\textwidth}{X} 有晚上，有早晨，這是第五日。 \end{tabularx} \\ \\ \relax
1:24 & \begin{tabularx}{0.7\textwidth}{X} 神說：「地要生出有生命之物，各從其類，就是牲畜、爬行動物、地上的走獸，各從其類。」事就這樣成了。 \end{tabularx} \\ \\ \relax
1:25 & \begin{tabularx}{0.7\textwidth}{X} 於是神造了地上的走獸，各從其類；牲畜，各從其類；和地上一切的爬行動物，各從其類。神看為好的。 \end{tabularx} \\ \\ \relax
1:26 & \begin{tabularx}{0.7\textwidth}{X} 神說：「我們要照著我們的形像，按著我們的樣式造人，使他們管理海裡的魚、天空的鳥、地上的牲畜和全地，以及地上爬的一切爬行動物。」 \end{tabularx} \\ \\ \relax
1:27 & \begin{tabularx}{0.7\textwidth}{X} 神就照著他的形像創造人，照著神的形像創造他們；他創造了他們，有男有女。 \end{tabularx} \\ \\ \relax
1:28 & \begin{tabularx}{0.7\textwidth}{X} 神賜福給他們，神對他們說：「要生養眾多，遍滿這地，治理它；要管理海裡的魚、天空的鳥和地上各樣活動的生物。」 \end{tabularx} \\ \\ \relax
1:29 & \begin{tabularx}{0.7\textwidth}{X} 神說：「看哪，我把全地一切含種子的五穀菜蔬和一切會結果子、果子裡有種子的樹，都賜給你們；這些都可作食物。 \end{tabularx} \\ \\ \relax
1:30 & \begin{tabularx}{0.7\textwidth}{X} 至於地上一切的走獸、天空一切的飛鳥，並一切在地上爬行的，有生命的動物，我把綠色植物賜給牠們作食物。」事就這樣成了。 \end{tabularx} \\ \\ \relax
1:31 & \begin{tabularx}{0.7\textwidth}{X} 神看一切所造的，看哪，都非常好。有晚上，有早晨，這是第六日。 \end{tabularx} \\ \\
[1ex]
\hline
\hline
\end{longtable}
$^{1}$(自動翻譯).
繼續聆聽祂的話.
我想起在疫症期間.
我們的生活大受影響.
所以有幾種情況出現了.
有很多人選擇留在家裡.
因為都是回應政府的呼籲.
減少社交活動.
但也有很多人選擇外出.
其中一個活動就是行山.
你看到社交媒體上.
很多照片和行山有關.
我一直看的時候都很羨慕.
當時想我已經很久沒有去行山了.
覺得很懷念很多我看過的很美的風景.
我又想起了一個片段.
這個片段我從昨天晚上到今天.
都在腦海裡盤旋著揮之不去.
想起很多年前.
和我的家人一起去北京旅行.
真的很多年前.
我和我姐姐一起爬上香山.
北京的香山.
北京的香山你說很矮肯定不是.
你說很高也不算是.
不過要花點力氣爬上去.
爬到頂的時候有兩個選擇.
一是走下來.
一是可以坐纜車.
當時我姐姐很想試新鮮的東西.
她總是感到很興奮.
她說不如坐纜車.
可以看看風景.
我說好啊坐纜車.
那個纜車就特別.
它是一張椅子.
不是密封式的椅子.
開放式的一張椅子.
然後只有一個安全圍欄.
在前面扎著保護我.

$^{41}$當時我們就一起坐上去.
我的心情就開始好不好.
因為我有一個很大的弱點.
我是畏高的.
很厲害的畏高.
香山也高的.
我們就坐上去.
完全是周圍沒有保護的.
只有前面的圍欄.
一直坐坐坐.
去到前面準備要衝出閘口.
要面對高山.
當時的心情越來越緊張.
可能是因為纜車已經很舊了.
很舊的速度不夠快.
於是工作人員為了增加我們的刺激感.
我們出閘前.
纜車出閘前.
工作人員在後面用力一推.
把車推出去.
我們就衝出閘.
不是人衝出去.
是纜車衝出去.
看著美麗的高山.
我記得那一刻.
我即時的反應是.
我真的很想死.
真的很辛苦.
因為高山很高.
我閉上眼睛完全不看.
我姐姐在旁邊就有完全另一個反應.
因為她很大膽.
她就說「嘩!很美啊!」.
大聲地喊.
「很美啊!」.
「你張開雙眼!張開雙眼!」.
「快點看著!多美啊!」.
她越叫我就越覺得很想死.
兩者有很大分別.
我必須承認.

$^{81}$真的很美.
後來過了一會.
我沒那麼害怕的時候.
我張開雙眼看.
很美.
那個大自然的景色.
沒辦法用文字去形容.
就算當時有一個多美麗的手機.
我總覺得拍不完這麼多美景.
眼睛所看到的那一刻是最美的.
我很同意我姐姐口中的呼喊.
當時姐姐是未信主的.
她已經發出了這種讚歎.
你說大自然有多麼美呢?.
今時今日很多香港的市民.
都去行山.
特別喜歡挑戰高難度的山.
我深信大家走到山上的時候.
都會有一個讚歎.
很美.
香港有很多很美的山.
大自然就是這麼美.
今天很想跟大家說一件事.
起初神創造天地.
就很想說.
那個創天造地的神.
究竟是怎樣的.
我們一直看大自然景色的時候.
原來這一切背後.
都有一位創造的主.
究竟是怎樣的呢?.
今天看這張經文其實非常豐富.
我至少看到三點.
很想跟大家分享.
這個創天造地的神.
至少有三大特色.
第一,他是超越萬有的.
全能的神.
超越萬有的全能的神.
起初神創造天地.

$^{121}$是整本聖經的第一句說話.
很震撼,很有豪氣的說出來.
起初這裡是指人類歷史開始的時候.
天地這裡是說宇宙.
在人類歷史開始的時候.
神創造了整個宇宙.
這是很令人感到震撼的宣告.
按照維基百科的資料.
什麼是宇宙呢?我也查了一下.
按照我們人類很有限的了解.
宇宙是指所有時間,空間.
和它所包含的內容,物質等等.
所構成的這一切為之宇宙.
宇宙包含了行星.
星系間的空間,次元子粒.
這些很深.
還有所有的物質和能力.
我深信從來沒有人可以量度宇宙有多大.
或者裡面的內容物質有多豐富.
沒有人可以量度.
沒有人可以告訴我們.
你頂著頭腦才可以估計.
但是那些估計也不實際.
而我們人在哪裡呢?.
是在地球上生活.
地球不代表整個宇宙.
是地球.
地球也不知道佔宇宙的幾分之幾.
可能佔很少.
一個很小的位置.
歷史歷代的科學家花盡了他們的心機和力量去研究地球.
得出來的都是一點點的知識.
但是這些一點點的知識對於全世界的人來說已經是很多了.
是很多.
但是你記住是一點點的知識.
直到目前為止.
只是說我們這個地球.
還有很多奧秘等待我們人類去發掘的.
所以才會出這麼多書.
世界有多大多大的奧秘.

$^{161}$在哪裡哪裡有什麼解不開的謎.
連國家也找不到答案.
如果地球有這麼多奧秘.
更何況是整個宇宙呢?.
整個宇宙浩瀚無邊.
宇宙的奧秘多不勝數.
起初神創造天地.
顯明了神是宇宙的創造者.
但是在這裡最重要最重要告訴我們.
神比這一切更大.
你覺得宇宙大.
神比這一切更大.
神超越萬有.
不單是這樣.
上帝更加是全能的神.
神藉著祂的話語去創造天地.
神的話語滿有權能.
一發出就會成就.
我舉一些簡單的例子.
例如第三節那裡這樣說.
如果你在家裡.
在哪裡都好.
你手上有聖經.
你可以打開創世記第一章.
跟我一起很快看看.
今天也多看經文.
第三節那裡這樣說.
神說要有光就有了光.
神說 神說.
第六節.
神說中水之間要有窮倉.
把水和水分開.
神就造了窮倉.
把窮倉以下的水.
和窮倉以上的水分開.
事就這樣成了.
神說 事就這樣成了.
第十四節又是神說.
神說話.
天上要有光體來分晝夜.

$^{201}$讓他們作記號.
定季節 日子 年份.
他們要在天空發光.
照在地上.
事就這樣成了.
每一節都是神說.
事就這樣成了.
或者神說.
然後就有了.
神要怎樣創造.
就怎樣創造.
沒有人給過他意見.
根本沒有人可以給他意見.
也沒有人可以攔阻他.
神是萬物的源頭.
萬物是因為神而來.
萬物都顯出.
神創造的奇妙.
和祂的大能大力.
神更加是我們生命的源頭.
生命是因為神而來.
每一顆生命.
無論是動物的生命.
還是人的生命.
都顯示出.
神創造的智慧和能力.
神是全能的.
祂超越萬有.
這是第一個觀察.
第二個觀察.
神是掌管宇宙的.
祂管理宇宙的.
因為這樣的神.
神很喜愛秩序.
很喜愛和諧.
在這裡很有趣.
很多時候我們只分.
創造我們和小朋友.
或者少年人說.
聖經故事都可以說.

$^{241}$第一天做什麼.
第二天做什麼.
第三天做什麼.
我們只顧著.
每天有什麼不同的事情.
物件.
我們有沒有留意一個原則.
很重要.
神很喜愛秩序.
神創造的.
不會混亂.
神創造天地的時候.
祂會將各樣物件.
擺放在一個合適的位置上.
而且各自有自己的功用.
例如說.
神創造光.
光就是生命成長.
一個不可以缺少的元素.
第四節.
神看光是好的.
於是神就把光和暗分開.
分開的意思是什麼.
就是將這些創造物.
擺放在合適的位置上.
令這些創造物.
顯示出有一個秩序.
神將光暗分開.
這樣就可以定晝夜.
可以分日子.
第六和第七節.
出現了分開的詞語.
神說.
種水之間要有窮倉.
把水和水分開.
神就造了窮倉.
把窮倉以下的水和窮倉以上的水分開.
這裡的分開.
是要為水訂立界限.
窮倉以上的水和以下的水.

$^{281}$會以不同的形態出現.
簡單來說.
窮倉以上的水是雲層.
和以下的水不同.
以下的水有不同的形態.
水是生命之源.
對動植物和人類都很重要.
這個分開.
顯明了神在裡面掌管.
神有自己的心意.
祂有祂的智慧.
祂就是這樣在管理.
在掌管.
我要達成神自己的心意.
除了分開這個詞語之外.
還有一個詞語很值得我們注意的.
神怎樣掌管呢?.
就是在這裡特別說.
「各從其類」.
第十一節說.
神說.
地要長出植物.
就是含種子的五穀菜蔬.
和會結果子.
果子裡有種子的樹.
在地上各從其類.
十二節.
於是地長出了植物.
含種子的五穀菜蔬.
各從其類.
會結果子.
果子裡有種子的樹.
各從其類.
我上述提到的蔬果.
都是各從其類.
各有各的特色.
它們有自己的功用.
其中一個功用是什麼呢?.
就是作為食物.
原來神就是這樣.

$^{321}$在這裡預備食物.
這些很重要.
這些菜蔬.
這些果子.
是為大自然的環境帶來很多好處.
第二十一節.
神創造了大魚和水裡之生的.
各樣活動的生物.
各從其類.
以及各樣有翅膀的鳥.
各從其類.
天空和水裡的動物.
各有各的特色.
也各自有它們在大自然裡的位置.
和角色.
透過分類.
神就可以對大自然萬物.
進行不同的管理.
達至祂的目的.
這個真的很重要.
你見到神在掌管.
神在掌管.
提到各從其類.
分類.
最近我終於下定決心.
要做一件事.
口說了很久要做的.
要收拾房子.
家裡非常混亂.
我決定要收拾房子.
說了很久了.
最近收拾房子.
收拾的時候覺得十分不容易.
因為物品很多.
很久沒有清潔了.
於是有很多塵.
但我一邊收拾.
腦海就出現了幾個字.
分類.
將它們分類.

$^{361}$分類就容易收拾好.
將它們分為同一類.
之後就容易處理.
房子就會整潔.
慢慢是這樣.
我真的這樣將它們分類.
分類就很好收拾.
其實神在創造的時候.
祂也將它們分類.
每個類別各有自己的特色.
不會誇個類別.
在裡面你見到神的管理能力.
管治的能力.
這個很重要.
聖經不單說神喜歡秩序.
在管理.
他強調他的能力.
他的掌管性是多麼強.
神會為不同物件命名.
這代表了神對宇宙萬物.
擁有絕對的管理權.
你為一件事命名.
通常你是它的主人.
你有權管理它.
你才可以為它命名.
你看到神對宇宙萬物.
是有管理權的.
例如第二節.
這裡這樣說.
神稱光為晝.
祂稱光為晝.
祂稱暗為夜.
祂稱暗為夜晚.
神給名不同的事物.
顯示出神在掌管.
神絕對掌管光暗.
掌管白晝黑夜.
掌管宇宙的一切.
後來神把命名的權給了人.
所以人可以為不同的動物起名.

$^{401}$這個叫什麼.
但在此之前.
神擁有絕對的權.
祂只不過把權下放給人.
你見到祂是宇宙真正的掌管者.
祂掌管整個宇宙.
不單在創造世界的那一刻.
上帝掌管這個世界.
今天祂掌管這個世界.
今天神沒有失去這個權.
今天神掌管全世界整個宇宙.
祂是最終極的主人.
將來直到永恆.
祂也是.
祂永遠在此坐著為王.
這一點值得我們記住.
第三個特色.
第三個特色就是.
這位創天造地的上帝.
是很愛人.
很愛大自然萬物.
祂的創造不是一個冷冰冰的創造.
也不是這樣.
只是欣賞自己的作品的創造.
是一個愛的創造.
裡面充滿了愛.
整個創世記的第一章.
實際上如果換個主題.
換個方法去講.
你可以說上帝怎樣為人.
預備居住的地方.
上帝怎樣為大自然萬物.
預備他們的公營.
可以從這個角度去看.
創世記一章二節提到.
地是空虛雲遁.
冤冕上面一片黑暗.
這個地和第一節的天地不一樣.
第一節天地是講宇宙.
但這裡講的地.

$^{441}$我們縮小範圍.
窄一點.
是地球.
地是冤冕黑暗.
這個地是地球.
人類在地球生活之前.
其實地球是怎樣的?.
空虛雲遁.
空虛雲遁是講沒有人住.
荒涼.
很亂.
沒有秩序.
其實是很可怕的地方.
地球上全部都是水.
冤是講甚麼?.
水.
冤冕很黑暗.
不只是一般的水.
是大水.
整個地球被大水包圍住.
大水上面是一片黑暗.
沒有人可以在這裡住.
這裡很荒蕪.
神逐步改變這個地方.
向著一個目標.
就是要讓這個地方成為一個很美好的地方.
成為一個人可以居住的地方.
首先針對地球很黑暗的情況.
神創造了光.
光的出現.
為生命的成長帶來了盼望.
提供了第一個很重要的要素.
神為光暗訂立界限.
讓時間的運轉變得很有規律.
接著針對地球上被大水所包圍的情況.
神為水分界限.
分開天空上的水和天空下的水.
剛才不斷地說水是生命之源.
當神為水訂立界限的時候.
神就可以讓水在地球上有規律地運作.

$^{481}$這樣更加適合生命的成長.
而針對地球荒涼的情況.
神一步一步地帶來生命.
剛才我們說神為水訂立界限.
在這個時候.
當神為水訂立界限.
就會讓乾地露出來.
有了乾地就有下一步.
可以長出植物.
這些植物包括五穀菜蔬.
結果子的樹等等.
這些植物不單止為地球帶來了生機.
也像剛才我說的.
成為了動物和人類的食物.
神很體貼萬物的需要.
體貼人的需要.
祂知道人所需要的.
不單止光,水,空氣,地.
人和自然萬物的生物一樣都需要食物.
所以神會去預備.
當然不單止食物.
其實我們人在生活裡.
很需要光體的幫助.
甚麼是光體?.
就是太陽,月亮,星獸.
這一切對我們的生活.
對自然萬物都有很大影響.
所以神就將太陽,月亮,星獸.
擺在天空.
照在地上.
管理就業.
分別光暗.
做不同的工作.
到了這一刻.
你會發覺地球開始很適合生物生存.
有光,有空氣,有水,有乾地.
很適合生物生存.
這時候神就開始做海裡的大魚.
水裡各樣活動的生物和天空的鳥.
神很愛他所做的生物.

$^{521}$所以經文說.
祂賜福給這些生物.
令他們繁殖增多.
神不單是愛人.
神也很愛這些生物.
祂看著這些生物.
祂喜愛他們.
祂憐惜他們.
祂做這些生物的時候.
同樣地球環境已經很適合他們生長的時候.
這樣去做他們.
之後神又創造了地上的走獸,牲畜.
和一切的爬行動物.
地面上不再那麼荒涼.
開始有生命氣息.
開始很熱鬧.
動物和靜物之間的關係.
又非常微妙.
互惠互利.
生生不息.
神欣賞自己的創作.
感到滿足快樂.
誰知道.
「神看為好的」這句話.
就這幾個字眼.
在這章經文出現了四次.
在頭章那裡.
「神看光是好的」又出現了一次.
即是神一直看自己做的事.
很好,很滿足,很欣賞.
就是神看到這些事.
每樣事物都那麼好.
他覺得很滿意.
這個時候.
他才開始做人.
他覺得每樣事物都很欣賞,很好,很美.
他覺得這個時候人出現.
他生活是非常適合.
他這樣才是做人.
所以你看到.

$^{561}$神對人那份的愛.
神對人的愛.
而且神創造人的時候.
是按照他自己的形象,樣式去做人.
這個我要留待下一次我的講道.
才慢慢講關於人.
被做那方面的事情.
這裡簡單來說就是.
人被創造出來之後.
人是很像神的.
就好像說子女.
好像爸爸媽媽一樣.
很有爸爸媽媽的形象,樣式.
這個也是一個愛的表現.
上帝做人和他做萬物是不一樣的.
其他的創造並沒有看到.
上帝按照自己的形象,樣式.
做植物,動物,做光,做這些.
沒有.
唯一講人是按照他的形象,樣式.
你看到他很愛我們.
我喜歡用一個比喻.
上帝愛我們好像爸爸媽媽愛子女一樣.
我們想想.
我經常在腦海裡出現了這個圖畫.
或者這些片段.
想像一下有一對夫婦.
結婚之後.
太太懷孕了.
準備生孩子了.
在這個時候.
你想想父母做多少預備.
如果條件許可.
首先預備一個房間.
這個房間將來等孩子長大了.
孩子可以擁有自己的房間.
房間裡有孩子的床.
最重要是房間的牆.
有柔和,漂亮的顏色.
將來讓孩子的眼睛睜大.

$^{601}$看到感到很快樂,很溫馨,很舒服.
就算不在也要換一個牆紙很適合他.
床也要很安全.
床上有玩具.
床又舒舒服服.
然後為孩子預備很多衣服.
又要想想將來孩子出生.
是沒有魚還是奶粉呢?.
奶粉是哪一種呢?.
很多很多很多東西要想.
想好了.
接著迎接孩子的出生.
其實神去做這個世界的時候.
他也是這樣想的.
做好所有的事情.
他很細心地去計劃.
很細心地去部署.
做好了.
覺得滿意了.
這樣他就做人了.
他想人來到世界.
是享受他所預備的.
正如爸爸媽媽希望.
孩子出生得到最好的一樣.
在這裡你看到神對我們那份愛.
神對萬物也有那份愛.
神派人去管理萬物.
那個管理不是剝削.
不是貪婪的剝削,殺害.
是神派人去管理萬物.
叫萬物在秩序裡好好地生活.
享福.
這是神的心意.
我們回看.
原來神愛人,愛大自然,萬物.
是很久以前的事.
由一開始神創造天地.
神就愛我們.
很久了.
很多時候我聽弟兄姊妹都是這樣說.

$^{641}$新約的神很慈愛.
很愛我們.
因為他派耶穌基督為我們在十誡上犧牲.
所以我喜歡新約的神.
祂很愛我們.
但是我很想說.
這位上帝從來都沒有變過.
這位上帝由一開始就愛我們.
祂的愛深到我們是不能夠明白的.
這份愛很久了.
是一份永恆的愛.
由一開始祂創造天地就愛我們.
一直愛愛愛很多世紀.
愛到今天都在愛我們.
直到永恆祂都愛我們.
沒有變過.
一直愛我們.
這份愛我們有多了解呢?.
既然我們剛才對上帝有這三個發現.
這三個發現帶來我們有什麼回應呢?.
值得我們去想.
第一件事.
剛才說那位創天造地的上帝是超越萬有的.
是全能的神.
我們知道這一點有什麼回應呢?.
我們的回應是.
大大的讚美神.
如果讓我再來一次.
我再去香山.
走上山.
下山的時候再坐一次纜車.
今時今日我的姐姐已經信耶穌了.
我們兩個再來一次.
經歷香山的纜車之旅.
再坐纜車下山.
我想做一件很瘋狂的事.
之前是姐姐大叫.
「好靚啊」.
我打算和她一起坐下來的時候.
這次一起大叫.

$^{681}$深信靠著神的恩典.
我現在應該有進步的.
我的畏高症.
她又信了主.
她那麼欣賞上帝的創造.
我也欣賞.
很想和她一起大大聲大叫.
「好靚啊」大自然萬物.
然後再加一句.
「上帝你好好啊」.
「上帝你好奇妙啊」.
「好想讚美上帝你」.
很想在纜車上和她一起大叫這一句.
因為上帝佩得.
上帝佩得我們去頌讚.
物理學家之父牛頓.
他深深意識到.
創造浩瀚宇宙者是多麼偉大.
他曾經說.
在沒有物質的地方究竟有什麼呢?.
他真的問這些問題.
太陽和行星的引力從哪裡來呢?.
宇宙為什麼那麼井然有序呢?.
動物的眼睛是根據光學原理設計的嗎?.
豈不是宇宙間有一位造物主嗎?.
雖然科學未能使我們立刻明白萬物的起源.
但所有這一切都引導我們.
歸向萬有的神面前.
牛頓在讚歎.
有很多科學家.
一邊探討世界的奧秘.
一邊讚歎神的偉大.
這些讚歎絕對可以帶來更多讚美的聲音.
星期日我們崇拜的時候.
我覺得很寶貴.
不單是聽上帝的話.
也是我們一同讚美我們的神.
我誠意鼓勵在家中的弟兄姊妹們.
即使你在家中.
不是在現場沒有實景.

$^{721}$我鼓勵你們唱詩歌的時候.
放聲歌唱.
你找一個環境可以放聲歌唱.
上帝佩德我們敬拜讚美.
用盡力氣,力量.
用心靈誠實去讚美祂.
佩德.
但我們的讚美不應該只留在星期日.
是每天都要讚美上帝.
每天都歌頌祂.
將榮耀歸還給祂.
我們可以有聲無聲.
有心靈的歌頌.
有嘴唇的歌頌.
上帝真是佩德.
第二件事.
既然這位創天造地的上帝.
祂掌管宇宙.
是一個很奇妙的大能的掌管者.
昔在,今在,以後永在.
祂真的一路掌管.
這一點帶給我們什麼思考呢?.
就是我們應該依靠祂.
我們有困難的時候.
第一時間不是去找人幫忙.
Peter,Mary.
發生一些很不開心的事.
第一時間就去傾訴.
有沒有想過.
第一時間先找上帝?.
上帝佩德我們去依靠祂.
在這個世界你告訴我.
有誰比上帝更厲害呢?.
上帝今天還在掌管這個宇宙.
不找祂找誰呢?.
為什麼我們找身邊的人呢?.
那些人和我們一樣都有限制.
那些人都有自己的需要.
我不是說完全不找人.
而是第一時間.

$^{761}$你最依靠的.
你找誰呢?.
我知道疫症令我們心情很低落.
我們很害怕,很沮喪.
我們很不開心.
我知道人生很多經歷令我們害怕.
知道人際衝突令我們害怕.
我們感到沮喪,灰心.
我們經常覺得神去了哪裡.
祂不會幫我們.
祂不知道,祂不了解.
對不起,我告訴你.
神知道,神了解.
神今天還在.
祂是最大的那個.
我們一起來找祂.
專心依靠祂.
我深信在新的一年.
當我們校正眼光.
看著這位上帝.
看清楚誰是最大的.
困難不是最大的.
憂傷不是最大的.
失敗,失意都不是最大的.
包括我們身體的疾病.
我們感到很痛苦的.
都不是最大的.
經濟的困難不是最大的.
上帝是最大的.
讓我們仰望祂.
第三點,最後一點.
既然這位創天造地的上帝.
祂很愛人,很愛萬物.
我們的回應是什麼?.
我們更應該要愛祂.
今天苦心自問.
問問我們自己.
我們最愛誰?.
我們什麼時候才會找神?.
有需要時找祂.

$^{801}$我們很喜歡把神當成阿拉丁神燈.
揉揉祂的神燈.
「出來啊,幫幫我」.
「我現在有困難」.
之後就放下祂.
很久很久很久沒有聯絡祂.
過了一陣子.
「人生又有困難」.
「這個時候又來揉揉神燈」.
「你快點出來啊」.
「我又有這個需要」.
「很大困難,你幫幫我」.
「很緊急」.
幫完之後.
又離開了.
這不是愛.
什麼叫愛?.
你日思夜想都是想著那個人.
你是為那個人的好處著想.
你想那個人的感受.
我們這樣對上帝絕不是愛.
是利用.
我們希望神為我們服務.
比較像我們自己是神.
祂不是神.
當祂是神的時候.
你真的把祂當成神去看.
不一樣.
你會敬畏祂.
你如果真的愛祂.
就不會這樣.
愛祂就會常常親近祂.
不是因為今天發生了很悲痛.
很悲慘.
很困難的事.
我才親近祂.
天天親近祂.
就像你的情人.
你天天要跟祂說情話.
上帝留了一本情書給我們.

$^{841}$就是聖經.
祂天天都想跟我們說祂的情話.
但我們很多時候都不想聽.
新一年讓我們天天聽上帝的情話.
天天跟上帝說情話.
天天都想了解上帝的心意是什麼.
天天都想知道.
怎樣令上帝的心得到滿足.
去愛祂.
去佩特.
我們一同禱告.
今天讓我們重溫《創世記》這章經文.
讓我們看回你的奇妙.
你的偉大.
你佩德榮耀.
尊貴.
稱重.
看回你的愛.
是這麼寶貴.
我們真的有很多事.
需要靜下來.
去思考.
我真的很需要思考.
今天我們對你的反應.
對你的回應.
究竟是怎樣的.
天父啊.
在過去的這段時間.
如果我們上上冷落你.
撇棄你.
上上藐視你.
背叛你.
就求你寬恕我們的罪惡.
願以耶穌基督潔淨我們的生命.
泡沫我們的罪惡.
聖靈啊.
求你充滿我們的心.
感動我們回轉.
感動我們真的回到上帝你的身邊.
敬畏你.

$^{881}$愛你.
讓我們一生不改變.
在新一年開始的時候.
願聖靈給我們力量.
給我們力量去立志.
給我們力量去改變.
奉耶穌基督的名字.
祈禱.
阿們.
(今天視頻就拍到這裡了,下次見).
\newpage



\section{提摩太後書 2:1-7}
\label{sec:hPE_EBTYs3o}
\textbf{2021.01.17《傳下去》提摩太後書 2\_1-7( 和修版)  張美薇牧師}
\newline
\newline
連結: \href{https://youtube.com/watch?v=hPE_EBTYs3o}{\texttt{https://youtube.com/watch?v=hPE\_EBTYs3o}} ~~~~ 語音日期: 2021-01-18
\newline
\newline
\hyperref[sec:kzUx2k7DXdQ]{\small{< < < PREV SERMON < < <}}
~
\hyperref[sec:index_chronic]{\small{[返順時目]}}
~
\hyperref[sec:index_scriptual]{\small{[返順卷目]}}
~
\hyperref[sec:sVUIck2trk4]{\small{> > > NEXT SERMON > > >}}
\newline
\newline
提摩太後書 2:1-7
\newline
\begin{longtable}{cl}
\hline
\hline
章節 & 經文 (和合本修訂版)\\
\hline
2:1 & \begin{tabularx}{0.7\textwidth}{X} 我兒啊，你要在基督耶穌的恩典上剛強起來。 \end{tabularx} \\ \\ \relax
2:2 & \begin{tabularx}{0.7\textwidth}{X} 你在許多見證人面前聽見我所教導的，也要交託給那忠心而又能教導別人的人。 \end{tabularx} \\ \\ \relax
2:3 & \begin{tabularx}{0.7\textwidth}{X} 你要和我同受苦難，作基督耶穌的精兵。 \end{tabularx} \\ \\ \relax
2:4 & \begin{tabularx}{0.7\textwidth}{X} 凡當兵的，不讓世務纏身，好使那招他當兵的人喜悅。 \end{tabularx} \\ \\ \relax
2:5 & \begin{tabularx}{0.7\textwidth}{X} 運動員在比賽的時候，不按規則就不能得冠冕。 \end{tabularx} \\ \\ \relax
2:6 & \begin{tabularx}{0.7\textwidth}{X} 勤勞的農夫理當先得糧食。 \end{tabularx} \\ \\ \relax
2:7 & \begin{tabularx}{0.7\textwidth}{X} 我所說的話，你要考慮，因為主必在凡事上給你聰明。 \end{tabularx} \\ \\
[1ex]
\hline
\hline
\end{longtable}
$^{1}$網上高的弟兄姊妹平安,現場我們也有幫手的弟兄姊妹平安,好像很久沒見了,很開心今天在這裡和大家見.
中國人有句話,富不及三代,意思就是說一個很有錢的人,將他那些有錢的情況,比較難去到他的孫子那一代或者繼續下去.
其實某個程度上相對知識或者信仰上其實都是一樣,我們知道信二代或者信三代流失是很嚴重的,就算在南韓的教會,全國基督徒數目都曾經達到25-30\%.
信二代的流失問題仍然是嚴重的,怎麼樣將信仰傳下去呢,今天老傳道給我的題目就是信仰的成全,我們要將信仰傳下去.
有人開玩笑說神只有子女,神沒有孫子沒有孫女,因為基督徒不能因著遺傳將我們的信仰傳下去,我們都知道的,其實我在學院裡面教會增長.
都知道我們報道或者我們去叫教會的人數增加其實是有三個途徑的,第一個就叫做biological growth,就是說我們有第二代的信徒,第二代的信徒就是爸爸媽媽是信徒,他自小在教會長大.
他當然知道第二代信徒其實都需要有成全的,如果不是的話他們都會流失,這個是第一個教會增長的方向,第二個是什麼,就是叫做transfer growth,就是說在第二間堂會過來的.
第三個就是我們經常說很重要的conversion growth,就是說報道而帶回來,所以當我們知道有子女在我們手下的時候,我們要幫他們將他們的信仰實現出來.
我看到很多外國的宣教士,他們可能是第六代第七代的信徒,我自己聽到心裡面很感恩,因為我自己是第一代的信徒,我爸爸媽媽都不是信徒,我是第一代.
現在上一代開始受我們影響,我們都會影響下一代,但是我們怎麼傳下去呢,這個都很重要,但是都不會失望的,因為我在神學院裡面都有個學生他原來是第六代的信徒.
中國的學生,他在國內來的,他已經是第六代的信徒,所以其實還有的,我們要明白我們是要想我們怎麼樣承傳下去,又或者在座很多的兄弟姐妹,你們可能會想承傳信仰的承傳是不是全道人做的呢.
或者一些教會領袖先進去做的呢,和我好像沒什麼關係,今天講道的內容和我都沒什麼關係,其實不是的,我們知道中國人有另外一句話,三人行必有我師,我們三個人在一起,無論你是去講也好,你不講也好,你都可以成為別人的老師.
所以很可能我們當中有些弟兄姊妹你是做教導的,你是教會裡面領袖,你當然要思想我們怎麼樣將我們的信仰傳下去,但是我們當中有很大部分的弟兄姊妹是沒有這些職位的,但是其實無論你一言一行.
其實都會影響在你附近的人,你的子女,或者你的學生,或者你的鄰居等等,他們看到你說,甚至你不講你的信仰,你走出來其實都很大聲的,你就算不需要開口,你走出來都是很大聲的.
因為人家會看到你怎麼做的時候,他們就會心中有數,如果你說我要時時喜樂的,而你又不喜樂的話,人家會看你怎麼走,看你怎麼講,很多時候你走出來是被說出來,是重不重視.
我們看到保羅寫提摩塔後書的時候,是他人生最後的階段,他在羅馬被囚的當中,可能離開他被判死刑的日子都不遠了,是他在新約聖經裡面最後所寫的一本書.
保羅在提摩塔後書第一章最後那裡,他提到,凡在亞西亞的人都離棄我,這是你知道的,其中有肥吉路和黑魔奇,意思就是說好像很多人都離開我了,都好像覺得他的信仰去到某一個點的時候.
他突然說風一轉,就轉到提摩塔後書第二章第一到第七節,大家有聖經的,你們翻開聖經,或者看我們這裡的經文,是會單單主力講這裡的經文.
保羅在提摩塔後書稱提摩塔為什麼,他在第一節說,我兒啊,你要在基督耶穌的恩典上肛強起來,他稱提摩塔是什麼,是他的兒子,我們剛剛說成全,我們想著成全只是背前指的.
但是其實我們都知道實際上保羅不是提摩塔肉身的爸爸,因為提摩塔我們都知道提摩塔的爸爸是外邦人,媽媽是游費,但是我們知道保羅和提摩塔可以說是屬靈的父親和屬靈的兒子.
所以我們會發覺保羅是很著緊很著緊提醒提摩塔,因為他知道保羅知道他的生命就快結束了,他希望提摩塔是持守著這個信仰,將這個信仰傳下去,所以他給了他三個命令.
在這一段經文在這七節的經文裡面,保羅給了提摩塔三個命令,第一個命令就是說他要提摩塔在基督的恩典上肛強起來,他說你要肛強,提摩塔你無論在新教或者延教的時候,你都要肛強,就是說要加一些力量變得有能力.
我們知道提摩塔是較為柔弱的,所以在一張七節裡面,保羅提醒提摩塔,提一張七節,提醒提摩塔說他因為神賜給我們的不是膽怯的心,但是肛強仁愛緊守的心.
意思就是說保羅提醒提摩塔不要膽怯,要肛強,因為神給我們一個仁愛緊守的心,保羅提醒他首先是要肛強,他要自己認同自己所教,你要自己信自己所教的,自己走自己所教的,教到才有能力.
我們走和我們說是要如出一轍的,不可以心口不一樣,除了他要自己知道之後,保羅提醒他他要肛強為什麼,因為他不一定得到學習者的認同,大家有沒有想想,當你教人家怎麼做的時候,對方不是一定認同的.
所以保羅說你要站硬,你要站穩,你所學習的,不一定得到當,不單不得到學習者的認同,更加不一定得到當權者的同意或者贊賞的,弟兄姊妹現在大家都在想,我們可以說現在這一百個圈叫威權臨近的時間.
那究竟我們所講的我們所教的,究竟政府或者中國的官方會怎麼去看呢,今天拆十字架,我們要跟黨走那些話很多時候都在,但是保羅提醒提摩泰,說你要肛強,但是他更加提醒提摩泰不單只是肛強,他要在基督的恩典上高肛強.
所以他提醒提摩泰不是靠自己去肛強,我們知道提摩泰比較怯懦比較柔弱,但是保羅提醒他,你要在基督的恩典上高肛強,他首先提醒提摩泰,所以他在前面都說過了,神給很多恩典他的,他的外祖母和他的母親有離基.
兩個都是信主的,兩個都是在他的生命當中給很多的恩典給提摩泰,他更加提醒提摩泰,怎麼可以肛強起來呢,不是靠自己的,他是要依靠和他聯合的基督,他在第一章第九節說,神救了我們,以聖詔教我們,不是按我們的行為,乃是按他的旨意和恩典.
這恩典是萬古之先,在基督耶穌里賜給我們的,所以保羅提醒提摩泰,神已經借助耶穌基督賜給我們有這個恩典,所以我們是要在基督的恩典上高肛強起來.
我們因為主給了很多的恩典,我看過一些書,梁家倫院長和我們說,他說每個人在自己的過去的時候,都可以找到很多神的恩典,當我們去數算神的恩典的時候,我們要好好記錄下來,而當我們心裡面有猶疑,有不是很肯定的時間,我們可以不斷的去重溫去當者.
同樣的弟兄姊妹,你們首先要肛強起來,但不是因為你們自己很強,不是因為你們自己很厲害,沒錯,教會過去是有很多成功或者很多成長的,但是我們要在基督的恩典上肛強,數算神恩,你可能去數算一下.
數算主恩是很重要的,這個是保羅給提摩泰的第一個命令,就是他要在基督的恩典上肛強起來,保羅跟著給第二個命令提摩泰,他說你要在很多見證人面前聽見我所教訓的,也要交托那忠心能教導別人的人.
意思就是說保羅要提摩泰將教訓交托忠心能教導別人的人,我們去看一下保羅在這裡說了多少代的信徒呢,大家去小心的去看,他說你在很多見證人面前聽見我所教訓的,也要交托忠心能教導別人的人,我們看看第一個是誰,第一個就是.
五代的教訓,第一個是我,是保羅,第二個就是在很多見證人面前,是吧,保羅在很多見證人面前去教誰,是去教你,去教提摩泰,然後跟著就是提摩泰就要交托忠心能教導的人,這是第四代了,是吧.
然後第五代是誰,第五代就是別人,保羅說你要在很多見證人面前聽見我所教訓的,也要交托忠心能教導別人的人,我們去看回這樣的話,如果這樣是五代的話,假如他每一代只教了四個人,五代下來你猜有多少人呢,我們可以看過,第一保羅,是吧.
第二很多見證人,第三提摩泰,第四提摩泰所教導的,第五就是提摩泰所教導的學生,當保羅是一個,跟著他教四個,四個每人又教四個,就是十六個,十六個再教四個,就是每人教四個,六十四個,另外也再教,每人教四個的話,加起來就是三百四十一.
所以不要小看一個教四個,一個教五個,教六個,教七個,是更加多的人會受到這些教導,所以保羅提醒提摩泰是要什麼,要忠心,他說你要在很多見證人面前聽見我所教訓的,也要交托忠心能教導別人的人.
首先他提醒提摩泰,他自己要忠心,另外他要去找一些忠心能教導別人的人去做這個教導的工作.
我們知道不單單自己要遵守主導,也要為教會下一代的需要設計,我們知道當中的陶輝,他的太太教小朋友,很重要的教小朋友,因為我們要培育真理見證.

$^{41}$其實不單單是小朋友,大人我們初信的,都一樣要好好的接受這些培育,所以一個忠心的人是一個受人相信,是一個可靠的人,所以絕對不向逼迫或者錯漠屈膝的.
忠心的人心必須集中在誰身上,集中在基督的身上,記住我們不是要求我們當中有人要忠心於我們的傳道,我們要忠心於我們的導師,我們要忠心於某某某,我們要忠心於誰,基督.
所以當我們要思想的時候,我們要搞清楚我們的思想,我們要忠心於基督的身上,任何危險都不能夠有到他離開忠心的道路,任何假教訓的誘惑或者別人花言巧語的誘惑,都不能夠勾引他離開對耶穌基督的忠心的道路.
我們的生活和思想一定堅定不移,我們看到這裡保羅說的很清楚,這個是忠心能教導別人的人,意思就是說是受教的,是受教的,我們很多時候當我們有人會希望教我們一些東西的時候,我們未必很受教的,是要受教的.
能夠去做到受教,他在這裡給我們看到,領受信仰其實我們知道是可以聽回來的,另外還是怎麼樣,是看回來的,是看到人的見證,所以當我們要傳出去的時候,我們可以用我們的口來傳出去的.
用我們話語上傳出去,但是更重要的是什麼,我們都要在見證上傳出去,最理想就是心口如一,但是假如我們所說的和我們所做的是不可以夾到的話,人家會信我們哪裡,人家就看我們怎麼走.
同樣的弟兄姊妹,保羅在這裡提出要把信仰傳下去,其實每一個信徒都有責任,每一個信徒都有責任,可以用口去傳,或者你可以用你的行為去傳,剛才都說過,說了很多次了,無論你開口或者不開口,其實你的行為都已經顯出你的信仰.
要明白,我們要瞭解,所以保羅再一次提醒,提摩太第二條的命令就是你要去將我所交托你的,能夠去交托給忠心能教導別人的人,這個是他給提摩太的第二個命令.
其實都不容易的,但是當他給提摩太第三個命令的時候,對香港的信徒來說可能更難,他給提摩太第三個命令是什麼,他說你要和我同受苦難,保障基督耶穌的正義.
我們知道基督徒成為天國的精兵,好像你一信主一生下來就是的了,剛才都說過,每一個都要將你的信仰傳下去,所以不是簡的,我們知道如果你要考特警部隊是要考的,要你跑得多快,多fit,做到多少個人體上升那些東西.
但是做基督的精兵其實每一個都是,所以保羅他首先勸提摩太不要為主做見證為止,也不要以保羅作為主被囚的為此.
因為保羅受這些苦難是因為他是福音奉派做傳道,做使徒,做師傅,所以保羅現在是受苦難,他在第一章第八節說,你不要以給我們的主作見證為此,也不要以我這為主被囚的為此.
總要按神的能力與我為福音同受苦難,然後他在第一章第十一節說,我為這福音奉派做傳道的,做使徒,做師傅,為了緣故我也受這些苦難,然而我不以為此,因為知道我所信的是誰,也相信他能保存我所交付他的,直到那日.
保羅很清楚的告訴提摩太,提摩太要預備為主受苦,好像保羅為主受苦一樣,其實對於香港的信徒來說,這個對我們來說是比較新的一個經歷.
我記得我從小在教會成長,每個人都覺得做基督徒是好事,就算他們不信主,他們都會覺得基督徒是好的,基督徒是好人,基督徒人家說笑借錢不用還之類的,那些是錯誤的想法,但是他們一般來說都覺得做基督徒是好的.
但是我們到今天大家開始有些想法,究竟我們做基督徒有沒有苦難呢,究竟苦難臨到的時候我們怎麼辦,於是有很多人想辦法,我要走我要離開這裡等等,沒錯的,我今天不是評論這件事.
但是我們要知道保羅在這裡他告訴提摩泰,他說你是要和我同受苦難,保羅受了什麼苦,我們看到其實他被棍打被逼迫挨餓被囚被拒絕等等,他現在還在等著被判死.
但是保羅會說你要和我同受苦難,好像基督耶穌的精兵,通常來說當我們說比喻我們只是看一點,比喻可以用m這麼多一點來對待,他說基督的精兵是會受苦,怎麼受苦,我們知道其實受困的時候會受苦,受困又會有負重的訓練.
很重的東西去一些荒野的地方,又會長時間的訓練,又或者難度高的訓練,受困一定是苦的,其實受困是初步來的,當你去打仗的時候會怎麼樣,做精兵肯定是受困,受困是預備打仗的,打仗的時候會怎麼苦,受傷是會苦的,死亡也是會苦的.
更加苦是被俘虜,如果被俘虜了更加苦,保羅就用這裡叫提摩泰思想,他會不會要預備好,好像保羅一樣為基督的緣故受苦,然後他就用了三個例子,記住我們看例子都是一樣,都是單體,其中一點.
好像基督耶穌的精兵一樣同受苦難,保羅給了三個例子提摩泰,第一個例子他用當兵,凡是軍中當兵的,不將世無前身,好叫招他當兵的人,豈願,意思是什麼,你要當兵就要專心,要專心不要將世無前身.
很苦的,當兵要離開家裡,要去到軍營里集中的一起生活,當兵要吃大鍋飯,沒有私做,是大家一起做的,我們都知道他不是當兵,運動員牛車神都說過,阿里維斯都說受訓要離開家裡,很少見到家裡人.
很多時候很多好吃的東西都不能吃,就算很想吃的雪糕菠蘿包都吃不到,是很苦的,因為我們知道當兵的不要將世無前身,但是其實我們知道神會公平,很多不同的例子.
在不同的工場裡面,在自己讀神學的時候,神都公平,但是更重要的一點,好叫招他當兵的人,豈願,當兵的人要討誰的喜悅,要討教官的喜悅,教官招他當兵.
當我們是耶穌基督徵兵的時候,我們都應該專心的去求主的喜悅,不是求我們自己的喜悅,我們求主的利益而不是求我們自己的利益,不是求我們自己的前途,士兵要習慣於服從,士兵是沒有責任問為什麼這麼做,他有責任什麼,他有責任是去執行傳承的計劃.
所有決策的東西都是招他當兵的負責,是他的上司的負責,我們知道要聽上帝的聲音,有時候我們都不會明白為什麼是這樣的,就是我們要知道我們要討招我們當兵的人,這個是保羅比和紅鱷派第一個例子.
他給他第二個例子就是人若在場上比武,非按規矩就不立十半年,我們知道做運動員都是慘的,都是很多的苦難,當然是有很多的榮耀的,當你拿到第一的時候,當你拿到金牌的時候,是很多榮耀的,但是其實這麼多的運動員當中,只有某幾個拿到金牌.
李慧詩都知道在訓練的時候是很苦的,我之前聽到大英姐們和我們說,他說他要去跑馬拉松,原來跑馬拉松都是要受訓的時候,某幾天只能吃飯,不能吃肉,某幾天只能吃肉,不能吃飯.
因為他要拿蛋白質,有些又要拿澱粉質等等,去預備好自己,所以知道特別我先生周牧師以前是打柔道的,他說他有時候體重在兩級之間,他可能因為要比賽,他要節食.
讓他的體重可以減到低到某一個水平才能出賽,又或者要增長到一個地步才能去上一級,所以其實是很苦的,但是這裡除了做運動員是苦之外,他說非按規矩就不能夠得冠冕,得冠冕是運動員一個很重要的目的.
但是他要按規矩才能得冠冕,按參加比賽的人要嚴格遵守比賽的規矩,當我們是屬於耶穌基督徵兵的時候,規矩是誰定的,規矩不是我們定,也不是我所跟的老師或者我所跟的導師,是誰定的.
是神定的,我們是要按神所定的規矩比賽,如果不是的話就會被淘汰,所以我們要找清楚基督給我們的規矩,然後我們就去跟,然後我們無論多辛苦我們都要這樣去比賽.
然後第三個例子,努力的農夫理多思得能成,我想今天來說不是那麼多人去做農夫的工作,我剛剛剛才來的時候看到紅人堂下面有個海報,你們現在有個地可以找人一起耕,這個是很難得的經歷.
農夫一定是努力的,農夫也要艱苦的去等,等收成,等收購,可能假如你種一些菜短時間一些,可能二十天左右就收成,但當你如果要種稻米等等,可能半年才有收成.
是要等的,日復一日的努力的去等,但是我那時候去湖南的時候去短孫發覺他們大部分都是做農夫的,他們平時的時候他們很早去崇拜的,他們平時是七點鐘崇拜的,禮拜天很早排練,跟他們一起去敬拜,你知道他們龍王的時候是幾點鐘崇拜的.
他們不是不崇拜的,他們龍王的時候是六點鐘崇拜,六點鐘崇拜完他們就會去田裡工作,他們可能今天就去A那個田裡工作,其他人就會一起幫忙上收成,明天就收完A那個田,明天就去B那個,他們就大家一起去幫忙.
但是他們是努力的,但是同時這裡更加告訴我們,是你當先得糧食,神是給我們有這個好處,讓我們好像龍妖一樣,我們先得糧食,如果你是熟糧食那是什麼,就是糧,我們是要努力的去讀好我們的聖經,努力的去明白神的心意.
我們自己得了靈壽,然後我們就可以傳下去傳給大力的,這個是保羅給這三個例子,第一個是做政兵,第二個就是做運動員,第三個就是農夫,當時來說很多人都會很明白得青就會變,然後保羅在最後這一段的說話.
我所說的話你要思想,因為凡事只必吸引聰明,請姐妹,善性信仰不是盲目的跟隨,善性信仰是要思想的,子也給我們清明,我們要思想我們是跟隨的,我們要思想我們是轉移的.
與新教和異教教導自己附近的人,與新教和異教教導自己的學生,自己的異類,自己下一代的,我說了很多年了,我都很喜歡用這一句來結束的,就是愚弄何乾.

$^{81}$我和大家說了大概38分鐘左右,說了這麼久,你會說對你有什麼關係,你有沒有聽到,你有沒有聽到保羅怎麼去向提摩泰國,你就對這件事有什麼感覺,大家一起去思考.
我們都會將我們的信仰傳下去,無論我們用口或者用言語,願你使用我們,願這個信仰可以繼續的成長,叫弘恩堂或者弘恩堂附近不斷的傳下去,以亮明,得以平高的,聽我們的祈禱,和耶穌基督的名求,阿彌陀佛.
\newpage



\section{出埃及記 3:1-15}
\label{sec:sVUIck2trk4}
\textbf{2021.01.24《再思上主的呼召》出埃及記 3\_1-15( 和修版)  曾裔貴牧師}
\newline
\newline
連結: \href{https://youtube.com/watch?v=sVUIck2trk4}{\texttt{https://youtube.com/watch?v=sVUIck2trk4}} ~~~~ 語音日期: 2021-01-26
\newline
\newline
\hyperref[sec:hPE_EBTYs3o]{\small{< < < PREV SERMON < < <}}
~
\hyperref[sec:index_chronic]{\small{[返順時目]}}
~
\hyperref[sec:index_scriptual]{\small{[返順卷目]}}
~
\hyperref[sec:tlLyIq5teaA]{\small{> > > NEXT SERMON > > >}}
\newline
\newline
出埃及記 3:1-15
\newline
\begin{longtable}{cl}
\hline
\hline
章節 & 經文 (和合本修訂版)\\
\hline
3:1 & \begin{tabularx}{0.7\textwidth}{X} 摩西牧放他岳父米甸祭司葉特羅的羊群，他領羊群往曠野的那一邊去，到了神的山，就是何烈山。 \end{tabularx} \\ \\ \relax
3:2 & \begin{tabularx}{0.7\textwidth}{X} 耶和華的使者在荊棘的火焰中向他顯現。摩西觀看，看哪，荊棘在火中焚燒，卻沒有燒燬。 \end{tabularx} \\ \\ \relax
3:3 & \begin{tabularx}{0.7\textwidth}{X} 摩西說：「我要轉過去看這大異象，這荊棘為何沒有燒燬呢？」 \end{tabularx} \\ \\ \relax
3:4 & \begin{tabularx}{0.7\textwidth}{X} 耶和華見摩西轉過去看，神就從荊棘裡呼叫他說：「摩西！摩西！」他說：「我在這裡。」 \end{tabularx} \\ \\ \relax
3:5 & \begin{tabularx}{0.7\textwidth}{X} 神說：「不要靠近這裡。把你腳上的鞋脫下來，因為你所站的地方是聖地」。 \end{tabularx} \\ \\ \relax
3:6 & \begin{tabularx}{0.7\textwidth}{X} 他又說：「我是你父親的神，是亞伯拉罕的神，以撒的神，雅各的神。」摩西蒙上臉，因為怕看神。 \end{tabularx} \\ \\ \relax
3:7 & \begin{tabularx}{0.7\textwidth}{X} 耶和華說：「我確實看見了我百姓在埃及所受的困苦，我也聽見了他們因受監工苦待所發的哀聲；我確實知道他們的痛苦。 \end{tabularx} \\ \\ \relax
3:8 & \begin{tabularx}{0.7\textwidth}{X} 我下來是要救他們脫離埃及人的手，領他們從那地上來，到美好與寬闊之地，到流奶與蜜之地，就是迦南人、赫人、亞摩利人、比利洗人、希未人、耶布斯人之地。 \end{tabularx} \\ \\ \relax
3:9 & \begin{tabularx}{0.7\textwidth}{X} 現在，看哪，以色列人的哀聲達到我這裡，我也看見埃及人怎樣欺壓他們。 \end{tabularx} \\ \\ \relax
3:10 & \begin{tabularx}{0.7\textwidth}{X} 現在，你去，我要差派你到法老那裡，把我的百姓以色列人從埃及領出來。」 \end{tabularx} \\ \\ \relax
3:11 & \begin{tabularx}{0.7\textwidth}{X} 摩西對神說：「我是甚麼人，竟能去見法老，把以色列人從埃及領出來呢？」 \end{tabularx} \\ \\ \relax
3:12 & \begin{tabularx}{0.7\textwidth}{X} 神說：「我必與你同在。這就是我差派你去，給你的憑據：你把百姓從埃及領出來之後，你們必在這山上事奉神。」 \end{tabularx} \\ \\ \relax
3:13 & \begin{tabularx}{0.7\textwidth}{X} 摩西對神說：「看哪，我到以色列人那裡，對他們說：『你們祖宗的神差派我到你們這裡來。』他們若對我說：『他叫甚麼名字？』我要對他們說甚麼呢？」 \end{tabularx} \\ \\ \relax
3:14 & \begin{tabularx}{0.7\textwidth}{X} 神對摩西說：「我是自有永有的」；又說：「你要對以色列人這樣說：『那自有永有的差派我到你們這裡來。』」 \end{tabularx} \\ \\ \relax
3:15 & \begin{tabularx}{0.7\textwidth}{X} 神又對摩西說：「你要對以色列人這樣說：『耶和華－你們祖宗的神，就是亞伯拉罕的神，以撒的神，雅各的神差派我到你們這裡來。』這是我的名，直到永遠；這也是我的稱號，直到萬代。 \end{tabularx} \\ \\
[1ex]
\hline
\hline
\end{longtable}
$^{1}$各位弟兄姊妹早晨.
再一次能夠代表部會.
來在雄恩和大家.
將神的話語去講解.
剛才.
曾姑娘讀經.
帶出到.
神的威嚴.
摩西.
一個戰經的心靈.
那個戰經的.
那個說話的聲調.
當然.
原文是希伯來文.
但是都反映到.
當時的情勢.
一個高高在上.
在荊棘樹.
沒有銷壞的一個.
異象神跡背後.
說話的這一位.
摩西知道是非同小可.
不是尋常的.
一個人跟祂說話.
因為能夠讓一個牧羊人.
在曠野.
看到一個大異象.
你就知道.
這一段聖經.
是相當重要的呼召.
今天很想藉著這段聖經.
不是跟一些初信弟兄姊妹說.
當然崇拜每一個人坐在這裡.
你都應該有領受.
不過.
有一個很重的負擔.
是跟一些事奉.
跟一些心靈已經失落的弟兄姊妹說.
包括在家中.
大家來到看著.

$^{41}$直播的弟兄姊妹.
又或者其他的教會.
偶然間進入到這個視頻.
將來成為一個記錄.
很想彼此勉勵.
我們的信心.
我們的事奉.
是不是已經失落了那個初心.
那個蒙召的.
那個服侍的力量呢?.
我們簡單有一個祈禱.
我們的主.
你是呼召的主.
你看到.
以色列人受痛苦.
他們的哀聲.
達到你的耳中.
我們坐落在這個社區.
其實同樣有很多人.
他們不認識你.
他們甚至乎都不知道自己.
在發出痛苦的哀聲.
我們很盼望.
我們不會因為.
我們自己內心的痛苦.
哀聲.
來遮蓋了周圍.
地方未蒙救恩的人.
他們的聲音.
求主幫助我們.
我們見到一個失落的老人家摩西.
如常每天.
做他喜歡重複重複的工作.
就是帶著羊群.
去到這個曠野.
去到這個荷列山.
親愛的主.
求你將這段聖經.
藉著聖靈的工作.
叫很多很多弟兄姊妹.

$^{81}$心靈再一次甦醒過來.
看到你是自有永有的主.
是亞伯拉罕的神.
以撒的神.
和雅各的神.
更加是摩西父親.
父親輩的神.
以致他們看到的.
不再僅僅是埃及的奴役.
而是看到.
有一個釋放.
等待他們去執行.
今天早上我們等待在你的面前.
禱告奉耶穌基督的名.
阿們.
有一位年輕人.
在台灣.
他讀書已經讀到研究院.
讀到碩士.
他找工作就有些障礙.
有些困難.
因為工作因為生活.
沒辦法.
他就去找一些.
做一些吊臂車.
起重的很粗重的工作.
人工很高.
折合台幣.
一個剛剛出來.
有六萬元的台幣.
和一般的文職.
都是兩三萬.
他就覺得有些介意.
我讀這麼多書.
已經是碩士研究生.
為什麼要做這些工作呢.
他有一次就有機會.
就問台灣的首富.
郭台銘先生.
郭台銘先生就很有智慧地.

$^{121}$來回覆他.
他說.
如果你夠聰明.
如果你對那個事情.
你有一個很強的心.
就算那個工作本身.
你都可以改變.
接著他說了一個例子.
他說中國大陸.
有一位當時的首富.
梁允根.
中國的三一重工集團的老闆.
他是一個普通工人.
但在他工作裡.
不斷發現那些機械.
可以在他做的時候.
可以改良.
不斷不斷.
他就嘗試不斷地改良.
慢慢慢慢就自己創業.
慢慢慢慢經過一段長時間之後.
他就成為了三一重工的董事長.
就成為了在國內.
一個很出名的.
由一個很窮的基層工人.
慢慢就致富.
郭台銘就用這個例子.
只要你整個人.
你的內心是.
他用英文.
你的內心夠聰明.
那些工作不會再是卑賤.
你有一個錯覺就是.
你讀那麼多書.
你現在做這些工作.
好像有一個錯配.
其實不會錯配.
弟兄姊妹這個例子.
給我們看摩西的時候.
摩西內心裡面有很重的抱負.

$^{161}$四十歲的時候.
不夠時間和大家讀很多聖經.
第二章就已經有這樣記載.
四十歲的時候.
他很想來到去拯救以色列人.
有一天見到一個埃及人.
和以色列人爭吵.
他就錯手打死了那個埃及人.
後來這件事被法老王知道.
他就很怕逃走.
所以當他見到這個異象.
神問他的時候.
他說我是何許人.
我怎麼可以回去見法老呢.
他自己知道.
這個在他內心的這條刺.
他是一個逃犯.
一個逃犯又怎麼能夠回去見法老王呢.
其實法老王當時已經死了.
是有新的王.
這裡告訴我們.
摩西曾經內心很想事奉神.
很想帶領自己的同胞.
來離開埃及.
不過一個人的力量.
怎麼能夠去改變整個埃及帝國呢.
他所做的無論他的時間.
和他的力量.
都不足以去改變.
由四十歲逃亡.
聖經就沒有再記載了.
第三章剛才主席讀了的經文.
就告訴我們.
摩西在這個時候已經八十歲了.
八十歲的老人家.
有說四十年的時間.
在曠野如常.
每天這樣牧放外府的羊.
每天的生活.
就是他可以來做的.

$^{201}$也是慢慢地.
似乎應該都可以隔兄就熟.
越來越在他的牧羊裡面.
來成為他餘生餘下的工作.
但這個是不是上主給他的呼召呢.
當然我們看完這段聖經就知道不是.
原來四十年是準備他的生命.
是讓他能夠將來為神做.
神要他做的工作.
也可以這麼說.
神的時間在摩西的熬煉裡面.
現在夠了.
神要來大大地使用他.
讓他成為一個可以帶領自己的同胞.
離開埃及的這一位領袖.
留意一下這些經文很重要.
就是第五節.
相信你和我如果在那個場景.
看到荊棘燒著.
曠野就從來都不特別的.
我們來不及清明.
同樣都會看到山火的.
因為有些人燒.
但是燒完就一定會救得息.
特別的地方就是.
為什麼荊棘燒著又不會燒毀呢?.
神要用一個這麼大的神跡.
不是普通的自然的現象.
是一個神跡的一個意象.
讓摩西要走過去.
當走過去看那些荊棘的時候.
荊棘有聲音出來.
這裡就開始重點了.
就是說原來那個意象.
只不過是要引發摩西的好奇.
好像要挑動他內心.
他裡面本來已經平靜如鏡的內心.
突然間湧起一個巨浪.
什麼原因?.
為什麼我會看到這麼特別的奇景呢?.

$^{241}$當他走過去的時候.
你要留意這裡很重要.
摩西原來立即站在神的面前.
那個地方其實是一個曠野.
是一個很平常不過的地方.
神的同在立即成為神聖.
摩西立即要驚覺自己是一個罪人.
是一個在神.
不是在法老面前這麼簡單.
不是他內心自己的罪疚感湧上來這麼簡單.
而是在神的面前.
他現在所站的地方.
就是一個神聖的聖地.
就是將來的聖殿一樣的聖地.
聖殿不是因為那個地方的華麗.
是因為有神的同在.
所以每一個教會.
有奉主名聚集的地方.
是立即成為神聖.
你回來是一個聖堂.
你坐在這裡.
你事奉.
是在服侍這位聖潔的造物主.
摩西立即發現自己.
置身在神的面前.
是神直接向他宣告.
從來沒有人知道.
這樣的一個環境.
突然間變成神聖.
好像一個時空隧道.
他去到天堂一樣.
一個這樣的心態.
你無法想像那個戰境.
是多麼可怕.
神就在你面前.
但是摩西脫鞋.
代表著他要學習謙卑信服.
在神的面前.
如果不是神.
很難大用你摩西.

$^{281}$讓他好像經驗到.
現在神有什麼.
要來去差遣我呢.
當然立即的推搪是很正常的.
我是何許人.
第十一節.
竟能夠去見法老.
就是內心湧出.
神啊.
你知不知道.
如果你是全知的.
你知不知道.
我四十歲的時候.
打死了埃及人.
我逃難去到這個曠野.
輾轉就認識了我的太太.
這個西波拉.
又得到她的外父.
葉特羅的信任.
我才可以現在牧羊群.
做一個牧羊人.
還有就是創世記四十九章.
埃及人是不喜歡牧羊的.
對牧羊的人是討厭的.
摩西都知道自己是一個牧羊人.
這麼卑微的人.
怎麼去見法老呢.
不過.
將你自己的生命.
放在上主的手.
在這個時候.
上主要來去溫暖摩西的心靈.
讓他看到.
不是藉著你.
是藉著我和你同在.
就帶來你的生命.
會有一個不一樣的改變.
這個就是摩西的呼召.
回到我們現代.
回到我們弟兄姊妹.

$^{321}$我們的事奉.
很多在教會的領袖.
有時慢慢慢慢在工作的失落.
困難.
生活面對的種種的挫折.
有時就會有一種的分割.
慢慢慢慢生活就還生活.
事奉教會的聚會就還教會的聚會.
禮拜日回來好像是一個透氣.
不是想著來到.
很認真站在聖地的事奉神.
神有什麼的差遣.
我要去做一些很困難的事.
你想想.
你要帶領你的族人.
二百萬人離開一個強大的帝國.
往後由第四章開始.
一直到十二章.
其實十三章就是他們過紅海.
你就會知道十災.
摩西是怎樣來的.
透過神去行神蹟.
讓法老慢慢慢慢低頭.
但是法老從來沒有低頭.
他的心硬如石頭.
然後還要在紅海追殺摩西.
最終就是他們的軍兵淹死了.
你可以想像得到.
原來摩西將要面對一個這麼強大的對手.
不過神沒有告訴他這麼大的困難.
不過摩西慢慢慢慢經驗到.
最寶貴是有神同在.
還有神繼續在他生命裡面的呼召.
讓他的內心再一次來到他知道.
原來他的人生要放在這個位置才是最好.
弟兄姊妹.
有沒有將工作和你的生活.
和你的教會分割了呢?.
有沒有在你的生活裡面.
你只是形形意意.

$^{361}$這些是世俗的工作.
和我的事奉無關.
我只是賺錢.
我只是做一個小世界.
就是這樣.
回到教會.
我是一個坐在這裡很乖.
聽道.
就是這樣.
我的生活是沒有人可以進入.
或者來干預的.
會不會有這樣的分割呢?.
會不會就是神給你的那個召命.
你應該在你的職業.
在你人生的職場.
你來投放.
那個地方就是神聖.
因為神跟你一起.
那個地方.
那個職業成為神聖.
有沒有這樣想自己的.
六天的人生的心靈呢?.
有沒有帶上神.
其實不需要你帶的.
神無所不在的.
不需要你帶著神去到哪裡的.
有沒有想過.
你星期一至六.
你那個職業的操守.
你內心的事奉神的心態.
有沒有將這個神.
是聖的神.
神聖的.
你所站的是聖地.
不只是這個.
呱呱的港台才是神聖.
是你和我在職場.
你開車的時候.
有沒有隨便就違反交通規則呢?.
雙白線就說過就過.

$^{401}$紅燈沒人見到照衝.
會不會這樣呢?.
抑或.
我所站的地就是聖地.
使徒保羅在提摩太后書.
他有很深的對自己的認識.
他說基督耶穌降世.
為要拯救罪人.
這話是可信可佩服的.
在罪人當中.
我是一個罪魁.
不只是這樣來背.
一些探索班金句.
更重要的是他自己的內心啟發.
然而我蒙了憐憫.
是因耶穌基督要在我這罪魁身上.
顯明他一切的忍耐.
給後來信他得永生的人作榜樣.
不介意說自己的過去.
其實是跟一個傳道人年青的提摩太說.
我是一個罪魁.
我曾經不認他.
迫害信他的人.
這樣的保羅.
所以他就發現自己現在的人生.
他每天就是在服侍一位這麼聖潔的主.
可以說他明白.
根本沒有一個地方是不神聖的.
沒有一個職業是不神聖的.
每一個人只要你事奉這位主.
你就是神聖的.
這樣理解自己.
但反過來.
如果你不是在事奉神.
你的人生大部份都是為你自己.
為你自己的人生的驕傲.
為你自己的野心理想.
你可能很成功.
不過神體不是他喜悅你的人生擺放.
這樣你就算站在一個很神聖的地方.

$^{441}$你都變得不神聖.
求主幫助我們每一位弟兄姊妹.
我再說.
今天這篇信息是向一些經歷過耶穌的愛.
亦事奉過主的弟兄姊妹說.
我們一起同心看到.
為何要呼召摩西.
對他特別情有獨鍾.
在萬民當中揀選了他.
其實他本來就應該浸死在尼羅河.
真的正在去的那個屠宿商.
是埃及王法老的女兒.
公主見到他就收養他.
否則他跟其他歲數的小朋友已經浸死了.
摩西這個名字就是從水中拉上來.
即是他的生命是在神的手中.
弟兄姊妹.
一起看到我們的生命.
因為神.
因為這位亞伯拉罕以撒雅各的神.
而變得尊貴.
是因為神才令你和我變得尊貴.
離開了神.
你所佔的一切都不是聖地.
是尋常的曠野.
但你在神裡面.
你所做的事.
你的工作就是神聖.
你的職業就是神聖.
因為你整個心態都已經純正.
當你六天都是為你自己生活工作的時候.
星期日就更加拼命來事奉神.
更加是為我們將福音傳給有需要的人.
記住我們不要只聽到自己的痛苦.
自己的哀聲.
摩西四十年天天做牧羊人.
不是他甘心.
但情勢令他繼續這樣.
不過原來神不是期望他這樣.
而是希望他能夠帶領族人離開埃及.

$^{481}$成就他往後在神眼中的身份.
這是摩西.
是不是你的故事呢?.
是不是你內心裡面很多很多的痛苦,哀聲.
而遮蓋了周圍真正需要福音的人的哀聲呢?.
這是第七節.
我確實看見我百姓在埃及所受的困苦.
我又聽見因受監宮苦袋所發的哀聲.
我確實知道他們的痛苦.
最知道的就是摩西.
最知道這件事的已經四十年.
天天他都知道.
但是手無寸鐵.
個人自己都無可奈何.
自己的內心好像他的軀殼和他的靈魂有些分割.
曾經幫過,但是越幫越忙.
幫倒忙,做不到任何事.
但是他難道不知道那個困苦的哀聲嗎?.
難道我們不知道信不到耶穌,上不到天堂的人.
他們要永遠沉淪嗎?.
知道的,但是有什麼辦法呢?.
這個時候我們需要再一次看回.
這位做物主完全沒有因為摩西的過去.
就將他定性了.
這一次的出現顯現.
正正就是告訴他.
摩西,你在我的手中.
你就可以改變.
因為我是自有永有.
我們需要可能是一些當頭棒喝的鼓勵.
讓我們不要再留戀我們自己.
我們認為的對錯.
我們應該回去看神.
怎麼看周圍的地方的需要.
嘗試一個例子給大家一個共鳴.
鄭丹瑞可能大家都認識.
最近在網上都有一些很正面的信息發佈給市民.
讓大家在疫症期間都有很多的鼓勵.
他講一個見證他自己很多年前.
來到去做DJ的時候.

$^{521}$他那時候在商台做得很出色.
後來就轉工去了港台.
香港電台.
這樣來到去安排了一個黃金時段.
早上的九點到十一點.
音樂彈這個節目.
他都知道以為自己是很受歡迎的.
這樣.
過了一段時間之後.
當時香港電台就要.
有一些的緊縮.
他就無意中有一天看娛樂版.
報紙.
那時候不會有那麼多iPad.
不會上網看新聞的.
看報紙娛樂版.
看一看.
為什麼報紙說香港電台會cut budget.
第一個cut就是音樂彈這個時段的節目.
他就很生氣.
為什麼我不知道.
而娛樂版就會登出來.
他就回去問高層.
覺得很生氣.
為什麼會這樣.
完全自己不知道.
輾轉.
高層就安撫他.
就是說現在要cut budget.
首先就來到去.
這個高層其實都欣賞他.
就幫他爭取.
爭取就是說還給你一點時間.
如果你做到成績.
就不來炒你了.
這樣.
就幫他爭取到.
在凌晨的兩點.
到六點.
做一個叫做輕彈淺唱不夜天的一個節目.

$^{561}$凌晨兩點.
各位.
由一個黃金時段.
九點至十一點的早上.
去到半夜.
這樣的一個落差.
如果你是阿旦.
你是鄭丹瑞.
你會有什麼心情呢.
他在一個見證裡面.
他都說.
敷衍做.
充滿憤怒地做.
簡直是一種對我的一個迫害.
這樣的心態.
你想想自然的成績.
照樣播.
那些歌毫無生氣.
毫無串連.
做的人.
大家都感覺到他充滿了一種敷衍.
那時候就有這個詞.
HEA.
敷衍真是一個很好的詞.
很反映到那個人的內心.
那個態度.
對那個工作的熱誠.
完全沒有.
叫做HEA.
有一天.
就有一個派信的職工.
阿錦叔.
就派給那些同事.
通常聽眾來信.
等等.
因為你是DJ.
那封信是給他.
他也奇怪.
為什麼會有信呢.
沒有抬頭.

$^{601}$當然信封面當然有.
但信的內容是沒有抬頭的.
只有一句說話.
你不想做就不要做.
你這樣HEA做.
我每晚.
就需要聽著音樂來工作.
這樣的一封信.
很噁心.
很沒禮貌.
但對鄭丹瑞來說.
他突然當頭棒喝.
是喔.
為什麼我不用心.
做好我現在的工作.
讓聽眾.
讓高層看到.
我不是一個第一時間.
要開刀cut budget的人呢.
就是這樣.
用心的.
在半夜.
做好他的工作.
接著開始有不同的聽眾來信.
多謝他.
但鄭丹瑞.
在他的見證裡說.
我等一封信.
大家都知道.
等一封信.
就是沒有抬頭.
罵他的人的信.
有一天過了一段時間.
他真的收到這封信.
他從筆跡認得這封信.
就來到.
這次有抬頭了.
阿旦.
你可不可以四點鐘.
幫我播這首歌.

$^{641}$他覺得是一個.
他自己生命的鼓勵.
重新讓人肯定.
他真的看到自己.
原來真的可以.
不過是透過一個.
這樣的手術.
一個落差.
親愛的弟兄姊妹.
大家有一個.
事奉主的心.
大家在人生的路.
總有不同的.
跌跌碰碰.
但我們來到神的面前.
首先我們要記住.
這位造物主.
祂才是唯一的聖潔.
我們去到祂的根前.
立即就要脫鞋.
這個要成為.
我們很強烈的意識.
如果不是我們.
不能事奉這位聖潔的造物主.
因為祂神聖.
到不得侵犯.
但當你願意.
放下的時候.
你就會看到.
聽到.
真正的哀聲.
其實不是你和我的內心.
我們不要太多.
不斷的無病呻吟.
我們要看到的是.
周圍的福音的需要.
看到有很多人.
從來沒有機會聽福音.
那種的痛苦.
那種的哀聲.

$^{681}$其實已經達到上主的耳中.
記住.
亞伯拉罕的神.
以撒的神.
雅各的神.
即是說神在四百年前.
已經有救世的計劃.
如果摩西繼續.
在他自己的自我裡.
不肯跳出來的時候.
他仍然是一個.
牧羊的摩西.
他仍然.
只會帶著他自己的痛苦.
過他餘下的人生.
但是.
神不希望.
服事主的弟兄姊妹.
這樣.
讓我們大家一起在這個聖地.
一起去服事.
這位聖潔的主.
就是將福音帶到更多的人群當中.
求主幫助我們每一位.
我們來到去學習.
聖經的真道.
在思想.
上主的呼召.
我們一起簡單祈禱.
再一次多謝主.
讓今早.
藉著這段聖經.
我們彼此的去勉勵.
親愛的主.
你是聖潔.
是尊貴的.
求你讓我們看到.
其實我們在每一個環境.
只要你的同在.
你的愛在我們內心.

$^{721}$就使我們能夠.
充滿了這一份對你的敬拜.
還有.
我們在事奉裡.
能夠將人帶到你的面前.
我們自己.
都感受到這一份神聖的尊貴.
在我們的生命裡.
不斷地會有.
那個成長.
就是越來越像我們的主耶穌.
有真理.
仁義和聖潔.
多謝主讓我們今早.
一同的來思想.
學習你自己的真道.
願你的工作.
在我們人的心裡.
再一次感動光照.
我們這樣禱告.
奉耶穌基督的名.
阿們.
\newpage



\section{耶利米書 1:1-10}
\label{sec:tlLyIq5teaA}
\textbf{2021.01.31《亂世中的召命》耶利米書 1\_1-10( 和修版)  楊穎兒導師}
\newline
\newline
連結: \href{https://youtube.com/watch?v=tlLyIq5teaA}{\texttt{https://youtube.com/watch?v=tlLyIq5teaA}} ~~~~ 語音日期: 2021-02-02
\newline
\newline
\hyperref[sec:sVUIck2trk4]{\small{< < < PREV SERMON < < <}}
~
\hyperref[sec:index_chronic]{\small{[返順時目]}}
~
\hyperref[sec:index_scriptual]{\small{[返順卷目]}}
~
\hyperref[sec:xeKLoIITcbQ]{\small{> > > NEXT SERMON > > >}}
\newline
\newline
耶利米書 1:1-10
\newline
\begin{longtable}{cl}
\hline
\hline
章節 & 經文 (和合本修訂版)\\
\hline
1:1 & \begin{tabularx}{0.7\textwidth}{X} 這些是便雅憫地亞拿突城的祭司，希勒家的兒子耶利米的話。 \end{tabularx} \\ \\ \relax
1:2 & \begin{tabularx}{0.7\textwidth}{X} 亞們的兒子猶大王約西亞在位第十三年，耶和華的話臨到耶利米。 \end{tabularx} \\ \\ \relax
1:3 & \begin{tabularx}{0.7\textwidth}{X} 從約西亞的兒子猶大王約雅敬在位的時候，直到約西亞的兒子猶大王西底家在位的末年，就是第十一年五月間耶路撒冷被擄時，耶和華的話也常臨到耶利米。 \end{tabularx} \\ \\ \relax
1:4 & \begin{tabularx}{0.7\textwidth}{X} 耶利米說，耶和華的話臨到我，說： \end{tabularx} \\ \\ \relax
1:5 & \begin{tabularx}{0.7\textwidth}{X} 「我尚未將你造在母腹中，就已認識你；你未出母胎，我已將你分別為聖，派你作列國的先知。」 \end{tabularx} \\ \\ \relax
1:6 & \begin{tabularx}{0.7\textwidth}{X} 我就說：「唉！主耶和華，看哪，我不知道怎麼說，因為我年輕。」 \end{tabularx} \\ \\ \relax
1:7 & \begin{tabularx}{0.7\textwidth}{X} 耶和華對我說：「不要說『我年輕』，因為我差遣你到誰那裡去，你都要去；我吩咐你說甚麼話，你都要說。 \end{tabularx} \\ \\ \relax
1:8 & \begin{tabularx}{0.7\textwidth}{X} 你不要怕他們，因為我與你同在，要拯救你。這是耶和華說的。」 \end{tabularx} \\ \\ \relax
1:9 & \begin{tabularx}{0.7\textwidth}{X} 於是耶和華伸手按住我的口，對我說：「看哪，我已將我的話放在你口中。 \end{tabularx} \\ \\ \relax
1:10 & \begin{tabularx}{0.7\textwidth}{X} 我今日立你在列邦列國之上，為要拔出，拆毀，毀壞，傾覆，又要建立，栽植。」 \end{tabularx} \\ \\
[1ex]
\hline
\hline
\end{longtable}
$^{1}$弟兄姊妹大家好.
帶我來到紅恩堂裡面去實習.
多謝紅恩堂.
特別是老傳道.
給我這個機會.
在這個主日.
和大家分享神的話語.
讓我們先低頭祈禱.
親愛的神.
求你使用孩子.
說出你的話語.
求你開會眾的心.
眼和耳朵.
讓你的話成為他們的幫助.
奉主明求.
阿們.
今天很想和大家.
思想什麼是召命.
召命是信仰核心的一個基本觀念.
上帝的呼召.
不只是限於聖功.
事奉或職業.
是包括整個人的生命.
即是當上帝.
向人發出呼召.
人要以生命來回應.
這就是召命.
召命是沒有年紀.
沒有次數.
神很可能在不同的人生階段.
去呼召我們每一天.
都以我們的生命.
去回應神.
今天我要和大家講的題目是.
亂世中的召命.
我相信如果用亂世.
來形容今天的情況.
大家都會認同.
怎樣才算亂呢?.
亂就是失去了一個應有的秩序.

$^{41}$失去了平衡.
好像在過往的兩年.
我們經歷了社會上的混亂.
世界性的混亂.
令我們失去了人.
渴望甚至追求.
那種安穩的秩序.
我們看看聖經.
似乎神要我們追求的.
不是安穩.
而是安穩的秩序.
我們看聖經.
我們追求的不是安穩.
幾乎整卷聖經.
都是在說.
在混亂的世代裡.
神有審判.
祂對子民有教導.
更多的是.
神的應許.
怎樣一次又一次與祂的子民同在.
拯救祂揀選的子民.
我們今天所講的.
耶利米書.
整卷書的背景.
都是在說神的審判將要來到猶大國.
越來越亂.
最後.
整個猶大國都要滅亡.
正正就在.
國家滅亡之前.
神父召耶利米成為先知.
他猜險.
去宣講.
神審判的訊息.
剛才我們也讀過經文.
我們會發現.
當神父召耶利米的時候.
耶利米用了一些藉口.
去回應他的召命.

$^{81}$同樣.
我們會否因環境的改變.
或眼前的混亂.
我們會否像耶利米一樣.
用藉口.
去迴避神在我們身上的召命呢.
讓我們一起.
從耶利米書第一章一至十節.
學習.
不要只看見世上的混亂.
也不要看見眼前.
自己的限制.
反而要轉移目光.
去明白神對我們的應許.
讓我們先進入.
經文的時代背景.
去看看當時的混亂.
經文說到第二節.
是猶大王約西亞在位的第十三年.
大概是公元前.
627年左右.
第三節.
說到猶大王西底迦在位的末年.
這裡大概有.
四十多年.
猶大王在世.
在世的第十三年.
這裡大概有.
四十多年的日子.
是橫跨三個王朝.
約西亞.
約瓦勁和西底迦.
這三代王朝裡面.
神的話語一直.
淋到耶利米.
可能在我們當中.
已經有人知道.
這裡不止三代王朝.
乃是五代王朝.
當中也包括了.

$^{121}$兩位只是做了三個月左右的君王.
不過經文沒有提及.
所以我們也不多說.
這兩位君王.
在當時的猶大國.
這幾代王朝裡面發生了甚麼事呢?.
公元前627年左右.
猶大國就是紫色的部分.
是一個小國.
周邊有亞述帝國.
埃及帝國.
巴比倫帝國.
他們的勢力也有限.
猶大國在沒有外憂的情況下.
相對來說.
是比較平靜.
於是約西亞在在位的期間.
實行了制治的改革.
包括廢除外族影響的.
敬拜的制治方式.
務求帶領整個民族.
回到利尾祭司那種敬拜的傳統.
到了609年.
約西亞試圖阻止法老.
去援助亞述人.
於是猶大國就陷入了.
與埃及人的戰爭.
最後他也在巴比倫的戰役中被殺.
而後來他的弟弟.
約瓦勁登上王位.
不久後就被擄到巴比倫去.
屈服在巴比倫王之下.
最後到597年.
約瓦勁的叔叔西底加上位.
但最終猶大國.
還是被巴比倫.
佔領了耶路撒冷.
在一代不如一代的統治之下.
國家最終是步向滅亡.
其實神興起的外邦勢力.

$^{161}$是對猶大國的審判.
因為猶大國的子民離棄了.
神最初訂立的制治傳統.
而神在這個時代背景開始.
先呼召耶利米.
要起來做先知.
宣講神將要審判的訊息.
喚醒子民對神.
去到中心的敬拜.
這是當時的時代背景.
當然我們今天去看經卷的時候.
我們是看了整個故事.
因為敘事者是從回顧歷史的角度去講.
但對於當時蒙召的耶利米.
沒錯就是當時蒙召的耶利米.
其實它是看不到那麼遠的.
它也沒有辦法可以猜測到未來的景況.
那麼耶利米是看到甚麼呢?.
我們又有甚麼可以從中去學習?.
今天我整理了三點.
和大家一同去思考.
第一點.
從生命的限制到神拯救的應許.
耶利米書一開始第一節.
描寫耶利米出生的背景.
她是出生自祭司的家庭.
第一節講到她是希臘家的兒子.
希臘家是被人稱為.
「希臘家」.
希臘家的意思是.
希臘家是被人稱為.
「希臘家」.
希臘家的意思是.
希臘家是被人稱為.
「希臘家」.
希臘家的意思是.
希臘家是變雅敏地亞特城的祭司.
甚麼是祭司呢?.
根據《新命記》.
祭司設立的目的就是要守護以色列人.

$^{201}$在全國各地設立聖地.
他們主要的工作是.
舉行祭祀儀式.
指導人去學習摩西的律法.
叫人去遵守律法.
從而避免冒犯神.
而先知和祭司是不同的.
先知是神所揀選授權的發言人.
他們要向子民宣告神的審判將要來.
做先知的要求比做祭司更高.
除了要清楚祭司的職責外.
先知還要向整個國家講話.
去揭露百姓的罪.
宣告神審判的訊息.
目的就是要呼召子民去悔改.
回到神面前重新立約.
這些忠言逆耳.
通常都不是人想聽的訊息.
換言之.
神不是要呼召耶利米.
像祂的父親一樣去做祭司.
即使祂一生在祭司家庭裡成長.
祂的經驗,祂的訓練.
好像不足夠讓祂去承擔一個更大的工作.
第五節更加提及祂要做列國的先知.
我們剛才也看過歷史.
一直以來.
猶大國的祭祀方式.
一直以來都是受到外族的影響.
以至約西亞在位期間.
要帶領猶大國進行祭祀的改革.
也就是當時的祭司.
忽略了他們應有的責任.
融合了外族的祭祀儀式.
這個祭司的傳統變得腐敗.
所以當神呼召耶利米做先知的時候.
小不免會引起祂家人極大的不安.
甚至祂的家庭會希望耶利米保持沉默.
但是神對耶利米有祂自己的心意.
耶利米要承擔一個更加艱鉅的任務.

$^{241}$第一,祂要向列國宣講審判的信息.
第二,要帶領猶大國去悔改回到申命記那種祭祀的傳統.
第三,祂要挑戰這些沒有盡忠而又腐敗的祭司.
因為這些祭司沒有按祖宗的規定去事奉神.
他們是耶利米生命的一大羞辱.
由此可見,在耶利米眼中.
更大,更艱鉅的可能是祂先要去面對自己整個身份的背景.
當我們明白了耶利米的背景.
我們明白了祂的難處和痛苦.
但是神是有應許的.
神說我早就選擇了你.
第四節提到.
我尚未將你做在母腹中,就已認識你.
認識這個字是遠遠超過認識祂的名字.
是一種更深層次的認識.
當中包含了神知道耶利米的生命.
祂認識耶利米的想法.
祂的過去等等.
前段時間,政府頒發好市民的獎項.
獲獎的人不是那天才被選出來的.
而是祂過去為社會做了一些事情.
可能是破案,可能是拯救過某些人而被認識.
選好市民的籌委.
祂不單認識這些人所做的事.
祂更加認識他們背後堅持的價值.
而這些價值很可能是來自他們的成長背景.
甚至是他們父母自小培養出來的價值.
神不是亂選耶利米出來的.
當聖經說神認識耶利米.
從祂的母腹之中,神就已經認識耶利米.
是神親自創造耶利米.
而耶利米很可能知道.
祂被呼召作列國的先知是一大重任.
既要拯救國家走向另一個屬靈的光景.
有很多未知之數.
任務是艱鉅的.
但是耶利米又有沒有意識到.
其實神是藉著這個召命.
而去拯救祂自己.
走出祂自己能力的限制呢?.

$^{281}$第九節說:因為我與你同在,要拯救你.
這個「你」字原文是指向耶利米.
就是神對耶利米的應許.
要耶利米知道,祂要衝破自己的限制.
祂要去學習聽明神的話語.
因為我與你同在,要拯救你.
神的應許不是單單要與祂同在.
更是要拯救祂.
拯救祂不要只是看到生命的限制.
而是要從祂自己的枷鎖裡釋放出來.
這個拯救不是只對耶利米獨特的應許.
而是對我們每一個人的應許.
或許我們的過去有很多生命的限制.
我們會有一些埋怨.
如果怎樣怎樣就好了.
我們很想重塑過去.
但我們知道我們沒有辦法去改寫我們的歷史.
今天我們成為了基督徒.
就是要明白我們的生命是神所創造的.
是神親自去揀選我與你進入歷史的裡面.
就像耶利米一樣.
神帶領他衝出他的命限秩序.
因為神從無縛之中創造他,認識他.
而我們也一樣被神不斷更新創造.
這是歷克胡哲.
相信你也聽過他的故事.
他天生四肢不完整.
但他絕對有更大的理由或藉口去說我做不到.
但他沒有放棄他的人生.
他更加去到不同的國家演講.
去分享神在他生命奇妙的見證.
去鼓勵人擁抱生命.
在人的眼裡,他的身體不完整.
但在神的眼裡,他的生命力戰勝了生命的限制.
豈不是在顯示神的同在呢?.
神更加是用歷克的生命.
成為其他擁有相同缺陷的人的祝福.
右邊的圖就是他和另一位小朋友.
他有同樣的缺陷.
那一刻他們在擁抱.

$^{321}$第二點想和大家分享的是.
從自身的不足到神同在的應許.
在耶利米追尋自己生命的限制的時候.
他發現了什麼呢?.
第六節,主耶和華說:.
「看啊!這個語氣透露了耶利米當時感到自己不能勝任神選擇的工作.
又或者當神呼召他做列國先知的時候.
他有所嚇倒.
當時耶利米所看見的是他不能做到的.
第六節:我不知怎麼說.
然後他再解釋:因為他年輕.
他所看見的是自己經驗的不足.
我們又回到經文的開始.
看看當時神是如何和耶利米互動的.
第二節:耶和華的話臨到耶利米.
第三節:耶和華的話日常臨到耶利米.
連耶利米他自己也說.
第四節:耶利米說:耶和華的話臨到我.
聖經連續三節用了三次「耶和華的話臨到」.
但是為什麼當耶利米被神呼召做先知的時候.
他還要用藉口去回覆神呢?.
主耶和華看,我不知怎麼說,因為我年輕.
難道他自身的限制竟然高過神的話語一直臨在他身上?.
當然今天的經文是關於耶利米被呼召做先知.
在座的各位不是人人都被呼召去做聖職的工作.
但是我們可以想想,在今天,在我們當中.
我們會不會也像耶利米一樣.
明明聽到神的話,但是又聽不明白神的話語呢?.
又沒有行出神的心意呢?.
情況就好像明明是做了靈修.
明明是看了聖經的了.
但是為什麼沒有得著呢?.
又或者我們會不會給自己很多藉口.
例如,我不知道,我不懂,我太忙.
我年紀小,我年紀大,我身體不好等等.
當神呼召我們去開始一個任務的時候.
神不會給我們一個很明確的路線圖.
你要這樣這樣這樣走.
反而神是要求他要和我們同行同在.
而這個同行同在是先於我們的意識.

$^{361}$我們一起再看看經文.
所用的字眼是「母腹中成形之前」.
「未生之前」.
這些都是聖經裡面和創造有關的言語.
神昔日如何立大地的根基.
去定海水的界限.
光從何而來.
星星是如何穩定地在天際.
動物是如何懂得覓食.
飛鳥是如何懂得因著氣候的變化而向南向北飛.
這些都是神創造的工作.
神創造的工作是不需要問我們的意見.
因為我相信當神問我們的時候.
我們都未必懂得回答.
但是神這一切的創造都是盛載著他的心意.
可能你會說.
「那神為何要選擇我經歷某些困難?」.
我相信我們可以將這些問題帶到神面前.
但是我們要知道我們這樣問.
是因為我們對未來沒有把握.
我們不知道將來會怎樣.
處於一個不知的裡面.
我們變得害怕.
就好像在一個黑暗的環境裡面.
我們不知道怎樣向前走.
我記得有一次我去行夜山.
是我人生唯一一次夜媽媽去行馬鞍山的山峰.
我發現電筒的光只能夠指引我前面那一步.
但幸好我有一班跟我一起行山的人.
我要靠著他們在前前後後說話的聲音.
我才知道要怎樣去行.
行的時候其實我也不知道何時到頂.
或何時下山有多遠.
但其他人知道.
因為他們有經驗他們行過.
他們可以帶著我行.
而我也可以放心去行.
所以在不知的黑暗裡面.
有同行者會使我們不再害怕.
現在我們知道.

$^{401}$很可能當時耶利米很害怕.
因為他意識不到神的話語一直與他同在.
即使他意識到神的話語是臨在他身上.
他聽到但他有沒有聽明白呢?.
早兩個星期.
我在教會與弟兄姊妹預備新年的裝飾.
就是你螢幕看到的金魚.
起初我們都覺得很難接.
有些弟兄姊妹更表示他們只是想接某一個工序.
因為對他們來說.
單單處理一個工序是比較容易.
而且在人的眼裡.
我們做一個流水式的作業.
效率亦都比較高,比較快.
但當時我看到我們的傳道人曾姑娘.
她花時間教每一位姊妹完成接一條魚.
我看到的時候.
我覺得這個圖畫很美麗.
因為她讓每一個人都去克服很難接的想法.
最後每個人都至少接到一條魚.
有些弟兄姊妹更加是學了之後.
還將這些魚,這些利是封帶回家再繼續接.
我們看自己的時候.
經常都會覺得能力不足.
看不到神為我們預備的天使在旁邊幫助我們.
第三點我想和大家分享的是.
將生命的限制交給神.
經歷神的能力.
再看回第六節.
主耶和華看,我不知怎麼說.
因為我年輕.
這一句看似是耶利米的藉口.
但其實我們看回原文的語氣.
就會發現耶利米向神表達自己那種痛苦的感嘆.
意思是耶利米將他的痛苦帶到神的面前.
耶利米將他的難處帶到神的面前.
然後第七節.
耶和華回應耶利米的難處.
他這樣說.
不要說我年輕.

$^{441}$然後神解釋了有兩個原因.
第一,我差遣你到誰哪裡去.
你都要去.
第二,我吩咐你說什麼話.
你都要說.
神重新一次向耶利米肯定.
神主動去揀選差遣.
帶領耶利米去宣講神審判的話語.
神所要的不是耶利米的能力.
不是要他在祭司家庭學到的祭祀禮儀.
而是要耶利米用行動去回應神的召命.
神要用耶利米的雙腳去服從他的差遣.
當神差遣他去哪裡他都要去.
又要耶利米用口去宣講神的吩咐.
當神吩咐他說什麼他都要說.
在耶利米去的途中.
神應許會有帶領,有吩咐.
就好像每一步都有神在旁邊.
有什麼可以跟這份同行相比呢?.
在宣講的途中.
神更加要拯救他脫離自己的過去.
重新去塑造一個全新的耶利米.
這不是祭司耶利米.
而是造先知的耶利米.
既然神命令耶利米要走這條路.
神就要負責.
有什麼好懼怕呢?.
神不單只是透過解釋去讓耶利米明白.
反而神自己有進一步的行動.
第九節.
「於是耶和華伸手按著我的口對我說.
看啊,我已將我的話放在你口中」.
這是一種充滿意象性的表達.
說的是神有他自己獨特的心意.
神要創造一個全新的耶利米.
是不需要向耶利米諮詢的.
好不好?.
因為神有他的心意.
然後神就將他將來要成就的工作.
表明給耶利米知道.

$^{481}$第十節.
「我今日立你在列邦列國之上.
為要拔出拆毀毀壞傾覆.
又要建立栽植.
在一個國家將要被傾覆.
被建立之前.
神先讓耶利米從他自己的生命.
去經歷那種拔出再被栽種.
這一種再被創造的過程.
以致到他將來出來做先知的時候.
他明白神曾經這樣.
在他生命彰顯過這種能力」.
不知道我們又會不會像耶利米一樣.
懂得將自己的不足向神說出來呢?.
其實你不說.
神也是知道的.
但為什麼我們要說給神聽呢?.
從經文之中.
我找到有兩個應用.
第一點.
回望過去.
但不是停留在過去.
我們的生命總會有一些創傷.
有一些痛苦.
是沒有人明白.
只有自己才知道.
但神從來不是叫我們自己一個去面對.
而是要明白其實神是知道的.
既然神是知道.
為什麼他不幫我呢?.
神給每個人的生命功課都不同.
正如神創造我們一樣.
我們每個人的樣貌.
我們的指紋.
我們的人生經驗都是獨特的.
縱然生命是有不完美的地方.
但我們就要在生活裡的形形異異裡去學習.
在主裡找到那種釋放.
讓我們的生命學會如何擺脫我們獨特的枷鎖.
因此神的拯救.

$^{521}$不只是在決志那一刻.
而是在生活裡.
神不斷地去拯救我們.
所以我們可以向耶利米學習.
學習去回顧我們的過去.
當中是會有痛楚.
甚至連自己都接納不了自己的身份.
但我們可以讓我們的過去和今天的我們共存.
學習像耶利米那樣.
向神說出自己的痛楚.
讓自己有意識去聽到這些痛楚.
然後去求神此時此刻.
讓這些痛楚這些限制給神知道.
讓祂的道.
聖經的話語成為醫治.
傳到我們的身上.
然後去經歷祂的同在.
祂的拯救.
但記住我們的專注點不是過去.
引伸到教會或社會甚至世界的場景.
有沒有什麼是我們眼前見到的.
我們經歷過的.
仍然留在我們的心裡.
成為一些包袱.
一些障礙物去遮蓋神的同行同在呢?.
這些會不會是一些根深蒂固的傳統.
或固有的觀念.
這些傳統觀念都是很有價值的東西.
但我們可以問神.
這些有價值的東西會不會成為我們的限制呢?.
我們的藉口會不會是過去都是這樣的.
來限制了神再創造的工作呢?.
然而我們也學習.
讓神透過祂的話語.
去選擇怎樣更新教會.
甚至社會整個群體.
所以這一點我們可以向耶利米去學習.
回望過去.
但不是停留在過去.
第二點.

$^{561}$展望將來.
接受生命被更新.
我們的題目是亂世中的照命.
相信大家都很好奇.
我們各人的照命是什麼呢?.
剛才說過.
照命就是我們要以人生.
整個生命去回應.
耶利米蒙召做先知.
是因為神親自跟他說.
但神未必有跟我們每一個說.
我們的照命是什麼.
我們怎知道呢?.
我們回到耶利米的生命.
當神未呼召耶利米的時候.
耶利米就在祂的祭司家庭裡成長.
在祂的環境裡生活.
這就是祂從出生到蒙召前的照命.
中世紀改革家馬丁路德曾經說過.
信徒皆祭司.
所有的基督徒.
不單是擁有一個共同的洗禮.
共同的信仰和福音.
也是屬於屬靈的產業.
這句意思就是指出信徒相同的地方.
不同的是.
他繼續說.
基督徒之間唯一的分別就是職務的區分.
那是為什麼呢?.
例如在我們熟悉的教會場景裡.
平時我們沒有限聚令的時候.
我們會有主席唱詩講道.
負責音響.
負責祈禱祝福教主日學等等.
還有會眾.
各人都有自己的位置和職份.
但是如果主席唱詩講道.
不好好準備.
會眾不好好的安靜坐好.
整個崇拜就不能夠獻上最好.

$^{601}$在社會裡我們都有不同的職業.
而且職業是無分貴賤的.
大家都是社會上的一份子.
在工作上.
如果你是一個做麵包的師傅.
你就用心去做好每一塊麵包.
因為你所做的麵包.
是用來餵飽肚子餓的人.
在家裡.
如果你是一個照顧者的角色.
你就需要用心去照顧你要照顧的人.
甚或只是好好照顧你自己一個.
每個人的生命.
在不同的階段.
都有不同的照明.
每一個時刻.
神所要的是要我們好好地.
活好今天這一刻.
這樣看來.
亂世中的照明.
其實沒什麼要擔心.
憂慮或者懼怕.
反而如果我們願意走這條信仰的道路.
神的能力就能夠在我們這麼多的.
生命的限制裡變得完全.
我想起哥林多後書第十二章九節.
神說:我的能力在人的軟弱上顯得完全.
這不是說神的能力不足夠.
而是如果沒有了我們的軟弱.
反而他的能力不能夠顯得完全.
就是正正因為我們的不足.
神才能夠完全顯出他的能力.
我們常說神永恆不變.
沒錯.
神愛的那種本質永恆不變.
但是神是一個活神.
活的神.
他的同在是會有變化.
如果我們看看聖經裡的每一個故事.
他描寫的是一個活神.

$^{641}$他有感情.
他有慈愛.
他有公義.
他有審判.
也有憐憫.
神此時此刻是這樣和我們同在.
但是下一刻.
我們的生命可能會有另一些的轉變.
他的同在也會有所變化.
意思就是這個同在不斷被神更新.
他不會永遠審判我們.
也不會永遠只安慰我們.
他是一個活的神.
願意我們一同學習.
用這樣的眼光去看神.
這是需要慢慢從生活中領會出來.
最後我想挑戰一下大家.
現在這一刻的你.
會是需要和自己的過去互動.
還是學習如何和我們的活神互動呢?.
兩件事情是沒有衝突的.
可以共存的.
而我們又會不會用養神.
在下一刻更新和改變我們呢?.
就好像耶利米那樣.
神在下一刻就將他的話語放在他的口裡.
為耶利米帶來更新.
總結今天的講道.
講了三個我們可以一起學習的地方.
第一.
我們從明白生命的限制.
去體會到神拯救的應許.
第二.
從意識自己的不足.
到領會神同在的應許.
第三.
我們將各種限制大膽地交給神.
經歷神的同在和拯救.
讓我們一起祈禱.
親愛的神.

$^{681}$我們從耶利米的出生.
蒙召,帶領.
到我們去尋索過去的人生.
和思想你創造的大能.
你從來沒有離開過我們.
願你透過你的話.
時常臨在我們的生命.
求你讓我們不要只看到自己生命的限制.
幫助我們擺脫我們的枷鎖.
釋放我們的生命.
讓我們去學習用心去信靠你.
並以行動進行你的話語.
經歷你的同在和拯救.
讓我們更加能夠擁抱今天你在亂世之中.
交託給我們的照明.
教導我們用生命去回應.
祈禱奉我主名求.
阿們.
(下集再見).
\newpage



\section{創世記 2:4-25}
\label{sec:xeKLoIITcbQ}
\textbf{2021.02.07《唯獨你是不可取替》創世記 2\_4-25( 和修版)  曾敏姑娘}
\newline
\newline
連結: \href{https://youtube.com/watch?v=xeKLoIITcbQ}{\texttt{https://youtube.com/watch?v=xeKLoIITcbQ}} ~~~~ 語音日期: 2021-02-10
\newline
\newline
\hyperref[sec:tlLyIq5teaA]{\small{< < < PREV SERMON < < <}}
~
\hyperref[sec:index_chronic]{\small{[返順時目]}}
~
\hyperref[sec:index_scriptual]{\small{[返順卷目]}}
~
\hyperref[sec:nPhRUxT71P8]{\small{> > > NEXT SERMON > > >}}
\newline
\newline
創世記 2:4-25
\newline
\begin{longtable}{cl}
\hline
\hline
章節 & 經文 (和合本修訂版)\\
\hline
2:4 & \begin{tabularx}{0.7\textwidth}{X} 這就是天地創造的來歷。在耶和華神造地和天的時候， \end{tabularx} \\ \\ \relax
2:5 & \begin{tabularx}{0.7\textwidth}{X} 地上還沒有田野的草木，田間的菜蔬還沒有長出來，因為耶和華神還沒有降雨在地上，也沒有人耕種土地。 \end{tabularx} \\ \\ \relax
2:6 & \begin{tabularx}{0.7\textwidth}{X} 但是，有霧氣從地上騰，滋潤整個土地的表面。 \end{tabularx} \\ \\ \relax
2:7 & \begin{tabularx}{0.7\textwidth}{X} 耶和華神用地上的塵土造人，將生命之氣吹進他的鼻孔，這人就成了有靈的活人。 \end{tabularx} \\ \\ \relax
2:8 & \begin{tabularx}{0.7\textwidth}{X} 耶和華神在東方的伊甸栽了一個園子，把所造的人安置在那裡。 \end{tabularx} \\ \\ \relax
2:9 & \begin{tabularx}{0.7\textwidth}{X} 耶和華神使各樣的樹從土地裡長出來，可以悅人的眼目，好作食物。園子當中有生命樹和知善惡的樹。 \end{tabularx} \\ \\ \relax
2:10 & \begin{tabularx}{0.7\textwidth}{X} 有一條河從伊甸流出來，滋潤那園子，從那裡分成四個源頭： \end{tabularx} \\ \\ \relax
2:11 & \begin{tabularx}{0.7\textwidth}{X} 第一條名叫比遜，它環繞哈腓拉全地，在那裡有金子。 \end{tabularx} \\ \\ \relax
2:12 & \begin{tabularx}{0.7\textwidth}{X} 那地的金子很好，在那裡也有珍珠和紅瑪瑙。 \end{tabularx} \\ \\ \relax
2:13 & \begin{tabularx}{0.7\textwidth}{X} 第二條河名叫基訓，它環繞古實全地。 \end{tabularx} \\ \\ \relax
2:14 & \begin{tabularx}{0.7\textwidth}{X} 第三條河名叫底格里斯，它流到亞述的東邊。第四條河就是幼發拉底。 \end{tabularx} \\ \\ \relax
2:15 & \begin{tabularx}{0.7\textwidth}{X} 耶和華神把那人安置在伊甸園，讓他耕耘看管。 \end{tabularx} \\ \\ \relax
2:16 & \begin{tabularx}{0.7\textwidth}{X} 耶和華神吩咐那人說：「園中各樣樹上所出的，你可以隨意吃， \end{tabularx} \\ \\ \relax
2:17 & \begin{tabularx}{0.7\textwidth}{X} 只是知善惡的樹所出的，你不可吃，因為你吃它的日子必定死！」 \end{tabularx} \\ \\ \relax
2:18 & \begin{tabularx}{0.7\textwidth}{X} 耶和華神說：「那人單獨一個不好，我要為他造一個配偶幫助他。」 \end{tabularx} \\ \\ \relax
2:19 & \begin{tabularx}{0.7\textwidth}{X} 耶和華神用泥土造了野地各樣的走獸和天空各樣的飛鳥，都帶到那人面前，看他叫甚麼。那人怎樣叫各樣的動物，那就是牠的名字。 \end{tabularx} \\ \\ \relax
2:20 & \begin{tabularx}{0.7\textwidth}{X} 那人就給一切牲畜、天空的飛鳥和野地各樣的走獸都起了名。只是亞當沒有找到配偶幫助他。 \end{tabularx} \\ \\ \relax
2:21 & \begin{tabularx}{0.7\textwidth}{X} 耶和華神使他沉睡，他就睡了；於是取下他的一根肋骨，又在原處把肉合起來。 \end{tabularx} \\ \\ \relax
2:22 & \begin{tabularx}{0.7\textwidth}{X} 耶和華神就用那人身上所取的肋骨造了一個女人，帶她到那人面前。 \end{tabularx} \\ \\ \relax
2:23 & \begin{tabularx}{0.7\textwidth}{X} 那人說：「這是我骨中的骨，肉中的肉，可以稱她為女人，因為她是從男人身上取出來的。」 \end{tabularx} \\ \\ \relax
2:24 & \begin{tabularx}{0.7\textwidth}{X} 因此，人要離開父母，與妻子結合，二人成為一體。 \end{tabularx} \\ \\ \relax
2:25 & \begin{tabularx}{0.7\textwidth}{X} 當時夫妻二人赤身露體，並不覺得羞恥。 \end{tabularx} \\ \\
[1ex]
\hline
\hline
\end{longtable}
$^{1}$(拍攝).
弟兄姊妹早晨.
早晨.
讓我們今天仍然有生命氣息.
來到祂的面前去敬拜祂.
下個星期就是農曆新年了.
大家會不會很期待呢?.
雖然因為疫症的緣故.
我們少了機會去見親友.
但是我深信大家會很珍惜這個機會.
因為是一個很好的假期.
我們可以好好休息.
下個星期除了農曆新年.
也是情人節.
絕對是一個很好的機會.
可以向身邊的人表達愛意.
說到情人節.
不知道大家最喜歡聽到的情話是什麼呢?.
如果有人跟你說.
唯獨你是不可取替的.
你有什麼感受?.
是否開心?.
這句話是不是覺得很熟悉呢?.
很可惜.
大部分弟兄姊妹現在都是在.
看我們的網上直播.
不是在現場.
你在現場我可能會叫你舉手.
有沒有聽過這句話呢?.
一舉手就暴露了自己.
是不是跟我同年代呢?.
這句話當然就是.
香港歌星鄭秀文一首.
很出名的歌的歌名.
那時候這首歌是很紅的.
我也很喜歡這首歌.
唯獨你是不可取替.
如果有人跟你說.
你的感受如何.
如果這句話是出自你的情人的口.

$^{41}$我相信你也會很開心.
因為在她心目中.
你是有一個很獨特的位置.
如果是神跟你說這句話呢?.
是神跟你說.
你又會覺得怎樣呢?.
或許弟兄姊妹就會說.
神會跟我們說這句話嗎?.
會呀.
原來在創造宇宙的時候.
神已經向我們表達了這份心意.
唯獨你是呀.
就是你.
就是人.
唯獨你是不可取替的.
神是很愛很愛我們的.
很多人看聖經的時候.
總是要將神分開.
舊約的神是不同的.
跟新約的神.
新約的神是很慈愛的.
很愛我們的.
耶穌基督為我們犧牲.
覺得舊約的神總是令我們想起了別的東西.
但我想跟你說.
其實不是的.
神從頭到尾都沒有改變過.
由祂創造宇宙萬物的時候.
祂就是這樣.
祂很想跟我們表達.
祂很愛很愛我們.
那份愛是很永恆的.
很久了.
不會改變的.
我們從哪些地方可以看到.
祂對我們的那份愛呢?.
我們透過今天創世記第二章.
去看這件事.
現在我們再一次有個禱告.
天主上大人懇求你.

$^{81}$今天開我們的眼睛.
讓我們看到你奇妙的作為.
更深的明白你的愛.
願意聖靈在我們的心裡面.
叫我們的心.
向上帝你的說話去打開.
我們明白.
就去回應.
讓上帝你的心意得到滿足.
奉耶穌基督的名字祈禱.
阿們.
今天我們如果大家可以的話.
都想大家手上可以翻開聖經.
去看創世記第二章.
我們都會看當中一些經文.
在哪裡看到神很愛很愛我們呢?.
第一件事可以這樣說.
神做人的方法非常奇妙.
凸顯出人的尊貴.
神做人的方法奇妙.
凸顯人的尊貴.
當時這樣說.
耶和華神是用地上的塵土做人.
將生命之氣吹進他的鼻孔.
這人就成了有靈的活人.
耶和華神用地上的塵土做人.
這個做字是指神不是亂來一通的.
隨便隨便隨便地做了人出來.
而是他經過精心的構思.
他用心去想去設計.
才將人創造出來.
就好像一個藝術家一樣.
他在開始進行設計他的作品之前.
他就會一邊想一邊構思.
花了很多心思才開始動手.
他期望做出來的是他的傑作.
神也一樣.
他一邊想一邊設計才做.
他很希望人是他手中的傑作.
結果做出來的時候.

$^{121}$人又真的是傑作.
如果大家有機會聽很多科學家的分享.
很多不同研究報告的分享.
你會留意整個人體的構造是非常精細.
我深信直到今天.
無論人用了多少時間和精力去研究.
對於人體的奧秘.
我們的認識仍然是非常有限.
神是掌管著一切奧秘的神.
他才知道裡面的奧妙.
你就知道原來人是不簡單的.
《創世記》一章二十六節說.
神說:我們要照著我們的形象.
按照我們的樣式做人.
使他們管理海裡的魚,天空的鳥.
地上的生畜和全地.
以及地上爬的一切爬行動物.
神不單是想過,度過,深思熟慮才做我們.
神是按照他的形象,樣式做我們.
這個形象,樣式是指理性,道德意識.
有知識,有能力等等.
既然神是按照他的形象,樣式做我們.
即是說我們很像神.
就像子女,父母一樣.
例如神會愛人,憐憫人.
我們也會.
神有語言能力,我們也會.
神有創造力,我們也有.
特別是地面上不同國家,地方的建築.
你會覺得很驚訝.
這些是人造出來的?.
是的,是人造出來的,很厲害的.
只不過我們可以這樣說.
雖然我們也有這些能力.
我們也有不同的知識等等.
但是說到級數.
神的級數肯定是高過我們很多倍.
我們沒辦法跟神比.
但是我們是像他.
在萬物當中.

$^{161}$只有人是按照神的形象,樣式做的.
所以人在神眼中是很特別,很重要的.
比萬物都重要.
不單是這樣.
神做人的過程也是很特別的.
以章七節這樣說.
「而我話神用地上的塵土做人.
將生命氣息吹進他的鼻孔.
這人就成了有靈的活人」.
神用地上的塵土做人之後.
其實人是沒有生命氣息的.
只不過是塵土的製作品.
但是當神親自將生命氣息吹進人的鼻孔之後.
人就有生命氣息了.
人的生命是直接來自神.
今天你看到世界上的人都覺得自己很了不起.
是不是因為我們覺得世界的科技很發達.
只要你有錢,什麼都可以做.
人很厲害,很多發明創造.
但是不要忘記了.
其實所有東西都來自神.
包括我們的生命氣息.
今天神將這個生命氣息拿走了.
其實我們什麼都做不到.
我們絕對不是神.
但是我們留意今天的人卻是驕傲.
忘記了這口氣.
這生命氣息是神給我們的.
而且將生命氣息吹進人的鼻孔這個動作.
是令人感到很溫暖,很親切.
好像什麼呢?.
好像親吻,好像親吻這麼親密.
我們平時看的時候也沒有留意這一點.
這個動作很親密.
神將生命氣息吹進人的鼻孔.
提醒了我們沒有了神,我們是不行的.
我們和神要緊緊結連在一起.
而這個親密的創造方法只用在人身上.
上帝不會用這個方法去做動物.
做其他的生物.

$^{201}$你從這裡看到.
神和人的關係很親密.
人在神心目中的位置絕對超越了世上的萬物.
單單從這個創造的方法我們看到.
人在神心目中的地位是獨一無二的.
神愛我們.
正正是因為神很愛我們.
所以我們的生命是很有價值的.
在神眼中,我們是很寶貴的.
好像鑽石一樣寶貴.
我這一生人直到這一刻為止.
暫時還沒有機會去看什麼鑽石展覽.
我總覺得這些和我相隔很遙遠.
我覺得是那些通常和這個行業有關.
或者在經濟上比較富裕的人.
他們會去看這些展覽.
但是也看過一些圖片.
通常這些展覽裡面的場館.
都是佈置得非常漂亮.
整個地方是很高級的.
但是重點不是場館.
不是場地的佈置.
而是當中很小的.
就是那個鑽石,價值不非過億的鑽石.
所以進去參觀的人.
我相信他們看的時候.
就會看這個佈置多漂亮.
所有周圍的東西怎樣.
看就看,這一顆很小很小的.
就是整個場館的中心.
價值不非的,幾億的鑽石.
其實我們生命都是這樣的.
在這個世界有很多很多被造的物件.
很多的生物等等.
但是我們的價值就像這顆鑽石那麼厲害.
我們有價值不是因為我們住豪宅.
開名牌的車.
穿名牌的服裝.
我們工作很有成就.
學業很有成就.

$^{241}$不是因為今天我和你很有錢.
我們生得好像模特兒那麼美女.
那麼英俊.
更加不是因為今天我人生際遇.
是否順利.
不是這一切決定我們的價值.
我們寶貴就是因為上帝創造我們的時候.
賦予我們有這個寶貴的價值.
而這個價值.
甚至可以這樣說.
鑽石都不足以形容.
是無價寶.
如果不是都不會.
是要耶穌基督用他的生命來換取我和你的生命.
我們的價值.
在這麼寶貴的時候.
但是今天我們怎樣看自己的價值呢?.
我們的自信心建立在哪裡呢?.
是世人的眼光.
世人的評論.
還是神的看法.
我們的生命很多時候出現很多問題.
都是因為我們不認識自己的價值.
而變得很自卑.
我們不開心.
我們很喜歡和別人比較.
我在人群當中經常聽到這些說話.
很多比較總覺得自己不可愛.
他又好.
不用做那麼多事.
你看他的皮膚多滑.
他又好,很白.
他又好.
整個站出來,無敵的身材.
你看我整個水桶.
很多東西比較.
他好,讀很多書.
又有錢.
連三餐都是一個困難.
很多很多很多.

$^{281}$在這裡很不開心.
很自卑.
如果有什麼場合.
最好就躲起來.
不要出現.
你出現了就覺得自己矮人一截.
很多比較,很不開心.
我們不愛自己.
我們不愛自己,就不會愛別人.
我們很容易批評別人.
看別人.
看那些比自己不太好.
他這樣.
他這樣比不上我.
那樣比不上我.
如果交朋友,不要找這些.
我們不懂得去愛.
其實這個自我價值是很重要.
我不是說我們要很驕傲.
怎樣去稱讚自己.
而是我們懂得從上帝的角度去看自己.
原來我是很有價值的.
神給我這個價值.
要學會愛自己,才可以愛別人.
我們教會有干理樓事工.
我自己也很欣賞這個事工.
我們會提供物資給經濟上需要的朋友.
每一次.
當這些干理樓的朋友.
來到我們教會當中的時候.
我看到他們看到什麼呢?.
我看到上帝的愛.
傾倒在他們的生命裡.
我看到神.
很看重他們每一顆生命.
其實他們和我們一樣寶貴.
不會因為我們今天是幫助者.
就比受助者更加珍貴.
我們和他們一樣.
在神眼中都是那顆鑽石.

$^{321}$有時候我們去短孫.
真的很有趣.
我們去短孫.
我們帶很多物資去到宣教的地方.
我們要為那個地方.
或為那個機構當地人.
提供很多服事等等.
我們很容易站在一個高的地方.
看上去低的地方.
覺得對方有很多需要.
感覺上我們的價值好像高一點.
但其實對不起.
無論幫助者或受助者.
我們在神眼中的價值是一樣的.
我們不會比別人高.
因為神賦予人價值.
從來不是因為外在的東西.
是他創造我們的時候.
比我們是一樣的.
所以這個過程很盼望.
我們弟兄姊妹一起彼此鼓勵.
從神眼中看自己.
不自卑也不驕傲.
愛自己.
也愛人如己.
第一點剛才說了.
神做人的方法很奇妙.
突顯人的尊貴.
第二點.
神細心體貼.
為人預備一切.
當我再看的時候.
從另一個角度看這段經文的時候.
我就很驚訝.
神是這麼好.
預備是這麼豐富.
二章八節說.
也說神在東方的伊甸.
栽了一個原子.
把所造的人安置在那裡.

$^{361}$伊甸園其實是一個果園.
有很多水果.
很豐富的水果.
原來在古時代的中東地方.
一般人只種菜園.
有分別的.
菜園的級數不及果園.
一般人只種菜園.
誰可以種果園呢?.
只有皇帝可以種果園.
層次是高一點的.
當時我們可以這樣說.
伊甸園的主人是誰呢?.
就是神.
神是各位君王.
下次再看這段聖經.
感覺應該會有些不同.
原來不是一個普通的果園.
神是君王.
是伊甸園的主人.
他派人去管理這個園子.
人其實是神的僕人.
但是也很有榮耀.
很有榮幸.
可以這樣做神的僕人.
去服侍這位神.
人可以在園裡住.
絕對很開心.
因為裡面應有盲有.
很豐富.
如果以今天買樓的角度來說.
這是超級大豪宅.
不只是有地方睡覺這麼簡單.
有一個很大的園子.
裡面的水果應有盡有.
今天應該是很貴的.
價值不菲的.
昔日人是不用付錢的.
我們昔日不用買樓.
亞當·夏娃很幸福.

$^{401}$住在這麼貴的地方.
是皇帝住的地方.
他們住.
你說多幸福呢.
到今天.
人不能住伊甸園.
但是神仍然為我們預備居所.
無論我們住哪裡.
都值得感恩.
我回想.
神這麼多年為我預備居所.
我不知道有沒有在崇拜這裡提過.
我其實在大陸出生.
然後去到中一的時候.
我就來香港.
初期來香港住什麼地方呢.
我是住過豬場的鐵皮屋.
養豬的地方.
豬場裡面的鐵皮屋.
裡面會有毛毛蟲在那裡滾來滾去.
就是這樣的.
夏天非常熱.
冬天又冷.
但是我想和你們說.
我不感到難過.
有地方住.
我就開心.
我仍然很懷念那時候在豬場生活的日子.
和家人一起推著豬車.
去餵豬的生活.
我還是記得.
我在過程中心裡沒有什麼特別要求.
慢慢地.
神很愛我們一家.
由豬場被我們搬到另一條村住.
開始住村屋.
一邊住村屋.
一邊從村屋搬屋.
搬到住一些丁屋.
一邊從一個地方搬到另一個地方.

$^{441}$每一次搬遷.
神都是恩待我們.
讓我們住得越住越好.
我經常都很喜歡和弟兄姊妹說這個見證.
神為我們預備是多過我們所求所想.
每個地方我都欣賞.
我就算最近和弟兄姊妹聊天.
弟兄姊妹都會和我說.
以前條件差的時候.
試過住在哪裡.
差到這樣.
由那裡搬過來紅福村住.
多感恩 多開心.
我也很感恩.
我替她開心.
我在那裡見到神的作為.
神的供應.
是一個很好的福氣.
所以今天弟兄姊妹都可以提醒居住那裡.
神有沒有給你什麼福氣.
怎樣為你預備.
這一切不是必然的.
我想說.
神為人預備水源.
以章十節說.
有一條河從伊甸流出來.
滋潤那源子.
水是生命之源.
植物的生長需要水.
人和動物都需要水.
沒有水萬萬不能.
神就是這樣透過水.
去滋潤那個源子.
很重要.
其實伊甸源本身的名字的意思.
就是水源充足的地方.
今天我們在香港說水源充足.
我們的感受不是很深刻.
但是如果你回到中東的地方去說.
水源充足.

$^{481}$很滋潤.
感受就很深刻.
因為真的不是每個地方.
都有這麼多水源.
但是上帝為人造的這個伊甸源.
是水源很充足.
以致慢慢地.
植物可以生長起來.
然後預備水之後.
就預備食物.
當時人是吃素的.
神怎樣為人預備食物呢?.
就以章九節這樣說.
神使各樣的樹.
從土地裡長出來.
可以越人的眼目好作食物.
就是這樣有水.
滋潤土地.
樹就會長出來.
就有水果.
有蔬果.
那時候可以作為人的食物.
人是不需要擔心.
每天亞當夏娃都吃.
水果自製.
想吃什麼水果.
就有什麼水果.
你不要輕看水果.
今天買水果非常貴.
你買多幾種.
已經要花很多錢.
現在說的是每天水果自製.
所以可以想像.
他們多麼健康.
有水有食物.
神為人預備工作.
以章十五節.
神把打人安置在伊甸園.
讓他耕耘看管.
耕耘看管.

$^{521}$這個詞語有不同的解釋.
有照字面直接翻譯.
如果是這種解釋.
原來人不是無所事事.
在伊甸園生活.
可以管理那個地方.
讓每件事都井井有條.
工作是快樂的.
工作不是苦差.
神給人工作.
是給人在裡面.
有一個很大的滿足感.
其實今天.
不知道大家從事什麼工作.
大家享受工作嗎?.
還是覺得天天捱世界?.
一路由星期一捱到星期五.
很不容易捱到星期六中午.
然後苦難結束.
去到週末.
過完星期日又再來過.
到底是捱世界還是享受工作?.
在工作裡找到你的使命.
發揮你的長處.
但如果不是耕耘看管.
另一個層面.
敬拜服從.
這樣很有意思.
原來神想我們生命的工作.
是神把體作工作敬拜服從.
是我們生命不可缺少的一部分.
神很喜歡我們敬拜祂.
很喜歡我們的生命去順從祂.
今天我想問弟兄姊妹.
我們的敬拜生活又如何?.
是不是星期日回來一會兒?.
或者我們在網上直播.
一會兒.
就說一會兒是敬拜.
平時和神沒有關係.

$^{561}$工作是恆常地進行.
不是那星期過一天.
而是恆常地進行.
神想我們恆常地敬拜服從.
我們有多麼的恆常?.
神為我們預備水,食物,工作.
剛才再早的就是居住的地方.
還有一樣很重要的.
神為我們預備配偶.
為人建立了婚姻的制度.
二章十八到二十五節.
「哪人單獨一個不好.
我要為他做一個配偶幫助他」.
神看到一個人這樣生活.
很孤單.
人是一個社會性的動物.
社會性的意思是.
我們不可以一個人生活.
需要有群體.
於是神已經在想.
要為人做配偶.
但很特別.
你看經文.
經文描述不是馬上說.
神叫亞當沉睡.
在取下他的一根肋骨.
經文不是這樣.
中間很特別插入了一段經文.
「神用泥土造了野地.
各樣的走獸.
天空各樣的飛鳥.
都帶到那人面前.
看牠叫什麼」.
這麼有趣.
神覺得人很孤單.
轉頭就說.
把動物帶到牠那裡.
好像很不相關.
為什麼這樣說話呢?.
原來有意思.

$^{601}$神要讓人自己意識到.
他需要配偶.
於是他讓那些野地的走獸.
天空各樣的飛鳥.
都帶到那人面前.
第一.
神給予權柄為牠們起名.
因為人將要管理動物.
起名之餘.
亦給予人去看.
人會看到什麼?.
這麼有趣.
動物一對一對.
成雙成對.
在我面前經過.
牠們都有盤.
那我呢?.
這個或那個.
在當中會不會有.
可以和我做盤的呢?.
有沒有和我是同類的?.
於是動物一路經過.
神就讓牠自己去意識去發掘.
一個一個經過.
一對一對經過.
都沒有.
都沒有一個和我合適的.
沒有我的同類在.
所以經文才這樣說.
「只是亞當沒有找到配偶幫助牠」.
神讓人自己去看.
自己去發現.
在這個時候.
當人發現沒有.
神就為牠做配偶.
「神使牠沉睡.
牠就睡了.
於是取下牠的一根肋骨.
有在遠處把肉合起來」.
神種下蛙是很特別的.

$^{641}$是在亞當的身體裡.
拿出肋骨.
去做下蛙.
做了下蛙之後.
就讓亞當和下蛙可以一起成為夫婦.
你看看.
他這樣說.
「這是我骨中的骨.
肉中的肉」.
用骨來形容.
女人是代表什麼?.
有沒有想過?.
骨給你的感覺是怎樣的?.
很堅強的.
骨是很堅強的.
有這麼堅強的人在身邊幫助你.
大家一定可以得到幫助.
肉中的肉.
肉是想起什麼?.
親密.
又堅強又親密.
而這裡說的.
男人和女人的關係.
是非常緊密.
骨除了說堅強.
還在說什麼?.
那種衛生度.
不是那麼容易被破壞.
對關係的衛生度.
是很堅持的.
是很親密的.
所以婚姻的關係很寶貴.
如果將來有一次.
再探討神對婚姻的想法.
可以再一次探討這段經文.
但神在這裡.
透過這個過程.
讓亞當了解到.
自己需要配偶.
神為他做配偶.

$^{681}$這個配偶很重要.
和他關係很緊密.
這個配偶會毀身在這段感情.
這段婚姻關係裡.
還有很重要.
配偶幫助他什麼?.
幾方面.
第一 我深信耕耘的管理.
不只是亞當去做.
夏娃也可以一同去耕耘.
去看管.
第二.
夏娃可以和亞當一起.
在各方面彼此配搭.
在工作生活不同層面.
她也可以成為亞當的支持.
亞當不再是一個人.
有人與他同行.
第三.
很直接的.
神很想人生養眾多.
遍滿地面.
夏娃的出現.
可以和亞當一起.
去成就上帝的旨意.
一路生養眾多.
一路開枝散葉.
他們的後代就可以管理大地.
神真的這麼好.
為人預備一切他所需要.
住屋 水 食物 工作 配偶.
多少 是不是?.
所有這一切.
是神為人預備的.
神細心到.
我們真的想像不到.
神並沒有為世界其他受眾物.
預備這麼多東西.
所以你看到神很愛 很愛我們.
人在他心目中的位置.

$^{721}$是獨一無二.
來到這裡.
值得我們問自己.
我們有沒有為自己所擁有的.
去感恩呢?.
我最怕去到類似感恩分享的場合.
一些聚會.
我會問大家.
大英姐妹.
你們有沒有什麼要感恩的?.
沒有.
我想不到.
不知道要為什麼感恩.
其實說真的.
仔細想想.
很多東西要感恩.
我不及這個住得舒服.
也不及那個那麼成功.
誰又給我美男美女.
這個又順利.
又有錢.
你覺得我應該為什麼感恩呢?.
我又經常不舒服.
我覺得我好像沒有什麼要感恩.
你們說吧.
不是的 大英姐妹.
我們現在不是要比較.
我不是要做全世界第一.
所有東西超越所有人.
那個人才可以感恩.
我想說我們所有東西.
得到的都不是必然的.
是福氣.
所有東西都是神給我們的.
就是為了所有這一切都值得感恩.
不需要比較.
就是要感恩.
大英姐妹我們多久沒有感恩呢?.
其實我們教會.
真的很需要有這種氣氛.

$^{761}$一起回來說感恩.
我做干理樓事工的時候.
有一個姐妹.
那個居民.
那個街坊信耶穌.
她真的給我很多激勵.
我這樣祈禱神.
這樣幫我.
她為我預備這一切.
我也很得激勵.
她看到神很多恩典.
我心裡也很為她歡喜快樂.
但願我和你都是這樣.
生命啊 我們常常因為神的工作.
而感到快樂.
常常感恩.
最後一點.
就是神不嫌棄人的卑微.
甘願與人同在.
這是最寶貴的.
伊甸園不是一個果園.
是皇帝的果園.
伊甸園象徵了神與人同在的地方.
神特別做了這個地方.
將人安置在裡面.
神想與人一起生活.
《創世記》三章八節這樣說.
天起了涼風.
那人和他妻子聽見.
耶和華神在園中.
來回行走的聲音.
第三章是說人犯罪的紀錄.
罪是如何進入世界.
三章八節是說.
當人犯罪之後.
他們聽到耶和華神在園中.
來回行走的聲音.
其實這個來回行走的聲音.
不只是這一天才出現.
是每一天.

$^{801}$每一天都出現.
人都聽到.
每一天神都會來伊甸園.
與人一起拍拖.
每一天神都會找人.
神與人一起生活.
是很親密的.
是何時開始.
人才沒有與神緊密.
就是當人犯罪.
人被趕出伊甸園開始.
這也是一個很重要的.
象徵性的時刻.
人與神的關係拉遠了.
罪令我們與神分隔很遠.
然後你會發覺.
人開始轉變.
以前在伊甸園很自然.
人與神一起生活.
我覺得很親密.
神做人的時候.
就是希望與人這麼親密.
人都習慣了很親密.
但犯罪之後就不再親密.
然後人開始有自己的世界.
在自己的圈子裡生活.
人開始慢慢覺得.
上帝不是我生活中的一份子.
開始慢慢覺得.
世界是沒有神的.
我習慣我們自己生活.
不要經常和我說神.
神是分開的.
我們還我們自己.
這是令神很痛心的.
其實由起初的時候.
神不是這樣想.
神想和我們一起.
但不嫌棄我們很卑微.
雖然我剛才說人很特別.

$^{841}$很尊貴.
但不要忘記.
神是用塵土做我們.
不是用泥做我們.
是塵土.
比泥還低的物質.
我們其實很卑微.
神不想我們驕傲.
但在過程中.
神不嫌棄我們卑微.
我們很想和神一起.
但我們人是否很想和神一起呢?.
有神同在.
我們什麼都不用怕.
不用怕疫症.
去年社會運動的時候.
其實有很多人.
心靈被攪動了.
很多人失眠.
其實很多人失眠.
不是因為現在疫症開始.
是去年社會運動的時候開始.
很多人憤怒,沮喪,害怕.
不知怎麼辦.
一直延續到今年有疫症的時候.
失業了,經濟不好.
很多很多.
很害怕.
這個疫症.
我們怕隨時都會感染.
入屋前噴滿全身.
由頭噴到腳.
我們很害怕.
我們做安全設施是對的.
但你會聽到很多對話.
很多擔心和害怕.
我想說有神同在.
什麼都不用怕.
不用怕疫症.
不用怕失業.

$^{881}$不用怕健康出現問題.
不用擔心.
人際關係出問題怎麼辦呢?.
家人健康出問題怎麼辦呢?.
我想告訴你.
一切都會發生.
人生不如意事.
十常八九.
是很自然的.
有不同的不開心的事.
困難的事會在我們生命出現.
我們避免不了.
只可以這樣說.
有神同在,什麼都不怕.
神會和我們一同經歷.
今天我們沒有在一個.
昔日的醫電園生活.
但有神同在.
就可以像昔日一樣.
阿登烏夏娃在醫電園.
和神一同甜蜜.
親近.
是這麼好.
鄭秀文那首歌.
《唯獨你是不可取替》.
當中的歌詞是這樣說的.
如果今天將失去眼前的一切.
剩下清風兩袖也不計.
唯獨你是不可求人.
是我生命裡的一切.
剩下清風兩袖也不計.
唯獨你是不可求人.
是我生命裡的一切.
當我們愛一個人.
好心好心的時候.
我們什麼都可以失去.
唯獨這個人.
是不可以失去.
其實對於神來說.
祂最愛的就是我們.

$^{921}$祂不願意失去我們.
所以我再次強調.
祂才會這樣派祂的兒子.
耶穌基督.
為我們犧牲在十字架上.
因為祂真的不願意失去我們.
對於神來說.
我們是不可取替的.
祂愛我們.
祂很想與我們同行.
祂每一樣東西都和我們一起.
但是今天我和你.
是不是很想和祂同行呢?.
但是今天我和你.
是不是很想和祂同行呢?.
祂很想和我們同行.
祂很想和我們同行呢?.
對於我和你.
神是不是不可取替呢?.
我們的生命最貴重的.
最寶貴的是誰呢?.
是你的配偶.
是你的孩子.
你的事業.
你的功名.
是誰呢?.
是什麼呢?.
你敢不敢說.
如果今天將失去眼前的一切.
承低.
清風兩袖也不計.
你對著神說.
唯獨你一個.
是不可取替的.
你敢這樣說嗎?.
我深信神配得我們說這一句.
祂為我們預備一切.
祂將一切都給了我們.
我們會不會這樣愛神呢?.
我們一起祈禱.

$^{961}$今天藉著這段經文.
提醒我們.
你多麼愛我們.
你的愛.
是很耐的一份愛.
從過去到現在.
一直到永遠都不會改變.
天父你為了愛我們的緣故.
今天求你激勵我們.
為愛你的緣故.
願意什麼都捨棄.
這個愛真的不容易說出口.
盼望你給我們力量.
讓我們走出來.
以行動去愛你.
求你恩待我們教會中的弟兄姊妹.
奉耶穌基督的名字.
祈禱.
阿們.
\newpage



\section{約翰福音 6:1-15}
\label{sec:nPhRUxT71P8}
\textbf{2021.02.14《年年有餘》約翰福音 6\_1-15( 和修版)  老耀雄牧師}
\newline
\newline
連結: \href{https://youtube.com/watch?v=nPhRUxT71P8}{\texttt{https://youtube.com/watch?v=nPhRUxT71P8}} ~~~~ 語音日期: 2021-02-16
\newline
\newline
\hyperref[sec:xeKLoIITcbQ]{\small{< < < PREV SERMON < < <}}
~
\hyperref[sec:index_chronic]{\small{[返順時目]}}
~
\hyperref[sec:index_scriptual]{\small{[返順卷目]}}
~
\hyperref[sec:t7FE7L_V2U0]{\small{> > > NEXT SERMON > > >}}
\newline
\newline
約翰福音 6:1-15
\newline
\begin{longtable}{cl}
\hline
\hline
章節 & 經文 (和合本修訂版)\\
\hline
6:1 & \begin{tabularx}{0.7\textwidth}{X} 這些事以後，耶穌渡過加利利海，就是提比哩亞海。 \end{tabularx} \\ \\ \relax
6:2 & \begin{tabularx}{0.7\textwidth}{X} 有一大群人因為看見他在病人身上所行的神蹟，就跟隨他。 \end{tabularx} \\ \\ \relax
6:3 & \begin{tabularx}{0.7\textwidth}{X} 耶穌上了山，和門徒一同坐在那裡。 \end{tabularx} \\ \\ \relax
6:4 & \begin{tabularx}{0.7\textwidth}{X} 那時猶太人的逾越節近了。 \end{tabularx} \\ \\ \relax
6:5 & \begin{tabularx}{0.7\textwidth}{X} 耶穌舉目看見一大群人來，就對腓力說：「我們到哪裡去買餅給這些人吃呢？」 \end{tabularx} \\ \\ \relax
6:6 & \begin{tabularx}{0.7\textwidth}{X} 他說這話是要考驗腓力，他自己原知道要怎樣做。 \end{tabularx} \\ \\ \relax
6:7 & \begin{tabularx}{0.7\textwidth}{X} 腓力回答他：「就是兩百個銀幣的餅也不夠給他們每人吃一點點。」 \end{tabularx} \\ \\ \relax
6:8 & \begin{tabularx}{0.7\textwidth}{X} 有一個門徒，就是西門‧彼得的弟弟安得烈，對耶穌說： \end{tabularx} \\ \\ \relax
6:9 & \begin{tabularx}{0.7\textwidth}{X} 「這裡有一個孩子，帶著五個大麥餅和兩條魚，但是分給這麼多人還算甚麼呢？」 \end{tabularx} \\ \\ \relax
6:10 & \begin{tabularx}{0.7\textwidth}{X} 耶穌說：「你們叫大家坐下。」那地方的草多，人們就坐下，男人的數目約有五千。 \end{tabularx} \\ \\ \relax
6:11 & \begin{tabularx}{0.7\textwidth}{X} 耶穌拿起餅來，祝謝了，就分給坐著的人，也同樣分了魚，都照他們所要的來分。 \end{tabularx} \\ \\ \relax
6:12 & \begin{tabularx}{0.7\textwidth}{X} 他們吃飽後，耶穌對門徒說：「把剩下的碎屑收拾起來，免得糟蹋了。」 \end{tabularx} \\ \\ \relax
6:13 & \begin{tabularx}{0.7\textwidth}{X} 他們就把那五個大麥餅的碎屑，就是大家吃剩的，收拾起來，裝滿了十二個籃子。 \end{tabularx} \\ \\ \relax
6:14 & \begin{tabularx}{0.7\textwidth}{X} 人們看見耶穌所行的神蹟，就說：「這真是那要到世上來的先知！」 \end{tabularx} \\ \\ \relax
6:15 & \begin{tabularx}{0.7\textwidth}{X} 耶穌知道他們要來強迫他作王，就獨自又退到山上去了。 \end{tabularx} \\ \\
[1ex]
\hline
\hline
\end{longtable}
$^{1}$各位紅英家的弟兄姊妹平安.
新星頭先和大家拜個年.
祝大家新一年身體健康.
福輝滿日.
每日都能夠與主同行.
靈命日日高.
今日講題就是年年有餘.
在中國人的過年的時候.
寫揮春或者祝賀人的一句祝賀語.
寓意每一年都有一靜.
永遠都沒有短缺這樣的意思.
每一年都有一靜.
的確是一件很開心的事.
用會計學的計法.
就是收入比支出大.
有盈餘.
用信仰的講法.
就是恩典滿滿.
我相信每一個信徒都希望恩典滿滿.
那聖經有沒有一些類似年年有餘的事件.
可以給我們從中吸取教訓呢.
當然有.
在四福音裡面有一個神蹟.
講到耶穌為五餅二魚.
祝謝了之後.
不單止可以餵飽五千人.
那些剩下的.
還可以裝滿了十二個男子.
從表面上去理解.
都可以說和有餘.
這個意思有關係.
那整個神蹟事件.
能不能夠說.
讓我們能夠獲得這個恩典滿日的秘訣呢.
今天我們就從約翰福音六章一至十五節.
去探討幾個得福的秘訣.
我們一起去祈禱.
我們可以一起去敬拜你.
去尊崇你.
去聆聽你的話語.

$^{41}$這個真是無比的恩典.
因為你是那位賜福的源頭.
你是那位滋潤人心靈.
又激勵我們生命的主.
我們願意萬口都稱頌.
萬失都跪拜你.
求你臨到我們的心靈深處.
打開我們的心.
讓我們每一個都盛載著你的話語離開.
我們的祈禱.
奉靠主耶穌基督得勝明智祈求.
阿們.
在進入經文之前.
我們先了解一下經文的背景.
在耶穌出來事奉的三年多時間.
耶穌施行過很多神蹟.
這些神蹟都被不同福音書的作者記錄下來.
成為福音書的內容.
不過約翰福音的作者和其他福音書的作者.
在取材方面有很大的差異.
例如約翰福音.
祂只記載了耶穌七個神蹟.
婦女福音就多很多.
這是量的差異.
另外.
約翰福音記載的七個神蹟之中.
有六個是約翰獨有的.
而婦女福音是完全沒有記載的.
這是質的差異.
而最重要的分別.
是婦女福音以權能來詮釋耶穌所施行的神蹟.
即是神蹟的出現.
是彰顯了耶穌作為神的兒子的能力.
但是約翰就以記號來詮釋耶穌所行的神蹟.
或者用另一個說法.
約翰是用一個標示去解釋耶穌的神蹟.
大家記住一點.
稍後進入經文的時候.
我再和大家補充.
而今天說到這個神蹟的事件.

$^{81}$也是四福音書中唯一一個神蹟.
所有福音書都有記載.
所以很值得我們專注地去看一看.
今次我想和大家從三個角度去檢視這個神蹟的內容.
第一個角度就是.
當日跟從耶穌的群眾.
他們怎樣去理解耶穌這個神蹟呢?.
在第二節很清楚地透露了群眾當時的一些情況.
有一大群人因為看見他在病人身上所行的神蹟.
就跟隨他.
原來大部分的跟隨者.
都曾經見過耶穌在病人身上所施行的神蹟.
我相信他不應該見過一次耶穌施行的神蹟.
在上一章.
耶穌在畢氏大祠中治好一個瘸腿患者的神蹟.
他們可能當時已經跟隨著耶穌.
甚至可能再早一些.
在第四章.
當耶穌醫治大神的兒子的時候.
可能他們已經在那裡.
他們一直跟隨著耶穌.
可以這樣說.
他們應該是一群耶穌基督的跟隨者.
當他們經歷了五餅二魚的神蹟之後.
他們怎樣去形容耶穌呢?.
經文在第十四節這樣說.
人們看見耶穌所行的神蹟.
就說.
這個真是那個要到世上來的先知.
那個真是要到世上來的先知.
究竟是什麼意思?.
究竟群眾所指的先知是哪一個先知?.
我剛才說過.
約翰福音的作者.
是以記號或者標示.
來詮釋耶穌所行的神蹟.
記號或者標示的作用.
就是一見到這個記號.
立刻就令人有一個特別的聯想.
讓他立刻想起一件事.

$^{121}$一件事或者一個物件.
如果我和大家說.
我們一定要有獅子山下的精神.
那你立刻會想起什麼?.
香港精神.
香港人肯拼的精神.
獅子山下好像一個記號.
讓我們立刻有一個集體回憶.
記起一些我們很熟悉的事件.
或者記得一個很重要的概念.
五餅二魚的神蹟.
究竟是什麼記號?.
這個記號.
又會令到當時的群眾想起什麼?.
令到他們會說到這句話?.
原來對於整個猶太民族來說.
饑餓和食物的供應的事件.
是一個很特別的記號.
這個記號很自然地令他想起.
猶太人先祖在曠野的四十年的經歷.
因為他們的先祖在曠野的時候.
整個民族同樣面對一個.
怎樣可以有得吃的困難.
在《文造記》十一章十三節.
摩西說.
我在哪裡拿些肉去給這些百姓吃?.
如一肥力.
好像摩西一樣.
他這樣說.
我們去哪裡買些餅給這些人吃?.
群眾所遇見的困境.
好像重演昔日一個關於食物的難題.
耶穌所施行的神蹟.
又好像重演一件昔日.
天賜嗎哪天賜鵪鶉的事件一樣.
類似一樣的困難.
類似同樣解決的方法.
猶太人自然就記得.
《申命記》十八章十五至十九節的一段經文.
耶和華你的神要從你弟兄中.

$^{161}$給你興起一位先知像我.
你們要聽他.
誰不聽從他奉我名所說的話.
我必親自向他追討.
在以前.
耶和華應許以色列人的先祖.
會興起一位先知.
好像摩西一樣.
帶領著以色列人.
要求以色列人要聽這位先知的話.
因為這位先知就代表了耶和華.
這也是舊約所說的彌賽亞應許.
因此這批領受了五餅二魚的群眾.
從這個食物供應的記號.
他即時會聯想到.
原來耶穌就是昔日耶和華.
所應許的會來的先知.
耶穌就像昔日摩西一樣.
帶領著他們離開這個困境.
所以群眾的回應就說.
那真是要到世上來的先知.
一個信仰的記號.
讓這批群眾記得昔日.
對他們整個民族的應許和恩典.
事實上.
約翰記載這個神跡的事件.
也是想讓當代的讀者.
能夠領受到.
耶穌不單止是昔日耶和華.
所應許會來的彌賽亞.
耶穌還是神的兒子.
三位一體的第二位.
聖子耶穌基督.
紅衣家的弟兄姊妹.
我們不是猶太人.
我們沒有曠野的經歷.
我們自然認不出這個信仰的記號.
不過我們每一個人的得救自身經歷.
都是獨特的.
這個獨特的經歷.

$^{201}$能不能成為你一個屬靈的記號.
讓我們可以有一個回憶.
記得主耶穌對我們的供應.
和對我們的救贖呢.
或者我們還記不記得.
當日我們受浸加入教會的時候.
我們曾經在會友面前宣認過.
耶穌是我們的生命的主.
耶穌是我們請求的主.
我們在受浸的時候.
立志願意放下主權.
讓耶穌管理我們的生命.
過一個以主為先.
以主為尊的信仰生命.
這個和神立約的時刻.
能不能成為我們每一個會眾的信仰記號.
讓我們一起參與其他人的浸禮的時候.
我們可以有一個反思.
今天我和神的關係.
究竟是一個什麼狀況呢.
是原地踏寶.
還是一個熟齡的嬰兒.
還是我們已經長大成人.
按照神的心意.
去過我們的信仰生活呢.
第二個進入神蹟的事件.
就是教我們如何去承傳神的恩典.
經文裡面有一個門徒.
他叫做肥力.
耶穌在眾人都飢餓的處境之下.
去試煉他.
叫肥力解決這個食物的問題.
其實耶穌當時可以有很多人去問.
為什麼偏偏要選擇肥力呢.
如果我們有看其他福音書.
我們就知道這個地方有名叫白菜大.
肥力正正在白菜大出生.
也就是說.
這裡是肥力的出生地.
是肥力的主場.

$^{241}$他最熟悉的地方.
耶穌問一個出生在白菜大的肥力.
問他我們去哪裡買餅給這些人吃呢.
肥力很直接回答.
不夠.
就算是二十兩銀子買到的餅.
都是不夠.
有色經家計算過.
二十兩銀子的價值.
是當時一個工人七個月的工資.
事實上沒有人會帶七個月的工資出門.
而且肥力也沒有回答耶穌的問題.
耶穌是問.
去哪裡買.
肥力說不夠錢.
買到都不夠.
可以這麼說.
肥力答非所問.
面對耶穌對肥力的挑戰.
他完全應付不到.
耶穌另一個門徒安德烈.
面對同一個問題.
他的反應很不同.
他說我這裡有.
不過很少.
不夠分的.
嚴格來說.
安德烈也沒有直接回答耶穌的問題.
不過安德烈有行動.
他帶了一個只有五個餅兩條魚的小孩.
將他們所有的一切都交給耶穌.
少少的五個餅兩條魚.
放在耶穌手裡.
經過竹蔗.
分出去.
就可以餵飽五千人.
《馬可福音》記載.
只是男人就有五千人.
所以有些人就估計到.
如果當時加上一些婦孺和小朋友.

$^{281}$估計當時全部的人數.
應該都超過一萬人.
大家可以進入當時的場景去想一想.
整個分遞過程.
究竟是怎樣進行的呢?.
你記住當時超過一萬人.
坐在草地.
一條很長很長的人鏈.
是不是很虛陷?.
耶穌是不是先將五個餅兩條魚.
變成一萬份的數量?.
然後開始傳.
如果是這樣的話.
越接近耶穌的人.
就會先接觸了那些食物.
不過他不是自己拿的.
他要分出去.
他要將拿到的東西.
分給後面的人.
所以越接近耶穌的人.
他是最後才能夠拿到食物的人.
而距離耶穌最遠的人.
是首先分到食物.
是不是這樣呢?.
耶穌先將五餅二魚.
變成一萬份的數量.
然後逐個分.
又或者會不會是這樣呢?.
每個人拿起自己的那份.
然後再分出去.
分的過程.
神就說.
分的過程.
神就發生了.
不是我全拿走.
不是我拿盡.
下一個沒有份.
我拿走.
留下我的.
分出去.

$^{321}$如果是這樣的話.
最接近耶穌的人.
最先拿到食物.
食物經過分到.
再遞出去.
這個簡單的動作.
讓神蹟不斷地發生.
你想一想.
每個人拿到自己的那份.
再分出去.
每個人拿到的那份.
又再分出去.
神蹟就是這樣不斷地發生.
這個神蹟有別於之前第一個.
第一個是.
耶穌先將五餅二魚.
變成一晚份.
然後分出去.
現在這個是.
一路遞.
神蹟不斷地產生.
經文沒有告訴你.
這個神蹟是怎樣發生的.
用什麼方式去發生.
但有一點可以很肯定的.
整個神蹟的事件.
除了耶穌有介入之外.
人都有回應和行動.
馬太福音這樣形容這件事.
耶穌出來.
見到一大班人.
就憐憫他們.
五個餅 兩條魚.
要滿足一萬個人的需要.
的確是杯水車薪.
但只要交在主耶穌的手裡.
主會使用.
而主所成就的.
超過我們所想所求.
超過一萬人.

$^{361}$在這個分遞的過程裡面.
享受著神對人的憐憫.
以及人與人之間的關愛.
經文給我們看到的.
不單止這一萬人可以吃得飽.
還有得剩.
剩下多少.
裝滿了十二個男子.
我們知道十二這個數字.
對猶太人來說.
是一個圓滿的數字.
耶穌的供應.
可以豐豐足足.
令到每一個人參與的.
都感覺到很豐盛.
絕對切合今天的講題.
是甚麼.
年年有餘.
耶穌基督所賜予給我們的恩典.
不單單是食物.
可以是我們對人的禱告.
可以是我們對人的探訪.
可以是我們對人的接納.
可以是我們對人的慰問和安慰.
無論是甚麼形式都好.
只要願意向別人付出.
哪怕只是一個微笑.
都可以成為別人的祝福.
我預備這篇講章的時候.
就想起我們上個月.
教會推動了一次.
在洪福村派福袋的行動.
雖然有限聚令.
但無礙弟兄姊妹的參與.
由派發之前的福袋包裝.
到當日上午和下午的派發.
弟兄姊妹都全程參與.
我們透過與人分享.
實踐信仰的精神.
福袋被神祝福了之後.

$^{401}$送到街坊的手裡.
神透過這個行動去工作.
讓街坊感受到甚麼叫愛.
神的名被高舉.
神的愛被傳遞.
洪恩家的弟兄姊妹.
成為傳遞愛的工具.
好像五餅二魚的群眾一樣.
群眾得到食物之後.
又將食物分享給別人.
洪恩家的弟兄姊妹.
得到神的恩典之後.
又將這個恩典傳遞給別人.
各位弟兄姊妹.
我們已經是神家裡的人.
我們享有豐盛的恩典.
我們願不願意將我們手上.
看似微小的東西.
交在主的手裡.
無論是多麼不起眼.
多麼微不足道.
但經過被主耶穌基督的祝福.
再分遞出去.
同樣可以祝福到很多人.
你對其他人的分享.
就好像今天的講題一樣.
可以年年有餘.
甚至我們每一個的會友.
我們都能夠成為.
耶穌基督這個分遞神蹟的其中一環.
透過我們雙手.
透過我們的禱告.
透過我們對人的關心.
透過我們對人的問候.
不斷地收到從神而來的恩典.
又不斷地將這個恩典傳開去.
讓我們同樣地能夠實踐到.
我們今天講到的題目.
年年有餘的信仰版.
「恩典滿人」.

$^{441}$第三個角度進入神蹟事件.
就是怎樣才是真正的.
耶穌基督的跟隨者呢?.
甚麼才是真正的.
耶穌基督的跟隨者呢?.
經文第十五節講到.
耶穌知道他們要來強迫他作王.
就獨自由退到山上去了.
群眾既然確認.
耶穌就是那個彌賽亞.
是以色列民族的拯救者.
他們想擁納耶穌做他們的王.
做他們的領袖.
但耶穌就避開群眾.
獨自走上山.
為何耶穌要避開眾人呢?.
為何那麼多人跟隨.
不是好的嗎?.
原來耶穌知道他們擁納他.
做王的目的是甚麼.
當群眾稍後再次找到耶穌的時候.
耶穌在第六章二十六節.
這樣跟群眾說.
他說:我實實在在告訴你們.
你們找我並不是因為見到神蹟.
而是因為你們吃餅吃飽了.
不要為了吃餅而去找我.
不要為那些會壞的食物操勞.
而是要為那些傳到永生的食物操勞.
耶穌說:你們跟從我.
只是因為我能夠供應食物給你們.
滿足到你們的肚腹之慾.
滿足到你們的需要.
你們之所以跟從我.
都是因為你自己.
你們不明白.
我施行神蹟背後.
的含義.
你們不明白.
神的心意.

$^{481}$單單只是關注你們自己的需要.
原來一個真正的.
耶穌基督的跟隨者.
不應該只關注自己的.
需要和缺乏.
耶穌說:不要為今生的食物操勞.
不要為今生的食物操勞.
而是要為永生的食物操勞.
今生的食物表示了什麼?.
我們是依食住行.
永生的食物表示了什麼?.
神的話語.
以及神在我們生命裡的帶領.
而所謂《馬太福音》六章的.
其中一段說話.
最能解釋這個道理和教訓.
祂說:所以我告訴你們.
不要為你為生命憂慮.
吃什麼,喝什麼.
或者為你的生命.
身體憂慮著什麼.
祂說:生命不勝於飲食嗎?.
身體不勝於衣裳嗎?.
你們要先求神的覺和義.
這些東西都要加給你們了.
所以不要為明天憂慮.
因為明天自有明天的憂慮.
一天的難處一天當就夠了.
經文說:不要憂慮.
我們吃什麼,喝什麼,穿什麼.
雖然依食住行很重要.
但是一個真正的.
耶穌基督的跟隨者.
不應該只是單單關注.
在物質上的依食住行.
經文說:你們要先求.
祂的覺和祂的義.
這些東西都要加給你們了.
一個主耶穌基督的跟隨者.
應該有上述的心態.

$^{521}$就是要追求祂的覺和祂的義.
什麼叫神的覺?.
耶穌說:天國近了.
你們當悔改.
一個改變了的世界.
一個改變了的教會.
一個改變了願意悔改的生命.
這個就是天國.
為了實踐天國的來臨.
我們要勇於向每一個人傳福音.
那個是悔改的福音.
你願不願意為天國的緣故.
向人宣講一個悔改的福音.
神的義是什麼呢?.
先知尼加向世人宣告.
只要你行公義,好憐憫.
全謙卑的心與你的神同行.
一個全憐憫,謙卑的心.
在你的生活裡面去實踐真理.
這個就是神的義.
你願不願意為天國的緣故.
在你的生活裡面實踐真理?.
對人多一份憐憫.
多一份公義的心.
而且追求一種與神同行的屬靈生命.
一個耶穌基督的跟隨者.
他應該以耶穌的心懷意念.
成為他的信仰使命.
保羅在《迦太書》說.
現在活著的不再是我.
乃是基督在我裡面活著.
並且我如今在肉身活著.
是因信神的兒子而活.
他是愛我,為我捨己.
因著耶穌對我們的愛.
我們甘心樂意地去跟隨耶穌.
因著耶穌基督對我們的捨己.
我們能不能夠甘心樂意地為耶穌而活呢?.
各位同性戀家庭姐妹.
在2021年的農曆新年.

$^{561}$我們能不能夠藉著今天的訊息.
為我們找到一個信仰的記號.
叫我們一見到這個記號.
我們立刻可以數算神的恩典.
我們能不能夠在神的恩典裡面.
成為這個恩典分隸的一部份.
在信仰的生命過程裡面.
不要只顧著自己.
不要只記著.
只是讓恩典在自己的生命裡面收走而沒有出.
只有入而沒有出.
而是有入又有出.
好像今天的講題一樣.
恩典在我們生命裡面.
透過我們的生命分享.
祝福更多人.
讓更多人感受到.
年年有餘,恩典滿日的人生.
最後,當我們宣認.
告訴別人.
我們是一個基督徒的時候.
我們願不願意改變我們的價值觀.
不再以自己的需要為優先.
不再單單追求那種衣食住行的需要.
而是將生命的優先次序.
改為學務神的話語.
追求一種與神同行的人生價值觀.
成為一個真真正正的.
耶穌基督的追隨者.
我們一起去祈禱.
主,我們的恩典從來都沒有斷絕.
你的恩典滿滿.
年年有餘.
主,我求你讓我們每一個的生命.
成為恩典傳遞的一部分.
好像一個流通的管子.
有恩典入.
也有恩典流出.
讓我們的生命能夠祝福到身邊的人.
讓你的榮美讓人感受得到.

$^{601}$聖靈求你幫助我們每一個.
因為我們真的很軟弱.
我們面對世界價值觀的衝擊.
我們很容易跌倒.
求你堅固我們的心.
讓我們能夠實踐你的話語.
好叫我們的行為.
是和我們的恩典相稱.
保守我們的心.
叫我們能夠以你為先.
以你為主.
甚至我們每一個屬你的.
都能夠被你所使用.
成為光.
成為炎.
讓教會成為一個燈台.
照亮這個地區.
奉主耶穌基督的聖名祈求.
阿們.
願主祝福大家.
\newpage



\section{腓立比書 2:12-18}
\label{sec:t7FE7L_V2U0}
\textbf{2021.02.21《天國的拓展》腓立比書 2\_12-18( 和修版)  黎光偉牧師}
\newline
\newline
連結: \href{https://youtube.com/watch?v=t7FE7L-V2U0}{\texttt{https://youtube.com/watch?v=t7FE7L-V2U0}} ~~~~ 語音日期: 2021-02-22
\newline
\newline
\hyperref[sec:nPhRUxT71P8]{\small{< < < PREV SERMON < < <}}
~
\hyperref[sec:index_chronic]{\small{[返順時目]}}
~
\hyperref[sec:index_scriptual]{\small{[返順卷目]}}
~
\hyperref[sec:PDRx0v0AuiQ]{\small{> > > NEXT SERMON > > >}}
\newline
\newline
腓立比書 2:12-18
\newline
\begin{longtable}{cl}
\hline
\hline
章節 & 經文 (和合本修訂版)\\
\hline
2:12 & \begin{tabularx}{0.7\textwidth}{X} 我親愛的，這樣看來，你們向來是順服的，不但我在你們那裡，就是我現在不在你們那裡的時候更是順服的，就當恐懼戰兢完成你們自己得救的事； \end{tabularx} \\ \\ \relax
2:13 & \begin{tabularx}{0.7\textwidth}{X} 因為是神在你們心裡運行，使你們又立志又實行，為要成就他的美意。 \end{tabularx} \\ \\ \relax
2:14 & \begin{tabularx}{0.7\textwidth}{X} 你們無論做甚麼事，都不要發怨言起爭論， \end{tabularx} \\ \\ \relax
2:15 & \begin{tabularx}{0.7\textwidth}{X} 好使你們無可指責，誠實無偽，在這彎曲悖謬的世代作神無瑕疵的兒女。你們在這世代中要像明光照耀， \end{tabularx} \\ \\ \relax
2:16 & \begin{tabularx}{0.7\textwidth}{X} 將生命的道顯明出來，使我在基督的日子得以誇耀我沒有白跑，也沒有徒勞。 \end{tabularx} \\ \\ \relax
2:17 & \begin{tabularx}{0.7\textwidth}{X} 我以你們的信心為供獻的祭物，我若被澆獻在其上也是喜樂，並且與你們眾人一同喜樂。 \end{tabularx} \\ \\ \relax
2:18 & \begin{tabularx}{0.7\textwidth}{X} 你們也要照樣喜樂，並且與我一同喜樂。 \end{tabularx} \\ \\
[1ex]
\hline
\hline
\end{longtable}
$^{1}$各位紅恩堂的弟兄姊妹.
新年祝大家.
新一年.
蒙上帝的福.
願上帝的恩惠.
平安與大家常在.
今日的講題是天國的拓展.
「拓」在中文字裡面有「推開」的意思.
好像甚麼呢?.
好像一團麵粉.
我們用木棍將它慢慢從四方八面推開一樣.
過程中你需要用力.
你需要重複地去滾動木棍.
最終使麵團能夠均勻地慢慢地展開.
於是麵團的範圍就會不斷地擴張.
同樣教會是天國的實踐.
我相信是這樣.
所以如果我們期待教會得到拓展的話.
我們同樣需要用力在各個領域上.
將呼音的種子推展得更遠.
我相信在座每一位弟兄姊妹.
你和我一樣.
我們作為天國的國民.
即使在現在艱難的環境下.
這個我相信仍然是我們的願望.
天國的拓展仍然是我們的宏願.
我要問下去.
究竟天國怎樣才算是被拓展了呢?.
史奴保羅是這樣認為的.
他說教會的拓展關鍵.
不在乎它發展的規模有多大.
它的經濟實力有多強.
乃在於教會的群體能否保持它獨有的影響力.
他在《爾忽所書》裡是這樣說的.
保羅說:從前你們是暗媒的.
但如今在主裡面是光明的.
幸是為人要像光明的子女.
意思是說在教會的人群裡面.
我們理應發揮它獨有的影響力.
好像在黑暗裡面發出點點光芒一樣.

$^{41}$讓人認識那個真光耶穌.
教會在什麼時候能夠發揮這種影響力.
就說明天國正正在那個地方裡面被拓展開去.
今天很希望透過《肥納比書》第二章.
跟大家去討論,去思考.
教會怎樣可以發揮這種影響力.
請聽我讀出以下經文.
《肥納比書》第二章十二到十八節.
聖經是這樣說的.
「我親愛的,且讓看來,你們向來是信服的.
不但我在你們那裡.
就是我現在不在你們那裡的時候更是信服的.
就當恐懼戰驚完成你們得救的事.
因為上帝在你們心裡運行.
使你們又立志又實行.
為要成就他的美意.
你們無論做什麼事.
都不要發怨言,起爭論.
好使你們無可指責,誠實無偽.
在這彎曲,缽縫的世代.
作上帝無瑕疵的兒女.
你們在這世代中好像明光照耀.
像生命的道表明出來.
使我在基督的日子.
得以誇耀我沒有白跑也沒有徒勞.
我以你們的信心為貢獻的祭物.
我若被囂墊在旗上也是喜樂.
並且與你們一同喜樂.
你們也要照樣喜樂.
並且與我一同喜樂.
讓我們同心禱告」.
上帝我們懇求你幫助我們.
在新一年的裡面.
讓我們恩著依靠上帝你的說話而得福.
讓我們今天所聽的道.
所解明的聖經.
不單止能夠相信.
更加讓我們有足夠的勇氣和信心實踐出來.
祈禱奉耶穌基督的名字.
阿們.

$^{81}$經文剛才我們所讀的一開始.
保羅說「這樣看來」.
說明保羅在前面已經說了一番話.
原來之前保羅是勸勉肥立給信徒.
讓我們一起看看.
他是這樣說的.
他說「你們信徒之間.
要意志相同.
愛心相同.
有一致的心思.
一致的想法.
要心存謙卑.
各人看別人比自己強.
各人不要單顧自己的事.
也要顧別人的事」.
保羅原來一開始.
先向肥立給的信徒說明.
教會合一.
信徒之間彼此信服是很重要的.
接著下一步.
才向他們呼籲三個實踐的行動.
這個實踐的行動.
是和我們今天講的主題很有關係的.
天國的拓展.
必須要先實踐以下這三個行動.
讓我們一一來看.
第一個行動是什麼呢?.
是合一個工作.
我們務必要做下去.
保羅呼籲的第一個行動.
他是這樣說的.
十二節裡面說.
他呼籲信徒.
你們要恐懼戰驚的.
完成你們得救的事.
什麼是完成得救的事呢?.
是不是肥立給信徒信得不夠深.
所以要投入一點.
我們才得救呢?.
我相信不是.

$^{121}$根據《上文下理》.
保羅所指得救的事.
是指教會合一的事.
保羅指出.
要教會得到拓展.
發揮它的影響力.
合一是關鍵.
而且還要繼續做下去.
所以因為這個緣故.
保羅對教會的合一.
有兩個提醒.
第一個提醒是什麼呢?.
是合一個工作.
是每一個人都需要經營的.
你們要完成你們自己得救的事.
每一個人都要經營這個合一的工作.
合一不等於沒有意見.
教會的合一不是.
你出小句聲當幫忙.
教會的合一也不是.
你喜歡怎樣做就怎樣做.
我沒有意見.
而是每一個人都要付出.
正如保羅所說.
他說信徒是需要有一致的心思.
一致的想法.
要尋求一致就要溝通.
要彼此同工.
彼此配搭.
我們才會有共識.
不單單是這樣.
保羅繼續說.
「各人不要單顧自己的事.
也要顧別人的事」.
信徒不是被動地去等待餵養.
而是主動地去顧教會的需要.
爭取事奉的機會.
第二個提醒是什麼呢?.
保羅繼續去呼籲門徒.
他說.

$^{161}$你們恐懼戰經.
完成你們得救的事.
這個恐懼戰經.
並不是指個人.
不是針對你要怕某一個人.
你要對某一個人戰經.
而是對上帝一種敬畏的態度.
意思是說.
我們要帶著敬畏的心.
認真去完成這個合一的工作.
為什麼要這麼小心謹慎呢?.
現實告訴我們.
因為在教會裡面.
信徒要認真去對待事奉的同時.
又要和其他人同心.
其實是很艱難的.
不知道你有沒有這個經驗.
很多夫婦.
很多時候他們爭拗.
他們有爭拗.
其中一個原因.
是基於他們對子女的教導方法不一樣.
不知道這是不是你的經驗.
我相信你都明白.
這不是因為作為父母.
他很不愛這個家.
他很有自己的風格.
他喜歡自己怎樣做就怎樣做.
不是這個原因.
我相信相反的是.
他們都重視這個家.
他們都重視他們的子女.
所以他們很希望用他們覺得最好的方法.
去對待他們的子女.
去教導他們.
結果平時夫妻之間.
其實沒有什麼爭拗.
什麼都無所謂.
大家都很同心.
不過因為緊張子女的緣故.

$^{201}$就會發生爭拗.
所以保羅並不是要去否認.
這些信徒他們對教會的愛心.
他們都相信.
這些信徒他們很愛教會.
他們好像我們剛才所說的父母一樣.
他們很緊張教會的發展.
不過他們也分享了他們自己的秘訣.
他說唯有帶著敬畏的心.
大家才可以暫時放下彼此的歧見.
去完成教會合一的工作.
因為保羅相信.
當信徒願意靠著上帝尋求合一的時候.
我們就不需要擔心了.
因為上帝會在我們當中工作的.
祂會促使我們合一的.
保羅怎麼說呢.
他說「當恐懼戰驚完成你們得救的事.
因為是上帝在你們心裡運行.
使你們又立志又實行.
為要成就祂的美意」.
因為教會合一是上帝的美意.
上帝期望教會合一.
弟兄姊妹.
我們回想一下.
在教會二千多年的歷史裡面.
合一的工作.
其實從來都是這麼重要.
我們看史陶保羅.
雖然他用整個人生.
很努力的去做傳福音的工作.
但在他的書信裡面.
你不難發現.
絕大部分談論的.
都是關乎到教會合一的問題.
寫一箱.
要教會保持合一.
一點都不容易.
因為每一個人對教會都有不一樣的期望.
所以才讓保羅這麼賣力.

$^{241}$圍著合一的緣故.
而講述和教導這麼多.
我相信教會總是經常會面對各樣的轉變.
而當中不同的肢體.
因為愛教會的緣故.
我們總是會帶著各種不一樣的期望.
或許有人會覺得.
教會的步伐太慢了.
一點都沒辦法回應大家的訴求.
與此同時.
可能有人會認為.
教會的轉變太快太多了.
心裡面覺得很不安.
不過無論如何.
我們要相信一件事.
只要我們不灰心.
我們不中途離場.
我們繼續懸疑.
尋求弟兄姊妹的共識.
致力教會的合一.
這樣.
上帝是會在我們當中工作的.
祂會調節我們.
祂會連結我們.
因為上帝喜悅教會的合一.
祂期望信徒之間能夠互相配合.
彼此信服.
弟兄姊妹.
今天.
你是被動地接受教會所面臨的轉變.
還是積極地尋求共識.
同心參與教會新一年的福事呢?.
讓我們繼續看下去吧.
接下來.
保羅進一步說教會應該如何合一.
他說合一的群體.
唯有保持合一.
才能夠為上帝發光.
保羅向菲勒比信徒.
呼籲第二個行動是.

$^{281}$他說.
你們無論做什麼事.
都不要發怨言.
起爭論.
以致信徒能夠成為.
在這個彎曲薄茂的世代.
成為上帝無瑕疵的兒女.
到底發怨言起爭論.
和成為上帝兒女有什麼關係呢?.
保羅這句話.
原來是引至於摩西的.
在摩西臨終的時候有這樣的教導.
他警告當時的以色列人.
他說從前這些以色列民.
當他們進入迦林地之前.
他們的上一代.
在怪僻彎曲的世代.
在當時.
在他們的上一代裡面.
他們行了敗壞的事.
因為他們的弊病.
因為他們對上帝的背叛.
以致他們不能夠成為上帝的兒女.
保羅回顧當時的以色列人.
他們有失敗的經驗.
是因為他們向上帝發怨言.
彼此之間經常起爭論.
最終他們喪失了上帝兒女的身份.
不能夠進入應許之地.
所以保羅提醒菲臘比的信徒.
我們既然已經成為上帝的兒女.
我們就要分別為聖.
所以無可指責.
誠實無偽.
而且我們要成為一個分別為聖的群體.
首要做的是什麼呢?.
不是各自修行.
不是行善積德.
而是立即壓止教會內部的紛爭.
要彼此接納.

$^{321}$要互相的相愛.
或許教會的信徒.
追求聖潔仍然有很漫長的路.
但是彼此相愛.
是教會成聖的第一步.
正如耶穌這樣教導.
祂說:你很熟悉的經文.
我賜給你們一條新命令.
乃是叫你們彼此相愛.
我怎樣愛你們.
你們也怎樣彼此相愛.
如果你們有彼此相愛的心.
眾人因此就能夠認出.
你們是我的門徒.
跟隨耶穌的門徒.
最主要的特徵.
不是聰明過人.
不是特別神聖.
而是他們成為一群有愛的群體.
弟兄姊妹我想問一下.
當初你是被什麼吸引而回來教會的呢?.
又是什麼事情吸引你.
以致你留下在這裡呢?.
對我來說.
教會一開始吸引我的.
不是教會有一些超凡入聖生活的人在裡面.
也不是因為有高深的屬靈知識.
有很多很厲害的聖經專家.
更加不是因為有什麼名人在教會裡面聚會.
我當初被教會吸引的.
是因為這個群體很真誠.
當時我看到.
有我的牧師師母.
還有很多比我年紀大的哥哥姐姐.
當時我回來教會.
年紀還是很小.
他們不單是彼此之間的關係很親密.
他們更加接納我這個小子.
沒有將我排他冤外.
我記得當時我很小就回來教會.

$^{361}$沒多久.
當回到教會的團契一段時間.
當時我回少年團.
我第一次要嘗試參加教會的主日崇拜.
於是就參加了.
完結之後.
這個崇拜就完結了.
當時少年團的導師.
他做了一個動作.
他帶我.
拖著我的手.
我還記得.
他帶我一一認識教會其他的弟兄姊妹.
一年之後.
教會沒有一個人是我不認識的.
我當時就認定教會.
這一間教會.
就是我屬靈的家.
因為我在這裡.
我感受到幸福.
弟兄姊妹.
我相信在座不少人.
你都嘗過這種幸福.
你都是因為這樣而留下在這裡.
你願意讓身邊的人.
都能夠擁有這種幸福嗎?.
教會要吸引人.
感受耶穌的愛.
是需要付上代價的.
不過.
這個代價.
是你付得起的.
而且.
你會覺得.
付這個代價.
是值得的.
怎樣付出呢?.
保羅就繼續跟信徒說.
他說.
我們在這個世代裡.

$^{401}$要好像明光照耀一樣.
將生命的道顯明出來.
有關於這兩節.
十五十六節.
我另外找了一個翻譯.
是更加貼近原來的意思.
我們一起去看新漢語譯本的翻譯.
聖經是這樣說.
這段聖經可以翻譯成.
你們要堅守生命之道.
藉此在他們當中發光.
好像宇宙中的光體一樣.
事實上.
保羅沒有低估當時社會的複雜性.
他指出.
當時教會的世代是扭曲的.
是跟上帝的心意背道而馳.
不過即使如此.
保羅仍然認為.
教會是可以發揮它的力量的.
保羅形容.
我們好像什麼呢?.
我們好像宇宙的星體一樣.
是中數.
不是一粒.
是很多粒.
在黑夜裡.
一粒星.
光芒是有限的.
它不能夠發揮太大的作用.
不過.
當很多星一起發光的時候.
就能夠照亮黑夜.
讓人能夠辨別方向.
我記得四年多前的暑假.
在我自己服侍的教會.
有幾位大專生和畢業生.
他們去了韓國參加青少年.
呼音工作的會議.
他們回來之後很興高采烈地.

$^{441}$分享他們的所見所聞.
我記得.
他們提到最後一晚.
有一個陪靈的聚會.
他們回憶當晚是這樣的.
在這些聚會中間.
還是開始的時候.
他們每一個人都被派發.
有一支蠟燭.
不過這個蠟燭是沒有點燃的.
當所有的年輕人都坐滿了場地的時候.
後面開始就有人點燃那些蠟燭.
一開始是一個人.
幾個人點著那些蠟燭.
然後他們就將那些火光.
一直傳遞到每一個青少年人的手上.
讓他們的蠟燭都被燃點.
你可以想像一下當時的畫面.
一開始不是很光.
不過到了最後.
所有人都拿著燭光.
整個場地都發出光芒.
弟兄姊妹,當教會要守住上帝生命之道的時候.
單靠牧者.
單靠某一個教會的領袖.
我相信是不足夠的.
教會需要對公.
就好像耶穌選立十二個門徒一樣.
他們需要彼此配合.
互相勉勵.
互相打氣.
才能夠讓更多人認識上帝.
你試想像一下.
如果教會有多一些的信徒.
他們願意站出來.
成為屬靈的領袖.
人就能夠更立體知道.
信耶穌是一回事.
因為這些領袖.
是有來自不同背景.

$^{481}$不同性格.
不同恩賜的信徒.
他們可以知道.
原來我.
可能這種性格的弟兄姊妹.
我有些不同.
但我可以和我相似性格.
恩賜的弟兄姊妹.
去效法他們來服侍上帝.
而事奉的領袖.
我相信也因為這個緣故.
他們不再感到.
他們是孤身作戰.
關鍵是什麼呢.
關鍵是.
有沒有人願意站出來呢.
有沒有人願意站出來.
成為那個領袖呢.
領子妹.
我們沒有一個人.
生出來就成為教會領袖.
我們都需要上師.
我們都需要學習.
但倘若我們連試.
我們都沒有去試的話.
這幅圖畫就很難實現.
教會也很難來說服.
讓人認識生命之道.
最後.
很想和大家做一個小小的總結.
我相信.
合一的成功.
是在乎弟兄姊妹.
甘心樂意的付出.
甘心的犧牲.
最後.
保羅向菲臘比信徒呼籲第三個行動.
他說.
你們也要照樣喜樂.
並且與我一同喜樂.

$^{521}$到底是什麼事情.
值得保羅和菲臘比信徒喜樂呢.
保羅是這樣形容自己的工作.
他說.
我以你們的信心為貢獻的祭物.
我若被囂獻在旗上.
也是喜樂.
並且與你們眾人一同喜樂.
什麼是囂獻在旗上呢.
原來古時以色列聖殿的獻祭.
有一種祭叫做囂酒祭.
古時當祭司獻上祭身.
原來在獻的同時.
或者獻的之後.
他們會將酒灑在祭壇上.
以致所獻的祭的香氣.
更加濃郁.
一同像煙一樣升上去呈獻給上帝.
保羅是菲臘比教會的成長.
就好像呈獻給上帝的祭身一樣.
當祭身被燃燒.
香氣的煙緩緩上升的時候.
雖然保羅當時因為坐牢的緣故.
他不能夠與這群信徒在同一處的地方.
但是他將他自己為福音受苦的服侍.
比喻作囂酒祭的祭身的酒一樣.
當菲臘比信徒彼此合一不斷成長的時候.
保羅亦在另一個地方.
甘心為教會作為福音的使者.
犧牲和付出.
最終與菲臘比的信徒.
一起將他們最好的事奉呈獻給上帝.
就好像囂酒祭的酒一樣.
兩種的香氣.
兩種的濃煙.
最終升到天上.
混在一起.
就一同蒙上帝的喜悅.
弟兄姊妹.
保羅要說的是什麼呢?.

$^{561}$保羅要告訴當時的信徒.
教會要合一.
是需要一同付出的.
而他亦都是身先士卒.
他率先的付出.
他很希望透過這種犧牲.
讓大家可以跟隨.
效法著他為教會.
為上帝的緣故.
付上他們的福事.
所以跟隨耶穌的人.
是要帶著這種奉獻的心.
甘心付上造門徒的代價.
勇敢站出來.
彼此信服的.
你都是服侍上帝.
這樣我們就能夠常到做基督徒的快樂.
就能夠見證到天國的拓展.
但願弘恩堂家.
因著大家的付出.
以致能夠成為同心的教會.
因著大家的合一.
以致我們能夠成為明亮的神聖.
照亮黑暗.
讓天國得以拓展.
阿們.
\newpage



\section{約翰福音 13:34-35}
\label{sec:PDRx0v0AuiQ}
\textbf{2021.02.28《新的命令》約翰福音 13\_34-35( 和修版)  林佩心姑娘}
\newline
\newline
連結: \href{https://youtube.com/watch?v=PDRx0v0AuiQ}{\texttt{https://youtube.com/watch?v=PDRx0v0AuiQ}} ~~~~ 語音日期: 2021-02-28
\newline
\newline
\hyperref[sec:t7FE7L_V2U0]{\small{< < < PREV SERMON < < <}}
~
\hyperref[sec:index_chronic]{\small{[返順時目]}}
~
\hyperref[sec:index_scriptual]{\small{[返順卷目]}}
~
\hyperref[sec:JLt4gSCI6Y8]{\small{> > > NEXT SERMON > > >}}
\newline
\newline
約翰福音 13:34-35
\newline
\begin{longtable}{cl}
\hline
\hline
章節 & 經文 (和合本修訂版)\\
\hline
13:34 & \begin{tabularx}{0.7\textwidth}{X} 我賜給你們一條新命令，乃是叫你們彼此相愛；我怎樣愛你們，你們也要怎樣彼此相愛。 \end{tabularx} \\ \\ \relax
13:35 & \begin{tabularx}{0.7\textwidth}{X} 你們若彼此相愛，眾人因此就認出你們是我的門徒了。」 \end{tabularx} \\ \\
[1ex]
\hline
\hline
\end{longtable}
$^{1}$今天第一次來到紅恩堂跟大家分享神的話語.
在開始之前我們一起祈禱.
親愛的天父愛多謝你.
你在一年的開始也帶領孩子來到紅恩堂.
來分享神你自己的信息.
主啊願你的靈在我們當中啟迪我們.
賜我們有屬靈的悟性.
有屬靈的耳朵.
使我們所聽的所領受的.
都是從神理而來的信息.
有潔淨孩子.
被我所說的.
都是神理的心意.
祈禱靠耶穌基督的德性的名.
阿們.
好.
2021年的新年.
所有新衣服.
新鞋.
新髮型.
來一個新的形象.
來一些新年的傳統.
在疫情之下.
不知不覺就失落了.
按照衛生防護中心的呼籲.
過年我們每個人都仍然需要減少一些家庭的聚會.
所以熱熱鬧鬧的傳統就失落了.
反而視像拜年.
不知不覺成為我們疫情中一項的新常態.
今天跟大家分享幾個經文.
是記載在《日行福音》的第十三章.
三十四至三十五節.
耶穌一開始就這樣跟門徒說.
祂說:我賜給你們一條新命令.
乃是叫你們彼此相愛.
我怎樣愛你們.
你們也要怎樣相愛.
在新的一年開始.
我們一起來思想.
耶穌這個新的命令.

$^{41}$究竟這個新的命令.
在什麼情況下來發出呢?.
當天其實是在預鑰節的晚上.
耶穌就知道.
祂跟門徒吃這一頓飯.
對祂來說就是祂在地上吃飯最後一頓飯.
當時祂的心裡面就很憂愁.
就跟門徒說.
我實實在在地告訴你們.
你們中間有一個人要賣我.
當人知道有人要出賣自己.
一般的反應就是很生氣.
按照人自我防衛的機制.
就揭發這個人的陰謀.
讓他在眾人面前慕地自容.
這是很人之常情的正常反應.
當耶穌知道猶大一定會出賣自己的時候.
祂反而很憂愁地跟門徒說.
祂說:我賜給你們一條新命令.
乃是叫你們彼此相愛.
我怎樣愛你們.
你們也要怎樣相愛.
我們一起思想一下.
在聖經裡面的一個愛的命令.
其實對以色列人來說.
他們從小到大.
在律法書中.
律法書裡面怎樣教導愛呢.
它說你要盡心盡性盡力.
愛耶和華你的神.
所以愛的誡命不是一個新的誡命.
而且律法書也有這樣教.
它說你是不可報仇.
也不可埋怨你本國的子民.
它說要愛人如己.
舊約裡面.
愛的誡命是對神一種完全的奉獻.
就是盡了.
盡是什麼呢.
盡就是盡.

$^{81}$多盡就去到多盡.
去到多盡就多盡.
盡自己心裡面的力量.
耶穌不可以面對要被人出賣.
祂對他們的門徒說.
我賜給你們一條新命令.
乃是叫你們彼此相愛.
我怎樣愛你們.
你們也要怎樣相愛.
這個新的愛的命令.
就是你們要彼此相愛.
這個不單單是從自己出發.
是要跟別人配合的.
因為彼此.
是要互動產生出來的一種果效.
另外單方面就不能夠成效的.
最少也有兩個人.
重點就是彼此.
彼此是怎樣的.
彼此要願意付出.
以及要願意接受愛.
真的不容易.
所以當一對情侶拍拖.
我們就要恭喜他.
為什麼呢.
因為真的不容易.
是很值得恭喜的.
結婚週年紀念真的要值得慶祝.
不要說老夫老妻結婚三十年了.
說我們珍珠婚.
我們金婚.
更加要慶祝.
為什麼呢.
當中經歷了每一天一起生活的磨合.
生活的點點滴滴.
包括很多的包容和接納.
不是一天.
不是一年.
甚至你過了十年二十年三十年也好.
每一天.

$^{121}$就算兩個人一起生活.
都會有新的挑戰.
兩個人都不是那麼容易.
耶穌現在對著他十二個門號說.
你們要彼此相愛.
對他們來說這不是愛的更新版.
簡直就是強化版.
加強版.
複雜了很多.
耶穌的呼召.
就是每一個人對耶穌的回應.
耶穌呼召我.
我對耶穌回應.
我祈禱.
我信服神.
我依靠神.
我事奉神.
我已經盡了愛神.
盡力愛神的力量了.
但是耶穌在這裡說.
在盡這個誡命上.
加上了彼此.
加上了一種人際的互動.
不是一個人單單可以做.
奉獻給神就夠了.
所以盡心盡性盡力愛神.
還要彼此相愛.
有一個目標的.
就是一同見證我們是耶穌基督的門徒.
是一個教會的見證.
這個愛的加強版.
帶來了一個新的挑戰.
第一.
我們要思想一種關係上的重整.
在約翰福音最後一章.
第21章.
約翰記下了.
他的其中一個同學.
彼得.
和耶穌有一段對話.

$^{161}$當時彼得指著約翰.
問耶穌.
主啊.
這個人將來如何?.
約翰說.
約翰將來如何?.
耶穌就對他說.
我若要他等到我來的時候.
與你何干?.
你跟從我吧.
在人際關係裡面.
被認同.
被肯定.
是一個很人性.
一個很基本的需要.
所以一個比較.
也有一個人性很自然的反應.
重點.
在於我們如何去看人際關係.
如果我們經常以一個宿敵的態度.
來看我們的同伴.
那麼我們在家庭生活裡面.
我們的兄弟姐妹全部.
我們一起爭寵.
媽媽疼誰多一點.
爸爸疼誰多一點.
爺爺經常買東西.
封利是裝一下.
其實我也試過.
有一次我真的.
有一點不開心.
我媽媽經常封.
一百塊利是給我.
我很開心的.
當我發覺我媽媽.
原來封給她的孫子.
就是我哥哥的兒子封五百塊.
我立刻嘴巴吐出來.
我立刻扁嘴.
這個很人性.

$^{201}$我是女兒.
一百塊是老爺孫.
又五百塊.
唉.
真是的.
那麼.
如果我們這樣看的時候.
同學.
很老友.
同學一起成長.
我們一起.
談一些項目.
做功課.
一起搞定一些事情.
但是呢.
如果我們當同學都是宿敵.
那麼我們就會在我們的成績.
和課外活動裡面.
經常在老師面前爭取表現.
是不是?.
嗯.
好了.
同事更簡單.
所以特徵職加薪.
是不是?.
多大都要博上位.
如果我們認為.
這些是我們人生裡面.
不能夠擺脫的競爭對手.
彼此相愛.
就是強化中的強化.
比強化玻璃更強化.
很難衝破的.
所以彼此相愛.
不是一種理想.
耶穌一路說.
是一個命令.
命令的意思.
就是要做的.
要實踐的.

$^{241}$耶穌這樣問彼得.
彼得說到彼得的將來.
為什麼耶穌會指著約翰韋.
他怎樣呢?.
原來耶穌之前有跟他說.
彼得你將來會怎樣的.
耶穌說.
彼得.
你為我養的羊.
我實實在在地告訴你.
你年輕的時候.
自己戳上帶子.
隨意往來.
大年老的時候.
你要伸出你的手.
別人要把你戳上.
帶你去.
到你不願意去的地方.
耶穌說這句話.
是指著彼得.
將來會怎樣.
死來榮耀神.
說到這裡.
耶穌對他說.
來.
跟從我吧.
彼得很盡.
他真的盡心盡性盡意.
愛著他的神.
說完之後怎樣?.
他繼續跟.
跟到底.
雖然後來他被耶穌說.
其實你心靈有軟弱的.
你會三思不認我.
但是那一剎那.
他真的盡了他的心理能量.
說要跟從耶穌.
耶穌告訴彼得.
他說你年輕的時候.

$^{281}$你做什麼都可以.
你去哪裡都可以.
但是你年老的時候.
你所遇到的事情.
跟你年輕時候那種.
隨意考自由.
去服侍主.
除非什麼?.
是成為強烈的對比.
你那種自由而獨立的性格.
因著主的緣故.
會被剝奪的.
你要被迫伸開雙手.
被人觸賞.
帶你去你不願意去的地方.
伸開雙手的意思.
就是只要彼得.
就要跟當時會被釘十字架的人一樣.
當時要被人釘十字架的人.
他要伸開他的手背.
背起自己的十字架.
走去行刑的地方.
彼得將來就是這樣.
殉道榮耀神.
彼得聽到自己將來的結局之後.
他什麼都不說.
他沒有說主啊.
你將這個苦杯挪開.
主啊.
你救我.
什麼都沒有說.
他很奇怪.
很奇怪就說.
那約翰怎麼辦?.
就是耶穌的回憶.
與你何干.
你跟從我吧.
從此可見彼得不僅是為了好奇心而知道這件事.
他的內心有一種很人性的.
比較.

$^{321}$哦.
那我這樣殉道.
那約翰怎麼辦?.
就是他.
就是怎樣?.
我好死一點.
他好死一點.
信征頭不喜歡聽.
不要緊我們做基督徒.
我們也是聖經這樣說的.
他內心是比較.
自己約翰將來的日子怎樣.
結局怎樣.
在跟隨主的人生裡面.
各人都有不同的經歷.
耶穌的吩咐是.
你來跟從我.
不需要比較.
你專心跟從神的帶領.
根據聖經的記載.
彼得成為耶路撒冷教會的領袖.
他在五旬節聖靈降臨的時候.
宣講大有力量.
厲害了.
宣告耶穌是基督是西人的救主.
神給他滿滿的果子.
三千人悔改信主.
Hallelujah.
暫時.
他即時成為耶路撒冷教會的領袖.
耶路撒冷教會.
就是這樣誕生的.
但是後來因為逼迫的緣故.
他沒有辦法.
他就被迫離開耶路撒冷.
到了羅馬傳福音.
究竟使徒約翰怎樣呢?.
根據聖經所說.
使徒約翰很厲害.
因為他得到神很大的啟示.

$^{361}$他寫了約翰一二三書.
內容教導信徒.
如何做一個合神心意的基督徒.
鼓勵他們在生活上如何表現彼此相愛.
過聖潔的生活.
在信仰上要堅持.
要戰勝異端的誘惑.
要堅持堅守真道.
但是後來他年老的時候.
他被放逐到拔摩島.
他的事奉還沒有完結.
人放逐他.
但是神就說.
人想不到.
在那個情況下.
讓他自己的僕人有一個很大的意象.
將來世界的終局的事告訴他.
他就能夠在放逐一個很寧靜的地方.
沒有人理會他.
他就可以將神給他的啟示一一寫下來.
寫成啟示錄.
而且他真的是十二個門徒裡面最長命的一個.
他是自然死去的.
不是別人釘他的十字架.
當耶穌面對要走上十字架的晚上.
他在撒克西勒瑪利園祈禱.
他帶著彼得,雅各和約翰.
他和他最親密的門徒.
就吩咐他們.
你們跟我一起在這裡警醒禱告.
彼得和約翰.
都是一起在耶穌面對十字架掙扎的時候.
陪伴他的門徒.
所以耶穌沒有比較他們這兩個門徒.
他們有不同的恩慈.
但是他們有共同的心志.
為神國付上自己的人生.
付上自己的生命.
盡他們所能.
去事奉主.

$^{401}$愛主愛神.
他們都是成為主所愛的門徒.
所以在愛的加強版.
彼此相愛的裡面.
第二方面我們要想一下.
我們要重整一下自己事奉的目標.
耶穌說:我怎樣愛你們.
你們也要怎樣彼此相愛.
這樣眾人就認出你們是我的門徒了.
耶穌基督的門徒.
就是承接著福音的棒.
要將耶穌基督的愛的見證,犧牲的見證.
救贖的見證繼續傳揚開去.
這是他們奉獻給神.
盡心盡性盡義盡力愛主你的神的目標.
所以耶穌基督犧牲十字架的愛.
修復了人和神的破裂關係.
這是神和人的一個新約.
是一個恩典之約.
恩典.
原本應該要死在過犯罪惡之中的人.
因為耶穌基督十字架犧牲的救贖罪得赦免.
使我們恢復神兒女的名分.
所以耶穌當祂面對.
要有人傷害祂.
要釘祂十字架的時候.
耶穌說你們要彼此相愛.
耶穌對門徒的吩咐.
彼此相愛.
對他們來說是一個很深的觀念.
因為你們要彼此相愛.
我怎樣愛你們.
你們也要怎樣相愛.
當時的門徒可能都知所以而不知所以而.
他們不太了解究竟是怎樣的.
當耶穌被釘十字架之後.
他們就知道了.
因為耶穌對他們說.
你們要彼此相愛.
我怎樣為你們犧牲愛你們.

$^{441}$你們就要這樣犧牲彼此相愛.
犧牲.
可能對於門徒來說.
是一個很深很深的觀念.
他們沒有聽過.
但耶穌就是這樣.
我們經歷過.
每一個相信耶穌基督的人.
我們都經歷過.
耶穌基督十字架.
犧牲的愛的無條件的救贖.
這是人際關係的更新版.
祂將愛的力量推到最高.
發揮愛最大的力量.
彼此以犧牲的愛.
來實踐基督的愛.
這個愛的加強版.
超越了人際關係一種很基本的互動.
史圖普羅這樣說的.
他說.
「因為我們還軟弱的時候.
基督就按所定的日期為罪人死.
為異人死是少有的.
為人人死或有感作的.
唯有基督在我們還作罪人的時候.
為我們死.
神的愛就在此向我們顯明」.
《羅馬書》五章六至八節.
去到盡.
這樣去愛神.
愛人.
是不錯的.
但是.
犧牲的愛比愛人如己.
在層次上.
不是更高一層樓這麼簡單.
耶穌對門徒說這是一個新的命令.
是要遵守的.
要實踐的.
要做的.

$^{481}$要成為你生命裡面.
你要走出來的.
這是你的生命的一部分.
唯有這樣.
這個世人未認識耶穌基督的人.
才認得出你們是我的門徒.
命令的目標很清晰.
為主作見證.
見證我們是世上.
被分別出來的群體.
就是那個跟從耶穌基督的門徒.
如何實踐彼此相愛.
我們首先看看耶穌是如何服侍.
在宣告一個新的命令以前.
耶穌那天是逾越節的晚餐.
耶穌原來為門徒洗腳.
這個工作本身是誰做的.
是家裡的婢女.
就是下人做的.
有一天我去親戚家裡吃飯.
裝飯.
我近距離看.
我的親戚說.
這些是下人做的.
就是家裡的傭人姐姐.
當時人們經常有這種階級的觀念.
其實這些工作是怎樣的呢?.
當一個很好的主人.
我們都知道巴勒斯坦地區.
我們不是穿我們現在的鞋子.
我們穿涼鞋.
一天走很遠的路.
整個腳都是灰塵.
很累.
腳都腫了.
又沙沙石石.
所以當人客來到到訪的時候.
通常那些屋子前面都有幾缸水.
婢女就會裝水.
幫人客洗乾淨腳.

$^{521}$不是怕踩進弄髒的地方.
而是讓他經過一天的行走.
腳都腫了.
沙沙塵塵.
黏著腳很不舒服.
讓他清潔了.
讓他清涼一下.
又舒服一點.
是一個很殷勤接待的禮貌.
是一個款待.
但那天是沒有婢女.
在門徒裡面.
你看我一眼.
誰做這件事.
誰做.
為什麼他們不做.
因為怕別人做.
12個門徒.
沒有一個想自己是門徒裡面最小的.
每個人都想考第一第二第三.
我是大門徒.
所以彼得會問耶穌約翰將來如何.
很仁性.
耶穌看他們不做.
就自己做.
為他們洗腳.
以身作則.
教導一種謙卑的服事.
耶穌經常說.
誰願為大.
就作眾人的僕人.
所以教會裡面.
應該不會這樣分這些是誰做的.
教會裡面.
就是看那樣做那樣.
每樣都做.
你覺得這是服事主的.
服事弟兄姊妹的.
不是積分他.
這些不是我做的.

$^{561}$這些是他的.
做他做得不好.
開會的時候說.
為什麼你做得不好.
為什麼他這樣.
你漏了什麼.
這樣完全不是耶穌的教導.
耶穌是怎樣的.
保羅說他本有神的形象.
卻不堅持自己與神同等.
反倒虛己.
取了奴僕的形象.
成為人的樣式.
既有人的樣子.
就謙卑自己.
全心信服以至於死.
且死在十字架上.
耶穌基督就是這樣倒空自己.
原來彼此相愛的.
怎麼能夠實踐.
第一倒空自己.
有首歌叫做亡妙我是誰.
不要管自己在公司裡.
自己的身份地位在家裡.
在新約教會裡.
保羅經常教.
自主的為奴的.
在耶穌基督新約教會裡.
不分彼此.
不應該這樣分.
這樣才能實踐彼此相愛.
謙卑服事.
我們一些不配的罪人.
就是耶穌基督為我們立下.
一個犧牲的愛的榜樣.
耶穌吩咐門徒.
要彼此相愛.
不是他們是很可愛的人.
相反在三年多的時間裡.
如果我們認真讀一下科幻書.

$^{601}$耶穌的門徒.
每個都有差不多的性格.
每個都有自己的硬的東西.
很難忍受的東西.
硬的東西.
耶穌全部都接納了他們.
不是硬吃.
而是愛和包容.
耶穌有十二個門徒.
每個都有優點.
不過都有缺點.
也都有軟弱.
耶穌當然很清楚.
十二個門徒之間.
彼此的關係.
所以彼得問他的時候.
你來跟從我.
我與你何干.
你們不要比較了.
你們每一個都是在耶穌基督的裡面.
在我裡面是我的門徒.
十二個門徒自己都很清楚.
彼此之間的關係.
當耶穌吩咐他們要彼此相愛.
他們可能都嘔了.
不懂得回應.
彼得後來就說.
你知道我會很愛你的.
耶主.
彼得自認自己會為主捨命.
耶穌也沒有批評彼得.
為什麼他們說得這麼快.
全知的神預告.
其實你心裡面有一種戰經.
叫做「熬底」.
我被人抓的時候.
你就「熬底」.
你三思不認.
很自然.
沒有說是負面.

$^{641}$有些事情我們沒有經歷過.
你會害怕.
耶穌就是這樣接納.
就是這樣包容.
因為他認識我們每一個.
彼得衝動的性格.
應該不是很可愛.
人家還沒說完他就插嘴.
但是耶穌仍然認為.
他在門徒裡面能夠實踐彼此相愛.
可愛和不可愛.
其實不是一個絕對的標準.
其實在我們日常生活裡面.
很多時候特別是我們中國人.
我們經常都會說關係.
由關係出發.
我的親人.
朋友.
同事.
甚至我的寵物.
我家裡的狗.
他不是很多缺點.
但他是我最愛的.
的狗.
應該不是絕對可愛.
每個人都有不同的缺點.
但是當他們是我們的親人.
他們是我的父母.
他們是我的兒女.
說這個是死黨.
是我同天的同事.
如果我們真的這樣看.
我們不會經常以宿敵的態度.
看我們周圍的人的話.
我們就會這樣說.
我們經常都會聽到人說.
不好意思.
他是這樣的.
不好意思.
其實他搞錯了.

$^{681}$為什麼你認錯.
他搞錯了.
為什麼你認錯.
因為你愛他.
你跟他包底.
是很大的包容和接納.
愛那些不可愛的人.
耶穌就是這樣愛我們.
犧牲的愛.
可以在具體的生活裡面實踐.
耶穌完完全全知道.
我們每一個人的限制.
我們的軟弱.
但是他以無條件的愛來愛我們.
而且當我們願意認罪.
我們願意悔改.
耶穌是會接納我們.
耶穌的教導就是怎樣.
他怎樣愛我們.
我們也要怎樣彼此相愛.
下法耶穌基督愛那些不可愛的人.
接納那些.
我們認為不應該接納的人.
要恕那些傷害我們的人.
並且為他們犧牲.
彼此相愛要成為耶穌基督教會.
耶穌基督的門徒.
耶穌基督教會有力的見證.
弟兄姊妹.
其實我們不單單在教會裡面.
可以要付出愛的力量.
剛才我說過.
是互動的.
是彼此的.
所以更加基本的.
我們要有被愛的勇氣.
有些人的性格比較強.
別人表示對你關心.
表示對你愛.
表示對你的支持.

$^{721}$你也可能.
不用了.
我自己可以了.
好像覺得不是.
好像覺得如果我接受別人.
這樣愛我.
我就好像很弱者.
耶穌基督不是.
我們要讓.
我們要來到神的面前.
為什麼我們可以愛人.
要犧牲.
我們要來到耶穌基督的面前.
經歷耶穌基督對我們生命的犧牲的愛.
他醫治我們有罪的生命.
使我們這個不完全的人生.
在耶穌基督裡面.
能夠與神和好.
修復和好的關係.
所以我們不要逃避.
面對耶穌基督犧牲的愛.
我們要承認自己是.
蒙恩不配的罪人.
我們只不過是在救了耶穌基督的恩典裡面被救贖.
他不喜歡我們的罪.
卻為我們這些罪人.
死在十字架上.
救贖我們.
願救主耶穌基督的愛.
是今天再一次成為我們的激勵.
使我們能夠活出彼此相愛.
成為愛的見證群體.
我們在這裡同心祈禱.
愛我們的主.
我們再一次來到面前.
獻上我們的感恩.
獻上我們的讚美.
因為你就是這樣.
為我們這些不配的罪人.
死在十字架上.

$^{761}$你以你寶貴的生命.
犧牲你的生命.
成為我們的救贖.
你在十字架上受的痛苦.
潔淨了我們.
除去我們的罪孽.
你的寶血成為我們罪的潔淨.
讓我們再一次來到主你自己的面前.
再次被你犧牲的愛激勵.
也讓我們有柔軟的心.
我們效法救主基督你.
我們學習.
我們願意謙卑自己.
去進入人群當中.
愛那些不可愛的人.
願意主你再次我們柔軟的心.
讓我們接納救主基督.
你每一天為我們的罪.
在十字架上釘死.
讓我們每一天願意接受救主基督.
你犧牲的愛.
成為我們的激勵.
願我們能夠接受弟兄姊妹的代禱.
接受弟兄姊妹的關心.
接受弟兄姊妹的支持.
我們也都在一個接受的裡面.
我們的經歷的裡面.
我們又能夠付出.
我們為我們的弟兄姊妹守望.
為我們的弟兄姊妹禱告.
我們為我們的弟兄姊妹獻上我們的關心.
成為他們的支持.
成為他們的同行者.
陪伴我們這樣.
我們在教會裡面.
我們在耶穌基督裡面.
主願意主你再次賜福給我們.
我們紅恩堂.
能夠成為一個愛的群體的見證.
在紅水橋這個地方.

$^{801}$高舉耶穌基督你.
十架犧牲的愛.
成為這個社區的祝福.
主願意你.
幫助我們的教會.
每一個弟兄姊妹.
在你面前的祈禱.
靠耶穌基督你得勝的名.
阿們.
\newpage



\section{撒母耳記上 16:1-13}
\label{sec:JLt4gSCI6Y8}
\textbf{2021.03.07《神怎樣看我們》撒母耳記上 16\_1-13( 和修版)  老耀雄牧師}
\newline
\newline
連結: \href{https://youtube.com/watch?v=JLt4gSCI6Y8}{\texttt{https://youtube.com/watch?v=JLt4gSCI6Y8}} ~~~~ 語音日期: 2021-03-11
\newline
\newline
\hyperref[sec:PDRx0v0AuiQ]{\small{< < < PREV SERMON < < <}}
~
\hyperref[sec:index_chronic]{\small{[返順時目]}}
~
\hyperref[sec:index_scriptual]{\small{[返順卷目]}}
~
\hyperref[sec:2_wlsYE_36Q]{\small{> > > NEXT SERMON > > >}}
\newline
\newline
撒母耳記上 16:1-13
\newline
\begin{longtable}{cl}
\hline
\hline
章節 & 經文 (和合本修訂版)\\
\hline
16:1 & \begin{tabularx}{0.7\textwidth}{X} 耶和華對撒母耳說：「我既厭棄掃羅作以色列的王，你為他悲傷要到幾時呢？你將膏油盛滿了角；來，我差遣你到伯利恆人耶西那裡去，因為我在他兒子中已看中了一個為我作王的。」 \end{tabularx} \\ \\ \relax
16:2 & \begin{tabularx}{0.7\textwidth}{X} 撒母耳說：「我怎麼能去呢？掃羅一聽見，就會殺我。」耶和華說：「你可以手裡牽一頭小母牛去，說：『我來是要向耶和華獻祭。』 \end{tabularx} \\ \\ \relax
16:3 & \begin{tabularx}{0.7\textwidth}{X} 你要請耶西來一同獻祭，我會指示你當做的事。我對你說的那個人，你要為我膏他。」 \end{tabularx} \\ \\ \relax
16:4 & \begin{tabularx}{0.7\textwidth}{X} 撒母耳遵照耶和華的話去做，來到伯利恆，城裡的長老都戰戰兢兢出來迎接他，有人問他說：「你是為平安來的嗎？」 \end{tabularx} \\ \\ \relax
16:5 & \begin{tabularx}{0.7\textwidth}{X} 他說：「為平安來的，我來是要向耶和華獻祭。你們要使自己分別為聖，來跟我一同獻祭。」撒母耳把耶西和他眾兒子分別為聖，請他們來一同獻祭。 \end{tabularx} \\ \\ \relax
16:6 & \begin{tabularx}{0.7\textwidth}{X} 他們來的時候，撒母耳看見以利押，就心裡說，耶和華的受膏者一定在耶和華面前了。 \end{tabularx} \\ \\ \relax
16:7 & \begin{tabularx}{0.7\textwidth}{X} 耶和華卻對撒母耳說：「不要只看他的外貌和他身材高大，我不揀選他。因為耶和華不像人看人，人是看外貌，耶和華是看內心。」 \end{tabularx} \\ \\ \relax
16:8 & \begin{tabularx}{0.7\textwidth}{X} 耶西叫亞比拿達從撒母耳面前經過，撒母耳說：「耶和華也不揀選他。」 \end{tabularx} \\ \\ \relax
16:9 & \begin{tabularx}{0.7\textwidth}{X} 耶西又叫沙瑪經過，撒母耳說：「耶和華也不揀選他。」 \end{tabularx} \\ \\ \relax
16:10 & \begin{tabularx}{0.7\textwidth}{X} 耶西叫他七個兒子都從撒母耳面前經過，撒母耳對耶西說：「這些都不是耶和華所揀選的。」 \end{tabularx} \\ \\ \relax
16:11 & \begin{tabularx}{0.7\textwidth}{X} 撒母耳對耶西說：「你的兒子都在這裡了嗎？」他說：「還有一個最小的，看哪，他正在放羊。」撒母耳對耶西說：「你派人去叫他來；他若不來這裡，我們必不坐席。」 \end{tabularx} \\ \\ \relax
16:12 & \begin{tabularx}{0.7\textwidth}{X} 耶西就派人去叫他來。他面色紅潤，雙目清秀，容貌俊美。耶和華說：「起來，膏他，因為這就是他了。」 \end{tabularx} \\ \\ \relax
16:13 & \begin{tabularx}{0.7\textwidth}{X} 撒母耳就用角裡的膏油，在他的兄長中膏了他。從這日起，耶和華的靈就大大感動大衛。撒母耳起身回拉瑪去了。 \end{tabularx} \\ \\
[1ex]
\hline
\hline
\end{longtable}
$^{1}$紅映街的弟兄姊妹平安.
中國人有一句說話叫做.
先敬羅衣後敬人.
意思就是在人際關係裡面.
大部份的人都以衣著作為一個衡量的標準.
穿得乾鮮一點.
抑或穿得又殘又舊.
給人的印象是會截然不同.
或者這樣說.
穿得乾淨整齊.
亦都是對人的一種尊重.
給人一種信心.
亦都留下一個很良好的印象.
方便日後交往的時候.
能夠有一個很好的合作和溝通的基礎.
當然先敬羅衣後敬人背後說的是一種包裝.
不單是衣著.
亦講求身份地位住址職業.
甚至乎是駕駛甚麼車.
以前的說法叫做批頭.
今日的說法叫做包裝.
雖然說法不一樣.
但其實都是一種形象工程.
以一個最好的形象展現於人前.
以達至到最好以及最吸引的效果.
從而得到最佳的效益.
今日的形象設計是最吃香的行業.
亦都是一種能夠點石成金的行業.
這種熟世的價值觀.
以及聖經裡面的屬靈價值觀.
究竟有甚麼分別呢.
今日很想透過撒母耳記上的十六章一至十三節.
去看看我們每一個信徒在神面前.
怎樣去建構我們一個生命的形象工程.
我們一起去祈禱.
求主你讓我們有從你而來的智慧.
去過一個合乎你心意的信仰人生.
我們亦都求聖靈對我們作出引導.
讓我們能夠謙卑地領受你對我們的提醒.
我們的祈禱是奉主耶穌基督得勝明治祈求.

$^{41}$阿們.
經文一開始就講了一件事.
就是耶和華已經厭棄了蘇羅作以色列的王.
厭棄的原因有兩件事是最關鍵的.
第一件事在十三章.
撒母耳要求蘇羅在吉格等他.
而蘇羅等了撒母耳七天.
撒母耳都還未出現.
當時非利士人亦都在吉格集結.
準備攻打以色列.
蘇羅當時就很心急.
他不再等撒母耳來到.
而私自獻祭.
希望能夠鞏固軍心.
後來蘇羅還埋怨.
撒母耳為何這麼遲都不出現.
和撒母耳去辯解.
他說所以我勉勉強強地獻上繁祭.
表示我這樣做.
勉為其難而已.
不是我想的.
撒母耳後來跟蘇羅說.
他說你做了糊塗事了.
沒有遵守耶和華.
李神吩咐的命令.
如果不是.
耶和華會在以色列中.
建立你的國度直到永遠.
在這件事.
耶和華沒有因為這件事而放棄蘇羅.
只不過就是他的國度.
不可以延續給下一代.
比較關鍵的第二件事.
在十五章.
耶和華要求蘇羅前去攻打亞馬力族.
並且要求滅絕他們所有族人.
和所有的生畜.
不過蘇羅卻留下亞馬力王阿格.
也沒有將所有的生畜殺盡.
最後要撒母耳在吉格裡.

$^{81}$將阿格殺了.
當撒母耳質問蘇羅.
為什麼你不遵守耶和華的命令.
耶和華就推了到百姓身上.
他說是百姓將這些老弱之物留下最好的.
是要獻給耶和華.
因此不關蘇羅事.
蘇羅第一次將責任推卸在撒母耳身上.
說他遲到.
所以他要私自獻祭.
第二次將責任推卸到百姓身上.
是他們留下的生畜.
是他們要獻給耶和華.
所以又不關蘇羅事.
所以最後撒母耳就向蘇羅宣告.
他說因為你厭棄耶和華的命令.
因此耶和華都厭棄你作王.
耶和華厭棄蘇羅的原因.
最主要是他不遵守耶和華的命令.
而且作為以色列君王.
他不願意承擔自己所犯的罪.
將責任推卸給別人.
當撒母耳還為蘇羅被耶和華廢除一事難過的時候.
耶和華就吩咐不要再難過了.
因為要前往伯利衡.
去高納耶和華所揀選的新王.
撒母耳即時心表示擔心.
他說我怎麼能夠去.
蘇羅一聽見就會殺我.
因為無論是真的高納.
還是意圖高納新王.
其實在蘇羅的眼裡都是一個叛國.
撒母耳和新王隨時都因為這樣而有生命的危險.
另外撒母耳怎麼進入伯利衡.
這也是一個問題.
因為如果讓遍雅門支派.
或者猶大支派的人知道.
他這個行動背後的目的的話.
必然會引起支派的內部分裂而爆發內戰.
耶和華沒有責怪撒母耳的信心.

$^{121}$或者為什麼要害怕.
反而提出一個可行的計劃.
讓撒母耳去完成使命.
就是他可以帶著牛族前去伯利衡獻祭.
並且安排耶西的家人一起去參與.
然後再高納一個新的王.
他說我會指示你當做的事.
就是讓撒母耳能夠安心前去.
對於撒母耳的來到.
伯利衡的長老很害怕.
戰戰兢兢地出來迎接.
他們不知道撒母耳此行的目的是來審判.
還是來祝福.
所以他們劈頭第一句就問.
你是不是為平安而來的.
因為他們覺得如果以審判來說.
他們就當災了.
反映當時撒母耳的身份和權柄.
是受到百姓的尊重.
撒母耳安撫他們.
為平安而來的.
你們放心.
耶和華獻祭有三個步驟.
第一個就是自結.
第二個就是獻祭.
第三個就是飲宴.
所以在眾人自結的過程裡面.
撒母耳就可以逐一逐一地去審視.
耶西的每一個兒子.
為他們行結敬禮.
撒母耳以為是高納耶西的長子.
以利亞.
但是耶和華提出了他揀選的標準.
他說不是只是看他的外貌.
和他身材高大的.
我不揀選他.
耶華說不像人去看人.
人是看外貌的.
但是耶和華是看內心的.
當耶西的七個兒子都被撒母耳審視完畢.

$^{161}$但是沒有一個人是耶和華所揀選的.
因此撒母耳就問耶西.
還有沒有兒子還沒有出現.
原來還有一個最少的正在放養.
耶西於是派人去叫大衛回來.
經文用他面色紅潤.
雙目清瘦.
容貌盡美來形容大衛.
大衛當時可能年紀還小.
身材可能比較矮小.
但是容貌是盡美的.
耶華對撒母耳說.
起來,高他,因為這就是他了.
撒母耳於是即時在大衛的兄長面前.
高納了大衛.
當時長子以利亞和大衛.
同樣都是容貌盡美的.
但是最後耶華就揀選了大衛.
因為耶華是看人的內心.
我們就問.
究竟大衛的內心是怎樣的.
大衛的內心是有什麼特殊.
同樣以利亞和他都是容貌盡美的.
為什麼不揀選以利亞而揀選大衛呢.
什麼因素讓耶華揀選大衛呢.
這件經文沒有透露.
我們只能夠在上下文去找答案.
大衛在第十七章的時候.
他曾經請纓去挑戰非利士人哥利亞.
他曾經和素羅說.
他說他為父親放羊的時候.
有時候會遇到獅子或者熊.
他就追趕他們.
擊殺他們.
將那些羊從獅子和熊的口中救出來.
因此大衛深信.
既然耶華說能夠救他脫離獅子和熊的爪.
今天耶華說同樣會將他從哥利亞手中拯救出來.
因此在薩姆爾見到大衛之前.
或者正確地說.

$^{201}$在薩姆爾未高立大衛之前.
大衛屬靈生命的經歷.
就只是在放羊的時候.
他遇見獅子或者熊.
而大衛面對這個情況.
他卻有一個完全倚靠耶華的心.
令他可以在殺獅殺熊的過程裡面.
建立對神的信心.
亦鞏固對神的倚靠.
這個就正正可以解釋了.
耶華在第七節裡面說.
耶華是看內心的.
因為耶華見到大衛的內心裡面.
對他有一種完全信靠的心.
這種完全信靠的心.
在面對危難的時候.
可以將整個生命完全交託.
無畏無懼.
不單單是這樣.
這個信靠耶華的心態.
亦激發了一種能夠為神爭戰的信心.
大衛面對哥利亞的感受.
就是神的名怎能夠被污辱.
這個未受國禮的非利士人哥利亞.
竟然夠膽咒罵耶華的軍隊.
大衛在面對哥利亞的時候.
他講了一個信仰宣言.
顯明他的出戰.
或者他要應付哥利亞.
並不是因為素來有甚麼獎賞.
而是他為耶華而戰.
他講了甚麼呢.
他說你來攻擊我.
是靠著刀槍和銅矛.
但我來攻擊你.
是要靠著萬軍的耶華之名.
就是你所辱罵帶領以色列軍隊的神.
事實上.
素來和當時所有的以色列人.
都懼怕哥利亞.

$^{241}$這些以色列的軍兵.
全部都擁有真實的作戰經驗.
真的在戰場上打過仗.
而大衛只不過是在牧羊的時候.
去擊退那些獅子和熊的經驗.
這種經驗又豈能和真實的戰場相比呢.
你試想一想.
一個毫無作戰經驗的大衛.
他不穿著一個作戰的軍服.
只是拿著一把垃圾機.
竟然有膽量.
有勇氣去挑戰哥利亞.
如果他不覺得.
耶華與他同在的話.
不是因為大衛有一個.
完全依靠神的心和交託的心的話.
如果不是大衛有一個.
無畏無懼的信心的話.
試問大衛又豈敢.
作出一個如此性命攸關的挑戰呢.
大衛出戰的信心來自耶和華.
他依靠都是來自耶和華.
蘇羅和以色列所見到的是.
哥利亞的勇猛.
而大衛見到的是.
哥利亞是耶和華的敵人.
耶和華必然會將他的敵人.
哥利亞擊敗.
因此我們就更加清楚知道.
神為什麼要捨棄蘇羅而揀選大衛.
蘇羅在第十三章師治憲制之前.
當時蘇羅也是同樣面對非利士人.
蘇羅所見到的是敵人有戰車三萬兩.
騎兵八千.
士兵好像海邊的沙一樣那麼多.
不過蘇羅見不到的就是.
耶和華是與他同在的.
他更加沒有想過依靠耶和華.
以至他急急地去憲制.
想挽回軍心.

$^{281}$然後以自己的能力去抵抗非利士人.
而在第十七章.
蘇羅見到的是哥利亞的巨大衛猛.
他見不到的是哥利亞是耶和華的敵人.
蘇羅依靠自己的能力去抵抗非利士人和哥利亞.
所以縱然他身邊有很多的士兵.
他仍然感覺到恐懼.
而大衛依靠的是萬軍的耶和華.
所以縱然他矮小.
他穿不起軍服.
只有他一個人.
但仍然可以無畏無懼.
大衛的信靠使他的屬靈視野.
能夠擺脫他自身的限制.
他有一種屬靈的洞見.
勇氣,無懼的心.
以及為耶和華爭戰的信心.
這正正就是蘇羅和大衛在屬靈生命和視野上的分別.
亦都是耶和華捨棄蘇羅.
而揀選大衛的原因.
當然有人會替蘇羅抱不平.
蘇羅縱然有犯錯.
但大衛沒有做錯嗎?.
為何耶和華對蘇羅不可以寬容一點呢?.
對大衛又處處寬容.
蘇羅私自獻祭.
以及滅絕亞馬力族這兩件事.
不單沒有依靠神的心.
甚至不願意承認自己所犯的罪.
只是一味地推卸責任.
大衛一生犯過很多罪.
最嚴重的我們所記得的一定是謀殺了烏尼亞.
不過大衛願意在耶和華面前認罪悔改.
承擔罪責.
這就是耶和華喜悅大衛的原因.
大衛人生不完美.
但大衛對神的那種倚靠和忠誠.
正正是大衛對神的倚靠.
他知道縱然在神手裡接受懲罰.
但在神手裡接受懲罰仍然有愛.

$^{321}$弟兄姊妹.
縱然我們在神面前認罪悔改.
接受那種的刑罰.
但在耶和華手裡的懲罰仍然有愛.
因此正正是這種屬靈的洞見.
使大衛在犯錯的時候.
他願意回轉認罪悔改求赦免.
這也是我們今天弟兄姊妹應該要學習的功課.
做錯要承擔責任.
但神的懲罰裡仍然有愛.
各位同學弟兄姊妹.
今天講的是神如何看待我們.
如果神看待我們是根據我們的內心世界.
我們心裡的動機.
而不是看待我們那種熟悉的身份地位能力以及財富的話.
我們真的需要停一停想一想.
我們今天所追求的.
是不是和神的要求有一種偏差呢?.
神要求我們的心靈深處有一顆單純的信靠的心.
而這種單純的信靠的心.
可以讓我們每時每刻遇到什麼困境危難的時候.
能夠有一種屬靈的視野.
可以無畏無懼有勇氣地去面對.
這才是神看待我們的標準.
這才是神揀選我們為祂做事事奉的標準.
我們又要再問.
究竟我們的心是像大衛嗎?.
因為我們的心是像素羅.
如果我們覺得自己像素羅.
凡事都只是靠自己的聰明才智.
而不是依靠神的話.
或者我們臨急抱佛腳.
有需要的時候求神去幫的.
那麼基督信仰對我們有什麼作用呢?.
或者我們再回顧一下.
素羅一個人生去做一個反思.
就是縱然你很甘願被神所使用.
你很想被神使用.
但是都沒有一個好結果.
一種不符合神要求的心態去事奉.

$^{361}$一種沒有神參與同行的事奉.
最終都不會被神所使用.
就好像素羅一樣.
如果我們好像大衛一樣.
很想為神的命被污辱而出頭.
為神而戰.
但是我們問一問自己.
我們有沒有大衛那麼勇敢.
我們不敢站出來在強敵面前不屈服.
還可以無畏無懼.
我們心裡面很想.
但現實上又做不到.
那怎麼辦呢?.
可能你會說.
Gary,世界不是那麼易圓的.
不一定要在素羅和大衛之間去選擇.
生命還有很多選擇.
是的,我們的人生不一定要像大衛一樣.
神給每一個人的人生都很獨特.
所以我們要打造我們自己的人生.
開始的引言我們不是說了嗎?.
我們要為自己打造一個形象的工程.
因為我們知道.
我們有一個好的形象.
我們就有一個好的開始.
關鍵是我們怎樣去打造.
我們自己的形象工程呢?.
因為我們知道.
神不看我們的外表.
神看我們的內心.
從大衛的生命我們知道.
他的屬靈生命是建基於信靠.
然後在信靠的經驗裡面.
產生勇氣.
產生信仰的眼界.
產生為神爭戰的信心.
因此信靠神是一個基礎.
在這個基礎上再打造其他屬靈的品格.
我們看看大衛.
大衛早期的屬靈生命.

$^{401}$是從每一日日常生活裡面建立.
那種信靠神的心.
大衛年輕的時候.
每日都要為父親去放羊.
其實放羊的生活是很枯燥的.
但是每當遇見獅子或者熊的時候.
就能夠激發那種信靠神的心.
大衛是從一種平淡和枯燥的生活裡面.
建立那種信靠神的心.
大衛被高立之後.
他的生命也不是順風順水的.
他雖然打敗了哥利亞.
但是在他漫長的人生裡面.
還有很多很多的挑戰.
由於素羅認為大衛會威脅他的王位.
於是追殺他.
大衛需要四處漂泊.
他有兩次的機會可以殺死素羅.
他仍然放棄.
因為他信靠神.
他認為素羅是神所高立的.
他沒有權去殺素羅.
他相信神會保守看顧他.
所以縱然要顛沛流離.
他都不肯去殺害神所高立的素羅.
大衛這段屬靈的之旅.
寫了很多首詩.
去描述他的心路歷程.
因此大衛生命裡面各種的屬靈果實.
都是從他每一日的生命一點一滴去累積而成.
而信靠神的心.
也必定是從我自身的經歷開始和累積的.
不能天降而來的.
也沒有任何的捷徑.
是要我們親自的經歷.
然後內化在我們生命裡面.
因此我們形象工程的包裝的用料.
不是華貴的衣服.
不是我們的身份.
不是我們的地位.

$^{441}$不是我們住在哪裡.
更加不是我們的財富.
因為這一切一切都不能夠叫我們.
在面對逆境的時候.
可以有一份無畏無懼的心.
事奉主的人從信靠神產生一種.
產生不同的屬靈果子.
為自己打造一個獨特的形象工程.
不同的屬靈果子產生信仰.
生命的香氣.
讓其他的人被這些香氣所吸引.
最後歸入主的名下.
一切的基礎都是從信靠主開始.
各位姐妹.
我們問一問我們自己.
我們在我們人生的經歷裡面.
我們有多少時候是信靠主.
我在一月尾被主高立成為一個牧師.
在寫這篇講章的時候.
我一路寫一路反思.
今天這個講題.
神怎樣看我.
究竟我是否能夠好像講題一樣.
我有沒有講到裡面那種信靠主的心.
我的屬靈香氣是怎樣的呢.
我自己那個獨特的形象工程究竟完成了沒有呢.
我很想和大家分享.
我信靠主的心有兩個階段.
第一個階段是在我信主的初期.
我記得我當時信主之前有很多壞習慣.
吸煙喝酒賭錢這些是我的常態.
有一段時間我很想戒掉.
用了很多方法去戒掉但戒不到.
而且我覺得我自己在享受那種最終之樂.
我根本不想戒掉.
也不知道多麼享受.
直至我信了耶穌.
我信主初期當然都好像一般的信徒.
參加崇拜參加主日學.
我還記得那時候我參加完崇拜之後.

$^{481}$等著回主日學.
我就在後樓梯吸煙.
不止吸一支還要吸兩支.
因為吸兩支才能治癮.
我以為別人不知道.
吸了煙就回主日學.
別人以為不知道其實別人好心不戳爆我.
我知道主不喜歡我這種人格分裂式的生活.
我在主面前祈禱.
我求主加能賜力給我.
讓我能夠戒除這個習慣.
就是這樣.
憑著這種單純去信靠主的心.
我就戒除了我的壞習慣.
成為一個新做的人.
對於當時一天吸兩包煙的我來說.
我視為是奇蹟.
一天吸兩包煙.
賭錢是我正職.
不賭錢是奇蹟.
這個經歷給我一個很大的信心.
讓我能夠在淑齡階段裡面有一個很好的成長.
第二個信靠主的時期是讀神學的階段.
當時只靠太太去做一份傳道人的薪酬.
初初做傳道人的薪酬都不會很高.
我會靠太太做傳道人.
去支持我們一家四口的開支.
不過這次信靠主的經歷不是我一個人面對.
而是我和太太一起去經歷.
我們一起見證我們的家庭.
在我們讀神學這四年期間.
神對我們的供應.
當時我的女兒曾經說過.
她說:爹地媽咪,你不是說家裡沒有錢嗎?.
但為甚麼我不覺得呢?.
為甚麼我覺得一切都有,一無所缺?.
單純的信靠,神是會供應的.
這種信靠主的心慰.
我產生了幾種淑齡的香氣.
我都頗確定,我自己有這種淑齡的氣質.

$^{521}$第一就是毀身.
我從讀神學開始就立志終身事奉祂.
所以是否安立我做牧師.
其實我都是會終身事奉祂.
不過當然安立了牧師的時候.
過程裡我更加知道.
我一生都要被祂所使用.
因為我深信神昔日既然去扶持我.
現在都在保守我.
為甚麼我要擔心將來神不會看顧我呢?.
神是昔日,昔在今在永在的主.
第二個我覺得我自己有的淑齡香氣就是忠心.
我知道我自己有很多缺點.
亦知道自己有很多限制.
但我相信我只要忠心的話.
神仍然會樂於使用我.
第三是憐憫.
其實認識我的人都知道.
我的性格很硬.
但神讓我在事奉裡面.
讓我能夠接觸很多人.
神讓我有機會進入別人的生命裡面.
有時當我聽到別人的生命故事的時候.
我心裡面有一種憐憫.
特別是在紅印格.
我有很多機會接觸過來留的新青人.
每一次跟他們聊天.
我都覺得我們一定要做這個事奉.
第四就是我是一個很有承擔的人.
我是一個不怕辛苦的人.
我覺得做事是不會做死人的.
所以我就願意去做一些開方的工作.
而開方的工作是需要很大的承擔.
你不能夠因為遇到一些挫折.
然後就說我放棄了.
或者我逃避.
那我的形象工程完成了沒有?.
當然還沒有.
我仍然在自己裡面繼續學習.
例如好像謙卑.

$^{561}$我不夠謙卑.
我不夠溫柔.
沒有那種同理心.
急躁.
要放慢腳步.
事實上形象工程就好像成聖之類一樣.
是需要一生去打造的.
甚願我日後的生命是為別人帶來祝福.
而不是令人跌倒,失腳.
或者是福音的難阻.
弟兄姊妹.
我們會不會一起去打造一個信靠主的心?.
為主結出一個雷雷的果子.
好不好?.
我們一起去祈禱.
你在罪裡面拯救了我們.
讓我們行在天國之裡的當中.
求主加能賜力.
讓我們能夠更加的信靠你.
從信靠中結出不同屬靈的果子.
一生被主所使用.
榮耀主名.
我們的祈禱是奉靠主耶穌基督.
聖名祈求.
阿們.
願主祝福大家.
\newpage



\section{約翰福音 1:1-28}
\label{sec:2_wlsYE_36Q}
\textbf{2021.03.14《教會存在於見證中》約翰福音 1\_1-28( 和修版)  陳韋安博士}
\newline
\newline
連結: \href{https://youtube.com/watch?v=2_wlsYE-36Q}{\texttt{https://youtube.com/watch?v=2\_wlsYE-36Q}} ~~~~ 語音日期: 2021-03-18
\newline
\newline
\hyperref[sec:JLt4gSCI6Y8]{\small{< < < PREV SERMON < < <}}
~
\hyperref[sec:index_chronic]{\small{[返順時目]}}
~
\hyperref[sec:index_scriptual]{\small{[返順卷目]}}
~
\hyperref[sec:oH_gH6L8AqQ]{\small{> > > NEXT SERMON > > >}}
\newline
\newline
約翰福音 1:1-28
\newline
\begin{longtable}{cl}
\hline
\hline
章節 & 經文 (和合本修訂版)\\
\hline
1:1 & \begin{tabularx}{0.7\textwidth}{X} 太初有道，道與神同在，道就是神。 \end{tabularx} \\ \\ \relax
1:2 & \begin{tabularx}{0.7\textwidth}{X} 這道太初與神同在。 \end{tabularx} \\ \\ \relax
1:3 & \begin{tabularx}{0.7\textwidth}{X} 萬物都是藉著他造的，沒有一樣不是藉著他造的。凡被造的， \end{tabularx} \\ \\ \relax
1:4 & \begin{tabularx}{0.7\textwidth}{X} 在他裡面有生命，這生命就是人的光。 \end{tabularx} \\ \\ \relax
1:5 & \begin{tabularx}{0.7\textwidth}{X} 光照在黑暗裡，黑暗卻沒有勝過光。 \end{tabularx} \\ \\ \relax
1:6 & \begin{tabularx}{0.7\textwidth}{X} 有一個人，是從神那裡差來的，名叫約翰。 \end{tabularx} \\ \\ \relax
1:7 & \begin{tabularx}{0.7\textwidth}{X} 這人來是為了作見證，是為那光作見證，要使眾人藉著他而信。 \end{tabularx} \\ \\ \relax
1:8 & \begin{tabularx}{0.7\textwidth}{X} 他不是那光，而是要為那光作見證。 \end{tabularx} \\ \\ \relax
1:9 & \begin{tabularx}{0.7\textwidth}{X} 那光是真光，來到世上，照亮所有的人。 \end{tabularx} \\ \\ \relax
1:10 & \begin{tabularx}{0.7\textwidth}{X} 他在世界，世界是藉著他造的，世界卻不認識他。 \end{tabularx} \\ \\ \relax
1:11 & \begin{tabularx}{0.7\textwidth}{X} 他來到自己的地方，自己的人並不接納他。 \end{tabularx} \\ \\ \relax
1:12 & \begin{tabularx}{0.7\textwidth}{X} 凡接納他的，就是信他名的人，他就賜他們權柄作神的兒女。 \end{tabularx} \\ \\ \relax
1:13 & \begin{tabularx}{0.7\textwidth}{X} 這些人不是從血生的，不是從情慾生的，也不是從人的意願生的，而是從神生的。 \end{tabularx} \\ \\ \relax
1:14 & \begin{tabularx}{0.7\textwidth}{X} 道成了肉身，住在我們中間，充充滿滿地有恩典有真理，我們也見過他的榮光，正是父獨一兒子的榮光。 \end{tabularx} \\ \\ \relax
1:15 & \begin{tabularx}{0.7\textwidth}{X} 約翰為他作見證，喊著說：「這就是我曾說：『那在我以後來的先於我，因為在我以前，他已經存在。』」 \end{tabularx} \\ \\ \relax
1:16 & \begin{tabularx}{0.7\textwidth}{X} 從他的豐富裡，我們都領受了恩典，而且恩上加恩。 \end{tabularx} \\ \\ \relax
1:17 & \begin{tabularx}{0.7\textwidth}{X} 律法是藉著摩西頒佈的；恩典和真理卻是由耶穌基督來的。 \end{tabularx} \\ \\ \relax
1:18 & \begin{tabularx}{0.7\textwidth}{X} 從來沒有人見過神，只有在父懷裡獨一的兒子將他表明出來。 \end{tabularx} \\ \\ \relax
1:19 & \begin{tabularx}{0.7\textwidth}{X} 這是約翰的見證：猶太人從耶路撒冷差祭司和利未人到約翰那裡去問他：「你是誰？」 \end{tabularx} \\ \\ \relax
1:20 & \begin{tabularx}{0.7\textwidth}{X} 他就承認，並不隱瞞，承認說：「我不是基督。」 \end{tabularx} \\ \\ \relax
1:21 & \begin{tabularx}{0.7\textwidth}{X} 他們又問他：「那麼，你是誰？是以利亞嗎？」他說：「我不是。」「是那位先知嗎？」他回答：「不是。」 \end{tabularx} \\ \\ \relax
1:22 & \begin{tabularx}{0.7\textwidth}{X} 於是他們對他說：「你到底是誰，好讓我們回覆差我們來的人。你說，你自己是誰？」 \end{tabularx} \\ \\ \relax
1:23 & \begin{tabularx}{0.7\textwidth}{X} 他說：「我就是那在曠野呼喊的聲音：修直主的道。」正如以賽亞先知所說的。 \end{tabularx} \\ \\ \relax
1:24 & \begin{tabularx}{0.7\textwidth}{X} 那些人是法利賽人差來的。 \end{tabularx} \\ \\ \relax
1:25 & \begin{tabularx}{0.7\textwidth}{X} 他們就問他：「你既不是基督，不是以利亞，也不是那位先知，那麼，你為甚麼施洗呢？」 \end{tabularx} \\ \\ \relax
1:26 & \begin{tabularx}{0.7\textwidth}{X} 約翰回答：「我是用水施洗，但有一位站在你們中間，是你們不認識的， \end{tabularx} \\ \\ \relax
1:27 & \begin{tabularx}{0.7\textwidth}{X} 就是那在我以後來的，我給他解鞋帶也不配。」 \end{tabularx} \\ \\ \relax
1:28 & \begin{tabularx}{0.7\textwidth}{X} 這些事發生在約旦河東邊的伯大尼，約翰施洗的地方。 \end{tabularx} \\ \\
[1ex]
\hline
\hline
\end{longtable}
$^{1}$各位兄弟姐妹,早安!很高興來到紅磡堂跟大家分享.
我也是第一次來到教會.
雖然很多時候都知道教會已經開了一段時間.
但因為很多原因都沒有機會來到當中.
來到的時候很開心.
開車來到洪水橋附近的時候.
看到很多不同的街坊.
看到很多茶餐廳,單車舖.
看到很多不同的附近的人.
很感恩教會在這個地方成為一間新的教會.
我以前也在宣德堂開展直堂.
所以也很特別,很體會一間新教會的存在.
或者是一個新意思,可能也不是十年以上.
一間教會在一個新的地方很有意思.
特別是疫情,在這一年我們來崇拜的時候.
讓我們更加思考教會在這個地方所做的事情的意義.
所以今天牧師選擇了一個有關教會的主題的時候.
我也用這段經文跟大家一起來思考.
《約翰福音》第一章一到十八節是一段神學引言.
很多人都建議初信主的人來讀《約翰福音》.
除了這一段,這一段很深奧.
這一段有很多神學的信息.
太初有道,對神同在道,徒至神.
說到永恆的上帝怎樣來猜險,讀生子.
上帝是一個完全的人,完全的神.
這些很深奧秘,三維一體,永恆之類的話題.
可能大家也未必這麼容易理解這段經文.
不過這段經文裡面其實是很特別的.
第六節是加插了一個人進去.
你會發現一段這麼神聖,關乎於上帝永恆,奧秘的經文.
竟然是加插了一個平凡的人在當中.
第六節裡面加入了施瑟若翰這個人的故事.
在整個神學引言裡面.
一個甚至是一個很粗鄙,很市井的男人.
滿口鬍鬚,然後他那個人基本上都不是怎麼穿衣服.
穿著一些駱駝毛的皮,到處走.
吃的東西是蝗蟲和野蜜.
一個很市井,很街坊的形式.
竟然是加插在一段這麼神聖,上帝請求的神學引言當中.
就好像港台裡面有一本聖經.

$^{41}$聖經裡面突然間翻開就發現夾了一本TVB雜誌的字.
好像很不配合.
一個這麼莊嚴,這麼神聖,太初有道的經文.
突然間加了一個人的故事在當中.
這個正是我們今天來思考教會的存在.
教會在洪水橋這個地方的意義.
當我們來看施瑟若翰的自我介紹的時候.
你會發覺更加知道這個重點.
第十九節開始一班的祭司和一班的利未人去找施瑟若翰.
當他們問一個問題.
你是誰?你是誰?的時候.
施瑟若翰有一個很奇怪的回答.
他用了一個很特別的方式來回答這個問題.
他不是回答他是誰.
他用了一個很否定的方法來回答.
他說:我不是基督,我不是耶穌基督.
很奇怪的,當別人問你是誰的時候.
你就回答自己的名字,其實他又不是.
他竟然回答:我不是基督.
就好像你認識一個新朋友.
你問他是誰,他就回答:我不是羅德華.
整件事不是很奇怪.
我知道你不是羅德華,但你是誰呢?.
所以若翰其實是故意地用了幾次都是這樣說.
去到後面有一節更加來強調再否定.
他說:我不是以利亞,不是先師.
所以在見到這個若翰的三次的自我介紹裡面.
他說:我不是基督,我不是以利亞,也不是先師.
三次的一個否定式的語句.
來作出一個自我介紹.
當然這個介紹是故意的.
若翰作為一個耶穌基督的見證人.
他的這個自我介紹.
是要去點出作為見證者一個最首要的身份.
那個最首要的身份不是他是誰,而是他不是什麼.
一個見證耶穌基督的人.
他首先要知道他不是什麼.
他不是耶穌基督.
這個是他最首要提醒外面的人.
也提醒了自己,一個這樣的見證.

$^{81}$這個見證正正來點出.
今天我們教會在地上.
無論在洪水橋,香港,普世都一樣.
我們教會的處境.
講一下我以前在德國宣教事奉的經歷.
基本上我們以前在德國事奉的時候.
德國是一個基督教國家.
不過當中基督教其實是很弱勢的.
那個勢很弱.
為什麼說勢很弱呢.
如果你看他的市集的時候.
有時候他見不到市集裡面有很多不同的人.
賣衣服,賣香腸,賣不同的東西.
印象很深刻的是.
那時候有一個男人,那時候是柏林冬天.
很冷,很暗暗的地方.
很記得一個男人.
有一個木打出來的攤檔.
上面就放了幾本福音冊子.
很記得那個男人的樣子.
沒有人理他,整個市集裡面.
很熱鬧,但是沒有人理這個人.
然後這個男人.
他好像連自己都不理自己.
他自己都很憔悴地站在攤檔那裡.
前面就放了幾本德文的福音冊子.
這個正正就是在德國裡面教會的形勢.
勢很弱,應該說是.
人們不回教會,人們都不太理教會的事情.
當時我們都很想在德國買一個十字架.
找了很久,原來在德國找十字架都不容易.
雖然是基督教國家.
香港很明顯是以前去宣傳書局,或者去基督教.
但是原來在德國找十字架都不容易.
我們找了很久都找不到.
今天的整個世界正正.
如果你去到更加遠的地方.
你會發現教會信仰的力量越來越減少.
我想香港也是這樣.
大家特別在這一年裡面.

$^{121}$你會發現教會在這個社會裡面.
往往是一個很渺小的情況.
不知道大家在洪水橋裡面會怎麼看.
當我們在這個社區裡面的時候.
我們這個地方是一個什麼樣的位置呢.
大家外面的街坊,外面的微信主人.
他們怎麼看這個地方呢.
可能他們會很關心茶餐廳,單車舖.
多於我們這個地方.
人們都不會特別去想教會這件事.
所以這個正正就是我們今天教會的情況.
就像約翰一樣.
在一個極大的曠野裡面.
不斷地大聲呼喊.
你們要悔改,你們要信耶穌.
不過對於很多人來說.
我們的聲音雖然聲色力竭.
但其實都是一個很微弱的聲音.
所以約翰第一句的自我介紹.
他說我不是基督.
這個點名出一個很重要的事實.
約翰不是基督.
雖然他用水幫人去施洗.
大聲叫人去認罪悔改.
但他不是基督.
雖然他滿口耶穌.
但他絕對不是基督.
教會也一樣.
教會雖然說的是永恆很大的道理.
但我們絕對比不上我們的道理.
我們的道理絕對比我們自己更加大.
我們承載的寶貴的雅氣裡的寶貝.
比我們本身更加大.
所以我們所說的永遠都大於我們自己.
我們沒有特殊的能力.
我們沒有特殊的神蹟.
我們是社會上的少數派.
但我們所說的是這樣的福音.
所以約翰更加清楚地指明自己的身份.
第23節裡他強調.

$^{161}$他開始說出一個比較正面的自我介紹.
第23節他說.
我就是在那曠野有人大聲喊著說.
修植著的道路.
正如先知以徐也所說的.
第23節的時候.
約翰終於說出一個比較正面的自我介紹.
他不是基督不是以利亞不是先知.
他是甚麼呢.
他說我就是一個在曠野當中大聲喊著說.
修植著道路的人.
甚麼是修植著道路.
其實如果看以徐也書的時候.
這個就是為主預備道路.
就是將道路修整和平衡.
簡單來說就是一個預備的功夫.
預備功夫永遠都是最難做的.
不知道大家有沒有這樣的感覺.
就好像一個演唱會裡面.
如果你看一些外國的演唱會裡面.
演唱會裡面有一個開場的樂隊表演.
就是開場裡面找一些其他人來熱身.
那個不是主角.
只不過是開場裡面熱身一下就走了.
或者更加貼切就是在婚禮裡面的伴郎.
我們英文叫做best man.
伴郎.
婚禮裡面的伴郎永遠都是這樣.
他稱之為best man.
最好的人.
他要全心全意盡力做到最好.
不過你永遠知道你不是主角.
你所做的事永遠都是做一個配角的角色.
當日裡面沒有人特別focus在你裡面.
不過你要做到最好.
他為新郎新娘預備一切.
然後就黯然地讓人去光輝.
所以約翰這樣說.
我必興旺,我必須微.
他必興旺,我必須微.

$^{201}$這個正正就是約翰自己作為一個教會.
作為一個見證基督人.
最首要的身份.
教會離開了他的使命.
他就一點都不是.
當教會離開了這個使命的時候.
他就一點價值都沒有.
今天我們不斷地去講耶穌.
在這個地方裡面向人們傳揚福音.
是有一些寂寞的感覺.
就好像一個曠野裡面不斷地去呼喊.
不斷地去推銷一個信仰的那樣的人.
所以你看到教會正正是一個面對著這樣的難題.
不知道大家在這幾年裡面.
在這個地方裡面去開教會.
去傳福音有什麼感受.
正正是一個我們很用力去做.
好像都得不到很大的回報一樣.
我自己也是一樣.
以前我在柏林,德國教會柏林的一間華人教會裡面事奉.
都是一個三十幾四十人的群體.
我們用盡很多的方法.
就是全心全意.
我們見到開放吃苦火熱.
去耗盡我們自己.
但發覺都是很不容易.
永遠都是在整個社會裡面很少數的人.
可能大家在香港裡面不覺得.
香港是只有4\%基督徒.
你想想4\%是什麼意思.
100裡面有4個人是基督徒.
其餘96人都不是基督徒.
大家可能不太覺得這個數字.
大家平時這裡都是100\%.
大家在這裡回來.
你的信主朋友永遠都多過你不信主人.
順便你覺得不是只有這麼少.
你身邊的,平時飲茶的.
你Facebook裡面的,你WhatsApp裡面的.
都是基督徒.

$^{241}$但其實原來那個數字只有4\%.
100個裡面其實只有4個人是基督徒.
其實我們永遠都是在一個.
不是永遠,就是一直都是在一個很邊緣的數字裡面.
香港這個社會對於基督教.
教會其實一直都是在一個很邊緣的數字裡面.
所以我們平時不覺得.
但我們原來從來都是在一個這樣的位置裡面.
在整個世界裡面.
我們一直都是一個很少數的群體.
但我們所傳揚的正正是對著這個世界.
以前我都做過幹事.
在書店裡面.
禮拜一幹事要做些什麼.
要做入數.
將教會的奉獻來到銀行.
去匯豐入數.
那時我去元朗大馬路的匯豐裡面.
將教會奉獻入數.
當我將教會神聖的金錢拿去匯豐的時候.
有很大感覺.
你知道禮拜一匯豐是一個很多人排隊的地方.
就是下午吃飯的時間.
很多不同的商舖.
都一起入數.
除了教會之外.
有街市的嬸嬸.
五金舖的叔叔.
很多不同的人都一起去匯豐入數.
當我排隊的時候有很大感覺.
我們教會的那筆錢.
那些支票的錢.
那些聖供的錢.
其實只不過是很多很多不同行業裡面.
其中一員.
當你這樣看的時候.
發覺原來教會是很渺小很渺小的.
在整個社會裡面.
我們只不過是其中一個一員裡面.
這份感覺正正是我們面對的處境.

$^{281}$雖然我們所見證的是一個世界裡面的真光.
但我們所面對的是一個極大的曠野.
當然故事並未完結.
故事去到最後的時候.
約翰福音第一章第26到27節.
約翰這樣來到.
他清楚地再一次認定.
他自己的身份.
他說我是用水去施洗.
但有一位站在你們中間.
是你們不認識的.
就是那在我以後來的.
我給他解鞋大一不配.
有一位站在你們中間.
是你們不認識的.
去到整段說話的最後.
約翰在這班人面前.
作出一個見證.
一個有關耶穌基督的見證.
一個告訴其餘90人聽.
基督就在你們中間的見證.
事實上.
這正正是約翰開口說話的原因.
也是他大聲呼喊的原因.
也是他存在在地上的唯一原因.
他要向世界作出一個見證.
耶穌基督正正就在我們中間.
正如第14節所說.
到成了肉身住在我們中間.
充充滿滿有恩典有真理.
這件事已經是出現了.
發生了.
我只不過是來告訴你.
去見證這件事.
一開始說到的時候.
我們問為何一個這麼神聖的神學隱言.
竟然會出現一個這麼市井普通平凡的人.
原因只有一個.
這個人正正是要為這個真光來作見證.
我們是世界上的光.

$^{321}$而基督的真光.
他自己是會土照這個世界.
我們所做的只不過是為這個真光.
來作一個很簡單不配的見證.
他不是那光.
乃是要為真光作見證.
這種這樣的身份.
這種這樣的狀態.
正正是我們教會群體要做的事情.
雖然今天我們面對著很多很多很困難的難題.
我想每一間教會.
都面對著或多或少差不多的難題.
在今年裡面.
崇拜連崇拜都回不了.
我們搞很多的班.
很多人都參與不了.
或者很多很多的不同的事工都停滯了.
我們好像各自寸步難行.
不過我們仍然要謹記.
我們正正生下來就是這樣的處境.
我們從來都不是那個真光.
我們都是一個很平凡的人.
但我們卻願意去見證這個真光的時候.
這個就是我們教會存在的唯一的目的.
所以我想洪恩堂或者這間教會.
或者任何地上的教會.
面對著這一年的難題.
2019年2020年面對著這麼多這麼多難題的時候.
其實讓我們更加可以去無忘初心.
來思考我們自己存在的目的.
我們的存在所謂基督徒.
從來都不是純粹一班不救群體.
基督徒這個字.
本身都是解著一個有召命的人.
基督徒正正是帶著一個使命.
而去存在的一班人.
如果我們沒有了我們的使命.
我們基督徒的字就只不過是illusion.
一個幻象.
所謂mission這個字.

$^{361}$所謂使命這個字.
很多人以為這個解作宣教.
或者是解作某些猜拳的活動.
但mission這個字正正是我們教會最首要的目的.
我們是帶著一份使命來做事情.
以前在德國裡面.
有些德國人將mission這個字解作mission.
這個字其實不是只解作宣教.
而是解作在地區裡面的關愛.
以前在德國裡面.
教會會幫那些露宿者.
去到幫他們有個地方讓他們睡覺.
一個避寒中心.
對於德國教會來講.
mission這個字不是單單解作跨文化的宣教.
而是我們在整個地區裡面.
所作出的最基本的見證.
我們有很多不同的部.
傳道部,猜拳部,社關,示威.
其實很多不同部門.
最根本最做的事情只有一個.
就是見證耶穌基督自己.
傳道部是本地的猜拳.
海外的社關就是關心社會的.
其實沒有分別.
我從來都是單單去做一件事情.
教會上面的事情都是做見證耶穌基督.
無論是近方的洪水橋的社區.
或者香港的社會.
或者遠方的不同地方.
我們仍然留心不同人的需要.
去見證耶穌基督.
這一年裡我們忘記了很多我們已有的做法.
去捉緊我們的初心.
紅恩堂或者這間教會.
正正是為了見證基督而存在.
很多的事工,很多的事情.
都是圍繞著這件事情去做.
讓我們能夠去捉緊這個任務.
任務更加是我們最首要的.

$^{401}$是我們自己心靈裡面的使命.
帶著這份使命去做基督徒.
永遠是首要的.
比我們所說要做的事工更重要.
讓我們被延起這個使命.
這個不是純粹教會的使命.
而是作為基督徒.
當你信了耶穌之後.
耶穌基督呼召你去完成見證的使命.
這個不是屬靈的高層次的事情.
而我們每個基督徒首要地.
就已經開始被耶穌基督呼召去做的事情.
這個才是基督徒的定義.
基督徒就是一班見證耶穌基督的人.
讓我們在這個難關裡面.
重新去思考這件事情.
去審視我們的架構.
審視我們的事工.
審視我們自己的思維.
好讓我們能夠去.
將見證基督這件事情.
放在洪水橋的地方裡面.
讓這間教會能夠成為一個.
被上帝使用的教會.
我們一起祈禱.
我們祈求你去幫助.
因為你自己是世界的真光.
讓我們能夠在地上.
能夠去迎迎的去見證你.
讓我們教會裡面每個人.
能夠去校正我們的方向.
忘記一些我們已經有的事情.
單單純粹去思考.
如何能夠在這個世代.
在這個年裡面.
去重新將你自己見證給其他人聽.
求主你這樣的幫助.
我們所做的不是單單圍內.
自己崇拜.
而是去到將你的福音.

$^{441}$將你自己見證給其他人.
求主你這樣去感動我們.
讓我們能夠每一個人.
就算在家中裡面崇拜的時候.
仍然記得我們怎樣.
在生活裡面去見證你.
讓我們能夠成為你的美好的見證.
奉主明求.
阿們.
\newpage



\section{以賽亞書 6:1-7}
\label{sec:oH_gH6L8AqQ}
\textbf{2021.03.21《聖哉!聖哉!聖哉!》以賽亞書 6\_1-7( 和修版)  曾敏姑娘}
\newline
\newline
連結: \href{https://youtube.com/watch?v=oH_gH6L8AqQ}{\texttt{https://youtube.com/watch?v=oH\_gH6L8AqQ}} ~~~~ 語音日期: 2021-03-25
\newline
\newline
\hyperref[sec:2_wlsYE_36Q]{\small{< < < PREV SERMON < < <}}
~
\hyperref[sec:index_chronic]{\small{[返順時目]}}
~
\hyperref[sec:index_scriptual]{\small{[返順卷目]}}
~
\hyperref[sec:7LBy_VBf8iw]{\small{> > > NEXT SERMON > > >}}
\newline
\newline
以賽亞書 6:1-7
\newline
\begin{longtable}{cl}
\hline
\hline
章節 & 經文 (和合本修訂版)\\
\hline
6:1 & \begin{tabularx}{0.7\textwidth}{X} 當烏西雅王崩的那年，我看見主坐在高高的寶座上。他的衣裳下襬遮滿聖殿。 \end{tabularx} \\ \\ \relax
6:2 & \begin{tabularx}{0.7\textwidth}{X} 上有撒拉弗侍立，各有六個翅膀：兩個翅膀遮臉，兩個翅膀遮腳，兩個翅膀飛翔， \end{tabularx} \\ \\ \relax
6:3 & \begin{tabularx}{0.7\textwidth}{X} 彼此呼喊說：「聖哉！聖哉！聖哉！萬軍之耶和華；他的榮光遍滿全地！」 \end{tabularx} \\ \\ \relax
6:4 & \begin{tabularx}{0.7\textwidth}{X} 因呼喊者的聲音，門檻的根基震動，殿裡充滿了煙雲。 \end{tabularx} \\ \\ \relax
6:5 & \begin{tabularx}{0.7\textwidth}{X} 那時我說：「禍哉！我滅亡了！因為我是嘴唇不潔的人，住在嘴唇不潔的民中，又因我親眼看見大君王－萬軍之耶和華。」 \end{tabularx} \\ \\ \relax
6:6 & \begin{tabularx}{0.7\textwidth}{X} 有一撒拉弗向我飛來，手裡拿著燒紅的炭，是用火鉗從壇上取下來的， \end{tabularx} \\ \\ \relax
6:7 & \begin{tabularx}{0.7\textwidth}{X} 用炭沾我的口，說：「看哪，這炭沾了你的嘴唇，你的罪孽便除掉，你的罪惡就赦免了。」 \end{tabularx} \\ \\
[1ex]
\hline
\hline
\end{longtable}
$^{1}$領子妹平安.
願意神讓我們今天大大的得著祂.
我們一同再次禱告.
讓我們仍然有生命氣息去親近你.
願意今天你開我們的眼睛.
讓我們看見你的榮耀.
看見你的奇妙與聖潔.
就讓我們將心歸給你.
願聖靈在我們每一個心裡面去工作.
感動我們 回應神.
奉耶穌基督的名字 祈禱 阿們.
今天的題目叫做「盛哉盛哉盛哉」.
這三個詞語是一個詞語的重複.
不停地說「盛哉盛哉盛哉」.
其實是一個三頌讚.
例子一段敬拜讚美的說話.
一個讚美的詞語重複了三次.
就是說很重要.
盛哉是指聖潔.
如果我們把它再轉一轉.
盛哉 盛哉 盛哉.
代表了這個盛哉是最高級別.
最盛哉的意思.
這個三重的頌讚是發自撒拉弗的口.
撒拉弗是指天使.
經文裡這樣說.
主坐在高高的寶座上.
這個寶座應該是在天上的聖殿裡.
有天使停留在神之上的空中.
在那裡侍候神.
天使向神發出了敬拜讚美.
一連三次讚美神的聖潔.
他們讚美神是最聖潔的那一位.
什麼叫做最聖潔呢?.
這些詞語我們在教會聽得多.
說得多 其實都很抽象.
我現在叫你形容一下.
什麼叫做最聖潔呢?.
那怎麼形容呢?.
我們試試拿一些形容詞.

$^{41}$是不是像一張紙一樣.
白白的 這樣叫聖潔.
但紙的白色.
都有很多不同的程度.
有時候感覺上看起來.
叫白紙 也不是那麼白.
又或者不行 用雪.
雪一樣的白.
聖潔就是這樣.
好像雪一樣的.
是不是這樣叫聖潔呢?.
又或者我們用太陽的光芒.
會不會近一點.
去形容神的聖潔 行不行呢?.
當太陽的光照著我們的眼睛的時候.
大家有沒有試過這個經驗?.
實際你的眼睛.
可不可以看著太陽?.
看不到的.
太陽的光真的直接照著你的時候.
你會覺得很刺眼.
看不到 避開.
或者閉上雙眼.
是不是這種感覺.
很光芒四射的感覺.
是不是聖潔呢?.
我想是很難的.
我們的詞語.
其實是很難表達給我們聽.
什麼叫最聖潔呢?.
怎麼想像呢?.
真是很難.
但是神的聖潔.
是讓我們看到神自己的榮耀.
在六章一到七節.
短短幾節的經文裡.
我們看到不少文字.
是形容神的榮耀.
這些文字也帶來了一個對比.
我們現在和大家看看這些文字.

$^{81}$六章一節這樣說.
當烏西亞王崩的那年.
我看見主坐在高高的寶座上.
烏西亞是南國猶大國的一代明君.
他正職彪炳.
在位期間成功收復了大部分的失地.
恢復了大衛和所羅門王時代南部的版圖.
他很努力發展國家的經濟.
令國家繁榮昌盛.
到了晚年的時候.
他因為驕傲.
強行進入聖殿燒香.
而得罪了神.
他最後的結局是.
生大麻風到死亡的那一天.
雖然烏西亞是晚節不保.
但他在以色列民心目中.
是一個很重要.
很有影響力.
很厲害的一個君王.
他的失敗和離世.
是令人感到很唏噓的.
我們會想到.
無論地面上的君王有多出色.
都難保自己不會失敗.
不會跌倒.
無論地面上的君王有多成功.
他的管治都會有結束的一天.
我們的神都是君王.
為何這樣說呢?.
經文有一個對比.
六章一節這樣說.
主坐在高高的寶座上.
這個寶座也是形容君王的寶座.
只不過是在天上.
六章一節繼續這樣說.
他的衣裳下擺.
這是修訂版.
遮滿聖殿.
衣裳下擺代表了神的掌權.

$^{121}$他的榮耀.
神完全掌管天上的聖殿.
他的榮耀充滿了這個地方.
六章五節又這樣說.
由於我親眼看見大君王.
萬君之耀和華.
他的榮光遍滿全地.
全地都是屬於上帝的.
不只是以色列國.
全世界每一個國家.
每一個地方都是屬於神的.
神在當中掌權.
神所掌管的地方遠遠超過地面上的君王.
而且這位大君王永遠不會改變.
由過去到現在.
直到將來他都掌權.
他不會失敗.
他不會跌倒.
他的管治是直到永遠.
以色列民一直都很渴望有君王帶領他們打仗.
去打敗外敵.
去發展國家經濟等.
但是上帝很想他們知道.
地面上的君王有很多限制不足.
其實他們是不能完全滿足以色列民的心願.
你看看天上的君王.
他滿有能力,慈愛,有智慧.
他帶領以色列民就好得無比.
上帝很想透過這些經文.
透過這個對比.
希望百姓真心真意跟從他.
以他為王.
這並不是代表我們不需要地面上的君王.
而是神要我們更加體重他自己.
體重天上的君王.
今天我和你又如何呢?.
我們最體重的是哪一個?.
我們覺得哪一個最能帶領我們.
走人生的路.
面對生命的挑戰.

$^{161}$我們所倚靠的是哪一個?.
在這裡我有一個經歷很想和大家分享.
最近全城熱話是說打疫苗.
有很多討論在網上.
不停說打不打.
打又怎樣不打又怎樣.
打這個又怎樣打那個又怎樣.
不停討論到處都是這樣說.
很特別的是本來我自己.
是想不太早打疫苗的.
但是家人就想快點打.
而且指定是某一個.
那就好我們共同進退.
無論如何我和他們一起打.
做了這個決定之後.
就告訴了我的朋友.
就有很多聲音.
你打這個?.
這個有這麼多問題.
你不知道嗎?.
數據是這樣的.
你知不知道有這些問題.
你有沒有想過不要打這個?.
你寧可不打.
還是打另一個.
我又覺得很困擾.
於是我又出去和人家聊天.
於是另一陣營的人就說.
誰說的?.
你這個疫苗有什麼問題?.
那個疫苗才有問題.
你要留意聽聽數據是怎樣的.
於是就說你想想.
你打了就沒有這個險了.
聲音很激烈.
當面談完還不夠.
繼續在信息上來往.
開始令到我的心搞動.
不停在想怎樣打不打呢?.
回家和家人討論.

$^{201}$在過程中我發覺.
好像沒有誰能說服我.
我發覺我繼續問人.
就會分開兩派.
可能甚至三派.
中間那派就不打那派.
於是就有很多看法.
有一天我決定了.
要好好問上帝.
我真的感到很困擾.
那天是我下班的時間.
這件事也有在祈禱會的時間.
和弟兄姊妹分享過.
我下班的時間.
去到天水圍站對出的巴士站.
我站在那裡心裡問上帝.
我一時看到有很多討論.
關於打哪種疫苗.
打不打的問題.
一時我怎樣看呢?.
我真的覺得很困惑.
我雙眼望上天空.
聽到聲音就很清楚了.
說了三大字.
你猜哪三大字呢?.
就是「我保護你」.
你放心.
上帝不會在天空中說哪種疫苗的名字.
祂不會告訴你哪種疫苗好.
哪種疫苗不好.
祂說「我保護你」.
夠了.
討論結束.
我知道了.
上帝會保護我.
打也好不打也好.
哪種疫苗好.
上帝與我同在.
祂帶領我保護我.
夠了.

$^{241}$我就安心了.
回去跟著家人去打第一針.
在打第一針的過程.
我真的經歷了上帝的保護.
看顧.
很平安.
打完之後.
身心靈平安.
家人也平安.
我們在這件小事上看到.
上帝很奇妙.
神不停提醒我們.
祂帶領我們.
今天我相信更多的靈性姐妹.
在想打不打疫苗.
打哪種疫苗.
我們不如試試.
每個人都問問神.
祂有一個答案.
祂不會有一個模型的答案.
祂只有一個答案.
那個疫苗名.
我告訴你你就這樣打.
不是的.
我相信神對我們有不同的帶領.
我們每一個自己尋求神.
我們面對經濟的壓力.
這段時間疫症期間.
我心上有不同的家庭.
有自己的困惑.
有自己的挑戰在裡面.
可能會有失業.
家人失業.
家裡不夠錢用.
那怎麼辦?.
我們可以和靈性姐妹分享.
和教會分享和朋友分享.
最重要是自己和神分享.
祂是我們生命的神.
祂帶領我們走.

$^{281}$祂一定有話想和我們說.
我們要靜下來聽祂的聲音.
還有可能有些.
不只是經濟上.
水上有困擾.
很長時間.
我們固然有弟兄姊妹.
有朋友有家人.
但我們還有神.
這是最幸福的.
求神幫助我們.
求神醫治我們.
洪恩家所需要的.
也是神自己.
我們有很多很多的挑戰.
我們有我們的需要.
現在不是誰夠強勢.
出來說這一派的意見.
我告訴你.
要做這些.
然後另一派出來說.
一定要這樣.
不是誰鬥大聲.
我們就改變局面.
我們教會屬於上帝.
等上帝來作主.
等上帝帶領我們.
我們所需要的.
就是禱告尋求.
我自己很老了.
在這段時間.
推動不同的禱告會.
個人禱告.
弟兄姊妹之間的禱告.
其實我們弟兄姊妹.
都很興奮.
特別在我們星期三晚.
Zoom的祈禱會.
弟兄姊妹彼此代禱.
我們一起見證.

$^{321}$上帝怎樣聽我們的禱告.
改變我們的生命.
幫助我們.
我心上的只是一個開始.
我們要更多禱告.
太少了.
我們的禱告.
我們將主權交出來.
等上帝帶領我們.
記住上帝比地面上所有領袖都大.
所有的帶領者.
上帝是最終極的.
今天不靠他.
靠誰呢.
在以塞瓦孫這段經文裡.
上帝就是要這樣.
深刻地將這件事.
放在我們的心中.
上帝是榮耀的大君王.
他遠遠超越地面上的君王.
他是最聖潔的那一位.
面對這個.
榮耀聖潔的神.
天使的反應是甚麼呢.
這是很有趣的.
《六章二節》這樣說.
上有撒拉弗時納.
各有六個翅膀.
兩個翅膀遮臉.
兩個翅膀遮腳.
兩個翅膀飛翔.
每一位撒拉弗.
即每一位天使.
他有六個翅膀.
這個天使的特色.
是和火有關.
這個字.
這個詞語本身的字根.
和焚燒有關.
所以這個天使和火很有關.

$^{361}$今天你看的經文裡.
也和火很有關係.
這個撒拉弗有六個翅膀.
有兩個翅膀是用來做甚麼呢.
很特別 遮臉.
遮臉.
遮臉做甚麼呢.
是否他的臉不好看.
不是.
在我們的臉裡.
有一個器官很重要.
在我們的眼睛.
遮臉最主要是遮一雙眼.
做甚麼呢.
原來天使.
經常都站在神的旁邊.
他們不能夠.
直望著上帝.
為甚麼呢.
因為神很聖潔.
很榮耀.
聖潔榮耀到一個地步.
他們不能夠直望著上帝.
有點像我剛才的比喻.
太陽光照著你的時候.
直照著你.
你不會這樣張開雙眼.
看太陽是怎樣.
你不會這樣.
看完之後兩眼分花.
很容易瞎.
就是那種光芒.
那種榮耀聖潔.
在撒拉弗要遮著雙眼.
天使是以一個.
很敬畏的態度.
來面對神.
除了有一個很敬畏的態度.
天使是會開聲.
敬拜讚美神.

$^{401}$天使讚美的聲音.
很震撼.
說到甚麼呢.
好像門檻的根基.
都震動.
殿裡充滿了煙雲.
在這個讚美裡.
我們看到.
天使的聲音多麼的厲害.
神同在.
在過去這麼多年.
我所經歷過.
最震撼的.
敬拜讚美的場面.
是葛福林報道大會.
當時我還是神學生的時候.
很多間教會.
神學院的師班.
一起在香港大球場.
嘗試過.
你未見過這麼大的師班.
是很震撼的.
你未到大球場.
遙遠.
就已經聽到.
他們的歌聲.
歌聲繞樑三日.
在人的心裡留下一個很深刻的印象.
不過我想告訴你.
即使如此.
都未震撼到.
可以震動.
香港大球場.
那些門的根基.
或者震動球場的根基.
未到煙雲.
充滿了.
大球場的上空.
沒有這種震撼的場面.
所以這一場.

$^{441}$在天上的敬拜讚美會.
不是你和我.
所可以想像的.
天使.
以一個很敬畏的心.
去敬拜讚美神.
今天我和你又如何?.
神在我們教會當中.
我和你是如何敬拜的?.
什麼叫敬畏呢?.
這麼抽象的.
什麼叫敬畏呢?.
我們能夠將.
從裡到外.
你真的將最好的.
要獻給神.
在你眼前.
你看這個上帝很佩德.
很崇高.
甚至乎有些戰戰兢兢.
但不只是那種怕.
你有對上帝的那種愛.
但你不慶佛你的態度.
你要想著.
你不是隨隨便便給上帝.
是最好的.
這叫敬畏.
敬拜的態度.
今天我們還未恢復.
實體崇拜.
我們還在看網上直播.
我們正在預備.
我們的心靈.
很期待疫症快點好一點.
我們就可以回到.
洪恩那裡.
有實體崇拜.
到回來那一天.
我誠意邀請靈靈姐妹.
想想什麼叫.

$^{481}$以一個敬畏.
將最好的給上帝.
在裡面.
是什麼?.
我們的心預備好了嗎?.
整個星期的生活.
非常忙碌.
你是家庭主婦也好.
你是上班的也好.
有很多活動也好.
你經常在家裡休息也好.
這個星期的生活如何?.
你和神的關係如何?.
有沒有得罪神?.
有沒有得罪人?.
有沒有做上帝要你做的?.
遠離上帝不要你做的?.
你這一刻的生命光景是怎樣?.
回到來的時候.
真的要有一個時間好好安靜自己.
預備自己.
你要將最好的獻給神.
崇拜是最好的.
不只是靠靈唱.
唱很震撼的詩歌.
講完.
講到有沒有趣?.
是不是震撼人心?.
主席是不是很觸動我的心.
帶我進入崇拜?.
不只是靠這些.
我們自己預備好了嗎?.
我很欣賞.
有些弟兄姊妹.
我見到他們.
我們很難得.
實體崇拜的時候.
整個星期沒有見.
一回來.
捉著這個捉著那個.

$^{521}$我之前聽你說.
那些事情如何?.
我很興奮.
發生什麼事情.
中午去哪裡吃飯.
一路說了大量.
其實這些溝通是要的.
回到教會就想這樣.
我想崇拜前.
準備好自己.
有些弟兄姊妹就靜下來.
不會和人聊天.
一路默默的禱告.
預備好自己.
準備迎見神.
你將最好的.
心最好的.
給神.
由裡到外.
外面是什麼?.
收拾一下.
用心就行.
不是.
我想問大家.
你所愛的人.
會不會收拾你的心?.
你會不會難生難逝.
蓬頭救命.
你知我份心.
我愛你.
你知我愛你.
我就來.
是這樣的.
你會收拾你的心.
因為你愛.
因為你敬畏.
不一定要.
每次都西裝.
不用穿什麼.
拿什麼很昂貴的東西.

$^{561}$怎樣去整.
化了很多妝.
不一定是這一種.
至少端莊.
美美的.
你最美的造型.
因為是裡到外最好的獻給神.
上帝喜歡美的.
上帝喜歡秩序.
上帝喜歡整潔.
我知道上帝喜歡這些.
我在上帝面前就預備最好的.
他配得.
他崇高.
他配得.
整個人準備好.
我記得我以前.
有八年時間.
做青少年侍工.
在崇拜.
有一個很大的教導.
我經常會教青少年.
講好衣著.
我們當然.
回到成年人的世界.
我們很少這樣.
但青少年會講明.
有什麼衣著在崇拜時.
不要穿回來.
因為青少年總是充滿年輕熱誠.
回來.
穿得有多短多短.
有多潮多潮.
最重要是鞋.
現在最有型的是拖鞋.
踢拖鞋回來.
整個造型.
又潮又性感.
又這樣那樣.
整個崇拜.

$^{601}$很吸引人去看他的樣子.
結果大家的焦點.
不是好好敬拜讚美上帝.
而是很特別.
整個景觀.
在過程中.
我一定會教導他們.
有什麼要注意.
是會教導的.
因為你不只是講你的心.
你整個生命獻情給神.
沒有什麼是非的.
而我們的外在.
也是一個姿態.
我要將什麼給神.
這是很重要.
值得我們留意.
期待有一天.
實體回來見到大家時.
我們每個人.
以最好的獻給上帝.
第二件事.
在面對榮耀聖潔的神時.
剛才提及天使的反應.
是很敬畏的.
敬拜讚美.
人的反應是什麼呢?.
人的反應很特別.
但這是一個很正常.
很好的反應.
我們看回.
《耳塞瓦書》六章五節.
他說.
「那時我說」.
這個「我是耳塞瓦」.
「那時我說:禍災,我滅亡了.
因為我是嘴唇不潔的人.
住在嘴唇不潔的民眾.
又因我親眼看見.
大君王萬君之耶和華.

$^{641}$又因我親眼看見.
大君王萬君之耶和華.
耳塞瓦.
看到榮耀聖潔的上帝的反應.
「好耶!上帝!我來了!.
你真好啊!」.
不是這隻王.
是很害怕.
「我有禍了!我死了!.
我滅亡了!」.
是吧?.
害怕.
害怕什麼?.
他看到自己的罪.
「因為我是嘴唇不潔的人.
因為我是嘴唇不潔的人.
我的嘴唇犯罪.
我的嘴唇犯罪」.
第一層意思像說嘴唇.
我們的嘴唇很容易犯罪.
是吧?.
有時會想.
本來是說一件事.
說某一個人.
好像說一些事實.
說著說著,我們會發覺.
我們說話會說多了.
會覺得弟兄怎樣怎樣.
姐妹怎樣怎樣.
加一些評語.
說著說著,A君說給B君聽.
B君又說給C君聽.
說著說著,就變成是非.
就很容易.
這樣傳出去.
嘴唇很容易犯罪.
說錯話,傷害人.
是吧?.
我們不是想傷害人.
和人吵架.

$^{681}$不知為何對方的反應會令我們衝口而出.
這樣就死了.
很容易.
我們說些說話.
令人跌倒.
我想這不是我們的原意.
嘴巴.
而作為以賽亞.
他.
亦很明白這一點.
嘴巴.
很容易不潔.
但其實嘴巴不只是指嘴巴.
是代表整個生命有罪.
在聖潔的神面前.
以賽亞看到自己有罪.
他亦看到他的群體有罪.
住在嘴唇不潔的民眾.
就是他自己的百姓.
是吧?.
他們也和我一樣.
嘴唇犯罪,得罪神.
但其實是整個生命.
我們住在這個群體裡.
整個群體我們都有犯罪,得罪神.
這是以賽亞的反省.
我們在這個群體裡.
反省.
可能很多人.
會這樣.
我作為以賽亞會怎樣呢?.
哎呀!上帝.
我看到你榮耀聖潔,就想起.
我所在的群體.
他們每一個都有犯罪.
其實我沒有問題.
不過我真的很委屈.
我住在這個群體裡.
每一個都有罪,每一個都有問題.
他們真的不好.

$^{721}$是吧?上帝你憐憫他們.
你寬恕他們.
你幫他們處理他們的罪惡問題.
很多人都會從這個角度.
覺得自己沒有問題.
別人才有問題.
我想起我剛剛.
加入教會的時候.
我在舞會的時候.
很厲害,我的嘴巴很厲害.
不需要別人教.
不知為何,看聖經又不算很熟.
一回到教會沒多久.
弟兄不屬靈.
姐妹不屬靈.
又不靈修,又不禱告.
我覺得他們有這個問題.
那樣問題.
我專門跟別人說.
怎樣怎樣有問題.
一回到教會沒多久,整個審判官.
全世界都有問題,只有自己沒有問題.
真的很容易.
不需要別人教.
我們只看到別人.
眼中有刺.
我們自己有良目.
每個人都壞人.
自己是好人.
最壞就是他們.
令我們這樣.
我在這裡就好了.
我記得我剛剛回教會也是這樣.
我記得我當時回一個青年團.
青年團的團友.
從小到大.
都在教會成長.
是教會的子弟.
他們之間的關係很親密.
他們的關係不單親密.

$^{761}$連繫也很強.
我是中途信耶穌.
大學時才信耶穌.
我加入他們.
自然需要很長時間才融入.
我就感覺,這麼久.
好像融入不了他們的世界.
他們經常都自己一個小圈子.
於是我心裡.
開始很憤怒.
我心裡一直責備他們.
一直說他們.
這班人很自私.
自己的.
不容納愛人.
我很不開心.
心裡罵了很多次.
有一次一直在罵.
我去到上水火車站.
我太記得這一刻.
我當時不太熟悉聖經.
但不知為何神就有一節經問到我.
既然夢照.
行事為人.
就當與夢照的恩相稱.
你已經得夢拯救.
你就不要這樣.
去對待人.
心裡不停地咒罵人.
心裡咒罵人.
你以為人家看不到你.
上帝看到的.
上帝提我,用聖經提我.
那一刻.
我知道神責備我.
就停下來.
過了不久.
我又開始在罵他們.
覺得他們很有問題.
指責他們有罪.

$^{801}$你搞笑圈子.
罵罵罵.
神就繼續用他的方法去管教我.
在這個過程.
如果我.
生命常常覺得自己沒有問題.
別人有問題.
我就會這樣.
這些就是嘴唇不潔的人.
我住在他們當中.
這些就是犯罪的人.
我被這些人.
玷污了我的生命.
我是受傷的人.
但以差鵝不是這樣的.
沒錯.
他看到整個以色列民.
這群人他們是犯罪.
但他才承認.
自己有罪.
我嘴唇不潔.
我死得了.
我認罪.
我真的不好.
這一年.
我自己.
很努力地學認罪功課.
現在才學認罪功課.
剛好信耶穌時.
要學這些.
現在學,信耶穌多久了?.
信耶穌多久.
都要好好學.
我去看我自己生命.
有什麼罪惡.
我向上帝認罪.
逐樣逐樣數出來.
向上帝幫助我對付我的罪.
然後.
我又為我們整個紅恩家一同禱告.

$^{841}$為我們整個群體.
我們求神憐憫.
求神寬恕.
你看到.
神正等我們這件事.
為什麼這樣說?.
六章六字.
有一撒拉夫向我飛來.
從紅的壇試用火鉗.
從壇上取下來的.
撒拉夫向我飛來.
這個動作是即刻的.
什麼即刻?.
以塞埃一認罪.
撒拉夫即刻飛過來.
感覺是什麼?.
上帝就在看著.
認罪了嗎?.
一認罪我就過來.
等這一刻等很久了.
上帝就想我們認罪.
祂真的在等我們.
等到這一刻.
拿著燒紅的壇.
會不會燙到嘴唇.
這麼可怕.
其實我們不是從這個角度.
這個壇是象徵.
潔淨.
是潔淨.
上帝會潔淨我們.
透過天使拿著壇.
去潔淨.
以塞埃的嘴唇.
潔淨整個生命.
讓他得潔淨.
他的罪惡得赦免.
今天的經文沒有說到.
大家如果有機會再看下去.
六章八節一路打後.

$^{881}$當以塞埃為自己認罪.
為百姓認罪之後.
神潔淨他之後.
神就用他.
差遣他.
然後就用他.
差遣他.
今天值得.
我和你.
去想一件事.
是不是神也在等我們認罪呢.
在等我們.
你說有什麼罪呢.
我們在神面前安靜下.
問下神.
我的生命是不是全然討你喜悅.
我有沒有得罪你的地方.
聖靈一定會光照我們.
除非我們心很硬.
否則聖靈一定會告訴我們.
有什麼問題.
今天神也在等我們.
正如昔日.
他在等以塞崖.
為自己為百姓認罪.
一等到.
他就去寬恕.
今天神也在等我們.
等我們認罪.
他就會寬恕我們.
潔淨我們.
然後就會用紅印鑑.
這是第一步.
如果我們不做好這一步.
我們又期望.
我要下一步做什麼.
這是我們人的想法.
在神裡面.
他想要聖潔的生命.
他要的.

$^{921}$那些生命的質素.
在這一刻.
我邀請大家安靜地一同禱告.
請大家問神.
這麼久以來.
特別是這段時間.
我們生命裡有什麼是神不起悅的.
我們先不要想.
身邊誰誰.
一想起就想那個人不對.
這個人不對.
想自己.
我們有什麼要向神認罪.
我們這一刻向神認罪.
請大家安靜地禱告.
天父上帝讚美你.
是聖潔榮耀的神.
如果你在場.
聽到撒拉弗.
怎樣稱頌你.
盛哉盛哉盛哉.
我相信.
我們的心會很震撼.
在你的面前.
我們一定會賤敬.
可能會和.
二賽亞一樣.
想叫死吧.
有我了.
因為我們真的犯罪得罪你.
求你寬恕我們每一個.
盡你的地方.
我願意耶穌基督.
你真的以你的寶血.
潔淨我們的罪惡.
主啊你說過.
輕憂傷痛悔的心.
你必不輕寒.
主讓我們為我們的罪.
憂傷痛悔.

$^{961}$你讓我們悔改.
你讓我們深深的.
經歷你的接納.
願主使用紅印格.
使用每一個生命.
成為傾覆於別人的器皿.
在我們要成為紅水橋的.
明亮登台之前.
願主.
全然改變我們的生命.
奉耶穌基督的名字.
祈禱.
阿們.
\newpage



\section{使徒行傳 10:34-35}
\label{sec:7LBy_VBf8iw}
\textbf{2021.03.28《誰是主宰》使徒行傳 10\_34-35( 和修版)  鄧泰康傳道}
\newline
\newline
連結: \href{https://youtube.com/watch?v=7LBy_VBf8iw}{\texttt{https://youtube.com/watch?v=7LBy\_VBf8iw}} ~~~~ 語音日期: 2021-03-29
\newline
\newline
\hyperref[sec:oH_gH6L8AqQ]{\small{< < < PREV SERMON < < <}}
~
\hyperref[sec:index_chronic]{\small{[返順時目]}}
~
\hyperref[sec:index_scriptual]{\small{[返順卷目]}}
~
\hyperref[sec:rYjqMvMSdYU]{\small{> > > NEXT SERMON > > >}}
\newline
\newline
使徒行傳 10:34-35
\newline
\begin{longtable}{cl}
\hline
\hline
章節 & 經文 (和合本修訂版)\\
\hline
10:34 & \begin{tabularx}{0.7\textwidth}{X} 彼得開口說：「我真的看出神是不偏待人的。 \end{tabularx} \\ \\ \relax
10:35 & \begin{tabularx}{0.7\textwidth}{X} 不但如此，在各國中那敬畏他而行義的人都為他所悅納。 \end{tabularx} \\ \\
[1ex]
\hline
\hline
\end{longtable}
$^{1}$頂智本早晨.
如果以港獨來說.
我也是第二次來到裡面.
上一次在這裡高朋滿座.
這裡已經坐滿了.
之後跟大家分享.
這次是一個直播.
在即將來臨的新消息.
昨天是第一次.
11月到現在第一次清零.
即是香港沒有本土的個案.
我們也收到一些消息.
會不會有機會.
在復活節之前.
可以崇拜重新再開放呢.
雖然人數可能會減少.
但是頂智妹.
很實在我是很期待很渴望的.
因為真的可以和頂智妹面對面接觸.
這種的港獨是很不同的.
其實我是一個很不喜歡對著鏡頭說話的人.
我很喜歡看到頂智妹的眼神.
我很喜歡看到頂智妹的回應.
所以我對自己的港獨只有一個要求.
就是我港獨的時候.
我希望不會看到頂智妹睡著了.
這就是我的要求.
我希望盡量在港獨裡面的時候.
讓我頂智妹能夠專心聽上帝的話.
所以在家中的時候.
我是有些擔心的.
因為有些頂智妹可能在被窩裡面聽.
有些可能是剛剛睡醒的狀態去聽.
所以我在這裡更加鼓勵.
萬一繼續要網上的話.
請頂智妹最少提早十分鐘.
做好一切的事.
放下一切的事.
之後專心坐在電腦面前.
如果是可以的話.

$^{41}$我更加鼓勵換一套衣服.
坐在這裡好好地去敬拜上帝.
整個狀態進入這個敬拜的狀態.
這是很重要的.
我今天跟大家分享這段經文的時候.
是整個小道行經第十章.
剛才我請主席讀的時候.
只有兩節的原因.
因為這裡也是一個比較頗長的故事.
所以我想用一點時間.
很快地跟大家交代這個故事.
是說什麼.
第十章的時候.
他說在凱撒尼亞一個外邦的地方.
有一個叫哥利留.
是意大利營的白夫長.
是一個軍人和軍官.
但聖經形容他就是一個虔誠人.
他全家都敬畏神.
而且他行出信徒的榜樣.
多多州際百姓上上禱告.
很明顯的就是.
他聽到上帝的時間.
然後他就信奉.
然後他就按照著聖經對他們的要求去做.
然後他就行出來了.
有一天的時候.
大概下午三點鐘.
他看到一個異象.
有使者來到跟他說.
他是哥利留.
哥利留就問主下什麼事.
然後天使就說.
你的禱告和你的州際已經去到神面前.
已經蒙月立了.
所以你現在打發人去到約柏的地方.
找一個叫彼得西門的人.
你去吧.
然後他和他的家人商量完之後.
就派了一個很虔誠的兵.

$^{81}$我想大家留意一下.
連派他的兵都是虔誠兵.
不是隨便派一個.
而是愛主的兵去到那裡.
然後他差不多去到的時候.
聖經說第二天.
彼得在樓頂中午去禱告的時候.
然後覺得肚子餓了.
他在家裡也預備煮飯.
他也看到一個異象.
這個異象就是天開了.
在上面的時候.
有一塊布他放下來.
攤開的.
上面有各樣的動物和昆蟲和天上的飛鳥.
然後聖經就和他說.
彼得起來.
趕快吃吧.
彼得就說不行主.
不可以.
因為這些是不潔的物.
我從來沒有吃過.
最重要的是第二把聲音.
神所潔淨的你不可以當作俗物.
聖經形容三次之後就收回上去了.
正在彼得心裡疑惑之間.
聖靈就提醒他.
你下去吧.
有人來找你.
你跟他們去吧.
是我差派來的.
然後彼得就下去見那些人.
問他們什麼事.
那些兵也和他們分享整件事發生什麼事.
然後他就說你放心.
我會和你們去.
因為這是上帝告訴我的.
到第二天的時候我們就出發.
去到凱撒利亞的時候.
哥利留一家人就已經很熱切的期待.

$^{121}$等他們來.
之後哥利留又重新再一次交代了整件事是什麼事.
然後彼得的回應就是剛才大家聽的那段經文.
我真看出神是不偏待人.
原來各國中那敬畏主行義的人為主所喜悅.
為主所悅立.
然後他就說了一篇道.
這篇道是非常短.
其實我很羨慕的.
彼得在說道的時候.
是很簡單直接了當.
很短的說話.
將主耶穌整個救恩交代了.
然後聖經最好玩的地方就是.
彼得還沒說完.
這篇道已經夠短了.
但是彼得還沒說完.
然後這些人就被聖靈充滿.
然後他們說了謊言去讚美上帝.
彼得就說既然這些人聖靈都充滿他們了.
我還怎麼可以禁止他們受洗呢.
然後彼得就為他們受洗.
所以頂子們.
我想告訴你們.
在這段經文的時候.
我真的覺得很驚訝.
我作為傳道人.
其實我是很羨慕的.
你想想如果有一天你請鄧傳徒上來.
然後我上來.
你試試讀一下那段經文.
如果就這樣說.
我相信是連兩分鐘都不用的.
我都還沒兩分鐘說完之後.
下面全部舉手.
我信耶穌.
這個很震撼.
我仍然很震撼.
但是這個也是很巧妙的地方.
就是整件事真的與彼得無關.

$^{161}$就是說這些人信耶穌是與彼得無關的.
整件事從頭到尾都是上帝的工作.
為什麼會有這件事呢.
有趣的地方就是我想大家留意.
其實彼得正在面對一個很大的矛盾.
這個矛盾就是.
他從出生一直堅持著的律法.
由律法告訴他什麼叫做不潔淨.
你說他隨便說就沒得說.
但他不是很認真.
我知道什麼不潔淨.
我從出生到現在都沒有吃過.
他跟上帝說我沒有吃過.
但是現在上帝跟他說.
我說潔淨.
我以前一直覺得不潔淨的東西.
上帝突然跟我說潔淨.
上帝怎樣啊.
用現代人的術語是怎樣.
你玩完啦.
潔淨也是你說.
不潔淨也是你說.
但是弟兄姊妹我想告訴你.
是呀.
是上帝決定的.
因為律法是上帝設立的.
我們跟上帝相處.
請大家不要以我們同一個人相處去相處.
你記住我們同一個上帝去相處.
跟上帝相處的時候.
是有很多的超越性.
我經常說上帝所知的.
其實是比我們所知的多很多.
所以你會發覺在這個時候.
上帝說是呀.
沒錯.
在以前的時候.
我跟你說了這些是不潔淨的.
但是我現在告訴你.
現在是潔淨的.

$^{201}$為什麼.
我想大家留意.
當日上帝去設立這些律法的時候.
其實有一些問題.
到今天你為什麼都解不通.
我想問問大家.
為什麼豬要比牛不潔.
你客觀解釋給我聽.
今天現代人.
很多解經的人嘗試用很多方法.
但我想問他們是廢話.
我曾經聽過一些最有自給的解釋.
豬是豬欄的.
大家有沒有豬欄.
豬欄很髒的.
又臭.
滾來滾去.
豬永遠出來都不是白色的.
永遠都是髒的.
所以不潔.
當然不是.
當然不是這樣.
我很清楚地告訴你.
這只是當時上帝對以色列人定了的標準.
即是我可以告訴你.
上帝同樣可以定一個標準.
豬是比牛潔淨的.
是同樣可以的.
意思是什麼.
一直以來的時候.
上帝要人去聽.
和納粵的準則都是.
你要聽明.
你不是自己去分辨什麼是潔淨.
什麼是不潔淨.
我想大家弄清楚.
因為如果你記得的時候.
當阿當夏娃開始的時候.
上帝對他們的挑戰那一棵就是分辨善惡樹.
他早就說了.

$^{241}$人是不應該自己去分辨善和惡的.
早就說了.
善和惡是上帝的權柄.
人可以去做的就是聽明.
所以當大家留意的時候.
最關鍵的那段經文的說話就是.
神所潔淨的你不可當作俗物.
留意.
神所認為是潔淨的你不可以當作俗物.
另外一個翻譯就是.
上帝已經潔淨的東西你不能稱為不潔.
你留意.
上帝說是上帝決定.
所以其實你去看哥利留的故事的時候.
我會這樣詮釋.
凡合乎上帝標準的.
才是潔淨.
所以你就會發現一件事.
究竟上帝是不是反來覆去.
上帝是不是在耍花樣.
不是的弟兄姊妹.
我想你弄清楚.
上帝在聖經裡面曾經有沒有試過反口.
有沒有?.
有.
記不記得約拿.
他叫約拿離開的時候是怎樣的.
你告訴他我一定會滅了你們.
然後出現的情況是甚麼.
約拿城的人全部由君王至動物.
全部披麻蒙灰.
是誇張到整個城.
然後上帝又怎樣.
後悔了.
上帝寧願讓你說他反口.
但他不滅.
出於甚麼.
出於愛.
出於憐憫.
上帝並不是憑著自己的血氣.

$^{281}$我說想怎樣就怎樣.
上帝不是這樣的人.
你弄清楚上帝不是這樣的神.
上帝是以善.
以愛.
以好為最高的標準.
只需要人聽命他的說話的時候.
就一定是好的.
所以凡是聽命令的人.
在神的眼中都是蒙越納.
為甚麼這麼重要呢.
我想大家留意.
原來第十章和第九章是連結的.
第八第九章發生甚麼事.
就是蘇羅第一次出場.
蘇羅是未來上帝將外邦人交給他的使徒.
預備要將這個福音傳給外邦.
但我想大家留意蘇羅是一個怎樣的人.
蘇羅其實我這樣形容他是.
法利賽人的影子.
正宗的血統.
很認真的學習聖經.
做事熱心.
但我想問你他弄了甚麼出來.
迫害上帝的僕人.
殘害教會.
為甚麼會這樣.
因為他們認為自己掌握了上帝的標準.
我已經很清楚上帝告訴我甚麼是潔淨甚麼是不潔淨.
我很清楚甚麼是彌賽亞.
所以我不接受主耶穌你是彌賽亞.
我不接受你是彌賽亞.
因為我們知道我們是君王.
因為君王有甚麼可能會在馬槽出世.
沒有可能.
我們的彌賽亞有甚麼可能會被十字釘死.
沒有可能.
所以我們不接受.
他們就將自己所學自己所堅持的東西.
成為了自己的標準.

$^{321}$而這個標準大到連上帝都拒絕.
認為上帝你都沒有資格去改.
因為你不會改.
你搞錯了.
原來一個人當他就算熱愛上帝的律法.
熱愛上帝的聖經.
但他沒有將生命的主權交給上帝的時候.
是這麼大件事.
但大家有沒有留意出現的哥利奈.
哥利奈其實我想告訴你.
第九章出現的人物其實可以不是一個外邦人.
是可以是以色列人裡面一個人.
大家有沒有想過.
他偏偏是一個外邦人.
對比強烈到不得了.
是一個外邦人.
擁有當時最高的權力.
記住是羅馬人.
他最重要.
但他卻是一個敬畏上帝.
行上帝的道去周濟百姓的人.
他卻是一個謙卑下來.
他不用他的權力.
他反而將他放下.
將他自己交給上帝.
然後行出上帝的道.
所以聖經就要告訴你.
凡是聽從上帝說話的.
都蒙上帝月立.
有沒有留意.
所以我想幫大家重新認識.
上帝對於潔淨的標準.
弟兄姊妹.
我們今天不是光有一個信徒的身份.
你就在聖結.
而是你要有信徒的生命.
其實你會發覺一樣東西.
就像我剛才所說的素羅一樣.
他的身份是正統的.
他是認真的.

$^{361}$但問題是.
他的熱心放在他自己認為的標準裡.
他沒有行出上帝的生命.
但這個哥利留.
在當時的世界.
他絕對不是一個信徒身份.
因為當時的外邦傳統還未開始.
但反過來.
他卻是滿有一個信徒的生命.
誰蒙上帝月立.
就是哥利留.
所以弟兄姊妹.
我想你們留意到一樣東西.
其實今天我們作為基督徒.
當然了.
記住這些是好事.
要平衡.
我們很誠實.
我們很榮譽於我們是一個基督徒.
因為這是尊貴的.
這個身份是上帝給我們的.
但同時.
請你不要只掛著一個基督徒的名份.
其實今天有很多人掛著基督徒的名份.
不等於生命的聖結.
你弄清楚.
所以有些弟兄姊妹.
我做婚前輔導.
或者做年輕人工作.
未結婚的多.
有些很多時候走來問我.
一些我認為是不需要問的問題.
很多餘的問題.
但一轉眼.
我將來找老公找老婆.
是不是要基督徒.
你說是不是佛教.
是不是要基督徒.
你說是不是廢話.
你問我的時候.

$^{401}$我當然要跟你說.
當然了.
我不會跟你說不用的.
當然了.
但我一定會補充一句.
你要找的不是基督徒.
而是一個很認真信耶穌的基督徒.
他是一個很有生命的基督徒.
你的婚姻就受到保障.
明白我的意思嗎?弟兄姊妹.
今天有很多人掛著的名字是基督徒.
之後其實生命是什麼都不是.
又回到婚姻.
我曾經見過一些很搞笑的.
要追求女孩的時候.
回應得怎樣.
我小時候曾經回教會.
我又見過有些時候.
外面有些人亂來.
然後來找教會幫忙.
我經常說這些.
就是世界仔.
先跟你裝朋友.
人家跟你說.
我小時候曾經回教會.
是真的.
你小時候曾經回教會.
那又如何.
代表什麼.
不要和上帝拉關係.
我是很願意幫你的.
但不要用俗世的方法.
同樣的就是今天.
有很多很多的人.
當你說自己是基督徒的時候.
上帝看的不是你這個名字.
而是看你的生命.
你的生命是如何.
究竟今天你的生命.
是否配合基督徒的生命.

$^{441}$這才是上帝認為是聖潔的.
弟兄姊妹.
第二個就是.
不要空有聖經律法的知識.
而是你要去遵行有多少.
今天這個香港的社會.
我可以說我們真的很幸福.
你要的資源.
要多少有多少.
神學院我們也有這麼多間.
我經常說笑.
除非是你不讀.
你不會沒有得讀的.
但今天整個過程裡.
大家都很追求聖經的研究.
你記住.
我不是叫大家不要讀聖經.
不要研究聖經.
你不要誤會我.
我們要很認真去學聖經.
去追求聖經.
不過更重要是你學了之後.
你有沒有放回你生命裡面.
去實行出來.
我坦白說.
我也認識不少的人.
他今天對聖經很認識.
甚至乎以教導聖經去事奉.
但很實在.
他回到家裡的時候.
跟他的父母.
跟他的配偶的關係.
是一塌糊塗的.
我作為今天.
我自己是一個傳道人.
我也相信弟兄姊妹.
同樣會把尺子放在我身上.
我今天正因為.
我有很多機會站在上面講道.
所以我知道弟兄姊妹.

$^{481}$對我們的要求會更高.
我不是只聽你講.
究竟平時鄧傳道.
你的生命是怎樣的.
我所認識的你.
是不是你口中所講的那種愛.
是不是你所講的那種善良.
是不是你口中所講的那種信服.
那種信心.
講很容易.
我可以告訴大家.
但你生命實行出來.
這才是上帝要的.
你容許我去感恩.
我在之前事奉的教會.
因為是一間小教會.
我的小朋友在那裡成長.
我的小朋友今天在大學畢業.
已經出來工作了.
他從小到大的時候.
去教會 回兒童聖經學校.
但我們的教會不夠資源.
去到中學初中.
他就要坐在聖人崇拜.
由初中就要聽我講道.
我們的教會請人不多.
我每一個月都有兩次講道機會.
你知道做傳道人.
永遠都有壞習慣.
因為你沒有話說.
你當然要講自己家.
講自己家當然會講兒子.
有很多傳道人的子女不喜歡.
你講道就講道.
你為甚麼要講我.
我也會這樣分享.
但有一次我問他.
我坐下來.
當時他比較大了.
高中了 再大一點.

$^{521}$我就問他.
其實你坐在那裡崇拜.
你聽爸爸講道.
你尷尬嗎.
或者你覺得.
為甚麼你會坐在那裡聽.
其實我有一點.
我有一點.
他在學校也有回過團期.
我說不如你回學校幫他教會.
會不會怎樣.
他說不會 我不會回來.
我問他為甚麼.
他說因為你講的道.
和你在家裡我見到的那個你.
是一樣的.
我不覺得有甚麼問題.
第二件事.
你從來分享我的事.
你都不是講我的壞事.
你是講感恩我的好事.
我聽完之後.
我恍然大悟.
因為由這個小朋友出生開始.
我就已經很清楚.
這個小朋友是上帝給我的產業.
不是我的.
我是要用上帝的方法.
用上帝的原則去教導他.
讓他去成長.
我不是他的主人.
我不是他的上帝.
我是要幫他.
讓他去發揮上帝在他生命的東西.
所以有些人會覺得很奇怪.
為甚麼我兒子這麼快就已經.
因為我兒子在台灣.
為甚麼他去台灣.
我可以告訴大家.
他中二那年就已經跟我說他去台灣.

$^{561}$我不會當他開玩笑.
因為他在當時.
他在網上認識了一班台灣的大學生.
大家談得很好.
他很嚮往台灣.
所以他很有趣.
他知道整個過程他自己就.
我要去台灣.
我對他努力.
我身邊他會支持你.
所以他整個過程很順利地.
他就去台灣讀完書.
然後去台灣畢業.
他現在很感恩.
他剛剛在當地找到工作.
他可以開始慢慢的生活.
上帝給他的軌跡.
我們今天在神的面前的時候.
連你的家人都一樣會看你是誰.
所以在神面前的聖潔.
不是你做了甚麼.
而是你的生命是如何的.
所以求神幫助我們去學習彼得的謙卑.
彼得在這裡的時候其實很奇妙.
我經常都說.
在這裡的時候.
我最震撼的地方就是.
整件事都不關彼得的事.
你想想.
彼得其實很固執.
他的固執是從哪裡來的.
他的名字是石頭.
後來才叫盤石.
原本就是石頭.
他很固執.
你們會發覺他以前的固執是怎樣.
當耶穌問他的時候.
將來你們會不認我.
彼得怎麼說.
我發誓我不會不認你.

$^{601}$然後這個人是怎樣.
我又發誓我不認識這個人.
後來軟弱到要做漁夫.
耶穌把他撈回來.
但我想告訴大家.
他後來出去講道的時候.
一次講道三千順耶穌.
我可以告訴你.
他很清楚這個不是他的能力.
但他的生命是著實改變.
因為問題是他是固執的.
但後來他謙卑下來.
他變成擇善而執.
他擇上帝的善而執.
上帝你說好.
我就堅持下去.
所以他改變了之後.
他對外邦人講的第一篇.
第一番說話是什麼.
上帝要我說的.
我不可以不說.
就算我不說.
石頭也會說.
多麼厲害.
所以你會發覺一件事.
原來他的生命真正的轉變.
是要謙卑下來.
去馴服聽命於這位.
超越的上帝.
剛才我提過.
他的話還未講完.
聖靈就替他工作.
其實原來我們為上帝工作的時候.
上帝一直強調.
我不是要你多厲害.
當然我上一期.
記住永遠都要平衡.
我們要預備好.
盡我們的能力去做好.
但上帝最後看你的不是你有多好.

$^{641}$上帝是看你有多忠心.
我就用一個例子.
我都頗欣賞師班的.
因為師班是有沒有獻師都好.
都每個星期要練習.
我們都要練的.
我們要準備好.
準備好之後.
我們就走去為上帝奉獻.
為上帝做好.
我們是要盡我們的能力.
但上帝看的不是你唱得多好.
而是你這班人又忠心又善良.
所以今天是上帝替我們去完聖所有的工作.
我們謙卑下來.
被上帝使用就夠了.
第二件事.
我不知道大家有沒有留意.
剛才我說第九章和第十章.
當時素羅剛剛出場.
彼得是未認識素羅的.
他亦都完全不知道.
究竟接下來的時候.
他會使用這個人.
但我想大家留意.
有沒有留意.
一切的事情都未發生.
上帝就親自疏通所有的問題.
因為他知道.
第一個要改變的是誰.
就是以色列人門徒的頭兒彼得.
他們的觀念一直都是.
我連和你相處都不恰當.
怎麼傳.
但現在上帝親自用這個方法.
讓他親自有一次.
要和這些人去受洗.
讓他們見到上帝的工作.
所以大家記不記得.
當素羅後來變成保羅.

$^{681}$走出來傳道的時候.
他遇到最大的問題是什麼.
就是被人罵到瘋了.
怎麼搞的.
你在叫人不要受割禮.
然後就被人抓了.
所以一直被猶太人迫害.
就是迫害這件事.
後來他就抓了去見彼得.
記不記得彼得當時的說話.
就是這些原因.
他就說回這件事.
是上帝親自將所有的計劃.
自己砌好.
我們今天身處其中.
我們都不會知道發生什麼事.
坦白說這兩年.
我們都不會知道發生什麼事.
疫情也好.
社會的情況也好.
或者現在經濟情況也好.
我們完全不會知道上帝在做什麼.
不過我想告訴你.
是呀你不知道的.
不過請你好像彼得記一樣.
謙卑下來擇善而執.
我經常都說.
這段日子.
我是一個小人物.
我改變不了什麼.
我做不了什麼.
但我想告訴大家.
總是有很多事情等著你去做.
對不對.
無論這個世界上是多混亂.
這個世界上是多麼有問題.
但仍然有很多人在水深火熱當中.
要你去關心他.
我們有沒有做到.
還是我們只是停留在這個困境裡面.

$^{721}$我們只是不斷在埋怨上帝.
上帝有沒有搞錯.
搞到我們現在這樣.
為什麼這個社會又這樣.
我們會不會只是停留在這裡.
而我們什麼都沒有做.
所以弟兄姊妹.
求上帝幫助我們擇善而執.
繼續在任何環境.
堅持上帝要我們做好工作.
盡心盡性盡意盡力去愛神和愛人如己.
我們今天整個人生裡面的總括就是這兩句話.
將來上帝去評論你.
也是用這兩句話.
所以求上帝給我們生活裡面的每一個小節.
是的沒錯.
我們的信仰是生活.
我們生活裡面應該是滲透著.
我們整個信仰上帝的愛和他的能力.
當你每一次和你的家人相處.
請你想一想.
上帝想你怎樣和他相處.
當你身邊有些人出現的時候.
上帝想你怎樣和他相處.
怎樣去幫助他.
怎樣和他同行.
每一樣東西和上帝一起去走.
我可以告訴你.
你的生命才是真正的豐盛.
所以到最後的時候請大家不要搞錯.
我們只是上帝的流通管子.
我們沒有什麼好誇張的.
譬如簡太太很好.
每個星期幫我們一起去彈光琴.
這是上帝所給你的恩賜.
是不是.
簡太太彈得很好.
但我要這樣說.
如果有一天.
上帝突然收回她.

$^{761}$她的手受了傷.
她在鋼琴這件事上就什麼都不是.
我們今天盡力去練.
盡力去自馮是對的.
但是我想提醒姐妹.
其實我們沒有什麼好誇張的.
我們今天在上帝面前.
只可以做我們最基本上的.
按恩賜給我們的事奉.
但我們沒有什麼值得誇張的.
我們有什麼好誇張的.
我們人與人之間.
誰高誰低.
沒有.
教會就是要把事情削平.
任何人來到神的殿裡面.
都是一樣的.
是沒錯.
我們和某些人相處.
性格上未必很合得來.
是對的.
所以要承認.
不要美化到.
不是犯人我也可以.
這是上帝給你的恩賜.
但是你不用和他做仇人.
你可以和他好好相處.
你不是和他去資深朋友.
但你可以和他一起走.
因為我們有同一個目標.
要他傳承福音.
所以一樣可以走得通.
是對的.
所以弟兄姊妹求上帝讓我們在當中.
去知道原來我們只不過是一個卑微的人.
我們在神面前的時候.
我們要學習.
你就是主宰.
我們看別人比自己強.
我們去看其他人在上帝眼中是可愛.

$^{801}$我們的包容接納程度就會大很多.
我們的紛爭就會少很多.
我們在神面前的時候.
我們每一個都尊心於.
在上帝給我們的生命當中的時候.
整個教會會美善很多.
所以求神幫助我們.
在神面前透過他的說話.
再一次提醒我們.
祂才是主宰.
由祂決定.
最後我也要和大家說.
我覺得這是一個.
我覺得是時空裡面很好玩的地方.
為什麼會說好玩呢.
因為上帝每一天.
讓出現在你身邊的人都不同.
你都猜不到.
其實每一個時候.
你只需要什麼.
上帝每一次都和你玩一個遊戲.
我將這個人帶到你身邊.
你怎樣做.
其實每一次都可以選擇.
我可以說什麼.
不好意思啊.
我很累啊.
很忙啊.
遲些再說吧.
你可以這樣.
你又可以很細心認真地聽祂說話.
你又可以去想想.
究竟上帝想你怎樣去同行.
你可以去很多地方去想.
而最精彩的地方是什麼呢.
上帝竟然透過我們這些.
這麼差勁.
這麼微小的人.
你可以讓祂得到幫助.
多震撼啊.

$^{841}$我做了什麼.
所以有時候我說完之後.
有些弟兄姊妹說.
鄧傳道我很喜歡聽你講.
多謝.
不是我說得好.
我覺得是上帝的恩典.
祂透過我可以去感動弟兄姊妹.
讓他們能夠專心聽到哪裡.
和上帝的聖靈親自去改變他們的生命.
所以求神幫助我們.
在教會裡面.
我們越來越多的彼得.
我們越來越多.
願意去改變生命的彼得.
我們願意去用我們的生命.
去回應上帝的彼得.
我們要一起去建立上帝這個盆石.
這個教會.
這個地方.
求神帶領我們.
我們一起同心祈禱.
上天為我們求主.
你自己親自帶領我們.
最後其實我們在你面前.
我們只可以去盡力而為.
其實每一次當我們去做每一件事.
我們都不會知道那個果效.
但我們總是相信上帝在當中.
你將最好的東西.
你已經給了我們.
所以我們將你賜給我們的恩賜.
你賜給我們的資源.
我們能夠放出來.
我們能夠和人分享.
這個就已經足夠.
我只是求主.
你讓我們在你面前.
我們會看到的.
我們是上帝你自己去和我們說滿意.

$^{881}$我們要求的也不是人的讚譽.
而是我們在神面前的時候.
我們求你讓我們能夠.
專心回到上帝的面前.
我們求你借親自的幫助.
讓我們一直跟隨你.
讓你成為我們一生的主宰.
基督教主耶穌.
得勝名至祈禱.
多謝.
咱們就這樣吧 咱們就這樣吧 咱們就這樣吧.
\newpage



\section{約翰福音路加福音 23:32-49}
\label{sec:rYjqMvMSdYU}
\textbf{2021.04.04《我們所做的,我們知道嗎?》路加福音 23\_32-49;約翰福音19\_30( 和修版)  老耀雄牧師}
\newline
\newline
連結: \href{https://youtube.com/watch?v=rYjqMvMSdYU}{\texttt{https://youtube.com/watch?v=rYjqMvMSdYU}} ~~~~ 語音日期: 2021-04-06
\newline
\newline
\hyperref[sec:7LBy_VBf8iw]{\small{< < < PREV SERMON < < <}}
~
\hyperref[sec:index_chronic]{\small{[返順時目]}}
~
\hyperref[sec:index_scriptual]{\small{[返順卷目]}}
~
\hyperref[sec:q2mVvzFJiBo]{\small{> > > NEXT SERMON > > >}}
\newline
\newline
$^{1}$大兄姐妹平安.
如果手頭上.
剛才主席在讀聖經的時候.
你有打開聖經的話.
千萬不要收起來.
因為待會講道的時候.
都需要我們大家一起看著聖經.
一起去領受.
最近正在看一本.
由日本作家.
遠藤周作所寫的翻譯小說.
它叫做《沉默》.
故事講了兩個葡萄牙神父.
去日本找他們師父的經歷.
因為這兩個神父.
收到教廷的消息.
說他們的師父在日本.
放棄了信仰.
他所指的放棄信仰.
即是叛教.
他們很不明白.
他們的師父為何要這樣做.
所以他很想.
去日本找他的師父.
求證這件事情的真偽.
這本小說講道.
這兩個神父去日本.
親身體會到當時日本的.
鎖國禁教的策略.
而當時的幕府政權.
執政官井上.
如何用酷刑.
去迫害信徒.
而信徒在信仰生活.
和現實迫害的對比裡面.
有很多的掙扎.
以及他們對信仰的回應.
兩個神父經過很多次的逃避.
追捕的過程裡面.
他見到當時的信徒.

$^{41}$為了救他們兩個而犧牲.
其中一個神父向神埋怨.
主啊,我恨你.
我恨你一直都在沉默.
不過有個聲音.
他向這個神父回應.
他說,我不是沉默.
我是和你一起受苦.
另外,當神父最後選擇.
用腳去踐踏聖像.
向政權表明.
他都要叛教的時候.
他聽見來自聖像的聲音.
他說,踏下去吧.
踏下去吧.
你腳上那種痛.
我是最清楚的.
你踩下去吧.
我是為了要讓你踐踏.
然後才出生在這個世上.
是為了分擔你的苦痛.
才背起這個十字架.
我們所相信的神.
在我們受苦的時候.
祂並不是沉默的.
是和我們一起受苦的.
我們所相信的耶穌.
是要分擔我們的痛苦.
而走上十字架.
今天是復活節.
我們一起透過.
路加福音23章32-49節.
藉著耶穌在十字架上的一片經歷和片段.
去反思耶穌在十字架上的救贖和復活.
對我們帶來甚麼意義.
我們一起有個祈禱.
因為你是為我們而來到這個世界.
你也是為我們的緣故走上十字架.
承擔著我們的生命.
祂讓我們每一天所言所行.

$^{81}$都能夠與你恩典相稱.
不要做一個污辱主的聖徒.
我們的祈禱是奉靠主耶穌基督.
得勝明智祈求.
阿們.
如果我們仔細一點.
去對比四卷福音書的話.
我們一定會發覺到.
耶穌在被釘上十字架之後.
一共說了七句說話.
在路加福音記載了三句.
約翰福音又記載了三句.
而馬太和馬可都同樣記載了一句.
所以如果我們綜合四卷福音書的內容.
我們就可以歸納出.
耶穌在被釘上十字架之後.
祂說了七句說話.
我們稱作十加七言.
如果我們手裡拿著一本.
和合本修訂版的聖經.
在路加福音第23章34節那裡有一個括號.
而這篇聖經的左下角有一個注釋.
寫著有古卷沒有括號裡的文字.
這是甚麼意思呢?.
即是說當聖經翻譯的時候.
它手上的譯本是有34節這段經文的.
但是有些考古學的發現.
發覺有些抄本是沒有這一節經文的.
翻譯的學者.
不肯定這篇經文究竟是有還是沒有.
而這篇經文又不是可有可無的.
不能夠隨隨便便就剪掉了.
所以在修訂的時候就加了注釋.
而這篇經文正正是耶穌被釘上十字架之後的一句說話.
「父啊,赦免他們,因為他們所做的.
他們不知道」.
如果聖經真的沒有這節經文的話.
那就只有十架六賢.
就不是七賢了.
另外我們還需要去問一個問題.

$^{121}$就算真的有這段經文.
在這節經文裡面.
「赦免他們,因為他們所做的.
他們不知道」.
這三個「他們」.
是指哪些人?.
是指同樣被釘十字架.
在耶穌左右兩邊的犯人?.
因為他們對於自己所犯的錯事.
他們不知道.
所以耶穌要求天父赦免他們.
抑或是嘲笑和戲弄他的羅馬長官和士兵?.
他們不知道.
其實他們是釘了神.
在十字架上而不自知.
因此耶穌要求天父赦免他們的無知.
抑或是耶穌所指的是另有其人?.
會不會是要求釘死他的猶太人.
和大祭司蓋亞法和他們的祭司擋雨?.
又或者明知道耶穌沒有罪.
卻又將耶穌交給猶太人.
和判他釘十字架的羅馬巡撫彼拉多?.
我們試試根據路加福音23章32-49節經文.
從耶穌在釘十字架期間圍繞耶穌的人.
對耶穌的看法和反應.
從而估計一下.
耶穌在34節所說的「他們」.
究竟是甚麼人?.
我們剛才叫大家不要收起聖經.
我們現在一邊看經文一邊繼續分享.
在經文的觀察可以發現一共有五批人.
在第35至37節的官長和士兵.
士兵說:你若是猶太人的王.
救救你自己吧!.
在耶穌左上方有一個牌.
寫著「這是猶太人的王」.
士兵對耶穌的認識.
是從掛在耶穌頸上的牌而來.
因為掛在耶穌頸上的牌.
寫著「這是猶太人的王」.

$^{161}$所以士兵是根據牌上的資料.
對耶穌作出嘲諷.
士兵嘲諷耶穌.
你既是猶太人的王.
那就救救你自己吧!.
從士兵對待耶穌的態度可以知道.
他們既不相信耶穌的身份.
亦不相信耶穌的能力.
這種對人的態度.
其實源於無知.
因為他們對人的判斷.
只是從眼睛所見到的.
一種非常膚淺的認知.
一個無知的人.
是不是耶穌要求天父赦免的對象呢?.
那長官呢?.
長官應該是羅馬人.
他沒有猶太社群的信仰背景.
但是長官嘲笑耶穌說.
他救了別人.
他若是基督是神所揀選的.
救救自己吧!.
我們要問.
那些官長為何會知道.
耶穌是基督這個稱謂呢?.
又如何知道耶穌是神所揀選的呢?.
長官稱呼耶穌為基督.
相對於士兵稱呼耶穌為猶太人的王.
長官這個基督的稱呼.
更加貼近耶穌的真實身份.
而且長官不單單對耶穌的身份有較貼切的稱呼.
他對耶穌過往的行為.
亦比士兵有多一些的了解.
他說:他救了別人.
羅馬的長官知道耶穌救過別人.
可能他聽過耶穌曾經行過的神蹟.
好像拉薩路.
或者管會堂艾魯的女兒.
或者那因城寡婦的兒子.
這些事件都是耶穌能夠使人復活.

$^{201}$救過別人的真實事件.
當然如果計算耶穌醫病.
趕鬼的事就更加多.
這些使人復活的神蹟.
在加利利猶太地被傳開了.
很多人都會談論.
長官可能聽過猶太人.
談論耶穌行神蹟的事件.
而且可能聽過一些猶太人.
對耶穌身份的猜測.
耶穌曾經問過門徒.
他說:人說我是誰?.
大家就說:你是以利亞.
有些就說:你是耶利米.
有些說:你是先知中的其中一位.
耶穌再問:你們說我是誰?.
當時彼得這樣回應.
他說:你就是神所立的基督.
基督是永生神的兒子.
耶穌是基督.
是神的兒子這個身份.
被門徒確認了.
也被傳開了.
因此長官對耶穌的能力.
對耶穌身份的認識.
應該是從別人的口中聽回來的.
雖然長官聽過耶穌的身份.
也聽過耶穌的能力.
聽過耶穌曾經救過很多人.
但是他們根本不尊重.
也不相信耶穌的身份和能力.
或者換另一個說法.
長官根本認為和他沒有關係.
這只不過是猶太人的信仰.
耶穌就算是神.
也只是猶太人的神.
而我是一個羅馬人.
我的神就是凱撒.
否則他不會譏笑耶穌的身份和能力.
羅馬長官是一個異教的信徒.

$^{241}$是不是耶穌要求天父赦免的對象呢?.
第二批人.
在第33節.
39節和43節之間.
那兩個凡人.
和耶穌一起被釘十字架的有兩個凡人.
一個在耶穌的左邊.
一個在耶穌的右邊.
其中一個說.
你不是基督嗎?.
救救你自己我們吧!.
這個凡人對耶穌的認知.
只是根據官長和士兵的口述.
官長說耶穌是基督.
救過其他人.
這個凡人就以嘲笑的口吻.
你不是基督嗎?.
救救你自己和我們吧!.
他可能連基督這個身份.
代表什麼都不知道.
更加不知道耶穌以往.
拯救人的事件.
當然.
他也不會相信耶穌.
真的能夠拯救他.
他只是一個人魂亦魂.
再加上一份輕視.
和不相信的心態的一類人.
一個口中說耶穌是基督.
要求耶穌拯救.
但心裡完全不相信.
耶穌是拯救主的人.
是不是耶穌要求天父.
去赦免的對象呢?.
另一個凡人就不同.
他說你是一樣受刑的.
還不怕神嗎?.
我們是應得的.
因為我們是自作自受的.
但是這個人.

$^{281}$卻沒有做過一件不對的事.
這段說話表達了.
凡人心裡的一些看法.
我們是應得的.
因為我們是自作自受.
表明我們是有罪的.
但這個人沒有做過一件不對的事.
表明.
這個凡人心裡.
耶穌是沒有罪的.
你是一個受刑的人.
還不怕神嗎?.
表明一個有罪的人.
在神面前.
應該害怕的.
而耶穌就是神.
為什麼在耶穌面前.
你仍然毫無悔意.
毫無敬畏之心?.
一個凡人輕看耶穌.
對自己的罪毫無悔意.
面對神的兒子耶穌.
毫無敬畏之心.
一個凡人知道自己有罪.
所以對耶穌說.
耶穌啊.
你進入你國的時候.
求你紀念我.
一個很想被神紀念的人.
他的結局會是怎樣呢?.
耶穌對他的回應就是.
我實在告訴你.
今天你要和我在樂園裡了.
其實你想想.
這個犯人.
正在十字架上受刑.
應該在他人生裡充滿失望.
犯人面對死亡.
內心應該充滿了恐懼.
但是因為他很珍惜耶穌在他身邊.

$^{321}$他對耶穌有另一個體會.
這種體會有別於另一個犯人.
也有別於羅馬的士兵和長官.
雖然他和耶穌相遇的時間很短暫.
場景也不是一個很莊嚴神聖的地方.
不過他把握著這個機會.
他向神所立的基督.
永生神的兒子耶穌.
作出了一個請求.
這個承認自己的罪.
將自己的生命交託給耶穌的請求.
帶來的是耶穌對他作出了一個.
充滿永恆的應許.
你將會永遠與我同在.
這個犯人的經歷.
對我們有很大的提升.
第一.
就是時間上.
無論多麼趕急.
環境多麼不理想.
只要我們肯向神去求.
神會回應我們.
第二就是.
無論我們做錯了什麼.
我們都可以在神面前.
坦然地尋求幫助.
一個願意在耶穌面前.
認罪悔改的人.
是不是耶穌要求天父赦免的對象呢?.
第三批我們在經文能夠看到的人.
是伯夫長.
在47節.
伯夫長應該是一個羅馬人.
他對耶穌的認識.
有別於其他羅馬的長官.
他看到這一切的事情.
他就歸榮要給神.
為什麼要歸榮要給神呢?.
經文說.
他知道耶穌是一個沒有罪的人.

$^{361}$但卻被釘在十字架上.
所以他說.
這個人真是一個義人.
義人.
是可以指.
無可指責的意思.
這個伯夫長對耶穌的評價是.
耶穌是一個無可指責的人.
弟子們你想想.
這個評價.
和羅馬官方的判斷.
是一個很大的落差.
對一個羅馬人.
一個小小的伯夫長.
在羅馬的政權底下.
他堅持做一個真誠的人.
他敢說出事實的真相.
其實.
是很需要無比的勇氣.
對一個擇善而行的羅馬伯夫長.
是不是耶穌要求天父赦免的對象呢?.
第四批.
第四批人可以在第48節裡面我們看到.
就是.
另外還有第35節.
和35節的百姓是不是同一批人呢?.
他們是什麼人呢?.
經文沒有說.
是不是有要求釘死耶穌的猶太群眾在內呢?.
又或者.
有份審判耶穌的祭司和聖殿的領袖呢?.
我估計.
不會是這一批人.
經文說.
這些人都隨著空回去了.
看到整個過程而有這些反應的.
應該是同情耶穌的人.
他們認為耶穌是無辜的.
但卻被如此殘暴地對待.
只是因為他們無權無勢.

$^{401}$他不可以做到什麼.
所以只能夠隨著空回去.
他們可能聽過耶穌的事跡.
也可能認為他是一個先知.
或者是一個神人.
但他們對耶穌的反應.
就是面對強權.
只能愛莫能助.
他們有良知.
但沒有能力作出回應.
或者正確地說.
他們不敢去回應.
因為他們是一群弱勢社群.
一群弱勢社群.
是不是耶穌要求天父赦免的對象呢?.
最後一批.
在第49節.
所有與耶穌熟悉的人.
是從加里利跟著他來的婦女們.
經文已經說了.
這群人是最熟悉耶穌的人.
也有一些是從加里利.
跟著耶穌進入耶路撒冷的婦女.
這批很熟悉耶穌的人.
遠遠地站著看這件事.
他們與耶穌應該有一個很親密的相處.
他們心裡想些甚麼呢?.
我們估計一下.
這群婦女有些甚麼人.
應該有拉撒路的姐姐.
瑪大的妹妹.
那個用純拿撻拉香膏.
去告耶穌的那位.
莫迪拉的瑪利亞.
她見證耶穌.
將拉撒路從死人中復活過來.
她用最好的獻給耶穌.
她揀選了上好的福分.
就是坐在耶穌身邊.
聽耶穌講解真理.

$^{441}$應該還有西比泰的兩個兒子.
雅各和約瑟.
兩個兄弟的母親瑪利亞.
她曾經叩頭請求耶穌.
在這個角裡面.
請讓我這兩個兒子.
一個坐在你右邊.
一個坐在你左邊.
她是一個偉大的母親.
她為兩個兒子.
求在天上的榮耀.
她絕對相信耶穌.
是有權柄和能力的.
雖然她不明白.
坐在耶穌的左邊和右邊.
真正的意義是甚麼.
她也不明白她所求的.
是需要付上代價.
但起碼.
她讓那兩個兒子.
去跟從耶穌.
也不阻止兩個兒子去跟從.
另一個肯定在現場的.
一定有耶穌的母親瑪利亞.
瑪利亞從聖靈懷孕.
去到耶穌出世.
東方博士來訪.
耶穌十二歲在聖殿走失.
之後耶穌出來事奉.
在迦納見證耶穌將水變為酒.
所有事情發生之後.
經文都用同一句說話去形容.
瑪利亞把這一切的事.
存在心裡.
作為一個記錄.
瑪利亞作為耶穌最親近的人.
見證耶穌的一言一行.
今天耶穌被釘十字架.
她有甚麼感受.
這一群和耶穌有親密關係的婦女.

$^{481}$是否耶穌要求天父赦免的對象.
在這五批人當中.
有對耶穌完全不認識.
完全沒有關係的.
有一些對耶穌有少許認識.
但是關係不深的.
有一些對耶穌有深入的認識.
並且有很親密的關係的.
天父赦免他們.
因為他們所做的.
他們不知道.
究竟哪些人.
是耶穌要求天父赦免的對象呢?.
在三十二至四十節裡.
還有一節經文.
可以給我們多一個參考.
在耶穌斷氣之前.
第四十五節說.
聖所裡的象脈.
從中裂開了.
這節經文的概念.
就是人再不受到這幅掛在至聖所的象脈所隔了.
人再不需要透過大祭司進入至聖所.
去與神相遇.
象脈從中裂開.
表示人可以直接來到神面前.
耶穌走上十字架的目的.
就是要所有人都可以與神相遇.
所以如果加上約翰福音.
十九章三十節那句「成了」.
我們就可以有一幅比較清晰的圖畫.
十加七言的最後一句就是「成了」.
「成了」代表了耶穌完成了天父交給他.
要在地上達成的使命.
這個使命就是要承擔世人所有的罪.
以致人能夠與神復和.
「成了」就是表示所有得罪神的人.
都藉著耶穌在十字架上的救贖完成了.
耶穌在世上所受的苦.
被人誤解,敵視,詛咒,侮辱.

$^{521}$被鞭打,被釘死.
一切一切都是源於人的罪.
罪進入了世界.
罪使人蒙蔽了自己的眼睛,良心.
罪讓惡勝善.
罪將藏在人的內心.
屬於神的那種真善美掩蓋了.
罪使我們以自己為中心.
以致人所做的,所說的.
既得罪神又得罪人而不自知.
保羅在《羅馬書》如何形容這批被罪蒙蔽的人呢?.
他說:「他們雖然知道神.
卻不把他當作神榮耀他.
他們的思想變為虛妄.
無知的心昏暗了」.
當時羅馬的士兵,官長,巡撫彼拉多.
猶太公會的領袖,聖殿的大祭司.
法利賽人,薩都該人.
這班人正正是保羅所說的「他們」.
他們雖然知道神.
卻不把耶穌當作神榮耀他.
耶穌在十字架上所做的一切.
就是讓罪不再克制我們.
讓我們能夠得釋放.
可以過一個自由的人生.
可以與神同行.
而不被神所擊打的人生.
而耶穌的復活.
卻是印證這個救贖所帶來的真實.
可以這麼說.
耶穌向天父祈求.
「父啊,赦免他們.
因為他們所做的他們不知道」.
裡面的「他們」.
就是指所有被罪所蒙蔽的.
包括那些士兵,長官,巡撫彼拉多.
猶太公會領袖,聖殿的大祭司.
法利賽人,薩都該人.
以及當時被罪惡所捆綁的一切人.
不單止當日.

$^{561}$就算是今天.
耶穌在十字架上的救贖.
對今天仍然被罪所捆綁的人.
仍然有效.
同一家弟兄姊妹.
今天的講題是.
我們所做的我們知不知道?.
透過路加福音所描述的五類人.
我們有沒有想過.
我們現在的信仰光景.
和哪一類人最相似?.
我們知不知道.
今天我們所做的.
我們所講的.
究竟是土神喜悅.
還是土神厭惡呢?.
我相信我們和耶穌的關係.
不會像羅馬長官和丁丁那樣.
起碼我們口中.
稱呼耶穌為主.
但會不會像另一個犯人那樣.
口中雖然稱呼耶穌為主.
但心中卻沒有神呢?.
心中對神是毫無敬畏的心呢?.
抑或我們是那種.
很想被耶穌拯救.
但卻不肯交出主權的人呢?.
如果從好的方面想.
我們最好就像.
梅大拉的瑪利亞.
我們能夠領受上好的福分.
用最好的香膏獻給耶穌.
抑或我們好像.
雅各和約瑟的兩兄弟的母親.
瑪利亞一樣.
我們不會拒絕.
我們的子女.
完全奉獻給主耶穌.
在今天.
有很多母親.

$^{601}$都不想子女.
全然奉獻.
你讀完大學.
就去再讀神學嗎?.
你在教會的薪水.
養不了家庭.
當今天我們處身在一個.
彎曲末謀的世代裡面.
我們會不會像一個白膚匠一樣.
在強權底下.
我們仍然願意做一個真誠的人.
去講出真理.
抑或我們苦心自問.
我們都是一個弱勢社群.
我們人微言輕.
對於發生在面前的事.
我們不聞不問.
不敢過問.
只是扔著胸口回家就算了.
我相信今天我參與崇拜的朋友.
正在看直播的弟兄姊妹.
你們大家都是已經相信了耶穌.
但我們有沒有反省過.
耶穌在十字架上為我們承擔的罪.
我們今天有沒有明知故犯地繼續犯下去.
我們所做的.
是神所喜悅的.
還是神所厭惡的呢?.
耶穌走上十字架的目的.
就是讓我們不被罪惡所捆綁.
今天我們問問自己.
我們有沒有被罪惡所捆綁.
為何我們今天仍然被罪.
被自己的私慾所捆綁.
我們需要很認真地反思一下.
我們在家庭,工作場所.
甚至在教會.
我們所做的.
有多少是從信仰出發.
有耶穌基督的心意.

$^{641}$有多少是被這個世界價值觀所影響.
有多少是被自己的私慾所控制.
為了去預備這個講章.
我就找了Facebook一些資訊.
在2004年.
米露吉遜拍攝了一套著名的宗教電影.
《受難曲》.
雖然這套電影有很多的教義.
都引發了很多爭論.
但是他有一個反省.
他說耶穌的犧牲能夠喚醒了容忍.
喚醒了愛和寬恕.
這些都是現今世界所需要的.
今天我們的生命.
有沒有彰顯出主耶穌那種容忍.
愛和寬恕呢?.
另外作為演員的卡維素.
他說每一次拍攝.
耶穌被釘上十字架的時候.
他就有一個很深的體會.
他說當我們每一次犯罪.
我們就好像有份將耶穌再交出來.
重釘十字架.
好像希伯來書四章六節這樣說.
我們每一次犯罪.
就將耶穌再釘多一次十字架.
但是每一次當我們去背叛耶穌.
耶穌卻不斷地叫我們重新起來.
呼求我們.
要求我們千萬千萬千萬不要放棄倚靠.
兄弟姐妹.
我們一起合上眼.
沉澱一下剛才的訊息.
我帶大家一個審查的時間.
主耶穌求你讓我們不單止在受難節.
或者復活節的時候.
我們先去審查你為我們所受的苦.
讓我們每時每刻.
我們都要紀念你為我們所做的一切.
主耶穌求你讓我們時刻記得.

$^{681}$你的恩典不是廉價的.
而是貴價的恩典.
這個恩典是以你的血和你的生命救贖回來的.
主耶穌求你讓我們記得.
當初我們說相信你的時候.
我們承認你是我們生命的主.
承認你是我們救贖的主.
我們實實在在需要將我們生命的主權交給你.
主耶穌讓我們深深地去明白.
信仰的根基要建立在盤石上.
而不是建立在沙土上.
因為風吹雨打水沖撞擊了房子.
防止腫不倒塌.
求你讓我們聽到而行到.
因此主耶穌今日此時此刻.
當我們反思我們的信仰生命的時候.
我們需要好像浪子一樣.
要和主耶穌說.
我們得罪了天又得罪了你.
今日此時此刻.
當我們反思我們所言所行所信的時候.
我們是否需要和天父說.
父啊赦免我們.
因為我們所做的我們不知道.
親愛的主你已經許過每一個屬你的兒女.
我們若認自己是罪.
你是信實的是公義的.
必要赦免我們的罪.
洗淨我們一切的不義.
蒼天造地的天父.
在十字架上為我們承擔一切罪責的主耶穌.
每日用使不出的嘆息為我們代求的聖靈.
我們就憑著你的應許.
坦然無懼的來到你面前.
承認我們每一個人的罪.
求你赦免我們.
求你堅固我們每一個願意認罪悔改.
願意跟從你的心.
縱然我們軟弱.
但主你知道.

$^{721}$我們愛你的心從來都沒有改變.
主耶穌我們愛你.
誠心所願.
阿們.
願主祝福大家.
\newpage



\section{約翰福音 6:7-9}
\label{sec:q2mVvzFJiBo}
\textbf{2021.04.11《一位平凡但重要的使徒--安德烈》約翰福音 6\_7-9( 和修版)  郭鴻標牧師}
\newline
\newline
連結: \href{https://youtube.com/watch?v=q2mVvzFJiBo}{\texttt{https://youtube.com/watch?v=q2mVvzFJiBo}} ~~~~ 語音日期: 2021-04-11
\newline
\newline
\hyperref[sec:rYjqMvMSdYU]{\small{< < < PREV SERMON < < <}}
~
\hyperref[sec:index_chronic]{\small{[返順時目]}}
~
\hyperref[sec:index_scriptual]{\small{[返順卷目]}}
~
\hyperref[sec:D7R35FEAwH0]{\small{> > > NEXT SERMON > > >}}
\newline
\newline
約翰福音 6:7-9
\newline
\begin{longtable}{cl}
\hline
\hline
章節 & 經文 (和合本修訂版)\\
\hline
6:7 & \begin{tabularx}{0.7\textwidth}{X} 腓力回答他：「就是兩百個銀幣的餅也不夠給他們每人吃一點點。」 \end{tabularx} \\ \\ \relax
6:8 & \begin{tabularx}{0.7\textwidth}{X} 有一個門徒，就是西門‧彼得的弟弟安得烈，對耶穌說： \end{tabularx} \\ \\ \relax
6:9 & \begin{tabularx}{0.7\textwidth}{X} 「這裡有一個孩子，帶著五個大麥餅和兩條魚，但是分給這麼多人還算甚麼呢？」 \end{tabularx} \\ \\
[1ex]
\hline
\hline
\end{longtable}
$^{1}$我們一起同心祈禱.
我們親愛的天父上帝.
我們要向你敬拜.
向你讚美.
因為你是神.
在我們人生的旅途裡面.
我們有很多不同的體驗.
有甜酸苦辣.
有順利有不順利.
我們向你禱告.
就是因為我們人是很渺小.
我們很需要依靠你.
有時我們用人的想法去解決問題.
但很多時候是解決不到的.
我們很需要你給我們有智慧.
我們很需要你給我們的內心.
有平安.
有滿足和喜樂.
有世界給不到的那種安全感.
我們就願你的說話.
成為我們生命的力量.
願你的靈在我們當中工作.
讓我們明白你話語的寶貴.
這樣的祈禱是奉耶穌基督的名.
阿們.
很感恩.
原來今天是洪恩嘉.
恢復一個現場實體崇拜的大日子.
我相信弟兄姊妹帶著一個期待的心回來.
我今天很榮幸.
剛好是一個這樣的日子.
和安排弟兄姊妹在元朗西鐵站接我過來.
真的非常好.
我自己這段時間經常都有看聖經.
看聖經對我來說是要很多考驗的.
為什麼呢.
我的問題可能以為自己很熟悉.
就是因為如果我以為自己很熟悉的時候.
我就會忽略很多東西.
很感恩.

$^{41}$我讀聖經.
大家應該聽過五餅二魚的神蹟.
你也很熟悉的.
我今天就是講這個題目.
這段經文.
但是我這次講的就跟我平時的發現有不同.
我說的話都是跟著聖經的.
不會亂來的.
但是你試一下聽一下.
我發現了什麼.
我也很有得著.
我的經文就選了約翰福音第六章第七到第九節.
其實講到的時候.
我用的經文是多過這兩三節的.
因為我恐怕主席讀經文是用很多時間去讀的.
是這樣的.
但是題目是怎樣呢.
是一個平凡但是重要的使徒.
安德烈.
你想一下十二使徒的名單.
你最熟悉的是哪幾個.
彼得.
約翰.
亞國等等.
其實安德烈你們熟不熟悉呢.
不熟悉呢.
這個就是有趣的地方.
我講到有三個部分.
第一個部分.
原來安德烈是耶穌第一批呼召的使徒或者門徒.
原來安德烈很資歷深厚.
第二個部分就是.
原來安德烈是五餅二魚神蹟的促成者.
當然你要知道這個世界沒有說誰不行的.
神做事不需要我們.
但是事實上五餅二魚的神蹟.
就是因為安德烈.
他的角色.
令到這件事實現了出來.
好了講了這些經文.

$^{81}$到了第三部分.
我就應用了.
我們去學習安德烈.
將人帶去耶穌基督面前.
這個就是我講到的結構.
首先你看看.
安德烈是耶穌基督第一批呼召的門徒.
我就嘗試看幾卷科音書.
因為不是只看一卷.
我先講馬可福音.
第三章十三到十九節.
有一個十二門徒的名單.
我先講名單.
排名呢.
安德烈就排在彼得亞國約翰之後.
你當好奇看看他怎樣排.
好了.
他排第四.
如果我們看馬太福音.
第十章第一到第四節.
那個十二門徒的名單.
安德烈就排在彼得之後.
厲害了排第二.
如果我們看路加福音第五章.
就講耶穌上了彼得的船.
意思是什麼.
只是提到彼得.
但沒有提到安德烈.
我們看路加福音第五章第十節.
他這樣講.
他的夥伴西比泰的兒子亞國約翰也是這樣.
然後耶穌對西門說.
不要怕.
從今以後你要得人了.
好了.
到路加福音第六章.
才開始提到安德烈這個人.
路加福音第六章十四節.
這十二個人有西門.
耶穌又給他起名叫彼得.

$^{121}$還有他兄弟安德烈.
這句話很有趣.
講完彼得就說.
還有他的兄弟安德烈.
又有亞國和約翰.
費力和巴多羅買.
就是說.
安德烈其實在這幾卷福音書的十二門徒名單裡面.
都不是很重要.
反而重要的是誰呢.
就是彼得.
他總是排第一.
但是當我們看約翰福音第一章三十七到三十九節.
你有另外一個發現.
我就用舊的和合本.
他說兩個門徒聽見他的話.
這兩個門徒是誰的門徒.
是施洗約翰的門徒.
這兩個門徒聽了施洗約翰講的話.
就跟從了耶穌.
你又想施洗約翰見到耶穌來了.
對他兩個門徒說.
這個就是主耶穌了.
這個是神的兒子.
耶穌轉身來了.
就見到他們兩個跟著.
就問他們說.
你們要什麼.
他們說拉比你在哪裡住.
耶穌說你們來看.
他們就去看他在哪裡住.
那一天就跟他同住.
那時約有新政了.
就是說這兩個門徒跟了耶穌.
去耶穌住的地方.
約翰福音記載.
原本是施洗約翰的兩個門徒.
現在就跟了耶穌.
究竟這兩個人是誰.
我們就要看約翰福音第一章四十到四十二節.

$^{161}$他說聽見約翰的話.
跟從耶穌那兩個人.
一個就是西門彼得的兄弟安德烈.
留心後面說什麼.
他先找著自己的哥哥西門.
對他說我們意見彌賽亞了.
你這樣看誰是第一個跟了耶穌.
起碼安德烈還有另外一個人.
聖經學者大部分都覺得那個人就是約翰.
就是說安德烈和約翰.
應該是耶穌第一批呼召的門徒.
也有聖經學者說.
其實可能安德烈是第一位.
我們不要斟酌他是第一還是第二.
總之是最早期的那兩個門徒.
馬太科音,馬可科音,路加科音的門徒名單.
他因為要凸顯彼得的領導地位.
他就不是按著時序去排列.
大家明不明白.
不是按著時序的發展去排列.
其實如果你跟著漢科音.
你會看得很清楚.
安德烈其實是耶穌第一批呼召的門徒.
當然你可以說.
他雖然是第一位或者第二位.
但是介紹他的時候都挺.
你可以這麼說.
有些人他有兩兄弟兩姐妹.
如果讀同一間學校的時候.
做老師的那些.
眼看就立刻有共鳴.
做老師的時候.
很可能最深印象的是哥哥或者姐姐.
接著哥哥姐姐就有個弟弟有個妹妹的時候.
接著就會這樣想.
你陳大文.
你就是陳大大的弟弟.
你永遠都是在你哥哥的影子底下.
或者是你永遠在你姐姐的影子底下.
人家認識你就是要這樣介紹.

$^{201}$現在安德烈就是這樣.
你看看他幾卷科目書都是介紹這樣的.
介紹他還要搬他哥哥彼得出來才認識他.
其實你是安德烈你會怎樣.
其實都不是很快樂.
但是如果你留心.
原來安德烈是耶穌最早呼召的門徒.
只不過他可能領導因此不是很強.
科目書的作者就是這樣記載他.
當然有另一次記載是厲害一點的.
馬克科音十三章第三節記載.
安德烈,彼得,雅各約翰跟耶穌上去橄欖山.
這樣都厲害了他排第四.
安德烈是核心三人組外面的第四位.
彼得,雅各約翰是核心三人組.
安德烈是第四位.
即是叫mung 車邊.
我這樣介紹安德烈出來給大家.
你就覺得他很平凡.
很普通的不是耶穌最近身得力的弟子.
不是彼得,雅各約翰就是.
但是安德烈是外面一點.
這個外面一點的人.
很容易被我們忽略.
但如果我再看第二部分我要講的五餅而喻的神蹟.
安德烈就是那個速成者.
你細心看你會發覺安德烈其實很厲害.
為什麼呢.
首先約翰福音第一章43到44節.
記載了耶穌呼召完安德烈.
和呼召完彼得之後.
呼召一個人叫肥力.
你記住彼得也是後面的.
即是他弟弟帶他信主.
在約翰福音一章43到44節就講了.
有次人耶穌想要去到加利利去.
就遇見到肥力對他說.
你來跟從我吧.
這個肥力是什麼人呢.
是白菜大人和安德烈和彼得是同城.

$^{241}$即是同鄉.
肥力都是較早被耶穌呼召的一位.
起碼都是前面那些.
不過始終彼得亞國約翰是最厲害.
現在我講到肥力.
一會兒就出來了.
經文提到肥力和安德烈和彼得是同鄉.
這樣就進入伍秉儒的神蹟的故事.
在這個記載裡面.
我想大家留心.
肥力和安德烈兩個人有什麼不同的反應.
耶穌叫他們做事.
但兩個人做事的反應是怎樣呢.
當時那班人相傳有五千人聽耶穌講到不願走.
都開始晚了.
耶穌就想想怎樣.
要不解散他們.
要不就想辦法給他們吃東西.
然後他就叫門徒去附近買東西.
這個是不是很高難度的動作呢.
很難的.
大家留心肥力怎樣回答.
剛才介紹了肥力和安德烈和彼得是同鄉.
肥力就說.
就是二十兩銀子的餅.
叫他們每人吃一點都不夠.
這個是肥力的反應.
然後後面那位有一個門徒.
經文的記載真是.
他這樣介紹安德烈.
有一個門徒.
他就是西門彼得的兄弟安德烈.
你可以明白安德烈也是阿四.
他就對耶穌說.
這裡有一個孩童.
他帶著五個餅兩條魚.
他都叫做提供資料.
不過他都補一句.
只是分給這麼多人吃算得了什麼.
肥力其實說得對.

$^{281}$首先是說要給五千人吃飽.
二十兩銀子.
等於當時是怎樣.
當時的工人兩百天的工資.
近七個月的人工.
你說要籌到那二十兩銀子.
其實都不容易.
就算你籌到回來.
你去買那餅.
買到了給你.
你夠不夠給五千人分呢.
分到了.
每人只有一點點.
仍然不夠吃.
肥力一開始已經分析過整個情況.
是沒什麼作用的.
經文就跟著說有一個叫安德烈.
安德烈就剛才說了.
有一個門徒.
他是彼得的兄弟.
他怕那些人不知道誰是安德烈.
要提一個有名望的彼得出來.
再要說明他是彼得的兄弟.
然後就緊接著說.
安德烈就對耶穌說.
這裡有一個孩童.
帶著五個大麥餅.
兩條魚.
只是分給這麼多人吃.
算得了什麼呢.
在這裡我想大家想一想.
經文就很簡短.
我想問大家.
你有沒有想過.
究竟為什麼安德烈知道.
有個小朋友有五個餅.
兩條魚呢.
他憑什麼知道呢.
很大的可能性.
他當然到處去問.

$^{321}$他不問怎麼會知道呢.
沒理由他有特異功能.
知道有個人有五個餅.
兩條魚.
很合理的.
你還要想一想.
他去問當然是先問大人.
如果讓你安德烈.
你會不會這樣.
你走去小朋友那裡問.
不會的.
你去問大人.
你有沒有餅.
有沒有魚.
有什麼食物.
但是現在的答案.
不是由大人提供的.
很有趣的.
是小朋友提供的.
為什麼小朋友提供呢.
真的不知道.
經文沒有提到.
但是你去想像一下.
其實如果安德烈.
去到這樣的情況.
你都很為難的.
如果小朋友說.
我有餅.
有魚.
我可以給你.
你都不是很想要的.
問一下父母.
你們自己處理他吧.
小朋友.
但是很奇怪的.
安德烈.
他照樣將這個消息.
告訴耶穌.
我想大家想一想.
換成是你.

$^{361}$你會不會覺得.
將這個消息告訴耶穌.
耶穌會不會覺得我傻傻的.
或許當我告訴耶穌的時候.
我周邊的門徒都說.
你是不是病了.
可能會覺得.
你是不是真的沒有腦子.
大人的東西你不要.
要小朋友的東西.
還有你拿五餅魚來幹嘛.
哪裡夠吃.
其實你有沒有想過.
安德烈要開口告訴耶穌.
有這個消息.
他都有壓力的.
是不是.
不過你可以這樣說.
安德烈的領導能力是一般般.
他想事情都是一般般.
耶穌叫我.
我問到什麼.
就告訴耶穌.
他真的這麼簡單.
他真的照樣說.
他就告訴耶穌.
這裡有個小朋友.
有五個餅.
兩條魚.
他很誠實.
不過分給這麼多人吃.
哪裡夠吃.
你有沒有留心.
安德烈和肥力兩個人.
有什麼不同反應.
肥力就很會想事情.
其實不夠.
你拿兩百.
二十兩銀子去買.
分給五千人吃.

$^{401}$根本吃不夠.
經文沒有說他有沒有做事.
是不是.
又沒有說他有沒有去問.
是不是.
但是經文就說.
安德烈走去.
起碼知道有個小朋友有五餅魚.
還有他走去告訴耶穌.
這個就是兩個人之間的分別.
很細微的.
不知道你有沒有發覺.
如果安德烈一早就說.
那兩個魚.
五餅沒用的.
那我不如就不要告訴耶穌.
那你想想.
這個五餅魚的神蹟會不會發生.
如果他不說.
那你可以說.
可能有第二個人說也說不定.
但起碼他不說.
他不會告訴耶穌.
所以我想.
整件事我會覺得.
一件事就是說.
其實安德烈他的反應.
其實他也經過考慮和掙扎.
但他最後告訴耶穌.
這件事就真的讚他.
我覺得在我自己看經文裡面.
安德烈他的才華不是很多.
但他有件事.
就說些東西給耶穌知道.
是不是.
這個信息.
其實按他的判斷覺得沒什麼用.
但他照樣告訴耶穌.
好了.
還有另一件事.

$^{441}$第三個應用.
就是學習安德烈將人帶到耶穌面前.
當我讀聖經的時候.
我發覺有另一段經文.
記載安德烈的事.
其實福音書介紹他的東西不多.
這件事是什麼事情呢.
他.
我覺得他厲害在什麼呢.
懂得轉告消息.
不是講閒話那些.
是轉告重要的消息.
在約翰福音第一章40到42節.
剛才我讀過一次.
聽見約翰跟從耶穌的兩個人.
一個就是西門彼得的兄弟安德烈.
他先抓住他的哥哥西門.
對他說我們遇見彌賽亞了.
這個你有沒有看見他轉告的恩賜.
他第一件事.
他遇見彌賽亞之後.
他回去跟他的哥哥說.
我見到彌賽亞了.
大家有沒有留心.
後來記載彼得是顯赫過安德烈的.
安德烈有件很厲害的事.
就是找到一個比他更厲害的人去跟耶穌.
你有沒有看到這件事.
其實這件事對我都很大提醒.
就是說.
我們可能會覺得自己都是普通人.
就是很平凡.
但是你總是覺得很多人很厲害.
在你身邊很多人比你厲害.
不要緊.
你就是將這個轉告的恩賜發揮.
找一些更厲害的人去跟耶穌.
大家有沒有想過這件事.
不知道有沒有想過.
有時會覺得.

$^{481}$為什麼神你做人做得這麼不公道.
那個人這麼厲害我這麼差.
會不會有這個想法.
有都不奇怪.
但是如果你看回聖經的時候.
你可能學一下.
其實厲害不厲害都不是我選的.
不過不要緊.
我可以學安德烈這樣.
去介紹一些比我更厲害的人去認識耶穌.
我將他帶去耶穌面前.
我都很厲害.
大家同不同意這個說法嗎.
你不要當他是阿Q.
而是安德烈的恩賜.
好了.
安德烈就轉告彼得他見到彌賽亞.
還有另一件事剛才都講了.
約翰福音六章八到九節就說.
有個門徒就是西門的兄弟安德烈.
就對耶穌說.
這裡有個孩童帶著五個大麥餅兩條魚.
只分給這麼多人吃算是什麼.
他轉告一件事給耶穌知道.
耶穌吩咐他去找東西吃.
但是他轉告給耶穌知道.
有個小孩有五餅二魚.
還有一件事.
我讀完這段經文.
我真的發覺很大得著.
約翰福音十二章二十到二十節.
就說那時候上來過節的人中間有幾個希利利人.
就是希臘人.
他們來到加利利的白菜大那個肥力那裡.
肥力又出來了.
就對他說.
先生我很願意去見耶穌.
肥力就去告訴安德烈.
安德烈就和肥力去見耶穌.
很簡單的一段經文.

$^{521}$但是當我再發現的時候.
其實人與人之間是有個密切關係的.
其實雖然肥力和安德烈和彼得是同鄉.
但是肥力有個老死.
講得俗一點.
有個friend.
那個friend是誰.
就叫拿彈業.
為什麼我這樣說.
在約翰福音一章四十五節.
肥力抓著拿彈業對他說.
摩西在律法上所寫同眾先知所記那位.
我們預見了就是約瑟的兒子.
拿撒勒人耶穌.
肥力都有個friend.
就是拿彈業.
我就問一個問題.
其實這件事有希臘人來說想見耶穌.
如果你是肥力會不會馬上叫你的friend.
拿彈業我們一起去見耶穌.
很正常的.
但是這次不是.
肥力他竟然找了安德烈.
以下那些是我的推論.
為什麼很簡單的道理.
因為安德烈已經有一個很眾人發現的因子.
就是轉告的因子.
如果我有希臘人要我帶他去見耶穌.
我當然不找拿彈業.
拿彈業不是這樣強項.
反而安德烈是.
我當然拉著他.
安德烈我們去見耶穌.
這裡有幾個希臘人.
所以我看完之後我發覺.
這個安德烈原來有些斤兩.
雖然他在十二門徒的名單裡面.
最厲害是有時候排第四有時候排第二.
其實他根本不是核心三人組.
但是他有一樣東西真的與眾不同的就是.

$^{561}$他可以見到耶穌之後.
他將這個消息告訴他的兄弟朋友知道.
或者耶穌叫他做事.
他就將那些消息又告訴耶穌知道.
或者有希臘人來想見耶穌.
他又去華耶穌那裡.
將這些人帶去耶穌面前.
所以我會覺得這個安德烈有些東西很值得學習.
我看回歷史資料.
原來安德烈大概在主後六十年殉道.
安德烈是一個非常忠誠的信徒.
他的見證是很值得我們學習的.
大家同不同意這件事.
他真的很特別.
彼得是一個偉大的領袖.
當然值得我們學習.
但是大家都明白.
我們不是每個人都像彼得那麼威風.
我們都是平凡人.
所以我讀到這裡的時候.
我就覺得安德烈他不是核心領袖.
但是他最厲害是忠心的轉告信息給有關的人.
無論是耶穌基督無論是他的朋友他的兄弟.
就令一件事情更大的事可以成就.
大家明不明白這件事.
安德烈所做過的事.
就令到更大的事成就.
他是一個配角.
或者俗話說的叫做搏腳.
他是這樣的.
我將他轉告的動作.
就指向一個方向.
就是他將人帶到耶穌基督面前.
所以我讀完這段經文.
這樣去處理完之後.
我發覺我得到很多東西.
我以前沒有.
我看五餅二魚的記載裡面.
我都沒有這樣的領袖.
其實我告訴大家.

$^{601}$如果一個傳道人.
看聖經有領袖是很開心的.
為什麼呢.
我們經常看.
日看夜看.
又看英文又看中文什麼都看.
其實看得多.
你講得不好聽.
有些人做那一行會厭那一行.
經常都看到悶.
但是當你看聖經有種收穫的時候.
你就會很感恩.
你明白我的意思嗎.
這件事我就覺得神有教導.
我們怎樣去學安德烈呢.
這個就是實際的問題.
我想我就會很直接地說.
安德烈他做的事情.
他不是將人帶到人的面前.
他是將人帶到耶穌基督的面前.
這句話希望大家慢慢咀嚼.
因為你將人帶到人的面前.
你只是用人的眼光.
人的方法去看那件事.
大家都很明白.
人一定會重複犯錯.
根據一些不準確的理論去處理問題.
結果就不斷在漩渦裡面打轉.
找不到出路.
這個是人間的事情.
人間一定是這樣的.
在新冠病毒的衝擊下.
其實教會面對很多考驗.
就是牧者和信徒領袖都沒有經驗.
怎樣去處理網上崇拜.
團契 主日學 祈禱會.
這些根本我們都不懂的.
我們現在是一個很艱難的年代.
你能夠守住目前的狀態已經是很厲害了.
哪裡還有力量去開拓呢.

$^{641}$我們很需要神賜下屬靈的分辨能力.
分析目前這個複雜的情況.
神學院有很多同工舉辦工作坊 講座 研討會.
最近有人提到教會的人力和財力的下跌.
這些事情.
大家又根據數據去分析等等.
有些牧者和我溝通.
就讓我學會很多不同的知道情況.
對當前的環境知道多一些.
今早魯牧師又跟我聊天.
又告訴我洪恩堂近年的發展.
我聽完很開心 為什麼呢.
那個社區的服務又很好.
也接觸了很多不同的家庭.
聽完之後我就為你們感恩.
其實你知道.
當教會恢復有現場實體的崇拜.
就好像學生考試一樣.
其實是派成績表的意思.
現在有很多人擔心.
恢復現場實體崇拜之後.
會流失多少人.
很多教會都在頭痛.
見到有些人還在.
其實是很感恩的.
不是說笑.
這是一個考試.
考試考什麼呢.
考試考你之前那段時間有沒有做關心的工作.
我認識一些傳道人.
真的粉墨登場.
首先有一個女傳道.
她的教會有二百個長者.
她的教會很大.
但那些長者不懂得上網看電視的崇拜.
不懂的.
那她怎麼辦呢.
她用手提電話認識那些長者.
手提電話有WhatsApp.
她就錄音錄影自己拍片.

$^{681}$60分鐘零收等等.
就在群組發放.
但還有些老人家連這個智能電話都不用.
那她怎麼辦.
她就剔出來有多少個沒看過我那些.
錄音影片那些.
親自打電話給她.
做到這樣.
很忙碌的.
有一個傳道人.
她的年紀都五十幾歲.
她沒什麼頭髮.
她本來做長者部.
但教會沒有了.
很小.
她的教會兒童崇拜都沒有了.
結果她粉墨登場.
就去拍片.
但因為她做長者侍工.
說話是比較慢一點.
比我更慢.
所以拍完之後.
弟兄姊妹就很有感動.
傳道人我來幫你拍.
你明白什麼意思.
會令到拍攝效果更好.
這樣都好.
兒童崇拜都有網上版.
所以我見到很多事.
神正在工作.
說了這麼多之後.
其實我講回我自己.
我寫下了.
我成長的時候.
在一間美國宣教士的教會成長.
那時都是十幾歲.
都不知道做什麼事.
那些西教士對我們青年人.
很平等.
我們叫他名字.

$^{721}$無大無小.
後來在教會去參與事奉.
所以我整個習慣上.
在信主之後就會事奉.
我不喜歡坐在那裡.
不知為什麼.
我不太能夠坐在那裡.
所以有很多事奉的機會.
我自己講近一點.
就是1995年的時候.
當時在德國.
不知為什麼有一個很強烈的感動.
要成為牧師.
那時我都經常講道.
不過我嫌教會很麻煩.
老實說.
我很年輕就知道教會很複雜.
我可以去講道.
但你不要叫我做教會牧師.
很頭痛.
但1997年回來香港教神學.
2000年開始幫教會去牧羊.
其實不知為什麼.
你見到那個tum 自己踩下去.
但很開心.
過去20年事奉的體驗.
我發覺牧者領受了一個屬靈的即時.
更加重要的是.
神的呼召.
還有很多比我厲害的人.
他放下了事業發展的機會去事奉神.
所以我自己的理解就是.
牧者和信徒領袖一起去領導教會.
以屬靈方向由牧者去帶領.
因為他更加敏感到神的工作.
我自己用傳道人的身份參與教會事奉.
其實我全職是神學院的.
大家都知道我在長洲住.
這個就是我自己的限制.
所以我只主力在講道,教道.

$^{761}$很少去行政開會.
我最怕那些.
我就問一個問題.
怎樣可以將人帶到耶穌面前呢.
其實我講回頭.
我受洗一年左右.
我1978年受洗.
1979年左右.
那時最少有兩個美國宣教士家庭.
後來更多.
他們都會講廣東話.
不過聊聊天就可以.
但開始崇拜要講道.
他們就找人翻譯.
我真的不明白為什麼.
他們找了我做其中一個翻譯.
這是對我很震撼的.
因為在我成長的那班人之中.
講英文比我大有人在.
很厲害的.
我就說為什麼你不找他們呢.
他們更厲害.
我直接當著他們的面.
他們更好.
他們那班人.
這些弟兄姊妹很抵打的.
行了 你行的.
我不知道他們是真的覺得我行還是不覺得我行.
總之就是你了.
所以我就開始去事奉.
就是這樣.
那時十幾歲就學習.
現在都六十歲了.
過了六十歲.
已經不是再衝刺的時候了.
我又覺得怎樣呢.
我都是在後面做一些配搭的東西.
就是搖旗吶喊.
舉牌那些人.
做這些事.

$^{801}$隨著人生的閱歷增加了.
我就大概可以看到教會有些什麼明顯的問題.
其實很多時候是人.
那個眼光不夠宏觀.
他限制了神的屬靈的工作.
我經常都是這樣.
我就經常見到很多弟兄姊妹.
他們對信仰很認真的.
我其實都認識.
我們今天崇拜主席Richter姐妹很久了.
我應該沒有記錯你的名字吧.
我第一次聽她們祈禱這麼詳細.
不用讀稿的.
我們在錦繡堂那些崇拜祈禱是讀照寫的.
照本宣科的.
我第一次聽到這麼流暢的祈禱.
這麼有內容的.
很特別.
所以我自己的看法就是.
我看到弟兄姊妹是很認真去事奉的.
你不要以為.
有些人去看任何的事奉崗位.
他都是很嚴肅的.
很認真的.
他們做主席也好.
領唱也好.
都很認真.
你去看是看得到的.
所以有些人對崇拜不是很著緊的.
這些你會分得到的.
但你會發覺原來有很多人很忠心去事奉的.
當我見到弟兄姊妹這麼熱心.
我想起我們要學安德烈.
安德烈真的不是很厲害的人.
為什麼我會這麼欣賞他.
因為我自己覺得我也是.
有種認同感.
我們要學他一樣東西.
就是將人帶到耶穌基督面前.
其實現在我接觸了一些人.

$^{841}$他有病.
他們反而覺得我剩下來的時間不多了.
我要更加把握將人帶到耶穌面前.
不知道你有沒有同感.
有些人是這樣的.
所以你可以說現在很多挑戰從四方八面而來.
牧者也是.
信徒領就是.
但如果我們都是一種很混亂的狀態又不行.
反而我們去學安德烈.
想想怎麼做.
未來還有很多挑戰.
還有移民潮.
還有很多事情.
其實我們沒有一套方法.
沒有什麼可以解決的.
你可以說見招拆招.
見步行步.
教會需要方向這句話永遠都對.
在任何一間教會都對.
只不過情況不同.
但對我來說.
我們學安德烈.
將人帶到耶穌面前.
我覺得這是我們任何一間教會.
今時今日應該學習的.
因為這才是最根本的東西.
我們有些事情是應付不了的.
解決不了.
但我們就要把握我們的青春.
用這個字眼青春.
將人帶到耶穌面前.
好像安德烈這樣.
其實安德烈是耶穌第一批呼召的門徒.
甚至第一位.
另一位就是約翰.
如果你覺得你信主很久了.
你未必像彼得那麼威風.
但不要緊.
你可以學安德烈.

$^{881}$安德烈不是門徒裡面的核心.
他是五餅二魚神蹟的促成者.
你想想.
五餅二魚的神蹟不是由彼得促成的.
是不是很公道.
上帝很公道.
不是最厲害的促成這件神蹟的.
安德烈因此將人帶到耶穌面前.
讓更大的事去成就.
在我們的想法之中.
我們能夠做的事是有限的.
我們一定是這樣.
但你將人帶到耶穌面前.
耶穌基督可以成就你想像不到的更大的事.
你用這個想法去想.
你就會發覺.
我將人帶到耶穌基督面前.
那個作用有多大.
很大.
很多時候.
譬如我的同學帶我回教會.
他沒有奉獻做傳道人.
我奉獻了做傳道人.
我不是說這件事厲害一點.
而是你可以想像.
將來你帶那個人信主.
那個人成為.
他更加親近神.
做更多的事也不一定.
是不是.
求主幫助我們學習安德烈.
你將人帶到耶穌基督面前.
一定可以成就更大的事.
如果你想是小人物為神做大事.
其中一個方法.
你就像安德烈這樣.
跟他哥哥說我遇見彌賽亞了.
因為他哥哥很厲害.
比特.
所以我想我們未必是直接做事那個.

$^{921}$但是我有份參與在神的工作那裡.
你會不會很開心呢.
好,我們一起同心祈禱.
我們的天父上帝.
因為聖經裡面很寶貴.
他不僅僅是有大道理.
他也有一些生活的智慧在裡面.
我們都是很平凡的人.
小人物.
做不到什麼.
沒什麼作用.
但是我們看到安德烈.
其實他也很坦白.
他說五餅二魚.
其實有什麼作用呢.
但是他最低限度都將這個消息.
告訴耶穌基督.
他最低限度都聽耶穌的話.
去問有沒有人有食物.
最低限度他都勇敢的將這個消息.
告訴耶穌基督.
就求你藉著安德烈這個行動.
他這種轉告的恩賜.
來去激勵我們.
讓我們把握我們在世的日子.
來去為神去做我們應該做的事.
我們真的很希望.
在一個比較混亂的狀態底下.
是你安定人心.
安定教會.
是你帶領我們走過不容易的日子.
這樣的祈禱是奉耶穌基督的名.
阿們.
\newpage



\section{羅馬書 5:1-11}
\label{sec:D7R35FEAwH0}
\textbf{2021.04.18《我們快樂是因為……》羅馬書 5\_1-11( 和修版)  曾敏姑娘}
\newline
\newline
連結: \href{https://youtube.com/watch?v=D7R35FEAwH0}{\texttt{https://youtube.com/watch?v=D7R35FEAwH0}} ~~~~ 語音日期: 2021-04-20
\newline
\newline
\hyperref[sec:q2mVvzFJiBo]{\small{< < < PREV SERMON < < <}}
~
\hyperref[sec:index_chronic]{\small{[返順時目]}}
~
\hyperref[sec:index_scriptual]{\small{[返順卷目]}}
~
\hyperref[sec:NkgRuGpNbys]{\small{> > > NEXT SERMON > > >}}
\newline
\newline
羅馬書 5:1-11
\newline
\begin{longtable}{cl}
\hline
\hline
章節 & 經文 (和合本修訂版)\\
\hline
5:1 & \begin{tabularx}{0.7\textwidth}{X} 所以，我們既因信稱義，就藉著我們的主耶穌基督得以與神和好。 \end{tabularx} \\ \\ \relax
5:2 & \begin{tabularx}{0.7\textwidth}{X} 我們又藉著他，因信得以進入現在所站立的這恩典中，並且歡歡喜喜盼望神的榮耀。 \end{tabularx} \\ \\ \relax
5:3 & \begin{tabularx}{0.7\textwidth}{X} 不但如此，就是在患難中也是歡歡喜喜的，因為知道患難生忍耐， \end{tabularx} \\ \\ \relax
5:4 & \begin{tabularx}{0.7\textwidth}{X} 忍耐生老練，老練生盼望， \end{tabularx} \\ \\ \relax
5:5 & \begin{tabularx}{0.7\textwidth}{X} 盼望不至於落空，因為神的愛，已藉著所賜給我們的聖靈，澆灌在我們心裡。 \end{tabularx} \\ \\ \relax
5:6 & \begin{tabularx}{0.7\textwidth}{X} 我們還軟弱的時候，基督就在特定的時刻為不敬虔之人死。 \end{tabularx} \\ \\ \relax
5:7 & \begin{tabularx}{0.7\textwidth}{X} 為義人死，是少有的；為仁人死，或者有敢做的。 \end{tabularx} \\ \\ \relax
5:8 & \begin{tabularx}{0.7\textwidth}{X} 惟有基督在我們還作罪人的時候為我們死，神的愛就在此向我們顯明了。 \end{tabularx} \\ \\ \relax
5:9 & \begin{tabularx}{0.7\textwidth}{X} 現在我們既靠著他的血稱義，就更要藉著他得救，免受神的憤怒。 \end{tabularx} \\ \\ \relax
5:10 & \begin{tabularx}{0.7\textwidth}{X} 因為我們作仇敵的時候，尚且藉著神兒子的死得以與神和好，既已和好，就更要因他的生得救了。 \end{tabularx} \\ \\ \relax
5:11 & \begin{tabularx}{0.7\textwidth}{X} 不但如此，我們既藉著我們的主耶穌基督得以與神和好，也就藉著他以神為樂。 \end{tabularx} \\ \\
[1ex]
\hline
\hline
\end{longtable}
$^{1}$弟兄姊妹早安.
果然站在這個角度.
看到弟兄姊妹的感覺是非常不同的.
現在很明白詠倫的感受.
在這裡看著弟兄姊妹.
還會想起認識在座當中的一些成員.
也和大家有不同的經歷.
有這麼多奇妙的經歷.
願意今天上帝的說話.
是一同的激勵我們.
我自己預備這個講章的時候.
感覺是很想和自己說話.
我和自己說完話之後.
我又想將這個和自己說的話.
和大家去分享.
我們快樂是因為什麼呢?.
什麼事情令我們感到快樂呢?.
是不是事業發展很順利.
有很多錢.
有房有車.
名成利就.
身體健康.
很少病痛.
有美滿的家庭關係.
有好的婚姻.
兒女聽話.
學業有成.
自己有美好的明星.
是不是這些呢?.
可能有弟兄姊妹馬上就會說.
怪不得我不開心.
因為這些這麼好的事情都沒有發生在我身上.
又或者有人就會說.
我曾經都擁有一些.
或者全部這些福氣.
現在又沒有了.
所以就不開心了.
相信這些事情都是好的.
我沒有說不好.
但是這些是不是代表.

$^{41}$就是我們生命最寶貴的.
最能夠令我們感到快樂的事情呢?.
希望今天的訊息.
都帶領我們再一次去思考這個問題.
我們一起去祈禱.
讓我們在疫情稍為緩慢的時候.
讓我們可以回到神的家裡敬拜你.
願意今天你對我們每一個去說話.
願聖靈在我們每一個心裡去動功.
讓我們聽得清楚明白.
我們又願意回應上帝你的聲音.
奉耶穌基督的名字祈禱.
阿們.
《流馬書》一到四章.
保羅花了很長的篇幅去談論.
什麼叫做「恩順清儀」.
原來世界上的每一個人.
我們都犯了罪得罪了神.
耶穌基督為了我們的罪死在十字架上.
第三天又從死裡復活.
成就了救恩.
當我們願意相信耶穌基督的時候.
神就稱我們為義.
神稱我們這一群罪人為義.
這個信仰的內容.
我深信弟兄姊妹聽過了.
特別在座大部分都是經過洗禮的弟兄姊妹.
對你來說.
對我來說都不陌生.
「恩順清儀」是我們整個信仰的核心內容.
我們經常掛在嘴邊.
我們也有很一部分成員.
經常都很想和別人分享福音.
這就是我們說的內容.
但是弟兄姊妹.
我們有多少時間會去思考.
「恩順清儀」這件事.
對我們生命的意義是什麼呢?.
今天這段經文看上去似乎很熟悉.
但是我卻發覺.

$^{81}$我們並不多在講壇上分享.
《羅馬書》五章一到十一節提到.
「恩順清儀」所帶給我們的影響.
當中出現了幾次.
和快樂有關的詞語.
例如第二節.
「並且歡歡喜喜盼望神的榮耀」.
第三節.
「就算在患難中也是歡歡喜喜的」.
第十一節.
「也就藉著他以神為樂」.
原來「為樂」這個詞語和「歡歡喜喜」.
在原文是同一個詞語.
都是形容歡喜快樂.
是大大的歡喜.
看來「恩順清儀」所帶給我們的.
其中一個影響.
是會令我們感到快樂的.
我們多不多從這個角度去看呢?.
為什麼我們「恩順清儀」之後.
會歡喜快樂呢?.
我們透過這段經文去找答案.
首先.
我們「恩順清儀」之後.
很重要的一點.
就是我們「同神和好了」.
並且我們在他的榮耀裡是有份的.
羅馬書五章一節這樣說.
「所以我們既恩順清儀.
就藉著我們的主耶穌基督.
得以與神和好」.
五章九到十節這樣說.
「現在我們既靠著他的血清儀.
就更要藉著他得救.
免受神的憤怒.
因為我們作仇敵的時候」.
原來在我們未信耶穌之前.
神視我們為仇敵.
其實我覺得這一點.
我們以往並不多提.

$^{121}$我們只說我們犯罪.
我們得罪了神.
我們要求神赦免等等.
你有沒有想過嚴重到.
神視我們為仇敵.
我們不是故意敵對神.
但因為我們犯罪的緣故.
我們激怒了神.
所以自己成為了神的仇敵.
我們很少思考神會憤怒這件事.
有不同的弟兄姊妹.
總是希望教會的講壇.
多些說上帝的慈愛.
上帝的愛怎樣溫暖我們的心.
擁抱我們的生命.
因為我們的生命有很多困難.
我們很需要神的安慰.
所以最好多說些神的慈愛.
但如果我們只說神的慈愛.
而不說神的聖潔公義.
這種教導其實是很不全面的.
神絕對是聖潔公義的神.
祂看到人犯罪就會感到憤怒.
神憤怒的時候是怎樣的呢?.
我覺得很難完全用言語去形容.
我想起很多年前的有一天.
當時我忘記了.
我是否在一個機構服務的時候.
有一天下班的時候.
非常厲害.
雷霆電閃.
下著傾盆大雨.
很多地方是水浸.
那天我不知道為什麼.
一定堅持不飛回家.
一定要走回家.
全身濕透了.
水浸到我小腿的位置.
很深很深.
所以我也走得不快.

$^{161}$雷聲震天.
我感覺每隔一會兒.
就好像有閃電在我很近的位置.
其實當時心裡面.
啪啪啪啪啪.
這樣跳.
很害怕很害怕.
一邊祈禱一邊回家.
心裡面有一種很大的懼怕.
我想起這一幕很深刻.
我相信上帝煩惱的時候.
比這一幕震撼得多.
真的不是我和你受得起.
如果真是的.
神將他的怒氣.
不斷傾倒在我們身上.
不是我和你受得起.
我們會害怕.
上帝視我們為敵人.
其實我們根本沒有辦法.
跟上帝作對.
敵人會打仗.
我們怎樣打得過呢?.
本來我們是死路一條.
那條死路是甚麼?.
那條死路是說通往的.
永遠死亡的那條路.
本來我們的光景是這麼淒涼的.
現在因為耶穌基督為我們成就救恩.
因為我們信耶穌基督.
我們就免受了神的憤怒.
跟神和好了.
我們跟神的關係和好重見.
對我們來說是天大的好消息.
很值得我們開心.
頂曉妹我們有沒有單單因為得到救恩.
而感到歡喜快樂呢?.
還是你一說到快樂.
一定要說到剛才.
世界上的人都很想說的事件.

$^{201}$我們的事業,學業,家庭,健康.
那些就歡喜快樂.
從來我講歡喜快樂.
我不會想到救恩這個層面.
但你有沒有想過.
其實這件事何等值得我們歡喜快樂.
甚至可以說.
差不多是最開心的一件事.
上帝不再生我們氣.
和我們和好了.
這是一件很開心的事.
當年我讀神學的時候.
曾經聽過一些同學的見證.
我不知道會不會有不同的頂曉妹.
或者朋友會覺得.
什麼人會去讀神學呢?.
一定是那些愛上帝吧.
會不會讀神學呢?.
沒有工作的.
沒有工資的.
這些人去讀神學.
要不就很有錢.
之前存了一大筆.
很多財產的.
不會休息.
於是就去讀神學.
這樣的.
要不就有個強而有力的後盾.
後台.
會撐著他們經濟上各方面的.
是不是這樣呢?.
我記得當年我讀神學的時候.
身邊有好一些同學.
在經濟上很缺乏的.
我真的會形容他們為窮的.
為什麼窮呢?.
經常都會說.
哎呀,快要交學費了.
快要交食宿費了.
當時我的那間神學院是要住宿的.

$^{241}$還會包三餐.
所以我們交又要住宿費.
又要吃的費用.
有同學就說.
快要到期交錢了.
我沒有錢了.
真的一分都沒有.
我說怎麼辦呢?.
我還沒想到.
不知道.
我還沒到最後.
不知道要怎麼處理.
就是沒得交.
一路如是者.
過到那個期限的時間.
很奇妙.
不知道為什麼.
同學就可以在信箱.
或者在哪裡.
收到一張支票.
剛剛好支付學費.
食宿費.
過完這一關之後.
過到下一次又來過.
哎呀,又來到期了.
哎呀,又沒有了.
那我又怎麼辦呢?.
我不知道.
我經過上一次.
我知道上帝會供應.
就是這樣的.
關關難過,關關過.
這是我們神學院的同學.
最喜歡說的那句.
我們有窮的人.
也有身體常常不適的.
經常生病的同學.
卻要去面對很繁重的.
神學的學習.
神學的實習等.

$^{281}$原來這一群.
又會生病.
又經濟上缺乏的人.
為什麼要去讀神學呢?.
就是因為他們自己.
經歷過福音的寶貴.
他們有一個很豐富的.
很滿足的生命.
當上帝呼召他們的時候.
他們會走去和別人分享.
這個福音.
怎樣帶給人真正的快樂.
怎樣令人和上帝復和.
是一個令人滿足快樂的福音.
所以他們義不容辭就去.
甘願為了這件事付代價.
我不知道在座的弟兄姊妹.
你是怎樣信耶穌的.
是不是也有傳道人.
神學生,宣教士.
和你傳福音的呢?.
其實這些神學生,傳道人,宣教士.
都是抱著一個心.
很想令大家和上帝復和.
有一個真正滿足快樂的人生.
今天我們活得如何?.
快樂嗎?.
有沒有活出這個福音.
帶給我們的那種快樂?.
那種開心興奮.
還是我們已經被現實生活打沉了?.
看著我們不堪的經濟,健康.
人際關係.
我們被打沉了.
即使聽過福音也不快樂.
甚至有時比未信耶穌的人更慘.
我有時很佩服.
未信耶穌的人.
他們經常說.
我們積極生活之道.

$^{321}$如何快樂起來.
他們會說很多積極之道.
我想我們信耶穌的人.
應該比他們更加憂醒.
因為我們所擁有的.
那些未信耶穌的人.
他們未有這樣的福氣擁有.
但今天我們是否真的這樣經歷呢?.
值得我們去思考.
除了與神和好之外.
我們也可以在神的榮耀上有份.
五章二節說.
「我們有藉著祂恩信.
得以進入現在所站立的這恩典中.
並且歡歡喜喜盼望神的榮耀」.
這個「神的榮耀」是指神榮耀的顯現.
在末日的時候.
神會顯出祂的榮耀.
而當神這樣做的時候.
原來上帝的百姓.
可以與神一起分享祂的榮耀.
原來恩信清二這件事.
可以令我們這些卑微的人.
有一份盼望.
在上帝的榮耀上有份.
這是多麼寶貴的事情.
是很值得我們開心的.
很久以前我曾經想過一個形容.
怎樣形容我自己的生命.
和上帝的生命是怎樣不同呢?.
我以前在想.
我去形容.
神就是在皇宮裡面.
的國王.
我所說的那種榮耀.
那種聖潔.
真是穿金戴銀.
光輝閃閃 很光亮.
皇宮裡面的生活是很富足.
很豐富.

$^{361}$很清潔 聖潔.
很漂亮.
很令人嚮往.
我呢 我的身份.
我想起我自己未信耶穌前是甚麼呢?.
好像一個乞丐.
在街上流浪.
骯髒骯髒 無父無母.
整身很骯髒.
這個國王很憐憫.
這個乞丐於是就放了進皇宮生活.
他令我洗白白 洗乾淨.
讓我可以在他身邊.
當我可以和上帝一起的時候.
這個昔日的乞丐.
竟然可以和國王有份.
其實這個差距很大.
但我深信上帝和人的差距更加大.
上帝讓我們在他的榮耀上有份.
這個真是我們不能明白.
所以不怪聖經裡面說.
要歡歡喜喜 盼望.
是否一種勉勵 一種福利.
要開心 開心地盼望.
除了歡喜盼望.
我們還會因為其他事情歡喜.
這個才是最刺激.
《流氓書》五章三到五節說.
不但如此.
就是在患難中也是歡歡喜喜的.
因為知道患難生忍耐.
忍耐生老練.
老練生盼望.
盼望不止落空.
因為他的愛已藉著所刺及我們的性靈.
囂冠在我們心裡.
我們因信清義.
同神和好之後.
生命不再一樣.
其中一件不可思議的事情.

$^{401}$就是我們竟然會因為患難而歡喜.
患難有很多種.
在這裡患難是說.
因為事奉 因為見證的緣故.
而帶給我們的苦難.
只是指這一種.
在患難中.
他在這裡說.
我們不是感到無奈 感到痛苦.
我們很消極.
我們就忍.
這個忍耐不是這個意思.
而是我們在患難裡都是歡歡喜喜的.
我們以積極的態度去面對.
因為知道患難生忍耐.
忍耐生老練.
老練生盼望.
我們絕對可以以一個很怨毒.
很憤怒的心去面對患難.
但亦可以以一個很欣然接受的心去面對.
兩者所帶來的結果是完全不一樣的.
如果我們的態度是後者.
我們可以生出忍耐的心.
即是我們以一個欣然接受的心.
以一個積極的心.
去面對我們的患難.
我們可以生出忍耐的心.
這裡的忍耐不是被動的.
是積極 是主動的.
是一種堅持 一種堅守.
當我們面對患難的時候.
能夠以一個積極的態度去面對.
就會生出一個堅忍的生命特質.
當我們繼續堅忍下去.
就會生出老練.
老練就是指經過考驗.
我們的生命質素就會合格.
是要經過考驗.
這種經得起考驗的生命.
在神裡面是有盼望的.

$^{441}$五章五字裡這樣說.
盼望是不至於落空的.
為何不至於落空呢?.
經文接著下來.
花了很大的篇幅去說.
為何不至於落空呢?.
經文裡這樣說.
我們一起去看回.
第五節.
因為神的愛已藉著所賜給我們的聖靈.
囂冠在我們心裡.
為何在神裡面的盼望不會落空?.
因為神的愛愛我們.
這份愛是怎樣來的?.
上帝賜給我們有聖靈.
而神就藉著聖靈.
將這份愛放在我們心裡.
因為這份愛.
因為神的愛.
這份盼望不會落空.
經文花了很多時間.
去描述神的愛是怎樣的.
神的愛是甚麼時候發動的?.
第六節.
我們患軟弱的時候.
基督就在特定的一時刻.
為不敬虔之人死.
神何時愛我們呢?.
是我們不好的時候.
我們軟弱的時候.
這個軟弱是說我們屬靈生命很軟弱.
我們和神隔得很遠.
我們在罪惡裡面.
這個軟弱就是下面的不敬虔.
我們不敬畏上帝.
我們沒有說哪一位.
在這個世界上是特別的敬虔.
哪一位就不敬虔.
羅馬書說世人都犯了罪.
我們每一個都犯了罪.

$^{481}$得罪神.
每一個軟弱的時候.
每一個都不敬畏上帝的時候.
在那個時候.
基督就在特定的時刻.
為我們死了.
為我們犧牲.
這份愛有甚麼特別?.
在我們還是軟弱.
我們不敬虔的時候犧牲.
有甚麼軟弱?.
你看看神的愛多好.
人的愛多有限.
第七字就說人的愛多有限.
為二人死是少有的.
為人人死或有感造的.
這個二人是說那些很遵守律法的.
行為上做得很好的.
不過他這一種是獨善其身.
我自己做得好.
不過我不一定要關心你.
我又不一定要幫助別人.
你和我沒有甚麼關係.
我做好我自己就好了.
這一種就叫二人.
這一種人因為不關心別人.
和別人沒有甚麼關係.
於是為二人死的人是少有的.
少有不是說可能有一兩個.
人數少.
少有的意思是很難找到人.
為這些人死的.
你問有沒有誰為這個.
叫做很守律法的.
行得正站得正的人犧牲呢?.
很少.
沒有.
很難找到.
沒有.
為人人死.

$^{521}$人人就是好的人.
我自己行得正站得正.
我又很關心別人.
我和別人分享.
我經常幫助別人.
我心裡牽掛著別人.
我們都會在人群當中找到這樣的人.
為這些人犧牲.
或者或者.
有感受.
你可能會找到人.
願意為他們犧牲.
可能而已.
不一定.
人的愛就這麼有限.
更何況那些又不行得正.
不站得正.
又不為人著想的.
你不用想了.
我們人都是這樣的.
會按照對方是怎樣的人.
決定幫不幫他.
愛不愛他.
要衡量一下他多值得愛.
我就愛了.
我付出多少.
我們都是這樣計算的.
這份愛隨時都可以改變的.
他以前很好的時候.
我就愛他多一點.
他現在沒那麼好了.
更加愛他少一點.
我們說到犧牲.
拿我的命出來.
我先想一想.
我最喜歡問這個問題.
跟別人傳呼音.
你覺得在這個世界上.
有誰會為你犧牲呢.
人們通常會想一想.

$^{561}$大家猜猜說得最多的答案是什麼.
爸爸媽媽.
當然了.
爸爸媽媽.
看這個是我自己的骨肉.
最大機會是這樣.
所以父母的愛很偉大.
配偶都想一想.
不一定.
很刺激地問他.
你會為誰犧牲呢.
先想一想.
為什麼要為別人犧牲呢.
人是這樣的.
我們的愛很有限.
但這裡他說.
唯有基督在我們還作罪人的時候.
為我們死.
我們遠約的時候.
我們不敬虔的時候.
我們犯罪得罪上帝的時候.
他就為我們犧牲.
如果回應這段經文.
一開始的時候.
我們做他仇敵的時候.
他就為我們犧牲.
今天我們再看聖經教導.
什麼叫為仇敵呢.
寫命呢.
這句話的感覺應該很不同.
我們很壞的時候.
我們犧牲.
我們很壞的時候.
我們得罪他的時候.
做頂心杉的時候.
他為我們死.
神的愛.
就是在這裡.
向我們顯明.
是因為這份愛.

$^{601}$我們有盼望.
我們生命不再一樣.
這份愛.
是何等寶貴呢.
今天呢.
剛才我們是領唱帶.
我們唱一首詩歌.
叫《奇妙的愛》.
你問我呢.
神的愛.
是多麼高深.
其實是很難回答的.
神的愛.
那種心.
其實他這些是很難做到.
對人來說.
很難做到.
上帝做到愛我和你.
因為這樣.
我們才有這麼多福氣.
第九節到十一節.
就做個總結.
做個總結.
他一路做一些比較.
既然是這樣就更加要.
如果我們是這樣的時候.
尚且是這樣.
就更加要怎樣.
不但如此.
還要怎樣.
透過連續幾句說話去做一個總結.
他說我們既然靠著.
耶穌基督的血.
因為耶穌為我和你.
犧牲.
我們信他.
神稱我們為義.
既然已經這樣.
這是很寶貴的事.
我們更加要藉著他得夢拯救.

$^{641}$免受神的憤怒.
讓神不要再生我們氣.
是不是.
既然是這樣就要做到這樣.
因為我們做仇敵的時候.
尚且藉著神兒子的死.
得以與神和好.
我們得夢拯救.
我們與神和好.
更加要因他的生得救.
這個「生」是說.
他復活的生命.
我們更加要因為.
耶穌基督復活的生命.
去得夢拯救.
這個拯救是說.
將來有一天.
我和你面見耶穌的時候.
那個徹底的.
一分百的拯救.
耶穌基督為我們成就.
奇妙的救恩.
不但如此.
我們既然藉著耶穌基督得以與神和好.
更加.
去到總結.
以神為樂.
以神為.
我和你的快樂.
這整段經文最想說.
以神為我和你的快樂.
就是因為神愛我們.
在我們很軟弱的時候.
我們不配得愛的時候.
我們本來要配得滅亡的時候.
他愛我們.
派他的兒子耶穌基督救我們的生命.
是因為信得清為宜.
同神和好.
我們的人生不再一樣.

$^{681}$在這裡.
我仍然想問那句.
剛才我開始問.
弟兄姊妹你快樂嗎?.
我問我自己.
是什麼拿走了我們的快樂.
我們的眼是看著什麼.
我們人一樣看著.
日常生活很想要的東西.
世人很想要的東西.
看著那些有就快樂.
沒有就不快樂.
但我和你信了耶穌.
我和你蒙上帝的愛.
我們的生命.
有沒有因為他感到快樂.
可以超越.
這個世界帶給我們的經歷.
每天上班帶著什麼心情.
我和你.
捱世界的心情.
磨爛的心情.
面對困難的時候.
是什麼心情.
是打敗仗的.
被徹底打垮的心情.
還是我得勝了.
我贏了.
其實我和你已經贏了.
因為神的愛.
充滿我們的生命.
上帝不再對我們憤怒.
從今以後還有什麼害怕.
是不是.
神的愛不只是那一下救贖我們的生命.
人這種和好的關係.
是由現在直到永遠.
神的愛都不離開我們.
試問有什麼.
是不可以面對的.

$^{721}$在大小事情上.
上帝都與我和你同在.
我記得很多年前.
去短孫.
那次也是我其中一次.
很深刻經歷的一次.
我去了緬甸短孫.
我在想什麼令我快樂.
那次其實心裡很大掙扎.
我就以這次短孫.
看我自己的人生.
這次短孫之後.
我要回答自己的問題.
究竟我要繼續服侍上帝.
還是不如索性.
不再走服侍的路.
出來找工作.
好好地賺錢養家.
雖然我還未結婚.
但也要養父母.
當時真的想以這個短孫.
去認識自己.
究竟是一回事.
那次去緬甸的第一天.
就很刺激.
第一天我遇到交通意外.
我們的車.
撞到一個當地人.
那位當地的女孩子.
受了重傷很危險.
司機很有良心.
很盡責.
他是神學院的校長.
他把短孫隊都折斷.
一群人站在路邊.
然後把傷者送到醫院.
我們也很支持校長這樣做.
那次很震撼.
因為發生交通意外.
車上的玻璃全碎了.

$^{761}$我們可以很重傷.
不單是車外的女孩子.
但看到上帝保守我們.
所以車內的人.
大部分都很平安.
我記得我們站在路上.
隊員們已經.
幫對方撥去.
我們身上的玻璃.
我就已經開始拿電筒.
因為當時天黑.
就照著一些貨本.
準備備貨.
預備未來幾天的事奉.
未來幾天.
我們都有很多的經歷.
這個事件要繼續處理.
政府應該說警察方面.
很想趁這個機會.
又賺一筆錢.
希望神學院.
要給錢警察那邊.
息事寧人.
這件事.
糾纏來糾纏去.
我們也很擔心被撞傷的女孩子.
但神很奇妙.
這件事.
逐一幫我們處理.
那個女孩子.
很奇妙地.
完全康復.
最後應該是.
警察也不能成功.
從神學院身上撥了一筆錢.
而短宣隊要做的事.
完全順利地完成.
那一次我就看到.
其實事奉上有很多困難.
困難.

$^{801}$作為宣教士.
作為傳道人.
很多困難.
不是那麼簡單.
剛才經文說的患難.
是否因為見證.
因為事奉帶來的困難.
是很實在.
但我卻感到.
莫名的快樂.
為甚麼呢.
因為原來這個傳呼音的責任.
傳呼音的使命.
是一件令人感到快樂的事.
無論你在本地傳呼音.
到外面宣教.
都是.
能夠有這個呼召.
被神所用.
是一個榮幸.
我從這件事.
看到上帝愛人.
很愛世界上.
每一個人.
愛我和你.
上帝很想拯救世界上.
我們每一個罪人.
他們不只是想拯救.
他們想我們.
從罪裡得釋放.
有生命真正的快樂.
我很盼望.
經過今天的講道之後.
大家回想一下.
甚麼是令我們快樂.
上帝的救贖恩典.
有沒有成為我和你的動力.
面對生活的挑戰呢.
頂尖會起來.
不要被生活打垮.

$^{841}$我特別要對.
紅衣人的姐妹們說.
不要被這裡種種的.
生活的困難.
打垮了自己.
耶穌基督給我們的.
是寶貴無比的.
是已經勝利了.
我們一同禱告.
天父上帝.
今天讓我們去看這段經文的時候.
看到是何等的寶貴.
原來我們.
曾經被你視為.
敵人.
本來我們.
是往死裡走.
但是你卻.
藉著你的愛子耶穌基督.
去救贖我們的生命.
讓我們和你和好.
你讓我們的人生.
是一個得勝的人生.
不會被打垮.
今天無論我們.
在生命裡有沒有擁有.
世人覺得很寶貴的人.
物,事,情也好.
上帝我們已經贏了.
我們不是自欺欺人.
上帝.
你給我們一個徹底的.
勝利的.
一個.
這樣的成就.
天父求你激勵我們.
起來,離開我們的憂傷.
離開我們的沮喪.
離開我們那種.
感到失敗的.

$^{881}$是沒有辦法的.
這樣的處境.
讓我們天天抓住你.
天天倚靠你.
天天以你為樂.
讓我們這裡.
是一個彰顯耶穌基督榮耀的群體.
我們將我們自己的生命.
交在你手中.
奉耶穌基督的名字祈求.
阿們.
\newpage



\section{路加福音 5:17-20-24-26}
\label{sec:NkgRuGpNbys}
\textbf{2021.04.25《信心的服侍》路加福音 5\_17-20,24-26( 和修版)  楊穎兒導師}
\newline
\newline
連結: \href{https://youtube.com/watch?v=NkgRuGpNbys}{\texttt{https://youtube.com/watch?v=NkgRuGpNbys}} ~~~~ 語音日期: 2021-04-26
\newline
\newline
\hyperref[sec:D7R35FEAwH0]{\small{< < < PREV SERMON < < <}}
~
\hyperref[sec:index_chronic]{\small{[返順時目]}}
~
\hyperref[sec:index_scriptual]{\small{[返順卷目]}}
~
\hyperref[sec:FhlMWYYtvXU]{\small{> > > NEXT SERMON > > >}}
\newline
\newline
路加福音 5:17-20-24-26
\newline
\begin{longtable}{cl}
\hline
\hline
章節 & 經文 (和合本修訂版)\\
\hline
5:17 & \begin{tabularx}{0.7\textwidth}{X} 有一天，耶穌教導人，有法利賽人和律法教師在旁邊坐著；他們是從加利利各鄉村、猶太和耶路撒冷來的。主的能力與耶穌同在，使他能治好病人。 \end{tabularx} \\ \\ \relax
5:18 & \begin{tabularx}{0.7\textwidth}{X} 這時，有些人用褥子抬著一個癱子，要把他抬進去放在耶穌面前， \end{tabularx} \\ \\ \relax
5:19 & \begin{tabularx}{0.7\textwidth}{X} 卻因人多，找不出法子抬進去，就上了房頂，從瓦間把他連褥子縋到當中，在耶穌面前。 \end{tabularx} \\ \\ \relax
5:20 & \begin{tabularx}{0.7\textwidth}{X} 耶穌見他們的信心，就說：「朋友，你的罪赦了。」 \end{tabularx} \\ \\ \relax
5:21 & \begin{tabularx}{0.7\textwidth}{X} 文士和法利賽人就開始議論說：「這個人是誰，竟說褻瀆的話？除了神一位之外，誰能赦罪呢？」 \end{tabularx} \\ \\ \relax
5:22 & \begin{tabularx}{0.7\textwidth}{X} 耶穌知道他們所議論的，就回答他們說：「你們心裡為甚麼議論呢？ \end{tabularx} \\ \\ \relax
5:23 & \begin{tabularx}{0.7\textwidth}{X} 說『你的罪赦了』，或說『你起來行走』，哪一樣容易呢？ \end{tabularx} \\ \\ \relax
5:24 & \begin{tabularx}{0.7\textwidth}{X} 但要讓你們知道，人子在地上有赦罪的權柄。」他就對癱子說：「我吩咐你，起來！拿你的褥子回家去吧。」 \end{tabularx} \\ \\ \relax
5:25 & \begin{tabularx}{0.7\textwidth}{X} 那人當著眾人面前立刻起來，拿了他所躺臥的褥子回家去，歸榮耀給神。 \end{tabularx} \\ \\ \relax
5:26 & \begin{tabularx}{0.7\textwidth}{X} 眾人都驚奇，也歸榮耀給神，並且滿心懼怕，說：「我們今日看見不尋常的事了！」 \end{tabularx} \\ \\
[1ex]
\hline
\hline
\end{longtable}
$^{1}$弟兄姊妹平安.
多謝魯牧師的安排.
讓我可以在實習期間.
跟大家分享聖經的訊息.
我相信「坦子」這個得醫治的故事.
大家都很熟悉了.
在禱告之後.
我決定講道的題目是「信心的服事」.
經文是來自.
《路加福音》第五章.
今天的講道經文.
會分三個部分.
第一個部分是.
回顧福音書關於「坦子」的故事.
第二個部分.
我們分析人物的說話.
和行動.
看看今天的經文.
裡面的獨特之處.
第三部分.
我們一起思想一下.
什麼是「信心的服事」.
在四福音裡.
關於「坦子」得到醫治.
這個神蹟的故事.
除了在《約翰福音》.
沒有提及之外.
其餘三卷福音書都有記載.
我們一起看看.
按照《馬太福音》的記載.
比較簡單.
只有一句.
有人把一個躺在床上的「坦子」.
帶到他那裡.
然後耶穌就看見他們的信心.
就是這麼多.
有多少人去抬「坦子」.
抬「坦子」的人.
是否有信心.
抬「坦子」的人遇到什麼困難.

$^{41}$他們用什麼方法去解決.
全部都沒有提及.
這是《馬太福音》的記載.
而去到《馬可福音》裡.
我們有多一些的內容.
經文提到有幾個重點.
第一.
就是作者很客觀地描述眼前的事物.
經文這樣說.
那時有一個人.
那時有人把一個「坦子」.
帶到耶穌那裡.
然後耶穌就看見.
有人把一個「坦子」.
帶到耶穌那裡.
是由四個人抬來的.
現在我們知道.
原來是有四個人去抬「坦子」.
第二點.
就是當時要面對的困難是什麼.
原來是人多.
他們沒有辦法將「坦子」.
抬到耶穌面前.
而第三點.
就是按當時的情況.
那四個人作出了一些策略.
他們就是將「坦子」.
調到耶穌所在的地方.
拆了房頂.
然後就將「坦子」.
連「玉子」都追了下去.
即是調了下去.
這就是《馬可福音》的記載.
到今天.
我們在《路加福音》的經文.
剛才.
《馬可福音》提到的幾個描述.
就是有人去抬「坦子」.
因為人多.
不能進去.

$^{81}$所以就提到這個解決的方法.
但是值得留意的.
就是作者在《路加福音》.
提到一個很獨特的地方.
就是第十七節.
「主的能力與他同在.
叫他能治病」.
這個「他」是指主耶穌.
主耶穌有這麼多的能力.
偏偏就是提到.
祂有醫病的能力.
其實作者在這段經文之前.
正正就是在說.
耶穌醫好一個麻瘋病人.
為甚麼不在醫治麻瘋病人的時候.
去說祂有治病的能力.
而要到現在「坦子」的故事.
才去提祂有治病的能力呢?.
我們看完下一句就知道了.
經文提到第十八節.
「要抬進去放在耶穌面前」.
新譯本的領導人.
新譯本的另一個解釋.
就是更加強調「要」的意思.
其實就是在說.
他們很想.
放「坦子」.
去到耶穌面前.
我們想想,這不是客觀的描述.
這是作者很主觀地.
告訴我們.
那幾個抬「坦子」的人的想法.
他們想抬「坦子」進屋.
放在耶穌面前.
這幾個人有這樣的想法.
他們就按照這個想法去做.
這個想法是甚麼呢?.
就是想把「坦子」送進去.
放在耶穌面前.
為甚麼呢?.

$^{121}$因為主的能力.
與他同在,叫他能治病.
現在我們明白了.
原來這一句.
是要為夏文埋下一個伏線.
然後到了二十節.
耶穌看見他們的信心.
就跟「坦子」說.
「你的罪赦了」.
我們就要問.
抬「坦子」的人.
他們有甚麼生命的質素.
以至耶穌看見他們的信心呢?.
我們看回《路加福音》.
經文的編排.
第一章至第四章.
是說耶穌的出身.
祂是「神兒子」的身份.
「你是神兒子」.
祂是神兒子的身份.
來到第五章的時候.
連續三個故事.
都是關於靈魂的需要.
我們一起看看.
第五章第一個故事.
是耶穌看見門徒心靈的需要.
但是門徒.
他們自己看不到自己的需要.
於是.
耶穌就呼召門徒.
「你們不要再去打魚了」.
「你要成為我的門徒,跟隨我」.
「要做一個得人的漁夫」.
接著到第五章.
第二個故事.
是說一個被世人排斥的麻瘋病人.
他知道自己靈魂有需要.
肉身有缺陷.
於是去尋求耶穌的醫治.
接著來到第五章的第三個故事.

$^{161}$亦都是我們今天看的經文.
「癱子」的故事.
我們就會發現.
不是耶穌去看見癱子的情況.
而是.
祂去看見.
癱子的情況.
不是癱子看到自己.
肉身靈魂需要的情況.
而是那幾個.
抬癱子的人.
他們看到癱子的苦況.
就想將他帶到.
耶穌面前.
我想大家看到漢見的層次.
第一個故事是門徒自己看不到.
第二個故事.
是患病的人自己看見.
第三個故事.
就是有人.
看到其他人的需要.
所以我們今天藉著「癱子」的故事.
我們都要去學習.
看到別人的需要.
特別是靈魂的需要.
影響靈魂的因素.
有很多可能是來自.
身體的需要.
情緒上.
生活上的需要等等.
很多需要都會影響到.
靈魂如何生長.
我們看看「癱子」的故事.
抬癱子的人.
所看到的是.
癱子肉身的缺乏.
因為他沒有了神創造的形象.
他沒有了人本身可以活動.
自由自如的形象.
於是他就帶.

$^{201}$「癱子」去到耶穌面前.
作者的寫作.
導向是很清晰的.
經文說.
主的能力與耶穌同在.
叫耶穌能治病.
經文是強調耶穌能夠治病.
這幾個人就以行動.
去表示他們相信的事.
相信神.
是與耶穌同在.
相信耶穌.
是擁有治病的能力.
然後耶穌看到他的信心.
那個朋友的罪就卸了.
所以這段經文是引導我們去看.
抬癱子的人.
是要他肉身得到醫治.
不只是要癱子.
去到耶穌面前.
去得到那種靈魂的得救.
可能抬癱子的人.
都想癱子靈魂得救.
但是我們留意這段經文.
是提醒我們.
靈魂的需要有很多.
不只是靈魂得救.
是不同的.
拯救靈魂.
和使靈魂得救.
是三一上帝的工作.
故事說到.
當人帶癱子來到耶穌面前.
是耶穌很清楚.
很清楚地看到.
癱子靈魂深處那種渴望.
因為當時.
社會往往都會將.
肉身那種缺憾.
和罪連上關係.

$^{241}$所以其實癱子.
他肉身的苦況.
牽連到他有很大的罪疚感.
這一種罪.
是更加讓他得不到.
靈魂的自由.
比他動不了更加痛苦.
於是耶穌.
先赦免他的罪.
讓他從裡面釋放那種.
被罪捆綁的痛苦.
使他靈魂得救.
很多時候.
我們一說到靈魂.
我們就很火熱.
我們很想傳福音.
因為我們看到還有很多人未得救.
所以我們無論去到哪裡.
我們都要傳福音.
叫更多人認識耶穌.
相信耶穌.
沒錯,傳福音是我們要做的.
但是今天坦展這個故事.
是提醒我們.
去學習去顧念別人.
一些很現實的需要.
就好像如果有個人.
去問你拿一杯水.
你連一杯水都不給.
只叫他信耶穌的話.
難免就會令人認為.
我們的信仰是.
如此冷漠無情.
所以今天我們要.
學習看到別人.
很實際,生活上.
靈魂上很實際的需要.
我們不要只顧自己的事.
也要顧其他人的事.
來到第二點.

$^{281}$學習相信主耶穌的能力.
我們首先要知道.
能力是從何而來.
這幾個台坦子的人.
就是看到神與耶穌同在.
耶穌擁有醫治的能力.
同時.
耶穌所說的話.
也是帶著神能力.
是神能力的出口.
我們看第24節.
耶穌說:我吩咐你.
起來,拿你的玉子.
回家去吧.
在這裡.
耶穌行了一個神蹟.
就是叫坦子.
起來行走.
這件事很有趣.
為什麼.
耶穌要坦子拿起.
他的床呢?.
為什麼要拿起床呢?.
他還需要床嗎?.
為什麼不直接叫他.
不用拿床了,你回家吧.
為什麼耶穌不敢說呢?.
耶穌正正就是.
要他去承認.
一些痛苦的經驗.
他要坦子去學習.
背起自己曾經不能夠.
活動自如的那種經驗.
因為這個經驗.
是很獨特.
是屬於坦子他自己一個.
獨有的經驗.
當耶穌去醫治一個人.
的身心靈的時候.
從來都不是要將他過去的人生.

$^{321}$經驗抹去.
反而.
耶穌是要將他的經驗.
成為生命的見證.
當其他人.
看到坦子被醫治.
就驚訝了.
他們就會想.
這個就是以前不能活動的坦子.
現在已經治好了.
同時.
我們也會發現.
耶穌從來都沒有叫他.
坦子起來,你相信我吧.
我滿有能力,你相信我.
就會康復的.
耶穌不是這樣的.
他不是將榮耀歸給自己.
而是將.
讚頌神的權利.
讚頌神的任務.
賦予給一個剛被.
醫治的坦子.
吩咐他起來,拿起你的床.
回家吧.
第25,26節.
說到.
那人立即當眾起來.
拿著他躺過的床.
還稱讚神回家去了.
眾人都.
驚奇.
讚頌神,並且十分.
懼怕,說.
我們今天看見了.
不平常的事.
神蹟就是.
這樣發生的.
透過某些人將坦子.
帶到耶穌面前.

$^{361}$耶穌醫治了他.
靈魂最深處的那種需要.
使得生命.
得到更新.
在最軟弱的生命裡面.
耶穌賦予他.
讚頌神的權利.
使用他的生命.
使他能夠開口.
讚頌神的榮美.
這個舉動.
其實就是顛覆.
在場所有人的一眾.
的想法.
就是他坦的.
他怎會醫得好.
但偏偏耶穌就是.
將這個神蹟發生在他身上.
今天我們也要去學習.
以信心去想.
然後我們就容讓.
上帝去成就.
我們看見信心的服事.
可以牽連得很廣.
經文說到.
抬坦子的人.
他們想抬坦子進去.
放到耶穌面前.
就是這樣.
憑著信心就帶了坦子.
去見耶穌.
在這個信心的服事裡面.
會遇到困難.
但經文說.
抬坦子的人.
他們是很積極地去面對困難.
亦留意他們其實是有備而來.
經文沒有說.
但你們可以想想.
怎樣可以撬起屋頂.

$^{401}$怎樣吊坦子下去.
這一切都是提醒我們.
其實我們要有準備.
起碼要有些工具.
要有些成熟.
所以對於我們今天來說.
當我們想帶人.
到耶穌面前的時候.
我們亦都要好好裝備自己.
當教會說要做關懷社區的工作的時候.
我們都需要預備有策劃.
在與人分享福音的層面上.
起碼我們自己都要對聖經有多些熟悉.
平日都儲一些信靠神的例子.
經歷可以與人分享.
所以求主教導我們.
開我們的眼睛.
當我們遇到一些.
未認識主的朋友的時候.
我們不要只想著他靈魂未得救.
教我們要貼備一些.
我們去看到他實際生活上.
情緒上的需要.
然後我們好好裝備自己去接觸他們.
這一切都是建基於我們今天.
有沒有相信主耶穌的能力.
我們有沒有相信神是與我們同在.
今天就祈求神按著賜給各人的信心.
讓我們更加去捉住祂的能力.
好像那幾位憑信心帶毯子去見耶穌的人一樣.
勇敢去想.
我們去想.
想了我們去做.
然後耶穌就在當中成就了整件醫治毯子的事情.
同時神的命也都得到稱頌.
另外我們要留意那幾位帶毯子去見耶穌的人.
當然他們很希望耶穌能夠醫治毯子.
但是我們看回經文.
經文提到他們帶毯子來到耶穌面前.
但是沒有提及他們的意願.

$^{441}$意思就是他們有這樣的想法.
但是經文作者選擇沒有將這個要求寫出來.
經文只是提到要抬進去.
放在耶穌面前.
為什麼呢?.
我想到有三點.
第一.
耶穌不是我們的神燈.
我們怎麼可以去指示耶穌要做什麼呢?.
第二.
神的拯救工作是有祂的心意.
有祂的主權.
有祂的時間表.
尤其是由我們去吩咐神要做什麼呢?.
第三.
正正就是我們不知道神會做什麼.
我們不能夠測透祂的心意.
我們不能夠超越祂的主權.
我們不能夠控制祂的時間表之下.
但是我們仍然願意這樣走.
單單是因為我們知道神沒有能力.
沒有事情是做不到的.
於是我們就願意來到神的面前.
我們去奉獻自己的時間去事奉.
我們去奉獻金錢去推動教會的事工.
讓神可以去做神自己心意的工作.
這就是信心的服侍.
今天我的題目是「信心的服侍」.
所揀選的經文都是想將這個焦點調教在信心這個功課上.
「坦子之所以得到醫治,不是坦子他自己有信心」.
而是那幾位抬坦子的人相信耶穌的能力去行動.
耶穌見到這幾位抬坦子的人的信心.
坦子的靈魂就得到赦罪的恩典.
經歷了肉身的拯救.
今天我想和大家分享一些早期宣教的故事.
今天你和我之所以能夠認識上帝.
能夠我們今早坐在這裡敬拜.
其實都是因為昔日有西方的宣教士.
來到中國大陸這片土地上去傳教.
其實早在唐宋元明清的時代.

$^{481}$已經有宣教士來到中國大陸去傳教.
但那時候的發展很慢.
亦都受到當時執政皇帝支不支持的因素.
所以整體來說都是很慢的.
大概到了二百年前.
馬禮信傳教士來到中國大陸這片土地上.
他見到華人那種靈魂的需要.
他甘心放下世上的享受.
千里迢迢地來到中國.
首先去學中文.
去做一些翻譯的工作.
而到了後來就有更多的宣教士.
來到中國這片土地.
特別是香港去建立學校.
讓平民甚至是婦女有機會去識字.
以致他們能夠去讀明聖經.
我們想想如果沒有馬禮信.
還有很多無名的宣教士.
如果不是他們憑著信心.
來到中國大陸這片土地去宣教.
我想信仰未必可以在香港這片土地上.
播下這個福音的種子.
正正是因為他們憑著信心.
領受了神的差遣.
望見那個福音的和祥.
靈魂的需要.
今日坐在這裡的你和我.
就有能夠進入神的救贖歷史裡面.
活在神的愛和恩典當中.
我自己都有一個去非洲宣教的經驗.
想和大家分享.
在我未入神學院之前.
我跟著一個警察去非洲宣教.
那時是伊波拉病毒最嚴峻的時間.
染病後的死亡率是二分一.
那時我們都不知道為甚麼.
我們整個短宣的團隊.
都有一種不怕死的信心.
我們深信如果主帶我去.
是生是死是福是禍.

$^{521}$我們都交在主的手裡.
那時在非洲有很多人自稱牧師.
但他們沒有受過神學裝備.
他們甚至連一本聖經都沒有.
為甚麼他們會是牧師呢.
原來他們有一個傳統.
就是如果他們向人傳福音.
領人信主.
他們就去了牧師的級別.
所以一直以來.
警察會安排很多士工.
去非洲教學培訓裝備信徒.
這個士工在非洲很重要.
目的就是讓他們能夠.
在領人歸主的士工裡.
傳達一種更正確的信息.
那次的短宣.
其實我負責翻譯好幾堂的講道.
同行的香港牧者.
就會用中文去講道.
然後在旁邊用英文來翻譯.
因為他們都聽得懂英文.
但到了一些地區.
他們真的聽不懂英文.
就好像圖片中.
還有一個非洲人.
他會將英文翻譯成當地的語言.
我記得有一次.
我們去到一個很偏遠的村莊.
那裡的人住在茅屋.
茅屋沒有窗沒有燈.
家裡沒有洗手間.
要去附近的公廁.
公廁又是一個茅屋.
茅屋裡面有一個洞.
就是這麼簡陋.
我們就在村莊進行會堂教學.
這些都是會堂.
但這個會堂很臭.
因為整個天花板.

$^{561}$你見到有些洞.
其實日間有很多蝙蝠在.
所以有時有些東西滴下來.
所以環境是非常差.
但我記得.
來到的人令我覺得很驚訝.
因為他們來學神的話語.
他們穿到盛裝出席.
衣服你看多鮮.
白色衣服是雪白的.
燙過的.
大家看看我身旁的位置.
有煲胎的.
沒有拍到他們的腳.
但其實他們可能是穿拖鞋.
而且他們需要走兩小時.
來到會堂學習聖經.
那次之後.
我看到非洲的靈魂.
那種需要.
在他們靈魂的深處裡面.
他們是多麼渴慕.
去聽神的話語.
他們在那段時間.
他們要很努力去上課.
去考試.
然後他們的獎品就是那本聖經.
所以自己對非洲都有些負擔.
希望將來可以再去非洲服侍.
最後我想帶大家去回顧一下.
洪恩家在過往一年.
在聖靈的帶領之下.
教會的領袖怎樣同感一靈.
定義要在洪水橋這個社區.
要透過關懷社區.
參與信心的服侍.
過了一年.
洪恩堂做了甚麼呢?.
教會推動了一個「哥利樓」侍工.
是一個食物支援的侍工.

$^{601}$由善美基金去贊助.
受惠的人士可以一年八個星期.
去領取食物援助.
至今已經派發了超過二千包米.
由今年年頭到現在.
我們亦都派發了超過一千四百份的福袋.
和一些食物包.
教會亦都會和紅孩兒合作一個機構.
會透過利用奶粉.
對低收入家庭做出一個支援的計劃.
接下來下半年我們亦都會派發.
近四百罐的奶粉給這些低收入家庭.
剛才提到的物資.
九成都是來自機構.
或者弟兄姊妹的贊助.
另外洪恩家在這一年來.
亦都提供了一些免費的義展.
疫情較早的時候.
都有些烹飪的活動.
和一些社福的機構去合作.
將關懷社區的侍工.
推到不同年齡層面的人士.
接下來都會有兩個機構.
分別是仁愛堂和天朗膳食坊.
因為我們一直有做這樣的服務.
他們看到我們的服務.
所以在接下來的日子.
亦都願意將更多的物資分派給我們.
透過我們教會去分派給更多有需要的人.
當教會的領袖有一個這樣的想法.
就將這顆心帶到上帝面前.
求上帝去閱辣帶領,供應.
這就是洪恩家在過往一年信心的服侍.
成為了這個社區的流通管子.
將福音傳遞出去.
我們單單看這些數字很興奮.
沒錯,扶貧這個侍工我們是要做的.
但是如果我們用扶貧侍工來視為一種手段.
或者成為一種讓人可以誇口的侍工.
反而會帶來一個很摧毀的後果.

$^{641}$甚至乎神是不閱辣的.
扶貧這個侍工我們是要做的.
但是我們要知道.
扶貧其實是無底深潭.
要做的話,沒完沒了.
神不是要我們這個信仰群體.
成為世俗裡面一般的慈善機構.
而且我們看回聖經.
其實他從來沒有打算去解決貧窮這個問題.
反而耶穌更加說.
貧窮的人就有福了.
所以作為一個基督的信仰群體.
第一,我們要深信這一切的扶貧侍工.
是神讓我們去認識到.
神就是那位供應我們所需的神.
這個供應的能力是來自神.
第二,我們更加去相信.
教會的扶貧侍工只是一個開始.
讓我們可以透過分派物資.
洗腦教會的童工,弟兄姊妹.
能夠去接觸,去明白.
這個社區的群體實際的需要.
第三,更加重要的就是.
神興起我們當中的弟兄姊妹.
他們要成為這些服侍的領袖.
就好像當時耶穌到城肉身那樣.
走進社區,在這個群體裡面去服侍.
而在當中,我們就去見證.
聖靈怎樣在他們的生命裡面.
去驅走黑暗.
最終我們都是要等待.
等待神親自去擴建我們洪恩家.
讓更多的人加入我們.
在生命中經驗到那種赦罪的恩典.
我們一起建立了一個.
滿有神形象,樣式的生命.
就好像今天的經文的譚子一樣.
去見證了自己怎樣被耶穌赦罪.
透過耶穌醫治的能力.
讓他可以起來.

$^{681}$耶穌亦都是賦予他.
讚頌神大能的權力.
讓他能夠走回人群裡面做見證.
叫人們見到這個神蹟的時候.
去驚訝神的作為.
總結今天的講道.
這個譚子的故事.
是要我們去明白.
在服侍.
不好意思.
這個譚子的故事是要我們去明白.
在服侍的裡面.
我們看到很多破碎的生命.
我們見到靈魂的需要.
我們很想將榮耀的豐盛帶給他們.
很多時候在禱告裡面.
我們都聽到他們的故事.
我們心裡很火熱.
我們很想為他們禱告.
這些禱告不是為了求上帝.
立即供應他們所需.
甚至使他們的生活有180度的轉變.
但是這些禱告是要我們謙卑下來.
教導我們去做那幾個譚子.
台譚子的人.
就好像他們.
其實在經文裡面.
他們連名字都沒有.
他們單單就是有信心.
將這個破碎的生命帶到耶穌的面前.
然後懷著信心去領受神的恩典.
而神要在他們生命裡面.
做一些甚麼的工作.
這個是神的工作.
神要親自去成就.
而且當神去成就的時候.
是遠遠超過我們所想所求.
今天我們看到上帝.
怎樣帶領紅恩家.
用信心去服事.

$^{721}$深信這個信心的服事.
是需要整個紅恩家一同去參與.
以致我們各人能夠親身去經歷到.
那種信心服事所帶出來的果後.
在這個時候.
我很想邀請大家.
按著你今天所聽到的訊息.
去回應神.
以下有一分鐘裡面.
我很想邀請在座的各位.
去向神禱告.
我們可以學習這個信心的功課.
我們怎樣用信心去回應.
無論我們是有感動去參與事工.
或者我們有感動去將這些事工.
放在禱告裡面紀念.
又或者願意以金錢獻上.
都想邀請你在這個禱告裡面.
讓神看到各人的信心.
如果當中有未明白的弟兄姊妹.
亦都邀請你祈求神.
繼續帶領教會的事工.
預備紅恩家所需.
賜這個信心給各人.
以下有一分鐘的禱告時間.
之後我就會帶領大家.
一起去禱告.
(錄影開始).
(錄影結束).
(錄影開始).
(錄影結束).
我們一起禱告.
親愛的主.
我們滿心的去讚美你.
因為你就是那位滿有能力的主.
是你愛我們.
揀選我們.
感動我們.
帶領我們來到紅恩家.
你一直都和紅恩家同在.

$^{761}$從來都沒有離開過.
我們不能夠測透你的心意.
但我們願意用信心去擺上.
主若然你要興起紅恩家.
去服侍這個社區.
在這個社區傳揚你的福音.
求你讓我們每一個都能夠參與.
見證你的大能.
願你垂聽.
看見並且閱立各人的心.
按你所賜給各人的信心.
喚醒我們各人去看見別人的需要.
感動我們用不同的方式去繼續去回應.
當你彰顯你的能力的時候.
願人都將這個榮耀歸給你.
禱告共主命求.
阿們.
\newpage



\section{約翰福音 5:1-15}
\label{sec:FhlMWYYtvXU}
\textbf{2021.05.02《你要痊癒嗎?》約翰福音 5\_1-15( 和修版)  曾敏姑娘}
\newline
\newline
連結: \href{https://youtube.com/watch?v=FhlMWYYtvXU}{\texttt{https://youtube.com/watch?v=FhlMWYYtvXU}} ~~~~ 語音日期: 2021-05-03
\newline
\newline
\hyperref[sec:NkgRuGpNbys]{\small{< < < PREV SERMON < < <}}
~
\hyperref[sec:index_chronic]{\small{[返順時目]}}
~
\hyperref[sec:index_scriptual]{\small{[返順卷目]}}
~
\hyperref[sec:k0ge2og3_Pw]{\small{> > > NEXT SERMON > > >}}
\newline
\newline
約翰福音 5:1-15
\newline
\begin{longtable}{cl}
\hline
\hline
章節 & 經文 (和合本修訂版)\\
\hline
5:1 & \begin{tabularx}{0.7\textwidth}{X} 這些事以後，到了猶太人的一個節期，耶穌上耶路撒冷去。 \end{tabularx} \\ \\ \relax
5:2 & \begin{tabularx}{0.7\textwidth}{X} 在耶路撒冷，靠近羊門有一個池子，希伯來話叫畢士大，旁邊有五個柱廊； \end{tabularx} \\ \\ \relax
5:3 & \begin{tabularx}{0.7\textwidth}{X} 裡面躺著許多病人，有失明的、瘸腿的、癱瘓的。等候水動， \end{tabularx} \\ \\ \relax
5:4 & \begin{tabularx}{0.7\textwidth}{X} 因為有天使按時下池子攪動那水，水動之後，先下水的人無論患甚麼病都能得痊癒。 \end{tabularx} \\ \\ \relax
5:5 & \begin{tabularx}{0.7\textwidth}{X} 在那裡有一個人，病了三十八年。 \end{tabularx} \\ \\ \relax
5:6 & \begin{tabularx}{0.7\textwidth}{X} 耶穌看見他躺著，知道他病了很久，就問他：「你要痊癒嗎？」 \end{tabularx} \\ \\ \relax
5:7 & \begin{tabularx}{0.7\textwidth}{X} 病人回答他：「先生，水動的時候，沒有人把我放在池子裡；我正要去的時候，別人比我先下去了。」 \end{tabularx} \\ \\ \relax
5:8 & \begin{tabularx}{0.7\textwidth}{X} 耶穌對他說：「起來，拿起你的褥子走吧！」 \end{tabularx} \\ \\ \relax
5:9 & \begin{tabularx}{0.7\textwidth}{X} 那人立刻痊癒，就拿起自己的褥子走了。那天是安息日， \end{tabularx} \\ \\ \relax
5:10 & \begin{tabularx}{0.7\textwidth}{X} 所以猶太人對那被治好了的人說：「今天是安息日，你拿褥子是不合法的。」 \end{tabularx} \\ \\ \relax
5:11 & \begin{tabularx}{0.7\textwidth}{X} 他卻回答他們：「那使我痊癒的人對我說：『拿起你的褥子走吧！』」 \end{tabularx} \\ \\ \relax
5:12 & \begin{tabularx}{0.7\textwidth}{X} 他們問他：「對你說『拿起褥子走』的是甚麼人？」 \end{tabularx} \\ \\ \relax
5:13 & \begin{tabularx}{0.7\textwidth}{X} 那治好了的人不知道那人是誰，因為那裡人很多，耶穌已經躲開了。 \end{tabularx} \\ \\ \relax
5:14 & \begin{tabularx}{0.7\textwidth}{X} 後來耶穌在聖殿裡找到他，對他說：「你已經痊癒了，不要再犯罪，免得你的遭遇更壞。」 \end{tabularx} \\ \\ \relax
5:15 & \begin{tabularx}{0.7\textwidth}{X} 那人就去告訴猶太人，使他痊癒的是耶穌。 \end{tabularx} \\ \\
[1ex]
\hline
\hline
\end{longtable}
$^{1}$請姐妹平安.
感激主讓我們回到神的家去敬拜祂.
這也是神給我們的福氣.
我們再次向上帝禱告.
讓我們有生命氣息.
仍然可以來到你的面前.
願意今天你對我們每一個人說話.
讓我們聽到你的聲音就歡喜.
讓我們的心去回應你.
奉耶穌基督的名字祈禱.
阿們.
今天的題目叫做「你要痊癒嗎?」.
你要不要痊癒?.
這句話是耶穌說的.
究竟耶穌對誰說呢?.
在什麼地方.
在什麼情況之下.
祂說這句話?.
為何祂要說這句話呢?.
我們今天逐一找出這些問題的答案.
「你要痊癒嗎?」這句話.
究竟祂對誰說呢?.
《約翰福音》五章五到六節說.
在那裡有一個人病了三十八年.
耶穌看見他躺著.
知道他病了很久.
就問他.
「你要痊癒嗎?」.
原來這句話是耶穌對著一個長期病患者說的.
一個人病超過十年.
已經是很長的時間.
如果病超過三十八年.
簡直是折磨和煎熬.
曾經在我身邊認識有弟兄姊妹.
病超過二十年.
真的是長期病患.
常常要到醫院接受治療.
身體出現很多傷痕.
我每一次看見她.
心裡都不敢輕易地說.

$^{41}$「求神給你更多信心」.
「行了,最重要是對神有信心」.
這些說話我不會輕易說出口.
我反而經常稱讚對方.
「你簡直是一個勇士」.
「很勇敢」.
「面對這個病這麼久」.
「換作是我」.
「我不知道自己能否做到更堅強」.
今天這個人病了三十八年.
你想想他的痛苦是多麼大.
他的絕望是多麼大.
簡直不是一般人可以明白的.
這個人是躺著的.
可能代表他是癱瘠.
又或者是身體太虛弱.
他沒有辦法站起來.
這個病人躺在不士大池旁邊的儲牢裡.
不士大池在耶路撒冷.
靠近陽門這個地方.
其實我們看經文.
覺得它是一個水池.
實際上你看它是兩個水池.
一個在北面.
一個在南面.
有四個儲牢圍繞著這兩個池子.
在兩個水池當中.
又有一個儲牢隔著.
所以經文裡會提到有五個儲牢.
這兩個水池的水是會動盪的.
當時的人一定知道這一點.
其中一個原因是甚麼呢?.
是他們的水源可能是來自間歇泉.
我上網看一看間歇泉是怎樣的呢?.
原來是地下水變成蒸氣.
然後很間歇地噴出來.
有條柱子.
所以有些東西是會動的.
噴出來的間歇泉會導致水池的水會動盪.
其中一個原因.

$^{81}$猶太人當中有一個傳說.
天使會按時候到水池.
會攪動水.
當水被攪動的時候.
誰首先到水裡.
就一定會得醫治.
無論你甚麼病也好.
多嚴重也好.
最重要是在水中第一個下去.
就行了.
很多猶太人都相信這個傳說是真的.
所以經文裡會寫到有很多失明的.
瘸腿的,癱瘓的.
都躺在不事大池旁邊的儲牢裡.
他們都等待機會進入水池裡.
整件事情引發我們很多思考.
我說的是整個不事大池.
引發我們很多思考.
這個傳說引起很多人的注意.
但我們要問的是.
上帝是否會用這個方法.
去醫治他的百姓呢?.
這個水池的事件.
究竟帶來的是祝福.
還是引發更多問題呢?.
我們來看看.
有很多病患者覺得.
他們的盼望就在水池裡.
最重要是你要成為第一個.
進入水池的人.
那你就有盼望了.
這個機會要怎樣呢?.
你要爭取的.
你不是就這樣躺在那裡.
坐在旁邊.
那你自然就得醫治.
你要爭取第一個進去.
在過程中.
其實其他人是否得到醫治.
是不重要的.

$^{121}$重要的是自己得到醫治.
什麼彼此關心互相扶持.
都是不重要的.
重要的是自己的利益.
為了爭取這個機會.
不惜和別人發生衝突.
在這個過程中我們看到.
人未得到醫治.
首先是心靈受傷.
首先是關係破裂.
在水池邊的世界.
是冷酷的世界.
只顧自己的世界.
我們要問.
上帝會用這個方法嗎?.
不單是這樣.
這個方法令人感到很吃力.
很費勁.
很沒有盼望.
首先你要等.
他沒有說明.
每一天天使會下來.
攪動水一次.
說明是傳說.
你怎知道天使何時來攪動水.
何時來攪動.
那就要等.
尋找機會.
然後你要和別人競爭.
你爭取第一個進水池.
你看看在水池邊的人.
要不是盲的.
懶的.
怎會有力氣進水池.
要花這個力氣.
你要有好朋友幫你.
將你抬進去.
又要等.
又要這樣下去.
很辛苦.

$^{161}$上帝是否真的會用這個方法?.
這個是否上帝用的方法?.
我肯定答案是不用質疑的.
上帝不會用這個方法.
各位醫治人的上帝.
他有很多方法.
其中我和你知道的.
神可以用地面上的醫生.
可以用地面上的藥物.
去醫治我們.
他可以自己直接行神蹟.
還有很多.
有很多方法.
上帝會用.
去叫人徹底康復.
不用用到這一種.
上帝是我們最終極的出路.
我們只要倚靠祂就好了.
但是我們在《約翰福音》這段經文裡.
我們看到的是什麼?.
在處郎裡的百姓.
倚靠的是那兩個水池.
不是上帝.
那個病了三十八年的人.
也是這樣.
他等什麼?.
他的眼睛看什麼?.
那個水池.
不是上帝.
所以當耶穌問他問題的時候.
他回答的是什麼?.
耶穌問他:你要不要痊癒?.
他回答是什麼?.
還是提到那個水池.
提到自己能不能夠下去.
那個就是他幫助.
其實我們看到.
上帝的百姓.
是這些人.
他們不認識神.

$^{201}$在耶路撒冷這一帶的人.
應該都知道.
上帝大概是怎麼樣.
對猶太的文化.
各方面都有了解.
對猶太人的信仰.
也都有了解.
但是你看到這裡.
這群人不認識神.
甚至是倚靠錯的對象.
今天的我們又如何呢?.
我們雖然踏足在教會裡.
都說自己是相信上帝的人.
但是我們認識神嗎?.
我們倚靠的是什麼?.
在我們的生命裡.
是不是也有一個不是大詞呢?.
而一個不是大詞.
可能是一個人.
或者一群人.
無論他們說什麼.
給什麼建議都是好的.
他們就是我們的出路.
在我們心中地位崇高.
每當我們生命遇到困難的時候.
第一時間想到的就是他們.
我們向神祈禱就先找他們.
他們是最高權威.
說什麼都是對的.
都是好的.
我都跟隨的.
有沒有這樣的人.
這樣一個人.
這樣一群人呢?.
在你生命裡.
我們會花很多時間.
跟這些我們覺得很重要的人去傾談.
我們抓住他們.
倚靠他們.
而不是倚靠上帝.

$^{241}$我在這裡.
預備這個講章的時候去思考.
我都很慚愧.
我發覺.
我都要向上帝認罪.
在我生命不同的階段.
都曾經出現過.
不同的不是大詞.
是真的.
在不久之前.
當這個疫情好轉的時候.
教會就開始思考.
我們怎樣重新開始.
不同的事工呢?.
我整天都在想.
那些事工怎樣.
找誰幫忙呢?.
怎樣做法?.
接著會遇到什麼情況?.
應該怎樣解決?等等.
與此同時.
我還要繼續手頭上的工作.
所以我常常都在想.
現在做什麼?.
下一刻要做什麼?.
怎樣做?.
怎樣運用時間?.
不停地在想.
腦海裡常常出現的字.
就是惡.
惡.
惡怎樣?.
惡這樣.
是不是整天都惡?.
我變得越來越倚靠自己.
我發現我成為自己的不是大詞.
其實這樣對我一點好處都沒有.
結果呢.
你知道不是大詞.
其實不是幫到別人的.

$^{281}$是我依靠自己也幫不到自己.
反而變得很憂慮.
我記得有一天早上.
坐巴士回教會的時候.
很特別.
那天早上.
神在車上.
很大聲地跟我說.
真的很大聲.
但是你放心.
身邊的人應該聽不到.
不要擔憂.
很清晰的四個字.
不要擔憂.
大聲地跟我說.
上帝知道我擔憂.
所以提醒我.
鼓勵我.
要依靠祂.
我的眼睛什麼時候離開上帝.
什麼時候看著自己.
我就憂慮了.
當我成為自己的不是大詞.
我就憂慮.
我很盼望.
永遠都要看著上帝.
不要看著自己.
今天我們看到這段經文裡.
那些處郎裡的病人.
那個病了三十八年的人.
就是看錯了對象.
看著不應該看的.
反而他要依靠上帝.
他們沒有看到.
面對這個病了三十八年的人.
耶穌做了什麼事情呢?.
祂做了三件事情.
值得我們留意.
第一個.
祂再次被這個病患者又盼望.

$^{321}$這個病人躺在池邊.
只是為了等待一個機會.
就是在水冷的時候.
首先進入水池裡.
得到醫治.
他等這個機會等了三十八年.
都沒有等到.
他對耶穌說.
水冷的時候.
沒有人把我放在池子裡.
我正要去的時候.
別人被我先下去了.
一次又一次.
他感到失望.
耶穌問他.
你要痊癒嗎?.
你是否要康復?.
如果我們是那位病人.
我們有什麼感受?.
我們會不會覺得.
耶穌你明知故問.
有點多餘.
我躺在這裡當然是想醫治.
為什麼躺在這裡這麼多年?.
你還問這個問題?.
其實這句話一點都不多餘.
為什麼不多餘?.
說真的.
這個患病者已經沒有盼望.
在這裡等什麼?.
是不是真的覺得自己.
還會有機會恢復健康?.
我心想.
他不是.
耶穌想讓他心裡再次興起盼望.
再去想起這件事.
你再去思考.
你是否想好?.
我們有沒有遇到一些事情.
是很困難的?.

$^{361}$無論我們試了多少次.
都找不到出路.
試完又試.
在裡面有沒有經歷過.
上帝仍然給我們有盼望呢?.
這是非常寶貴的.
我記得很久之前.
我在舊屋裡有一個很尷尬的經歷.
但時間比較短.
我經歷了絕對的失望,沮喪.
到有盼望.
有一次在家裡.
這是舊屋.
我在家裡洗澡的時候.
浴室的玻璃門被截住了.
當我洗澡準備出來的時候.
完全開不了.
那天晚上我嘗試很努力.
拍攝玻璃門.
拍攝任何東西.
我希望喊醒爸爸媽媽.
因為爸爸媽媽已經睡著了.
希望他們幫我開門.
無論怎樣拍攝,怎樣叫.
原來浴室是非常噘聲的.
聲音完全傳不出去.
我喊了很久.
真的很久.
都沒有人回應.
當時我落到一個很絕望的光景.
我打算在浴室過一整晚.
直到明天天亮.
爸爸媽媽進來的時候才發現.
原來我整晚困在裡面.
我真的試了很久.
試了任何方法都沒有辦法打開浴室門.
那種絕望是很深.
很多事情也在胡思亂想.
想想會怎樣.
我記得那一晚.

$^{401}$我仍然有祈禱.
雖然感到絕望.
但仍然在祈禱.
一直等.
也不是一會兒.
那扇門就打開了.
我就出去了.
一直等.
後來差一點就想.
算了.
等就等吧.
等到天亮.
是不是也很短.
對比這個人這麼長時間.
我算得了什麼呢.
但不是.
不知道為什麼.
有一次祈禱之後.
那扇門就突然打開了.
我就從絕望裡走出來.
我想起整個過程.
想起來真的呼吸大氣.
想想如果不是.
整晚會怎樣過呢.
是真的有盼望.
上帝會給我們.
有盼望.
那時間很短.
將這個時間放大.
放長.
三十八年不容易.
但上帝仍然會給我們盼望.
第二件耶穌做的事.
就是耶穌透過實際行動.
告訴病患.
盼望的源頭在哪裡.
五章八到九節.
耶穌對他說.
「起來,拿起你的玉子走吧」.
「那人立刻痊癒」.

$^{441}$「就拿起自己的玉子走了」.
這個病患病了三十八年.
就憑耶穌一句說話.
立刻好了起來.
三十八年和立刻這兩個時間.
有一個很強烈的對比.
三十八年很長.
顯示了病的嚴重性.
立刻這個時間性.
顯示這個很短.
立刻就這樣就行了.
不拖.
顯示了耶穌的能力.
是何等偉大.
短短的時間就改變了一切.
治好了三十八年的病患者.
藉著這個神蹟.
耶穌告訴病患者.
盼望在他身上.
看就要看著他.
不是看著他不是大慈.
耶穌是透過他的說話.
去醫治病患者.
所以你可以見到.
耶穌的說話充滿了權能.
弟兄姊妹.
今天我們願意花多少時間.
去閱讀神的話語呢?.
我們越多去閱讀神的話語.
就會越認識神.
越能夠從祂那裡得到力量.
每一天.
上帝都是要透過祂的說話.
告訴我們很多事情.
所以盼望我們願意花這個時間.
還有在裡面.
我們不單止明白上帝的心意.
最重要的是在裡面.
經歷上帝的大能.
第三件事.

$^{481}$耶穌向這個人所做的是什麼?.
耶穌向這個患了38年病患者所做的.
是提出愛的責備.
很特別吧?.
他之前被他盼望.
讓他知道盼望的源頭是自己.
第三件事.
要向他提出愛的責備.
在不是大慈旁邊有五個處郎.
裡面躺著很多病人.
有生命的渠腿的貪玩.
耶穌只選了這個病人.
大家有沒有留意?.
其實他可以去到.
每個人都幫他.
你得醫治了.
你得痊癒了.
你又康復了.
你的眼睛可以看到了.
你又可以站起來.
他可以的.
他有能力的.
他每個人都可以去幫他們.
他沒有.
他只選了一個.
選了這個.
患了38年病的人.
為什麼要這樣選呢?.
因為經文沒有說.
我們只可以說.
這是出於上帝的主權.
他的揀選.
還有他的憐憫.
耶穌憐憫這個38年的病患者.
醫治他.
給他徹底的康復.
其實這個病患者應該覺得很感恩.
這麼多人躺在處郎那裡.
都沒有得到醫治.
是他得到醫治.

$^{521}$是不是應該很感恩.
很快樂.
是不是?.
但是你看這個病患者.
對耶穌的反應.
是不是呢?.
很奇怪.
五章九字這樣說.
那人立刻痊癒.
就拿起自己的玉子.
走了.
就這樣走了.
沒有一句多謝的說話.
多謝你.
醫治我.
說一句都可以吧.
沒有.
他很開心.
走了.
好像他覺得這一切都是應得的.
其實這一切神的恩典.
沒有什麼是我們應得的.
後來當猶太人責備這個康復者.
說他幹什麼干犯了安息日.
這裡是什麼一回事呢?.
原來在神的教導裡.
就說到我們要守安息日.
在安息日不可以作供.
作供是什麼意思呢?.
原來只是說.
你不要去做那些.
習慣性的.
有積分的工作.
是不是?.
如果你上班的.
你平時做什麼行業的.
你去到安息日.
你要守安息日.
你就不要上班了.
說得白一點.

$^{561}$原來只是這樣.
但是呢.
後來猶太的拉比.
將它發展到無限大.
有39項事情.
都叫做干犯了安息日.
例如其中一件是很好笑的.
你如果在安息日.
搬東西.
將一個物件.
由一個地方搬到另一個地方.
你就干犯安息日.
所以在安息日.
大家不要搬椅子.
我們在禮堂裡.
不要搬來搬去.
也不要拿東西.
你只要由這個地方.
走到那個地方搬東西.
我們已經犯法了.
有39項.
是人定出來的.
當時猶太人見這個人得到醫治.
他本來是躺著的.
他不能夠站起來.
無論是因為癱瘓的緣故.
身體虛弱的緣故.
本來要替他開心.
他痊癒了.
讚美主也可以.
正常反應是這樣.
他們不是.
你犯法了.
你干犯了安息日.
你想想這些人是什麼心腸.
當猶太人這樣責備.
這個康復者.
你干犯了安息日的時候.
這個人即時反應是什麼.
是那個人叫我.

$^{601}$拿起我的玉子走.
我才干犯安息日.
是那個人叫我拿起.
我就干犯了.
即是什麼.
退卸責任.
他連累我的.
他叫我拿起我才走.
你以為我很想犯法嗎.
第一時間自保.
說他不對.
這個人簡直沒有心肝.
別人問他.
那個人叫你拿起你的玉子走.
叫什麼名字.
真的不知道.
沒有心肝.
我真的很感激你剛才幫了我.
你貴姓.
你也問人家姓甚名誰.
沒有.
走了就走了.
二號就二號.
我應份的.
我應得的.
他沒有心肝.
後來他知道了耶穌的名字.
他知道了耶穌的名字之後.
馬上告訴猶太人.
剛才那個叫我拿起玉子走的.
就是耶穌.
你不要想他出於感恩.
我真的要感激他.
他對我很好.
所以就這樣說耶穌的名字.
仍然是為了自保.
你說就說耶穌好.
是耶穌叫我拿起玉子走.
你不要說我.
這些人.

$^{641}$忘恩負義的人.
即使他經歷上帝的恩典.
他沒有去感恩.
但耶穌仍然是愛他.
在當下.
耶穌是曾經避開了.
當猶太人在責備.
這個康復者的時候.
那一刻耶穌是避開了.
後來他又再回到聖殿裡.
去找這個康復者.
他知道他忘恩負義.
他知道他不感恩.
但他仍然要對他提出愛的責備.
是責備.
他找到他之後.
第一就強調你已經痊癒了.
你康復了.
提醒他.
他其實已經經歷了上帝的恩典.
第二.
你以後不要再犯罪了.
免得你的遭遇更壞.
這個遭遇更壞的意思.
不是說他會有更嚴重的病.
睡得更久.
而是說他將會面對將來的審判.
甚至面對永遠的死亡.
你現在不要再犯罪了.
否則將來的後果更加嚴重.
這裡經文留了空白.
留了空間.
我們可能想問.
耶穌治好這個人的身體.
有沒有寬恕他之前的罪呢?.
你沒有寬恕他之前的罪.
怎可以叫他之後不要再犯罪呢?.
會不會是耶穌寬恕了他之前的罪.
治好他身體.
他為了愛這個人.

$^{681}$他提醒他.
你以後不要再犯罪了.
否則你要知道你有後果.
讓我遇到一個這麼不感恩的人.
我能不能有這份愛呢?.
這個責備也是出於愛.
耶穌很想這個人好.
很想他得夢成就.
今天看到耶穌所做的這三件事的時候.
值得我們有些反思.
第一件事.
當耶穌問你要痊癒嗎?.
你要不要痊癒呢?.
如果他問我們這個問題.
我心想我們會想回答.
我們想痊癒.
想痊癒要怎樣呢?.
那就要看著耶穌.
耶穌是真正能夠醫治我們.
幫助我們那個.
讓我們依靠他.
第二件事很重要.
要有感恩的心.
當耶穌幫助我們之後.
要有感恩的心.
不是覺得什麼都是應當的.
在這裡我想起.
其實神對我一生的恩典很多.
儘管我在這裡數一點.
我希望我這個數算主因是一個開始.
我鼓勵台下弟兄姊妹.
你們自己也去數算主因.
其實神對我們每一個人一生很多恩典.
由我還在沒有福氣的時候.
在那時候已經保護了我.
我後來怎知道.
神在一個很偶然的場合.
讓我思想到.
原來我能夠在媽媽的肚腹裡面.
這麼多個月平平安安.

$^{721}$不是必然.
可以有很多事情發生.
我可以在媽媽的肚腹裡面已經沒有生命.
但上帝從那裡開始保護我.
看顧我.
讓我來到這個世界.
祂供書教學.
這麼多年給我沒有缺乏.
祂養育我.
到今天.
弟兄姊妹你看到台上的我.
肥肥白白.
你就知道上帝有多少恩典.
我不是生在一個富豪的家庭.
給我一無所缺.
吃得好.
住得好.
給我有很多的機會.
不單在香港讀書.
也有機會去海外.
有一些見識.
有一些學習.
是救贖我的生命.
我不覺得.
我比起聖經裡說的那些罪人.
要聖見.
我覺得自己充滿了罪惡.
甚至比他們更壞.
但上帝這麼多年.
救贖我的生命.
洗淨我的罪惡.
還揀選我做他的僕人去服侍他.
這個實在是無比.
在這麼多年的服侍裡.
我經歷了很多高高低低.
上帝從來沒有離開過我.
祂一直建立我的生命.
掃著我.
在我受傷的時候.
醫治我.

$^{761}$讓我呼吸.
很寶貴的是.
上帝有很多愛我的人.
在我身邊.
我直到今天都認定.
為什麼有這麼多人愛我.
不是因為我可愛.
我也很三尖八角.
我也很有性格.
我記得有一次.
我和一個姐妹聊天.
姐妹說我的說話.
很有刺激性.
是的.
我的說話很有刺激性.
某個程度上.
也像我的性格.
我這麼三尖八角的性格.
上帝也愛我.
還感動很多人愛我.
那些人愛我之先.
是因為上帝愛我.
我知道所有這一切.
都不是必然的.
直到最近.
我還在不停地經歷祂的恩典.
每一件大小事情上.
祂都與我同在.
為我開路.
我只是覺得.
我這個這麼軟弱的罪人.
是不配得這一切恩典.
我很盼望.
我這個分享.
只是一個開始.
弟兄姊妹你們回去.
數算著這個恩典.
當我們有一個感恩的心的時候.
生命才有真正的平安.
喜樂,盼望.

$^{801}$你的眼睛才會看著上帝.
才會經常得到力量.
否則的話.
我經常覺得我生命裡.
什麼都是應得的.
上帝欠我的東西多過祂祝福我的東西.
如果我這樣的生命.
一定會充滿懼怕,痛苦和沮喪.
而且在這裡.
更有第二個提醒.
除了要感恩之外.
我們經歷神的恩典之後.
要遠離惡事.
否則後果會更加嚴重.
就好像耶穌對這個病了38年的人所說.
信了耶穌之後.
不代表我們不會犯罪.
因為有很多軟弱.
弟兄姊妹.
盼望我們一起在上帝面前去審查.
看看我們有什麼地方得罪神.
真的要遠離惡事.
我舉個例子.
如果以前.
未信耶穌前我拜開偶像.
信耶穌後.
要徹底離開一切這些拜偶像的行為.
是徹底的離開.
我剛才在主日學也說.
上帝創造天地.
起初神創造天地.
全世界都由祂而來.
祂是唯一的主宰.
所以我們跟隨這個上帝.
沒有別的.
所有拜偶像的行為.
一定要徹底離開.
如果我在未信耶穌前.
做很多壞習慣.
很多惡劣的行為.

$^{841}$信耶穌後全部都要去脫離.
這個患了38年病的人.
究竟之前犯了什麼罪.
經文裡沒有說.
但肯定很嚴重.
否則不會病了38年.
耶穌警告他.
之後不要再犯罪.
今天耶穌也要透過這段經文.
警告我和你.
以後不要再犯罪.
今天就在神面前安靜一下.
我們有什麼地方得罪上帝.
我們是不是真的很敬畏上帝.
按照上帝的教導而行.
有得罪上帝的地方.
讓我們認罪.
真是悔改.
好讓上帝的心意得到滿足.
你要全愈嗎.
你的答案是什麼.
我們一起圖告.
藉著這段經文提醒我們.
你很愛我們.
我們的生命很差很軟弱.
你也愛我們.
你很想我們改變.
你不只是想我們經歷你的恩典.
而是你要我們改變.
成為一個真正的.
夢你賜福的人.
一個真正的.
感到生命滿足的.
快樂的人.
天父你很願意醫治我們的生命.
求你讓我們的心也願意被你醫治.
你幫我們離開罪惡.
幫我們數算你的恩典.
將榮耀歸還給你.
奉耶穌基督的名字祈禱.

$^{881}$阿們.
\newpage



\section{馬太福音 20:20-28}
\label{sec:k0ge2og3_Pw}
\textbf{2021.05.09《一位母親的要求》馬太福音 20\_20-28( 和修版)  老耀雄牧師}
\newline
\newline
連結: \href{https://youtube.com/watch?v=k0ge2og3_Pw}{\texttt{https://youtube.com/watch?v=k0ge2og3\_Pw}} ~~~~ 語音日期: 2021-05-10
\newline
\newline
\hyperref[sec:FhlMWYYtvXU]{\small{< < < PREV SERMON < < <}}
~
\hyperref[sec:index_chronic]{\small{[返順時目]}}
~
\hyperref[sec:index_scriptual]{\small{[返順卷目]}}
~
\hyperref[sec:yXf3anXpfOw]{\small{> > > NEXT SERMON > > >}}
\newline
\newline
馬太福音 20:20-28
\newline
\begin{longtable}{cl}
\hline
\hline
章節 & 經文 (和合本修訂版)\\
\hline
20:20 & \begin{tabularx}{0.7\textwidth}{X} 那時，西庇太兒子的母親和她兩個兒子上前來，向耶穌叩頭，求他一件事。 \end{tabularx} \\ \\ \relax
20:21 & \begin{tabularx}{0.7\textwidth}{X} 耶穌問她：「你要甚麼呢？」她對耶穌說：「在你的國裡，請讓我這兩個兒子一個坐在你右邊，一個坐在你左邊。」 \end{tabularx} \\ \\ \relax
20:22 & \begin{tabularx}{0.7\textwidth}{X} 耶穌回答：「你們不知道所求的是甚麼。我將要喝的杯，你們能喝嗎？」他們對他說：「我們能。」 \end{tabularx} \\ \\ \relax
20:23 & \begin{tabularx}{0.7\textwidth}{X} 耶穌說：「我所喝的杯，你們要喝。可是坐在我的左右，不是我可以賜的，而是我父為誰預備就賜給誰。」 \end{tabularx} \\ \\ \relax
20:24 & \begin{tabularx}{0.7\textwidth}{X} 其餘十個門徒聽見，就對他們兄弟二人很生氣。 \end{tabularx} \\ \\ \relax
20:25 & \begin{tabularx}{0.7\textwidth}{X} 耶穌叫了他們來，說：「你們知道，外邦人有君王作主治理他們，有大臣操權管轄他們。 \end{tabularx} \\ \\ \relax
20:26 & \begin{tabularx}{0.7\textwidth}{X} 但是在你們中間，不可這樣。你們中間誰願為大，就要作你們的用人； \end{tabularx} \\ \\ \relax
20:27 & \begin{tabularx}{0.7\textwidth}{X} 誰願為首，就要作你們的僕人。 \end{tabularx} \\ \\ \relax
20:28 & \begin{tabularx}{0.7\textwidth}{X} 正如人子來，不是要受人的服事，乃是要服事人，並且要捨命，作多人的贖價。」 \end{tabularx} \\ \\
[1ex]
\hline
\hline
\end{longtable}
$^{1}$主席.
洪映嘉的弟兄姊妹平安.
今日是母親節.
在此祝在在的母親.
母親節快樂.
有人說天下的母親都是偉大的.
她們一生通常都不是為自己.
大部分都是為家庭的成員.
好像為自己的丈夫.
為自己的子女.
甚至是為自己的相親.
似乎在一個家庭裡面.
最願意為家人作出犧牲的.
大部分都是母親.
當然這個不是必然的.
不過統計上.
普遍上.
母親為家庭所犧牲的都是最多.
歷史上有很多見證.
可以找到很多偉大的母親的見證.
在中國.
我們最耳熟能詳.
家傳戶曉的.
就是孟母三遷.
孟子的母親.
本來住在墳場附近.
但是見到兒子學人家拜祭鬼神.
當然不想兒子耳濡目染.
於是就搬走了.
一搬就搬到圖牆附近.
誰知兒子學人家去宰豬.
母親也不想.
於是孟母又開始搬了.
這次搬到太廟旁邊.
孟母這次就很開心.
因為兒子可以有一個很好的讀書環境.
以便日後可以勤力讀書.
將來她可以金榜提名.
為國家盡一分力.
而在人品的品格方面.

$^{41}$也有人會即時想起一個人.
就是岳飛的母親.
岳飛的母親叫什麼名字呢?.
很多人留意,也沒有人記得.
大部分只記得岳飛背後那四個字.
就是什麼呢?.
精忠報國.
岳飛一生精忠報國.
就是他母親姚大夫人的教導.
母親對兒子的影響.
世界任何一個國家都可以找到.
不單單是中國.
在聖經裡面有很多不平凡的母親.
今天我們透過《馬太福音》20章.
20至28節去看一個不平凡母親的生命.
我們一起去祈禱.
你讓我們有家庭這個制度.
讓我們可以透過父母對子女的愛.
去體驗到創天造地的主.
同樣都是愛我們.
求你讓我們對地上的父母的愛.
好像你對我們的愛一樣.
是沒有條件的.
是不離不棄的.
願意為他們作出犧牲的.
將神的愛完完全全地去表達出來.
也讓我們每一個的家庭.
都能夠見證神的愛.
在我們當中彰顯.
我們的祈禱.
奉靠主耶穌基督.
聖名祈求.
阿們.
我們首先了解一下經文的背景.
這段經文是在耶穌第三次向門徒預言.
他將要走上十字架.
被殺第三天復活之後.
雖然門徒已經聽過兩次.
但事實上門徒對耶穌要走的十字架路.
他還未消化得到.

$^{81}$不過相比第一次和第二次.
耶穌說的時候.
門徒的反應已經好了很多.
另外還有一個要處理的問題.
這段經文同時記載在.
《馬可福音》的十章三十五至四十五節.
不過在《馬可福音》的記載裡面.
沒有了西比太兩個兒子.
母親向耶穌的請求.
即是說在《馬可福音》的記述裡面.
這個母親沒有向耶穌請求過什麼.
而向耶穌請求的.
就是西比太那兩個兒子.
我們稍後再處理經文差異的問題.
耶穌第一次預言將要受苦和復活的時候.
當時彼得就立刻斥責耶穌.
他說主啊千萬不要這樣.
這件事絕對不可以臨到你身上.
為什麼會這樣呢?.
因為彼得已經確認了.
耶穌就是基督.
耶穌就是一個永生神的兒子.
所以在彼得的心裡面.
他想像不到一個神的兒子.
竟然會走去受難和死亡.
所以耶穌就回應了.
他說撒旦退到我後面去.
你是我的絆腳石.
因為你不體會神的心意.
而是體會人的心意.
顯示出當時的門徒.
對耶穌這個使命毫無認識.
不過門徒就開始思考.
耶穌所說的每一句話.
耶穌在第二次預言受苦和復活的時候.
門徒當時的反應就是很憂愁.
不過他們憂愁完之後.
就走去問耶穌.
耶穌耶穌天國誰是最大呀?.
門徒這個問題的背後.

$^{121}$反映著門徒已經在關注.
究竟他們有什麼獎賞.
因此當耶穌第三次向門徒.
再預言自己將要受苦和復活之後.
馬太二十至二十一節這樣描述.
他說那時西比太的兒子母親.
和她的兩個兒子上前來.
向耶穌叩頭求他一件事.
耶穌問他你要什麼?.
他對耶穌說.
在你的角裡面請讓我這兩個兒子.
一個坐在你右邊.
一個坐在你左邊.
經文的焦點固然放在西比太兒子母親.
為她兩個兒子所求的事.
就是指這個母親很想為兒子.
在耶穌將來的角度裡面.
安排坐在耶穌的左右兩邊這個位置.
能不能夠預留給她兩個兒子.
雅各和約翰呢?.
這個母親究竟是誰?.
這段經文沒有提供任何的蛛絲馬跡.
只是說她是西比太的老婆.
如果我們看完馬太福音二十七章.
和馬可福音十五章作一個比較的話.
我們就知道這個女子的名字叫撒羅米.
撒羅米又是誰呢?.
我們再對比一下約翰福音十九章.
原來撒羅米就是耶穌母親瑪利亞的姐妹.
那不就是耶穌的姨媽嗎?.
因此約翰和雅各就是耶穌的表兄弟.
當我們知道了這個關係之後.
對這段經文會不會有更多的體會呢?.
撒羅米是不是以一個長輩的身份.
姨媽的身份.
要求耶穌的兒生為她兩個表兄弟提供一個特權呢?.
這個特權就是坐在耶穌的左右兩邊.
兩個家庭有著這麼親密的關係.
都被人感覺到可能會有一些特權.
有人會認為撒羅米在拉關係.

$^{161}$以大壓小,要求耶穌偏私.
尤其是當耶穌其他那十個門徒.
一聽到撒羅米和耶穌有這樣的要求的時候.
經文怎麼說呢?.
就對雅各和約翰很生氣.
這種反應都很正常.
證明覺得不公平.
不是因為你有親戚關係就這樣的.
他們覺得撒羅米在拉關係.
為他們兩個兒子爭上位.
我想在這裡為撒羅米說一句公道的話.
他不是在拉關係.
或者提出一個不合理的要求.
事實上在撒羅米心裡.
他只不過是為他們兩個兒子將來可以得到的東西.
作出一個較好的安排.
為什麼這樣說呢?.
因為耶穌在十九章二十八節.
他曾經對這十二個門徒說過.
他說什麼?.
「我實在告訴你們.
你們這些跟從我的人.
到了萬物更生.
人子坐在他榮耀寶座上的時候.
你們也要坐在十二個寶座上.
審判以色列十二個支派」.
因此撒羅米是按照耶穌對十二個門徒的應許.
而作出這個請求的.
因為按撒羅米的理解.
他們兩個兒子雅各和日漢的確是跟隨著耶穌.
從加利利到現在從來沒有離開過.
因此他覺得他們兩個兒子將來必定會坐在十二個寶座上.
不過還沒有確實的位置.
所以作為母親的他.
在這一刻就想請求耶穌.
可不可以預早訂位.
預早安排了.
坐在你左右兩邊可以嗎?.
這個要求過不過分?.
我們知道.

$^{201}$我們會去接受一個宴會.
坐在主家席上.
跟主家的秘書說.
可不可以坐在主家席左右兩邊.
我一定坐在主家席上.
安排了我坐在主家席上.
不過位置還沒有訂.
可不可以先訂位.
另外還有一個問題還沒有解決.
這個故事同樣記載在馬可福音.
不過在馬可福音35節這樣說.
西比太的兒子雅各和約翰進前來對耶穌說.
老師我們無論求你什麼願你為我們做.
原來在馬可的記載裡.
向耶穌提出請求的不是撒羅米.
而是耶穌的兩個表兄弟.
這件事跟撒羅米一點關係都沒有.
究竟真相是誰向耶穌提出這個請求?.
是撒羅米.
因為他們兩個兒子.
《色經學者》提出.
馬可的成書日期是早於馬太.
因此在馬太的資料裡大部分都是參考馬可的.
為什麼馬可沒有記載撒羅米這個請求.
而在馬太裡有撒羅米的出現呢?.
最大機會是馬太和馬可的作者.
是根據不同的版本來源去編輯的.
事實上有很多《色經家》都說過.
四福音裡不單單只有兩個版本.
而是有四個不同的版本.
即是叫做四元說.
如果這樣我們就解釋到.
四福音裡雖然記載著同一件事.
但會有不同的描述和處理方式.
當然也有一些不能夠查證的解釋.
這個不能夠查證的解釋就是.
馬太的作者不想雅各和約翰.
這兩個門徒擔上污名.
所以就把撒羅米擺上台.
讓撒羅米做這個醜人的角色.

$^{241}$你知道很多時候都是找媽媽.
做醜人的角色.
不過這個說法聽了就算了.
你不要認真地去和人說.
因為沒有證實到.
你不要說Gary說的.
Gary只是小道消息.
聽一下就算了 只能笑.
事實上雅各是最早殉道的門徒.
在仕途行.
撒羅米向耶穌的請求.
既然是在耶穌之前對十二個門徒的應許.
因此他跟進這個應許.
是不是應該很合理?.
關鍵是他不理解這個.
整個應許背後的精神.
撒羅米所想的是兩個兒子將來的榮耀.
也是他應該得到的獎賞.
他兩個兒子從加利利開始.
已經拋下一切跟從耶穌.
已經付出很大的代價.
不過他沒有想過.
他要得到這個獎賞.
還要付上一個更大的代價.
更加沒有想過.
究竟他兩個兒子有沒有能力.
和願不願意付出這種代價.
耶穌對於撒羅米的請求沒有直接回應.
反而轉而詢問撒羅米那兩個兒子.
所以作為母親的撒羅米.
她為兒子求賞賜.
因而詢問耶穌.
耶穌也在處理.
所以在這一刻.
撒羅米的任務已經完成了.
聖經也將他消失在這個舞台之上.
當耶穌向雅各和約翰提出.
我將要喝的杯你們能喝嗎?.
耶穌的意思就是.
將要為天父而作出犧牲生命.

$^{281}$他們的回應是.
我們能.
兩兄弟很清楚知道撒羅米的請求的內容是什麼.
是關於將來的角度.
耶穌對他們兩兄弟的提問就是一個考驗.
測試一下他們是否願意為將來的獎賞.
而作出一個受苦的心理預備.
因為在這一刻.
其實耶穌還沒有走上十教路.
他說我將要.
兩兄弟也不知道將來真實的苦路究竟是怎樣走.
但是他們應該知道.
耶穌要走的一定不是一條很舒服的路.
他們要跟隨耶穌.
要得將來的獎賜.
就必須要付出一個代價.
耶穌要他們確認.
你們願不願意付出代價.
我將要喝的杯.
你們能喝嗎?.
根據聖經的記載.
在耶穌被捕之後.
門徒就四散了.
只有彼得跟隨著耶穌.
看看發生了什麼事.
但我們要知道最後是什麼呢?.
連彼得也三次不認主.
不過在十字架下.
我們看到約翰的蹤影.
證明他沒有離棄耶穌.
事實上雅各也是最早殉道的門徒.
在《少行傳》十二章就記載了雅各被希律王所殺.
而約翰雖然沒有為主殉道.
但晚年被羅馬皇帝多米田放逐到拔摩島.
之後回到以弗所.
寫下啟示錄,約翰書信和約翰福音.
兩兄弟在這一刻對耶穌所作出的「我們能」這個承諾.
日後證實他們真的做得到.
耶穌回覆這兩兄弟.
「我所喝的杯你們要喝.

$^{321}$不過坐在我的左右就不是我可以施的了.
而是我為誰預備就為誰預備了」.
耶穌要走的路.
縱然當時門徒未必知道要付出什麼代價.
但只要立定心意去跟隨耶穌.
其實最後都要好像主耶穌一樣.
為了見證信仰而要作出犧牲.
關鍵是犧牲多少.
撒羅米提出了一個連耶穌都答不到的問題.
因為將來坐在寶座上的左右兩邊.
不是由耶穌決定.
是由天父決定.
從撒羅米向耶穌作出這個請求.
我們可以再追問多一樣東西.
究竟是撒羅米自己的意思.
還是兩個兒子商量過.
而代兒子去提出.
你明白我的意思嗎?.
這是媽媽自己一廂情願去問的.
還是跟兒子商量過.
覺得媽媽你去比我好.
這樣呢?.
經文沒有交代過.
不過我們從馬可福音的記載.
或者從經文的脈絡去看到.
兩兄弟和媽媽一起去.
想去跪在耶穌面前.
所以提問的人是不是兩兄弟.
還是撒羅米呢?.
不重要.
但最重要的是.
兩兄弟對這個問題都很有興趣.
很想知道.
可能三個人商量過.
最後由撒羅米提出都未定.
不過從撒羅米跪在耶穌面前.
真誠的求問.
證明他不是以長輩的身份.
不是以姨媽的身份.
想以大壓小地去提出請求.

$^{361}$另外還有一樣很肯定的.
就是他們對耶穌的說話很留心去聽.
尤其是當他們知道.
耶穌已經是基督神的兒子的身份.
已經確認了.
耶穌在之前沒有直接回應到.
天國裡誰是最大的問題.
他只是說要謙卑像小孩子.
才可以進入天國.
三個人心裡面仍然有一個疑問.
因此今次這個請求.
可以一次過解決.
天國裡誰是最大的問題.
坐在耶穌的左右兩邊就是最大.
是不是?.
耶穌安排我兒子坐在.
他的寶座的左右兩邊.
那還不是最大?.
所以撒羅米的請求.
觸怒到其他門徒是可以理解的.
因為其他門徒同樣都是這樣想.
他們同樣都很想.
不過他們沒有像撒羅米那樣.
直接去問耶穌.
所以當他們聽到撒羅米向耶穌這樣說.
他們就很生氣.
耶穌也知道門徒對天國誰最大的問題.
很關注.
所以最後對所有的門徒.
都作出一個回應.
他說你們中間誰願為大.
就要作你們的用人.
誰願為首.
就要作你們的僕人.
門徒所關心的是.
天國裡誰是最大.
是大與小的問題.
是一種地位.
是一種世界價值觀的制度.
耶穌告訴你.

$^{401}$天國沒有這種制度.
天國沒有大與小.
天國的價值是講服事.
為別人服事.
才是門徒最需要關心的事.
而不是在天國的地位.
耶穌最後也指出.
門徒應該以祂作為榜樣.
正如人子來不時要受人的服事.
乃是要服事人.
並且要捨命作多人的贖價.
因此耶穌的提問.
不單單指向雅各和約翰.
亦向當時所有的門徒提問.
我將要喝的杯.
你們能喝嗎?.
耶穌向十二個門徒問.
我將要喝的杯.
你們能喝嗎?.
我將要為天父而作出的犧牲.
你們能不能夠一樣?.
同樣的提問對今天的信徒.
對今天在座每一位相信主耶穌.
願意跟隨主耶穌的信徒.
同樣有一個提問.
耶穌昔日所喝的杯.
今天我們願不願意和主一起喝?.
弟子們.
請你們在心裡回應.
主耶穌向你們的提問.
剛才的聖餐禮.
你們喝的葡萄汁.
不是純粹喝了葡萄汁.
耶穌昔日所喝的杯.
我們願不願意同樣.
好像耶穌一樣.
一起喝.
一起為天父而作出犧牲.
弟子們.
在撒羅米爾的提問裡.

$^{441}$我們可以有兩個反思.
我們了解到門徒對神的應許是有盼望.
當耶穌應許門徒.
將要坐在十二個寶座上的時候.
門徒就時刻地在想.
究竟我將來究竟坐在哪裡?.
我是否可以坐在耶穌的左右兩邊?.
我是否天國裡面最大?.
雖然仍然未發生.
但對門徒來說是一種盼望.
從好處出發.
盼望可以產生動力.
推動我們成長.
我們從門徒日後的生命可以見到.
十一個都為主而犧牲了.
沒有人有退縮.
全部都很慷慨地為主犧牲了.
從壞處出發.
對盼望作出一個錯誤的理解.
就會造成一種衝突.
因為其他門徒知道撒羅米爾的提問.
就很不開心.
他們的錯誤在哪裡?.
他們的錯誤是認為天國的獎賞要有競爭.
或者天國裡面和世界一樣.
是有階級和制度的.
這是他們的錯誤的理解.
弟兄姊妹,我們對永恆的盼望.
我們今天在座的每一個信徒.
都一定有一個永恆的盼望.
是嗎?.
我們一定有一個永恆的盼望.
這個永恆的盼望對於你.
對於我,是否有一種積極的推動力.
讓我們能夠活得更加精彩.
讓我們活得更加開心.
是因為我們有一個永恆的盼望.
這個永恆的盼望能不能帶給我們一個積極的價值觀.
抑或我們好像昔日的門徒一樣.
我們有一個錯誤的理解.

$^{481}$我們以為天國是有地位之分的.
有金銀銅的冠冕.
我要做得比別人好.
我要做得比別人厲害.
我要受到別人的讚賞.
然後我才能夠拿到一個金冠冕.
得到更大的榮耀.
我們有一個錯誤的理解.
我們要競爭才能得到一個金冠冕.
令門徒將來可以得獎賞.
這是耶穌對他們的應許.
沒錯,這是耶穌對他們的應許.
不過在他們得到應許之前.
是要付代價的.
這個代價就和耶穌一樣.
要飲苦杯.
我們應該記得耶穌說過.
學生不能夠高過老師.
僕人不能夠高過主人.
如果耶穌要走的一條苦路.
作為耶穌基督的門徒.
怎麼可能走一條安舒又舒服的道路.
嚴格來說.
門徒真正為耶穌而飲的苦杯.
以及所要付的十字架.
仍然未出現.
或者這也是我們今天所犯的毛病.
這可能也是我們今天所犯的毛病.
我們相信耶穌.
想去拯救我們.
因此我們願意承認耶穌是我們的拯救主.
是不是?.
耶穌在十字架上已經成就了這個應許.
耶穌在十字架上已經做了這個拯救.
耶穌應許我們如果我們願意認罪.
罪就可以得到赦免.
耶穌做了.
我們願不願意?.
我們願意.
我們願意因著我們的罪求耶穌赦免.

$^{521}$以致我們能夠擺脫罪對我們捆綁.
可以過一個自由的人生.
不過罪得赦免不表示我們可以進入永恆.
是不是?.
我們能夠進入永恆.
還要確認一件事.
就是耶穌是我們生命的主.
我們整個生命要屬於耶穌.
才可以進入這個永恆神的角度.
即是我們要承認我們的生命的主權是屬於耶穌的.
我們的生命主權不屬於我們自己的.
一個屬於神的生命.
是由神去塗造.
神去使用.
以及踐行神的信仰.
如果沒有一個聖潔.
沒有一個神喜悅的生命.
我想問.
我們怎樣進入永恆?.
如果我們的生命污穢.
如果我們的生命充滿污穢.
我們口裡充滿謊言.
怎樣進入永恆?.
我們會不會有一個錯誤的期盼.
我們以為舉手決志.
我相信耶穌.
不需要付任何代價.
這樣就可以得到耶穌的應許.
這是不是我們錯誤的期盼?.
實際上我們要飲的杯.
我們所要付的十字架.
根本未出現.
我們會不會認真去想一想.
我們今天回來這個教會.
參與這個崇拜.
我們願不願意真的放下主權.
背負主耶穌基督.
為我們而設的十字架.
縱然那個十字架仍然未出現.
耶穌要喝那個苦杯.

$^{561}$我們願不願意喝?.
我們願不願意和主耶穌一起.
走上為天父而犧牲的苦路.
如果我們說我不願意.
我不願意.
這個苦路太苦了.
我不願意.
那怎樣可以得到耶穌的應許?.
撒羅米對耶穌.
向十二個門徒所作出的應許.
是有期盼的.
因此她以真誠的態度.
跪在耶穌面前去求問.
她很想兩個兒子.
在耶穌坐在榮耀的寶座上的時候.
可以得到更大的賞賜.
這個就是作為母親.
撒羅米認為跟從耶穌的掌上.
作為母親沒有為自己求.
她只是為她兩個兒子.
真誠地向耶穌祈求.
而她所求的.
是按照著耶穌之前所說的應許.
撒羅米知道永恆國度的重要.
她要確保她的兒子能夠參與其中.
很多人看這段經文的時候.
對撒羅米有一種負面的感覺.
覺得很貪心.
是不是?.
很貪心.
要兩個兒子坐在耶穌的左右兩邊.
你貪不貪心?.
你會不會太貪心?.
覺得是妄求.
是不是?.
你兩個兒子這樣.
可以坐在耶穌的左右兩邊.
我想要為她說一句公道話.
可能是.
我們覺得她很貪心,很妄求.

$^{601}$不過我們看看撒羅米的屬靈生命.
當耶穌在加里利夫,趙雅各和日漢的時候.
兩兄弟就放下工作.
向父親告別.
跟從耶穌.
撒羅米也隨著兩個兒子.
一起跟著耶穌.
途中服侍她兩兄弟.
供應眾人日常所需.
撒羅米雖然沒有被稱為門徒.
但實際上.
她過著門徒的生活.
時時刻刻跟著耶穌.
和耶穌一起生活.
甚至耶穌被釘十字架的時候.
就算門徒四散.
撒羅米也和三個瑪利亞.
仍然守著在耶穌身邊.
撒羅米從加里利.
到耶穌被釘十字架那一刻.
她跟了耶穌三年多.
當耶穌在十字架那裡.
將母親瑪利亞.
交給她所愛的門徒的時候.
很多的《釋經》家都說.
這個耶穌所愛的門徒.
就是撒羅米的兒子日漢.
作為日漢的母親撒羅米.
在這一刻.
會不會感到一份戀意和光榮呢?.
就是自己的兒子日漢.
在這一刻.
耶穌釘十字架這一刻.
我的兒子仍然被主所使用.
而且撒羅米也是第一批.
去到空墳墓見證主耶穌復活的婦女.
當耶穌呼召她兩個兒子.
跟從她的時候.
作為母親的她從來沒有反對過.
甚至耶穌自己的親生兄弟.

$^{641}$都不相信耶穌的時候.
她和她兩個兒子.
繼續陪著耶穌.
她甘心樂意.
兩個兒子被耶穌所使用.
她最大的心願就是.
她兩個兒子能夠在天國裡面得榮耀.
她最想的就是.
她兩個兒子能夠進入永恆.
耶穌所施行過的神蹟.
她參與過.
她見證耶穌是神的兒子.
耶穌所說的每一長道.
她聽過.
她見證耶穌就是基督.
這個就是作為母親撒羅米的屬靈生命.
她從來沒有為自己求過什麼.
她整個生命都是見證耶穌的.
從耶穌的事奉的開始.
直到耶穌的死亡和復活.
她都是耶穌基督的生命的見證人.
今天是母親節.
如果我們作為孩子的父母.
如果我們有小孩.
我們會為孩子求什麼.
我們多數都是為孩子求什麼.
一生幸福.
豐衣足食.
工作暢順.
一帆風順.
婚姻美滿.
很正常.
是不是.
不過真的很可惜.
真的很可惜.
耶穌沒有應許過這件事.
反而耶穌說.
狐狸有洞.
天空的飛鳥有窩.
人子卻沒有枕頭的地方.

$^{681}$對於人在今生的要求.
耶穌只是這樣說過.
你們不要求吃什麼.
喝什麼.
也不要憂慮.
因為這一切都是世上不信的人所尋求的.
你們的父原知道你們需要這一切.
你們只管求他的國.
這些東西都必加給你們.
耶穌要求我們求的是求他的國.
他的義.
作為地上的父母.
我們應該怎樣去為孩子祈求.
再正確一點.
我們應該為孩子求什麼呢.
我們會不會像撒羅蜜一樣.
為孩子能夠在天國裡得榮耀而祈求呢.
我們會不會為孩子的屬靈生命去祈求呢.
我們會不會為孩子能夠回歸神懷抱裡去祈求呢.
我們會不會為孩子日後能夠奉獻給主耶穌而祈求呢.
我們會不會為孩子日後能夠成為宣教士去祈求呢.
十月金國母親節.
我們作為兒女的.
能不能夠為我們父母的屬靈生命祈求.
能不能夠為父母能夠在天國裡得榮耀而祈求.
更加為那些未信耶穌的父母.
他們能夠回轉歸向神而去祈求.
下個月.
六月的第一周.
洪恩嘉就會舉行報道會.
我們就為了這個報道會.
在三樓裡有一棵祈禱樹.
讓那些很想很想那些未信朋友那些人.
你為他們寫一個無花樹的果子.
貼在樹上.
讓我們教睦同工能夠每一天上班的時候.
能夠為他們祈禱獻上禱告.
願不願意呢.
你身邊有沒有一些未信主的朋友呢.
你願不願意為他們獻上禱告呢.

$^{721}$將他們的名字寫在無花果樹上.
貼上去.
為他們獻上禱告.
讓聖靈去軟化他們剛硬的心.
讓每一個父母兒女家人.
都能夠在永恆的角度裡有份.
這是不是我們的心願.
一個作為母親節.
我們作為父母的.
我們想家人信耶穌.
我們作為子女的.
我們想家人朋友信耶穌.
這是不是母親節的心願.
我們一起祈禱.
主求你讓我們追求進入在你的角度裡.
求你讓我們懇切地為我們的家人祈禱.
讓每一個家庭的成員都能夠願意成為.
一生跟隨你的門徒.
甘心樂意被你所使用.
成為你在地上美好的見證.
我們祈禱.
奉主耶穌基督的聖名祈求.
阿們.
願主祝福大家.
\newpage



\section{利未記 25:1-24}
\label{sec:yXf3anXpfOw}
\textbf{2021.05.16《神悅納人的禧年》 利未記 25\_1-24 ( 和修版)  楊穎兒導師}
\newline
\newline
連結: \href{https://youtube.com/watch?v=yXf3anXpfOw}{\texttt{https://youtube.com/watch?v=yXf3anXpfOw}} ~~~~ 語音日期: 2021-05-17
\newline
\newline
\hyperref[sec:k0ge2og3_Pw]{\small{< < < PREV SERMON < < <}}
~
\hyperref[sec:index_chronic]{\small{[返順時目]}}
~
\hyperref[sec:index_scriptual]{\small{[返順卷目]}}
~
\hyperref[sec:HojbxuWCThI]{\small{> > > NEXT SERMON > > >}}
\newline
\newline
利未記 25:1-24
\newline
\begin{longtable}{cl}
\hline
\hline
章節 & 經文 (和合本修訂版)\\
\hline
25:1 & \begin{tabularx}{0.7\textwidth}{X} 耶和華在西奈山吩咐摩西說： \end{tabularx} \\ \\ \relax
25:2 & \begin{tabularx}{0.7\textwidth}{X} 「你要吩咐以色列人，對他們說：你們到了我所賜你們那地的時候，地要休耕，向耶和華守安息。 \end{tabularx} \\ \\ \relax
25:3 & \begin{tabularx}{0.7\textwidth}{X} 你們六年要耕種田地，六年要修整葡萄園，收藏地的出產。 \end{tabularx} \\ \\ \relax
25:4 & \begin{tabularx}{0.7\textwidth}{X} 第七年，地要守完全安息的安息年，就是向耶和華守安息。你們不可耕種田地，也不可修整葡萄園。 \end{tabularx} \\ \\ \relax
25:5 & \begin{tabularx}{0.7\textwidth}{X} 不可收割自然生長的莊稼，也不可摘取沒有修剪的葡萄樹上的葡萄。這年，地要完全安息。 \end{tabularx} \\ \\ \relax
25:6 & \begin{tabularx}{0.7\textwidth}{X} 地在安息年所長出的，要給你和你的奴僕、使女、雇工，以及寄居在你那裡的外人作食物。 \end{tabularx} \\ \\ \relax
25:7 & \begin{tabularx}{0.7\textwidth}{X} 所有的出產也要給你的牲畜和你地上的走獸作食物。」 \end{tabularx} \\ \\ \relax
25:8 & \begin{tabularx}{0.7\textwidth}{X} 「你要計算七個安息年，就是七個七年。這就成為你的七個安息年，一共四十九年。 \end{tabularx} \\ \\ \relax
25:9 & \begin{tabularx}{0.7\textwidth}{X} 七月初十，你要大聲吹角；這是贖罪日，你要在全地吹角。 \end{tabularx} \\ \\ \relax
25:10 & \begin{tabularx}{0.7\textwidth}{X} 你們要以第五十年為聖年，在全地向所有的居民宣告自由。這是你們的禧年，各人的產業要歸還自己，各人要歸回自己的家。 \end{tabularx} \\ \\ \relax
25:11 & \begin{tabularx}{0.7\textwidth}{X} 第五十年要作為你們的禧年。你們不可耕種，不可收割自然生長的莊稼，也不可摘取沒有修剪的葡萄樹上的葡萄。 \end{tabularx} \\ \\ \relax
25:12 & \begin{tabularx}{0.7\textwidth}{X} 因為這是禧年，是你們的聖年；你們要吃地中自然生長的農作物。 \end{tabularx} \\ \\ \relax
25:13 & \begin{tabularx}{0.7\textwidth}{X} 「這禧年，你們各人的產業要歸還自己。 \end{tabularx} \\ \\ \relax
25:14 & \begin{tabularx}{0.7\textwidth}{X} 無論你賣甚麼給鄰舍，或從鄰舍的手中買甚麼，彼此不可虧負。 \end{tabularx} \\ \\ \relax
25:15 & \begin{tabularx}{0.7\textwidth}{X} 你要按照禧年後的年數向鄰舍買；他要按照可收成的年數賣給你； \end{tabularx} \\ \\ \relax
25:16 & \begin{tabularx}{0.7\textwidth}{X} 年數越多，價錢就越高；年數越少，價錢就越低，因為他賣給你的是收成的數量。 \end{tabularx} \\ \\ \relax
25:17 & \begin{tabularx}{0.7\textwidth}{X} 你們彼此不可虧負，只要敬畏你的神，因為我是耶和華－你們的神。」 \end{tabularx} \\ \\ \relax
25:18 & \begin{tabularx}{0.7\textwidth}{X} 「你們要遵行我的律例，謹守我的典章，遵行它們，就可以在那地上安然居住。 \end{tabularx} \\ \\ \relax
25:19 & \begin{tabularx}{0.7\textwidth}{X} 地必出產果實，你們可以吃飽，在那地上安然居住。 \end{tabularx} \\ \\ \relax
25:20 & \begin{tabularx}{0.7\textwidth}{X} 你們若說：『看哪，第七年我們不耕種，也不收藏農作物，我們吃甚麼呢？』 \end{tabularx} \\ \\ \relax
25:21 & \begin{tabularx}{0.7\textwidth}{X} 我必在第六年發令賜福給你們，地就長出三年的農作物來。 \end{tabularx} \\ \\ \relax
25:22 & \begin{tabularx}{0.7\textwidth}{X} 第八年你們要耕種，也要吃陳糧；等到第九年農作物收成的時候，你們還有陳糧吃。」 \end{tabularx} \\ \\ \relax
25:23 & \begin{tabularx}{0.7\textwidth}{X} 「地不可以賣斷，因為地是我的；你們在我面前是客旅，是寄居的。 \end{tabularx} \\ \\ \relax
25:24 & \begin{tabularx}{0.7\textwidth}{X} 在你們所得為業的全地，要准許人有權將地贖回。 \end{tabularx} \\ \\
[1ex]
\hline
\hline
\end{longtable}
$^{1}$弟兄姊妹平安.
多謝洪恩堂老牧師的安排.
讓我可以在實習期間.
有機會再和弟兄姊妹去分享聖經的信息.
今天我們看的經文是來自利未記第25章.
剛才主席讀出第一至二十四章.
但是稍後我們會看多一點的經文.
都是來自第25章的.
第25章所提到的核心價值.
是包著神的公義,恩典和慈愛.
在禱告之後.
我決定講到的題目是.
「神閱立人的欺憐」.
我們會嘗試從經文去明白.
安息憐,欺憐的意義.
然後從而去學習.
活出土神所喜悅的樣式.
我們一起先低頭禱告.
親愛的恩主.
求你以你的話語.
引導我們去認識你.
去學習如何去討你喜悅.
我們同心合意地去祈求.
阿們.
我今天的講題是.
「神閱立人的欺憐」.
其實這一句話.
「神閱立人的欺憐」.
是耶穌說過的.
在我們進入今天的經文之前.
我想和大家一起去看看.
路加福音第四章.
曾經提到.
耶穌一出來傳道的時候.
祂就在會堂讀出這一段經文.
主的靈在我身上.
因為祂用羔羔我.
叫我傳福音給貧窮的人.
差遣我宣告.
被擄的得釋放.

$^{41}$失明的得看見.
受壓迫的得自由.
然後第十九節.
宣告神閱立人的欺憐.
根據路加福音這個技術.
神的靈是叫耶穌去傳福音給貧窮的人.
耶穌的工作就是要去傳福音.
但是當我們細心去回想.
耶穌所傳的福音.
其實就是向那些被鬼附的.
他們被擄的得釋放.
我們去.
耶穌是向失明的去宣告.
他們現在可以再看見東西.
向從前受壓迫的.
無論是肉身還是心靈.
如今他們都得到自由.
這些很多耶穌的事蹟.
我們都記得.
我們都明白.
然後耶穌就說.
宣告神閱立人的欺憐.
問題就來了.
甚麼是欺憐呢?.
甚麼是人的欺憐呢?.
似乎欺憐這個概念.
或者當中的價值.
是包含著神喜歡我們做的一些事情.
而且是和傳福音給貧窮人有關.
而當我們去做這些事情的時候.
神就會對人產生一種喜悅之情.
閱立人所做的事.
我相信我們在座的每一個.
都很想做一些討神喜悅的事情.
所以今天很想和大家一起去思想一下.
甚麼是欺憐.
簡單來說.
欺憐就是當時神設立的一個經濟體系.
而這個經濟體系.
是和我們今天一般認識的運作.

$^{81}$有些不同.
要明白這個經濟體系.
首先我們就需要知道.
神是很關心這個大地的.
從亞伯拉罕以撒雅各.
神對他們的應許.
從來都沒有離開過這片大地.
而且我們知道.
這些列祖以往都是以耕種牧羊為業.
因此土地就是他們的資本.
財富的來源.
人要有地.
藉著耕種,藉著牧羊.
人才能夠有工作.
有修成,能夠自給自足.
經文的背景就是.
摩西在帶領以色列民出了埃及.
在抗日未進入迦南美地之前.
神預先為這班以色列人.
制定好一套經濟的政策.
內容都是環繞著這片大地.
為甚麼呢?.
神親自解釋了.
第23節.
地不可永邁,因為地是我的.
你們在我面前是客裡,是寄居的.
在這裡.
神強調土地的主權是屬於他的.
是他允許他的子民暫時住在這個地上.
去耕種,去享用地的出產.
意思就是住在地上的人.
都只不過是地上的管家.
管家要做甚麼呢?.
如果大家記得創世記一開始的時候.
教導過我們.
我們要看守,修理,管理這個大地.
讓這個大地時刻都展現著神美好的創造.
神讓人能夠在這片大地上工作.
同時神對地的優色都有指引.
第3節.

$^{121}$你們六年要耕種田地.
六年要修整葡萄園.
收藏地的出產.
第7年.
地要守完全安息的安息年.
就是向耶和華守安息.
你們不可以耕種田地.
也不可以修整葡萄園.
不可以收割自然生長的莊稼.
也不可以摘取沒有修剪的葡萄樹上的葡萄.
這一年地要完全安息.
因為耶和華所關心的是這片大地的狀況.
所以在這裡他吩咐以色列人.
你們六年就要耕種田地.
去修理葡萄園.
去收割地裡的出產.
但是到了第7年.
地就要向耶和華守安息.
人可以耕種,不可以修整.
不可以收割,不可以摘取.
因為耶和華要地有完全的安息.
留意耶和華是要地有完全的安息.
不是要人有完全的安息.
人耕種不修整,不收割,不摘取.
都是要配合地要向耶和華守安息這個條件.
經文特別提到修整葡萄園.
如果大家有種植過一些欄藤類的植物.
你就知道.
好像葡萄樹,它要結果子的話.
它一定要修剪那棵植物.
一般來說,每年要修整兩次.
冬天一次,夏天一次.
就要剪掉前一年沒有修整的枝節.
當雨季一來的時候.
葡萄樹就會按時結出果子.
然後就可以將這些果子剪下來去收藏.
雖然第7年是不能耕種.
但是田地裡的葡萄園.
仍然會自己生長,有莊稼.
園主就不可以耕種,又不可以收割.

$^{161}$因為這些收成是要留給低下階層的人.
和一些動物,一些走獸去食用的.
目的是要讓一切在大地上生活的生物.
都可以有食物吃.
我們看看經文,經文是這樣說的.
第6節,地在安息年所長出的.
要給你和你奴僕,侍女,故宮.
以及寄居在你那裡的外人食用.
所有的出產也要給你的畜生和你地上的走獸作食物.
在這裡,奴僕,侍女是指長期在家裡工作的工人.
故宮和奴僕,侍女又不同了.
故宮通常是指散工.
今天來說可能是做燈,水喉.
搬一次貨物的職業.
在當時來說,故宮就不像僱主那樣有足夠的儲備.
他們每一天都要靠微薄的工錢來過活.
所以僱主在僱用故宮的時候.
要很敏感,故宮是金錢上的需要.
如果僱主將這些薪金在第二天才發給故宮的話.
故宮那一晚可能就沒有東西吃了.
而寄居的人是指暫時來到居住的人.
他們要依靠當地的人給予他們工作的機會.
或者要依靠一些社會的政策來維生.
在今天來說,這些寄居的人可能是一些家庭傭工,難民.
一些去旅行工作的年輕人.
甚至是一些外勞.
神很顧念這片大地上的勞工階層.
經文提到,這個安息的福祉是擴展到跟人有關係的動物.
由此可見,安息年背後的精神不是叫人一直很辛勞地工作的人.
他們在安息的年代就去休息,甚麼都不做.
做一個懶惰的人.
不是這樣的.
反而神是要指紋去實踐,去照顧其他人的需要.
讓生活在這片大地的每一個人都體會到神的豐盛.
以至萬物都是因著神自然的創造而能夠從土地得福.
在安息年,地要有完全的安息.
神要讓這些土地慢慢恢復它原有的養分.
地上的地主會因著地要有完全的安息而停止工作.
目的其實是要叫這些地主,他們有些累積財富的慾望要截斷.
神容許這些地主無止境地開發,採摘,剝削甚至破壞.

$^{201}$有多少賺多少.
對這些地主來說,如果第七年又不可以耕種又不可以收割.
那安息年吃甚麼呢?.
我們看看經文.
「我必在第六年法令賜福給你們,地就長出三年的農作物.
第八年你們要耕種,要吃陳糧.
等到第九年農作物有收成的時候,你們還要吃陳糧」.
陳糧就是指之前儲下的糧食.
我們在這裡看到神看顧人無微不至.
第六年有三年的收成.
是供給第六,第七,第八甚至第九年都有得吃的.
是神親自去祝福這片大地.
所以安息年亦提醒我們.
人所需要的,神會親自去負責.
而人就儘管安然地生活在神的愛裡.
在安息年裡經歷神與人同在的應許.
第十八節.
「你們要遵行我的律法,我的律例,謹守我的典章.
遵行它們,就可以在那地上安然居住.
地必出產果實,你們可以吃飽.
在那地上安然居住」.
神很強調,祂很想人在這個大地上安然居住.
這是安息年的精神.
當我們明白甚麼是安息年的時候.
我們就看看甚麼是欺年.
第八節.
「你要計算七個安息年,就是七個七年.
這就成為你的七個安息年,一共四十九年」.
第十一節.
「第五十年就是聖年,也就是欺年」.
原來第五十年就是欺年.
「五十年一個循環」.
欺年的核心價值其實與安息年是一脈相承的.
第十一至十二節.
「在欺年,你不可以耕種,不可以收割,不可以摘取.
因為這是欺年,是你們的聖年.
你們要吃地中自然生長的農作物」.
那這個欺年與安息年有甚麼分別呢?.
這個欺年有甚麼特別呢?.
第十節,經文說.

$^{241}$「在全地,要向所有居民宣告自由.
這是你們的欺年.
各人的產業要歸還自己.
各人要歸回自己的家」.
這裡提到各人的產業要歸還自己.
各人要歸還自己的家.
是甚麼意思呢?.
原來第25章,其實今天沒有讀到的經文.
其實提到一個社會的現象.
我很簡略地歸納了一些經文.
原來有一些弟兄在這個社會上漸漸貧窮.
他們面對的都是一些經濟的困難.
經文沒有解釋他們為甚麼貧窮.
只是提到他們漸漸貧窮.
貧窮到一個地步.
他們需要賣掉自己土地的房產.
他們的資產.
甚至他們要賣掉自己去做地上的奴隸.
然後到了欺年.
神就宣告他們「你們自由了」.
各人都要取回自己原有的產業.
而且他們都要回到自己的家.
這個欺年就是神要向全地的人宣告你們自由了.
在這個欺年的經濟體系之下.
我有三個討神喜悅的價值.
很想跟大家去思想.
第一,就是欺年使人有重新開始的機會.
也就是說設計這個欺年.
是要人重新開始.
重新起來.
過一個新的生活.
這是我們的神.
在這個他設計的經濟體系之下.
他希望人最終得到自由.
試想想.
50年.
一個勞碌了50年的工人.
曾經經歷了漸漸貧窮.
需要變賣自己的資產.
失去了人生的自由.

$^{281}$而去換取生活的所需.
但現在到了第50年.
欺年就是強調.
重新來過了.
你得回你自己的土地了.
你有回你自己的自由了.
你不用你過去發生的事.
你重新再來過.
由此可見.
欺年其中一個的核心價值.
其實不是要解決短期的溫飽這個問題.
而是要讓窮苦的人.
在長遠來說可以重新開始.
50年一個循環.
到了第50年.
這個欺年就完全改變了.
從前是勞工階層的宿命.
一切可以重新開始.
我們今天聽起來.
很不可思議.
我相信今天.
在地上都沒有一個國家.
會實行這個50年一次的經濟大赦.
但這就是我們的神.
對人生活的心意.
充滿了他的恩典和慈愛.
縱然地上的制度.
說它不容許我們去實踐.
這個欺年這種實際的運作.
但並不代表當中的價值.
不能夠由我們每一個去實踐.
欺年是強調對弱勢社群的關懷.
《新命紀》第15章10至11節提到.
因為地上的貧窮人.
永遠都不會斷絕.
意思就是神不會完全解決地上貧窮的問題.
試想想.
對神來說.
要解決貧窮的問題.
是會幫助人去經歷神多一點.

$^{321}$還是要神的子民.
去學習關懷弱勢社群.
從而去經歷神多一點呢?.
如果神去解決了貧窮的問題.
我們還有甚麼可以去實踐,去參與呢?.
所以安積年,欺年的意義.
不是要我們去休息.
不是要我們去放大價.
不是要我們用我們的努力.
在這一年再賺得更多.
而是要透過養地去守這個安積.
從中去學習分享我們有餘的.
我們生活上有餘的.
跟別人分享.
使地上所有人都可以有飽足.
吃得飽.
這是神要我們在頭腦上.
和實際去實踐的時候.
這兩方面我們都要成熟起來.
滿有神待人處事的標準和樣式.
而當神看到我們這樣做的時候.
神就喜悅了.
遵行我的律例.
謹守我的典章.
遵行它們.
就可以在這個地上安然居住.
和大家分享一下.
在我家一入門口.
我放了一個小豬的前鞍.
這個前鞍其實是提供給我們一家人.
誰去完街.
可能去超市買完東西.
吃完飯.
家中還有些散銀的時候.
就進去.
這個除了將自己有餘的.
將來可以存起來.
之後捐給有需要的群體之外.
其實是無時無刻提醒我們自己.
上帝每日所賜的.

$^{361}$都是充充足足.
都是有餘的.
所以有時回到家看到這個小豬.
就很開心.
又有些錢進去.
第三個價值想和大家分享的就是.
不可以虧虧人.
原來神設立的這個體系.
是要去維持一種公平的原則.
而且神是要求有權有勢的人.
去學習放下擁有財富.
要學習去放手.
那個要求就是.
每五十年土地就要歸還原主.
目的就是要避免有人可以漸漸囤積一些土地.
漸漸有人可以不勞而獲.
而造成更大的貧富懸殊.
意思就是避免富有的後代.
慢慢變成富二代.
貧窮的後裔就慢慢變成貧窮的二代.
在這個過程中.
神對有資產有財力雄厚的人有要求.
第十七節.
「你們彼此不可虧虧.
只要敬畏你的神.
因為我是耶和華你們的神」.
原來不虧虧人就是敬畏神.
換句話說.
虧負人就是不敬畏神.
在這裡經文很強調你們的神.
就是跟那些以色列的子民說.
你們的神.
耶和華是要提醒子民.
我耶和華是你們的神.
昔日我耶和華你們的神.
跟你們做過甚麼事.
是我拯救你們那些圍爐子民出埃及.
意思就是要提醒那些有財力的人.
你們也做過別人的奴隸.
你們今天不要反過來欺壓人.

$^{401}$第十八節.
「我是耶和華你們的神.
曾經領你們出埃及地.
為要把迦南地賜給你.
要作你們的神」.
其實邏輯就是.
雖然子民中有富貴之分.
但是無論哪一種人.
今天他們的身份地位是怎樣.
他們的起始點都是奴隸.
這班子民是曾經在埃及做過寄居的.
做過奴僕的.
今天有能力財富.
完全是耶和華昔日的拯救他們.
如果當日耶和華沒有拯救他們.
他們今天可能仍然是奴隸.
是耶和華使他們從奴役的生活中.
拯救了出來.
使他們脫離了奴隸的身份.
可以得到自由.
所以今天這班有財有勢的人.
他們仍然要效法耶和華.
去幫助一些弱勢社群.
讓其他人都能得到這份自由.
這其實是耶和華對強者的提醒.
大家都是在這個起始點之上.
來抗論.
就是曾經在埃及寄居做過奴隸.
因此大家不應該強調彼此有強弱之分.
有貧富之分.
反而是要努力的去將這個差距減少.
第五十五節.
因為以色列的人.
就是我領出來的人.
都是我的僕人.
他們是我的僕人.
是我領他們從埃及出來的.
沒有分別的.
我是耶和華你們的神.
怎樣才不是虧負神呢?.

$^{441}$我們看看第十四至十六節.
神舉了一個例子.
第十四節.
無論你賣甚麼給論社.
或者從論社手中買甚麼回來.
彼此不可虧負.
要按著欺年後可以收成的年數去定地價.
意思就是距離第五十年欺年越遠.
你的地價就可以賣得貴一點.
因為那五十年裡面.
那個收成是可以多一點的.
越接近欺年.
那個地價就越來越低了.
因為所剩下的收成不多.
就著神所設立的經濟原則.
我們看看這個大地發生了甚麼事.
現在在印度的疫情.
我們知道.
人們需要買一些氧氣.
黑市的一瓶氧氣原來是四百多元.
即是四百多元.
但現在其實是賣到五千元一瓶.
其實是坐地起價.
我們看到.
神創造的空氣要用錢買嗎.
我們又回到本港.
我們去看看一些tong 房的情況.
我們都知道.
它租的呎價其實是比你租豪宅還要貴.
我們感恩政府.
接下來有一些tong 房.
去制定一些tong 房的政策.
我們都求天賦賜.
官員有智慧.
希年是強調一個公平的交易.
這些例子公不公平.
一開始已經提過了.
希年的大前提就是地是屬於耶和華的.
我們是地上的管家.
看守 修理.

$^{481}$在這個過程中.
有沒有見到有人虧負了人.
有沒有人不敬愛神.
總括今天的訊息.
希年是提醒人.
神如何恩待過子民.
現在子民就要在頭腦上擁有這些價值.
並且要用生命去實踐.
去回應其他人生命的無常.
在我們的程序表裡.
其實都列了很多的禱告事項.
很鼓勵大家回到家裡.
為教會弟兄姊妹去守望.
特別想和大家介紹一個.
社會關懷及宣教的部分.
這個袋土事項.
其實是根據發生在這片大地之上的事情去編寫的.
當中有宣教的需要.
有國際新聞.
以及本地的需要.
除此之外.
除了這份東西大家可以拿回去之外.
我想問問大家有沒有人看新聞的習慣.
舉手.
有沒有人看新聞的.
有呀,很多人看新聞.
其實很鼓勵大家平時在家裡.
看新聞的時候.
去向神祈禱.
就算你在新聞裡面看到每天發生的事情.
我們去求神去顧念那些不同階層.
去失去自由.
無論他們是患病.
心靈受痛苦.
受到迫激的人.
我們都去紀念他們.
這是一個屬靈的操練.
因為我們透過我們的土局.
我們讓神去知道.
我們關心大地上發生的事情.

$^{521}$相信我們這樣做的時候.
神就會閱立我們的心.
願我們都能夠一同學習.
一起低頭祈禱.
親愛的恩主.
多謝你藉著今天的經文.
教導我們安息年.
熙年的精神.
你所設立的體系.
是要讓人有重新開始的機會.
透過我們各人對身邊的人.
去施予憐憫幫助.
去改變這個社會.
去建立人與人之間的尊重.
那份不分彼此的關係.
敬拜從來都不是單單.
我們星期日回來教會去敬拜你.
但願你親自幫助我們.
使我們在所行的一切事上.
都不做虧負人的事情.
以表示我們對你真誠的敬拜.
主啊 求你開闊我們的眼界.
讓我們看到世界上人的苦難.
在禱告之中與人休戚與共.
我們如此的禱告.
乃是奉愛我們.
為我們犧牲寫命的主的名而求.
阿們.
\newpage



\section{帖撒羅尼迦前書 1:1-10}
\label{sec:HojbxuWCThI}
\textbf{2021.05.19《「逆/疫」是有信,有望,有愛》 帖撒羅尼迦前書 1\_1-10 老耀雄牧師}
\newline
\newline
連結: \href{https://youtube.com/watch?v=HojbxuWCThI}{\texttt{https://youtube.com/watch?v=HojbxuWCThI}} ~~~~ 語音日期: 2021-05-19
\newline
\newline
\hyperref[sec:yXf3anXpfOw]{\small{< < < PREV SERMON < < <}}
~
\hyperref[sec:index_chronic]{\small{[返順時目]}}
~
\hyperref[sec:index_scriptual]{\small{[返順卷目]}}
~
\hyperref[sec:0Rk0UlP9K50]{\small{> > > NEXT SERMON > > >}}
\newline
\newline
帖撒羅尼迦前書 1:1-10
\newline
\begin{longtable}{cl}
\hline
\hline
章節 & 經文 (和合本修訂版)\\
\hline
1:1 & \begin{tabularx}{0.7\textwidth}{X} 保羅、西拉、提摩太寫信給帖撒羅尼迦在父神和主耶穌基督裡的教會。願恩惠、平安歸給你們！ \end{tabularx} \\ \\ \relax
1:2 & \begin{tabularx}{0.7\textwidth}{X} 我們為你們眾人常常感謝神，禱告的時候提到你們， \end{tabularx} \\ \\ \relax
1:3 & \begin{tabularx}{0.7\textwidth}{X} 在我們的父神面前，不住地記念你們因信心所做的工作，因愛心所受的勞苦，因盼望我們主耶穌基督所存的堅忍。 \end{tabularx} \\ \\ \relax
1:4 & \begin{tabularx}{0.7\textwidth}{X} 神所愛的弟兄啊，我知道你們是蒙揀選的； \end{tabularx} \\ \\ \relax
1:5 & \begin{tabularx}{0.7\textwidth}{X} 因為我們的福音傳到你們那裡，不僅在言語，也在能力，也在聖靈和充足的確信。你們知道，我們在你們那裡，為你們的緣故是怎樣為人。 \end{tabularx} \\ \\ \relax
1:6 & \begin{tabularx}{0.7\textwidth}{X} 你們成為效法我們，更效法主的人，因聖靈所激發的喜樂，在大患難中領受了真道， \end{tabularx} \\ \\ \relax
1:7 & \begin{tabularx}{0.7\textwidth}{X} 從此你們作了馬其頓和亞該亞所有信主的人的榜樣。 \end{tabularx} \\ \\ \relax
1:8 & \begin{tabularx}{0.7\textwidth}{X} 因為主的道已經從你們那裡傳播出去，你們向神的信心不只在馬其頓和亞該亞，就是在各處也都傳開了，所以不用我們說甚麼話。 \end{tabularx} \\ \\ \relax
1:9 & \begin{tabularx}{0.7\textwidth}{X} 因為他們自己已經傳講我們是怎樣進到你們那裡，你們是怎樣離棄偶像，歸向神來服侍那又真又活的神， \end{tabularx} \\ \\ \relax
1:10 & \begin{tabularx}{0.7\textwidth}{X} 等候他兒子從天降臨，就是神使他從死人中復活的那位救我們脫離將來憤怒的耶穌。 \end{tabularx} \\ \\
[1ex]
\hline
\hline
\end{longtable}
$^{1}$第一姐妹平安.
平安.
剛才陳牧師說.
我很支持新界西與牧師公.
其實在我在錦宣服侍的時候.
已經開始在新界西與牧師公服侍.
所以當然來到紅欣更加需要.
無論我和新界西與牧的關係.
或者是那個場地.
今天是一個假期.
你們不來這裡也是空了.
不如借給你們吧.
很高興能夠和大家分享神的話語.
如果有人問你.
教會是一個什麼人來的地方.
大部分人都會說.
教會當然是聚人聚集的地方.
但如果我們再深入一點再問.
聚人聚集在教會有什麼目的.
我想每一個人都有不同的答案.
這些細微的答案.
最後都會指向同一個方向.
聚人需要在教會裡面更加認識神.
去被牧養.
讓他們的屬靈生命要不斷地成長.
而最終的目的.
就是要見證神.
榮耀神.
將神的救恩講全.
一個能夠見證神生命的人.
仍然是一個土神喜悅的信徒.
我相信在座這麼多.
新界西縣牧事工的義工.
你們參與新界西縣牧的每一個服事.
其實我深信.
你最終都很想去榮神益人.
做一個土神喜悅的信徒.
今天我的講題就是.
在逆境裡面的信望愛.
在初期教會的歷史裡面.

$^{41}$究竟有沒有像今天香港一樣.
或者好像全球都發生的疫情呢?.
聖經似乎沒有提過.
當時有這麼大規模的疫情.
但是對於信徒處身在逆境之中的經歷.
聖經就有很多的描述.
我們從帖撒萊迦前書一章一至十節的經文裡面.
透過保羅對帖撒萊迦教會的教導.
讓我們學習怎樣在逆境裡面.
去實踐我們的信仰的關鍵.
我們一起有個祈禱.
就是因著我們求你臨到在我們當中.
開我們的心.
讓我們將你的話語栽種在內心深處.
成為承托我們信仰生命的根基.
好叫我們的信仰的根基穩妥.
不會隨風飄去.
隨水漂流.
成為一個無根的追求者.
我們的祈禱是奉教主耶穌基督.
得勝明智祈求.
阿們.
開始的時候.
我們不如先讀一下經文.
先讀一下那十節的經文好不好.
如果你們.
我知道耳目事公已經將一個PDF檔案給他們.
如果不是你可以看一下螢光幕.
我們一起先讀那十節經文.
好預備一二三.
保羅西拉提摩泰寫信級貼撒羅尼加再父神和主耶穌基督裡的教會.
願因為平安歸於你們.
禱告的時候提到你們.
在神我們的父神面前.
不住的紀念你們因信心所作的功夫.
因愛心所受的勞苦.
因嚮往我們主耶穌基督所傳的忍耐.
被神所愛的弟兄.
我知道你們是無揀選的.
因為我們的福音傳到你們那裡.

$^{81}$不獨在乎言語.
也在乎權能和聖靈.
並充足的信心.
正如你們知道.
我們在你們那裡.
為你們的緣故是怎樣為人.
並且你們在大能之中.
蒙了聖靈所賜的喜樂.
領受真道就效法我們.
也效法了主.
甚至你們作了馬其頓和阿該亞.
所有信主之人的榜樣.
因為主的道從你們那裡已經傳揚出來.
你們向神的信心.
不但在馬其頓和阿該亞.
就是在各處也都傳開了.
所以不用我們說什麼話.
因為他們自己已經報明.
我們是怎樣進到你們那裡.
你們是怎樣離棄偶像.
歸向神.
要服事那又真又活的神.
等候他而止.
從天降臨.
就是他從死裡復活的.
那位教我們脫離將來憤怒的耶穌.
開始的時候.
我想交代一下.
帖撒萊加前書的背景.
這間教會是保羅在第二次宣教旅程的時候建立的.
當時他和一切的同工.
還有西拉和提摩泰.
帖撒萊加前書的成書其實很早.
估計耶穌被釘十字架之後.
二十年左右就已經書寫成功.
根據《使徒行傳》第十七章.
一至九節記載.
記錄了什麼呢?.
他說保羅在帖撒萊加城的一個愛邦信徒.
耶孫的家裡.

$^{121}$至少有連續三個安息日的時間.
向愛邦人進行傳福音的工作.
這個傳福音工作.
做成一個結果.
就是有些猶太人.
以及一些希臘人信主.
當中也包括了一些不少的婦女.
不過很快就引來不信的猶太人迫害.
引起暴動騷亂.
保羅和西拉差點就被人抓了.
所以當時的信徒.
就安排他們立刻離開帖撒萊加.
保羅當時很想要傳福音.
但是面對很多的迫害.
在一個逆境的情況下.
保羅和西拉就只能夠離開了.
縱然他們肯離開.
但是這批不信的猶太人.
仍然不放過保羅.
就算保羅離開了帖撒萊加.
去到比利亞.
他們仍然跟到那邊.
而且去挑動眾人去抓保羅.
最後保羅又只能夠留下西拉和提摩泰.
自己一個人.
走到雅典去避開他們.
是一個我們叫做扯進扯退的環境.
其實很多人都會放棄的.
但是保羅沒有.
他仍然很掛念這些福音的果子.
保羅離開了帖撒萊加.
建立教會的牧者都走了.
當初剛才所說的那些信主的猶太人.
希臘人和他們的婦女.
誰去牧養?.
因此帖撒萊加的教會.
是由一群初信主的信徒所組成.
他們缺乏了牧者埋身的牧養.
而且面對很多的挑戰.
所以不單單保羅面對這個逆境.

$^{161}$就算這批初信主的信徒.
他們同樣面對一個逆境.
相信沒有人會說不同意.
牧者因為逆境的緣故.
離開帖撒萊加.
到最後要走到雅典避.
不是叫避難.
避過被人拘捕.
那批初熟的果子沒有了牧者.
面對信仰的疑惑.
面對社會文化的衝擊.
同樣都是一個逆境.
因為他們的信仰生命隨時會夭折.
實在令人擔心.
所以保羅在雅典和西拉和提摩泰會合後.
他立即打發提摩泰到帖撒萊加.
了解教會的情況.
這很正常.
後來提摩泰在哥林多再和保羅會合.
就這事帶來帖撒萊加成了近況.
保羅在這樣的環境下.
去寫帖撒萊加前書.
建立教會的牧者保羅.
和這批帖撒萊加的信徒.
同樣都處身在一個逆境裡.
究竟他們怎樣去面對呢?.
保羅寫的教會書信裡.
一般開始的時候.
都會有一段感恩或祝福的圖文.
帖撒萊加前書有三段感恩圖文.
今天講到的經文.
就是保羅第一段感恩的圖文.
我們開始進入經文的內容.
第一,就算在逆境裡.
我們仍然可以有一個信望愛的生命見證.
在第一章二節.
保羅透露他時時刻刻為帖撒萊加信徒.
在第三節就詳細講明了.
在第三節裡有一個動詞就是.
紀念,不住的紀念.

$^{201}$在保羅未收到提摩太向他匯報的時候.
保羅心裡當時已經紀念著這批信徒.
紀念他們甚麼呢?.
究竟他們有沒有跌倒呢?.
有沒有軟弱呢?.
有沒有被異端迷惑呢?.
這是一個牧者的心腸.
紀念他的羊.
因此當保羅聽到帖撒萊加的信徒.
能夠展現他的生命見證的時候.
在保羅所寫的經卷裡.
有很多地方都講到信望愛.
例如在哥林多前書十三章.
我們經常聽到的.
信望愛之中的最大就是愛.
在羅馬書五章裡面.
保羅所講的盼望是歡歡喜喜的.
在歌羅西書一章裡面.
提到信心的對象就是耶穌基督.
甚至在加拉太書的五章裡面.
講到信心是要連著人愛和果孝.
雖然不同的經卷都有講信望愛這個主題.
但是每一個的重點都不一樣.
而保羅對帖撒萊加的信徒所講的信望愛.
有他的獨特之處.
什麼獨特之處?.
就是他的信望愛.
是對應一個實際的行動.
我們看第三節.
因信心所做的工作.
這裡所講的信心是和工作有關的.
是什麼工作?.
就在第九節.
你們是如何離棄偶像歸向神來服侍那有真有活的神?.
你們是指什麼人?.
就是那些希臘人.
他們不是在猶太的會堂裡面信耶穌的.
那些外邦信徒信了主之後.
就離開了他們先前的信仰.
願意毀身服侍神.

$^{241}$如果用香港的例子.
很多人都相信民間信仰.
家裡有拜祖先.
當他信了耶穌之後.
他就能夠將那個偶像拆了.
這些就是他所講的信心和工作的關係.
信了耶穌.
就能夠將他以前所敬拜的偶像拆毀了.
保留到這班初信主的信徒感恩.
因為他們接受了主耶穌之後.
就以信心去毀身去服侍.
另外因愛心所受的勞苦是什麼呢?.
勞苦的意思就是身體上的勞苦工作.
有辛勞至身體疲乏的含義.
這裡所講的愛心.
原文和三章十二節的彼此相愛的心.
愛眾人的心.
以及第四章九節裡面手足之情.
是同一個字根.
整本書都是互相呼應.
貼紗的信徒.
以愛心去幫助有需要的人.
去到一個地步.
甚至自己身體疲乏的地步.
即是我們服侍不眠不休.
去探病.
可以由早上探到晚上.
面對著躺在病床裡面的病友.
不斷的安慰.
不斷的鼓勵.
不斷的為他代禱.
那種捨身.
忘記了自己由早上已經開始做了.
那種愛心不斷的去實踐.
為的是什麼呢?.
為的是服侍那批有需要的人.
這種透過愛心去服侍.
亦都是對應一章九節裡面.
要服侍那又真又活的神的一種實踐.
信與牧師公去探病人.

$^{281}$不是因為你認識那個病人.
你是很想服侍神.
你透過探望一個你不認識的人.
去服侍神.
這個就是服侍.
那個對方躺在病床的那個人.
你完全不認識.
但你走到病房.
你見到他身淫.
你見到他呼求.
你就走過去為他祈禱.
你告訴他.
已靠耶穌了.
就是以這種的愛心去關心.
幫助你身邊的每一個人.
第三節最後又講到.
因盼望我們的主耶穌所傳的忍耐.
這裡所指的盼望.
與他一章二章三章四章五章所講的.
完完全全一脈相承的.
所指的盼望只有一樣東西.
就是主在內的盼望.
而在這個盼望裡面.
信徒就需要有忍耐和行為和心態.
這個盼望的動力.
使到貼薩萊加信徒的生命裡面產生忍耐.
貼薩萊加信徒忍耐甚麼?.
沒有目的了,保羅都走了.
那他可以有甚麼盼望?.
主在內.
主在內對他們成為一種盼望.
今天我們都是.
我們面對香港有一個疫情.
剛才跟幾姑娘說.
我們可以快些疫情可以回轉.
可以脫口罩.
我們都不知道何時.
是嗎?.
我們作為一個信徒.
我們最終的盼望.

$^{321}$不是盼望何時脫下口罩.
而是盼望主在內.
主在內為我們帶來一個新天新地.
主在內讓我們能夠可以.
得到一個永恆的生命.
是嗎?.
就算地上給了榮華富貴給你.
給多少層樓給你.
戶口裡面有多少個零都好.
這個都是短暫.
作為信徒我們的盼望.
永遠都是主究竟何時來.
甚至我們在人生在世裡面.
總會有時被人冤枉.
是嗎?.
有理說不清.
講都講不到.
誰為我們伸冤.
可能地上的法律幫我們不到.
是嗎?.
我不相信任何一個國家沒有冤案.
多公義多公正的國家.
都總有錯判的情況.
誰為那些受屈者伸冤.
主在內是我們的盼望.
各位新界西語目的弟兄姊妹.
在保羅的土文裡面.
是因為貼誠的信徒能夠.
以信心去服事神.
以愛心去愛身邊的人.
以主在內的盼望.
去忍耐面前的困難.
今日我們在我們服事的歲月裡面.
找不找到這些信望愛的生命見證.
我們一起出隊.
我們能否見到隊友所展現的溫柔.
他所展現的愛心.
他所展現的為神毀身的品格.
我們能否見到.
我們能否讓人在我們生命裡面見到耶穌.

$^{361}$其實我們問一問自己.
我們是否想做一個土神喜悅的信徒.
我相信大部份人都想.
如果我們想的話.
我們要去反思一下.
在我們服事的過程裡面.
在我們生活的過程裡面.
我們有沒有這種信望愛的生命見證呢.
今日那個社會面對很嚴重的撕裂.
人與人之間很缺乏信任.
和一個人溝通.
溝溝下就吵架收場.
有沒有這樣的經驗.
有啊.
在過往這一兩年很多.
敵裡分明.
很多陰謀論.
不相信.
不信任.
當對方指鹿為馬.
扭曲事實的時候.
我們怎樣去彰顯真理呢.
我們心裡面焦急.
當我們面對以上處境的時候.
我們怎樣去回應呢.
我們能否謙卑呢.
我們能否對不同政見的人.
能否仍然心存愛心.
少一份批判.
這個是我們的淑寧生命.
因為真理不在對方的手.
真理也不在自己的手.
真理在三一神的手裡面.
唯有三一神才是擁有真理.
我們沒有權.
我們也不知道真理的全部.
我們只知道真理的一部份.
一切的政見.
一切的學說.
一切的見解.

$^{401}$都在神的真理之下.
神的真理超越一切.
而我們需要聖靈引導我們去認識真理.
這個就是我們的淑寧生命的方向.
當我們面對前面的路.
感到灰心,失望,難過的時候.
我們能否有主在來的盼望.
而產生一份忍耐和堅持.
香港人普遍都性急.
眼光又短視.
退收能夠讓我們安靜一下.
給我們一些反思.
讓我們能夠提升我們淑寧的眼界.
讓我們站高一點.
我們可以望遠一點.
然後我們的目光.
能夠和永恆的主相稱.
是的,真的.
我們要站高一點.
站高一點才能夠望遠一點.
我們永遠都站在這裡的情況下.
我們的眼界是開不了的.
我們要多一點默想操練.
靜心問一問.
究竟自己這樣做的動機是什麼.
我的動機是榮耀神.
因為榮耀自己.
坦白一點去面對.
第二個我們可以在這個.
土文裡面可以反思的就是.
就算在逆境之中.
我們都能夠可以彰顯一種叫毀身的生命.
剛才所說的.
信望愛的生命見證.
但仍然還有一種毀身的生命.
在第四節.
他說神所愛的弟兄.
我知道你們是夢簡選的.
簡選這個字是一個授格名詞.
即是完全由神去做主動的.

$^{441}$我們是被動的.
是神簡選了貼實羅尼加的信徒.
作為教會一員.
貼實羅尼加教會.
怎樣可以成為神所愛夢簡選的呢.
第五節一開始就因為這個字.
帶出貼成信徒能夠被神所愛的一個原因.
第一個原因.
第一個原因就是因為保羅.
將福音帶到貼實羅尼加.
經文說我們的福音傳到你們那裡.
如果沒有保羅將福音傳入貼實羅尼加.
貼實羅尼加的信徒怎能夠信主.
又怎能夠建立教會.
所以傳福音很重要.
當我們進入病房的時候.
那個病人可能是沒有盼望.
因為他不認識神.
他們對死亡可能很恐懼.
因為不知道死後會是怎樣.
當你去探望他的情況之下.
你不是單單給了他心靈一個慰藉.
而是你有機會將福音傳給他.
這個是很重要的.
沒有人傳福音.
誰能夠信耶穌.
沒有保羅進入貼實羅尼加.
貼實羅尼加又如何能夠有信徒.
當一個病人在生命的最後階段裡面.
他固然很想見到他的親戚朋友.
他固然很想見到他的至親.
但是你用屬靈的眼光看一看.
想一想.
如果他和耶穌擦身而過.
他有什麼後果?.
如果你帶著耶穌.
去讓那個病人和耶穌相遇.
他和你就有一個永恆的福樂.
這個何其重要?弟兄姊妹.
我們知道保羅有三次宣教旅程.

$^{481}$貼實羅尼加教會在保羅第二次宣教旅程當中.
在菲律賓建立了教會之後.
然後再去到貼實羅尼加.
然後建立的.
保羅在《哥林多後書》十一章裡面.
曾經說過他對外邦人傳福音的一個經歷.
我們聽聽他怎麼說.
他說「我被他們多受勞苦,多下監獄.
受鞭打是過重的.
誣告死是屢次都有的.
被猶太人鞭打了五次.
每次四十砍去一.
被棍打了三次.
被石頭打了一次.
遇著損壞三次.
一晝一夜在深海裡.
又旅行遠路.
遭江河的危險.
盜賊的危險.
同族的危險.
外邦人的危險.
城裡的危險.
曠野的危險.
海中的危險.
賈弟兄的危險.
這麼多的危險.
受勞碌,受困苦.
多次不得睡.
又飢又渴.
多次不得食.
受寒冷赤身怒體」.
一般人聽到這些經歷.
會怎樣?.
打退堂鼓吧.
還早?.
是不是?.
回家好過了.
是不是?.
去得不得救關我什麼事?.
我得救就行了.

$^{521}$保羅不是這樣想的.
保羅說自己蒙召.
毀身向外邦人傳福音.
在整個過程裡面.
保羅付出了很大很大的代價.
正正是由於保羅願意毀身.
貼實來家的信徒.
才有機會聽到福音.
繼而信著.
我們願不願意為福音而受苦?.
我們願不願意毀身傳福音.
縱使受苦.
為福音緣故而受苦.
我都甘心樂意.
接承教會能夠被神所愛.
除了剛才那個原因之外.
還有第二個原因.
固然保羅有願意傳福音的毀身.
但是這裡說.
他說福音本身是大能.
從保羅.
從保羅從聖靈而來的信心.
同樣在第五節裡面說到.
不僅在言語也在能力.
也在聖靈和充足的確信.
新譯本也更加準確.
他說不僅是藉著言語.
也是藉著能力.
藉著聖靈和堅定的信念.
信不信耶穌.
聖靈要很大的幫助.
堅定的信念.
不是說保羅和福音的信心.
而是指保羅對聖靈的能力和信心.
你進去病房見到病人.
你說到連樹上的小鳥也掉下來.
哄到下來.
但是聖靈不參與也沒有用.
你進去病房.
你感受到聖靈與你同在.

$^{561}$就算你說的話你覺得很笨拙也好.
那個人也會相信耶穌.
所以為什麼我們進入病房的時候.
我們要求聖靈與我們同在.
我們不是憑血氣去傳福音.
我們憑著聖靈的大能去傳福音.
所以這個意思就是說.
福音的傳遞.
不僅是透過傳福音者所說的言語.
不僅是重要的.
但不是唯一.
都是透過福音本身的能力.
福音本是神的大能.
而且傳福音帶著從聖靈而來的信心.
我們每一次和別人宣講福音的時候.
你感不感受到聖靈與你同在.
因為你只覺得我可以的.
我可以的.
我辯論組我冠軍的.
我去街市買菜講價我就第一名.
是不是靠著這些去傳福音.
不是.
靠著神靈.
靠著聖靈與你同在.
確信那個聽到者是蒙神請求.
所以保羅來到帖撒羅來加傳福音.
是帶著聖靈而來的信心.
帖撒羅來加教會能夠被神所愛的第三個原因就是.
保羅帶著一個無愧的生命來做見證.
第五節下半節.
他說你們知道我們在你們那裡為你們的緣故是怎樣為人.
經文本身這裡沒有講到他怎樣為人的內容.
不過從經文的上下文就可以找到答案.
第二章的第十節保羅說.
我們向你們信主的人是何等聖潔公義.
無可指責.
有你們在那裡做見證.
也有神做見證.
所以我們知道保羅在帖撒羅來加的行事為人是什麼.
聖潔的,公義的,無可指責的.

$^{601}$而這些行為都是帖撒羅來加的信徒可以見證的.
可以這樣說.
這個是無偽的生命見證.
帖撒羅來加教會之所以被神所愛.
是因為保羅從沒有計較付出.
沒有從聖靈而來的信心.
而且實踐著一個無愧的生命見證.
將福音帶來帖承.
才能夠令到帖承教會有如此美好的結果.
保羅這一種忠於照明的生命.
我們豈會不羨慕呢?.
我們羨慕嗎?.
我們羨慕.
我們願不願意效法呢?.
我們都願意效法.
《新界西演武》的侍工是不是被神所使用的?.
是.
是不是合乎神心意?.
是.
所以每一個參與的都有責任.
不只是那個阿頭才有責任.
不是紀姑娘,陳牧師才有責任.
每一個都有責任.
每一個只要活出信仰.
侍工就更加被神所使用.
如果每個信徒都不願意毀身的話.
侍工不可能復興.
也不可能被神大大使用.
或者我們可以這樣說.
我們沒有保羅那麼能捱苦.
我被人鞭打兩下就死了.
我們不行了.
浸我一天,皮都皺了.
我們不行的.
我們沒有他那麼強的心智.
但我們只要有一個願意毀身的心.
你在神面前說,主我願意.
你加能賜力給我.
每一個人都可以用你自己.
甘心情願的方法去服侍神.

$^{641}$我們不需要像保羅那樣.
當然保羅是值得我們去讚賞.
但我們只是小薄餅.
我們每一個都是小薄餅.
但我們只要願意.
我們願意被神使用就足夠了.
只要我們願意,神就會大大使用我們.
最後一點.
就算在逆境裡.
我們都能夠以效法主作為我們的使命.
在第六節.
並且在大能之中.
蒙了聖靈所賜的喜樂.
領受真道就效法我們也效法了主.
新漢語將經文的次序改了一點.
它說你們效法了我們,又效法了主.
在大患難裡.
懷著聖靈所賜的喜樂領受真道.
這個譯法強調了效法我們.
又效法了主的訊息.
他們如何效法保羅和主耶穌呢?.
就是他們在大能中.
帶著聖靈的喜樂接受了真道.
領受真道即是接受了上帝的道.
即是福音.
貼撒羅尼加的信徒.
在兩種狀態之下去接受福音.
第一種接受福音的狀態是大能.
大能中.
這裡的大能是甚麼呢?.
這裡所說的大能.
這裡所說的大能未必是指基督徒在信仰上.
高受到羅馬政府的大迫害這麼嚴重.
事實上羅馬皇帝尼祿.
大規模的迫害基督徒的時間.
大約在公元64年左右.
而保羅在貼撒羅尼加傳福音的時候.
就大概在公元53年左右.
距離真正的大迫害相差近十年.
即是未是那個時候.

$^{681}$我們知道保羅在貼撒羅尼加傳福音的時候.
就說有些猶太人和希臘人信了耶穌.
包括了婦女.
這些希臘的新信徒轉信了另一個陌生.
而且備受社會懷疑的宗教的話.
如果他們參與.
即是拒絕參與以往的祭祀或拜鬼活動.
即是等於舉個例子.
一個圍村人現在信了耶穌.
然後圍村一個太公婚豬肉.
拜神裝香.
很大壓力.
爸爸叫你裝香.
族長叫你叩頭.
你又說我信了耶穌.
不拜行不行呢.
這樣就很大壓力.
這個大能.
可能就是這樣就得罪了人.
另外好像剛剛信主的爺孫.
一信主沒多久就被人拉去警署.
要向地方官保證保羅不是貼撒羅尼加在這裡生事.
這種都是大能.
是不是.
你傳福音然後被人抓.
你害不害怕.
害怕一點.
警察問你是不是傳義端.
你是不是邪教.
你都要慢慢解釋.
你覺得這種是不是逼迫.
同樣有些傳統猶太教背景的信徒.
他們的家族成員認為.
你不再維持祖先的傳統.
你終於和這個家庭不關心不負責任.
令猶太信徒左右做人.
所以那個大能中的意思可能是指.
信徒日常生活所遇到的衝突.
他們被排擠甚至被陷害.
我相信我估計大家都沒有這種情況.

$^{721}$但對一個初信主的人.
他在吃飯祈禱都有壓力.
是不是.
剛剛信耶穌.
然後牧師說你吃飯祈禱.
他周圍的同事都不信耶穌.
然後他吃飯要祈禱.
有壓力嗎.
有壓力的.
他們在這種情況仍然信耶穌.
你說對於我們新界西演目.
在這裡有什麼應用.
其實你想一想.
那躺在那裡的病人.
他們是否在大能中.
他們在大能中.
治得好固然可以離開醫院.
有機會治不好.
剛才和海美談.
她說青山醫院有些病人.
一輩子住在那裡.
是不是在大能中.
在大能中怎能有盼望.
靠誰.
靠我們.
是我們才能有這樣的福分.
可以進入醫院探望他們.
事實上真的.
不是阿珠阿九可以進入醫院探望他們.
經過訓練.
是不是.
要經過蒙照.
要經過毀身.
神給了你恩賜.
給了你身份.
給了你一個遺憤.
然後你才能進入醫院探望那個人.
他和你三不識七.
醫院都怕你去賣傳銷.
你病死了不如吃這個藥.

$^{761}$起死回生.
是不是.
第二種接受福音的狀態.
是蒙了聖靈所賜的喜樂之後.
這段經文除了可以理解為.
聖靈所帶來的喜樂之外.
亦都可以解釋為.
聖靈與信徒同在成為喜樂之源.
所以第六節可以這樣說.
即是這樣翻譯.
他說貼薩雷加的信徒成為了.
使徒和主效法者.
是在於他們在患難裡面.
帶著聖靈的喜樂接受了福音.
患難裡面都能夠可以接受福音.
在患難裡面.
因為福音本是神的大能.
而第七第八節.
他說從此你們作了馬其頓和阿閣亞.
所有信主的人的榜樣.
因為主的道已經從你們那裡傳播出去.
你們向神的信心.
不只在馬其頓和阿閣亞.
就是在各處也都傳開了.
所以不用我們說甚麼話.
這兩節是對貼薩雷加信徒.
在第六節的表現.
一種肯定和稱讚.
當時他們身邊沒有牧者.
保羅西拉提摩泰都不在他們身邊.
但是貼薩雷加信徒能夠成為.
馬其頓和阿閣亞信主的人的榜樣.
是因為他們在生命上願意效法主.
如有一日.
有人說.
新界西爾牧師公主真是厲害.
真是值得我們效法.
他們的信往愛的精神.
真是淹死人.
這就是他們情況差不多.

$^{801}$貼薩雷加的人.
可以成為馬其頓和阿閣亞信主的人的榜樣.
香港島那邊也有些遠目部.
香港島遠目部.
有一日找幾姑娘說.
我們收到消息你們很厲害.
我們想來考察一下.
聽聽你們隊員的心路歷程.
讓我們可以借鏡學習.
就是這樣.
想不想.
生命能夠影響生命.
貼薩雷加的信徒能夠為馬其頓和阿閣亞的人.
展現了另一種生命力.
這種生命力讓人讚歎.
所以保羅說.
所以不用我們說什麼.
因為信徒的生命比傳道者的言語.
更加能夠展現信仰的真實.
貼薩雷加的信徒.
這班初熟的教子.
他們雖然沒有仕途.
沒有傳道人的名分.
但他們同樣都是福音的使者.
同樣新界西遠目的義工.
你們都不是傳道人.
你們都不是牧師.
但同樣你們都是信徒.
你們都是牧師.
但同樣你們都是福音的使者.
同樣你們都是能夠讓人生命改變的人.
因為他們不單單自己改變了.
你們的生命改變了.
也透過你們的生命見證.
可以改變很多的人.
我相信就是這麼多弟兄姊妹.
過來新界西遠目的義工.
與新界西遠目的義工的目的.
我們領受了恩典.
我們很想別人都領受這個恩典.

$^{841}$我們的恩典讓我們的生命改變.
我們都很想那些人能夠領受這個恩典.
讓他們的生命都改變.
所以弟兄姊妹.
在三個逆境裡面展現出來的信仰.
可以歸納為甚麼呢.
第一種能夠體驗信望愛的生命.
第二種就是能夠帶著照命.
願意毀身的生命.
第三個是甚麼呢.
能夠效法主.
在大患難裡面仍然有喜樂.
並且能夠成為別人生命的榜樣.
當每一個信徒.
能夠展現到任何上述一種的生命特質.
其實都可以被神所大大使用.
不過當然我們明白.
每一種生命的特質.
都需要我們持續不斷的屬靈操練.
然後才能夠結出果實.
所以為甚麼今天要搞退休.
今天就是讓我們的屬靈生命.
能夠有一個操練.
透過安靜,默想,分享.
然後讓我們的屬靈生命能夠提升.
我們願不願意花多點時間.
做這種屬靈操練.
可能你跟幾個人說.
一年一次太少了.
這種增強屬靈生命的操練方法.
可能真的要多一點.
如果你反映的話.
我想幾個人都會想一下.
我對帖撒萊加的背景.
我想再分享多一點.
在《仕途行傳》的記載裡面.
沒有說到保羅究竟在帖撒萊加.
實際停留了多久.
最準確的資料.
只是記載了他連續三個安息日.

$^{881}$向外邦人進行傳福音的工作.
之後就說到一班不信的猶太人.
召集了一些壞分子.
在城中引起騷動.
他們衝入耶穌的家守護保羅.
所以保羅要立刻離開帖撒萊加.
但是根據學者的估計.
相信保羅在帖撒萊加.
牧養信徒的日子不超過半年.
即是留下了半年的時間.
即是說帖撒萊加的信徒.
由開始信耶穌.
再接受保羅的培訓.
最多都是半年左右.
接著就沒有了牧者牧羊.
在這樣的環境下.
帖撒萊加的信徒竟然能夠表現上述的生命特質.
你會不會有興趣再問下去.
究竟保羅是怎樣培訓信徒的?.
這麼厲害?.
半年而已.
半年就走了.
然後他們可以成為亞該亞信徒的榜樣.
半年而已.
他們的生命就像脫胎換骨.
可以有順望愛的生命見證.
作為牧者我就很想知道.
你教教我吧.
然後在紅欣那裡翻兜.
半年翻一翻.
都不知道多好.
究竟保羅的培訓內容的核心是什麼呢?.
有什麼核心價值.
可以推動這些信徒.
有這麼強的生命特質呢?.
在二章十二至十三節保羅這樣說.
他說要叫你們行事.
對得起那召你們進他國.
得他榮耀的神為恥.
因你們聽見我們所傳神的道就領受了.

$^{921}$不是以人的道乃以為是神的道.
這個道實在是神的.
並且運行在你們信主的人心裡.
保羅想說什麼呢?.
他說行事為人是與召命有關.
這個召命是由神所發出的.
在二弗所書四章一節.
保羅也曾經說過.
既然蒙召行事為人.
就與當蒙召的恩相稱.
如果用整個新約來說.
在彼得前書二章二十一節.
他更加說到你們蒙召原是為此.
因基督也為你們受過苦.
給你們留下榜樣.
叫你們跟隨他的腳蹤行.
如果耶穌基督已經成為.
貼實羅尼加信徒的榜樣.
而信徒想要效法主的話.
他們能夠在苦難中有喜樂.
能夠在召命之中體現.
信望愛的行為.
那就一點也不奇怪.
如果這個訊息就是我們要找的答案的話.
只要我們新界西語目的每一個義工.
我們能夠以效法主的心態.
不是效法人.
不是效法牧者.
不是效法一個屬靈領導.
是效法主耶穌.
然後去回應神給我們的呼召.
這樣我們就能夠在逆境裡面.
去展現出這三種生命的特質.
弟兄姊妹我們實在需要去問自己.
今天我們為誰而活呢?.
我們為自己?當然了.
我們不是為了生活?.
我們為豐衣足食?.
因為我們是為耶穌而活.
當我們進入病房去探訪病人的時候.

$^{961}$當我們向他們發出一個問候安慰的同時.
能不能夠展現出這種信望愛的屬靈特質.
如果我們的生命能夠散發出基督的愛.
病患者就可以在我們生命裡面見到主耶穌.
我們的身份不單單是基督徒.
我們的身份更加是主的使者.
代表著一個神聖的使命.
神將我們分別出來.
讓我們在新界西演武部事奉.
不是要我們安安處處開開心心.
而是要為新界西演武部.
揭出福音的果子.
讓這些福音的果子.
能夠好像一棵樹一樣.
屢屢的果子.
然後將榮耀歸還給神.
我們願不願意去毀身回應.
這個主耶穌對我們每一個人的呼召呢.
我們一起去祈禱.
因為是你將我們從罪裡面拯救出來.
你將我們每一個人的生命都改變.
你讓我們能夠成為你的兒女.
讓我們能夠進入神的角度.
我們有君尊的祭司.
我們有身份無比的尊貴.
是你 是主.
你將我們救拔的.
是你將我們從糞堆中拯救過來的.
今天我們成為何等樣式的人.
我們需要去回應主你這個給我們的恩典.
這個不是一個廉價的恩典.
這個是貴價的恩典.
這個是主耶穌用你的生命寶血.
將我們買熟過來的恩典.
所以主求你叫我們好好地去回應.
好好地去回應主你對我們的呼召.
好叫我們明白.
我們今天我們作為一個信徒.
我們作為你的福音使者.
我們的生命如何能夠被你所使用.

$^{1001}$以至你的名被高舉.
以至更多人因著我們的生命而得福.
讓更多的人能夠成為天國的一員.
我們同享永恆盼望的福祿.
願主賜福給新界西每一個弟兄姊妹.
我們的祈禱.
奉教主耶穌基督得勝名至其球.
阿們.
好 願主祝福大家.
\newpage



\section{馬可福音 2:18-22}
\label{sec:0Rk0UlP9K50}
\textbf{2021.05.23《新皮袋》馬可福音 2\_18-22( 和修版)  吳景霖傳道}
\newline
\newline
連結: \href{https://youtube.com/watch?v=0Rk0UlP9K50}{\texttt{https://youtube.com/watch?v=0Rk0UlP9K50}} ~~~~ 語音日期: 2021-05-25
\newline
\newline
\hyperref[sec:HojbxuWCThI]{\small{< < < PREV SERMON < < <}}
~
\hyperref[sec:index_chronic]{\small{[返順時目]}}
~
\hyperref[sec:index_scriptual]{\small{[返順卷目]}}
~
\hyperref[sec:JhzI3zqDYAI]{\small{> > > NEXT SERMON > > >}}
\newline
\newline
馬可福音 2:18-22
\newline
\begin{longtable}{cl}
\hline
\hline
章節 & 經文 (和合本修訂版)\\
\hline
2:18 & \begin{tabularx}{0.7\textwidth}{X} 那時，約翰的門徒和法利賽人都禁食。他們來問耶穌說：「約翰的門徒和法利賽人的門徒禁食，你的門徒卻不禁食，這是為甚麼呢？」 \end{tabularx} \\ \\ \relax
2:19 & \begin{tabularx}{0.7\textwidth}{X} 耶穌對他們說：「新郎和賓客在一起的時候，賓客怎麼能禁食呢？只要新郎和他們在一起，他們不能禁食。 \end{tabularx} \\ \\ \relax
2:20 & \begin{tabularx}{0.7\textwidth}{X} 但日子將到，新郎要被帶走，那日他們就要禁食了。 \end{tabularx} \\ \\ \relax
2:21 & \begin{tabularx}{0.7\textwidth}{X} 「沒有人把新布縫在舊衣服上，若是這樣，所補上的新布會撕破舊衣服，裂口就更大了。 \end{tabularx} \\ \\ \relax
2:22 & \begin{tabularx}{0.7\textwidth}{X} 也沒有人把新酒裝在舊皮袋裡，若是這樣，酒會脹破皮袋，酒和皮袋都糟蹋了。相反地，新酒要裝在新皮袋裡。」 \end{tabularx} \\ \\
[1ex]
\hline
\hline
\end{longtable}
$^{1}$好,OK.
好,大型節目,祝你們平安.
清不清楚,後面聽到我這個音量.
OK,好,大家好,我是吳景霖傳道.
叫我不傳道,但其實我是會傳道的.
所以大家不用擔心.
也少來到紅人堂跟大家講道.
可能大家戴著口罩不認得我.
也很高興,特別是看到這麼多人能夠回來.
坐在你旁邊的人,你看一看他.
他還在你旁邊,他還沒被人捉去隔離.
其實也是值得感恩的一件事.
能夠回來,不是必然.
今天講的經文,其實這段經文已經很熟了.
我們很熟悉,講新舊的問題.
通常人們會新年去講.
因為新一年就講新事物.
但是新舊這件事,我想問大家.
是否固定的?有沒有一些是歷久常新,永遠都新的?.
未必的,新舊有時候是一個相對.
或者是一個流動的概念.
為什麼這樣說?舉個例子,譬如手提電話.
你出了一部手提電話,用得好好的.
我的iPhone11也很新.
但是出了iPhone12之後.
它原來只不過是一個舊電話.
原來所謂新的東西,新只有一年.
那個旗艦機,原來下一年已經不是旗艦了.
現在的事情很有趣.
但是,是不是有些東西永遠都是舊的呢?.
舊了就永遠都是舊的,是不是?.
有時候未必的,如果你發現有些考古的新發現.
挖到一個古墓,挖到一個很完整的恐龍花石.
挖到一些古代的文獻,譬如死海古卷.
我們找到一卷新的東西.
是不是真的新的?幾千年了,幾萬年了.
骨頭幾億年了.
但是對我們來說就是新的了.
因為是一些新的發現,是不是?.
有些東西可能是一些舊的知識.

$^{41}$不過你以前從來沒有聽過.
於是當你聽到,你學到的時候.
對你來說就是一種新知識.
新和舊,你會發現是一個流動的概念.
是一個相對的概念.
沒有一個固定的,沒有一個絕對的在裡面.
所以今天說的經文.
耶穌對所謂新和舊的情況有一點批判和教導.
我們首先看18到20節.
第18節說:那時約翰的門徒和法利賽人都禁食.
他們來問耶穌說:.
約翰的門徒和法利賽人的門徒禁食.
你的門徒卻不禁食,這是為什麼呢?.
這群「他們」指什麼人呢?.
就是指約翰的門徒和法利賽人.
法利賽人我想大家都知道.
他和耶穌的關係不是很好.
甚至是到了這段經文的時候.
他們已經是有心來找耶穌麻煩.
他當然不是來誠心問問題.
他是來找耶穌麻煩的.
但是約翰的門徒就不同了.
照道理他們應該對耶穌有些認識.
記不記得在耶穌的門徒裡面.
其實有兩個原本是約翰的門徒.
其中一個就是彼得的弟弟安德烈.
所以照道理約翰的門徒對耶穌是有一定的認識.
他們應該知道耶穌是什麼人.
他不應該來找他麻煩.
他們的目的是不同.
他們問了什麼問題.
約翰的門徒和法利賽人的門徒禁食.
你的門徒卻不禁食,這是為什麼呢?.
大家有沒有操練過禁食?.
有禁食吧?都有試過禁食.
禁食是為了什麼?.
為了什麼呢?.
佩華禁食是為了什麼?.
認得你,所以就唯有問你.
(為了什麼禮儀?).

$^{81}$為了一些重要的節日.
或者是一些重要的事情.
例如有些很大的事情.
例如身邊有些朋友的肢體患了重病.
我們為他禁食禱告求上帝的醫治來到.
或者是有一些情況.
就像當時的文化.
禁食是代表認罪悔改.
痛悔.
於是他禁食.
這是一個很普遍的宗教行為.
法利賽人一個星期禁食兩次.
所以禁食是很普遍的.
他們為了什麼要認罪悔改呢?.
那時候無緣無故認罪悔改什麼呢?.
其實是在說以色列和猶太.
當時已經亡國了.
那時候耶穌的時代.
他們已經在被統治的情況下.
他們亡國.
猶太人相信.
是因為自己的祖先得罪上帝.
他們得罪上帝.
所以上帝懲罰他們.
於是以色列和猶太就亡國了.
這是他們很明白的一件事.
他們一直都過著一種痛悔的生活.
他們希望上帝的拯救會來到.
他們希望彌賽亞會來拯救.
所以他們經常保持著一種很敬虔.
或者他們很持守律法.
他們很認真在信仰裡面.
約翰的門徒和法利策人.
雖然是兩個不同的群體.
但是他們同樣是猶太人.
所以他們有實踐禁食這個行為.
是很正常的.
不會有什麼意見分歧.
或者不會有什麼吵架的問題出現.
就算是現代我們不同宗派也好.

$^{121}$你也不會說有太大的分歧.
對於祈禱禁食這些東西.
問題就是為什麼耶穌那幫人反傳統呢?.
不禁食.
一些做慣做熟的事情.
全部猶太人都是禁食的.
就是你這個耶穌帶著那幫人不禁食.
對於他們來說他們會擔心的.
得罪上帝那怎麼辦?.
他們不知道耶穌是彌賽亞.
你要知道.
你這個人帶著那幫人不禁食.
得罪上帝那怎麼辦呢?.
法利賽人來找耶穌麻煩.
但是約翰的門徒可能真的要來搞清楚這件事.
因為不搞清楚可能會引起很大的問題.
就是一個信仰的危機.
好了,那耶穌就解釋了.
第十九節.
耶穌對他們說新郎和賓客在一起的時候.
賓客怎能禁食呢?.
只要新郎和他們在一起.
他們不能禁食.
但日子將到新郎要被帶走.
那一天他們就要禁食了.
這個新郎是指誰呢?.
誰是新郎呢?.
是耶穌吧.
耶穌指自己.
那為什麼說自己是新郎呢?.
因為是男人.
不說自己是新娘.
新郎是娶老婆回去的人.
是不是?新郎.
好了,其實他不是隨口說的.
我們要看《河西亞書》.
就是舊約裡面的《河西亞書》.
河西亞的時代.
大概是以色列國開始走下坡.
那時候被列強圍住.

$^{161}$被亞述和埃及圍住.
於是就和亞述和埃及簽訂了和約.
同盟.
但問題是.
這些同盟引致一些外邦的信仰.
入侵了以色列.
以色列人開始.
他們的信仰走下坡.
他們不再只是相信耶和華.
而是開始有一些混雜的信仰出現.
信仰走下坡.
於是乎.
這個《河西亞書》裡面.
以色列被上帝形容他們是.
一個淫亂的婦人.
就是不忠.
不忠.
他們對誰不忠呢?.
對上帝不忠.
以色列人對上帝不忠.
於是他們開始拜一些假神偶像.
上帝就追討他們的罪.
之後以色列和猶太亡國.
其實就是因為上帝的懲罰.
淋到.
但上帝是否就這樣算數呢?.
懲罰完就算數.
不是的.
上帝也有預言會拯救他們.
是會拯救的.
雖然有懲罰.
但會拯救.
在《河西亞書》二章十九到二十節這樣說.
我必聘你永遠歸我為妻.
聘就是娶回來的意思.
以公義公平慈愛憐憫.
聘你歸我.
也以信實聘你歸我.
你就必認識耶和華.
聘就是娶老婆的意思.

$^{201}$我們結婚有聘禮.
娶他回來.
雖然以色列人對上帝不忠.
但上帝會以公義公平慈愛憐憫.
娶他們回來.
因此耶穌說自己是新郎.
不是隨口說的.
有舊約的根據支持.
他形容自己就是新郎.
來娶你那群以色列人回到上帝那裡.
他代表了上帝.
剛才說過.
猶太人禁食.
是向神代表認罪痛悔.
希望彌賽亞尼拯救.
彌賽亞也來了.
來了的時候還敢吃什麼.
新郎也來了.
還等什麼呢.
這個時候應該是歡喜快樂的時間.
要珍惜的時間.
要回到祂身邊的時間.
所以絕對不是禁食的時間.
新郎上在身邊的時候.
是不需要禁食.
不過同時間耶穌也預言.
日子將到新郎要被帶走.
那天他們就要禁食.
通常看聖經.
特別是看科幻書的時候.
說的日子指向耶穌上十字架的日子.
就是說耶穌其實已經預言.
祂會被帶走.
當祂被帶走的時候.
你就要禁食了.
因為你需要重新再去等候.
等候耶穌基督的再臨.
耶穌被帶走.
釘死後來復活之後升天.
我們相信耶穌基督會再來.

$^{241}$剛才連使徒信經也說了.
會再來.
所以耶穌基督會再來.
這個時候禁食就變得合宜了.
所以大家剛才點頭說有禁食.
是對的.
你不用擔心.
禁食是合宜的.
其實耶穌一方面是在解答問題.
另一方面是在表達自己的身份.
不過祂是意猶未盡.
祂只是說完也不夠.
還要加一個比喻在後面.
再多說幾句.
21節.
沒有人把新布縫在舊衣服上.
若是這樣所補上的新布.
會撕破舊衣服.
裂口就更大了.
也沒有人把新酒裝在舊皮袋裡.
若是這樣酒會橡破皮袋.
酒和皮袋都糟蹋了.
相反地新酒要裝在新皮袋裡.
其實這個道理好像很顯淺.
也不是十分複雜.
他說沒有人會把新布補在舊衣服上.
為什麼呢?.
這個時代其實較少人補衣服.
通常如果爛一點也好.
爛得大就扔掉.
或者有些人特意弄破.
牛仔褲穿了一個洞.
我阿麻試過幫我補了牛仔褲的洞.
不知道是哭還是笑.
補衣服這個時候較少人去做.
因為衣服也不是太貴.
但是以前的世代物質沒有這麼豐富.
衣服是很貴的一件事.
通常人們可能只有一件起兩件止.
你看見耶穌上十字架的時候.

$^{281}$那些兵兵搶他的衣服.
要抽籤爭他的衣服.
所以當時衣服其實是貴的.
所以要修修補補.
補衣服是很平常的.
但其實衣服穿久了會變得霉.
布質也會變得霉.
我也試過有些布放了很久.
拿上手輕輕一撕.
整個布就爛開了.
越舊的布就越霉.
如果你補新的布上去.
新布硬朗.
舊布就變得霉.
於是補衣服的過程可能一拉.
就會將舊布拉爛.
還有縮水.
原來舊衣服縮水不太多.
新衣服就比較容易縮水.
補完上去一洗.
新布縮水比較多就會拉爛舊布.
所以當時的人都會很明白.
新布是不可以補在舊衣服上.
同樣道理.
酒也不可以放在舊皮袋裡面.
因為酒是會發酵的.
越新的酒發酵的時候.
釋放的氣體越多.
於是在舊皮袋裡面放新酒.
就會把酒釋放.
因為舊皮袋已經脆弱了.
新酒釋放的氣體多.
就會把酒釋放.
其實是同一個道理.
耶穌想說什麼呢?.
祂是否在教什麼生活冷知識呢?.
就像你看電視.
今天先說冷知識.
教你如何保存酒和衣服.
當然不是.

$^{321}$一定有屬靈教訓在背後.
一定有意義在背後.
今天說的是什麼呢?.
我們最重要明白的是.
其實耶穌暗示什麼是新.
什麼是舊.
什麼是新 什麼是舊.
不需要禁食.
是一種新的情況.
不過你要明白.
這只是現象.
你眼睛看到的叫做現象.
你見到祂沒有禁食.
但是一個真正新.
是在說背後的屬靈現象.
就是原來彌賽亞已經道成肉身.
去到人的中間.
這才是真正的新.
以前未見過 未出現過.
上帝拿了肉身來人中間.
這件事未見過.
這才是新.
耶穌基督的來到.
象徵著一個新時代的降臨.
祂來更新這個世界.
祂要改變這群人的思維.
要在舊的制度裡釋放出來.
我們要明白.
舊制度 說得好像舊東西不好.
是不是真的不好呢?.
耶穌是不是來廢掉舊的制度?.
不是的.
你記不記得勝經怎麼說?.
耶穌說.
「不要以為我來是要廢掉律法,我才知,.
我不是要廢掉,而是要成全.」.
耶穌來不是說舊的制度有問題.
不是說律法有問題.
有問題的是什麼?.
你怎麼守.

$^{361}$你怎麼用.
你怎麼使用那件事.
很多時候我們要明白.
其實一些法例不是有問題.
是執法的怎麼做.
怎麼用才是一個問題.
馬太福音15章1至6節.
我沒有給經文.
所以聽我讀.
「那時有法利卓人和文士.
從耶路撒冷來見耶穌.
說:你的門徒為什麼違反古人的傳統?.
因為他們吃飯時不洗手.」.
現在吃飯不洗手是很大的事情.
疫情大家記得要洗手.
「耶穌回答他們:你們為什麼因你們的傳統.
而違反神的誡命?.
神說:當孝敬父母.
又說:咒罵父母的必須處死.
你們賭說:無論誰對父母說.
我所當貢獻的已作了奉獻.
就可以不孝敬他的父親.
這就是你們藉著傳統廢了神的話.」.
這裡說到那些法利卓人.
竟然說如果你給父母那份家用.
如果拿了來奉獻.
你就不需要給家用了.
這件事怎麼說都不通.
在現代怎麼說都不通.
奉獻給神這件事有沒有問題?.
奉獻給神沒有問題.
但是那些法利卓人為了自肥.
現在潮流流行說自肥.
便宜了自己.
奉獻去聖殿.
最終都是進了宗教領袖的袋子.
但是父母是你的.
他當然不理你那麼多.
他為了自肥.
於是要求將貢獻父母那份奉獻了他.

$^{401}$其實就是美其名.
要進自己袋子.
法例本身沒有問題.
奉獻本身沒有問題.
有問題是法利卓人.
他怎麼去守律法.
他怎麼用律法.
這個才是最大問題.
同一個道理.
禁食本身是沒有問題.
但是怎麼禁法.
怎麼禁法就是一個問題.
你記不記得.
曾經耶穌說過一個比喻.
法利卓人和瑞利.
瑞利低頭雷空遁足.
去求神可憐他.
但是法利卓人是怎樣的.
他自言自語.
其實自己在說.
說給別人聽.
我不像別人.
其實是說瑞利.
勒索不義奸淫.
也不像這個瑞利.
我每週禁食兩次.
凡我所得的都獻上十分之一.
他每週禁食兩次.
你有沒有禁食兩次.
他比較屬靈還是你比較屬靈.
他每週禁食兩次.
但問題是.
他只是表現自己.
一種宗教表現.
一種宗教的行為.
我做出來好像我特別屬靈一點.
這個其實是自我滿足.
我滿足自己.
或者讓人看到覺得我很好.
其實這個也是滿足自己的虛榮.

$^{441}$神又怎會喜悅呢.
你不要說每週禁食兩次.
你每天禁食兩次都沒有用.
所以其實耶穌是處理他們的問題.
第一.
約翰和那些門徒和法利賽人.
他們不明白.
耶穌就是那個彌賽亞.
於是他們就不懂得去尋求耶穌.
他們唯有用舊的方法.
就是繼續禁食.
繼續去認罪悔改.
但其實你向耶穌認罪悔改.
你不是禁食.
你的對象已經來到你前面.
第二就是.
其實法利賽人的心態是有問題的.
他們用錯誤的心態.
去實踐禁食.
其實就是舊衣和舊皮帶.
舊衣和舊皮帶.
他們的舊的宗教行為.
是不可以勝在耶穌帶來的更新.
我要講清楚.
有問題的不是制度和律法.
而是心態和怎樣去遵守.
是心態和你怎樣去守.
怎樣去用.
時至今日.
對我們來說.
其實這是一個很大的提醒.
我們會想我們是基督徒.
有沒有人不是基督徒.
應該都是.
如果不是也好.
你可以在這裡聽聽福音.
我們是基督徒.
於是我們相信耶穌基督.
即是彌賽亞.
我們很容易認為.

$^{481}$在這段經文裡面.
我們就是在信新的東西.
猶太教那些才是信舊的東西.
我們就是在信新的東西.
如果是這樣的話.
其實我們這段經文講到這裡就完了.
我不用繼續講下去.
是不是.
很容易我們就代入了這個位置.
但問題是正如我一開始所講.
新和舊的概念不固定.
如果舊皮袋.
就是用錯誤的心態.
去持守舊有的律法.
就算我們信了耶穌.
你也有可能落入這種光景.
你也可能是在用舊皮袋的思考模式.
用錯誤的心態去遵守聖經的教訓.
不知道大家會不會覺得很大衝擊.
或者很驚訝.
信了耶穌.
不是信了獨一的真神嗎.
已經不用變了嗎.
我們的信仰會不會變呢.
其實我們的信仰不斷地在變.
很坦白地說.
不知道大家有沒有留意到.
其實今年是馬丁路德的宗教改革504年.
早幾年慶祝過宗教改革500年.
宗教改革是什麼呢.
知不知道什麼是宗教改革.
有些人不知道.
你知不知道你信的基督教只有500年歷史.
500年4年.
由耶穌時代到初期教會.
一直是羅馬天主教.
一直到馬丁路德的時代.
即是500年前.
因為教會的腐敗.
所以馬丁路德貼了95條論綱.

$^{521}$即是有95條覺得教會有問題的東西.
扔在教堂門口.
於是引爆了宗教改革的浪潮.
從此基督教和天主教分了價.
所以你要知道我們的信仰是變了.
我們的信仰是變了.
時至今日你會知道天主教和基督教是不同的.
但500年前其實是一家的.
有很多弟兄姊妹其實是不知道這件事的.
你可能以為.
很有趣的想法就是耶穌保羅到我們教會.
中間的歷史對我們來說完全沒有真空.
但這些東西是要知道的.
所以我要提醒大家.
特別是一些信徒內的肢體.
其實我們是要反省我們的信仰.
我們的信仰是有東西要變.
我們的信仰是有東西要更新.
不是永遠守著一套舊皮袋.
我舉一個例子.
可能大家已經過渡了.
但我信耶穌信了二十多年的時候.
我都看著改變的發生.
那地崇拜裡面很習慣有音樂有詩歌.
我很習慣這件事.
其實二十多年前信耶穌的時候.
有些教會還是非常不接受.
有一套爵士鼓.
有沒有搞錯?這麼不屬靈.
我不是說不屬靈.
這樣怎會行?.
當然要唱生命聖詩.
有三國琴是屬靈的.
最好有風琴.
唱生命聖詩最屬靈.
你唱現代詩歌那些.
那些不能接受.
有些是這樣的.
但其實又是講宗教改革.
五百年前有一個很重要的人物叫慈雲里.

$^{561}$不是慈雲山,慈雲里.
他很喜歡音樂.
你不知道他是誰也不要緊.
總之你知道他很厲害.
他很厲害就行了.
他很喜歡音樂.
但他反對崇拜裡面有音樂.
他覺得音樂令人不能聚焦在上帝的道裡面.
於是他的崇拜是沒有音樂的.
沒有聖詩的.
只是講的,聽道的.
另外有一個很重要的人物叫韓家二汶.
你不記得他是誰也不要緊.
只要你記得他很厲害就行了.
他對音樂的看法是正面一點.
在崇拜裡面可以唱詩.
但不准有樂器.
即是說你彈三國琴.
不屬靈.
影響崇拜.
他們是這樣想的.
其實你想想對於我們來說.
可能有些東西你覺得很傳統.
彈琴,生命聖詩你覺得很傳統.
但是在五百年前你離經半道.
竟然唱生命聖詩這樣的詩歌.
在五百年前你是很反叛的.
其實我們現在是同樣的想法.
你會覺得唱現代詩歌打樂隊那班人很反叛.
在台上彈結他.
你覺得他們很反叛.
接受不到.
在崇拜裡面不應該出現.
但是原來只不過我們的眼光比他們遲.
你曾經做的事情在以前的人眼裡.
你也是離經半道.
有很多事情在不斷地變.
以前的傳道人全部都像老牧師一樣.
西裝骨骨最好沒有頭髮.
像我這些這裡剷得這麼青.

$^{601}$我經常都被人說.
我剛剛出來牧會的時候我紮著馬尾.
我有沒有試過紮馬尾來港獨.
我讀神學院的時候.
這裡剷青了紮著馬尾在那裡走來走去.
有些人看到我都覺得.
以前的童工不可以染頭髮.
現在的女童工每天染頭髮都沒有人說她.
我想遲些還可以紋身.
我不是說可以.
我跟大家說.
我不是說一定可以.
我只是想提出一件事就是世界在變.
我想提出的就是這個世界在變.
我想大家去思考.
如果我們像我剛才那樣.
我保持著五百年前的心態.
我保持著五百年前的態度.
其實你就是在用舊皮袋.
你不願意用新皮袋去承接新的東西.
你說有多大的損失.
如果我們保持著.
就是不可以用音樂不可以唱詩.
在詩歌裡面其實大家領受了很多東西.
你有沒有覺得有時有些詩歌是在塗鑿你的生命.
你有沒有覺得有時有些詩歌是幫助你去認識這個信仰.
但如果我們不願意用新皮袋的話.
我們就損失了很多很激勵人心的詩歌.
我們損失了很多寶貴的屬靈經驗.
其實同樣道理.
今天很多的基督教音樂的樂隊.
他們用很現代的方式.
有些說唱 有些不同的音樂方式.
他們都做得很成功.
已經有些基督教音樂其實入到屋.
去到一些未信的人的耳裡面.
有些就算未信的都聽過.
其實都可能有點嘔.
鄭秀文的《不如驚動愛情》.
那些歌入到屋.

$^{641}$其實對於非信徒來說已經是一個認識教會的機會.
我們如果停留在一種.
還要唱生命聖詩.
沒有人聽.
我不是說生命聖詩有問題.
也是崇拜唱是很好的.
我講的是你失去了一個去接觸其他人的機會.
又再舉多一個例子.
在疫情期間這個世界變了很多.
這個世界真的變了很多.
很多東西變了可以網上進行.
明天都有一個網上報道會.
我們很多東西變了網上.
你可以網上上課.
你可以網上開會.
你可以網上開主日行.
你可以網上崇拜.
很多東西都變到可以網上.
世界變了能不能回頭.
有些東西回不了頭.
你不可以說到了疫情過了之後.
不可以報道會就一定要現場.
不可以直播.
那時候特事特辦.
不准直播一定要現場.
一定要回來.
還一定要訂大球場.
可能是不可行的.
既然是這樣.
你還需要停留在一種舊皮袋的時代.
其實可能是上帝透過疫情.
在改變我們教會的生態也不一定.
這個時候我們就要思考一下.
我們是否應該要用新皮袋.
去迎接這個新走.
是否應該要用新的思維.
去迎接這個改變了的世界.
我不是說甚麼都可以改.
我最後要戴上頭盔.
不是甚麼都可以改.

$^{681}$亦都不可以隨便改.
但是不可以有一種心態.
就是甚麼都不可以改.
如果你覺得甚麼都不可以改.
這個才是問題.
你是會和時代脫節的.
你有認真思考過.
最終你有一個論述去說出來.
為甚麼不可以改.
我覺得這件事完全可以.
我要強調.
有些東西我沒有給稿子.
翻譯的就辛苦了.
例如崇拜是否可以傳到網上.
完全不回來可以嗎.
我剛才說過你可以上網.
但崇拜完全不回來可以嗎.
不可以的.
我們的信仰.
我們說的上帝是耶穌基督.
是道成肉身的.
他是肉身在人中間的.
我們肉身的交往.
人與人之間的交往.
是不能夠被網上取代的.
特別是聖餐.
聖餐你真的拿起餅.
大家同手一個餅.
同手一個杯.
這個是有神學意義在背後的.
你是不可以隨便改了它.
亦都不可以就這樣改變了它.
廢棄了它.
你經過思考.
經過神學的思考.
有神學的判斷.
然後不改.
是沒有問題的.
是可以的.
問題是千萬不要想都不想.

$^{721}$我就說不可以改了.
是這樣的.
我們做了幾十年都是這樣.
動都不動.
這樣就很大問題.
所以我們一同去摸索.
我們一同去思考.
不要與這個時代脫節.
我們要想想.
有沒有一些舊皮袋.
束縛了我們.
阻礙了我們做一些新的事工.
阻礙了耶穌基督.
阻礙了教會的發展.
我們一同去摸索.
怎樣去做一個基督徒.
我們同心祈禱.
我們相信你的教導.
是永遠都是神的.
是歷久常新.
聖靈是會透過你的話語.
不斷地去提醒我們.
讓我們能夠在這個時代裡面.
繼續作你的見證.
成為你的炎 成為你的光.
求你保守我們.
經常都有一種新的思維.
讓我們敏感這個世界的改變.
讓我們知道這個世界怎樣發展.
好讓我們能夠把握不同的機遇.
去傳福音.
去擴展上帝的國.
求你的靈與我們同在.
我們祈禱交託奉耶穌基督的名號.
阿們.
\newpage



\section{尼希米記 4:7-14}
\label{sec:JhzI3zqDYAI}
\textbf{2021.05.30《內外受敵,卻不被困住》 尼希米記 4\_7-14 ( 和修版)  陳一華牧師}
\newline
\newline
連結: \href{https://youtube.com/watch?v=JhzI3zqDYAI}{\texttt{https://youtube.com/watch?v=JhzI3zqDYAI}} ~~~~ 語音日期: 2021-05-30
\newline
\newline
\hyperref[sec:0Rk0UlP9K50]{\small{< < < PREV SERMON < < <}}
~
\hyperref[sec:index_chronic]{\small{[返順時目]}}
~
\hyperref[sec:index_scriptual]{\small{[返順卷目]}}
~
\hyperref[sec:11ItqReN_00]{\small{> > > NEXT SERMON > > >}}
\newline
\newline
尼希米記 4:7-14
\newline
\begin{longtable}{cl}
\hline
\hline
章節 & 經文 (和合本修訂版)\\
\hline
4:7 & \begin{tabularx}{0.7\textwidth}{X} 參巴拉、多比雅、阿拉伯人、亞捫人和亞實突人聽見耶路撒冷城牆正在修造，破裂的地方開始進行修補，就非常憤怒。 \end{tabularx} \\ \\ \relax
4:8 & \begin{tabularx}{0.7\textwidth}{X} 大家同謀要來攻打耶路撒冷，使城混亂。 \end{tabularx} \\ \\ \relax
4:9 & \begin{tabularx}{0.7\textwidth}{X} 然而，我們向我們的神禱告，又因他們的緣故，就派人站崗，晝夜防備他們。 \end{tabularx} \\ \\ \relax
4:10 & \begin{tabularx}{0.7\textwidth}{X} 但猶大有話說：「扛抬的人力氣衰弱，瓦礫太多，我們自己不可能建造城牆。」 \end{tabularx} \\ \\ \relax
4:11 & \begin{tabularx}{0.7\textwidth}{X} 我們的敵人說：「趁他們不知道，看不見的時候，我們進入他們中間，殺了他們，使工作停止。」 \end{tabularx} \\ \\ \relax
4:12 & \begin{tabularx}{0.7\textwidth}{X} 那靠近敵人居住的猶太人十次從各處來見我們，說：「你們必須回到我們那裡。」 \end{tabularx} \\ \\ \relax
4:13 & \begin{tabularx}{0.7\textwidth}{X} 我叫百姓站在城牆後邊低窪的空處，使百姓各按宗族站著，拿刀、拿槍、拿弓。 \end{tabularx} \\ \\ \relax
4:14 & \begin{tabularx}{0.7\textwidth}{X} 我察看了，就起來對貴族、官長和其餘的百姓說：「不要怕他們！當記得主是大而可畏的。你們要為你們的弟兄、兒女、妻子、家園爭戰。」 \end{tabularx} \\ \\
[1ex]
\hline
\hline
\end{longtable}
$^{1}$大姐妹早安.
早安.
我們今時今日可以有實體的去敬拜.
感不感恩?.
感恩.
很感恩.
尤其是三十四度高溫.
走到教會都平平安安.
感不感恩?.
感恩.
所以能夠敬拜主.
一切都是恩典.
牧師今早跟你們分享的題目.
是內外受敵卻不被困住.
希望這個情況不是在你的身上.
不過有時候也避無可避.
有時候我們生活當中.
真的遇到什麼?.
內外都有困難.
所以我們如何去面對這些困難呢?.
所以今天選用的經文是比較冷門一點的.
《離希米基》.
有沒有讀過《離希米基》?.
有些人讀過.
但可能很多人未聽過《離希米》.
其他米可能是吃過.
《離希米基》是什麼呢?.
都未知道.
今天的經文.
亦有時候有點苦悶.
因為有些歷史背景.
剛才一開始有兩個人物.
有參巴拉.
有多比亞.
這些人都是屬於阿拉伯人.
又是阿滿人.
又是阿薩達人的領袖.
他們要去攻擊誰呢?.
攻擊離希米.
為什麼要攻擊離希米呢?.

$^{41}$因為離希米是帶頭.
要將耶路撒冷的城牆重建起來.
為什麼要建城牆呢?.
是因為大家都知道.
猶太人曾經亡國.
亡國後被人擄去外地.
一個朝代一個朝代.
到後來.
七十年後.
上帝的應許實現了.
讓猶太人可以回歸祖國.
結果猶太人有三次回歸.
這些行動當中.
第三次就是離希米帶隊.
離希米那時候.
是在朝廷做官.
因為他在波斯王亞達舍西的朝代當中.
做官.
經常見到皇帝.
有一次皇帝見到他.
哎呀.
愁眉苦臉.
愁眉心鎖.
即是頭凹凹的那種.
皇帝當然覺得他有事.
就問他.
其實離希米真的要抹一抹汗才行.
為什麼呢?.
原來在皇帝面前頭凹凹的.
行不行的?.
不行的.
你怎麼可以在皇帝面前愁眉心鎖呢?.
因為離希米有心事.
心事重重.
皇帝一問.
他知道皇帝都什麼.
都看出他.
不過都戰經的.
不知道皇帝會不會一生氣.
將他人頭斬下來.

$^{81}$所以他就祈禱.
他要回答皇帝這個問題的時候.
他就祈禱.
他心裡面祈禱.
他當然不敢公開.
心裡面祈禱.
因為皇帝問他這個問題.
為什麼你的樣子這麼難看.
於是他就跟皇帝很坦白說.
他說皇帝.
我的祖國.
城牆倒塌.
城門失火.
我的祖家.
團緩敗瓦.
不堪入目.
你明白我的意思嗎?.
在這樣的情況下.
我的心情就很沉重了.
所以皇帝現在明白.
為什麼他會愁眉心鎖.
皇帝很好.
看出他.
你有什麼要求?.
他就說.
皇啊.
如果你給我機會.
我想告大架.
我想拿大架.
做什麼?.
我想回到自己的祖家.
能夠重建城牆.
皇帝批不批給他?.
批給他.
不僅批給他.
他說.
你可以回去.
還帶猶太人回去.
他就跟皇帝說.
如果我在你眼前蒙恩.

$^{121}$皇帝可不可以給我一些建築材料呢?.
你知道.
建城牆要不要材料?.
要建築材料的.
皇帝給不給?.
給.
這個皇帝真是很好人.
疼愛他.
其實.
你知不知道還有一個要求.
他說皇帝.
既然你給我材料.
又讓我回去.
你可不可以派兵.
護送我回去.
這個真是越來越大膽的要求.
但是皇帝批不批?.
都批.
可見.
尼希米在皇帝眼前蒙恩.
當然.
大家都知道.
他之所以這麼順暢.
其實背後有天上的神.
有這位創造的主.
以色列人.
猶太人的上帝.
在背後掌管著一切.
令到皇的心都願意配合.
答應尼希米的要求.
這是一個小小背景.
讓大家知道.
為何他要建城牆.
為何氣氣下.
現在就是了.
氣氣下就發生大事了.
為甚麼?.
原來回到自己祖家的時候.
四周的外邦人.
甚麼阿滿人.

$^{161}$阿薩達人.
阿拉伯人.
那些外國人.
外族人.
看見他們在這裡開工.
要建城牆.
他們的心.
快不快樂?.
當然不快樂.
一定是甚麼?.
一定是心情很緊張.
虎視眈眈.
看你們建城牆.
做甚麼?.
是否要造反?.
是否要跟他們對抗?.
你明白我的意思嗎?.
如果一旦讓你們成功建了城牆.
你便可以放肆.
做甚麼都可以.
所以這班人.
真的不容易放過尼希米和猶太人建城牆.
你知道這單工程.
遇到的困難真的很大.
遇到的艱難很大.
所以.
我們一開始看這段經文之後.
你便發覺到.
這些經文剛才主席已經幫我們讀了.
很清楚.
我們看下去.
當我們遇到困難的時候.
在前一章.
他們靠著祈禱.
靠著堅持的態度.
要繼續甚麼?.
完成這個工程.
但是之前那班人.
不過是奚落一下.
不過是嘲笑他們.

$^{201}$但是來到這段經文就不同了.
噩夢是甚麼?.
開始了.
為甚麼?.
不只是口說.
原來這班外族人.
外邦人.
就要包抄他們.
要開始用行動去攻擊他們.
你看看有甚麼行動.
原來有三個行動.
來攻擊猶太人和尼希米.
大家有留意到嗎?.
第一個行動.
就是這幾個外族人.
大家合謀.
大家同謀.
一起要甚麼?.
要阻擋猶太人.
尼希米.
不讓他們得逞.
這個就是同謀合謀.
第二.
他們就說要偷襲他們.
即是說攻其不備.
當猶太人沒甚麼的時候.
我們要.
趁他們一個不留神.
就殺光他們.
令他們的工程全部停頓.
這個就是攻擊他們.
第三招就是重告招.
第三招是甚麼?.
是心理戰術.
就是散播謠言.
在他們當中散播無中生有的謠言.
令猶太人聽到的時候.
心裡很怯.
每個人都沒有鬥志.
令他們開始感覺到.

$^{241}$這些謠言都不知道真假.
你知道很多時候.
人們聽一點不聽一點.
人很有趣.
有時候你說真話沒甚麼人聽.
你說假話很多人聽.
你說正當話.
人們耳邊風.
你說邪話.
他們都相信.
有沒有這樣的人?.
很多這樣的人.
所以這些心理戰術.
三招令猶太人在那個時候.
就好像一條車胎.
那條車胎做甚麼?.
就插穿了.
洩氣了.
本來想著一鼓作氣.
去起城牆.
但現在遇到這三招.
他們的猶太人.
面對那種逆境.
面對別人攻擊的時候.
就好像車胎沒了氣.
車胎都爆了.
在這樣的光景下.
就出現了四個現象.
哪四個現象呢?.
這四個現象真是令到.
尼希米作為領袖.
真是很煩惱.
怎樣面對.
不要說攻城.
怎樣面對外來的敵人.
怎樣防備外族人.
已經夠頭痛了.
但又看回現在自己國家.
起城牆的工程.
搞到甚麼呢?.

$^{281}$搞到不湯不水.
即是說.
第一.
聖經描述.
當時警方是甚麼?.
是灰土上多.
他們有四失.
第一個失是甚麼呢?.
就是失掉了意志.
就是失去了意志.
為甚麼?.
為甚麼見到他們失去了意志呢?.
因為他們的工程.
起城牆.
周圍都很多甚麼?.
灰土,沙土.
大家知道這樣的情況.
有沒有人開工?.
沒有人開工.
大家有時都聽過爛尾樓.
有沒有聽過爛尾樓?.
希望你不要買到這樣的樓.
有時爛尾樓就是這樣.
起到兩層,三層.
然後還要再起嗎?.
就不再起了.
旁邊很多泥沙.
很多鋼筋.
被人劈到一邊.
這種情況就好像當時.
起城牆.
第一件事.
就是灰土上多.
他們看到這些沙土,灰土.
越看就越閉眼.
就好像今天.
弟兄姊妹.
你們有時平日生活當中.
非常忙碌.
很多事要做.

$^{321}$很多家務.
很多文件未處理.
很多報告未做好.
有沒有遇到這種情況?.
很煩惱.
為甚麼工夫.
一來就不來.
有時工夫一來.
就像排山倒海.
根本應付不來.
有時生活很混亂.
甚至不知道自己最近做甚麼.
最近自己生活.
都不知道為甚麼而生存.
有時人就落在一個很混亂的地方.
這種情形就好像.
剛才我所描述.
猶太人就失掉了.
失去了意志在當中.
好.
第二個失是甚麼呢?.
就是失掉了力量.
為甚麼呢?.
因為猶太人.
失掉了力量.
因為第六節告訴我們.
這班猶太人.
起牆起氣.
就覺得很疲累.
或者心情很惡劣.
你明白嗎?.
聽到這麼多謠言.
聽到這麼多說話.
影響自己的心情.
你知道.
當你的心情受影響.
心情差的時候.
自然力氣都失去了.
為甚麼呢?.
因為都不想做了.

$^{361}$連力量都失去了.
難怪工程最大的關卡.
是甚麼時候呢?.
就是當你做到一半的時候.
是嗎?.
你想想.
就算做生意.
開店鋪也好.
做甚麼工作也好.
起初你都很熱切.
很熱誠.
很熱情.
一定要甚麼?.
一定要去做成.
但誰不知一開工.
一開店鋪.
一開工程.
去到一半就最辛苦了.
為甚麼呢?.
因為創業容易.
下半句你知道是甚麼.
是嗎?.
真是怎樣埋尾呢?.
怎樣守下去呢?.
各位頂尖的朋友們.
我對陶儀輝這句說話.
是很容易甚麼?.
打擊一個人.
令他失去力量.
再向前行.
這是第二個失.
第三個失是甚麼呢?.
這是一個例子.
如果你行山.
你喜歡行山的弟兄姊妹.
你發覺第一個小時不難行.
第二個小時開始抽筋.
第三個小時就不行.
或者找人抬你.
所以有時真是甚麼?.

$^{401}$去到一半是甚麼?.
是最大考驗.
我們看看.
第三個失是甚麼?.
失去了信心.
因為猶太人.
整整下.
你和我說.
我和你說.
都是不行的.
都是起不成的.
都是不可以做的.
做都沒有用.
你明白我的意思嗎?.
這些就是.
在他們遇到.
外面的攻擊的時候.
那種狀況.
失去了甚麼?.
信心.
這是很敝屣.
很倫敘的地方.
所以你發覺.
做領導.
帶領工作的.
做頭.
容易做嗎?.
很容易做的.
因為猶太人真是面對.
外憂又面對來還.
所以.
最大的危機.
就是失去了信心.
你做甚麼都好.
如果你被人打沉了你的信心.
你就很難.
再繼續做下去.
這是我們時時都要提醒自己的.
上帝給你有甚麼的事奉.
有甚麼的工作.

$^{441}$記得.
不要.
打沉了信心.
因為一失去信心.
你就會半途而廢.
你就做不下去.
所以這個要很警惕.
很警醒.
有時魔鬼撒旦.
被你的攔阻.
去到一半.
你就做不下去.
沒有信心.
我們發覺第三.
最可惜就是失去了信心.
為甚麼呢?.
就是失去了安全.
為甚麼?.
因為那些外族人.
大家都夾定了.
要攻其不備.
等到猶太人做事.
或者收工的時候.
晚上的時候.
殺入去.
乘虛而入.
將這班猶太人清除.
他們就不可以再.
起城牆.
在這樣的情況下.
猶太人的處境.
是甚麼?.
每個人都覺得.
很岌岌可危.
在這樣的情況下.
真的失去了安全感.
所以你發覺到.
當你沒有安全感的時候.
你心更加慌.
害怕.

$^{481}$不想去做事.
在這樣的情況下.
敵人就得逞了.
所以各位電視節目.
幸好聖經的經文沒有停在這裡.
原來還有下文.
我們真的很期望.
看看下文是甚麼呢?.
原來.
在這樣.
內憂外患下.
來希米做阿頭.
做帶領的.
真的很辛苦.
有時候.
在我們.
最近的生活當中.
要面對疫情.
你知道.
疫情.
都對很多人.
造成了很多打擊.
有些在經濟上.
受到打擊.
有些做中小企的.
又遇到難處.
有些.
維持不了.
又要動身.
又要減薪.
甚至裁員.
有時候你都發覺.
在疫情下你都會遇到.
遇到困難.
在我們的生活當中.
有時候都出現了.
一些攻擊.
令到你內外.
都受包圍.
內外都受敵.

$^{521}$你都不知道怎樣去處理.
在這樣的光景下.
你認識這位老人家.
他做阿頭.
辛苦嗎?.
他最初做阿頭就黑頭髮.
做了七年之後.
變成了這樣的頭髮.
眉毛都白了.
因為做阿頭帶領.
是很辛苦的.
所以我們多些要為我們的領袖.
祈禱.
求上帝加能賜力給他們.
好.
面對一個這樣的內外.
受敵的時候.
黎希米就用了四招.
來化解.
這場的困難.
剛才我們遇到有四失.
現在我們看看.
黎希米怎樣去.
化解這場的危機.
他有四個的做法.
第一.
他的行動是什麼呢?.
就是不住的祈禱.
懇切的禱告.
因為在.
這樣的危險.
以及在.
現在.
最脆弱的時候.
最軟弱的時候.
猶太人又不能夠.
一致同心的時候.
他唯一可以做的.
就是鼓勵大家去祈禱.
鼓勵大家一起去禱告.

$^{561}$仰望上帝.
透過祈禱.
來重獲他們的信心.
能夠將那些謠言.
能夠去化解.
因為我們大部分人.
都不再化解.
因為我們大家的注意力.
專注力就不再放在.
其他人的身上.
也不放在閒言閒語的身上.
大家都將焦點.
放在正.
放在什麼?在神的身上.
所以弟兄姊妹們.
原來在困難的當中.
遇到有疫情.
遇到有攻擊的時候.
當你好像走投無路.
又不知道怎樣去.
化解這場危機的時候.
我們不妨學一下.
尼熹米.
尼熹米做得很好.
他在第九節.
他怎樣說呢?.
願意我們向我們的神禱告.
向我們的神禱告.
所以這個行動.
是整個化解危機的第一步.
就是將困難.
帶到神的面前.
求上帝去幫助.
求上帝.
去給他們有智慧.
求上帝引領他們.
怎樣再去踏出第二步呢?.
這個是很重要的.
當我們遇到有困難的時候.
我們要學習尼熹米.

$^{601}$這個很重要的一個行動.
就是懇切的禱告.
一起的來禱告.
來禱告倚靠上帝.
我們再看.
第二個.
尼熹米做的.
就是他和所有的猶太人.
來講說話.
告訴他們.
在現在這個情況.
我們要守望相助.
要大家同心團結.
特別在第十三節的時候.
他就怎樣說呢?.
尼熹米說.
我叫所有的百姓集中.
告訴他們.
每個人都在你們家住的附近的地方.
來去把守住.
在城牆後面低窪的地方.
又要拿什麼呢?.
拿所有的工具.
不單是開工.
還要什麼呢?.
還要抵抗敵人.
所以在這個情況裡面.
第二個行動是很重要.
就是他們願意聽從領袖的帶領.
大家每一個人都要來到.
都要在自己的崗位上.
自己所住的地方.
盡我們的本份.
盡我們的責任.
來做好我們把守崗位的工作.
這個真的有時候.
你知道是很不容易的.
當你發號施令是容易的.
叫大家一起上班.
一起工作.

$^{641}$但是叫了出來之後.
可能反應是不同的.
有些人說不好意思.
我要去買菜煲湯.
我來不了.
有些人說不是.
開工了.
我的工廠剛剛開了.
我現在要回去工作.
各位電視節目的觀眾.
如果當我們遇到.
有任何困難的時候.
我們大家都顧著自己的利益.
你明白我的意思嗎?.
大家都顧著自己的需要.
而忘記了整體的需要.
而忘記了上帝的心意.
這樣就淪盡了.
為什麼?.
因為這樣你就會做自己個人的事.
但是反過來.
當你在這些危機困難當中.
你能夠把焦點放遠一點.
放在其他人的身上.
放在你身邊的人的身上.
看回整體的需要的時候.
那你就不同了.
為什麼?.
因為你的焦點不再看自己.
而是看你身邊所發生的事.
身邊的人.
怎樣來到大家可以互相配合.
在這樣的情況下.
就真的發揮到守望相助.
大家來到站穩崗位.
堅持要去做領袖所吩咐他們做的事.
我覺得這種同心合意.
重不重要?.
很重要.
這種一同去承擔困難.

$^{681}$一同去面對困難.
這樣的心智是非常非常重要.
所以當我們要合一的時候.
弟兄姊妹們.
記得我們要放下個人的利益.
放下個人的需要.
我們要反過來去看一看.
我們整體的需要是什麼.
我們身邊的人的需要是什麼.
我們現在落在什麼處境裡面.
所以在這樣的情況下.
我們更加要大家團結同心去做.
無論是你家裡發生事也好.
又或者你的工廠發生困難也好.
或者在教會裡出現困難也好.
弟兄姊妹們.
記得這是逃避不了的.
但是重要的是什麼呢.
就是上帝是掌管一切.
我們繼續去仰望神.
去禱告.
大家繼續互相配搭.
團結同心.
一起去事奉神.
守望相助.
這樣猶太人就可以渡過這場危機.
免招致殺身之禍.
否則他們很快就會被外族人攻入城牆.
所以我們看到.
第二個行動.
守望相助是非常需要的.
第三個行動是什麼呢.
第三個行動就是要專注靠神.
為什麼呢.
因為尼希米向他們宣告.
他說我們要紀念我們的主.
我們的神.
一起說是什麼.
大而可畏.
大有能力.

$^{721}$為我們征戰的神.
原來我們能夠做成任何的工作.
記得不是我們自己的能力.
是可以做得到的.
很多時候你發覺.
當你無能為力的時候.
剛剛就是神開始動弓的時候.
有些弟兄姊妹告訴我.
他說他和人傳福音.
很有成功感.
跟隊去派單張傳福音.
他很有成功感.
為什麼呢.
因為他真的做到很多這些工作.
但是當他說向爸爸媽媽傳福音的時候.
他就有很大的失敗感.
為什麼呢.
因為每次說都被人罵.
你明白我的意思嗎.
每次說都被人罵.
所以他一向家人傳福音.
就投降了.
有沒有這樣的情況.
有的.
各位弟兄姊妹們.
這就證明一件事.
原來困難的事情.
不可能成就的事情.
在人的確會是.
但是當你專注靠上帝.
整幅圖畫就不同了.
因為上帝是一位無所不能的神.
尼希米說我們的神是大而可畏的.
大有能力為我們爭戰.
還需要怕這麼多嗎.
不需要.
沒有怕了.
為什麼呢.
因為我們這次共渡危機困難.
不是靠自己.

$^{761}$不是靠我們的才剛.
我們的經驗.
我們的能力.
我們這次能夠渡過這個危機.
衝破內外受敵.
是因為什麼呢.
因為神在我們當中.
因為神與我們同在.
因為神是幫助我們.
當我們知道神是幫助我們的時候.
仇敵的攻擊就算不得什麼了.
所以各位弟兄姊妹們.
這個行動是很重要的.
我不得不將這句金句送給大家.
做你的座右銘.
當你們做什麼功夫都好.
參與什麼的時逢都好.
遇到什麼的困難都好.
記得不要看自己.
不要靠自己.
而是什麼.
黎希米怎麼說.
紀念我們的主是什麼.
大而可畏.
大有能力為我們爭戰.
當我們知道上帝為我們爭戰的時候.
我們何懼之有呢.
所以我就很快想起.
以前以色列人.
他們不聽話的時候.
上帝給他們一些教訓.
上帝怎麼給他們教訓呢.
就是用外邦人來壓制他們.
所以就出現了時思的年代.
有沒有聽過時思記.
其中有一個時思.
大家都很熟悉他的名字.
叫做基澱.
有沒有聽過基澱.
聽過基澱.

$^{801}$你知不知道基澱.
怎麼打羊那場仗.
怎麼打聖米殿人.
阿瑪利人你記不記得.
是不是因為他們人多就打勝.
是不是因為他們用人海戰術.
不是的.
上帝說你一吹雞吹得不錯.
一吹就三萬二千人就舉手跟你.
不過我覺得人多過頭.
你告訴他們.
誰心情害怕就不要來了.
不要參加打仗了.
一吹就吹了二萬二千人.
厲不厲害.
剩下一萬人.
上帝說一萬人夠不夠.
一萬人都多過頭.
為什麼.
因為他們將來打贏仗.
就容易驕傲.
為什麼.
這次我們能夠打贏仗.
能夠打贏那些米殿人.
是因為我們人多勢眾.
你明白我意思嗎.
我們人強馬壯就容易驕傲.
所以上帝說一萬人都太多.
叫他們每個人去喝水.
如果你讀過《時詩記》.
是很有興趣的.
最後.
你知不知道上帝揀選多少人.
上帝揀選三百人.
去打好像蝗蟲這麼多的米殿人.
有沒有打.
人體有沒有打.
結果.
機殿和那三百人.
大勝米殿人.

$^{841}$而且海船而歸.
各位頂智妹妹.
因為什麼.
因為上帝為以色列人去征戰.
這是一個最關鍵的重點.
今天我們遇到任何的困難.
頂智妹妹.
要學習離棄米給我們的教導.
上帝是什麼.
大而可畏.
我們要敬畏上帝.
不要自以為是.
我們要聽神的話.
當你們聽神的話.
走在神的旨意當中的時候.
頂智妹妹.
神就會為我們去征戰.
但當我們沒有聽上帝的話.
自把自為.
做自己的事的時候.
上帝離棄了我們.
我們就不能夠作神.
所以要記得.
離棄米能夠扭轉這個形勢.
是因為他對以色列人的教導.
是很清晰.
我們再看離棄米的第四招是什麼.
第四式是什麼.
第四式.
就是要堅持下去.
他在聖經第十四節.
告訴我們.
他怎麼說呢.
他說你們要為你們的弟兄.
兒女.
妻子.
家園.
你們去征戰.
換句話來說.
我們每一個人.

$^{881}$都要盡我們的本份.
當上帝願意幫助我們的時候.
我們自己都需要盡我們的努力.
不是說翹起手.
什麼都不做.
叫上帝來幫我們打勝仗.
不是的.
我們都要盡我們的本份.
所以當這一班的以色列人.
猶太人.
他們願意各盡其職.
盡他們的本份的時候.
堅持做上帝喜悅的事情.
結果他們就能夠將這一次的危機.
變成了一個契機.
所以各位弟兄姊妹們.
今天也是.
當你發覺在你生活當中.
遇到有什麼難題的時候.
不要忘記.
不要忘記親近神.
不要忘記讀聖經.
不要忘記靈修.
不要忘記回團契.
回崇拜.
不要忘記跟人傳福音.
當你做上帝喜悅的事情.
當你去做上帝要你去做的工作的時候.
你可以放心.
你所遇到的難題.
遇到的困難.
上帝一定會幫助我們.
將它們解決.
因為我們的神.
祂是無所不能的神.
所以各位弟兄姊妹們.
今天我們在這個訊息裡.
我們大家互相的鼓勵.
互相的提醒.
我們發不得以禮欺迷.

$^{921}$給我們這次四個行動.
也都成為了我們一個操練的行動.
以致到在我們遇到.
有受困難攻擊逆境的時候.
我們都能夠可以將它.
將危機變為一個契機.
我們一起低頭我們祈禱.
透過一段我們比較陌生的經文.
再一次給我們看到.
我們在遇到困難的時候.
我們怎樣能夠可以去克勝這些困難.
我們真的需要更加同心團結.
一起有一個合一的心智.
來聽上帝的說話.
來仰望依靠上帝的能力.
以致到我們在我們的征戰裡.
經歷到你自己的祝福.
你的同在.
多謝天父.
叫我們個人的生命.
我們的家庭.
我們的教會.
我們的事奉.
我們的社會.
都能夠依靠上帝.
仰望上帝.
來遵行你的命令.
做你所吩咐我們做的事情的時候.
我們就能夠在種種的困難當中.
看到原來有曙光的出現.
原來上帝是幫助我們.
叫我們更加能夠去讚歎.
去感恩你自己奇妙的作為.
祈禱感恩.
奉耶穌的名求.
阿們.
謝謝大家.
\newpage



\section{路加福音 4:18-19}
\label{sec:11ItqReN_00}
\textbf{2021.06.06《釋放與捆綁》 路加福音 4\_18-19 ( 和修版)  雷子健牧師}
\newline
\newline
連結: \href{https://youtube.com/watch?v=11ItqReN_00}{\texttt{https://youtube.com/watch?v=11ItqReN\_00}} ~~~~ 語音日期: 2021-06-06
\newline
\newline
\hyperref[sec:JhzI3zqDYAI]{\small{< < < PREV SERMON < < <}}
~
\hyperref[sec:index_chronic]{\small{[返順時目]}}
~
\hyperref[sec:index_scriptual]{\small{[返順卷目]}}
~
\hyperref[sec:hNtKC6DPyIM]{\small{> > > NEXT SERMON > > >}}
\newline
\newline
路加福音 4:18-19
\newline
\begin{longtable}{cl}
\hline
\hline
章節 & 經文 (和合本修訂版)\\
\hline
4:18 & \begin{tabularx}{0.7\textwidth}{X} 「主的靈在我身上，因為他用膏膏我，叫我傳福音給貧窮的人；差遣我宣告：被擄的得釋放，失明的得看見，受壓迫的得自由， \end{tabularx} \\ \\ \relax
4:19 & \begin{tabularx}{0.7\textwidth}{X} 宣告神悅納人的禧年。」 \end{tabularx} \\ \\
[1ex]
\hline
\hline
\end{longtable}
$^{1}$我讀今日講道的經文給大家聽.
是路加方向四章十八到十九節.
請聽我讀便可以.
路加方向四章十八到十九節.
主的靈在我身上.
因為他用告告我.
叫我傳福音給貧窮的人.
差遣我報告被擄的得釋放.
乞眼的得看見.
叫那受壓制的得自由.
報告神月立人的欺憐.
請我們同心祈禱.
神啊我們多謝你.
我們來到你的面前要感恩.
因為你能夠釋放我們.
無論我們大罪也好小罪也好.
在你來講沒有分大小.
你都能夠釋放我們.
只要是我們親近你.
神啊你就充滿我們.
主我們感恩多謝你.
我們聽到Anna姐妹的分享.
她的生命的經歷.
真的要來到你面前要頌讚你.
因為你拯救她.
主我們多謝你.
主我們更加感恩.
你的救恩臨到我們每一個人的身上.
將紅恩堂每位弟兄姊妹.
特別新朋友我們交托給主.
主啊這一刻充滿他們.
聖靈要臨在他們當中.
主啊讓我們同感一靈領受上帝的說話.
並且我們行出上帝的說話.
知道你來到這個世界.
釋放每一個罪人.
多謝主.
聽我們同心祈禱.
奉基督明求.
阿們.

$^{41}$好了.
牧師也要介紹一下自己.
牧師也是一位過來人.
我14歲的時候跟黑社會.
16歲的時候吸毒.
不知數了多少年.
現在你們看我的照片就知道我小時候.
在門上那副.
弄著腰的那副.
那時候只有十幾歲.
那次就是第一次坐牢.
剛剛出差回家拍的那張照片.
很奇妙.
神的改變.
神的臨在在我的身上都是很奇妙的.
我怎麼會想到我是一個吸毒的人.
我坐過幾次牢.
神竟然在我身上發生作為.
祂呼召我.
祂並且讓我成為牧師.
大家看到旁邊那張照片.
就是案立的時候拍的那張照片.
如果看清楚日期是2001年.
今年案立了20年.
我值得感恩.
其實我的出身也不是一個好的出身.
我爸爸是黑社會人.
我姓雷.
你知不知道我爸爸的花名叫什麼.
雷公就是雷公.
哈哈哈哈.
我小時候住在七層寨子區的東頭村.
雷公這個名字是很厲害的.
沒有人敢靠近他.
他都很害怕他.
他就是名副其實.
所以他是一個很可怕的人.
我媽媽也不是好人.
我媽媽開了一家大店.
兩個人就黑鬥黑.

$^{81}$在家裡就.
我總共有五個兄弟姐妹.
我對上有三個是不同爸爸生的.
所以我媽媽帶了三個不同爸爸生的兄弟.
兩個哥哥一個姐姐就跟了我爸爸.
也不是正式結婚住在一起.
我之下有個弟弟.
所以我五個兄弟姐妹.
後來我出生幾歲的時候.
爸爸和媽媽就分開了.
媽媽就帶了兩個哥哥一個姐姐走了.
我就跟了阿嬤.
我最小的弟弟跟著我爸爸成長.
我到了十一歲的時候.
我阿嬤死了我就回到爸爸家裡住.
一個超過十年沒有相處的家庭.
回到家裡很陌生.
這個爸爸也是很陌生.
我爸爸是一個酒鬼.
就是酗酒.
所以家裡也是很複雜的.
我爸爸早死了.
在我十四歲那年他就死了.
他喜歡開快車.
他撞車死的.
他開快車和載車.
就死了.
所以只剩下我一個人在西市區長大.
很容易在行差搭菜.
如果你看到.
我很多個哈哈笑公仔遮住的那幅.
我那時候就是去了狼騎的福音戒毒.
來了戶外中心.
那個時候是二十五歲.
在最後一次坐牢.
其實我坐過幾次牢.
我最後一次坐牢在大欖監獄.
有一年聖誕節.
戶外的弟兄姊妹就進去傳福音.
到現在都是.

$^{121}$其實我們每一年都進去大欖監獄傳福音.
就在聖誕那天.
過去那年因為疫情我們進不去.
不知道今年能不能進去.
所以我就聽到福音.
當時幾百個囚犯在禮堂那裡.
我坐在最前的那一行就聽到福音.
我聽到有人講見證.
我心很感動.
我說主啊.
台上高高這麼壞的人.
好像我這麼壞.
神都能夠改變他.
我相信神你也可以改變我.
我那時候就是這樣想.
我說主啊.
我都願意去試一下跟隨你.
我不敢說是跟隨.
試一下信一下.
結果有一個牧師講道.
他說如果你相信耶穌要跟隨耶穌.
請你舉手.
我想請問你們猜猜我有沒有舉手信耶穌.
有些人說沒有.
有些人說有.
答案是有或沒有.
都是兩個.
我告訴你.
可以說有又可以說沒有.
因為我心裡面我很想舉手.
就是我想信.
就是我信.
但我沒有舉手.
因為三十多年前坐牢.
現在好一點.
現在我進去大欖監獄.
一說信耶穌.
很多人舉手.
黃先生也試過.
現在開放程度又不同了.

$^{161}$以前很慘的.
你說信耶穌.
搞鬼你.
取笑你.
你不能生活.
在坐牢的時候.
所以.
我們有Madam在.
知道這些行情.
三十多年前.
我當然不舉手.
我舉手之前.
我想舉手.
我就看後面有沒有人舉手.
我當然有盤才舉手.
一看.
看到有兩個人舉手.
幾百人之中有兩個人舉手.
一個叫阿痴.
痴線那個痴.
他有些不正常.
真的有些不正常.
一個叫阿美.
他的頸不知道為什麼這麼歪.
他這樣走的.
我心裡想.
一個又痴又歪.
歪來歪去的人.
才信主的.
因為全部都沒有人舉手.
所以牧師當時放棄了.
舉手這個機會.
舉手決戰.
但我心裡都很心動.
動得很厲害.
結果.
在過了幾個月之後.
我出監.
出監俗稱出賊.
你們要懂的.

$^{201}$我想你們都懂的.
叫出賊.
出賊一個吃了白粉.
吸毒九年的人.
一出賊就上什麼.
死人尋舊路.
找一些朋友.
沒什麼的.
找些朋友.
我走回以前.
以前我蒲的地方.
找些朋友.
朋友就找不到.
一個都沒有.
結果在我迎面而來.
就碰到.
那次在聖誕節.
入監獄.
全福音站在台上.
彈結他那個人.
我認到他.
我跟他聊天.
他說我是戶外的童工.
我給你一張卡片.
如果你上去這裡.
我們一定幫到你.
結果我有沒有上去.
當然有了.
不然今天搞成這樣.
哈哈哈哈.
結果上了去.
心裡很大掙扎.
其實因為當時我沒有毒癮.
坐監哪有毒品吃.
我一個沒有毒癮的人.
我為什麼要去福音大道.
其實我想生命去改變.
生命改變就是大前提.
我記得我入村那天.
我什麼都沒有.

$^{241}$村的中心就買了兩條內褲給我.
和一些替換的.
兩件很簡單的衣服.
我進去村裡面生活.
家人又不原諒我.
所以一分錢都沒有給我.
我全部都是.
吃神住.
住神給我們的地方.
我就是這樣的環境裡面去成長.
今天.
我所有的親戚都原諒我.
所有的親人.
我媽媽在臨死之前都信耶穌.
我姑姐在臨死之前信耶穌.
很感恩.
今天雖然什麼都沒有.
我出差的時候.
出戶外什麼都沒有.
神就給我老婆.
神又給我兒子.
神又給我讀神學.
其實讀神學也是很奇妙的.
我讀書不多的.
我去過神學院.
神學院不收我的.
不收我的原因是.
我讀書不多.
讀書.
讀神學.
入學都要有條件.
基本的條件.
中學畢業都沒有.
但結果神學院都給我.
你試讀一下.
你試讀一下行不行.
所以我第一個學位是96年畢業.
我總共讀了十年書.
前前後後讀了十年書.
讀了一個Diploma.

$^{281}$讀了一個學士.
讀了一個碩士.
讀了十年書.
十年兼職讀書.
十年窗下有沒有人問.
當然有.
讀了十年書.
今天我又在讀書.
都為我祈禱.
我寫一個博士論文.
我希望神在我身上.
都幫我來成就.
只要我們努力.
只要我們的心向著神.
神會不會扶你.
一定扶你.
Amen.
這個就是我簡單的人生的介紹.
神在我生命中.
他中賢罪是捆綁我.
捆綁了很辛苦.
九年的吸毒裡面.
很大的痛苦.
經常都很想不吸毒.
但又不行.
又要去吸.
我打針的.
我就是要插自己.
拿針去插自己這麼蠢.
在九年裡面就是這樣.
很想爭脫.
很想不吸毒.
但又不行.
過了很久.
一會又坐牢.
一會又跟人吵.
一會又被人打.
很痛苦.
我最後睡到睡街.
沒有一個親戚原諒我.

$^{321}$睡街.
但今天神.
怎樣.
釋放我的罪.
神釋放我的人生.
神將本來很重的罪.
他拿開.
今天我們看聖經.
剛剛我讀那段經文.
就是耶穌回到自己的家鄉.
有執事第一卷書給他.
他打開來讀.
就是讀了這段經文.
這段經文就是以賽亞書61章1節.
簡稱611.
就是這段經文.
他讀了這段經文.
主的靈在我身上.
因為太陽高高.
叫我全方位給貧窮的人.
參與我報告被虜的得色狂.
下眼的漢奸.
叫我受下的得自由.
報告神以立民的欺孽.
讀完之後就把這本書還給他.
他讀出這段經文.
簡單來說.
這段經文就是.
神要他做的事.
他來到這個世間.
一定要做的這幾件事.
就是他的使命.
算額.
如果大家去看四福音.
耶穌很貧窮.
他到處去傳福音.
他最主要做三件事.
就是口說.
宣講神的話.
用口來宣講神的話.

$^{361}$第二就是講鬼.
第三就是醫病.
他最主要做這三件事.
在他使命宣言.
他就落實告訴大家.
他要做的事.
第一件事.
貧窮人得福音.
沒錯.
在耶穌的年代.
貧富很懸殊.
貧窮人得福音.
耶穌傳福音.
大部分時間都在貧窮人當中.
當然有時候會在有錢人.
例如瑞利.
瑞利很有錢.
但瑞利的人貧不貧窮.
他是有錢.
不過他心靈貧窮.
耶穌最主要來到這個世界.
就是去到一些心靈貧窮.
和實際上真的貧窮的人.
他說你們要得福音.
福音的原文解釋就是好消息.
你們每一個人都要聽好消息.
你們每一個人都要得著神的好處.
那個好消息就是.
罪得赦免和得永生.
信耶穌的人是不是可以得永生.
是不是可以罪得赦免.
這個就是神要告訴他們.
我讀這句話就是告訴他們.
我實際上就是要來做這件事.
沒有一個人.
能夠自己可以解開自己的重擔.
你試一下.
你游水淹到.
你想像一下.
你游水抽筋.

$^{401}$這個人救你的.
你難道自己拔自己的頭髮.
拔起它.
說我沉下去.
我拔起自己.
行不行.
不行的.
罪就是這樣.
一路跌下去.
但你神不拔起你.
是不行的.
這個就是罪得赦免.
貧窮人要得到福音.
是他要做的重要的第一件事.
第二件事.
基督徒說他是被擄得釋放.
什麼人是被擄.
一些叫做奴隸.
他是被擄的.
有沒有自由.
在耶穌時代做奴隸很慘的.
主人叫他做什麼就做什麼.
限制住他.
你不睡也要他完成.
他說他是奴隸.
不值錢的.
一文錢不值的.
做奴隸你永遠都要跟著我.
做到死為止.
但耶穌告訴他們.
被擄得釋放.
最重要的是被罪去擄掠.
我們有些事是做不到的.
有些事想掙脫都不行.
我們那些人.
我們戒毒中心的人.
最主要進來就是戒毒.
我昨天去荔枝角拘留所.
見一個年輕人三十歲.
他有個小孩.

$^{441}$很小的幾個月.
他吸毒犯了事.
有少量的毒品.
他等報告.
等感化官報告.
看判進來我們那裡.
還是判去政府坐牢的戒毒.
我去見他.
見完了.
所有東西見完.
好啊牧師我去你那裡.
好好啊你那裡.
住一年對的.
對的不停的up 頭.
然後我就想走.
談完就走了.
牧師慢著.
你那裡可不可以吸煙.
哈哈哈哈.
我就告訴他.
我們學做好人就要學到底.
煙都不吸.
簡單點說.
不要解釋那麼多.
簡單點說.
他說對.
我們那些人.
很多時候戒完毒出來.
第一時間先吸煙.
喝杯酒.
通常都是這些.
但這些是不是陷阱.
絕對是陷阱.
你吃那些.
來來下下就吃了毒品.
所以呢.
這也是一個捆綁.
很多人戒了毒.
戒完毒就戒不到煙.
那麼慘.

$^{481}$在我們那裡住了一年.
明明什麼都戒了.
出來都要吸煙.
我們這麼蠢.
這麼辛苦戒完毒.
你出來又要吸煙.
這些都是捆綁.
壞習慣的捆綁.
有些是說髒話的捆綁.
我們在村子裡經常遇到.
楚南傳說.
洗洗戶人.
那個人馬上爆幾句髒話出來.
哈哈哈哈.
但我們值得原諒.
剛剛進村子沒多久.
有時候發怒起來.
弟兄吵架.
姐妹也是.
因為他們.
以前做慣了.
如果我們不是真正.
靠著神能不能釋放.
是釋放不了的.
鼻孔的能釋放.
第三黑眼的.
如果你看四福音.
有很多盲眼的.
病人來到神的面前.
能不能得到漢見.
得漢見.
這些是實質的盲眼.
但耶穌要表達.
不單是實質的盲眼.
是心裡的黑眼.
你的心都不打開.
你看不到最重要的東西.
你看不見最重要的.
神的恩典和救恩.
就是這樣.

$^{521}$黑眼的要漢見.
你蒙蔽了你的心眼.
我們的人也是.
我們蒙蔽了心眼.
我們的心眼經常都S兩洞.
多賺點錢.
我們的弟兄.
完成了一年戒毒.
出來第一時間問我.
牧師.
我做這個行業好不好.
我說你做這個行業幹什麼.
他說多賺點錢.
我最近跟一個弟兄聊天.
他說我們現在.
在做一個實習童工.
他說我不做了.
我說你不做幹什麼.
我出去賺第一桶金.
我出去賺多點錢.
那我們賺多點錢幹什麼.
娶老婆.
那我想想.
他們現在連女朋友都沒有.
他只看到.
多賺點錢.
那我問他你做什麼職業.
飲食業.
那我頭痛了.
飲食業星期天不用上班嗎.
你怎麼上到教會.
死定了.
你不親近神就死定了.
Anna都說了.
出來沒有教會生活就死了.
我們這些人不親近神就死定了.
黑眼的.
我們的眼睛.
經常想著錢.
想著這個想著那個.

$^{561}$原來我們缺少神.
是不行的.
今天.
在疫情的緣故.
可能有很多人在線上.
看直播.
都是慘的.
我覺得都是一個很大的危機.
因為看慣了線上的直播.
就線上的.
經常坐在家裡.
經常坐在家裡看.
就看了很多教會的崇拜都可以.
但最重要的是什麼.
自己教會.
我們基督教的信仰.
不是看著螢幕.
我們要和弟兄姊妹同行.
如果你不同行.
是成長不了的.
沒有激勵沒有激素.
沒有鼓勵的.
你經常看著螢幕那怎麼行呢.
所以.
黑眼的德漢.
有很多東西我們眼睛看的是表面.
但心眼是避開的.
耶穌說給人聽.
你的心眼要打開.
還有.
受壓制的.
得自由.
在耶穌時代.
在教會懸殊.
還有很多特別的人.
他們是受壓制的.
如果你看聖經.
文氏和.
法利賽很多時候都攻擊耶穌.
他說.

$^{601}$說瑞利和罪人.
你們親近這些人.
文氏和法利賽不是這樣說的.
瑞利不是人嗎.
瑞利都需要福音的.
最出名的.
有兩個瑞利信耶穌是哪兩個.
撒該.
和利微.
撒該和利微.
厲害嗎.
撒該是瑞利長.
利微更厲害.
十二使徒之一.
馬太.
耶穌愛他們.
瑞利和罪人.
罪人就是妓女.
妓女來到耶穌面前.
打破真那達.
香膏末主.
這件事震撼了.
人家每個人有錢拿去別的地方.
但他拿了他畢生.
一年賺的錢.
他來到耶穌的腳前.
打破香膏末主.
受壓制的.
得自由.
本來困住.
被人看低.
被人小看.
壓住他.
但神宣告.
我來第四件事.
要讓這些人得自由.
所以.
耶穌一定要做的.
有這四件事.
接下來我講幾個真人真事.

$^{641}$給大家聽.
當然之前講的都是真人真事.
監倉的閘門.
懲教署.
委任了我.
我做懲教署牧師.
我有張證.
可以自由出入懲教所探人.
我進去.
不用警察帶我周圍走.
我想去哪裡探誰.
都可以.
可以周圍去.
有一次我去赤柱監獄.
赤柱監獄.
你們的印象是什麼.
重犯才在那裡.
高度設防的.
重犯.
受監禁在那裡.
有一次我去赤柱.
探訪.
我去到一個屋裡.
第幾組.
其中一組.
他們那組人.
大約有二十幾人.
因為他們分開.
很多個組.
因為他們不想混在一起.
那麼多人混在一起我怎樣處理.
每一日都跟著那組人生活.
他們工作在隔壁房.
住在樓上.
走來走去.
又有一個房間給他們玩.
我那次是星期日去.
大約有二十幾人在一個地方玩.
角色奇色.
有些人看電視.

$^{681}$有些人彈結他.
有些人做運動.
有些人閒聊.
我和其中一個人聊天.
有一個人在後面拍我.
問我認不認得我.
重犯拍我.
我認不認得他.
我轉頭.
看一看這個人.
有個輪廓.
有些少印象.
他告訴我.
他說牧師.
你三十幾年前和我一起坐牢.
三十幾年前.
我最後一次坐牢是三十幾年前.
他說你認不認得我.
我再看清楚.
真的認得.
他告訴我.
牧師.
我因為殺了人.
他是越南人.
他在難民營殺了一個朋友.
所以終身監禁.
今天已經坐了三十年.
終身監禁有否可放棄.
其實沒有.
他今天已經坐了三十年.
二十幾歲坐到五十幾歲.
坐了三十年.
我們三個小時在家.
都不行.
如果入禁閉營.
二十一天都不行.
租一間酒店給你十四天.
你去外國旅行都不行.
都很辛苦.
他坐了三十年.

$^{721}$三十年坐牢.
接著他告訴我.
牧師我很開心.
坐了三十年監禁都開心.
我很開心.
因為我十年前信了耶穌.
還洗禮.
我聽到的時候當然替他開心.
但我奇妙.
因為沒有自由的一個人.
受限制的.
受克制的一個人.
原來他在那個地方.
只有釋放.
罪放.
是讓他開心.
他在那裡.
釋放.
他坐了三十年監禁.
今天告訴我他很開心.
因為罪的釋放.
他信了主.
洗了禮.
很開心.
他內心最重要的開心.
在裡面發出來.
不是外面的燕落.
他最開心.
他在裡面發出來.
他信了主.
最近一次.
我問他.
他告訴我.
有機會見Bot.
提早解釋.
他有機會.
當然見了才知道.
在他人生.
雖然犯過錯事.
在外人來看.

$^{761}$殺死人.
很慘.
不過二十多歲.
太過重手失手.
但神釋放了他們的生命.
這是第一個.
真實的見證.
神願意釋放我們.
無論什麼處境.
今天你面對任何狀況.
只要神在你生命釋放你裡面.
你什麼都能得到自由.
第二個見證.
這裡拍攝一個男人.
背後黑漆漆.
前面有太陽的曙光.
我教會.
其實我也是街坊.
我住在屯門碼頭附近.
所以和華姐經常碰見.
同一個屋村.
華姐一來到.
老院已經叫.
牧師.
整家人知道我是牧師.
老院已經叫.
我教會在我家隔離.
一間叫.
神字會家治中學.
我是神字會屯門堂的牧師.
今天也是一個顧問牧師.
我教會在那裡.
有一個男人.
太太先信耶穌.
太太就回教會.
始終都是.
所有教會都是.
姐妹多過弟兄.
今天一眼看都是.
可能兩夫妻.

$^{801}$兩夫妻.
可能兩夫妻都是太太先信耶穌.
女性的心容易被觸動.
主動很多.
所以女性.
交朋友聊天比較豪爽.
聊得暢談.
男人多數不出聲.
是不是這樣的感覺.
當然有些男人.
很多話題.
多數男人不是很出聲.
藏在心裡面多些.
太太先回教會.
信主很積極.
信主洗禮後.
很積極.
參與教會.
各樣大小事項.
她都很熱心參與.
但男人又怎樣.
男人很少出聲.
但他都很積極.
跟著老婆回教會.
但他不承認她信耶穌.
他回了二十年教會.
都說自己不信耶穌.
很難得.
二十年跟著老婆回教會.
又坐在下面聽道.
除了上班時間難度.
他都要回來.
因為他開輕鐵.
有時星期日要上班.
他除了上班.
全部教會都回到家.
他懂得修理電器.
所以教會壞了一點.
他都不理會.
他不理會.

$^{841}$他就不理會.
教會壞了什麼.
全部都是他修理.
他自動自覺幫我們.
唯一一樣我不信耶穌.
在幾年前.
他承認自己信耶穌.
甚至幾年前.
是我幫他施洗.
其實很難得.
我問他.
事後我跟他聊天.
問他為什麼這麼大的突破.
二十年來都一點承認都沒有.
這麼大的突破.
他說牧師.
我用了二十年來看你教會怎樣.
我看我老婆.
為什麼這麼迷了.
為什麼要搞到這樣.
所以我要看清楚.
我才信.
他看清楚.
還是決定信.
神在他生命裡.
有他的作為.
沒錯.
他在眾人面前是個好好先生.
他真是個好好先生.
他又不吸煙又不喝酒.
又不說髒話.
老老實實地上班.
供養自己.
供養自己家裡養大子女.
全部都是個好好先生.
但好好先生有沒有罪.
聖經說世人都犯了罪.
他有罪性.
他一樣有罪.
神願意.

$^{881}$來到他面前.
他捆綁他釋放他.
所以他承認信耶穌.
第三個真事.
一個媽媽的手.
我們有個家長義工.
她的先生.
在大陸工作.
有外遇.
有第三者.
他很不開心.
極度的不開心.
加上.
他有兩個女兒吸毒.
兩個女兒吸毒.
弄得他頭都暈了.
很慘.
她告訴我.
有一天她買了砒霜.
想煮一鍋糖水.
整家人毒死.
自己死了.
結果不成事.
神出手阻止她.
在她的生命裡.
沒有了很多不好.
很多難處.
很多困難.
很多不幸在她身上發生.
老公有外遇.
兩個女兒吸毒.
我們那裡戒毒.
有一個女兒死了.
死了一個女兒.
今天仍然有一個女兒吸毒.
她所有事都不如意.
沒錯.
她今天到現在.
仍然整個家庭都不如意.
有人還在罪.

$^{921}$但這位家長.
很醒目.
你有沒有聽過.
我決定信耶穌.
是我個人的決定.
我要接納耶穌.
為我個人的救主.
有沒有聽過這句話.
神是你的神.
別人不可以代表你的.
神是我的神.
沒有一個人代表我的神.
我的決定是我的決定.
這個姐妹.
她誠心來到神的面前.
沒錯.
家裡全部都還沒搞定.
她不會說.
等我女兒搞定了.
我才信耶穌.
她不會這樣.
她自己先信耶穌.
還很誠心很熱心地信耶穌.
洗禮加入教會.
熱心參與事奉.
出席祈禱會.
今天是我們家長會的主席.
她不理任何的事.
沒錯.
彷彿很多不幸在我身上.
但我不理那麼多.
因為神是我的神.
今天可能.
在我們當中.
遭遇很多的事.
很多的困難.
甚至有綑綁在我們身上.
一罪的綑綁.
阻止我信耶穌.
很多的事.

$^{961}$我都是想著想著.
我拜什麼那個比較靈.
什麼那個比較靈.
那些是偶像.
怎麼會靈.
真神你都不會拜.
找一個真神.
就像我以前.
我坐了幾次牢.
有幾次都是犯了打劫罪.
我和兄弟.
打劫之前.
去拜關公.
保佑我一間.
拿把刀都厲害一點.
馬上人們拿錢出來.
後來被人抓回CID房.
那裡又有一個關公.
有沒有搞錯.
關公又幫CID又幫我.
亂來的.
偶像.
信就信個厲害一點.
信真神.
可能在我們生命.
遭遇到有這樣的事.
但神願意.
來到我們的生命裡面.
去釋放我們.
去給我們.
有一個新的生命.
今天走出去堂堂正正.
今天在這裡.
我們好好地在耶穌裡面.
去成長.
今天我可以保證.
牧師可以保證.
真神只有一個.
耶穌基督.
我請司琴幫我彈音樂.

$^{1001}$慢慢的音樂.
我們閉上眼睛.
先去為自己.
禱告.
我們閉上眼睛.
去為自己.
主啊我全心擺在你的面前.
主啊.
我願意去信服你.
去倚靠你.
主啊在我們的生活裡面.
可能有很多不如意.
阻滯.
但主啊.
願意緊靠你.
耶穌來到這個世上的目的.
就是將福音.
帶給我們.
這個好消息.
就是永生的信息.
就是最得赦免的信息.
很容易給我們.
很願意給我們.
每一個人.
牧師又去問.
在座呢.
有哪位未信耶穌的.
在這一刻很願意.
我願意相信耶穌.
我願意.
釋放我.
讓我得自由.
在座位裡面試試舉手.
有沒有這樣的朋友.
在座有沒有這樣的朋友.
又或是.
你的生命很被捆綁.
今天主啊.
你給我自由.
釋放我.

$^{1041}$見到一個伯伯舉手.
我們感恩.
好像後面那裡.
也有兩位新來的朋友.
你願意相信耶穌嗎.
願意嗎.
你不要學我.
在監房裡不停地看別人舉不舉手.
你自己的事.
願不願意.
洗了禮.
在座還有哪一位.
我今天有個決定.
願意.
相信耶穌.
我請各位伯伯.
可不可以走出來.
中間.
你陪他走出來.
慢慢走出來.
這麼大年紀.
願意相信耶穌.
不簡單的一件事.
神的福氣.
臨在伯伯身上.
慢慢走出來.
我請老牧師為這位伯伯.
祝福禱告.
除了這位伯伯之外.
沒有人.
是神去感動你.
你是願意這一刻.
去相信耶穌.
接受神的祝福.
如果有的可以一起出來.
\newpage



\section{路加福音 2:36-38}
\label{sec:hNtKC6DPyIM}
\textbf{2021.06.13《使命人生》 路加福音 2\_36-38 ( 和修版)  老耀雄牧師}
\newline
\newline
連結: \href{https://youtube.com/watch?v=hNtKC6DPyIM}{\texttt{https://youtube.com/watch?v=hNtKC6DPyIM}} ~~~~ 語音日期: 2021-06-15
\newline
\newline
\hyperref[sec:11ItqReN_00]{\small{< < < PREV SERMON < < <}}
~
\hyperref[sec:index_chronic]{\small{[返順時目]}}
~
\hyperref[sec:index_scriptual]{\small{[返順卷目]}}
~
\hyperref[sec:T6XL5VlhAWY]{\small{> > > NEXT SERMON > > >}}
\newline
\newline
路加福音 2:36-38
\newline
\begin{longtable}{cl}
\hline
\hline
章節 & 經文 (和合本修訂版)\\
\hline
2:36 & \begin{tabularx}{0.7\textwidth}{X} 又有位女先知，名叫亞拿，是亞設支派法內力的女兒，年紀已經老邁，從童女出嫁，同丈夫住了七年， \end{tabularx} \\ \\ \relax
2:37 & \begin{tabularx}{0.7\textwidth}{X} 就寡居了，現在已經八十四歲。她不離開聖殿，禁食祈求，晝夜事奉神。 \end{tabularx} \\ \\ \relax
2:38 & \begin{tabularx}{0.7\textwidth}{X} 正當那時，她進前來感謝神，對一切盼望耶路撒冷得救贖的人講論這孩子的事。 \end{tabularx} \\ \\
[1ex]
\hline
\hline
\end{longtable}
$^{1}$吾人間係弟兄姊妹平安.
如果有一日你去見一份工.
個老闆同你講好jor 所有聘任ge 條件.
大家都好滿意啦.
咁之後呢你第一日返工.
老闆同你講.
你有一個好重要ge 任務.
咩任務呀.
你就係office裡面呢.
等一個好重要ge 人物.
當呢個好重要ge 人物黎到ge 時候.
你就將呢個人物呢.
介紹俾每一個人識啦.
咁於是呢.
你就每一日都期待呢個好重要ge 人物出現.
但呢個好重要ge 人一路都唔出現.
咁你就問老闆.
究竟佢幾時黎架.
我等jor 好耐.
個老闆就話.
你等啦.
佢實出現架.
你一等就等jor 十年.
十年又十年.
人生有幾多個十年呢.
不過你一等就等jor 六個十年啦.
呢個重要ge 人物仲未出現.
係咁ge 情況.
你咪問你自己.
你心裡面諗D 咩呢.
原來等待都係一種學問黎架喎.
從信仰角度等.
即係從信仰角度去睇.
等待係一個屬靈ge 操練黎架.
聖經就一個咁ge 人物.
佢等jor 一個人差唔多六十年啦.
咁我地就透過路加福音.
二章三十六至三十八節.
從阿拿ge 生命學習信仰ge 兩個態度.
開始ge 時候我地有一個祈禱.

$^{41}$創天造地ge 主.
我地知道你創造我地每一個.
都有你美好ge 心意同計劃.
但願我地每一個係成聖ge 旅程入面.
我地能夠明白你對我地gwo 種獨特ge 心意.
我地又願意行係你ge 心意裡面.
能夠與你同行.
係我地ge 恩典同福氣.
我地ge 祈禱奉靠主耶穌.
基督聖名祈求.
阿們.
信仰第一個ge 態度.
就係忠於自己ge 使命.
經文一開始就提及阿拿ge 身份同出處.
阿拿屬於阿徹之派.
阿徹係雅各ge 第八個仔.
亦都係拉潔ge 使女.
色柏為雅各所生ge 第二個仔.
摩西係臨終ge 時候.
曾經為以色列十二個之派祝福.
而佢為阿徹之派祝福ge 時候.
佢加多jor 一個應許.
佢話你ge 日子如何.
你ge 力量也必如何.
呢句說話可以理解為.
你點樣用你ge 時間.
係同你ge 生命力成正比ge .
你點樣用你ge 時間.
係同你ge 生命力成正比.
阿拿ge 一生.
究竟點樣去用佢ge 時間.
我地一齊再睇落去.
阿拿係一個寡婦.
經文話佢年界八十四歲.
猶太人一般都唔會太遲結婚.
當阿拿十六歲結婚.
經過七年ge 婚姻生活.
就係佢二十三歲ge 時候.
佢就離開jor 呢個世界.
咁你識計數.

$^{81}$計一計阿拿獨身jor 六十一年.
如果我地手頭上有和合本.
聖經ge 話呢.
佢會多jor 一個ge .
即係呢.
全息ge 三考.
佢話呢句說話可以將.
亦為呢阿拿寡居jor 七八十四年.
即係話唔係阿拿八十四歲.
而係話佢寡居jor 八十四年.
如果佢寡居jor 八十四年ge 話呢.
咁佢就差唔多一百歲.
不過無論佢係寡居jor 八十四年好.
亦或佢只係八十四歲好.
六十一年同八十一年.
都係一個好長好長ge 日子.
經文話阿拿係一個女先知.
先知係聖經ge 解釋.
就係神ge 代言人.
佢能夠領受從神而來ge 說話.
然後將神ge 心意.
轉述比神ge 子民聽.
所以阿拿作為一個女先知.
為神發聲.
係佢ge 使命黎ge .
係聖經ge 記載裡面.
我地知道好多舊約ge 先知.
佢地都行遍整個猶太地.
我地數得到ge 有以利亞.
有以利沙.
佢地四圍去.
將神ge 話語講比百姓聽.
更加有D 先知唔單單係猶太地.
仲去到外邦異族.
去先講神ge 話語ge .
記唔記得邊個呀.
若拿.
若拿就走到去亞述尼尼尾城.
去先講神ge 說話.
就算係新約.

$^{121}$同阿拿同時期出現ge 施洗約漢.
佢都係四圍去ge .
去到曠野先講神ge 道.
不過經文就話阿拿.
只係留jor 係聖殿裡面服侍.
佢冇離開過聖殿.
所以阿拿所領受ge 意象.
唔係四圍去.
而係聖殿裡面等待.
等待咩呀.
等待神所應許的彌賽亞出現.
然之後.
佢先至可以發聲.
為彌賽亞作見證.
一日彌賽亞未出現.
佢都係度等.
一等就六十幾年.
或者八十幾年.
經文就冇話比我知.
阿拿係幾多歲領受呢個意象.
亦都冇話比我聽佢究竟.
領受jor 意象之後等jor 幾耐.
但係阿拿好清楚知道自己ge 使命.
就係係聖殿裡面等候彌賽亞出現.
然之後為佢作見證.
而唔係出去四圍咁去搵彌賽亞.
然之後為佢護駕保航.
我地知道孫子ge 身份係舊約好尊貴.
好受百姓ge 擁戴.
先至每到一個地方.
都會有猶太人作出對佢ge 款待.
因為先知就係神ge 代言人.
佢ge 說話就係神ge 說話.
不過阿拿ge 使命與別不同.
佢要為彌賽亞作見證.
要彌賽亞出現jor .
佢先至可以傳神ge 話語.
佢唔能夠四圍去.
佢只能夠守係聖殿裡面等待.
等待一個唔知幾時.

$^{161}$先至會出現ge 彌賽亞.
去入黎聖殿.
不過呢個亦都係.
呢個先知等待彌賽亞ge 唯一ge 方法.
阿拿就用jor 佢一生ge 時間.
好忠心咁樣去實踐.
呢個神交託比佢ge 使命.
所以經文話.
正當那時.
就係直到約翰同瑪利亞.
帶同耶穌到達聖殿ge 時候.
阿拿當時第一個反應.
唔係即刻去實踐使命.
係同人講耶穌就係彌賽亞.
實實在在係阿拿內心裡面ge 情感ge 流露.
因為使命將要完成.
因為一個漫長ge 等待終於結束.
因為以色列ge 百姓ge 盼望終於黎到.
係因為個擔子可以卸下來.
阿拿雖然未見過耶穌.
但係佢亦都係最認識耶穌ge 其中一個.
佢知道耶穌ge 身份.
佢知道耶穌黎到呢個世界ge 目的.
佢知道耶穌就係以色列百姓ge 盼望.
因為整個事奉ge 生命.
就係為耶穌作見證.
證明耶穌就係大家所等待gwo 個彌賽亞.
曾經睇過一本書.
叫做天地一沙鷗.
有冇人睇過呢本書呀.
作者係一個空軍上校.
尼察巴哈.
書裡面講jor 一隻海鷗.
叫做鄂納山.
呢隻小小ge 海鷗.
佢唔甘心做一隻係沙灘上搶食細魚.
或者跟住gwo D 漁船追逐麵包碎ge 海鷗.
佢相信生活ge 意義.
絕對唔單單係要填飽個肚.
佢要追求更高ge 生活理想.

$^{201}$為jor 追求呢個理想達到目標.
佢忍受嘲諷.
忍受痛苦同埋孤獨.
即使被趕出去呢個海鷗群.
佢都唔放棄去學習更高境界ge 飛行技巧.
後來成功jor .
成功jor 但係冇忘記到佢ge 同類.
所以佢帶住一個使命.
將佢ge 技術傳遞比下一代.
讓其他海鷗都體會到.
其實我地可以有一個自由無限ge 觀念.
唔係只係追逐係個漁船後面gwo D 麵包碎.
我地能夠超越海鷗ge 極限.
成為一個海鷗ge 傳奇.
使命讓到呢隻小海鷗企得高望得遠.
去搵到生命ge 終極意義.
生命可以不平凡.
關鍵係在於.
我地搵唔搵到gwo 個屬於我地ge 使命.
我地能唔能夠搵到神比我地ge 使命.
然後將我地終其一生去實踐呢個使命.
各位紅映街ge 弟兄姊妹.
我地每一個雖然係一個好平凡ge 信徒.
我地都唔係一個先知.
我地冇呢個身份.
但係唔代表我地冇使命.
耶穌係馬太福音28章19至20節.
向所有ge 門徒咁講.
所以你地要去使萬民作為我門徒.
奉父子聖靈ge 名.
給他們施洗.
凡我所吩咐你地ge .
你都教導佢地遵守.
看下我天天與你地同在.
直到世代ge 終結.
只要我地係信徒.
我地就有使命.
將福音講傳.
係我地其中一個ge 使命.
我地已經領受jor 呢個救恩.

$^{241}$可唔可以好似呢個海鷗.
鱷立山一樣帶住一個使命.
將救恩傳比gwo D 仲未得救ge 朋友呢.
我地唔使好似阿娜咁樣.
要等60年.
我地唔係好似阿娜咁樣.
要守係教會裡面ge .
我地可以真係走出去.
去到我地所認識ge 群體.
用我地ge 言語.
我地ge 行為.
我地ge 生活.
去見證我地所相信ge 主耶穌.
呢個係咪你ge 使命.
信仰ge 第二個態度.
我地需要生命ge 操練.
去建立同神有更深ge 關係.
我地知道阿娜由領受使命開始.
佢就放棄jor 一切.
佢用jor 一生ge 時間去實踐呢個使命.
阿娜能夠咁堅定咁去執行使命.
佢有咩秘訣.
61年.
甚至80年.
去持守一個使命.
經文話阿娜禁食祈求.
就夜事奉神.
我地知道食物係我地身體所需要ge 其中一種物質.
食物吸收我地身體吸收唔到呢種食物.
所釋放出黎ge 能量ge 話.
嚴重ge 話呢我地會導致死亡.
所以基梗都係30.
基梗90就唔係個個都得架啦.
我地需要食物.
去保持我地ge 身體ge 能力.
但禁食ge 操練.
就可以將我地放低jor 我地身體ge 需要.
食物係必須ge .
但我地將佢降低少少.
縱使呢種需要係咁重要.

$^{281}$重要到生死攸關.
但阿娜要係信仰裡面去不斷去問自己.
究竟肉體ge 需要重要D .
因為神ge 心意重要D .
透過禁食ge 操練.
讓阿娜更加專注.
更加清心去測驗神ge 心意.
最後阿娜願意將神ge 心意.
放係佢生命ge 首位.
使佢有能力承擔呢個使命.
又忠於呢個使命.
我地每一個人都有好多ge 需要.
呢D 需要可以好重要.
但可唔可以將你ge 使命.
ge 優先次序排高D 呢.
阿娜六十幾年ge 事奉生命.
一定有好多ge 引誘同埋試探.
你話咩係引誘同試探呀.
你諗下23歲就冇jor 丈夫.
年青ge 時候會唔會有情慾同婚姻ge 試煉呢.
有架.
越等待越唔耐煩.
會唔會有懷疑神ge 試煉呢.
十年又十年仲要等幾耐呀.
外面世界咁大.
係咪應該出去闖一闖呢.
我係先知黎架.
我受人尊敬架.
我行出去受人接待架.
會唔會有呢種私慾ge 試煉呢.
不過阿娜將神ge 心意放係首位.
呢個就係佢抵抗引誘.
剋性試探ge 致勝關鍵.
我心尊主為大.
我地ge 心能唔能夠尊主為大.
主係我地一切.
我為主可以放棄一切.
或者我地唔好咁偉大.
唔好話放棄一切.
我地可唔可以將我地ge 慾念.

$^{321}$我所追求ge ye 稍微放輕一D .
弟兄姊妹.
我地察唔察覺到.
原來世界不斷向我地發出唔同ge 挑戰同試探.
物質ge 生活.
同信仰ge 實踐.
令到我地感受到一種好大好大ge 張力.
有多初訓班ge 同學.
好多時都問我一個問題.
星期要返黎崇拜.
我冇jor Family Day.
我平時呢個時間我同屋企人飲緊茶.
我平時呢段時間同我ge 配偶行緊山.
你過每個禮拜都返黎崇拜.
可唔可以唔返.
引誘挑戰.
有D 年青人問我.
老闆話升我職.
升職好.
不過升jor 職之後.
星期日就冇得返教會.
唔升職.
星期日就可以返.
升jor 職就冇啦.
咁我升唔升職呀.
張力.
係咪呀.
有邊個唔想自己升職.
有邊個唔想自己加人工.
但係升職加人工.
你冇jor 一個崇拜心.
社會ge 文化.
同信仰ge 價值觀南轅北轍.
咩叫社會文化呀.
我記得呢幾年裡面好輕講一個贏在起跑線.
如果你有個細路.
要讀幼稚園.
你會報幾多間幼稚園.
五間.
少得滯呀.

$^{361}$起碼都十間.
係咩呀.
執書行頭.
贏在起跑線.
我地唔係要交托ge 咩.
我地信仰gwo 個交托去jor 邊.
點解要報五間幼稚園報十間幼稚園.
我地驚輸呀.
社會ge 文化同我地有衝擊.
不過如果我地知道.
有D ye 俾我地肉體ge 需要.
有D ye 俾我地物質ge 享受.
更加重要同有永恆價值ge 時候.
我地就可以作出一個取捨.
今生一切都可以無.
我地冇永恆ge 盼望.
透過淑寧ge 操練.
讓我地知道生命咩係叫最重要.
如果肉體ge 需要都可以放低.
咁物質ge 需要就更加容易放低.
我地ge 信仰生命需要我地去操練.
提供一個反思ge 平台.
比我地能夠更加清心咁樣去.
察驗神ge 心意.
其實有陣時我地唔係唔知道神ge 心意.
不過我地軟弱呀嘛.
我地唔想付呢個代價呀嘛.
經文提到阿拿除jor 禁食之外.
仲有祈禱.
祈禱讓阿拿同神ge 關係更加親密.
祈禱係信徒同神溝通一個最直接ge 方法.
透過祈禱我地可以認識神.
亦可以認識自己.
我地越認識神.
就越覺得神gwo 個偉大.
神gwo 個無限.
我地越認識自己.
就覺得自己有限制.
自己有多ge 軟弱.
透過祈禱阿拿將自己ge 軟弱交託俾神.

$^{401}$讓神去兼顧佢ge 信心.
透過祈禱阿拿可以將自己ge 限制同神分享.
讓神加添佢ge 能力.
透過祈禱阿拿同神ge 距離拉近jor .
並且同神建立一個深厚ge 關係.
呢種親密ge 關係能夠促使阿拿能夠承擔使命.
又忠於使命.
今日我要向大家問一問大家.
你今日同神ge 關係係一個咩ge 光景.
今日你同神ge 關係係咪一種可有可無ge 關係.
有又得.
無亦都得.
阿拿唔係天生出黎就有咁ge 耐性同埋毅力.
61年81年唔係生出黎就有.
佢要等待彌賽亞出現已經等jor 60幾年.
如果等待ge 期間.
無神ge 安慰臨到.
無神ge 鼓勵臨到.
無神加能賜力俾佢.
以及增添佢ge 信心.
我真係唔相信以人ge 能力係可以做得到.
我地ge 人ge 能力唔好高估.
可能我地一年都守唔到.
如果我地無神與我地同行.
一年我地都頂唔順.
弟兄姊妹.
如果祈禱能夠讓我地同神ge 關係拉近.
使我地同一位又無限又大能ge 神去連結.
如果我地ge 信仰生命能夠體驗過.
我地可以同永恆ge 主結連ge 話.
我地仲有咩ye 係可以唔願放棄ge 呢.
個關鍵係.
我有冇體驗過.
我有冇體驗過我同神ge 親密關係.
如果體驗過ge 話.
我地就可以作出選擇.
但係如果我地從來冇體驗過ge 話.
我地就唔識選擇.
我地有冇錯過神俾我地係禱告裡面ge 恩典.
每個信徒都可以禱告.

$^{441}$因為主耶穌已經成為jor 中堡.
我地透過主耶穌我地同神ge 關係復和.
我地同神ge 關係建立親密ge 關係.
呢個係我地作為信徒ge 權利.
我地有冇錯過jor 我地呢個權利.
抑或我地根本唔知.
我地原來可以係禱告裡面可以操練.
可以係個禱告裡面操練到我地同神gwo 個關係可以咁親密.
另外經文雖然冇直接提到阿娜經歷緊一個等待ge 操練.
但事實上佢用jor 一生ge 時間去學習呢個功課.
等待ge 操練促使佢同耶和華有一個親密ge 關係.
因為佢要確定.
呢個就係佢一生要承擔ge 使命.
每個人都有.
你搵唔搵到.
如何搵到.
搵到之後點去承擔.
呢個係我地信仰ge 成聖之旅.
等待ge 操練加深jor 去認定自己ge 使命.
因為一定要確信.
呢個係唯一實踐使命ge 方法.
唔係四圍去.
唔係周圍講.
唔係周圍問.
等待係你唯一ge 方法.
要明白使命ge 內容.
以及點去實踐.
唯一ge 方法.
就係同神有一個好親密ge 關係.
呢個係唯一ge 法門.
你問唔到個牧者.
因為牧者唔係創造你.
神先至係創造你gwo 個.
你要直接搵神.
問返神.
你創造我ge 目的是咩.
你擺係我生命裡面ge 恩賜係咩.
你要我為實踐ge 信仰.
生命使命係咩.
唯有神先至可以回應呢個答案.

$^{481}$弟兄姊妹.
如果我地唔願意係神面前等候.
亦都唔去等候.
我地只會一條路.
就係按住自己ge 心意走.
你永遠都聽唔到神對我地ge 說話同心意.
如果我地ge 信仰世界只係自己得一個.
我地ge 信仰只係按住自己ge 心意而行ge 話.
我地叫做自說自話.
咁信仰對我地有咩意義.
就先聖餐禮裡面我地講.
我地係浸禮ge 時候我地咪確認.
神係我地生命ge 主咩.
邊個先係我地生命ge 主.
我地唔可能按住自己ge 心意行.
我地一定係按住主ge 心意行.
我呢度再介紹多一個人物比大家識.
可能我之前講過.
如果有睇過烈火戰車ge 話一定識佢.
佢就係蘇格蘭一個運動員.
叫Eric Henry.
中文譯名叫李愛瑞或者李達.
戲裡面講到係1924年.
巴黎舉辦ge 夏季奧林匹克運動會裡面.
Eric因為星期日要參加主日崇拜ge 緣故.
參加主日崇拜.
就放棄jor 佢最有機會lor 金牌ge 100米短跑比賽.
你話Gary係咪呀.
主日崇拜有幾咁大件事呀.
lor 奧運金牌喎.
lor 奧運金牌唔重要.
對於Eric黎講.
返崇拜先係最重要.
佢就係唔參加呢個100米短跑比賽.
而選擇jor 一個佢唔擅長ge 400米.
因為呢個400米比賽唔係主日裡面去進行.
咁個個人都覺得.
你都唔係跑400米呀.
你跑100米咋嘛.
你都係輸架啦.

$^{521}$就係一片咁唔好ge 情況之下呢.
佢最後lor jor 個400米金牌.
呢套戲呢.
就以佢lor 到呢個400米金牌ge 奧運金牌ge 結束.
戲就完jor 啦.
不過.
佢真正ge 人生.
佢真正ge 事奉人生先至arm arm 開始.
運動會之後.
當然國家好想佢拎jor 呢個金牌之後為國家服侍啦.
佢有名有利啦.
不過佢冇咁做.
Eric跟jor 佢爸爸.
去jor 中國ge 天津宣教.
開始服侍中國人.
嘩.
天同地呀.
係咪呀.
一個奧運金牌係國家幾咁受尊重呀.
你走jor 去中國天津冇人識你架.
1937年中日戰爭爆發.
1943年Eric比日本.
即係日本軍隊俘虜.
成為戰俘.
關jor 係山東ge 集中營.
係集中營ge 期間呢佢冇停止事奉.
繼續為gwo D 集中營ge 細路仔.
開始主日學.
係個期間.
日本同英國曾經協議.
有一個交換戰俘ge 計劃.
Eric就係呢個交換名單裡面ge 頭一批.
如果你係Eric.
你開唔開心呀.
可以重見天日啦.
交換戰俘.
你可以返返英國.
不過Eric放棄jor .
放棄jor 呢個從惡自由ge 機會.
佢讓jor 位比人.

$^{561}$點解呀.
佢認為係個集中營裡面.
有更多人需要去照顧.
gwo D 人需要去照顧.
所以從不從惡自為對佢黎講唔重要.
重要ge 就係佢ge 使命係乜.
佢個使命就係服侍中國人.
我去到天津做宣教士就係服侍中國人.
係呢個關口我都係要服侍.
而唔係呢離開呢個國家.
Ericge 使命唔係係運動場上lor 好多金牌.
lor 金牌唔能夠滿足佢對上帝事奉ge 心.
當生命遇到逆境ge 時候.
能夠離開集中營.
能夠重獲自由.
都唔係佢覺得生命裡面最重要ge .
生命最重要ge 照顧佢身邊ge 弱小ge 兒童.
讓佢感受到神ge 大愛.
唔放棄上帝交托比佢ge 使命.
呢個先至係Eric生命中最重要ge 核心價值.
唔放棄神交比我ge 使命.
我永遠都唔放棄.
命都冇啦唔重要.
佢用佢宣教ge 生命黎榮耀上帝.
咁你會話Gary佢咁勁je .
我冇佢咁勁.
係真ge .
我地每一個都冇佢咁忠心.
但係願唔願意將我地ge 生命交係主ge 手裡面.
讓主加能賜力比我地.
讓主ge 恩典托住我地.
今日ge 講題係使命人生.
人生在世總有D 目標要完成.
係咪呀.
我記得早幾年有D 人講話.
買唔到層樓.
不如做鹹魚好過啦.
佢ge 使命咪就買樓囉.
冇壞ge 係咪先啦.
冇信仰ge 人人生目標可以係升官發財.

$^{601}$事業有成.
世ge 目標可以係買一層樓.
安居樂業.
家庭ge 目標可以望子成龍.
望女成鳳.
婚姻美滿.
呢D 都可以促進我地ge 動力.
有盼望咁去生活.
冇壞ge .
不過如果我地帶起一副信仰眼鏡去睇.
一副信仰眼鏡去睇.
我地就一定要發援到.
呢D 目標只能夠滿足我地一時.
卻唔能夠滿足我地心靈深處gwo 個內心個渴求.
呢個渴求需要我地超越今生直到永恆.
基督信仰話比我地知.
唯有耶穌先至能夠滿足到我地呢個新心靈gwo 個需要.
係耶穌基督裡面.
我地先至有一個永恆價值ge 使命.
將與神和好的福音傳開去.
將我地能夠被釋放ge 生命與人分享.
使到更多人能夠體驗神ge 愛.
同享有永恆ge 生命.
如果上個禮拜你有返過黎崇拜.
你有聽過Anna姐妹ge 見證.
你一定發覺到.
佢ge 生命真係好精彩.
點樣由毒海裡面能夠被釋放出黎.
我地唔使好似佢咁樣.
我地唔使好似佢咁樣走去吸毒.
然後去經歷呢樣ye .
我地每一個人或多或少總有D ye 係被克制.
我地或多或少總有D ye 係未釋放.
Anna話比我地聽.
我地主gwo 個大能.
我地主gwo 個奇妙.
使命透過與人分享.
經過實踐產生ge 意義.
Anna就係一個好例子.
佢經過分享.

$^{641}$佢實踐佢ge 信仰.
產生jor 個意義.
呢個分享成為我地心裡面一個好大ge 鼓勵.
分享生命比獨自享受更有意義同價值.
我地每一個人都有得救見證.
你ge 得救見證埋係個心裡面.
唔肯同人分享有咩意義.
冇意義.
唯有你ge 生命.
你ge 得救生命.
你gwo 個與神同行ge 生命.
你gwo 個能夠可以.
即係呢黑性gwo 個捆綁ge 生命.
同人分享.
人就得到鼓勵同安慰.
雖然我地要領受使命.
需要實踐.
但我地心知道係實踐ge 路上.
我地經常跌倒.
甚至跌倒jor 之後起唔返身.
係咪.
我地好多時候初信耶穌gwo 陣時候滿腔熱誠.
跟住沖完一輪之後.
借jor .
借jor 就借jor .
潛jor 水就潛jor 水.
潛水艇落jor 去河底.
都唔知何年何日先升返上黎.
唔知.
我地要有個屬靈操練.
要建立一種同神親密ge 關係.
可以令到我地生命能夠強壯D .
有能力有勇氣去接受挑戰同試探.
屬靈操練亦都俾我地呢.
遇上天國ge 一D 能道.
與主同在.
唔係只係末世裡面天國出現.
耶穌黎到啦然之後先出現.
係今生.
此時此刻我地都可以.

$^{681}$好似係天堂一樣.
享受與主同在.
gwo 種ge 深刻經歷同體驗.
其實屬靈操練唔單單只係禁食祈禱.
仲有好多.
有D 操練可以幫助我地面對物質ge 牽引.
例如十一奉獻.
有D 操練係想使我地呢.
係使用金錢ge 主權交返比主.
有D 操練係讓我地呢.
將個時間交返比主.
有D 操練又唔涉及物質ge .
譬如返祈禱會去操練我ge 禱告.
一班人同心合意ge 祈禱.
肯定更加能夠加強我地係禱告上ge 信心.
所以我鼓勵大家.
你能夠有時間ge .
星期三晚返返去祈禱會.
一班人同心合意.
係禱告上面去終身ge 事奉.
等於只係我地願唔願意過一個.
有永恆價值ge 使命人生.
係我地ge 屬靈生命上.
有冇一D 軟弱位.
一篤我就死得架啦.
一篤就死得.
就需要我地持續ge 操練.
讓呢D 軟弱變得剛強D .
好叫我地能夠承擔使命又忠於使命.
我就好想大家安靜係座位裡面.
瞇埋眼.
將剛才我聽到ge 訊息沉澱一下.
然後我帶大家有一個默想.
我深信神賜給信徒ge 使命.
並唔限於昔日ge 阿那.
亦包括一切在座相信主耶穌ge 信徒.
換言之亦都包括你同我.
包括坐係呢個聖殿裡面.
每一個屬主ge 子民.
不過主耶穌.

$^{721}$我地每一個都好軟弱ge .
我地有時會感到冇能力去承擔.
主耶穌我知道你就係我地身邊.
你願意成為我地每一個人ge 幫助.
主耶穌我知道.
當我地願意承擔ge 時候.
你會加能賜力比我地.
並且你會陪我地.
一齊行呢條信仰ge 旅程.
今日我地每一個都確信.
聖靈已經內在我地心裡面.
今日我地都確信.
主已經與我同在.
成為我地隨時ge 幫助.
今日此時此刻.
我地都需要承擔.
主耶穌託付比我地ge 使命.
我而家就奉主耶穌ge 名.
向在座ge 每一個弟兄姊妹作出邀請.
你願意向你未信主ge 親戚朋友.
見證主耶穌嗎?.
如果你願意ge .
請你舉一舉手.
讓耶穌知道你的心意.
你願意係未來日子去付代價.
去同耶穌建立一個親密ge 關係嗎?.
如果你願意ge .
請你舉一舉手.
我會為你祝福.
各位願意一生跟從耶穌ge 弟兄姊妹.
你願意用你的生命去見證主耶穌嗎?.
你願意同耶穌有親密關係嗎?.
請你在座位舉一舉手.
讓耶穌知道你的心意.
好,請放下.
我地一齊去祈禱.
親愛的恩主.
你知道我地每一個ge 心意.
今日有好多弟兄姊妹願意為你作見證.
每一個都願意係你面前毀身.

$^{761}$佢地舉手立志.
願意係生命上流露出屬你ge 香氣.
我地好想其他人.
唔認識你ge 人.
可以係我地ge 生命裡面.
見到主耶穌你自己.
主呀求你兼顧我地ge 心.
求你使用我地每一個ge 生命.
叫我地與你同供.
成為一個福音ge 使者.
親愛的恩主.
有好多弟兄姊妹願意親近你.
我地想贖領生命有成長.
我地想飲你ge 活水.
食你ge 靈量.
求你進入我地ge 心靈深處.
滋潤我地ge 渴求.
使我地身心靈都同樣得到飽足.
親愛的恩主.
求你閱立每一個真誠渴慕你ge 心.
讓我地信仰ge 生命.
同你所賜ge 恩典相稱.
讓你名被高舉.
我地祈求感恩.
奉主耶穌聖名祈求.
阿們.
願主賜福俾大家.
\newpage



\section{詩篇 139:13-18}
\label{sec:T6XL5VlhAWY}
\textbf{2021.06.20《無可取替的愛》詩篇 139\_13-18( 和修版)  曾敏姑娘}
\newline
\newline
連結: \href{https://youtube.com/watch?v=T6XL5VlhAWY}{\texttt{https://youtube.com/watch?v=T6XL5VlhAWY}} ~~~~ 語音日期: 2021-06-21
\newline
\newline
\hyperref[sec:hNtKC6DPyIM]{\small{< < < PREV SERMON < < <}}
~
\hyperref[sec:index_chronic]{\small{[返順時目]}}
~
\hyperref[sec:index_scriptual]{\small{[返順卷目]}}
~
\hyperref[sec:M84h3HizR5w]{\small{> > > NEXT SERMON > > >}}
\newline
\newline
詩篇 139:13-18
\newline
\begin{longtable}{cl}
\hline
\hline
章節 & 經文 (和合本修訂版)\\
\hline
139:13 & \begin{tabularx}{0.7\textwidth}{X} 我的肺腑是你所造的，我在母腹中，你已編織我。 \end{tabularx} \\ \\ \relax
139:14 & \begin{tabularx}{0.7\textwidth}{X} 我要稱謝你，因我受造奇妙可畏，你的作為奇妙，這是我心深知道的。 \end{tabularx} \\ \\ \relax
139:15 & \begin{tabularx}{0.7\textwidth}{X} 我在暗中受造，在地的深處被塑造；那時，我的形體並不向你隱藏。 \end{tabularx} \\ \\ \relax
139:16 & \begin{tabularx}{0.7\textwidth}{X} 我未成形的體質，你的眼早已看見了；你所定的日子，我尚未度一日，都在你的冊子寫上了。 \end{tabularx} \\ \\ \relax
139:17 & \begin{tabularx}{0.7\textwidth}{X} 神啊，你的意念向我何等寶貴！其數何等眾多！ \end{tabularx} \\ \\ \relax
139:18 & \begin{tabularx}{0.7\textwidth}{X} 我若數點，比海沙更多；我睡醒的時候，仍和你同在。 \end{tabularx} \\ \\
[1ex]
\hline
\hline
\end{longtable}
$^{1}$各位弟兄姊妹平安.
讓我們可以在神的家裡.
去敬拜祂.
去認識祂.
真是很多無比.
今天是父親節.
我相信有不少家庭會把握今天的機會.
和爸爸一起慶祝.
送禮物.
一起吃飯.
在大部分的家庭裡.
爸爸是很重要的.
是主要的經濟支柱.
供應一家人的需要.
爸爸又是大力士.
是搬搬抬抬.
更加是英雄.
是保護一家人的安全.
爸爸的愛對於孩子來說.
是非常寶貴.
有了爸爸的愛.
孩子才可以健康地成長.
地面上的爸爸.
令我想起天上的爸爸.
他的愛更加寶貴.
甚至可以說是無可取替.
他的愛勝過地面上的爸爸.
沒有東西可以取替他的愛.
今天我們來看看.
天父的愛是怎樣的寶貴.
我們會看詩篇139篇13到18節.
我們開始的時候再一同禱告.
讓我們今天聚集在一起.
去認識你.
天父求你開我們的眼睛.
讓我們認識你是怎樣的神.
讓我們明白你的愛.
願聖靈充滿我們每一個的心.
去感動我們.
賜我們智慧悟性去明白.

$^{41}$也讓我們今天懂得去回應你.
願在裡面天父你顯出你的榮耀.
奉耶穌基督的名字祈禱.
阿們.
大家看到這幅圖.
這個人就是大衛.
稍稍介紹詩篇139篇.
這個詩篇是大衛初期作王的時候.
所寫的作品.
大衛是以色列第二任的君王.
他能征善戰.
很會打仗.
但他打仗不是靠自己.
比較多的時候都是靠著上帝的恩典去打敗敵人.
與此同時他也是一個非常感性的詩人.
寫了很多詩篇.
表達自己對上帝的看法和感受.
在詩篇139篇裡.
他讚美神三件事情.
就是神是無所不知.
無所不在.
和無所不能的神.
13到18節.
今天我們要看的經文.
是強調神是創造我們的神.
是無所不能的神.
他的愛是無可取替的.
我們從三方面去看他的愛.
第一方面.
神親手創造我們.
賦予我們無比珍貴的價值.
13到14節這樣說.
13節說.
我的肺腑是你所造的.
我是無服眾.
你已偏激我.
14節.
我要稱謝你.
因我受造奇妙可畏.
你的作為奇妙.

$^{81}$這是我心心知道的.
肺腑是人身體裡很重要的器官.
例如什麼呢?.
神臟.
其實猶太人在這裡提到肺腑.
說到神臟.
是用來比喻我們內心的情緒和感受.
原來不單是我們人的身體.
就連我們人的情緒和感受.
都是上帝所造的.
換言之.
我們整個人.
由身體到情緒到感受.
都是神所造的.
神在哪裡造我們呢?.
13節這樣說.
我在無服眾.
你已偏激我.
15節說.
我在暗中受造.
在地的深處被塑造.
地的深處是指很隱密的地方.
是很隱密的地方.
我想先上去.
(上鏡時).
不要緊.
這裡說地的深處是隱密的地方.
再進一步說子宮的深處.
偏激是說連結.
當我們還在媽媽的肚子裡的時候.
神就將我們不同的部分連結在一起.
我們每個人雖然是經過父母的結合.
經過媽媽的肚子來到這個世界.
但實際來說.
仍然是出於上帝的工作.
完全是上帝的工作.
以至我們才可以出現.
神製造我們每一個生命.
所以我們都是獨一無二的.
每一個可以說是手工製造的.

$^{121}$獨立製作的.
不是像工廠一樣.
有一個模具.
吸完一個又一個.
工廠的模具是逐個吸出來的.
所以全部都是一樣的.
但神是逐個逐個製造我們.
所以我們獨一無二.
剛才很快地飛了一個powerpoint.
下面看到2017年的時候.
當時的全球人口已經達到75億.
我相信今天是不止的.
我還沒有最新的數字.
但在當時的75億裡面.
沒有一個人的DNA是一樣的.
就算在不同的家庭裡面.
有些孿生的姐妹,兄弟.
他們都是不一樣的.
所以我們每一個人.
是神特別的,獨立的製造出來的.
我們不單止獨一無二.
還可以這樣說.
我們奇妙無比.
所以詩人才這樣說.
我受造奇妙可畏.
我們的身體是怎樣奇妙可畏呢?.
我找了一些資料.
大家不妨聽聽.
我們平時比較少從這個層面去看.
例如是甚麼呢?.
我們眼睛原來.
如果當中有讀生物的弟兄姊妹可能會知道.
眼睛有六條肌肉.
我們眼睛的六條肌肉是要互相協調.
才可以幫助你的眼球.
轉到一個正確的目標上.
原來這個運作不簡單.
我們平時很隨意.
動得你覺得沒甚麼特別.
原來是很特別的.

$^{161}$還有你的瞳孔.
你要調較好.
一個適當的光圈.
這樣看東西才是剛剛好的.
最直接的就是.
我們要留意鼻孔是向上還是向下呢?.
還用說嗎?.
當然是向下.
有人說.
向下才好.
下雨的時候.
雨水才不會進入.
這是說笑的.
都是事實上的.
但還有很多好處.
我們的器官.
有不同的數目.
每一樣東西剛剛好.
我們的腦非常複雜.
腦控制了我們整個身體的運作.
腦壞了.
整個身體就不能正常運作.
這些研究.
我想是這麼多年來的科學家.
生物學家一直研究的時候.
只能夠認識其中一部分.
我和你平時都比較思考這些問題.
是更加難明白.
甚至感覺很奇妙.
做我們很特別.
上帝愛我們.
不單止賦予我們生命.
也賦予我們很寶貴的價值.
我們有沒有從神的角度去愛自己.
和愛其他人呢?.
這個世界的價值觀是扭曲的.
為什麼我說扭曲呢?.
這個世界覺得.
你一定要有一個很漂亮的外表.
很聰明的頭腦.

$^{201}$你要健康,口才,金錢,名譽,地位.
你很有權利,很有成就.
你才是值得愛的.
還記得我在中學的時候.
我曾經因為自己的身形.
那時候已經是一個肥妹仔.
在青春時期.
整個肥妹仔.
因為每個同學都很瘦.
我就顯得很突出.
我當時真的覺得自卑.
還有開始長了很多青春痘.
一長青春痘就很想不要出門.
躲起來,不要見人.
當時的健康也不是很好.
經常會有病.
還有自己也很介懷.
我過去的成長背景等等.
我感到很自卑.
但當我成長以後.
特別是信耶穌以後.
神一步一步地改變我.
讓我對自己多了很多信心.
神告訴我們.
不論我們外在的條件如何.
祂賦予我們很寶貴的價值.
我們不需要自卑.
反而要好好地愛和珍惜自己.
也要好好地愛和珍惜別人.
今天你們弟兄姊妹都問問自己.
你真的好好地愛自己嗎?.
還有什麼令你感到很自卑的呢?.
你真的能從神的角度看自己的價值嗎?.
這一點是很重要的.
第二點.
神愛我們.
所以祂一直保護和看顧我們.
讓我們活在祂的恩典當中.
13節說:我在母腹中,你已偏激我.
原來這個偏激還有另一個解釋.

$^{241}$就是「奉庇」和「保護」.
當我們還在媽媽的肚子裡.
神已經保護我們.
這是一件很寶貴的事情.
我們在媽媽的肚腹時.
可以遭遇不同的事情.
以至我和你未必可以生存下去.
例如什麼呢?.
可能是媽媽的健康出現問題.
或者是遭遇意外的緣故.
以至媽媽發生了流產的事件.
胎兒失去了生命.
也有可能是因為人為的緣故.
例如墮胎.
以至胎兒無法生存下去.
這裡也給大家一些數據.
在2011年1月將墮胎合法化後.
直到2005年還未到今天.
美國有4500宗墮胎個案.
全球個案有多少呢?.
當時是2005年的數字.
全球墮胎的數字是13億宗.
是不是很誇張?.
現在全世界每一年.
都有好幾千萬的孩子被墮胎.
比任何的戰爭和天災死亡人數還多.
今年我又找過.
新冠疫情的死亡人數超過了188萬人.
已經高過去年全年總和.
但是其實188萬是不是超過了呢?.
我剛才說墮胎是多少呢?.
一年幾千萬.
其實這幾千萬.
是因為我們墮胎的原因.
比新冠病情的原因還要嚴重.
只不過沒有人提出.
我們夢神的保守.
可以順利來到這個世界.
是一件很值得感恩的事情.
當我們出生來到這個世界以後.

$^{281}$上帝繼續保護我們.
今天仍然有很多原因.
可以導致我們失去生命.
例如交通意外.
例如有疾病.
現在最多說的是新冠病毒.
其實也不只是這些.
每天有很多人因為疾病的緣故離世.
被人謀害,天災等等.
在不同的情況下.
都是上帝逐一保護了我們.
以致我們才可以活到今天.
這一切不是必然的.
我們有沒有為了這樣而獻上感恩呢?.
保護我.
讓我可以順利出世來到這個世界.
多年來我也很想分享一些例子.
說明上帝多次保護我.
遠離不同性質的危險.
我知道如果說到這些例子.
其實也不是最驚險,刺激的.
在在的弟兄姊妹.
你經歷的事可能比我深刻得多.
可能更加刺激,更加大的危險.
我只是拋磚引玉.
希望引起大家獻上感恩.
遠離被蛇咬的危險.
很特別的,我講這件事.
我住在一個很接近上水的村子裡.
那條村子給我最深刻的印象.
就是有很多蛇.
因為有很多田,樹,草地.
有不同種類的毒蛇在那裡.
也有沒有毒的蛇.
總之就是很多蛇.
有不止一次,我自己是很近距離的看到蛇.
有時蛇就在旁邊,其實我不知道.
因為我身體重.
我腳步聲很重的時候,就驚動了蛇.
牠一動的時候,我才發覺牠在旁邊.

$^{321}$每一次都嚇得腳都軟了.
遠離他們的傷害.
其實可能有些人說.
蛇不是每一次都有攻擊性.
那倒不是,我住在那條村子的時候.
也試過有一次有蛇攻擊我家人.
那次當然是最後把蛇打死了.
牠是有攻擊性的.
也試過有蛇攔著村子的路口.
我又進不去.
但我感恩上帝保護我.
遠離被狗追的危險.
有一次我經過工地的門口.
有很多隻狗追著我.
當時我知道被追到的話.
一定會被咬.
這個傷害就很嚴重了.
我平時跑得不快.
但過了一天,我很感恩.
上帝給我一雙摩打腳.
忽然跑得很快,那種死亡力量.
於是就跑得很快,遠離了那批狗.
牠就追咬了.
很感恩,特別多是.
上帝保護我離開交通意外的危險.
我曾經在港島說過.
我去緬甸短孫的時候.
第一天就發生很嚴重的交通意外.
當時擋風玻璃是碎了.
不停地向我們飛過來.
是上帝真的很厲害.
很快我們已經可以投入.
在那次短孫的生活裡面.
前一段時間,這反而令我更加深刻.
在一個星期日的早上.
其實不算是很久之前.
我飛的士回教會.
因為那天我想早點回來.
預備事奉上的事情.
當時我的的士正準備要乘車.

$^{361}$在不遠的地方有一個大貨車.
看到我們準備要乘車.
其實他應該要收慢制.
等我們乘車後,轉過去向前行.
他才跟著我們.
但這貨車很特別,他看到的.
但他沒有打算停下來.
而且是以很快的速度向我們衝過來.
是攔腰衝過來.
我看著的時候,那種感覺.
是害怕,我控制不了.
在的士上大叫一聲.
很感恩,是神為我預備了.
一個非常厲害的司機.
他安慰我,停下來,不用怕.
一下子他就向前飛奔.
過了對面的空地.
就一瞬間,貨車在我們後面衝過去.
我們事後討論這件事.
為何他看到我們都要衝過來.
我當時驚魂未定.
後來回到家才想.
如果那天我撞車.
應該回不了紅印堂.
我現在應該在醫院休養.
今天應該不會跟大家說.
我那天乘的士時不會想.
我會遇到這樣的事情.
神就是這樣保護了我.
在我不知道的時候成為我的幫助.
我很感恩.
還有很多晚上,我很晚回家.
我以前住的地方.
要經過一些很黑的.
危險的公園地帶.
其實晚上一個人這樣走.
是危險的.
所以我被家人罵了很多次.
一直聽我禱告,保護我.
經過黑夜.

$^{401}$否則有很多事情可以發生.
我知道恩典不是必然.
希望這一切的分享.
引發我們去想.
上帝是怎樣保護我們呢?.
有些人還在媽媽的肚腹裡.
到今天,他在哪些地方保護了我們?.
願意我們獻上感恩和讚美.
第三點.
上帝愛我們到一個什麼地步呢?.
他的愛是完全關注我們.
他為我們的人生定下美好的旨意.
一章十六節.
說到他很關注我們.
我未成形的體質,你的眼早已看見了.
未成形的體質是胚胎.
當我們還在胚胎的時候.
神已經看到我們了.
神不需要透過超聲波.
在胚胎的時候,他看到我們清清楚楚.
這個體見不單代表上帝知道我們的情況.
還代表他很關心我們,很愛我們.
好像什麼呢?.
好像孩子出生後被放在育嬰房.
父母就在育嬰房的玻璃裡.
看著自己的孩子.
很想快點看到他.
覺得他這個孩子和其他孩子不一樣.
他覺得很愛他.
漢奸就是這種感覺.
他看著肚子裡的我們.
關注我們.
深深愛著我們.
這麼久了,真的很久了.
還有,在哪裡看到他關注呢?.
十六節,你所定的日子,我尚未到一日.
都在你的冊子寫上了.
連未過的時間,未過的日子.
都寫下來了.
我尚未經歷這些事.

$^{441}$我尚未經歷這些階段.
神都寫下了.
這方面是因為.
神很有能力.
祂知道我們的未來.
今天我們不知道的事,上帝都知道了.
將來的事,我們不知道的.
但神都知道了.
更重要的是,上帝自己為我們的人生.
訂下美好的旨意.
所以十七節,詩人這樣說.
你的意念向我何等寶貴.
神的旨意不單是美好這麼簡單.
還很豐富.
所以詩人繼續這樣說.
祈禱何等眾多.
祈禱海沙這麼多.
上帝對我們的人生.
有美好而豐富的旨意.
是很多很多的.
我想起.
我在大陸.
一個城市出生.
我出生時.
從小就在那裡成長.
我開始很投入.
在那裡的生活.
我記得我的小學.
我的小學是師兄師姐們.
他們是撐磚,撐泥.
他們是一點一滴慢慢地建的.
我的小學雖然不富有.
但也有游泳池.
因為是師兄師姐們很厲害地建的.
我很喜歡那裡很可愛.
很單純的學校生活.
在那裡的我.
並不知道原來有一天.
我會出來香港.
來到香港也去過不同的中學讀書.

$^{481}$那一刻又不知道有一天.
我會信耶穌.
又不知道將來真的會考上大學.
到了大學的時候.
自己又在想,我希望我的人生.
跟其他的同學一樣.
畢業後找工作,賺錢養家.
過我自己美好的人生,很想是這樣.
又不知道原來上帝.
一直對我的人生有計劃.
祂想我去做傳道人.
很多的不知道.
每一個步,每一個階段.
上帝都讓我遇到不同的人.
我經常回想,那時候的位置.
我又沒有想過今天會遇到這些人.
會跟這些人做朋友或同工.
但這一切都已經寫在上帝的冊子上.
你去想想你的人生.
我們尚未到一日.
上帝都寫上.
祂有美好的旨意.
祂很想賜福氣給我們.
再回想剛才我們說的那幾件事.
上帝的愛為何無可取替呢?.
有很多事是地面上的爸爸做不到的.
上帝做得更加盡,更加徹底.
祂的能力很超越.
第一是祂做我們給我們生命.
賦予我們價值.
祂一直保護我們.
沒有人可以24小時保護我們.
只有天父可以.
天父對我們人生有奇妙的旨意.
是直到永遠.
就這幾件事讓我們看到.
上帝的愛無可取替.
既然上帝的愛這麼寶貴.
我們有沒有好好地愛祂呢?.
我們聽到父親節.

$^{521}$我們如何對父親孝順等等.
其實天上的爸爸.
我們又如何待他呢?.
我們如何才愛到他呢?.
我深信就是要聽他的話.
遵行他的旨意.
那就是愛他了.
你問我們在現在的生命裡.
我們是否真的聽他的話.
遵行他的旨意.
我們一起來禱告.
讓我們看到你的作為是何等奇妙.
可能當我們還不太懂得愛自己的時候.
你已經很深深地愛著我們.
願意讓我們明白你的愛.
能夠從你的角度去愛自己.
去愛別人.
天父同樣讓我們看到你的保護是何等珍貴.
我們就上上去獻上感恩.
說說你的恩典.
還有天父求你幫助我們去倚靠你更多.
最後的是.
天父你對我們一生有美好的旨意.
願你讓我們有開放的心去聆聽你的心聲.
明白你的旨意.
我們有恩賢的踏上.
做你所喜悅的人.
願你賜福給洪恩家.
讓洪恩家的弟兄姊妹更深地明白你的愛.
也因為你的愛的激勵.
而活出一個精彩的人生.
我們這樣的禱告是奉耶穌基督的名字.
阿們.
\newpage



\section{歌羅西書 1:16-32}
\label{sec:M84h3HizR5w}
\textbf{2021.06.27《不要自欺欺人》 歌羅西書 1\_16-32 ( 和修版)  鄧泰康傳道}
\newline
\newline
連結: \href{https://youtube.com/watch?v=M84h3HizR5w}{\texttt{https://youtube.com/watch?v=M84h3HizR5w}} ~~~~ 語音日期: 2021-06-29
\newline
\newline
\hyperref[sec:T6XL5VlhAWY]{\small{< < < PREV SERMON < < <}}
~
\hyperref[sec:index_chronic]{\small{[返順時目]}}
~
\hyperref[sec:index_scriptual]{\small{[返順卷目]}}
~
\hyperref[sec:F4zrvJPl9fU]{\small{> > > NEXT SERMON > > >}}
\newline
\newline
歌羅西書 1:16-32
\newline
\begin{longtable}{cl}
\hline
\hline
章節 & 經文 (和合本修訂版)\\
\hline
1:16 & \begin{tabularx}{0.7\textwidth}{X} 因為萬有都是在他裡面造的，無論是天上的、地上的，能看見的、不能看見的，或是有權位的、統治的，或是執政的、掌權的，一概都是藉著他為著他造的。 \end{tabularx} \\ \\
[1ex]
\hline
\hline
\end{longtable}
$^{1}$(主持人: 各位觀眾早安).
(主持人: 大家很精神).
我最感恩的是每一次來到大家當中的時候.
大家都是帶著很有精神的一起參與在神的面前.
為什麼呢?我經常說每天回來崇拜的時間.
最重要的是要預備好.
我經常覺得如果每一次崇拜之前都能夠早一點回來的時候.
你就能夠有更好的精神去聽上帝的話.
跟著跟聖靈一起合作.
讓聖靈在你心裡面去感動你.
然後行出神的生命.
這才是真正崇拜的裡面我們能夠達至.
所以我很開心的.
特別是在大堂的時候.
大堂講完之後我都不太聽得到.
特別是去到清善崇拜的時候.
我是很要求的.
我講第一次的時候通常沒有人回答我.
然後就不斷講早晨的時候.
讓他們回答我為止.
我之後才開始.
上次跟大家已經是兩三個月前見面了.
今天我再跟大家分享的時候.
其實我心裡面一直感覺到.
有一件很重要的事情.
就是什麼呢?.
我們經常說今天我們教會在地上的時候.
為上帝去作見證,去傳福音.
但其實我想告訴大家.
我聽的時間為上帝去傳福音.
最重要的不是技巧.
我想再強調不是技巧.
因為香港是一個太知識型的地方.
什麼都是.
我不懂就不做.
每一樣東西我不懂的時候.
我就要學.
學到夠了我才做.
同樣的就是.
你有沒有聽過很多人說.

$^{41}$你跟他傳福音.
之後你很想他相信.
之後很多人就跟你說.
我不可以回教會住.
我今天早上聽郭文前牧師講道的時候.
他講的一個例子很有趣.
他說有一個伯伯.
他跟他去傳福音.
之後他很想他相信.
接著伯伯就說.
遲一點吧.
為什麼呢?.
因為伯伯說我不想早上回教會.
下午就去馬會.
所以在這樣的情況下.
有很多人就覺得.
不是的,我要OK的.
我OK而已,我才可以回教會.
但那位牧師說了一件事很有意思的.
你搞錯了.
你完全不了解我們的信仰.
我們的信仰就是正因為你不好.
所以你相信了耶穌之後.
你靠著神的能力去變好.
是不是?.
所以同樣的就是.
今天我們在神面前.
我們為上帝去傳福音.
最重要的不是你有什麼技能.
哪怕你能夠很暢順地將三福.
將所有傳福音的工具.
福音橋.
你跟別人很流利地說完.
但是如果你沒有一個屬於上帝的生命.
我可以告訴你.
你是一點作用都沒有的.
為什麼呢?.
因為我當你是很厲害.
你就將那些人帶回教會裡面.
但我經常有句話.

$^{81}$小心,教會前面的門開多大.
當你不小心,後門就開多大.
因為那些人一回來之後.
我經常說人與人的關係很實在.
一接觸的時候就見真章.
我和你不熟的時候.
坦白說.
你每次見到鄧全道的時候.
鄧全道的事多好.
當然好,我每次都在港版國安法.
到你和我接觸多的時候.
就真的挑戰了.
鄧全道平時為人怎樣呢?.
我和你說要靈修.
鄧全道都不靈修的.
我和你說愛的時候.
鄧全道凶神惡煞的.
我和你說愛的時候.
我都沒有耐性.
就見真章.
所以最重要的就是.
今天我們的生命.
所以我今天很想回到聖經裡面.
和大家回到一個很基本的問題.
是今天很多時候.
當教會慢慢成長.
我們相信住了一段日子之後.
我們慢慢會忽略了.
就是關於罪的問題.
羅馬書是一卷非常重要的書卷.
我很鼓勵弟兄姊妹.
你無論如何都要讀懂它.
為什麼呢?.
因為羅馬書是整個聖經裡面.
講罪講得最清楚.
和講得最徹底的一卷書.
第一至十一章的時候.
其實它分享到.
罪對人生命裡面的影響.
律法和清義之間的關係.

$^{121}$和人依靠自己在罪裡面的時候.
根本是沒有出路.
其實保羅在裡面的時候.
我想告訴大家.
我們一直都覺得.
保羅是一個很厲害的信徒.
很厲害.
但是我想告訴大家.
他在第七章的時候.
他很實實在在地告訴大家.
他面對罪的時候.
他的掙扎和你的掙扎是一模一樣的.
大家記不記得那段經文?.
立志為善,由得我.
只是行不出來,由不得我.
接著是什麼?.
我真的很想做好.
但我裡面有一個律.
令到我沒辦法可以做好.
接著最後有一句歎息的說話.
就是什麼?.
神啊!.
我這個這麼取悉的身體.
誰能夠去救我?.
我想告訴大家.
保羅在這裡做了一個.
很重要很重要的示範.
今天你千萬不要在別人面前的時候.
你記住.
不是假裝.
假裝我是一個.
頭上有一個光圈.
我是一個怎樣的聖人.
而是我們今天面對罪.
我們都在掙扎.
我們同樣在掙扎.
但是保羅有一個說話很重要.
他不是跟你說了他的困難.
接著他說了一個很重要的結尾.
是什麼?.

$^{161}$靠著我們的主耶穌基督.
我就能夠得勝.
記住這句話.
保羅今天這麼厲害的緣故.
不是他變成一個聖人沒問題.
而是他穩住了一個很重要的方法.
原來我們靠著主耶穌基督的話.
我們就能夠得勝.
弟兄姊妹.
今天我們每一個信耶穌的人.
都一樣是這樣.
今天你能夠帶別人信耶穌.
你不是說在別人面前的時候.
多好多好.
我挺肯定的.
因為教會要成長.
坦白說有一個一定走不出的程序.
就是由你的生命.
去吸引其他人來到神的面前.
接著帶他們去認識上帝.
是不是?.
所以在教會裡面.
如果大家沒有關係的話.
根本不會傳到福音.
對嗎?.
但最重要的就是.
不是你有多好.
你要搞清楚.
所以有些人說你真好.
鄧傳道你真好.
我經常提醒自己.
我就會告訴他.
不是我好,是誰好?.
上帝好.
是耶穌好.
是他怎樣救我.
是他怎樣令我這個這麼骯髒的罪人.
我能夠得勝.
對嗎?.
所以弟兄姊妹.

$^{201}$今天我的主題叫做.
不要自欺欺人.
因為罪永遠都在影響著我們.
不要掉以輕心.
不要以為我信了耶穌之後.
接著我就有了保護罩.
我不會犯罪.
我相信是騙人的.
罪對人的影響.
永遠是超過了我們的想像.
就舉個例子.
我以前去事奉的教會.
是一間教會和戶外中心有合作.
戶外中心裡面的時候.
是一些過來人去吸毒.
對嗎?.
當你和他們聊天的時候.
他就會告訴你.
其實大部分吸毒的人.
都有一個同樣的特徵.
就是他認為這粒小小的毒品.
是不可以控制我的.
因為每一個人的時候.
通常的渠道都是這樣的.
有一個朋友去吸毒.
接著時間.
說實話,我問大家.
你吸過毒嗎?.
當然不吸了.
接著那些朋友又怎樣?.
你沒有膽量?.
你這麼差勁?.
接著又怎樣?.
我要證明我有膽量.
我要證明我不差勁.
接著我要吸.
吸完之後.
這粒沒事的.
你最少吸,沒事的.
接著時間,第二次沒事.

$^{241}$之後一直下去.
行不行?.
控制不了.
由他去控制回來.
每一個染上毒癮.
倒癮的人.
都認為自己可以控制不了.
大家有沒有.
以前還沒回教會.
有沒有去過澳門島的經驗?.
應該有吧?.
你去澳門不玩兩手?.
你只是吃牛肉乾.
我不相信.
OK?.
但我想告訴大家.
我先戴上頭盔.
如果有一天的時候.
當大家在澳門.
你看見鄧全導在島上出現.
你不用覺得奇怪.
因為我和我太太.
以前去澳門.
澳門除了吃玩樂.
沒什麼好做的.
我和我太太.
以前有一個習慣.
很喜歡去賭場.
看人們怎樣賭.
我們很喜歡站在後面.
翹起雙手.
看人們怎樣賭.
之後我們發覺.
我記得有一次.
有一個經驗.
有一對年輕拍拖的人.
去澳門玩玩.
他們在我們旁邊.
聽到他們的對話.
我們不是偷聽.

$^{281}$接著那個男生和女生說.
我們玩玩.
他們說好啊.
五百塊.
贏也好.
輸也好.
我們都走了.
我想問.
是不是大部分都是這樣的.
是吧.
玩玩嘛.
那又不知道為什麼.
他們去賭的時候.
那個男生很運氣.
給他很短時間.
半個小時之內.
贏了幾千塊回來.
走不走啊?.
大家很肯面對自己.
當然不走.
如果你告訴我.
走的那些很厲害.
真的很厲害.
大部分都不走.
這麼好運氣都走.
接著怎樣.
繼續賭.
大家知道結果了.
我記得以前何先生有句話.
我不怕你贏.
我最怕你不賭.
接著他賭很久了.
慢慢慢慢輸下去了.
輸著輸著的時間.
那五百塊都輸了.
走不走啊?.
為什麼不走?.
剛才我說了.
我們輸五百塊就走.
為什麼?.

$^{321}$剛才我這麼好運氣.
這麼容易就贏幾千塊回來.
為什麼走啊?.
我一定贏的.
那怎樣?.
繼續賭.
接著我和我太太後面.
真的去吃花生看電影.
看著他們在玩.
看著他們在起起跌跌.
由很好關係開始吵架.
我們看到那個女生開始扯著男生.
走啦走啦.
不走不走.
到最後怎麼樣?.
接著之後.
全部去澳門的旅費都輸了.
接著之後.
那個女生跟他說.
我們酒店還沒訂.
回香港的船票還沒買.
怎麼辦?.
那時候不是說了嗎?.
五百塊走啊.
我誰知道啊.
結果大家很開心去旅行.
變成什麼?吵架.
每一個人都會認為自己可以控制自己.
第三方面.
有分愛情的人.
都認為自己的聰明智慧.
是可以處理兩段感情.
白面玲瓏.
如果有自知之明的人.
是不會有分愛情的.
就是那些人認為自己很聰明.
不會的.
我和這個的時候.
我不會讓我太太知道.
我做什麼的時候.

$^{361}$我太太不會清楚.
認為自己可以去控制.
到最後通常怎麼樣?.
很奇怪的.
天網恢恢.
疏而不漏.
又不知道什麼時候.
接著你那張坐過的車的車票被人看到.
接著之後.
你曾經怎麼匯錢的那張單.
又被人不知道為什麼會看到.
你會發覺一件事.
原來你發覺裡面都有一個特性.
每一個人都自以為自己聰明.
但弟兄姊妹.
我想告訴你.
你不要以為上帝不知道.
今天最奇怪的地方就是.
我們犯罪的時候.
永遠都在一個黑暗的環境當中.
我們以為沒有人知道.
弟兄姊妹.
上帝知道.
上帝清清楚楚.
知道你在做什麼.
所以弟兄姊妹要小心.
其實剛才那段經文裡面.
一開始就告訴你.
上帝透過萬事萬物告訴你.
祂存在.
祂主動向你顯示祂的存在.
而聖經裡面說.
你是心裡就知道.
其實人有一種天性.
我不知道大家知不知道.
大家有沒有去過北京天壇.
認不認識天壇.
天壇用來做什麼?.
BBQ?.
不是吧.

$^{401}$天壇是皇帝祭天的地方.
你弄清楚.
不是拜神.
是祭天.
因為連皇帝都認為.
有一個高高在上的力量.
或者神明是比祂高.
祂知道的.
所以祂自稱是天子.
我是以神明的力量.
賦予我這個身份.
所以你們全部的人.
都要聽我的話.
所以其實人的心裡.
本來就知道.
所以當你了解的時候.
你會發覺全部很原始的信仰.
都不是有偶像的.
他們都是拜石.
拜山.
拜樹.
看到嗎?.
香港也有.
我想問問大家.
如果你以前.
你未信耶穌之前.
如果你是拜偶像的話.
你有沒有看到在街上.
通常有些被人標記的偶像.
是放在哪裡的?.
樹底.
還有呢?.
石頭邊的.
有沒有留意?.
其實你會發覺.
他們心裡很奇怪.
他們認為這些東西.
是有一種神靈的力量.
大家知不知道土地公.
原來就是石頭?.

$^{441}$原來他們覺得這個世界.
是有一種力量.
我想大家留意.
這是聖經告訴你的.
是上帝創造你的時候.
放在這裡的.
讓你去追求.
原來這個世界是有一個神的.
你不知道為什麼.
你出生就已經有這種感覺.
所以當你去了解地球的時候.
你就會發覺.
你的眼見.
首先是心裡知道.
接著透過你的眼見.
你會看到地球是有設計的.
地球如果大家讀一點點科學.
你都知道.
地球是在轉的.
為什麼你會日出日落?.
大家記住.
不要像以前舊的信仰.
我們不動的.
是太陽這樣走.
其實是什麼?.
不是.
是我們在轉.
是上帝創造的.
地球原來是斜斜的.
為什麼要斜?.
大家知不知道?.
原來和海的潮汐漲退有關.
所以為什麼有時候會亂七八糟.
就是差一點點.
所以很多時候.
我們傳福音的時候.
我會和別人說一個.
很常用的例子.
有一位物理學家.
霍先生.

$^{481}$他說這個世界是爆出來的.
這個地球是爆出來的.
很多理論解釋這個世界是爆出來的.
但我想問大家.
我很喜歡用的例子就是.
我拿著一隻錶.
我想問這個錶有沒有人設計的?.
肯定?.
為什麼你這麼肯定?.
沒理由的.
又有弧面.
裡面又有電子.
又會跳.
但是.
接著我說不是.
沒有人設計的.
爆出來的.
在世界裡面一直爆出來.
最初的時候.
世界創造.
爆了爆了.
爆了石頭.
再震了震了震了.
震了碎沙.
再一直在世界轉變的時候.
融了融了.
不知道為什麼加了氣溫.
就變成零件了.
零件震了震了.
又不知道為什麼震了一起.
剛好就變成一隻錶.
有沒有可能?.
如果連這隻錶你都覺得沒有可能.
這個地球.
你說沒有一個創造主?.
但很奇怪.
就是一個這麼具名的科學家.
告訴你.
沒有的.
這個世界是爆出來的.

$^{521}$人就是這麼奇怪.
自以為聰明.
聖經說人不可以否定.
但是人就喜歡挑戰.
人就是喜歡挑戰權威.
就好像最初的亞當和夏威夷一樣.
上帝明明跟你說了.
你不要吃那棵分辨善惡樹.
到後來.
就算找個藉口我也不聽.
我也要吃.
我也要挑戰.
上帝你說了.
吃了的話一定會死.
然後之後.
透過蛇的引誘說什麼?.
不一定死.
人總是喜歡跟上帝違反.
跟他頂撞.
不聽他的話.
人視而不見.
然後就將自己心裡面.
所想的化成偶像.
剛才我說的很具體.
人喜歡將自己變成上帝.
但又不喜歡告訴別人我是上帝.
為什麼我會這樣說呢?.
聖經裡面說了一樣東西.
人對造物主應該有一個追求和敬畏.
是以身俱來的.
就是說你不對這個造物主敬畏.
偏要自己做主.
我舉個例子.
剛才說的那些山和天.
變成什麼呢?.
天變成了什麼?.
天庭.
上面有個玉皇大帝.
天君天將.
我想問大家是什麼?.

$^{561}$你想清楚.
那些是地上的體制.
為什麼會這麼實在呢?.
我心裡想的.
地上是這樣.
上面是這樣.
山變山神.
石就變土地公.
所有萬事萬物的時候.
都有一個神明去掌管著.
我想告訴你.
不是東方.
就算連西方都是一樣.
所以現在就算看電影.
雷神.
宙斯.
全部都是這些東西.
大家還沒相信耶穌之前.
有沒有去拜一下偶像?.
起碼拜一下祖先吧.
我想問問大家.
有沒有見過有人燒紙紮的那些?.
我小時候很頑皮.
我小時候經常去紙紮.
大家有沒有見過紙紮的那些很大間的屋子?.
我很喜歡裝進去裡面看看.
看看放什麼呢?.
有沒有東西呢?.
我很喜歡看一看.
但如果大家有機會.
你會燒一個包括什麼呢?.
紙紮的屋子.
電視.
還有什麼?.
車.
燒兩個工人進去跟他們打麻將.
是不是?.
現在再新潮一點.
要什麼?.
要蘋果電話.

$^{601}$是不是?.
還有什麼?.
你的名牌.
我在地上買不起的那些.
你燒給他吧.
OK?.
我想問問大家一個問題.
你放一部電視機進去.
為什麼不燒發電廠給他呢?.
你明白我的意思嗎?.
你想把車給他.
你又不給他油.
他怎麼走路呢?.
我想從這裡告訴你.
其實人就是那種自我滿足.
人就是那種.
我認為是這樣的.
所以有時候你聽到一些.
那些真的是靈異事件.
你跟我說很冷.
然後怎麼辦呢?.
燒衣服給他.
燒被子給他.
你真的相信燒下去.
明明燒下去之後會發一件衣服.
你真的覺得燒下去之後會變回一張被子?.
為什麼呢?.
因為這是你的真實感受.
弟兄姊妹.
我想告訴你.
人很奇怪.
自己自以為神.
但他被放在一些東西裡面.
然後.
去掌管這些.
但反過來又被這些東西.
核制自己.
舉個例子.
在街上很多時候.
不要偶像劑的石頭邊.

$^{641}$我會很頑皮.
我以前頑皮到.
我打那些不是斗槍.
是氣槍.
我用什麼來練靶.
當然是那些.
整個山頭都是.
用來練靶.
然後我媽媽就跟我說.
你將來會死的.
人家那些你都敢動.
以前有沒有見過.
不要踢到香爐灰.
肚子痛.
鄧先生今天在這裡.
跟大家有一個好見證.
我正正齊齊.
沒有任何災病.
因為這個緣故.
但我想說.
上帝製造了一些東西.
但又被這些東西控制了自己.
人是不是很矛盾.
但你會發覺.
上帝很奇妙地告訴你.
我們經常誤會了.
我們經常說.
舊約的上帝比較兇.
大家有沒有發覺.
又要裂開地.
又要降火雷燒整個城.
這個上帝很兇.
但我想告訴你.
如果上帝真的用舊約的標準.
對待你的話.
我想告訴大家.
今天你和我都不在了.
因為本身.
我們的罪惡一直都在.
只不過.

$^{681}$上帝以更大的恩慈.
去代理.
然後在聖經.
告訴你.
神真正對人的懲罰.
是什麼呢.
是任憑.
任憑.
有時候你都誤會了.
為什麼那些殺人的.
金腰帶.
修橋補路.
沒有屍鞋.
為什麼那些作惡多端的.
他們很地道.
弟兄姊妹你放心.
原來在神的眼中的時候.
不是這樣的.
在神的眼中的時候.
其實祂說既然你們是這樣的時候.
我就任憑你們去做.
而聖經在裡面的時候.
因為當時的世界是針對希臘文化.
希臘文化裡面有一個很特別的地方.
就是什麼都和性.
連在一起.
如果大家有機會.
你熟悉電腦.
你寫在電腦上打希臘文化.
看完之後你會毛骨悚然.
因為希臘文化裡面.
它告訴你.
在男人的角度.
怎樣看女人的呢.
每一個娶回來.
每一個屬於我的女人.
都有她的功能和作用.
娶老婆回來做什麼.
去管理家庭.
妓女用來做什麼.

$^{721}$用來給我享樂.
我得到的奴隸.
女奴.
又幫我做事.
和給我享樂.
他們是這樣的觀念.
他們會覺得.
任何的事情.
我們今天只需要快樂.
就可以了.
我們追求今生的快樂.
所以你會發覺.
他們的神明.
都是淫亂的.
為什麼會這樣呢.
同一個道理.
他們心裡好色.
令他們的神明偶像.
都好色.
所以聖經裡面說.
所以神就任憑他們心裡的情慾.
行污穢的事.
以致彼此沾污自己的身體.
他們將神的真實變為虛方.
去敬拜事奉受造之物.
不敬奉那造物之主.
主乃是可稱頌的.
直到永遠.
神是任由他們去做的.
然後聖經裡面說.
又任憑.
我想大家留意.
下面這麼短的經文.
說了很多次任憑.
因此神就任憑他們放縱可恥的情慾.
去將他們女人的順性用處.
變成逆性用處.
我不詳細再讀.
在這裡的時候.
今天特別是現世的裡面.

$^{761}$我不知道為什麼.
人很喜歡高舉自己的智慧.
包括基督教裡面.
他用自己的方法.
去解釋.
同性戀沒有問題.
甚至他用很多.
的色經告訴你同性戀沒有問題.
但弟兄姊妹在這裡的時候.
這段經文是清清楚楚.
告訴你.
女人有順性.
逆性.
男人有順性逆性.
是上帝創造下來的時候.
清清楚楚的.
是不是啊?弟兄姊妹.
是沒得爭辯的.
雖然今天有很多人.
我想強調.
我們不是去推開.
或者拒絕這些朋友.
不是,你搞清楚.
在我的角度裡面.
我覺得每一個人都有難處.
我以前的時候.
都會和很多朋友聊天.
我不會一和他聊天的時候.
就在說你對我錯,我對你錯.
我不會說這些的.
反過來,我會聽他講他的故事.
最初一定是.
我可以告訴你,十個十個都告訴你.
我從很小的時候就開始覺得.
自己是想做女生的.
或者很小就想做男生.
十個十個都這樣說.
但你常常弟兄姊妹.
用你的耐性.
聽她說話.

$^{801}$陪她,不要先害怕.
了解她.
慢慢當你的愛.
你的關心融化之後.
我可以告訴你.
百分之差不多一百.
都會將她的故事告訴你.
我可以告訴你.
我所總結的.
基本上有幾個原因.
第一個,我都勸勉.
這裡的弟兄姊妹.
如果你做爸爸媽媽.
覺得他們很可愛.
做什麼呢?.
有沒有見過其他人.
他們的男生.
很小的時候,男生會給他們裙子.
然後紮一條辮子.
你很可愛.
很可愛.
然後拍照.
從小到大,當他們像女生.
去玩.
其實你知不知道.
影響很大.
因為有很多人跟我說.
小時候.
每次媽媽.
給我穿裙子.
她就會稱讚我.
她說我漂亮,我可愛.
但每次我穿回男裝.
她都會罵我.
你為什麼這麼頑皮?.
為什麼搞這個搞那個?.
慢慢成長之後.
她接受了女性的身份.
是好的.
她說她從小開始.

$^{841}$認為自己是女生.
第二個.
很真實.
特別是現在的社會.
原來.
不少人,包括男人和女人.
在他們的成長過程中.
他們都被人.
否定了他們的性別.
例如.
一個男生.
比較溫柔.
說話音聲細氣.
大家怎樣?.
那型.
是吧?.
男生不是這樣的.
我想告訴大家.
我有一個很好的例子.
我很感激我沒有變成女人.
我小時候.
其實我是一個很容易哭的男生.
六歲之前,我什麼都會哭.
見到螞蟻死了.
又哭,見到蟑螂死了又哭.
我很容易哭.
所以我小時候有個綽號叫哭包.
但我每次被人說哭包.
我一直長大後.
哭是不好的.
之後我就不再哭了.
我去到十幾歲.
我就不再哭了.
我去到變態度,那是什麼情況呢?.
我見到朋友去世了,我都不哭.
我就不斷告訴自己.
我是一個很理性的人.
我是一個不會哭的人.
在我讀完神學的時候.
我被醫治了.

$^{881}$我發現原來上帝創造了我.
原來我是有理性和感性兩邊的.
為什麼男生不能哭呢?.
為什麼男生一定是流血不流淚呢?.
我想問問大家.
這是什麼邏輯呢?.
男生就沒有眼淚?.
女生哭就很可憐?.
為什麼要小心這些事情呢?.
有很多人在成長的過程中.
女生被否定,為什麼呢?.
因為豬扒.
是不是?不漂亮.
然後你就去傾向一個.
什麼叫漂亮的呢?.
慢慢的.
那些不能進入流淚的就變成什麼呢?.
變成男生.
你有沒有發覺?很極端的.
就把全部頭髮剪短.
像男生一樣,出來之後反轉.
你很帥.
做男生.
頂尖妹原來全部.
都有他們的轉化過程.
但我想告訴頂尖妹.
在當日的社會更加恐怖.
我根本不用跟你說這個轉化.
我想告訴你,在裡面的世界很簡單.
我喜歡就行了.
我不喜歡跟女人,我喜歡跟男人.
我認為怎樣就怎樣.
最慘的就是.
他們用他們自己信仰的.
當日他們自己的.
去解釋他們的行為.
頂尖妹.
但上帝說,任憑你們.
按你們的慾火攻心.
最後的時候.

$^{921}$去得著這個妄為當得的報應.
我形容的時候就是.
上帝讓人在我們人與人之間.
任憑你們去做的時候.
大家在你們的人際裡面.
受夠了自己生命裡面的傷害和破損.
我形容的是什麼呢?.
了倦之還.
如果大家剛才有留意那段經文的時候.
聖經從來沒有試過.
這麼詳細.
這麼細緻.
將裡面的所有的東西.
裡面的邪惡.
不義,惡毒,貪婪.
恨父母,什麼的,全部都出來了.
他很想告訴你.
原來一個在罪裡面的人.
理心裡面充滿著的.
就是這些東西.
頂尖妹,請你記得.
原來這是一個對比.
我想問問大家.
今天我們是一個在上帝裡面的基督徒嗎?.
你怎樣證明自己是一個在上帝裡面的基督徒?.
我什麼時候舉手決了戰?.
我是基督徒.
是,沒錯.
但問題是.
客觀地我怎麼知道?.
我和上帝知道就行了.
這些是藉口.
聖經有標準的,大家知道嗎?.
記不記得.
聖靈的果子.
人愛喜樂和平.
恩慈善良節制.
聖經裡面是有標準的.
我所說的不是說.
你有這些標準之後你才行.

$^{961}$我不是在說這個問題.
這是告訴你.
你究竟生命裡面有沒有從上帝裡面長出來的果子.
有沒有.
如果當你有的時候.
你在聖靈當中.
但如果你心裡面自己問.
不是,我心裡面是這一堆.
邪惡多一點.
請你小心.
證明你在罪當中.
為什麼不要自欺欺人?.
今天在教會裡面.
最糟糕的就是.
我們回到整個氣氛環境.
令我們要假裝.
我們是基督徒.
我舉個例子.
你明明還沒戒煙.
你也不敢在教會裡吃.
是吧.
但接著怎樣?.
第一時間衝出門口.
先說兩句.
是吧.
在這裡裡面的時候.
你變得的時候.
明明裡面是這樣,你也假裝了.
但弟兄姊妹,我想告訴你.
其實不是要假裝.
是要面對.
是要處理.
如果你今天.
仍然在這個狀態裡面的時候.
我想告訴你.
你自己也覺得很累.
你自己也覺得很辛苦.
上帝在裡面的時候.
不斷逼使你.
你要了卷之還.

$^{1001}$第十八節說.
原來神的憤怒從天上顯明在一切不乾不易的人身上.
就是那些行不義.
阻擋真理的人.
上帝要我們.
在裡面的時候.
人際關係當中受損.
破損,然後才能回到神的裡面.
我舉個例子.
我是一個信了耶穌.
從十幾歲到現在.
已經這麼多年了,我也不想再說.
我信了耶穌多少年了.
很少去教會.
你告訴我,在教會裡沒有人際關係的問題.
有沒有可能?.
沒有可能.
但很多時候.
我們的處理就是.
避開.
還有一件事.
自己騙自己.
上帝說要憐憫.
我不跟祂吵.
你的心裡面.
是不是真的不跟祂吵?.
你想多一句,不如少一句.
我不是叫大家衝突.
我不是說這個意思.
我是想說.
如果當你有這些問題的時候.
請你好好回到神的面前.
禱告.
求上帝憐憫我們兩個.
我們在神的面前.
其實大家都是罪人.
聖經從來都說.
整個教會裡面.
我們都是蒙恩的罪人.
我們一直都是蒙上帝的恩典.

$^{1041}$而能夠走到今天.
我們有什麼這麼厲害?.
我們有什麼這麼了不起?.
我們有什麼比別人高?.
回到上帝的面前.
回到夫妻的面前.
夫妻裡面的艾煞.
我經常說.
一隻手是拍不響的.
永遠都是.
大家有事才能夠拍響.
雙方裡面.
都有一些事情做得不好.
做得不OK.
所以求上帝憐憫我們兩個.
然後求上帝.
給我們一個好的時機.
可以去處理.
這個這麼不容易的問題.
當然裡面需要很多的勇氣.
需要很多裡面上帝的憐憫.
需要很多我們的謙卑.
需要很多我們裡面.
去在神面前求幫助.
但我想告訴你.
這個才是真正上帝要我們教會做的事.
在神裡面的時候.
我們能夠回到上帝的面前.
因為沒有一個人.
不可以在神面前得到醫治.
在以前的時候.
我曾經認識一位弟兄.
這位弟兄很小的時候就開始回教會.
中學時代就開始回教會.
但他是一個信得很差的人.
為什麼我會這樣說呢.
他在教會是一個人.
出道是另外一個人.
他就是那種早上回教會.
下午就去打麻將的人.

$^{1081}$但是我們小時候.
我們那些年代.
很小的時候就受裝備.
很快就叫你去做事奉.
大家有沒有聽過.
如果你聽過事奉的職位.
證明你跟我一樣大.
有沒有聽過一個事奉職位叫靈修.
我們長大了都有.
靈修做什麼的.
就是你預備好一些靈修的資料.
給弟兄姊妹靈修.
但你可不可以想像一個從來都不靈修的人.
去預備資料給別人靈修.
當時就是這樣.
這個弟兄從來都不靈修.
交功課.
每次的時候都會去找些書.
扔給別人拿.
完成了.
這樣的信仰生活.
其實說真的.
到最後一定會死的.
所以到了中午的時候.
這個弟兄就說教會都沒有愛的.
你知道年輕人很喜歡小圈子.
又不成熟.
自己玩又不叫你.
他跟上帝祈禱.
他說上帝啊.
我知道有你在.
但我不需要你.
因為我信你和不信你是沒有分別的.
然後他就離開教會.
讀書又不成熟.
出去之後工作.
很奇怪.
離開教會那五年是他人生裡面.
最豐盛的那五年.
為什麼.

$^{1121}$他想怎樣都可以.
然後呢.
在八十年代.
大家知不知道.
我們都要體諒這班年輕人.
八十年代.
在我成長的年代.
是一個上班.
四點半就可以洗手.
預備下班的時候.
我們的家生.
當時是怎樣呢.
不用努力.
隨隨便便.
做完之後每人都有百分之十.
每年百分之十的家生.
好一點的百分之二十三十.
所以當時很多人跳槽.
跳的時候你又再跳.
所以在那個年代.
由做文員升做高級文員.
高級文員之後行政助理.
行政助理之後升做副主任.
副主任升做主任.
順順利利.
一片光明.
上帝看你沒有你更好.
但大家知不知道為什麼會好.
我想問一下.
在那個年代那麼多人要升職.
為什麼要升你.
只有兩個可能性.
在我們成長的年代.
只有兩個可能性.
一是你真的很厲害.
另外一個.
就是很會擦鞋.
經常陪老闆吃飯.
經常陪老闆這樣那樣.
這是一群後者.

$^{1161}$所謂升職升得很快.
在那個時候亦都結婚.
是他結婚的時候.
結婚後滿不滿.
買車買樓.
人生裡面那麼多兒子.
還未生孩子.
齊齊了.
很好.
但我想告訴你.
他才發現一件事.
原來人心是不滿足的.
當他一路升上去之後.
他認為自己很厲害.
然後娶了一個老婆在旁邊的時候.
然後有人看著他.
黃臉婆.
為什麼我公司裡面.
每一個都比你漂亮.
為什麼.
上班的時候大家會不會很窩心.
不會 起碼你都梳頭.
但我想告訴你.
如果家裡的老婆每天都化妝.
對著你 你會不會很害怕.
其實是很正常的.
但又怎樣.
他又覺得不入流.
男人有句金句.
不要為了一棵樹.
就放棄一個樹林.
然後怎樣.
想離婚.
但又不想做陳世美.
那方法是怎樣.
就是不斷推動老婆.
煮東西煮得不好.
家裡那麼混亂.
你都不知道怎樣.
很多的埋怨.

$^{1201}$每天都在轟炸老婆.
然後你要感恩.
幸好最後老婆並沒有離婚.
過了日子之後.
這邊搞不定了 離不了婚.
在公司裡面又沒有朋友.
你都踩著人家的照片.
誰當你朋友.
然後就想是不是做的壞事多.
就出去做義工.
但做義工又被人陷害.
我真的想做好事.
但又被人抄襲.
這樣的情況下.
我想告訴你.
人很脆弱.
讀書又不擅長.
又沒有得到好成績.
家裡又不開心.
那最後怎樣.
就想死.
不想生存.
但我想告訴大家.
很奇怪 弟兄在他想死的時候.
突然間很奇怪.
五年裡面上帝沒有再出現.
突然那一刻心裡面.
自己問自己.
你的神在哪裡.
然後鼻涕淚流.
到老婆下班回來後.
已經哭了.
跟著老婆說我想回教會.
老婆說.
如果能幫到你.
我陪你回去.
因為當時兩夫妻已經沒什麼感情.
如果對你好就陪你回去.
一回去之後.
不得了.

$^{1241}$弟兄回到教會後.
認識了傳道人.
之後180度轉變.
他真正看到自己那種邪惡的罪.
離開了神那麼多年.
他很可怕.
到最後他完全改變.
半年之後他很投入教會.
大家知不知道有三福.
讀過吧.
有三福最難的地方就是找報道對象.
那弟兄.
第一個找報道對象是誰.
當然是自己老婆.
但他老婆知道他不相信.
因為他老婆後來已經沒回教會.
回到教會後說我要做實習.
你聽我說.
然後就在那裡說了一輪三福.
最後順理成章問他.
你信不信.
他老婆回應.
我信.
然後他弟兄說你不要玩了.
你都不回教會.
你不要敷衍我.
他說我認真信.
他問他為什麼.
他說這半年雖然我沒回教會.
但我看著你變.
他說我以前用了那麼多方法.
哭又哭過.
求又求過.
但你都沒有轉變.
但這半年裡你的上帝真的很厲害.
你的上帝真的很厲害.
我從你身上看到這個上帝是真的.
然後他信耶穌.
之後這個姐妹信得非常好.
一家人裡.

$^{1281}$就一直去事奉上帝.
很開心.
為什麼我那麼清楚這個弟兄.
因為那個就是我.
有很多人問我.
為什麼你會做傳道人.
我很簡單回答你.
是上帝幫我拯救了我的家庭.
拯救了我的生命.
拯救了我整個人生.
我還有什麼不可以給他.
弟兄姊妹.
究竟你在神的裡面.
你有沒有經歷過上帝拯救的大能.
你知夠壞了沒有.
我經常問我自己.
夠壞了沒有.
在神面前的時候.
我們去受這些苦.
夠了沒有.
我想告訴大家.
跟大家讀第十六節.
我不以福音為始.
這福音本是神的大能.
要救一切相信的先是猶太人.
後是希臘人.
因為神的義正在這福音上顯明出來.
這義是本於信.
以至於信如經上所記.
二人必因信得生.
我想告訴大家.
這福音本是神的大能.
在一個譯本裡的意思是正在運行的大能.
弟兄姊妹.
從初期教會.
由聖靈降臨在教會.
到現在其實一直運行.
從來沒有變的.
他的能力昔日可以改變人.
今天同樣可以改變人.

$^{1321}$請你不要看低聖靈的能力.
請你不要相信耶穌之後.
我都感覺不到聖靈.
你當然感覺不到.
你都不肯讓聖靈住在你裡面.
你怎麼感覺呢.
什麼叫讓他住在裡面.
在他面前承認自己.
是多麼無能.
在他裡面的時候.
原來我們人不聰明.
在他裡面承認原來我們.
真的源源似水.
只可以依靠你.
我們就能夠迎接聖靈住在你裡面.
這樣你就能夠謙卑下來.
去聽上帝在你身上.
怎樣去改變你.
在第十七節那段經文的時候.
有兩個翻譯是這樣的.
這福音告訴我們.
上帝怎樣使我們在他眼中成為二人.
而這自始至終都是靠著信心來完成.
正如聖經所說.
二人藉著信心才能夠得到生命.
弟兄姊妹,我們的信仰從最初到現在都沒有變過.
我們一直都在神面前.
我們對上帝的信心.
什麼信心?.
上帝,我不可以,但你可以.
弟兄姊妹,這是你今天的救命符.
為什麼我要這樣說呢?.
在這幾年裡,我想大家應該都完完全全清楚.
變幻原是永恆.
每一分每一秒都在變.
我們很多時候眼睛一直看著這個世界的變化.
看著這個疫情,看著這個政局,看著這個環境.
看著所有的事情,我們心裡是怎樣的?.
跟著它,像股市一樣,一時上一時下.
一時上一時下.

$^{1361}$看著這個時間,那些經濟好一點,我又開心一點.
經濟不好,我又差一點.
剛才我聽到有弟兄姊妹說.
有些弟兄姊妹可能因為這個疫情的緣故而失去工作.
然後就很徬徨.
弟兄姊妹,但我想告訴你.
如果你相信上帝是可以的.
你相信上帝的話.
你就安心在上帝的裡面.
你去相信上帝在他的時間.
他一定會祝福一個在他裡面的義人.
只需要你保持你自己在神的裡面.
你一直維持你自己和上帝的關係.
在無論任何環境裡面.
你堅持去行善.
堅持去和別人建立好關係.
堅持繼續將這個福音傳給其他人.
我想告訴你.
上帝一定會按著他的說話.
他的應許.
他會將一切所需要的都會加給大家.
為什麼我這麼肯定呢?.
我經常說,記住弟兄姊妹,不要怕困難.
我在讀神學到現在.
其實上帝是針對著我來做的.
我這個人,如果你問我.
其實我對於經濟是很沒有信心的.
我對於經濟的事情是很擔心的.
可能我有一種開懷的責任感.
你養家的,你錢都不夠怎麼養家呢?.
我記得有一段日子裡面.
特別是我太太以前是病的那段日子裡面.
我們要賣掉房子,然後去醫.
醫到最後的時候.
我曾經試過一段時間.
連戶口裡面只剩下幾百塊.
你不知道怎麼搞的.
其實你不會想到什麼方法的.
但是我想告訴弟兄姊妹.
我不想這麼神化這件事.

$^{1401}$當時我們的童工是好的.
每一個月都有一次大家相聚.
我們就會請他們為我們經濟祈禱.
就這樣吧.
但是有一天的時候.
我又很失驚無神地收到一封信.
你收到這封信,我上去一看.
如果你收到的話,你一定很害怕.
因為這封信的整個款式.
是一封標心的信.
那些綁票的.
為什麼呢?.
全部在前面的那些.
那些郵票機.
但是全部的字都不是寫的.
全部是剪下來的,一格一格的字.
鄧太空全部是剪出來的.
那個地址是剪出來的.
接著我收到很害怕.
我說什麼事?.
這麼恐怖的.
但是我想告訴你,我一打開的時候.
裡面有一張支票.
一萬塊的支票.
接著時間上又是同一個格式.
用一張很小的書籤寫著.
弟兄,我們一起努力.
上帝為你預備.
弟兄,怎會這樣?.
如果我不經歷在這個困難當中.
我不會經歷到上帝為我預備.
我的人生裡面就是.
不斷不斷不斷經歷上帝的預備.
經歷上帝很多的化險為夷.
經歷很多裡面的時候.
雖然上帝沒有在那一刻幫我.
但是祂讓我平安度過.
所以在我們人生當中.
我們要有力量.
我們不是我們有多少.

$^{1441}$我們不是有什麼.
而是你的生命當中有沒有上帝與你同在.
所以求上帝幫助我們.
去相信上帝的大能.
所以今天我鼓勵.
如果你是信了耶穌的.
你應該是弟兄姊妹的.
讓我們好好回到神面前.
檢視我們的生命.
我們究竟今天有什麼卡住了.
讓我們生命當中.
沒有辦法彰顯上帝的大能.
我們有什麼卡住了.
令我們不可以很有效地.
去傳揚福音.
有什麼卡住了.
令我們經常都消消沉沉.
總是不開心的狀態.
什麼卡住了我們.
如果我們見到自己的問題.
請你好好回到神面前.
求你憐憫我們.
如果今天還沒有相信耶穌.
當中有朋友的話.
我鼓勵你.
其實當你看到的時候.
上帝在你身邊擺了很多的見證人.
我今天告訴你.
原來上帝.
祂所給你的生命是很精彩的.
聖經裡面說了.
你要得生命.
並且得的更豐盛.
不會假的.
所以求上帝幫助我們.
好不好?.
我好相信,特別是紅恩堂在這個社區裡面.
我每一次來的時候.
我都會為這裡禱告.
因為我以前也有去過宣德堂.

$^{1481}$做直堂的經驗.
上帝每一次擺教會.
在那個地方都有祂的意思.
對面有一些中產.
後面有基層.
附近也有一些基層的地方.
其實上帝在裡面的時候.
很多時候就透過我們的教會.
祂能夠將這麼大能的福音傳給他們.
有很多人正在需要當中.
但當你生命裡面.
你都沒有力量.
你怎樣將福音傳給他們呢?.
但求上帝幫助我們.
我們沒有什麼值得去高舉.
我們都是軟弱的人.
但是我們的上帝是大能的.
求上帝保守我們.
一個簡單的祈禱.
我們一起讀上帝的福音.
我們求你幫助帶領我們.
很願意我們來到你面前的時候.
我們知取你的力量.
最重要的是你告訴我們.
我們來到你面前.
我們需要謙卑下來.
我們需要跪在你面前.
我們知道上帝是你自己.
願意我們承認自己的軟弱.
承認自己的罪.
你已經赦免了我們.
你已經將你最寶貴的恩典賜給我們.
很願意我們經歷你的愛.
經歷你的大能.
因為我們看到仍然有很多在仇人當中.
有需要的人在當中.
我們求你親自去保守.
祈禱奉教主耶穌.
得勝名至其球.
多謝大家.

\newpage



\section{列王記上 19:1-18}
\label{sec:F4zrvJPl9fU}
\textbf{2021.07.04《你在這裡做甚麼?》 列王記上 19\_1-18 ( 和修版)  老耀雄牧師}
\newline
\newline
連結: \href{https://youtube.com/watch?v=F4zrvJPl9fU}{\texttt{https://youtube.com/watch?v=F4zrvJPl9fU}} ~~~~ 語音日期: 2021-07-04
\newline
\newline
\hyperref[sec:M84h3HizR5w]{\small{< < < PREV SERMON < < <}}
~
\hyperref[sec:index_chronic]{\small{[返順時目]}}
~
\hyperref[sec:index_scriptual]{\small{[返順卷目]}}
~
\hyperref[sec:LNNbZTaQW6E]{\small{> > > NEXT SERMON > > >}}
\newline
\newline
列王記上 19:1-18
\newline
\begin{longtable}{cl}
\hline
\hline
章節 & 經文 (和合本修訂版)\\
\hline
19:1 & \begin{tabularx}{0.7\textwidth}{X} 亞哈把以利亞一切所做的和他用刀殺眾先知的事都告訴耶洗別。 \end{tabularx} \\ \\ \relax
19:2 & \begin{tabularx}{0.7\textwidth}{X} 耶洗別就派使者到以利亞那裡，說：「明日約這時候，我若不使你的性命像那些人的性命一樣，願神明重重懲罰我。」 \end{tabularx} \\ \\ \relax
19:3 & \begin{tabularx}{0.7\textwidth}{X} 以利亞害怕，就起來逃命，到了猶大的別是巴，把僕人留在那裡。 \end{tabularx} \\ \\ \relax
19:4 & \begin{tabularx}{0.7\textwidth}{X} 他自己在曠野走了一日的路程，來到一棵羅騰樹下，就坐在那裡求死，說：「耶和華啊，現在夠了！求你取我的性命吧，因為我不比我的祖先好。」 \end{tabularx} \\ \\ \relax
19:5 & \begin{tabularx}{0.7\textwidth}{X} 他躺在羅騰樹下睡著了。看哪，有一個天使拍他，對他說：「起來吃吧！」 \end{tabularx} \\ \\ \relax
19:6 & \begin{tabularx}{0.7\textwidth}{X} 他觀看，看哪，頭旁有燒熱的石頭烤的餅和一壺水，他就吃了喝了，又再躺下。 \end{tabularx} \\ \\ \relax
19:7 & \begin{tabularx}{0.7\textwidth}{X} 耶和華的使者回來，第二次拍他，說：「起來吃吧！因為你要走的路很遠。」 \end{tabularx} \\ \\ \relax
19:8 & \begin{tabularx}{0.7\textwidth}{X} 他就起來吃了喝了，仗著這飲食的力走了四十晝夜，到了神的山，就是何烈山。 \end{tabularx} \\ \\ \relax
19:9 & \begin{tabularx}{0.7\textwidth}{X} 他在那裡進了一個洞，在洞中過夜。看哪，耶和華的話臨到他，說：「以利亞，你在這裡做甚麼？」 \end{tabularx} \\ \\ \relax
19:10 & \begin{tabularx}{0.7\textwidth}{X} 他說：「我為耶和華－萬軍之神大發熱心，因為以色列人背棄了你的約，毀壞了你的壇，用刀殺了你的先知，只剩下我一人，他們還要追殺我。」 \end{tabularx} \\ \\ \relax
19:11 & \begin{tabularx}{0.7\textwidth}{X} 耶和華說：「你出來站在山上，在耶和華面前。」看哪，耶和華從那裡經過。在耶和華面前有烈風大作，山崩石裂，耶和華卻不在風中；風後有地震，耶和華也不在其中； \end{tabularx} \\ \\ \relax
19:12 & \begin{tabularx}{0.7\textwidth}{X} 地震後有火，耶和華也不在火中；火以後，有輕微細小的聲音。 \end{tabularx} \\ \\ \relax
19:13 & \begin{tabularx}{0.7\textwidth}{X} 以利亞聽見，就用外衣蒙臉，出來站在洞口。聽啊，有聲音向他說：「以利亞，你在這裡做甚麼？」 \end{tabularx} \\ \\ \relax
19:14 & \begin{tabularx}{0.7\textwidth}{X} 他說：「我為耶和華－萬軍之神大發熱心，因為以色列人背棄了你的約，毀壞了你的壇，用刀殺了你的先知，只剩下我一人，他們還要追殺我。」 \end{tabularx} \\ \\ \relax
19:15 & \begin{tabularx}{0.7\textwidth}{X} 耶和華對他說：「去吧，從原路回去，往大馬士革的曠野去。到了那裡，你要膏哈薛作亞蘭王， \end{tabularx} \\ \\ \relax
19:16 & \begin{tabularx}{0.7\textwidth}{X} 又膏寧示的孫子耶戶作以色列王，並膏亞伯‧米何拉人沙法的兒子以利沙作先知接續你。 \end{tabularx} \\ \\ \relax
19:17 & \begin{tabularx}{0.7\textwidth}{X} 將來逃過哈薛之刀的，必被耶戶所殺；逃過耶戶之刀的，必被以利沙所殺。 \end{tabularx} \\ \\ \relax
19:18 & \begin{tabularx}{0.7\textwidth}{X} 但我在以色列中留下七千人，是未曾向巴力屈膝，未曾親吻巴力的。」 \end{tabularx} \\ \\
[1ex]
\hline
\hline
\end{longtable}
$^{1}$對一個事奉主的信徒來說.
一生所追求的其中一個恩典.
必然是事奉得力.
事奉者能夠不斷地從主領受事奉的力量.
事奉才不會乏力.
在教會我們經常聽到一個term.
就是一個人burn out.
即是說事奉到沒有動力.
由事奉得力到事奉無力.
中間必定經歷過一個過程.
怎樣能夠時時都事奉得力呢?.
或者怎樣才能夠不會事奉無力呢?.
這個都要有一些智慧.
今天我想引用《列王記》上十九章一至十八節.
引用以利亞的一段故事.
去和大家有三個信仰的反思.
去了解一下事奉者為何會陷入一個事奉的低谷.
以及怎樣可以重新得力的一個秘訣.
開始的時候我們一起去祈禱.
「創天道地的主,我們願意俯伏在你面前.
求你與我們親近.
用你的話語去洗滌我們每一個的心靈.
好叫我們能夠遵行你的旨意.
做一個行為與恩典相稱的基督跟隨者.
我們的祈禱是奉教主耶穌基督.
得勝名自祈求」.
在進入經文之前.
我都先講一下經文的背景.
當大衛王國分裂了.
大衛王朝分裂的時候.
不過以色列陷入一個偶像當道的年代.
整個以色列國.
由國王阿哈,皇后耶洗別.
以至整個民族.
都敬拜巴力和亞舍拉的偶像.
認為這些偶像能夠讓以色列風調雨順.
生產眾多.
因此他們就離開了耶和華.
先知以利亞領受了一個異象.
他對以色列說.

$^{41}$「我指著所事奉永生的耶和華.
以色列的神起誓.
這幾年如果我不禱告的話.
必不降露水.
也不下雨」.
以利亞先知代表了耶和華.
向整個以色列國發出了警告和審判的信息.
究竟誰才是掌管宇宙的神.
究竟誰最值得以色列百姓去敬拜.
究竟是耶和華.
還是你們所敬奉的巴力呢?.
於是由那天開始.
天真的不下雨了.
整個大地都落在旱災和饑荒之中.
以色列王阿哈於是派人出去尋找以利亞.
希望找到以利亞尋求一個解決的方法.
不過一直都找不到.
三年之後.
以利亞再領受異象.
要他去見阿哈.
以利亞對阿哈說.
請你派四百五十個巴力的先知.
以及四百個亞瑟拉的先知.
一起上加密山.
到時看看誰才是真神.
以利亞選擇在加密山是有原因的.
因為加密山是巴力的聖山.
以利亞讓巴力先知.
佔盡天時地利和人和.
也讓這些先知不能推搪.
當八百五十個巴力和亞瑟拉先知.
到達加密山之後.
以利亞就對以色列的百姓說.
你們心持異議要到何時呢?.
如果耶和華是神.
就應該順從耶和華.
如果是巴力的話.
就當順從巴力.
以利亞就和八百五十個假先知比試.
以誰能夠降火.

$^{81}$在祭壇的祭僧那裡.
誰就是真神.
因為巴力是雷神.
作為一個雷神.
降火應該是巴力的強項.
巴力和亞瑟拉先知.
由早上到晚上一起去祈禱.
都不能夠令火從天降下來.
以利亞當時還嘲笑他們.
說他們事奉的巴力是不是睡了.
然後以利亞就向耶和華禱告.
耶和華就降下火.
燒盡祭壇的勢僧.
當以色列的百姓見到火從天降下來.
燒盡一切祭僧的時候.
百姓就臉伏於地說.
耶和華是神.
以利亞就將八百五十個假先知殺了.
之後耶和華順從承諾.
以利亞再禱告的時候.
天就下雨了.
在這一刻.
耶和華和巴力的比試正式結束.
耶和華大勝.
以色列回轉.
重新敬拜耶和華.
以利亞在加密山的經驗到.
耶和華的同在.
此時此刻.
可以這樣說.
以利亞是達到她屬靈生命的高峰.
很可惜.
當皇后耶洗別.
向以利亞發出追殺令之後.
以利亞就開始逃命了.
從加密山一直走到荷列山.
荷列山即是西乃山.
是耶和華的聖山.
不過以利亞只是走了一天.
她就已經要求死了.

$^{121}$以利亞從一個人.
打敗八百五十個假先知的屬靈高峰.
一下子.
跌到要求死的屬靈低谷.
即是我們說的burn out.
究竟以利亞為何會burn out.
以及她如何處理自己burn out的情況呢?.
今天的講道經文.
就是以利亞求死.
和與耶和華相遇的一段對話.
作為我們一個信仰的反思.
第一個信仰反思就是.
事奉真的不要跟別人比較.
也不要以效果.
作為我們衡量事奉的標準.
在以利亞大勝巴力先知之後.
以利亞為何會陷入一個屬靈的低谷呢?.
是否因為耶洗別要追殺她呢?.
一般人都會認為.
以利亞當時是很擔心自己的安危.
尤其是十九章三節.
它這樣說.
以利亞害怕就起來逃命.
她害怕耶洗別的追殺.
以致她要逃跑.
不過如果你看一下.
應該這樣說.
我們手上沒有和合本修訂版的聖經.
如果你能有一本和合本修訂版的聖經.
你會在五十頁的左下角裡面.
看到一行細字.
上面寫著.
害怕是根據一些古卷和七十事譯本的.
原文是漢見.
所以以利亞逃跑的另一個原因.
不是因為她害怕.
而是她看到一些東西.
其實說以利亞害怕.
這個理由很牽強.
因為她之前已經隱藏了三年.

$^{161}$整個以色列都找不到她.
如果她不是因為受到異象.
要主動去見阿蝦的時候.
其實阿蝦仍然找不到以利亞.
所以如果以利亞要再刻意隱藏自己.
去逃避追殺的話.
耶洗別都應該和阿蝦一樣.
找不到以利亞.
而且作為一個先知.
她知道耶華會保護她.
耶穌經過加密山的經歷.
她明白耶華的大能.
如果耶華要定義保護她.
耶洗別的追殺令.
根本對她來說是毫無作用的.
所以以利亞不是因為害怕.
不是害怕被追殺.
不是害怕沒命而逃跑的.
那什麼令到以利亞陷入了屬靈的低谷呢?.
是因為以利亞看到一些東西.
正確地說.
不是以利亞的眼睛看到一些東西.
而是以利亞的屬靈視野.
令她有一種感覺.
而最後作出求死的反應.
在以利亞在羅騰樹下求死的時候.
她說了一句話.
這句話是.
耶華,現在夠了.
求你取我性命吧.
因為我不比我的祖先好.
以利亞認為自己不及她的祖先好.
以利亞所指的祖先究竟是誰?.
以利亞究竟和哪個祖先去比較?.
比較什麼?.
如果真的要比較的話.
就要有相類似的背景和經歷.
才可以有得比較.
如果是這樣.
相信以利亞是和摩西作出一個比較.

$^{201}$因為大家都是以耶華的名.
向一群外邦群體施展神跡旗幟.
摩西以十災向埃及的法老.
展示耶華是以色列的神.
以利亞以旱災,降火和降雨.
向巴力和亞瑟拉先知展示耶華是真神.
另一個相類似的背景就是.
摩西和以利亞同樣要帶領以色列民族歸回耶華.
摩西帶領以色列離開埃及.
做一群跟隨耶華的子民.
以利亞要帶領以色列離開偶像崇拜.
做一群敬虔的耶華子民.
可以這樣說.
以利亞求死是假的.
求死只是一種情緒發洩的表現.
事實上他想抱怨耶華是真的.
他想抱怨以耶華的大能.
為什麼不能夠像昔日一樣.
用大能一手讓整個民族回歸.
而且復興?.
以利亞心底想說昔日耶華協助摩西.
以旱災對付法老.
法老最後屈服.
整個以色列民族跟隨摩西離開埃及.
現在耶華同樣用旱災,天火,降雨的神蹟.
展示他的大能.
雖然以色列民願意回轉.
但是阿蝦和耶洗別仍然是無動於衷.
而且無視耶華的審判.
依然怪惡.
還派人追殺自己.
以利亞認為.
同樣是以神的名帶領以色列民.
為什麼摩西能夠超乎想像的結果.
神的大能能夠令到法老恐懼.
令百姓信服.
而我卻沒有呢?.
這就是以利亞認為自己不及祖先的緣故.
也是他在屬靈低谷裡尋死的第一個理由.
以利亞認為神與摩西同工.

$^{241}$可以有這樣的效果.
現在神與自己同工.
是不是應該都有相同的效果呢?.
很可惜.
以利亞看到的並不是這樣.
以利亞一旦與摩西比較.
就跌入了這個試探裡.
或者用另一個說法.
以利亞燃燒的其中一個原因.
就是他想看到的效果沒有出現.
當他見不到整個以色列.
包括阿蝦和耶洗別都願意回轉.
他見不到全民復興的情況.
他就與摩西比較.
也因為他與人比較.
以效果作為他事奉的核心.
他就對神抱怨.
對人埋怨.
以至他由一個屬靈的高峰.
跌落一個屬靈的谷底.
紅英家弟兄姊妹.
每一間教會都倡導全民皆兵.
每一間教會都希望每一個會友都投入事奉.
教會都很期望每一個弟兄姊妹.
都投入一個事奉的行列.
不過在事奉的過程裡.
千萬不要與人比較.
也不要用效果作為事奉成績.
是否好與不好的一個標準.
事奉者所追求的.
應該是多元而合一的事奉觀.
事奉者應該關注的.
是應該如何與神同工.
或者是他與神的關係.
而不是與人比較.
事奉一旦以效果作為我們的指標.
我們很容易跌入一個屬靈的陷阱.
以致我們不斷與人比較.
或者不斷以最好的成績作比較.
失卻了事奉者應有的態度.

$^{281}$那個態度就是.
以主耶穌作為我們事奉的中心.
主耶穌才是我們事奉的焦點.
我們一旦偏離了這個焦點的話.
事奉者就會走錯了一個跑道.
迷失了在事奉的十字路口裡.
洪恩嘉正值在疫情開始緩和的時候.
我也不斷向不同的弟兄姊妹邀請.
期望他們參與我們教會不同的事奉崗位.
我也很希望洪恩嘉每一個弟兄姊妹.
都能夠投入事奉.
但是我也在這裡作出勸勉.
當你踏出了這個事奉的時候.
千萬不要與人比較.
神給每一個人有不同的恩賜.
神給每一個人的恩賜有不同的發揮.
不要覺得別人比你好.
或者不要覺得我比別人好.
我們的事奉焦點永遠都是聚焦在主耶穌身上.
第二個信仰的反思就是.
事奉不應該只見到自己.
而見不到其他人.
以利亞其中一個尋死的理由.
就是他以為只有自己一個人在事奉.
也只有他一個人為神大發熱心.
他分別在第十節和第十四節.
重複地對耶和華說.
我為耶和華萬軍之神大發熱心.
因為以色列人背棄了你的約.
毀壞了你的壇.
用刀殺了你的先知.
只剩下我一人.
他們還要追殺我.
以利亞為神大發熱心是好事.
每個跟隨主耶穌的信徒都應該為神大發熱心.
但是以利亞認為.
只剩下我一人.
那就未必正確.
以利亞似乎忘記了兩件事.
第一就是在第十八章.

$^{321}$當他遇到阿哈的管家俄巴底的時候.
俄巴底已經跟以利亞說.
我將一百個先知藏起來.
每五十個藏在一個山洞裡.
拿餅和水供養他們.
難道俄巴底的行為不是為神大發熱心嗎?.
難道那一百個先知雖然藏起來.
就不是先知?.
就不是神的僕人?.
第二.
以利亞也似乎忘記了在加密山.
幫他殺那八百五十個先知的人.
那八百五十個先知不是以利亞一個人去做的.
是誰幫他去做的?.
所以以利亞認為只剩下我一人.
或者以利亞認為只有他一個人.
才是為神做事.
其他人就不是為神做事.
以利亞似乎自視過高.
或者可以這麼說.
以利亞在這一刻患了一種屬靈驕傲症.
最後神在十八節裡面親自回應他.
給他一個更新的視野.
他說:但我在以色列人中.
留下七千人.
是未曾向巴力屈膝.
未曾親吻巴力的.
現代中文修訂版是這樣翻譯的.
他說:但我要在以色列留下七千人.
這些人都效忠我.
沒有跪拜過巴力.
也沒有親吻過他的偶像.
現代中文修訂版強調這七千人.
是效忠耶和華.
沒有跪拜巴力.
耶和華想以利亞知道.
不只是只有你一個人效忠於我.
不只是只有你一個人去事奉神.
另外和合本.
不是修訂版.

$^{361}$和合本的譯法就是.
但我在以色列人中為自己留下七千人.
是未曾向巴力屈膝的.
未曾與巴力親嘴的.
和合本強調耶和華為自己留下七千人.
耶和華想以利亞知道.
不只是只有你一個人可以用.
可以被神使用的有七千人.
耶和華所顯示的人數.
是俄巴底的一百人.
多了七十倍.
耶和華想以利亞知道.
在你眼中只有你一個人在事奉.
只有你一個人效忠於我.
不是.
在神眼中.
在神的心意中.
還有很多人在為神事奉.
神要用誰都可以.
不是只有你,以利亞一個.
事奉不要看自己看得那麼重.
不要覺得只有你一個人在為神事奉.
不要覺得只有你一個人是為神大發熱心.
另外耶和華又將一個新的視野給以利亞看.
當以利亞將自己一個人.
和整個神國的僕人比較的時候.
神讓他有另一個角度去看.
其實在神的國裡.
還有很多人是未曾離棄神.
還有很多人沒有向巴黎屈膝.
而這些人全部都是神的僕人.
當以利亞看到整個以色列國都背棄了神的約.
都離棄神的話.
神卻讓以利亞知道.
還有七千個以色列人沒有離棄過神.
事情並不是好像以利亞所見到的一樣.
無論是俄巴底還是耶和華.
都看到有其他人同樣在事奉神.
可能事奉的內容不一樣.
事奉的角色不一樣.

$^{401}$好像俄巴底那樣.
他雖然是阿蝦的管家.
但是他秘密地去做.
去為神保護其他的先知.
又或者那批未曾向巴黎屈膝的先知.
同樣都是效忠耶和華的.
他們沒有放棄信仰.
可能他們沒有以利亞那麼被重用.
但是他們仍然都是神的僕人.
他們同樣都是在事奉神.
或者我可以這樣說.
以利亞固然是被神所重用.
這是毋庸置疑的.
但是被神重用.
就不等於其他人是沒有用.
被神重用.
並不是等於其他人是沒有用.
重用也不等同神不用其他人.
神可以重用你.
神也可以重用其他人.
以利亞在事奉過程裡.
這一點他只看到自己.
他看不到其他人.
他覺得自己比其他人更加重要.
甚至覺得自己所做的才是大發熱心.
因此一旦覺得事情不如意的時候.
就跌落了屬靈的谷底.
這就是他求死的另一個原因.
各位洪恩嘉弟兄姊妹.
我們很想在這裡提醒大家.
以利亞這一刻的事奉觀.
他只是以自己為中心.
無視其他人.
無視其他事奉的同路人.
我們要引以為教.
教會是一個基督的身體.
每一個都是屬基督的.
在基督裡面.
每一個肢體都很重要.
沒有誰高些.

$^{441}$沒有誰低些.
不是牧師就很厲害.
也不是會友就很差.
在教會裡面不是這個標準.
我們縱然事奉的方式不一樣.
我們事奉的範疇不一樣.
但我們同樣都是神所使用的僕人.
無分彼此.
而且我們都需要彼此的守望和待土.
第三個信仰的反思就是.
如何能夠從谷底裡面走出來的秘訣.
這就是回到神的身邊.
無論是以利亞主動的.
走去荷列山找耶和華.
還是耶和華透過天使引導.
以利亞去荷列山.
都是同一個目的.
那個目的就是要回歸到神身邊.
要在神的聖山裡面與神相遇.
唯有這樣.
以利亞才可以重新得力.
再次起步.
其實你看看以利亞.
很軟弱的.
他只是走了一天.
四十天的路程.
走了一天.
就要在羅藤樹下求死.
如果不是天使的幫忙.
提供餅和水.
我恐怕以利亞不能夠在曠野裡面.
走完餘下的三十九天.
神總是會在我們最軟弱的時候.
為我們提供不同的恩典.
托住我們.
如果說天使的幫忙是恩典的話.
耶和華刻意與以利亞相遇.
這是一個恩典中的恩典.
耶和華在荷列山第一次詢問以利亞.
他怎麼說呢?.

$^{481}$他說: 以利亞.
你在這裡做什麼?.
以利亞如何回應?.
他就說: 我為神大發熱心.
因為以利亞人背棄你的約.
毀壞你的壇.
用刀殺了你才知道.
只剩下我一個人.
你還要追殺我.
他們還要追殺我.
仍然在抱怨.
以利亞遇到神.
與神相遇.
仍然在抱怨.
對於以利亞的抱怨.
耶和華沒有正面去回應.
只是他站在山上.
站在耶和華的面前.
然後耶和華就在那裡經過.
在耶和華面前有烈風大作.
山崩石裂.
耶和華卻不在風中.
風後有地震.
耶和華也不在其中.
地震後有火.
耶和華也不在其中.
耶和華不在烈風中與以利亞相遇.
也不在地震中與他相遇.
也不在火中與他相遇.
究竟耶和華以甚麼方式.
與以利亞相遇呢?.
耶和華與以利亞相遇的情況.
就好像《出埃及記》33章.
耶和華與摩西相遇的情況.
一模一樣.
摩西說:求你顯出你的榮耀給我看.
耶和華說:我要顯示我一切的美善.
在你面前經過.
並要在你面前宣告耶和華的名.
我要恩代誰就恩代誰.

$^{521}$要憐憫誰就憐憫誰.
昔日耶和華在摩西面前經過.
此時此刻.
耶和華從以利亞所站的地方經過.
既然以利亞要與摩西比較.
既然以利亞求死時說:我不如先知.
不如早先.
耶和華就以與摩西相遇的地方.
和形式與以利亞相遇.
你這麼喜歡比較.
我就比較到底.
昔日我如何與摩西相遇.
今天我就用那個形式與你相遇.
恩典中的恩典就是.
耶和華對以利亞的一種做法.
是體貼入微.
耶和華在以利亞最軟弱的時候.
為她加能,賜力.
然後耶和華再以微小的聲音再問以利亞.
以利亞你在做什麼?.
耶和華重覆詢問以利亞.
有深層的意義.
就好像主耶穌在彼得三次不認主之後.
當主耶穌再遇到彼得的時候.
主耶穌三次問彼得.
你愛我比這些更深.
主耶穌要彼得再一次確認.
你是否真的跟隨我?.
重覆去問一個問題不是多餘.
而是讓人再一次去思考.
再一次去思想.
我們與神的關係.
我們是否真的跟隨主?.
我們是否真心願意跟隨主?.
今天的講題是.
你在這裡做什麼?.
當一個事奉者身心靈疲乏的時候.
當一個跟從主耶穌的僕人.
去燃燒的時候.
當一個信徒心灰意冷的時候.

$^{561}$能不能回到聖殿裡.
好好思考一下.
我究竟在事奉的崗位.
想做些什麼?.
你在這裡想做什麼?.
是大發熱心.
想自己做自己的事奉?.
還是以造就別人.
為你的事奉目標?.
是否以主的心意.
為主耶穌而去事奉?.
我2013年神學畢業.
在錦宣家事奉了七年多.
如果加上實習期.
我事奉了十年.
神學畢業的時候.
神學院院長提醒我們.
他說事奉的頭五年是高危期.
很容易就會心灰意冷.
離開職場.
過了五年會好一點.
應該適應得到.
如果說這十年的事奉裡.
你每天都這麼開心的話.
毫無疑問我一定是說謊.
不會的.
我一定有些事情是不開心的.
但當我遇到不開心或遇到困難的時候.
我必然會回到神身邊.
我多數會選擇默想或做良心審查.
當我默想到.
為什麼我要在這裡事奉的時候.
即是今天的講題.
我在這裡做什麼的時候.
我就會想起.
在我修讀神學的過程裡.
我曾經在主耶穌面前的立志和承諾.
這個立志和承諾就是.
我願意忠心去牧養.
主耶穌你所交託給我的羊.

$^{601}$當我再一次這樣去認定的時候.
錦宣家的會友.
就是主耶穌交託給我去牧養的羊.
我就知道我要承擔責任.
我就要處理好我不開心的那種情緒.
在去年四月十五日.
錦宣家正式差派我下來紅欣服侍.
我再一次問主耶穌.
你所交託給我的羊在哪裡.
我的答案很清楚.
你們就是我的羊.
你們每一個都是我的羊.
主耶穌說.
我問主耶穌.
我怎樣帶領我的羊.
主耶穌就說.
與聖靈同工.
用禱告托住你們.
與我同工.
用真理去建立你們.
讓他們的信仰能夠有根有機.
讓他們每一個都能夠活出信仰.
這個就是我來到紅欣的時候.
所領受的現象.
你們就是我的羊.
我用禱告托住你們.
我用真理建立你們.
因此我很清楚知道為什麼我要站在這裡.
也知道我為什麼要在這裡服侍.
以利亞你在這裡做什麼.
不是問以利亞在這裡想做什麼.
而是問以利亞.
你來到這裡的目的是什麼.
你來這裡是求死?.
你來這裡是抱怨?.
還是你來這裡與神相遇.
是再重新一次在神面前立志.
繼續為祂事奉?.
你是為主事奉.
而不是為自己事奉.

$^{641}$以利亞似乎不明白.
不明白.
不明白耶和華的一番心意.
所以我這次用另一個方式去引導以利亞.
十五節.
去吧.
從原路回去.
這句話可以譯為.
回到你原初的路上.
也就是說.
以利亞你回到你原先的事奉崗位上.
回到你原先事奉的道路上.
以利亞回到你以往事奉的崗位.
繼續為我事奉.
如果說在哪裡跌倒.
在哪裡站起來的話.
耶和華就要以利亞繼續為祂事奉.
不過是繼續按照神的吩咐去事奉.
因此耶和華吩咐以利亞去告三個人.
以利亞似乎明白了.
他不再去問.
或者他再沒有糾纏於自己不明白的地方.
他沒有再去發怨言.
沒有再自以為是.
他只好順服地按照耶和華的吩咐去做.
同樣的弟兄姊妹.
以利亞的經歷.
或者以利亞那種屬靈的信仰經歷.
其實對我們應該有一個很大的提醒.
我們事奉不要跟別人比較.
不要以果孝作為事奉的成敗標準.
更加不要以為只有自己才是為主而事奉的.
其他人對事奉不熱心.
以上一切都可能令我們遭遇一個死胡同.
最後會遭遇一個屬靈的低谷.
當然最好在我們還沒有跌入低谷之前.
懂得回到神身邊.
重新去思考.
我在這裡做什麼.
各位弟兄姊妹.

$^{681}$當你在洪恩家事奉的時候.
你問問你自己.
你在這裡做什麼.
你是追求一種自我感覺良好.
還是覺得你是在事奉神.
你是為造就別人的生命而做這個事奉.
還是為自己的榮耀而事奉.
當我們知道一切事奉的果孝和力量.
都是來自神的時候.
我們只需要按照神的心意.
就必能夠重新得力.
將神所交託給我們每一件事.
能夠榮神益人.
其實耶和華真的很愛這個以利亞.
在荷列山.
耶和華以微星和以利亞相遇.
在《列王記》下二章.
記載了他們以另一種相遇的方式.
耶和華要用旋風接以利亞升天.
這次耶和華是用旋風和以利亞相遇.
洪恩家弟兄姊妹.
我很深信主耶穌很願意.
向每一個渴求祂的信徒相遇.
主耶穌很願意與你們相遇.
很想每一個願意跟隨主耶穌的人相遇.
這個信徒是否你?.
你有沒有想過.
你用什麼方式與我們的主耶穌相遇?.
有沒有想過?.
或者你試過沒有?.
你嘗試過與主耶穌相遇沒有?.
你嘗試過擁抱主耶穌沒有?.
或者你被主耶穌擁抱過沒有?.
你想用主耶穌用哪種方式與你相遇呢?.
主耶穌給了我們一個應許.
尋找的就找到.
叩門的就給他們開門.
我們一起去祈禱.
你讓我們在以利亞的宿命生命裡.
曉得什麼叫做事奉的真義.

$^{721}$曉得我們的能力一切都是來自主.
一切都是恩典.
我們能夠事奉你是恩典.
叫我們聚焦在主耶穌你的身上.
而不是聚焦在自己的心意身上.
讓我們成為基督身體的一部份.
我們與其他的肢體是一個合一的肢體.
而不是一個個別的肢體.
我們不是一個人事奉.
我們是一群被召的人一起事奉.
為的是要建立主耶穌基督的身體.
也就是教會.
也就是我們洪恩家.
主求你恩待我們每一個.
讓我們每一個所做的.
都按著主你的心意而做.
以獻上感恩.
祈禱交託奉教主耶穌基督.
得勝名字其合.
阿們.
願主祝福大家.
\newpage



\section{彼得前書 1:1-6-7-22-25}
\label{sec:LNNbZTaQW6E}
\textbf{2021.07.11《奇妙的道》 彼得前書 1\_1,6-7,22-25 (和修版) 陳一華牧師}
\newline
\newline
連結: \href{https://youtube.com/watch?v=LNNbZTaQW6E}{\texttt{https://youtube.com/watch?v=LNNbZTaQW6E}} ~~~~ 語音日期: 2021-07-13
\newline
\newline
\hyperref[sec:F4zrvJPl9fU]{\small{< < < PREV SERMON < < <}}
~
\hyperref[sec:index_chronic]{\small{[返順時目]}}
~
\hyperref[sec:index_scriptual]{\small{[返順卷目]}}
~
\hyperref[sec:5qNMTaHZqr4]{\small{> > > NEXT SERMON > > >}}
\newline
\newline
彼得前書 1:1-6-7-22-25
\newline
\begin{longtable}{cl}
\hline
\hline
章節 & 經文 (和合本修訂版)\\
\hline
1:1 & \begin{tabularx}{0.7\textwidth}{X} 耶穌基督的使徒彼得寫信給那些被揀選，分散在本都、加拉太、加帕多家、亞細亞、庇推尼寄居的人， \end{tabularx} \\ \\ \relax
1:2 & \begin{tabularx}{0.7\textwidth}{X} 就是照父神的預知，藉著聖靈得以成聖，以致順服耶穌基督，又蒙他血所灑的人。願恩惠、平安多多地賜給你們！ \end{tabularx} \\ \\ \relax
1:3 & \begin{tabularx}{0.7\textwidth}{X} 願頌讚歸於我們主耶穌基督的父神！他曾照自己的大憐憫，藉著耶穌基督從死人中復活，重生了我們，使我們有活的盼望， \end{tabularx} \\ \\ \relax
1:4 & \begin{tabularx}{0.7\textwidth}{X} 好得到不朽壞、不玷污、不衰殘、為你們存留在天上的基業， \end{tabularx} \\ \\ \relax
1:5 & \begin{tabularx}{0.7\textwidth}{X} 就是為你們這些藉著信、蒙神大能保守的人，能獲得他所預備、到末世要顯現的救恩。 \end{tabularx} \\ \\ \relax
1:6 & \begin{tabularx}{0.7\textwidth}{X} 雖然你們必須在百般試煉中暫時憂愁，你們要為此喜樂 ， \end{tabularx} \\ \\ \relax
1:7 & \begin{tabularx}{0.7\textwidth}{X} 使你們的信心既被考驗，就比那被火試煉仍然能壞的金子更顯寶貴，可以在耶穌基督顯現的時候得著稱讚、榮耀、尊貴。 \end{tabularx} \\ \\ \relax
1:8 & \begin{tabularx}{0.7\textwidth}{X} 雖然你們沒有見過他，卻是愛他；如今雖看不見，你們卻因信他而有說不出來、滿有榮光的喜樂， \end{tabularx} \\ \\ \relax
1:9 & \begin{tabularx}{0.7\textwidth}{X} 因為你們得到信心的效果，就是靈魂的得救。 \end{tabularx} \\ \\ \relax
1:10 & \begin{tabularx}{0.7\textwidth}{X} 論到這救恩，那預先說你們要得恩典的眾先知已經詳細地搜索查考過， \end{tabularx} \\ \\ \relax
1:11 & \begin{tabularx}{0.7\textwidth}{X} 查考在他們心裡的基督的靈預先證明基督受苦難，後來得榮耀，是指甚麼時候，甚麼樣的情況。 \end{tabularx} \\ \\ \relax
1:12 & \begin{tabularx}{0.7\textwidth}{X} 他們得了啟示，知道他們所服事的不是自己，而是你們。那藉著從天上差來的聖靈傳福音給你們的人，現在將這些事傳給你們；這些事連天使也都切望察看呢！ \end{tabularx} \\ \\ \relax
1:13 & \begin{tabularx}{0.7\textwidth}{X} 所以，要準備好你們的心，謹慎自守，專心盼望耶穌基督顯現的時候帶給你們的恩惠。 \end{tabularx} \\ \\ \relax
1:14 & \begin{tabularx}{0.7\textwidth}{X} 作為順服的兒女，就不要效法從前蒙昧無知的時候那放縱私慾的樣子。 \end{tabularx} \\ \\ \relax
1:15 & \begin{tabularx}{0.7\textwidth}{X} 但那召你們的既是聖潔，你們在一切所行的事上也要聖潔； \end{tabularx} \\ \\ \relax
1:16 & \begin{tabularx}{0.7\textwidth}{X} 因為經上記著：「你們要成為聖，因為我是神聖的。」 \end{tabularx} \\ \\ \relax
1:17 & \begin{tabularx}{0.7\textwidth}{X} 既然你們稱那不偏待人、按各人行為審判人的主為父，就當存敬畏的心，度你們在世寄居的日子。 \end{tabularx} \\ \\ \relax
1:18 & \begin{tabularx}{0.7\textwidth}{X} 你們知道，你們得以從你們祖先傳下來虛妄的行為中救贖出來，不是靠著會朽壞的金銀等物， \end{tabularx} \\ \\ \relax
1:19 & \begin{tabularx}{0.7\textwidth}{X} 而是憑著基督的寶血，如同無瑕疵、無玷污的羔羊的血。 \end{tabularx} \\ \\ \relax
1:20 & \begin{tabularx}{0.7\textwidth}{X} 基督是神在創世以前所預知，而在這末世才為你們顯現的。 \end{tabularx} \\ \\ \relax
1:21 & \begin{tabularx}{0.7\textwidth}{X} 你們也因著他而信那使他從死人中復活、又給他榮耀的神，好讓你們的信心和盼望都在於神。 \end{tabularx} \\ \\ \relax
1:22 & \begin{tabularx}{0.7\textwidth}{X} 既然你們因順從真理而潔淨了自己的心靈，能真誠愛弟兄，就該以清潔的心彼此切實相愛。 \end{tabularx} \\ \\ \relax
1:23 & \begin{tabularx}{0.7\textwidth}{X} 你們蒙了重生，不是由於會朽壞的種子，而是由於不會朽壞的種子，是藉著神永活常存的道。 \end{tabularx} \\ \\ \relax
1:24 & \begin{tabularx}{0.7\textwidth}{X} 因為「凡血肉之軀的盡都如草，他的一切榮美像草上的花；草必枯乾，花必凋謝， \end{tabularx} \\ \\ \relax
1:25 & \begin{tabularx}{0.7\textwidth}{X} 惟有主的道永遠常存。」這話就是傳給你們的福音。 \end{tabularx} \\ \\
[1ex]
\hline
\hline
\end{longtable}
$^{1}$弟兄姊妹早安.
今天早上我們在崇拜當中.
無論我們所唱的詩歌.
又或者我們所聽到的獻唱.
或者是聖餐的過程當中.
大家都很感受到主的靈.
是感動我們的心.
特別是我自己坐在前面的一排.
感受到那份內心的震撼.
跟大家一起去敬拜.
今早牧師跟大家分享的題目.
就是《奇妙的讀》.
大家讀第一節的時候.
是否能讀出一些地方名字.
或者我們試試看第一節的經文.
大家.
這是最後的.
還沒這麼快說完.
還有一段時間才說完.
第一節經文記載了幾個地方名字.
這些都是很陌生的名字.
例如本都.
加拉泰.
加帕多加.
亞細亞.
比推尼.
其實你們的程序表也有寫出來.
這些名字第一.
我們很三普.
第二.
讀起來是否「甩不甩咳」.
都挺「甩咳」的.
到底這些是什麼地方來的呢?.
原來讀這些名字的地方.
就是集中在小亞細亞的地方.
這個小亞細亞的地方.
就是現在土耳其的版圖裡面.
所以這樣說你就明白多一點.
接著可能你會覺得很奇怪.
為什麼可以無緣無故寫一封信.

$^{41}$去給那些地方的人.
就是那些信徒.
其實也不是無緣無故的.
是有原因的.
為什麼呢?.
因為聖經告訴我們.
當地的信徒.
面對一個百般的試煉.
有沒有讀到這個字?.
一個如火的試煉.
這些字是描述當時的信徒.
在那些地方.
遇到有極艱難的生活.
極大的困苦.
所以叫做百般的試煉.
如火的試煉.
可能你會覺得奇怪.
為什麼那些人會這樣呢?.
原來是有背景的.
主要原因.
收信的那一群基督徒.
那一群信徒.
原來他們原先不是住在那裡的.
是移民移到那裡.
大家明白這個意思嗎?.
原本他們應該住在以色列的地方.
猶太人聚居的耶路撒冷.
或者巴勒斯坦的地方.
但是他們移民到那裡去了?.
去了小亞,細亞那邊.
問題浮現了.
這一群信徒.
第一.
他們做了僑民.
僑居在那裡.
第二.
他們的人數應該是少之又少.
為什麼呢?.
因為他們不是本地人.
他們是做僑民,做僑胞的.

$^{81}$所以人數是少之又少的.
第三.
那些人離鄉別井.
真的要適應新的環境.
是很不容易的.
就像以前.
我們移民的時候.
都有很多問題要處理.
很多的適應.
好了.
這一群猶太人.
還有一個情況.
他們是信猶太教的.
不過.
收信這一群猶太人.
原來他們放下了自己的猶太教.
改信了基督教.
你有沒有留意?.
所以就成為了基督徒.
如果這樣一層一層分析.
抽絲剝繭地看.
你就很快會知道.
原來那一群猶太信徒在當地生活.
難怪這麼辛苦.
難怪這麼多衝擊.
我再舉一些實例.
你會明白多一點.
大家知道我們香港地.
真的很有福氣.
因為我們可以請外傭.
是不是?.
又有非籍的.
又有印籍的.
有些像.
現在我們敢宣邀請了.
帕斯特爾芬迪.
或者帕斯特卡羅.
他們是在外地請回來的牧者.
目仰什麼?.
他們本身的信徒.

$^{121}$為什麼?.
因為印尼的.
這一群的外傭.
來到香港人數都是少的.
第二.
印尼.
大家都知道.
他們的宗教背景.
他們是回教的.
他們現在擺了回教.
改信了什麼?.
基督教.
所以你會發覺.
帕斯特爾芬迪.
或者帕斯特卡羅.
他們目仰的那一群對象的人.
在香港地來說.
是不是見少買少?.
是不是越來越少?.
本身已經少了.
現在還是更少.
你現在明白了.
為什麼彼得要寫信.
鼓勵他們.
跟他們打打氣.
是有原因的.
原因就是因為當時.
這一群信了耶穌的猶太人.
在當地的生活.
是非常艱苦的.
很簡單而已.
舉個例子來說.
找工作容易嗎?.
一點都不容易.
更加不要奢望找到高薪厚職.
什麼高薪良準.
找不到工作嗎?.
很難找到.
因為好的工作.
都被當地人做了.

$^{161}$是不是這樣說?.
有沒有一些油位?.
沒有.
豬頭骨就可能有.
所以你會發覺到.
當時這一群信徒.
工作上都遇到很多困難.
可能要做一些義務性的工作.
第二.
在身份上.
也出現問題.
因為他們僑居到別人地方.
我猜他們的身份.
最多頂籠的.
都是三等公民.
你認不認同?.
因為第一等公民.
是屬於羅馬籍的公民.
第二等公民.
就是當地的人.
第三等公民.
就是這些僑居.
在當地做僑民的猶太人.
所以他們的身份.
也不是很優勝.
所以我猜如果派一萬元.
不,不要一萬.
派五千元.
你猜他們有沒有份?.
可能未必有份.
因為住不夠七年.
所以弟兄姊妹們.
你會發覺到.
這群信了耶穌的猶太人.
無論在工作上.
在國籍方面.
在生活上.
其實也困難.
不過這些也是.
這些也是捱得住的.

$^{201}$都是支撐得到的.
最難是什麼呢?.
原來最難.
就是在他們的宗教生活方面.
原因是.
他們這群信徒.
已經不是當地.
一方面要適應當地的文化宗教.
但另一方面.
他們又脫離了猶太教.
現在信了耶穌.
所以在宗教信仰生活上.
這才是最艱辛的地方.
就是被其他信仰宗教的人.
來壓迫他們.
或者歧視他們.
甚至難聽一點.
可能逼迫他們.
為什麼呢?.
因為跟他們不相同.
所以你現在開始明白.
為什麼彼得要寫這封信.
給住在小亞西亞的猶太信徒.
是事必有因.
你現在明白了這個背景之後.
我們就會問第三個問題.
第三個問題是.
這樣嗎?.
這麼淒涼.
這麼辛苦.
這樣的環境.
到底他們有沒有出路呢?.
這是我們很喜歡小敏的問題.
他們有沒有希望呢?.
什麼是他們的出路呢?.
說到這方面.
原來他們是有出路的.
有盼望的.
出路是什麼呢?.
出路就是神的道.

$^{241}$所以彼得就用上帝的話.
神的道來鼓勵他們.
告訴他們.
你們千萬不要小看自己.
千萬不要洩氣.
千萬不要閉眼.
為什麼呢?.
因為你們每個人都擁有什麼呢?.
神的道.
你們信了主.
為什麼你們要信耶穌呢?.
是因為你們已經被神的道改變了你們.
所以.
神的道第一.
就是改變了他們的生命.
使得這些猶太信徒.
他們的人生觀.
他們的價值觀.
他們的世界觀.
都不同了.
是不是?.
差很遠的.
如果你的觀念改變了.
你看東西.
就會看得比較.
看得比較豁達一點.
大家明白我的意思嗎?.
如果你的觀念不改.
經常拿著以前的陳年舊物思想.
經常跳不出傳統的僵化思想.
你自己就會有很多的重擔.
所以神的道.
第一.
改變了他們的生命.
第二是什麼呢?.
第二就是提升他們的素質.
讓他們屬靈的生命的素質提升了.
因為他們現在是神的兒女.
所以在他們內在的素質方面.
他們能夠展現.

$^{281}$他們是有信心的.
有盼望的.
就算是辛苦.
困難.
或者是黑暗的環境.
都有明天的.
有信.
有望.
再加上有什麼呢?.
有愛.
原來他們的生命素質.
已經擁有信,望,愛的人生.
你說.
弟兄姊妹們.
他們是否找到出路了?.
有出路的.
原來上帝的道.
是這麼奇妙的.
上帝的道是這麼有能力的.
然後彼得就詳細告訴他們.
原來上帝的道.
神的道.
自低限度有四大功能.
今早牧師跟大家分享這四個功能.
上帝的道有哪四大功能呢?.
第一個功能是二十二節.
二十二節聖經怎麼說呢?.
聖經說.
因為你們順從真理.
就是神的道.
就潔淨了你們的心靈.
各位弟兄姊妹們.
原來上帝的說話這麼有能力.
可以將我們內心的問題.
處理得好.
大家有沒有留意到.
一個人生活得不開心,不快樂.
實境不多不少.
與心有沒有關係?.
有的.

$^{321}$很簡單.
如果這個人心胸窄.
每件事都很計較.
這個人何時得罪我.
那個人何時開口盯著你.
這個人何時不懂得做人.
如果記著這些東西.
我想問大家他開心嗎?.
不開心.
為什麼呢?.
心胸窄.
記著.
記著這件事.
有沒有辦法晚上睡得好呢?.
有沒有辦法睡得開心呢?.
為什麼呢?.
因為記著太多不好的事情.
這個就是心胸窄.
你說不單是心胸窄.
有些人是心地差.
心地差更麻煩.
為什麼呢?.
日想夜想都想到損人利己的事情.
你明不明白我的意思?.
怎麼睡得著呢?.
經常想著好處.
經常想著佔人便宜.
如果一個人心地這麼差.
他快不快樂?.
都不快樂的.
他自己知道自己的事.
他很怕有報應.
在他身上.
所以.
你說心胸窄也不好.
心地差也不好.
如果心骯髒.
更加倫潤.
如果一個人心骯髒.
就很倫潤.

$^{361}$為什麼呢?.
日想夜想.
都是怎樣做不好的事情.
做壞事.
如果一個人的內心.
充斥著這麼多問題.
試問.
怎麼可以活得輕散.
怎麼可以活得快樂呢?.
是不是?.
所以.
唯有什麼?.
唯有我們仰望上帝.
神的道.
來潔淨我們的心靈.
所以.
大家都知道.
人心比萬物都鬼咋.
壞到極處.
誰能識透呢?.
再下一幅.
麻煩你.
再下一幅.
大家知道.
現在我們在香港生活.
都要高度警惕.
為什麼?.
因為人心比萬物都鬼咋.
你知不知道.
原來在香港有很多騙子.
大家知不知道?.
所以.
有時候我們發覺自己聽電話都要很謹慎.
是不是?.
不要隨便回覆別人.
是不是?.
問你拿新聞證,副本.
好不好給他?.
不要了.
問你拿電話號碼好不好給他?.

$^{401}$不要了.
為什麼?.
原來香港有很多電話騙案.
是不是?.
那麼.
因為有電騙.
我現在才知道原來香港人這麼有錢.
特別是那些富婆.
隱形的.
原來這麼有錢.
被人騙了才知道這麼有錢.
各位頂尖的朋友.
人心比萬物都什麼?.
鬼咋.
壞到極處.
怎麼可以快樂呢?.
不可能的.
是不是?.
所以.
我們唯有仰慕什麼?.
神的道來潔淨我們.
原來上帝的道有一個洗滌的功能.
比你現在在超市買什麼什麼補品.
更加強力強勁.
你有神的話.
就將你心裡面的骯髒骯髒的東西洗淨乾淨.
難怪詩人說.
我用什麼潔淨自己的行為呢?.
有印象嗎?.
答案就是.
是要遵守你的道.
所以.
詩人說.
我將你的話藏在心裡.
免得我什麼?.
得罪你.
是不是?.
你的話是我腳前的燈.
是我路上的光.
是不是?.

$^{441}$所以頂尖的朋友.
原來有上帝的話在你的心中.
你的心靈才得到潔淨.
才活得輕散.
是不是?.
所以我們今時今日.
在一個複雜的社會裡面.
我們希望自己的內心更加清潔.
更加單純.
更加需要上帝給我們有清潔的心.
讀完這節聖經.
大家都可以放心了.
不一定要去換心臟.
或者要換血.
都不需要了.
我們主要是有上帝的話.
追求上帝的話.
在心中得到潔淨.
好了.
上帝的道.
神的道.
第二功能是什麼呢?.
就是真實相愛.
真實相愛.
原來我們人.
真的很少懂得.
去愛別人.
因為我們人的本性.
天性.
都是愛自己多一點.
是不是?.
都是很少去愛別人的.
這是一個天性.
再加上環境.
再加上環境的因素.
你會發覺到.
有時候遇到一些人.
說是很有愛心.
但原來那些愛心.
都不是正統的愛心.

$^{481}$不是合神心意的愛心.
為什麼呢?.
因為如果是合上帝心意的愛心.
就應該是來自羅馬書的十二章告訴我們.
那個模式.
就是說.
愛人不可虛假.
我要什麼?.
義務,善要親近.
有沒有印象?.
這才是上帝給我們一個.
真實相愛的模式.
是愛人不可虛假.
但很可惜.
今時今日.
有些人很少有愛心對別人.
因為愛自己多.
另一方面.
就算你去愛人.
可能帶有目的也未定.
為什麼這個人這麼關心我?.
為什麼這個人經常給我電話?.
為什麼這個人虛寒問暖?.
才發覺原來他要我買這個買那個.
你說是不是也有些失望?.
或者他對我這麼好.
原來他希望我日後會做.
對他好一點.
有時候也有些目的在當中.
所以你看看世界社會的潮流.
愛心也有很多不好的動機在當中.
所以我們不要要這些模式.
我們要上帝的道來改變我們.
叫我們學習真實相愛.
上帝的道給我們力量,能力.
可以真誠地去愛其他人.
這才是合神的心意.
有時候我被人問一條問題.
問得口啞啞.
問什麼呢?.

$^{521}$他說牧師.
為什麼你們教會裡面.
好像很多康復者?.
我聽聽是什麼意思呢?.
他說你們教會為什麼總是有這麼多.
要麼什麼不舒服.
要麼什麼有病.
要麼精神康復.
你明白我的意思.
為什麼這麼多這樣的人呢?.
有時候問這些問題.
真的不知道怎麼回答他.
可以將我的實況告訴他.
我告訴他.
是呀,沒錯.
教會之所以多這些康復者.
是正反映社會的實況.
今天我們的社會.
其實人情是很冰冷的.
如果你有什麼缺陷.
有什麼問題.
很多人很快將你拒諸門外.
或者不想接觸你.
甚至不想給你機會.
試問這些人.
在哪裡找到人間的友情呢?.
這些人在哪裡找到溫暖.
找到關心,找到支持呢?.
所以不覺得當有人帶這些康復者.
回教會之後.
這些康復者很投入教會的生活.
為什麼呢?.
因為他體會到教會這個家.
真的有神的愛.
教會這個家有溫暖的愛.
他們真的從心裡面來關愛,關心人.
這就是真實的相愛.
所以弟兄姊妹們.
這就是神的愛第二個功能.
我們再看第三個.

$^{561}$第三個功能是什麼呢?.
就是「重生了我們」.
原來神的道很奇妙.
神的道可以讓我們每一個人.
能夠死過活生生.
你說厲不厲害?.
真的.
還記不記得.
尼哥底姆問耶穌一條問題.
他說人怎麼可以重生呢?.
耶穌說.
「人若不重生就不能進神的國」.
於是尼哥底姆說.
我都已經老了.
怎麼可以再回媽媽的肚腹再生出來.
重生呢?.
耶穌就開始糾正他的觀念.
耶穌說.
「人若不是從水和聖靈生的.
就不能進神的國」.
水就是象徵到神的話語.
當我們有機會聽到神的道.
在教會聽到上帝的說話.
聖靈是感動我們的心.
於是我們就願意謙卑自己.
感覺到上帝的道加聖靈的工作.
讓我們看到自己的真面目.
所以我們的心就開始軟化.
願意謙卑接受主耶穌做救主.
所以這才是重生的意思.
各位弟兄姊妹們.
原來神的道.
是令我們每一個人.
可以擁有一個新的生命.
就好像一個新造的人.
人若不重生就不能進神的國.
這讓我們更加體會到.
原來神的道是這麼厲害的.
無論你過去有什麼背景.
或者你過去曾經是一個.

$^{601}$有不好習慣的人.
或者你過去心地不好.
思想不好.
有什麼問題也好.
頂至無謂也不需要擔心.
為什麼呢?.
因為上帝的道.
是可以將我們從這些罪惡深淵裡.
救拔出來的.
請看下一幅.
不知道大家認不認識這位牧師呢?.
這位牧師可能有些人會認識.
有些人未必認識他.
認不認識他?.
雖然他滿頭白髮.
但他很厲害的.
問一下有多少人認識他.
認識他舉手吧.
這位牧師為什麼這麼厲害.
我要提他的名字呢?.
他姓蕭的蕭牧師.
是因為他透過神的話.
神的道改變了很多人的生命.
令到很多在罪惡裡打滾的人.
能夠重獲新生.
你說厲不厲害?.
所以這位蕭牧師.
做了很多福音界賭的工作.
厲不厲害?.
賭徒走到山窮水盡的時候.
找蕭牧師.
蕭牧師就用神的道栽培他.
改變了他的生命.
所以現在很多賭徒都變成了教徒.
他不單可以幫助賭徒.
甚至有些人吸毒去找他.
他都可以幫助他.
這就叫做福音界賭.
他不單是福音界賭.
他還可以由福音界賭.

$^{641}$這個蕭牧師很厲害.
如果有些人沉迷在色情裡.
他還可以用福音界識.
你說厲不厲害?.
這位牧師真的透過神的道.
轉化了很多人的生命.
難怪這位蕭牧師.
有一句很精簡的口頭禪.
他怎麼說呢?.
他說:黃賭毒有得戒.
耶穌開你眼界.
厲不厲害?.
厲害.
原來神的道這麼有能力.
無論我們的生命過去怎樣敗壞.
不要緊.
只要你回到神的道裡.
接受神的話.
在你的心中.
你人生就可以得到重生.
所以我們的重生.
不是因為能壞種子.
而是不能壞種子.
就是神活潑常存的道.
最後看第四個功能.
神的道.
第四個功能.
原來是有一個永恆的價值在當中.
各位頂尖的姐妹們.
真的很感恩.
這個世界上.
有什麼東西是永恆不變的呢?.
你看看東,西,左,右.
你會發覺.
這個世界上所有你看到的事物.
都會改變.
但是唯有什麼呢?.
聖經說.
唯有神的道是永不改變.
耶穌說:天地要廢去.

$^{681}$我的說話一點一滑也不能廢去.
耶穌說:天地要改變.
我的說話是永不改變.
難怪以賽亞先知這樣說.
以賽亞先知這樣說.
他說:雨水如何由天而降並不返回.
同樣.
我的說話一出.
也不會徒然返回.
因為神的話是永遠安定在天.
各位頂尖的姐妹們.
當你擁有神的道.
你的內心的生命.
就有一個永恆的價值.
你的人生.
就有一個永恆的盼望.
為什麼呢?.
因為神的道是永遠不變的.
所以你看看這位詩人.
這位皇帝.
他做皇帝.
又寫了一些詩句.
在詩篇裡面.
這位是什麼皇呢?.
就是大衛皇.
或者下一章.
大衛皇.
他在做皇的時候.
是很顯赫一時的.
大家都知道.
他可以擁有很多東西.
不過大衛最後的總結.
就是詩篇23篇第一節.
是什麼呢?.
耶和華是我的牧者.
我必不自缺乏.
為什麼呢?.
因為他知道世界上的東西.
他擁有的.
無論是他的財富.

$^{721}$又或者是他的王位.
又或者是外邦的神.
或者是一些打贏仗.
這些都是會改變的.
不能夠存到永遠.
唯有神自己是永恆的.
所以他寧願擁有永恆的神.
做他的牧者.
他知道他人生.
必定是一個豐盛的人生.
一個永恆的人生.
所以各位電視節目.
什麼是真正有永恆的價值呢?.
原來就是神的道.
我們中國人很喜歡看一些歷史的書籍.
我相信在座當中.
很多人都看過三國志.
看過三國演義.
喜歡三國志,三國演義的.
舉舉手.
你有沒有喜歡看這一部份?.
你看三國志,三國演義的時候.
裡面最出名的詩句.
莫過於什麼?.
滾滾長江.
東逝水.
浪花途盡英雄.
是非成敗.
轉頭空.
是告訴我們.
在這個世界上.
你如何權傾天下.
你如何叱tak風雲.
其實你都是一剎那.
短暫就過去了.
所以各位電視節目的各位朋友.
什麼是永恆的價值呢?.
就是神的道.
所以今早牧師向大家推薦.
這個奇妙的道.

$^{761}$讓大家知道.
原來今天我成為基督徒.
成為神的兒女.
就好像彼得寫信.
給小小的猶太信徒.
讓他們知道.
因為擁有神的道.
即管他們遇到百般的試煉.
也不需要懼怕.
不需要擔心.
所以最後在結束的時候.
可能你會問.
既然上帝給了我們這麼奇妙的道.
祂有什麼期望呢?.
祂對我們有什麼探望.
我們去做的事呢?.
我相信有兩點是很真實的.
上帝將奇妙的道賜給我們每一個人.
讓我們經歷當中.
上帝奇妙的功能.
原來有兩點是上帝想我們做的.
第一.
上帝不想我們獨食.
因為獨食難肥.
所以上帝想我們將祂的道.
要拿出來.
要分享出來.
要拿出來跟別人分享.
所以我看到現在大家弟兄姊妹.
都很懂得用社交媒體.
差不多早上你當然收到一些金句.
收到一些靈修小品.
是很好的.
如果我們能夠有什麼屬靈的心得.
有什麼神的話.
是很激勵自己.
你可以拿出來跟別人分享.
互相勉勵.
你記得要將神的道傳開去.
讓多些人都聽到神的道.

$^{801}$以致生命得到改變.
第二個目的.
我相信天父都想我們做的.
就是希望我們將神的道.
活現在我們的生命當中.
換句話來說.
上帝很希望你將祂的道.
能夠將祂生活化.
在我們每人的工作裡面.
與人相處的當中.
活現了神的道.
讓人一看到你.
就很明顯看到.
你這個人是與別不同的.
為什麼?.
因為你有神的話語.
在你的生命當中.
這個見證是一個最有效的見證.
我們希望弟兄姊妹.
都將這兩個目的.
上帝的心意操練出來.
我們低頭祈禱.
親愛的主.
多謝你透過彼得寫這封信.
給小亞,細亞,猶太的信徒.
再一次提醒我們.
原來上帝你的話語.
是何等的寶貴.
真實.
大有能力.
所以我們今天的生活當中.
我們必須要將你的話.
藏在心中.
實踐在我們的生活裡面.
與人分享.
以致我們每一天.
都能夠知取到.
從神的道裡面.
可以得到生命的餵養.
生命的力量.

$^{841}$聽我們的禱告.
讓我們能夠.
活得更加快樂.
活得更加豐盛.
是因為我們有神的道.
湧現在我們的心中.
多謝耶穌.
我們的祈禱.
奉耶穌的名求.
阿們.
\newpage



\section{腓立比書 2:25-30}
\label{sec:5qNMTaHZqr4}
\textbf{2021.07.18《隱藏的事奉》 腓立比書 2\_25-30 (和修版) 曾敏姑娘}
\newline
\newline
連結: \href{https://youtube.com/watch?v=5qNMTaHZqr4}{\texttt{https://youtube.com/watch?v=5qNMTaHZqr4}} ~~~~ 語音日期: 2021-07-20
\newline
\newline
\hyperref[sec:LNNbZTaQW6E]{\small{< < < PREV SERMON < < <}}
~
\hyperref[sec:index_chronic]{\small{[返順時目]}}
~
\hyperref[sec:index_scriptual]{\small{[返順卷目]}}
~
\hyperref[sec:eRsLDyLgpD0]{\small{> > > NEXT SERMON > > >}}
\newline
\newline
腓立比書 2:25-30
\newline
\begin{longtable}{cl}
\hline
\hline
章節 & 經文 (和合本修訂版)\\
\hline
2:25 & \begin{tabularx}{0.7\textwidth}{X} 然而，我想必須差以巴弗提到你們那裡去。他是我的弟兄、同工和戰友，是你們差遣來供應我需要的。 \end{tabularx} \\ \\ \relax
2:26 & \begin{tabularx}{0.7\textwidth}{X} 他很想念你們眾人，並且極其難過，因為你們聽見他病了。 \end{tabularx} \\ \\ \relax
2:27 & \begin{tabularx}{0.7\textwidth}{X} 他真的生病了，幾乎要死。然而神憐憫他，不但憐憫他，也憐憫我，免得我憂上加憂。 \end{tabularx} \\ \\ \relax
2:28 & \begin{tabularx}{0.7\textwidth}{X} 所以，我更要盡快送他回去，好讓你們再見到他而喜樂，我也可以減少憂愁。 \end{tabularx} \\ \\ \relax
2:29 & \begin{tabularx}{0.7\textwidth}{X} 故此，你們要在主裡歡歡喜喜地接待他，而且要尊重這樣的人， \end{tabularx} \\ \\ \relax
2:30 & \begin{tabularx}{0.7\textwidth}{X} 因他為做基督的工作不顧性命，幾乎至死，為要補足你們供應我不夠的地方。 \end{tabularx} \\ \\
[1ex]
\hline
\hline
\end{longtable}
$^{1}$弟兄姊妹平安.
平安.
讓我們可以回到祂的家裡.
親近祂.
敬拜祂.
我們開始的時候再有一個禱告.
顧念我們的卑微.
讓我們可以有這樣的福氣.
去敬拜讚美你.
這一刻.
求你真的監察我們每一個的生命.
看看在我們裡面.
我們有沒有得罪你,得罪人的地方.
如果有的話.
懇求你赦免潔淨.
免得這一切的罪惡.
難了我們,去到你的根前.
聖靈啊.
今天願意在我們每一個心裡面動功.
讓我們聽到.
上帝要我們聽到的說話.
也感動我們.
作出悉切的回應.
我們今天的禱告.
是奉耶穌基督的名字.
阿們.
今天的經文.
其實比較少在講壇上.
被人去講的.
但是這段經文.
是很值得我們去學習的.
在這裡.
我想說.
因為最近我們教會.
都在重整不同部門的.
事奉.
大量的招募人手.
所以也是一個非常好的時間.
我們去說說.
關於事奉方面的教導.

$^{41}$我相信很多靈性姐妹.
在思考.
我想超越什麼崗位的事奉呢.
多少個崗位呢.
怎樣是好的呢.
我的恩賜又如何等等.
關於這一系列的問題.
我都會有不同的講道的.
我希望都可以.
涉獵到這些不同的層面.
今天首先要分享的是.
究竟.
土神喜悅的事奉.
是怎樣的.
土神喜悅的事奉.
是怎樣的.
事奉現在不是說我要做一些事情.
是不是.
我要做一些事情.
有份工作.
有些事情我要做.
我要做一些事情.
現在是事奉.
我們就去想想.
怎樣去討上帝的喜悅呢.
在回答這個問題之前.
我們首先要問.
什麼是事奉呢.
提醒一下靈性姐妹.
有沒有很快的.
可以回答我呢.
什麼是事奉呢.
這個名詞在教會裡.
有沒有人嘗試回答.
教會工作.
好的開始.
有沒有其他的分享.
什麼是事奉.
還有沒有.
靠著主的能力.

$^{81}$做祂想我們做的事情.
靠著主的能力.
做祂想我們做的事情.
都是很好的一個發揮.
大家繼續想的時候.
我又嘗試去理解一下.
事奉的原文翻譯.
就是服務.
為人服務.
或者是為神工作.
都可以稱為事奉.
有時候事奉這個詞語.
也可以翻譯為敬拜.
原來我們敬拜神的時候.
是在事奉.
今天我和你一起敬拜神.
我們都是在事奉.
事奉的內容.
非常豐富.
有剛偉性的.
有非剛偉性的.
無論我們是服務人.
是為神工作還是敬拜.
我們都有一個中心.
唯一一個中心就是上帝.
我們是為了神而服務人.
為了神而工作.
也是為了神而敬拜.
敬拜的對象.
更加只有神一位.
所以簡單來說.
我們事奉的對象就是神.
為什麼我們要事奉神呢?.
希伯來書十二章二十八節這樣說.
所以既然我們得了.
不能被震動的國度.
就要感恩.
照著神所喜悅的.
用虔誠.
敬畏的心事奉神.

$^{121}$因為我們的神.
是吞滅的火.
這個不能被震動的國度.
是指神的國度.
是指神掌權的所在.
我們能夠相信耶穌進入神的國度.
是神所賜的恩典.
和福氣.
不是我們所配得的.
所以當我們得到這麼大的福氣.
這麼大的恩典的時候.
我們就應該感恩.
而感恩就有行動的.
除了口說.
為我成就了這麼多恩典之外.
是我們用行動.
我們用虔誠和敬畏的心去服侍神.
這是第一點.
因為要感恩所以要服侍.
第二點,彼得前書四章十一節說.
若有人服侍.
他要按著神所賜的力量服侍.
好讓神在凡事上.
因耶穌基督得榮耀.
第二點,原來我們服侍神.
是為了榮耀神.
不是為了榮耀自己.
我們服侍不是為了榮耀自己.
我們有了這些很重要的條件.
說服侍是怎樣的時候.
接著我們就開始探討.
究竟土神喜悅的事奉是怎樣的呢?.
今天是看其中一點.
我們透過一個人物去學有關的功課.
這個人物的名字叫.
爾巴忽提.
爾巴忽提是出現在.
菲勒比書裡面的.
菲勒比書寫作的時間.
大約是主後62年.

$^{161}$當時保羅.
保羅就是菲勒比書的作者.
寫這本書的人.
當時在羅馬.
當時在羅馬坐牢.
他就寫信給菲勒比教會.
這間教會是他在歐洲建立的.
第一間教會.
保羅很厲害.
你可以這樣稱呼他.
他是一個文學家.
他寫了13卷書信.
他也是一個很偉大的.
牧者,一個宣教士.
沒有人不認識保羅.
信耶穌.
一段時間都會知道保羅是誰.
是很出名.
而這個菲勒比教會.
是他在歐洲建立的第一間教會.
在菲勒比書的二章.
保羅稱讚兩位弟兄.
為事奉者的典範.
其中一位是提摩泰.
提摩泰也是聖經裡.
很出名的一個人物.
同樣大家都非常熟悉.
另外一位就是爾巴忽提.
反而這一位.
就不是那麼多人認識他.
今天.
我只想說這個.
不是很多人認識.
不是很被人認識的.
爾巴忽提.
他的生命.
是很值得我們欣賞的.
爾巴忽提.
是誰呢?.
保羅用三個詞語去形容.

$^{201}$他和自己的關係.
是我的弟兄.
是我的同工和戰友.
究竟爾巴忽提在教會裡.
是一位領袖.
還是普通的信徒.
我們是沒有充足的資料.
去了解的.
我們只可以這樣說.
他不是使徒.
保羅形容他是弟兄.
是要表達出.
他們在主裡面的關係.
引起大家.
有同一個信仰.
所以大家的關係.
是緊密的.
信仰將他們結連在一起.
保羅再進一步說.
他是同工.
可能是因為.
爾巴忽提在福音的工作上.
曾經和保羅合作過.
他們一同在.
傳福音的工作上.
服侍上帝.
彼此有一個共同努力的目標.
所以很能夠.
明白對方.
這是非常寶貴的.
所以.
被稱為同工.
同工的關係.
比弟兄姊妹的關係再進一步.
在傳福音的過程.
我們都知道.
這也不簡單.
因為要從撒旦的手中.
搶救靈魂.
就像打仗一樣.

$^{241}$是非常厲害的.
所以保羅也稱呼.
爾巴忽提是戰友.
大家在.
搶救靈魂的事工裡.
是互相支持的.
互相保護對方.
撐著對方.
不能夠失去對方.
在這個過程.
戰友的關係.
比同工更緊密.
更重要.
一層一層上.
在這三個描述裡.
我們可以看到保羅和爾巴忽提的關係.
是非常緊密的.
爾巴忽提.
對於保羅來說.
是很重要的.
雖然是這樣.
我們會發現.
爾巴忽提在聖經裡.
是一個不太紅的人物.
用世俗的字眼說.
不是很紅.
不是很出名.
牽涉的文字不多.
保羅是紅很多.
從這個角度去看.
是很出名的.
基本上來說.
爾巴忽提是一個很低調的人物.
他雖然低調.
卻是上帝很欣賞的人物.
為什麼上帝.
這麼欣賞他呢.
我們來了解一下.
爾巴忽提.
實際的身份是什麼呢.

$^{281}$我們只能了解到.
他是菲律賓教會所差遣的.
所以他是菲律賓教會的代表.
在這段.
短短的經文裡.
提到他事奉的內容.
他事奉的內容是什麼呢.
原來當時保羅.
在羅馬坐牢.
很需要教會在經濟上的支持.
而當時的羅馬政府.
是不會為囚犯提供什麼物資.
去滿足他們的需要.
所以囚犯.
都很依賴親友.
去供應.
供應什麼呢.
包括了.
有日常的食物.
有日常的生活用品.
保羅當時的處境.
是這樣的.
沒有人供應.
就沒有東西了.
就沒有滿足.
有人供應.
生活才能過得去.
所以當時的菲律賓教會.
就派伊巴忽提.
去羅馬的地方.
將教會裡的.
那些的饋贈要送給保羅.
或者可以這樣說.
教會裡的一些奉獻.
送給保羅.
去供應他生活所需要的.
另外伊巴忽提有個任務.
是要留下.
陪伴和服侍保羅.
我們要記得.

$^{321}$當時保羅正在坐牢.
如果伊巴忽提.
來就是要這樣陪伴他.
陪伴坐牢的保羅.
要服侍在監獄當中的保羅.
很明顯.
我們聽得出.
伊巴忽提的服侍是比較後方的.
隱藏的.
不是很前線的.
如果不是.
保羅刻意提到他.
沒有人會留意他.
正如今天.
電影姐妹.
大家回來崇拜.
還有.
正在看直播的電影姐妹.
那些感覺應該更加深刻.
看直播的電影姐妹.
看到什麼?.
只是看到鏡頭前.
看到的人.
領唱.
和獻唱的呼音樂曲隊.
主席.
講員.
報告和祝福.
就只有這些.
其他都看不到.
這些可以說是.
比較前線的.
很顯露的事奉.
大家看到的.
在場的電影姐妹會看到多一點.
但是.
有很多事奉人員.
默默做了很多事.
大家看不到的.
除非.

$^{361}$我們刻意提起他們.
舉個例子.
我手上有一個程序表.
大家手上都有程序表.
這個程序表.
原來都經過不同人的手.
不同的人負責不同部分.
但從來裡面不會寫下名字.
讓你聽.
這裡是這樣,那裡是這樣.
到你手上的時候.
你只看到程序表.
你不知道背後的事奉者.
大家今天看到.
台上有一盤花.
每個星期都有人插花.
我們的影音.
也有不同的人輪流.
還有很多.
其實.
有很多人的事奉.
在背後默默做很多事.
我們不一定知道他們的名字.
但每個星期.
他們都在做.
很低調的,後方的.
隱藏的.
是不是這些人.
比不上在前線的.
前方的.
顯露的人.
好像有分開的,這些看不見的人.
他們沒有那麼重要.
看得到的人是重要的.
有時我都會聽一些弟兄姊妹這樣說.
我只是走來走去.
隨便收拾一下.
我只是閒人.
不重要的.
好像分開了,某些事奉特別重要.

$^{401}$某些事奉不重要.
其實沒有分開的.
你是後方的.
別人看不見你的.
你是前方的,別人看得見你的.
經常看得見你的.
是一樣的重要.
是一樣的有價值.
西人的眼睛總是覺得.
顯露的東西是最寶貴的.
最有價值的.
在上帝眼中從來不是這樣看.
上帝看我們.
是一視同仁.
我最近有兩個星期六.
坐在教會,都有種很感動的感覺.
我星期六就見到弟兄姊妹在預備.
有很多不同方面的預備.
我見到那些工種.
很細緻的.
各自在做自己的角色.
我為每一個弟兄姊妹感恩.
我很想說.
你們未必認識他們.
未必知道他們在做什麼.
但他們在努力.
他們就像耳包忽提.
默默在做很多事情.
在這裡很重要.
很重要.
我們去看看.
這個後方隱藏的事奉.
究竟有什麼.
令上帝去欣賞他.
是否耳包忽提的態度.
在態度上有什麼令人欣賞呢.
有三點很值得去看.
第一點,他繼續保持低調.
他本身的事奉就是低調.
他繼續低調.

$^{441}$耳包忽提沒有到處跟人說.
我將要將教會的大筆數.
送去保羅那裡.
他沒有到處說.
如果不是因為健康的緣故.
因為他健康不好.
保羅要差遣他回到菲律賓去.
如果不是因為保羅刻意提起這件事.
相信我們對耳包忽提的事情了解得很少.
聖經並沒有說.
於是耳包忽提的事跡.
傳遍了哪裡哪裡.
大家對耳包忽提越來越認識.
可能我們會說,那個年代沒有手機.
也沒有使用那麼多社交媒體.
所以沒有辦法拍那麼多照片.
也沒有辦法使用社交媒體.
IG,Facebook等等.
也沒有辦法拍廣告.
於是就沒有人認識他.
也不是的.
沒有社交媒體.
一樣有方法可以令自己很紅.
我見人就說.
我和你說.
我現在在家.
很重要的一個使命.
我要拿這些奉獻去哪裡哪裡.
說完之後見下一個又說.
輝哥,我告訴你我接下來要做什麼.
我每一個都說.
一個全十十全百.
很厲害的,可以的.
如果那個是一個很高調的.
很想紅的人一定有方法.
但你見到耳包忽提沒有這樣做.
他就安靜地做.
他要做的事.
我想問今天我們是不是很甘於.
參與後方的事奉呢.

$^{481}$我們不是甘於參與比較隱藏的事奉呢.
我們心裡有沒有比較.
為什麼那些人.
只會某某某某.
為什麼他們.
只能站出來事奉.
我也很想站出來事奉.
讓人認識我,行不行呢.
有沒有這樣的心態呢.
有沒有問過上帝.
我們想參與什麼事奉.
是不是前方,是不是後方.
是上帝為我們安排.
上帝看什麼是適合.
上帝安排前方.
就是前方,是不是.
站在人前,上帝安排後方.
隱藏的,就專心做後方的事.
在參與的時候.
我們有沒有安靜地做我們要做的事呢.
還是心裡經常比較.
經常不開心呢.
總是不甘心呢.
這些就是在後方.
做隱藏事奉的人.
對於那些參與前方的.
比較顯露服侍的人.
那需不需要保持低調呢.
要的.
那你說我站在前方.
全世界都看著我還用低調.
都知道我了.
不是的,低調和不低調有分別.
我想讀給大家聽一段經文.
是很經典的.
《馬太福音六章》一到六節是這樣說的.
你們要謹慎.
不可故意在人面前表現虔誠.
叫他們看見.
若是這樣.

$^{521}$就不能得你們天賦的賞賜了.
所以你施捨的時候.
不可叫人在你前面吹號.
將那假冒為善的人.
在會堂裡和街道上所做的.
故意要得人的稱讚.
我實在告訴你們.
他們已經得了他們的賞賜.
你施捨的時候.
不要讓左手知道右手所做的.
好像你隱密地.
隱蔽地施捨.
你負在隱蔽中察看.
必然賞賜你.
你們禱告的時候.
不可將那假冒為善的人.
外站在會堂裡和十字路口禱告.
故意讓人看見.
我實在告訴你們.
他們已經得了他們的賞賜.
你禱告的時候.
要進入內室.
向那在隱蔽中的父祖告.
你負在隱蔽中察看.
必將賞賜你.
我們做什麼都好.
其實呢.
說的是這一班.
其實是一班領袖.
很喜歡在人面前.
故意做一些事情.
讓人看到就可以稱讚自己.
本身他的位置已經很顯露.
很前方.
還要故意顯露自己.
再顯露多一點.
讓人稱讚自己.
讓自己得榮耀.
心態決定最後的生命是怎樣.
如果我的心態很喜歡.

$^{561}$自己得榮耀的時候.
我一定會令自己.
敲鑼敲鳴鳴.
全世界都知道我在做什麼.
但我低調的時候.
就算我在前方.
我都可以做到很低調.
我曾經見過不同的牧者.
是真的很低調.
除了他要上台講道.
除了他要去做一些服事的時候.
其他時候他真的很謙卑.
很低調.
很安靜.
不會到處表露自己怎樣怎樣.
以致有時候你會有種感覺.
他坐在那裡.
如果他不出聲.
你還以為他是路人甲乙丙丁.
你不知道.
原來是某某牧者.
他真的不覺得.
所以可以做到很低調.
做前方顯露事奉的人.
可以保持低調.
忠心,安靜,做好我們的事情.
為什麼?.
剛才我一開始的時候說.
事奉是為了什麼?.
因為我們要感恩.
事奉是為了榮耀誰?.
榮耀上帝.
不是榮耀自己.
所以無論今天我在前方.
在後方服事.
保持低調.
以巴弗提在後方服事.
他繼續低調.
所以上帝很欣賞他.
第二件事.

$^{601}$很重要的,以巴弗提的服事.
得人欣賞的地方,他是在做神的工作.
不是做人的工作.
保羅在這裡三十節.
保羅讚以巴弗提.
做基督的工作.
不顧聖名,幾乎致死.
對於人來說.
以巴弗提的工作.
是要供應保羅的需要.
要陪伴他,要服事保羅.
這個很明顯.
是做人的工作.
但是在神的眼中.
這是做基督的工作.
因為如果沒有了以巴弗提的供應.
保羅的生命都難以維持下去.
他就不能繼續傳福音.
保羅不坐牢的時候.
他還可以自己出去找兼職.
可以織帳篷維持自己的生活.
但是現在他在坐牢.
他不可以做兼職.
是完全靠別人的供應.
在這個時候.
上帝就透過以巴弗提的供應.
去支持他的生活.
讓他可以繼續服事.
無論在監獄裡.
或者將來他出監獄.
都可以繼續服事.
所以以巴弗提的服事.
都是在做基督的工作.
上帝委派他.
這不是單單服事保羅那麼簡單.
而以巴弗提服事保羅.
不只是看自己和保羅有多友好.
多老友.
我和他是朋友.
服事他好一點,做多一點.

$^{641}$你不要以為.
基督徒沒有這樣的層面.
有的.
我們基督徒都看的.
我和弟兄姊妹和那個是不是比較友好.
我做的事做得不一樣.
這個沒那麼友好.
沒那麼好,做的時候沒那麼賣力.
都會的.
他不是,以巴弗提不是看這些.
雖然保羅形容他和自己的關係很緊密.
但是他看的時候.
是一個服事,是上帝的工作.
服事上帝.
無論怎樣都會做得很好.
今天我和你的服事又如何呢.
究竟最終極的.
我們是服事人.
還是服事耶穌基督.
這個有很大的分別.
如果純粹是出於人情.
為了人的緣故.
在神的標準裡.
還未達標.
我們應該是為上帝的緣故.
為神做到最好.
這是第二點,以巴弗提.
是令上帝很欣賞他的地方.
他是做神的工作.
不是人的工作.
第三點.
就是以巴弗提不稀罕.
這些後方的.
隱藏的服事.
他拼盡全力.
27節這樣說.
他真的生病了.
幾乎要死.
30節,引他為做基督的工作.
不顧性命,幾乎致死.

$^{681}$應該按照當時的情況.
他應該是運送.
這一筆奉獻.
或者我們叫做.
去羅馬的途中.
是生病了.
他病得很厲害.
大家就不要想.
只是起義病.
今天我們誰沒有病過.
只是很平常.
病在這個年代就很平常.
看醫生,吃藥.
甚至有些人的抵抗力很好.
自然痊癒.
病沒什麼大不了.
但當時來說.
醫學不是那麼發達的時候.
病是一件很嚴重的事情.
更何況他這個病病得很厲害.
就更加的嚴重.
有生命的危險.
換作其他人是他.
病得很厲害.
他在途中.
你猜他會做什麼事呢?.
回頭吧.
不要繼續去.
我病了.
我病得很嚴重.
我想回頭要好好照顧自己的身體.
我不去羅馬了.
我想回菲律賓.
我想見我的親人.
我很思念他們.
經文裡有說他很思念家鄉的人.
我回頭吧.
其實也不差.
保羅的貢獻.
我回去之後.

$^{721}$再找別人拿.
遲一點而已,沒問題.
一定到的.
我身體要緊.
我現在家鎮真的很想死.
很大問題.
就回去吧.
換作其他人就已經走了.
回到菲律賓.
回到親人身邊.
應該照顧好一點.
現在在路途上.
這樣不是很好.
他不需要冒著生命的危險.
經文說他.
幾乎要死.
不顧性命.
病得很厲害,快要死了.
他堅持要去羅馬.
就算死也要去.
這筆錢.
不可以拖延.
一定要給保羅.
他的生活不可以受到虧損.
上帝的工作不可以受到影響.
病也要去.
死也要去.
是不是.
那種認真.
服侍.
以至於死那麼厲害.
所以你可見.
他那種認真.
多麼厲害,用盡全力.
別人去看.
你那些.
是甚麼工作,運錢.
供應.
你做那些很偉大的.
報道會.

$^{761}$一些很重要的聚會.
你是一些.
很大影響性的事.
你就說重要的事.
你就認真.
如果是不顧性命.
也差不多.
你這些.
用不著,好像有些小題大做.
他不是.
他真的不流汗.
這個態度.
是很值得.
很值得我們去欣賞.
我深信上帝就看在眼裡.
記在心中.
很欣賞這個弟兄.
其實這些位置.
是沒有人看到的.
我都說,如果不是保羅.
說出來,誰知道他不顧性命.
我們有沒有.
為這些.
小小的,隱藏的.
背後的事奉,傾盡全力.
沒有人看到的.
沒有人會讚你的.
會不會呢?.
為什麼他會傾盡全力.
就是剛才第二點說的,服事神.
我不分大小.
服事神就是這樣傾盡全力.
有沒有人看不重要.
我都會傾盡全力.
我不是要人的讚賞.
我不是要得人的榮耀.
我就是要.
為上帝做到最好.
今天我和你.
又如何呢?.

$^{801}$我記得很多年前.
我當時還是一個機構.
去服事,我不知道大家有沒有.
聽過遠東廣播這個名字.
是一個.
福音性的機構.
它就是在網上.
透過做節目.
向我們祖國的.
同胞們傳福音.
這個.
福音機構.
其實都是很多人的.
分幾個部門的.
有節目部,當時我是隸屬.
節目部,節目部是做什麼呢?.
通常是做節目,可能是主持人.
就是做監製.
寫稿等等.
接著就有.
工程部.
工程部就是稍為.
後方一點的,幫你剪接.
節目錄音錄影.
等等,做很多事情.
一些混音的事奉.
接著還有.
是回覆信的.
就是有些.
聽眾就會寫信來.
我想問一些信仰問題.
或者我生命有很多困擾.
我想知道是怎樣怎樣.
於是.
那個部門就負責回覆信件.
本來最好的做法.
是節目主持人回覆.
回覆所有信件.
但是因為節目主持人.
每天就很忙,做很多.

$^{841}$節目,經常都要趕.
Deadline,所以根本沒有.
時間去回覆聽眾的.
來信,於是就有一個部門.
是回覆信的.
我說了三個部門,是不是?.
每個部門都各有特色.
只有一個部門在前方.
就是節目部.
聽眾只是認識.
節目主持人.
最紅的都是那些人.
但是後面,無論是工程部.
回覆信的.
是沒有什麼人認識的.
這些都是很隱藏,很後方的.
這些同工們全部都很.
盡心盡力地服侍.
從來都沒有說.
因為我都不會紅了.
永遠都不會在那個流行榜排第幾.
收信第幾高.
我不會在那個榜.
但他們全部都很努力地服侍.
但是更加有一個人.
三個部都不入的.
是更加努力的.
那個人三個部都不入.
是什麼人呢?.
就是負責我們當時的機構清潔的.
那個阿嬸.
她的工作.
很簡單.
就是負責掃地.
清潔,還有.
她會幫同工們洗杯.
把地方弄得乾乾淨淨.
其實她是.
三個部都不屬於.
更加很少人.

$^{881}$過問她做什麼.
留意她做什麼.
但是這個同工.
她默默地做很多事情.
在她自己的.
領域裡面.
是很忠心.
她不稀罕.
她的工作,她的事奉.
她做到最好.
我經常以我自己觀察.
她是經常都很主動.
那些事情不一定是她的範疇.
她主動爭取去做.
去幫忙.
有些時候做完那些事情.
大家就覺得很開心.
誰幫忙做了.
這麼好呢.
大家回來,每天都覺得那個地方很乾淨.
很舒服.
拿個杯喝水的時候都覺得不同.
總之整個環境.
令到大家很享受.
在那裡事奉.
還有這個姐妹.
這個真的是一個姐妹.
她是會.
分享她自己生命的一些見證.
她也會默默地聽同事的分享.
有很多的鼓勵.
但是她從來都不張揚.
她那個位置.
是最沒有人留意的.
但是你有沒有想過.
原來她的事奉.
默默地一點一滴.
留在我們每個人的心裡.
最沒有人留意的.
好像是.

$^{921}$最不前方的.
卻是最值得欣賞.
我記得有一年.
我們機構會選.
最喜歡.
最出色的童工.
我記得是.
每一次都選有五位童工.
她從來都沒有想過.
那一年她拿了獎.
成為一個.
最出色的童工.
她拿這個獎的時候.
大家是很服的.
真的覺得她很配得.
你不要看.
她沒有人理會,沒有人注意.
好像做的事情.
你覺得節目.
比較暴力就最重要.
但其實她那個.
很值得欣賞.
她絕對是我們的弟兄.
姐妹,是我們的童工.
在事奉路上.
不可能沒有她.
沒有她.
我們那個事奉就不完整.
她也是我們的戰友.
沒有她.
你想像一下我們整個機構.
多骯髒,多不舒服.
回到家裡可能經常生病.
也不奇怪,就有她在.
令到整個事奉.
不一樣.
她拿了獎.
同樣的,以巴弗提.
我深信有一天在神面前.
她會拿一個很大的獎.

$^{961}$在經文裡面.
其實保羅是稱讚她的.
稱讚她.
怎樣稱讚她呢?我們去看看.
當這個.
以巴弗提病了.
病得很嚴重.
上帝因待她,醫治她.
那段經文說到.
27節,「願以神憐憫她」.
原來這個憐憫.
不單是那種.
情感上的.
那種憐憫,心意上的憐憫.
而是實際真的醫治了她.
是吧?.
以致她康復.
所以不單是憐憫她.
其實是憐憫保羅.
保羅是憂上加憂.
因為她有病而擔憂.
神醫治她,讓她康復.
接著保羅.
要送她回去.
因為這個以巴弗提.
病的消息傳回.
她家鄉裡的人.
家鄉的弟兄姊妹.
很擔心她.
很掛心她.
所以保羅就很想.
她康復,快些送她回去.
讓大家都可以歡喜.
不要再憂愁.
在過程中,保羅如何知道.
欣賞她?你們要在主裡.
歡喜地接待她.
29節.
而且要尊重這樣的人.
這個人很有價值.

$^{1001}$沒錯,以巴弗提要提早回來.
因為健康的緣故.
但你們要尊重她.
而且要歡喜地.
接待她,她很佩得.
她的生命很有價值.
佩得你們接待.
佩得你們尊重.
保羅欣賞她.
那你說她做了什麼?.
是不是?.
是不是?說了幾堂道.
傳了很多福音,叫很多人信耶穌.
嘩!真是厲害.
拯救多人靈魂,是不是?.
是不是這樣?不是.
她只是拿一些捐獻去.
她只是陪伴保羅.
服侍她,就是這樣.
但她的整個生命.
是佩得上帝的喜悅.
所以你們要在裡面歡喜地接待.
尊重這樣的人.
你看看這些人.
為了服侍上帝是怎樣?不顧性命.
幾乎致死.
為了補足你們.
供應我不足夠的地方.
整個菲律賓教會的弟兄姊妹.
根本就去不了羅馬.
陪不了保羅.
他們都沒有辦法去.
將這些奉獻帶過去.
他們沒有辦法.
做這個供應的角色.
是有不足夠的地方.
但是.
以巴弗提.
就代表了教會.
做了這件很重要的事.

$^{1041}$將兩者緊緊結連在一起.
所以以巴弗提的服侍.
是不是只是在.
她自己所做的工作?不是.
她讓保羅和.
菲律賓教會結連在一起.
低調的人.
都可以做到這個角色.
後方隱藏的角色.
那個事奉都可以這樣.
令不同的教會的部分.
不同的弟兄姊妹.
結連在一起.
在主裡合一.
這個是上帝體寵.
這個是真正有價值.
今天去到最後.
我去看回的時候.
未來的日子.
我深信弟兄姊妹有很多思考.
我申請參與哪個部門的服侍?.
做些什麼?.
我只是想.
邀請弟兄姊妹去禱告.
求神告訴你.
祂想你在哪個部門服侍.
我們思考的時候.
不是只是想.
這個位置可能多些人注意我.
那個位置那麼隱藏.
是不是?.
我們思考這些.
上帝想我擺在哪個位置.
到你事奉的時候.
我們的心態是保持低調.
前方,後方都好.
記住我一開始時說.
服侍是為了上帝.
不是為了誰.
服侍也不是為了榮耀自己.

$^{1081}$是榮耀上帝.
所以什麼位置前後都不要緊.
為什麼要低調?.
因為我們要榮耀上帝.
得上帝的喜悅.
好像伊巴忽提一樣.
能夠做到這樣就很好了.
我們一同禱告.
藉著這個很低調的人物.
提醒我們.
要怎樣去服侍你.
我們很想.
得榮耀.
我們很想得榮耀.
你會欣賞我們的服侍.
我們的服侍.
才會有價值.
求天父給我們.
弟兄姊妹將來的日子.
也是這樣.
要有一個討上帝喜悅的.
服侍的生命.
當我們好像伊巴忽提.
這樣服侍的時候.
我們也可以讓教會.
不同的部分.
彼此結連在一起.
在主裡面可以去合一.
主求你.
叫到我們的生命復興.
教會復興.
我們每個人有一個很美的.
服侍的生命.
整間教會一起有一個很美的.
合一的見證.
願你賜福給我們.
奉耶穌基督的名字祈禱.
阿們.
\newpage



\section{馬太福音列王記下 5:1-14}
\label{sec:eRsLDyLgpD0}
\textbf{2021.07.25《活著為主》 列王記下 5\_1-14;馬太福音22\_36-40 (和修版) 高淑芬傳道}
\newline
\newline
連結: \href{https://youtube.com/watch?v=eRsLDyLgpD0}{\texttt{https://youtube.com/watch?v=eRsLDyLgpD0}} ~~~~ 語音日期: 2021-07-27
\newline
\newline
\hyperref[sec:5qNMTaHZqr4]{\small{< < < PREV SERMON < < <}}
~
\hyperref[sec:index_chronic]{\small{[返順時目]}}
~
\hyperref[sec:index_scriptual]{\small{[返順卷目]}}
~
\hyperref[sec:G34IvTirgOo]{\small{> > > NEXT SERMON > > >}}
\newline
\newline
$^{1}$(日語).
不知道為什麼一提到泰國.
或者有機會講泰文.
就好像回到故鄉一樣.
在16年在泰國裡面的宣教.
因為我是一個人.
我是在曼谷北邊.
一個叫巴吞塔尼的地方服侍.
就算我們的警察會也沒有派其他的宣教士和我一起.
所以我的泰文會突飛猛進.
因為我沒有機會講廣東話.
當時很開心,短孫一來.
原來廣東話也是一個很喜歡講的語言.
短孫一來,當然他們也幫我做了很多爆炸的工作.
我這個女孩就在這裡撿回一些零碎.
在泰國裡面做了一些服侍神的工作.
聖經就是一個最大的陪伴.
在泰國裡面.
裡面有很多安慰,很多鼓勵.
很多令你繼續振作的話.
其中有一個小女子.
她到今時今日也不知道她叫什麼名字.
國內有些人就叫她小猴.
你從哪裡來的?你怎知道這個小女子?.
我查了很久也找不到她叫小猴.
中文就是小好.
原來一個女,一個子加起來就是個好.
所以說看見了,她就是小好.
當然這是一個笑話.
雖然聖經裡面沒有說她叫什麼名字.
她多了一些背景.
其實你剛才讀的也只有一兩句是形容她.
但她做了很偉大的事情.
而且我們也可以知道.
當時這個小女子的背景是一個很混亂的時代.
我相信我們香港人過去不是很理政局.
但這個小女子她生長的以色列國裡面.
就是一個政治非常動盪的時候.
經常都會打仗.
其實大家看新聞在中東國家你也會看到.

$^{41}$整天都在打仗.
在當時來說同樣都是這樣.
當時我們現在說的這個以色列國.
已經不是我們啟舊約開初的時候.
以色列人一定要立王的第一位皇帝.
就是素羅,大衛,所羅門.
不是那個時代了.
不是整個十二個支派一起的那個以色列大國.
而在小女子這個時代.
是在所羅門的兒子羅波安.
他那個很強制的政權底下.
結果這個以色列的大國就分了兩個國家.
分了兩個國家一個就叫做猶大國.
一個就叫做以色列國.
而這個小女子就是在這個以色列國裡面成長.
小女子她就在這個以色列國裡面.
是一個整天都會打仗.
是一個很混亂很戰亂的時代.
以色列國附近的國家.
整天都會有很多這些的政治衝突.
而且動不動就會打仗.
而小女子成長的時候.
以色列國的國王就是若蘭王.
是在位的.
若蘭王又是誰呢?.
他的父母都很出名的.
就是阿哈王和耶洗別.
他就是她的兒子.
雖然他不像她的父母那樣.
是拜巴力.
大家都記得早三個星期左右.
魯木斯有講到.
以利亞先知.
當時大鬥法的巴力.
就是阿哈王和耶洗別皇后的事跡.
若蘭王就是他們的兒子.
雖然若蘭王沒有像她的父母那樣.
拜巴力這麼厲害.
但是他也不是一個完全敬畏上帝.
一個為主而活的皇帝.

$^{81}$他不是這樣.
當時以利亞這位先知.
在以色列國裡面.
宣講上帝的話語.
但是之後被上帝接去.
他的徒弟以利沙.
剛才大家讀的經文裡面.
以利沙就是這位以利亞的徒弟.
他接住了他施父的先知的使命.
在以色列國裡面做一個先知.
他也是一個很著名的先知.
有很多神蹟奇事.
都是在他的身上.
神透過他來彰顯出來.
小女子很不幸.
當我們剛才讀的經文裡面.
我們會看到她已經在一個戰爭裡面.
被亞蘭閣.
亞蘭閣可以說是現在的敘利亞.
俘虜了這個小女子.
去服侍乃蠻.
乃蠻就是亞蘭閣的一個.
一人之下萬人之上的一個大元帥.
去她家裡面服侍.
她認識的她的家人.
她認識的親戚朋友.
不知道究竟是在戰亂裡面身亡了.
受傷.
還是和她一樣做了俘虜.
她可能都不會太清楚那個情況.
這個小女子.
面對在人生這麼大的逆境裡面.
我很被她吸引.
就是她仍然是尊主偉大.
在這幅圖裡面.
這個就是.
介紹一下.
這個是在巴吞塔尼教會裡面的崇拜.
很懷念.
大家看到後面剛才是很多樹木.

$^{121}$就是在一個這樣的地方裡面.
過去是在那裡服侍.
這個小女子.
她是尊主偉大.
她真的好像新的聖經裡面描述.
她是盡心盡性盡義地去愛主.
所以我很覺得.
她很多時候.
你知道事奉或者我們的人生都不是一帆風順.
雖然我們不至於像小女子.
在一個戰亂裡面.
被迫要和我們的家人.
我們所親愛的人分開.
但是問題是.
都是非常不容易的事情.
但是這個小女子.
她在五章三節裡面.
她仍然說.
她向她的主母.
即是奶孟的太太說.
我的主人就是奶孟.
如果她可以去撒瑪利亞那棵樹.
去找我們的先知.
她一定得到一個意志.
而且她的重點是.
她可能得到意志.
或者可以.
她是說到必定能夠.
如果你不是一個對上帝有信心的人.
你怎會說一定可以.
必定可以.
所以你會看到.
縱然她落在一個這麼逆境的裡面.
但是你會看到她對上帝的信心.
是非常令人難以理解.
為何她會這麼有信心呢?.
她現在不是很慘嗎?.
她正在做奴隸.
不能像以前那樣和家人一起生活.
但是她竟然在她的主母面前.

$^{161}$這麼勇敢來做這個見證.
你不要以為說一句話很輕鬆.
在她那個年代.
如果她說這句話不成立的時候.
她隨時可以死亡.
或者遭受很多的刑罰.
大家還記不記得勝經裡面.
以斯帖皇后.
她身為皇后要去見國王.
都是要拼了一條命的.
何況是一個完全沒有身份地位.
更加她不是雅蘭國的人.
她是一個俘虜.
所以大家看到.
她能夠說這句話.
你看到她縱然自己在一個逆境裡面.
但是她知道上帝仍然是上帝.
她所信的神仍然是全能的父上帝.
小女子為何可以這麼尊主偉大呢?.
大家要留意一下.
所以我們看聖經很有趣.
我很鼓勵大家不要怕了舊約.
其實舊約很好看.
很多很多的事蹟.
很多的故事.
以色列人的傳統家庭教育是很重要的.
做父母的都是一個家庭裡面的老師.
不是現在的家長只要他找中英數.
為了學校.
他是一個除了生活上的教導之外.
聖經的教導,律法的教導是不能少的.
所以當我們看到.
如果在聖經裡面說小女子.
即使她的年紀是很輕很輕的.
一個很年輕的女子.
竟然她可以這麼明白.
上帝是真真實實.
這麼有權柄的神.
你可以看到她的家庭教育.
是好好地教導她.

$^{201}$上帝的律法.
就好像十誡裡面.
除了神以外.
你不可以有別的神.
甚至愛我守我誡命的.
我一定要向他施以慈愛直到萬代.
這個在律法書裡面.
是不斷地重複的.
小女子為何這麼肯定.
上帝會透過撒瑪利亞的先知.
來治好乃曼的大麻瘋.
而且大家都知道.
大麻瘋是要去到1940年.
才有一些治療.
可以治療到大麻瘋這個疾病.
在過去裡面.
這個簡直是治療不了的.
治療不了的病.
竟然一個年輕女子.
為何她可以這麼大膽.
對神有這麼大的信心.
這個真是令我很佩服.
她為何找撒瑪利亞的先知.
這裡要交代一下背景.
因為當以色列大國.
分了兩個國家的時候.
當然猶大國的首都.
就在耶路撒冷.
以色列國.
他們誓死都不會去耶路撒冷敬拜.
那他們如何讓那十個支派.
這班以色列國分了的子民去敬拜呢.
他們就在撒瑪利亞山那裡.
買了這個地方.
撒瑪利亞這個位置就是他們的首都.
就是以色列國的首都.
所以去找撒瑪利亞的先知.
就是找以利沙.
因為在原文裡面.
這個形態是一個單數.

$^{241}$不是隨便找一個先知.
是找以利沙.
他又怎知道以利沙這麼厲害呢.
你去看到.
以利沙其實有很多神跡.
是在以色列境裡面出現過的.
如果你有興趣回去再去看前面.
甚至乎看以利亞的先知的奇妙事.
你都會看到.
昔日是透過先知做了很多很多的事情.
而以利沙他相信.
當時的事就沒有現在的傳媒那麼發達.
但是問題是一傳十 十傳百.
就是個個都知道.
就算是那麼發達的傳播網絡裡面.
我們都會看到.
以利沙的出名.
是他們以色列國的人認識他的.
知道這個先知是大行其事的先知.
雖然這個年輕女孩.
她已經被俘虜了.
現在是在亞蘭閣裡面.
但是她的眼睛不是在她的困難裡面.
不是在她的悲慘當中.
不是在她的受苦裡面.
她的眼光是定睛在上帝的面.
我們今天有沒有遇到很多很多的事情.
會不會定睛的仰望我們的神呢?.
去尊主為大呢?.
她超越了她自己的傷痛.
淪落成為一個別鄉的一個工人.
她很難得的就是.
她不只是在她的困難裡面.
她是盡心盡性愛著.
她很珍惜現在眼前.
現在面前她要面對的人.
和要做的事.
她面對逆境.
她是會傷感.
但是她沒有停留在傷感裡面.

$^{281}$不斷抱怨.
今天我們遇到很多事情.
例如現在戴口罩這件事.
差不多快要接近兩年了.
我們是不是只是定睛在這個逆境裡面.
或者我們香港的氣候都變了.
是不是就是這樣抱怨呢?.
我們懇求主幫助我們每一個.
讓我們知道這個仍然是天賦的世界.
我們的眼睛是要望準這位親愛的上帝.
好像這個小女子一樣.
專注為大.
其實一個女孩子.
她是不情願的.
因為她被俘虜.
去到一個地方其實是很不容易的.
可以幫我繼續下去第十頁.
她是被俘虜的.
而我去泰國的時候.
我不是被俘虜去的.
我真的是上帝的差遣.
其實我還沒有去的時候.
我在宣道會荃灣堂做了七年的傳道人.
做傳道人不是我自己說的.
自己去做傳道是上帝呼召.
同樣其實已經七年了.
我在宣道會荃灣堂也是我的母會.
已經是很熟悉.
但是上帝竟然要我出去.
當時也知道要聽神的話.
所以2001年就是我的差遣禮.
去一個泰國.
大家現在看我又說泰文.
就以為我泰國當然有很多親戚.
對不起.
在泰國沒有親戚.
2001年我只會說「你好」.
你們也會吧.
所以可以去下一幅了.
起初在巴吞他尼的日子裡.

$^{321}$這一書也不是我的會友.
只有一個左右.
全部都是一些短孫.
當時我剛剛去的時候.
這個地方是有女傳道和女幹事的.
我那時候在曼谷學泰文.
學完泰文之後.
想不到一個女幹事去讀神學.
現在也成為女傳道.
但是女傳道就辭職.
又是剩下自己一個人.
以為有同工.
剩下一個人.
下一幅.
但是很感恩.
蒙著神的恩典.
剛才你看到的地方.
在頭三年已經因為白蟻侵蝕.
沒有再租那個地方.
繼續做實的工作.
所以第一年的述職.
我很慘的.
你知道嗎.
我被我妹妹笑.
人家述職就說怎樣去見教會.
你述職呢.
你那間教會就.
因為我要回來香港前.
我還要把教會的東西.
快點搬去我家裡.
因為還沒有找到地方.
你述職就是.
沒有收檔.
在教會收檔.
接了.
有得去說笑.
但是真的是蒙著神的恩典.
剛才你看到那個地方.
我們由一個單位變成兩個單位.
都是租的.

$^{361}$結果人數是越來越多.
後來甚至神的恩典.
給我們有自己的地方.
那輛車就是我在這麼多年.
裡面開的車.
就是這輛車.
最厲害的其實是七人車.
現在多一些.
在香港我看到這些車.
我就很懷念.
以前回來香港我沒有見過.
但現在越來越多見到.
不過後面沒有弄一個.
我們泰國那裡叫Carry Ball.
後面都可以裝到有冷氣.
人都可以塞進去坐.
所以當有試過.
我的車又要去催一些比較遠的弟兄姊妹.
或者有時候我們要出去一些事奉.
都是要用這輛車.
特別最高峰期.
就是有短孫隊來.
記得有一次短孫隊聖誕節來的時候.
我們好像玩東華三院籌款.
這輛車其實不是籌款.
是塞車.
大家塞車.
你猜塞了多少人?.
20人.
多謝.
18人.
不過20人都對.
為什麼呢?.
因為聖誕節你去做的節目.
全部我們都要拿著很多禮物.
單張.
所以你真的猜不到這輛車.
塞了18人這麼多.
後來這個地方.
你看到我們泰國很特別的文化.

$^{401}$就是崇拜完.
下面那一幅是崇拜.
崇拜完之後.
我們就會一起有外賢的時間.
好,接下來.
大家留意.
我不是被俘虜.
但是我要第一年在泰國那裡.
你是什麼都做不到.
雖然你已經是有七年的年資的一個傳道人.
但是你不會泰文.
你做什麼呢?.
就在學習.
學到你都會覺得.
你很低B.
我們泰國傳播就明白了.
因為泰文有一個字.
真的好像B一樣.
從小到大我們寫B是怎樣的呢?.
你當你前面有個黑板.
你寫個B.
這樣.
然後這樣.
是不是?.
你知道老師教你.
還要教你寫字.
你回去寫.
我們要寫到這些.
這個也不是一個字母.
這個是音節旁邊的一小部分.
每逢看到這個.
我當然是這樣去.
樣子是一樣的.
和老師比的那些.
但是全部畫了交叉.
為什麼呢?.
原來泰文字的秘訣.
就是一定要畫圓圈.
如果你畫了圓圈.
才寫這個B.

$^{441}$你會怎樣寫呢?.
你寫給我看吧.
你們真的很有這個.
這個潛質.
你們應該去泰國.
不是我去.
我就很低B地去了.
我低B.
就是.
老師.
我真的不明白.
因為一模一樣.
和你這麼大個人.
寫字都寫錯.
你會覺得.
你一個這麼大的人.
都會覺得你很低能.
我問老師為什麼錯呢?.
他就跟我說.
你一定要畫圓圈.
才上去.
所以他會看得到.
你的筆順不對.
所以錯了.
但是你也要適應當時的文化.
泰國的文化語言是其中之一.
有很多東西要適應.
過去七年的傳道裡面.
因為我有父母照顧.
家人.
真的可以專心傳道.
但是來到泰國.
你不可以.
像過去那樣.
你的生活有人照顧.
我那時候回來香港.
三年之後.
就回來一次述職.
我第二次述職.
你都不知道我這麼大個人.

$^{481}$竟然說營養不良.
問我你在泰國吃什麼.
我一直說都說不到飯.
因為我覺得煮飯又要煮菜.
是很麻煩的.
因為很多事情要做.
我最多是吃蝦子麵.
蝦子麵和一些菜.
我覺得自己很有營養.
因為泰國最多是很多水果.
不同的月份都有不同的水果.
我經常都會吃水果.
還有.
我學了泰國人.
很多葉.
菜.
豆角.
那些都可以生吃.
真的很方便.
因為我去到哪裡都要開車.
我洗乾淨放在旁邊.
我就拿一條來吃.
拿一個紅蘿蔔.
這樣吃.
我真的料不到.
這麼有維他命C的.
裡面.
還有維他命D.
那些陽光.
竟然說我營養不良.
你會看到.
當我們去事奉神的時候.
其實有神的誡命其實很好.
你會看到.
第一.
當然是盡心盡性去愛我們的神.
但其次.
很重要的就是.
愛人如己.
我們很容易一跳就跳了.

$^{521}$愛別人.
其實你要留心看到.
你要愛人好像愛自己.
你自己都不愛得到.
你愛別人就不夠完全.
所以.
很感受到.
真的要有平衡的生活.
我當時就營養不良.
後來我的身體一直變差.
所以我就不能去泰國.
不過也很感恩.
竟然在這樣的條件下.
蒙神的恩典.
足足在泰國16年的時間.
但是在再大的病痛裡面.
更加感受到.
過去真的很忽略.
上帝創造這個世界給我們.
我一面倒.
很事工型地去做很多不同的事情.
但是上帝告訴我.
其實你都要愛.
愛你自己.
生活上都要有平衡.
然後.
第15幅圖.
我們要愛自己.
愛自己是做好自己.
除了在生活上.
我們要有自己的規律.
好像傳道師說.
「睡有時,吃有時,玩樂有時」.
原來這三個部份.
是我們生活上要有的東西.
人在不容易的情況下.
我們當然.
剛才在生活上要做好.
在這裡.
為什麼會有.

$^{561}$《Wonder》.
《Wonder》中文是.
《奇蹟男孩》的片.
這個小朋友.
他一生出來的面.
是和其他不同的.
樣子是有些扭曲了.
是真人真事改編的.
但是.
他一路從小到大.
父母在家庭裡教他讀書.
但他到初中的時候.
父母都覺得他應該要.
回到正常的學校裡.
來受教.
當然.
你看看他樣子.
他單是很害怕.
父母都幫他做了很多心理準備.
但他的確回到學校.
真的被人家.
.
我現在.
因為上午教完滿份.
我還有些泰文在裡面.
看低他.
看不順眼他.
所以.
他被人欺凌.
但他是靠著.
父母的鼓勵.
還有自己.
做好自己.
自己不要沉在一個.
只是一個被欺凌.
自己很慘的裡面.
他是跳出來.
他是慢慢慢慢.
因為他的學業很好.
旁邊的同學都是.

$^{601}$他又待人又好.
於是他慢慢.
由這個欺凌的裡面.
被人欺負的裡面.
反而可以結交很多朋友.
旁邊那幅.
就是一個黎智偉.
我想我很想表達我們愛自己.
縱然我們去到一個不好.
很患難的裡面.
或者家家有本難唸的經.
都可能有著不容易的處境.
這個黎智偉.
他本來是一個攀山的名人.
但是很可惜.
有一次交通意外.
他的腳撞了.
不可以再用雙腳.
來爬山.
但當他康復完之後.
他竟然.
因為他的兒子的緣故.
他會覺得我要再爬.
他坐著輪椅去爬.
早幾個月.
他好像在魚深.
在荃灣那座大樓裡面.
很想做一個紀錄.
爬到底達大樓裡面.
但後來失敗了.
下次再來過.
但你會看到.
其實我們身邊有些人.
他在不容易的裡面.
但我們不要沉在不容易的裡面.
你要愛自己.
因為你.
特別我們坐在這裡.
我們大部分.
我們都是相信主的人.

$^{641}$我們都是用上帝的寵駕.
賣熟回來.
你是很有價值.
你是很寶貴.
你不可以就這樣.
中了撒旦的圈套.
只沉在一些逆境裡面.
你會看到.
其實我們身邊.
這只不過是舉例.
身邊有很多很多的人.
你會看到他們都是.
要做好自己.
例如這套電影.
這套電影.
讓你更加了解到.
小女子的情況.
因為大長今.
大家在我的年代.
應該看大長今.
或者再看阿信.
即是一些很刻苦.
但又怎樣.
然後可以好好地.
來到很有成功的人生.
這個大長今.
大家知道.
後來她很小的時候.
因為有一次.
有一個機緣際遇.
她的二父二母.
賣走給皇宮.
結果就讓她.
小時候就進了皇宮裡面.
做廚房.
學做廚房.
你會看到.
去到這些大戶人家.
小女子就是大戶人家.
大長今就是皇宮裡面.

$^{681}$是不容易讓你.
一去就去到這麼高.
我這裡很想帶出給大家看.
經文一兩句就說到.
她被擄之後.
從以色列第二節那裡說到.
以色列以前的亞蘭人.
一群人撞了出去.
在以色列擄了一個女子.
這個女子就是服侍賴孟的妻.
很簡短地交代了這件事.
但我要告訴大家.
短短的篇幅裡面.
經歷過很多事.
就好像大長今.
我想要洗碗.
做一些很下流的事.
才慢慢看到她的工作表現.
她的誠信.
才會慢慢去到今天.
她去服侍主人的太太.
你不會讓一個俘虜去做的.
說一些國家機密.
怎知道她會不會是間諜.
你這樣看.
我們看得太輕鬆.
但其實裡面已經經過了很多不容易.
這個小女子給我覺得很深刻的印象.
就是她已經是尊主偉大.
但同時她真的.
好像《新約聖經》裡面的說話.
她真的是一個愛自己的人.
無論在什麼處境下.
她仍然是愛.
在這裡帶進自己.
是有一個好的前面.
人的生命就是這樣.
一路一路去學.
學一步一步去.
我相信她一定是有一些.

$^{721}$有些特殊的氣質.
出眾的表現.
她才會成為了.
奶萬元帥的太太的傭人.
去服侍她.
而且她所說的話.
是被重視,聽.
聽了之後.
你知道她是跟奶萬的太太說.
太太跟奶萬說.
奶萬跟著要跟皇帝說.
你會看到這件事.
是來自一個敵國的俘虜說的.
源頭在那裡.
你那個源頭大一點也好.
就是這麼小的一件事.
所以你會看到.
這個小女子在不容易的裡面.
年紀輕輕的她.
她也懂得靠著神去做好自己.
人的生命真的很脆弱.
可以活著一天.
我們就真的要好好珍惜眼前的一切.
我今天也跟門訓的姐妹說.
讓我站在台門.
不過站在台門突然間腳軟.
因為這段身體.
當我回來後有些大病的時候.
我覺得我們的身體,我們的精神是判斷不到的.
我發現這段時間.
我也會突然間腳軟.
突然間為甚麼腳會這樣不能動.
第一.
我來紅欣塘.
其實最難倒我的就是那些樓梯.
我第一天.
其實我不認識這裡.
我沒有怎麼來過紅水.
第一天來的時候.
為甚麼牧師不坐在那裡跟我談.

$^{761}$還要走上去.
一路一路走到三樓.
我就在那裡說.
上帝啊.
如果這裡是三層樓.
你要讓我好好思想.
上到三樓就不要下到二樓.
二樓就不要去三樓.
不要上來上去.
不要上上下下.
因為當時我真的在看我的膝蓋.
醫生都說我不能醫治.
但是很痛的.
其實四月中我們才實體崇拜.
我一月份就已經和我的泰國姐妹.
團契的姐妹一起去探訪的時候.
他們都知道我很不舒服.
其實我身體也不是很痛.
但是.
但是我去到五月很感恩.
我每個月都在祈禱.
上上下下沒有問題.
五月份的時候.
上下下好多了.
但是身體一會兒好一會兒不好.
都是不能夠自己控制的.
雖然這個是你的身體.
你是祂的主人.
但是都控制不了.
所以我們更加要求神幫助我們.
做好我們自己.
生活上都要平衡.
你又不可以一面倒.
只要事奉你的家庭.
你的父輩.
如果你有子女.
你的子女.
很多事情都要平衡去做.
而且無論你在哪一個處境裡面.
你記住.

$^{801}$你是一個重價買回來的.
你一定是有著很寶貴的價值在當中.
所以.
做好自己.
無論落在一些不好的情況下.
好,接下來.
給大家看一下.
接下來那個應該是十六.
十六的圖.
第十六.
我其實在泰國裡面.
其實做事很不容易.
因為這段時間.
我現在是泰國皇帝.
但他頒給我的時候.
他只是一個泰國皇子.
我就拿到一個金牌.
其實我覺得神是透過這件事去鼓勵我.
因為那段時間.
我是非常低落.
為什麼呢?.
我拿金牌的緣故.
因為我是去癌症病人那裡.
去探訪那些病人.
是持之以恆去做的.
和我們一起去的.
是整個巴吞帕尼的不同的教會.
都有人一起去受了醫院的訓練.
三天這麼多.
教你怎樣見到病人.
要怎樣.
基本常識.
開頭的時候很開心.
因為每逢星期四去.
A教會,B教會,C教會,D教會的人一起.
是很開心的.
甚至我自己教會.
我都會叫大園姐妹一起去.
他們去是很好的.
我們現在去香港.

$^{841}$最多去屯門醫院等.
你去.
我去開車到你家.
除非他有車.
但大部分都沒有車.
因為我們泰國姐妹知道.
巴吞帕尼在泰國裡面.
是一個.
我現在回來才知道.
幸好在那裡不知道.
沒那麼害怕.
是一個黑名單.
罪惡的.
因為它接近曼谷.
所以什麼混毒.
非法的東西.
在那裡是很多的.
怪不知我有一次試過開車.
我好好的開車.
那個人可能他想爬我.
結果爬不到.
幸好那時候有個姐妹.
短孫在那裡.
「高姑娘好像拿著一支槍」.
我說是嗎?.
你知道我應該向前望.
真的不知道.
因為泰國是很容易買槍.
他竟然對著槍.
好像是上帝.
讓我快點走.
不要理他.
很多事情.
我想日後有機會可以再慢慢說.
這樣就看到.
很不容易.
總之這個時段.
譬如你去.
我就會開車去到你家.
一起去到醫院.

$^{881}$受訓完之後.
就送你們回家.
結果之後的探訪都是.
我們那時候就是逢星期四一次.
去探這些癌症病人.
都是這樣.
但是去著去著.
四個變成三個.
三個變成一個.
就剩下自己了.
每個人都說沒空了.
不去了.
如果他們一起和我去.
堅持到底.
就不只是我一個人去拿金牌.
我告訴你.
其他的教會都是這樣.
沒有人.
因為我們後期去.
只剩下我們教會去.
都談得很好.
和大英姐妹堅持到底.
堅持到一個.
但是真的沒有人告訴你.
你這樣做有個金牌.
我自己都不知道.
只不過我覺得.
去探癌症病人.
我覺得神開了這個門.
因為在這間醫院.
是有和尚去做院務的.
和尚去探.
但是他難得給你基督教.
人可以去.
所以我覺得.
神開了這個門.
不可以關了.
我就去.
所以剩下一個人也好.
我就去.

$^{921}$去到有一個家.
最後拿到獎.
我之外還有一個家庭.
那個家庭的太太.
是曾經癌症.
所以她有心同感受.
就叫了她的親戚.
她的妹夫.
她的丈夫.
就是妹妹.
四個.
他們有四個人都拿到金牌.
其餘的都看不到.
所以當我拿了這個金牌的時候.
神就說你繼續堅持下去.
結果我就可以堅持下去.
有幾秒鐘就上了電視.
那一晚.
這個金牌.
但是弟兄姊妹知道.
他們知道我怎樣拿的時候.
其實上帝是親自給他們一個教訓.
其實你繼續去.
他們是和我一樣可以拿到.
應該去十七頁那裡.
除了我們要愛自己的時候.
愛人如己.
我們要去愛別人.
這個小女子.
她沒有懷著心仇大恨.
她仍然很懂得去愛其他人.
特別是愛敵人.
你看到她.
她在聖經裡面描述.
她和主母巴不得.
在和合本裡面.
巴不得主人去見撒瑪利亞的先知.
巴不得在和修本裡面說希望.
意思是.
真是很想主人去見撒瑪利亞的先知.

$^{961}$是真心很希望.
你看到嗎?.
她不是隨便就這樣.
不像若拿.
若拿是一個先知.
上帝派她來微城宣講.
隨便說說.
是上帝的角度.
這樣就走人了.
她不是.
她是很誠懇.
很想你去.
很想很想你去.
還有.
必定自好.
自好在那個自裡面.
就是收集了這個東西並且帶走.
或者是收集了你那些麻瘋.
那些病菌.
有的那些病帶走.
是一個對神有信心之餘.
她對一個敵人.
真是秉承了.
新約裡面我們說要愛你的敵人.
什麼要輸他七個七次.
她是完全做到.
所以我就很感受到.
她這個愛別人.
是非常之誠懇.
她其實可以面對著.
這件事情可以閉口不說.
但是她卻是這麼勇敢.
這麼愛.
這麼肉緊.
希望她可以好起來.
可以繼續下去第十八頁.
我在泰國的第十八頁.
開頭我就是開荒牧會.
但是也是一些機緣的巧合下.
上帝竟然給我有機會認識了.

$^{1001}$在巴吞塔尼的穆斯林的長老.
我就有機會去他們的學校.
去他們的會堂.
來跟他們建立關係.
其實之前在宣道會你會看到我的樣子.
有照片登.
但我後期的三年因為做穆斯林的工作.
我就沒有了.
沒有我的樣子.
你只看到高淑芬.
因為泰國人不認識我的名字叫高淑芬.
我竟然可以認識她.
在很多學校裡.
知道巴吞塔尼除了是一個罪惡城市之外.
巴吞塔尼是一個屬靈真正很大的地方.
為什麼呢.
原來穆斯林的總部.
我如果不是認識他們我也不知道.
以為他們的總部一定在南部.
因為泰國的南面是最多穆斯林的人.
你經常聽到他們扔炸彈去軍營 去警察局.
甚至乎後來炸一些學校.
因為他們很想獨立.
他們那幾個 泰國有77個府.
在南部的那些府.
特別是穆斯林勢力大的那些.
他們希望獨立成為一個國.
竟然知道原來總部在巴吞塔尼那裡.
但是叫做「接讓著曼谷」.
由我們教會去到30分鐘路程.
是很近的地方.
所以除了穆斯林.
還有一個佛教的一個大的組織.
是國際性的.
都是在巴吞塔尼裡面的.
那個叫「法新寺」.
所以看到是很不容易.
但是神竟然再下一幅.
神竟然給我有機會去認識這班人.
跟著的十九.

$^{1041}$還有因為我是他之前的一幅.
這個說一說反正有了.
這個是我那次回來香港的時候.
這個長老帶了他的.
那時候還是女的.
五個女的.
來一起穿著盛裝去送我上飛機.
這裡你看到.
你真的看不到這些圖畫.
我開頭以為沒有什麼.
但是後來我參加一些.
穆斯林的一些課程.
那個老師是馬來西亞人.
經常都跟我說.
「Man to man, woman to woman」.
因為我經常告訴他.
我跟男人對著長老.
男的男 女的女.
不可以一起的.
因為馬來西亞是這樣.
還有一起受訓練的同學們.
他們是中國的回民.
或者是中東那些國家更加.
你說你做宣教士.
你沒有一個男人帶著你.
你都不用在那裡做.
神就知道我沒有男人.
竟然在泰國裡面.
做一些這麼自由的穆斯林.
起碼他們很尊重我.
開會跟他們.
全部男人.
女人有的在外面做飯.
我是唯一一個不用做飯.
跟他們在那裡開會的人.
其實是做得很好.
跟他們的關係很好.
所以希望日後有機會.
我回去繼續探望他們的時候.
我希望可以帶你們一起去.

$^{1081}$有沒有人願意去?.
當然了.
如果去到要坐牢.
坐酒店牢.
我們當然不去.
因為時間去到坐牢是不行的.
真正開放的時候.
我很希望.
窮恩堂可以跟我一起去.
神在宣教裡面讓我知道.
我要愛神之餘.
我要愛自己.
更加愛其他不同種族的人.
因為今天的時間.
其實是最後一課了.
想不到他們剛才說的長老家庭.
他一家只有一個太太.
跟著那五個女兒.
他可能想追逐.
結果最後第六個女兒.
已經是兒子了.
不過已經是我回到香港的時候了.
他們很好.
我從來沒有試過.
被人這樣在泰國送機.
在香港他們都穿著城服.
他們在機場那段時間.
每個人都跟你拍照.
覺得他們穿得很漂亮.
又很可愛.
還每一個人.
你看我手上的玫瑰.
是每一個人從小到大排隊給你.
好像學校.
想不到.
我隨便告訴你.
我那天要走.
上午我就走了.
他沒有告訴我他真的要去.
你知道很多時候.

$^{1121}$對不起 泰國大英姐妹.
我在泰國裡面.
經常被人放飛機.
他說去.
其實不去的.
他說三點來.
可能六點才來.
或者夜一夜才來.
以為他不來.
三個小時等到人資如燼.
走了他就來了.
在泰國裡面.
很多這些.
要經常收拾.
所以我以為他沒有來.
結果他還早來.
我還在巴丹塔來.
他已經說.
我已經在機場了.
你在哪裡.
我說我還在家裡.
為什麼你不載我.
我在心裡.
為什麼你不載我走.
當然他不知道我在哪裡.
我只試探他.
沒有讓他來我家.
但問題在我心裡.
真是可惜.
但想不到.
他竟然將我們的相片.
用了他的.
因為我當時.
還可以和他聯絡.
為什麼你拿了這個來上片.
所以很感恩.
上帝給我有機會.
在過去十六年在泰國服侍.
還有列王記下.
我們還未說完.

$^{1161}$現在時間已經是12點40分.
所以大家記住.
剛才我說的歷史背景.
給你們一點功課.
8月15日.
我一樣會說列王記下.
你今天看到我.
忍住不說其他人.
只說小女子.
因為小女子對我來說.
是很大的鼓舞.
很大的借鏡.
因為這個不是故事這麼簡單.
是真人真事.
我們在下一次.
會透過這段經文.
再去看宣教的現象.
今天希望神祝福我們.
祝福我們每一個.
真是尊主偉大.
我們紅印堂的弟兄姊妹.
每一個是尊主偉大.
無論我們香港社會.
變成怎樣.
我們都要愛我們自己.
做好自己.
無論你現在正處於一些.
很難念的經驗.
很多事情不容易.
但是你可以做好自己.
另外求神也給我們力量.
不單單是愛自己.
我們愛身邊其他人.
特別一些過去傷害過你的人.
或者一些不容易的人.
因為這些功課.
人與人做事.
是會碰到的.
教會也是一個小社會.
也是很不容易.

$^{1201}$因為我們都是一個聚人.
進來這裡.
聚在一起.
所以大家很多不同.
很多的磨合.
我們都要懂得去包容其他人.
上帝愛世界上每一個人.
特別是紅印堂.
神很賜福給大家.
你們這裡有泰語團契.
不然我就不來了.
沒有機會來了.
來認識大家.
我們很奇妙.
我真的沒想到.
我來做泰國團契.
竟然買一送一.
我的團契要有英文.
所以現在我做分.
我們有香港的弟兄姊妹.
所以每一次課程.
我都要有三個語文.
泰文,中文和英文.
但我覺得買一送一.
也是鍛鍊我的時間.
因為上帝真的愛世界上每一個人.
每一個族群.
特別現在奧運會.
大家看到很多名字.
是沒有聽過的.
會不會一邊看的時候.
這個什麼國家.
不只是奇怪.
而是神給他們這個國家.
也有人相信你.
是萬失敗敗的一個國家.
我們很感恩.
這裡有兩位非洲裔弟兄.
在我們當中.
我們做他的朋友.

$^{1241}$你們可以跟他打招呼.
跟我們的泰文姊妹.
說廣東話.
我們願意.
我們紅欣塘是一個充滿愛的國家.
讓人一看到我們.
就知道我們是上帝的兒女.
願神祝福大家.
\newpage



\section{創世記 12:1-20}
\label{sec:G34IvTirgOo}
\textbf{2021.08.01《絕不離地的信仰》 創世記 12\_1-20 (和修版) 老耀雄牧師}
\newline
\newline
連結: \href{https://youtube.com/watch?v=G34IvTirgOo}{\texttt{https://youtube.com/watch?v=G34IvTirgOo}} ~~~~ 語音日期: 2021-08-04
\newline
\newline
\hyperref[sec:eRsLDyLgpD0]{\small{< < < PREV SERMON < < <}}
~
\hyperref[sec:index_chronic]{\small{[返順時目]}}
~
\hyperref[sec:index_scriptual]{\small{[返順卷目]}}
~
\hyperref[sec:_2_W0T2LR08]{\small{> > > NEXT SERMON > > >}}
\newline
\newline
創世記 12:1-20
\newline
\begin{longtable}{cl}
\hline
\hline
章節 & 經文 (和合本修訂版)\\
\hline
12:1 & \begin{tabularx}{0.7\textwidth}{X} 耶和華對亞伯蘭說：「你要離開本地、本族、父家，往我所要指示你的地去。 \end{tabularx} \\ \\ \relax
12:2 & \begin{tabularx}{0.7\textwidth}{X} 我必使你成為大國，我必賜福給你，使你的名為大；你要使別人得福。 \end{tabularx} \\ \\ \relax
12:3 & \begin{tabularx}{0.7\textwidth}{X} 為你祝福的，我必賜福給他；詛咒你的，我必詛咒他。地上的萬族都必因你得福。」 \end{tabularx} \\ \\ \relax
12:4 & \begin{tabularx}{0.7\textwidth}{X} 亞伯蘭就遵照耶和華的吩咐去了；羅得也和他同去。亞伯蘭離開哈蘭的時候年七十五歲。 \end{tabularx} \\ \\ \relax
12:5 & \begin{tabularx}{0.7\textwidth}{X} 亞伯蘭帶著他妻子撒萊和姪兒羅得，以及他們在哈蘭積蓄的財物、獲得的人口，往迦南地去。他們就來到了迦南地。 \end{tabularx} \\ \\ \relax
12:6 & \begin{tabularx}{0.7\textwidth}{X} 亞伯蘭經過那地，直到示劍地方，摩利橡樹那裡；當時迦南人住在那地。 \end{tabularx} \\ \\ \relax
12:7 & \begin{tabularx}{0.7\textwidth}{X} 耶和華向亞伯蘭顯現，說：「我要把這地賜給你的後裔。」亞伯蘭就在那裡為向他顯現的耶和華築了一座壇。 \end{tabularx} \\ \\ \relax
12:8 & \begin{tabularx}{0.7\textwidth}{X} 從那裡他又遷到伯特利東邊的山，支搭帳棚；西邊是伯特利，東邊是艾。他在那裡又為耶和華築了一座壇，求告耶和華的名。 \end{tabularx} \\ \\ \relax
12:9 & \begin{tabularx}{0.7\textwidth}{X} 後來亞伯蘭漸漸遷往尼革夫去。 \end{tabularx} \\ \\ \relax
12:10 & \begin{tabularx}{0.7\textwidth}{X} 那地遭遇饑荒。亞伯蘭因那地的饑荒嚴重，就下到埃及，要在那裡寄居。 \end{tabularx} \\ \\ \relax
12:11 & \begin{tabularx}{0.7\textwidth}{X} 將近埃及，他對妻子撒萊說：「看哪，我知道你是美貌的女人。 \end{tabularx} \\ \\ \relax
12:12 & \begin{tabularx}{0.7\textwidth}{X} 埃及人看見你會說：『這是他的妻子』，他們就會殺我，卻讓你活著。 \end{tabularx} \\ \\ \relax
12:13 & \begin{tabularx}{0.7\textwidth}{X} 所以，請你說你是我的妹妹，使我可以因你得平安，我的性命也因你存活。」 \end{tabularx} \\ \\ \relax
12:14 & \begin{tabularx}{0.7\textwidth}{X} 亞伯蘭到達埃及時，埃及人看見那女人極其美貌。 \end{tabularx} \\ \\ \relax
12:15 & \begin{tabularx}{0.7\textwidth}{X} 法老的臣僕看見了她，就在法老面前稱讚她。那女人就被帶進法老的宮中。 \end{tabularx} \\ \\ \relax
12:16 & \begin{tabularx}{0.7\textwidth}{X} 法老就因她厚待亞伯蘭，給了亞伯蘭許多牛、羊、公驢、奴僕、婢女、母驢、駱駝。 \end{tabularx} \\ \\ \relax
12:17 & \begin{tabularx}{0.7\textwidth}{X} 耶和華因亞伯蘭妻子撒萊的緣故，降大災擊打法老和他的全家。 \end{tabularx} \\ \\ \relax
12:18 & \begin{tabularx}{0.7\textwidth}{X} 法老召了亞伯蘭來，說：「你向我做的是甚麼事呢？為甚麼沒有告訴我她是你的妻子？ \end{tabularx} \\ \\ \relax
12:19 & \begin{tabularx}{0.7\textwidth}{X} 為甚麼說『她是我的妹妹』，以致我把她接來要作我的妻子呢？現在 ，看哪，你的妻子在這裡，帶她走吧！」 \end{tabularx} \\ \\ \relax
12:20 & \begin{tabularx}{0.7\textwidth}{X} 於是法老吩咐人把亞伯蘭和他妻子，以及他一切所有的都送走了。 \end{tabularx} \\ \\
[1ex]
\hline
\hline
\end{longtable}
$^{1}$各位弟兄姊妹平安.
在《聖經創世記》十二至五十章裡.
記載了以色列民族幾個族長的故事.
包括了亞伯拉罕,以撒,雅各,約瑟.
族長時代有甚麼特點呢?.
其實它是一個父系的社會.
父親是一家之主.
負責照顧族人各方面的需要.
而家人亦都需要父親的照顧.
所以如果一個家庭沒有了父親的話.
整個家族就陷入了經濟上和生活上的困難.
簡單來說.
在當時沒有一個社會或政府的架構裡面.
父親就是一個法官,一個祭司.
甚至是一個主人.
家人成為他們的財物.
他可以全權支配.
權力是毫無限制的.
不過雖然族長擁有無上的權力.
但都有無比沉重的負擔和壓力.
因為他要負責整個族群的生計和發展.
有人說過.
如果他想知道一個人真正的品格.
可以從他面對逆境的態度和行為.
就可以知道得一清二楚.
今天我們透過《創世記》的十二章一至二十節.
在《喀巴拉罕》初期的兩個生命片段裡面.
去學習一下信仰實踐的兩個提醒.
我們一起有個祈禱.
主你透過聖經向我們說話.
讓我們透過聖經認識你.
求你讓我們將你的話語藏在心裡面.
成為我們一生的幫助.
叫我們一生都跟隨你.
我們祈禱奉靠主耶穌基督.
得勝神明治祈求.
阿們.
開始之前我們先了解一下經文的背景.
在《創世記》第一章到第十一章.
除了記載了世界和人類的起源之外.

$^{41}$還描述了人類屢次犯罪而引致被神懲罰.
首先.
亞當和夏娃不聽神的吩咐.
被趕出伊甸園.
大地因此被咒了.
人類雖然繼續繁險眾多.
但罪惡一樣多.
人類的邪惡.
最終引致神用洪水來毀滅世界.
經文似乎重覆一個主題.
就是咒詛受罰.
但《族長》記載一開始.
就提到上帝的祝福.
兩者成為一個很強烈的對比.
《創世記》的作者.
很希望讀者通過這個對比.
能夠看見上帝的心意和本性.
上帝的心意是要賜福人類.
而不是要懲罰人類.
另外亞伯蘭在第十二章出現時.
他還未改名.
還是叫亞伯蘭.
要到第十七章.
他九十九歲時.
神才為他改名.
叫做亞伯拉罕.
我們知道這些背景後.
我們開始進入經文.
我們信仰實踐的第一個提醒就是.
我們能夠以信服,專注為大的心.
去回應神對我們的呼召.
亞伯蘭的夢召過程.
讓我們感受到神人關係的親密.
神與亞伯蘭說.
你要離開本地.
本族,父家.
往往要指示你去的地方.
本地是甚麼意思?.
指他的國家.
本族指他的故鄉.

$^{81}$即是他與親戚居住的地方.
父家是指他所住的屋.
他的祖家.
就是他與直系親人.
包括他的父母,兄弟姐妹同住的地方.
耶和華要求亞伯蘭.
離開一個他最熟悉,最有安全感.
最有歸屬的地方.
人離開自己的國家已經不容易.
離開一個土生土長的故鄉更加艱難.
與自己的父母,兄弟姐妹道別.
更加痛苦.
上帝吩咐亞伯蘭.
與自己的同胞,親人的親告別.
割斷與親人,社團的聯繫.
目的是要帶他去另一個新的地方.
去哪裡呢?.
在這一刻.
神沒有提供一個正確的目的地.
只有一個方向.
上帝只是將方向顯示給亞伯蘭知道.
甚至連細節都沒有提出.
不過亞伯蘭心裡想.
如果方向是正確的話.
其他的問題就應該很容易解決.
雖然他不知道要去哪裡.
但是亞伯蘭知道.
吩咐他去哪裡的那個是誰.
那個是耶和華.
為何上帝吩咐亞伯蘭.
要離開自己的家鄉呢?.
有學者認為是基於吾爾的宗教.
與他的道德很墮落.
亞伯蘭要離開這個地方.
才可以與將來他的後裔.
遵守上帝的道.
秉行公義.
這個說法其實不是很準確.
因為迦南的宗教與道德.
並不會比吾爾優勝.

$^{121}$那麼上帝的吩咐.
為何要亞伯蘭離開吾爾呢?.
答案在下一節的經文.
「我必使你成為大國.
我必賜福給你.
使你的名為大.
你要使別人得福.
為你祝福的.
我必賜福給他.
詛咒你的.
我必詛咒他.
地上的蠻族.
都必因你得福」.
如果亞伯蘭聽從上帝的命令.
離開了國家,故鄉.
前往上帝指示的地方.
就會得到很多上帝給他的應許.
這些應許包括了.
成為一個大國.
你想想.
亞伯蘭當時是一個人.
要成為一個大國.
很難想像.
「我必賜福給你.
叫你的名為大.
他也要叫別人得福.
為你祝福的.
就必賜福給你.
地上的蠻族.
都要因你得福」.
如果今天耶和華跟你說.
Gary.
你要因你而得福.
有很多人因著你的生命而得福.
我們會怎樣回應?.
我們會想.
我們只是一個小蝦米.
但因為接觸我的人而得福.
你想想我是否一個很蒙福的人?.
或者換另一個說法說.

$^{161}$上帝給人類的祝福.
由亞伯蘭一個人開始.
然後一路擴張開去.
直到萬國萬族.
關鍵是亞伯蘭願不願意.
踏出這個信心的第一步.
亞伯蘭怎樣回應?.
亞伯蘭就遵照耶和華的吩咐去了.
羅德都與他同去.
亞伯蘭離開哈蘭的時候.
七十五歲.
根據經文記載.
亞伯蘭離開的時候所攜帶的東西.
有財物.
亦有獲得的人口.
而且亞伯蘭當時七十五歲.
按照亞伯蘭在一百七十五歲.
離開世界的計算.
亞伯蘭當時是壯年.
這樣的年紀.
又薄有家財.
在當時的生活.
應該都相當捨異和休閒.
不過亞伯蘭完全吩咐.
耶和華的吩咐.
離開本地本族.
夫家.
帶同他的財物.
家人前去耶和華所指示的地方.
亞伯蘭在這一刻.
完全信服在耶和華的心意裡面.
亦對耶和華投入一個信心的回應.
對於當時的亞伯蘭.
其實他不是身無一物.
他不是單身.
無牽無掛的.
他有家人.
有親戚.
還有財富.
他其實可以留在自己的家鄉.

$^{201}$過一個悠閒的生活.
對於神賜給他的福.
其實是將來的事.
要放棄目前的生活.
去追求將來的福氣.
是一個對神信心的表現.
而要表達出這個信心的回應.
就是要完完全全地信服的行為.
亞伯蘭最後踏上一個信心的旅程.
去回應神對他的呼召.
當亞伯蘭到達神所指示要去的地方.
即是迦南事件的時候.
耶和華再次向亞伯蘭應許.
他說我要將這個地賜給你的後裔.
這個應許就相當具體了.
是一塊地.
一塊可以世代相傳的地.
一塊可以落葉歸根的地.
於是亞伯蘭即刻在事件.
為耶和華築了一座壇.
之後亞伯蘭搬到伯特利附近的時候.
他再次為耶和華築了一座壇.
求告耶和華的名.
築壇以及求告耶和華的名.
就是公開地去敬拜耶和華.
以及在迦南地.
向迦南人正式公開他和神的關係.
這個公開的宣告是很重要的.
除了表示一個敬拜的禮儀之外.
他還表示亞伯蘭全心全意地去事奉耶和華.
當亞伯蘭要祝福.
當耶和華要祝福亞伯蘭.
使他的名為大的同時.
亞伯蘭也同樣高舉耶和華的名.
專主為大.
當亞伯蘭在迦南地去實踐他的信仰的時候.
亦是耶和華透過亞伯蘭.
去祝福他身邊的人的開始.
亞伯蘭回應神的呼召.
彰顯了一種完全信服專主為大的態度.

$^{241}$如果我們在這一刻去檢視亞伯蘭的品格.
我們可以說他是取得很高分的.
雖然處身在光輔的生活當中.
他仍然肯回應神對他的呼召.
是否取得很高分呢?.
各位弟兄姊妹.
神昔日呼召亞伯蘭跟隨他.
今日神對每一個信徒都發出同樣的呼召.
他說:來吧!跟隨我.
我們如何回應呢?.
昔日亞伯蘭願意離開安書的生活和環境.
今日安書的生活會否成為我們跟隨主的半腳石呢?.
昔日亞伯蘭在迦南地公開宣認他與神的關係.
今日我們願不願意在我們的職場.
我們的生活圈子公開我們基督徒的身份呢?.
昔日亞伯蘭全心全意地去事奉神.
今日亞伯蘭的經歷能否成為我們全心全意去事奉神的榜樣呢?.
甚至亞伯蘭的生命能否成為我們每一個.
跟從主耶穌的鼓勵和幫助?.
信仰實踐的第二個提醒.
就是活出生命就是最好的信仰實踐.
亞伯蘭雖然充滿信心.
但並不是表示沒有軟弱或沒有試探.
事實上危機即刻就發生在他面前.
經文說:那地遭遇饑荒.
糧食短缺不單單只關乎亞伯蘭一個人.
是關乎整個族群的生與死的考慮.
亞伯蘭作為族長.
肩負整個家族的安危.
他於是決定去埃及.
要在那裡寄居.
亞伯蘭想的如意算盤就是.
埃及是一個大國.
物資豐盛.
應該可以彼得過饑荒.
而且我們只是去寄居.
我們又不是在牆上住.
饑荒過後我們才回到迦南.
不單單如此.
亞伯蘭深思細密.

$^{281}$想到到了埃及之後會發生什麼事.
他估計到有兩件事.
他知道自己的老婆撒萊很漂亮.
埃及人一定會被撒萊的美貌所吸引到.
埃及人會因為垂涎撒萊的美色而搶奪她.
如果知道亞伯蘭是她的丈夫.
一定會下手加害她.
亞伯蘭有這樣的想法不是沒有道理的.
因為傳說在埃及的法律裡面.
禁止搶奪有夫之婦.
於是埃及人想到一個辦法.
既然喜歡別人的老婆.
就將她的老公殺了.
逃避了法律的制裁.
於是亞伯蘭想到一條絕世好主意.
他自己是撒萊的哥哥.
因為當時的風俗.
如果有人喜歡一個女孩.
而她父親不在場的話.
哥哥就可以全權辦理妹妹的婚姻大事.
如果到時真的有人想娶撒萊打她主意.
任何求親的人都要先跟亞伯蘭交涉.
亞伯蘭就可以施展拖延的方法.
然後尋求脫身之計.
的確是一條可以保命的好方法.
而事情的發展就一如亞伯蘭所料.
他們一到埃及.
埃及人就真的覺得撒萊很美.
很想打她主意.
唯一是亞伯蘭始料不及的.
就是看中他老婆的人.
不是一般平民百姓.
而是埃及法老王.
法老的神僕還將撒萊帶入宮.
亞伯蘭的拖延政策完全發揮不到.
他現在可以怎樣做呢?.
亞伯蘭自以為聰明絕頂.
現在卻束手無策.
最後關頭還是要神出手.
神介入在事件裡面.

$^{321}$降大災在法老和他家人身上.
令法老要查明原因.
他召見亞伯蘭.
向她質詢.
你向我做的是甚麼事呢?.
為甚麼要向我講大話呢?.
為甚麼不告訴我撒萊就是你老婆呢?.
為甚麼騙我.
說她是你妹妹呢?.
亞伯蘭無言以對.
經文沒有提到亞伯蘭有甚麼辯解.
最後法老沒有懲罰亞伯蘭.
他只要求亞伯蘭和撒萊.
以及他們全部一切.
所有的東西.
離開埃及.
法老不想再和亞伯蘭有任何的聯繫.
他不想惹禍上身.
聖經將亞伯蘭這兩個生命片段放在一起.
兩段生命的實踐之間有很大的落差.
上一段.
亞伯蘭對神很有信心.
完全信服神的心意.
下一段.
他沒有提到亞伯蘭和神的關係.
至少饑荒來到的時候.
他沒有問耶和華.
究竟應不應該去埃及.
上一段.
神祝福他.
要使別人得福.
下一段.
他的行為令到法老和家人受到大災.
不單祝福不了人.
還令到人有苦難.
法老甚至不想和他有任何的關係.
在上一段.
他和妻子.
一起走在神的心意裡面.
下一段.

$^{361}$他和妻子要隱瞞身份和關係.
隱藏了他們和神的關係.
一切的落差.
源於亞伯蘭的害怕.
他害怕在饑荒裡面不能夠生存.
他害怕透露了他和撒萊的夫婦關係之後.
會被殺害.
亞伯蘭心裡面有擔心.
有害怕.
有恐懼.
以致他走歪了.
甚至可以說.
他走在神的面前.
而不是走在神的後面.
緊緊跟隨他.
今天我們可以馬後炮地說.
亞伯蘭沒有信心.
是嗎?.
不過亞伯蘭的害怕.
那種感覺是很真實的.
真實到我們不能當他不存在.
我們不能夠抽離現實.
一味說只要有信心.
這樣就行.
事實上.
如果亞伯蘭不去埃及.
有機會整個家族可能會餓死.
如果亞伯蘭不認撒萊做妹妹.
可能他真的會死在埃及裡面.
而撒萊就真的可能被埃及人搶了去.
亞伯蘭所提出的方法.
可能在當時來說是一個最好的方法.
也未必的.
也可能是唯一一個可以保存性命的方法.
也未必的.
信仰如果真的要實踐的話.
真的不能夠空談的.
也不能夠站在道德高地裡面指指點點.
當我們說要實踐信仰的同時.
其實我們要預埋.

$^{401}$我們有十足十的代價.
負代價的準備.
而這個負代價可能就是我們連命都要給他.
在這裡.
現在這個疫情還沒有過.
但我想起上一次那個疫情.
十八年前.
當沙士來到的時候.
有一段這樣的報導.
在2003年的3月.
他說嚴重的急性呼吸道症候群.
疫情在香港爆發.
屯門醫院接收了三個沙士病人.
院內的空肺的專科醫生是不足夠的.
當時的謝宛文醫生自願.
由內科病房轉去沙士病房.
是自願的.
當時情況危急.
謝宛文與另一個殉職的男護士劉永佳.
先後親自為病人插喉.
懷疑因此感染致病的病毒.
4月3日留院治療.
4月15日轉入深切治療部.
中文大學沈祖堯和香港大學的袁國勇.
曾經為他聯診.
5月13日凌晨四時搶救無效.
而逝世,享年35歲.
一個我們很熟悉的故事和人物.
35歲,與亞伯蘭一樣.
正值盛年,黃金歲月.
他當時可以不舉手.
為何要自願?.
他不需要自願走入沙士病房.
他也可以不與病人插喉.
因為插喉後,不保證能夠救回.
他可以獨善其身,保存性命.
繼續做一個專業的內科醫生.
不過,他選擇了用生命去實踐他的信仰.
當年胡說呼召亞伯蘭的時候.
第一件事就是要亞伯蘭離開他的安居區.

$^{441}$神母保證亞伯蘭事事順風順水.
他沒有保證到了迦南後.
仍然需要面對饑荒.
亦沒有保證到埃及,不會被人殺害.
當生命源於一線的時候.
當害怕的感覺很真實,很淋度的時候.
我們不能單單說要有信心.
這樣就可以跨過去.
要多些甚麼去支持我們呢?.
除了說我們的信心之外.
信仰的核心不是單單解決罪的問題.
不是單單在神面前認罪後.
就甚麼事都可以自然解決.
信仰亦不是翻了多少年的主日學.
或是我們整本聖經讀了多少次.
有深厚的聖經知識固然不能夠少.
但信仰還要說和神的關係.
與主同在,與主同行的生命.
如何實踐在生命裡面呢?.
可能對某些人來說很陌生.
甚麼叫與主同在?.
甚麼叫與主同行?.
可能從來都未經歷過.
生命的見證都是信仰的核心的一部份.
神人關係是雙方面的.
神對人有美好的心意.
神會祝福保守每一個跟從祂的人.
對於神,祂永遠都不會失信.
只是我們人失信.
人對神的信心不單單是建立在知識層面裡.
即我們所說的,不單單是聖經的認知.
我們還需要建立在經歷裡面,關係裡面.
我們與神的關係是怎樣?.
我們有沒有經歷過神?.
如果我們的實踐由累積回來.
一路一路地累積.
能夠有很多與主相交,與主同行的經歷.
而產生出來的信心的話.
這一種的經歷,這一種的信心.
才能夠幫我們面對危機.

$^{481}$當我們面對生與死的時候.
我們仍然能夠跨過去.
各位弟兄姊妹.
透過亞伯蘭的兩個生命的經歷給我們反省.
當神呼召我們的時候.
我們能不能夠以一個信服的心.
尊主偉大的心去回應.
而擁有與主同行,與主同在的經歷.
就能夠活出生命.
這也是我們最好的信仰實踐的方法.
也是建立我們的信心的最好方法.
今天的講題是「絕不離地的信仰」.
「絕不離地的信仰」是說什麼呢?.
是說經過信仰的高峰和低谷的亞伯蘭.
他的生命經歷告訴我們.
信仰從來都不會離開我們的生活.
或者這麼說,信仰跟我們的生活是分不開的.
如果我們強行要將信仰跟生活分割開.
星期日回來,我們講的是信仰.
離開教會,我們講的是生活的話.
信仰對於我們來說是毫無意義可言.
我們可以說,亞伯蘭在這次的事件裡.
有軟弱的部分.
但是不代表他整個生命都是軟弱和失敗.
他仍然可以在日後的生命當中繼續成長.
只要他願意與神同工,與神同行.
我們的信仰生命其實跟亞伯蘭沒什麼分別.
我們都是有血有肉的.
我們可能在初信主的時候.
我們跟神的關係是很好的.
經過一段時間,我們可能又冷淡了下來.
我們殷純了.
我們可能在初信主的時候.
我們可以有很多靈修的生活.
但是今天我們冷淡了.
甚至我們連看聖經那種渴求都沒有.
無論我們的生命是開心還是傷心難過.
只要我們願意被神介入在我們的生命裡.
那就足夠了.
我們的生命能夠有神介入.

$^{521}$自然就有神的恩典承托著我們整個生命.
只要我們有與神同在的經歷.
才培養得到堅定不移的信心.
而這種信心能夠幫助我們.
當我們面對害怕的情況的時候.
有害怕的感覺的時候.
我們仍然可以很勇敢地去面對和承擔.
當生命遇到挑戰的時候.
我們仍然能夠挺起胸膛.
昂首闊步地去面對.
我們一起去祈禱.
我們從亞伯蘭的生命我們見到.
神對亞伯蘭是不離不棄的.
祂呼召亞伯蘭.
在祂信心最高的時候.
祂沒有離開他.
但是當亞伯蘭軟弱的時候.
信心跌到谷底的時候.
神同樣都沒有離開他.
我們今天可以稱亞伯拉罕為信心之父.
這個信心之父是怎樣來的?.
是用他一生的經歷去累積而來的.
我們的人生都一樣.
我們的生命經歷過起起伏伏.
高高低低.
無論我們信心高的時候.
神沒有離開我們.
我們信心低的時候.
神都沒有離開我們.
甚至我們每一個.
我們都能夠緊緊地跟隨著主耶穌.
過那種與神同在.
與神同行的信仰生命.
因為唯有如此.
我們和神的關係.
才能夠緊緊地結連在一起.
我們既然被神去救贖.
將我們從罪裡面拯救出來.
我們就要在信仰裡面.
我們能夠好好地珍惜.

$^{561}$這種與神復和的關係.
正正是我們已經和神復和了.
我們就能夠享受那種與神同工.
與神同行那種的甘甜.
但願我們每一個弟兄姊妹.
都能夠經歷這種和神親密的關係.
以致我們的生命能夠被神所大大的使用.
願主賜福給我們.
我們祈禱交托.
奉教主耶穌基督.
得勝名自祈求.
阿們.
願主祝福大家.
謝謝大家.
\newpage



\section{以弗所書 4:11-16}
\label{sec:_2_W0T2LR08}
\textbf{2021.08.08《同心合意建立基督的身體》 以弗所書 4\_11-16 (和修版) 曾敏姑娘}
\newline
\newline
連結: \href{https://youtube.com/watch?v=_2_W0T2LR08}{\texttt{https://youtube.com/watch?v=\_2\_W0T2LR08}} ~~~~ 語音日期: 2021-08-11
\newline
\newline
\hyperref[sec:G34IvTirgOo]{\small{< < < PREV SERMON < < <}}
~
\hyperref[sec:index_chronic]{\small{[返順時目]}}
~
\hyperref[sec:index_scriptual]{\small{[返順卷目]}}
~
\hyperref[sec:mw0QSsEsPEU]{\small{> > > NEXT SERMON > > >}}
\newline
\newline
以弗所書 4:11-16
\newline
\begin{longtable}{cl}
\hline
\hline
章節 & 經文 (和合本修訂版)\\
\hline
4:11 & \begin{tabularx}{0.7\textwidth}{X} 他所賜的有使徒，有先知，有傳福音的，有牧者和教師， \end{tabularx} \\ \\ \relax
4:12 & \begin{tabularx}{0.7\textwidth}{X} 為要裝備聖徒，做事奉的工作，建立基督的身體， \end{tabularx} \\ \\ \relax
4:13 & \begin{tabularx}{0.7\textwidth}{X} 直等到我們眾人在信仰上同歸於一，認識神的兒子，得以長大成人，達到基督完全長成的身量。 \end{tabularx} \\ \\ \relax
4:14 & \begin{tabularx}{0.7\textwidth}{X} 這樣，我們不再作小孩子，中了人的詭計和欺騙的法術，被一切邪說之風搖動，飄來飄去。 \end{tabularx} \\ \\ \relax
4:15 & \begin{tabularx}{0.7\textwidth}{X} 我們反而要用愛心說誠實話，各方面向著基督長進，連於元首基督， \end{tabularx} \\ \\ \relax
4:16 & \begin{tabularx}{0.7\textwidth}{X} 靠著他全身都連接得緊湊，百節各按各職，照著各體的功用彼此相助，使身體漸漸增長，在愛中建立自己。 \end{tabularx} \\ \\
[1ex]
\hline
\hline
\end{longtable}
$^{1}$去敬拜祂 親近祂.
我們開始的時候再有一個禱告.
天父我們來到你的面前.
知道你深深的認識我們每一個.
在這一刻.
我們在你的面前無所遁形.
裡面的好與不好.
都顯露在你的眼中.
天父我求你再一次赦免.
我們裡面一切得罪你 得罪人的地方.
很願意耶穌基督以你的寶血.
潔淨 遮蓋我們.
免得這些罪惡攔阻我們到你的面前.
又願意聖靈在我們眾人心裡面去工作.
讓我們今天有一個很開放的心靈.
聽上帝你的聲音.
我們就去切切地回應你.
求你賜福.
奉耶穌基督的名字祈禱.
阿們.
這段時間我在喪禮當中很多體會.
其中一個喪禮是我們教會籌備的.
已經完成了.
我看到在座的一些弟兄姊妹.
在過去一段時間都和我一起.
去籌備一個喪禮.
是天婆婆的安息禮拜.
沒有多久之前.
在當中我看到上帝透過弟兄姊妹.
將愛傾倒在家人的心裡面.
在那一天分別有三個不同的群體.
在安息禮拜裡面獻屍.
過程中很大轉外.
感恩那一天有人幫我們做一些協調的工作.
所以幾個師班坐的位置剛剛好.
是更加方便獻屍的服侍.
而那一天除了獻屍.
還有主禮,祈禱,讀經,講道.
還有之前預備程序表.
還有陪伴喪家的弟兄姊妹.

$^{41}$很多人一起都有美好的配搭.
大家一起努力將喪禮搞好.
這一個就是神的家.
是神的家.
我心裡面覺得是很感動,很溫暖.
另一個喪禮是我爸爸的喪禮在籌備當中.
雖然我來到洪恩家不到一年.
但是弟兄姊妹很關心我.
給我很多慰問和關懷.
有些人很專心在背後為我禱告.
有些人要在我爸爸的安息禮拜裡幫忙.
在這裡我很感受到眾弟兄姊妹對我的那份愛.
更加是上帝的愛.
是很滿滿的愛.
在這裡衷心地多謝大家.
也多謝我是天上的父神.
祂讓我們可以在神的家裡生活.
在神的家裡生活很開心.
神的家就是基督的身體.
是非常寶貴和榮耀的.
上帝希望我們怎樣去對待基督的身體呢?.
這正是今天我很想和大家分享的.
第一件事很明顯的.
就是神很希望我們同心合意地建立基督的身體.
是同心合意地建立基督的身體.
在這裡聖經說道.
上帝首先把五類人.
是和上帝的說話有相關恩賜的人.
賜給了教會.
他們分別是使徒,先知,傳福音的牧者和教師.
使徒很快地說是有廣義和合義之分的.
廣義是說受到教會委派的代表.
例如聖經裡有些人物是提多,以巴弗提.
我記得我在這裡說過以巴弗提這個人物.
又或者是一些巡迴的宣教士.
而合義的意思就是夢神呼召和委任.
對於基督的復活作見證和傳福音的.
包括了十二使徒和保羅等人.
這些人是見過復活的主.
是主自己親自委派去服侍的.

$^{81}$這些是合義來說一班的使徒.
先知就是神說話的傳講者.
不只是在舊約裡有些先知.
到了新約的時候也是有的.
而使徒和先知是基督奧秘的領受者和傳揚者.
是具有權威性的.
牧者的職責比較明顯一點.
是慰養,關懷,照顧,保護教會的群揚.
好像今天教會裡的牧師,傳道等等.
他們要帶領和治理教會.
而提到慰養.
就一定是和教導神的話語有關的.
教師都是專心做教導聖經的.
教導真理工作的人.
所以你看這五類人.
全部都是和上帝的說話很有關的.
這些人的服侍是很重要的.
但並不代表整間教會.
上帝將擁有這五類恩賜的人賜給教會.
是為了兩件很重要的事情.
第一件事情.
要裝備聖徒.
做事奉的工作.
為了要裝備聖徒.
我們這裡就要想一想.
要裝備聖徒要裝備多少個呢?.
教會裡有沒有數的數目呢?.
是一個?.
兩個?.
還是怎樣?.
在經文裡很明顯不是一個,兩個這麼少.
而是每一個.
這些聖徒本身.
都從上帝領受了不同的恩賜.
以弗所書四章八字這樣說.
所以有話說.
他聖上高天的時候.
「勞掠了俘虜,將各樣的恩賜賞給人」.
耶穌基督勝過死亡.
從死裡復活.

$^{121}$祂勝過仇敵.
祂擁有無比的權柄去處理這些恩賜.
祂將它賞賜給所有人.
是所有人.
是每一個屬於祂的人.
所以弟兄姊妹我想說.
我們每一個人都是有恩賜的.
今天就坐在洪恩家台下的每一個弟兄姊妹.
你看看你自己.
你看看你身邊前後左右的弟兄姊妹.
每一個都是有恩賜的.
既然是這樣.
我們每一個人都應該按照自己的恩賜事奉.
不應該由一小撮人.
做完教會大部分的工作.
今天我們回看洪恩家的時候.
是不是每一個弟兄姊妹都有事奉呢?.
或者不要只說洪恩家.
不同的教會.
教會裡面是不是.
真的每一個弟兄姊妹都有事奉呢?.
還是普遍出現一個現象.
就是一小撮人.
他們承擔了大部分的事奉呢?.
你會留意那些人.
又做主席.
又做領唱.
又做師班.
又做影音.
又要教主人學.
是同一批人來來去去.
但是還有很多的弟兄姊妹又如何呢?.
聖經給我們看.
每一個.
神都給我們有恩賜的.
你的恩賜有沒有發揮出來呢?.
這裡值得我們去想一想.
經文這裡說.
當這個聖徒領受了恩賜之後.
神也會透過.

$^{161}$剛才我提過的那些五類的.
有特別恩賜的領袖.
去裝備聖徒.
以致這些弟兄姊妹就可以參與事奉的工作.
事奉的工作是非常多的種類.
是很多很多.
有些人是容易看得見.
有些人是看不見的.
這個是第一件事.
好了.
第二件事是做什麼呢?.
上帝將擁有五類剛才說的那些恩賜的.
那些領袖賜給教會.
除了裝備聖徒做事奉的工作.
還有和那些聖徒一起.
建立基督的身體.
這一個就是今天的重點.
要去建立教會.
去建立教會的時候.
不是某一個人.
保羅自己很有恩賜.
保羅自己的恩賜是很獨特的.
他寫了很多卷的福音書.
為上帝做了很多的福音工作.
建立教會等等.
就算一個這麼有恩賜的人.
他不會只是想著自己一個去服侍.
他要和弟兄姊妹一起.
上帝賜給教會有這麼多.
和話語服侍有關的領袖在這裡.
不是只是要他們幾個.
或者某些人去服侍教會.
而是要他們裝備其他人.
大家一起去建立基督的身體.
少一個都不可以.
在這裡是非常特別的.
我們可以問一個問題.
就是我們同心合意.
建立基督的身體要到什麼地步呢?.
這個是以弗所說很想說的.

$^{201}$今天究竟上帝對於紅恩堂的心意是什麼.
是不是只是想我們每個星期回來聚會.
看到弟兄姊妹很開心.
打一下招呼.
接著聽一下道.
唱一下詩歌.
感受一下自己今天開不開心.
是不是這些呢?.
我肯定這個不是上帝的心意.
上帝的心意想我和你同心合意的.
建立基督的身體要到什麼地步呢?.
第一件事.
我們要在信仰上同歸於一.
在真道上要合一.
我們在這裡領受了真道.
我們明白上帝要我們相信的是什麼.
要我們追隨的是什麼.
要我們跟從的是什麼.
但是我們要在這些真道裡面.
要一起合一.
不是你有你的跟從.
我有我的跟從.
而是我們要一起的.
我們這個群體不分彼此.
這是第一件事.
第二件事我們要認識神的兒子.
我們要知道耶穌基督是神的兒子.
上帝將祂差派來到這個世間.
耶穌基督是神的兒子.
這是一個很重要的教義.
有很多異端邪說.
說起耶穌基督的身份.
耶穌基督的作為的時候.
並不會這樣去認同.
耶穌基督真是神的兒子.
是說了很多很多不同的東西出來.
原來這個我和你.
一定要達至一致.
我們要認定耶穌基督真是神的兒子.
我和你的教義是一致的.

$^{241}$在這個裡面我們是要合一的.
建立基督的身體要達到這些.
信仰在真道裡面合一.
在這個教義上合一.
第三件事很重要.
長大成人.
經文這樣說.
長大成人原文的意思.
就是要成為成熟的男性.
其實這是一個比喻.
我們建立基督的身體.
什麼叫成為成熟的男性呢.
我們這裡有弟兄有姐妹.
意思是在靈性上要成熟.
要很穩定.
在真理上要成熟.
要穩定.
穩定到什麼地步.
是達到基督完全長成的新娘.
是要完全像基督的.
能夠展現基督完美的品格.
今天我想問弟兄姊妹.
我們這個群體.
今天我經常說群體.
不是說我們個人.
我們這個群體成熟嗎.
我們穩定嗎.
在靈性上我們是怎樣呢.
會不會我們很受人影響.
人有變動.
我跟著變動.
人不變我就不變.
我自己能不能夠站立得穩.
我的屬靈生命經不經得起風浪.
風浪一來就立刻倒塌.
左搖右擺.
是不是呢.
這就好像基督.
人們看我們整個群體.
看不看得到耶穌基督.

$^{281}$不是說某一個人.
某一撮人.
是整個群體.
今天當上帝看著大家.
在敬拜祂的時候.
你覺得神會給我們一個什麼評分.
外面有人進來.
看到我們敬拜上帝的時候.
覺得我們像不像耶穌基督.
有多像.
有多能夠展現基督完美的品格.
在經文裡.
就將成熟和不成熟做了一個對比.
成熟的男人和小孩子有很大的差別.
小孩子是很不穩定的.
屬靈的小孩子是不成熟的.
經文說到會因為.
邪說之風而動搖飄來飄去.
邪說之風這裡是說一些很新奇的教義.
假教義.
或者聽起來好像是真的.
但實際是說謊的.
異端邪說等等.
小孩子.
會被這些影響.
不成熟的人會被這些影響.
不成熟的教會都會被這些影響.
成熟的教會不會.
會站立得穩.
我們今天的光景又如何.
我記得洪恩堂剛剛在這裡.
設立堂祠的時候.
那時候經常提議.
要很小心附近的異端.
我看到當中有人點頭.
起初真的是.
我知道我們附近有很大間的.
異端的教會.
他們都很努力傳他們的信息.
當時說我們要互教.

$^{321}$要守護弟兄姊妹.
不要受到這些異端邪說的影響.
當時是這樣.
其實今天也是.
今天他們仍然很努力傳他們的信息.
我看到在不同的地方.
除了我們本身視為正統派的教會.
很努力傳福音之外.
也有很多異端教會很努力傳他們的信息.
在過程中.
我們懂得分辨嗎.
今天別人跟我們分享信仰.
你猜我們夠不夠膽跟別人說.
他那本所謂的正典.
他從頭到尾滾瓜爛熟.
我們那本聖經呢.
我叫你今天說幾節經文.
你能說多少節呢.
是不是.
不容易.
是不是.
他們也說到.
屬靈的小孩子.
還有一個問題.
就是不成熟的人.
他會中了人的詭計.
和欺騙的法術.
這裡其實不只是說人的詭計.
欺騙的法術.
原來這些是來自於惡者.
惡者會透過這些去搞擾人心.
欺騙人.
叫人走在歪的路上.
成熟的人可以分辨得到.
不成熟的人分辨不到.
成熟的教會可以抵擋這些.
不成熟的教會.
就會落在網羅裡面.
所以這裡他透過小孩子這個例子.
作為一個負面的教材.

$^{361}$提醒我們.
我們一起建立基督的身體.
就是要成熟.
最重要是要站得穩.
屬靈上很穩陣.
今天我們給自己多少分呢.
是不是很穩陣.
接著很重要.
怎樣同心合意建立基督的身體呢.
有幾樣東西值得成為我們的提醒.
第一.
要用愛心說誠實話.
這裡的誠實話其實是在說真理.
不只是說真話還是假話.
是不是.
哎呀曾姑娘.
你好肥呀.
是不是.
說真話這樣.
或者你覺不覺得又多了一些皺紋.
又或者是.
覺不覺得今天走路又慢了.
這樣的誠實話.
不是說這一種.
而是真的說真理.
呼音的真理.
聖經的真理.
信仰裡面的真理.
在教會裡面.
你怎樣去建立基督的身體呢.
要有真理.
有時候真理說出來令人不舒服.
因為那個照到我們內心不對的地方.
照到我們生命軟弱的地方.
我們不想別人說我們.
是不是.
最好你不要說.
最好見面說些好聽的話.
曾姑娘.
是不是.

$^{401}$你好好.
或者什麼這樣.
最好令到我舒服就好了這樣.
但上帝要我們在這裡說真理.
在我們的群體裡面要彼此提醒.
這件事是很需要有勇氣的.
我在和弟兄姊妹的聊天裡面.
試過不止一次.
我說這件事不對的.
其實是要提出來的.
我們自己不提.
希望別人提.
是不是.
這個是不行的.
上帝就想我們自己提出來.
你說我們自己提.
是不是比外人提到我們更不對呢.
別人說我們就很醜.
是不是.
失去見證.
教會內部就是要彼此提.
不對的地方.
但我們就很怕得罪.
我們怕得罪人卻不怕得罪上帝.
其實應該要先敬畏上帝.
不要怕得罪人.
但是說出來的時候.
又要在愛裡面.
所以用愛心說誠實話.
有另一個極端.
有些人是不怕的.
不怕得罪人的.
是不是.
我告訴你.
應該怎樣是更加好.
是對的.
如果按照聖經的教導.
說起來就是一把刀捅進別人的心臟裡.
是不是.
直說就插進去.

$^{441}$哇.
痛得別人不得了.
真是.
要在愛裡面.
你真的愛你的弟兄.
你的姐妹.
你愛他出於愛跟他說.
你希望挽回他.
你希望提醒他.
讓他改變.
在愛裡面.
愛是一個很重要的.
《已乏所恕》裡面的一個主題.
失去愛是不行的.
原來你提醒他是為了愛.
你不是想徹底毀滅對方的生命.
我明白了.
當初你提醒我這樣.
我覺得很不舒服.
是不是.
現在我知道你是為我好.
在愛裡面將真理說出來.
第一樣東西.
怎樣同心合意建立基督的身體.
第二樣東西很重要.
少了一定不行.
我們有同一個方向.
有同一個動力.
就是要向著基督長進.
連於元首基督.
靠著祂全身都連接得緊湊.
要有耶穌基督.
這個連於和連接都是現在的事態.
表示這種連結要一直堅持.
你不可以在一分鐘的熱度.
不斷不斷地堅持.
過去堅持.
現在堅持.
將來堅持.
你要建立一個基督的身體.

$^{481}$就是要有這種堅持.
什麼堅持.
和耶穌基督結連.
讓耶穌基督作我們的頭.
這樣教會才可以在主裡面保持合一.
才能夠成長.
這個值得我們探討.
向著基督長進的意思.
就是越來越似耶穌基督.
經文裡面說是各個層面都是.
每樣東西都越來越似耶穌基督.
你永遠評分都是這樣想.
那我像不像耶穌基督呢.
這部分像不像耶穌基督呢.
我們可以逐步的.
每一部分的服事都會去想的.
你去想我們教會有多少個部服事.
每一個部是不是像耶穌基督.
不只是部.
教會有不同的元素.
每個元素像不像耶穌基督呢.
向著方向.
「連於元首」.
耶穌基督是我們的頭.
人不是教會的頭.
記住沒有任何人是教會的頭.
耶穌基督是教會的頭.
整個身體要靠耶穌基督結連.
是祂將我們黏在一起.
這裡第一件事不是某個人.
有時候我們很努力追求.
我們屬靈的生命.
是某一些人.
某一個人在追求.
是不可以的.
是要整個群體.
所以我今天說這個訊息.
是對著台下每一個弟兄姊妹說.
是我們每一個都要緊緊黏著耶穌基督.
緊緊以祂為我們的頭.

$^{521}$每一個都要.
你不要只是說.
行了 祂做就行了.
祂們做就行了.
不關我的事.
我又沒有什麼事奉崗位.
剛才我說了.
每個人都有恩賜.
每一個人都要事奉.
每一個人都要一起建立這個身體.
每一個人都要靠耶穌基督.
什麼叫和耶穌基督結連呢?.
這個就很有趣了.
是不是我經常讀經.
靈修祈禱.
就已經是結連呢?.
我想問問大家.
是不是基督徒每一天.
靈修讀經祈禱就是結連呢?.
不一定的.
我們是要讀經靈修祈禱.
但我們的心有沒有到?.
我們是不是真的和耶穌基督.
有一個真正的關係呢?.
是不是真的明白耶穌基督的心意呢?.
知道耶穌基督喜歡什麼.
不喜歡什麼.
知道之後是不是跟呢?.
是不是真的走出來呢?.
還是每一天打開聖經.
看完了.
行了 結束了.
我都叫做靈修了.
就這樣祈禱.
我最喜歡.
我教主的時候很喜歡這樣的.
第二堂就問了.
我上一堂講過那樣是什麼.
就是什麼.
就問了大家.

$^{561}$然後我發覺.
已經把那些東西都給我了.
如果問你.
哪天你自己靈修.
那段經文講什麼呢?.
對你有什麼幫助呢?.
是不是?.
你是不是真的在你生活裡面.
實踐出來呢?.
你怎麼回答我?.
主要是學時覺得好聽.
講故事很多東西聽.
聽完之後.
下一堂回來.
其實我都不知道.
不知道發生了什麼事.
對你生命有什麼意義呢?.
是不是真的.
晝夜思考上帝的話.
你真的活出來.
有時候明明.
我們心裡面聽上帝的聲音.
是很清楚的.
上帝跟我們說.
這樣不是很好.
是不是?.
我聽了未必跟.
我就是不跟了.
這樣是不是結年呢?.
這樣的生命會不會成長呢?.
真的黏著.
結年.
結年是會有信服在的.
他怎麼教我們.
就要怎麼跟.
你宣告.
你以行定宣告.
不是你把口說.
你是我的主.
你說是這樣.

$^{601}$就這樣吧.
就算我未必很想.
我都聽你講.
信服.
我信靠.
沒有別的了.
是要黏著你.
我才有辦法.
是要黏著你.
才有出路.
我不會去找ABCDE.
這麼多人.
全靠那些人.
你黏著他.
不要分開.
不要一個半個黏.
如果今天玩遊戲.
我邀請台下所有弟兄姊妹一起.
就是要全部人用盡力量.
一起黏著耶穌基督.
是不是?.
如果耶穌基督.
今天現身在我們當中.
玩的遊戲就是要一起.
黏在他身上.
一起.
這樣聖經才說.
教會基督的身體.
就會慢慢地成長.
是要一起.
少一個都不行.
還有一個過程.
第三點很重要.
除了靠耶穌基督.
就是我們要什麼的配合.
百節各按各職.
照著各體的功用.
彼此相助.
這個在聖經上.
有一些難題.

$^{641}$百節各按各職.
實際的原文是說.
是每一個支持的人帶.
每一個支持的人帶.
直接將不同的部分.
結連在一起.
這些人帶很重要.
它將不同的身體的肢體.
結連.
這個人帶可以解釋為.
剛才我們說.
上帝賜給教會的.
那些與說話恩賜有關的.
教會領袖.
這些教會領袖.
再加上後面的.
各體的功用.
就是不同的弟兄姊妹.
不同的弟兄姊妹.
都有領受.
從上帝而來的.
那些恩賜.
教會領袖.
與弟兄姊妹一起.
彼此相助.
互相幫助.
這樣去建立基督的身體.
我自己在洪恩家服侍的.
大半年.
都很深刻的體會.
我會在這裡.
有我自己要參與的.
服侍的部分.
但我發覺我需要很多弟兄姊妹.
我是不可以這樣的.
是不是.
或者老牧師自己都不可以一個人的.
一人組在這裡服侍.
我們需要很多人.
我講一個笑話給大家聽.

$^{681}$我想影音組是最熟悉的.
就是.
每隔一陣子.
我就會有一些.
虛幻的東西.
有一次.
就應該是.
我想將USB.
插在手機上.
傳檔案的.
但應該是插不實.
所以手機是讀不出來的.
當時插不實自己又不是很懂.
證明自己又不是很熟悉手機.
不熟悉電腦的東西.
一直都很著急.
我記得那天梁傳道.
要帶整個影音組一起來看.
這個場面真的很震撼.
大家就一起研究.
是什麼問題呢.
導致我讀不出檔案.
後來一試.
原來是沒有插實.
我相信大家當時很想笑.
原來是這樣的.
之後又一次試過.
跟電腦有關.
又是沒有做好某一個小小的步驟.
結果又是搞不定.
這個過程我真的認定.
自己的電腦是不太行的.
對手機又不算是很熟悉.
所以我才需要弟兄姊妹.
今天大家看到台上的花很漂亮.
你叫我去插花.
其實我應該是沒有這樣的美感的.
真的不簡單.
每次看他們.
為什麼可以這樣呢.

$^{721}$這麼厲害的.
可以做到這麼漂亮.
很有層次感.
很有藝術感.
真的是神給他們的恩賜.
還有很多人是我很欣賞的.
一些電腦設計的人.
嶺昌.
有很多等等.
在過程中.
其實我真的發覺我們彼此需要對方.
還有一些人關心人.
是特別窩心很溫暖的.
我們跟著他們一起都學到東西的.
全部人都可以從弟兄姊妹身上學到東西的.
只是自己一個.
只是小小的童工.
能夠做到多少呢.
整間教會需要弟兄姊妹起來.
在這個過程我們各人發揮自己的恩賜.
發揮自己的長處.
這樣互相幫助.
我需要你.
你需要我.
然後.
基督的身體就會成長.
教會就會成長.
我們就在愛中建立自己.
在這裡我們就經歷到愛.
正如我剛才一開始的時候.
沒錯.
我在說的是喪禮.
但是我提到的.
在教會裡面經歷愛.
我在別人的喪禮經歷到愛.
在自己家人的喪禮經歷到愛.
不只是這些.
是在教會事奉裡面.
我們常常都可以經歷到愛.
這個就是上帝的心意.

$^{761}$想我們在這裡彼此相愛.
在愛裡面就建立自己.
我們建立這個身體.
也都建立自己的生命.
我很盼望將來有幅很美的圖畫.
紅恩家.
是每個弟兄姊妹起來服侍.
其實我們可以有崗位的服侍.
也可以有非崗位的服侍.
是不是.
崗位的服侍很多的.
很多.
你看得見的.
今天在這裡會場裡面的.
不同的弟兄姊妹在參與的.
那些是崗位的服侍.
還有很多什麼非崗位呢.
很重要的.
例如禱告.
禱告.
如果今天在座這麼多個.
都同心合意禱告.
我們的教會會很興旺.
一定會復興.
弟兄姊妹在這裡會經歷到愛.
在這裡服侍都會成長.
上帝會將更多的人.
帶到我們的當中.
上帝也都要叫到我們.
自己成長.
你不要說這裡坐飽.
將來紅恩堂有機會.
變成一個婚堂.
你有沒有想過這幅圖畫呢.
禱告.
每個人都禱告.
還有很多.
我們都可以做關懷的工作.
你回去問阿撒.
給你什麼恩賜.

$^{801}$我希望今天.
你回去就做這個禱告.
我們一起的祈禱.
很愛我們.
你給我們有無比的榮幸.
你將傳道人牧者賜給教會.
你將那些.
有教導說話恩賜的.
領袖賜給教會.
裝備信徒.
然後一同建立神的家.
你一同建立神的家.
你想我們成熟.
賜禮.
是個過程.
你要我們依靠你和你結連.
你讓我們在信仰上站立得穩.
你讓我們.
無論傳道同宮頂.
姐妹彼此結連.
彼此的配搭.
最後我深信.
上帝要透過這樣.
榮耀你自己的名.
也都要透過這樣.
讓我們在愛裡面去生活.
去成長.
是這樣愛著我們.
我們不配得.
很願意你叫弘恩家.
不斷的復興.
成為你喜悅的教會.
將來這個教會在這裡.
是一盞很大的燈台.
不是為了榮耀我們任何人.
只是單單的為了榮耀上帝你.
讓上帝你的心意得到滿足.
奉耶穌基督的名字祈禱.
阿們.
\newpage



\section{列王記下彼得後書 3:9}
\label{sec:mw0QSsEsPEU}
\textbf{2021.08.15《宣教的異象》 彼得後書 3\_9;列王記下 5\_1-27 (和修版) 高淑芬傳道}
\newline
\newline
連結: \href{https://youtube.com/watch?v=mw0QSsEsPEU}{\texttt{https://youtube.com/watch?v=mw0QSsEsPEU}} ~~~~ 語音日期: 2021-08-19
\newline
\newline
\hyperref[sec:_2_W0T2LR08]{\small{< < < PREV SERMON < < <}}
~
\hyperref[sec:index_chronic]{\small{[返順時目]}}
~
\hyperref[sec:index_scriptual]{\small{[返順卷目]}}
~
\hyperref[sec:h6Nm0N8GuaU]{\small{> > > NEXT SERMON > > >}}
\newline
\newline
$^{1}$你好.
今天我們會講這個列王記下的很美麗的圖畫.
還沒有到使徒時代.
你就會看到上帝讓這件事.
讓一個愛邦人.
一個高官在位大能的勇士.
貼服服地相信馬耶和華這一位的真實.
在這本書講這個異象的時候.
宣教這件事.
甚至這個事件裡面.
你會看到有很多個人的質素.
讓這位大能勇士去相信主.
今天也會順道教大家一些泰文.
好不好?.
好.
那你留心了.
有很多個「濟」.
你要記住這些「濟」.
第一個「濟」就是「偉濟」.
「偉濟」泰文的「濟」就是心的意思.
「偉濟」在第十一幅那裡.
「偉濟」的意思就是可以很信任你.
「濟偉」就是擺的意思.
我可以信任你.
我的心可以擺在.
就是很放心.
也有些像我們廣東話.
放心下去.
泰文就是這樣的意思.
「偉濟」.
我們可不可以做到一個呢?.
「吞係倌圓偉還濟呆」.
就是讓別人都可以很信任你.
我想在宣教裡面.
這個信任是很重要的.
無論是一個警察會.
支持的教會.
支持的人士.
你一定要有一份信心.
就是給一個信任感.

$^{41}$給前面那個人.
那個人去了.
好像我一樣.
我去了泰國.
我是沒有親戚在那裡.
所以我皮膚很黑.
人們都以為我是.
你是不是泰國華僑.
我不是.
我去那裡是舉目無親的.
我真的很信任後面的團隊.
警察我去之後.
都會遙遠距離照顧.
所以這份信任是很重要的.
當然你自己也要做好自己.
你去那裡.
是不是每天都去布吉島呢?.
當然我們那裡去布吉島.
真的是山長水遠的.
是去不了的.
在這一書我們會看到.
「偉濟」.
「偉濟」就是信不信任.
你有沒有這份信任.
在這個人的身上呢?.
整件事.
「乃萬佢信主」.
「他相信上帝是真的神」.
你會看到.
是有很多人.
那份信任連結在一起的.
「亞蘭王」這個元帥.
不要忘記.
他的主人的事.
是在「亞蘭王」的面前.
「亞蘭王」就是「變哈達二世」.
「亞蘭」這個國家.
就是敘利亞的國家.
他是一人之上.
萬人之下的一個神父.

$^{81}$是一個很權重的人.
但是在聖經裡面第一節有說到.
上帝是藉著「乃萬」.
來讓他得勝的.
透過他的手來教訓以色列人.
他是一個大能的勇士.
但是很可惜.
在未到可惜之前.
他這麼大能的時候.
我想我們如果有些長官.
或者在學校裡面.
有些校長,老師.
他做得好,教得好.
我們都很喜歡和他聊天.
或者和他拍照.
都很好.
那麼這位大能的勇士.
其實他得到人民的愛戴.
他有他的榮華富貴地位.
什麼都有.
很可惜人生裡面.
不是每樣都有齊.
總是有一些缺陷.
他很遺憾地.
他有了這個麻風.
這個麻風.
大家都知道是去到1940年.
才可以能夠治好一些麻風病.
換句話說.
1940年之前.
麻風是一個人見人驚的事.
當然它不像我們今天.
我們這個COVID-19.
傳染性很快.
它是有傳染.
不過它的傳染沒有這麼快.
比較慢一些.
但是它的潛伏期.
最糟糕是你開頭.
中了招都不知道.

$^{121}$5到20年.
這個是很慢很慢的情況.
但是問題是.
因著不同的病人.
它會侵略神經系統.
呼吸道.
皮膚 眼睛.
甚至乎皮膚上會出現一些肉瘤.
慢慢侵害神經系統.
連痛楚 感受都可以慢慢消失.
所以你會見到麻風病人.
手和鼻子都脫離了.
他們不覺得痛.
不過就是消失了.
慢慢連肌肉都可以沒力.
你試想一個大能的勇士.
他要帶著他的部隊去打仗.
他的肌力慢慢的失強.
視力慢慢的模糊.
他的巔峰.
就可能是慢慢慢慢的衰退.
所以這一位大能的勇士.
雖然是很勇猛.
但是都敵不過病菌的侵襲.
就好像我們今天這樣.
這個病菌襲擊我們.
甚至乎一路一路的變種.
我們看到這個病菌的底下.
無論你是有錢有權位的人.
我們都看到有人過了生.
或者去到基層的人士同樣都會有.
所以無人倖免.
那個病菌不會認識你是誰就給你面子.
不會的.
所以在這個麻風.
臨到乃曼身上的時候.
他很享受著現在他有這一份榮華富貴.
和他真的用生命拼回來的這一種的尊貴.
打勝仗.
但是很可惜.

$^{161}$潛伏著他的就是這些細菌.
一路一路的去侵襲他.
我相信他在家的時候.
夜間人或者.
總之不容易在外面的人.
去露出他這一份威脅.
大麻風對他的那種威脅.
但是我相信他都會有他那種擔憂.
憂傷.
顧慮在當中.
直到有一天.
他的太太竟然跟他說.
我聽到以色列的小女子這樣跟我說.
就是說你去找我們薩馬利亞的先知.
你一定得救.
這個的.
如果你是乃曼將軍你會怎麼想呢.
我是一個勝國來的.
我是打贏仗的那個.
你這個小女子都是.
那一次打贏了你們國家.
擄你回來的一個俘虜.
你做事做得很好.
今天你就是照顧我太太的人.
但是你的說話信不信得過呢.
我相信.
乃曼裡面都有著很多的矛盾.
或者有人都會說.
這個可能都是他已經窮途末路了.
我相信他在皇宮裡面如醫名醫.
一切民間的一些.
那個醫得好.
都會走去.
我想全部.
總之可以救他的他都會去做.
但問題是.
這個是他打敗了的一個國家.
哪裡會有好東西出呢.
因為這個觀念.
是在當時的舊約時代.

$^{201}$普遍的觀念.
我今天A國.
我打贏了你B國.
我的A國是強大的.
我A國所拜的神都是強大的.
你這個B國輸定了.
因為你B國的神不行.
他們是這樣想的.
但是.
乃曼他都去試.
我相信不只是去到.
沒有路了.
平一平.
因為這個在面子.
在什麼都過不去.
我相信他願意去一試的時候.
是因為小女子.
他這個.
我上一次都有講到.
小女子他去到.
照顧乃曼太太這個位置.
我相信是.
是有很多他的.
做事的交不交代.
他這個人品如何.
他的態度是怎樣.
是有時間去看到.
這個人.
真的可以委祭的.
是可以委祭的.
是可以信任他的.
他才來考慮.
我相信他的太太.
他告訴乃曼的時候.
也是這樣重覆重覆去想.
應該信得過.
去.
這樣.
以至乃曼去信.
都是一連串.

$^{241}$你會看到是一連串的信任.
信任小女子.
信任他的太太.
以至到.
好.
我們.
他就去了.
你會看到.
人與人之間的委祭.
無論是在警察團裡面.
這個是很重要.
你讓你的教會.
讓你的家人.
警察會.
他放心.
你在遠方宣教.
你這個在遠方宣教的人.
你又要放心.
我在那裡生活到.
我在那裡.
是在事奉的.
很多東西的平衡.
那個信任很重要.
我在這裡很感恩.
我已經有好幾次.
一想到這件事.
覺得非常之.
就是.
你們很多謝教會的信任.
給了我一個地方.
我可以工作.
以至到.
我們泰語團契就知道.
開頭我都告訴你.
不是.
開頭的時候.
我沒有筆記給他們.
因為我沒有這些工作的場地.
當我有張台.
有個電腦去工作的時候.

$^{281}$他們從此之後.
就有筆記了.
這個.
他們問的時候.
我都說多謝教會.
給了我一個地方.
我可以去工作.
更加的就是.
我有教會的鎖匙.
你們怕不怕.
我都很受寵若驚.
拿著鎖匙.
所以我就想像到.
給得不容易.
給得拿著天國的鎖匙.
我拿著鎖匙.
第一天關門的時候.
外面又在下雨.
我多害怕.
我的意思是.
我在想三層樓要關的東西.
我有沒有關好呢.
我就拿著廢柴.
對不起.
我在說泰文.
我拿著的廢柴是.
是電筒.
廢柴就是電筒.
因為我一個人在這裡的時候.
他們都很好.
樓梯都不會開燈.
我覺得我只是用房間的電就夠了.
所以差不多.
真的沒有人.
我就拿著廢柴.
不好意思我們這些人.
語文都調轉了.
所以我下次說廢柴.
我是在找電筒.
因為我真的有個電筒.

$^{321}$在我的房間.
我那個廢柴.
我就拿著電筒.
巡一下房間.
因為我覺得.
既然教會這麼信任我的時候.
我都應該.
你明不明白人家信任你的時候.
你都應該有一個信任的回應.
你信任我.
我都是信任得過的.
在那個責任的範圍裡面.
你會希望做好一點.
我就會巡下來關燈.
關上他們很友善的樓梯燈.
我覺得可以的.
我有支廢柴可以看到東西.
可能我在泰國都習慣了.
我們閃一閃電.
我們的燈就沒有了.
所以有一次曾姑娘都笑.
你幹嗎這麼有趣拿著廢柴.
我說不要緊.
習慣了.
那裡一會閃一閃電就沒有了.
東西就什麼都沒有了.
所以我已經很久沒有用電腦來工作了.
因為是不容易的.
還有我發現.
有了這個位置的時候.
去工作的時候.
原來我是失去了一個安靜工作的辦公室.
這麼多年來.
原來我已經忘記了.
當我坐在那裡工作的時候.
我覺得很感恩.
因為原來我去到泰國之後.
很多事情都是即步即介.
你開了桌子.
聚會來了快點收桌子.

$^{361}$每一次都要設置好.
然後才做.
特別我們都不能坐定定.
都是要拿著車.
所以短孫隊來到就知道.
我的車你們在上次的片見過.
很大部.
很多東西的.
我隨時都是.
塞車的時候就用腦去想事情.
眼睛就好好地看著你開車.
就是因為整個運作就是要到處走.
塞車不知道什麼時候.
所以我很感恩.
這一份信任.
我相信.
這個是我們宣教的素質裡面所需要.
也是我們每一個.
遵從上帝的大使命的人.
要有的一個很重要的素質.
就是偉濟.
你是一個偉濟的人.
你是一個值得我信任的人.
當我們去.
你說你去傳福音給人.
特別我在泰國.
人家說三點鐘來.
我三點鐘等.
雖然他可能四點五點.
但是總之他說的點數.
我就會在那裡出現.
是這麼久以來.
就是這份信任而建立到.
我這個人不是說說就算了.
會做的.
所以求神都看顧給我們.
給我們一個偉濟的人.
這個已經是很中文的.
就是燈濟的濟.
那個濟就好像濟字的濟.

$^{401}$就是信心的人.
以後我們雄恩堂有些.
有些密語.
有些暗號.
你偉不偉濟?.
偉濟.
這樣.
然後.
其實我告訴你.
為什麼要記住廢柴.
因為我覺得.
香港話就不是很好.
什麼廢柴.
但是我覺得.
在泰文裡面.
其實我下台後.
我都好像你們一樣.
一樣要照顧爸爸媽媽.
很多一樣的人.
我經常覺得.
雖然在過去上帝呼召我做傳道.
後來做宣教士.
但是都是一個.
和大家一樣生活在這個世上的人.
廢柴給我.
特別當我在泰國認識.
這個叫廢柴.
的時候我覺得上帝真是很有用.
因為剛才我都說.
泰國整天.
閃一閃就沒電.
特別我在車裡.
因為太黑.
就算開了這個燈.
都很黑.
我都是用廢柴來照著.
去找我的東西.
你可能覺得自己.
是一個廢柴.
但是你不要小看我.

$^{441}$那個廢柴很小.
但是都不知道多有用.
會照亮著我要走的路.
今天可能我們會看自己.
我都不行的.
這樣不行.
那樣不行.
對不起.
在這個這麼黑暗的時代裡.
很需要.
即使你覺得你這個廢柴.
很小很小都好.
很有用.
剛剛前天.
去買一些東西.
是那些日用品.
店員拿不到.
不是不是.
在裡面.
他說卡住了.
我說是呀是呀.
你的手指卡住了.
他說是呀你知道.
七月份我要帶這些.
這些是銀.
七月份銀.
我都要定一定時.
不是定一定神.
他在說什麼.
我知道了.
你要信耶穌了.
因為我和他不可以太久相處.
他拿完東西給我.
就走了.
他說你要信耶穌了.
我唯有這樣回答他.
你想想那個很年輕的.
真的很年輕的少女.
不是很害怕.
不是一個害怕的人.

$^{481}$你表面看不到.
但她竟然這麼害怕.
七月份要帶銀在手上.
是呀.
就是這樣.
你就會看到.
原來我們周圍.
還是一些很黑暗的地方.
求神給我們.
就算做廢柴.
都做得很精彩.
在這些黑暗的地方.
都可以點一點.
你試一下.
很黑很黑的地方.
很小的光都幫到你.
很大很大的用.
在乃萬將軍.
去接納這個少女子的意見的時候.
你會看到.
是一連串的信任下.
甚至乎.
我們可以去第十二福那裡.
甚至乎你會看到.
連亞蘭王.
你可以去.
都沒有想過他會做義務.
你是不是去.
投奔以色列國.
背棄亞蘭國呢.
他不是.
完全很信任.
好吧,我幫你寫信給以色列王.
你會看到.
這些建立.
是君臣之間.
他們都有這份信任.
你知道嗎.
如果得到別人的信任.
你在家裡.

$^{521}$你在公司.
你在很多人際相處裡面.
你是會舒服很多.
還有你的發揮會大一點.
所以這份信任.
求神給我們.
雄恩家委祭.
是一個令人放心的.
一個這樣的屬靈的家.
接著.
那一幅我們就要學一下.
登祭.
真的音好像是登祭.
但下一幅.
就是專注的意思.
在這件事裡面.
雖然我們看到乃萬.
他是很專注.
很登祭的.
就去.
去什麼.
去這個以色列國.
他就是專注於.
去面對這件事情.
就好像我們亞運完了.
不是,奧運完了.
你看到我們的運動員.
雖然他可能腳受傷了.
雖然他可能.
已經在落後.
但是他專注在比賽裡面.
專注.
一分一分地打.
專注這個素質.
對於我們每一個都很重要.
特別我們說.
我們專注在神的身上.
專注地.
跟神有一個好的關係.
這個都是很重要的事.

$^{561}$特別大家知道.
現在的社會越來越變.
很多人都在笑.
我們遲一點會不會.
連一個.
教師信賴的機構.
政府都不信賴.
遲一點會不會輪到.
輪到我們呢.
輪到教會呢.
倘若真的有這樣的日子的時候.
我們更加要專注.
我們是上帝的兒女.
你的專注的心.
有沒有放在上帝的裡面呢.
乃萬呢.
他知道.
要去以色列國裡面.
去找.
去求醫.
他真的是一個大能的勇士.
他真的按部就班.
一件一件地去做.
他一人之上.
萬人之下.
他就要去以色列王那裡.
去告訴.
他告訴蘭王.
我要去以色列國.
蘭王就很好.
很有信任地寫了一封信給.
乃萬去.
去這個以色列國那裡.
我們看到.
他真的是大能.
真的懂得高懂得低.
又是.
擺放在他的處事裡面.
甚至乎呢.
你看到.

$^{601}$他是預備了很多.
很多的禮物.
帶過去給.
這個.
以色列國的.
這一處除了帶了一封.
王信.
亞蘭王寫的信.
他還拿了.
很多的錢.
在經文裡面.
我們看到.
他拿了.
六千的.
他帶了十他年得的銀.
六千赦夏勒的金紙.
他年得.
赦夏勒.
這些都是他那些.
道量行的東西.
這樣.
十他年得.
其實我算了一個數.
但是我的數再參考.
其他人算完的數.
我的數好像少了一點.
一會兒再告訴你少了多少.
這樣.
但是我想.
不要這樣搞了.
不如搞一下在聖經裡面.
怎樣讓我們看回視野.
究竟他帶了多少多少的錢去.
在《催拉及記》.
21章32節那一段說到.
給我們聽.
銀紙.
三十赦夏勒.
就可以.
買了一個人.

$^{641}$那些是奴僕.
婢女.
所以你想想.
小女子更加是.
俘虜回來.
更加不值錢.
連三十赦夏勒都沒有.
這個.
乃萬所帶的十他年得.
就有.
三萬個赦夏勒.
他又可以買到一千個奴隸.
這麼巨大的錢.
好了.
甚至乎.
這個十他年得的銀.
如果再.
在《列王記》上.
十六章再前一點的時候.
以色列當南北兩個分開的時候.
暗利王.
他只是用二他年得的銀.
就已經可以在.
撒瑪利亞那裡.
買了整個山頭.
做他們的成都.
這樣.
是二他年得.
所以說.
他帶的錢.
可以買了五個.
撒瑪利亞山的價錢在那裡.
更加.
你會看到.
是很多.
很多的錢在當中.
更加.
你會再看到.
好.
然後.

$^{681}$再下一幅.
麻煩.
六千赦夏勒.
是金.
大概是六十八公斤.
所以為什麼.
最後剛才你讀經的時候.
哈基西這個僕人.
他拿了一些東西.
要乃萬的僕人幫他抬.
抬回去.
是很重的.
很重禮的.
這個六千赦夏勒.
可以是怎樣呢?.
想當年.
這個也是想當年.
相比現在.
乃萬這個朝代.
因為是大胃王的時代.
他們只是用六百赦夏勒.
的金.
就已經可以買了整塊地.
整個土地.
但是現在.
乃萬為了治好自己的大麻風.
帶了六千赦夏勒的金.
去找.
這個.
以利沙這個先知.
你就會看到.
他是.
很重很重的住.
我剛才說我自己計算.
我是少於別人一點點.
十他年得.
大概是我們的港幣多少呢?.
513萬.
都很厲害.
接著.

$^{721}$六千赦夏勒.
就是這個.
我的數字真的不行.
是否.
599萬.
加起來也有成千萬.
但是這個算是一個比對.
我是參考人那些.
我這樣算是少了一點點.
但是怎樣也好.
也很多.
對於我們來說.
有這麼多.
我們.
這麼多.
你會看到乃萬的專注.
他願意不惜所有.
因為他是一個.
大能的勇士.
他去到別人的國家.
這是戰敗國.
他不想.
有一種.
好像我欠了你一樣.
下次打仗我就要留守.
所以清清楚楚.
帶足夠的銀兩.
多到不得了.
是很多.
相比當時前後.
就算有通脹也不用這麼多.
足夠有餘的錢.
你醫好我.
我給你醫生費.
這樣.
所以你看到.
乃萬他完全這樣.
很專注.
他希望得到醫治.
甚至乎給醫生費.

$^{761}$重金禮聘.
已經帶過去.
甚至乎有亞蘭王的信.
還有權威.
這樣.
非常之好.
有這些東西去到.
很可惜.
若果.
在這個專注裡面.
我們看到.
他開始發脾氣.
因為帶著這麼多東西.
去見以色列王的時候.
當然想著找以色列.
找最大的.
最大的.
你就會找先知來醫我.
誰知以色列王就怕得.
面青口唇發.
怕得不得了.
真是很悲哀的.
還要怕.
我想我們人生都會遇到.
一些很怕的懼怕.
我很怕買好衣服.
買好衣服是小強.
姐妹.
不要這麼喜歡買好衣服.
在泰國千萬不要.
因為那些是買好十.
那個音應該是買好十.
但我有一次教過短孫.
他說買好衣服.
那些泰國人懂得聽.
從此之後.
挺有趣.
買好衣服可以叫小強.
這樣.
那一次我完了聚會之後.

$^{801}$我還要開車去.
那時候我們還是去舊的機場.
因為泰國有兩個機場.
我就是去舊的.
舊的就叫.
在楞士那裡.
我們去的地方不算太遠.
不過問題是那段時間.
晚上.
我們那時候.
不是現在這個教育.
早期在學生宿舍.
不是說有一個單位變兩個單位.
那個位置.
外面.
一出門口這裡.
很大的水喉.
不是,是坑渠蓋很大的.
那段時間是很多蟑螂.
群起.
出來.
我的車就這樣停在那裡.
我已經很快.
因為在這樣的地方.
你已經要想好.
一上車要快點關上門.
又不要夾到手指.
這樣.
我以為我搞定.
我準備要開車.
走去機場那裡.
去接人.
一邊開車的時候.
我覺得我腳.
有點不妥.
但是你開車.
膝蓋後面.
好像有些東西.
我呢.
怕是買了衣服.

$^{841}$我又要鎮定.
在宣教裡面.
經常教我遇到驚慌也好.
也要鎮定一點.
何況我開車.
我在那裡祈禱.
其實在車裡面除了我辦公室.
也是我和上帝.
經常交往的地方.
上帝就死了,我很害怕.
如果他再.
向前走快一點.
走到再近距離.
再近距離.
我肯定會撞車.
會失控.
不知道會飛去哪裡.
上帝讓我鎮定.
一邊開車一邊鎮定.
這樣.
慢慢.
開到旁邊.
也算是安全的地方.
我呢.
非常之.
立刻.
停下來.
快點飆出去.
大力地.
真的掉下來了.
我立刻把它扔掉.
因為我怕他會跟著我上車.
很害怕.
就是這樣.
但我仍然要專注.
專注地上車.
專注地泊車.
泊好位置.
處理這些害怕的事情.
在.

$^{881}$這個乃曼.
她的專注裡面.
我們看到她是一路.
專注.
但會不會有時候我們也像乃曼.
遇到一些卡位.
我們就停了步.
幸好.
幸好她有些僕人.
就在這裡.
我的父啊.
人家先知如果叫你.
做一件大大的事.
你也會做.
因為乃曼在想.
這個先知真的不懂得抗拒.
我是一個大能的勇士.
我是.
敘利亞的元帥.
歡迎也沒有.
也沒有露面.
就這樣寫一個僕人.
告訴我去.
約旦河湖.
泡七次.
我沒有去過.
但看了資料就知道.
約旦河這個湖很小.
水也不清.
很濁.
如果在他們大馬式的兩條河.
聖經說的兩條河.
水質也比這裡漂亮.
他開始在想.
是不是水質可以.
醫治好他的大麻風.
如果用水質.
來說真的沒有比例.
但是.
這個已經讓他.

$^{921}$生氣.
想退縮.
因為之前.
好好的見你這個以色列王.
叫你找人醫我.
他怕成這樣.
怕得像什麼.
他一個錯一個不對.
再來一個不對.
以致他.
就是.
都消沉了.
人生都會有消沉的時候.
但是很需要一些.
委祭,可以委祭的人.
委祭是什麼.
信任的人.
去扶你一把.
這些僕人是他所信任.
這樣說的時候.
又說得合乎中道.
如果這個才知道.
根據他們的古代文化.
可能要你在這裡跳舞.
泰國.
泰國我們曼谷.
那裡有個四面佛.
有很多明星去拜.
然後要去那裡.
還神諸如此類.
有些人說.
要在那裡.
跳舞.
他們都會做.
這些是迷信.
他們覺得做一些事.
就可以贏取一些醫治.
他們就用.
這些僕人用.
因為他們專注.

$^{961}$我們來是希望元帥.
你好.
現在只差一點點.
為什麼不試他.
浸七次.
就是這樣.
以致乃萬斯放下.
預設好的東西.
願意.
下去浸七次.
所以.
周圍的人.
雖然有時候我們會擋住.
但周圍的人提醒我們.
我們可以再在那件事.
專注.
接下來有個「濟」.
剛才我們的「濟」叫什麼.
「登濟」.
是專注.
「濟」的意思是.
我很肯定.
很確實的知道.
這個在第十六頁.
「濟」.
很確實的知道.
很清楚的知道.
當乃萬.
其實很不容易.
在約旦河七次.
他不是漸進式好返.
我浸了下去.
好了一點.
第二,第三,第四次,第五次.
很踴躍下去.
但不是這樣.
一次浸完想回去.
而且他是一個大名人.
周圍的人怎會看住.
他真的.

$^{1001}$放下了他的尊嚴.
下去這條.
他看不起的那條河.
來去洗.
一次,兩次.
可能第三次已經算了.
但我們看到他.
又很好.
很登濟地繼續去洗.
洗了七次.
七次之後神跡先來.
很多時候.
我們信神.
我們有很多事.
都好像關注了神的作為.
求神幫助我們.
七次的洗.
一上來的時候.
奶孟的皮膚.
我們真的非常學悟.
來到現在的時候.
肉復原.
復原已經很開心.
一個大麻風的人.
復原不止.
超乎他所想所求.
好像小孩子的肉.
已經復原.
復原.
小孩子的肉.
今天有個小女孩.
來探望我們.
如果不是的話.
待會你去街上留意一下小孩子的肉.
或者回去看看你的孫子.
小朋友.
小孩子的肉.
表面.
可能他們的僕人都會看到.
奶孟你真的完全不同.

$^{1041}$那些白色的.
麻風造成的東西.
都沒有了.
奶孟的皮膚.
變了有彈性.
滑,沒有皺紋.
總之皮膚很漂亮.
你看看她的肉.
好像小孩子的肉.
她不只是皮膚,是肉.
那些肌肉.
她感受到,奶孟一定會感受到.
那些肌肉.
真的.
她還未認識上帝的時候.
上帝都要透過她做大事.
她經歷上帝之後.
我相信上帝一樣會繼續.
讓她成就大事.
當她經歷到.
上帝這個恩典的時候.
奶孟.
當然去報恩.
去找那個先知.
因為從頭到尾她都沒有見過先知.
終於找到他,好回來了.
你會看到.
她很肯定.
她知道現在她所信的.
是獨一的真神.
她知道以色列國的神.
是一個救星.
是一個救主.
她也知道.
我雖然有權位.
物力.
我是不能夠.
去買到.
神的一些恩典.
這個恩典不是買得到的.

$^{1081}$她很感受到.
這個上帝是真的時候.
她很願意.
給她帶來的寵禮.
給了以尼沙.
但是以尼沙.
她說我向神發誓.
我不會要你的東西.
然後.
奶孟她真的是一個大能的將軍.
她真的很有.
高低.
還有輕重.
她說到.
當我.
因為她有責任在身.
她是一人之上.
萬人之下.
她也差不多是皇帝的保鏢.
皇帝去臨門廟.
臨門廟不是臨門一腳.
不過很好記.
就是這樣記到她的廟.
臨門廟當時.
是一個拜雷神.
總之去拜廟的時候.
我就要屈一屈身.
扶她的時候屈身拜.
但是我在這裡.
好像劃清界線.
我是履行我的職責.
我不是拜祂.
我從今以後.
我要向這位神.
我要向這位.
你會看到.
她這個神跡.
在她身上的時候.
她有很大很大的轉變.
甚至乎.

$^{1121}$她連這件事都.
報了在案.
雖然以尼沙不要她的禮物.
但是她也是.
根據過去.
過去.
舊有的文化.
就是.
她不可以.
她不會移民.
因為她應該要回去.
忠於她的敘利亞國.
但你可不可以給你的土地.
因為他們的相信是這樣的.
我要拜你的神.
我要不就在你的土地裡面.
去拜你的神.
要不我就要拿你的土地回去.
做一個壇.
以至到我去拜你的耶和華神.
以至到以色列的神.
她連這件事都很願意去做.
這件事.
在.
在新約福音.
耶穌有提過.
就是在路加福音.
第四章第27節.
祂說到先知.
以尼沙的時候.
以色列裡面有很多長大麻風.
但是裡面.
除了敘利亞的拉曼.
之後.
沒有一個.
得到潔淨.
這件事是事實.
在聶王記下.
五章一到二十七節我們看到這件事.
但是.

$^{1161}$那些文士.
律法師.
祭司他們聽到的時候.
他們非常之.
憤怒.
他們怎樣.
當時的聽眾.
想置耶穌於死地.
想推祂下山.
因為靠近山崖.
但是.
耶穌沒有讓他們做到.
你會看到.
為什麼他們會.
憤怒到要殺耶穌.
因為在他們.
猶太人的.
觀念裡面.
上帝是我們的.
外邦人.
他們看低我們外邦人.
不值得看他們.
我們是低等一點.
你們的外邦人.
是不可能的.
耶穌還要說.
在當時.
上帝醫好敘利亞的.
拉曼元帥.
他更加.
覺得不可思議.
不願意福音.
外流.
他是唯威位在我們以色列人的那裡.
但是我想告訴你.
上帝的心腸.
不是這樣.
上帝的心腸.
是不願一人沉淪.
我們求神讓我們有齊.

$^{1201}$這個素質.
我們需要建立.
我們繼續.
去第十七頁.
就好像一個茶壺.
茶壺裡面.
沒有東西在裡面.
你倒什麼出來給人.
所以你一定要好好地.
增進自己.
給自己內涵.
素質.
這些東西可以向上帝求.
你看到有些人說.
孫教聖甚麼都不怕.
我剛剛上星期.
打針.
打針之前我一直拜託.
姐妹們為我祈禱.
因為我平時打針都很怕.
抽血又怕.
我每次去醫生都說緊張.
你很緊張.
你一定要抽血.
今次打針.
還要有這麼多新聞說打完針.
又這樣又那樣.
真的很害怕.
不過都要打.
但是.
可以.
鍛鍊.
甚至我們不是一個人去練.
是我們一班人一起.
洪恩家一班人一起去練.
就很不同了.
我們會看到我們的IT.
IT群組.
不是星期天早一點來.
這麼簡單.

$^{1241}$他們平時都要來.
不斷練習不斷受訓.
你也知道他們.
不懂電腦.
但可以訓練成這樣.
很厲害.
我們是互相的配搭.
一起去配搭.
我們在崇拜的時候.
有個姐妹就要很努力.
將我的廣東話.
變成英文.
去翻譯給我們有.
尼哥利日利亞的一位弟兄.
知道這位利日利亞的弟兄.
原來是在洪恩堂洗禮.
然後參加泰國團契.
是我們的一份子.
你看不看到洪恩堂很特別.
就好像舊約聖經.
不是去.
傳福音.
叫他們來.
是他們不知怎樣.
上帝就讓他們擺在這裡.
洪恩堂也有這個.
一個也好.
你看不看到.
我們真的要有人.
來配搭才能做到.
這個造就的功夫.
就算在團契裡面.
大家也要忍耐.
因為一篇課程我們又要中文.
泰文,英文.
所以我歡迎你們.
來探索一下.
不過要忍耐.
三次.
這個就是.

$^{1281}$洪恩堂已經有的.
跨文化工作已經給了你們.
你們有泰語的團體.
有利日利亞英語的群眾.
我知道還有菲律賓的.
不過他們不是常來.
你看到我們互相.
我們洪恩家真的.
我們自己就是一個茶會.
茶會裡面有甚麼東西.
我們不可以只用盡它.
不時也要加些東西.
加些盡心.
接著下一幅.
我們要承認.
有時我們回到小社會.
我們這個舊約裡面.
下一幅是第十八.
第十八那幅.
我們會看到.
別人有些.
都不是很委濟.
委濟是甚麼.
他做事.
都不是很登濟.
又是甚麼.
接著.
他是否太厲害.
他相信的是甚麼.
他太厲害是甚麼.
很確定.
你要知道.
我們都是每一個人.
我自己都是.
工程進行中.
還未完成.
我有心之人.
這個人我還是工程進行中.
不知道你們同不同意.
大家都在.

$^{1321}$目標上.
繼續去行.
但是工程進行中.
我們求主給我們有自省.
原來我們.
待人處理不太好.
改進.
我希望大家不會像.
以利沙那個.
那個僕人.
基哈西那樣.
他要去到.
自己有大麻風.
他才知道.
究竟他相信的是誰.
他已經跟著.
以利沙.
以利沙是雙倍的零.
下去的一位先知.
他跟著這位先知.
這麼久.
竟然他都不懂.
他都不知道.
他究竟服侍.
我的心意是怎樣.
我們主禱文.
今天歌裡面有.
我們主禱文說我們.
願遵行父的旨意而行.
我們知不知道上帝的旨意.
上帝的旨意.
跟著下一幅.
在彼得後書三章九節那裡說.
主所應許的尚未成就.
有人就以為.
他是擔言.
其實不是擔言.
是寬容我們.
人人沉淪乃願人人悔改.
上帝願意萬人得救這個心.

$^{1361}$在舊約時代.
還言那一班服侍神的人.
法理錯人.
文士.
他們都不曉得.
不是那麼明白.
甚至到今天.
猶太人仍然是.
自己覺得.
我們的利益.
是我們自己的利益.
我們自己的利益.
是我們自己的利益.
所以我們有.
警察會.
專門向.
以色列人傳福音.
告訴他們.
利塞亞已經.
來就是耶穌基督.
很多他們本來.
是上帝揀選的子民.
但他們都不肯定.
都不確定.
他們所信的是什麼.
我們求神.
求神給我們.
我們不斷改進的時候.
我們知道.
我們的神的心意.
更加摸索到上帝的心腸.
願意紅恩堂已經有好的開始.
有好的開始.
有跨文化的工作.
在這裡.
願意我們繼續.
讓周圍的人.
多人信主.
願意神的心願.
同樣的成就.

$^{1401}$求主祝福每一位紅印堂弟兄姊妹.
\newpage



\section{使徒行傳 16:6-10}
\label{sec:h6Nm0N8GuaU}
\textbf{2021.08.22《使命必達》 使徒行傳 16\_6-10 (和修版) 謝靄薇牧師}
\newline
\newline
連結: \href{https://youtube.com/watch?v=h6Nm0N8GuaU}{\texttt{https://youtube.com/watch?v=h6Nm0N8GuaU}} ~~~~ 語音日期: 2021-08-25
\newline
\newline
\hyperref[sec:mw0QSsEsPEU]{\small{< < < PREV SERMON < < <}}
~
\hyperref[sec:index_chronic]{\small{[返順時目]}}
~
\hyperref[sec:index_scriptual]{\small{[返順卷目]}}
~
\hyperref[sec:65pDMB41cZg]{\small{> > > NEXT SERMON > > >}}
\newline
\newline
使徒行傳 16:6-10
\newline
\begin{longtable}{cl}
\hline
\hline
章節 & 經文 (和合本修訂版)\\
\hline
16:6 & \begin{tabularx}{0.7\textwidth}{X} 因為聖靈禁止他們在亞細亞講道，他們就經過弗呂家、加拉太一帶地方。 \end{tabularx} \\ \\ \relax
16:7 & \begin{tabularx}{0.7\textwidth}{X} 到了每西亞的邊界，他們想要往庇推尼去，耶穌的靈卻不許。 \end{tabularx} \\ \\ \relax
16:8 & \begin{tabularx}{0.7\textwidth}{X} 他們就越過每西亞，下特羅亞去。 \end{tabularx} \\ \\ \relax
16:9 & \begin{tabularx}{0.7\textwidth}{X} 夜間，有異象向保羅顯現。有一個馬其頓人站著求他說：「請你過來，到馬其頓來幫助我們！」 \end{tabularx} \\ \\ \relax
16:10 & \begin{tabularx}{0.7\textwidth}{X} 保羅既看見這異象，我們就立即設法往馬其頓去，認為神呼召我們傳福音給那裡的人。 \end{tabularx} \\ \\
[1ex]
\hline
\hline
\end{longtable}
$^{1}$各位弟兄姊妹平安.
很久不見了.
代表阿錦書問候大家.
很羨慕阿堅然.
希望大家都是一個.
常住喜樂的心.
我們去到神的面前.
你會青春常駐.
好像姊妹一樣.
我也很希望像她一樣.
我們一起祈禱.
當我們來到你的殿裡的時候.
你的靈與我們同在.
你在當中帶領我們的心.
當我們有什麼不夢理喜悅的.
求主剔除.
以致我們能夠歡喜快樂的.
專一的去仰望你.
得到幫助.
與主你幫助孩子在你面前.
口中的言語.
心中的意念是夢理悅立.
我們這樣祈禱交代.
奉主耶穌基督的名求.
今天我們看這張的時候.
我們想起原來.
有些事情.
你計劃好的.
但不可以如期進行.
但你會發現.
可能你做了錯誤的決定.
或者.
是因為天氣影響你.
甚至因為疫情.
很多事情都改變了.
婚期也改變了.
人們要結婚.
本來要選好日子的.
已經過了一年.
很多事情都在改變.

$^{41}$我們今天看保羅.
他在敘利亞和基烈加.
那地方.
來到他第一次宣導旅程.
差不多在迦太南部的地方.
那一次.
是弄在前面的.
還沒到.
(記者:第一點).
再前一點.
還沒到第一點.
我第一點還沒開始.
第二張圖.
敘利亞和基烈加.
我們看回畫面.
這是第三點.
我們不要理他.
我們繼續.
他去探訪.
以前曾經.
去和他們宣教的地方.
這個城市.
所以這次的目的.
不是要建立教會.
而是他想要培植.
和兼顧過去設立的教會.
他在特比.
如果看回聖經的地圖.
特比,盧斯德.
然後去到以干念.
比西底的安提阿.
原來安提阿有兩個.
一個在敘利亞這邊.
然後去到中間那裡.
是比西底的安提阿.
然後他在那裡過了一段日子.
這個地區.
很接近亞西亞的邊界.
很自然朝著西南去走.
他們沿著巴斯特大道.

$^{81}$經過哥羅西.
再去伊弗所.
原本的目標.
是想在亞西亞.
所以他們朝著那方向.
走了一段路.
但不知道什麼原因.
聖靈禁止他們.
不讓他們在亞西亞廣道.
保羅他從東面.
來的時候.
自然是從特比.
然後去盧斯德.
原來提摩泰是盧斯德人.
保羅是親自帶信主的.
他的媽媽.
是一個信主的猶太人.
外祖母.
也是信主的.
也是信耶穌的.
鄰近地方的人.
信徒很讚賞提摩泰.
所以保羅很留意他.
很希望他可以加入他們的行程.
然後協助宣導工作.
有一個身份很特殊.
爸爸是希利尼人.
不是猶太人.
猶太人來看.
是外邦人.
因為他們不是猶太人.
他們信耶和華神.
但希利尼人.
他們不是.
所以他們未受國禮.
但將來在宣教時宮.
他們一定會遇到猶太人.
所以保羅就和他們受了國禮.
行了國禮之後.
保羅和他的同伴繼續他們的行程.

$^{121}$這次議決.
在耶路撒冷大會.
決定了一些事情.
外邦人.
他們不用受國禮.
但提摩泰.
他是猶太人.
他和外邦人不同.
所以你會看到.
他一定要將會議.
決定的事情.
告訴其他地方的人.
讓他們知道外邦人不用受國禮.
雖然提摩泰.
和外邦人沒有關係.
但他都要.
將開會.
決定的事情.
要和大家宣佈.
去平息猶太人和外邦人.
之間的衝突.
保羅由敘利亞的安提阿.
出發.
但不知什麼原因.
被聖靈禁止.
不讓他們在亞西亞.
去港島.
其實在聖經.
剛剛讀了.
第十節前.
寫史達林傳的.
作者就是路加.
路加.
之前寫了他們.
但到第十節.
他寫了我們.
我們發現.
他去特羅亞時.
已經加入行程.
我們看第一點.

$^{161}$傳遞使命者.
首要條件.
必須要.
信服聖靈的帶領.
第一.
傳遞使命的人.
一定要信服聖靈的帶領.
他們說的地方.
不是羅馬人.
管治的地方.
他們嘗試向北走.
進入比推尼.
就是黑海南部的.
一個部分.
又被耶穌禁止.
耶穌的靈.
也是聖靈.
但不清楚.
耶穌的靈.
如何禁止呢?.
路加沒有交代.
沒有解釋.
很多時候我們遇到難阻時.
是合一的內在見證.
不讓他們去.
或者是.
外在的疾病.
甚至是猶太人反對.
或者是.
法律上.
嚴禁他們.
甚至是透過.
基督徒的先知.
這裡有一個先知.
就是西拉.
所以他們也有.
聖靈去感動他們說話.
可能透過他們說的話.
他們知道.
不可以向前進.

$^{201}$我們也看到.
李文斯津.
這個宣教士.
他想去中國宣教.
但是神.
拒絕他去非洲.
很奇妙的.
他現在很清楚.
看到耶穌的靈禁止他們.
應該是第四張圖.
既然向西南.
向北的路也不通.
唯一正下向西北方向.
就是越過.
美西亞.
去到特羅亞的時候.
有一天晚上.
保羅看到一個異象.
看到一個馬其頓人.
好像跟他說.
你快點過來.
馬其頓幫助我們.
好像是請求的意思.
好像在跟你打招呼.
招手.
叫你快點過來我們這裡幫忙.
古時候用夢.
因為夢就是.
神指引他們.
譬如.
如果你有記性.
記得.
在史太人傳第九章.
曾經.
有一個人.
叫做亞娜尼亞.
他看到一個異象.
異象叫他去直街.
去猶大的家.
找一個大數人叫做蘇羅.

$^{241}$叫他去找.
而蘇羅人正在祈禱.
他又在異象裡面.
看到一個人.
叫做亞娜尼亞.
跟他按手.
令他看到.
很奇妙.
所以保羅和同伴.
馬上覺得.
原來.
夢是神的呼聲.
於是馬上.
他們就.
向馬其頓進發.
他們覺得.
是神的呼聲.
於是信服聖靈的大靈.
既然不可以去亞西亞.
現在就在馬其頓進發.
來到一個地方.
這個地方就是.
菲臘比.
菲臘比這個地方.
頂章塘.
似乎那裡是沒有會堂.
因為那裡有經文.
會堂裡面.
發現.
以前猶太人的會堂.
十個男人就可以成立一個會堂.
但不見了.
於是他們就去到.
甘寨提斯河.
那裡看到.
有一個祈禱的地方.
那祈禱的地方有很多婦女.
於是保羅有機會.
跟他們傳播音.
在頂章塘那裡就說到.

$^{281}$他們由特羅雅開船.
直駛到.
莫特拉.
去到尼亞波利.
然後那裡.
去到菲臘比.
大家很熟悉的菲臘比.
菲臘比其實是馬其頓這個地方.
一個重要的城市.
是一個駐防城.
羅馬的駐防城.
他們在那裡住了幾天.
在那裡祈禱的地方.
保羅有機會.
跟他們傳播音.
在那裡你看到有一個.
很有影響力的商人.
他叫做呂帝雅.
他是賣紫色布的.
是推亞推拉城的人.
他是一個敬畏神的人.
他的心已經準備好了.
他很留心去聽.
聽保羅在那裡講道.
所以他明白到.
保羅在講耶穌的真理.
他聽了.
他就相信了.
保羅不是對他一個人講的.
當時有很多婦女.
但他特別留心.
他聽他很渴望.
主的真理.
所以主就開導他的心.
他相信了.
頂子有試過.
當聖靈光照你的時候.
你的心就會受感動.
我們願意跟從主的吩咐去做.
去行出真理.

$^{321}$所以你看到.
他會順從聖靈的帶領.
會順服這本聖經的真理.
去實行.
去跟從主的腳步.
所以你看到.
一步步你的靈情裡面.
你會願意奉獻自己的身心.
去給主.
呂帝雅.
她和她一家人.
都洗禮了.
洗禮之後.
她還會接待她的同伴.
你看她影響她的家人.
她的家人也相信了主.
還洗禮.
你看她信主就洗禮.
我們的過程可能要等很久.
但現在你看到.
她不單信主.
還洗禮.
洗完禮之後.
她還有行動.
有愛心的行動.
很明顯的她用愛心去款待神的僕人.
去接待他們.
而保羅.
他也使用了信息.
讓他去說得明白.
可以成就神的心意.
一路來看,很順利的.
他當了目的敵.
有機會去講信息.
很順利的.
來到馬其頓的地方.
在腓立比.
你會看到.
可能這張完成的時候.
這班人.

$^{361}$在腓立比.
很自然的.
可能這張完成的時候.
之後才建立了.
這個腓立比教會.
下一張圖.
就是看到十六節.
十六節就是.
他再去祈禱的地方.
發現一些很奇怪的事.
看到一個.
被占卜的零附身的使女.
迎面走來.
她使用法術.
使主人發大財.
她還跟著保羅.
和他一起.
這個行程的人.
所以這裡就說.
他跟隨保羅和我們.
哭著說.
這些人是至高神的僕人.
對你們傳講救人的道路.
這個女人一年幾日都跟著他們.
每天都這樣叫.
你想想占卜的零.
她不是姓零.
她和姓零是敵對的.
她怎會贊成.
神的工作呢.
十六節保羅這次.
去到腓立比.
遇到一個屬零的爭戰.
在祈禱的路上.
遇到一個.
被占卜的零附身的.
一個女子.
還迎面來.
她透過占卜.
看著主人賺很多錢.

$^{401}$因為主人當她是一個工具.
去騙錢.
奇怪的是.
她跟著保羅.
經常這樣說.
這些人是至高神的僕人.
對你們傳講救人的道路.
一年幾日.
都這樣呼喊著.
你當然很高興.
有人意外的跟你宣傳.
起初你不覺得危險.
但要小心.
慢慢的你以為.
他跟你是同一陣線.
以為是自己人.
危險了.
你沒有防備.
但保羅在這裡說.
他心中厭煩.
後來看出他被邪靈支配.
他奉主耶穌的名.
趕他出去.
叫他從屎裡的身上出來.
所以靈就出來了.
奉耶穌基督的名.
吩咐你從他身上出來.
靈出來之後.
怎樣呢?.
主人就沒有希望了.
因為本來是靠他發達的.
現在靈出來了.
就沒法發財了.
於是.
這兩個人.
保羅西就被他抓去見官.
第二點.
講到一個傳遞使命的人.
他必須.
要經歷考驗.

$^{441}$在考驗中.
他要.
忍耐等候.
第八章.
經文就給他看了.
去到官長的面前.
這個人怎樣說呢?.
這個主人.
他說的話.
他的控告.
他的原因是什麼呢?.
他說這些人是騷擾我們城市的人.
他們是有意的.
他們是猶太人.
竟然傳一些盧曼人.
不可以接受.
不可以遵守規矩.
說得姑勉堂皇.
去控訴他的原因.
其實是不是這樣呢?.
他引起一些群眾.
在當中.
人們本身都對基督徒不滿意.
對猶太人.
現在火上加油.
原來信主.
未必是一帆風順.
會坐牢.
被抓.
進去還未問.
就打他們.
用棍打他們.
然後放在監牢裡.
還叫人看守著他們.
嚴謹地看守他們.
控告他們.
說他們騷擾城市.
說一些不可接受的規矩.
這樣控告的時候.
引起大家.

$^{481}$很混亂.
引起暴動.
還引進一些.
說他們傳異教.
很嚴重的事.
一般來說.
盧曼人.
他們不可以信奉一些.
他們國家未批准的異教.
但習慣上.
其實可以這樣做.
只要他們的異教沒有違反.
羅馬生活的法律.
或者規條.
只要他不牽涉到任何政治.
或者社會的罪行.
其實沒有人理會.
但二十一節.
主人.
控告的原因是包裝得很漂亮.
完全看不到他憤怒的原因.
因為錢.
他不能賺錢.
就是這樣仇恨他.
所以要抓他們見官.
你看到.
不信的人.
更加.
對你不滿的時候.
很惡劣.
你看到.
他們被人對待.
多麼坎坷.
被人脫了衣服.
被人打.
在心路.
最近過去二百年.
有首位.
來中國的宣教士.
一個殉道者.

$^{521}$死在中國土地.
這個人就是.
劉理華宣教士.
他是屬於.
美國長老會的宣教士.
他很有學問.
他是美國十大名校.
普林斯頓大學畢業.
是很年輕的一個學者.
大概一百三十年.
大約二十四歲.
那時候.
有馬利信聖經出版.
他自己.
很想翻譯聖經.
有一次他有事情.
要坐船由上海去寧波.
中間經過.
杭州灣的時候.
他看到海盜.
逼近他們.
他們撐著艇.
有八個海盜.
接近他們的船.
那些人拿著矛槍.
試圖打開他們的行李箱.
看看有什麼值錢的東西.
這個宣教士.
他很從容的.
他看到他.
他拿鑰匙給他.
免得他叫.
我給你鑰匙.
你那本聖經不要動.
其他的任你拿.
他那本聖經很值錢的.
因為他要翻譯聖經.
有希伯來文.
有希臘文還有當時的英文版.
巴克斯特版本.

$^{561}$因為他要做翻譯的工作.
那這八個人.
準備離開的時候.
有人覺得不妥.
因為覺得有外國人在.
將來.
他們如果去領事館.
告發我們的時候會有麻煩.
於是他們八個人就商量.
將他扔到海裡.
用矛槍.
黏著他.
這個人.
臨死之前他不斷掙扎.
要將這本聖經扔回船裡.
為什麼?.
因為他希望聖經可以流傳下去.
他自己沒命.
但他也要扔這本聖經.
雖然他已經沒力氣了.
但他也要扔這本聖經上船.
希望可以流傳聖經.
你看到.
他就是這樣的.
被弄死了.
那時候.
他才28歲.
你看到他.
宣教的生涯.
只是四五年的時間.
很短.
但他覺得.
聖經可以流傳是很重要的.
你看到宣教時.
是用血汗.
這樣獻身給神.
是用血汗.
用他們的生命去見證.
我們有沒有公敬的.
去讀這本聖經呢?.

$^{601}$如果我們不讀聖經.
我們是不是對不起他們呢?.
有很多人花了很多心機.
去翻譯聖經.
但我們今天只是去看.
我們識字的人.
我們只是看,很輕鬆的.
但我們都不看的時候.
是不是對不起他們呢?.
看完聖經.
也要花很多時間.
但他翻譯.
是要多少年的時間.
但今天如果我們不去.
認真地看聖經.
我們真的對不起他們.
我們真的要在神面前.
我們要悔改.
好好地去公敬.
去讀神的話語.
神的話語就是我們腳踝的燈路上的光.
去指引我們.
我們可以見到一個服事主的人.
他走的路.
不是那麼平坦.
危險是會有性命的.
第三點.
一個傳道.
要傳遞使命的人.
他的堅持是怎樣的呢?.
他為了要忠於主的託付.
所以在第十個圖.
裡面.
就是講到.
25節,在半夜的時候.
保羅和西拉.
他們是祈禱.
他們唱詩讚美神.
囚犯都要側耳聽.
留意一下.

$^{641}$如果有人在講話的時候.
側耳聽的時候.
是什麼態度呢?.
就很想聽聽他們在講什麼.
側耳去看他們.
是不是有些問題呢?.
被人拉進來.
被人打得那麼慘.
還會唱詩,還會讚美神.
是不是很奇怪.
他們被抓去坐牢.
沒有自由了.
不信主的時候是好好的.
信了主,反而.
要面對這樣的遭遇.
你也要考慮一下.
我要不要回來教會呢?.
我要不要跟你們一起信主呢?.
但是保羅和西拉.
他們受到那麼大的困難.
被人打.
都還沒問他們什麼罪.
才把他們打了.
把他們放在監獄裡.
但是他們呢.
沒有詛咒.
他們反而是唱詩讚美神.
所以難怪這些囚犯.
要側耳去聽.
為什麼他們那麼開心呢?.
半夜.
突然間獄警醒了.
為什麼醒了呢?.
因為突然間地大震動.
地大震動.
地震嘛.
大家都害怕.
地大震動.
連門都開了.
囚犯的鎖都開了.

$^{681}$不是你跟他開的.
是震的時候開了.
但是旅行的人.
沒有死,沒有受傷.
保羅有機會帶囚犯去信主.
這是很奇妙的經歷.
但是獄警.
他醒過來的時候.
看到這樣的情形.
門又開了.
人家的鍊子又鎖了.
他很慌了.
為什麼這些囚犯都走了呢?.
我要背負很大的責任.
當然要被處死.
我要負責的.
他很害怕.
他拿起刀來就要.
自殺.
但是保羅叫著他.
叫他不要.
千萬不要死.
我死你也不要死.
他很害怕.
太害怕了.
但是保羅人比他還害怕.
因為他不死.
他也要受刑罰.
如果他死了,刑罰更重.
所以當然叫他不要死.
這個人.
很有趣.
忽然間.
跪在保羅人西欖面前.
他說:二位先生.
我必須做什麼才可以得救.
留意紅字.
二位先生.
我必須做什麼才可以得救.
有沒有看到.

$^{721}$這個獄長.
他知道他需要什麼.
弟兄姊妹,你知不知道你需要什麼.
他知道他需要什麼.
他是一個獄長.
這樣的行動是不容易的.
跪下,你是囚犯,跪在你面前.
不容易的行動.
有沒有看到.
但是保羅人和西欖直接告訴他.
當信主耶穌,你和你一家都必得救.
大家很熟悉的金句.
當信主耶穌,你和你一家都必得救.
很直接的.
這個福音.
他這樣告訴他.
教他怎麼做.
你要信耶穌了.
而這個獄長.
對保羅西欖.
有感恩的心.
所以他立刻的行動.
和他們.
擦傷口.
因為獄長和他家人.
都信了主.
你信主也不容易.
要令你全家都信主.
是更加困難.
你會看到.
很奇妙的.
在這個時候,獄長.
和他家人都信了主.
還洗禮.
全家的人都洗禮.
獄長還帶他們.
去他家.
保羅和西欖是.
囚犯,有沒有想到.
他是.

$^{761}$很重的囚犯.
在地牢.
獄長帶他們.
去他們的地方.
離開他們的房間.
離開房間是有罪的.
可能死刑.
你怎麼可以私下.
帶犯人離開監房.
這麼大膽.
而且他們在地牢.
你怎麼可以私自帶走他們.
可見這個.
獄長信心很大.
而他家人信了主.
可能他們就是菲律賓教會.
最早期的信徒.
他們擺上飯.
去接待他們,請他們吃飯.
他和全家的人.
信了主.
都滿心喜樂.
你看到.
信了主的時候.
很自然,心中有.
平安,有喜樂.
所以,到天亮的時候.
官長就打發.
差役來,說要放他們.
放他們好不好?.
要放保羅和西拉.
如果你是保羅,開心嗎?.
但很奇怪.
這個獄長當然祝福他們.
平平安安的.
回去,祝福他們.
但是.
37節.
下一張圖.
就看到保羅怎麼說.

$^{801}$他說我們是羅馬人.
羅馬人.
有保障的.
並沒有定罪.
他們竟然在公眾.
打我們,將我們放在監獄裡.
現在要私下.
趕我們走?.
暗地裡放我們走.
不就開不了了之.
你又不用認罪.
不用跟我們道歉.
不用道歉.
我不可以接受.
很奇怪.
我們放他們,他們說不可以接受.
我是羅馬人.
他說不行.
叫他們自己來.
領我們出去.
很大膽.
差役就照樣把話說給官長聽.
官長一聽到.
他說羅馬人就害怕.
聽到羅馬人就害怕.
羅馬人有保障的.
他都沒有.
問他什麼,就先打了他.
就把他放在監獄裡.
等於虐待.
在工作上.
其實羅馬人在法律上.
是有他們的權利的.
是得到保障的.
除非他本人不同意.
他本人不同意的時候.
你不可以.
用普通地方的法律.
去限制他.
甚至審訊他都不行.

$^{841}$他可以上訴.
但不可以虐待他.
不可以殺他,不可以毒打他.
羅馬人有很多權利.
他上訴的那段時間.
你都不可以捆綁他.
所以保羅這樣做.
他就是想討回公道.
他希望藉這個機會.
可以見到官長.
然後當面跟他解釋清楚.
另外就是.
將來有.
比特幣教會存在.
他都要保障.
這班人不會受到官方.
歧視他們.
還有避免.
將來他都要回來探望他們.
不想有不良的紀錄.
所以.
他這樣做.
就是他的原因.
希望他將來.
可以再回來探訪比特幣教會.
當這個.
羅馬人.
這個明清.
官長知道之後.
他見到.
令到他們.
很不開心.
很害怕,要親自來.
跟他們道歉.
低聲下氣.
叫他們走.
為什麼要叫他們走?.
很怕那些人.
會搞事.
那些人就因為你.

$^{881}$差不多要暴動.
你快點走,怎樣道歉.
我都會跟你殺個lum .
他希望你快點走.
所以保羅都要做完.
他的事情.
他出了監之後.
他都去探訪那些弟兄姊妹.
勸慰他們一番.
離開那個地方.
各位弟兄姊妹.
今天我們都領受了.
主給我們的使命.
你有沒有信服.
使你帶領.
肯聽從主的吩咐.
去行.
當你遇到考驗的時候.
你還記不記得.
主與你同行.
你可不可以忍耐等候.
你可不可以堅持到底.
去忠於主的託付.
我們一起祈禱.
親愛的主耶穌.
你讓我們今天可以成為你的兒女.
我們知道走主的道路.
在這個彎曲勃謬的世代裡.
我們都要走.
走主的道路.
在這個彎曲勃謬的世代裡.
是不容易的.
只求你幫助我們.
可以勝過試探.
幫助我們真的能夠信服聖靈的帶領.
能夠堅持到底.
直至見你面.
你都可以稱呼我們.
有良善.
有忠心的僕人.

$^{921}$求主引領我們的心.
去保守我們.
能夠忠心為主.
能夠堅持到底.
我們這樣祈禱.
教導是奉主耶穌基督的名教.
阿們.
\newpage



\section{詩篇 150:1-6}
\label{sec:65pDMB41cZg}
\textbf{2021.08.29《當來讚美神》 詩篇 150\_1-6 (和修版) 張美薇牧師}
\newline
\newline
連結: \href{https://youtube.com/watch?v=65pDMB41cZg}{\texttt{https://youtube.com/watch?v=65pDMB41cZg}} ~~~~ 語音日期: 2021-09-01
\newline
\newline
\hyperref[sec:h6Nm0N8GuaU]{\small{< < < PREV SERMON < < <}}
~
\hyperref[sec:index_chronic]{\small{[返順時目]}}
~
\hyperref[sec:index_scriptual]{\small{[返順卷目]}}
~
\hyperref[sec:xUmmDjoz5Ds]{\small{> > > NEXT SERMON > > >}}
\newline
\newline
詩篇 150:1-6
\newline
\begin{longtable}{cl}
\hline
\hline
章節 & 經文 (和合本修訂版)\\
\hline
150:1 & \begin{tabularx}{0.7\textwidth}{X} 哈利路亞！你們要在神的聖所讚美他！在他顯能力的穹蒼讚美他！ \end{tabularx} \\ \\ \relax
150:2 & \begin{tabularx}{0.7\textwidth}{X} 要因他大能的作為讚美他，因他極其偉大讚美他！ \end{tabularx} \\ \\ \relax
150:3 & \begin{tabularx}{0.7\textwidth}{X} 要用角聲讚美他，鼓瑟彈琴讚美他！ \end{tabularx} \\ \\ \relax
150:4 & \begin{tabularx}{0.7\textwidth}{X} 擊鼓跳舞讚美他！用絲弦的樂器和簫的聲音讚美他！ \end{tabularx} \\ \\ \relax
150:5 & \begin{tabularx}{0.7\textwidth}{X} 用大響的鈸讚美他！用高聲的鈸讚美他！ \end{tabularx} \\ \\ \relax
150:6 & \begin{tabularx}{0.7\textwidth}{X} 凡有生命的都要讚美耶和華！哈利路亞！ \end{tabularx} \\ \\
[1ex]
\hline
\hline
\end{longtable}
$^{1}$主席.
高聲地讚美您.
我們更加可以一起來聆聽您的說話.
求主開我們的心,開我們的眼.
讓我們的心柔軟.
能夠明白您的話語的能力.
也能聽懂,聽明白.
也能幫助我們,不僅聽懂,聽明白.
更加幫助我們,讓我們有能夠有行道的手和腳.
去將您的話走出來.
我們願意用我們的生命來榮耀主.
聽我們的祈禱,奉耶穌基督的名求.
阿們.
今天想和大家分享詩篇第150篇.
「當來讚美神」.
讚美.
其實每一個民族都有不同的民族性.
大家同不同意嗎?.
我自己小時候.
可能這樣就出賣了我的年紀.
我小時候很少得到大人的讚.
可能年輕一點的媽媽,爸爸經常讚你.
年紀大的就不是.
通常父母都是罵的多.
讚的少.
做得好像是怎樣?.
是英訓的,大家都知道.
我聽說有人因為考不到第一,考第二.
某一科的成績只有98分.
被媽媽追著問.
你坐在哪裡?.
為甚麼那兩分被人扣了?.
但其實在美國文化是不一樣的.
我以前在美國讀書的.
所以知道.
原來在美國文化,父母通常都會讚子女.
經常說.
Good,做得好,做得好,做得好.
Well done.
在中國文化裡面是不會的.

$^{41}$甚至我們很多時候對外人會說.
這是個蠢仔,還有乞丐仔.
通常在孩子的自尊就得不到見.
所以我們是要學的.
當我們小時候沒有被人讚的時候.
我們都不是很懂得讚人.
所以詩篇第150篇裡面是說讚美.
我們可以看一看.
我將詩篇第150篇六節都抄了出來.
和和修本第一句是甚麼?.
和修本第一句是Hallelujah.
其實詩篇在和合本裡面.
就將它譯了出來.
你們要讚美耶和華.
其實這首詩篇很短的.
如果是計和合本來說.
翻譯完是大概133個字.
還包括標點符號.
133個字.
你知不知道說了多少次.
讚美他,讚美神,讚美耶和華?.
總共說了13次.
13次即是多少個字?.
即是如果你數一下.
有45個字是和讚美他,讚美耶和華,讚美神.
133個字裡面有多少?.
有45個字,多不多?.
其實是很多很多.
所以我們去學讚美神.
這一首是很好學的一首詩篇.
有人就說.
其實雖然我們小時候.
我小時候,你們小時候未必是一樣.
很少被人讚.
但我們人其實是很喜歡別人讚的.
大家想想小孩子很喜歡別人讚他甚麼?.
讚他聰明.
讚他乖.
還有讚他靚.
現在我兒子還喜歡別人讚他甚麼?.

$^{81}$讚他Q.
即是他Q就可以.
他說我不用Lag.
他說我Q就可以.
那些是很重要.
政治人物其實都喜歡別人讚的.
讚他甚麼?.
是人民的大救星.
是好的領袖.
大家最近派那兩千元.
花光了沒有?.
都可能花光了差不多.
大家會不會發覺.
派出來的時候.
陳茂波那一朝.
每間間電台都在拍著他.
四處去請人吃蝦餃.
然後就算我們的特首.
都請人吃蘋果.
是怎樣?.
他想人家去讚他.
讚政府.
但我們又看看.
其實.
如果你看報紙.
是彈多過砸的.
是嗎?.
我有時都不是很明白的.
因為派消貴券.
我不知道你們有沒有這樣去拿.
他們說如果你在OK店那裡拿.
就給你二十元.
每一張是五元.
於是怎樣?.
排長龍.
於是網上的人說.
政府真是擾民.
這樣做.
令人排長龍.
排到這麼長.

$^{121}$其實關政府什麼事?.
其實是OK店的做法.
但人們就經常去說.
很多時候我們就去彈.
我們這樣看下去.
其實世上最值得.
我們最需要讚的是誰?.
是神.
但我們就不是很懂得去讚神.
所以在詩篇第150篇裡面.
很好.
他教我們要讚美神.
他第一句和最後一句.
都是說.
你們要讚美耶和華.
弟兄姊妹.
有時我們可以開口讚美神.
但假如你現在的環境不是太暢順.
又或者你遇到一些困難.
你容易讚美神嗎?.
不是這麼容易.
所以有個詩人就說.
假如神的旨意.
是要我處於黑暗當中.
我要怎樣?.
我讚美神.
假如神的旨意.
是要我居於光明當中.
我都讚美神.
你如果賜我安慰.
我讚美神.
如果你加我痛苦.
我怎樣?.
我同樣讚美神.
所以這是我們去讚美神的一個很重要的基本做法.
就是我們無論是什麼情況.
我們都要讚美神.
在世間.
我們如果真的不願意讚美神.
或者不懂得讚美神的時候.

$^{161}$到天家就論盡.
因為到天家.
所有天使,所有聖徒都只會做一件事.
就是讚美神.
所以我們要學.
所以今天我們會用一個查經的方式.
就是六合.
大家有沒有聽過?.
六合方法來學我們怎樣讚美神.
第一個.
原來講完了.
我們讀第二節.
要因祂大能的作為讚美祂.
按著祂極美的大德讚美祂.
意思就是說.
為什麼要讚美神?.
我們很喜歡問為什麼.
我們知道任何人要獲得讚美.
決不應該單憑他的身份.
或者他的地位.
就得人讚.
這是假的.
譬如說.
因為我是張牧師.
你們要讚我.
沒理由的.
一定不是這樣.
所以也不會說.
因為林鄭是特首.
因為她是特首.
我們就一定讚她.
不是的.
我們一定要有.
她一定要有好的作為.
好的表現.
我們先政是會讚美她的.
所以我們第一樣.
剛才那節經文我們看到.
我們要因神的大能來讚美祂.
剛才唱的詩歌是很對的.

$^{201}$我給大家選詩歌.
詩歌會說什麼?.
原來神的大能是很厲害的.
因為神是創造天地的神.
天地是神造的.
地和地中的一切都是造的.
而天地顯出神的大能.
祂也做人類.
創造出你.
創造了什麼?.
創造了我.
所以神是很厲害的.
我們又看舊約聖經.
我們回來教會讀經.
看舊約聖經.
神又選了亞伯拉罕.
在亞伯拉罕的生命裡面.
做了很多事情.
神又厲害到帶領以色列人出埃及.
又帶領他們入迦南地.
當以色列人被擄當中的時候.
神也保護他們.
更厲害的是.
以色列人亡國兩千年.
神卻在第二次世界大戰之後.
讓他們復國.
所以神的大能是很厲害的.
但可能你會說.
你說這些都是歷史.
好像和我很遠很遠.
有沒有看到.
其實神在洪恩當中.
也有很多大的作為.
大家在洪恩有一段日子了.
大家有沒有去聽一下洪恩的歷史.
我想我們當中如果要問.
輝哥就很熟悉.
可以和我們去講.
神其實在很多地方當中.
讓錦宣家選擇了洪水橋直塘.

$^{241}$更加讓業主願意租這間地方給我們.
其實也不簡單.
因為他可以做很多不同的事情.
神又願意讓錦宣家.
讓我們有一個這麼大的地方.
有沒有想過人家直塘會有三層樓.
地面那層就小一點.
第二樓和三樓都很大.
很足夠的地方.
還有神又讓錦宣家願意.
直塘和礦塘一起做.
我們第一次直塘礦塘的時候.
做得很辛苦.
那時候大概是2001年.
我們那時候建了錦宣的座塘.
我們又直了宣德出去.
2002年那時候很辛苦.
所以當我們後來說要直塘的時候.
我們本來就說我們要分開兩件事.
我們直塘和礦塘要分開.
結果我們卻在2011年的時候.
又直塘然後又礦塘.
這個是很特別.
然後神又讓一群弟兄姊妹.
願意跟著李牧師一起出來直塘.
不簡單的.
我如果看回.
我問寶湖.
寶湖就給了數字我.
是有35位弟兄姊妹.
一起願意出來直塘.
我想在座當中.
都有這些直塘勇士在我們當中.
回想起來都很興奮.
回想起當時.
其實還有.
你們有沒有看到.
這張圖.
其實是我問寶湖拿的.
就是將逐年在洪恩,洗禮.

$^{281}$和轉會的弟兄姊妹的數字列出來.
你看看不簡單的.
在過去九年多.
總共有107人在這裡信主.
有35個是直塘勇士.
也有30人轉會.
神的能力是不是很厲害.
當我們不是這樣數的時候.
我們都不會發覺.
原來神是很厲害.
我想你們都是在當中的.
一是直塘勇士.
一是在這裡信主.
一是轉會.
又或者現在才出來.
現在又有魯牧師和曾姑娘.
繼續的接力下去.
神又感動各行委,領袖,弟兄姊妹.
一起來經營.
大家在這裡坐得悶悶的.
很好.
其實神的大能是很不簡單.
我們更加可以看到.
我們還有合作的機構.
剛剛看到報告.
有兒童的合作機構.
也有每月的派福袋,行動等等.
我們看到洪恩堂的未來.
神是用你們.
神是大能地用你們.
在這個社區裡面.
去改變這裡的社區.
神會使用洪恩.
改變這裡的社區.
改變洪水橋.
改變元朗.
是否感受到神的大能.
其實除了神的大能之外.
神也有良善.
在和學本裡面.

$^{321}$是大德.
而在和修本.
因德極偉大.
去讚美神的大德.
甚至於十字架.
成就了耶穌在普世的救贖工作.
不單救贖了以色列人.
也延伸到萬國,萬民.
神對他所做的人的慈愛.
神降雨給好人.
但壞人呢?.
神一樣降雨給壞人.
白天照好人的.
白天同樣照壞人.
神的醫治.
神醫好人.
神也同樣醫不好的人.
神有公認.
神有看顧.
其實我們相信.
在洪恩的匯種當中.
也有很多很多改變生命的見證.
讓你們去發掘.
讓你們去宣講.
我自己在大概二十年前.
即是十九年前.
我們出席了宣德堂.
宣德堂有一個很特別的例子.
他們每五年.
慶祝五週年,十週年,十五週年.
他們就會出一本書.
將匯種裡面的見證.
就會講出來.
在匯種當中.
做了很多很多.
改變生命的見證.
我記得我自己初信的時間.
也是的.
我們那時一班教會的弟兄姊妹.
我們教會很小.

$^{361}$我們教會只有四,五十人.
但是我們的團契.
每個星期我們都回去分享.
我們當時的團長.
那一年就定了一個題目叫我們分享.
每次我們都要分享兩件事.
第一件事就是.
我們的靈修心得.
靈修心得都比較容易.
如果有.
即是說.
即是說一些壞事給你們聽.
如果我那個星期.
做得不是那麼好.
靈修.
我星期六早上回去分享.
我星期五一定會早點去找一段經文.
然後隨便都會做到一些東西出來分享.
但是第二個題目就更加難.
就是說.
你要去看看.
神在你附近.
在你的身上.
或者在你附近.
有些什麼大的作為.
祂做了什麼好事.
最初那時真的很難.
不習慣.
很難去看.
但是慢慢當我習慣了去看到的時候.
很多時候星期一就已經等.
怎麼那麼久還沒有到星期六.
我可以回去分享.
這些神的作為.
很簡單的.
我記得有個弟兄.
他是長老.
很年輕.
我想他那時剛出來工作.
他就說.

$^{401}$他有一次.
美國那些叫Garage Sale.
不知道大家知不知道.
就是人家車庫賣東西.
他說.
他那天就分享.
我買了一個一元的煲蓋給教會.
所以很開心.
我為教會省了一筆錢.
然後他分享的時候突然說.
咦.
為什麼我買一元的煲蓋給教會.
但是我買給自己.
就是一個新煲.
他說.
咦.
我可能是有一點點的問題.
我們就看到他在分享裡面.
他看到自己的需要.
發現原來神是很良善的.
去提醒他.
所以弟兄姊妹.
我們是要去學習.
去學習.
去讚美神的大能.
讚美神的良善.
有沒有東西要讚美?.
有.
所以其實是很多很多東西要讚美.
另外.
接著的問題就是.
我們在哪裡讚美呢?.
我們在第一節裡面就說了.
他說你們要讚美耶和華.
在他的聖所讚美他.
在他顯能力的窮倉讚美他.
聖所.
聖所是什麼呢?.
聖所是人造的地方.
是用來什麼的?.

$^{441}$是用來敬拜神.
所以.
這一首詩篇的對象.
是神的子民.
而神的子民聚集在一起的時候.
他們就建築了聖所.
來敬拜神.
來讚美神.
所以當我們會發覺.
我們在神的聖所.
我們可以看到世界各國.
他們未必是認識神的.
他們都會在他們建立最美最美的地方.
來敬拜神.
他們的神.
這個是什麼人?.
什麼的敬拜?.
這個是回教旨,是嗎?.
這個是印度教旨.
而這個是埃及的廟宇.
而這個就是一些比較落後的民族.
可能是原住民的圖騰.
他們都是用這些東西.
來敬拜他們所認為的神.
中國人一樣,是嗎?.
我們天壇.
都是用來敬拜中國人所認為的神.
但當神的子民認識了神之後.
我們也會在我們今天的世界.
我們會有不同的敬拜神的地方.
可能是一些很宏偉.
很大很大的禮拜堂,是嗎?.
又或者是一些很小的教堂.
尖頂.
大家看到這個是尖頂的教堂.
面無意外,可能是一些下雪的地方.
所以要這麼斜的.
讓雪可以下來.
又或者.
其實這個是一個在泰國的.

$^{481}$很簡陋的敬拜神的地方.
簡陋,或者好像剛才第一張這麼美.
不要緊的.
重點是什麼?.
我們在那裡是敬拜神.
又或者今天大家可以在這裡.
大家到一起是敬拜神.
所以我們知道教會是信徒組成.
這個才正重要.
有人說禮拜堂不重要的,是嗎?.
禮拜堂不是最重要.
但禮拜堂也不是不重要,是嗎?.
一樣是很重要.
我記得對敬拜神的地方.
我們一定要重視.
所以大家會看看.
地方都幾乾淨,是嗎?.
幾整齊.
我記得藤牧師跟我們說.
藤近輝牧師.
他說他剛剛出來教會事奉的時候.
他說有一次星期六早上.
他進到教會的禮拜堂.
想去拿些東西.
他就發覺有幾個老人家.
有些講上海話的婆婆.
他們本來家裡是幾有錢的人.
但他們都自己回來打掃好教堂.
預備星期日的敬拜.
他說他看到他們的服侍.
他很敬畏的心.
他很靜靜地退出去.
因為他不想打擾到他們.
弟兄姊妹.
我們對敬拜神的地方.
一定要尊敬,尊崇.
我們今天對網上的弟兄姊妹說.
你有機會的話.
回來崇拜.
在教堂裡面.

$^{521}$在教會裡面去敬拜神.
跟弟兄姊妹一起崇拜.
去敬拜神.
我們知道在末後的日子.
黃金街,碧玉城,玻璃海.
在三層天上高的正所裡面讚美.
而且是面對面的.
是讚美神.
我們知道第一個去讚美神的地方在哪裡?.
在禮拜堂.
我們今天說.
第二個他說可以在窮倉.
在他顯能力的窮倉裡面.
讚美神.
神是造整個世界.
整個宇宙.
彰顯神的榮耀的地方.
我們知道.
我們教會裡面是讚美神.
但在整個宇宙裡面.
任何一個地方.
我們都可以讚美神.
我們亦都需要.
在我們去旅行的時候.
我們去到不同的地方的時候.
我們同樣地向神獻上讚美.
我們在天地當中讚美神.
因為天地可以彰顯神的榮美.
所以當我們見到.
早幾天天藍得很漂亮的時候.
我們會怎樣?.
我們會讚美神.
但當天陰行雷閃電的時候.
我們都看到它彰顯了神的榮美.
彰顯了神的大能.
我們同樣要讚美神.
就算在一個風暴當中的時候.
我們同樣要讚美神.
又或者在雨後彩虹的時候.
我們都一樣要讚美神.

$^{561}$所以我們會看到.
大如天空的星體.
小如電子,原子等粒子.
其實規律都是一樣的.
叫人驚訝神的奇妙.
在人看的細物.
亦都不是科學所能完全明白的.
連人的眼眉毛,腳趾尾.
都有一定的用處.
所以窮倉不斷地讚美神.
怪不得勝經都有說.
如果人不開口讚美神的時候.
石頭都會開口來讚美神.
所以整個宇宙是一個很大的聖所.
所以人應該在神顯能力的窮倉.
去讚美祂.
所以無論天是晴或暴風雨.
我們都讚美神.
深挖.
就算我們在監獄裡面.
我們都要學習讚美神.
記不記得.
當我們讀到保羅書信的時候.
保羅在監獄裡面.
在困境當中的時候.
他都一樣的讚美神.
直至到他在菲律賓監獄裡面.
唱歌讚美神的時候.
地大震動.
然後監門大開.
弟兄姊妹.
你現在在順境當中.
還是逆境當中.
我們都要學習讚美神.
我們要讚美神.
如果你處於逆境當中.
或者人生低潮當中.
我們要開口的讚美神.
因為神是不改變的.
無論何境況.

$^{601}$祂仍然愛你.
仍然愛你.
祂不會撇下你.
祂會不斷的幫助你.
所以我們在不同的境況.
我們都要讚美神.
你看多療癒化香.
但不是經常都是這樣.
所以我們要學習讚美神.
有人就會問.
我知道我們為甚麼要讚美神.
我們知道我們在哪裡讚美神.
然後我們要看怎樣去讚美神.
用甚麼來讚美神呢?.
這裡有三節經文.
要用覺聲讚美他.
鼓瑟彈琴讚美他.
擊鼓跳舞讚美他.
用詩源的樂器.
和簫的聲音讚美他.
用大響的筆讚美他.
用高聲的筆讚美他.
其實有幾方面的.
第一方面就是用各種樂器.
為甚麼?.
因為詩源說.
其實我們每一個都要讚美神.
所以我們有不同的樂器.
我們在這裡都看到.
我們這裡有鼓.
我們這裡有琴.
我們有時還有結他.
最簡單是甚麼?.
是用我們的嘴巴.
我們有各種樂器.
有各種特殊的美妙的聲音.
當我們同時一起演奏的時候.
合起來就是一首多姿多彩的協奏曲.
所以亦都包括我們不同種族的人.
我們知道我們這裡.

$^{641}$剛才都有說.
我們有些東西是用廣東話去說的.
有些東西是用普通話去唱的.
我們有時又可以用英文去唱.
等等等等.
所以我們是不同種族.
各樣的樂器.
還有更加是說.
各種的姿勢.
有些是用口吹的.
有些是用手撥的.
有些是用擊打的.
即是用敲的.
有些是用手彈的.
有些是用全身來跳舞的.
有些是單單用口的.
還有大小不同的聲音.
有些聲調高的.
有些聲調低的.
有些聲調是大的.
我們是演出一個讚美的人生.
其實雄恩也有很特別的一個詩班.
有兩個詩班.
一個叫做傳統的詩班.
我們唱傳統的詩歌.
另一個是甚麼.
福音樂曲.
我們唱樂曲來讚美神.
用中國人文化的音樂.
同樣去敬拜神創造天地的獨一真神.
用我們所能去讚美神.
這比我們一生的遭遇.
其實是變化無常的.
但因為神的恩佑.
這些不同的遭遇.
都成就了不同美妙的音樂.
有時我們德性好像吹角一樣.
打仗的時候.
是很雄壯的吹角.
有時我們很苦惱.

$^{681}$好像是鼓失一樣.
有時我們快樂好像彈琴.
有時我們爭戰如擊鼓.
有時我們軟弱卑微.
好像撥弄思言的樂器.
有時我們也是卑虎,卑傷.
好像吹簫一樣.
有時我們好像興奮.
用大響的笛等等.
這些人生高高低低的經驗.
經歷放在神的面前.
就成為了我們一首多姿多彩的生命奏曲.
弟兄姊妹.
用你可以用的聲音,姿勢.
來歌頌讚美神.
最後一個問題.
最後一個問題就是.
我們哪一個人要讚美神.
凡有氣息的都要讚美耶和華.
你們要讚美耶和華.
第一個是誰要讚美神.
你們.
你們即是誰.
即是神的子民.
即是信徒.
屬神的人.
信徒首先要帶領人去敬拜.
所以如果你想加入詩班.
想加入福音樂曲.
想加入鼓班.
想加入鋼琴.
又或者想加入結他等等.
去找可以教你的人.
去幫你去學習.
一起去敬拜.
一起去開聲.
剛才靈昌說得很好.
他說一會兒不要只有我的聲音.
你的聲音要比我的聲音更加深.
一起去敬拜.

$^{721}$在神的聖所讚美神.
屬神的人每週的崇拜.
是先決的條件.
是要出席的.
因為讚美神是很重要.
很重要的一件事.
其實家庭的崇拜也重要.
有時我們一家大小.
特別你家有小朋友的話.
一家大小一起去敬拜神.
又或者大家可以分享禱告.
分享大家關心的事情.
個人的靈修也是很重要的事情.
我們一起來每一個信徒都讚美神.
但我們看到不只是信徒.
這裡說什麼?.
「凡有氣息的」.
即是誰?.
即是所有人.
即是有些是未認識神的人.
凡有氣息的都應該讚美神.
即是神的創造是每一個人的.
所以不單單是已經認識神的人.
凡有氣息的人都要讚美神.
凡是活人都要讚美神.
我們知道世上每一個民族.
每一個種族.
其實都會讚美他們的神.
我們知道印尼未開發之前的獵頭族.
他們都會讚美.
美洲的印第安人,非洲的土人,回教徒,佛教徒,印度教徒.
他們都會了解到人所未能控制的事.
他們會敬畏天,敬畏宇宙,敬畏大自然.
敬畏他們的神.
他們都會做聖所來敬拜他們認為的神.
但我們知道什麼?.
只得哪一個?.
只得我們所認識的獨一真神.
祂才是神.
今天的世界.

$^{761}$世人不單沒有讚美做他們的耶和華.
反而是背叛,轉向敬拜偶像.
違反了神的要求.
但當我們問.
究竟他們怎樣才可以敬拜這個真神呢?.
羅馬書說了什麼?.
沒有人去傳.
誰會去聽呢?.
沒有人去聽,就不會明白.
沒有人去明白,就不會去求.
沒有人去求,就不會去信.
所以我們有責任是怎樣?.
我們屬神的人.
我們要去到凡有氣息的人面前.
向他們宣講這個神.
這個獨一的真神.
我們不單用口去說.
我們還要用什麼?.
我們要用我們的見證.
在他們當中.
去活出一個屬天的見證.
我們哪有奉差去宣講神的大能和大德.
還有這個全呼音的使命.
神把這個全呼音的使命交託給我們.
神要我們完成祂的使命.
要我們凡有氣息的人來讚美神.
其實這是我們的使命.
我們要在這裡讚美神.
我們更加要離開這裡.
去到我們住的地方.
去到我們所熟悉的人當中.
甚或去到我們不熟悉的人當中.
去讚美神.
紅恩堂快要慶祝十週年.
看到老牧師開始和大家說宣教的信息.
他在紅恩的網頁加了一些猜拳的資料.
也在你們的周刊裡面.
今天也有猜拳專訊.
他想我們每一個人都去明白.
因為誰來讚美神.

$^{801}$凡有氣息的都要讚美耶和華.
你們要讚美耶和華.
我們認識了神的人.
我們就要有責任出去.
將這個好消息說出去.
沒錯,你們將會猜揮楊姑娘出去.
你們都可以祈禱,可以奉獻.
但我們都要親身去.
延續你們以前去國內去做荒野曲.
或者想去歐洲.
去傳福音的心志.
將福音傳出去.
詩篇一開始第一篇是說到兩種人生.
他勸我們一生要敬愛神.
做一個有福的人.
而詩篇最後一篇,150篇.
他以歡樂,讚美作結束.
這是一個信徒一生有福的最後寫照.
他一生是一首詩歌.
是越唱越美.
因為詩歌的內容充滿了神的榮美.
而且也是榮美的神親自幫我們去寫成.
弟兄姊妹,我講道通常都是問這一句.
聽了45分鐘.
你會問「與我何干?」.
有沒有幫助?跟你有沒有關係?.
有關係.
我們要在這裡開始讚美神.
我們更加要將讚美帶到地極.
讓地極的人同樣可以聽到神的故事.
讓他們能夠歸向神.
也讓他們能夠讚美神.
我們低頭祈禱.
天父上帝我們雖然我們現在.
可能有時我們的情景未必太理想.
但是我們知道.
無論是黑暗或者是光明.
主你都是掌管一切.
你看到你過去的大德.
你過去的大能在我們當中彰顯.

$^{841}$你也將會繼續用你的大能來扶植我們.
我們就願你親自幫助紅人家.
讓我們能夠成為一群大能的子民為你所用.
聽我們的祈禱.
奉耶穌基督的名及.
阿們.
(影片完結).
\newpage



\section{彼得前書 4:7-11}
\label{sec:xUmmDjoz5Ds}
\textbf{2021.09.05《在苦難中的信仰實踐》 彼得前書 4\_7-11 (和修版) 老耀雄牧師}
\newline
\newline
連結: \href{https://youtube.com/watch?v=xUmmDjoz5Ds}{\texttt{https://youtube.com/watch?v=xUmmDjoz5Ds}} ~~~~ 語音日期: 2021-09-07
\newline
\newline
\hyperref[sec:65pDMB41cZg]{\small{< < < PREV SERMON < < <}}
~
\hyperref[sec:index_chronic]{\small{[返順時目]}}
~
\hyperref[sec:index_scriptual]{\small{[返順卷目]}}
~
\hyperref[sec:oCu8Aak_B2s]{\small{> > > NEXT SERMON > > >}}
\newline
\newline
彼得前書 4:7-11
\newline
\begin{longtable}{cl}
\hline
\hline
章節 & 經文 (和合本修訂版)\\
\hline
4:7 & \begin{tabularx}{0.7\textwidth}{X} 萬物的結局近了。所以你們要謹慎自守，要警醒禱告。 \end{tabularx} \\ \\ \relax
4:8 & \begin{tabularx}{0.7\textwidth}{X} 最要緊的是彼此切實相愛，因為愛能遮掩許多的罪。 \end{tabularx} \\ \\ \relax
4:9 & \begin{tabularx}{0.7\textwidth}{X} 你們要互相款待，不發怨言。 \end{tabularx} \\ \\ \relax
4:10 & \begin{tabularx}{0.7\textwidth}{X} 人人要照自己所得的恩賜彼此服事，作神各種恩賜的好管家。 \end{tabularx} \\ \\ \relax
4:11 & \begin{tabularx}{0.7\textwidth}{X} 若有人講道，他要按著神的聖言講；若有人服事，他要按著神所賜的力量服事，好讓神在凡事上因耶穌基督得榮耀。願榮耀和權能都歸給他，直到永永遠遠。阿們！ \end{tabularx} \\ \\
[1ex]
\hline
\hline
\end{longtable}
$^{1}$鴻升街弟兄姊妹平安.
近期聽到最多的話題是甚麼呢?.
一定離不開移民.
一個月內可能要吃幾餐離別的聚餐.
當然無論是個人也好,家庭也好.
他們選擇移民也有很多原因.
未必只是有一個單一的原因.
不過如果我們感覺到有壓迫.
或者有苦難的來到的時候.
也可能是我們其中一個很想離開的原因.
面對苦難的臨到.
有人選擇離開.
有人選擇留下.
對於非信徒來說.
安全感可能是一個最大的考慮.
對於信徒來說.
他們的安全感是從何而來呢?.
苦難的來臨.
可以給我們有甚麼信仰的反思呢?.
今天我想透過彼得前書四章七至十一節.
去看看在苦難裡面.
信仰可以給信徒一個甚麼的安全感.
我們一起去祈禱.
因為你救贖了我們.
讓我們在這個急速幻變的世代裡面.
能夠作你的子民.
你揀選我們生活在這個時代裡面.
去為你作見證.
求你兼顧我們的心.
好叫我們在這個逐世洪流裡面.
我們能夠仍然站立得住.
榮耀你的名.
我們的祈禱.
是奉教主耶穌基督得勝明志祈求.
阿們.
在我們進入經文之前.
首先看看經文的背景.
根據彼得前書一章一節的內容.
經文差不多透露了作者是誰.
而且在彼得前書裡面.

$^{41}$也有幾段經文.
是對於耶穌受苦的描述.
似乎都是出自於.
一些目擊證人的手筆.
很多證據顯示.
作者和耶穌有一個很親密的關係.
因此初期教會的教父都認為.
彼得前書的作者.
應該就是跟從耶穌的十二使徒之一.
也是第一批被耶穌呼召的門徒彼得.
我們對彼得最深刻的印象.
或者這是其中一個很深刻的印象.
必然就是耶穌被祭司長捉拿之前.
已經預言了彼得將會在雞叫之前.
三次不認主.
而事實亦一如耶穌所料.
彼得三次不認主的詳細事件.
很清楚地記錄在福音書裡面.
彼得的信仰生命固然有一個軟弱的一面.
但也有愛主的一面.
在主耶穌復活之後.
主耶穌在提比利亞海邊.
和彼得有一次面對面的單獨對話.
耶穌三次問彼得.
你愛我比這些更深嗎?.
彼得當時的回應就是.
你知道我愛你.
彼得是經過一段.
耶穌經歷過死而復生的信仰反思之後.
再一次去回應主耶穌對他的呼召.
你知道我愛你.
或者今天我們都需要再經歷一次.
耶穌基督的死而復生.
然後我們再一次單獨回應耶穌.
究竟你愛的是哪一個?.
然後主耶穌三次吩咐彼得.
你牧養我的羊.
這是耶穌升天之前.
親自交託給彼得的使命.
彼得領受了這個使命之後.

$^{81}$一直都忠心事奉主.
直到主後六十五年左右.
被尼祿王殺害而殉道.
彼得前書寫書信的年期.
大概是主後六十年左右.
即是彼得殉道前的五年.
而彼得前書的對象.
就是散居在小亞細亞一帶的外邦信徒.
這班主所交託給彼得牧養的小羊.
所居住的社會被排擠.
在信仰上受到迫害和苦難.
如果用今天的說法.
他們是一批弱勢社群.
在所居住的社區.
受到邊緣化的對待.
他們不是面對苦難的臨到.
而是苦難已經很真實地.
發生在他們身上.
彼得在這樣的背景之下.
寫信給這批人.
作出鼓勵和安慰.
在第一章六節.
彼得已經說了.
雖然你們必須在百般試煉中暫時休受.
百般試煉中.
是指他們所面對的試煉.
是逃避不了的.
既然逃避不了.
就要改變我們的心態去面對.
在第三章十四節.
彼得對信徒說.
即使你們為義受苦.
都是有福的.
不要怕人的威嚇.
也不要驚慌.
經文顯示出.
信徒已經身處在苦難當中.
第四章一節.
彼得又說.
既然基督在肉身受苦.

$^{121}$你們都應該將這個心智作為兵器.
彼得勸勉信徒要學孝主耶穌.
主耶穌所走過的這段苦路.
信徒都要有受苦的心智.
信徒雖然面對迫害.
但彼得在第四章第五節說.
他們必須在那位將要審判活人死人的主面前交仗.
迫害信徒的人.
最終都要面對主耶穌的審判.
申冤在主.
信徒正處身在苦難中.
他們在迫害與苦難面前應該如何自處呢?.
因此就帶出了七至十一節的訊息.
我們開始進入經文.
第七節一開始.
有一個前提.
就是萬物的結局近了.
結局就是指主耶穌再回來的事件.
近了.
是指時間上已經來臨了.
這個結局就是主耶穌將要再回來.
追討那些迫害信徒的人.
審判活人死人.
信徒面對一切的苦難都有終結的一天.
這一天就是主耶穌再回來的那一天.
這個日子何時臨到呢?.
不知道.
但這一天必定會臨到.
亦成為當時信徒的盼望.
彼得勸勉信徒在一個等待盼望臨到的期間.
信徒可以有什麼信仰回應呢?.
彼得提出了幾個提醒.
第一個.
要謹慎自守.
要警醒土局.
謹慎和警醒的原文是一個命令式的動詞.
兩個命令.
有一個意思.
就是頭腦清醒.
兩個命令都主要一個信仰的行為.

$^{161}$就是禱告.
在主耶穌再回來的日子之前.
信徒首先的.
第一件事是要頭腦清醒.
我們還記不記得在格拉森這個地方.
格拉森這個地方.
就有一個人被二千隻烏鬼附身的故事.
這個人不穿衣服.
也不住在家裡.
而是住在墳場.
他更加被人用鐵鏈和腳鏈鎖住.
因為被二千隻烏鬼附身.
耶穌幫他趕走那些烏鬼.
他就穿著衣服.
心裡明白過來.
這個心裡明白過來.
就是這一節經文所說的.
謹慎,警醒的意思.
因此信徒的禱告不是無關痛癢的.
信徒是要很清楚明白.
我為什麼要禱告.
以及禱告的目的是什麼.
他們需要很清楚的知道.
他們這一班信仰群體.
雖然現在分散在本都,加拉泰,多帕多加,亞細亞.
不同的地方.
而且彼此分隔得很遠很遠.
未必能夠經常有接觸有交流.
但是他們可以經過禱告.
透過禱告彼此守望彼此紀念.
他們要知道為什麼要活在這個世上.
也要知道活在這個世上的使命是什麼.
他們要知道苦難是有盡頭的一天.
因此今天就算為易受苦.
都仍然可以忍受.
耶穌不是在登山寶訓那裡已經說了嗎.
他說為易受逼迫的人有福了.
因為天國是他們的.
當時地上的羅馬政府不重視他們.
不要緊.

$^{201}$但是天國重視他們.
他們要知道有主為他們伸冤.
因此就算地上得不到公義的伸張.
但仍然可以繼續努力行善.
主是審判的主.
主會為受迫害的人伸冤.
還他們一個公義.
他們要知道邪不能勝正.
因為他們知道主耶穌才是那個得勝的王.
當時地上的羅馬政權是一個盛世的年代.
他們需要有一個很清楚的流露.
因為知道羅馬盛世是會過去的.
主耶穌才是永恆的主.
用禱告為自己將要面對的一切去守望.
今天我們去基督教的書室.
我們一定會看到很多屬靈偉人的自傳.
必定會發覺他們同樣有一個共同的屬靈特質.
就是很喜歡禱告.
他們分享到如果沒有禱告的生活.
根本就沒有動力.
所以有神庚祈禱的習慣.
甚麼是神庚祈禱?.
用今天的說法就是四點鐘.
四點鐘起床祈禱.
我相信在座的各位.
我講我自己.
四點鐘中有些人在做夢.
但這些屬靈偉人就說四點鐘就要去祈禱.
因為我根本沒有動力去面對我當日一天的工作.
在每天工作開始之前都必定祈禱.
讓生命被注入一種動力.
那種動力包括了怎樣愛,怎樣忍耐,怎樣恩慈以及怎樣謙卑.
以致我們可以面對每一天的挑戰和磨練.
弟兄姊妹,我很想跟大家分享.
禱告不是真的向空氣說話.
而是透過禱告向神知取能力去面對一切的挑戰.
禱告是信徒和神溝通的不二法門.
我很感恩我們洪恩家的童工團隊.
從星期二到星期四的上午.
我們必定有童工的祈禱會.

$^{241}$我們就在禱告裡面去知取力量.
然後我們才有屬靈的力量去牧養,去事奉神.
因此我鼓勵大家.
星期三的晚上回來教會參與禱告上的操練.
彼此同心合意向上帝禱告.
將讚美的祭,感恩的祭,仰望的祭.
不同的祭都在禱告裡面呈獻給主.
也讓我們更加敏銳.
神在我們面對試探的時候的挑戰.
當然我知道香港是一個很繁忙的都市.
有很多人住得很遠.
不能夠每一個星期都來.
但我在這裡向你發出一個挑戰.
你縱然不能夠每一個星期都來.
你能不能夠在神面前承諾.
我起碼一個月來一次.
一個月來一次.
如果都不行的話.
你在神面前祈禱.
兩個月來一次.
不要不來.
不要跟這個祈禱絕了緣.
這個祈禱是我們與神保持關係的不二法門.
練武功不單單需要內在的心法.
也需要武功的招式.
如果禱告是一個武功心法的話.
那招式是什麼呢?.
信徒可以從主的生命裡面.
汲取一種動力成為他們的內功心法.
他們需要在他們生活的地方.
彼此實踐相愛的信仰生命.
彼此相愛不是一個口號.
是一個信仰的實踐.
基督教不單單是一個講愛的宗教.
也是一個實踐愛的宗教.
彼得認為彼此相愛有一個目的.
就是愛可以遮掩許多的罪.
這個真理對於彼得來說.
應該感受很深.
從彼得跟隨耶穌開始.

$^{281}$直到耶穌復活升天之後.
我相信彼得很清楚.
也明白自己有多軟弱.
他所犯的罪.
在聖經裡很清楚地被記錄下來.
而沒有記錄的可能更多.
但在彼得裡.
他沒有活在罪咎裡.
因為他經歷了主耶穌對他的愛和寬恕.
然後在事奉裡.
他經歷了其他人對他的包容.
在保羅寫迦太書的時候.
保羅指責彼得因為割禮事件.
而跟外邦人疏遠.
保羅如何形容彼得呢?.
保羅說彼得裝假.
與福音的真理不合.
你想一想.
彼得當時已經是.
然後殺難的屬靈領袖.
他受到一個初出茅廬的彼得去批評.
不過彼得和保羅沒有因為這件事而分裂.
他們仍然都一樣忠心事奉.
彼得所犯的錯.
被寬恕.
被包容.
保羅性格的直率衝動.
也被接納.
被包容.
這就是愛的力量.
彼此相愛不單單是信仰的真理.
也是信徒在信仰群體裡.
能夠建立的親密關係的關鍵.
我們明白到沒有一個人是完美的.
神才是最完美.
但神沒有要求我們是一個完美的人.
每個人都有自身的優點和缺點.
彼此相愛就是欣賞別人的優點.
包容別人的缺點.
我們星期六晚有彪竿人生的課程.

$^{321}$在彪竿人生第十八課裡.
他提到彼此相愛的其中一個實踐的方法.
就是彼此的饒恕.
彼得饒恕保羅對他的冒犯和衝動.
保羅饒恕彼得一時的軟弱.
在彼此饒恕當中.
彼此都遮掩了對方的罪.
不再看對方的缺點.
原本罪本來想引發出來的.
那種破壞信徒合一的力量.
沒有了.
罪最喜歡就是破壞我們的合一.
彼此歡恕沒有了.
因此相愛有很多種體現方法.
歡恕和你心的接納.
彼此認罪.
真誠的分享.
都是相愛的一種表現.
每一種相愛的表現.
都能拉近你和別人的距離.
讓彼此的關係更加成熟和親密.
弟兄姊妹.
彼此相愛最適當的實踐場地在哪裡?.
一定是在教會.
在教會和弟兄姊妹相交過程裡.
我們都是三尖八角.
我們彼此都可能有機會有衝突.
無論是有心還是無意.
如果不盡早去處理.
就會成為心病.
藏在我們心裡.
彼此都有芥蒂.
這種芥蒂影響到我們生命的交流.
彼此生命的建立.
如果我們彼此能夠坦誠.
在神面前和解和復和.
這幅是一個很美麗的圖畫.
所以我很想鼓勵大家.
如果你真的和某一個弟兄姊妹有些嫌隙.
你很想和他和解.

$^{361}$我們就實踐彼此寬恕.
不過你會說:牧師,很尷尬的.
我和他有些嫌隙,很尷尬的.
你找我吧.
我幫你做一個和事佬.
透過我們三個人,包括上帝.
我們一起在信仰裡實踐彼此的要恕.
另外,弟兄姊妹.
我們有沒有一個彼此相愛的團契?.
你有沒有一群彼此能夠守望的同路人?.
無論你今天是否面對苦難.
是否面對迫害,你都需要這個同路人.
如果只是你一個人走你的信仰路.
你很容易走錯路,甚至走歪路.
很容易迷失方向.
因為只有你一個人.
你很容易落入魔鬼撒旦為你而設的陷阱裡.
在當時的社會文化.
在彼此相愛的基礎上.
我們可以看到彼得有幾個盡奴提醒.
他說要互相款待,不發怨言.
彼得所牧養的對象分散在不同的地方.
無論是關係的建立,還是彼此的守望.
其實都需要刻意的經營,要主動.
對於當時的羅馬世界,找女管不是很普遍.
有些從遠方而來的信徒.
他們就需要有人去接待.
這種接待可能包括.
任飲任食,好像食自助餐一樣.
包三餐,包食包住.
不過這都是進行主耶穌的教訓.
因為主耶穌在馬太福音25章說過.
祂說,這些事你做在我弟兄中一個最小的身上.
就是做在我身上.
究竟接待一個從遠方而來的信徒.
可以給他吃多久,住多久.
如果有些濫用別人的愛心的話.
在那裡白吃白住,濫吃濫住.
那怎麼處理呢?.
經文說,不發惡言,不發怨言.

$^{401}$其實這個可能反映到.
有些信徒其實濫用別人的愛心.
而引致的後果.
在初期教會有一份叫做.
十二使徒遺訓的手冊.
其中一段是這樣說.
若是來的人是一個旅客.
你們當盡量幫助他.
但他應該只停留在你們的住處.
不多過兩天.
或者必要的時候三天都可以.
倘若他想長住在你們裡面.
就要有一技之長.
他應該自食其力.
如果他不會首邁.
你們就按照你們的了解為他預備.
不要任一人為了他是一個基督徒的緣故.
而閒住在你們中間.
但若他不願意這樣做的話.
他就在售賣基督.
你們就要防備這些人.
這段教訓就是要告訴當時的信徒.
怎樣適當去接待外地的信徒.
亦都防止一些信徒濫用這個款待.
短暫的接待兩三天就足夠了.
長期的接待就要他自己自力其食了.
當然套用在今天的環境可能完全不適用.
今日從遠方的信徒可以住在酒店.
香港這麼方便.
無需要在家裡接待他.
但款待都可以是開放家裡的.
成為弟兄姊妹可以聚會的一個場所.
所以弟兄姊妹.
我們怎樣去實踐這個款待呢?.
今日我們很講究私隱.
我們未必願意開放家裡的.
就算我願意開放家裡.
我配偶或家人成員接不接受呢?.
有些人跟我說:木芝你不要來啦.
你一來我就要打掃三個鐘頭.

$^{441}$下樓茶餐廳就算了.
你一上來我就很煩.
是的,沒錯的.
但教會是一個屬靈的家.
教會可以開放給會友使用.
你可以帶一個朋友來.
回到教會.
讓他們可以有一個舒適的地方.
去彼此分享,彼此祈禱,彼此守望.
可以嗎?.
你家裡不方便.
帶他回來教會.
教會舒舒服服.
訂一間房間.
大家可以聊天,唱歌,分享.
帶朋友回教會.
在一個又寧靜又乾淨的房間裡面.
聊天,參觀教會.
介紹給我認識.
讓人知道教會是一個甚麼場地.
很溫馨,很溫暖,很歡迎你們.
是嗎?.
這也是一種事奉.
是嗎?.
我記得有一個很不好的見證.
有一次我問一個人.
如果你相信耶穌.
你想不想他帶回教會?.
他說:當然不想.
真是死定了.
你很想朋友相信耶穌.
但朋友相信耶穌.
你不想他回到你的教會.
是嗎?.
我不希望是這樣.
我希望你的朋友相信耶穌.
你很願意帶他回教會.
甚至未相信耶穌.
你也很願意帶他回教會.
帶他回教會,參觀教會.

$^{481}$跟牧者聊天.
讓他感受肅靈的家的溫暖.
彼得提出.
在彼此相愛的基礎下.
除了款待對方.
還可以彼此服侍.
經文說.
人人要照自己所得的恩賜彼此服侍.
作神各種恩賜的好管家.
彼得對恩賜的講論很少.
除了這一節.
沒有在其他經文提及.
不過有兩個希臘文.
都是用來形容恩賜的.
第一個意思是.
屬靈的事.
或者是關於靈的事.
這個字強調恩賜的本質和來源.
即是說恩賜不是天之才幹.
不是與生俱來.
而是從聖靈而來.
第二個字的意思是恩典的次語.
這個字強調恩賜是神恩典的次語.
恩賜既不是一種與生俱來的能力.
而是神給信徒的一種恩典.
即是說神給你就有.
不給你就沒有.
肅靈的恩賜.
可以用一個很簡單的定義.
是一種恩典的次語.
比較精細的解釋可以這樣說.
是賜給基督身體裡面的肢體.
一種特別的事奉能力.
換句話說.
恩賜只有信徒才有.
非信徒是沒有恩賜.
非信徒可以有不同的才幹和能力.
可以造福世界.
但不可以叫恩賜.
信徒的恩賜.

$^{521}$從神而來.
焦點是讓人在肅靈上得福.
而不是在肅世界上得福.
因此彼得說的恩賜.
即是恩賜的彼此服事.
不單單涉及個人信仰實踐.
亦涉及如何造就和建立教會.
信徒不同的肅靈恩賜彼此服事.
目的是要造就教會.
讓每一個進入教會的人.
能夠在肅靈上得福.
因此彼得提出.
若有人講道.
要按著神的聖言講.
我們知道保羅的宣教旅程.
每到一個地方.
都會先到猶太會堂傳福音.
後來因為會堂不方便.
所以就去信徒家裡聚會.
所以教會的初形.
就是從會友的家裡開始.
尤其是彼得書信的對象.
既然是外邦信徒.
根本不可能被猶太會堂所接納.
所以很常見.
都是以小群教會的形式去聚會.
彼得提出講道時.
需要有恩賜是很刻意的.
因為彼得後書已經提到.
當時有很多假教師的出現.
他們在聚會裡.
宣講一些違反基督教信仰的道理.
引誘信徒來經半道.
當時的講道就不像現在這樣.
有個講台可以讓傳道人去講.
而在聚會裡有人說有感動.
這樣就可以發言.
所以彼得才提出.
按著神的聖言講.
這個感動是不是來自神.

$^{561}$保羅亦都在提摩太後書.
勸勉提摩太.
按著正義講解真理的話.
初期教會有很多異端.
可以傳播得這麼快和廣泛.
就是很多假教師.
他們不是按著神的聖言去宣講.
而是按著自己的心意去講.
其實今天很多時候.
有些教會都被人批評.
有些上了台.
就當是一個大炮向著別人發射.
究竟他們是不是按著神的話語去講.
還是按著自己的心意去講呢?.
當然在教會裡.
信徒除了有講道的服事之外.
還有其他恩賜可以服事.
因此彼得也提出.
若有人服事.
他要按著神所賜的力量服事.
能夠服事的範圍可以很大.
亦都包括了不同的恩賜.
不過無論是講道也好.
還是其他服事也好.
其實一切的目的.
都是要建立教會.
讓耶穌基督得榮耀.
這是每一個事奉者應該要有的心態.
我們上星期.
舉辦了一個崇拜神學的講座.
就是讓每一個在崇拜參與的過程裡.
我們知道為何我們要參與崇拜.
就是讓每一個進入教會.
參與崇拜的每一個弟兄姊妹.
他們都可以在屬靈上得福.
可以來到教會裡與神相遇.
在崇拜裡能夠與神相遇.
是否一個福氣?.
這就是我們每一個事奉神的人.
應該要有的態度.

$^{601}$香港其實現在還有一些教會.
我們叫做小群教會.
可能只有一個傳道人.
其他事奉崗位都是信徒承擔.
主日學是屬靈領袖去教的.
去編排的.
團契團長會友負責帶茶經.
教會的財務是由會計恩賜的人負責.
甚至程序表,海報設計.
全部都是會友去做.
教會是全民皆兵.
為何可以這樣?.
是因為每一個肢體都當教會是一個家.
你想一想.
當家庭有需要的時候.
你就會自動自覺去幫忙.
回教會除了拿神的恩典.
還要彼此服侍.
信徒的恩賜願意彼此配搭.
教會就能夠更加被神所使用.
今天講題是.
在苦難裡的信仰實踐.
給得怯面信徒.
萬物的結局近.
末日臨到.
縱然我們身處在苦難裡.
我們仍要堅守信仰.
要有一個清楚明白的心去禱告.
要為自己要面對的一切.
為自己,為別人去禱告守望.
如果你今天仍然能夠屹立得到.
你就為那些軟弱的肢體祈禱.
你不要為你自己能夠屹立不倒.
而沾沾自喜.
你要為那些軟弱跌倒的人.
你要流淚祈禱.
在苦難裡怎樣可以活出信仰.
彼此相愛.
因為愛能夠遮掩彼此的罪.
我們不能夠這麼口響說我們沒有罪.

$^{641}$我們是一個有罪的人.
我們的罪除了主耶穌去赦免之外.
還可以有弟兄姊妹彼此寬恕.
彼此認罪.
彼此成為信仰的同路人.
信徒在犯錯裡.
每一個人都有錯.
但卻可以經歷彼此的寬恕.
要接待別人.
讓別人在苦難裡.
或者困境裡.
在逆境裡.
仍然感覺到人間有情.
因為接待別人就等同於接待耶穌.
是一種事奉.
我們每個星期五都有哥利留時宮.
我們派一些物資給有需要的人.
每一個有需要的人.
他來到我們教會.
接受這些食物.
這是一個開始.
他們需要的不是單單是食物.
他們需要的是主耶穌的愛.
我們接待他.
不是單單給他物資.
我們接待他.
是要展現耶穌基督對人的那種愛.
盼望我們弟兄姊妹能夠多參與.
教會的不同活動.
因為教會的每一個活動.
其實都是很想展現我們的信仰.
這個信仰的核心就是愛.
不單單是個別去服侍別人.
還要彼此服侍.
為的是要建立教會.
讓基督的身體更加榮美.
透過教會.
讓人可以認識基督.
透過教會.
可以引領迷路人走向光明.

$^{681}$教會能夠招聚一班聚人.
去讓信仰展現於人前.
這就是基督得榮耀的時刻.
我再說一次.
教會招聚一班聚人.
而這班聚人.
卻可以將信仰展現在人的面前.
人們會覺得.
這個人我認識他的時候.
嘩!粗口濫竊.
為什麼今天可以這麼溫文爾雅.
原來他信了耶穌.
基督得榮耀的時刻.
明白嗎?.
我們每一個都是聚人.
但當人在我們生命裡見到我們的改變.
他會問.
為什麼他的生命會改變?.
原來信了耶穌.
這個耶穌真是厲害.
當他說這個耶穌真是厲害的時候.
耶穌就得榮耀.
紅映家弟兄姊妹.
當我們說.
還在想.
究竟苦難是不是來的?.
什麼時候來的?.
彼得告訴我們.
萬物的結局近了.
我們不需要害怕苦難.
因為信仰可以讓我們承擔和面對.
當我們想.
究竟是留還是走的時候.
主耶穌告訴我們.
趁著白日.
我們必須做猜我來的那位的工.
我們每一個在地上仍然有使命.
我們每一個都有.
我們只要忠於我們的使命就足夠了.
無論是黑夜將來.

$^{721}$還是光明白晝.
信徒之間仍然要實踐彼此相愛.
彼此服侍.
因為唯有愛才能流傳到永恆.
信徒是為永恆而活的.
而不是為今生而活的.
我們一起去祈禱.
讓我們能夠在教會裡.
和眾聖徒去實踐和學習.
我們能夠彼此相愛.
以及彼此服侍.
都需要主你親自的保守和帶領.
甚至我們每一個.
坐在這裡的每一位.
我們都能夠在順境.
或者逆境裡.
我們能夠緊緊的跟隨你.
甚至我們每一刻.
我們都能夠感受到.
你與我們同在.
與我們同行.
作我們的安慰.
好叫我們能夠一生的跟隨你.
我們的祈禱.
是奉教主耶穌基督.
得勝的明智祈求.
阿們.
願主祝福大家.
\newpage



\section{哥林多後書 4:7-15}
\label{sec:oCu8Aak_B2s}
\textbf{2021.09.12《無懼患難的人生》 哥林多後書 4\_7-15 (和修版) 曾敏姑娘}
\newline
\newline
連結: \href{https://youtube.com/watch?v=oCu8Aak_B2s}{\texttt{https://youtube.com/watch?v=oCu8Aak\_B2s}} ~~~~ 語音日期: 2021-09-16
\newline
\newline
\hyperref[sec:xUmmDjoz5Ds]{\small{< < < PREV SERMON < < <}}
~
\hyperref[sec:index_chronic]{\small{[返順時目]}}
~
\hyperref[sec:index_scriptual]{\small{[返順卷目]}}
~
\hyperref[sec:fusOcxn2XoI]{\small{> > > NEXT SERMON > > >}}
\newline
\newline
哥林多後書 4:7-15
\newline
\begin{longtable}{cl}
\hline
\hline
章節 & 經文 (和合本修訂版)\\
\hline
4:7 & \begin{tabularx}{0.7\textwidth}{X} 我們有這寶貝放在瓦器裡，為要顯明這莫大的能力是出於神，不是出於我們。 \end{tabularx} \\ \\ \relax
4:8 & \begin{tabularx}{0.7\textwidth}{X} 我們處處受困，卻不被捆住；內心困擾，卻沒有絕望； \end{tabularx} \\ \\ \relax
4:9 & \begin{tabularx}{0.7\textwidth}{X} 遭受迫害，卻不被撇棄；擊倒在地，卻不致滅亡。 \end{tabularx} \\ \\ \relax
4:10 & \begin{tabularx}{0.7\textwidth}{X} 我們身上常帶著耶穌的死，使耶穌的生也在我們身上顯明。 \end{tabularx} \\ \\ \relax
4:11 & \begin{tabularx}{0.7\textwidth}{X} 因為我們這活著的人常為耶穌被置於死地，使耶穌的生命在我們這必死的人身上顯明出來。 \end{tabularx} \\ \\ \relax
4:12 & \begin{tabularx}{0.7\textwidth}{X} 這樣看來，死是在我們身上運作，生卻在你們身上運作。 \end{tabularx} \\ \\ \relax
4:13 & \begin{tabularx}{0.7\textwidth}{X} 但我們既然有從同一位靈而來的信心，正如經上記著：「我信，故我說話」，我們也信，所以也說話； \end{tabularx} \\ \\ \relax
4:14 & \begin{tabularx}{0.7\textwidth}{X} 因為知道，那使主耶穌復活的也必使我們與耶穌一同復活，並且使我們與你們一起站在他面前。 \end{tabularx} \\ \\ \relax
4:15 & \begin{tabularx}{0.7\textwidth}{X} 凡事都是為了你們，好使恩惠既藉著更多的人而加增，感恩也格外顯多，好歸榮耀給神。 \end{tabularx} \\ \\
[1ex]
\hline
\hline
\end{longtable}
$^{1}$願至聖聖三一保佑我們每一個人.
各位弟兄姊妹平安.
讓我們可以回到神的家去敬拜祂.
開始的時候我們再一次有一個禱告.
讓我們有這樣的福氣.
我們可以去親近你.
我們這麼卑微的人竟然可以來你的面前去敬拜你.
是何等的美好.
天父今日求你對我們每一個人說話.
你讓我們的心向你開放.
願意聖靈在我們心裡面去工作.
讓我們聽得明白.
我們就去跟從和回應.
奉耶穌基督的名字祈禱.
阿們.
今日我會說的是無懼患難的人生.
其實這個題目並不是太多會在講壇上說的.
我印象中講壇比較多是說.
我們現在經歷很多困難.
我們的生活有很多的需要.
耶穌基督怎樣幫助我們面對我們現在生活的困難呢.
心裡面真的有很多難處.
通常講這類比較多.
我們都很渴望聽到很安慰的訊息.
有我們很需要被安慰.
這些訊息都是很寶貴的.
但今日這件事我是很想挑戰大家.
回到我們信仰很核心的層面.
今日是一個非常好的場合.
是聖餐.
聖餐主日通常是比較多弟兄姊妹回來.
紀念耶穌基督的生.
祂的死.
應該先紀念祂的死.
然後紀念祂的復活.
這個是我和你信仰的基礎.
今日我講這堂道都很想回到這個基礎那裡.
最近當我們留意世界新聞的時候.
大家覺得這個詞是甚麼.
阿蓋達.

$^{41}$塔利班.
無論是怎樣都是我們很震撼的消息.
在阿富汗政權的改變.
自從塔利班在阿富汗掌權之後.
國內外都有很多的不安.
國內的基督徒更加擔心自己.
陷入更大的壓迫和傷害之中.
可能隨時都有生命的危險.
這個不是開玩笑的.
所以世界各地的人都很關心阿富汗的基督徒.
不知道他們在患難當中是否能夠撐得住呢?.
我們其實也曾經在我們教會的禱告會裡.
為他們禱告.
來到這裡可能你會說.
你今日這堂道不就應該對阿富汗的基督徒說.
無懼患難的人生.
跟我們有甚麼關係呢?.
我們在香港很安全.
我們在教會生活.
但求平安快樂.
無風無浪.
閒時三午好友相約去玩.
這樣就好了.
這個是我的人生.
所以這個患難的人生好像跟我很遙遠.
但我並不是這樣認為.
提摩太后書三章十二節這樣說.
凡立志在基督耶穌裡敬虔道入的.
也都將受迫害.
這個沒有分你在哪裡生活.
你是很自由的地方.
你是很不自由的地方.
沒有分的.
當我們是很忠心跟從主.
為主而活.
努力服侍主的時候.
我們都會經歷患難.
這是很自然的事情.
所以我今日說這堂道.
是很想去預備大家的心.

$^{81}$亦都很想挑戰大家去敬虔道入.
願意為了服侍主而受患難.
而且在面對患難的時候.
都不要害怕.
究竟我們應該怎樣去面對患難呢?.
我們可以從保羅的身上去學這個屬靈的功課.
在保羅的事奉生涯裡.
保羅經常是經歷患難的.
但他的生命並沒有因為這樣而被打垮.
保羅得性的秘訣在哪裡呢?.
我們來看以下兩點.
第一點:認識我們自己的本相.
大家在這裡看到甚麼呢?.
一個很特別的瓦器.
瓦器裡面會閃的.
會發光的.
在經文第七節裡這樣說.
「我們有隻寶貝放在瓦器裡.
要顯明這莫大的能力是出於神.
不是出於我們」.
保羅用一個暗語去形容自己.
他形容自己是一個瓦器.
瓦器是用泥土製造的.
價值較低的.
微不足道.
亦都比較容易碎的.
你一撞到地上.
它就會碎的.
所以你見到這幅圖.
瓦器裡面有些裂紋.
是很脆弱的.
保羅看自己的生命都是這樣.
在患難當中.
他更加看到自己的微小脆弱.
在我和你的眼中.
我們看保羅是一個很有牧者心腸的傳道人.
是一個勇敢的宣教士.
去了不同的地方宣教.
帶領不同的人信耶穌.
他的業績非常標榜.

$^{121}$他的思辨能力很強.
是一個很偉大的神學家.
他寫的書信有13卷被列入聖經當中.
他的影響力無法估計.
一位這麼強的.
在我們眼中看為一個很偉大的事奉者.
他竟然形容自己是一個瓦器.
我們其他人還可以說些什麼呢?.
我相信我們都會同意.
其實我們的生命一樣.
都是這麼脆弱,這麼卑微.
在患難面前,在很多考驗面前.
我們很難說自己很強,很厲害.
靠著自己的力量就可以勝過.
很難.
保羅自己也說自己只是一個瓦器.
得性的認識自己本相之後就要想想.
得性的關鍵在哪裡呢?.
就是閃閃發光的部分.
那個寶貝.
我們有隻寶貝放在瓦器裡.
為了顯明莫大的能力是出於神.
不是出於我們.
得性的關鍵是寶貝.
這個寶貝是指保羅一直傳揚的福音.
這個福音說的是.
耶穌基督為了拯救世人的生命.
而在十字架上犧牲了自己的生命.
他被埋葬.
第三天從死裡復活.
聖天為世人承受了救恩.
我心想今天如果走到台下.
問弟兄姊妹.
我相信你也很容易回答這個答案.
我們都知道福音是什麼.
如果你經常傳福音還會上口.
我想跟你說這個福音極為珍貴.
好像寶藏一樣.
就有這個寶貝在瓦器裡.
瓦器就不再一樣.

$^{161}$瓦器的價值不再一樣.
同樣有福音進入了保羅的生命.
保羅也不再一樣.
他就能夠從神那裡得到力量.
去勝過苦難.
是因為這個寶貝.
是因為這個福音.
他可以從上帝那裡得到力量勝過苦難.
人在軟弱當中更加能夠經歷和見證上帝的大能.
當我們有福音在我們生命裡的時候.
我們就能夠從神去支取力量勝過苦難.
接著經文就跟我們分享.
我們的生命因為這個情況會有什麼經歷.
保羅就用四個對比告訴我們.
福音的大能和來自神的力量會造成什麼呢?.
我們看四個苦難的經歷.
第一,處處受困,卻不被困住.
我們無法確保人生不會遇到壓迫,困難,沮喪,挫折.
我們無法確保,也無法避免.
但是處處受困也好.
也不會真的沒有出路.
這個處處受困是說我們面對很多壓迫.
很多不同的壓力,困難.
是纏著我們,圍著我們,好像一面圍牆.
感覺是要讓我們窒息,好像走不出去.
在這樣的處境也好,你仍然會有出路.
你不會被困住.
福音的能力,來自上帝的能力.
可以令我們在這種無比的壓迫裡.
仍然有盼望,看到有路可走.
這是第一個經歷.
第二個經歷,內心很困擾,卻沒有絕望.
雖然我們內心感到迷失,混亂,不知道該如何處理.
但是就算再差,也未到達絕望的境地.
是因為什麼?.
有福音在你生命裡.
是因為有神的能力,發出來.
就算你內心的情感是這樣,你的情緒是這樣.
有很多事情你控制不了,你也未到達絕望.
又出路.

$^{201}$第三,遭受迫害,卻不被撇棄.
遭受迫害,這裡形容一種像野獸被追捕的情形.
很緊迫.
這裡形容保羅為了傳福音的緣故.
被敵人一直追捕.
雖然我們會被敵人一直追捕,一直趕著去.
但是卻不被撇棄.
被誰撇棄呢?.
不會被上帝撇棄.
人撇棄我們不是最可怕.
最可怕是上帝撇棄我們.
原來有福音,有上帝的力量在我們裡面.
有神在我們裡面.
無論你遭遇什麼,被人追捕,被迫害.
上帝也不會撇棄你.
這是最寶貴的.
第四個經歷.
這個應該是最深刻,也是最困難的.
擊倒在地,卻不致滅亡.
這個擊倒在地,簡直是在說你的肉身.
被擊倒在地,面臨死亡,傷痕纍纍.
但是我們是不會滅亡的.
我們仍然會存活.
就是因為有福音在你生命裡面.
就是因為有上帝的力量.
你外在的肉身,你經歷什麼都好.
你不會滅亡.
這四個經歷,都是保羅經歷過的.
他大大小小不同的經歷.
所以在哥林多後書11章23到27節.
他這樣說.
他們是基督的傭人嗎?我說句狂話.
我更是,我被他們忍受更多勞苦.
坐過更多次監牢.
受過無數次的鞭打,上上冒死.
我被猶太人鞭打五次.
每次四十減去一下.
被軍打了三次.
被石頭打了一次.
遭海難三次.

$^{241}$一晝一夜在海裡掙扎.
我又屢次行遠路.
遭江河的危險,盜賊的危險.
同族人的危險,外族人的危險.
城裡的危險,曠野的危險.
海中的危險,假弟兄的危險.
我勞碌困苦.
上上失眠,又饑又渴.
忍饑內寒,赤身老體.
很多患難,保羅經過了.
如果今天大家一齊走過來.
說說人生有多苦.
我經歷過困難多大.
保羅有很多東西可以跟我們說.
但這一切都沒有將他打倒.
因為有福音在他裡面.
有救恩在他裡面.
有上帝的能力發出來.
保羅不被打倒.
而是讓神的能力在他裡面彰顯出來.
保羅多了很多機會去經歷上帝.
這一切都讓保羅能夠逐一去勝過那些患難.
今天我和你說起來很輕易.
好像就這樣數白欖.
但當日對於保羅來說很困難.
你今天叫我去經歷這一切.
我真的不夠膽說.
我覺得這個挑戰實在太大.
所以你可以看到上帝的能力和恩典是多麼的浩大.
無論保羅經歷的患難有多大.
他都可以從神之所有力量去經過.
今天我想問問弟兄姊妹你們又如何呢?.
福音對你們來說有什麼影響?.
是不是我相信上帝.
有福音在我生命裡我只是拿著護照.
拿著護照上天堂呢?.
只是這樣?.
有這個福音在你生命裡.
有沒有令你能夠從上帝之取力量去面對你的患難呢?.
這正正是我要問你的.

$^{281}$可能有人會問.
這些患難對保羅有什麼意義呢?.
很有意義.
第十節這樣說.
我們身上尚帶著耶穌的死.
使耶穌的生也在我們身上顯明.
保羅經常經歷患難.
他就在患難裡參與了耶穌的死.
我再說一次.
保羅經常經歷患難.
就在患難裡他參與了耶穌的死.
正正是因為這份參與.
就令他可以經歷耶穌基督復活的大能.
有死就必有生.
在耶穌的死裡有份.
就能夠經歷耶穌復活的大能.
所以他是可以很有力地跨過重重的難關.
第十節.
保羅進一步解釋這個意義.
他說因為我們這個活著的人.
尚為耶穌備至於死地.
使耶穌的生命在我們這個必死的人身上.
顯明出來.
當日耶穌被門徒出賣.
被交付在人的手裡.
一步一步地邁向死亡.
今天我們常常為了耶穌的緣故.
而面對患難和死亡的威脅.
其實我們所遭遇的.
或者當日保羅所遭遇的.
就好像主所遭遇的一樣.
正正是因為當日的保羅.
在患難裡參與了主的受苦和死亡.
所以耶穌基督復活的大能.
也彰顯在保羅的生命裡.
當我們在患難裡.
參與了主的受苦和死亡.
耶穌基督復活的大能.
也會彰顯在我和你的生命當中.
所以這是很大的盼望.

$^{321}$我們不怕受苦.
因為會經歷耶穌基督復活的大能.
特別是為主而受苦的時候.
參與主的死和復活.
讓我們更加明白貼近主.
耶穌為什麼需要受苦,受死和復活?.
今天是聖餐主日.
每個月第二週都回來紀念.
是為了什麼?.
耶穌為誰受苦,受死和復活?.
就是為了今天坐在禮堂裡.
看直播的每一位弟兄姊妹和同工.
祂為我們受苦,受死和復活.
今天當我們都是為主的緣故.
為了服侍祂,為福音緣故.
去受苦的時候.
經歷患難的時候.
我們就更加明白耶穌.
進入到耶穌的受苦,受死和復活裡.
這個關係是很緊密的.
我們才明白耶穌.
真正與祂有一個很緊密的結連.
這種關係不是一般的.
是很容易明白的.
是很寶貴的.
亦給予我們很大的動力.
去面對我們的生活.
因為我們會發覺主是很明白.
很明白我們的受苦.
很明白我們要經過的.
我們才發覺.
原來受苦的時候才知道.
耶穌的經歷是這樣的.
這個關係才真正的緊密.
同樣的.
這樣的患難參與.
都會令我們與其他人有奇妙的結連.
十二節這樣說.
「這又看來,死在我們身上運作.
生卻在你們身上運作」.

$^{361}$保羅經歷患難是為了什麼?.
是傳福音給哥林多後書的受眾.
讓他們信耶穌.
生命得夢拯救.
他要經歷患難.
好讓人們得著拯救.
他面對死亡.
人們就得著生命.
所以說.
生就是在人們的身上運作.
這是非常寶貴的.
這裡有幾幅圖畫.
經歷耶穌基督的死和生.
很深刻的.
就在這一次的行程裡.
大家看到.
這個地方很明顯就是中東地方.
我不說那個國家在哪裡.
當時我還在讀神學院的時候.
我本身是讀跨文化系的.
曾經有想過將來要去宣教.
在一個暑假.
與老師及同學去了一個中東國家裡.
去試一下體驗一下.
如果我在那個地方宣教.
生活會是怎樣呢?.
當時也有八個星期.
在裡面也做了很多事情.
例如學習語言.
學習阿拉伯語.
也要接觸當地人.
參與不同的事工.
認識一下當地的教會.
去講英文的教會.
也去看華人教會.
也去了不同的地方.
那次的行程.
那八個星期非常豐富.
令我很深刻的感受.
其實不只是阿富汗這個地方.

$^{401}$會有白基督徒.
我去這個國家也會有.
這個國家最主要的信仰是伊斯蘭教.
所以進去後.
不可以這樣說.
我們是一群來自香港的.
還要說是神學院的同學.
第一時間會將你趕出去.
你進去後.
如果你到處張揚.
跟別人說.
我們是一群來自香港的.
神學院的一群神學生.
一定會被人抓住.
很多事情要處理等等.
甚至乎有機會有生命危險.
在過程中.
所以去到這個地方要很小心.
老師也教了我們很多東西.
怎樣跟別人說話.
怎樣說自己的身份.
怎樣處理.
衣著如何等等.
所以你看到.
其實我們的衣服是比較長的.
長的衣服,長的褲子,裙子等等.
我記得我去到沒多久.
我就聽到一件事.
有一位當地的牧師.
被人在三樓那裡.
是丟下來了.
當場死了.
但是當局並沒有查清楚這個案件.
是誰殺死這個牧師的呢?.
但肯定這個牧師不是自己.
在那麼高的地方跳下來.
自己解決自己的生命.
不是的.
肯定是有人將他丟下來.
他的生命就這樣結束了.

$^{441}$這件事不了了之.
然後很快.
我又聽到有一個見證.
有一個女孩子.
年輕的女孩子.
本身是信奉伊斯蘭教的.
後來有人帶她信了耶穌.
她回教會聚會.
她感到很快樂.
敬拜讚美上帝.
真的很好.
真的在過程中.
感受到耶穌的愛.
她覺得很寶貴.
她很喜歡.
她真的善這個位.
她得到這個福音.
是個寶貝來的.
她就一直珍惜回教會的時間.
家人發覺了.
發覺了之後.
最主要是她有一個弟弟.
弟弟就來逼迫她.
囚禁她.
禁錮她在家裡.
不讓她出外.
不讓她回教會.
又恐嚇她.
會殺死她.
這些不是只是嚇嚇.
在當地來說.
如果背叛了伊斯蘭教.
家人可以殺死你.
不是犯法的.
當時家人真的有這個打算.
想恐嚇她.
看她會不會改變.
然後就請他們的阿公.
就是那些律法的教師.
走過來遊說她.

$^{481}$你快點放棄你的基督教信仰.
快點回來.
不要再做叛徒了.
不准回教會.
一直鎖住她.
很多的恐嚇.
我就看到.
是會受苦難的.
是真的會為了耶穌而被交於死地.
要面臨死亡.
我很深信.
那位姐妹在那一刻.
很經歷到.
什麼叫患難.
耶穌基督的死是一回事.
是很不容易.
當時這個姐妹.
其實有一個選擇.
是可以逃跑的.
教會也有朋友跟她說.
你要不要逃跑.
我可以安排.
從今以後.
只要逃跑離開這個家.
就不可以再回去.
她想過之後.
她決定.
就算死.
我也不離開家.
因為我離開了以後.
就不用再和家人傳福音了.
那誰會和家人講信仰呢?.
我救自己是容易的.
我在這裡可能會犧牲.
我會被交於死地.
我要面對死亡.
好像到處都沒有出路.
我是四面受敵.
但是.
這是為了耶穌的緣故.

$^{521}$我會留下.
但我不會走.
就是這樣.
教會知道她的意向.
尊重她的意向.
當時我們也有為她禱告.
這個生命.
後來是很奇妙.
上帝真的給她很奇妙的經歷.
她在後面.
竟然經歷很大的轉變.
神讓她的弟弟.
竟然在其他地方找到工作.
弟弟沒有辦法.
不可以再看著姐姐.
阻礙著她.
以致弟弟要離開家工作.
姐姐就有了自由.
但我相信.
在這個實際的處境改變之前.
是耶穌基督復活的大能.
在她心裡工作.
給她很大的力量.
盼望.
去勝過這一切外在的環境.
也是的.
好像我今天所說.
我講起來好像很輕描淡寫.
好像很容易.
但當天聽的時候.
我們的心情也是很沉重.
在那段時間.
我很感受到.
原來死亡是很真實的.
原來我們是會很容易.
會為耶穌受苦難的.
需要的.
有一次更加特別的.
但是在這些圖片裡.
我沒有留下我同學的圖片.

$^{561}$也沒有留下其他.
只留下我自己.
和一些當地人的圖片.
是為了保護的緣故.
我記得有一次更加直接.
上帝透過過程去挑戰我的思考.
你是否真的可以?.
如果有一天你要為我犧牲.
你是否真的可以?.
有一天我們就一起要做一個功課.
我們要入清真寺.
入這些回教寺.
我們有一個體驗.
一群神學院的同學.
分開不同的組合.
進去後.
我們不可以暴露我們的身份.
也不可以說他們去做什麼.
但那天我和我的同學.
有一個男同學被人抓住了.
抓住了.
有一個在清真寺裡.
是一個很重要的人.
他想知道我們做什麼.
很兇地問我們.
你們是誰?.
你們來這裡做什麼?.
一直很兇地問我們.
那一刻.
我立即停下來.
我想一下應該怎樣回答他.
純粹說是遊客.
我們來看看.
覺得很好奇.
對於這個清真寺裡.
是這樣回答嗎?.
還是什麼呢?.
對方一直在追問.
追問了很多次.
很兇的.

$^{601}$你是誰?.
你在這裡做什麼?.
我想如果很誠實地回答.
是不是應該.
應該有下一步的經歷.
是不是這樣?.
一直有很多思考.
那一刻.
不是現在那麼平安.
當時的環境.
令我感受到.
是有威脅的.
有危險的.
我的同學很期待.
當時他真的是一個.
很快的男同學回答.
我們是基督徒.
我心想.
你什麼事?.
他用英文回答.
我們是基督徒.
我想你用不著這樣回答.
誰知道對方的反應.
是很有趣.
我相信對方的英文.
應該不是很好.
如果不是正常.
應該會抓了我們.
對方就.
嗯,OK.
繼續放過我們.
我心想你沒理由.
我真的不敢回答.
我肯定是.
誤會了.
英文程度很差.
所以我同學都覺得很奇怪.
我這樣回答.
他竟然都放過我.
在那個時候.

$^{641}$我發覺.
哎呀,很大挑戰.
是不是?.
雖然對方還沒有.
做什麼傷害我.
只是那一下.
忽然間.
如果有一天真的.
我真的被抓了.
就是問你.
問你的信仰.
問你要不要.
放棄你的信仰.
如果真的有一天是這樣.
或者你放棄你的事奉.
你為什麼要傳福音?.
是不是?.
行不行?.
那一刻我發覺.
是一個很大的挑戰.
死亡不是很遠.
但是對於寶來說.
這些是經常經歷的.
經常為耶穌.
被放於死地.
經歷耶穌的死.
然後經歷耶穌基督復活的大能.
多麼親近.
不是開玩笑的.
所以我在你面前.
我不敢去誇口.
免得自己顯得很虛假.
因為我覺得不容易.
真的不容易.
但是今天我講這些堂道的時候.
我自己也在這裡.
很想挑戰一下自己.
挑戰一下大家.
會不會進入這種生命當中.
捨得為耶穌基督.

$^{681}$犧牲寫記.
付代價呢?.
保羅和古代的詩人.
都是有同樣的信心.
13節這樣說.
「但我們既然有從同一位靈而來的信心.
正如經常記著.
我信故我說話.
我們也信所以也說話.
我信故我說話」.
是來自詩篇116篇10節.
那位詩人受了極大的困苦.
他向耶和華呼求.
上帝就聽他禱告.
拯救他.
保羅和詩人的經歷很相近.
都是遭遇患難.
最後都經歷了上帝的拯救.
古代的詩人是因為信而說話.
保羅亦是因為這份信而分享福音.
這是非常寶貴的.
保羅不但在患難的生活中.
經歷了耶穌基督復活的大能.
他還相信他的肉體將來一定會復活.
14節這樣說.
「因為知道那使主耶穌復活的.
也必使我們與耶穌一同復活」.
這是在說將來.
「並使我們與你們一起站在他面前」.
保羅相信將來他死了之後.
父神會讓他復活過來.
會和復活的主在一起.
不單這樣.
上帝亦會令哥林多的信徒復活.
上帝會令保羅和信徒一起復活.
然後一起來到主的面前.
這份相信帶給人很大的盼望.
地面上的苦難很短暫.
會過去.
但有一天我們都會復活.

$^{721}$所以我們不需要害怕.
這時的苦難會是怎樣.
這個盼望是很大.
弟兄姊妹.
你信耶穌那麼久.
這個福音不只是起初時成為你的激勵.
成為你的幫助.
而是每一天.
我也很願意這個福音成為你生活的動力.
每一天我們都應該靠福音繼續生活下去.
從那裡支取力量有盼望.
而不是只是停在某一個位置.
最終的時候.
保羅經歷苦難.
最終的目的是甚麼呢?.
15節.
凡事都是為了你們.
好使因為既藉著更多的人而加增.
感恩也格外顯多.
好歸榮耀給臣.
凡事是指保羅所宣講的內容.
以及他所經歷的一切苦難.
凡事都是為了你們.
即是保羅所說的任何事.
以及他所經歷的苦難.
都是為了你們好.
為了哥林多教會的好處.
因為是在說神的恩典.
包括了神所賜下的救恩.
祂給信徒的眾多的恩惠等等.
保羅很努力宣教.
令到很多人信耶穌.
經歷到神的恩典.
這樣神的恩典也會更加越發地增加.
信徒經歷神的恩典就會感恩.
最後歸榮耀給神.
在這裡最後有兩點值得我們留意.
保羅最終受苦的目的是為了別人.
為了哥林多的信徒.
而最後最後為了榮耀神.

$^{761}$不是為了榮耀自己.
今天我挑戰大家.
忠心跟從主.
也挑戰我自己.
忠心跟從主.
努力服侍主.
挑戰我們一起為主受苦.
你要知道我們這樣受苦.
會祝福很多人的生命.
我想說香港為什麼這麼多教會是怎樣來的.
大家有沒有想過.
今天我和你很舒服吧.
坐在這個冷氣室裡.
很平安.
坐在一個很舒服的椅子上.
敬拜讚美上帝.
我們有自由.
大大聲地嘗試著讚美上帝.
很快樂.
很想上帝得榮耀.
你有沒有想過.
有多少人為了大家可以坐在這裡.
享受這個敬拜.
付出了多少代價呢.
香港這麼多教會.
這麼多信徒.
享受這些福氣.
因為什麼呢.
有很多宣教士來香港.
付上很大的代價.
遠離自己的家鄉.
經歷很多事.
是為了祝福別人.
是因為有很多很多人.
經歷苦難受苦.
付上很大代價.
我和你.
我們香港的信徒才可以享受這一切.
我們一直下去.
是否只停在這裡.

$^{801}$繼續享受三五好友.
閒時好好互相探望.
吃喝玩樂.
這是我們的人生呢.
還是我們都要去.
為別人的好處.
付代價呢.
這是我挑戰大家想的.
還有呢.
有沒有想過榮耀上帝.
為主受苦.
一定榮耀神.
我們一同去禱告.
天父上帝在這裡.
今天藉著寶來的這段經文.
讓我們去思考受苦的意義.
為什麼我們不怕受苦的人生.
因為受苦很重要.
就是在經歷苦難裡.
我們要更加認識.
更加明白.
更加知道你的用心.
我們也被挑戰.
要為你的緣故去受苦.
好像寶來一樣.
天父求你今天激動.
我們紅人家中的影子們的心.
你要我們起來.
離開我們的安樂窩.
要為了服侍你.
付代價.
我們不一定每個人.
拋頭顱灑熱血.
但至少要肯起來.
將我們的恩賜擺上.
將我們的時間精力金錢擺上.
要帶領人去信主.
要按上帝你的心意去服侍.
為上帝你的緣故去付代價.
做上帝你喜悅的人.

$^{841}$我們經歷患難.
就在你的死裡有份.
你的復活大能.
也會在我們的生命裡彰顯.
今天願意耶穌基督的復活大能.
每天與弟兄姊妹同在.
以致我們很有能力.
面對我們的患難.
以致我們都被激勵.
更多更忠心的跟從主.
祈禱交托.
奉耶穌基督的名字.
阿們.
\newpage



\section{馬太福音 12:18-21}
\label{sec:fusOcxn2XoI}
\textbf{2021.09.19 講道《壓傷的蘆葦》 馬太福音 12\_18-21 (和修版) 王楚男傳道}
\newline
\newline
連結: \href{https://youtube.com/watch?v=fusOcxn2XoI}{\texttt{https://youtube.com/watch?v=fusOcxn2XoI}} ~~~~ 語音日期: 2021-09-23
\newline
\newline
\hyperref[sec:oCu8Aak_B2s]{\small{< < < PREV SERMON < < <}}
~
\hyperref[sec:index_chronic]{\small{[返順時目]}}
~
\hyperref[sec:index_scriptual]{\small{[返順卷目]}}
~
\hyperref[sec:_ayvrCyPl7U]{\small{> > > NEXT SERMON > > >}}
\newline
\newline
馬太福音 12:18-21
\newline
\begin{longtable}{cl}
\hline
\hline
章節 & 經文 (和合本修訂版)\\
\hline
12:18 & \begin{tabularx}{0.7\textwidth}{X} 「看哪，我所揀選的僕人，我所親愛，心所喜悅的；我要將我的靈賜給他，他必將公理傳給外邦。 \end{tabularx} \\ \\ \relax
12:19 & \begin{tabularx}{0.7\textwidth}{X} 他不爭吵，不喧嚷，街上也沒有人聽見他的聲音。 \end{tabularx} \\ \\ \relax
12:20 & \begin{tabularx}{0.7\textwidth}{X} 壓傷的蘆葦，他不折斷，將殘的燈火，他不吹滅，直到他使公理得勝。 \end{tabularx} \\ \\ \relax
12:21 & \begin{tabularx}{0.7\textwidth}{X} 外邦人都要仰望他的名。」 \end{tabularx} \\ \\
[1ex]
\hline
\hline
\end{longtable}
$^{1}$各位弟兄姊妹,早安.
今天有機會來到當中.
代表母愛來問候.
也介紹蘇怡的先生.
他來了,站起來給大家看看.
剛才已經看了.
阿倫.
很奇妙,神的恩典很大.
今天跟大家分享一段經文.
聖經可以打開.
在你手上的是《馬太福音》.
《馬太福音》十二章.
十八到二十一節.
剛讀過.
雖然版本不同.
可能不更清楚所講.
在1990年的時候.
我在伊利沙巴醫院有機會去事奉.
在這條遠目當中一起事奉.
在十年的時間我跟上帥都問了一個問題.
我問神為甚麼我們人有那麼多苦難呢.
為甚麼有那麼多苦難呢.
因為每次見到都好多令人都.
好像浪蓋過來是受不住的.
我覺得心裡面常常問的.
特別是遇過三兄弟連續患老癌.
一個一個離世.
我就心裡不期然問主人究竟發生甚麼事.
我那時如果心水清.
1990年的時候我都很年輕.
三十歲左右.
我問究竟人生為甚麼呢.
為甚麼有那麼多苦難呢.
在這裡經文講到.
阿斯利康的蘆葦.
我直接講.
蘆葦在地利斯蘭的時候.
是一些看不上眼的.
但它可以造成一些可以用來吹的笛子.
當中可以奏起一些很漂亮的樂章.

$^{41}$我記得有間教會有一個弟兄吹一些陶笛.
同樣是那個聲音漂亮到不得了.
我相信笛子有它的特別聲音.
同樣在摩西很小的時候.
有一個小的嬰兒.
當時不能夠養育嬰兒.
唯有扔到海裡.
但他用一個蒲草.
即是當時的蘆葦草.
造成一些籃子放在裡面.
並且當時有些油燈.
油燈用麻鹽油做一些燃點.
原來麻鹽油燃點之後.
吹熄的時候都有一陣味道.
在60年代70年代.
我們一代人就不忘記那種味道.
就是火水味.
當我們熄了火水爐.
那種味道就出來了.
但現在當然沒有了.
當時這些情況.
小小的背景說到.
原來這段經文的出處.
其實是在《爾作亞書》42章1-9節.
在《僕人之詩》.
本來還有.
總共是有不止這麼少.
有四首的.
在42篇裡面.
之後就有四首的《僕人之詩》.
今天只不過拿了第一.
《僕人之詩》的第一小小的部分.
跟大家分享.
這個情況.
第一點.
壓傷的心.
我想容易就明白.
壓傷的心.
就好像這裡說.
壓傷的爐位不截斷.

$^{81}$燈火不熄滅.
我在想這個過程的時候.
其實好像很簡單.
但當我想的時候.
原來可以簡單分兩大類的分開.
第一類就是一個.
我們會在內或外面來造成壓力.
這個就是由外.
一個就是由內造成的壓力.
好像很簡單.
但當說是由外的壓力來的時候.
原來壓傷的爐位不截斷.
為甚麼不截斷呢.
其實是.
原來裡面是一件.
壓傷這兩個字.
很簡單.
壓傷.
如果是外力的時候.
很自然是甚麼.
有可能大自然.
或者是一些甚麼.
人為.
或者是天生出來.
造成的一些情況.
這個就是在一些這樣的情況造成的.
其實壓傷的爐位.
有沒有人覺得很可惜.
其實沒有的.
我相信你走過草.
經過甚麼地方.
壓傷的爐位.
沒有人在乎的.
沒有人在乎.
沒有人理會.
其實原來爐位可以做很多好事.
當你說沒有用的時候.
原來它就勝在一條生命.
剛才剛剛說的.
摩西就是靠著爐位.

$^{121}$造了一艘蒲草的小船.
成為甚麼.
成為王子.
如果沒有爐位載著他.
這個故事就不成了.
很重要的.
很重要.
其實在.
如果深水清的人看到馬太夫音.
原來馬太夫音前一點點在這張裡面.
他上文.
十二章一到十七節.
是在說人子是甚麼.
人子是安息的主.
安息是甚麼.
安息本來是為人而設的.
為人而設的.
當時耶穌在.
他去到每一個地方的時候.
很多人都將這件事倒轉來說.
讓人受了很多壓力.
真正.
原來耶穌就是.
他想去告訴人聽.
扭轉回.
安息是為人而設.
耶穌在心裡面.
是看人.
看每一個生命是甚麼.
是天大地大的事情.
所以.
在這裡說到.
壓傷的爐位.
他不截斷.
張纏的燈火.
他不熄滅.
耶穌可不可以在這麼艱難的裡面.
去扭轉這些局勢呢.
基督的時代也是一個殖民的時代.
但是好像香港一直轉變一樣.

$^{161}$背景很相似.
但耶穌在這個情況下很特別.
我看到近來很熱門的情況.
外體雖然殘.
但是內心卻很勇敢.
很能夠看到他們的生命力.
這班人在裡面.
雖然殘而不廢.
並且為香港人爭光.
我們都很開心看到他們為香港人爭光.
但是.
有些人有負面的情緒.
雖然為了國家奪取很多獎牌.
但是因為抑鬱,負面.
覺得自己沒用.
就會有另外一個想法.
他覺得自己不行.
最後還是選擇了結自己的生命.
原來很多時候我們得和失.
都是一念之間.
一念之間怎麼去面對呢.
原來主耶穌視我們每一個生命很寶貴.
寶貴到一個地步.
祂願意去安慰.
願意去接觸.
願意赦免我們一切的罪.
祂撫摸,祂醫治.
在四福音裡面的時候.
耶穌常常叫人所做的事出其意表.
祂叫那些小朋友來祂的裡面擁抱祂.
甚至祂說的故事.
一個浪子的故事.
祂就好像父親一樣.
當有人願意悔改.
願意擁抱他.
其實祂好像天父一樣.
擁抱耶穌就是一個這樣的心腸.
祂要救我們脫離這一切艱難.
這些都是往常有名的人.
其實我們都一樣.

$^{201}$當你回轉我們心願意回頭.
在艱難裡面起來的時候.
就可以有很大的改變.
我們在浪茄每次走過一個叫做.
下山.
我們叫做環山路.
環山路那棵樹.
在幾年前有一次很大風.
山竹事件.
那棵樹扭斷了.
但很奇妙.
那棵樹每次我走過.
那棵樹好像是手刮一樣.
每次見到它都很硬.
原來很奇妙.
在這裡裡面.
在神像我們經歷那麼多苦難.
似乎好像有少許出路.
祂叫我們更加在越困難的地方.
越生得硬.
要留意所有的樹木.
它扭過來或斷過來.
再扶好再生的時候.
它的栗子.
是很硬的.
比之前更加硬.
不知道硬了多少倍.
原來神很奇妙.
這些人.
你也有份.
你就是那些人.
你就是我們的弟兄姊妹.
終於要改變過來.
相反.
我們在呂村.
我聽到.
看到.
這個不是真的上來的.
例如.
我們不可以出街.

$^{241}$呂村有一次有很多PO.
很多的錦法官都去探村.
其中一個就說自己見證.
我浪費了人生16年的時間.
他做什麼.
他跟著一個人.
經常踢他打他.
打到他頭破血流.
但是他都要跟著他.
聽的時候那個PO都流眼淚.
他還幫那個男人生了兩年.
他都不知道.
他每次醒來就吃.
醒來就睡.
醒來就吃.
很淒涼.
醒來之後就在村裡面說.
這件事情.
每一個錦法官都.
每一個都拿了很多紙巾出來哭.
我們那次不是一個.
一直數.
數了很多的.
生命講出自己的難處.
我在當中的時候.
我都忍不住來哭.
甚至在裡面.
叫每一個生命被改變.
不好意思.
我太急了.
另外在.
下一張就是.
不行.
就是.
我課了.
不要緊.
下一張就是講到一個.
養魚的一個朋友.
他就.
我輸了.

$^{281}$養魚的朋友.
他經過一些風暴.
所有魚都反了.
他的所有東西一無所有.
我睡著了.
我太急了.
所有的魚都睡著了.
他其實都很.
當時精神混亂.
很慘.
要進到青山裡面接受治療.
但是這個人一樣是.
當時我見到他就在QE認識他的.
他後來信了耶穌.
然後他關心人.
他很開心.
原來這些就是來自外在的壓力.
第二點就是來自內在的壓力.
無論心靈情緒等等.
聖經說在裡面講到.
原來.
將塵燈火拍不催滅.
原來在一些社福機構.
有這樣的說話.
我看完都很詫異.
老人不是漸漸變老的.
是突然變老的.
原來人生當遇到困難的時候.
其實我們.
上次我姐妹講完.
她說一把年紀.
看上去都不算老.
其實是不會變老的.
只是突然間你會在.
心中突然間出現很多的東西.
很多負面的東西.
很多的情況.
突然間才變老.
我覺得我都挺贊同這句話.
原來我們就是很奇妙的.

$^{321}$沒有反應.
不好意思.
我看回前後.
可以嗎?.
可以啊.
好.
原來不單單這樣.
當人風促殘年的時候.
特別是他講到將塵燈火.
就好像很微弱的燈芯.
我們以前過去用油燈的時候.
風促殘年就好像甚麼呢?.
就好像一些老人家.
他感受到.
本來家人是很好的.
但是當被家人的遺棄.
或者老人家想.
我不想連累人.
所以他有時候覺得很絕望.
縱然他進了老人院也好.
其實有很多人照顧.
但是他內心是甚麼呢?.
內心是覺得.
老人院好像是很感到被遺棄.
所以我有一段時間.
在院舍裡面事奉工作的時候.
我都見到一個家人去探望他的兒子.
但是不試探望他的兒子.
他放下了糖餅.
怎樣呢?.
他不是用兩重閘.
都是禁閉式的.
當他去開一個小窗.
他說我放下東西就走了.
我說你進來吧.
進來看看你的兒子.
他說不用了.
當然他不來也沒辦法.
但是在我心裡面.
我都有點難受.

$^{361}$如果我是他的兒子.
我慘了,明白嗎?.
為甚麼家人不來呢?.
你明白我說甚麼嗎?.
為甚麼他不來呢?.
當然他的智商可能不是很高.
但是你猜他知不知道?.
他知道的.
他知道,他很難受.
同樣的.
我說的處境其實是很事實的.
更加是我們戶外裡面.
我們星期天.
就會給弟兄打電話回家.
弟兄打電話回家的時候.
我記得一個很深刻的印象.
特別是弟兄回家的時候.
他媽媽打電話跟他聊天.
兒子你為甚麼還不回家呢?.
因為他媽媽開始有一個腺亡症.
他講著講著就忘記了.
最淒涼的就是當弟兄畢業的時候.
剛剛這個疫情下.
他媽媽又安排到老人院住.
他已經忘記了是誰來的.
他又探不到他.
當時不讓探,你明白嗎?.
那種很為難.
大家想見.
但是媽媽可能完全忘記了他是誰.
那種很無奈.
其實在這裡想說.
無論怎樣艱難也好.
神也是.
耶穌很愛我們.
耶穌好像剛才唱歌.
耶穌是說.
神不會截斷.
神不會摧滅.
無論我們是多艱難.

$^{401}$無論我們在甚麼環境也好.
原來神將你一切艱難.
神都知道.
聖經說壓傷路位.
將塵埃登火.
可能令到我們.
無論對方或自己也好.
心靈受創.
但是耶穌將你.
銘刻在心上.
銘刻在你的掌上.
這個心上就是耶穌.
將你的名字釘在十字架上.
將你釘在他的心裡面.
那個痕跡.
神是知道.
我相信無論怎樣也好.
神知道你愛他.
你發出禱告神也知道.
在九十年代的時候.
我有機會去到.
不單單.
那時候參加一些跟太太.
很想關心一些殘疾的朋友.
一些傷健的朋友.
他們傷健的朋友.
如果有機會去接觸他們的時候.
你是捨不得走的.
他們的團體是很真實的.
每個人都說話.
很艱難才說出他們的話.
知道嗎?.
只說自己的名字.
說了多久.
我叫什麼名字.
原來他們很艱難.
所以他們對你的出現.
是很開心.
還有很想說真話給你聽.
這個就是我們當時.

$^{441}$我看到你們.
其實我心裡很尊重他們.
每一個人出現.
我對他們很尊重.
每次都花很多人力物力.
當時不是很流行的復康車.
我們已經可以用復康車.
借別人的教會去聚會.
我自己當時初出茅廬.
就不太知道.
我就想神可不可以.
將來他們有一間教會.
有一間叫回聲相見協會的教會.
不是挺好嗎?.
當時也被人禁止.
我們不可以將教會變成.
全部都殘疾.
但相見是相和見.
都可以.
在2000年之後.
神的作為很奇妙.
就成了一間相見協會的教會.
並且裡面有很多人信主.
我看到神的手很奇妙.
所以我們怎麼說呢?.
請不要輕看你.
不要輕看甚麼?.
不要輕看別人.
將來可能是神所重用的人.
我們不要輕看身邊任何人.
不要輕視他們.
神是愛他們的.
第二點.
主治的手.
很簡單的說.
當時我一句經文.
在想甚麼事呢?.
甚麼將真理傳開呢?.
第十二節.
十二章十二二十節.

$^{481}$等他施行公理.
叫公理得勝.
他施行公理.
他施行公理就是堅持正義.
在《熱愛中文》譯本.
他要堅持正直.
余文說他憑誠實將公理傳開.
在《熱愛創下書》更加這樣說.
但要相信神必傳揚公理.
按公義審判.
原來是說到一個情況.
神要將公理傳開.
怎樣呢?.
過去神已經在《熱愛創下書》預言了.
神耶穌基督的降臨.
現在已經成為事實.
耶穌在新約的時候.
已經來到世界上.
修補人的罪問題.
只要相信他的人.
就得到拯救.
並且當他心裡的創傷.
壓傷或破碎.
神要修補.
如何將公理傳開呢?.
怎樣呢?.
叫公理傳開.
叫公理得勝.
就是必傳揚正義.
直到最後的正義得勝.
這就像一種橫幅.
好像有人在舉著東西.
遮住了,先回頭.
原來叫公義得勝.
好像一種橫幅.
《熱愛創下書》就說到.
按公義審判.
最後都要得勝.
就是直著主耶穌基督.
來恩訊我們都要稱義.

$^{521}$因此我們就像這裡.
我們去到最後一個.
神的審判.
神愛的審判.
這裡.
下一張.
原來神愛到多盡呢?.
我.
要這樣說.
耶穌不是用權勢.
威脅來.
那是什麼呢?.
那是以愛和憐憫.
來施行公義和正義.
公理.
正義公理.
祂連自己都不顧.
為了怎樣顧呢?.
祂將生命放在.
釘身在十字架.
犧牲自己.
這個就是不顧.
不顧法就是.
簡單說.
祂為我們去厭棄.
被苦待.
被鞭傷.
被釘死.
被埋葬被復活.
都是什麼呢?.
為了你.
神為了我們.
為公理為正義來擺上.
為了成就舊恩來擺上.
這個就是苦惱十四站.
神就是為我們這樣擺上.
其實在四福音裡面.
很清楚.
有說到.
在.

$^{561}$一切神.
去醫治所有的.
被鬼附的.
被疾病的.
要生命被改變.
生命就是這樣一個一個.
主動地.
憑著神的恩典來改變.
耶穌主動地將公理傳開.
這是事實.
包括什麼呢?.
包括你和我.
原來這個公義的傳開是包括誰呢?.
是我們.
我們每一個人的生命.
剛才蘇儀的分享.
我們將生命擺上.
當有一段時間.
可能四五年前.
我要等候工場.
好像你們一樣.
弟兄姊妹就等.
找工作.
等候.
轉工的時間.
等工場.
等了一個月.
等了半年.
等了一年.
兩年.
害怕嗎?.
等了兩年.
再等你很害怕.
我好寒.
你不害怕我害怕.
太太就說.
不如去讀補健課程.
學一下你以前的東西.
我想也好.
以前在醫院做過很多.

$^{601}$涼血壓.
弟兄姊妹就說.
你試試去量吧.
那時候.
啪一聲就有涼血壓.
高血壓低血壓.
那時候也不知道什麼事.
但當我在等候的時候.
我開始明白.
明白什麼?.
原來弟兄姊妹很辛苦.
才找到工作.
那時候就開始.
之前也沒有緊會的.
變成在等候的過程就明白.
弟兄姊妹很艱難.
很多眼淚.
我心裡就很明白.
當我真的聽太太說.
去讀補健課的時候.
我又要.
我也不知道.
我要去執藥.
原來我去執藥的地方.
下面.
下面是什麼?.
下面是一個糞區.
是一個糞池.
我執了一段時間.
當然我才知道.
你為什麼要把東西全都拔出來.
原來下面是糞池.
原來人生就是這樣.
你才知道原來.
有很多東西要學習.
我那時候帶一些大診.
帶到晚上.
你也知道現在急症很久了.
我看到晚上.
明天就要下班.

$^{641}$我還沒下班.
在那裡等.
我又不喜歡.
我又怕其他下一家人辛苦.
就沒有去.
變成有時候真的很辛苦.
最艱難的一次.
我在等工場的時候.
有一段時間.
還沒到補健之前.
補健的學習的時候.
我去做一些消防員的.
美國消防員也要煮飯.
有人說我會煮這些焦糖.
我說好啊你試試.
當我去見了工的時候.
我不是很會煮.
但是他又請我去做.
我到處去.
有一次去到一個地方.
我買了很多菜.
車上有很多一抽一扔的菜.
上了巴士.
去到一個地方.
在過程裡面.
突然間收到一個電話.
以前在天水圍一間教會.
單親家庭的一個姊妹.
App過來問我.
記不記得我是誰.
原來她坐在我旁邊.
我也不知道.
App過來.
我真的記得她是誰.
跟她說.
談了很久.
當我想和她打招呼.
知道的時候.
她就站在閘門說拜拜.
原來我想告訴大家.

$^{681}$我也是一個人.
我都愛我的家人.
我真的沒有意識.
當時我心裡面.
好像很難受.
什麼都被你看光了.
我做到這麼多工作.
很奇妙.
原來我們都是淪落人.
我一個人只不過在不同的角色.
做不同的事情的時候.
你就反而多一點體會到不同的人.
不同的處境.
這個就是那個情況.
其實神很愛我們每一個.
神很愛我們每一個.
於是我們一起為主去做見證.
一起為主擺上我們的生命.
當我看到這段經文.
在那裡說.
要使神功利.
叫功利得勝.
我開始知道.
我們的見證.
就能夠叫生命.
叫功利.
叫神的話活潑起來.
我們不怕去.
好像保羅說什麼.
保羅說的都是一些失敗見證.
保羅就是這樣.
我們都是這樣.
在《浪劇》我們看到.
每一個生命或者女村.
每一個都是一個.
從糞堆被提拔出來.
神太愛我們.
我一個人懂得出來.
做人可以這麼說.
真的不懂得做人.

$^{721}$但是看到的時候.
神是這麼愛我們.
祂們能夠都為主發光.
最後.
我不說這裡了.
神是愛我們的.
不離不棄.
神是愛我們.
不會忘記我們.
神愛我們.
放在神的心裡.
我先轉過來.
不要緊.
我們始終是連在一起的.
不要緊.
我想有一點總結.
不管我們今天怎樣都好.
神都在祝我們每一個的生命.
第二點就是.
留意耶穌的愛和憐憫.
祂安慰抱緊我們.
第三就是.
無論我們在順境逆境.
天父最明白你的處境.
神很明白我們.
無論今天我們在紅恩堂.
不單純是紅恩堂.
整個香港的教會.
都面臨著很多的改變.
很多的牧者離開.
其實紅恩堂很有恩典的.
所以不需要比較任何人.
可以在這裡成長就成長.
生命就是這樣可以改變.
讓我們在神裡面打我們的生命.
讓我們生命在神裡面.
茁壯成長.
這個就是我們要將功利傳開.
分享到這裡.
\newpage



\section{路加福音 5:1-11}
\label{sec:_ayvrCyPl7U}
\textbf{2021.09.26 講道《與神同工》 路加福音 5\_1-11 (和修版) 楊明麗姑娘}
\newline
\newline
連結: \href{https://youtube.com/watch?v=_ayvrCyPl7U}{\texttt{https://youtube.com/watch?v=\_ayvrCyPl7U}} ~~~~ 語音日期: 2021-09-29
\newline
\newline
\hyperref[sec:fusOcxn2XoI]{\small{< < < PREV SERMON < < <}}
~
\hyperref[sec:index_chronic]{\small{[返順時目]}}
~
\hyperref[sec:index_scriptual]{\small{[返順卷目]}}
~
\hyperref[sec:BYuafMUFtfY]{\small{> > > NEXT SERMON > > >}}
\newline
\newline
路加福音 5:1-11
\newline
\begin{longtable}{cl}
\hline
\hline
章節 & 經文 (和合本修訂版)\\
\hline
5:1 & \begin{tabularx}{0.7\textwidth}{X} 耶穌站在革尼撒勒湖邊，眾人擁擠他，要聽神的道。 \end{tabularx} \\ \\ \relax
5:2 & \begin{tabularx}{0.7\textwidth}{X} 他見有兩隻船靠在湖邊，打魚的人卻離開船，洗網去了。 \end{tabularx} \\ \\ \relax
5:3 & \begin{tabularx}{0.7\textwidth}{X} 有一隻船是西門的，耶穌就上去，請他把船撐開，稍微離岸，就坐下，在船上教導眾人。 \end{tabularx} \\ \\ \relax
5:4 & \begin{tabularx}{0.7\textwidth}{X} 他講完了，對西門說：「把船開到水深的地方下網打魚。」 \end{tabularx} \\ \\ \relax
5:5 & \begin{tabularx}{0.7\textwidth}{X} 西門說：「老師，我們整夜勞累，並沒有打著甚麼。但依從你的話，我就下網。」 \end{tabularx} \\ \\ \relax
5:6 & \begin{tabularx}{0.7\textwidth}{X} 他們下了網，圈住許多魚，網險些裂開， \end{tabularx} \\ \\ \relax
5:7 & \begin{tabularx}{0.7\textwidth}{X} 就招手叫另一隻船上的同伴來幫助。他們就來，把魚裝滿了兩隻船，船甚至要沉下去。 \end{tabularx} \\ \\ \relax
5:8 & \begin{tabularx}{0.7\textwidth}{X} 西門‧彼得看見，就俯伏在耶穌膝前，說：「主啊，離開我，我是個罪人。」 \end{tabularx} \\ \\ \relax
5:9 & \begin{tabularx}{0.7\textwidth}{X} 他和一切跟他一起的人對打到了這一網的魚都很驚訝。 \end{tabularx} \\ \\ \relax
5:10 & \begin{tabularx}{0.7\textwidth}{X} 他的夥伴西庇太的兒子雅各、約翰，也是這樣。耶穌對西門說：「不要怕！從今以後，你要得人了。」 \end{tabularx} \\ \\ \relax
5:11 & \begin{tabularx}{0.7\textwidth}{X} 他們把兩隻船靠了岸，就撇下所有的，跟從了耶穌。 \end{tabularx} \\ \\
[1ex]
\hline
\hline
\end{longtable}
$^{1}$各位早晨.
今天感恩我有機會在這裡一起聽泰語姐妹的分享.
美雲.
我想我們大家都很受到鼓勵.
我們一起低頭祈禱.
然後我們一起聆聽神的話語.
親愛的天父上帝.
是你恩典.
主耶利撒讓給我們有這個機會.
在雄恩這個地方.
一起敬拜你.
一起事奉你.
有福音的工作.
有彼此建立的工作.
求主紀念我們.
主耶利撒讓我們在今天早上能夠一起聽到上帝的說話.
主耶利撒讓我們無論說的聽的都得著你的恩典.
得著你的話語.
祈禱仰望你奉主耶穌基督的聖名.
阿們.
今天這一段經文可能很多電影姐妹也很熟悉.
也對於這段經文好像有一首歌.
知不知道?.
開到水深的地方.
開到水深之處.
這一處我們會有一首歌.
我應該是對著哪個方向呢?.
是.
不要緊,你幫我下一張吧.
我們在神的恩典當中有聖經.
而聖經當中我們有這一段.
這一段是我們很熟悉的.
很熟悉的時候,第一樣是講到關於彼得的一些經歷.
彼得將船開到水深之處.
一開始的時候.
其實彼得和耶穌的接觸.
不是去到耶穌找他.
跟他說:你來跟從我吧.
做得人的漁夫吧.
不是在那裡開始的.

$^{41}$原來一開始是因為施洗約翰有一群門徒.
其中有一個可能就是西門彼得的弟弟安德烈.
安德烈認識了耶穌之後.
覺得這個耶穌一定是神要猜他的.
於是就介紹自己的兄弟給耶穌認識.
他介紹西門彼得給耶穌認識的時候.
那時候西門彼得的名字.
如果我們看回這一書.
他的名字還叫西門.
他還沒有叫彼得.
就是那一次和耶穌基督接觸的時候.
耶穌基督一見到他.
就馬上幫他改了名字.
就說你是約翰的兒子西門.
但是你的名字應該叫做石頭.
基法或彼得的意思就是石頭.
原來石頭這個名字是耶穌基督親自為他改的.
弟兄姊妹.
你有沒有印象小時候叫什麼名字?.
記不記得小時候身邊的人.
家裡的人.
同學仔,朋友仔,同村的人怎麼叫你?.
那個可能就是你最深刻對自己的認識.
或者是別人對你最深刻的認識.
但是耶穌一見到彼得就幫他改了一個名字.
名字就叫做彼得.
叫他做石頭.
我們又看一看.
接著下去.
《路加福音》第五章的經文.
第五章第一節.
就說耶穌去到格尼撒勒這個地方.
(請梁傳道).
去到格尼撒勒這個湖邊的時候.
很多人就黏著耶穌.
黏到一個地步.
直接就推著他向一個方向.
因為太多人推著他.
推到真的去到海邊.
去到岸邊的時候.

$^{81}$他就覺得其實這樣.
是又做不到什麼.
又不安全的.
推著推著可能整群人就下水了.
於是耶穌就做了一件事.
他就自己上了一個人的船.
為什麼呢?.
上了船的時候.
他就可以向這群很想聽到.
很想黏著他問話,說話的人.
來到港島.
而這一群很想聽到的弟兄姊妹.
就可以安全地站在岸邊也好.
坐在石頭地,沙地也好.
去聽耶穌基督講道.
耶穌基督就上船.
上船的時候他就上了哪個船呢?.
其實他對彼得已經有一些認識.
彼得對他也有一些認識.
他又幫彼得改了名字.
另外他又再早前一點.
去過彼得家裡.
或者叫他岳母家裡.
治好他岳母的病.
彼得認識他.
知道一點點.
對於耶穌是一個很特別.
又能夠治病的人.
這兩句都很明白.
耶穌上了他的船.
之後就叫他撐開一點點.
就出去外面.
撐開一點點出去外面的時候.
其實他那時候是甚麼時候呢?.
我覺得彼得很特別.
他是很不辭勞苦的.
原因就是他是捕魚的.
捕魚是晚上天黑了.
吃完晚飯才出去開工.
出去開工之後就天亮下班.

$^{121}$回來就洗魚網.
洗了魚鱗.
看看魚網有沒有哪裡需要補.
補完之後掛乾魚網.
他就休息一下.
睡一睡.
晚上再開工.
應該他就是回到夜更.
但是就看到耶穌來了.
很多人就想聽耶穌講道.
誰知道耶穌就上了他的船.
耶穌又沒有訂到艇.
又沒有下訂.
又沒有預約.
我問一下.
其實你們回來事奉有沒有預約?.
沒有.
沒有的吧?.
沒有預約你們就回來事奉了.
老牧師一定會編排一下.
主席呀.
哪位主席.
哪位領師呀.
師琴不用編排的了.
我們永遠都有一個.
專業師琴.
專業鼓手.
是.
吃了早餐了沒有?.
有沒有吃早餐才來事奉?.
有沒有很趕.
要馬上走過來?.
彼得那時候可能是的.
等下做完事了.
他想著下班的時候.
耶穌上了他的船.
又沒有問他吃了飯沒有?.
又沒有問西門彼得.
吃了早餐沒有?.
我帶個麵包給你.

$^{161}$你幫我撐船.
沒有.
耶穌就上了船.
撐出一點點.
彼得可能是很仰慕耶穌.
很想聽道.
他很乖很馴服.
撐出一點點.
撐出一點點的時候.
就去聽道.
聽完道了.
想做什麼?.
可能大家都會的.
晚上睡不著就會聽道.
彼得聽完道.
可能想做某一件事.
就是回去休息一下.
誰知道耶穌說完.
就對西門彼得說.
你撐船出去湖中心水深的地方.
原來湖中心水深的地方.
深度大概四十米.
岸邊你只能走腳下去一點點.
再下去一點點就不行了.
然後就一直深下去.
那個地方特別深.
為什麼叫水深之處呢?.
那裡的深度比隔離附近的地方.
深了將近一倍.
比起岸邊深了幾倍.
所以可能很多魚.
美魚也在那裡.
其實大家都知道.
他們捕魚.
耶穌就對西門說.
你撐船出去水深的地方.
然後就在那裡捕魚.
西門彼得就說.
他叫耶穌做老師.
其實我們剛開了一整晚工.

$^{201}$其實現在挺累的.
不過既然你說道.
老師我很尊重你的.
你既然說道的時候.
我們就撐出去.
就捕一次魚.
他的員工都很聽話.
他的員工就照著水所說的.
跟這個沒有預約的老師.
耶穌撐出去.
彼得很謙卑.
又能夠帶領到這一班人.
這一班他的員工.
又無怨言地聽到這一位.
無專業資格的耶穌.
教他們怎樣去捕魚.
耶穌的專業資格.
是做拉比.
是講聖經的.
捕魚不是他的專業.
我們在在有哪一位強項是捕魚.
我們最多喜歡去釣魚.
我們未至於到強項是去捕魚.
捉一捉魚過來.
但是他們去做.
他真的很天秤床服.
如果有一天有一個人.
又不是很懂電腦的東西.
突然之間去跟梁傳道說.
你的電腦應該這樣做的.
音響影音你全部要這樣做的.
梁傳道又是很謙卑.
他都會哈哈哈哈.
跟著他就會繼續工作了.
是不是.
變了樣子.
(笑).
這個西門彼得就開了船出去.
跟著就看一看魚.
你猜結果是怎樣.

$^{241}$你們每個都很熟悉的.
結果就是他一看魚之後.
就很多魚進了魚網.
很多魚進了魚網.
這個彼得他是很特別的.
他的特別就是在於.
他很機智的.
他很機智就是不是死了.
或者哎呀那怎麼辦這麼多魚.
因為魚的份量是他駕船不夠裝的.
他的網都綁不住的.
不知道是他開船出去的時候.
他另一對很朋友的兄弟.
這個雅各和約翰.
就跟著船出去.
還是他們打到招呼.
見到他的情況就馬上去幫忙.
他就這一架船.
有一架船去到他身邊的.
估計這一架就是雅各和約翰的船.
因為接下來就有提到他們了.
有提到他們一起對耶穌基督所做的事.
很有反應.
這兩架船過去之後.
發生了什麼事呢?.
這兩架船過去之後.
就是將魚.
他們將魚搬上船的時候.
裝滿兩架船.
弟兄姊妹.
他們有一架船出去.
但是最後是有兩架船去拿魚的.
為什麼叫做拿魚呢?.
原來他們當時捕魚的時候.
告訴你考古已經找到一架.
是第一世紀在格尼薩拉湖.
當年代的船的化石.
找到那架船出來.
長度是二十多呎.
闊度是七呎多.

$^{281}$高度就是深度是四呎.
大概每一次可以有十四五個人在那裡.
裝多些人又不是很好.
然後他們可以裝的魚都是有限的.
但是不知道為什麼.
另外有一架船就來幫忙了.
於是他們將魚.
他們不是整個網收上來的.
他們是將魚逐條逐條.
或者一堆撥到船上.
然後等魚在船上.
網才慢慢收上來.
魚的份量是一架船裝不完的.
一架船裝不完.
於是要兩架船.
兩架船能不能裝得完呢?.
原來剛剛好.
可能多多個人或者多多條魚.
船就會沉了.
聖經說的.
船差點要沉.
不過沒有沉.
是差點要沉.
不過沒有沉.
即是滿了.
網差點要爛.
不過沒有爛.
沒有掉.
你領受恩典是不是每次都是這樣?.
很豐富.
豐富到差點收不完.
但是剛好收得完.
很多.
但是不會帶來損失的.
不會船沉了.
爛了.
網裂了.
那就要搞很久.
不是.
他們仍然可以回家洗好網.

$^{321}$掛好網.
今晚再繼續去網魚.
但是魚的份量就剛剛好.
耶穌可能數過.
就兩架船.
裝得完.
嘩 多好啊.
你的感覺是怎樣?.
當是這樣的時候.
你的感覺是怎樣?.
那個恩典剛剛好.
足夠你的需要.
很豐富.
如果是我.
第一件事可能是感恩.
感恩.
又跟別人說天父很恩典.
但是給的反應不同.
給的反應.
他就看自己.
他覺得.
他自己是一個罪人.
他第一件事就是向神.
向主耶穌基督說.
我是一個罪人.
你不要這麼靠近我.
我是罪人 我骯髒的.
我的心靈是污穢的.
我是一個罪人.
記不記得一開始.
我們說耶穌上那架船的時候.
叫他下網的時候.
他怎樣稱呼耶穌?.
他叫他做老師.
是對的.
他一直叫他做老師.
可能很多人都叫耶穌做老師.
但是這一次的時候.
他不是叫耶穌做老師.
他說主啊.

$^{361}$你離開我遠一點.
我是一個罪人.
我和你不匹配的.
我站在你旁邊是不對的.
我不配站在你旁邊的.
甚至乎可能是想說.
我不配去幫你撐船的.
弟兄姊妹.
你回教會這麼久.
你怎樣稱呼耶穌?.
耶穌在你心中.
是你子女都會回來祈禱說的神.
是你覺得祂都像是神.
還是祂真是你的神?.
是你的主?.
是你會向祂祈禱.
是你的事情交給祂的嗎?.
見到祂的時候.
我們求恩典是多一點的.
感恩都會有的.
認罪我們好像少一點.
可能我們現在神聖化了.
沒有犯那麼多罪了.
我們少一點認罪.
讓祂領受了恩典之後.
竟然第一件事.
祂是.
我是一個罪人.
我不配.
我真是一個罪人.
我不配事奉你.
這是彼得的反應.
弟兄姊妹.
在這件事裡面.
你對耶穌基督的工作.
和彼得和祂的關係.
你有什麼認識?.
除了他們之外.
你說有另外一艘船.
就是在這裡告訴我們.

$^{401}$可能這艘船就是亞國約翰的船.
因為他們都在這件事上.
他們覺得很驚訝.
可能他們就是.
另一艘裝滿了魚回來的船.
接著之後.
彼得的反應.
耶穌說了一句話.
耶穌和彼得說.
祂說不用怕.
從今之後.
彼得你要做得人的漁夫了.
你要去得人了.
你不要只是去.
做一些和魚有關的工作.
你現在是要去做和人有關的工作.
接著兩艘船就洗了岸.
就是說剛才的對白.
那些驚訝的人.
他們真的就在湖中心附近.
然後那艘船就慢慢洗了岸.
然後.
誰跟了耶穌?.
原來不是西門彼得.
撇下所有的.
跟從了耶穌.
他們把兩艘船.
洗了岸.
他們就撇下所有的.
他們就跟從了耶穌.
是的 他們跟從了耶穌.
剛才我們說.
彼得也是一個很有影響力的人.
一個很有影響力的彼得.
他能夠帶領整艘船的人.
其實我們下班要回家.
見一下老婆子女的時候.
現在要我們再撐船出去.
再釣一次魚.
回來又要再洗那個岸.

$^{441}$多了很多功夫.
但是那些做他員工的.
都很聽他的.
他是一個很好的領袖.
除此之外.
他這一次的經歷.
一個不在行的耶穌.
教了他去捕魚.
而捕了很多魚的時候.
他發現了自己是一個罪人.
他認識自己多了.
而另一方面.
他對神 對主耶穌.
那個感覺不同了.
耶穌是很特別的.
祂不單只是一個老師.
彼得稱呼祂說.
你是主.
最後他放下所有去跟耶穌.
為什麼呢?.
他相信耶穌可以供應他.
我想在他.
他應該不少年紀.
起碼當他二,三十歲.
他也不是第一年去捕魚.
他也不是第一代去捕魚.
他可能三幾代都捕魚了.
他可能沒試過捕滿兩船的魚.
網打兩船的魚.
而沒有任何損失.
船好好的.
旁邊的船也好好的.
網也沒事.
他經歷了一些.
可能他覺得爸爸,爺爺.
在他的工作上.
都沒有經歷過這樣的神蹟.
這個時候他相信.
不是他可以供養自己多少.
他做多久有這一船的魚呢?.

$^{481}$可能打一整晚.
都未必有一船的.
半船都挺好的.
多一點還好一點.
但這一次是兩船的魚.
他明白到.
為什麼昨晚自己什麼都打不到.
而今天早上一網.
在一個不適當的環境.
天光之下.
可以網到這麼多魚.
不是因為什麼時機好運.
是因為耶穌.
因為耶穌叫他去網魚.
因為耶穌在船上.
耶穌給的.
如果靠他去養家.
他真的在養.
但不是時時都一定.
有這麼好的成果.
物資從哪裡來呢?.
誰能夠令他一定養到家.
養到這班員工呢?.
可能是耶穌.
因為他一句話.
因為聽他一句話.
接著就有這樣的結果.
彼得他明白多了.
他放下所有的.
因為他不需要擔心.
其實開船出去的時候.
也未必不擔心.
有什麼擔心?.
有沒有人出過海網魚?.
出過海網魚的話.
出海都擔心風浪.
他們不是沒遇過風浪.
格陵薩勒湖經常有風浪.
風浪不喜由是可.
他一氣就出了.

$^{521}$所有脾氣都出了.
是很難去控制的.
除了安全之外.
是不是魚在湖上.
一定要被捕捉到?.
他都不知多想享受一下自由.
享受一下他的生命.
為什麼要做你今晚.
碟子上的那隻魚?.
是耶穌給他的.
他很明白是誰給他的.
恩典是從哪裡來的.
他發現了一件事.
我只是一個罪人.
耶穌你和我一起.
影低了你.
我是不配的.
但是當耶穌叫他跟隨的時候.
他又跟了他.
因為他明白.
供應是從哪裡來的.
他接受了耶穌基督的呼召.
似乎有一群人跟著他.
這群人是不是經常跟著他們.
出去傳福音呢?不知道.
不過久不久他們會出現的.
就是這樣.
這四個人.
彼得,安德烈,雅各,約翰.
跟著耶穌.
跟著了耶穌.
他們做了什麼事呢?.
我想他們做了一件事.
就是那天早上.
聽到他們認識了耶穌多一點.
不過是駕駛船去湖中心的時候.
去到水深之處的時候.
他們真的認識了耶穌.
他們認識了耶穌.
是他們可以跟的對象.

$^{561}$他們真的相信耶穌.
就是要來的彌賽亞.
不過還有一件事.
他們認識了自己.
彼得跟耶穌說.
他說:主啊.
你離開我遠一點.
我是個罪人.
但你看到一個是主的時候.
就是他是主的.
他是主就應該是他決定的.
他站在哪裡.
和你站在哪裡都不是自己決定的.
是他決定的.
但是彼得就說.
你走遠一點.
我不配跟你一起.
這個是.
你有沒有這樣的經歷啊.
有沒有一些時候.
你有這種的.
很有矛盾.
又很有感受的反應啊.
可能有些時候.
魯牧師或曾姑娘.
找你去參與一些事奉的時候.
你又會覺得.
是啊 你找我.
我都真的很開心的.
魯牧師叫我做事.
曾姑娘叫我做事.
不知道多開心.
不過我不行的.
你找別人吧 曾姑娘.
我做不到.
我做得不好的.
你叫別人吧.
可能我們都會有這樣的反應.
可能我們都會覺得.
有機會事奉神是很好的.

$^{601}$但是我呢.
我都可以的嗎.
是不是真的我都可以的.
耶穌跟彼得說.
是啊 就是你啊.
就是你要去得人啊.
可能彼得他對自己的認識就是.
我隨時亂說話的.
我突然之間衝一句話出來的.
是不是 大家都認識彼得.
我適不適合去做一個得人的人啊.
阿弟姐妹.
無論你是哪個人.
什麼人.
如果曾姑娘找你.
魯牧師找你.
請你幫忙一些事奉的話.
他們認識你的.
你要相信他們真的認識你的.
是你的牧者.
他們找你幫忙一些事奉的時候.
他們覺得你很適合那個位置.
也就像主耶穌基督找彼得一樣.
彼得做到多少呢.
是主耶穌基督一直都會幫助他.
魯牧師和曾姑娘.
也都會在背後.
一直找你事奉的時候.
會給你訓練.
會支持你.
會關心你做得怎麼樣.
會陪同你一起去做.
「紅因」.
我們剛才唱第一首詩歌.
就好像很特別.
那首詩歌好像是為「紅因」而寫的.
裡面有「紅因」兩個字.
「紅因」是為誰而存在呢?.
「紅因」也是為神而存在.
也是為在座每一位而存在.

$^{641}$也是讓每一位在座的.
可以為「紅因」得人而存在.
在你的生活裡面.
可能你還有很多要牽掛的事.
就像彼得一樣.
彼得要牽掛他的船.
牽掛他的父母.
他的岳母.
他的太太,妻兒.
還有他的十幾個員工.
他不是要找一個人的生活費.
他應該要找十幾個家庭的生活費.
做老闆的你知道.
你要不要繼續做生意.
你的生意,你的公司要怎樣去處理.
不單止是一個人的事.
是你每一個員工都受到影響.
是的.
彼得看到.
他相信耶穌可以照料他們這麼多人.
當神藉著老牧師.
藉著曾姑娘找你去事奉的時候.
他相信你能夠在這裡付出你的能力.
也可以像彼得一樣.
放下他的牽掛.
還有他覺得的生活需要.
他知道神會幫助他.
讓他在家庭上可以平衡得到.
讓他在生活上有神來照料.
是的,原來是可以這樣.
我們可以有些東西.
天父幫助我們去處理.
不過我們有一樣東西要認得.
就是認得基督是誰呢?.
基督是主.
彼得提我們基督是主.
我們通常做什麼呢?.
不知道大家有沒有犯過這些小毛病.
我也會犯過這些毛病.
就是什麼呢?.

$^{681}$我們一開始祈禱.
很敬重地稱呼神之後.
我們就說了很多代祖事項.
然後就說天父主耶穌基督聖靈.
我這十幾項的餐單就交給你了.
你有空就幫我做吧.
這樣就能夠押門了.
其實誰是主呢?.
誰是僕人呢?.
誰是主誰是僕人呢?.
有時候我們真的很心急地去做一些事的時候.
我們有時候把主和僕人調轉了.
彼得就認得.
耶穌基督是主.
他是回應耶穌基督的呼召.
去做耶穌基督的僕人.
所以有些事原來不是我們去決定的.
彼得是不是繼續要去做漁夫.
不應該是彼得決定的.
雖然耶穌是很客氣地說來吧.
你跟著我去做得人的工作吧.
耶穌沒有鎖著他.
沒有拉著他要走.
很客氣地說.
但是耶穌才是主.
是耶穌去決定彼得應該做什麼.
然後他應該去什麼地方得人.
都是耶穌去決定.
什麼時候去哪裡.
是耶穌基督會下命令的.
我好像告訴他.
現在撐離岸一點.
就說到.
然後撐到水深之處就得魚.
不是這位捕魚的專家彼得的決定.
而是主耶穌基督的決定.
因為他是主.
他才是主.
我們是僕人.
讓我們藉著禱告.

$^{721}$我們都告訴他.
他是可以差遣我們.
他是可以藉著教會的教務傳道.
讓我們在教會當中.
我們有份去參與.
一起建立雄恩.
我們一起低頭祈禱.
參見天父上帝.
你為我們預備救恩.
你讓我們能夠得救.
主耶穌基督是我們的主.
我們要跟隨祂.
我們是你的僕人.
是基督的僕人.
是聖靈的僕人.
我們需要你對我們的差遣.
你對我們的鼓勵.
讓我們在你預備給我們的工作上.
我們能夠有一個正確的安排.
在你預備給我們的工作上去努力.
天父有時候我們真的有種矛盾.
我們很想好好事奉你.
但有時候我們很有自己的意思.
有時候我們又很不敢走上去.
求主你恩代我們幫助我們.
感動我們.
藉著教會的傳道人牧師來鼓勵我們.
讓我們能夠勇敢地在紅恩堂中.
一起辛苦 一起學習 一起努力.
做你的僕人.
主願意你紀念看顧我們.
賜恩典給我們 指引我們.
賜能力 賜恩賜給我們.
去做你的僕人.
我們祈禱仰望你奉主耶穌基督的聖名.
阿們.
\newpage



\section{列王記上 2:1-9}
\label{sec:BYuafMUFtfY}
\textbf{2021.10.03 《活在神的應許中》 列王記上 2\_1-9 (和修版) 老耀雄牧師}
\newline
\newline
連結: \href{https://youtube.com/watch?v=BYuafMUFtfY}{\texttt{https://youtube.com/watch?v=BYuafMUFtfY}} ~~~~ 語音日期: 2021-10-08
\newline
\newline
\hyperref[sec:_ayvrCyPl7U]{\small{< < < PREV SERMON < < <}}
~
\hyperref[sec:index_chronic]{\small{[返順時目]}}
~
\hyperref[sec:index_scriptual]{\small{[返順卷目]}}
~
\hyperref[sec:P_diqejYp0w]{\small{> > > NEXT SERMON > > >}}
\newline
\newline
列王記上 2:1-9
\newline
\begin{longtable}{cl}
\hline
\hline
章節 & 經文 (和合本修訂版)\\
\hline
2:1 & \begin{tabularx}{0.7\textwidth}{X} 大衛的死期臨近了，就吩咐他兒子所羅門說： \end{tabularx} \\ \\ \relax
2:2 & \begin{tabularx}{0.7\textwidth}{X} 「我要走世人必走的路了。你當剛強，作大丈夫， \end{tabularx} \\ \\ \relax
2:3 & \begin{tabularx}{0.7\textwidth}{X} 遵守耶和華－你神所吩咐的，照著摩西律法上所寫的行耶和華的道，謹守他的律例、誡命、典章、法度，好讓你無論做甚麼，不拘往何處去，盡都亨通。 \end{tabularx} \\ \\ \relax
2:4 & \begin{tabularx}{0.7\textwidth}{X} 耶和華必成就他所說關於我的話，說：『你的子孫若謹慎自己的行為，盡心盡意憑信實行在我面前，就不斷有人坐以色列的王位。』 \end{tabularx} \\ \\ \relax
2:5 & \begin{tabularx}{0.7\textwidth}{X} 你也知道洗魯雅的兒子約押向我所做的事，他對付以色列的兩個元帥，尼珥的兒子押尼珥和益帖的兒子亞瑪撒，殺了他們。他在太平之時，如同戰爭一般，流這二人的血，把這戰爭的血染了他腰間束的帶和腳上穿的鞋。 \end{tabularx} \\ \\ \relax
2:6 & \begin{tabularx}{0.7\textwidth}{X} 所以你要照你的智慧去做，不讓他白髮安然下陰間。 \end{tabularx} \\ \\ \relax
2:7 & \begin{tabularx}{0.7\textwidth}{X} 你當恩待基列人巴西萊的眾兒子，請他們常與你同席吃飯，因為我躲避你哥哥押沙龍的時候，他們親近我。 \end{tabularx} \\ \\ \relax
2:8 & \begin{tabularx}{0.7\textwidth}{X} 看哪，在你這裡有來自巴戶琳的便雅憫人，基拉的兒子示每。我到瑪哈念去的那日，他用狠毒的言語咒罵我。後來他卻下約旦河迎接我，我就指著耶和華向他起誓說：『我必不用刀殺死你。』 \end{tabularx} \\ \\ \relax
2:9 & \begin{tabularx}{0.7\textwidth}{X} 但現在你不要以他為無罪。你是有智慧的人，必知道怎樣待他，使他白髮流血下陰間。」 \end{tabularx} \\ \\
[1ex]
\hline
\hline
\end{longtable}
$^{1}$有人說:人之將死,其言也善..
當然,人的生命已經去到盡頭,.
亦毋須說謊..
不過,對於一個重要領袖去臨終遺言,.
不單說的是要真,亦要有善意..
因為作為一個領袖,.
需要向所帶領的群體.
說出一個未來的展望和預測,.
讓下一代能有一個清晰的藍圖,.
作為對下一代的指引和發展方向..
在聖經裡面有很多重要人物的遺言,.
往往都是涉及未來的重要提示..
如雅各的遺言,.
涉及以色列十二支派的未來;.
摩西的遺言,.
涉及以色列人入迦南的未來;.
約書亞的遺言,.
涉及十二支派在迦南分地後的未來;.
撒母耳的遺言,.
涉及以色列立王後的未來;.
大衛的遺言,.
涉及大衛王朝的承傳未來..
今日很想透過《列王記》上二章一至九節,.
從大衛向所羅門的遺言中,.
和大家探討兩個信仰的重點,.
我們一起去祈禱..
作為我們生命的陀,.
亦有聖靈指引著我們生命的方向,.
你的憐憫,保守,.
叫我們不會偏離航道,.
你的提醒,.
亦都讓我們能夠按著你的心意而行,.
求你繼續施恩典,憐憫,.
好叫我們的屬靈生命,.
每日都能夠成長,.
被你使用,榮耀主名..
我們在祈禱,奉主耶穌聖名祈求..
阿們..
在大衛對所羅門的遺訓中,.
大致分為兩個部分,.

$^{41}$在第二至第四節中,.
內容與神和以色列王國有關,.
第五至第九節內容與人有關..
我們先看看大衛遺言的第一部分,.
在第一部分,.
是說遵守神的話語,.
作為承傳給下一代的教導,.
我們再看看第二至第三節,.
「我要走世人必走的路,.
你當剛強,作大丈夫.」.
遵守耶和華,.
你神所吩咐的,.
照著摩西律法上所寫的.
行耶和華道,.
謹守祂的律例,誡命,典章,法道,.
好讓你無論做什麼,.
不驅往何處去,盡都亨通..
我將這裡的內容隱藏起來,.
這兩節經文的重點,.
其實與約書亞一章七節,.
摩西勸勉約書亞的三個重點,.
差不多完全一樣..
在約書亞一章七節,.
他說:只要剛強,大大壯膽,.
謹守遵行我僕人摩西所吩咐你的一切律法,.
不可以偏離左右,.
使你無論往那裡去,都可以順利..
第一重點是剛強,.
剛強作大丈夫..
在聶皇記上和約書亞裡面的.
「剛強」這個希伯來文,.
不只有剛強的一個意思,.
是有兩種的向度,.
包括了力量和勇氣..
力量是指外在的力,.
勇氣是指內在的能力,.
是一種由內至外面對世界的力量..
第二個重點,.
是要謹守耶和華向摩西頒布的.
《律例誡命典章法道》,.

$^{81}$這裡其實是指摩西的律法..
第三個重點,亦是結果,.
就是不拘往何處去,盡刀亨通..
簡單來說,.
就是在神的律法或真道裡面,.
要有力量和勇氣,.
意念與行為一致,.
這樣去踐行出來,.
這樣的話,.
無論去到哪裡,.
都會一切順利..
這不單單是大衛對所羅門的遺訓,.
其實亦都是摩西向約書亞所講的,.
不單單是這樣,.
其實還可以追溯到.
摩西與約書亞的關係,.
還可以追溯到神與以色列的關係..
在《申命記》30章16至20節,.
摩西怎樣說呢?.
他說:.
「我將生死禍福來形容神與以色列的關係,.
如果揀選耶和華,.
遵守他的誡命律例短章,.
耶和華就應許被列祖的地賜給以色列人.」.
所以大衛遺訓這一點,.
一開始就是要所羅門記得.
神對整個以色列民族的要求,.
不過第二至三節所涉及的,.
都只是神對整個以色列民族的應許,.
沒有涉及到神對大衛王朝的應許..
所以大衛在第四節,.
就提到神對他自己的應許..
在第四節他說:.
「你的子孫若謹慎自己的行為,.
盡心盡義,憑信實,.
行在我面前,.
就不斷有人在以色列的王位.」.
原來耶和華曾經對以色列的王朝.
有過這樣的應許..
在最後一節他說:.

$^{121}$「這書要存在他那裡,.
他一生的年日要重讀,.
好使他學習敬畏耶和華他的神,.
謹守遵行這律法書上的一切話,.
和這些律例,.
免得他的心向弟兄高傲,.
偏離了這個誡命,.
或向右或向左,.
這樣他和他的子孫.
就可以長久作王,治理以色列.」.
不過耶和華對大衛王朝最直接的應許.
不是這一段,.
而是在撒母耳記下的七章十二至十六節..
神透過先知拿丹.
對大衛直接說出了這一段說話:.
「當你受歲滿足,.
與你祖先同歲的時候,.
我必使你身所生的後裔接觸你,.
我也必堅定他的國,.
他必為我的名建造奠宇,.
我必堅定他國道的王位,.
直到永遠.」.
因此大衛對所羅門的遺訓.
其實是根據神對以色列的列祖的應許,.
對以色列這個民族的應許,.
以及以色列王朝的應許,.
當中包括了個人,.
包括了民族,.
以及包括了大衛王朝..
因此神所頒佈的律法.
不單單只是針對以色列的人民,.
也針對了大衛王朝以及他的繼承者..
對於大衛來說,.
他最關心的是什麼呢?.
就是整個民族以及王朝的承傳..
以色列這個民族能不能夠承傳下去呢?.
大衛王朝能不能延續下去呢?.
關鍵是他們所做的能不能夠遵守耶和華的律法..
同一家弟兄姊妹,.
大衛的遺訓的前半部對於我們有什麼反省?.

$^{161}$大衛心裡面所繫的最重要的,.
固然是以色列民族將來是否能夠盡刀亨通呢?.
以及王朝是否能夠延續呢?.
這個涉及大衛的承繼人和耶和華的關係..
一個與耶和華有親密關係的人,.
才能夠遵守耶和華的誡命律例典章..
一個能夠遵守神律法的王,.
他才能夠帶領著整個以色列民族.
去遵守耶和華的律法..
而整個民族能夠遵守耶和華的律法群體,.
才能夠得到神的祝福..
不只是王去遵守神的律法,.
而是百姓都需要遵守神的律法..
只有百姓遵守神的律法,.
王不遵守,不可以..
只有王遵守律法,百姓不去遵守,都不可以..
如果你說這段遺訓對我們洪恩家.
可以有什麼反省?.
我就會想,.
洪恩家的墓者,.
他的屬靈生命是不是聖潔?.
墓者能不能夠肯遵行神的心意?.
這是決定了洪恩家的發展息息相關,.
沒有一個聖潔的生命,.
神不喜悅的..
如果墓者偏行己路,.
只會走上一言堂,.
或者獨裁的路..
反過來,.
墓者如果沒有了會眾的支持,.
教會都只是孤軍作戰,.
失敗是必然,.
因為教會不會只有一個人去做..
洪恩家如果要盡刀亨通,.
關鍵是墓者和全體的會眾.
都和神有一個好的關係..
也就是說,.
墓者和會友的屬靈生命.
都要貼近神的心意..
這反映到墓者和會眾.

$^{201}$能不能夠在這個教會裡,.
毀身服侍..
弟兄姊妹,.
如果今天我們輕視神的話語,.
我們又不去實踐,.
我們的屬靈生命.
怎麼能夠去見證神?.
下一代又怎麼能夠從我們身上.
去有什麼效法的機會?.
或者用另一個說法,.
我們要問問自己,.
我們想要成全一個什麼信仰.
給我們下一代?.
是一個充滿信望,愛的信仰,.
還是一個可有可無的信仰?.
這個下一代.
可以包括我們血肉相連的子女,.
這個下一代.
也可以是我們信仰的下一代,.
我們教會的第二代,第三代,第四代..
我們想教會.
成全一個什麼信仰給下一代?.
信望,愛的信仰.
就不是只有一個「港」字,.
要人看到這間教會.
真的有信望,愛的果實,.
才能夠成全..
還是一種可有可無的信仰?.
什麼叫做可有可無的信仰?.
你回來也可以,.
不回來也可以,.
事奉也可以,.
不事奉也可以,.
奉獻也可以,.
不奉獻也可以..
這種就是可有可無的信仰..
今日的講題是《活在神的應許》裡面,.
大衛對所羅門臨終的遺訓.
就是要他遵守耶和華的《誡命律例典章》.
即是摩西的律法.

$^{241}$以至一切盡都亨通.
盡都亨通.
就是耶和華對以色列人的應許.
也是耶和華對大衛的應許.
大衛很想所羅門繼續活在神的應許裡面.
今年是洪恩嘉踏入十週年的時刻.
能不能有一個新的開始?.
十年重新起步.
重新立志.
重新出發.
牧者很期待每一位信徒.
包括我自己.
都能夠活在這個應許裡面.
而能夠活在這個應許裡面.
關鍵的條件.
就是一切.
一切去遵行神的《誡命律例典章》.
剛才主席在祈禱裡面說到.
我們能不能夠一切.
盡心,盡意,盡性,盡力.
愛我們的主.
不是可有可無.
不是那種可有可無的信仰.
而是盡心,盡意,盡性,盡力.
愛我們的主.
在第二部分.
不是以人的智慧取代神的應許.
大衛的遺訓下半部.
是關乎三段與人的關係.
一段是知恩.
兩段是念神.
先說知恩的部分.
在第七節.
「你當恩代基列人巴西萊的眾兒子.
請他們常與你同席吃飯.
因為我躲避你哥哥阿沙隆的時候.
他們親近我」.
事件在阿沙隆發動叛變期間發生.
由於當時很多人都歸向阿沙隆.
所以大部份人都不看好大衛.

$^{281}$一般人都是一尋百踩.
有多少人會雪中送炭呢?.
巴西萊是其中一個雪中送炭的.
他帶著被褥,盤,瓦器.
還有小麥,大麥,麥餅.
烘熱的穀碎.
豆子,紅豆,炒豆,蜂蜜.
奶油,綿羊,奶羊,奶餅.
供給大衛和跟隨他的人吃.
因為他們想.
百姓在抗疫中必定又飢餓又疲乏.
這就是雪中送炭.
巴西萊在大衛最需要支援時提供糧食.
並且加上關心.
大衛在平息叛變後.
回到耶路撒冷途中.
曾經對巴西萊作出一個承諾.
王對巴西萊說.
「你跟我一起渡過去吧」.
「我要在耶路撒冷身邊奉養你」.
這就是大衛對他的邀請.
巴西萊如何回答呢?.
他說「僕人送王過約旦河只是小事一樁」.
「王何必用賞賜來報答我呢?」.
「請讓我回去死在我本城」.
「葬在我父母的墓旁」.
「於是眾百姓過了約旦河」.
「王也過去了」.
「王親吻巴西萊為他祝福」.
「巴西萊就回自己的地方去了」.
巴西萊沒有應約.
只是回家就可以了.
小事一樁.
施恩莫忘部.
所以大衛要求所羅門好好善待巴西萊.
和他的眾子.
作為報答巴西萊當日雪中送炭之恩.
報恩是大衛想所羅門為他做的其中一件事.
作為基督徒.
我們如何報主耶穌對我們的恩.

$^{321}$主耶穌說「小事一樁」.
真的小事一樁?.
主耶穌在十字架上為我們死.
對於兩件念仇的事件.
第一件在第五節.
「你也知道.
在路亞的兒子約翰向我所做的事.
他對付以色列的兩個元帥.
尼爾的兒子阿尼爾和益帖的兒子阿瑪撒.
殺了他們.
他在太平之時如同戰爭一樣.
流著兩人的血.
把這戰爭的血染了他腰間.
足的大趾和腳上穿的鞋」.
約翰是誰?.
約翰是大衛的元帥.
亦是大衛的外甥.
阿尼爾是誰?.
阿尼爾是當時的素羅軍隊的元帥.
即是說兩個元帥是死對頭.
是素羅的叔尼爾的兒子.
當時大衛和素羅在交戰.
所以約翰是大衛的元帥.
阿尼爾就是素羅的元帥.
阿尼爾在素羅死後.
繼續服侍素羅的兒子伊斯波切.
戰爭本來就會一觸即發.
不過彼此透過外交手段.
阿尼爾後來和大衛講和.
並且準備帶領以色列的民眾歸降.
去投靠大衛.
如果是這樣的話.
就避免了一場內戰.
接著就發生了這件事.
阿尼爾回到希伯倫.
約翰領他到城門中間.
要與他私下交談.
就在那裡刺穿了他的土腹.
他就死了.
因為他流了約翰兄弟阿薩克的血.

$^{361}$對於以色列第一個王素羅的元帥.
阿尼爾的歸降.
對於約翰的地位有相當的威脅.
什麼叫一山不能藏二虎.
就是這個意思.
因為他極有可能取代了自己的地位.
因此約翰殺死阿尼爾的用意.
就是防止對方取代他的地位.
也為他自己的兄弟阿薩克報仇.
所以經文說.
約翰是在太平的時候殺死阿尼爾.
而不是在雙方交戰的戰場殺死他.
在第五節也提到另一個元帥亞瑪薩.
亞瑪薩是誰呢?.
亞瑪薩就是大衛的兒子.
阿沙隆的元帥.
當阿沙隆叛變的時候.
就立了亞瑪薩做元帥.
後來兵變事敗.
阿沙隆身陷險境.
約翰沒有聽從大衛的命令.
施治將阿沙隆殺了.
大衛在平息了阿沙隆的叛變之後.
曾經在眾人面前預告.
將會立亞瑪薩做元帥.
取代約翰.
約翰知道自己地位不保.
然後又發生了這件事.
亞瑪薩沒有防備.
約翰手裡拿著一把刀.
約翰用刀刺入他的肚腹.
他的腸子留在地上.
約翰沒有再刺.
他就死了.
對於這兩個元帥的品格.
在《列王記後》《灑海經文》.
有這樣的描述.
耶和華必使約翰的血.
歸到他自己頭上.
因為他擊殺兩個.

$^{401}$比他又公義又善良的人.
就是尼爾的兒子.
以色列的元帥阿尼爾.
和益帖的兒子猶大的元帥亞瑪薩.
用刀殺了他們.
我父親大衛卻不知道.
如何形容這兩個元帥.
又公義又善良.
而他殺了他們的時候.
大衛說不知道.
經盟約這兩個元帥稱為.
又公義又善良.
約翰就在.
阿尼爾和亞瑪薩歸降之後.
才殺了他們.
約翰之所以去殺.
阿尼爾以及亞瑪薩.
都是以自己的利益出發的.
因為兩次事件.
都不在一個戰爭狀態之下.
蘇羅和阿沙隆已經死了.
戰事也已經結束了.
雙方都在一個外交的交接裡面.
不過約翰就為了自己的地位.
殺害了兩個以色列的元帥.
其實在戰爭的過程裡面.
這些事都是很常見的.
大衛在位的時候.
他沒有親自去懲罰約翰.
為何臨終的時候.
才吩咐所羅門去處理呢?.
有人說大衛不想.
再犯上一個殺人的罪名.
不過不像的.
因為大衛不像一個要避開殺人的罪名.
事實上他以前也曾經因為犯下殺人的罪名.
就好像殺害烏尼亞一樣.
先知拿丹已經宣判了大衛殺人的罪名.
我們中國人很喜歡說一次污兩次衛.
殺了一個不再殺第二個.

$^{441}$大衛要所羅門處理約翰.
其實是想幫所羅門鞏固他自己的王位.
在第一章所羅門的哥哥亞多尼亞.
已經密謀篡位了.
而作為大衛的元帥.
當時的約翰是支持亞多尼亞的叛變.
約翰是一個很能幹的軍人.
而且幫了大衛這麼久.
在軍隊裡面有一定的影響力.
如果他要叛變.
必定會為所羅門帶來軍事上的不穩定.
所以大衛就要所羅門除去這個障礙.
大衛用自己的智慧去幫助所羅門鞏固王位.
就是以人的智慧去取代約翰的保守.
與他遺言的第一部分的精神完全相反.
另外一段念仇的事件.
記載在第八節.
「看!在你那裡有來自巴胡林的汴雅文人.
基拉的兒子氏妹.
我到馬哈念去的那一天.
她用狠毒的言語咒罵我.
後來她卻下若淡河迎接我.
我便指著向她起誓說.
我必不用刀殺死你」.
氏妹是甚麼人?.
氏妹是素羅家族的一個僕人.
在阿沙隆叛變期間.
大衛逃走了.
去到馬哈念的時候.
大衛遇到氏妹.
那氏妹怎樣對待他呢?.
他一面走一面咒罵.
又向大衛王和王的眾臣掟石頭.
眾百姓和勇士都在王的左右.
氏妹這樣詛咒他.
「你這個好流人血的無賴.
走滾吧滾吧.
你流了素羅全家的血.
接續他作王.
耶和華把這罪歸在你身上.

$^{481}$耶和華將這國交在你兒子阿沙隆的手中.
看!你夠由自取.
因為你是好流人血的人」.
一塵百踩就是這些人.
氏妹就是那種一塵百踩的人.
當時大委身邊的將士看不過眼.
就想去割了氏妹的人頭.
大衛怎樣呢?.
大衛就這樣說.
他說「我親生兒子都想要我的命.
何況現在這個汴南人.
由得他罵吧.
因為這是耶和華吩咐他的」.
大衛都相當豁達.
後來事件平息了.
阿沙隆的叛變之後.
大衛又經過這個地方.
回到耶路撒冷的時候.
變色龍就是變色龍.
氏妹又再走出來.
這次不是詛咒.
而是向大衛賠罪.
他說「我主王離開耶路撒冷那天.
僕人行了半億的事.
現在求我主不要因為此加罪於僕人.
不要記得也不要放在心上.
僕人明知自己有罪看.
約瑟全家之中.
今日我首先下來迎接我主王.
認罪」.
於是王怎樣說呢?.
當日大衛就對氏妹說.
「你必不死」.
這是王對她的喜誓.
當時大衛向氏妹喜誓.
承諾不會殺她.
為何到現在臨死的時候.
就要跟她計算得這麼清楚呢?.
究竟大衛當日的寬恕是真情還是假意?.
大衛要對付氏妹.

$^{521}$其實都是為了王朝的鞏固.
因為氏妹是素羅家族一個很大的影響力.
如果氏妹都叛變的話.
是一個很大的威脅.
殺了氏妹.
就是一種殺雞儆猴的方式.
告訴素羅這批變雅敏之派的人.
凡是親素羅反大衛的.
必不寬容.
因此無論他對付約瑟.
還是對付氏妹.
大衛要等到臨終的時候才處理他們.
最大的理由.
就是為了他的兒子所羅門的王位.
清除一切的障礙.
大衛想怎樣處理這兩段恩怨呢?.
對於約瑟.
他說:所以你要照你的智慧去做.
不讓他白髮安然下陰間.
對於氏妹.
但現在你不要以他為無罪.
你是有智慧的人.
必知道怎樣待他.
使他白髮流血下陰間.
不讓他白髮安然下陰間.
使他白髮流血下陰間.
其實都是同一個意思.
就是叫他們不得好死.
大衛提醒所羅門.
要用智慧去處理.
什麼叫智慧?.
在這裡說的智慧是什麼智慧?.
大衛沒有說.
只能夠意會.
即是今天說的「執生」.
但結果都是一樣.
兩個被提名的人.
最終都是要死.
不過諷刺的是.
大衛在第三章.

$^{561}$真的向耶和華求智慧.
怎樣說呢?.
神回應所羅門的時候.
和大衛所想的有一個很大的對比.
神對他說.
你既然求這件事.
即是求智慧這件事.
不是自己求壽.
不是自己求富.
也不求滅絕你仇敵的性命.
大衛是否要求什麼?.
大衛要求滅絕.
所羅門求的智慧.
不是要求滅絕.
只是求明辨可以聽從.
神回應的其中一點.
不求滅絕你仇敵的性命.
這種做人做事的智慧.
和大衛要求所羅門所做的事情.
是完全相反.
是截然不同的做法.
大衛要求所羅門滅絕虐壓.
滅絕這個仕妹.
他用這種智慧去鞏固這個王朝.
神回應所羅門.
你不求滅絕你仇敵的性命的智慧.
我會賜給你.
神給所羅門的智慧.
不是一個復仇報復的智慧.
也不是一個鞏固王朝的智慧.
而是一個明辨是非的智慧.
明辨是非.
如果今天讓我們去求.
我們要求一個什麼智慧呢?.
如果我們沒有信仰.
當然是求智慧.
看我們怎樣可以賺錢.
炒兩炒.
或者怎樣可以極速上位.
做CEO,CFO.

$^{601}$一個能夠明辨是非的智慧.
對於一個信徒來說有什麼幫助?.
如果我們今天去求智慧.
求一個明辨是非的智慧.
有什麼幫助?.
當大衛用他的智慧.
取代神的保守和應許的時候.
我們就明白.
神所賜給所羅門明辨是非的智慧.
就是要分辨什麼是神所喜悅.
什麼是神所厭惡.
以至所羅門能夠走在神的心意當中.
如果我們是一個信徒.
我們今天說.
向神求智慧.
這個智慧就是怎樣去明白神的心意.
神喜歡什麼.
神不喜歡什麼.
紅映加的弟兄姊妹.
在大衛的遺訓裡面.
我們看到兩個部分.
第一個部分要求所羅門遵守神的律法.
顯示大衛在列祖和整個以色列民族裡面.
和耶和華的關係出發.
這個關係就是要所羅門.
謹守著耶和華的誡命律例典章.
需要按著神的心意而行.
也是以色列列祖所承傳下來的.
今天紅映加的會友.
能否像按照耶穌和大公教會的傳統一樣.
又像主席剛才的祈禱一樣.
我們能夠盡心盡性.
盡意盡力愛主理神.
其次就是愛人如己.
去踐行我們的信仰.
說是很容易的.
做不是那麼容易.
第二部分.
大衛要求所羅門紀念巴西來的眾子.
以及要清除虐壓和煽誤.

$^{641}$雖然是為了幫助所羅門去鞏固王朝.
以及提高王位的穩定性.
不過大衛使用的是人間的智慧.
而不是依靠神的應許.
耶和華既然應許過大衛.
只要你遵守耶和華的律法.
他就會讓你的子孫世世代代做王.
盡刀亨通.
但是大衛沒有完全依靠這個應許.
仍然以自己所想的.
為所羅門去鞏固這個王位.
今天.
神對我們每一個信徒都應許過.
無論我們處於一個什麼的光景.
神都與我們同在.
無論我們在什麼的逆境裡面.
我們都可以呼求主.
去幫我們面對這個困難.
為什麼我們仍然用我們的智慧.
聰明去面對我們的逆境呢?.
這個有時候我也覺得.
我們的信仰不能夠深化在我們的生命裡面.
被信主之前.
我們每一個人都是靠著自己的.
努力,關係,智慧.
去面對我們的困難.
很正常而已.
不過今天我們信了主.
為什麼仍然用舊的模式.
去面對我們的信仰生活呢?.
我們為什麼不能夠交托.
或者為什麼不能夠依靠主呢?.
是不是我們的信仰沒有深化在我們的生命裡面?.
在大衛遺訓裡面.
給我一種反思.
很多時候人們說大衛是神最合適的僕人.
或者大衛是神最合心意的僕人.
我相信在這份遺訓的下半部.
就加入了一些遺憾.
成為一個很不完美的結局.

$^{681}$聖經沒有為大衛的一生粉飾大平.
他有功亦有過.
他有得到耶和華喜悅的時候.
但也有得罪神的時刻.
一生的功過.
由神去評價.
或者用另一個說法.
人生不盡是完美.
亦都不需要追求完美.
不過在一個不完美的人生裡面.
我們能不能夠求神賜給我們一個明辨是非的心.
去測驗神所喜悅.
或者神所厭惡的事.
能夠一心一意.
為到神所喜悅的事.
而為主擺上呢?.
我們一起去祈禱.
因為你對我們每一個信徒都有應許.
求主你讓我們真的活在你的應許裡面.
而不是活在自己的聰明才智上.
因為你的應許.
是讓我們能夠有一個豐盛的人生.
你的應許.
是說過與我們同行.
與我們同走這個信仰的路.
甚至我們真的能夠珍惜.
主你給我們每一個信徒的應許.
讓我們每一個都能夠在這個應許裡面.
活出一個豐盛的人生.
我們祈禱仰望.
奉靠主耶穌.
得勝名自祈求.
阿們.
願主祝福大家.
\newpage



\section{哥林多前書 12:12-27}
\label{sec:P_diqejYp0w}
\textbf{2021.10.10 《蒙神喜悅的事奉》哥林多前書 12\_12-27 曾敏姑娘}
\newline
\newline
連結: \href{https://youtube.com/watch?v=P-diqejYp0w}{\texttt{https://youtube.com/watch?v=P-diqejYp0w}} ~~~~ 語音日期: 2021-10-12
\newline
\newline
\hyperref[sec:BYuafMUFtfY]{\small{< < < PREV SERMON < < <}}
~
\hyperref[sec:index_chronic]{\small{[返順時目]}}
~
\hyperref[sec:index_scriptual]{\small{[返順卷目]}}
~
\hyperref[sec:9yhz2ddk1BQ]{\small{> > > NEXT SERMON > > >}}
\newline
\newline
哥林多前書 12:12-27
\newline
\begin{longtable}{cl}
\hline
\hline
章節 & 經文 (和合本修訂版)\\
\hline
12:12 & \begin{tabularx}{0.7\textwidth}{X} 就如身體是一個，卻有許多肢體，身體的肢體雖多，仍是一個身體；基督也是這樣。 \end{tabularx} \\ \\ \relax
12:13 & \begin{tabularx}{0.7\textwidth}{X} 我們無論是猶太人是希臘人，是為奴的是自主的，都從一位聖靈受洗成了一個身體，並且共享這位聖靈。 \end{tabularx} \\ \\ \relax
12:14 & \begin{tabularx}{0.7\textwidth}{X} 身體原不只是一個肢體，而是許多肢體。 \end{tabularx} \\ \\ \relax
12:15 & \begin{tabularx}{0.7\textwidth}{X} 假如腳說：「我不是手，所以不屬於身體」，它不能因此就不屬於身體。 \end{tabularx} \\ \\ \relax
12:16 & \begin{tabularx}{0.7\textwidth}{X} 假如耳朵說：「我不是眼睛，所以不屬於身體」，它也不能因此就不屬於身體。 \end{tabularx} \\ \\ \relax
12:17 & \begin{tabularx}{0.7\textwidth}{X} 假如全身是眼睛，聽覺在哪裡呢？假如全身是耳朵，嗅覺在哪裡呢？ \end{tabularx} \\ \\ \relax
12:18 & \begin{tabularx}{0.7\textwidth}{X} 但現在神隨自己的意思把肢體一一安置在身體上了。 \end{tabularx} \\ \\ \relax
12:19 & \begin{tabularx}{0.7\textwidth}{X} 假如全都是一個肢體，身體在哪裡呢？ \end{tabularx} \\ \\ \relax
12:20 & \begin{tabularx}{0.7\textwidth}{X} 但現在肢體雖多，身體還是一個。 \end{tabularx} \\ \\ \relax
12:21 & \begin{tabularx}{0.7\textwidth}{X} 眼睛不能對手說：「我用不著你。」頭也不能對腳說：「我用不著你。」 \end{tabularx} \\ \\ \relax
12:22 & \begin{tabularx}{0.7\textwidth}{X} 不但如此，身上的肢體，人以為軟弱的，更是不可缺少的； \end{tabularx} \\ \\ \relax
12:23 & \begin{tabularx}{0.7\textwidth}{X} 身上的肢體，我們認為不體面的，越發給它加上體面；我們不雅觀的，越發裝飾得雅觀。 \end{tabularx} \\ \\ \relax
12:24 & \begin{tabularx}{0.7\textwidth}{X} 我們雅觀的肢體自然用不著裝飾；但神配搭這身子，把加倍的體面給那有缺欠的肢體， \end{tabularx} \\ \\ \relax
12:25 & \begin{tabularx}{0.7\textwidth}{X} 免得身體不協調，總要肢體彼此照顧。 \end{tabularx} \\ \\ \relax
12:26 & \begin{tabularx}{0.7\textwidth}{X} 假如一個肢體受苦，所有的肢體就一同受苦；假如一個肢體得光榮，所有的肢體就一同快樂。 \end{tabularx} \\ \\ \relax
12:27 & \begin{tabularx}{0.7\textwidth}{X} 你們是基督的身體，並且各自都是肢體。 \end{tabularx} \\ \\
[1ex]
\hline
\hline
\end{longtable}
$^{1}$多謝各位收看 再見!.
各位電影姐妹平安!.
因為十週年同慶的緣故.
我相信我們會經常說「事奉」.
我們不停地說過去的事奉.
將來的事奉.
都是一件很重要的事.
事實上最近我們的教會.
都開始熱鬧起來.
多了人參與不同部門的事奉.
我們要為此獻上感恩.
我們的事奉是為了神.
不是為了自己.
究竟我們整個群體.
不單是我們個人.
整個群體的事奉.
怎樣才是討神喜悅呢?.
這是很值得我們去探討的.
因為洪恩家是屬於神的.
不是屬於我們哪一個人.
所以在這裡我們更加要問.
怎樣去服侍神.
才是祂的喜悅呢?.
我們一同來禱告.
讓我們今天能夠以你的話語.
能夠彼此鼓勵.
求你今天開我們的眼睛.
讓我們去看明白神的心意.
讓你的耳朵真的打開.
聽清楚我們的聲音.
願聖靈在我們每一個人裡面動功.
感動我們.
回應神理自己的教導.
奉耶穌基督的名字.
祈禱.
阿們.
好.
夢神喜悅的事奉.
基本上可以說是有兩大點的.
第一點.

$^{41}$夢神喜悅的事奉.
我們說群體.
既是合一的.
也是多元化的.
什麼是合一的呢?.
什麼又是多元化呢?.
哥林多前書十二章十二節這樣說.
就如身體是一個.
它有許多肢體.
身體的肢體最多.
仍是一個身體.
基督也是這樣.
從身體和肢體的比喻裡.
我們看到.
身體就是指教會.
而肢體就是形容信徒.
身體上的肢體各有特色和功能.
就如教會裡的信徒.
各有不同的恩慈.
大家會參與不同的事奉一樣.
我嘗試數過.
在我們洪恩家裡.
有些什麼不同類型的事奉呢?.
我試一下數.
我相信我有機會是數漏的.
證明我們的事奉可以是很豐富.
我們洪恩家裡有什麼呢?.
有私事招待.
坐在在座當中.
是不是也有私事招待呢?.
試一下舉手.
私事招待是誰呢?.
原來我們有不同的事奉人員在這裡.
有領唱的,有師班的.
領唱師班的要舉手.
有的,有不同的位置.
有教主,學的在哪裡?.
舉手,教主,學的.
有不同的人.
有負責兒童事奉的.

$^{81}$可能現在不在這裡.
在上面努力中.
有影音的.
我們的影音不同了.
多了一些姐妹參與.
有堂位.
今天我們也有堂位.
有負責講訊息的.
講訊息不單是現在這裡大堂崇拜.
也有上面兒童事奉那裡.
有祈禱的.
祈禱這個位置.
有這麼一個崗位嗎?.
有些是沒有崗位的.
但是默默在背後的.
做很多禱告的工作.
是的.
有負責清潔的,搬東西的.
這些有位置的嗎?.
大家有沒有試過參與過清潔.
搬東西的事奉?.
有關顧的,探訪的.
有傳福音的.
還有今天這裡.
大家看不到.
有一個我在講話.
他也在同時講話.
就是我們這個.
崇拜的翻譯.
我知道昨晚.
我很辛苦時.
他也很辛苦的.
有參與社會關懷的.
是不是?.
在過程中.
剛才我只是講師班.
其實我們也有福音樂曲.
在這個洪恩家.
很豐富的地方.
我可能會有少數.

$^{121}$還有很多.
背後的,默默的.
有些事奉沒有崗位的.
但是你在這裡參與.
你會看到的.
原來我們是這麼豐富的.
弟兄姊妹的服侍是多元化的.
因為弟兄姊妹的恩賜.
都是多元化的.
雖然信徒很多.
大家有不同的恩賜.
但是要記住一件事.
我們同樣都是屬於.
耶穌基督的教會的.
是屬於神的家的.
所以我們是需要.
在教會裡合一.
正如肢體需要.
在身體裡合一.
談到教會的合一.
保羅這段經文.
有以下兩方面的.
很重要的提醒.
第一個.
我們合一的基礎究竟是什麼?.
13節那裡這樣說.
無論是猶太人.
是希臘人.
是遺奴的,是自主的.
都從一位聖靈受洗.
成了一個身體.
並且共享這位聖靈.
當我們相信耶穌以後.
神就讓聖靈住在我們裡面.
聖靈會在我和你的生命裡工作.
改變和更新我們.
就讓我們成為了教會的肢體.
是聖靈的工作.
令我們成為一個合一的群體.
所以無論我們是來自什麼種族.

$^{161}$好像上面所說的.
是猶太人,是希臘人.
在我們當中.
我們不單止有香港的弟兄姊妹.
我們還有來自哪裡的弟兄姊妹呢?.
大聲一點.
是啦,是太極的.
太極的姊妹們.
還有呢,剛才有人說的.
非洲的弟兄們.
我們都有不同的種族.
或者我們是來自什麼社會階層都好.
在文字裡提到有遺奴的,有自主的.
在當時的年代還有做奴隸的.
所以他就分了遺奴的,自主的.
說的是兩個很不同的社會階層.
無論是哪一種.
我們在主裡面都是屬於一個大的群體.
我們不分彼此.
《已忽所書》四章三字裡這樣說.
以和平彼此聯繫.
竭力保守聖靈所賜的合一.
這個合一.
聖靈已經賜給我們.
我們只需要保持就很好了.
但保持的時候要竭力.
很用力的.
當我們要保持合一的時候.
又要注意以下的事情.
這個就是一個很重要的提醒.
第一件事要合一.
就不要比較.
自卑甚至妒忌.
不要比較自卑甚至妒忌.
十四節說身體原不只是一個肢體.
而是許多肢體.
十五節假如腳說「我不是手」.
所以不屬於身體.
它不能因此就不屬於身體.
十六節假如耳朵說「我不是眼睛」.

$^{201}$所以不屬於身體.
它也不能因此就不屬於身體.
教會有不同的弟兄姊妹.
各自有不同的恩慈.
難免會出現比較的情況.
剛才那段經文描述.
是肢體之間有比較.
比較之後就不開心.
有自卑.
甚至有機會妒忌.
作為腳.
大家看看我們的身體.
今天經常說身體.
手手腳腳 眼耳口鼻.
我們的腳可能會覺得.
我是身體的下半部分.
人都是留意手多一點.
又覺得手做的事情比腳多.
所以人會比較看重手.
於是腳就感到不快樂.
甚至想離開身體.
是我的自卑.
怎麼像手呢.
手就好了.
做這麼多事情.
算了 我們這些腳.
我們還是離開吧.
同樣的 論到頭上的器官.
耳朵也可能會覺得.
比不上眼睛.
這麼受到重視.
通常這個人.
你見到他.
很多時候會看什麼.
看他的眼睛.
你形容人時怎麼說.
他眼睛大大的.
你很少這樣形容.
他耳朵大大的.
很少這樣說.

$^{241}$所以感覺眼睛比較重要.
你形容眼睛是人的靈魂之窗.
那耳朵是什麼呢.
有沒有什麼形容呢.
感覺是吧.
所以耳朵可能會有機會感到不爽.
又自卑.
妒忌.
我想離開身體.
普羅話我們不需要這樣看.
即使腳覺得自己比不上手.
耳朵覺得自己沒有眼睛這麼重要.
他們都不能夠離開身體.
因為他們是身體的一部分.
沒有了他們.
身體就不能夠正常運作.
你只是看.
聽不到.
你整個人馬上就很不舒服.
你真的有很多很多的限制.
你發覺自己.
你只有手.
你的腳不能走.
你都會很不開心.
任何一部分都不可以少.
那些感覺自己沒有這麼重要的.
都是不可以少的.
我們覺得自己不重要.
神不是這樣看.
我們整個教會都不是這樣看.
每一個都是重要.
不可以少的.
所以在這裡不要比較.
不要自卑.
也不要妒忌.
第二件事.
一定要按神所賜的恩賜去服事.
要按神所賜的恩賜去服事.
我們的恩賜是怎樣來的呢?.
十二章四節說.

$^{281}$「恩賜有許多種,卻是同一位聖靈所賜.」.
十二章十一節都說.
「這一切都是由唯一的.
同一位聖靈所運行,.
隨著自己的旨意,.
分給各人的.」.
恩賜不是我們買回來的.
不是我們找回來的.
不是我和你爭取回來的.
不是我們想要就能要到的.
是聖靈按照自己的旨意分給我們的.
不是我們這個人決定要什麼恩賜.
是神決定將什麼恩賜賜給我和你.
賜恩賜的主權是在神的手裡.
而且神會希望我們能夠將祂所用的恩賜.
所給我們的恩賜.
用在一個很合適的崗位上.
十八節這樣說.
「但現在神隨自己的意思.
把肢體一一安置在身體上了.」.
究竟今天我和你是做眼睛一耳朵.
手一腳還是其他什麼器官呢?.
這要視乎神給我們的恩賜是什麼.
但我想問靈子們一個問題.
今天我們的事奉.
是否按照神所給我們的恩賜.
我們是否在合適的崗位上.
發揮自己的恩賜.
你說沒問題.
你憑著愛心.
你憑著對上帝的愛.
你想做什麼就做什麼就行了.
上帝一定喜悅的.
不是.
這間教會.
每一間地方的教會.
這個世界的耶穌基督教會.
只有一個頭.
是耶穌基督.
上帝是教會的掌管者.

$^{321}$祂給我們恩賜.
祂就有對我們的心意.
現在不是我們和你去決定.
我想怎樣事奉.
我們想怎樣事奉.
是神決定我們.
所以是要問的.
是否真的按恩賜事奉.
是合適的崗位.
不合適的崗位有很多問題.
對教會是負面影響的.
一定會.
也很奇怪.
如果不是按恩賜事奉.
第一自己會很辛苦.
別人也很辛苦.
我們不能建立基督的身體.
所以這是一個很大的影響.
17字說.
假如全身都是眼睛.
聽覺在哪裡.
假如全身都是耳朵.
嗅覺在哪裡.
有時候我們有一種心態.
很渴望自己能夠擁有某些恩賜.
說明是自己渴望.
有時候見到別人.
真是很好.
他可以這樣事奉.
很好.
我羨慕.
我希望自己也可以有這些恩賜.
我也想這樣事奉.
於是我們就一起.
參與這個事奉.
就正如我們也想成為眼睛.
因為眼睛是很重要的器官.
今天不停提眼睛.
是很被人重視的器官.
但如果全身都是眼睛.

$^{361}$沒有其他器官.
例如沒有耳朵.
沒有聽覺.
會是怎樣呢?.
你有沒有想像.
只是頭.
全部都是眼睛.
沒有其他器官.
這是什麼頭呢?.
其實這肯定不是頭.
真是完全不正常.
不能運作.
不能生活下去.
這裡甚至有說.
如果耳朵都很重要.
全身都是耳朵.
又怎樣呢?.
又沒有鼻子.
沒有嗅覺.
又怎樣呢?.
如果只有一個器官.
一種器官.
身體根本就不存在.
這不是一個身體.
正如十九節所說的.
假如全部是一個肢體.
身體在哪裡呢?.
假如今天我們都覺得.
某一個部的事奉很重要.
教會裡每一個人.
都參與這個部的事奉.
沒有其他事奉了.
這間教會肯定沒有辦法.
好好地運作下去.
不會拓展.
不會吸納更多的人.
也不可以好好的.
去拓展神的國度.
我舉個例子.
全部的人.

$^{401}$我們很認定.
做領唱的事奉是很美好.
很重要.
從今天開始.
台下所有人都只是領唱.
沒有其他東西了.
我們不會有人.
做一個程序表.
不會有人關心探訪.
傳福音.
不會有人講道.
不會有人做影音.
只是有人站在這裡.
帶大家唱歌.
這間教會會發生什麼事?.
我深信大家一定會忍受不住.
這個哪是教會呢?.
所以如果我們只是想按照自己的喜好.
而不是按照神賜給我們的恩賜.
和神的心意去服侍.
第一,我們的生命仍然是很自我的.
你從這裡看到誰是主.
我是主還是上帝是主呢?.
是上帝是主.
那我要明白我是由什麼恩賜.
應該如何事奉.
如果自己是主.
就自己作主.
我想這樣這樣.
我可以讓人看到我很熱心.
這個熱心.
是不是神想要的?.
是不是對教會是最好的呢?.
值得問.
如果教會要合一.
我們就要一起放下自我.
在這裡有一個很重要的.
需要分享的地方.
我如何知道神賜給我的恩賜是什麼呢?.
這個我也聽不少.

$^{441}$我不知道我有什麼恩賜.
那是做什麼呢?.
是不是?.
特別是剛剛信耶穌.
剛剛回教會沒多久.
好像在這方面對自己認識不多.
我簡單說幾個方法.
其實是很多的.
第一,你直接的.
你可以禱告問神.
神會回應你.
「啊,神啊,我現在加入教會.
你想我參與什麼服事啊?」.
「我不知道神給了我什麼禮物呢?」.
恩賜真的是禮物.
我可以如何為你.
去用上我的恩賜.
神你自己去定.
你可以問.
第二,我深信聖靈也會在你裡面提醒你.
總是有些很特別的感動.
經常覺得我要去插花.
很感恩.
現在我們整個插花的團隊.
我們教會是一個很壯大的團隊.
很多弟兄姊妹在當中去參與.
聖靈可能一直在提醒.
為什麼不去參與這個事奉.
為什麼不去參與那個呢?.
第三,你自己參與事奉的時候.
身邊的人也會給你某個程度的印證.
「啊,我發覺你很懂得關心人.
你去探訪會很好,帶給人安慰.
人們會很喜歡,你會很鼓勵人.
可能不止一個.
弟兄姊妹,甚至道師,傳道人.
也會這樣跟你說.
這個也是一個初步的對自己的了解.
第四,你自己參與某項事奉的時候.
你自己會感到有能力.

$^{481}$你感到有滿足感,是開心的.
你發揮到的.
如果有一樣東西.
你做來做去都不是很行.
你很努力地學習.
你付出很多時間都做來做去不行.
你要想一想.
你問一下神.
「哇,這個東西不知為何我做了很久.
我學了很久,還是不上手」.
是不是你恩賜呢?.
你要繼續尋求.
這樣東西你絕對可以找傳道人聊天.
去尋求指引.
現在是很好的.
坊間出了一本書.
其實那本書已經很久了.
是說恩賜的測試表.
一直做,做,做,做,做.
有些問題可以這樣.
一直做,幫你了解自己.
做完之後,原來我有這些恩賜.
都可以的.
原來有很多方法.
但你不一定要用這本書.
這是其中一個方法.
可以令你更加認識自己.
還有很多.
突然之間會值得花時間.
你去問一問神.
你去認識了解.
我的恩賜是什麼.
上帝擺我在這裡.
一定有原因.
上帝想我在哪裡服事.
你要按神心意去服事.
這是第二件事.
在合一裡一定要注意.
第三件事在合一裡要小心的.
就是不要驕傲.

$^{521}$你要堅持相信.
肢體之間彼此需要對方.
不要驕傲.
你要相信肢體之間彼此需要對方.
二十節這樣說.
但現在肢體最多.
身體還是一個.
眼睛不能對手說.
我用不著你.
頭也不能對腳說.
我用不著你.
眾多的肢體合成一個身體.
既然是屬於同一個身體.
就彼此有關連的.
要互相配搭.
很有趣.
這裡說的是很特別的關係.
眼睛和手的關係.
在這段文字裡.
眼睛是比手高的.
眼睛可以指揮手去工作.
頭也比腳高.
頭也可以指揮腳去行動.
在某個程度上代表了兩個不同的群體.
在教會裡做領袖的.
帶領者的.
和一些被帶領者的角色.
做帶領者的.
不能夠去興旱.
不可以去驕傲.
覺得看不起被帶領者.
帶領者與被帶領者.
一起彼此需要對方.
如果眼睛不需要手.
又或者頭不用腳.
其實眼和頭.
都不能夠為身體.
發揮很大的功用.
可以說是唇亡齒寒.
你不用他們.

$^{561}$不和他們配搭.
自己也成為無用的.
正如在教會裡負責帶領工作的.
無論是弟兄姊妹同工.
如果不能和被帶領者配搭去服侍.
一同合一.
也不能有效地建立教會.
剛才說的合一.
很重要.
「夢神喜悅的事奉」既是合一.
亦是多元化.
「夢神喜悅的事奉」還有一樣東西.
是「滿有愛」.
二十二節這樣說.
「不但如此.
身上的肢體.
更意為軟弱的.
更是不可缺少的.
身上的肢體.
我們認為不體面的.
與法給他加上體面.
我們不雅觀的.
與法裝飾得雅觀」.
在事奉的過程當中.
我們不能忽略一些.
條件不是那麼好的.
有缺陷的.
人看為很軟弱的肢體.
保羅說.
這些肢體.
是不可缺少的.
在我們的群體裡.
他們亦是我們的一分子.
正正他們有缺陷.
條件上有限制.
我們要與法給他加上體面.
與法裝飾得雅觀.
什麼意思呢?.
我們要關注他們多一點.
比其他人要多.

$^{601}$我們要照顧他們多一點.
比其他人要多.
教會是學習愛的地方.
在服侍的過程中.
我們要彰顯出神那份.
很豐厚的愛.
我想起好幾年前.
我在帶一些青少年活動時.
有一次我玩一個遊戲.
覺得很深刻.
有一次我要一班門訓成員.
那班門訓成員當時還是中學生.
是高中生.
我要他們一起去學.
完成一個任務.
要用繩織一個網.
由一個地方將一個人.
將他們小組當中一個人.
運去另一個地方.
要記住繩子是很柔軟的.
他們要想辦法織一個網.
要承載一個人的體重.
然後他們要由一個地方運去另一個地方.
當時那個門訓小組很特別.
不是每一個成員都很精靈.
很聰明.
很青春活潑.
大部分都是.
有一個在智力上.
有少少差一點.
但我真的很欣賞這個.
智力上差一點的一個姐妹.
她對神的心很好.
我記得她洗禮時.
她讀她的《偏見證》.
她很堅持讀完一次又一次.
她自己本身連揭聖經都不行.
讀聖經每次我教她的時候.
教洗禮班.
一定要有其他弟兄姊妹陪她一起.

$^{641}$陪她一起讀才能讀到.
閱讀困難.
很多方面困難.
但去她洗禮那天.
她很用心.
讀完一次.
自己之前練習.
練了很多次.
去洗禮的時候.
就要將最好的獻給神.
平時她是很認真的.
她不會遲到.
出席最多.
但她在身體上有些缺陷.
那次小組其他.
大部分都是跳舞.
青春活力.
結果那次我很感動.
整班成員.
很快就談.
怎樣做網.
應該這樣.
過程中她們並沒有忽略這個姐妹.
她們很懂得做.
將她納入裡面.
即使織網的時候.
其實這個姐妹很困難.
好像幫不到什麼.
但她們很聰明.
教她.
你負責拿著繩頭.
我們會給你繩子.
穿來穿去.
姐妹覺得都有份.
她們很懂得照顧.
最後織好.
她們決定.
你躺在上面.
我們負責抬你.
姐妹很聽話.

$^{681}$其實感覺很好.
因為她好像被眾人所注意.
被眾人照顧.
她們真的很疼她.
我很感動.
到最後那次.
最大收穫是見到那份愛.
其實在教會裡的事奉.
都是要有這份愛.
不是因為條件好.
特別要強調條件好的.
自己一堆.
如果在群體裡有這樣的成員.
要將他們一起納入.
我們就以愛彼此結連.
在愛裡一同服侍.
二十四節說.
我們雅觀的肢體.
自然用不著裝飾.
單純配搭者身子.
把加倍的體面.
給賴有缺欠的肢體.
免得身體不協調.
總要肢體彼此照顧.
神將加倍的關懷和關注.
給了那些有缺欠的肢體.
是為了調配身體.
讓它們能夠協調.
整個身體要達至.
一個很合一的境界.
很和諧的一個境界.
是平衡的.
很明顯.
雅觀的肢體.
會付出多一點.
去照顧缺欠的肢體.
以達至那種平衡和諧.
最重要的是肢體之間.
真的互相照顧.
彼此相愛.

$^{721}$這是神的心意.
在事奉裡不是你還你的.
我還我的.
你不要搞我.
其實我也不想理會你.
各自完成就好了.
我們很忙的.
時間有限的.
不是這樣.
是要互相照顧.
星期日回來.
好像打仗一樣.
大家有沒有留意.
哇 影音部很緊張.
領唱在這裡正在努力.
如果有師班就努力練習.
福曲隊是一樣的.
都要練習準備好.
教主要學又主要學.
又有師事招待.
又準備好了.
去歡迎新朋友.
看著整個過程.
很多很多.
怎樣叫彼此相顧.
彼此相愛.
如果我發覺這個部有問題.
不是我的部.
可以幫的 都一起幫.
我的心也掛念.
那份愛和情在這裡.
不是一塊塊的.
各自自己搞定就算了.
是一切.
如果有一天.
我們做到的效果.
回來的時候.
感到很溫暖.
雖然另一個部的事.
與我無關.

$^{761}$但我很關心他們.
我也為他們禱告.
有事幫忙一定出.
同樣的.
如果我們有需要.
別人也一起過來.
感覺很美好.
究竟我們要愛到什麼程度呢?.
剛才Rita姐妹有分享.
26節說.
假如一個肢體受苦.
所有肢體就一同受苦.
假如一個肢體得光榮.
所有肢體就一同快樂.
對於肢體受苦.
我們不會冷漠無感受.
我們會感同心受.
這個我曾經在港島裡分享過.
在我爸爸的安息禮拜裡.
就是我們一家人.
面對傷痛的時間.
與我們一同去感受那份傷痛.
其實這個是很無比.
很寶貴.
不是說.
「曾姑娘,你自己爸爸的傷痛.
你們家裡慢慢經歷吧」.
「不關我們事」.
不是,是一起.
直到如今.
我心裡也很感動.
為弟兄姊妹的這份愛.
如果我們在座當中.
有弟兄姊妹.
在她事奉的過程裡有很多難處.
有自己身體上的需要.
很多方面的需要.
我也想弟兄姊妹記住.
你不是一個人.
是有整個群體.

$^{801}$與你一起.
你不是一個人.
整個群體都關心,都在乎.
因為神想我們.
一起的.
去經歷那份傷痛.
一起承擔那份苦.
同樣的.
有人一個肢體得光榮.
所有的肢體就一同快樂.
有一個肢體.
無論是得獎也好.
他做得很好.
我們就不要妒忌.
有時我們的心也很有趣.
我好,別人稱讚我就舒服.
他好.
有些不舒服.
很容易妒忌.
要提醒的是.
他得光榮.
我們一起得光榮.
欣賞他.
稱讚他.
這個團隊是這樣的時候.
整個教會是這樣的時候.
我心想每個人都很想事奉.
因為我想做得好的時候.
我知道弟兄姊妹是支持我的.
喜歡的.
我有不開心.
我有憂傷.
弟兄姊妹會跟我一起分擔.
愛就是這樣.
事奉是要這樣的.
不是分開了.
是不同的部分.
這是一幅很美麗的圖畫.
肢體能夠這樣相愛.
27字說明了什麼.

$^{841}$這是基督的身體.
我們在教會裡面.
各自是肢體.
這些肢體是結連在一起.
耶穌基督是我和你的頭.
我們是要黏在一起.
任何一個肢體.
發生什麼情況是不理想的.
其他人一切都不理想.
任何一個肢體得榮耀.
很好.
我們一切都很有榮耀.
所以在過程中.
很願意在未來的日子.
我們再思考.
我們的事奉是怎樣的.
夢神喜悅的事奉是合一.
是多元化.
夢神喜悅的事奉是沒有愛.
特別是剛才提到的那些位置要小心.
我們不要比較自卑妒忌.
我們要按神所賜的恩賜事奉.
更加不要驕傲.
堅持需要對方.
我們一同禱告.
愛我們.
給我們恩賜.
給我們事奉.
是我們不配得的.
天父求你給我們回應你的愛.
我們就好好地按你的心意去事奉.
你愛洪恩家.
你給我們在這裡合一.
我們不同部門的弟兄姊妹.
在事奉當中彼此配搭.
彼此相愛.
讓更多的弟兄姊妹被激勵要起來事奉.
這裡就成為一個很溫暖的地方.
能夠吸納更多未信的人進來這裡.
他們可以認識主.

$^{881}$主願意你去建立我們的生命.
復興我們的生命.
在未來的日子.
我們很快就會迎接十週年的慶典.
願意在裡面.
你叫到我們整個群體.
更加有上帝你的愛.
更加去明白什麼是你的大靈.
讓我們跟著你的大靈.
去迎向未來的更多的十年.
奉耶穌基督的名字祈禱.
阿們.
\newpage



\section{希伯來書 12:1-3}
\label{sec:9yhz2ddk1BQ}
\textbf{2021.10.17 《得力秘訣》 希伯來書 12\_1-3 (和修版) 陳秀賢牧師}
\newline
\newline
連結: \href{https://youtube.com/watch?v=9yhz2ddk1BQ}{\texttt{https://youtube.com/watch?v=9yhz2ddk1BQ}} ~~~~ 語音日期: 2021-10-18
\newline
\newline
\hyperref[sec:P_diqejYp0w]{\small{< < < PREV SERMON < < <}}
~
\hyperref[sec:index_chronic]{\small{[返順時目]}}
~
\hyperref[sec:index_scriptual]{\small{[返順卷目]}}
~
\hyperref[sec:PriyvxgsE30]{\small{> > > NEXT SERMON > > >}}
\newline
\newline
希伯來書 12:1-3
\newline
\begin{longtable}{cl}
\hline
\hline
章節 & 經文 (和合本修訂版)\\
\hline
12:1 & \begin{tabularx}{0.7\textwidth}{X} 所以，既然我們有這許多見證人如同雲彩圍繞著我們，就該卸下各樣重擔和緊緊纏累的罪，以堅忍的心奔那擺在我們前頭的路程， \end{tabularx} \\ \\ \relax
12:2 & \begin{tabularx}{0.7\textwidth}{X} 仰望我們信心的創始成終者耶穌，他因那擺在前面的喜樂，輕看羞辱，忍受了十字架的苦難，如今已坐在神寶座的右邊。 \end{tabularx} \\ \\ \relax
12:3 & \begin{tabularx}{0.7\textwidth}{X} 你們要仔細想想這位忍受了罪人如此頂撞的耶穌，你們就不致心灰意懶了。 \end{tabularx} \\ \\
[1ex]
\hline
\hline
\end{longtable}
$^{1}$各位弟兄姊妹早晨.
我已經退休了.
剛剛七月退休.
我代表迎風堂向大家問安.
很開心老牧師邀請我來分享神的話.
很開心和你們敬拜.
很開心聽到你們的樂曲的伺服器傳福音.
我們一起祈禱.
主啊天國君萬國君.
主啊你是色在永在的傳能者.
主啊我們眾之體.
今天奉你的名在當中聚集.
去敬拜一位這麼偉大的神.
天國君萬國君.
你是災雷的君王.
所以主啊你配得我們的敬拜和讚美.
主啊求你幫助我們.
去讀你的話的時候.
要明白你的真理.
明白你自己的救恩.
明白你自己認識你自己.
要讓我們繼續去跟隨一位這麼偉大的主.
禱告奉主名求.
阿們.
好.
我不知道大姐妹.
你心靈上有沒有感到一些擔憂.
懼怕.
或者是徬徨.
或者是疑惑.
或者是很愧疚.
很憤怒.
或者是很疲乏.
很多新聞.
很多香港的事.
很多世界的事.
充斥在我們的腦子.
我們會有這些感受.
我有一個朋友.
他也六十多歲了.

$^{41}$早幾個月他跟我分享一件事.
他就像這個公仔一樣.
就是很抑鬱.
甚至要吃精神藥.
其實他是一位很熱心事奉神的人.
他很熱心很熱心事奉神很多年.
到他經過這兩年香港所發生的事.
或者是疫情.
他人很低落.
很抑鬱.
很情緒.
還要照顧一個老人家.
九十多歲的老人家.
所以不知道為什麼突然陷入一個很抑鬱的狀態.
很不開心.
但是因為他是一個這麼熱心事奉的人.
不知道怎麼跟別人說他心裡的懼怕.
還有他失眠了.
很長時間失眠了.
不知道怎麼說.
他有一天就打給我.
我聽完他說話之後.
我心裡有一種很不忿的感覺.
就是那種不忿氣.
我們基督徒服侍這麼熱心事奉神.
真的很不忿氣.
不可能要吃精神藥.
不可能抑鬱到這樣的狀態.
所以我很不忿氣.
我就跟他說我們要去敬拜和讚美.
以及倚靠耶穌基督.
所以我就跟他說你不要躲在家裡.
你不要躲在家裡.
你出來吧.
跟我們一起有聚會.
跟一群基督徒熱心事奉的基督徒.
敬拜神,讚美神.
以及認識神的神蹟和奇妙的作為.
所以他就出來了.
就是很奇妙.

$^{81}$跟他相處了兩個月.
彼此的鼓勵,敬拜,祈禱.
他就吃了一些藥.
在兩個月之後.
全部康復了.
精神醫生說你不用吃藥了.
你已經很好了.
各位觀眾原來我們人生裡面.
我看起來我們人生都是很有年齡的人.
在世界上你做了幾十年人.
其實你所面對的這兩三年的事.
是令我們不知道怎樣去處理的.
我們的心情不知道怎樣去處理.
身為一個基督徒.
我們不知道怎樣去處理.
剛才有機會認識.
羅木斯的太太斯母.
她很有氣質.
她分享她有去一些服侍.
不過現在疫情已經停擺下來.
我們原來在這個情況下.
我聽到粵曲的弟兄姊妹.
來了這裡八年.
但在外地也去了很多次.
因為疫情又去不了.
我們活在一個這樣的困局困境裡.
今天我很想和我們一起鼓勵.
我們一起讀希伯來書十二章一至第三節.
預備請.
所以既然我們有許多見證人.
如同雲彩的環繞我們.
就該卸下各樣重擔和僅僅前來的罪.
以堅忍的心奔那擺在我們前頭的路程.
仰望我們的信心的創始.
成中的耶穌.
他因擺在前面的喜樂.
興旱羞辱.
(聖經).
我們看到希伯來書這段聖經教我們.
我們如何走人生十字架路.

$^{121}$繼續服侍神.
雖然我們年紀大.
機器有少少壞.
但我們走十字架路.
我們如何得力.
彼此去勸勉.
希伯來書的作者提醒我們.
其實有許多見證人.
如同雲彩的環繞我們.
我們就要卸下各樣重擔.
僅僅前來的罪.
以堅忍的心奔那擺在前頭的路程.
甚至要仰望我們創始成中的耶穌基督.
在希伯來書裡面.
有許多偉人描述.
亞伯拉罕,阿葛,以撒,摩西.
這些偉人你都不陌生.
弟兄姊妹.
這些偉人.
例如基澱,巴拉,參孫,耶佛他,大衛,撒母耳等等.
這些人都在世上認識神.
並且他們為神捱很多苦楚.
特別我高舉了一些字.
紅色的字.
有人是忍受刑罰的.
有人忍受戲弄和鞭打.
以及被困所,監獄,金戒各種磨練.
他們甚至被人用石頭打死.
被割死,被刀殺.
他們遇到許多困難.
今早我早上在我自己教會迎風堂崇拜.
聽了一個中東的宣教士分享.
這位女士也是年輕的.
是我們宣道會的分堂的一位傳道人.
堂主任.
她三年前去了中東一個地方叫索米城.
其實她去的時候.
剛剛在伊斯蘭國非常蓬勃的時候.
你聽過ISIS.
她是一個小女子.

$^{161}$在三年前去了這個地方被警察傳福音.
有沒有困難?.
她今早說有困難.
因為中東那個地方有戰亂.
有火爆.
很多殺害.
很多毀壞.
神呼召她去索米城這個地方.
索米城是一個代名.
去到中東的地方傳揚救主.
在世上原來有很多偉人.
蒙召去服侍主.
我們的神.
許多時候這些信心偉人.
在不同的時代.
在不同的際遇.
表現出他們信心的果效.
他們願意為主作見證.
我們應該跑在這個路途上.
無論遇到什麼困難和危難.
我們都可以靠主得勝有餘.
我們看看聖經.
如何提醒我們跑十字架路.
而得勝.
第一要放下各樣的重擔.
放下.
卸下各樣的重擔.
你就會想.
是不是放下我們照顧子女.
照顧孫.
照顧家人的責任.
特別60歲的弟兄姊妹.
我們的家人都80歲90歲.
甚至有100歲.
我們是不是放下照顧老人家的責任.
可能有時事奉很多重擔.
是不是放下我們事奉的重擔.
是不是放下我們工作的重擔.
這叫卸下重擔的意思.
即是不用做.

$^{201}$不用做.
聖經教我們不是放下重擔.
不用做.
而是你做的時候有秘訣.
即是有很多重擔在你的肩膀.
但你是有秘訣的.
耶穌基督是教我們提醒我們.
如何可以輕省背負的擔子.
我們一起讀聖經.
預備請.
「凡勞苦擔重擔的人.
到我這裡來.
我要使你們得安息.
我心裡又和謙卑.
你們當負我的壓.
向我.
(聖經).
耶穌基督沒有叫我們不用做.
即是放下所有的重擔.
不是.
而是祂要呼籲我們.
「凡勞苦擔重擔的人.
如何做到我這裡來.
你就可以得安息」.
耶穌說我心裡又和謙卑.
耶穌其實是神.
但祂都謙卑.
去依靠耶和華.
祂說你要當負我壓.
學我.
你心裡就得安息.
我很想說「壓」字.
我不知道鄧英姐妹有沒有耕田.
這裡有沒有人耕田的.
舉手來看看.
咦這麼少.
只是副主席嗎.
這麼少人耕田.
所以你在教會也要耕田.
就是這麼耕田.

$^{241}$其實你看到牛是負壓的.
壓是木造的.
放在牛那裡.
其實牛是耕田.
但它是用手腳耕的嗎.
因為它不能用手腳耕.
它是走路的.
但它如何去耕田呢.
就是它負壓在肩膀裡.
人們就要控制壓.
它就會動.
木的壓在它的肩頭上.
就可以幫它爬泥.
它一邊走.
壓一邊拉著.
後面就幫它鬆一些泥.
耶穌基督說.
你可以到我舍內來.
我使你們得安息.
他說你當乎我壓.
壓的意思是.
大家知不知道耶穌基督的職業是做什麼的.
祂的爸爸.
木匠.
他的爸爸約瑟是做木匠.
他爸爸做木匠.
你猜耶穌會不會鬥木呢.
他會不會鬥這些傢俬呢.
會.
耶穌應該會.
因為約瑟是做木匠.
所以耶穌基督說.
農夫.
牛的壓.
是負在他的肩頭上.
是最後的.
好的木匠.
他會做一個櫃子.
好的木匠會做一個很適合的壓.
給牛的.

$^{281}$因為如果他做得好的話.
做好的牛就不會那麼辛苦.
壓就在牛的肩膀上.
他就可以走得很順利.
但如果壓過重.
他會壓死牛的.
是不是.
或者壓本身有刺.
都會壓死牛的.
同意嗎.
所以那個壓是很重要的.
那個壓要量身訂做的.
好的木匠.
木工.
他會量身訂做的壓.
放在牛的肩頭上.
那隻牛.
是托著那個壓的時候.
他怎樣.
他會很輕散地走.
他走的時候.
他可以爬到田裡.
爬到下面的泥.
所以耶穌就說.
你要負我壓.
這個壓是神量身訂做.
為那隻牛而做的.
所以你今天所事奉.
你所擔當的東西.
其實神是為你量身訂做的.
神是.
如果你唱歌唱得好.
有些人五音不全.
唱不到.
你可以又唱到.
又平衡.
什麼的.
你可以服侍神.
其實你的壓是輕散的.
為什麼.

$^{321}$是量身訂做的.
你很獨特的.
上帝是獨特為你而做.
獨特給你的.
弟兄姊妹.
你做招待.
我正好坐在這裡.
一個招待.
給我一個擁抱.
我說不用了.
你很有愛心.
我不用一個擁抱.
我放在椅子上.
其實每一個人.
他都是神量身訂做.
給他服侍的.
是不是.
我很敬佩.
正好.
做主席.
或者打鼓.
弟兄姊妹.
你不同的人.
有不同的職份.
但耶穌說.
壓是我量身訂做給你的.
所以你要靠耶穌基督.
耶穌說.
你負我的壓.
並且你學我.
耶穌基督說.
我們為什麼要學耶穌.
他說他柔和謙卑.
耶穌是神.
他也什麼.
他也要依靠神.
我們的上帝.
所以我們就要.
效法耶穌基督.
去依靠我們偉大的上帝.

$^{361}$我們不要靠己力去服侍神.
我們不要靠己力.
去負我們的重擔.
我不知道你在職場裡.
或者在教會裡.
你負著什麼壓.
我不知道你遇到什麼壓力.
什麼事奉的壓力.
主告訴我們.
到我借你來.
我給你安息.
我給你的壓.
是容易的.
是輕散的.
第二就是脫去容易纏累我們的罪.
希伯來斯十二章四節.
他說你們與罪惡要爭鬥.
說沒有抵抗到流血的地步.
原來我們走十字架路裡.
你承不承認.
其實有很多東西.
都是纏累我們的罪.
有人曾經在深海裡.
被樹葉海草纏住他的腳.
他怎樣.
他就算很厲害.
游泳很厲害也好.
被樹葉纏住.
樹枝纏住.
他都會拉他下去.
他會死.
我們做教徒.
就算愛了十年.
二十年.
三十年.
四十年.
我們都會有些罪纏住我們的腳.
令我們跨下去.
所以求主祝福我們.
耶穌在聖經裡教導我們.

$^{401}$我們要脫去.
緊緊纏累我們的罪.
並且我們要與罪搏鬥.
搏鬥到一個要流血.
好像還未到流血的地步.
這麼重要.
在保羅的哥林多後書十章.
四至六節裡說.
我們征戰的.
本來不是屬血氣的.
我們要憑著上帝給我們的能力.
能夠攻破堅固的邢侶.
我們要攻破這些計謀.
各樣難阻人認識上帝的邢侶.
將人的心奪回.
跪向忌妒.
我個人很少唱粵曲.
但聽到你的萬世軍戰國軍.
很強烈的打仗氣勢.
我們要將人.
將他們.
因為我們未信耶穌的時候.
我們有很多邢侶.
而這些邢侶難阻我們.
去跪向基督.
這些邢侶非常堅固在我們生命當中.
叫我們看不到耶穌的光輝.
所以我們要攻破.
我們一生要打仗.
要對罪那種抵擋.
好像要流血的地步.
你這樣抵擋.
你看到這些邢侶.
人生很多邢侶.
這些邢侶難阻我們.
去跪向基督.
所以我們要像一個.
屬靈的基督徵兵.
要打仗,常常都要打仗.
要為神打美好的仗.

$^{441}$我們要思想一下.
我們有什麼邢侶.
或者纏著我們的嘴.
我們要留意.
弟兄姊妹.
我早前聽到.
我自己的WhatsApp錄音.
原來某間教會.
不知為何轉錯了錄音給我.
錄音的內容是.
他是基督徒.
他說我打算要離婚.
兒女們誰供養誰供養.
我就很驚訝.
為什麼有這個錄音.
立即有人說.
對不起,我轉錯了錄音給你.
但他轉錯了錄音給我.
我又不開心.
他說他是基督徒.
他說他離婚.
他說他的兒女怎樣分化.
原來你不要以為.
信了耶穌之後就很正當.
因信清義.
但其實我們都要.
繼續去面對我們的罪.
面對我們.
撒旦給我們很多誘惑.
我們要抵擋.
可能你看看.
你看看你的Facebook.
你可以留意一下你自己.
生命現在怎樣的狀態.
你敢不敢將你的Facebook的記錄.
即是你上網的記錄.
你上網的Google的記錄.
上YouTube的記錄.
你敢不敢讓人看呢?.
突然間牧師說.

$^{481}$你給我手機看看.
看看你YouTube的紀錄.
你最近看了什麼.
你給我手機看看.
你上Google有什麼紀錄.
你其實尋找什麼資料,資訊.
很多弟兄.
我估計姐妹都有份.
上網看色情的故事.
色情的一些刊物.
但有人在.
我們要留心.
這些是什麼纏累了我們的罪.
我們自己不知道.
我們可能一心好,熱心事奉.
但如果有罪是纏累我們.
轟炸我們的時候.
我們要留心.
要脫去這些罪.
還有我們的心要被奪回.
歸回耶穌基督的身上.
走十字架路.
要堅忍的奔跑十字架路.
弟兄保羅提醒我們.
你奔跑十字架路.
我不是以自己已經得著.
我只是一件事.
忘記背後,努力面前.
向著標竿直跑.
要得著上帝在耶穌基督裡.
從上面照我們來的獎賞.
如果我要奔跑這條路.
我要很有耐性去跑.
並且我要忘記背後.
我要努力地去跑.
並且要向著目標去跑.
擺在我面前的目標去跑.
要忍耐,要堅忍.
堅忍這個字其實是一個堅持.
我們越來越忍耐很差.

$^{521}$有時我去超市買菜.
我會留意.
哪條路的收銀員會比較快.
我會知道誰比較快.
我會排隊.
因為收銀員很慢的.
會很快.
不是吧.
他們超市的收銀員很慢.
逐樣地掃複.
我們很沒耐性.
有時我自己去櫃員機按機按錢.
很害怕,自己也有壓力.
有三,四個人在排隊.
他們會說.
你知道的.
你又要按了你的號碼.
按了號碼又要轉戶.
轉戶的奉獻.
轉戶的奉獻有時又回答不了.
又要再按一次.
按完之後又要轉戶.
轉戶對了.
看清名字就轉戶.
這樣也要一點時間.
有時眼不靈,手不靈.
有時慢了,後面就.
其實你也會被人壓.
你看人家後面.
你看前面的人.
有十張支票這樣進去.
你就說.
搞錯了.
十張支票你去銀行.
在櫃員機這樣進去.
我們很沒忍耐.
糟了,我們現在走十字架路.
幾十年.
直到見主面.
如果你長命百歲一百年.

$^{561}$已經要忍耐去跑.
所以希伯來書教我們.
你要看著目標跑.
就好像馬拉松的人.
很有耐性.
他們很有耐性去練.
很有耐性去跑.
他們不一定要跑第一名.
幾百人怎麼跑第一名.
第一名只有一個.
他們的原意是什麼.
要跑完.
跑完半馬也好,全馬也好.
總之要跑完.
跑完就很開心.
所以我們走十字架路.
就是這樣跑.
而且要忍耐心去跑.
跑得來也有力量.
我們要仰望創始而成中的耶穌基督.
仰望我們信心而創始的耶穌基督.
因為祂擺在前面的喜樂.
前面很開心.
就忍受今天的羞辱.
忍受十字架的苦難.
祂現在已經坐在上帝的寶座右邊.
我們要細心仔細想想耶穌.
這位忍受罪人.
而且被人頂撞的耶穌基督.
你就不至灰心失意了.
我們跑這條路的時候.
我們要看著耶穌基督跑.
因為十字架的路.
其實不是順風順水的路.
同意嗎.
不是一定成功的路.
不是一定發達的路.
但這條路可能是很艱辛的.
基督徒在香港有邊緣的數字.
三十多萬.

$^{601}$很少的.
所以我們經常說.
現在怎樣.
現在的情況怎樣.
我們香港教會怎樣.
是不是很擔心.
很詭異.
所以我們走的路一點也不簡單.
我們弟兄姊妹在不同的職場工作.
都會有壓力的.
有困難的.
一點也不簡單.
但是我們要跑的時候.
我們要一直看著.
你不要看左右的風景.
你不要左顧右盼.
你只看著耶穌跑.
看著耶穌跑.
甚至你看到耶穌.
他受今天的羞辱.
他看到前面的喜樂.
因為他頂十字架.
他忍受這個苦難.
他頂十字架.
就帶領億萬的人進入神的天國.
所以現在的苦楚.
其實不是苦楚.
你看到這個苦楚.
但是你看到前面的喜樂的時候.
耶穌就是甘心的為我們受苦.
而作出犧牲.
你要思想.
你要忍受.
有時候我們會被人質問.
傳福音.
你有沒有被人質問過.
有沒有.
我都見到有人點頭.
如果你從來沒有被人質問過.
你都會有些問題.

$^{641}$你不可能沒有被人質問過.
你有沒有受過苦.
我會回應.
如果你說不是.
我很悠閒.
我每天早上都喝早茶.
和我的朋友喝早茶.
沒有人質問我.
因為你沒有傳過.
你沒有說過.
你沒有受過苦.
剛才聽到兩位姐妹說.
他們有時候都會去內地.
去坐車.
去坐船.
去偏遠的地方.
或者去城市.
去傳揚耶穌.
他們傳揚耶穌.
每個人都很好.
在這裡拍手掌.
不是的.
有些人會罵你.
有些人就不專心.
有時候你會覺得.
原來傳福音不是經常都這麼成功.
或者很多人信儲.
可能有人瀏覽你.
但是他說你走十字交道的時候.
你要忍受別人這樣頂撞你.
三位一看真純的耶穌基督.
祂尚且被人頂撞.
那麼我是誰呢?.
我是誰呢?.
我是被人責罵.
是很如常的事.
被人罵是很如常的事.
你不被人罵才好奇.
特別現在這個世代.
我們做基督徒.

$^{681}$有時不同的時代.
基督徒在香港是很被人讚賞的.
有時候不是被人讚賞的.
現在基督徒.
一時之間.
我不知道現在是怎樣.
但是我們要知道.
在不同的時代.
我們要傳揚主耶穌基督的時候.
是有困難的.
不過你走十字交道的時候.
你看著耶穌.
我是誰呢?.
耶穌也是甘心遇到這些羞辱.
祂看到前面的喜樂.
祂就很開心.
各位弟兄姊妹.
洪水橋這十天八天.
經常在新聞上聽到的.
聽到嗎?.
本來人們都不知道哪裡是洪水橋.
現在看到這個施政報告.
就知道洪水橋在哪裡.
就是洪水橋.
而且是發展北區.
我已經替你們很興奮.
其實我唯一認識的教會.
在洪水橋就是你們.
我已經替盧牧師很興奮.
因為你們將來有很多事情要做.
在這個區的發展.
地盤工人.
住進來的人.
隨意喝茶.
你會看到你們有很多朋友.
帶他們來.
跟他們分享神.
關心他們.
靠近幫助他們.
跟他們做朋友.

$^{721}$很多很多.
有時候也會遇到困難.
但聖經告訴我們.
看著耶穌基督.
(等待中).
聖經在羅馬書.
哥林多後書.
說過.
聖經.
他說.
今天我們所受的苦楚.
其實是極重無比.
是自贊自輕的苦楚.
相對將來是極重無比.
永遠的榮耀.
各位兄弟姐妹.
你們知不知道.
很多人說將來的喜樂.
就是北區發展.
你們這個區很大的發展.
洪水橋.
這個是喜樂.
看到這麼多人進來.
很開心.
所以要經常擔擔子.
有時候會覺得有些困難.
或者不開心.
但你今天所受的苦.
跟你傳方.
分享耶穌的時候.
別人拒絕你.
或者逼迫你.
對你不好.
其實你所面對的苦難.
是自贊自輕的.
到末日回來.
主耶穌.
天國君.
萬國君.
回來的時候.

$^{761}$我們會得到極重無比.
永遠的榮耀.
所以這是聖經教導我們.
我最近看了一本書.
追隨基督.
是潘霍華牧師寫的.
潘霍華牧師.
我又說一下這位牧師.
其實他已經過世了.
大家知道.
德國希特勒的時候.
第二次世界大戰.
希特勒自己要做王.
叫很多人跪拜他.
說他是王.
他也發動了很多.
將猶太人在集中營謀殺.
那時候的時代.
希特勒殺了六百萬人.
在那個時代.
同樣的時代.
潘霍華就是那個時代.
潘霍華是神學家.
也是柏林大學的教授.
很有知識.
很有地位.
其實很顯赫.
他覺得不贊成人這樣.
去跪拜希特勒.
希特勒和歐洲宣戰.
希特勒也叫德國人.
出來打仗.
叫年輕人出去幫忙打仗.
潘霍華牧師拒絕.
他不打仗.
他不參戰.
他不拒絕.
所以他拒絕打仗.
所以他們就抓了潘霍華牧師.
去集中營那裡住.

$^{801}$在1945年.
他們就把潘霍華牧師.
送上殿頭台.
死了.
這樣就殉道了.
在他死了幾天之後.
第二次世界大戰停止.
因為盟軍已經贏了希特勒.
所以.
但是人們就說.
潘霍華真是不值.
早一點.
盟軍贏了.
早幾天就不用死了.
潘霍華牧師.
39歲就死了.
是很出色的牧者.
他曾經說了一句話.
潘霍華牧師說.
當基督呼召一個人的時候.
他是呼召他來.
為他死.
是的.
他說.
基督呼召一個人的時候.
他是呼召他來死.
他就殉道了.
其實有一段時間.
他有一個美國朋友.
他還沒死之前.
有一個美國朋友就說.
其實你不要再回德國了.
你不要回德國.
我們幫你去美國.
他就.
想了一段時間之後.
他就說我也要回德國.
回自己的國家.
他真的這樣說.
他說假如現在.

$^{841}$現在.
假如我現在.
不跟我的同胞一起分擔.
當前的事年的時候.
將來我是沒有資格.
跟他們一起重建.
復興.
復元後的.
德國基督徒生活的.
如果我今天走了.
我要去美國.
他走了.
這麼顯赫的人.
是神學家.
又這麼出名.
一定很多人會抱他.
他就說我要回自己的國家.
如果我今天走了.
將來我沒有資格.
再回來德國.
跟基督徒重建.
基督徒復興的生活.
所以最後.
在集中營裡面.
去殉道.
今天我分享的訊息.
在十字架路上.
得力的秘訣就是.
要放下重擔.
不是不做.
而是.
耶穌給你的壓.
是輕生的.
我們要效法基督.
怎樣去依靠耶華的能力.
第二.
你留心.
你每天.
現在所面對的生活.
是否有什麼罪.

$^{881}$是在把你潛入海底.
第三.
你要忍耐.
要堅忍的.
走完十字架道路.
不要羞辱基督.
好好看看.
六十歲.
七十歲.
八十歲.
你仍然愛主.
最後.
你跑路不孤單.
要仰望.
創始.
而成鐘的耶穌基督.
最後送這幅畫給你.
漂亮嗎?.
我畫的.
其實我退休的時候.
我的會友.
不知為何.
其實他叫我畫個公仔出來.
不知為何.
你退休了.
做個紀念品.
我又不知為何.
我就畫個公仔出來.
這個公仔很開心.
太陽花.
他仰望神.
他有兩個袋子.
一個袋子就是.
麥克風.
我很喜歡唱歌.
所以你只要有些粵曲.
我都會跟著唱.
其實我未學過.
筆很喜歡寫東西.
和畫東西.

$^{921}$他是一隻.
企鵝.
我很喜歡企鵝.
肥一點點.
很像我.
有一個字.
就是要仰望耶穌.
創始成中的耶穌.
是的.
而且是開心心的去仰望.
神母應許我們十字架.
這條路困難.
不困難.
沒有說過.
不過我們真的要仰望.
創始成中耶穌基督.
我們一起祈禱.
多謝耶穌.
你揀選我們成為你的子女.
多謝你讓我們成為天國的子民.
多謝你讓我們跟著耶穌基督.
行十架的路.
多謝你創始而成中.
我們就是看著你這樣跑.
我們人生的目標.
人生的道路.
事奉的道路.
所以祝福紅印堂的弟兄姊妹.
月發的愛你.
月發愛主.
祝福紅印堂的弟兄姊妹.
道路都不偏離你.
最後祝福紅印堂的弟兄姊妹.
在十字架道路上.
努力的向祂周邊的人.
分享神的愛.
講述神的愛.
求主祝福紅印堂.
多謝主聽我們禱告.
奉主耶穌基督名而求.

$^{961}$阿們.
\newpage



\section{創世記 5:1-5}
\label{sec:PriyvxgsE30}
\textbf{2021.10.24 《耶和華眼中的亞當》 創世記 5\_1-5 (新譯本) 楊穎兒傳道}
\newline
\newline
連結: \href{https://youtube.com/watch?v=PriyvxgsE30}{\texttt{https://youtube.com/watch?v=PriyvxgsE30}} ~~~~ 語音日期: 2021-10-27
\newline
\newline
\hyperref[sec:9yhz2ddk1BQ]{\small{< < < PREV SERMON < < <}}
~
\hyperref[sec:index_chronic]{\small{[返順時目]}}
~
\hyperref[sec:index_scriptual]{\small{[返順卷目]}}
~
\hyperref[sec:L2pB3hPJTOs]{\small{> > > NEXT SERMON > > >}}
\newline
\newline
創世記 5:1-5
\newline
\begin{longtable}{cl}
\hline
\hline
章節 & 經文 (和合本修訂版)\\
\hline
5:1 & \begin{tabularx}{0.7\textwidth}{X} 這是亞當後代的家譜。當神造人的日子，他照著自己的樣式造人。 \end{tabularx} \\ \\ \relax
5:2 & \begin{tabularx}{0.7\textwidth}{X} 他造男造女。在他們被造的日子，神賜福給他們，稱他們為人。 \end{tabularx} \\ \\ \relax
5:3 & \begin{tabularx}{0.7\textwidth}{X} 亞當活到一百三十歲，生了一個兒子，形像樣式和自己相似，就給他起名叫塞特。 \end{tabularx} \\ \\ \relax
5:4 & \begin{tabularx}{0.7\textwidth}{X} 亞當生塞特之後，又活了八百年，並且生兒育女。 \end{tabularx} \\ \\ \relax
5:5 & \begin{tabularx}{0.7\textwidth}{X} 亞當共活了九百三十年，就死了。 \end{tabularx} \\ \\
[1ex]
\hline
\hline
\end{longtable}
$^{1}$弟兄姊妹平安.
平安.
多謝洪恩堂的堂主任魯牧師的安排.
讓我能夠有機會在實習完之後.
仍然可以回到洪恩堂.
與弟兄姊妹分享聖經的信息.
剛才讀出的經文是來自《創世記》第五章一至五節.
事實上第五章整章都是亞當的家譜.
在禱告之後.
我決定今天的講道題目是.
耶和華眼中的亞當.
亞當的故事.
可能大家蓋上聖經都能背得出.
那還有什麼好說呢?.
盼望我們今天可以藉著這個經文.
讓我們回望亞當的一生.
有些片段之餘.
還有更多的發現.
所以我在今天的講道裡面.
所用的經文都會是與亞當有關的.
讓我們一起先禱告.
親愛的天父.
求你來到我們的中間.
你張開孩子的口.
打開我們的心.
讓我們各人都聽得明你的道.
看得見你的慈愛和信實.
祈禱奉主名求.
阿們.
今天我的題目是.
耶和華眼中的亞當.
我不知道大家對亞當的印象是怎樣的呢?.
讓我先簡單介紹一下他.
在聖經裡面.
他是神創造的第一個人.
他也是第一個.
也是唯一一個.
曾經居住在伊甸園的男人.
他也是.
他也是第一個有工作的人.

$^{41}$耶和華神將那人安置在伊甸園裡面.
使他修理,看守那個園子.
他和夏娃建立了第一個家庭.
因為耶和華說那人獨居不好.
我要為他做一個配偶幫助他.
於是耶和華神就用那人身上所取的肋骨.
造成了女人.
之後亞當夏娃犯罪的記錄就開始了.
那個女人見到那裡的果子好作食物.
也悅人的眼目.
且是可喜愛的.
能使人有智慧.
就摘下果子來吃.
又給他丈夫.
他丈夫也吃了.
他們二人眼睛就明亮了.
才知道自己是赤身露體.
拿無花果樹的葉子為自己編造裙子.
亞當和夏娃一起吃了這個禁果之後.
改變了整個人類的命運.
耶和華神即時的咒詛.
他們更是被逐出伊甸園.
聖經這樣說.
耶和華神便打發他出伊甸園去.
耕種他所自出之土.
由此可見除了即時的審判之外.
也有世世代代的影響.
對於女人來說.
耶和華神說.
「我必多多加增你懷胎的苦楚」.
「你生產兒女必多受苦楚」.
「你必戀慕你丈夫」.
「你丈夫必管轄你」.
對男性來說.
耶和華神對亞當說.
「你既聽從妻子的話」.
「吃了我所吩咐你」.
「不可以吃那棵樹的果子」.
「你必為你的緣故受咒作」.
「你必終身勞苦才能得到地裡吃得吃的」.

$^{81}$當他們被逐出伊甸園後.
他們生了戈隱和阿伯.
戈隱因為耶和華看中了阿伯和他的禮物.
而沒有看中戈隱和他的禮物.
他非常憤怒.
垂頭喪氣.
最後更殺了阿伯.
如果我們問問上面兒童主一學的小朋友.
我相信他們都能說出.
亞當,夏娃和他們兩個兒子的故事.
既然如此,還有甚麼可以說呢?.
有的,我們可以再深入一點.
我們去說說亞當在神眼中的意義.
先跟大家說說亞當的名字.
為何聖經要將第一個人的名字命名為亞當呢?.
原來在希伯來的聖經裡.
亞當,亞丹在舊約裡出現了560次.
在絕大多數的情況下都是解作「人」或「人類」.
所以在希伯來文背景成長的人.
去讀聖經時.
一看到亞當就會意識到亞當是指人類.
所以在《創世記》第一章.
我們看到神創造了這個世界和其中的一切後.
神創造了亞當.
亦都寓意了最後一個被造的種類就是人類.
當亞當和夏娃吃了這個果子後.
他們就能夠像神一樣具備與神同等的智慧.
他們的罪就是來自他們想超越神和人的界線.
當亞當和夏娃吃了這個果子後.
他們的眼睛就明亮了.
發現了自己赤身露體需要找地方去躲藏.
耶和華神就在這個伊甸園中尋找他們.
呼喚著「你在哪裡?」.
剛才說過了.
亞當亦即是人類的意思.
我們可以看看今天.
仍然有很多人想做神.
很想用自己的經驗,自己的智慧去主宰自己甚至別人的生命.
並且要賺取全世界.
甚至要超越神和人的界線.

$^{121}$今天又有人因違反了神的話語而逃避神.
而神仍然在大地上不斷向人發出呼喚的聲音.
「你在哪裡?」.
神的呼聲就是從創世記開始.
神的心意就是要人類.
包括不同的萬族,萬民.
回歸到祂的身邊.
要人類就像亞當一樣.
承認自己的過犯.
願意在神的教導裡.
生命得到重整.
得著神為人預備的福.
神真的那麼在意祂所創造的人類嗎?.
按照今天的經文.
我有三個方面跟大家分享.
第一是創造人類的心意.
在第五章亞當的家譜裡的第一節這樣說.
神做人的時候是按照自己的樣式去做的.
人類在所有的受造物裡是獨一無二的.
即使我們的身體有許多方面都像動物一樣.
但是唯獨人類能夠作為一個群體.
共同參與在神國的建立.
和耶穌一起同工.
並且在聖靈裡面能夠有一個合一的團契.
這是神創造的心意.
當神創造世界的時候.
祂已經有這個計劃.
祂透過我們的身體.
我們的生活與神與人同在.
有一位宣教學者Overglasser這樣說.
人類不僅僅是困在肉體中的靈魂.
而且是擁有完整心理社會的機體.
儘管所有其他的生物都要依賴他們的創造者.
但是人類接受了神的氣息.
就能夠與神進行一個智能的互動.
擁有服從或不服從祂的自由.
因此人類為了與造物主建立特別的關係.
而被造出來的.
因此人類就是要盡心盡性盡義盡力愛他們的神.
經文提到.

$^{161}$神是按照祂的樣式做人的.
這個樣式又是甚麼呢?.
神學家加爾文在十六世紀已經說過.
人類是非物質的本性.
人類會使用語言,理性去辨識一些事物.
人類又懂得用技術製造工具和使用工具.
同時人類又有社會性.
能夠起草一些法律建立一個團體.
人類又有歷史性.
會關注累積的傳統.
人類更加有審美的眼光.
能夠創造,裝飾一些事物.
去享受一些事物.
人類又有道德倫理的價值觀.
人類又能夠投入宗教去接觸那個看不到的世界.
最後他提到人類有認識神的能力.
而且在人類裡面.
他們的本質就是有一種對宗教信仰的渴望.
希望和超自然的事物建立關係.
渴望在某個程度上去反映神的本性和相似之處.
這些非本質的本性都是指向神的善.
這些樣式就好像種子一樣.
埋藏在我們裡面.
有待時日結出果子.
更強調這些樣式是因為我們每一個人.
都可以有機會做神的代表.
當我們有了這些樣式的時候.
我們就會映照出神的光輝.
就好像月亮,星星這些受眾物一樣.
他們自己本身不懂得發光.
唯有太陽的光照落他們的身上.
月亮,星星就顯出太陽的光.
並且在漆黑的天空中.
變成明亮的天燈.
所以神創造的心意是要人有神的樣式.
從神那裡映照出光.
充滿在地上.
第二點我想和大家分享的是.
管理大地的任務.
神按照他的樣式做人是有必要的.

$^{201}$創世記這樣說.
我們要照著我們的形象.
按照我們的樣式做人.
使他們管理海裡的魚,空中的鳥.
地上的畜生和全地上爬的一切昆蟲.
在神的創造計劃裡.
原來人是有任務的.
因為神將修理,看守的工作給了人類.
意思就是要人去管理神.
被神創造的東西.
所謂的管理不是容許人任意妄為.
卻是要在生活上,工作上都要顯得出神的樣式.
意思就是在我們的頭裡,內裡有這些樣式.
然後我們就要將這些樣式轉化成為行動.
在身體以外的地方展示出來.
所以我們會透過工作,生活裡.
讓我們的生命展示出神的樣式.
去愛人,愛神,愛人.
要有神的美善.
由忍耐,寬恕,憐憫等等的特質.
有很多信徒從星期一打拼到星期六.
星期日回到教會.
就是在那一堂課要吸取所有的力量.
為的是要回到家庭,社會裡繼續發亮,發光.
但事實上有時我們面對日常的種種困難.
其實真的很疲倦.
我們會有心力交碎的時候.
有時我們會覺得與神相近.
有時覺得與神相拒.
完成了一個星期的擔子之後.
又變得很疲累.
好像沒有甚麼空間去承載神的話語.
沒錯,這是很正常的.
因為我們都是人.
我們在一個有限的區隔裡生活.
但神就是如此真實地賜予我們身體.
讓我們不斷在生活裡學習尋求神.
去求神引導,操練我們的靈性.
我們去學習將我們軟弱的交給神.
我們這個身體要成為聖靈的電.

$^{241}$要藉著我們不斷循環去學習.
神透過萬事萬物不斷呼召我們回到祂美善的樣式.
就是在與神相交的過程裡.
人漸漸就會自發地去按照神的話語.
去服從祂的旨意去生活.
漸漸與神的樣式就越來越相似.
下一張.
第三節.
阿當一百三十歲的時候.
生了一個兒子.
樣式和形象都和自己相似.
就給他起名「失特」.
阿當生了該人和阿伯之後.
其實之後就生了第三個兒子.
阿當就看著這個嬰兒.
這個嬰兒其實甚麼都沒有做過.
為甚麼阿當會看著這個嬰兒.
就看到那個樣式,形象是和他自己相似呢?.
阿當活了一百三十年.
意思就是他經歷了人生很長的時間之後.
他明白了這種樣式和形象是從神而來的.
是由神所賜的.
是神造人類的心意.
顯現在生命裡面.
而不是這個嬰兒做了甚麼事情.
阿當明白了樣式和形象原來是與生俱來.
是神造人的心意是美善的.
而當阿當看見這一點的時候.
阿當正正就是開始和神的心意相近.
大大話話阿當去明白神創造的這個心意.
都要用一百三十年.
何況是我們呢?.
神的這個心意是要把這個祝福延到後代的.
第五章第二節.
「他創造了一男一女」.
「在創造的時候,神賜福給他們,稱他們為人」.
這裡提到「做男做女」.
當然我們很快會知道他在回應.
神認為乃人獨居不好.
要為他做一個配偶去幫助他.

$^{281}$神建立這個家庭的體系.
目的是要人透過生育去延續下一代.
這是其中一種神賜福的方式.
但是在神眼中.
在創造人的一刻.
其實已經和這個福有關.
就是從創造人的一刻.
神吹了一口氣給人類.
然後阿當吸了.
阿當就有了這個靈氣.
就由本來沒有生命的人.
變成有靈的活人.
下一章.
《創世記》二章七節這樣說.
耶和華神用地上的塵土做人.
將生氣吹在他的鼻孔裡.
他就成了有靈的活人.
名叫阿當.
神是用他自己的說話去創造世界.
他說有光就有光.
要水分開,水就分開.
唯獨是人類.
他不單只是用口說.
並且是用地上的塵土去做人.
並且加入自己的生命氣息.
去創造人類.
這個生命氣息.
就是神賜的福.
而且這個生命氣息.
是每個人都有的.
遍滿全地.
當我們走在街上的時候.
我們見到的每一個人.
他們都有著神這一口氣.
只不過能夠活著.
而又明白這個福氣的.
並不是每個人都做得到.
不知道大家是否留意自己的呼吸.
我們睡覺的時候.
通常是很平穩的.

$^{321}$一呼一吸.
大家做完運動.
跑完步.
吸了很多空氣的時候.
我們就會感覺到.
空氣更新了我們裡面的養分.
我們做完運動之後.
就會感覺渾然一新.
大家有沒有試過.
趕巴士,趕上班.
我們會很跑.
我們會急速,緊張,煩亂.
這個呼吸的步伐又不同.
跟跑步和睡眠的呼吸節奏.
又有分別.
其實我們一呼一吸的節奏.
是對應著我們生活的每一個步伐.
如果我們透過意識.
我們呼吸的節奏.
我們學習如何不急忙.
不慌亂在世上的事情.
我們懂得將這個憂慮交給神的話.
我們一呼一吸的氣息.
都會跟神的步伐接近.
神創造世界是有規律的.
是有晚上,有早晨.
有了這一種規律.
祂才會去做人.
晚上是休息的時間.
我們就去休息.
白天是活動的時間.
我們就去活動.
其實我們是按著神給我們的規律.
用我們的身體去親近我們的神.
每一天休息有時.
活動有時.
透過調節我們生活的節奏.
一呼一吸.
看到我們能夠無時無刻.
其實我們都感受到.

$^{361}$符在我們裡面.
我們又在符的裡面.
越來越接近神.
無論在社會上,在家庭上.
我們都是透過呼吸去連繫於神.
所以為甚麼我們牧者們都會建議.
信徒在早上做靈修.
就是要將整天的節奏先放慢.
回歸到神的話語裡面.
去得到力量.
調教好每天的節奏.
然後與神同行.
最後說說亞當一生的年日.
第五章第四節.
亞當生失特以後還活了八百年.
並且生了其他兒女.
第五節.
亞當共活了九百三十歲就死了.
其實亞當是否活了九百三十歲呢?.
我們真的不用那麼斟酌.
在今天來說究竟是多久的日子.
意思就是一段很長的日子.
下一章.
詩篇九十篇第十節這樣說.
我們一生的年日是七十歲.
若是強壯可以到八十歲.
但其中所經過的不過是受煩.
擲老虎受煩.
轉眼成空.
我們便如飛而去.
當人生經歷了好些年日.
回望過去.
有些甚麼值得記載呢?.
今天的家譜.
這個經文是亞當的家譜.
家譜通常是一堆已經死去的人的名字.
這個家譜是神為亞當所編寫的.
是記錄了他從出生到死亡.
然而當中有神的創造.
有神的賜福.

$^{401}$又有一些延續性.
大家可以細心去看看.
第五章一至五節.
是上帝親自為亞當所寫的.
可能你會說.
整本聖經都是神所密使的.
這有甚麼特別呢?.
在《創世記》第一,二章.
是講神的創造一個無罪的世界.
第三章是講人的墮落.
人想越過神和人的界線.
最終人要承擔死亡的後果.
第四章是聖經第一宗的殺人事件.
經過一連串的罪惡.
恐怖的殺人案.
神極度憤怒的咒詛之後.
我們揭到第五章.
是神為亞當親自提筆.
寫出亞當一生的事情.
有趣的是神沒有記錄亞當的罪.
或是亞當兒子的惡行.
神為亞當寫的家譜.
卻是神在亞當身上作的美事.
有創造,有賜福,有延續.
這是神所關心的.
事實上,整本聖經都是講人如何不斷墮落.
而神在當中又不斷創造.
不斷願意將福賜給人.
又要人承接這一切的福.
一代又一代地延續下去.
神在亞當的故事中.
的確要亞當負上他的罪責.
然而在我預備經文時.
我自己很深的感受.
就是我看著這個家譜.
930年過去了.
亞當已經歸於塵土.
所留下來的家譜.
就是要我們用神的眼光.
去回望亞當的一生.

$^{441}$所看見的不是亞當犯過甚麼彌天大罪.
而是神充滿了豐盛的慈愛.
神赦免罪惡和過犯.
神不以罪惡來定人的一生.
因為神已經透過主耶穌基督.
被釘於十字架的救贖事件.
耶穌基督已經用祂的寶血.
洗淨了所有人的罪.
以致人類得以與神和好.
亞當一個人犯的罪.
全人類都要負上這個罪責.
但同時耶穌基督.
一人背負了全人類的罪.
人類就從神那裡得到這個恩典.
在亞當這個家譜的短短的五節裡.
正正就是因為他沒有寫過去的罪.
讓我更加看到.
思想到耶穌的作為.
亞當亦是人類不用再背負著這個罪責.
在我一開始講到的時候.
我提到亞當就是指人類.
神在人身上作的美事.
就是要賜福給人.
使人從原本不明白的.
到明白.
到活出有神形象的樣式.
這件事原來是更加值得.
存留在神的記憶裡.
死亡是人生命的一部分.
我會想如果有一天.
自己回到天家的時候.
神會怎樣為我這個墓碑.
去撰寫一些文字呢?.
不知道神會寫些什麼呢?.
盼望藉著今天的講道.
大家都去整理一下自己生命的樣式.
同時調節一下我們與生俱來的.
這種呼吸的節奏.
讓我們懂得每時每刻都有神的形象樣式.
更加重要的是我們不要忘記.

$^{481}$要花點時間坐下來.
數算一下神賜給我們各人的福.
讓我們一起祈禱.
親愛的天父.
我們滿心歡喜的來讚美你.
因為你就是那位創造天地萬物.
使全地都充滿著你生命氣息的神.
我們讚美你.
在你的創造心意裡面.
你加入了人類.
我們每一個都是極之的微小.
但是你卻願意不斷呼召我們.
回到你的身邊.
做你的兒女.
學習你的形象樣式.
你也在這個世界上.
不斷呼召還未認識你的人.
你不斷在說你在哪裡.
聖靈就求你教導我們各人.
使用我們.
使我們不只是星期日才做信徒.
乃是我們每一天.
每一分.
每一秒.
我們的一呼一吸.
都連繫於你.
時刻與你同工.
我們如此的禱告.
乃是奉那位為我們犧牲.
捨命的主的名而求.
阿們.
(開機中).
\newpage



\section{哥林多後書 4:1-18}
\label{sec:L2pB3hPJTOs}
\textbf{2021.10.31《不喪膽的事奉》 哥林多後書 4\_1-18 (和修版) 林誠信牧師}
\newline
\newline
連結: \href{https://youtube.com/watch?v=L2pB3hPJTOs}{\texttt{https://youtube.com/watch?v=L2pB3hPJTOs}} ~~~~ 語音日期: 2021-11-01
\newline
\newline
\hyperref[sec:PriyvxgsE30]{\small{< < < PREV SERMON < < <}}
~
\hyperref[sec:index_chronic]{\small{[返順時目]}}
~
\hyperref[sec:index_scriptual]{\small{[返順卷目]}}
~
\hyperref[sec:zUN5t6YF3yU]{\small{> > > NEXT SERMON > > >}}
\newline
\newline
哥林多後書 4:1-18
\newline
\begin{longtable}{cl}
\hline
\hline
章節 & 經文 (和合本修訂版)\\
\hline
4:1 & \begin{tabularx}{0.7\textwidth}{X} 所以，既然我們蒙憐憫受了這事奉的責任，就不喪膽， \end{tabularx} \\ \\ \relax
4:2 & \begin{tabularx}{0.7\textwidth}{X} 反而把那些暗昧可恥的事棄絕了，不行詭詐，不曲解神的道，只將真理顯揚出來，好在神面前把自己推薦給各人的良心。 \end{tabularx} \\ \\ \relax
4:3 & \begin{tabularx}{0.7\textwidth}{X} 即使我們的福音被遮蔽，那只是對滅亡的人遮蔽。 \end{tabularx} \\ \\ \relax
4:4 & \begin{tabularx}{0.7\textwidth}{X} 這些不信的人被這世界的神明弄瞎了心眼，使他們看不見基督榮耀的福音。基督本是神的像。 \end{tabularx} \\ \\ \relax
4:5 & \begin{tabularx}{0.7\textwidth}{X} 我們不是傳自己，而是傳耶穌基督為主，並且自己因耶穌作你們的僕人。 \end{tabularx} \\ \\ \relax
4:6 & \begin{tabularx}{0.7\textwidth}{X} 那吩咐光從黑暗裡照出來的神已經照在我們心裡，使我們知道神榮耀的光顯在耶穌基督的臉上。 \end{tabularx} \\ \\ \relax
4:7 & \begin{tabularx}{0.7\textwidth}{X} 我們有這寶貝放在瓦器裡，為要顯明這莫大的能力是出於神，不是出於我們。 \end{tabularx} \\ \\ \relax
4:8 & \begin{tabularx}{0.7\textwidth}{X} 我們處處受困，卻不被捆住；內心困擾，卻沒有絕望； \end{tabularx} \\ \\ \relax
4:9 & \begin{tabularx}{0.7\textwidth}{X} 遭受迫害，卻不被撇棄；擊倒在地，卻不致滅亡。 \end{tabularx} \\ \\ \relax
4:10 & \begin{tabularx}{0.7\textwidth}{X} 我們身上常帶著耶穌的死，使耶穌的生也在我們身上顯明。 \end{tabularx} \\ \\ \relax
4:11 & \begin{tabularx}{0.7\textwidth}{X} 因為我們這活著的人常為耶穌被置於死地，使耶穌的生命在我們這必死的人身上顯明出來。 \end{tabularx} \\ \\ \relax
4:12 & \begin{tabularx}{0.7\textwidth}{X} 這樣看來，死是在我們身上運作，生卻在你們身上運作。 \end{tabularx} \\ \\ \relax
4:13 & \begin{tabularx}{0.7\textwidth}{X} 但我們既然有從同一位靈而來的信心，正如經上記著：「我信，故我說話」，我們也信，所以也說話； \end{tabularx} \\ \\ \relax
4:14 & \begin{tabularx}{0.7\textwidth}{X} 因為知道，那使主耶穌復活的也必使我們與耶穌一同復活，並且使我們與你們一起站在他面前。 \end{tabularx} \\ \\ \relax
4:15 & \begin{tabularx}{0.7\textwidth}{X} 凡事都是為了你們，好使恩惠既藉著更多的人而加增，感恩也格外顯多，好歸榮耀給神。 \end{tabularx} \\ \\ \relax
4:16 & \begin{tabularx}{0.7\textwidth}{X} 所以，我們不喪膽。雖然我們外在的人日漸朽壞，內在的人卻日日更新。 \end{tabularx} \\ \\ \relax
4:17 & \begin{tabularx}{0.7\textwidth}{X} 我們這短暫而輕微的苦楚要為我們成就極重、無比、永遠的榮耀。 \end{tabularx} \\ \\ \relax
4:18 & \begin{tabularx}{0.7\textwidth}{X} 因為我們不是顧念看得見的，而是顧念看不見的；原來看得見的是暫時的，看不見的才是永遠的。 \end{tabularx} \\ \\
[1ex]
\hline
\hline
\end{longtable}
$^{1}$主席.
各位兄弟姐妹早安.
首先我在此恭賀錦宣堂.
你們四十週年實在是一個不短的日子.
我在此恭賀曾牧師,陳映華牧師.
還有你們教會整個執事團隊.
所有團契小組的組長,領袖們.
我相信過去這四十年神的恩典和慈愛信實.
是伴隨著大家的.
為什麼我們要在這個地方建立一個教會呢?.
我相信早期神已經將這個感動.
很清晰地放在眾領袖的心中.
願意藉著這個地方一個群體.
你們彼此相愛也都衷心去傳揚福音.
建立門徒.
剛才我很享受在你們當中一起的敬拜.
我看見弟兄姊妹.
特別師班還有台下伴奏的一班姐妹們.
我真的大開眼界.
你們吹的樂器叫什麼來的?.
Aerophone.
非常好.
奏出來的音色配搭我們的傳統鋼琴,風琴.
你會發覺也有一個以別一身的感受.
其實上帝是配得我們這樣尊崇的.
還有你們不同的語言的弟兄姊妹.
菲律賓裔的姐妹們剛才的舞蹈.
我看見神在你們教會所作奇妙大而可畏的事.
每年都有新信主的新抱.
我相信上帝在前面的四十年.
若主未回來已先.
我深信祂要激勵我們每一個弟兄姊妹.
今天為什麼我們說我們心裡有時會喪膽呢?.
會沮喪呢?.
你看見題目寫著「毫不沮喪的服侍」.
其實在聖經裡面.
我很想今天透過保羅他自身的一個生命的見證.
跟大家分享.
在哥林多後書第四章.
如果我們整章聖經來看.

$^{41}$來回顧一下保羅他的人生.
你就可以發覺他裡面的生命是一些秘訣.
聖經說到他真的不沮喪.
不喪膽.
其實這是我們中文聖經的譯本的不同.
和合本就是譯做不喪膽.
新譯本是譯做不沮喪.
但新漢語譯本是譯做不灰心.
英文兩個最普遍的譯本.
NIV或者是NASV.
就說「We do not lose heart」.
弟兄姊妹你今天是不是心有時沉了下來.
或者沮喪了下來.
或者你覺得我今天的事奉還有價值嗎.
或者你覺得我今天身處在這個時代裡面.
我怎樣可以進一步去伏線呢.
其實我們先問這個問題.
我們要問一問我們自己.
我們清不清楚今天我們所事何事.
我們為何要去事奉.
原來我們今天每一個人的人生都是短暫的.
我們必須要珍惜.
在這個過程裡面.
讓我們不要看我們自己.
或者單是在看別人.
或者看這個環境.
當我們忘記是誰呼召我們去事奉的.
是誰差遣我們去服侍的.
我們就會很容易跌在這個陷阱裡面.
所以我先用幾個見證和大家稍為分享.
約伯記第五章七節講到.
人生在世必如患難如同火悲騰.
我們不要以為奇怪.
其實信了耶穌的人.
不代表我們的人生就順風順水.
很多時候都不是.
你都知道信耶穌的人一樣會失業.
甚至可能會有癌症.
甚至有不好的際遇.
但正在這些困難裡面.

$^{81}$我們如何去回應上帝的呼召和揀選.
我們今天始駐不渝.
我們看一看在歷史裡面.
甚至在近代裡面.
都有很多如同雲彩的見證人.
大家可以看到這位姊妹Joni Erikson.
她是在美國一個很出色的輪椅上的畫家.
她十七歲的時候因一次跳水意外.
導致到她頸部以下的手腳都癱瘓.
經過了一段痛苦深刻的信仰掙扎.
她靠主的恩典成為傑出的口含畫筆的畫家.
勵志作家和演說家.
成為多人生命的祝福.
旁邊那位我相信可能有些弟兄姊妹會熟悉的.
她是一個神學家.
Martha Dawn Tomoha.
是加拿大溫哥華維珍神學院的靈修神學教授.
她提出的軟弱神學是國際知名的.
她不單止有學術和聖經的研究.
她自己個人也有很真實的一個經歷.
她曾經是癌症的病人.
後來病好之後再確診.
患有乳癌.
她甚至一隻眼已經失明.
她曾經也坐在輪椅上.
香港學生分組團契曾經邀請她來香港分享.
同工問她為何願意揭露自己身體和關係上的傷痛.
她說你寫的時候不痛嗎.
Martha教授她怎樣回答呢.
她說痛是痛的.
但若果不分享這些.
我的讀者就不會明白.
我相信的神其實在人的痛苦裡面仍然是信實的.
Tomoha姐妹.
她在痛苦之中仍能持守信仰.
而且她很努力地書寫神學的著作.
因為父母從小就教導她去學習聖經.
神的話早已在她的心裡成為她生命的力量.
讓她感受從天上而來的愛的傳源.
以致她有勇氣去觸摸自己的傷痛.

$^{121}$亦都能夠很真誠地和其他人分享.
我們也看看我們熟悉的.
我們中國人的教會裡面.
蘇欽沛姐妹.
她少年的時候已經患了甲狀腺癌.
有沒有人認識蘇欽沛姐妹?.
有多少姐妹認識她?.
有點奇怪.
可能大家都是比較年輕一點.
所以不是很認識.
欽沛姐在80年代在香港是非常出名的一個作家.
她在文字服務工作裡面激勵了很多.
特別是向年輕人呼喚.
她成為突破機構的創辦人之一.
在她的鼓勵下.
蔡元雲醫生放下她的專業.
投身服務年輕人.
她所有的文學作品和劇作.
都反映出她的人格和信仰的反省.
她的遺作《死亡別狂傲》.
成為歷久不衰激勵人心的一個上街的書籍.
欽沛姐在1982年的復活節已經安息了.
她成為一個時代很特別的見證.
她率先是發動幫助失落的青少年.
去抗衡消費文化.
她提出減薄生活.
反思女性身份的角色.
對祖國作出承擔.
如今她雖然死了.
卻因著信營救說話.
旁邊這位姐妹比較現代一點.
這位大家看到她的名字叫廖志.
不知道有沒有人認識她呢.
竟然有啊 非常好.
這位廖志姐妹.
她的故事也很特別.
她在2008年.
大家知道四川發生什麼事嗎.
是一個閩村的大地震.
在當時5月14日下午.

$^{161}$她原本是會出去跟人有一個活動.
但她留守在家裡.
陪著她的小女兒.
沒有人預計到原來一個大地震在那時發生.
她們在家裡的這三個人.
她和她的女兒.
和她的奶奶.
不幸撞身於這些壓力當中.
最後她是倖存者.
她能夠倖存.
不過她要截肢.
兩隻腳都要鋸走.
所以你看到這幅相就是真實的.
這個不是動畫.
不是那些所謂的P圖.
這個是真實的一個見證.
當時她的姐妹.
她的生命失去了她最重要的雙腿.
因為她的工作就是一個舞蹈的教師.
一個喜歡跳舞的人.
失去了自己的雙腿.
你說她的生命還有什麼意義呢.
在當時她真的想過自殺.
甚至她的丈夫嫌棄她.
當時也拋棄她.
如果你是她.
你猜你的生命會怎樣.
我想你會真的和我一樣.
我想這個打擊太大了.
太大了.
但是感恩就是她後來遇到有一班基督徒.
在查經班裡面.
她認識了上帝的恩典慈愛.
結果她來到重拾這種生活.
繼續要好好地為神而活.
這種的鬥心.
她生命這種意志和堅定.
是源於她信仰有一份新的動力.
後來她遇到她的丈夫.
是一個台灣出生.

$^{201}$但是在美國成長長大.
我們俗稱就是ABC.
American Born Chinese.
她在當中的過程.
這位丈夫其實是專做義肢的.
在過程裡面.
他們在上海相識.
然後後來代入愛河.
他們在幾年前結婚.
現在有兩個子女.
這個就是她的生命的見證.
非常感人.
甚至最近在中央電視台.
竟然在一個大型的節目裡面.
她來到表演她的舞蹈.
她述說她的見證.
我覺得這種比我們看見.
今日我們的信仰不是個人性的.
我們的信仰是在整個世界.
甚至在不同的舞台上.
我們要向世人宣告.
就是上帝不是缺席.
上帝在人間各樣的困難.
上帝不是不見了.
往往只是我們人的信心.
我們失落了.
我們迷失了.
大家可以看到這位弟兄.
是來自澳洲的Nick Rudgevich.
他自己出世的時候.
是沒有手沒有腳的.
但他父母充滿著神的愛.
讓他自幼認識主.
在他成長的過程.
他受到很多人的歧視.
他曾經試過自殺.
但基督的信仰.
讓他重新毫無怨恨.
而且他重尋生命的意義與目的.
耶穌基督給他內心.

$^{241}$有無比堅定的力量.
堅忍地克服各樣的障礙.
他學有所成.
獲得兩個大學學位.
亦都懂得游泳與踢球.
現在他到全球不同的地方.
見證主耶穌的恩典.
他曾經亦都進入過.
共產和回教的國家.
去述說他的生命故事.
旁邊那幅相.
大家可能亦都會.
不知道有沒有人認識.
是來自台灣的一個盲人的合唱團.
是由一班盲人的弟兄組成的.
我們教會曾經邀請過他.
來到我們當中.
來服侍見證神的恩典.
我亦都曾經.
跟他們一起同台食飯.
我看到其實這班弟兄很可愛.
他們內心其實比很多.
是沒有我們眼睛.
沒有盲到的人其實還清心.
他很敏銳.
他很留意到世界上.
發生的每一件的事.
我們聽到他們.
生命的見證的時候.
就很體會我們今天所信的神.
就是又真又活的.
保羅確實讓我們明白.
我們今天我們的信仰.
要有一個毫不沮喪的服侍的話.
我們需要學習這幾方面.
我們來看這張聖經.
他提到第一節和第十六節.
都有提到不喪膽.
就是因為我們蒙了這份憐憫.
就了這個事奉的責任就不喪膽.

$^{281}$在第一節裡面其實是承接前文.
第三章六節的主題.
提到我們是這個新約的執事.
我們有一個榮耀的積分.
我們今天.
我們知道我們的身份是誰.
我們是神的兒女.
是蒙他所愛為他所揀選的.
這個身份.
是我們成為我們蒙召的基礎.
今天神揀選我們.
並非因為我們有多少知識學問.
也都不是在乎我們在社會上的地位.
高與低的問題.
在神的眼中.
他要我們很清楚知道.
我們是蒙了他的憐憫.
我們是蒙了他的愛與及揀選.
所以我們不會沮喪和喪膽.
我們能夠成為這個新約的執事.
我們好好地去傳揚福音的時候.
不是來將一些自異常的.
一種的律例條文去跟人分享.
我們要將一種活潑的生命的見證.
就是耶穌基督這種帶給我們新造的人.
這種出死入生的一種經歷.
我們與世人去分享.
在結束這張聖經的最後.
他亦都提到我們不喪膽.
是因為我們有這個復活的指望.
所以讓我們一起看看這張聖經裡面.
我試一下和大家用三方面來看.
就是頭六節是講到一個很清晰的照明.
而中間由第七到第十二節.
就是這個生命的能力.
我們要能夠持續不疲倦不灰心地去服侍主.
是因為耶穌在我們這個雅氣裡面成為那個寶貝.
而最後在十三到十八節.
是講到一個實在的信心.
這個實在的信心有三方面.

$^{321}$是因為我們有同一位的聖靈的見證.
帶給我們有相同的信心.
亦都經歷這位相同的基督給我們的復活.
以及一個榮耀的願景.
這個願景就是來自父神給我們看見.
在他的榮耀的寶座面前.
我們可以看得很遠.
原來將來永恆的家鄉裡面.
在那裡再沒有任何眼淚.
痛苦 疾病 死亡.
這一切都要過去.
原來在基督耶穌裡面.
我們是有這一種是打不死亦都是打不倒的.
這一種韌力我們去服侍主 服侍人.
讓我們一齊先看頭六節裡面.
清晰的照明.
讓我們一起讀.
所以既然我們蒙尼文.
受了者事奉的責任就不喪膽.
反而擺那些暗媒可恥的事.
棄絕了 不行詭詐 不曲解神的道.
又將真理顯揚出來.
好在神面前擺自己.
善給各人的良心.
即使我們的福音被遮蔽.
那只是對滅亡的人遮蔽.
這些不信的人.
被這世界的神明弄黑了心眼.
使他們看不見基督榮耀的福音.
基督本是神的像.
我們不是傳自己.
而是傳耶穌基督為主.
並且自己因耶穌作你們的僕人.
那分副光從黑暗裡照出來的神.
已經照在我們心裡.
使我們知道榮耀的光.
站在基督的臉上.
好我們在這一段經文.
我們可以再一次看見.
神給我們這個呼召裡面.

$^{361}$是一個很清晰的.
都是出自祂的憐憫.
我們既然有這一份來自神的憐憫的時候.
讓我們不要沮喪 不要灰心.
我們要忠於上帝.
託付給我們的真理.
聖經這個真道.
在末世的末期.
是更加會有很多的假先知假師傅.
其實每個時代都會有.
但是在耶穌再回來以先.
這個情況是會越來越更加混亂.
我們不要曲解神的道.
我們要好好彰顯神的真理.
正如我們今天.
我們很好的來和人分享救恩.
我們不是要譁眾取寵.
我們要帶領人去信主.
是認真地讓他明白思考道.
我們每一個人為甚麼有需要呢.
因為人人都確是一個罪人.
我們今天作為神的兒女.
不是說我們站在一個道德高地.
我們去體輕別人 不是.
我們是告訴別人.
我們都是一個蒙了憐憫的一個人.
我們是蒙了上帝的恩典和拯救.
所以當我們和人分享救恩的時候.
並非說我們好像來體輕對方.
我們要用真正的愛.
既有真理.
但也都用真正的憐憫.
我記得很多年前.
當我在外國生活的時候.
我們在星期日下午完崇拜之後.
我們會帶領姐妹去商場.
我們會有福音報道.
當我們去帶人信主.
在這個過程有一次.
有一個小隊他和我說.

$^{401}$牧師我們遇到有一個太太.
她似乎有興趣想再進一步聽.
我們不知道怎樣和她說.
可不可以和我們一起去家訪.
我說好.
結果我和他們一起去探訪這個太太的時候.
她告訴我們原來她還有一個丈夫.
我說你丈夫為什麼不出來見我們.
原來他丈夫有重病.
她說他在房裡.
他不想見客.
我們不強人之難.
我們繼續在客廳.
和這位太太分享我們的主耶穌.
當我們越來越分享的時候.
我們覺得好像很迫切.
甚至想邀請她來決志.
但她都說還不行.
她說要等一等.
我們就不勉強她了.
唯有告辭.
怎知道之後她又再來和我們說.
她想我們再去一次她家.
我們覺得是不是發生了什麼事.
原來在我們走了不久.
有另一個小隊.
也在商場接觸過這位太太.
你知不知道那個小隊是屬於什麼教會.
原來他們是耶穌,基督,末世,聖徒,教會.
每個字其實在聖經都找到.
不過他們是基督教的一個大異端.
有錢有勢的摩門教.
她說的耶穌和我們的耶穌是差天動地.
結果我們很奇怪.
為什麼她又會叫我們再回去呢.
原來她的丈夫雖然沒有接客.
沒有來和我們見面.
第二批的這些摩門教徒和他們的長老.
和這位太太去傳他們的福音的時候.
她丈夫在這個房間聽得一清二楚.

$^{441}$聽完之後送了客走了.
他和他太太說.
你叫回第一班那些人來.
你信第一班那些人.
原來病人,患癌症的人.
其實他的內心很清晰.
他每一句都聽到.
而他們所說的耶穌是又真又活的.
相反摩門教的耶穌是有分等級的.
他說天父最大的.
耶穌是第二級.
聖靈是第三級.
他們是多神主義的.
當我們來到他們講.
他們都嘗試拿著他們的魔門經.
就說這本魔門經.
十九世紀是近代.
比這本聖經更超越了.
但我們說這些人是扭曲聖經的真道.
我們不可以和他有妥協的餘地.
我們必須要今天仍然堅守.
在神給我們的托付裡面.
我們忠心和人分享.
我們也看到今天.
其實我內心裡覺得最傷痛的一件事.
就是整個世界好像真的黑了眼.
今天有很多人說這個世界是.
後真理,後真相的時代.
人們不在乎這些東西是真是假.
他們說我喜歡是否真.
我喜歡就行了.
真的很恐怖.
最近我不知道大家有沒有留意到一宗新聞.
不是很久之前.
在十月中.
在美國Pennsylvania的一個城市裡面.
就在一個地鐵站.
竟然有一個女士.
是被一個三十五歲的流浪漢.
性侵犯.

$^{481}$我是說得比較斯文一點.
不是真是普通的非禮.
這個說的是強姦.
在哪裡呢?.
在一個地鐵車廂裡面.
足足八分鐘.
竟然周圍是有人的.
沒有一個人施予援手不入耳.
沒有人打911去通知警察.
甚至有人只是拿著手機直播.
like 放上facebook.
很不恐怖嗎?弟兄姊妹.
為甚麼今天我們這個世界會變得.
這麼沒有人性.
甚至我們見死不救呢?.
後來為甚麼警察會找到上門呢?.
是因為地鐵車廂都有閉路電視.
當那些職員發覺不對勁.
車廂發生一些很嚴重的事情.
結果立即通知警察.
就在另一個車站截停了.
跟著上車.
立即當場逮捕這個流浪漢.
我很想告訴弟兄姊妹.
整個世界都再不以上帝為神的話.
就因為人心裡頭黑了眼.
我們為自己著做很多偶像.
我們為自己來到不是挑剔偶像.
我們自己變了上帝.
我們自己做了一切標準的判斷.
很恐怖.
如果我們今天不悔改.
不回到神的面前的時候.
我們就沒有辦法去忠誠回應上主這個清晰的召命.
我們看見我們今天要忠於主所託.
我們就要放膽無懼.
我們要回想你幾時信主.
你幾時在神的面前再一次經歷祂的恩典.
上一次你自己打開聖經.
你去靈修.

$^{521}$你去祈禱.
幾時聖靈觸動你的內心?弟兄姊妹.
幾時你回到神的家.
你聽到師班唱詩.
你聽見弟兄姊妹帶我們去歌頌神的時候.
你心裡頭是有一種很強烈的感動.
你的眼裡面會充滿著淚水.
是因為你知道神的靈.
跟著祂的話.
藉著不同的思覺.
你在感動我們.
我們今天我們回想一下.
我們不是全自己.
我們是全由耶穌基督為主.
我們同作是耶穌基督的僕人.
浸信會很強調.
是使徒皆濟私.
我們宣道會.
我相信大家也很重視的.
就是整個的群體.
我們一起無論在宣教.
在醫治.
在這種的社會關懷裡面.
在傳福音.
我們都要好好完成神給我們的託付.
我們播道會一樣.
我們很著重傳福音和社會關懷.
我們今天神給我們各樣的恩典.
在基督的身體裡面.
我們是互為之體.
雖然我們屬於不同的宗派.
但我們在基督耶穌裡面.
我們可以合而為一.
我回想我昔日夢照的時候.
是我當我自己將要進入大學.
來讀我的工程系.
當時那個書.
我每天如常翻開聖經靈修.
讀到這個歷代至上28章.
是大衛王對他的兒子所羅門一個祝福.

$^{561}$好像一個臨終的託付.
他說:我兒啊所羅門.
你要誠心樂意的事奉耶和華.
因為神揀選了你.
要建造他的奠宇.
你要好好地來跟隨神.
如果不是的話.
神亦都會凋棄你.
是有應許有警告.
我當時看見這段經文.
我自己內心有很大的一個感動.
其實我還很記得.
在那天靈修之前那一晚.
是我去了多倫多的一個宣道會.
當晚是一個猜拳的陪靈會.
講完就是藤劍輝牧師.
其實當時我沒有回應那個呼召.
但是我想不到.
原來上帝竟然在第二天早上.
祂要我很清楚地.
回應祂自己的揀選.
我自己很懼怕.
我覺得上帝.
我自己不是能言善辯的人.
我覺得我沒有口才.
我也很不習慣.
在這麼多人面前去講說話.
我是比較屬於一些.
所謂的怕醜的人.
但是聖靈不斷在我內心.
有一個很清楚的感動.
祂說Peter.
我現在不是問你有多少能力.
我是問你願不願意.
Are you available?.
這個呼召來得很突然.
但是亦都很清晰.
當時我在這個感動底下.
我自己說上帝.
你可能揀錯了人.

$^{601}$應該不是我.
我以為自己.
在從小上教會的那些.
基督徒應該明白我所講.
我們好像屬於一些乖乖的小朋友.
好像覺得今天早上靈修.
同神有個交往.
其實這個不是正確的靈修觀念.
靈修不是你自己早上打開聖經的那段時間.
靈修是我們整個人的生命.
我們每一日生活24小時.
你都可能要有這種交往.
當時我自己看完這條聖經.
有感動.
但是我想消滅聖靈的感動.
我就蓋上這本聖經.
我就跟著想做一件什麼事.
我就想祈禱.
多謝天父帶領我.
我想好像交完功課一樣.
你知不知道我發生什麼事.
當我這樣做的時候.
竟然我想開口說話.
我一句話都說不到.
是真的.
我一句話都說不到.
我啞口無言.
然後神的靈繼續在我內心.
很大的催迫和感動.
直到我看見.
我說上帝這個若然是出於你的感動.
我願意求主幫助我.
我信心不足求你幫助我.
當我這樣說的時候.
我淚如雨下.
我簡直心裡像一個大石鑼開.
好像被上帝的靈充滿著.
是一個愛的環繞.
一個很真實的見證.
和一個感動.

$^{641}$我回想這個已經是發生很多年之前.
我不告訴你多少年.
免得你猜我的歲數.
但確實這件事是很震撼.
對於我來說.
我人生我現在已經是.
由神樂園畢業到現在已經32年去傳道.
但在這個過程裡面我回看.
每一次遇到任何困難.
叫自己不要灰心.
不要沮喪.
不要所謂的懼怕.
我回想神在最原初那一點裡面.
祂是怎樣呼召我.
怎樣來給我一個肯定.
我要感恩.
我要和弟兄姊妹繼續竭力.
去完成上帝給地上教會的大使命的託付.
我們看中間這個段落的時候.
由第7到第12節.
我們一起來讀.
我們有這寶貝在眼裡.
為要顯明莫大的能力.
是出於神.
不是出於我們.
我們處處受困.
卻不被困住.
內心困擾卻沒有絕望.
遭受迫害卻不被拋棄.
的倒在地卻不致滅亡.
我們身上常帶著耶穌的死.
使耶穌的生也在我們身上顯明.
因為我們這活著的人.
常為耶穌被置於死地.
使耶穌的生命在我們這必死的人身上顯明出來.
死是在我們身上運作.
生卻在你們身上運作.
我們在這個段落裡面.
我們都很清楚看見.
保羅和我們分享.

$^{681}$他事奉主的艱難.
但他怎樣能勝過呢.
當我們看第7到第9節的時候.
我們要明白保羅當時在第一世紀的背景.
其實教會是充滿著內憂外患.
教會裡面有賈師傅在傳別的福音.
但教會外邊亦有猶太教的信徒.
在挑唆眾人去逼迫使徒.
有些人想抹黑保羅的身份和權柄.
但保羅和其他使徒.
能夠堅忍前行服事.
既打不死亦打不倒.
是因為他用這個比喻.
讓我們看見這個啞氣.
這個身體裡面是有耶穌基督.
作為生命寶貝的這種能力.
這種莫大的能力.
就是出於神.
以至我們不怕任何的迫害.
保羅在第一次宣教旅途裡面.
我相信大家都記得.
在《使徒行傳》記載.
他在路斯德.
他遇到猶太人煽動群眾.
用石頭打他.
打到他真的面流很多的血.
人們以為他真的死了.
拖到城外.
誰知保羅當時傷勢雖然嚴重.
但他醒過來之後.
他不是走去休息.
他竟然進入城裡面.
繼續去堅固門徒的心.
保羅這一份傳福音和建立.
所牧養鼎子妹的心志.
再再表明他愛神愛人.
這一份牧者的心腸.
可以稱為一個為父的心.
鼎子妹今天這個知識爆炸的時代.
其實是作師父的有千千萬萬.

$^{721}$但為父的卻不多.
今天你願不願意成為教會裡頭.
一個屬靈的父親.
屬靈的母親嗎.
你願不願意去關心我們的下一代.
一些的新生代.
好好用心機.
用神的話.
用愛心去扶植他們.
來成長起來.
這個也是我很想帶給大家的一個挑戰.
在第十到第十二節.
這裡所說的是一個.
先死後生的屬靈原則.
唯有當人願意天天背起十字架.
去跟隨主.
每日對罪和世界看自己是死的.
我們才能夠活出耶穌基督復活的大能.
同樣得著這一份生命的祝福.
當我們看到經文最後這個段落.
在第三段.
我們一起來.
或者先讀了這一段聖經.
我們一起來讀.
經上記著說.
我信所以我說.
我們既然有同樣的信心也就信.
所以也說話.
因為知道主耶穌也與我們一同復活.
並且擺我們和你們呈現在他的面前.
這一切都是為了你們.
好使因為既然恩著許多人而增多.
感恩的心更加增多了.
使榮耀歸給神.
所以我們並不沮喪.
我們外面的人雖然漸漸扭壞.
但裡面的人卻日日更新.
因為我們短暫輕微的患難.
是要為我們成就極大無比永遠的榮耀.
我們所顧念的不是看得見的.

$^{761}$而是看不見的.
因為看得見的是暫時的.
看不見的是永遠的.
頂詞目我們最後看這段經文的時候.
讓我們都來明白.
保羅他說我們需要有實在信心.
首先就是我們需要有同一位的聖靈指教我們.
他這裡引用的是舊約裡面詩篇116篇.
所說的.
詩人說死亡的繩索纏繞我.
陰間的痛苦抓住我.
我遭遇患難受苦.
主啊!你救我的命了.
免了死亡.
救我的眼免了流淚.
救我的腳免了跌倒.
我要在耶和華面前行活人之路.
我因信.
所以如此說話.
我受了極大的困苦.
保羅說在詩篇116篇.
詩人曾經歷過聖靈所說的這番說話.
我同樣和你們既然有同樣的信心.
我們就相信不移.
我們今天要學習的.
就是我們以前.
等於四十年前.
錦宣堂這裡.
每一位是創會的牧者.
和教會的領袖.
他們都是因為神給他們的意象體見.
弟兄姊妹我們需要有同樣的這些所謂先賢的信心.
效法他們的至死忠心.
去跟隨主.
服侍主這一份的心志.
我們懷著同樣的信心.
看見一件事情就是.
耶穌基督.
我們所相信的是同一位基督.
而我們要復活的時候.

$^{801}$是同一班的弟兄姊妹.
與基督是一同復活.
我們是帶著這一份榮耀的盼望.
在最後這個大結局的時候.
我們會看見.
就是神給我們看見這個遠像.
這個遠像就是外面的人和裡面的人.
那個不同之處.
保羅在他結束整章經文的時候.
他要我們顧念.
顧念那個原文的意思就是.
叫我們定睛誓看.
用心靈信心的眼睛.
看見屬靈的榮耀和真實.
如果一個人.
他是真正看到永恆榮耀的真實遠像.
他就會甘心樂意.
毫無保留地去毀身事奉主.
所以在最後這一段經文裡面.
出現的這五個對比.
是很值得我們去了解的.
第一世紀的時候.
希羅文化裡面的文學.
常常以為是將人分為靈魂和身體.
這樣義分.
甚至有些哲學家.
是看靈魂是被困於卑賤的身體當中.
保羅在這一段經文要肯定的.
不是這件事.
他要肯定的是生命的優先次序.
不是要將身體的價值貶低.
因為神在原初的創造裡面.
一切的物質世界的生命.
都是視為好的.
可惜人類的始祖亞當夏娃犯罪之後.
罪就進入.
以及去污染了整個受造界.
聖經的人觀是完整的.
是包含著靈魂體整個的結合.
所以我們今天看這段經文最後的段落.

$^{841}$同樣是很重要的.
弟兄姊妹我們要看到這個榮耀的遠像.
不是一個抽象不可觸摸的屬天世界.
而是我們看見這位道成肉身的耶穌基督.
祂曾經為人.
祂在歷史上按照神的心意.
祂來為我們眾人死.
也都為我們眾人復活.
為我們眾人升天.
祂說祂一定要再回來.
就是因為這一種屬靈的透視力.
我們今天要懷著一個剛強壯膽的心.
我們要效法保羅.
這一份的基礎.
讓我們可以看透.
是因著我們已經與耶穌基督.
有同一復活的指望.
所以我們這個身體雖然會衰殘會扭壞.
但是我們可以耐心一天生持一天.
在後文五章一節裡面.
我們最後讀這一節聖經.
我們也知道.
如果我們在地上的帳盤拆毀了.
我們必得著從神而來的居所.
那不是人手所做的.
而是天上永存的房屋.
求神幫助我們.
讓我們好好來認清.
我們今天不要喪膽,沮喪,灰心.
我們要努力的去服侍主,服侍人.
帶著這一份清晰的使命.
生命的能力和實在的信心.
讓我們可以成為別人生命的幫助.
最後我將一個運動員的見證放在大家面前.
這個發生在1968年墨西哥的奧運會裡面.
來自坦桑尼亞的這位運動員.
他是走一個所謂馬拉松的長跑.
當時他跑的時候跌傷了.
膠也掉了.
他要紮住繃帶.

$^{881}$有人勸他應該退出比賽.
當時其實是有一班的運動員.
其實都是和他一樣.
可能因為受傷緣故.
他們沒有完成那個賽事.
但是這位John Avary.
他沒有放棄.
他走到終點的時候.
比跑過終點的那個冠軍.
已經遲了一個多小時.
那些人以為已經散了.
但是竟然見到他在日落西斜的時候.
竟然帶著一個這樣的受傷的身軀.
來到的時候.
那些記者訪問他.
他說John.
為什麼你還要完成這個賽事.
還要跑回來.
其實你棄權就可以了.
那麼辛苦做什麼呢.
你猜他怎麼回答他們呢.
他說我的國家派我來這裡.
由五千里以外來到這裡.
不是來開始比賽的.
而是要完成這個比賽.
弟兄姊妹.
今天我們人生都一樣.
上帝不是派我們來這個世界上.
只是出世.
過幾十年就算了.
神是要我們好好的完成.
我們在地上這個所謂生命的賽道.
我們有沒有好像保羅一樣.
我們好好持守我們的真道.
要跑的路我們已經跑過了.
要打的仗我們打過了.
我們問心無愧.
到將來我們見主面的時候.
我們可以來到去.
是完全不是帶著任何遺憾.

$^{921}$或者羞愧去見主.
到時我相信我們每一個人.
都可以跟我們的主說.
是神的恩典救我們用.
因為祂的能力.
是在人的軟弱上顯得完全.
我們都好像保羅一樣.
要跨我們的軟弱.
好叫基督的能力去奉庇我們.
今天全世界的教會都出現一個.
很真實的問題.
就是高靈化的現象.
我不知道大家同不同意.
我見有些弟兄姊妹已經笑笑口了.
我就去不同教會講道.
我都發現.
真是很多教會.
包括我們自己教會一樣.
有青少年人流失.
但我們過去這五年已經竭盡所能.
我們要堵塞這一切的漏洞.
我們真的要懷抱.
關心我們的下一代.
所以過去這兩年.
我們特別不單止教會要年青化.
我就強調一樣東西.
我們要跨代同行.
跨代同行的意思就是.
我們的年長一輩的弟兄姊妹.
不要看自己.
你以為我一把年紀了.
退休了.
其實我們每一個都可以做一個退而不休.
是因為你的內心是充滿著神的愛.
神的恩典.
你繼續可以和年青人一起去服事.
最近我在教會推動一件事.
就是我們要拍短片.
去將我們教會裡面一些的弟兄姊妹.
生命的見證.

$^{961}$講好耶穌的故事.
就是這個系列.
我們現在有一個標題叫「神奇故事」.
歡迎大家上網去參觀.
這個神奇的奇就是祈禱的奇.
第一集已經拍了.
第一集的主角.
是要述說耶穌在一位姊妹身上的奇妙的工作.
我現在口上戴著的口罩.
叫聖靈果子的口罩.
是第十週.
我們由仁愛,喜樂,和平.
一直數到這個就是節制.
最後那個.
聖靈的果子有多少個呀?弟兄姊妹.
不是九個.
聖靈果子是一個.
不過是九種特性.
這個第十週就九種特性都在這裡.
我們將這個姊妹.
為何她明明是在做生意.
轉行做口罩廠呢?.
所為何事呢?.
是神給她甚麼感動,甚麼啟示呢?.
當她來做這個口罩.
背後都是一個榮譽.
就是這些都是聖經的說話.
這麼有趣的.
我們之前有一個節制.
就好像青蛙那種綠色的.
有人問我.
牧師為何好像一隻青蛙那樣.
不是青蛙王子.
我說這個就是.
好比當我們有時見到人家.
你自己心裡頭牢嘈.
你自己煩躁.
你想發脾氣的時候.
聖靈又要提醒自己.
你停一停.

$^{1001}$上一上.
你不要像青蛙那樣.
跳上別人.
跳上不同的物體.
我們很多時候以為用我的手段.
用我的方法.
你好不好安靜下來.
停一停.
被聖靈來感動我們.
平靜我們的情緒.
我們同人的關係.
很多事都可以復和的.
不是需要這樣來撕裂的.
弟兄姊妹.
求主幫助我們.
施恩於我們.
好不好呀?弟兄姊妹.
跨代同行.
不要說自己年紀老邁.
現在世衛六十五歲.
才踏入中年.
是啊.
所以我們當中.
我見到很多都是年青人.
你們都是六十五歲或以下.
六十五到七十五歲才是中年.
所以我們要有健壯的身體.
有好的靈性.
前面的四十年.
若主未回來的時候.
我們都要忠心去事奉祂.
我最後用一首詩去做結束.
這首詩就是靈實醫院的院長.
梁志達醫生.
我和他都有很好的交往.
他也是一個很好的朋友.
他翻譯做英文.
不過我將他的中文放在大家面前.
他說在神眼中.
老人永不會沒用.

$^{1041}$外體雖然會衰退.
但內心卻日日更生.
越來越明白他話語的真實.
家多住木永恆.
老年人仍可以在信望愛中被建立.
在默想聖言和禱告的甜蜜中被造就.
仍可以同奔天路.
向著彪光直跑.
繼續敬拜讚美和服侍我們的恩主.
我們愛神的人永不會太老.
在他的角度裡永不會沒用.
讓我們彼此勉勵.
翹首永恆.
好不好 弟兄姊妹.
我們都是神的教徒的年青人.
我們要跨代與我們的下一代同行.
願主賜福你們每一個.
賜福曾牧師和你們整個團隊.
都能夠讓教會走起神的喜悅和心意.
當眾我們一起結束祈禱.
天父上帝我們很感恩.
你過去四十年怎樣使用錦宣堂.
我們很相信在前面的日子裡.
你要大大繼續賜福和使用他們每一個.
讓他們忠心在你的聖言上.
讓他們能夠每一個弟兄姊妹.
好發揮神賜他們的不同恩賜.
在神的國 在這個社區.
在香港整個社會.
甚至在普世的差傳宣教裡.
都能夠注著上帝你自己美好的計劃和心意.
繼續好好去祝福各族各方各民.
能夠成為主的門徒.
被主稱讚.
我們同心合意祈禱.
奉耶穌基督得勝的名堂.
阿們.
\newpage



\section{詩篇 23:1-6}
\label{sec:zUN5t6YF3yU}
\textbf{2021.11.07 《好牧者之歌》 詩篇 23\_1-6 (和修版) 曾敏姑娘}
\newline
\newline
連結: \href{https://youtube.com/watch?v=zUN5t6YF3yU}{\texttt{https://youtube.com/watch?v=zUN5t6YF3yU}} ~~~~ 語音日期: 2021-11-08
\newline
\newline
\hyperref[sec:L2pB3hPJTOs]{\small{< < < PREV SERMON < < <}}
~
\hyperref[sec:index_chronic]{\small{[返順時目]}}
~
\hyperref[sec:index_scriptual]{\small{[返順卷目]}}
~
\hyperref[sec:TujTTaedR4Y]{\small{> > > NEXT SERMON > > >}}
\newline
\newline
詩篇 23:1-6
\newline
\begin{longtable}{cl}
\hline
\hline
章節 & 經文 (和合本修訂版)\\
\hline
23:1 & \begin{tabularx}{0.7\textwidth}{X} 耶和華是我的牧者，我必不致缺乏。 \end{tabularx} \\ \\ \relax
23:2 & \begin{tabularx}{0.7\textwidth}{X} 他使我躺臥在青草地上，領我在可安歇的水邊。 \end{tabularx} \\ \\ \relax
23:3 & \begin{tabularx}{0.7\textwidth}{X} 他使我的靈魂甦醒，為自己的名引導我走義路。 \end{tabularx} \\ \\ \relax
23:4 & \begin{tabularx}{0.7\textwidth}{X} 我雖然行過死蔭的幽谷，也不怕遭害，因為你與我同在；你的杖、你的竿，都安慰我。 \end{tabularx} \\ \\ \relax
23:5 & \begin{tabularx}{0.7\textwidth}{X} 在我敵人面前，你為我擺設筵席；你用油膏了我的頭，使我的福杯滿溢。 \end{tabularx} \\ \\ \relax
23:6 & \begin{tabularx}{0.7\textwidth}{X} 我一生一世必有恩惠慈愛隨著我；我且要住在耶和華的殿中，直到永遠。 \end{tabularx} \\ \\
[1ex]
\hline
\hline
\end{longtable}
$^{1}$弟兄姊妹平安.
今天神讓我們在這裡有福氣.
來敬拜親近她.
我們開始的時候再有一個禱告.
真的很愛我們.
我們這麼卑微的人.
能夠來到你的面前.
實在是很不配.
藉著美娟的禱告提醒我們.
我們生命裡有很多軟弱.
有很多罪惡.
很願意神光照我們.
寬恕和潔淨我們.
願意我們這一刻.
就是守潔心清的來到你的面前.
你對我們每一個說話.
願意聖靈在我們每一個人的生命裡工作.
讓我們真的聽到上帝你的聲音.
我們又能夠作出回應.
奉耶穌基督的名字祈禱.
阿們.
11月份就是我們的堂慶月.
這句話今天已經不是第一次說了.
我們會紀念紅人堂十週年的堂慶.
這是大日子.
過去這十年的時間.
有很多人在紅人堂這個地方.
留下了他們的故事.
或者我們的故事.
很多人應該包括弟兄姊妹和全道同工.
全道同工是我們的牧者.
對我們的生命影響非常深遠.
我不知道有哪位牧者.
在你心中留下了很深刻的印象.
十週年是一個時間去回顧.
除了地面上的牧者.
我們還有一位在天上的牧者.
這一位牧者對我們來說是最重要的.
因為他牧養我們的時間.
是從過去到現在.

$^{41}$直到永遠.
是最長最長.
這一位天上的大牧者.
絕對是一位好牧者.
他有多紅呢?.
今天就讓我們透過這首寶貴的詩篇.
讓我們去思考.
詩篇23篇被稱為大衛的詩.
有不少學者認為是大衛逃避他的兒子.
壓沙龍追殺的時候寫的.
在這裡讓我們稍稍說一下.
大衛的生平事蹟.
讓電影姐妹可以去重溫.
伯利恆人耶西的兒子.
就是大衛.
大衛對上有七個哥哥.
很厲害.
大衛原先是放羊的.
他三個最大的哥哥.
就跟著當時第一任.
以色列第一任的軍王.
索羅出去打仗.
當時一般人都覺得.
打仗的比較有出色.
放羊的呢?.
如果對比打仗的.
就比不上打仗的那麼威風.
所以大衛的哥哥.
當時的人會覺得.
他們是比大衛更加有出色的.
但是上帝卻偏偏揀選了.
大衛作為以色列第二任的軍王.
當時非利士人哥利亞.
向以色列盟罵戰.
如果有看過聖經裡面.
這些故事的電影姐妹就會知道.
哥利亞真是一個巨人.
當他罵戰的時候.
索羅王和以色列眾人.
都感到很害怕.

$^{81}$沒有人願意出去迎戰.
大衛就勇敢地去挑戰哥利亞.
將她打敗和殺死.
婦女唱歌稱讚大衛.
更甚於索羅.
大家還記不記得那首歌是怎樣說的?.
誰千千誰萬萬?.
真是很雷.
原來大家很熟悉大衛的故事.
索羅殺死千千.
大衛就殺死萬萬.
真是這一下不得了.
招惹了索羅的妒忌.
過了一段時間之後.
索羅追殺大衛.
讓大衛進入人生第一次的逃亡當中.
神一直保護大衛.
索羅沒有辦法去捉拿大衛.
索羅死了之後.
過了一段時間.
大衛成為了以色列的君王.
他統一全國.
有一天.
他睡到黃昏才起床.
當他在皇宮的平頂散步的時候.
他看到一個容貌非常美麗的女人.
在那裡沐浴.
後來大家都知道了.
故事發生就是.
他和這個女人發生了關係.
這個女人其實是大衛一個非常忠心的手下.
烏莉亞的太太.
女人或者人.
大衛很想隱瞞這件事情.
希望不斷地用一些方法造成局勢.
令人覺得這個孩子.
是女人和她的丈夫烏莉亞所生的.
結果他不成功.
他就借刀殺人.
讓烏莉亞戰死沙場.

$^{121}$神非常不喜悅這件事.
就讓先知拿丹去責備大衛.
最後拿丹說.
刀劍永遠都不會離開他的家.
神也會從大衛家裡興起災禍攻擊他.
結果大衛的家就接二連三.
發生了一些很悲慘的事情.
大衛的孩子彼此傷害.
厄沙隆就是大衛的兒子.
他很想做王.
所以他起來背叛和追殺大衛.
詩篇23篇.
就是大衛在逃避厄沙隆的時候所寫的.
大衛在一個特別的時勢.
正正要透過詩篇.
去描述神和他之間的關係.
他的眼目再次望著上帝.
再提醒自己.
這個神是一個怎樣的神.
他要去依靠祂.
在詩篇裡面我們可以看到.
神是一位好牧者.
詩篇23篇.
大衛用了兩幅圖像去表達他想說的內容.
第一個圖像是一到四節.
就是牧羊人和羊的圖像.
第二個圖像是五到六節.
是君王和戰士的圖像.
但我想跟大家說.
其實這兩個圖像是相通的.
並不是分開的各有各講.
其實在舊約裡面.
他們去看君王.
經常會想起牧羊人和牧者.
所以君王和牧者這個角色.
不會說是彼此分開.
基本上大家差不多會走在一起.
在這個過程中.
我們透過這六節經文去看.
好牧者是怎樣的.

$^{161}$第一件事.
好牧者和羊有一個很親密的個人關係.
第一節.
耶和華說牧者必不辭缺乏.
牧羊人和羊有一個很親密的個人關係.
去認識牠的羊.
在外國.
有些牧羊人會為羊起名.
不知道這裡有沒有.
弟兄姊妹去外國旅行.
有見過這樣的情況.
牧羊人為羊起名.
例如多利 瑪利 比特 約翰.
比特 約翰也是嗎?.
會的.
因為你記住Peter這個名字.
也是很多人叫的名字.
那隻羊也叫Peter.
當他呼喚羊到他身邊的時候.
就說多利 多利 多利 過來呀.
Peter 過來.
我們看羊很像樣子.
很難分辨.
這隻和那隻好像差不多.
但是牧羊人不會弄錯羊.
也不會弄錯名字.
因為他認識牠們.
知道牠們每一隻的特色.
例如.
他很明白牠們的手腳.
眼耳口鼻.
毛色.
羊的性格有什麼不同.
不單是這樣.
牧羊人掌管羊的一生.
牧羊人帶領著羊.
有不同的經歷.
那羊又怎樣呢?.
牠們又認識牧羊人.
當牧羊人叫牠們的名字的時候.

$^{201}$John.
牠們就會跟著.
牠們會聽的.
牠會認住這把聲音.
牠會跟著牧羊人走.
對於這些羊來說.
牧羊人就是牠們唯一的主人.
一生的幸福都在牧羊人的手上.
跟著牠們可以吃.
跟著牠們就有好的羊生.
牧羊人和羊的關係是親密的.
耶和華神完全認識大衛.
很認識牠的樣子.
因為是上帝造牠的.
很認識牠的性情.
牠的優點和缺點.
牠做得好的地方.
牠在什麼地方跌倒.
牠都知道.
神很明白大衛的心在想什麼.
祂掌管大衛的一生.
而大衛也認識神.
他認識神的性情.
神是怎樣的神.
他知道大衛有一樣東西.
很討神喜歡.
他經常想著神.
很願意跟從神.
他稱神為他個人的牧者.
就好像他擁有這個牧者一樣.
我的牧者.
我擁有他.
我們和神的關係是怎樣呢?.
這是值得我們思考的.
神一定認識我和你.
因為是祂造我和你的.
祂知道我們的一切.
我們在祂面前.
我用這個形容詞.
簡直是透明的.

$^{241}$不要以為只有你自己知道.
上帝是完全知道.
祂掌管我們的一生.
那我和你又如何呢?.
我們認不認識這個神呢?.
這正是我預備這篇講章時的一些思考.
認識不只是頭腦上的認知.
我們上教會上了很久.
很容易說.
別人問你上帝是怎樣的.
我知道.
上帝是慈愛的.
他做了甚麼.
12345678.
上帝也是公義聖潔的.
說這些話很順口.
因為這些話上了教會很久.
就懂得說了.
我們認識不只是頭腦上的認知.
而是這個頭腦上的認知會帶來實際的回應.
例如當我知道神是愛的時候.
我們就很渴望去親近祂.
當我們知道神是聖潔和公義的時候.
我們在祂面前就要存著一個敬畏的心.
要遠離罪惡.
我們的心會不會麻木了呢?.
連我們自己陷在罪惡裡都不知道?.
我們有沒有常常反省自己的生命.
在神面前認罪悔改.
真正的戰經來到神面前?.
這個值得我們問.
如果頭腦上的認知沒有帶來實際的回應.
這個不是真正的認識.
我們的心有沒有像大衛那樣.
常常想著神呢?.
我們願不願意跟從神呢?.
我們是否真的擁有神呢?.
神認識我們不代表我們關係就親密.
我們要同樣認識神.
貼近神.

$^{281}$這樣才會親密.
我們是否有需要重建我們跟神的關係呢?.
在十週年的慶典裡.
沒錯,我們會有很多的表演.
在過去帶來我們事工的發展.
弟兄姊妹,這正是一個好的機會去思想.
我跟神的關係是怎樣.
不要進了寶山空手而回.
我們親近神,來到神面前這麼久.
得著什麼呢?.
你是不是得著神呢?.
你是不是被神得著呢?.
這十年來,我們只有活動嗎?.
如果我們要去說.
我們跟上帝的關係,我們能說些什麼呢?.
說真正的關係.
可以說什麼呢?.
這是值得我們思考的.
未來還有很多的十年.
我們想怎樣走下去.
是更多的活動?.
還是我們想重建跟神的關係?.
這是值得我們去思考的.
好牧者和他的羊有一個親密的個人關係.
第二點,好牧者是怎樣的呢?.
很眷顧他的羊.
這是詩篇23篇說得最多的一個位置.
好牧者很眷顧我們.
23篇一字說.
耶和華是我的牧者.
我必不自缺乏.
一般來說,缺乏這個動詞後面.
都會跟一些名詞的.
例如缺乏食物,缺乏衣服.
衣食住行等等.
但這裡沒有跟任何名詞.
即代表我們什麼都不缺乏.
無論是物質的,心靈的.
還是靈性上的都沒有缺乏.
木羊人眷顧羊,讓他羊一無所缺.

$^{321}$就正如神眷顧大胃.
讓大胃一無所缺.
這裡我們要看看.
木羊人是怎樣眷顧羊的呢?.
詩人從三個層面去解釋.
神的眷顧,是木羊人的眷顧.
第一,木羊人是眷顧羊身心靈的需要.
23篇二節這樣說.
他使我躺臥在青草地上.
領我在可安歇的水邊.
木羊人帶領羊去尋找青草的地方.
讓羊吃得飽足.
我們留意這節經文的重點不是說吃.
重點是說躺臥.
不是說吃,是躺臥.
他已經吃飽了.
吃飽之後,原來木羊人會讓羊躺臥.
在草地上享受休息.
休息是很重要的.
休息不到,我們就會感到疲累.
情緒也會很受困擾.
休息是香港人一個很重要的議題.
每隔一段時間,香港就會公佈數字有多少人失眠.
每一晚,全港的燈都熄得差不多.
有多少人眼睛瞪.
眼睛光,睡不著.
直到第二天,感到非常疲累.
就算人睡不著.
心靈也睡不著.
有多少人?.
休息有多重要?.
木羊人想羊休息.
上帝想我們休息.
而且是享受休息.
這很重要.
木羊人知道羊需要喝水.
所以帶羊到水邊.
這不是一般的水邊,而是安靜的水邊.
因為木羊人知道羊的一個秘密.
就是羊怕死.

$^{361}$一點聲音也會嚇到牠們.
如果帶牠們去急流的水邊.
羊就會很害怕.
所以帶牠們去安靜的水邊.
讓羊感到.
就算喝水也不用被驚嚇.
23篇的三節這樣說.
上半節.
他使我的靈魂蘇醒.
靈魂在這裡不一定是我們一般所說的.
我們人的靈魂.
而是我們自己.
去蘇醒我們自己.
神復蘇我們的生命.
這非常寶貴.
我們生活在地面上.
有時大家有沒有感覺.
好像死了一樣.
整個人.
太痛苦了.
太不開心,好像死了.
神復蘇我們整個生命.
這個含義是非常深和廣.
牧羊人很努力滿足我們身心靈的需要.
他的細心.
令我們感到快樂.
上帝也是這樣供應大衛的需要.
令他感到很滿足.
可能在這一刻你會說.
通常應該是那些比較富裕的人.
生活沒有憂慮的人.
什麼都有的人.
比較容易感受到詩篇23篇這裡說.
牧羊人滿足羊.
事實不一定是這樣.
可能是有缺乏的人.
更加能夠經歷到上帝的供應.
大衛經歷過人生的低潮.
當中包括了兩次逃難的經歷.
剛才大家還記得我提到他人生的故事.

$^{401}$第一次就是被掃羅追殺的時候.
第二次就是被他的兒子亞沙隆追殺的時候.
在第一次逃難的經歷當中.
大衛曾經向羅伯的亞希米勒祭司.
和在加密的大富翁.
這個拿伯討食物.
由此可見.
他是曾經經歷過缺乏.
應該不止一次.
最後神很奇妙.
透過不同的人供應了他.
讓他渡過難關.
在這裡我想說一個.
我很欣賞的屬靈偉人的經歷.
大家有沒有聽過.
喬治·穆勒的故事.
穆勒.
是,穆勒的故事聽過吧.
當中有些反應.
他是一個很好的屬靈偉人.
是一個德國人.
在1805年出生.
剛開始的時候.
他也是一個傳道人.
到了32歲的時候.
上帝就感動他.
去設立孤兒院.
他就按照這個感動去做.
但是他有些原則.
他不宣傳.
他不會告訴別人.
我現在在一個孤兒院.
各位善長仁翁.
快點慷慨解囊.
拿你的錢出來.
我們孤兒院有很多需要.
講明孤兒院.
我們是很貧乏的.
很缺乏的.
給錢吧.

$^{441}$有沒有感動.
給錢吧.
他不做這些宣傳.
他不會說這些.
可能你說.
他不說是因為他有錢.
是吧.
所以他就這麼信心.
沒有人資助他都可以.
又不是.
穆勒本身自己就是一個窮人.
他很貧窮.
他沒有什麼米飯幫主在後面.
沒有一個強而有力的後台.
沒有.
他只是專心禱告上帝.
以靠神的公營.
他剛剛開辦孤兒院的時候.
院內只有20個孤兒.
但是經過了61年的時間.
你猜這麼多年.
孤兒院收留了多少孤兒呢.
有沒有人想嘗試猜.
億多.
那又太多了.
太多了.
多少.
二千個.
誰說一萬的.
一萬接近了.
好東西.
一萬接近了.
本身開始的時候.
20個孤兒.
經過61年.
可以收容到9744個.
這個不是穆勒唯一的工作.
他還有建立一些貧窮的學校.
他有的.
做了很多這些.

$^{481}$他這個人很窮.
他會去做這些.
不是因為我有錢去做.
窮我都會去做.
為什麼.
專心依靠上帝.
他在事奉上是很深刻經歷.
神的公營.
他也見證著上帝賜福給那9744個孤兒.
還有貧窮學校的學生.
他不需要自己有個銀行去做這些.
他的銀行就是上帝.
上帝有源源不絕的錢.
他就專心依靠上帝.
耶和華是穆勒的牧者.
他也是孤兒的牧者.
他滿足他們每一個的需要.
今天.
神也會滿足我們的需要.
有錢的.
沒錢的.
上帝會滿足我們的需要.
因為耶和華神.
是我們洪恩家真正的牧者.
他牧養我們.
主禱文有一句這樣說.
我們日用的飲食.
今日賜給我們.
耶穌基督就教我們這樣祈禱.
所以我們都可以向上帝這樣禱告.
我們學穆勒這樣禱告.
無論我們什麼事工也好.
無論我們關心什麼有需要的人也好.
我們向上帝禱告.
下一次我們再見面的時候.
你不要跟我說.
我沒有錢.
找些有錢的人來幫助.
那些有需要的人.
有錢沒錢.

$^{521}$我們都有上帝.
夠了.
有上帝就夠了.
在這裡我們看到.
剛才說的那一點.
好牧人去眷顧羊的需要.
第一點就是眷顧羊身心靈的需要.
第二件事.
牧羊人和羊同在.
牧羊人會帶領羊.
走在正確的道路上.
23篇3字下這樣說.
為自己的名引導我走二路.
二路是正確的路.
羊的方向感不是太好.
很容易迷路.
但只要牧羊人在身邊.
就不用怕.
因為牧羊人會帶領羊.
走在正確的道路上.
我們人都是很軟弱的.
我們很容易迷失.
我們生命的方向.
走著走著就不知道.
走了去哪裡.
大家有沒有這樣的經歷呢?.
不知道走到哪裡.
不只是這樣.
甚至我們會陷在罪惡的網絡裡.
大衛是一個土神喜悅的君王.
我經常很欣賞他對神的心.
真的覺得很好.
覺得要學習也不容易.
他很願意跟從神的教導而行.
也很倚靠神.
但他也有軟弱.
大家記不記得他有什麼軟弱?.
我剛才說他的背景.
記不記得?.
一開始我說他的背景.

$^{561}$其實他也有犯奸淫.
他也有借刀殺人.
這不是他生命唯一的軟弱.
但他也有跌倒的遺憾.
但神沒有放棄他.
還是跟他一起.
神派先知拿丹去責備他.
讓大衛知道自己犯了罪.
得罪神了.
從而要認罪和悔改.
上帝帶領大衛回到正確的道路上.
弟兄姊妹.
今天我們需要的.
不只是不停有人呵護我們.
你很不開心,是吧?.
我抱抱你.
我們需要什麼?.
也有人提醒我們.
喂!.
我們好像離開了上帝的道路.
小心啊.
這樣做是否神不喜悅?.
要回到上帝的身邊.
我們也需要有好牧者.
提醒我們走回正確的道路上.
我們不只是聽呵護的話.
需要這種帶領.
這位好牧者就是我們的上帝.
祂會帶領我們.
牧羊人也會保護和管教羊.
《詩篇》23篇4節這樣說.
我雖然行過死任的幽谷.
也不怕遭害.
因為你與我同在.
你的象,你的肝都安慰我.
牧羊人帶領羊去尋找青草地的時候.
經常會經過幽谷.
雖然幽谷是很可怕和危險的.
例如有猛獸.
其實有些學者也有提過.

$^{601}$在幽谷也有機會有毒草.
這個比較少提及.
通常是提猛獸多.
但是羊是不需要害怕.
因為牧羊人與牠在一起.
牧羊人會保護羊.
用象打走猛獸.
我們之前提到大衛在人生當中.
有兩段時間是經歷逃亡的.
一段是被掃羅追殺的時間.
另一段是被他的兒子追殺的日子.
每一段時間都不容易過.
每一段時間都有生命的危險.
但是神保護大衛.
不將他交在敵人的手裡.
到最後的時候.
神更加讓大衛勝過他的仇敵.
弟兄姊妹想問問大家.
死任的幽谷對大家來說是指什麼呢?.
每個人都不同.
我們不用跟別人比較.
對我們每個人來說.
死任的幽谷是什麼呢?.
是否一些很危險的處境.
困難的處境.
患病.
身體的.
情緒的.
做大手術.
陷入經濟的困難.
家庭關係的破裂.
很多的危機.
不知道怎樣生活下去.
是我們受傷害.
在這些位置.
我們有沒有經歷到神的保護呢?.
你有沒有感受到.
神圍著我們.
抱著我們呢?.
神絕對願意保護我們.

$^{641}$神也有能力保護我們.
我們有沒有依靠祂呢?.
除了保護以外.
牧羊人也會管教羊.
當羊不聽牧羊人的話.
偏行己路的時候.
牧羊人會用長桿拍打羊.
等牠返回正路上去.
因著大衛犯奸淫.
借刀殺人的緣故.
神非常不喜悅.
他除了透過拿丹責備大衛之外.
還讓管教臨到大委身上.
正如神讓拿丹說的.
刀劍必不離開你的家.
神要從大衛的家中興起和幻攻擊他.
兒子阿沙隆追殺父親.
是管教之一.
這個管教非常痛.
我深信大衛很深刻.
在兒子阿沙隆追殺父親之前.
大衛的家裡是做什麼的?.
兒女互相傷害.
大家如果有機會.
看回這段記載的時候.
你也會覺得很心痛.
神真的在他家中興起和幻.
去攻擊他.
神的管教很嚴厲.
但是神所愛的.
神必管教.
這是真的.
我想說這麼多年來.
我經歷過神大大小小不同的管教.
被神管教的時候.
我雖然感到不開心.
有時真的哭著說話.
很不開心.
但心裡明白是我自己做錯了.
我犯罪得罪神.

$^{681}$我記得很久之前.
我去短孫.
那次我很深刻.
我記得自己腦海裡.
當時有些不好的心思意念.
結果那次我經歷.
神的管教是我跌倒了.
我不知道為什麼.
追趕某些事情就跌倒.
跌得很痛.
那一刻就想.
神管教一定是我想些不好的.
真的非常不好.
跌得很痛.
痛到手有些傷.
腫.
那晚我記得女傳道很有愛心.
捉著我的手.
因為當時沒有冰敷.
捉著我的手去沖凍水.
痛啊痛痛.
一邊痛一邊思考.
都是自己不好.
自己不好.
這只是其中一個例子.
一邊跌一邊想.
在不同階段都經歷過神的管教.
但我感恩神管教我.
至少神覺得我還不是一個.
要被放棄的一個.
所以祂管教我.
在認罪過後.
就嘗試努力去悔改.
然後繼續走前面的路.
神管教過後.
我們的生命一定會有所成長.
弟兄姊妹們.
我們是否知道.
神一定會管教我們呢?.
是一定會的.

$^{721}$我們是否懂得分辨.
神的管教和一般普通的經歷呢?.
可能我們被管教也不察覺.
不知道神在管教.
也不知道自己錯了.
當我們經歷神的管教的時候.
我們有沒有反省自己的生命呢?.
在哪個位置得罪神?.
有沒有認罪?有沒有悔改?.
如果我們執意不悔改.
要知道後果很嚴重.
我今天說這番話.
是對坐在教會裡的弟兄姊妹說的.
不是對未信主的人說.
信耶穌的人.
也要常常在神面前警醒.
要小心自己的生命.
是否有得罪上帝的地方.
否則的話.
坐在教會裡的弟兄姊妹.
我們都要面對.
要面對神.
我不覺得我們在教會裡面.
面對神一定面對得起.
因為我們的神是聖潔的神.
公義的神.
祂要管教我們的時候.
亦都可以很嚴厲.
所以在這裡要彼此提醒.
第三件事.
好牧者眷顧羊.
除了滿足我們身心靈需要.
保護和管教羊.
好牧者也會讓羊過得性.
和豐盛的生活.
詩篇23篇5到6節.
轉換了另一個圖像.
另一個就是君王和戰士的圖像.
剛才已經說過了.
在舊約聖經裡.

$^{761}$君王的形象和牧者的形象是互通的.
君王令我們想起了牧者.
君王牧養百姓.
就好像牧者照顧群羊一樣.
詩篇23篇5到6節這樣說.
在我敵人面前.
你為我擺設筵席.
你用油告了我的頭.
使我的福杯滿溢.
我一生一世必有因為慈愛隨著我.
我且要住在耶和華的殿中直到永遠.
君王在敵人面前.
為戰士舉行慶功宴.
預祝他的勝利.
既然說是預祝.
就代表這場仗尚未打完.
但君王已經特意在敵人面前.
為他的戰士預祝勝利.
君王這樣做.
其實說真的.
不是因為這個戰士很厲害.
不是啊.
不是因為他的能力很高超.
他一定能夠打敗敵人.
是因為這個君王對他的戰士很有信心.
為什麼他對他的戰士很有信心.
因為君王自己會與這個戰士同在.
是君王的能力高超.
是君王令到這個戰士得勝.
所以他會這樣舉行這個慶功宴.
大衛打了很多場仗.
都贏了.
其實不是因為他很會打仗.
當然聖經描述.
他打仗的戰場上經常得勝.
最主要是因為他懂得依靠神.
每一次他打仗.
如果大家有機會看回那些經文.
描述大衛打仗是很豐富的.
每一次打法都不一樣.

$^{801}$神會給他不同的聰明智慧.
在過程中是上帝令他得勝.
上帝將仇敵交在他的手裡.
所以我在這裡深信.
在大衛打仗之前.
神已經知道大衛的勝利.
因為神自己是得勝的君王.
祂會令到大衛勝利.
今天神都會和我們擺這個慶功宴.
我們還沒打完仗.
但是上帝會宣告我們得勝.
不是因為我們很強.
是因為我們的神很強.
祂會令到我們勝利.
有祂在身邊的時候.
我們就沒有什麼懼怕.
無論今天你的戰爭是什麼.
面對你的健康挑戰.
經濟挑戰.
家庭關係破裂的挑戰.
情緒困擾各方面都好.
上帝說我們一定贏.
因為祂有這個能力.
絕對會使我們勝利.
我們吃這頓飯.
這餐的外賢是很歡喜快樂的.
所以弟兄姊妹打起精神來.
我們是那個勝利者.
因為背後有一位.
直到永遠都勝利的神.
除了得勝以外.
牧者也會叫他的羊.
過一個很豐盛的生活.
詩篇23篇5到6節這樣說.
使我的福杯滿日.
我一生一世必有因為慈愛隨著我.
我且要住在耶和華的殿中直到永遠.
君王給戰士的福氣是滿舍的.
他真的很愛他.
就正如牧者傾福於他的羊一樣.

$^{841}$最特別的詞語是隨著我.
這三個字不是隨著.
不是普通的跟.
而是形容那些野獸.
追著獵物的情況.
是怎樣追的?.
是慢慢跑的?.
是施施然的?.
是不會追到的.
是很急著的.
是很熱切的.
一直追我直到吃到為止.
他在形容上帝的福氣是怎樣?.
追著我們來.
這個很誇張.
是真的追著我們來.
神的恩典追著我們都是這樣的.
這個很寶貴.
詩人要住在神的殿中直到永遠.
住這個詞是形容軍隊回到營裡.
住在神的殿裡便有平安穩妥.
昔日打仗的時候是這樣的.
譬如今天有一場仗.
由早上一直打到黃昏時間.
如果你去到黃昏時間.
竟然可以回到你的營裡.
代表什麼?.
第一,你還沒死.
第二,你沒有被人捉.
那你真的贏了.
你還沒死,沒有被人捉.
那你便很平安了.
所以這裡說我要住在耶和華的殿中.
我平安了.
我回來了.
我勝過了這一切的困難.
我沒有被這一切的困難綑綁住.
我勝利了.
是上帝令我得勝.
我回到來繼續享受神同在的美好.

$^{881}$整幅圖畫是很美的.
在這裡我想起.
在我人生裡有很多福氣.
因為我預備這個講章的時候.
都在想.
神的恩典怎樣一生追著我去呢?.
是怎樣多發呢?.
哇,真的很奇妙.
記得最近有一晚.
跟家人分享爸爸一些往事.
因為前一段時間爸爸離世.
我們有時候會隔一陣子說.
他那時候這樣,我們那時候這樣.
很喜歡回顧.
回顧的時候.
我記得有一個位置.
很感動地說.
我說我們家裡最有福氣的位置.
全部都相信了耶穌.
這是最有福氣的.
生命有一個好牧者.
帶著我們走人生的路.
有什麼會缺乏呢?.
是得到最好的.
一家都有這樣的福氣.
有什麼需要擔心呢?.
我真的覺得我們很蒙福.
還沒有計算.
每一天的生活裡.
在大小的細節.
神都賜福.
有很多安排.
有時候有些很細微的位置.
我們會覺得.
啊,這麼奇妙.
上帝這樣預備了.
啊,神有這樣的工作.
真的很開心.
每一天都很多.
你數來數去.

$^{921}$是數之不盡.
我們最近.
我知道香港發放了消費券.
很多人馬上.
大家很有反應.
很開心.
是開心的.
都是神的恩典.
有消費券.
神的恩典絕對比消費券多.
何止幾千呢?.
這麼少嗎?.
多很多.
每一天你試著數算.
還有呢.
有人說.
最貴的是養育小朋友.
說要幾百萬.
去栽培一個小朋友生命.
你不要說小朋友.
我們一生都花很多錢.
有多少人曾經照顧過我們.
花了很多錢.
我們的生命只是這麼少錢嗎?.
神給我們這麼少嗎?.
肯定比這一切多.
多很多.
弟兄姊妹.
讓我們一同去數算神的恩典.
每天去數算.
看看.
感受一下.
什麼叫神的恩典追著我們去.
感受一下在神的家的平安.
最後一個總結.
今天我們說.
好牧者之歌.
唱一首歌.
讚美好牧者有幾樣東西.
第一樣東西.

$^{961}$好牧者和我們有很親密的個人關係.
很親密.
第二樣東西.
牧羊人.
好牧者和我們同在.
這個同在.
應該說他很眷顧我們的生命.
他眷顧我們的生命.
包括了他眷顧我們身心靈的需要.
他與我們同在.
帶領我們走對的路.
他亦都保護管教.
到最後最後.
好牧者很重要.
就是讓我們過一個.
得勝豐盛的生活.
很盼望這一首詩篇.
是真的繼續在生活裡.
成為我們的提醒.
我們的幫助.
我們一同禱告.
成為我們的好牧者.
你給我們的實在太豐富了.
我們無以為報.
就算今天我們做多少事.
有多少個奉獻.
都不足以去回應你的恩怜.
天父求你激勵我們.
帶著感恩的心.
經歷每日的生活.
在困難裡面.
又常常依靠你.
對你充滿信心.
亦都當上帝你提醒我們.
我們是有犯罪.
得罪你的地方的時候.
我們就很認罪.
很悔改.
讓我們很敏銳.
聖靈的提醒.

$^{1001}$離開罪惡.
主合給我們這一生.
繼續的去跟從你.
繼續的在人群當中.
去訴說你奇妙的作為.
願你在十週年的時間.
透過洪恩家.
你真的得到一切的榮耀.
還要讚.
在未來的日子.
你都繼續拖大.
整個洪恩家的發展.
成為洪水橋.
一盞明亮的燈台.
奉耶穌基督的名字.
祈禱.
阿們.
\newpage



\section{以斯拉記 6:13-22}
\label{sec:TujTTaedR4Y}
\textbf{2021.11.14《 神的恩典 》以斯拉記 6\_13-22 (和修版) 老耀雄牧師}
\newline
\newline
連結: \href{https://youtube.com/watch?v=TujTTaedR4Y}{\texttt{https://youtube.com/watch?v=TujTTaedR4Y}} ~~~~ 語音日期: 2021-11-15
\newline
\newline
\hyperref[sec:zUN5t6YF3yU]{\small{< < < PREV SERMON < < <}}
~
\hyperref[sec:index_chronic]{\small{[返順時目]}}
~
\hyperref[sec:index_scriptual]{\small{[返順卷目]}}
~
\hyperref[sec:sbhfqkg2NGI]{\small{> > > NEXT SERMON > > >}}
\newline
\newline
以斯拉記 6:13-22
\newline
\begin{longtable}{cl}
\hline
\hline
章節 & 經文 (和合本修訂版)\\
\hline
6:13 & \begin{tabularx}{0.7\textwidth}{X} 於是，河西的達乃總督和示他‧波斯乃，以及他們的同僚，急速遵行大流士王所頒的命令。 \end{tabularx} \\ \\ \relax
6:14 & \begin{tabularx}{0.7\textwidth}{X} 猶太人的長老因哈該先知和易多的孫子撒迦利亞的預言，就建造這殿，凡事順利。他們遵照以色列神的命令和波斯王居魯士、大流士、亞達薛西的諭旨，建造完畢。 \end{tabularx} \\ \\ \relax
6:15 & \begin{tabularx}{0.7\textwidth}{X} 大流士王第六年，亞達月初三，這殿完工了。 \end{tabularx} \\ \\ \relax
6:16 & \begin{tabularx}{0.7\textwidth}{X} 以色列人、祭司和利未人，以及其餘被擄歸回的人都歡歡喜喜地為神的這殿行奉獻禮。 \end{tabularx} \\ \\ \relax
6:17 & \begin{tabularx}{0.7\textwidth}{X} 他們為這神殿的奉獻禮獻了一百頭公牛、二百隻公綿羊、四百隻小綿羊，又照以色列支派的數目獻十二隻公山羊，作以色列眾人的贖罪祭。 \end{tabularx} \\ \\ \relax
6:18 & \begin{tabularx}{0.7\textwidth}{X} 他們派祭司按著班次，利未人也按著班次在耶路撒冷事奉神，正如摩西律法書上所寫的。 \end{tabularx} \\ \\ \relax
6:19 & \begin{tabularx}{0.7\textwidth}{X} 正月十四日，被擄歸回的人守逾越節。 \end{tabularx} \\ \\ \relax
6:20 & \begin{tabularx}{0.7\textwidth}{X} 祭司和利未人一同自潔，他們全都潔淨了。利未人為被擄歸回的眾人和他們的弟兄眾祭司，並為自己宰逾越節的羔羊。 \end{tabularx} \\ \\ \relax
6:21 & \begin{tabularx}{0.7\textwidth}{X} 從被擄之地歸回的以色列人，並所有歸附他們、除掉這地外邦人的污穢、尋求耶和華－以色列神的人，都吃這羔羊。 \end{tabularx} \\ \\ \relax
6:22 & \begin{tabularx}{0.7\textwidth}{X} 他們歡歡喜喜地守除酵節七日，因為耶和華使他們歡喜。耶和華又使亞述王的心轉向他們，堅固他們的手，去做神－以色列神殿的工。 \end{tabularx} \\ \\
[1ex]
\hline
\hline
\end{longtable}
$^{1}$洪恩堂弟兄姊妹平安.
平安.
今日是洪恩堂十週年的堂慶崇拜.
我在兩個月之前已經在想要講一篇甚麼訊息給大家聽.
我想起傳道書裡面一句經文.
它說拆毀有時 建造有時.
究竟洪恩家過往的歷史裡面.
有沒有被神拆毀過.
有沒有被神建造過.
如果是拆毀的 神拆毀了甚麼.
如果是建造的 神又建造了甚麼.
我覺得今日的經文很適合洪恩堂的堂慶崇拜的訊息.
我想用以色列記錄章十三至二十節.
和大家分享幾個在洪恩家事奉的神學反省.
我們一起去祈禱.
主的恩典數之不盡.
我們需要在不同的恩典裡面不斷地反思.
讓我們的心意能夠更生而變化.
測驗何謂神的善良 純存可喜悅的旨意.
每一個恩典都可以讓我們生命有所成長.
每一個恩典都可以令我們心意更生而變化.
最終就是要知道主的心意是甚麼.
好讓教會能夠按著你的心意而行.
讓我們可以成全初心 再次起步.
我們的祈禱是奉教主耶穌基督得勝名字祈求 阿們.
我們進入經文之前一起了解一下當時的一些背景.
我們知道自從所羅門死了之後.
以色列就分裂為南北兩國.
不過以色列在耶羅幫安領導底下脫離了大衛家的統治.
聯盟了北部十個支派 成立了獨立王國 建都在事件.
不過所有的王都行耶羅華眼中看為惡的事.
最後順藉著亞述攻陷北角 將百姓擄去亞述.
北角正式滅亡.
南國稱為猶大 屬大衛家的猶大支派所統治.
猶大的王好壞參半 其中出現過希西加的改革運動.
不過由於嗎哪西所犯的罪太惹神厭惡.
耶羅華不願意赦免他的罪 要將猶大趕出應許之地.
最後神也藉著巴比倫攻陷耶路撒冷.
將百姓擄去巴比倫 南角亦正式滅亡.
整個以色列民族開始經歷局破家亡之苦.

$^{41}$當時先知耶利米曾經預言過.
神應許他們七十年後可以重返耶路撒冷.
隨著亞述被巴比倫吞併 巴比倫又被波斯打敗.
以色列百姓不斷被不同的統治者遷移到不同的地方居住.
在異地漂泊.
經過七十年後 當時的波斯王古列大帝採取開放政策.
攀下玉子 讓被擄的民可以歸回原居地.
重建自己的家園和奠宇.
當時神就感動所羅巴伯帶領了四萬九千八百九十七人.
重返重歸耶路撒冷 重建聖殿.
亦是以色列民族第一次回歸.
當他們重修聖殿的根基之後.
就受到很多當地人和外邦人的攻擊.
他們進行賄賂誣告 甚至上奏到阿達薩西王.
於是阿達薩西王下令停工.
聖殿重建計劃受阻十五年.
直至大劉氏王第二年.
他重新確認古列的玉子.
於是批准以色列帶同五千個百姓第二次回歸耶路撒冷.
他返回耶路撒冷.
用了四年的時間將聖殿重建完成.
距離以色列人第一次回歸已經有二十一年.
而距離聖殿被毀已經超過九十年.
今日講到經文就是以.
以色列為聖殿獻聖殿禮的過程.
我們從經文裡面去反思一個信徒.
對神應有的態度.
第一 神藉先知對以色列人的提醒和勸勉.
在第十四節.
他說猶太人的長老恩哈該先知.
和亦多的孫子撒迦利亞的預言.
就建造寫殿凡事順利.
他們遵照以色列神的命令.
和波斯王居老氏大劉氏阿達薩西的預旨.
建造完畢.
經文提到兩個先知.
就是阿哈和撒迦利亞.
他們對聖殿重建的事.
向以色列的長老發出了預言.
我們先看看哈該的預言.

$^{81}$他在哈該書一章四字這樣說.
他說這個殿方糧.
你們自己還住天花板的房屋嗎?.
現在萬軍的耶和華如此說.
你們要審查自己的行為.
原來聖殿重建的工程已經停建了十五年了.
百姓擇以為常.
他們忙於修葺自己的房屋.
忘記了聖殿.
亦都忘記了耶和華的事.
忘記了他們為何要返回耶路撒冷.
所以先知提醒他們.
你們要審查一下你們自己的行為.
是否與回歸耶路撒冷的原意一致.
哈該在一章九至十節.
再一次警告百姓.
他說你們盼望多德.
看那所得的卻少.
你們修道加忠.
我就催去.
這是為甚麼呢?.
因為我的殿方糧.
你們各人卻只為自己的房屋奔走.
這是萬軍的耶和華說的.
所以因你們的緣故.
天不降甘露.
地不出土產.
耶和華警告百姓.
由於你們只顧著自己.
不是以神為大.
不是以神為先.
因此你們所做的一切.
都是徒然的.
如果你再這樣過你們的生活.
我也不會再賜福給你們.
這是哈該寫給他的百姓.
撒迦利亞又說了甚麼呢?.
在撒迦利亞書八章九至十二節.
這樣說.
這段經文也長一點.

$^{121}$萬軍的耶和華如此說.
你們的手要堅強.
這些日子你們已聽見先知的口了.
在萬軍的耶和華的殿的根基立定.
聖殿建造日子所說的話.
他說這些日子以前.
人得不著功架.
生畜也無人僱用.
且因敵人的緣故.
出入不得平安.
因我使人與人互相攻擊.
但如今.
我對這愚民必不像先前日子.
寫至萬軍的耶和華說的.
因為他們要平安殺種.
葡萄樹要結果子.
土地必有出產.
天也必降甘露.
我要使這愚民享受這一切.
甚麼叫這些日子之前?.
就是指聖殿未重建的時候.
百姓不會享有平安.
不過聖殿一旦能夠重建完成.
一切都不會像以前一樣.
重建之後.
葡萄樹要結果子.
土地必有出產.
天也必降甘露.
而要百姓享受這一切.
前一種是缺乏的日子.
後一種是豐足的日子.
以色列的百姓如何選擇呢?.
總結兩個先知的預言.
就是要這些回歸的百姓去想一想.
你們為甚麼要回到耶路撒冷?.
你們回到耶路撒冷的目的是甚麼?.
如果回到耶路撒冷只是收集你的房屋.
如果你回去只是因為生活安書.
那你真的不需要回來了.
只需要留在波斯.

$^{161}$或者巴比倫就可以了.
今日面對的缺乏.
探望將來的豐足.
百姓應該如何去選擇呢?.
如何才是一個正確的選擇呢?.
因此百姓聽從先知的教導.
開始上下一心.
重新去建造這個聖殿.
以四年的時間.
就將聖殿重新建造起來.
因此經文說.
「凡事順利」.
洪恩家的弟兄姊妹.
神的家如果荒涼.
你會如何回應?.
你會記掛自己的事?.
還是記掛神的家的事?.
今日洪恩家不需要在硬件上重建.
任何一個人.
今日進入洪恩堂.
都會覺得企企理理.
不過教會不是看外表.
還要看會眾的屬靈生命.
教會注重甚麼.
比教會的外表漂亮不漂亮.
重要得多.
如果問洪恩家現在最缺乏的是甚麼呢?.
我相信很多人都很清楚.
就是斷層.
教會缺乏兒童和青少年的肢體.
這個斷層現象如果不解決.
不處理.
教會就不可能承傳下去.
洪恩家的荒涼.
就是沒有了第二代.
亦沒有第三代,第四代.
因此如果洪恩家要重建的話.
就必定要重建兒童及青少年的事工.
在疫情之前.
即是去年年初.

$^{201}$兒童主日學仍然都未停止.
兒童事工都只有十個.
現在兒童事工未完全恢復.
半恢復當中.
回來的都只有兩個.
青少年事工就由頭到尾都沒有人.
完全斷層超過一段時間.
除了青少年的流失.
亦都沒有傳道童工負責.
上個月開始.
教會和新福事工合作.
每個周六都有大概二十八個小朋友.
由小一至中三的兒童.
回來教會參與一些興趣班.
教會開始和他們建立關係.
我們很期望這些興趣班完結之後.
他們會留下.
成為我們教會的一員.
因此教會現在需要的人手.
由傳道童工到義工都需要.
我們需要招聘青少年傳道.
但亦都要跟從區聯會的指引.
教會需要有一定的人數聚會.
大概這個指引就是去到一百四十人.
我們就可以申請去聘請第三個傳道童工.
所以請大家為此代討.
如果大家有看教會的財務報告.
我們就知道我們每一個月奉獻.
其實不足夠給這個新童工的薪酬.
亦都不足夠去發展青少年的事工.
我們的奉獻收入.
就只能夠維持現狀.
維持現在的活動.
因此我們不單需要金錢的奉獻.
去支持教會的青少年事工.
亦都需要更多的人參與教會的崇拜.
讓我們有足夠的數據去聘請童工.
我在這次的堂慶崇拜裡.
很想讓大家知道教會來年的方向.
而且亦都讓大家知道教會的優先次序.

$^{241}$就是我們很想全力去解決教會斷層的危機.
我重申.
我不是在這裡捉大家的數.
我是想邀請大家.
一起去參與洪恩家未來的重建工程.
第二個我在神學裡面的反思.
就是在神的殿中敬拜的喜悅.
十六節.
以色列人,祭司和利未人.
以及其餘被擄歸回的人.
都歡歡喜喜地為神的殿行奉獻禮.
經文以「歡歡喜喜」來形容他們的心情和態度.
這種心情和態度有什麼具體事件可以證明?.
他們為這個奉獻禮獻上.
一百隻公牛.
二百隻公綿羊.
四百隻小綿羊.
和十二隻公山羊.
全部都是他們獻給神的祭物.
要屠宰和清潔超過七百隻的動物.
想一想就知道需要多少的人手和時間.
他們可能要很早就預備.
甚至半夜就要開始工作.
這些人不是利未人.
也不是祭司.
只是普通的百姓.
反映到當時每一個歸回的人.
都很願意參與這場奉獻禮的盛況.
眾志聖城.
應該可以用來形容他們的狀況.
事奉人員又參與了什麼?.
十八節.
他們派祭司按著班次.
利未人也按著班次在耶路撒冷事奉神.
正如摩西律法書上所寫的.
祭司和利未人按著班次事奉神.
表示什麼呢?.
表示事奉人員是有偏濟的.
是輪流的當值的.
而這種事奉的模式.

$^{281}$是根據摩西律法的內容而執行的.
他們從被擄到現在.
已經有九十年沒有聖殿的敬拜.
其實在所羅巴伯建造聖殿的根基時.
已經有百姓大聲哭號.
亦有人大聲歡呼.
這是相當複雜的情緒表達.
有些回歸者可能從未在聖殿敬拜過.
他們可能在異邦出世.
聖殿敬拜的狀況.
只是從長輩口中聽聞.
亦有一些是真的曾經在被擄之前.
有參與過聖殿的敬拜.
他們從流乃爾密之地.
被擄到亞述,巴比倫,波斯.
不同的地方.
九十年來從未想過可以有聖殿的敬拜.
現在能夠再次在聖殿敬拜神.
心裡面是何等的歡樂.
這種感受是何等的恩典.
靜姐妹,我們有沒有相同的經歷?.
我記得在疫情期間.
教會停止了現場的崇拜.
只有網上直播.
亦有些人不懂得上網.
他們完全沒有崇拜的生活.
在疫情緩和之後.
當你能夠重返教會的心情.
還記不記得呢?.
我記得當日主席潤暢.
當日潤暢做主席的時候.
他那種雀躍的表現.
我今天還記得.
大部份的會長都感到很開心.
能夠回來教會崇拜很開心.
能夠見到那些會友很開心.
我們這段不能回到實體崇拜的經歷.
只不過是大半年.
但猶太人當時是經過90年的等待和仰望.
這種失而復得的感受.

$^{321}$應該不單用歡歡喜喜的心態來形容.
應該可以說是感恩.
再感恩的原意.
聖殿敬拜停了90年.
當時的聖殿禮儀還有沒有人記得呢?.
如何按照律法書上所寫的去執行呢?.
估計一定有敬虔的猶太人.
或是大祭司保存著聖殿禮儀的資料.
雖然聖殿被毀.
但仍心存盼望.
有望日後當聖殿恢復敬拜的時候.
祭司和利未人可以按著那些刺事奉神.
我們看到他們雖然失望.
但是沒有絕望.
因為他們相信神仍然看顧著他們.
我還記得上個月開始.
洪恩家恢復傳奉獻袋的環節.
對於這個以往每一個星期都會做的環節.
原來教會是沒有具體的指引保存下來.
因此就要詢問一些資深的會友.
究竟你們以前是怎樣做的?.
相關的流程是怎樣的?.
今天的奉獻環節是經過一些改革.
也符合奉獻的核心意義.
不過最重要的是.
今天已經有具體的文字檔案保存下來.
成為教會的資產.
昔日耶路撒冷的聖殿被拆毀.
百姓被擄,四處漂泊.
但是他們沒有放棄.
神都沒有放棄他們.
最終聖殿經過九十年.
又可以重建.
肅神的子民又可以在聖殿裡敬拜神,事奉神.
弟兄姊妹.
洪恩家曾經被神拆毀過.
無論在社會事件所造成的分裂.
也在團契分裂了.
疫情的限聚隔離,不能回崇拜.
慢慢地就不回崇拜了.

$^{361}$有些會友不懂得怎樣去看直播.
超過大半年沒有崇拜生活.
甚至沒有人事的變動.
對每一個肢體或多或少都有情感上的傷害.
有人失望過.
有人有疑惑.
有人離開.
也有人堅持留下.
我在這裡勸勉大家.
不要將焦點放在人身上.
因為人是很軟弱的.
而不要將焦點放在制度上.
因為地上沒有一個完美無瑕的制度.
我們要將焦點放在神身上.
因為唯有神才可以拆毀之後再建立.
也唯有神才可以為我們帶來盼望.
今天,我深信神開始重新建立我們洪恩家.
神賜了一個合一的團隊給教會.
神賜了一個很明確的意象給教會.
神賜了你們每一個給教會.
神愛洪恩家.
神按著他的心意去拆毀.
神也按著他的心意去建築.
或許我們曾經有失望過.
但千萬不要絕望.
因為神是那位次盼望的神.
今天是洪恩家十週年的唐慶崇拜.
我們雖然沒有宰殺到七百隻動物作為祭物.
但我們可以獻上我們每一個肢體為祭.
無論我們回到洪恩家有多久.
無論我們今天的屬靈生命是一個什麼狀態.
我們都可以將我們的新心靈呈獻給主.
作為今天十週年唐慶活動的禮物.
如果我們大家願意將我們自己的新心靈呈獻給主的話.
我們一起說一聲阿們好不好?.
阿們.
深信神必定閱立我們獻給他的禮物.
第三個.
在這段經文的神學反省.
我們就是說.

$^{401}$事奉者需要自潔.
在憲聖殿禮完畢之後.
他們就守逾越節.
在第二十節.
祭司和利未人一同自潔.
他們全都潔淨了.
利未人為被老歸回的眾人.
和他們的弟兄眾祭司.
並為自己在逾越節的羔羊.
祭司和利未人一同自潔.
他們全部都潔淨了.
事奉人員在事奉之先.
必須要潔淨自己.
免得被罪所污染.
一個不聖潔的事奉人員.
他們所做的.
都不會被神所閱立.
他們所獻的祭.
都是沒有果效的.
舊約隨罪的方法.
就是行贖罪祭.
以祭身代替自己去承擔罪責.
一個被罪污染了的贖罪祭.
神又怎會接納呢?.
舊約的信仰核心是隨罪.
新約的信仰核心.
亦都是要隨罪.
基督教的信仰核心.
就是如何去免除人類所犯的罪.
我所認識的神.
既慈愛又公義.
今天我們知道.
我們需要在主耶穌面前認罪悔改.
去滿足神的公義.
因為神必追討我們的罪.
然後主耶穌就施行祂的慈愛.
赦免我們的罪孽.
今天下午有一個感恩分享會.
對教會的事奉.
他們有些在洪恩家事奉了十年.

$^{441}$亦有些肢體很忠心地事奉.
教會不但對他們加以肯定.
不過我仍然在這裡提醒.
我們所有正在事奉的群體.
不單單是他們.
而是我自己都要被提醒.
作為一個事奉的人員.
我們需要好像昔日的祭司和利未人一樣.
我們要自潔.
我們要審查自己的罪.
要在主面前守潔心清.
否則我們所做的一切.
都不會被主所熱納.
而且更重要的.
還是有機會讓人跌倒失腳.
這個才是最重要.
不自潔最大的影響.
就是不能夠合一.
你有你的事奉.
我有我的山頭.
各自為政.
你不要理我的事.
我也不會給你意見.
教會一旦落入這個光景.
就會好像昔日南北兩國一樣.
當昔日南北兩國都有敬拜耶和華.
都有獻屬罪祭.
不過從君王到大祭司.
祭司利未人都過著犯罪的生活.
他們關注自己多於關心神.
沒有審查過自己.
過著一種自以為是的信仰生活.
主就直先知跟他們說過.
污穢的祭.
耶和華不會熱納.
以色列對神所做的一切.
無論獻多少隻牛,多少隻羊.
都是徒然的.
聽命聖於獻祭.
順從聖於供羊的知方.

$^{481}$神喜悅我們聽從祂的話.
按祂的心意而行.
而不是按自己的心意而行.
教會的主是主耶穌.
教會的王也是主耶穌.
教會的頭依然是主耶穌.
牧師只是主的僕人.
牧師只是一個器皿.
牧師只是一個.
仍然都是一個罪人.
需要神的挽回.
神的赦免.
牧師需要按著主的心意帶領教會.
牧師需要與會眾同心合意地.
履行神的使命.
牧師更加需要.
會眾用禱告托著和守望.
讓牧者不要偏離神的道.
牧師與會眾沒有高低之分.
因為大家都是基督的身體.
牧師需要認識你們.
你們也需要認識牧師.
今日是洪恩堂十週年的堂慶崇拜.
牧師願意在此向神立志.
與大家再次起步.
建立一間神所閱立的教會.
牧師願意以謙卑和忠心的心態.
與大家一起重建洪恩家.
今日的講題是「神的恩典」.
昔日的先知提醒我們.
關心神的家.
神的家興旺.
才是恩典.
唯有神的家豐盛.
信徒才能從豐盛的家得福.
要讓神的家有糧.
亦要讓神的家有工人.
唯有洪恩家盛載滿滿的恩典.
才可以將這福氣帶給每一個進入教會的人.
以色列的百姓提醒我們.

$^{521}$能夠回到聖殿敬拜.
是莫大的恩典.
能夠有敬拜的生活.
不是必然的.
能夠在聖殿裡與神相遇.
領受神的話.
更加是恩典之中的恩典.
每一個基督的跟隨者.
都要珍惜每一次的敬拜.
因為每一次的敬拜.
都能夠增強我們與神的關係.
建立我們的信心.
讓我們的屬靈生命更加成熟.
甚至洪恩家每一周的敬拜.
都能夠洗刷我們每一個的心靈.
讓我們更加貼近主.
祭司和利未人提醒我們.
能夠事奉是一個恩典.
事奉者要有聖潔的生命.
否則一切的事奉都是徒然.
事奉者需要專注於神的話.
亦要與神有美好的關係.
只要每一個事奉者都能夠彰顯神的屬性.
就能夠讓人在我們生命裡面見到主耶穌.
讓主得到榮耀.
最後,如果說神拆毀了洪恩家甚麼呢?.
我會說,神拆毀了洪恩家的紛爭和是非.
所有的誰是誰非都應該結束.
放下我們的不滿,憤怒.
以及對人,對教會的指責.
在主的恩典底下,我們重新起步.
教會不應該停留在過往的歷史裡面.
教會可以檢視過往的得失.
但不是被過往的歷史所纏繞.
而令到教會停滯不前.
如果說神建立了甚麼呢?.
我會說,神建立了洪恩家與母堂的關係.
以往洪恩家與母堂都有很多誤解.
洪恩家以為母堂對洪恩諸多流難.
母堂又以為洪恩家自把自為.

$^{561}$我可以肯定地說,母堂是全力支援我們.
不過我們也要爭氣.
我們也要做一些成績給別人看.
而且在今天的經濟狀況,逆境底下.
教會也需要學習量入為出的理財方案.
我們不能夠每樣東西都攤大手擺,問人拿.
我們也需要在各方面有進步的成果.
甚麼叫進步的成果?.
我們有多少人參加主一學?.
我們有多少人學武神的話語?.
我們有多少人參與祈禱會?.
這些都是一個很具體的指標.
弟兄姊妹,今天是洪恩家一個很重要的里程碑.
藉著洪恩家大學十週年.
我們能否在主面前,我們一起去承諾.
我們上下一心,同心合意地去重建洪恩家.
我們一起獻上禱告.
親愛的恩主,洪恩家在你的看顧之下.
恩典滿門,數之不盡.
當我們去檢視過去的時候.
我們要審查我們的不足,軟弱.
承認我們的錯,求主你赦免我們.
當我們展望將來的時候.
我們要謙卑在你面前,求主你帶領我們.
讓我們可以與你同工同行.
不走在你前面,只走在你後面.
求主你堅固我每一個弟兄姊妹的心.
讓我們每一個都緊緊地跟從你.
又讓我們有一個合一的心.
好叫我們能夠按著你的心意而行.
主啊,過往曾經受過傷害.
心靈有創傷的.
求主你親自去醫治憐憫.
因為你是醫治的主,也是最好的醫生.
求主讓我們每一個都在你的重建計劃裡面有份.
也讓我們每一個都能夠與你同行同工.
甚至我們每一個都能夠以你為尊,以你為大.
以你的旨意成為我們行事為人的最高的優先次序.
求主你閱立我們同心合意的禱告.
是奉主耶穌基督,得勝名自祈求.

$^{601}$阿們.
願主祝福大家.
\newpage



\section{路加福音 10:38-42}
\label{sec:sbhfqkg2NGI}
\textbf{2021.11.21 《接待上主之道》 路加福音 10\_38-42 (和修版) 趙崇明博士}
\newline
\newline
連結: \href{https://youtube.com/watch?v=sbhfqkg2NGI}{\texttt{https://youtube.com/watch?v=sbhfqkg2NGI}} ~~~~ 語音日期: 2021-11-25
\newline
\newline
\hyperref[sec:TujTTaedR4Y]{\small{< < < PREV SERMON < < <}}
~
\hyperref[sec:index_chronic]{\small{[返順時目]}}
~
\hyperref[sec:index_scriptual]{\small{[返順卷目]}}
~
\hyperref[sec:S_q0D_6qEuY]{\small{> > > NEXT SERMON > > >}}
\newline
\newline
路加福音 10:38-42
\newline
\begin{longtable}{cl}
\hline
\hline
章節 & 經文 (和合本修訂版)\\
\hline
10:38 & \begin{tabularx}{0.7\textwidth}{X} 他們繼續前行，耶穌進了一個村莊。有一個女人，名叫馬大，接他到自己家裡。 \end{tabularx} \\ \\ \relax
10:39 & \begin{tabularx}{0.7\textwidth}{X} 她有一個妹妹，名叫馬利亞，在主的腳前坐著聽他的道。 \end{tabularx} \\ \\ \relax
10:40 & \begin{tabularx}{0.7\textwidth}{X} 馬大伺候的事多，心裡忙亂，進前來，說：「主啊，我的妹妹留下我一個人伺候，你不在意嗎？請吩咐她來幫助我。」 \end{tabularx} \\ \\ \relax
10:41 & \begin{tabularx}{0.7\textwidth}{X} 主回答說：「馬大，馬大，你為許多的事操心煩惱， \end{tabularx} \\ \\ \relax
10:42 & \begin{tabularx}{0.7\textwidth}{X} 但是不可少的只有一件。馬利亞已經選擇了那上好的福分，是沒有人能從她奪去的。」 \end{tabularx} \\ \\
[1ex]
\hline
\hline
\end{longtable}
$^{1}$各位弟兄姊妹平安.
想不到李姊妹會幫我做廣告.
我今天選了一個題目.
是接待上主之道.
我盼望你聽完我講到之後.
你明白我這個題目其實是.
我改這個題目是語帶相關的.
我希望你聽完之後你會明白.
怎樣語帶相關.
這段經文我相信除非你是剛剛回教會.
否則你一定很熟悉這段經文.
不過我的印象不知道有沒有錯.
我總覺得不少基督徒可能會誤解這段經文.
最容易的誤解就是我們很多時候以為耶穌.
只是肯定瑪利亞所做的事.
而否定馬大所做的事.
我覺得這是一個很大的誤解.
很多時候我們怎樣去比較.
人很喜歡和人比較.
馬大代表的只是為屬世的事操心.
關心的只是物質的世界.
而瑪利亞關心的是屬靈的事.
嚮往的是屬靈的世界.
或者這麼說.
馬大代表的是一種外向型的性格.
只會死做爛做.
做事就厲害了.
工作狂也不止.
瑪利亞代表的是一個很內斂很靜的人.
喜歡靜修喜歡靈修.
喜歡物觀喜歡夢想.
又或者瑪大代表一個英文是aggressive.
aggressive是自己能譯的.
他有一些正面但有一些負面的意思.
正面的意思就是一個人很主動很進取.
但又不是這麼正面.
aggressive就是過度主動過度進取.
而帶來其他人的壓迫感.
你看聖經就這樣說.
在廚房做事做到一頭煙.

$^{41}$看到妹妹還在坐著.
主人快點為何不叫她進來.
起碼給妹妹有些壓迫感.
瑪利亞就不是.
代表一個很被動沒有他那麼主動.
一個被動一個很順服很聽從的人.
甚至我聽聞有人說.
認為瑪大代表靠行為得救.
瑪利亞代表恩順清義.
說大了你都會笑了.
說大了我真的聽說是有人這樣說的.
於是結論就是.
是一種很二元對立.
一個這樣的思維影響.
不要說笑.
我覺得我們還教很多基督徒.
是被這種思維影響得很深.
於是耶穌不欣賞瑪大.
否定瑪大所做的事.
耶穌只肯定和稱讚瑪利亞.
背後就是屬靈屬世的餘分.
屬世的不好.
我最初上教會時有個感覺.
屬世的是邪惡的.
只是追求屬靈的東西.
你不要小看這種.
我們經常被這種思維影響.
我們解經時看很多東西.
要了解一段經文.
譬如《福音書》.
一定要了解經文當時的故事.
所講的背後.
是一個甚麼社會處境.
一個甚麼文化處境.
大家知道當時羅馬和猶太社會.
明顯仍然是一個男權社會.
好像我們中國傳統社會.
這一代好一點.
我祖母那一代仍然是男人大哂.
男尊女卑.

$^{81}$男主外女主內.
這些是我們華人很熟悉的東西.
當時的社會也是這樣.
猶太人的社會也是這樣.
羅馬的社會也是這樣.
當時社會對兩性的角色.
是有一些特定的規範.
有一些特定的期望和要求.
社會期望甚麼.
就像我剛才所說的.
像中國人一樣.
女人在家中扮演重要的角色.
一切家務由女人打點.
洗衣服煮飯.
一腳踢.
有客人到訪的時候.
當然是婦女去接待.
是份內的事.
所以如果你知道當時.
也是一個這樣的背景的時候.
馬大對耶穌很辛勞的款待.
是沒有問題的.
符合當時社會傳統對女人的要求.
你不是這樣做你才會死.
所以馬大正在做你所當年的事.
甚至可以說是當時社會中的女性典範.
你再看看聖經.
剛才讀完經文.
其實耶穌從頭到尾沒有否定馬大所做的事.
沒有.
你找一句告訴我.
哪一句.
否定馬大所做的事.
沒有.
事實上如果熟悉科林書.
你知道我們的神很有趣.
沒有一個宗教.
就算那個宗教說有神.
沒有一個宗教像我們這樣的.
我們的神是什麼.

$^{121}$就是神成為人.
道成肉身.
成為人.
道成肉身.
有肉身的身體來到這個世界.
所以我不知道你知不知道.
我們教會經常說.
耶穌是救世主.
這個大家很熟悉.
你知不知道耶穌也是創造主.
知道就好了.
我怕你真的不知道.
《猶人徽音》很多書信.
歌羅西書.
都是耶穌是創造主.
耶穌有份創造這個物質的世界.
然後他來到這個有份創造的物質世界.
他用肉身來當中成為肉身世界的一部份.
以馬來尼.
即將到來.
聖誕節.
以馬來尼是什麼意思.
神與我們同在.
神與我們同在很特別.
在我們教會很特別的就是.
他用肉身與我們同在.
你知道肉身的同在很重要嗎.
現在有實體崇拜很重要.
以前疫情比較嚴重的時候.
大家在網上崇拜.
肉身崇拜很重要.
肉身一起.
那個關係很不同.
很簡單.
譬如你知道某一個弟兄姊妹.
她很不開心.
很難過.
最近她很傷心.
如果大家肉身同在的時候.
你可以過去拍她的肩膀.

$^{161}$同性的可以抱她.
你在網上怎麼抱.
不能抱的.
但抱的擁抱勝過千言萬語的安慰.
所以我們說.
我們不可能人與人之間的關係.
我們不可能不是一種肉身的關係.
肉身一起同在的關係.
耶穌就是這樣.
我們的神就是這樣.
祂來到這個世界上.
還有.
祂來到這個世界上又如何呢.
耶穌很有趣.
聖經形容祂是.
貪吃好酒的人.
你記得嗎.
很喜歡吃東西.
祂經常與人吃飯.
不過很有趣的.
聖經記載祂多數與什麼人吃飯.
與罪人,妓女,稅吏.
當時社會上一些不太喜歡的人.
不太歡迎的人.
或者是一些弱勢社群.
耶穌很喜歡與他們吃飯.
貪吃好酒.
真的.
祂臨死都要與門徒一起吃飯.
吃完飯我才走.
才可以去.
在最後的完結.
最後晚餐.
入聖餐.
聖餐都是與吃東西有關.
其實整部聖經.
66卷聖經.
我幾年前在神學院的一個完新課程.
我開了一課叫.
聖經有講飲講食.

$^{201}$是嗎.
是呀.
聖經經常講飲講食.
這位主如果是這樣的時候.
他沒有理由反對.
馬大在廚房煮東西給他.
馬大煮飯.
招待耶穌是一件美事.
耶穌沒有理由反對.
所以我估計在耶穌心目中.
馬大和瑪利亞不是兩個對立的角色.
我們常常說.
屬靈.
比屬世好.
屬靈和屬世.
心靈和物質.
靜態和動態.
沉思和行動.
靜修和工作.
主動和信服.
行為與信心.
這些不是對立的.
不是有你沒我.
而是兩者共存.
有時我真的很害怕.
我們基督徒.
我們常常說.
講耶穌講耶穌.
全副音就是講耶穌講耶穌.
只有一個「講」字.
講就是天下無敵.
做就是甚麼.
無能為力.
軟弱無力.
總之是不會做的.
亞國書說過.
有信心沒有行為是甚麼.
死的.
只有一個「講」字.
講到多大都可以.

$^{241}$行動很重要.
行動關心一個人.
關心他肉身的需要.
這個很重要.
你不是只跟你說耶穌.
你之後靈魂得救上天堂.
只有靈魂.
我們只說靈魂得救.
肉身不重要嗎.
物質不重要嗎.
對耶穌來說同樣重要.
因為祂創造了這個世界.
這個物質的世界是祂創造的.
祂看著是好的.
我很相信.
我絕對相信.
耶穌很期望.
你跟我.
既是瑪利亞.
又做馬大.
我很肯定.
不過.
你剛才讀了一段經文.
畢竟.
你讀經文的時候.
又好像.
讀到耶穌對馬大所做的事.
他沒有否定他煮飯.
但真的有些批評.
於是你讀上去的時候.
你會覺得.
好像不及瑪利亞.
因為耶穌說過最重要的.
他選擇了一個上好的福分.
沒有人能夠奪取.
這句很關鍵.
如果我們說.
做事,行動,煮飯.
沒有問題.
問題在哪裡呢?.

$^{281}$耶穌提醒他些什麼呢?.
耶穌提醒馬大些什麼.
以至耶穌今天提醒你和我些什麼呢?.
這是我們要想的.
究竟是些什麼呢?.
我相信.
關鍵在41節.
41節.
他說:馬大.
馬大.
你為.
留意四個字.
許多的事.
馬大很多事要做.
很忙.
許多的事.
可能是問題.
但不一定是問題.
你很忙.
可以是問題.
但忙不一定又是問題.
你說什麼呢?朱先生.
因為還有下面那四個字.
操心煩惱.
這是和修本的翻譯.
如果用和修本的翻譯.
就是思慮煩擾.
其實是這四個字.
如果我們很忙的時候.
令我們操心煩惱.
或者思慮煩擾.
就有問題.
如果你很忙.
很多事要做.
但是沒有思慮煩擾.
未必有問題.
你明白我說什麼嗎?.
重點就是.
操心煩惱.
煩惱.

$^{321}$什麼意思呢?.
大家一定很熟悉.
在你的日常生活中.
你的經驗中.
很多事要做.
千百樣事在你身邊圍繞著你.
很容易將你拉扯.
是嗎?.
令你分心.
是嗎?.
香港人習慣忙.
忙得很厲害.
有時日間工作.
現在很難分開工作.
下班回家.
我們不工作.
很多時不是這樣.
我不知道你有沒有忙到這樣.
你白天很多事要做.
回家後.
跟家人一起吃飯.
一家人很開心.
太太也好.
就算傭人煮飯也好.
煮了很美味的食物.
一家人應該品嘗那些.
很美味的食物.
然後大家聊天.
但可能你很忙.
又要回覆WhatsApp.
又要回覆這樣和那樣.
一邊吃飯.
你繼續忙著工作.
沒有好好享受飯桌.
沒有跟家人好好交流.
很多事令我們分心.
忙的時候已經是這樣.
我看下去.
不知道有沒有猜錯.
不知道當中會不會有些人退休.

$^{361}$退休還很忙.
不是的.
有些人退而不休.
很難停下來.
又或者這樣.
可能你很難看.
很多事要做.
仍然很忙.
因為我們習慣了.
香港的生活太忙.
習慣了.
令我們為很多事操心煩惱.
思慮煩擾.
沒辦法.
正常.
真的挺正常.
香港活了這麼久.
不正常也不行.
你不是在農村社會.
不過最麻煩.
我覺得.
我們的社會仍然是這樣.
不過不是最近.
現在還多了這個.
誰沒有帶這些東西在身.
請舉手.
真的少.
你竟然可以不用帶這些東西.
你不用的可以.
在上面.
忘記帶在身上.
不可能不帶在身上.
你不帶錢.
不帶現金出街.
也沒所謂.
你不可以.
這個真的越來越厲害.
這是小型電腦.
什麼都靠它.
是不是.

$^{401}$我不用說這件事有什麼好處.
一定數到很多好處給我們聽.
方便.
它最厲害的.
最厲害的是什麼.
已經打破時空的隔膜.
是不是.
幸好這個世界有這個科技.
不然整個疫情來的時候.
真的怎樣崇拜.
真是太好了.
可以在網上崇拜.
是不是.
多好.
就算現在我不知道你們有沒有移民.
在很遠很遠的地方.
其實都可以崇拜.
一起崇拜.
這是好處.
不過.
這個東西.
其實我們已經活在一個.
電子網絡的世界裡面.
其中一個不好處就是上癮了.
我不叫你舉手.
都不用怎麼想.
大大小小都上癮了.
上癮了.
手不離肌.
另外一個問題.
你知不知道.
我們活在一個這樣的文化裡面.
其實不知不覺地展現了一件事.
我們活在一個英文叫Multitasking.
就是同一時間可以做很多事.
在一個這樣的生活模式裡面.
這樣就好了.
這樣就省時間了.
就是它厲害的東西.
多好.

$^{441}$同一時間可以做很多事.
我希望我沒有說中你.
那時候.
完全在Zoom的時候.
在家裡網上崇拜的時候.
其實你們應該都很晚了.
十一點半.
我很少去教會十一點半.
第一堂.
都很晚的.
不過你習慣了.
我女兒也是.
特別年輕人.
十一點半.
剛剛起床.
可能你遲起床.
崇拜怎樣.
做個早餐.
一邊吃早餐一邊崇拜.
我都說我希望我沒有說中你.
不僅如此.
因為太忙了.
過去一星期太忙了.
忙到我根本沒有收拾房子.
網上崇拜.
嘩.
然後看總.
起碼說完就說到.
起碼說半個鐘至三十五分鐘.
夠時間收拾房子.
開大電視.
聽而已.
廣東話聽的不用看.
那個講完的趙崇明又不是帥哥.
所以不用看他的樣子.
聽到就行了.
於是拿吸塵機回來.
一邊聽到一邊吸塵.
你手快的時候.
整個房子可以搞定.

$^{481}$嘩.
同一時間又聽到.
又做完.
我聽到的.
我每個字都聽到的.
你猜上帝看到你.
你猜上帝會不開心嗎.
當然不開心了.
不過更重要的一點.
慢慢塑造我們不能專注.
弟弟姐妹.
分心是現代人的一個問題.
本身工作忙碌已經令我們分心.
再加上這些.
我們已經上了癮.
有時候你回來教會.
你現在只是聽到.
我怎知道你怎樣.
下面這樣.
一間餐廳.
去哪裡吃飯.
敬拜都不能專注.
我希望我沒有說中你.
這個也不是最大問題.
也不是沒有問題.
不過我接著想說另一個.
從一個屬靈的層面.
其實耶穌透過對馬大的提醒.
很想提醒你和我.
今天特別處於文化的你和我.
要小心.
從屬靈的層面來說.
我覺得專注是什麼.
專注不單單是.
你能不能夠在同一時間.
只做一件事.
集中精神做一件事這麼簡單.
從屬靈層面來說.
專注是你能不能夠.
經常保持生命只有一個焦點.

$^{521}$你能不能夠經常保持.
你生命只有一個焦點.
你知道那個焦點是誰.
是上帝.
你知道的.
你和我知道的.
這裡知道的.
不過做不到.
分心.
我們上帝是不上心.
我們很多東西已經佔據了我們生命.
祂成為我們的焦點.
哪樣東西上心.
那個就是你生命的焦點.
很容易查的.
在拍拖的時候.
當然是男朋友或女朋友.
是我最上心的.
他一言一語.
真的很上心.
經常想念他.
做夢也想念他.
結婚之後當然是配偶.
生了孩子.
但跟著子女.
做爺爺奶奶.
當然是孫子.
外婆.
當然是孫子.
外孫.
那些是最上心的.
你經常想念他們.
最上心.
可能是你的事業.
很多東西令我們很上心.
現在我女兒最上心的是甚麼.
其中一樣.
MIRROR.
她可以做很多事.
她是姜糖.

$^{561}$姜濤出現.
好像整個人很興奮.
好像生命.
淑齡粉輕.
她有偶像.
不只是年輕人有.
你和我有.
你和我一定有偶像.
以致我們的生命.
不是一個焦點.
是上帝.
不是其他那些.
很多.
如果你有很多偶像就慘了.
思慮煩擾.
不能專注.
其實耶穌已經提醒我們.
耶穌說的那句話是金句.
你們不能又事奉上帝又事奉晚門.
如與熊掌不可兼得.
我們很貪心.
又聽道又打人家頭細毛.
我們很擔心.
不可兼得.
耶穌也不是叫我們選擇.
不能選擇.
耶穌那句話叫我們不要將生命的焦點放在財利上.
然後耶穌說完那句.
你知道接著說甚麼嗎.
生命為民天憂慮的問題.
憂慮甚麼.
吃甚麼穿甚麼.
你的工作是甚麼.
搵食.
就是搵食.
財利是甚麼.
就是搵食.
當然我們可能多些搵食的東西.
除了搵食之外.
在那裡說我自己得到自我的滿足感.

$^{601}$覺得自己很棒.
回到原來我一個這麼有能力的人.
在那裡找回我自己的.
我多有成就.
全部都是.
我.
弟兄姊妹.
專注在甚麼.
單單以上帝為生命的焦點.
如耶穌所說.
我們單單事奉上帝.
那你說那不用工作嗎.
不用打你家頭細毛嗎.
不是 當然不是.
而是我們做一切世上的事的時候.
上帝.
以上帝為焦點.
然後做那些事.
但很多時候你和我不是.
很坦白分享.
我在神學院事奉.
全世界事奉.
每天都教這些東西.
每天都和聖經離不開.
應該很專注.
有時未必.
我承認有時未必.
因為.
很多時候我們.
人性.
我們有時也會討人喜悅.
我們所事奉的是討人喜悅.
多於討神喜悅.
有時我會這樣.
當我討人喜悅的時候.
我所做的事.
表面上都是在事奉.
但其實.
最上心的是那些人.
最想討誰的喜悅.

$^{641}$領錢會.
這個很重要.
另外現在想說說.
我經常去到哪裡.
我一定去教會.
拿一本實體聖經.
你們現在沒有了.
現在沒有手提電話.
你拿手提電話出來.
我相信你不會到處找.
打開經文.
因為實體聖經是甚麼.
實體聖經很重要.
我們對一段經文的了解是上下文.
我們不能離開上下文.
了解一段經文的意思.
所以實體聖經是最容易了解上下文.
很有助於我們解經.
你沒有實體聖經.
如果你不介意.
你拿手提電話出來.
如果你有下載經文.
就去到路加芬第十章.
我和你說說.
這段聖經.
這段聖經記載在路加芬第十章最後.
第十章最後的部分.
38至42節第十章最後的部分.
你不可以只抽這段經文來理解.
你想更深入理解這段經文的意思.
你看前一點.
原來這段經文是源於25節.
源於25節.
第十章25節說甚麼.
你都很熟悉.
我告訴你.
這段經文記載.
有一個律法師.
他來試探耶穌.
甚麼是律法師.

$^{681}$他是律法的專家.
對舊約的律法非常純熟.
然後他來試探耶穌.
他問.
夫子我做甚麼可以承受永生.
耶穌就跟他說.
律法上寫的甚麼.
他是律法師.
於是就馬上回答.
你要盡心盡性盡力.
盡意愛主你的神.
又要愛人死而同我自己.
100分.
耶穌說你答對.
當然是對.
他是律法專家.
熟讀律法.
律法的總綱.
律法的精神.
律法的精義.
他講到了.
然後耶穌說.
你回答的是.
你知識有100分.
不過你怎樣行.
所以剛才我說.
你知道是沒有用的.
很多事你都知道.
你有沒有可以行呢.
你要行才可以永生.
然後.
你記得耶穌說了一個故事.
這個故事是虛構的.
就是好撒瑪利亞人的比喻.
行嗎.
好撒瑪利亞人的比喻是甚麼.
你不用說.
這個故事你很熟悉的.
就是有一個人受傷了.
兩個宗教人士走過.

$^{721}$不理會他.
見死不救.
然後有個撒瑪利亞人見到.
做足100分有得.
這個故事就是回應剛才.
那個律法的總綱的下半部分.
盡心盡性盡義盡力.
愛鄰舍如同愛自己.
行嗎.
這個故事就是回應那部分.
去行動.
很重要.
好了.
完結的故事.
就是我們今天所讀的經文.
就是馬大瑪利亞這段故事.
所以這個故事.
講的是甚麼.
是回應經文的律法總綱的上部分.
盡心盡性盡義.
愛主你的神.
剛才說甚麼叫對神專注.
就是「盡」字.
弟兄姊妹.
你有沒有盡心盡性盡力.
盡義愛主你的神.
抑或在你和我的生命裡.
只是撥10\%給神.
90\%就是我剛才說的其他東西.
你就很難盡.
關鍵就是「盡」字.
無所分心.
搶奪了我們對神那種.
愛的盡.
愛到盡.
弟兄姊妹.
耶穌稱讚瑪利亞.
瑪利亞不單止在聽道.
她在享受和主耶穌同在.
那種團契交流.

$^{761}$那種美好的共融的關係.
她在聽耶穌說話.
這就是耶穌說的.
上好的福分.
她很專心聽耶穌說話.
所以你就明白.
路加福音十一章一開播說甚麼.
「祈禱」.
明白嗎.
聖經的作者編輯聖經.
是有他的道理.
他是藉著這些一段一段的訊息.
所以我說看聖經容易看到.
是用PowerPoint的一兩句.
這些都是.
弟兄姊妹.
你和我.
有沒有這種上好的福分.
你看不看重這種上好的福分.
其實.
這一段經文.
剛才那一大段經文.
由二十五節到四十二節.
說的主題就是「接待」.
今天我的題目是.
「接待上主之道」.
接待.
剛才耶穌說的虛構故事.
是在說好撒馬林.
肉身的接待.
那個傷者.
接待的英文是甚麼.
Hospitality.
和Hospital.
醫院的字很接近.
醫院.
我在說香港.
我們香港的醫療也不差.
我在說公立醫院.
不是說私家醫院.

$^{801}$公立醫院是甚麼.
見到一個人有需要.
就要入急診實.
不論他是甚麼社會階級.
就算是乞丐.
骯髒邋遢.
他有受傷的時候.
有需要的時候.
就一定入來.
也不會說你先給錢.
給錢我就給.
私家醫院可能會.
先進來.
慢慢才付錢.
甚至你窮的時候.
就申請不用.
這個就是醫院.
醫院就是一個接待.
無條件地接待.
見到有人肉身的需要.
就接待.
這個就是好撒馬林的故事.
盡地用愛淪陷.
肉身的接待.
去到我們今天說的.
經文馬大.
瑪利亞.
瑪利亞其實是.
耶穌藉著.
祂做這件事.
來說.
為接待加個意思.
耶穌沒有否定物質的接待.
不過.
當瑪利亞坐在主的腳前.
聽主講道的時候.
瑪利亞很安靜地聽到.
表面上是沒有做事.
什麼都沒有做.
瑪利亞做到一頭煙.

$^{841}$但瑪利亞表面上沒有做事.
但其實.
如果用我們中國道家思想.
好像很無謂.
但其實大有所謂.
她正正在接待主.
她接待主就是聆聽主的話.
她接待主.
她就是用聆聽主的話.
這個行動接待主.
原來接待這件事.
接待主也好.
接待人也好.
不否認物質的接待.
那些很重要.
但同樣.
如果你要在肉身裡.
懂得怎樣接待有需要的人.
先後次序不可以倒轉.
就等於律法還說.
你要先要盡心盡心.
愛神.
然後愛靈.
然後愛自己.
這個先後次序不可以倒轉.
所以當耶穌.
用這個故事就說.
瑪利亞是很專注接待.
聆聽耶穌的道.
當我們常常聆聽耶穌的道.
知道耶穌的教導的時候.
我們在那裡不單止是知識上.
而是我們經驗到剛才說.
那種和主席.
就是我們平時聆聽.
和一個這麼親密.
這麼愛的關係.
以至我們領受到耶穌的愛.
耶穌的道主要是他的愛.
而且這個道.

$^{881}$生生當日在瑪利亞面前.
成為肉身的道.
甚至將我們相信十字架.
這個愛的上帝.
當你感受到愛的時候.
你才有力量.
有資源來愛.
接待其他人.
這個先後次序不可以倒轉.
否則外面的人.
都在做很多事情.
和我們基督徒做的事情.
分別在哪裡呢.
就是在這裡.
還有你想想.
你想像一下.
耶穌是一個客人.
來到瑪大利亞的家.
假如瑪大利亞出生.
為什麼我女兒.
不是我女兒.
為什麼不進來幫我.
如果瑪利亞真的進來幫的時候.
只剩下耶穌自己一個.
在客廳裡.
孤伶伶的獨自外坐.
難道這樣是好的待客之道嗎.
真正好的款待客人的道.
就是坐在這裡.
和他聊天.
如果不認識他.
聽聽他的故事.
其實這個就是好的待客之道.
我的題目與大相關就在這裡.
接待上主這裡.
什麼叫做好的款待客人之道.
最好的待客之道是什麼.
就是接待.
耶穌就是上帝的道.
所以最好的待客之道.

$^{921}$如果那個客人是耶穌.
就是接待耶穌本身.
各位姐妹.
最近我們香港社會.
經歷了一個很大的移民潮.
我不知道你們教會有沒有移民.
很多教會很多人移民.
其實移民不是開玩笑的.
非不得已.
一般來說非不得已都不想走.
我有一個同學認識了很多年.
他最近去了英國.
他跟我們說他被迫移民.
去到那裡你會有很多適應.
一個新的環境.
新的文化.
很多的適應.
有些社會未必真的很歡迎中國人.
可能有些歧視.
諸如此類.
很容易被人冷落.
甚至乎排擠.
道成肉身的故事.
其實都是一個被人冷落的故事.
《楊富安第一章》十節十一節這樣記載.
這個世界本來是耶穌做的.
這樣記載.
祂在世界.
世界也是藉著祂做的.
世界卻不認識祂.
祂到自己的地方來.
自己的人都不接待祂.
剛才那個很撒瑪利亞人的故事.
大家記住.
那個虛構的故事.
你看盧家鋒揭前幾章.
真正的撒瑪利亞人是不接待耶穌的.
那個記載了.
現實上撒瑪利亞人是不接待耶穌的.
那個故事.

$^{961}$不知道是不是耶穌說的故事.
是專門給撒瑪利亞人.
留意《楊富安第一章》十節十一節的經文.
將不認識和不接待.
能上關係.
事實上是的.
農村社會.
家門是門上開的.
你走過去看看.
沒問題的.
農村社會.
但現代人的社會是怎樣的.
門上關的.
而且緊緊關上.
在洞裡有人按鐘.
在洞裡撞.
不認識的.
坐著也傻.
你會不會不認識的.
我給你一條鑰匙進來.
住幾天.
不會.
一定不會.
正常一定不會.
我也不會.
因為不認識.
所以不接待是正常的.
不過你想想.
因為不認識.
所以不接待是正常的.
但因為你不認識.
所以才不接待他.
你永遠不接待.
你怎樣認識他.
即是永遠都不會接待他.
你明白我的道理嗎.
所以你要接待一個人.
當然接待其實是.
那個很散漫的故事.
那個受傷的是陌生者.

$^{1001}$你接待一個熟的人.
譬如Gary.
認識了很多年.
當然要接待他.
但不認識的.
我不認識你的.
我今天見到你的.
你說來.
你知道我去旅行.
我放一條樹給你.
當然不行.
我不認識你的.
接待是要冒險的.
愛人是要冒險的.
是沒有一個.
你不認識他們而信任的險.
今日我們社會.
人際關係越來越疏離.
越來越不信任人.
似乎令到我們教會裡.
都可能不信任.
弟兄姊妹.
這些東西要學.
重新再學.
人為甚麼不接待上帝呢.
你有沒有試過不接待上帝.
當然有.
甚至你今時今日.
有時你和我都.
沒有心理準備.
或者不接待上帝.
為甚麼呢.
通常是因為我們的罪.
你和我生命裡面.
有一些特別的.
很污穢的罪.
一些很黑暗的罪.
即是不見光的罪.
那些罪通常不會和人說.
最近沒有靈修.

$^{1041}$都會和人說.
有些罪你不會和人說.
你最親人都不會說.
有時上帝都不想.
在上帝面前提.
那些就是阻礙我們.
和人的關係.
和和上帝之間的關係.
所以我們不接待上帝.
封閉自己.
封閉自己的心門.
不歡迎上帝進入.
所以耶穌教導我們.
教導罪人.
要重新打開心門.
接待他.
接待一位.
耶穌是真光.
照亮人心黑暗的真光.
將我們罪人生命裡面.
將我們生命的內心裡面.
以致.
他和我們真真正正同在.
藉著聖靈和我們同在.
這個故事.
有一個很特別的地方.
表面上是人接待上帝.
那兩姐妹.
是罪人接待上帝.
不過其實.
首先是上帝接待我們.
上帝差遣了.
獨生子耶穌基督.
來到我們.
我們愛.
因為神先愛我們.
誰先是神先.
耶穌.
上帝差遣耶穌.
來接待我們.

$^{1081}$接納我們.
事實上.
當瑪大很忙碌.
在廚房預備一個.
豐富的筵席給耶穌的時候.
他預備接待耶穌的時候.
瑪利亞.
都是剛才說.
都是直接聽到.
接待耶穌.
其實耶穌已經預備了一個.
屬靈的.
很好的筵席.
接待瑪利亞.
所以那是上好的福分.
不能奪去.
各位姐妹.
我想你都同意.
我希望你都有感受.
我不管你的政見是什麼.
我希望你都應該感受到.
香港現在是和以前.
很不同.
我們現在是面對一個.
我們不是很知道將來會怎樣.
一個相對比較艱難的時候.
很多時候我去教會講道.
或者是特別講座完結.
很多時候.
弟兄姊妹問我.
甚至有些傳人問我.
趙先生你對教會將來.
怎樣走下去.
怎樣看.
我說沒得看.
見步行步.
沒有人知道會怎樣.
不過時勢艱難了.
時勢艱難的時候怎樣.
裝備自己.

$^{1121}$除了知識上的裝備自己.
當然人讀聖經.
不用我提.
環教會一直都很重視聖經.
讀多些歷史.
香港教會的歷史.
中國教會的歷史.
開始讀一下.
真的要讀一下書.
讀一下聖經.
除了這些知識上的裝備.
另一個裝備就是.
重新檢視自己的生命.
以前在安書的日子.
已經訓練我們.
有時我覺得安書日子.
未必是好事.
剛才我們有些文化訓練.
我們其實很不專注.
上帝在我們生命裡.
其實很多時候不賞心.
另一個裝備就是.
現在我們要重新檢視自己的生命.
怎樣繼續走下去.
一個人是不行的.
一班人.
互相鼓勵.
互相扶持.
互相提醒.
彼此大家在屬靈生命裡.
再一次提醒.
我們跟神結連.
然後真的重新檢視自己的生命.
究竟上帝在我們的生命裡.
擺了甚麼位置.
特別是當中已經退休.
你比一些很忙的人.
更加有條件.
容易來將上帝.
擺在你的生命裡佔據多些時間.

$^{1161}$好好來裝備自己.
以致在我們每一個崗位上.
日常生活裡.
真的作美的見證.
我們不知道將來會怎樣.
沒有人知道.
外在的環境我們掌握不到.
但是人的心.
我們可以學習.
可以修煉.
當然不是靠我們自己.
靠上帝.
我們一起低頭祈禱.
愛我們的三一上帝.
愛我們的主耶穌.
你真的來到這個世界上.
你進入了一個苦俗的世界.
我們要紀念你的誕生.
主一求你幫助我們.
以致我們可以學習.
專注於你.
我們願意盡心.
學習盡心盡性.
盡力地愛你和愛其他的人.
多於愛自己.
不過我們承認.
我們很多時候.
人性是我們愛自己.
我們很軟弱.
求主幫助.
求聖靈帶領.
我們這樣祈禱交托.
感恩和轉名給下一個.
\newpage



\section{撒母耳記下啟示錄 12:22-23}
\label{sec:S_q0D_6qEuY}
\textbf{2021.11.28 《 我信永生 》撒母耳記下 12\_22-23;啟示錄 21\_3-6 (和修版) 陳麗蟬牧師}
\newline
\newline
連結: \href{https://youtube.com/watch?v=S-q0D_6qEuY}{\texttt{https://youtube.com/watch?v=S-q0D\_6qEuY}} ~~~~ 語音日期: 2021-11-29
\newline
\newline
\hyperref[sec:sbhfqkg2NGI]{\small{< < < PREV SERMON < < <}}
~
\hyperref[sec:index_chronic]{\small{[返順時目]}}
~
\hyperref[sec:index_scriptual]{\small{[返順卷目]}}
~
\hyperref[sec:foX2kGxGGAY]{\small{> > > NEXT SERMON > > >}}
\newline
\newline
$^{1}$「擴寬Touchout」.
很開心見到群玉在這裡做見證.
我記得我初初去錦繡堂的時候.
我就認識了群玉.
除了一方面她是很親切之外.
另外就是因為她的名字跟我自己舞堂的.
一位很老友的姊妹是一模一樣的.
所以我很快就記得她的名字了.
還有有時候搭地鐵走的時候.
也會跟她同車過.
今天聽到她分享.
很感恩 很感恩.
我們在同行全訪的時候.
有時候我們就會遇到一些問題.
就是當我們跟她說.
「你要信耶穌」.
「如果你信了耶穌之後」.
「你就會有這個永生的了」.
但是人們又會說.
「我現在沒有空住」.
「等我將來有空一點」.
「我退休了」.
「或者到我將來臨終之前」.
「我才信」.
其實想想曾珍是這樣的.
臨終之前信耶穌也挺好的.
又不用上崇拜.
又不用參與事奉.
省下很多時間.
還有又不用做奉獻.
又可以省下很多錢.
但是呢.
在我們去讀《使徒信經》.
原來呢.
我很久沒有來洪恩講道了.
原來洪恩現在也在讀《使徒信經》.
你發到他.
在裡面講道.
「我信上帝」.
「全能的父」.

$^{41}$「創天造地的主」.
然後去到最後一句.
就是「我信永生」.
「我信永生」.
在整個《使徒信經》裡面.
就是在講我們的信仰.
在它開始的時候就說.
創造天地.
上帝創造天地.
但是這是以前的事.
是不是和現在沒有關係呢?.
上帝究竟有沒有繼續創造呢?.
到最後那句是.
「我信永生」.
這是不是將來的事.
又和我今生沒有關係呢?.
如果是這樣的話.
我用什麼來讀它呢?.
信耶穌.
是不是就是買了一張.
天堂的入場券.
那今生我可以喜歡怎麼做就怎麼做.
然後有什麼呢?.
臨終才改變.
好啊.
這個「我信永生」.
對於今生來說有什麼意義呢?.
在《使徒信經》裡面說到.
「我信上帝」.
「全能的父」.
「創造天地的主」.
當我們說這個信.
我相信的時候.
這個「相信」這個字的意思.
就是承認,同意.
並且對我所說的話.
有一個堅定的信念.
為什麼會有這個《使徒信經》出現呢?.
原來在初期教會的時候.
大約在一,二世紀的時候.

$^{81}$教會面對很多的異端.
人們在討論.
究竟耶穌基督是人還是神呢?.
究竟耶穌基督在十字架上這麼痛.
如果祂是神的時候.
祂不應該痛.
那不就是在裝模作樣嗎?.
但是也有些異端.
他就說了.
特別是當時很流行的.
「諾斯底主義」.
他就說.
在這個宇宙裡面.
有很多很多神.
就像我們現在中國人一樣.
信很多神.
滿天神佛的.
他們也是這麼想的.
那麼耶穌也是神.
不過耶穌是最低等,最低等的神.
所以當時對耶穌基督的身份.
很多人有很多不同的意見.
教會為了澄清這個人的信仰.
也叫所有信徒.
都很清楚知道自己在信什麼.
所以他就將我們整個信仰.
列成了一些條文.
讓信徒來背誦.
特別是洗禮的人.
一定要懂得背誦.
才能夠去接受這個洗禮.
「我信永生」.
這是《使徒信經》最後的一句.
是信徒.
我們每一個作為耶穌信徒.
一個最終最終的盼望.
但是究竟有沒有說永生呢?.
剛才主席為我們讀出.
華娟學位為我們讀出《撒母耳記》第十二章的時候.
她就說到背景說到大衛的兒子死了.

$^{121}$她很難過.
為什麼她的兒子會死呢?.
就是當大衛做了以色列王之後.
她非常的驕傲.
別人的將領在外面打仗.
打生打死.
她在家裡睡到黃昏白夜.
心裡的驕傲就出來了.
還要在一個很安逸的日子裡.
放縱她的肉體絲肉.
接著她看到一個美女.
看到對面屋的美女.
於是她就用自己皇帝的權力.
搶了別人的太太.
和別人發生了關係.
之後又殺死了別人對方的丈夫.
她就做了一件令上帝很生氣的罪.
上帝就說了.
你現在姦污了別人的妻子.
她所懷孕生的兒子必須要死.
因為上帝在告訴她這個事實的時候.
她已經說了這個審判.
但是戴維很不捨得.
所以中賢從這個先知裡面知道.
她的兒子必定要死.
不過她就覺得.
我試試去多點的祈禱.
又禁食祈禱.
日夜的祈禱.
通宵的祈禱.
她希望求上帝憐憫她.
以致這個兒子不需要死.
不過當上帝定了這個審判的時候.
她最終都要這個寶寶.
去回天家.
當時這個戴維很難過.
知道自己的祈禱上帝沒有垂聽.
但是她另一方面又知道.
神將這個寶寶帶回上帝自己國度裡面.
所以當她知道自己的兒子死了的時候.

$^{161}$她就這樣說.
寶寶都死了.
我何須要再禁食呢?.
我不能夠叫這個寶寶.
回到地上.
但是我知道.
我將來一定會去到某一個地方.
我會和這個寶寶一起的.
所以奇妙的是.
原來在這個舊約裡面.
上帝已經讓人知道.
這個人死了之後.
不會回到地上.
相反.
所有屬上帝的人.
將來都會在某一個地方.
可以再一次相聚.
這個地方就是在新約的時候.
我們來講的永生.
就是上帝自己國度的地方.
很有趣.
在舊約或者新約.
當我一講到彌賽亞.
在講耶穌基督的救贖的時候.
就很清楚告訴我們.
耶穌基督的救贖.
是和永生很有關係的.
我們有一節很熟悉的經文.
在約翰福音三章十六節裡面.
都是這樣講的.
我們一起讀這節的經文.
1,2,3.
神愛世人.
甚至將他的獨失兒子賜給他們.
叫一切信他的人不致滅亡.
反得永生.
這裡就是說.
只要你信耶穌基督的時候.
你就不需要去受這個永死.
反而你可以得到永遠的生命.

$^{201}$但這個永生是怎樣的呢?.
這個永生是否說.
我是持續永遠永遠都不會死的.
如果是永生的時候.
我們的樣子.
我們最擔心的就是.
我們的樣子是怎樣的呢?.
如果我們持續都是.
十幾二十歲皮膚很滑的時候.
就真的沒有死.
但如果說.
如果我永生的時候.
我變成一個千歲人魔.
皺紋皺得不得了.
這麼醜怪的時候.
我們就可能想過.
會不會去這個永生.
但當我們想想.
在這個永生的時候.
如果我和一些.
我很討厭.
我很不喜歡的人.
走在一起的時候.
我會怎樣呢?.
有一位姐妹跟我聊天.
她說她和她丈夫吵架.
吵到在一個很衝動的情況下.
她和她丈夫說.
我很討厭見到你.
如果有第二世的時候.
我一定不會再嫁給你的.
其實這樣說很傷心的.
但想想.
如果在永生裡面.
和一些我們很討厭.
很不喜歡.
很討厭的人.
永遠永遠走在一起.
苦不苦啊?.
比下地獄還要慘.

$^{241}$試過一次在搭西鐵的時候.
人都很擠迫.
我都站在那裡.
有一個.
我想她不知道是不是.
剛剛看完醫生.
那些衣服又不是很整齊的人.
可能我覺得她有一點抗議.
抗議沒有人讓位給她.
她就故意把衣服拿開.
有一個傷口在那裡.
有沙布貼在那裡.
那些沙布又有血在那裡.
她就故意彎腰.
然後露出傷口給你們看.
露出這裡有血在流.
然後又弄出很慘的樣子.
我不知道她是真還是假.
不過仍然沒有人讓位.
但是我們站在那裡的人.
我就在看.
下一個站開始有人下車.
我也在想我下不下車.
因為跟她再坐幾個站.
不知道會怎麼樣.
如果在永恆裡面.
經常跟一些這麼骯髒.
或者一些這麼令我不喜歡的人.
那時候怎麼辦.
怎麼辦是永恆的.
在啟示錄裡面的經文就告訴我們.
到將來的時候是全新的.
整個警方是全新的.
他說我又看見一個新天新地.
先前的時.
先前的天和地都已經過去了.
海也都不會再有.
我又看見星辰.
新耶路撒冷由神那裡從天而過.
預備好了.

$^{281}$這個新耶路撒冷就好像新娘一樣.
打扮整齊等候丈夫.
其實是等到我們教會.
講到教會就好像一個新娘一樣.
等候耶穌基督.
所以大家放心.
在將來的時候.
整個的情況是很開心很高興的.
是一個很華麗.
一個很輝煌的場景.
到那時候你不用怕了.
不會再有一些你很討厭.
或者很骯髒的一些事情.
或者一些令你很嘔心很激氣的事情發生.
更加不需要怕.
那個樣子不會變成千年人魔.
在聖經那裡講到.
不再有眼淚.
不再有死亡.
不再有悲哀.
不再有哭號.
不再有痛苦.
你放心.
上帝是那位獨行其事的神.
祂自有祂自己的辦法.
好了.
為什麼可以這樣呢.
在這個啟示錄裡面講.
看啊.
神的帳幕在人間.
祂要與他們同住.
他們也要作祂的指紋.
神親自到來的.
跟他們同在.
在這裡.
祂說.
神的帳幕在人間.
這個帳幕.
就是講到在舊約的時候的回望.
講到在當時的宗教中心.

$^{321}$有這個回望的時候.
象徵著上帝與他們一起.
只要上帝與他們一起的時候.
我們知道上帝是美善的源頭.
在他們裡面.
一切不好的事情都會變為好.
上帝是能力的源頭.
就算我們多麼軟弱.
但是在上帝的裡面.
我們都可以變為剛強.
總之在上帝的裡面.
一切都會變為美好的.
如果是這樣的話.
這個永生.
究竟是不是將來的事情.
與現在無關呢.
在剛才我們讀了的那段經文.
講到信耶穌與永生.
是緊緊相連的.
是緊緊相連.
只要信耶穌就不自滅亡.
反得永生.
所以當我們說.
我讀《使徒信經》說.
我信的時候.
我信上帝.
是全能的主.
創造天地的主.
不單只信上帝是過去的神.
也是今天的神.
也是將來的神.
我們現在要怎樣呢.
我們就要活出.
這個昔在今在以後永在的神.
我們就要活出.
這位神一切的美善.
曾經看過一本書.
很好看的.
可能在座當中也有人看過.
在天堂裡面遇見的五個人.

$^{361}$是講到這個作者.
是一個舅父.
舅父臨死之前.
可能我們有時也會遇過.
我爸爸也是.
當他臨死的時候.
他會見到很多以前的親戚.
這個舅父就是.
這個舅父叫艾迪.
這個作者就藉著.
他舅父臨終的時候.
見到很多的親戚.
他就來寫了這本小說.
他的舅父艾迪.
他是一個性情很古怪.
很古癖的那種.
經常都埋怨的.
每樣事情他都不喜歡.
煽得樹葉都掉下來.
但是他是一個很忠心的人.
他在一個遊樂場裡面.
去做這個維修員.
有一天.
當他在這個遊樂場周圍.
去看看哪裡要修.
哪裡要修的時候.
他突然間看到.
他修那個摩天輪.
有其中一個車廂.
在搖搖欲墜.
他就想起了.
他在維修摩天輪的時候.
其中一個車廂.
有些螺絲還沒有上實.
但是開始在運行.
他就一直看到.
其中一個車廂開始浮了出來.
正向整個車廂.
撞向一個女孩.
他想想這個女孩很危險.

$^{401}$於是他就飛身一撲過去.
想救那個女孩.
然後他就不省人事了.
就不知道發生什麼事了.
不知道過了多久.
他就醒了.
醒來的時候.
他發現自己仍然在遊樂場那裡.
但是整個環境很寧靜.
很舒暢.
他也覺得整個遊樂場.
他已經在那裡幾十年了.
以前覺得他很殘舊.
但是今天.
他看到整個遊樂場.
他自己做完摩天輪.
變得非常燦爛.
那個色彩很舒服.
很明亮.
最奇妙的.
就是他原本的身體.
周圍又肩膀痛又腰骨痛.
因為幾十年了.
在遊樂場裡面做維修.
又要搬很多重物品.
所以周圍都身體痛.
但是那個時候.
他發覺一切都沒事了.
最奇妙的就是.
他原本是跛的.
他原本是有一隻腳.
在打仗的時候.
被人射中了.
他跛了.
但是不知道為什麼那時候.
他突然間整隻腳可以走路.
一點都不痛.
所以覺得真的很奇妙.
為什麼會撲一撲救人.
可以變成這樣呢?.

$^{441}$但是同時他也察覺到.
原來他自己死了.
他可能已經去到天堂了.
突然就有個人出現告訴他.
他說.
你一會兒將會遇到五個人.
這五個人.
都會跟你在生平的事情裡面有關.
他會跟你聊天的.
於是他就去遊樂場走路.
走著走著的時候.
他就看到了.
他在年輕的時候.
他打仗的時候.
和他一起打仗的其中一個朋友.
這個朋友當時是做隊長的.
他很欣賞的地方.
當時隊長看到他的時候.
就跟他說.
艾迪很久沒見了.
一直跟他說起.
在打仗的時候的事情.
但是說著說著的時候.
他發覺原本這個地方.
我剛才說了.
很光明,很燦爛的.
但是突然間所有的色彩.
變成了灰色.
那些樹葉都枯了.
葉都掉下來了.
隊長告訴他.
因為他在戰場上.
他都同樣犧牲了.
他死了.
難怪變成這樣.
但是隊長跟著告訴他.
告訴艾迪.
你知不知道.
你那隻跛腳.
是我射到你跛的.

$^{481}$這個艾迪很生氣.
你為什麼無端射我的腳.
以致我跛了.
我沒有前途了.
蹲在這個遊樂場做了幾十年.
這個艾迪很生氣.
生氣到他拿著拳頭.
就上來打這個隊長.
但是這個隊長告訴他.
為什麼我當時開槍打你呢.
就是因為你的情緒非常激動.
叫你走,不肯走.
你如果不走下去的時候.
你會害死全隊的兵.
所以我唯有出此下策.
但迪森仍然覺得很生氣.
為什麼你害得我一生沒有成就.
然後這個隊長又問艾迪.
喂,你當日受傷了.
醒來的時候.
你有沒有再見過我.
艾迪就想.
是呀,沒有再見過他.
隊長就告訴他.
在當時你暈倒了.
我們繼續逃亡.
去到我們地方.
知道敵人就安裝了地雷.
隊長他自己自願.
自己下去走走.
測試一下哪裡有地雷.
他走著走著走著的時候.
他突然聽到天上有直升機來.
救援隊來救他的時候.
他很興奮.
所以突然興奮的時候.
他就疏忽了.
腳就不經意踏中了地雷.
結果這個隊長就粉身碎骨.
所以我死了.

$^{521}$然後慢慢跟艾迪解釋.
你當日不肯走.
我就寧願犧牲了你的腳.
但我保住全隊人的性命.
你一輩子為了自己沒有了腳.
一輩子都在憤怒的生活.
一輩子又變成了古叔.
什麼都埋怨.
值不值得?.
你是不是要將這個憤怒.
帶到永生呢?.
隊長接著就說.
說句真話.
我都犧牲了.
我為了全隊人.
我自己主動去測試地雷.
究竟在哪裡.
但就是因為我走錯了一步.
我粉身碎骨.
我不能夠再見到.
我在世界上的妻子.
見不到我自己的兒子.
也見不到我父母.
不過我犧牲了自己一條命.
我可以救了整隊人的生命.
有時我們好像失去了一些很珍貴的東西.
但如果我們失去了這些珍貴的東西.
但換來其他人得到很多益處的時候.
這個是不是值得呢?.
我當時也很感動.
現在才明白了.
為什麼隊長要射他的腿.
我很感動.
於是伸出手來抱緊隊長.
這時候他又發現.
整個遊樂場.
原本那些變成了很灰暗的顏色.
突然又光明又燦爛了.
原本那些樹又枯了,又落葉了.
現在竟然又出了樹葉.

$^{561}$又出了花,又出了果子.
整個遊樂場變得很有色彩.
我發現原來.
當人處於一個.
肯要恕,肯彼此接納的時候.
整個人生周圍的環境.
就充滿了色彩,充滿了光明.
但當整個人生都處於一個埋怨.
在一個憤怒的時候.
周圍的顏色都是灰暗的.
我只說其中一個人.
如果說完五個人就要分五堂課講.
原來當我們在地上生活的時候.
我們將來上帝說.
不再有眼淚,不再有痛苦,不再有憤怒.
不再有不好的事情的時候.
如果在今生的裡面.
我們能夠將這個妖樹.
接納,活出來的時候.
這個不用去到永生,都可以拿到.
所以永生不是長命百歲.
不是將日子不斷地而生.
而是以馬來利.
神的帳幕在人間.
是上帝與人永遠同在.
上帝的人愛,上帝的妖樹.
接納時刻與我們同在.
而我們就要將上帝的人愛妖樹.
接納自己來走出來.
天國也不是一個地方.
一無是事,無所事事.
週日就唱歌,Hallelujah,Hallelujah.
是一個歡樂天地,不是的.
而是我們與這位永恆的上帝.
可以很親密地走在一起.
在當中有喜樂,有平安,有福祿.
但不是上帝說他是一個很獨行其事的神.
他沒有難勝的事.
所以原本有這個妖樹.
原本有仇恨,憤怒.

$^{601}$突然間止一止,全部變得沒有了.
不是,而是上帝要我們在今生的裡面.
就要走出來.
所以千萬不要將這麼多的憤怒.
這麼多的積怨積在我們的生命裡面.
變得我們整個生命是憂憂愁愁,古古叔叔.
所以永生不是將來的事.
作為信徒的你和我.
我們就要與耶穌基督同死.
將我們舊有的生命趕走.
不再作罪的路步.
將一切的仇恨,疾度,苦毒去致死.
唯有這樣我們才可以在今天裡面.
活出一個復活的新生命.
將將來那份沒有眼淚,沒有憂愁.
沒有痛苦,沒有仇恨的永生的福樂.
在今生的裡面活現出來.
我們一起禱告.
讓我們知道在你的裡面我們有永恆的福樂.
在你的裡面我們就去領受到.
享受到那份被遙恕,被接納.
在地上高,我們實在遇到很多人.
我們真的很討厭,我們很不喜歡.
主但讓我們有耶穌基督那份的遙恕.
有耶穌基督那份的接納.
主我們知道不容易做的.
但你懇求聖靈,你加力給我們.
加這個愛給我們.
以致到我們可以能夠努力.
我們可以能夠做得到.
天父上帝求主工作.
奉靠主耶穌基督得勝靈治而求.
阿們.
\newpage



\section{撒母耳記上 13:5-14}
\label{sec:foX2kGxGGAY}
\textbf{2021.12.05 《事奉者的眼界》 撒母耳記上 13\_5-14 (和修版) 老耀雄牧師}
\newline
\newline
連結: \href{https://youtube.com/watch?v=foX2kGxGGAY}{\texttt{https://youtube.com/watch?v=foX2kGxGGAY}} ~~~~ 語音日期: 2021-12-08
\newline
\newline
\hyperref[sec:S_q0D_6qEuY]{\small{< < < PREV SERMON < < <}}
~
\hyperref[sec:index_chronic]{\small{[返順時目]}}
~
\hyperref[sec:index_scriptual]{\small{[返順卷目]}}
~
\hyperref[sec:ds7i5doUxn0]{\small{> > > NEXT SERMON > > >}}
\newline
\newline
撒母耳記上 13:5-14
\newline
\begin{longtable}{cl}
\hline
\hline
章節 & 經文 (和合本修訂版)\\
\hline
13:5 & \begin{tabularx}{0.7\textwidth}{X} 非利士人集合，要與以色列人作戰。他們有戰車三萬輛，騎兵六千，士兵像海邊的沙那樣多。他們上來，在伯‧亞文東邊的密抹安營。 \end{tabularx} \\ \\ \relax
13:6 & \begin{tabularx}{0.7\textwidth}{X} 以色列人見自己危急，軍隊被圍攻，百姓就藏在山洞、叢林、巖隙、地窖和深坑中。 \end{tabularx} \\ \\ \relax
13:7 & \begin{tabularx}{0.7\textwidth}{X} 有些希伯來人過了約旦河，逃到迦得和基列地。掃羅還在吉甲，所有的人都戰戰兢兢地跟隨他。 \end{tabularx} \\ \\ \relax
13:8 & \begin{tabularx}{0.7\textwidth}{X} 掃羅照著撒母耳所定的日期等了七日。但是，撒母耳還沒有來到吉甲，百姓就離開掃羅散去了。 \end{tabularx} \\ \\ \relax
13:9 & \begin{tabularx}{0.7\textwidth}{X} 於是掃羅說：「把燔祭和平安祭帶到我這裡來。」掃羅就獻上燔祭。 \end{tabularx} \\ \\ \relax
13:10 & \begin{tabularx}{0.7\textwidth}{X} 他剛獻完燔祭，看哪，撒母耳就到了。掃羅出去迎接他，向他問安。 \end{tabularx} \\ \\ \relax
13:11 & \begin{tabularx}{0.7\textwidth}{X} 撒母耳說：「你做了甚麼事啊？」掃羅說：「因為我見百姓離開我散去，你又不照所定的日期來到，而且非利士人已在密抹集合； \end{tabularx} \\ \\ \relax
13:12 & \begin{tabularx}{0.7\textwidth}{X} 我說：『現在非利士人已經下到吉甲來攻擊我，可是我還沒有向耶和華禱告。』所以我就勉強獻上燔祭。」 \end{tabularx} \\ \\ \relax
13:13 & \begin{tabularx}{0.7\textwidth}{X} 撒母耳對掃羅說：「你做了糊塗事了，沒有遵守耶和華－你神吩咐你的命令。不然，耶和華會在以色列中堅立你的國度，直到永遠。 \end{tabularx} \\ \\ \relax
13:14 & \begin{tabularx}{0.7\textwidth}{X} 現在你的國度必不長久。耶和華已經尋著一個合他心意的人，立他作百姓的君王，因為你沒有遵守耶和華所吩咐你的。」 \end{tabularx} \\ \\
[1ex]
\hline
\hline
\end{longtable}
$^{1}$各位弟兄姊妹平安.
在信仰來說.
我們每一個都是罪人.
而且每一個都是軟弱的人.
就算信了主也好.
我們仍然都是一個很軟弱的罪人.
信徒面對不同的困難和挑戰.
都有不同程度的軟弱.
有信徒可以跨過.
亦有信徒是軟弱跌倒.
在我們乘勝的旅程上.
有沒有一些方法.
能夠令我們不容易軟弱跌倒呢?.
答案當然是有.
這亦是我們每一個星期回來崇拜.
去聆聽神話語的目的.
神藉著祂的說話.
讓我們信心堅固.
信徒亦可以藉著神的提醒.
避免犯上相同的錯誤.
今日我想和大家探討一個題目.
就是「事奉者的眼界」.
一個事奉神的信徒.
他可能會遇到一些盲點.
我希望從《摘武耳記》上.
十三章五至十四節.
和大家一起思考一下.
在事奉上我們需要有一個甚麼態度.
我們一起有個祈禱.
可以在洪水橋這個地區服侍.
透過與不同需要的群組接觸.
可以將福音廣傳.
求你讓每一個參與服侍的弟兄姊妹.
都能夠按著神的帶領.
忠心地做好主所交託的即時榮耀主名.
我們祈禱事奉主耶穌基督.
得勝名字祈求.
阿們.
開始的時候我們先了解一下.
經卷的背景.

$^{41}$《摘武耳記》上和下.
在聖經的分類裡.
都是屬於歷史書的系列.
歷史書主要是講.
撒母耳先知,素羅王和大衛王.
三個人的故事.
歷史書所關注的是.
過去的真實世界的人和事.
而且很多的內容.
都是以故事的形式去呈現的.
既然是講故事.
必然是一個講故事的人.
因此故事有些情節的佈局.
氣氛的營造.
對白的安排上.
都有講故事的人.
很想表達的一個訊息.
而且一個好的故事.
亦離不開有幾個要素.
分別就是情節,角色,場景.
和敘事觀點.
甚麼叫情節?.
情節就是故事裡面的衝突事件.
一個故事如果沒有衝突的事件出現的話.
故事就是一潭死水.
很平淡.
而衝突事件的高潮.
必然是需要讀者去關注的一個地方.
在這段經文裡面.
衝突事件就是一個獻祭禮儀的執行.
素羅王是否應該審時度勢.
去獻繁祭和平安祭呢?.
抑或他要最終都要等到撒母耳來到.
由撒母耳去主持大局呢?.
在角色的刻畫裡面.
今次的主角是素羅和撒母耳.
而以色列人和非利士人都是Catafari.
透過素羅和撒母耳的對話和行為.
讓我們知道整件事的來龍去脈.
而場景就是角色行動的空間.

$^{81}$故事背景是素羅準備和非利士人打仗.
雙方都在集結兵力.
非利士人有戰車三萬兩.
騎兵六千.
士兵好像海邊的沙一樣多.
真的可以用來勢洶洶.
不能輕敵.
而以色列那邊呢?.
百姓有些躲在山洞,叢林,深坑裡面.
有些還過了約旦河.
離開了交戰的地方.
很明顯.
這個對比讓我們知道.
以色列人不夠齊心.
就算跟著素羅那些.
都是甚麼?.
驚驚呻呻.
不單止這樣.
素羅一方面的軍隊.
見撒母耳還未到.
已經開始有人打退堂鼓.
想散水.
素羅在這一刻應該怎樣處理呢?.
繼續等撒母耳來到.
才獻祭.
任由百姓離開.
當時的兵力已經很懸殊.
如果百姓都離開.
沒有人怎麼打仗?.
抑或是自己先去主持梵祭.
獻了梵祭和平安祭.
藉此.
藉一個獻祭的禮儀去穩定軍心呢?.
可以這樣說.
這個故事是甚麼呢?.
是一個危機處理的場景.
素羅身為以色列王.
他要帶領百姓去面對這個危機.
敘事者透過一個危機處理的場景.
藉著素羅和撒母耳的對話.

$^{121}$向讀者表達出一個訊息.
我們現在一起去發掘一下.
這個訊息的內容.
首先.
敘事者向我們表達出.
事奉者應有的第一種態度.
就是要遵從神的吩咐.
素羅在一個雙方交戰的危機裡.
他沒有等到撒母耳.
而親自獻梵祭和平安祭.
如果你說他錯了.
那錯在哪裡呢?.
一般人認為.
素羅本身是一個君王.
但他不是一個祭司.
他不應該代撒母耳去獻祭.
他似乎是僭越了祭司這個職份.
這個很多時候都.
普遍人認為.
這個是素羅所犯的最基本的錯誤.
當然有很多人會為素羅不值.
因為.
你撒母耳遲到.
你遲遲都不來.
百姓都走了.
素羅這樣做都是要穩定軍心.
他想留住那些人.
危機處理.
就是要靈活多變.
就算是錯.
是不是都情有可原呢?.
經文在第十三節.
透過撒母耳和素羅的對話.
他做了一個總結.
撒母耳怎樣說?.
他說你做了糊塗事.
你沒有遵守神所吩咐的命令.
罪事者認為素羅獻祭.
是一件沒有遵守神命令的回號事.
我們立刻要問.

$^{161}$為什麼素羅的罪名.
由僭越祭司職份的罪.
變成沒有遵守神的命令呢?.
變得這麼快?.
如果我們再看前一點.
在第十章的八節裡面.
我們就知道了.
當時撒母耳高立了素羅之後.
就已經即時跟他說.
你當在我耳先下到吉甲.
我也必下到那裡.
獻繁祭和平安祭.
你要等候七日.
等我到達那裡.
指示你當行的事.
撒母耳作為一個先知.
他是按照神的旨意去高立素羅的.
也是按照神的吩咐.
叫素羅去吉甲.
叫他去吉甲等候神的吩咐.
簡單來說.
素羅去到吉甲.
最重要就是等命令.
等待神進一步的指示.
而不是代替撒母耳去獻繁祭.
所以素羅心急.
在撒母耳還未到的時候.
私自獻祭.
就是不遵從神的吩咐.
素羅作為一個被神揀選的王.
神的僕人.
在第一次領軍對抗非利士人的戰事裡面.
犯了這個錯誤.
戰事還未開始.
就已經不遵守神的吩咐.
不去等候神的指示.
如果罪事者想我們以史為鑒.
我們就要謹記.
作為一個事奉者的首要心態.
就是要遵從,聽從神的吩咐.

$^{201}$洪英佳的弟兄姊妹.
我說的事奉者.
就不單單只是參與一些有形事奉崗位的會友.
即不是剛才的主席,領唱,詩禽,詩侍.
當我們回教會崇拜.
不遲到.
不構成別人的滋擾.
這個可不可以算是一種事奉?.
當我們在座位裡面.
專心地等待與神相遇.
不去與人聊天.
不去滋擾其他人.
可不可以是一種事奉?.
當我們同心去讚美.
獻上咀唇的祭的時候.
每一個人都獻出咀唇的祭.
可不可以是一種事奉?.
因此每一個參與敬拜的信徒.
嚴格來說都是一個事奉神的人.
你與我聚集在教會.
其中一個目的.
就是要聽神的吩咐.
針對個人來說.
就是要我們盡心盡性盡意盡力愛主你的神.
其次就是愛人如己.
對於這個吩咐我們應該一點都不陌生.
但如果針對一個群體呢?.
如果針對一個教會來說.
我們就很少去討論.
神對教會的其中一個吩咐.
就是要讓教會成為地區的燈台.
就是讓人認識耶穌.
讓人認識真理的地方.
每一個人來到教會.
除了要敬拜神.
聆聽神的話語之外.
也要對身邊的人.
或者社會近期發生的事.
能夠予以回應.
或者是予以施以援手和代禱.

$^{241}$所以教會崇拜是一個敬拜神的信仰行為.
我們不應該有一種有空就回.
沒空就不回的態度.
教會的崇拜也不是為了滿足信徒.
追求感覺的活動.
而是一個聆聽神話語.
和接受神差遣的聚會.
所以更加不應該有一種.
感覺開心我就回.
感覺我不開心我就不回的態度.
教會崇拜更加不是一個為了滿足信徒.
實踐消費主義的地方.
而是實踐一個甚麼?.
行公義,好憐憫的地方.
所以不應該有一種.
聽到聽得開心就回.
聽得不開心就不回的態度.
各位弟兄姊妹.
我們每一個.
你和我.
都有份參與洪恩家的事奉.
無論你是有一個實質的事奉崗位.
抑或你只是坐在那裡參加這場崇拜.
你都是一個事奉者.
所以請大家謹記.
敬拜神和聆聽神的話語.
實踐神的命令.
才是大家為何回到來.
坐在那裡的焦點.
而不是因著我們的喜好.
我們的感覺.
或者我們是否有空.
才回到崇拜.
甚願洪恩家能夠成為洪水橋這個社區的燈台.
神的光,神的話.
是照耀這個社區.
事奉者應有的第二種態度.
就是要耐心等候.
不要走在神的前面.
在講故事裡.

$^{281}$角色之間的對話.
亦都是罪事者.
很想表達訊息的一種很重要的工具.
在素羅和撒母耳的對話裡.
我們可以觀察到.
素羅對於獻祭一件事.
他不認為自己是有錯.
所有責任應該是撒母耳去承擔.
經文第八節說.
素羅照著撒母耳所定的日期等了七天.
撒母耳還沒有來到吉格.
第十一節.
素羅回應撒母耳的提問.
他說:因為我見百姓離開我散去.
你都不照所定的日期來到.
而且非利士人已經聚集在墨沒了.
素羅認為.
正當大軍壓境的時候.
百姓又離我而去之時.
你啊撒母耳還未到.
也不照所定的日期來到.
於是觸發了整件事的因.
所以十二節.
素羅說:我就勉強獻上繁祭.
這是事件的果.
經文以勉強來形容.
素羅獻祭時的心態.
原文有強迫自己的意思.
所以《詩高聖經》亦為迫不得已.
素羅似乎想說.
釀成這個局面.
不是我的錯.
是你撒母耳遲到.
你應該承擔一切責任.
如果撒母耳不遲到.
素羅就不用迫不得已去獻繁祭和平安祭.
一切都是你這個遲到大王.
撒母耳遲到才搞成這樣.
是不是呢?.
罪事者在第十節說.

$^{321}$素羅剛剛獻完祭.
撒母耳就到了.
這個就像電視的情節.
獻完祭就到.
在時間上作出這麼巧合的安排.
這麼戲劇化.
罪事者在時間上.
為什麼要用這個剛剛的表達手法呢?.
他是有意還是無意的?.
如果是有意的話.
他想我們留意些什麼?.
如果要留意的話.
我就很想知道.
究竟素羅在吉格所等待的七天.
是不是已經完成了.
就是那七天完全結束了.
獻祭的時間是第八天.
素羅才迫不得已去獻祭呢?.
抑或素羅還在第七天的中途.
那天還未完結的時候.
他就已經急不及待呢?.
如果我們再看第十五節.
撒母耳來到之後.
他就馬上離開了.
而素羅還有時間去數點人數.
所以估計其實第七天還未完.
罪事者認為撒母耳沒有遲到.
撒母耳在第七天當天出現的.
不過素羅太心急了.
在危機和壓力面前.
素羅選擇了運用自己的智慧去解決.
而不理會.
神叫他.
你要等待.
面對很強大的壓力.
我們往往都很心急.
我們很想盡快解決這個問題.
平息這件事.
我們往往忘記了作為一個事奉者.
應該注意的其中一個心態.

$^{361}$就是我們不應該走在神的前面.
我們總要跟在神的後面.
素羅用自己的智慧和經驗去處理一個危機.
他不願意等待神進一步的指示.
這成為我們一個反面的教材.
大家都聽過龜兔賽跑的故事.
不過我今天不是講賽跑.
我想以龜和兔做一個比喻.
烏龜雖然走得慢.
用慢的字就最適合.
但只要他方向正確.
他就跟目標越來越近.
縱然是慢.
但只要方向正確.
他每過一天.
他都跟目標越接近.
兔子跑得很快.
但方向不正確.
他越跑得快.
距離和目標是甚麼?.
越遠.
所以關鍵不在於快和慢.
關鍵在於方向是否正確.
只要方向正確.
多慢也好.
你都是跟目標越來越接近.
跑得快沒有用.
方向只要錯.
越跑得快.
距離和目標越遠.
離神的心意越遠.
如果我們不願意在神面前等候.
然後才確定睛確方向才起步的話.
而喜歡走在神的前面.
按著自己的心意而行.
按著自己的心態.
說做就做的話.
我們就有機會成為.
走錯方向的兔子.
越跑得快.

$^{401}$距離和目標越遠.
事奉上徒勞無功.
侍沛而功半.
最後走到筋疲力盡.
還質問神.
為何我這麼熱心事奉?.
為何我這麼委身?.
你都不保守我.
你都不幫我.
最後受傷離開.
從此不再事奉.
這就是每一個事奉者.
如果他的事奉走在神的前面.
最後的結果.
事奉不力.
事奉到燒傷.
事奉到離開了職場.
離開教會.
離開信仰.
這永遠都不是.
合神心意的事奉.
各位事奉著的弟兄姊妹.
我總希望你能夠停一停.
想一想.
你事奉的方向是否對準神?.
你事奉的目標是否榮耀神?.
甚願我們不要做那種.
既走錯方向.
又跑得快的兔子.
事奉者應該有的第三種態度.
就是不能夠心中保神.
在這件憲制事件上.
很多人都會替這個訴羅不值.
因為訴羅雖然在這件事上有犯錯.
縱然是那種大錯.
但所受的懲罰.
就很令人吃驚.
而且覺得不合理.
犯錯而已.
用不著死罪.

$^{441}$撒母耳記第十四節說.
「你的王位必不長久」.
「而說已經沉著一個合他心意的人」.
「立他作百姓的君」.
難道一犯錯就要受到這麼嚴重的懲罰?.
訴羅不單失去了君王的身份.
就是連神的僕人.
事奉神的身份都失去了.
讀者看到這裡.
都不禁要問問自己.
在事奉的生涯裡.
總會有機會犯上.
不聽從神的吩咐吧?.
總會有機會犯上.
走在神的面前的錯誤吧?.
任何一個事奉者都有機會犯.
我自己也在想.
作為一個牧者.
是否也有機會被褫奪.
事奉神的身份?.
罪事者沒有交代.
究竟為什麼訴羅會被褫奪君王.
這件事件的深層次原因.
表面的理由就是不聽從神的吩咐.
但我們看看.
相比聖經其他人物.
其實也有犯過同樣的錯誤.
你看看先知約拿.
同樣不聽從神的吩咐.
神叫他去東,他就去西.
但約拿的下場比訴羅好很多.
如果我說.
神有主權.
神要憐憫誰就憐憫誰.
神要因對誰就因對誰.
你們敢向神強嘴?.
耶和華昔日可以輕而訴.
重亞國.
現在也可以輕訴羅.
重約拿.

$^{481}$神有絕對主權.
不行嗎?.
行….
如果你利用神的絕對主權這張王牌.
打這張大牌出來.
我當然要閉嘴.
不過這場戲不順.
神是有主權的.
不過你這個決定我不順.
金庸寫武俠小說.
很多人都看過.
我是其中之一.
一看就收不到掣.
一看就不放下.
因為情節很吸引.
很想追下去.
想知道結局.
講故事最高的技巧.
就是要吊著你.
不告訴你.
結局先不揭曉.
等續者繼續追下去.
我估罪事者都有這樣的想法.
他沒有即時說出終極的原因.
但是在四章之後的第十七章.
他用了同一個對比的手法.
以同樣的場景.
主角同樣需要面對危機.
透過對話.
讓讀者自己找出答案.
在第十七章.
大委身處的情況是甚麼呢?.
是在比他強大很多倍的.
非利士人哥利亞面前.
情況比塑羅所面對的危機更加大.
因為塑羅在吉甲的時候.
只是雙方軍隊集結的階段.
集結.
即是還在接班.
未正式交鋒.

$^{521}$距離正式開戰還有一段時間.
但是大衛當時是面對面.
直接和哥利亞賣身肉搏.
當時的以色列人包括塑羅在內.
全部感到很害怕.
不過大委身處其中.
他就說出了一個事奉者.
所必須要有的心態.
而這個必須要有的心態.
正正就是塑羅所沒有的.
在第十七章四十五節.
大衛向非利士人說.
「你來攻擊我.
是靠著刀槍和銅擊.
我來攻擊你.
是靠著萬軍的耶和華之名.
就是你所怒罵帶領以色列軍隊的神」.
當時塑羅還是以色列的王.
一個國家的領袖.
他和以色列百姓所看到的.
是一個強大的敵人.
塑羅和百姓因著對手的強大.
所以很害怕.
他看不到神究竟在哪裡.
而大衛看到的.
是打這場仗的主角.
主角並不是以色列.
也不是大衛自己.
主角是萬軍的耶和華.
這場戰事是由耶和華所掌控的.
作為一個神的僕人.
一切的力量,能力的來源.
都是來自背後的主人.
應該倚靠的是神.
並不是自己.
而這個信仰的核心.
正正就是塑羅.
他作為一個屬靈領袖.
在相同的處境裡面.
他所缺乏的心態.

$^{561}$作為一個事奉者.
如果心裡沒有神.
不覺得神掌管一切.
凡事都是依靠自己的知識和能力.
他怎樣去事奉神.
他怎樣去進行神的命令和吩咐.
如果這個正正就是罪事者.
很想讀者明白的.
一個事奉者必要有的心態.
當讀者回望塑羅所得到的懲罰.
就是要褫奪他作為一個王.
一個帶領者的身份的話.
那就變得相當合情合理.
他用自己的能力.
他帶領以色列.
他沒有一個屬靈的眼界.
他告訴民眾.
我帶領你們打這場仗.
而不是耶和華帶領你們打這場仗.
事奉者心裡沒有神的意思.
不是說他認知上不認識神.
不是.
而是他沒有讓神成為他生命的主.
他仍然以自己的心意.
以自己的智慧和能力去事奉.
所以他只能夠以自己的眼界.
去看待這個世界.
他沒有一個屬靈的眼界.
去看待身邊的人和事.
神有沒有參與其中.
神在當中的角色是一個甚麼人.
各位紅映家弟兄姊妹.
今日的講題就是事奉者的眼界.
事奉者的眼界就是叫我們.
不要以自己的眼光作出發點.
事奉的目標是神.
就應該以神的眼光作為出發點.
以耶穌基督的心為心.
因此事奉者有三個應有的態度.
就是要聽神的吩咐.

$^{601}$不要走在神的面前.
最不好的就是.
當事奉是一個工作.
要心裡面有神.
是神作王而不是自己.
如果每一個信徒都是事奉者的話.
而事奉者對教會的責任.
就是要建立基督的身體.
最簡單最簡單的一件事.
由今日開始.
認識一下你前後左右的朋友.
你知不知道你前後左右的人叫甚麼名.
你知不知道?.
不知.
或者知.
你問問他.
有甚麼可以為你守望.
有甚麼可以為他代禱.
你就在這個禮拜為他代禱.
你為你身邊的人一起翻崇拜.
你認識與否認識的會友.
你為他祈禱守望.
然後下次見面的時候.
你就可以彼此分享.
不久後.
你們可能成為一個土伴.
一起為教會土伴.
這是一個最簡單又極之有效的方法.
去建立大家的屬靈關係.
你不是單獨一個人來參與崇拜.
你是同一群屬神的子民.
一起參與崇拜.
這個就是我們《鑰匙條信經》裡面所講的.
「聖徒相通」.
你連隔離的那個都不知道他是甚麼.
通甚麼呢?.
無法通.
你怎樣通得到呢?.
彼此守望代禱.
可以彼此服侍對方.

$^{641}$《聖經》說彼此相愛.
你回來這裡.
你一個人參加崇拜.
怎樣彼此相愛?.
如果我們每兩個禮拜.
不斷循環去實踐彼此代禱.
彼此祝福的功課.
我很肯定.
我真的很肯定.
洪恩嘉一定能夠成為社區的祝福和加油站.
不過你願不願意?.
你願不願意回應這個事奉者的呼召?.
還是你不覺得自己是一個事奉者?.
我在此邀請大家在座位上靜默一下.
沉澱一下我們內心的感受.
聆聽一下神對我們有甚麼提醒.
如果你對剛才我提出的建議有回應的話.
你也可以用禱告告訴神.
你的想法.
最後我們有一個結束的祈禱.
秋下因主我們求你月立每一個會友.
剛才在座位裡面.
向你的回應和祂的內心的心意.
求你加能設立給每一個願意服侍你的信徒.
讓他們都能夠實踐剛才在你面前的承諾.
秋下因主.
讓他們能夠真是做一個合你心意.
討你喜悅的事奉者.
一起高舉你的名.
讓洪一家真真正正成為這個社區的祝福和燈台.
而不是只是一個口號.
讓我們這間教會.
真是照亮整個社區.
讓人覺得.
如果洪水橋沒有了洪恩這間教會.
洪水橋就會失色.
而不是有和沒有都無關痛癢.
沒有了都不可惜.
我們的祈禱是奉主耶穌基督得勝明智祈求.
阿們.

$^{681}$願主祝福大家.
\newpage



\section{列王記上 22:14}
\label{sec:ds7i5doUxn0}
\textbf{2021.12.12《只為真道爭辯》 列王記上 22\_14 (和修版) 郭鴻標牧師}
\newline
\newline
連結: \href{https://youtube.com/watch?v=ds7i5doUxn0}{\texttt{https://youtube.com/watch?v=ds7i5doUxn0}} ~~~~ 語音日期: 2021-12-14
\newline
\newline
\hyperref[sec:foX2kGxGGAY]{\small{< < < PREV SERMON < < <}}
~
\hyperref[sec:index_chronic]{\small{[返順時目]}}
~
\hyperref[sec:index_scriptual]{\small{[返順卷目]}}
~
\hyperref[sec:yEoqVyieILA]{\small{> > > NEXT SERMON > > >}}
\newline
\newline
列王記上 22:14
\newline
\begin{longtable}{cl}
\hline
\hline
章節 & 經文 (和合本修訂版)\\
\hline
22:14 & \begin{tabularx}{0.7\textwidth}{X} 米該雅說：「我指著永生的耶和華起誓，耶和華向我說甚麼，我就說甚麼。」 \end{tabularx} \\ \\
[1ex]
\hline
\hline
\end{longtable}
$^{1}$我們一起同心祈禱.
我們親愛的天父上帝.
我們知道你是我們生命的主.
你知道我們的優點和缺點.
你知道我們的限制和軟弱.
你就是不斷賜給我們機會.
賜給我們恩典的神.
我們願意你的說話.
成為我們生活的力量.
讓我們面對各種考驗和挑戰困難的時候.
我們都能夠看到你一路的人靈.
你是一路教導我們.
建立我們.
求你藉著今天講道的時間.
祝福我們每一位弟兄姊妹.
祈禱奉耶穌基督的名.
阿們.
邀請我跟大家分享聖經的說話.
最近在一位學生的小組課堂的報告上.
聽到有兩節經文.
很吸引我的注意.
第一段經文.
就在尤大書的第一章第三節.
他說.
親愛的弟兄.
我想盡心寫信給你們.
論我們同基督教的時候.
就不能夠不寫信勸你們.
要為從前一次交付給聖徒的真道.
你要去竭力的爭辯.
第一就是要為真道爭辯.
到第二段經文.
就是提摩太后書第二章第十四節.
他說.
你要使眾人回想這些事.
在主的面前祝福他們.
不可為言語爭辯.
這是無益處的.
只能夠敗為聽見的人.
我聽見這個報告有兩節不同的經文的時候.

$^{41}$我就問一個問題.
究竟在什麼情況下我們可以去爭辯.
什麼情況下是不可以去爭辯呢?.
有一次聽到的時候.
我又聽到一段經文.
就是列王記上第二十二章.
講到阿哈王.
以及以色列王阿哈.
就是北面的.
就跟南面的猶大王約沙伐的故事.
另外有一些領袖.
可能大家不知道這個故事是什麼.
我一會兒會跟你們解釋一下.
由於我講到的經文其實是很長的.
為了方便.
我只選了列王記上第二十二章第十四節.
講到的題目就是.
我們只是為真道去爭辯.
講到的內容我怎麼處理呢?.
我就分兩個部分.
第一就是經文的解釋.
第二就是生活的應用.
首先我們看看經文的解釋是怎麼樣的.
大家應該聽過以色列王阿哈有關的兩件事.
可能你在主日學,查經班也聽過.
第一件事就是.
以利亞和四百五十個巴力先知比試的故事.
這件事記載在列王記上的第十八章.
由第十八章到我今天講到的經文.
二十二章還有幾章.
故事是在延續下去的.
第二件事就是耶駛別.
大家聽過這個人名.
他用假見證去陷害一個人.
叫做拿伯.
原因就是要奪取了他的葡萄園.
這件事記載在列王記上的第二十一章.
這兩件事都讓我們看到.
跟阿哈王有關係.
阿哈是一個怎樣的人.

$^{81}$我們有了這個基本的概念.
就可以進入列王記上的第二十二章.
列王記上的二十二章第一節.
講的經文就是.
亞蘭國和以色列國三年沒有戰爭.
列王記上的記載.
就是阿哈娶了一個西頓王的女兒.
叫做耶駛別.
其實是另外一個國家的.
耶駛別他就大力推動一種叫做.
二幫的宗教崇拜.
即是他將他鄉下的宗教帶過來.
不過以色列在政治上是相當成功的.
起碼三年跟隔離的亞蘭大國沒有打過仗.
你明白嗎.
這個經文告訴我們.
阿哈娶了一個耶駛別.
耶駛別就是將他鄉下的二幫宗教帶去以色列.
但是這個危機還沒有爆.
不過另一方面的政治環境是相當不錯.
因為亞蘭沒有跟他打仗.
你明白背景.
到列王記上的二十二章第三節.
就講到這件事了.
這麼穩定當然要想一些東西.
阿哈王就希望奪回.
基列的拉門.
即是要拿回他失去的地.
列王記上二十二章第四節就記載了.
以色列王阿哈就問猶大王約沙伐.
你會不會一起去打基列的拉門.
所以他想著自己一個不夠打.
就找南面的約沙伐合作.
猶大王約沙伐就說好.
我和你出兵.
但是要先求問上帝.
很正常.
以色列是這樣的.
要問耶和華.
列王記上的二十二章第六節就記載.

$^{121}$阿哈王就找了四百個先知.
這批先知都認為應該出兵.
到列王記上二十二章第七節.
接著猶大王約沙伐就說話.
他說這裡不是還有一個耶和華的先知.
我們可以去求問他.
你明白吧.
阿哈王他叫了很多先知.
只有四百個.
但是他有一個不叫.
為什麼不叫呢.
一會兒故事就來了.
列王記上的二十二章第八節就記載了.
以色列王阿哈對約沙伐和.
是沒錯還有一個.
他叫欽拉的兒子米該雅.
我們可以托他去求問耶和華.
不過後面那句.
只不過我是很恨他.
這個字用得很盡.
我真的痛恨他.
因為他指著我所說的二元.
不說吉語.
就是不說好東西.
單說空語.
約沙伐說.
王不必這樣說.
這篇經文你看到沒有.
裡面約沙伐知道有一個先知叫米該雅.
不過恨他入骨.
所以這次召集先知來問問題.
就不會找他.
他就說這個先知米該雅.
說的都不是好東西.
都不適合他聽.
所以就不找他.
但你有沒有留心到.
這篇經文最後那句.
南國的約沙伐.
就和北國的以色列王阿哈說.

$^{161}$王你不用這樣說.
你明不明.
阿哈王要約沙伐出兵一起打仗.
約沙伐就說.
問米該雅.
其實阿哈都不想的.
不想之中他都要做.
因為要靠別人一起打仗.
好了.
猶大王約沙伐就應該知道有米該雅這個先知.
他不會問這個問題.
你想想當四百個先知都說可以出兵的時候.
他問這裡還有一個先知.
我們可以請他問上帝.
其實問這個問題有什麼意思呢.
大家想一想.
既然四百個先知都有同一個答案.
多一個有什麼意義.
如果和這個米該雅都出兵.
他來和不來沒有分別.
是意料之內.
好了,反過來他說不要出兵.
一個人米該雅說有什麼用.
你明白約沙伐問這個問題.
其實按常理來說.
都沒有什麼必要.
但是在沒有必要的情況下.
神就有些事情做了.
很有趣.
當我們讀《熱狼紀》上二十二章的時候.
我們就應該謹記.
經文記載這件事是有教導作用的.
他不是白記在這裡的.
你想想剛才那件事這樣發生.
有什麼原因有什麼要學呢.
究竟這件事裡面是教什麼呢.
上帝需要米該雅這個先知.
說從上帝而來的話.
中史其他先知有不同的看法.
你想一想這件事就很爆炸性.

$^{201}$四百個先知都說同一番話.
你說另一番話.
你想想對著皇帝.
你想想這件事多麼爆炸性.
我很希望大家注意.
以色列王阿克是知道米該雅存在.
在傳召先知的時候.
就是偏偏不找他.
阿克對猶大王約沙伐說.
他很恨這個米該雅.
心裡面那句說了出來.
一個忠於上帝的先知.
可以令其他人討厭他.
你想一想.
恨他的.
不讓他參與的.
不給他機會的.
邊緣化他的.
架空他的.
以色列王記上的22章第9節記載.
以色列王阿克就派人去接米該雅來.
因為你答應了約沙伐要找他.
以色列王記上的22章第9節.
就記載阿克的使者.
就是他的馬仔.
就跟米該雅說.
米該雅.
眾先知都一口同音.
向王說吉言.
你不如跟他們說一樣的話.
都說吉言.
你明不明.
阿克的伙計就跟他說.
你識做些.
這次米該雅有機會參與.
阿克王的使者又提醒了米該雅.
其他先知都說了好話.
你都不如說一份.
米該雅的回答對我很大提醒.
在以色列王記上的22章第14節.

$^{241}$即是今天我們講到的經文.
米該雅就說.
我指著永生的耶和華.
來去.
起誓.
就是耶和華對我說些什麼.
我就說些什麼.
你看了整個經文的發展.
你就明白這一句.
列王記上22章第14節.
他教我們什麼道理.
你看下去你發覺原來很多東西要學.
大家覺得米該雅的態度.
是不是正確呢.
如果你問.
如果你是先知米該雅.
你又會怎樣回答.
你會不會.
咦 還是像米該雅那樣.
我又看看米該雅怎樣去.
落場對應.
在社會上其實大家都知道.
很多人都會看風洗理.
作為上帝的先知.
可能不能夠討好人.
不過只能夠忠於上帝.
列王記上22章第10節.
他就說所有先知都在以色列王阿克.
和約沙法王的面前說語言.
個個都出去做事.
列王記上22章第11節就記載.
有些人名出來了.
有個叫基那那的兒子.
叫西底加.
就是這些祭司級的.
做了兩個鐵角就說.
耶和華就這樣說.
你要用這個角抵觸阿難人.
直到將他們滅絕.
有高級數的先知就說.

$^{281}$我做了兩個角出來.
靠著這些就可以打敗阿難人.
基那那的兒子西底加.
就是以上帝的名字來支持出兵.
其他四百個先知都說.
是呀 是呀 是呀.
要呀 出呀.
你明不明.
就是說這個是四百個先知的頭.
基那那的兒子西底加.
究竟米該雅面對著四百個先知.
都支持出兵的時候.
他又怎樣回答呢.
我們看回列王記上的第22章第15節.
就說米該雅去到王的面前.
王就問她.
米該雅.
我們上去攻取基列的拉末.
可以不可以.
她就回答.
可以.
我只是配音.
因為她是經文上.
可以上去.
必然得勝.
你上吧 贏定了.
耶和華必定將城交給你.
有一點點語氣.
是講反話的.
從字面上來說.
米該雅是同意出兵.
不過你讀下去就知道她是反對的.
接著到列王記上的第22章第16節.
王就說話.
王就對她說.
我當祝福你多少次.
你先奉耶和華的名向我講實話.
阿哈王就說.
你別在那裡玩了.
我知道你說什麼.

$^{321}$以色列王阿哈就知道.
米該雅同意出兵那句話.
其實是反諷式的說話.
然後米該雅就說出她真正從上帝而來的信息.
那個就在列王記上的第22章第17節.
米該雅就說.
我就看見以色列的眾民.
散在山上.
好像沒有牧人的羊一樣.
耶和華就說.
這些民是沒有主人.
他們可以平平安安的各歸各家去.
你想想一個國家出兵.
士兵不可能在山上散羊.
不可能的.
你還說沒有牧人的羊群一般.
這裡還像兵一樣.
散兵游勇.
上帝就說.
這些民是沒有主人.
意思是什麼.
首領都不見了.
可能失蹤或者死了.
諷刺的就是.
首領不見了.
大家可以平平安安的各歸各家去.
你明白嗎.
米該雅在領受上帝的信息說出來.
就是列王記上的第22章第17節.
你明白嗎.
米該雅就說出真的說話.
再看第18節.
記載阿哈對著猶大王約沙伐說.
米該雅這個人就是專門說空話.
不說吉語的.
我都老早說了.
劇情就是這樣下去.
然後去到第19節.
米該雅就說.
你要聽耶和華的話.

$^{361}$我看見耶和華坐在寶座上.
天上的萬軍.
立在他左右.
你想想米該雅說出這番話.
她是聽到耶和華的話.
看見耶和華在寶座上.
天上的萬軍立在他左右.
你想想米該雅是不是亂說話.
大家明白嗎.
看這段列王記上的時候.
米該雅不是隨口說的.
到第20節.
耶和華就說.
誰去引誘阿哈上去激烈的拉末去陣亡呢.
這段經文說.
阿哈會打仗會死的.
這個就這樣說.
那個就這樣說.
當時.
你知百忌都有.
天上有天庭.
這個我們都不知道怎樣.
沒辦法想像.
但這段就說.
上帝就問誰去做這件事.
誘惑阿哈上去激烈打仗.
然後他就打仗死.
這個就這樣說.
那個就這樣說.
要配合上帝的計劃.
請大家注意.
上帝是要以色列的阿哈王上激烈的拉末陣亡.
但經文是沒有交代原因的.
為何走這一步.
不過以色列王阿哈.
還有一件事.
他縱容了450個巴力先知.
你記不記得.
和以利亞比試.
還有他令以色列人陷入異邦崇拜.

$^{401}$肯定要罰的.
以色列王阿哈在耶駛別用假見證陷害拿伯之後.
奪取葡萄園是有悔意的.
列王記上21章27節說.
阿哈聽見這個話就撕裂他的衣服.
禁食新村麻衣.
睡臥也穿著麻布並且緩緩而行.
即是說他有些悔意有些良知.
到21章29節記載.
上帝說臨到以利亞.
阿哈在我面前這麼自卑.
你看不看見.
因為他在我面前的自卑.
他還在世的時候.
我就不降和.
到他兒子的時候.
我就必降和.
在他家.
在上帝眼中阿哈都不壞.
還知道有些人性.
耶駛別陷害拿伯.
奪取了葡萄園.
他都會有些悔意.
但是以色列王阿哈.
雖然記得有上帝.
還有些自責.
但他這種宗教行為是沒有辦法彌補他之前所犯的錯誤.
譬如他縱容耶駛別那裡有巴力450個先知.
他引進異邦崇拜.
這件事上帝一定會罰他.
所以我們不要以為上帝不會罰人.
上帝有他的時間表.
我自己看到這裡.
我有一點都要提醒一下自己.
你怎麼知道你自己做的事.
上帝會怎麼看你呢.
以色列王記上下22章21節.
記載了.
隨後有一個神靈出來.
就走在耶和萬面前.

$^{441}$我去誘惑他.
故事就講到有一個僕人出來做事.
經文沒有交代這個神靈是誰.
無名無姓.
希伯來文就是Haruah.
英文聖經的NASB的本就是A Spirit.
總之是一個零.
你當他是上帝僕人就可以了.
總之上帝是僕人去做事.
就誘惑阿哈上去激烈的拉梅打仗.
接著又來看頭燈.
22章的22節.
上帝問那個靈.
耶和華問他.
你用什麼方法.
你這麼聰明就出來自告奮勇.
這個靈說.
我去要在他眾先知口中.
作謊言的靈.
耶和華說.
這樣你一定會引誘到他.
你去做吧.
上帝問他用什麼方法.
他說我用謊言的靈.
讓那些仙子都講謊言.
謊妙的話.
虛謊的話.
上帝說.
這樣應該行得通的.
行了你做吧.
烈皇紀上的22章23節.
記載米該亞對以色列王阿哈說.
現在耶和華就使用謊言的靈.
進入了你這群仙子的口.
並且耶和華已經命定降和給你.
你想想.
米該亞是不是很大膽.
阿哈你知道嗎.
耶和華使用謊言的靈.
進入了你這群小馬.

$^{481}$那400個小馬.
那你在那400個小馬會怎樣反應呢.
那就要看下面了.
大家記住米該亞只有一個人.
基拿拿的兒子西底加有400個.
加起來400個仙子是他的小馬.
米該亞的意思是.
基拿拿的兒子西底加為首的400個仙子.
全部都是被謊言的靈控制.
你想想.
基拿拿的兒子西底加聽到的時候.
你會不會火滾呢.
你有400個小馬在那裡.
你一個這麼大膽.
你會怎樣呢.
在烈皇紀上的22章24節就記載了.
基拿拿的兒子西底加就走前來了.
打了米該亞的臉.
一巴掌就升過去.
就說耶和華的靈從哪裡離開我來跟你說話.
基拿拿的兒子西底加說.
耶和華的靈什麼時候都在我這裡.
它什麼時候離開我來跟你說話.
你是誰啊.
你明白嗎.
在烈皇紀上的22章25節.
米該亞就說.
你從你進入嚴密的屋子.
躲藏的那一天.
你就一定看見了.
他不是跟他鬥氣.
他說基拿拿的兒子西底加.
到那一天你躲起來的時候.
你就看得見了.
你明白嗎.
他講這些東西.
米該亞的意思就是.
那時候你就知道了.
就是這樣的說話.
在烈皇紀上的22章27節就記載了.

$^{521}$以色列王阿哈就說.
王就如此說.
將這個人收入監獄.
使他受苦.
吃不飽 喝不足.
等候我平平安安的回來.
阿哈王就說.
抓他進去 關著他.
不要給他吃東西.
不要給他吃飽.
不要給他喝夠.
等我打完仗回來.
打贏仗才收拾他.
阿哈王講得很聰明.
接著你由心米該亞.
在烈皇紀上的22章28節講了一句什麼的話.
他就說.
如果你能夠平平安安的回來的話.
那就是耶和華沒有藉著我說話.
那就說.
中文你們都要聽.
你看這個米該亞多厲害.
你明白嗎.
當阿哈王說我要關你進監獄.
不給你吃飽.
不給你喝夠.
等我回來的時候好好對待你.
米該亞就說.
你回來再說吧.
你明白嗎.
經文你讀下去會發覺很有味道.
他說你回來平平安安的話.
那就是耶和華沒有藉著我說話.
就這麼簡單.
你想想一個米該亞對著四百個.
號稱是耶和華的先知.
他夠膽這樣說話.
是吧 你想想.
結果就是阿哈王死了.
我不再說這麼長了.

$^{561}$那個故事很長.
結果就是這樣.
當我讀這段經文的時候.
我都想了很多東西.
究竟這段經文教了我們什麼呢.
那就到了生活的應用.
《列王記上》22章14節.
米該亞就說.
我指著永生的耶和華起誓.
耶和華對我說什麼.
我就說什麼.
這個是不是很高難度.
你明白嗎.
這個很高難度.
《列王記上》記載.
耶洗別有450個巴力先知.
阿哈有基娜娜和兒子西底加帶領著四百個先知.
這四百個人名義上都是耶和華的先知.
名義上.
另外阿哈他的家宰.
有一個叫俄巴底.
他做了一件很獨特的事.
就收藏了一百個真正耶和華的先知.
你明白嗎.
皇帝手下他悄悄收藏了一百個先知.
還有這個米該亞先知.
金庸.
如果我不說你不知道有個米該亞金庸.
但我們就知道有個以利亞.
其實這一班人.
在那段時間都是很重要的人物.
耶洗別那450個巴力先知肯定不是上帝的先知.
基娜娜的兒子西底加帶領四百個先知.
名義上是披著上帝的名.
但他從來沒有聽到上帝的話.
只是討好以色列王阿哈.
阿哈說要什麼就什麼.
當然不是真先知.
在熱王紀上哪些是真先知呢.
阿哈的家宰俄巴底收藏了一百個先知.

$^{601}$米該亞.
以利亞.
所以讀到這裡的時候.
我就學到一個功課.
我如何從聖經裡面.
去學習分辨哪些是真先知哪些是假先知.
你看到這個故事你學到一些東西了.
阿哈的家宰俄巴底收容了一百個耶和華的先知.
這一班是沒有妥協的先知.
這一百個先知未必有米該亞或者以利亞這樣的勇毅.
其實聖經沒有記載他們的事跡.
但是他們都是不願意同流合污的真先知.
你明白嗎.
你看到聖經的故事給我們很多東西學.
米該亞先知是敢於面對四百個披著上帝的名.
但是沒有聽到上帝的聲音.
來講話的假先知.
就是他沒有聽到上帝講話.
不過自己去講話.
以利亞是敢於面對四百五十個巴力先知的人.
這個時代的牧者和信徒領袖.
面對著很多內憂外患.
世俗社會的價值.
是影響著我們對基督信仰和教會的看法.
這個挑戰是非常之大的.
這段時間我們經歷的事.
一個世俗的人可以掛著基督徒的名.
用屬靈的消費者的態度去參加教會的活動.
我們都見到的.
很多的教會都是這樣的.
面對著外面對教會的威脅.
教會的內部是不是有足夠的屬靈的力量去抵抗呢.
這幾年就看得很清楚了.
教會是不是真的裡面有足夠的屬靈力量呢.
有些信徒對聖經是一知半解的.
他可能拿一兩節經文就以篇蓋全.
加上按著社會潮流文化的想法.
他自己帶進聖經裡面.
他自己去解釋就可以了.
有些信徒在網絡世界吸收了很多關於聖經的知識.

$^{641}$就覺得自己都進步了.
我又反問自己.
是不是教會真的做得太少.
做得不足夠呢.
我也會問這些問題.
由於這樣去讀聖經法.
我是有些感觸的.
作為牧者 作為神學院的老師.
我也覺得有責任去撥亂反正.
在眾說紛紜的當中提出正確的聖經解釋.
不過我會集中在神學院裡面訓練那些未來的牧者.
變成在教會就比較少這樣講的.
不過目前的需要是很大.
在變幻的時代裡面.
其實大家越來越渴望聖經的答案.
究竟是怎樣的才行.
這樣又可以 這樣又可以.
哪一種才是.
我一定要很嚴正的指出.
有些人一生去研究聖經 教會歷史 神學.
他都不敢說自己擁有所有的答案.
他不敢說自己擁有所有真理.
他只是說我越來越接近真理.
我自己也會這樣看.
我越來越接近.
但是越來越接近的時候.
那個矛盾就很大.
不知道大家可不可以明白我以下所講的.
你越接近真理 你越看到光明.
那就是好像什麼呢.
好像一塊鏡子那樣看到自己.
你越看到真理.
你是越相反相對的.
就看到自己的那種真我.
我的真面目.
那個就是很矛盾的 你明不明.
一個基督徒最矛盾的地方就是.
我越來越明白上帝.
但我越來越明白自己不像上帝.
你明不明.

$^{681}$這種的落差是很痛苦的.
不知無惹的 你明不明.
所謂不知由自可.
我不需要太過介意.
我知道了我才知道自己的問題.
如果我不知道自己有問題.
我不知道多快樂 是不是.
所以這個才是我們很多基督徒.
在這個過程之中.
我要怎樣去面對這件事.
我又想靈靈更新.
但我又不想這麼煩惱 你明不明.
我知道了我自己不配.
我很煩惱 你明不明我的意思.
我很難面對自己.
這個才是我們最難的事.
所以我們在這裡希望彼此互勉.
我們一定經歷這個歷程.
看到自己的不足都不再講.
我們除了看到自己的不足.
我們還看到很多人的不足.
你明不明.
一定是這樣.
所以我自己都在想.
很多時候義端見到教會有很多迷洋.
所以他們千方百計奪取了牧者的話語權.
去分化教會和分裂教會.
所以你會見到近年.
其實這個現象是很普遍的.
當然牧者和信徒領袖.
都有可能迷失方向.
亦都有可能是牧者和信徒領袖.
造成對教會的傷害.
這個亦都有.
有人有不同的性格.
有些人是和平型的性格.
即是什麼呢.
好好先生好好女士.
有些人的原則性很強.
是就是不是就不是.

$^{721}$據理力爭.
不會為了和諧.
不講出自己不同的意見.
每種性格都有好處和壞處.
優點缺點各都有.
我們不能夠說好好先生好好女士就是最好.
因為很多事情是影響很多人的.
如果我們知道有問題.
卻不提醒其他人.
其實這個是不負責任的態度和做法.
即是你想深一層就是了.
我們又不能夠說據理力爭的就是最好.
因為有些事情各不相讓.
互相對立矛盾.
衝突分裂就一定出現.
同樣哪種性格是屬靈呢.
好好先生好好女士就屬靈.
據理力爭就不屬靈.
是不是這樣當然不是.
因為舊約聖經的先知都是據理力爭.
誰敢說先知不屬靈呢.
這樣據理力爭就屬靈.
好好先生好好女士就不屬靈.
當然又不是.
因為新約聖經保羅在羅馬書十二章十八節教我們.
若是能行總要盡力與眾人和睦.
我看完聖經的那種幅度的時候.
我發覺高難度不得了.
有人說聖經很難解.
我們怎樣能夠明白呢.
究竟我們應該據理力爭還是凡事包容呢.
這個問題真的非常難答.
當我讀到列王記相的時候.
我內心是有很多感觸的.
因為我是一個凡事忍讓的人.
從小到大家裡就教.
哎呀.
一人小句.
或者是.
哎呀.

$^{761}$不要強出頭.
我們這些傳統文化.
要隱惡揚善.
我的訓練就是這樣.
我一直都以為識事令人.
多一事不如少一事是對的.
在這麼多年事奉裡.
其實剛才我回來的時候.
我都見到.
有些我認識了二十年.
真的不是開玩笑的.
孫大英.
我們相識很久的弟兄姊妹.
我們的性格是完全透明的.
你騙到誰呢.
尤其是各位在社會上.
經歷這麼多人生歲月的時候.
你講一句.
挾起尾巴都知道你做什麼.
是不是這樣說.
所以我自己都很感恩.
我自己真的沒有敵人.
我想熟悉我的弟兄姊妹知道.
在教會裡.
我是沒有跟任何一個牧者.
或者信徒領袖衝突的.
你可以說我不做事.
所以就沒有衝突.
你可以這樣說都可以.
可能我真的不下水.
或者不踩到別人的地雷.
但是我就會做一個眾人的支持者.
你誰做頭.
我都是這麼支持你的.
是不是沒什麼用呢.
我可以跟很多人合作.
什麼人都可以合作.
我覺得這是神的恩典.
我是來幫忙的.
我又不是什麼.

$^{801}$但是當我見到教會有問題的時候.
問題越來越大的時候.
我也會問一個問題.
為什麼會發展到這樣的地步呢.
我都會問的.
不是不問的.
最近我就靈修讀到以賽亞書.
第21章第6節.
主對我如此說.
你去設立守望的.
使祂將所看見的術說.
上帝設立先知.
就要守望上帝的子民.
要將上帝從上帝那裡看見.
聽見的去宣講.
就是打奸佬.
先知做什麼.
危險啊.
最近你有沒有看電視.
或者看新聞.
長洲出現了野豬.
有的.
我沒有看到.
有人拍了照片給我.
我鄰居看到.
還有一段片子很有趣的.
就是野豬游水去長洲.
真的.
不是開玩笑.
游得很快.
不是狗仔式.
是野豬式.
有人說好像青衣那隻很大隻.
你明白嗎.
原來我們要看見就要守望.
我們同事就開玩笑.
因為我們長洲有野狗.
那些野狗有時候.
晚上你知道.
我們有些夜歸的時候.

$^{841}$牠不是一兩隻.
是十隻八隻.
有人開玩笑.
今次野豬來.
一定霸佔了野狗的地盤.
好像很開心.
但你想想晚上回家.
一時可能遇見野豬.
一時可能遇見野狗.
很大件事.
所以昨晚我女兒回來.
我也要留意去接她.
因為我路上來上山的時候.
不知道有什麼.
你會發覺原來守望很重要.
有危險你要打更.
你告訴人.
你有沒有見到蛇.
我們長洲蛇.
這陣子很多蛇.
你一定會告訴大家.
哪裡有蛇.
有野狗.
有野豬.
同樣我們都是在上帝的指紋裡面提醒.
有些危險.
我們要警醒.
《爾創亞書》第21章第8節說.
他像獅子嘲叫說.
主啊.
我白日常在守望樓上.
整夜立在我守望所.
這很重要.
即是說什麼.
不單止是牧者.
信徒領袖.
我們中弟弟兄姊妹.
其實都有個祈禱守望的職責.
你要去提醒大家.
有些事不可以輕忽.

$^{881}$你明白嗎.
可能你會覺得說了出來.
人家會有些不開心都不一定.
舉一個簡單的例子.
我最近有預備另一篇講章.
講崇拜.
我就回想以前的年代.
崇拜之前是有祈禱會的.
我不知道為什麼那麼乖.
那時候連崇拜之前的祈禱會都參加.
因為那時候認為.
我們祈禱就會支持到崇拜.
屬靈的力量氣氛就會強一些.
也有人說.
教會的人回到祈禱會的比例.
是一個探業針.
你沒有人去祈禱.
教會就沒有屬靈的力量.
因此我就看了一下.
教會的祈禱.
不是那些奸教會.
總之有些教會.
祈禱會好像都越來越少人.
我就會想怪不得.
你做足了所有的程序.
但你沒有屬靈的力量.
好像一會兒你們唱那個.
願那靈火復興我.
那首歌其實都是我的意思.
所以變成你要有人提出來.
好像祈禱會沒有什麼人來參加.
你要有人說才行.
沒有人說就當沒有事發生.
站在守望樓的意思.
你是一方面要觀察外面的世界的變化.
另一方面你要回到聖經裡面.
你去領受上帝的說話.
這個就是時代的先知.
很簡單的道理.
舊約也是這樣.

$^{921}$新約也是這樣.
現在都是這樣.
我們有沒有這個心智呢.
我們要有心理準備.
任何時代都有假先知.
他們的人數勢力可能很大.
大家說一些相安無事的東西.
讓你快快樂樂.
變成沒有事的.
但這個是不是上帝的聲音.
是不是上帝的遺像.
這個是我們要去問的問題.
很多時候我們是可以有寬容.
但是在真理的事上.
其實是不能夠妥協的.
大家聽完我說這番話.
都覺得很不容易拿捏.
在很多事上我們是有寬容的.
但在真理的事上.
是不能夠妥協的.
先知被呼召.
就是為真理去爭辯.
只是為真理.
就不是為義氣.
你明不明白.
在這裡我最後要說一件事.
其實我們任何一個人.
我們想的事.
做的事.
都是希望能夠造福教會.
我們都是要這樣想.
用這個想法去指導我們的思想和行動.
無論我們是什麼職位.
無論我們是什麼崗位.
我們都沒有人願意.
看到任何一個教會.
是走下坡的.
我們沒有人願意這樣的.
我們的心願就是每一間教會.
它都是充滿屬靈的力量的.

$^{961}$大家都是這樣想的.
所以我們當知道有任何一個教會.
任何一個茶經班.
任何一個團契.
它裡面有些問題.
或者有些難處.
有頂之會有失望.
或者是各種表現的時候.
其實我們.
你想一下我們都不會開心的.
另一個更不開心的.
我們可以想像一下.
主耶穌基督是更加不開心.
等於我們剛才唱.
領我到足留地的那首歌一樣.
同樣的道理.
所以在這裡我想.
最後都是彼此互勉.
希望藉著這段經文.
可以互相勉勵.
究竟我們是不是預備自己.
只是為真道去爭辯.
不只是要有這個心智.
還要天天讀經.
剛才你聽到我說.
我聽同學報告.
我就有些領受.
再聽人講道又收到一些屬靈的訊息.
不是這樣哪有這麼多東西說.
神是透過你每天的靈修.
祈禱.
供應一些屬靈的糧食給我們.
你才喜歡聽講道.
熱愛讀屬靈書籍.
或者是讀神學去裝備自己.
其實這個是我們可以互勉的.
神是很願意將活潑的屬靈生命的糧食給我們.
但是我們要放時間才可以.
而我自己很感恩.
因為你知道我們經常有一句話.

$^{1001}$做那行厭那行.
其實如果我們經常講道.
可能講道是一個厭惡性動作.
又要講了.
都不一定的.
但如果我們講了幾十年之後.
都好像很有興趣.
又要再找找找.
其實很感恩很難得.
譬如說你是專家做餅做了幾十年.
你要有新的一個動力.
令你再創高峰才可以.
如果不是你可能又要做了.
不過沒辦法了.
開店.
這店又只有我一個人做.
你明不明白那種心態是很不同的.
但如果我們天天讀經.
天天靈修.
熱愛聽道.
熱愛讀屬靈書籍.
經常覺得神跟你講話.
你會很開心.
究竟我們怎樣辨別是真假才知道.
哪些才是不隨從世俗的價值觀.
講他的說話是沒有個人的利益的.
不是為了自己的知名度.
他是願意成為上帝流通的管子.
將上帝的話完完整整的宣講.
那些才是真才知道.
最近我還學到另一樣東西.
就是說我們願意做一個很平凡的人.
做很平凡的事.
在弟兄姊妹中間.
真真正正的關心人.
你明白我的意思嗎.
這是很重要的事.
所以我也是學.
其實我不懂得關心人.
我這樣講給你聽.

$^{1041}$但最低限度你懂得為人代禱.
代禱.
你知道他有難處.
你為他代禱.
那你才不只是做一些動作.
不是做一些事.
你真的對人有這樣的關切.
你好像耶穌那樣關心人.
面對變幻的世代.
我們是承擔著很重的責任.
那你願不願意成為上帝指引的守望者.
你不要想著只有牧者才是守望者.
每一個弟兄姊妹他都可以是一個守望者.
他都可以是其中一個.
留心一下外面有什麼變化的人.
我們願神賜福各位.
我們一起祈禱.
生日天父上帝我們向你禱告.
因為我們實在要聽到你的聲音.
願你親自向我們說話.
讓你的道,你的經文成為我們的安慰.
成為我們的指引,成為我們的力量.
讓我們能夠明白我是怎樣配合你.
怎樣在一個複雜,充滿挑戰.
甚至充滿混亂的環境.
我們能夠跟隨你.
我們的主,我們的神.
我們知道人性有很多軟弱.
但我們很希望.
我們的軟弱不會妨礙你的工作.
我們能夠將我們很微小的那部分呈獻給你.
就讓你不介意去使用.
你可以將我們這麼有限.
你可以變得更加成為多人的祝福.
願你紀念我們眾弟兄姊妹各樣的需要.
我們的禱告是奉耶穌基督的名.
阿們.
\newpage



\section{馬太福音 2:1-12}
\label{sec:yEoqVyieILA}
\textbf{2021.12.19《來敬拜我們的王》 馬太福音 2\_1-12 (和修版) 曾敏姑娘}
\newline
\newline
連結: \href{https://youtube.com/watch?v=yEoqVyieILA}{\texttt{https://youtube.com/watch?v=yEoqVyieILA}} ~~~~ 語音日期: 2021-12-19
\newline
\newline
\hyperref[sec:ds7i5doUxn0]{\small{< < < PREV SERMON < < <}}
~
\hyperref[sec:index_chronic]{\small{[返順時目]}}
~
\hyperref[sec:index_scriptual]{\small{[返順卷目]}}
~
\hyperref[sec:_NykIK_k97E]{\small{> > > NEXT SERMON > > >}}
\newline
\newline
馬太福音 2:1-12
\newline
\begin{longtable}{cl}
\hline
\hline
章節 & 經文 (和合本修訂版)\\
\hline
2:1 & \begin{tabularx}{0.7\textwidth}{X} 在希律作王的時候，耶穌生在猶太的伯利恆。有幾個博學之士從東方來到耶路撒冷，說： \end{tabularx} \\ \\ \relax
2:2 & \begin{tabularx}{0.7\textwidth}{X} 「那生下來作猶太人之王的在哪裡？我們在東方看見他的星，特來拜他。」 \end{tabularx} \\ \\ \relax
2:3 & \begin{tabularx}{0.7\textwidth}{X} 希律王聽見了，就心裡不安；耶路撒冷全城的人也都不安。 \end{tabularx} \\ \\ \relax
2:4 & \begin{tabularx}{0.7\textwidth}{X} 他就召集了祭司長和民間的文士，問他們：「基督該生在哪裡？」 \end{tabularx} \\ \\ \relax
2:5 & \begin{tabularx}{0.7\textwidth}{X} 他們說：「在猶太的伯利恆。因為有先知記著： \end{tabularx} \\ \\ \relax
2:6 & \begin{tabularx}{0.7\textwidth}{X} 『猶大地的伯利恆啊，你在猶大諸城中並不是最小的；因為將來有一位統治者要從你那裡出來，牧養我以色列民。』」 \end{tabularx} \\ \\ \relax
2:7 & \begin{tabularx}{0.7\textwidth}{X} 於是，希律暗地裡召了博學之士來，查問那星是甚麼時候出現的， \end{tabularx} \\ \\ \relax
2:8 & \begin{tabularx}{0.7\textwidth}{X} 就派他們往伯利恆去，說：「你們去仔細尋訪那小孩子，找到了就來報信，我也好去拜他。」 \end{tabularx} \\ \\ \relax
2:9 & \begin{tabularx}{0.7\textwidth}{X} 他們聽了王的話就去了。忽然，在東方所看到的那顆星在前面引領他們，一直行到小孩子所在地方的上方就停住了。 \end{tabularx} \\ \\ \relax
2:10 & \begin{tabularx}{0.7\textwidth}{X} 他們看見那星，就非常歡喜； \end{tabularx} \\ \\ \relax
2:11 & \begin{tabularx}{0.7\textwidth}{X} 進了房子，看見小孩子和他母親馬利亞，就俯伏拜那小孩子，揭開寶盒，拿出黃金、乳香、沒藥，作為禮物獻給他。 \end{tabularx} \\ \\ \relax
2:12 & \begin{tabularx}{0.7\textwidth}{X} 因為在夢中得到主的指示，不要回去見希律，他們就從別的路回自己的家鄉去了。 \end{tabularx} \\ \\
[1ex]
\hline
\hline
\end{longtable}
$^{1}$主持人:好,弟兄姊妹平安.
這也是神給我們的福氣.
我們開始的時候再有一個禱告.
不厭棄我們的卑微.
讓我們來到你的面前去敬拜讚美你.
我們原是不配.
這一刻我們知道我們的生命在你的面前是無所遁形.
求你寬恕我們一切的罪惡.
我們嘴巴的罪惡.
心靈的,行事為人的.
天父全然的求你寬恕和潔淨.
讓我們守潔心清的來到你的面前.
今天願意聖靈在我們當中去工作.
對我們每一個說話.
讓我們聽到上帝你的聲音.
我們就在裡面作出回應.
奉耶穌基督的名字祈禱.
阿們.
下個星期五就是平安夜.
所以今天崇拜裡面有不同的詩歌.
都令我們想起聖誕節.
當然星期六更加就是聖誕節的當日.
你們知不知道弟兄姊妹最期望的是什麼呢?.
你最期望什麼呢?.
是參與聖誕派對.
收禮物,吃大餐.
還是思想耶穌基督的降生.
對於我們來說有什麼意義呢?.
今天最想問的是.
在我們的心中我們是怎樣看待耶穌的呢?.
你說「還用問的,回教會那麼久了,這些問題還問?」.
不是的,就是回教會久了才更加要問.
我們和耶穌的關係又如何呢?.
今天很願意透過這段經文再次思考.
我們和耶穌的關係.
我們怎樣看待耶穌呢?.
就決定了我們會怎樣對待耶穌.
我再說一次這一句.
我們怎樣看待耶穌.
就會決定我們怎樣對待耶穌.

$^{41}$在《馬太福音》第二章1到12節記載了三類人.
他們對耶穌的看法是完全不同的.
導致他們對耶穌有不同的態度.
你試試數一下哪三類人.
第一類是哪一類呢?.
剛才大家都有讀經文的.
第一類是哪一類呢?.
請姐妹們試試回答.
大聲一點.
博學之士.
博學之士,真是好了.
多謝孫生.
第一類就是博學之士.
這一類呢,博學之士呢.
他們是尊耶穌為王.
所以他們是會敬拜耶穌的.
博學之士尊耶穌為王.
所以他們敬拜耶穌.
在二章一到二節這樣說.
在希律作王的時候.
耶穌生在猶太的伯利恆.
有幾個博學之士從東方來到耶路撒冷說.
那生下來作猶太人之王的在哪裡?.
我們在東方看見他的星.
特來拜他.
經文提到耶穌出世之後.
有幾個博學之士從東方來到耶路撒冷.
傳統的主要學說.
有三個博士從東方而來.
其實經文並沒有說三.
不是三個.
只是博士將三份禮物.
獻給小孩子耶穌.
所以人們才覺得是三個博士.
實際這裡說的是幾個.
是幾個博士.
博學之士是說那些對於解夢,占星和魔術.
感興趣的人.
所以他們不是上帝的子民.
他們是外邦人.

$^{81}$他們絕對不是像我們今天這樣.
坐在教會裡已經是常常敬拜上帝的人.
不是.
但對於這一類的博學之士.
我們可以這樣看.
不是所有人都一樣.
有些博學之士.
他們是很認真的.
很想了解真相是什麼.
但也有很多人.
只是想.
呃呃拐拐.
做流氓.
找一碗飯吃.
馬太福音所提到的博學之士.
不是這種呃呃哄哄.
想找一碗飯吃的人.
反而這些人.
他們可能曾經研究過舊約聖經和各類猶太書籍.
所以他們很清楚關於猶太人王的事情.
他們在東方看到那顆星.
那顆星指向猶太人的王.
於是他們特別走到猶太這個地方.
去找這個小孩子.
當然.
其實這群人只知道方向.
大概的位置.
他們不知道實際上那個小孩在哪裡.
他們一想.
這個小孩應該在最重要的城市耶路撒冷.
所以他們先到了這個地方.
當他們到了耶路撒冷時.
就問人.
生下來作猶太人之王的在哪裡呢?.
生下來作王和生下來成為王.
有很大分別.
我再說一次.
生下來作王.
和成為王.
有很大分別.

$^{121}$我生下來作王的意思是.
由他出生那一天.
他就是一個君王.
他一開始就有這個身份.
有這個位置.
他不是.
他君王的身份.
不是後面的人加給他的.
《馬太福音》第一章就說到耶穌基督的家譜.
從肉身上來說.
耶穌是大衛的後人.
大衛大家都知道.
是以色列的第二任君王.
現在說的是大衛王後面的家譜一直下.
耶穌基督在那裡而出.
就是說他是君王的後裔.
但耶穌基督不單是地面上君王的後裔.
他更加是神國度裡的君王.
他是萬王之王,萬主之主.
如果大家有機會去看啟示錄的時候.
這是更加明顯,更加突出地說出來.
當然我想說博士們還沒有到這個位置.
看得這麼深.
真的在想永恆上帝.
他的國度裡的君王.
並不是的.
他停留在地面上.
這個君王的層面.
但在博學之士的心中.
這個小孩子的地位非常尊貴,崇高.
他出生的那天就是君王.
他這個身份不是後人加添給他的.
他天生就是這樣.
從一開始就是這樣.
是超越其他人的.
他很配得人去尊崇.
因為博學之士視耶穌為一個這麼重要.
這麼偉大的君王.
所以他們對耶穌的態度有以下三點.
第一是有渴慕之心.

$^{161}$博學之士他們很渴慕見到這個新生王.
所以他們從很遙遠的東方來到耶路撒冷.
他們願意為了見耶穌.
付上時間,金錢,精神和體力的代價.
由東方來到耶路撒冷.
不是很短的時間.
一定要錢.
路上的路費都要.
食物,住宿等等.
一定要精神,長途不捨.
就算騎著生畜過來也好.
也要體力,是要付代價.
當希律王派他們去伯利恆找小孩子時.
他們就去了.
第九節這樣說.
忽然在東方所看到的那火星.
在前面引領他們一直走到.
小孩子所在地方的上方停住了.
他們看見那火星就非常歡喜.
那顆星星一直引領著博士.
去到小孩子的區域.
這件事就是這件事.
令到博學之士很歡喜.
為什麼歡喜呢?.
因為他們很想見到這個小朋友.
很渴望見到他.
有一個很渴慕的心.
第二件事.
這些博學之士會獻上禮物.
十一節這樣說.
「遵了房子,看見小孩子和他母親瑪利亞.
就俯伏拜那小孩子.
揭開寶盒,拿出黃金,乳香,抹藥.
作為禮物獻給他」.
當博學之士見到耶穌這個小孩子後.
他們就獻上三樣禮物.
我經常說是耶穌小孩子.
是因為到這個時候.
耶穌已經不是一個嬰孩.
為什麼我知道呢?.

$^{201}$因為後面當希律下命令.
要殺死小朋友的時候.
不是說初生嬰兒.
是兩歲以下.
其實這些已經是小朋友.
博學之士見到這個小孩子.
獻上三樣禮物.
其中一樣是黃金.
在當時是非常貴重的.
其實今天黃金的價格也很貴.
當時更加不用說.
一般人不容易擁有的.
乳香是什麼呢?.
有沒有人知道?.
有沒有人知道乳香是什麼呢?.
樹什麼呢?.
樹枝,是吧?.
樹膠.
樹膠是會發光的.
還有香味的.
而末藥就是香水和香料.
用在屍體上有防腐的作用.
這三樣東西的特色就是貴.
很貴重.
一般人來說很少送這三樣禮物.
除了黃室成員.
我相信比較普遍會有這三樣東西.
一般人是少的.
但是博學之士覺得.
這個在他們心目中是新生王.
是很值得他們送上這麼名貴的禮物.
所以他們就送上這份禮物.
我深信其他人會覺得奇怪.
為什麼送禮物給這個很普通的小孩子.
值不值得呢?.
博學之士覺得值得.
他們有學務之心.
他們會送上禮物.
第三樣東西.
敬拜耶穌.

$^{241}$11節.
「進了房子看見小孩子和他母親瑪利亞.
就俯伏拜那小孩子」.
幾位博學之士.
看到小孩子和他的母親.
就敬拜他.
敬拜小孩子.
獻上最大的敬意.
耶穌基督是佩得人的敬拜和讚美.
事實上,他不單是地面上皇室的後人.
更加是整個神的國度裡的萬王之王,萬主之主.
是佩得人的敬拜和稱頌.
我們在這裡會佩服幾位博士.
佩服他們什麼呢?.
我一開始是怎樣形容他們的?.
他們不是那些經常坐在教會上敬拜上帝的人.
是不是?.
他們不是神的百姓.
他們是外邦人.
那些占星的.
喜歡魔術的.
是這一類的人.
但是他們對耶穌卻有一份很深的敬意.
說真的,他們對耶穌的認識比不起上帝的百姓.
但是那些認識已經足以令他們很想去敬拜耶穌.
他們的反應是很特別的.
所以在他們心中如何去看待王的身份呢?.
不簡單.
這是第一種人.
第二種人是什麼人呢?.
大家又說說.
希律,沒錯了.
希律王是第二種很典型的人.
他視耶穌為威脅.
所以想盡辦法剷除耶穌.
耶穌是在希律作王的時候出生的.
我曾經在講道的時候也有介紹過希律.
這個希律是大希律王.
本來他是以土買人安提柏特的兒子.
是一個非常出色的人.

$^{281}$在主前四十年的時候.
羅馬元老院任命他為猶太王.
所以這是後天加上去的.
他絕對不是生來就是王的人.
由於他很聰明.
歷屆羅馬的君王也很喜歡他.
但是大希律這個人的特色是他很喜歡權力.
他會加重在百姓的身上.
百姓也覺得他.
你這個人也不是名正言順坐那個位子的.
你也是透過篡位才得到權位.
所以百姓也看不起他.
所以對這一點他感到很憤怒.
但是他很喜歡他的權位.
一旦得到就要抱著他.
到了晚年的時候.
他的行為變得很殘忍.
因為憤怒因為妒忌的緣故.
殺死了他的親信.
他的老婆和至少兩個兒子.
你可以想像這是一個怎樣的人.
一個非常殘忍的人.
經文說.
希律王聽見了就心裡不安.
耶路撒冷全城的人也不安.
希律王不安是因為他很擔心.
這個傳聞中.
人們每個人都說.
生來就是猶太人王.
這個人會影響他.
他的出現.
是否代表自己的王位不保.
他來就要趕走我.
我就不能做猶太人的王了.
這個很實際.
是自己的自身利益受到影響.
他很明白.
正如剛才所說.
他自己作王是後天的事.
人家加給他.

$^{321}$不是天生就擁有這個權柄.
所以耶穌基督的出生.
令他感到很不舒服.
希律不安.
全城的人都感到不安.
這些人感到不安.
不是因為這個英雄的出現.
影響我們所有人的利益.
這些人不安.
是他們知道.
希律不安.
我們就沒有好日子過了.
你怎知道希律王又做了什麼不好的事.
我們就民不聊生.
所以百姓都不安.
由於希律覺得耶穌基督是一個威脅.
所以他對耶穌的態度.
又有以下三點.
第一點.
起了殺機.
他動了殺機.
他召集了祭司長和民間的文士問他們.
基督該生在哪裡.
他們說在猶太的伯利恆.
希律想知道那個新生王.
是在哪裡誕生的呢.
做什麼呢.
去追尋他.
當他知道原來這個新生王.
是在伯利恆出生的.
他就去照這個博學之事來問他.
查問那個星是什麼時候出現的.
希律想知道.
這個新生王大概現在多大了.
什麼時候出現了.
什麼情況.
如果要找人去找他.
都要有一個譜和範圍.
不然怎麼找呢.
多大的人.

$^{361}$他要鎖定一個範圍.
希律做這麼多事.
都是因為他心裡起了一個殺機.
在他心裡沒有半點對耶穌的渴慕和愛.
第二點.
虛偽一番.
第一先起了殺機.
第二他又虛偽了一下.
希律派了博學之事去伯利恆.
對他們說.
你仔細去尋訪一下那小孩子.
找到了就來報信.
我也很去拜他.
說得很好聽.
其實我和你們一樣.
我也很想敬拜這個新生王.
我也很有這份心.
其實這只是一個謊言.
說這個謊言.
使博學之事相信他.
沒有了戒心.
日後他們得到資訊.
就會願意和他分享.
告訴他這個新生王在哪裡.
這是第二點.
虛偽一番.
第三點是格殺勿論.
博學之事在十二節記載.
他們因為在夢中得到主的指示.
不要回去見希律.
他們就從其他路回自己的家鄉去了.
十六節.
希律見自己被博學之事愚弄.
極其憤怒.
警察將伯利恆城裡和四景所有的男孩.
根據他向博學之事仔細查問到的時間.
凡兩歲以內的都殺盡了.
希律見自己沒辦法透過博學之事.
知道耶穌的所在.
就唯有殺死所有年紀相若的小孩子.

$^{401}$確保這個新生王沒辦法生存下去.
他希望殺害耶穌.
這個想法去到這一刻.
是成為了事實.
不是再想了.
他轉化成為一個行動.
他要做出來.
這個就是希律王.
這個就是想剷除耶穌的這一代人.
這一批人.
我們聽到這裡心裡會不會想.
這一類人和我們最遠.
我們怎會想著作王呢?.
我們不說這些.
我們敬拜上帝.
想著作王嗎?.
是希律這些人才想作王.
他會不會想這些對付耶穌呢?.
其實也不是的.
我們人的天性.
我們自己很想作王.
我們作王不一定像古代那樣.
真的有個王國.
在裡面我是最高的領導人.
我都可以在我所在的地方作王.
我們不想耶穌作主也會.
耶穌在這裡作主.
我們每次都要求求祂.
我作王就很順暢.
我想怎樣就怎樣.
所有利益都是屬於我的.
如果現在每次都要求求耶穌.
影響我的利益.
我也很不喜歡.
我們人會不會這樣想呢?.
會的.
教會裡面的人都可以這樣想.
所以這一類人不是離我們很遠的.
這是第二類.
第三類是哪一種人呢?.

$^{441}$你說是嗎?有第三種人嗎?.
哪一類呢?.
看來看去都不知道看到哪一類.
哪一類是第三類人呢?.
這是最不明顯的.
其實在經文裡面提到一種.
是宗教領袖.
宗教領袖對於新生王誕生的事情.
感到非常冷漠.
他們都不打算去尋找耶穌.
宗教領袖暫時提到的.
只是提到祭司長和民間的民事.
這些人很熟悉神的律法.
很熟悉神的話語.
他們知道舊約聖經不止一次提到大衛的後裔.
以色列最終極的君王一定會出現.
他要統管以色列民.
統管萬邦.
他要拯救他的百姓.
這些宗教領袖絕對是很熟悉.
所以當希律問他們的時候.
基督該生在哪裡呢?.
他們說很容易知道.
在猶太的伯利恆.
因為有先知記著.
猶大地的伯利恆.
你在猶大諸城中並不是最少的.
因為將來有一位統治者要從你那裡出來.
募養我以色列民.
這個整個第六節其實是來自.
小先知書尼加書.
小先知書尼加書五章二節那裡.
記載了一段相當接近的文字.
這班宗教領袖絕對熟悉聖經.
亦很懂得用聖經回答希律王的話.
但說真的.
聖經只是停在那裡.
沒有記載這班人如何渴望耶穌.
他們亦沒有興趣認識這個.
在他們心目中所謂的新生王.

$^{481}$這是否真的要來拯救我們以色列民.
他們有很多懷疑.
這班人很冷漠.
為何會這麼冷漠呢?.
這班認識聖經,譯聲神話,語語人.
他們是教人的.
可以教人聖經.
為何會這麼冷漠呢?.
其實我們值得問一個問題.
這班人的上帝的話.
是否真的能進入他們的生命裡呢?.
曾經有一個人這樣說.
他說你不要以為.
一定是那些神學院的教授.
最懂得勝經.
可能隨時來一個比賽.
誰能解釋聖經更好呢?.
可以是一個不信耶穌的人.
一個學者.
相反,他上台時.
這堂課可能會覺得眼前一亮.
這段經文是可以這樣解釋的嗎?.
很精彩,真的會拍爛手掌.
可能會的.
他不信耶穌也可以說得很精彩.
那你說聖經對他的意義是甚麼呢?.
一個學術.
他解釋得很厲害的一段文字.
他用他的學術去解釋這段文字.
教書地教你.
但這段文字對他來說沒有意義.
這班宗教領袖也可以這樣.
很熟悉聖經.
我知道上帝將來會這樣安排.
等等等等.
不過這些有沒有意義呢?.
可以沒有意義的.
所以就很冷漠.
這才是最大的諷刺.
不是真正認識上帝的人.

$^{521}$或認識神不是這麼深的人.
渴慕耶穌.
反而這些所謂的.
說自己認識神的人.
卻是不渴望耶穌的.
這段經文是充滿了諷刺.
所以我想提醒大家一件事.
今天不是說.
我上了很多主日學.
我聽了很多唐道.
我自己每天看聖經.
從早到晚我已經很熟悉聖經.
我就很熟悉了.
不是.
我想說.
你有沒有讓上帝的話.
真的進入你裡面.
這是否真正建立你和神的關係.
你真的明白神?.
你真的將你的心給神?.
如果我讀完這麼久.
學完這麼久.
我都不明白祂.
我也不打算體貼祂的心意.
那麼讀和學的意義在哪裡?.
到最後原來我還是很冷漠.
那麼這個生命.
究竟價值在哪裡?.
是否很可惜呢?.
這種生命.
我每年聖誕節都慶祝.
我參加所有聚會.
什麼事情都要做.
只有驅魔.
你的生命裡面沒有真正改變過.
我仍然問這個問題的意義在哪裡?.
到最後我發現.
我其實不認識神.
那麼你看這麼多聖經.
這些是為了什麼?.

$^{561}$我很喜歡人看聖經.
我很喜歡教人看聖經.
但是我更加喜歡.
你看完之後.
你和上帝有一個真實的關係.
是真實的關係.
不是虛假的關係.
很值得大家思考.
最後應用的時候.
值得我們去想一想.
我們究竟屬於以上哪一種人呢?.
我由第三種開始說起.
剛才說到冷漠的這種人.
我們都覺得自己很認識聖經.
很多參與教會聚會等等.
我們是否認識神?.
和神有一個好的關係?.
有沒有問過自己.
為什麼我們的心這麼冷漠?.
為什麼?.
我們有沒有靠著聖靈.
活出神的話語呢?.
其實我想每一次.
你聽完道之後.
開禮堂的時候.
假如你能夠捉到一句神的話.
然後在那個星期活出來.
我覺得已經是非常成功了.
就不是下個星期我問你.
上個星期我說什麼呢?.
我已經忘記了.
不用下個星期.
可能一會兒吃飯時問.
其實我都不太記得.
你剛才說了什麼.
是不是?.
那一刻就覺得聽了一些東西.
之後就不記得了.
那是沒有意思的.
所以在這裡值得大家思考.

$^{601}$我們是不是這一種人?.
第二種人.
就是剛才我們說的.
好像希律王這一種.
我做王.
你不要影響我做王.
不只是不相信耶穌的人才會這樣.
信耶穌的人都會這樣.
我很想做我生命的王.
我可以說一些很淑女的東西.
但我想做王.
所以到最後你不要干涉我.
你走到一邊.
耶穌.
我們是不是這一種人?.
如果有.
問一下.
為什麼會這樣呢?.
我們什麼時候開始成為耶穌的敵人呢?.
和上帝爭這個位置.
其實都很容易聽得出來.
我們的生命是不是讓耶穌做王.
你說一番話.
很快就能看見.
那個人是在說自己.
自己想怎樣怎樣.
還是真的讓上帝做王.
很快就聽得出來.
這是第二種.
還是我們是屬於.
不是的.
我告訴你.
我們是屬於敬拜耶穌那種.
我們學武耶穌的.
敬拜耶穌.
那我就要問一個問題.
在你心目中.
你覺得耶穌是怎樣的呢?.
耶穌是不是王呢?.
耶穌真的是王?.

$^{641}$是哪一種王呢?.
我發覺.
我們看耶穌的觀念是很不一樣的.
有人看耶穌是阿拉丁神燈.
我擦擦擦擦擦擦擦擦擦.
那個阿拉丁神燈.
叫你出來.
我有幾個願望.
我想這樣想這樣.
那個神很信服.
好,我為你達成.
你真是好.
下次有需要.
又擦擦擦擦擦擦.
他又出來成就我們心願.
真是好.
只要有一次.
擦那個燈.
你怎麼沒有出來.
沒有達成我心願.
這個是不是王呢?.
這是我們的僕人.
不是王.
上帝聽你的話.
什麼時候才知道.
你是不是視他為王呢?.
就是當我有不信你的時候.
我是怎樣看神就知道.
我那個時候最容易.
因為不信你.
人生很痛苦.
我就很想將矛頭指向上帝.
你令我很痛苦.
這樣那樣.
你欠我的.
你欠我太多.
你那個時候就知道.
原來祂不是王.
祂只不過是我的僕人.
祂真是服侍得我不周到.

$^{681}$祂可能是王大仙的神.
我們看祂是.
我求籤你就答應我.
為什麼我求籤你不答應我呢?.
我們是不是看祂是這樣呢?.
是王是怎樣的?.
你好尊崇祂.
不是祂服侍你.
是你服侍祂.
是我們服侍祂.
我們一起服侍祂.
這樣就叫做王.
這樣就祂說了算.
就不是我們說了算.
是不是?.
今天的題目是.
來敬拜我們的王.
聖誕節更是一個好的時間.
思想耶穌基督.
祂是不是王才行呢?.
在我們心中.
如果我已經很久.
坐在我自己生命的寶座上.
已經把耶穌趕到不知去哪裡呢?.
這個聖誕節.
是不是應該把自己這個位置.
讓出來?.
不是要祂服侍我們.
而是我們去服侍祂.
去敬拜祂.
博士會獻上很珍貴的禮物.
黃金,玉香,末藥.
是很珍貴的.
我們獻給神是很珍貴的.
不是將我們不要的.
有剩的就給你吧.
上帝說很好的.
我給你.
有很多人不給你.
哇!神啊!.

$^{721}$我好稀罕啊!.
上帝說我擁有天地萬物.
我好稀罕你這些剩下的.
不要的,殘缺的東西.
我們是怎樣對上帝?.
是否將最好的獻上?.
去配得?.
我們有沒有很好的敬拜祂呢?.
有一首聖詩叫《尊主為王》.
是一首很舊的詩歌.
是英國著名的牧師羅勒德的作品.
羅勒德是18世紀英國.
屬靈復興運動的領袖.
他寫的詩歌都有一個特色.
充滿了聖靈的能力.
特別是這一首《尊主為王》.
我很想有機會找回那些歌譜.
或者那些聲音.
因為比較舊.
試試找回.
這首歌是他的不老之作.
這首歌影響了很多人的生命.
其中有一個小故事是這樣的.
有一個海外的宣教士.
他同樣是很有名的.
叫做斯國脫(Scott).
這個斯國脫牧師有一次.
進入了印度荒山.
一個野蠻的部落傳福音.
當時他隨身帶著一把小提琴去山區.
他很勇敢.
忽然間有一次.
他被一群的土人包圍了.
他們是很原始的.
而且是很野蠻的.
殺人也是等閒之事.
他們就高舉長矛.
準備要殺他.
他心中這個時候.
開始默默祈禱他的上帝.

$^{761}$在這個時候他祈禱他的上帝.
而且很特別.
在一個快要死的場景.
他做了一件什麼事呢?.
就是敬拜上帝.
敬拜神為他的王.
他拿起他的小提琴.
閉上雙眼.
一邊拉一邊唱這首詩歌.
很特別吧?.
我不知道如果我們真的有這樣的處境.
有人將會殺死我們.
那一刻毫無準備.
我們會打算做什麼呢?.
應該也會祈禱吧?.
但可能情緒上太慌亂了.
已經想不到想做什麼.
但他那一刻就是敬拜上帝.
一邊唱這首詩歌.
尊主為王唱了兩次.
他閉上雙眼.
接著不要緊.
他敬拜完上帝.
如果要離世就離世吧.
他準備好了.
但是很特別.
當他張開雙眼的時候.
很特別他看到什麼呢?.
敬拜神之後.
他看到周圍的人.
那些土人放下了他們的長矛.
改變了他們很敵對的態度.
反而接受了施覺脫為他們的朋友.
還接待他為上賓.
這個真的很特別.
在敬拜上帝後.
上帝有他奇妙的工作.
改變人心.
在裡面將整件事.
化為一個很大的祝福.

$^{801}$施覺脫也趁這個機會.
將耶穌基督的福音介紹給他們.
就是這樣.
他一直留在土人當中.
服侍他們.
兩年半之後.
他因為健康的問題沒有辦法.
他要回到美國.
那時候已經帶了很多人信耶穌.
這一班人.
曾幾何時.
都是在那裡拿著長矛殺人的.
很野蠻的這班人.
信了耶穌.
經歷上帝的愛.
這一班人很不捨得施覺脫離開.
他們一直送施覺脫走.
送了多久呢?.
送了40里的路.
很長的路.
一直陪著很不捨得.
但是到了臨離開.
施覺脫自己.
反而很不捨得.
被這一班土人.
深深的感動.
最後留下來繼續服侍他們.
當中這一首詩歌.
尊主為王.
成為了很美的一幕.
我相信這個位置.
不只是那首詩歌.
由這個詩歌提到這個人的生命.
是一個敬拜上帝的生命.
他也將他一生.
獻給神.
上帝帶領他去哪裡.
他就去哪裡.
不是他自己說了算.
上帝說了算.

$^{841}$上帝要這樣用他的生命.
他就這樣獻上給神.
是一切.
這首詩歌.
只不過是他生命的一個寫照.
今天弟兄姊妹.
都很想邀請大家.
是敬拜我們的王.
這個聖誕節.
我們都很想.
邀請我們的王.
這個聖誕節.
再一次看看你生命的寶藏.
是誰坐在這裡.
如果是我們.
自己請我們下來.
讓耶穌基督坐在那裡.
耶穌基督不是聖誕老人.
不是阿拉丁那種神.
他是王.
而且是萬王之王.
在這個聖誕節.
讓我們的生命來一個重整.
要去敬拜他.
把他當得的榮耀歸還給他.
還有將我們的生命的主權.
真的切實交在他手裡.
我們一同禱告.
今天透過這段經文.
透過不同的生命.
你讓我們去思考.
今天我們是什麼的光景.
我們總是會覺得.
聖經裡的人物.
是別人的故事.
其實在裡面.
也有我們自己的影子.
可能我們的生命就是這樣.
我們想自己作王.
我們不想耶穌在這裡管著我們.

$^{881}$又或者我們覺得.
我們的生命這一刻很冷漠.
我們不是實際很關心你的事.
如果是的話.
我求你寬恕我們.
求你寬恕我們.
我們看完聖經.
但我們的心沒有真正向你開放.
天父求你寬恕我們.
求你潔淨我們的生命.
求你幫助我們.
真的謙卑下來.
將你的寶座讓給你.
願你坐在裡面.
坐著為王.
願你興起我們的敬拜.
將榮耀歸還給你.
是可以的.
主啊,不只是那些不信主的人.
他們渴望你.
外邦人渴望你.
昔日我們都是外邦人.
今天我們進到這個門來.
你讓我們繼續渴望你.
而且我們的生命.
要被你搞動.
要被你改變.
主啊,求你得著我們的心.
奉耶穌基督的名字祈禱.
阿們.
\newpage



\section{路加福音 2:8-14}
\label{sec:_NykIK_k97E}
\textbf{2021.12.26《聖誕節的意義》 路加福音 2\_8-14 (和修版) 曾裔貴牧師}
\newline
\newline
連結: \href{https://youtube.com/watch?v=_NykIK-k97E}{\texttt{https://youtube.com/watch?v=\_NykIK-k97E}} ~~~~ 語音日期: 2021-12-31
\newline
\newline
\hyperref[sec:yEoqVyieILA]{\small{< < < PREV SERMON < < <}}
~
\hyperref[sec:index_chronic]{\small{[返順時目]}}
~
\hyperref[sec:index_scriptual]{\small{[返順卷目]}}
~
\hyperref[sec:yWaJULCl_Mc]{\small{> > > NEXT SERMON > > >}}
\newline
\newline
路加福音 2:8-14
\newline
\begin{longtable}{cl}
\hline
\hline
章節 & 經文 (和合本修訂版)\\
\hline
2:8 & \begin{tabularx}{0.7\textwidth}{X} 在伯利恆的野外有牧羊人，夜間值班看守羊群。 \end{tabularx} \\ \\ \relax
2:9 & \begin{tabularx}{0.7\textwidth}{X} 有主的一個使者站在他們旁邊，主的榮光四面照著他們，牧羊人就很懼怕。 \end{tabularx} \\ \\ \relax
2:10 & \begin{tabularx}{0.7\textwidth}{X} 那天使對他們說：「不要懼怕！看哪！因為我報給你們大喜的信息，是關乎萬民的： \end{tabularx} \\ \\ \relax
2:11 & \begin{tabularx}{0.7\textwidth}{X} 因今天在大衛的城裡，為你們生了救主，就是主基督。 \end{tabularx} \\ \\ \relax
2:12 & \begin{tabularx}{0.7\textwidth}{X} 你們要看見一個嬰孩，包著布，臥在馬槽裡，那就是給你們的記號。」 \end{tabularx} \\ \\ \relax
2:13 & \begin{tabularx}{0.7\textwidth}{X} 忽然，有一大隊天兵同那天使讚美神說： \end{tabularx} \\ \\ \relax
2:14 & \begin{tabularx}{0.7\textwidth}{X} 「在至高之處榮耀歸與神！在地上平安歸與他所喜悅的人！」 \end{tabularx} \\ \\
[1ex]
\hline
\hline
\end{longtable}
$^{1}$鼎姐妹主女們平安.
再一次代表教會,母堂,感宣家.
來問候大家聖誕節快樂.
牧師今天能夠在大家當中.
來分享神的真道.
是一件很開心很開心的事.
我們一起有個簡單祈禱.
多謝主.
你是那一位創造的主.
更加降生成為人的主.
我們知道嬰孩在馬槽.
絕對不是一件浪漫的事.
是很坎坷,很狼狽.
還有主,你甘願卑微.
降生成為貧窮.
使我們因為你而得到心靈,生命的付足.
我們獻上感恩.
讓我們每一位的聖徒.
看到我們自己是多麼的蒙恩.
我們的生命在神眼中是多麼的尊貴.
讓我們願意在地上的生活.
去榮耀你的聖名.
我們能夠像年提.
明年的年提.
大家都毀身來事奉.
因為是上主所喜悅的.
求主帶領.
今早早.
當我們再一次去思想聖誕節的意義的時候.
求你讓我們的心靈都能夠明白.
心靈裡面去振奮.
一起去追求.
這樣感恩禱告.
奉耶穌基督的名.
阿們.
剛才這段聖經一點都不陌生.
有一個記號.
那個記號就是嬰孩在馬槽.
馬吃東西的那個盤子.
大家都知道那個背景.

$^{41}$因為當時帝國.
羅馬帝國.
一個皇帝.
這個蓋薩奧古士都.
或者叫做屋大維.
他統治了整個羅馬帝國.
他是第一位的君王.
他要做什麼呢.
全個帝國要做一次.
大規模的人口普查.
所有人要回去自己的原居地.
那就有根有據.
你是什麼戶籍.
為什麼.
因為要收稅.
所以聖經的描述.
約瑟和他太太瑪利亞.
這個他所聘取的妻子.
太太.
當時已經有新人.
當然就不是他和太太.
來結合生的.
是聖靈使到瑪利亞懷孕的.
就去到大衛的城.
伯利恆.
就很多人.
所以那些旅店沒有地方.
很狼狽地.
有一個地方就是馬槽.
我們的主耶穌.
就在馬槽那裡出生.
包著一塊布.
很坎坷.
一點都不浪漫.
但是重點就是天使.
有一個天使.
向在伯利恆野外的牧羊人.
當時是一些勞苦大眾階層.
夜間當值.
夜間.

$^{81}$就是當夜班.
有沒有在座的弟兄姊妹.
你是當夜的.
要通宵等等.
夜班就通常沒什麼人做的.
很難的.
我們有個弟兄.
7-11做通宵.
好像多了20.
時薪多一點點.
但是很辛苦的.
打亂了節奏.
那些牧羊人是分班次的.
當夜班的牧羊人.
就是說可能都沒什麼人做的.
他就去做.
但是很奇妙地.
天使是不是去錯了地方.
是不是那個GPS定位定錯了.
跑去找這些牧羊人.
還是夜班的牧羊人.
大家有沒有留意.
在《馬太福音》第二章第一節.
當希律王的時候.
有幾個博士來到猶太.
去找希律王.
就問希律王.
那個生下來做猶太人的王在哪裡.
我們要去見他.
去拜他.
還有將那些禮物給他.
那些博士不是去錯地方.
既然生下來做猶太人的王.
當然要去皇宮找.
真的去皇宮找希律王.
就問希律王.
但是為什麼.
那個天使是不是去錯了地方.
找錯人了.
為什麼找一些當夜的牧羊人.

$^{121}$這是一個很值得我們去想的問題.
當時牧羊這個階層.
是被社會歧視的.
他們是沒有社會地位的.
根據一些拉比的文獻說.
當時代.
這些牧羊人.
如果要做出庭的法庭的見證人.
法庭是不受理的.
因為他們那個明星.
他們平時不是那麼得到.
在上的統治階層所接納的.
這個就是牧羊的階層.
真的是勞苦大眾.
基層.
但是天使向他顯現.
是不是找錯了地方.
前兩天準備平安夜.
下午吃完午飯.
飯盒沒理由放在辦公室.
自己拿出門口扔.
外面的垃圾桶可以快點清除.
走出門口沒多久.
對面有個太太.
牧師 牧師 牧師.
快點來 快點來.
什麼事.
先扔下飯盒.
快點進來.
進去.
原來他看不到無線電視.
追不到台.
快點叫我幫他追台.
追來追去都追不到.
不知道什麼原因.
然後他說.
你不是哪個議員嗎.
我說不是.
我是教會牧師.
我是曾牧師.

$^{161}$我不認識.
原來他以為我是議員.
但是叫我叫牧師.
叫我快點進去幫他.
我說你找錯人了.
我不認識.
好像很好笑.
但是原來他是找錯人了.
認錯人了.
但是天使沒有認錯.
天使去找牧羊人.
告訴他們一個大喜的訊息.
今天已經是26號.
稱為Boxing Day 拆禮物.
很多人已經慢慢回復正常生活.
當然今天還是假期.
明天還是假期.
但是在主耶穌降生之後.
你想一下牧羊人.
他們完了.
回到自己的崗位.
自己的生活階層.
是不是突然之間.
因為他們可以在野外捕獲了聖嬰.
他們突然間身份提高呢?.
相信不會.
很快連約瑟和瑪利亞.
都會被希律王追殺.
所以有天使提醒這對夫婦.
趕快帶著你的小朋友.
即是耶穌.
逃難去埃及.
逃避希律王的追殺.
你可想而知.
這個降生的消息.
沒有帶來整個耶路撒冷.
猶太的人.
他們賞心.
他們雀躍.
因為為我們生了一個彌賽亞聖嬰.

$^{201}$整個消息沒有帶來整個城的轟動.
帶來的是人心的不安.
因為希律王感到受威脅.
所以要殺兩歲以來的嬰兒.
就是這樣很平鋪直敘.
告訴我們.
聖誕在當時的人的心裡.
如果沒有聖靈.
如果沒有信心.
來調和這個節慶.
只會帶來.
不過是一個普通的小朋友出世.
不過好像有些博士說.
他是猶太人的王.
是他說的.
不是真正可以影響到世界.
各位弟兄姊妹.
你有沒有想過.
什麼是聖誕真正的意義呢?.
世界的人.
可以的就是出國旅遊.
出不了就要出外吃東西.
或者住酒店.
也可以散心.
或者買禮物.
沒有信仰的人.
就會這樣去過了長的假期.
這兩天就發現.
24日平安夜我們有聚會.
有晚堂的平安夜的崇拜.
有很多平時不來教會的人.
來教會.
很多平日不回來的.
都來教會.
25日也是.
很多人平時都不來的.
都來的.
他們感覺到就是.
我來到禮拜堂聽聖詩.
我就很舒服了.

$^{241}$我的心靈就會有一種滿足了.
但是過了之後.
剛才我來之前.
11點鐘.
我招呼弟兄姊妹回來崇拜.
和他們打招呼.
少了很多人.
天氣一來冷.
二來就是這兩天都有聚會.
當然也有直播.
少了很多人.
沒有平時午堂那麼多.
相信弟兄姊妹在在.
都很穩定地來到聚會.
原來我們只是一種節慶的感動.
沒有真正的對這位救世主耶穌.
有一個真正的認識.
而這個認識帶來的.
就是你的信心會有一個提升.
生命有一個改變.
各位.
你想想過了耶穌的出世之後.
首先就是有幾個人.
我們要想想他們的後來是怎樣的.
聖經表面就沒有說太多.
都已經要我們自己去思想.
他們往後的人生應該是怎樣.
牧羊的人.
你猜一下他們的階層會不會突然提升.
我們組織一個公會.
我們牧羊公會.
我們曾經去探訪過聖嬰.
所以我們是特別神聖.
有天使曾經向我們這一班人顯現.
所以我們要爭取這個社會的地位提高.
相信沒有.
繼續是受歧視的老虎階層.
不過在上帝的眼中.
選這一班人成為報喜的第一批人.
第一批知道這個福音.

$^{281}$關乎萬民的福音.
因為今天在大衛的城裡.
為你們生了救主.
就是主基督.
這個聖嬰帶來的就是人類.
在地上因為他而得到生命的平安.
所以牧羊人.
他們當時一定不會有電話.
又拍下照片.
立刻放到社交群組.
Instagram或者是Facebook.
一定不會.
但是他們這一次的經歷.
沒錯.
社會上一般人.
不會對他們特別有一個改變的態度.
但是他們在上主的眼裡.
原來是尊貴.
原來是上主不會對他們有一個歧視.
他們對這個信息.
他們看到這件事的真確的應驗.
真的有一個嬰兒.
包著一塊布.
在馬槽裡.
那個就是神的兒子.
相信對這一班的牧羊人來說.
他們的生命的改變.
不是因為看到一個奇蹟.
是因為他們看到上主對他們的愛和接納.
弟兄姊妹.
我們要盡快在禮拜堂這裡.
不只是透過一些的感動.
一些的擺設.
一些的詩歌.
來有一種的享受.
是盡快要跟神.
跟這一位的耶穌基督.
有你個人的生命的接觸.
那種信心的成長.
以至無論你落在什麼的光景.

$^{321}$困難.
耶穌基督對你的牧羊關愛是不改變的.
在網上最近看到一篇短片的視頻.
很感人的.
有一對夫婦.
丈夫是一個戰鬥機師.
太太是一個醫生.
韓國瑜的醫生.
他們有一個女兒.
在美國的.
五歲.
只是五歲.
兩歲的時候就患了一個絕症.
神經性的一些問題.
那些的退化症.
會慢慢那些器官會退化.
慢慢就會離開世界.
但是一個五歲的小朋友.
兩歲已經開始.
逐個機能一路一路衰退了.
去到差不多.
真的去到那個階段沒有了.
很多的器官都已經損壞了.
媽媽跟女兒說.
女兒叫Julianne.
Julie.
她姓Slo.
雪那個字.
那個姓.
叫史洛.
跟這個女兒說.
就快又要進醫院了.
女兒說我不想再進醫院了.
我想在家裡.
媽媽就很清白的跟她說.
如果你不在醫院.
你很快就要去天堂了.
女兒說.
是的,我知道.
接著.

$^{361}$父母說了一句話.
很令人感動.
爸爸媽媽暫時都不可以跟你一起去天堂.
女兒怎麼說.
一個五歲.
就快要離世的小朋友.
她說爸爸媽媽不要緊.
因為天父在那裡.
我先去.
你們還沒來,我知道.
一個五歲的小朋友.
這個視頻你今晚上去看一看.
很感人的.
這樣的信心.
你想一下你五歲在做什麼.
她五歲要面對生和死.
自己一個離開.
還有她那個器官的損壞.
不是令她的人很難看樣子.
是一個很漂亮的女兒.
五歲就要跟父母說再見.
但是那個信心.
她知道她在天父的手中.
她現在是回到天家.
等待父母將來來.
一點都不覺得孤單.
五歲.
有沒有弟兄姊妹在這裡已經超過五年了.
你的信心有沒有這種程度.
原來不是計算你幾歲.
是計算你跟耶穌的關係.
有沒有那種信心.
所以牧羊人.
他們的生命很大改變.
為什麼會向我們顯現.
我們是勞苦階層.
當野的.
你試想一下.
如果你是做看更的.
你當野的話很辛苦.

$^{401}$沒什麼地位.
但是天使向他們顯現.
還要有一大隊的天君.
為他們唱歌.
至高之處榮耀歸於真神.
地上的平安歸於人.
天君唱的.
天君.
為什麼牧羊人會見到他們有分別.
天使和天君是不是真的武裝有槍.
不知道.
但是他們分辨到.
你可想而知.
為這些牧羊人.
來作見證.
就是說原來神愛每一個階層.
你會不會覺得你自己的階層.
比較低下呢.
或者你覺得你自己的人生.
很普通.
普通到你都覺得自己很普通.
神的眼中.
每一個弟兄姊妹.
都是他所愛.
所寶貴.
所以在路加福音.
你會很容易就看到.
很多的描述.
都是神向這些窮苦人.
來顯現的.
譬如當時那些血流寡負.
十二年的.
耶穌就會向他來施恩.
還有其他的.
各位弟兄姊妹.
聖誕的意義.
如果我們藉著這一段聖經.
還有剛才馬太福音第二章.
那段經文來做一個比對的時候.
我們就會發現.

$^{441}$聖誕的意義不在於.
那個氣氛.
不在於那個佈置.
其實當時怎會有佈置呢.
當時其實是一個很狼狽的場景.
聖誕是告訴我們.
神愛每一個人.
神愛你.
不是因為你的身份.
不是因為你是一個怎樣的人.
神愛你.
因為他認識你和我.
因為他是造物的主.
因為他創造了我們.
我們每一個人在他眼中.
是寶貴的.
是尊貴的.
我們每一個人在神的眼中.
都是有價值的.
這個價值不是用人間去衡量的.
如果用人間去衡量.
其實就可以分了很多不同的階層.
所以盼望我們能夠在神的裡面.
來看到上主對我們每一個人的接納.
聖誕的意義.
是在於我們和主耶穌建立的信心.
主耶穌自己就是這樣面對.
在地上最艱難的時候.
生下來自己決定不了睡哪張床.
媽媽要將他放在馬的槽.
吃東西的地方.
誰媽媽想自己的寶寶.
放在馬槽裡面呢.
我們當年那麼窮.
都在醫院裡面.
廣華醫院裡面出生.
回到家裡都有張嬰兒床.
誰呢.
如果你想一下現在這個時代.
任何的小朋友都是最重要的.

$^{481}$無論他出生的預備.
為他預備的各樣東西.
佈置整間嬰兒房.
還有其他的公書教學.
都是最好的.
為他想到最好最好的.
這個就是現代人對小朋友的安排.
古代怎麼會整個政權.
會維繫在一個嬰兒身上呢.
但是我們的主耶穌就是這樣.
由他出生到兩歲.
爸爸媽媽要和他一起逃難.
去到埃及.
直至等待到當時的希律王死.
才敢回來.
還有天使提醒他們回來.
回來都不敢再住在伯利恆.
就回到拿撒勒.
一個完全不是當時代出名的地方.
很窮的地方.
再加上那個地方有很多羅馬的駐軍.
所以變相是猶太人看不起的地方.
那耶穌就變成了拿撒勒人.
各位親愛的弟兄姊妹.
來到今天26號.
牧師很想和大家一起去思想.
我們的救贖主.
他是怎樣來到去成為貧窮.
使我們的生命能夠在神裡面富足.
是生命上面的富足.
是生命因為耶穌而帶來富足.
以致我們能夠在神的面前.
來得到生命的恩典.
就是寶貴的福音.
很盼望我們能夠在聖誕.
在這個節慶.
我們想到的就是.
你的生命有沒有和主耶穌有一個緊密的關係.
以致你能夠真的在神的面前.
你得到的就是神給你的永生的福氣.

$^{521}$願主幫助我們.
很盼望我們真的大家一起.
在這個地方.
神的家.
我們去將福音傳開.
讓更多人得到生命的豐富和建立.
願主祝福大家.
我們一起有簡單祈禱.
多謝主.
讓我們今早能夠透過聖詩.
透過神的話語.
來認識你.
祝你有美好的生命相交.
多謝你揀選我們.
多謝你將我們拯救.
更加多謝你無條件地救贖我們世人.
使我們得到這個永恆的生命.
很盼望我們能夠像牧羊人.
像弱失瑪利亞.
一樣地來到.
雖然是卑微.
但是在神眼中.
我們是尊貴的.
讓我們看到我們的身份.
不是用人的角度去衡量.
而是在神的裡面.
看自己的價值.
求主祝福.
叫我們每一位.
洪恩家的弟兄姊妹.
在神的恩典裡面.
繼續的剛強.
長進.
有力.
信心是增長.
一直的來得著.
上主美好的祝福.
這樣的感恩禱告.
奉耶穌基督的名.
阿們.

$^{561}$祝福各位.
\newpage



\section{歌羅西書 1:15-23}
\label{sec:yWaJULCl_Mc}
\textbf{2022.01.02《尊貴無比的基督》 歌羅西書 1\_15-23 (和修版) 老耀雄牧師}
\newline
\newline
連結: \href{https://youtube.com/watch?v=yWaJULCl-Mc}{\texttt{https://youtube.com/watch?v=yWaJULCl-Mc}} ~~~~ 語音日期: 2022-01-05
\newline
\newline
\hyperref[sec:_NykIK_k97E]{\small{< < < PREV SERMON < < <}}
~
\hyperref[sec:index_chronic]{\small{[返順時目]}}
~
\hyperref[sec:index_scriptual]{\small{[返順卷目]}}
~
\hyperref[sec:VMnlnDfbr9Q]{\small{> > > NEXT SERMON > > >}}
\newline
\newline
歌羅西書 1:15-23
\newline
\begin{longtable}{cl}
\hline
\hline
章節 & 經文 (和合本修訂版)\\
\hline
1:15 & \begin{tabularx}{0.7\textwidth}{X} 愛子是那看不見的神之像，是首生的，在一切被造的以先。 \end{tabularx} \\ \\ \relax
1:16 & \begin{tabularx}{0.7\textwidth}{X} 因為萬有都是在他裡面造的，無論是天上的、地上的，能看見的、不能看見的，或是有權位的、統治的，或是執政的、掌權的，一概都是藉著他為著他造的。 \end{tabularx} \\ \\ \relax
1:17 & \begin{tabularx}{0.7\textwidth}{X} 他在萬有之先；萬有也靠他而存在。 \end{tabularx} \\ \\ \relax
1:18 & \begin{tabularx}{0.7\textwidth}{X} 他是身體（教會）的頭；他是元始，是從死人中復活的首生者，好讓他在萬有中居首位。 \end{tabularx} \\ \\ \relax
1:19 & \begin{tabularx}{0.7\textwidth}{X} 因為神喜歡使一切的豐盛在他裡面居住， \end{tabularx} \\ \\ \relax
1:20 & \begin{tabularx}{0.7\textwidth}{X} 藉著他，神使萬有與自己和好，無論是地上的、天上的，都藉著他在十字架上所流的血促成了和平。 \end{tabularx} \\ \\ \relax
1:21 & \begin{tabularx}{0.7\textwidth}{X} 從前你們與神隔絕，心思上與他為敵，行為邪惡； \end{tabularx} \\ \\ \relax
1:22 & \begin{tabularx}{0.7\textwidth}{X} 但如今，他藉著他兒子肉身的死，已經使你們與他自己和好了，把你們獻在他的面前，成為聖潔，沒有瑕疵，無可指責。 \end{tabularx} \\ \\ \relax
1:23 & \begin{tabularx}{0.7\textwidth}{X} 只要你們持守信仰，根基穩固，堅定不移，不致動搖，離開了你們從前所聽見的福音的盼望；這福音也是傳給天下一切被造之物的，我—保羅作了這福音的僕役。 \end{tabularx} \\ \\
[1ex]
\hline
\hline
\end{longtable}
$^{1}$各位姐妹平安.
平安.
今天是2022年的第一個主日崇拜.
恭祝大家新年快樂.
身心靈健康.
中國人有一句話.
他說一年之計在於春.
一日之計在於辰.
一生之計在於勤.
勤力總是中國人的美德.
也是香港人的核心價值.
我還記得在2018年山竹颱風吹襲香港.
當時天文台懸掛十號風球.
風球下了之後.
香港人在那些倒塌的樹木.
仍然都未移除的情況之下.
走去搭車那種狼狽樣.
當時有一張海報就這樣說.
香港人就算面對重重的困難.
都要去上班.
這句話究竟是讚譽還是諷刺呢?.
對於基督徒來說.
勤力的重點並不是放在工作上.
而是放在言讀神的話語上.
信徒要對聖經有全面的認識.
才是一個成熟的基督的跟隨者.
牧師甚願2022年.
同恩家的每一個會友.
能夠勤於靈修.
勤於言讀聖經.
往年的講道.
我都是以主題的形式去宣講神的話.
今年除了主題的形式講道之外.
也都加入了以經卷的形式去講道.
所以今年的第一季.
就會以三次連續的形式.
去做一個經卷講道.
去帶出整卷歌羅西書的核心信仰.
讓大家能夠對歌羅西書有一個很整全的了解.
今日我就是以歌羅西書一章.

$^{41}$十五至廿三節去探討基督論的三個向度.
就是基督的本質.
基督的旨意.
以及信徒的回應.
我們一起去祈禱.
你藉著你忠心的僕人.
將信仰的核心書寫在不同的經卷裡.
承傳給我們.
求主讓每一個聽道的會眾.
能夠將你的話語藏在心裡.
成為我們腳前的燈.
路上的光.
光照著我們的生命.
讓我們每一個都能夠傳言成聖.
我們祈禱奉主耶穌基督.
明智祈求.
阿們.
由於是一個經卷式的講道.
我就會用較多少少的時間.
讓會眾先了解歌羅西書的背景.
歌羅西書在新約裡面是一封很簡短的書信.
只有四章.
沒有相關的資料可以知道.
哥羅西教會是怎樣建立的.
事實上就算你翻查《士道行傳》.
都沒有提到哥羅西這個地方.
也都沒有資料顯示.
保羅曾經去過這個哥羅西教會.
所有有關哥羅西教會的資料.
只能夠在歌羅西書裡面去找.
在內政裡面.
即是歌羅西書的內容裡面.
我們知道.
直接建立哥羅西教會的人.
不是保羅.
而是以巴弗.
以巴弗本身是哥羅西人.
他很可能在保羅在以忽所事奉的時候.
由保羅帶領他去決志信主.
他信了主之後.

$^{81}$經過信仰的栽培.
然後就回到他老家.
哥羅西去傳福音.
之後就建立了教會.
而哥羅西教會被建立的時間.
估計在主後53至55年.
亦都是在保羅在以忽所事奉的那段時間.
所以哥羅西教會是保羅傳福音給以巴弗.
由以巴弗所結出的一個福音的果子.
由於哥羅西教會的信徒大部份都是外邦人.
因此由外面引入教會的異端.
就在哥羅西教會出現了.
雖然書信裡面沒有直接講出異端的教義.
但從內政去看.
仍然可以勾畫一些輪廓出來.
這個異端的理論.
是很有具有日後諾斯底主義的元素.
或者現在只是初型.
但對於教會來說.
是要即時處理的.
因為諾斯底主義認為物質本身是邪惡的.
這種想法.
引致他們徹底否認.
上帝在基督裡面的那種道成肉身.
即是否定耶穌的道成肉身.
他們認定.
即是諾斯底主義認定.
神是不可以與一個人的軀體結合在一起.
諾斯底主義從兩種途徑去推翻.
神道成肉身的真理.
第一種論說就是否定耶穌的真實人性.
如何去否認?.
他們說耶穌的身體只是一個幻影.
是假象.
看起來像是真的.
這種就叫做幻影說.
第二種論說就是.
完全否定耶穌的神性.
以推翻道成肉身.
即是說耶穌只是一個人.

$^{121}$兩種傾向在《歌羅西書》裡面都有提及.
不過保羅最主要去辯解的.
是針對第二種論說.
就是對於否認耶穌具有神性那一面.
這二端不單止有諾斯底主義的思想.
其實也有基督教的元素.
因此道學說的中心.
是基督教和異教思想的混合物.
而這些思想雖然沒有否認基督.
但是卻將主耶穌從寶座裡面拉下來.
他仍然留給基督一個位置.
不過不是一個最高的位置.
其實在香港現在一種叫五教合一.
正正是這種輪廓.
五教合一當中有耶穌.
即是說耶穌不是最高那個.
耶穌只是與其他等同.
所以這種以基督教為招來的手法.
就更加危險.
因為謬誤的思想會在教會裡面孕育出來.
貶低了基督的榮耀.
這樣對於那些明目張膽的挑戰.
更加令人防不勝防.
因此歌羅西書的主題很清晰.
就是宣告主耶穌基督的絕對崇高的地位.
以及他是唯一能夠滿足我們贖罪的要求.
歌羅西書是屬於監獄書信之一.
即是說保羅寫信這封信的時候.
他是在坐牢的.
保羅的監獄書信裡面.
總共有三篇講述基督論的經文.
歌羅西書一章十五至二十節是其中一篇.
另外兩篇記載在《以弗所書》和《肥納比書》裡面.
對於歌羅西書的熱端.
最少為教會帶來三種危險.
第一,這種論述將恩典得救的真理歪曲了.
他認為單是相信基督不足以叫人得救.
即是說他不相信恩信清義.
第二,他誤解了基督徒的生活.
走向兩個極端.

$^{161}$哪兩個極端呢?.
第一,苦行式的禁慾.
一是放縱情慾.
不過最危險的都是第三種.
第三種就是他貶低了耶穌基督.
對於這些熱端來說.
基督不是一個得勝的主.
他只不過是天上地下其中之一.
因此保羅這段基督論的經文.
亦都是特別針對這些熱端來寫的.
讓哥羅西的信徒對耶穌基督有一個正確的認識.
我們研讀歌羅西書可以讓我們再一次確認和評估我們的信仰.
我們所相信的主耶穌基督.
是一真神裡面的第二位.
他具有神人異性的位格.
亦都是唯有他在十字架上的救贖.
才能夠滿足到天父所要求的義.
換句話來說.
天上人間除他以外並沒有另一條得贖的途徑.
我們開始進入經文.
剛才主席沒有幫我們逐一讀經文.
我們在程序表裡面看回我們的經文.
基督的本質是尊貴無比的.
在第十五節上.
他說:愛子是那看不見的神的像.
這句話說明了基督和神的關係.
甚麼叫神的像?.
神的像可以譯為神的形象,真象.
就是神的自己.
因此基督耶穌不是次一等的靈體.
祂本身就是神.
新普及的譯本譯得比較好.
它是怎樣譯呢?.
它說:基督是不可見的上帝那可見的形象.
神本來是看不見的.
我們沒有一個人能夠看見神.
但藉著耶穌基督的道成肉身.
就可以讓人能夠見到神.
因為基督是神的像.
從基督的身上可以反映出神的本質.

$^{201}$我們看到原本的經文是愛子.
但也有譯本譯成基督.
為甚麼呢?.
和學本是譯愛子的.
因為愛子的原文是一個關係代名詞.
就是那一位.
所以有些譯者將那一位譯為愛子.
也有些譯者譯為基督.
和修本就譯為基督.
上帝是不可見的.
但基督是可以見的.
不但可見.
祂還來到人間.
很近距離地將神的形象彰顯出來.
如果我有讀過《約翰福音》.
我就聽過這句話.
祂說:從來都沒有人見過上帝.
只有在父懷裡的獨生子.
才能夠將祂顯露出來.
因此人可以透過可見的基督.
去認識神的本質.
那神的本質是甚麼呢?.
在15節下解釋了.
是首先的.
是一切被做的以先.
基督不是被做的.
是首生的.
首生是一個形容詞.
是在用一個時間.
在一切被做之先.
即是說基督耶穌不是被做的.
而是在時間上先於一切的創造物.
亦是指耶穌基督的身份.
高於一切的受造之物.
因為基督的本質.
就是宇宙和世界的創造者.
一切都是基督所創造.
16節上再加以解釋.
因為萬有都是在祂裡面做的.
在祂裡面是甚麼意思呢?.

$^{241}$是指上帝藉著基督創造了一切.
基督是一切創造的來源.
16節下,祂這樣說.
祂說無論是天上的,地上的.
能看見的,不能看見的.
說是有權位的,統治的.
說是執政的,掌權的.
一切都是藉著祂,為著祂做的.
天上的,地上的是指空間.
能看見的,不能看見的.
是指視覺,也就是我們的感官世界.
有權位的,統治的,執政的,掌權的.
是指國家,政權或者體制.
因此一切都是藉著祂,為祂做的.
這句話有兩個意思.
第一個意思是藉著祂.
藉著祂的意思就是指由祂所創造.
第二個意思是為著祂.
為著祂的意思就是指創造的目的.
是為了基督.
17節,祂說「祂在萬有之先.
萬有也靠祂而存在」.
因此結論就是基督就是創造者.
一切的創造都是因基督而存在.
而一切創造的目的都是為了基督.
所以萬有能夠存在.
是依靠著耶穌基督.
沒有基督.
萬有根本不能夠獨立於,存留於宇宙裡面.
如果萬有的存在都要依靠耶穌基督.
那我想問.
人類究竟有什麼能耐.
可以超然於耶穌基督.
人類又怎麼可能不依靠耶穌基督而活呢?.
你可以說.
不是呀,我現在生活得很好.
沒有耶穌我也是這樣生存.
你也可以這樣說.
但耶穌的永恆性你有沒有呢?.
當你今天去信耶穌.

$^{281}$你有沒有想過信耶穌.
你能夠有一種永恆的生命.
你不信耶穌.
你何來有一個永恆的生命呢?.
在《士徒行傳》十七章二十八節就說過.
他說我們生活,行動,存在都在於他.
或者現實一點說一句.
我們今天只可以坐在這裡.
是因為基督.
是因為基督賦予我們生命和氣息.
也是基督賦予我們一切的價值.
有信徒在發生禍患的時候質問神.
為什麼發生意外的時候是我呢?.
有信徒去問神.
為什麼我會有這樣的禍患.
有神學家就這樣回應.
每天都有人死.
為什麼死的不是你?.
當信徒質問.
為什麼我有禍患的時候.
神學家就這樣回應.
每天都有人死.
為什麼死的不是你?.
當我們遇到不好的遭遇的時候就質問神.
神賦予我們每一個生命氣息以及價值.
馬克理穆斯在美國建立了一間.
二萬三千人的馬鞍峰教會.
寫了一本銷量達到五千萬本的書.
叫《彪竿人生》.
《彪竿人生》的一書就指出.
我們被造是為了基督.
我們活著都是為了基督.
保羅一開始就說出基督和神的關係.
接著又說出基督是創造者之外.
他更加說出基督和教會的關係.
也就是基督和我們的關係.
基督在創造這個世界之後.
又參與其中.
第十八節他這樣說.
祂是身體教會的頭.

$^{321}$祂是原始.
是從死人中復活的守生者.
好讓祂在萬有中居首位.
《和修本》有一個括弧.
括著教會兩個字.
顯得很特別.
是指身體就是教會嗎?括弧.
《和合本》沒有這個括弧.
《和合本》是譯為.
祂是教會全體之首.
原文教會和身體是兩個獨立的名詞.
比較好的譯法就是.
基督也是身體的頭.
這身體就是教會.
教會既然是基督的身體.
教會自然應該有基督的屬性.
十八節下就解釋了這個屬性.
祂是原始.
是從死人中復活的守生者.
好讓祂在萬有中居首位.
十八節的守生和十五節的守生是同一個原文.
不過十五節的守生.
所指的是時間上.
即是基督在創造之先就存在了.
十八節的守生是指他的歷史性.
是指基督是第一個從死裡真正復活的.
如果我們看四福音裡面.
都記載過很多死人復活的事件.
最出名的一定是拉薩路的復活.
拿恩城寡婦的兒子都復活了.
甚至保祿和彼得曾經讓死人復活.
但這種復活不是永恆.
復活之後肉身最終都會死亡.
但基督的復活就不同了.
基督的復活是永恆的.
不會再死.
正正因為基督裡面有永恆的生命.
而屬基督的子民.
即是你和我.
就因為能夠與基督合一.

$^{361}$因此亦可以享有這種永恆的生命.
為何基督是首位從死裡復活.
為何他在萬有中居首位呢?.
基督不單單在舊約的創造裡面居首位.
因為他是萬有之先.
是一切生命的創造者.
但當罪進入世界之後.
我們知道罪的功架就是死.
因此死就是我們人類的終局.
如果我們不認識耶穌.
死就是我們的終局.
但基督的救贖.
就是我們人因著相信主耶穌而有永生.
當一個罪人再被賦予永恆生命的時候.
基督就成為一個新的創造者.
明不明白呢?.
再簡單一點說一句.
創造之先.
基督是一個創造者.
創造萬有.
所以他居首位.
直到新約.
或者可以這樣說.
直到耶穌基督降生之前.
人都是絕望.
人沒有永恆生命.
人的最後結局都是死.
耶穌基督出現了.
將這一批死的人.
將要死的人.
能夠賦予一個永恆生命.
他為我們人類再賦予一個創造者.
居首位的明星.
基督作為創造者.
他在萬有之先.
所以他居首位.
但基督作為教會這個身體的頭.
而他亦都是從死裡復活的首一個.
第一個.
他使所有屬於教會的信徒.

$^{401}$都有永恆的生命.
他從創造.
從歷史教會.
都是居首位.
可以這樣說.
世界不能夠沒有了基督.
沒有了基督的世界.
是一個敗壞.
以及沒有盼望的世界.
教會不能夠沒有了基督.
沒有了基督的教會.
只是一棟建築物.
沒有與主同行.
沒有與主同在的屬靈生命.
以及沒有了這個永恆的盼望.
一切都是謊言.
教會裡面充滿了罪人.
一切的罪人.
都要面對永遠的死亡.
沒有一個可以倖免.
不過.
基督是從死人中復活的首生者.
基督既能夠擺脫死亡.
亦都表示一切屬於他的.
即是包括你和我.
都能夠藉著基督的恩典.
我們都可以擺脫死亡這個綑綁.
這個亦都是我們人類.
永遠沒有辦法擺脫的一個死穴.
不相信耶穌基督.
沒有永恆的生命.
這個是唯一真理.
弟兄姊妹.
2022年.
第一個崇拜.
牧師想問你.
你屬於基督嗎?.
你在基督裡面嗎?.
你的生命是否連於基督?.
是.

$^{441}$我們一定說是.
我不相信你們說不是.
不過我們的回應.
不單止要表現在我們頭腦上的認知.
說是.
而是表現在我們的言行.
感神的恩典和我們的行為.
相不相稱?.
當我們說我是連於基督的時候.
當我們說我屬於基督的時候.
當我們說我們在基督裡面的時候.
我們問一問自己.
我們的行為與這個恩典相不相稱?.
這個是我們在2022年.
第一個的信仰反省.
第二,基督的旨意.
在第十九節.
因為神喜歡使一切風聖在裡面居住.
神的風聖是什麼呢?.
就是神的一切所有.
神的一切真善美.
以及神一切的本質和權柄.
神的形象造人.
因此人在被創造的時候.
已經擁有了神的真善美.
我們有神的真善美.
不過很可惜.
犯罪了之後.
神的形象就失落了.
我們只有殘餘的神的形象.
在《讓福音》主耶穌曾經說過.
他說我來了.
是要揚得生命並且得得更豐盛.
神不單止將他的風聖賜給基督.
也很想賜給我們.
神將萬有的主宰.
教會的頭都賜予給基督.
是有一個目的的.
二十節,直著去.
神使萬有與自己和好.

$^{481}$就是要萬有.
直著基督與神和好.
這個就是福音.
是耶穌基督的福音.
讓人與神的距離拉近.
是耶穌基督的福音.
使信徒與神的永恆有份.
這種萬有的恩典.
包括在天上,在地上.
即是說這個不單止人類.
是包括了神所創造的一切.
都可以直著基督在十字架上所流的寶血.
促使了復和的實現.
保羅在《二弗所書》一章八節說過.
他說要照著所安排的.
在時機成熟的時候.
使天上的,地下的.
一切所有的.
都在基督裡面同季如日.
這種復和不單單是人類.
都是萬物的.
萬物與神復和.
帶動萬物與萬物的復和.
人不是單獨與神復和的.
都是要與人與萬物復和的.
人展現出一個彼此相愛的行為.
就能夠彰顯在基督裡面同季如日.
很多時候我們的信仰誤解了.
我們與神的關係好就行了.
我個人與神的關係好就行了.
所以這麼多網上崇拜.
我與神的關係好就行了.
我有崇拜生活嗎?.
不是的.
你不單單要與神復和.
你要與人復和.
你要與萬物復和.
我與人如何復和?.
教會裡面這麼多紛爭.
這麼多勾心鬥角.

$^{521}$這麼多決裂.
就藉著耶穌基督去復和.
不然教會怎麼能夠叫教會呢?.
教會只能夠變成教會.
就因為我們有彼此相愛.
在彼此相愛裡面.
人與人能夠復和.
人對神所產生的萬物.
產生一種愛護的心.
而不是摧毀的心.
這個就是萬物.
這個就是與萬物產生一種復和的關係.
我們不單單愛神.
愛人.
還要愛我們身邊的人和事.
弟兄姊妹.
我們的信仰是一個復和的信仰.
首先是我們與神復和.
然後我們與其他的人復和.
最後是與萬物復和.
用信仰的說法.
與神復和是什麼?.
就是與神建立一個關係.
你與神有什麼關係?.
如果拍拖都知道.
你拍拖的時候.
無時無刻不想著對方.
無時無刻不想和對方聊天.
熱戀中的男女.
如果你說熱戀中的男女只顧著工作.
不顧著對方.
那就不是熱戀.
你與神的關係.
就是一個熱戀的男女.
你想不想和他聊天?.
你想不想聽他話語?.
你想不想和他相會?.
如果你說沒有.
沒有.
那我們就說.

$^{561}$這個只是軀殼.
我們說什麼叫掛名的基督徒.
就是這一類.
掛名的基督徒.
我們與人是否能夠建立一個和好的關係?.
我們心裡是否仍然對別人心存怨憤.
忿忿不平?.
能否以一個寬恕的心.
與別人和好.
建立一個真善美的關係?.
一個外人去看教會.
不是看那個建築物.
而是看教會裡面的肢體.
究竟如何去相交.
一個外人進入教會.
見到肢體相交都是面左左的.
吸引不到他留在教會.
一個外人來到教會.
見到肢體之間能夠彼此相愛.
這一種的相愛.
是吸引人留在教會.
教會不能夠實踐彼此相愛.
其實不可能叫做教會.
極其量都是宗教活動.
第三點,人的回應.
第二節,他說從前你們與神隔絕.
心思上與他為敵,行為邪惡.
當人未認識神的時候.
罪將人與神隔絕了.
罪亦都將人與神為敵.
心思意念都是邪惡.
其實我們想一想我們自己就知道.
未信耶穌的時候.
我們心裡面想的是什麼.
或者你可以說我信了耶穌之後.
我心裡面想的都是一樣.
我們沒有給信仰去轉化我們的生命.
第二節,他說但如今.
他藉著他兒子肉身的死.
已經使你們與他自己和好了.

$^{601}$把你們獻在他的面前成為聖傑.
沒有瑕疵,無可指責.
藉著基督在十字架上的代贖.
他流出的寶血.
洗淨了我們的罪,救贖了我們.
使一切相信主耶穌的人與神和好了.
在神面前成為一個義人,無可指責.
有很多悉經家或者神學家.
對這段經文用了很多方法去解.
沒有瑕疵,無可指責,行不行的呢?.
我們一個信徒能不能夠做到.
沒有瑕疵,無可指責呢?.
如果做不到的.
這個要求是不是過高?.
對於人來說.
信了耶穌之後,是不是真的可以沒有瑕疵,無可指責?.
有些悉經家就說.
這個沒有瑕疵,無可指責.
是在神面前,對神.
在神面前你能夠完完全全的毀身.
完完全全的將自己的生命獻上.
神就看你為無瑕疵,無可指責.
如果我們認為自己是沒有罪的.
如果我們認為能夠靠自己得救.
如果我們認為靠自己的行為.
可以得到神的喜悅.
這個正正是褻瀆和侮辱主耶穌基督.
因為沒有一個人能夠藉著我們的行為.
靠著自己,認為自己沒有罪.
當我們信了耶穌,人的本質改變了.
人的罪被赦免了.
人在神的面前能夠成為聖潔.
這個都是在神眼裡.
這一切都不是我們可以做得到.
這一切都是基督的恩典.
當我們接受了這個恩典.
我們怎樣回應呢?.
23節,他說.
只要你們持守信仰,根基穩固,堅定不移,不致動搖.
離開了你們從前所聽見的福音的盼望.

$^{641}$怎樣可以持守信仰?.
這個也很老生常談.
當信仰成為你生活習慣的一部分.
成為一種習慣的時候.
就能夠持守了.
我相信在座每一位.
今天早上起床都會有刷牙的,是不是?.
但從來都不會問你為什麼還要刷牙.
因為刷牙成為了你的生活的一部分.
你不會再問為什麼.
就好像吃飯祈禱一樣.
你不會不做.
也不會問為什麼要這麼做.
怎樣根基可以穩固?.
建樓要打樁,就要地基建穩.
信仰要穩固,就要在神的話語裡紮根.
讓根部能夠牢牢地抓住泥土.
讓神的話語成為我們生命的根基.
坦白說,這個老生常談的說法就是經常靈修.
你不靈修,就不能抓住泥土.
第一節我今天要講的題目是尊貴無比的基督.
基督的尊貴.
是因為祂是一個獨一的真神.
祂配得一切的榮耀.
基督的尊貴.
因為祂是創造者.
一切都是祂的創造.
也是為祂而被造.
基督的尊貴.
因為祂是我們與神之間的中寶.
藉著基督,我們才能夠與神和好.
基督的尊貴.
因為祂是永恆的.
只要是屬於祂的.
都能夠被帶入永恆.
正正因為基督的尊貴.
當我們擁有了這份寶貴的恩典的時候.
我們就要持守信仰.
根基穩固.
堅定不移.

$^{681}$不致動搖.
剛才我說過.
我們有一個反省.
就是我們的行為與恩典相不相稱.
如果我們不想成為一個掛名的基督徒.
我們就要勇於面對我們自己.
和要去實踐.
今年洪恩嘉的年提是什麼?.
記不記得?.
「毀身事奉」.
正正就是要每一個基督的跟隨者.
作出一個具體的回應.
「毀身」就是要全人的投入.
「事奉」就是以神給我們的恩賜.
去服侍神,服侍人.
那你說我有什麼恩賜?.
沒有人是沒有恩賜的.
每一個都有.
每一個都有神給你的恩賜.
你願不願意發圭?.
你願不願意被神所使用?.
你要想一想.
在2022年.
你要想.
我可以在教會哪個崗位上高事奉呢?.
教會實在有很多事奉的崗位.
等候著大家「毀身」的參與.
我們一起同心祈禱.
親愛的主,我們獻上感恩.
因為你尊貴無比.
卻為我們甘願道成肉身.
你為我們走上十字架.
為我們每一個跟隨你的.
獻上你的生命.
為我們贖罪,捨命.
讓我們與你有份.
讓我們有永恆的盼望.
求你讓我們在乘勝之旅的過程裡.
堅守信仰,根基穩固.
堅定不移,不致動搖.

$^{721}$求你加能設力給我們.
讓我們能夠回應你的呼召.
就是「毀身」事奉,服侍神,服侍人.
好叫「紅恩家」.
能夠傳言的被你所使用.
讓你的旨意能夠成就.
讓更多未得之民.
因著「紅恩家」而能夠認識主.
歸入你的名下.
你的名被高舉.
我們完全祈禱.
奉教主耶穌基督.
得勝名字祈求.
阿們.
願主祝福大家.
祝福大家.
\newpage



\section{路加福音 7:36-50}
\label{sec:VMnlnDfbr9Q}
\textbf{2022.01.09 《以愛還愛》 路加福音 7\_36-50 (和修版) 曾敏姑娘}
\newline
\newline
連結: \href{https://youtube.com/watch?v=VMnlnDfbr9Q}{\texttt{https://youtube.com/watch?v=VMnlnDfbr9Q}} ~~~~ 語音日期: 2022-01-09
\newline
\newline
\hyperref[sec:yWaJULCl_Mc]{\small{< < < PREV SERMON < < <}}
~
\hyperref[sec:index_chronic]{\small{[返順時目]}}
~
\hyperref[sec:index_scriptual]{\small{[返順卷目]}}
~
\hyperref[sec:biU1pOyNb74]{\small{> > > NEXT SERMON > > >}}
\newline
\newline
路加福音 7:36-50
\newline
\begin{longtable}{cl}
\hline
\hline
章節 & 經文 (和合本修訂版)\\
\hline
7:36 & \begin{tabularx}{0.7\textwidth}{X} 有一個法利賽人請耶穌和他吃飯，耶穌就到那法利賽人家裡去坐席。 \end{tabularx} \\ \\ \relax
7:37 & \begin{tabularx}{0.7\textwidth}{X} 那城裡有一個女人，是個罪人，知道耶穌在法利賽人家裡坐席，就拿著盛滿香膏的玉瓶， \end{tabularx} \\ \\ \relax
7:38 & \begin{tabularx}{0.7\textwidth}{X} 站在耶穌背後，挨著他的腳哭，眼淚滴濕了耶穌的腳，就用自己的頭髮擦乾，又用嘴連連親他的腳，把香膏抹上。 \end{tabularx} \\ \\ \relax
7:39 & \begin{tabularx}{0.7\textwidth}{X} 請耶穌的法利賽人看見這事，心裡說：「這人若是先知，一定知道摸他的是誰，是個怎樣的女人；她是個罪人哪！」 \end{tabularx} \\ \\ \relax
7:40 & \begin{tabularx}{0.7\textwidth}{X} 耶穌回應他說：「西門，我有話要對你說。」西門說：「老師，請說。」 \end{tabularx} \\ \\ \relax
7:41 & \begin{tabularx}{0.7\textwidth}{X} 耶穌說：「有兩個人欠了某一個債主的錢，一個欠五百個銀幣，一個欠五十個銀幣。 \end{tabularx} \\ \\ \relax
7:42 & \begin{tabularx}{0.7\textwidth}{X} 因為他們無力償還，債主就開恩赦免了他們兩個人的債。那麼，這兩個人哪一個更愛他呢？」 \end{tabularx} \\ \\ \relax
7:43 & \begin{tabularx}{0.7\textwidth}{X} 西門回答：「我想是那多得赦免的人。」耶穌對他說：「你的判斷不錯。」 \end{tabularx} \\ \\ \relax
7:44 & \begin{tabularx}{0.7\textwidth}{X} 於是他轉過來向著那女人，對西門說：「你看見這女人嗎？我進了你的家，你沒有給我水洗腳，但這女人用眼淚滴濕了我的腳，又用頭髮擦乾。 \end{tabularx} \\ \\ \relax
7:45 & \begin{tabularx}{0.7\textwidth}{X} 你沒有親我，但這女人從我進來就不住地親我的腳。 \end{tabularx} \\ \\ \relax
7:46 & \begin{tabularx}{0.7\textwidth}{X} 你沒有用油抹我的頭，但這女人用香膏抹我的腳。 \end{tabularx} \\ \\ \relax
7:47 & \begin{tabularx}{0.7\textwidth}{X} 所以我告訴你，她許多的罪都赦免了，因為她愛的多；而那少得赦免的，愛的就少。」 \end{tabularx} \\ \\ \relax
7:48 & \begin{tabularx}{0.7\textwidth}{X} 於是耶穌對那女人說：「你的罪都赦免了。」 \end{tabularx} \\ \\ \relax
7:49 & \begin{tabularx}{0.7\textwidth}{X} 同席的人心裡說：「這是甚麼人，竟赦免人的罪呢？」 \end{tabularx} \\ \\ \relax
7:50 & \begin{tabularx}{0.7\textwidth}{X} 耶穌對那女人說：「你的信救了你，平安地回去吧！」 \end{tabularx} \\ \\
[1ex]
\hline
\hline
\end{longtable}
$^{1}$弟兄姊妹平安.
雖然今天很多弟兄姊妹都在家裡敬拜神.
但是我們的心是一體的.
開始的時候我們一同去禱告.
讓我們這麼卑微的人.
可以去敬拜你.
你知道我們的軟弱.
知道我們心裡所想的一切.
嘴裡所說的每句話.
你知道我們的行事為人.
天父在這裡求你徹底寬恕我們一切的罪惡.
免得這一切難了我們去親近你.
去聽你的聲音.
求你今天潔淨.
孩子的嘴,孩子的生命.
潔淨他們每一個弟兄姊妹的生命.
讓我們能夠擁抱你的說話.
願聖靈在我們當中工作感動我們.
讓我們知道神的心意.
你激動我們的心.
讓我們向著你所喜悅的方向去走.
我們這樣的禱告是奉耶穌基督的名字.
阿們.
今天這段經文.
你說很多人說又不是.
你說很少人說又不是.
曾經有人試過說今天的經文.
試過從不同的向度.
我今天很想嘗試從一個向度.
就是以愛去還愛.
我們今天說愛這個題目.
我們2022年開始沒多久.
我們第二個星期.
原來就有新的疫情.
我們的生活又要改變了.
大家會多了留在家裡.
少了社交活動.
教會很多聚會都改為網上的直播.
很多人是感到不安的.
因為不知道疫情會怎樣發展下去.

$^{41}$我也知道有人跟我說.
我們最重要是要為疫情禱告.
是的.
我們真的可以為這個疫情去禱告.
但是我們有沒有試過.
從另一個角度去看現在這件事.
看現在這個疫情.
上帝讓我們在疫情裡.
可以停下來.
靜一靜.
本來我們覺得很可惜.
有很多活動,很多事情.
很多人我們要見,要做.
但實際上是讓我們停下來.
不是一個半個人可以停下來.
是大家一起停下來.
靜一靜.
我們去思考有關神的事情.
是好的.
今天讓我們透過.
這段《路加福音》七章.
三十六到五十節這段經文.
去思考上帝的愛.
和我們怎樣去回應祂的愛.
《路加福音》七章三十六到五十節.
記載一個法利賽人.
邀請耶穌和他吃飯.
法利賽人是宗教領袖.
很熟悉律法.
也很熟悉舊約的聖經.
他們很樂意為人講解神的律法.
他們也會看著人.
「你!你!你們有沒有遵守?」.
「你!聖經這樣說,是嗎?」.
「神的律例這樣說」.
「你遵守了嗎?」.
他們很喜歡這樣想.
他們被人視為很高尚的人.
在猶太人的社會裡.
大家很尊重他們.

$^{81}$但這些人非常重視階級的觀念.
所以這個法利賽人.
肯定覺得耶穌至少和他比較近階層.
所以才會邀請他和他一起吃飯.
耶穌接受法利賽人的邀請.
去他家裡吃飯.
一起吃飯是很重要的.
我們無論什麼國家的人.
一起吃飯是很重要的.
對猶太人來說.
一起吃飯一定表示了接納.
那種肯定.
「我願意和你做朋友」.
「我才會和你一起吃飯」.
耶穌願意吃這頓飯.
是代表他也接納了這個法利賽人.
耶穌接納法利賽人.
不是因為他來自一個很高尚的階層.
也不是因為這個人很好.
所以他要和他吃這頓飯.
事實上在福音書裡.
耶穌曾經責備過法利賽人.
以假冒為善.
他們很刻意地在人面前.
去表現自己是多麼虔誠.
事實上他們是不認識神.
不明白神的心意.
也不能切實地.
真正地進行神的教導.
他們會指責別人.
他們離神的道很遠.
他們的生命並不高尚.
但是當有法利賽人邀請耶穌的時候.
耶穌願意應約.
因為耶穌愛每一個罪人.
無論這個人來自什麼階層.
他都會應約.
有錢的,窮的.
被人看為高尚的.
被人看為卑微的.

$^{121}$他都會去,他不分.
在言籍當中出現了一個沒有被邀請的人.
一個有罪的女人.
在37節裡這樣說.
「那城裡有一個女人是個罪人」.
如果從原文來說.
可以有另一個看法.
有一個更加詳細的看法.
這句話可以這樣翻譯.
那個城裡有一個女人.
她素常是個罪人.
又或者說.
她一直都是一個罪人.
這個女人在城市裡是很有名的.
不是因為她有什麼貢獻.
或者她個人的條件是多麼優厚.
而是因為她曾經犯罪.
可能是比較嚴重的罪.
有些學者認為她可能是妓女.
但經文沒有直接說.
這裡漏了一個空白的位置.
我們只可以說城市裡的人.
都知道她犯了什麼罪.
在人的眼中.
這個女人地位非常卑微.
而且因為罪的緣故.
她是被人看不起的.
這個女人出現在筵席裡.
不是因為她是被邀請的.
純粹因為她要找耶穌.
整段經文讓我們看到一個對比.
法理賽人和那個有罪的女人.
對耶穌的態度有很大的分別.
首先我們再要問一個問題.
我們看回剛才那段經文.
我們覺得這個法理賽人.
她是否認識耶穌呢?.
她是否知道她真正的身份是誰呢?.
她會稱呼耶穌為老師.
因為耶穌會教導人聖經.

$^{161}$他會教導人神的話.
當那個有罪的女人.
後面的經文我們會看到.
這個有罪的女人.
在耶穌身上做了一連串的事情.
這個法理賽人心裡.
就會這樣說一句.
在39節.
「這人若是先知.
一定知道摸他的是誰.
是個怎樣的女人.
他是個罪人」.
先知是神的代言人.
神會透過他去說話.
會說預言的.
神也會告訴先知很多事情.
所以法理賽人認為.
如果耶穌是先知.
你一定會知道那個女人是罪人.
她這樣說.
如果耶穌是先知.
如果.
如果就會知道.
大家覺得她認為耶穌是否先知呢?.
如果耶穌是.
你便知道他是罪人.
實際上她覺得耶穌是否先知呢?.
不是的.
覺得你一定不是.
你知道他是罪人.
你應該很想和他分開.
他骯髒.
你教導聖經.
大家不應該分開.
不同階層.
這個法理賽人叫做西門.
到了這個時候.
耶穌開始稱呼他的名字.
其實我們知道.
耶穌一定是先知.

$^{201}$他不單知道那個女人是罪人.
他也知道這個法理賽人西門在想什麼.
當然他不是一個普通的先知.
他是尼賽亞先知.
上帝揀選了那個.
要去拯救他的百姓.
在這裡.
48字我們還看到.
耶穌對那個女人說.
你的罪都赦免了.
耶穌宣告那個女人所有的罪惡.
都獲得赦免了.
其實只有神有赦罪的權柄.
換言之.
耶穌不單只是先知.
他還是神.
我們知道耶穌基督是神的兒子.
他就是神.
所以他有赦罪的權柄.
這些法理賽人完全不知道.
他是無知的.
他不認識真正的耶穌.
除了法理賽人不認識耶穌的身份.
他不認識耶穌是哪位之外.
法理賽人也沒有很熱情的款待耶穌.
44到46節記載耶穌評價西門.
他說:我進了你家.
第一,你沒有給我水洗腳.
是很自然的.
當時的人去款待客人.
他們走很多路.
腳非常骯髒.
也是因為他們的鞋子比較特別.
很容易碰到灰塵.
所以他去到一個家庭裡.
如果那個家主很喜歡款待客人.
他至少會給水洗腳.
西門沒有.
第45節:你沒有親我.
大家的關係這麼友好.

$^{241}$可以親嘴問安.
表示一種友好.
西門也沒有.
好像有點冷漠.
46節:你沒有用油抹我的頭.
這是對於上賓的款待.
如果你認識對方.
覺得對方的位置很崇高.
這個人是特別值得你尊敬的.
你至少要用橄欖油抹他的頭.
西門也沒有.
他沒有熱情的款待耶穌.
他純粹是請他吃飯.
他的態度表明了他怎樣看待耶穌.
其實耶穌在他心目中.
不是特別值得人尊敬的.
很配得人去服侍的.
他也不是很愛耶穌.
他的態度顯出了冷漠.
這是法理賽人西門的態度.
對比之下.
那個女人的態度很不一樣.
44至46節記載了.
耶穌對於那個有罪的女人的描述.
她說這個女人用眼淚滴濕了我的腳.
又用頭髮擦乾我.
西門沒有給水讓耶穌洗腳.
但這個女人用她的眼淚洗了她的腳.
當然我們稍後會知道.
為何她的眼淚會滴在她的腳上.
用眼淚洗她的腳.
直接用頭髮擦乾.
這個女人的位置是很卑微的.
很謙卑的一個位置.
也是一個很尊崇的表達.
還有45節.
這個女人從我進來就不住地親我的腳.
西門不會和耶穌親嘴問安.
表達友好.
但這個女人就親她的腳.

$^{281}$是表達友好.
也是一種很大的尊重.
一種敬意.
很體重她.
還有第46節.
這個女人用香膏抹我的腳.
這個香膏非常貴.
比橄欖油更貴.
是很貴重的香料.
西門沒有用橄欖油抹耶穌的頭.
這個女人卻用香膏.
很貴重的香膏.
抹耶穌的腳.
又是一種尊重.
一種景仰.
她很重視.
也是表示一種更大的友好.
這令我們看到.
兩個人的分別很大.
我們再回頭看.
為什麼這個女人的態度這麼特別.
會做這三件事情呢?.
這個女人的眼淚滴濕耶穌的腳是因為什麼?.
這是感激的眼淚.
很感激.
應該是較早前的時間.
耶穌寬恕了她的罪.
雖然我們看經文後面.
耶穌宣告你的罪被赦免.
好像那一刻才赦免她的罪.
其實不是.
是較早前.
耶穌寬恕了她的罪.
她很感激耶穌.
一直以來她的罪.
她感到很羞恥,很難過.
你想想整個城市的人.
記得的是什麼?.
不是她的好.
不是她的尊貴.

$^{321}$整個城市的人記得.
就是她是一個犯罪的女人.
她犯了什麼罪.
整個城市的人都知道.
人們口中就在討論.
都是她的不好.
多羞愧.
這生活多難過.
還有自己的罪疚感.
她很辛苦.
她根本沒辦法在別人面前抬起頭來.
也過不了自己這一關.
她背著的是一個重的罪的重擔.
她的人生是沒有出路的.
如果是這樣下去.
她也很清楚在永恆裡面.
等著她的是什麼.
這個.
我和你同樣清楚.
耶穌寬恕了她的罪惡.
是耶穌寬恕她.
她感受到從未有過的輕散.
和自由.
她感受到耶穌的愛和接納.
她很感動.
從來沒有人這樣對她做.
但是耶穌會這樣做.
所以她一定要去法利賽人.
西門的家找耶穌.
她要多謝他.
她要回報他.
在這個女人的心中.
耶穌是很崇高.
很配得愛和尊重.
她的回報是要有實際的行動的.
在這個回報裡面.
剛才是說了三件事.
我們再來看看.
38節說.
她站在耶穌的背後.

$^{361}$靠著他的腳哭.
眼淚滴濕了耶穌的腳.
就用自己的頭髮擦乾.
她這是一個很深的情感表達.
很感激.
深深的感激耶穌.
她真的覺得自己很不配.
她說:你知道嗎?.
我過去的人生是這樣的.
現在我可以被寬恕.
你知道我多開心.
多感動.
有人愛我,接納我,寬恕我.
真的無法形容這種激動.
那是很深很深的.
真正的情感表達.
接著第二件事.
也是動作.
用嘴巴連連親她的腳.
謝謝她.
真的要謝謝她.
也是一種尊重.
一種欣賞,肯定,回應.
第三件事.
她用聖滿香膏的玉瓶.
把香膏抹在耶穌的腳上.
抹在她的腳上會有什麼效果?.
很香,很滋潤.
去尊崇她.
她配得最好的東西.
去獻上她最好的東西.
在這裡.
我看回天淵之別.
一個是非常非常愛耶穌.
很感動,很感激.
會有行動的.
去做出來.
一個是很冷漠的.
什麼都不做.
只是吃一頓飯.

$^{401}$可能她覺得吃飯已經很好.
我給你面子.
我請你吃這頓飯.
我可以不請你吃.
兩個人來自不同的社會階層.
為什麼有這麼大的分別?.
耶穌一定看到這個分別.
祂自己親身感受到.
兩個人對祂不同的回應.
所以在這裡.
耶穌說了一個比喻.
耶穌說這個比喻.
也是要回應西門.
因為西門心裡.
在說話.
以為耶穌應該不知道.
我心裡在說話.
耶穌如果你知道.
這個人是誰.
他是一個罪人.
下面有很多.
「…」.
可能在想.
你今天會打算怎樣對待他?.
你知不知道這個人怎樣?.
可能有很多事情要繼續說.
他覺得自己說這些.
沒有人會聽到.
耶穌完全知道.
完全知道.
所以在這裡.
祂要回應他.
要跟他說.
第四十節.
耶穌回應他.
西門,我有話要對你說.
西門說:老師,請說.
有兩個人欠某一個債主的錢.
一個欠五百個銀幣.
一個欠五十個銀幣.

$^{441}$一個銀幣就是一天.
當時人的工價.
其實五百個銀幣是很多錢的.
五十個銀幣.
其實也不少.
五十天的工錢也是很多的.
重點是什麼?.
第四十二節.
因為他們無力償還.
無論你欠債是少還是多.
你都無力償還.
你欠債少.
也是還不了的.
你欠債多.
也是還不了的.
結局是一樣的.
當然有分別.
差距是有的.
欠多一點.
欠少一點.
也是還不了的.
於是債主開恩赦免了他們倆人的債.
耶穌問:那兩個人.
哪一個是更愛他呢?.
更愛的債主呢?.
債仔更愛債主呢?.
西門很懂得回答.
全世界都懂得回答.
我想是那個多得赦免的人.
那個欠了五百個銀幣.
現在債主說:行了.
不用還了.
哇!震開心.
我欠很多.
我現在不用還.
我真的一身輕.
是自由了.
人生多好.
那個欠五十個銀幣.
欠得沒有那麼多.

$^{481}$你寬免我的感覺.
好像沒有五百個銀幣那麼開心.
耶穌接著在說.
所以我告訴你.
他有很多的罪.
在說這個有罪的女人.
很多的罪都赦免了.
因為她的愛多.
在這裡小心看.
意思不是說.
因為這個女人很有愛心.
所以我就赦免她的罪.
好像剛才那個經歷.
就是因為她為我做那麼多事.
又用眼淚洗我腳.
用頭髮抹我腳.
又用膏去高抹我的腳.
又親我的腳.
因為她做了那麼多事.
於是我寬恕她.
不是.
反過來.
這個女人那麼有愛.
證實了.
她真的有很多罪.
得忘赦免.
因為她經歷過罪的赦免.
她很激動.
她就會去表達她的愛.
因為她經歷赦免.
她感受到自己被赦免.
所以她要表達她的愛.
這是重點.
「少得赦免的,愛的就少」.
當那個人表達愛少的時候.
證實了什麼?.
好像她經歷的赦免是少.
其實不是.
那個債主.
不願意去寬免別人的債.

$^{521}$是有原因的.
這個什麼叫做赦免少.
愛的就少呢?.
如果用舊恩的層面來說.
實際是說什麼?.
其實一個人表達的愛少.
證明了什麼?.
我們好像很少經歷神的赦免.
不是神不赦免我們.
是我們沒有好好將我們的罪.
交在上帝的手裡.
去經歷祂的赦免.
上帝是願意徹底.
百分百赦免我們的罪.
但是我們有多看到自己罪的可怕.
認罪.
交給上帝.
如果是這樣的時候.
我們經歷的赦免也不會是少的.
是徹底的.
很多時候.
在這個位置.
其實這個例子.
不單是說錢,銀方面的欠債.
是免債的問題.
更加說我們在舊恩的路上.
上帝寬恕我們.
寬恕我們的罪.
我們的回應.
是顯示出我們經歷了多少上帝的寬恕.
如果我們好像這個女人.
我們真的願意認自己犯了罪.
是這麼可怕的.
徹底將自己交在上帝手裡的時候.
上帝比她的寬恕大.
法利賽人覺得我有什麼問題.
我很高尚的.
我教別人學聖經.
我不知道多熟悉律法.
你來問我.

$^{561}$從尾巴背回來我都可以.
我經常看著你們.
有沒有很好的遵守神的律例.
我當然很高尚的.
其實感受不是很深.
他不會覺得.
哇!耶穌你真的很寶貴.
你赦免我的罪.
我不會覺得很寶貴.
因為我覺得我好像沒有經歷過.
那種赦免的寶貴.
今天我想問一件事.
這是重點.
我們有沒有經歷過赦罪的寶貴.
我們覺得自己是哪個人呢?.
欠債50個銀幣的人.
還是欠了500個銀幣的人.
問心那句.
我相信我們很多時候會覺得自己是.
欠了50個銀幣的人.
我們哪有這麼多罪.
我沒有去偷,搶,殺人,放火.
我做很多好事.
別人都很讚我的.
你知道我的人很好.
我信耶穌的時候.
想來想去都不知道自己有什麼問題.
有什麼不好.
所以別人說赦罪.
想了很久.
感覺很難入罪.
我記得我剛剛信耶穌的時候.
我也有這個問題.
那時候覺得自己在學校都算乖乖女.
一直很努力讀書.
我沒有什麼很大的.
在家裡很活躍的東西.
一直在家裡乖乖女.
在學校好像也不錯.
信耶穌初期我記得我自己.

$^{601}$心態一定是覺得.
我只是欠50個銀幣而已.
但當我一直看聖經.
越看越多的時候.
原來這些都是罪.
原來驕傲也是罪.
原來還有這些.
很多很多.
說是非,貪心,妒忌.
很多,逐樣逐樣看.
越看越覺得自己很不妥當.
原來我是很多罪的.
然後在我裡面的聖靈會光照我.
喂,你看看.
你那時候這樣跟別人說話.
你這樣想的.
你覺得別人不知道,是嗎?.
你看你生命裡面是這樣的質地.
一直上帝開著我的眼睛去看.
每一個階段上帝都會開著我的眼睛去看.
你的生命有這些你要處理.
是一個又一個又一個階段.
每個階段上帝都給我對自己有新的認識.
所以我在講道裡面.
我從來不敢說.
我覺得自己真的很好.
你們沒理由不欣賞我,是嗎?.
真的覺得自己很根深的說話.
我很怕說這些話.
因為我知道.
我一樣是一個罪人.
只不過是一個蒙恩的罪人.
我絕對覺得自己是欠五百個銀幣的人.
有時候在神面前.
面對自己的罪惡的時候.
都覺得真的很羞愧.
不要這樣做就好了.
不停地想這些事情.
會很難過.
我絕對是那個欠五百個銀幣的人.

$^{641}$有時候看聖經人物.
看到他們有一些軟弱.
有不完美的地方.
我不敢說.
他們真的壞到這樣.
我就不會了.
我一定會好好地告訴你.
我心想.
換了我在那個位置,我更壞.
我敢這樣說.
換了我在那個位置,我更壞.
我一定是那個欠五百個銀幣的人.
我不覺得自己是欠五十個銀幣.
很寶貴.
沒錯,在人生不同的階段.
我對上帝的回應是不同的.
但是我越來越覺得.
是的,我真的是一個蒙恩的罪人.
我欠上帝太多.
實在太多.
這裡說得最好.
因為他們無力償還.
今天不管我是犯罪少還是犯罪多.
我都還不了給上帝.
我們欠上帝就永遠欠他.
是上帝說開恩寬恕我.
我才得到這個自由.
一個有罪的女人很明白這個道理.
很能夠愛耶穌.
很深的去愛耶穌,很激動.
但是這個法利賽人.
口說不清,說聖經不明白耶穌.
今天我去想一件事.
很重要的一件事.
今天我們回到教會.
怎樣看自己和上帝的關係.
我們是否覺得自己是一個很蒙愛的人.
上帝愛我們.
上帝愛我們不是說人生有很多錢.
我沒有經濟上的挑戰困難.

$^{681}$我一定沒有病.
我一定沒有人際上的困難.
我不會遇到沮喪.
我什麼都很好的,很順利的.
如果上帝是這樣,你就愛我了.
這樣其實和皇大先識的信仰沒有什麼分別.
我們看神的愛看得太膚淺.
上帝最愛我們的地方是祂.
派祂的兒子耶穌基督為我和你犧牲.
上帝寬恕我和你的罪.
是耶穌基督為我和你的罪犧牲.
所以我們得蒙救贖.
其實這種輕散我們有沒有感受到呢?.
還是我們覺得太普通了?.
為什麼我們覺得太普通了?.
是不是我們看不到自己的罪呢?.
我不覺得自己有什麼罪.
所以我不覺得這個世面有什麼特別.
如果我看到我的罪.
我知道這個世面這麼寶貴的時候.
我應該很開心.
就算我回到來.
我的生命有很多困難.
錢的問題,身體的問題.
家庭的問題,人際的問題.
各方面人生的挫折等等.
我都還是很開心.
因為上帝給了我們最寶貴的東西.
而且上帝向我們顯明祂毫無保留的愛.
耶穌基督很愛那個有罪的女人.
其實祂也愛西門.
祂愛每一個人.
在福音書裡面到最後.
祂給我們看到那個愛到極點.
就是我把我的命給你.
我要向全世界的人宣告.
我愛你們每一個人.
我把我的命給你們.
我什麼都毫無保留.
有創造天地萬物的主宰.

$^{721}$向我們宣告這份毫無保留的愛.
我和你還有什麼是不開心的呢?.
所以我都在思考這個問題.
為什麼我們常常不開心?.
是不是我們不明白這份愛?.
我們不知道自己罪的重.
因為我們不明白這份愛.
其實我們都很難有持久的動力.
去愛慕神的話語.
如果我們不停地在講壇說.
「神啊,你要返主教學校」.
那個感覺很外在.
「來,弟兄姊妹,返主教學校」.
一定要學慕神的話語.
是好的,這個呼籲是重要的.
但是如果你在裡面.
你很學慕上帝.
我就想多看一些聖經.
我想知道上帝怎麼說.
我想認識這個神是怎樣.
你自己主動就會去學慕.
我們走在一起.
你知不知道我們的探討題目就不一樣.
我們會講一下聖經.
但是你有沒有留意我們走在一起.
我們講什麼?.
我們講聖經.
我們講什麼?.
我們容易講是非.
我們有時候站在法理賽人的位置.
去看.
你覺不覺得那個是誰?.
我都說他不是很好.
你覺不覺得?.
他不是很好.
我看他.
很激氣.
很激心.
他這樣.
你知不知道.

$^{761}$他這樣,我們嘴巴就用來講這些.
彼此之間很不和諧.
我們又不明白上帝想我們在這裡做什麼.
他那麼愛我們.
他又希望更多人得到這份愛.
但我們自己都沒有感受到這份愛.
我們怎會有這樣的動力.
再叫更多人來擁抱耶穌基督的那份愛呢?.
我們都不太認識耶穌的心意.
所以在過程中.
很願意在新的一年.
我們更深認識自己.
是一個怎樣的人.
我們認罪.
好像欠債的.
欠債的.
有罪的女人.
她這樣去愛上帝.
我們去愛耶穌.
我們都要傾盡自己的一切.
去愛耶穌.
我們學會親近耶穌.
自己裡面發出來.
我們離開罪惡.
將自己有的東西盡力擺上.
今年教會的主題是「毀身是奉」.
講壇都會呼籲.
弟兄姊妹你們想好了.
你們有什麼恩賜擺上.
其實我覺得更加寶貴.
是弟兄姊妹你覺得.
上帝太好了.
我欠上帝很多.
那我今年要擺什麼給上帝?.
金錢.
時間.
恩賜.
不是出於我們一直推啊.
推啊.
是你裡面發出來.

$^{801}$你想將這些給上帝.
更加重要的是.
我們學武神.
回來彼此相愛.
上帝愛我們每一個.
我們真的不完美.
我們有很多問題.
走在一起一定有衝突.
預料到了.
家裡人都會有爭吵.
神的家裡的人都會有爭吵.
爭吵之後要靠上帝的愛.
走在一起.
讓上帝的心滿足.
這是更加重要的.
在結束的時候.
很想講一個很寶貴的經歷.
今天我們有師班獻師.
是很寶貴.
我們都很愛詩歌.
其中其實我記得.
我不知道這麼多年.
師班有沒有機會唱一首歌.
《二恩典》.
這首歌非常.
大家很熟悉.
無論是甚麼年齡層的弟兄姊妹都很熟悉.
作者約翰·牛頓.
其實有很深刻的經歷.
他的媽媽是一個很敬虔的基督徒.
他在牛頓小時候.
盡力將他所有的靈異遺產.
傳授給他.
教他看聖經.
不過那時候牛頓還小.
他在牛頓七歲的時候離世.
媽媽很早就安息主懷.
去到十一歲.
牛頓跟著爸爸踏上航海的生活.
他很快就染上了水手那些.

$^{841}$很惡的,很壞的習慣.
去到二十歲.
一個非常年輕的時間.
他成為一個相當冷血的奴隸販賣者.
他專門喜歡販賣非洲的黑奴.
他為奴隸船定下一個目標.
他自己應該是有經營的.
奴隸船的目標是甚麼?.
在一個很小的空間裡面.
要塞滿那些黑奴.
我管得了他們是甚麼光景.
他們有沒有食物.
他們有沒有足夠的空氣,環境.
不重要.
我要賺錢.
總之在一個很小的空間裡塞滿了.
這些人在他心目中不是人.
是真的物件.
賺錢的物件.
很小的空間可以塞滿二百五十個人.
那些黑奴的手腳都上了枷鎖.
那些倉庫充滿了排泄物,汗臭,發酸的食物.
還有腐爛的屍體.
由非洲去到英國的南卡羅萊納港.
這一段行程有四分之一的奴隸會死在途中.
死了之後他們會拋入海裡.
他不要把這些人當人.
但上帝的管教臨到.
在一次的厄運當中.
他成為了.
他很諷刺.
他成為了黑奴的奴隸.
為什麼呢?.
因為當時有一個非洲的情婦.
竟然虐待他.
鎖住他.
他那次的經歷是很淒涼的.
最後他的父親要派船去救他.
整個的光景被人看到了.
他是沒有面子的.

$^{881}$他也是很慘的.
辛森寧在那次經歷了一次的管教.
在他回英國的行程當中.
他讀到一本書叫《效法基督》.
他在想自己過去的生活充滿了罪惡.
他開始看到自己的罪.
那種黑暗醜陋.
他求神憐憫他.
求神救他離開這種黑暗的日子.
過了一段時間之後.
有一次徹底扭轉了他的人生.
他經歷了風暴.
當時他的船在海上經歷了風暴.
那次船已經破爛了.
他面對生命的危險.
那一次他徹底回轉歸向上帝.
他跟上帝說.
上帝憐憫我們.
如果能上岸.
還能生存.
我要過不一樣的生活.
他就徹底改變.
他歸向耶穌基督.
上帝賜他一條生命.
他之後改變.
他寫了一首很出名的歌《奇異恩典》.
講述自己整個過程.
他如何犯罪.
如何經歷上帝的寬恕.
那種拯救.
之後他很用力反奴隸的反賣.
這種行動.
以前他反賣奴隸.
信耶穌後他就要反這件事.
去改變.
將他的人生交託給上帝.
第一步.
認清自己是罪人.
深深的犯罪得罪神.
沒錯.

$^{921}$我們犯的罪是不一樣的.
不是所有人像約翰·牛頓那樣犯的罪.
但我剛才告訴你.
我在上帝面前看自己的時候.
上帝都會光照我.
我並不讓其他人聖潔高尚.
我都是一個罪人.
如果上帝不救贖我.
甚至上帝不揀選我.
今天我不會站在這裡去講道.
在這裡我們一起禱告.
透過這段經文提醒我們.
要認識自己的不配.
我們都是罪人.
求上帝開我們的眼睛.
讓我們看到自己的生命.
有多少的污穢.
每個人的罪都不一樣.
但上帝知道我們.
我們每個人都無法還這個罪債.
正如經文所說.
我們欠五百銀幣.
五十銀幣都還不了.
今天我們每個人都無法處理我們罪的問題.
等待著我們的下場.
我們都沒有辦法.
我們沒有辦法.
處理我們罪的問題.
等待著我們的是永遠的死亡.
永遠的和上帝分開.
今天我們擁有這個新生.
很值得歡喜快樂.
天上為我們拍掌.
天父求你給我們歡喜快樂.
每個星期回來我們帶著感恩的心.
回來聚會.
求你給我們學會你的華語.
經常被我們看得不配.
我們就甘心樂意.
將我們的金錢,時間.

$^{961}$我們的恩慈.
用在你的身上.
我們回到這裡.
愛你所愛.
你愛每一個熟你的人.
讓我們愛對方.
讓我們愛你的心儀.
\newpage



\section{啟示錄 21:1-8}
\label{sec:biU1pOyNb74}
\textbf{2022.01.16 《神的帳幕在人間》 啟示錄 21\_1-8 (和修版) 蔣文忠博士}
\newline
\newline
連結: \href{https://youtube.com/watch?v=biU1pOyNb74}{\texttt{https://youtube.com/watch?v=biU1pOyNb74}} ~~~~ 語音日期: 2022-01-16
\newline
\newline
\hyperref[sec:VMnlnDfbr9Q]{\small{< < < PREV SERMON < < <}}
~
\hyperref[sec:index_chronic]{\small{[返順時目]}}
~
\hyperref[sec:index_scriptual]{\small{[返順卷目]}}
~
\hyperref[sec:0IxLPtCeZ88]{\small{> > > NEXT SERMON > > >}}
\newline
\newline
啟示錄 21:1-8
\newline
\begin{longtable}{cl}
\hline
\hline
章節 & 經文 (和合本修訂版)\\
\hline
21:1 & \begin{tabularx}{0.7\textwidth}{X} 我又看見一個新天新地，因為先前的天和先前的地已經過去了，海也不再有了。 \end{tabularx} \\ \\ \relax
21:2 & \begin{tabularx}{0.7\textwidth}{X} 我又看見聖城，新耶路撒冷由神那裡，從天而降，預備好了，就如新娘打扮整齊，等候丈夫。 \end{tabularx} \\ \\ \relax
21:3 & \begin{tabularx}{0.7\textwidth}{X} 我聽見有大聲音從寶座出來，說：「看哪，神的帳幕在人間！他要和他們同住，他們要作他的子民。神要親自與他們同在。 \end{tabularx} \\ \\ \relax
21:4 & \begin{tabularx}{0.7\textwidth}{X} 神要擦去他們一切的眼淚；不再有死亡，也不再有悲哀、哭號、痛苦，因為先前的事都過去了。」 \end{tabularx} \\ \\ \relax
21:5 & \begin{tabularx}{0.7\textwidth}{X} 那位坐在寶座上的說：「看哪，我把一切都更新了！」他又說：「你要寫下來，因為這些話是可信靠的，是真實的。」 \end{tabularx} \\ \\ \relax
21:6 & \begin{tabularx}{0.7\textwidth}{X} 他又對我說：「成了！我是阿拉法，我是俄梅戛；我是開始，我是終結。我要把生命的泉水白白賜給那口渴的人喝。 \end{tabularx} \\ \\ \relax
21:7 & \begin{tabularx}{0.7\textwidth}{X} 得勝的要承受這些為業；我要作他的神，他要作我的兒子。 \end{tabularx} \\ \\ \relax
21:8 & \begin{tabularx}{0.7\textwidth}{X} 至於膽怯的、不信的、可憎的、殺人的、淫亂的、行邪術的、拜偶像的和一切說謊話的人，他們將在燒著硫磺的火湖裡有份；這是第二次的死。」 \end{tabularx} \\ \\
[1ex]
\hline
\hline
\end{longtable}
$^{1}$(自動翻譯).
各位弟兄姊妹.
平安.
很感恩在在此大家中間一起分享神的話.
這次是一個很特別的安排.
弟兄姊妹應該都在家中.
去和大家崇拜.
但無論如何.
我相信神的靈和我們同在.
我們心歸向神的時候.
無論我們處於一個什麼位置.
我們都以神為中心.
一起以神的道為我們教會群體的指引.
教導為我們的中心.
能夠聚集.
能夠一起去敬拜神.
這也是神的恩典.
拍攝今天的訊息.
能夠成為大家的安慰和幫助.
我不知道大家在這些日子當中.
有沒有什麼聖經會為大家多讀.
或者對自己的信仰有更深的體會.
我這段時間裡.
很多時候都會有機會去讀啟示錄.
特別在工作裡.
和大學生查經的時候.
其實我們都是去查啟示錄.
問他們查完幾章之後.
去到一個學期的完結.
再開始的時候.
其實他們仍然都說很想繼續查下去.
都是一個很特別的.
對我來說的一個體會.
可能在啟示錄裡面.
特別是它的聖經背景.
聖經的訊息.
都可能在我們今天.
香港的處境當中.
會成為我們的安慰.
和大家分享的這段經文是.

$^{41}$在啟示錄最後的總結.
啟示錄的總結是二十一和二十二章.
而這一段一至八節.
就是總結的開頭的一段經文.
如果你再想闊一點.
其實啟示錄是聖經裡面最後一段書.
可能都是整段新舊約的傳書的總結.
上帝要告訴我們.
告訴祂的子民.
有什麼訊息.
要祂常存心裡.
要祂繼續存在教會當中.
繼續保存.
繼續相信.
和繼續去把持著.
去過我們的生活.
這也是上帝可能在啟示錄最後的時候.
為我們預備的心意.
求主幫助讓我們今天的訊息.
都是讓我們能夠.
存著開放的心去領受.
我們先再有個禱告.
然後我們開始進入上帝的話語當中.
慈悲天父多謝你給我們今天.
雖然我們不能夠實體的見面.
實體的聚集.
但你的靈和我們同在.
你的話吸引我們.
以致我們雖然.
我們的肉身.
是身處在不同的地方.
但我們的靈.
可以是同一.
指向歸向你.
讓你的道.
在我們生命當中.
去得著.
讓我們的心靈.
謙卑.
開放.

$^{81}$去聆聽你的說話.
讓你親自與我們同在.
提醒我們.
祝福.
鼓勵我們.
聽我們同心向你禱告.
奉主耶穌.
得勝.
明哲而求.
阿們.
[音樂].
講一點背景的時候.
可能大家都會知道.
就是.
究竟啟示錄是一卷講.
將來的書.
還是講現在的書呢.
我想.
在新約的研究當中.
都有不同的.
變化.
或者有不同的.
就是時候.
有不同的重點.
去提出的.
特別是當我們.
發現啟示錄的成書.
或者背景.
是.
就是都比較肯定.
在第一世紀末的時候.
那個一個.
是羅馬帝國迫害信徒的.
一個背景之下.
而他最後成書的時候.
可能很多人都覺得.
是.
其實啟示錄都是一卷.
就是對於當代.
當下的信徒.

$^{121}$是一個.
鼓勵.
和一個安慰.
當時是.
羅馬帝國有一個.
皇帝的名字叫多米田.
是一個.
很殘暴的君王.
去殘害了很多基督徒.
也一樣.
令到很多基督徒都真的.
去到一些地下教會.
就是地下教會真的在.
地面下面的教會.
在當時就開始出現.
他們堅持信仰.
繼續聚集.
繼續去敬拜.
以及繼續彼此的.
相交.
這個是當時的信徒很多時候.
面對很多的.
就是對信仰的迫害的時候.
他們能夠繼續堅持下去的時候.
而他們就.
嗯.
有啟示錄的信息.
來幫助他們.
所以.
你可以說啟示錄既是在說一個.
終末的世界.
新天新地.
將來.
上帝怎樣完全的.
去掌權等等.
但是呢.
同時間這個信息也都會對.
當下今天的信徒.
其實是有他們的意思.
有他們的領受.

$^{161}$有他們的幫助的.
我們其中一個很重要的主題就是.
主必快來.
是不是?.
就是大家從頭到尾.
都有說主會快來.
既是在說一個.
終末的大審判.
也是一個信徒迫切的禱告.
就是.
盼望現在那種黑暗.
或者現在那種迫害.
終必過去.
能夠可以.
在他們眼見的世界當中.
可以有一種改變.
這個也是信徒.
很期盼上帝在當中的工作.
這個很重要.
就是說我們的盼望.
不只是對於今世無意思.
就是不只是被動的等候.
我們忍耐到底就可以了.
這個世界.
其實是.
上帝好像不在當中去改變.
或者不在當中去作王.
或者去掌權.
而是覺得上帝他的能力.
仍然會在這裡有他的彰顯.
所以呢.
這個終末的盼望.
是會讓信徒對於今天.
有一份力量.
有一份堅持下去的信心和盼望.
所以這個的.
對於這群繼續堅持信仰的信徒是很重要的.
他們沒有放棄.
他們希望可以繼續敬拜.
繼續去讀上帝的話.

$^{201}$繼續作見證.
繼續去在受著一些不同的迫害的時候.
不同的壓力的時候.
他們仍然很忠心見證上帝.
這個是.
當時或者其實裡面仍然是有一個這樣的向度.
去提醒我們.
這份盼望.
不是什麼都不用做.
而是提醒我們.
正正主是會作王.
正正主是最終會得著世界的時候.
我們今天.
特別作為他的教會群體.
我們需要繼續持守我們的信仰.
去等候迎接主.
他自己的臨在.
這個是在基督徒裡面提醒我們的信息.
到了這兩章的時候.
我現在很快地跟大家一個概括.
然後我們就會下去今天的經文.
這裡有幾重要的信息.
這是幾點跟大家點題去提出的.
上帝的榮耀.
遍佈聖城.
大家剛才讀經的時候也讀到.
我們史圖約翰看到.
聖城.
新耶路撒冷.
這個也是一個很重要的意象.
對於猶太人來說.
耶路撒冷是他們的命脈.
他們的皈依.
所以.
在新天新地的時候.
仍然用回新耶路撒冷這個意象.
來表述.
但是這個表述.
大家要留意的是.
他已經不是在歷史上的那種表述.

$^{241}$是有一種超越了歷史的情況.
而是有上帝永遠的掌權之下.
當中有一種新的意思.
所以特別是.
譬如他在21章第23節.
他說神的榮耀偏半聖城的時候.
他說當時是沒有黑夜.
你看到紀水路線讀的經文裡面.
很多東西都是不再有了.
在一個新的時代.
在一個新的上帝徹底掌權的時候.
其實很多過去.
舊的時代的東西他都不再有了.
而在21章第23節就說沒有黑夜.
很有趣吧.
沒有黑夜.
為什麼沒有.
以前我們的世界裡面有黑夜有白晝.
這個也是創造當中.
七天創造裡面一個秩序.
但是到了新聖城的時候.
他說沒有黑夜.
是高揚親自作為指紋的燈照耀著他們.
神的高揚在人間.
然後好像明亮的燈台照耀他們.
大家可能想.
那就好了.
不用交電費了.
教會又不用開電掣了.
就是說.
是不是?.
上帝高揚永遠照耀.
大概大家未必明白.
不是好像就這樣按著字面去解釋.
是說上帝永遠帶領他們.
光照他們的生命.
他們.
有時候我們是不是?.
都會很迷惘.
都會有黑暗.

$^{281}$都會有失落的時候.
不知道會怎麼尋求上帝.
但是就是說到了那個情況的時候.
他永遠照耀著我們.
永遠提示著我們.
我們不會迷失.
不會失去了方向.
是一種這樣的.
跟神很緊密的關係.
然後接著就說上帝與子民永遠掌權.
這個就很特別了.
我們知道在.
就是新天新地裡面.
上帝跟我們分享了一切.
但是來到這裡.
特別說到在22章第5節.
就說.
跟我們一起去作王.
這個很特別.
連上帝的王權.
都跟我們分享.
在那個.
聖城的時候.
整個社群的生活.
他是很.
就是.
honor(榮譽).
他是很.
就是對於他的子民是有一種.
很大的權柄賜予他們.
讓他們一起去管理.
這個屬於他的地方.
這個是到了.
新城.
聖城又殺難的時候.
一種很特別的賦權.
給我們的子民.
不僅跟他分開恩典.
分開他的權柄.
給我們一起去管理.

$^{321}$第三點就是大家都.
可能剛才也說了很熟悉的就是.
主必快來.
成就他的旨意.
一個亞南文.
可能大家以前的詩歌都聽過.
"Mawen lava"(馬丁拉法).
那個字.
其實就是主必快來.
就是.
當時剛才說了.
在地下教會.
他真的刻了這些亞南文.
在牆壁上面.
去成為他們的提醒.
雖然他們不能夠.
是不是?.
就是面對著很多的困難.
很多的阻隔.
就是他們不能夠如常地.
在一個普通的自由公民的生活狀態當中.
可以下到一個地下去做這件事情.
但是呢.
他提醒自己.
就是主必快來.
這個必快來是在說.
既是相信這個局面是會改變.
而且相信.
唯有上帝才會永遠掌權.
而不是當時的羅馬君王.
成為他們的掌權.
我們進入到經文.
就是今天這三節的.
這一段的經文.
我會分三段跟大家去分享.
一節三節.
神要親自跟我們同在.
這裡是講一份關係.
雖然在啟示了一個中末的經文裡面.
他是強調一個天地的改變.

$^{361}$但是大家留心.
其實他最終的焦點.
都是放在神跟他的子民的關係裡面.
一節三節.
是強調著神親自與我們同在.
是強調著這份關係.
然後四至五節.
神要更新他的一切.
這個更新一切.
不是純粹是說世界發生的事被更新.
而是更新我們自己生命的經歷.
我們人生當中.
可能很多起跌.
很多不同類型的傷痕.
或者很多不同的狀況需要醫治.
因為我們生命可能是在人世裡面.
經歷了很多碰碰撞撞.
很多的傷害,被傷害之下.
一切在主的當中都被更新.
第三段就是.
不單止這樣.
我們會被賜予生命的滿足.
因為最終的結果.
上帝是要我們生命經歷一份.
跟他永遠同在之後.
生命感到滿足.
感到喜樂.
這是在這一段的一至八節當中.
他給我們的訊息.
我會逐一跟大家分享下去.
在一至三節裡面.
剛才大家讀了這段經文.
我再讀一次新漢語譯本給大家聽.
我看見一個新天新地.
因為先前的天和先前的地都過去.
海也不再有了.
我又看見聖城新也又剎那.
從神那裡至天上降下來.
已經預備好了.
將新娘為她的丈夫裝扮好了.

$^{401}$我又聽見有響亮的聲音.
從寶座發出來.
說漢那神的將王與人同在.
他要與他們同在.
他們要作他的子民.
神要親自與他們同在.
作他們的神.
很清楚.
是不是?.
就是說見到一個新天新地.
先前的天.
先前的地都過去了.
過去了.
是的.
這是啟祖上常用的字眼.
好像已經是屬於過去的意思.
留在那裡.
沒有帶著.
沒有隨著.
上帝的掌權跟隨著.
那些已經是放下了.
有這樣的意思.
大家留心.
有天也有地.
我們有時候信仰裡面.
以為到了.
中末的時候就上天堂.
是不是?.
覺得呢.
就好像是離開了人世間.
去到一個.
純粹在天上和上帝一起享樂的狀態.
但是如果我們仔細看的時候.
新的情況是一個.
天和地都在一起.
特別是他們先前的天.
先前的地都沒有.
是一個新天新地.
是很清楚的.
天和地都被更新.

$^{441}$大家留心.
就是這個.
是從天而降的.
是不是?.
就是不在地上自己演變出來.
而是上帝次下來主動.
再有一個新的創造.
安置在地面的一個.
或者整個的創造界.
這樣看.
是不是?.
我們記得.
創世的第一章.
是否起初神創造天地.
似乎在最後這裡的經文裡面.
是在回應這個.
所謂舊的天和舊的地.
就是上帝第一次的創造.
是一種這樣的回應.
呼應著的一種這樣的味道.
不是說舊的天地就這樣不要.
而是他在當中.
在裡面見得到一些新的景象.
在舊的當中.
見到一些新的景象.
有些舊的就過去了.
不再出現.
但是裡面.
有一些從天上降下.
或者是從上帝而來的.
一些新的景象.
其實這麼說的時候提醒我們.
我們是腳踏實地.
在這個地方裡面.
去存活.
為上帝作見證.
我們這個當下的處境.
我們這個生活.
其實都是在上帝這個保守.
照顧的裡面.

$^{481}$那種新的景象.
是更新了我們自己.
的一路以來的生活.
但是裡面可能有很多的.
不同的問題.
不同的那些罪的那種.
帶來的那種結果.
就放下了.
過去了.
而我們繼續有一種.
更新和演變.
一種這樣的.
身身不適的.
一種這樣的情況.
這個是在.
新天新地裡面我們見到.
一會兒我們去到中間一段.
更加明顯.
有這樣的說法.
而且這段經文裡面.
有不同的意象.
是不是?.
就是.
這個聖城.
新陽剎那.
同天而降.
是什麼意思呢?.
是不是整座城.
就是天空之城這樣降下來呢?.
其實我們要留心.
他接著說的意象.
其實就是.
給我們一個很明確的方向.
去想想.
其實那個意象.
是在說什麼.
他說.
好像新娘.
是不是?.
為她的丈夫裝扮.

$^{521}$等候.
是一個很關鍵的.
就是.
新郎和新娘是一種關係.
是一種愛情的關係.
現在說什麼時刻呢?.
就是他最.
是不是?.
就是應該裝扮得最漂亮的時候.
一個新娘預備出嫁.
就是應該是最美麗.
而且最等待.
是不是?.
去出嫁.
或者等待她另一半.
她的丈夫.
去迎接她.
是這樣的一個時間.
一種這樣的關係.
一種這樣的期盼.
是.
這個是上帝子民的關係.
所以.
當他用一個.
聖城這樣的圖像.
然後畫這個這麼美麗的城.
剛才大家說了.
記住耶路撒冷是他們的敬拜中心.
是他們一個很重要的標記.
到什麼時候用這個標記.
來去.
表示.
好像是.
預備了給.
就是新娘.
預備好自己.
迎接等候她的新郎.
和他之間.
一起.
一起去.

$^{561}$進入一個新的生活.
這個就是那個.
新娘出嫁的意象的意思.
我們要抓得緊.
這個意象才能夠明白那個新城.
那個意思在哪裡.
第二個就是更加特別.
就是用「將莫」或者「會莫」這個字.
就是.
我今天說到也說了題目.
「神的將莫在人間」.
我們大家都很熟悉.
那一節經文.
那個字是.
Tabernacle(天堂).
可能大家聽過.
這個很特別的字.
其實那個字.
是從舊約開始用的.
到新約.
都是用同一個意象的字.
舊約就在說那種.
就是當以色列人.
去四十年在曠野飄泊的時候.
大家記得那個住棚節.
就是住在那些棚.
其實不是真的搭棚.
都是一個Tent(室)而已.
是那些帳幕.
上帝和他一起走的.
他那個會莫.
也是一個Tent(室).
一個帳幕.
不過上帝住.
我們給他一個特別的名字.
叫做會莫.
人住的我就叫做住棚.
這樣去形容.
這個是當時他們.
上帝陪伴著他們四十年.

$^{601}$去一起飄泊流離的時候.
沒有離棄他們.
你去到會莫當中找到我.
都是一個.
我是在你們中間.
一種這樣的意思.
大家留心就是.
所以.
究竟來到啟示錄的時候.
剛才說了.
他怎樣轉化那個意思呢?.
是不是自己是在重複舊約的意思.
還是來到這裡.
有新的意思呢?.
我們要參考一節.
另一節經文也挺特別的.
就是在十三章第六節.
是另一個地方有說.
帳幕.
會莫這個字.
我讀給大家聽.
或者那個情況是怎麼樣呢?.
當時是有海獸.
是.
興起.
然後就褻瀆神的話.
在說的時候他說.
他不僅褻瀆神.
還要褻瀆神的名.
還要說褻瀆他的帳篷.
就是這個會莫的字.
然後那個字.
然後後面還有一句短句的.
他說就是那些在天上.
跟他一起的子民.
那麼如果你看回不同的譯本.
有一個很重要的說法就是.
其實那兩句之間的關係.
就是說那就是.
當聖經說到他褻瀆神.

$^{641}$沒有問題.
但是如果只是說一個建築物.
或者物質的時候.
為什麼都要說明說.
這隻海獸.
要褻瀆神的會莫呢?.
其實原來重點不是那些物質的東西.
原來說.
那就是上帝在天上.
跟你一起的子民.
就是說其實.
你可以想像到.
如果他打那一幫人的時候.
他除了打那個大佬之外.
他還打那些同黨.
他不僅要褻瀆神自己.
還要褻瀆跟從他一起的那一黨的人.
就是那一幫子民.
就是教會.
這個就是當時.
為什麼用這樣的.
海獸褻瀆神的一個圖像.
來參與或者牽涉會莫的事情呢?.
讓我們明白就是.
其實這個會莫不是只是說一個地方.
而是一個群體.
你再這樣想的時候.
其實他第三節的經文是很特別的.
就是說看神的帳幕在人間的時候.
他要跟他們同住.
他們要做他的子民.
神親自要跟他們同在.
做他們的神.
是不是?.
這個很熟悉.
在舊約裡面好像是上帝跟子民立約.
互相去說我是你的神.
你要做我的子民.
那種的關係.
現在是一個圖.

$^{681}$用會莫的圖畫.
神帝會莫在人間的圖畫.
來帶出這個上帝與人同在.
在這群人的中間的意思.
這個會莫象徵著上帝在人中間的意思.
就是說這個群體.
是以神為中心.
你看到會莫.
看到這群人的時候.
你就看到上帝.
因為神透過這群人.
去表示他自己出來.
神親自與人同在.
不需要再隔著會莫.
不需要再隔著萬子.
是不是?.
而是神透過這群人.
真的忠心敬拜聚集的時候.
神就在這裡了.
神就在這群人的中間.
你透過這群人.
就看到神.
這個是第三節裡面一個.
很特別的意思.
他不是在說一件事.
是在說這群人.
看到他們.
就看到神的會莫.
看到神的會莫.
就看到上帝在中間.
當我再想的時候.
有什麼可以更加例子.
或者是實在地.
將這件事表達出來呢?.
我想起一個在歷史當中.
其實也不是發生了很久.
發生在七八年前的一件.
就是一個在歐洲烏克蘭的一個問題.
一個軍事的衝突.
可能大家都留意.

$^{721}$在網上都看到這些照片.
那個地方就是他們國家的中心.
幾乎首都的人民廣場.
好像類似這樣的地方.
當時那段時間.
大概是2014年的1月2日至3月的時間.
常常都有一些穿著教會穆斯普的人在中間.
當大家留心那個細節.
就是說其實這群人.
因為烏克蘭東正教會.
大家明白.
好像一個國家宗教那樣.
但是原來他們有很多不同的.
小的或者其他不同的宗派.
當然天主教會也有.
但是也有一些福音教會是認識的.
有浸信會在當中.
他們描述到就是說.
當時這段時間.
教會也沒有分到宗派.
大家這群的牧者領袖.
主教啊.
牧師啊.
走在一起.
聯合在一起.
站在這些.
好像軍隊和人民的中間.
站在這裡.
這是什麼場景呢?.
這就是一個.
就是有些軍事衝突.
當然有些傷亡.
有很多的生命.
就是會被淘汰的事情發生.
這群的牧者就在當中.
在那個地方裡面.
舉行一些悼念的祈禱會.
他無論是什麼人.
無論他是軍事的軍人.
或者是平民百姓也好.

$^{761}$總之有人的生命傷亡的時候.
他也為這些的靈魂.
為這些生命去哀悼.
去安慰那些失去了至親家人的.
至親人的家人.
去帶給他們的安慰.
帶給他們一些上帝而來的.
那種對生命的接待.
然後讓他們能夠在生命裡面得到平安.
是不是我們在神.
代表著神的帳幕在人間裡面.
是要帶來一種另外的一種.
如果這裡好像是.
我相信當時的氣氛很緊張.
可能很多的仇恨.
很多的傷害.
但是教會就在當中.
賜下了另一種的角色.
是不是?.
就是播下了人愛.
播下了對生命的重視.
播下了珍惜生命.
這些這樣的元素.
在當中.
其實未必能夠解決問題的.
但是卻在當中.
見證給人知道.
上帝在中間.
上帝是有一個角色.
在我們面對著很多不同的.
就是人和人之間的衝突裡面的時候.
我們作為上帝的見證.
我們教會的時候.
是應該將一份的愛.
一份的和平.
去帶給一些.
是充滿著可能仇恨啊.
傷害啊.
掙扎的人的生命當中.
神的帳幕要很實在.

$^{801}$在我們身處的地方當中.
去顯明出來.
我們今天教會.
處於一個地方的時候.
香港是不是還有很多疫情啊.
或者很多不同的壓力啊.
是不是?.
就是人很多的生活上.
都很艱難的.
我們就問自己.
究竟我們也可以怎樣.
代表著上帝.
將這個帳幕.
是能夠展露在.
人.
人民的生活.
社群當中的裡面呢?.
求主去幫助我們.
這個的親自同在.
是跟這份關係.
是一個.
上帝和人子民之間很緊密.
是不是?.
保持一份尊貴.
子民很委身.
而且彼此愛惜的一份關係.
神親自.
不需要經過不同的.
就是以前的舊約裡面的制度啊.
憲制等等.
很多不同的.
就是.
步驟.
不需要.
神親自.
和他的子民同在.
去堅持一份.
和子民.
去立約的關係.
我們去到第二段的經文.

$^{841}$中間的第.
四至五節.
我說的就是神要.
更新.
一切.
的經歷.
這個是一個新漢日本.
不知道在在頂尖.
可以跟我一起.
去將第四節.
一起讀出.
好不好?.
預備.
一.
二.
三.
神必擦去.
哈門每一火眼淚.
不再有死亡.
也不再有悲哀.
哭號.
痛苦.
因為先前的事都過去了.
我不知道大家讀這個.
新漢日本和讀和修本.
有沒有什麼不同.
有沒有一些字眼不同?.
是不是看不看到?.
聽到一點聲音.
B字.
和修本是.
是怎樣的?.
和修本是一樣的.
擦去.
他們一切的眼淚.
是的.
一切.
是不是?.
這裡呢.
每一火.

$^{881}$是的.
就是.
如果你看清楚英文也是這麼說的.
是.
heavy tear.
每一滴的眼淚.
就是那個是.
挺特別的.
是的.
就是.
這個位置.
其實當我們再看第五節的時候.
講到說.
在座寶座上說一切都更新.
你下面也有提到.
是不是?.
就是第二點那裡.
其實那個的.
時態呢.
其實是一個不住.
作更新工作.
這樣的.
就是如果你看看英文.
它就說I am making all things new.
是一個現在.
進行式.
就是說現在還在做.
的一種這樣的時態.
而並非是一種完成了.
我已經做完了.
或者去到最後的時候.
我會做完.
那件事.
一種.
就是.
表示一個.
就是齊全.
或者是.
做完了.
那種的味道.

$^{921}$不是這樣的.
因為當我們看啟示錄的時候.
想到最後的時候.
我們很容易就會想.
有一個傾向就是.
啊 上帝.
就是會一鋪清袋.
處理完所有事情.
是不是?.
所有眼淚.
是不是?.
就全部都更新了.
就是有一種這樣成了.
這樣的.
味道.
但是你讀細心一點看.
第四第五節的時候.
其實它是有一種很細微.
對於每一個生命的經歷.
都很重視.
這樣的意思.
其實沒有衝突的.
就是說神仍然都是所有的.
眼淚.
所有東西都更新.
但是祂不是純粹按一個按鈕.
或者是在這裡.
就是.
好了 我一次過.
好像神仙棒那樣.
一叮.
所有東西都沒事了.
就是突然間.
本來很差.
突然間變好.
而是祂是.
很細微.
很.
很看顧我們每一個生命裡面的經歷.
而每一樣東西.

$^{961}$甚至每一滴的眼淚.
每一滴眼淚表示我們每一個經歷.
都去細心.
去察看.
而且當中都有一種的更新.
這樣的改變.
上帝是一個.
就是.
能夠.
作成一切事情.
而且祂是一個.
無所不看顧的上帝.
甚至是.
你記得勝經也說.
就是連麻將.
都看顧.
我們每一個人的頭髮都數過.
這些都在新約裡面.
其他的經文裡面有描述過.
是很細心.
我們.
不要.
輕看自己.
或者不要以為上帝.
行了 我都不是一個什麼人物.
所以上帝都未必會記得起我.
其實我們不需要有這樣的.
對上帝的想法.
這裡已經說完了.
就是說.
我們每一滴的眼淚.
每一顆火眼.
祂都在著.
祂都去察看.
我們每一個人生的經歷裡面.
起起跌跌.
每一個傷害.
每一個的經歷裡面.
心裡面有挫敗也好.
心裡面有很多不同的難過也好.

$^{1001}$上帝都會在當中.
祂會處理.
都會看顧.
不再有.
悲哀.
哭號.
痛苦.
這些經歷.
神是不住在我們生命當中.
去更新.
直到去終末的時候.
這個更新的工作.
在我們繼續.
今天活著的時候.
就已經開始了.
聖靈在我們生命當中.
不住運行.
去改變.
更新我們的生命.
我想跟大家分享一個.
的生命故事.
這個是偶然在.
電視裡面看到.
這位.
Karen Low 的女士.
她的經歷.
我不知道大家有沒有留.
見過她的生命的經歷.
她現在是一個.
就是都.
幾.
出色的一個.
叫做什麼.
就是.
運動的心理學家.
是的.
一個就是幫不同的人.
去處理.
就是.
運動可能你以為只是體能.

$^{1041}$不是的.
原來心理學的.
可能大家有沒有留意到.
就是奧運那些.
是不是有.
選手.
臨場都退出的.
就是你發覺原來很多年輕的.
當然也不一定是年輕的.
但是我們看到很多都是年輕的.
選手.
其實他們去到.
準備好了.
體能上都準備好了.
但是心理是過不了的.
這個是我相信大家明白.
現在我們.
很多的人.
我們都是面對很大的心理壓力.
其實他自己也有很多的經歷.
雖然今天.
他成為一個運動心理學家.
其實他原來是一個.
游泳.
健將來的.
他小時候.
11歲已經代表香港隊.
去出戰.
去到.
12歲的時候很小的時候.
是不是.
他是.
就是囊括了所有香港的.
那個分靈的紀錄的.
包括所有的泳蝕.
包括所有的.
就是泳蝕還有100米200米的.
不同的長度.
他都還數了.
你想想多厲害.

$^{1081}$所以當時的教練看到他這麼厲害.
就已經預備他.
去參加.
可能15,16歲.
16歲好像是.
還是17歲左右.
預備呢.
參加那個的.
奧運.
是雅典奧運.
但是.
看他自己的一些的網誌.
或者他自己的分享裡面.
其實他一直.
那個心理壓力.
隨著成績越來越好的時候.
或者別人.
對他期望越大的時候.
心理壓力就越來越厲害.
是不是.
無數的.
就是那種的.
緊張的心情.
肚子痛.
失眠.
吃不了東西.
情緒困擾.
流過很多.
不知道怎麼.
就是自己都控制不了的眼淚等等.
這個是他自己很深的經歷.
所以很多的賽事裡面.
都是想.
去到本來.
準備好了.
但是呢.
就去到.
就是預備.
去出賽的時候.
就是一種心理.

$^{1121}$情緒的困擾呢.
使得他不能夠.
去參賽.
到最後.
他都不能夠.
就是如他教練或者他自己所願.
代表香港去參加.
當時的那個奧運.
他說.
心理壓力.
天天都上升.
而且當時他十六七歲開始會考了.
他說當時會考一顆張近.
就是使得他.
作一個很痛心.
忍痛就決定呢.
要放棄游說.
專心去避.
那個考試.
他這麼說.
他經常擔心達不到別人的期望.
教練又不了解他的心情.
當時覺得游說呢.
他的心好像已經到了盡頭了.
所以呢.
他就覺得呢.
很多人都可以.
就是將運動或者是.
學業兼顧.
但是似乎他自己.
就不能夠了.
他就只能夠充滿著失落.
和可惜.
很多人都能夠兩邊兼顧.
但是自己卻做不到.
後來就是他偶然之下.
戶口之後另一個是.
他又上網看看.
怎麼可以有什麼出路呢.
他就偶然上網就這樣找.

$^{1161}$就看到有一個叫做.
就是運動心理學的專業.
於是他就.
下定決心.
就預備自己呢.
將來呢.
要成為一個.
就是運動心理學家.
接著就報大學.
就在香港先讀心理學系.
然後就預備自己呢.
去到美國讀Master.
因為美國才有一個.
專業的資格去讀.
完成他的碩士課程.
回到來.
香港之後.
他開始就在不同的地方.
去推廣.
就是回到.
運動心理學的重要性.
或者那種運動員.
特別年輕的運動員.
那種心理.
那種健康的重要性.
我好像偶然看到一個片段.
就是他.
他都感動的.
向著一群.
就是年輕的運動員.
應該就是當時.
青年的.
香港的游泳隊.
他們做很多的練習.
他們將他們的心聲畫出來.
然後分享.
這些大家可能都.
試過的一些.
心理表達.
或者是.

$^{1201}$關心的一些.
工作坊.
或者一些.
一些這樣的一些.
練習.
他自己說呢.
他說.
他突然間發現了自己一個.
更加重要的角色.
他說以前.
是自己.
由.
自己拿到多少的獎牌.
自己的成就歸於自己.
他說當他見到一群年輕人的時候.
接著他說.
很希望他們能夠.
就是個個都很.
就是不像他那樣.
能夠可以突破到.
那些的心理壓力和關口.
如果他們能夠.
繼續的.
就是去努力.
突破到自己的時候.
他們會贏到更多更多的獎牌.
拿到的成就.
會更加更加超過他自己一個人.
所能夠達到的目標.
他說當他這樣想的時候.
他覺得自己呢.
更加能夠呢.
可以將自己發揮到.
繼續去將自己.
供在這個.
就是游泳的.
這個運動的項目.
事項的裡面.
為什麼我會分享他的故事呢?.
我覺得.

$^{1241}$是不是我們.
人生裡面.
可能都好像他一樣.
不是那麼順利的那些.
或者.
是很必然的.
但是我們能不能夠.
在那些經歷當中.
不斷地去.
整理.
不斷地去調節自己.
更加重要的就是.
我們有沒有不斷地.
讓上帝.
在當中.
去修整我們自己.
或者讓上帝.
去幫我們修整自己.
或者讓上帝.
的靈.
在當中.
更新.
我們自己的生命.
讓我們能夠可以.
繼續地.
在上帝國裡面.
去為祂.
做更大的事.
再不是停在自己.
以為可以.
得到的成就.
而是將自己放在上帝的手中.
讓祂.
不住地塑造我們.
為祂國度.
發揮更大的果效.
一個甚至是超乎你.
想像.
想求.
的果效.

$^{1281}$的位置.
的中間.
求主去幫助我們的生命.
都不住被上帝.
更新.
我們一個.
有眼淚的生命.
被上帝拆光的時候.
我們同樣.
可以繼續.
去為上帝去工作.
而也同樣可以.
有更大的.
力量.
去為主去作見證.
去到.
最後一段的時候.
在諾茲第八節.
諾茲第八節.
接著去第.
五節的經文.
當我們再讀.
「我是阿拉伯.
我是俄羅斯人」的時候.
是不是?.
大家都知道.
這個希臘文.
那個頭尾.
α ΩΓ的字母.
一切成了.
然後再說.
我要把生命的.
泉水.
無償地.
賜給那些口渴的人.
所以大家可以再對這個所謂.
α ΩΓ有一個新一點的提護.
我們.
就是最初.
或者我們以前的想法就是.

$^{1321}$這個.
可能就說上帝整個.
工作.
創造工作.
有始有終這樣.
但是如果接著去第四第五節說的時候.
其實是在說一些.
每一個上帝.
在我們生命裡面.
更新的工作.
上帝也是.
創始成終.
上帝開始得.
他就會有終結.
上帝是一個信實的上帝.
上帝的工作不會爛尾的.
就是上帝如果以我們生命當中.
去改變我們.
更新我們的時候.
上帝就是那個α ΩΔ.
上帝就是那個ΩΓ.
這個是.
可以說.
在啟示錄裡面.
上帝成了.
那件事情.
除了說一個整個.
大的.
圖畫.
上帝要更新改變之外.
其實裡面.
到了第四第五節再接著下來的時候.
那個意思可以是一個.
對於我們每一個生命.
上帝都有一個計劃.
上帝都有一個心意.
上帝都會在我們生命當中有開始.
上帝都會在我們生命當中有它結束的時候.
他的工作.
是信實不變.

$^{1361}$特別重要是賜給我們.
更新我們生命之後.
他賜給我們活水.
讓我們學了.
永遠不去再學.
然後.
他就說到.
得勝的.
就會承受他自己的國度.
如果他是一些.
膽怯的.
不信的等等.
他就會.
是背棄上帝的人.
就會有他自己.
上帝給他們的結果.
我們能夠得勝.
不同世界妥協的時候.
堅持上帝為首位的時候.
要承受一切.
上帝賜下的.
亦都是成為神的兒女.
亦都是上帝賜給我們的.
一份的福分.
一份的滿足.
活水很重要.
我們每一天.
都需要有水.
都需要活水.
活水帶給我們.
生命的滋養.
生命的.
那種在重重複複.
形形異異的裡面.
我們能夠見到上帝.
的一些亮光.
一些生命的滋養.
讓我們能夠重新得力.
這就是上帝的活水.
上帝的道.

$^{1401}$對他們生命的意思.
一個神學家莫特曼.
他有一本書叫做.
《上帝的來臨》.
是寫一個終末論.
我覺得他的.
說法.
和啟示了最後這段經文.
都是一個很大的共鳴.
莫特曼說.
上帝的來臨.
其實重要的.
不是結束.
不是要結束一些.
舊世代的問題.
而是一個新的開始.
所以說是一個新的創造.
一個新的創造.
當然.
沒有罪,沒有痛苦,沒有痛.
那些.
好像啟示所所說的.
都是重要的.
但是.
我們更加不可以忽略的就是.
當剛才說了那種.
擦乾眼淚.
一切都更新的時候.
生命是.
經歷一個重新的創造.
生命.
的獨特性.
和生命的尊嚴.
最終在上帝的.
這種的.
來臨裡面.
是被更新.
來臨不單止是要解決.
我們.
過去一直.

$^{1441}$面對的問題.
那些困難.
那些掙扎.
上帝的來臨更加重要的是.
讓我們生命得到滿足.
得到力量.
讓我們能夠再出發,再上路.
去走下去.
所以可以說其實這個的.
上帝的來臨.
上帝的那種的.
終末的更新.
對生命的再次的創造.
其實不是一個終點.
我們很難去想像.
但是呢.
我只能夠體會到就是一種.
我們原來生命是不住地重新得力.
重新看見到新的可能.
新的方向.
新的目標.
這個就是在啟示道到最後.
或者莫特曼在書裡面所說.
讓我們去再對終末.
一種重新的體會.
到了結語的時候.
拍拍我們.
等於節目我們在地上.
作為上帝的帳幕.
做人間的時候.
能夠在世上成為.
一個很確實的.
有神同在的群體.
看見我們.
就看見上帝.
在當中.
在人的中間.
神母停止了她的工作.
甚至是不住地.
更新一切.

$^{1481}$在我們生命當中.
在我們跟人接觸的裡面.
在我們身邊的人的生命當中.
祂會繼續.
當中祂有祂自己的作為.
祂沒有停止過.
也不會停下來.
祂不會覺得疲乏.
不會困厥.
我們每一顆的眼淚.
祂都會顧念.
祂都會擦乾.
祂都會看顧.
我們.
上帝是一位很細心.
無微不至的上帝.
同時間我們不能夠.
或者不應該被罪惡所剩.
這也是啟示了當時.
這群信徒堅守的信仰.
就是堅守忠信.
最終上帝賜給他們.
承受生命的冠冕.
而且我們主觀地.
能夠享受到.
生命的滋養.
和滿足.
這就是上帝賜給我們的活水.
讓我們生生不息.
繼續重新得力.
那個狀態.
以致我們成為祂.
一個新的創造.
生命會繼續有新的可能.
新的嘗試.
發揮更大的效用.
更大的作用.
特別是在上帝的國度裡面.
弟兄姊妹.
我們面對著不同的衝擊.

$^{1521}$不同的問題.
我們如何再為我們自己定位呢?.
有時候.
有時候被那個世界.
改變著我們.
或者是.
在動我們.
有時候我們好像很被動.
好像在應接不同的.
挑戰和考驗.
但是我們回到上帝裡面.
其實我們是能夠.
在上帝裡面自己.
為自己定位.
能夠處變不驚.
去面對一切.
從外而來的挑戰.
我們教會群體.
在當中有一個什麼的使命呢?.
我們以臣為忠心.
堅持這種生活.
和這樣的群體.
有什麼意思呢?.
盼望啟示裡面的訊息.
特別是最後這章.
這段經文的訊息.
能夠給予我們一個亮光.
讓我們體會自己.
就是一個忠目的群體.
祭上帝的帳幕.
要在今天香港這個社會當中.
讓我們能夠成為了.
人看見就看.
神的帳幕.
就在人的中間.
我們一起在這裡禱告.
【天下無二人】.
七天會多謝你.
你有沒有撇下過.
你的兒女.

$^{1561}$我們的生命.
我們的生活.
在這裡堅持的時候.
你都是一一的看顧.
一一的保守.
更加重要的就是.
你說透過我們.
就好像是你的帳幕一樣.
能夠在人世間當中.
彰顯你的榮耀.
顯明你那種.
對生命的重視.
將你的愛.
將你的憐憫寬恕.
帶到去.
可能充滿著很多.
仇恨.
艱難的.
人的中間.
願你親自的.
去.
讓我們今天.
藉著啟示徒經文裡面.
給予我們一個.
心的安慰.
你拆去我們每一顆的眼淚.
你根深不住.
根深我們的生命.
以致我們能夠在你當中.
得到生命的滿足.
得到你的活水.
你讓我們成為德性的一群.
主要求你當中親自的.
讓我們能夠.
很深的體會.
並不是靠著我們的力量.
而是你親自與我們同在之下.
我們這群人.
特別在今日這個.
幻變的社會當中.

$^{1601}$我們怎樣可以堅持.
堅守信仰.
亦都處變不驚.
不會過分的恐懼.
過分的擔憂.
心裡面.
相信你是那位.
創始成中.
你試過.
你試阿拉伯.
你試阿摩加.
你成就你一切的心意.
你不會說出你的話.
它是不成就你所說的一切.
求你親自的.
恩代.
憐憫.
寬恕.
祝福我們每一個.
祝福紅印堂的群體.
聽我們同心向你禱告.
逢救主耶穌.
得勝榮辭而求.
阿們.
\newpage



\section{約翰福音 1:6-8-15}
\label{sec:0IxLPtCeZ88}
\textbf{2022.01.23《見證》 約翰福音 1\_6-8,15 (和修版) 曾裔貴牧師}
\newline
\newline
連結: \href{https://youtube.com/watch?v=0IxLPtCeZ88}{\texttt{https://youtube.com/watch?v=0IxLPtCeZ88}} ~~~~ 語音日期: 2022-01-26
\newline
\newline
\hyperref[sec:biU1pOyNb74]{\small{< < < PREV SERMON < < <}}
~
\hyperref[sec:index_chronic]{\small{[返順時目]}}
~
\hyperref[sec:index_scriptual]{\small{[返順卷目]}}
~
\hyperref[sec:lmTkPW75UmA]{\small{> > > NEXT SERMON > > >}}
\newline
\newline
約翰福音 1:6-8-15
\newline
\begin{longtable}{cl}
\hline
\hline
章節 & 經文 (和合本修訂版)\\
\hline
1:6 & \begin{tabularx}{0.7\textwidth}{X} 有一個人，是從神那裡差來的，名叫約翰。 \end{tabularx} \\ \\ \relax
1:7 & \begin{tabularx}{0.7\textwidth}{X} 這人來是為了作見證，是為那光作見證，要使眾人藉著他而信。 \end{tabularx} \\ \\ \relax
1:8 & \begin{tabularx}{0.7\textwidth}{X} 他不是那光，而是要為那光作見證。 \end{tabularx} \\ \\ \relax
1:9 & \begin{tabularx}{0.7\textwidth}{X} 那光是真光，來到世上，照亮所有的人。 \end{tabularx} \\ \\ \relax
1:10 & \begin{tabularx}{0.7\textwidth}{X} 他在世界，世界是藉著他造的，世界卻不認識他。 \end{tabularx} \\ \\ \relax
1:11 & \begin{tabularx}{0.7\textwidth}{X} 他來到自己的地方，自己的人並不接納他。 \end{tabularx} \\ \\ \relax
1:12 & \begin{tabularx}{0.7\textwidth}{X} 凡接納他的，就是信他名的人，他就賜他們權柄作神的兒女。 \end{tabularx} \\ \\ \relax
1:13 & \begin{tabularx}{0.7\textwidth}{X} 這些人不是從血生的，不是從情慾生的，也不是從人的意願生的，而是從神生的。 \end{tabularx} \\ \\ \relax
1:14 & \begin{tabularx}{0.7\textwidth}{X} 道成了肉身，住在我們中間，充充滿滿地有恩典有真理，我們也見過他的榮光，正是父獨一兒子的榮光。 \end{tabularx} \\ \\ \relax
1:15 & \begin{tabularx}{0.7\textwidth}{X} 約翰為他作見證，喊著說：「這就是我曾說：『那在我以後來的先於我，因為在我以前，他已經存在。』」 \end{tabularx} \\ \\
[1ex]
\hline
\hline
\end{longtable}
$^{1}$主席 各位弟兄姊妹 主裡面平安 再一次在疫情雖然都是比較嚴峻 都代表著母堂和所有弟兄姊妹都送上在主裡面的祝福和文安 盼望我們大家身體心靈都是健壯 和大家一起有一個祈禱 跟著開始我們的正道 再一次我們等候主的聖靈.
來將聖經的話語教導我們 今天很想透過施洗約翰的見證和他怎樣去做一個見證人.
讓大家彼此得到免禮 更加盼望洪恩我們的子堂能夠在這個社區成為主的見證人 一個明亮的登台 叫人得到因為主而帶來的那個幸福感 而帶來的真的幸福 求主幫助 我們這樣的祈禱 奉耶穌基督的名 阿們.
相信不同的年代 大家都會有追求不同的偶像 明星 現在青少年這一代你們知不知道的 我也是剛剛才知道的 原來那個MIRROR很瘋癲的 有些小朋友12,13歲就已經學到很厲害.
他們跳的舞跟MIRROR的鏡粉跳得很像 我們這些年代 可能現在就是陶大宇的倒轉地球 或者我們的年代就是譚詠麟 張國榮那些的追求 不同的時代都會有不同的追求 約翰福音很特別.
剛剛主席已經讀了幾節 其實在這一章裡面 作者一開始一到第五節就說了主耶穌其實就是神的道 就是神藉著祂去創造萬有 祂本身就是神來的.
還有這個道就快跟人有接觸 就是來這個世界 但是很特別 第六第七第八節是突然間加入了一個人 第六節說有一個人名叫約翰 是從神那裡差遣而來的.
這個人來只有一個目的 就是要作見證 幫誰作見證 就是為光 就是上文所說的這位道 為光作見證 叫眾人因為他的見證就可以相信了 相信光就是這個道.
接著第八節 這個是我之前給教會的經文給漏了 作者 寫約翰福音的作者特別強調 特別強調 他不是那個光 不是那個主角 不是那個焦點.
所以千萬不要追求錯了他 他不是那個焦點 要追求的是他介紹的那個光 就是當然是耶穌 這裡告訴我們 祂不是拿光 乃是要為光作見證.
這裡三節聖經已經說了這位施洗約翰 他出現在這個世界 還有他有一個很清楚自己的身份 目的就是要為光作見證 就是說他是那個見證人 很重要了.
接著第十五節 其實我抽開了 給大家先有一個小小的了解 然後我等一下會再詳細跟大家說一說 第十五節就是他的信息了.
他說 約翰為他 為上文講的真光 這位道 這位主耶穌 他說約翰為他作見證 大聲的喊著 當時在曠野 在這個白大利約旦河外的地方來作見證.
待會我就會告訴大家一些背景 這個就說 也就算我曾經說 那在我以後來的 反成了在我以前的 因他本來就是在我以前了.
他不斷的說 有一位將要來的 我會來幫他施洗 但是其實他的身份在我以前的 所謂在他以前 根本是約翰 這個施洗約翰都還沒有出世 他已經存在了.
這樣來講這個好像很矛盾的道理 其實不是的 這裡告訴我們 待會我們要 第二個方面要留意 作一個見證人 去到最極致 最極致就是要發現你自己什麼都沒有.
你自己是要將這個主帶出來 讓人去得到他 那就可以得到真正的福氣 福音 所以這一大段聖經 不知道大家有沒有覺得 剛才我們沒有讀整段.
其實第六第七第八節 是特別的加上去的 因為其實 如果你看那個經文的上下文 一直都在講這個道 很快讀給大家聽 今次是沒有第六第七第八節的.
太初就有道 道與神同在 道就是神 道太初與神同在 再一次用另一個角度 講道與神的關係 萬物是藉著他做的 凡是被做的 特別強調 每一樣都是藉著他做的.
沒有一樣不是藉著他做的 即是反面的說 第四節 生命在他裡頭 這個生命就是人的光 接著光照在黑暗裡 即是說這一位的主 與物質世界的關係 原來黑暗是拒絕接受光的.
接著我們撇開第六第七第八節 不讀了 留意第九節 立即就講這個真光 光是真光 照亮一切生在世上的人 他在世界 世界都是藉著他做的 再一次講第三節那裡.
但是強調的就是 他創造的萬有已經拒絕他 不認識他 甚至這個拒絕是一種黑暗的拒絕 這裡告訴我們世界卻不認識他 接著第十一節.
將世界的範圍收窄到一個地步 就不是一個普通的泛指的世界 是講他自己的地方 即是猶太人的地方 他到自己的地方來 自己的人倒不接待他.
接著有一個盼望和一個很重要的一個信息 就是這位作者說 但不是的 有一些人會接待的 所以這裡其實就是回應六七八節所講的 必須有見證人 見證真光.
所以第十二節 凡是接待這個光的人 就是相信他明的人 他就賜他們權柄作神的兒女 而且還強調這個作神的兒女 那種命定神的旨意.
這等人 即是這些接待耶穌的人 不是憑血氣而生的 不是從情慾而生的 情慾即是男歡女愛 即是肉身的關係 也不是從仁義去生 而是從神而生.
即是說如果你帶領一個人去接待耶穌 你就是那個終生的見證人 你就可以令到他從神而生 成為神的兒女.
然後十四節就講道成了肉身 所以到底是接待誰呢? 就是接待耶穌 道成了肉身 住在我們中間 充充滿滿的有恩典 有真理 我們都見過他的榮光 正是負獨生子的榮光.
所以這一段聖經 原來很重要的焦點 整大段一到十四節都是講主耶穌 這位的道 怎樣成為耶穌這個人.
但是很特別 很特別 就是第六 第七 第八節 居然插入了一個人 透過他傳遞給人知道這個光 或者這個道的重要 人必須要接待他 就可以得到這個生命.
所以今天很想從這個角度 這個焦點 和大家思想一下 教會存在在地上 我們是見證人 我們要有一個見證 個別的信徒 我們都是一個見證人 我們去見證神的真實和大愛.
有一個這樣的故事 馬雲 阿里巴巴的創辦人 其中一個創辦人 他們有十八個創辦人 後來當然很成功 他們開始需要有很多組織裡面的培訓.
他們請了一些訓練師 可能是一些促銷 去訓練他們的僱員 或者培訓班 道師有一次在班裡面 可能是一些經理級等等 去培訓.
就說 現在我給你們一個題目 你們怎樣去做一個推銷 現在有一把梳 譬如我們的工廠是製造梳的 你怎樣可以成功賣給和尚 賣給和尚 他賣梳 馬雲經過聽到他這樣 我不知道是真是假 好像是真的.
就馬上解僱了這個導師 不可能的 你帶給他們的創意不是我們公司的文化 我們公司的文化不是這樣的 不應該將一些沒有需要的來製造成為需要 那個和尚根本就不需要外梳 為什麼你要賣梳給他 是錯誤的一個解釋.
所以這個命題是不對的 如果這樣 你朝這個方向訓練我的僱員 可能他們就會走向旁門左道 即是說務求達到目的 什麼都要做 馬雲我覺得這一點很有意思.
沒有需求你製造需求 這個就不是真正的營銷 即是他們的價值就是那個地方有需要 我們提供到那個服務的需要 他就能夠得到那個福氣 這個就是阿里巴巴的價值觀 這個就是真正的銷售.
所以這個故事在網上都很火 一個這樣的見解 兄弟姐妹 今天我們自己感覺到在耶穌裡面 我們自己有那個真正的幸福 我們自己有那個心靈的平安.
而我們要將這個在世人可能還未覺得是需要的那個心靈 我們提供給他 當然不是叫銷售 我們要去見證 讓他們去認識 怎樣能夠讓這個幸福感帶給其他未信主的人呢.
這裡的應許就是說凡是接待這個道 這一位的耶穌 就是相信他明的人 即是說耶穌這個寶貴的名字 如果人明白了 求告就可以得到神永恆的生命 神兒女的身份 可以是這麼大這麼嚴肅的事來的.

$^{41}$所以我們要製造這種屬靈的需要給未信的朋友 首先就是我們自己要有這個真幸福 約翰很認識他自己 作者說有一個人 不是神化他 是普通人來的 是從神那裡差遣而來的.
奉差遣 奉差遣 約翰其實他是大祭司撒迦利亞和以利沙白的兒子來的 其實他一出生就是大祭司 他是在耶路撒冷那些宗教的群體那裡成長的.
他沒有需要自己走到那個曠野來去在那裡施洗 叫人來去悔改 和他的身份 即是生下來的身份是不匹配的 但是他居然自己走到曠野來做這個施洗者 為什麼呢?.
因為有一個人是由神那裡差遣的 感覺到自己這個身份是有一個新的使命 所以就放下自己那個祭司的身份 放棄祭司的尊榮的階層 來將自己放在一個傳道者的見證人的一個這樣的身份 自己就去到曠野.
在其他的福音書 馬太 馬可 盧迦都有記載 他就地在那裡來生活 穿的是駱駝毛的衣服 吃的是蝗蟲野蜜.
就地的很簡樸的生活 穿著要見證那個快要來的那一位的道 那個真光 所以這個是創造別人的幸福 自己放下自己的那個安書 尊貴的祭司身份 來自己走到去有需要的地方.
神安排他的那個地方 就是他的工場 或者叫教會 或者是他的港台 他就是要見證就快有一個人 我要幫他洗 幫他施洗 但是他是在我以前的 我幫他解鞋帶都不配的 這樣的施洗約翰.
所以教會 洪恩的弟兄姊妹 特別是在線上的每一位弟兄姊妹 我們要創造別人的幸福 我們自己不能夠製造 但是我們是一個見證人 見證耶穌基督在其他人的生命裡面 當有結合 求告 接待耶穌的時候.
是的 全世界都是不信耶穌 拒絕的 來到這個世界 世界卻不認識他 所以肯定福音是不受人歡迎 不是一個人見人愛的 因為人還沒看到背後自己的需要和神的偉大.
所以我們這個見證人 去將耶穌讓人認識 製造那個幸福感 我在網上看到一個故事 都很感動 有一間公司有300個員工 在美國的 300個員工 每年的聖誕節 他們都會有一個這樣的規矩 傳統 就是每人拿10元美金出來.
然後寫下自己的名字 當然那個作為一個儲備 就放進去那個抽獎箱 寫自己的名字 到時候因為聖誕聯歡 就抽獎 抽到誰呢 因為300個員工 每人10元美金 那就有3000元 那個中大獎的就可以獲得3000元美金.
那都很開心的 在聖誕節聯歡 那個故事告訴我們 真人真事 有一個男同事 知道最近有一個清潔的醫醫 他的兒子進了醫院 很急需要錢來幫他.
這個男生一直想著想著 雖然那個抽中的機會很渺茫 他都覺得總是有一個機會 就將他自己的名字寫了那個清潔醫醫的名字 那就有2/3的機會 他覺得都很少的了 沒可能的了 不過盡力幫一下.
就將自己的名字改了是他這個清潔醫醫的名字 希望有奇蹟出現了 去到抽獎的時候 相信大家都會知道那個主持一抽就宣讀那個名字 正正就是這個清潔醫醫的名字.
他很開心 當場就流眼淚 他說這3000元美金對我來說是很重要 這個聖誕節我經歷到神蹟 大家都很開心 那個年輕人都很感動 有沒有可能3/2會抽到他呢 沒多久就結束了 去看看那個抽獎箱.
拿了一張 拿了一張 很多都是那個清潔醫醫的名字 原來很多的同事都覺得都是很想那個醫醫能夠拿到這個安慰和這個獎金 就是這個抽獎的傳統的福利.
這些都是人為的感動 可以帶來好像是一個神蹟 人製造的神蹟 但是我們作為見證人 我們在教會裡面 我們自己必須真的經歷主耶穌在我們生命裡面帶來的幸福 帶來的生命改變.
我們告訴別人 我們真的經歷過耶穌的愛 祂的拯救 祂的醫治 祂對我們的饒恕 和我整個生命的建立 有因為我是神兒女而有的尊榮 尊榮這份honor 有時候會令一個人有很大的改變.
我又見過一個這樣的故事 有一個推銷原子筆的推銷員 一路他的業績不好 在街上很惆悵 怎麼辦呢 我的業績這麼差 剛好有一個有錢的做生意的人經過 見到他這麼落魄 就問他什麼事 我是推銷原子筆的 但是我一直都賣不到.
這個人就說20塊一支 我幫你拿著吧 給你20塊 然後就走了 這個人就很開心 走了沒多久 做生意的那個有錢人就回頭了 我給你20塊 我忘記拿回原子筆了.
然後就說了一句話 你和我都是做生意的 你都會是一個很成功的商人 就拿了原子筆就走了 這個推銷員被這個有錢人 一個路邊不認識的人 一番的安慰勉勵的說話 說你是一個成功的商人.
這樣就令到他往後的努力奮鬥 過了一年之後 這個推銷員居然在一個小小的商場裡面 開了一家店鋪 就是說他已經由一個很老的推銷員 慢慢因為他的努力 去建立了自己的小小的事業.
剛好一年後 這個做生意的人又走過了 這個推銷員這次認得他了 就跟他說 你知不知道你當年 就是去年 你不是給了我20塊 你是給了一份尊重我 其實我知道我自己根本就不成功.
但是你這一番的說話 你說我和你一樣 都是兩個做生意的人 你是會成功的 哇 這樣的一個勉勵 就帶給這個推銷員 慢慢慢慢越來越有力量的去建立自己的自信 而帶來有一種自己的生命的改變.
親愛的弟兄姊妹 我們在教會裡面 我們來對別人有沒有能夠隨時說 造就別人的安慰的話呢 叫聽見的人 不是那些阿諛奉承 不是那些很假的說話 是叫聽見的人得到益處 隨時說一些可以造就他的說話.
造就可以是善意的勸勉 可以是 譬如我們長輩對年輕人 如果是熟的 有關係的 可以直接的去指出 你這樣做可能會對你有影響的 那這就是善意的隨時說造就的好話 是叫聽見的人得到益處 就是說他自己 是喔 他可以反思.
他可以來到有一個發現 然後他有一個改變 那就會帶來那個人會有一個新的轉機 我們很需要求神 讓我們能夠在教會 看到我們能夠行使 運用的那個權柄來自三一上主 還有我們能夠可以帶給很多人 得到從神那裡而來的真正的幸福.
弟兄姊妹 這個世界很少有一個地方 是這麼密集 每個禮拜大家聚在一起 來自不同的背景 不同的階層 大家可以都是有同樣的神的兒女的生命 大家可以彼此的來到去唱詩歌 團契 講道 彼此的勸勉 又有交易的時間 很少有這樣的群體的.
很少大家是一信耶穌 大家都是一個很單純的心靈 這樣來到去相交 是沒有任何的企圖的 你想一下有哪個地方的人會是這樣呢?.
所以教會是神安排我們的一個場景 我們要為主作見證 將人帶進那個真正的幸福裡面 我們要將人帶進那個神的救恩的福樂裡面.
我們看看下面 19節之後 有一個調查團下來調查 為什麼你這麼受歡迎 為什麼你在那個地方這麼荒謬的曠野 來做一個施洗者 到底你是不是基督呢?.
約翰很清楚 第20節 他就明說 並不隱瞞的明說 我不是基督 這個不斷否認 很少的 你賣一件產品 不斷都要說你自己是怎樣怎樣的好 這樣才多人買 銷量才會高 不斷吹噓.
但一個真正在耶穌裡面的見證人 他做到最極致的就是 他認定自己什麼都不是 耶穌就真的是那一位 應許的彌賽亞 你去到他那裡 你就會得到生命的改變.
所以你會看到後面 他居然叫他的門徒跟耶穌 不要跟他 例如安德烈 例如這個肥力 拿丹葉 來跟耶穌 大家可能不知道那個背景.
當時施洗約翰 根據馬太福音說 耶路撒冷的人 由祭司的宗教階層到平民百姓 甚至一些稅吏 一些官長 兵丁 都去找約翰施洗 找他接受他的洗禮.
你想一下 是很厲害的一種熱潮 但是約翰居然間否認自己的身份是救世主 是基督 他只是那個聲音 他只是那一點光 真正的大光 真光 是這一位的耶穌.
是我將要幫他施洗的這一位耶穌 他是神的兒子 他是神的羔羊 他是那一位救主 所以這個真的很寶貴.
透過今天我和大家去彼此的勉勵 施洗約翰他的見證 可以帶給人有一個真正的幸福和生命.
所以他完全不介意他的忠心門徒 一個一個跟隨這一位的耶穌基督 因為他正正就是要為這一位的耶穌作見證.
所以他自己能夠說 他要越來越興旺 我約翰要越來越衰微 衰微到一個地步 最好就是盡快被人忘記.

$^{81}$這個是一個很高的情操 一個人不會將自己的功勞歸給自己 更加不會將那個名聲 將人帶到自己的面前.
所以很盼望無論錦宣家 無論是我們的洪恩家 或者我們的宣道會 在香港眾多的宗派 其中一個宗派.
我們都不是將人帶到自己這個宗派 或者自己的面前 而是好像約翰那樣 將人帶到耶穌那裡.
甚至乎要見證的 就是那個真光 就是這個道 就是這一位主耶穌.
所以很盼望在新的一年開始 牧者我們有機會和大家來彼此的勉勵.
我們現在在 我們剛剛堂董會開完會 我們都定了一堂三點.
很希望能夠在資源的共享上面 期望洪恩家能夠未來可以更加的茁壯成長.
一路步向那個獨立 能夠大家繼續可以在這個社區為主的真道去發光.
我們要很準確的見證的是主耶穌帶來那個人的真幸福.
不過首先都真的要檢視一下 我們自己是不是都感覺到 我們自己是在享用神給我們的福氣.
以致我們自己都真的有這個產品 就是耶穌基督 我們去介紹給人.
可以得到從主那裡而來的新生命 很盼望這個階段和大家彼此的勉勵.
我們一起祈禱 多謝主 今早能夠透過施洗約翰這一位忠實的見證人.
來將耶穌基督的身份 救贖來傳開 帶給當時的時代一個很大的震撼.
這一位的救主居然成為人 居然在當時的約翰的年代來出現.
成為人可以見到的其實就是神這一位造物主.
雖然現在疫情全球這樣爆發 但是我們仍然相信這個是天賦的世界.
造物主仍然掌管著萬有 叫我們大家有信心 有盼望.
繼續去將這個救人的福音 叫人得到好處 得到真的幸福.
\newpage



\section{約翰福音 4:3-26}
\label{sec:lmTkPW75UmA}
\textbf{2022.01.30 《耶穌領人歸順的奧妙之處》 約翰福音4\_3-26 (新漢語譯本) 楊穎兒傳道}
\newline
\newline
連結: \href{https://youtube.com/watch?v=lmTkPW75UmA}{\texttt{https://youtube.com/watch?v=lmTkPW75UmA}} ~~~~ 語音日期: 2022-02-03
\newline
\newline
\hyperref[sec:0IxLPtCeZ88]{\small{< < < PREV SERMON < < <}}
~
\hyperref[sec:index_chronic]{\small{[返順時目]}}
~
\hyperref[sec:index_scriptual]{\small{[返順卷目]}}
~
\hyperref[sec:aSFgvj6BhUY]{\small{> > > NEXT SERMON > > >}}
\newline
\newline
約翰福音 4:3-26
\newline
\begin{longtable}{cl}
\hline
\hline
章節 & 經文 (和合本修訂版)\\
\hline
4:3 & \begin{tabularx}{0.7\textwidth}{X} 他就離開猶太，又回加利利去。 \end{tabularx} \\ \\ \relax
4:4 & \begin{tabularx}{0.7\textwidth}{X} 他必須經過撒瑪利亞， \end{tabularx} \\ \\ \relax
4:5 & \begin{tabularx}{0.7\textwidth}{X} 於是到了撒瑪利亞的一座城，名叫敘加，靠近雅各給他兒子約瑟的那塊地。 \end{tabularx} \\ \\ \relax
4:6 & \begin{tabularx}{0.7\textwidth}{X} 雅各井就在那裡；耶穌因旅途疲乏，坐在井旁。那時約是正午。 \end{tabularx} \\ \\ \relax
4:7 & \begin{tabularx}{0.7\textwidth}{X} 有一個撒瑪利亞婦人來打水。耶穌對她說：「請給我水喝。」 \end{tabularx} \\ \\ \relax
4:8 & \begin{tabularx}{0.7\textwidth}{X} 因為那時門徒進城買食物去了。 \end{tabularx} \\ \\ \relax
4:9 & \begin{tabularx}{0.7\textwidth}{X} 撒瑪利亞婦人對他說：「你是猶太人，怎麼向我一個撒瑪利亞女人要水喝呢？」因為猶太人和撒瑪利亞人沒有來往。 \end{tabularx} \\ \\ \relax
4:10 & \begin{tabularx}{0.7\textwidth}{X} 耶穌回答她說：「你若知道神的恩賜，和對你說『請給我水喝』的是誰，你早就會求他，他也早就會給了你活水。」 \end{tabularx} \\ \\ \relax
4:11 & \begin{tabularx}{0.7\textwidth}{X} 婦人對耶穌說：「先生，你沒有打水的器具，井又深，哪裡去取活水呢？ \end{tabularx} \\ \\ \relax
4:12 & \begin{tabularx}{0.7\textwidth}{X} 我們的祖宗雅各把這井留給我們，他自己和兒女以及牲畜都喝這井裡的水，難道你比他還大嗎？」 \end{tabularx} \\ \\ \relax
4:13 & \begin{tabularx}{0.7\textwidth}{X} 耶穌回答，對她說：「凡喝這水的，還要再渴； \end{tabularx} \\ \\ \relax
4:14 & \begin{tabularx}{0.7\textwidth}{X} 誰喝我所賜的水，就永遠不渴。我所賜的水要在他裡面成為泉源，直湧到永生。」 \end{tabularx} \\ \\ \relax
4:15 & \begin{tabularx}{0.7\textwidth}{X} 婦人對他說：「先生，請把這水賜給我，使我不渴，也不用到這裡來打水。」 \end{tabularx} \\ \\ \relax
4:16 & \begin{tabularx}{0.7\textwidth}{X} 耶穌對她說：「你去，叫你的丈夫，再到這裡來。」 \end{tabularx} \\ \\ \relax
4:17 & \begin{tabularx}{0.7\textwidth}{X} 婦人回答，對耶穌說：「我沒有丈夫。」耶穌說：「你說沒有丈夫是對的。 \end{tabularx} \\ \\ \relax
4:18 & \begin{tabularx}{0.7\textwidth}{X} 你已經有過五個丈夫，你現在有的並不是你的丈夫。你這話是真的。」 \end{tabularx} \\ \\ \relax
4:19 & \begin{tabularx}{0.7\textwidth}{X} 婦人對他說：「先生，我看你是一位先知。 \end{tabularx} \\ \\ \relax
4:20 & \begin{tabularx}{0.7\textwidth}{X} 我們的祖宗在這山上敬拜神，你們倒說，應當敬拜的地方是在耶路撒冷。」 \end{tabularx} \\ \\ \relax
4:21 & \begin{tabularx}{0.7\textwidth}{X} 耶穌對她說：「婦人，你要信我。時候將到，你們敬拜父，既不在這山上，也不在耶路撒冷。 \end{tabularx} \\ \\ \relax
4:22 & \begin{tabularx}{0.7\textwidth}{X} 你們所敬拜的，你們不知道；我們所敬拜的，我們知道，因為救恩是從猶太人出來的。 \end{tabularx} \\ \\ \relax
4:23 & \begin{tabularx}{0.7\textwidth}{X} 時候將到，現在就是了，那真正敬拜父的，要用心靈和誠實敬拜他，因為父要這樣的人敬拜他。 \end{tabularx} \\ \\ \relax
4:24 & \begin{tabularx}{0.7\textwidth}{X} 神是靈，所以敬拜他的必須用心靈和誠實敬拜他。」 \end{tabularx} \\ \\ \relax
4:25 & \begin{tabularx}{0.7\textwidth}{X} 婦人對他說：「我知道彌賽亞—就是那稱為基督的—要來；他來了，會把一切的事都告訴我們。」 \end{tabularx} \\ \\ \relax
4:26 & \begin{tabularx}{0.7\textwidth}{X} 耶穌對她說：「我就是，正在跟你說話呢！」 \end{tabularx} \\ \\
[1ex]
\hline
\hline
\end{longtable}
$^{1}$弟兄姊妹平安.
多謝洪恩堂堂主任魯牧師的安排.
讓我可以在洪恩堂和弟兄姊妹分享聖經的訊息.
每一次回到洪恩堂都很有回家的感覺.
剛才讀出的經文是來自約翰福音第四章三至二十六節.
如果大家家裡有聖經的話.
請大家去預備.
我在預備經文的時候.
我發現在經文裡面.
耶穌和婦人的對答當中的內容是非常豐富.
能夠幫助弟兄姊妹去領略傳揚福音.
令人歸主.
參與主在地上復和的工作.
很有幫助.
所以我決定今天的講道題目是.
耶穌領人歸順的奧妙之處.
讓我們先禱告.
親愛的天父.
求您來到我們中間.
張開孩子的口.
打開我們各人的心.
讓我們明白你的話語.
學習耶穌如何領人歸順.
豐富我們傳揚福音的技巧.
祈禱奉主並求.
阿們.
今天我們要看的故事.
是耶穌傳福音給一個撒瑪利亞婦人的故事.
這個故事的結局是很震撼的.
很多人聽了婦人的見證.
以及耶穌的教導之後.
因而相信耶穌是救世主.
但是今天的講道的重點並不是這個結局.
因為一個人歸順主耶穌與否.
是三一神的工作.
今天我們會看看耶穌帶領這位撒瑪利亞婦人.
歸順祂的過程.
婦人與祂的對答過程.
然後我們在當中去明白主要我們學習的功課.
撒瑪利亞這個地方是位於耶路撒冷.

$^{41}$與猶太人之間.
由猶太人向北走到耶路撒冷.
最快要經過撒瑪利亞.
但是通常來說猶太人是不會經過撒瑪利亞的.
他們會沿著東面的約旦河而走上去.
就算猶太人走著走著.
他們遇到撒瑪利亞人.
他們也會繞過撒瑪利亞人.
甚至不讓自己踩到撒瑪利亞人的影子.
原因是猶太人認為撒瑪利亞人的血統不純正.
這個問題就要回到歷史裡面.
要追溯到猶大王阿哈斯的年代.
大概是公元前八世紀.
當時亞述王攻擊以色列.
佔領了撒瑪利亞.
亞述王就遷移了一些外族.
住在撒瑪利亞的城邑.
與原有的以色列人同住.
意思是這些撒瑪利亞人.
要與外族人一起共同生活.
撒瑪利亞人不單止血統不純正.
連信仰也是不純正的.
根據《列王記下》第十七章記載.
亞述王遷移到撒瑪利亞的外族.
他們都是為自己製造神像.
並安置在撒瑪利亞人所造的幽壇.
即是在一些殿裡.
在撒瑪利亞的中間也有一些祭司.
亞述王去差派他們.
在撒瑪利亞人當中的外族裡.
如何敬畏耶和華.
所以我們知道他們有外族.
有撒瑪利亞人 又有祭司.
在他們信仰的過程裡.
其實他們是不遵守耶和華.
吩咐雅各后裔如何聽從這個律法.
他們的信仰一方面.
他們認識耶和華.
又懼怕耶和華.
但另一方面他們又事奉自己的神.

$^{81}$有時他們去到哪裡居住.
就拜到哪裡.
就好像有拜錯.
沒有放過一樣.
這是撒瑪利亞人的背景.
約翰福音的故事很特別.
作者經常用一些記號去寫.
因此我們有時要明白.
約翰福音時就要像解碼一樣.
拆掉這些記號.
經文提到耶穌接觸了一個.
曾經有過五個丈夫的婦人.
這個婦人簡直是世間難尋.
多難才找到一個婦人.
嫁了五次呢?.
我們也會想想這個婦人是怎樣的呢?.
這個婦人的道德觀念又是怎樣的呢?.
可能我們會立即聯想到.
她必定是一個道德敗壞.
不忠於她丈夫的婦人.
沒錯,這就是經文要表達的意思.
如果我們看回舊約.
經常會發現.
經文會描寫到行淫邪的婦人.
她會背著丈夫去行淫.
婦人又會對丈夫不忠等等的異警.
而我們去理解這個婦人肉體上的不忠.
她道德的敗壞是其中一個解釋的方法.
另一個解釋其實是指.
她們靈性上的不忠.
所指的是人在信仰裡背棄了神.
去崇拜別神.
在舊約裡面.
神經常去審判背棄神的子民.
他們去拜其他別神.
當神看到這些人去拜別神的時候.
發現這些人會花費精神和時間.
所有資源去討好這些別神或偶像.
神看到這些人的興奮程度.
那種投入的程度.

$^{121}$神就好像看到自己所愛的人類.
和其他人進行不道德的行為一樣.
所以在今天的故事裡.
耶穌指撒瑪利亞婦人曾經有五個丈夫.
婦人這個特徵其實是代表著她不忠的記號.
不單止是婦人的不忠.
亦是整個撒瑪利亞民族的不忠.
為何是五個丈夫而不是六個七個丈夫呢?.
原來在聖經的歷史發展裡.
我們可以找到這個代表著撒瑪利亞民族的五個丈夫.
在《列王記》下第十七章二十四節所寫.
亞述王從巴比倫,古他,阿瓦,哈瑪和西伐瓦欽而來.
這些都是外族.
安置在撒瑪利亞的城邑.
與以色列人同住.
他們各自建立神像和祭壇.
祭師在這些壇上獻祭.
所以這五個丈夫其實是指這五個外族.
巴比倫,古他,阿瓦,哈瑪和西伐瓦欽.
如果到經文去看.
其實他們拜很多別神.
有不同的名字.
撒瑪利亞人原本都是以色列人.
而且他們的先祖是雅各.
代表著他們與猶太人出自同一條根.
但經過公元七八多年與外族同住的歷史.
撒瑪利亞人的信仰已經被外族同化了.
所以一直以來猶太人對撒瑪利亞人都存有極大的偏見.
而撒瑪利亞人亦長期住在那裡.
長期受壓,被邊緣化.
他們亦不會主動與猶太人有甚麼來往.
所以婦人中午出去打水的時候.
可以指她是一個不合乎道德的婦人.
要在沒有人打水的時候才出來避開眾人.
又或者我們可以說.
這個婦人是寓意著撒瑪利亞是一個當時被社會排擠的民族.
講完記號.
很想跟大家看看這位撒瑪利亞的婦人與耶穌的對答.
我們會發現這個婦人不是一個普通的婦人.
第一,這個婦人對她一直以來所信奉的.

$^{161}$其實她是真心覺得很清楚的.
為何這樣說呢?我們來看看經文.
在撒瑪人的理解他們的信仰的時候.
他們是有一些祭祀的儀式.
所以當耶穌去跟婦人說.
你有沒有.
所以耶穌去跟婦人說.
祂有活水可以給這個婦人的時候.
這個婦人就立刻留意到耶穌沒有帶工具.
怎樣打水呢?.
意思就是其實耶穌在傳福音的時候.
這個婦人問祂:你裝香拜神的工具都沒有.
你哪裡有這些福音,福氣給我?.
不僅如此.
婦人還要證明她自己所信的是雅各留下來的傳統.
這個傳統是正統的信仰.
於是她更加跟耶穌說.
雅各自己,雅各的子孫,他們的生畜.
全部人都是這樣做的.
都是喝這個水的.
由此可見.
這位婦人對她一直以來所信奉.
她很熟悉的.
而且很堅持她所信的信仰是真的信仰.
第二.
婦人不僅知道她自己信奉什麼.
她還知道猶太人所信奉的習俗.
所以當耶穌問她拿水的時候.
第九節她向耶穌說.
因為猶太人和撒瑪利亞人不會用共同的器皿.
意思就是婦人知道她所信奉的習俗.
和猶太人是不同的.
大家有不同的寺廟,儀式,祭祀的方法.
就好像今天來說.
有些宗教會裝香.
有些宗教會點蠟燭一樣.
婦人的意思就是.
你們還你們.
我們還我們.
是她先說的.

$^{201}$你們猶太人.
你們說是耶路撒冷.
耶路撒冷才是人要去敬拜的地方.
第三.
婦人在哪裡去敬拜.
都很有了解.
婦人在第20節提到.
他們的祖先在這個山上敬拜.
這個山是指基利心山.
這個山的歷史很悠久.
昔日摩西分庫以色列人到了迦南地之後.
就要在這個基利心山為百姓祝福.
基利心山才是蒙神賜福的福地.
婦人的意思就是說.
自我們先祖開始.
我們已經在這個山上敬拜.
現在耶穌來到.
你知道是猶太人.
所以她更加跟耶穌說.
我都不明白為什麼你們猶太人要去耶路撒冷敬拜.
用今天的語氣來說.
其實婦人都很尊重耶穌.
就說我都知道你是一個先知.
但你跟我說敬拜.
你都不知道過去.
無論是敬拜的方式.
歷史的承傳.
祭祀的地點.
婦人都很能夠跟耶穌理論.
甚至提起這件事跟耶穌理論.
由此可見這個婦人.
很可能在當時她自己的民族宗教來說.
她很熟悉.
甚至她是敬虔的一群.
所以後來第39節說到.
當婦人聽完耶穌說話之後.
她進到城裡去講見證的時候.
城裡的人都很相信她的話.
如果這個婦人是代表著.
整個撒瑪利亞民族的宗教認知.

$^{241}$我們可以推斷.
整個城裡的撒瑪利亞人.
都很熟悉這些過往的歷史.
因此當這個女人去到城裡去講見證.
當這些人去見到耶穌的時候.
知道過去的歷史.
都願意去相信耶穌.
原來真是他們一直等待的救主.
不過雖然這個婦人.
很虔誠於她自己信奉的宗教.
也並且引以為傲.
然而當耶穌叫她帶她現在的丈夫來的時候.
這個婦人竟然說沒有.
這個婦人自己很熟悉自己民族的信仰.
但為什麼會說沒有呢?.
原來是她一回到自己心靈的深處的時候.
她不能夠將她所信奉的信仰.
連繫到她的靈魂.
中世紀的神學家奧古斯丁.
曾經這樣去形容約翰福音這個撒瑪利亞人.
提到這個撒瑪利亞人的頭腦.
他對信仰的理解是很充足.
甚至支配了他內心的靈魂.
使得他的靈魂根本不能夠.
從他平時的敬拜的模式來得到滿足.
於是.
第一,耶穌會糾正這個婦人對敬拜的理解.
指出敬拜的重心不在於地方.
不是基里森山.
不是耶路撒冷.
第二,耶穌又從歷史傳承的角度去告訴這個婦人.
指出撒瑪利亞人敬拜的神.
是歷史遺留下來的神.
是外邦的神.
撒瑪利亞人又怎會認識他們呢?.
認識的話,只不過是頭腦的知識.
缺乏靈裡面的敬拜.
第三,耶穌繼續說到救恩不是來自摩西或雅各.
但是現在要相信是耶穌自己.
真誠的敬拜是要以靈,以真理去敬拜.

$^{281}$耶穌就是看到撒瑪利亞人心中的空虛.
他很關心婦人心中沒有靈的那種信仰.
心是多麼的空虛.
以靈,以真理的原文.
就是形容關於聖靈.
就是強調聖靈的工作.
是聖靈去引領人歸順於祂.
同時,敬拜神的人要在聖靈裡面去認識耶穌基督.
當耶穌說到要以靈,以真理去敬拜的時候.
這正是撒瑪利亞人內心最需要的東西.
耶穌親身展示了祂如何去引領其他有宗教信仰的人歸順於祂.
我發現了有四方面可以和大家一起在傳福音或宣教領域上一同學習.
第一,以靈,以真理敬拜.
第24節是這樣說的.
神是靈,敬拜的人必須以靈,以真理來敬拜.
剛剛提到,靈,真理就是強調聖靈的工作.
有時我們在不同的譯本裡.
會看到他們是用了「以心靈和誠實去敬拜」.
這關乎到在譯聖經裡.
譯者對經文有不同的理解.
有不同的翻譯過程.
但我們有時去理解心靈誠實的標準在哪裡呢?.
我們何時才可以肯定自己是誠實的?.
或是我們如何用自己去解釋才能理解到心靈的誠實呢?.
如果是這樣的話,每個人的敬拜都很視乎敬拜者的生命質素.
所以我今天揀選了漢語新譯本.
它用「以靈,以真理敬拜」.
所說的是強調聖靈的工作.
關鍵不是我們有多努力.
關鍵不是我們在哪個地方.
不是基里森山.
不是耶路撒冷.
敬拜也不是關乎我們一個星期有多少天很努力地回來教會.
敬拜當然不是我們像薩瑪尼亞民族那樣.
要靠一些祭祀的工具,祭祀的儀式去敬拜.
敬拜乃是我們有沒有認定神的靈.
是與耶穌同工.
此時此刻,在我們的生命裡.
我們有沒有與神連上關係.
我們的神是不斷很想去尋回不認識祂的人.

$^{321}$因此在領人歸主的事上.
我們都要清楚知道.
神是有祂自己修葛的時間.
不是我們有多努力去做.
不是我們在哪個地方.
而是我們默默共同的.
以一個群體去尋求聖靈的帶領.
例如,如果神在你的心裡.
讓你去記掛著某一個人的時候.
你就去打個電話.
去發個訊息給他們.
問問他們近來怎樣.
或者我們去花時間跟他們聊天.
或者像剛才我們雲鳳姐妹.
在她的短訓課程裡學習.
我們先聽別人說話.
然後慢慢再去回答他們.
我們也可以在跟他們聊天的過程裡.
跟他們祈禱,聆聽他們的內心.
如果他們是不相信耶穌的.
我們也求主能夠讓福音去到他們的心中.
第二,突破地域界限.
第三節提到.
他是耶穌,他離開猶太再次進入加利利.
他必須經過撒瑪利亞.
剛才說過了.
如果從猶太去到耶路撒冷.
是要經過撒瑪利亞的.
這是最快的路途.
但是我們知道.
撒瑪利亞人是被孤立被歧視的一群.
很多猶太人都會妖道而行.
然而耶穌特意選擇一定要經過.
必須經過撒瑪利亞.
是他親自想去接觸這些撒瑪利亞人.
因此他必須要經過.
宣教很明顯是要進入其他地區去服事.
去進入那個地區去接觸那裡的人.
同樣在本地傳福音的事公上.
我們也有很多機會可以跳出自己的安書區.

$^{361}$跳出自己的領域.
不單是突破了地域的界限.
而且可以跨種族,跨語言,跨文化,跨宗教去服事.
就好像當耶穌向撒瑪利亞婦人取水的時候.
婦人很驚訝.
她說:你是猶太人.
你怎會向我這個撒瑪利亞人取水?.
偏偏就是藉著耶穌的舉動.
去喚醒了這個婦人.
去展開了一輪對話.
婦人去提出了一些提問.
從而幫這個婦人再去想多一些.
去深入討論她所相信的信仰.
第三,認識其他的信仰.
剛才在講道的時候也提到.
有些人會離棄神去崇拜別神.
為何這些人會離棄神.
去依靠別神.
其實背後每一個原因都不同.
又或者他們根本沒有聽過信仰.
我們說他也會遇到像撒瑪利亞的婦人.
她自己已經有自己的信仰.
而且她更加會見到像耶穌的猶太人.
信徒去挑起理論信仰的對話.
我也曾經遇過不同信仰的人.
約我去討論宗教.
我發現他們約我其實是有備而來的.
特別是一些無神論者.
因為他們要否定我們相信的神是怎樣.
因為他們要否定一些東西.
所以他們就知道有神的宗教究竟是怎樣去講這個神.
然後他們才可以否定這個神.
因此我們除了要學習聖經的知識之外.
我們也要認識其他的宗教信仰.
我們就可以更加有效地.
當我們遇到有這樣的機會.
我們可以幫助身邊的人.
去使他們的心靈得到平安.
去認識我們所相信的耶穌.
就是那位賜人重生的主.

$^{401}$所以很鼓勵大家.
無論在本地傳福音.
或者海外宣教方面.
都嘗試去認識其他的信仰.
是非常重要的.
除了認識之外.
第四,我們也要接納其他的信仰.
接納不代表認同.
有時我們很想去傳福音.
但我們也不知道從何而起.
但耶穌正正示範給我們看.
祂不僅叫了他的門徒去買食物,去做服事.
祂也邀請在信仰上不同的人.
去幫助耶穌自己.
第七節說到.
有一個撒瑪利亞婦人前來打水.
耶穌對她說.
「請給我水喝」.
因為她的門徒已經進城買食物去了.
我們留意到耶穌是如何接近這個女人的呢?.
文中記載.
當時耶穌是打發了門徒去買食物.
而且剛剛是第六時,即是正午的時候.
烈日當空.
耶穌也有祂的人性.
也會口渴.
於是祂就將祂的需要放在這個婦人面前.
希望這個婦人給祂水喝.
耶穌用這個方法.
會不會比一靠近就傳福音更有效呢?.
因為當我們去接觸一些非信徒.
或者接受他們進入我們的生命的時候.
我們就能夠有機會去互相認識,交流.
從而可以發現他們的心靈是怎樣.
也讓他們可以看到我們.
有耶穌基督在生命的特質究竟是怎樣的.
我們再看看.
耶穌是如何和這個女人去展開對話.
耶穌不是問這個婦人.
「婦人,你需要甚麼?」.

$^{441}$而是耶穌讓她知道耶穌的需要.
而由這個需要去決定對話的起點.
耶穌去求她,幫助她.
耶穌肯定了她的尊嚴.
有時人最深切的需要.
就是要讓別人肯定她的尊嚴.
人想感受到被別人尊重自己.
知道自己是一個有價值的人.
其實我們也可以邀請一些未信的朋友.
或者非信徒去參加我們教會的活動.
例如幫忙修理教會.
做一些義工幫忙.
讓他們能夠在這個群體中.
感到尊重和被認同.
聖靈向我們展示了這個很深刻的真理.
就是耶穌的愛.
耶穌願意用她的肉身.
這個身軀.
走到另一個地域.
去到一個撒瑪利亞婦人的面前.
為這個撒瑪利亞婦人帶來一線的希望.
並且移除了種族的障礙.
因著耶穌基督身體的需要.
耶穌向撒瑪利亞婦人傳福音.
打破了猶太人和撒瑪利亞人之間.
多年來的種族障礙.
甚至仇恨.
今天我們教會都是耶穌的身體.
通過我們.
我們就可以在社區那裡.
嘗試好像耶穌昔日.
打破和撒瑪利亞婦人之間存在的一些障礙.
所以領人歸順主的目的.
其實是要人和神得以復和.
家庭和家庭之間得到復和.
社區和社區之間得到復和.
民族與民族之間得到復和.
甚至宗教與宗教之間得到復和.
這是神所關心的.
讓我們一起祈禱.

$^{481}$親愛的天父.
我們滿心歡喜的來讚美你.
你就是那位蒼天造地.
滿有大愛的神父.
我們讚美你.
在你的創造心意裡.
你賜了耶穌給我們.
約翰福音一章九節說道.
那光是真光.
來到世上照亮每個人.
沒有人能夠躲避這個光.
這個光照遍大地.
照好人也照壞人.
照信你的人也照未認識你的人.
願你時常提醒我們.
讓我們生命為你作出見證.
聖靈啊.
求你內住在我們的生命裡.
教導我們各人.
去順從你的帶領.
聖靈你是生命的氣息.
自由的流動.
貫穿萬物.
讓我們明白.
我們的生命是時刻.
與你與不同的人一起連結.
我們這樣的禱告.
求主幫助.
乃是為我們犧牲生命的主而求.
阿們.
\newpage



\section{歌羅西書 2:6-23}
\label{sec:aSFgvj6BhUY}
\textbf{2022.02.06 《尊貴無比的基督信仰》 歌羅西書 2\_6-23 (和修版) 老耀雄牧師}
\newline
\newline
連結: \href{https://youtube.com/watch?v=aSFgvj6BhUY}{\texttt{https://youtube.com/watch?v=aSFgvj6BhUY}} ~~~~ 語音日期: 2022-02-06
\newline
\newline
\hyperref[sec:lmTkPW75UmA]{\small{< < < PREV SERMON < < <}}
~
\hyperref[sec:index_chronic]{\small{[返順時目]}}
~
\hyperref[sec:index_scriptual]{\small{[返順卷目]}}
~
\hyperref[sec:_xVz_wL7W7s]{\small{> > > NEXT SERMON > > >}}
\newline
\newline
歌羅西書 2:6-23
\newline
\begin{longtable}{cl}
\hline
\hline
章節 & 經文 (和合本修訂版)\\
\hline
2:6 & \begin{tabularx}{0.7\textwidth}{X} 既然你們接受了主基督耶穌，就要靠著他而生活， \end{tabularx} \\ \\ \relax
2:7 & \begin{tabularx}{0.7\textwidth}{X} 照著你們所領受的教導，在他裡面生根建造，信心堅固，充滿著感謝的心。 \end{tabularx} \\ \\ \relax
2:8 & \begin{tabularx}{0.7\textwidth}{X} 你們要謹慎，免得有人用他的哲學和虛空的廢話，不照著基督，而是照人間的傳統和世上粗淺的學說，把你們擄去。 \end{tabularx} \\ \\ \relax
2:9 & \begin{tabularx}{0.7\textwidth}{X} 因為神本性一切的豐盛都有形有體地居住在基督裡面； \end{tabularx} \\ \\ \relax
2:10 & \begin{tabularx}{0.7\textwidth}{X} 你們在他裡面也已經成為豐盛。他是所有執政掌權者的元首。 \end{tabularx} \\ \\ \relax
2:11 & \begin{tabularx}{0.7\textwidth}{X} 你們也在他裡面受了不是人手所行的割禮，而是使你們脫去肉體情慾的基督的割禮。 \end{tabularx} \\ \\ \relax
2:12 & \begin{tabularx}{0.7\textwidth}{X} 你們既受洗與他一同埋葬，也就在此禮上，因信那使他從死人中復活的神的作為跟他一同復活。 \end{tabularx} \\ \\ \relax
2:13 & \begin{tabularx}{0.7\textwidth}{X} 你們從前在過犯和未受割禮的肉體中死了，神卻赦免了你們一切的過犯，使你們與基督一同活過來， \end{tabularx} \\ \\ \relax
2:14 & \begin{tabularx}{0.7\textwidth}{X} 塗去了在律例上所寫、敵對我們、束縛我們的字據，把它撤去，釘在十字架上。 \end{tabularx} \\ \\ \relax
2:15 & \begin{tabularx}{0.7\textwidth}{X} 基督既將一切執政者、掌權者的權勢解除了，就在凱旋的行列中，將他們公開示眾，仗著十字架誇勝。 \end{tabularx} \\ \\ \relax
2:16 & \begin{tabularx}{0.7\textwidth}{X} 所以，不要讓任何人在飲食上，或節期、初一、安息日等事上評斷你們。 \end{tabularx} \\ \\ \relax
2:17 & \begin{tabularx}{0.7\textwidth}{X} 這些原是未來的事的影子，真體卻是屬基督的。 \end{tabularx} \\ \\ \relax
2:18 & \begin{tabularx}{0.7\textwidth}{X} 不要讓人藉著故作謙虛和敬拜天使奪去你們的獎賞。這等人拘泥在所見過的幻象，隨著自己的慾望無故地自高自大， \end{tabularx} \\ \\ \relax
2:19 & \begin{tabularx}{0.7\textwidth}{X} 不緊隨元首；其實，由於他全身藉著關節筋絡才得到滋養，互相聯絡，靠神所賜的成長而成長。 \end{tabularx} \\ \\ \relax
2:20 & \begin{tabularx}{0.7\textwidth}{X} 既然你們與基督同死而脫離了世上粗淺的學說，為甚麼仍像生活在世俗中一樣，去服從那「不可拿、不可嘗、不可摸」等類的規條呢？ \end{tabularx} \\ \\ \relax
2:21 & \begin{tabularx}{0.7\textwidth}{X} 見上節 \end{tabularx} \\ \\ \relax
2:22 & \begin{tabularx}{0.7\textwidth}{X} 這些都是根據人的命令和教導，論到這一切都是一經使用就都敗壞了。 \end{tabularx} \\ \\ \relax
2:23 & \begin{tabularx}{0.7\textwidth}{X} 這些規條使人徒有智慧之名，用私意崇拜，自表謙卑，苦待己身，其實在克制肉體的情慾上毫無功效。 \end{tabularx} \\ \\
[1ex]
\hline
\hline
\end{longtable}
$^{1}$各位網友平安.
雖然農曆新年過了幾天.
但今日仍是年初六.
所以仍在新年假期.
牧師在此為大家拜個年.
祝大家今年靈命的日長.
愛主的日心.
在新的一年.
滿有主的恩典.
過一個土神喜悅的人生.
上個月的講道我跟大家講過.
今年第一季我會用經卷形式去講道.
以三次的講道內容.
跟大家對歌羅西書有更深入的認識.
在上一次的講道裡面.
我提過哥羅西教會所面對的問題.
就是信徒受到異端的引誘.
讓他們偏離了基督的信仰.
尤其是耶穌基督的身份.
有些異端完全否定了耶穌獨特的神聖.
耶穌只不過是眾多神明之中的其中之一.
有些否定了耶穌的神聖.
認為耶穌只是一個人形的軀殼.
神只不過是依附在耶穌的肉體裡面.
異端對耶穌基督的神人異性.
產生了不同程度的質疑.
令到有些信徒最終離開了基督的信仰.
因此保羅在歌羅西書第一章.
講出耶穌基督的超越性.
祂是宇宙的開始.
亦是世界的終結.
世界被創造之先.
耶穌基督已經存在.
世界是被基督所創造.
一切都是為了基督而被創造.
因此神的本質尊貴無比.
而神的心意就是要我們活在祂的風城裡面.
來到第二章.
保羅對異端所提出的不同主題.
逐一地辯解.

$^{41}$讓信徒不再陷入被誘惑的裡面.
重新回到基督的信仰裡面.
我們進入經文之前一起祈禱.
你賜下聖經.
讓我們可以在裡面明白真道.
認識你更加多.
甚願我們每一個跟隨你的.
都在你的話語上高扎根.
成為我們信仰上美好的根基.
我們祈禱奉主耶穌基督.
得勝名字祈求.
阿們.
我們開始進入經文.
異端當時有四種.
與基督信仰背道而馳的理論.
是哪四種呢?.
就是一般的世俗學說.
律法主義.
敬拜天使.
以及撫修主義.
我們一個一個探討.
他們對基督信仰的影響.
在古代的社會裡面.
哲學是一種很廣泛的學說.
不單單作為一種探討.
與神有關的知識性交流.
甚至一切與世界.
和生命意義的內容.
都包括在哲學的理論層面裡.
因此哲學是當時.
知識分子溝通的一種媒介.
保羅不是反對所有的哲學.
而是反對當時.
異教所提出與神有關的理論.
第八節是這樣說.
他說差點哲學和虛空的廢話.
不照著基督.
而是照人間的傳統.
和世上粗淺的學說.
虛空的廢話.

$^{81}$就是指到自然界的靈體.
而人間的傳統.
和世上粗淺的學說.
就是指逐世的價值觀.
自然界的靈體.
是泛神論的其中一種.
認為神只不過是被萬物所充滿.
世界上有不同的靈體.
充滿在地上每一個角落.
負責著不同的工作.
此外人間的傳統和粗淺學說.
就是指當時世俗的價值觀.
保羅認為這兩種學說.
都不是跟著基督的.
就是說兩種觀念背後的核心價值.
是違背了基督信仰的精神.
基督信仰所說的是一神論.
神是獨一的.
是宇宙和世界的創造者.
世界是被神所創造.
而不是神是世界的一部分.
正正是神是一切的創造者.
因此神是尊貴.
並不是滿天神佛.
也不是與其他靈體.
享有同一個地位.
至於逐世的價值觀.
也只是一種按當時的處境.
而產生的處世態度.
是屬世的.
不能夠將人帶入永恆.
甚至是粗淺的.
例如教人如何發達致富.
這種就是所說的粗淺.
各位姐妹.
你覺不覺得今天.
有沒有一種異端邪說在我們身邊呢?.
當然有.
「範神論」是現今最普遍的流行學說.
尤其是我們很熟悉的中國民間宗教.

$^{121}$是「範神論」的表表者.
天堂有玉皇大帝.
掌管著其他者的神.
而其他神也各有功能.
甚至無有個別的神.
是會受歡迎的程度.
比玉皇大帝更高.
有沒有聽過?.
例如觀音.
例如關公.
但對入屋的神.
莫過於每個家門口的土地.
祂是家庭的神.
只要你敬拜祂.
你供奉祂.
祂就會保護著你一家平安.
這種「範神論」.
就是讓人可以各有所需.
你喜歡哪一個.
你就供奉哪一個.
哪一個可以帶來心理的利益.
那就敬拜哪一個.
但我們從來都不會反思一下.
他們所敬拜的神.
究竟有沒有違格呢?.
如果天上的神.
完全沒有主權.
只不過是人的一種滿足工具.
這種神究竟有沒有神性呢?.
或者用另一個說法.
他們有沒有一種超越性呢?.
能夠帶領敬拜祂.
或供奉祂的人.
可以超越祂本身的限制和綑綁的呢?.
所以這種信仰.
除了是一種利益的交換之外.
也只是一種心靈的慰藉和寄託.
完全不能帶給我們一個永恆的超越性.
對於俗世的價值觀.
西方世界現在最流行吸引的.

$^{161}$其實就是叫做新紀元運動.
好像他們說.
萬物歸一.
一切宗教最終都會歸於一統.
一切都有神性.
人性即是神.
相信宇宙進化.
包括人的意識.
以及人性的進化.
而這種進化使世界最終都是樂觀的.
信念創造現實等等.
新紀元是甚麼理論呢?.
其實是吸收了東方和西方的古老宗教傳統.
而且將很多觀念和現代科學的觀念融合在一起.
特別是心理學和生態學.
新紀元思想也涉及世界各大宗教的靈感.
最普遍的就是猶太教.
基督教.
伊斯蘭教.
印度教.
中國民間傳統.
例如儒釋道.
特別是西洋的神秘學.
諾斯底主義.
由於他們將所有的理論滲入在不同宗教的思想裡面.
所以一般信徒是不容易辨識的.
我舉個例子.
很多信徒覺得參加儒家.
玩一下儒家班.
做拉筋.
有甚麼所謂呢?.
這種包裝成養生或心靈平靜的活動.
是會強調它的功效.
但是儒家或超覺靜坐背後的印度教思想.
它就不會再跟你說了.
再舉另一個例子.
信念創造現實.
Life Dynamic的課程現在在外面很流行.
它們目前最流行的一種理論.
發掘你的潛能.

$^{201}$發掘你的小宇宙.
創造你的未來.
人是自己生命的主宰.
這種以人為本的理論很吸引.
就算是信徒也不容易察覺有甚麼不妥.
基督教的信仰是以神為本.
神才是我們一切的倚靠.
神才是我們一切得力的來源.
而不是我們自己.
弟兄姊妹.
今天世界有很多觀念.
很多理論.
我們如何去界定這些理論和學說.
是有違信仰.
我們不能夠跟進呢.
保羅在第七節提醒我們.
在他裡面生根建造.
我們的信仰生命是需要有根的.
而根源是需要來自甚麼.
神的話語.
最後才能夠在這個基礎上去建造.
最後才能夠以神的眼光去看這個世界和人生.
因此信徒不能夠只吃靈奶.
不去.
經常被人餵.
我們要去吃粥.
要吃飯.
讓我們整個生命都能夠成熟.
將生命建造在神的話語裡面.
才能夠生根建造.
否則我們只不過是一個在信仰上吃奶的嬰兒.
第二方面值得我們去批判的就是律法主義.
在摩西律法裡面對於潔淨和不潔淨的食物有嚴格的規定.
因此規範了信徒可以吃甚麼.
有甚麼是不能吃的.
而且猶太教亦有很多節期和禮節.
例如鑄鵬節.
逾越節.
五旬節.
而安息日更加是他們必定要守的一個條例.

$^{241}$對於當時的愛邦信徒來說.
他們面對一班非常頑固的保守的猶太信徒.
指責這批愛邦信徒.
說他們不遵守摩西所定的節期安息日等等的條例.
其實我們知道在《賜堂行傳》裡面記載過.
保羅曾經為了愛邦人是否要行國禮.
而上去耶路撒冷請示過彼得.
雅各和約翰等等的教會領袖.
最後得到教會領袖的認可和同意.
愛邦人信主之後是不需要行國禮的.
其他猶太教的節期和禮節.
愛邦信徒是否同樣不需要遵守呢.
因此愛邦信徒就受到一些頑固的猶太信徒攻擊.
認為他們不守摩西的律法.
是對摩西的一種藐視.
保羅在第十一至十三節十四至十五節提出.
既然你們都在基督裡面受了不是人手所行的國禮.
甚麼叫做不是人手所行的國禮呢?.
就是我們的受浸.
神就赦免了你們一切的過犯.
使你們與基督一同活過來.
我們是一個新造的人.
抹去了在律法上所寫的敵對我們.
束縛我們的自居.
拍打切去釘在十字架上.
一切的律法在十字架上都已經成就了.
基督既將一切執政的掌權解除了.
就在凱旋的行列中將他們公開示眾.
仗著十字架跨城.
今天我們每一個耶穌基督的跟隨者.
我們倚靠耶穌的.
都是倚靠著這個十字架而得勝.
保羅的意思就是說.
主耶穌已經是一個得勝的主了.
既然可以將執政者的權勢解除.
自然可以解除律法上對信徒的一切捆綁.
所以十六節他就說.
不要讓任何人在飲食上或節期.
初一,安息日等事上評論你們.
因為這些律法和條例.

$^{281}$都是未來的預表.
但是現在這些預表的東西都已經來到了.
還要守這些預表的法律嗎?.
舊約的禮儀,律法和條例.
都是向人描述並且指明基督的預表和象徵.
信徒要留意的重點應該集中在.
哪一個才是真正要來的基督.
現在基督來了.
信徒若再堅持遵守這些禮儀,律法和條例.
就是寫本逐末了.
沒有認清這個禮儀和條例的功用.
今天我們也是.
我們不要被那些律法所規綁.
成為我們律法的奴隸.
我們已經被主耶穌的救贖所釋放.
不過我們仍然都需要注意.
我們雖然是自由.
但我們卻不能夠放縱和濫用主賜給我們的自由.
當我們一旦濫用主賜給我們的自由的時候.
我們就走向另一個極端.
今天仍然有信徒問我.
拜過神的雞,能吃嗎?.
我們能不能進廟參拜呢?.
問這些問題的信徒.
表示他們仍然在吃奶.
問這個問題本身一點問題都沒有.
有問題就要問.
不過如果他們的生命不再在這裡生根建造的話.
再過十年.
他們仍然會為這些問題而煩惱.
保羅在哥林多前書說過.
凡事都可行,但不都有益處.
凡事都可行,但不都造就人.
在主裡面我們什麼都可以做.
不過如果對人對己都沒有益處.
又不能夠造就別人的話.
縱然可以做.
但我都可以選擇不去做.
或者我舉多一個例子.
去說明什麼叫真正的自由.

$^{321}$能夠有真正的選擇.
才是真自由.
沒有真正的選擇的選擇.
其實是捆綁.
我用打機為例.
你可以選擇什麼時候打.
打多久.
什麼時候收手.
最後如果完全是按照你決定的範圍進行的話.
這是你的自由.
但你不想打機.
一旦接觸到這部遊戲機.
你就控制不了.
由什麼時候打.
打多久.
什麼時候停.
你都完全不受控制.
你無論作出任何的選擇.
其實都只是一種結果.
就是繼續打機.
而且不能夠控制.
什麼時候終止打機的時間.
這種選擇不是自由.
是捆綁.
類似的例子就是.
吸煙,上網,看電視.
你以為自己可以有選擇.
但其實是被他們捆綁.
被他們克制.
弟兄姊妹.
當我們面對罪的誘惑.
我們能不能戰勝他們呢?.
基督已經賦予了每一個跟隨他的人.
可以選擇向罪說.
不.
我們是否有真正體驗過.
這種真正的罪惡.
有沒有真正向罪惡說.
不.
保羅對異端的第三種批判.

$^{361}$就是對敬拜天使的批判.
十八節.
不要讓人藉著故作謙虛和敬拜天使.
奪去你們的賞賜.
這等人的驅離在所見過的幻象.
隨著自己的慾望無故地自高自大.
當時有一種異端認為.
神高高在上.
人不能直接敬拜他.
必須以一種謙卑的態度.
經過敬拜天使來討好神的喜悅.
另外他們也認為自己可以修煉到一種高超的境界.
能夠看到一般世人看不到的異象.
因此他們自覺高人一等.
主耶穌甘心情願道成肉身.
走上十字架.
成為人神之間的忠補.
是出於對信徒的愛.
而不是想高高在上.
受人的景仰.
如果天使能夠成為人神之間的忠補.
不單單貶低了主耶穌.
和天父的救贖心意.
也抬高了天使的身份.
因為天使只是一個服侍的靈.
並沒有這麼重要的責任和角色.
如果主耶穌不能直接接受人的敬拜和親近.
那將主耶穌視為一個高高在上.
不惜人間煙火的神.
對人毫不親近.
毫無憐憫.
對主耶穌的屬性就有嚴重的偏差.
其實初期天主教的信徒敬拜瑪利亞.
也有一種視瑪利亞為忠補的心理作用.
當時的信徒認為透過耶穌的母親去請求主.
主就會因為母親的緣故而不會拒絕.
這種思想就將主耶穌的忠補身份奪去了.
後來教廷在梵蒂岡第二次會議修正了瑪利亞的身份.
指出瑪利亞並沒有任何神聖.
最多只是一個聖人.

$^{401}$她的行為值得信徒仿效.
但卻不具有神與人之間忠補的重要角色.
因此由於當時的異端過份強調他們的神秘經驗.
讓他們自覺高人一等.
我們曾經聽說.
保羅在哥林多的前書中提及.
他自己也有三層天的經歷.
或許是這樣導致其他信徒.
都誇張地表達個人的神秘經歷.
認為是一種與別不同的優越感.
讓他們產生了一種驕傲.
今日個別的靈恩教會都強調方言.
很強調先知預言的恩賜.
認為有方言才顯明與主耶穌的親密關係.
有預言或圖像才顯明他們能夠從神領受特別的啟示.
當然我們很相信.
方言和預言都是神的恩賜.
但當信徒過份強烈一種個人與主的神秘經驗.
就會令其他信徒覺得他們與別不同.
這也是一種信仰的偏差.
弟兄姊妹.
神學家加爾文提出過.
信仰裡有四個唯獨.
哪四個呢?.
唯獨是恩典.
我們得救是一個恩典.
唯獨是聖經.
唯有聖經是神的話語.
唯獨是基督.
唯獨是基督拯救我們.
以及我們唯獨榮耀神.
基督作為神人之間的忠報是不能代替的.
我們也唯有透過主耶穌才能進入婦那裡.
我們也唯有透過敬拜主去表達我們對神的讚美.
保羅對二端的第四個批判.
就是對他們的苦修主義.
二十三節.
這些規條使人們徒有智慧之名.
用私義崇拜.
自表謙卑.

$^{441}$苦待己身.
其實在克制肉體的程度上毫無功效.
由於苦修者認為.
他受到二端所影響.
他視肉體是邪惡的.
故此認為誠信之道.
是禁制肉體一切的意欲喜好.
並且將他的需要減到最低.
對苦修者來說.
肉體是應該被懲罰和打擊的.
應該將它當作敵人般看待.
在主教第三世紀和第四世紀.
苦修主義極為普遍.
有些人甚至將自己的身體禍害.
搞得很受削.
結果忽視了身體的護理.
保羅對苦修主義作出一個回應.
在二十至二十一節說.
既然你們與基督同死.
而脫離了世上粗淺的學說.
為何仍然將生活在世俗中一樣呢?.
為何還要服從那些.
不可以拿,不可以嘗,不可以摸的規矩呢?.
基督徒既然已經與耶穌同死.
超越了地上一切的限制.
為何仍然被這些俗世的規條所影響呢?.
苦修制肉的條例.
不單是人為的.
而且對控制肉體本性的放縱.
是毫無功效的.
不可以拿,不可以嘗,不可以摸.
是異教徒的飲食戒條.
得到救恩的門路.
一部分的規條可能是源於.
改頭換面的摩西律法.
而其他的必然是源自異教的節育主義.
我們看動詞的表達裡.
便知道它逐少地減少範圍.
不可以拿,不可以嘗,不可以摸.
其實是一種自我設限的思維.

$^{481}$弟兄姊妹.
今天我們未必會被苦修主義所影響.
去苦待我們的身體.
但我們對身體的觀念.
仍然深深被扭曲的大眾文化所影響.
香港自從有選美活動的開始.
便已經將女性身體的價值觀扭曲了.
入圍的參加者.
必定要符合某種對身體的要求.
才不會被淘汰.
參選者穿著泳衣去展現不同的體態.
成為一種大眾文化對美麗身體的標準.
電視台將女性美好身體的意識形態.
透過大型電波.
我們今天所說的豬扒.
便成為一個醜女的標籤.
因此也造就了日後的千體文化.
無論是否參加選美.
女士們都將身體塑造成某一種標準.
改善我們的麒麟臂,大象腿.
成為千體廣告的字眼.
以符合大眾文化覺得這是美麗的要求.
因此無論是增肥或減肥.
甚至是消脂整容.
都不是因為身體健康的緣故.
而是為了滿足社會文化對美和醜的界定.
對於今天的女性來說.
這是否一種捆綁?.
弟兄姊妹今天說的講題是.
尊貴無比的基督信仰.
第一章的教禱告訴我們.
我們所相信的基督是尊貴無比的.
第二章的教禱告訴我們.
我們接受了基督.
祂帶給我們的信仰.
都是尊貴無比的.
我們的信仰從對生命感到絕望.
最後可以達到豐盛,滿有盼望.
我們的信仰超越捆綁.
讓每一個屬神的都能得著自由.

$^{521}$脫離捆綁的人生.
我們的信仰提升生命質素.
讓我們的老我成為願意謙卑服侍的人生.
我們的信仰毫不孤單.
因為有主同行同在.
我們的信仰因著主的應許.
我們能夠超越今生.
直到永恆.
既然我們所相信的信仰是如此尊貴無比.
我們可否像第七節所說的.
心裡存著信心堅固.
信心的堅固不在於我們的行為是多麼渺小和軟弱.
而在於我們所相信的主是多麼愛我們.
不在於我們回應的行動是如何.
而在於我們的愛.
而在於我們愛主的心是如何.
我們愛主的心越大.
我們就越願意回應主對我們的呼喚.
今年教會的年提.
大家都很清楚.
應該是「毀身事奉」.
神對我們洪恩家已經發出了一個呼召.
弟兄姊妹我們如何去回應呢?.
甚至我們都能夠以我們的經文.
生根建造.
信心堅固.
我們一起獻上祈禱.
因著你的寶貴.
也讓我們每一個跟從你的都顯得尊貴無比.
我們原本不配的人生.
因著認識你而改變.
我們的身份改變.
我們的宗教也改變.
我們的生命因此而改變.
一切都是因著神賜給我們的恩典.
求主讓我們每一個願意跟隨你的.
都心存感恩的心去過我們每一天.
讓我們生命充滿著你的恩典.
好叫我們所遇到的每一個人.
都因著我們所實踐的信仰而得福.

$^{561}$讓更多人能夠歸入你的名下.
也讓主的名在洪水橋這個地方能夠得著榮耀.
我們祈禱.
是奉教主耶穌基督得勝名字祈求.
阿們.
願主賜福給大家.
\newpage



\section{出埃及記 17:8-16}
\label{sec:_xVz_wL7W7s}
\textbf{2022.02.13 《耶和華尼西》 出埃及記 17\_8-16 (和修版) 曾敏傳道}
\newline
\newline
連結: \href{https://youtube.com/watch?v=-xVz-wL7W7s}{\texttt{https://youtube.com/watch?v=-xVz-wL7W7s}} ~~~~ 語音日期: 2022-02-13
\newline
\newline
\hyperref[sec:aSFgvj6BhUY]{\small{< < < PREV SERMON < < <}}
~
\hyperref[sec:index_chronic]{\small{[返順時目]}}
~
\hyperref[sec:index_scriptual]{\small{[返順卷目]}}
~
\hyperref[sec:jmOEpWc4LtY]{\small{> > > NEXT SERMON > > >}}
\newline
\newline
出埃及記 17:8-16
\newline
\begin{longtable}{cl}
\hline
\hline
章節 & 經文 (和合本修訂版)\\
\hline
17:8 & \begin{tabularx}{0.7\textwidth}{X} 那時，亞瑪力來到利非訂，和以色列爭戰。 \end{tabularx} \\ \\ \relax
17:9 & \begin{tabularx}{0.7\textwidth}{X} 摩西對約書亞說：「你為我們選出人來，出去和亞瑪力爭戰。明天我要站在山頂上，手裡拿著神的杖。」 \end{tabularx} \\ \\ \relax
17:10 & \begin{tabularx}{0.7\textwidth}{X} 於是，約書亞照著摩西對他所說的話去做，和亞瑪力爭戰。摩西、亞倫和戶珥都上了山頂。 \end{tabularx} \\ \\ \relax
17:11 & \begin{tabularx}{0.7\textwidth}{X} 摩西何時舉手，以色列就得勝；何時垂手，亞瑪力就得勝。 \end{tabularx} \\ \\ \relax
17:12 & \begin{tabularx}{0.7\textwidth}{X} 但摩西的雙手沉重，他們就搬一塊石頭來放在他下面，他就坐在上面。亞倫與戶珥扶著他的手，一個在這邊，一個在那邊，他的手就穩住，直到日落。 \end{tabularx} \\ \\ \relax
17:13 & \begin{tabularx}{0.7\textwidth}{X} 約書亞用刀打敗了亞瑪力和他的百姓。 \end{tabularx} \\ \\ \relax
17:14 & \begin{tabularx}{0.7\textwidth}{X} 耶和華對摩西說：「你要把這事記錄在書上作紀念，又念給約書亞聽：我要把亞瑪力的名字從天下全然塗去。」 \end{tabularx} \\ \\ \relax
17:15 & \begin{tabularx}{0.7\textwidth}{X} 摩西築了一座壇，起名叫「耶和華尼西」。 \end{tabularx} \\ \\ \relax
17:16 & \begin{tabularx}{0.7\textwidth}{X} 他說：「我指著耶和華的寶座發誓，耶和華必世世代代和亞瑪力爭戰。」 \end{tabularx} \\ \\
[1ex]
\hline
\hline
\end{longtable}
$^{1}$天父,今天求你開我們的眼睛.
讓我們認識你更加深.
也讓我們聽到你的聲音.
要明白上帝的心意如何.
就讓我們有力量.
面對現在的處境.
面對我們的生命.
面對不同的問題.
天父,求你賜福給我們.
奉耶穌基督的名字祈禱.
阿們.
這個講題我們應該比較少講.
耶和華尼西這個題目.
但卻是非常非常寶貴.
今天我們要看一場戰爭.
這場戰爭是很特別的.
不是單純的一個層面.
我們先來看一下.
耶和華尼西是耶和華的「禎奇」.
究竟是什麼意思呢?.
希望今天在這堂道裡.
讓大家深深地明白.
也知道耶和華尼西這個名字.
對於我們的意義在哪裡.
今天我們看的那段經文.
是《春海及記》十七章八節打後的一段經文.
當時阿瑪力來到利菲定這個地方.
與以色列爭戰.
這個圖大家看.
或許不是十分清晰.
利菲定在圖上.
現在你看到的畫面.
右手邊比較下一點的位置.
你看到有四個字叫西乃曠野.
上一點的位置就是利菲定.
在那個端端那裡.
就在那個地方.
有一場戰爭.
這個是以色列民出埃及後.
在曠野的第一場戰爭.

$^{41}$阿瑪力究竟是什麼人呢?.
值得我們認識一下.
在《創世記》三十六章十到十二節那裡這樣說.
「以素子孫的名字如下.
以素的妻子雅大生以利法.
亭納是以素兒子以利法的妾.
她為以利法生了阿瑪力」.
所以我們知道阿瑪力是以素的後人.
阿瑪力人與以色列人是親屬.
他們有關係的.
因為以素與雅各是兩兄弟.
所以他們的後人也是親屬的關係.
但是我們很肯定地說.
他們的關係並不友好.
甚至可以這樣說是宿敵.
這個值得我們去思考.
為什麼親屬到後來又發展成為敵人.
並不友好呢?.
我們很肯定一件事就是相當不起於阿瑪力人.
因為他們的手段非常卑劣.
可以這樣說他們欺負以色列民.
特別記載這一段.
以色列民出埃及的路上.
他們遇到阿瑪力的這一段經歷.
聖經在其他地方有記載.
神對這一群人的評價是什麼.
申命記的二十五章十七到十九節這樣說.
你要記得你們出埃及的時候.
阿瑪力在路上怎樣對待你.
在路上迎擊你.
趁你疲乏困原的時候.
擊殺所有在你後面軟弱的人.
並不敬畏神.
到後面上帝說.
你要把阿瑪力的命從天下逃去.
其實在這個時候.
以色列民並沒有主動挑動戰爭.
沒有.
但是阿瑪力卻要欺負他們.
而且卑劣的地方就是.

$^{81}$趁他們累了就要取他們的命.
他們非常累.
很疲乏很軟弱的時候.
他就要攻擊他們很軟弱的人.
真的想要取他們的命.
這個如果大家記得.
出埃及的場面是浩浩蕩蕩.
不只是有男人在裡面都有婦孺.
也肯定有長者.
在這個時候在走難當中.
這樣被人攻擊.
是很不容易去面對的.
今天值得我們問.
誰是我們的阿瑪力人.
阿瑪力人是以色列的仇敵.
那誰是我們的阿瑪力人呢?.
或許我們和弟兄姊妹會有衝突.
我們不喜歡某個人.
或者某個人不喜歡我們.
但是我可以很肯定地說.
他們絕對不是我們的敵人.
他們不是我們的敵人.
所以我們的矛頭.
沒有指向我們的弟兄姊妹.
是因為我不喜歡他.
我覺得他這樣.
他覺得我這樣.
所以我們彼此為敵.
不是的.
他們不是我們的敵人.
真正的敵人只有一個.
那惡者撒旦.
彼得前書五章八節這樣說.
務要謹慎 要警醒.
因為你們的仇敵魔鬼.
如同咆哮的獅子.
走來走去.
尋找可吞吃的人.
務要謹慎 要警醒.
你們的仇敵魔鬼.

$^{121}$這段話不是對著.
未信主的人說的.
是對著信耶穌的人說的.
上帝想我們要去小心.
是這個仇敵.
不是我們身邊的人.
我們甚少去說.
究竟這個仇敵.
會怎樣對付我們的生命呢?.
原來有至少幾個方式.
第一樣很自然的.
是引誘我們犯罪.
當我們陷在罪惡當中的時候.
我們自然的心就會遠離神.
我們跟神的關係很疏遠.
其實惡者無非想做一件事.
最好就是遠離神.
最好就是沒有親近神.
跟神有這麼遠 得這麼遠就最好了.
我破壞了這個關係就成功了.
所以所有事情都向著這個向度.
第二樣就是當我們遇到.
人生的風浪和巨大的困難的時候.
那個惡者都會趁機影響我們的信心.
讓我們對神失去信心.
我們開始懷疑神.
一定是上帝作弄我.
一定是上帝對我不好.
為什麼要這樣呢?.
在疫情期間一定會有很多人問.
上帝你知道疫情很厲害.
你見不見到.
每天排隊去強檢.
你見不見到我們每天一千多宗.
有機會上到一萬宗.
上帝為什麼不出手?.
上帝是不是你要整頓我們?.
我們開始懷疑神.
我們埋怨神.
一定是祂對我們不好.

$^{161}$我們甚至遠離神.
這個也是惡者的攻擊.
一些詭計.
令我們覺得對,懷疑.
一定是祂不好.
祂對你好,怎麼會有疫情呢?.
這些時候要小心.
第三件事.
當我們彼此.
弟兄姊妹當中.
我們肢體當中發生衝突的時候.
那個惡者也在當中搞破壞.
加深大家不愉快的感覺.
造成教會分裂.
分黨分派.
自然大家就沒辦法只為主作見證.
大家都很憤怒.
我很想對付那個對我不好的人.
每天用嘴巴說是非.
就要說對方怎樣不好.
那惡者就贏了.
在不知不覺之間.
他已經勝利.
因為我們已經不同心.
我們也沒有這樣的力量.
令人歸主.
第四件事也是很卑劣的.
就是他可以直接攻擊我們的身心靈.
令我們經歷很多患難和痛苦.
我想說一個事實.
就是如果我們經常聽見見證.
宣教士和傳福音的勇士.
他們經歷的屬靈攻擊是很多的.
我很肯定這一點.
因為他們要搶救靈魂.
所以一定很直接地.
會跟那個惡者碰到正.
還有凡納至在基督耶穌裡.
過敬虔生活的.
也要受到逼迫.

$^{201}$要面對很多屬靈的攻擊.
這個一定是真的.
除非我今天.
我不打算敬虔道人.
否則你要知道.
那惡者是一定不會放過你的.
我最近去預備這篇講章的時候.
去思考的時候.
我發覺不是.
不只是這些人面對很多屬靈的攻擊.
不只是他們面對那惡者的攻擊.
其實其他人也會.
那個惡者就是這樣.
透過不同的位置.
一直就是這個好像咆哮的獅子.
走來走去.
找不同的空間去攻擊不同的人.
我們平時很少思考這些問題.
我們是否在自己的罪惡裡.
是否在很多關係的破裂裡.
根本沒有這個空間去思考.
但今天藉著這場仗要想.
我們的亞馬力人是誰.
弟兄姊妹想清楚.
不是你的弟兄.
你的姊妹.
而是那惡者.
究竟我們要怎樣去打這場仗呢.
一會兒我會繼續說.
其實我們生命裡有什麼仗.
我們要去打呢.
要去說.
我們說回.
又看回那場仗.
很明顯亞馬力人是有備而來的.
他們很善於打仗.
以色列人就不是了.
說得再下一點.
以色列人是很集排軍.
集排軍沒有訓練的一批人.

$^{241}$他們不善於打仗.
但在這個位置.
以色列人沒有選擇.
他們並不好戰.
他們不是想去打別人.
別人打到過來.
你不去迎戰沒有辦法.
所以被迫一定要去打.
我們可以從三個層面去看這場仗.
這是非常精彩.
其實很少在聖經裡.
讓我們看這層面是這麼清晰.
這三場仗.
但我想說.
事實上我們在生命裡.
在生活裡.
在事奉裡打仗.
都會經歷這三個層次.
哪三個層次呢.
第一個地面上的爭戰.
肉搏戰一定會.
第九節摩西對約斯亞說.
你為我們選出人來.
出去和亞馬力爭戰.
十節.
於是約斯亞照著摩西對他所說的話去做.
和亞馬力爭戰.
你一定要有前方的戰士出去打.
但有一個特別的位置值得我們注意.
就是你要選出人來.
原來不是每個人都打.
要選人.
但是.
聖經沒有說得很清楚.
你的準則是什麼呢.
因為這不是這段經文的重點.
但我知道要選人.
我不知道準則是什麼.
但知道一定是選對的人.
上帝覺得對的人.

$^{281}$他才會選去打仗.
這個值得我們思考.
上帝覺得對.
不是我覺得自己對.
不是別人覺得我對.
上帝覺得我對.
是元帥約斯亞覺得那個人對.
適合這個軍隊.
我就選他來打.
弟兄姊妹.
我和你的生命對上帝嗎?.
我和你的生命.
是否一個被揀選的生命.
我們是否預備好裝備好.
上帝會選我們去打仗呢.
合不合用.
是值得問這個問題.
這個地面上的仗.
什麼情況.
是否有多激烈.
過程是怎樣.
約斯亞拿著什麼武器衝過去.
遇到什麼敵人.
對方的反應是怎樣.
其實經文是沒有說的.
因為這不是重點.
但知道約斯亞要去打.
選人去打.
他一路打是很艱鉅.
打了很久.
不是十幾二十分鐘.
由早上打打打.
打到下午.
一直打到很長時間.
只知道他打了很久.
在打的過程中.
原來有人會很影響戰局.
第二個層面就是山上的爭戰.
第九節這樣說.
摩西對約斯亞說.

$^{321}$明天我要站在山頂上.
手裡拿著神的杖.
神的杖是代表神的權能.
他拿著神的杖去宣告.
上帝的同在.
上帝的權能在這裡.
在這場仗中.
上帝會彰顯他的權能.
上帝和我們一起.
摩西要做這件事.
摩西不是一次拿著神的杖.
如果大家記得.
神的權能是怎樣.
我們看看.
在《春拉及記》十四章十五節.
跟摩西說.
你吩咐以色列人往前走.
你舉手向海伸杖.
將水分開.
原來叫紅海分開那一刻.
是摩西拿著那支杖.
宣告上帝的權能.
上帝和我們一起.
當他伸杖.
紅海就分開.
上帝的權能是這麼大的.
有上帝同在是這麼奇妙的.
我深信過紅海.
這個特別的經歷.
是以色列人當日才會經歷到.
是這麼奇妙.
今天摩西就在山上.
第二個層面就是.
他在山上去舉手.
去舉手宣告上帝的權能.
上帝和以色列同在.
在那裡撐著那場戰爭.
很多時候我們都會想.
他舉著那支杖.
就是這樣舉著那支杖.

$^{361}$很多時候我們會聯想到.
他會禱告.
他就是一路舉手.
一路在那裡山上禱告.
托著整個以色列.
托著約書亞整個軍隊.
這是這段經文很重要的位置.
那裡這樣說.
《猜疑及記》17章到11節說.
摩西,亞倫和,胡爾都上了山頂.
摩西何時舉手.
以色列就得勝.
何時垂手.
阿瑪力就得勝.
原來不可以只是這樣.
約書亞在下面打.
約書亞應該很能打的.
他年輕,他力壯.
但是我想和你說.
經文並沒有說.
哪個約書亞真厲害.
他一帶著軍隊打仗就贏了.
戰無不勝.
哪天從早到晚.
約書亞就是哪個得勝的將軍.
沒有這樣說.
反而說摩西何時舉手.
他就贏.
何時垂手.
他就輸了.
原來摩西是會影響他的.
他的禱告.
他宣告上帝的權能這件事.
是很重要的.
沒有摩西這個禱告.
沒有摩西一直全程托著他.
根本就不行.
他托多久.
是不是托一兩分鐘呢.
不是托一兩分鐘.

$^{401}$不是托幾個字.
是他打多久,他就站多久.
今天我們很多時候.
我們都說.
弟兄姊妹你有需要.
我們教會有需要.
我們開始前站一會就好.
一會就好.
真的一會.
很厲害的.
我沒見過.
去到一個小時的祈禱都沒有.
幾個字就已經很厲害了.
我為弟兄姊妹祈禱.
就幾個字.
你們知不知道這場仗.
是他打多久就站多久.
幾時垂手他就輸了.
你的禱告是要一同配合的.
一同是一切的.
是要這樣的.
沒有了這個.
他在山下多勇猛都不行.
其實在過程某個程度.
你都看到一件事.
說真的.
以色列文約書亞這邊的軍隊.
有多厲害呢.
沒有了禱告.
更加重要.
一會兒第三個層面.
沒有了那樣.
其實以色列是會輸的.
說真的.
在能力上不是真的那麼大.
沒有禱告怎麼行呢.
當然這不是純粹.
肉身上的那場仗.
也是一個屬靈的爭戰.
在裡面不是禱告.

$^{441}$一定不能夠贏.
我們看到這個很震撼的場面.
摩西要去舉仗.
去舉仗才能贏.
宣告上帝的權能.
痛哉.
有人在他身邊的.
經文讓我們很清楚的看到.
摩西也不是一個人.
在爭戰的道路上.
我們需要有同行者.
因為會累的.
一個人都撐不了那麼久.
你說摩西祈禱很厲害.
就他一個人.
他從頭到尾一個人.
一直撐到黃昏日落.
都不行的.
他累了.
所以需要亞倫和會義.
去扶他的手.
這裡看到經文說.
摩西的雙手沉重.
你舉手舉了多久.
上帝要他這樣舉著.
宣告權能.
上帝痛哉.
我們曾經玩過遊戲.
就是舉手的遊戲.
原來不是舉了很久.
很快手就酸軟了.
就算沒有拿著東西.
手也不能舉很久.
酸軟之後就想放下.
還要像摩西一樣.
舉著像這樣.
從早到晚.
不能撐很久.
所以他也累了.
雙手沉重.

$^{481}$於是亞倫和會義就搬了一塊石頭.
放在他下面讓他坐下.
祈禱整天這樣.
受不了.
還要扶著他的手.
一個在這邊.
一個在那邊.
他的手就穩住.
直到日落.
手累不累.
撐足一天.
累了.
沒有人扶著你.
也不用質疑.
剛才說什麼.
他什麼時候舉手就贏.
一垂手就輸了.
所以亞倫和會議多重要.
撐著他.
否則怎麼行呢.
在真正的路上.
我和你都需要有同行者.
你不可能單打獨鬥.
因為會累.
我們真的是人.
有很大的限制.
摩西一個這麼偉大的領袖.
上帝大大重用他也好.
他也需要同行者.
最重要.
耶和華的「尼西」是什麼意思.
耶和華的「禎奇」代表什麼.
上帝在這裡.
上帝同在.
上帝的權能.
上帝的保護在這裡.
上帝在這裡爭戰.
今天沒有上帝的爭戰.
其實整場戰爭會徹底失敗.
所以關鍵在於上帝自己.

$^{521}$到最後摩西要捉一座壇.
起名叫「耶和華尼西」.
他指著耶和華的寶座.
發誓耶和華必世世代代.
和亞馬力爭戰.
在那一天.
原來第三個層面是上帝和亞馬力爭戰.
德性的關鍵在這裡.
如果在那一天的戰爭裡.
上帝是不會出現的.
摩西即使高舉權杖.
上帝不在.
就會輸.
他舉手也好.
他放手也好.
要輸在下面都會輸.
因為上帝不在.
摩西的禱告才有意思.
他才可以托著整個以色列民.
整個以色列的軍隊.
所以那座壇.
不是紀念什麼?.
不是紀念約書亞.
我們那位偉大的將軍約書亞.
他打仗真好.
所以我建了一座壇紀念約書亞.
世世代代都要記住約書亞打仗多厲害.
不是.
也不是紀念摩西.
大家記住摩西在山上舉手.
他舉手多重要.
撐著整個軍隊.
沒有了摩西.
這場仗就輸了.
不是.
只是紀念耶和華.
是與以色列民同在的上帝.
是他的貞旗.
他保護以色列民.
他同在.

$^{561}$他顯出他的權能.
他才是德性者.
他令這場仗勝利.
所以要紀念是紀念他.
這三場仗缺一.
這三個層面缺一不可.
沒錯.
上帝是主宰勝負.
他也是最大的保護.
但是上帝都會透過摩西去禱告.
宣告他的權能同在.
上帝會用摩西.
上帝不會這樣撇開摩西.
上帝自己在這裡.
約書亞你打吧.
不是,他用摩西.
在過程中.
上帝自己這麼厲害.
自己一個打吧.
你用什麼摩西祈禱.
用什麼約書亞打.
簡直多此一舉.
你自己來打.
一下子就打.
聖經都有記載.
在以色列戰爭的歷史裡.
有一場仗.
根本以色列民沒有出過手.
上帝就殺敗敵人.
很多很多的人.
有這樣的記載是可以的.
但是上帝不想這樣.
上帝想我們參與.
所以他給約書亞有機會去打仗.
帶領軍隊.
去經歷上帝.
他也給摩西去禱告.
給每一個在裡面經歷.
原來是上帝令我們得勝的.
這個是非常非常重要.

$^{601}$今天我們也有爭戰.
我們有爭戰.
我們有個人生命的層面.
我們個人生命有經歷很多的困難.
很多的傷害,痛苦.
我們很容易灰心,失望.
可能根本在生命上.
生活上走不下去.
更不要說為主作供.
如果我們要按照神的心意.
這樣去生活下去.
又要為上帝去打仗的話.
其實我們整個生命要起來.
其實我們也要為自己去爭戰.
今天有一個很特別的見證.
很想跟大家說的.
但是這個見證牽涉到第二點.
第二點是我們的教會也在爭戰當中.
教會要為主作炎作光.
要帶領更多人歸向耶穌基督.
要去守護更多的生命.
要醫治更多的生命.
其實也需要爭戰.
今天大家看回.
其實我們洪恩家也需要有很多方面去重建.
但我們也知道我們有很多挑戰.
我們今天也需要打仗.
我深信弟兄姊妹是深深明白這一點的.
我們有很多很多很多的需要.
必須要透過打仗.
將這一切重奪回來.
歸於耶穌基督的掌管下.
第三點是我們為了洪水橋區.
靈魂的得救去爭戰.
在這裡有很多人未信主.
我試過有一次我們教會出去.
有街頭的報導.
我記得那天跟一位姊妹一起拍檔.
排單張講福音.
那天吃了很多閉門羹.

$^{641}$遭受了很多很多拒絕.
但我覺得很感恩.
上帝也讓我看到.
這一區的人現在的光景如何.
需要怎樣.
原來要拯救這麼多人的靈魂.
是需要打仗.
不簡單.
不是直接走過去跟他說.
耶穌愛你.
講個福音.
他就會相信.
這裡有很大的需要.
我們更加要為了香港去爭戰.
在這裡過去兩年發生了很多事情.
我們很想讓香港重新歸於.
耶穌基督的掌管下.
讓上帝在這裡掌王權.
這個也不容易.
是要爭戰.
在這麼多的爭戰裡.
在這個過程裡.
我們要思考我們的位置在哪裡.
今天我一早很特別.
我在Facebook上看到一個很奇妙的見證.
是講爭戰的.
令我心靈為之一振.
講到有一個外國的姐妹.
叫做薇琪佐拉迪.
她剛剛跟丈夫默克.
慶祝完她六十歲的生日.
她是一位老師.
她一路身體都很健康.
在六十歲生日那一年.
她也拿了一個最佳老師的獎.
所以她的事業非常成功.
當然我深信這對夫婦都很愛神.
正在上教會.
她丈夫也很敬畏上帝.
剛才說到她很健康.

$^{681}$怎樣健康呢?.
應該是不需要怎樣請病假的.
她累積病假的時間有幾百天.
不需要請病假.
很好.
但是過完生日之後.
她就患了一個很急性的疾病.
情況急轉直下.
有一天她的生命有很大的危險.
在未來的二十四小時.
都是非常關鍵.
她有機會離開這個世界.
她丈夫默克在禱告當中呼求主.
「主啊,如果薇琦要活下去.
除非你嫁人.
否則她真的有生命的危險」.
這位先生聯絡了牧師,小組長,家人和朋友.
一起來禱告.
那天早上也有五十名婦女聚集在教會.
一路為這個姐妹禱告.
薇琦的情況起起伏伏.
一會兒好一點.
一會兒又陷入低谷.
醫生覺得已經束手無策.
覺得如果她的生命要結束.
就這樣吧.
就放手吧.
其實對於醫生來說已經無能為力了.
但是作為丈夫的默克就沒有放棄.
她一路召集人去禱告.
雖然薇琦自己已經陷入半昏迷的狀態.
但她仍然心裡禱告.
這是很特別的.
上帝給她很大的力量.
有一個話題是.
丈夫在隔離房間和家人禱告.
唱詩歌的時候.
薇琦自己也在房間裡.
心裡默默地禱告.
接著說她見到一些很不尋常的東西.

$^{721}$她見到很不尋常的東西就是.
她抬頭一看見到什麼呢?.
這真是一場爭戰.
她看到一群天使.
是很巨大的天使.
在一場巨大的戰鬥中.
抵禦邪惡和死亡.
是天使在爭戰.
是上面的天使在爭戰.
接著她見到下面的天使也在爭戰.
我想這指的是她的家人們.
她的家人們,弟兄姊妹們.
朋友們在為她迫切地禱告.
在向天父禱告.
接著她自己也在禱告.
她的內心很平靜.
她不害怕.
你在這裡看到的是一個三層的戰爭.
她為自己去爭戰.
她去禱告,去爭戰.
她的家人們,朋友,親戚們在禱告.
就好像摩西在山上禱告一樣.
一直為她禱告,代禱.
然後你看到上帝在天使的最上面那層.
在最頂層一直在爭戰.
最後得勝.
本來醫生已經想宣告.
這位女士的生命已經到了盡頭.
但是很奇妙.
在醫生放棄的時候.
一個小時之後.
她是奇蹟般地康復.
在這場爭戰裡她康復起來.
她令所有人感到很驚訝.
雖然之後很特別.
她的手和腳都出現問題.
需要截肢.
但是上帝卻要成為她的手和腳.
因為上帝是這樣跟她說的.
上帝也使用她的生命.

$^{761}$去展開了截肢的福音工作.
她繼續為上帝去搶救靈魂.
帶領人歸向神.
這是一個很好的例子.
我們為我們自己去爭戰.
原來我們為自己去爭戰.
不只是自己去打.
也需要有代禱者.
更加需要上帝的同在.
我們為教會爭戰.
我們要重建教會.
要讓教會走在上帝的心意當中.
我們要突破種種的困難.
在過程中我們弟兄姊妹一起來禱告.
一起來禱告.
我們背後也需要有摩西.
繼續為我們禱告.
更加需要上帝與我們同在.
需要上帝與洪恩家同在.
我們為香港.
我們為洪水橋區的福音工作去禱告.
不只是我們禱告.
我們背後也要找人.
繼續為我們禱告.
也更加需要上帝自己與我們同在.
這樣.
這個爭戰才會得勝.
所以我們的爭戰是有前方的戰士.
要有代禱者.
我們需要上帝的同在.
今天我們作為前方的戰士.
我需要問一個問題.
再問約書,揀人去打仗.
今天我與你的生命.
我剛才問過.
我們裝備好了嗎?.
我們預備好了嗎?.
我們是不是神很合用的器皿呢?.
是上帝用得著我們的嗎?.
如果上帝用不著我們.

$^{801}$我們如何去打仗呢?.
我們不可以只是自己去打仗.
教會也不可以只是自己去推示宮.
你就這樣走出去.
洪水橋區說傳福音.
走吧!憑著一股熱誠,勇氣.
就傳吧!.
我肯定你一個是不行的.
我記得那時候.
即使我還沒有做傳道人.
去短孫中心.
去學一些傳福音的課程.
他們也說你需要有土伴.
你需要有人為你每次出隊去禱告.
你是在說打仗.
不是兒戲.
沒有人為你禱告.
沒有一個摩西.
有亞倫胡爾在那裡.
你如何可以打這場仗呢?.
靠自己的力量嗎?.
靠自己的血氣嗎?.
是需要的.
我們需要代禱者.
其實我記得.
應該是前天.
我參加安息禮拜的時候.
我真的聽到有間教會提到.
他們教會有代禱士工.
代禱士工應該不只是某一個人.
某幾個人就在代禱.
應該是一個團隊.
當教會有不同弟兄姊妹有需要.
教會有不同需要.
就會給一班弟兄姊妹.
整個團隊為教會禱告守望.
我心裡很感動.
這就是我們需要.
我們需要有代禱士工.
不只是某一個人做摩西.

$^{841}$是整個團隊.
團結就是力量.
合一就是力量.
弟兄姊妹.
我們有沒有準備好自己.
也做一個代禱者.
現在疫情.
我們沒有辦法只有實體的祈禱會.
但是我們有Zoom的祈禱會.
弟兄姊妹有沒有準備好自己.
回Zoom的祈禱會呢?.
不只是為了我們個人需要禱告.
有沒有為了教會.
為洪水橋區的福音工作.
為香港禱告呢?.
實在有太多人有需要.
弟兄姊妹.
我們的代禱是怎樣的?.
一兩分鐘的代禱.
幾分鐘的代禱.
十幾二十分鐘的代禱.
有沒有行前地.
一直一直一直堅持禱告?.
是祝福別人的生命嗎?.
值得我們去思考.
第三樣東西是神的同在.
要有神的同在.
我們所做的.
所求的一切.
都要合乎上帝的心意.
今天值得我們問的.
我們的一切.
是不是合乎神的心意呢?.
我們自己有多少個親近上帝.
明白上帝的心意是甚麼呢?.
我們是否按照上帝的心意求?.
我們是否按照上帝的心意去行呢?.
行在神的心意當中.
上帝才會與我們同在.
我今天真的很願意.

$^{881}$講完這堂道之後.
很渴望鼓勵弟兄姊妹.
去經歷上帝的同在.
經歷上帝的大能.
沒錯,我剛才說過洪海的經歷.
只有以色列民當日有.
但是我深信上帝.
可以以其他的形式.
不同的方法.
讓我們經歷祂的大能.
祂的同在.
祂的保護.
祂的看顧.
也會說尼西這個名字非常寶貴.
很願意我們天天宣告這個名字.
這個名字成為我和你.
很大的鼓勵.
很大的幫助.
我們一同去禱告.
讓我們今天重溫這場仗.
你讓我們看到.
你是多麼的愛我們.
你要我們參與征戰.
要我們有份去參與.
你不要我們只是做旁觀者.
但是上帝你也要我們的生命準備好.
以致我們才可以成為被揀選的.
今天求你操練我們的生命.
讓我們上讀經,禱告,親近你.
認識你.
上上認罪,悔改,離開罪惡.
上上敏銳聖靈的聲音.
你讓我們整個生命真的準備好.
是很適合被上帝你用的.
天父求你.
你也興起我們代禱的事工.
求你為我們預備代禱的人.
但也興起我們每一個做代禱者.
不是幾分鐘的禱告.
是恆切的,恆常的禱告.

$^{921}$天父求你幫助我們.
你興起洪恩嘉成為禱告的教會.
願意感動更多人回來禱告.
在疫情裡我們無法實體的回來.
你讓我們在Zoom裡認罪.
我們去禱告.
不單是為了個人的需要.
是為了整個教會.
為洪水橋,為整個香港.
天父上帝在這裡.
求你更讓我們看到你的權能.
讓我們看到你是這麼奇妙的神.
你勝過萬有.
你征戰的時候你必然得勝.
你讓我們有信心.
讓我們靠著你就有信心.
讓我們靠著你什麼都不怕.
今天在這個疫情底下.
求你讓我們看到.
要怕的不是疫情.
今天你讓我們看到.
更加需要的是重建我們和你的關係.
更加要讓我們明白你的心意.
讓我們走在你的心意當中.
天父上帝我懇求你去幫助我們.
紅衣家中弟兄姊妹.
很深的經歷你的權能.
你的同在.
你大大的復興我們.
你真的讓我們在洪水橋區.
為你作明亮的燈台.
在這裡只願意人認識你的榮耀.
而不是看到人的榮耀.
天父上帝願你的心意得到滿足.
奉耶穌基督的名字祈禱.
阿們.
\newpage



\section{詩篇 1:1-3}
\label{sec:jmOEpWc4LtY}
\textbf{2022.02.20 《有福的人》 詩篇 1\_1-3 (和修版) 張美薇牧師}
\newline
\newline
連結: \href{https://youtube.com/watch?v=jmOEpWc4LtY}{\texttt{https://youtube.com/watch?v=jmOEpWc4LtY}} ~~~~ 語音日期: 2022-02-22
\newline
\newline
\hyperref[sec:_xVz_wL7W7s]{\small{< < < PREV SERMON < < <}}
~
\hyperref[sec:index_chronic]{\small{[返順時目]}}
~
\hyperref[sec:index_scriptual]{\small{[返順卷目]}}
~
\hyperref[sec:PKTsSbC_DIE]{\small{> > > NEXT SERMON > > >}}
\newline
\newline
詩篇 1:1-3
\newline
\begin{longtable}{cl}
\hline
\hline
章節 & 經文 (和合本修訂版)\\
\hline
1:1 & \begin{tabularx}{0.7\textwidth}{X} 不從惡人的計謀，不站罪人的道路，不坐傲慢人的座位， \end{tabularx} \\ \\ \relax
1:2 & \begin{tabularx}{0.7\textwidth}{X} 惟喜愛耶和華的律法，晝夜思想他的律法；這人便為有福！ \end{tabularx} \\ \\ \relax
1:3 & \begin{tabularx}{0.7\textwidth}{X} 他要像一棵樹栽在溪水旁，按時候結果子，葉子也不枯乾。凡他所做的盡都順利。 \end{tabularx} \\ \\
[1ex]
\hline
\hline
\end{longtable}
$^{1}$「弟兄姊妹平安」.
在踏入2022年農曆新年第一個月裡.
在這裡祝賀基督教宣導會.
紅恩堂的弟兄姊妹.
新年蒙福 主恩更多.
我去看了一個四格的花生漫畫.
主角伴茶李在新年的時候.
他見到Lucy.
他說 Lucy 祝你新年快樂.
很自然的.
誰知Lucy並沒有同樣地回祝他新年快樂.
她反而問他.
喂 你說了就是了嗎.
然後他說 是不是你說了新年快樂.
就叫我在新的一年裡面真的快樂.
是不是你單單說了事情就成了.
這是不是一個保證.
是不是 是不是 是不是.
伴茶李沒有辦法.
只可以用一句.
哦 我的天來回應.
當我剛才祝賀大家新年蒙福的時候.
大家未必一定會和Lucy有同樣的反應.
不過不知道大家心裡所想到的有福.
是甚麼意思呢.
我最近在網上看到這幅圖畫.
是百感交集的.
香港的老弱竟然要這樣日以繼夜的.
在醫院門外等候入醫院.
新年快樂的願望對他來說.
可能是願望能夠好像過去一樣.
想入院就入院.
又或者上個星期.
在網上看到這幅圖是一位姐妹的女兒.
她要趕在2月10日之前.
去理髮師處理她的頭髮.
現在因為不知道甚麼時候.
可以再去理髮店那裡理髮.
原來隨意去剪頭髮.
也是一個快樂,是一個福氣.

$^{41}$所以,各位姐妹,福氣是甚麼呢?.
是否要靠別人的保證呢?.
是否要靠別人容許的呢?.
今天我們一起來看.
詩篇第一篇的第一至三節.
原來我們在詩篇第一篇裡面看到.
福氣是某個程度上是自己可以控制的.
我們看回詩篇第一篇第一至三節.
「不從惡人的計謀,不站罪人的道路,不坐洩萬人的座位.
為喜愛耶和華的律法晝夜思想,這人便為有福.
他要像一棵樹栽在溪水旁.
按時候結果指,葉子也不枯幹.
凡他所作的,盡都順利」.
我們看到詩篇第一篇.
原來我們立志做有福的人.
是可以自己控制的.
我們看回這裡.
雖然在和合本的聖經裡面.
有福在第二節尾才出現.
但是原來在預文.
或者我們不看預文.
我們不懂得預文.
我們看英文翻譯本的聖經.
其實有福是在第一節開頭的時候出現的.
而第一節是說.
有福的人有些事情是不做的.
而第二節是說.
有福的人有些事情是做的.
然後第三節是用一個比喻.
來告訴我們.
有福的人是怎樣的呢?.
所以讓我們慢慢逐節逐節去看.
有福的人不做,不做什麼呢?.
原來我們可以再說.
有福的人是要遠離惡人的.
有福的人要遠離惡人.
原因是用三個否定句來講.
不從惡人的計謀.
不暫聚人的道路.
不坐舌萬人的座位.

$^{81}$希伯來文的比較方式是這樣的.
他重複.
如果他不喜歡一件事.
就說我不喜悅.
如果他真的不喜歡.
就說我不喜悅.
我不喜悅.
如果他真的不喜歡.
就說我不喜悅.
不喜悅.
所以三個不.
不.
表示私人極之重視.
語重心長的提醒聽眾.
不要親近.
而是要遠離惡人.
在這三個字裡面我們看到.
其實是完成式.
已經過去的.
是過去式.
是強調.
《荊棘夜話》的人.
從來都不會涉及任何.
這裡講到的一件事.
我們首先可以從.
這三個形容詞來看.
就是惡人,聚人,舌萬人.
惡人其實是從他的計謀.
顯出惡人的愚昧.
他在心裡盤算著這些惡.
什麼叫做聚人呢.
聚人就是將思想裡面的惡.
付諸行動.
所以就說他們整個道路.
或者整個生活的方式.
都是敗壞的.
第三個舌萬人是什麼呢.
舌萬人就是那個人.
無視神和他的誡命.
他誤會應神的吩咐.

$^{121}$卻以他的無慢去引發爭鬥.
亦時時嘲笑.
不將惡看為惡.
更變本加厲地誇耀他的罪行.
而且輕視公義.
我們看完形容詞之後.
我們就看動詞.
有三個動詞.
不從,不站,不坐.
不從我們可以翻做不走.
最初開始的時候.
是思想上的問題.
想想惡行和惡事.
或者思想上跟隨惡人去行惡.
之後就要求他不要站在罪人的道路中.
站的意思就是將心裡的惡念.
付諸實行.
跟惡人一起行惡.
我們可以說跟惡人做同路人.
站的焦點就是在行為上的模仿和參與.
然後到最後就坐下.
坐下的話就即是跟惡人為無.
或者跟惡人為伍.
意思就是在幕後策劃.
甚或是指揮手下執行惡事.
我們再從名詞來看.
有三個名詞.
計謀,道路,座位.
其實剛才都說了.
意思是我們做福的人.
我們要不從惡人的計謀.
不站在罪人的道路上.
也不坐在蝕萬人的座位上.
我問你.
意思是我們要遠離惡人.
我想細蚊仔都知道.
不用大人教.
細蚊仔都知道.
我們要遠離惡人.
但容不容易做到呢?.

$^{161}$其實是不容易做到的.
為什麼呢?.
因為惡人的額頭沒有字.
就在那裡說我是惡人.
沒有啊.
沒有的話.
你怎知道他是惡人.
而要遠離他呢?.
所以第二節就教導我們了.
他給我們一個標準.
標準我們要喜愛耶和華的律法.
晝夜思想.
所以第二節就告訴我們.
有福的人要追求親近神.
為喜愛耶和華的律法.
晝夜思想.
這人便為有福.
第二節的文法上高.
是將第一節的完成式.
記不記得.
第一節是完成式.
第二節就變成不完成式.
現在時時都仍然要進行.
或者叫進行式.
表達喜愛耶和華的律法.
晝夜思想.
是要不斷進行.
是習慣性的.
習慣性的思想神的話語.
習慣性的在生活當中.
反映出神的話.
這是熱切的言讀神的律法.
被神的旨意充滿.
而且遵行神的命令.
我們可以去看看.
這裡的耶和華的律法.
這裡不單是講十誡.
不單是講摩西五經.
不單是講舊約或者整本聖經.
可以指是和神而來的訓誨.

$^{201}$啟示是來自神的.
而目的是要幫助人.
過一個符合神的生活.
符合神的旨意的生活.
在每一個時代.
敬虔的人都要照著神的啟示而過活.
並且以之為樂.
詩人在這裡不是在談論.
律法致死的影響.
我們知道律法是有兩面的.
一面是具有威嚇和懲罰性.
我記得當時我還是孫德堂的堂主任.
我是徐澤林的校監的時候.
我在元朗特別小心.
特別過馬路的時候.
明明前面的紅公仔已經亮了.
看來看去都沒有人.
但我不會走的.
為什麼呢.
因為如果被警察抓了.
你在行人紅燈的時候過馬路.
這樣不就是論準嗎.
因為律法是會有威嚇的.
你犯了律法就會被懲罰.
但其實這裡不是談論律法的威嚇性.
這裡就是重視律法賜人生命的面向.
信徒由按著律法的生活.
按著神的旨意.
新生命喜樂的去回應神的默示.
和回應聖靈的帶領.
信徒的喜樂不單止是認識.
言讀背誦神的話.
尤其是以遵行神的旨意為樂.
而不是受惡人的欺騙.
這裡更加提到我們要喜愛耶和華的律法.
我們要表達我們要愛神.
我們很多時候都將愛神.
好像今天的情感化.
我們說我們要愛一個愛人.
或者愛的伴侶.

$^{241}$愛神就是一個這樣.
是有情感的,有感性的愛.
沒錯 我們要學習用感性來愛神.
但是我們更加要表達.
我們要喜愛,敬畏耶和華.
是敬畏神的律法上.
聖經人物參孫是一個很好的例子.
我們每個小朋友都認識.
所以我們做過大人,教過主一學.
我們都很熟悉參孫的故事.
參孫在庭那遇見一個非理士的女子.
他喜歡她,愛她.
願意娶那個女子為妻.
當時的人都會問父母的意見.
父母當然不贊成.
因為她是外邦的女子.
參孫說什麼呢?.
我喜悅她呀!.
結果他就娶了那個女子.
這已經注定參孫的失敗.
後來因為他喜歡了另一個非理士的女子.
他最後就失去了他的靈力.
剪了他的頭髮.
落在仇敵的手中.
參孫深愛外邦女子.
他的腳就跟著那個外邦女子.
去了外邦裡.
心裡所愛就決定他的腳zoung .
人如果愛世界,愛情慾,愛物質,愛罪惡.
就決定他的腳zoung 要走世界的道路.
要走情慾的道路.
走物質的道路.
走罪惡的道路.
決定他將來的結局.
什麼是你和我所愛的呢?.
你愛神嗎?.
你愛神的話語嗎?.
你怎樣愛神的話語呢?.
我們看看下一道.
這裡說要晝夜思賞神的話語.

$^{281}$一般人很難擁有聖經.
在古時的時候.
他們只可以背誦聖經.
而時時都思賞.
但我們知道到今天.
我們已經不知道你有多少部聖經.
就算在我的辦公室裡.
起碼有超過十本不同版本的聖經.
但默想神的話語仍然很重要.
背誦神的話語,默想神的話語.
然後思賞神的話語.
是很重要的.
記不記得我們說.
從第一節的完成式.
去到第二節的未完成式.
或者我們叫現在進行式.
在現在進行式裡.
我們要不斷的思想.
不斷的喜愛.
我們要以切然讀律法.
不同的默想.
叫我們認識神的話更加深.
好像讀經.
我不知道你們讀完整本聖經多少次.
我記得我信主的時候第一年.
我去讀聖經.
我是3月29日信主的.
我去到4月初左右.
就去到出埃及記.
就困在那裡.
出埃及記有很多東西.
很難很難明白.
特別去到利未記.
困在那裡.
很難.
但後來我強行衝過了.
然後就繼續讀下去.
到最後又去到以賽亞書.
又困在那裡.
但我都讀完它.

$^{321}$當我第二年再讀第二次的時候.
原來我第一年讀到的資料.
幫助我去明白第二次讀聖經的時候.
一些的內容.
當我讀完第二次.
讀第三次.
讀第四次.
讀第五次的時候.
越讀就會越清晰.
其實記在心中都可以的.
幫助我們需要的時候.
我們對場景可以有即時的.
適當的回應.
結果我們就會.
當我們時時去默想神的話的時候.
我們就會被祂的旨意充滿.
而且我們時時都可以出來.
行祂的命令.
所以默想不單單是早上靈修.
或者晚上靈修.
而是我們在日常生活的過程當中.
不斷的反覆思想神的話.
相信大家都很熟悉.
《申命記錄》第七節.
無論你坐在家裡.
走在路上的屋子.
躺在床上起來.
都要談論.
談論什麼呢?.
談論神的律法.
所以當摩西去教導以色列人的時候.
他們要在行走站立坐的時候.
都思想上主.
要將遵行神的旨意.
變成一個習慣.
好不好熟悉?.
和第一節不一樣.
我不走在錯誤的當中.
我不站在罪人的道路上.
我又不坐在涉罪人的上.

$^{361}$但我們是做什麼呢?.
我們就站在,坐在,行在神的律法上.
應用神的話.
是需要個人的努力的.
但其實我們都需要弟兄姊妹的支持和鼓勵.
我不知道紅磡堂的弟兄姊妹.
會不會在團契或者會眾的層面.
推行一些靈修運動.
彼此的提醒.
記得有一個團契的做法是這樣的.
彼此的協議.
在傍晚的時候.
有弟兄或姐妹會致電提醒.
喂!你今天做了靈修了沒有?.
我們怎樣做一個有福的人呢?.
喜愛耶和華的律法.
在這裡指出有福快樂.
是必須有所付出,有所行動的.
在消極方面.
我們不做這件事.
我們要遠離惡人,遠離惡行.
但相反在積極方面.
我們要一心一意地跟隨神.
並且喜愛祂的誡命律例.
時時在心裡反覆思想.
所以敬虔的人.
在他們行走,躺臥,起來的時候.
都熱愛神.
所以無論在家中或者在路上.
在他們所走的路上.
在他們所有的活動當中.
他們都要遠離不敬虔的人.
免得受到他們的影響.
他們在家庭事業和社交關係上.
小心保護自己.
設立自己的人際關係和一個界線.
我以前看了一本書.
叫做《In His Steps》.
即是跟隨他的腳踏中行.
這本小說裡面有一個問題.

$^{401}$每一次當你需要一個決定要做的時候.
他的問題是.
What would Jesus do?.
即是W.W.J.D.
不知道大家有沒有這樣的手鏈.
如果耶穌在的時候.
祂會怎樣做呢?.
我們會思想.
我們去問.
如果我們知道答案又怎樣呢?.
就去做.
所以我們會知道.
原來有福的人.
是不做一些事情.
然後再做一些事情.
這就是我們可以控制.
我們當然知道.
最後有沒有福氣臨到呢?.
其實都需要有另外神的幫助.
所以第三節就會說到.
有福的人好像一棵樹.
他說一棵什麼樹呢?.
他說他好像一棵樹栽在溪水旁.
我在網上找到這樣的圖畫.
按時候結果指.
葉子也不枯幹.
凡他所作的盡道順利.
這裡是很好的.
他讓我們看到.
原來有福的人好像一棵樹.
一棵樹是怎樣呢?.
又結果指.
另外樹葉又常青.
當我們看到一棵樹會遇到什麼困難.
什麼挑戰呢?.
我們知道不是每棵樹都能夠結果指.
不是每棵樹,樹葉都常青.
我們又看到聖經有另外一個地方.
告訴我們.
原來樹會面對什麼挑戰.

$^{441}$在耶利米蘇第十七章第七到八節說.
「倚靠耶和華,以耶和華為可靠的,那人有福了.
他別像樹栽於水旁,在河邊紮根.
炎熱來到並不懼怕.
葉子仍必清脆.
在乾旱之年,毫無掛慮.
而且結果不止」.
所以告訴我們.
當我們可以用耶利米蘇這裡.
讓我們更加清晰地看到.
「要像一棵樹栽在溪水旁」.
意思是什麼呢?.
意思是這棵樹是有主人的.
是因為什麼呢?.
是被栽在溪水旁.
變成一個被動的魚.
我們發覺第一,第二節都是主動的.
第一節是過去式.
即是我們現在完全不會和罪人,惡人,褻瀆人相來往.
第二節就是不斷進行式.
現在仍然是進行的.
我們喜愛耶和華的律法.
晝夜思想.
我們就為有福.
而第三節告訴你.
這是一個被動魚.
這棵樹是被栽在溪水旁.
或者我們看呂鎮忠的譯本.
他說是移植的.
是被移植到溪水旁.
意思是移植到一個地方在水邊.
和一般在野地凹陷的地方種植.
或者在田野的樹木是不一樣的.
因為那裡仍然距離水源遠.
但這裡說是一個有福的人.
他就好像一棵樹被栽在溪水旁.
樹會遇到什麼難處呢?.
我們會看到.
第一就是炎熱.
其實春夏秋冬四季是運行不息的.

$^{481}$春天之下就夏天.
夏天之下就秋天.
秋天就冬天.
無論怎樣都會有些時候是比較炎熱的.
夏天來到的時候.
無論多熱.
但當那棵樹還在水邊的時候.
他們怕不怕?.
不怕.
因為多熱都沒有關係.
因為葉子是靠著根達到水邊那裡.
吸到水來滋養著.
所以葉子是不枯乾.
而第二個天敵是什麼?.
是乾旱.
他說在乾旱之年毫無掛慮.
弟兄姊妹我們今天都知道.
我們今天是乾旱之年.
我們大家都遇到很多的困難.
很多的難處.
疫情是很厲害.
乾旱之年不是經常有的.
是偶發的.
好像人生一樣.
我們有純景.
有逆景.
當乾旱臨到的時候.
雖然天沒有下雨.
但是因為有樹的根可以伸到溪水的旁邊.
水的供應仍然源源不絕.
所以仍然能夠結果不止.
有一個很真實的例子.
在倫敦的皇宮裡面.
有一棵葡萄樹.
初栽種的時候.
每年都很茂盛的葉子.
很漂亮的.
但是每年都結不到果.
園丁當然是很失望.
突然有一年結了很多的果子.

$^{521}$那時候園丁就很奇怪.
要研究一下是什麼原因.
令到這個栽種在倫敦的皇宮裡面的葡萄樹.
可以結這麼多果子.
於是他就挖開泥土.
發現原來這棵樹的根.
越長越遠.
去到太吾士河那裡.
在太吾士河的邊.
就發現有活水去吸收.
去供應新鮮的活水.
永留不息的活水.
所以它的果子可以是很多.
如果我們時時在主的恩典裡面.
當然就是一個蒙福的人.
在世上所蒙的福分.
沒有一樣比時時在基督裡面更加穩妥.
所以我們所作的都是順利.
這是一幅又美麗又快樂又安心又順利的圖畫.
所以我們想一想.
我們今天讀了這三節經文.
很簡單的三節經文.
我專門選這三節.
讓大家可以背.
有福的人是怎樣呢?.
不從惡人的計謀.
不暫聚人的道路.
不助洩萬人的作為.
我們不做一些事.
為喜愛耶和華的律法.
晝夜思想.
這人便能有福.
所以我們每個人都可以做.
我們不做一些事.
然後我們喜愛耶和華的律法.
晝夜思想.
但我們知道不是單靠我們自己的才能和能力.
而是靠著神.
我們靠著神的時候.
就好像一棵樹被栽在溪水旁.

$^{561}$安事後結果子.
葉子也不富幹.
凡他所作的盡都順利.
講完的時候.
大家也可以背這三節經文.
我想舉個例子和大家去看.
我知道其實紅春堂.
你們也有很多宣教上的一些人物.
你們也很想去看.
在你們的周刊裡.
時時都有一些宣教的例子.
今天和大家分享一個宣教士.
他叫做博恩德.
博恩德是十九世紀末的人.
出生年是1887年.
他生在美國的伊利諾州.
一個很有錢的家人.
他媽媽大概在他七歲的時候信主.
然後到他高中畢業.
大概十六歲的時候.
家裡給了他一份生日禮物.
大概是一九零幾年.
他父母送了一部聖經給他.
然後給了錢他去環遊世界.
今天大家都很想.
我到今天都沒有試過環遊世界.
我不知道大家有多少人環遊世界.
他一個高中畢業生.
就已經環遊世界.
他遊歷了亞洲,歐洲,中東.
很多的城市.
他看到很多人很富麗堂皇.
但也看到很多人很窮乏的生活.
在生活裡沒有盼望.
於是神就感動他.
他立志要做一個宣教士.
他要將福音傳給這群這麼大需要的人.
他更加想將福音傳給.
當時仍然在世界上很大的宣教工場.
中國.

$^{601}$他更加想傳給中國的穆斯林.
成為宣教士對他來說是不容易的.
他要放棄一個原本奉意足食的生活.
很多人都勸他不要放棄這麼好的前程.
經過一段時間和內心的掙扎之後.
他就寫信告訴父母.
我將我全部的生命.
要進入宣教工場作預備.
同時他為了堅定自己的奉獻心志.
在他爸爸送給他的聖經上.
親筆在最後一頁寫了.
No Reserve 即是沒有補償.
當然他只是中學畢業.
還要做很多的預備.
所以他回到美國之後.
他就進入了當時很知名的耶魯大學.
我想大家都有聽過.
現在仍然是很出名的.
大家都很想.
如果我兒子可以去耶魯大學讀就好了.
很快他進去之後.
他就覺得這個名校對他來說很失望.
因為他覺得學生的道德很低落.
整個校園都充斥著各種罪惡.
但博德恩沒有因為這樣就袖手旁觀.
相反他開始和幾個好朋友約定.
在早餐之前一起讀聖經.
一起祈禱.
很奇妙.
越來越多人聽到這個祈禱的消息.
並且選擇加入小型的祈禱會.
祖禱會就日漸擴大.
成為全校的運動.
當博德恩大一結束的時候.
祖禱會的每週聚會人數.
你猜到多少人?.
150人.
不簡單.
是到150人.
然後會繼續到第二年.

$^{641}$第三年.
第四年.
到四年級到畢業的時間.
你猜有多少人一起祈禱?.
每個星期的祈禱會.
一起祈禱的人數達到1300人.
我們看到.
其實神是使用他.
他也很懂得去動員他附近的人.
神的呼召一直在他的心裡.
仍繞不去.
他心中很明白.
神要他去中國的穆斯林.
將福音傳給他們.
他大學畢業之後.
其實有很多大公司和學校.
爭相用高薪請他.
但他拒絕了.
他堅守起初神給他的感動.
他申請進入普林斯頓大學裡面的神學院就讀.
因為當他知道他的申請被接納的時候.
他就在聖經最後那頁.
寫下他第二個.
他第二個就說.
No retreat.
我不會退縮的.
我不會退後.
他不保留.
不退縮.
完成了普林斯頓神學的訓練之後.
他開始為他預備去中國做準備.
由於當時.
其實現在都是的.
穆斯林共同的語言是阿拉伯語.
穆斯林是不會將他們的聖經翻做外語的.
所以一直都是用阿拉伯語來做他們的聖經.
所以他決定首先去到埃及去學這些阿拉伯語.
然後才去到中國的甘肅向穆斯林宣教.
去到1913年初的時候.
伯德恩來到埃及的開羅.

$^{681}$希望在這裡學好他的阿拉伯白文.
並且親自向陰霧很久的茨維謀.
這位向回民宣教的前輩去學習.
他每天在開羅的街上都會分發福音單張.
特別留意到當中有沒有中國人的回教徒在當中呢.
有一天他在日記裡面提到.
終於在回教世界的最高學府.
即是埃及是一個回教的國家.
愛茲哈爾大學見到一位來自甘肅省的中國留學生.
他很開心很興奮.
不過很可惜他去了埃及不久就患了腦膜炎.
一個月之後就離世.
當年他只有25歲.
1887年去到1913年.
伯德恩逝世的消息很快傳回美國.
所有的報紙都刊登了這一則消息.
因為他是一個很有錢的家族.
很有前途的年輕人.
但他卻願意放下他這些一切.
然後踏上宣教的路.
但是他竟然出師未捷身先死.
一般人看來是一宗很大的悲劇.
是一個很大的遺憾.
但當伯德恩的父親收拾兒子遺物的時候.
他翻閱很多年前給兒子的聖經的時候.
那裡寫了「不保留不退縮」.
原來在「不保留不退縮」之下.
伯德恩更親筆地留下了「不後悔」.
No regret 不後悔的三個字.
長話短說 伯德恩的身家是很多的.
他有一百萬美金.
記住是1913年時代的一百萬美金.
是很多錢.
他全部捐出來 他的家人幫他全部捐出來.
其中有四分之一 即是25萬.
捐給向中國傳福音為宗旨的中國內地會.
即是今天的OMF.
內地會就用這一筆錢.
在中國的甘肅省蘭州市成立了一間.
服務穆斯林為主的醫院.

$^{721}$名為「伯德恩福音醫院」.
後來成為今天的蘭州市第二人民醫院.
伯德恩的死也成了激勵了很多的基督徒.
前往中國宣教.
為他毫無遺憾的人生畫下了很完美的句號.
因為不單是他去.
他更可以一代一代有人跟著他.
其中想和大家介紹一個內地會的宣教士.
海春心就是一個例子.
海春心在1916年.
記不記得伯德恩是何時死的.
1913年.
即是說他在1916年4月的一個晚上.
他正在參加普告會.
正在專注祈禱的時候.
突然在他意識當中有個聲音問.
無論什麼地方你都願意為我去嗎?.
海春心就說.
是呀 主呀.
什麼地方我都願意為你去的.
然後主立刻讓他的思想上到.
中國西北部的穆斯林的地方.
然後他就說我知道這是主所指示的.
無論什麼地方.
後來海春心才知道.
原來那一天晚上剛好是伯德恩逝世三週年.
就是在他逝世三週年的紀念晚會上.
慈惠謀和伯德恩的媽媽.
還有一些朋友聚在一起.
在紐約一起祈禱.
求神呼召一些年輕人.
繼承伯德恩的未竟之志.
向中國西北穆斯林宣教.
海春心來到中國西北的西陵.
去服侍回民.
正是神回應了他們的禱告.
一棵樹栽在溪水旁.
安事後結果子.
葉子也不枯幹.
一代一代一代的傳下去.

$^{761}$弟兄姊妹.
相信從2021踏入2022的時候.
大家都收到很多WhatsApp的訊息.
不知道有沒有收到這個呢.
好像大家都覺得2021真的有很多困難.
很多不開心.
很多不好的事.
好像大垃圾是不好的.
我們把它扔掉.
我們等待2022年來.
所以我送給弟兄姊妹的訊息裡面.
我送了這幅圖畫.
由2021跨過去到2022的時候.
我們有什麼呢.
我們主內平安.
恩典滿日.
永享恩福.
皆主同行.
意思就是說不是轉了年就可以了.
而是我們要有主.
我在網上看到另外一幅圖畫.
很有趣的那幅圖畫.
就是說新的一年.
並不會帶來新的希望.
又真是.
2022和2021進來的時候.
或者2020年進來的時候.
或者2019年快要來的時候.
都是大家一樣.
都是對新的一年.
好像有很多的盼望.
很多的希望.
但是呢.
當我們去看清楚的時候.
什麼東西會給我們希望呢.
只有那位做新事的神.
會給我們希望.
因為祂會在新的一年裡面.
給我們新的恩典.
新的同在.

$^{801}$新的盼望.
所以就好像我每次講道.
我都會在這樣的結束.
你聽了大概三四十分鐘的道.
對你有什麼關係呢.
你想做一個有福的人.
原來我們可以自己控制的.
怎樣控制呢.
我們不從惡人的計謀.
不暫聚人的道路.
不助洩萬人的座位.
為喜愛耶和華的律法.
晝夜思賞.
這人便為有福.
當我們去按著這樣去做的時候.
神就會賜福給我們.
我們就好像一個被栽在溪水旁的人.
或者一棵樹被栽在溪水旁.
安事後結果子.
孽子也不富幹.
凡他所作的盡都順利.
我們只要做了我們的部分.
神就會用祂的部分來賜福給我們.
我們低頭祈禱.
讓我們在一年之初的時候.
能夠重新再次思想.
我們有什麼是不應該做的.
有什麼是做的.
我們就等待你給我們的同在.
你給我們的恩典.
我們現在面對一個很嚴峻的環境.
疫症橫行.
我們不知道第五波的疫情.
什麼時候才會跌下去.
但我們卻知道.
我們只要做好我們的部分.
你就會做你的部分.
在你裡面.
一切都有你的恩典.
有你的旨意.

$^{841}$願你的旨意成就.
聽我們的祈禱.
奉耶穌基督的名求.
阿們.
\newpage



\section{詩篇 27:1-14}
\label{sec:PKTsSbC_DIE}
\textbf{2022.02.27 《基督光--燃亮幽暗心》 詩篇 27\_1-14 (和修版) 陳麗蟬牧師}
\newline
\newline
連結: \href{https://youtube.com/watch?v=PKTsSbC-DIE}{\texttt{https://youtube.com/watch?v=PKTsSbC-DIE}} ~~~~ 語音日期: 2022-02-27
\newline
\newline
\hyperref[sec:jmOEpWc4LtY]{\small{< < < PREV SERMON < < <}}
~
\hyperref[sec:index_chronic]{\small{[返順時目]}}
~
\hyperref[sec:index_scriptual]{\small{[返順卷目]}}
~
\hyperref[sec:jiG98Pa_BV4]{\small{> > > NEXT SERMON > > >}}
\newline
\newline
詩篇 27:1-14
\newline
\begin{longtable}{cl}
\hline
\hline
章節 & 經文 (和合本修訂版)\\
\hline
27:1 & \begin{tabularx}{0.7\textwidth}{X} 耶和華是我的亮光，是我的拯救，我還怕誰呢？耶和華是我生命的保障，我還懼誰呢？ \end{tabularx} \\ \\ \relax
27:2 & \begin{tabularx}{0.7\textwidth}{X} 那作惡的就是我的仇敵，前來吃我肉的時候就絆跌仆倒。 \end{tabularx} \\ \\ \relax
27:3 & \begin{tabularx}{0.7\textwidth}{X} 雖有軍隊安營攻擊我，我的心也不害怕；雖然興起戰爭攻擊我，我仍舊安穩。 \end{tabularx} \\ \\ \relax
27:4 & \begin{tabularx}{0.7\textwidth}{X} 有一件事，我曾求耶和華，我仍要尋求，就是一生一世住在耶和華的殿中，瞻仰他的榮美，在他的殿宇裡求問。 \end{tabularx} \\ \\ \relax
27:5 & \begin{tabularx}{0.7\textwidth}{X} 因為我遭遇患難，他必將我隱藏在他的帳棚裡，把我藏在他帳幕的隱密處，將我高舉在磐石上。 \end{tabularx} \\ \\ \relax
27:6 & \begin{tabularx}{0.7\textwidth}{X} 現在我得以昂首，高過四面的仇敵。我要在他的帳幕裡歡然獻祭，我要唱詩歌頌耶和華。 \end{tabularx} \\ \\ \relax
27:7 & \begin{tabularx}{0.7\textwidth}{X} 耶和華啊，我呼求的時候，求你垂聽我的聲音；求你憐憫我，應允我。 \end{tabularx} \\ \\ \relax
27:8 & \begin{tabularx}{0.7\textwidth}{X} 你說：「你們當尋求我的面。」那時我的心向你說：「耶和華啊，你的面我正要尋求。」 \end{tabularx} \\ \\ \relax
27:9 & \begin{tabularx}{0.7\textwidth}{X} 求你不要轉臉不顧我，不要發怒趕逐你的僕人，你向來是幫助我的。救我的神啊，不要離開我，也不要撇棄我。 \end{tabularx} \\ \\ \relax
27:10 & \begin{tabularx}{0.7\textwidth}{X} 即使我的父母撇棄我，耶和華終必收留我。 \end{tabularx} \\ \\ \relax
27:11 & \begin{tabularx}{0.7\textwidth}{X} 耶和華啊，求你將你的道指教我，因我仇敵的緣故引導我走平坦的路。 \end{tabularx} \\ \\ \relax
27:12 & \begin{tabularx}{0.7\textwidth}{X} 求你不要把我交給敵人，遂其所願；因為妄作見證的和口吐凶言的都起來攻擊我。 \end{tabularx} \\ \\ \relax
27:13 & \begin{tabularx}{0.7\textwidth}{X} 我深信在活人之地必得見耶和華的恩惠。 \end{tabularx} \\ \\ \relax
27:14 & \begin{tabularx}{0.7\textwidth}{X} 要等候耶和華，當壯膽，堅固你的心，要等候耶和華！ \end{tabularx} \\ \\
[1ex]
\hline
\hline
\end{longtable}
$^{1}$有一個故事,不知道你們聽過沒有,就是講到一隻鬼王,他吩咐小鬼,去到某一個城市裡面,他要用這個病毒去殺死二千人..
但是那個小鬼出去了五天,竟然殺死了五千人,他回來向鬼王匯報,他告訴鬼王,他說不關我的事,我只是用病毒去殺死了二千人,其他的三千人是他們自己太過擔心,自己嚇死自己..
有人說這個憂慮是現代人的病,也有醫生說這個憂慮令人心靈裡面產生了懼怕,產生了沮喪的心理,其實是健康的大敵..
工作忙,工作很有壓力,不會殺死人,最能夠殺死人的就是這個憂鬱..
憂鬱症的成因是怎樣的呢?就是生活或者工作上遇到很大的壓力,例如失業,或者感情受到創傷,婚姻又失敗,學業又失敗,人際關係又出現這個問題..
這一切一切的事令人感到很沮喪,很擔心,甚至很憤怒,所以種種的負面情緒就會出現,但當我們不懂得怎樣去處理這些情緒的時候,很容易就患上這個憂鬱症..
剛才主席為我們讀出詩篇27篇,作者是大衛,在他寫這個作品的時候,其實他都經歷了很多很多的艱難,他經歷過被這個掃羅的追殺,又經歷過被自己的兒子阿沙隆的追擊,.
還有很多很多次的逃亡,很多很多次的打仗,但每一次他都靠著倚靠耶和華的時候,他就能夠經歷過了..
大衛的人生經歷裡面,可以說一輩子都很慘,如果他不懂得處理他的情緒,不懂得去仰望上帝的時候,他一定是一個很憂鬱的人..
但我們從詩篇27篇裡面,我們看到的大衛是一個很堅忍,是一個對上帝充滿了信心的人,是我們要學習的榜樣..
我們先看第一方面,就是他人生的經歷裡面,是怎樣令到他的心靈是多麼的幽暗..
在這個詩篇27篇裡面,講到他的困難有多少呢?.
在第二節他說,那作惡的就是我的仇敵,前來要吃我的肉..
我很怕狗的,所以有時候去探訪的時候,那些大狗撲過來的時候,我會大叫救命的..
雖然主人告訴我,他們是來歡迎我的,但是我有一個感覺,他們就撲過來要抓我,很害怕的..
但何況大衛呢,他那些的仇敵,真的想將他置諸死地..
第三節又說,他的仇人多到怎樣呢?好像軍隊一樣,來攻擊他..
第五節又說到,他遭遇很多的患難..
到第十節,他說到他被父母離棄..
其實大衛沒有被他的父母離棄的,他只不過想形容當他被蘇羅追殺的時候,.
他要逃難,要打仗,在逼不得已之下,他將他的父母安置在非利士人的營裡面..
其實是一個大恥辱,他自己照顧不到父母,這只是一個形容..
形容真是最愛的父母都保護不了他..
在第十二節又說到,當中有些人是作一些假見證,甚至說一些很惡的很大的聲音,來去控他,來去攻擊他..
大衛接二連三的受到攻擊,受到重傷,可以說是一波大過一波..
心裡面的壓力,憂悶,厭煩,可想而知,都不知道什麼時候才能夠去停止..
我想現在香港人,或者全世界人多少都體會得到,在這個疫情底下,又沒有公開,很多時候都不准時出糧,又不准旅行,又不准多人聚會..
現在又說出街吃飯,又要在六點鐘之前吃完,還有一桌子又不可以多過兩個人,還要一定要去打針,然後才可以去到這些公共場所..
真是令人很擔心..
有位姐妹跟我分享,她很怕打針,但又要去打,所以她去到打針的中心裡面,在牆外邊,她就不斷地徘徊,到時間都不敢進去,直至她的兒子催她,媽,你還不進去就過了時間..
進到中心的時候,對齊了文件,排隊準備去打了,她又頻頻地去洗手間,等了十幾二十分鐘,去了五六次洗手間..
她心裡想著,原來她不自覺,一直想逃避,一直怕打針,不知道會帶來她身體有什麼副作用..
有一位導演,張堅庭導演,他很多時候都做喜劇,我想大家比較容易知道的,有一套叫做《表姐,你好ye 》..
他講到在他人生裡面,他自己都經歷了一場笑劇,就在2001年的時候,他和成龍在泰國拍戲,他發覺他的食指不受控制地動來動去,他就開始懷疑自己是不是患了柏金遜症..
他非常擔心,所以一個月之內,無論在泰國還是回到香港,他看了六七個專科醫生,有腦科的醫生,有精神科的醫生,有免疫系統科的醫生..
這六七個專科醫生都告訴他,都是同一個答案,我們查過很多東西,都證實不了你有柏金遜,不過有可能是柏金遜的先兆..
所以張堅庭就覺得,有沒有搞錯,個個都查不到,他繼續陷在恐懼裡面,他很怕自己患上這個症,.
他過分的擔心,他對著自己平時很喜歡吃的東西,對著很珍貴的食物都覺得淡而無味,更慘的是他每晚只能睡兩三個小時,.
最嚴重的一次,就是他只能睡兩三個小時,一晚竟然發惡夢醒了十次,他完全睡不著覺,結果在三個月裡面,他瘦了三十多磅,還怕,不知道是不是有這個癌症呢?.
因為張堅庭導演,他以前無論在他讀書,或者他的工作,婚姻各方面都是一帆風順的,但到現在他的健康不知道什麼事,.

$^{41}$以致他陷在一個很低谷的裡面,他自己想自殺死了就算了,於是向他的家人,也都向他的朋友鄭丹瑞,後來也有說,他和他的朋友告別了,再見了..
人生裡面最具毀滅性的情緒,就是憂慮,最難克服的就是害怕..
昨晚在電視上我們都看到,俄羅斯和烏克蘭快要打仗了,當看到人民都拿好行裝要離開,又看到坦克車在冰天雪地裡面行走,軍隊又在整頓槍的時候,看見心都寒..
但在這個擔心驚惶的當中,我們可以做些什麼呢?.
我們在這個幽暗的心靈裡面,我們有什麼出路呢?.
詩人在第一節裡面,他已經用信心來宣告他對耶和華的信仰,第一節,耶和華是我的亮光,是我的拯救,我還怕誰呢?耶和華是我性命的保障,我還懼誰呢?.
在這一節的經文裡面,詩人一口氣用了三個對上帝的稱呼,第一個就是講到耶和華是他的亮光,這個亮光能夠驅走這個黑暗,能夠照亮我們前面的路程..
所以無論我們在那個日子是怎麼困苦,怎麼危險,因著有上帝的光來照預示的時候,我們就有盼望,就有出路..
我們可以想像一下,當我們被困在一個黑漆漆的房間裡面,完全沒有光,我們會害怕的,十分鐘,二十分鐘,三十分鐘的時候,我們會害怕的,因為我們不知道房間裡有什麼東西,心靈裡面慢慢就會不安的了..
但只要有人打開房門一點點,有一些光能夠透進來的時候,人就會開心了,人就覺得有盼望,有出路的了..
所以耶和華是我們生命的亮光,就是給我們在這個幽暗,困苦的生活裡面,給我們有出路..
第二個就是講到耶和華是我們的拯救,耶和華是我們的拯救..
這個拯救對以色列人有一個特別的獨特的理解,因為祂的祖先曾經在埃及裡面做過奴隸的,但是上帝拯救他們,就呼籲了摩西帶領他們出紅海..
一出了紅海,他們就唱詩歌頌耶和華,說我要向耶和華歌唱,因祂大大戰勝,將馬和騎馬的都投在海中..
耶和華是我的力量,是我的詩歌,也都成了我的拯救..
在文訴記裡面就記載了,在以色列人出埃及的時候,有多少人呢?.
他說不算這個小孩子,不算老人家,不算婦女,不算利未俗的人,都已經有六十多萬了..
如果統統計算的話,估計有上百多二百萬的人..
哇,這麼多人,耶和華都可以照顧得到,都可以救得到,何況大衛一個人呢?.
怎麼會救不到呢?所以大衛很有信心地說,耶和華是他的拯救,他還怕什麼呢?.
第三個,他對上帝的稱呼,就是耶和華是大衛生命的保障..
這個保障這個字呢,又可以解為力量和避難所..
這個避難所是一個最安全的地方..
小時候,我經常覺得我的避難所最安全的地方就是國際..
因為我們一家八口住在一個房間,除了屋架床之外,還有一個最頂的角落..
所以有時候激怒媽媽,知道媽媽快要用藤條來打我的時候,我就會爬上最高角落,.
還要去到角落那裡,媽媽一定不會追上來的..
雖然遇到這麼多的攻擊,但是耶和華成為了他的力量,成為了他的避難所..
所以他對一切的東西,一切想傷害他的東西,他都不懼怕..
大衛一生經歷無數的患難,經歷無數的仇敵的攻擊,.
甚至自己至親的兒子的遺棄,背叛..
每天有無數的危險,患難,.
但是因為他有上帝三重的盾牌,.
亮光,拯救,保障,你去做這個支持..
斯人就沒有避到他人生裡面所有這些惡劣的環境來嚇倒..
相反,他可以很有力量,很有膽量,這樣來去面對..
因為他堅定的相信,他所信靠的上帝會給他隨時的幫助..
剛才我也說過,我很怕狗的,.
所以我去探訪的時候,我就很害怕,.
我一定會說,麻煩你了,麻煩你了,拉緊這隻狗..

$^{81}$當我安全進入了要探訪的人家裡的時候,.
我就會和那隻大狗玩一下,.
我會說,你好,我知道你很歡迎我,.
不過我很害怕的..
為什麼我可以和他玩呢?.
就是因為我知道這個門或者這個鐵閘隔開了我和那隻狗,.
於是我就可以二文二道地和他玩..
這個鐵閘就成為了我的拯救,成為了我的保護..
但大衛他經歷了就是夜禍化,成為了他的拯救,成為了他的保護..
所以就因為大衛他知道這個三重的盾牌,.
亮光,拯救,保障,.
所以大衛可以很得意的唱出這個第六節的歌詞,.
他說:「現在我得以昂首,高過四面的仇敵,.
我要在他的帳幕裡歡言獻祭,.
我要唱詩歌頌夜禍化.」.
我很喜歡這句的,.
他說:「我現在得以昂首,高過四面的仇敵.」.
為什麼我剛才說我進了門之後可以二文二道地和那隻狗玩呢?.
就是因為我知道這個門是一個很堅固的保障..
大衛同樣一樣,他可以很得意的..
高昂這個人,我經常想起在六十年代的時候的樣板戲,.
那些宣傳的人在戲院門外,.
你看到每個人都很醒目的,.
高一點點頭,很醒目的,.
好像告訴別人我會努力,我會奮鬥的..
在人生裡面遇到這麼多的攻擊,.
不如意,失敗,失業,失戀,欠債,疾病,.
孤單,沮喪,甚至憤怒,.
或者這次的疫情都給我們一種感覺,無可奈何,.
因為避無可避..
但我盼望天父所給我們的三重銅牌,.
甚至我說這是天父給我們的三寶,.
亮光,拯救,保重,只要我們拿著這三個盾牌的時候,.
我們就可以很醒目的面對人生任何的遭遇..
人為什麼會懼怕呢?.
就是因為找不到支持,找不到力量,.
找不到安全的地方..
人怎麼可以能夠脫離懼怕呢?.
就是找到最大,最穩妥的靠山,.
知道上帝一定在患難當中來去答救我們..

$^{121}$我們又說一下剛才說的張堅庭導演,.
他真的找了很多很多的專科醫生,.
看完都不知道發生什麼事,.
他住在一個極度極度的抑鬱底下,.
有一天晚上他起來去洗手間,.
他突然間記起他小時候去過教會,.
那時候那些導師教過他們,.
有什麼困難的時候可以向上帝祈禱..
於是他就在洗手間裡面跪下,.
他跟上帝說,.
上帝啊,你救我吧,因為他沒有辦法了,.
他自己就在想,.
如果我患了這個柏金遜,.
我就不是一個成功的爸爸了,.
因為我照顧不了我的小朋友,.
照顧不了我太太了,.
那我的家人,小朋友,我太太怎麼辦呢?.
但是心靈有一種聲音告訴他,.
他不用怕,地上高的爸爸不行,.
有餅嘛,那就可以找天上高的爸爸了,.
他想到的時候突然間,啪一聲,.
好像發現了一些東西,.
接著那個星期天,.
他就開始帶他的小朋友上主一學,.
帶他的太太一起去教會崇拜,.
參加教會裡面各種的講座,.
他繼續去找醫生,.
但是後來他發現,.
他患的不是柏金遜症,.
他說他的手指只是骨移了位置,.
不知道怎麼扭傷了,.
以致他的手指不自動地來去動,.
但也因為他過分擔憂身體健康,.
以致他患了中度抑鬱症,.
他覺得這是人生裡面一場大喜劇,.
虛驚一場,.
他也覺得上帝藉著他人生裡面,.
給了他一場這麼大的笑劇,.
但他可以能夠重拾在上帝裡面的新生命,.
在2003年的4月12日,.

$^{161}$他就受浸,加入上帝自己的教會,.
我們從張展廷導演的經歷裡面看到,.
這個擔憂,疑慮,壓力,.
會拖垮了人的,.
一個小小的病而已,.
但竟然想到這麼大,.
然後令自己落入一個中度抑鬱症裡面,.
我們看到這個憂鬱,.
這個擔憂,.
是一個惡者,.
來剝奪人的信心,.
剝奪人對上帝的信心..
最後一方面,.
我們看看,.
在這個充滿憂患,戰爭和不安的環境裡面,.
大衛最大的心願是什麼呢?.
他最大的心願就是說,.
一生一世住在耶和華的殿中,.
瞻仰他的榮美,.
在他的殿宇裡面去求物..
他當時因為大衛時期要去逃難,戰爭,.
所以他很久都不能夠在這個會幕裡面去敬拜,.
所以他最大的心願,.
就是能夠在耶和華的殿裡面,.
當時在這個會幕裡面去敬拜,.
就是說在耶和華的殿中,.
是能夠和上帝有一個親密的關係,.
我們知道上帝愛我們,.
上帝也有足夠的能力來幫助我們,.
在這個疫情的底下,.
我們都很體會到,.
原來能夠回教會團契,.
能夠回教會崇拜,.
是多麼幸福的事情..
但現實裡面,.
我們真的回不了教會,.
我們真的就算一起聚集的時候,.
都只能有兩個人,.
我們有時看到這樣的環境,.
我們真的覺得很苦悶,.

$^{201}$但我們盼望我們有私人大衛那種心態,.
有上帝成為了我們生命的亮光,.
有上帝成為了我們生命的保障,.
有上帝成為了我們拯救的時候,.
我們不再過分的擔心和恐懼現在的環境,.
不再過分的擔憂我們前面的路程,.
我們心靈可以能夠有滿足,有平安..
就正如在最後一節裡面,.
私人大衛這樣說,.
是要等候耶和華,.
當壯膽堅固你的心,.
要等候耶和華..
所以,弟兄姊妹,.
雖然這個疫情,或者周圍不如意的事情,.
好像是沒時沒了,.
不知道什麼時候才有結果,.
但是要等候耶和華,.
在等候期間,我們要壯膽,.
要有一個堅強的心,.
等候耶和華,.
祂一定成就,.
祂一定幫助我們的,.
我們同心討高..
天父上帝,我們謝謝你,.
因為你的話語就成為了我們的亮光,.
成為了我們的拯救,成為了我們的保障..
我們謝謝你,.
祝福我們,.
祝福我們紅恩堂的每一位弟兄姊妹,.
在這個疫情裡面,.
雖然會遇到不少的艱難,.
有很多的擔心,.
在經濟,在生活裡面,.
我們都不知道會怎樣的,.
但是,讓我們專心的仰望天父上帝的時候,.
我們就可以從上帝的話語裡面,.
我們能夠有力量,有盼望..
我們在這裡都求你,.
求你用你自己大能的聖手,.
去將這些病毒快快地趕走,.

$^{241}$讓他們親自到來的,幫助每一位病患的朋友,.
他們在等候醫治的,.
他們灰心的,他們孤單的,.
求主同在,求主醫治..
為了所有的醫護人員,.
他們很吃力的,.
這樣一直在工作的,.
求主加力,.
求主保護他們..
我們又為了俄羅斯和烏克蘭,.
他們現在的那個的境況,.
我們求主賜下平安,.
求主賜下和平,.
也都叫雙方的領袖,.
能夠有這個的智慧,.
能夠處理各樣的事情,.
求主工作,我們做公敬,.
將現在整個世界的情況,.
交給主耶穌基督來,.
你是世界的主宰,.
求你憐憫,求你工作,.
多謝你垂聽禱告,.
奉靠主耶穌基督,.
得勝名字而求..
阿們..
\newpage



\section{歌羅西書 3:1-17}
\label{sec:jiG98Pa_BV4}
\textbf{2022.03.06 《尊貴無比的基督徒身份》 歌羅西書3\_1-17 (和修版) 老耀雄牧師}
\newline
\newline
連結: \href{https://youtube.com/watch?v=jiG98Pa_BV4}{\texttt{https://youtube.com/watch?v=jiG98Pa\_BV4}} ~~~~ 語音日期: 2022-03-08
\newline
\newline
\hyperref[sec:PKTsSbC_DIE]{\small{< < < PREV SERMON < < <}}
~
\hyperref[sec:index_chronic]{\small{[返順時目]}}
~
\hyperref[sec:index_scriptual]{\small{[返順卷目]}}
~
\hyperref[sec:RMhyNLglQvM]{\small{> > > NEXT SERMON > > >}}
\newline
\newline
歌羅西書 3:1-17
\newline
\begin{longtable}{cl}
\hline
\hline
章節 & 經文 (和合本修訂版)\\
\hline
3:1 & \begin{tabularx}{0.7\textwidth}{X} 所以，既然你們已經與基督一同復活，就當求上面的事；那裡有基督，坐在神的右邊。 \end{tabularx} \\ \\ \relax
3:2 & \begin{tabularx}{0.7\textwidth}{X} 你們要思考上面的事，不要思考地上的事。 \end{tabularx} \\ \\ \relax
3:3 & \begin{tabularx}{0.7\textwidth}{X} 因為你們已經死了，你們的生命與基督一同藏在神裡面。 \end{tabularx} \\ \\ \relax
3:4 & \begin{tabularx}{0.7\textwidth}{X} 基督是你們的生命，他顯現的時候，你們也要與他一同在榮耀裡顯現。 \end{tabularx} \\ \\ \relax
3:5 & \begin{tabularx}{0.7\textwidth}{X} 所以，要治死你們在地上的肢體；就如淫亂、污穢、邪情、惡慾和貪婪—貪婪就是拜偶像。 \end{tabularx} \\ \\ \relax
3:6 & \begin{tabularx}{0.7\textwidth}{X} 因這些事，神的憤怒必臨到那些悖逆的人。 \end{tabularx} \\ \\ \relax
3:7 & \begin{tabularx}{0.7\textwidth}{X} 當你們在這些事中活著的時候，你們的行為也曾是這樣的。 \end{tabularx} \\ \\ \relax
3:8 & \begin{tabularx}{0.7\textwidth}{X} 但現在你們要棄絕這一切的事，就是惱恨、憤怒、惡毒、毀謗和口中污穢的言語。 \end{tabularx} \\ \\ \relax
3:9 & \begin{tabularx}{0.7\textwidth}{X} 不要彼此說謊，因為你們已經脫去舊人和舊人的行為， \end{tabularx} \\ \\ \relax
3:10 & \begin{tabularx}{0.7\textwidth}{X} 穿上了新人，這新人照著造他的主的形像在知識上不斷地更新。 \end{tabularx} \\ \\ \relax
3:11 & \begin{tabularx}{0.7\textwidth}{X} 在這事上並不分希臘人和猶太人，受割禮的和未受割禮的，未開化的人、西古提人、為奴的、自主的；惟獨基督是一切，又在一切之內。 \end{tabularx} \\ \\ \relax
3:12 & \begin{tabularx}{0.7\textwidth}{X} 所以，你們既是神的選民，聖潔、蒙愛的人，要穿上憐憫、恩慈、謙虛、溫柔和忍耐。 \end{tabularx} \\ \\ \relax
3:13 & \begin{tabularx}{0.7\textwidth}{X} 倘若這人與那人有嫌隙，總要彼此容忍，彼此饒恕；主怎樣饒恕了你們，你們也要怎樣饒恕人。 \end{tabularx} \\ \\ \relax
3:14 & \begin{tabularx}{0.7\textwidth}{X} 除此以外，還要穿上愛心，因為愛是貫通全德的。 \end{tabularx} \\ \\ \relax
3:15 & \begin{tabularx}{0.7\textwidth}{X} 你們要讓基督所賜的和平在你們心裡作主，也為此蒙召，歸為一體。你們還要存感謝的心。 \end{tabularx} \\ \\ \relax
3:16 & \begin{tabularx}{0.7\textwidth}{X} 當用各樣的智慧，把基督的道豐豐富富的存在心裡，用詩篇、讚美詩、靈歌，彼此教導，互相勸戒，以感恩的心歌頌神。 \end{tabularx} \\ \\ \relax
3:17 & \begin{tabularx}{0.7\textwidth}{X} 你們無論做甚麼，或說話或行事，都要奉主耶穌的名，藉著他感謝父神。 \end{tabularx} \\ \\
[1ex]
\hline
\hline
\end{longtable}
$^{1}$主席.
各位同學,各位姐妹平安!.
《歌羅西書講道》的系列已經到了第三講,也是最後一講了..
在第一章,我們講的是尊貴無比的基督,因為基督就是神的真相,祂是守生的,也是教會的頭..
祂在萬有之先,萬有也靠祂而存在,一切的豐盛都在基督之內..
因此,我們所相信的基督是尊貴無比的..
在第二講,我們講的是尊貴無比的基督信仰,這個信仰不是地上的知識以及價值觀,因為這些價值觀都會因著環境的改變而有變化的..
也不是那種不可以拿,不可以嘗,不可以摸的死硬的律法,只是按照條例的字面解釋而沒有了解背後的核心精神,更加不是那種禁慾或苦行的主義,透過一些極端的行為,以為可以就這樣滿足神的需要..
基督的信仰是使人得自由的,以及得釋放的福音,是與神復和,恢復與神有親密的關係,以及生命連結的屬靈狀態..
今天來到第三講,也是第三章,我們講的是尊貴無比的基督徒身份..
當我們信了耶穌,再得到赦免,身份就會改變..
這個改變了的身份,對我們每一個信徒來說,有什麼寶貴呢?.
這個新的身份對我們行事為人有什麼影響呢?.
因此,我們就以第三章的1至17節,去看看一個新的身份背後的屬靈價值,以及新的身份所應該做和不應該做的..
我們一起有個祈禱,開始!.
你讓每一個相信你的,都有一個尊貴的身份,也都以這個尊貴的身份去踐行我們的信仰,以及實踐你所交託給我們的使命..
願每一個屬你的,都能夠被你使用..
我們同心合意的禱告,是奉靠主耶穌基督,得勝明智祈求..
阿們..
一個人相信了主耶穌之後,他的本質和身份就因著神的超越性而發生了變化..
在第17節的經文裡,保羅提出了四個本質或身份的改變..
在第一節,他這樣說:既然你們已經與基督一同復活,.
經文表示,這件復活的事情已經發生了,我們就有些疑惑了,.
我們現在仍然在世上生存著,我們未經歷過死亡,何來已經復活呢?.
因為經文所指的復活並不是指我們的身體肉身的復活,而是指一個以死的心靈被神所喚醒的復活,.
一個被喚醒的生命,也是指一個從死裡復活的新生命,.
這個新生命已經被神所喚醒,也是重生..
神學家馬丁路德說過,一個靈魂未甦醒的人,他們所做的都是一個有軀殼卻沒有靈魂的人生,.
因為他們所作的任何選擇都是傾向於惡,因為未甦醒的靈魂被污染到近乎死亡,.
失卻了神在創造時的真善美,也就是失卻了神的形象..
因此,凡相信耶穌基督的人,藉著靈裡面的重生,生命就得以改變了..
這個改變的可貴之處,就是可以過一個嶄新而豐盛的人生,而不是過一個被罪所污染的人生..
另外在第三節又提到,生命與基督一同藏在神裡面,.
一個相信主耶穌的新生命,是與基督的生命相連的,.
而這種相連的生命在神裡面,也就是說,我們的生命是與神有份的,.
可以分受到神的聖潔,永恆,屬性和超越性..
一個污穢的人生,內心的污穢在神的光照之下,.
可以逐步逐步得以潔淨,成為一個聖潔的人生..
也唯有聖潔的信徒,才能彰顯到神的真善美..
一個有限的生命,會因著身體不斷的扭壞,最後我們會死亡..

$^{41}$由於神的永恆性賜予了子民一個不扭壞的身體,.
去承載著復活之後的生命,讓每一個屬主耶穌的,都可以有一個永恆的生命..
我們想一想,一個軟弱無力的人,可以在主裡面成為一個有用的器皿..
我們的本質,原本都是自私,驕傲,貪心的,.
但是神注入了愛在我們的生命裡面,.
讓這個愛成為我們與人相處的一個潤滑劑,.
不單單改變了我們的品格,也讓我們可以跟別人復活..
這個改變的寶貴,超越了我們自身的限制..
我們以往的生命有很多不可能的,.
現在因著能夠與這個無限的神有連結,一切都可以變為可能了..
經文第十節也說到,.
「這新人照著做他的主的形象,在知識上不斷地更生..
這個新做的人能夠不受罪的咒詛綁和搏戳,.
而是照著神原初創造人的形象與樣式而活,.
也都是活出一個真善美的人生.」.
一個真善美的人生所獲取的知識,不是屬世的知識,.
而是渴求認識神的一切屬性的知識,.
藉著不斷地認識神的屬性和喜好,.
以及生命不斷地更生和成長,.
讓信徒可以像主耶穌的品格一樣,.
在世上活出真正的信仰精神,.
將天國的價值帶到人間..
這個改變的寶貴,.
是在於我們可以忠於自己的本性而活,.
以真我去面對世界的人和事,.
不需要戴著不同的面具去面對不同的人.」.
最後,經文第十二節說到,.
「你們既是神的選民,.
也就是神的子民,.
是屬於神的,.
信徒是被神所揀選進入神的角度裡面的..
因此,作為神的選民,.
活在神的角度當中,.
就不是為自己而活,.
而是為神而活,.
也是為神的角度而活.」.
這個身份改變的寶貴,.
是在於我們永遠都站在一個得勝者的一方,.
其實人生在世,.
我們不知道自己的選擇何時是對,何時是錯的,.

$^{81}$但當神揀選了我們進入他的角度的時候,.
我們就不需要再怕去選擇了,.
我們只需要緊緊地跟隨和信服,.
就已經可以了..
無論是一個靈魂被喚醒的生命,.
還是一個與神結連的生命,.
還是一個可以忠於自己的生命,.
都是同屬於神的子民,.
都是同樣活在神的角度裡面..
弟兄姊妹,.
你對於現在所擁有的新的身份,.
以及改變了的屬靈本質,.
你有多珍惜呢?.
你有沒有想過,.
你要用什麼方法,.
讓你已經擁有的身份和本質,.
不會失去的呢?.
還是你覺得,.
既然已經擁有了,.
那就永遠都不會失去的呢?.
一個屬於主的信徒,.
我們自然要知道,.
神的厭惡是什麼..
神的厭惡可以分為兩個階段,.
一個是未信主之前,.
另一個就是信了主之後..
在第五節,.
經文這樣說,.
要致死你們在地上的肢體,.
就如淫亂,污穢,邪情,惡慾和貪婪..
貪婪就是拜偶像..
經文所說的淫亂,污穢,邪情,惡慾,貪婪,.
都是神所厭惡的..
這都是一個象徵和代表,.
是指我們的身體受到自身的世俗,.
私慾,操控而產生的一種惡念和罪行..
這些惡念和罪行,.
成為我們要接受終末審判的罪證..
當世界被罪的網絡所包圍,.
人類就只能夠在這個罪裡面生活了..

$^{121}$當人按著這些意念,.
在不同的階段,.
結出了不同的惡果,.
我們最終必定被神追究我們所犯的罪,.
沒有人能夠豁免..
因此,活在罪裡面的人的終局,.
最後就是滅亡..
保羅說過,.
罪的功架乃是死..
這種死並不是指我們人類肉身的死,.
而是永恆的滅亡..
這種永恆的死亡,.
也是永遠與神隔絕..
因此,縱然我們今生所犯的罪,.
可能在我們肉身死亡之前被隱藏了,.
沒有人發現,.
但卻在主再來的時候,.
祂會站在審判台前,.
向每一個人追討..
沒有一個人能夠僥倖地逃避,.
也沒有一個人能夠豁免..
人要為自己所犯的罪,.
付上這個永死的代價..
今天,很多人都忽視了一個真理,.
就是我們無視神的審判,.
以一種自以為是的態度,.
去面對自己所犯的罪行,.
認為只需要沒有人知道,.
不被揭發,.
就可以安然無事了..
不過,神的審判必定會臨到,.
只不過是時間的問題..
可以這麼說,.
人如果沒有了神的恩典臨到,.
其實活著都是枉然,.
因為縱然是活著,.
但他們卻不知道,.
自己的將來結局就是永死..
也因為不知道將來的結局,.
以至今生追求的是一種虛空的人生,.

$^{161}$一種只追求物慾的滿足,.
而無視將來審判的人生..
因此,人需要救恩,.
改變被審判的結局..
人雖然信了主,.
一切的罪行,.
雖然都得以赦免,.
可以重新做人,.
有一個新生命的樣式,.
不過,很多時候,.
老我依然在拖我們的後腿,.
老我如果沒有成長,.
不願意更生,.
仍然會促使信徒繼續犯罪..
在第八節,九節,.
經文這樣說,.
「你們要棄絕這一切的事,.
就是惱恨,憤怒,惡毒,.
誹謗和口中污穢的言語,.
不要彼此說謊,.
因為你們已經脫去舊人和舊人的行為..
已經脫去的舊人,.
理應是聖潔的..
但仍然可以有.
惱恨,憤怒,惡毒,誹謗,污穢的言語,.
以及說謊,.
這些惡念和罪行..
或者,我們可以這樣說,.
這些罪行可能沒有淫亂,污穢,.
邪情,惡慾,貪婪那麼嚴重,.
不過都是罪,.
罪是沒有分大小的,.
他們同樣都是惡種子所結出來的惡果實,.
同樣都需要被信徒去棄絕的..
或者我們不會去犯這些那麼嚴重的罪,.
但我們問一問,.
為什麼我們仍然刻意去犯罪,得罪神呢?.
有兩個可能性,.
第一,我們仍然是享受著最終之樂,.
我們不捨得放棄這種樂趣..

$^{201}$第二,我們對舊恩的內容仍然未全面的認識,.
或者曲解了舊恩,.
認為我們的信仰是不需要付代價,.
只需要相信就足夠了..
那怎麼可以棄絕這種種的惡念呢?.
就是讓聖靈不斷地提醒,引導,督責,.
然後認罪,悔改,求赦免,.
每日心意更生而變化,.
效法基督,.
就能夠將我們的老我,.
這件衣服脫去,.
換上聖潔的新衣..
我相信大部分的信徒不是不明白,.
而是不願意去實踐,.
那我們如何去實踐呢?.
作為一個被主拯救,進入神國的信徒,.
如何去實踐信仰,去取悅神呢?.
就是時刻渴求能夠達到神的期望..
經文在第十二節這樣說:.
「你們既是神的選民,聖潔蒙愛的人,.
要穿上憐憫,恩慈,謙虛,溫柔和忍耐..
作為神的選民,.
首要的就是要有聖潔的生命,.
其他都是聖潔而結出的果,.
神是聖潔的,.
因此與神生命相連的生命,.
都需要是聖潔的..
聖潔亦都是指一種分別為聖,.
討神喜悅的生命,.
因此聖潔的生命,.
就是追求一種與主同行,與主同在的生活.」.
而十三至十七節,.
就是一些討神喜悅的特質..
在第十三節,經文這樣說:.
「倘若者人與那人有嫌隙,.
總要彼此容忍,彼此饒恕..
主怎樣饒恕了你們,.
你們也要怎樣饒恕人.」.
既然信徒的生命是與主結連的,.
就需要展現到主耶穌基督的特性,.

$^{241}$主耶穌要求我們以祂為榜樣,.
主怎樣歡恕了你們,.
你們也要怎樣歡恕人..
其實人很容易被憤怒,.
報復的情緒所支配,.
有人到死那刻仍然放不下仇恨,.
仍然抱怨別人怎樣對不起他,.
仍然被這些情緒所牽動,.
無法走出這個死胡同..
歡恕是讓人走出這個死胡同的最好方法,.
歡恕不單讓彼此得到內心的釋放,.
也能夠醫治受創的傷口..
完全的歡恕更能夠將個人的創傷.
化為一種生命見證,.
可以讓別人得福..
另外,第十四節也說過,.
除此以外還要穿上愛心,.
因為愛是貫通全德的..
基督教最喜歡說愛的宗教,.
因為愛是基督教信仰的核心..
基督的愛促使他放下尊貴,.
威嚴的身份,.
道成肉身,.
來到罪人的中間..
基督的愛也促使他背負十字架,.
受人誣衊,嘲笑,唾罵,凌辱,鞭打..
基督的愛促使他以自己為屬罪制,.
承擔人的罪,以至於死..
這種愛是甘願降悲,.
甘願沉默不語,.
甘願犧牲,.
甘願為人而活的生命..
信徒就是需要這種愛,.
貫穿著整個信仰的生命,.
來見證主耶穌基督的真實..
我們會問,.
一個原本是自私,貪心,驕傲的生命,.
何來愛的動力呢?.
神就是愛,.
神的本質就是愛..

$^{281}$只要信徒能夠真正地與神有生命的結連,.
信徒才能夠有愛的動力,.
否則信仰也只是一個口號,.
毫無用處,.
或者只是一種心靈的慰藉,.
只能夠滿足自己的需要而已..
第十五節又提到,.
你們要讓基督所賜的和平在你們心裡作主,.
又為此蒙照,歸為一體,.
今日,我們生活在一個動盪的世代,.
世界或者社會有和平嗎?.
促使世界和平的可能是威嚇,.
可能是武器,.
或者外交,.
還是一些虛無縹緲的彼此信任..
以上種種都是世界所使用的方法,.
但我們知道並不有效..
神所賜的和平就是一個「復和的福音」,.
讓宣揚「復和的福音」成為人生的使命..
信徒蒙照的其中一個使命,.
就是要同心合意,.
將「復和的福音」傳給人..
神的角度裡,.
每一個人都是被「復和的福音」所承托住,.
每一個人都可以在主裡得到平安,穩妥,.
這種平安不單單在現世可以獲得,.
也延續到永恆,.
無分各籍,.
無分各界,.
都是同屬於主,.
都是同一個角度..
在第十六節,.
經文說:.
「當用各樣的智慧,.
把基督的道,.
豐豐富富地藏在心裡,.
用詩篇,贊美詩,靈歌,.
彼此教導,互相勸勉,.
以感恩的心歌頌神.」.
信仰的生活態度就是將神的話語傳在心裡,.

$^{321}$讓神的話語栽種在我們心田裡,.
成為結出一百倍的好土,.
也成為我們信仰的根基..
我們要無時無刻,.
以各種方法去歌頌,讚美神..
敬拜不單單只是在星期日的崇拜裡進行,.
敬拜是一種生活態度,.
是日常生活的表現..
信仰更不是將我們的生活切割,.
在教會之內我們就是信徒,.
在教會之外就不是信徒,.
或者星期日是屬於主的,.
其餘的時間是屬於自己的..
信仰就是將敬拜神的生活態度展現出來,.
人看到你怎樣去敬拜神,.
就看到你所信的神有多真實..
一個沒有敬拜生命的信徒,.
怎樣可以將所信的神去傳揚出去呢?.
我們想一想,.
一個沒有喜樂生命的信仰,.
怎樣令人相信你所相信的神是可以賜予喜樂的呢?.
一個沒有平安的信仰生命,.
怎樣令人相信你所信的神是可以賜予平安的呢?.
一個沒有見證的信仰,.
怎樣能夠令人相信你所信的神是真實的呢?.
因此一個土神喜悅的生命,.
就是一個能夠時刻彰顯到敬拜神,.
以及活出信仰的生命..
去到最後那一節,第十七節,.
「你們無論做甚麼,或說話,或行事,.
既然我們信徒整個生命都是屬於主的,.
因此我們無論說甚麼,做甚麼,.
是對信仰的真實回應..
並且以生命去見證,.
才是取悅神的最好和最直接的方法.」.
弟兄姊妹,.
今日的講題是「尊貴無比的基督徒身份」,.
我們擁有了神子民的身份,.
我們與神有份,.
我們分受了神的聖潔,永恆,.

$^{361}$屬性和超越性,.
因著這個尊貴無比的身份,.
因此我們要過一個土神起源的人生,.
我們避免去做神所厭惡的行為,.
神所厭惡的就是人活在罪惡裡面,.
讓欲念所衍生的罪行展現出來,.
讓老我不斷的更生,改變,.
寧願不願成長,.
而神所起源的是甚麼呢?.
就是能夠與他的生命結連,.
在生命裡面能夠彰顯愛,.
並且在愛裡面活出信仰..
今日在疫情困擾下,.
你有沒有印象信仰的緣故,.
使你在逆境裡面仍然充滿起立,.
縱使不能回教會,.
屬靈的生命仍然可以不斷成長,.
在不能與人接觸的情況下,.
你仍然願意為別人禱告守望,.
甚至每一個相信主的人,.
我們都能好好珍惜這個基督徒的身份,.
因為這個身份裡面所承載的,.
是我們根本不配能夠獲得的恩典..
最後我想和大家分享一個.
屬靈導師的見證,.
在二十世紀有一個著名的靈修作家,.
他叫盧雲,.
他曾經在耶魯和哈佛神學院裡面,.
擔任神學的教授,.
對於一般人來說,.
成為一個長春神大學的教授,.
是一種成就,.
是一份夢寐以求的即時,.
不過對於盧雲來說,.
他覺得能夠聆聽基督,.
與主同行,.
才是真正的呼召..
他就在他的人生下半場裡面,.
就離開了哈佛的教席,.
他加入了一個專門服務弱智人士的組織,.

$^{401}$叫方舟團體,.
並且在黎明之家裡面去事奉,.
去照顧的是一班弱智的人士..
你想想有多大的差別,.
一個就在神學院裡面教書,.
一個就照顧弱智人士..
當中有一個住在黎明之家的院友,.
他叫阿當,.
阿當是一個嚴重傷殘的人,.
他不能夠用言語去表達自己,.
也不能夠料理自己的起居生活,.
盧雲每一天大概花兩個小時去幫他整理,.
包括餵飯,沖涼等等..
這段時間也造就了盧雲一段的靈修之旅,.
盧雲分享到他體會到天主的愛,.
深深地隱藏在阿當的生命裡面..
他說阿當雖然說不出話,.
但是他帶給人的平靜,.
他身體雖然殘缺,.
但是卻在這個殘缺裡面見到他的溫柔..
阿當外表上好像很無能,.
需要人幫助,.
但是在這個無能的身體裡面,.
卻顯出上帝在他身上的恩賜..
在阿當身上他見到人的和睦,平靜和真愛,.
天國就透過這樣的給予和捨棄展現出來..
盧雲認為阿當是他生命的導師,.
你會不會覺得諷刺呢?.
一個神學院的教授,.
卻認為這個嚴重的弱智的阿當是他生命的導師..
他從阿當的經歷裡面寫了一本書,.
叫做《無能者的大能》..
他當時最大的感受就是,.
他一個以為是去幫助別人的強人,.
卻成為一個被幫助的弱者..
弟兄姊妹,.
今天當我們說他愛人,.
好像盧雲所說的,.
好像一個強者去幫助一個弱者那樣,.
不過盧雲的經歷告訴我們,.

$^{441}$神可能藉著一個看似是弱者的人,.
去幫助我們都未必,.
信仰就是要首先承認自己的軟弱,.
首先承認自己的無能,.
然後才能夠在神裡面去支取這個力量,.
在生活裡面不斷地更生和實踐,.
這樣才是一個神所喜悅的人生..
我們一起有個祈禱,.
你不單給我們生命氣息,.
你更加讓我們在生命裡面與你有份,.
可以享受到你的豐盛,.
祂讓每一個願意跟隨你的,.
都經歷到這種從軟弱到剛強的體驗,.
一生都被你的恩典所盛載,.
過一個土神喜悅的人生,.
為主打美好的仗,.
我們的祈禱仰望,.
是奉靠主耶穌基督得勝明智祈求,.
阿們..
願主祝福大家..
\newpage



\section{約翰福音 14:27}
\label{sec:RMhyNLglQvM}
\textbf{2022.03.13 《願主賜你們平安》 約翰福音 14\_27 (和修版) 曾敏姑娘}
\newline
\newline
連結: \href{https://youtube.com/watch?v=RMhyNLglQvM}{\texttt{https://youtube.com/watch?v=RMhyNLglQvM}} ~~~~ 語音日期: 2022-03-14
\newline
\newline
\hyperref[sec:jiG98Pa_BV4]{\small{< < < PREV SERMON < < <}}
~
\hyperref[sec:index_chronic]{\small{[返順時目]}}
~
\hyperref[sec:index_scriptual]{\small{[返順卷目]}}
~
\hyperref[sec:bW0AAWns91U]{\small{> > > NEXT SERMON > > >}}
\newline
\newline
約翰福音 14:27
\newline
\begin{longtable}{cl}
\hline
\hline
章節 & 經文 (和合本修訂版)\\
\hline
14:27 & \begin{tabularx}{0.7\textwidth}{X} 我留下平安給你們，我把我的平安賜給你們。我所賜給你們的，不像世人所賜的。你們心裡不要憂愁，也不要膽怯。 \end{tabularx} \\ \\
[1ex]
\hline
\hline
\end{longtable}
$^{1}$弟兄姊妹平安.
仍然可以一同去敬拜神.
雖然今天我見不到大家.
但相信上帝讓我們的心緊緊走在一起.
我們再一次地向祂禱告.
在今天仍然賜給我們生命氣息.
去敬拜你親的你.
願意今天你對我們眾人說話.
願意聖靈開啟我們的心門.
讓我們明白.
也給我們力量去回應.
我們這樣的禱告.
是奉耶穌基督的名字.
阿們.
今天只有一節經文.
但非常寶貴的一節經文.
是願主賜你們平安.
這幅圖是不是很美呢?.
是一個很舒服很光明的圖畫.
但今天現時在香港或世界發生.
至少有兩件事很觸目.
不單在香港發生.
在世界各地也發生.
大家看到這幅圖畫.
是我們香港做的強檢.
很多人排隊.
我相信頂展會當中.
有不少人曾經排隊.
很多人排隊幾小時.
有些人在天氣非常冷的日子排隊.
大家的感受一定很深刻.
就算不是.
我們每天看新聞也很厲害.
鋪天蓋地不停去報導.
我們每天都有記者發佈會.
說說每天有多少人確診.
多少人離世等等.
是非常多的.
所以我相信我們每一個人.
是可以離開這個話題.

$^{41}$大家在思考在聽.
亦有弟兄姊妹跟我分享過.
也會影響心情.
也會擔心.
其實這也是很自然的.
我相信在這段時間.
沒有人不思考這件事.
我們的心情或多或少受到影響.
另外一件事.
亦是很厲害的.
我相信在基督徒圈子裡.
不停傳出很多祈禱事項.
我們為烏克蘭這個地方去祈禱.
現在正在打仗.
俄羅斯和烏克蘭之間.
大家亦有很多討論.
縱然我們比較多被香港疫情影響.
但亦沒有忘記在遠方的戰爭.
亦有很多思考.
思考生命的起起跌跌.
生死的問題.
究竟平安在哪裡?.
怎樣可以得著平安呢?.
我相信不同的人都在思考這個問題.
亦都在想方法.
怎樣可以得著平安呢?.
我亦相信神給我們一個很清晰的答案.
在哪裡得著平安?.
今天我們看回這段經文.
約翰福音14章27節.
我再讀一次給大家聽.
我留下平安給你們.
我把我的平安賜給你們.
我所賜給你們的.
不像世人所賜的.
你們心裡不要憂愁.
也不要膽怯.
這一句話是耶穌基督所說的.
耶穌基督說.
祂要留下平安給他的門徒.

$^{81}$經文背景.
在約翰福音.
我們可以看回13章.
經文是14章27節.
我們看回13章背景是什麼呢?.
原來耶穌知道自己離世歸附的時候到了.
祂快要離開門徒.
快要離開這個世界.
祂也知道自己即將踏上受苦之路.
和在十誡上犧牲.
所以祂也告訴門徒.
自己即將被出賣.
在這樣的情況下.
祂提到平安.
平安究竟是什麼呢?.
平安在希伯來文裡是Shalom.
在猶太人見面和道別的時候.
很習慣說這句話.
這句詞語.
在這裡耶穌基督真的和門徒說再見.
準備和他們說再見的時候說.
當然稍後的時間.
在另一個場合.
耶穌基督再說平安.
是很不同的.
那不是說再見.
是和他們說見面.
平安有三個層面.
哪三個層面呢?.
第一個層面是個人的平安.
我們自身的平安.
就算外在環境如何煩亂.
外在環境如何惡劣.
很個人,很鎮定,很平靜.
我的整個情緒就是這樣.
我可以消除那種恐懼.
所以外在環境和我個人自身的情緒不一樣.
這是第一個層面.
個人的平安.
第二個層面是人和人之間的關係.

$^{121}$平安也可以解釋為和平.
這裡是說我們和別人的關係.
是否一個和平的狀態.
其實這絕對有很大影響.
特別是在職場,不同層面.
我們經常說.
工作不會造死人.
是甚麼令我們最辛苦?.
人際關係.
不只是職場.
回到家裡.
在朋友圈當中.
在不同層面.
甚至在弟兄姊妹當中.
甚麼最大影響?.
工作造不死人.
是那種關係.
平安這個層面是說人和人之間的關係.
是否能夠有一種和平呢?.
我和其他人是否有一個好的結連呢?.
第三個平安.
是說我們和上帝之間的關係.
是否復和?.
是真正的復和.
大家如果再回到創世記的時候.
會記得那種很深刻的畫面.
上帝創造了亞當夏娃的時候.
神和人的關係是非常親密的.
很甜蜜的.
每一天都拍拖.
上帝很愛人.
人也是在那裡回應神,愛神.
但是當亞當夏娃犯罪之後.
那個關係就破裂了.
神和人的關係之間就有了隔閡.
一代一代下去的時候.
你會發覺神和人之間有很多的隔閡.
我們神和人之間是否能夠有一個復和的關係呢?.
這是很重要的關鍵.
原來平安牽涉到這個最重要的層面.

$^{161}$(自我自身的情緒).
最特別的是.
在三個平安裡面.
自身的情緒上的鎮定,平靜,安穩.
和我和人的關係.
和我和上帝的關係.
這三個層面裡.
其實最重要的是我們和上帝的關係.
當我們和神復和的時候.
其他兩樣東西都可以做到.
我們和人可以復和.
有一個美好的關係.
我們自己的心靈也有一個很大的平安.
但當我們不能和上帝復和.
(這是一個真正的復和關係).
當我們不能和神復和.
我們和人的關係也不好.
心靈裡很難有一個真正的平安.
所以最重要最重要是和神的關係.
究竟誰可以令我們和上帝復和呢?.
我相信大家心裡一定有答案.
只有耶穌基督.
只有耶穌基督能夠令我們和天父復和.
耶穌基督為了令我們和天父復和.
為我們成就救恩.
耶穌基督說:我就是道路,真理,生命.
若不藉著我.
沒有人能到父那裡去.
耶穌基督就是那條路.
是道路,真理,生命.
要藉著祂.
我們才可以和父神復和.
這些絕對對我們來說是很熟悉的道理.
我們真的聽了很多年.
我們真的這樣信耶穌.
但今天再提的時候.
我們得到平安的最基本最重要的基礎.
沒有了這一點就沒有了真正的平安.
而耶穌基督為了成就救恩.
需要付上沉重的代價.

$^{201}$我相信大家仍然記得.
他付上沉重的代價.
就是即將在約翰福音14:27後將會發生的事情.
耶穌已經說了自己將會被出賣.
是被其中一個門徒出賣.
被出賣後.
他被捉拿.
要經歷不公義的審判.
被戲弄,被鞭打.
最後走上十字架.
被釘死.
這個代價是非常非常大.
裡面有很多苦,很多痛.
不是人所能承受的.
耶穌基督一定要這樣成就救恩.
才可以令我們和天父復和.
我們才可以得到最重要的平安和和平.
很快的時間.
當耶穌說這番話.
說要將他平安留給門徒的時候.
說不久他就要經歷這些事情.
大家可以想像這一刻.
他的心情是怎樣.
但這一刻你看到的耶穌.
卻是心中充滿了平安.
一點都不像那個.
快要踏上受苦之路.
快要被釘死的耶穌基督.
他的平安是非常非常大.
而且他還說要將他的平安留給門徒.
他可以去祝福門徒的生命.
在《約翰福音》14:27.
再看一次.
「我將我的平安賜給你們.
我所賜給你們的不像世人所賜」.
耶穌基督賜的平安和世人是不一樣的.
其實平安這兩個字對世人來說.
是非常吸引的.
世人都很想得到平安.
都會用很多很多的方法去追求得到平安.

$^{241}$我深信在不同的宗教信仰裡.
人都會有很多的祈求.
祈求什麼?.
祈求不同的一些很大能的靈界去保護自己.
祈求誰誰去幫助自己.
一定要有平安.
最好遠離那些患難.
遠離一切的困難.
最好一切都順利平安.
最好就是這樣.
好像什麼圖畫呢?.
一會兒就給大家看看.
我想先給大家看一幅圖畫.
好像這幅圖畫.
祈求什麼?.
就好像這樣.
就算有白雲也好.
都是藍天在上.
下面平靜如水的湖面.
有很漂亮的花.
有美麗的景色.
人生是一片美好.
世人很想渴望.
最好最好不要有風浪.
不要有困難.
不要有任何事.
就是這樣就最好.
如果有困難.
有什麼事最好避開.
逃離.
怎樣都好.
總之不要傷害到自己.
求來求去都是這些.
我想這也是人性.
我們很自然.
我們都很想.
沒有人想求.
最好苦難淋到我生命裡.
讓我受到這個熬煉.
沒有人會想這樣求.

$^{281}$每個人都想求.
不要發生患難.
不要發生困難.
一起平安平靜.
每天的湖水都這麼平靜.
不要起漣漪.
不要起波浪.
一起順利.
由出世那刻直到離世那刻都順利.
這樣就好了.
但是人生呢.
是不是這樣一回事呢.
很明顯人生不是.
人生就好像這幅圖畫.
風浪.
閃電.
很多挑戰.
很多的苦.
很多的難.
所以中國人最喜歡說的那句話.
人生不如意事.
就十常.
是常.
十常八九.
人生不如意事是多的.
避不開.
我們避不開.
在這個時候.
其實我們所需要的一份平安.
是超越一切實際的處境.
是超越艱難困苦失意的.
是沒有人可以奪去的.
無論外在環境如何.
都不受影響的一份平安.
這份平安絕對不是世人所可持的.
所以耶穌基督才這樣說.
我所給你們的.
不是好像世人所持的.
這份平安為什麼和世人不似呢.
因為.

$^{321}$是一個根基和上帝復和而來的.
和上帝有一個關係上的連結而來的平安.
這份平安無論什麼外在的事情都不可以奪走.
無論是什麼處境都不會影響到.
就是在這裡.
可以在生命裡發出來.
這段關係.
是有這樣的力量在.
世人不可以.
世人是用自己的力量.
求不同的.
他們所覺得是神明的.
他覺得是神明的.
去幫助自己.
但我們都知道那些不是真的.
不是能夠幫助到他們的.
但今天耶穌基督.
祂是神的兒子.
就是神自己.
我們三位一體獨一的真神.
唯有祂才能夠幫助到我們.
而這份結連是無可取締.
這份平安.
剛才說過是超越實際的處境.
艱難困苦失意.
沒有人可以奪去.
耶穌基督自己就做了一個很美的示範.
剛才說過.
祂即將踏上受苦之路.
祂即將要踏上十字架.
那份艱難何等大.
那份痛苦何等大.
祂要承載全世界人的罪惡.
這是什麼感受.
不是我們真正可以明白的.
祂要踏上這一切.
要承受這麼多.
但耶穌那一刻是很平安.
祂這份平安就是超越了外在這一切.
就是這麼奇妙.

$^{361}$而且耶穌基督知道.
只有這份平安可以幫助到屬祂的人.
所以祂說祂將祂的平安留給門徒.
好像財產一樣.
原來這是最寶貴的財產.
祂將祂的平安留給門徒.
讓門徒都可以有這份平安.
勝過外在實在的處境.
門徒將來亦都要經歷很多患難.
迫伯,傷害.
在耶穌的門徒裡.
我今天沒有機會讓大家看看.
他們的結局.
有很多是殉道的.
是經歷很大的苦難.
但深信耶穌基督的那份平安與他們同在.
所以他們絕對可以去勝過.
這是一點也不簡單.
在這裡我想起自己兩人的經歷.
也是很想和大家分享.
這麼多年在神裡有大大小小的不同經歷.
去感受到耶穌那份平安的寶貴.
感受到和上帝結連而有的力量.
我記得很多年前.
有一次是暴風雨的日子.
當時我還住在舊屋的地方.
那個地方距離.
我記得當時還在錦宣事奉.
由村巴站走回我家.
很遠的路程.
大約二十分鐘.
那次是暴風雨.
還要閃電打雷.
是非常厲害.
感覺那些雷很近.
閃電很近.
就在旁邊的位置.
一直在打雷.
很厲害.
雨水很厲害.

$^{401}$我的雨傘是沒有用的.
在這麼厲害的情況下.
我已經全身濕透了.
雨水還要在雨傘的骨頭那裡一直流下來.
我已經發覺很多地方水浸.
水浸到我的腳踝.
一直走路.
其實很危險.
那天的走法比較好.
我的坐車是比較好的.
但我選擇一直走路.
一直走路的時候.
整個情況很惡劣和危險.
只要一下雷劈中了.
我就要安息主懷.
那天一直走路的時候.
其實很害怕.
一直祈禱.
耶穌基督幫助我.
保守我.
保守我遠離雷.
其實是最害怕的雷.
不怕雨水.
一直祈禱.
心裡面就越來越有力量.
一開始時不害怕.
但又不是.
也會害怕.
一直濕透了.
一個人在走路.
其實很危險.
完全可以成為被雷電擊中的目標.
但我一直祈禱.
我記得中間有一個位置.
感覺到.
我不知道是電流還是什麼.
有一點點在雨傘頂下.
有一點暖和熱力.
但到我的手指位置就停下了.
我知道是危險.

$^{441}$但在過程中.
上帝一直保守.
就算閃電在很近.
我旁邊的位置.
打雷在很近的位置.
其實上帝都保守我.
沒有被打中.
最重要是心裡面的力量.
出來支持我.
讓我一直回到家.
上帝一直保護.
一直給我平安.
回到家一定給家人.
一頓飯.
這麼危險.
是否不想不顧生命.
但在過程中.
回想起來.
是一個很特別的經歷神的機會.
這是其中一次.
另外有一次更加深刻.
在我去教會事奉前.
其實我在一個基督教機構裡事奉.
在遠東廣播裡做節目.
有一次機構派我去廣西.
接觸少數民族.
跟他們錄音.
我去了.
我租了一間酒店.
其實是想幫機構節省一點.
不要花太多錢.
那次租的酒店比較普通.
衛生環境也不算十分理想.
那次我跟少數民族一起錄音.
要照顧他們生活.
就出街買菜.
我買了一些非常辣.
非常油膩的食物回來.
因為他們很喜歡吃辣.
吃完之後.

$^{481}$我就非常不妥.
真的十分感到腸胃很疼痛.
一直忍到晚上.
真的工作完.
才盡情去洗手間.
過程不得了.
那一晚我去了非常多次洗手間.
整個人虛脫.
然後我發覺自己發高燒.
過程中.
有少數民族知道我病得很厲害.
甚至第二天我做不到事.
真的不OK.
但他們可能基於不同原因.
沒有過來看我.
變成我自己一個人.
由頭到尾都是自己一個人.
燒得很厲害.
很辛苦.
整個身體完全不在狀態.
在我酒店附近也沒有很好的醫院.
如果想買什麼藥.
藥房的藥物也有限.
所以也很辣.
再加上我經常要黏著洗手間.
是比較困難.
結果發高燒燒了一段時間.
其中一件事就是祈禱.
當然我身邊也有些藥物.
我自己也有些藥物.
吃了少許.
但藥物在這麼惡劣的情況下.
應該不足以扭轉我的身體.
但我也吃了少許.
我便祈禱.
很記得有一個祈禱很深刻.
我當時很辛苦.
真的辛苦到有個衝動.
想打電話給香港的家人.
跟他們說我十分不舒服.

$^{521}$可不可以飛過來接我回去.
十分不可以.
但後來想也沒有辦法.
其實他們飛過來也幫不了我什麼.
我便一路祈禱.
自己一個人祈禱.
沒有叫任何人為我祈禱.
就是祈禱.
其中一個祈禱是.
上帝我知道你是大能的.
你是全能的.
其實你要醫治我.
你絕對可以.
但神如果不用這個方式.
我也很信服.
我便將自己交在你的手裡.
由你去帶領.
我做完這個祈禱.
我很記得.
當時第一個最大的反應.
是自己的心靈很得釋放.
很記得那份平安.
本來整個人的新心靈.
是非常不可以的.
但祈完的肚子很釋放.
然後便很奇妙.
很奇妙.
上帝便大大出手.
在這個禱告之後.
上帝大大醫治我.
令我每件事都康復起來.
然後我記得大概兩三天後.
我徹底康復了.
我坐飛機回香港的時候.
完全是很OK.
好像沒有事發生.
真的很感恩.
我很記得那份平安.
不是康復之後有平安.
而是禱告完後.

$^{561}$我的心靈得釋放.
是那份平安.
帶著我經過.
到最近的時候.
我回想身邊有很多討論.
說說疫情如何.
說說多少人確診.
身邊認識的人確診.
或有認識的人安息主懷等等.
我們的心情受影響.
但我經常堅持一件事.
很相信一件事.
只要上帝保守我們的時候.
我們一定可以渡過這個難關.
我經常問同工一個問題.
最喜歡問.
有沒有想過如果我確診會怎樣.
好像經常都會有客人.
如果我確診會怎樣.
大家打算怎樣.
過程不是說我很渴望確診.
而是我將一切交投給神.
確診也好.
不確診也好.
上帝在我們裡面掌管.
上帝與我們同在.
那份平安會給我們力量去勝過.
無論今天身邊的環境變得多惡劣.
現在的情況多惡劣.
上帝與我們同在.
一定帶領我們渡過這個難關.
無論未來政府如何安排.
我們未來的日子會經歷甚麼.
其實神都是與我們同在.
不會離開我們.
今天第一節我想說.
上帝已經把這份平安給了你.
給了我.
因為我們與祂復和了.
有祂與我們一起.

$^{601}$天天都不用害怕.
天天上帝都保護帶領我們.
但我知道我們心裡仍然有很多擔心.
我們受環境影響.
我們有很多掛慮.
可以怎樣做呢?.
在《非拉比書》四章六到七節.
有一段寶貴的經文.
「應當一無掛慮.
將你們所有的高素神.
神所賜那超越人.
所能了解的平安.
必在基督耶穌裡.
保守你們的心懷意念」.
禱告.
凡事禱告.
不單止禱告.
不單止祈求.
還要感恩.
經常數算上帝恩典.
正如我們能夠回來教會服事.
就感恩了.
我覺得是神給我好好的福氣.
每個星期上帝還給我生命氣息.
還有健康.
做這做那.
就要感恩.
弟兄姊妹我們都可以一同感恩.
一路感恩.
我們就將眼光放回上帝那裡.
看到神的恩典還伴隨著我們.
神還沒有離開我們.
還在帶領我們的人生.
還有我們有憂慮的時候.
有害怕的時候.
就禱告.
一路禱告.
是一路的.
凡事都去禱告.
將我們所有的告訴上帝.

$^{641}$上帝給我們的平安.
在這裡說.
是超越人所能了解的.
是超越我們所能了解的.
在基督耶穌裡面保守我們心懷的意念.
在這裡願耶穌基督賜福給大家.
賜給大家有那股力量.
在這時勢仍然是.
滿有力量勝過.
無論你看到甚麼,聽到甚麼.
耶穌基督.
直到永遠都在我們生命裡面掌管.
我們無懼風雨.
無懼困難.
我們一同禱告.
有我們這麼寶貴.
在疫情底下可以去思想.
甚麼是平安.
天父原諒你就是這麼疼我們.
你藉著耶穌基督.
為我們成就救恩.
讓我們藉著耶穌基督.
與上帝復和.
當我們復和之後.
就非常好.
上帝讓我們有無窮無盡的力量.
勝過外界一切困擾.
一切艱難,一切苦.
今天求你讓我們看到.
神是這麼好.
神是這麼寶貴.
讓我們捉著你.
讓我們的眼睛看著你.
無論我們繼續聽到甚麼新聞.
看到甚麼事情.
上帝讓我們看到.
你超越了這一切.
你也會幫助我們超越這一切.
我們會一直很有力量.
走每一天的路.

$^{681}$向你賜福給弟兄姊妹.
讓我們在這樣的時勢裡.
仍然能夠好好休息.
好好吃東西.
但更重要的是.
好好親近你.
每天從你得力.
每天要認識你更深.
更加要明白.
上帝要我們在疫情裡.
如何生活.
能夠更加討你的喜悅.
求你興起基督徒.
在疫情當中.
我們作美好的見證.
要讓身邊的人看到這個見證.
是深深被上帝吸引.
我們要他歸向你.
他們也要一同歸向你.
求主賜福保守.
奉耶穌基督的名字.
祈禱.
阿們.
\newpage



\section{哥林多前書 12:4-11}
\label{sec:bW0AAWns91U}
\textbf{2022.03.20 《聖靈隨己意分給各人的恩賜》 哥林多前書 12\_4-11 (和修版) 何智雄牧師}
\newline
\newline
連結: \href{https://youtube.com/watch?v=bW0AAWns91U}{\texttt{https://youtube.com/watch?v=bW0AAWns91U}} ~~~~ 語音日期: 2022-03-20
\newline
\newline
\hyperref[sec:RMhyNLglQvM]{\small{< < < PREV SERMON < < <}}
~
\hyperref[sec:index_chronic]{\small{[返順時目]}}
~
\hyperref[sec:index_scriptual]{\small{[返順卷目]}}
~
\hyperref[sec:atdLdwQVuag]{\small{> > > NEXT SERMON > > >}}
\newline
\newline
哥林多前書 12:4-11
\newline
\begin{longtable}{cl}
\hline
\hline
章節 & 經文 (和合本修訂版)\\
\hline
12:4 & \begin{tabularx}{0.7\textwidth}{X} 恩賜有許多種，卻是同一位聖靈所賜。 \end{tabularx} \\ \\ \relax
12:5 & \begin{tabularx}{0.7\textwidth}{X} 事奉有許多種，卻是事奉同一位主。 \end{tabularx} \\ \\ \relax
12:6 & \begin{tabularx}{0.7\textwidth}{X} 工作有許多種，卻是同一位神在萬人中運行萬事。 \end{tabularx} \\ \\ \relax
12:7 & \begin{tabularx}{0.7\textwidth}{X} 聖靈彰顯在各人身上，是要使人得益處。 \end{tabularx} \\ \\ \relax
12:8 & \begin{tabularx}{0.7\textwidth}{X} 有人藉著聖靈領受智慧的言語；有人也靠著同一位聖靈領受知識的言語； \end{tabularx} \\ \\ \relax
12:9 & \begin{tabularx}{0.7\textwidth}{X} 又有人由同一位聖靈領受信心；還有人由同一位聖靈領受醫病的恩賜； \end{tabularx} \\ \\ \relax
12:10 & \begin{tabularx}{0.7\textwidth}{X} 又有人能行異能，又有人能作先知，又有人能辨別諸靈，又有人能說方言，又有人能翻方言。 \end{tabularx} \\ \\ \relax
12:11 & \begin{tabularx}{0.7\textwidth}{X} 這一切都是由惟一的、同一位聖靈所運行，隨著自己的旨意分給各人的。 \end{tabularx} \\ \\
[1ex]
\hline
\hline
\end{longtable}
$^{1}$各位紅人黨的弟兄姊妹,你們好.
我是孫克航的學務司儀.
原本今天會跟大家分享.
但知道疫情不可能.
所以要在辦公室跟大家分享聖經.
我知道背景不太理想.
所以才容許我到分享的PowerPoint.
首先要代表孫克航所有弟兄姊妹.
大家民安,希望平安.
平安這幾個字說就容易.
但現在實際上能夠經歷到.
能夠感受到確實很困難.
這個世界是變化萬千.
今天是不知明日事.
怎可以有平安.
我們所相信的主是永不改變的主.
祂的愛是不會離開我們的.
所以弟兄姊妹,這段時間我們更需要捉緊著上帝.
捉緊著祂給我們的說話.
以不變迎萬變.
以捉住我們不變的上帝.
我們去面對這個世界.
求主幫助我們.
讓我們先低頭一起祈禱,好嗎?.
今天的早上,讓我們一齊來到你面前.
我知道你愛我們.
求你閱立我們在你面前的奉獻.
在你面前的敬拜.
也求主你賜個給我們.
是你所應許的.
我們知道我們一切的好處不是你而我.
求你賜下你的恩典給我們.
主啊,現在在這個世代裡,我們真的很需要你的恩典.
以致我們才能夠好好地生活.
所以我們恭敬的將以下時間交上.
求你與我們同在.
同心合意在你面前一齊禱告.
奉耶穌基督的名字祈求,Amen.
我今天可以跟大家分享.
知道你們這半年裡有一系列的講道.

$^{41}$今天講的是關於恩慈.
我特別選了哥林多前書第十二章第四到十一節.
我的題目是「聖靈隨己而分給個人的恩慈」.
如果大家看這篇詩篇.
愛的篇的真諦的時候.
相信大家都認識的.
不會很陌生.
因為當你去我們《大英之母》的婚禮.
你可能會看到,甚至聽到這個愛的詩篇.
愛是永恆,能夠忍耐,也有恩慈.
愛是不嫉妒,不自誇,不張狂.
不作害羞的事,不求自己的益處.
不輕易發怒,不計算人家的惡.
不喜歡不義,只喜歡真理.
凡是包容,凡是相信,凡是盼望,凡是忍耐.
愛是永不自食.
一個很美麗的詩篇.
但是,大英之妹.
你知不知道,詩篇的愛.
並不是講關於夫妻之間的愛.
並不是講關於男女之間的愛.
在某程度來說.
我們有時是用錯了聖經.
為何這樣說呢.
我們今天打開聖經的前書.
一起去看看.
先看看前書裡.
其實保羅想跟我們講的是甚麼.
之後我們慢慢回到第十二章和第十三章.
我們就深入一點去看.
究竟保羅用第十三章的愛的真體.
其實他很想跟當時的大英之妹說甚麼.
今天我們再說,他跟我們說的是甚麼.
原來我們可以在少人主我們看的時候.
保羅記熟了他在生命中有三次的轉生.
去到亞洲不同的地方.
甚至去到歐洲.
在第二次的時候他跟同工.
他先去由安提柯出發.
之後去之前的教會探望他們.

$^{81}$其實保羅很想在亞細亞的範圍裡.
甚至在伊布的地方去傳福音.
去見到教會.
但聖經上說得很特別.
聖靈不容許他們這樣做.
他想去大聖靈在當中禁止他們.
於是他們一路向西退.
去到一個地方叫特洛亞.
保羅就在這個位置.
特洛亞,大家如果看地圖.
上面是馬其頓的地方.
大家記得是馬其頓呼聖.
保羅在這裡聽到他們說.
你可不可以過來幫我們.
當我們去看聖經,特別是西洋的聖經.
看少人傳的時候.
我們真的需要認識地圖.
看的時候一定要打開地圖.
不然會一起看.
因為裡面有很多地名.
有很多不同的方向.
還有,如果能夠在地圖上看到這些訊息.
就會更加明白.
其實在這次小人轉到這裡的時候.
傳福音是一個很大的突破.
福音由亞洲進入到歐洲.
就是這一次.
是一個突破性的旅程.
保羅帶著同工到歐洲才去到菲律賓.
之後到特洛尼加.
之後到比利亞.
為何他們會一直轉呢?.
其實他們的傳福音是很有果效的.
保羅很快就帶了人信耶穌.
不過可惜當地的猶太人很妒忌.
很不喜歡他們留在那裡.
因為他們說的是耶穌基督的復活.
耶穌基督的救恩.
他的好消息.
是一群猶太人不能夠接受的.

$^{121}$於是就一直迫害他們.
由菲律賓,到特洛尼加,然後到比利亞.
保羅就跟同工一路走.
到那裡之後.
保羅就放下了提摩太太.
西拉他們.
之後就自己到雅典.
然後就到哥林多.
在哥林多就住了一年半.
就建了一間教會.
教會裡面是很感恩.
裡面除了猶太人外.
還有很多外邦人.
但是保羅在那裡一年半之後就離開了.
他就去了醫學院.
再停留了一段時間.
之後他就回到耶路撒冷.
開了耶路撒冷大會.
在過程當中.
當然保羅是很想念.
哥林多裡面的弟兄姊妹.
因為他是親手帶出來的.
當他聽到.
哥林多教會裡面有很多問題.
很多困難,很多情況出現的時候.
於是他就寫信.
當他第三次退出的時候.
他就寫了一封信.
給哥林多前書和後書.
鼓勵他們.
幫助他們.
還有裡面有些錯誤的想法.
他就幫他們改變.
讓他們不要偏離真道.
當然說保羅就是.
哥林多這個地方.
其實發生了很多信仰上的問題.
當時是希臘一個很大的城市.
羅馬帝國阿波拉省的首都.
是州府.

$^{161}$裡面的文化是很高的.
但品行很複雜.
是一個世界性的城市.
很多人都來到這裡.
但另外.
裡面的道德是非常敗壞.
道德淪落.
甚至有妙技.
淫亂的事也有很多.
正因為這些緣故.
所以令到哥林多教會裡面.
都有很多問題.
因為都是那些人.
不過他們上了耶穌之後.
他們進了教會.
但裡面有很多的舊我.
其實一直都在裡面.
因為這些緣故.
所以他們的教會裡面.
就衍生了很多困難.
問題.
在聖靈的感動下.
保羅就寫了前書和後書.
特別是前書.
哥林多前書是比較長的.
大家看新約裡面.
不同教會的聖經.
是比較長的.
因為裡面的問題.
比較多.
和比較長一點.
要去多一點的去教導.
一班弟兄姊妹.
所以保羅就寫了這個.
很簡單來說.
我們看看哥林多前書的大綱.
第一章.
因為保羅去看見.
聽見哥林多的教會時.
其實裡面很感恩.

$^{201}$保羅.
因為裡面有很多很好的信徒.
很有恩慈.
他們做得很好.
但裡面有些困難出現.
在第一章第十節到第四章二十一節裡面.
講到教會是分裂.
分黨派.
我是信保羅的.
我是信基督的.
我是信阿波羅的.
人們就是這樣分門別類.
分黨分派.
令到教會有很多困難.
在合一上有很多困難.
在第五章第一節到第六章第二十節.
講到關於淫亂.
社會上的充斥問題.
其實都發生在教會裡面.
而第七章第一節到第十五章第五十八節.
保羅就講出.
教會裡面有很多不同的困難.
很多的問題.
拜偶像.
聖餐.
耶穌復活.
娶嫁.
受人恩賜.
愛.
女人的蒙頭.
先知講道.
婦女的講道等等.
令到整個教會裡面.
因為沒有了頭.
保羅不在.
很多的問題都解決不了.
沒有一個很好的指導.
令到教會裡面很混亂.
保羅就是這樣.
逐個逐個的問題.

$^{241}$都跟我們講應該怎樣做.
第十六章就是結餘.
回到我們今天的章節.
我們特別要講第十二和第十三章.
是講到關於屬靈的恩賜.
第十二章和十三章的關係在哪裡呢?.
仔細地分析一下.
第十二章第一到第十一節.
是講到關於屬靈的恩賜.
第十二至三十一節是講到一個身體.
保羅做一個比喻.
用身體去做一個比喻.
一個身體裡面有不同的肢體.
不同的肢體需要彼此配合.
才能夠有一個完整的身體.
但是要彼此配合的時候.
保羅提及很多教會裡面.
他們更加需要求一個更加大的恩賜.
保羅說更奄妙的道.
這個道是甚麼呢?.
就是愛.
第十三章裡面是講愛.
一到三節是講愛.
而四到七節是今天的經文.
主席讀了的經文就是愛的身體.
第十三章第八節.
愛是永不止息.
先知的講道始終都會沒有.
講方言的能力始終都會停.
知識都歸於無有.
唯有一件事永遠長存的是甚麼呢?.
就是愛.
神就是愛,他就是愛.
而我們都很需要愛.
保羅用愛去指出一個很重要的問題是甚麼呢?.
其實保羅,我剛才都從大綱裡面看.
多倫多教會是一個很有恩賜的教會.
他們的講道很強.
他們講完很有知識.
很有傲備,知道內容.

$^{281}$能夠引人興奮.
大家都知道,如果講道能夠這樣.
突然間每一節都會很得著.
他們有信心.
有人行神蹟,有異能.
有些病的時候,有土地就得到醫治.
有窮人的時候.
當中都可以得到信徒的熱切幫助.
他們肯幫助,甚至犧牲.
如果身處這樣的教會.
其實是很幸福的事.
我希望你和我都希望.
張學堂和洪恩堂都是這樣.
但其中一個令這教會烏煙瘴氣.
一無是處.
沒錯,有很多很多恩賜.
不過,卻少了一件事.
就像一道菜式.
色,香,味都有.
但當少了鹽去調味的時候.
這道菜可能不好吃.
頂尖都不享受.
但如果多了鹽去調味.
可能整道菜都不一樣.
恩賜都是這樣.
只有恩賜,沒有其他.
是沒有意義的.
在恩賜加上愛的時候.
心靈就很有意義.
原來保羅在用愛去形容.
去跟恩賜拉上關係.
剛才我讀的經文.
第13章的四到八節.
愛是恆久忍耐,又有恩賜.
愛是不疾度,不自誇,不張嘴.
不作海修的事,不求自己益處.
不輕而發怒,不計算人的惡.
不喜歡不義,只喜歡真理.
凡事包容,凡事相信,凡事拋妄.
凡事忍耐,愛是永不完成.

$^{321}$這些的鹽,這些恩賜.
加上剛才大家讀第十二章第四到十二節.
就很不一樣.
哥林多前書第十二章第四節.
因此原有分別,聖靈卻是一位.
即是也有分別,與卻是一位.
公用也有分別,神卻是一位.
在眾人裡面運行一切的事.
聖靈顯在個人身上.
是叫人得益處.
這人蒙聖靈賜他智慧的言語.
那人也蒙這位聖靈賜他知識的言語.
又有一人蒙這位聖靈賜他信心.
也有一人蒙這位聖靈賜他醫病的恩賜.
又有一人能行異能.
又叫一人能作先知.
又叫一人能作辨別諸靈.
又叫一人能夠講方言.
又叫一人能夠翻方言.
這些一切都是聖靈所賜.
隨己分級個人.
十二章講完之後.
十三章做一個很重要的補充.
愛.
每個人都有不同的恩賜.
每個人都可以好好發揮他們.
但是各自發揮,沒問題.
但如果大家走在一起.
要彼此配搭.
彼此去造就.
去建立一間教會的時候.
如果缺乏了愛.
就有很多很大的問題.
就好像保羅體建.
哥林多教一樣.
沒錯,有人講得很厲害.
有人醫治,有人其他.
有人信心,等等都有.
但最後見到的是甚麼呢.
分黨分派.

$^{361}$不合一.
合一就成為了.
哥林多教的一個很重要的問題.
每人不能彼此配搭.
保羅在十二章的時候.
剛才講到十一節之後.
他講到一個比喻.
就是身體.
整間教會變成身體一樣.
聖靈給每個人都有不同的恩賜.
就好像肢體當中有不同的.
手,眼,耳,口,鼻,手,腳等等.
每個人都不一樣.
恩賜都不同.
但是都不能夠說.
眼說我不是腳.
我不和他在一起.
我和他比他更好.
因為他不懂得做我做的事.
我可以看見東西.
但腳告訴我.
我看不見,但我可以走路.
大家有各有不同的恩賜.
有不同的長處.
如果各自分滿別類.
各自分黨分派.
我是屬基督的.
我是屬保羅的.
我是屬阿波羅的.
如果在這樣的情況下不合一.
就出現很大問題.
我們都是屬於基督的.
基督需要成為我們的元首.
基督是我們的頭.
是他一直帶領著,掌管著.
以至整個人身體才能夠有很好的發揮.
保羅很好.
保羅用一個很適當,很好的比喻.
令到我相信當哥林多教會看官信時.
他們都會明白,知道.

$^{401}$現在的時代,大家知道疫情很嚴重.
確診人數真的很多.
沒想過香港現在的情況是這樣的.
最近聯會也很好.
他們也準備了快速檢測包.
提供給教會.
以至教會能夠給予一些需要的人.
大家想想整個過程其實是不容易的.
需要很多的.
看似很簡單,一件事.
但裡面需要很多不同的,有恩賜的人.
彼此配合才能做得到.
我再多一點分析一下.
先是一個異象,看見一個圖畫.
心裡想著如何在現在的疫情中見證耶穌.
需要一些方面比較強的恩賜才能看得到.
只看見是不行的.
需要一些人有恩賜能夠定一些目標出來.
讓他們一步一步如何聯絡供應商.
或一些肢體願意奉獻.
如何將快速檢測包送給他們.
如何取貨,如何分派給堂會.
如何做物流的工作.
很多很多的.
每個堂會都要取得檢測包.
取得之後要了解元朗的情況.
大埔不同,香港島有不同的堂會.
有不同的情況,各種事物都不同.
他們要了解裡面有甚麼人,有甚麼需要.
然後看他們如何派發給他們.
現在的情況是要很小心.
不可以有多人聚集.
派發時要戴口罩.
很容易會在過程中感染.
一說就容易.
但由開始想到這個方法.
到最後有需要的人能夠取得檢測包.
整個過程中,如果仔細想想.
其實需要很多人,很多的恩賜.
很多突然之間彼此配搭才能做到.

$^{441}$否則這件事是不能成功的.
我簡單地將行動或這件事解釋出來.
大家一想就很明白.
中間真的需要合一.
如何才能合一呢.
每個突然之間,我們心中都有一個主耶穌.
就像我們一幅圖畫的拼圖一樣.
大家看到,這是香港的拼圖.
我相信要拼圖的時候.
必須先拼一個香港夜景的圖畫.
之後就可以逐張拼在一起.
我們也是一樣.
上帝給我們不同的恩賜.
每個人都有的,不會沒有的.
因為聖經說得很清楚.
聖名是按照自己的己意分給各人.
這個各人是每一個基督徒.
我們都有從上帝而來的恩賜.
當我們按照自己的恩賜去做的時候.
會做得特別好.
有力量的,果效特別大.
因為擅長一定做得好.
所以如果各人都將自己的擅長合在一起.
大家一起做,而且配搭得好.
整整有條理.
最後就好像拼一幅香港島的拼圖一樣.
逐塊逐塊拼,拼得好,拼得清楚.
明白這個位置應該在哪裡.
最後就有一幅很美的香港夜景.
在我們眼前出現.
當我們今天面對疫情的時候.
我們每個人都需要多走一步.
每個人都需要多走一步.
那一步是怎樣走呢.
我們需要問上帝.
究竟,比如說,我今天活在元朗.
我今天在這裡,我可以做些什麼.
現在和以往的情況是否一樣.
剛才我說過,大家都知道.
現在我們面對八年都沒有一遇的疫情.

$^{481}$從來都沒有發生過.
是你和我.
所以從兩個角度看.
為什麼我們這麼不幸.
活在今時今日的情況.
這麼慘,這麼沮喪.
這麼令人不安,這麼困難.
你可以從這個角度去理解現在的情況.
但我們也可以從另一個角度去看.
我們其實是面對一個黃金機會.
為什麼是黃金機會呢.
最近我在街上走的時候.
我聽到很多不同的人,有不同的反應.
突然之間我聽到有人「啊」.
然後他們自己說話.
有些人,我看到在街上.
兩個人,可能中間發生了一些事.
應該不是很大的事.
然後他們說了很多不好聽的話.
甚至粗暴.
比以前多了很多.
為什麼呢.
正因為人心裡很不安.
很暴躁,很有情緒.
不知道該如何發洩.
所以一些小問題就引發很大的迴響.
人在這裡面正面對一份不安全.
無奈的感覺.
誰能幫到他們.
如果我們從傳播音的角度.
現在真的是一個很好的黃金機會.
耶穌說,犯了苦難重擔的人.
可以來到我這裡.
我就可以得到他們的安置.
這個說話如果以前說的時候.
人都沒事,很平安.
上班,看電影,吃飯,都沒問題的時候.
對於這句話,他們的感覺上是不大的.
今時今日.
為何要說「懦夫囚犯」呢.

$^{521}$重擔不是很輕的.
每個人都很重.
當你說這句話的時候.
突然人們會豎起耳朵.
聽多一點.
是不是真的.
最近大家知道.
在烏克蘭和俄羅斯打仗.
我收到那邊的消息.
在Galatoo的消息.
有一間神學院.
最近因為這樣的情況.
打開神學院.
讓居民進入去.
暫時躲避.
他們發現很多人都願意一齊祈禱.
甚至星期日,多了一些人.
當然這是在之前的情況.
現在的情況是在之前.
很多人願意來教會.
為甚麼.
因為心裡很不安.
不知道怎麼辦.
他們很喜歡,很需要這些安慰的說話.
誰能夠給他們.
我聽到一些消息.
就是一些牧者.
神學院.
他們繼續留在那裡.
他們不走,為甚麼.
因為他們很相信.
如果這段黃金時間.
他們能夠在這段時間.
給到安慰訊息給他們.
幫到失望的人.
給他們一些盼望.
一些受傷的人.
可以得到心靈上的醫治.
這個機會.
過了就沒有了.

$^{561}$我記得.
有一個很年輕的神學家.
潘佛華.
他身處在第二次世界大戰的時候.
其實他可以去美國.
他不需要回德國受苦.
他完全可以去.
但是他心裡很想.
他有一個意象.
如果他現在在大戰的時候.
不能夠和他的國民.
一起去受苦.
他將來就沒有資格.
和他們一起.
打完仗之後.
和他們一起去享受和平的時間.
是因為這個原因他回去了.
最後他被希特拉抓了.
在第二次世界大戰和平之前.
剛剛好之前.
他就被殺害.
殉道.
現在是一個怎樣的時間.
沒錯.
我們覺得很困難.
全世界都覺得很困難.
但作為基督徒.
作為我們的時候.
其實我們靠著耶穌.
我們是可以比其他人好一點.
好一點.
而當其他人來的時候.
我們可以將耶穌的福音.
那個好消息.
以前都不覺得是好消息.
但現在就不同了.
我們可以捉緊這個機會.
和別人去說.
耶穌說得很清楚的比喻.
一般人只要我們做一個小事.

$^{601}$在他身上.
縱使是一杯涼水.
我們給他.
其實上帝是看在眼裡面.
大家知道《馬太福音》第24章.
講到問公打問國.
要公家國.
其實這些說話.
我們之前打開聖經的時候.
我們都看了.
以往我們看完就算了.
但今天不同.
我們真的在經歷這件事.
我們教會是需要走在一起.
按照自己的恩賜.
合一成為一個見證基督的群體.
有人看到我們的盼望.
我們不是說我們能夠.
把這個疫情變成武.
把自身馬上變成和平.
我們沒有這個能力.
上帝是有.
但我們卻有能力靠著耶穌.
得到那份平安.
這個是很吸引的.
粉絲們,我也鼓勵大家.
我們一起努力.
我們也知道很多不同的教會.
在不同的地區都在做支援工作.
我們會給予一些.
我們會派給一些有需要的人.
我們會知道有些人有需要.
我們會幫他們.
有些人不能外出.
我們會幫他們去買.
用安全的方法把物資給他們.
雖然我們做的事不是很多.
但是大家都記得五餅二魚.
五個餅一條魚.
最後餵飽了多少人呢?.

$^{641}$是五千人.
誰去餵呢?.
十二個門徒.
當耶穌來到面前.
當耶穌把他們捉住後.
把五個餅的一部分給他們的時候.
如果他們只看著眼中的那小小的餅.
他們敢不敢去派呢?.
那麼少,在眼前.
那麼多,怎麼派呢?.
但他們對耶穌有信心.
既然耶穌叫他們去的時候.
他們就去做.
於是就拿去派.
嘩!竟然不單止這麼多人吃飽.
還有剩下十二個門徒.
靜姐妹.
今天的故事.
這個神蹟仍然可以發生.
仍然可以發生.
問題是你和我.
是否也將我們的生命.
將我們自己再次交給耶穌.
和耶穌說.
主啊,在今時今日.
我應該要做什麼.
你給我的恩賜,我應該怎樣發揮.
我們在教會裡.
大家如何彼此合一.
走在一起.
才能讓在這些時間.
有人真的看見.
耶穌的大體活在他們當中.
合一很重要.
多不多教會.
雖然有很多恩賜.
靜姐妹有很多不同的恩賜.
但他們分黨分派.
不合一.
不合一正正就是.

$^{681}$當時教會的很大問題.
求主幫助我們今天.
當我們看那麼多前書.
當我們看恩賜的時候.
讓我們都能夠成為一間合一的教會.
不單止是紅恩堂.
紅恩堂和黑恩堂都是合一的.
我們都是一個主.
一信一誓.
一個聖靈.
紅恩堂可能會在洪水橋這裡.
做一些工作.
我們在元朗市中心會做一些工作.
大家都是向著主.
大家都是求民主.
在這一刻,我們在這裡的地方.
我們應該要做些什麼.
我們不一定要做些很偉大的事.
但我們必須要做的就是.
人家要我們去做今天要做的事.
很簡單.
伸開.
打開你的手.
將你的物資.
分給身邊的人.
教會裡面的一些.
信是給我們的一些恩典.
要分給我們身邊的人.
其實這樣就足夠了.
求主幫助我們.
讓我們靠著上帝的力量.
把握著這一刻的時間.
我們已經經歷了上帝的愛.
已經經歷了上帝給我們恩典的時候.
我們將這些恩典流傳出去.
以致更多的人.
在這段黃金時間.
能夠經歷耶穌.
我們低頭一起祈禱.
因為你愛我們.

$^{721}$你不會放棄我們.
沒錯,我們每一天都面對很多困難挑戰.
但是主你告訴我們.
我與你同在.
主你的名字本身就是這樣.
神與我們同在.
所以我們就靠你.
在這段時間.
讓我們先去經歷你.
先去經歷從上帝而來那份.
出生以外的平安.
我們就捉緊這份平安.
和我們身邊的人分享.
帶領他們來到你的面前.
認識你.
求主你潔淨.
求主你使用.
多謝你.
我們同心合意在你面前一起.
祈禱奉靠耶穌.
這個名字祈求.
Amen.
多謝大家.
\newpage



\section{哥林多後書 5:17}
\label{sec:atdLdwQVuag}
\textbf{2022.03.27 《新造的人》 哥林多後書 5\_17 (和修版) 曾裔貴牧師}
\newline
\newline
連結: \href{https://youtube.com/watch?v=atdLdwQVuag}{\texttt{https://youtube.com/watch?v=atdLdwQVuag}} ~~~~ 語音日期: 2022-03-27
\newline
\newline
\hyperref[sec:bW0AAWns91U]{\small{< < < PREV SERMON < < <}}
~
\hyperref[sec:index_chronic]{\small{[返順時目]}}
~
\hyperref[sec:index_scriptual]{\small{[返順卷目]}}
~
\hyperref[sec:LGE2dbGZ_40]{\small{> > > NEXT SERMON > > >}}
\newline
\newline
哥林多後書 5:17
\newline
\begin{longtable}{cl}
\hline
\hline
章節 & 經文 (和合本修訂版)\\
\hline
5:17 & \begin{tabularx}{0.7\textwidth}{X} 所以，若有人在基督裡，他就是新造的人，舊事已過，都變成新的了。 \end{tabularx} \\ \\
[1ex]
\hline
\hline
\end{longtable}
$^{1}$各位弟兄姊妹 主女兒們平安.
再一次將舞堂的祝福帶到大家當中.
大家能夠一起在神的家裡敬拜.
以及來聆聽神的話語.
都盼望神祝福我們.
我們一起有祈禱.
再一次我們將大家的心靈.
仰望我們私人的主.
你向我們顯現.
讓我們明白你的心意.
讓我們看到什麼才是永恆.
什麼才是真的值得我們去追求.
也讓我們看到世界的本質.
那個有限.
以及那個在你眼中的軟弱無力.
叫人 特別我們屬主的兒女.
都能夠真的可以來找到這個永恆.
我們在你面前等候禱告.
奉耶穌基督的名 阿們.
今天選取的經文和題目都相當簡單.
我們是一個創造人的地方.
創造新人類的地方.
如果你看英文.
新造的人是一個新的創造.
保羅很清楚.
造物主創造萬有.
六日創造世界.
有猶太背景的.
特別保羅這些拉比.
他們很明白.
舊有的創造不能夠承載神的舊恩.
必須是要有新的創造.
所以在保羅的眼中.
他看到不是他能夠做到這些事情.
十八節說一切都是出於神.
而且這個新的創造.
是要將人帶到神的面前.
這裡告訴我們.
是神自己藉著基督.
使得我們和他和這位造物主和好.

$^{41}$使得我們能夠在神面前.
不單止我們的罪得到赦免.
而且到了第21節.
神使得這些無罪的.
當然就是這位主.
看到這位無罪者替我們成為罪.
以致我們可以在神的面前.
得到的是神的義.
不只是無罪這麼簡單.
得到的是神的義.
在保羅的眼中他明白.
所以今天很想藉這個題目.
五章的十七節.
其中重要的地方就是.
舊事已過.
創造的新人類.
新的族類.
教會負起這個責任.
教會要在這個地上.
告訴世界的人.
什麼才是最最重要.
早幾天星期四.
下午做完事.
出去元朗吃飯.
見到很多人.
很久沒見到這麼多人排隊.
在大會堂.
我走過去看看.
問問那些人排什麼.
看看有沒有我自己的好處.
排了很多人.
那些歲數.
那些年齡層.
那些背景.
全部是基層.
而且是五六十歲.
沒有年輕人.
排什麼呢?.
原來等拿救濟.
急不急切?.

$^{81}$你看到每一個男男女女.
他們的眼神.
充滿了憂慮.
充滿了那種困苦.
很多是低學歷的.
你看到那個隊就知道.
原來是排隊拿一張表格.
填那個一萬元失業救濟.
那就和我無關.
我就不用排了.
但離開的時候.
我作為教會的牧者.
有什麼我可以多走一步.
幫到這些人呢?.
為這些人祈禱.
教會可以做什麼呢?.
那個思緒在這一大段時間.
包括老牧師.
我們大家交通都是.
教會可以多走一步.
做什麼呢?.
我們開同工會不斷問.
教會可以做什麼呢?.
烏克蘭的戰爭.
我們又可以做什麼呢?.
我們就呼籲頂展會奉獻.
去為這個戰爭來祈禱.
和金錢.
我們在今天最後收集那些獻金.
我們就會傳授.
去一些NGO的組織.
去發放出去.
有些事是我們可以做的.
但是.
有什麼更加更加重要.
是教會才能夠做.
一般的地上的群體.
沒有信仰的群體.
甚至乎政府.
都做不到的.

$^{121}$什麼是我們可以做.
必須要做.
是要經常去動員.
頂展會要做的呢?.
我相信.
必然就是這一節聖經.
上下文所涵蓋的.
創造新人類.
讓人進入神的和好裡面.
帶領人的生命的改變.
讓他在地上可能很短暫的.
有些可能很長的.
譬如很小的時候信主.
很大年紀才離開世界.
他的人生就很長.
在教會的人生就很長.
他可以有很多年的教會生活.
他可以有很多的成長.
有些很短暫的.
譬如在病床的泥流.
你能夠做的.
就是帶領他信主.
清楚自己得救.
他就可以得著永恆生命.
但是.
這是一個很重要的訊息.
今天很想真的有一個呼召.
我們有多少人看到.
我們的生命.
是要為神去做呢?.
這裡告訴我們.
舊事已過.
哪些人舊事已過呢?.
必須有一個條件.
要在基督裡面.
你才可以舊事已過.
你的一切.
耶穌為你承擔了.
這裡說的是罪.
那是必然的.

$^{161}$因為下面.
剛才都引述了.
但是這一段聖經.
其實要放在保羅所說的第四第五章.
當然再遠一點就去到第十一章.
你就會更加準確地聽到.
保羅他內心裡面.
所要去擔心的.
哥林多教會他們的情況.
整段聖經.
其實是有一個對象的.
那個對象.
那個對頭.
就是那些假使徒.
在哥林多教會.
要來去核製信徒.
要令到他們歸向他們.
跟隨他們.
來疏遠保羅.
因為保羅是使徒.
他們就要在這個教會裡面.
當保羅離開之後.
就在這裡去慫恿那些人.
就說你們不守那個國禮.
你們不做猶太人做的事情.
你們不跟那些律法.
你就不行了.
那個保羅為什麼他說來又不來.
這個一章那裡所說.
好像他是自然而又飛.
說了說到好像很兇的.
但是來到又很婉轉.
又好像沒有什麼能力等等.
這樣來去批評.
論斷使徒的積分.
接著還要誇耀的就是.
我們是希伯來人.
第八天就行國禮了.
這個是十一章.
保羅那裡所描述的.

$^{201}$這些假使徒的特徵.
所以保羅就說.
就著這些事情.
我都要誇耀的.
如果那些人可以誇耀.
我都要誇耀的.
第十六節.
十一章的十六節.
我再說.
你們不可以看我是愚妄的.
縱然如此.
也要把我當作愚妄人.
來接納.
叫我可以略略的自誇.
我說這些話.
不是奉主明說的.
乃是好像一個愚妄的人.
這樣放膽的自誇.
留意十八節.
既有好些人.
憑著血氣自誇.
我也要自誇了.
你們既是精明人.
用一些很像很.
很威風的寫作技巧.
來讓哥林多教會.
在神面前有一個反思.
到底誰才是真正愛你的使徒.
你是跟那些假使徒.
來走錯了真理.
來跌落一個.
越來越離開了基督的狀態.
譬如十九節.
你們既是精明人.
就能夠甘心忍耐愚妄的人.
你們如果這麼有智慧.
這麼精明.
第二十節.
假若有人強你們作奴僕.
或侵吞你們.

$^{241}$或擄掠你們.
或舞蔓你們.
或打你們的臉.
你們都能夠忍耐他.
我說者說.
是羞辱自己.
好在我們從前是軟弱的.
然而人在何事上勇敢.
我也勇敢.
他們是稀白來人麼.
我也是.
開始在這裡做一個比較.
就是說.
自律行間.
一直由第四章開始.
就已經告訴我們.
保羅要面對的是.
假使徒對哥林多人的影響.
而且這個影響是令他們離開基督的.
所以保羅就不能夠.
來到去.
不出聲.
要來面對這些假使徒.
他們自稱是稀白來人.
就是說.
哥林多你們是外邦的人.
你們不夠我們猶太人的優越.
我們是稀白來人.
接著保羅說.
他們是以色列人嗎.
我也是.
他們是亞伯拉罕的後裔嗎.
我也是.
他們是基督的僕人嗎.
你想想.
一路一路越來越近.
說這班人.
他們是如何誇耀自己的身份.
來歷.
背景.

$^{281}$等等.
接著保羅就說.
我說一個狂話.
我更加是基督的僕人.
大家感受到那個氣氛嗎.
感受到保羅的那個情詞迫切嗎.
感受到.
哥林多人的愚昧.
跟隨著一些假使徒.
離開了基督.
所以保羅在這裡告訴我們.
回到五章這段經文的時候.
我們就看到保羅在這裡告訴我們.
有一個很重要的訊息.
原來.
只要你在基督裡面.
一切的過去.
都過去了.
在神的眼中.
在基督的救贖裡面.
我們每一個都是更新的新人.
其實這個新人這個字.
在新約出現了兩次.
另外一次就是加拉太書的六章十三節.
保羅這樣做一個結論.
他說受不受割禮.
都無關緊要.
不關事.
無關緊要.
受割禮對猶太人來說是很重要的.
不受割禮.
你是外邦人當然是沒問題.
不受割禮都無關緊要.
最要緊是什麼.
什麼是最要緊.
保羅說.
最要緊就是要做新做的人.
就是新的人類.
這裡就是告訴我們.
保羅看這個生命的改變.

$^{321}$在基督裡面的生命的改變.
是看得很透徹.
所以這裡的舊事.
值得我們大家去想一想.
剛才引述了的經文十一章.
再回到我們這裡的時候.
五章的十二節.
就有一點很近的上下文.
就告訴我們.
保羅面對著.
在教會裡面的一個很大的分化.
這個分化是讓一些.
沒有背景的人.
他們不斷感覺自己.
好像是在神面前.
不能夠達到一個水平.
而且還有一種.
正在建立的一種歧視文化.
在教會裡面.
第十二節.
我們不是向你們再舉薦自己.
乃是叫你們因我們有可誇的地方.
留意一下.
好對了.
憑外貌.
不憑內心誇口的人.
有言可答.
原來教會裡面.
有些憑外貌看人.
不憑內心去看人.
可能多教會.
就因為可能種種的原因.
當然更大的.
是先是從那個真理上的出現了問題.
就會對人就會有一個的錯覺.
所以可能這些的假使徒.
保羅這樣說的時候.
相當明顯他們是猶太背景.
他們有一定的律法的修養.
也都是跟隨基督.

$^{361}$但是他們就形成了一種敵對.
形成了一種分門結黨的這樣的情操.
在教會裡面.
導致信徒就受到一個很大的影響.
救恩不可以改變一個人的整個的生命.
還要經過可能重重複複的手續.
不同的階梯.
你爬到上去.
你就可以達到我們的標準.
人的標準.
例如受國禮.
所以有很多外邦人.
很擔心自己沒有受國禮.
你受了國禮.
你不是第八天受國禮.
你就是說不是純正的猶太人.
你沒有做到一些儀式.
你就不能夠好像我們一樣的優越.
這種的種子.
在哥林多教會就很嚴重了.
其實前書已經很白了.
我是屬保羅的.
我是屬基法的.
我是屬基督的.
已經是這樣很清楚.
來直接的去責備哥林多教會一些的領袖.
很明顯就是有這樣的.
對真理的誤解.
而帶來的問題.
所以回到第十七節.
弟兄姊妹.
做了一點點簡單的釋經之後.
我們有一個很寶貴的結論.
原來在耶穌的眼中.
神藉著他的愛子.
將我們去拯救.
我們就能夠在他的裡面.
一起更新.
我們成為了這個新的族類裡面的其中一個.
再沒有分高低.

$^{401}$當然在工作上.
因此上是有的.
臨前十三章已經說了.
十四章也都說.
很清楚的去說了.
其實由十二章的臨前.
一直到十四章.
就是說因此.
不過夾雜著十三章.
就說愛.
在教會裡面.
有不同的分工.
聖靈隨他的己意.
來分給不同的人.
有因此.
但是在生命上.
這一節聖經.
很清楚的告訴我們.
只要我們在基督裡面.
舊的事情是過去.
保羅這樣理解每一個弟兄姊妹.
保羅這樣去看每一個人.
你應該怎樣去在新的生命裡面.
你為誰而活呢?.
所以你要留意到第十六節.
五章的十六節.
當然可以再看多一點.
十四節.
原來基督的愛激勵我們.
因我們想一個人.
既替眾人死.
眾人就都死了.
並且他替眾人死.
是叫那些活著的人.
這個新的族類.
不再為自己活.
只要你和基督連上這個生命關係.
救贖主已經住在你的心靈裡面.
你就是那個新人類.
舊的事情已經過去了.

$^{441}$包括你負面的你的罪惡.
你的個性.
你的各方面.
你在社會上帶來的種種的.
那種人性的錯誤的理解的那些階級.
民族.
政治理念.
學問.
男女.
各方面.
在基督裡面的新人類已經沒有了這一切.
完全沒有了.
要我們教會就是這樣去實踐.
踐行這個新的族類.
創造新人類.
在保羅的眼中.
我們要替這位死而復活的主活.
第十五節.
不再為自己活.
所以這裡可以有一個很重要的訊息.
我們會不會覺得自己是輸在起跑線.
怎樣努力.
你都不會有成功.
和怎樣理解.
怎樣去釐定成功這個定義呢.
錦宣都是一個中產的教會.
普遍的影子會有一套房.
甚至有兩套三套.
隨著整個香港世界的經濟量化寬鬆.
有資產有物業的.
只會隨著經濟不斷的升值.
沒有資產的.
就活在一個很大的困境裡面.
這是已經很基本的整個社會的觀察.
我想問一個問題就是.
是不是這樣就是贏在起跑線呢.
是不是這樣沒有資產的人就是輸在起跑線呢.
誰去釐定那個輸還是贏呢.
這是一個很大的課題.
我們要問自己和教會要回答的問題.

$^{481}$誰是真正的輸.
誰是真正的贏呢.
相信這裡我們有答案了.
第十五節剛才說了.
接著第十六節.
所以從今以後不再憑著外貌去認人了.
雖然憑著外貌我曾經認過基督.
如今卻不再這樣認他了.
不再認基督從一個表面的拿撒勒人.
不是我們的救世主彌賽亞君王.
過去的保羅.
現在因為聖靈光照他.
大馬士革路上面大光照向他.
他的生命改變了.
他看基督完全不同了.
他再重新看舊約.
原來一直那個救恩是為我們預備.
不單止為猶太人.
是為外邦人的.
所以他是外邦人的使徒.
第十七節這個延續.
保羅也這樣看其他人.
不再分任何的因為種族.
因為學識背景.
因為你的經濟能力.
你是否一個女性.
猶太人他們很深的.
對這件事他們的歧視是很強的優越.
因為三次祈禱.
不是女人.
我們識字的.
即是對那些不識字的.
特別外邦人沒有律法的.
很大的鄙視.
這樣祈禱的.
未信的猶太人.
保羅對女性是非常的歧視的.
但是你可以看到他在加拉賴書告訴我們.
如今在基督裡面.
不再分任何的種族男女各方面的背景.

$^{521}$所以我們為這個新人類創造新人類.
我們可以去到多盡呢.
我們可以做多少的事情.
來為主而活呢.
是不是那個起跑線輸贏.
只是你一生有多少在你手上可以擁有的事物呢.
你擁有多少的為主而活.
是主眼中算數的.
那些真正的屬靈資產呢.
不是要否定人類我們大家的努力.
工作.
勤奮.
積蓄.
致富.
不是要否定這些事情.
每個人有你的職場.
有你的跑道.
是神為你安排的.
但這段聖經告訴我們.
你要怎樣令到其他人看到.
包括你自己看到.
這句聖經不是為你自己的.
是為你看到的人.
他怎樣能夠看自己.
是在這個創造的新人類裡面.
再沒有任何的歧視.
有些人為自己的職業.
可以去到很盡的.
大家認識的叫林成斌.
我就不是太認識他.
口才非常好.
樣子又有趣.
很多人喜歡他.
他說話好笑.
一個很出色的電視主持.
原來他還沒做主持之前.
他在新城電台.
就進去大學畢業之後.
就進去做.
但又不是開咪去做那些DJ.

$^{561}$是進去做營業部推銷廣告.
那個commission很高的.
8\%.
如果推銷一個產品.
譬如100萬.
你就有8萬元.
他說他試過.
通常就快掛八號風球.
那些營業員就躲起來看電影.
喝下午茶.
因為快掛八號風球.
你沒理由去別人公司.
最多你打電話.
他說他可以在差不多掛八號風球的時候.
有一個客人差不多簽單了.
他說他特意走到公司樓下.
然後還想一些辦法.
買半打蛋撻.
有手信.
快下大雨了.
特意走到街上淋濕自己.
很匆忙的上去.
按鐘.
剛剛走過.
有蛋撻很香.
放下.
然後就想走了.
真的走了.
那個老闆娘就說.
拿張單出來吧.
不是啊.
我走開附近的.
不是啊.
那個老闆娘說.
你不要裝了.
拿張單出來吧.
他到處找.
真的有.
當然有.
因為他真的有計劃.

$^{601}$想簽好那張單.
可以這麼盡.
可以這樣為自己營銷.
很多人星期日放假.
就玩.
他就星期日坐在新城電台.
準備下個禮拜的單.
誰是差不多可以簽的.
就做好那些準備.
這樣為自己.
所以他幾個月.
他的營業額已經很高了.
這樣為自己打算.
是很有計劃地.
為自己的人生打拼.
不信主的人.
都會這樣為自己.
如果我們.
怎樣去看自己的人生呢.
我們怎樣去看.
在這句聖經裡面.
所涵蓋的震撼性.
誰人在基督裡面.
只要他在基督裡面.
他就立即成為新人類.
拯救一個人.
得著永恆的生命國度.
上幾個禮拜.
星期四晚.
完了查經班.
突然收到一個電話.
弟兄說他的公司的.
女同事的媽媽.
現在離樓在QE.
九點鐘.
準備開車回家.
你可不可以出來.
就二話不說.
其實都未吃東西.
立即開車出去.

$^{641}$都差不多要半小時.
未泊車的車.
行錯了一段路.
要去到紅磡兜回頭.
接著未泊車在停車場.
就已經收到短訊.
剛剛離開世界.
遲了十分鐘.
仍然上去為家屬祈禱.
後來輾轉一路辦喪事.
家屬就說.
原來有另外的同工.
已經在他離樓之前.
來跟他決志.
那個生命就得到拯救.
是多麼寶貴的一件事.
有時趕不及去.
都有一種遺憾.
有一個好像很可惜.
差一點點為何不行呢.
但是感恩.
有同工已經可以在那裡做了.
以致他真的已經得著這個永恆生命.
親愛的弟兄姊妹.
我們怎樣去讓人在基督裡面.
得著這個永恆的國度生命呢.
保羅很清楚告訴我們.
不是他能夠做到什麼.
如果從第四章第七節.
一直看到後面.
一切是出於神.
再一次說.
十八節.
一切都是出於神.
他藉著基督使我們與他和好.
又將勸人與他和好的積分賜給我們.
這個積分已經賜了給教會.
我們去向人.
來勸說.
你要趕緊的與神和好.

$^{681}$有什麼可以改變到一個生命呢.
相信是基督教的教育.
基督教可以改變人的生命.
我們能夠做的.
不僅僅是給一個快速檢測.
其實很多社會人士都在做.
很多立法會議員都在做.
很多人都在做.
我們要做的.
哪些是教會才能夠做到.
我們要認真去做.
就是勸人與神和好.
保羅用勸字.
如果我們能夠用手可以降火.
你信不信.
不信就馬上降火.
每個人都相信.
因為你行神蹟.
每個人都馬上跟隨.
教會就坐滿了人.
不是的.
不然保羅所說的就是.
勸人與他和好.
這個積分.
是勸.
好像我們很對不起他.
好像我們欠了你.
好像我們說的話.
只有我們自己相信.
每個人都看為傻瓜才會相信.
勸.
虎口婆心的虧勸.
我們沒有那種能力.
可以令到人.
由零變一來相信.
保羅告訴我們.
四章第十六節.
我們不喪膽.
外體雖然毀壞.
他已經說得很客氣了.

$^{721}$被人打到半死.
在十一章說他使徒的身份.
內心卻一天生持一天.
我們這些自暫自輕的苦楚.
要為我們成就極重無比永遠的榮耀.
因為我們不是顧念所見.
是顧念所不見的.
因為所見的是暫時.
所不見的是永遠.
最後我用一個例子.
來結束今天的講道.
保羅班德醫生.
是一位.
外科的醫生.
他特別擅長.
針對麻瘋病人的矯治.
由他發現了.
麻瘋病人的誤判.
以為是過濾性病毒.
令到手腳出現問題.
其實不是.
是由他發現了.
麻瘋病人最大的問題.
就是病毒侵害神經系統.
所以手腳已經沒有任何感覺.
以致他們在街上或工作時.
是不會痛的.
不斷流膿,但不會痛.
沒有痛覺神經.
以致他有時會撞得很厲害.
他也不覺得自己有事.
原因是失去了痛覺神經.
因此扭轉了整個醫學界.
對麻瘋的判斷和診治.
帶來翻天覆地的改變.
他的父母是宣教士.
他在印度的醫院.
有一個病人.
非常麻煩.
一來是面容.

$^{761}$因為撞損,撞傷,扭曲和潰爛.
令到他不喜歡自己.
也很恨自己會有麻瘋病.
所以就帶來整個醫院很麻煩.
製造麻煩.
人們看到他也害怕.
因為麻瘋是很恐怖的.
有些人就歧視他.
令他感覺到自己已經不喜歡.
更加製造麻煩.
輾轉.
保羅班德醫生的母親.
一位在印度很多年宣教的姊妹.
對這個病人另眼相看.
來循循善誘她,教導她,幫她.
慢慢地她相信了耶穌.
和洗禮.
這個病人突然有一天問保羅班德醫生.
他從來沒有去過教會.
他知道教會的人對他的歧視.
因為他一直以來.
由他有麻瘋之後.
他就受盡歧視.
他就問保羅班德醫生.
我可不可以跟你星期天去教會.
非常為難.
為難不是因為保羅醫生不喜歡他.
而是為難是怕教會對他有一個.
令他的內心更加充滿憤怒.
保羅班德醫生請求很為難.
他就先在平日去找教會.
跟他們的傳道人說.
這個病人其實已經沒有嚴重的傳染性.
但是他的樣子和他的行為.
是因為他過去的各樣情況.
所以才會這樣.
我保證他真的改變了.
他已經是信主了.
那個星期是聖餐.
印度的教會是同喝一個杯.

$^{801}$我們先不要理會那個潔淨問題.
一起喝那個杯.
教會思考了一段時間.
讓這位弟兄來聚會聖餐.
輾轉在那個星期天.
保羅醫生和這個弟兄去到教會.
前面有一個長老在準備.
早知道這個弟兄會來聚會.
會聖餐.
這個弟兄很凝重.
包括保羅醫生都很凝重.
到底這位在聚會裡的弟兄.
怎樣看這個麻風的新信主的弟兄呢.
很擔心.
那個弟兄因為習慣了受歧視.
被人瞪眼.
沒有人敢看他.
正面都不敢看他.
樣子太醜了.
這位長老走過去搭他的肩膀擁抱.
歡迎他參加今天的崇拜.
稍後有聖餐.
這樣的一個接納.
對這個弟兄來說.
帶來的是一個徹底生命的改變.
從來沒有被接納的.
現在一個教會接納他.
他知道自己在基督裡面是一個新人類.
舊事已過.
包括他自己對自己的懷恨.
對別人的懷恨.
他對自己的人生的坎坷.
一筆勾銷的舊事已過.
輾轉過了一段時間.
這個弟兄在一個工場裡面.
成為一個最出色的工人.
得到那個工廠很多的家長.
就是因為這個擁抱.
就是因為他在基督裡面.
而帶來的一個改變.

$^{841}$是相當不容易.
保羅告訴哥林多人.
你知不知道你們在基督裡面已經是新人類.
你不需要是猶太人.
你不需要受割禮.
你不需要是誰誰的影子.
你就是基督裡面的新人類.
所以求主幫助教會.
無論感宣.
我們在弘恩一起去服事.
我們創造新人類.
讓人進到基督的救恩裡面.
他的生命會改變.
聖靈會在他裡面教導他.
他的生命會有改變.
求主祝福幫助我們.
疫情很快要過去了.
教會準備迎接這個實體的崇拜.
但是我們用什麼來招呼新來的朋友呢.
若有人在基督裡.
他就是新人類.
舊的一切已經過去.
在基督裡面.
他有嶄新的基督賦予他的身份.
我們一起祈禱.
來思想你的話語.
很盼望教會真的準備好.
成為迎接未信朋友的地方.
創造他們的新生命.
進到基督裡面.
改寫他們的命運.
我們這樣感恩禱告.
奉耶穌基督的名.
阿們.
\newpage



\section{馬太福音 26:36-46}
\label{sec:LGE2dbGZ_40}
\textbf{2022.04.03 《願你的旨意成全》 馬太福音 26\_36-46 (和修版) 曾敏姑娘}
\newline
\newline
連結: \href{https://youtube.com/watch?v=LGE2dbGZ-40}{\texttt{https://youtube.com/watch?v=LGE2dbGZ-40}} ~~~~ 語音日期: 2022-04-04
\newline
\newline
\hyperref[sec:atdLdwQVuag]{\small{< < < PREV SERMON < < <}}
~
\hyperref[sec:index_chronic]{\small{[返順時目]}}
~
\hyperref[sec:index_scriptual]{\small{[返順卷目]}}
~
\hyperref[sec:ra6IMFq2wMc]{\small{> > > NEXT SERMON > > >}}
\newline
\newline
馬太福音 26:36-46
\newline
\begin{longtable}{cl}
\hline
\hline
章節 & 經文 (和合本修訂版)\\
\hline
26:36 & \begin{tabularx}{0.7\textwidth}{X} 耶穌和門徒來到一個地方，名叫客西馬尼。他對他們說：「你們坐在這裡，我到那邊去禱告。」 \end{tabularx} \\ \\ \relax
26:37 & \begin{tabularx}{0.7\textwidth}{X} 於是他帶著彼得和西庇太的兩個兒子同去。他憂愁起來，極其難過， \end{tabularx} \\ \\ \relax
26:38 & \begin{tabularx}{0.7\textwidth}{X} 就對他們說：「我心裡非常憂傷，幾乎要死；你們留在這裡，和我一同警醒。」 \end{tabularx} \\ \\ \relax
26:39 & \begin{tabularx}{0.7\textwidth}{X} 他就稍往前走，俯伏在地，禱告說：「我父啊，如果可能，求你使這杯離開我。然而，不是照我所願的，而是照你所願的。」 \end{tabularx} \\ \\ \relax
26:40 & \begin{tabularx}{0.7\textwidth}{X} 他回到門徒那裡，見他們睡著了，就對彼得說：「怎麼樣？你們不能同我警醒一小時嗎？ \end{tabularx} \\ \\ \relax
26:41 & \begin{tabularx}{0.7\textwidth}{X} 總要警醒禱告，免得陷入試探。你們心靈固然願意，肉體卻軟弱了。」 \end{tabularx} \\ \\ \relax
26:42 & \begin{tabularx}{0.7\textwidth}{X} 他第二次又去禱告說：「我父啊，這杯若不能離開我，必須我喝，就願你的旨意成全。」 \end{tabularx} \\ \\ \relax
26:43 & \begin{tabularx}{0.7\textwidth}{X} 他又來，見他們睡著了，因為他們的眼睛困倦。 \end{tabularx} \\ \\ \relax
26:44 & \begin{tabularx}{0.7\textwidth}{X} 耶穌又離開他們，第三次去禱告，說的話跟先前一樣。 \end{tabularx} \\ \\ \relax
26:45 & \begin{tabularx}{0.7\textwidth}{X} 然後他來到門徒那裡，對他們說：「現在你們仍在睡覺安歇嗎？看哪，時候到了，人子被出賣在罪人手裡了。 \end{tabularx} \\ \\ \relax
26:46 & \begin{tabularx}{0.7\textwidth}{X} 起來，我們走吧！看哪，那出賣我的人快來了。」 \end{tabularx} \\ \\
[1ex]
\hline
\hline
\end{longtable}
$^{1}$願意今天你開我們的眼睛,讓我們看到上帝的作為,開我們的耳朵,讓我們聽到你的聲音.願意聖靈充滿我們眾人的心,在我們裡面去工作,讓我們明白神你的旨意.我們這樣的禱告是奉耶穌基督的名字.阿們..
今天已經是四月的第一個星期日,時間過得非常快.我們在四月份有一個很重要的日子,就是壽苦節和復活節.都很想在四月第一個星期日開始講有關的訊息,以至弟兄姊妹可以開始去思考了..
今天有一句說話是「願你的旨意成全」.我相信大家剛才有用心去讀經的話,會發現這個句子出現過.「願你的旨意成全」.這個人是誰來的呢?.
這個肯定是耶穌基督了.但是這樣伏在這裡,我們好像看不清楚.其實這個是耶穌基督的祈禱來的.剛才我們看到出自祂的禱文.耶穌基督在什麼情況下作出這個禱告呢?當時祂的心情如何呢?.
這個禱告是很特別,值得我們去看一下.這個是什麼地方呢?字體原來非常小,這個地方就是赫西瑪尼..
耶穌和門徒去到這個地方,經文裡面說,耶穌說你們坐在這裡,我去那邊禱告,祂和門徒說的.於是祂留了好幾個人在一邊,祂只是帶著彼得和西比泰兩個兒子一起走前一點去禱告..
在這段時間,我們看到經文說到耶穌的心情是祂憂愁起來,極其難過.這個描述在聖經裡面是非常少的.平時我們看到的耶穌是很盡心盡力,很用心去服侍人群的..
你看到的祂是不會受到任何環境的影響.我們一直在看的耶穌,我們可以覺得祂是一個非常勇敢的人,一直情緒各方面都很正面..
雖然祂有恩著人的犯罪,祂有祂憤怒的事,例如在聖殿裡面,祂看到人在買賣,將聖殿變成賊窩的時候,祂為上帝的緣故,祂有祂憤怒的事..
祂也有恩著在拉薩路死亡復活的整個事件裡面,為了人心,耶穌有難過,有哭過..
但是整體來說,要說祂很憂愁,極其難過,是沒有什麼的.就在這裡,我們看到是描述祂的情緒,很深刻..
忽然間我們眼前睜大了一眼,這是在形容耶穌?是的,是在形容耶穌,而且耶穌自己還要說,祂說我心裡面非常憂傷,幾乎要死..
憂傷到差不多死,很深很深的痛苦,這一刻很難受,這是耶穌心底的話,耶穌會的嗎?耶穌會憂傷到差不多死的嗎?.
接著,你看到祂再離開彼得和西比太兩個兒子的一段距離,大概一塊石頭扔出去的位置,祂就跪下禱告,祂說父啊,如果你願意,求你將這杯切去..
而而不是,照我的意願,以至要成全你的旨意.祂就作了這個禱告..
這一段經文,我抽自路加福音,為什麼要抽路加福音呢?原來在這個禱告裡面,我們看到耶穌的情緒是怎樣的呢?.
剛才說祂說我憂傷,憂傷到差不多死,祂自己也說我很憂傷,在這裡說祂憂傷是怎樣呢?.
就是說祂當時是非常痛苦焦慮,這是路加福音的描述,祂禱告的時候更加的懇切,祂的汗好像大血點滴在地上,很誇張,是很大滴,好像血一樣,祈禱更加的懇切,真的很焦慮,很痛苦,真的很難想像..
祂為什麼會痛苦,會這麼焦慮呢?是什麼令祂憂傷痛苦呢?.
神的兒子耶穌基督,也是人的兒子耶穌基督,這一刻在經歷什麼呢?這一刻耶穌基督是以祂作為人的兒子的身份,在經歷一件很重要的事,但是祂感到很憂傷,很痛苦,我們要從祂祈禱的內容去看,究竟因為什麼事而憂傷痛苦?.
祂稍往前走的時候祈禱,我父啊,如果可能,求你將這個杯離開我,然而不是照我所願,而是照你所願的,祂的祈禱就是這樣,將這個杯離開我..
究竟這個杯是指什麼呢?這個杯令祂非常憂慮痛苦,原來這個杯是一個很豐富的含義,不只是指著患難,苦難和受死,而是指著上帝的憐怒,憤怒..
上帝對著什麼憤怒呢?對著人犯罪很憤怒.我們很多時候都會說一個信息,就是上帝就是愛,我最喜歡說上帝擁抱我,上帝呵護我,照顧我,我只喜歡說上帝的慈愛,教會最喜歡說上帝的慈愛,一天都說..
我們很少提到上帝的憐怒,上帝是憤怒的,在哪裡看到上帝的慈愛,你才要由他的憤怒開始說起,你才看到他的慈愛有多深.上帝看到人犯罪,看到人犯罪,很憤怒.特別是人是不斷累積上帝的憤怒,因為人會不斷犯罪的..
就算基督徒信耶穌之後的人,我們都還在犯罪,在過程中我們在累積上帝的憤怒.在這裡說,上帝的憤怒在這個杯子裡.這個杯子裡還有上帝的憤怒不是對著一個人,兩個人,上帝的憐怒是對著我們全世界所有的人..
是很多的,多憤怒.你說從亞當夏娃開始犯罪,得罪上帝開始,罪惡進入世界之後,一直到今天,人做了什麼事,人犯了什麼罪..
其實在過程中,我自己很深刻,不止一次,我去看人的基督教,我自己出去傳福音,我心裡很深刻,或者我思考自己的罪惡,我都很深刻..
我試過有出去跟人傳福音的時候,我初時帶著一個心,今天又有機會傳福音,很興奮.而對方已經說,我很想信耶穌,我就去跟他講福音,跟那個對象講..
講著講著,我也很想認識一下這個人,他是怎樣的呢?對方就跟我說他人生經歷過的事情,一講之後,我就發覺,如果不是因為我信耶穌,真的會去接納犯罪的人..
我相信,我一聽到他講這些事情,我是很不喜歡這一類型的人,因為他所犯的那類罪,是我很討厭的,真的覺得是對人的家庭,對人很大的傷害..
我一邊聽的時候,我心想,我平時一定會跟這些人分開,如果我沒有信耶穌,我不會跟他交朋友,我更加不會做什麼..
但是今天我信了耶穌,上帝接納了我這個人,上帝也想我接納其他人,所以在這個時候,我繼續跟他傳福音,到最後我帶他信耶穌的時候,我要跟那個人說,其實剛才一直說,原來那些那些那些那些是得罪上帝的..
但我事後再想,其實那次那個人是信了耶穌,後來他也接受栽培,慢慢慢慢一直成長..
但這個過程,我想回頭,我真的想不可思議,如果是一個沒有信耶穌的我,我真的不可以接受這種人..
接著,最近我自己想起來,我聽到一些人的事蹟,有些人生前做了很多很多傷害別人的事情,但到最後還是信了耶穌,上帝的憐憫臨到,上帝真的很愛人..
在這個過程中,我不停不停地想,上帝真的很愛,很愛那些人,這樣的人也信耶穌..
但回頭想起我自己,我又不是一樣是一個罪人,我很好嗎?我一直看著那些人,你說別人怎樣壞怎樣壞,這樣也信耶穌,那我又不是,我是一個很壞的人,我這麼壞的人也可以信耶穌,上帝也憐憫我..
其實上帝對這麼多人,對所有的世人都是很愛,但我記得,我一樣,在我出生,到我未信耶穌之前,甚至我信耶穌後,也是一樣,一直都是要認罪的,一直很多罪,我的罪孽是一直激怒上帝,我也在累積著神的憤怒..
我們全世界的人,不要說誰犯罪多誰犯罪少,我們全世界的人在做這件事,是使得上帝非常憤怒,這一切的憤怒就要傾倒在耶穌基督的生命裡面,是這一切的憤怒..
然後就是我們非常熟悉的苦難,耶穌基督將要面對的,是一連串的苦難,他將要被他的門徒出賣,被捉拿,經歷不公義的審判,被戲弄鞭打,最後被釘在十字架上,一連串的..

$^{41}$原來這個杯是這麼不容易的,這麼艱難的,其實這個杯不是一般的杯,我想這個杯只有一個人可以承受的,就是耶穌基督,只有他有這樣的資格,因為只有他完全沒有犯罪..
我們全世界的人都犯了罪,全部得罪上帝,我們全部都不可以領受這個杯,不可以為別人承受,其實這個要為別人承受..
本來我們要受的就是經歷上帝的憤怒,面對這一切憤怒,最直接的後果是永遠的死亡,本來我們要面對的是直接的,就是永遠的死亡,本來的苦難我們要去承受,今天耶穌基督要代我們去承受這一切..
在這裡我想給大家看一幅圖畫,在後面的《馬太福音》27章45節是說耶穌基督被釘的那一天的情況,從正午到下午三點鐘,遍地都黑暗了..
這是何等令人感到心酸的一個很難過的場面,耶穌基督要去經歷,在裡面耶穌基督要去喝這個杯..
其實當他要面對上帝的烈怒的時候,還有一樣東西他要經歷,這個杯帶來的很大的痛苦是什麼呢?就是下午三點鐘的時候,耶穌大聲地高呼說:.
「以彌以彌拉瑪撒巴國大尼」,就是說:我的神,我的神為什麼離棄我?原來當他承載著全世界所有人的罪惡的時候,他要喝下這個杯的時候,上帝掩面不看他,在那一下子上帝要離棄他..
聖父聖子一直都在一個很緊密的很親密的團契裡面,是非常非常親密的關係,沒有什麼是難了他們之間的團契,但是就在這裡,他為了承載全世界的人的罪,為了這些罪去付出代價,就在那一下子上帝要離棄他,上帝掩面不看他..
這個在神學上當然引起很多的討論,是怎樣的離棄呢?是怎樣的掩面不看呢?但是我們知道那一下子耶穌所承受的痛苦是最大最難忍受的,他其實完全沒有犯罪,就是為全世界的人去承受這個罪的後果,為了這樣,要去面對一個很大的痛苦,和上帝分離..
這個根本是受不了的..
耶穌在早上九點的時候被釘在十字架,一共有六個小時在十字架上,到了下午三點左右,在這個前段的時間,他一直經歷很多很多不同的苦難,被人在下面嘲笑,經常恨他,你快點離世吧,譏笑他等等..
後面他經歷的是和上帝和父神的分離,上帝的烈怒傾倒在他裡面,他沒有犯罪,但是上帝的烈怒傾倒,因為他為世人承受這一切..
這是很難以忍受的.耶穌基督所要去承受的,是遠比殉道更加艱難..
我記得很多年前有機會認識一個姐妹,這個姐妹跟我說上帝呼召她,她說我跟你說,呼召的過程是怎樣呢?.
話說有一次我參加一個聚會,我看到大會播放一些影片,說到有一個國家,說到有一個國家,你去到那裡傳呼音,你將會面對以下的情況..
什麼情況呢?就是看到一個媽媽和女兒的圖片,是怎樣的?就是她們被肢解了,女兒很小,頭,手,腳都被斬開了..
我當時看到這個照片,我的心跳得很快.我知道上帝要呼召我去這個地方宣教,但是我的眼淚不停地流..
原來在這個國家裡,她們就是這樣對待基督徒,就是這樣傷害他們,拿他們的性命,是會殉道的..
但是上帝就要動我,要感動我,問我要不要去,就要去這個國家.換言之,這個姐妹如果將來答應去這個國家,也有機會會殉道..
其實這只是其中一個例子,殉道的例子很多,犧牲的方法有很多類,各式各樣不同的犧牲方法..
很多時候說起這些例子都是可歌可泣的,令人很感動的說,哇,殉道是多麼慘烈的事情,這個犧牲很大,這個痛苦很大,一直說..
但是其實我想說,耶穌基督那一刻要承受一定比殉道更加艱難..
殉道沒有錯,是為福音的緣故,為上帝的國度,是慘烈的,是痛苦的..
但是我想說,殉道不是為了全世界人的罪,不是承載全世界人的罪,為全世界人的罪去受這個死刑,是永遠的這麼大的痛苦..
不是,耶穌基督為我們承受這個刑罰,是有分別的,全世界只有祂一個可以這樣做,只有祂一個,所以祂的痛苦是獨一無二的..
你想想,全世界人的罪都壓在了他的生命裡面,壓在了他的生命裡面,上帝的憤怒傾倒在他的生命裡面,他要承受和天父的分離,有誰試過呢?.
殉道者不是這樣,殉道者不是為全世界所有人的罪惡,不是為他們而犧牲,不是,殉道者也做不到這個角色,唯有耶穌基督..
耶穌基督的死也是獨一無二的,縱然這個世界有很多死,何國可入,很震動人心,但是不及耶穌基督的死..
但是這一刻還沒有去到那個位置,耶穌基督在克西瑪利元禱告的時候,是在受苦和釘十字架之前的,在經歷和上帝分離之前的那一下,但是祂知道祂將會這樣,在知道祂將會這樣那一下,祂心裡面就很痛苦很憂傷..
其實這個痛苦和掙扎才是真正的辛苦,當你知道你將要去面對一個很大很巨大的痛苦,你知道你將要這樣的時候,之前的心情,那種心理裡面的痛苦是很大很大,你會發覺自己非常的孤單..
沒錯,在另外的經文裡面也有說,天使有加添給他力量,但是即使加添給他力量,也是他自己要去面對的,那些痛苦和分離所承受的一切,也是要他自己去面對..
所以在這一刻,耶穌真的很真心地說,我很憂傷,我很難過,我們有多明白?.
結果,他為了這個情況,進行了三次的禱告,三次的禱告都是在他的掙扎裡面很真實地,他很想求天父挪走這個杯..
經文沒有描述天父有什麼回應,只是說耶穌進行了三次的禱告,每一次都是這樣,求神挪走這個杯..
但是他也會說,如果這個杯不能夠離開我,必須要我去喝,那我就願你的旨意成全了..
如是者,他一直這樣向上帝禱告,向上帝挪走這個杯,他一直在掙扎.其實說真的,他可以選擇的,他可以選擇不要這個杯..
在創造世界之前,他和天父在一個很甜蜜的很緊密的團契裡面,他享受一切的豐盛,為什麼他要受這些苦呢?.
他來到這個世界,他作為人的兒子,他沒有犯過罪,他為什麼要為全世界的人去喝這個苦杯呢?還要面對天父的烈怒,要分離,要這麼多痛苦,是不是可以選擇不信服呢?.
他可以選擇不去喝這個杯,可以這樣選擇,連祈禱也不祈禱也可以,就這樣選擇,就去吧,不喝這個杯了,可以嗎?.
如果這樣的話,會發生什麼事呢?我想問,我不要這個杯,可不可以?.

$^{81}$最後你看到他是選擇信服的,他選擇信服,他以信服天父旨意為最重要的,三次禱告之後,就是真正願你的旨意成全,他信服天父..
我可以這麼說,其實赫西曼利元那一晚很重要,是一個掙扎的時間,是一個選擇的時間,有兩條路可以選擇,直接就不掙扎,就不喝這個杯,所有這一切我都不承受,是不是?.
我享受由創世之前一切的豐盛,平安,第二個信服天父旨意,就是要走下去,如果耶穌基督所選擇的第一種就是我不會去喝這個杯,我們會發生什麼事呢?.
就像這幅圖畫,不是說十字架那一邊,是另一邊,其實那條路很闊的,那條路是走向滅亡的路很闊,很多人行在其中,耶穌基督那一晚在那個掙扎,赫西曼利元的掙扎,如果他選擇放棄,他不要喝這個杯的時候,我和你今天肯定不會在紅印堂..
我深信教會不會出現,其實我們全部都是繼續背著我們的罪惡,我們的罪的擔子,直走向永遠滅亡,我們人生沒有盼望,人沒有盼望在永遠的黑暗裡面生活,如果那一晚在赫西曼利元,如果耶穌的選擇是這樣,.
有多少人會走向永遠的滅亡?我們今天也不需要在這裡說一些同道,我們根本不存在在教會裡面,教會也不存在,但是就是因著耶穌基督信服天父的旨意,在之後走下去,去喝了這個苦杯,經歷過一切一切的時候,.
這個十字架的橋樑,這個生命的橋樑搭建起來,我們才可以回到天父的身邊,我們才可以得蒙拯救,生命又盼望..
原來是要這樣,付一個這麼大的代價,但是你有沒有想過在赫西曼利元那一晚,耶穌多麼痛苦,知道即將要面對這一切,心裡面的擔子有說不出來的重,願你的旨意成全,其實這句話不是這麼輕易說的,隨便說,沒所謂,你說怎樣就怎樣,不是這樣的,那句話很重..
願你的旨意成全,是不是能夠說得出口?.
在這裡,當我們重溫這一晚的片段的時候,重溫赫西曼利元那一晚,耶穌的掙扎,耶穌的痛苦的時候,我想問問弟兄姊妹,也問問我自己,我們是不是真的明白耶穌基督為我們付出了什麼代價?.
是不是真的知道,真的知道祂這麼苦,每一年的受苦節,復活節是不是很恆常地,很隨意地渡過,就這樣,隨便說說這些信息,聽過一下,就過去了,我們是不是真的明白,真的知道,有沒有回應呢?.
當我們知道耶穌原來是承受這麼大的痛苦,為我和你付出這麼大的代價的時候,我們願不願意愛上帝更多,愛耶穌基督更多?.
愛不是在嘴裡說愛,是有行動的,有行動,我想挑戰弟兄姊妹,在我們接下來的四月份,整個月裡面,除了思想耶穌基督的受苦和復活,我挑戰弟兄姊妹想,你打算怎樣愛耶穌?.
你今天在你生命裡面,耶穌是佔第幾位的呢?.
是不是我們首先排列好家人們,好友們,因為耶穌排列得很厚的呢?.
還是我所愛的事業,興趣,一切都排列得很先,耶穌排列得很厚的呢?.
在這個四月份,受苦節,復活節之後,你願不願意重新思考你和耶穌的關係?.
我們自己去回應神,我們一起禱告..
你讓我們明白,耶穌真的為我們承受了很多..
你的受苦和你的死,都是獨一無二..
你選擇了去信服天父的旨意..
正正因為這樣,之後你非常勇敢地面對每一個苦難和傷害..
主耶穌,因著你所擺上的,今天我們有何等的盼望?.
因著你為我們付出沉重的代價,我們得到了一個新的人生..
我們和天父復和,每天經歷天父的愛..
在這裡,耶穌基督很求你,真的感動我們,去愛你更深..
你配得,你配得我們為你擺上一切..
你配得我們將你擺在我們人生的第一位..
你配得我們去尊崇敬拜,讚美你..
願在這個受苦節,復活節,你給我們的人生有一個很大的改變..
你挑戰我們一群基督徒,一群跟從你的人,生命要去成長..
阿們..
\newpage



\section{創世記 1:26-28}
\label{sec:ra6IMFq2wMc}
\textbf{2022.04.10 《因為愛,所以創作》 創世記1\_26-28;2\_18-23;3\_7,21 (新普及譯本) 陳朗欣傳道}
\newline
\newline
連結: \href{https://youtube.com/watch?v=ra6IMFq2wMc}{\texttt{https://youtube.com/watch?v=ra6IMFq2wMc}} ~~~~ 語音日期: 2022-04-10
\newline
\newline
\hyperref[sec:LGE2dbGZ_40]{\small{< < < PREV SERMON < < <}}
~
\hyperref[sec:index_chronic]{\small{[返順時目]}}
~
\hyperref[sec:index_scriptual]{\small{[返順卷目]}}
~
\hyperref[sec:DcwFkXHzJa4]{\small{> > > NEXT SERMON > > >}}
\newline
\newline
創世記 1:26-28
\newline
\begin{longtable}{cl}
\hline
\hline
章節 & 經文 (和合本修訂版)\\
\hline
1:26 & \begin{tabularx}{0.7\textwidth}{X} 神說：「我們要照著我們的形像，按著我們的樣式造人，使他們管理海裡的魚、天空的鳥、地上的牲畜和全地，以及地上爬的一切爬行動物。」 \end{tabularx} \\ \\ \relax
1:27 & \begin{tabularx}{0.7\textwidth}{X} 神就照著他的形像創造人，照著神的形像創造他們；他創造了他們，有男有女。 \end{tabularx} \\ \\ \relax
1:28 & \begin{tabularx}{0.7\textwidth}{X} 神賜福給他們，神對他們說：「要生養眾多，遍滿這地，治理它；要管理海裡的魚、天空的鳥和地上各樣活動的生物。」 \end{tabularx} \\ \\
[1ex]
\hline
\hline
\end{longtable}
$^{1}$各位紅印堂的大兄姐妹,平安.
很高興可以再次在線上見面.
同工再次向大家問安.
我今天的講題是「因為愛,所以創作」.
我們聚焦於創世記的故事.
來看愛,創作,甚至管理等關係.
我們先聽一段經文.
可否到我的簡報?.
起初,上帝創造天地.
地是空虛混沌,冤綿黑暗.
上帝的靈,運行水面之上.
上帝說,要有光,就有了光.
上帝看著光是好的,就將光暗分開了.
上帝稱光為日,稱暗為夜.
有晚上,有早晨,這是第一天.
上帝說,諸水之間要有,事就這樣成了.
上帝稱空氣為天,有晚上,有早晨,是第二天.
上帝說,天下的水要,事就這樣成了.
上帝稱旱地為地,稱水的聚集處為海.
上帝看著是好的.
上帝說,地要生出青草,事就這樣成了.
於是地生出了青草和結種子的菜蔬,各從其類.
和結果子的樹木,各從其類.
果子都包著雲,上帝看著是好的.
有晚上,有早晨,是第三天.
上帝說,天上要有,事就這樣成了.
上帝看著是好的.
有晚上,有早晨,是第四天.
上帝說,水要多,上帝看著是好的.
上帝就賜福給這一切,說.
滋生繁多,充滿海中的水,雀鳥亦要多多生在地上.
有晚上,有早晨,是第二天.
上帝說,地要生出活物,各從其類.
生出昆蟲,野獸,各從其類.
事就這樣成了.
上帝看著是好的.
上帝說,我們要照著我們的形象.
按著我們的樣式做人.
使他們管理海裡的魚,空中的鳥,地上的生畜和泉地.
和地上所爬的一切昆蟲.

$^{41}$上帝就照著自己的形象做人.
是照著他的形象做男,做女.
上帝就賜福給他們,又對他們說.
要生養眾多,遍滿地面,治理大地.
亦要管理海裡的魚,空中的鳥,和地上各樣行動的活物.
上帝說,看看,我做了這麼多事就這樣成了.
上帝看著一切所做的都非常好.
有晚上,有早晨,是第六天.
是第六天.
主上帝用地上的塵土做人.
將生氣吹入他鼻孔裡.
他就成了有靈的活人.
主上帝在伊甸的東邊做了一個園地.
將所做的人安置在裡面.
主上帝又洗地,生出各種樹木.
他們有又好看又美味的果實.
園地中間有生命樹和分別善惡樹.
有一條河由伊甸流出來,灌溉園地.
那裡又分為四條支流.
主上帝將那個人安置在伊甸園裡面.
讓他看管園地的一切.
主上帝就這樣安置在伊甸園裡面.
如果我們有機會去加拿大.
假設你那裡有朋友.
你一家人去探望他們一家人.
你們想像坐了一場飛機.
然後要處理所有檢測的事情.
你一定會很累.
然後你朋友在機場接你.
然後送到你們家.
一打開門,整間屋都是暖氣.
有熱辣辣的朱古力.
然後等你休息一下.
享用熱辣辣的晚餐.
很豐富的現食.
還有客房已經預備好了.
一切都預備好妥當.
只待你來到和一起享受共聚的時光.
其實上帝的創造都是一樣的.
六日裡面的創造像詩一樣優美.

$^{81}$有對稱的結構.
有重複的關鍵詞.
並且在創造中間的高潮.
其實就是神做了人.
整個故事的情節.
又好像一對父母預備小朋友出生一樣.
名字都改好了.
然後預備好嬰兒一切要用的東西.
無論是床,衣服,奶粉等等的用具.
安排好照顧嬰兒的人.
無論是陪月員或是安排外母,姐姐等等.
所有東西都預備好.
只待嬰兒來到.
為什麼呢?.
因為父母愛他.
將最好的為他預備.
神做人的故事其實都是一樣.
他預備最好的給人.
整個地球和宇宙創造好了.
還有適合人類和一切就造物居住的環境.
空氣,水,陽光.
還預備了裝飾.
想想日出的天空是甚麼顏色.
日落的天空是甚麼顏色.
晚上的天空有甚麼.
就是滿佈天空水.
好像沙一樣多的星星.
是星雲.
還有很多數不完的不同的動植物.
又漂亮又可以做食物的水果.
預備了不同的地貌.
還有一個花園給人居住.
那花園還有四條河.
河下面有不同的金屬寶石.
這是甚麼呢?.
這是愛.
在人類未來之前.
神已經預備好了.
在我們未出生之前.
已經為我們預備好了.

$^{121}$為你籌謀.
這是愛.
人在被愛的環境下就造.
在未出現之前已經是被愛的對象.
無論我們身處甚麼地方.
身為甚麼民族.
我們有甚麼文化.
神很愛我們.
我又想問問你.
誰又是你愛的對象呢?.
你又會為他而預備的呢?.
你會為他而創作花心思的呢?.
我們作為基督家庭的一份子.
誰又是我們愛的對象呢?.
在這個雄恩群體裡面.
在這個社區.
在這個城市.
或是這個世界.
誰是我們作為基督家庭的一份子.
愛的對象呢?.
我們剛才看了.
上帝因為愛而做人.
接著我們看看.
上帝給人的召命.
其實都是和創意和創作物不可分.
人被上帝委任的第一件工作就是管理.
上帝說我們要照著我們的形象.
按著我們的樣式去做人.
然後讓他們管理世界所有的受助物.
上帝就這樣按著他們的形象.
做了男做了女.
並且賜福給他們.
之後又跟他們說.
你們要治理,管理受助界.
在二章七節上帝.
聖經形容上帝將塵土造成一個人.
然後將生命的氣息吹入他的鼻孔.
這個人就成了有靈的活人.
同樣地他用泥土造動物.
但他沒有將生氣吹入他的鼻孔.

$^{161}$這就是人和所有受助物不同之處.
人有上帝的靈.
氣息其實是上帝的靈的意思.
萬物之靈是沒有錯的.
我們在地上作為上帝的代表去管理大地.
很有趣.
管理這個字很多時候都會覺得是高級人員.
古代甚麼人要管理呢?.
是國王才要管理.
他要管理他的臣民,土地.
但當我們在歷史當中發現很多管理的例子.
一點都不算好.
他用他的權力去轄制人.
君要臣死 臣不得不死.
但在聖經中我們看到一個角度.
是跟愛和創意不可以分開的.
當上帝委任人類去管理大地的時候.
而聖經記載阿當的第一個工作是甚麼呢?.
是改名.
是有權柄的人物和其他人改名.
正如父母有權柄和子女改名.
這是一個權柄.
父母有權柄和受眾物改名.
上帝也會改名.
但很有趣.
仔細想想這是甚麼工作?.
是改名.
你不會隨便改名.
你叫1,2,3,4.
上帝我做完了.
不會的.
其實你會用心去做.
而且是一個創作.
他叫甚麼名字好呢?.
因為他不是廣東人.
他不會用廣東話改名.
我們香港人用廣東話也會明白到.
例如看到一隻大鼠.
我們叫牠袋鼠.
看到一隻馬有斑紋.

$^{201}$我們叫牠斑馬.
看到一隻鹿的頸很長.
我們叫牠長頸鹿.
其實你會按照牠的特徵.
觀察牠之後改名.
這是要花時間的.
花時間的事其實是甚麼?.
背後是你要用心.
你要資源.
其實是一份愛.
你不隨便做其實是愛.
但你要做的工作是甚麼?.
是創意.
是創作的工作.
即是說管理的工作.
都是跟愛和創意分不開.
管理這個位份.
管理大地這個職份.
是跟愛和創意分不開.
我再問.
我們作為基督家庭的一份子.
誰又是我們愛的隊長呢?.
無論是這個群體.
這個社區.
甚至面對著城市.
甚至作為世界的受造界的代表.
面對著不同的民族.
誰是我們愛的隊長呢?.
我們知道在這個月.
就是穆斯林的齋戒月.
他們是一個很大的.
未得之民的群體.
每一年闖開的門.
這個機構都會公佈.
全球首望名單.
就是全世界頭五十個.
最迫不及待的基督徒的地方.
過去頭一位是北韓.
今年下跌了.
變成阿富汗是第一位.

$^{241}$北韓是第二位.
你想像到那裡其實是有教會.
有基督徒的.
但他們受著很嚴峻的歧視.
或者迫不及待.
無論是工作不順利.
甚至人生資源有限.
甚至生命有危險.
或者人生有安全的危險.
我們來看看.
這個在伊朗的一家人的故事.
(電話鈴聲).
(悲傷的音樂).
(阿富汗語).
(悲傷的音樂).
(阿富汗語).
(悲傷的音樂).
我想上帝吩咐我們去愛倫社.
其中一個很重要的元素就是.
你見到便去做.
你知道的便去做.
剛才那段影片是《闖開的門》準備的.
有程式,網頁,中文版本.
我們可以了解世界各地.
不同的弟兄姊妹發生的故事.
而《闖開的門》的創辦人安德烈弟兄.
其實本身他和我們一樣是基督徒.
在50年代他去了波蘭.
當年是一個鐵幕國家.
他參加全球共產主義青年大會.
他不是共產黨員.
他是基督徒.
他想了解一下發生什麼事.
當他去到波蘭.
亦有機會去探訪波蘭的教會.
在聚會的時候.
他還未說話的時候.
當他分享完後.
當地的波蘭牧師說.
我們以為全世界遺忘了我們.

$^{281}$但你來到.
你不用說話.
我們都感受到神的愛和同在.
他發現了原來在鐵幕國家背後.
很多基督徒被忽略,遺忘.
他便開始了支援鐵幕國家背後的基督徒教會.
他運送聖經給他們.
給他們資源和各種需要的東西.
很有趣.
他當年駕駛一輛甲蟲車.
他運送聖經和基督教的單張資源.
他每一次出入關口都會祈禱.
祈求神當年將盲人開眼.
現在求神令不應該看到的東西.
官員的眼睛被他弄盲.
讓他們看不到不想看的東西.
他每一次都成功出入境.
把聖經運送到共產鐵幕國家背後的教會支援他們.
到今時今日他們的工作.
發展到關注全球很多國家.
例如新冠肺炎的時候.
在初期兩年前.
在利比利亞的基督徒婦人.
因為她是基督徒身份.
雖然丈夫過世了.
是牧師被人殺了.
但她很辛苦等到政府可以派救濟物資.
政府不派給她.
因為她是基督徒.
結果坐她的門的同工支援他們.
讓他們渡過難關.
我們是一個神呼召的群體.
很特別.
在這個地上代表他.
我們有政治的本質.
因為我們是天國的群體.
無論我們身處香港或世界各地.
都有他的心意.
他讓我們成為香港人.
有不同的創意和才幹.

$^{321}$是有意思的.
我們一起去努力.
無論是利用創意去管理所有資源.
無論是時間,才幹,甚至情感.
我們去看看神給人的第二個召命.
是愛裡面創作.
神給人的第二個召命是生子.
上帝叫亞當夏說你們生多些.
披滿全地.
其實男女結婚在婚姻裡面再有生命.
是一個很奇妙的事.
彷彿體會再造生命的能力.
很有趣.
就算一個家庭裡面的不同小朋友.
都是獨特的.
每一個都是一樣.
性格,外貌都是獨特的.
在第二個任命.
生養眾多.
延伸到創作.
由生子生養眾多.
到從事多些.
和創意有關的事.
無論是衣服的配搭.
試試新的髮型,配搭.
行地路試試新路線.
煮菜試試新菜式,煮法.
甚至表達思想,過去的思維.
試試新的.
試試這樣看或學習去說.
其實在創作學習裡面.
你會不斷感受到很多可能性.
很多快樂.
亦是認識自己的過程.
認識自己,接納自己的過程.
這是藝術創作可以得到的.
愛是可以令人創作的.
我們記得小朋友.
我們小時候都自發地畫張卡.
寫「爸爸,我愛你」.

$^{361}$很簡單的動作.
是因為愛.
你會覺得文字,語言都說不完.
你要用藝術品去做.
做一些東西去表達你的愛.
剛才看的這段影片.
是在去年下半年.
我有機會在宣導會的南區一間教會.
深恩堂.
很榮幸地擔任藝術傳道的角色.
並獲得南區中學的校長邀請.
我們在他的校園裡有一個壁畫.
我邀請了一批學生.
我們跟他一起創作.
構思如果這面白色的牆.
要畫一個打卡位.
有學校的校徽.
和突出你們學校的特色.
你會畫些甚麼.
我們一起創作.
結合幾十個學生的畫稿.
結合成一幅完整的壁畫.
你會看到我們幫手潤色的藝術家.
背後有顆雙手捧著的大珍珠.
大珍珠上有學校的校徽.
這個設計象徵著這間學校的上上下下.
好像上天手上的掌上明珠一樣寶貴.
畫裡面還有很多東西.
每一樣東西都有意思.
正如聖餐或禮儀.
每一樣東西都有意思.
但我作為主要的統籌人.
當我畫的時候.
我感受到我將所有對學校和學生的祝福和愛.
都畫在上面.
例如筆日光.
筆日光的產生是地球的磁場和太陽粒子.
遇上的時候遇上圈內的筆日光.
其實都想勉勵他們.
人生有時機遇來到.

$^{401}$你就會發出光芒.
還有很多東西.
全部都是一些創作.
結合我們對他們的祝福和愛.
這樣放在上面.
當有個電話生產商找ERROR.
12個男生拍攝很大型的手機廣告.
我們教會有個傳道.
用二次創作.
用同一個姿勢.
邀請了12個弟兄.
拿著12本聖經.
去打Ekklesia.
被呼召出來的群體的希臘文翻譯.
其實教會都是這樣的.
之前政府說要限聚.
宗教的住所關閉.
私人地方都不可以多於兩個家庭.
於是我們教會的傳道.
和我們的設計師想了一個貼文.
表達一家人.
教會有多少家人.
是兩家還是N家.
還是主羅一家.
去表達這些反思.
在這個大齋期.
我們有一個很特別的信仰特別行動.
是用以色列先知的藍本.
以行動體會信仰的堅持.
男士是40天不太梳.
兼不剪頭髮.
女士是40天有一個特別的髮色.
所以這40天大部分時間.
我之前病了.
總之病好之後.
我就是不斷夾不同的髮型.
都是一個學習對信仰的堅持.
由生活上一些小的行為.
再去體會對信仰的堅持.
這位姐妹其實都是我妹妹.

$^{441}$她是孫獨煒 沙葛台.
她是一個浸大創意媒體畢業的年輕導演.
現在她是從事殯儀的.
很有趣.
對我們來說殯儀就是賺錢的行業.
過去殯儀都是被不同的群體壟斷了.
這個行業.
這個女生入行是帶著她的照明.
第一 她體會到每一個喪禮都是一個創作.
是一個服侍喪家的創作.
很想藉著她的專業.
無論是拍攝 設計和創意.
每一個度身訂做去服侍喪家.
另外她很想轉化這個行業.
她發現行家的自我形象都比較低.
很多人都會問她.
你讀完大學為何要來做殯儀.
但她不是這樣看.
她不是覺得殯儀一定要沒有讀過書的人才做.
我覺得很欣賞神對她的呼召而她活出來.
最近她在三月.
幾個月前在我媽媽的安息日.
對不起.
我媽媽的死忌.
她回到天家的日子.
她送了一束花給我.
她說這是送給你媽媽的.
放在我媽媽的殺灰紀念花園.
我體會到她很用心和我曾經喪失親人的人同行.
由籌備喪禮到禮儀.
到過了哀悼期.
她仍然會記住這份創意和愛和專業.
其實我一直經歷.
原來藝術也可以傳福音.
再慢慢認識藝術更加闊.
我們透過藝術敬拜.
神幫助我作為一個讀藝術的人.
去整合藝術和信仰.
最近有一個新的事奉項目.
獲得浸大的老師邀請.

$^{481}$讓我在校園有一個藝術項目.
也是回應那位弟兄和我說.
Alison你在教會剛分享的所有東西.
其實都可以在我學校發生.
包括藝術產生對話.
有創意等等.
有一個團隊.
我們在這個時間的展覽舉行了.
大家可以有空間去深水Po 大南街Parallel Space.
大南街202號.
這個地方實體有一個關於栽種和平的展覽.
作為我的初心是藝術產生對話.
在對話中我們人與人生命相遇.
我作為有基督信仰的世界觀的人.
我可以透過創作和我的生命與你對話.
我們有交流.
無論你有信仰或沒有信仰.
然後大家去看一些生命的重點.
重溫一些生命已經有的元素.
或者可以探索的元素.
在這個展覽以栽種和平為題.
為何改名為「海的偶處」.
類似羅亞方舟.
羅亞在水未退不斷放鳥仔.
由烏鴉到百鴿.
以致牠知道洪水退去.
我們可以落地.
其實這個過程都是一個等待.
漫長的等待.
就正如我們面對現實.
很多不平安,戰爭,人與人的衝突.
甚至內心有很多糾結.
我們如何由國際到自己心中的層面.
我們如何面對,自主.
彷彿平靜的景況很遠.
我們想藉著創作這個展覽帶給大家看.
和平不單是很遠.
也可以很近.
是可以實踐的.
大家也可以共同經歷.

$^{521}$包括不平安的糾結和對盼望的等待.
所有參展的藝術家其實都是基督徒.
有少許展品給大家看.
如果你來到我會招呼你.
我們一起努力創作.
藝術好像很難.
日常生活亦很難.
甚至作為代人接物.
很多東西都可以嘗試.
我們再看看.
上帝因為愛而創造了人的配偶.
上帝說一個人不好.
要做一個配偶幫助他.
上帝做人,現在為人做配偶.
阿當也覺得.
他改甚麼都沒有配偶幫助他.
但配偶這個字不是等於低級.
而是大家同等的.
同等的形象,同等的能力.
彼此配搭.
男人做到的,女人做不到.
但男人做不到的,女人可以捕捉.
其實是互相捕捉的意思.
上帝做了一個外科手術家.
令他沉睡.
然後用一條條肋骨縫合起來.
然後做了一個女人.
帶到他那個人的面前.
然後這個時候.
阿當一見到夏娃就作了一首詩.
中文譯作「骨中骨,肉中肉」.
又押韻又靚.
還有他說.
你可以稱她為女人.
因為她是從男人而出的.
她說的那個語言.
其實男人和女人是押韻的.
在她說的那個語言.
我們聽聽這首歌.
請播放.

$^{561}$這就是張國榮的熱情的沙漠.
其實就是這樣.
他見到女人.
簡直是很棒.
很棒.
簡直是言語都表達不到.
她內裡龐大的愛意.
然後就作了一首詩.
要用藝術去表達她的那份愛意.
因為愛所以創作.
所以我們一起努力.
多點創作.
我們這篇道理.
基本上已經說完了.
因為愛去做人.
因為愛去為人做配偶.
在結構上.
最重要的是人的回應.
人如何用愛和創意去回應.
管理這個照明.
還有伸向眾多.
延伸到創作多一點的照明.
我們可以再留心.
我有一部壓軸篇.
和大家分享.
這個壓軸篇某程度是一個記號.
是一個有序.
但又有恩典覆蓋的記號.
在三章記載.
人做了一件事.
讓他開了眼.
發現原來我沒有穿衣服.
他做了一件神吩咐.
一定不可以做的事.
你做了會死的.
他就做了.
然後發現自己沒有穿衣服.
於是他們就創作.
他們用無花果葉.
去做了一件事.

$^{601}$他們用無花果葉.
去做了一群給自己.
很有趣.
上帝在園裡呼叫他們.
躲避上帝.
我們看到.
罪帶來梳理.
還有帶來蒙蔽.
他們失去了智慧.
因為無花果葉.
是有輕微的毒性.
在皮膚上會刺痛和癢.
即是說.
人做了一些.
有了罪性.
他的影響是.
蒙蔽了自己.
沒有足夠的智慧去分辨.
他仍然有神的形象.
但他有扭曲.
和有些地方蒙蔽了.
缺少了智慧.
但在.
三章二十一節.
聖經記載.
上帝為人做衣服.
他用獸皮.
去為人做衣服.
呂鎮忠.
亦用皮做了皮掛紙.
上帝繼承.
為創作藝術家.
外科醫生.
還有.
做一個很有智慧的時裝設計師.
他用獸皮.
去做衣服.
是否比無花果葉更耐用舒服.
其實除了耐用.
舒服之外.

$^{641}$還有一個很深層次的含義.
就是獸皮.
是取自動物的生命.
其實獸皮.
動物的死.
其實是代表了人的死.
動物可能不是一隻.
可以做到皮掛紙.
可能要幾隻動物.
但他們某程度充當了.
祭身的角色.
而我們知道.
在二千年前.
耶穌基督亦親自成為.
永遠的祭身.
贖罪祭.
去代替人的罪.
而成為二.
即是說.
人是有罪的.
但上帝再做一件創作.
去遮蓋我們的罪性.
讓我們成為二.
這是一個罪裡面.
有因點或因點蓋罪的記號.
即是說三.
三是一個因點蓋罪的記號.
創作和人生.
其實都是一個過程.
是一個不斷被上主根深.
不斷有它的因點蓋罪.
軟弱因點交織的旅程.
正如《耳粉所書》說.
我們是上帝手中的傑作.
《女真宗》亦稱.
《媳婦及翼本》亦稱傑作.
英文是masterpiece.
我們是上帝手中的傑作.
上帝不斷根深我們.
在耶穌基督裡重新創造我們.

$^{681}$讓我們能夠做.
祂早已安排的善事.
我相信這亦是.
剛才我們唱的那首《活出愛》.
我們如何活出愛.
需要不斷有因點.
亦不斷用創意和愛去履行.
願主祝福大家.
(影片完結).
\newpage



\section{馬可福音 15:33-39}
\label{sec:DcwFkXHzJa4}
\textbf{2022.04.15 《為甚麼離棄我》 馬可福音 15\_33-39 (和修版) 老耀雄牧師}
\newline
\newline
連結: \href{https://youtube.com/watch?v=DcwFkXHzJa4}{\texttt{https://youtube.com/watch?v=DcwFkXHzJa4}} ~~~~ 語音日期: 2022-04-18
\newline
\newline
\hyperref[sec:ra6IMFq2wMc]{\small{< < < PREV SERMON < < <}}
~
\hyperref[sec:index_chronic]{\small{[返順時目]}}
~
\hyperref[sec:index_scriptual]{\small{[返順卷目]}}
~
\hyperref[sec:lmsk0116lHM]{\small{> > > NEXT SERMON > > >}}
\newline
\newline
馬可福音 15:33-39
\newline
\begin{longtable}{cl}
\hline
\hline
章節 & 經文 (和合本修訂版)\\
\hline
15:33 & \begin{tabularx}{0.7\textwidth}{X} 到了正午，全地都黑暗了，直到下午三點鐘。 \end{tabularx} \\ \\ \relax
15:34 & \begin{tabularx}{0.7\textwidth}{X} 下午三點鐘的時候，耶穌大聲呼喊：「以羅伊！以羅伊！拉馬撒巴各大尼？」（翻出來就是：我的神！我的神！為甚麼離棄我？） \end{tabularx} \\ \\ \relax
15:35 & \begin{tabularx}{0.7\textwidth}{X} 旁邊站著的人，有的聽見就說：「看哪，他叫以利亞呢！」 \end{tabularx} \\ \\ \relax
15:36 & \begin{tabularx}{0.7\textwidth}{X} 有一個人跑去，把海綿蘸滿了醋，綁在蘆葦稈上，送給他喝，說：「且等著，看以利亞會不會來把他放下來。」 \end{tabularx} \\ \\ \relax
15:37 & \begin{tabularx}{0.7\textwidth}{X} 耶穌大喊一聲，氣就斷了。 \end{tabularx} \\ \\ \relax
15:38 & \begin{tabularx}{0.7\textwidth}{X} 殿的幔子從上到下裂為兩半。 \end{tabularx} \\ \\ \relax
15:39 & \begin{tabularx}{0.7\textwidth}{X} 對面站著的百夫長看見耶穌這樣斷氣，就說：「這人真是神的兒子！」 \end{tabularx} \\ \\
[1ex]
\hline
\hline
\end{longtable}
$^{1}$洪恩嘉弟兄姊妹平安.
如果我們有詳細看過四卷福音書的話.
我們就會發覺到.
耶穌在被釘上十字架之後.
一共說了七句話.
在路加福音記載了三句.
約翰福音又記載了三句.
而馬太和馬可又記載了一句.
所以如果我們綜合四卷福音書的內容.
我們就可以歸納出.
耶穌被釘十字架之後.
一共說了七句話.
我們稱為十加七言.
今天我就以馬可福音第十五章.
三十三至三十九節.
耶穌在十字架上的一段經文.
去反思我們和神的三個關係.
我們開始的時候有個祈禱.
「主耶穌,我求你在恩典.
臨到每一個屬你的兒女.
讓我們在受苦節的崇拜裡.
思考到你對我們所付出的.
是多麼何等的長闊高深.
好讓我們能夠傳一個感恩的心.
去回應你對我們的愛」.
我們祈禱是奉主明求.
阿們.
我們第一個關係就是.
無論我們處於甚麼環境.
甚麼狀態.
神仍然是我們的倚靠.
馬可福音第十五章三十三節記載道.
由正午十二點到下午三點.
全地都黑暗了.
一個違反了自然規律的現象.
維持了差不多三個小時.
一個應該是光明.
視野清晰的狀態.
現在反常了.
變得很暗昧.

$^{41}$視野模糊.
黑暗代表甚麼呢?.
黑暗可能代表一種.
黑暗的勢力.
不單單是指沒有方向.
沒有出路.
亦象徵邪惡勢力的籠罩.
撒旦正在橫行霸道的時候.
耶穌正在這個處境之下.
被釘在十字架上.
在黑暗勢力籠罩.
撒旦掌權.
一切顯得絕望的一刻.
耶穌在十字架上大聲呼喊.
耶穌的呼喊.
代表了耶穌心靈深處的一種感受.
這一種從心底裡發出的情感.
需要大聲呼喊出來.
根據一個猶太的傳統.
祈禱有十個名稱.
排第一位的名稱.
就是呼喊.
當我們看不到跡象.
顯示我們和神站在同一個陣線.
當我們覺得神沉默不語.
當我們傷心欲絕.
當我們感覺到自己完全被遺棄的時候.
失敗的感覺.
可能誘惑我們放棄對神的那種倚靠.
我們應該怎樣做呢?.
即使被最大的苦難所淹沒.
即使覺得神不在我們身邊.
我們仍然要堅持抓住神.
我們需要向神大聲呼喊.
耶穌在十字架上的呼喊.
正正顯示一種仍然最天賦那種倚靠的心.
其實昔日很多屬靈的領袖.
同樣有這種體驗.
摩西大令以色列人在曠野.
遇到沒有水沒有食物.

$^{81}$最困難的時候.
他呼喊過.
約書亞大令以色列人.
面對阿摩尼人迦南人攻擊的時候.
他呼喊過.
約伯耶利米在他們面對絕境的時候.
同樣呼喊過.
當我們好像耶穌.
面對撒旦的張牙舞爪.
當我們好像摩西.
面對沒有水沒有食物的逆境.
當我們好像約書亞.
面對強大敵人的挑戰和攻擊.
當我們好像約伯.
面對眾叛妻離的那種孤獨.
甚至當我們好像耶利米.
面對國破家亡.
感到軟弱無力的時候.
每個人都需要向神呼喊.
因為祂是我們唯一的倚靠.
向神呼喊並不是一件羞恥.
不見得剛的事.
相反.
呼喊是我們承認我們軟弱無力.
我們需要依靠神的一種表現.
能夠向神呼喊.
表示我們生命裡面仍然有神.
我們不要害怕去向神呼喊.
因為呼喊反映我們不是依靠自己.
而是依靠神.
呼喊也可以說是信心的一種表現.
在今天早上受苦節的崇拜當中.
我們思考一下.
我們上一次呼喊神究竟是何時.
抑或在我們的信仰生命裡面.
從來都沒有向神呼喊過.
神從來都沒有介入過在我們的生命裡面.
我們有沒有這些體驗呢?.
是因為我們沒有遇過一些苦難困境的經歷.
如果是.

$^{121}$我們要為你感恩.
因為生命沒有困境並不是必然的.
不過生命如果只是有順境.
而沒有經歷過逆境的試煉.
生命就只有一種顏色.
而不是色彩繽紛.
弟兄姊妹請留意.
我不是鼓勵大家去追求苦難.
而是想請大家不需要刻意去逃避苦難.
因為苦難逆境都是在神的心意當中.
苦難逆境可以為我們生命添上另一種色彩.
一種與神同行.
與神同在的色彩.
生命不是要展現個人的成功.
生命是要展現出.
神如何在我們生命裡活著.
我們願意向神呼喊.
表示願意讓神介入在我們的生命裡.
生命中如果充滿了與神同行的色彩.
我們就越能夠被神所使用.
第二種的信仰關係就是.
主耶穌就是我們的拯救.
在第三十四節記載了耶穌呼喊的內容.
「以羅伊,以羅伊,拉瑪撒巴,葛大利」.
這句話是亞蘭文的原文.
它其實譯出來就是.
「我的神,我的神,為什麼離棄我?」.
呼喊的內容引用詩篇第二十二篇的第一節.
詩人向神呼喊.
「我的神,我的神,為什麼離棄我?」.
「為什麼遠離不救我,不聽我的呻吟?」.
詩人開始的時候相當絕望.
他認為神離棄了他.
遠離他,不聽他的呼喊.
詩人身處於絕望,無助,孤單的境地.
耶穌表明他現在的處境.
與詩人的處境一模一樣.
不過詩篇二十二篇的終結是這麼說的.
「驚慰耶和華的人要讚美他.
雅各的後裔要榮耀他.

$^{161}$以色列的後裔要懼怕他.
因為他沒有藐視厭惡受苦的人.
也沒有轉念不顧他們.
那受苦之人呼求的時候.
他就垂聽」.
詩人的呼喊最終都得到神的垂聽和安慰.
神是那個垂聽呼喊的神.
詩人亦以讚美神作為結束.
有人說過.
人的盡頭就是神的開始.
詩篇二十二篇印證了這個道理.
由於亞蘭文爾羅伊的發音.
與以利亞很相似.
所以我們能夠解釋到.
在第三十五節.
為何士兵會說.
「看那!他叫以利亞」?.
士兵以為耶穌呼叫以利亞先知來救他.
於是士兵再嘲諷耶穌說.
「且等看!看以利亞會不會來把他放下來?」.
原來士兵聽錯了耶穌呼喊的內容.
其實耶穌呼喊的對象不是先知以利亞.
耶穌呼喊的對象是我的神.
我的神是指天父.
是那位創天造地的耶和華.
是那位與聖子,聖靈同為一體的聖父.
是那位三位一體的聖父.
按著三位一體.
從第二位位格的聖子的角度去看.
聖父從來沒有離開過他.
因為三位一體的神是不可能分開的.
從昔日,今日到將來.
聖父,聖子,聖靈.
三位一體的神從來沒有分開過.
不過耶穌固然是一百分之一的神.
但他也是一百分之一的人.
從肉身耶穌的角度去大聲呼喊.
我的神.
為什麼離棄我?.
因為祂知道即將要面對死亡.

$^{201}$祂要代替人類走上這條待贖之路.
由於人類所犯的罪.
需要人類自己去背負.
但是人類充滿了罪.
怎麼能夠去背負呢?.
因此耶穌就以一個沒有罪的人生.
去代替我們這個有罪的人生.
耶穌以人的身份和代表.
背負整個人類的罪責.
耶穌為所有犯罪得罪神的人類.
以自己的身體獻上贖罪祭.
這不僅是天父的心意.
這也是三位一體.
聖父,聖子,聖靈的心意.
也是道成肉身的耶穌的心意.
神對人的心意就是要施行拯救.
讓我們可以回轉.
神對耶穌的心意.
就是要祂作為拯救人類的中寶.
藉著耶穌的死.
我們可以與神復和.
藉著耶穌的死.
我們可以罪得赦免.
藉著耶穌的死和復活.
我們可以進入神的國度.
同樣的大兄姐妹.
我們所犯的罪.
理應由我們自己去承擔.
不過很可惜.
就算我們肯承擔.
也承擔不了.
一個污穢的生命.
不可以在審判台前.
有任何脫罪的機會.
對於耶穌.
祂不需要付出這麼高的代價.
耶穌曾經向天父求過.
求你使者杯離開我.
不過祂甘願順服天父的旨意.
所以祂繼續求.

$^{241}$然而不是照我所願的.
而是照你所願的.
耶穌是按照天父.
拯救世人的心意.
而走上十字架.
今天我們能夠成為神的兒女.
天國的子民.
完全是藉著耶穌.
在十字架上為我們所成就.
詩篇二十二篇.
以讚美神作為結束.
我們也應該要讚美神.
同英家的弟兄姊妹.
耶穌走上十字架的過程.
被鞭打,被唾罵,被凌辱,被冤枉.
耶穌一生所受的苦.
全都是為了我們每一個罪人.
也是為我們每一個願意跟隨主的信徒.
我們要怎樣去回應.
耶穌為我們所做的一切呢?.
耶穌願意順服天父的心意.
甘心樂意放下自己的主權.
對於今天高舉個人自由的我們.
我們會不會向神祈禱.
然而不是照我們所願的.
而是照你所願的呢?.
耶穌向我們展現出.
信仰是一個要付代價的旅程.
雖然要受苦.
但最後能夠成就神的心意.
我們在教會事奉.
有沒有一種為主受苦的心?.
我們的事奉是要別人遷就自己.
抑或我們願意犧牲自己去苦咒別人呢?.
耶穌拯救我們.
不是叫我們去過一個安逸的信仰生活.
而不是叫我們享受一個.
不需要付代價的信仰人生.
不嫁恩典.
絕對不是主耶穌呼召我們.

$^{281}$要去走的成聖之旅.
我講的第三個信仰關係就是.
一個與神靈距離的關係.
一個與人合而為一的關係.
耶穌為世人.
為天父所做的.
在第38節已經說了.
經文說.
「殿的萬子從上到下列為兩半」.
「至聖所的萬子象徵神與人的左隔」.
「大祭司要進入至聖所朝見神」.
「必定要揭開萬子才進入得到」.
「亦唯有大祭司才有這個資格」.
這個萬子的左隔.
不單單是神與人在距離上的左隔.
亦不單單是神與人在空間上的左隔.
亦都是神與人在關係上的左隔.
現在萬子裂開了.
一切的左隔都沒有了.
神與人的距離沒有了.
我們可以藉著耶穌與神溝通.
神與人的空間沒有了.
主耶穌可以與我們同在.
《約翰福音》十五章說.
「常在我裡面的,我也常在他裡面」.
神與人的關係恢復了.
我們能夠被稱為神的兒女.
遷到神的角度之內.
洪顯嘉弟兄姊妹.
昔日我們沒有可能與主相遇.
昔日我們與神的關係破裂.
昔日我們連朝見神的資格都沒有.
但是今天我們與神零距離.
今天我們可以與主同在.
今天我們可以透過敬拜朝見神.
不過我們有沒有珍惜這一切的恩典呢?.
最後三十九節說.
「一個站著的白膚長看見耶穌這樣斷氣.
就說:這人真是神的兒子」.
對於耶穌的身份.

$^{321}$彼得曾經宣認過.
彼得在《馬太福音》十六章.
曾經宣認過耶穌的身份.
他說:你是基督.
是永生神的兒子.
對於猶太人來說.
耶穌就是那位要來的彌賽亞.
是耶和華對整個猶太民族的應許.
耶穌基督這個神的兒子.
耶和華所應許的彌賽亞.
今天要來拯救整個以色列的民族.
不過現在宣認耶穌的身份的是一個白膚長.
嚴格來說是一個羅馬人.
一個與神為敵的愛邦人.
一個執行羅馬暴政的異族人.
因著他宣認和相信.
耶穌就是神的兒子.
因著福音的大能.
也因著神對世人的恩典.
因此不單止猶太人得到拯救的恩典.
一切凡相信耶穌是神的兒子的人.
都可以領受這個恩典.
聖經在《皆太書》三章十六節這樣說.
不再分猶太人或希臘人.
不再分遺奴的自主的.
不再分男的女的.
因為你們在基督耶穌裡都成為一了.
就在神的恩典和應許.
今天我和大家可以在這裡敬拜神.
讚美神.
因為凡一切相信耶穌是神的兒子的人.
都在基督裡合而為一了.
我們在這裡可以互相稱呼為弟兄姊妹.
同樣的弟兄姊妹.
今早的受苦節崇拜.
我希望大家透過上述的經文.
讓我們明白.
苦難是有價值的.
即是耶穌在苦難當中.
仍然可以呼喊神.

$^{361}$我們願意呼喊.
就表示神可以介入在我們生命裡.
讓每一個屬天父的兒女.
都可以無時無刻地呼喊神.
耶穌的救贖不是廉價的.
耶穌是付上整個生命來救贖我們的.
祂一信服.
成就天父對我們救贖的旨意.
我們願意信服.
亦能夠成就主耶穌在我們生命的旨意裡.
讓每一個屬主的都甘心樂意.
信服主耶穌對我們的呼喚.
我們和耶穌的關係.
是可以完全沒有阻隔的.
唯一阻隔我們的原因.
是我們不願意和耶穌建立關係.
甚至我們每一個信徒都能夠體驗.
耶穌對我們的應許.
就是我在你裡面.
你在我裡面.
與主同在的今天.
我們每一個屬主耶穌的信徒.
都需要合一.
因為唯有大家有合一的心.
聖靈才在我們心裡面工作.
甚至每一個主耶穌基督的跟隨者.
我們都能夠敏銳聖靈對我們提醒.
督責,引導.
好叫我們的生命與恩典相承.
最後我邀請大家一起站起來.
我們一起祈禱.
今天主耶穌已經去過我們.
無論我們在哪裡.
你都與我們同在.
因此縱然我們感覺不到你在我們身邊.
縱然我們覺得你不回應我們的祈禱.
縱然我們向你呼喊.
你為什麼要離棄我.
但是我們仍然都要歌頌你.
要讚美你.

$^{401}$我們要因你的名歡喜.
因你的名快樂.
因為人的感覺.
不會高得過你的應許.
你的應許直到地極.
你的應許直到窮滄.
你的應許直到永恆.
既然有你信實的應許.
我們就不需要畏懼了.
一切全是你的恩典.
願我們每一個屬於你的.
都能夠合一地去事奉你.
讓你的名被高舉.
讓你的名響遍全地.
我們誠心獻上禱告.
奉主耶穌基督的聖名祈求.
阿們.
會眾請坐.
願主祝福大家.
\newpage



\section{約翰福音 20:24-29}
\label{sec:lmsk0116lHM}
\textbf{2022.04.17 《信心的轉化》 約翰福音 20\_24-29 (和修版) 老耀雄牧師}
\newline
\newline
連結: \href{https://youtube.com/watch?v=lmsk0116lHM}{\texttt{https://youtube.com/watch?v=lmsk0116lHM}} ~~~~ 語音日期: 2022-04-18
\newline
\newline
\hyperref[sec:DcwFkXHzJa4]{\small{< < < PREV SERMON < < <}}
~
\hyperref[sec:index_chronic]{\small{[返順時目]}}
~
\hyperref[sec:index_scriptual]{\small{[返順卷目]}}
~
\hyperref[sec:O3rEhqjezKM]{\small{> > > NEXT SERMON > > >}}
\newline
\newline
約翰福音 20:24-29
\newline
\begin{longtable}{cl}
\hline
\hline
章節 & 經文 (和合本修訂版)\\
\hline
20:24 & \begin{tabularx}{0.7\textwidth}{X} 那十二使徒中，有個叫低土馬的多馬，耶穌來的時候，他沒有和他們在一起。 \end{tabularx} \\ \\ \relax
20:25 & \begin{tabularx}{0.7\textwidth}{X} 其他的門徒就對他說：「我們已經看見主了。」多馬卻對他們說：「除非我看見他手上的釘痕，用我的指頭探入那釘痕，用我的手探入他的肋旁，我絕不信。」 \end{tabularx} \\ \\ \relax
20:26 & \begin{tabularx}{0.7\textwidth}{X} 過了八日，門徒又在屋裡，多馬也和他們在一起。門都關了，耶穌來，站在當中，說：「願你們平安！」 \end{tabularx} \\ \\ \relax
20:27 & \begin{tabularx}{0.7\textwidth}{X} 然後他對多馬說：「把你的指頭伸到這裡來，看看我的手；把你的手伸過來，探入我的肋旁。不要疑惑，總要信！」 \end{tabularx} \\ \\ \relax
20:28 & \begin{tabularx}{0.7\textwidth}{X} 多馬回答，對他說：「我的主！我的神！」 \end{tabularx} \\ \\ \relax
20:29 & \begin{tabularx}{0.7\textwidth}{X} 耶穌對他說：「你因為看見了我才信嗎？那沒有看見卻信的有福了。」 \end{tabularx} \\ \\
[1ex]
\hline
\hline
\end{longtable}
$^{1}$足夠法定人數.
主席.
香港復活節假期.
香港大小小的商場都會佈置得美輪美奐.
而且通常都會安排一隻復活套或吉祥物.
在大派禮物.
他們好像聖誕老人一樣.
在聖誕節派禮物一樣.
將歡樂的氣氛帶給市民.
商場將復活節的主角.
變成了復活節復活蛋和復活套.
不過以往在香港.
人在復活節所關心的.
其實是在哪裡的復活節自助餐最抵食.
或者復活節假期去哪裡玩.
似乎復活節的意義.
就是一個提供給人消費的日子.
香港人似乎忘記了復活節的真正意義.
今年的復活節.
再沒有這種歡樂的氣氛.
大家都受到疫情的影響.
盡量減低外出活動.
避免感染.
疫情由2020年開始到現在.
已經超過17000人.
死於新冠狀病毒.
香港人對於何時可以回復.
以往有序的生活.
沒有人可以回答.
世人或不明白.
或不清楚復活節的意義.
但作為信徒的我們.
我們應該知道.
復活節是一個關乎信仰的日子.
主角一定是耶穌.
而不是代表消費的復活套.
因此我們今天.
用《約翰福音》第20章24至31節.
去探討一下耶穌的復活.
帶給我們信心的三個轉化階段.

$^{41}$我們一起去祈禱.
因著你的生與死.
讓每一個屬你的.
都能享受你賜予的恩典.
甚至我們每一個.
都在你的恩典之下.
過一個豐盛的人生.
我們祈禱奉教主耶穌.
得勝名字祈求.
阿們.
開始講道之先.
我先交代一下經文的背景.
約翰在20至21章.
用了五個故事.
去描述耶穌復活的事件.
第一個故事是.
彼得和約翰去到空墳墓.
他們沒有看到復活的耶穌.
只看到細麻布和果頭根.
不過他們就相信了.
第二個故事是.
瑪利亞看到復活的耶穌.
耶穌不讓瑪利亞摸他.
理由是耶穌還沒有升上去.
見他的父.
第三個故事是.
耶穌在門徒之間顯現.
除了差遣門徒之外.
還被他們領受了聖靈.
第四個故事就是今天講道的經文.
第五個故事是.
耶穌在提比利亞海邊顯現.
還三次問彼得.
你愛我嗎?.
其實第三個和第四個故事的場景是相同的.
都是在門徒聚集的屋裡.
不過人物的數量不同.
第一次耶穌顯現的時候.
其他門徒都在場.
唯獨多瑪不在.

$^{81}$第二次耶穌顯現的時候.
多瑪就在其中.
所以耶穌第二次在門徒中間顯現的目的.
就是為了多瑪.
作者透過多瑪的經歷.
讓我們知道信心的三個轉化的階段.
第一個階段.
見不到耶穌的信心.
是沒有永恆價值的.
經文在第24節.
講明了耶穌第一次向門徒顯現的時候.
多瑪不在.
第25節裡面.
門徒向多瑪說.
我們已經看見主了.
門徒這句說話表示了一個宣認.
就是門徒見證了耶穌已經復活.
而耶穌已經復活這件事.
代表了耶穌之前所說的所有說話都是真的.
可信的.
真的會應驗的.
耶穌在第16節曾經預言自己的死和復活.
所以如果耶穌已經復活.
是被證實的話.
那表示耶穌之前所說的所有應許.
承諾都同樣會應驗和實現.
約翰就在20章31節講明.
他說要藉著耶穌已經復活這個事實.
是要叫你們信耶穌基督是神的兒子.
並且叫你們信了祂.
就可以因祂的名得生命.
相信耶穌是神的兒子和得生命.
是約翰寫福音書的其中一個目的.
而得生命就是相信耶穌是神.
的一個具有永恆價值的應許.
簡單一點說.
相信耶穌已經復活這個事實.
就會相信耶穌是神.
而相信耶穌是神.
就可以因耶穌的應許而得生命.

$^{121}$不過多瑪對於耶穌已經復活這個事實的反應就是.
不信.
其實多瑪的反應是很矛盾的.
一方面見證耶穌已經復活這個事.
是多瑪所認識的門徒.
多瑪沒有理由認為門徒會說謊.
他理應相信門徒真的見到耶穌已經復活這個事實.
但是另一方面.
耶穌的復活是超出了多瑪的想像.
也都是超出了多瑪的經驗.
所以多瑪的信心不能單靠別人的說話而確立得到.
縱然所說的人是誠實可靠的.
為了滿足這個超越想像和經驗的信心.
多瑪要求的.
是不單要看見復活的耶穌.
他還要親自觸摸耶穌手上的釘痕和他的肋行.
然後才會再相信.
多瑪要親自觸摸耶穌還有一個意義.
就是懷疑門徒所見到的耶穌究竟是否真實.
是否一種日有所思夜有所念的幻象呢?.
門徒所見到的其實不是一個真實的形體.
只是一個幻覺.
所以如果多瑪能夠親自觸摸到這個形體.
就能夠證實門徒所見到的不是一個幻象.
也都證實耶穌是真的復活了.
在這個階段.
多瑪對耶穌的信心.
並不能夠依靠第三者的說話和見證.
而是要依靠自己親身的經歷.
才可以建立得到.
多瑪信心的狀況是什麼呢?.
是要見不到,摸不到,就不相信.
換另一個說法.
多瑪信心是需要什麼呢?.
需要看到並且需要驗證過.
才能夠建立信心.
對於多瑪來說.
這一刻他不相信耶穌已經復活.
也就是說.
他對耶穌之前的承諾,應許,盼望.

$^{161}$同樣都是沒有信心的.
多瑪抱著一個懷疑的態度.
可以這麼說.
多瑪這一刻對耶穌的信心.
是沒有價值,沒有意義的.
更加說不上什麼永恆的價值.
本來曾經說過.
如果基督沒有復活的話.
我們的信就是謊言.
同樣如果我們今天不相信.
耶穌已經復活這個事實.
我們的信都是謊言.
我記得以前和弟兄姊妹.
去傳福音的一個片段.
有一次我們應邀向一位弟兄傳福音.
去到的時候大家一起分享.
這位男士一直都沒有什麼反應.
不過當我們一說到.
耶穌就是神的時候.
他就馬上答嘴.
他說如果耶穌是神的話.
那不如出現在我面前.
如果祂真的出現在我面前的話.
我就相信了.
這位男士好像多瑪一樣.
他不相信別人的見證.
只是相信自己眼睛所看到的東西.
和自己所經驗到的東西.
他一天見不到耶穌.
他總是不相信.
我們很多時候都會說.
這些人很硬心.
需要神慢慢軟化他們.
一個人如果不相信耶穌已經復活.
自然不會相信耶穌就是神.
更加不會相信.
耶穌對人類的應許是會實現的.
所以這些人的信是沒有盼望的.
也沒有永恆價值.
那怎樣的信心才有永恆價值呢?.

$^{201}$第二個階段.
看到耶穌的信心是有盼望的.
對於多瑪要見到耶穌.
並且要親自觸摸到耶穌的要求.
耶穌怎樣去回應呢?.
如果根據《馬可福音》的記載.
耶穌向梅底拉的瑪利亞顯現的時候.
其實耶穌吩咐瑪利亞向門徒和彼得說.
我會先去到加利利.
叫門徒往加利利去見他.
所以如果門徒相信瑪利亞的話.
他們應該立刻趕去加利利.
而不是留在耶路撒冷.
不過門徒的信心實在太脆弱了.
以至於耶穌要在耶路撒冷多次顯現.
才能夠建立門徒的信心.
經文在二十六節透露.
「過了百日,耶穌來站在當中說:.
「願你們平安!」.
耶穌願意回應多瑪的要求」.
在這八天裡.
耶穌再次在眾門徒當中顯現.
很明顯耶穌這次顯現的目的.
就是要鞏固多瑪對耶穌已經復活的信心.
多瑪在這八天裡.
究竟想了些什麼呢?.
他可能正在經歷一個信心的危機.
他沒有把握耶穌會回應他的要求.
他只能夠等待.
他很渴望像其他門徒一樣.
可以見到耶穌.
但是如果耶穌不再顯現呢?.
那他會怎樣呢?.
可能他整個生命都會改寫.
或者他的信心會崩潰.
所以在二十七節.
耶穌介入了多瑪的生命裡.
耶穌不單止讓多瑪見到.
他還邀請多瑪觸摸他的手和肋行.
這兩個部位正正就是多瑪在二十五節所提出的.

$^{241}$他很想觸摸的.
耶穌這個邀請不是要挑戰多瑪.
而是想圓滿多瑪對信心建立的要求.
多瑪的要求就是.
要看見.
並且需要經過驗證才會相信.
耶穌就滿足多瑪的要求.
耶穌沒有輕看多瑪的信心.
耶穌是憐憫多瑪.
耶穌在多瑪還沒有行動之前.
加多一句勸勉對他說.
「不要疑惑,總要信」.
疑惑這個字和信這個字同樣都是一個形容詞.
呂鎮忠譯本譯作.
「別無信心了,務要有信心」.
司高譯本譯為.
「不要作無信的人,但要作有信得的人」.
原文可以直譯為.
「不要變得無信,要信」.
這是由信心從無到有的轉化.
神的介入本來是要滿足多瑪兩個要求.
就是看到,摸到.
不過這一刻多瑪已經覺得足夠了.
透過和耶穌的對話.
也透過耶穌對他的勸勉.
多瑪的信心已經被建立了.
他不需要真真正正地捉住耶穌的手和立行才相信.
耶穌在他面前的顯現已經足夠了.
多瑪在二十八節說出.
「我的主,我的神」.
是多瑪從心底裡相信耶穌的信仰表白.
主和神是對耶穌的宣認.
主是多瑪對耶穌的關係性宣認.
神是多瑪對耶穌神性身份的宣認.
除了表示多瑪相信耶穌已經復活之外.
也承認耶穌是神的身份.
是對耶穌的神性最清晰和最簡單的宣認.
既然多瑪承認了耶穌已經復活這個事實.
更加宣告耶穌是神的身份.
自然就建立了耶穌是神的信心.

$^{281}$也肯定了耶穌的應許必定會實現.
這個「應許」就是三十一節所說的.
「叫你們信了他,就可以因他的名得生命」.
所以多瑪這一刻的信心是充滿盼望的.
而這個盼望也是充滿了永恆價值的.
神憐憫多瑪.
並且直接介入在多瑪的生命裡.
幫助他信心的建立.
以至多瑪對耶穌已經復活這個事實.
又存而轉做徹底的相信.
多瑪的事件讓我們知道.
我們對神的信心的的確確需要耶穌的介入.
才可以建立得到.
弟兄姊妹.
同樣,如果我們今天對神失去了信心.
或者我們對神的「應許」失去了盼望.
或者我們正經歷一個信心危機的時候.
我們需要學多瑪一樣.
祈求神的幫助.
只要耶穌介入在我們的生命裡.
我們的生命就必定能夠由懷疑.
沒有信心.
轉化為一個肯定和有盼望.
第三個階段就是.
看到耶穌在人生命中的信心是超越的.
我們知道耶穌沒有否定多瑪的信心.
不過這種信心是建基於多瑪在物質上.
世界上見到耶穌.
又能夠摸到耶穌.
如果耶穌不介入.
多瑪的信心是不會被建立的.
多瑪固然是被耶穌憐憫和祝福.
不過,如何滿足每一個.
都想見到耶穌已經復活的要求呢?.
是不是每一個想的話.
耶穌都在那個人面前顯現.
滿足每一個的要求呢?.
耶穌既然可以憐憫多瑪.
我相信耶穌也可以憐憫其他.
很想見到耶穌已經復活的人.

$^{321}$耶穌和多瑪的對話沒有停止.
耶穌繼續說出信心建立.
是可以超越到另一個視野.
他說:你因看見了我才信.
那沒有看見就信的有福了.
耶穌這句話說明了有另一種信心.
而這種信心是多瑪和其他門徒所沒有的.
耶穌不是想將兩種信心作出比較.
去分辨哪一種信心比較好.
而是想說明兩種信心建立的方法是有不同.
一種是見到肉身的耶穌的信心.
而另一種是無需要見到.
那個肉身耶穌就有相信的信心.
耶穌所宣告的超越信心是怎樣的呢?.
是賜給那些就算見不到.
耶穌肉身的顯現.
但仍然相信耶穌已經復活的人.
在這裡其實耶穌是指向多瑪.
和使徒圈子以外的信徒.
也就是指向教會的世界.
那些透過使徒的見證.
而相信耶穌的人.
他們無需要親眼見過肉身復活的耶穌.
而是透過門徒的生命裡面.
能夠見到耶穌.
從多瑪的事件我們知道.
耶穌已經復活的信心.
是有兩個條件的.
就是耶穌的介入.
以及耶穌的被看見.
超越的信心仍然需要這兩個條件.
耶穌的介入是必須的.
但耶穌的被看見.
就不再需要親眼見到耶穌肉身的復活.
而是只需要看到耶穌.
在信徒生命中的顯現.
就是耶穌在物質世界中的被看見.
轉化成為耶穌在別人生命中的被看見.
這種被看見的轉化.
也為我們建立了.

$^{361}$耶穌已經復活的信心.
這種超越信心.
就是讓那些能夠在物質性的看見.
轉化成在屬靈上看見的人.
他們能夠在別人的生命裡面見到耶穌.
同時生命經過轉化.
以及被耶穌介入.
其他人又可以從他的生命裡面.
見到復活的耶穌.
復活的耶穌.
可以在每一個信徒的生命裡面.
被人見得到.
耶穌說這種超越的信心是有福的.
有福起碼有幾種屬靈的含義.
第一.
有福是不需要等待的.
有福是不可以馬上見到耶穌.
我們如果回望一下多瑪.
他在等待耶穌顯現的第八天.
才能見到.
他是多麼的焦急.
如果他有這份超越的視野.
他就不需要等這一百天.
第二.
有福是確定基督徒生命之中有耶穌.
因為生命要被耶穌介入.
我們的生命才能彰顯到耶穌.
然後才能被人見得到.
第三.
有福是因為我們能夠成為恩典的其中一個管道.
透過耶穌介入信徒的生命.
信徒生命活出耶穌.
而令第三者看得見.
並且相信耶穌.
將神的祝福賜予別人.
我們成為耶穌已經復活的中介和見證者.
剛才我提到傳福音的例子.
還未說完.
那位弟兄是一個會友的丈夫.
姐妹在結婚之後才信主.

$^{401}$他透露他老公很暴躁.
一喝酒就打人.
所以兩夫妻經常打架.
打得連隔壁的人都聽到.
後來姐妹信了主.
開始回主一學.
主一學老師教她謙卑.
不要跟老公吵.
凡事要忍耐一下.
姐妹覺得不太合理.
她認為又不是自己錯.
為什麼不出聲呢?.
又不還手?.
這樣就很吃虧了.
不過姐妹最後都願意信服.
她說那我就不還手了.
如果真的打得我厲害.
我就走開避開.
所以姐妹很想有教會的人.
向她老公傳福音.
希望她老公信了主.
改變一下她的性格.
不過那次她老公根本聽不進去.
有一次她老公又喝醉酒.
不過這次沒有出手打人.
如果在以前.
姐妹會怎樣做呢?.
拿一盆冷水兜頭淋過去.
等他醒過來.
不過這次姐妹沒有這樣做.
還用毛巾幫他敷面.
弄杯深茶給他喝.
她老公迷迷糊糊.
被她照顧了一晚.
但老公醒過來.
他說昨天好像有個人服侍我.
好像是天使一樣.
他問天使是誰呢?.
姐妹說哪有天使呢?.
除了我還有誰肯服侍你呢?.

$^{441}$老公開始有點不好意思.
問老婆為何轉性呢?.
為何肯照顧我呢?.
姐妹說如果不是信了耶穌.
我真是懶得理你.
她老公雖然沒有表示.
但開始對老婆關心了.
還說上次得罪了教會的人.
不好意思.
很主動地想約我們去吃飯.
說是道歉.
我們又再上去.
這次氣氛很好.
我們再次問他信不信耶穌.
他這次說.
好呀!試試吧.
我問他為何這次會信呢?.
他說自從老婆回到教會後.
真的看到他改變了很多.
他真的覺得.
耶穌在他老婆的生命裡.
他也想試試被耶穌改變一下.
他老公是在老婆身上.
看到神的工作.
這個看見吸引了老公去認識耶穌.
我們的生命能否讓人看見耶穌.
在我們生命裡活著呢?.
如果我們的生命能夠讓人看到.
耶穌在我們生命裡活著.
我們就成為一個見證耶穌的流通管子.
透過耶穌向多瑪顯現的八折經文.
我們知道多瑪初時的信心.
是沒有意義的.
沒有永恆價值.
這個永恆價值就是.
多瑪不能因著信耶穌而得生命.
後來經過耶穌的介入.
多瑪看到耶穌.
有盼望的信心被建立起來.
不過耶穌更加宣告一種超越的信心.

$^{481}$這種信心只需要看到別人生命的耶穌.
就可以建立了.
這種信心是要在別人的生命裡見到耶穌.
同樣可以建立得到.
同樣擁有這種信心的人.
也可以得到神所應許的永恆生命.
弟兄姊妹.
今天是復活節的主日.
我們反思一下.
耶穌已經復活了.
我們相信耶穌是神.
也相信耶穌真的復活了.
我們已經有一個永恆的價值.
我們的信有永恆的價值.
不過我們能不能成為那個流通的管子呢?.
我們能不能讓人在我們生命裡見到耶穌呢?.
我們生命可不可以成為別人的祝福呢?.
我們洪恩家很多時候都說.
我們要服侍洪水橋的街坊.
當我們說的時候.
我們能不能反思一下.
我們可以為這批街坊做些什麼呢?.
當我們說洪恩家很想去傳福音的時候.
很想多人歸信主耶穌的時候.
我們能不能反思一下.
我們為這批未得之民可以做些什麼呢?.
我相信信仰不是一個口說的空話.
信仰需要我們實踐.
在實踐的過程裡.
我們願意按神的話語去謙卑自己而被主使用.
如果我們願意的話.
我鼓勵大家實踐今年教會的年提.
「為身事奉」.
參與教會不同的事奉崗位.
讓你所接觸到的人.
能夠在你的生命裡見到耶穌.
讓每一個屬主的生命.
能夠顯現出耶穌的屬性.
成為更多人的祝福.
我們一起有個祈禱.

$^{521}$你介入在我們的生命裡.
使我們認識你.
歸向你.
求你使我們對你的信心每天加添.
從肉身的看見.
能夠轉化為屬靈上的看見.
不單在別人生命裡見到你.
也使我們的生命.
讓別人能夠見到你.
使我們的生命能夠成為一個流通的管子.
讓更多人認識你.
學悟你.
既然我們的生命是屬你的.
就讓我們甘心樂意地被你使用.
好叫你的名被高舉.
榮耀指屬利.
我們的祈禱仰望.
是奉教主耶穌基督得勝名至其球.
阿們.
願主祝福大家.
拜拜.
\newpage



\section{彼得前書 1:1-3}
\label{sec:O3rEhqjezKM}
\textbf{2022.04.24 《主復活的盼望》 彼得前書1\_1-3 (和修版) 謝靄薇牧師}
\newline
\newline
連結: \href{https://youtube.com/watch?v=O3rEhqjezKM}{\texttt{https://youtube.com/watch?v=O3rEhqjezKM}} ~~~~ 語音日期: 2022-04-25
\newline
\newline
\hyperref[sec:lmsk0116lHM]{\small{< < < PREV SERMON < < <}}
~
\hyperref[sec:index_chronic]{\small{[返順時目]}}
~
\hyperref[sec:index_scriptual]{\small{[返順卷目]}}
~
\hyperref[sec:D5wSPpsbnpU]{\small{> > > NEXT SERMON > > >}}
\newline
\newline
彼得前書 1:1-3
\newline
\begin{longtable}{cl}
\hline
\hline
章節 & 經文 (和合本修訂版)\\
\hline
1:1 & \begin{tabularx}{0.7\textwidth}{X} 耶穌基督的使徒彼得寫信給那些被揀選，分散在本都、加拉太、加帕多家、亞細亞、庇推尼寄居的人， \end{tabularx} \\ \\ \relax
1:2 & \begin{tabularx}{0.7\textwidth}{X} 就是照父神的預知，藉著聖靈得以成聖，以致順服耶穌基督，又蒙他血所灑的人。願恩惠、平安多多地賜給你們！ \end{tabularx} \\ \\ \relax
1:3 & \begin{tabularx}{0.7\textwidth}{X} 願頌讚歸於我們主耶穌基督的父神！他曾照自己的大憐憫，藉著耶穌基督從死人中復活，重生了我們，使我們有活的盼望， \end{tabularx} \\ \\
[1ex]
\hline
\hline
\end{longtable}
$^{1}$各位大家卿卿 平安.
今天很開心.
放鬆.
聚集的時候.
第一次見到你們.
很開心 盼望神.
在當中帶動我們的心.
更加活躍.
像以前一樣更加愛主.
愛神 愛人.
我們一起祈禱.
當我們來到你殿中的時候.
我們心中真的歡喜快樂.
願你真是你的靈魂與我們同在.
也願我們心中有什麼.
不滿意的.
求你來開.
以致我們能夠專一地.
去仰望你.
仰望你在當中對我們說話.
以致我們心裡.
能夠有你給我們的.
力量 我們如何去奔跑.
前面的路 也願孩子.
口中的言語.
心中的意念滿意如納.
我們這樣祈禱交託.
你所記得名頭 阿們.
(回應觀眾).
當我們在疫情之下.
不知道有什麼盼望呢?.
當然學生就開心.
希望他們能夠.
快點見到同學.
可以實體地去學習.
不容易的.
就算我們老人家.
你常常叫他們在家裡.
困在家裡.
也是很悶的.

$^{41}$他們也要出去走走.
不然會有憂鬱的.
有些人.
可能很想.
快點結婚.
但總是.
限制兩個人.
很慘的.
不可以見到更多的朋友.
所以在這些情形之下.
有很多很多的盼望.
但耶穌死後.
祂曾經預言.
祂第三天復活.
其實.
門徒.
不是相信的.
他們覺得婦女跟他們說的時候.
是胡言的.
你是不是亂說的.
其實婦女也不太相信.
但不相信的人.
又記得耶穌說.
祂第三天會復活.
於是他們邀請.
彼拉多.
叫人看守住那個地方.
說要好好地看守.
兩個門徒.
他們去到墓裡.
其中一個就是今天.
我們所看到的.
寫了彼得書信的彼得.
究竟主復活了.
彼得有什麼盼望呢?.
當時的信徒有什麼盼望呢?.
今天你和我.
有什麼盼望呢?.
今天我們來看看.
原來在.

$^{81}$第一節.
說到.
有蒙簡選的.
我們的盼望是蒙簡選的.
第一節的經文說到耶穌基督.
的使徒彼得.
被簡選.
分散在.
本都,加拿大.
加帕多加,亞細亞.
比推尼寄居的人.
在第一節說到.
當時期.
在第一世紀.
的.
後半節的時候.
基督徒已經遭遇到迫迫.
這裡有一個字眼.
就是分散.
分散寄居的.
各地的選民.
分散指當時彼得.
他寫的信.
給一些已經信了主的.
基督徒的.
猶太人.
猶太人很重視.
他的國籍.
雖然他們已經成為基督徒.
但是他們仍然.
不會放棄他們猶太人的傳統.
好像你信了主你都不會放棄.
你的國籍.
同時他們在羅馬帝國的官吏.
去統治之下.
這些猶太人都有他們.
享受的特權.
在主效62年.
其實基督教和.
猶太教已經無法.

$^{121}$避免分裂.
基督徒受到迫迫.
猶太人不喜歡基督徒.
亦都傷害他們.
煽動羅馬政府.
去攻擊他們.
你看當年保羅.
也是.
他未認識耶穌基督的時候.
他都去捉拿基督徒.
你看到.
這些信徒.
受迫迫.
散居去.
羅馬帝國各處.
越來越受迫迫.
原來這個信仰.
這個福音可以傳到.
羅馬帝國的各個地方.
整個羅馬帝國.
都有福音傳到.
主耶穌的兄弟雅各.
遭遇到殉難.
他死後兩年.
基督徒擁有的特權.
已經取消了.
被稱為一個不合法的團體.
在主後六十四年.
更加反基督教的情況.
越來越厲害.
彼得.
大概在大悲不之前.
寫了一封書信.
大概在主後六十三年.
下半年的時候.
已經完成了這本書.
當時.
這些神的選民.
已經分散在小亞細亞.
以北的各個地區.

$^{161}$這些寄居的地方.
當時.
各地寄居的基督徒.
大部分是奴隸.
或者是外邦人.
其實.
彼得很強調.
他們是被父神所揀選的.
他們是鼓勵讀者.
從前只有以色列人.
他們自稱自己是神的選民.
但今天你直接耶穌基督.
你信耶穌的人.
無論你是外邦人.
或者猶太人.
你都是屬於神的.
今天我們可以.
蒙請教.
蒙神的眷顧.
完全是神的憐憫.
神的揀選.
我們無論遇到怎樣的患難.
和逼迫.
都不可以奪去神給你的永生.
其實保羅在這封書信裡.
正正是提醒.
告訴我們.
在這個世界上有苦難.
不是每一次都順利.
一定會有難處.
有首詩歌.
講到「這世界非我家」.
我們是屬於神的.
但神為我們預備的家鄉在哪裡呢?.
不是在地上.
因為更美的家鄉.
是屬天的.
所以我們留在地上.
在這個世界.
我們是屬天的.

$^{201}$但在地上.
在這個世界.
神有使命給我們.
所以我們為神做見證.
等待有一天.
主耶穌基督回來的時候.
這群屬神的子民.
在末日的時候會被焦聚起來.
為甚麼.
為甚麼要選擇我們呢?.
有甚麼原因呢?.
在第二節.
在前書第一章第二節.
講到「就是照父神的預知.
藉著聖靈得以成聖.
以致信服耶穌基督.
有蒙他血所灑的人.
願因為平安多多地賜給你們」.
原來這是神的計劃.
神在創世之先.
祂預備了耶穌基督.
去救贖我們.
所以這是神的旨意.
一開始已經定了.
在創立世界以前.
已經有的.
揀選了我們.
另一段經文.
是在二弗所書的一章四節.
因為祂從創世以前.
在基督裡揀選了我們.
使我們在祂面前成為聖潔.
沒有瑕疵.
沒有愛心.
原來耶穌基督.
祂為我們而死.
祂是沒有罪的.
反而說我們罪人而死.
藉著聖靈.
使我們成為聖潔.

$^{241}$使我們得到救恩的福氣.
可以與世人分別出來.
所以將我們分別為聖.
結證我們.
可以去事奉祂.
耶穌基督的血.
成為神和人的納入的憑據.
所以有一個憑據.
我們接受的人.
一起納入的人.
是需要信服的.
所以主說我們捨命.
祂用重價買贖我們.
我們更加要信服主.
祂那麼愛我們.
第三方面.
就是讓我們看到我們的盼望.
是因為我們無了重生.
在第三節.
在第三節說.
願眾贊歸於我們主.
耶穌基督的父神.
祂曾照自己的大憐憫.
直到耶穌基督.
從死人中復活.
重生了我們.
使我們有活的盼望.
重生.
即是再生一次.
給我們一個新的生命.
對於父神.
是如何重生我們呢?.
其實重生就是我們個人.
與神產生生命的關係.
這個關係是出於神的憐憫.
直到耶穌基督.
從死人中復活.
所以人可以夢重生.
是出於神的大能.
神的大能就是.

$^{281}$看到耶穌基督.
神叫他從死人中復活.
這個大能.
重生的生命是一個新的.
是一個活潑的生命.
信徒直到耶穌基督.
從死人中復活.
得到一個新的生命.
耶穌基督死而復活的生命.
所以第四方面.
我們看看這個死而復活的生命.
是如何的呢?.
記不記得耶穌基督.
在復活之前.
祂曾經.
叫那恩成.
寡婦的兒子.
復活的.
本來很傷心的那隊.
送殯的行列.
但耶穌基督.
叫寡婦的兒子.
坐下來.
祂立即可以說話.
另外這個.
以前叫做.
艾魯.
現在是.
翻譯成葉魯.
葉魯的女兒.
古回堂的女兒.
她快死了.
她來求耶穌.
去她家.
但還沒去到的時候.
她已經死了.
但耶穌說.
不要怕.
不如拍子去信.
於是去到她家.

$^{321}$叫這個女兒起身.
她立即可以起身.
你看到.
耶穌在當中.
有叫她們復活.
另外這個.
你更加熟悉了.
馬大瑪利亞的弟弟.
埃塞路.
已經埋葬了.
耶穌叫他們出來.
由墳墓出來.
他們的身體.
裹著屍體的布.
還包著.
耶穌立即叫他們.
幫他們把布拿走.
這幾個人.
他們都是死了.
復活了.
但他們還有原來的身體.
他們復活.
不是得到靈性的身體.
所以他們都是死了.
但耶穌基督.
復活之後.
他們的形象有改變.
當他們出現在門徒之前.
門徒.
嚇了一跳.
覺得他們.
看到了什麼呢?.
他們不用.
自己開門穿牆.
耶穌叫他們.
摸摸我的手和腳.
當主人.
向他們顯示的時候.
他們才知道是主耶穌.
馬大瑪利亞也是.

$^{361}$本來以為耶穌是圓丁.
但當耶穌叫他們的名字的時候.
他們忽然間.
才知道.
這是耶穌基督.
那時候才知道.
是夫子.
拉波尼.
還有兩個.
在二萬五思路上的.
門徒.
耶穌跟他們說話.
說了很久.
他們都不知道是耶穌.
當耶穌在他們面前.
打餅來.
給他們的時候.
他們恍然大悟.
這幾個.
原來是神向他們啟示.
他們才知道.
他們知道耶穌基督.
在門徒當中.
祂的寧靜身體.
是要向人啟示.
他們才知道.
但祂不是只是暈.
暈是沒有骨.
沒有肉的.
但祂問他們有什麼吃的.
他們說有一片.
烤魚.
祂在他們面前吃了.
你仍然可以認出祂的身體.
不是叫他們看丁行的手.
仍然有祂的身體在.
但祂那時候是一個靈性的身體.
主耶穌接受.
這個死是勝過了死亡.
耶穌既然復活.

$^{401}$有一天我們每一個信主的人.
你都會復活.
復活之後你就會得到這個.
靈性的身體.
這和我們.
天然的身體是不同的.
靈性的身體是不會.
扭壞不會死亡的.
可以完全適應.
榮耀.
屬靈的生活.
但天然的身體是會扭壞的.
會死亡的.
我們每一個信徒就是盼望.
有一天我們可以復活.
去得到這個靈性的身體.
這個榮耀的身體.
所以你看到.
屬血氣的生命.
和屬靈的生命.
是很大的對比.
到主耶穌回來的時候.
死人要復活.
如果耶穌基督回來的時候.
你還在地上的話.
你的身體要改變.
當日耶穌可以.
叫一些睡了的人去復活.
今天.
祂使得死人復活是更加容易的.
第五方面.
重生的生命.
在前書第一章.
第三節的後半段.
那裡說到耶穌基督.
從死人中復活.
重生了我們.
使我們有活的盼望.
復活真的說到.
我們生命的改變.

$^{441}$復活之後除了形體的改變.
有一個靈性的身體.
復活的人得到一個.
屬靈的生命.
在哥林多前書的.
第十五章.
五十節那裡說到.
弟兄們我要告訴你們的是.
血肉之軀不能承受.
神的國.
必扭壞的也不能承受不扭壞的.
這裡看到.
復活的人.
得到屬靈的生命.
所以復活的生命是屬靈的.
要有屬靈的生命.
才可以承受神的國.
我們要承受神的國.
我們一定要有屬靈的生命.
至於重生.
生命就很重要.
這裡說到.
信主之後我們生命的改變.
我們的內在.
應該有改變.
信了耶穌之後.
舊事已過都變成新的.
你是一個新造的人.
大家很明顯.
看到浪子的比喻.
那裡說到浪子.
他本來.
他爸爸還沒離開世界.
就問他要遺產.
還用光了.
接著他的光景.
很落魄.
要和豬爭食.
最後他想一想.
我落了這樣的光景.

$^{481}$回去做父親的僕人就好了.
當他回去的時候.
爸爸每天都在等他.
當他回去的時候.
他醒覺.
他和爸爸說什麼呢?.
父親我得罪了天.
又得罪了你.
從今以後我不配稱為你的兒子.
你看到一個人.
最怕的是什麼呢?.
他不知道自己錯了.
他知道自己的兒子會悔改.
他看到自己的問題.
他願意悔改.
得罪了天又得罪了你.
我們不得罪了燈.
得罪人還會得罪我們天父.
他看到了.
他爸爸很開心.
馬上叫人拿袍子給他.
給他戴上戒指.
承認身份.
穿上鞋子.
僕人不穿鞋子.
他的兒子.
給他穿上鞋子.
給他身份的象徵.
承認是他的兒子.
還要和他慶祝.
他說.
爸爸說什麼呢?.
他說.
我這個兒子是死而復得的.
死而復活.
過去.
真是.
太荒唐了.
但今天.
有改變.

$^{521}$我們今天認識耶穌基督的人.
我們生命也有改變.
記不記得在耶穌基督.
死後.
有一個阿利瑪泰的藥室.
他要求彼拉多.
領耶穌的身體.
這個阿利瑪泰.
這個省可能.
是一個城.
可能是辯雅敏境內的拉瑪.
在猶太省.
猶太省裡面的.
這個藥室.
可能在耶穌基督.
在被釘十字架斷氣的時候.
他在場的時候.
來到那裡.
你見到他很富有.
但他一點也不驕傲.
他是猶太公會的義士.
他見彼拉多很容易.
他曾經是暗暗的.
做主的門徒.
但因為他愛耶穌.
這個時候他也不避嫌.
也不怕別人說他閒話.
他要去.
領耶穌的屍體去埋葬.
你見到他買細麻布.
是新的.
去將耶穌取下.
由十字架拿下來.
一個屍體.
還要起釘.
但在聖經裡面說.
凡是摸死屍的.
會不潔淨七天.
摸的人或物.
也不潔淨到晚上.

$^{561}$但在約翰福音裡.
說到這兩個人.
說到亞利瑪泰藥室.
和尼哥底姆.
以前叫尼哥底姆.
現在叫尼哥德姆.
在翻譯上.
雖然他們不潔淨.
但因為他敬愛耶穌.
他也有犧牲的心.
尼哥底姆.
也是一個財主.
他也是猶太公會的義士.
他也曾經晚上去找耶穌.
但可能.
這個藥室的行動.
令他更壯膽.
他帶了一百斤的煤藥.
和陳香.
這些可能是末常粉.
加在一起.
是要去找耶穌的.
一百斤等於.
三十三公斤.
是皇室人員.
他們埋葬用的份量.
你看到他.
很多.
可能一世麻布.
但無論如何.
他的心.
是給耶穌用的.
他很敬愛耶穌.
他去.
表示他的心意.
看到這兩個人.
可以看到他們.
沉睡的靈已經蘇醒過來.
他們做了一些.
不可思議的事.

$^{601}$另外一段經文.
洗禮的人.
可能會有印象.
在羅馬書第六章.
三到五節.
難道你們不知道.
我們這受洗歸入基督耶穌的人.
就是受洗歸入他的死嗎?.
所以我們直著洗禮歸入死.
和他一同埋葬.
是要我們行事為人.
都有新生的樣子.
像基督.
直著父的榮耀.
從死人中復活一樣.
我們若與他合一.
經歷與他一樣的死.
也將經歷.
與他一樣的復活.
提到我們.
要與主同死.
同埋葬.
同復活.
我們要看信徒的得救.
不是看洗禮.
當然洗禮是一個過程.
是需要做的.
最重要是看你.
生命有沒有改變.
如果生命沒有改變.
洗禮很多次也沒有用.
你一樣.
不可以承受神的國.
承受神的國.
是要屬靈的生命.
所以你檢視一下.
如果你今天信了主很多年.
你還是依然如故.
或者有些不信的人.
覺得我比你還要好.

$^{641}$為什麼要來教會呢?.
我們真的要檢討一下.
我們到底有沒有.
高舉耶穌呢?.
有沒有讓祂去掌管我們的生命呢?.
我們需要有.
耶穌基督給我們.
這個屬靈的生命.
我們才可以承受神的國.
我們透過耶穌基督.
救贖了我們.
我們可以與耶穌基督一起復活.
今天基督徒最大的盼望.
最喜悅的是什麼呢?.
就是等候主榮耀的再臨的時候.
我們這個.
敗壞的身體.
我們要復活.
死人要復活.
耶穌回來的時候.
如果你還在這個世界.
你的生命會改變.
你的身體要經歷改變.
所以如今我們在地上.
我們如何與主同死.
同埋葬同復活呢?.
我們才要致死.
我們這個老我.
過去我們這個舊人.
我們要致死了他.
我們要埋葬了他.
所以聖經說.
耶穌叫我們天天背起.
自己的十字架要跟從耶穌.
耶穌提醒我們.
天天的時候.
我們要為主復活.
為主去活.
我們的生命.
是不斷要更新的.

$^{681}$耶穌升天之後就派聖靈來.
聖靈會經常提醒你的.
你覺得不對的話.
就要去祈禱.
假怪安的孫教士.
他22歲的時候.
已經很清楚知道.
神父召他去中國.
他原本.
不是很明白聖靈.
當他看到一個.
舊的園館.
彎曲變形.
但是他可以湧流出.
很乾淨很晶瑩的水.
令他忽然間.
茅塞.
頓開.
他一醒來.
原來這樣也可以嗎?.
原來聖靈的真理.
就是這樣的.
信我的人如經上說.
從他的福中.
要流出活水的江河.
我們要信徒.
自己要蒙了恩.
下一步.
不是坐著享受.
要帶人去信主.
使別人跟你一樣.
蒙受神的恩.
聖靈就會在你裡面.
像活水的江河.
如願不絕的.
永流不息.
所以信靠耶穌的人.
不單是靈性的饑渴.
得到滿足.
之後.

$^{721}$就像是全元在你裡面.
永流不息.
所以他向神祈禱.
就是告訴神.
我像一個圓館.
變形了.
彎曲了.
這樣的圓館.
這樣的生命.
求聖靈充滿我.
讓我流到別人身上.
就像水流過圓館一樣.
當然我們不知道神.
如何去成就祈禱.
但沒多久就去了中國.
他首先在揚州.
接受一些婦女.
語言訓練.
在學校裡.
然後就去了.
鑊綠園.
在1900年6月.
就有權匪之亂.
這個時候.
他和一對夫婦.
清桂蓮.
夫婦是在一起的.
他們收到一個消息.
他們被迫逃到.
一個廟裡.
住了六天.
然後又再去農村住了一個月.
但權匪.
打聽他們的消息.
他們很緊張.
然後他們再逃到一個山洞.
那個山洞.
形容得很潮濕.
在那個地方.
很美麗.

$^{761}$當他們就要被發現的時候.
忽然間.
這個姐妹.
有一個神的話語.
很強烈的.
浮現在她心上.
所以有千人.
撲倒在你旁邊.
萬人撲倒在你右邊.
這災卻不得臨近你.
她看到這句神的話.
她很清楚知道.
神與她同在.
所以她就不怕.
終於他們被發現了.
權匪很狠的.
把他們推到山洞裡.
當他們.
走回到城裡的時候.
真的很狼狽.
很多人圍觀.
但不敢出聲.
她看到一個劉嫂.
她就很大聲的跟她說.
神與我們同在.
我們不怕.
後來他們被送到衙門.
但那些官員很怕事.
把他們踢到.
一個更危險的路上.
後來他們就.
被迫坐船.
船不知道會帶他們去哪裡.
船開了不久.
在一個偏僻的地方.
船夫就放下他們.
原來那些權匪.
告訴他們.
你們去到沒有人的地方.
解決了他們.

$^{801}$但船夫不忍心.
所以就放下他們.
於是他們就躲在草叢裡.
在草叢裡.
他們就等有沒有機會.
有其他船夫經過.
可以載他們.
這個時候.
清桂蓮夫婦.
清桂蓮孫女士.
就問她太太.
又問賈姑娘.
賈姑娘.
今天上帝有沒有.
說話給你們呢.
賈姑娘.
忽然間有一個很奇怪的想法.
她就說.
很想小鳥.
送一封信給我們.
沒多久.
有人通風報信.
知道他們在那裡.
於是又被那些權匪.
綁住他們.
於是就把他們.
押送到.
一個廟裡.
那個廟裡非常骯髒.
有蒼蠅,老鼠.
很臭的.
在那個環境下.
他們整晚.
這三個人.
睡不著.
他們內心充滿著.
好像看到耶穌多少受難.
那個很深刻的經歷.
耶穌正要.
走去各個攤位.

$^{841}$那個光景.
他們在那個情形下.
雖然被人侮辱.
但他們住了一個星期.
忽然間.
有一天下午.
賈姑娘被人在夢中.
吵醒.
原來清貴年.
這位宣教士.
他拿到一個字條.
有人在外面.
把字條拿了一塊.
他打開來看.
原來.
有人給了字條.
告訴他.
不要害怕.
權費已被中國和外國軍隊掃清了.
我會到天津通知你們.
國家的軍隊.
來保護你們.
這個時候.
這個.
賈桂康姑娘.
她真的不知道要笑還是要哭.
她說.
她看到這個消息.
就是這件事.
然後.
她雖然看到這個消息.
但不是那麼快.
脫離.
他們還住了幾個星期.
被人俘虜的情形下.
住了幾個星期.
清貴年.
宣教士有一個五歲的女兒.
因為她患了理疾.
所以她死了.

$^{881}$他們非常傷心.
到了十月底的時候.
英軍法軍來救她們.
有一個有錢的中國人.
他很有智慧.
想了一個辦法.
去幫他們.
去救他們.
雖然各方的謠言.
以為他們死了.
他們已經停止.
取消了計劃去救他們.
在這次行動裡.
有189個宣教士.
犧牲了.
其中有58個成員.
是內地會的.
賈姑娘很深深的體會到.
原來神一次又一次的.
救她們.
是有她的目的.
表示她們有使命在身.
當她.
回到自己的國家休教不久.
她又回到中國.
回到和樂縣.
繼續和清貴年夫婦.
宣教士一起工作.
那時已經是.
1901年.
她已經30歲.
她就專心的.
在農村裡.
報道工作.
她大部份時間.
都在一個中國人的家庭裡居住.
她可以了解到.
中國人的言語.
風俗.
社會隱藏的毒瘤.

$^{921}$或者是有什麼秘密的組織.
她非常了解.
都幫到她日後.
她去做報道工作.
是有很大的影響.
她專做婦女的工作.
向婦女去講道.
所以她可以一針見血.
講出她們的問題.
因此很多的婦女.
她們的生命得到改變.
1926年.
她曾經寫給朋友一封信.
她說15年.
警察會總是.
滿心喜樂的.
派她去某個地方工作.
然後又滿心喜樂的.
接她回來.
現在的情況已經告一段落.
她去過15個省份.
舉行過183場報道會.
在那裡跟隨耶穌的.
太太們.
或者是小姐們.
有5342名.
她很清楚.
神的確為她.
開了這裡.
婦女福音工作的大門.
她讚美神.
她說神從未違背.
從天上來的意象.
當福音大門一打開.
她就毫不猶豫地進去工作.
當日她說.
今天有很多門.
福音大門已經關了.
要進去的地方.
幾乎是不可能的.

$^{961}$各位弟兄姊妹.
今天你還有自由.
可以傳福音.
你還有健康.
你還有生命.
你願不願意為主而活呢?.
你的靈蘇醒了沒有呢?.
你願不願意回轉.
歸向神呢?.
你肯不肯聽神的話?.
她真的盼望聖靈充滿你.
因為聖靈會引導你.
帶領你去完成主的使命.
我們一起去祈禱.
你讓我們可以成為你的兒女.
今天我們看到很多人.
他們都願意.
很努力地放上自己.
今天我們.
是不是落後了呢?.
你說在前的要在後.
在後的要在前.
求你聖靈提醒我們.
鞭策我們.
何時我們違背了你的意思的時候.
求你幫我們能夠走回正路.
我們聽你的旨意.
以致我們可以.
走在你的道路當中.
去榮耀你.
求你帶領我們的心.
我們這樣祈禱教託.
奉主耶穌基督的名求.
\newpage



\section{馬可福音 8:27-37}
\label{sec:D5wSPpsbnpU}
\textbf{2022.05.01 《你願意付甚麼代價 ?》 馬可福音 8\_27-37 (和修版) 老耀雄牧師}
\newline
\newline
連結: \href{https://youtube.com/watch?v=D5wSPpsbnpU}{\texttt{https://youtube.com/watch?v=D5wSPpsbnpU}} ~~~~ 語音日期: 2022-05-06
\newline
\newline
\hyperref[sec:O3rEhqjezKM]{\small{< < < PREV SERMON < < <}}
~
\hyperref[sec:index_chronic]{\small{[返順時目]}}
~
\hyperref[sec:index_scriptual]{\small{[返順卷目]}}
~
\hyperref[sec:8C2XRS_C_0M]{\small{> > > NEXT SERMON > > >}}
\newline
\newline
馬可福音 8:27-37
\newline
\begin{longtable}{cl}
\hline
\hline
章節 & 經文 (和合本修訂版)\\
\hline
8:27 & \begin{tabularx}{0.7\textwidth}{X} 耶穌和門徒出去，往凱撒利亞‧腓立比附近的村莊去。在路上，他問門徒：「人們說我是誰？」 \end{tabularx} \\ \\ \relax
8:28 & \begin{tabularx}{0.7\textwidth}{X} 他們對他說：「是施洗的約翰；有人說是以利亞；又有人說是先知中的一位。」 \end{tabularx} \\ \\ \relax
8:29 & \begin{tabularx}{0.7\textwidth}{X} 他又問他們：「你們說我是誰？」彼得回答他：「你是基督。」 \end{tabularx} \\ \\ \relax
8:30 & \begin{tabularx}{0.7\textwidth}{X} 於是耶穌切切地囑咐他們不可對任何人說起他。 \end{tabularx} \\ \\ \relax
8:31 & \begin{tabularx}{0.7\textwidth}{X} 從此，他教導他們說：「人子必須受許多的苦，被長老、祭司長和文士棄絕，並且被殺，三天後復活。」 \end{tabularx} \\ \\ \relax
8:32 & \begin{tabularx}{0.7\textwidth}{X} 耶穌明白地說了這話，彼得就拉著他，責備他。 \end{tabularx} \\ \\ \relax
8:33 & \begin{tabularx}{0.7\textwidth}{X} 耶穌轉過來看著門徒，斥責彼得說：「撒但，退到我後邊去！因為你不體會神的心意，而是體會人的意思。」 \end{tabularx} \\ \\ \relax
8:34 & \begin{tabularx}{0.7\textwidth}{X} 於是他叫眾人和門徒來，對他們說：「若有人要跟從我，就當捨己，背起自己的十字架來跟從我。 \end{tabularx} \\ \\ \relax
8:35 & \begin{tabularx}{0.7\textwidth}{X} 因為凡要救自己生命的，必喪失生命；凡為我和福音喪失生命的，必救自己的生命。 \end{tabularx} \\ \\ \relax
8:36 & \begin{tabularx}{0.7\textwidth}{X} 人就是賺得全世界，賠上自己的生命，有甚麼益處呢？ \end{tabularx} \\ \\ \relax
8:37 & \begin{tabularx}{0.7\textwidth}{X} 人還能拿甚麼換生命呢？ \end{tabularx} \\ \\
[1ex]
\hline
\hline
\end{longtable}
$^{1}$各位姐妹平安!經過了四月的壽苦節以及復活節之後,我期望大家的屬靈生命能夠被聖靈喚醒,真真正正去思考今年教會的年提,委身事服.教會的年提不是一個口號,而是教會帶領會眾在今年的信仰反思和實踐的具體..
年提的真正目的,不是想會眾參與教會做義工,為教會做事.如果只是參與教會的事工,而不了解事奉的意義,這個屬靈生命的反省,只是留於事工化,不是教會設立這個年提的真正目的..
年提的真正目的,是要想會有在屬靈的層面裡面去回應你和耶穌的關係.當大家去思考你和耶穌的關係的時候,你必然會面對一個處境,你怎樣去回應耶穌為你所做的呢?你怎樣去回應耶穌對你的愛呢?.
受苦節和復活節的崇拜,就帶給大家一個很深入的思考.因此,你的回應就是一種與神關係的反思,你可以為主耶穌做些什麼呢?.
亦都是今天的講題,你願意付一個什麼代價?亦是你屬靈生命邁向成長的一個必然過程.因此,今天我想透過馬福音8章27至37節,去探討作耶穌的門徒的使命和代價.我們開始的時候有一個祈禱..
親愛的主,求你開開我們的心,讓我們明白你的話,將你的話語藏在我們心裡面,成為我們腳前的燈,路上的光,指引我們前面的路.我們的祈禱,奉教主耶穌基督,得勝名字祈求..
先講到我們先了解一下經文的背景.在馬福音裡面,耶穌曾經三次預言過自己的身份和使命,三次的模式大致都一樣.每當耶穌講明他要受苦,要死亡,以及復活的使命之後,門徒都是不明白的,包括第一次彼得對耶穌的斥責..
第二次門徒彼此爭論誰是大誰是小,以及第三次雅各和約翰的母親要求他們兩個兒子坐在耶穌的左右兩邊.最後耶穌就向門徒講明作門徒的代價,品格和條件..
而在耶穌三次講論自己的使命的經文的前後,都有醫好黑眼的神蹟作為一個啟認,包括在白菜大醫好一個盲人,以及在耶利哥醫好盲人巴底買,黑眼被醫好是一個隱喻..
如意門徒的屬靈奉建,由看不到到看到耶穌的真正使命.我們今次講到的經文,就是耶穌第一次預言自己的身份和使命的一段過程..
作為耶穌的門徒的使命,就是要按著天父的心意而行.耶穌在27-29節曾經兩次主動向門徒講出一個身份的詢問,他說:「人說我是誰?」以及:「你們說我是誰?」.
耶穌首先問完其他人對自己的看法,然後再問門徒對自己的看法,當中有什麼意義呢?.
可以是一種漸進式對耶穌的了解,由一般人的觀點作為開始,然後再進入到門徒的觀點,由淺入深地了解耶穌的身份.又或者是一種比較,凸顯出一般人和門徒之間在觀點上的差異,表明門徒和耶穌的親密關係..
不過,無論是怎樣的理解,有一點是可以很肯定的,就是耶穌很關注門徒對自己身份的看法.門徒已經跟了耶穌一段日子了,耶穌現在正是朝著耶路撒冷出發,承擔天父交託給他的使命,或者是時候要求門徒向自己表一表態了..
彼得回答說:你是基督.彼得在眾多的門徒之中,他首先能夠確認耶穌是基督的身份.如果我們對照一下《芙蕾福音》,我們會有很多的資料作為補充..
《馬太福音》對耶穌的身份有多一些的描述.在《馬太福音》十六章十六節,彼得說:你是基督,是永生神的兒子.永生神的兒子,彼得說出了耶穌和神的關係..
另外在《路加福音》也有補充,在《路加福音》九章二十節,彼得說:是神所立的基督.說明了基督的身份是由神所命立的.耶穌是基督,耶穌是永生神的兒子,耶穌是神所立的,通通都是指向舊約所指的彌賽亞..
這就是彼得對耶穌身份的宣認了.耶穌在《馬太福音》更加說到彼得是有福的,因為他能夠說出耶穌的身份,完完全全是天上的父所指示的.彼得能夠知道耶穌的身份的確是一個好的開始,畢竟彼得也是耶穌第一批呼召的門徒之一..
之後三十節,耶穌就吩咐門徒不要把自己的身份傳開去,其實這個情況不單單是在這次發生,在一章四十四節,耶穌醫好大麻風的病人,他也吩咐大麻風不要告訴別人,你只需要把身體給祭司擦汗就可以了..
在三章十一至十二節,當污鬼認出耶穌的時候,哭著說你就是神的兒子了,耶穌又再三叮囑他不要把自己顯露出來..
在七章三十五及三十六節,耶穌醫好耳聾不能說話的人,也吩咐他不要告訴別人.耶穌為什麼不想公開自己的身份呢?讓別人知道他是彌賽亞不是更好嗎?這是馬可筆下彌賽亞的秘密,我稍後再和大家解釋..
當門徒已經知道耶穌的身份之後,耶穌就開始向門徒講他的使命了..
第三十一節,人子必須受許多苦,被長老,祭司獎和民事戲拙,並且被殺,過三天復活.受許多苦,被殺,復活,這就是耶穌所要走的路,也是耶穌在世上的使命..
第三十二節,經文說,耶穌是「明白地說者」,即是說,明白地是一種強調的用語.呂鎮忠認為,耶穌坦然無懼地說這些話,即是說,耶穌是毫不含糊,很認真地去表明他的使命..
然後,彼得就拉住他,責備他.和合本的譯法,就是叫彼得拉住他,勸他.和修本的譯法,就是叫責備他.責備相比起和合本的勸他,顯示出彼得當時對耶穌的態度強烈很多..
勸耶穌和罵耶穌,是兩碼子的事.證明彼得是極之抗拒耶穌將要受苦死亡的事實.換言之,彼得不認同耶穌作為基督,以色列所等待的彌賽亞,你所承擔的使命,竟然要走的是一條受苦和死亡的路..
雖然耶穌最後說,三天後就會復活,我相信彼得已經聽不進去.彼得看著門徒,責備彼得說,撒旦,退到我後面去.因為你不體會神的心意,而是體會人的意思..
這句話也不夠精準,中文譯本也很好.撒旦,走開,你所想的不是上帝的想法,而是人的想法.耶穌看著門徒,對著彼得說這番話,有兩點可以很肯定的..
第一,耶穌很想和彼得劃清界線.第二,耶穌不想彼得影響到其他門徒,即使其他門徒也可能和彼得有相同的想法..
一個能夠認出耶穌是基督身份的彼得,竟然被耶穌斥責為撒旦,要和他劃清界線,原因就是彼得作為耶穌的門徒,他只能以自我中心去看耶穌的使命,他不能以神的角度去看耶穌的使命..
這也是馬可不下彌賽亞的秘密,耶穌不想太早公開自己身份的原因,因為他不想跟隨他的人,錯誤去理解他的使命..
以色列百姓以為彌賽亞的來臨是要幫他們推翻羅馬政府,重新統治百姓,彼得也認為耶穌要走的路不應該是一條受苦和死亡的路,彼得更加不接受..
基督為什麼要受苦?這怎麼可能是上帝計劃的一部分?門徒雖然知道耶穌是基督,是神的兒子,是神所立的彌賽亞,但卻錯誤地理解耶穌的使命,以至他們跟隨耶穌期間還要爭論誰才是最大..
雅各和約翰又要求坐在耶穌的左右兩邊,他們不明白當中的意義,門徒始終不明白耶穌的使命,因為他們仍然是以自己的心意去取代了耶穌的使命..
所以當耶穌很清楚地向門徒說明,祂要走一條受苦被殺然後復活的路的時候,彼得就勸了祂,責備耶穌..
我們知道彼得是最早被耶穌呼召的門徒之一,當時他和安德烈是即時放下慾望去跟隨耶穌的,門徒曾經按著耶穌的吩咐並且領受權柄,兩個人一組去傳道,醫病,趕鬼..
他們經歷過耶穌如何餵飽五千人五餅二魚的神蹟,他們見過耶穌所行的神蹟,他們對耶穌的認識比其他人更加清楚,不過他們仍然是一個屬靈的盒子,他們需要耶穌的醫治,然後才能夠看清楚耶穌在世的使命..
彼得在當時的的確確不能夠看清楚耶穌的使命,甚至我們知道他日後三次不認主,不過我們都知道彼得最後都知道耶穌的使命,並且為了回應耶穌對他的呼召,最終走上殉道者的行列..
今天我們很清楚知道耶穌已經完成了他在地上的使命,不過耶穌仍然好像昔日呼召彼得一樣,向我們同樣發出呼召,要我們參與他的使命..
我們這一班屬主的兒女,是否都可以像彼得一樣,願意回應主的愛,放下我們安恕的生活,去承擔傳福音的使命呢?.

$^{41}$今天我們不會像彼得一樣,走上殉道者的行列,但是我們能否像彼得一樣,願意去回應耶穌呢?縱使軟弱,縱使得罪主,三次不認主,但是沒有辦法不說的是彼得對主的回應真誠,是真心的..
洪恩家的弟兄姊妹,我想向大家發出一個問題,我們認不認識耶穌?.
認識耶穌!.
沒錯,我們每一個人都會說認識耶穌,但是耶穌對我們的生命有什麼意義呢?.
我們是否需要像彼得一樣,要耶穌醫好我們的眼睛,以致我們有一個熟齡的眼睛,看清楚我們自己今天身處的處境,以及我們能夠為主承擔一個什麼使命呢?.
聖經說過,學生不能夠大過先生,如果我們承認耶穌是我們的主,我們的使命都應該像耶穌一樣,是按著神的心意而行,而不是按著自己的心意而行的..
那麼神的心意又是什麼呢?.
在馬福音後一點的十二章,三十至三十一節,耶穌說:你要盡心盡性盡義盡力,愛主你的神,其次就是愛人如己..
我們問問自己,我們對神,我們對身邊的人,尤其是那些不認識主的人的朋友,我們可以為他們盡力嗎?.
第二個作耶穌門徒的代價就是寫記,背起自己的十字架,跟從主..
耶穌在第三十四節很明確地跟門徒說,若有人要跟從我,就當寫記,背起自己的十字架來跟從我..
現代中文譯本是這樣譯的,如果有人要跟從我,就要捨棄自己,背起自己的十字架來跟從我..
寫記和背起兩個字的原文是一個命令式的語法,即是你無法選擇的,是耶穌的命令..
耶穌是以命令的語氣跟門徒說,如果你要跟從我,你就必須要捨棄自己,背起自己的十字架..
捨棄自己是一種心態,背起自己的十字架是一個行動.捨棄自己含有放棄個人意願的心態,也是一種不以自己的喜好為準則的態度..
簡單地說,寫記就是放下你自己的自我中心,不以自己的心意為標準..
背起是一個行動的象徵,背起自己的十字架就是一種具體的行動.十字架是羅馬帝國處置死囚的一個極刑..
這個刑罰會首先將死囚的衣服脫下,然後鞭打他,然後由死囚自己背起組成刑具的長刑木條,.
然後由他自己背起這個長刑木條走去刑場,然後被釘死在自己背起的十字架上而死..
十字架代表著受人的侮辱,受苦,以致死亡.背起自己的十字架就是背起自己的刑具..
每個門徒的十字架都不一樣,有些人重一點,有些人輕一點,但是都是向著同一個目標..
門徒所做的一切犧牲都是必須要為耶穌而做的.他們所背的十字架必須是由耶穌的十字架所決定的..
而門徒一旦背起自己的十字架,就必須有一個心理準備,就是面對受苦,甚至死亡..
也就是說,想跟隨耶穌,就不一定是安逸的生活,而是放棄一種自我中心,並且甘心去走一條受苦,甚至死亡的路..
神學家潘博華在一本追隨基督的書裡面是這樣寫的,他就當基督呼召一個人,.
他就在呼召他,為他去死.為什麼要付這麼大的代價去追隨耶穌?為什麼要付這麼大的代價去死亡,去作為耶穌的門徒?.
耶穌在三十五節裡面有解釋的,他說:「凡要救自己生命的,必喪掉生命;凡為我和福音喪掉生命的,必救了生命.」.
這句話很有弔詭性,如果用中國人說的話,會有一點佛偈,色即是空,空即是色.究竟是救生命,還是喪失生命呢?.
經文這兩句部分是平衡的,是從正和反兩方面去解釋的道理,這個道理就是為了事奉耶穌,為了福音的緣故,犧牲生命的,這個生命最終是會得到拯救..
生命可以指肉身的生命,或者是指屬靈的生命.我們知道肉身的生命是有限制的,我們每一個人都會經歷生老病死這個階段,生老病死這個過程是不會改變的..
如果有人很想利用世界所提供的科技方法,用以確保自己能夠永恆地生存,保住自己的生命,是不會成功的,因為肉身的生命最終是會死亡..
那什麼叫屬靈生命呢?屬靈生命就是一種與神相交,與神聯合的生命,也是聖經裡面所說的神在你裡面,你在神裡面的生命..
所以如果你是為了神,或者你是為了福音的緣故,肉身的生命縱然被人殺了,甚至死亡,但卻可以與神建立一個屬靈生命的聯繫,而這個屬靈生命卻是永恆的..
屬靈生命超越了肉身生命最終都會死亡的限制,也超越了今生的空間,我們的生命直接就可以進入永恆..
這就是「凡為我和福音喪掉生命的必救了生命」的道理,原來門徒要寫記背起十字架,與永恆的生命是一個因果關係,寫記背起十字架是因,得到永恆的生命就是果..
所以三十四和三十五節的理解,你可以簡單地說,如果有一個人要跟從耶穌,就要放下自我,甘願走一條受苦的道路,就算因為這樣而失去生命,但最終都會得到一個永恆的屬靈生命,與自己建立一個關係..
這個因果關係很清楚地說明,想要有一個永恆的生命,你必須要付代價..
耶穌最後在三十六節,加多了一句充滿價值判斷的話,他說人就是賺得全世界,賠上了自己的生命,有什麼益處呢?.
人還可以拿什麼來換生命呢?耶穌以賺取全世界,與屬靈生命作一個比較,讓人思考一下,哪一種的生活方式對你最有益最有利呢?.
然後用一句疑問句,人還能拿什麼來換生命呢?來帶出作者的答案,你願意拿什麼來換取屬靈生命呢?.

$^{81}$很多人都引用過三十六節這段經文,來說明生命可貴,不可以任意糟蹋,鼓勵世人要珍惜生命..
不過經文背景所指的生命,是指屬靈的生命,是上帝與人建立關係之後所賦予人的生命,這種屬靈生命不是世俗世界可以提供得到的..
你不認識主,你不會有屬靈生命,你這個世界多成功,多有成就,多有地位,你都沒有這種屬靈生命,你都不會有這種永恆的盼望..
因為基督徒即使事事亨通,生活安舒,贏得了全世界,但卻沒有遵行作門徒的必需條件,就是寫記,背起自己的十字架,.
這就等於他賠上了自己的屬靈生命,他現在所擁有的,對於他的永恆來說,有什麼益處呢?.
紅春家弟兄姊妹,今天講題是「你願意付什麼代價?」.
如果我們願意作主的門徒,作主的僕人,我們願意回應主的話,主已經對我們發出了一個很明確的命令,就是要寫記,背起自己的十字架..
我們的生命只要有耶穌,一切都值得了,除了有耶穌和我們一起走這條事奉的路之外,我們教會還有一群同路人,有這麼多屬靈的同行者陪著我們,我們就不會孤單了..
剛才美蓮姐妹分享了關古部的事奉經歷,這種事奉就是一種生命影響生命的事奉,是一種讓人可以在我們身上看到耶穌的事奉,是一種從心裡謙卑讓自己與神同行的事奉..
弟兄姊妹,神的角度需要你的參與,神的永恆計劃需要你的參與,洪恩家更加需要你的參與,每一個事奉崗位都有不同的困難,但是神的恩典永遠夠我們用,我們就在困難裡仰望那個為我們創始成中的救主耶穌,祂為我們所成就的,祂所賜給我們的,一定超過我們所想所求..
但願我們每一個屬主的兒女,將來在審判台前,我們可以這樣跟耶穌說:那美好的仗我已經打過了,當跑的路我已經跑盡了,所信的道我已經守住了,我們這一生都能夠按照你的命令,做到捨己,背起自己十字架來跟從你,成為一個無愧的工人..
我們一起去祈禱..
七位恩主,我們知道你的身份,我們也知道你的使命,我們知道要回應你的愛和呼召,我們知道不要按著自己的心意而行,而是按著你的心意而行,當我們越去追求今生的慾念,我們就越遠離你,當我們越體貼自己的私慾,就越得不到你的恩典..
因此求主,堅固我們的心,甦醒我們的靈,好叫我們甘心樂意地去回應你的呼召,成為你合用的器皿,一生被你所用,讓你的名被高舉,榮耀傳歸你..
我們的祈禱仰望是奉主耶穌基督,得勝名字祈求..
阿們..
願主祝福大家..
\newpage



\section{撒母耳記上 14:6-15}
\label{sec:8C2XRS_C_0M}
\textbf{2022.05.08 《轉敗為勝的秘訣》 撒母耳記上 14\_6-15 (和修版) 陳一華牧師}
\newline
\newline
連結: \href{https://youtube.com/watch?v=8C2XRS-C-0M}{\texttt{https://youtube.com/watch?v=8C2XRS-C-0M}} ~~~~ 語音日期: 2022-05-09
\newline
\newline
\hyperref[sec:D5wSPpsbnpU]{\small{< < < PREV SERMON < < <}}
~
\hyperref[sec:index_chronic]{\small{[返順時目]}}
~
\hyperref[sec:index_scriptual]{\small{[返順卷目]}}
~
\hyperref[sec:A3yUdy51iFc]{\small{> > > NEXT SERMON > > >}}
\newline
\newline
撒母耳記上 14:6-15
\newline
\begin{longtable}{cl}
\hline
\hline
章節 & 經文 (和合本修訂版)\\
\hline
14:6 & \begin{tabularx}{0.7\textwidth}{X} 約拿單對拿兵器的青年說：「來，我們過去到那些未受割禮之人的駐軍那裡，或者耶和華為我們施展能力，因為耶和華使人得勝，不在乎人多人少。」 \end{tabularx} \\ \\ \relax
14:7 & \begin{tabularx}{0.7\textwidth}{X} 拿兵器的對他說：「隨你的心意做吧。你上去，看哪，我一定跟隨你，與你同心。」 \end{tabularx} \\ \\ \relax
14:8 & \begin{tabularx}{0.7\textwidth}{X} 約拿單說：「看哪，我們要過去到那些人那裡，在他們那裡展現我們自己。 \end{tabularx} \\ \\ \relax
14:9 & \begin{tabularx}{0.7\textwidth}{X} 他們若對我們這麼說：『站住，等我們到你們那裡去』，我們就站在原地，不上他們那裡去； \end{tabularx} \\ \\ \relax
14:10 & \begin{tabularx}{0.7\textwidth}{X} 但他們若這麼說：『上到我們這裡來吧』，我們就上去，因為耶和華把他們交在我們手裡了。這就是我們的憑據。」 \end{tabularx} \\ \\ \relax
14:11 & \begin{tabularx}{0.7\textwidth}{X} 二人就讓非利士的駐軍看見。非利士人說：「看哪，希伯來人從躲藏的洞穴裡出來了！」 \end{tabularx} \\ \\ \relax
14:12 & \begin{tabularx}{0.7\textwidth}{X} 站崗的士兵對約拿單和拿兵器的人說：「上到這裡來，我們有一件事要告訴你們。」約拿單就對拿兵器的人說：「跟我上去，因為耶和華把他們交在以色列人手裡了。」 \end{tabularx} \\ \\ \relax
14:13 & \begin{tabularx}{0.7\textwidth}{X} 約拿單手腳並用爬上去，拿兵器的人跟隨他。非利士人仆倒在約拿單面前，拿兵器的人跟著他，殺死他們。 \end{tabularx} \\ \\ \relax
14:14 & \begin{tabularx}{0.7\textwidth}{X} 約拿單和拿兵器的人第一次擊殺的約有二十人，都在一畝 地的半犁溝之內。 \end{tabularx} \\ \\ \relax
14:15 & \begin{tabularx}{0.7\textwidth}{X} 於是在軍營、在田野、在眾百姓中，人心惶惶，駐軍和突擊隊都戰兢；地也震動，這是從神那裡來的驚恐。 \end{tabularx} \\ \\
[1ex]
\hline
\hline
\end{longtable}
$^{1}$各位弟兄姊妹,早安!.
今天我們可以有實體敬拜,感不感恩?.
天父真的很疼愛我們,讓我們有機會一齊來敬拜.
今天很特別,因為又是母親節.
祝天下間的母親身體健康,主恩更多.
日子如何,力量也如何.
另外今天又是聖餐.
剛才魯牧師帶領我們有一個聖餐親近主的時刻.
我們的心裡是不是也得到很大的鼓勵?.
所以我們今天就要思想一個題目.
就是轉拜為聖的一個秘訣.
大家都知道香港人比人都很忙.
所以你有沒有聽人說香港人很忙碌.
每天都好像打仗一樣.
跑來跑去,騰上騰落.
好像沒有機會靜下來.
好像沒有機會休息.
因為工作,生活都很忙碌.
不單是這樣,我們信了耶穌做神的子女.
我們不單只是生活,有時候好像一般人說打仗.
更加我們會多一樣就是屬靈的戰爭.
大家認同嗎?原來我們屬神的兒女.
我們都會面對在屬靈上的爭戰.
當我們在地上來去事奉上帝.
去傳福音,去關懷人的時候.
有時候我們發覺都不是一帆風順.
都不是那麼容易的,為什麼呢?.
因為我們都會遇到魔鬼,撒旦.
有時候會攻擊我們,攔阻我們.
叫我們不要那麼愛神.
叫我們不要那麼關心人.
叫我們做一些上帝不喜悅的事情.
所以這些都是屬靈的爭戰.
今早牧師和大家看看.
如何我們可以在屬靈的爭戰裡面.
能夠轉敗為勝呢?.
我們今天就看看一個很典型的例子.
就是若拿丹.
大家有沒有留意到?.
那個背景是怎樣的呢?.

$^{41}$那個背景就是當時以色列人和非利士人打仗.
大家知道兩隊人馬是雙形見絕的.
為什麼呢?.
原來非利士人人頭湧湧.
非常壯觀.
但是以色列人就很淪盡.
原來他們只有六百人去打仗.
淒不淒涼?.
人丁單薄.
不單止這樣.
原來他們手上的武器也很落後.
因為當時在以色列的境界.
是沒有鐵匠的.
沒有人製造鐵器的.
所以以色列人的陣地只有兩把刀.
你說是不是很淒涼?.
只有兩把刀.
面對這麼龐大的非利士人軍隊.
真是雙形見絕.
當他們落在這個困局的時候.
似乎形勢一面倒.
在一個危險的時候.
又被約拿丹索羅的兒子想到一個辦法.
想到什麼方法呢?.
與其被人圍住.
沒有出路.
不如這樣吧.
索性我們突圍而出.
如何突圍而出呢?.
於是約拿丹帶著幫他拿兵器的侍從.
兩個人突圍而出.
於是整個形勢逐漸出現變化.
我們今天來看看如何在這場戰爭中.
以色列人能夠轉敗為勝.
我們發現有秘訣.
我們來看看秘訣是什麼.
第一個秘訣.
第一個秘訣在第六節.
不如我們一起說第一個秘訣.
預備起.

$^{81}$具正確觀念.
第一個秘訣是什麼呢?.
是因為約拿丹有一個很正確的觀念.
這是一個很正確的屬靈觀念.
是什麼觀念呢?.
他對侍從說.
「耶和說使人得勝不在乎什麼.
人多人少」.
這個觀念重不重要呢?.
很重要.
因為約拿丹知道自己以色列人數很單薄.
如何能夠面對這麼龐大.
人多勢眾的非利士人呢?.
他現在知道.
原來上帝使人得勝不在乎什麼.
人多人少.
各位頂尖的觀眾.
這個觀念很重要.
原來打一場仗.
能夠贏能夠輸.
不在乎什麼.
人多人少.
是在乎什麼.
上帝有沒有介入當中.
上帝有沒有親自動功.
如果上帝動功親自介入的時候.
形勢就可以扭轉過來.
我記得香港.
也問問大家.
香港年初三的時候.
通常人們是怎樣的.
年初三.
一說到年初三.
香港人就跑馬.
如果不是跑馬.
就車工.
所以有一群基督徒的牧師.
發覺年初三.
要麼跑馬要麼車工.
不是辦法.

$^{121}$我們年初三應該怎樣做.
我們年初三應該要向真神祈禱.
去呼求.
以致當我們年初三.
一個新年的日子.
我們就能夠向真神禱告呼求.
這樣就更加為香港人祈福.
求平安.
重要嗎.
重要.
所以這一群牧師.
就發起一個年初三的祈福大會.
但誰不知道.
我們想做這麼正確的事情.
但是信不信你.
不信你.
我們去借一個地方開祈禱大會.
誰知那間教會說.
不好意思借不到給你.
因為我們有職員確診了.
借不借到.
借不到.
然後就想辦法去找第二間教會.
找到的時候.
以為可以了.
誰知一個星期後.
就告訴他們.
不好意思.
我們教會地下.
這一區地下的污水.
呈現陽性.
都不知道哪裡有這樣的疫情.
所以要看清楚.
結果又等多一個星期.
等報告出來.
原來教會的範圍沒有人確診.
地下的污水和他無關.
於是才借到地方舉行祈福大會.
各位節目.
其實那時候的環境.

$^{161}$不是像現在這麼好.
有你們一群會眾在.
我記得年初三的時候.
我們在台上舉行這個祈福大會.
台下是沒有觀眾的.
你明不明白我的意思.
因為是不可以有聚會的.
只有工作人員.
各位觀眾.
雖然我們面對一二三的困難.
我們這個祈福大會.
在網上直播播出的時候.
順利嗎?.
很順利.
結果有二千多三千人.
去看這個祈福大會.
比現場的實體更加多人.
讓節目記住.
原來當我們做一件事情.
是合乎上帝的心意.
是神所喜悅的.
有一個正確的屬靈觀.
我們不要怕什麼?.
不要怕有困難.
不要怕有什麼障礙.
為什麼呢?.
因為耶和華說.
使人得勝不在乎什麼?.
人多或人少.
如果我們有神的同在.
上帝就會祝福我們的事奉.
這是第一個秘訣.
我們一起看第二個秘訣.
請大家一起讀.
第二個秘訣是什麼?.
「具同心的人」.
原來當時能夠看贏仗.
是因為約拿丹有一個很好的侍從.
這個年輕人怎樣跟約拿丹說?.
他說:我一定跟隨你.

$^{201}$跟著什麼?.
與你同心.
你說:寶不寶貴?.
很寶貴的.
原來能夠打贏一場仗.
一定要有同心的人.
或者現在來說.
如果我們要做一件事情.
就要眾志聖城.
大家要有同一個心志.
是很重要的.
如果沒有同心合意的人.
很難發揮最成功的效果.
以一個簡單的例子.
我們新界西宇牧師宮.
在疫情期間.
其實都遇到很多挑戰.
大家知不知道.
我們新界西宇牧師宮.
要負責多少間醫院?.
我們要負責五間醫院.
由屯門,青山,小欖,博愛.
再加上天水圍.
你明不明白?.
五間醫院的病人.
病床都非常非常多.
因為在疫情期間.
家人可不可以探病?.
不可以探病.
所以我們每一天.
接到的求助電話非常多.
每一天我們也都接到.
病房轉介病人給我們探訪.
當作是一個恩恤探訪.
也是很多.
疫情期間的姑娘,醫生,護士.
辛不辛勞?.
很辛勞.
他們特別是最淒涼.
2月末,3月初的時候.

$^{241}$病人都沒有病房給他們.
躺在帳篷上.
大家可否看到那一幕情景?.
醫生姑娘根本很難休息.
直踩幾間醫院.
都要去做這些救護工作.
各位頂智慕們.
面對這麼龐大的病人需要.
家人需要,員工需要的時候.
神讓我們新界西語牧師宮.
有一個很團結,同心的團隊.
以致我們能夠在這個疫情最危險,嚴峻的時候.
我們可以發揮到對人的關懷,牧養.
以及去支援前線的醫護人員.
各位頂智慕們.
這個行動如果沒有了大家的團結,同心.
能否做到?.
很難做到的.
為什麼?.
因為大家的標準都不同.
大家信心的程度都不同.
有些人是心不害怕.
有些人什麼都害怕.
大家明白我的意思嗎?.
所以我們的團隊.
大家團結,同心起來的時候.
就能夠去關懷更多有需要的人.
這是第二個原因.
我們再看第三個秘訣.
第三個秘訣.
為什麼能夠轉敗為勝呢?.
一起讀.
預備起.
一句尋求印證.
原來約拿丹打這一場仗.
他謹不謹慎呢?.
他也很小心的.
他不敢隨便亂來.
他知道上帝會與我們同在.
但他自己也不敢老率或老莽輕率去做事.

$^{281}$於是他跟侍應說.
不是侍應,是侍從.
他跟侍從說.
你記住.
如果對方叫我們站好.
不要動.
我們就不要動.
如果對方叫我們過來.
這就是神的印證.
證明什麼呢?.
證明上帝要把這些非利士人交在我們的手裡.
你明白我的意思嗎?.
如果對方不動.
我們就不要上去.
為什麼呢?.
因為上帝不是這個時候要我們打這個仗.
如果對方說動,上來吧.
我們就不要客氣.
就上去.
去尋求上帝的印證.
重不重要?.
很重要的.
各位頂尖的姐妹們.
在我們每天的生活裡.
每天的工作裡.
人生歷程裡真的有很多決定要做.
而且有些決定.
是很重要的決擇.
影響你一生一世.
你認不認同?.
我舉了兩個例子.
你們看到了.
是否影響你一生一世?.
如果做錯了決擇.
做錯了決定.
你就會很後悔莫及.
所以我們做任何事情.
記得.
我們要尋求上帝的心意.
不要自作主張.

$^{321}$我們要多些祈禱.
安靜.
等候.
聽聽上帝是否真的要我做這件事.
有沒有經文給我一些支持.
是否真的上帝要我做這個決定.
各位頂尖的姐妹們.
我們真的不能輕率魯莽.
我們要謹慎.
好像約拿丹那樣.
尋求上帝的恩證.
清楚了.
然後我們去做.
就一定能夠打贏這場仗.
我們再看第四個秘訣.
能夠轉敗為勝是什麼呢?.
請讀.
具信心勇氣.
是的,沒錯.
誰不知.
非利士人果然說.
你們上來吧.
來到我們這裡.
當他們這樣說的時候.
約拿丹清楚知道這是什麼.
神的心意.
這個時間.
神會將他們.
以色列人交在我們的手中.
所以拿到印證之後.
就要憑著信心和勇氣.
去到敵人的陣地那裡.
大家明白嗎?.
因為對方叫你上來.
你知不知道對方的陣地.
有什麼陷阱.
你知不知道?.
你不知道的.
對方的陣地.
有什麼地雷.

$^{361}$有沒有地雷在.
你都不知道的.
總而言之.
危機就是什麼?.
是四伏.
在這樣的情況下.
弟兄姊妹們.
當你去到一個陌生的環境.
去到敵方的陣地的時候.
其實這個信心和勇氣.
是很需要的.
如果你沒有了信心和勇氣.
你不敢走前一步.
因為你對前景.
是完全不知道發生什麼事.
對敵方的陣地.
一無所知.
所以.
我經歷到.
信心和勇氣.
對我們每一個人的行事.
是很重要的.
大家都知道.
疫情在香港最肆虐的時候.
就是在三月初的時候.
你記不記得?.
那時候每天都五萬多宗.
你有沒有留意到?.
在一個這麼危險的情況下.
竟然有一群基督徒的弟兄姊妹.
和牧師.
他們有三百多位.
大家聯署.
去登報紙.
呼籲全港基督徒.
懇切為疫情祈禱.
因為當時三月初.
大家都不知道發生什麼事.
只知道疫情天天都飆高.
飆高到什麼時候停止都不知道.

$^{401}$在這樣的水深火熱的時候.
這群牧師和基督徒的領袖.
作出呼籲.
我們發覺不是辦法.
不可以靠自己.
我們真的在最危險的時候.
要依靠神.
所以呼籲全港基督徒.
一起來懇切的禱告.
呼籲在四大報章登出來之後.
很多教會,機構,團體,小組.
很懇切向上帝禱告.
因為發覺這個形勢很嚴峻.
天天都五萬多的時候.
都不知道怎麼辦.
你明白我的意思嗎?.
在這個時候.
上帝聽不聽禱告?.
神真的聽祈禱.
當神的兒女懇切,同心祈禱的時候.
上帝就親自介入當中.
結果由3月4日的五萬宗登報紙.
到3月5日已經是三萬多宗.
一直到今天都是直線下降.
到昨天多少宗?.
都是二百多宗.
昨天還好.
有沒有人死亡?.
沒有人死亡.
我們很希望沒有人死亡.
弟兄姊妹們.
我們需要祈禱嗎?.
需要祈禱.
我們要不要懇切的祈禱?.
要.
千萬不要鬆懈.
所以我希望洪恩龐,弟兄姊妹們.
你們更多的為香港祈禱.
以致我們能夠過一個正常的生活.
也減少人命的損失.

$^{441}$所以信心和勇氣是很重要的.
雖然你不知道前面發生什麼事.
但你信得過神掌管一切.
你信得過神坐著為王的時候.
你就能夠有勇氣向前行.
正如約拿丹.
我們再看他成功的第五個秘訣是什麼.
一起讀.
預備起.
行動.
沒錯.
當約拿丹和他的侍應.
真的去到對方陣營的時候.
那個環境是不是很平坦,很方便?.
不是的.
原來敵方的陣地是很陡峭的.
很險峻的.
所以約拿丹就要什麼?.
手腳並用.
用什麼?.
爬上去.
你說這樣的環境.
是不是真的不容易向前行?.
但是約拿丹和侍從沒有放棄.
他們手腳爬爬進入到敵方的陣地.
各位聽眾.
當他們願意付諸行動.
聽上帝的話.
憑信心,勇氣向前行的時候.
果然見到奇蹟發生在他們眼前.
果然他們兩個人一把刀.
都可以殺退這些非利士人.
因為上帝在當中.
令非利士人的陣腳大亂.
各位電視節目的觀眾.
你看到嗎?.
原來我們付諸行動是什麼?.
是很重要的.
有時候當你清楚上帝要我們去做.
關懷的工作.

$^{481}$如果我們沒有去做.
就是錯失良機.
大家明白我的意思嗎?.
雖然香港第五波的疫情.
沒有人喜歡.
你和我都不喜歡.
不過既然來到.
我們發覺疫情之下.
也有些人情味出來了.
因為有人有很大的需要.
所以教會弟兄姊妹很熱心.
來關懷人.
將一些物資,物品,食物.
送給有需要的人.
為什麼呢?.
因為當時疫情很嚴峻的時候.
很多工作都開不了.
很多人失業.
很多人手停口停.
但是在這一刻的時候.
教會沒有忘記身邊的人.
教會沒有忘記社區的需要.
很感恩我知道.
洪恩堂也是不甘後人.
在最貧乏的時候.
缺乏的時候.
教會很主動,很積極地.
將物資送給有需要的人.
關心身邊的人.
雖然不認識他們.
但仍然關心他們.
這個見證重不重要?.
很重要.
這個就是我們付諸行動之後.
我相信疫情之後.
其實有很多基層的人.
都開始認識教會.
也對教會的態度,觀感改變了.
為什麼呢?.
因為在他們人生最缺乏的時候.

$^{521}$艱難的時候.
教會就去關心他們.
支援他們,支持他們.
這是很重要的行動.
如果我們不採取這些行動.
其實我們真的錯失良機.
疫情過後.
每件事都順風順水之後.
人們就忘記了.
所以我們今天.
在每天的生活當中.
我們怎樣能夠在屬靈的增進裡.
可以轉敗為勝呢?.
我希望今天我們每一位.
都記住這五個秘訣.
這五個秘訣成為了我們今天.
在屬靈的增進裡.
一個很重要的提醒.
頭主幫助我們.
祝福我們.
使用我們.
我們低頭祈禱.
慈愛的天父.
我們人生裡.
真的遇到很多的挑戰.
甚至我們都會遇到很多.
自己預料之外的一些難處.
一些困境,逆境的出現.
因為你是我們天上的爸爸.
你時刻都會看顧我們.
你也是在我們最困難的時候.
你和我們在一起.
你與我們一起去面對.
人生一些大大小小要去面對的戰爭.
叫我們能夠靠主得勝.
天父多謝你.
我們香港這個地方.
來興起很多屬你的兒女.
向你禱告,向你呼求.
天父你憐憫我們.

$^{561}$你幫助我們.
你令到這個疫情逐漸壓制下來.
我們獻上感恩.
求主讓我們每一天.
我們服事當中,生活裡.
都是緊緊的依靠你.
好讓我們的生命能夠體會到.
靠著耶穌.
我們能夠過得勝的生活.
我們的禱告奉耶穌名求.
\newpage



\section{馬可福音 16:14-20}
\label{sec:A3yUdy51iFc}
\textbf{2022.05.15 《宣道人 . 宣教人》 馬可福音 16\_14-20 (和修版) 游志豪傳道}
\newline
\newline
連結: \href{https://youtube.com/watch?v=A3yUdy51iFc}{\texttt{https://youtube.com/watch?v=A3yUdy51iFc}} ~~~~ 語音日期: 2022-05-17
\newline
\newline
\hyperref[sec:8C2XRS_C_0M]{\small{< < < PREV SERMON < < <}}
~
\hyperref[sec:index_chronic]{\small{[返順時目]}}
~
\hyperref[sec:index_scriptual]{\small{[返順卷目]}}
~
\hyperref[sec:PjWDhON898w]{\small{> > > NEXT SERMON > > >}}
\newline
\newline
馬可福音 16:14-20
\newline
\begin{longtable}{cl}
\hline
\hline
章節 & 經文 (和合本修訂版)\\
\hline
16:14 & \begin{tabularx}{0.7\textwidth}{X} 〔 後來十一使徒坐席的時候，耶穌向他們顯現，責備他們不信，心裡剛硬，因為他們不信那些在他復活以後看見他的人。 \end{tabularx} \\ \\ \relax
16:15 & \begin{tabularx}{0.7\textwidth}{X} 他又對他們說：「你們往普天下去，傳福音給萬民聽。 \end{tabularx} \\ \\ \relax
16:16 & \begin{tabularx}{0.7\textwidth}{X} 信而受洗的必然得救，不信的必被定罪。 \end{tabularx} \\ \\ \relax
16:17 & \begin{tabularx}{0.7\textwidth}{X} 信的人將有神蹟隨著他們：就是奉我的名趕鬼；說新方言； \end{tabularx} \\ \\ \relax
16:18 & \begin{tabularx}{0.7\textwidth}{X} 手能拿蛇；若喝了甚麼毒物，也不會受害；手按病人，病人就好了。」〕 \end{tabularx} \\ \\ \relax
16:19 & \begin{tabularx}{0.7\textwidth}{X} 〔 主耶穌和他們說完了話以後，被接到天上，坐在神的右邊。 \end{tabularx} \\ \\ \relax
16:20 & \begin{tabularx}{0.7\textwidth}{X} 門徒出去，到處傳福音。主和他們同工，藉著伴隨的神蹟證實所傳的道。〕 \end{tabularx} \\ \\
[1ex]
\hline
\hline
\end{longtable}
$^{1}$大家好,我是洪振英.
我想問問你們,為什麼會回來宣道會教會?.
又或者回來洪恩堂教會?.
是不是朋友叫你回來?.
又或者不經意走過,這裡有一間教會?.
不如上來看看吧.
但是為什麼是宣道會教會,而不是其他宗派呢?.
可能有很多原因.
不過我想神帶大家來到洪恩堂或宣道會教會,一定有祂的心意.
因為宣道會教會有一點點特別,有一點點不同的地方.
原來宣道會教會有宣道會自己的基因.
所謂的DNA基因在裡面.
有人說宣道會教會比較特別,因為他們很著重宣教和猜拳.
不知道大家有沒有發覺.
孫迅博士建立宣道會的時候.
其實他想組織一個宣教的組織.
一直發展猜拳的工作.
後來就成為了宣道會.
所以大家回來宣道會的時候.
都很有猜拳的感覺.
都要出去做一些宣教的服事.
當然福音傳福音也要.
或者遠方的宣教也要.
你有沒有懷著.
有宣教的DNA基因在你裡面.
你是不是一個宣教人呢.
我給大家看一幅圖畫.
這個超人就很厲害了.
不要問出處了.
十多年前有位弟兄.
就成為了宣教人.
就叫做猜拳人.
這個猜拳人就很厲害了.
你看看他左邊和右邊.
他有很多的裝備.
譬如簡單說.
他有頭盔有通話器.
他和宣教士就要保持聯絡.
他又有劍.
他又有消費卡.

$^{41}$消費卡很重要.
因為猜拳需要金錢.
所以定期要奉獻猜拳基金.
還有他的鞋子也很重要.
因為我們要去宣教.
要去探訪宣教士.
走到腳底也很遠.
要有對好的鞋子去保護.
我們在想.
一個猜拳人.
你是不是一個猜拳人呢.
你衡量一下自己.
你是不是一個猜拳人.
你有沒有頭盔又有劍又有鞋子.
常常保護你.
常常做足所有的事情.
我在想可能會難一點.
有些就已經是剛回來.
有些就做不足這麼多的時候.
我不是猜拳人.
我在想猜拳人是不是一定要做足所有的事情.
你才算是做猜拳呢.
當然答案是不是的.
為什麼呢.
今天我們會透過經文.
去看看一群門徒.
他們是怎麼樣.
想想這個門徒.
他們究竟有什麼特質在裡面呢.
當主耶穌升天之後.
我們知道一群門徒聚集.
他們的聖靈降臨在那群門徒身上.
他們心裡火熱.
四處傳道.
四處傳福音.
但是那群門徒是什麼人呢.
你想想.
耶穌那時候去到湖邊.
呼召彼得.
彼得也是打魚.

$^{81}$馬太也是碎利.
其實那群門徒.
就是一群烏合之眾.
什麼人都有的.
組合在一起.
不過當我們看剛才的經文.
你也看到那群門徒很厲害.
到處傳揚福音.
心裡火熱.
宣教人.
一定是宣教人.
但是為什麼他們能夠這麼厲害呢.
有三個重點.
今天拍網我講完之後.
記得三個重點就可以了.
他們有.
信心.
第二.
他們有行動.
他們不是坐著.
他們是出去的.
第三 有聖靈.
挑撥他們.
信心 行動 聖靈.
拍網我們也有這三樣東西.
你們逐一去看一看.
我們先讀一讀經文.
這是另一個版本.
我們一起去讀.
預備起.
不久之後.
耶穌又在.
十一個門徒一起吃飯的時候.
向他們顯現.
責備他們.
靈魂不信.
因為他們不相信.
那些在他復活以後.
見過他的人.
他們不相信.

$^{121}$見過他的人.
他又告訴他們.
你們要到世界各地去.
向所有人.
傳揚福音.
凡是相信.
並接受洗禮的.
必定得救.
拒絕相信的.
必被定罪.
這些神跡會伴隨著相信的人.
他們要供我的命.
趕鬼.
說新的.
軟弦.
能用手捉蛇.
而安然無恙.
即使喝了有毒的東西.
也不會受傷害.
按手在病人身上.
病人就康復了.
主耶穌.
向他們說完.
這些話.
就被接到天上去.
坐在上帝的右邊.
永遠不離.
上帝的右邊.
榮耀的位置.
門徒到處傳揚福音.
主直直.
在他們而作作.
並用許多神跡.
來證實他們所傳的道.
好 謝謝.
剛才就是.
新普及的版本.
他們就用另外一些的句語.
去演繹了.
那幾段的經文.

$^{161}$難不難呢 不難的.
都是一個故事.
講到門徒.
如何出去傳福音.
短短幾節的經文.
講到.
耶穌向門徒顯現.
叫他們出去.
後來.
門徒真的出去.
做福音的工作.
還有許多神跡.
顯現在他們的生活.
裡面.
我們這裡就看到.
三個大的主題.
包括信心 行動.
還有.
忘記了 聖靈的工作.
我們先看一看.
第一個.
信心的基礎.
哪幾個經文這樣說呢.
我們看到在第十四節.
主耶穌.
和他們說什麼呢.
祂說祂責備他們.
說他們信心不足.
不過經文也很有趣.
到了最後的第十九節.
耶穌.
被接上升天.
坐在神的右邊.
這裡也是強調.
使得門徒有信心.
能夠確認.
他們所信的.
真是拯救他們的神.
第二.
行動在哪裡呢.

$^{201}$行動在第十五第十六節.
耶穌和他們說.
他們要出去.
天下傳福音.
經文的第二十節.
又說得很有趣.
門徒真的能夠出去傳福音.
第十節上面說.
門徒出去到處.
傳福音.
經文是這樣說.
第三個就說到神蹟.
神蹟就說到聖靈的神蹟.
騎士.
第十七第十八節.
耶穌和他們說 信的人就有神蹟伴隨.
第十節門徒.
出到外面的時候.
有沒有神蹟呢.
經文說有神蹟伴隨著.
這群門徒.
我們逐點去看看.
今天讓我們去思想.
學習.
當中.
信心的基礎.
第一就說到信心是什麼.
後來十一個使徒.
坐直的時候.
耶穌就向他們顯現.
責備他們心裡.
不信.
因為他們不信.
那些在他復活之後.
看見他的人.
糟了.
耶穌責備那群門徒.
那群門徒你們真的.
信心少了.
或者不信了.

$^{241}$有人見過耶穌復活.
跟著跟他們說.
但是他們信不信.
所以耶穌在當時.
責備他們信心不太好.
心裡剛硬.
堅固.
頑固.
不信.
這就是耶穌所責備他們.
其實在整個馬可福音.
由頭至尾.
都有說到那些門徒.
信心不太好.
不知道大家信心信得好不好.
不知道大家信了這麼久.
或者回教會這麼久.
面對著耶穌的時候.
信心不足.
還是信心很強.
不過在其他的.
福音書其實都有.
罵這群門徒.
不知道大家有沒有記得.
大使命.
說什麼.
大使命馬太福音第28章.
當一開始的時候.
耶穌責備那群門徒.
心裡不信.
他們心裡仍然有疑惑.
馬太福音就說.
那群門徒心裡有疑惑.
可能他們疑惑.
就是耶穌復活了.
究竟是不是真的.
究竟有沒有這件事.
他們就心裡去疑惑.
不過我想問一個問題.
當那群門徒.

$^{281}$心裡仍然疑惑的時候.
仍然不信的時候.
或者信得不好的時候.
主耶穌有沒有.
不頒佈使命給他們.
似乎沒有.
當那群門徒還沒準備好.
還沒做得很好的時候.
主耶穌已經.
將大使命.
馬上頒佈給他們.
可能他們還沒做得太好.
不過主耶穌都給了他們使命.
叫他們出去做.
很有趣的.
有時候我們在大學的時候.
我們會有一個.
大使命.
很有趣的.
有時候我們猜拳.
傳福音.
有些人就覺得我還沒做到.
我剛剛回來教會.
我剛剛初信.
我還是不做傳福音了.
自己就在猜拳.
去不同的教會.
做一些猜拳教育.
聽說.
有些教會.
靜悄悄說.
有些宣導會.
都說不做猜拳.
不知道大家知不知道.
有些說我不做猜拳.
或者說我很少做猜拳.
其實.
有幾個原因.
他說我們堂會很小.
堂會只有20人.

$^{321}$30人做不到猜拳.
又有些說.
我剛剛才成立堂會.
不做猜拳.
有些就說.
我們本地.
本地都還沒做完.
傳福音的工作.
怎麼做猜拳的工作呢.
有些會這樣說.
不多的.
但是有這樣的教會.
他們覺得不做了.
不過當然.
我相信我們不是這樣.
主耶穌當我一群的門徒.
還沒信得很好的時候.
就已經頒布了.
這個使命給他們.
叫他們去做猜拳的工作.
在這裡.
很明顯說到.
一個信心的問題.
那群門徒不信.
第一個要提醒我們.
就是一個信心的問題.
在裡面.
剛才也提到了.
在經文第19節就說.
主耶穌和他們說完話之後.
被接上天.
坐在神的右邊.
作者這樣說.
雖然那群門徒.
信心不足.
但是你想想當時的情景.
他們信心不好.
但是他們看到主耶穌.
被接.
升上天的時候.

$^{361}$他們能夠確信.
耶穌再一次告訴他們.
祂就是拯救的主.
你能夠將你的生命.
交給祂.
主耶穌在神的右邊.
有祂的權能.
有祂的大能和權柄.
能夠翻回祂的身份.
回復祂的地位.
那群門徒是眼睜睜.
看著主耶穌基督.
得著.
升天.
回到祂的位置.
對他們來說.
信心是一個很大的提升.
和幫助.
今天.
宣教是甚麼呢.
我自己也說過.
我自己做宣教或猜拳.
的工作.
究竟宣教是甚麼呢.
有些人會說.
如果問神那麼厲害.
其實不用我們做宣教.
他派天使.
就可以做到.
宣教的工作.
或者他做很多奇蹟.
就可以做到宣教的工作.
不用堂會.
不用我們參與.
不過有位宣教士.
提出過一件事.
神給了我們一個很寶貴的機會.
去參與.
猜拳.
讓我們能夠.

$^{401}$和神一起同工.
其實當我們.
做宣教或猜拳的工作.
不知道大家有沒有試過傳福音的工作.
但是當你.
傳福音給一個.
朋友或家人的時候.
當他信主的時候.
你會怎樣.
你會很開心很喜樂.
你不知道為甚麼會這麼開心.
他信主.
不關你的事.
不過你心裡會覺得很開心.
就算你去到.
短宣.
有一個可能是不同文化的人.
可能跟你沒有甚麼關係.
但是他信了主.
你心裡也會很大的開心.
很大的開心.
主要我們能夠參與.
猜拳.
其實是同樣分享他這一份.
喜樂的內容.
讓更多人能夠認識他.
萬族萬民.
都可以敬拜他.
今天.
如果你問你對神的信.
有沒有當時的門徒.
這樣不信.
還是充滿了所有的力量.
在裡面.
你很相信神.
他給了你意象.
給了你盼望.
給了你去做甚麼.
這個不知道大家.
之前有沒有看過.

$^{441}$都是說到我們開創人.
就是這個.
孫迅博士.
當時原來宣道會.
講一下歷史.
宣道會在初期.
其實宣道會博士.
他也有一個很特別的意象.
他的意象.
就是說他在晚上.
他曾經有一晚.
做夢.
他在裡面.
醒了.
他心裡很熱.
很火熱.
心裡有一種很特別的感覺.
他不斷地.
想著夢的裡面.
他就說.
他看到有一個很大的.
教會.
他坐在後面.
他眼前看到很多不同的人.
很多不同的人.
他說有幾百萬的人.
在裡面.
他說看起來很像.
中國人.
一個西方的宣教士.
但是他做夢.
他夢見裡面是中國人.
他說那些中國人沒有出聲.
一聲也沒有.
但是很痛苦地.
挽著自己的手.
很痛苦地坐在那裡.
他不能忘記他們的表情.
然後.
孫迅博士後來醒了.

$^{481}$他沒有想過.
是中國人.
他沒有想過神給他的.
意象是中國人.
他說聖靈感動他.
他就跪下.
他就跟神說.
神啊我願意去.
一個西方的宣教士.
神給他意象.
去中國.
做宣教的服事.
剛才說.
第一個宣教的DNA.
是信心.
是我們的信仰.
在我們個人的生命裡面.
你信神.
信成怎樣呢.
你信得很堅固.
還是很虛浮呢.
神有沒有給你.
一個意象.
好像孫迅博士一樣.
要你參與在.
猜拳裡面.
為他去做一點事情呢.
我記得有一個.
在一個大會的工場.
有一個故事.
有一個姊妹.
大家聽不聽得到.
她參與在短孫裡面.
沒什麼特別.
很多人都會去短孫.
近兩年就去不到.
但她過去就去短孫.
這個姊妹.
很特別.
她特別的地方就是.

$^{521}$她已經70歲.
這位70歲的.
姊妹.
她參與在緬甸的短孫.
她去緬甸.
當然是落後一點的地方.
她70歲.
一直每年去一次短孫.
去到80歲.
去了10年的短孫.
當時.
工場跟我們分享.
這個很厲害.
當然她的身體可以.
她的精神可以.
如果不是就去不到.
但她有這樣的魄力.
有這樣的意象.
神叫她能夠參與.
她就參與在短孫裡面.
今天可不可以呢.
可以的.
我們信神也帶領我們.
參與在主的.
福士裡面.
可能神叫你做一點點的事.
但我們也可以參與在當中.
裡面.
神給你什麼意象呢.
什麼宣教的意象呢.
有沒有宣教的心.
或者傳福音的心.
放在你的心的上高.
以致你很想身邊的人.
去認識神呢.
這是第一樣.
第二個就說到.
行動的裡面.
剛才說到行動是什麼呢.
那些門徒.

$^{561}$真的出去傳福音.
他說.
耶穌對他們說.
你們往普天下去.
福音給萬民聽.
信而受洗的必然得救.
不信的必被定罪.
這個經文呢.
剛才說到這個故事.
來到這裡的時候.
讓我們想到另一個.
就是宣教的行動.
猜拳的行動.
裡面.
有些人說這個.
馬太福音28章大使命.
所說的就是.
你們要去.
萬民作我的門徒.
是吧.
奉主的名和他們施洗.
其實這裡說到.
普天下萬民.
聽福音呢.
萬民就是所有的人.
不單是你身邊.
那個.
記住喔,不單是你認識的那個.
甚至你還沒認識.
和你不同文化的人.
其實都是說到萬民.
所以我們都說宣教.
說猜拳的裡面.
其實是說遠一點的.
不單是自己的範圍.
不單是這個社區.
再遠一點的萬民.
我們都要關心.
我們未必做到那麼多.
不過我們能夠去關心.

$^{601}$不過這個經文同樣告訴我們.
我們要去.
不是你們要停.
不是你們要回來.
而是你們要出去.
出去是一個行動.
不是坐在那裡等.
而是真的走一點.
做一點.
去看看神怎樣開路.
所以在宣教的.
DNA或者宣教的.
元素,第二就是.
一個行動是很重要的.
路加福音.
這樣說,人呢.
人要奉他的名.
全說的是悔改.
赦罪的道.
信的就必然得救.
不信的就必會被定罪.
好像這裡這樣說.
信而受洗必然得救.
不信的必被定罪.
所以今天我們出去.
傳福音的時候.
其實我們的信息是很重要的.
我們告訴他耶穌的時候.
他如果不信的.
他就已經被定罪.
但是那個人.
如果他能夠.
相信主的時候.
他就會得到那個拯救.
所以今天.
我們每個人都有很重要的.
角色在裡面.
去傳講這個.
很重要的福音.
羅馬書說如果沒有人.

$^{641}$去傳.
那裡會有人聽得見呢.
天父.
怎樣猜險主耶穌.
主耶穌也怎樣.
猜險我們出去.
所以今天.
耶穌是猜險我們出去.
出去可能很遠的地方.
有些宣教士去到很遠的地方.
你想像也想像不到那麼遠.
但是有些可能是近一點的.
國內.
或者東南亞.
甚至再收窄一點.
在你的範圍之內.
你的社區.
上帝也是同樣猜險你.
去到身邊.
和身邊的人去傳福音.
裡面.
剛才那個經歷.
剛才那個經文.
第20節不知道大家還有沒有印象.
提到那些門徒.
耶穌和他們講完.
傳福音之後那些門徒有沒有出去.
那些門徒.
有行動.
那些門徒.
門徒出去.
到處傳福音.
那些門徒其實不是.
再坐在那裡.
望天打卦 當時他們門徒.
祈禱.
祈禱之後是不是還坐在那裡.
等呢.
不是.
在萬惡福音所提到的.

$^{681}$到處傳揚福音.
隨處隨傳.
他們去到哪裡的時候.
就講耶穌.
離開了加利利 離開了耶路撒冷.
他就去傳揚.
福音 見到什麼人.
他就繼續去講.
他的福音裡面.
門徒是這樣.
我們是坐在那裡.
這裡真是很漂亮.
又有冷氣又舒服.
還是我們都是願意.
走出去.
去到身邊的時候.
關心身邊的人.
將福音帶給他們.
這個地方.
不知道大家認不認識.
我不認識這個地方.
我在網上找來的.
這個聖奧斯本.
在英國的地方.
和我們有個英國的宣教士.
有些拉姦.
因為我們在英國有一個宣教士.
他去到倫敦的地方.
做宣教.
話說這個宣教士.
看到就很想去做福音的工作.
他就和.
在一個代討信這樣說.
你知不知道原來有很多人.
香港有很多人去了英國.
他說用BNO申請去英國.
2021年就有12至15萬人.
去了英國.
預計會有差不多30萬人.
去到英國.

$^{721}$很大.
很多人去到英國.
他就說他和兩位宣教士.
在倫敦附近不斷去.
看看哪些地方.
可以做福音的工作.
他開車的時候.
開車到這個地方.
聖奧斯本的地方.
他就說看看.
這裡究竟有多少.
香港人呢.
他就問一問.
原來在這個地方有200個.
香港的家庭.
在這裡住.
200個香港的家庭.
如果一個人是4個人的時候.
就說8100.
有一個很大的數字.
香港來的人.
去了英國.
他就在想.
問當地的人.
在這200個家庭裡面.
大約有50個是信主的.
即是說.
起碼有150個家庭.
還有很多福音的機會.
裡面有很多.
有很多新移民.
有些大學區.
所以很多人就會來這裡.
每年都很多人.
但是那個宣教士再看看.
有沒有教會呢.
有英國人的教會.
但是一間華人的教會都沒有.
不知道你聽到的時候.
如果是你會在想些什麼.

$^{761}$不過宣教士就會想到一句經文.
他就說要收的莊家多.
作供的人少.
這句我們都懂.
但是我們會不會套入了.
去這裡看.
宣教士就心裡很迫切.
就會覺得150個家庭.
是不是應該在那裡.
去做福音的工作呢.
直堂.
或者是一些小組.
那些童工就在那裡.
逐家拍門去做一些福音的工作.
不過最感動我的就是.
他寫了幾句說話.
他說說真的.
在這個地方.
他一個人都不認識.
他說真的不知道從何著手.
唯有在禱告裡面.
求主人.
求弟兄姊妹紀念.
我聽到都有點唏噓.
很想去做福音的工作.
但是又不知道怎樣做.
連一個人都不認識的時候.
唯有可能逐家拍門.
不過宣教士就看到.
作供的人是少.
但是裝嫁卻是多.
宣教是需要行動的.
當我們看到的時候.
其實就不是坐在那裡.
而是真的出去去接觸.
今天在我們的宣教生命.
或者在我們的人生裡面.
我們常常都.
我們是不是馴服神呢.
馴服神的時候.

$^{801}$馴服神的帶領 他的吩咐.
讓我們走出去.
去做一些關懷的工作.
我們都願不願意.
突破自己的限制.
為主去做一些工作呢.
不知道這個社區是怎樣.
我們可以想的時候.
宣教的工作.
我相信團天區有很多.
少數民族.
或者是一些不同文化的人.
這些其實我們都可以去關心.
至少我們都可以.
為他們去禱告.
因為現在不用去到這麼遠的地方.
我們都知道他們已經來到.
我們家門口裡面.
我們自己教會.
就有一個很特別的印尼團契.
印尼的姐姐.
通常我們就會.
星期日會去到那些地方.
搞一些活動給他們.
不過我們教會.
搞這個就在星期五的.
晚上 印尼團契.
你問問那些姐姐.
你怎麼回來 星期五晚上.
不是要煮飯 洗衣服嗎.
原來呢.
我們就和那些基督徒的.
僱主說.
你都很想你的印尼姐姐.
能夠信主.
不如你星期五晚上.
放他們假.
你要做一個有使命的僱主.
當然僱主聽到.
是喔 真是很想.

$^{841}$那些姐姐信主.
可能煮飯.
就放他們假.
他們就七點八點來到教會.
有小組 有學習.
雖然我們不是.
因為教會裡面基督徒的僱主不多.
所以都不是很多姐姐.
四五個 五六個.
不過我們維持了.
這個小組就幾年的時間.
都很感恩.
因為有基督徒的僱主.
他們有心.
給他們的姐姐.
放他們假回來.
都是一個猜拳的行動.
你放你姐姐的假.
其實就是讓他們.
可以去接觸信仰.
今天有什麼我們可以做.
或者我們做到的呢.
未必去到那麼遠.
未必走出去.
去到其他地方.
但可能已經在我們家裡.
都可以做到一些猜拳的服事.
不過主題我們就是要行動.
不是坐在這裡.
第三個就是神跡的伴隨.
第十六和第十七節這樣說.
信的人將會有神跡隨著他們.
就是奉我的名.
趕鬼 說語謊言.
手能拿蛇.
如果喝了什麼有毒的.
都會康復的.
很有趣.
剛才說到要出去傳福音.
但突然間經文又說.

$^{881}$你們會見到 遇到神跡.
告訴我們一個很特別的提醒.
在宣教傳福音猜拳裡.
不是經常做的.
我們都知道我們要做很多事.
不過他告訴我們.
我們都會經歷神跡.
經歷聖靈的帶領.
經歷神和我們同在.
那份甘甜.
神跡是一個sign.
即是一個表示.
表示到我們在神裡面.
正在做他的工作.
剛才說到說謊言 拿蛇.
其實很多已經在思路行轉.
出現過的神跡奇事.
可以這樣說.
其實今天我們宣教也有很多神跡奇事.
在香港比較少聽.
但當我們聽宣教的時候.
我們會聽到很多落後的地方.
或者部落的地方.
很多時候他們因為醫藥比較落後.
有宣教士告訴我們.
他有病 但沒有藥醫.
他唯有怎樣.
他說我都不知道怎樣去應付.
就是給他兩粒止痛藥就算了.
然後回去就說.
神啊 醫不到 因為沒有藥.
不過求你彰顯你的神跡.
很感恩.
他吃了兩粒止痛藥.
又康復了.
他們在缺乏的地方.
正在經歷神的醫治.
有很多神跡的奇事.
盧家福音說.
一群門徒等候聖靈的能力.

$^{921}$聖靈是宣教的靈.
聖靈落入他們的心.
他們就火熱.
他們就做起宣教的工作.
所以聖靈在我們心裡.
催促我們去傳福音.
我們心裡覺得很火熱.
很想去做福音工作.
因為聖靈在我們心裡.
催促我們.
所以第三個元素.
這裡說到神跡.
我就用另一個字眼.
就是聖靈的工作.
聖靈的大能在我們心裡.
第三個DNA就是.
有聖靈在我們心裡.
和我們同行.
那群門徒.
能不能經歷聖靈呢?.
在剛才20節的最後.
不知道大家還要回想一下.
想想我的經文.
當耶穌說你出去.
做福音的工作.
那群門徒就出去了20節.
他們就說主和他們同工.
藉著伴隨的神跡.
證實所傳的道.
那群門徒正在經歷神跡.
神跡不是主角.
也不是焦點.
神跡是告訴他們.
他們正在做神的工作.
神和他們同行.
證實他們所傳的道.
今天我們信不信神跡呢?.
在香港很文明很理性的.
我們比較少.
但是今天仍然有神跡的存在.

$^{961}$每一天我們都會經歷很多神跡.
不知道大家有沒有短孫的經驗.
可能兩年前或是多少年前.
都有短孫的經驗.
話說有幾次去泰國.
柬埔寨短孫.
我們幾十人坐一輛大旅遊巴.
一上車的時候.
不知道大家有沒有試過.
那一次坐一輛孫教士開的車.
一上去的時候.
不知道大家有沒有印象.
孫教士會做些什麼呢?.
又聽到祈禱.
然後孫教士說停一停.
我們先祈禱.
然後我們想我們繞個彎.
走兩個彎就停.
不用祈禱.
保守安全.
真的兩個彎就到了.
下次回到家.
又停一停.
又先祈禱.
我想大家如果有短孫經歷過.
孫教士開車.
或者是旅遊巴.
開車之前.
孫教士就說先不要開.
先祈禱.
為什麼會這樣說呢?.
我有時都會問.
為什麼每次都祈禱.
巴士有沒有祈禱.
沒有.
我們很少.
因為香港的危險意識不太高.
但是在工場裡.
聽他們說.
他們可以很危險.

$^{1001}$他們開著車.
可能電單車在旁邊彈出來.
就會撞車.
可能會發生很多意外.
尤其是長途車的時候.
可能會有很多意外.
在佛法裡.
他們常常提醒自己.
真的要神保護.
我們祈禱.
保護整個旅程.
在他們裡.
在工場裡.
孫教士特別.
經驗到.
神和他的同在和神蹟.
不是說香港.
香港比較少.
因為我們比較沒有這樣的需要.
或者是我們.
有時候工作很繁忙.
但是在工場裡.
他們真的這樣做.
今天我們會嘗試.
每天都經驗.
神的神蹟奇事.
聖靈的大能.
這是真實的.
因為他們這樣說.
你做主的工.
神的神蹟就會證實你全的道.
會伴隨著你一切.
這裡就是.
一個叫恩園的地方.
是警察會.
有一個愛主的孫教士的家庭.
把元朗.
元朗錦田.
有一幅地.
大約幾萬呎的地.

$^{1041}$農地.
給了警察會.
我們在想.
但是可以做些什麼呢.
過往在社運.
下神都很保守.
警察會的奉獻.
都不差很遠.
我們感恩.
我們面對著有這個地.
給了一幅地.
這是舊的建築物.
我們怎樣去用呢.
我們打算建一個.
孫教中心.
還有一個孫教的驛站.
可以給孫教士休息.
在這個地方重建.
但是原來在香港.
重建房子是很難的.
我不知道大家有沒有.
這樣的經驗.
因為它是農地.
所以你要和政府商討.
轉變它做建屋.
又要改地契.
又要補地價.
很多很多事情.
我們就和一些顧問公司開會.
他們就幫我申請.
轉變規劃.
進去這個規劃處.
規劃處有一個職員.
和我們說.
其實他們都不太知道.
我們申請做退休中心.
他們都不太知道.
什麼叫退休中心.
可能他們不是基督徒.
他們不懂這些字眼.

$^{1081}$你是不是做住宅.
是不是做tong 房.
是不是拿出來將來就租出去.
我們就說不是.
我們就想這次真的難搞.
他們就說其實都難一點.
都是當你住宅.
我們心裡都想.
都想很難去成事.
唯有交託給神.
去禱告.
如果是神的開路.
他就會帶領.
去到尾的時候是批准了.
馬上說到尾批准了.
現在就差不多.
做一些入職的地方.
但那個情況是什麼.
當我收到問那個顧問公司.
為什麼後來會批准了.
因為疫情的關係.
退後了一段時間.
拖拖拉拉.
原來原先在這裡.
做規劃的官員.
因為政府會調的.
那個男生過了兩年就調走了.
調走了另外一個女的官員.
他是基督徒.
真的.
然後就.
因此我不知道是不是.
我想都是公正的.
但他至少明白.
什麼叫退休中心.
輾轉之下.
批准了.
那個顧問公司.
他都不是基督徒.
都很奇妙.

$^{1121}$因為都不是很多例子.
當然可能地方不是大.
我們申請那裡都很小.
所以他都會批准.
我想很多時候.
在工場或者在我們生命的點滴.
都會發現了.
原來神真的有神跡和我們同行.
不過視乎我們怎樣去相信.
我們是否仍然都相信.
主和我們同行.
當我們.
一路去做神的工作的時候.
聖靈和我們同工.
證實祂所行的道.
說了三樣東西.
信心.
行動.
和聖靈.
其實是三個很連著的.
一件事.
我們對神有信心.
堅固我們的信心.
我們不是停在這裡.
我們是行動.
去做祂的福音工作.
去做猜拳的工作.
當我們一路行動的時候.
聖靈就作供.
讓我們正在經歷.
神在我們生命裡的真實.
當我們越經歷的時候.
我們的信心又會增強.
我們繼續不斷不斷不斷.
去運神.
在我們心裡.
剛才提及那群門徒.
其實都是烏合之眾.
耶穌罵他們.
這群不信的門徒.

$^{1161}$靈還不靈.
不過.
經文說他們最後很厲害.
火熱到處.
傳揚福音.
到處有神跡伴隨.
你說他們不信.
但到最後.
他們卻是很厲害的一群門徒.
其實宣教人.
就是一群相信的人.
我們信主的.
其實已經是一群宣教的人.
信的就是宣教的人.
我們有這樣的信心.
繼續去行.
我們就正在經歷神的大靈.
我們盼望.
同樣地.
我們同樣地.
對神有信心.
堅固地繼續去行神的道.
我相信.
神的奇蹟.
一直發生在堂會裡.
個人也好.
堂會也好.
不斷經歷神.
不斷帶領你們.
去經歷真實的神.
當然日後的困難.
越來越多的挑戰.
不過我們都能夠繼續去依靠主.
最後.
我都有一個影片.
給大家看.
關於剛才的地.
先看看.
宣教士在前線工場.
征戰連年.

$^{1201}$身心疲累.
因此回港述職的時候.
實在需要一個休息.
重整.
重新得力的歇息居所.
2013年.
宣教士在前線工場.
征戰連年.
身心疲累.
2013年.
有一個外族的家庭.
捐出他們的先父母.
留給他們位於錦田的農地.
給警察會.
警察會後來議決.
把這塊地皮發展成多功能宣教中心.
當中包括.
退休宣教士居所.
退休場地.
訓練中心.
以及跨文化教會等等.
在諮詢完政府各部門之後.
由於這塊地皮.
暫時仍然沒有足夠的配套.
去做支援.
所以發展得非常困難.
2018年.
我們打算分拆.
這個多功能宣教中心部分的功能.
先發展其中的一個部分.
地段287.
作為整修用途.
這個細小的部分.
有一間名為恩園的舊樓房.
這個樓房.
昔日就在翁立德.
和他的父母住的.
這個舊樓房.
昔日就用來做報道所.
和查經的地方.

$^{1241}$並向村民分享福音.
雖然今天只是一間破舊的樓房.
但是感恩.
日後將會是我們回港宣教士的.
中途驛站.
一個為宣教士加油.
給宣教士整合得力的家園.
2020年.
我們得到神的印證.
在重重困難之下.
政府委員會批准我們改變.
土地用途的申請.
在有附帶條件之下建設退收中心.
現在.
這個退收中心.
已經正式命名為靜修恩園.
我們和工程顧問公司.
已經簽訂合約.
相信可以盡快.
進行補地價的示意.
以及展開建築的工程.
靜修恩園.
一個宣教士的靜修中心.
在面積5000多平方呎.
兩層樓的建築裡.
包括了退收房間.
一個跨文化敬拜禮堂.
淑齡導引室.
輔導室.
以及多用途活動房間.
祈禱室等地方.
我們相信.
這個靜修中心.
可以為宣教士提供一個退收安靜的地方.
成為宣教士的驛站.
一個整合得力的家園.
以致.
他們可以安靜得力.
重新展翅上騰.
弟兄姊妹.

$^{1281}$靜修恩園整項工程費用.
總共需要.
籌募港幣六千多萬.
當中包括工程顧問費.
建築費.
行政費.
以及必須繳交給政府的費用等等.
我們誠邀你的參與.
一同共建.
我們一同禱告.
主佑天上的阿爸父.
讓我們能夠經歷您自己.
也能夠信心堅穩.
能夠出去.
與主一同做福音的工作.
盼望近處遠處.
都得著福音.
在這個困難的時代.
求您更加堅固我們.
帶領我們 讓我們見到異象.
讓我們能夠繼續與您同行.
我們禱告 奉主耶穌基督的名字.
其他事項.
\newpage



\section{馬太福音 7:21-23}
\label{sec:PjWDhON898w}
\textbf{2022.05.22 《我從來不認識你們》 馬太福音 7\_21-23 (和修版) 曾敏姑娘}
\newline
\newline
連結: \href{https://youtube.com/watch?v=PjWDhON898w}{\texttt{https://youtube.com/watch?v=PjWDhON898w}} ~~~~ 語音日期: 2022-05-25
\newline
\newline
\hyperref[sec:A3yUdy51iFc]{\small{< < < PREV SERMON < < <}}
~
\hyperref[sec:index_chronic]{\small{[返順時目]}}
~
\hyperref[sec:index_scriptual]{\small{[返順卷目]}}
~
\hyperref[sec:5T528Sf9V9Y]{\small{> > > NEXT SERMON > > >}}
\newline
\newline
馬太福音 7:21-23
\newline
\begin{longtable}{cl}
\hline
\hline
章節 & 經文 (和合本修訂版)\\
\hline
7:21 & \begin{tabularx}{0.7\textwidth}{X} 「不是每一個稱呼我『主啊，主啊』的人都能進天國；惟有遵行我天父旨意的人才能進去。 \end{tabularx} \\ \\ \relax
7:22 & \begin{tabularx}{0.7\textwidth}{X} 在那日必有許多人對我說：『主啊，主啊，我們不是奉你的名傳道，奉你的名趕鬼，奉你的名行許多異能嗎？』 \end{tabularx} \\ \\ \relax
7:23 & \begin{tabularx}{0.7\textwidth}{X} 我要向他們宣告：『我從來不認識你們，你們這些作惡的人，給我走開！』」 \end{tabularx} \\ \\
[1ex]
\hline
\hline
\end{longtable}
$^{1}$我們開始的時候,再有一同禱告的時間.
給我們這麼卑微的人.
可以去親近你.
去敬拜你.
天父,當這一刻.
我們是來臨在你的面前的時候.
求求你.
給我們更多的機會.
讓我們更多的機會.
去親近你.
去敬佩你.
在你的面前的時候.
求你光照我們的生命.
讓我們看到.
在我們的生命當中.
我們有什麼地方是得罪你的.
是得罪人的.
願你親自寬恕潔淨.
幫助我們守潔心清.
去親近你.
這一刻,願意聖靈.
充滿了我們每一個的心.
對我們每一個說話.
願意聖靈.
去幫助我們.
讓我們的耳朵.
真的張開.
眼睛是打開的.
我們看到上帝,我們聽到上帝的教導.
然後我們的心.
就樂意作出回應.
將以下的時間交上.
奉耶穌基督的名字.
祈禱,阿們.
這一段經文.
應該是比較少講的.
一段經文.
在當中是一個很重要.
很重要的一個.
一個願勉.

$^{41}$甚至乎是一個警告.
我們來看看這句說話.
這句說話很特別.
我從來不認識你們.
什麼叫從來呢.
不是現在我不認識你們.
以前認識你們.
現在就不知道為什麼不認識了.
或者忽然之間不認識了.
不是,是我從那時候到現在.
一直以來我都不認識你們.
這是一個很深的.
一個拒絕的表達.
我完全.
完全的不認識你們.
由以前到現在.
都是.
是誰說的呢.
我聽到.
有些聲音,主耶穌.
是耶穌說的.
這句說話是耶穌說的.
出自耶穌的口.
是不是很特別呢.
那個背景.
就是整個馬太福音.
七章十三到二十七節.
這裡有很豐富的內容.
但是.
它前後的經文.
是非常吻合.
很配合.
今天那段經文.
前後究竟在說什麼呢.
這個像是什麼呢.
大家看下去.
這個猜不到.
這個是閘門.
很小很小的.
一道門,閘門.

$^{81}$原來經文七章十三節.
那裡一直說.
走進入.
閘門,是走閘路.
這裡是說.
一條很特別的路.
是閘的門,是閘的路.
才是通往永生的.
反而那條.
是通往滅亡的路.
是很闊的.
那道門通往.
滅亡的門也是很闊的.
進去的人也多.
進去閘門的人.
卻是少的.
這是開始的那段內容.
接著之後.
這個是什麼呢.
披著羊皮的狼.
是不是.
是一個狼.
說的是假先知.
經文提到.
要防備假先知.
他們披著羊皮.
我們不是很分辨到.
以為是羊,實際上是狼.
他進入羊群當中.
無非就是.
為了傷害羊,擔就羊.
讓羊.
走在歪的道路上.
作者叫我們要去.
這個其實.
是耶穌基督說的.
叫我們要去防備假先知.
你由一棵樹的果子.
去分辨出.
那棵樹是一棵好的樹.

$^{121}$還是一棵不好的樹.
你由果子.
去看出.
那個究竟是真先知.
還是假先知.
我們要小心防備.
接著就到我們.
今天這段經文.
今天這段經文就是說.
哪一個可以.
進入到天國.
裡面去呢?.
其實我們表面上是在問.
哪一個可以進去呢?.
實際上耶穌提出了一個很大.
很大的警告.
很大的一個勸勉.
不要以為.
今天我按照.
我的理念,我就一定.
進去,我們一會兒再看這一段.
這一段你看上面.
一直在說,究竟作為一個.
基督徒,耶穌的.
跟從者是怎樣生活呢?.
你要走紮路.
進入紮門.
還是走向滅亡呢?.
真正跟從耶穌的人.
怎樣選擇?.
跟從耶穌也要懂得.
去防備那些假先知.
要堅心.
堅持去跟從耶穌基督.
還有要小心.
真的要清楚想清楚.
我是否那個可以進天國的人.
之後呢.
他再說的時候.
有很特別的.

$^{161}$有兩間屋.
這個是蓋房子的比喻.
這個比喻裡面說明了.
兩種生命的根基.
今天我們怎樣選擇.
究竟是我聽了神的話語.
去進行.
還是我聽了神的話語,我不走.
前者一定是.
聽了神話語去走.
是很好的.
這個房屋是非常穩固.
後者就說.
我聽了神話語,我不會走.
聽就聽吧.
我不會進行.
這時候我們的根基.
是非常薄弱.
那間屋是很容易倒塌的.
耶穌基督也是問我們.
怎樣選擇?.
你想怎樣選擇?.
究竟今天.
耶穌不停不停.
在問.
你打算怎樣選擇.
你的人生.
你想.
去走向永生.
去得著永生.
還是走向滅亡.
你將你人生的根基建立在哪裡?.
是上帝的話語.
還是總之聽了都浪費力氣.
我不會跟的呢?.
你肯定.
自己可以進天國?.
你說.
為何說一堂這樣的道?.
我回教會這麼久.

$^{201}$信耶穌華洗二十年.
三十年,還跟我說這些?.
不是喔.
兄弟姐妹.
這堂道不是對著.
未信主的人說的.
是要對著信主的人說的.
我今天就要問你,怎樣選擇?.
是一個選擇.
你怎樣選擇是有後果的.
甚麼人可以進天國?.
就是今天想說的.
在經文裡面.
看到.
七章二十一節那裡.
很明顯說到有不少人.
是會說.
主啊,主啊,經常提到.
耶穌基督是我主啊.
你回教會信誰?.
當然是耶穌.
耶穌是我所跟從的.
是不是主啊,主啊.
的人就可以進天國?.
還有.
你以為我是小事?.
只是回教會?.
我不是只是回來崇拜.
主學的.
我告訴你,我有事奉.
我有很多事奉.
我有私事.
招待,我是關顧的.
我唱詩班領唱.
做主席,做影音的.
是不是?.
我很熱心的.
我還不是那個人?.
我就應該是那個人進天國.
這裡說的.

$^{241}$事奉更厲害.
七章二十二節這樣說.
在那天必有許多人對我說.
主啊,主啊,我們不是奉你的名傳道.
奉你的名趕鬼.
奉你的名行許多義能嗎?.
說到趕鬼.
傳道,這些都是昔日耶穌基督.
參與的事奉.
很尊貴的事奉.
很重要的.
你知不知道.
我們還不是.
我跟從耶穌.
所以我傳道.
我趕鬼.
還有奉耶穌的名行許多義能.
不是我.
一定是我.
我一定是那個進天國的人.
是不是?.
有疑惑的嗎?.
主耶穌這樣說.
祂在這裡向我們發出一個很嚴厲的警告.
七章二十二節說.
在那天必有許多人對我說.
主啊,主啊,我們不是奉你的名傳道.
奉你的名趕鬼.
奉你的名行許多義能嗎?.
主的警告是什麼?.
我要向他們宣告.
我從來不認識你們.
你們這些作惡的人.
給我走開.
這樣也行.
是不是?晴天霹靂.
很大的警告.
我有沒有聽錯?.
我去教會.
我事奉.

$^{281}$我不是.
你不是在說.
我恩順清義嗎?.
你現在在說另一件事嗎?.
這真是耶穌說的.
我從來不認識你們.
你們這些作惡的人.
給我走開.
為什麼?.
原來很重要的一件事.
我的嘴巴在說.
主啊,主啊.
耶穌基督,你是我主.
我不跟耶穌.
我不跟誰.
所以經常去教會.
我做很重要的事.
我有很多行動.
很多活動.
我事奉的.
主啊,你知不知道?.
你看不到嗎?.
你的嘴巴就這樣說.
你有些外在的行動.
不過,你整個生命的內在.
是不配合.
這裡應該說.
不單是行動.
是聰明.
我真正的生命是不配合的.
耶穌基督說.
不是每一個.
稱呼我主啊,主啊.
的人都能進天國.
唯有遵行我天父旨意的人.
才能進去.
你不是只有嘴巴說.
你跟耶穌的.
你也有很多行動,外在的.
但你整個心改變了嗎?.

$^{321}$你整個生命改變了嗎?.
耶穌問我們.
有沒有遵行天父的旨意?.
嘩,你說這麼抽象.
什麼叫天父旨意呢?.
我事奉不是嗎?.
我們今年的教會年提是什麼?.
沒有人夠膽回答我.
不夠膽說,是吧?.
說完之後,周姑娘怎麼回應我?.
毀身事奉.
是嗎?.
不是嗎?.
耶穌基督要問我們.
有沒有遵行天父旨意?.
生命有沒有改變?.
我們天父旨意.
是否真的要問一下?.
好像很抽象.
你的生命有沒有真正改變?.
有沒有悔改離開罪惡?.
有沒有和主建立.
一個愛的關係?.
這個我一定要問.
如果洪恩嘉將來.
要有一個好的發展.
我們要承載.
洪水橋區的人的生命.
今天就要好好地.
面對自己.
問自己一個問題.
不要覺得我已經進入天國了.
你在說.
我是信耶穌的.
我也做很多事.
耶穌現在問你.
你改變了嗎?.
悔改了嗎?.
真的離開罪惡了嗎?.
你真的和主建立.

$^{361}$一個愛的關係嗎?.
什麼叫改變呢?.
是嗎?.
不再拜偶像了嗎?.
有形的.
無形的.
我曾經.
和一個基督徒聊天.
這個不是我們教會的人.
有一個基督徒.
那次應該是.
很久之前.
我接觸他,代表教會關心他.
打給他的時候.
和他聊天.
這個基督徒和我說.
他也會去旅行.
去旅行途中.
難免我們一定會進入廟宇.
你進入廟宇幹什麼?.
燒香.
我另外也認識.
一些基督徒朋友.
和我一起.
進入廟宇燒香.
玩玩而已.
我心想,你拿這些來玩?.
你真的相信耶穌嗎?.
你相信耶穌.
還這樣拜拜?.
這個很容易的.
以前在我們的家庭裡.
我們家人是拜拜的.
一定有很多東西.
傳統家族的信仰.
我們可能是大女兒.
或是大兒子.
要幫忙燒一些東西.
於是就一路拜.
我不是相信的.

$^{401}$我一定要幫忙.
於是我就一路拜.
爸爸媽媽一定有些事情要求.
我就求一下.
這麼多年都是這樣.
你放心,我不是真的相信.
這個習慣還是存在.
不清脆的.
有時候.
想踏一隻腳,一腳踏兩船.
這一刻.
如果耶穌是靈的.
就跟耶穌.
如果我的祖先.
或是我的家人.
拜開也好.
他靈的,我就跟他那個.
總之不要跟教會說就行了.
不要跟熱心的弟兄姊妹說.
誰知道你做什麼.
我回到家裡又不是這樣.
你怎麼知道我家裡在做什麼.
弟兄姊妹告訴你.
你不要這樣想.
上帝只是在教會.
上帝無處不在.
你在哪裡做什麼.
你覺得教會沒有人知道.
你覺得上帝不會知道嗎.
所以不要覺得.
這些說到別人的事情.
在我們當中,如果有弟兄姊妹.
直到今天還抱著這些東西.
願意上帝寬恕我們.
我們真的要認罪.
要撇下這些東西.
有形的.
無形的重敬.
你問問自己.
今天在你生命裡最重要的是誰.

$^{441}$是什麼.
你敢跟我說是耶穌.
不是你的孫子.
你的子女,你的配偶.
你的事業.
你喜歡的興趣.
你敢跟我說耶穌基督是最重要的.
你每個星期看你的時間表.
你花多少時間在什麼事情上.
花在什麼人身上.
你的心思意念.
是放在誰身上.
你很快就知道誰是你生命的上帝.
你不需要有形的東西.
無形的東西才更有殺傷力.
我試過有一次.
真的很震撼.
真的去探訪的時候.
進到一間屋.
這個其次很鮮明.
那個廳.
很厲害.
那個廳四面.
包圍.
浪得很高.
一疊疊.
一棟棟的.
包圍了什麼呢.
是孩子的玩具.
很厲害.
那種愛孩子的心很迫切.
所有事情.
你跟一個家庭聊天.
話題最重要是繞著孩子.
所有事情都是講孩子.
很厲害.
鋪天蓋地.
我們愛孩子是應該的.
愛家人應該.
愛朋友應該.

$^{481}$愛弟兄姊妹也是應該的.
我們生命是要有興趣.
有事業.
要擁抱他們.
愛是對的.
但愛到一個地步.
鋪天蓋地.
超越上帝的時候.
對不起,那些都是你的偶像.
你問問自己最愛什麼.
名譽,地位.
成就感.
我們抱著這些.
怎會有很大的空間.
去擁抱上帝呢.
我們怎會很渴慕知道.
上帝對我們的生命有什麼旨意呢.
是不是因為我們已經.
用盡精力和時間.
去擁抱這些.
還有空間給上帝說話呢.
身心靈各方面的罪惡呢.
這麼久以來.
是吧.
嫉妒的,貪心的.
什麼惹起你的根.
什麼令你.
眼睛發涼.
很想要或者什麼.
憎恨.
憎恨到很多種形式.
講是非,講謊.
講骯髒的說話.
心裡有很多不好的意念.
等等.
千萬不要覺得.
不要覺得.
只是我們自己知道.
或者我身邊和我接觸的人知道.
上帝無處不在.

$^{521}$怎會看不到,聽不到.
怎會不知道我們在做什麼呢.
你想想.
其實我們整個群體當中是很大影響的.
不要緊.
其他人不是這樣就行了.
怎會呢.
一個人影響整個群體.
特別我們每個人都這樣想的時候.
彼此的影響是非常負面的.
很厲害的.
群體的影響.
我們想想.
如果我們擁抱罪惡.
怎會擁抱耶穌基督呢.
耶穌基督是聖潔的主.
看到我們這麼多的罪惡.
祂會很心痛.
我們有沒有離開這些罪惡.
還是被這些罪惡.
完全包圍著自己呢.
我們有沒有和主.
建立一個愛的關係.
做主所喜悅的事情.
我想問.
弟兄姊妹每個禮拜.
你回來崇拜的時候.
你最想做什麼呢.
真的很想敬拜上帝.
還是.
整個禮拜都沒見過朋友.
可以聊聊天.
一直聊.
崇拜前又聊.
崇拜後又聊.
還有沒完的話.
總之很渴慕很開心.
愛弟兄姊妹是對的.
但你對上帝的心又如何.
你回來最渴慕的.

$^{561}$是不是上帝.
是不是最想親近祂.
你說要和上帝有關係.
那你將祂擺在什麼位置呢.
你真的最渴慕祂.
我也在問我自己.
我最渴慕上帝.
最想見到祂.
回來崇拜.
什麼都不要緊.
最重要的是敬拜上帝.
我記得.
大概兩年前.
疫情在香港.
是很厲害的.
兩三年前.
我記得那時在錦宣事奉.
有一次.
是很好笑的.
剛剛疫情爆發的時候.
那一次.
爆發的時候.
當時家人很喜歡回大陸.
我們一家人.
回了大陸.
回了大陸之後.
疫情越來越厲害.
於是區聯會就有個指引.
誰回了大陸就要申報.
還有回到香港.
自己要去居家隔離.
不是政府的要求.
我一收到這個指引.
就想.
我們就.
回到香港之後.
就不能回教會一段時間.
要在家居家隔離.
我回到香港.
應該有兩星期.

$^{601}$不能回教會.
崇拜.
在那個過程.
我就發覺非常不習慣.
以往.
可以這樣回教會.
敬拜上帝真的很多無比.
雖然你說在家都可以.
敬拜上帝的意義是不一樣.
就在那兩個星期.
就知道原來.
是很想.
快點回去.
真的很不習慣.
在家是很方便.
你可以不停地睡覺.
有很多的享受.
很多的休息.
但都不及回教會那麼開心.
我想問弟兄姊妹.
我們有沒有真正.
渴慕上帝的心呢?.
與主建立一個愛的關係.
是要做主所喜悅的事情.
我想問.
我們是否知道.
主喜悅我們做什麼事情呢?.
又很抽象.
現在又見不到主.
誰知.
主喜悅我們做什麼事情呢?.
有什麼方法呢?.
大家有沒有想過?.
我聽到說.
《聖經》.
沒錯.
《聖經》是其中一樣.
要勤讀《聖經》.
才知道主喜悅我們做什麼事情.
如果每個星期.

$^{641}$純粹是星期日.
回來聽堂道.
有些人再多一點.
星期六上主的課.
星期日上主的課.
又或者.
任何時候有少許讀經.
都不足夠.
而是每一天.
自己都要有讀經的時間.
每一天自己都要有靈修的時間.
去親近神.
你才知道主喜悅什麼.
在那裡.
上帝要跟我們說話.
否則的話.
我們都會覺得很迷茫.
我怎知主喜悅我做什麼呢?.
我亦都不會.
自然都不會有行動.
不會行動.
怎可以跟主建立一個.
愛的關係呢?.
你要去愛主.
就要做主所喜悅的事.
做主所喜悅的事.
的過程.
上帝亦都會給我們很深的經歷.
祂的愛.
這是非常寶貴的.
在這裡.
我很想說一個很好笑的例子.
譬如我要說一個人物.
是我們在在的弟兄姊妹.
應該是比較熟悉的.
如果我說起張學友的名字.
在在有沒有人.
不認識張學友呢?.
我相信我們差不多年代.
我們不怕說.

$^{681}$沒有人說.
不認識張學友.
我們大家都知道張學友.
是一位歌星.
是我們年代的.
我相信有不少人.
喜歡聽他的歌.
張學友亦都很紅.
在我們那個年代.
經常強調.
如果要問起.
大家說.
你認不認識張學友呢?.
大家怎樣回答?.
大家很有反應.
你認識他的甚麼?.
唱歌.
可以逐首歌說一下歌名.
甚至乎.
對他的家庭都略知一二.
因為我們會很關心的.
張學友的婚姻狀況.
家庭是怎樣呢?.
大家在報道都留意一下.
你覺得張學友.
認不認識我們呢?.
不認識,為甚麼呢?.
他不是回來教會嗎?.
還有呢?.
為甚麼不認識我們?.
沒有接觸.
對了,非常好的回應.
沒有接觸.
所以.
如果今天張學友.
舉行一個派對.
邀請不同的人.
去他的家.
總之是開放.
一個很大型的派對.

$^{721}$每個人都可以進去.
我們就走去.
派對.
有一個各敲門.
如果張學友來應門.
你又能說出我們的名字嗎?.
完全說不出.
但你又能說出.
張學友,我知道你們.
舉辦一個派對,所以我現在進來了.
張學友說,你是哪位?.
我想問我們認識耶穌.
是否都是這樣?.
對了,都聽過.
耶穌不同事跡.
耶穌厲害,耶穌很厲害.
行神蹟奇事.
都聽過耶穌為我們.
成就救恩.
真厲害,他是萬王之王.
萬主之主,教會常說.
唱詩歌也有這樣唱.
耶穌說不認識你們的意思.
不認識我們的意思是甚麼?.
用同樣的道理.
我們和耶穌沒有關係.
沒有交流.
真的沒有交流.
主喜悅我們.
每天去親近祂.
向祂傾心圖義.
又聽祂說話.
主喜悅我們每天要明白.
祂想我們做甚麼,我們去做.
主又給予回應.
我們是這樣來來去去.
一路相處,建立愛的關係.
是有關係的.
不只是星期日.
回來聽堂道,去上主的學.

$^{761}$現在主說.
從來不認識我們.
從來沒有關係.
這就是無來無往.
我們說自己信了耶穌.
生命又不改變.
一路流於那種光景.
很長時間都是這樣.
以前未信主前.
是怎樣行事為人.
現在信了主後也是一樣.
我就要嘰嘰嘰嘰.
我感覺自己也差不多.
信主後也是差不多的光景.
主說.
從來都不認識我們.
現在這番對話.
是一個警告.
是說將來.
在哪天,主再回來的時候.
就會有很多人這樣說.
主啊,主啊,我又怎樣.
事奉你,又怎樣怎樣.
主就說,將來的日子.
祂會告訴這些人,祂不認識他們.
今天就警告我們.
祂是否想將來這樣一天.
主跟我們說這番話.
祂不認識我們.
從來都沒有關係交流.
說自己信耶穌這麼久.
生命又沒有改變.
人完全一模一樣.
我沒有按照主所喜悅的生活.
我沒有成長過.
我還是我.
耶穌就是耶穌.
不過去到那一天.
我就忽然間很想.
去到那一天跟耶穌一起.

$^{801}$這個是不可以的.
其實如果到那一天.
我們才醒來,是很可憐的.
是很淒涼的.
你想想,去到那一天,我們歡天喜地.
主啊,進來吧,我跟你一起.
主啊,我不認識你.
你哪位啊?.
你不認識我?.
你知不知道我回教會二十年,三十年了.
我從來都不認識你.
不是吧?.
你不認識我?.
你跟我有什麼關係?.
弟兄姊妹.
在這裡挑戰我自己.
挑戰你們.
從今天開始.
生命要改變.
如果今天我們說出去傳福音.
帶人回來.
我們自己.
還沒有搞好的時候.
怎樣去跟那些新朋友.
或者將來.
我們很想發展兒童事工.
青少年事工.
今天我們就要做個板給他們看.
你信耶穌是怎樣的?.
你怎樣跟上帝建立關係的?.
我們有沒有.
這樣的信心給人做一個好的榜樣?.
不可以只是這樣.
做那個兒童事工.
青少年事工.
他們做好了.
他們也想看一下.
那些回洪恩堂這麼久的弟兄姊妹.
我想知道他們怎樣生活.
怎樣經歷上帝.

$^{841}$怎樣跟上帝有互動.
你覺得他們會看到我們什麼?.
你想他們看到什麼?.
今天就要改變.
今天是一個時候改變.
今天我們每一個.
我很想這一刻.
去講訊息的時候.
我想我們弟兄姊妹有一個.
很好的認罪悔改的時間.
在最近這段時間.
我自己也有很多認罪悔改.
為了我自己生命的問題.
我誠意邀請弟兄姊妹加入我的行列.
向上帝認罪悔改.
這麼久以來.
在生命裡有什麼是纏繞我們的?.
跟上帝的關係有什麼難了?.
有沒有每天真的.
將你的生命交出來.
你說主啊 主啊.
祂是你生命的主人 是我生命的主人.
我們是不是聽祂的話呢?.
祂說遵從主旨意.
就是從今以後.
不是為自己而活 不是我作主.
而是主作主.
那有多少人.
是聽祂的話呢?.
所以我不是.
只是對弟兄姊妹說.
也對自己說 我們一起去認罪.
將生命的主權.
交出來.
讓上帝在這裡.
掌管.
還有很多的罪.
跟上帝之間的.
那些問題.
我們犯很多不同的罪 剛才也列舉了一些.

$^{881}$還有跟人的衝突.
在這麼多年.
在我們教會生活裡.
有大大小小的事情發生.
有很多的衝突在我們身邊.
今天.
在我們生命裡.
還有沒有.
這些裂痕?.
有沒有這些憎恨?.
我有沒有不喜悅.
什麼弟兄姊妹.
覺得她總是在我眼中.
得罪我 如果說仇敵.
一定是她 對不起.
我們的弟兄姊妹 永遠都不是我們的仇敵.
你會有不開心 我們會有不開心.
但她不是我們的仇敵.
那惡者才是.
在這裡會不會將我們.
跟人的關係都交上.
是要做主.
所喜悅的事.
你真是.
每天都很想知道.
我們的主是怎樣的呢?.
他喜悅我做什麼呢?.
你不要等到.
星期天才聽到.
每天都要親近主.
你家裡的聖經是不是封了塵呢?.
如果封了塵就拿手機.
手機有下載程式.
都可以讀經的 現在有很多選擇.
我就想我們今天改變.
起來 悔改.
為我們整個群體.
向上帝禱告.
現在我們一同.
有一個一分鐘的安靜時間.

$^{921}$然後我帶領大家一同.
向上帝禱告.
(安靜的呼吸).
天會上帝.
我們在這裡.
將我們的生命交上.
將孩子自己的生命.
將弟兄姊妹的生命.
交在你的手裡.
我很肯定在過去的日子.
我們有很多得罪你的地方.
可能我們有拜有形的偶像.
也有拜無形的偶像.
在我們生命裡有些東西.
我們總是很渴慕的.
有些人我們很渴慕的.
遠超過對上帝的渴慕.
每個星期日回來崇拜.
是否最渴慕上帝你.
都不一定.
天父在這裡求你寬恕我們.
直到如今.
可能我們擁抱著很多的罪惡.
身心靈的.
關係上的.
天父你看到我們滿身的罪惡.
求你寬恕我們.
我願意你是耶穌基督.
你以你的寶血傳言.
塗抹遮蓋我們的罪污.
洗淨我們的罪惡.
聖靈啊.
求你光照我們的生命.
你讓我們看到.
我們真的這麼骯髒.
我們真的很需要.
上帝你的憐憫.
需要離開罪惡.
聖靈求你讓我們有那股力量.
真的在你面前謙卑.

$^{961}$認罪悔改.
讓我們整個群體.
都離開罪惡.
要去得著上帝你的寬恕.
為了討上帝你的喜悅.
去悔改.
聖靈求你充滿我們的生命.
也讓我們真的很想認識.
主所喜悅的事.
究竟是什麼.
讓我們去禱告.
讓我們看聖經.
讓我們天天親近上帝你.
我們天天去聽.
上帝你的聲音.
天天亦都向上帝傾心圖義.
讓我們就從現在.
直到永遠.
和上帝都是這樣親近.
以致到主啊.
我們越來越認識你.
主你亦都會和我們說.
你認識我們.
我們不會成為被撇棄的.
一層人.
主啊在這裡我求你.
復興我們洪恩家.
讓我們洪恩家是一個.
渴慕追求聖潔的群體.
讓我們離開罪惡.
讓我們追求聖潔.
讓我們明白真道.
主啊求你.
在這裡興起.
敬拜讚美的聲音.
是真的.
以心靈誠實來敬拜讚美你.
願你在這裡得著榮耀.
主啊求你在這裡.
去堵塞我們一切的破口.

$^{1001}$讓我們弟兄姊妹.
是很同心的彼此禱告審望.
讓我們和上帝復和.
讓我們彼此又復和.
過去的種種不開心.
主啊求你安慰我們.
醫治我們.
包紮我們的傷口.
主啊我們同樣在這裡懇求你.
懇求你去看顧我們.
去讓我們的生命.
更加的廣闊.
讓我們天天經歷你的愛.
我們就有能力.
愛你愛人.
主啊求你預備我們整個群體.
是將來.
足以去承載洪水橋區.
的人的生命.
也足以去承載.
一群年輕人的生命.
主啊我很深信.
當我們準備好的時候.
神就會將人帶到我們當中.
我們所需要一切的供應.
你都會賜給我們.
大大賜福保守我們.
很願意今天聖經說了一句.
警戒的說話勸勉的.
警告的說話不會發生.
在我們身上.
我們仰望你的憐憫.
奉耶穌基督的名字祈禱.
阿們.
\newpage



\section{使徒行傳路加福音 24:44-49}
\label{sec:5T528Sf9V9Y}
\textbf{2022.05.29 《宣教任務( 一) :務要 傳 道》 路加福音24\_44-49;使徒行傳1\_6-11(新漢語譯本) 楊穎兒傳道}
\newline
\newline
連結: \href{https://youtube.com/watch?v=5T528Sf9V9Y}{\texttt{https://youtube.com/watch?v=5T528Sf9V9Y}} ~~~~ 語音日期: 2022-05-31
\newline
\newline
\hyperref[sec:PjWDhON898w]{\small{< < < PREV SERMON < < <}}
~
\hyperref[sec:index_chronic]{\small{[返順時目]}}
~
\hyperref[sec:index_scriptual]{\small{[返順卷目]}}
~
\hyperref[sec:abVWbFOM8VQ]{\small{> > > NEXT SERMON > > >}}
\newline
\newline
$^{1}$我又可以來到大家當中分享神的話語.
在基督的宗教裡面有一個教會的年曆.
這個年曆除了會幫助弟兄姊妹在不同的節期裡面去閱讀經文之外.
最重要的是引領教會群體去紀念主耶穌的一些重要的事件.
可能大家也聽過在過去的壽苦節和復活節之前的七個星期裡.
大概四十天,這個叫做大齋期.
就是叫教會群體一起默想耶穌基督進入壽苦的路.
同樣在復活節之後的七個星期,大概四十天.
這段時間叫做復活期,就是紀念主耶穌從死裡復活.
並且在地上再顯現了四十天,才被雲彩接去升上天.
所以在我預備講章的時候,我發現今天剛好就是復活期的最後一個主日.
所以禱告之後也很想和大家藉著這個機會.
去紀念主耶穌昔日復活之後,仍然繼續在信徒的中間.
去猜念他們完成抗戰福音的使命.
當中我們看的經文,耶穌對門徒有很多豐富的吩咐.
今天我們又怎樣可以按著主耶穌的教導,繼續去接續傳揚福音的使命呢?.
所以我決定今天講到的題目是宣教任務一,無要傳道.
宣教任務一,有一就是有二.
希望將來可以繼續和大家分享關於宣教使命任務的訊息.
亦都很想從不同的層面去認識宣教,猜傳,傳福音.
這些都是和耶穌的使命有關.
讓我們一起祈禱.
親愛的天父,你親自猜猜你的愛子主耶穌基督降臨在這個世界上.
去按著你的拯救計劃,讓世人都去認識你.
願我們今天再一次去紀念主耶穌的復活.
並且在世上繼續與人說話.
並且教導我們怎樣去完成主耶穌的使命.
求你使用孩子的口去分享你的話語.
由預備個人的心去聆聽你的說話,你的教導.
奉主耶穌的名求,阿們.
在進入經文之前,想和大家分享一個經歷.
我有一個六歲的兒子,他叫做子倫.
他經常問我一個問題.
「天堂是怎樣的呢?」.
每一次他問我,我就透露一點給他聽.
我就跟他說,天堂就是永遠同住一起的地方.
那裡沒有眼淚,那裡是一個很喜樂的地方.
每次聽完,他都很興奮.
而且告訴我,馬上很想去天堂.
他今天也在場,聽我說,他覺得很尷尬.

$^{41}$有一天,子倫問我,究竟怎樣才能去到天堂.
如果我要把整個福音,救音都告訴他.
按照他六歲來說,似乎太複雜了.
但如果我跟他說,你死了之後,你上天堂了.
你又有永生了,但是很不完整.
而且很誤導性.
那怎樣好呢?.
於是我就跟他說,子倫.
天堂在這裡,在我們的心裡.
有聖經根據的.
因為傳道書是這樣說的.
傳道書,三章十一節.
神造萬物,各安其事,成為美好,又將永生安置在人心裡.
我跟他說,在這裡的.
子倫聽了很開心.
他一直很好奇,很想擁有的天堂原來一直在他裡面.
之後,子倫就沒有再問了.
他沒有問我,但我就繼續告訴他.
你可以將這個天堂放在你的心裡面.
或者你可以將這個天堂跟別人分享.
子倫又更開心了.
原來還有後續的.
我說你可以跟別人分享愛.
你可以關心別人,跟別人祈禱.
你不要說謊話,不要作弄別人.
有很多教導的.
你還可以跟主耶穌一起.
一起將天國帶到地上.
到你所去的地方,就好像天堂一樣.
那次談完之後.
子倫沒有再問我天堂的事.
一句都沒有問.
我相信他明白了.
當然,子倫現在六歲.
他還有很多時間,要一生去學習.
並且活出主耶穌的教導.
小朋友就是這麼單純,這麼學武.
這麼有毫無懷疑的心去迎接天國.
主耶穌曾經這樣說.
除非你回轉,變得好像小朋友一樣.

$^{81}$否則絕對不能進入天國.
今天我們所學習的宣教任務.
其實都是要人回轉悔改.
回到神的身邊.
今天的經文是來自《路加福音》和《仕途行傳》.
這兩本書就像上策和下策.
作者都是老家醫生.
他是跟隨著保羅時期去宣教.
就像第二代的信徒.
我特意將這兩段經文放在一起.
是想看看老家作者.
在上策和下策記錄耶穌升天日發生的事情.
作為今天我們要學習的經文.
很多時候我們有意無意都會覺得.
上策是講主耶穌.
下策是關於保羅和其他的信徒.
但其實我們閱讀聖經的時候.
應該以整本聖經都是以基督為中心.
甚麼是基督為中心?.
意思是我們整本聖經.
每一章每一節每一個段落.
其實都是跟耶穌基督有關.
不單止跟耶穌基督有關.
而且是跟耶穌基督拯救的任務有關.
今天我們看的《路加福音》經文.
短短的六節.
是耶穌自己重新再講一次.
從過去現在到將來在地上的救恩任務.
關於過去的部分.
第44節.
耶穌對他們說.
「我從前和你們在一起的時候都跟你們說過.
就是摩西的律法.
先知書和詩篇上所記.
關於我一切的都應驗了」.
當時還沒有我們現在看的聖經.
我們現在手上看的那本.
那時候的信徒不是看這本.
他們看的那本叫做《希伯來聖經》.
有三個部分.

$^{121}$就是摩西的律法,先知書和一些寫作.
就好像一些詩篇一樣.
所以耶穌說.
從前希伯來.
就是那時代的人看的那本《希伯來聖經》.
所說的話都應驗了.
耶穌道成肉身來到這個世界上.
然後講現在.
第46節.
他對他們說.
「經上記住.
耶穌必須受難第三天從死人中復活.
並且人要奉祂的名傳講悔改和赦罪的道.
從耶路撒冷開始直到萬邦」.
你們就是這些事的見證人.
耶穌來到世上.
就是要聚集更多的人.
和耶穌一起去作見證.
然後耶穌又說到將來.
耶穌說.
「漢娜,我要把我父所應許的差遣到你們那裡.
但是你們要留在城裡.
直到你們得著從天上而來的能力」.
如果我們有了這個框架去閱讀聖經的時候.
其實整本聖經都可以以耶穌基督為中心.
當我們打開聖經的時候.
我們很想閱讀神的話語.
我們想與主相遇.
我們很想學習耶穌.
我們在聖經所學的教導.
我們會學習去遵從.
我們想與主耶穌親近.
耶穌所思想的事.
我們都會掛念.
耶穌的教導.
我們會去學習.
我們會去學愛耶穌愛的人.
當我們經常想著耶穌所想的事的時候.
我們整個行為表現.
言行舉止.

$^{161}$內心產生的那種憐憫.
使我們越來越像主耶穌.
而主耶穌在他升天之前.
一個這麼重要的時刻.
他最後會說什麼呢?.
你想想.
如果你離開家人一段時間.
去旅行或是去公幹.
你會跟家人說什麼呢?.
可能是小心.
照顧自己.
照顧家人.
鎖好門.
這些聽起來很日常的用語.
都是由心底而出的一些吩咐.
耶穌說什麼呢?.
在我今天整合的路加福音.
和《使徒行傳》的經文裡面.
平衡來看.
有三個重點或方向.
我們一同可以學習.
第一.
得著能力去傳道.
路加福音.
二十四章四十九節這樣說.
但你們要留在城裡.
直到你們得著從天而來的能力.
同時在《使徒行傳》又是這樣說.
聖靈會臨到你們身上.
你們必得著能力.
這個能力當然是神的靈.
聖靈的帶領.
這個靈在路加福音裡面.
不是現在才出現的.
不是耶穌臨要離開才出現的.
早在施洗約翰.
未出生之前.
聖靈已經感動了他的母親撒加尼亞.
他的母親被聖靈充滿.
到瑪利亞懷孕的時候.

$^{201}$他也被聖靈感動.
到施洗約翰預言.
耶穌會來.
耶穌會以聖靈去更新百姓.
這個聖靈也在.
當耶穌受洗的時候.
聖靈降臨在他的身上.
而到耶穌出來傳道的時候.
聖靈又引領他到曠野.
去受魔鬼的試探.
到耶穌第一次在會堂的講道.
主的靈也臨到他的身上.
要他去拯救一些貧窮.
受迫迫的人.
所以我們不要誤會.
聖靈只是在耶穌升天之後.
和信徒一起去承接耶穌時代的工作.
這個靈其實早就在了.
經文提到.
信徒需要等候神的靈.
降臨在他們的身上.
這個能力是用來做什麼呢?.
《路加福音》第47節提到.
這個能力是使人.
能夠奉他的名去傳講.
悔改和赦罪的道.
無論我們是參與宣教.
猜拳 傳福音.
這些方向.
其實都是需要有聖靈的帶領.
得著聖靈的能力.
然後聖靈感動我們.
能夠奉他的名去傳講悔改赦罪的道.
中文的翻譯是傳.
好像傳講 傳道.
好像局限了.
是不是只有牧師 傳道人.
才可以去做這件事呢?.
但原來我們看回原文的時候.
其實它有一個宣告的意思.

$^{241}$那種熱情就是.
你可以站在哪裡都可以.
向人宣告.
耶穌在我身上的作為.
宣告神在我們身上的事實.
發生了的事實.
所以我們去宣講的時候.
不是說.
你回來教會吧.
教會好 教會開心.
教會有活動玩.
這些都不是吸引人.
是吸引人來教會的事.
其實這個能力.
是要讓我們能夠去宣告.
神怎樣去拯救我們.
這個能力.
不只是滿足我們的生活.
我們的慾望.
我們能夠一帆風順.
這個能力.
就是要我們去懂得.
去傳講昔日.
神怎樣令我們悔改.
令我們歸順他.
是神的道再次吸引人.
次盼望給人.
讓他們能夠脫離罪惡.
悔改.
讓神進入他們的生命.
去改變他們.
盼望我們都能夠.
祈求主次聖靈的能力.
讓我們去傳講悔改的道.
第二.
宣教任務的方向.
可能你都會留意到.
其實當你閱讀路加福音.
《士同行傳》的時候.
很強調主和聖靈的旅程.

$^{281}$他們一起去了哪裡.
去做了什麼事情.
所以我們看看以下兩節.
都給了我們一個很清晰的方向.
耶穌給門徒有一個很清晰的.
宣教傳福音猜拳的方向.
路加福音24章47節.
從耶路撒冷.
起直傳到萬邦.
《士同行傳》一章八節說.
並且要在耶路撒冷.
整個猶太,撒瓦利雅各地.
直到地極.
大家有沒有留意到.
兩個經文都很強調.
這個福音的方向.
這個旅程都是在耶路撒冷開始.
不是由地極回來耶路撒冷.
而是由耶路撒冷開展出去.
當時門徒聽到耶穌說耶路撒冷.
就好像我們聽到人們說.
獅子山就等於香港人的精神.
太平山是景點.
所以當門徒一聽到耶路撒冷.
這是一個很重要的地方.
總在他們閱讀的希伯來經卷裡.
其實耶路撒冷已經是一個聖殿的地方.
到了耶穌的時期.
耶穌在耶路撒冷的哈西勒瑪利園禱告.
在耶路撒冷和門徒有最後的晚餐.
也在耶路撒冷受死被釘.
所以耶路撒冷其實就是一個.
門徒和主相遇的地方.
與主相遇的地方.
可能是在教會.
不一定是在教會.
可能是一些經歷.
不知道在大家的生命裡.
是什麼情況.
什麼光景去與主相遇呢?.

$^{321}$我和大家分享一個.
我曾經與主相遇的經歷.
很多年前我加入了一個青年團契.
團契裡有導師.
團友有五個男五個女.
很平均的.
在導師教導我們聖經的學習裡.
提到年輕人最好找一個基督徒的配偶.
這是很好的教導.
因為我們有共同的信仰.
可以同行天路.
但我慢慢發現我參加的團契裡.
弟兄姊妹拍拖的對象.
都是在這個團契裡.
後來我到外國參加一些職稱的團契.
我又發現一些很類似的情況.
就是很想在教會裡找一些配偶的對象.
後來我跟主說.
我覺得不一定要在教會裡找.
如果我找到一個未信的弟兄.
我又跟他一起經歷信主的過程.
不是更加能夠親眼體驗神的能力嗎?.
沒多久我就蒙召.
我記得蒙召的那一刻站起來.
很感動.
主給了我一個任務.
就是我要找一個配偶跟我一起走這條路.
後來沒多久我又認識了我先生.
他從不認識福音.
到現在在時空的路上.
他都很支持我.
在這十年來.
從帶主的經歷是很不容易的.
但在不容易當中.
我又看不到出路.
身邊的弟兄姊妹可能會覺得.
為什麼會選擇這個配偶.
以為我遠離了主.
這就是我跟主相遇的一個片段.
也是我人生朝聖旅程的其中一幕.

$^{361}$回頭看的時候.
如果今天沒有了我先生.
我都看不到我今天怎樣可以再走下去.
經文說在耶路撒冷.
整個猶太 撒瑪利雅各處直到地極.
這是我們傳福音的方向.
這是一個旅程.
可以解作地理性旅程.
但當我看的時候.
這也是一個人生的旅程.
所以我想說.
宣教傳福音.
這些耶穌的使命.
我們要繼續接續下去的時候.
不單單是說我們今天在這個地方.
我們要本地傳福音.
我們有宣教士去外國做宣教工作.
這是一個地理性的理解.
但原來當我們去擁抱耶穌的使命.
本質的時候.
其實要用我們每個人的生命.
自己去經驗 去參與.
我們每個人的生命都經歷過主.
與主相遇.
這個朝聖的旅程是很獨特的.
是你才獨有.
所以我們去到哪裡.
其實我們也可以帶著這個旅程.
這個經歷.
向人宣告.
悔改的福音究竟是怎樣可以.
吸引到其他人更多去認識主耶穌.
第三 作我的見證人.
路加福音第24章47節.
又這樣說.
你們就是這些事的見證人.
這事的見證.
第四圖行傳一章八節.
也是作我的見證人.
耶穌說了很多事情.

$^{401}$告訴他們很多事情都應驗了.
都是從希伯來的聖經已經說到.
他還未道成肉身的時候.
經卷已經說有關他的預言.
他來到之後又有很多的教導.
到現在對他們的門徒.
都要等候聖靈再去引領他們去繼續工作.
為的就是要世世代代的人.
都認識神的作為.
要為神作見證.
就是要講實這個主是多麼真實.
從前主耶穌怎樣開我們的眼睛.
讓我們明白福音.
我們是需要講出來的.
神真的給我們很多獨特的經歷.
上帝讓我年輕的時候去到教會.
讓我看到.
我不用局限了選擇一個信主的弟兄做配偶.
所以他才會有這樣的任務給我.
原來感動我去找一個未信的弟兄.
最終不單止成為信徒.
而且今天我看著我先生.
無論去到哪處他都會和人分享福音.
他不單止叫人信耶穌.
他花更多的時間去陪伴一些不認識主的人.
和他們走過人生一些很難過的難關.
然後帶他們來神的面前.
叫他們悔改.
為的是盼望這些人都得到和他一樣.
那份新生命的禮物.
我覺得他展現的這種精神.
無論去到哪處他都能將福音帶給人.
由耶路撒冷與主相遇的地方.
去到猶太各處.
到地籍.
我相信今天你和我也有很多獨特的經歷.
有些經歷是你才有我沒有的.
所以我們真的要很珍惜.
去整理一下我們自己的經歷.
不是說我們的經歷有多特別.

$^{441}$或神是怎樣愛我們.
而是這些經歷是見證著神的作為.
神在我們身上的能力.
祂的慈愛.
讓我們都能夠盼望繼續與聖靈同工.
一同各自在自己的崗位那裡.
去完成耶穌吩咐門徒.
繼續要做的宣教.
猜傳 傳福音的任務.
所以我都說了今天三個學習的方向.
第一就是會得著聖靈的能力.
第二就是宣教任務的方向.
其實都是我們很鼓勵大家.
去回想自己最初與主相遇的那一刻.
就是整個生命的朝聖的旅程.
主是怎樣帶領你的.
這些都是我們時刻可以去思想整理去感恩.
無論我們將來會以禱告.
或親身去到不同的角落去傳福音宣教.
都是帶著這一份精神.
要為主作見證.
為的證明不是神怎樣祝福我們自己.
而是祈求神能夠使用我們的生命.
將這個祝福再推到給其他身邊的人.
去年五月尾我在紅印堂完成了最後一個月的實習.
到現在又是五月了.
一年過去了.
很想藉這個機會和大家分享一下.
我這一年與主走過的什麼路.
我記得去年和大家在紅印堂的祈禱會.
自己也領受到很愛非洲這個群體.
也曾經在主日的講道裡.
提到過去去過非洲短孫.
也很想再次服侍非洲的群體.
很奇妙地畢業之後.
神又帶領我到現在的教會.
九龍幽靈堂去服侍.
主要我都是服侍一些來自非洲這片土壤的弟兄姊妹.
有來自非洲各國.
還有階級的信徒.

$^{481}$階級其實都是屬於非洲的土壤上.
他們因著不同的政治迫發.
宗教的迫發來到香港.
來到香港他們的資源很困乏.
過的生活也需要政府的資助.
或者有心人的資助.
所以在這個服侍的裡面.
其實很多時候我都會去探訪他們.
其實很開心分享這些相片.
跟他們探訪的時候.
他們會將自己國家的事告訴我.
或者從他們的國家.
他們以前怎樣生活.
文化.
很多故事可以從他們的口中聽到.
就好像自己也去了非洲一樣.
我也會在教會有些團契的活動.
這次的團契活動就是跟他們做一些勞作.
一家人也有來.
有不同的桌子.
做一些活動.
我也會跟他們一起去唱詩.
他們有一個詩班.
久不久就會出去出隊.
去唱一些非洲的歌曲.
跟不同教會的弟兄姊妹去互動.
這是我的第二個小朋友.
當她在教會上班的時候.
她也接觸到其他來自非洲的小朋友.
其實如果你只是看相片.
其實我好像去了非洲一樣.
除了你知道她在香港之外.
近期在疫情之下.
因為他們得到的資源也很困難.
很多東西都封鎖了.
他們得到的資源有限.
物價也上漲.
很開心跟其他機構做了派送隊.
我們有機構捐贈一些食物.
水果蔬菜物資.

$^{521}$綠色衣服的是一個非洲人司機.
他在港可以工作.
他負責開這輛車.
我們去全港九新界有幾個點.
一個下午去幾個點.
派發一些物資給難民群體.
這是我們去到深水Po 派物資的時候.
跟他們拍了照.
他們有很多聚居在深水Po 土瓜灣的位置.
這是在元朗.
元朗也有很多非洲群體.
其實在我參與派送的時候.
發現這個群體很大.
我接觸到的只是少數.
有很多人聚居在錦田,元朗.
有些人未必接聽過福音.
有些是回教徒.
在我心裡也在想.
會否有一天可以跟他們講福音的訊息.
這個可以介紹一下.
自2008年開始.
九龍幼齡堂開展了難民事公.
除了鼓勵難民朋友參加崇拜和團契之外.
我們還會做很多落地的服事.
提供日常物資,生活津貼等等.
幫助改善他們實際的生活需要.
另外,九龍幼齡堂也不時透過主日廣道,團契領會.
合唱團出隊,外展見證分享,難民廚房等等活動.
幫助難民朋友更加懂得擁抱他們作為難民的身份.
同時,九龍幼齡堂也藉著這個事公.
提供一個跨種族,跨信仰的平台.
讓社會各界能夠關懷和支援寄居在港的難民朋友.
在種種的生命互動交流裡.
參與建立天國的事公.
我也很開心能夠跟大家分享.
我在九龍幼齡堂的服事.
也是一個見證.
由紅印堂服侍非洲人.
到現在每天接觸他們的時候.
都是想跟大家見證的經歷.

$^{561}$也是我朝聖旅程的一部分.
盼望今天復活期的最後一個主日.
我也用真命的見證去演繹了這個經文.
更加盼望的是紅印堂的弟兄姊妹.
你們很珍惜與主相遇的故事.
大家一起起來為主作見證.
又求主開了個人的心.
就是「未信的朋友」的心.
讓他們能夠在悔改之後.
不單止停留在認罪的那一刻.
更加能夠擁抱一個新的生命.
讓我們一起祈禱.
親愛的天父.
我們滿心歡喜的來讚美你.
因為你就是那位賜生命給我們.
賜經歷給我們的那位神.
我們讚美你.
主耶穌求你呼喚我們個人的心處.
呼喚紅印家.
呼喚我們在這個朝聖的旅程裡.
願你繼續的拖大人靈.
時常提醒我們要讓人明白悔改的道.
和為你作見證.
聖靈也求你臨到我們的生命裡.
叫我們懂得去等候你的來臨.
叫我們得著能力.
叫我們可以隨著你的帶領.
我們一起宣告主耶穌復活的大能.
叫更多人悔改歸你.
我們這樣的祈禱.
是奉我們復活的主名而求.
阿們.
\newpage



\section{啟示錄 2:1-7}
\label{sec:abVWbFOM8VQ}
\textbf{2022.06.05 《重回信仰生命的起跑線上》 啟示錄 2\_1-7 (和修版) 老耀雄牧師}
\newline
\newline
連結: \href{https://youtube.com/watch?v=abVWbFOM8VQ}{\texttt{https://youtube.com/watch?v=abVWbFOM8VQ}} ~~~~ 語音日期: 2022-06-07
\newline
\newline
\hyperref[sec:5T528Sf9V9Y]{\small{< < < PREV SERMON < < <}}
~
\hyperref[sec:index_chronic]{\small{[返順時目]}}
~
\hyperref[sec:index_scriptual]{\small{[返順卷目]}}
~
\hyperref[sec:qiH_9wRKv74]{\small{> > > NEXT SERMON > > >}}
\newline
\newline
啟示錄 2:1-7
\newline
\begin{longtable}{cl}
\hline
\hline
章節 & 經文 (和合本修訂版)\\
\hline
2:1 & \begin{tabularx}{0.7\textwidth}{X} 「你要寫信給以弗所教會的使者，說：『那右手拿著七顆星，在七個金燈臺中間行走的這樣說： \end{tabularx} \\ \\ \relax
2:2 & \begin{tabularx}{0.7\textwidth}{X} 我知道你的行為、勞碌、忍耐，也知道你不容忍惡人。你也曾察驗那自稱為使徒卻不是使徒的，看出他們是假的。 \end{tabularx} \\ \\ \relax
2:3 & \begin{tabularx}{0.7\textwidth}{X} 你能忍耐，曾為我的名勞苦而不困倦。 \end{tabularx} \\ \\ \relax
2:4 & \begin{tabularx}{0.7\textwidth}{X} 然而，有一件事我要責備你，就是你把起初的愛心拋棄了。 \end{tabularx} \\ \\ \relax
2:5 & \begin{tabularx}{0.7\textwidth}{X} 所以你要回想你是從哪裡墜落的，並且要悔改，做起初所做的工作。你若不悔改，我要到你那裡去，把你的燈臺從原處挪去。 \end{tabularx} \\ \\ \relax
2:6 & \begin{tabularx}{0.7\textwidth}{X} 然而你還有一件可取的事，就是你恨惡尼哥拉派的行為，這種行為也是我所恨惡的。 \end{tabularx} \\ \\ \relax
2:7 & \begin{tabularx}{0.7\textwidth}{X} 凡有耳朵的都應當聽聖靈向眾教會所說的話。得勝的，我必將神樂園中生命樹的果子賜給他吃。』」 \end{tabularx} \\ \\
[1ex]
\hline
\hline
\end{longtable}
$^{1}$主席.
各位紅映家的弟兄姊妹主內平安.
在我未讀神學的時候.
有一段時間我是做市場策劃的工作.
周不時整個團隊人一起去想一些方案.
寫好詳細的計劃書之後就交給老闆去審議了.
很多時候老闆彈回來的意見.
除了批評計劃寫得不夠全面之外.
他有一個更大的提醒.
就是一個這麼長遠的計劃.
為什麼在過程中竟然沒有一個檢查點呢?.
原來在一個標準的企業發展計劃裡.
一定會設立很多這樣的監察點.
這些檢查點可以讓團隊停一停.
想一想重新檢視一下.
在運作過程中有沒有地方走錯了.
有沒有地方需要修正.
好的地方會保留.
不好的地方要停止.
調整好之後再向目標繼續進發.
檢查點的設立是想大家不要只一味衝.
大家還要留意一下.
現在走的方向有沒有錯.
然後才繼續走.
如果一間公司或社會.
在發展和運作過程中.
需要有這樣的檢查點的話.
去審視有沒有走錯路.
一間教會或信徒的生命.
在乘勝的旅程過程中.
是否都要加入一些checking point.
這樣的監察機制.
去審視一下我們自己的屬靈生命.
有沒有偏離神的方向和心意.
今天講到的經文是.
在一間初期教會期間.
神要求其中一間教會.
對自己的屬靈狀況進行一個檢視.
好叫他們能夠調整信仰的方向.
再重新出發.

$^{41}$我們今天一起在啟示錄.
二章一至七節當中.
檢視一下信仰的三個過程.
我們一起去祈禱.
七一恩主你是教會的元首.
你帶領著每一個跟隨你的信徒.
甚願我們這群屬你的子民.
能夠專心聆聽你的教導.
信服你的指引.
並且實踐在我們每一天的生活裡.
讓人在我們的生命裡見到你.
我們的祈禱是奉主耶穌基督的聖名祈求.
阿們.
在進入經文之前.
我們一起去了解一下.
這本經文的資料.
啟示錄有別於聖經其他經卷的體材.
因為它不是詩歌體材.
不是敘事體.
也不是歷史.
更加不是書信.
它是屬於一種天啟文學.
這種文學的體材.
在主前二世紀猶太地方相當流行.
它的特點就是.
以一些象徵符號和數字去見稱.
符號的重點不在於符號本身.
而在於符號所指向的方向.
即是符號的象徵意義.
大於符號實質的意義.
在這次講道的經文裡.
有七個星.
七個金燈台的符號和數字.
不過約翰在一章二十節.
已經為我們解釋了符號的意義.
七個星是指七個教會的使者.
而七個燈台是指七間教會.
有關地理方面.
爾忽所城在啟示錄七個城市之中.
是一個很重要的城市.

$^{81}$因為爾忽所不單是羅馬法庭的所在.
亦是地方總督和政府的所在.
每一個總督上任.
都必須進入爾忽所去行省.
爾忽所是一個宗教的中心.
女神亞底米的崇拜風氣相當盛行.
所以爾忽所城有一個很尊貴的頭銜.
它叫做神廟看守者.
而號稱世界七大奇觀之一的大神廟.
亦都坐落在爾忽所.
有學者從地理結構去看.
去記錄這七間教會.
由於爾忽所城距離不摸道最近.
所以約翰以順時針的方向.
按著地理的次序.
逐間逐間地提名七間教會.
這是一個看法.
亦有學者從另一個看法.
就是從形式方面去看這七間教會.
因為七間教會有著同一個格式.
這個格式是甚麼呢?.
有宣告者,有嘉許的地方.
有責備的地方.
有警告提醒的地方.
有呼召的地方.
亦都有應許的地方.
如果從內容和教會的情況來看這七間教會.
第三和第四,第五間教會的情況都相似.
都是很壞參半.
第二和第六間教會是最好的.
都是沒有受到責備的教會.
而第一和第七間教會的情況就最危險.
所以審判是最嚴厲的.
如果從內容和教會的情況來看.
以弗所教會正正是屬於第一類別的教會.
亦都是屬於一間很危險的教會.
我們開始進入經文.
一起看看以弗所教會究竟如何危險.
第一個檢視信仰的過程就是回想起初的呼召.
當以弗所教會回想以往的事奉歲月的時候.

$^{121}$都不是完全一無是處.
起碼有兩個範疇是給神稱讚的.
第一是以弗所教會的屬靈生命.
經文在第二節這樣說.
「我知道你的行為,勞碌,忍耐」.
以及第三節.
「你能忍耐,曾為我的名,勞苦而不困倦」.
第二節所指的行為並不是指工作上的行為.
而是指事奉上的行為.
「神稱許以弗所的信徒在事奉上勞心勞力」.
即是說做到抽筋.
筋疲力盡.
至於忍耐是指他勇於接受苦難.
讓苦難成為神的榮耀.
亦都藉著苦難讓生命更加成熟.
「神稱許以弗所教會」.
他很願意為生.
在事奉上勞心勞力.
肯犧牲不計較.
在面對羅馬政府的迫迫.
他能夠忍耐沒有離棄信仰.
可以這樣說.
神稱讚以弗所教會的信徒的靈命是成熟的.
不是那種還在吃奶的屬靈BB.
第二個被神稱讚的是他的教義.
亦都是他行在真理之上.
經文第二節這樣說.
「你也曾察驗那自稱為使徒卻不是使徒的.
看出他們是假的」.
以及第六節.
「你還有一件可取的事.
就是你恨惡尼哥拉派的行為」.
這種行為都是我所厭惡的.
在初期教會期間.
自稱是先知或使徒.
這些假使徒混入教會.
去博取信徒的信任.
然後傳遞一個錯誤的福音.
事實上保羅在牧養以弗所教會的時候.
已經提醒過信徒有關這些事情.

$^{161}$我們可以在保羅在使徒行傳.
二十章二十九至三十一節.
對以弗所的信徒這樣說過.
「我知道在我離開之後.
必定有兇暴的豺狼進入你們中間.
他們不顧適陽群.
就是你們中間.
也必有人起來.
說謊的話.
要引誘門徒跟從他們.
所以你們要警醒.
紀念我三年之久.
晝夜不斷地流淚勸誡你們各人」.
從保羅說這番說話.
我們可以推斷當時假使徒已經出現了.
早在保羅時期已經活躍在以弗所這個城市.
到了日漢寫啟示錄的時候.
假使徒的問題.
已經實質地影響當時的教會.
而且也影響了信徒對真理的實踐.
另外第六節也提到尼哥拉一黨的教訓.
有什麼教訓呢?.
是指一種唯利是圖.
主張妥協自由放蕩的思想.
如果用今天的說法就是.
向淺漢個人主義抬頭.
而以弗所又是一個商業發達.
迷信盛行的城市.
當時羅馬皇帝迫害教會.
有些信徒就嫌教會太過拘謹.
很執著.
尼哥拉一黨所提倡的就是妥協信仰.
妥協信仰.
什麼意思呢?.
不要緊.
無所謂.
做了不用死.
只要回來就行了.
他和侵入別加摩教會的巴蘭異端.
以及推亞推拉的呂先知耶駛別.

$^{201}$都有一個共同的特點.
就是認為信徒可以拜偶像.
即是你信耶穌.
OK呀好呀.
不過拜觀音也無妨.
初三是否去拜車公廟?.
也無所謂.
神也不會生你的氣.
類似這種思想.
所以以不鎖教會能夠分辨真假仕途的身份.
沒有被假仕途將錯誤的福音影響到教會.
同時恨惡而端.
在教義上真理上能夠堅守信仰.
所以避諸稱讚.
說到這裡.
其實我們也應該很羨慕.
他這種優點.
是不是?.
他是一個信徒靈命成熟的教會.
願意事奉.
委身.
願意克服和內怒.
不過回想的過程裡面.
在checking point裡面.
固然有好的地方.
也有不好的地方.
以不鎖教會有一樣東西是被主所責備的.
就是第四節.
你把起初的愛心拋棄了.
我們想知道.
以不鎖信徒起初的愛心是甚麼呢?.
我們可以看回.
保羅在以不鎖書一章十五至十六節.
怎樣形容他們的愛心.
保羅這樣說.
我們看到以不鎖信徒起初是甚麼呢?.
對主耶穌有信心.
對眾信徒有愛心.
也就是以不鎖教會起初是一間.
愛神愛人的教會.

$^{241}$現在神提醒以不鎖教會.
這種愛神愛人的心.
消失了.
甚麼原因使他們消失呢?.
甚麼原因使他們那種.
愛神愛人的信心消失呢?.
我們要進入.
檢視信仰的第二個步驟.
第二個步驟就是.
尋找不足並且悔改.
在尋找不足上有兩個原因.
就令到以不鎖教會失落了這種.
起初的愛心.
第一.
以不鎖信徒雖然熱心事奉.
但沒有合一的心.
保羅在以不鎖書四章一至三節.
提醒過他們.
勸勉過他們.
保羅怎樣說.
我為主作酬勞的勸你們.
既然蒙召.
行事為人就要與你們所蒙的呼召相稱.
凡事要謙虛,溫柔,忍耐.
用愛心互相寬容.
以和平彼此聯繫.
竭力保持聖靈所賜的合一.
如果事奉的群體.
只是著重侍宮的果效.
多於彼此關係的話.
事奉的重點.
只是對焦於侍宮身上.
只顧成果.
而不是對焦在神的心意身上.
沒有體諒到其他成員的感受.
整個事奉.
必然沒有合一的心.
即是我們說甚麼呢?.
侍宮化.
教會是有侍宮的.

$^{281}$但侍宮化只是看目標.
而不看過程中.
大家是否有同心.
大家只看結果.
做到就行了.
舉辦報道會.
只要有人信主就成功了.
但沒有想過.
在舉辦的過程中.
大家是否合一.
大家有沒有為那人祈禱?.
第二.
雖然謹守教義.
但過猶不及.
信徒在非核心信仰上.
彼此針鋒相對.
堅持己見.
排除異己.
而不所有信徒固然恨惡.
罪惡.
但同時也恨惡罪人.
視他們為仇人.
忘記了.
大家都是罪人.
我們要傳一個彼此寬恕的心.
已不所信徒在繁忙的事奉中.
失去了焦點.
沒有一個合一的心.
失去了對神對人的愛.
在力爭教義之際.
堅守真理的同時.
忘記了關顧信徒的需要.
再次失去了對人的愛.
作者約翰在其他書信.
對愛神愛人有這樣的提醒.
他說:人若說我愛神.
卻恨他的弟兄.
這就是謊言.
不愛漢德傑的弟兄.
就不能愛漢不見的神.

$^{321}$愛神的都要愛弟兄.
這是我們從神所受的命令.
神是看不到的.
我們說愛神.
我們愛神.
你愛的弟兄.
就表示你愛神.
你不能在口中說愛神.
卻恨你的弟兄.
你每個星期回來教會.
遇到的每一個都是弟兄姊妹.
我真的想問.
你認不認識他們?.
你有沒有跟他們打過招呼.
你有沒有跟他們聊過天?.
你知不知道他們叫什麼名字?.
他們住在哪裡?.
我們要承認.
我們不能只用口說愛神.
只用口去愛神.
卻不去愛我們的弟兄姊妹.
愛神愛人.
就好像同閉的兩面.
兩面都不能分割.
我們不能說我愛神.
卻不愛我們的弟兄姊妹.
既然尋找到失落愛心的原因.
就要有悔改歸回的心.
信仰生命裡總會有起起跌跌的.
正所謂勝不驕.
敗不累.
神知道我們每一個都很軟弱.
不過祂也曾經應許我們.
祂說過什麼呢?.
祂說我們若認自己是罪.
神是信實的.
是公義的.
是會赦免我們的罪.
洗淨我們一切的不義.
我們每一個人都會犯錯的.

$^{361}$犯錯不是罪無可寬恕的.
罪無可寬恕的就是.
我們犯了錯之後.
都不肯去承認我們的錯.
承認自己的錯.
向神認罪悔改.
重新回到神的懷抱.
也重新回到信仰的起始點.
重新再來過.
我曾經參加一個區聯會的講座.
其中有一個講者.
也是一間教會的堂主任.
他在這裡分享了一個經歷.
在教會的一次執事會的會議過程裡.
會議裡有一個很重要的議題.
要去討論.
是有關教會一個相當重要的發展.
有人贊成.
有人反對.
有人贊成.
有人反對.
這個都很常見的.
他說當時氣氛很僵.
會議討論了很久.
兩三個小時.
有些執事就不耐煩.
就說主席不如表決.
堂主任就是當原主席.
就宣佈不如開始投票.
投票結果六對四.
就是六票贊成.
四票反對.
按照議會的規定.
是什麼?.
少數服從多數.
主席應該宣佈議案應該獲得通過.
那四個反對的又如何?.
服從議會的規定.
走到外面.
你都不能說你不反對.

$^{401}$你一定說你贊成.
是不是?.
這個是遊戲規則.
不過這個堂主任當時說了一番話.
他說既然這個議案這麼重要.
對教會有這麼大的影響.
如果大家不能夠同心的話.
就算通過了.
我都不能夠坦白地告訴全會眾.
這個是我們全部執事會的一致決定.
因為有分期的.
六對四.
他接著提出一個臨時動議.
這個臨時動議就是.
押後這個決定.
雖然通過了六對四.
他說押後決定.
再給這四位反對的執事.
說一下他的反對意見.
再給贊成的執事.
對反對意見作出一個回應.
會議又再次為這個重大議題.
進入一個更加深入的討論.
不過當然氣氛好了很多.
大家都很具體地說出自己的憂慮.
反對的人的憂慮.
兩個小時之後.
大家為議案祈禱.
祈禱之後再次表決.
最後全體通過.
唐主任最後分享到.
他見到一個很大的關鍵.
就是信徒合一.
是教會的核心價值.
他寧願事公延遲執行.
都不想見到大家不同心.
為什麼呢?.
他說如果會友見到這個執事團隊合一.
會友會信服.
會友見到執事團隊不和.

$^{441}$會影響到會友的靈命成長.
他覺得你們這些屬靈領袖.
都不合一.
都不同心.
那我教會會怎樣呢?.
我跟誰呢?.
每個人都是屬靈領袖.
這個唐主任說完之後.
我當時立刻有一個反思.
唐主任沒有使用議會的權力.
去令到反對的執事團隊轉為同意.
事實上是6比4.
那4票的一定要同意.
這個是遊戲規則.
因為議會是集體負責制.
反對的執事團隊要服從遊戲規則.
出來面對會友.
你都要說我支持.
雖然你內心裡面反對.
唐主任分享到.
他不是想事工能不能通過.
他是想執事團隊需要有一個合一的心.
唯有一個合一的心.
才能夠承擔神給他們教會的意象.
這麼重要的議案.
不同心如何去承擔.
如何去執行.
一個少數服從多數的決定.
是不能夠讓教會承擔使命和意象.
唯有教會合一.
才能夠承擔神所交託給我們的意象.
唯有教會合一.
才能夠實踐愛神愛人的信仰生命.
我希望我們紅英家都一樣.
我們能不能夠有一個合一的心.
放下我們的權利.
就算六對四.
都可以放下.
我們是一間教會.
檢視信仰第三個過程就是要行動.

$^{481}$再繼續實踐信仰.
既然我們是一間教會.
既然在回想裡面能夠察覺到不足.
自然要向神認罪求赦免.
但並不表示認罪後就算了.
而是有悔改的行為.
很多信徒不明白.
或者很多信徒扭曲了我們的信仰核心.
經常都覺得.
信了耶穌便可以了.
我們學習的不是信耶穌得永生嗎?.
我信了耶穌.
我應該得永生.
信是要有悔改的行動.
不是你信.
信以及悔改.
你信了.
有沒有悔改?.
以前所犯的罪你有沒有重覆犯?.
還是你以前所犯的罪.
你悔改了之後.
認罪了之後.
你立志不犯.
這是我們的信仰核心.
這個悔改的行為就是回復.
以往所做的事.
現在不再做了.
以往所犯的錯.
悔改.
認罪悔改.
以後不再犯.
以前的日子做得好.
現在被淡化了.
神要求這福所教會.
再重拾昔日的熱情.
用行動去證明他們仍然是一間愛神愛人的教會.
換句話說.
不要讓屬靈生命再後退.
其實我們的屬靈生命是一個逆水行舟的屬靈生命.
你問問你自己.

$^{521}$一個星期回來教會多久?.
一個崇拜大概一個半小時.
一個星期被外面世界的影響.
影響了多久?.
你只是一個星期回來教會.
一個半小時.
然後就回到這個世界.
這個一個半小時.
要頂住你一個星期的其他時間.
是不是逆水行舟?.
世界對你影響大一點.
還是教會對你影響大一點?.
你能夠回來教會.
在這個一個半小時裡.
所能夠吸收到的東西.
所能夠領受的東西.
讓你可以抗衡這個星期.
這個世界對你的挑戰和引誘,試探.
是不是逆水行舟?.
所以你的屬靈生命能夠成長.
一定是恩典來的.
不是你的能力.
是神的恩典托住你.
神的恩典為什麼要托住你?.
是你愛祂.
才能夠是這個恩典.
你不回教會.
神的恩典怎麼托住你?.
你回來崇拜黑眼訓.
神的恩典怎麼托住你?.
你回來教會還看手機.
神的恩典怎麼托住你?.
究竟主席的祈禱是怎麼說的?.
聽道行道.
神的恩典才能夠托住你.
能夠真正悔改和行為自然有獎賞.
俗話說.
德勝的.
我必將神樂園中生命樹的果子賜給你吃.
經文所說的德勝.

$^{561}$含有戰爭的意思.
今天的思綱是萬世戰爭.
戰爭的敵人有兩個.
一個是魔鬼撒旦.
我們需要依靠神的大能大力.
才能戰勝魔鬼.
另一個敵人.
我們很熟悉的.
是誰?.
另一個敵人就是自己.
我們願不願意放下主權.
全心信服.
然後才能夠進行主導.
要做一個至死不渝的信徒.
一生忠心侍主.
這些人才能夠進入神的樂園.
而獲得生命樹的果子吃.
我再重複一次.
不是信了耶穌就得永生.
而是要信而悔改.
你沒有悔改的行為.
永生離你很遠.
主席祈禱的時候說什麼?.
你的信仰生命.
和你的恩典相不相稱?.
樂園原本是一個波斯的字.
是用來形容帝王的花園.
猶太人相信二人死了之後.
會到一個好像《創世記》二章八節所記載的.
伊甸園的地方.
啟示樂園代表一個沒有罪惡污染的環境.
在那裡人和神可以暢通無阻.
彼此感通.
有獎賞自然有懲罰.
教會和信徒若不悔改有什麼後果呢?.
約翰在第五節說出了後果.
把你的燈台從懸柱挪去.
什麼意思呢?.
表示教會實在有可能被主挪去燈台.
燈台的失去.

$^{601}$顯明了一間失去了愛心的教會.
一間失去了愛心的教會.
只是苟延殘喘.
名存實亡.
是主的一種審判.
一間實踐不到愛神愛人的教會.
它的光芒已經暗淡了.
它的熱力減退了.
剩下的價值和社會上任何一個團體沒有分別.
任何人都看不到.
感受不到.
這間教會有什麼信仰來行.
他回到教會.
不察覺到信徒正在敬拜.
只察覺到這群人只是在參與一個宗教聚會.
他怎麼能夠吸引到他們學務主呢?.
一間實踐不到愛神愛人的教會.
根本就不是教會.
只是一群人按著習慣.
聚集在一間建築物之內而已.
他們參加的不是一個敬拜.
不是回應神對他們的愛.
他們只是參與一個宗教活動.
一個令自己自我感覺良好的活動.
信徒說我有回教會.
我有回崇拜.
這是自我感覺良好.
不管他們在裡面笑得多辛苦.
有多忍耐.
他們參與的事工有多龐大.
有多少資源.
甚至他們的建築物有多華麗.
有多少空間.
他再不能夠向世界說.
他是一間代表主耶穌基督的教會.
沒有.
沒得說.
因為他沒有愛神愛人的特質.
同樣是一間教會的事奉只是為了責任.
為了點綴教會有這些使命.

$^{641}$奉獻只不過是另一種變相的廉價贖罪券.
禱告只是空洞公式化的內容.
或者將至高至大的神.
當成是一種燈神去祈求.
這間還是不是神要放在這個地區.
發光發亮的燈台呢?.
不是.
我的引言說到這所教會很危險.
我相信一間面臨被神挪走的教會.
不單止可以用危險來形容.
更加可以用甚麼?.
生死邊緣來形容.
被神挪開.
被神挪走.
神對這所教會的提醒.
是否都可以作為我們洪恩家的一個信仰反思呢?.
我們說很想洪恩家成為洪水橋的燈台.
這是我們很想的.
我們很想是否就可以成真呢?.
我們的想是否可以成真呢?.
不可以的.
我們的想是不可以成真的.
唯有我們淺行.
我們淺行.
我們要淺行神的真理.
我們淺行愛神愛人的真理.
我們才真正能夠成為一個燈台.
不被神挪開.
洪恩家的弟兄姊妹.
今天的講題是甚麼?.
重回信仰生命的起跑線上高.
在香港這個節奏這麼快速的社會.
我們很多時候都在衝衝衝.
最高目標當然是為我們的理想而奮鬥.
我們衝得這麼勤.
可以為我們家庭為子女為將來而奮鬥.
是對的.
不過很多信徒其實都是為生活為戶口為賺兩餐.
他們實在很難有機會停一停去想一想.
究竟我們做基督徒的終極意義是甚麼?.

$^{681}$甚麼叫基督徒的終極意義?.
作為一個信徒我一定要問.
我為甚麼要信?.
我的生命有沒有改變?.
我的生命變成怎樣?.
變得快還是慢?.
我是一個熟齡的嬰兒.
還是一個可以吃飯的成人?.
我可以為主做甚麼?.
我可以討主哪裡的喜悅?.
我深信忙碌的生活已經令我們不知如何反省.
你一年前的信仰根基是怎樣?.
有沒有反省過一年前你一個怎樣的熟齡狀況?.
我深信做基督徒這麼久就更加需要回顧我們的信仰歷程.
好像年年驗身一樣.
我們需要為我們的熟齡生命做身體檢查.
神對耶福教會的提醒就是.
讓我們學習到信仰.
檢視信仰的三個過程.
第一回想.
想一想我以前是怎樣?.
我現在是怎樣?.
第二認罪悔改.
第三行動.
完成了這三個步驟.
我們就可以告訴自己.
我今天基督徒的生命去到哪一點?.
我是否仍然是一個熟齡的嬰兒?.
我正在吃奶.
正在吃烏仔.
正在吃飯.
還是人可以在我身上見到基督?.
這是我們的目標.
我們很想人在我們的生命裡見到耶穌.
那麼人在你的生命裡見不見到耶穌?.
還是人見到你.
他也相信耶穌?.
他呀.
他說他是基督.
我就不做了.

$^{721}$有沒有聽過這樣的說法?.
上個月第三個主日.
關國部有一次大探訪的行動.
我知道洪恩嘉以前對關國部的探訪很熱心.
我聽過很多會友說.
以前關國部很興旺.
不單有很多人願意參與.
也有很多的果子結了出來.
讓教會成為一個祝福別人的平台.
能夠真正實踐作鹽作鋼.
今天我們能否找回這個初心?.
剛才惠玲說過.
我們關國部又再reform(改革).
你會不會參與?.
關國部未來的部署.
是會在每個月的第三周.
或是一些不定期的外出探訪.
昨天我們部長美蓮.
和我們當值的委員劉世雄.
探訪了一個會友.
探訪完之後馬上有果效.
Edmond告訴我.
這個會友今天有回來參加主日學.
你鼓勵他們吧.
這些探訪試探訪甚麼人呢?.
探訪一些很久沒有回教會的肢體.
探訪一些新朋友.
探訪一些對福音有興趣的朋友.
不過六月份第三周是父親節.
所以我們改為第四周.
你會否參與?.
或者你參與不了.
其實我都很忙.
你會否推介一些朋友給我們去探訪?.
有些人對福音有興趣.
我叫他回教會他又不肯.
不如叫我們教會的人去探訪他吧.
你可以為教會做很多事.
而最大的事就是你的心.
你的心.

$^{761}$無論你想參與或介紹探訪對象.
我們都鼓勵你去聯絡美蓮部長.
美蓮在不在?.
美蓮,站起來讓人認識你.
不然我想找美蓮,但不知道她是誰.
主耶穌曾經在馬太福音說過.
因為我餓了,你們讓我吃.
渴了,你們讓我喝.
耳人回答.
主啊,我們甚麼時候見你餓了.
讓你吃,渴了,讓你喝.
王回答.
我實實在在告訴你們.
這些事你們做在我弟兄中一個最少的身上.
就是做在我身上.
我們見不到主.
當你說我愛主,我愛神.
你做在弟兄身上.
弟兄姊妹,愛神愛人不是用口說的.
而是用行動去實踐.
你有沒有這份感動.
在教會裡,在弟兄姊妹當中.
為一個你認識或不認識的朋友.
遞上這杯涼水.
以我的經驗來說.
這是常態.
在教會裡,有一個選定的位置.
你坐在這邊,就不會坐在那邊.
你坐在前排,就不會坐在後排.
所以每個星期,你左右的人.
理論上都是離去那一批人.
不會有甚麼態度去改變.
你上一年,都是離去那一批人.
如果你坐在這邊,你認不認識那些人.
你坐在後排,你認不認識前排的人.
我們教會都不大,我們不是那些church.
可不可以認識一下他們.
為他們遞上一杯涼水,可以嗎.
不為甚麼.
只是因為服事的和被服事的.

$^{801}$都是神所愛的兒女.
服事的和被服事的.
都是神所愛的.
你就為他們遞上一杯涼水.
是否真的這樣遞上.
不是,你關心一下他們.
有甚麼需要我為他們代禱.
祈禱不一定要說出口.
在心裡,為他們默默祈禱都可以.
重要的是你有沒有這個心.
一年前,你坐在這裡.
你不認識他.
能否一年後.
無論你坐在哪裡都好.
你都彼此認識呢.
這個就是家.
這個就是屬靈的家.
是不是.
你回到這裡十年.
你和他都不認識.
你只知道這個會有.
叫甚麼名字.
住在哪裡.
這就是.
這是否你的弟兄姊妹.
我們一起去祈禱.
七恩主,我們今天領受了你的恩典.
成為你的兒女.
也是天國的子民.
我們有君尊濟私這個身份.
這一切都是我們不配得的.
但是你卻一次過給了我們.
你想我們怎樣回應呢.
聖靈,我求你進入我們每一個人的內心深處.
你向我們說話.
亦聆聽我們向你那種真實的回應.
今天早上.
今時此刻.
願你的愛傾倒在我們的生命裡面.
亦願我們可以踏出一步.

$^{841}$接受你的愛.
亦都實踐你的愛.
讓我們用行動去實踐那種愛神愛人的生命.
我們的祈禱是奉教主耶穌基督.
聖名祈求,阿們.
願主祝福大家.
\newpage



\section{民數記 24:1-25:3}
\label{sec:qiH_9wRKv74}
\textbf{2022.06.12 《豐厚?缺乏?》 民數記 24\_1,25\_3 (和修版) 曾敏姑娘}
\newline
\newline
連結: \href{https://youtube.com/watch?v=qiH_9wRKv74}{\texttt{https://youtube.com/watch?v=qiH\_9wRKv74}} ~~~~ 語音日期: 2022-06-12
\newline
\newline
\hyperref[sec:abVWbFOM8VQ]{\small{< < < PREV SERMON < < <}}
~
\hyperref[sec:index_chronic]{\small{[返順時目]}}
~
\hyperref[sec:index_scriptual]{\small{[返順卷目]}}
~
\hyperref[sec:L3GT0EvH0Vg]{\small{> > > NEXT SERMON > > >}}
\newline
\newline
民數記 24:1-25:3
\newline
\begin{longtable}{cl}
\hline
\hline
章節 & 經文 (和合本修訂版)\\
\hline
24:1 & \begin{tabularx}{0.7\textwidth}{X} 巴蘭見耶和華喜歡賜福給以色列，就不像前兩次去求法術，卻面向曠野。 \end{tabularx} \\ \\ \relax
24:2 & \begin{tabularx}{0.7\textwidth}{X} 巴蘭舉目，看見以色列人照著支派紮營。神的靈就臨到他身上， \end{tabularx} \\ \\ \relax
24:3 & \begin{tabularx}{0.7\textwidth}{X} 他唱起詩歌說：「比珥的兒子巴蘭說，眼目關閉的人說， \end{tabularx} \\ \\ \relax
24:4 & \begin{tabularx}{0.7\textwidth}{X} 聽見神的言語，得見全能者的異象，俯伏著，眼睛卻睜開的人說： \end{tabularx} \\ \\ \relax
24:5 & \begin{tabularx}{0.7\textwidth}{X} 雅各啊，你的帳棚何等華美！以色列啊，你的帳幕何其華麗！ \end{tabularx} \\ \\ \relax
24:6 & \begin{tabularx}{0.7\textwidth}{X} 如連綿的山谷，如河畔的園子，如耶和華栽種的沉香樹，又如水邊的香柏木。 \end{tabularx} \\ \\ \relax
24:7 & \begin{tabularx}{0.7\textwidth}{X} 水要從他的桶裡流出，種子要撒在多水之處。他的王必超越亞甲，他的國必要振興。 \end{tabularx} \\ \\ \relax
24:8 & \begin{tabularx}{0.7\textwidth}{X} 神領他出埃及，為他有如野牛的角。他要吞滅那敵對他的國，壓碎他們的骨頭，用箭射透他們。 \end{tabularx} \\ \\ \relax
24:9 & \begin{tabularx}{0.7\textwidth}{X} 他蹲如公獅，臥如母獅，誰敢惹他？凡為你祝福的，願他蒙福；凡詛咒你的，願他受詛咒。」 \end{tabularx} \\ \\ \relax
24:10 & \begin{tabularx}{0.7\textwidth}{X} 巴勒向巴蘭怒氣發作，就緊握拳頭。巴勒對巴蘭說：「我召你來詛咒我的仇敵，看哪，你竟然這三次為他們祝福。 \end{tabularx} \\ \\ \relax
24:11 & \begin{tabularx}{0.7\textwidth}{X} 如今你趕快回本地去吧！我想使你大得尊榮，看哪，耶和華卻阻止你得尊榮。」 \end{tabularx} \\ \\ \relax
24:12 & \begin{tabularx}{0.7\textwidth}{X} 巴蘭對巴勒說：「我不是對你所差遣到我那裡的使者說： \end{tabularx} \\ \\ \relax
24:13 & \begin{tabularx}{0.7\textwidth}{X} 『巴勒就是把他滿屋的金銀給我，我也不能違背耶和華的指示，隨自己的心意做好做歹。耶和華說甚麼，我就說甚麼。』 \end{tabularx} \\ \\ \relax
24:14 & \begin{tabularx}{0.7\textwidth}{X} 現在，看哪，我要回到我的百姓那裡。來，讓我告訴你這百姓日後要怎樣對待你的百姓。」 \end{tabularx} \\ \\ \relax
24:15 & \begin{tabularx}{0.7\textwidth}{X} 他就唱起詩歌說：「比珥的兒子巴蘭說，眼目關閉的人說， \end{tabularx} \\ \\ \relax
24:16 & \begin{tabularx}{0.7\textwidth}{X} 聽見神的言語，明白至高者的知識，看見全能者的異象，俯伏著，眼睛卻睜開的人說： \end{tabularx} \\ \\ \relax
24:17 & \begin{tabularx}{0.7\textwidth}{X} 我看見他，卻不在現時；我望見他，卻不在近處。有星出於雅各，有杖從以色列興起，必打破摩押的額頭，必毀壞所有的塞特人 。 \end{tabularx} \\ \\ \relax
24:18 & \begin{tabularx}{0.7\textwidth}{X} 以東將成為產業，西珥將成為它敵人的產業；但以色列卻要得勝。 \end{tabularx} \\ \\ \relax
24:19 & \begin{tabularx}{0.7\textwidth}{X} 有一位出於雅各的，必掌大權，他要除滅城中的倖存者。」 \end{tabularx} \\ \\ \relax
24:20 & \begin{tabularx}{0.7\textwidth}{X} 巴蘭看見亞瑪力人，就唱起詩歌說：「亞瑪力是諸國之首，但它終必永遠沉淪。」 \end{tabularx} \\ \\ \relax
24:21 & \begin{tabularx}{0.7\textwidth}{X} 巴蘭看見基尼人，就唱起詩歌說：「你的住處堅固；你的巢窩造在巖石中。 \end{tabularx} \\ \\ \relax
24:22 & \begin{tabularx}{0.7\textwidth}{X} 然而基尼族必被吞滅，直到何時亞述把你擄去？」 \end{tabularx} \\ \\ \relax
24:23 & \begin{tabularx}{0.7\textwidth}{X} 巴蘭又唱起詩歌說：「哀哉！若神做這事，誰能存活呢？ \end{tabularx} \\ \\ \relax
24:24 & \begin{tabularx}{0.7\textwidth}{X} 有船隻從基提邊界來到，要壓制亞述，要壓制希伯；他也必永遠沉淪。」 \end{tabularx} \\ \\ \relax
24:25 & \begin{tabularx}{0.7\textwidth}{X} 於是巴蘭起來，回本地去；巴勒也回他的路去了。 \end{tabularx} \\ \\ \relax
25:1 & \begin{tabularx}{0.7\textwidth}{X} 以色列人住在什亭，百姓開始與摩押女子行淫。 \end{tabularx} \\ \\ \relax
25:2 & \begin{tabularx}{0.7\textwidth}{X} 這些女子請百姓一同為她們的神明獻祭，百姓吃了祭物，跪拜她們的神明。 \end{tabularx} \\ \\ \relax
25:3 & \begin{tabularx}{0.7\textwidth}{X} 以色列與巴力‧毗珥聯合，耶和華的怒氣就向以色列發作。 \end{tabularx} \\ \\
[1ex]
\hline
\hline
\end{longtable}
$^{1}$主持人:各位弟兄姊妹平安!.
每一次能夠回來教會敬拜神都是一件很美好的事情,我們一起禱告:.
而且你藉著耶穌基督去救贖我們的生命,讓我們永遠屬你,能夠和你復和..
今天很願意你親自對我們每一個說話,特別願意聖靈在我們眾人心裡面去動功,.
讓我們聽得到上帝你自己的聲音,又願意聖靈去感動我們,回應神你自己的教導,奉耶穌基督的名字祈禱..
阿們..
今天看這個題目,看經文,好像不太知道想說什麼,這麼奇怪,這兩節好像沒有什麼相關的經文,.
大家一路聽下去就會很明白,完了之後你應該會找到答案,什麼是豐厚,缺乏什麼呢?.
這段經文,教會或我們紅印堂不算是很多看的經文,但是裡面的屬靈教導是非常豐富,.
如果說真的,如果我們有時間,應該是由民數記的22章,一路看下去,看到25章,內容是很豐富,.
裡面有很多故事說,第一件事,什麼是豐厚呢?我們就看看,民數記說的是兩次,.
以色列民數點文素當中發生的事,這兩次數點文素都是以色列民在曠野發生的,.
以色列民在曠野漂流了40年時間,前後分別有前後各一次的數點文素,.
中間發生了什麼事呢?裡面有很寶貴的教導,特別是說以色列民對上帝的信,.
信與不信究竟是怎樣呢?其中一段是很特別的,說巴蘭和巴勒的故事,.
當中還有一段說奴的故事,一隻會說話的奴,這個故事我就不詳細說了,.
反而說以色列民在經歷什麼呢?以色列民一直在曠野遊走,直到現在這裡,摩亞平原這個地方,.
其實距離他們要進入迦南地的地方不遠了,很近了,但是在這個時候,他們又面對另一個挑戰,.
我相信他們是沒有什麼心理準備的,甚至有機會在事情發生之前,他們根本不知道會這樣發生,.
為什麼呢?說的是一大群以色列民去到摩亞平原的時候,摩亞王一看到,.
哇,好多人啊,以色列民非常多,摩亞看到這麼多百姓都非常害怕,.
那個背景就是原來之前上帝幫助以色列民打敗了阿摩利王西弘,原來這個阿摩利王西弘其實是很有名氣的,.
原來阿摩利人都很厲害的,在那個地方稱王稱霸,竟然以色列民可以打敗他們,那以色列民就更加厲害了,.
以色列的上帝就更加厲害了,所以大家對於這個以色列民是有些避忌的,這麼大群人在摩亞平原,會發生什麼事呢?.
當時他們就很擔心,大家聽我說就行了,因為我沒有記錄這麼多的經文在這裡,在民數記22章3字那裡這樣說,.
摩亞因為以色列百姓這麼多,非常懼怕,甚至是憂懼,摩亞對於米顛的長老這樣說,.
他說現在這一群人簡直要舔盡我們四處的一切,好像牛舔盡田間的草一樣,他們來吃光我們的東西,他們來總之就是我們的威脅,.
對我們來說是有威脅,是很不安全的,於是這個摩亞王巴勒,他就去召一個外邦的先知巴蘭尼,要去詛咒以色列民,.
很特別,在這裡你看到,這個摩亞王不是堂堂正正的,好好的帶著他的軍隊打仗,你可以的,帶著軍隊打敗以色列民就可以了,.
他不是,他就找人詛咒他們,因為害怕,覺得對方應該很有實力,就這樣打,硬碰硬未必會贏,詛咒他們令民族崩潰,軟弱,這樣就可以勝過對方,.
他當時是這樣想,所以就找人詛咒,找巴蘭尼,你一聽他講他是什麼先知,大家很容易就想,這個一定是屬於上帝,又不是,.
他不是屬於上帝的人,他本身是一個外邦人,是幼法拉底河的一個被奪的地方的人,他擅長什麼呢?用法術,用占卜之術去祝福和咒咒,.
應該他過往的歷史,他的祝福和咒咒對人是有一定影響力的,所以他就很出名,那些人如果想去祝福或者咒咒人,就會去找巴蘭,所以現在這個巴勒王就是這樣做,.
如果我要去對付這個以色列民,就是要找一個他很厲害的法術方面的人,這樣過來幫忙,大家記住,他那個巴蘭的能力不是出自上帝的,你不要被先知的名影響了,.
他有能力,但不是來自上帝,這個肯定是上帝敵對陣營那邊的能力,那照他來去詛咒以色列民的時候就很有趣了,因為經文很長,我把它簡單分了幾個段落,.
巴勒當然是開了很豐厚的條件,去給巴蘭一些過濾就可以了,我就要給很多很多這些東西,甚至差不多可以說是可以開一個空頭支票,你想要多少我都給你,巴蘭當然心動了就去了,.
一共有四次很特別的,巴蘭的發言,有四次他對著以色列民要說一些話,四次本來是要去咒咒的,本來巴蘭看著以色列民就要詛咒他們了,.
但是那四次裡面上帝都轉化了他的嘴巴,去說出一些祝福的話,在第一次的時候很特別,我想帶大家看一下,第一次的時候,他獻上他的祭物,.
看著一個地方,然後就去咒咒,但是很特別,經文去形容上帝將他的說話放在巴蘭的口裡面,這個是非常特別的,你記住巴蘭本身不是一個跟從上帝的人,.
他一向都是用其他方法,占卜的方法,是跟上帝作對的方法,跟上帝不喜悅的方法去做事,所以這個人不是屬於上帝,但是上帝都有方法將他的說話放在巴蘭的口裡面,.

$^{41}$本來他要咒咒,到他一說出來的時候,是一番祝福的話,那番祝福的話很簡單,就是說了一件事,就是上帝沒有咒咒的,巴蘭都不可以去咒咒,所以今天上帝是要祝福以色列民,巴蘭是不可以咒咒以色列民,.
還有,他說到誰能夠數點雅各的塵土呢?其實是說上帝要祝福雅各,祝福整個以色列民的人數眾多,你要好像塵土那麼多,你數到塵土,你便數到以色列民有多少,在這裡看到,神反而要去祝福他的百姓,.
好了,這一次完了之後,這個祝福完了之後,巴勒,就是摩訶王就很生氣,有沒有搞錯?我請你回來,不是要你去祝福他們的,要我去咒咒,幫我打敗他們,這個位置不是很好,可能上帝不喜歡我在這個位置去咒咒他們,我們換個位置,做同樣的事,獻祭,然後你咒咒他們,.
可能換了位置就可以,於是就換了位置,於是就第二次,第二次就很有趣,第二次又是,上帝又向巴蘭顯現,將說話放在他的口裡,又是很特別的,.
在這裡,上帝表明,他定義要賜福給以色列民,他定義賜福的是,人不能夠扭轉,在裡面的福氣很多,包括了什麼呢?是上帝的大令,上帝的同在,沒有人可以傷害到以色列民,.
還有就是在那個過程裡面,上帝特別提出,沒有任何占卜可以傷害以色列,所以今天巴蘭要召人去傷害以色列民,去咒咒,用占卜,上帝就是不讓他這樣做,不讓這幫人這樣做,上帝說不讓,這幫人就不會受到傷害,.
所以在裡面你看到,神其實是要保護以色列民,這是第二次,到了第三次的時候,是,很搞笑,經過兩次之後,又是這樣,他們覺得好像不太對,要換位置,又換第二個位置,希望這次巴蘭又可以咒咒,都是不成功的,.
這一次,上帝又透過巴蘭的口去祝福以色列民,用他們的,他說是雅各的帳篷,他們整個居所,他們所在之地,是有很豐盛的生命,那些圖畫非常漂亮,你看到,是很美麗的生命氣息在裡面,還有,以色列民是對付敵人很有力量,.
這幫以色列民根本就沒有心理準備,人家要找人去咒咒他們,去對付他們,傷害他們,沒有心理準備,上帝已經出手保護,將這些咒咒的說話轉化成為很大的祝福,然後到了第四次的說話的時候,就有一點特別了,.
上帝特別去透過巴蘭去說,將來有這個權柄,有這個政權會在以色列興起,其實這個是說將來尼賽亞的來臨,尼賽亞真正在所有的國家當中去統領,真正的我們整個宇宙的一個大君王,.
這個說將來,也說其他民族的情況,每一次都是對以色列民好,都是只有祝福,沒有咒咒,結果巴勒的計謀是不能夠成功的,整件事我強調以色列民是沒有心理準備,不知道人家要這樣害他們,上帝就保護了他們,.
還有巴蘭本來就是很想要巴勒的利益,太吸引了,是很豐富的,不只是錢這麼簡單,是有那種尊榮,誰不想要,他就特別想要,他很想要的了,在過程中竟然又被上帝這邊陣營的人影響了,.
在過程中我們值得有些問題去思考,第一件事,你透過上帝定義要去祝福以色列民看到什麼,神對他的百姓的心意是什麼,要賜給他們很豐厚的恩典,這是第一個問題,豐厚些什麼,上帝的恩典很豐富,上帝是要祝福他的百姓,.
而且不只是給一點點,是那個恩典可以傾出來,我們懂不懂得看到那個恩典的豐厚,還是覺得,應份的啊,上帝,我今天跟你,你不保護我,誰來保護我,我應份的,我得到任何從你而來的福氣都是應份的,.
對不起,沒有應份這回事,上帝沒有欠我和你,其實在過程中,你有沒有看到這個恩典是這麼豐厚,剛才那一句,他講到24章第一節,巴蘭建爺說喜歡賜福給以色列,上帝是喜歡給他的百姓福氣,.
那種喜歡是出面的,人家是看到的,外邦才看到,可想而知他多愛他的百姓,以色列人有沒有看到這份愛呢?有沒有看到自己得到的恩典是很寶貴,很豐富,是自己不配得的呢?有沒有呢?.
今天同友問我和你,我們得到上帝的什麼恩典呢?有沒有好好去數算,有沒有看到神給我們每一個都是很豐厚,他不只是祝福我們當中某位弟兄,某位姐妹,是祂祝福我們每一個,而那些恩典都是我和你不配得的..
我們試試去看回我們人生,我們人生從小到大經歷了什麼,當然不會每天都很開心,每天都很順利,我們有高低起伏,有歡樂的,歡笑的時間,也有眼淚的日子,但是你看回頭,我們經歷了什麼?.
在裡面如果給今天大家有一個很長的時間,整個星期給你數算上帝的恩典,你會怎樣去看?.
你說整個星期?可能我幾分鐘就說完了,我說整個星期太長了,我不會用整個星期去數算上帝的恩典,如果你這樣看的時候,看不到上帝過去在你生命裡面為你預備了多少的恩典..
我很感恩上帝給我有一個很寶貴的家庭,就是爸爸媽媽,姐姐和我,我們人少少,我們也很平凡,絕對不是社會上什麼出名的人士,也沒有什麼特別的技能,也不出眾,但是我很感恩上帝給我們這個家庭很豐富的恩典..
一直想起過去這麼多年,上帝給我們很多很美好的回憶,生活的共同經歷,看到上帝在我們的家庭裡面,給每個人都得到很多的恩典,不只是供書教學,不只是生活的東西,是還有很多上帝自己給我們的那些特別的經歷..
如果你讓我去數算,幾分鐘絕對是不足夠的,是真的可以,至少用一個星期慢慢慢慢說,我們整個家庭經歷了什麼,是非常寶貴的..
在過程中,我覺得最豐厚的恩典就是上帝給我們整個家庭得到救恩,能夠一個家庭一起,分別是在不同時間,我們信耶穌,得蒙救贖,這個是無比的,上帝是很愛我們一家..
今天我說的不是我們有什麼大富大貴,不是說明利,也不是說什麼很出眾,但我就是在說上帝愛我們一家,救贖我們一家,這個已經是最寶貴的禮物,是最好..
在這裡,我就是很想說,就是為了這一切,值得我和你去感恩,我想起這些,心裡也很歡喜快樂,至少我知道上帝很愛我的一家..
弟兄姊妹,你們有沒有想過,你們是不是真的深深感受到上帝很愛你們,你們真的覺得上帝的恩典很豐厚,還是不是?.
總之,我一看著隔壁的人,我就覺得上帝欠我的東西,總之我們心裡很不平衡,我看著他,為什麼我不是有那些,或者看著另一個人,又沒有那些..
上帝不公平,我數了ABCDEFG這麼多東西,你如果全部給我,那你就很愛我了,那你猜上帝怎麼回答你?.
其實,兒子女兒,ABCDEFG要給你有多難,但這些是不是對你最好?我給你的東西比這些還要好,你們有沒有看過?.
我們團契聚會,弟兄姊妹走在一起,有多少講感恩的事項呢?.
今天很多詩歌都是講感恩,大家有沒有留意?.
我們有多少講感恩的事項?數算上帝的恩典,上帝給我們這些,給我們那些,很快樂的,特別是上帝給我們救恩,今天是聖餐主日,.
有沒有為了上帝救贖我們,和我們的家人而歡喜快樂呢?值得大家想一想,這件事和後面是有關連的,.
如果我覺得上帝給我的恩典很豐厚,我很感恩的時候,後面的那件事應該是可以避免到,盡量避免不要發生,.
雖然很難確保人不會軟弱犯罪,但仍然有機會避免到,為什麼呢?後面發生什麼事呢?.
由22到24章的時候,是講上帝怎樣保護以色列民,離開巴蘭的詛咒,扭轉他們的嘴,.
在這件事上可以看到上帝是多麼的厲害,上帝掌管萬民,不只是掌管熟悉的人,不只是掌管基督徒,上帝也會掌管非基督徒,.
其實上帝真的超越萬有,在眾人當中不斷工作,所以今天大家想想自己生活的環境,.
不要只覺得,我肯定耶穌只是在我們教會裡面工作,只是在基督徒當中工作,.

$^{81}$對於那些未信耶穌的朋友,上帝不會在裡面工作,對不起,一定不是的,神也會這樣工作..
後面說回那件事,那件事反映出我們是不是真的覺得上帝給我們的恩典很豐厚呢?.
在民數記25章的時候,就是另一句話,那裡發生一件非常悲慘的事,悲慘不是以色列民受傷害,而是他們犯罪了..
之前還好好的,在經歷上帝的恩典,很豐厚的恩典,上帝的保護,其實他們應該很感恩,應該很快樂才行,.
是不是?自己的上帝簡直是保護自己,離開別人的詛咒,經歷這麼大的事情,要詛咒幾次,幾次都被上帝保護了,.
本來應該感恩,但是很快樂,轉過頭來,經歷了之後,25章的時候,你看到以色列就跌倒了,跌倒了是為什麼呢?.
25章是說到以色列人住在什廷這個地方,百姓就開始和摩訶女子行淫了,.
這些女子請百姓一同為他們的神明獻祭,百姓吃了祭物,跪拜他們的神明,於是就去到剛才那一句,.
以色列與巴力比爾聯合,耶和華的怒氣就會向以色列發作,這個聯合的意思就是立約,.
之前還好好的經歷上帝,下一刻就跌倒,以色列人拜偶像,不只是拜偶像,是和摩訶女子行淫,淫亂,是雙重的,.
身體方面的淫亂,也是靈性方面的淫亂,是這個令上帝非常憤怒,為什麼呢?.
我們明明是在經歷神的恩典,是很興奮,很雀躍,下一刻就會犯罪,這麼容易?為什麼?.
你問一問,以色列人是否真的覺得上帝的保護這麼好?是真的覺得上帝的恩典這麼豐厚?還是覺得沒所謂,你給我就要,覺得我應份,很特別嗎?.
加上你的保護,可能心底裡就是這種心聲,不懂得珍惜,不懂得感恩,看到上帝這麼好,所以下一刻很快就犯罪,.
在這次拜偶像的整個情況,行淫亂的情況,以色列人是很慘的,因為上帝是聖潔的神,一定會施行審判和管教的,.
結果這次上帝令他們經歷瘟疫,因為瘟疫而死的有二萬四千人,所以以色列人在這件事情上是一個很大的跌倒的情況,.
缺乏感恩,我們不是缺乏物質,不是缺乏福氣,也不是缺乏祝福,缺乏感恩的心,缺乏一個心以愛還愛,去回應上帝的愛,.
我們沒有那份心,不感恩,也不覺得很寶貴,總之就是覺得應份,你給我的是應該的,所以就會這樣跌倒,生命是會這樣軟弱,.
弟兄姊妹,我們問問自己,我們得到多少恩典,我們是否覺得很珍貴,是以愛還愛的呢?還是我們很容易陷入罪惡裡面,很容易跌倒,很容易背叛上帝,令上帝傷心呢?.
我舉一個例子,在這段時間,我很想邀請弟兄姊妹多點為我們的崇拜去禱告,崇拜是我們很重要的一部份,每個星期回來我們敬拜讚美我們的上帝,上帝是唯一的主角,回來我們獻上頌讚,.
由主席的頌讚禱告開始,到詩歌敬拜,一直到講完的訊息,每一部份都要令人將心歸給上帝,將榮耀歸給上帝,弟兄姊妹回來是應該要預備好自己的心靈,是要令上帝喜悅的,整個崇拜是這樣,.
而不是反過來的,這個崇拜是否令我感覺到很快樂呢?我有沒有聽過一堂很精彩的道,最好要動督笑,每隔幾分鐘就有一個爛笑,我笑完之後,你講得很精彩,前面各方面的事奉人員表現得很好,好像一場show一樣,.
我今天回來很滿足,滿足不是靈裡面和上帝親近的滿足,而是感受到我被服侍的快樂,不是這一種,崇拜不應該是這一種,崇拜應該是以神為中心,.
而背後就是你覺得上帝很佩得你敬拜讚美,不是因為有人叫你要這樣敬拜讚美,一直拉著你,扯著你,是你自己覺得上帝很佩得,我經歷了很多恩典,回來不讚美上帝簡直是對上帝不滿,大大聲敬拜讚美,.
不用別人叫的,全程投入,準備好自己整個人,不是好像剛剛起床,整個人都不是狀態,這樣來迎見上帝,同樣道理,感恩的心,你以你的愛回應上帝,你認識上帝是這樣佩得你敬拜讚美的,.
今天我們有沒有感恩的心,我們有沒有缺乏了這一份心,為什麼我總覺得好像沒有什麼經歷上帝恩典,有沒有很好的數算,盼望我們教會有這麼好的文化,數算神的恩典,看到上帝恩典的風口,.
神的恩典這麼好,以至我整個人生命的狀態,我都要將最好的給上帝,我也要遠離罪惡,我們就以以色列民在十庭的失敗經驗引以為鑒,我不覺得我們比他們好很多,.
我們可能說,我一定不會的,我告訴你,什麼凡間人,我不知道多麼聖潔,我一定不會拜偶像的,真的,他們很壞,我們很好的,盼望他們的經歷,令我們可以引以為鑒,.
也要看自己,為什麼沒有那份熱情,為什麼沒有那份熱情,我們需要有那份對上帝的熱情,我們需要重新起神,我們一起禱告..
令上帝很生氣,很難受,但天父你仍然要愛這班百姓,你定義要賜福氣給他們,只可惜百姓看上帝你的恩典,沒有看到這麼豐厚,以致他們很容易跌落罪惡裡面..
主下,今天求你讓我們引以為鑒,讓我們看到我們自己是一個怎樣的人,你開我們的眼睛,讓我們看到上帝你的恩典是何等豐富,讓我們看到神對我們多好,你是怎樣救贖我們的生命,我們是何等不配,.
就讓我們很被激勵,要去愛你,很被激勵要離開罪惡,很被激勵總是要預備好自己,要去回來將最好的獻給你,特別我知道上帝在崇拜裡面,總是在等我們,求神興起我們,賜給我們有那種真正愛你的熱情,.
讓我們回來的時候,將最好的獻給你,主下不單止這樣,你幫助我們忠於你,我們遠離身體的淫慾,也遠離我們屬靈裡面的淫慾,讓我們成為討上帝你喜悅的生命..
我將同恩家交在神你的手裡,求你賜福給我們,你讓我們整個群體有一個很大的復興,神你得著我們,你也讓我們得著你,從今以後不再一樣,主願你在我們當中工作,奉耶穌基督的名字祈禱,阿們..
阿們..
(影片完結).
\newpage



\section{約書亞記創世記 5:21-24}
\label{sec:L3GT0EvH0Vg}
\textbf{2022.06.19 《我和我一家必定事奉耶和華》 創世記 5\_21-24;約書亞記 24\_15 (和修版) 高淑芬傳道}
\newline
\newline
連結: \href{https://youtube.com/watch?v=L3GT0EvH0Vg}{\texttt{https://youtube.com/watch?v=L3GT0EvH0Vg}} ~~~~ 語音日期: 2022-06-22
\newline
\newline
\hyperref[sec:qiH_9wRKv74]{\small{< < < PREV SERMON < < <}}
~
\hyperref[sec:index_chronic]{\small{[返順時目]}}
~
\hyperref[sec:index_scriptual]{\small{[返順卷目]}}
~
\hyperref[sec:PeueLB8OHRY]{\small{> > > NEXT SERMON > > >}}
\newline
\newline
$^{1}$大家都知道教會全年的主題是「毀身事奉」.
在這書中,我覺得毀身事奉是需要在座的爸爸,弟兄.
在家中的很重要的爸爸,如果你說「我和我一家必定事奉耶和華」.
這句話是更加有威力的.
今天祝大家父親節快樂.
因為在父親節中,我們會想起天上的爸爸.
還有我們肉身的爸爸.
在《箴言》也有教過我們.
在十章說道「至畏諸子使父親歡樂」.
不知道我們有沒有令到天上的爸爸歡樂.
有沒有令到我們肉身的爸爸同樣的歡樂呢?.
在這段時間,我要教一些剛來香港的泰國小朋友.
一個八歲,一個九歲,來自兩個不同的家庭.
剛好在一班.
我說「這個星期日是父親節」.
因為你不告訴他們,他們不知道.
泰國的父親節是在12月5日,是第九任皇帝生日的日子.
所以是不同我們香港的.
我只是給這個題目,順便教了四個字.
就是「我愛爸爸」.
中文字是他們寫的.
但題目是「我愛爸爸」.
他們連媽媽也加進去.
我覺得他們非常好.
因為他們的觀念很正確,我很欣賞.
因為一個家庭真的要有爸爸,要有媽媽.
我們做子女的,我們要愛他們.
這個很特別.
「Family」這個字是由爸爸和媽媽的字加上「I love you」的字頭.
就變成了一個家庭.
彭若,今天你仍然是一個健全的家庭.
是一個感恩的福氣.
已經是上帝給我們在世的福氣.
當我們思想到,我們作為子女的時候.
我們怎樣對我們的父母.
因為這是一個家庭裡面兩位很重要的人物.
這張照片是我拍的.
因為這段時間經常要由天水圍搭車去九龍城.
我發現爸爸或者媽媽.
當小朋友很累的時候.

$^{41}$爸爸的手,媽媽的肩膀.
或者有時候是手抱.
都是讓小朋友好像躺臥在青草地.
每一個小朋友見到自己的爸爸媽媽.
就覺得這個世界很好.
不知道你有沒有感受到.
如果沒有感受.
或者星期六來體驗一下.
因為星期六教會有很多爸爸媽媽在.
因為他們要等他們的子女放學.
很多爸爸媽媽在.
子女見到自己的爸爸媽媽.
他們很開心.
你就會看到昔日我們的父母.
怎樣由我們小時候培養我們.
成為現在的大人.
在聖經裡面我找到八個原則.
怎樣讓我們的父母.
聖經教我們不要讓我們的父母有憂傷.
當然不要偷父母的東西.
不要輕慢父母.
因為到我們長大的時候.
小朋友在十歲前.
大部分都很愛父母.
但當他去到青少年的時候.
就開始喜歡自己的老友.
自己的朋友.
我們都經歷過這些時間.
有時候父母不明白就輕慢.
小看父母.
或者我們的教育程度已經比父母好.
我們又小看.
其實我們的父母.
我們是不得輕慢.
當然不要打父母.
不要咒詛.
我們要孝敬我們的父母.
我們要在主裡面聽從.
在主裡面尊敬我們的父母.
因為每一個我們都會有我們的父母.

$^{81}$我們求神給我們.
我們特別是上帝的兒女.
我們更加求神給我們.
上帝很著重給我們有育身的家庭.
也讓我們好好地做子女.
今天是父親節.
所以很想和大家介紹一位很經典的父親.
這位經典的父親叫做爾諾.
爾諾你會覺得他很厲害.
在聖經裡面沒有經歷過死亡的只有兩位人物.
他是其中一個.
就是爾諾.
他沒有經歷過死亡.
他就直接被上帝接去.
我們不知道當天何時來.
我們在座還是有機會.
常存盼望.
因為不知道何時我們會被提.
爾諾就經歷過被提.
還有昔日的先知爾利亞.
先知.
聖經裡面只有這兩位人物.
是沒有經歷過死亡的.
被主接了去.
另外爾諾亦是上帝親口說.
他是喜悅的人物裡面.
在聖經裡面只有兩位.
一位就是主耶穌基督.
當祂洗禮的時候.
天上來的聲音.
「這是我的愛子.
我所喜悅的」.
你想想.
天上的父親.
為我們這些不配的罪人.
做得多徹底.
把祂所喜悅的獨生子.
為了我們.
讓祂降生.
釘十架.

$^{121}$贖罪.
為了我們這個不配的人.
耶穌固然是祂喜悅的人.
但想不到爾諾同樣.
是天父上帝喜悅的其中之一位.
所以我們會看到.
這位爾諾很經典.
但經文不多.
只是說他在舊約裡面.
剛才我們的主席也幫我們讀了出來.
但是在新約裡面.
他同樣會排行在希伯來書的偉人榜裡面.
爾諾有提名.
上帝也密視了.
當時寫聖經時也有紀念他.
甚至在尤大書裡面.
同樣有詳細紀念他.
究竟爾諾與神同行時.
遇到些甚麼.
究竟爾諾是誰呢?.
他是阿當.
很早期的.
在《創世記》裡便會看到.
是阿當的第七代孫子.
阿當和他一起在地上生活過.
因為我計算過.
阿當享受了有九代同堂的時刻.
我不知道大家能否看到我的數字.
照著聖經裡面的歲數.
他何時生了誰.
很好看,很好玩.
這樣計算便會知道.
他們有多少代同堂.
我們知道爾諾是阿當的第七代孫子.
在這個時候.
我們會看到.
當時.
雖然他們是第一代的人.
是上帝親手造的人.
他是第一代.

$^{161}$但想不到在這個短短的家譜裡.
我們會發現到當時的人已經有很多罪惡.
或者有很多不好的事情發生.
爾諾究竟何時才開始和神同行.
聖經裡面剛才有讀到.
他已經65歲.
我們的書有些已經到65歲.
有些還未到.
所以我們要想想自己究竟.
我們開始了同神同行沒有.
每一個都需要.
特別是我們的父親,我們的弟兄.
他的轉捩點.
我們看到爾諾的轉捩點.
在《創世記》五章二十一節裡說到.
當他生了馬土撒拉的時候.
他就開始與神同行.
不知道會不會像最近新聞裡說.
有些快要做爸爸了.
他去學習怎樣培育.
因為他們覺得他們要有些付出.
不要只有太太一個人承擔所有.
另外一個就是.
他不想子女只貼著媽媽.
不貼著爸爸.
所以他們就轉化子孫去學習這件事.
我們看到馬土撒拉.
他出生之後.
竟然成為了爾諾要跟從神的方向.
甚至與神同行了三百年.
他的壽數是三百六十五年.
如果在阿當的家譜.
這九代的家譜裡去看.
他屬於比較年輕.
過世了.
但是他受到上帝的喜悅程度是高的.
甚至他是一位與神同行的爸爸.
我相信我不知道我們的弟兄.
你做爸爸的時候.
你的感受是怎樣呢?.

$^{201}$很緊張還是怎樣?.
要做些什麼?.
每一個宗賢.
雖然是太太生孩子.
但是爸爸的心情.
不知道會不會像片段一樣.
都很緊張.
但是緊張只是一剎那.
我們的父親.
我們的爸爸有沒有堅起.
做這個爸爸的責任呢?.
付諸下來的就是那種責任.
有一天.
我問問我爸爸.
爸爸你覺得做爸爸容易嗎?.
他就點頭了.
我以為容易就完了.
我問他答案.
他就說做爸爸很容易.
總之有兒子有女兒.
生了就是爸爸.
不過要做一個負責任的爸爸.
就不容易了.
他又說他的陳年往事.
因為我爸爸有很多親戚在印尼.
有一次他去了印尼.
那間屋有張很大的海報.
好像一比一般大.
我爸爸看著.
這個男士是誰呢?.
以為是.
就問這個是誰來的?.
怎知道印尼的親戚說.
你是不是香港來的?.
這樣都不認識?.
這個是你們的Andy Lau.
劉德華.
我爸爸真的不認識.
因為我爸爸就說.
如果我真的認識那些.

$^{241}$你們就不會帶了.
因為他要工作.
掛念去謀生.
沒有娛樂沒有看電影.
我怎麼認識這些明星.
所以他要犧牲.
過去可能沒有子女的時候.
怎樣都可以.
起碼你賺的薪水.
都是你自己去享用.
但是當你有了子女的時候.
其實我有時覺得父母很偉大.
很多東西他不捨得用.
都要給女兒給兒子.
特別是星期六.
你會看到有些學跆拳道.
有些學畫畫班.
有些學芭蕾舞.
我看到是很願意付出.
為子女.
你想想.
你辛辛苦苦賺的錢.
就是為了你的子女.
這個就是有責任的爸爸.
有責任的媽媽所會面對.
原來做一個有責任的父親.
是很難的.
也真是的.
在我們人生裡.
學生有學生手冊.
當我們不懂得用這個喇叭的時候.
都會有一個使用手冊.
我想問在座的爸爸.
你有沒有收過一本叫做爸爸手冊呢?.
或者媽媽手冊呢?.
沒有.
我們都是一邊做爸爸一邊做媽媽.
一邊去摸索怎樣走.
當以諾生了馬土撒拉的時候.
因為在當時.

$^{281}$舊約時代他們很緊張.
究竟我要幫我的兒子.
我的下一代.
叫做甚麼名字呢?.
上帝竟然啟示叫做馬土撒拉.
知道嗎?.
馬土撒拉在聖經裡有多少個?.
如果我們說約書亞.
你就會看到聖經裡有很多個約書亞.
叫做彼得.
或者到現在我們很多人都叫彼得.
叫David 大衛.
很多這些名字.
馬土撒拉是不會有人要這個名字.
也不會改這個名字.
馬土撒拉的意思就是死亡的意思.
馬圖死亡.
撒拉就是將會有件事發生.
即是你的兒子死了的時候.
就會有事發生.
而當以諾再深深地.
受到上帝的提示的時候.
更加知道這是一個審判.
如果你聽到這個信號的時候.
你是否更加要好好地跟著神?.
因為怎樣去養育這個兒子呢?.
不要讓他輕易地死.
三歲死 接著審判就來了.
所以很奇妙.
馬土撒拉是長壽的.
969歲.
他剛剛在羅亞出生.
羅亞還要是五百歲的時候才生兒子.
接著就生了三個兒子.
在羅亞六百歲的時候.
馬土撒拉.
你看那些數字.
就是剛剛好.
他969歲的時候.
當洪水降臨地面的時候.

$^{321}$馬土撒拉就死了.
你會看到上帝的審判.
就是在那時候臨到這個泉地.
羅亞得到以諾的家族.
一直這樣.
你會看到聖經裡面很多的養兒育女.
養兒育女.
如果在翻譯原文.
意思是成為一個爸爸.
不是一個名銜.
是一個有位份 有責任.
裡面當然有權柄.
是有一個爸爸本身的位份.
是沒有人可以代替的.
昔日以諾恩著這個啟示.
他就好好地.
真的不知道怎麼做.
怎麼走前面的路呢?.
他選擇了與神同行.
我告訴你.
雖然他是第七代.
當之後第七代的人類.
但是在眾言.
在這個世界裡面.
第七代的人類也好.
人就好像聖經裡面所說.
人比萬獸更加惡.
人的罪已經在第七代裡面亂七八糟.
第二代大家都知道.
已經發生了這個謀殺事件.
第一代就說謊.
推卸責任.
所以以諾要與神同行.
其實不是一件容易的事.
今天其實上帝都會透過聖經.
透過一些事情.
透過一句話.
讓你去轉變.
你不要再靠自己了.
你不要聽世界對你說的聲音了.

$^{361}$你要轉過頭來.
跟從我.
我知道今天坐在這裡.
你是相信上帝.
但是你有沒有與神同行呢?.
剛才我們領師的姐妹帶來的那首歌.
很好聽的.
在花園裡.
你有沒有經歷過呢?.
我們香港這麼嘈雜.
哪裡有個花園.
花園都人山人海.
對不起.
當你真的與神同行的時候.
在花園裡面.
你真的看到.
現在我不是這麼早起床.
但是你真的下完雨.
那些水點在花葉裡承托住.
雖然外面很嘈雜.
很亂.
但這是一個靜觀的操練.
當你真的有神在裡面.
跟你一起行的時候.
可以與上帝一起去行.
例如去花園.
屯門花園.
或者元朗花園.
我們香港有很多花園.
或者你種的盆栽等等.
你從這些上帝所創造的.
你都會經歷到.
在花園裡.
上帝與你一起.
並且要告訴你.
你是祂的好朋友.
你是天父的愛子愛女.
今天我們都說.
我們是相信神.
求神給我們更加的跟從神.

$^{401}$更加的聽著主.
讓我們知道.
我們下一步是怎樣去行.
或者這些事情是怎樣去面對.
真的要問一下自己.
以諾他六十五歲.
與神同行.
你呢?.
你開始了沒有?.
你開始了與神同行沒有?.
我們說要全然偽生.
很重要的.
我想我們教會不是要大家來做.
幾個小時的兼職.
幾個小時做一下這個.
做一下那個.
不是的.
與神同行.
與上帝一起.
這個就是全然的偽生.
一個很基本的第一步.
求神給我們紅恩堂的弟兄姊妹.
開始.
開始的就繼續了.
未開始的就快快要開始了.
與神同行.
無論我們香港的節奏是很繁忙.
你是可以與神同行的.
我們遇到很多棘手麻煩的事.
你同樣都可以與神同行.
讓神成為我們一起的同路者.
以諾又怎樣與神同行呢?.
我要與神同行.
再過多三年我就退休了.
我退休了.
我就可以.
不要說了.
全世界與神同行.
我退休了.
你退休之後.

$^{441}$可能你的孫子.
可能很需要你.
又或者你開始有很多.
過去拼搏的人生.
未完成的事.
你可能在那段時間就完成它.
你還可以有很多理由.
告訴自己.
還不是時候.
還不是時候的.
我告訴你.
以諾的時候.
那個時間.
那個社會的氣氛.
不是像我們現在去避世那樣.
大家都有經歷過.
去退休的日子.
退休的日子是很開心的.
那幾天如果是靜修營.
物觀營.
什麼都好.
你會覺得很好.
真的好得無比.
不用理家裡的事.
真的好好地去享受詩歌.
享受聖經.
享受與神的獨處.
是好得無比的事.
但是一下山.
一出營.
我又要做一個凡人.
對不起.
以諾是做凡人的時候.
一個普通的爸爸.
他就可以與神同行.
同樣我們要等到那一刻的時候.
我們不知道能不能等到.
但是今天.
我就可以好好地與神同行.
讓上帝成為.

$^{481}$今天活著的不在時.
我平時都可以唱歌.
可以經文.
今天就是這樣.
我們求神幫助我們.
以諾他是一個信心的偉人.
他在一個烏煙瘴氣的世界裡.
他如何與神同行呢?.
他的信很重要.
就好像希伯來書十一章五到六節說.
以諾是因著信被接去.
不至於見死.
人都找不到他.
因為神已經將他接去了.
只是他接去而先.
已經得到神所喜悅他的名證.
人非有信就不能得見神的喜悅.
因為到神面前來的人.
是必須信有神.
而且信他上次他那尋找他的人.
聖經裡.
上帝真的很愛我們.
給了我們很多很多的應許.
給了我們很多鼓勵.
在當時以諾的世界裡.
你看看.
很早的創世記.
耶和華已經看到地上的罪惡很大.
中日所思想的都是惡.
但是以諾.
你會看到到羅亞的時候.
他的真孫.
你會看到進去方舟的.
只是羅亞的一家.
不多的.
你就發現到.
以諾他生存在社會群體當中.
就像我們今天.
我們坐在這裡好像人多多.
但是當我們離開紅印堂.

$^{521}$出紅水橋.
出元朗屯門天水圍.
回到家裡附近的時候.
我們變得很少數.
有很多的聲音.
結了婚都不要生孩子.
自己賺的錢怎麼會給子女用呢.
一會兒他們未必會孝順我.
我倒不如好好地自己賺錢.
自己享受.
或者有些聲音說.
這個世界這麼多苦難.
生孩子?.
一會兒他們也一樣這麼苦難.
不要生孩子.
越來越多人結婚之後.
不願意生孩子.
因為他們沒有你們這麼有信心.
一個孩子四百多萬.
其實這是宣傳的.
是不是真的要這麼多錢.
才養一個孩子呢.
如果是這樣.
生孩子是不是變成有錢人的專利呢.
不是這樣的.
所以我們會看到.
以諾的世界裡面.
其實不容易.
就像我們這樣.
是不容易的.
他還是要在需要裡面.
依靠神的供應.
我想當你面對一個家庭.
每一個家庭都有一本很難讀的書.
很不容易的事情.
不能夠跟別人說的話.
但是你可以向神說.
在關係裡面.
當你做過父母.
你是父母的時候.

$^{561}$你就會知道.
你的開心很多東西.
都是子女帶給你.
或者孫子帶給你.
你就很開心.
他做得好.
你會替他開心.
很開心.
但是都是成對比的.
他有些事情.
他人家做.
不會傷到你的.
他做.
如果你的子女.
你的孫子.
做了一些這樣的事情.
你就會覺得.
那個心好像被刀割一樣.
在這個時候.
你會不會想到.
天父創造我們的時候.
其實都是這樣.
我們是他所創造的人.
但是我們對神.
是不是時刻都讓上帝得著歡喜.
得著喜悅呢.
當在無助的裡面.
我們有沒有堅固的.
相信神掌權.
掌握著呢.
昔日約書亞.
跟著摩西.
後來摩西死了之後.
他要成為領袖.
要帶著人們過約旦河.
也是一件不容易的事.
打了一場大仗.
很多很多場的戰爭.
他經歷到上帝是真的.
由他由十二個探子.

$^{601}$要去探路的時候.
他都看到.
我們靠的是上帝.
是的.
雖然外面的聲音告訴你.
這些又大又難的事.
但是在上帝裡面.
沒有難成的事.
只不過.
上帝的時間.
不是我們所想的時間.
弟兄姊妹.
主權在上帝的裡面.
在打完仗得到迦南地.
根本都沒有超過三十年.
以色列人又背叛神了.
又是各人.
行他所以為正的事.
以為其實是惡的事.
你會看到.
約書亞身為一個領袖.
他出來要開聲.
現在擺在你面前的.
就是.
你們要麼跟從神.
要麼你們繼續走你現在這條路.
但是上帝過去這麼多的恩典.
在我約書亞來說.
約書亞說他才有這個威力.
約書亞說.
不單是我一個.
如果我一個人事奉神.
很容易的.
但是是我和我一家.
以至到我後後代代.
好像以諾那樣.
他都要傳承與神同行的這件事.
去到他的後裔.
以至他說.
我和我一家都必定的事奉耶和華.

$^{641}$今天.
我們說與神同行.
在泰國裡面.
開車很塞車.
非常塞車.
我還要開一輛棍波車.
你會知道.
如果開棍波車.
你就知道要按壓一點.
走一厘米的路.
因為塞車.
我經常覺得塞車這段路.
好像在考驗我怎樣跟從神.
前面這個是神.
我要跟著祂.
但是又不可以碰到祂.
又不可以碰到祂.
但是又不可以被我旁邊左右的煞車.
是我的比喻.
他們不是真的煞車.
他們會過來.
因為當你有一厘米的空間.
他們就會過來.
然後你又要再排隊.
他們的車是塞得很厲害.
所以你要逆步逆隨.
又不可以碰到前面的車.
然後要貼一貼.
貼到不要有那一厘米.
即是半厘米.
不要被人家塞過來.
這樣跟從神.
在跟從神裡面.
我感受到.
我們很感恩.
同恩同.
我們都叫做有弟兄.
特別我們泰國祖族.
都有一位尼哥弟母.
但是你知道尼哥弟母.

$^{681}$來自於尼日利亞.
在他們的文化傳統裡面.
男性是高一點的.
沒辦法.
這是上帝給的遺訓給男士.
所以我們這裡的弟兄.
你知不知道.
上帝已經給了這麼好的遺訓給你.
你知道什麼叫遺訓.
或者角色.
或者什麼.
是不同的.
例如我爸爸媽媽.
不會聽我們說什麼.
叫他記得留心吃藥.
要做運動.
他不會怎麼聽.
醫生這樣說.
叫你這樣做.
他就會聽.
或者你的孫子子女.
做這個不做.
但是老師說要做.
老師說要買一個筆盒.
不是筆袋.
他會聽老師的話.
那就是那種角色帶來的權柄.
一個家庭父親的角色.
是母親無法代表.
沒有人代表得了.
在泰國裡面.
我曾經服侍一個叫孤兒院.
但我覺得他們都不是孤兒.
有時候放假.
這個探爸爸.
一會兒去到的時候.
看不到這個.
這個又探媽媽.
他們有爸爸媽媽的.
當然了.

$^{721}$不然他怎麼會出生呢.
不過問題是.
不理會他.
是沒有責任的父母.
以致他們成為孤兒.
有一次.
有一對的短孫女.
是親子團.
即是爸爸媽媽和小孩.
小孩很小.
都是十歲以下的小孩.
去到這個孤兒院.
親子團我們都不會要求太多.
我們要求.
不要嫌小孩黑漆漆.
牽手.
因為當時沒有疫情.
牽手就可以了.
你們唱這首歌牽著他們.
一開始是叫小孩牽著小孩.
誰知.
特別是去到男孤兒院.
因為通常孤兒院裡面.
做的都是女士.
他們很少見到男士.
還有這麼多個爸爸.
他們的小孩主動.
不是牽著很精靈.
很白皙的.
因為泰國人很喜歡白皙皮膚的.
不是牽著小可愛.
他們是靠近牽著爸爸的手.
泰國裡面.
很多很多變性人.
特別是男變女.
多於女變男.
為什麼呢.
因為他們沒有父親的形象.
所以今天.
我們的父親.

$^{761}$你不要和一些人去比較.
我們有上完毅班.
完毅治療班.
每個學生的材料都一模一樣.
但是出來的製作.
每個不同.
幾次的東西都是.
很多次的製作.
每個不同.
所以我告訴你.
每一個的爸爸.
每一個的家庭都不同.
但是上帝給你的位份.
絕對是當中.
可能.
這是高傳道猜測.
我真的不太行.
做燈泡的不是我.
是我老婆.
有蟑螂我找我老婆來.
可能是這樣.
但不重要.
你會覺得我真的不太行.
我真的好像廢柴.
對不起.
泰文廢柴.
是很重要的.
好像剛才你的歌.
求天父使我生活如光.
是一個廢柴.
你知道是什麼意思嗎.
泰文廢柴是一個電筒.
所以今天你覺得.
你自己廢柴的時候.
請你記住.
這個世代是黑漆漆的世代.
你有沒有試過.
很黑很黑的地方.
很小的電風扇.
去到最後一格.

$^{801}$有盞燈很小.
但你很黑很黑的時候.
因為我們泰國整天都沒有電.
那盞燈.
都可以救了我一個晚上.
很小很小的.
雖然你覺得你是一個大廢柴.
也好小的廢柴也好.
在這個黑暗的世代裡.
很有用.
所以.
可能姐妹你都覺得.
自己都是廢柴.
也不要緊.
你是上帝的子女.
廢極有限.
當你認識了上帝.
已經是智慧的開端.
所以你要想泰文的想法.
廢柴很重要.
是殿堂.
我們說為身事奉.
就是給上帝多踩一些蛾.
我們的姐妹知道甚麼意思.
就是用多些蛾.
踩就是用.
你願不願意給神踩.
願不願意.
我們的弟兄呢.
大聲一點.
你不是廢柴.
你是當管.
願意.
你願意要給神用.
特別.
我很想提.
你不是單單在教會裡給神用.
在你的家庭.
縱然你是一個默不作聲.
即是在家裡很安靜的人.

$^{841}$你在這裡已經成為家裡.
很重要的一個角色.
所以求神祝福.
我們雄恩家裡的弟兄.
因為上帝.
亦派了助手.
有你的太太.
有姐妹在.
我們一起.
事奉神.
給神踩.
這個踩不只是做甚麼.
做甚麼.
是我們的心.
願意與神同行.
我們不只是在教會裡一個樣子.
是我們的家庭.
在我們平時工作裡.
昨天有個姐妹很好.
她有些製成品.
她說要擺在公司裡.
擺在公司裡為甚麼呢?.
因為我家人都相信主了.
擺在公司裡.
人們看到這個製品.
這個製作.
他們就會有興趣.
可以問我.
我便可以跟你講福音.
是的.
有些人.
有個小朋友說.
爸爸.
我很喜歡你在教會裡的樣子.
為甚麼你們會笑?.
我還未說.
你們知道了.
爸爸.
這次真的廢柴了.
他都不知道.

$^{881}$即是怎樣?.
我在家裡和他不一樣.
我也是這樣的樣子.
他以為這個樣子.
這個樣子是沒有問題的.
不是的,爸爸.
你在教會裡.
你很有禮貌.
牧師.
傳道.
姐妹們.
叫你做甚麼.
你馬上去做.
在家裡.
我叫你幫我修理書包.
我扣爛了.
我講了一個月了.
你都未做.
還有.
你的聲音.
在教會裡.
很有禮貌.
但你有時經常都學媽媽.
我們求神幫助我們.
當與神同行的時候.
求神讓我們知道.
上帝在逆境裡面.
在黑暗的當中.
仍然和我們一起行.
願意我們都經歷到上帝.
即是說.
洪恩堂的弟兄姊妹.
你們是上帝所喜悅的.
我們如此事奉他.
是理所當然.
這個事奉是一齊回來敬拜.
也是的.
與神多些.
與神緊密溝通.
這也是本位的事奉.

$^{921}$再多些.
你可以參與一些事工.
一起參與.
我很希望我們的弟兄.
很雄壯.
只是一次.
不能預演.
我們的弟兄可否站起來.
給些鼓勵.
我們的姐妹做得很厲害.
我很想你.
說一句.
要雄壯一點.
要很實在地說.
就是這句.
好不好.
1 2 3.
我們一起祈禱.
我們親愛的天父上帝.
我們知道上帝創造我們每一個人.
我們的角色每一個人都不同.
因此主幫助我們.
當我們在一起的時候.
我們都是一個罪人.
叫我們不要比較.
叫我們讓上帝給我們一種獨特.
在當中去榮耀你.
主我們站在這裡的弟兄.
他可能已經成為父親.
又或者還沒有.
但是上帝給他一個男性的角色.
主啊.
無論是在他肉身裡的養育兒女.
或者在靈命裡.
去培育一些弟兄姊妹.
主啊.
你給他有做父親的權柄.
有做父親那種肯承擔一些責任.
主啊.
他願意的就是跟從你.

$^{961}$主啊給我們這群姐妹.
我們更加的去向著.
因為上帝你是我們的頭.
主啊.
你給到男士.
也都是成為家裡頭的時候.
主啊.
叫我們因著服從你而服從.
就是在地上裡的一家之主.
求主幫助我們.
讓我們每一個都是成為你所喜悅的子女.
主啊.
更加願意我們已經成為父親.
因為肉身的父親.
或者靈性的父親.
更加的每天.
每一天.
每時每刻與主同行.
讓我們每一個出去.
縱然是很細小.
很微小.
但是在這個黑暗裡.
哪怕我們只不過是很小的光線.
好像泰文說的廢柴.
一個小廢柴也好.
但是我們是很有光芒.
可以照到這個黑暗.
求主幫助我們每一個.
幫助我們的紅恩堂.
當牧師又再帶領我們.
因為神你給牧師是這裡的一個代表.
當祂說我和我一家都要事奉耶和華的時候.
就是我們全部.
我們這個屬靈的大家庭.
我們就一起的.
要與主同行.
我們多謝主.
多謝主讓我們有在一起的時候.
我們這樣禱告.
奉主耶穌基督的名字.

$^{1001}$阿們.
請坐.
好.
\newpage



\section{詩篇 103:1-5}
\label{sec:PeueLB8OHRY}
\textbf{2022.06.26 《拒絕忘記》 曾裔貴牧師 詩篇 103\_1-5 (和修版)}
\newline
\newline
連結: \href{https://youtube.com/watch?v=PeueLB8OHRY}{\texttt{https://youtube.com/watch?v=PeueLB8OHRY}} ~~~~ 語音日期: 2022-07-01
\newline
\newline
\hyperref[sec:L3GT0EvH0Vg]{\small{< < < PREV SERMON < < <}}
~
\hyperref[sec:index_chronic]{\small{[返順時目]}}
~
\hyperref[sec:index_scriptual]{\small{[返順卷目]}}
~
\hyperref[sec:zIFZXKCsIdk]{\small{> > > NEXT SERMON > > >}}
\newline
\newline
詩篇 103:1-5
\newline
\begin{longtable}{cl}
\hline
\hline
章節 & 經文 (和合本修訂版)\\
\hline
103:1 & \begin{tabularx}{0.7\textwidth}{X} 我的心哪，你要稱頌耶和華！凡在我裡面的，都要稱頌他的聖名！ \end{tabularx} \\ \\ \relax
103:2 & \begin{tabularx}{0.7\textwidth}{X} 我的心哪，你要稱頌耶和華！不可忘記他一切的恩惠！ \end{tabularx} \\ \\ \relax
103:3 & \begin{tabularx}{0.7\textwidth}{X} 他赦免你一切的罪孽，醫治你一切的疾病。 \end{tabularx} \\ \\ \relax
103:4 & \begin{tabularx}{0.7\textwidth}{X} 他救贖你的命脫離地府，以仁愛和憐憫為你的冠冕。 \end{tabularx} \\ \\ \relax
103:5 & \begin{tabularx}{0.7\textwidth}{X} 他用美物使你的生命得以滿足，以致你如鷹返老還童。 \end{tabularx} \\ \\
[1ex]
\hline
\hline
\end{longtable}
$^{1}$弟兄姊妹早安.
代表武唐錦繡堂來祝福大家.
剛才主席Victor姊妹都說了.
其實我們是很需要彼此去分享的.
所以當我們一路跟Gary傳道.
對不起Gary牧師.
來到一路去關乎到雄恩的發展.
我們都很想有不同的支援.
都是一些大家彼此的分享.
所以由下個月開始.
逢第四週.
當我們有牧者下來講道.
我們都會有兩位領唱的姊妹.
趙小玲姊妹和梓玲姊妹.
都會過來幫忙.
跟大家一起彼此的分享.
在他們的事奉裡面.
他們的聖樂來幫助大家.
跟大家一起提升對神的敬拜的學習.
這個都是大家彼此的分享.
所以主來一家.
都很希望雄恩堂的弟兄姊妹.
在未來的日子更加的靈明.
茁壯成長之後.
就可以朝向一路的獨立.
我們開始聚會講道之前.
有一個簡單禱告.
再一次多謝主.
我們打開聖經.
盼望藉著大衛的詩篇.
讓我們去看到.
這個很重要的主題曲.
我們要敬拜我們這位造物主宰.
而且是我們從心靈裡面.
和在整個人.
我們的身體.
這個物質的身體.
我們能夠用盡最好的去敬拜主.
今早我們等候你的訓誨.
我們這樣禱告奉耶穌基督的名.

$^{41}$阿們.
拒絕忘記.
特別是用到很重的字眼.
大衛說不可以忘記祂的恩惠.
但我用的題目拒絕忘記.
和大家一起去思想忘記這個課題.
大衛強迫自己.
不可以對神.
在你生命裡面所做的一切.
來忘記.
他一直說.
他的赦免.
和他用很多美物.
美好的.
無論是人事物.
去讓大衛的人生豐富.
如果這篇詩篇是他造王.
如果是後期的時候.
一點都不為過.
他有很多人生的閱歷.
他的生命很豐富.
有很多的經歷.
但是他選擇對神的恩典不忘記.
剛剛過去的24號.
上個星期三.
因為我都是幼稚園的校董.
我們感受幼稚園.
我們就在劇院.
元朗劇院.
一個很大的畢業禮.
不知道你有沒有參加過幼稚園的畢業禮.
你的孫子,孫女.
或者你的小朋友.
或者你認識的.
你有沒有去過呢?.
其中有一個很重要的環節.
當然固然有一些老師.
他們是師兒.
大的師兒.
你想想.

$^{81}$在整個元朗劇院.
很大型的.
他們有其中一個.
有三個小朋友.
都是畢業班的.
他們做小師兒.
你想想.
小師兒.
他們一定要有那些卡紙.
有那些稿.
但是小朋友四歲.
他怎會懂得寫字呢?.
靠什麼?.
靠機器.
沒錯.
一開始的時候.
輪到這幾個小師兒出來.
穿著很漂亮的衣服.
拿著咪高峰.
準備走出來.
誰知道半分鐘.
完全忘記了.
三個都忘記了.
站在那裡越來越緊張.
我們那些是家長.
還有那些嘉賓.
我們不是很緊張.
而是我們很開心.
很好笑.
還沒說話.
每個人都在想.
不知道說什麼.
後來一開始.
第一個順.
連續就順了.
你很難想像他們的記憶力.
他可以記得那些嘉賓.
是什麼身份.
他有什麼銜頭.
你想想他們背了多久.

$^{121}$你就可以想像到.
他們小朋友的記憶力.
是非常非常強的.
這是上主.
給我們小時候.
有這種很特別的記憶.
將來我們去到晚年.
小朋友的時代.
歷歷在目.
我相信你和我都有共鳴.
小時候我們發生過的事.
很深刻的.
反過來近期的.
我們就會忘記了.
這是大衛.
他告訴我們.
一開始的時候.
其實一到第五節.
已經可以將這篇詩篇.
做了一個很重要的總論.
因為第六節之後.
就開始描述.
他所認識的耶和華.
是一位怎樣的上帝.
一到五節.
他自己個人的經歷.
我的心.
你要去稱頌耶和華.
我用的是舊的和合本.
因為我帶來的聖經.
凡在我裡面.
我整個的內心.
都要去稱頌這位主的聖名.
我的心.
你要稱頌耶和華.
千萬不要忘記.
他一切在你生命裡面的恩惠.
然後就說了.
第四第五節.
他有什麼恩惠在大委身上呢.

$^{161}$赦免一切的罪孽.
大衛犯過什麼罪呢.
一個君王.
為什麼有這樣的膽量.
向他的臣民.
承認自己的罪孽呢.
大家都知道.
還有幾天.
特首林鄭月娥女士.
就要卸任了.
記者就問一些很尖銳的問題.
你有沒有什麼需要向香港人道歉.
有什麼做錯.
我不會為香港人.
我所做的事道歉.
我為我自己對我的家庭.
不夠時間去照顧.
來道歉.
不會為自己所做的事道歉.
大衛這位君王.
他第一樣想到的.
這位的臣.
不只是會給你福氣.
不計較你的道德行為.
這位的臣.
大衛第一個想到的.
就是他內心裡面.
和臣之間.
沒有罪的阻隔.
他能夠得到臣.
那個赦罪的福氣.
還有醫治你一切的疾病.
固然我們明白.
不是說每一個疾病.
都得到醫治.
大衛自己有一個兒子.
都無名的.
沒有改名的.
他拼命為這個兒子祈禱.
最終這個兒子都要被臣擊打而死.

$^{201}$他明白的,他知道的.
他痛的.
但很快臣就給他一個.
第二個兒子.
就是所羅門.
甚至叫拿丹仙芝.
去和這個大衛說.
你的兒子要改一個名字.
叫做耶底底亞.
希伯來文來的.
就是神喜悅這個兒子.
名字叫做所羅門.
就是我們所認識的所羅門王.
這裡告訴我們.
他救贖你的命脫離死亡.
以仁愛和慈悲為你的冠冕.
大衛戴的是最好的皇冠.
因為他的國度.
以色列當時是最強大的.
以後都沒有那麼強大.
後繼人都沒有那麼強大.
頂多是希西家王.
但都沒有那些地圖.
都因為有亞述的入侵.
後來那些君王就更加的沒落.
大衛有冠冕.
但他看見的最寶貴的.
是神給他的.
用仁愛.
用慈悲.
為他所編織的冠冕.
每年都有一些調查.
剛剛2022年.
就調查全世界.
他用了149個國家.
哪個國家是全世界最開心最快樂的.
他有些準則.
就是那個國家的國民.
是否可以自由選擇自己的生活.
那個國家的收入.

$^{241}$是否都穩定和高.
還有那個國家的國民.
如何去理解對內對外的政府貪污.
各方面大概有6個準則.
又不是很主觀地去問.
而是用一個相對的.
首先設立一個虛擬的國家.
將這149個國家放進去.
用相對來看.
最差的.
可能大家都猜到.
全世界最不開心的國家.
就是裡面的國民.
就是阿富汗排首位.
這麼多年都排首位.
然後是非洲那些很窮.
和經常內戰貪污的國家.
前三名全部是歐洲的國家.
芬蘭已經是第二年連續.
瑞典.
然後就是有荷蘭等等.
相對於其他這些這麼窮.
這麼戰亂的國家.
全世界他們認為最開心的.
國民很有自己個人的身份認同.
全世界最開心.
其實你想一想.
最近在美國有兩個很重要的判案.
最高法院的判案.
如果大家有留意新聞就知道.
美國人21歲以上.
是絕對有憲法的保障.
可以帶槍出街.
個人可以有槍攜帶.
剛剛才立了法.
現在18至21歲.
是比較嚴格去審查你.
可不可以擁有槍械.
不是代表你不可以.
比較嚴格.

$^{281}$如果你的分數不夠高.
就是說你有潛在可以射人的危險.
暫時不讓你有槍械.
另外一個最近大家都知道.
最高法院推翻了50年.
一直在美國墮胎.
是沒有任何理由的合法化.
只要你想墮胎.
你就可以墮胎.
因為美國人認為身體是自己的.
我要怎樣就怎樣.
所以地上表面上是全世界最開心.
最快樂的國家.
當然是相對於戰爭.
很貪污.
很多內戰.
很多疾病的國家.
表面上是開心.
表面上是個人有很強的自主的生活空間.
但實際上.
假如我們離開了對造物主的敬拜.
對他的稱頌.
對他的愛.
和他在我們生命裡面.
所來工作的.
尤其是射罪的恩典的工作.
而我們根本沒有去計算在這個類別裡面.
那個所謂的開心.
那個所謂的個人的自由.
其實是放縱.
其實是離開神越來越遠.
其實是很大的一個問題.
有時反而最簡單的.
才是最重要.
大衛很想告誡他的國民.
我們這個民族是有福的.
我們這個民族是因為敬拜這位獨一的真神.
所以他要勉勵他自己.
其實是一篇公開的詩篇.
所以每個以色列人都需要從心裡面學習.

$^{321}$基督教是一個敬拜的群體.
我們來到神的面前.
剛才領唱的弟兄姊妹帶領我們敬拜.
其實就是邀請我們一起.
暫時放下地上的身份.
以神的兒女的身份.
在這個地方和眾聖徒一起敬拜神.
一起將我們在神裡面所得到的恩典.
來去重揚我們這位主.
不要忘記上主恩惠.
剛才那些小朋友.
他背得很熟.
其實他背出來的.
他大概不太知道說什麼.
他們一直這樣念.
老師帶他們背.
背了很久.
所以有一點點的東西卡住.
他們就完全停頓了.
不能夠去深化所讀的.
所說的東西.
基督徒不同了.
隨著我們人生有不同的閱歷.
和人的相處.
人在我們生命裡面所留下的痕跡.
恩典.
嚴規.
不同的事情.
你要記哪些東西呢?.
為什麼我們要拒絕一些東西.
我們不能夠忘記.
但是有些事情.
你真的要選擇.
堅決的選擇.
你要忘記它.
印度有一套很出名的電影.
如果你認識一個叫印度的劉德華.
阿米爾漢.
他拍了很多套電影.
是反映印度民間很嚴重.

$^{361}$深層次的結構問題.
例如教育的問題.
有一套電影叫打死不離三兄弟.
聽過看過.
很好看的.
很正面的講教育.
印度的填鴨式.
另外又有一套是講到他作為父親.
要很強的訓練.
令他的女兒希望成為摔角手.
另一套.
今天我想講的一套電影.
這套電影很沉重.
這套電影叫未知死亡.
中文的翻譯.
是揭露印度販賣兒童.
販賣婦女的一個現實.
當然用電影的橋段去包裝.
讓人很深刻.
印度出現的問題.
帶來的痛苦.
主角當然是阿米爾漢.
一出場.
他是一個很健碩的.
因為他訓練自己.
要很健碩.
全身紋身.
那些紋身全部是仇人的名字.
為什麼呢.
原來這個主角.
本身是一個在印度很成功的生意人.
電信公司的總裁.
為人很低調.
但怎知偶然之間.
有一個女廣告推銷員.
誤認了這個主角阿米爾漢.
是她的男朋友.
由這樣的橋段.
引發出大進的劇情.
就是這個女主角.

$^{401}$有一次坐火車.
無意中揭露了.
黑社會的黑幫.
販賣這些婦女.
這些女孩.
無意中被她發現.
她就輾轉揭發.
釋放解除了這些小朋友的危險.
就遷怒於這些黑幫.
追殺這個女主角.
輾轉.
電影告訴我們.
黑幫將這個女孩.
來到岩山打死.
接著主角被打到失憶.
但很特別.
有15分鐘的記憶.
超過15分鐘就完全忘記了.
所以他要不斷寫紙.
那些卡紙提醒自己.
每一次都要拍張照片.
寫下他做過什麼.
整套電影大部分時間.
就是他不斷記得自己.
很深的記憶.
就是他的女朋友.
被殘暴殺害.
他要追殺.
報復這些黑幫.
整個故事就這樣發展.
所以很沉重.
反映整個社會.
一些結構的問題.
透過電影去包裝.
但有一個很深的教訓.
他自己因為很深的仇恨.
他自己都出不來.
不斷的要報復.
最終被黑幫打到半死.
後來結局當然是結束了.

$^{441}$那些黑幫都被抓了.
但他自己因為已經重傷.
後來因為有一個女實習醫生.
本來是腦科實習的學生.
他對這個案子.
亞米爾漢這個主角.
這個失憶.
就很有興趣.
就問教授可不可以研究這個人.
一路由這個女實習醫生.
就一路追隨著亞米爾漢.
一路看著他去報仇.
最後.
這個女實習醫生.
帶他去到一個地方.
就是一個幼稚園的地方.
當然這是一個故事.
虛構的故事.
亞米爾漢最終.
他整個人變回一個正常人.
沒錯,他的記憶都沒有.
但他正常是一個很開心.
很有愛的地方.
來到他的生活裡.
不再是一個凶神惡煞.
要不斷練到自己很健碩.
要報仇雪恨的一個人.
在他內心.
我想說的就是.
已經沒有了仇恨.
已經沒有了要復仇的內心.
那種痛和那種恨.
換來的是很單純的愛.
在一個環境受保護.
各位.
這是不是我們信耶穌.
在耶穌的愛裡.
我們就是主的羊呢?.
今天我們選擇的是什麼呢?.
記得對你不好的人.

$^{481}$因為你選擇的.
記住神在你生命裡的恩典.
和不同的人.
在你生命的成長裡.
加給你的愛.
這個很重要的.
人越大.
那個記憶力越衰退.
但不開心.
你不喜歡的東西.
就越來越深刻.
對人有時多了挑剔.
少了欣賞.
這是一個通病.
每個人都會的.
不關你的學問.
不關你的什麼.
所以在座的應該跟我差不多一代.
開始童年的東西很深刻.
近期的東西就經常忘記.
不過再要學多一個功課.
一些真是聖經不容許的.
一切的惱恨.
讓鬧.
那些事情.
都要在我們中間被忘記.
不要再記了.
要記的就是.
神的恩賜.
別人在我們生命裡面留下的.
那些愛.
那些祝福.
甚至你要轉化它.
從一個另外的眼光.
去看別人.
我們整個的自己.
整個群體.
我們就是開心的群體.
這個開心.
不是因為我們有衣有食這麼簡單.

$^{521}$是因為人和人之間.
是有一份的信任.
是有一份的相愛.
今年.
剛剛.
諾貝爾和平獎的得主.
是俄羅斯人.
穆拉托夫.
我還沒有展示照片給大家看.
今晚你回去搜一下.
他將他的和平獎的.
那個諾貝爾獎的金牌.
拿去拍賣在紐約.
你猜拍賣了多少錢.
一億三百五十萬美金.
即是大概百億多的港幣.
全數一毛錢都不要自己.
全數捐給俄羅斯.
烏克蘭的兒童基金.
因為已經超過五百萬的兒童.
離開自己的家園.
流離失所.
沒有了自己的家園.
這個諾貝爾和平獎的.
其實是俄羅斯人.
他要帶頭製造這個社會.
要有和平.
弟兄姊妹我們需要求主.
讓我們能夠選擇.
好像大衛那樣.
不可忘記他一切的恩惠.
我們要在神的愛裡面.
去記我們自己所得到的.
這個時間是我們需要.
在神的恩典裡面.
去反思我們自己的生命.
求主祝福大家.
簡短的祈禱.
多謝主.
讓我們能夠透過剛才.

$^{561}$大衛這一篇詩篇.
其中一段的經文.
讓我們去思想人生.
大衛足夠有資格.
他可以在做王的時候.
可以開心.
可以快樂.
可以做他自己喜歡做的事.
的確他真的曾經在神意外.
做了一些對人.
對自己都不好的事.
但是神啊.
你有恩典.
有寬恕.
有拯救的恩惠.
在大衛的身上.
但願我們能夠在主的裡面.
我們能夠不忘記.
你一切的恩惠.
我們在神的恩典裡面.
繼續的學習.
將一些需要忘記的.
特別是一些可能是人和人之間的嫌隙.
或者一些不開心的事情.
我們能夠彼此的寬恕.
互相的在神裡面.
這樣來挽回.
以致神的恩典.
不會受到我們個人內心的狹隘.
來限制.
我們這樣禱告.
奉耶穌基督的名.
阿們.
好,願主祝福大家.
謝謝大家.
\newpage



\section{啟示錄 3:14-22}
\label{sec:zIFZXKCsIdk}
\textbf{2022.07.03 《自以為是的屬靈生命》 老耀雄牧師 啟示錄 3\_14-22 (和修版)}
\newline
\newline
連結: \href{https://youtube.com/watch?v=zIFZXKCsIdk}{\texttt{https://youtube.com/watch?v=zIFZXKCsIdk}} ~~~~ 語音日期: 2022-07-05
\newline
\newline
\hyperref[sec:PeueLB8OHRY]{\small{< < < PREV SERMON < < <}}
~
\hyperref[sec:index_chronic]{\small{[返順時目]}}
~
\hyperref[sec:index_scriptual]{\small{[返順卷目]}}
~
\hyperref[sec:EJ_GoCr8XKk]{\small{> > > NEXT SERMON > > >}}
\newline
\newline
啟示錄 3:14-22
\newline
\begin{longtable}{cl}
\hline
\hline
章節 & 經文 (和合本修訂版)\\
\hline
3:14 & \begin{tabularx}{0.7\textwidth}{X} 「你要寫信給老底嘉教會的使者，說：『那位阿們、誠信真實的見證者、神創造的根源這樣說： \end{tabularx} \\ \\ \relax
3:15 & \begin{tabularx}{0.7\textwidth}{X} 我知道你的行為，你也不冷也不熱；我巴不得你或冷或熱。 \end{tabularx} \\ \\ \relax
3:16 & \begin{tabularx}{0.7\textwidth}{X} 既然你如溫水，也不冷也不熱，我要從我口中把你吐出去。 \end{tabularx} \\ \\ \relax
3:17 & \begin{tabularx}{0.7\textwidth}{X} 你說：我是富足的，已經發了財，一樣都不缺，卻不知道你是困苦、可憐、貧窮、瞎眼、赤身的。 \end{tabularx} \\ \\ \relax
3:18 & \begin{tabularx}{0.7\textwidth}{X} 我勸你向我買從火中鍛鍊出來的金子，使你富足；又買白衣穿上，使你赤身的羞恥不露出來；又買眼藥抹你的眼睛，使你能看見。 \end{tabularx} \\ \\ \relax
3:19 & \begin{tabularx}{0.7\textwidth}{X} 凡我所疼愛的，我就責備管教。所以，你要發熱心，也要悔改。 \end{tabularx} \\ \\ \relax
3:20 & \begin{tabularx}{0.7\textwidth}{X} 看哪，我站在門外叩門，若有聽見我聲音而開門的，我要進到他那裡去，我與他，他與我一起吃飯。 \end{tabularx} \\ \\ \relax
3:21 & \begin{tabularx}{0.7\textwidth}{X} 得勝的，我要賜他在我寶座上與我同坐，就如我得了勝，在我父的寶座上與他同坐一般。 \end{tabularx} \\ \\ \relax
3:22 & \begin{tabularx}{0.7\textwidth}{X} 凡有耳朵的都應當聽聖靈向眾教會所說的話。』」 \end{tabularx} \\ \\
[1ex]
\hline
\hline
\end{longtable}
$^{1}$主席.
各位姐妹平安.
以前香港的其中一個核心價值.
知不知道是甚麼.
香港最大的核心價值就是買樓.
如果買不到一層樓似乎人生都是失敗的.
很多香港人節意足食.
終於買了一層樓.
然後就結婚生子.
為了供這一層樓可能再陷入經濟危機.
很擔心失業.
人要有信仰為目標奮鬥.
這都可以是人生的盼望.
沒有信仰的人.
一般以金錢名譽地位.
作為他們奮鬥的目標和盼望.
一個有信仰的信徒.
他們的人生目標和盼望.
又應該是怎樣的呢.
上個月我講過啟示錄裡的已忽所教會.
今次是講老弟家教會.
大家還記不記得我曾經講過.
如果從內容和教會的情況來看.
這七間教會.
第三間,第四間和第五間教會的情況都很相似.
都是有好有壞的.
第二和第六間的教會就最好了.
都是沒有受到責備的教會.
而第一和第七間的教會.
情況就最危險.
因為審判是最嚴厲的.
第一間教會就是上次講過的已忽所教會.
他們失去了起初的愛心.
如果再不改善的話.
神就會挪走他們.
而第七間教會就是今天要講的老弟家教會.
老弟家教會又面臨甚麼危機呢.
如果老弟家教會不改善的話.
下場又會是怎樣呢?.
今天我們試試從啟示錄三章十四至二十二節.

$^{41}$透過神對老弟家教會的訊息.
去反思我們信仰的三個態度.
我們一起有個祈禱.
各位色在,金在,永在.
坐在天上大寶座的大君王.
我們聚集在你的殿裡敬拜你.
因為你配得我們對你的讚美.
當我們能夠享受你賜給我們的恩典.
又滿有盼望等待著你再來臨的時候.
求你讓我們在等待得贖的過程裡.
我們能夠有信心,有忍耐.
並且能夠遵行你向我們所啟示的真理和教訓.
好叫我們所信的,所行的,所講的.
都和你的恩典相呈,榮耀你的名.
我們在你面前仰望禱告.
我們奉主耶穌基督得勝名字祈求.
阿們.
信仰的第一個態度.
不要以為自己現時的屬靈狀態很好.
就不需要做信仰反省.
經文在第十七節說.
「你說我是富足的,已經發了財,一樣都不缺」.
老弟家教會認為自己一樣都不缺.
並且是富足的.
如果這真是事實的話.
應該值得為他們感恩.
事實是不是呢?原來不是.
因為經文接著說.
「卻不知道你是困苦,可憐,貧窮,赫硬,赤身的嗎?」.
這種困苦,可憐的狀態和富足相比.
落差實在太大了.
原來只是老弟家教會一種自以為是.
自我感覺良好的狀態.
而且不單單有所缺乏.
還有其他方面有虧欠.
在第十五和第十六節.
「我知道你的行為,你又不冷又不熱.
我巴不得你或冷或熱.
既然你如溫水,又不冷又不熱.
我就要從你口中把你吐出來」.

$^{81}$老弟家的行為被形容為「不冷也不熱」.
究竟是指什麼呢?.
什麼叫「不冷也不熱」呢?.
原來「不冷也不熱」是一個比喻.
比喻老弟家的行為是一無是處.
可有可無.
沒有了教會應該有的特質.
教會不成教會.
原來距離老弟家六里之外.
有一個地方叫希拉波納.
那裡盛產溫泉.
溫泉的溫度高達華氏95度.
而溫泉含有碳酸鈣.
泉水在高溫之下就有醫療的作用.
所以主說「我巴不得你熱」.
因為在這麼熱的溫度下.
才能發揮到醫療價值.
而當泉水流到哥羅西.
溫度開始下降.
雖然沒有了醫療價值.
但泉水流動的過程.
就將水中的礦物質全部沉澱.
泉水變得冷而清澈.
可以給人飲用.
所以主說「我巴不得你冷」.
因為冰凍的涼水已經沒有了雜質.
可以供人飲用.
但魯迪加的地理卻是.
在希拉玻納和哥羅西中間.
流向魯迪加的泉水.
既沒有了高溫的醫療價值.
亦未完全沉澱了泉水中的雜質.
可以供人飲用.
所以主對魯迪加的責備就是.
「你的行為不冷也不熱」.
就是說他的行為好像泉水一樣.
對一個口渴的人來說.
泉水充滿雜質.
喝了也會吐出來.
根本是一無是處.

$^{121}$全無價值可言.
當魯迪加教會.
很自滿地認為自己已經很富足.
發財了,一樣都不缺.
他們陶醉在一種自我感覺良好的狀態下.
殊不知這種感覺卻是建築在物質條件上的滿足和自信.
而引用出來的屬靈的驕傲.
以至導致屬靈上的失明.
因此主就點出他們的危險.
其實魯迪加教會是一個困苦.
可憐,貧窮,黑眼,赤身.
而且一無是處.
主用了五個形容詞來形容魯迪加教會的真實情況.
最可惜的是.
他們竟然察覺不到自己的不足之處.
用今天的術語來形容魯迪加教會.
就是他們是一個屬靈的核子.
而且窮得只剩下錢.
窮得只剩下錢.
就是魯迪加教會的真實情況.
所以魯迪加教會所犯的錯.
不是對人很冷漠.
也不是屬靈熱度不足.
從他們自誇一樣都不缺的心態.
我們深信這間教會在外表來看是不錯的.
可能好像其他六間教會.
曾經被人稱讚過.
是一個宗教活動頻密的魯迪加城.
他們懷有這種自大自滿的態度.
不可能在宗教上是一個完全沒有作為的群體.
因此他們的問題並不是沒有宗教的活動.
正如魯迪加城的問題不在於沒有水.
他們有水.
但水含有雜質不能飲用.
魯迪加教會也有他們的宗教活動.
但因為有雜質.
以至這些活動全沒有屬靈的效力.
這才是魯迪加教會真正的問題.
他們有宗教活動.
但沒有生命.

$^{161}$他們有宗教活動.
不活潑的.
就像一潭死水一樣.
沒有質素.
沒有影響力的屬靈生命.
在主的眼裡.
就是全無價值.
一無是處.
如果我們換上信仰的詞彙去形容.
什麼叫不冷又不熱?.
這句話就是既沒有好的見證去影響身邊的人.
又沒有好的靈性去滿足內心對永恆的渴求.
對人對自己.
對整個社區.
一點幫助都沒有.
在這裡和大家說一個諷刺式的笑話.
話說在俄國革命之前.
俄國革命很久了.
坊間有一個流行的笑話.
話說有一天.
就像今天一樣崇拜日.
天使走過教會門口.
她遇到一個魔鬼.
天使很緊張地問魔鬼.
你來這裡做什麼?.
這裡每個人都是來敬拜上主的.
沒有人會是你的信眾.
魔鬼說:是嗎?.
未必的.
她只是來教會.
未必真的用心敬拜神的.
天使就不相信.
於是魔鬼就說:不如這樣吧.
我們一起坐在這裡.
一起數一下.
如果有人真的進了教會.
是誠心祈禱敬拜的.
天使就將他的名字寫在你的生命冊上.
但是如果進了去.
不是真心祈禱敬拜的.

$^{201}$那個人就歸我了.
天使和魔鬼就一起.
在教會門口一起數.
既然這是諷刺式的笑話.
最後屬於魔鬼的人數.
大幅拋離天使.
這個笑話想說什麼呢?.
想諷刺什麼呢?.
是諷刺當時俄國的教會和信徒.
他們不單沒有信仰的生命.
也對其他人完全沒有影響力.
他們雖然身處在教會裡.
但是所表現出來的信仰.
完全發揮不了信仰的核心價值.
這個笑話諷刺他們.
就算身處在教會裡.
你也表現不了一個信仰的精神.
更何況你離開了教會之後.
更加沒有人知道你是一個信徒.
明白了嗎?.
如果今天天使和魔鬼在雄恩堂數點人數.
不知道結果會是怎樣.
各位雄恩家庭的姐妹.
今天我們的信仰生命.
是不是覺得很舒適呢?.
是不是覺得很自由自在呢?.
是不是覺得自我感覺良好呢?.
是不是覺得我們現在很好呢?.
一樣都不缺.
老弟家教會的情況也一樣.
我們是不是也需要反思一下.
我們自己所信的究竟是什麼呢?.
什麼才是我們的信仰核心價值呢?.
是我們的祈禱靈不靈?.
是我們的祈禱有求必應?.
信仰的第二個態度就是.
承認自己的罪悔改並且接受主的管教.
經文在第十八節說.
我勸你向我買從火中鍛鍊出來的金紙.
洗你腹足又買白衣纖上.

$^{241}$洗你赤身的羞恥不露出來.
又買眼藥抹你的眼睛.
洗你能看見.
老弟家教會當時最自持.
最驕傲最不可一世的.
就是他們是一個財經中心.
成衣製品和藥物的中心.
但想不到他們最需要改變的.
就是正正是這一點.
所以主給的勸勉就是.
向我買火鍊的金紙.
白衣和眼藥.
中國古籍淮南子一書.
在漢朝被認為是收集了朱子百家的其中一部著作.
裡面有一段說話.
他說自誇善射者死於此.
善戰者死於兵.
善泯者溺於水.
這個正正是我們認為最強的地方.
就是我們最輕視和忽略的地方.
我們的屬靈驕傲就是這樣來的.
我們最認為.
那樣東西是最厲害的.
就成為我們的屬靈驕傲.
向我買從火中鍛鍊出來的金紙.
洗你腹足.
主提醒老弟家教會.
我們的屬靈生命如果沒有經過一些試煉.
只是一個溫室裡面的嬰兒.
一旦暴風雨來到.
就會全然軟弱跌倒.
毫無抵抗的能力.
唯有經過火一般的試煉.
生命才會堅強和茁壯.
如果信徒的信仰生命經過試煉.
一定能夠為主而結出屬靈的果子.
並且在生活裡面能夠造就別人.
和建立別人生命的質素.
這些腹足.
豈不是比地上的金銀財寶.

$^{281}$名譽地位更有永恆價值嗎?.
我們剛才聽到玉芬的分享.
苦不苦?.
苦啊.
她女兒的病仍然在病.
不是因為做了信徒之後.
她女兒就能夠病得痊癒.
不需要她照顧.
她仍然要照顧她女兒.
但她學會如何與神同行.
香港的離婚率很高.
信徒也一樣.
離婚對一個婦人來說.
或者一個女士來說.
苦不苦?.
苦啊.
向了多少年才走出死暗的幽谷?.
七年.
今天她能夠站在這裡講見證.
是信仰承托了她的生命.
我們信徒的生命很喜歡一帆風順.
我們不喜歡有逆境.
但逆境就像一個火煉的金一樣.
燒盡我們.
讓我們能夠建立一種肅靈的果子.
經文說.
「有買白衣穿上.
使你的赤身羞恥不露出來」.
主要提醒老弟家教會.
我們所犯的過錯.
最終在神面前必定會顯露出來.
當一個充滿罪惡的人.
站在聖潔的神面前的時候.
一定會覺得很羞恥.
感覺到無地自容.
如果我們看回《創世記》的書.
人類的始祖亞當和夏娃.
一吃了分辨算角樹的果子.
他們就看見自己赤身露體.
他們怎樣做呢?.

$^{321}$他們就編織了無花果樹的葉子.
為自己造成裙子.
當神呼喚他們的時候.
他們怎樣回應呢?.
他們說「我在圓中聽見你的聲音.
我就害怕了.
因為我赤身露體.
我就藏起來了」.
他們感到羞恥.
感到害怕.
還有藏了起來.
讓我們知道.
逃避是不能解決他們所犯的罪.
今天我們也一樣.
我很喜歡逃避.
犯了罪,犯了錯.
不承認,逃避.
算了.
很快就不記得了.
算了.
沒有大不了.
一件白色純潔無瑕的衣袍.
是代表神的寬恕和憐憫.
當我們能夠自覺犯罪.
得罪人,得罪神的時候.
我們記得在面前.
去求這件白衣.
求神以他的寬恕和憐憫.
去遮蓋和赦免我們的罪.
這是唯一的方法.
在整個舊約以色列的歷史.
我們也看到.
神對以色列百姓的寬容和赦免.
主同樣要求老弟家教會.
向神求取寬恕.
去遮蓋他們所犯的罪.
教會在這個社區的.
其中一個目的是什麼呢?.
是要建立一個聖潔的群體.
進來後.

$^{361}$進來後每一個信徒.
他要自覺人生污穢.
覺得自己犯罪,得罪神.
來到教會的時候.
求主赦免,求主寬恕.
以至他能夠守潔心清.
可以安然在這個聖殿裡敬拜.
我再說一次.
如果心裡面深存污穢.
你的敬拜是沒有意義的.
神是不會喜悅的.
你唱得多大聲也好.
神也不會喜悅的.
唯有你.
在神面前認罪悔改.
你存一個清潔的良心.
你在神面前敬拜.
神就喜悅了.
聖人說:有買眼藥.
抹你的眼睛.
使你能看見.
能夠使我們打開淑寧的眼睛.
唯一的方法.
就是與神有親密的相交.
當我們與神的關係進入一步.
一種叫什麼呢?.
你在我裡面.
我在你裡面的狀態.
這就是約翰福音所說的.
葡萄樹與枝子的關係.
我們要連於葡萄樹.
我們就能夠從神裡面.
能夠有質量,健康,強壯地成長.
也唯有我們將生命重整.
淑寧生命有成長.
我們才能夠提升我們的淑寧視野.
使我們以前所看不到的.
今天我們能夠看見.
一個很明顯的分別.
你未信耶穌的時候.

$^{401}$你看見所有人.
那些都是人.
你信了耶穌之後.
你再看見那些.
所有都是主所創造的人.
那些都是主所愛的人.
這就是淑寧視野的分別.
你信耶穌之後.
你看不見.
你看不見那些是神所愛的人.
你信了耶穌之後.
有神的眼光.
你看見每一個.
你知道每一個都是神所創造.
每一個都是神所愛.
然後你就會有連吻.
這個也很重要.
我們有沒有神的眼光.
去看這個世界.
雖然老弟教教會的情況是這麼惡劣.
但是主的管教卻又是主動的.
在十九至二十節這樣說.
他說凡我所愛的.
我就責備管教.
所以你要發熱心.
也要悔改.
我站在門外叩門.
若有聽見我聲音而開門的.
我要進到他那裡去.
我與他.
他與我一起吃飯.
老弟教教會縱然犯了很多錯.
淑寧情況很差.
但是他們仍然都是主所愛的.
仍然都是主所眷顧的一間教會.
主在門外叩門.
只要聽到主的聲音.
呼喚,邀請.
願意接受.
願意接受被主管教.

$^{441}$打開心扉的主.
就會與我們一起.
今天我們願不願意打開我們的心.
讓主進入我們那個污穢的靈裡面.
其實我們心裡面真的很污穢.
以前我曾經聽過一個這樣的比喻.
就是我們想起耶穌.
心房裡面有一百間房.
主耶穌.
邀請主耶穌進入其中一間房.
這間房.
主耶穌你可以在這間房裡.
怎樣都可以.
但是那九十九間房.
你就不要理會.
那九十九間房.
我才可以進去.
主你不要進來.
黑暗是會逃避光的.
我們也是一樣.
當我們心裡面污穢的時候.
我們都很想逃避主.
進入我們生命裡面.
只要肯回轉.
只要肯認罪悔改.
恢復回屬靈生命的活躍.
老弟家教會仍然能夠與主.
建立那種親密的關係.
以至最後被主所縱容.
神對人最大的懲罰.
不是要用最嚴厲的方法去管教他.
而是要用最嚴厲的方法.
而是完全放手不理.
讓你任意而為.
讓你任意去放縱你自己.
這個才是最大的懲罰.
主肯管教你.
仍然是愛你.
主不理你.
任由你怎樣做.

$^{481}$這是一個最大的懲罰.
當人完全滿足在自己的心意的同時.
正正就是與神的關係.
最疏遠.
以及在距離上.
最遠離神的時候.
在春秋的時候.
孔子的學生曾三.
曾經很勤奮好學.
很深得他孔子的喜愛.
有些同學就問他.
為什麼你進步得這麼快.
他就說.
吾日三醒吾生.
為人謀而不忠乎.
與朋友交而不信乎.
存不擇乎.
意思就是.
我每日都會多次問自己.
能夠替別人辦事.
是否盡力了.
與朋友交往.
有沒有不誠實的地方.
所教的學生.
有沒有教不好.
如果發現教得不好.
就立即改正.
三醒吾生.
就是能夠隨時都要作出.
一個自我反省.
看看有什麼可以改善和修正.
信仰的三醒吾生.
又是怎樣的呢.
我們可以每一晚都問問自己.
神所交託的事.
我做好了沒有.
為別人的守望和待討.
完成了沒有.
聖經的教訓.
和對神的認識.

$^{521}$有沒有認識多一些呢.
我相信如果我們每一天.
都做一次三醒吾生.
信仰反思.
我們一定可以進步神速.
洪英嘉弟兄姊妹.
更加有信心.
更生自己的屬靈生命.
悔改的心.
和神親密的關係.
就是主對老弟家教會的勸勉.
我深信.
也是對洪英嘉.
眾信徒的一個勸勉.
所以只要我們能夠反省到.
自己的不足.
願意悔改.
重新起步.
和主建立一個親密的關係.
我們的生命.
仍然可以被主所使用.
榮臣益人.
這不是我們信徒.
很想要做的事嗎.
榮臣益人.
信仰的第三個態度就是.
生命的終極.
需要有得勝的目標和盼望.
經文在21節這樣說.
得勝的我要賜他.
在我的寶座上與我同坐.
就如我得了勝.
在我父的寶座上.
與他同坐一般.
主要是要他.
主透過約翰.
寫給啟示錄的七間教會.
都透過不同的應許方式.
表達了他對教會的期望.
對這所教會.

$^{561}$就是重拾起初的信心.
也可以是繼續忍受患難.
又可以是不再向這個世界妥協.
不過主最後給了老弟家這封信.
主耶穌是以自己作為一個榜樣.
作為他對老弟家教會的期望.
和他對得勝的定義.
主耶穌對天父至死不渝的信服和信心.
使得他從死裡復活.
與神與父神同坐寶座.
老弟家教會的信徒如果要得勝.
就要好像主耶穌信服天父一樣.
他們同樣要信服主耶穌的命令.
和對主耶穌有充足的信心.
對於猶太人來說.
他們一生的終極盼望.
就是能夠與主一同坐寂.
彌賽亞的筵席是他們的盼望.
現在主對老弟家的應許.
不單可以與主在宴會上.
宴席上一起坐.
不單單是這樣.
更加是與主一起坐在寶座上.
這一定比一起坐在宴席上更加尊貴.
與主坐在寶座上.
這是一個很尊貴的應許.
你想一想將來有一天.
你在天上.
你坐在主給你的寶座上.
坐在寶座上.
你旁邊的就是主.
你會怎樣?.
你的心情會怎樣?.
這會不會是你的盼望?.
今天大家坐在聖殿裡敬拜主.
你知道主與你同在.
將來你坐在寶座上.
主耶穌就坐在你旁邊.
有沒有想過自己的信仰目標和盼望?.
能夠與得勝的君王一起坐在寶座上.

$^{601}$會不會就是一個信徒終極的盼望?.
如果我們的目標.
真的很想過一個得勝的人生.
我們就必定要審視一下我們的方向.
究竟我們的方向能否走到這一步.
我們的方向是否可以去到與主一同坐在寶座上.
詩篇九十篇.
剛才玉芬說很喜歡讀詩篇.
今天又有一個詩篇九十篇.
第十節這樣去形容我們現在的人生.
他說我們一生的年日是七十歲.
如果是強壯的.
可以讀八十歲.
但其中可以驚誇的.
不過是.
勞苦受煩.
轉眼即逝.
我們就如風飛去.
有什麼好厲害呢?.
讓你讀到八十歲,九十歲又如何呢?.
勞苦受煩.
詩人說到人生苦短.
勞苦受煩.
轉眼即逝.
如果人世間的現象是這樣的.
我們又怎可以過一個得勝的生命?.
這麼苦的人生.
怎樣過一個得勝的生命?.
有人類歷史以來.
都不知道有多少聰明的人.
或者哲學家.
都同樣問一個問題.
人生的意義和目標究竟是什麼?.
我們為什麼而活?.
為什麼而生存?.
為那一口飯而生存?.
為那一層樓而拼搏?.
今時今日.
有很多不同的理論.
亦有很多不同的宗教呈現在我們面前.

$^{641}$任我們去選擇.
不過基督教所提出的答案就是.
不是人去找神.
不是.
你找不到.
而是神主動去找我們.
神透過啟示向我們.
講明祂的心意.
聖經和主耶穌的道成肉身.
就是對人類最直接的啟示.
你認為.
你在報道會裡面.
你舉手決志.
是你去選擇神嗎?.
不是.
沒有聖靈在你的生命的工作.
你永遠都找不到神.
我們可以透過聖經裡面的神人關係.
可以說是去認識神的屬性.
亦可以透過耶穌在世的經歷和見證.
讓我們的屬靈眼睛能夠見到.
見到甚麼?.
見到原來生命可以有另一種向度.
有另一種向度.
我們的生命不只是要買樓.
不是要名譽.
不是要財富.
可以是一種向度.
但基督教的價值觀告訴你.
可以有另一種選擇.
而這種生命的向度.
不單單可以脫離.
施人所講的人生苦短.
勞苦受煩.
而是可以帶領我們進入永恆.
在世上展現一種甚麼叫得勝的生命.
這種正是神為我們每一個.
相信主耶穌的信徒而設計的生命.
主耶穌為我們而設計的生命.
得勝的生命.

$^{681}$讓我們去尋找.
上主為我們設計的人生.
恢復上主創造我們的心意.
我們被主喚醒了的自由意志.
去選擇一條主所喜悅的生命歷程.
你今天是有這種選擇的.
你回到伊甸園裡面.
神給你一個自由意志.
然後你選擇.
你遵從上主的話而活.
還是你跟隨那條蛇的引誘.
抗拒主耶穌的話而活.
生命裡面能夠彰顯神的屬性.
讓每一個信徒.
能夠在生命裡面見到主耶穌.
能夠為主而活.
用生命見到主耶穌.
用生命見到主.
好像玉芬所說.
在困境裡面與主相遇.
在困境裡面.
倚靠主而走過這個困境.
用生命去見到主.
然後昂首闊步地走入永恆.
這就是一種得勝的生活.
不是被苦難勝過你.
不是被情緒勝過你.
不是被逆境勝過你.
很多人自怨自艾.
好像他所說的自憐.
為什麼是我.
自怨.
為什麼我這麼苦.
不是,不是,完全不是.
一種得勝的生命.
當他再說出他的見證的時候.
能夠幫助別人.
我走過這條困苦的路.
我靠著主得勝.
你也可以.

$^{721}$神給我們有生命,有恩賜.
難道只是.
就這樣過一個為自己而活的人生.
為買樓,為升職.
為累積財富的人生.
還是我們可以用神賜給我們恩賜.
去過一個為神而活的人生.
在二千多年前.
主耶穌來到這個世界.
為我們親身示範過.
怎樣走這條為神而活的人生.
主耶穌得勝的道路.
是一條十架的道路.
祂說過.
學生不能夠高過老師.
僕人也高不過主人.
今天如果我們說要勝過這個世界.
我們說要過一個得勝的生活.
我們也需要像主耶穌一樣.
走這條十架的道路.
就是一條信服天父心意的人生.
我接著就問.
你知不知道天父的心意?.
如果你不知道天父的心意.
你怎樣信服他的心意?.
當我們完成了這個信仰人生之後.
我們所得到的.
不單單是得到永生這麼簡單.
很多時候我們向神說.
信耶穌吧 信耶穌得永生.
不單單是這樣.
我們以一個勝利者的身份.
和主一起坐在神的寶座上.
你想過沒有?.
一個受造物.
能夠和一個創造者.
一同坐直.
是最高的榮耀.
我們是最高的榮耀.
昔日我們的始祖阿當.

$^{761}$他為了要與神同等.
他甘願犯罪得罪神.
違背神的命令.
去吃那個分辨善惡樹的果.
就是為了要與神同等.
今天我們不再需要.
好像走回阿當的舊路.
也不需要走自己的路.
只需要按著天父的心意.
走主所走過的路.
活出信仰就可以了.
不知道大家有沒有聽過.
一個溫水煮蛙的故事.
話說在美國.
康納爾大學的科學家.
做過一個溫水煮蛙的實驗.
把一隻青蛙放在一個熱水裡面.
青蛙碰到熱水.
就立刻彈走了.
但是如果將青蛙放在熱水裡面.
青蛙放在一個冰冷的水裡面.
青蛙就好像如常地暢游.
它不會彈走.
然後慢慢將鍋裡面的水加溫.
加到水很熱.
慢慢熱 慢慢熱 慢慢熱.
青蛙習慣了.
它就不可以跳出這個鍋.
最後就死在熱水裡面.
實驗想證明一件事.
青蛙由於對漸變的溫度.
產生適應性和習慣性.
適應性和習慣性.
漸漸失去了警惕和反抗的能力.
為什麼說這個例子.
我們對罪.
會不會適應了它.
會不會習慣了它.
習慣了.
一件污 兩件威.

$^{801}$有什麼所謂.
許美蔭犯罪.
每個人都是罪人.
有什麼問題.
我們每個人都是罪人.
犯多一件 少一件.
有什麼所謂.
適應性和習慣性.
這個就是對我們的溫水著畫.
當我們的生命沒有警醒的心.
沒有一種對敏銳世界對我們的影響.
我講過的.
你每一個星期回來教會.
逗留多久.
一個半小時.
離開教會面對這個世界多久.
是六天.
三百六天.
即是七天減少了一個半小時.
你想想.
世界對你的影響多.
還是教會對你的影響多.
正正因為這是神的話語.
才可以讓你們逆水行舟.
如果你不能領受到神的話語.
神的話語不能在心裡.
成為那種腳前的燈.
路上的光.
你如何抗衡七天減少了一個半小時的影響.
沒有什麼可能.
世界不斷影響我們的時候.
我們是否察覺到.
還是我們已經習慣了.
習慣了.
習慣了我們現在的舒適感.
作為一個沒有信仰的人.
他們的盼望和目標.
可以是很短暫的.
如飛而去.
但作為一個神的兒女.

$^{841}$天國的子民.
君尊的祭司.
我們有沒有一個永恆的盼望呢.
一個得勝生命的目標.
如果信仰.
三省吾生是有盡心板的話.
我們要問.
我們的生命有沒有被主所使用.
我們的生命有沒有對其他人.
有美好的影響.
我們的生命和神的心意.
是否配合呢.
昔日老弟家教會.
主身在一個安定繁榮的社會之中.
信徒所倚靠的是自己的能力.
而且自滿而不自知.
當主提醒他們說.
指出他們一無是處的時候.
唯一的回應.
就是認罪悔改.
再一次讓神去管教.
認定主就是超乎一切以上的元首.
主耶穌是洪恩家的主.
也是洪恩家的王.
當我們能夠在世實踐一個得勝的生命.
將來在永恆的角度裡.
就能夠與主一同坐在寶座上.
享受一個永恆的尊榮.
我們一起去祈禱.
透過你對老弟家教會的教導.
你叫我們不要自滿.
叫我們要有認罪悔改的心.
還要有一個得勝生命的盼望.
求主兼顧我們每一個會有的心.
讓你的話語栽種在我們心裡.
好讓我們能夠活出信仰.
榮耀主名.
我們去祈禱.
奉教主耶穌得勝明治祈求.
阿們.

$^{881}$願主賜福大家.
\newpage



\section{哥林多前書 13:13}
\label{sec:EJ_GoCr8XKk}
\textbf{2022.07.10 《愛的呼喚》 哥林多前書 13\_13 (和修版) 陳一華牧師}
\newline
\newline
連結: \href{https://youtube.com/watch?v=EJ_GoCr8XKk}{\texttt{https://youtube.com/watch?v=EJ\_GoCr8XKk}} ~~~~ 語音日期: 2022-07-14
\newline
\newline
\hyperref[sec:zIFZXKCsIdk]{\small{< < < PREV SERMON < < <}}
~
\hyperref[sec:index_chronic]{\small{[返順時目]}}
~
\hyperref[sec:index_scriptual]{\small{[返順卷目]}}
~
\hyperref[sec:2v_ng9IARWc]{\small{> > > NEXT SERMON > > >}}
\newline
\newline
哥林多前書 13:13
\newline
\begin{longtable}{cl}
\hline
\hline
章節 & 經文 (和合本修訂版)\\
\hline
13:13 & \begin{tabularx}{0.7\textwidth}{X} 如今常存的有信，有望，有愛這三樣，其中最大的是愛。 \end{tabularx} \\ \\
[1ex]
\hline
\hline
\end{longtable}
$^{1}$請教授和學生們進行聖祭禮儀.
我們可以敬拜天父.
感不感恩?.
感恩.
雖然最近天氣熱.
可能你正在走路.
大巷踏細巷.
回來教會.
但我相信你的心.
都很渴想親近天父.
我今早早一點來到.
看到有不少弟兄姊妹.
很早回來.
預備心來敬拜.
這實在是很好的事情.
回來教會.
你又可以唱詩.
又可以讀聖經.
又可以聽到神的話語.
亦可以聽到師班獻唱.
更加難得.
還有聖餐來紀念主.
所以一切都是恩典.
今日牧師和大家讀聖經.
和思考經文.
很簡單.
一節聖經.
你們都很熟悉.
不如我們一起讀一次.
預備.
一二三.
如今常傳的.
有信有愛.
有愛這三樣.
其中最大的是愛.
這本經文雖然熟悉.
但亦很挑戰性.
因為不是很容易實踐.
所以我們今天.
來思考這節經文.

$^{41}$近來這幾個月.
國際新聞最熱門的是甚麼呢?.
我相信大家都知道.
就是俄烏戰爭.
這場戰爭.
可以說是沒有贏家.
這場戰爭亦令到.
很多人流離失所.
逃離家園.
甚至死傷不少人.
在這情況下.
我們很多時都會問.
為何要去打仗呢?.
實際上大家都知道.
如果我們用武力去征服人.
其實這不會是長久的.
不過是曇花一現.
這個國家不會千秋萬代.
反過來可能是遺臭萬年都未定.
所以大家都明白到.
用武力去征服人.
不是一項強項的事情.
相反如果我們用愛心去征服人.
這才能夠令人長久的心悅誠服.
大家認同嗎?.
以耶穌基督為例.
耶穌基督的大愛.
耶穌基督的愛心.
能夠令到這二千多年來.
很多人都得到感染,激勵.
所以樂意效法耶穌的榜樣.
來關愛其他人.
所以唯有主耶穌的愛.
祂的表現.
才能夠令到我們人的心裡.
更加願意去學習跟隨祂的榜樣.
所以你發現到.
原來愛的威力,動力是非常之大的.
為什麼呢?.
因為愛包含了兩個很重要的元素.

$^{81}$所以今早牧師跟大家分享.
愛怎麼包含兩個很重要的元素呢?.
第一個元素.
愛是付諸行動的.
大家看到這張照片.
有沒有印象在當中呢?.
其實這張照片不是拍了很久.
如果你有印象.
是在上個月的父親節.
不謀而合幾份報紙都登了這段新聞.
可能因為父親節.
你看到這張照片的時候.
你有什麼感想呢?.
一個老人家揹著一個年輕人.
是一個七十多歲揹著一個四十六歲的兒子.
是爸爸揹著兒子.
而爸爸只有五十四公斤.
兒子只有七十多公斤.
做什麼呢?.
就是在一條天橋樓梯的地方.
上上落落.
原因是兒子的腳癱瘓了.
爸爸很心急.
也很有愛子之心.
就揹著兒子上下樓梯.
希望兒子試試自己的腳.
扶著樓梯的時候.
可不可以把腳動一動呢?.
他就帶著這樣的目的.
每一天都跟兒子上下樓梯.
各位頂層的朋友們.
當我看到這張照片的時候.
我心裡很感動.
為什麼呢?.
爸爸愛兒子.
可以說是付諸行動.
是有實質的行動.
就算他自己受虐也好.
為了兒子的康復.
他都願意揹著兒子.

$^{121}$來跟他做一個運動.
希望幫助兒子.
原來是真的.
愛是什麼呢?.
是一個付諸行動.
就好像主耶穌愛我們世界上的人.
他就付諸行動.
他放棄了自己天上的榮耀的寶座.
放下了他自己的尊貴.
來到地上.
道成肉身.
在馬槽裡出生.
而且還生活在我們一個.
充滿罪惡的人世間.
試問一下.
有哪一位是願意.
將自己的尊貴身份放下.
而願意來到人世間.
這個行動是由最高跌到最低.
是表示什麼呢?.
表示耶穌因為愛.
他就付諸行動.
所以愛能夠激勵到很多人.
平時不想做的事.
或者平時做不到的事.
因為你有愛心在背後推動.
你就能夠做得成.
大家相信嗎?.
曾經有位媽媽.
因為幫兒子.
兒子身體有肝病.
媽媽想幫兒子.
但是因為自己身體不太好.
所以她要操練自己的身體.
每天都走十公里.
練好身體.
希望將來能夠幫助兒子.
各位頂智母們.
很多時候你覺得很難做.
但是因為背後有愛.

$^{161}$你就願意付諸行動.
第二個元素是什麼呢?.
第二個元素就是作出犧牲.
頂智母們.
爸爸為兒子犧牲大嗎?.
很大.
他失去了自己私人的時間.
也放棄了自己工作的時間.
也放棄了自己社交的時間.
他將自己所有的時間都給兒子.
為了他的康復.
報紙說這個老人家.
做這個動作做了多久?.
做了二十年.
厲不厲害?.
二十年如一日.
來背著兒子上天橋.
下樓梯.
希望操練到兒子能夠走路.
頂智母們.
這個爸爸為兒子.
很樂意很甘心犧牲他所有的東西.
目的是希望兒子能夠.
有一個康復的機會.
同樣主耶穌也是.
因著愛你和愛我的緣故.
耶穌就走上了十字架的道路.
祂被釘身在十字架上.
流出祂的寶血.
洗淨我們的國泛.
正如今天我們領聖餐.
老牧師提醒我們.
為什麼我們透過聖餐.
會更加思想到和主的關係呢?.
為什麼我們在聖餐裡.
更加容易反思到自己的靈命.
好像有時停滯不前呢?.
是因為每一次你和我都感受到.
來到聖餐桌前的時候.
我們會被主耶穌犧牲的偉大的愛.

$^{201}$來激勵自己.
因為耶穌基督的犧牲.
祂的無比的愛.
激勵我們的緣故.
所以我們在主耶穌的面前.
我們更加願意追求成長.
更加追求有豐盛的人生.
弟兄姊妹們.
這兩個元素是非常重要的.
所以牧師今天早上很想鼓勵大家.
每一位.
如果說我也很想有愛心.
來關懷別人.
關愛別人的時候.
你就要時時都問自己.
到底我有沒有擁有這兩個元素呢?.
我對別人的關愛.
有沒有實際的行動呢?.
有沒有付諸行動呢?.
我對別人的關愛.
對別人的關心.
有沒有願意作出犧牲呢?.
這兩個元素是很重要的.
大家都知道近日.
香港社會的氣氛都很喜氣洋洋.
為什麼呢?.
因為我們慶祝回歸25週年.
上星期日.
如果你還記得.
新政府成立的典禮裡面.
國家主席和我們的新特首.
都是不約而同.
在他們的演詞當中.
是強調關愛的重要.
有沒有留意到?.
他們很想政府成為一個關愛的政府.
這個目標方向是很正確的.
是很重要的.
因為我相信你和我都很希望.
新一屆的政府.

$^{241}$新的特首.
新的官員.
他們能夠要做的事.
不單止是要解決香港新層次的問題.
不單止是關注到房屋的問題.
或者是民生的問題.
或者青少年向上流的問題.
這些都是需要關心的.
又或者我們不單止希望新的政府.
這一屆的政府.
不單止著眼於如何做好KPI.
或者做好結果為目標.
我們不單止期盼.
它可以達成這些方向.
在這些行動的背後.
我猜你和我是最想.
希望這些行動的背後.
我們的政府.
能夠擁有一種關愛的心態.
大家認同嗎?.
當我們的政府能夠擁有關愛的心態.
作為它的背後動力的時候.
它就會能夠想市民所想.
就會急市民所急.
就會成為一個父母官.
能夠愛護自己的子民.
能夠在他們的執政期間.
推出更多的為民的德政和人政.
為什麼?.
因為愛心的動力是很重要的.
不是單單達到目標.
而是背後需要有愛心的動力.
能夠推駛到我們每一個人.
不單止是政府.
各位電視節目們.
今早你和我都是.
原來上帝讓你成為祂的兒女.
都是希望你在生活當中.
和其他人有些分別.
你有沒有記得耶穌在馬太福音說什麼?.

$^{281}$耶穌在馬太福音說了一句話.
很重要的.
祂說:你們是世上的什麼?.
光,是吧?.
這個光是怎樣的?.
「成造在山上,是不能」.
「隱藏的」.
我們的光放在檯面上.
也能夠照亮一家的人.
這些經文大家有印象嗎?.
怎樣能夠令基督徒的生命.
令信耶穌的人.
他的見證是很特別的呢?.
其中一個最大的秘訣.
就是我們的生活.
我們的工作.
我們和人相處當中.
我們有沒有關愛的心態在裡面?.
如果我們能夠擁有這種.
關愛的心態的時候.
你是很不同的.
你和人聊天,說話是很不同的.
你的態度會比較謙卑.
會比較溫和.
你不會覺得自己是高高在上.
你不會覺得自己很了不起.
你很差.
不會,為什麼?.
因為關愛的心態.
令你將自己變得更謙卑.
當你和人聊天的時候.
你會想多一點別人的需要.
多於想自己的需要.
為什麼?.
因為你有耶穌那種關愛的心.
所以你和弟兄姊妹相處的時候.
或者和弟兄姊妹一起合作.
去配搭,去事奉的時候.
你那種著眼點.
你會放多一點.

$^{321}$在別人的需要那方面.
不是單單看自己的需要.
這是很不同的出發點.
大家認同嗎?.
所以你發覺到.
原來有這種關愛的心態.
是非常非常重要的.
大家記不記得.
在很多年前.
香港有個島叫做喜寧洲.
你們有沒有聽過這個島?.
這是一個孤島.
沒有人住.
但很多年前.
喜寧洲這個孤島住了什麼人?.
住了麻風的病人.
這班麻風病人.
在去到荔枝角醫院之前.
全部困在喜寧洲的孤島上.
弟兄姊妹們.
這個孤島有一個故事.
就是在這個孤島上.
原來有位牧師姓梁.
竟然每個星期都去喜寧洲.
探望這些麻風病人.
關心他們.
服侍他們.
有時和他們洗傷口.
有時和他們聊天.
更加重要的是.
這個牧師去到.
和他們講道.
你明白我的意思嗎?.
他一做.
你知不知道他做了幾年?.
一做就做了三十年.
在這個島上.
你說這位牧師.
令不令人敬佩?.
很令人敬佩.

$^{361}$一個孤島.
一個荒島.
交通很不方便的地方.
但這位牧師.
他不忘他的初心.
因為他知道上帝呼召他.
就是要關心這些麻風病人.
各位弟兄姊妹們.
德蘭修女的故事.
你們都聽過很多.
可能你會問.
牧師.
你說一些人物.
都好像很大人物.
好像很重要的人物.
我不過是一個.
很小的人物.
我很細微.
我很微小.
我做不到什麼大事.
其實不是.
各位弟兄姊妹們.
你是可以做到的.
為什麼?.
因為只要你擁有.
主耶穌那份愛.
就是你經歷了.
天父那份無比的愛.
感動了你之後.
你就能夠.
擁有這份關愛的動力.
去關心有需要的人.
所以你是可以做到的.
很簡單.
可能你去.
和一個人在電話裡聊天.
如果你有關愛的心.
你就能夠在電話裡聊天.
鼓勵到對方.
為什麼?.

$^{401}$因為你聽到對方的聲音.
是很低沉.
是很沒有心機.
是一個人很灰心.
但因為你有愛的緣故.
你就能夠在電話聊天中.
去關心他.
去鼓勵他.
和他同行.
這是什麼呢?.
這是很簡單.
在生活上很細微的事.
但你能否做到呢?.
你是可以做到的.
結果被你關心的那位朋友.
就感覺到你真的是一個.
很不同的朋友.
再久一點.
他知道因為你有耶穌基督的愛.
來關愛他.
所以你就吸引了他.
來主耶穌的面前.
又或者可能你在一些.
很細微的動作.
你每天都會做的.
你打開手機.
經常都會在群組裡.
你發訊息給對方.
如果你的內容是很正面的.
很鼓勵對方.
是很勵志的時候.
你雖然有一個短短的訊息.
但收到的人都得到什麼呢?.
很大的鼓勵.
各位頂尖的朋友們.
可能這些是很輕微.
很平常.
很簡單的動作.
其實你都可以做到的.
只要你擁有天父那份無比的愛.

$^{441}$感動到你.
你就能夠付諸行動.
也能夠作出犧牲.
以致你將耶穌基督的愛.
與人分享的時候.
你自己都很大的喜樂.
來充滿你的心.
也都有一種滿足感.
因為耶穌基督的愛.
令你能夠關心有需要的人.
我們今天結束的時候.
可否再讀一次今天的金句.
很簡單.
一節的金句.
預備起.
如今常傳的.
有信有望有愛.
最大的是愛.
哥林多前書十三章十三節.
我們一起低頭祈禱.
今天我們從一節熟悉的經文當中.
再去反思我們的生活.
是否真的好像世界的光.
來光照我們身邊黑暗的地方呢?.
我們是否好像一盞燈台.
擺在檯面上.
能夠光照全家的人呢?.
天父求你常常提醒我們.
可能一些平時很細微的動作.
很不起眼的一些行動.
但是因為背後有耶穌基督的關愛.
就能夠變成不一樣.
能夠在我們的生活當中.
可以實踐愛的功課.
愛的操練.
也都能夠去見證.
上帝你就是愛的源頭.
我們屬你的兒女.
也都成為你的器皿.
能夠去將福音傳開.

$^{481}$求主祝福帶領我們弟兄姊妹.
能夠享受到這一份.
從天父而來的愛.
能夠激勵我們的生命.
我們感恩禱告奉耶穌的名及.
阿們.
(下集待續).
\newpage



\section{路得記 1:1-22}
\label{sec:2v_ng9IARWc}
\textbf{2022.07.17 《 忠心:是禍抑是褔? 》路得記 1\_1-22 梁子勇牧師}
\newline
\newline
連結: \href{https://youtube.com/watch?v=2v_ng9IARWc}{\texttt{https://youtube.com/watch?v=2v\_ng9IARWc}} ~~~~ 語音日期: 2022-07-22
\newline
\newline
\hyperref[sec:EJ_GoCr8XKk]{\small{< < < PREV SERMON < < <}}
~
\hyperref[sec:index_chronic]{\small{[返順時目]}}
~
\hyperref[sec:index_scriptual]{\small{[返順卷目]}}
~
\hyperref[sec:rVzIB96PGyI]{\small{> > > NEXT SERMON > > >}}
\newline
\newline
路得記 1:1-22
\newline
\begin{longtable}{cl}
\hline
\hline
章節 & 經文 (和合本修訂版)\\
\hline
1:1 & \begin{tabularx}{0.7\textwidth}{X} 士師統治的時候，國中有饑荒。在猶大的伯利恆，有一個人帶著妻子和兩個兒子往摩押地去寄居。 \end{tabularx} \\ \\ \relax
1:2 & \begin{tabularx}{0.7\textwidth}{X} 這人名叫以利米勒，他的妻子名叫拿娥米；他兩個兒子，一個名叫瑪倫，一個名叫基連，都是猶大伯利恆的以法他人。他們到了摩押地，就住在那裡。 \end{tabularx} \\ \\ \relax
1:3 & \begin{tabularx}{0.7\textwidth}{X} 後來拿娥米的丈夫以利米勒死了，剩下她和兩個兒子。 \end{tabularx} \\ \\ \relax
1:4 & \begin{tabularx}{0.7\textwidth}{X} 兩個兒子娶了摩押女子，一個名叫俄珥巴，第二個名叫路得，在那裡住了約有十年。 \end{tabularx} \\ \\ \relax
1:5 & \begin{tabularx}{0.7\textwidth}{X} 瑪倫和基連二人也死了，剩下拿娥米，沒有丈夫，也沒有兒子。 \end{tabularx} \\ \\ \relax
1:6 & \begin{tabularx}{0.7\textwidth}{X} 拿娥米與兩個媳婦起身，要從摩押地回去，因為她在摩押地聽見耶和華眷顧自己的百姓，賜糧食給他們。 \end{tabularx} \\ \\ \relax
1:7 & \begin{tabularx}{0.7\textwidth}{X} 她和兩個媳婦就起行，離開所住的地方，上路回猶大地去。 \end{tabularx} \\ \\ \relax
1:8 & \begin{tabularx}{0.7\textwidth}{X} 拿娥米對兩個媳婦說：「你們各自回娘家去吧！願耶和華恩待你們，像你們待已故的人和我一樣。 \end{tabularx} \\ \\ \relax
1:9 & \begin{tabularx}{0.7\textwidth}{X} 願耶和華使你們各自在新的丈夫家中得歸宿！」於是拿娥米與她們親吻，她們就放聲大哭， \end{tabularx} \\ \\ \relax
1:10 & \begin{tabularx}{0.7\textwidth}{X} 對她說：「不，我們要與你一同回你的百姓那裡去。」 \end{tabularx} \\ \\ \relax
1:11 & \begin{tabularx}{0.7\textwidth}{X} 拿娥米說：「我的女兒啊，回去吧！為何要跟我去呢？我還能生兒子作你們的丈夫嗎？ \end{tabularx} \\ \\ \relax
1:12 & \begin{tabularx}{0.7\textwidth}{X} 我的女兒啊，回去吧！我年紀老了，不能再有丈夫。就算我還有希望，今夜有丈夫，而且也生了兒子， \end{tabularx} \\ \\ \relax
1:13 & \begin{tabularx}{0.7\textwidth}{X} 你們豈能等著他們長大呢？你們能守住自己不嫁人嗎？我的女兒啊，不要這樣。我比你們更苦，因為耶和華伸手擊打我。」 \end{tabularx} \\ \\ \relax
1:14 & \begin{tabularx}{0.7\textwidth}{X} 兩個媳婦又放聲大哭，俄珥巴與婆婆吻別，但是路得卻緊跟著拿娥米。 \end{tabularx} \\ \\ \relax
1:15 & \begin{tabularx}{0.7\textwidth}{X} 拿娥米說：「看哪，你嫂嫂已經回她的百姓和她的神明那裡去了，你也跟你嫂嫂回去吧！」 \end{tabularx} \\ \\ \relax
1:16 & \begin{tabularx}{0.7\textwidth}{X} 路得說：「不要勸我離開你，轉去不跟隨你。你往哪裡去，我也往哪裡去；你在哪裡住，我也在哪裡住；你的百姓就是我的百姓；你的神就是我的神。 \end{tabularx} \\ \\ \relax
1:17 & \begin{tabularx}{0.7\textwidth}{X} 你死在哪裡，我也死在哪裡，葬在哪裡。只有死能使你我分離；不然，願耶和華重重懲罰我！」 \end{tabularx} \\ \\ \relax
1:18 & \begin{tabularx}{0.7\textwidth}{X} 拿娥米見路得決意要跟自己去，就不再對她說甚麼了。 \end{tabularx} \\ \\ \relax
1:19 & \begin{tabularx}{0.7\textwidth}{X} 於是二人同行，來到伯利恆。她們到了伯利恆，全城因她們騷動起來。婦女們說：「這是拿娥米嗎？」 \end{tabularx} \\ \\ \relax
1:20 & \begin{tabularx}{0.7\textwidth}{X} 拿娥米對她們說：「不要叫我拿娥米，要叫我瑪拉，因為全能者使我受盡了苦。 \end{tabularx} \\ \\ \relax
1:21 & \begin{tabularx}{0.7\textwidth}{X} 我滿滿地出去，耶和華使我空空地回來。耶和華使我受苦，全能者降禍於我。你們為何還叫我拿娥米呢？」 \end{tabularx} \\ \\ \relax
1:22 & \begin{tabularx}{0.7\textwidth}{X} 拿娥米從摩押地回來了，她的媳婦摩押女子路得跟她在一起。她們到了伯利恆，正是開始收割大麥的時候。 \end{tabularx} \\ \\
[1ex]
\hline
\hline
\end{longtable}
$^{1}$各位兄弟姐妹早安.
高興能夠和各位兄弟姐妹在這個地方來敬拜.
洪水橋這個地方慢慢越來越出名了.
在整個北部都會.
我在新界東粉嶺區.
你們在新嶺西這邊.
整個北部都會.
不知道每月將會有多少人住?.
250萬人.
是250萬人.
所以求神祝福洪水橋堂.
洪恩堂在這裡繼續能夠有更加.
蒙神賜恩的發展.
所以今天能夠和大家一起去敬拜.
我們施香神教華語之師我們一起祈禱.
你的恩典承托我們的生命.
讓我們能夠活著.
不是枉然.
求你來施恩.
藉著你自己的華語.
光照我們所當行的路.
你是那位永活的上帝.
我們一餐一宿都滿有你的恩典.
最後求你來施恩.
讓我們不枉然領受你自己的恩典.
唯有你連戳.
唯有你敦促.
我們能夠快跑的來跟從你.
願你來施恩.
祝福我們今天的敬拜.
讓我們在中間與你相遇.
叫你的華語甦醒我們的心靈.
叫你的臨在.
能夠叫我們哪怕光暗的心都軟化.
俯伏.
為要尊崇你.
祝福我們中間每一位你所愛的.
誰聽我們在面前的祈禱.
奉耶穌聖名祈求.
阿們.

$^{41}$路得記其實是一卷.
很百感交集.
但是又容易錯解的聖經.
在看路得記的時候.
對應著今天我們的社會.
移民好還是不移民好呢?.
一罵我真的千車萬馬.
不過他又回流.
這段經文對你和我有什麼啟迪呢?.
在早前的時候認識一位姐妹.
她說在一年的時間裡.
兩個自己的兒子.
相繼離世.
白頭人送黑頭人.
是很大的悲痛.
剛剛這個星期我聽到.
有一位事奉神將近退休的路姐.
他的兒子離世.
很大的悲痛.
對經文來說.
那奧米的丈夫離世.
她兩個兒子又離世.
路德這一段婚姻.
還未能生兒時.
丈夫又離世.
三個寡婦走在一起.
一個年輕的跟一個年長的.
回到伯利恆這個地方.
就是路得記裡面的那段經文.
如果我想問你.
今天說忠心.
你做誰呢?.
你想做男或女還是想做路德呢?.
選擇什麼時間.
最好沒有禍患的時間.
男或女能夠最終有個孫子.
當乾了是他的孫子.
這樣就安樂了.
或者路德再遇到一個有錢人家.
那個時間就好了.

$^{81}$我們人生有些時候.
我們很想選擇某一個片段.
如果一生單單是這樣就好了.
路得記不是這樣去通訊.
在整個過程裡面.
神容讓他自己的選民.
失去自愛和倚靠.
其實充滿眼淚悲情的一個記載.
如果我們細心去掌握體會.
其實很難去明白.
究竟路得記是教我們做什麼呢?.
是不是教有錢人.
你懂得找.
跟就一定要跟.
這個幫下師後面.
這樣就有好處了.
這個婆婆也是這樣教你.
你聰明一點吧.
有些時候我們可能以為這樣是路得記.
不是這樣的意思.
絕對不是.
是不是奶奶就.
哇!我媳婦能夠改嫁了.
我將來就無憂了.
到了路得記最後就好像這樣.
還要是什麼?.
王的後裔.
可能我們中間有些弟兄姊妹.
未必看過路得記這段經文.
剛才讀第一章.
其實他想說什麼呢?.
路得記其實用一個愛邦的女子.
來羞辱神的選民.
容許我講是用羞辱.
路德勝過那奧米.
不單止這樣.
他以路德的忠貞.
路德這個字是解.
友情,友誼.
他的名字是有意思的.

$^{121}$來羞辱什麼人?.
那奧米.
作為神的選民.
他的小信.
一個愛邦女子.
對一個人的忠貞.
比你對上帝的承諾還要聖.
又或者這樣說.
其實同時之間.
是以一個乞丐的手藥.
乞丐來的.
路德的手藥.
來羞辱什麼?.
是時思記.
時思記的強者的放任.
如果你留意路得記.
其實和時思記是分不開的.
看一段聖經最好是什麼呢?.
你最好看前一節.
不要只看那一段.
路得記的前一節是什麼呢?.
時思記最後那一節.
時思記最後那一節是怎樣的?.
我們看看.
這一節沒有提及大家.
最後那一節.
"克雅舍以色列終沒有王".
"各人" 怎樣呢?.
"任意而行".
這個就是面貌.
所以帶出什麼呢?.
在路得記接續的第一節.
第二節就是那種情況.
大頭的老爺.
老爺叫什麼名字呢?.
"爾利米勒".
就是拿奧米的丈夫.
這個老爺的名字"爾利米勒".
你知不知道.
在希臘白萊文裡是有意思的.

$^{161}$是什麼意思呢?.
很有趣的.
什麼意思呢?.
"神是王".
是不是很諷刺呢?.
神是王.
明明是時思記最後那一節是怎樣的呢?.
"以色列終沒有王".
"爾利米勒"的老爺幫他取名字取得很好.
"神是王".
在你的生命裡神是王.
不過接續在第一第二節.
就讓我們看到是什麼呢?.
名字就叫"爾利米勒".
不過誰來作主呢?.
就是自己.
是任意而行.
所以這個最大頭的大男人.
戴瑜的家.
就離開伯利恆.
就去了摩押地.
頂尖無疑.
我想和你看看.
在整卷的路得記裡.
其實他想說什麼呢?.
整卷的重點是什麼呢?.
不是你和我牽掛著地上的明星財富或安穩.
焦點只有一個.
找應許.
不過這個字很弔詭.
什麼是應許呢?.
神你答應我給我安穩.
神你答應我要怎樣怎樣.
神答應過你嗎?.
沒有啊.
神有沒有答應過你情快快過去啊?.
我們有祈禱啊.
神你說你聽祈禱啊.
我又聽到你祈禱.
我什麼時候答應過你.

$^{201}$按你的心意去做呢?.
沒有啊.
神應許了什麼?.
我們要很細心地看聖經.
神答應過什麼?.
不要隨便用"應許"這兩個字.
清清楚楚是神應許的.
有時候有些頂尖無疑.
至於我和我家.
我們也試過這樣.
至於我和我家呢.
我們一定會得救的.
什麼時候有這麼說過應許是這樣的?.
沒有啊.
求神連續幫助我們.
找應許是什麼呢?.
今天跟大家思想中心的時候.
我想跟大家看看.
在這個中心的裡面.
其實我們很需要仰靠.
神的指引.
這些我剛才也提到.
是中心仰靠神的指引.
神的帶領是怎樣?.
伯利恆這個字.
知不知道為什麼?.
又有意思的.
原來這個字是什麼呢?.
糧倉.
行不行?.
是糧倉.
很吊詭的一件事.
在《路得記》裡面.
劈頭就告訴我們.
伯利恆遭遇什麼?.
饑荒.
神不對辦的.
你要知道.
《時史記》在什麼年代裡面.
是要進入迦南地.

$^{241}$是差派以色列人離開埃及.
過完紅海.
攻了耶利哥.
然後進入迦南.
滅了迦南七族.
整個大地.
神就說給你的.
現在會什麼?.
流奶與麥.
這班人選擇了在伯利恆定居.
還是糧倉.
結果是怎樣?.
饑荒.
以你物同南河米.
發腳踢.
不知是否學了伯蘭.
你知道我說誰嗎?.
伯蘭是韓國人.
遇到饑荒又怎樣?.
就走.
遊民族.
當然是去有草的地方.
丙姐妹.
神的引領是怎樣的?.
很有趣的.
有些時候.
原來福中是會有禍的.
正正在中間.
其實.
以你物同南河米.
怎樣去學習在信心中馴服?.
伯利恆這個地方.
糧倉之一.
不是純粹流奶與麥.
丙姐妹.
你也記得伯利恆是一個什麼地方.
將來主耶穌出世的地方.
糧倉的意思不是因為佛.
不是.
是因為這是君王出生的地方.

$^{281}$那種豐盛是屬靈的豐盛.
是從這位王帶來給神的子民的恩澤.
是這個糧倉.
所以神選民一切的豐盛.
是離不開什麼?.
這位葡萄樹.
這位君王.
對不對?.
我們所有的好處.
是不在神以外的.
所以伯利恆這個名字說得很清楚.
為什麼路德基對焦的是這個?.
你的寄望是什麼?.
你有沒有看上帝帶領是怎樣的?.
但是你看眼前的環境.
今天我們很容易看眼前的環境.
你在香港這個地方可能已經有30年50年.
甚至60年.
這三年將你翻一翻.
我是不是要留在這裡?.
離開不要緊.
路德基不是叫你離開還是不離開.
不是說這一點.
上帝在你生命的帶領.
你有沒有尊這位神為你的主?.
為你的王.
這才是焦點.
所以不要搞錯.
今天牧師帶路德基是叫我離開還是不離開.
沒有說過這件事.
忠心於誰?.
是上帝那份人靈.
以至你要得應祂之靈.
你要相信信耶穌是得著豐盛的生命.
你要一生來依靠祂.
經文裡面其實原來忠心是會有考驗的.
奈何米和路德其實真的很悲情.
孤寡窮.
有時最眼淵的是什麼呢?.
看見別人很好.

$^{321}$每個人都好.
我們兩個寡母頗醉眼淚洗臉.
晚上睡覺不知怎樣過.
明天就去做乞丐.
同胞很好過.
弟兄姊妹.
所以難怪奈何米在中間.
經歷這樣.
你留意經文裡面單單第一章.
他講了多少次.
好像耶和華在中間怎樣.
所有苦都是彷彿出自耶和華.
加在她的身上.
弟兄姊妹.
神有時很連述.
仍然在中間幫助奈何米.
讓她在當下有路德陪伴.
她原本打算一個人回去.
如果沒有路德每天幫她撿墨水.
她都不知吃什麼.
上帝原來有時在不同的考驗裡面.
仍然有關心你,安慰你.
與你同行的人在中間.
我看見弟兄姊妹.
紅春堂即將.
預備參加港九培練大會.
在教會裡面一起看直播.
弟兄姊妹你記不記得兩年前.
郭文慈牧師在中間去港島.
有這樣的聚會.
你知不知道郭牧師在家發生什麼事.
你可能不是特別留意或聽完.
其實郭牧師在過程中分享過.
他的孫子在泰國曼谷.
發現第三期的淋巴牙.
正正是他承擔培靈會講完的時間.
嘩!真是很致命.
你讓我安靜一下行不行?.
你給我考驗.
我同時也不需要將自己放在別人面前.

$^{361}$我後面還要掛心那個孫子.
如何面對手術如何面對化療.
嘩!所以弟兄姊妹很不簡單.
你要來就要走.
留下來就留下來.
搞到我這麼尷尬幹什麼?.
你明白嗎?.
你可以不是這樣的.
嘩!弟兄姊妹.
有機會聽回.
十天內郭牧師的分享.
又或者大白今次.
我所了解李思敬博士.
其實早前也曾經有重病.
他要揭下來一段時間.
放下做院長的職份.
弟兄姊妹.
上帝讓考驗淋到中間.
讓你和我不因為面對這些事.
我們就搖擺.
我們就退後.
弟兄姊妹.
哪怕有些事是上帝容許要管教我們的.
我們有些事很錯解.
管教是什麼呢?.
打兒子.
打兒子是什麼意思?.
頑皮.
你頑皮我就打你.
不過你不要搞錯.
不是這樣的.
希伯來書很清楚說.
為什麼要管教你?.
因為你是我兒子.
就是這麼簡單.
其實希伯來書十二章裡面說管教.
首先說管教的是誰呢?.
是主耶穌.
是讓他在苦難當中經歷以至完全.
所以十二章希伯來書的經文裡面.

$^{401}$叫你看眼看誰?.
看耶穌.
那忍受罪人這樣頂撞的.
你們要思想.
管教不是以為.
我都沒有做錯.
你為什麼這樣對我?.
我們經常這樣想.
做什麼事我會有不舒服的.
是不是我有罪.
是不是我哪些事情做錯.
你有這樣的反省是好的.
但是另一個角度就是.
不是單單是這一方面.
為什麼要有教練?.
幫助你操練整個過程.
無論你是彈琴.
又或者好像我們科音樂曲班.
我相信你們練了很久.
對不對?.
一定要操練的.
是不是?.
那一定要有教練.
要有師班長.
是不是?.
一起去做才行.
其實上帝的管教是這樣的.
不過我們不是很喜歡.
我們不是很喜歡.
除非什麼?.
我參加.
是不是?.
弟兄姊妹.
你參加科音樂曲班.
你要想清楚.
我們這些時間練歌.
我只是出席可以嗎?.
我只是你獻唱的時候出席可以嗎?.
當然不可以.
是不是?.

$^{441}$當然不可以.
上帝今天為什麼叫你和我玩活著?.
是要操練你和我的生命.
我們教會有些姊妹.
她每逢安息禮拜都會出席.
不過她很獨特.
她坐輪椅.
占仰為榮的時候.
其實全場基本上沒有人坐輪椅.
你留意安息禮拜真的不會有坐輪椅出席.
在占仰為榮的時候.
你知不知道.
無論信主或未信主的人看到的場景.
不是虛色地說.
有心.
你真的有心.
你知不知道去殯儀館的路.
很不簡單.
特別有些路.
斜坡很斜的.
怎樣上?.
是坐著輪椅上和下.
上帝有時候逐一不同的人.
經歷管教.
我講的是正面來毫數他的生命.
以至他能夠成為美好的見證.
神所使用的那種器皿.
像我剛才所說.
用路德修育納何米.
所以看到這位的肢體這樣的時候.
我們一班弟兄姊妹.
我們心想我們做了什麼.
遠遠及不上.
你明白嗎?.
遠遠及不上.
其實納何米正正要這樣去學習.
認識路德.
以至體會上帝在整個管教的過程中.
叫他滿滿的出.
空空的回.

$^{481}$這些扔掉不要了.
你就是抱著這些.
是不是? 亭姐妹.
他仍然很記掛.
我沒有了.
我沒有這些東西了.
我滿滿的離開伯利恆.
現在空手回來.
你知道上帝管教你嗎?.
拿走這些東西.
亭姐妹.
你多喜歡的東西都好.
有一天都要拿走的.
有哪樣東西不會拿走的?.
沒有.
是沒有的.
小心我們年長的亭姐妹.
我們有時候說什麼呢?.
最重要是健康.
這句話要小心.
你跟了我微信的人都在說這句話.
跟微信的人有什麼分別呢?.
你說最重要是健康.
你一定沒有健康.
這句話很矛盾.
不是說健康不重要.
我今早都有運動.
我不是說運動不重要.
我不是說健康不重要.
你要健康來做什麼呢?.
是服侍上帝嘛.
這才是.
不是最重要是健康.
教會說不回來了.
我要去做運動.
我要行山.
我跟這個那個去.
你真的最重要是健康.
你不用擔心.
你一定沒有.

$^{521}$是不是?.
你不用憂慮.
是真的.
亭姐妹.
求神幫助.
你看到這種管教.
他是要管教你和我的.
不是因為你頑皮.
是因為他恨鐵不成鋼.
他看重你和我.
他想造就你和我.
他想祝福你和我的生命.
不要在中間閒懶.
在那裡嘻.
在那裡以為安然的度日.
不是這樣的.
我這個年底退休了.
不過亭姐妹我想告訴你.
我兩年前告訴眾亭姐妹.
我做了三十年.
今天也是一樣.
正因為這樣我要逼自己要運動.
我每次都不想運動.
為什麼要弄到滿身汗.
這麼累.
但是亭姐妹.
你和我還活著幹什麼?.
怎樣去珍惜神給你和我的生命.
去事奉上帝.
讓你我.
今受管教的時候.
神就會連恤.
很有趣的.
他甘心面對羞辱.
他甘心然然回到伯利恆.
是最羞辱的一個地方.
他然然去到中間.
這是值得讚賞的一個地方.
他乞求神連恤我.
上帝不僅預備了老德陪他.

$^{561}$神也預備了幫阿施.
從人的角度來說.
這個好像是有錢人.
多少都可以有墨水.
你又看到.
剛剛那個時間是十墨水的時間.
亭姐妹.
神的連恤上上.
是你和我好像缺乏的時間.
有時候你不經意.
原來那些時間是你最經歷上帝的時候.
亭姐妹.
你回望你自己屬靈的過程.
這個旅程的時候.
你最經歷神觸動你生命是什麼時候?.
可能是你最挫敗.
最低谷.
最有罪疚.
最患難的時間.
忽然間你發現.
原來神最連恤你.
就是那些時刻.
所以整個老德記的場景.
原來十墨水是最好的.
有墨水十,你就十吧.
你就十吧.
求神幫助你和我看到.
神的預備.
叫你今天和我還能夠活著.
我在教會裡和亭姐妹說敬拜.
其實敬拜是什麼?.
你知不知道你回來禮拜堂.
第一樣東西應該看到什麼?.
如果用一個比方.
就應該看到一個行刑的工具.
如果在紅欣堂.
假設.
你不要學,也不要怪我這樣說.
有一個斷頭台在紅欣堂門口.
你進來還是不進來?.

$^{601}$疫情嚴重一點我也不進來.
有一個斷頭台.
玩什麼呢?.
不過亭姐妹.
你知不知道聖殿是這樣的?.
聖殿的頭是什麼?.
是安陽的地方.
宰殺的地方.
為什麼要搞這樣的東西?.
你去見神.
你什麼是新鮮蘿蔔皮?.
你憑什麼去見神?.
第一個字應該怎樣?.
該死.
行不行?.
是該死的.
無論是罪是該死.
又或者你獻情你自己.
當繁祭.
都是什麼?.
該死.
所以死是第一個字.
返教會.
死.
好.
我死也返.
這樣才是返教會.
說真的.
你看回舊約的聖殿.
其他地方你不能進去.
那些是給利未人進去.
祭司進去.
什麼時候輪到你進去?.
你進去就是進去什麼?.
這個祭壇.
他那裡的祭壇就是這麼多.
所以去聖殿是什麼?.
死.
殺了羊不關我的事.
什麼不關你的事?.

$^{641}$你就是說關你的事才來的.
我死.
我應該死.
可不可以聽姐妹?.
今天你還活著出去.
是連續的.
全是連續.
求神幫助我們對敬拜有更加深入的認識.
敬拜不是聚會.
遠遠大於聚會.
敬拜是你與神相遇.
是相遇.
真正的主角不是聚會.
是神自己.
唯有你和我能夠看得到.
我們才知道什麼是真實去敬拜上帝.
以至其實整個路得記最大的恩慈是什麼?.
就是歸屬.
歸屬什麼?.
不是幫阿斯.
不是因為大衛王.
歸屬了大衛好還是不好?.
其實有很久都沒有.
那個是什麼?.
他的祖父.
俄比德.
可以嗎?.
然後才是耶西.
接著才是大衛.
那麼歸屬大衛好還是不好?.
你想清楚.
大衛的家族很多人死了.
祭司是橫死的.
還要被兄弟殘殺.
他的大兒子暗亂被誰殺了?.
四兒子阿沙隆殺了他.
阿沙隆後來也不好.
是不是?.
你發現歸屬大衛有什麼好?.
其實也不好.

$^{681}$如果你的後裔是這樣的.
你也擔心.
糟了.
如果我的後裔將來是這樣.
所以歸屬大衛有什麼好?.
看路得記你要看到.
不單單看片段.
歸屬大衛有什麼好?.
是因為歸屬.
神的選民.
神的國度.
神的君王.
我們這些弟兄姊妹.
你今天是不是真的歸屬基督?.
你今天是不是已經受了洗?.
我這個星期的時間.
有弟兄姊妹慣常都會這樣.
轉會也好.
不見了洗禮證也好.
牧師.
你幫我寫個證明書.
我真的在洗禮.
我心想.
什麼是洗禮證?.
什麼是洗禮?.
是不是證明你和神的關係.
是仍然常新的?.
是這樣去毀身.
去馴服.
去跟從主.
還是說我洗了禮.
沒什麼意義.
拿著洗禮證.
你和這位上帝的關係是怎樣的關係?.
這才是最關鍵最重要的.
你告訴別人有什麼意義?.
幫到你什麼?.
弟兄姊妹.
今天.
你和我受了洗.

$^{721}$你和我是那位上帝收納成為我們.
做祂的兒女.
讓你我能夠看到.
其實在馬福音裡面很清楚.
很清楚地說.
當你認順耶穌是基督的時候.
他就告訴你.
你不要搞錯.
跟這個君王.
這個君王做什麼的?.
大家讀一讀.
31節.
31節準備請.
從此他教訓他們說.
人子必須受許多的苦.
被長老.
祭司長王.
民事氣絕.
並且被殺.
過三天復活.
你知不知道我是一個這樣的君王?.
你認順是.
是神所差來的君王.
這位基督.
認順完劈頭立刻說.
你又要跟從我.
我們讀一讀接著那段經文.
34節至36節.
於是叫眾人和門徒來.
對他們說.
若有人要跟從我.
就當捨己.
背起他的十字架來.
跟從我.
因為凡要救自己生命的.
必喪掉生命.
凡為我和福音喪掉生命的.
必救了生命.
人就是賺得全世界.
賠償自己的生命.

$^{761}$有什麼益處呢?.
我們是不是真的選擇一條.
火裡火裡.
水裡水裡.
都跟從著的道路.
是這樣的一個跟從.
求神幫助我們看到.
我們有時很容易錯解什麼是忠心.
請你留意.
忠心不是對一份工作.
忠心.
我幾十年都是這間公司.
那不是.
不是這裡說忠心.
我很顧家的.
我最顧家的.
那些書.
全部都是我關心我照顧的.
不是說這個忠心.
欣賞你.
不過.
上帝所說的忠心.
不是這種忠心.
有時我們對崗位的忠心.
又或者.
我們對一些即時忠心.
我是三十年裡.
超過三十年.
是同一間堂會事奉的.
不過沒有意義.
外面風大雨大.
當然不要走.
有什麼意義呢?.
看東西兩邊看.
你明白嗎?.
沒有特別意義.
有些忠心.
我對聚會最忠心.
我未缺過崇拜.
未缺過祈禱會.

$^{801}$有沒有用?.
沒有用的.
頂尖的.
不是這種忠心.
也不是對洪恩堂的忠心.
不是柴台.
不是同桌.
也不是對宣道會的忠心.
我是宣道會的人.
洪恩堂以前的舞堂.
或者不是以前.
現在的舞堂在哪裡?.
感受堂.
我對感受堂忠心.
我對華人.
又或者.
甚至對我們的理想.
對制度的忠心.
無論是政見.
又或者是事後.
運動都可以很忠心.
是不是? 弟兄姊妹.
我和這個球隊.
弟兄姊妹不是這種忠心.
《聖經》裡面所說的忠心是什麼?.
忠心是關係.
是你和神的關係.
是創造主的關係.
如果你忠心的話.
不要受到奉獻這麼少.
行不行? 弟兄姊妹.
因為你和我十分之十都是誰的?.
是上帝的.
你十一奉獻就說忠心嗎?.
那十分之九去了哪裡?.
我也要養妻活兒.
我知道.
是因為上帝要你這麼做.
行不行?.
不是十分之九隨你自己怎麼洗.

$^{841}$是不是?.
要清楚才行.
不要以為我十一奉獻就很忠心.
沒有這樣的事.
是關係.
所以同時是救贖主和你的關係.
今天你蒙潔靜.
我們經常很想赦免.
赦免了你和我做什麼?.
安於繼續犯罪.
是不是這樣?.
不是吧?.
你和我潔淨就是要給神用.
是給他用.
為什麼要洗了杯子?.
要倒水.
要用.
是這樣的.
忠心就是忠心.
永生上帝.
所以.
活著還是死了.
死了還是主的人.
是主的人.
他是永生神.
但是我想你知道.
忠心不是單是對神的.
包括什麼?.
對群體.
今天我們在路德基的時候.
為什麼說伯利恆?.
其實那是神的選民.
不是抱著地方.
不是.
是神的選民.
今天你和我不用去伯利恆.
但是我們要緊扣對教會的承擔.
我們忠心於神的託付.
神的國度裡面.
馬化帝兄姊妹.

$^{881}$路德基完了的時候.
其實他們還沒有得到應許.
神的國是不是真的完全臨到?.
他還沒有完全擁有.
還沒有得到皇權.
這是大胃王.
而且接下來還會有戰亂.
弟兄姊妹.
讓你和我.
仍然等候那份應許.
應許的成就不是甘心.
不是地上的國.
是天上的國度.
讓你我在忠心裡面.
是一種對神關係的那種忠心.
以致或左或右.
我有心有所屬.
一生去毀身.
不是單單是和神的關係.
和他的聖教會.
和他眾多子民.
能夠一起去承擔.
伯莫弟兄姊妹.
你和我能夠剛強壯膽.
事奉這位永生上帝.
何等的母恩.
我們一起同心祈禱.
主一切都是出於你的廉恕和智慧.
又唯有你的護蔭.
你的保守.
讓我們可以剛強.
可以壯膽.
主願意你施恩.
扶立我們中間每一位你所愛的.
你多年來.
為要造就和建立我們.
讓我們不輕看這個還活著.
可以給你使用的時間.
願主你兼顧我們.
幫助我們每一位.

$^{921}$能夠成為你手中的器皿.
給你使用的瀑被.
高舉主你自己的名.
聽我們的祈禱謝恩.
奉耶穌聖名祈求.
阿們.
\newpage



\section{出埃及記 12:40-41}
\label{sec:rVzIB96PGyI}
\textbf{2022.07.24 《榮耀的軍隊》 出埃及記 12\_40-41 (和修版) 曾敏姑娘}
\newline
\newline
連結: \href{https://youtube.com/watch?v=rVzIB96PGyI}{\texttt{https://youtube.com/watch?v=rVzIB96PGyI}} ~~~~ 語音日期: 2022-07-25
\newline
\newline
\hyperref[sec:2v_ng9IARWc]{\small{< < < PREV SERMON < < <}}
~
\hyperref[sec:index_chronic]{\small{[返順時目]}}
~
\hyperref[sec:index_scriptual]{\small{[返順卷目]}}
~
\hyperref[sec:pNmpdKyDl3E]{\small{> > > NEXT SERMON > > >}}
\newline
\newline
出埃及記 12:40-41
\newline
\begin{longtable}{cl}
\hline
\hline
章節 & 經文 (和合本修訂版)\\
\hline
12:40 & \begin{tabularx}{0.7\textwidth}{X} 以色列人住在埃及共四百三十年。 \end{tabularx} \\ \\ \relax
12:41 & \begin{tabularx}{0.7\textwidth}{X} 正滿四百三十年的那一天，耶和華的全軍從埃及地出來了。 \end{tabularx} \\ \\
[1ex]
\hline
\hline
\end{longtable}
$^{1}$請各位靈性的朋友回到家中.
以色列人住在埃及共四百三十年.
正滿四百三十年的那一天.
耶和華的全軍從埃及地出來了.
這裡有提到軍隊.
在說的是出埃及的時間.
耶和華的全軍從埃及地出來了.
全軍是指甚麼呢?.
即是整個軍隊.
整個軍隊包括了甚麼呢?.
是所有出埃及的.
屬於耶和華的子民百姓.
大家如果對於出埃及記略有所認識.
就知道在這個之前.
其實以色列民一直在面對埃及人很大的逼迫.
埃及的法老覺得.
以色列民這麼多.
會成為我們埃及國的一個很大的威脅.
於是就要他們做苦工.
逼迫他們.
甚至在想有初生的男嬰.
那些以色列民的男嬰.
就是要殺死他們.
總之就不要讓他們長大成人.
要限制和控制著以色列民的人數的增長.
在這種環境下.
以色列民是過得非常辛苦的.
他們又要做苦工.
又要面對法老的種種逼迫.
叫苦連天.
他們去哀求的時候.
上帝是聽到他們的聲音.
於是上帝就要差派摩西亞倫.
要去帶領以色列民出埃及.
這個過程非常艱鉅.
摩西亞倫要去見法老.
要跟法老說.
放我們的百姓出埃.
法老的心很硬.
大家還有印象.

$^{41}$上帝是走了十災.
走了十災.
最後是殺長子之災.
之後.
法老才終於說.
好了,走吧.
你們快點走吧.
就是這樣.
這一次.
就是這個時間.
以色列人住在埃及.
共430年的這一天.
耶和華的全軍.
從埃及地出來了.
出來的是全軍.
所有的子民百姓.
不只是那些在前面.
拿刀拿槍的人.
大家印象很深刻.
摩西亞倫是稍後.
摩西亞倫帶領以色列民.
過紅海的時間.
但背後就是這一群百姓.
是非常浩浩蕩蕩.
很多人.
裡面包括了男人和女人.
年輕的.
中年的,老年的.
老人,婦女,孩子.
都是其中的一份子.
神清所有這些人.
是耶和華的全軍.
大家覺得特別嗎?.
其實經文可以這樣說.
換另一種說法.
就是正滿430年的那一天.
耶和華的子民百姓.
所有的子民百姓.
從埃及地出來了.
都可以.

$^{81}$完全講得通.
但經文不是這樣說.
就是要耶和華的全軍.
從埃及地出來.
每一個.
無論在外表上.
看來他們能不能打仗.
實際看起來.
能力是怎樣.
上帝眼中看他們.
都是軍隊的一份子.
我們去看.
究竟.
是不是呢?說真的.
如果真的有戰爭發生.
真的每一個子民百姓.
都去到前方.
真的拿刀拿槍去打?.
是不是呢?.
我們看過.
我曾經在講道裡說過.
在出埃及之後.
以色列人曾經在尼菲定這個地方安營.
亞馬力人來跟他們打仗.
大家記得嗎?.
那一節我曾經說過.
當時在十七章九節裡說.
摩西對約書亞說.
你為我們選出人來.
出氣和亞馬力爭戰.
原來這段經文有沒有說.
所有人都會出去打.
選人去打.
是不是?.
那時候是選人去打.
其實經文並沒有詳細說.
約書亞怎樣選人.
是選一些人.
大家還記得.
在山下.

$^{121}$在地面上的時候.
就是約書亞.
帶領他所選的人.
在那裡打仗.
在山上摩西在那裡.
禱告.
那場仗是這樣的.
有人在前方打仗.
有人在山上禱告.
接著你會問.
沒有被揀選去打仗.
又不在山上禱告的.
在那裡做什麼呢?.
上帝稱整個百姓.
所有子民百姓.
是他的軍隊.
上帝是這樣看的.
那其他人在那裡做什麼呢?.
經文沒有說.
會不會你們在想.
前面有人打仗.
上面有人祈禱.
於是其他人就想.
我又不是祈禱.
我又不是打仗的.
所以玩吧.
我們自成一角.
我們自己玩.
我們休息.
我們休養生息.
前面的戰場.
和我們後方.
是不一樣的.
我們有我們的世界.
這樣是不是軍隊呢?.
我看到弟兄姊妹搖頭.
怎麼會是軍隊呢?.
軍隊會這樣嗎?.
其他人.
一定不敢這樣.

$^{161}$因為這段時間只是打仗的時間.
打仗的結果絕對.
是關乎大家的安危.
所以大家都很緊張.
其他人可能在那裡做什麼呢?.
默默地為前方的士兵打氣.
專注地留意戰況的進展.
其實那一天.
約書亞這場仗.
都不容易打.
摩西什麼時候舉手.
以色列什麼時候得勝.
大家記不記得?.
他什麼時候手垂下來累了.
他們就在輸了.
所以一直起起伏伏.
起起伏伏的那天很難打.
我相信其他人.
看著的時候不敢.
在那裡聊天.
玩耍,休息.
不敢.
看著輸了.
贏了,加油啊.
可能就是背後在做這些.
甚至他們不在山上.
不是摩西那樣禱告.
他們也可以在下面禱告.
軍隊的一分子.
是我們不分彼此的.
是一齊的.
這場仗不是牽涉到前方的人.
所有人都有份.
是牽涉在裡面的.
既然臣事所有的人.
為軍隊的成員.
也就是說不在前方打仗的成員.
都要有所參與.
軍隊裡面所有的人.
都是同心合意.

$^{201}$向著同一個目標進發的.
不會這樣的.
在前方打仗的.
祈禱的那些.
你們自己想想.
你們想怎樣就怎樣.
我們後方有我們的想法.
不會的,同一個目標.
希望這場仗要贏.
簡單的說.
他們前方努力.
我們後面撐著他們.
我要讓他們看到.
他們不是孤單的.
在那裡打仗.
是很同心的.
從另一個角度去看.
摩西帶領百姓.
過紅海的時候.
後面的群眾.
印漢.
耶和華的百姓子民.
你看看出埃及的那一群人.
剛才說的,有男的,女的.
不同年齡層的.
有婦孺的,有老人家的.
有孩子的.
有精壯的男士.
一群人,有不同的狀況的人.
可能外面的人看.
耶和華的百姓是一群.
很軟弱的人.
都不像軍隊.
能打嗎?.
怎會能打呢?.
都不專業.
應該一群人都不能打.
這麼軟弱.
這些上帝的百姓.
上帝看他們百姓.

$^{241}$是怎樣的?.
榮耀過軍隊.
上帝看和人看.
不一樣.
他們不是軟弱的人.
這些人不是不能夠.
參與的,不是不能夠貢獻的.
不是只分出一群人.
那一群就是精兵.
其他就是.
沒有價值的,不是的.
整個群體是軍隊.
在上帝眼中.
是榮耀的.
為何我會說是榮耀呢?.
在哪裡有個字眼說榮耀呢?.
為何上帝看是榮耀呢?.
因為上帝自己.
是他軍隊的元帥.
上帝自己是榮耀.
而他是這個軍隊的元帥.
就不一樣.
今天要看誰領兵.
要看誰領兵.
如果你看回.
世界歷史,看過不同國家的歷史.
去看打仗.
不是每個人都能打仗.
有些人逢打仗必輸.
有些人一出來打仗就很厲害.
一出來就說.
這個會打仗的人.
上帝更加厲害.
他是絕對勝利的.
傳言榮耀的軍隊的元帥.
這個元帥.
在聖經裡面.
有個稱呼.
是萬軍之耶和華.
我們讀的時候.

$^{281}$很快就讀到萬軍之耶和華.
其實你知不知道.
萬軍之耶和華是甚麼意思呢?.
是統領萬軍.
在地面上.
軍王或軍隊的元帥.
再厲害也好.
他的軍隊是有限量的.
沒有一個.
可以有這麼榮耀的稱呼.
是萬軍.
去統領的,所以他是最大的.
只有耶和華神.
有這個能力.
有這個權柄,有這個威榮.
是萬軍之耶和華.
我們有沒有想過呢?.
平時讀聖經.
因為讀得太慣了,太上口.
好像覺得沒甚麼特別.
在這個世界,這個宇宙.
最厲害的.
最厲害的元帥.
他一出,沒有人可以跟他比.
在這裡.
看看.
《倚創亞書》六章一到三節.
是一段很寶貴的經文.
當烏西亞王崩的時候.
當時.
先知看見.
主坐在他的寶座上.
講到他那種榮耀.
講到我們的主.
他的衣裳下擺.
遮滿聖殿.
然後講到撒拉夫的.
敬拜讚美.
一路讚美他.
聖哉!聖哉!聖哉!讚他甚麼?.

$^{321}$萬軍之耶和華.
他的榮光.
遍滿全地.
從來沒有一個元帥.
在地面上擁有一種這樣的榮耀.
唯有上帝可耀.
因為他.
統領萬軍.
這位統領萬軍的元帥.
他戰無不勝.
一定贏的.
我們很快.
之後你看到出埃及記.
是十四章.
是這章之後的兩章.
你看到神就帶領.
以色列民過紅海.
這個過程.
多厲害.
說真的.
當中.
不是以色列民自己去打.
純粹是元帥.
自己出手.
就是元帥自己出手.
他完全勝過.
埃及的軍隊.
當時埃及的軍隊.
是很有名聲的.
他有很精良的裝備.
他的武器很好.
很多戰車.
軍隊有嚴密的結構.
很善於打仗.
這些人一定.
不會說他是烏合之眾.
以色列民反而有機會.
人們看到會覺得.
這些烏合之眾.
他們不是的,很懂得打仗.

$^{361}$很多裝備在.
這個這麼強的軍隊.
一直浩浩蕩蕩的追著.
以色列民的時候.
你看到上帝為以色列民開路.
一直保護他們.
拖大他們.
讓他們順利的過紅海.
紅海分開.
以色列民.
走在乾地上.
一路平安.
以色列民過了紅海之後.
紅海才合起來.
上帝用他很大能的作為.
讓整個埃及的軍隊.
全軍覆沒.
萬軍之驕.
一出手就贏.
他根本不需要人.
去幫他.
這一場仗更加厲害.
在以上亞書有一場仗.
37章記載.
當時亞述王派軍隊.
攻擊猶大國.
而且他辱罵.
以色列的神.
他覺得根本沒有一個國家的神.
可以救他的百姓.
離開自己的手.
當時亞述.
國勢非常強盛.
軍隊也很強盛.
他當時被人認為.
戰無不勝.
所以他去哪裡一定贏.
所有國家都要投降.
你以色列民算得了什麼.
你以為上帝可以救你.

$^{401}$離開我的手嗎.
他真的這樣辱罵上帝.
結果神派使者.
去亞述營當中.
殺了18萬5000人.
在這一場仗裡都是一面倒.
上帝根本不需要.
地面上的軍隊.
他自己出手.
完全.
完成了.
他要去贏.
一點也不難.
是戰無不勝,是甘有能力.
就是這位萬軍之耶和華.
戰無不勝的.
耶和華神.
是以色列軍隊的元帥.
有機會大家可以重溫.
這場這麼厲害的戰爭.
還有.
為什麼說這軍隊是榮耀的軍隊.
除了元帥.
是非常.
來頭很強.
還有.
榮耀是因為什麼.
元帥帶領他的軍隊.
離開黑暗.
遺奴之地.
讓他們尊崇自己.
成為榮耀.
以色列民在埃及.
過了很多年的苦日子.
經歷被逼白.
面對生命的威脅.
很多年.
在埃及生活.
特別是後期的日子.
一直做奴隸.

$^{441}$捱世界.
沒有自由,沒有尊嚴.
被逼白.
但是耶和華神卻將他們.
帶離.
埃及.
埃及對於他們來說.
是黑暗之地.
上帝救他們.
給他們自由.
給他們尊貴的身份.
從今以後.
這群白人是專屬於自己.
上帝自己是榮耀的.
這群白人屬於自己.
就是榮耀.
而這群人神愛他們.
所以視他們都為榮耀.
是因為神愛他們.
所以他們成為.
一個很榮耀的軍隊.
原來如此.
我們看回.
整個以色列民.
一直經歷.
上帝如何帶領他們得勝.
重點是.
不單是昔日.
上帝看他的百姓.
是他的榮耀軍隊.
今天我想說.
教會也是軍隊.
今天我們洪恩嘉.
是上帝軍隊的一步伐.
要為上帝做什麼?.
要為上帝征戰.
搶救靈魂.
去拓展上帝的國度.
我想昔日.
錦繡堂差派我們.

$^{481}$來到洪水橋這席.
堂就是有.
神的美意在.
想在這裡拓展上帝的國度.
傳播福音.
其實換一個角度看事.
也是為上帝征戰.
我們是軍隊.
雖然上帝會挑選弟兄姊妹.
承擔不同職責.
有不同事奉.
但可以肯定.
每個弟兄姊妹.
都有份.
在軍隊裡.
我們是有份的.
我們不要說.
上帝一定選了這班弟兄姊妹.
他們負責這些部份.
上帝一定選了這些.
他們是這樣參與的.
你就放過我吧.
我就差不多了.
總之我們分開.
兩大陣營.
這些是屬於他們的.
很舒服的生活.
我們彼此.
不一樣的.
不可能的,全部一樣.
職責是有不同.
但你記住我們是軍隊的一部份.
我們是軍隊的一部份.
我們有份的.
我想說.
有弟兄姊妹參與.
前方的事奉.
我和你.
不是局外人.
你記得嗎?.

$^{521}$約書亞那場仗.
有人打仗,有人在山上禱告.
我們是怎樣的?專注的.
留意進展,我們關注.
我們支持的.
我們不會這樣,教會發生什麼事.
我不知道.
某些人知道,你問那些知道的人.
不是的,應該整間教會一起知道.
我關心的.
不只是傳道人在上面.
在那裡禱告.
每個人都可以祈禱.
今天紅印加.
我看下去,有幾十人在這裡聚會.
每個人一起祈禱,力量多大.
而且是同心合意的祈禱.
團結就是力量.
不應該分開.
一堆堆,各有各做.
我們彼此不了解對方.
甚至我找不到我的位置.
在哪裡,那是做什麼呢?.
其實我也不知道要負責什麼.
我這些做不到,那些又做不到.
不行了,你放過我吧.
我回到家,見到朋友,我已經很開心.
聊聊天.
享受我的人生.
這樣就很開心了.
對不起,上帝不是這樣看.
上帝不是這樣看我們.
是局外人.
兄弟姐妹今天不要看著別人.
我和你們說話.
和你們每一個人說話.
你是一份子.
你有份的.
你是榮耀軍隊的一部分.
你是不可或缺的.

$^{561}$少了你的支持.
少了你的參與,是不行的.
這個軍隊就不完整了.
上帝看我們,是很不完整的.
上帝看我們,是可以參與的.
因為聖靈在我們當中會給我們恩賜的.
這裡沒有人沒有恩賜的.
每個人都有恩賜的.
聖靈給我們的.
一定在這裡有位置.
我問你,你在這裡找不到上帝給你的使命呢?.
我在團體裡面.
很喜歡挑戰弟兄姊妹做一件事.
我們未必能夠跑跑蹼蹼.
走很多地方.
可能體力勞動很強的一些事情.
未必做到.
但我們很厲害的.
我們成為禱告兵團,是不是?.
我不懂得如何祈禱.
可以學習的.
所以就要多點回來祈禱.
坐在這裡不會學到祈禱.
我們姐妹一起去祈禱.
一邊去祈禱,一邊去學習.
自己回去又操練.
自己又多點祈禱.
回來又多點祈禱.
就很厲害了.
你有份的.
你的祈禱是可以祝福很多人的生命.
例如在家裡.
我們今天整個教會的群體.
我們的年齡是比較成熟的.
但我們可以祝福很多人.
很多人的.
我們祝福我們的兒女.
為他們禱告.
為他們的信仰禱告.
讓他們回教會.

$^{601}$我們帶他們回教會.
為我們的孫兒禱告.
為我們的鄰居禱告.
我們都會出去做運動的.
認識很多街坊.
很開心的跟他們聊天.
我們不單止跟他們分享生活上的事.
我們都可以跟他們說.
「回教會」.
「我可以為你祈禱」.
有很多很多的事可以做.
今天就在這個層面.
我們都可以一起幫忙.
拓展上帝的國度.
可以很忙碌的,很精彩的.
弟兄姊妹,你找不找到你的使命?.
很想挑戰大家去想.
你是軍團的一部分.
上帝看我們是榮耀的.
有能力的.
因為元帥的緣故.
因為我們上帝的緣故.
教會裡所有的人.
都是同心合意向同一個目標進發.
不會你有你的.
我有我的.
我喜歡什麼就什麼.
你喜歡什麼就什麼.
不是,是同一個目標.
弟兄姊妹,我想問洪恩堂的目標是什麼?.
有人想嘗試回答.
你覺得我們同一個目標是什麼?.
大聲一點說.
敬拜上帝.
非常好.
我們同一個目標,要敬拜讚美上帝.
傳福音.
傳福音,是吧?.
洪水橋有很多人未信耶穌.
我和你都有份去傳福音.

$^{641}$不只是一小撮人.
大家都有份,有很多方法.
傳福音在這個世代有很多方法.
祈禱.
我發覺每個星期.
最了解教會發生什麼事.
通常是星期二,三,四早上.
參與我們網上祈禱會.
或者星期三晚.
參與我們Zoom.
祈禱會的弟兄姊妹.
其他真的要了解的話.
也不是那麼容易.
除非身邊的人.
和你分享.
所以在過程中再想.
我們可不可以每個人都了解.
教會的事多一點.
多一點的禱告.
你問上帝多一點.
我在這個軍團裡.
這個軍隊裡,可以怎樣參與.
我的角色是什麼?.
上帝想洪恩嘉向哪裡發展.
是不同的.
今天的我們是否同心呢?.
問問自己.
神體教會是榮耀的.
洪恩嘉也是榮耀的.
因為是屬於他的.
以色列民是榮耀的.
因為耶和華神是他們的元帥.
耶和華神也救贖他們的生命.
今天上帝看教會是榮耀的.
因為我們是屬於他的.
我們的元帥是耶穌基督.
他是榮耀的.
他是萬王之王,萬主之主.
有一天耶穌基督要再回來.
要審判萬民.

$^{681}$我們也要看到耶穌基督的榮耀.
而教會的元帥就是耶穌基督.
還有,耶穌基督是全然得勝的主.
他絕對可以帶領他的軍隊.
去勝過仇敵.
教會很重要的是什麼?.
為上帝爭戰.
拯救靈魂.
拓展上帝的國度.
在過程中絕對有仇敵.
就是那惡者.
他絕對不想上帝的國度拓展.
但是今天上帝.
是我們的王,我們的主.
他就會帶領我們得勝.
耶穌基督會帶領我們得勝.
要去勝過仇敵.
所以教會是榮耀的.
還有,這個元帥耶穌基督.
也帶領我們離開罪惡的綑綁.
脫離永遠的死亡.
去救我們的生命.
讓我們專屬他.
成為榮耀的.
所以.
神體教會是榮耀的.
我們有沒有讓人看到.
在我們的生命.
看出上帝軍隊的榮耀呢?.
大家有沒有膽量去問.
別人.
覺得我們這個軍隊.
是否很榮耀.
我們整個生命的見證.
是怎樣的.
上帝看我們很榮耀.
上帝也期望我們是這樣的.
成為一個榮耀的軍隊.
但是我們被別人看到.
是否一個榮耀的軍隊呢?.

$^{721}$我們被別人看到的是甚麼?.
問問自己.
這一刻.
是否別人會看到.
你們是榮耀的軍隊,你們屬上帝.
你們信耶穌是特別不同的.
是否這樣呢?.
還是看到另外一些東西.
和榮耀相反的.
和德性相反的.
是甚麼呢?.
值得我們去反省.
這一下要想.
上帝看我們是榮耀.
我們怎樣活出這個生命.
有沒有定精在主的身上.
看著他.
他是我們的元帥.
所以我們要聽他的聲音.
最怕的就是.
元帥就很厲害.
元帥全有能力.
帶領我們去打那場仗.
不過我們沒有聽他的聲音.
我們沒有聽他的聲音.
你有你說.
我有我生活.
這樣就分開了.
這就不是軍隊.
士兵和元帥分開了.
我不太知道元帥想怎樣.
我也不打算問他.
因為我太習慣.
我有自己的生活.
我自己就是這樣繼續下去.
但感覺上又在做很多.
好像很好的事.
這個軍隊.
不會是榮耀.
也不會是特性.

$^{761}$因為沒有了元帥.
其實不是他不存在.
而是我們沒有把他當元帥.
還有.
你聽他的聲音.
會不會去跟.
我很喜歡挑戰一件事.
最近有和第一節目討論.
開會文化.
教會.
我指的是香港的教會.
也很喜歡開會.
透過開會解決問題.
不開會怎樣解決.
沒理由自然而然.
事情就能夠.
完成,辦好.
要透過開會.
沒錯,是需要開會.
但在過程中.
在之前,是否有很多禱告.
其實可以有很多東西.
可以脫勾.
不是貨幣.
或者經濟體.
是脫勾.
是與元帥的關係.
可以脫勾.
其實耶穌基督.
有對教會的很多看法.
但我們不去問他.
我們就想.
憑我們的聰明智慧.
我們透過開會.
談好他.
ABCDEFG就會出來.
那就搞定了.
感覺有多好.
好像頭頭是道.
但你記住那些是自己的想法.

$^{801}$不是我們真正的.
我們的王,主,基督的看法.
這樣下去.
這就不是軍隊.
是否元帥要怎樣發號施令.
有沒有人聽呢.
接著可能是我們自己談了.
這一定是上帝的心意.
我知道你沒有反對.
因為我們一向都是這樣.
我們談完之後.
我知道上帝你也沒有反對.
就是這樣走.
其實.
這樣是不能真正.
去得勝.
也不是.
真的可以去拓展上帝的國度.
也不會讓人看到.
上帝的榮耀.
上帝帶領.
和人帶領是兩回事.
是否在過程中.
如果你說真的.
要去活出榮耀的生命.
整個榮耀的團隊.
是要去聽.
上帝的聲音.
你要聽就去禱告.
去等,不是一個半個人在祈禱.
在等.
是整個教會一起.
在未來的日子.
我深信.
上帝要去.
大大的發展洪水橋區.
無論是我們的交通.
很多的配套.
我相信都會改變.
今天我們看到這裡.

$^{841}$和我們將來.
看到的應該不一樣.
我相信我們有更多.
傳呼音的機會.
但是否.
我們就是這樣自己衝出去.
還是在過程中.
打仗前真的要問清楚.
我們的元帥.
耶穌基督.
是怎樣打的.
他這樣說.
就跟了.
就這樣打了.
我很盼望未來的日子.
我們是有不同的.
我們的生命是有不一樣.
是活出榮耀的樣式.
是活出一個得勝的樣式.
讓上帝的心意滿足.
亦都是過程中.
我們會很興奮.
當我們聽上帝的聲音.
跟著他的教導去.
以上帝讓我們.
經歷他的得勝.
我們會很歡喜快樂.
我們和上帝一起.
歡喜快樂.
和我們教會的主耶穌基督一同歡喜快樂.
在這裡願神賜福給我們.
我們一同禱告.
天母上帝.
今天透過這個.
以色列民出埃及的這段經文.
你讓我們思考教會的身份.
讓我們再次去思考.
究竟這支軍隊.
是否真的有力.
是否真的榮耀.

$^{881}$好像上帝眼中所看的.
主耶穌求你.
開我們的眼睛.
讓我們看到現在.
我們這一刻的光景.
你讓我們看到.
我們和你的關係是怎樣.
我們是否真的.
去求問你的心意.
真的專注聽你的聲音.
又願意去跟從.
主我們整個群體.
有沒有真的讓你在我們裡面.
作王.
顯出你的大能.
我們一同去經歷你的得勝.
在此求你讓我們思考.
讓我們活出來.
也讓我們從今以後.
我們的生活不再一樣.
主幫助我們來到你的面前.
好好地向你禱告.
整個生命的態度.
不再一樣.
我們真正地將我們生命的主權.
教會的主權.
交出來.
歸還給你.
由你去掌管.
願主帶領我們.
奉耶穌基督的名字祈禱.
阿們.
\newpage



\section{路加福音 10:1-11-16-20}
\label{sec:pNmpdKyDl3E}
\textbf{2022.07.31 《宣教任務( 二) :熟讀這篇武林秘笈》 路加福音 10\_1-11,16-20 ( 漢語新譯本)  楊穎兒傳道}
\newline
\newline
連結: \href{https://youtube.com/watch?v=pNmpdKyDl3E}{\texttt{https://youtube.com/watch?v=pNmpdKyDl3E}} ~~~~ 語音日期: 2022-08-04
\newline
\newline
\hyperref[sec:rVzIB96PGyI]{\small{< < < PREV SERMON < < <}}
~
\hyperref[sec:index_chronic]{\small{[返順時目]}}
~
\hyperref[sec:index_scriptual]{\small{[返順卷目]}}
~
\hyperref[sec:MN9at1kI2W8]{\small{> > > NEXT SERMON > > >}}
\newline
\newline
路加福音 10:1-11-16-20
\newline
\begin{longtable}{cl}
\hline
\hline
章節 & 經文 (和合本修訂版)\\
\hline
10:1 & \begin{tabularx}{0.7\textwidth}{X} 這些事以後，主另外指定七十二個人，差遣他們兩個兩個地在他前面，往自己所要到的各城各地去。 \end{tabularx} \\ \\ \relax
10:2 & \begin{tabularx}{0.7\textwidth}{X} 他對他們說：「要收的莊稼多，做工的人少。所以，你們要求莊稼的主差遣做工的人出去收他的莊稼。 \end{tabularx} \\ \\ \relax
10:3 & \begin{tabularx}{0.7\textwidth}{X} 你們去吧！看！我差你們出去，如同羔羊進入狼群。 \end{tabularx} \\ \\ \relax
10:4 & \begin{tabularx}{0.7\textwidth}{X} 不要帶錢囊，不要帶行囊，不要帶鞋子；在路上也不要向人問安。 \end{tabularx} \\ \\ \relax
10:5 & \begin{tabularx}{0.7\textwidth}{X} 無論進哪一家，先要說：『願這一家平安。』 \end{tabularx} \\ \\ \relax
10:6 & \begin{tabularx}{0.7\textwidth}{X} 那裡若有當得平安的人，你們所求的平安就必臨到那家，不然，將歸還你們。 \end{tabularx} \\ \\ \relax
10:7 & \begin{tabularx}{0.7\textwidth}{X} 你們要住在那家，吃喝他們所供給的，因為工人得工錢是應當的；不要從這家搬到那家。 \end{tabularx} \\ \\ \relax
10:8 & \begin{tabularx}{0.7\textwidth}{X} 無論進哪一城，人若接待你們，給你們擺上甚麼食物，你們就吃甚麼。 \end{tabularx} \\ \\ \relax
10:9 & \begin{tabularx}{0.7\textwidth}{X} 要醫治那城裡的病人，對他們說：『神的國臨近你們了。』 \end{tabularx} \\ \\ \relax
10:10 & \begin{tabularx}{0.7\textwidth}{X} 無論進哪一城，人若不接待你們，你們就到大街上去，說： \end{tabularx} \\ \\ \relax
10:11 & \begin{tabularx}{0.7\textwidth}{X} 『就是你們城裡的塵土黏在我們的腳上，我們也當著你們擦去。但是，你們該知道神的國臨近了。』 \end{tabularx} \\ \\ \relax
10:12 & \begin{tabularx}{0.7\textwidth}{X} 我告訴你們，在那日子，所多瑪所受的，比那城還容易受呢！」 \end{tabularx} \\ \\ \relax
10:13 & \begin{tabularx}{0.7\textwidth}{X} 「哥拉汛哪，你有禍了！伯賽大啊，你有禍了！因為在你們中間所行的異能，若行在推羅、西頓，他們早已披麻蒙灰，坐在地上悔改了。 \end{tabularx} \\ \\ \relax
10:14 & \begin{tabularx}{0.7\textwidth}{X} 在審判的時候，推羅和西頓所受的，比你們還容易受呢！ \end{tabularx} \\ \\ \relax
10:15 & \begin{tabularx}{0.7\textwidth}{X} 迦百農啊，你以為要被舉到天上嗎？你要被推下陰間！」 \end{tabularx} \\ \\ \relax
10:16 & \begin{tabularx}{0.7\textwidth}{X} 耶穌又對門徒說：「聽從你們的就是聽從我；棄絕你們的就是棄絕我；棄絕我的就是棄絕差遣我來的那位。」 \end{tabularx} \\ \\ \relax
10:17 & \begin{tabularx}{0.7\textwidth}{X} 那七十二個人歡歡喜喜地回來，說：「主啊，因你的名，就是鬼也服了我們。」 \end{tabularx} \\ \\ \relax
10:18 & \begin{tabularx}{0.7\textwidth}{X} 耶穌對他們說：「我看見撒但從天上墜落，像閃電一樣。 \end{tabularx} \\ \\ \relax
10:19 & \begin{tabularx}{0.7\textwidth}{X} 我已經給你們權柄可以踐踏蛇和蠍子，又勝過仇敵一切的能力，絕沒有甚麼能害你們。 \end{tabularx} \\ \\ \relax
10:20 & \begin{tabularx}{0.7\textwidth}{X} 然而，不要因靈服了你們就歡喜，而要因你們的名記錄在天上歡喜。」 \end{tabularx} \\ \\
[1ex]
\hline
\hline
\end{longtable}
$^{1}$洪恩堂的弟兄姊妹平安.
平安.
多謝洪恩堂老牧師的安排.
讓我能夠回來洪恩堂.
在大家當中去分享神的話語.
我們先祈禱.
親愛的父啊.
來到你的殿中去聽你的話語.
求你真的將聖靈.
藏在每一個人的心裡面.
開他們的耳朵.
去聆聽你的話語.
亦預備孩子的口去宣講你的話語.
我們這樣的禱告.
奉主明求.
阿們.
先給大家看一幅圖畫.
可能大家都認得這個畫面.
這個就是粵語武俠片《如來神掌》.
當中由曹達華飾演龍劍飛.
當時是1964年.
有沒有人看過?.
有.
又有一幅畫.
當中有弟兄姊妹都看過.
就是2004年周星馳的電影《功夫》裡面.
都拿了一本武林秘笈出來.
就是《如來神掌》.
有沒有人看過?.
都看過.
在武俠小說裡面.
其實武林秘笈的設計.
是遵循著幾個特性的.
首先當然要厲害.
一出掌就厲害.
不可以每個人都懂得.
要很隱密.
它很奇妙地可以做到一些.
很奇妙的效果.
你一推出去.

$^{41}$掌就去到別人那裡.
亦都很靈.
而且它很正宗.
有門有派.
亦都要合乎道德.
還有這個武林秘笈.
其實是要傳承下去的.
在香港的電影世界裡面.
我們有《如來神掌》的武功要傳下去.
同樣在我們的信仰裡面.
宣教裡面.
都有一個秘笈是要傳承下去的.
剛才都讀過經文.
就是耶穌親自去傳授.
這個全福音宣教的武林秘笈.
給七十二個弟子.
關於當差派門徒去傳講福音的時候.
是門徒要跟隨和要注意的事情.
留意我今天的講題是.
宣教任務二.
熟讀這篇武林秘笈.
讓我們一起來看看這篇.
由耶穌親自去教授門徒.
信徒.
宣教秘笈有什麼值得我們反思的.
第一節.
「這些事已後」.
什麼事呢?.
我們要看看前文.
發生了什麼事呢?.
以至耶穌要傳授這篇武林秘笈給他的門徒.
我們很簡單地看回路加福音.
其實在路加福音第一至八章.
是講耶穌的出生.
呼召門徒一起在加里里海和加伯隆一帶.
去傳道行神蹟.
醫病,趕鬼和貧窮人傳福音.
叫「被擄得釋放,乞眼得自由」.
「受逼迫得自由」.
「乞眼看見,受逼迫得自由」.

$^{81}$到第九章.
主耶穌去差遣十二個門徒傳道.
那時候主耶穌早已講了.
第九節.
耶穌對十二個門徒說.
什麼都不用帶.
不用帶手杖,行囊,食物和銀錢.
什麼都不用帶.
你們無論去到哪一家.
就住在那裡.
直到離開那個地方.
無論誰不接待你們.
你們離開那個城的時候.
就撕掉那些塵土.
去指控他們的不信.
意思是.
其實你們不用帶東西去.
因為有人會接待你們.
你帶很多東西去的時候.
別人哪有機會接待你呢?.
要專心去傳道.
有人不接待你就算了.
因為第九章四十八節這樣說.
無論誰接待我.
就是接待猜我來的那個.
意思是.
誰不接待你.
就是不接待猜你來的那個.
但那時候門徒的反應很纏繞.
他不接待我那怎麼辦呢?.
門徒甚至去問耶穌.
第九章第五十三節這樣說.
當時有人接待他們.
而門徒雅各和約翰.
看到有人接待他們就問.
主啊.
你要我們去吩咐火從天上降下來.
燒了他們嗎?.
耶穌是不是被這些門徒氣死了?.
有人不接待耶穌.

$^{121}$不用燒別人的屋子.
耶穌的行程是要去耶路撒冷.
他要順服上帝的旨意.
他要上十字架.
去為世人贖罪.
拯救世人.
但那時候門徒很不像樣.
怎麼辦呢?.
當地人不接待耶穌.
不用燒了別人的屋子.
到了第十章.
即是今天的經文.
我再一次向七十二位門徒.
即是信徒.
都是說些相似的話.
第十章.
又是不用帶錢,袋子,鞋子.
一路上.
難道要問候任何人嗎?.
因為問候任何人.
別人又會說話.
而會分擾了你專心傳道的心.
所以甚麼都不用帶.
有時帶太多東西去宣教.
一去就被搶走.
一帶去.
別人又說你不合用.
所以你不帶東西去.
才有機會讓當地人幫你.
大家還記得那時候.
保羅去宣教.
他離開馬其頓.
遇到一個做生意的女人.
她叫做呂底亞.
他帶她順著住在他家.
後來又認識了一些官員.
這些官員又帶領當地人.
去幫助保羅.
另外.
耶穌亦吩咐他們.

$^{161}$不要隨便從一家搬到另一家.
意思是叫他們不用急.
不用趕.
慢慢做好一家才一家.
而這個被截.
其實是有一個指標的.
那聽從你們的.
就是聽從我.
那拒絕你們的.
就是拒絕我.
那拒絕我的.
就是拒絕差遣我來的那一位.
所以這個指標是說.
有兩個可能性.
就是你去到當地的時候.
全部都是這樣.
兩個可能性.
一是接待你.
一是不接待你.
所以不接待你的話.
那怎麼做呢?.
就說無論你去到哪座城.
人若不接待你.
你們就走到街上.
即使是你們城裡的塵土.
有黏在我們的腳上.
我們都要拆掉.
還給你.
雖然是這樣.
但你應該知道.
神的國近了.
其實他不接待你不要緊的.
我們拍一拍腳的塵.
但我們仍然要告訴他們.
神國近了.
你們要悔改.
遠離惡行.
記得嗎?.
無論誰接待我.
就是接待差遣我來的那一位.

$^{201}$不接待的就走.
去下一家.
亦都不用死纏爛打.
亦都不用燒了別人的房子.
耶穌對十二門徒是這樣說.
對七十二位門徒也是這樣說.
但不同的是.
原來到路加福音.
即是他寫的時候.
大概是主升天之後三十年.
當時那裡的信徒.
已經傳福音到希臘.
羅馬一帶.
他們當時面對的.
是一個社會充滿著不同種族的人類.
是一個多元的社會.
當時的社會已經不再只是猶太人.
有流散的猶太人.
有希臘人.
有羅馬人.
所以在這篇秘笈.
和之前給十二門徒的秘笈.
有什麼不同呢?.
原來是耶穌加了一些跨文化的貼士.
第八節.
無論你們進那座城.
人若接待你們.
給你們擺上什麼.
你們就吃什麼.
其實這是宣教必須要做的.
當你去到當地.
你要吃別人的東西.
別人才會覺得.
他親手做了一些東西給你吃.
他又接待你.
你又欣賞他的食物.
另外.
無論你們去到哪一家.
都要說平安.
願平安臨到這個家.

$^{241}$很簡單.
平安.
跟別人說一聲.
你們進屋也要叫人.
平安.
如果那裡有人配得平安.
你們所求的平安就會臨到那個人.
如果那個人不配得平安.
那個祝福就無效了.
其實如果看原文的意思.
其實不是無效.
原文是說那個祝福會反彈給我.
這樣很好.
我祝福你.
如果你不配得平安.
反彈給我.
這樣很好.
而猶太人當時的想法.
那個平安.
其實不單是出入平安.
他們有很圓滿.
和平.
和諧.
安全.
安穩.
一個發展得很好.
一個盛況.
甚至是一個救贖.
拯救.
一個很圓滿的平安.
就好像.
我記得我們.
一百年前.
當宣教士來到中國地方宣教的時候.
他們都要入窗隨俗.
他們都要穿上唐裝.
學兩句中文.
是不是.
另外就是.
你們住到哪裡.

$^{281}$就吃他們提供的食物.
因為工人得報酬是理所當然的.
那時候.
記不記得保羅去到哥林多的時候.
他遇見散居的猶太人.
阿基拉.
妻子伯基拉.
保羅就和他們一起織帳篷.
和他們一起做工.
去融入當地的文化.
當地的生活.
所以如果我們今天請到耶穌來.
再和我們講一下這個秘笈.
其實精要都是差不多的.
第九節.
要治好城裡的病人.
對他們說神的國近.
去到宣教的地方.
其實有很多問題.
甚至是我們在本地.
自己生活.
經驗不到的問題.
例如現在我服侍的.
一些難民群體.
他們都是一些.
從外國流散到香港的地方.
他們住的環境很小.
他們的房間甚至缺乏空氣.
tong 房都很窄.
所以經常說這裡痛那裡痛.
所以我覺得.
不是說我不是醫生.
我不可以治好他們很多的病.
但是至少我會告訴他們.
不要忘記耶穌基督.
怎樣拯救了我們.
耶穌基督怎樣在我們的生命.
顯了祂的大能.
所以我們所傳的是.
耶穌基督上帝兒子降臨在世上.

$^{321}$為眾人死.
這個如此反智的道理.
就讓世人知道.
讓世人遠離他們現在的罪惡.
而走向真理.
說完了.
這個孫教的武林秘笈.
它達不達到武林秘笈的特性?.
高超?.
厲害了 耶穌教的.
隱密?.
隱密嗎?.
隱密吧 藏在家裡.
在家裡福音.
不說大家也未必知道.
正宗?.
耶穌傳的正宗.
有沒有道德元素?.
也有的.
你們要和別人有禮貌.
平安.
不要隨便從一家搬到另一家.
打完齋不要和尚.
出去實戰比武.
不帶秘笈.
所以熟讀了這篇秘笈.
可以放在一旁.
有三方面值得我們反思.
第一點.
孫教官要更新.
第九章.
猜顯了十二個門徒.
第十章主要猜顯了七十二個門徒.
可能你會說.
這些數字怎麼解釋?.
對我來說.
這十二個門徒.
是指那時候.
舊約.
有十二支派代表著舊約.

$^{361}$七十二個門徒.
是指新約.
這些外邦信徒.
當我們稍為玩玩數字.
我們就會發現.
原來猜顯了八十四個.
其實就是門徒.
十二.
乘以七.
其實七是一個很圓滿的數字.
其實.
聖經是否說.
猜顯信徒的時候.
是全民皆兵?.
在主耶穌的心意裡.
今天在座的每一位.
都是天國的精兵.
都是能夠參與天國的.
孫教的事工.
很多時候.
我們的孫教官.
還留在一百年前.
孫教士千里迢迢.
搭半年船.
又運船浪來到新大陸.
然後又感染風土病.
又要學語言.
又要換上一件唐裝衣.
服侍了一段時間.
孫教士就死了.
沒錯.
今天教會仍然是要.
差派孫教士出去外地.
去將這個福音.
帶到美德之民裡面.
今天的美德之民.
大概集中在.
日本,中印地區.
東亞,中東.
中間這個地方.

$^{401}$當中的信徒.
有回教,佛教.
有些民間宗教.
這些人民.
都很需要福音.
在我服侍的地區.
我們都派物資的供應.
經常都會見到.
不同的.
來自不同地區的.
世界各地的人.
來取物資.
很多都是.
有其他信仰的.
透過這些服侍.
我能夠接觸到他們.
和他們分享信仰.
甚至是分享愛.
說一句平安.
其實大家已經是.
一個很好開始的交流.
與此同時.
我們也看看.
今天的處境.
今天我們面對全球化.
我們因為全球化.
已經很方便.
我們很容易去經商.
很容易去貿易.
很容易搬到另一個城市.
以哈比為例.
哈比是因著宗教的迫害.
來到一個陌生的地方.
今天我們也面對科技的發展.
很多年輕人.
在網上的世界.
在元宇宙的世界裡.
他們不只是玩遊戲機.
而是學校教授的東西.
也由這些世界裡.

$^{441}$灌輸給他們.
所以我們今天面對的.
不再是平面地域界限的宣教.
反而我們要想一個.
多元空間的宣教工作.
留意今天的經文.
其實沒有很指明地說到地點.
原文就是說.
這些門徒很主動地.
要去到城市裡.
宣教不再局限在.
地理上的理解.
我們在這個空間裡.
反而是我們去到哪裡.
就起來做.
如果今天主要在.
這些空間裡去預備人心.
我們在這個城市裡.
我們在做什麼的猜拳工作呢?.
對於西方的國家來說.
他們的社會比較.
接觸多一些難民的需要.
因為在中東地方.
一逃到西方國家比較近.
但在香港.
其實有很多社會大眾.
都不知道他的故事.
所以今天很開心.
哈比來到當中.
跟大家分享宗教迫害.
是多麼真實.
他們整個群體.
一起去受苦.
就像耶穌基督一樣.
與耶穌基督連結在一起.
在香港他們被貼了太多的標籤.
他們很多時候都會被標籤為.
「你來拿福利」.
但其實如果他們可以選擇.
他們也不會想離開.

$^{481}$那是他們的家.
讓我在散居的群體裡服侍.
在今天的世界裡.
宣教不再局限在去哪裡.
而是無論我們去到哪裡.
我們都要帶著宣教心.
而今天我們也看到.
仍然有很多人因為戰爭.
饑荒,迫害.
帶到不同的國家.
帶到我們的城市裡.
就像現在的疫情一樣.
他們很多人被困在城市裡.
所以今天我們要更新宣教觀.
本地跨文化的事工.
是必須要做的.
是宣教的一部分.
第二點.
「莊家多,工人少」.
當時耶穌指出了一個問題.
第二節.
他對門徒說.
「莊家多,工人少」.
「所以你們要求莊家的主」.
「趕緊派警察工人去收割他的莊家」.
我們從耶穌的話裡.
感受到那種緊迫的語氣.
第三節.
「你們去吧」.
「看,我差遣你們出去」.
「就像把羊羔送進狼群中」.
我們大多數人現在都不是以農為業.
所以我們可能不太理解這裡全部的意思.
對於任何一個田主來說.
在收割的時候.
有兩個重點是需要關注的.
一,有多少收成?.
二,有很多收成.
但是趕不趕得及收呢?.
假如一個田主的農作物都長得很好.

$^{521}$是時候要收割了.
但是收割不及的話.
農作物就會腐爛.
結果就是損失.
所以我們想像到.
經文說莊家多不夠工人去收割.
收割不及.
其實是田主很大的憂慮.
留意經文是在說.
去派工人.
不是在說請工人.
不是在說去裝備工人.
大家聽得出分別了嗎?.
工人是有的.
為什麼不出去呢?.
可能工人在磨鋤頭.
建大一點的狼瘡.
負責計劃一下.
待會要怎麼收割呢?.
或者看看旁邊有沒有田還可以耕.
要做神的工人.
又要等呼召.
又要讀神學.
又要裝備.
又要讀宣教課程.
宣教視野心.
又要參加警察會的奉獻團契.
做完了.
讀完了.
要先目睽睽.
學習如何目睽睽教會.
抗堂等等.
所以今天我們真的要求莊家的主.
趕緊派警察工人去收割莊家.
你們去吧.
我警察你們去.
今天香港警察全市工聯會.
在2019年做了一個調查.
發現當年宣教視的人數有624位.
是香港警察的宣教視.

$^{561}$當年新增的有58位.
這個近一點.
新增的有58位.
但退休回來有33位.
即是說那一年的增長只有25位.
香港有1300間教會.
25位即是2\%都不夠.
如果真的這麼難警察出去的話.
我們教會.
我們信徒要更加爭氣.
為宣教出一分力.
你記不記得我剛才說過.
84位門徒.
12個加72個.
全民皆兵.
外邦的信徒.
要接觸這個福音的棒.
繼續宣教的事工.
在這個空間.
任何時間我們都要有宣教的心.
但又有人說.
今時不同往日.
沒錯.
世界是不斷變化的.
第三節.
你們去吧.
我警察你們出去.
就像把羊羔送進狼群中.
把羊羔送進狼群.
這是主的吩咐.
也是秘笈的精髓所在.
容易做就不用主.
自己靠自己就可以.
反而越大挑戰.
逼不越大.
工人才更加要依靠主的幫助.
宣教的工作.
修葛的工作.
上帝比我們更緊張.
田也是祂的.

$^{601}$水也是祂的預備.
物資也是祂供應.
錢也是祂提供.
送你出去.
當然是上帝看著.
叫你去見證神的工作.
你都不去嗎?.
或者你會說.
有些猜出去的宣教士會死的.
那又怎樣呢?.
生命不是又回歸到上帝的愛裡嗎?.
第三點.
宣教士齊見證.
第十七節.
那七十二個人歡喜地回來.
說主啊.
因著你的名.
就連鬼也服從了我們.
耶穌對他們說.
我看見撒旦像閃電一樣從天上墜落.
作供的是主.
這班人見證了.
因著耶穌的名連鬼也服了他們.
但是大家有沒有留意.
這班人回來.
還沒有說主的作為.
他們只是說主的名.
主自己也要補一句.
我看見撒旦像閃電一樣從天上墜落.
今天打的這場仗是屬靈的戰爭.
我們其實是看不到撒旦像怎樣從天上墜落.
這不是我們空間能夠看到的事.
這場仗是主耶穌親自去打的.
我已經賜給你們權.
使你們能踐踏蛇和蠍子.
並勝過仇敵的一切能力.
絕對不會有甚麼能夠傷害你們.
神國道的工作包括醫治,驅魔,報道.
一直都是和宇宙的力量打交道.
目的是要戰勝這些邪惡的力量.

$^{641}$蛇和蠍子就是代表這些邪惡的力量.
得勝的根據不是來自那七十二個人.
不是來自那十二個人還是八十二個人的能力.
而是來自耶穌自己的能力.
能夠戰勝敵人的一切力量.
同時神國的工人是能夠得到被神保護的.
今天神就是差遣我們.
兩個兩個參與在他的事工.
我們有甚麼能力可以參與在神的事工.
我們去到也是做見證的.
但不要驕傲.
不要因為邪靈服從了你們的喜樂.
卻要因為你們的名字已經記錄在天上而喜樂.
不知道大家是否想得到這個喜樂.
其實我們真的沒有能力.
沒有甚麼恩賜.
可以參與在神的宣教事工上.
但神就是那麼愛我們.
邀請我們去加入他.
去參與他的事工.
然後我們就可以很驕傲地說.
上帝我們是倚靠你.
我們是倚靠你的引領.
無論你帶我們去哪裡.
我們都成為你的光.
無論你帶了誰進來我們教會.
我們都是喜樂的服侍.
讓我們一起禱告.
三一的上帝我們向你的禱告.
我們滿心讚美你.
你就是那位創天造地.
擁有全地權柄的主.
全地都充滿著你的大愛.
你是我們宣教的主.
我們盼望你的國度能夠臨在地上.
在你的心意裡面.
你們要我們做的就是憑著信心出去.
甚麼都不用帶.
帶著你的福音.
帶著你的基督的寫生.

$^{681}$這個事實.
我們相信你帶我們進去狼群中.
平平安安地帶我們回來的主.
聖靈求你內注在我們的生命裡面.
除去我們的膽怯.
教導我們各人.
去專心順從於你的帶領.
若似各人宣教的心.
讓你許可的話.
讓我們在洪水橋的地區.
繼續去忠心的服侍.
讓你繼續預備人心.
得聞福音.
我們這樣的祈禱.
我主耶穌基督明而求.
阿們.
\newpage



\section{啟示錄 4:1-11-20:11-15}
\label{sec:MN9at1kI2W8}
\textbf{2022.08.07 《世界的終局》 啟示錄 4\_1-11,20\_11-15 (和修版) 老耀雄牧師}
\newline
\newline
連結: \href{https://youtube.com/watch?v=MN9at1kI2W8}{\texttt{https://youtube.com/watch?v=MN9at1kI2W8}} ~~~~ 語音日期: 2022-08-09
\newline
\newline
\hyperref[sec:pNmpdKyDl3E]{\small{< < < PREV SERMON < < <}}
~
\hyperref[sec:index_chronic]{\small{[返順時目]}}
~
\hyperref[sec:index_scriptual]{\small{[返順卷目]}}
~
\hyperref[sec:lBPL13HHJ8A]{\small{> > > NEXT SERMON > > >}}
\newline
\newline
$^{1}$主持人:待會兒的弟兄姊妹平安.
在今年的一月至三月.
我跟大家分享了《歌羅西書》的講道系列.
到了六月.
我就以講道系列的形式.
跟大家分享《啟示錄》.
《啟示錄》講道的系列裡面.
我跟大家講過《以弗所教會》.
和《老弟家教會》的信息.
今天進入第三講.
也是最後一講.
天上的敬拜和世界的忠告.
基督教信仰的世界觀是有開始的.
內容記載在《創世記》.
神創造宇宙.
聖獸,山川河流,天地萬物.
一切的生物包括人類.
都是神所創造.
並賦予生命氣息.
神託管著萬有.
讓一切被造物.
享有神的供應和恩典.
神有藉著亞伯拉罕.
從一個人.
伸延到一個民族.
一個國家.
賜福給世界上的萬族萬國.
今天我們正正是在享受.
神所賜給我們的福氣.
基督教信仰的世界觀.
是有終結的.
內容記載在《啟示錄》.
世界的時間是有期限的.
生命是有限制的.
而且短暫.
同會扭壞.
地上的一切物質.
都沒有永恆價值.
最終都會變為無有.
而且在世界終結的時候.

$^{41}$是有審判.
每一個人都要對自己所做過的事.
所犯的罪.
負上責任和代價.
而神就是終極的審判者.
可以這麼說.
人從出生到死亡.
就是一個生命探索和成長的旅程.
一生之中遇到什麼事.
最終可以去到哪裡.
最後可以得到什麼.
一切的答案.
都要視乎每個人的經歷和際遇.
以及如何去回應.
我們有沒有想過.
我們的生命之旅.
可以為我們帶來什麼結局呢?.
是一個豐盛的人生.
還是一個虛空的旅程呢?.
今天我們從啟示錄兩段經文.
看看在世界終結的時候.
會發生什麼事.
讓我們今天的人生.
能夠作出一個.
即時的信仰反思和回應.
我們一起去祈禱.
因為你是創始成中的主.
你既是開始.
也是終結.
萬有都歸你.
萬有也因你而立而生.
我們的生活.
存留.
氣息.
一切都在乎你.
求主讓我們每一個.
都按著你的心意而活.
也為著我們所領受的恩典.
甘心樂意為主你作美好的見證.
我們祈禱.

$^{81}$奉靠主耶穌基督得勝名字祈求.
啟示錄的作者.
是仕途約翰.
也是主耶穌最愛的其中一個門徒.
約翰當時被羅馬政府.
流放到拔摩島.
但主耶穌仍然與他同在.
讓約翰領受從神而來的意象.
讓他看到將來在天上的敬拜.
以及世界的終局.
經文第一節這樣說.
「理想者裡來.
我要把此後必須發生的事指示你」.
耶穌基督啟示仕途約翰.
讓他有機會看到將來的事.
約翰也將自己所見到的.
記錄在啟示錄裡.
讓後世的信徒.
也有機會感受和經歷到.
神給他們領受的啟示.
在意象裡.
約翰看到天上有一個寶座.
而坐在寶座上的是誰呢?.
作者以碧玉,紅寶石來形容神的樣子.
事實上.
地上沒有任何文字.
可以描述到神的模樣.
因為神的形象是超越的.
因此約翰只能夠以地上.
最耀目璀璨的寶石來描述.
而圍在寶座四周的.
也只能夠以類比的彩虹,綠寶石來形容.
因此約翰所見到的.
其實是一個極其威榮,尊貴的景象.
坐在寶座上的是擁有極其尊貴的榮耀.
是超越地上一切所有.
亦表示坐在寶座上是超越一切的.
除了寶座之外.
最接近寶座的是有二十四個座位.
上面坐著二十四位長老.

$^{121}$身穿白衣.
頭上戴著金冠面.
二十四位長老.
有學者認為是表示.
舊約的十二支派.
以及新約的十二使徒.
二十四位長老是新舊約的代表.
亦代表著一切屬神的人.
是神所揀選的子民.
每一個被神揀選的信徒.
都包括其中.
他們身穿白衣.
代表著每一個都聖潔無瑕.
被神接納.
他們頭戴金冠面.
代表榮耀和勝利.
是德性的一群.
亦都是被神賞賜的一群.
在寶座前有閃電,聲音,雷轟發出.
代表有極大的權能.
從寶座發出.
神的七靈.
即是聖靈.
神透過聖靈施展他的大能.
寶座前有玻璃海.
能夠揭示世界一切.
以及終局的信息.
一切的奧秘.
都可以從玻璃海裡面知道.
想一想.
弟兄姊妹想一想.
如果我們身處其中.
我們會有甚麼感受?.
是否一種興邁不得的感覺?.
到時我深信.
我們只能夠.
只有一個動作.
就是俯伏在寶座前.
甚至連頭都不敢抬起來.
寶座的四位.

$^{161}$寶座周圍四邊有四個活物.
遍體前後都長滿了眼睛.
長滿了眼睛代表了他的透視能力和智慧.
四活物的描述.
和以西結書所提到的天使基路伯很相似.
不過以西結書所提到的天使.
只有四個翅膀.
反而和以賽亞書所描述的薩拉忽.
就相似了.
同樣有六個翅膀.
第一個活物好像獅子.
第二個像牛軸.
第三個像獵人獵.
第四個像飛鷹.
獅子是野獸之中最強大的.
牛軸是加速中最強壯的.
人是萬物之中最有靈性的.
而飛鷹就是鳥類之中最兇猛的.
四種最優秀的特點.
同時集結在這個四活物身上.
四活物就代表著整個創造界的一切生物.
包括天上的天使.
他們都在寶座之前.
晝夜不停地讚美神.
聖哉!聖哉!聖哉!.
主全能的神.
昔在,今在,以後永在.
我們有沒有想過.
晝夜不停地去讚美神?.
詩篇第一篇有說過.
他為喜愛耶和華的律法.
晝夜思想祂的律法.
這個人就是有福了.
晝夜思想神的律法.
以及晝夜讚美神.
都是神所喜悅的.
我們閱讀了上述的經文.
就有一個印象.
這是一個敬拜,讚美的景象.
天上的敬拜就是這樣的.

$^{201}$當二十四位長老俯伏敬拜的時候.
他們將戴在頭上的冠冕脫下.
放在寶座前.
然後才說出讚美的說話.
當他們身處在榮耀,尊貴,威嚴的主面前.
就算已經賞賜給他們的冠冕.
他們都願意甘心樂意地.
放在寶座前.
因為與尊貴的主相比.
個人頭上的冠冕算不得什麼.
弟兄姊妹.
使徒約翰所見到的.
是將來在天上的敬拜.
是眾信徒,眾天使.
以及眾被造物一起敬拜的情景.
整個情景都顯得莊重,威嚴,輕慢不得.
今天我們回到教會.
這個聖殿裡面敬拜.
有聖經學者甚至神學家都說.
這個就是將來天上敬拜的預演.
我們要問.
如果這個真是預演.
我們今天地上敬拜的預演.
以致我們將來在天上敬拜的時候.
可以實踐的時候.
我們要問自己.
我們有沒有一份對主耶穌尊崇的態度.
我們會不會遲到呢?.
我們會不會不專心呢?.
我們會不會不預備自己的心靈去敬拜呢?.
我們有沒有為了這一天.
這一天主日崇拜.
為主獻上最好的呢?.
什麼是最好的?.
在主耶穌聖經上的回應.
我們可以知道這個答案.
什麼叫做最好.
就是你要盡心,盡性,盡意,盡力去愛祂.
我們盡了沒有呢?.
我們盡了什麼呢?.

$^{241}$今天我們對敬拜有沒有另一種更深的感受.
就是崇拜的焦點.
不再是去批評那個領唱.
唱到走音.
也不是批評主席的公土很湧長.
也不是說講完的講到很沉悶.
因為敬拜讚美的焦點從來都不是領唱.
從來都不是主席.
也從來都不是講員.
崇拜不是要滿足我們的感覺.
以及自然.
而是要討那位坐在天上大寶座的主的喜悅.
我們要透過敬拜,讚美.
透過我們的真誠的心.
將榮耀歸還給上帝.
我們今天應該好好地反思.
我們有沒有這一種心態.
去參與我們每一個禮拜這個預言.
這個天上敬拜的預言.
這個敬拜是不是可有可無的?.
如果這個敬拜真的可有可無的話.
我很擔心.
你能不能夠出現在天上的敬拜的場景裡面.
將來神所賞賜給我們的.
戴在頭上的冠冕.
尚且要除下.
今天我們在地上還要爭什麼呢?.
我們還要跟別人比較什麼呢?.
所有的牧者,同工,屬靈領袖,部長,信徒.
就算我們在地上的事奉很成功.
給別人很多的稱讚.
但是在聖潔,威嚴,有權能的主面前.
每一個,每一個都需要謙卑自傑.
彼此合一服侍和守望.
因為我們都很盼望.
我們每一個雄恩家的肢體.
都能夠在天上敬拜中有份.
你想一想.
將來你在天上敬拜.
你看到旁邊那個.

$^{281}$卻不是今天你所見到的那個.
你會有點失望.
你會期望將來你在天上見到.
你坐在旁邊的那個.
就是現在你坐在你旁邊的那個.
如果是的話.
你可以為你旁邊那個做什麼?.
你說我很有信心.
我一定可以上去.
坐在你旁邊的那個能不能上去?.
我們不能夠只有一個人說.
我們可以為他做什麼?.
如果我們雄恩家的肢體.
都很想將來的敬拜.
我們整個雄恩家每一個肢體都參與其中.
都一定有份的.
你可以為他做什麼?.
主耶穌在給約翰領受的另一個意象.
就是世界的終局.
對非信徒來說.
他們不相信有審判.
所以非常輕視末日的審判.
但是對於信徒來說.
有些人都認為.
最後的審判是遙不可及的.
不過無論是主耶穌.
使徒保羅,彼得,約翰.
卻相當重視.
因為這一天.
彌臨的時候大而可畏.
沒有人可以逃彼得到.
經文說.
我又看見一個白色的大寶座.
和那坐在上面的.
天和地都從他面前逃避.
再也找不到他們的位置了.
天和地都從他面前逃避.
是什麼意思?.
有一派的學者.
認為天和地都要被神所毀滅.

$^{321}$因為要毀滅了之後.
才能夠有新天新地的彌臨.
不過也有另一派持相反的意見.
神那麼慈愛.
神是不會毀滅天地的.
而天和地都退隱了.
要使死人顯露.
接受審判.
不過無論是哪一種的解釋.
都表明了一件事.
就是地上所有的物質都會改變.
我們的房屋.
我們的財富.
一切所擁有的物質.
將再也看不見了.
如果今天我們以我們的財富.
以我們住在哪裡.
作為我們喜樂的來源的話.
到那一天.
我們就不會再喜樂了.
到那一天.
我們應該關心的.
不是我們的房屋.
以及財產.
而是我們要怎樣面對審判.
經文說.
「於是海交出其中的死人.
死亡和陰間也交出其中的死人.
去到照各人所行的受審判」.
有人相信.
人死如燈滅.
人死了之後一了百了.
了無牽掛.
沒有什麼需要付代價的.
沒有責任可言.
基督教信仰不同意這個觀點.
基督教信仰認為.
人要按著自己所犯過的罪.
要去接受審判.
來顯出神的公義.

$^{361}$因此人縱然已經死去.
但是神卻要求.
海要交出其中的死人.
死亡和陰間.
都要交出其中的死人.
人最大的恐懼就是死亡.
人在死亡面前顯得很渺小.
也顯得軟弱無力.
不過經文說.
死亡和陰間.
都服膺在主的主權之下.
神才是最終的勝利者.
因此死亡和陰間.
最終都要被神.
「瘟盡火狐」中.
人能夠勝過死亡和陰間的.
就只能夠倚靠主耶穌.
中國人有一句話.
叫做什麼呢?.
「殺人放火金腰帶」.
「修橋保路無屍骸」.
是形容做惡人卻有好的下場.
反而做好人卻沒有好的下場.
因此佛家就回應了一句.
「善有善報,惡有惡報」.
不事不報,時辰未到.
佛學講因緣.
它是同意人最終都會有報應的.
但是什麼時候呢?.
不知道,不過一定有.
今世無,便來世.
教會很確切地告訴我們.
報應的時間就是末日的審判.
無論人是否死去.
或是到時再生.
都會在審判的日子裡作一個了結.
而且是由神負責審判的過程.
人的一生就像一部紀錄片一樣.
每一個事件的過程.
每一個內心的動機.

$^{401}$都會呈現在神面前.
無法隱藏.
你可以隱瞞世人.
在世沒有人知道.
甚至你連自己都可以騙.
但你永遠隱瞞不了.
能夠透視人心的神.
在審判台前.
一切心思意念.
一切惡行都會呈現出來.
你如何在神面前辯解呢?.
你說我不知道.
不知道是否一種解釋.
就可以過得關呢?.
一旦審判的結局出來的時候.
你會怎樣面對?.
聖經又說.
聖經用了四個字.
來形容當時的情景.
聖經說了什麼呢?.
「哀哭切此」.
哀哭切此就表示要後悔,要補救.
都已經太遲了.
一切都不能夠還回.
哀哭切此.
弟兄姊妹.
不想讓自己陷在哀哭切此的境況之下.
應該怎樣做呢?.
應該盡快在神面前認罪悔改.
求主寬恕和赦免.
否則後悔已經太遲了.
經文說.
凡明字沒留在記載生命冊上的人.
就被冰進火壺裡.
這火壺就是第二次的死.
人在肉身上的死是無法避免的.
因為人的肉身是會扭壞的.
摧殘的.
而且也是人犯罪的代價.
不過這個代價仍然未完全抵消人所犯的罪.

$^{441}$保羅在《羅馬書》說過.
罪的功架乃是死.
死有雙重意義.
第一就是肉身的死亡.
第二就是靈性的死亡.
人對所犯的罪要付的終極代價是第二次的死.
也是永遠的靈性死亡.
是永恆的.
也是基督教信仰所說的.
在地獄裡的沉淪.
地獄未必是一個地方.
可以是一個象徵.
象徵著一個極度痛苦的狀態.
而人之所以極度痛苦.
就是沒有神在他身邊.
也沒有愛與關懷在他身邊.
極度痛苦是因為不能夠與主同在.
永遠與主隔離.
因此陷入一個極度痛苦.
絕望又沒有盼望的狀態.
無論是肉身還是靈性.
神都不再介入了.
這就是一個永死的狀態.
那麼人如何可以避免第二次的死亡呢?.
答案就是.
明智在生命冊上.
沒有人可以看到.
生命冊裡面是不是有自己的明智.
那麼又如何確保.
自己的明智在生命冊之內呢?.
主耶穌曾經差遣過十二個門徒.
前往各城替人醫病.
趕鬼和傳福音.
當他們回來的時候.
很興高采烈地說.
主啊!因你的名.
就是鬼也服了我們.
很開心.
因為奉耶穌的名.
鬼也服了我們.

$^{481}$不過耶穌如何回應他們呢?.
不要因名服了你們就喜歡.
要因你們的名.
記錄在天上歡喜.
可見主耶穌對門徒的名.
是否記在天上.
讓他們能否戰勝惡魔.
更加關心.
因為能夠有永恆價值.
比任何事物都更加重要.
今天很多信徒都高舉恩信稱義.
認為只要信了耶穌.
一切都穩妥了.
這種穩妥就可以上天堂.
得永生了.
這種說法太過粗疏.
將信仰廉價化.
恩信稱義的重點是想說明.
人的行為與得救無關.
免得人因著自己的行為而自我誇耀.
你做些什麼.
對你得救是沒有關係的.
人能夠得到救恩.
是因著對神的信心.
是神賜予的恩典.
不是人的行為.
而人對神的信心是需要持久的.
是一生之久.
《雅各書》也說過.
你信神只有一位.
你信得很好.
不過連鬼魔也相信.
不過鬼魔就很害怕.
連魔鬼也相信耶穌.
連魔鬼也認同耶穌的身份.
那我想問.
你的相信與魔鬼的相信.
分別在哪裡?.
魔鬼雖然相信主耶穌是神的兒子.
是三一的真神.

$^{521}$魔鬼沒有否認.
但是魔鬼卻抗拒神的主權.
他拒絕服侍神.
魔鬼拒絕神的吩咐.
也拒絕認罪悔改.
因為魔鬼要與神同等.
他要搶奪神的榮耀.
與神為敵.
人要相信主耶穌是生命的主.
拯救的主.
願意放下主權.
一生服侍主.
並且要真誠地在神面前認罪和悔改.
立志成為聖潔的人.
不再犯罪.
所言所行都期望神得到當得的榮耀.
這就是魔鬼的信.
與得救的人信心的分別.
得救與否不是信的內容.
不是因為你相信耶穌是神.
不是因為你相信耶穌是三一的神.
不僅僅是信的內容.
而是對神的主權.
身份行動的回應.
魔鬼是拒絕和抗拒.
而得救的人是接納和馴服.
在《馬太福音》.
主耶穌曾經說過.
不是每一個人稱呼我主呀主呀的人.
都能夠進天國.
他說唯有進行我天父旨意的人.
才能夠進去.
在那天必有許多人對我說.
主呀主呀.
我們不是奉你的名傳道.
奉你的名趕鬼.
奉你的名行許多義能嗎?.
我要向他們宣告.
我從來不認識你們.
你們這些作惡的人.

$^{561}$給我走開.
耶穌認為進天國的重點.
不是對主耶穌的稱呼.
或者曾經為主做過的事.
而是對神的吩咐.
是否願意一生遵守和踐行.
因此我們的命.
是否在生命冊內.
不是一次在報道會裡的決志信主.
不要認為一次得救.
永遠得救.
也不是單單相信主的能力和身份.
因為魔鬼也相信.
更不是口裡雖然稱呼主耶穌為主.
但對主的呼召或提醒毫不回應.
與我無關.
而是你能否按著主的心意.
在地上踐行信仰.
無論你遇到什麼困難.
你都讓生命連結在主的生命裡.
讓主耶穌介入在我們生命的每一個部分.
讓主的恩典在我們生命裡流露.
讓人可以在我們生命裡見到主耶穌.
使主得到當德的榮耀.
當我們願意一生忠心地跟隨祂.
按著主耶穌的心意而行的話.
我們的命志就應該在生命冊上有份.
弟兄姊妹.
當我們明白今天的崇拜.
就是將來天上敬拜的預演的時候.
我們有沒有一份興奮的心.
興奮.
我們正在做一個彩排.
將來敬拜的彩排.
我們有沒有盡力在每一次的敬拜裡.
做好自己的角色.
當領唱唱歌的時候.
只是他唱.
因為我們一起唱.
當主席祈禱的時候.

$^{601}$我們有沒有感受到祂的禱文.
然後同聲說「阿們」.
當我們明白我們是否得救.
關鍵是在命志需要在生命冊上的時候.
我們有沒有一份警醒的心.
警醒的心.
就是我們的屬靈生命.
以及我們和耶穌的關係.
我們的屬靈生命.
是否按著主耶穌的心意而行.
一切一切都在影響我們的命志.
是否記錄在生命冊上.
最後我引用我們宣道會創辦人.
孫迅博士的一段屬靈旅程.
作為我們這次講道的結束.
孫迅博士在《非我為主》一書裡提到.
在他的生命裡有三段.
改變他一生的屬靈經歷.
第一段經歷.
是在十五歲的時候.
他接受了主耶穌重新得救時的感受.
他感受到.
藉著我們的主耶穌基督.
得以與神和好.
與神和好是進入信仰的第一步.
也是關鍵的一步.
能夠與神和好.
也是與主耶穌建立親密關係的重要基礎.
不能夠與主有一個親密的關係.
信仰是沒有意思的.
鬼魔也信.
你所信的.
與魔鬼信的沒有什麼分別.
要分別的就是你與主耶穌的關係.
魔鬼與主耶穌的關係是敵對的.
信徒與主耶穌的關係.
是連結的.
第二段經歷.
是在他三十一歲的時候.
孫迅被按納牧師已經十年了.

$^{641}$他感到事奉很乏力.
當他尋求更深的屬靈指引的時候.
聖靈提醒他.
我只能在基督裡成聖.
然後孫迅就明白到.
我必須要先將自己奉獻給主.
然後才能夠成聖.
成聖至到完全.
乃是從神而來的聖潔.
並不是人靠自己自我改造而得到的.
聖潔的主.
是不會設立一個污穢的人.
作為他的僕人.
無論僕人有多少能力都好.
所有事奉都是徒然的.
你事奉能力多好都好.
如果你的生命是污穢的話.
你所做的一切都沒有意思.
神是不會預立的.
因此孫迅不斷地請求聖靈提醒.
讓他能夠知罪,認罪.
第三段經歷.
在他身體很軟弱的時候.
醫生也絕望地宣告沒有辦法了.
他就領受到神的旨意.
乃作我完全的救主.
不但救我的靈魂.
也醫治我的身體.
作我的大醫生.
所以當他得到神的醫治的時候.
他明白神的救恩是全避的.
甚麼是全避呢?.
這個全避的救恩包括了.
我們的靈,魂,體.
是一個全人的拯救.
不單單是醫治你的病.
是你的心靈也醫治了.
我們知道宣道會的信仰特點.
就是四重福音.
四重福音是說甚麼呢?.

$^{681}$基督是拯救的主.
這是他第一次經歷中領受到的.
第二個是基督是成聖的主.
就是宣信在他第二段經歷中領受到的.
第三個就是基督是醫治的主.
也是宣信在他第三段經歷中領受到的.
最後,基督是再來的王.
是宣信博士在他整個信仰旅程中領受到的.
四重福音的特色.
和宣信博士的熟能經歷息息相關.
甚至可以說.
宣信的熟能經歷.
引用出四重福音的信仰宣言.
我們洪恩家是宣道會的一份子.
所以我們信徒是很有需要地認識四重福音的核心意義.
所以每一個轉會的申請者.
他都要上兩堂.
第一堂就是四重福音.
第二堂就是宣道會的歷史.
當我們每一個都願意進行神的旨意.
當我們在地上踐行神的道.
活出信仰的時候.
我們的名字就很自然地會寄在生命冊上.
我們一起去祈禱.
你不單給我們有生命氣息.
你更加讓我們的生命與你有份.
可以享受你的豐盛.
我們要歌頌你.
要讚美你.
要尊你的名為「聖」.
求你讓每一個願意跟隨你的.
我們都能夠經歷到那種從軟弱到剛強的體驗.
我們一生都能夠被你的恩典所盛載.
過一個土里喜悅的人生.
我們能夠活出信仰.
為主打美好的仗.
我們祈禱仰望奉教主耶穌的名聲求.
阿們.
願主祝福大家.
謝謝大家.

\newpage



\section{約翰福音 15:1-8}
\label{sec:lBPL13HHJ8A}
\textbf{2022.08.14 《活出結果子的生命》 約翰福音 15\_1-8 (和修版) 曾敏姑娘}
\newline
\newline
連結: \href{https://youtube.com/watch?v=lBPL13HHJ8A}{\texttt{https://youtube.com/watch?v=lBPL13HHJ8A}} ~~~~ 語音日期: 2022-08-18
\newline
\newline
\hyperref[sec:MN9at1kI2W8]{\small{< < < PREV SERMON < < <}}
~
\hyperref[sec:index_chronic]{\small{[返順時目]}}
~
\hyperref[sec:index_scriptual]{\small{[返順卷目]}}
~
\hyperref[sec:J30380FlPik]{\small{> > > NEXT SERMON > > >}}
\newline
\newline
約翰福音 15:1-8
\newline
\begin{longtable}{cl}
\hline
\hline
章節 & 經文 (和合本修訂版)\\
\hline
15:1 & \begin{tabularx}{0.7\textwidth}{X} 「我就是真葡萄樹，我父是栽培的人。 \end{tabularx} \\ \\ \relax
15:2 & \begin{tabularx}{0.7\textwidth}{X} 凡屬我不結果子的枝子，他就剪掉；凡結果子的，他就修剪乾淨，使枝子結果子更多。 \end{tabularx} \\ \\ \relax
15:3 & \begin{tabularx}{0.7\textwidth}{X} 現在你們因我講給你們的道已經潔淨了。 \end{tabularx} \\ \\ \relax
15:4 & \begin{tabularx}{0.7\textwidth}{X} 你們要常在我裡面，我也常在你們裡面。枝子若不常在葡萄樹上，自己就不能結果子；你們若不常在我裡面，也是這樣。 \end{tabularx} \\ \\ \relax
15:5 & \begin{tabularx}{0.7\textwidth}{X} 我就是葡萄樹，你們是枝子。常在我裡面的，我也常在他裡面，這人就多結果子，因為離了我，你們就不能做甚麼。 \end{tabularx} \\ \\ \relax
15:6 & \begin{tabularx}{0.7\textwidth}{X} 人若不常在我裡面，就像枝子被丟在外面，枯乾了，人撿起來，扔進火裡燒了。 \end{tabularx} \\ \\ \relax
15:7 & \begin{tabularx}{0.7\textwidth}{X} 你們若常在我裡面，我的話也常在你們裡面，凡你們想要的，祈求，就給你們成全。 \end{tabularx} \\ \\ \relax
15:8 & \begin{tabularx}{0.7\textwidth}{X} 你們多結果子，我父就因此得榮耀，你們也就是我的門徒了。 \end{tabularx} \\ \\
[1ex]
\hline
\hline
\end{longtable}
$^{1}$我們開始時再一同禱告.
讓我們有這樣的福氣.
來到你的面前.
向你獻上嘴唇和心靈的祭.
也讓我們有福氣.
去聆聽你的話語.
讓我們有福氣.
去聆聽你的話語.
讓我們有福氣.
去聆聽你的話語.
讓我們有福氣.
去聆聽你的話語.
讓我們有福氣.
這一刻,求你再一次.
去監察我們的生命.
看看我們當中有沒有得罪你.
得罪人的地方.
如果有,請你赦免潔淨.
免得這些罪惡.
難了我們,去到你的根前.
聖靈啊.
求你今天在我們眾人心中.
工作.
讓我們聽到上帝的聲音.
讓我們就作出回應.
奉耶穌基督的名字.
祈禱.
阿們.
今天這段經文.
我相信是曾經有很多的講員.
講過的一段經文.
是非常寶貴.
很願意今天和大家分享的時候.
我們互相勉勵.
這段經文.
講真葡萄樹的.
這段經文.
講真葡萄樹的.
先看看.
真葡萄樹的隱喻.

$^{41}$其實究竟.
上下文大概.
在一個什麼位置呢?.
原來耶穌基督知道.
自己當時就要.
被出賣.
要被捉拿.
離開這個世界.
回到天父那裡.
祂為門徒洗腳.
祂預言.
自己將會被出賣.
祂也命令門徒.
彼此相愛.
還講了很多很多不同的教導.
今天這個.
隱喻是其中一個.
的教導.
是很重要.
在當中我們逐樣.
逐樣去看.
剛才講到.
這個是真葡萄樹.
的一個隱喻.
在舊約裡.
不只一次去看到.
葡萄樹.
很多時候葡萄樹.
是和以色列民掛勾的.
是用來形容.
以色列民.
上帝和以色列民立約.
所以以色列民.
是有一個很尊貴的身份.
他們是.
屬於聖潔的群體.
上帝選了他們.
所以他們覺得自己.
是很特別的.
選擇的意思就是.

$^{81}$好像差不多.
是唯我獨尊.
上帝首先選了我們.
我們和上帝有份的.
你說是多麼寶貴.
所以他們經常都很自豪.
對比那些外邦人.
總覺得他們.
不是被選擇的.
自己是被選擇的.
所以自己是不同的.
但很可惜舊約.
和以色列民.
很多時候不是一個.
很正面的描述.
是比較負面的.
甚至裡面.
帶有審判的意味.
我要稍稍和大家看看.
因為和我們今天的經文.
都有些相關.
這個字體稍為有些小.
但應該都可以看到.
以賽亞書五章一到七節.
那裡這樣說.
是一個葡萄園之歌.
他說.
我要為我親愛的唱歌.
我所愛的.
他的葡萄園之歌.
我親愛的有葡萄園.
在肥沃的山崗上面.
他刨滑原子.
清除石頭.
栽種上等的葡萄樹.
在園中蓋了一座樓.
又作出走池.
指望他結葡萄.
反倒結了野葡萄.
本身這個葡萄樹.

$^{121}$是上等的.
是很好的.
栽種的人是有個期望的.
期望他是結好果子的.
但結果出來.
野葡萄就不是好果子的意思.
所以是很負面的.
這裡當然是形容.
以色列民的生命.
是非常破爛.
是誤討上帝的喜悅.
所以在這裡.
經文接著.
在這裡.
發出審判的訊息.
其實經文本身.
很快就說.
究竟葡萄樹.
是指甚麼?.
就是指耶路撒冷的居民.
就是指猶大人.
在過程裡.
上帝說到.
我指望他是結葡萄的.
希望這個群體結出好的果子.
但為甚麼.
結果不是好的果子呢?.
我現在告訴你.
我要怎樣做.
我要切去那些泥巴.
令它被銷毀.
這是一個審判的訊息.
拆毀的圍牆.
令它被踐踏.
我要令它荒廢.
甚至收剪都不再收錢.
我不會除草.
不會令它長得更好.
任由荊棘擲泥在這裡生長.
我也在這裡.

$^{161}$叫到.
可以分庫蜜雲.
我要棄置.
這個葡萄樹.
不要它.
在這裡.
代表對以色列民.
很大的不喜悅.
一個很大的審判.
以色列民.
覺得自己.
很尊貴.
很幸福.
我被揀選.
但生命結不出好的果子.
上帝感到非常不喜悅.
要撇棄它們.
這個警告的訊息.
是非常嚴厲.
在這裡你也看到.
第七節有解釋.
耶和華的葡萄園.
就是以色列家.
上帝期望在當中.
見到有公平公義.
但結果卻是滿滿.
寫寫的罪惡.
流血.
有冤聲.
所以審判會臨到這一刻.
來到約翰福音.
這裡就很特別.
在這裡的舊約背景.
在上下文之下.
耶穌基督有這個教導.
關於真葡萄樹的隱喻.
想說什麼呢?.
耶穌說.
我就是真葡萄樹.
耶穌基督.

$^{201}$是真以色列.
不是以前的以色列.
不是現在的以色列.
是真以色列.
不是以前的以色列人.
他遠遠勝過.
舊約的以色列民.
他是上帝所揀選的.
不只是.
聖約的子民那麼簡單.
他是上帝的兒子.
上帝的國度.
很重要的主宰.
他擁有真正的生命.
他能夠結實累累.
耶穌基督是不一樣的.
他說這句話.
是很有豪華的.
他說這句話是很有豪氣的.
耶穌基督是真葡萄樹.
那麼誰是知子呢?.
信徒就是知子.
這個關係.
非常緊密.
葡萄樹本身.
是擁有生命.
當他和知子連結的時候.
他的生命.
會進入知子裡面.
他會藉著知子.
去結果子.
知子本身.
是沒有生命的.
他一定要.
和葡萄樹連結.
他才可以得到生命.
不單是得到生命.
而且可以.
結果子.
他不能夠獨立存在.

$^{241}$葡萄樹卻可以.
自己去存在.
我們再看.
這個非常.
索果累累.
很美的豐收時間的圖畫.
耶穌基督.
擁有真正的生命.
當他和信徒結連的時候.
他的生命.
進入信徒裡面.
信徒經歷他的同在.
信徒就可以.
結出果子.
耶穌基督也藉著.
信徒去結出果子.
去榮耀上帝.
這個關係.
這個互動.
在裡面.
這個門徒在主裡面.
是一種很寶貴的關係.
很重要的一件事.
必須強調的是.
對於葡萄樹與知子.
栽培的人.
是有期望的.
大家如果剛才用心.
去讀經.
大家告訴我.
栽培的人是誰?.
啊!.
一分鐘前讀的經文.
我忘記了.
栽培的人是誰?.
耶穌.
耶穌又是葡萄樹.
又是栽培的人.
哎!.
天父.

$^{281}$一分鐘前我們才讀經文.
這個栽培的人.
天父.
是有期望的.
期望知子結出果實.
很正常的.
如果大家.
去種一棵果樹.
你沒理由只是希望.
有些葉.
有些樹在你家.
是不是?.
你當然期望.
我種芒果樹,它結出芒果.
如果是龍眼樹.
它能夠結出龍眼.
葡萄樹更加.
因為葡萄樹.
如果不結果子,是沒有其他功用的.
為什麼呢?.
會這樣說呢?.
以西結書15章.
有一個很特別的描述.
人子.
葡萄樹比一切.
其他的樹.
就像樹林裡種樹木的樹枝.
有什麼長處呢?.
可不可以.
從其中取木料來造功呢?.
可不可以用來造釘呢?.
掛東西在那裡.
很好用的.
不出果子也可以造釘.
或者.
造一個木架.
掛東西,可不可以呢?.
你看這幅圖.
都知道是不可以的.
其實沒有什麼.

$^{321}$不結果子就沒有其他功用.
最重要.
就是要結果子.
否則就燒掉它.
做燃料.
燒掉它.
燒掉它的意思是.
不要.
簡單一點,要不就結果子.
要不就不要它.
所以栽種的人.
是有期望的.
天父去栽培.
葡萄樹.
枝子的時候,很簡單.
一定是期望結果子.
這個含義.
是什麼呢?.
在這裡後面.
含義是什麼呢?.
今天上帝對我和你.
有個期望.
我們信耶穌之後.
我們的生命與耶穌基督.
緊緊的結連.
不只是拿著護照.
回到天家.
就夠了.
夠了,我有護照就行了.
我的生命.
如常地繼續進行.
以前如何吃喝玩樂.
我就繼續.
我以前如何享受人生.
我現在繼續,我以前如何做人.
繼續,是沒有變化的.
我不會讓任何東西改變我的生命.
我的生活節奏.
我的喜好,我所有的東西.
我如何用時間和金錢.

$^{361}$完全不會改變.
信耶穌之後,一模一樣.
只是多了護照.
對不起,天父不是這樣想.
這個期望是必然的.
不是說可以.
你或者可以結果子.
或者,你可以也說結果子.
你想也說結果子.
不是.
或者.
或者,可以.
是一定要的.
弟兄姊妹,有沒有聽到這句話.
一定要結果子.
這是唯一的期望.
如果不結果子會怎樣.
你告訴我.
燒掉它,剪除它.
所以呢.
我每一次聽到.
你跟我說.
行了.
我相信耶穌已經很好.
你放過我吧.
我人生想繼續.
我以前如何生活,現在如何生活,你任由我吧.
我想以前這樣做.
現在也是這樣.
我不想有改變.
算了,你留下那些事奉人員.
你留下那些.
站出來的.
他們改變就好了.
你不要搞我.
對不起.
你和耶穌基督結連的那一刻.
你的生命不是屬於你的.
栽培的人是期望的.
是他在掌管.

$^{401}$他期望.
一定要結果子.
那什麼是果子呢.
這段話.
其實我們綜合了上下文.
去看.
究竟果子是什麼呢.
很空泛.
不知道結什麼呢.
這裡有兩個不同的範疇.
第一個範疇是我們生命的質素.
簡單來說.
第一件事.
很重要,上下文其中一件事.
特別強調.
在主裡面.
上上在主裡面.
什麼在裡面呢.
聽上帝話.
葡萄樹與枝子.
枝子是要依靠葡萄樹.
才能夠生存.
才能夠結果子.
所以這是一個很強烈的.
一個重塑的關係.
枝子不可以自己決定.
葡萄樹.
你現在聽我說.
我現在告訴你.
我想這樣生活.
我作為枝子就想這樣.
你聽不聽我說.
你跟著我.
對不起,沒有這件事.
你只能夠跟葡萄樹.
所以很重要的一件事.
是信服.
這樣就在主裡面.
不是一刻的信服.
是上上的信服.

$^{441}$在主裡面.
不是很空泛地信服.
是信服上帝的話語.
所以經文裡面.
也有提到.
如果上上的話語在我們裡面的話.
我們要去.
認識上帝的話語.
要去信服.
上帝的話語.
在他裡面.
這是其中一個質素.
是很重要的.
我和你要去結出來.
喜樂.
什麼喜樂.
你叫我開心.
我自然就要開心.
這個喜樂有點不同.
是經歷耶穌基督而有的喜樂.
我們信耶穌以後.
主在我們裡面.
耶穌基督在我們生命裡面.
我們經歷祂的同在.
在生命.
事奉.
各方面.
上上經歷耶穌.
有很大的喜樂.
但是這個根源.
最大的.
最起初的喜樂是什麼.
是我們經歷被救贖的喜樂.
耶穌基督救贖.
我和你的生命.
離開永遠的死亡.
就很開心.
我們.
那是首先.
應該有的喜樂.

$^{481}$我們有沒有.
這份喜樂呢.
我經常都.
強調救贖的恩典.
很寶貴.
但是你有沒有發覺.
很多時候我們會看這件事.
太平常.
一來救恩.
我們經常都提到.
好像很平常.
第二.
我們看不到救恩有什麼.
這麼寶貴.
因為看不到我們罪惡的可怕.
看不到我們自己的罪惡的心重.
看不到耶穌基督.
救贖我們這麼寶貴.
所以我們不覺得喜樂.
回到教會.
本來應該每個人都很雀躍.
每個人都被拯救的人.
自由的人,是不是很歡喜快樂.
不是的.
我們不覺得這些是喜樂.
我們要找一些很外在的東西來喜樂.
例如我們要有大型的活動.
玩得開心嗎.
盡不盡興.
有沒有禮物.
有沒有見到朋友.
開心嗎,聊得來嗎.
人生是不是每天都有不同的事.
安排得一般般.
這樣就叫很豐富的人生.
我們本末倒置.
最大的喜樂.
應該就是信耶穌那一刻.
我們得到了.
但不應該變成很丟淡.

$^{521}$要找一些新的刺激.
為什麼教會生活這麼悶.
真悶.
沒有刺激.
沒有新意.
你看外面人的生活多豐富.
多姿多彩,回到社區中心.
還好玩.
幾十八個班.
是不是這樣的.
基督徒的喜樂不是這樣的一回事.
我們重想.
究竟我們因為什麼喜樂.
我們是不是失落了.
甚至乎都不覺得.
信耶穌是一件這麼喜樂的事.
但這裡說的是.
經歷耶穌的喜樂.
是我們的生命果子.
基督徒是應該有這種氣質.
人家走到你那裡.
感受到.
弟兄,姐妹.
你這麼開心.
為什麼.
每天有人送禮物給你.
玩得開心.
有什麼活動,我也想去參加.
不,弟兄,姐妹告訴我.
信耶穌很開心.
你知道耶穌拯救我的生命.
很開心.
你能不能有這種表達.
不能的時候.
要想一想,為什麼.
永遠都在想.
教會和社區中心有什麼不同.
如果我們和社區中心沒有分別.
的時候.
是很可惜的.

$^{561}$是很可悲的.
最寶貴的.
耶穌基督在我們生命裡.
是沒有價值的,救恩沒有價值的.
那這個是不是真的要相信呢.
要問自己.
但這個生命的質素.
是要顯示出來.
那種喜樂,不止.
救贖的恩典.
上帝給我們每天都有數不盡的恩典.
哇,經歷耶穌的喜樂.
實在太大.
我們彼此相愛.
這個要有這個質素.
還是彼此.
相殺.
我毒你,你毒我.
我看你不順眼,你又看我不順眼.
遙遠見到你,我就想走開.
然後我們就不停.
在話題上說,這個這樣的.
那個這樣的.
我真是數之不盡,因為.
港人的是非,真是整本書.
要這麼厚.
這個不是.
我們的質素.
耶穌基督希望我們.
彼此相愛,你看回.
經文上下文.
是不停提這件事.
這個是我們要結出的.
果子.
這個第一個層面.
生命的質素,其實應該不止的.
如果你說.
我們結出.
那個屬靈的.
果子,你記得有不同的.

$^{601}$那個一個屬靈的果子.
有不同的質素,其實那些都應該.
是我們生命應該.
要結出來的,這個果子.
另外就是我們第二件事.
很重要,我們要成為美好的見證.
所以他去強調.
是要讓天賦得.
榮耀,我們的.
生命出來是一個美好的見證.
耶穌基督.
派我們出去都要做見證.
讓人看到.
天賦的榮耀.
我們的生命就是一個活的板.
我一出去.
究竟被人看到什麼呢?.
我最怕.
聽到的關於基督徒的評價.
就是.
哎呀,我不知道.
那個人信耶穌嗎?.
直到某一次我才發覺.
哇,他做這個.
無間道做了這麼久.
我終於發覺了他信耶穌.
哎呀,你覺不覺得有點奇怪呢?.
我從來都不覺得他像.
基督徒,這樣,這個是最失敗的.
人家在你生命看不到耶穌基督的.
那你怎樣見證他呢?.
當你開口.
跟人家說,你信耶穌.
是不是?我先跟你說說.
我們的救恩是怎樣的.
一說下去,說完之後.
人家最後一句,馬上就會叫你.
閉嘴.
你為什麼不覺得呢?.
你的生命呢.

$^{641}$真的這麼像基督徒.
你又說信耶穌很好.
為什麼這麼久我不覺得你改變了呢?.
你跟.
信耶穌之前都差不多.
哇,那怎麼說呢?.
後面還怎麼說下去呢?.
怎麼榮耀天父呢?.
其實這裡後面呢.
有很多年代的.
一些思考.
當你榮耀天父的時候呢.
其實你才更有能力.
去結福音的果子.
我很深信.
結福音的果子這裡是一個很.
豐富的內容.
裡面一定包含了福音的果子.
天父期望我和你的生命呢.
在信耶穌之後.
有質素的.
不一樣的.
要脫離過去的價值觀.
脫離被罪捆綁的生命.
是要有一個很.
好的質素,像祂的.
每樣東西都像祂的.
是榮耀的,是見證到祂的.
我們是可以帶人信耶穌的.
令到更多的資質.
很特別的,這個比喻.
令到更多的資質可以黏著葡萄樹.
這一生我們的生命是豐富有力的.
這個是天父的期望.
記住,不是說.
可能吧,你得了就結果子.
你想你就結果子.
不得也不要緊.
你不想就由得你.
不是,是一定要.

$^{681}$一定要.
好了,結果子和不結果子.
有什麼後果呢.
結果子的生命.
是會被栽培的人.
會有修展的.
經文說得.
非常清晰.
不是某某人.
才經歷他的修展.
十五章二節說.
凡屬我.
所有屬我的人.
無一例外.
凡屬我的.
凡結果子.
無一例外.
凡結果子.
就是那些在上帝眼中.
體會他覺得.
是真的可以繼續栽培.
的生命,結果子的.
凡所有.
所有結果子的.
無一例外,他就會修展.
乾淨.
這個修展.
就是一個管教.
我們的生命會結果子的時候.
上帝真的覺得.
這個生命.
是一個.
他會喜悅的生命.
但是他不會.
停在這裡,救了你.
結果子,我放過你.
他會修展.
要管教.
修展的時候,舒不舒服.
聖經裡說.

$^{721}$管教是怎樣的?不舒服的.
是不是?.
被人管教,怎會舒服?.
上帝管教.
我們有很多方式.
有時都哭了.
哎呀,我知錯了.
我知錯了.
一而再,再而三,不斷地管教.
直到.
知錯為止,直到改為止.
修展.
不暢快.
不開心.
有時很怕.
但是這是愛的修展.
真的證明.
上帝很愛我們.
這個生命.
可以要.
令我們做什麼?.
是要結出更多的果實.
讓我們的生命更豐盛.
更加可以見證上帝.
每一個.
都會這樣經歷,所以不用.
在這裡想,你放心.
我通常只聽旁邊的人.
經歷上帝管教.
是不會淋到我身上的.
如果想想.
經歷在別人身上,自己不會.
那你想想,上帝也特別愛那個人.
上帝如果愛我,一樣會管教我.
是每一個都會.
是嗎?.
這是屬於上帝.
上帝所愛的,他也會這樣做.
是要令我們結果子.
更多.

$^{761}$你也可以想想.
之前所經歷管教的.
情況是怎樣.
當時是怎樣呢?.
在這裡我們看到.
真的會修剪.
修剪的時候,真的.
結果實會更多.
會更加好.
在這裡.
凡結果子,修剪乾淨.
使枝子結果子更多.
不結果子的生命.
是怎樣呢?就是要面臨被剪除.
其實我們要面對.
上帝的審判.
上帝的一個很嚴厲的.
很嚴厲的.
對我們的結果.
就是剪除.
不要.
這個果子.
這個枝子為甚麼會這樣呢?.
黏著葡萄樹.
這麼久都沒有果實出.
沒有價值.
讓我們做那個.
園丁農夫也會.
哇!種那個葡萄樹.
真的沒有其他用途.
它的木不能用來.
種甚麼東西.
應該不能用葡萄葉.
來煲湯.
有些藥用的價值,如果有的話.
你事後來提醒我,可能我的見識淺薄.
沒有的時候,真的只有結果子.
但是那個枝子.
黏葡萄樹這麼久.
都沒有結果子.

$^{801}$不要啦!.
剪了它更好.
不要阻礙.
多些空間,那些好的枝子.
就生長得更好.
結更多果實.
也不要讓它浪費了位置.
所以會剪除.
所以它說.
凡屬我不結果子的枝子.
它就剪掉了.
人若不常在我裡.
就像枝子被雕在外面.
枯乾了.
人撿起來,瀕盡火雷.
燒了.
剪了它.
燒了它.
沒有用的.
上帝不要這些生命.
其實你想想.
正常一般來說.
枝子黏葡萄樹.
從樹之取生命.
它是會結果子的.
但是枝子.
一直純粹黏在那裡.
一直不結果子.
極有可能.
它本身也枯乾了.
本身也死了.
是沒有生命的.
今天枝子和葡萄樹結連.
就像我們信了耶穌一樣.
我們也叫.
屬上帝.
上帝也稱我們屬祂.
但是如果那個關係.
不是真正在它裡面的關係.
我們是不能進入.

$^{841}$耶穌基督的生命裡.
不能取得.
耶穌基督的生命.
不能取得耶穌基督的生命.
耶穌基督的生命.
不能進入我們裡面.
我們的生命.
就算你拿了護照.
回天家.
其實是一個.
死亡的生命.
這不是一個活著的.
基督徒的生命.
是沒有生命.
一個沒有生命的.
這樣的人.
上帝怎麼看我們.
想不想要我們.
這個絕對不是開玩笑的.
其實我們很多時候.
太過安恕了.
坐在這裡.
每個星期我們回來崇拜.
只是感到很舒服.
看看.
今天的領唱.
唱得好不好.
師班唱得怎麼樣.
講完有沒有菜.
整堂自我感覺.
是否良好.
御使者很快.
一個多鐘就完了.
回家與上帝生活.
好像和上帝沒有.
什麼關連在裡面.
如果是這樣.
生活的時候.
御使者過很多年,很多天.
我的生命.

$^{881}$又不會改變.
我完完全全沒有顯出.
上帝期望我有的質素.
我也不會結果子.
我不會.
去帶人去信耶穌.
不會榮耀天父.
一個這樣的生命.
大家怎麼看.
其實這個真是.
沒有了真正的生命.
沒有耶穌基督的生命.
在我們當中.
好像是否基督徒.
也不是很大分別.
但其實.
我們是基督徒.
基督徒也不是很大分別.
值得我們反思.
今天我們是一個.
什麼樣的光景.
我很想說.
很多時候.
我們回來可能會想.
因為傳道人會看著.
我們團契有些領袖會看著.
我們不要做壞事.
我們不要做不好的事.
但一出門口.
就不同了.
例如.
一出門.
可以想想.
是否買六合彩.
有誰知道呢?沒有人知道.
出去就如常地做自己.
好像以前未信耶穌時一樣.
我買六合彩.
我想怎樣就怎樣.
滿口粗言穢語.

$^{921}$一生氣的時候.
其實我跟外面說粗口的人.
沒有什麼分別.
你怎知道.
我出了門,難道你跟著我去?.
不是的.
出不出門.
上帝也與你同在.
不需要攝錄機.
上帝全部記錄在案.
這個生命完全是.
一個沒有質素的.
沒有任何.
完全沒有.
耶穌基督生命在我們裡面.
一個這樣的人.
上帝.
怎樣看我們?.
記住.
結果是被剪除.
被冰盡火雷.
上帝絕對不喜悅我們.
如果去到最後.
很重要的是.
我不想有這樣的生命.
我要活出一個結果的生命.
我要改變,那怎樣做?.
經文說.
我就是.
葡萄樹,你們是枝子.
上在我裡面的.
我也上在它裡面.
如果要活出一個結果子的生命.
很重要的就是.
上在耶穌基督裡面.
上在祂裡面.
就是.
經常親近主.
不是只拿護照.
而是親近主.

$^{961}$進入關係裡面.
你真的和耶穌基督有關係.
在裡面,你聽上帝的說話.
聽耶穌基督的聲音.
跟隨祂.
跟著祂走.
你知道耶穌喜歡甚麼.
你做祂所喜歡的.
是要這樣有關係的.
以致不是這樣的.
你有你生活的.
不是耶穌有耶穌的.
你是認識耶穌的.
耶穌也認識你.
就算我們好意說.
耶穌也認識.
但現在你也認識祂.
你跟著祂,關係是這樣.
經常靠著祂,黏著祂.
祂叫你怎樣走.
你跟著祂走.
不夠力的時候.
也是靠祂,捉著祂.
總之,這個枝子和葡萄樹的關係.
是緊緊相近.
是在對方的裡面.
我們是很經歷到上帝的同在.
耶穌的同在,有力的.
我們尚在主的裡面.
祂就會尚在人的裡面.
這個人就多結果子.
離了我.
你們就不能做甚麼.
更加重要.
除了尚在他裡面.
是上帝的說話很重要.
我的話也尚在你們裡面.
在十五章.
十五章的頭裡有說.
耶穌基督的道.

$^{1001}$令我們的生命變得潔淨.
神的話對我們有影響力.
令我們成為聖潔.
好像一盞燈.
照著我們.
走在哪裡,走在一個正確的方向.
離開黑暗,離開罪惡.
跟著祂.
走去對的地方.
耶穌基督的話語令我們成為潔淨.
今天上帝的話.
常常在我們裡面.
我們就跟著走.
你的生命整個人的質素.
就是耶穌基督想要的質素.
天父想要的質素.
一路成長.
很有上帝想要的氣質.
整個人是不同的.
除了跟從祂的話.
還有.
凡你們想要的.
祈求就給你們成全.
你說好.
原來葡萄樹和枝子這麼好.
但你跟著耶穌基督.
想怎樣祈禱,祂都答應你.
所以我要一大筆錢.
我跟上帝說.
你給我,你說的.
葡萄樹和枝子,凡你們想要的.
祈求就給你們成全.
我現在要一大筆錢.
我要一間大屋.
你不給我.
你應該會成全的.
按照《肉汗福音》說的.
這不是說這一點.
而是因為我們常常在主的裡面.
和主的關係親密.

$^{1041}$依靠,信服.
跟從.
上帝的說話.
也常常在我們裡面.
我們認識神.
按祂的心意生活.
這樣祈求是不一樣的.
這樣的祈求.
最終會按照.
上帝的心意祈求.
所以.
這樣的話,祈求就給你們.
成全.
這個是真正的結果子.
生命最後的質數.
祈禱不是按照我所想的.
而是我知道.
上帝想我怎樣.
我的祈禱就按照祂所想的祈禱.
上帝一定成全.
是這樣的意思.
你們多結果子的時候.
我父就因此得榮耀.
我們可以見證到.
我們天父的榮耀.
你們也就是我的門徒了.
門徒和信徒.
有什麼分別?.
我想問大家.
門徒和信徒有什麼分別?.
耶穌基督不是說.
你們就是信徒了.
你們是相信我的人.
你們是人.
讀經.
聽不到.
門徒和信徒有什麼分別?.
很大分別.
不只是單純的信.
我真的按照耶穌基督所教導的.

$^{1081}$去生活,活出來的.
這樣就是門徒.
耶穌基督要的是門徒.
不是單純的信徒.
是要門徒.
是榮耀天父.
能夠按他心意而生活.
知著他.
是能夠有質素.
是能夠結出福音的果子.
那個秘訣.
就是經常要親近主.
聽祂的話.
經常要看聖經.
跟從祂的話.
在裡面的意思.
永遠都有跟從的意思.
是很重要的.
然後祈求.
是按照上帝心意祈求.
是要祈求.
所以今天我們很重要的.
就是親近,依靠上帝.
看聖經.
按照上帝心意禱告.
這樣就會多結果子.
但是我想問一件事.
究竟我們今天有多看重.
這些東西的價值呢?.
如果讓你選擇.
你的錢是有限的.
你寧可.
吃一餐自助餐.
還是拿來.
買一本好一點的聖經.
讓自己看得舒服一點.
多一點看.
鼓勵自己讀經呢?.
我知道.
我曾經在《大英展望》當中生活.

$^{1121}$我知道大部分人選擇的是什麼呢?.
我也很不客氣地說.
吃自助餐,吃好東西.
不捨得.
拿去買聖經.
但你要吃一餐,你知道有多快樂嗎?.
是不是?.
但我想想.
你吃完那些,最後去了哪裡呢?.
就去了廁所那裡.
但是上帝的話語.
卻可以留在我們生命裡.
讓我們生命有力量.
我們自己不會從聖經支取力量.
是親近神支取力量.
一有出事的時候.
最厲害的就是打電話給弟兄姊妹.
最重要不是打給一個.
打給很多個.
打完一個又一個.
講述那件事是何等悲慘.
講完一次之後.
不行,還是解決不了.
第二,我再找下一個,那個熟悉一點.
講完一次之後,不行.
再打給第三個,第四個,第五個,第六個.
你打完這麼多個之後.
行不行?.
還不行.
你見到第七個,第八個,繼續講.
我告訴你那件事,你問我多麼悲慘.
我就想問.
我們信了耶穌.
耶穌的位置在哪裡呢?.
我們和一個不信的人有什麼分別呢?.
為什麼我們要和這麼多人講?.
因為我們必須承認一件事.
人是有限的.
他不能夠真的.
切實的,徹底的幫助我們生命.

$^{1161}$唯有耶穌基督.
我不是說不找弟兄姊妹.
弟兄姊妹是重要的.
教會的緣故,有團契.
各方面互相扶持.
但是最重要的耶穌基督.
不抓住他.
在那個時候,完全沒有力量.
就這樣找人.
如果那時候.
還要和他講.
聖經哪裡有力量呢?.
不行.
我完全不知道哪裡打哪裡.
有什麼意思?對我來說.
是沒有意思的.
那就很可惜了.
上帝的說話是很有力量的.
但是我們就錯過了.
這些寶貝.
真的錯過了.
所以我鼓勵大家再次想想.
自己的讀經生活是怎樣的.
我們的祈禱是為了什麼而祈禱的?.
我本來今天.
想做一件事,逐個逐個問.
我們當中信耶穌多久了?.
是嗎?.
誰信耶穌一年?.
五年?十年?.
我們怎樣去看祈禱呢?.
誰讀完聖經一次?.
我最不敢問的問題就是讀完聖經多少次?.
有沒有讀完整本聖經?.
我真的很害怕.
我又很尷尬,你又很尷尬.
那祈禱又很難問.
我們有沒有怎樣去祈禱呢?.
依靠上帝.
會不會我信耶穌十年呢?.

$^{1201}$其實我也不太懂得怎樣去祈禱.
為什麼我信耶穌十年也不太懂得祈禱呢?.
所以你說這麼多年.
和耶穌基督的關係是怎樣的?.
有生命嗎?.
我真的活著嗎?.
好好地問自己的問題.
我是不是活著呢?.
在主裡面活著,真的嗎?.
這麼多年也不是活著的.
在我們接下來.
我們在計劃十一周年.
堂慶的時候.
很想趁著今天的.
講道挑戰弟兄姊妹.
做一個結果子的生命.
過了這麼多年.
無論信主多少年也好.
過去怎樣也好,今天就要重新來過.
要做一個結果子的生命.
不要再死死地掉下來.
隨時拿刀一割.
它就會斷.
隨時被人拿去燒的一個這樣的生命.
我們的生命可以很有力.
榮耀天父.
今天洪恩嘉是上帝所愛的.
弟兄姊妹,你們是上帝所愛的.
你們可以活得更加精彩.
我們一起來祈禱.
藉著我們再一次重溫.
葡萄樹與枝子的隱喻.
讓我們再思考我們的生命.
究竟信主多年.
今天的光景如何.
我們是不是那些.
都要被剪除的生命.
要被瘟盡火裡的生命.
如果是的話,求你切切寬恕我們.
這麼多年.

$^{1241}$活著都好像死了一樣.
主求你讓我們認清自己的罪.
認罪悔改.
讓我們從哪裡跌倒.
從哪裡重新開始.
我們要親近你.
要在主你的裡面.
要去讀聖經.
好好禱告.
去跟從主.
從上帝你得著力量.
去黏著你.
生命去成長,去賜你.
能夠結福音的果子.
榮耀我們天上的父神.
主,求你賜福給我們洪恩嘉.
我知道你愛我們.
我知道你不會想放棄我們.
所以求你.
用這份血激勵我們的心.
奉耶穌基督的名字.
祈禱.
阿們.
\newpage



\section{加拉太書 5:13-26}
\label{sec:J30380FlPik}
\textbf{2022.08.21 《聖靈的果子》 加拉太書 5\_13-26 (和修版) 黎光偉牧師}
\newline
\newline
連結: \href{https://youtube.com/watch?v=J30380FlPik}{\texttt{https://youtube.com/watch?v=J30380FlPik}} ~~~~ 語音日期: 2022-08-22
\newline
\newline
\hyperref[sec:lBPL13HHJ8A]{\small{< < < PREV SERMON < < <}}
~
\hyperref[sec:index_chronic]{\small{[返順時目]}}
~
\hyperref[sec:index_scriptual]{\small{[返順卷目]}}
~
\hyperref[sec:D4j0WaSCINU]{\small{> > > NEXT SERMON > > >}}
\newline
\newline
加拉太書 5:13-26
\newline
\begin{longtable}{cl}
\hline
\hline
章節 & 經文 (和合本修訂版)\\
\hline
5:13 & \begin{tabularx}{0.7\textwidth}{X} 弟兄們，你們蒙召是要得自由；只是不可把這自由當作放縱情慾的機會，總要用愛心互相服侍。 \end{tabularx} \\ \\ \relax
5:14 & \begin{tabularx}{0.7\textwidth}{X} 因為全部律法都包括在「愛鄰如己」這一句話之內了。 \end{tabularx} \\ \\ \relax
5:15 & \begin{tabularx}{0.7\textwidth}{X} 你們要謹慎，你們若相咬相吞，恐怕要彼此消滅了。 \end{tabularx} \\ \\ \relax
5:16 & \begin{tabularx}{0.7\textwidth}{X} 我說，你們要順著聖靈而行，絕不可滿足肉體的情慾。 \end{tabularx} \\ \\ \relax
5:17 & \begin{tabularx}{0.7\textwidth}{X} 因為肉體的情慾和聖靈相爭，聖靈和肉體相爭，這兩個彼此敵對，使你們不能做所願意做的。 \end{tabularx} \\ \\ \relax
5:18 & \begin{tabularx}{0.7\textwidth}{X} 但你們若被聖靈引導，就不在律法之下。 \end{tabularx} \\ \\ \relax
5:19 & \begin{tabularx}{0.7\textwidth}{X} 情慾的事都是顯而易見的；就如淫亂、污穢、放蕩、 \end{tabularx} \\ \\ \relax
5:20 & \begin{tabularx}{0.7\textwidth}{X} 拜偶像、行邪術、仇恨、紛爭、忌恨、憤怒、自私、分派、結黨、 \end{tabularx} \\ \\ \relax
5:21 & \begin{tabularx}{0.7\textwidth}{X} 嫉妒、醉酒、荒宴等類。我從前告訴過你們，現在又告訴你們，做這樣事的人必不能承受神的國。 \end{tabularx} \\ \\ \relax
5:22 & \begin{tabularx}{0.7\textwidth}{X} 聖靈的果子就是仁愛、喜樂、和平、忍耐、恩慈、良善、信實、 \end{tabularx} \\ \\ \relax
5:23 & \begin{tabularx}{0.7\textwidth}{X} 溫柔、節制。這樣的事沒有律法禁止。 \end{tabularx} \\ \\ \relax
5:24 & \begin{tabularx}{0.7\textwidth}{X} 凡屬基督耶穌的人，是已經把肉體與肉體的邪情私慾同釘在十字架上了。 \end{tabularx} \\ \\ \relax
5:25 & \begin{tabularx}{0.7\textwidth}{X} 我們若靠著聖靈而活，也要靠著聖靈行事。 \end{tabularx} \\ \\ \relax
5:26 & \begin{tabularx}{0.7\textwidth}{X} 不要貪圖虛名，彼此惹氣，互相嫉妒。 \end{tabularx} \\ \\
[1ex]
\hline
\hline
\end{longtable}
$^{1}$各位紅印堂弟兄姊妹平安.
首先再一次多謝教會的邀請.
可以在這裡分享上帝的說話.
我收到的講題是聖靈的果子.
加拉太書是這樣說的.
聖靈的果子就是人愛,喜樂,和平,忍耐,恩慈,良善,信實,溫柔,節制.
這些事沒有律法可以禁止.
以上的經文或許對你來說已經是非常熟悉了.
當想起基督徒應該具備什麼樣的品格.
我們很自然就會想起這九條.
曾經何時我們會買聖靈的果子一式九款的書籤.
送給團友或組員.
去勉勵教會的肢體鼓勵他們要追求人愛,喜樂,和平.
不能否認基督徒當然要注重品格.
這九種的美德也是我們所嚮往的.
不過根據《尚文》下里.
我們會發現修煉這些品格並不是使徒保羅寫這封信的目的.
原因有兩個.
第一 保羅說這九種美德是聖靈所結的果.
也就是說這種美德的流露是結果.
而不是修煉的起點.
就好像讀書時代每一年結業禮一樣.
我們或許會為每位得獎的同學而鼓掌.
期待自己都可以成為台上領獎的人.
但是學年已經完結了.
我們再不能夠時光倒流改變這個結果.
所以保羅向當時的信徒羅列出這些美德.
並不是要列一張屬靈成長的清單.
告訴當時的信徒.
哪些地方已經很好.
哪些地方要努力一點.
你要學會人愛,喜樂.
可能你的忍耐要多一點.
並不是這樣的我們相信的目的.
第二 保羅說這些是聖靈所結的果.
也就是說這也不是我們可以修煉出來的.
而是聖靈修成正果.
不是我們修成正果.
所以保羅當時要求信徒關注的.
並不是要找出修煉品格的方法.

$^{41}$因為這是聖靈的工作.
信徒理應關注的是什麼呢?.
信徒理應關注的是.
當我們相信主耶穌.
上帝的靈已經和我們同在.
我們如何可以讓聖靈在我們身上發揮作用.
以至教會當我們在群體相處的裡面.
能夠流露出聖靈帶來的美德呢?.
弟兄姊妹.
當大家經歷了三年的抗疫.
當生活的節奏反覆被中斷又再度恢復的時候.
無論是身邊的朋友.
還是教會的弟兄姊妹.
或許都有這種同感.
覺得這幾年的變化.
自己的心態也改變了.
當體會到人生無常.
人生的時間並不能掌握在自己手上的時候.
我們發現我們更加珍惜和家人朋友相處的機會.
我們期望修補建立某些關係.
我們期望趁著還有時間.
可以做幾件造福人又能夠令自己無悔的事.
事實上作為基督徒.
建立愛的關係.
尋求做更有價值的事.
這也是主對我們的期望.
不過,如何著手開始呢?.
可能我們很快就會發現.
單靠我們個人的努力.
是無法完成的.
我們需要彼此.
我們需要同行的人.
這也是主耶穌設立教會的目的.
當然,保羅當時也清楚知道.
教會並不是完美的.
他同樣要面對各種挑戰和困難.
不過,保羅仍然憧憬.
當信徒認清自己的境況.
當我們懂得如何依靠上帝的時候.
聖靈就會帶來教會截然不同的局面.

$^{81}$信徒可以在人愛,喜樂,和平,忍耐,恩慈,良善,信實,溫柔,節制的底下.
不單止在這個氣氛底下.
能夠建立彼此相愛的關係.
更加能夠攜手完成上帝的託付.
活出不一樣的人生.
這樣,我們就要問.
如何才能依靠聖靈的幫助達到這個局面呢?.
在講解以先,讓我們同心禱告.
因為你愛我們.
你期待我們在屬靈上可以有成長.
你期望我們可以建立愛的關係.
在教會裡不單止給予彼此的支援.
更加我們可以攜手一起去做上帝要我們做的事.
讓我們履行使命.
完成上帝給我們的託付.
求你幫助我們以下的時間.
讓我們知道如何在你的話語裡得到力量.
我們能夠因著明白上帝的話語.
知道如何進行你.
祈禱奉耶穌基督名字.
阿們.
讓我們回到加拉太書.
了解保羅究竟在什麼情況下.
要提出聖靈所結的果子.
他期望當時的信徒學到什麼功課呢?.
讓我們先從加拉太書第五章第一節開始了解.
保羅一開始勉勵他們說.
基督釋放我們.
為使我們得自由.
所以要站穩了.
不要再被奴隸的欺騙脅制.
保羅強調.
得自由是主耶穌請求我們的目的.
不過這裡自由的意思並不是自由奔放.
想做什麼就做什麼.
保羅指的自由是不被箝制.
到底當時加拉太信徒被什麼所箝制呢?.
歷代的人雖然有不同的見解.
但至少我們可以排除兩個方面.
第一.

$^{121}$箝制他們的不是摩西的律法.
雖然保羅強調我們得救是上帝的恩典.
但他從來都沒有反對人守律法.
如果你看加拉太書第五章.
當你看到第十四節的時候.
你會發現保羅告訴當時的信徒.
總要用愛心互相服侍.
因為全部律法都包括在愛倫如舉這句話之內.
保羅並沒有否定律法的重要.
所以不是因為當時信徒守律法.
他們感覺被箝制失去了自由.
第二.
箝制他們的不是當時迫害他們的猶太宗教領袖和法利賽人.
因為從他們信耶穌開始.
就已經受到猶太宗教人士的反對.
保羅毋須在這裡重複說明.
這樣是否因為他們不再信耶穌.
他們改信猶太教.
以致他們不再自由呢?.
也不是.
因為在整卷的加拉太書信裡.
保羅似乎沒有提及過他們要放棄信仰.
我們有理由相信.
保羅認為箝制他們的不是教會以外的人.
而是教會裡面的自己人.
他們當中有部分是猶太人.
基本民族信仰.
他們自小就行割禮.
也有一些渴慕猶太文化的外邦人.
他們以不同的方式去接受割禮.
割禮是什麼呢?.
割禮就是當時他們按摩西律法.
將男性部分的包皮割去.
行割禮對教會群體不會造成什麼影響.
不過.
倘若這些信徒不但以自己的割禮為傲.
覺得我們做割禮.
我們遵行摩西律法.
我們感覺高一些.
並且他們覺得不只是高一點這麼簡單.

$^{161}$他們還感覺自己優於其他信徒.
他們甚至譴責那些未受割禮的肢體.
認為他們不夠敬虔.
將他們排他於外.
結果.
有些從來未受割禮的外邦信徒.
他們內心就很掙扎.
到底我不願意去接受割禮.
是不是我們愛主不夠深呢?.
當我們有了這個背景.
我們就更加能夠去理解.
保羅為什麼要這樣去回應.
保羅強調.
我保羅告訴你們.
五章二節.
你們若受割禮.
基督就對你們無益了.
這個無益並不是割禮本身.
而是那些主張行割禮的人.
背後那種優越感.
他們並不是為了尋求律法的真義而行割禮.
他們只是純粹靠表面的割禮禮儀.
來提升自己在群體內的地位.
這些人喜歡以自己的道德標準來批評別人.
用這種方法來操控其他人.
勉強別人去跟從他們的做法.
保羅指出.
這些人不能夠在基督身上得到任何的好處.
因為真正提升我們的是聖靈的工作.
而不是割禮的儀式.
正正因為我們承認我們的無能.
才能夠靠上帝的恩典成為義人.
就算我們遵守整套的律法.
也不能夠為我們的屬靈生命帶來半點的提升.
相反.
保羅說.
這些要靠割禮稱義的人.
其實是與基督隔絕.
他們並不是提升.
相反.

$^{201}$他們是在恩典裡向下墜落了.
最後.
保羅勸勉當時的信徒.
真正能夠使到屬靈生命成長的.
其實是信心.
他說.
我們是靠著聖靈.
憑著信心.
才能達到我們之後所盼望的意義.
而更重要的是.
怎樣才能夠表現出我們的信心呢?.
保羅說.
唯獨使人發出人愛的信心才有功效.
第六節.
什麼叫做發出人愛的信心呢?.
好像很複雜似的.
《新漢語》的譯本有一個更接近原文的翻譯.
是這麼說的.
只有用愛表現出來的信心.
才有用處.
真正的信心不是勇於去割.
而是勇於去愛.
放下屬靈的優越感.
用愛了解弟兄姊妹的需要和難處.
嘗試與其他人同行.
信得過當我們這樣做的時候.
上帝會改變我們.
這種信心.
才能帶來屬靈的益處.
我們才有機會成長.
弟兄姊妹.
或許我們都曾經.
不自覺地陷入被人優勝的試探中.
這種優越感.
可以是來自成就.
有時我們會以.
是否成功來判斷我們是否做了好的見證.
變相我們否定了那些過著平凡日子.
甚至他們曾經遭遇挫折的人.
其實他們的堅持.

$^{241}$都是美好的見證.
有時我們會強調自己信仰的年資.
認為我們因為信主信得久.
我們在教會最久.
所以我們才是最了解教會的人.
我們不自覺地將後來加入的信徒排他於外.
有時我們會不自覺地.
將我們自己化成為教會的批判者.
以高的角度去批判教會內外的人和事.
而將自己置身事外.
怎樣才能突破這個處境.
可以放下這種優越感呢.
原來我們可以從珍惜開始.
怎樣珍惜呢.
或者甚麼是珍惜呢.
最近我有機會看到一段電台節目的報導.
節目訪問了近期一齣電影的女主角毛順君.
和她一起參演的有一個新晉偶像姜濤.
在訪問內.
毛順君提到.
有人問過我.
如果拿姜濤和當日的張國榮相比.
你會怎樣比呢.
她這樣回答.
她說.
我覺得失去的東西永遠都是最好的.
但今日我們更加要珍惜的是我們的年輕人.
甚麼是珍惜呢.
珍惜不是單單去懷緬過去寶貴的回憶.
正是因為過去我們在主裡面經歷過快樂的回憶.
以致我們更加珍惜今天的弟兄姊妹.
我們期望守望他們.
我們期望他們可以在教會裡面.
可以擁有同樣美好的經歷.
我們不能顧及教會所有人的感受.
但至少我們先學習.
不要將所有焦點都放在自己的身上.
體諒尊重身邊肢體的需要.
今日在教會裡面.
我們正正是需要這種珍惜.

$^{281}$或許來到這裡.
可能會有人有這樣的想法.
既然我們都是蒙恩的罪人.
在主裡面沒有誰特別優勝.
也不要批評別人.
那我們就不如甚麼都不要做.
從今以後我們就繼續教會的生活.
你就繼續你事奉的盲鹿.
總之大家以後就不要多管閒事.
做好自己就算了.
我肯定這不是保羅期望的局面.
他不單要信徒放下自我的優越感.
更加要當時的信徒走多一步.
甘心實踐愛.
不單止珍惜我們眼前人.
更加進一步知道怎樣去愛這些眼前人.
我們回到經文裡面.
保羅繼續說.
弟兄們 你們夢照是要得自由.
只是不可把握自由.
當作放縱情慾的機會.
總要用愛心互相服侍.
到底這裡的放縱情慾是甚麼意思呢?.
放縱情慾原來的意思是慾體.
這個解釋似乎不能讓我們更理解經文過中的意思.
讓我們用另一個方法.
我們先看看到底保羅所說的放縱情慾.
是相對甚麼來講的呢?.
經文告訴我們.
當保羅說不可把這自由當作放縱情慾的機會之後.
接下來他講甚麼呢?.
保羅接下來不是說.
所以你們應該要約束自己.
不可以放縱情慾.
所以你們要約束自己.
消除你們的慾念.
保羅沒有這樣說.
相反 他說你們不可以把這個自由當作放縱情慾的機會.
然後他就說我們應該要做甚麼呢?.
我們應該要愛心服侍.

$^{321}$所以這裡的放縱情慾並不是指奸淫擄掠 縱情走色.
而是信徒.
這個放縱情慾是信徒不再顧及身邊的人.
從此只為自己而活.
純粹按著自己的喜好去建立自己喜愛的生活方式.
英文聖經The Message 譯本.
將這一句翻譯得很全神.
中文大概的意思是甚麼呢?.
大概的意思就是要確保你不再以自由為藉口.
為所欲為 以致破壞你的自由.
你應該要將你的自由用在愛裡彼此服侍.
這就是自由增長的方式.
當保羅批評部分的迦拉太信徒.
藉著國禮來提升自己之後.
或許保羅也擔心當時會不會有部分人有這個想法呢?.
既然我們在主的恩典下.
做甚麼也沒有甚麼特別意義.
不如我們就甚麼也不做吧.
在教會裡躺平就算了.
總之大家從此互相尊重.
不要再騷擾對方就行了.
結果保羅作了進一步的補充.
指出在教會躺平.
從此只按自己的想法來過信仰生活.
這不是屬靈成長的出路.
保羅提醒當時的信徒.
屬靈的生命如果要再進一步成長.
我們仍然有律法需要學習和進行.
保羅怎樣說呢?.
因為全部律法都包括在愛倫如己這句說話之內.
你們要謹慎.
如果你們相咬相吞.
恐怕要彼此消滅了.
加拿大書第十四第十五節.
可能你會問.
大家各自為政又怎會相咬相吞呢?.
試想像一下.
當教會之體的生活變得各顧各的時候.
可能大家就會慢慢發現.
教會很多事情都不是稱心滿意的.

$^{361}$為什麼教會可以任由小孩走來走去.
我自己安安樂樂地去崇拜.
為什麼有些人搞著我們呢?.
為什麼我們要遷就初信的朋友呢?.
為什麼外賢一定要等那些遲下班的信徒.
我們才可以開飯呢?.
時間越長.
矛盾和衝突就會慢慢浮面.
有人出言指責.
有人置之不理.
即使教會開始的時候多麼有愛.
也會慢慢被信徒這種只顧自己所吞噬.
保羅說.
信徒如果執意過自我為中心的生活.
這種自以為是.
除了帶來淫亂.
污衊.
荒蕩.
拜偶像.
行邪術之外.
還會帶來仇恨.
紛爭.
忌恨.
憤怒.
自私.
分派.
結黨.
和嫉妒.
這個可以說是歷代很多教會衰敗的原因.
保羅叮囑他們要以此為鑒.
弟兄姊妹.
昔日保羅領受主耶穌的吩咐.
傳福音給外邦人.
我們想像下.
其實他大可以按本子辦事.
到各地將福音傳開.
然後就自由自在過自己的日子.
不過.
歷史告訴我們.
保羅到處傳福音後.

$^{401}$他不單一次.
甚至第二,三次.
回到他到過的城市.
去兼顧眾信徒.
去幫助他們.
讓他們在信仰上更能著地.
從保羅的書信你會發現.
其實他大部分時間.
是處理教會成長的問題.
保羅運用上帝給他的自由.
選擇服侍教會眾人.
保羅在歌羅西書中怎樣形容自己.
他說我是成為教會的僕役.
他知道靠他一個人.
不可能關顧所有人.
所以他又呼籲其他的信徒.
加入他的行列.
放下自己的優越感.
一同用愛心服侍其他人.
同樣.
今天我們單靠我們幾個人.
是不可能為信仰群體大力改變.
但願在疫情稍緩的底下.
當教會逐步恢復之後.
我們可以重新出發.
我們用甘心彼此相顧.
去表明我們對主耶穌的信心.
由始至終.
人愛喜樂和平並不是憑努力訓練出來的.
而是當我們甘心樂意照顧彼此的時候.
當我們願意去做先顧別人那個的時候.
聖靈就會幫助我們.
讓我們結出美好的關係.
透過我們的嘗試.
讓教會流露出人愛喜樂和平.
忍耐恩慈良善信實溫柔節制.
阿們.
\newpage



\section{提摩太前書 6:1-19}
\label{sec:D4j0WaSCINU}
\textbf{2022.08.28 《假道理與真財富》 提摩太前書 6\_1-19 (和修版) 曾裔貴牧師}
\newline
\newline
連結: \href{https://youtube.com/watch?v=D4j0WaSCINU}{\texttt{https://youtube.com/watch?v=D4j0WaSCINU}} ~~~~ 語音日期: 2022-09-01
\newline
\newline
\hyperref[sec:J30380FlPik]{\small{< < < PREV SERMON < < <}}
~
\hyperref[sec:index_chronic]{\small{[返順時目]}}
~
\hyperref[sec:index_scriptual]{\small{[返順卷目]}}
~
\hyperref[sec:qTaSFAQFmRY]{\small{> > > NEXT SERMON > > >}}
\newline
\newline
提摩太前書 6:1-19
\newline
\begin{longtable}{cl}
\hline
\hline
章節 & 經文 (和合本修訂版)\\
\hline
6:1 & \begin{tabularx}{0.7\textwidth}{X} 凡負軛作奴隸的，要認為自己的主人配受各樣的尊敬，免得神的名和教導被人褻瀆。 \end{tabularx} \\ \\ \relax
6:2 & \begin{tabularx}{0.7\textwidth}{X} 奴隸若有信主的主人，不可因他是主內弟兄就輕看他們，更要越發服侍他們，因為得到服侍的益處的正是信徒，是蒙愛的人。 \end{tabularx} \\ \\ \relax
6:3 & \begin{tabularx}{0.7\textwidth}{X} 若有人傳別的教義，不符合我們主耶穌基督純正的話語與合乎敬虔的教導， \end{tabularx} \\ \\ \relax
6:4 & \begin{tabularx}{0.7\textwidth}{X} 他是自高自大，一無所知，專好爭辯，擅於舌戰，因而生出嫉妒、紛爭、毀謗、惡意猜疑， \end{tabularx} \\ \\ \relax
6:5 & \begin{tabularx}{0.7\textwidth}{X} 和心術不正與喪失真理的人不停地爭吵，以敬虔為得利的門路。 \end{tabularx} \\ \\ \relax
6:6 & \begin{tabularx}{0.7\textwidth}{X} 其實，敬虔加上知足就是大利。 \end{tabularx} \\ \\ \relax
6:7 & \begin{tabularx}{0.7\textwidth}{X} 因為我們沒有帶甚麼到世上來，也不能帶甚麼去； \end{tabularx} \\ \\ \relax
6:8 & \begin{tabularx}{0.7\textwidth}{X} 只要有衣有食，我們就該知足。 \end{tabularx} \\ \\ \relax
6:9 & \begin{tabularx}{0.7\textwidth}{X} 但那些想要發財的人就陷在誘惑、羅網和許多無知有害的慾望中，使人沉淪，以致敗壞和滅亡。 \end{tabularx} \\ \\ \relax
6:10 & \begin{tabularx}{0.7\textwidth}{X} 貪財是萬惡之根。有人因貪戀錢財而背離信仰，用許多愁苦把自己刺透了。 \end{tabularx} \\ \\ \relax
6:11 & \begin{tabularx}{0.7\textwidth}{X} 但你這屬神的人哪，要逃避這些事；要追求公義、敬虔、信心、愛心、忍耐、溫柔。 \end{tabularx} \\ \\ \relax
6:12 & \begin{tabularx}{0.7\textwidth}{X} 你要為信仰打那美好的仗；要持定永生，你為此被召，也已經在許多見證人面前作了那美好的見證。 \end{tabularx} \\ \\ \relax
6:13 & \begin{tabularx}{0.7\textwidth}{X} 我在那賜生命給萬物的神面前，並在向本丟‧彼拉多作過那美好見證的基督耶穌面前囑咐你： \end{tabularx} \\ \\ \relax
6:14 & \begin{tabularx}{0.7\textwidth}{X} 要守這命令，毫不玷污，無可指責，直到我們的主耶穌基督顯現。 \end{tabularx} \\ \\ \relax
6:15 & \begin{tabularx}{0.7\textwidth}{X} 到了適當的時候都要顯明出來：他是那可稱頌、獨一的權能者，萬王之王，萬主之主， \end{tabularx} \\ \\ \relax
6:16 & \begin{tabularx}{0.7\textwidth}{X} 就是那獨一不死、住在人不能靠近的光裡，是人未曾看見，也是不能看見的。願尊貴和永遠的權能都歸給他。阿們！ \end{tabularx} \\ \\ \relax
6:17 & \begin{tabularx}{0.7\textwidth}{X} 至於那些今世富足的人，你要囑咐他們不要自高，也不要倚賴靠不住的錢財；要倚靠那厚賜萬物給我們享受的神。 \end{tabularx} \\ \\ \relax
6:18 & \begin{tabularx}{0.7\textwidth}{X} 又要囑咐他們行善，在好事上富足，甘心施捨，樂意分享， \end{tabularx} \\ \\ \relax
6:19 & \begin{tabularx}{0.7\textwidth}{X} 為自己積存財富，而為將來打美好的根基，好使他們能把握那真正的生命。 \end{tabularx} \\ \\
[1ex]
\hline
\hline
\end{longtable}
$^{1}$各位弟兄姊妹早安.
早安.
再一次能夠代表務會.
過來和弟兄姊妹.
在神的話語上.
彼此一齊去追求.
更加能夠看到大家.
在崇拜裡面.
那份專心投入.
我們先一個簡單的祈禱.
多謝主.
我們能夠真的活在這個世界.
是很實在的.
我們每天的生活.
都離不開柴米油鹽.
但是我們知道.
我們屬主的兒女.
我們每一個信徒.
其實我們有永恆的天家.
是光明.
是神的居所.
是我們每一個信徒.
將來都可以永遠在那裡.
更加有我們親愛的.
救贖主耶穌.
和我們一起.
在主裡面的.
無論是親戚朋友.
和主裡面的弟兄姊妹.
將來我們都會在天家.
再相聚.
所以幫助我們在地上.
我們能夠在主的真道上.
去追求.
去長進.
今早我們事立.
等候在主的面前.
求你的聖靈.
向我們說話.
叫我們明白.

$^{41}$我們這樣祈禱.
祈求你的聖靈.
我們這樣祈禱.
奉耶穌基督的名.
阿們.
有少許抱歉.
因為待會一完港道.
我就要趕回舞堂.
因為有外來的港員.
因為如果太長時間.
我就好像不太方便.
因為我要招呼他.
待會一起用飯.
所以要抱歉.
我一完就立刻走.
這一段聖經其實是由第一節.
你可以去到第十九節.
其實講了四類人.
第一類當然是講第一節的奴隸.
當時在羅馬世界裡面的奴隸.
奴隸應該怎樣來生活呢?.
怎樣來在信了主之後.
表達他們內心的敬虔呢?.
去到第十七節.
講在那個時代.
保羅就叫提摩泰.
要勸一些有錢的人.
有能力運用他的資源的人.
有能力去累積財富的人.
怎樣來生活.
表達他們.
雖然有錢.
但是敬虔的有錢呢?.
這樣就很簡單了.
就是要甘於的去施捨.
樂意的去供給人.
使得自己.
這些有錢人.
能夠為自己.
積好了根基.

$^{81}$預備的是將來的永恆.
就是說.
原來錢不單止不是邪惡.
錢是為這些有錢人.
可以積成美好的.
自己的那個淑齡根基.
當然在這個分享的過程.
就表達到內心的敬虔.
第三類人.
就是第三到第五節.
有些傳假道理的人.
可能要讀讀第三節.
因為都頗為重要的.
三到第五節.
如果有人傳別的教義.
不符合主耶穌基督純正的話語.
和合乎敬虔的教導.
敬虔是穿梭在整張聖經的.
這些人都講敬虔.
但是你看看他們這個敬虔.
到底是他們內心裡面的態度.
那個本質.
因為認識主耶穌之後的內心的改變.
還是敬虔原來是一個手段.
為了達到的一個目的就是.
要騙錢.
成為自己那個自私的享用.
留意到第四節.
他說他是傳假道的人.
他是自高自大.
一無所知.
又自大又什麼都不知道.
這樣的人.
如果在你旁邊就很慘了.
又自大.
什麼都不會聽別人說.
但是原來自己很多事情都不懂.
不懂相信是在講淑齡方面.
對神的認識.
他們就會有一種自我保護的機制.

$^{121}$就會出現了.
尊好爭辯.
善於竊賤.
因而生出.
很自然的.
就是嫉妒.
紛爭.
毀謗.
惡意的猜疑.
和心術不正的.
與喪失真理的人不停的爭吵.
果然真的是一個手段.
以敬虔為得利的門路.
第六節之後.
保羅就教提摩泰.
其實錢在這個世界應該怎麼用呢?.
第六到第十節.
就是錢財的正確觀念.
不過留意第四類人.
第十一節.
但是你提摩泰.
你是什麼人?.
你是屬神的人.
你應該怎麼去表達.
你在神面前的這個慰焚.
這個呼召.
這個榜樣.
就是第十一節.
這些真的是聖經很重要的金句.
特別是對傳道人.
永遠都適用.
要逃避上文所說的這些事情.
要追求的是公義.
敬虔.
信心.
愛心.
忍耐.
溫柔.
全部是內心裡面.
因為主耶穌的救恩.

$^{161}$而帶來的生命改變.
對一個信徒.
尤其是傳道者.
更加要保羅很仔細地去描述.
敬虔的生命.
是包含什麼特質呢?.
有公義的.
也有敬虔.
有對神,對人的信心.
有愛心.
有對一些事情的忍耐,包容.
表達出一份內心很真實的一種溫柔.
溫柔不是那些音聲細氣.
溫柔是對人的等待.
不會是一種硬梆梆的.
我回想我初出來傳道.
講道很強.
有些長輩說不要那麼強.
你做到的事.
別人未必馬上做到.
應該要溫柔.
但不懂得.
但過了很多年.
二十多年了.
你就會很明白.
原來很多事情是要等待.
要用一個溫柔的心靈.
去勸導弟兄姊妹.
四類人.
尤其是在第三到第五節.
保羅說在教會裡.
必須要提防這樣的人.
因為他們不單止自己.
要得到那個財利.
所以就會用很多的手段.
更可怕,更嚴重的.
就是那個手段就是敬虔.
以敬虔為得到利益的門路.
但接著講的其他三種人.
就很不同了.

$^{201}$就是真正的敬虔是怎樣表達呢?.
很寶貴的地方就是第一節.
奴隸.
保羅沒有在他傳道的人生.
試圖將整個羅馬帝國.
一些維繫社會的階層.
他要因為信了耶穌.
就要立即推翻這些不公平.
對人不健康的壓制.
保羅沒有試圖這樣做.
反而聖經新約給我們的教導.
反而能夠維繫到.
那個社會的穩定之餘.
更加凸顯到.
原來不單止是改變那個制度這麼簡單.
立即可能用一種革命.
推翻這種不公平的制度.
為什麼你要看其他人.
是比你低下的奴隸呢?.
保羅沒有試圖這樣做.
你想想教的那些奴隸.
你信了主耶穌.
你在神眼中.
你就是很尊貴的弟兄姊妹.
不過在你的工作裡面.
你信了主的老闆.
同樣也是很尊貴的弟兄.
所以你就要用真誠.
用你的勤奮.
用你的內心改變了的品格.
去對付他.
對付就是恭敬他.
用愛心尊重他.
我在舞堂也說過一個例子.
其實學得不好.
端午節之前.
我媽媽就和我的姐姐.
包了很多種.
過百種的.
逐家逐戶派.

$^{241}$差不多有一天晚堂結束了.
就和我兒子一起.
開車去媽媽家裡.
去拿我家的種.
訂了多少種.
接著媽媽的菲傭就下來.
拿幾袋來.
順便派給我弟弟妹妹.
當時又晚又趕.
第二天又趕到.
就做錯事了.
大聲地喝菲傭.
快點放下吧.
很兇.
接著.
她很害怕.
不關我的事.
媽媽叫我給你.
我也不知道什麼事.
開車了.
我兒子的乘客就說.
曾牧師.
剛才你的說話不對.
你對菲傭這樣的態度.
也不關她的事.
馬上要說對不起.
接著再和菲傭說對不起.
弟兄姊妹.
保羅教導奴隸.
他說.
凡是在負額作奴隸的.
負額可以是因為社會問題.
可以是你的城邦.
被羅馬人吞併之後.
你就成為階下囚.
你就成為奴隸.
或者你的父母不知道什麼原因.
將你賣掉.
就成為奴隸.
就是說你一出生.

$^{281}$可能就已經是奴隸了.
凡是負額作奴隸的.
要認為自己的主人.
配受各樣的尊敬.
免得神的名聲和神的道理.
被人褻瀆.
一個信了主的奴隸.
重點是信了主帶來的身份.
在神面前改變了.
成為平等的弟兄姊妹.
而且在神眼中是尊貴的身份.
這是很重要的.
保羅就強調.
這個身份要活出來.
就是將你內心的敬虔.
仍然是一個奴隸.
穿著的衣服.
做的工作.
都是一個奴隸的工作.
但是內心裡面敬虔的品格.
就流露出來了.
透過哪裡去反映呢?.
服侍主人.
而且是你的弟兄.
那種美麗.
不需要用刀刀槍槍.
就可以某程度.
將這個不公平的對人的制度.
不知不覺間就打破了.
就有一種美.
在其中來流露.
所以就不要變成神的名字.
神的道理.
就被那些外邦人褻瀆.
所以很重的說話.
第二節就說奴隸.
如果有信主的主人.
不可以因為他是主裡面的弟兄.
就輕看他們.
更加越發地要服侍他們.

$^{321}$因為得到服侍益處的信徒.
是蒙愛的.
讓他感覺是蒙愛.
讓他得到服侍的益處.
讓這個主人看到自己家裡的奴隸.
是信主的.
所以大家回家.
你快點帶你的菲傭信主.
她會有很大的改變.
她會反過來.
會恭敬地對待你.
因為她信了主.
她的生命改變.
當然你加一點工資也更加好.
不過真的.
保羅告訴我們.
提摩太你要明白.
你有一個很重的身份.
以及職責.
就是要教導敬虔.
教導人去認識耶穌.
教導敬虔的生命.
所以到了第十七節.
第二類人.
至於那些今世富足的人.
當時有一些很權貴的階層.
彼得前書說那些人.
白晝宴樂.
方宴醉酒.
白晝宴樂.
你有沒有試過.
早上十點.
你在大酒樓裡面開五十席.
沒有人這麼做的.
一定是晚上.
下班後.
通常星期六晚才會這樣.
但是當時的貴族.
是白晝宴樂.
不過信了主耶穌的有錢人.

$^{361}$很不同了.
那些資源不是惡的.
資源當然是金錢.
因為整個上下文是在說錢.
在這些假的師父.
假教師的面前.
成為了一種追求的目標.
但是這些不需要追求.
本身已經有的有錢人.
第十七節.
至於這些今世富足的人.
你就要祝福他們.
不可以自高.
不可以依賴靠不住的錢財.
要依靠厚賜萬物.
給我們享受的神.
要回到那個源頭.
為什麼你能夠賺到錢呢?.
為什麼你能夠有餘暇.
去享用這些錢呢?.
為什麼你能夠運用這些金錢呢?.
你想想如果這樣的時候.
他能夠明白自己獲得的.
那就很不同了.
他就可以將他的金錢.
去幫助更多的人.
所以他的內心.
如果是一個敬虔的有錢人.
他的祈禱就是想賺更多的錢.
然後服侍更多的人.
這樣就真的是一個.
真真正正聖經所說的貨如輪轉.
真真正正的可以幫助更多人.
你留意第十八節.
不是令自己更加的自我中心.
更加的自高.
第十八節就是祝福他們要行善.
要在好的事情上付足.
甘心的施捨.
樂意的分享.

$^{401}$施捨和分享.
是一個敬虔的付足人的內心.
早一兩個月.
跟一些的牧者聊天.
劉恩德牧師.
他未做牧師之前.
就是安儲中學的校長.
很多年了.
聊天的時候.
他說他們學校.
有時候會請一些色囚.
坐過牢的人.
來做校工.
他說很多時候.
他們做著做著.
有時候都會有自己的人.
因為他們吸毒.
就會不做了.
有時候又會跟同事關係不好.
他說不是.
偶然間這麼多年裡.
試過有一兩個是非常好的.
因為信了主.
生命改變.
真的會帶來那個學校很大的祝福.
我也認識有香港連鎖的零售的老闆.
他說他們有時候都會請一些妓女.
但沒有人知道.
只有幾個高層知道.
讓他們有正當的工作.
在零售裡面做.
原來有錢人.
可以有這麼大的社會資源.
能力去幫助真正有需要的人.
所以這個樂意的分享.
和甘心樂意的施捨.
是一個有錢人.
要在神面前上上.
去求主憐憫.
為何保羅特別提到至高呢?.

$^{441}$有錢就是一個有錢人的神.
錢就是有錢人的神.
很容易掛勾.
但如果錢不是有錢人的神.
如果做物主.
耶穌基督真是你內心的神.
那有錢人就有很大的力量.
可以去甘心的施捨.
樂意的去分享.
就會帶來當時羅馬世界的一些信徒.
會有很大的祝福.
你看看保羅.
奴隸階層又可以教導他們.
有錢的階層.
保羅又可以叫提摩太去教導他們.
所以今天的講道題目.
假的道理會帶來很大的破壞.
所以教會很需要有敬虔的傳道者.
將純正的教導真理.
來教導弟兄姊妹.
所以第十一到第十六節.
就相當重要了.
第三類人就是在神眼中被選照的傳道者.
講聖經.
但理者屬神的人.
你要追求的是公義,敬虔,信心,愛心.
忍耐,溫柔.
你要為信仰打美好的仗.
要持定永生.
你是為這件事來被照的.
也在許多的見證人面前.
作了美好的見證.
保羅就在這裡再一次祝福提摩太.
我在那次生命級萬物的神面前.
並再向本雕彼拉多.
作過美好的見證的基督耶穌面前.
祝福你.
在神的面前.
這一位是創造萬物的神.
在耶穌的面前.

$^{481}$這一位曾經在他最後的晚上.
被拉多的公堂.
來受審的耶穌面前.
來祝福他什麼呢?.
要守這個命令.
這個命令就是傳永生的真道.
可以改變人.
讓人的生命得到那個敬虔.
守這個命令.
毫不玷污.
無可指責.
直至到主耶穌的顯現.
相信我們不能夠否定.
耶穌隨時會回來.
如果大家有留意新聞.
你就知道現在英國.
甚至全世界的旱災.
有科學家說是五百年.
五百年就是馬丁路德的時間.
五百年的旱災.
現在歐洲的大火沒有停止熄滅.
長江已經差不多有些地方枯乾.
萊茵河快要斷流.
下了一個很嚴重的.
好像近乎世界末日的光景.
但當然這些不可以作為一個印證.
主耶穌就是今晚回來.
但這是聖經新約不斷告訴我們.
使徒保羅來到緬泥提摩泰.
已經過了二千年了.
主耶穌還沒有回來.
但是保羅很肯定.
主耶穌一定會回來.
祂會顯現.
當然更快的是我們每一個人.
離開這個世界之後.
我們就進入冥行.
去見到我們的救贖主.
所以第十五節說.
到了適當的時候.

$^{521}$都要顯明出來.
祂是那一位可稱重獨有權能者.
萬王之王 萬主之主.
就是那獨一不死.
住在人不能夠靠近的光裡面.
是人未曾看見.
也是不能看見的.
好像將一個人的卡片.
不斷逐個銜頭讀出來.
祂是哪一家哪一家的會長.
祂是哪家公司.
祂是哪一位博士學位.
好像逐樣地讀出來.
不過這是主耶穌.
不知道你的袋子裡.
有沒有主耶穌的卡片.
這個卡片逐樣地說.
不過全部是在榮耀裡面.
永恆裡面的身份.
就是我們的主耶穌.
對提摩泰有什麼鼓勵作用呢?.
保羅不說他自己.
當然在其他地方他有說.
譬如提摩泰後書.
他有說自己.
勉勵這一位提摩泰.
提摩泰.
你不要以為我是一個正人君子.
基督耶穌降世.
說要拯救罪人.
在罪人當中.
其實你的師父我是一個罪魁.
不過然而我蒙了憐憫.
傳道者.
特別是一個內心比較柔弱的提摩泰.
很需要神的憐憫.
他能夠慢慢地.
自己可以獨當一面.
能夠慢慢地.
好像掛著P牌一樣駕車.

$^{561}$慢慢地自己有信心.
自己可以在保羅不在的時候.
他可以真真正正肩負起.
這個全陽永生的責任.
所以保羅來面對提摩泰.
你不斷地思想主耶穌.
他如何在彼拉多面前.
默默無聲地去受痛苦.
他現在是怎樣的身份.
時時要戴著卡片.
看著主耶穌原來是這樣的身份.
我現在事奉的就是他.
我就不怕了.
這個需要大家慢慢地學習.
當然第三到第五節.
是第四類.
就是聖經最不想出現的人.
也就是說提摩泰聽.
教會也可能會有這樣的人.
提摩泰很需要大眼識人.
有些人的敬虔原來是一個手段.
不是內心真正的態度.
如果有人傳別的教義.
不符合我們主耶穌基督純正的話語.
和合乎敬虔的教導.
這些人就會自高自大.
一無所知.
爭辯.
竊戰.
就會產生嫉妒.
紛爭.
然後就會毀謗.
就會猜疑.
然後就會不停地爭吵.
因為他們是以敬虔為得利的門路.
社會上有這樣的人一點也不奇怪.
在2009年的時候.
香港也有一宗比較大的新聞.
有一位姓卓的女孩.
31歲而已.

$^{601}$是一個恆生銀行的一個經理.
很出色的.
她說當時月入的月薪.
已經十幾萬一個月.
她的權限可以隨時調動客戶的資產.
當時有好幾個城中的明星.
邱淑貞.
馬榮成.
很多的資產她可以調動.
但是因為她自己投資失利.
擅自將這些資產調動.
最終過了一段時間.
因為通常明星都不理會.
信任這些經理.
調動的錢接近差不多一億.
那些明星的損失很大.
後來就知道了.
就報警.
這個女孩.
判了80個月監.
她很痛苦.
講述她自己的過程.
原來一直就是那個貪.
感恩啊.
在監裡面.
她信了耶穌.
沒什麼事做.
讀讀神學.
就真的很有心將來出來傳道.
當然現在我不知道她的情況如何.
這麼奇妙的經歷.
是很難的.
但是求主幫助我們要小心.
我們很容易就會有.
對於世界.
對於錢的試探.
會帶給我們.
保羅講的都很重的.
心術不正.
和喪失真理.

$^{641}$傳道變成是一些她的教義.
所以是別的教義.
不是純正的聖經的教導.
所以她會爭的是為自己來爭.
因為她本身的內心.
就是自高自大.
和敬虔.
原來只是一個裝扮.
是一個手段.
但是另外那三種人.
就很不同了.
就在生命裡面遇見主耶穌.
被祂的大愛改變.
而真心實意的.
願意在地上.
就算你自己現在的身份改變不了.
你是一個奴隸.
在座應該不會有吧?.
但是你有沒有覺得.
自己的身份很低微呢?.
有沒有對自己的出生.
成長家庭.
很多很多的厭惡呢?.
還是你信了主耶穌之後.
你越來越對自己的生命.
神在不同的人生帶領你.
你充滿了感恩.
要這樣來醫治自己的過去.
要在主耶穌的眼中.
看你自己的過去.
那你就會蒙福了.
如果你本身就已經是.
譬如在世界上讀書.
又能夠賺到錢.
又能夠積聚很多的財富.
哇!神給你的責任更大了.
你就可以祝福更多的人.
當然更寶貴.
如果能夠有更多的.
年輕的提摩太.

$^{681}$可以看到神的真道的寶貴.
可以改變別人的生命.
我聽到老牧師說.
你們的主日學現在很多人上.
我們牧者就很開心.
因為能夠大家對真道的認識.
渴慕和追求.
這樣就可以讓教會.
在神的真理裡面.
來扎根.
有更多的提摩太.
可以來到他起來.
所以求主幫助我們今天.
和大家在這幾段聖經裡面.
這樣來分享.
一起去學習.
我們做一個敬虔的人.
真正的財富.
就是你內心的敬虔.
對主耶穌的熱愛和追求.
和很想讓其他人.
都得到神的恩典.
我們一起祈禱.
多謝主.
你的話語何等的寶貴甘甜.
使我們能夠得著提醒,教導.
願主幫助.
聖靈在我們大家的內心.
燃點這一份對主耶穌的認識.
對真正的敬虔生命的渴慕.
以致我們每一位弟兄姊妹.
都一起去追求我們的主.
這樣感恩禱告.
奉耶穌基督的名.
阿們.
\newpage



\section{約翰福音 9:1-7-13-17-30-41}
\label{sec:qTaSFAQFmRY}
\textbf{2022.09.04 《屬靈的看見》 約翰福音 9\_1-7,13-17,30-41 (和修版) 老耀雄牧師}
\newline
\newline
連結: \href{https://youtube.com/watch?v=qTaSFAQFmRY}{\texttt{https://youtube.com/watch?v=qTaSFAQFmRY}} ~~~~ 語音日期: 2022-09-04
\newline
\newline
\hyperref[sec:D4j0WaSCINU]{\small{< < < PREV SERMON < < <}}
~
\hyperref[sec:index_chronic]{\small{[返順時目]}}
~
\hyperref[sec:index_scriptual]{\small{[返順卷目]}}
~
\hyperref[sec:u7aPgc4mo0Y]{\small{> > > NEXT SERMON > > >}}
\newline
\newline
約翰福音 9:1-7-13-17-30-41
\newline
\begin{longtable}{cl}
\hline
\hline
章節 & 經文 (和合本修訂版)\\
\hline
9:1 & \begin{tabularx}{0.7\textwidth}{X} 耶穌往前走的時候，看見一個生來就失明的人。 \end{tabularx} \\ \\ \relax
9:2 & \begin{tabularx}{0.7\textwidth}{X} 門徒問耶穌：「拉比，這人生來失明，是誰犯了罪？是這人還是他的父母呢？」 \end{tabularx} \\ \\ \relax
9:3 & \begin{tabularx}{0.7\textwidth}{X} 耶穌回答：「既不是這人犯了罪，也不是他的父母，而是要在他身上顯出神的作為來。 \end{tabularx} \\ \\ \relax
9:4 & \begin{tabularx}{0.7\textwidth}{X} 趁著白日，我們必須做差我來的那位的工；黑夜來到，就沒有人能做工了。 \end{tabularx} \\ \\ \relax
9:5 & \begin{tabularx}{0.7\textwidth}{X} 我在世上的時候，是世上的光。」 \end{tabularx} \\ \\ \relax
9:6 & \begin{tabularx}{0.7\textwidth}{X} 耶穌說了這些話，就吐唾沫在地上，用唾沫和了泥抹在盲人的眼睛上， \end{tabularx} \\ \\ \relax
9:7 & \begin{tabularx}{0.7\textwidth}{X} 對他說：「你到西羅亞池子裡去洗。」（西羅亞翻出來就是「奉差遣」。）於是他去，洗了，回來就看見了。 \end{tabularx} \\ \\ \relax
9:8 & \begin{tabularx}{0.7\textwidth}{X} 他的鄰舍和素常見他討飯的人，就說：「這不是那從前坐著討飯的人嗎？」 \end{tabularx} \\ \\ \relax
9:9 & \begin{tabularx}{0.7\textwidth}{X} 有的說：「是他」；又有的說：「不是，卻是像他。」他自己說：「是我。」 \end{tabularx} \\ \\ \relax
9:10 & \begin{tabularx}{0.7\textwidth}{X} 於是他們對他說：「你的眼睛是怎麼開的呢？」 \end{tabularx} \\ \\ \relax
9:11 & \begin{tabularx}{0.7\textwidth}{X} 那人回答：「有一個名叫耶穌的，他和了泥抹我的眼睛，對我說：『你到西羅亞池子去洗。』我去一洗，就看見了。」 \end{tabularx} \\ \\ \relax
9:12 & \begin{tabularx}{0.7\textwidth}{X} 他們對他說：「那個人在哪裡？」他說：「我不知道。」 \end{tabularx} \\ \\ \relax
9:13 & \begin{tabularx}{0.7\textwidth}{X} 他們把以前失明的那個人帶到法利賽人那裡。 \end{tabularx} \\ \\ \relax
9:14 & \begin{tabularx}{0.7\textwidth}{X} 耶穌和泥開他眼睛的那一天是安息日。 \end{tabularx} \\ \\ \relax
9:15 & \begin{tabularx}{0.7\textwidth}{X} 法利賽人又問他是怎麼得看見的。他對他們說：「他把泥抹在我的眼睛上，我一洗，就看見了。」 \end{tabularx} \\ \\ \relax
9:16 & \begin{tabularx}{0.7\textwidth}{X} 於是法利賽人中有的說：「這個人不是從神來的，因為他不守安息日。」另有的說：「一個罪人怎能行這樣的神蹟呢？」他們之間就產生了分裂。 \end{tabularx} \\ \\ \relax
9:17 & \begin{tabularx}{0.7\textwidth}{X} 於是他們又對那盲人說：「他開了你的眼睛，你說他是怎樣的人呢？」他說：「他是個先知。」 \end{tabularx} \\ \\ \relax
9:18 & \begin{tabularx}{0.7\textwidth}{X} 猶太人不信他以前是失明，後來能看見的，等到叫了他的父母來， \end{tabularx} \\ \\ \relax
9:19 & \begin{tabularx}{0.7\textwidth}{X} 問他們說：「這是你們的兒子嗎？你們說他生來是失明的，現在怎麼看見了呢？」 \end{tabularx} \\ \\ \relax
9:20 & \begin{tabularx}{0.7\textwidth}{X} 他的父母就回答說：「他是我們的兒子，生來就失明，這是我們知道的。 \end{tabularx} \\ \\ \relax
9:21 & \begin{tabularx}{0.7\textwidth}{X} 至於他現在怎麼能看見，我們卻不知道；是誰開了他的眼睛，我們也不知道。他已經是成人，你們問他吧，他自己會說。」 \end{tabularx} \\ \\ \relax
9:22 & \begin{tabularx}{0.7\textwidth}{X} 他父母說這些話，是怕猶太人，因為猶太人已經商定，若有宣認耶穌是基督的，要把他趕出會堂。 \end{tabularx} \\ \\ \relax
9:23 & \begin{tabularx}{0.7\textwidth}{X} 因此他父母說「他已經是成人，你們問他吧」。 \end{tabularx} \\ \\ \relax
9:24 & \begin{tabularx}{0.7\textwidth}{X} 於是法利賽人第二次叫了那以前失明的人來，對他說：「你要將榮耀歸給神，我們知道這人是個罪人。」 \end{tabularx} \\ \\ \relax
9:25 & \begin{tabularx}{0.7\textwidth}{X} 那人就回答：「他是不是個罪人，我不知道；有一件事我知道，我本來是失明的，現在我看見了。」 \end{tabularx} \\ \\ \relax
9:26 & \begin{tabularx}{0.7\textwidth}{X} 他們就問他：「他給你做了甚麼？是怎麼開了你的眼睛？」 \end{tabularx} \\ \\ \relax
9:27 & \begin{tabularx}{0.7\textwidth}{X} 他回答他們：「我已經告訴你們了，你們不聽，為甚麼又要聽呢？難道你們也要作他的門徒嗎？」 \end{tabularx} \\ \\ \relax
9:28 & \begin{tabularx}{0.7\textwidth}{X} 他們就罵他：「你是他的門徒，而我們是摩西的門徒。 \end{tabularx} \\ \\ \relax
9:29 & \begin{tabularx}{0.7\textwidth}{X} 神對摩西說話是我們知道的，可是這個人，我們不知道他從哪裡來。」 \end{tabularx} \\ \\ \relax
9:30 & \begin{tabularx}{0.7\textwidth}{X} 那人回答，對他們說：「他開了我的眼睛，你們竟不知道他從哪裡來，這真是奇怪！ \end{tabularx} \\ \\ \relax
9:31 & \begin{tabularx}{0.7\textwidth}{X} 我們知道神不聽罪人，惟有敬奉神、遵行他旨意的，神才聽他。 \end{tabularx} \\ \\ \relax
9:32 & \begin{tabularx}{0.7\textwidth}{X} 從創世以來，未曾聽見有人開了生來就失明的人的眼睛。 \end{tabularx} \\ \\ \relax
9:33 & \begin{tabularx}{0.7\textwidth}{X} 這人若不是從神來的，甚麼也不能做。」 \end{tabularx} \\ \\ \relax
9:34 & \begin{tabularx}{0.7\textwidth}{X} 他們回答他說：「你完全是生在罪中的，還要來教訓我們嗎？」於是他們把他趕出去了。 \end{tabularx} \\ \\ \relax
9:35 & \begin{tabularx}{0.7\textwidth}{X} 耶穌聽說他們把他趕出去，就找到他，說：「你信人子嗎？」 \end{tabularx} \\ \\ \relax
9:36 & \begin{tabularx}{0.7\textwidth}{X} 那人回答說：「主啊，人子是誰？告訴我，好讓我信他。」 \end{tabularx} \\ \\ \relax
9:37 & \begin{tabularx}{0.7\textwidth}{X} 耶穌對他說：「你已經看見他，現在和你說話的就是他。」 \end{tabularx} \\ \\ \relax
9:38 & \begin{tabularx}{0.7\textwidth}{X} 他說：「主啊，我信！」他就拜耶穌。 \end{tabularx} \\ \\ \relax
9:39 & \begin{tabularx}{0.7\textwidth}{X} 耶穌說：「我為審判到這世上來，使不能看見的看見，能看見的反而失明。」 \end{tabularx} \\ \\ \relax
9:40 & \begin{tabularx}{0.7\textwidth}{X} 同他在那裡的法利賽人聽見這些話，就對他說：「難道我們也失明了嗎？」 \end{tabularx} \\ \\ \relax
9:41 & \begin{tabularx}{0.7\textwidth}{X} 耶穌對他們說：「你們若是失明的，就沒有罪了；但現在你們說『我們能看見』，你們的罪還在。」 \end{tabularx} \\ \\
[1ex]
\hline
\hline
\end{longtable}
$^{1}$(自動翻譯).
各位姐妹平安.
平安.
我們都知道人都是有限制的.
好像我們身體的感官系統一樣.
限制了我們接觸另一個領域和空間.
我們的眼睛沒有鷹看得那麼遠.
也不像蒼蠅一樣.
可以有360度的視野.
我們的耳朵沒有狗一樣靈敏.
它們的聽力可以聽到超音波的聲音.
甚至狗的嗅覺都比人優勝.
起碼我們聞不到毒品.
我們的腦也沒有導航系統.
不像是後鳥一樣.
它們可以跨越大西洋,太平洋,印度洋.
而不會迷路.
而蝙蝠更加可以用超音波回聲定位的訊號.
去尋找食物,探測距離,確定目標.
上述生物所擁有的能力.
都是人所沒有的.
當我們人類只能夠接觸到長,闊,高的三維世界的時候.
它們只能夠領受從這個有限的空間所發出來的訊息.
它們不能夠感受到更加遠,更加高,更加深的世界.
我們所見到的.
就只是一個有限度的距離.
究竟人類能不能夠超越這些限制呢?.
人類的腦袋比較奇妙.
它發明了很多東西.
科技上不斷有新的發明.
以不同的機器,晶片,電腦.
去彌補了我們一些的不足和限制.
不過沒有一樣的發明.
能夠讓我們可以探索到生命的價值.
我們為甚麼要生?.
我們為何而活呢?.
我們將要何處呢?.
至今都沒有任何一種科技.
可以提供答案給我們.
人類仍然都要為這些生命的問題.

$^{41}$去尋找答案.
或者尋找生命的途徑.
並不能夠用我們有限的感官去尋找.
而是需要戴上一副屬靈的眼鏡.
才能夠超越我們的限制.
今天我想透過《約翰福音》第九章.
一個核子的故事.
去探索一個生命的故事.
我們以祈禱開始今天的講道.
我們人類雖然有很多的限制.
但你沒有限制我們渴求你的心.
你也沒有限制我們用心靈去感受和敬拜你.
你更加沒有限制我們一生去跟從你.
因此求主你開開我們的眼睛.
讓我們能夠看得見你的榮耀.
看得見你的美善.
看得見你是怎樣去愛我們.
我們獻上感恩祈禱.
奉靠主耶穌基督得勝明智祈求.
阿們.
《約翰福音》的作者.
用了整整一章.
四十一節的經文.
去講述這個核子被醫治的故事.
這個神跡的醫治故事.
是要讓人知道耶穌的身份,能力.
以及生命成長的過程.
經文一開始.
滿圖就一個生來失明的核子.
向耶穌問了一個問題.
拉比這個人生來就失明了.
究竟是誰犯了罪?.
是這個人犯罪?.
還是父母犯罪?.
原來猶太人的傳統.
一直都認為一個人的殘缺.
是因為他們犯罪得罪了神.
而被神懲罰的後果.
所以猶太人認為.
無論是核子.

$^{81}$或者是瘸腿的耳聾的人.
他們身患殘疾.
是因為他們有罪.
而且是罪有應得的.
不過這個核子是生來就失明的.
所以滿圖就產生一種疑惑.
究竟犯罪是核子本身?.
還是它的失明是它自身的惡果?.
還是是核子的父母犯罪?.
所以父母的罪.
使嬰孩在母胎裡面.
已經成為一個殘疾?.
因此孩子一出世就已經失明了.
滿圖其實向耶穌問了一個神學的問題.
什麼神學問題?.
就是罪是否會承傳到下一代?.
耶穌就回應了.
既然不是這個人犯罪.
也不是他的父母.
而是要在他身上顯出神的作為.
耶穌沒有回應到.
核子和他的父母是否有犯罪的問題.
而是回應了.
如何可以在人的罪裡面.
顯出神的榮耀這個問題.
既然人有罪.
而朝至有缺陷.
人當然沒有辦法解決這個問題.
因此人在如何解決罪的問題上.
是完全無能為力的.
不過如果能夠將罪所帶來的一切惡行.
以及對人的捆綁.
都能夠得到赦免和釋放的話.
這樣就能夠顯出神的能力和榮耀了.
因為人沒有辦法單獨去性覺惡.
但是靠著神的恩典與赦免.
人就可以得到釋放.
耶穌再加以補充.
我們必須做差我來的那位一工.
黑夜來到就沒有人能做工了.

$^{121}$我在世上的時候是世上的光.
原來為神作功以及顯出神的榮耀.
也都需要爭取時間.
因為黑暗勢力來到的時候就做不到了.
因為神榮讓黑暗勢力在地上.
可以掌權一段時間.
耶穌向門徒說明.
人的罪能不能夠被赦免和釋放.
是會有黑暗勢力出來作出抗衡的.
而這種黑暗勢力往往是我們忽略的.
雖然我們知道黑暗.
最終是不能勝過光明.
但我們不能夠因此而輕看這一場屬靈爭戰.
如果我們今次有看過.
今年培靈賢經大會最後一堂的奮興會.
鮑維奎牧師就形容.
神的子民就是神榮耀的軍隊.
神榮耀的軍隊的爭戰對象.
就是敵黨神的撒旦.
而神榮耀的軍隊要拯救的對象.
就是失喪的靈魂.
因此教會存在的其中一個目的.
就是讓更多的罪人得到赦免和釋放.
而我們爭戰的其中一種武器.
就是傳講復和和釋放的福音.
讓所有被撒旦捆綁的人.
能夠回歸神的角度.
這樣神就可以得到當德的榮耀.
撒旦國就會被擊敗.
如果信徒不去傳福音.
那就令撒旦國擴張.
令神的角度萎縮.
也令神得不到榮耀.
信徒屬靈的裝備當然也包括祈禱.
祈禱是戰勝撒旦的關鍵元素.
如果信徒輕看了祈禱的能力.
認為祈禱只是向空氣說話.
就表示我們已經放下了屬靈的武器.
只是用血氣與惡魔爭戰.
這樣慘敗是必然的結果.

$^{161}$教會想成為福音的平台.
也只是空談.
耶穌說「必須做差外來的那位的工.
必須」.
這個原文是一個命令的時態.
不是那種可有可無的.
必須讓人生命得釋放.
然後才得到榮耀神.
因此這不單單是耶穌的命令.
也是每一個神子民的使命和責任.
弟兄姊妹.
你們有沒有這份讓神得榮耀的心?.
我們傳福音的時候.
有沒有輕看撒旦的攻擊?.
既然知道有攻擊.
我們有沒有裝備自己.
然後才出去打這場仗?.
抑或我們只是憑血氣之勇.
然後撞到焦頭爛額回來?.
上個月教會推薦了七個肢體.
參加了青年歸主協會舉辦的成長八課導師課程.
課程讓學員對核心信仰有更深入的認識.
有更多的信仰反省.
七個肢體願意先接受信仰的裝備.
經過自己對課程內容的消化.
然後才建立別人的生命.
這一份是很認真的事奉態度.
因為我們不想空槍上陣.
而是先接受裝備.
然後才最後被神所使用.
另外也有三位姊妹.
接下來的星期二.
將要參加短宣中心的傳福音課程.
為期三年.
她們都是只想裝備自己.
然後被神使用.
也有一位姊妹.
去讀宣教的課程.
希望她們對宣教的認識裡.
更能夠得著那種神的力量.

$^{201}$各位姊妹.
我們真的需要知道.
我們不只是坐在這裡叫句口號.
這就叫做毀身事奉.
而是需要身體力行.
去打這場屬靈的爭戰.
如果教會推薦你去接受裝備.
你會不會接受?.
你願不願意在黑夜來臨之前.
在白晝為主作供.
抑或你只是願意坐在你的安書區裡.
繼續過你一個人的信仰生活呢?.
核子得到耶穌的醫治之後.
他的生命有甚麼改變呢?.
核子的鄰居問他.
「你的眼睛是怎樣開的?」.
核子就回應.
「有一個名叫耶穌的.
他把泥沫磨在我的眼睛.
叫我到西羅亞池子去洗.
我一洗就看見了」.
核子當時只是知道耶穌的名字.
他只是知道耶穌的名字.
他只是知道是耶穌讓他能夠看見.
這是他對耶穌的認識的一個開始.
回到塞爾就問核子.
「他開了你的眼睛.
你說他是甚麼人呢?」.
核子說「他是個先知」.
核子從知道耶穌的名字.
進而確認耶穌的身份.
核子知道耶穌不是一個普通人.
起碼耶穌是一個大有能力的人.
從對耶穌名字的認識.
到確認耶穌的身份.
是核子的一個成長過程.
之前法系賽人曾經詢問過核子的父母.
「你們說他生來是失明的.
現在如何看見呢?」.
核子的父母因為很怕被法系賽人趕出會堂.

$^{241}$所以他們就回應.
「他現在如何能夠看見我們不知道啊.
他已經是成年人了.
你們問他吧」.
核子的父母雖然知道耶穌的大能.
也知道是耶穌治好他的兒子.
不過卻害怕被趕出會堂.
受到社會所排斥和輕視.
而不願意為耶穌作見證.
我們看到人的軟弱.
縱然得到耶穌的恩典.
卻因著個人的利益.
而不願意為耶穌作見證.
法系賽人又對核子說.
「他給你做了甚麼呢?.
是甚麼開你的眼的?」.
又說「你張榮耀給神.
因為我們知道這個人不是一個罪人」.
核子針對法理賽人認為耶穌是罪人.
作出了這樣的回應.
他說「他是否一個罪人我就不知道了.
不過有一件事我就知道.
就是我本來是失明的.
現在我看到了」.
核子以現實的結果來回應法系賽人.
他從失明到能夠看見.
就是一種生命獲得釋放的過程.
甚至是生命轉化的過程.
耶穌這種能力.
可以是一個罪人可以成就的嗎?.
核子又針對耶穌如何開他的眼睛.
以及耶穌的身份作出回應.
他說「他開了我的眼睛.
你竟然不知道他從哪裡來?.
唯有敬奉神.
遵行他旨意的神才會聽的.
從創世以來.
未曾聽見有人開了生來就失明的人的眼睛.
這個人若不是從神而來.
甚麼都做不到」.

$^{281}$核子不像父母一樣軟弱.
他不單單願意為耶穌申辯.
也願意為耶穌作見證.
他明白耶穌的能力是從神而來的.
是神透過耶穌醫治他.
既然耶穌是從神而來的.
那麼耶穌就不會是一個罪人了.
這麼簡單的道理.
為甚麼法利賽人會不明白呢?.
不過核子雖然勇於為耶穌作見證.
不過他仍然要承擔.
他父母極想逃避的考驗.
核子最終被法利賽人趕出會堂.
核子為了耶穌而作了見證.
也為此付出了代價.
當耶穌知道核子被法利賽人趕出會堂後.
他就找到核子.
對核子說「你信人子嗎?」.
核子回應「主啊!人子是誰呢?」.
核子說「讓我信他」.
耶穌對他說「你已經見到他了」.
「現在跟你說話的就是他」.
於是核子說「主啊!我信」.
他就拜耶穌.
核子從認識耶穌的名字開始.
到確認耶穌的身份.
願意為耶穌作見證.
當核子稱呼耶穌為主的時候.
他對耶穌的信心達到圓滿.
這一刻他真真實實地遇到他生命的主.
他的眼睛不但被醫好.
他的心靈也得到耶穌的牧養.
是耶穌親自找到核子的.
是耶穌親自向核子呼召的.
核子在未見到耶穌之前.
他從來沒有盼望過那隻眼睛能夠醫好.
核子在未見到耶穌之前.
他是被排斥的.
被孤立的.
是一個被歧視的弱勢社群.

$^{321}$被社會認定是一個得罪了神的罪人.
而且是應有此報的.
甚至核子首次遇到耶穌.
都是耶穌作出主動去醫治他的.
不過核子把握著每一個機會.
當耶穌將唾沫和泥.
抹在他的眼睛上的時候.
核子沒有抗拒.
當耶穌要求他前往西雷亞池洗一洗的時候.
核子又願意順服.
我們想一想.
難道耶穌不可以用說話.
就能夠將核子立刻醫好嗎?.
耶穌是有這樣的能力的.
既然耶穌可以用說話就能夠醫好核子.
為什麼還要核子做出這樣的行為呢?.
因為這兩個行為的回應.
都是對核子的一個考驗.
他是否抗拒?.
他是否願意順服呢?.
這成為他對耶穌信心的一個基礎.
當耶穌知道核子被趕出會堂.
耶穌是主動去找回這個核子.
然後向核子作出呼召.
你已經看見他了.
現在跟你說話的就是他了.
核子也把握著這個機會.
回應耶穌.
主啊,我信.
對於核子來說.
他以往的人生是一個不被重視.
也沒有盼望的人生.
因為他無論在生活上.
宗教上都參與不了.
可以說他過著一個被社會遺棄的人生.
然而從他遇到耶穌的那一剎那開始.
他把握著每一個認識耶穌的機會.
讓他的屬靈生命一步一步地成長.
這個成長不單單有頭腦上的認識.
認同耶穌的身份.

$^{361}$也有對耶穌的信心和倚靠.
以至核子的生命改變和轉化.
讓他最終可以在社會上.
宗教上都獲得認同.
核子生命的改變.
最終成為見證神的榮耀.
弟兄姊妹.
核子被耶穌醫治之後.
他沒有畏懼這個權勢.
就算有機會被趕出會堂.
他仍然都要向法利賽人宣認.
這個人若不是從神而來.
甚麼都做不到.
核子對信仰的回應.
是一個勇於為主作見證的人生.
我們是否願意為主作見證?.
今日我們為主作見證.
不會被人趕出教會.
起碼我不會趕你出教會.
亦都不會被排斥.
如果我們有留意.
我們每一個禮拜的程序表.
都有生命見證的分享.
弟兄姊妹將他們的生命見證寫出來.
這個也是對神的一個回應.
當我們十一月堂慶的時候.
將這些生命見證收集成雲彩集的時候.
眾多肢體的見證.
就是成為洪恩家.
呈獻給主的一份禮物.
這份禮物有沒有你的份?.
你願意參與其中?.
將你的見證說出來.
你願意將你和神的故事述說出來?.
願榮耀歸還給我們天上的父.
核子屬靈生命的圓滿.
本來已經是故事的結束了.
它已經被耶穌治好了.
誰知耶穌在結尾的時候.
再加一句.

$^{401}$「我為審判到這世上來.
使到不能看見的看見.
能看見的反而失明」.
這句話很吊詭.
我們明白什麼叫做.
使不能看見的看見.
就好像核子一樣.
它本來看不見的.
但藉著耶穌的大能.
核子最終都能夠看到.
但是能夠看見的反而失明.
這是什麼意思呢?.
是讓那些看不見的人.
不是.
是讓那些看見的人成為核子嗎?.
耶穌在說論一個人.
能不能看見這句話的時候.
也有一個目的.
就是「我為審判到這世上來」.
因此一個人能夠看到.
是審判的結果.
一個人看不到.
也是審判的結果.
既然看得見是這麼重要.
我們就要問.
究竟看見什麼呢?.
看見什麼內容.
這個更加重要.
看見什麼呢?.
就是核子所看見的東西.
就是看到神的臨在.
也看到神的救恩.
核子看到的耶穌.
不單單是一個肉身的耶穌.
他也看到耶穌的神聖.
他知道耶穌如果不是從神而來的話.
他就什麼都做不到.
因此核子看到神的臨在.
他領受到從神而來的.
耶穌為他的醫治.

$^{441}$他看見耶穌的真實.
神聖和恩典.
他把握著每一個機會.
他看見站在他面前的耶穌.
他就知道耶穌就是神.
因此他對耶穌宣認.
「主啊,我信」.
在神的審判下.
核子這個宣認.
他的罪被得赦免.
因為他得到神的拯救.
法利賽人也看到耶穌.
但看不到耶穌的神聖.
以往耶穌多次在眾人面前行神蹟.
法利賽人也在場.
但法利賽人只看到.
耶穌在安息日裡治病.
他們只關注安息日的條例.
看不到耶穌有醫治的權柄.
法利賽人看過耶穌趕鬼.
他們認為耶穌是靠著鬼王趕鬼.
他們只相信魔鬼有權能.
卻看不到耶穌也有神的權能.
法利賽人看到.
耶穌的門徒不洗手,不禁食.
他們只看到潔淨條例和傳統被破壞.
他們看重的是禮儀.
以及恩秦的文化.
卻看不到耶穌要更新律法的精神.
法利賽人有眼睛.
卻看不到對生命最重要的救恩.
法利賽人沒有失明.
卻看不到道成肉身的耶穌的身份.
他們看到的只是破壞律法的耶穌.
是那些說褻瀆話的耶穌.
因此法利賽人聯合了民事,祭司合謀.
想殺害耶穌.
於是就印證了耶穌所說的.
「能看見的反而失明」.
法利賽人不是肉眼的失明.

$^{481}$而是屬靈上的失明.
他們看不到主耶穌的神聖.
抗拒救恩.
在神的審判下就注定被定罪.
真正能夠看見神的作為.
並不是用肉眼的眼睛去看.
而是用心靈的眼睛去感受.
我們要像核子一樣.
用心靈與神相遇,相交.
用心靈去察驗神的尾線.
亦要用心靈與主建立關係.
最後我們要問問自己.
我們是否想擁有屬靈的眼睛.
與神相遇,相交呢?.
今天講題是「屬靈的看見」.
生命的本質.
因為我們的肉眼真的有限制.
要用我們的屬靈生命去察看.
體驗和領受.
我們的生命被神從罪中釋放出來.
要經歷破碎到成熟的不同階段.
每一個階段都需要時間.
才能夠被神建立和栽培.
屬靈的操練是必須的.
沒有捷徑.
亦沒有不勞而獲.
關鍵是我們願不願意付上代價.
我們不單要從理性上去認識耶穌.
去體驗耶穌的真實.
亦要用屬靈的眼睛去看身邊的人和事.
以及這個世界.
這不是理性上可以感受得到的.
最後我要提醒我們自己.
我們的生命是有限的.
我們要把握時間.
用生命去為主作見證.
為主作供.
有人說.
我不知道.
我不明白.

$^{521}$我沒有聽過.
我們說這些人.
無知.
因為沒有聽過.
不知道,不明白.
那就是無知.
但如果聽過.
明白.
但不去行.
我們就不能說他是無知.
他是明知而固乏.
神不喜悅.
或神不悅立這種知而不行的信仰人生.
如果你這樣做的話.
你的罪就已經定了.
想和神建立.
想被神建立和塗造.
是需要我們信服和遵行.
信服是心態.
遵行是行為.
唯有心口如一.
恩典與行為相稱的信仰生命.
才是神所閱立的生命.
而生命最大的塗造工程.
就是讓我們有屬靈的洞見.
這種屬靈的洞見.
能夠超越我們的限制.
可以察驗到.
甚麼是神所喜悅的心意.
又感受到我們每一個人的內心深處.
神所厭惡的一種慾念.
以及我們的心思意念.
好像昔日的先知一樣.
我們能夠有這種屬靈的洞見.
就能夠看到神給我們的遺像.
讓我們的生命的遺像.
成為我們信仰生命的動力.
為主跑過標竿人生.
作美好的見證.
甚至我們每一個弟兄姊妹.

$^{561}$都能夠擁有屬靈的眼睛.
看到耶穌為我們所作的.
祂與我們同行同在.
祂叫我們以感恩的心去回應.
我近日看了一本書.
是我神學院老師趙崇明副教授.
他對祈禱有以下的分享.
一般的信徒.
對祈禱好像.
不太認識.
我讀完他這個分享之後.
我希望大家可以對祈禱.
有一個重新的反省.
你慢慢咀嚼.
這個祈禱的內容.
祈禱不是以我為主詞.
而是以上帝為中心.
祈禱不是以人的言語.
為主要的內容.
而是以上帝的言語為本.
祈禱不是始於人的開口.
而是始於人的閉口不語.
祈禱不是向上帝喋喋不休.
而是在靜默中聆聽上主的微聲細語.
祈禱不是人去索求上帝的垂聽.
而是上帝要求人的聆聽.
祈禱不是不斷地向上帝呼喚.
而是領受被上主的呼喚.
祈禱不是要改變上帝.
而是讓上帝改變我們.
祈禱不是對著空氣自說自語.
而是經驗三一上主臨在的交談.
祈禱不是要與眾星圈蛙的世界隔絕.
而是在眾星圈蛙中學習心神專注.
祈禱不是強求上帝成全我的意思.
而是要求自己按上主的意思行事.
祈禱不是首先顧念自己的事和益處.
而是先求上帝的覺和上帝的義.
祈禱不是一篇詞粗華麗的演說.
而是側樸無華的真誠心底話.

$^{601}$祈禱不是只為昨天懊悔及為明天憂慮.
而是為今天所領受的生命禮物而感恩.
祈禱不是懷著消費的心態去妄求.
而是祈求有神就一無所缺的知足.
祈禱不是向上帝報告日常所事.
而是閣下日常事務與上帝閒談.
我邀請大家在座位裡用心靈向上主祈禱.
然後我再帶大家做一個結束祈禱.
(結束).
秋下主我們每一個人的心思意念.
你都知道.
求你閱立眾人在你面前無聲的祈願.
並按你的旨意一一去成就.
求主你開我們的眼.
讓我們能夠見到感受到.
你的的確確臨在我們當中.
你與我們實實在在的同在同行.
你的臨在是我們得力量.
得安慰.
得勇氣.
得激勵.
以及一切得福的來源.
願我們懷著感恩的心.
我們的祈禱是.
「凡教主耶穌基督得勝名自祈求」.
阿們.
願主祝福大家.
謝謝大家.
\newpage



\section{馬太福音使徒行傳羅馬書 28:19-20}
\label{sec:u7aPgc4mo0Y}
\textbf{2022.09.18 《承擔宣教的使命》 馬太福音 28\_19-20;使徒行傳 1\_8;羅馬書 10\_14-15 (和修版) 余家寶傳道}
\newline
\newline
連結: \href{https://youtube.com/watch?v=u7aPgc4mo0Y}{\texttt{https://youtube.com/watch?v=u7aPgc4mo0Y}} ~~~~ 語音日期: 2022-09-21
\newline
\newline
\hyperref[sec:qTaSFAQFmRY]{\small{< < < PREV SERMON < < <}}
~
\hyperref[sec:index_chronic]{\small{[返順時目]}}
~
\hyperref[sec:index_scriptual]{\small{[返順卷目]}}
~
\hyperref[sec:QekXUaNsm08]{\small{> > > NEXT SERMON > > >}}
\newline
\newline
$^{1}$魯牧師 紅印堂改弟兄姊妹 早安 主脈平安.
我是儒家堡 係錦繡堂 差傳部支持的宣教士.
現在和荷蘭Send the Light華人基督教會合作.
在盧修甫堡做宣教的工作.
今次可以和大家分享一下宣教的事復.
說到宣教的使命.
我們可以回到《馬太福音》第28章最後那段經文.
看看主耶穌對我們有什麼的教導.
當時主耶穌已經完成了世上救贖的使命.
在升天之前向門徒說了一個很重要的信息.
因為當時門徒看著主耶穌遇難離開他們.
跟隨了一段時間的領袖魯斯突然離開他們.
大家都不知道如何應變.
突然而來的打擊令到大家都茫無頭緒.
前面的路都不知道如何繼續走下去.
就在這個時候主耶穌透過很簡單的指示.
告訴他們前面繼續應該怎樣走.
他們需要從過往跟隨主耶穌.
做門徒一個被動的角色.
一個學習的角色.
轉變成一個主動做領導,教導,帶領人傳福音的角色.
所以主耶穌的離開也是門徒踏入一個新局面的開始.
主耶穌在升天前的40天.
不斷向門徒顯現.
解釋給他們聽有關天國的道理.
最後叮囑門徒要耐心等候.
等候神所應許賜給他們的聖靈.
然後就隨著一朵雲彩升天.
領受主耶穌的使命.
很可能門徒都不是完全明白裡面的意思.
去?去哪裡呢?.
使萬民作耶穌的門徒.
萬民是什麼人?.
怎樣使他們成為耶穌的信徒?.
奉父子聖靈的名.
給他們施洗.
我們又不是施洗的漢.
哪有這樣的權柄為人施洗呢?.
教導他們遵守主耶穌的教訓.
其實主耶穌的教訓.

$^{41}$可能連他們都不能完全明白.
然後還要去到世界的活了.
什麼時候才是世界的活了?.
得來人的思想,理解和能力.
真是怎樣能夠做到.
去完成這些奧秘的事情呢?.
所以主耶穌跟他們說.
你們要等候.
等候聖靈.
那聖靈又是什麼呢?.
主耶穌跟他們說.
我有跟他們說過聖靈.
我們可以回到.
《密約開恩》第十四章十六節.
那裡耶穌提到一個保衛師的名稱.
保衛師.
或者在細節上寫.
奮衛師.
那保衛師是什麼呢?.
我們可以從原文的用字.
去理解裡面的意思多一點.
這個字的意思包含了.
很有力量的.
這個力量是能夠實現出來的.
第一個意思.
就是一個幫手的意思.
幫手就不是助手.
而是可以在困境裡.
去幫忙人.
去解決問題.
第二.
是一個安慰者的意思.
這個安慰是有一個支持著.
人可以站起來.
重新面對前面的力量.
也有鼓勵者的意思.
這個鼓勵者.
特別是形容到一些危難.
在戰爭的情況裡.
他不斷地要鼓勵這個人.

$^{81}$讓他有勇氣勇往直前.
第四.
是包含了辯護者的意思.
當人面對神的審判的時候.
他會做人的辯護者.
為人去爭辯.
最後還有一個.
智慧者的意思.
智慧者.
不是給意見那麼簡單.
而是告訴你.
一種有智慧可以解決問題的途徑和方法.
我們再看第十四章26到27節的時候.
更加解釋了.
聖靈可以幫助人.
去明白神的道理.
明白真理的奧妙.
所以當門徒在五巡節.
領受了聖靈的時候.
他們就開始明白.
一直以來疑惑不明白的地方.
但是聖靈就正式啟動了.
他們的生活也很奇妙地展開了不同的變化.
在《使徒行傳》第一二章那裡描述.
當使徒聽了主耶穌跟他們說.
要等候聖靈降臨.
然後他們就會得到力量.
得到力量之後.
有些事情他們要做.
就是將福音傳到更遠的地方.
當門徒領受了聖靈的時候.
其實同一時間他們也領受了聖靈.
所賜給他們的力量.
所以他們有能力去改變自己的生活方式.
他們開始一個新的生活方式.
他們突破了自己的能力.
也很開開心心地期待著為主而活.
當人聽了神的道理.
願意悔改.
願意去改變自己的思想和行為.

$^{121}$大家可以互相接納.
可以彼此照顧.
大家將自己擁有的東西.
拿出來大家繁文公容.
互相去支持.
和諧共處.
開開心心誠實地生活.
無論大家聚集在一起.
還是在家裡.
都很開心地唱詩歌去讚美神.
去祈禱.
去渴望去聽使徒解釋.
神的動力.
神的奧妙.
開開心心面對著自己每天的生活.
很多未信主的人.
都是因為這些人.
他們的行為見證.
和使徒所做的神蹟奇事.
而相信了神.
而福音就這麼奇妙地.
在耶路撒冷傳開了.
跟著隨著人群離開耶路撒冷.
也有些人因為信了主耶穌.
他們受到政治的迫害.
需要離開自己的地方.
或者好像保羅那樣.
日後去到很多不同的旅程.
去傳授福音.
神的動力.
主耶穌的教恩.
就這樣傳開了.
由耶路撒冷.
由大全地.
撒瑪利亞.
去到更遠的地方.
直到地極.
正如《使徒行傳》一章八節.
我們這樣說.
福音要由耶路撒冷.

$^{161}$由大全地.
撒瑪利亞.
傳到直到地極.
主耶穌的使命.
就這樣開始和延續下去.
當人願意去回應.
用行動去表達出來的時候.
聖靈又會加給他們力量.
在他們身上發揮出聖靈的工作.
在《使徒行傳》的記載.
更加讓我們看到.
有另外一幅更加清晰的圖畫.
神怎樣去成就你的大使命.
其實初期教會的信徒.
他們都是在不知不覺裡.
將福音帶到更遠.
不同的地方.
很多時候他們都不是.
透過什麼計劃.
或者知道自己要去什麼地方.
就去做出來.
而是跟隨著聖靈的帶領.
就好像在耶路撒冷的信徒.
當他們聚集和組織穩定之後.
聖靈又帶領他們開始新的工作.
支持著保羅和其他的信徒.
去到不同的地方.
做傳福音的工作.
保羅帶著幾個信徒.
離開耶路撒冷.
去到不同的地方去傳福音.
隔一段時間.
他們又會回到耶路撒冷.
在門徒的燈中.
講解給他們聽.
所發生的事.
所遇到的事.
就將他們所遇到的難題.
帶回來.
跟大家一起去分享去討論.

$^{201}$就例如.
外邦人需不需要接受國禮.
或者對於拜偶像的食物.
我們去問他們可不可以吃.
透過這些日常生活的點滴.
一些很真實生活的基礎.
讓大家看到聖靈怎樣工作.
聖靈是怎樣帶領著他們.
將在《少林寺第一章》第八節.
所講的話.
一步一步實踐出來.
神承就大使命.
有不同的方式.
其中一種是聖靈直接的帶領.
一個特別的帶領.
另一種.
就是透過教會穩定之後.
和宣教士互相結練.
互相配合.
互相激勵.
一同造就.
一起成長.
時至今日.
傳福音的方式已經有很多不同.
有很多不同的變化.
但基礎也很雷同.
就好像過往海外的宣教士.
他們和教會或組織合作.
來到亞洲,香港.
去到不同的地方.
做不同的福音工作.
按照支持的教會或組織.
和宣教士的恩賜.
每一天開始.
有不同的事奉.
例如醫院.
學校,兒童院.
做一些父母的工作.
一些社會的工作.
然後慢慢建立教會.

$^{241}$去實踐主耶穌基督的教導.
所以無論我們在本地.
去做傳福音,靜堂.
或者海外宣教的工作.
我們的宗旨都是一樣.
目標都是一致.
向未信主的人傳福音.
讓他們懂得福音的好處.
其實我們做本地的福音工作.
也不容易.
也遇到很多困難.
但因為我們在這裡的資源.
比較快,比較多.
所以很快又可以繼續和發展下去.
有時我們也有一個工作.
有時我們也有一個疑問.
既然本地的福音工作.
有這麼多.
做的也不完.
為什麼我們還要去做一些.
海外更加難的.
傳福音的工作呢?.
這個疑問.
當年我們也曾經向那些.
來到我們華人世界.
傳福音的海外宣教士.
發出過.
在羅馬書第十章.
十四到十五節.
普羅就這樣說.
人未曾相信神.
又怎可以向神求告呢?.
人未曾聽過神.
又怎能相信祂呢?.
沒有傳道的又怎能聽呢?.
沒有人奉差遣.
又怎能有人去傳道呢?.
這是普羅他自己個人的領袖.
向未信主的人.
未聽過福音的人去傳福音.

$^{281}$就在那時.
耶路撒冷的領袖彼得.
也從聖靈領受了.
一個新的意象.
開始支持.
外邦人的傳福音工作.
我們要知道.
當時猶太人.
要接受向不潔淨的外邦人.
傳講他們所相信的神.
是一個很難接受的事情.
打破了他們的思想.
和行為.
但這就是聖靈的工作.
聖靈的工作連結了教會.
和宣教士.
讓他們成為夥伴.
大家互相支持.
互相配搭.
成為前線後勤.
一個很好的配搭.
就正如普羅和耶路撒冷教會一樣.
聖靈的工作.
很多時候不是我們全部都能夠明白.
理解得到.
很多時候.
即使解釋了給我們聽.
我們都理解不到.
甚至接受不到.
所以當教會和宣教士.
大家都在聖靈裡面.
領受了相同的意象.
大家願意一起互為配搭.
一起可以去實踐傳福音的時候.
福音的大使命.
就可以走得更加遠.
去到更加不同的地方.
神的心意.
是人埋人得救.
無論在本地或海外的福音工作.

$^{321}$都願意我們信徒.
不斷去傳福音.
有些教會.
他們領受支持海外的猜拳工作.
這個意象和呼召比較早.
所以他們會有猜拳的事工.
有些教會.
他們去領受和參與比較遲.
但無論什麼情況下.
都是有神自己的心意和安排.
2011年.
是我加入.
感受堂祖傳道同工的時候.
當時教會只是密室一個地方.
希望可以發展直堂的工作.
教會弟兄姊妹都很開心.
大家一起祈禱.
雖然不知道在什麼地方可以做直堂.
但都是滿心開心的去祈禱.
求密神.
有些弟兄姊妹也很有心.
拿了在元朗附近.
日後將會發展的地區.
大家一起祈禱.
一起求密.
當我們去接觸到.
紅恩堂.
現在這個堂祠.
大家聚會的地方的時候.
大家都很開心.
一個教宗這麼方便.
有這麼大的地方.
不過在高興之餘.
又發現其實我們面對很大的難題.
例如業權的問題.
經濟的問題.
人手的問題.
人流的問題.
等等等等.
當時錦繡堂的弟兄姊妹.

$^{361}$因為經歷過在過去三十年.
神怎樣帶領錦繡堂的弟兄姊妹.
去建立教會.
現在我們數回四十年前.
錦繡花園是一個很荒蕪.
沒有交通工具可以直達的地方.
但是當時神就藉著.
宣教士在這裡聚集上.
幾個人開始.
大家組織小組.
然後一起去辦幽士院.
一起去建立教會.
因為大家經歷過神的帶領.
經歷過傳福音的好處.
得到神的恩典和祝福.
於是大家憑信心.
我們一起去發展紅一代.
在紅一代正式成立啟用禮之前.
我還記得帶著幾個年輕人小組.
來到這裡.
在樓下門口追著氣球.
在街邊派單槍.
邀請小朋友來參加活動.
唱詩歌.
玩遊戲.
吃茶點.
除了我們之外.
還有團體的弟兄姊妹.
來到教會裡面打碎地方.
大家很希望這個地方可以發展.
早日可以傳福音.
雖然都是一些默默.
微小的工作.
但是大家都很開心地事奉神.
十年過去了.
紅印堂在神的祝福.
牧師 傳道同工 執事.
和各位弟兄姊妹的努力付出之下.
走到今天.
有很多神的恩典.

$^{401}$有很多神的祝福.
無論在人數上.
或者在資源上.
如果十年是一個階段.
在下一個階段.
相信紅印堂會繼續努力.
尋求神的心意和更多儀仗.
除了在本地繼續努力傳福音.
會不會大家一起又會再踏一步.
打開支持海外傳福音的門呢?.
回顧在疫情之前.
我有這樣帶著一個小隊的弟兄姊妹.
在街頭派單槍.
在公園和小朋友 家長.
一起唱詩歌 講聖經故事.
一起玩遊戲.
只不過地方是一個語言和文化.
都不同的另外一個地方.
2018年我成為.
感受堂差全部支持的宣教士.
因為以前我曾經在德語的國家學習過.
所以我又回到德語地區.
開始做一些華人的福音工作.
而我的宣教工場.
就在歐洲中部一個很小的國家.
布申堡.
2020年代.
我們傳福音的方式.
真的沒有想過.
又好像回到以前一樣.
我第一次在人群之中派單槍.
靜地而坐.
拿著一張大的詩歌紙.
一起唱詩歌.
講聖經故事 講福音故事.
是沒有麥克風大聲地講的時候.
就是很多年前.
我讀神學.
每一個星期都跟著一個美國的宣教士.
去到當時香港的越南難民營.

$^{441}$做神學實習的時候.
幾十年過去了.
無論在香港 洪水橋 盧森堡.
或者其他不同的地方.
都還有很多沒聽過福音的人.
等待著我們去傳福音給他們聽.
雖然方法是神救.
資源又不夠.
但是神都有你自己的工作.
就拿我們盧森堡這裡的宣教工作為例.
我們開方隊只有幾個家庭.
幾個弟兄姊妹.
在我們都還沒有接觸到福音對象之前.
弟兄姊妹一起參與 一起準備的時候.
已經經歷到神給我們很大的喜樂.
很大的動力.
大家很開心 一起準備 事奉神.
神更加燃燒起.
弟兄姊妹在盧森堡生活了二十幾三十年.
從未有過對生命的意義和價值.
大家開開心心.
在有限的資源裡面去為神傳福音.
就連住在護理院.
長期病患 行動不方便的弟兄姊妹.
都很想和我們一起參與 做探訪的工作.
去到其他院舍.
探訪同樣都是長期病患.
要住院的院友.
鼓勵他們 向他們傳福音.
所以資源不夠.
我們的難度也有限.
甚至我們的奉獻.
連支持每一個星期聚會的.
場地中心和支出都不夠.
但是這樣有什麼問題呢?.
我們都知道在能力有限的裡面.
我們仍然可以做的事情.
這些就是聖靈的工作.
聖靈的工作真的不是人.
可以傳言 可以理解.

$^{481}$只能夠領受其中的一二.
我們相信只要跟著聖靈的帶領.
大家盡力發揮神所賜給我們的恩賜和資源.
開開心心努力事奉神.
神就會承受他的工作.
盧森堡的宣教工作.
困難比我們想像中的更多.
雖然盧森堡是歐洲最富裕的國家其中的一二.
但是很多資源都不是外國人入籍後可以使用的.
僑居在這裡的華人和下一代.
因為語言 文化 背景的限制.
在生活上都出現了很多的問題和困難.
我們不說其他了.
但是就說語言的問題.
在盧森堡我們有三種語言.
盧森堡文 德文和法文.
這個也是小朋友在學校必定要學習的語言.
在現實生活裡面.
在不同的場景.
在日常生活不同的場景.
這幾種語言都會同時間或者交替出現.
好像當你離開家 踏出家門.
盧森堡文的鄰居或外國人.
當我們看醫生的時候.
醫生只會說德文.
而護士只會說法文.
所以在很基本的日常生活的溝通裡面.
我們都要面對很大的挑戰.
很大的難處.
甚至是一個解不開的結.
就因為這些艱難.
就因為這些困難.
我們弟兄姊妹經歷了處境裡的無助.
那種孤單.
不知道怎麼做 無可適從.
所以更加覺得很需要有神的幫助.
有聖靈教給我們的能力.
可以繼續積極努力生活下去.
困難越多.
弟兄姊妹更加看到.

$^{521}$人對福音的需要越大.
雖然我們的人數不多.
資源也有限.
但我們都有意象和心智去事奉神.
我們希望首先在盧森堡穩定我們的聚會.
建立一個穩定的全福音基礎.
然後將我們的福音帶到鄰近的國家.
德國和法國的華人.
我們也求神給我們這樣的心智.
有一天將我們在盧森堡所建立的全福音基礎.
可以帶到附近其他地方.
不單傳給華人和下一代.
也傳給華人和外國人通婚的下一代.
起初保羅親自從聖靈接過宣教的棒.
無論他隻身上路或與教會合作.
福音就是這樣傳了出去.
宣教使徒的福音由耶路撒冷傳到歐洲.
傳到美洲 傳到亞洲 傳到非洲.
我們不知道什麼時候會傳到地球.
能夠建立穩定的教會.
這是神的恩典 固之然是好.
能夠再下一步與宣教士一起合作.
將福音傳到更遠的地方.
這是神的心意 也是神的計劃.
是神交給我們傳福音的大使命.
宣教的議長可以是個人的理事.
也可以是教會的議長.
大家能夠共同合作.
互相連結 互相配搭.
前線 後防 互相支持.
是一件合神心意的美事.
我們看回初期教會的弟兄姊妹.
他們的信仰其實都很艱難.
他們都面對過大迫害的時候.
就好像保羅在未信主之前.
他都曾經大力迫害相信主耶穌的人.
但是有天誰想得到.
甚至連他自己都想不到.
有一天他自己會成為.
傳福音一個最大的宣教士.

$^{561}$所以神的工作不是我們可以理解.
也不是我們可以去想像.
而是神透過聖靈的工作.
帶領著個人 帶領著教會.
互相合作 成就祂的大計劃.
成就祂的大使命.
前線傳福音的工作.
很需要後防 教會在土共.
在經濟各方面的支持.
兩者互為配搭.
更加是經歷神給我們.
共同承擔宣教使命.
保羅傳福音的經驗非常豐富.
他去到每一個不同的地方.
都在那裡傳福音.
帶領人信主 建立教會.
其中他去過哥林多這個地方.
在那裡也建立了教會.
在哥林多前書第一章那裡.
他這樣說.
他說我純純為了哥林多的弟兄姊妹.
因為主耶穌基督賜給他們的恩典.
非常之多.
哥林多的信徒非常之富足.
而且他們知識厚重.
樣樣全密.
甚至每一樣他們都不及其他人.
但是在這樣的情況下.
保羅也要再提醒他們.
我們傳福音.
我們建立教會.
是一個不可或缺的要素.
就在哥林多前書第十三章.
在那裡這樣說.
保羅這樣說.
他說我若能說萬人的方言.
天使的華語.
給滿有愛.
我就成了明的樂.
響的筆一般.

$^{601}$他說我若有先知講道之能.
也都明白各樣的奧秘.
各樣的知識.
而且有全密的信.
叫我能夠以心.
但是如果沒有愛.
我也都不算不得神愛.
所以無論傳福音.
宣教.
建立教會.
保羅都提醒我們.
我們要將主耶穌基督的愛.
成為我們的基礎.
我們要傳揚的.
是主耶穌基督的救恩.
藉著我們的生命.
藉著我們的行為.
去表達神的愛.
去將神的愛彰顯出來.
讓聖靈在我們身上工作.
改變我們.
幫助我們.
這個就是主耶穌給我們的.
大使命的內容.
我們去傳揚主耶穌基督.
不是只是用口去講.
而是用我們實際的行為.
要用主耶穌給我們的愛.
要用聖靈在我們裡面.
激發我們的能力.
讓人看到.
這個就是相信主耶穌基督的生命.
所以今天無論我們.
在哪裡傳福音.
去哪裡宣講神的訊息.
我們都不要忘記.
用愛來表達神的使命.
最後再一次多謝老穆星爺.
邀請我在這裡有機會.
和大家分享宣教的事服.

$^{641}$今天的時間不是很多.
能夠和大家分享我們.
在盧森堡宣教的工作.
也不是很詳細.
希望日後有機會.
可以和大家更詳細.
分享我們在這裡的事服.
也容許我在這裡.
邀請大家為我們在這裡的.
宣教工作來禱告.
希望藉此.
讓大家知道.
我們在這裡宣教的意象.
主有許可.
或許有一天.
大家會成為我們.
一起合作的宣教夥伴.
我們一起在這裡.
做宣教的配搭.
求神帶領我們.
幫助我們.
願主耶穌基督的平安.
恩惠.
慈愛.
常與大家同在.
恩運 慈愛 常與大家同在.
\newpage



\section{阿摩司書 5:1-17}
\label{sec:QekXUaNsm08}
\textbf{2022.09.25 《尋求耶和華者必存活》 阿摩司書 5\_1-17 (和修版) 陳麗蟬牧師}
\newline
\newline
連結: \href{https://youtube.com/watch?v=QekXUaNsm08}{\texttt{https://youtube.com/watch?v=QekXUaNsm08}} ~~~~ 語音日期: 2022-09-29
\newline
\newline
\hyperref[sec:u7aPgc4mo0Y]{\small{< < < PREV SERMON < < <}}
~
\hyperref[sec:index_chronic]{\small{[返順時目]}}
~
\hyperref[sec:index_scriptual]{\small{[返順卷目]}}
~
\hyperref[sec:bLMot_knMH4]{\small{> > > NEXT SERMON > > >}}
\newline
\newline
阿摩司書 5:1-17
\newline
\begin{longtable}{cl}
\hline
\hline
章節 & 經文 (和合本修訂版)\\
\hline
5:1 & \begin{tabularx}{0.7\textwidth}{X} 以色列家啊，聽我為你們所作的哀歌： \end{tabularx} \\ \\ \relax
5:2 & \begin{tabularx}{0.7\textwidth}{X} 「以色列民跌倒，不得再起；躺在地上，無人扶起。」 \end{tabularx} \\ \\ \relax
5:3 & \begin{tabularx}{0.7\textwidth}{X} 主耶和華如此說：「以色列家的城派出一千，只剩一百；派出一百，只剩十個。」 \end{tabularx} \\ \\ \relax
5:4 & \begin{tabularx}{0.7\textwidth}{X} 耶和華向以色列家如此說：「你們要尋求我，就必存活。 \end{tabularx} \\ \\ \relax
5:5 & \begin{tabularx}{0.7\textwidth}{X} 不要往伯特利尋求，不要進入吉甲，也不要過到別是巴；因為吉甲必被擄走，伯特利必歸無有。」 \end{tabularx} \\ \\ \relax
5:6 & \begin{tabularx}{0.7\textwidth}{X} 要尋求耶和華，就必存活，免得他在約瑟家如火發出，焚燒伯特利，無人撲滅。 \end{tabularx} \\ \\ \relax
5:7 & \begin{tabularx}{0.7\textwidth}{X} 你們這使公平變為茵蔯，將公義丟棄於地的人哪！ \end{tabularx} \\ \\ \relax
5:8 & \begin{tabularx}{0.7\textwidth}{X} 那造昴星和參星，使死蔭變為晨光，使白晝變為黑夜，召喚海水、使其傾倒在地面上的，耶和華是他的名。 \end{tabularx} \\ \\ \relax
5:9 & \begin{tabularx}{0.7\textwidth}{X} 他快速摧毀強壯的人，毀滅就臨到堡壘。 \end{tabularx} \\ \\ \relax
5:10 & \begin{tabularx}{0.7\textwidth}{X} 你們怨恨那在城門口斷是非的，憎惡那說正直話的。 \end{tabularx} \\ \\ \relax
5:11 & \begin{tabularx}{0.7\textwidth}{X} 所以，因你們踐踏貧寒人，向他們勒索糧稅；你們雖建造石鑿的房屋，卻不得住在其內；雖栽植美好的葡萄園，卻不得喝其中所出的酒。 \end{tabularx} \\ \\ \relax
5:12 & \begin{tabularx}{0.7\textwidth}{X} 我知道你們的罪過何其多，你們的罪惡何其大；你們迫害義人，收受賄賂，在城門口屈枉貧窮人。 \end{tabularx} \\ \\ \relax
5:13 & \begin{tabularx}{0.7\textwidth}{X} 所以智慧人在這樣的時候必靜默不言，因為這是險惡的時候。 \end{tabularx} \\ \\ \relax
5:14 & \begin{tabularx}{0.7\textwidth}{X} 你們要尋求良善，不要尋求邪惡，就必存活。這樣，耶和華－萬軍之神必照你們所說的與你們同在。 \end{tabularx} \\ \\ \relax
5:15 & \begin{tabularx}{0.7\textwidth}{X} 要恨惡邪惡，喜愛良善，在城門口秉公行義；或者耶和華－萬軍之神會施恩給約瑟的餘民。 \end{tabularx} \\ \\ \relax
5:16 & \begin{tabularx}{0.7\textwidth}{X} 因此，主耶和華－萬軍之神如此說：「在一切的廣場上必有哀號的聲音；在各街市上必有人說：『哀哉！哀哉！』他們叫農夫來哭號，叫善唱哀歌的來舉哀； \end{tabularx} \\ \\ \relax
5:17 & \begin{tabularx}{0.7\textwidth}{X} 各葡萄園都有哀號的聲音，因為我必從你中間經過。」這是耶和華說的。 \end{tabularx} \\ \\
[1ex]
\hline
\hline
\end{longtable}
$^{1}$(影片開始).
弟兄姊妹平安.
平安.
感恩 今天又可以跟大家一起.
以上帝的說話來彼此勉勵.
我在灣仔住.
在地鐵站對面那條街很有名的.
叫做玩具街.
有很多玩具賣的.
有一個情況我留意了好幾年了.
街尾那裡有一個攤檔.
那個是賣人字拖鞋的.
我記得幾年之前我走過的時候.
我看到那個老闆娘也挺兇的.
整天在罵人.
有一次我又聽到她大聲罵一個主婦.
她說買鞋買鞋試了就不要換了.
因為你要換就不要星期天早上來找我.
星期天早上我要回教會.
要換就星期天下午才來換.
講價講這麼多.
我們信耶穌的人怎麼會騙人呢.
所以跟我買東西不要跟我講價.
我看了一下.
她在那個攤檔旁邊又放了一張摺台.
鋪滿了很多中信的月刊.
寫著是免費取月的.
我覺得.
啊.
最初我看錯人了.
覺得這個人兇得要命的樣子.
誰知道原來她信耶穌.
可能她不懂得傳福音.
又放一些中信月刊在那裡.
讓人隨意拿.
我覺得這個老闆娘也不錯.
然後過了幾個星期之後.
就是10月了.
我又走過那個攤檔.
這次嚇了一跳.

$^{41}$因為10月是什麼月呢?.
萬聖節.
於是整個攤檔就不是賣人字拖了.
就是賣了萬聖節的怪怪的用品.
所以今天要講這篇道.
我前幾天又去看了.
還沒踏入10月已經在賣了.
我想一下.
沒錯.
基督徒也是要賺錢的.
上帝不想我們成為窮光蛋.
但是呢.
當我們用我們的信仰.
你成為了一個裝飾.
你們最好不要來跟我講價.
因為我信耶穌.
我不騙人的.
我覺得.
但是你又另一方面.
又去買這些東西的時候.
我覺得你又在做一些失現性的事情.
其實我們整個信仰裡面.
我們需要的.
需要的一件事情.
就是要言行都要一致的.
因為當我們所講和我們所信的事情.
都是不同的時候.
神一定會追討.
一定會來審判的.
剛才讀了長長的17節的阿摩斯書.
我們先認識一下阿摩斯.
阿摩斯的時代.
已經分了北國和南國.
北國就叫以色列國.
南國就是猶太國.
當時北國是一個挺繁盛的地方.
阿摩斯才知道他住在南邊南國的地方.
但是上帝不知道.
很有趣.
就猜他去到北國裡面.

$^{81}$跟他宣告審判.
所以其實有難度在這裡.
因為當時的以色列國挺繁榮的.
但是先知來就罵罵罵.
其實挺難做的.
不過在這個社會裡面.
充斥了很多不公義的事情.
那些惡人是很惡的.
有錢人也去欺凌很多窮人.
阿摩斯才知道.
他在第一節就說了.
他說他要以色列要聽他所作的哀告.
到了第十六節.
剛才我們有讀經文.
我們聽到說.
哀哉哀哉.
周圍寬闊之處廣場的地方.
都有人在哭.
甚至覺得不夠人哭.
連原本去耕種的農夫也拉過來哭.
說到全國都陷入一個很淒涼的境況.
一般來說.
不是一般來說.
是直接的.
哀歌是什麼時候唱的?.
人死了之後才唱的.
是不是?.
但是在這個時間.
阿摩斯就已經說了.
他要為以色列去唱哀歌.
他是在說預言.
他說這個國家將要滅亡.
還要慘的是.
他形容到以色列國.
就好像一個很年輕的女孩.
還沒結婚的女孩.
她要面對死亡.
甚至她要跌倒在地上.
沒有人去幫她.
他在描述將來以色列國.

$^{121}$他說以色列國.
因為他被上帝所揀選的.
是上帝跟他們立約的一個特別的國度.
所以本來以色列民是一個很有光彩.
很有盼望.
很有力量的民族.
但是現在呢.
將會哀聲四起.
為什麼呢?.
因為上帝在出埃及的時候.
在西賴山的腳底.
他就跟這群以色列民納了約.
在出埃及裡面這樣描述.
他說你們若真的聽從我.
就說回以色列民.
遵守我的約.
你們就在萬民當中作我的指紋.
因為全地都是我的.
你們聽我的話.
遵守我教你們的東西的時候.
你們就可以成為我的指紋.
你們成為我的指紋的時候.
上帝就會保護他們.
而這群指紋又如何呢?.
為上帝屢作祭司的國度.
有一個特別的身份.
榮耀的身份.
為上帝作一個神聖的國民.
但是呢.
這群以色列人沒有這樣做.
他們得罪了上帝什麼呢?.
第一方面.
他們在人與人之間.
他們沒有守在第十誡裡面.
第五至第十誡裡面.
上帝所教導他們的.
第五誡一直到第十誡.
當孝敬父母.
不可殺人.
不可姦淫.

$^{161}$不可偷渡.
不可作假見證.
不可貪戀你.
淪捨一切的所有.
但是呢.
在剛才我們讀了大段經文.
大家都不知道讀了之後有沒有印象了.
就才在罵.
罵那群以色列人什麼呢?.
在城門口.
在這裡.
這個城門口就是在舊約的時候.
他們來去做審判的地方.
在城門口.
就是這個法庭裡面.
他們本應就一定要講誠實話.
講真話的.
但是呢.
竟然不准人講真話.
那就是法庭也沒用了.
是不是?.
還有.
他們要踐踏這個貧民.
那麼他們辛辛苦苦種了那些東西出來.
他們就去拿了它們.
他們又收受賄賂.
在這個城門口.
又是說在城門口.
就冤枉那些窮人.
他們對上帝的話完全沒有遵守.
又剝削窮人.
是不是?.
又收賄賂.
這幾天.
我們都看到這個新聞.
就說.
在前北京的公安局長.
在2005年.
到2021年.
總共16年.

$^{201}$你知不知道他非法貪污了多少錢?.
1.17億.
我也不懂得數.
1.17億.
昨晚看報紙.
又是在國內的.
另一個人收了多少億呢?.
6億.
不是說他賺錢.
而是單單收賄賂.
哇.
這麼多億.
在當中.
多少人受害.
這幫人.
全部都是判死刑.
不過就緩期執行.
還有在香港.
香港這幾天又有了.
我們看到.
現在很喜歡打針.
我不想打針.
那怎麼辦?.
就去排隊.
排通宵隊.
拿錢就可以買到免針紙.
是不是?.
本來那些醫生.
應該是很尊貴的.
賺錢也賺很多.
但是呢.
為了賺多一點的時候.
就亂發免針紙.
我們看到整個世界.
這個世界就是這樣.
當時以色列國的時代更厲害.
但是上帝說.
上帝一定會審判的.
審判怎麼樣呢?.
他說.

$^{241}$你們那些人自己為自己建這個堅固的房屋.
做了也沒得住.
因為戰爭快到了.
還有.
他建的房子很漂亮的.
通常一般平民百姓就用泥磚來建的.
那些有錢人.
才用石頭來做的.
建了也沒得住.
還有.
整個國家很慘.
出去打仗的.
出一千.
就只有一百人回來.
出一百.
只剩下十個人回來.
所以呢.
整個國家是因為你們自己沒有守上帝的約.
沒有守上帝的約.
所以上帝就懲罰了.
上帝就不保護了.
除了這樣之外.
他也沒有守到.
十誡裡面第一至第四誡.
講到人和神之間的關係的約.
在第一至四誡裡面講到.
除耶和華之外.
不可有別臣.
也不可以為自己雕刻偶像.
不可以忘稱耶和華的名.
當紀念安息日.
好了.
那我才提醒他們.
提醒他們.
他們呢.
沒錯你們有敬拜.
你們有獻祭.
但是呢.
在這個敬拜.
在這個獻祭裡面.

$^{281}$他們加入了很多那些愛邦神的拜祭裡面.
那些淫亂的儀式.
所以呢.
在第五節他說.
你們不需要去伯特利裡面尋求.
也不需要去吉格.
也不需要去別士巴.
不需要去那些地方裡面敬拜.
因為呢.
這三個地方.
在舊約的時代.
都是一些很神聖的地方.
是不是?.
就好像我們現在.
你說一說去拜神.
那些人第一就去哪裡呢?.
黃大仙.
是一些很神聖的地方.
在那時候.
那些舊約的人.
想想覺得神聖一點的地方.
伯特利.
就是這個.
當這個雅各.
被他哥哥追殺的時候.
是不是?.
有一天.
他睡覺.
在曠野睡覺.
他就做夢.
看到有一條天梯.
有使者在那裡.
又上又下.
上帝來去建立他.
上帝來保護他.
所以呢.
他就說有一天.
我有命回來這個地方.
我就會在這裡.
捉壇來敬拜.

$^{321}$後來他真的平安回來了.
他就在那裡捉了這個壇.
所以以色列人很重視這個地方.
第二個就是吉格.
吉格就是約書人.
他就帶領以色列人.
在曠野地.
走了40年.
進入迦南地.
就好像那時候出埃及的時候.
經紅海.
這次就過約旦河.
過約旦河的時候.
因為當時差不多有200多萬人過約旦河.
雖然約旦河很小.
好像元朗那條.
元朗河仔.
但是200人過.
還有水漲時間是很危險的.
但是當那些祭司.
抬著玉櫃的祭司.
一伸進這個約旦河.
那些河水就上下分開了.
又是那些以色列人.
很深很深經歷的一個地方.
所以又是一個很出名的地方.
別氏巴.
別氏巴就是以撒.
因為挖到水漲.
那些非利士人就愛上他.
好幾個都不行.
後來到了別氏巴這個地方的時候.
上帝就向他顯現.
告訴他.
上帝一定會保護他.
所以別氏巴裡面.
他挖的水漲.
非利士人再沒有來搶了.
並且他和非利士的領袖阿比米勒.
來簽訂了和平的條約.

$^{361}$又在那裡祝壇敬拜了.
所以這三個地方很神聖.
他們很喜歡這個地方.
用來去獻祭.
敬拜上帝的.
但是上帝說.
你來這些地方.
你做足儀式.
但是沒有用的.
只有外表.
你一手拿著十個十字架.
另一手就拿著這些偶像.
繼續在行為上參與他們淫亂的敬拜.
所以上帝一點都不越立他們所做的一切.
我教中學的時候.
有一個男同學.
鍾義.
水上人.
已經習慣了.
久不久他就不上學了.
因為他要和家人出海打魚.
有一次他乖乖的.
之前就來跟我告假.
Mr. Chan.
我下個星期要告假一個星期.
又和家人去打魚.
他說這次不是.
我去長洲.
我會帶著十字架.
帶著蒜頭.
我要去長洲捉鬼.
我說.
你都還沒有信耶穌.
你捉什麼鬼?.
他說.
感恩你平安回來.
所以鬼沒有找你.
或者我們知道有時候.
當我們說要去泰國旅行.
有沒有聽過?.

$^{401}$去泰國旅行你進酒店.
看到酒店有本聖經.
一蓋上你就不用怕.
一打開你就要怕了.
因為一定是很猛的.
我說.
他們那些人不信的人覺得.
蓋上聖經那些字就發不了力量.
但是打開的時候.
那些力量就會威猛一點出來.
是不是?.
我就說.
所有這些東西.
都是物品.
無論你十字架,聖經什麼都好.
當你不信.
當你不是屬於耶穌基督的人的時候.
這些東西.
全部都沒有這個屬靈果效出來的.
所以.
阿摩斯說.
你不需要去任何這些高級,屬靈的地方.
去找我.
去敬拜我.
最重要的.
是行出我給你們的教導.
所以.
我說我們都是屬於上帝的子民.
我們需要去學上帝的聖經.
進行上帝的教訓.
然後這樣才能得以存活.
我們是信基督.
不是信基督教.
所以剛才我老闆娘.
我不知道她是真的信耶穌還是假信耶穌.
我覺得她只是有宗教傳福音的行動.
但是沒有實際的屬靈的生命.
所以.
這一幫以色列人.
被先知罵得很慘.

$^{441}$整張都是哀歌.
牠帶牛,帶羊去獻祭是沒有用的.
他們這樣的行為.
反而污穢了上帝的聖潔.
我們需要.
在我們的信仰和生活上.
一定有所結連.
千萬不要聽一套.
說一套.
而做的又是另外一套.
第二方面.
上帝說.
要怎麼做呢?.
要做什麼呢?.
在第八節裡面.
要尋求那做.
我的鄉下話不太行.
「卯星」.
要尋求那做卯星和參星.
使死蔭變為神光.
使白日變為黑夜.
命海水來囂在地上的神.
耶和華是他的名.
在這裡.
就是說要去找.
找回做卯星和參星的上帝.
因為以色列人.
他就跟隨了亞述人.
去拜卯星和參星.
他們覺得這些星神.
是掌管著整個宇宙的運行.
所以他們就拜這些星座.
不過.
阿摩斯就提醒他們.
不是拜這些星座.
是拜創造這些星.
這些星座.
那位上帝.
尋求的意思就是說.
他們要回去.

$^{481}$從原本拜這些星座.
他們回去拜上帝.
所以尋求的意思就是說.
他要回轉.
回去倚靠創天造地的上帝.
因為有耶和華.
才是真正的創造者.
不過他不僅創造宇宙萬物.
他也創造我們人類.
他也管理我們人的心靈.
他的能力有多大呢?.
創造卯星和參星的上帝.
他的能力有多大呢?.
他使得死刃變為神光.
當我們在黑漆漆的地方.
我們感到毫無盼望.
我們感到死了.
沒有希望了.
但是.
使得人在失望和絕望的時候.
他能夠帶他走出黑路.
只要他回轉.
去倚靠創天造地的上帝的時候.
我們的人生.
能夠從絕望.
沒有機會.
變為盼望.
充滿了光明的生命.
他的能力就是使白日變為黑夜.
使海水能夠淹沒整個地.
這裡是在說什麼呢?.
就是在說他的審判.
說到原本是光天化日的.
突然間黑沉沉.
突然下大雨.
淹沒整個大地.
在說舊約.
哪裡的審判呢?.
羅亞.
就是當時的人.

$^{521}$每個人都完全不聽話.
每個人都喜歡做什麼就做什麼.
上帝不是叫羅亞去做方舟嗎?.
光天化日.
白日.
陽光普照.
傻子不斷在打大船.
怕有大雨.
是不是?.
但是呢.
上帝就是這樣.
那些人有信心去依靠的時候.
他就得逞夠了.
他就說上帝能夠使原本是陽光普照的.
就變為黑夜.
然後用洪水來審判大地.
在第五章.
阿摩斯·先知.
三次去說.
你們要尋求.
在第四和第六節他也這麼說.
耶和華向以色列加諭此說.
你們要尋求我.
就必存活.
第四和第六節都這麼說.
到了後面.
到了第十四節.
他更加說多一點.
你們要求善.
求量善.
不要求惡.
就必存活.
這樣.
耶和華萬君之臣.
必照你們所說的.
與你同在.
就是說.
我們去尋求祂.
迴轉祂.
依靠祂的時候.

$^{561}$我們就必得存活.
就必得存活.
求量善.
量善是什麼呢?.
量善就指我們的上帝.
祂是一個很高尚的品格.
是很純潔的.
還有.
所有的力量.
所有的智慧都集中在那裡.
其實一句話說完.
所有好的東西.
都在上帝裡面.
所有的好東西.
都在上帝那裡.
只要我們去尋求祂.
我們去找祂.
去依靠祂的時候.
這些東西.
都會傳過來給我們.
但是.
當我們擁有上帝的量善的時候.
我們又要將這些行善.
走出來.
因為上帝擁有這個量善.
祂自己走出來的時候.
就是去差派祂的兒子.
來到世間.
去救贖我們.
同樣.
我們得著這個量善的時候.
我們都要走出來.
去幫助其他人.
在這個中信第724期.
剛剛出了一兩期左右.
其中有一篇叫做.
《出黑暗入光明》.
這一位作者叫做陳卓君.
他是一個家庭.
一個重男輕女的家庭裡面去長大.

$^{601}$可能跟我差不多年紀的人.
很多家庭都是這樣的.
我經常都說.
你問一下.
特別是潮州人.
我說潮州人家裡是重男輕女的.
舉手.
腳都舉起來都可以.
是不是?.
哇.
我們那時候經常都記得.
小時候我很不服氣的.
我舅父給利是分三級制的.
第一級就是我哥哥.
第二級就是我的姐姐.
第三級就是我兩個弟弟和我.
永遠年年都是最低級的.
很不服氣的.
又或者出去喝酒.
我們永遠都不能去喝酒的.
只有我哥哥可以去喝酒.
覺得很不喜歡.
但是這位陳卓君更慘.
他就說.
他不僅如此.
他父母不讓女孩子讀書.
他不能讀書.
但是他很喜歡讀書.
所以他寧願白天的時候.
幫他父母做事.
晚上怎麼都堅持去讀書.
上帝讓他讀得很好.
而且成績比他弟弟好.
他竟然有這樣的父母.
他不喜歡女兒讀書.
比兩個兒子厲害.
有這樣的人.
然後出來工作.
他又做得很好的成績.
一直升職.

$^{641}$他說.
他在1992年的時候.
已經在一家保險公司裡面工作.
每個月三萬多塊的薪水.
那時候1992年三萬多塊.
很厲害的.
他媽媽又不喜歡.
可以這麼說.
總之不能比兒子厲害.
他說.
最後他的弟弟欠了很多錢的時候.
他父母又要逼他幫他還錢.
幫小孩還錢.
所以他心裡面.
對父母很生氣.
他雖然結了婚.
跟丈夫後來沒多久又離婚了.
他整個生命裡面.
充滿了很多怨氣.
也有很多的苦毒.
覺得為什麼就算是父母也好.
為什麼會這樣對親生女兒.
所以他太多負面的事情.
出來做保險這些事情.
我相信在當中很多的.
大家鬥,搶客各方面的那些.
他很辛苦.
他患了嚴重的抑鬱症.
再加上一次車禍裡面.
他差不多死在ICU那裡.
他就經常都發惡夢.
就覺得自己不知道去到哪裡.
那裡黑漆漆的.
很害怕很害怕.
但是有一晚.
上帝跟他說話.
他說我的恩典夠你用.
因為我的能力.
在人的軟弱上高.
顯出完全.

$^{681}$我的恩典夠你用.
因為我的能力在人的軟弱當中.
顯出完全.
當他聽到這個在小學還是中學的時候.
在教會學校裡面讀到的一句聖經.
他覺得被支持.
過往父母都不疼他,不支持他.
說出來的那些同事是大家勾心鬥角的.
但是在當中他覺得有那份的惡.
所以當他後來出院之後.
他就開始回教會.
他就慢慢將過往很多的仇恨.
爭權.
那些生命令他整個很辛苦的事情.
他慢慢來去改.
因為直至透過一直讀聖經的時候.
他發覺到過往.
他又覺得這個人又壞.
這個人又欺負他.
那個人又自私.
那個人又不知道怎樣.
於是他覺得每個人都壞人自己最好.
但當聖經來光照他的時候.
他發覺到他自己是一個很嫉妒的人.
是一個很驕傲的人.
也是一個很貪心的人.
上帝就將他自己那個幽暗的生命.
他不單止覺得人壞.
他原來化到自己比人更差.
上帝就在那裡光照他.
他回轉,他倚靠.
之後他洗禮了.
他洗禮之後上帝再光照他.
上帝就光照他了.
在剛才的十誡.
但是當孝敬父母.
你做了多少呢?.
我覺得爸媽那時候那麼壞.
經常逼我.
又不讓我讀書.

$^{721}$我又那麼慘.
他以前對我那麼苛刻.
你叫我現在對他好一點.
很難吧?.
但就上帝的話.
一路光照他.
然後他帶了他父母信主.
很感恩.
我想他的生命就經歷了.
在這個原本的死暗裡面.
黑暗的裡面.
上帝就一路藉著話語.
來將他轉化.
變為這個光明.
變為這個有盼望的生命.
就好像在哥林多後書裡面說.
哪裡有主的聖靈.
那裡就有自由.
過往他只落在一個麻煙.
只落在詛咒其他人.
但到現在.
他就釋放了.
他跟人的關係.
特別是跟他自己父母的關係.
他釋放了.
所以他的生命不再像以前那樣.
污穢穢.
充滿了這個苦毒.
充滿了這個仇恨.
他在裡面.
在主的裡面.
有這個自由.
他經歷了生命那份的奉承.
所以這個.
阿摩斯才說.
要尋求主的良善.
不要尋求這個惡.
就必存活.
這樣.
耶和華萬軍之神.

$^{761}$就會與你們同在.
其實阿摩斯才說.
我們要得到上帝的良善的時候.
我們就要主動出擊.
同樣.
好像上帝那樣.
成為一個.
你去建立他人.
你去關心其他人.
耶穌基督給我們的就是.
你要盡心盡性.
盡義愛主你的上帝.
其次就是愛倫如己.
很喜歡一句話.
他說生命是神所給我們的禮物.
我們多謝上帝.
但我們怎樣去生活.
怎樣活出上帝的良善.
就是我們給上帝的禮物.
願我們生命上帝愉樂.
我們一起禱告.
天父上帝我們多謝你.
多謝主你的恩典.
主你給我們得著耶穌基督的救恩.
以致我們生命.
可以能夠從死刃地裡面.
我們去轉為白日.
我們多謝你.
主多謝你.
主給我們藉著耶穌基督的榜樣.
我們可以.
我們可以能夠去學習.
怎樣去愛周圍的人.
主祝福我們每一個人的生命.
我們說我們去尋求上帝的時候.
我們就活出上帝你的真善美.
加力量給我們.
多謝你垂聽禱告.
奉靠主耶穌基督得勝的名字而求.
阿 面.

\newpage



\section{馬可福音 5:21-43}
\label{sec:bLMot_knMH4}
\textbf{2022.10.02 《一個怎樣的信心》馬可福音 5\_21-43 (和修版) 老耀雄牧師}
\newline
\newline
連結: \href{https://youtube.com/watch?v=bLMot_knMH4}{\texttt{https://youtube.com/watch?v=bLMot\_knMH4}} ~~~~ 語音日期: 2022-10-03
\newline
\newline
\hyperref[sec:QekXUaNsm08]{\small{< < < PREV SERMON < < <}}
~
\hyperref[sec:index_chronic]{\small{[返順時目]}}
~
\hyperref[sec:index_scriptual]{\small{[返順卷目]}}
~
\hyperref[sec:KSHhb73QdyA]{\small{> > > NEXT SERMON > > >}}
\newline
\newline
馬可福音 5:21-43
\newline
\begin{longtable}{cl}
\hline
\hline
章節 & 經文 (和合本修訂版)\\
\hline
5:21 & \begin{tabularx}{0.7\textwidth}{X} 耶穌又坐船渡到對岸，有一大群人聚集到他身邊；他正在海邊。 \end{tabularx} \\ \\ \relax
5:22 & \begin{tabularx}{0.7\textwidth}{X} 有一個會堂主管，名叫葉魯，也來了，一見到耶穌，就俯伏在他腳前， \end{tabularx} \\ \\ \relax
5:23 & \begin{tabularx}{0.7\textwidth}{X} 再三求他，說：「我的小女兒快要死了，求你去為她按手，使她痊癒，可以活下去。」 \end{tabularx} \\ \\ \relax
5:24 & \begin{tabularx}{0.7\textwidth}{X} 耶穌就和他同去。有一大群人跟隨他，擁擠著他。 \end{tabularx} \\ \\ \relax
5:25 & \begin{tabularx}{0.7\textwidth}{X} 有一個女人，患了經血不止的病有十二年， \end{tabularx} \\ \\ \relax
5:26 & \begin{tabularx}{0.7\textwidth}{X} 在好多醫生手裡受了許多苦，又花盡了她所有的，一點也不見好，反而更重了。 \end{tabularx} \\ \\ \relax
5:27 & \begin{tabularx}{0.7\textwidth}{X} 她聽見耶穌的事，就夾在眾人中間，從後面來摸耶穌的衣裳， \end{tabularx} \\ \\ \relax
5:28 & \begin{tabularx}{0.7\textwidth}{X} 因她想：「我只摸到他的衣裳，就會痊癒。」 \end{tabularx} \\ \\ \relax
5:29 & \begin{tabularx}{0.7\textwidth}{X} 於是她的流血立刻止住，她覺得身上的疾病好了。 \end{tabularx} \\ \\ \relax
5:30 & \begin{tabularx}{0.7\textwidth}{X} 耶穌頓時心裡覺得有能力從自己身上出去，就在眾人中間轉過來，說：「誰摸我的衣裳？」 \end{tabularx} \\ \\ \relax
5:31 & \begin{tabularx}{0.7\textwidth}{X} 門徒對他說：「你看眾人擁擠著你，還說『誰摸我』呢？」 \end{tabularx} \\ \\ \relax
5:32 & \begin{tabularx}{0.7\textwidth}{X} 耶穌周圍觀看，要見做這事的女人。 \end{tabularx} \\ \\ \relax
5:33 & \begin{tabularx}{0.7\textwidth}{X} 那女人知道在自己身上所成的事，就恐懼戰兢，來俯伏在耶穌跟前，將實情全告訴他。 \end{tabularx} \\ \\ \relax
5:34 & \begin{tabularx}{0.7\textwidth}{X} 耶穌對她說：「女兒，你的信救了你，平安地回去吧！你的疾病痊癒了。」 \end{tabularx} \\ \\ \relax
5:35 & \begin{tabularx}{0.7\textwidth}{X} 耶穌還在說話的時候，有人從會堂主管的家裡來，說：「你的女兒死了，何必還勞駕老師呢？」 \end{tabularx} \\ \\ \relax
5:36 & \begin{tabularx}{0.7\textwidth}{X} 耶穌不理會他們所說的話，就對會堂主管說：「不要怕，只要信！」 \end{tabularx} \\ \\ \relax
5:37 & \begin{tabularx}{0.7\textwidth}{X} 於是他帶著彼得、雅各和雅各的弟弟約翰同去，不許別人跟著他。 \end{tabularx} \\ \\ \relax
5:38 & \begin{tabularx}{0.7\textwidth}{X} 他們來到會堂主管的家裡，耶穌看到一片吵鬧，並有人大聲哭泣哀號， \end{tabularx} \\ \\ \relax
5:39 & \begin{tabularx}{0.7\textwidth}{X} 就進到裡面，對他們說：「為甚麼大吵大哭呢？孩子不是死了，是睡著了。」 \end{tabularx} \\ \\ \relax
5:40 & \begin{tabularx}{0.7\textwidth}{X} 他們就嘲笑耶穌。耶穌把他們都趕出去，帶著孩子的父母和跟隨的人進了孩子所在的地方， \end{tabularx} \\ \\ \relax
5:41 & \begin{tabularx}{0.7\textwidth}{X} 就拉著孩子的手，對她說：「大利大，古米！」翻出來就是說：「女孩，我吩咐你，起來！」 \end{tabularx} \\ \\ \relax
5:42 & \begin{tabularx}{0.7\textwidth}{X} 那女孩子立刻起來走動—她已經十二歲了；他們就非常驚奇。 \end{tabularx} \\ \\ \relax
5:43 & \begin{tabularx}{0.7\textwidth}{X} 耶穌切切地囑咐他們，不要讓人知道這事，又吩咐給她東西吃。 \end{tabularx} \\ \\
[1ex]
\hline
\hline
\end{longtable}
$^{1}$主席.
分享經文就是剛才這一段.
和你們分享這段經文是有原因的.
因為我知道.
福音樂曲隊很想將這段血漏婦人的故事.
演繹成一套福音報道的話劇.
不過由於疫情的關係.
始終未能成事.
所以我特別選用這段經文去鼓勵他們.
讓他們更加明白故事想表達的神學意義.
在他們日後有機會演出的時候.
能夠更有效地呈現出故事裡的信仰精神.
不過由於上次分享的時間不太長.
而我覺得經文還有很多訊息.
可以讓我們反思.
於是我寫成了講章.
讓經文的屬靈意義.
能夠更加全面地呈現出來.
我們一起去祈禱.
讓我們能夠從你的話語裡面.
可以學習不同的屬靈教導.
從你的教導裡面.
可以明白你的心意更多.
好叫我們能夠按著你的心意而行.
讓你的名被高舉.
我們祈禱.
是奉教主耶穌基督得勝明智祈求.
阿們.
我們先了解一下經文的背景.
剛才大家都讀出了二十三節的經文.
《血漏婦人》的故事其實很短.
只有十節.
由二十五節開始到三十四節.
其實就完結了.
不過另外《血漏婦人》的故事.
是由兩組經文所包著.
分別是二十一節至二十四節的上文.
以及三十五節至四十三節的下文.
所以《血漏婦人》這個故事.
本身不是一個獨立的故事.

$^{41}$而是依附著上下文的一連串的事件.
色影學者將這種的寫作體材.
稱為三文治式的結構.
所以看《血漏婦人》.
這個故事不應該單獨去看.
還要看上下文的故事.
否則就錯過了.
作者這個三文治式結構的心思.
故事的上文是一個管會堂的業佬.
他再三請求耶穌去救他的女兒.
因為他的女兒就快死了.
在耶穌趕去業佬家中的途中.
就發生了血漏婦人的事件.
最後血漏婦人被治好了.
下文就說.
婦人被治好的同時.
業佬的女兒就死了.
耶穌應該不需要前去了.
不過耶穌沒有理會.
繼續前去.
他進入業佬的家趕出所有人.
最後將業佬的女兒從死裡拯救出來.
整個故事涉及了兩組人物.
兩件事情.
兩件事情都涉及了信心和行為.
無論是婦人的血漏病被治好.
還是女孩死而復生.
很明顯地讓讀者知道.
兩件事情都是一個神跡醫治.
業佬的行為很明顯.
他在眾人面前向耶穌再三請求.
顯示出他對耶穌醫治的信心.
婦人的行為是隱藏的.
他靜靜地走到耶穌後面.
沒有人知道婦人的企圖和信心.
兩件事情都指向不同的性別和身份.
業佬是一個男性.
有姓名.
是管會堂的.
是一個有身份.

$^{81}$有地位.
受人敬重的人物.
婦人呢.
是一個女性.
連名字都沒有被記載.
更加沒有身份和地位.
兩人的性別,身份,地位.
截然不同.
可以這麼說.
一個是有身份,有地位,有名聲.
受人尊敬.
另一個沒有身份,沒有地位.
根本是一個被社會遺棄的弱勢社群.
兩件神蹟被捆綁在一個沒有間斷的故事裡面.
究竟作者想我們明白一個甚麼訊息呢?.
我們開始進入經文.
管會堂的業佬.
他的女兒生病了.
雖然經文沒有說這個女孩生了甚麼病.
病了多久.
但經文說她快要死了.
女孩的生命與死亡相當接近.
一個有身份,有地位的業佬.
他應該找過很多醫生去醫治他的女兒.
不過始終都是徒勞無功.
他的女兒的病.
全無起色.
在這一刻.
業佬有一份絕望的感受.
他不知道可以為他的女兒做些甚麼了.
當他知道耶穌在海邊.
他又走過去要見耶穌.
雖然經文沒有交代到.
業佬為甚麼對耶穌那麼有信心.
但當我們看完三文治結構裡面.
中間那層的經文.
我們稍後就知道了.
他一見到耶穌.
俯伏在他腳前.
再三地求他.

$^{121}$我的小女孩快要死了.
求你去為她按手.
使她痊癒.
可以活下去.
業佬對女兒的病情的無助感.
束手無策.
這種感受推動到業佬去見耶穌.
當他見到耶穌的時候.
就請求耶穌的幫忙.
耶穌成為他在絕望裡面的唯一一個盼望.
耶穌對於業佬的請求也都欣然接受.
業佬的女兒雖然仍然未被醫好.
但業佬這一刻應該是充滿信心和盼望的.
因為耶穌答應他.
並且與他一起前往.
於是他們一起起程.
回到業佬的家裡.
就是想醫治業佬的女兒.
不過故事的發展卻被另一件事打斷了.
原來一個患了血漏的婦人.
亦都在人群裡面.
她患了血漏病十二年了.
十二年.
看過很多醫生.
很可惜都醫不好.
花盡了她所有的.
不但不能夠好轉.
病情反而重了.
我們想像一下這十二年.
婦人曾經有個盼望.
她每見一個醫生.
都有一個盼望.
希望那個醫生能夠醫好她.
所以她願意花錢.
她盼望有一個醫生可以醫好她的病.
很可惜.
一次又一次的失望.
財富沒有了.
病情卻越來越嚴重.
不單止這樣.

$^{161}$血漏是一個不潔的象徵.
所有跟她接觸過的人或物件.
都需要潔淨.
因此沒有甚麼人願意跟這個人接近.
婦人在日常的社交.
或者宗教生活裡面.
都是被排斥的.
可以這樣說.
她是一個被社會遺棄的弱勢社群.
婦人這十二年.
在這個情況下.
應該都是陷於一個絕望的心態裡面.
經文說.
她聽見了耶穌的事.
《馬可福音》第一章.
記載了耶穌醫好患麻風的病人.
在第二章.
耶穌又醫好了一個癱子.
第三章.
耶穌又醫好了一個.
枯萎了手的人.
到第四章.
耶穌平靜風和海.
第五章.
耶穌在格拉森趕出了二千隻污鬼.
因此耶穌的名聲被傳開了.
耶穌這個名聲傳開了之後.
婦人就聽見了耶穌的事.
葉老同樣也聽見了耶穌的事.
他們都是從別人所傳來的訊息.
或者故事裡面.
知道耶穌這個人.
他們知道耶穌是一個很大有能力的人.
也曾經醫治過很多絕望的人.
你想想.
一個癱子的人.
一個患麻風的人.
手脫皮的人.
被二千隻鬼附著的人.
他們同樣好像血流婦人和葉老一樣.

$^{201}$沒有盼望.
是絕望.
他們也很希望藉著耶穌的能力.
病患可以得醫治.
就好像那些患麻風的.
癱子的.
被鬼附著的人一樣.
最終能夠得到耶穌的憐憫和醫治.
在這一刻.
婦人就在耶穌的人群裡面.
兩個人的距離是這麼近.
經文對婦人所想的.
作出了一個很具體的描述.
經文是怎麼說的呢?.
他說:我只需要摸一摸她的衣服.
我只需要摸一摸她的衣服.
我就會痊癒了.
這是一個什麼樣的信心呢?.
我只需要摸一摸她.
摸一摸她的衣服.
我的病就好了.
這是一個什麼樣的信心呢?.
這個信心不需要耶穌的請求.
她不需要向耶穌請求.
這個信心也沒有經過耶穌的同意.
這個信心只是.
摸一摸就行了.
於是這個婦人真的從後面摸一摸耶穌的衣裳.
神蹟就這樣發生了.
經文說:婦人的血流止住了.
疾病就醫好了.
困擾了十二年的病.
洗盡她的家財.
困擾了十二年的心靈壓抑.
在這一刻.
釋放了.
對於婦人來說.
這一刻.
是何等的興奮.
想一想.

$^{241}$是何等的興奮.
這個信心超越了一般人的想像.
婦人認為摸到衣裳就可以了.
她沒有經過耶穌的同意.
婦人就是這種信心去支持她.
去摸耶穌衣裳的行為.
這個信心成就了婦人所期望.
這個信心也印證了耶穌的大能.
縱使耶穌處於一個不知情的狀態之下.
耶穌在事件開始的時候是完全被動的.
婦人沒有經過耶穌同意就主動去摸耶穌的衣裳.
不過耶穌縱然是被動的.
但無礙她這個恩典的施予.
當耶穌發覺有能力從自己身上出來的時候.
耶穌就馬上轉為主動.
祂就問.
問眾人:誰摸我的衣裳?.
門徒認為在這個擠擁的場合.
被人摸到是一件很自然不過的事.
為什麼耶穌還那麼介意,在意誰摸我?.
不過耶穌就是定義要知道.
祂定義要知道是誰做這件事.
婦人沒有逃避,沒有逃跑.
於是她就俯伏在耶穌面前.
經文形容婦人那時候的感覺是什麼?.
恐懼見驚.
恐懼見驚地將實情告訴耶穌.
婦人為什麼要感到恐懼見驚?.
因為她真的很害怕被耶穌罵.
事實上她沒有得到耶穌的同意.
就去摸祂的衣裳.
她也很害怕.
害怕什麼?.
害怕耶穌收回這種能力.
讓她的病不能夠醫治.
耶穌的回應不是對婦人的責備.
而是讚賞和安慰.
祂說:女兒,你的信救了你.
平安地回家吧.
你的疾病痊癒了.

$^{281}$我們還記得這個婦人是什麼?.
無名無姓.
在社會上是被忽略的一個.
但現在她有一個尊貴的身份.
她是女兒.
耶穌稱呼她為女兒.
她有一個很尊貴的女兒的名份.
她以前是很隱藏地走到耶穌後面.
不想被人發覺.
很不想和別人接觸.
不過這一刻.
她能夠光明正大地.
在耶穌和眾人面前.
很從容地離開.
她的病不單單因著耶穌的信心能夠被醫治.
她的信心也令她的罪能夠得到赦免.
得到拯救.
因為耶穌不是單單對婦人說.
你的病好了.
不是這麼說.
而是說:你的信救了你.
婦人的信心讓她獲得拯救.
這個拯救不單單是血流的病獲得拯救.
也是一個生命的拯救.
在猶太人的傳統.
身體上的殘缺.
通常歸類為犯罪.
得罪神的一種懲罰.
唯有罪得赦免.
身體才能得到康復.
因此耶穌對婦人的赦罪.
就讓婦人的身心靈也能得到康復.
婦人的身心靈被醫治.
不再感到恐懼.
她的人生應該充滿感恩和盼望.
也充滿讚美.
她在絕望裡.
帶著一個隱藏的信心.
走到耶穌面前.
然後俯伏在耶穌面前.

$^{321}$恐懼戰驚地向耶穌講述她的故事.
不過最後.
她帶著神的恩典.
和歡欣的心去離開.
血流婦人的事件結束.
卻是另一個信心故事的開始.
有人從業主的家裡.
跟業主說:你的女兒死了.
何必還勞駕老師呢?.
這是什麼意思呢?.
就是有人指出.
女孩已經死了.
耶穌也毋須要前來.
因為一切都已經完結了.
死亡將一切的盼望都終結了.
在這一刻.
業主有什麼反應呢?.
這位有身份有地位有名聲.
管理會堂的人.
在這一刻.
鴉雀無聲.
一言不發.
也不知道該說什麼.
甚至在這個三文治式的結構裡.
業主在下半節的故事裡.
連講一句台詞的角色也沒有.
在下文裡.
你完全找不到業主.
一句話也沒有.
好像完全消失在信仰故事裡.
耶穌主動跟業主說.
不要怕,只要信.
耶穌說這句話反映了什麼呢?.
其實就是反映了業主當時的心態.
就是非常的怕.
以及完全失去信心.
女兒的死去.
令他感到害怕.
女兒死去的消息.
令他失去了一切的信心.

$^{361}$眼前的耶穌.
是不是真的能夠令到他的女兒起死回生.
他沒有信心.
這就是他在下半場故事裡.
完全沒有角色的一個原因.
信心失去了.
耶穌沒有理會到業主的反應.
他帶著三個親密的門徒.
就是彼得,雅各和約翰.
繼續前往業主的家.
當時家裡的人.
還為這個女孩的死而難過.
眾人都哀哭飲泣.
耶穌的反應就是.
為什麼大吵大鬧呢?.
孩子不是死,只是睡著了.
眾人就嘲笑耶穌.
耶穌就把他們趕出去.
進入女孩的房間.
然後對女孩說.
女孩我吩咐你,起來.
這句話充滿了權柄.
叫人不得不聽從.
不單單人要聽從.
連死亡都要聽從.
女孩就立即起來.
恢復生命氣息.
正如耶穌所說.
她只不過是睡著了.
現在睡醒了.
馬上可以起來行走.
好像沒有病一樣.
經文透露.
他們很驚奇.
這個女孩已經十二歲了.
他們的驚奇.
當然是因為女孩能夠走動.
沒有死去.
也因為耶穌的權柄和大能.
竟然可以令到一個死去的人.

$^{401}$能夠恢復生命氣息.
他們因為見證到一個神蹟而驚奇.
不過他們除了驚奇之外.
是不是應該對耶穌有一份感恩之情.
是不是對耶穌的身份.
有一份敬畏的心呢?.
對於葉老來說.
會不會因為自己失去了.
對耶穌的信心.
而感到慚愧和懊悔呢?.
血肉婦人患了病十二年.
這個女孩也是十二歲.
這是巧合.
因為作者另有心意.
經文沒有解釋.
不過無論是婦人這十二年的.
患病的苦痛日子.
還是女孩十二年無憂無慮的青春歲月.
她們都是同樣在絕望裡面遇到耶穌.
在她們遇到耶穌之後.
她們的生命都經歷了一個很大的變化.
婦人十二年的病患被釋放.
聚得赦免.
女孩的病被醫好.
生命得以延續.
繼續過去那個美好的人生.
她們不單單同樣被耶穌視為女兒.
有著尊貴的身份.
同樣在未來的歲月裡面.
都會有一段難忘和不一樣的人生.
我們應該感恩.
我們不是在絕望裡面遇到耶穌.
我們回想一下我們信主的經歷.
我們不是那個血流婦人.
我們也不是那個像快要死去的女孩一樣.
但我們會否同樣都像經文裡面.
我們會有一個不一樣的人生.
當我們被耶穌拯救之後.
當我們的生命被耶穌釋放了之後.
我們是否也可以有一個不一樣的人生?.

$^{441}$在兩個故事裡面.
葉老主動去尋找耶穌.
耶穌也主動去回應.
就算旁人認為女孩已經死了.
事件已經告一段落了.
就算葉老在知道女兒死訊之後沒有進一步的請求.
甚至信心全失.
就算被葉老的家人嘲笑.
認為這是風人風語.
人死了還當是睡著.
但是耶穌對女孩的憐憫.
耶穌對女孩的憐憫.
沒有因為別人的誤會.
或者別人的無知.
而阻礙了耶穌對女孩的請求.
也不會因著葉老喪失信心而放棄.
血肉婦人也主動去尋找耶穌.
但卻隱藏著她的信心和行為.
令到耶穌被動地回應.
就算婦人沒有得到耶穌的同意或許可.
就算摸耶穌的衣裳.
這個行為是偷偷摸摸地進行.
但是對於耶穌要請求的人.
耶穌對人的憐憫超越了人一般的信心.
因此無論是耶穌主動.
或者被動地去回應人的請求.
無論是人是否失去了信心.
還是信心的行為沒有得到耶穌的同意和認許.
對於需要被請求的人.
對於耶穌所憐憫的人.
耶穌仍然樂意賜予豐厚的恩典.
我們是否應該感恩?.
因為我們都是耶穌定義地去揀選我們.
耶穌定義要去憐憫的人.
無論我們的信心有多大.
因為我們根本毫無信心.
無論我們是主動還是被動.
只要是耶穌定義去憐憫和請求的人.
我們同樣可以得到豐盛的恩典.
這就是我們所相信的.

$^{481}$以及我們所跟隨的三一真神的屬性.
不是因為我們.
不是因為我們的行為.
不是因為我們的信心.
不是我們所做的任何一件事.
而只不過是神的恩典與憐憫.
我們是否應該感恩?.
弟兄姊妹.
今天講題是一個怎樣的信心?.
從上述兩件神蹟醫治.
我們看到一個有身份有地位的業佬.
在他感到絕望的時候.
他主動去找耶穌.
想耶穌為他的女兒醫病.
他俯伏在耶穌面前.
他的信心大到他覺得.
只要耶穌為他的女兒按手.
女孩就能夠痊癒.
業佬向眾人顯示出他的信心和誠意.
有時候我們也會.
我們很想我們有這樣的信心.
不過當有人說他的女兒已經死了.
你的女兒死了.
在那一刻業佬的信心就全失.
但無礙耶穌請求的恩典.
業佬縱然從有信心到失去信心.
但神的恩典仍然臨到.
這就是我們所相信的主.
我們又看到一個沒有身份沒有地位.
連名字也沒有的婦人.
婦人很絕望.
用盡她所有的大病情卻更加嚴重.
她主動地尋求耶穌.
很想耶穌醫好她的血漏病.
她隱藏著動機.
靜靜地跟在耶穌後面.
她的信心很特別.
她不需要耶穌的同意.
也不需要耶穌的應許.
她認為只要摸到耶穌的衣裳就可以.

$^{521}$最後她的信心得以成就.
婦人的信心和行為.
雖然未被耶穌認可.
但神的恩典仍然臨到.
兩種信心.
可否作出比較.
婦人被醫治的過程裡.
信心未被挑戰過.
起碼沒有被死亡挑戰過.
相反.
耶穌的信心卻被女兒的死亡挑戰過.
以致她的信心全失.
我們一般人.
我們摸一隻薯條.
一般人一定是這樣想的.
讓死人服務的信心.
是否應該要比病得醫治的信心更加大.
這樣才合理呢?.
如果死人服務的信心要大一點.
事實上.
在這兩段經文裡.
或者這一連串的故事裡.
耶穌施行神跡醫治.
完全不在乎人的信心大小.
而在乎耶穌主動的憐憫.
耶穌主動的憐憫.
如何去應用這段經文的屬靈意義呢?.
我自己在這裡講章的時候.
我覺得真的很困難.
教會一般的教導.
我們一定要對神有信心.
我們要對神有信心.
然後神就會幫我們.
但我們看看在葉老的事件裡.
他在女兒死去之後就喪失了信心.
我們是否可以在人死去之後.
仍然很有信心宣告.
這個死去的人可以再一次站起來.
我們應用不了.
我們沒有這樣的信心.

$^{561}$婦人信心的行為呢?.
她未經耶穌的同意.
是否意味著我們的信心行為.
都可以不需要主耶穌的同意.
就可以成就呢?.
我想就可以了.
血柳婦人就想想.
我摸到耶穌的衣裳就可以了.
她不需要耶穌同意.
那我們呢?.
我們是否可以不需要耶穌同意?.
憑著我們的信心就可以讓信心達成.
兩種信心在現實的應用上都遇到困難.
那我們怎樣去應用?.
或者可以這樣說.
如果我們沒有信心的話.
神一定不會幫到我們.
如果一點信心都沒有的話.
神怎能夠在一個微弱的信心裡幫我們呢?.
信心是必須的.
但是要有多大的信心才能夠成就.
卻不是一個必然的條件.
信心越大.
神當然會越恩賢和越立.
但是信心的微小.
神的憐憫和應許.
卻可以補足我們信心的不足.
關鍵仍然是.
神有沒有對我們施予憐憫.
今天我們如果要尋求耶穌.
不論我們的身份,地位,性別.
耶穌同樣接納我們.
當我們要尋求耶穌的幫助的時候.
我們仍需要對耶穌有信心.
不過耶穌為我們所成就的事.
卻和我們的信心大小沒有必然的關係.
而是指耶穌的憐憫.
只要是主所揀選的.
只要是主所定義的.
主所定義要拯救的.

$^{601}$神所作的功必定能夠成就.
我們知道神是一個完美的神.
但是祂卻沒有要求信徒同樣完美.
而是要求信徒盡心,盡性,盡義,盡力.
每一個信徒所能盡的能力和效果都不一樣.
所以我們不能夠用成果作為標準.
而是以信徒所盡的心思和行為.
信徒是否盡了這種心思和實踐?.
主耶穌知道.
業佬以及血流的婦人.
在他們絕望的時候.
他們真的盡了他們的信心和能力.
對於他們來說,這是他們的極限.
因為他們有限制.
或許業佬最後喪失了信心.
或許婦人的行為未經主耶穌的同意.
但耶穌看到他們那種盡的心.
耶穌就憐憫他們.
補足了他們的信心.
因此今天我們不是要說信心的大小.
不是我們信心大一點,神就一定應許.
信心少,神不應許.
不是.
而是我們要看測.
審查我們自己對主的愛.
我們是否真的能夠愛主多一點.
我們愛主的心是否已經盡了.
我們一切的心思意念?.
是否已經盡了?.
我們回來崇拜.
是否已經盡了我們的心思意念?.
是否能夠準時來到?.
是否有足夠時間讓我們的心思醞釀進入敬拜?.
是否有精神去聽道?.
是否已經盡了?.
神知道,神一定知道.
最後我曾說過.
這篇道是為了福音樂曲隊而寫的.
所以最後我祝願每一個福音樂曲隊的成員.
在他們的事奉神的過程裡.

$^{641}$無論信心如何.
他們都能夠盡上一切的心思意念.
為的是要高舉主的榮耀.
願榮耀還讚歸給三一的主.
我們一起去祈禱.
三一的上主我們要稱頌你.
因你的慈愛憐憫.
充滿我們.
你的愛彌補我們信心的不足.
你的愛遮蓋了我們所犯的罪.
你的愛讓我們生命全然成聖.
你的愛讓我們可以與你復和.
也可以與別人復和.
你的愛讓我們能夠進入神的致勝所.
與你相遇.
在你裡面.
我們的生命可以有信.
有望.
有愛.
甚至我們每一個屬你的.
都在你的愛裡面茁壯地成長.
成為一個合你心意的僕人.
被你使用.
讓你的名被廣傳.
被高舉.
讓我們能夠一生地跟從你.
我們的祈禱是奉主耶穌基督得勝名至其球.
阿們.
願主祝福大家.
大家.
\newpage



\section{馬太福音 25:31-46}
\label{sec:KSHhb73QdyA}
\textbf{2022.10.09 《作一隻蒙天父賜福的綿羊》 馬太福音 25\_31-46 (和修版) 曾敏姑娘}
\newline
\newline
連結: \href{https://youtube.com/watch?v=KSHhb73QdyA}{\texttt{https://youtube.com/watch?v=KSHhb73QdyA}} ~~~~ 語音日期: 2022-10-15
\newline
\newline
\hyperref[sec:bLMot_knMH4]{\small{< < < PREV SERMON < < <}}
~
\hyperref[sec:index_chronic]{\small{[返順時目]}}
~
\hyperref[sec:index_scriptual]{\small{[返順卷目]}}
~
\hyperref[sec:rTQz8Qod6Ek]{\small{> > > NEXT SERMON > > >}}
\newline
\newline
馬太福音 25:31-46
\newline
\begin{longtable}{cl}
\hline
\hline
章節 & 經文 (和合本修訂版)\\
\hline
25:31 & \begin{tabularx}{0.7\textwidth}{X} 「當人子在他榮耀裡，同著眾天使來臨的時候，要坐在他榮耀的寶座上。 \end{tabularx} \\ \\ \relax
25:32 & \begin{tabularx}{0.7\textwidth}{X} 萬民都要聚集在他面前。他要把他們分別出來，好像牧人分別綿羊、山羊一般， \end{tabularx} \\ \\ \relax
25:33 & \begin{tabularx}{0.7\textwidth}{X} 把綿羊安置在右邊，山羊在左邊。 \end{tabularx} \\ \\ \relax
25:34 & \begin{tabularx}{0.7\textwidth}{X} 於是王要向他右邊的說：『你們這蒙我父賜福的，可來承受那創世以來為你們所預備的國。 \end{tabularx} \\ \\ \relax
25:35 & \begin{tabularx}{0.7\textwidth}{X} 因為我餓了，你們給我吃；渴了，你們給我喝；我流浪在外，你們留我住； \end{tabularx} \\ \\ \relax
25:36 & \begin{tabularx}{0.7\textwidth}{X} 我赤身露體，你們給我穿；我病了，你們看顧我；我在監獄裡，你們來看我。』 \end{tabularx} \\ \\ \relax
25:37 & \begin{tabularx}{0.7\textwidth}{X} 義人就回答：『主啊，我們甚麼時候見你餓了，給你吃；渴了，給你喝？ \end{tabularx} \\ \\ \relax
25:38 & \begin{tabularx}{0.7\textwidth}{X} 甚麼時候見你流浪在外，留你住；或是赤身露體，給你穿？ \end{tabularx} \\ \\ \relax
25:39 & \begin{tabularx}{0.7\textwidth}{X} 又甚麼時候見你病了，或是在監獄裡，來看你呢？』 \end{tabularx} \\ \\ \relax
25:40 & \begin{tabularx}{0.7\textwidth}{X} 王回答他們說：『我實在告訴你們，這些事你們做在我弟兄中一個最小的身上，就是做在我身上了。』 \end{tabularx} \\ \\ \relax
25:41 & \begin{tabularx}{0.7\textwidth}{X} 「王又要向那左邊的說：『你們這被詛咒的人，離開我！進入那為魔鬼和他的使者所預備的永火裡去！ \end{tabularx} \\ \\ \relax
25:42 & \begin{tabularx}{0.7\textwidth}{X} 因為我餓了，你們沒有給我吃；渴了，你們沒有給我喝； \end{tabularx} \\ \\ \relax
25:43 & \begin{tabularx}{0.7\textwidth}{X} 我流浪在外，你們沒有留我住；我赤身露體，你們沒有給我穿；我病了，我在監獄裡，你們沒有來看顧我。』 \end{tabularx} \\ \\ \relax
25:44 & \begin{tabularx}{0.7\textwidth}{X} 他們也要回答：『主啊，我們甚麼時候見你餓了，或渴了，或流浪在外，或赤身露體，或病了，或在監獄裡，沒有伺候你呢？』 \end{tabularx} \\ \\ \relax
25:45 & \begin{tabularx}{0.7\textwidth}{X} 王要回答：『我實在告訴你們，這些事你們沒有做在任何一個最小的弟兄身上，就是沒有做在我身上了。』 \end{tabularx} \\ \\ \relax
25:46 & \begin{tabularx}{0.7\textwidth}{X} 這些人要往永刑裡去；那些義人要往永生裡去。」 \end{tabularx} \\ \\
[1ex]
\hline
\hline
\end{longtable}
$^{1}$(主持人:弟兄姊妹平安).
(主持人:我們一同再有一個禱告).
(主持人:願意今天你親自對我們每一個說話).
(主持人:完聖靈充滿了我們的心 讓我們的心真的向上帝你打開).
(主持人:我們明白上帝你所要對我們說的).
(主持人:你讓我們的心感到紮心 我們就回應你).
(主持人:奉耶穌基督的名字祈禱 阿們).
(主持人:今天這段經文也是我們當中不多去說的經文).
(主持人:很多時候一說到這段經文會令人想起扶貧時宮).
(主持人:還不是 說到餓了 赤身老體 病了等等).
(主持人:感覺好像在說扶貧時宮 其實不是這回事).
(主持人:今天讓我們重新來看這段經文).
(主持人:盼望到最後我們真的明白什麼叫做).
(主持人:做一隻望天父賜福的羊 而且是綿羊).
(主持人:我們來看看 其實這裡是一個比喻).
(主持人:是綿羊和山羊的比喻).
(主持人:整個比喻的場景是什麼情況呢?).
(主持人:原來這個場景就是有一天耶穌基督要再回來).
(主持人:祂要回來審判世界的時候的一個場景).
(主持人:當人子在榮耀和眾天使來臨的時候).
(主持人:要坐在榮耀的寶座上).
(主持人:在這短短的一節裡出現了兩次榮耀).
(主持人:原來有一天主耶穌基督要再回來的時候).
(主持人:是非常榮耀的).
(主持人:最近我在教啟示錄).
(主持人:我們會說到耶穌基督有一天要回來的時候).
(主持人:是什麼樣的場景呢?).
(主持人:就是非常榮耀的).
(主持人:耶穌基督第一次來到這個世界的時候).
(主持人:大家記得祂的形象是非常卑微的形象).
(主持人:祂出生之後不能夠住在客棧裡).
(主持人:要被放在馬槽裡).
(主持人:到後來祂長大的時候).
(主持人:當有一天祂要進入耶路撒冷的時候).
(主持人:是祂騎著的小驢子進入耶路撒冷).
(主持人:祂的整個形象是和平的君王).
(主持人:是非常卑微的).
(主持人:但是到第二次).
(主持人:耶穌基督要再回來這個世界).
(主持人:要審判萬民的時候).

$^{41}$(主持人:祂的形象不一樣).
(主持人:在短短這一節).
(主持人:有兩個榮耀的字眼).
(主持人:耶穌回來就是這麼榮耀的).
(主持人:祂回來的時候).
(主持人:這一次不是騎著小驢子).
(主持人:而是和眾天使一起來).
(主持人:在這個場景裡).
(主持人:你也看到萬民都要面對耶穌基督的審判).
(主持人:所以萬民要聚集在祂的面前).
(主持人:然後耶穌要按著審判所定的).
(主持人:將祂們分別出來).
(主持人:將綿羊和山羊分開).
(主持人:綿羊放在祂的右手邊).
(主持人:山羊放在祂的左手邊).
(主持人:通常右手邊是比較重要的位置).
(主持人:是權柄,榮耀的右手邊).
(主持人:在這裡).
(主持人:左右的分別是右手指好的一方).
(主持人:左手指不好的一方).
(主持人:是這樣分開的).
(主持人:審判萬民是包括了信徒和非信徒).
(主持人:是各個民族,所有的人).
(主持人:在審判的日子卻要分開綿羊和山羊).
(主持人:這幅圖片很明顯是法庭的場景).
(主持人:因為耶穌要審判萬民,好像法庭一樣).
(主持人:分開綿羊和山羊).
(主持人:綿羊是上面那一種).
(主持人:看上去體型比較大,肥胖一些).
(主持人:毛又多一些).
(主持人:羊肉價值也很多,從皮到肉各方面).
(主持人:山羊比較瘦削,角比較直).
(主持人:也比較聰明,很機警).
(主持人:可以跳來跳去,甚至上樹也很容易).
(主持人:但我想說很特別).
(主持人:在這裡用綿羊和山羊作為比喻).
(主持人:以往在舊約聖經裡面).
(主持人:總會提到有些動物被視為不潔).
(主持人:以色列文是不可以吃的).
(主持人:有些是潔淨的動物是可以吃的).

$^{81}$(主持人:但我想說綿羊和山羊).
(主持人:都是被視為潔淨的動物).
(主持人:所以獻祭的時候可以獻綿羊和山羊).
(主持人:在這兩個層面,其實是一樣的).
(主持人:不是用這個比喻的原因).
(主持人:因為山羊是不潔淨的).
(主持人:綿羊是潔淨的,所以可以這樣分辨).
(主持人:但經文是有意思的).
(主持人:要將兩種這麼接近的動物).
(主持人:放在一起作一個比較).
(主持人:綿羊和山羊的共通性是甚麼呢?).
(主持人:都是潔淨的動物).
(主持人:還有一個很簡單的問題).
(主持人:共通性是甚麼呢?).
(主持人:羊,都是羊,潔淨的).
(主持人:可以獻給上帝).
(主持人:但又要在裡面區分).
(主持人:為甚麼呢?).
(主持人:是有原因的).
(主持人:其實看上去很相似的動物).
(主持人:性質也很相似).
(主持人:但在上帝眼中要將它們區分).
(主持人:這個分別很大).
(主持人:綿羊被視為比一變好).
(主持人:是夢上帝賜福的).
(主持人:但山羊被視為不好).
(主持人:要被詛咒的).
(主持人:結果是這麼大分別).
(主持人:我們一般看著綿羊和山羊).
(主持人:是不會覺得有這麼大的分別).
(主持人:但在這個比喻中).
(主持人:在神眼中看是有很大分別).
(主持人:究竟是代表什麼呢?).
(主持人:實際上要形容什麼呢?).
(主持人:它們相似).
(主持人:但審判時卻是不一樣的).
(主持人:綿羊和山羊究竟是什麼呢?).
(主持人:是什麼不同呢?).
(主持人:當我們看回後面的經文時).
(主持人:你會發覺).

$^{121}$(主持人:聖經會形容這些綿羊是義人).
(主持人:在第37節義人就回答).
(主持人:是義人).
(主持人:義人並不是說).
(主持人:這些是做好事的人).
(主持人:有很多善行的人).
(主持人:不是這個意思).
(主持人:這些義人是說).
(主持人:他們是相信上帝的人).
(主持人:是真正相信上帝的人).
(主持人:不是掛名是相信上帝的人).
(主持人:這些義人就是綿羊).
(主持人:另外那些).
(主持人:究竟是否真的相信神?).
(主持人:他們看上去像是相信神的人).
(主持人:但卻是被詛咒的).
(主持人:這些生命).
(主持人:我們要打的問號).
(主持人:不是真真正正相信上帝).
(主持人:人相信上帝不是只是說).
(主持人:我現在承認我犯罪).
(主持人:我得罪上帝).
(主持人:我相信耶穌基督是神的兒子).
(主持人:我邀請祂進入我的生命).
(主持人:成為我的救主和生命的主).
(主持人:我很清楚應信的).
(主持人:我就是這樣).
(主持人:是很重要的).
(主持人:這個應信是重要的).
(主持人:是一個開始).
(主持人:但真正的信).
(主持人:是你之後的生命是一個徹底的改變).
(主持人:真正的義人是生命改變的).
(主持人:你真是跟從耶穌的).
(主持人:你的生命是活出耶穌基督的).
(主持人:這樣叫義人).
(主持人:同樣進入教會的門內).
(主持人:每個人都是羊).
(主持人:但當中有綿羊有山羊).
(主持人:你看下去).

$^{161}$(主持人:每個人都說自己信了耶穌).
(主持人:每個人都缺了字).
(主持人:我每個星期都上崇拜).
(主持人:我上團契).
(主持人:我又事奉).
(主持人:我又這樣又那樣).
(主持人:你看下去).
(主持人:看下去都是羊).
(主持人:但上帝眼中看就分開了).
(主持人:有綿羊和山羊是不一樣).
(主持人:而且分別很大).
(主持人:真正的活出上帝).
(主持人:活出耶穌基督和不是有分別).
(主持人:真正的跟從耶穌生命改變).
(主持人:和不是跟從耶穌基督).
(主持人:生命沒有改變).
(主持人:是有很大很大的分別).
(主持人:結果是很大的分別).
(主持人:而到耶穌基督回來的時候).
(主持人:去審判的時候).
(主持人:原來不只是面對不信主的人).
(主持人:也會面對信徒).
(主持人:要面對信徒).
(主持人:原來我和你的生命).
(主持人:去到那一刻要面對耶穌).
(主持人:今天我們一起).
(主持人:紅恩家這裡).
(主持人:這麼多弟兄姊妹坐在這裡).
(主持人:一起崇拜).
(主持人:是很美的一幅圖畫).
(主持人:大家要記住這批人).
(主持人:將來要一起面對上帝).
耶穌的評價是什麼?判斷是什麼?.
有分別的.
分別就來了.
分開一批人去了右邊.
就是綿羊.
真正的義人.
跟從上帝的.
活出耶穌基督的.

$^{201}$生命有改變的.
這批人.
耶穌要稱讚他們.
你們這些是夢我.
付賜福的.
可以來承受.
創世以來為你們預備的國.
這裡說的是舊恩.
說的是永生.
是上帝的賞賜.
上帝很大的恩典.
這個所預備的國.
這份永生.
是神預備了很久的.
創世以來就預備的.
不是到後來.
不知哪一刻.
忽然間上帝就說.
我要見到人們.
為你們預備的國.
不是.
是很久很久的.
上帝愛我們是很久很久的事情.
創世以來已經愛我們.
已經知道人.
沒辦法去救贖自己的生命.
所以祂為我們預備最寶貴.
是很久的預備.
原來.
去審判的時候.
主要將這份禮物.
給那些綿羊.
給那些義人.
究竟這些義人做了什麼呢?.
因為我渴了.
你們給我吃.
渴了你們給我喝.
我流浪在外.
你們留我住.
我赤身露體.

$^{241}$你們給我穿.
我病了你們看顧我.
我在監獄裡.
你們來看我.
一連串很多的行動.
這些行動.
是做在誰的身上呢?.
哪怕是做在.
我弟兄中最少的一個身上.
換言之這些行動.
是做在.
屬上帝的人身上.
所以我說的.
不是扶貧侍工的意思.
這些行動.
不是說做在.
全世界所有的.
餓的,病的.
流浪在外的人身上.
不是.
做在弟兄.
屬於上帝的人身上.
耶穌基督的門徒身上.
耶穌基督的門徒.
都有很多的經歷.
很多的患難.
很多的困苦.
他要特別強調.
哪怕是做在.
最少的弟兄身上.
很多時候.
人去看人.
看事情.
都容易.
分階級.
無論我們承不承認.
在人群當中.
總會看.
這個人的位分高.
他是一個重要的人.

$^{281}$他影響很多人.
我對他必恭必敬.
我對他的態度是不同的.
說話又不同.
尊敬多很多.
比他特別多.
這個卑微很多.
這個影響力不大.
甚至乎沒有人知道.
他叫什麼名字.
他是哪位.
沒有那麼重要.
所以我對他的態度.
又差很遠.
在教會外是這樣.
在教會裡面也會.
我們必須承認.
是這樣.
有時你會發覺.
我們對不同的人.
態度不一樣.
在教會裡也是.
但耶穌基督說.
就算你做在一個.
最少的弟兄身上.
也要計算.
這個最少的弟兄.
可能就是.
沒有地位.
不被人重視.
不被人看重.
甚至乎人家.
知道他叫什麼名字嗎.
都不知道.
但耶穌基督.
看他為寶貴.
耶穌一定不會.
輕看這個弟兄.
最少的弟兄.
不會輕看他.

$^{321}$他經歷這一切.
患難困苦的時候.
耶穌基督與他認同.
這裡很特別.
很特別.
我們後面會看到.
不只是教會的.
耶穌基督的門徒.
經歷這些患難困難.
他們經歷患難困難的時候.
耶穌與他們同在.
耶穌與他們認同.
他們經歷患難.
就好像耶穌自己經歷患難一樣.
他們飢餓就像.
耶穌飢餓一樣.
他們病了就像.
耶穌病了一樣.
他們坐牢就像.
耶穌坐牢一樣.
耶穌完全將自己.
與他們認同.
耶穌完全的.
是身同感受.
就在他們的生命裡.
是不能夠分割.
耶穌是這樣看.
所以當這些義人.
去做這些行動的時候.
給飢餓的人吃東西.
給飢渴的人喝東西.
流浪在外的留他們住.
赤身奴隸的給他們穿衣服.
病了的去看顧他們.
去探監獄裡的人.
原來在耶穌眼中.
這些人就做在他們身上.
這些人正在探的.
其實不是直接幫的.
他們也不是幫耶穌基督的.

$^{361}$直接來說.
但是耶穌已經會說.
就是做在耶穌身上.
因為耶穌與這些人認同.
在這個過程是很特別.
原來我們這樣看顧我們的弟兄.
真正的愛弟兄.
第一件事是會有行動的.
不只是我口中說我很愛你.
你看著你前後左右的人.
每個星期都回來崇拜.
我們說完之後.
或者是牧者很喜歡這樣說.
互相問候.
彼此祝福.
不單止是這樣.
真正的愛不只是在舌頭上.
是在行動上.
你紀念對方的需要.
你心同感受明白.
他的痛苦.
你進入到別人的生命.
在你可行的裡面.
你付出代價.
這裡其實說真的.
他付的代價.
又不是大得這麼厲害.
是行動.
但未去到盡.
耶穌稱讚不是說.
你徹底盡力.
還要爆燈.
你的付出.
例如我舉個例子.
餓了就給我吃.
但有沒有說.
我請他吃自助餐呢?.
有沒有去到這麼盡?.
沒有吧?.
我學了你給我飲.

$^{401}$你真的給我滿足.
你留我住.
這個很寶貴.
現在香港很難做到.
很多我們的顧忌.
赤身露體給我三象.
病了你看我.
但沒有說在醫院陪他過夜.
是否一天去探.
一天去陪他.
去到盡.
在監獄你們來看我.
其實監獄來看.
可以再盡力一點.
以前可以陪坐監.
以前不是說現在.
我說那種善行.
那種好的行動.
關心的行為可以去到更盡.
但耶穌的要求不是去到盡.
你去到這個位置.
你見到別人有需要.
你付出.
神已經欣賞.
已經欣賞.
但你至少是由行動活出來.
而這個行動.
不是出於任何的機心,企圖.
不是因為我知道.
我將要拿很多禮物.
如果我做這件事.
我就會拿很多很多禮物.
於是為了那些禮物就做.
日後沒有禮物就不做了.
或者要做些什麼.
其實做人是很自然的.
從心發出的.
就是這樣發自內心的.
甘心樂意地做.
做的時候沒有想過.

$^{441}$原來耶穌基督這樣與那個人認同.
原來我正在做就是在服侍耶穌.
他也沒有這樣想過.
就發自內心地去做.
在背後有一份很大的愛.
如果不是愛.
是不會做得很久.
會很辛苦.
如果不是因為愛.
這是很特別的.
主是稱讚他們.
在做的時候也不知道.
所以當主稱讚他們的時候.
我餓了 你讓我吃.
我渴了 你讓我喝.
他們就會覺得很驚訝.
什麼時候?什麼時候做在你身上?.
我都不知道.
原來你做在打拍一個弟兄身上.
原來做在我身上.
主耶穌很愛.
在我們群體當中.
這些很有需要的肢體.
很愛 很愛他們.
自己的生命就與他們一起.
原來我們服侍這些弟兄 姐妹.
就是在服侍耶穌基督.
這個服侍是甘心命底的.
出乎愛的 毫無譏心的.
這樣做的時候.
就做在耶穌的身上.
其實這件事很特別.
我們之所以這樣做.
其實這反映了我們與上帝的關係.
與耶穌基督的關係.
當我很愛主的時候.
我真的很知道主的心意.
真正愛主的人就知道主的心意.
主就是愛這些肢體.
愛這些弟兄 姐妹.

$^{481}$當我知道主的心意的時候.
因為我愛主 我就去做.
就是要回應上帝的愛.
我不需要出於被迫.
不像守規條.
規條第一 必須這樣做.
規條第二 必須這樣做.
這些人並不是為了守規條.
沒有人強迫他們這樣做.
是自己甘心主動的.
因為愛神.
就算他不知道這樣做是服侍神.
也是出於他與上帝的關係.
出於他與上帝的愛的關係.
去愛神.
耶穌基督就是這樣照顧這些.
有病的 有飢餓的 口渴的.
流浪在外的人.
耶穌基督就是這樣愛這些人.
愛坐牢的.
愛群體中軟弱的 困苦的.
我愛耶穌基督的時候.
就是我愛耶穌基督所愛的.
當我這樣做的時候.
出於愛的時候.
就得到主的讚賞.
你也不知道是什麼時候.
有一天面對主 主就會讚我.
但之前我並不是為了什麼利益.
這就是二人的獎賞.
在過程中.
我們再看回.
另一批是關於山羊.
王向左邊的人說.
你們這些被詛咒的人.
離開我.
結果很多人害怕.
要進入魔鬼和使者所預備的永火.
這是地獄 是永死.
剛才是永生 是救恩.

$^{521}$上帝要獎賞給真正愛他 跟從他的人.
另一方面是詛咒.
很大的分別.
要詛咒那些不是真心跟從他.
而是真心愛他的人.
是非常非常嚴厲.
地獄本身是為魔鬼和使者所預備.
現在不是真心跟從他的人.
會有這樣的結局.
兄姐們有沒有想過呢?.
有沒有想過這件事?.
因為我餓了 你們沒有讓我吃.
渴了 你們沒有讓我喝.
我流浪在外 你們沒有留我住.
我赤身無體 你們沒有讓我穿.
我病了 我在監獄.
可能我們會說 很多人看著你.
為什麼要我來看著你呢?.
很多人照顧這些有需要的人.
我看到他們 我知道他們有需要.
我知道 不用我了.
我是否可以選擇不做呢?.
為什麼要我做呢?.
雖然說不去也很盡.
但也要付代價.
不做有什麼問題?.
我真的很相信群體這麼大.
有人做就行了.
他做完我讚他.
你做得很好.
好像弟兄姊妹是屬於那些人.
我沒有份 不關我的事.
我過好我的人生.
最重要是拿著護照去永生.
我可以的.
我很想享受我的人生.
不要緊 總有人做.
放心 教會這麼大.
這是很大的冷漠.
事不關己.

$^{561}$和我有什麼關係呢?.
我感受不到赤身露體的痛苦.
他病了 不是我病.
他說很辛苦 我怎麼知道呢?.
他餓了 我不知道.
現在還有人餓嗎?.
我即使知道別人這樣.
我不會切身處地站在他的處境.
感同身受他的痛 他的需要.
其實我覺得我們很容易進入這樣的狀態.
真真正正感同身受別人的需要.
真正愛別人.
不容易.
還有 香港是很忙的.
我很多事情忙.
交給時間多的人.
我忙這些 你讓我忙那些我想忙的事.
我有空我過來.
有空與否的問題是.
其實我講堂道 預備堂道.
我都要問自己 我忙什麼呢?.
我都覺得自己很忙.
是否忙就是一個理由.
不去感同身受.
別人的痛苦和需要.
我的弟兄 我的姐妹 是否我沒有份呢?.
主將我放在這個群體.
要我真真正正愛我的弟兄 我的姐妹.
我究竟做了什麼?.
可能某一處很有愛心.
能感受別人困難的弟兄姊妹 特別多付出.
那我呢?.
我們不要問別人 問問自己.
我在別人的身上做了什麼?.
我覺得很震撼.
在這裡 主是很體重.
主自己就是這樣去愛人.
愛每一個熟悉的子民百姓.
很深很深的愛.
但我們是否這樣活出這份愛?.

$^{601}$我是否真的明白主的心意?.
是否這樣愛人呢?.
還是我是很難忘?.
與我無關.
生命的需要 我估計不止這些.
不止餓 渴 流浪在外 赤身無體.
是更多.
可能我們在座的弟兄姊妹 各自有不同的需要.
生命有自己的困苦 有自己的難.
在這裡 是否經歷了愛?.
在這裡 我們是否成為一個愛別人的人 真心地愛?.
不是出於任何的機心 甘心樂意地愛 是不是呢?.
在這裡 主責備他們 主責備這群山羊.
你沒有真正地去愛這些困苦中的弟兄姊妹.
沒有真正地去愛.
這群山羊都很驚訝 是嗎?.
我何時沒有?.
就算見到這些人有需要 都是這些人.
是主你來的嗎?.
主說就是我 是不是?.
就是我.
你沒有照顧這些人 你沒有愛這些人 就是沒有愛我.
你沒有做在他們的身上 就是沒有做在我身上.
我深深地愛著這些人.
但你沒有我那份心腸去愛這些人.
其實最後真的不愛我.
所以這仍然反映了我們和上帝的關係.
我們是否真正地愛上帝?.
到最後想問這件事.
我真正地愛弟兄姊妹嗎?.
有多愛?.
他可愛的時候 他很好的時候.
他會陪我吃東西 特別好 我們很友善.
說話這樣那樣.
總之他對我有很多好處 我就愛他.
這樣是否真正的愛?.
如果他不是這一種.
他真的有生命很大的困苦 很大的需要.
我會放上多少時間?.
我會付出多少?.

$^{641}$這都是我問自己的.
是否真正的愛?.
是多愛?.
是否真心諾言 自然地愛?.
在裡面發出來的?.
是否因為愛上帝的緣故?.
真正愛上帝.
就會很自然地發出來.
不需要勉強 不需要因為規條.
有多愛?.
問問大家.
真正的愛神是怎樣的?.
就是有上帝的心腸.
與痛苦的人同在.
真正愛神的人.
就是服事上帝所服事的.
上帝在哪裡 我們就在哪裡.
沒有人看著我們的時候 我們都做.
因為不是我們自己做.
不是為了讓別人看.
我自己很感動.
在我們弟兄姊妹當中.
真的有人很有愛心.
一直默默地做很多事情.
其實很大部分弟兄姊妹都不知道.
因為我們有些探訪的對象.
未必回到我們教會.
但他們都有信耶穌.
我們有弟兄姊妹栽培他們的生命.
去探訪他們的身體.
心靈都有很多需要.
日常生活的照顧是很多的.
我們有弟兄姊妹煮飯給他們吃.
陪診.
是很多很多時間和金錢的擺上.
他們不會大聲地告訴大家.
今天我去探訪.
我今天煮飯.
我想讓全世界知道我做了多少.
不是.

$^{681}$就是這樣做在弟兄姊妹身上.
可能大家都不認識弟兄姊妹.
甚至乎她是否存在呢?.
大家都不太認識.
在什麼時候知道弟兄姊妹對她的感情呢?.
例如喪禮.
如果是弟兄姊妹.
其實根本大家不熟悉.
喪禮時有多少人會出席呢?.
我不是說所有喪禮我們都要出席.
但你也看到.
大家有多關心這個生命.
真正的愛是怎樣呢?.
我們每個星期回來崇拜.
前後左右的人你認識嗎?.
你知道人生命有什麼重擔嗎?.
有沒有想過勝在人生命的重擔.
真正的愛.
還是我們未曾愛之前.
很容易被一些是非.
彼此的爭執影響.
我們看到對方的不好.
很多問題.
我們這邊說一下.
那個不好.
那邊又說一下.
這些又是不好的.
於是整個群體就一直是是非非.
我們很不開心.
我們覺得自己的生命的困難.
是完全兩幅圖畫.
是兩幅圖畫.
不是愛.
真正的愛是這樣的.
所以弟兄姊妹.
我挑戰大家去想.
是否真正的愛弟兄姊妹.
山羊綿羊的比喻是說.
愛弟兄姊妹.
是否真正的愛神.

$^{721}$你愛神.
別人不用跟你說任何事.
你就會想做一些.
令神歡喜快樂的事.
神愛什麼人.
你就愛什麼人.
神在哪裡工作.
你就在哪裡工作.
神愛教會.
神很愛教會.
我今早說啟示錄.
說到神要藉著約翰寫信給七間教會.
說教會裡.
有什麼不好的地方.
也要說教會要認罪悔改的地方.
神很愛教會.
看著教會.
而且教會是在上帝的掌管裡.
在耶穌基督的掌管裡.
今天耶穌基督很愛洪恩家.
我們愛洪恩家.
不要說我很愛.
你愛.
就會心甘情願.
按洪恩家的需要付出.
這裡有很多事奉上的需要.
我們有沒有發揮我們的恩賜呢.
去服侍神.
不是一小撮人服侍.
是每個弟妹都要找回上帝給你的恩賜.
起來擺上.
不是被迫的.
不是拉拉扯扯的.
你愛就會做.
你想想你所愛的人.
你為他付出多少.
你心甘情願.
很多事主動做.
甚至多於對方要求.
因為你愛.

$^{761}$你愛教會.
愛神所愛的教會.
你也會主動付出.
甚至多於所需要.
其中一件事我要挑戰大家想.
愛教會.
你覺得教會每個月的支出大嗎.
我看到很多頭在發呆.
教會在什麼地方需要支出呢.
我今天夠膽說錢.
我們十世都不說.
教會要支出.
什麼地方要支出呢.
我們很開心.
坐在這裡有冷氣炭.
如果為了節省電費.
我們全部關掉冷氣.
今天大家一路搶著搶著.
搶著搶著撥著.
這裡可以省很多冷氣費.
冷氣費也不止.
燈油,火蠟全部要支出.
童工的薪津.
還有侍工的費用.
每件事都要支出.
我們平時甚少在講台上說.
我真正愛上帝的時候.
我不需要別人說.
就奉獻.
窮寡婦錢不多.
但她將養生的全部丟下來.
遠遠超過上帝希望她丟掉的.
因為她真是愛上帝.
那兩個錢可能在別人眼裡.
兩個錢而已.
不值錢啊.
她是富豪.
一丟就丟一百萬下來.
這樣才像樣.
不是.

$^{801}$她的心就是愛教會.
但我沒有這麼多.
我真的沒有.
但出於我的愛.
我覺得這些就是這麼多.
我就給到盡.
今天我和你怎樣看呢.
教會是需要更多更大的奉獻.
否則怎樣傳福音呢.
經常說我們走出去洪水橋區.
和別人傳福音.
帶領更多人進來.
這裡天天燈油火蠟.
在這裡花錢.
但卻遠遠比我們所需要的.
是差很遠的距離.
我最喜歡說的.
今天我和老牧師還是屬於感受.
是不是.
說得笑一點.
有一天感受把我們兩個召回去.
這裡就完全沒有傳道人了.
為什麼是感受呢.
感受是發薪水給我們.
所以在這個過程裡.
要更大更多的努力.
一起同心.
愛教會就丟上了.
為上帝懸故.
不是為我們懸故.
為福音的懸故.
你真的想傳福音.
所有事情都需要花錢.
怎樣帶新朋友回來.
我們舉辦活動.
不好意思.
我們想想有什麼經費.
我們又不行了.
等等等等.
這裡不行那裡不行.

$^{841}$你想福音拓展.
還有長遠想發展青少年士工兒童士工.
這些就是丟本的.
不斷丟錢.
從哪裡來呢.
老師出去打工.
我們各人賺了一份.
星期天才回來廣道.
接著半條命.
是不是這樣呢.
丁姐妹一起來.
愛神.
就愛神的家.
愛神所愛的.
愛神.
生命就要改變.
你成為他喜歡的樣式.
我最近經常問這個問題.
一年你聽多少堂課.
52堂.
不止.
有些人看阿培連經大會.
早午晚幾堂.
有些人喜歡聽聽那個講完.
外面又不知道是誰講完.
我很想聽聽他們講.
聽完一定超過52堂.
你聽了這麼多堂.
我最喜歡問的問題.
你記得哪一堂課到.
52堂到都很厲害.
我不相信主題重覆.
如果全部主題都不重覆.
52個教道.
哪怕每次只講一點.
你都有52個教道.
我想問丁姐妹.
你真的愛上帝.
你有沒有活出上帝的話語所提醒你的東西.
有沒有.

$^{881}$52個教道.
如果都活出來.
那已經很厲害了.
有些道可能完全沒有聽進去.
有很多大部分都忘記了.
意義在哪裡呢.
上帝不斷透過講完跟你講說話.
轉眼就沒有了.
問問自己.
是不是真的愛上帝.
萬望丁姐妹的生命不再一樣.
你聽了到就要改變.
你真的愛上神就有行動.
真的愛人都有行動.
不要聽完就轉頭回頭.
就講別的東西.
我們過我們的生活.
好像從來沒有參加過崇拜.
我們一起禱告.
我們一起祈禱.
天上帝求你寬恕我們.
當我們看這個比喻的時候.
可能我們覺得那些山羊很差.
是的.
他們真的差.
所以上帝說他們是要被詛咒的一群.
他們的生命好像有信徒的樣子.
有信徒的樣式.
但是沒有實質.
主大大恐怕我們今天的生命都是這樣.
我們有信徒的樣子.
我們過著信徒的生活.
但是沒有實質.
我們不愛你.
我們的耳朵發癢.
耳朵起疹.
不太聽到你的聲音.
不要聽入心.
也不會活出來.
不會照著你喜悅的樣式去活.

$^{921}$我們的心對著弟兄姊妹很硬心.
不是真真正正地愛她們.
所以今天我相信聖靈在我們裡面.
也會替我們擔憂難過.
主我求你提醒我們.
讓我們真的要認罪悔改.
每一次我們聽到都讓我們責心.
認罪悔改.
真的去成長.
主求你幫助我們.
幫助我們五在一讓.
我們的生命真的以行動去見證.
我們是愛你的.
是愛大人的.
是愛弟兄姊妹的.
主我很相信.
當有一天洪映嘉真正地愛你.
愛弟兄姊妹.
這裡整個面貌就會改變.
這裡就會真正是一個很溫暖的地方.
是會吸引更多人去歸向你.
也會吸引更多人去尋找你.
主求你幫助我們.
叫我們五在一讓.
讓我們有真正的悔改.
也讓我們真的做一隻.
上帝你要賜福的.
你喜悅的綿羊.
無論有沒有人看著我們的生命.
你都讓我們為你而做.
為你而活.
奉耶穌基督的名字祈禱.
阿們.
\newpage



\section{以弗所書 4:7}
\label{sec:rTQz8Qod6Ek}
\textbf{2022.10.16 《基督所賜的恩》 以弗所書 4\_7, 11-16節 (和修版) 張美薇牧師}
\newline
\newline
連結: \href{https://youtube.com/watch?v=rTQz8Qod6Ek}{\texttt{https://youtube.com/watch?v=rTQz8Qod6Ek}} ~~~~ 語音日期: 2022-10-18
\newline
\newline
\hyperref[sec:KSHhb73QdyA]{\small{< < < PREV SERMON < < <}}
~
\hyperref[sec:index_chronic]{\small{[返順時目]}}
~
\hyperref[sec:index_scriptual]{\small{[返順卷目]}}
~
\hyperref[sec:ngXOvQd6nwI]{\small{> > > NEXT SERMON > > >}}
\newline
\newline
以弗所書 4:7
\newline
\begin{longtable}{cl}
\hline
\hline
章節 & 經文 (和合本修訂版)\\
\hline
4:7 & \begin{tabularx}{0.7\textwidth}{X} 我們每個人蒙恩都是照基督所量給每個人的恩賜。 \end{tabularx} \\ \\
[1ex]
\hline
\hline
\end{longtable}
$^{1}$大家好,大約半年就有機會來一次紅恩堂跟大家分享神的話語.
很高興見到大家,見到很多老朋友.
特別今天見到福音樂曲隊獻唱.
又見到輝哥和姊妹一起用音樂來敬拜.
我知道你們都有師班.
其實是三個不同的方向.
可以用樂曲,可以用敬拜讚美.
又可以用傳統的師班,傳統的詩歌來敬拜神.
真是很多元,亦看到不同的弟兄姊妹.
有不同的恩賜,都可以一起配搭來事奉.
真是叫神的命得成長.
今天老牧師叫我和大家分享恩賜.
所以我選了這一段經文.
來到才發覺原來你們今年的年提是毀身事奉.
大家用大家的恩賜來一起合作事奉.
今日希望和大家分享一個簡單的訊息.
希望大家可以聽得入去,亦可以走出來.
我們低頭祈禱.
天主上求你使用孩子今日的講道.
去幫助會中,叫他們有一個能聽的耳朵.
有一個能夠把你的話語藏在心中的心.
亦有一個明白你的說話的頭腦.
更加明白之後,用他們的手和腳走出來.
以致你的命可以被高舉.
你的教會得以建立.
聽我們的祈禱,奉耶穌基督的名求.
剛才說過,恩賜其實有很多種.
所以很多時候恩賜會帶來教會一個很豐盛.
很好的成長.
但如果處理得不好,會怎樣?會吵架.
所以我們今天看第一個重點.
就是會說恩賜的合一根基.
經文其實在第七節裡面.
說我們各人蒙恩,都是照基督所量給各人的恩賜.
所以首先說一點.
恩賜是從哪裡來的?.
恩賜不是我從小就生出來的.
又或者是從不同的地方學回來的.
這裡說得很清楚,恩賜的來源是基督.
可以看到是一個很合一的根基.

$^{41}$因此我們可以看.
恩賜的來源是基督.
在《以弗所書》第四章第四至第六節裡.
告訴我們原來來源是有一個合一的根基.
因為第四章第四節裡說.
身體只有一個,聖靈只有一個.
正如你們蒙召,同有一個指望.
一主一信一使一神.
就算眾人的父超乎眾人之上.
貫乎眾人之中,也住在眾人之內.
所以我們可以看到原來恩賜來源是基督的時候.
它有七個合一的根基.
今天重點不是說這裡,而是簡單地說.
身體只有一個,教會是基督奧秘的身體.
在《以弗所書》前面第一章已經看到.
身體只有一個,但有很多不同的肢體.
有眼耳口鼻,手腳等等.
肢體雖然多,但卻仍然只有一個身體.
身體是不可以分割的.
所以這裡告訴你,我們身體只有一個.
也告訴我們,聖靈只有一個.
因為這是身體的成因和內容.
聖靈只有一個位置.
所以身體和身體只有一個是不可以分割的.
所以其實基督透過聖靈給我們很多不同的東西.
同時我們也同有一個指望.
在《以弗所書》第一,第二章裡有提及.
我們夢照的時候.
聖靈告訴我們將來會有基業.
我們在神裡面會有得著.
直到我們得熟進入榮耀.
這是夢照的指望.
基於基督奧秘的身體.
我們在靈裡面更加明白指望的事.
我們知道末後的時候.
我們大家都有同一個指望.
就是我們知道末後每個人都會有我們的基業.
我們還有一主.
大家都知道我們只有一位主人.
一位救世主和一位生命的主.

$^{81}$就是主耶穌基督.
所以我們在生活上,在教會的屬靈裡.
也只有一位主.
我們只有一個信,因為恩信稱義.
我們不是相信行為.
我們是因為信.
我們相信耶穌基督是神的兒子.
祂為世人的罪被釘死在十字架上.
所以我們都是同一信念.
還有一洗.
雖然不同教會有不同的洗禮方式.
有些教會可能是嬰兒洗.
有些教會是成人洗.
有些教會可能是灑水禮.
有些教會是浸禮.
但是形式不是最重要.
最重要的是祂所表明的實體意義.
實質的意義.
我們同有一個洗禮.
最後第七個根基.
我們只有一位神.
一位神.
就是眾人之父.
超乎眾人之上.
貫乎眾人之中.
也住在眾人之內.
所以保羅看來.
其實基督徒不能夠合一.
是不可以理解的.
但是我們稍後會說到.
有很多恩賜的時候.
我們會發覺不能夠合一的東西.
已經存在.
我們很多時候都會有困難.
但是保羅在這裡說.
我們要保守聖靈所賜合而為一的心.
是需要保守的.
因為這個合一已經賜給我們.
是根基.
我們接著再去看.

$^{121}$剛才那段經文.
我們各人蒙恩.
都是照基督所量給每一位信徒的恩賜.
所以很重要的一點.
恩賜是基督量給信徒的.
每一位信徒都有.
因此是基督賜給每一位信徒.
基督是賜恩的主.
恩賜是基督量給我們的.
量就是為我們度身訂造.
人人都有恩賜.
就是說我們人人都可以用我們的恩賜來事奉.
我們每個人都有一份很獨特的恩賜.
是基督量給我們的.
經文中的「各人」是一個強調的說話.
每一個人.
我們每一個都是神要用的.
我們每一個都已經有恩賜.
我們每一個都要用我們的恩賜來服侍神.
服侍人.
意思是在神家裡是沒有閒人的.
我不知道弟兄姊妹.
你們在這裡恩賜有多少的發揮呢?.
或者從另一個角度來看.
你們有多少是有服侍.
或者是事奉的崗位或者機會的呢?.
一般來說教會裡面可能只有百分之二十的人服侍.
百分之八十又如何呢?.
回來聚會來享受別人的服侍.
其實已經有百分之二十的服侍已經很好了.
但是不夠的.
因為這裡說什麼呢?.
恩賜是基督照人的需要.
人的特性.
量給「各人」的恩賜.
每一個都有的.
每一個都有.
所以每一個信主的人都要各盡其職.
建立基督的身體.
主雖然說有說不盡的恩賜.

$^{161}$卻是照我們最適宜的量給我們的.
如果我們想要更多的恩賜.
我們也可以的.
我們可以向主求的.
所以當我們有恩賜的時候.
我們不需要互相比較的.
譬如我們當中有些弟兄姊妹.
很有唱粵曲的恩賜.
但有些弟兄姊妹就不是.
她是喜歡唱《經拜讚美》的.
不需要比較的.
不需要彼此說誰厲害.
誰好.
誰重要.
我們是要去明白.
要去接納.
譬如我自己.
不知道大家知不知道.
我自己的音樂恩賜是很差的.
我曾經想過去學一下.
我嘗試參加過師班.
不過當然那些弟兄姊妹是很好的.
他們沒有說我退出師班是最大的貢獻.
但我唱著唱著就知道不是.
那怎麼辦呢?.
我當然可以求神給我.
但結果我沒有求神給我.
我求神給一個有恩賜的弟兄姊妹跟我配搭.
那又是.
神給了一些很有恩賜.
很有音樂恩賜的弟兄姊妹跟我配搭.
當我去出外短暫的時候.
有姊妹幫我選歌.
也幫我領唱.
也跟我說的信息很配合.
所以我們不需要羨慕別人.
我們也不需要去強求.
因為神會在適當的時候.
賜下適當的恩賜.
或者賜下適當恩賜的弟兄姊妹去幫我們.

$^{201}$我記得我自己成長的時間.
我第一間教會.
差不多每一個弟兄姊妹都會去教小朋友.
因為那時候在美國.
每一間教會有兩個車輛.
他們會去唐人街載一班小朋友回來.
因為我們在唐人街那兩個教小朋友的幼稚園老師.
都是我們的姊妹.
所以那班學生就一起上來上課.
我初信的時候就去上過一次課去做助手.
那次我不知道大家是否明白.
我坐在那裡抱著一個小朋友坐了一個小時.
因為那個小朋友會吵.
以後我就不再教小朋友了.
後來成長了.
做了傳道人.
我去了中部去直課.
那時候教會.
我去到的時候.
教會那時只有十多人.
只有一個家庭有四個小朋友.
我們說不行.
教會一定要有小朋友的侍工.
於是我們招募.
一招募就有二十多個小朋友回來.
但沒有人教.
於是我就要去教.
其實是硬著頭皮.
因為我從來都沒有去教小朋友.
因為那次的經歷.
但我向神求的時候.
神又給我恩賜.
我就可以去做到.
後來慢慢有弟兄姊妹的加入.
讓教會那群小朋友得到很好的教導.
所以是真的.
是基督量給人的.
所以每一個信徒都有恩賜.
在基督裡面.
人人都有事奉的功能.

$^{241}$因為各人都得著一份獨特.
而且是量給我們的恩賜.
所以我們要發掘我們的恩賜.
我們要用我們的恩賜.
發揮我們的恩賜.
我們更加要發展我們的恩賜.
看看可以怎樣.
把我們的恩賜用得更多.
更深更遠.
所以我們會看到.
沒錯.
我們恩賜的根基是合一的.
我們恩賜的根基是耶穌基督.
為我們每一個人量給我們的.
然後我們就落到第二點.
第二點除了我們有合一的恩賜以外.
我們的恩賜是多元和豐盛的.
因為在第四章第十一節.
「他所賜的有仕徒有先知.
有傳福音的有牧師和教師」.
在這裡說了.
有人說是四種.
有人說是五種.
我們今天不要詳細的去悉經.
究竟牧師和教師是一種恩賜還是兩種恩賜.
很簡單來說.
「仕徒」比較窄的定義來說.
是基督的仕徒.
是基督所揀選見過耶穌的人.
他們今天已經不存在了.
但是這個在廣義來說.
「仕徒」就是說有聖靈給予他能力.
他是奉神的呼召.
猜險到去不同的地方.
去開荒建立教會的人.
今天廣義來說「仕徒」是這樣做的.
「先知」就是有聖靈的感動.
接受神的真理啟示.
能夠傳講神的心意.
去做就,安慰,勸勉信徒.

$^{281}$在神的話語上是很有能力的.
「仕徒」就是開荒的.
「先知」就是在神的話語上很有能力的.
傳福音就是將福音傳給一些沒有聽過福音的人.
好像今天的報道者或者宣教士一樣.
大家如果年紀差不多.
都會認識葛培里牧師.
葛培里牧師是否一個很有報道恩賜的報道者.
我記得有一次他.
人家說那天下大雨.
他在報道會裡面.
他一出來到台上就對會眾說.
朋友們其實你們都預備好了.
這麼大雨我不說了.
你們誰預備信主的就出來吧.
你猜有沒有人出來?.
結果那些人就湧出來.
好像跟他平時一路說的時候一樣.
當然我們知道.
因為他們整個報道團的預備工作做得好.
他們可能一年之前就已經開始.
去到他們要去的地方去負責這些工作.
但我同時看到他是很有這方面的恩賜的人.
另外牧師和教師就是一個牧養,教導.
教師就是教導神的話.
牧師就是關心牧養.
有些人會說是一次的恩賜.
有些人會說是兩樣的恩賜.
他這裡就說有這幾樣.
但其實聖經裡面都有其他幾個地方.
都有說不同的恩賜.
很簡單和大家說.
他說:這人蒙聖靈賜他智慧的言語.
那人也蒙聖靈賜他知識的言語.
又有一人蒙聖靈賜他信心.
還有一人蒙聖靈賜他醫病的恩賜.
又叫一人能行義能.
又叫一人能作先知.
又叫一人能辨別諸靈.
又叫一人能說方言.

$^{321}$又叫一人能翻方言.
在哥林多前書第十二章八到十節裡面說到.
所以除了這幾樣恩賜之外.
更加有智慧的言語,知識的言語,信心,醫病.
行義能,先知,辨別諸靈,說方言,翻方言的恩賜.
還有在羅馬書十二章六到八節裡面說.
按我們所得的恩賜各有不同.
或說預言就當照著信心的程度說預言.
或作執事就當專一執事.
或作教導的就當專一教導.
或作勸化的就當專一勸化.
施捨的就當誠實.
自理的就當殷勤.
憐憫人的就當甘心.
更加有講預言的,有作執事的,有作教導的.
有作勸化的,有施捨的,有自理的,有憐憫人的.
更加多了一點.
就是各人的恩賜都不一樣.
但是我們要用好我們的恩賜.
所以說預言就要照著信心的程度去講預言.
做執事就專一執事.
不要理會其他東西.
作教導的就專一教導.
勸化就專一勸化.
施捨就要誠實地施捨.
因為當施捨的時候可能跟前菜有關.
自理的就要當殷勤.
憐憫人的就當甘心.
當我們有這麼多元的恩賜的時候.
我們就要用好我們的恩賜.
不單止這樣.
還有在比特前書4章10到11節上的時候說.
「個人要照所得的恩賜彼此服事.
作臣百般恩賜的好管家.
若有講道的,要按著臣的聖言講.
若有服事人的,要按著臣所賜的力量服事.
叫臣在凡事上都因耶穌,基督得榮耀」.
所以我們剛才很快地讀了這幾段經文.
已經看到有二十多樣恩賜.
但很重要的一點就是.

$^{361}$有沒有人有所有的恩賜?.
又有沒有哪一種的恩賜是所有人都要有的?.
都沒有說.
所以有很多種的恩賜.
有很多不同的配搭.
在這裡說.
有這麼多種恩賜.
我想沒有誰會所有恩賜都有.
也沒有說某一種恩賜.
所有人都要有.
而最後重要的是.
恩賜是用來彼此服事.
是叫臣得榮耀.
就像比特前書第十一章說.
「要按著臣所賜的力量服事.
叫臣在凡事上因耶穌,基督得榮耀」.
所以我們每個人都有耶穌賜給我們的恩賜.
我們的恩賜是不同的.
有人有這樣,有人有那樣,有人有其他的.
但是我們每個人都要用我們的恩賜.
去用我們的恩賜.
所以下面更加地解釋清楚.
在《以弗所書》第四章第十一至十二節.
再重覆說.
《以弗所書》說.
「他所賜的有仕徒,有先知,有傳福音的.
有牧師和教師.
為要承傳聖徒.
各盡其職.
建立基督的身體」.
所以他提到.
恩賜的目的是做什麼呢?.
是要建立基督的身體.
然後他在第十二至十六節裡.
告訴我們.
我們可以怎樣做.
才可以建立基督的身體.
基督的身體是什麼?.
就是教會.
他要做的是什麼呢?.

$^{401}$這裡就說第一件事.
他說.
「直等到我們眾人在真道上同歸於日.
認識神的兒子.
得以長大成人.
滿有基督長成的新娘」.
所以他第一件事.
我們怎樣去承傳聖徒.
各盡其職.
建立基督的身體呢?.
第一件事就是.
在真道上同歸於日.
直等到我們眾人在真道上同歸於日.
所以他所說的.
承傳聖徒各盡其職.
是牧師傳道的人要做.
教會的領袖要做.
還是每一個信徒都要做呢?.
每一個信徒都要做.
所以我們每一個都要做.
剛才說過20\%是不夠的.
我認為雄恩是比較好的.
我認為你們不止20\%.
但是有沒有100\%呢?.
我想你們未必有100\%.
但是你們要明白.
每一個信徒都要長大成熟.
這裡的長大成熟.
就是說在真道上長大成熟.
他說在真道上同歸於日.
認識神的兒子.
得以長大成人.
滿有基督長成的新娘.
是不是單單在神的話語上成長呢?.
其實是不夠的.
我看到其實有些人的恩賜.
是神的話語.
所以有些人信主.
是因為聽到神的話語而信主.
但有些人的恩賜是愛心.

$^{441}$大家想想有沒有一些很有愛心弟兄姊妹在我們當中呢?.
如果你發覺有些新朋友.
或者有些朋友病了.
很多時候會去探望他.
會幫他.
用基督的愛心去鼓勵他.
結果他就信主.
有沒有呢?.
都有的.
他都信主.
有些人是信心.
有些頂尖是很有信心的.
祈禱很有信心.
很多時候都會用他的信心.
去幫助一些朋友.
以致那些朋友就會信主.
但我們不只是停留在我們信主的那一刻.
假如我們是神的話語而信主的話.
我們也要在神的愛心上去長大.
因為耶穌也有愛心.
耶穌也會給我們信心.
所以當我們有神的話語的時候.
我們要在愛心上增長.
也要在信心上增長.
或者大家是由神的愛心而信主.
加入教會的話.
你也要在神的話語上受造就.
我們要在真道上受造就.
我們不可以只是在愛心上成長就夠了.
我們更加要在信心上都要成長.
又或者有些人是以信心信主.
我們同樣要在愛心上成長.
要在神話語上成長.
這樣大家就慢慢去想像.
我畫圖畫得不好.
畫得不好就未必能做到我想做的.
大家就會想像.
神的話語上慢慢擴大.
然後愛心慢慢擴大.
信心慢慢擴大.

$^{481}$到最後就三個合在一起.
中間有一部分三個都有.
這個就是我們各人都要這樣的時候.
在真道上同歸於一.
長大成熟.
去認識神的兒子.
認識神的兒子.
認識神的話語.
認識神的愛.
更加認識神的信實.
我們怎樣可以用信心去跟從神.
去幫助我們.
我們都要用這些恩賜.
去幫助教會所有的弟兄姊妹.
讓每一個都是這樣去長大成熟.
是有耶穌各方面給我們的恩賜.
沒錯,我們有些可能是強過其他人.
但我們都要有所認識.
我們都要有所明白.
所以這個就是在真道上同歸於一.
長大成熟.
這個就是讓我們成全聖徒.
各盡其職的其中的第一點.
第二點就是第十四節.
「使我們不再作小孩子.
眾了人的詭計和欺騙的法術.
被一切異教之風搖動.
飄來飄去就隨從各樣的疑端」.
剛才說了正面我們要怎樣.
在各方面都要認識耶穌基督的人.
以致我們可以在真道上.
我們大家都是同樣歸於一體.
第二個就是會說到.
我們要在信仰上.
不再像小朋友一樣.
眾了人的詭計.
以致被一些異端或異教.
叫我們跌倒.
然後被吸引離開正統.
所以一樣我們要在真道上成長.

$^{521}$我們亦要很小心.
有愛心的弟兄姊妹.
亦要幫助弟兄姊妹成長.
有信心的弟兄姊妹.
都要幫助.
這樣我們才可以讓弟兄姊妹.
不隨從各樣的疑端.
第一點就是在真道上.
同歸於一長大成熟.
第二點就是不隨從各樣的疑端.
正面的我們在耶穌上學習.
學習他的話語.
學習他的愛心.
學習他的信念.
然後我們不隨從各樣的疑端.
第三點就是.
「唯用愛心說誠實話.
凡事長進.
連於元首基督」.
顏色是.
我記得.
「全心都靠祂聯絡得合適.
百節各按各職.
朝著各體的功用彼此相助.
便叫身體漸漸增長.
在愛中建立自己」.
所以第三點.
我們如何用恩賜的目的呢?.
就是各按各職.
用愛心去建立身體.
說了這麼久恩賜.
但是愛心亦很重要.
為什麼呢?.
恩賜會令我們驕傲.
你有沒有見過.
有人很熟悉神的話.
他會很驕傲.
看不起一些不太熟悉的人.
有些人在某一方面很有恩賜的話.
他也會因為這個緣故.

$^{561}$而輕忽了其他人.
我想大家都有見過.
所以這是很重要的.
「各按各職」.
然後要用愛心說誠實話.
有恩賜還要有愛心.
然後更重要的是.
「連於元首基督」.
凡事長進.
「連於元首基督」.
我們有一個目的.
目的就是要讓基督的身體強大.
無論你的恩賜是哪方面.
剛才列了二十多種恩賜.
相信還有很多恩賜在聖經裡沒有說.
但是我們要很明白.
恩賜是用來建立基督的身體.
第三點更加是用愛心來建立基督的身體.
我想大家都聽過五千兩的比喻.
耶穌說的比喻.
主人派五千兩給他.
有五千兩的人是怎樣的?.
他們說:真好,我這麼多.
我就看不起其他人.
不是,他們很努力地用他們的五千兩.
結果他們是怎樣?.
賺多了五千兩.
有些人是二千兩.
二千兩沒有說:真是論盡了.
我只有二千兩這麼少.
我不知道怎麼辦.
沒有,他們一樣很努力地賺多了二千兩.
所以無論你的恩賜有多少.
你有哪一樣的恩賜.
你都要好好地用它.
用盡它.
也要幫忙讓基督的身體能夠一起成長.
我想你們都可能聽過.
有些小朋友去看眼科醫生.
就發覺他有Lazy eye.

$^{601}$有沒有聽過?.
甚麼叫Lazy eye?.
即是有一隻眼可能經常都不用.
明明沒事的.
但他經常都是用某一隻眼.
第二隻眼就沒有用.
Lazy eye.
醫生的醫治方法是甚麼?.
就是找一個黑的蓋.
蓋著他的眼.
蓋著他的甚麼眼?.
好的眼.
蓋著好的眼.
那怎麼辦?.
就讓不好的眼去用.
所以假如你去到一個情況.
你剛剛開始成長.
你有那個恩賜.
你要去發揮去用.
你就要好像那隻Lazy eye一樣.
要努力地去用它.
但假如你又是那一件事.
在團隊裡.
你是很精的.
很好的.
你都要被蓋著.
為甚麼?.
去讓你團隊裡面.
另外的一部份.
可以更加去運用到他的恩賜.
以致彼此將可以一起配搭.
一起成長.
所以Lazy eye是有弱的那隻.
但亦有強的那隻都要配合.
所以強和弱一起配合的時候.
他就怎樣?.
「唯用愛心說誠實話.
凡事長進.
連於元首基督.
全身都靠他聯絡得合適.

$^{641}$百節各按各職.
照著各體的功用彼此相助」.
結果怎樣?.
「叫身體漸漸增長.
在愛中建立自己」.
弟兄姊妹.
這個就是用我們多元的恩賜的方法.
神給我們每一個人都有恩賜.
是給我們糧食.
而我們要用它.
我們用的時候.
要用愛心去用.
要彼此的配搭.
是一起在教會裡面.
是讓教會來成長.
更加是用教會.
以愛心的成長.
結束的時候.
跟大家說一個故事.
這個故事其實是真的.
是十九世紀的時候.
教會歷史上的一個例子.
當這個小朋友.
當時他十二,三歲的時候.
他比較窮.
他每個週日都走路回教會的時候.
發覺教會前面的路很濕.
所以人們很難走.
很泥濘.
因為濕的泥路.
十九世紀的時候.
石屎的路不是那麼多.
於是他每個星期.
用少少錢買幾塊磚.
去砌一塊.
第二個星期再買幾塊磚.
又再砌多一塊.
他砌了幾個星期之後.
牧師看到.
牧師看到之後就問他.

$^{681}$為什麼你會這樣做呢?.
他就跟牧師說.
我想砌好那條路.
但沒有錢.
所以我每個星期砌少少.
牧師聽到之後就很受感動.
然後在崇拜的時候.
就去為他感恩祈禱.
然後在講道的時候.
就將他的例子講了出來.
結果怎樣?.
結果會友就說.
原來一個這麼缺乏的小朋友.
他也願意為教會.
為神的家有這樣的貢獻.
於是很快人們就籌到款.
那條路就建了一條很好的路.
變成大家去教會.
再不用像慣來慣去的.
又或者會積到整個腳都是泥.
那麼笨拙.
這個小朋友長大之後.
他也成為一個很好的基督徒.
他那時候十幾歲.
到了七十幾歲才離世.
他十七歲就開始做主一學的導師.
他十幾歲到十七歲.
他就做主一學的學生.
然後十七歲開始.
他就做主一學的導師.
一直到他七十歲離世的時候.
因為他一生的座右銘.
就會說我要先求神的國.
和神的義.
所以他不是主一學的學生.
就是主一學的導師.
不是主一學的導師.
他又會回去做主一學的學生.
後來他做生意做得很好.
很成功.

$^{721}$他在費城開了一間.
美國很大的百貨公司.
他有很漂亮的管風琴.
但是當他要預備開張的時候.
他才發覺原來慕迪先生.
在費城開報道會.
於是他就延遲他的開幕.
讓慕迪先生可以借用他的商場.
來開報道會.
然後他才開幕.
所以我們看到他一生的.
是為神奉獻.
一生的將他所有的.
用出來.
我講道通常都會.
最後是這樣說.
「於我何干」.
剛才你們都聽了.
大概四十五分鐘左右.
這篇道對你有什麼關係呢?.
你是否懶惰要掩蓋那一種.
還是你是懶惰要努力增大那一種呢?.
你會不會去發掘.
耶穌基督給了你什麼恩賜.
然後你發揮.
然後繼續為你的恩賜去發展.
假如你不是很清楚.
你有什麼恩賜.
你可以找魯牧師傾談.
他一定可以為大家開一個講座.
有些恩賜調查.
讓你去找回自己的恩賜.
願意每一個信徒.
都能夠成為主所能用.
因為恩賜是為了裝備每一位信徒.
做自奉的工作去建立基督的身體.
我們低頭祈禱吧.
天父上帝願你使用洪恩嘉的弟兄姊妹.
幫助他們.
讓他們都可以找到自己的恩賜.

$^{761}$也發揮自己的恩賜.
最後也可以為自己的恩賜去繼續發展.
讓他們成為一間蒙愛的教會.
教會能夠長大成熟.
能夠在你的命上被高舉.
也可以去榮耀你的名.
聽我們的祈禱.
奉耶穌基督的名頭.
阿們.
(影片完結).
\newpage



\section{提摩太前書 3:14-16}
\label{sec:ngXOvQd6nwI}
\textbf{2022.10.23 《屬靈的家》 提摩太前書 3\_14-16 (和修版) 謝靄薇牧師}
\newline
\newline
連結: \href{https://youtube.com/watch?v=ngXOvQd6nwI}{\texttt{https://youtube.com/watch?v=ngXOvQd6nwI}} ~~~~ 語音日期: 2022-10-27
\newline
\newline
\hyperref[sec:rTQz8Qod6Ek]{\small{< < < PREV SERMON < < <}}
~
\hyperref[sec:index_chronic]{\small{[返順時目]}}
~
\hyperref[sec:index_scriptual]{\small{[返順卷目]}}
~
\hyperref[sec:KbMX2DPfTSA]{\small{> > > NEXT SERMON > > >}}
\newline
\newline
提摩太前書 3:14-16
\newline
\begin{longtable}{cl}
\hline
\hline
章節 & 經文 (和合本修訂版)\\
\hline
3:14 & \begin{tabularx}{0.7\textwidth}{X} 我希望盡快到你那裡去，所以先把這些事寫給你； \end{tabularx} \\ \\ \relax
3:15 & \begin{tabularx}{0.7\textwidth}{X} 倘若我延誤了，你也可以知道在神的家中該怎樣做。這家就是永生神的教會，真理的柱石和根基。 \end{tabularx} \\ \\ \relax
3:16 & \begin{tabularx}{0.7\textwidth}{X} 敬虔的奧祕是公認為偉大的：神在肉身顯現，被聖靈稱義，被天使看見，被傳於外邦，被世人信服，被接在榮耀裡。 \end{tabularx} \\ \\
[1ex]
\hline
\hline
\end{longtable}
$^{1}$知道你們快要十一周年了.
預祝你們在這個地方.
洪水橋的地方是很有作為的.
繼續為主成為燈台.
為主發光.
代表舞堂問候大家.
大家一起彼此守望.
讓神的家更加的興旺.
我們一起一起祈禱.
每有恩典 每有慈愛的主.
每當我們來到您的殿的時候.
我們就歡喜快樂.
因為您就在當中.
您的靈眼當中運行.
與主您這個時候帶領我們的心.
若然我們有什麼是不蒙您喜悅的.
求聖靈在當中剔除.
以致我們能夠看到這個話語的時候.
我們就明白您的心意.
求主親自的祝福.
又為到孩子口中的言語.
心中的意念能夠蒙您悅立.
我們這樣祈禱.
教託是奉主耶穌基督的名求.
阿們.
相信大家應該是很久沒有去旅行了.
不過現在隔離的次數是零加三.
應該是很有機會了.
聽到很多人都已經準備.
或者11月份開始有很多人準備去旅遊了.
相信無論你去哪裡玩.
就算你住五星級的酒店也好.
但你總覺得家裡最好.
所謂的農場都不如你的狗窩舒服.
所以家是帶給人很溫暖的.
很安息的地方.
很安全感.
很輕鬆的.
不少年輕人他們很喜歡在教會.
表面上看來好像很好.

$^{41}$其實很多時候我們會勸年輕的導師.
你要看看他們到底和家人的關係是怎樣.
會不會趁機不想做家務就跑來教會呢?.
以致可能他經常回教會.
但是家人很反感的.
特別是過時過節.
都不可以忽略了做子女的本份.
又或者因為教會是我們的家.
所以很多人都會把家裡不要的東西.
就給教會.
但是教會又不是收次貨的.
我們自己都不要的.
我們為甚麼要給神呢?.
所以雖然我們去想這些的時候.
但是我們想一想.
其實教會是我們的家.
其實應該是充滿神的愛和溫馨的.
所以我們要將這份愛和人分享.
今天我會和大家分三方面去想.
我們擁有這個屬靈的家.
到底和教會有甚麼關係呢?.
在提摩太子主三章15節上半部就說.
倘若我認了.
你也可以知道在神的家中.
該怎樣做.
其實保羅原本是要去伊弗所的.
途中有一個特洛菲摩病了.
他在美麗都的行程稍為停留在那裡.
但是他就在那個時候被捕.
就沒辦法和提摩太子見面了.
提摩太子當時很需要保羅去教他.
當時他們可能有新的長老.
執事去成立.
或者當時有些異教的人.
說了一些誹謗的話.
他很需要保羅這個有使徒身份的人.
有權威的話.
可以去對付這些誹謗的根據.
因此保羅雖然要去伊弗所.
但他仍然先寫了一封信去提提摩太子.

$^{81}$就在14節那裡說.
我希望盡快到你那裡去.
所以才把這些事寫給你.
他可能有預感.
可能沒辦法見到提摩太子.
所以他才寫了提摩太子的前書.
這家就是永生神的教會.
真理的儲石和根基.
其實這裡也有意思.
叫你回想一下我們舊約.
舊約的聖殿多宏偉.
作為神的家.
是這麼堅固的建築物.
曾經有傳說人說了一句標語.
他說基督是我們的生命.
教會是我們的生活.
教會的生活是什麼呢?.
是不是一個星期回來崇拜.
用一個半小時.
或者上班就星期三.
或者在週間的時候去練習.
或者星期三又去祈禱會.
等等.
加起來.
是不是這些就是生活呢?.
教會的生活呢?.
潘博華在《團契生活》這本書裡.
就說到互相扶持的關係.
他說這應該是一切基督徒團契的目標.
就是要互相扶持的關係.
其實神的旨意就是要在人的口中.
和在弟兄姊妹的見證裡.
去見到神的道.
這是最重要的.
當人灰心沮喪的時候.
有人去關心他們.
去和他們做見證.
免得他們對真理有誤解.
而跌倒.
這也很重要.

$^{121}$教會是什麼呢?.
教會是神的家.
其實是被神呼召他們.
聚在一起的一群人.
就成為教會了.
所以教會不是一座建築物.
是被神呼召的一群人.
走在一起.
是有耶穌基督的生命的.
這群人才叫教會.
哥林多前書十二章裡說到.
基督徒是從一位聖靈受洗.
成立一個身體.
我們看來教會是一個很有生命力的.
是一個生命的體.
楊尼雅牧師曾經以團契來說教會的實體.
用團契來形容.
團契就是一種關係.
是一個活生生的人際關係.
不是只是一個組織.
不是一個群體這麼簡單.
夏東堅牧師也說.
聖經裡對教會並沒有一個很明確的定義.
但是有很多象徵.
讓我們明白教會的圖樣是怎樣的.
例如一個身子.
彼此做肢體.
大家也看到經文裡.
叫我們活石要建造聖靈宮.
既然大家同為肢體.
就不可以單獨存在.
一個人成為基督徒.
也不可以以個體的方式去存在.
就像你的耳朵.
你的眼睛.
不能獨立去生存.
要在耶穌基督的身體裡去成長.
我們每個弟兄姊妹.
無論你是手還是腳.
我們也不可以說.

$^{161}$我的腳走得快一點.
我怎樣去表達.
我厲害一點.
不可以的.
你不可以離開身體.
大家要一起合作.
所以教會裡.
不是看我要安排甚麼節目.
到時候要做甚麼內容.
而是怎樣使教會裡的弟兄姊妹.
怎樣在生活上產生關聯.
這是很重要的.
就是所謂的…….
為甚麼好像沒有反應.
幫我弄下一張.
就是要成為彼此相愛的團體.
正如大家很熟悉的.
《約翰福音》十三章三十四到三十五節.
我賜給你們一條新命令.
乃是叫你們彼此相愛.
我怎樣愛你們.
你們也要怎樣彼此相愛.
你們若彼此相愛.
眾人因此就認出.
你們是我的門徒了.
彼此相愛很重要.
曾經有一個姐妹搬家.
由於聯絡上出現了問題.
她的丈夫不在身邊.
那怎麼辦呢?.
我立即打電話給她的組長.
組長就是教會裡的長老.
很快就找到組員.
快到她家.
很快就幫她把所有行李搬上車.
之後長老就帶大家祈禱.
祈禱之後.
很奇妙的事發生了.
原來搬家的司機.
看到他們在祈禱.

$^{201}$於是他很感動.
他一方面拿手巾抹眼淚.
一方面拿錢出來.
給長老說.
「我很久沒有回教會了」.
「我現在要立志重新回教會」.
「那些錢你幫我帶到教會裡」.
「當作我的奉獻吧」.
「你看到原來」.
「在屬靈的分享裡」.
「也包括了物質」.
「在當中分享的」.
「你關心弟兄姊妹的需要」.
「原來在神眼中是很寶貴的」.
「很有價值的」.
「也能吸引別人」.
「去到神的面前」.
因為《雅各書》第二章14-16節.
那裡怎麼說呢?.
「有信心」.
「沒有行為是死的」.
看到弟兄姊妹的需要.
有缺乏的時候.
你不是說「平平安安去吧」.
但你也沒有什麼表示.
「我為你祈禱」.
但你自己好像很冷漠.
好像與你無關.
我們好像很熟悉的.
聖經說「你平平安安去吧」.
但原來我們也沒有心去為他祈禱.
也不是有些行動.
想想我們怎麼可以.
急忙去幫他呢?.
另外就是「彼此相顧」.
「彼此相顧」就好像.
臨前十二章那裡說到.
肢體之間有些是內臟.
有些是外面看到的.
眼耳口鼻那些.

$^{241}$聖經特別說到.
原來神配搭身子.
是加倍體面那些.
有缺陷的肢體.
免得身體不協調.
總要肢體彼此照顧.
正如保羅當年.
用馬其頓教會的捐獻.
做一個例子鼓勵弟兄姊妹.
因為馬其頓教會他們很熱心.
但他們當時的環境很缺乏.
他們在極窮的情況下.
他們知道別人有缺陷.
他們也很慷慨的.
所以參與了一份.
希望保羅給他們一份.
你看看能力這麼困難.
他也要求希望可以參與.
如果能力可以的.
是不是更加要放在自己.
可以做到的.
去保護不足呢?.
所以他說.
趁現在的富裕.
去保護他們的不足.
將來他們的富裕.
又可以保護他們的不足.
有人曾經說地上的人際關係.
會消失.
但主要是肢體的關係.
是永遠不會消失的.
第二方面.
我們去看看.
原來耶穌基督所救贖的.
奧秘的身體.
在《波蘭多傳說》的三章十六節.
很明顯地看到.
他說:敬虔的奧秘.
是公認為偉大的.
神在肉身顯現.

$^{281}$被聖靈稱義.
被天使看見.
被傳於外邦.
被世人信服.
被接載榮耀裡.
我們兩章讀了這麼多.
我們稍後也會再講.
最重要的是.
認識一下.
什麼叫敬虔的奧秘.
其實如果你看看《以弗所書》.
它會讓你讀一次.
可能你會有一點感覺.
可以認識更多.
關於奧秘.
《以弗所書》的第三章第五節.
到十一節是這樣說的.
「這奧秘在以前的世代.
沒有讓人知道.
像如今聖靈向他的聖使徒.
和先知啟示一樣.
就像外邦人在基督耶穌裡.
接受福音得以同為後製.
同為一體.
同為蒙應許的人.
我作了這福音的伯役.
是照著神的恩賜.
是照著他運行的大能.
賜給我的」.
「雖然我比眾使徒中.
最少的還少.
但如此我這恩典.
有我把這基督.
那差不多的豐富.
傳給外邦人」.
「由此眾人都明白.
什麼是歷代以來.
隱藏在創世萬物之神裡的奧秘」.
在這一節你會看到.
保羅的心態.

$^{321}$自從在大馬色路上.
他見到耶穌基督之後.
他知道自己很不配.
他立刻很謙卑.
在世俗裡他是最小的.
但神有一個恩典給他.
讓他向外邦人傳播音樂.
所以這也是一個很大的奧秘.
「為要在現今.
接受教會是天上.
執政的,掌權的.
知道神百般的智慧.
這是照著神在我們主.
基督耶穌裡所完成的.
永恆的計劃」.
我們看到.
第五節到十一節裡.
他講到不斷地形容.
我們看到原來以前.
以色列人覺得他是神的選民.
原來神的目的是要讓全世界.
不只是神的選民去認識主.
是外邦人.
今天你和我.
我們都是外邦人.
我們可以去認識主.
而今天這個十六節.
是更加特別的.
因為敬田奧秘其實是講到.
在創世之前.
已經隱藏了神的計劃.
是要透過耶穌基督的一生.
還有耶穌怎樣對神的敬田信仰.
還有他的生活方式去表達.
怎樣去啟示奧秘.
而十六節就很重要了.
十六節可能是初期教會裡.
一首相當普遍的詩歌.
保羅就拿出來強調.
所有基督徒都應該相信.

$^{361}$和肯定這個信仰.
所以十六節.
他就將福音真理的啟示.
分成六個階段.
他就按照這個順序.
逐個逐個讓我們知道.
「神在肉身顯現」.
這是第一個.
大家很清楚.
耶穌降生.
「道成肉身」成為人.
你見到就是.
他講到耶穌基督的「道成肉身」.
「神在肉身顯現」.
祂沒有犯罪.
但祂甘願為我們.
這個「義」的代替不義.
為我們死在十字架上.
去受死.
聖靈怎樣使祂從死裡復活.
第二就是被聖靈稱義.
你見到就是.
天使都見到祂復活.
見到耶穌復活.
按照聖靈.
這個神聖的聖靈講.
因從死人中復活.
用大能去顯明.
耶穌基督是神的兒子.
你看看耶穌基督被判死刑.
而聖靈使祂復活.
這樣就是從羅馬的控訴裡.
被稱為義.
本來是死刑的.
但復活了.
被稱為義.
耶穌基督復活證實了神的權能.
神的能力.
第三就是被天使看見.
天使報佳音.

$^{401}$開始的時候.
向牧羊人告訴人.
救主降生.
開始的時候.
天使已經看見.
然後耶穌基督.
在黑西瑪利園祈禱.
天使從天上顯現.
加添了力量.
在路加福音22:43.
是這麼說的.
然後.
你會看見.
耶穌勝過死亡.
復活的主.
所以.
耶穌基督復活了.
同時天使也預言.
有一天.
祂升天.
祂會駕雲降臨.
祂升到城裡.
然後.
被傳於外邦.
基督復活.
傳於外邦.
使今天不是以色列人.
不是神的選民.
都可以聽到福音.
世人聽到傳講的奧秘.
這麼大的奧秘.
外邦人都可以領受.
教會也有一個功用.
就是宣教.
被世人信服.
我們今天每個人.
都可以聽到福音.
可以得救.
這是神很大的恩典.
最後就是.

$^{441}$被接在榮耀裡面.
全部都是講耶穌基督的.
祂的一生裡面.
祂的工作.
耶穌基督升天了.
祂現今在天上.
祂擁有一個.
「全然勝利而無比榮耀」的位置.
基督是世人的救主.
祂也是主宰著我們.
將來耶穌基督要回來.
祂回來的時候.
就要進入永遠榮耀的國度.
我們這些同蒙應許.
同為一體.
在神的家裡面.
在耶穌基督的身體裡面.
每一個基督徒.
我們怎樣去承擔這個使命呢?.
我們要忠於主的託付.
剛才我們讀過.
「這家就是永生上帝的教會」.
我們清楚了.
我們特別講到.
真理的儲錫和根基.
耶穌基督就是真理.
而祂所買熟的教會.
因為教會就是耶穌基督的身體.
我們每個肢體.
就在這個教會裡面.
是被祂買熟的.
每一個就是儲錫和根基.
在裡面.
在屬靈的一家裡面.
我們每一個成員.
聖經說.
你們作為活石.
要被建造成為靈工.
或者現在我們和修館.
就是成為屬靈的殿.

$^{481}$講到活石.
就是說.
你是會動的.
活石.
你記得你是活石.
所以你就要參與事奉.
你不要坐在那裡.
好像大爺一樣.
有人招呼你.
可能你剛開始來的時候.
你什麼都不懂.
不要緊.
但是慢慢地你會懂的時候.
你就不是大爺了.
你就是家中的一份子.
你要去做事.
去服侍別人.
你要活石.
去參與事奉.
在這裡.
洪水橋.
你看那些房子.
真的很有發展.
將來很多人會來到這裡.
你要好好裝備.
去服侍別人.
去教別人.
令到他也像你一樣.
成為一個活石.
可以參與事奉.
所以在這裡.
就要看看我們事奉的態度.
房子裡.
有些很重要的支柱.
柱石和根基是不能動的.
是不是?.
很多時候.
人們貪漂亮.
砍了它.
就彎了.

$^{521}$柱石和根基是不能動的.
但是房子裡怎樣動都可以.
怎樣設計.
怎樣美化都可以.
但是.
重要的柱石和根基是不能動的.
不能拆掉.
很重要的.
因此教會就有責任.
去維護.
真理的存存.
有責任.
不要讓世俗.
一些不好的.
一些處事的方法.
去污染了.
很重要.
所以弟兄姊妹.
你對真理的維護.
堅決的態度.
就好像柱石和根基.
是這樣的.
穩固的.
我們環顧也好.
但是就要這麼固執.
去維護真理.
還有呢.
要各盡其職.
當你的屬靈生命成長了.
成長了之後.
你要成熟的.
不要成長了.
頭大身小都不行.
你要成熟.
成熟之後.
你又要裝備.
裝備最辛苦的.
要付代價的.
我們不喜歡記事物的.
但是沒辦法.

$^{561}$都要記一些方法.
怎樣去和人溝通.
或者怎樣將福音告訴別人.
很簡單的.
福音你要去熟習一下.
你又不可以學了就算了.
學了會很容易忘記.
特別是我們年紀大了.
會忘記的.
你要用.
越用就越有.
好了.
你用了就各盡其職.
然後呢.
神又給你不同的財果.
不同的恩賜.
你就去用.
就像剛才阿寧唱.
他們可以用聲音.
怎樣去讚美神.
我們又要去訓練.
訓練神給你這麼好的恩賜.
就不要浪費了.
你不用就沒有了.
不用了神就拿了給別人.
所以我們要用.
要發揮你的功用.
然後你可以和全部人一起去配搭.
一起去事奉.
這個身體.
就可以建立得很好.
很穩定.
當你身體強壯的時候.
那你就可以有更多人有能力去傳福音.
得救人數就增加了.
不單在裡面很扎實.
外面也不斷有人進來的時候.
或者你去宣教的時候.
這樣就有增長.
所以是很重要的.

$^{601}$還有呢.
我們要去迎接挑戰.
迎接挑戰的時候.
我們要看看我們的態度是怎樣.
我們的態度.
是有謙卑的心.
我們信主很久的時候.
忘記了我們的身份.
我們大家都是一樣的.
都是神.
因為我們是罪人.
大家都沒有什麼分別的.
無論你是多高的職位也好.
但是在神體我們是罪人.
但是神是拯救了我們.
我們進來這個群體的時候.
我們又不要擺盤.
我們大家都一樣的.
不分彼此的.
因為我們都是不配的罪人.
但是神不念.
我們有機會成為祂的兒女的時候.
我們就要學耶穌的謙卑.
然後去進入人群裡面.
在人群裡面去服侍.
因為人群裡面有很多不同的人.
祂不認識你的.
但是你不可以.
我以前是做經理的.
但是今天你去人群裡面.
去傳播音的時候.
你都要好好的跟祂說話.
你不可以說我等你招呼我.
不可以的.
那你怎麼可以有人跟你溝通.
所以你都要學耶穌的謙卑.
耶穌是天上的君王.
大家要共悲來到地上.
我們都要學祂.
怎麼虛己.

$^{641}$然後走進人群裡面去傳播音.
所以認罪的挑戰都很重要.
特別是在疫情之下.
是很難傳播音的.
當你有機會的時候.
你已經裝備好的時候.
隨時有人可以靠近的.
可以做事的.
所以我們要好好裝備自己.
也是要存有一個謙卑的心.
我們才可以做事的.
還有呢.
要勇於冒險.
冒險是一個不容易做的功課.
我們都很怕死的.
冒險.
但是馬達福音很真實的.
他說呢.
所以你們要去使萬民做我的門徒.
奉父子成人的名.
給他們施洗.
凡我所吩咐你們的.
都要教導他們遵守.
看啊.
我天天與你們同在.
直到世代的終結.
始終我們吩咐門徒.
要他們出去.
你們有沒有出去.
出去去哪裡呢.
有沒有聽耶穌的吩咐呢.
要去.
要走入人群當中.
去靠近他們.
還要影響他們.
不是給他們影響.
要去影響他們.
要記得去影響他們.
通常那些人.
好像是做保險的.

$^{681}$找自己人.
那你剛剛傳播音.
個個不懂.
不知道怎麼辦.
找自己人.
可以給他們練習.
找自己人.
慢慢你就有膽量.
自然你由認識的人開始.
然後你會見到.
你只要去.
就有一個影響.
接著你又不可以懶惰.
你要去跟他們傳播音.
要去教他們.
使他們成為門徒.
你要清楚地去傳播音.
去幫助他們.
如何克服信仰上的障礙.
最後主就應許.
要我們同在.
我們做了這麼多事.
他就應許要我們同在.
當我們聽他說.
直到世代的終結.
我們有沒有盡我們本份.
去完成這個大使命呢.
主耶穌向我們保證.
當我們去做這個計劃的時候.
他要我們同在.
我們將計劃付諸行動.
我們要去.
這樣我們就會對這個世界.
造成很大的影響.
要將這個消息傳出去.
只要你肯做.
事情就發生了.
當日.
耶穌基督和他的門徒.
門徒有十二個.

$^{721}$不過有一個就不爭氣.
掉頸了.
這十一個.
耶穌和基督的經文裡面.
當時對這十一個人說的時候.
你會覺得怎樣.
對十一個人說.
今天何恩堂有多少人呢.
如果經文對你說.
你會說浪費力氣.
我和旁邊的人都說不定.
但當日耶穌對這十一個人說.
你怎樣去影響這個世界.
你可能說.
天方夜譚.
當年.
蘋果電腦的創辦人.
大家很熟悉的.
喬布斯.
他的公司成長得很快.
他要找一個人.
可以幫他去管理.
如果一個領導的人幫他去管理.
於是他就找了.
約翰·斯卡利.
這個人當時在百事可樂工作.
他找他吃飯.
吃完飯之後.
這個人都沒有什麼反應.
好像沒有什麼興趣.
但他不忿氣.
他帶他去中央公園的大廈頂樓.
然後和他談.
去遊說他.
你去加入我們蘋果電腦吧.
然後他最後.
就看著他.
有一個很嚴肅的問題.
他和他說.
你到底要將你的餘生.

$^{761}$放在糖水裡.
還是你希望有機會去改變這個世界.
後來斯卡利的書裡就寫到.
當他聽到這麼有挑戰的話之後.
他就一語驚醒蒙眾人.
於是他就離開了百事可樂.
去加入蘋果電腦.
各位弟兄姊妹.
今天神也給你這份渴望.
祂希望你能夠去改變這個世界.
當然電腦影響世界的程度.
就不會像領人和基督建立的關係這麼強烈.
當人被引到神的恩典裡.
當一個孤獨的人.
他經驗到耶穌基督這麼愛他.
那種豐盛.
當人的良心受到責備.
一個罪人.
他的罪得到赦免.
他心裡重新找回清潔的良心.
他可以重拾人生的目標.
這是一個影響.
一個很強烈的連鎖反應.
你看看我們的疫情.
由兩個人開始.
數目飆升到十萬.
然後慢慢降低.
都是四個數字.
多厲害.
但今天.
你有沒有懸疑的心呢?.
你有懸疑的心.
你就離開你的安書區.
冒險吧.
不要怕死.
不畏縮的.
去為主做見證.
你願不願意呢?.
今天你跟神說.
靠著耶穌的幫助我願意去.

$^{801}$大家試著說吧.
我告訴神.
靠著耶穌的幫助我願意去.
請你跟旁邊的人說.
靠著耶穌的大能.
這個就不同了.
你跟旁邊的人說.
靠著耶穌的大能.
你要改變一些人的生命.
靠著耶穌的大能.
你要改變一些人的生命.
現在洪恩堂.
要慶賀十一週年.
你們長大了.
你們成熟了.
要起來了.
不要再坐在這裡了.
要起來了.
去影響一些人.
我相信這個地方的發展.
都是神的心意.
希望我們不要看.
我們現在這麼少人.
當年十一個人.
現在有多少人信主呢?.
我們現在不是十一個人.
我們要聽主的吩咐.
你們要去.
使萬民做我的門徒.
做耶穌的門徒.
奉父子之名.
跟他們廝駛.
神就跟我們同在.
成為世代的終結.
所以我們要願意聽耶穌的話.
祂就天天與我們同在.
我們一起祈禱吧.
\newpage



\section{出埃及記 13:17-22}
\label{sec:KbMX2DPfTSA}
\textbf{2022.10.30 《曠野中的更新》 出埃及記 13\_17-22 節 (和修版) 董家驊牧師}
\newline
\newline
連結: \href{https://youtube.com/watch?v=KbMX2DPfTSA}{\texttt{https://youtube.com/watch?v=KbMX2DPfTSA}} ~~~~ 語音日期: 2022-10-30
\newline
\newline
\hyperref[sec:ngXOvQd6nwI]{\small{< < < PREV SERMON < < <}}
~
\hyperref[sec:index_chronic]{\small{[返順時目]}}
~
\hyperref[sec:index_scriptual]{\small{[返順卷目]}}
~
\hyperref[sec:WZ8nJrtvmJs]{\small{> > > NEXT SERMON > > >}}
\newline
\newline
出埃及記 13:17-22
\newline
\begin{longtable}{cl}
\hline
\hline
章節 & 經文 (和合本修訂版)\\
\hline
13:17 & \begin{tabularx}{0.7\textwidth}{X} 法老放百姓走的時候，非利士人之地的路雖近，神卻不領他們從那裡走，因為神說：「恐怕百姓遇見戰爭就後悔，轉回埃及去。」 \end{tabularx} \\ \\ \relax
13:18 & \begin{tabularx}{0.7\textwidth}{X} 神領百姓繞道而行，走曠野的路到紅海。以色列人出埃及地，都帶著兵器上去。 \end{tabularx} \\ \\ \relax
13:19 & \begin{tabularx}{0.7\textwidth}{X} 摩西把約瑟的骸骨一起帶走；因為約瑟曾叫以色列人鄭重地起誓，對他們說：「神必定眷顧你們，你們要把我的骸骨從這裡一起帶上去。」 \end{tabularx} \\ \\ \relax
13:20 & \begin{tabularx}{0.7\textwidth}{X} 他們從疏割起程，在曠野邊上的以倘安營。 \end{tabularx} \\ \\ \relax
13:21 & \begin{tabularx}{0.7\textwidth}{X} 耶和華走在他們前面，日間用雲柱引領他們的路，夜間用火柱照亮他們，使他們日夜都可以行走。 \end{tabularx} \\ \\ \relax
13:22 & \begin{tabularx}{0.7\textwidth}{X} 日間的雲柱，夜間的火柱，總不離開百姓的面前。 \end{tabularx} \\ \\
[1ex]
\hline
\hline
\end{longtable}
$^{1}$弟兄姊妹早安.
非常的抱歉我不會講廣東話.
很抱歉我不會講廣東話.
今天是錦繡堂四十一的周年慶.
今天是錦繡堂四十一年的堂慶.
我在這裡代表世界華府中心跟弟兄姊妹問安.
我在這裡代表世界華福向弟兄姊妹問安.
剛剛我們一個一個來到台前領受聖餐的時候.
剛才我們一個一個來到台前領受聖餐的時候.
小時候每次想到教會要領聖餐.
小時候每次想到要去教會領聖餐.
我的第一個反應是今天的聚會會拖到時間.
我的反應就是這個聚會會拖到很長的時間.
但長越大的時候我越體會到聖餐的重要.
但是長大之後我越來越體會聖餐的重要.
在當中我們學習等待.
也等待彼此.
為了要活出我們在基督裡合一的真實.
為了要活出我們在基督裡面的真實.
四十六年前上帝感動一群華人教會的領袖.
四十六年前神去感動世界華人的領袖.
決定要連結全球各地的華人教會回應宣教使命.
決定要連結所有的教會去做這個教會的使命.
這也是一個不容易的事情.
這也是不容易的事情.
但是當年上帝把這個種子放在一群領袖的心中.
但是當年神將這個種子放在一群領袖的心中.
如今我們也看到它發芽.
而這個運動的開始地其實就是在香港.
這個運動的開始地其實就是在香港.
我相信香港的教會仍舊在未來會扮演全球華人教會當中重要的角色.
我相信香港仍然在以後的日子裡邊都會扮演重要角色.
今天是四十一週年慶.
今天是四十一週年的唐慶.
我常在想講四十週年很容易.
我在說要講四十週年的事情很容易.
但是第四十一年要講什麼.
但是四十一年要講什麼呢.
四十年不是就要進入到應許之地去了嗎.
四十年就要進入應許之地.

$^{41}$那四十一年不就在應許之地裡面很快樂了嗎.
那四十一年就應該在應許之地很快樂的時間.
我想到摩西四十一歲的時候.
摩西的人生也進入一個新的階段.
他的人生都進入一個新的階段.
什麼階段呢.
曠野的階段.
過去四十年摩西在埃及所學的一切知識都是有用的.
過去四十年摩西在埃及裡面所學的東西都是有用的.
但是當上帝真的要使用摩西的時候.
但是上帝要用他的時候呢.
在第四十一歲的時候.
在四十一歲的時候.
上帝把摩西帶入曠野.
我想正因為摩西曾在曠野待了四十年.
我想摩西在曠野待了四十年.
以至於他將來上帝使用他把以色列人帶到曠野的時候.
以至於神要使用他將這些以色列人帶到曠野的時候.
他知道怎麼帶領以色列人在曠野生活.
他知道怎麼帶領以色列人在曠野裡面生活.
這讓我們看到一個非常重要的原則.
就讓我們看到一個很重要的原則.
上帝首先要在我們身上工作.
上帝要首先在我們身上工作.
才透過我們工作.
然後透過我們去工作.
我們很多時候很想為上帝做大事.
我們很多時候想為主做大事.
我們想為上帝做很多很多的事工.
為上帝做很多很多的事工.
但若上帝首先不在我們身上做工.
如果上帝首先不在我們身上做工.
若我們的生命沒有被改變.
如果我們的生命沒有被改變.
我們所做的可能都是虛工.
但是當上帝在我們身上工作的時候.
我們需要的是耐心.
摩西一呆就是四十年在曠野.
摩西在那裡一去就是四十年了.
但是這四十年絕對不是浪費.

$^{81}$這四十年絕對不是浪費的.
而是在預備摩西這個人.
是在預備摩西這個人.
預備他帶領上帝的百姓離開埃及進入應許之地.
剛剛我們讀了這段經文.
講到說當法老釋放以色列人之後.
講到法老釋放以色列人之後.
以色列人要從埃及進到上帝應許他們的迦南地.
以色列人要從埃及進到迦南地.
如果你去看地圖.
你會發現從埃及到迦南地.
由埃及到迦南地.
曠野是必經之地.
換句話說曠野是在我們生命旅程當中不可少的.
換句話說曠野在我們生命旅程當中不可少的.
我們很喜歡講進入應許之地.
當我還是二十幾歲的傳道人的時候.
很多長輩看著我說.
太好了約書亞世代.
太好了是約書亞世代.
是約書亞的世代.
那個世代.
所以那個時候就會覺得.
對我們這一代要起來進入應許之地.
那時候就說對我們要進入應許之地.
如今已經四十一歲了.
有時候我在想是不是我們都不是約書亞世代.
我們其實都是曠野的世代.
其實我們都是曠野的世代.
為什麼呢.
因為聖經很清楚的講.
新天新地是從天而降的.
新天新地是由天而降.
我們都要進入那一個上帝所賜的安息.
我們都要進入神所賜的安息.
而當我們從曠野的角度重新來理解我們生命的時候.
一方面我們就明白為什麼在世上有苦難.
一方面我們就了解為什麼我們在世上會有苦難.
我們就有能力來回應世上的苦難.
但另外一方面我們又不失去盼望.

$^{121}$另一方面我們又不失去盼望.
我們不喪膽.
因為知道有一個應許之地是我們要進入的.
因為有一個應許之地是我們要進入的.
曠野是什麼地方.
是從埃及進到應許之地中間的地方.
以色列人離開了埃及.
我們可能覺得他們很開心.
他們的確應該很開心.
因為終於不用當奴隸了.
但是你知道當了一輩子奴隸的人.
你知道當了一輩子奴隸的人.
突然之間自由了.
那是多恐怖的一件事.
是恐怖的事.
過去他們所熟悉的生活不復存在.
過去他們所熟悉的生活不再存在.
以前當奴隸的時候.
有人告訴他們什麼時候要起床.
告訴他們要工作.
做完工作給他們吃的.
做完工作就給他們吃的.
明天再來一次.
明天再來.
一週工作七天.
一個禮拜工作七天.
跟我們有時候生活好像.
我們的生活都好像.
但突然之間他們從埃及離開了.
突然間他們從埃及離開了.
舊的生活方式已經過去.
可是又還沒有進到迦南地.
但是還沒有進入迦南地.
這是一個極度混亂的時刻.
是一個極度混亂的時候.
我想今天不論是面對全球的局勢.
我想現在無論是面對全世界全球的局勢.
還是教會堂會的發展.
或者是教會堂會的發展.
還是全球宣教策略.

$^{161}$或者是全球的宣教策略.
或者是神學教育的未來.
我們都在面對這樣典範轉移中間的曠野階段.
我們都是面對一個典範轉移的一個.
曠野.
曠野之中.
過去熟悉的秩序漸漸的瓦解.
過去熟悉的次序是瓦解了.
可是又還在.
但是又找到.
新的秩序又還沒來到.
新的秩序又沒來.
在這兩個中間我們到底該怎麼辦.
在這兩個中間.
兩個秩序中間.
秩序當中.
秩序當中我們怎麼辦.
但如果我們回顧整個教會歷史.
如果我們回顧整個教會歷史.
我們會發現往往這樣的時刻.
往往這個時候.
正是教會經歷更新的時刻.
正是教會重新認識自己是誰的時刻.
正是教會重新認識自己的時刻.
上帝把摩西帶入曠野.
預備了他.
要祝福將來的以色列人.
今天上帝把教會.
好像帶到一個曠野的處境當中.
今天上帝好像將教會.
帶到一個曠野的處境當中.
再一次的提醒我們.
願上帝在我們身上工作.
因為我們將來要服侍這個世代.
也許現在就開始了.
但是當上帝先更新我們的時候.
但是當上帝更新我們的時候.
我們才有能力來服侍這個世代.
我們需要面對曠野的階段.
因為曠野是從埃及到應許之地的必經之路.

$^{201}$第二個我們看到.
在這段經文裡面.
上帝沒有帶以色列人走最近的路.
這個很不符合我們華人的習慣.
好像不符合我們華人的習慣.
特別在大城市.
香港 新加坡.
也包括台北.
我們都喜歡效率.
我們都要效力.
今年暑假我到馬來西亞.
在那邊服侍.
跟當地的很多教會領袖聊天.
我們可以聊一兩個小時.
聊得非常開心.
我們可以聊一兩個小時.
聊得很開心.
這可能是當地的文化.
這些可能是當地的文化.
馬來西亞結束之後.
我到了新加坡.
馬來西亞結束之後.
我到了新加坡.
馬來西亞結束之後.
我到了新加坡.
新加坡有新加坡的文化.
我一去就感受到了.
因為我一天的行程就有六個約.
因為我一天的行程裡面.
就有六個約會.
早餐一個約.
早上一個約.
中餐一個約.
中餐.
午餐的時候一個約.
下午再一個約.
晚餐又一個約.
有時候晚餐結束還有一個約.
可能晚餐.
晚餐結束之後還有一個約.

$^{241}$非常講求效率.
而且見面也不廢話.
坐下來閒話家常兩三分鐘之後.
就開始講主題.
通常對方就會問.
David,今天你有什麼計畫.
就問你編排怎麼樣.
我想這是不同地方的文化.
如果我們把這個文化套入到.
今天這段經文當中.
如果我們把這個文化套入到.
今天的經文當中.
就會覺得上帝很沒效率.
因為這裡講到.
非利士人之地的路雖近.
因為這裡講到.
非利士人的路雖近.
但上帝卻不領他們從那裡走.
上帝帶領以色列人繞路.
上帝帶以色列人繞著那走.
走一條比較遠的路.
為什麼呢.
因為上帝知道以色列人還沒預備好征戰.
因為上帝知道以色列人還沒準備好征戰.
在一個講求效率的年代.
在一個講求效率的時代.
我們常會對上帝的帶領產生懷疑.
上帝你為什麼不這樣做.
我們會想上帝為什麼你不是這樣做呢.
你為什麼不那樣做.
你為什麼不是這樣做呢.
為什麼你不帶我走最近的路.
卻要帶我走遠路.
但是因為上帝真正關心的.
不是做工的效率.
而是做工的人.
是做工這個人.
上帝真正關心的是你與我.
上帝真正關心的就是你與我.
我們如何能夠在他裡面成長.

$^{281}$在我們今天這個時代.
其實我們活在極度的焦慮當中.
其實我們活在極度焦慮當中.
我們的過去的身分認同逐漸失去.
我們過去的身分認同逐漸失去.
一百年前.
你生在一個讀書人的家裡.
你這輩子就注定要做讀書人.
一輩子就注定要做讀書人.
哪怕你不會讀書.
就算你不會讀書.
你也要走這條路.
如果你生在一個種田的家裡.
你就是一個要種田的人.
你一輩子就要種田.
你走在大街小巷.
你在大街小巷.
在你的村莊裡面.
在村莊裡面.
大家一看就知道你是誰.
但在今天這個高度都市化的時代.
走在路上沒有人認識你是誰.
有一次我跟一個很有名的牧師見面.
我說牧師你在外面吃飯的時候.
你說牧師你在外面吃飯的時候.
你不會害怕嗎.
我跟牧師說你在外面吃飯.
你不會害怕嗎.
怕被別人認出你嗎.
怕別人認出你嗎.
他說家華.
基督徒的比例很少.
在路上沒有人會認得我.
就算有也很少.
就是表示很少.
我們活在一個.
其實走在路上沒有人認得我們的時代.
我們活在一個.
走在路上沒有人認識我們的時代.
我們很焦慮自己的身分到底是什麼.

$^{321}$我們的價值到底是什麼.
我的尊嚴到底是什麼.
為什麼今天我們有那麼多人.
包括我在內.
我們有的時候會進入工作狂的型態.
為什麼我們那麼多人.
包括我都是進入一個工作狂的狀態.
因為在今天的時代.
工作不再只是工作.
因為在今天的時代.
工作不再只是工作.
工作變成我們尊嚴的來源.
至少我們是這樣認為的.
至少我是這樣形容的.
工作變成我們被接納的來源.
所以我們很害怕.
一旦失去了我們的工作.
所以我們很害怕.
一旦失去了我們的工作.
我們很害怕一旦我們失去工作的時候.
一旦失去了我們的頭銜.
一旦失去了我們的頭銜的時候.
我就沒有價值了.
就沒有人會愛我了.
所以我們緊緊抓住我們的工作.
緊緊地抓住我們的頭銜.
但是耶穌基督有話要說.
祂真正關心我們的不是我們的工作.
而是我們這個人.
而很有趣的是.
唯有當我們不把身份意義放在我們的工作上.
我們才有可能好好的工作.
只有可能真正好好的工作.
喜樂的工作.
用工作來關顧我們的靈社.
用工作來關愛我們的靈社.
不再是為了我做什麼讓祂能夠回報我.
也不是我做什麼以至於祂更尊敬我.
因為我們知道我們最寶貴的身份是天父的兒女.
因為我們最寶貴的身份是天父的兒女.

$^{361}$有什麼身份比這個更重要呢.
有什麼身份比這個更有價值呢.
當我還是個二十幾歲的年輕傳道人的時候.
其實我很享受講道的服侍.
我覺得每一次的講道都是一個根神的創作.
其實我每一次講道都好像同神的創作.
我的主任牧師當時很有智慧.
他看得出來我很享受講道的事奉.
他就把我叫到他的辦公室.
他就叫我去他的辦公室.
他說嘉華.
你可能有講道的恩賜.
你有沒有想過一件事.
如果有一天你不再能夠講道了.
如果有一天你不再能夠講道.
你的人生還有意義嗎.
你的人生還有價值嗎.
那對我來講非常的震撼.
他這樣講的時候令我非常震撼.
因為這是非常有可能的.
有一天若我老了.
講不出話了.
頭腦不清楚了.
或者就是發生了一個意外.
或者發生了一個意外.
讓我說不出話來.
讓我沒得講話.
當我不再能夠講道的時候.
我這個人就沒有價值了嗎.
我就沒有尊嚴了嗎.
上帝呼召他的教會.
成為一個獨特的群體.
這個群體首先我們領受上帝的恩典.
這個群體是首先領受上帝的恩典.
我們是神的兒女.
以至於我們去愛彼此.
不是因為對方的工作.
不是因為對方的頭銜.
而單單的因為對方也是天父的愛子愛女.
因為對方都是神所愛的愛子和愛女.

$^{401}$我們都是天父家裡的人.
這是教會對這世界能夠給的最美的見證之一.
上帝在預備他的百姓.
在預備我們.
或許在這最困難的時候.
就是上帝在更新我們的時候.
事實上你仔細去想想聖經裡面.
其實我們只是想一下聖經裡面.
像以大衛為例子.
以大衛為例子.
大衛的一生許多的時間是在曠野漂泊.
在曠野當中.
不是大衛戰勝了曠野.
而是上帝戰勝了大衛得到他的心.
而是上帝戰勝大衛得到他的心.
曠野或許不是我們要去克服改變的.
而是在當中我們被上帝改變.
我們被上帝得著.
曠野是上帝培育我們的地方.
第三個.
曠野是我們學習順服的地方.
我們看到很有趣.
上帝帶百姓繞道而行.
我們看到上帝帶百姓繞道而行.
繞道去哪裡呢.
去紅海邊.
如果你再去把這段聖經讀下去的話.
如果你再把這段聖經讀下去的話.
你就會發現以色列人之所以會面對紅海.
以色列人之所以會面對紅海.
後面有埃及的追兵.
陷入到一個絕境當中.
是上帝的帶領.
是上帝刻意的帶領.
上帝精心的策劃.
是上帝精心的策劃.
而面對紅海.
這個看似難以跨越的困難.
看到這個是很難跨越的困難.
上帝讓以色列人學習的一件事情.

$^{441}$唯一的出路就是順服上帝.
順服不是一個抽象的原則.
而是有清楚對象的.
順服的對象是上帝.
當以色列人面對紅海抱怨的時候.
上帝親自讓以色列人看到.
沒有什麼困難是上帝不能夠戰勝的.
沒有什麼困難是上帝不能戰勝的.
唯一的困難是我們的想像力有限.
每一個教會我們都有一個裝飾品.
就是十字架.
剛剛在唱詩歌敬拜的時候.
我看著後面的玻璃.
那就是一個十字架.
就是一個十字架.
十字架就是提醒我們.
十字架是什麼記號.
它本來是失敗的記號.
它本來是羞辱的記號.
本來是羞辱的記號.
因為被釘十字架是全身發光釘在上面.
因為被釘十字架是全身曝光的.
被釘在上面.
人們路過的人看著你在上面赤裸的掙扎.
每個人經過就看到你在赤裸的掙扎.
一個失敗羞辱的記號.
也是人類最大錯誤的記號.
事實上想一想我們還可以犯什麼錯比這個更大的.
事實上我們還可以犯什麼錯比這個更大的.
把那來拯救我們的上帝自己釘死在十字架上.
但是三天之後耶穌基督從死裡復活.
但是三天之後耶穌從死裡復活.
如今我們看到十字架.
現在我們看到十字架.
我們看到的不是失敗.
而是上帝的得勝.
我們看到的不是羞辱.
而是上帝的榮耀.
我們看到的不是人類的失敗.
而是上帝的勝利.

$^{481}$這再一次的提醒我們.
再一次提醒我們.
沒有任何的錯誤是上帝無法翻轉的.
沒有任何的失敗是上帝不能夠轉化的.
你我生命當中沒有任何的羞辱.
是上帝不能夠以祂的榮耀來遮蓋的.
你的羞辱不是神的榮耀可以遮蓋的.
這是十字架的大能.
這是我們所宣講的福音.
這也是我們所信服的耶穌基督.
一個多禮拜前.
我在泰國遇到韓國一個很大教會的牧師.
我在泰國遇到一個韓國很大的牧師.
這位牧師跟我想像中的韓國教會大牧師感覺很不一樣.
這位牧師跟我想像裡面的.
在韓國很大的牧師是不同的.
我心中想像的韓國Mega Church的牧師都是很有威嚴的.
我想像中韓國的牧師是很有威嚴的.
但這位牧師很溫柔.
他帶著他的女兒跟我聊天.
因為他講韓文.
他女兒翻成英文.
因為他講韓文.
他女兒翻英文.
我講英文.
他的女兒翻成韓文.
我講英文.
他女兒翻韓文給爸爸聽.
聊到後來彼此都很有安全感的時候.
當大家聊得很有安全感的時候.
我就問他.
牧師,韓國教會過去四十年發展得非常快.
我就問他牧師.
牧師過去四十年韓國教會發展得非常快.
各地的華人教會都向你們學習.
你們的教會真的是大.
在台灣超過一千人的教會叫大教會.
在韓國叫做中小型教會.
你們所拆開的宣教士.
是全世界第二多的.

$^{521}$僅次於美國.
如果按比例來講你們超過美國.
但是過去十年整個韓國教會衰落得那麼快.
但是過去十年韓國教會衰弱得那麼厲害.
整個社會對教會的印象從正面到非常的負面.
整個韓國裡面對教會的印象是很負面的.
我就問這位牧師說.
我就問這位牧師.
牧師如果能夠時光倒流.
你回到二十年前.
如果時光倒流回到二十年前.
你會提醒二十年前的韓國教會什麼.
你會提醒二十年前的韓國教會是什麼.
我在問這個問題的時候完全沒有想到會得罪他.
他也沒有被得罪.
可是每當我跟別人講這件事情的時候.
人們都捏一把冷汗.
就很緊張,很為我緊張.
這位牧師的回答讓我非常的感動.
這位牧師讓我非常的感動.
他說他會提醒二十年前的韓國教會.
他說他會提醒二十年前的教會.
當我們把任何的事情看得比順服主更重要的時候.
我們就把那件事情給當作偶像來敬拜了.
我們就等於把那件事情當作偶像來敬拜.
當我們把發展堂會看得比順服主更重要.
當我們把發展堂會看得比順服主更重要的時候.
當我們把宣教事公看得比順服主更重要.
當我們把宣教事公看得比順服主更重要的時候.
這一切都變成偶像.
當這一切變成偶像的時候.
當一切變成偶像的時候.
就生出了競爭.
比較.
嫉妒.
妒忌.
衝突.
他說這是我會提醒二十年前的韓國教會.
這個是我會提醒二十年前的教會.
他說嘉華.

$^{561}$你要提醒華人教會.
沒有什麼是比順服基督更重要的.
以色列人之所以能夠過紅海.
不是依靠自己的武力.
而單純的是順服上帝的工作.
而今天我們特別需要回到這個最根本的原則.
今天我們需要回到這個最根本的原則.
每天我們有時間來聆聽上帝的聲音嗎.
在這麼繁忙的生活節奏當中.
我們有時間來閱讀上帝的話語嗎.
而當我們聽了和讀了.
當我們聽了和讀了.
我們有安靜下來讓神的話進到我們內心深處嗎.
我們有安靜下來讓神的話進入我們內心深處嗎.
而當我們有了領受之後我們有順服回應嗎.
當我們領受之後我們就順服去回應.
在曠野當中.
我們唯一的出路.
就是相信而順服上帝.
最後我想要講一點.
在曠野當中.
我們需要重新學習一件事.
就是依靠上帝的行動.
在第21節這裡.
在21節裡面.
在曠野當中上帝怎麼引導他們呢.
日間用雲柱引導他們的路.
夜間用火柱照亮他們.
日間的雲柱,夜間的火柱.
總不離開百姓的面.
日間的雲柱,夜間的火柱.
總不離開百姓.
教會往前行的路.
我們需要有計劃.
我們需要有討論.
但從來不是依靠我們的計劃跟討論.
我們去看初代教會的宣教.
並不是彼得,雅各,約翰聚在一起討論的宣教策略.
然後去推廣到安提亞教會.
推廣到伊佛索教會.

$^{601}$我絕對不是在否定計劃的價值.
我們需要做計劃.
這是對士工的尊重.
不然我們是隨隨便便.
不然我們就是便便.
但是我們在計劃之餘.
我們要不斷問自己.
我們依靠的到底是什麼呢.
我們依靠的是有一套完美的計劃.
我們依靠的是完美的計劃.
還是那位隨時能夠打亂計劃工作的上帝.
而是那個隨時可以打亂計劃的上帝.
也能夠使用我們的計劃成就他工作的上帝.
在今天現代生活當中.
其實我們的生活是被切割成很多不同的領域的.
包括基督徒在內.
我們都被訓練相信一件事.
就是其實不太需要上帝.
其實不太需要上帝.
我們也可以搞定自己的事情.
我們也可以把世界變得更美好.
而不知不覺之間.
當我們把上帝真實的工作放在一邊的時候.
我們以為可以靠自己來改變這個世界的時候.
我們就開始用這個世界的方式來衡量我們的事工.
當上帝的同在不再是我們行動的基礎的時候.
剩下的就是哪些做法可以帶給我們更多的人力.
物力跟財力.
在曠野當中.
我們唯一的依靠是上帝的雲助火助.
而這絕對不是跟計劃衝突.
這些絕對跟計劃是沒有衝突的.
因為你想像一下.
你想像一下.
兩百萬的人.
什麼時候上帝引導他們前進.
他們就需要趕快把帳篷收起來.
他們就要快點把帳篷收起來.
然後列隊要前進.
然後列隊前進.

$^{641}$什麼時候停下.
他們立刻就要把所有的行李都拿出來.
他們要把所有的行李都拿出來.
這需要多少的計劃.
這些需要多少的計劃.
這需要多少的紀律.
所以計劃跟紀律絕對不是跟依靠上帝是一個對比的.
但是我們再一次被提醒什麼是最重要的.
但我們要認定哪樣是最重要的.
我們可以做計劃.
我們可以訓練有紀律.
但如果失去了雲柱火柱的引導.
我們會一直在曠野當中繞路.
我們會一直在曠野繞路.
但是當我們順服上帝雲柱火柱的引導.
有一天我們都要在上帝賜給我們的應許之地當中有份.
最後我要分享一句話.
是我今年八月在馬來西亞的時候遇到一位主教.
他跟我分享的.
在今年八月我遇到一個主教跟我說的.
他是馬來西亞的Bishop.
馬來西亞衛理工會的Bishop 華勇.
他是馬來西亞教會的一個主教.
他是在馬來西亞當中非常受敬重的一位領袖.
不論是長輩還是年輕一代都非常敬重他.
無論是長輩或是年輕人都很敬重他.
我去他家旁邊的咖啡店看他.
我們是網友一兩年了.
第一次面對面見面.
我們聊到中間的時候.
我們聊到一半的時候.
他突然很認真的看著我.
然後對我說.
David.
你我都有博士學位.
因此我們更應該大聲疾呼.
教會不是建立在學位上面.
不是建立在建築物上面.
不是建立在產業上面.
不是建築在產業上面.

$^{681}$而是上帝的話語.
聖潔的生活.
跟真實的門徒訓練.
願上帝帶領我們.
在41週年的時刻.
我們在曠野當中.
被上帝更新.
以至於在下一個時代當中.
上帝也透過我們.
成為許許多多人的祝福.
我們一起禱告.
天父上帝謝謝您.
我們在這裡歡慶過去40年你的作為.
而面對未來.
我們面對未來.
求你在我們生命當中工作.
復興你的教會.
更新你的教會.
使我們成為靈社的祝福.
使我們成為萬族萬民的祝福.
奉耶穌基督的名禱告.
祈禱.
阿們.
\newpage



\section{尼希米記 10:28-39}
\label{sec:WZ8nJrtvmJs}
\textbf{2022.11.06 《重新立約》 尼希米記 10\_28-39 (和修版) 老耀雄牧師}
\newline
\newline
連結: \href{https://youtube.com/watch?v=WZ8nJrtvmJs}{\texttt{https://youtube.com/watch?v=WZ8nJrtvmJs}} ~~~~ 語音日期: 2022-11-13
\newline
\newline
\hyperref[sec:KbMX2DPfTSA]{\small{< < < PREV SERMON < < <}}
~
\hyperref[sec:index_chronic]{\small{[返順時目]}}
~
\hyperref[sec:index_scriptual]{\small{[返順卷目]}}
~
\hyperref[sec:VN5NAedAhq4]{\small{> > > NEXT SERMON > > >}}
\newline
\newline
尼希米記 10:28-39
\newline
\begin{longtable}{cl}
\hline
\hline
章節 & 經文 (和合本修訂版)\\
\hline
10:28 & \begin{tabularx}{0.7\textwidth}{X} 其餘的百姓、祭司、利未人、門口的守衛、歌唱的、殿役，所有與鄰邦民族分別出來、歸服神律法的，以及他們的妻子、兒女，凡有知識、能明白的， \end{tabularx} \\ \\ \relax
10:29 & \begin{tabularx}{0.7\textwidth}{X} 都隨從他們貴族的弟兄發咒起誓，要遵行神藉他僕人摩西所賜的律法，謹守遵行耶和華－我們主的一切誡命、典章、律例。 \end{tabularx} \\ \\ \relax
10:30 & \begin{tabularx}{0.7\textwidth}{X} 我們不把我們的女兒嫁給這地的居民，也不為我們的兒子娶他們的女兒。 \end{tabularx} \\ \\ \relax
10:31 & \begin{tabularx}{0.7\textwidth}{X} 這地的民族若在安息日，或甚麼聖日，帶了貨物或糧食來賣，我們必不買。每逢第七年必不耕種，凡欠我們債的必不追討。 \end{tabularx} \\ \\ \relax
10:32 & \begin{tabularx}{0.7\textwidth}{X} 我們又為自己定例，每年各人捐獻三分之一舍客勒，作為我們神殿之用： \end{tabularx} \\ \\ \relax
10:33 & \begin{tabularx}{0.7\textwidth}{X} 為供餅、常獻的素祭和燔祭，安息日、初一、節期所獻的祭和聖物，以色列的贖罪祭，以及我們神殿裡一切工作之用。 \end{tabularx} \\ \\ \relax
10:34 & \begin{tabularx}{0.7\textwidth}{X} 我們的祭司、利未人和百姓都抽籤，每年按父家定期將奉獻的木柴帶到我們神的殿裡，照著律法上所寫的，燒在耶和華－我們神的壇上。 \end{tabularx} \\ \\ \relax
10:35 & \begin{tabularx}{0.7\textwidth}{X} 每年我們又將地上初熟的土產和各樣樹上初熟的果子，都奉到耶和華的殿裡。 \end{tabularx} \\ \\ \relax
10:36 & \begin{tabularx}{0.7\textwidth}{X} 我們又照律法上所寫的，將我們頭胎的兒子和首生的牛羊都奉到我們神的殿，交給在神殿裡供職的祭司； \end{tabularx} \\ \\ \relax
10:37 & \begin{tabularx}{0.7\textwidth}{X} 並將初熟麥子所磨的麵和舉祭、各樣樹上的果子、新酒與新油奉給祭司，收在我們神殿的庫房裡，又把我們土地所產的十分之一奉給利未人，因利未人在我們一切城鎮的土產中當取十分之一。 \end{tabularx} \\ \\ \relax
10:38 & \begin{tabularx}{0.7\textwidth}{X} 利未人取十分之一的時候，亞倫的子孫中當有一個祭司與利未人同在。利未人也當從十分之一中取十分之一，奉到我們神的殿，收在庫房的倉裡。 \end{tabularx} \\ \\ \relax
10:39 & \begin{tabularx}{0.7\textwidth}{X} 因以色列人和利未人要把禮物，就是五穀、新酒和新油，帶到收存聖所器皿的倉裡，供職的祭司、門口的守衛、歌唱的都在那裡。我們絕不會不顧我們神的殿。 \end{tabularx} \\ \\
[1ex]
\hline
\hline
\end{longtable}
$^{1}$弟兄姊妹平安.
紅恩家兼赴武唐浸浴堂的差派.
在洪水橋開荒植唐的使命.
今年已經達到十一年.
實在是一個值得慶賀的日子.
甚願我們今天的唐慶崇拜.
能夠獻情給上主.
因為祂是至高的神.
今年唐慶籌委會開會時.
也討論過會否像去年那樣.
有午餐到會或辦晚宴.
可以吃一餐開心一下.
不過現在的感染個案仍然持續在五千宗.
雖然比過往一萬宗時回落了.
而且酒樓的限聚人數也放寬了.
不過教會也不想鬆懈.
因此集體進食活動盡量避免.
或許可以配合一些過往沒有的元素.
如果今天有回到教會.
你會發現有幾個特色.
我們有幾個唐慶佈置是往年沒有的.
教會門外掛了很多彩帶.
有慶祝的氣氛.
禮堂有個打卡位.
很搶眼的.
可以讓我們拍照留念.
後面有一個感恩的十字架.
我們可以寫一些對教會的祝福和期望.
然後呈獻給主.
這裡有一個十一字的剪版.
貼上教會過往的相片展示出來.
作為這十一年的集體回憶片段.
我們也將明年的年提寫在講壇後面.
讓大家對明年的年提有更深的印象.
他們想出了很多新的構思.
讓今年的唐慶與別不同.
今天的唐慶崇拜也加入了一些新的元素.
我賣個關子讓大家自行發掘一下.
去年十週年的唐慶講道.
我選用了以色列記第六章.

$^{41}$聖殿奉獻禮的其中一段經文.
就像潤唱弟兄所說.
其實去年的年頭.
大家都經歷過教會沒有實體的崇拜.
只能夠網上崇拜的情況.
因此去年就很想會有思考一下.
能夠回教會在聖殿裡面敬拜神.
不是必然的.
是一個恩典.
今年我想再延續聖殿奉獻禮之後的安排.
就是在以色列從早上到中午讀經會之後.
尼希米帶同眾百姓.
以及事奉神的祭司利未人.
在神面前重新立約的一段經文.
因此我今天選用了尼希米記的十章28至39節.
將這一段昔日被老歸回神的子民.
在神面前立約的經文.
去引發我們今天洪恩家的眾知體.
在這個十一週年唐慶崇拜日子裡面.
我們會不會一起在神面前.
也同樣有一個與神立約的時間呢?.
我們一起去祈禱.
愛我們的主.
你帶領著洪恩家走過十一個年頭.
我們去數算的恩典真是多疑又多.
數之不盡.
不過我們仍然需要你的寬恕和憐憫.
因為我們不聽你話的地方.
比聽你話的地方更多.
主耶穌你既是我們的主.
也是我們的神.
就求你繼續管教我們.
求聖靈繼續引導,督責,指引.
提醒我們每一個.
讓我們能夠專心的仰望你.
能夠好像明年的年提一樣.
學務真理.
遵行主導.
讓我們每一個屬你的.
都能夠心口如一.

$^{81}$活出信仰.
讓我們所領受的恩典.
與我們所實踐的行為能夠相稱.
讓我們身邊的親人,朋友.
能夠在我們的生命裡見到主耶穌.
成為一個美好的見證.
讓你的名被高舉.
讓你的榮耀被頌讚.
我們祈禱奉教主耶穌基督.
得勝明治祈求.
阿們.
我們進入經文之前.
都再重溫一下當時的歷史背景.
我們知道.
北角以色列以及南角猶大.
分別被亞述和巴比倫滅亡之後.
猶太人正式成為一個.
割破家亡.
而且被擄到異地生活的民族.
不過神沒有放棄他們.
神透過耶利米先知.
預言他們七十年年滿之後.
被擄的人可以回到耶路撒冷.
重建家園.
而以色列體驗過三次回歸的經歷.
第一次是由索羅巴伯帶領.
不過在重建聖殿的過程裡.
不斷受到敵對勢力的攻擊.
最終聖殿只建立了根基.
停止了重建的工程.
第二次是由以斯拉帶領.
藉著他的鼓勵.
聖殿真的能夠重建完畢.
回歸的百姓.
終於可以在聖殿裡敬拜神.
第三次是由尼希米帶領.
他主要是收拾耶路撒冷的城牆.
讓耶路撒冷城.
不會輕易被敵人進入和攻擊.
當城牆收拾完成.

$^{121}$尼希米帶領了.
回歸耶路撒冷的猶太人.
在耶路撒冷與神重新立約.
體現了猶太人在神面前合一的心.
因此立約不單單是涉及.
個人對神的承諾.
也是屬神指紋.
一個合一的體現.
今天的經文就是立約的內容.
我們一起進入經文.
這個立約的禮儀涉及甚麼人呢?.
28節.
他提到有百姓,祭司,利未人.
門口的守衛,歌唱的,電役的.
所有與倫邦民族分別出來.
「歸服神律法」的.
以及他們的妻子兒女.
雖然經文列出了很多種類的人.
但卻不是所有住在耶路撒冷的居民都有份.
因為是有條件的.
就是「歸服神律法」的.
這是表示甚麼呢?.
這是表示只有願意接受摩西律法的價值觀.
以及願意在摩西律法底下生活的人.
他們才有資格與神立約.
同樣.
如果洪恩堂的肢體很想與神立約.
是否所有人都可以呢?.
是否來者不拒呢?.
或者只是相信神就可以了.
還是同樣有條件呢?.
的確是有條件的.
就是相信主耶穌.
以及願意「歸服神律法」之下的人.
即是接受主耶穌為他生命的主.
以及救贖的主.
並且願意按照神的話語.
去活出信仰的肢體.
如果我們只是單單相信.
而不願意按照神的心意去活出信仰的話.

$^{161}$那又何必勉強自己去立約呢?.
立約是表明我們願意付出代價.
去踐行這個信仰.
如果我們不願意付出.
那我們就不要勉強自己去立約了.
昔日的百姓用什麼形式去立約呢?.
二十九節.
「法就起誓」.
「法就」是一個情緒的用詞.
是表達一種義無反顧的心態.
因為毀約者會按照約書的規定去受罰.
是用來堅定雙方的協定.
實際上這個約是以誓言的形式去確立.
表示出一種認真毫不反悔的態度.
對於不認真的誓言.
中國有一句話.
知不知道是說什麼?.
「法誓當食生財」.
是吧?大家都聽過了.
希望我們不是用這種態度去同神立約.
不要覺得「法誓當食生財」.
誓言的內容是什麼呢?.
二十九節.
「要遵行神籍他僕人摩西所賜的律法.
謹守遵行我們主的一切誡命典章律例」.
誡命律例典章是一個大綱.
具體內容是什麼呢?.
包括了三點.
第一點.
是一個分別為性的生活.
我們一起去讀第三十節.
第三十節.
一二三.
「我們不但不敬禮大吉者利君子.
但為我們自己最慚愧的罪孽」.
這個要求不是一個新的要求.
在以色列進入迦南的時候已經提出過.
以色列不可以與迦南七族的人通婚.
目的就是不想被異族的信仰滲透.
不覺以色列人阿哈就是一個最好的例子.

$^{201}$阿哈娶了西頓王的女兒耶洗別.
於是整個不覺就開始敬拜巴力和亞瑟拉.
敬拜別神是耶和華最憎恨的罪.
也是以色列人經常招惹神憤怒的一個原因.
今天如果信徒與神立約.
是不是表示不與異族通婚呢?.
當然不是.
當日不可以與異族通婚是有特殊的歷史背景.
今天信徒當然可以與不同膚色不同種族的人結婚.
但我們需要留意.
我們會不會被配偶的信仰所影響.
當一個信徒與一個非信徒進入婚姻生活.
必然會有很多張力.
因為不同的價值觀.
可以在不同的生活環節裡造成很多的摩擦.
因此很多信徒為了不與配偶產生衝突.
他就選擇了妥協.
這個妥協最後帶來的是什麼呢?.
輕的信仰的展現就大打折扣了.
與神的關係可以很疏離.
更加不要說什麼「毀身事奉」.
重的信仰就會被邊緣化.
可有可無.
甚至可以離開教會.
再作為肢體.
在目測上.
應該大部份人都經歷過婚姻的階段.
在經文裡面對我們有什麼提醒作用呢?.
不是的.
婚姻只是一個選項.
還有很多事情要引誘我們離棄神.
因此這段立約的經文的核心精神就是.
「分別為性」.
分別為性的生活是在基督的價值觀之下生活.
任何與信仰相違背的價值觀.
我們都要敢於去洩露.
因為我們是被神所揀選出來的.
我們有一個很獨特的身份.
是神的兒女.
我們是屬神的.

$^{241}$因此我們所持守的生活.
應該是一種與神的價值觀相融合的生活.
而不是追求大眾文化.
或者世界的價值觀.
當一些價值觀與基督信仰相違背的時候.
我們就要勇敢地說.
NO.
弟兄姊妹.
當我們願意與神立約.
我們願不願意在未來的日子裡面.
過一個分別為性的生活呢?.
第二個就是守安息日的意義.
我們一起讀31節.
1,2,3.
當時由於很多外邦人都為了經濟的利益.
經常在安息日那天.
帶著貨物走到耶路撒冷.
打算去做交易.
有些以色列人就被利益引誘.
就和外邦人進行了買賣交易.
觸犯了安息日的條例.
因此禁止安息日買賣.
是進一步落實安息日的條例.
確保他們不會干犯律法的條例.
休耕的目的也是讓土地能夠安息.
土地可以自由地生產土產.
這些休耕之下而所產生的土產.
可以讓一些貧窮的人自由地去取用.
讓社會低下階層的人都能夠得到溫飽.
豁免債務的條例是基於豁免年的定例.
有些猶太人因為欠債而成為奴隸.
因此律法規定在第七年.
就要釋放家裡猶太人的奴隸.
讓他們得到自由.
安息日的核心精神.
就是要在安息日裡休息.
享受神的供應以及對神的感恩.
在神的創造下.
不單單人要休息.
地也要休息.

$^{281}$而人在安息日最大的釋放是甚麼?.
就是避免所有的債項.
讓人得到自由重新的生活.
就好像我們信了耶穌一樣.
我們不再被罪所克制.
我們生命得釋放.
今天有信徒投訴.
他投訴甚麼?.
他說星期日比平日還辛苦.
因為要回來教會事奉.
很辛苦.
其實星期日是一個.
可以讓信徒身心靈休息的日子.
信徒星期日回教會.
是一種享受.
享受生命被神創造.
可以有生命氣息.
享受神的供應.
一粟三餐仍然可以無所掛慮.
享受神在生命裡彰顯奇妙的作為.
讓我們可以數算恩典.
然後透過敬拜讚美.
去讚頌和回應神救贖的恩典.
在安息的一刻.
我們享受與神相遇.
領受神的話語.
接受神的祝福.
也接受神的差遣.
崇拜整個過程.
讓我們能夠洗滌心靈.
讓我們可以回歸到神的身邊.
思考我們和神的關係.
如何讓我們盡心盡性盡義盡力.
去愛我們的主.
安息日並不是甚麼都不做.
安息日是要放下一切.
專心地去學求神.
唯有與神相遇.
才是我們整個生命最寶貴.
弟兄姊妹.

$^{321}$如果你們回到教會崇拜.
卻享受不到.
在安息日賜給你的福份.
就好像一個.
如入寶山空手回的人一樣.
白白浪費了.
虛耗了你們崇拜的時間.
弟兄姊妹.
我們願不願意與神納若.
在每一個星期日的早上.
我們回到教會的時候.
都能夠放下一切重擔和掛慮.
以一個學武的心去親近主.
以一個敬虔的心去敬拜主.
以一個真誠的心去與主相遇呢?.
第三個.
就是奉獻的意義.
我們一起讀三十二至三十三節.
一,二,三.
我們有為自己敬禮.
每年各人捐獻三分之一的養幸.
作為我們神殿之用.
為了祂的病.
常年的素祭和限祭.
我們昔日.
昨日.
前期所宴.
祈求神聖恆.
以聖經的贖罪聖聖.
以及我們神殿裡的齊工作之用.
百姓為了這個聖殿的需要.
願意作出一個信心的應獻.
每年各人捐獻三分之一的錫克勒.
購買聖殿獻祭時.
所需要的祭物.
這個信心應獻.
不是摩西律法所要求的.
是因為百姓看到.
聖殿有這個需要.
他就自願為此納入.

$^{361}$信心應獻.
是一種常廢以外的奉獻.
也是信徒對神.
豐富供應的一種回應.
我們再讀三十四節.
一二三.
(三十四節).
以聖帝的慣例.
通常都是用抽籤去解決事情.
這是一貫經常的方式.
不過祭司利未人和百姓一起參與.
體現了一種人人平等的律法精神.
摩西五經從來沒有提到.
奉獻祭壇用的柴.
只是規定祭壇的火不可以熄滅.
要維持祭壇的火不斷燃燒.
就自然需要大量的木柴.
因此神的指紋.
要恢復聖殿的敬拜.
首先確保祭壇上有柴.
這樣才可以獻祭.
利希米讓聖殿人員和會眾輪流.
去供應木柴.
以保證祭壇上的火不斷燃燒.
而會眾也能夠以此為榮.
因為人人都有機會.
來到祭壇上獻祭.
敬拜就融入了民眾的生活裡.
今日教會的鮮花奉獻的精神都一樣.
鮮花能夠讓講壇加添了色彩.
在崇拜裡帶出生命的朝氣.
你看看就知道了.
很漂亮的.
會友可以不懂插花.
但是鮮花奉獻體現了會友對崇拜的參與.
會友不是一個過客.
也不是一個旁觀者.
因為聖壇的佈置和擺設.
他是有份參與的.
就好像以色列人一樣.

$^{401}$祭壇上的火.
是用他奉獻的柴來燃燒的一樣.
我們一起讀三十五節三十六節吧.
(三十五節).
在一個農耕的社會裡.
初熟的土產和果子.
以及首生的牛羊.
都是一種很重要的收成.
但是在收成的分配優先次序下.
神是首先的享用者.
因為神是一切的供應.
沒有神的供應.
就沒有所謂的收成.
而投胎的兒子.
也是代表家庭最重要的產業.
所以要獻給神.
因為神是應得的.
納入的精神就是.
將神賜給我們一切的收穫與產業.
都將最好的獻給神.
各位姐妹.
今天.
當我們說將最好的獻給神.
是指甚麼呢?.
是否指在事奉裡更加精益求精.
還是參與更多的事奉.
奉獻更多的錢.
不是.
這一切或許都是好的.
但不是我們生命中最好的.
我們生命中最好的是甚麼呢?.
是一個回歸了.
回轉了一個聖潔的生命.
是一個聖潔的生命.
是我們生命中最好的.
一個被神赦免了的生命.
一個懂得感恩的生命.
一個願意信服神的生命.
一個願意信服神帶領的生命.
所以我想大家思考一下.

$^{441}$甚麼是你生命中最好的?.
你如何將你生命中最好的獻給主?.
我們一起讀37至38節.
1,2,3.
您將在所屬物質所愛的任何對象.
各樣思想的果子.
尋找與信仰.
共同得真事.
收藏在我們身邊的祖堂裡.
救馬我們土地所產.
十分之一.
奉獻給美人.
因您美人在我們一切成就.
到九三中.
當最十分之一.
您美人最十分之一的十分之一.
大聖帝至此中.
當我們在世.
與您美人共聚.
您美人也從十分之一.
到我們十分之一.
奉獻給我們十分之一.
收藏在我們身邊的祖堂裡.
十分之一的奉獻是神聖的.
傳歸耶和華.
摩西的制度是.
人要將土產帶到神面前.
不過尼希米.
將這個制度改為由.
您美人親自去郭城鎮.
向民眾收取.
聖殿主動收取奉獻.
是一個制度上很大的突破.
當您美人走到百姓那裡.
收撒奉獻的時候.
那些百姓是不能推諉的.
如果教會每個月.
都派人來向您.
喂 十一奉獻.
那您會怎樣呢?.

$^{481}$您會有很大壓力嗎?.
奉獻是甘心樂意的.
神從來不勉強人奉獻.
教會也不需要主動要求會有奉獻.
您美人收取奉獻的時候.
要由亞倫子孫的祭司.
在旁監督.
這是一個管理制度.
以防有人中飽私囊.
虧空供款.
今天教會需要有兩個人去數點獻金.
以及刊登徵信錄.
都是同一個道理.
要確保會有的奉獻.
是真正歸在教會.
您希米明白.
聖殿是敬拜和信仰的中心.
要讓聖殿有效運作.
以及發揮信仰承傳的作用.
就需要有足夠的奉獻.
去維持各項事奉的開支.
因此.
由信心奉獻.
十一奉獻.
獻上最好.
以至將奉獻融入敬拜當中.
其實.
您希米的心.
都是很想聖殿能夠持續有效運作.
因此您希米要求百姓.
為此納藥.
確保聖殿能夠按照神的心意.
去展現那種信仰精神.
新約沒有提到.
信徒要受一奉獻.
但在羅馬書十二章裡.
本來要求信徒.
將身體獻上當作活祭.
是聖潔的.
是神所喜悅的.

$^{521}$你們如此事奉乃是理所當然.
當信徒願意將整個生命奉獻給主的時候.
不需要再提及十一奉獻.
您整個生命都奉獻了.
弟兄姊妹.
教會需要有穩定的奉獻.
才能夠推廣聖功.
這是不爭的事實.
我們願不願意與神納藥.
持守十一奉獻呢?.
您願不願意獻上最好的呢?.
您願不願意今年有信心奉獻呢?.
今日的講題是重新納藥.
在座每一個受浸加入教會的肢體.
在我們受浸的時候.
都曾經與主耶穌納過藥.
納藥的內容是.
確保了主耶穌是我們生命的主.
拯救的主.
以致我們每一個都能夠稱為神的兒女.
得享神賜予給我們永恆的恩典.
之後我們在淺行信仰的過程中.
我們有沒有再與神納藥呢?.
或許有或許沒有.
不過無論有或沒有.
現在都是一個很重要的時刻.
在今日.
紅映加撒週年的唐慶崇拜中.
我有機會可以再次與神納藥.
在此之間.
我帶大家有一個納藥的默想.
然後將大家在神面前所納的心願.
呈獻給上主.
我們一起安靜閉上眼.
親愛的主.
當我們回顧今天講到的訊息的時候.
請你進入我們的心裡.
引導我們去作出一個回應.
今日無論我們願意納至過一個.
分別為聖的生活.

$^{561}$或者願意每個週日都以學務的心.
回教會親近神.
領受神的話語以及猜險.
又或願意將生命中最好的獻給主.
在金錢上去持守受一奉獻.
或者定義今年有一個信心奉獻.
抑或你有自己獨特的納藥內容.
無論是哪一種.
無論是哪一種的納至.
哪一種的承諾.
都是心甘樂意.
而不是出於勉強.
我們每一個人都是按著聖靈的感動.
不需要懼怕.
我們以信心去回應.
因為神與我們同在.
又與我們同行.
神願意與我們一起走這條信心.
以及誠信之路.
神昔日曾經拯救過我們.
今日祂與我們同在.
明日也一樣會保守看顧我們.
因為我們所相信的主耶穌.
是一個色在,金在,永在的主.
祂的應許從今時直到永恆.
今日我們能夠與一個永恆的主.
再重訂一個納藥的關係.
是我們的榮幸.
也是我們的福氣.
因為神是信實的.
祂永遠都不會失約.
所以這一刻.
我們安靜地去感受一下.
神在哪一段經文.
或哪一個教導.
觸動了我.
讓我擇心.
讓我願意為主獻上什麼.
我想與神納一個什麼的藥.
我們只需要用心靈.

$^{601}$告訴神就可以了.
親愛的主,我們要讚美你.
因為你沒有因我們的過犯而離棄我們.
你接納我們每一個的缺點.
你以溫柔的心與我們同行.
等待著我們每一個的成長.
親愛的主,與你納藥.
是會戰戰兢兢的.
因為我們每一個都是軟弱的.
我們很怕會失信於你.
因此求主你額外的加能設給我們.
讓我們在納藥之後.
有足夠的信心和能力.
去踐行我對你們的承諾.
親愛的主,今天是一個.
每一個在你面前納藥的.
都是寶貴的.
是我們誠心誠意的心願.
也是在聖靈引導下.
我們對你的回應.
深信神會預納我們每一個愛你的心.
求主建立我們的心.
也在我們乘勝的旅程裡.
與我們同在同行.
願我們整個生命都能夠被你使用.
讓主的名被高舉.
主的名被廣傳.
最後納藥也體驗了.
我們雄恩家信徒合一的心.
求主繼續帶領著雄恩家.
在往後的日子.
能夠為主結出累累的果子.
我們誠心仰望祈禱.
奉教主耶穌基督得勝明智祈求.
主祝福大家.
謝謝大家.
\newpage



\section{馬太福音 7:24-27}
\label{sec:VN5NAedAhq4}
\textbf{2022.11.13 《建立堅不可摧的生命根基》馬太福音7\_24-27(和修版) 曾敏傳道}
\newline
\newline
連結: \href{https://youtube.com/watch?v=VN5NAedAhq4}{\texttt{https://youtube.com/watch?v=VN5NAedAhq4}} ~~~~ 語音日期: 2022-11-13
\newline
\newline
\hyperref[sec:WZ8nJrtvmJs]{\small{< < < PREV SERMON < < <}}
~
\hyperref[sec:index_chronic]{\small{[返順時目]}}
~
\hyperref[sec:index_scriptual]{\small{[返順卷目]}}
~
\hyperref[sec:bh_7j2paksc]{\small{> > > NEXT SERMON > > >}}
\newline
\newline
馬太福音 7:24-27
\newline
\begin{longtable}{cl}
\hline
\hline
章節 & 經文 (和合本修訂版)\\
\hline
7:24 & \begin{tabularx}{0.7\textwidth}{X} 「所以，凡聽了我這些話又去做的，好比一個聰明人把房子蓋在磐石上。 \end{tabularx} \\ \\ \relax
7:25 & \begin{tabularx}{0.7\textwidth}{X} 風吹，雨打，水沖，撞擊那房子，房子總不倒塌，因為根基立在磐石上。 \end{tabularx} \\ \\ \relax
7:26 & \begin{tabularx}{0.7\textwidth}{X} 凡聽了我這些話而不去做的，好比一個無知的人把房子蓋在沙土上。 \end{tabularx} \\ \\ \relax
7:27 & \begin{tabularx}{0.7\textwidth}{X} 風吹，雨打，水沖，撞擊那房子，房子就倒塌了，並且倒塌得很厲害。」 \end{tabularx} \\ \\
[1ex]
\hline
\hline
\end{longtable}
$^{1}$各位弟兄姊妹平安.
敬拜祂都是一件很蒙恩的事情.
我們一同再禱告.
上上要藉著你的話語去提醒我們.
今天都很願意你開我們的心眼.
讓我們看見你.
開我們的耳朵.
讓我們聽到你的聲音.
願意聖靈在我們心裡去動功.
讓我們明白神的心意.
我們就去回應你.
奉耶穌基督的名字祈禱.
阿們.
我們剛剛上星期是一個很歡喜快樂的日子.
我們的十一週年的堂慶.
弟兄姊妹在裡面都有很多的參與.
無論在聚會的當日.
或者之前的籌備工作.
我看見弟兄姊妹真的很愛教會.
現在其實大家仍然可以再一次重溫.
上星期堂慶的一些很開心的情況.
大家可以上去三樓的瓷畫牆.
貼上大家的心心.
還記得當日大家在心心的紙上寫下心願.
我希望教會怎樣.
我希望自己的生命怎樣成長.
你看著這張心心紙.
就應該想起堂慶當時的快樂.
在這些心心紙裡寫下我們很多的心聲.
有很多弟兄姊妹很願意在新一年.
讀經,祈禱,傳福音.
每樣都成長.
亦都很想教會有生命的成長.
有人數的成長.
甚至能夠早點獨立等等.
在這裡反而我要提出一件事.
我相信這些心願都是很寶貴的.
我們亦都要將這些心願放在神的手中.
看看上帝是怎樣帶領我們.
除了對於整間教會.

$^{41}$我更加要問一個問題.
我們自己這個人準備好了沒有?.
在未來的一年或未來的幾年.
我相信我們會面對很大的挑戰.
亦都會有很多的考驗.
在這裡要問自己的問題.
我們準備好了沒有?.
我們生命的工程.
我這個人.
我們的生命的工程.
是否能經得起考驗呢?.
我們生命的根基是建立於甚麼之上?.
這個根基是不堪一擊的?.
一擊即散的?.
還是堅不可摧的?.
很穩固的呢?.
這些問題都要問的.
心願歸心願.
考驗挑戰就在前面.
但要看的是我們自己準備好了沒有?.
如果問你這些問題.
你打算怎樣回答呢?.
可以怎樣回答?.
所以今天很想和大家回到.
《馬太福音》七章二十四至二十七節.
就是剛才大家讀的經文裡面.
但其實這段經文是一個更大的段落裡面.
是七章十三至二十七節.
裡面的內容是環環相扣的.
原來不可以單單看那幾節.
我們稍稍看看前面的經文.
重點是今天二十四至二十七節.
在這段七章十三至二十七節之前.
其實是另一個大段落.
就是登山寶訓.
耶穌基督在那裡對於門徒.
對於一群會眾.
有很多很多的勸勉.
講述天國的子民應該有的生命質素.
講述了很多非常豐富.

$^{81}$那裡的教導如果要看的話.
應該要講述很多唐道.
或者上一個很豐富的主日學.
去到這裡的時候.
來到今天.
在登山寶訓之後.
耶穌提出四個警告.
祂教了很多很多的東西.
然後祂要用一個很嚴肅.
很認真的態度去提出警告.
要勸勉門徒真正的委身在信仰裡面.
教完了.
我們的回應是甚麼.
這個回應是很重要的.
而這些回應是很特別的.
是有影響的.
我怎樣回應.
就會有怎樣的結局.
第一個警告.
其實就是一個窄門小路.
對比闊門大路的警告.
窄的門令人想到不容易進入.
小的路亦不容易走.
在這裡其實.
耶穌要讓我們看到.
窄門小路是甚麼呢.
原來就是一條.
我和你都必須要踏上去的路.
都必須要經過的門.
為甚麼門會窄路會小呢.
因為有很多的挑戰考驗困難.
而在經文裡面說的是.
當時的教會是面對迫不得.
面對患難.
所以走這條路不容易.
門很窄路很小.
但你要知道.
這條路這門去到哪裡.
將會帶人進入永生.
是永遠和上帝一起.

$^{121}$我們在上帝的國度裡.
原來要經過窄門小路.
要經過迫不得,患難,困難.
要付代價.
是不容易的.
但只能夠這樣選擇.
耶穌都勸我們.
要這樣選擇窄門小路.
但另一個對比就是.
要對比闊門大路.
在這裡我們不需要面對迫不得.
患難,困難,很舒服.
甚至很大的吸引力.
這條路很闊很容易走.
門又很闊可以容納很多人.
所以很多人會經過闊門.
很多人會走在大路上.
但經文很清楚地說.
闊的門帶領我們走向滅亡.
大的路亦都是走向滅亡.
所以縱然是這麼容易,舒服,有吸引力.
是不需要經過困難的.
耶穌叫我們不是這樣選擇.
不要選擇走向滅亡的路.
不要進入走向滅亡的門.
反而要經過窄門小路.
這是一個很嚴肅的警告.
選錯了結果就不一樣.
這是第一個.
看回這個圖片都很大分別.
窄的路和闊的路很大分別.
窄的路和窄的門的確很少人能進去.
因為要付代價.
很多人都不想付這個代價.
這幾個警告其實不是對著非信徒說的.
是對著信徒說的.
對著門徒說的.
意思是挑戰我們.
你選擇的時候.
想好了沒有?.

$^{161}$是否真的為信仰付代價?.
是否忠心於我的信仰?.
知道要這樣付代價也要走下去?.
耶穌要挑戰我們.
這是第一個警告.
第二個警告是好樹結好果子.
壞樹結壞果子的警告.
一個這樣的對比.
我們很自然會想到.
好的樹自然有好的結果.
壞樹不可能結出好的果子.
特別讓我們看到.
壞樹結壞果子.
原來耶穌要我們在這裡防備假先知.
這些假先知.
他們就是這樣.
看起來像很純良的羊.
這些人都說自己要跟從上帝.
要好好服侍上帝.
甚至乎他都說要出去傳福音.
要為神做很多很多事情.
嘴巴說得非常響亮.
非常漂亮.
很好.
就像羊一樣.
但實際裡面就是狼.
是一個殘暴的狼.
實際裡面是敵黨真理.
是敵黨福音.
他的目的是拆毀傷害.
這些假先知會傷害群體.
會令人走上歪路.
耶穌要我們防備這些.
那你怎樣分辨得出他是羊還是狼呢?.
看他的見證.
好樹結好果子.
這個果子是生命的見證.
你從見證看這個生命是怎樣的生命.
從果子看這棵樹.
是不是一棵好的樹.

$^{201}$我們就從這些人的生命見證看.
他是真的跟從上帝呢?.
是善良的?.
還是傷害人生命.
破壞群體的.
令人走歪路的呢?.
就透過他生命的見證.
去分辨得出來.
耶穌要我們防備他們.
一來不要受他們影響.
二來也不要成為他們的一份子.
因為這些結壞果子的壞樹.
會有他的結局.
是面對火的審判.
後果非常嚴重.
正如剛才走去邁向閘門.
走小路會得到永生.
永遠的生命.
但如果走向闊門.
走上大路.
就會面向死亡.
是永遠的死亡.
這裡說的是這批人.
將會面對火的審判.
這是第二個.
第三個就是我曾經說過.
有一次我說過.
我從來不認識你們.
這堂道其實是來自這段經文.
是說對比表面.
跟從主的人.
對比真正遵從天父旨意的人.
這是第三個警告.
第三個對比.
是很重要的.
有些人總是會跟主說.
我為你傳福音.
我傳你的道.
去趕鬼等等.
我為你做多少事.

$^{241}$你看看我外在的行為是多麼好.
但主責備他們.
說他們沒有真正遵從天父的旨意.
這班人沒有真正將自己的生命.
放在耶穌基督的權下.
真的聽上帝的聲音.
去遵從.
他們跟上帝沒有真實的.
確切的關係.
所以只有外表的行動.
主責備這些人.
甚至說他們是惡人.
這班人.
耶穌基督不認他們.
如果我們要做這些.
表面跟從主的人.
結果都會跟他們一樣.
主會不認他們.
主視他們為惡人.
他們在主的內部是沒有份的.
哪怕在外面做了很多事.
對比真正遵從天父旨意的人.
不一樣.
真正遵從天父旨意的人.
明白上帝的心意.
真的會去跟從.
跟上帝是有關係的.
生命是會悔改的.
是不一樣的.
這些人有很大的對比.
來到這個警告的時候.
經文不是停在那裡.
耶穌基督不只是停在那裡.
接著來到今天我們要說的那段經文.
是一個很重要的.
最後的重錘出擊.
甚至可以說是一個總結.
這個第四個警告.
是建基於之前所說的一切.
不單止是之前的三個警告.

$^{281}$還有再上提過的登山部分.
那一大堆上帝自己的教導.
聽了耶穌的教導之後.
我們的回應是什麼.
這個引入第四個警告.
我們的回應將會決定我們的結局.
現在不是說我們有很多選擇.
你說得這麼絕對.
要麼你回應.
要麼你不回應.
你只有兩個選擇.
我中間行不行呢.
其實沒有中間的.
是沒有中間的.
你是回應.
要麼就不回應.
沒有中間的.
而這個回應將會決定我們的結局.
是一個很嚴重的影響.
接著經文就直接提出.
回應.
上面說了一大堆.
很多很多很多教導.
甚至耶穌提出警告.
去到這樣的地步.
這一刻我的回應是什麼.
如果我聽了耶穌這一切的話.
我又去做的話.
我的回應是.
我聽了.
我會做的.
我有回應.
那麼這個結局是什麼呢.
這裡用了一個比喻.
它說好比一個聰明人.
將房子蓋在盆石上.
聽了耶穌這麼多的教導和警告.
你聽了.
你去做.
你就是聰明人.

$^{321}$聰明人是怎樣的.
其實不是盲目的.
如果用聰明來形容.
這個人有智慧.
有思想.
我聽.
我是聽得明白的.
我聽得明白.
我思想.
我有智慧.
接著我什麼.
我馴服.
我做出來.
這樣叫聰明人.
這種回應.
耶穌說.
就是聰明人了.
將他的房子蓋在盆石上.
盆石給人的印象.
其實是非常堅固.
不會塌的.
真的很堅固.
你將房子建在盆石上.
怎樣都不會塌.
所以經文繼續說.
風吹雨打水沖.
怎樣撞擊房子.
防止總不倒塌.
因為根基穩固.
根基決定了房子.
上面的結構.
它的工程.
是否經得起考驗.
根基穩固.
上面也穩固.
根基不穩固.
上面就不穩固了.
聰明人建房子.
就會找些堅固的根基.
所以盆石是最好的選擇.

$^{361}$無論外面經歷什麼風浪考驗.
房子是不會塌的.
就像這間.
看上去非常穩固.
堅固.
無論什麼情況.
屹立不倒.
聽了耶穌的話.
去做的人.
代表了我們生命的根基.
是不被動搖.
我們生命的工程.
也得以堅固.
這是很重要的.
對比另一種人.
另一種人就是.
耶穌說了一大堆教導.
他也有提出警告.
到現在這一刻.
我聽了.
我不回應.
我不去做.
人是有自由意志的.
上帝也給我們選擇.
上帝不是強行壓低我們.
不聽.
不會.
你怎麼選擇.
我選擇不聽.
不聽.
我又不去做.
我入了耳.
但不會對我的生命有任何影響.
我不做.
那怎麼辦.
不做.
耶穌說.
這些是無知的人.
沒有見識的人.
沒有智慧的人.

$^{401}$聽了也不明白.
又不懂思考.
我又不信服.
也不會進一步行動.
那你是不是無知的人.
把房子建在沙土上.
想也想不出來.
沙土是什麼質地.
完全不行.
風吹雨打水沖.
一撞就沙土.
捱也捱不了.
一撞就倒了.
經文還要說.
還要倒塌得很厲害.
我們的根基.
這麼不堪一擊.
這麼軟弱.
你覺得上面的工程.
怎麼會立得住呢.
正如你今天去沙灘.
在沙上建一些東西.
很快海水一過來.
一衝就沒了.
風一大一點.
一吹就沒了.
是一樣的.
這房子的工程不行.
因為根基不行.
我們聽上帝的話.
不做就是這樣.
就好像這房子.
你在上面建多少都沒用.
人家看這房子很宏偉.
但是考驗一來到.
就倒塌了.
為什麼.
因為下面的根基不行.
聽了耶穌的話.
不去做.

$^{441}$我們的生命根基.
不堪一擊.
我們的生命工程.
會全然崩壞.
這是一個很大的提醒.
弟兄姊妹在這裡.
其實重點是什麼.
說了這麼久.
如何建立一個.
堅不可摧的根基呢.
最重要的是.
聽了耶穌的話.
就去做.
我們要做聰明人.
做聰明人.
把房子建在盆石上.
聽耶穌的話.
就去做.
曾經我在這裡提過.
一年我們聽52唐道.
有些講員是重複的.
但是也有不少的講員.
可能是單次性的邀請他過來.
或者有些並不是.
經常在這裡講道.
其實在這裡也有很多很多訊息.
52唐道.
我們聽了多少.
我這星期說.
如果我下星期問你.
還記不記得.
除了記不記得.
那你有多少唐道是活出來的.
你說聽了耶穌的話就去做.
那你聽了就要去做.
那有多少做了出來呢.
聽了祂的話.
左耳入 右耳出.
這個意義在哪裡呢.
聽了不去做.

$^{481}$對我的生命也沒有幫助.
我的生命根基不會穩固.
我的生命工程也不會很堅固.
所以你說接下來12年的唐慶.
13年 14年唐慶.
每一年我們都可以搞很多活動慶祝.
但是我這個生命還沒有準備好.
如果我的生命建築在沙土上.
上面砌了很多很漂亮的東西.
別人看起來很像樣.
不用的 很快的.
一個風浪來就倒了.
一個考驗來就倒了.
我可能就不想回教會了.
我不想再事奉.
我不想打開聖經看.
可以馬上倒.
就是這麼簡單.
其實這樣很難再繼續砌上去.
很難說.
哎呀 上帝我要在這裡拓展你的國.
上帝我要帶領更多人來到教會.
我自己的豆腐渣工程.
說得白一點.
豆腐渣的工程.
怎樣去承載別人的生命呢.
我記得上個星期.
去到我們有一個遊戲的環節.
很寶貴.
突然姐妹出來說很多感恩的事.
我覺得很珍貴.
在當中真的很珍貴.
看到他們的心聲.
為他們獻上感恩.
但是不足夠.
我們要繼續邀請人進來我們當中.
繼續拓展神的國.
我們的生命是怎樣.
我亦記得潤暢有分享過.
我們是否能承載別人的生命呢.

$^{521}$豆腐渣的工程.
可以承載多少人.
你不要說別人進來.
未曾進來我們才會倒.
所以今天挑戰我們.
建立一個堅不可摧的根基.
聽神的話.
就要去做.
有一堂崇拜應該是感受堂慶.
崇拜那天.
我很有感受對於講員的講道.
完了之後.
我都有在不同的群組分享.
我對於他的講道的回應.
心裡面很有感動.
真的很想.
接下來我們去參與教會的事奉.
參與不同的活動.
我都很想是這樣.
我們以信服神為重點.
講員講得很好.
信服神.
你可以開很多會.
你可以有很多策劃.
但不能沒有重點.
就是信服神.
甚麼是信服神呢.
上帝的心意是甚麼.
我怎樣知道上帝的心意呢.
禱告吧.
你不問上帝.
怎樣知道上帝的心意呢.
禱告 等候主 看聖經吧.
上帝在聖經裡面對我們說話.
不只是對一個人說話.
是對整個群體說話.
我們是否一起看聖經呢.
大家告訴我.
一個星期大家花多少時間看聖經.
最多是否今天呢.

$^{561}$大家笑吧.
上帝在十時後就煮了一盒.
我們一定會看聖經.
然後十一點半我們有崇拜.
講員一定要講聖經.
離開這間教會之後.
大家還有多少時間看聖經呢.
不敢回答我.
多少時間呢.
兩小時吧.
如果是兩小時怎樣計算呢.
每天一點點.
每天一點點算起來.
就是兩小時.
有人回答我.
還有人回答我嗎.
每個星期讀多少.
怎樣明白上帝的心意呢.
怎樣去信服呢.
怎樣知道祂要我們走去哪裡.
你說聽祂的話.
就要去做出來.
我都沒有聽祂的話.
怎樣做呢.
所以我們要讀聖經更多.
更多去禱告.
明白上帝就這麼簡單.
但要做出來.
我希望將來的時間.
看到我們的生命.
真的成長.
經得起風浪.
不要兩下打風就倒了.
不要豆腐渣工程.
要上帝眼中看.
經得起考驗的生命.
我們一起祈禱.
你要提醒我們.
聽到 行到.
真正的回應你.

$^{601}$其實我們好懶惰.
其實我們好少去看聖經.
都不知道你的教導是甚麼.
每個星期講完講完到之後.
其實很快我們就已經忘記了.
談不上甚麼行到.
主求你寬恕我們.
寬恕我們可能.
我們將更多的時間放在.
看電視 看手機.
那些可以花上很多很多個小時.
都不覺得悶.
祂不會將時間預留給我們.
天父求你寬恕我們.
正正因為這樣.
所以我們生命無力.
主耶穌基督求你幫助我們.
重新調教我們生命的優先次序.
賜我們渴慕你話語的心.
賜我們渴慕聆聽你的心.
讓我們想聽你聲音.
想知道你的旨意.
讓我們想信服 想跟從.
做出來.
主要讓我們聽完到.
不是就這樣算 而是有行動.
願你賜福給我們.
讓我們將來的生命更加健壯.
真真正正成為一間.
建立在盤石上的教會.
一間勝在別人生命的教會.
真正在洪水橋區.
能夠拯救多人的生命.
求主帶領我們.
奉耶穌基督的名字祈禱.
阿們.
\newpage



\section{但以理書 2:18-6:10-9:3}
\label{sec:bh_7j2paksc}
\textbf{2022.11.20 《讓但以理的生命勒著你和我》但以理書2\_18,6\_10,9\_3(和修版) 高淑芬傳道}
\newline
\newline
連結: \href{https://youtube.com/watch?v=bh-7j2paksc}{\texttt{https://youtube.com/watch?v=bh-7j2paksc}} ~~~~ 語音日期: 2022-11-25
\newline
\newline
\hyperref[sec:VN5NAedAhq4]{\small{< < < PREV SERMON < < <}}
~
\hyperref[sec:index_chronic]{\small{[返順時目]}}
~
\hyperref[sec:index_scriptual]{\small{[返順卷目]}}
~
\hyperref[sec:Hx9S3W_wWOQ]{\small{> > > NEXT SERMON > > >}}
\newline
\newline
$^{1}$Sawadee ka.
Sawadee dong chow.
哈哈哈哈.
Sawadee dong chow就是早晨.
剛才Sawadee ka.
你整天都可以用的.
今天這個題目.
因為是唐慶月.
正如剛才福音隊.
樂曲福音隊也有提及.
這個整個月都是我們的唐慶月.
不知道大家有沒有走上三樓.
看板你們的心心.
很多大部分都是希望.
我們雄恩堂裡面是越來越多人.
大家都可以一直去成長在靈明.
不知道你們做好準備沒有.
如果人多.
即是說你隔離的位置不可以放袋子.
隔離的位置就變成人.
或者.
走廊不夠大.
可能還要加椅子.
很多東西我們要準備.
正如今天的崇拜.
有沒有發現有些東西已經變更了.
你們頭上的球球變成什麼顏色了.
上個星期綠色就叫「屎橋」.
上個星期還是綠色.
還是「屎橋」.
現在變成粉紅色了.
其實我們教會每一個主日.
都有很多不同的.
有沒有欣賞一下這些花.
上帝所創造的花有沒有帶領你.
更加在敬拜裡面.
更加歡喜快樂地來誇獎神呢?.
心裡面更加感恩.
當看到唐慶教會已經有十一週年的時候.
很開心.

$^{41}$我們很希望神繼續帶領我們向前行.
甚至越來越多人在這裡一起去敬拜神.
我就想起「但以你的生命」.
「但以你的生命」我希望.
祂的生命可以「勒著你和我」.
「勒著」的意思是.
這個火是不熄滅的.
繼續是有火的.
我們經常說不要忘記了你起初對神的愛.
這個愛就好像那個火一樣.
不可以熄滅的.
有時候跟弟兄姊妹說.
我都很怕我會冷淡.
但是我們團在一起的時候.
我們就可以彼此勒著.
今天很想透過「但以你」.
祂有三個祈禱來成為.
因為「但以你」書其實全卷書是十二章.
其實是很好看的.
兒童主一學都是用來做教材.
如果我說「獅子坑」.
記不記得?.
又記不記得「但以你」的三個朋友.
進入了火坑.
但是「獅子坑」也好.
「火坑」也好.
生命沒有受損的.
繼續可以榮耀神的.
「但以你的生命」.
不是一路路坑通的.
但是祂的生命很值得我們去看.
以致這個生命.
勒著你跟我一起.
更加來合乎神的心意.
「但以你」的名字.
意思就是神是我的審判者.
這是祂的名字的意思.
祂出生是猶大王族的貴族.
不是你和我這樣平凡人.
我們是平民.

$^{81}$但是祂是貴族.
其實對於貴族這件事的觀念.
因為從小到大都在香港.
我不是怎樣領會這件事.
但是當我在泰國裡面的時候.
發現地上的王.
在泰國裡面.
這麼多國家裡面的一個國家.
泰國.
王是有著他的遺訓在當中.
還有是說.
看泰國人.
透過你的姓.
就知道你出生於平凡的一般人.
還是貴族.
又怎樣貴.
看你的姓就知道.
所以「但以你」不要小看祂的出生.
祂的出生是一個王族.
但是你會看到.
祂的生命不是這麼一帆風順.
雖然已經陷入了一條金鎖匙出世.
但是在祂十多歲的時候.
「但以你」已經成為俘虜.
被人用鎖鏈拉著祂走去巴比倫國.
祂本來是猶大國.
即是以色列民的一份子.
一個貴族.
一個人上人.
但是在公元前的607年.
你就會看到在約瓦敬王的第三年.
以色列國已經沒有了.
最後的一個君王.
都要變成一個俘虜.
王也是.
「但以你」貴族也是.
都是這樣鎖鏈鎖著去巴比倫國.
當時十六歲的「但以你」進入了巴比倫.
當時應該是一個懷柔政策.
皇帝,貴衛.

$^{121}$巴比倫國.
把他們.
即是善待他們.
可以說是軟禁他們.
以致他們的人民都會聽話一點.
除了這樣的時候.
巴比倫王.
更加把他們一些有潛質的年輕人.
年輕人.
十多歲.
因為希望他們繼續有貢獻.
就像現在香港一樣.
要吸納一些人才.
在這個徵選.
這班奴隸.
這班俘虜裡面.
還要選擇.
不是說有知識就可以.
也不是說聰明.
還要外表.
沒有殘疾.
我們這些就不行了.
戴眼鏡都是殘疾之一.
還要相貌俊美.
好像選美一樣.
更厲害.
智慧與外貌並重.
你會看到他們就是要有這些俘虜.
成為在王宮裡面.
來侍候的一班僕人.
於是.
但以理也成為了選上的一個代表.
但是他的名字以後不是叫做但以理.
改為巴比倫文的一個文字.
意思是要保住王的名.
你現在接下來的生命.
你就是要保住我巴比倫王的生命.
建立我這個巴比倫王國.
但以理.
大家很熟悉.

$^{161}$他剛剛開始被人俘虜的時候.
要受訓練的時候.
訓練了三年.
在這三年裡面.
他們得到很好的招待.
吃盡佳餚美酒.
因為王喝的東西.
吃的東西.
他照樣可以吃.
但是.
記不記得但以理吃什麼.
吃菜.
吃素菜.
但以理吃素菜.
喝水就行了.
為什麼會這樣呢.
因為在他們的以色列傳統裡面.
他們是不可以去吃一些獻祭的食物.
所以你看到但以理和他三個朋友.
他們一起都不要皇帝的珍饈百味.
他們選擇吃素菜和喝水.
但他是一個有情有理的人.
他看著他那個.
說不可以的.
你們一會變瘦到像排骨一樣.
皇帝就要找我來算帳.
但是在這個時間裡.
但以理說給我們十天試試.
給我們試試.
結果他們吃素菜.
這群人吃素菜.
就這樣喝水.
但是他們的面貌仍然很好.
在這裡有很多特別姐妹.
我都要像但以理一樣.
吃素菜好了.
千萬不要.
這是神蹟.
這已經是在但以理身上的一個神蹟.
因為你只吃素菜.

$^{201}$吃水你的營養是不均衡的.
你是要有蛋白質在裡面的.
所以你看聖經要看上文下理.
看他的歷史背景等等.
在這個過程當中.
你會看到但以理在被俘虜的時候.
其實人生都很慘.
自己的國家沒有了.
如果給了你和我的時候.
自己的國家亡了.
都是家破人亡的.
不知道在打仗裡面.
但以理的家人有沒有被殺害.
有沒有更加慘的遭遇.
不知道.
但是但以理.
他雖然亡國.
他仍然敬畏創造他的耶和華上帝.
無論生命是去到劣境也好.
他仍然堅持著上帝仍然是他的上帝.
在他的生命裡面不惹小的.
他經歷過兩個國家.
一個是巴比倫的國家.
還有波斯的帝國.
他經歷的亡國都有四任的皇帝.
是很特別的一個經歷.
我在第一個分題裡面.
他勇敢呼喚他的同伴一起去祈禱.
在但以理二章十八節.
我請匯仲一起中讀.
(念句).
因為當時尼布甲尼撒王就做了一個夢.
但是他又不告訴別人夢境是怎樣.
要那些聰明的人.
那些哲士文士官星學家博士.
要他們解釋這個夢的意思是甚麼.
每個人都站在那裡.
真的不知道怎麼辦.
王你講一下夢來聽.
我就是不講你們就要解釋給我聽.

$^{241}$於是沒有人解釋到.
他很生氣.
他生氣到尼布甲尼撒王.
所有這些聰明人全部幫他殺了.
或者全部殺了人才.
這個很生氣的裡面.
你就會看到人憤怒的時候.
殺傷力是很大很大的.
如果真的巴比倫國裡面的所有聰明人殺了.
這個國家還有沒有得剩.
沒有了.
所以當我們真的面對別人憤怒的時候.
而你仍然是清醒的.
仍然是很鎮定的時候.
請你在心裡靠著上帝.
來面對這個憤怒的人.
當然如果這個憤怒的人有暴力傾向.
你要保護自己.
但是憤怒.
所以聖經裡面特別講到.
憤怒不可憤怒到何時.
日落前你不要再憤怒了.
憤怒一段時間.
我們不要再將憤怒放在心裡.
如果但以利真的含著憤怒.
這個巴比倫國滅了我們的國.
現在又俘虜我們.
裡面是可以有很多很多的憤怒.
是可以有理由的.
但以利沒有這樣做.
而且我們看這裡就知道.
尼泊爾格尼撒王都不是蠢的.
不過因為憤怒所講出的語言.
所做出的行動.
一切都是很蠢的.
所以求主祝福我們每一個憤怒.
一段時間好了.
還要很短很短的時間.
日落前好收手.
一路禱告求神去幫我們.

$^{281}$但以利很有勇氣地.
其實與他無關.
因為但以利是一個俘虜.
在這個時間裡面.
是一個寂寂無聞的俘虜.
和一班聰明的人在一起工作.
可能他做的事情.
都不是太重要的事情.
因為始終他是敵國的一個俘虜.
他不是巴比倫的子民.
但當但以利聽到這個消息的時候.
他勇敢地去見皇帝.
請求寬限.
跟著說我的神會將這個夢顯明給我.
讓我知道.
跟著我會告訴你.
但以利有這份勇氣的時候.
他不只是自己祈禱.
他知道他有三個屬靈的同伴.
和他一起的.
他們一起試過不吃皇帝的真羞白味.
於是他就邀請他的三位同伴.
和他一起的祈禱.
去勇敢地去接受這個挑戰.
其實他可能不用死.
因為他是外邦人.
俘虜可能未必會殺到他.
但他有這樣的勇氣.
竟然抱自己在面前.
陷入這個危機的裡面.
或者是轉機.
危機會變成轉機.
這樣很奇妙.
神真的使他知道這個夢的解釋.
當但以利去見皇帝的時候.
他不是立刻心急地告訴尼布甲尼撒王.
你的夢是這樣這樣.
他是很強調.
我再次鼓勵大家去看《聖經》.
但以利書很好看.

$^{321}$一到十二章裡.
一到六章絕對好看.
七到十二章也不是不好看.
不過可能有些異象.
但我們還是可以接近的.
整本《聖經》我們可以接近的.
正如我們剛剛信主.
《聖經》有些東西我們看完不明白.
不明白不要緊.
特別你可以回來要求.
我不明白.
我們開一班但以利書做主好不好.
我相信老牧師一定額頭答應的.
所以大家如果一直看《聖經》.
有些東西不明白.
我們會不會看這一卷.
在團契在主一學裡.
或者我們一起可以再繼續看下去.
再深深的研究下去.
在這一段的說話.
就是說只有一位神.
是天上的神.
可以幫你將你的奧秘顯出來.
在這裡的經文裡.
更加再強調一件事.
不是因為我這個但以利.
不是因為我個人的智慧.
來幫助你.
而是我背後.
我心裡的上帝耶和華神.
去解答你的問題.
當時當但以利說完夢的解釋的時候.
尼布甲尼撒王.
是一個已經將以色列國.
全部俘虜戰勝的國王.
他竟然一聽完解釋之後.
向俘虜以色列人下拜.
並且像品臣一樣.
供物,香品當他神般拜.
你看到嗎?.

$^{361}$今天我們說我們很想人多一點.
多一點來得到福氣的時候.
你和我更加要掌握著.
上帝給我們的福氣.
將這份福氣讓我們身邊的人.
我們的家人.
我們認識的人.
都覺得你的神真是.
記不記得這個怎麼說?.
泰文是勁XX的神.
是一個很厲害的神.
在這裡上帝是勁XX的神.
以至尼布甲尼撒.
真的很自然地.
太感動了.
我這個夢沒有人能解釋.
你竟然知道我這個夢.
他完全不說這個夢是甚麼.
根本不知道你在做甚麼夢.
不知為甚麼.
但是竟然耶和華神可以藉著.
但以利來表達這件神蹟.
到今時今日.
你相不相信上帝同樣在你的生命裡面.
是像一個橋樑一樣.
將上帝的福氣留給其他人.
同樣透過你的生命.
可以讓人看到上帝就是那位.
施行神蹟的神.
在但以利被擄之後.
在這一刻裡.
他真的成為了宰相.
一人之下 萬人之上.
是一個很大的神蹟.
一個被俘虜的人.
竟然可以去到這個位置.
但以利也要求王.
讓他三位的好朋友.
靈命裡面的同伴.
一起在王宮裡面作事.

$^{401}$但以利就在朝廷.
即是最近皇帝.
我們看到剛才詩歌.
我們所唱的歌.
我們都看到我們所信的神.
是一個很厲害 勁勁的神.
而且在一個萬眾神觀念的.
舊約世界裡面.
他更加被巴比倫尼布格尼撒王.
譽為萬王之王.
或者是萬神之神.
如果有很多神的話.
耶和華神就是萬神之神.
是最厲害的神.
你今天和我有沒有這樣的認同呢?.
我們今天禮拜主日當中.
我們去敬拜神.
我們就是在敬拜一位萬神之神.
求主幫助我們.
讓我們知道我們所信的神.
實實在在是勁勁的神.
在這裡你會看到.
但以利他真的透過祈禱.
是個人的祈禱.
和同伴們的祈禱.
弟兄姊妹.
當你有團契 你有教會.
有你的主一學.
你一定在教會裡面有一個群體.
甚至如果沒有.
可能你只是回來崇拜.
我都強烈地建議你.
你參加一個團契.
參加主一學.
有一班是你的同伴.
建立一個靈明的同伴.
以致我們彼此的勒著.
不是勒著地憤怒.
彼此的靈明.
互相擦出繼續燃點的火花.

$^{441}$我的舞會有很多八十多九十多歲的長者.
在我還沒有去泰國的時候.
我要服侍他們.
我們有早午課.
晚課他們就不會去.
早午課他們多數是早課.
因為長者起得早.
他們經常說我們老了沒用了.
不過他們年輕的時候很厲害.
教會要買地方.
總之有些錢需要的時候.
我們都是平民百姓.
沒有人是貴族.
辛辛苦苦在沙廠做到.
有條鏈.
拔了條鏈出來.
奉獻.
他們說現在不行了.
我們老了.
但是我就經常鼓勵他們.
因為其實弟兄姊妹告訴我.
有時候弟兄姊妹.
年輕的都會懶惰.
或者偷懶.
五十二個星期.
一兩個星期請假.
但是當他們看到.
這班長者的弟兄姊妹.
很辛苦.
有些坐著輪椅.
以前的輪椅沒有電動.
要有人推.
有些滾著拐杖.
推著走路走得不穩.
都要來崇拜.
你知道嗎?.
這班長者的弟兄姊妹.
坐著在上去.
星期天來教會敬拜.
就是一個美好的見證.

$^{481}$同樣已經勒著那些.
打算下星期請假的那些.
回來.
所以今天你和我.
雖然還沒有到八十多九十幾.
但是你們回來.
你們在這裡敬拜.
你們同樣會勒著我.
你想想有一天.
剛好大家都找下星期做放假.
不來了.
這裡真的沒有什麼人.
只剩下那三個.
你猜那三個還能勒得起來嗎?.
勒不起的.
哇!真是!.
很多的疑問在當中.
所以你今天回來.
教會去敬拜神.
是一種彼此建立彼此的生命.
的一件事.
我們求主繼續.
記不記得「彼主用」的泰文是什麼?.
「踩」.
你願不願意被神踩?.
求主用我們的生命.
繼續的被神去踩.
被人踩.
如果神用得著.
就算了.
你就踩吧.
踩就是被神用了.
我們求神幫助我們.
禱告也同樣成為了.
但以理的生命的一部分.
在這個環節裡.
但以理已經不是在巴比倫了.
因為這個時段.
但以理在巴比倫的國家.
同樣被波斯帝國侵滅了.

$^{521}$現在這個皇帝.
是波斯王的大利烏王.
記住要看詳細.
很簡單一天一定會看完.
一至六章是很好看的.
在裡面.
你會看到大利烏王.
在這個情況下.
雖然他是俘虜.
還要是前一批的俘虜.
竟然去到波斯帝國.
仍然用這個但以理.
而且深得大利烏王的厚愛.
這個大利烏王很重用他.
以致到周圍的人非常妒忌.
你是甚麼人?.
這個但以理為甚麼可以.
比我們高的位置.
於是就設了一個.
連皇帝都察覺不到的詭計.
真是糞策.
糞策是綠色.
設了一個糞策.
出來陷害這個但以理.
禁令是怎樣說的呢?.
在這三十天裡面.
全國的人民.
除了大利烏王.
我們要拜祭之外.
其他的神.
其他的一切都不可以拜.
否則你就入.
扔你下去.
獅子坑.
你看到.
但以理是知道這個命令.
也知道這件命令.
如果違反了這條法例.
就是入獅子坑.
但以理仍然繼續.

$^{561}$他不是脈討.
如果給我機會.
我可能會走蜘蛛臉.
上帝我都是祈禱的.
我經常說.
我在泰國裡面.
我開車就是祈禱的.
因為操練了太忙了.
在心裡面繼續和神傾談.
神給我.
車不要撞我.
我又不要撞人.
因為我試過在泰國.
開車睡著了.
自從那次之後.
我連自己都不相信.
加上近期.
我29日就會去泰國.
探訪那班認識的.
特別是穆斯林群體.
我竟然買機票那一刻.
我的護照還在我手裡.
然後在上個月底.
我找到我的護照去了哪裡呢?.
我找了一個星期.
於是我又做夢.
做了好幾天.
夢裡面都是找護照.
到今時今日.
我還是沒有找到.
不過在那次的星期日.
主日完結.
星期一.
那星期日特別做夢.
做夢就是繼續在那裡找.
找到骨頭都痛了.
我還是不要找了.
快點去做吧.
去保齡.
那天打算去問問.

$^{601}$要甚麼資料.
要查查那裡有.
原來在元朗有.
很近的元朗有.
可惜電話打不通.
我不知道應該帶甚麼去保齡.
那好吧.
快快就隨便走走.
打算去找資料.
怎知道去到後.
人又多.
問到的時候.
都要三至四個星期.
而且你是不見了.
還要審核.
可能會久一點.
我一聽到日子.
我馬上要做了.
現在還可以做嗎?.
可以.
我給你一個手.
當時拍了一張.
全世界.
在我人生裡面.
最恐怖的相片.
眼睛不對焦.
總之.
我以後還未拿那本拍檔.
明天就拿了.
我想我都不敢看.
感恩.
但就會看到.
有一刻可能連自己都信不到.
但但以利.
他仍然在這個關頭裡面.
他知道後果是怎樣.
他仍然是一日三次.
還要跟隨他們以色列的傳統.
向著耶路撒冷這個城.
這個方向.

$^{641}$打開窗門.
來到一日三次的祈禱.
他本著一個大哥大的精神.
甚麼是大哥大呢?.
又是泰文.
大哥大就是死就死吧的精神.
今天從事了基督徒.
不知你多少年.
你有沒有在某些時段.
都試過有這種大哥大的精神呢?.
有沒有?.
我舉個例子.
你應該有的.
那個可能叫炒就炒的精神.
剛剛聽到一個姐妹說.
她以前做工.
每年有一次.
要全公司都要去葬香.
她是一個基督徒.
心裡面很心痛.
那怎麼辦呢?.
好了炒就炒吧.
炒就炒的精神.
就是一種大哥大的精神.
死就死吧.
這個真的是泰文.
大哥大在大義理書沒有寫這個大哥大這個字.
但是在以斯帖皇后.
當她要見皇帝的時候.
就是用這個字大哥大.
她說死就死吧.
這個心來見皇帝.
所以今天.
當我們說我們很想.
紅印堂越來越多人的時候.
有時候.
你要在我們神的旨意下.
這個立場當中.
你有沒有.
或者你願不願意.

$^{681}$附上這份大哥大的精神呢?.
未至於大.
大就是死.
可能炒就炒.
蝕就蝕.
輸蝕了一些.
這樣的精神.
意志到讓神的名氣高舉.
結果.
你知道當晚.
當大義理被奸人抓了.
因為他真的犯了法.
抓了之後.
皇帝其實很心愛他.
但是被奸臣.
這條例那條例原來就是想抓大義理.
他又不知道.
吸了鞋子.
皇命不能改.
大利烏王.
他說你的神應該可以救你.
反而他和大義理這樣說.
那一晚.
由大義理扔到獅子坑之後.
大利烏王一個王.
我要多少聰明人就多少聰明人.
但是他可以肺沉黃瘡.
他完全是吃不下.
睡不著.
原來宮殿裡有些音樂.
可以讓他開心一點.
不要.
他完全很想快點天亮.
一天亮了.
馬上跑到獅子坑那裡.
大義理.
你的神有沒有救你.
結果大義理無事.
沒事.
看到大義理無事的時候.

$^{721}$結果怎樣.
結果就是.
把大義理救回上來.
那些陷害他的人.
不止他那個人.
是他背後的家族.
他的子女.
他老婆父母.
全部扔到獅子坑.
對不起.
那些獅子真的很餓.
那天不吃大義理.
不是因為太飽.
是因為神職.
神封住了獅子的口.
到這一群人.
這麼多人.
扔下去.
聖經形容.
他還沒下去.
已經被獅子抬高咬著吃.
那些獅子飢餓到這樣.
大義理的禱告生活.
成為他生命的一部分.
哪怕要死.
他都這樣做.
你和我呢.
你和我有沒有好好地.
用這份權柄呢.
今天我們不是救藥時代.
你可以直接透過主耶穌基督.
我們可以去到父神面前祈求.
不只是祈求.
有時看到美好的花.
美好的事物.
讚美.
有時被人冤枉.
你很受委屈.
可能很多人都不能說.
但你可以和神說.

$^{761}$但有一件事.
我都很鼓勵大家.
你在你的團契裡.
或屬靈的朋友當中.
有時事情的詳細.
你不能說下去.
但起碼你可以說出題目.
因為記住.
雖然我不知道詳情.
但上帝知道.
對不起.
我不是在說你們.
我認識很多人.
有什麼代禱事.
其實就是要說.
我們又不是醫生.
如果是醫生.
就很詳細.
什麼什麼.
不用.
說完我們都不知道.
就是說出題目已經足夠.
如果你覺得那件事難以啟齒.
弟兄.
姐妹.
你為我.
我現在真的很傷心.
為我這件事祈禱.
都可以的.
不需要對方.
或你自己.
要說得很細節.
因為上帝知道什麼.
我現在需要的是你為我禱告.
所以有時候.
遠方的宣教士.
可能忙到沒辦法寫宣教信.
代禱信的時候.
你就為這個人禱告.
特別在香港的傳道人.

$^{801}$牧師.
不會像我們宣教士.
我們宣教士都會有一封代禱信.
沒有什麼代禱信.
你就為牧師.
為傳道.
為事奉人員.
到福音樂隊.
他們下一次再唱一些歌來勉勵我們.
剛才唱的歌.
就是每天我們行事為人.
是不是跟神一樣去走呢?.
我們香港回歸到中國.
我們的心.
是不是在上帝的裡面呢?.
若果不是.
我們就要回歸.
上帝同樣可以透過這些歌詞.
讓我們走近神的事.
求神幫助我們.
讓我們真的禱告.
和上帝建立讀經.
祈禱,靈修.
這些一切.
和回來敬拜神.
是一個和神保持關係密切的一件事.
在這件事裡面.
你看到兩件事.
兩個皇帝.
他們對於但以利這個俘虜.
是另眼相看.
但以利可能是很英俊,很厲害.
但他如果不是特別.
不是靈明裡面那麼漂亮的時候.
是不會有這件事出現.
這裡撇開了.
這裡仍然是但以利.
不過我這次沒有說.
很想你回去看第五章.
有一次你可能會記得但以利這個故事.

$^{841}$千人宴.
不是尼布甲尼撒.
尼布甲尼撒的孫子.
他那時候就上位了.
他就搞了個千人宴.
但沒有但以利.
我們都試過.
有時候功就由你做.
然後慶功宴就沒有你去慶.
於是你就很生氣.
可能會是這樣的處境.
可能你和我都經歷過.
有些事我做就做.
但沒有人提你的名字.
但記住.
這些事情你為上帝去做的時候.
絕不徒然.
在千人宴裡面.
沒有人記得但以利.
突然牆上出現了怪的手指.
在寫著字.
出現了手指就很奇怪.
但寫的字其實不奇怪.
現在學者學的文都讀得到.
但當時是尼布甲尼撒的太太.
現在的皇太后.
因為她很害怕.
害怕到在聖經裡面.
會看到很像卡通片或畫劇般.
很好看.
皇帝膝蓋都發抖了.
很害怕.
沒有人可以解答她的疑難.
那個時候千人宴裡面沒有但以利.
是皇太后上去.
你可以找但以利.
其中讚賞了但以利有六大點.
你回去看.
五章十二節全部在但以利書裡面.
第一點是美好的靈性.

$^{881}$第一點讚賞但以利.
就是有美好的靈性.
美好的靈性.
你和我有沒有?.
未夠高分.
好像我的名字高分.
也可以去求.
總之有一分得一分.
美好的靈性千萬不要低分.
要向著高分而行.
向著高處而行.
我們一生人.
要一直祈求美好的靈性.
求主幫助我們.
這個同樣是但以利.
這個特點讓巴比倫王.
和波斯王封他.
特別是波斯王去到大利烏王.
獅子坑那件事之後.
去頒佈全國.
頒佈全國的詳細你當然要看聖經.
當然他講到大家應該要在神的面前.
戰驚恐懼.
和合本裡面講這四個字.
你要在神.
全國的人民.
他不是叫全國人民只拜耶和華神.
但是你要在耶和華神面前.
戰驚恐懼.
我們現在講的是否敬畏神.
比敬畏神更高.
你要戰驚恐懼.
今天我們和神是否太親近.
以致我們不怕神.
你缺席一天不上班一天.
你都在搞醫生信請假.
很多程序.
今天不上班就不上教會.
不上團契就不上團契.
你明不明.

$^{921}$我的意思是.
我們所信的神.
應該是值得我們戰驚恐懼.
敬畏神.
當然上帝是有愛去愛我們.
但是神仍然是敬畏的神.
我們會看到.
但以利是很謙卑的祈禱.
其實落獅子坑.
已經是八十多歲的老齡.
落獅子坑沒有跌傷骨頭.
都是神蹟.
還要沒有被獅子咬.
這段經文.
為何神仍然維持每天三次禱告.
祈禱甚麼呢.
祂是希望以色列國可以復國.
因為上帝給了很多異象.
要經過多少就復興.
祂希望耶路撒冷復興.
今天你和我都很想洪恩堂復興.
你和我又會擺上多少禱告的力量.
在這件事裡.
求主給我們彼此的勉勵.
在但以利的禱告裡.
祂很明白是我們以色列人犯罪.
以致得罪神 以致亡國.
雖然在但以利的一生裡.
我們看不到以色列何時復國.
但以色列終於有復國.
是多少年了 大家記不記得.
1948年的時候.
因為我們的但以利.
仍然很清楚上帝是守施慈愛的神.
有時我們禱告.
可能我們的對眼看不到.
上星期我們請了一個弟兄.
他玩遊戲玩到很瘋狂.
媽媽在基督徒有為他祈禱.
但媽媽身為癌症過世.

$^{961}$他過世時兒子還未改變.
但在他媽媽過世後.
兒子的形容突然醒了.
覺得不能沉迷在玩遊戲機的生活.
於是在疫症期間.
因為他媽媽也是在2019年疫症期間過世.
仍然戴著口罩.
去探訪那些很有報道的心智.
這份報道的心智.
是他自己買口罩或應用品.
給弱勢的一群人.
於是他個人做的.
上帝感動他做.
結果由一個玩遊戲機的人.
勒著教會的牧師.
照樣不止弟兄做.
整間教會一起.
在疫症期間.
飯盒等需要用品.
來幫助社區.
所以生命可以彼此勒著.
你和我都想人多些在這裡.
除了準備好旁邊的位置不要放袋子.
我們的生命準備好了沒有.
求神幫助我們每一個.
上星期珍姑娘說的題目.
還記不記得.
建立堅不可摧的生命根基.
上星期一聽我已經跟珍姑娘說.
我下星期就說彈以利.
彈以利就是我心目中建立了一個.
堅不可摧的生命根基.
願彈以利的生命勒著你和我.
彈以利不是一帆風順.
你和我都不是一帆風順.
生命有挫折.
生命有不明白的時刻.
生命有被委屈的時候.
同樣生命也有喜樂開心的時刻.
但無論是怎樣的情況下.

$^{1001}$求神祝福你和我一樣.
可以透過我們的生命.
來勒著別人的生命.
更加去勒著一些未信主的人的生命.
因為這樣的時候.
他們才得到上帝的福氣.
上星期也說了.
有一份閱歷.
我經常覺得我們說是基督徒.
但連本基督聖經.
聖經我們都沒有看.
真的很過份.
每天.
不是每天.
逢星期三,六.
你走來走去搭輕鐵的時候.
你會看到輕鐵旁邊有馬會.
馬會附近.
你看看那些人.
拿著東西.
拿著筆.
沒有坐椅子.
在哪裡.
蹲也好.
怎樣也好.
在那裡做功課.
刨經刨得多厲害.
你看看那些布攤.
現在在賣的是甚麼經.
馬經.
馬經的報紙原來有很多份還未倒閉.
這就很代表那些馬迷.
他們很著重那本經.
我們呢.
你是基督徒.
跟著神的人.
你有沒有去看呢.
我覺得很好.
這本書上星期介紹的時候.
可以說這本月曆.

$^{1041}$就是洪恩堂十一周年的心血結晶.
好像一個嬰兒出世.
你會不會跟著這些細節去看約書亞呢.
記得有一次說到父親節.
你和你一家都要事奉神.
在約書亞記出版.
究竟約書亞是甚麼時候說這句話呢.
這次好了.
這次你看約書亞記.
你就知道了.
千萬不要怕難明的事.
不知道就做甚麼.
有些人可能不太識字.
我今天預備了一份禮物.
先不說這份禮物.
因為我們很多都是用手機.
手機裡面有很多APP.
有時我們看東西.
現在我都進入了老花.
看東西都要脫下眼鏡才看到.
都很麻煩.
但是聽.
如果你看有問題.
聖經有大字版.
有小字版.
有大字版.
你可以看大字版.
如果你連看聖經都有問題.
那就聽吧.
聽電話裡面有很多APP.
它是讀出來的.
讀出來有好處.
約書亞記一張你可以讀幾次.
一邊聽一邊聽.
例如我們姐妹們.
我們在家裡可能要做一些.
不用腦的事.
例如掃地.
洗碗.
這些事.

$^{1081}$你可以一邊聽一邊.
是很好的.
我就是經常這樣做.
因為我平時要照顧父母.
很多事要做.
一邊聽對我的益處很大.
重複再聽.
聽不到的再重複去聽.
這個在聖經裡面.
在你的電話APP裡面可以做到.
甚至乎你還可以買以前的產品.
要聽約書亞記.
就要聽.
這些可能你應用電話不方便.
現在誰先舉手.
誰可以拿走.
好呀.
你過來拿吧.
你真是.
真是要聽.
謝謝.
在這裡我真的很希望.
如果你今天還沒有這本書.
你其實可以偷步.
你真的覺得我怕我比別人慢.
你先偷步.
現在先看剔了.
靈命的事不怕你偷步.
最怕你停步.
一停步就不能勒著別人.
所以真是求神祝福我們.
讓但以你的生命.
真的可以勒著你和我.
我們要記住.
無論你的處境是怎樣也好.
上帝依然是那位.
勁默默的神.
還有.
縱然去到一個處境.
你可能要有些大哥大的精神.

$^{1121}$都求神祝福你.
求神不要去到那個地步.
我過去都很害怕.
以前都很喜歡看宣教士.
或者有些國家.
例如密室這本書.
我會看.
總之是說.
打仗了.
德國人怎樣收起以色列人.
我就在想.
如果那些人民高淑芬.
蟑螂就在你前面.
如果你信神.
蟑螂就會黏在你的臉上.
或者推到你的房間.
全部都是蟑螂.
蟑螂的泰文名字很好聽.
叫什麼.
買靚衫.
但我其實很怕買靚衫.
你知不知道.
我其實很怕這些東西.
但上帝讓我去泰國.
泰國都不知道多少買靚衫.
真的很害怕.
有一次回去教會.
鐵閘.
買靚衫好像開派對的閘.
比這個位置還要幼.
在穿梭.
我站在門.
應該要準備場地.
然後我們要教聖經.
我站在門.
不敢進去.
然後兩個大學生來到.
阿贊.
阿贊即是傳道人.
你做什麼.

$^{1161}$買靚衫.
他們的書.
輕描淡寫.
就幫我開門.
所以如果到一天.
真的有些東西恐嚇你.
要知道你不信主.
可能我是第一個.
我經常都很害怕.
因為買靚衫就搞定你.
但我求神憐憫我.
如果有一天是這樣.
求主真的赦免我的罪.
憐憫我.
我真的太膽小.
真的面對不了.
所以求神幫助我們.
縱然那一刻.
大哥大的精神也好.
求神給我們.
我們的生命可以蠟著.
其他人的生命.
讓這裡現在還有空置的座位.
成為別人的祝福.
讓祂一起來.
和我們一起敬拜神.
願神祝福每一位.
拜拜.
謝謝大家.
\newpage



\section{希伯來書 10:19-25}
\label{sec:Hx9S3W_wWOQ}
\textbf{2022.11.27 《讓我們定睛十字架》希伯來書 10\_19-25(和修版) 曾裔貴牧師}
\newline
\newline
連結: \href{https://youtube.com/watch?v=Hx9S3W_wWOQ}{\texttt{https://youtube.com/watch?v=Hx9S3W\_wWOQ}} ~~~~ 語音日期: 2022-11-29
\newline
\newline
\hyperref[sec:bh_7j2paksc]{\small{< < < PREV SERMON < < <}}
~
\hyperref[sec:index_chronic]{\small{[返順時目]}}
~
\hyperref[sec:index_scriptual]{\small{[返順卷目]}}
~
\hyperref[sec:qxsO_vpno54]{\small{> > > NEXT SERMON > > >}}
\newline
\newline
希伯來書 10:19-25
\newline
\begin{longtable}{cl}
\hline
\hline
章節 & 經文 (和合本修訂版)\\
\hline
10:19 & \begin{tabularx}{0.7\textwidth}{X} 所以，弟兄們，既然我們靠著耶穌的血得以坦然進入至聖所， \end{tabularx} \\ \\ \relax
10:20 & \begin{tabularx}{0.7\textwidth}{X} 是藉著他給我們開了一條又新又活的路，從幔子經過，這幔子就是他的身體。 \end{tabularx} \\ \\ \relax
10:21 & \begin{tabularx}{0.7\textwidth}{X} 既然我們有一位偉大祭司治理神的家， \end{tabularx} \\ \\ \relax
10:22 & \begin{tabularx}{0.7\textwidth}{X} 那麼，我們該用誠心和充足的信心，同已蒙潔淨、無虧的良心，和清水洗淨了的身體來親近神。 \end{tabularx} \\ \\ \relax
10:23 & \begin{tabularx}{0.7\textwidth}{X} 我們要堅守所宣認的指望，毫不動搖，因為應許我們的那位是信實的。 \end{tabularx} \\ \\ \relax
10:24 & \begin{tabularx}{0.7\textwidth}{X} 我們要彼此相顧，激發愛心，勉勵行善； \end{tabularx} \\ \\ \relax
10:25 & \begin{tabularx}{0.7\textwidth}{X} 不可停止聚會，好像那些停止慣了的人，倒要彼此勸勉，既然知道那日子臨近，就更當如此。 \end{tabularx} \\ \\
[1ex]
\hline
\hline
\end{longtable}
$^{1}$弟兄姊妹主女們平安.
做一次代表母會.
特別都是感恩節.
在十一月份.
都能夠帶著一份平安.
希伯來文叫SALOM.
來將這份祝福.
帶給紅恩的弟兄姊妹.
讓我們在講道之前.
一起有祈禱.
多謝主.
我們要打開聖經.
這一段聖經.
是邀請我們.
去到天上的致聖所.
其實就是去到神的面前.
而且是要有充足的信心.
再無懼怕.
這樣來與神親近.
剛才幾首的詩歌.
其實都已經講了.
十字架是唯一的原因.
我們可以藉著耶穌.
已經流出寶血的.
這一位主耶穌.
來在十架上的犧牲.
將我們帶到與神和好.
這個很重要的生命關係.
幫助我們.
我們中國人.
香港人.
未必對舊約的.
那些聖殿的條例.
會有一些很深的感受,體會.
因為我們未試過.
將牛羊帶去聖殿獻祭.
我們試過可能是爸爸媽媽.
來祭祖.
或者是在一些大的節日.
來燒香.

$^{41}$我們比較少.
好像猶太人那樣.
這段聖經.
真的幫助我們去想想.
我們怎樣可以.
很有平安的.
甚至乎應該很大膽的.
去到天上的致聖所.
今早我們仰望.
等候主的聖靈在聚會當中.
來激發我們大家的信心.
這樣的感恩禱告.
奉耶穌基督的名.
阿們.
帶了一本.
是新約的聖經.
但是我完全看不明白.
任何一個字.
因為這本聖經是.
希伯來文的新約聖經.
首先有兩個特別.
新約聖經是用希臘文寫的.
原文是希臘文寫的.
是不會用希伯來文寫的.
另外一個特別就是.
猶太人是不會接受新約的.
猶太人只有舊約的聖經.
由《創世記》到《馬拉基書》.
猶太人是不相信新約說的耶穌.
就是他們的彌賽亞.
這本新約聖經.
我們很有關聯.
因為是錦宣的一位弟兄.
在2019年.
來跟以色列的聖經公會合作.
印制.
印了12萬還是14萬本.
我是得到其中一本.
因為我們也有為這個時宮祝福去禱告.
這本聖經我完全一個字看不明白.

$^{81}$因為是用原文的希伯來文來寫.
這段聖經其實是很想給當時修信的讀者.
其中很大的部分是猶太的信徒.
要有一個很強的信心.
要擺脫舊有的舊約憲制的觀念和思想.
要將整個信心放在耶穌的裡面.
為什麼我會突然帶這本聖經來呢?.
昨天下午.
在辦公室一路預備.
想一想.
其實有沒有一些真正的猶太人.
他們有這些情意結.
現代的猶太人.
有沒有這些困難.
他們不能夠接納.
這一位的耶穌.
我們沒有什麼問題.
我們信了耶穌.
我們是外邦人.
我們信耶穌.
一信就得救.
沒有這些包袱.
但是猶太人有沒有呢?.
一路去想.
預備了很久.
腦袋也開始有點累了.
由我的辦公室走回去禮堂.
大家如果去過錦宣就知道.
有一個廣場和露天的茶座.
看到兩位弟兄在聊天.
其中一位就是支持這個事工的弟兄.
和另一位弟兄在聊天.
走過去打招呼.
和他說一聲.
很難得有機會星期六下午見到他們兩個.
聊聊天.
接著我就走進辦公室簽文件.
簽簽下.
弟兄就進來叫我.
有空嗎?.

$^{121}$我就立刻放下文件.
出去看看發生什麼事.
原來這一位弟兄.
馮弟兄.
同梗爾息烈一位傳道人.
來視像對話.
現場的視像對話.
他們是當時的早上.
我們是下午四點鐘左右.
弟兄拿起一本聖經.
就是這一本黑皮聖經.
他說他用英文.
相當簡單的英文.
他說他自小就在拉比學校成長.
他的父親是一個拉比.
後來輾轉有機會離開了以色列.
有人就送了一本希伯來文的新約聖經給他.
就是這一本.
他來到這裡讀.
一邊讀一邊有很多的掙扎.
為什麼這本聖經所說的耶穌.
和我小時候在拉比學校所說的耶穌是不同的呢?.
為什麼新約這本聖經會這樣說這位主耶穌.
是那個彌賽亞呢?.
他有很多的掙扎.
其實他很小的時候就已經對耶穌.
因為他們有讀他們的歷史.
對基督教所說的耶穌.
他是非常的憎恨.
因為不可能一個人說自己是神的兒子.
他自小就在拉比學校成長.
對舊約都很熟悉.
這一位中年的傳道者.
輾轉讀完這本書.
信了耶穌.
他現在在北加利尼一個地方.
那裡是唯一一間猶太人的基督教教會.
所以他要面對很多難處.
但是他的教會現在有64人.
很快就增長到144人.

$^{161}$所以很希望能夠有支持.
他們建造一間新的地方.
昨天他們就在對話.
準備支持他們數萬元的美金.
來建造一個地方.
就是這一本聖經.
新約的聖經.
很奇妙地.
我在想有什麼例子.
可以幫助我解讀這段聖經的時候.
多一點點的資料給大家.
居然間.
一個這樣的寒圈打個招呼.
就有一個以色列的傳道人.
自小就在拉比學校成長.
來跟我們的弟兄.
香港的弟兄在對話.
講他的基督教見證.
各位.
這段聖經.
其實有三個邀請.
就從21節開始.
或者應該22節.
居然間我們心中的天良.
就是我們的良心的虧欠.
已經洗乾淨了.
就應該怎樣呢?.
存著一個誠實的心.
和充足的信心.
去到神的面前.
剛才我也說了.
整本書是跟一些有信仰.
希伯來的傳統信仰猶太教的背景.
的弟兄姊妹.
來講.
勉勵他們.
不可以再透過物質的牛羊的血.
或者地上的大祭司.
這個體系.
來親近神.

$^{201}$那應該怎樣親近呢?.
你自己用信心直接親近神.
這是很大逆不道的.
不要祭司.
不要獻祭.
自己的心就是去到神面前.
但是聖經真的這樣講.
裡面包含的是接納.
裡面包含的是一個寬恕.
包含的是那個信徒.
要排除他過去的舊有觀念.
以信心.
那個信心是建立在主耶穌的犧牲.
就是說你對主耶穌有多少的認識.
你就會得到神對你有多少的接納寬恕.
你想想一開始的時候.
剛才弟兄已經讀了.
我們是可以坦然.
去到致聖所.
你想想這個是真真正正的禁地.
猶太人由出埃及的神.
設立了的回眸之後.
致聖所是沒有人可以進去的.
大祭司都是一年去一次.
就是他們的贖罪日.
七月初十.
才可以去一次.
去到才要將燒滿香爐的香煙.
放在裡面.
讓裡面充滿煙.
什麼都看不到的時候.
你不知道哪個方向.
進去之後.
總之你不要管了.
因為看不到.
就彈血.
彈七次血.
彈完就很快走.
不可以不走.
不走.

$^{241}$出不來.
可能就是被神擊殺死.
這個就是致聖所.
真真正正的禁地.
猶太人連想都不敢想.
他們可以坦然地進入致聖所.
但是作者現在在這裡.
告訴我們.
有一個不是一次一天的開放日.
是整個生命開放給每一個人.
你只要靠著.
這一位主耶穌的血.
就可以坦然進入致聖所.
所以去到二十二節.
就告訴我們.
我們就應該要帶著這一份誠實.
和充滿了信心.
充足的信心.
來到神的面前.
這是第一方面.
作者的勉勵.
讓我們帶著這一份信心.
這一份蒙月臘.
神對我們的寬恕的接納.
嘗試用一個例子.
可能勾起我們大家彼此的思考.
有一個律師.
他姓林.
林正江律師.
他說一個故事.
真人真事.
他說他有一個客戶.
第一次見他的時候.
十六歲.
因為在學校偷東西.
被人抓了.
他作為一個辯護律師.
自然一個少年人.
偶然間犯錯.
固然應該要幫他.

$^{281}$幫他辯護.
過了幾年.
又見到這個年輕人.
這次不是偷東西這麼簡單.
是鵝騙.
這個律師已經有點不開心.
為什麼你又這樣犯錯.
然後.
來到去.
算了.
仍然幫他.
希望他不要誤入歧途.
再繼續泥足深陷.
到第三次的時候.
律師.
當著他的面.
來罵他.
這次是搶劫.
已經又過了幾年.
律師說.
我不一定再幫你辯護.
因為你一而再這樣犯罪.
這個年輕人.
突然間很震驚.
為什麼律師這麼重的說話.
我請你幫我辯護.
你應該盡力幫我辯護.
這樣來到去.
義正嚴詞地罵他.
他突然間愕然之後.
他再想一想.
然後.
跟律師說.
跟林律師說.
你說得對.
我是應該.
注定要在監獄裡面.
向這位林律師道歉.
鞠一個躬.
然後離開.

$^{321}$過了一段時間之後.
這個年輕人.
差不多出監了.
寫信給林律師.
跟他說.
我差不多走了.
但是在監獄裡面.
輔導的人員.
有一次播了一首歌給我們聽.
是記載一件真人真事.
他說在美國.
同樣有監獄裡面的一個人.
犯了事.
坐了很久的牢.
他差不多釋放.
就寫信給自己的太太.
就跟太太說.
如果你肯接納我.
寬恕我.
過去的缺失.
我現在出監了.
我會回來.
你在村口那棵大樹.
你綁一條絲巾.
如果我見到這條絲巾.
證明你寬恕我.
閱立我.
他離開監獄.
就坐巴士回到他的村.
他沒有勇氣去望大樹.
他跟巴士裡面的乘客.
來講他自己的故事.
他說希望太太接納他.
就吩咐和乘客.
巴士司機幫忙.
望望差不多到達的時候.
那棵大樹有沒有一條絲巾.
他沒有勇氣抬頭去望.
過了一段時間之後.
車差不多到達那個站.

$^{361}$接著就很多人拍手掌.
因為那棵大樹裡面.
綁了很多的絲巾.
這個年輕人就寫信給林律師.
我差不多要出監了.
林律師你知不知道.
為什麼我會成為一個積犯.
我第一次犯事.
接著不是很嚴重的那個懲罰.
我回到學校.
學校的教職員和同學.
再看我是一個異類.
好像我的額頭寫著.
不是好人.
有學校誰不見了.
學校的職員訓導主任.
第一時間就去找我的袋.
找到沒有都不說道歉.
輾轉輾轉.
我好像在學校裡面是一個異類.
接著離開學校之後.
去找工作.
又會被人歧視.
更嚴重的就是.
原來在自己的社區.
有失事的案.
警察第一時間.
在這附近就會上來我的家.
來詢問我有沒有偷過東西.
所以家人看我是一個負累.
是對家門的明星的損毀.
所以我成為一個另類的人.
但更加諷刺的就是.
他和他的朋友.
大家都是犯事的.
反而感覺舒服.
所以慢慢慢慢.
他就是這樣來到去.
在沒有人去寬恕,接納.
和認同下.

$^{401}$一步一步地成為了積犯.
最後他說林律師.
我都很想離開這個監獄.
我見到有黃師大.
我見到有人接納我.
親愛的弟兄姊妹.
教會最寶貴的是什麼呢?.
我由2017年到今天.
在錦宣來承擔牧羊的職責.
最令我痛.
最令我睡得不好的.
是很多關係上.
因為一些很小的事.
或者可能有人看為很大的事.
就決.
決裂.
誰贏.
誰輸.
決裂.
決裂之後.
是真的大家不再用一種愛相交.
其實性和負之間.
是不是應該更加.
教會裡面應該有一個圓滿呢?.
是應該有一個彼此的寬恕接納呢?.
你看看作者如何勉勵當時希伯來的信徒.
他說既因為耶穌的血.
即是代價已經付出了.
不需要人和人之間爭誰勝誰負.
應該就是看到神寬恕每一個罪人.
其實我們每個人都是不是好人.
不過我們沒有人看到.
神看到我們大家的心靈.
所以我們就應該.
作者說我們要傳一個誠實的心靈.
和充足的信心.
去到神的面前.
即是你有多少對主耶穌的認識.
你就有多大的信心去到神的面前.
第二方面就說.

$^{441}$第三面就說23節要堅守我們所承認的指望.
不自搖動.
因為不是我們有信心.
是那一位.
應許我們的.
他是信實.
大家看完日本對德國那場球賽.
日本反勝2比1.
接著網上馬上怎樣呢?.
有很多道歉的聲明.
很多人帶頭說.
我們要向日本的教練森保一道歉.
因為未踢之前.
每個人都憎他死.
帶領那隊日本隊很差的戰績.
大家都沒有眼睛看.
就說這個人其實應該換一個教練.
不是換那隊球.
所以一贏了球之後.
尤其是下半場.
差不多完結.
還在輸1比0的時候.
總教練森保一.
毅然地用盡換人的限額.
四個前鋒一起出去.
拼命地進球.
贊成.
終於反勝2比1.
成為全球足球界的大事.
接著就向森保一認錯.
不再堅持對他的負面.
當然這是開玩笑的.
這樣去宣認.
但是這告訴我們.
無論在什麼情況下.
都要堅守這個指望.
就是說我們對永恆要有指望.
我嘗試用一個例子.
證明神的存在.
是我們想像不到的奇妙.

$^{481}$讓我們沒辦法相信.
這個世界的獨一真神.
他的存在是如此真實.
有一位醫生.
他姓張叫張煒.
他在北角浸信會洗禮.
這個年輕人.
其實是一個遼寧人.
在遼寧的醫學院.
拿了碩士學位.
在骨科很厲害的一個專家.
後來國內政府派他去外面.
一些保外的緩醫.
去幫一些弱勢的社群.
特別是非洲.
他也做得很出色.
這個年輕的醫生.
有一次休假.
就去跳水.
去游泳.
一跳.
原來下面有一些石.
弄到脊椎.
導致中身.
由他的肩膊以下.
完全不能夠再有任何的活力.
類似MIRROR演唱會的信賢.
李成能牧師的兒子.
腦部完全沒事.
但是肩膊以下完全不能夠再有活力.
生不如死.
一路康復.
一路復健.
越來越身體有意識.
但意識就是坐輪椅.
沒有了.
沒有什麼可以做了.
但腦部是很敏銳.
完全沒事.
跟媽媽說.

$^{521}$媽媽.
我長期睡病床.
那些浴瘡.
令我很辛苦.
很煎熬.
不如你去到懸崖.
推我下山就算了.
不要再留我在世界上.
這樣來受苦了.
一個二十多歲的年輕人.
在國內回到遼寧養傷.
輾轉有一天在機場轉機.
有一個日本的教練.
走過看到這個年輕人.
推著輪椅.
內心有感動.
這個日本人叫做恥口令子.
很感動.
接著就在他的袋子裡拿了一本書.
其實大家不認識的.
給他一本書.
叫做《輪椅上的畫家》.
鍾離也是美國一個.
很多年前的殘疾人.
也是相同的遭遇.
下半身完全不能夠有任何的動作.
用他的嘴巴來畫畫.
送給他.
你拿著一本書來贈慶吧.
我已經這麼慘了.
你還給我一本這樣的書來看.
叫媽媽放在抽屜裡.
不要讓他看到.
但是時間實在太長了.
全天坐在輪椅上.
不坐在輪椅上.
就躺在床上.
可以做什麼呢?.
記住腦子是很清醒的.
很精神的.

$^{561}$叫媽媽拿出一本書來看吧.
因為他根本動不了.
媽媽還要揭開給他看.
越看越特別.
為什麼他跟我一樣的遭遇.
為什麼他又散發出一種生命力.
一種這麼大的信心呢?.
越看越奇妙.
用了45天.
將這本書翻譯成中文.
是由媽媽幫他做.
他講,媽媽幫他做.
然後在這個過程裡.
兩母子信了耶穌.
然後神感動他.
要為殘疾人傳福音.
這位張煦醫生.
曾經在中央電視台.
全國的聯播.
講他為什麼這麼有生命力.
可以組織一個NGO.
來幫助很多的殘疾人.
他自己在朕黎的見證.
講了一番話.
在北角浸信會.
來接受洗禮.
香港的北角浸信會.
他說,不管環境多麼的黑暗可怕.
即使在沒有神思想分途的環境裡.
我們的創造者.
並沒有雕下我們不顧.
是他自己在朕黎的見證所講.
親愛的弟兄姊妹.
你的生活.
有沒有因為這幾年的疫情.
帶給你有很多的困難呢?.
心靈的.
生活的.
我當然改變不了你的即時的困難.
因為我也面對我的困難.

$^{601}$你的困難.
每個人都有我們自己的困難.
但是不要忘記.
我們有共同的指望.
不致搖動.
重要的是.
因為應許我們這一位.
他是信實的.
在沒有神環境分途的地方.
一本已經是幾十年前的書.
可以改變了一個活生生的年輕人.
讓他知道原來神沒有撇下任何一個人.
他是不看顧.
願主祝福大家.
也盼望洪恩嘉繼續茁壯成長.
有更多弟兄姊妹.
在這個社區成為那條絲帶.
成為別人的信心的一個榜樣.
我們一起禱告.
多謝主.
但願你自己的說話.
成為我們的支持.
我們的力量.
這樣的感恩禱告.
奉耶穌基督的名.
阿們.
\newpage



\section{哥林多前書 3:10-15}
\label{sec:qxsO_vpno54}
\textbf{2022.12.04 《建立蒙主喜悅的生命工程》哥林多前書 3\_10-15(和修版) 曾敏姑娘}
\newline
\newline
連結: \href{https://youtube.com/watch?v=qxsO_vpno54}{\texttt{https://youtube.com/watch?v=qxsO\_vpno54}} ~~~~ 語音日期: 2022-12-08
\newline
\newline
\hyperref[sec:Hx9S3W_wWOQ]{\small{< < < PREV SERMON < < <}}
~
\hyperref[sec:index_chronic]{\small{[返順時目]}}
~
\hyperref[sec:index_scriptual]{\small{[返順卷目]}}
~
\hyperref[sec:OuveUNRd_9o]{\small{> > > NEXT SERMON > > >}}
\newline
\newline
哥林多前書 3:10-15
\newline
\begin{longtable}{cl}
\hline
\hline
章節 & 經文 (和合本修訂版)\\
\hline
3:10 & \begin{tabularx}{0.7\textwidth}{X} 我照神所給我的恩典，好像一個聰明的工頭，立好了根基，別人在上面建造；只是各人要謹慎怎樣在上面建造。 \end{tabularx} \\ \\ \relax
3:11 & \begin{tabularx}{0.7\textwidth}{X} 因為，那已經立好的根基就是耶穌基督，此外沒有人能立別的根基。 \end{tabularx} \\ \\ \relax
3:12 & \begin{tabularx}{0.7\textwidth}{X} 若有人用金銀、寶石，草木、禾秸，在這根基上建造， \end{tabularx} \\ \\ \relax
3:13 & \begin{tabularx}{0.7\textwidth}{X} 各人的工程必將顯露，因為那日子要將它顯明，有火把它暴露出來，這火要試煉各人的工程怎樣。 \end{tabularx} \\ \\ \relax
3:14 & \begin{tabularx}{0.7\textwidth}{X} 人在那根基上所建造的工程若能保得住，他將要得賞賜。 \end{tabularx} \\ \\ \relax
3:15 & \begin{tabularx}{0.7\textwidth}{X} 人的工程若被燒了，他將損失，雖然他自己將得救，卻要像從火裡經過一樣。 \end{tabularx} \\ \\
[1ex]
\hline
\hline
\end{longtable}
$^{1}$大英姐妹平安.
令我們有很多思考.
我們再一次一同禱告.
真的很愛我們.
給我們有數不盡的福氣.
今天你更加要藉著你的話語.
去教導我們.
願意你開我們的眼睛.
開我們的耳朵.
開我們的心門.
讓我們完全明白你的心意.
願意聖靈在我們當中去工作.
感動我們.
聽到上帝的聲音的時候.
我們就去回應.
求你帶領以下的時間.
奉耶穌基督的名字祈禱.
阿們.
創風的分享令我有很多的思考.
一會兒我相信在講道當中.
都很想在那些位置.
我們彼此去提醒.
今天我們講的是一個.
叫做建立蒙主喜悅的生命工程.
今天已經進入到12月了.
真的很快.
2022年還有幾個星期就結束了.
不知道大家回想的時候.
這一年我們過得如何.
或者接下來2023年的時候.
大家又打算怎樣呢?.
麻煩你去我的powerpoint那裡.
這裡說的是建立蒙主喜悅的生命工程.
是指什麼樣的生命工程呢?.
其實這裡說的是教會裡.
整個群體的.
整個群體的生命工程.
是很值得思考的.
特別是去到我們一年最後的這個月份.
值得我們想一想.

$^{41}$有些事情我們去思考之後.
在新一年在未來的日子.
都值得我們在上帝面前好好禱告.
這個生命工程是指教會.
屬上帝群體的一個生命工程.
經文的背景是怎樣呢?.
我們稍稍看看.
這個字有一點小.
經文背景是哥林多前書三章五到七節.
其實這裡說的是有兩個人物.
分別是阿波羅和保羅.
保羅我們比較熟悉的.
他是一個很偉大的報道家.
也是一個非常出色的宣教士.
一個很有愛的.
很有牧者心腸的一位牧者.
上帝所重用的僕人.
他每到一個地方就很努力的傳福音.
帶領人去信主.
去建立教會.
他就周遊列國.
因為他是一個宣教士.
最主要的對象是對外邦人.
在這個過程中.
你會發覺.
事奉裡總不是一帆風順.
除了傳福音,事奉,牧養之外.
原來在牧養的過程中面對一些挑戰.
例如人們很喜歡和別人比較.
很喜歡比較.
在哥林多教會做什麼比較呢?.
我們也有一個阿波羅.
這個阿波羅的口才真的厲害.
他講話真的很大聲.
我喜歡聽阿波羅講話.
你這個保羅是怎樣一回事呢?.
我們來比較一下.
阿波羅是怎樣的.
你現在跟誰?.
你跟阿波羅還是保羅?.

$^{81}$我跟阿波羅.
我跟阿波羅派.
阿波羅派他出來.
現在另一個就說.
不是,我跟保羅的.
我們是保羅派的.
和你們阿波羅是分開的.
於是教會講著講著就變成了分裂.
分黨分派.
彼此比較.
很大的傷害.
其實我相信.
在阿波羅和保羅當中.
根本沒有這種競爭.
保羅很清晰.
今天他和阿波羅的關係.
是彼此配搭.
大家有不同的角色.
使命.
一個傳福音報道.
一個就是栽培.
大家多好.
不存在這種.
有我就沒有你.
有你就沒有我.
講清楚.
在哪個陣營.
但是教會很容易是這樣.
我之前這個傳道人.
我支持那個傳道人.
你們這些走在一邊.
我們什麼派.
蘋果派.
西瓜派.
我們互相不能走在一起.
一講起蘋果.
我心裡就湧流著熱血.
我很支持我的牧者.
那個又說.
想起我們的西瓜派.

$^{121}$真的十分之好.
我好好支持他.
你們這些蘋果派很膚淺.
與事者是非常愚昧.
這個不是神的心意.
神一定會興起不同的牧者.
各有不同的長處特色.
一起要配搭.
服侍建立神的家.
拓展上帝的國度.
這個本身就是神的心意.
但是人的情況.
就是令到他變成了鬥爭.
這是哥林多教會出現的情況.
保羅出來講的時候.
其實大家根本有不同的使命.
有不同的使命.
我們做的事有不同的.
彼此的配搭.
很明顯保羅是做災眾的人.
他傳福音.
殺種.
帶人歸向耶穌基督.
阿波羅囂冠.
要去栽培.
去陪靈.
但你們要記住.
最終是上帝令到這些生命成長.
是上帝自己的工作.
所以背景就是這樣.
而教會就是田地.
唯有神讓生命成長.
所以最重要的是.
神是教會真正的主宰.
真正影響教會的是神.
教會是屬於神的.
不是任何一個牧者.
就算是任何一個人出來.
哇,這個人講得很好.
我真的要擁戴他.

$^{161}$我跟著他,簡直是明星.
我記得以前有一次.
去參加一個陪靈的大會.
很厲害的.
那次的陪靈大會.
是以前遠東有一個舊同事邀請我去參加.
那個牧者很偉大.
真的很偉大.
他旗下的牧羊羊有幾十萬.
他一進去大球場的時候.
哇,很厲害.
前呼後擁.
那些人圍著他.
我感受到這是明星的一個講員.
真的很厲害.
但是我想說.
我相信這個不是上帝的心意.
上帝有他大大重用的僕人.
但是真正的明星.
真正的,那個很閃亮亮的.
是我們的上帝.
一定不是那個人.
那個人是要榮耀上帝.
彰顯上帝的奇妙.
所以今天教會無論如何.
真正的主宰是我們的耶穌基督.
是我們的上帝.
好了,這裡很重要.
趁著講這段經文的時候.
我想說教會是屬於神的.
每一樣東西都和耶穌基督有關.
這段經文是歌羅西書一章17到18節.
這樣說.
祂在萬有之先.
萬有也靠祂而存在.
祂是身體,就是教會的頭.
祂是原始.
是從死人中復活的守生者.
好讓祂在萬有中居首位.
這個在萬有中居首位的耶穌基督.

$^{201}$也是教會的頭.
耶穌基督是教會的頭.
教會的主.
是很榮耀的.
擁有所有的權柄和尊貴.
好了,在這個過程裡.
你看到很特別.
教會的頭就是耶穌基督.
教會的根基呢?.
這段經文雖然字體小.
但應該可以看得到.
保羅說:我照神所給我的恩典.
是上帝給他的呼召.
上帝呼召他成為一個宣教士.
一個傳福音的使者.
他要向外幫人傳福音.
上帝給他這個恩典.
他就做了一個很聰明的公頭.
做了一個很聰明的公頭.
立好了根基.
剛才我說了.
他負責傳福音的.
大人信主建立教會.
他就建立了這個很重要的第一層.
就是那個根基.
別人在上面建造的.
這個根基是什麼呢?.
十一節有說.
你看看你的程序表.
十一節有說.
因為那已經立好的根基.
就是耶穌基督.
教會就是那間屋.
做個比喻.
保羅做報道.
根基是耶穌基督.
阿波羅做陪靈.
栽培生命.
是很重要的.
所以耶穌基督不單是教會的頭.

$^{241}$耶穌基督也是教會的根基.
聰明的僕人.
就懂得將根基.
立在基督之上.
本來所有的神的僕人.
都應該將教會.
建立在耶穌基督的身上.
但你發覺不是的.
你發覺不是所有神的僕人.
都將教會建立在耶穌基督的根基上.
都有很多種的情況.
有義端的教會是怎樣的.
義端的教會扭曲了信仰.
除了說耶穌基督的救恩.
說起來也很奇怪.
很多字字節季扭曲.
他又會強調地面上有偉大的僕人.
這些僕人領受了信息.
所以要記住跟著這些偉大的僕人.
大家就將教會建立的根基.
不只是在耶穌基督的救恩上.
或者在耶穌基督的身上.
是有很多其他東西.
你說這些就是義端會做的.
我們正統教會不會做的.
不是的.
正統教會都很容易受到.
世俗化的影響.
很容易受到世俗化的影響.
商業世界怎樣我們又怎樣.
商業世界原則是這樣.
我們就又這樣.
所以到最後你會發覺有很多雜質.
你不是單純建立在耶穌基督的身上.
是有很多雜質.
很多利益.
很多世界的用的方法.
世界的人看重的東西都進來了.
所以教會根基不穩的時候就動不動的.
就是這樣的情況.

$^{281}$但耶穌基督是我們的頭.
我們的主.
祂是我們的根基.
而且這個根基是很穩固.
很可靠.
有祂做根基.
你將教會建立在耶穌基督的身上.
是一定不會倒的.
很多時候倒是因為.
不是單純建立在耶穌基督的身上.
太多的雜質.
太多扭曲的東西.
在這裡很快的給大家看看.
保羅作為神的僕人.
做了很多呼音的工作.
這是第二次的步道之旅.
如果大家有看《仕途行傳》的時候.
是看到他一直傳呼音.
一直傳到哪裡呢?.
後面紅色的那個位置.
就是哥林多教會.
今天經文所說的那個群體.
那個地方.
就是這裡發生了爭執.
也在這裡.
保羅很想去教導弟妹.
看看教會究竟是怎樣一回事.
生命的工程究竟是怎樣一回事.
我們看到上帝給我們福氣.
有榮耀去服侍他.
讓我們在教會裡是有份的.
在事奉裡是有份的.
所以有人建立了根基.
這個也是神給我們的福氣.
可以傳呼音的.
可以帶人歸主的.
建立教會的根基.
這個神給我們的.
是福氣來的.
是榮耀來的.

$^{321}$之後在上面再有人建立上去.
今天這個經文就花了很多篇幅說.
在上面再建立的時候是怎樣.
保羅去提醒那些人.
提醒弟兄姊妹去思考.
根基很穩固的了.
因為建立在耶穌基督的身上.
你再在上面建立的時候.
要小心 謹慎.
自己是用什麼材料去建立.
有人立好根基.
其他人就可以在上面建造.
其實建造的人都是很多的.
不只是教會的傳道人.
不只是那些頭羊.
還有很多牽涉建立生命當中.
實在是太多了.
所以弟兄姊妹你也有份的.
你有份去建立教會.
你有份去建立屬於基督的群體.
所以這番說話都是對大家說.
我們用什麼材料.
來建立上面的工程.
會帶來不同的果效.
說的是工程.
這裡說立好的根基是耶穌基督.
之後十二節就說了.
有不同的選擇.
如果有人用金銀寶石.
草木和街.
再在根基上建造.
各人的工程必將顯露.
因為那日子要將它顯明.
有火把它暴露出來.
火要試煉各人的工程怎樣.
這裡一共有六種材料.
一共有六種材料.
是可以用來建造工程的.
六種,但是你說.
看標點符號都看不到六個分別.

$^{361}$是不是?.
就不是的.
是金銀貼在一起.
應該是金鈍號銀鈍號.
寶石一個草鈍號木鈍號.
然後和街.
所以是六個材料.
六個材料很值得大家去思考.
給大家看看.
先看看.
金,厲害了.
金直到今天仍然是很值錢.
但是重點不是說它是否值錢.
金銀寶石屬於第一類材料.
金銀寶石.
銀,非常漂亮.
寶石,是吧?.
這三種材料是屬於很貴重的.
重的,重的.
但你看它不是很大份量.
不會是很龐大的組織.
是一個很小的,已經是很貴重的價值.
重和硬.
這裡要特別對比.
後面將會有火的試煉.
有火的試煉.
火,你想它燒的那些金.
燒的銀,寶石.
不是那麼容易著火.
因為你看是一個對比.
對比後面.
這些是很大件的.
前面那三種不是那麼大件.
但價值不非.
又不容易著火.
最重要是不容易著火和堅固.
重點是堅固.
後面這些是大大件的.
很顯眼的.
例如木,其實木不是說碎木.

$^{401}$木是說木柱.
很厲害,一看就很顯眼.
木,其實應該是木先.
但它說的是草木,實際應該是木先.
木之後就是草.
草之後就是和楷.
如果按次序.
一定是最容易著火的是和楷.
接著到草,接著才到木.
所以你數上去的時候.
到寶石,到銀,到金.
前三樣東西就很堅固.
不容易著火,有價值.
後面就是大大件的.
但越來越不堅固.
很容易著火.
而且你看得出沒有什麼價值.
其實這六種材料.
說真的,大家都看到的.
大家都知道價值如何.
知道它是多麼堅固,還是不堅固.
大家心中有數.
但你會發覺人仍然會選擇.
用草木和楷.
不是每個人都選擇金銀寶石.
你有沒有留意?.
為什麼不是每個人都選擇金銀寶石?.
我明知道那些東西很堅固.
我又知道那些東西很有價值,很貴重.
我還是要選擇木材,草木和楷.
很容易被著火,為什麼呢?.
這一點很值得我們思考.
在這裡說的是.
不同的人用不同的材料.
建造教會生命的工程.
我們用不同的質素.
參與教會的事奉,建立教會.
到有一天,上帝審判我們的時候.
就要看我們生命的工程.
是否能經得起考驗.

$^{441}$有審判的嗎?.
不是只給未信主的人嗎?.
如果有想我,啟示錄的弟兄姊妹是很深刻的.
啟示錄的第一個意象裡.
是說得很清晰.
教會應該說是神透過約翰.
告訴七教會.
上帝怎麼看待那七教會.
那七家教會有什麼做得好,有什麼做得不好.
上帝要稱讚什麼,提醒什麼,警告什麼.
要教會悔改.
啟示錄一開始的意象就說這件事.
有一天主要回來.
不只是審判其他未信主的人.
也要審判教會.
教會如果不悔改.
到那天有很嚴重的後果.
非常嚴重.
所以在這裡其實很特別.
原來不只是啟示錄說審判的信息.
在福音書裡,或者在這些保羅書信.
都是滲透這些信息.
是什麼情況呢?.
回到剛才那段經文,我想跟大家看.
因為哪一天.
哪一天要顯明主耶穌再回來.
剛才我記得靈祥一開始的時候.
有一個很大的盼望帶著我們.
因為聖誕節.
我們紀念耶穌基督的降生.
是很寶貴,真的很好.
在裡面說到我們盼望主的回來.
是的,我們屬主的人.
盼望主的再回來.
但是他再回來那一刻.
要問自己,我和你準備好了嗎?.
有火的試煉就是有審判.
上帝不只是審判外邦人.
不只是審判未信主的人.
主耶穌將會是審判者.

$^{481}$祂要審判我和你.
我和你準備好了嗎?.
將會有火的審判來考驗我們.
我們所建立的生命工程.
是否經得起考驗?.
經過火的考驗之後.
全部燒了,你說有些什麼?.
即是沒有了.
我的生命是經不起考驗的.
我所建立的東西.
全部在上帝眼中是沒有價值的.
完全不值錢的.
不單止沒有價值.
我自己的生命也像在火中經過一樣.
我的生命也有虧損.
這個審判是很嚴厲的.
但不是的,主耶穌再回來.
審判世界,審判教會的時候.
經過火的試煉,原來.
我的工程還屹立不倒.
是金銀寶石.
主耶穌就要稱讚我們.
所以我問你.
今天準備好了嗎?.
為什麼我們明知材料會帶來不同的效果.
我們仍然會這樣選擇呢?.
很明顯,第一件事.
我們是不懂得分辨.
什麼叫不懂得分辨?.
我不知道在上帝眼中的金銀寶石是什麼.
上帝眼中的寵物和皆是什麼.
可能我覺得好的東西.
和世界覺得好的東西一樣.
但原來上帝眼中的寵物和皆.
我最看不起的,最沒有耐性做的.
就是上帝眼中的金銀寶石.
我不會去選擇.
這是我們的眼光問題.
我不懂得什麼是上帝眼中的金銀寶石.
所以我可能選擇了上帝眼中的寵物和皆.

$^{521}$所以我一直在教會裡.
我都說自己很努力服侍.
我做了很多事.
到最後主宰回來的時候.
原來是沒有價值的.
這樣不是很慘嗎?.
但這裡說得很清楚.
火要試煉各人的工程如何.
不只是試煉一個,半個.
我們有份一起參與在教會的事工裡.
一起去建立神的家.
今天就這個生命的工程.
有一天要面對上帝.
我和你有沒有信心.
這些是經得起考驗的.
在根基上所建造的工程.
如果保得住的話.
他就要得獎賞.
如果被燒了,他將會損失.
但這個不影響救恩.
免得我們又去到極端.
哎呀,那我就不得救了.
不是的,我仍然得救.
但我這個工程沒有了.
那就得不到上帝的獎賞和肯定.
這絕對是一個很大的損失.
這真的很特別,值得提醒我們.
我想說,我都思考了很久.
上帝眼中的金銀寶石是什麼.
其實經文沒有一個很清晰的.
去說出來.
但我想起,耶穌基督是教會的頭.
耶穌基督也是教會的根基.
那中間的材料.
是不是完全和耶穌基督脫勾呢.
這不是很奇怪嗎.
我相信中間所用的材料.
去建立工程.
一定不會和耶穌基督脫勾.
這些材料一定是主所喜悅的.

$^{561}$主會重用的,看為寶貴的.
那我們就要花時間去求問上帝.
主啊,你體重是什麼.
你眼中的金銀寶石是什麼.
我們就要用這些去建立教會.
今天洪恩家也是.
我們事奉不是憑著我們的經驗.
我告訴你,我過去搞很多大型活動.
我經常都可以這樣跟別人說.
我以前,孫生在這裡都知道.
這麼多年我做青少年工作.
我搞多少營地,幾十八個活動.
對著多少青少年人.
我都可以完全不需要求問上帝.
一來到就行了.
我就憑著我的經驗.
12345678這麼多營地.
你要我搞活動,幾十八個都可以.
什麼都可以.
這些是不是金銀寶石呢.
主要這些,耶穌基督要這些.
我有沒有問過我們教會的頭.
耶穌基督,耶穌基督,你想我們教會.
怎樣進行呢.
怎樣去建立我們生命的工程.
有沒有問過祂.
我們的根基真是耶穌基督.
還是我們很多雜質.
有沒有問過上帝喜悅什麼.
耶穌基督喜悅什麼.
將主權給祂.
是不是.
憑著我們的經驗.
我可以很多很多活動.
那些活動可以像煙花一樣.
在天空綻放的時候.
很漂亮,煙花.
再來,可以很多個.
放完之後,沒有了.
煙花比這些草木和街還要差.

$^{601}$渣拿都沒有.
在2022年的最後一個月份.
我想我們回想.
如果是這樣,在新一年.
我們怎樣去事奉呢.
是不是就出來了.
我們最重要的是.
你想想有什麼經驗.
問問他,問問他.
全部都問了.
我們就不問耶穌基督.
不是專心尋求耶穌基督.
我想跟你說,專心尋求要時間.
要等候.
這些我們現代人.
特別是香港人是沒有耐性的.
香港人最重要快.
現在不能過關.
少過關.
我們通往國內.
以前我記得過關的時候.
看到香港人衝到去.
一直衝衝衝.
我心裡總是問一個問題.
你衝到哪裡去.
你的終點站你衝到哪裡去.
衝得這麼快幹什麼.
站在電梯上.
人們也是,電梯了.
扶手電梯,你就站在那裡.
不行的,還要衝上去.
什麼都講求快.
所以要停下來.
靜下來,祈禱.
等神,問神.
真的給神,說了算.
這些是很艱難的.
所以我們很喜歡選擇草木和街.
就是這個意思.
浪費了很多力氣.

$^{641}$新一年,是不是立志呢.
我剛才說回.
我聽創風.
分享的時候,心裡有一個很深的感受.
創風說.
昔日他回教會.
很久之前.
他覺得教會在那個位置.
好像幫不到自己.
其實這個呢.
都是.
好一類朋友的心聲.
回到教會.
不是這麼簡單.
我和你說.
我們等於姐妹彼此相愛.
如是者崇拜聚會.
行了,完結了.
下一個呢,我們又怎樣呢.
一起出去吃飯.
真是好了,我們在這裡.
我為你祈禱,你為我祈禱.
事情都是這麼簡單的.
很多時候,我們的生命.
是很多的創傷.
很多的痛苦,很深的痛苦.
創傷.
很多的甚至是.
捆綁.
眼淚.
我們進來教會.
很想找到幫助.
那怎樣呢.
行了,我們全部找老牧師.
老牧師交給你.
我們最重要是出去吃飯.
我們彼此聯誼,開心一下.
行行行,交給曾姑娘.
或者有誰.
創風負責.

$^{681}$做輔導,我們交給創風.
是不是這樣呢.
我想問,如果真是.
有很多.
我們洪水橋的街坊.
進來尋求幫助.
那不去教會.
去哪裡呢.
我們是否真的承載了.
是不是我們全部.
選擇那些生命沒有憂傷的.
生命沒有困難,沒有痛苦的.
很舒服的.
坐在這裡,說你好.
我好,那些,只是選擇.
這些來牧養.
我想問.
我們教會是否真的準備好承載.
用什麼去承載.
用我們人的意思.
用我們人的方法.
用我們人的愛.
會不會將來都會有更多.
分享,我想進到教會.
尋求幫助,但結果我尋求不到.
所以我挑戰大家想想.
那個生命的工程不簡單.
為何要.
回到上帝那裡.
不回到上帝那裡.
怎會有金銀寶石呢.
每一次都好像套餐.
ABC套餐.
不是A就B就C.
我們教會只有這樣.
你吃不吃.
怎會是套餐呢.
每個人的生命都不同.
現在說的是建立群體.
需要更大的力量.

$^{721}$所以在這裡.
想彼此挑戰.
不只是挑戰大家.
也挑戰我自己.
因為我記得我預備這堂課.
第一刻心裡有點痛.
我會想起.
以前我去牧養青少年群體的時候.
我的生命工程.
是否經得起考驗.
心裡都痛.
以前青少年的群體.
有些生命現在都很好.
在教會成長.
但很多都流失了.
很多的生命工程.
我自己想起.
我經不起考驗.
那些生命流失了.
很多生命都覺得.
在當中得不到幫助.
我自己都一直去反省.
是否不停地搞活動.
搞啊搞啊搞.
我做了很多寵物和街的事.
沒有真真正正回到上帝那裡.
去求這些金銀寶石.
用這些材料.
是合乎耶穌基督心意的材料去建立教會呢?.
這是我先提醒我自己的.
盼望和你勉勵.
我們不要離開了這個門口之後.
就當沒有這件事.
因為給自己一個納藥.
在2023年的時候.
一個全新的開始.
捨得付代價.
去上帝那裡求.
等神.
去幫神.

$^{761}$既然頭是耶穌基督.
根基都是他.
那我更加應該要用.
耶穌基督眼中的金銀寶石.
建立洪恩家.
我們一起禱告.
天父上帝.
今天透過你的話語.
去提醒我們.
提醒我們去思考.
今天我們是用什麼材料.
去建立.
我們的家.
我們的教會.
主啊,如果今天我們用的是.
草木和街.
求你真的光照我們.
你讓我們知道的時候.
我們的心是柔軟的.
肯認罪,肯悔改的.
我們回到你那裡,求你.
讓我們知道.
金銀寶石是什麼.
我們要用這些材料去建立洪恩家.
因為我知道.
你很愛我們.
你深深的愛我們.
要我們這間教會在洪水橋區.
成為明亮的燈台.
所以求你興起我們.
當我們走出這個門口的時候.
你帶領我們.
讓我們親近你.
讓我們將.
心,傳言.
交給你.
主啊,求你在這裡.
彰顯你的大能.
你讓我們洪恩家將來.
是一個很堅固.

$^{801}$很強壯的一間教會.
是一個.
很明亮的燈台.
在這裡照亮很多.
在黑暗中的人.
也有很多人進來尋求幫助的時候.
主啊,我們整個群體.
是盛載到人們的生命的.
而這個工程.
是經歷起考驗.
有一天面對你的時候.
主你會稱讚我們.
肯定我們是你所喜悅的教會.
我們就這樣.
將我們自己交在你的手中.
奉耶穌基督的名字.
祈禱,阿們.
\newpage



\section{約翰福音 14:16-17-25-26-15:26-27-16:5-15}
\label{sec:OuveUNRd_9o}
\textbf{2022.12.11 《認識聖靈》 約翰福音 14\_16-17,25-26,15\_26-27,16\_5-15 (和修版) 講員 : 陳秀賢牧師}
\newline
\newline
連結: \href{https://youtube.com/watch?v=OuveUNRd-9o}{\texttt{https://youtube.com/watch?v=OuveUNRd-9o}} ~~~~ 語音日期: 2022-12-12
\newline
\newline
\hyperref[sec:qxsO_vpno54]{\small{< < < PREV SERMON < < <}}
~
\hyperref[sec:index_chronic]{\small{[返順時目]}}
~
\hyperref[sec:index_scriptual]{\small{[返順卷目]}}
~
\hyperref[sec:pCQk0pIrnYA]{\small{> > > NEXT SERMON > > >}}
\newline
\newline
$^{1}$各位女士早安.
很感恩和大家一起敬拜和守聖餐.
紀念主耶穌基督.
上一次我們講到也是戴口罩.
現在也在戴.
所以三年來我們也面對一些困難.
不過我們的主是永恆不變的.
我們一起有祈禱.
主耶穌我們今天在聖餐裡紀念耶穌基督.
你為我們降生和犧牲.
因為你先愛我們.
所以我們能夠在疫情裡.
仍然持續敬拜和認識一位獨一的真神.
耶穌基督.
耶穌你是我們的救主.
是我們萬人的救主.
多謝耶穌你是醫治之主.
因為你是醫治之主.
所以我們懇請你醫治阿路穆斯.
因為看到他皮膚也有問題.
主你是醫治之主.
你施行拯救.
你醫治他的皮膚.
讓他的皮膚好過來.
主耶穌醫治我們所敬重所愛的童宮.
阿路穆斯.
求主恩待他.
主你是在來的君王.
所以我們要繼續認識你.
等候你在來的日子.
多謝你.
因為我們擁有一個不能震動的國.
所以主我們看到.
我們的弟兄姊妹努力的敬拜.
弟兄姊妹做.
很用心的服侍神.
主我們懷著敬畏你的心去服侍你.
因為神你先愛我們.
我求聖靈在我們當中運行.
叫我們明白聖經真理.

$^{41}$我們祈禱交托.
奉主名求.
阿們.
我很多時候來的時候.
我看到有些弟兄姊妹.
他們很努力的去按PowerPoint.
我呢就很烏龍.
我在試PowerPoint的時候.
按不到按不到.
其實我自己很烏龍.
都還沒有按安制.
所以弄得大家有一點緊張.
所以很欣賞大家的服侍.
今天我選了一個主題.
就是認識聖靈.
我們有時候在我們的成長.
其實我在宣道會都幾十年了.
有時候我們會.
很認識聖靈.
或者不是很認識聖靈.
我們不認識聖靈.
所以有時候會提聖靈.
但有時候我們又見到一些教會.
太高舉聖靈.
忽略耶穌基督.
我們怎樣從聖經裡面.
取以平衡.
究竟神是怎樣提聖靈.
教我們怎樣去知道.
我今天就去分享認識聖靈.
使徒信經我本來想和大家讀.
但原來你們很好.
程序表那裡也一起讀了.
在使徒信經那裡.
其實是有提到.
聖父,聖子,聖靈.
而聖靈是因聖靈感應.
耶穌基督是.
我們信聖靈.
我們信聖兒宮之教會.

$^{81}$我們信聖徒相通.
所以在使徒信經那裡.
是有提及聖靈的.
在宣道會的信仰.
可能魯木斯教你們的洗禮班.
一定是有教宣道會的信仰.
第一條.
可否讀一次給我聽.
預備請.
信神是無瑕.
完全唯一真神.
聖父,聖子,聖靈.
三而為一.
同位同權.
同榮.
永世無疆.
講第一條信令.
我們所信的.
聖父,聖子,聖靈.
三而為一的神.
祂是同位同權.
同榮同尊.
換句話說.
我們尊榮聖父.
我們要尊榮聖子.
與此同時.
我們要尊榮聖靈.
不要偏心.
我們要.
因為他們是同位同權.
同榮.
我們要尊榮.
三位一體的真神.
第三條.
宣道會的信仰.
都講聖靈.
可否讀一次.
預備請.
信聖靈是三一神的第三位.
受耶穌基督的差遣.

$^{121}$而居於信徒心中為保衛師.
領導信徒進入真理.
使世人為罪.
為義.
為審判.
自己責備自己.
第三條信條.
講到聖靈會帶動我們.
為罪.
為義.
為審判自己.
責備自己.
所以我們都很明白.
聖靈的工作.
我早來的時候.
與師母聊天.
原來她22日要講報道會.
我今天下午都要講報道會.
其實講報道會.
我們都要依靠聖靈.
因為聖靈是使世人為罪.
為義.
為審判自己.
責備自己.
我們都做不到什麼.
是聖靈工作.
我很想透過約翰福音.
剛才主席也跟我們讀過經文.
我們親自對約翰福音.
有些講述.
和認識聖靈.
約翰福音十四章.
講到真理的聖靈.
與我們同在.
讓我們明白聖經真理.
稍後我會詳細講解.
現在我只是跳過.
約翰福音十五章.
真理的聖靈來到這個世界.
與我們同在.

$^{161}$他是保衛師.
約翰福音十六章.
講到保衛師來到的工作和功能.
已經跟我們讀了一次.
我第一很想講.
聖靈是有位格的.
什麼叫有位格呢.
是有意志,有反應,有行動,有感情.
剛才主席帶我們唱歌時.
是否有感情.
人是有感情,有行動,有反應的.
我們有意志.
聽到也要有意志.
否則我們會睡覺.
我們要鬥智去聽到.
有意志.
我們有反應.
人有反應.
聖靈也是這樣.
而聖靈不是單一種能力.
可能我們說.
聖靈是一種能力.
我們求能力.
不單單是能力.
有時我們形容聖靈是如火或風.
其實聖靈也是有位格的.
好像聖父,聖子,聖靈一樣.
聖靈有位格.
即是祂會聽話.
所以祈禱祂會聽到.
你這麼激情去唱詩歌.
祂會聽到.
祂會說話.
祂會作見證.
又叫人去馴服基督.
聖靈又會榮耀基督.
會帶領和命令.
命令我們.
禁止我們做一些事情.
是給予人的幫助.

$^{201}$所以原來聖靈有位格.
意思是很有畫面.
聖靈不是一個死物.
雖然我們聖靈是看不到.
耶和華從舊約裡面.
你會看到如何帶領以色列人出埃及.
繪形繪聲很多神蹟奇事.
耶穌基督降世到城域.
祂是醫治人.
好像是人性化一點.
知道一些事情.
但其實聖靈也有祂的位格.
祂也有祂的位格.
我們要了解.
聖經例子裡面.
聖靈有什麼工作呢.
大家很熟悉的.
馬太方28章19節.
奉聖父聖子聖靈的命.
給他們施洗.
以及在伊弗所書4章30節.
你不要叫聖靈擔憂.
各位弟兄姊妹.
其實你的子女犯錯.
你會擔憂.
你的子女不生性.
你會擔憂.
聖靈一樣.
如果我們做錯事.
不生性.
聖靈也會擔憂.
我們不可以褻瀆聖靈.
不可以褻瀆聖靈.
我怕你們有阻礙.
不要褻瀆聖靈.
大家知不知道.
《四郎傳》第五章.
有一對夫婦.
他們也很奸詐.
就是阿娜莉亞和薩菲拉.

$^{241}$他們兩個奸詐.
買了田產.
又說已經買了.
其實那些錢奉獻了.
其實他們不是.
他們騙聖靈.
我們不可以消滅聖靈的感動.
聖靈有沒有感動你呢.
譬如唱詩歌的時候.
古舊十教很感動.
因為耶穌基督釘在十字架上.
我們很感動.
我們不要消滅聖靈的感動.
或者聖靈今天提醒你.
其實你今天得罪人了.
你應該向人道歉.
聖靈也會經常感動我們.
你要在意.
當你相信耶穌.
聖靈住在我們裡.
祂會感動我們.
如果我們洩漫聖靈.
我們不可以輕慢聖靈.
聖靈會禁止一些事情.
大家記得在《賜堂傳》.
阿保羅他們一班人.
宗教團隊.
本來想去亞細亞.
後來被聖靈阻止.
他們去了馬其頓.
聖靈會引導我們.
聖靈用說不出來的嘆息.
為我們祈禱.
各位姐妹你們有沒有試過.
憂傷到程度.
都不想祈禱.
有壓力到.
但又不能祈禱.
但聖經告訴我們.
聖靈用說不出來的嘆息.

$^{281}$為我們祈禱.
真是很叮噹.
聖靈會帶著肥理.
去傳福音.
在《賜堂傳》第八章.
聖靈會帶著肥理.
去安提阿伯.
即是太監那裡.
跟他說福音.
去帶順主.
這個太監也是管錢財.
很威風.
我們要留意.
聖靈有很多工作.
你能否感應到.
聖靈在你裡面.
做什麼工作.
我相信你們能感應到.
因為我們都是基督徒.
大家知道疫情.
其實在這三年裡.
我們都好似很困難.
帶人去信耶穌.
不過在疫情裡.
我都帶了很多人去信耶穌.
其中一個是神感動我.
向我舅父.
帶他去信耶穌.
我舅父都是帶我.
幾年左右.
六十幾歲.
話說在前年.
他患了胃癌.
大家知道胃癌.
要整個胃割了.
所以他吃東西.
做完手術之後.
都幾大劑.
腸很辛苦.
吃東西.

$^{321}$我有時都探望他.
跟他說福音.
不過出院之後.
有一次很精靈.
跟我喝茶.
他跟我說.
秀賢.
你以後不要再跟我說耶穌.
一盆冷水潑過來.
有些吃檸檬的感覺.
我就放下了這件事.
過了兩三個月之後.
我舅母打給我.
她說.
舅父有點不適.
我問她為何有點不適.
她說應該醫生說.
他今天快要死了.
我問她快要死了.
大家知道疫情.
有不同的時段.
有不同的情況.
剛剛那個時段.
醫院很緊張.
不讓我們去探望.
讓一兩個家人去探望.
很緊張.
大家知道.
那個時段就很緊張.
我跟舅母說.
舅母就跟我說.
如果舅父死了.
我就在天堂見不到他.
我跟她說.
舅母你上教會的嗎.
我這麼多年都沒聽過你上教會.
但無論如何.
既然你都想他跟你一起上天堂.
拜託你.
只有你去醫院探望他.

$^{361}$你可不可以跟他說.
帶他去耶穌.
我舅母就說.
我不懂得帶人去耶穌.
我就WhatsApp那個決志內文給他.
我說你行了.
拜託你.
你進去醫院的時候.
我照樣叫他跟你.
如果他相信耶穌.
聖靈都不斷帶動這件事.
誰知道.
我都衝去醫院探望他.
因為醫生說不行.
去到走廊.
我的表哥表弟表妹都在走廊.
他們很哀傷.
因為大家知道疫情.
探望不到家人.
第二他就會這樣過世.
他們很哀傷.
我問你們探望不到.
你們都探望不到.
他們真的很意興闌珊.
很失望離開醫院.
但現在我表弟就說.
表姐其實我們可以視像.
跟爸爸.
我可以視像.
現在在什麼視像呢.
在屯門醫院的大堂視像.
一視像.
我舅父就真的聽電話.
我立即視像.
你知道時間有限.
我立即視像.
舅父你真的要相信耶穌.
耶穌愛你.
你同不同意.
同意.

$^{401}$同不同意.
同意.
耶穌基督復活.
你信不信他.
信.
你現在快點跟我決志.
他就跟我決志.
決志之後.
你知道.
視像很短時間.
決志之後.
先收起來.
祝福你.
我的表哥表弟舅母就在我後面哭.
很感動.
他相信了.
那晚我舅父昏迷.
昏迷了.
就是那段時間.
我昏迷了一段時間.
又出院.
我就幫他洗禮.
在家裡洗禮.
我幫他辦喪禮.
當時在喪禮的時候.
我有幾十個親戚都是水上人.
大家知道水上人拜什麼.
很多東西都拜水上人.
我一次在喪禮講道.
什麼仇都報了.
可以一次過講給我的親戚聽.
很爽.
我很想講這件事.
其實是聖靈帶動我和感動我的.
因為當初我跟舅父說.
我吃檸檬.
叫我不要再講.
就不講了.
但是當有些事情發生.
聖靈就說快點吧.

$^{441}$把決志給我舅母.
聖靈一邊帶動.
就帶了我舅父信耶穌.
各位弟兄姊妹.
我們不要被撒旦騙我們.
就覺得.
這個不是得時.
這個時辰不是很順境.
帶不了人上耶穌.
其實不是的.
無論得時不得時.
我們都可以帶人上耶穌.
你會不會因為聖靈帶動你呢.
我有一個朋友.
他很淒涼的感覺.
他的朋友退休了.
前兩個月和他踩單車.
大家都想著退休很開心.
我也退休了.
他很開心.
以為有很多東西可以玩.
就和他踩單車.
誰知一踩.
這個朋友就跌在地上.
他親自送去醫院.
就死了.
我們就.
譬如我們紅春堂.
可能接觸過不同的退休人士.
你有些很好的朋友.
你求聖靈充滿你.
好像帶領阿肥你一樣.
帶著聖靈感動阿肥.
和全方位去.
所以其實你敏銳於聖靈.
在你生命裡的工作.
你可以看到很多果子.
第二,我想說聖靈的角色.
聖靈的角色在約翰福音14章16至17節.
提到保衛師.

$^{481}$Another Counselor.
第一個保衛師是誰呢.
如果聖靈是另一個保衛師.
第一個保衛師是誰呢.
就是耶穌基督.
約翰福音1:2.
他說:我小子們.
我將這些話寫給你們.
是叫你們不犯罪.
如果有人犯罪.
在乎那裡.
我們有一位忠報.
就是那二者耶穌基督.
他為我們的罪索了贖罪制.
不但為我們的罪.
為普天下的罪.
原來第一位忠報保衛師.
就是耶穌基督.
是耶穌基督.
而另一位保衛師是聖靈.
耶穌基督是他自己.
去面對離開世界的時候.
他和那些門徒說.
那些門徒都很害怕.
因為他經常和他們說.
預言他會釘十字架.
但是聖靈.
就是耶穌說.
我會派另一位保衛師.
保護你們.
將真理告訴你們.
那就是聖靈.
那聖靈的意思是.
保衛師是什麼意思呢?.
Counselor是什麼意思呢?.
是安慰人和輔導人.
你發覺有些弟子們不開心.
不知為何有很多安慰的話.
從你的口而出.
安慰到人.

$^{521}$聖靈幫助我們.
聖靈也是幫助我們.
支持我們.
也是一個顧問.
聖靈也是一個辯護者.
聖靈也是我們的盟友.
很多個角色.
第三.
我想和大家分享聖靈的工作.
聖靈的工作是什麼呢?.
真理的聖靈去教導信徒.
教導他們一切事.
讓他們想起.
耶穌基督所說的一切的話.
大家知道耶穌升天.
所以和他的徒弟說.
你們不用擔心.
因為我會派保衛師.
和你們一起.
所以聖靈的工作.
是會讓我們想起主的說話.
真理的聖靈教導我們的時候.
不是憑自己說.
而是將他所聽見的.
說出來.
告訴我們.
他是受於天父和耶穌基督.
所以他受於這些說話的時候.
他說出來讓我們明白.
這是很真實的.
有時候小時候讀了很多金句.
有沒有這個經驗?.
小時候背誦金句.
但當你有時疾病.
不開心的時候.
聖經會跳出來.
除了你的記憶跳出來之外.
其實是聖靈告訴你的.
這樣讓你想起.
有時候我們年紀大.

$^{561}$試用也很辛苦.
有些吃力.
但可能主又讓你想起一些說話.
加油啊.
我們要謹慎地做我們的傳道工作.
這樣就被鼓勵.
所以好像有些東西.
在我們心裡浮起.
提醒我們.
這就是聖靈教導我們.
想起耶穌基督.
聖經所說的話.
讓我們明白.
聖靈很厲害.
他要揭示世人的罪.
他來叫世人為罪為義.
為審判自己.
責備自己.
為罪.
是因為他們不相信我.
耶穌說.
聖靈會責備他們.
因為他們為罪.
因為他們不相信耶穌基督.
為義.
因為我枉負那裡是你們就不見我.
為義是什麼呢.
好像法理賽人的義.
是一種破爛的義.
以為自己義.
法理賽人的義.
自以為義.
我們很多朋友.
未信主都是自以為義.
我沒有坐牢.
我沒有殺人.
我沒有姦淫.
我沒有擄掠.
有很多人是這樣.
做社工 做老師.

$^{601}$這些是破爛的義.
因為這種破爛的義.
他們都不信神.
並且法理賽人和猶太人.
這種破爛的義.
是將基督釘在十字架.
耶穌基督以義者代替不義.
他已經枉負那裡.
他將會枉負那裡.
所以聖靈來.
讓世人知道自己破爛的義.
其實他們就將耶穌釘在十字架.
為審判.
因為這個世界王受了審判.
世界王是誰呢.
是撒旦.
撒旦他已經受審判.
所以耶穌基督說.
聖靈來.
這個世界王將會受審判.
所以聖靈來.
去指責這些一切的事.
這個圖不是我做的.
是幫我做的.
是帕邦的弟兄做的.
他好像一個律師.
聖靈好像一個律師.
一個檢控官.
指控世人的罪.
所以各位弟兄姊妹.
當我們去帶人信耶穌.
你必須要將這個經文.
自己背一次.
主要我帶人信耶穌.
都不是我自己的口才.
又不是我自己的能力.
聖靈你可不可以.
叫那個人認罪.
為罪為義.
為審判自己.

$^{641}$責備自己呢.
所以你們都有經驗的.
舉辦大型報道會.
會不會同步祈禱呢.
用這支經文去祈禱.
求聖靈光照那些罪人.
求聖靈光照我們.
以致我們見到基督的光輝.
第三就是.
聖靈為主作見證.
和榮耀耶穌.
經文告訴我們.
聖靈來是為主耶穌作見證.
聖靈來是為主耶穌榮耀.
如果我們太高舉聖靈.
是錯的.
因為我們要高舉基督.
要給耶穌基督的名字.
被舉起.
吸引萬人歸向他.
所以聖靈來是為主作見證.
和榮耀耶穌.
有一個很出名的神學家.
不知道你們有沒有看過他的書.
叫做巴赫.
他講過一個例子.
他寫了一本書.
《活在聖靈中》.
巴赫這個神學家.
他講的例子.
聖靈是怎樣呢.
一座建築物.
你看見建築物很光亮.
因為建築物裡面有燈.
所以就光亮.
所以聖靈會光照耶穌基督.
可以嗎.
將耶穌基督顯露出來.
這個景我不知道你們有沒有看過.
你們知道在哪個區.

$^{681}$這個區在我家樓下.
就是屯門.
在柏嶼商場.
他們新建的.
他們有很多樹.
如果晚上燈都關了.
遠看.
你從涼亭走過去.
就看不到樹的存在.
但當他晚上.
他有幾千顆小燈.
纏著樹.
他一有人按掣.
燈就亮了.
一亮就纏著樹.
所以就看見樹形.
所以就算路遠在涼亭村.
搭輕鐵走過.
都看見這些亮閃閃的樹.
這個例子都會想起巴赫這個例子.
其實聖靈是甚麼呢.
是Spotlight.
他將耶穌基督Spotlight出來.
讓基督的名字.
是被人知道和明白.
大概我講完了.
我可以做一些結論.
聖靈.
第一.
你記住.
他不是死物.
雖然我們看不到.
但你要知道他活著的.
他Living.
他有Person.
有位格.
他是另一位保衛師.
所以我們要尊榮聖靈.
剛才主席帶我們唱甚麼.
聖在聖在聖在.

$^{721}$聖父聖子聖靈.
所以就能夠達到.
同位同權同榮.
同尊.
所以這首歌是很classic.
和很永垂不朽.
因為唱對了.
第二.
聖靈.
我結論就是.
聖靈光照我們.
引導我們一切.
讓我們想起一切的說話.
詩篇說.
求你指教我.
遵行你的旨意.
因你是我的神.
你的靈本為善.
求你引導我平坦之路.
聖靈經常光照我們.
就算我們幾十歲.
我們也做錯事.
我們說話時.
得罪孫子.
得罪丈夫.
得罪後輩.
聖靈光照我們.
你不用這麼厲害吧.
你年紀大.
不用像老賣老.
你應該謙虛點.
和我們相處.
聖靈光照我們的時候.
我們就去下台街.
買些東西給他們吃.
或者說.
Sorry Sorry.
聖靈會光照我們.
我們做錯事的時候.
我們必須去.

$^{761}$去依聖靈去做.
可以嗎.
我一進來的時候.
見到老牧師.
為什麼情況這麼嚴重.
很想念他.
我們想念他的時候.
教會就要托著他.
為他祈禱.
希望他快快.
慢慢慢慢好.
正如我們.
等於姐妹有病的時候.
聖靈也提醒我們的時候.
我們也為我們等於姐妹祈禱.
將我們等於姐妹交給神.
所以聖靈光照我們.
想起很多黃道真理.
我們行真理去做.
第三就是.
我們經常留意聖靈的光照.
我們看看裡面有沒有黑暗.
有沒有黑暗的權勢.
或者黑暗的說話.
有沒有苦毒.
有些黑暗的東西.
有沒有黑暗的東西.
最近我們有沒有很多苦毒.
很多怨毒.
一開口就有很多怨毒.
留意一下.
求聖靈光照我們.
以致我們靠著上帝的恩典.
有悔改.
有改變.
神會幫我們.
我們的命會越省越靚.
我們要謙虛.
聖靈會光照我們.
看我們有什麼錯.

$^{801}$有沒有黑暗的面.
有沒有說話傷害人.
有沒有傷害人的眼神.
當聖靈光照我們的時候.
要收斂一點.
攻克一點.
攻克多深.
第四就是辨別真理的靈.
和謀望的靈.
我們讀一下聖經.
預備請.
凡另認耶穌基督.
是成了肉身來的.
就是出於神的.
從此你們不可以認出神的靈.
凡另不認耶穌.
就不是出於神.
這是那滴基督的靈.
你們從前聽見他要來.
現在已經在世上了.
有些人很害怕.
我認識聖靈.
會不會有邪靈上身.
或者我認識聖靈.
那些是不是邪靈.
或者你見到一些神蹟奇事.
發生在一些弟兄姊妹身上.
或者有些人.
很聖靈充滿.
你就很擔憂.
我們不要追求聖靈.
不要認識聖靈.
會不會有鬼魔上身.
那你就不用擔心.
因為一書說.
其實很容易認.
那個弟兄姊妹是被聖靈充滿.
還是邪靈充滿.
很容易認.
好像在少林寺初期教會.

$^{841}$那些人是.
聖靈充滿他們.
他們就帶很多人上耶穌.
是不是.
你跟隨他的果子.
會認到他的生命.
所以如果有人.
微信招人進來教會.
他有很多古怪的東西.
你就驗驗他的靈.
叫他跟你祈禱.
跟你唱詩歌.
通常如果裡面有邪靈.
根本他唱不到.
「勝在,勝在,勝在」.
唱不到這些歌.
因為是妙術神和敬拜神.
所以如果裡面有邪靈.
根本他跟不上你唱讚美歌.
當你去讀聖經的時候.
他也跟不上.
所以你可以辨別.
那個人是否被鬼附.
還是聖靈充滿.
是不是.
但因為當我們不認識聖靈.
或者害怕的話.
我們就會超越聖靈的工作.
和他的角色.
所以求主幫助我們.
在我們生命當中.
聖靈已經入住我們心裡.
我們很平安.
並且我們每一天都多認識他.
和了解他.
與他建立關係.
你帶人上耶穌.
求聖靈給你有能力.
求聖靈充滿你.
正如使徒行主一樣.

$^{881}$當神的靈充滿我們.
我們就必得著能力.
並要在耶路撒冷.
地極作神的見證.
今天下午.
我講完這篇道理後.
我就會坐的士去粉嶺.
去粉嶺的報導會.
我其實早已答應老穆斯.
來這裡的.
過了一段時間.
國安教會叫我去粉嶺的報導會.
我期待了很久.
我也去的.
很想去.
因為粉嶺宣導員.
他們借了宣導員.
大家知道宣導員是宣導會的地方.
在附近.
你們可能聽過一條村.
叫皇后山村.
是一個新屋村.
這間教會.
活石堂很想帶人上耶穌.
所以今天下午.
邀請一些街坊來宣導員.
聽福音.
所以我也願意為主效力.
走多理路.
將人帶到主面前.
我看上去也有年青的弟兄姊妹.
也有年長的.
我求主祝福我們.
雖然我退休.
但也有很多神感動我.
服侍的工作.
聖靈也感動我.
輔導,探訪,傳福音.
雖然我們這副主意.
是差的.

$^{921}$沒有以前那麼適合.
但當聖靈充滿你.
很奇怪.
你是有力量的.
你有喜樂的.
你有活力的.
去服侍神.
因為我不靠自己的力量.
是聖靈充滿我們.
祂賞賜喜樂給我們.
賞賜平安給我們.
我們一起祈禱.
聖父,聖子,聖靈三位的真實.
我們感激你.
我們有幸認識你.
讚美你.
主,我們願意在我們人生裡.
尊榮聖父,聖子,聖靈三位的真實.
並且我們效忠於你.
讓我們在生命裡繼續認識你.
也是將更多人帶到你的面前.
認識我們活著的神.
祝福紅印堂的弟兄姊妹.
無論我們的年紀如何.
但我們的心智是被你眺望的.
還有聖靈充滿我們.
我們就有能力.
也有恩賜.
繼續服侍一個不能振動的國.
禱告奉主名求.
阿們.
\newpage



\section{使徒行傳 16:6-10}
\label{sec:pCQk0pIrnYA}
\textbf{2022.12.18 《順著聖靈而行》使徒行傳 16\_6-10 (和修版) 鄧泰康傳道}
\newline
\newline
連結: \href{https://youtube.com/watch?v=pCQk0pIrnYA}{\texttt{https://youtube.com/watch?v=pCQk0pIrnYA}} ~~~~ 語音日期: 2022-12-23
\newline
\newline
\hyperref[sec:OuveUNRd_9o]{\small{< < < PREV SERMON < < <}}
~
\hyperref[sec:index_chronic]{\small{[返順時目]}}
~
\hyperref[sec:index_scriptual]{\small{[返順卷目]}}
~
\hyperref[sec:IOHR3FgR_rc]{\small{> > > NEXT SERMON > > >}}
\newline
\newline
使徒行傳 16:6-10
\newline
\begin{longtable}{cl}
\hline
\hline
章節 & 經文 (和合本修訂版)\\
\hline
16:6 & \begin{tabularx}{0.7\textwidth}{X} 因為聖靈禁止他們在亞細亞講道，他們就經過弗呂家、加拉太一帶地方。 \end{tabularx} \\ \\ \relax
16:7 & \begin{tabularx}{0.7\textwidth}{X} 到了每西亞的邊界，他們想要往庇推尼去，耶穌的靈卻不許。 \end{tabularx} \\ \\ \relax
16:8 & \begin{tabularx}{0.7\textwidth}{X} 他們就越過每西亞，下特羅亞去。 \end{tabularx} \\ \\ \relax
16:9 & \begin{tabularx}{0.7\textwidth}{X} 夜間，有異象向保羅顯現。有一個馬其頓人站著求他說：「請你過來，到馬其頓來幫助我們！」 \end{tabularx} \\ \\ \relax
16:10 & \begin{tabularx}{0.7\textwidth}{X} 保羅既看見這異象，我們就立即設法往馬其頓去，認為神呼召我們傳福音給那裡的人。 \end{tabularx} \\ \\
[1ex]
\hline
\hline
\end{longtable}
$^{1}$電視節目早安,你夠不夠我的精神?.
我經常說我很留意上帝的時間.
因為其實我很久沒有來到當中跟大家一起分享.
每一次我能夠來,我都很開心.
我便在祈禱上帝快點康復,但很奇妙的.
我發了兩條燒之後,基本上我原來是沒有事的.
但我一直都不轉陰,一直都還是陽.
直到剛剛,其實我是星期五才再次回到市中心.
所以其實上帝是很好的.
所以我祝大家聖誕快樂之餘,身體健康.
但亦都不需要懼怕,因為在過程中上帝一定會與我們同行.
今天跟大家分享這段經文的時候.
其實是一段很普通的記載.
但其實這段經文是很重要的.
為甚麼呢?.
這段經文裡面很有趣的地方就是.
其實是保羅事奉的一個很重要的分水嶺.
為甚麼呢?因為這段經文第十六章一開始的時候.
有一個很重要的人物出現了,就是提摩泰.
如果大家知道,提摩泰是保羅一個很重要的徒弟.
但這個人特別的地方就是提摩泰的爸爸是希臘人,媽媽是猶太人.
但他又是一個很虔敬的人.
在他身上亦都發生了一個原因,就是他又有行國禮.
我不知道大家是否記得,其實保羅在他的事奉裡面最掙扎的地方.
其實正好就是這件事.
就是他一路做一路做,他自己很清楚上帝所給他的領袖.
就是他要將福音傳給任何人,不應該只限於猶太人.
但猶太人就不斷要去強迫外邦人,你要行國禮,你要怎樣怎樣.
這是他們一直有很大的衝突,所以後來他才再回到耶路撒冷去匯報.
而來到這裡的時候,你就會發覺這個重要的人物出現了是重要的.
為甚麼呢?因為他之後就會帶著提摩泰一起進入這個行程當中.
而你也知道有很多很重要的教會是留下了提摩泰在裡面.
然後由他去牧養,其實是一個很重要的傳承.
第二件事,你會發覺這段經文這麼短,有兩個字很有趣.
一開始是說「聖靈既然今之,他們」在亞細亞.
突然間去到後來的時間就是「保羅既看見者而著,我們隨即」.
你會發覺原來就是這個作者路加正正式式親身加入這個團隊.
其實為甚麼這段經文這麼重要呢?.
在我還未說之前,我想透過這張圖告訴大家,很有趣的.
聖經在《小徒行傳》裡面,我們經常說「聖靈行傳」.

$^{41}$我們一直在說聖靈帶領著整件事.
這段經文裡面其實很少會出現兩次.
就是「聖靈既然禁止他們」在那裡說到.
然後「聖耶穌的靈卻不許」,兩次.
我想大家留意這幅圖裡面,有趣的地方就是.
你會發覺一開始是在特肥出發.
在那裡認識了提摩太之後,然後他們一直上去.
其實去到就是那個叫做安提柯的地方.
其實這個行程真正的起點是從安提柯.
大家知道安提柯是他們做宣教很重要的一個起點.
然後我們發覺很有趣,聖經就是說.
其實保羅原本想從安提柯去到加拉太,繼續向東走.
但是聖靈沒有讓他這樣做,然後他就轉了彎.
轉了彎之後就進了佛雷加.
然後佛雷加他又想著要向上面那個比推尼那裡去.
他想著要去北,但是耶穌的靈又不讓.
然後他又繼續向前走,經過美西亞去到特朗亞.
如果大家知道,去到特朗亞之後就是坐船.
然後他們收到馬其頓的呼聲之後.
就在那裡坐船去到馬其頓的地方.
這段經文有趣的地方是甚麼呢?.
其實當大家看完之後,這個很強調的就是.
聖靈介入了在保羅的工作當中.
而在剛才的圖裡面,大家很清楚.
保羅是有他自己的計劃的,有他自己想做的事.
但是你會發覺很奇怪,剛才說為甚麼路加的加入這麼重要呢?.
如果大家知道,路加福音和小太陽傳都是路加作者.
而路加是醫生,如果大家一起去學色經你都知道.
路加福音的特點就是很多很詳細的考究.
很多詳細的資料,因為他是一個科學家.
他有很多考證,他有很多的東西是寫得很清清楚楚的.
所以你會發覺一件事,就是在這裡反而出現了一個很有趣的地方.
他只是跟你略略帶過了聖靈介入,然後保羅就轉向兩次.
所以如果你是正常的話,你看聖經的時候,你一定會問這個問題.
究竟聖靈如何介入呢?.
大家想不想知道?其實我們很想知道.
為甚麼我會這麼肯定呢?因為在底英哲會我作為傳道人.
底英哲會經常來找我問我的時候,讓傳道說我這樣做是不是神的心意.
你明白我的意思嗎?其實我們是很想得到答案.
我都好像很想保羅一樣,我都很想聖靈直接告訴我,我就懂得怎樣走.

$^{81}$這個都是很多信徒,你最好告訴我,我就走.
但你會發覺路加故意不說,為甚麼呢?.
其實你會發覺路加是很聰明的.
聖靈的介入的時候,是保羅他們這個團隊很清楚領受了之後才會轉向.
但他不告訴你,因為他不想你的信仰是一個格式化的信仰.
甚麼是格式化信仰?其實格式化信仰是很危險的.
舉個例子,即是說,你是這樣,就是聖靈的意思.
你不是這樣,就不是聖靈的意思.
弟兄姊妹,但我想告訴大家,真正的聖靈的帶領,不是格式化的.
上帝和每一個人說話都不同.
我用一個例子,這次我中了covid,是不是上帝的心意?.
在我個人的領受裡面,我覺得上帝是我的恩典.
上帝在裡面的時候,為甚麼我會這樣說呢?.
就是在我中了之後,在我目的領受的弟兄姊妹,是多了很多人中的.
然後我是第一身的,可以即時提供很多資訊.
怎樣訂那些診所,怎樣訂那些扣離的士,吃甚麼,吃完藥之後有甚麼反應,怎樣怎樣.
我覺得上帝是給我的一個目的領受的工具.
但是說得不好聽,Joanne中了,上帝給她,是不是這個意思?不一定.
所以弟兄姊妹,我們今天其實如果你將你的信仰,去成為一個格式化的話.
是一件很危險的事,為甚麼我會這樣說呢?.
因為你就會有很多很多,其實是你自己給自己的理由,去行你自己的人生.
我記得有些弟兄姊妹,我是做傳達人,我聽過很多弟兄姊妹,當她要去讀神學的時候.
會有很多所謂的蒙召見證,我聽過很多的.
有些人,他們就怎樣呢?我打開那段經文,很奇怪,那段時間經常出現某一段經文.
然後出現了連續三次,所以我就覺得這個是神的呼召.
又或者反過來,某某某跟我說了某一些群體,連續幾次都見到的時候.
然後我覺得這個就是上帝自己的呼召.
我不會去排斥或者反對這些見證,因為這個是你個人領受.
但是如果你問我,這個是不是就等於是這樣呢?.
我不會這樣回答你,我說是的,因為為什麼這樣說呢?.
我記得在我以前,我還沒有回來感受之前,我在一間小學堂會,我去做堂主任.
我有很多弟兄姊妹想讀神學,其中有一個要過堂主任那一關.
那我要寫推薦信給你的,我在那間教會十二年裡面的時候.
我很記得我寫了兩封推薦信,你不覺得你是弟兄姊妹讀神學嗎?不是的.
有很多弟兄姊妹,她很興奮的走來跟我說.
我見到上帝怎樣怎樣呼召我讀神學,但是其實我後來去問她.
我說在我平時跟你相處的時候,你的生命當中很多問題你沒有去處理過.
我又不覺得你特別喜歡看聖經,我又不覺得你靈修奠定好.
我又不覺得你跟弟兄姊妹之間的關係是非常表現出來的,很有基督的生命.
然後你告訴我上帝呼召你去讀神學,我是有保留的.

$^{121}$所以很多時候我跟弟兄姊妹說,更加實際的是.
當你想去事奉上帝的時候,請你好好的檢視好自己的生命.
我今天整個生命出來的時候,是建立到其他人,還是我在拆毀其他人.
我跟人與人之間的人際關係裡面的時候.
是我更願意放下自己去成就上帝的工作,還是我堅持己見要去做我自己的工作.
可能這些會更加客觀和精準.
當然我還是那句,我是唐朝人,我不是上帝.
我不可以去攔阻,我不可以告訴你,這個不是上帝的呼召,我不敢這樣說,因為我不是上帝.
但是在我自己作為唐朝人,我要去守這一關的時候,我會很謹慎.
我想強調原來一個格式化的情況,會帶我們每一個弟兄姊妹進入一個危險當中.
剛才我提過,我們應不應該去問上帝的心意?.
要,但是如果你問我的話,我就會更加肯定一件事.
是基於保羅與上帝他自己的關係很密切,而去知道上帝的彼成.
我想大家都應該不會質疑.
保羅所呈現給大家看的時候,聖經是他很熟悉的.
我也相信他的禱告是不會少,他會不斷教導大家,你要堅持禱告,不要灰心禱告.
你會發覺他在裡面的時候,他似乎是很清楚的.
你會發覺你會從很多細微的裡面,你會知道上帝喜歡不喜歡,而你能夠作出這個決定.
所以當我要和你分享我的領受的時候,第一件事,我們怎樣去尋求那是不是神的心意呢?.
第一件事,這件事有沒有明顯地違反聖經的原則?.
聖經其實上帝已經和你說得很清楚,有些是應該,有些是不應該.
我用一個例子,年輕人最喜歡的,鄧炫東說,我最近認識了一個女孩.
你猜這個是不是神的心意呢?.
然後我第一句問他,他相信主了沒有?.
通常問得我就未相信主,相信主是不會問我的.
他很想我給他答案,這個就是神的心意.
然後我問他,他相信主了沒有?他說,未相信主.
那麼上帝在聖經裡面的教導是怎樣的?.
要同父恩厄,是不是?.
我想強調,他也不是同父恩厄,首先在這裡已經很肯定不是神的心意.
所以我就經常套下一步,你是不是不可以喜歡?不是.
你不要不急,不用一開始就拍拖,做朋友才行不行?.
你是很清楚的,因為我們這個信仰有價值,你去帶他回教會.
讓他認識這個,做朋友,讓他認識這個信仰.
讓他相信主,大家真的有同父恩厄,同一個價值觀.
然後你再跟他拍拖行不行?行的.
現在等於問我,鄧炫東,最近有一間公司我想轉工.
你猜是不是神的心意?我經常覺得自己是一個廟祝.
你最好告訴我清楚答案,我又問同一個問題.
你轉了工之後,你是可以更好回教會,還是你沒有時間回教會?.

$^{161}$很有趣,通常也是這樣的,為什麼問我?就是因為影響了回教會.
我就問他,聖經的教導是怎樣?以什麼為優先?以上帝為優先.
你連崇拜都沒有辦法去到,你怎樣以上帝為優先?.
如果你要問我,是不是神的心意?我首先回答,不是.
但倒過來,我可不可以轉工?可以,但你不要問我這是不是神的心意.
你明白嗎?所以我想強調,原來有很多東西是很明顯,究竟是不是聖經原則.
所以學聖經,讀經,靈修是很重要的.
你不熟悉上帝的心意,你怎樣去知道上帝是否喜歡呢?.
第二,究竟我做這件事,我是以上帝為先,還是以你個人為先?.
你問是否神的心意,就是你要先尋求神的國和神的義.
其實這個就要問你自己,我也不能澄清.
我只能帶領弟兄姊妹,一起重新摸一摸自己的心.
究竟我是以什麼為出發?.
我做這件事是想做更多上帝的事,能夠成就.
還是只是想達到我自己的目的?.
第三,當你細心地去尋求過之後,又倒轉了.
第一個,聖經原則沒有問題,第二個,我是以上帝為先的時候.
第三,你就不要怕錯了.
有時候當你尋求完之後,會不會有機會錯呢?會不會呢?會.
保羅這次就是一個很明顯的例子.
他有自己的想法,他想做神的工作,所以他才想向東.
但不等於上帝叫你一定要向東,你明白我的意思嗎?.
上帝會用他的方法去改變你的軌跡.
所以我們只需要繼續謙卑,留意著聖靈的引導.
隨時讓上帝修正你就行了.
會不會錯呢?會錯的.
所以弟兄姊妹不要怕錯.
不要因為怕錯的緣故,站在那裡猶豫不決,原地踏步.
我們要不斷在人生中尋求上帝的方法.
所以你會發覺這段經文展示什麼呢?.
就是無論保羅的計劃是怎樣也好.
只要上帝禁止他,他就會作出改變,配合上帝的計劃.
但這段經文給大家最大的挑戰是什麼呢?.
就不是你不知道神的心意.
而是你知道你是否願意順服去走.
這是我多年來的事奉中,一個很深刻的領袖.
有很多事情你明知上帝不喜歡.
但問題是你還是要走.
你不願意放下自己.
例如在教會的場景中.

$^{201}$紅恩堂有紅恩堂的環境.
佳恩樓有佳恩樓的環境.
錦繡堂有錦繡堂的環境.
其實每一個環境都不同.
上帝是不斷透過新變的東西告訴你.
我要讓你去牧養什麼群體,在做什麼人群.
而我們作為教會的帶領的時候.
或者我們參與弟兄姊妹的時候.
我們願不願意因為上帝告訴我們這些事情.
而改變呢?.
我們願不願意配合上帝你自己的工作去改變呢?.
很多時候是人不願意放下.
我用一個例子.
例如佳恩樓.
我們一直發展的時候,我們去尋求.
我們很多時間去禱告,我們一起去傾談.
到最後我們是一個很清楚的領袖.
所以佳恩樓的核心價值是什麼呢?.
我們要建立一個以家來做核心的崇拜.
什麼叫做家?.
就是一家大小可以一起崇拜.
意思就是在過去我們發覺一件事.
教會在過去叫做分離牧養.
一到崇拜的時候就叫做家散,不是人亡.
家散.
意思是什麼?.
大人大人崇拜,小朋友小朋友崇拜,青年青年崇拜.
大家各散東西.
但我們的領袖不是這樣的.
神的家裡面不是這樣的.
所以我們就發覺一件事.
有些事情我們就會改動.
改動什麼呢?.
其中一個是很不習慣的.
特別是我們有一班弟兄姊妹.
以前是在錦繡,錦繡是一個座堂.
座堂很多時候就是整個環境很寧靜.
你明白我的意思嗎?.
有時候有一點聲音的時候就覺得很大騷擾.
所以我們就會發覺.

$^{241}$其中一個特色就是.
很多教會都會的.
弄一間嬰兒房.
就是將全部有嬰兒的都推到房間裡面.
有聲音都不要妨礙你.
大家討論過你們都知道.
教會都會討論這個問題.
在嬰兒房裡面,那些人是不是崇拜的?.
不是,是講照顧小朋友的.
大家在那裡很吵的.
但不是的.
我們會發覺不是這樣的.
接著我們做了一個舉動.
就是我們想容許小朋友.
這班兒童.
他們可以在敬拜唱詩的時候.
讚美的時候.
我們可以一起去讚美唱詩.
所以我們容許小朋友在那段時間.
大家走回大堂那裡一起去唱.
最初成年人有沒有不習慣?.
當然有.
那些小朋友走來走去.
他們是否正常?.
正常的.
但我想大家留意的是.
當我們這樣做的時候.
我們看到一些很漂亮的圖畫.
就是小朋友看著你們爸爸媽媽.
你們在敬拜的時候.
他們跟著你們.
最奇妙的就是.
之前我們都一起去讀那個.
使徒信經.
有一個小朋友很大力的.
背了整個使徒信經出來.
我們一起舉手祈禱的時候.
我們看到小朋友在前面一起舉手祈禱.
一起唱詩歌.
他們原來是學我們一起去敬拜的.

$^{281}$我們在裡面的時候其實很開心.
我們不介意小朋友在前面走來走去.
我們不介意.
我們不覺得這樣叫混亂.
而是我們覺得在裡面的時候.
我們一起去讚美我們的神.
最近我們升級了一個.
升級版了.
我記得當時我還沒有回來.
在早期的監牢.
有些弟兄姊妹就說.
喂,不行的.
小朋友怎麼可以走上台上去玩樂器.
在那裡搗亂.
而最近我們很感恩.
有些小朋友真的很有恩慈.
打鼓呢.
真的很厲害.
然後我們去邀請他們的爸爸.
跟他們的小朋友一起唱歌.
跟他們的小朋友一起參與敬拜.
讓小朋友打鼓.
是打得又文又老的.
是整個過程裡面.
我們一起,大家都很感恩的.
小朋友的成長.
起不來就是.
應該在一個這樣的環境裡面.
他一直成長嗎.
所以你就會發覺.
原來在裡面的時候.
我們作為大人.
我們需要很多的改變.
我們要配合.
就等於回到監牢.
上帝不知為什麼.
讓我們那裡來了很多說英文的小朋友.
大家平時聊天都是說英文.
那怎麼辦.
難道你叫他改回中文嗎.

$^{321}$不會的.
那我們就求上帝為我們準備.
我們會不會有一個青少年的英文崇拜.
是可以用英文跟他們溝通.
或者叫英文團契.
行不行呢.
我們求上帝為我們準備.
所以弟兄姊妹.
最重要的就是.
究竟我們今天.
知道上帝要你扭轉的時候.
你願不願意放下.
什麼叫順著聖靈而行.
這個我個人領受.
我們在教會裡面.
我們一班的領袖.
我們一班的牧者.
我們一班弟兄姊妹.
可能更多的.
我們不是要開很多的會議.
怎樣做怎樣做.
不是.
而是我們真的更多的時間.
我們需要一起的禱告.
我們去尋求上帝.
對於在我們這間堂會裡面.
上帝啊.
你想我們怎樣.
當你願意去這樣尋求的時候.
弟兄姊妹.
你的生命.
你的眼目會改變的.
為什麼.
因為你就會自然.
你就會看著你附近的社區.
上帝有什麼人.
你交來給我們.
是不是.
當你今天走過.
隔壁的村的時候.

$^{361}$隔壁的洪福村的時候.
你見到很多不同背景.
不同的人.
他們在出現的時候.
上帝啊.
你想我們為他們做些什麼.
我們要做些什麼.
才能夠讓他們來到當中.
是不是.
我想強調的就是.
其實上帝擺每一個地方都不同的.
所以剛才譬如.
樂曲隊.
他們一起獻詩的時候.
其實我就想.
為什麼上帝要將樂曲隊.
加在這裡.
有它的目的.
而可能樂曲隊.
他們都需要去尋求.
除了獻詩之外.
這裡附近.
是不是有一些人.
透過樂曲.
能夠更加去接觸到這個信仰.
我們是不是願意去改變.
我們願意有更多的平台.
去讓我們可以去順服著上帝去走.
所以頂尖很多阻礙了我們更多的.
可能是我們自己的內心.
有很多我們既有的東西.
我們不願意去拆開.
我們不願意去走.
所以我們沒有辦法去改變.
我用一個例子.
這個很爭議性.
譬如很多時候教會的傳統.
是不是.
我們教會.
當然要透過弟兄姊妹的奉獻為一個主導.

$^{401}$主力的我們.
仰望神.
我們這樣叫仰望神.
弟兄姊妹.
我不知道大家有沒有發覺一件事.
其實我們身邊上帝給你的很多的人.
他的恩賜是很獨特的.
在監牢裡面的時候.
我可以告訴弟兄姊妹.
其實我們原來數著數著.
我們有很多弟兄姊妹.
是做營商宣教的.
然後我們最近的時候.
我想說是上帝給了我們很多的奉獻.
舉個例子.
我們在前一段時間.
當疫情停止了聚會.
我們要在網上做一些講座.
有一些.
我不知道大家有沒有聽過.
共享空間.
有沒有聽過.
有一些弟兄姊妹去做共享空間.
做得很好.
然後之後在外面.
原來外面發生很多奇妙的事情.
就是他們在外面.
透過將教會的地方去改變.
讓很多的年青人.
他們有這種發揮的潛能.
然後他們可以在教會.
用教會這個平台.
去建構這些有福音性的東西.
我舉個例子.
譬如一些福音性的商品.
有福音性的可能性的一些服侍.
可能是一些講座.
一些教學.
然後在教會那裡是可以發生的.
而那時候其實我形容是雙贏.

$^{441}$為什麼.
年青人說一句.
他們是.
我想說.
我們回來教會.
我們都是要生活的.
他們都需要去創作的.
而接著的時候.
教會.
我不知道大家.
這裡我不知道怎樣.
我們今年去思考的同樣都是.
監牢很多時候地方都空置了.
我們除了班組之外.
很多地方都沒有用到.
我們今年的目標是.
希望監牢的地方的使用率提高40\%.
我們希望可以更多的地方.
可以讓一些人在裡面.
去做一些有意義的事情.
而去使用上帝的地方.
而我不知道奇妙的地方是什麼.
可以上帝到最後.
被監牢所出現的就是.
我們不是只是單單弟兄姊妹的奉獻.
可以上帝透過弟兄姊妹的恩賜.
而去成為我們去做福音工作.
更多的資源.
而另外一個.
最近我們去探訪一間國際的幼稚園.
很好.
是他們主動聯絡我們.
我想問問大家.
你們身邊有沒有家庭有SEN的小朋友.
大家知道什麼叫SEN吧.
其實有很多的家庭是需要的.
小朋友是很需要的支援的.
而這一對的宣教士.
因為國內他不能做福音工作.
他回流了回香港.

$^{481}$其實他是專門訓練教會的人.
的弟兄姊妹.
只需要你有心.
你就可以去報名.
然後他建議什麼呢.
是可以在你的教會裡面.
開闢一個房間.
專門去服務這些SEN的家庭.
而這個人是可以收費的.
你可以自給自足的.
而令到對方得到這個服務.
而你又可以去傳播.
所以弟兄姊妹.
我想告訴弟兄姊妹.
當上帝聖靈去的工作發生的時候.
只需要你不要去局限你的眼光.
不要局限你自己的思想的時候.
很多事情都可以發生.
因為上帝在過程裡面.
他不斷地告訴我們.
跟我們說我要你是怎樣.
所以你會發覺.
在這裡的時候.
正因為保羅這麼馴服的緣故.
所以後面的經文就出現了一個.
大家都很熟悉.
有些人叫做馬其頓的異象.
或者叫做馬其頓的呼聲.
我剛才說了.
這個聖靈很清楚介入裡面.
完全扭轉了.
改變了保羅的事奉方向.
因為他從來沒有想過要去到那裡.
因為他有一個包袱.
他是猶太人.
他還是法理賽人.
你記住.
他一直在做的時候.
其實他很清楚.
他以為自己仍然要將福音傳給猶太人.

$^{521}$但是上帝不斷地扭轉.
扭轉到現在出現了這個異象.
很清楚有一個馬其頓人.
你過來幫我們.
他就去那裡.
你會發覺聖靈很奇妙.
當你願意的時候.
原來聖靈是很願意.
很樂意介入你的生命當中.
去告訴你.
想告訴你你可以怎樣做.
但是最重要的是.
你有沒有保羅這種這樣的生命.
我想大家留意.
這段經文裡面這樣說.
保羅既看見這異象.
我們隨即就要往馬其頓去.
依為神召我們全方位去那裡的人聽.
既看見.
隨即.
弟兄姊妹.
今天.
你看見了嗎?.
你有沒有隨即?.
今天我看到更多人是很多猶疑.
和考慮.
當你看見一些東西的時候.
我用最簡單的.
譬如你們現在這裡是一個很近.
譬如隔壁的屋村.
很多需要.
我可以肯定很多需要.
家庭很多需要.
婚姻.
剛才我提過.
SEN的爸爸媽媽.
貧窮.
當我們作為一個普通的人.
一看到這些議題我們會怎樣?.
我想告訴弟兄姊妹.

$^{561}$你會退縮的.
我一個小薯仔.
小薯仔.
我做到什麼了?.
弟兄姊妹不是這樣.
我想再強調.
保羅是一個普通人.
保羅和你一樣.
都是一個人.
但是他和你的不同.
就是他願意.
回應上帝所給他的膚質.
他願意一看見就除夕.
弟兄姊妹.
你今天見到身邊有這些需要的時候.
你有沒有去思考.
我可以做些什麼?.
我經常說教會的工作.
不是要打鑼打鼓.
開報道會.
很震撼的一場.
一百個人信耶穌.
其實最近.
我們一些的牧者或領袖.
都不斷思考這個問題.
弟兄姊妹.
你不發覺香港近幾年少了很多大型報道會嗎?.
因為你會發覺一件事.
就是他們都很清楚一件事.
我們只可以將信息給對方.
然後跟不上.
有很多人究竟有沒有回教會.
他不知道.
但是其實在教會裡面.
最紮實的工作是什麼?.
是關係報道.
是你認識這個人.
你跟他建立關係.
你曾經幫助他.
然後他因為你的緣故.

$^{601}$而去認識基督.
所以回到你的人.
剛剛回來的人.
說真一句.
請大家先不要高高在上.
所以你要尋求神.
你不要跟人.
不是的.
你先做好你自己.
他剛剛來到是跟人的.
正常的.
我是看你的頭的.
然後就透過你這種門徒生命.
去帶他成為另一個門徒.
是由你一直帶他出生的.
就好像保羅一樣.
當他呼召提摩太太之後.
你搞清楚.
提摩太太是跟著他一起去周遊列國.
做了很多工作之後.
才將他留下在那間教會.
就不是呼籲提摩太太.
行了,你去吧.
不是的.
他有保羅這個師父.
弟兄姊妹.
你準備好要在神國裡面.
去做神國裡面.
弟兄姊妹的師父了嗎.
當我通常做這種呼召的時候.
我相信弟兄姊妹心裡即時是甚麼.
肅殺.
我這麼差勁.
祈禱又差勁.
靈修又差勁.
聖經又差勁.
弟兄姊妹.
你不要這樣對著全人類說.
因為我經常說.
一年裡面教聖經.

$^{641}$講道理都不少篇.
當日教你們教了十年之後.
你們跟我說我還很差勁.
即是我很差勁.
我會問自己.
每個星期我講崇拜國是甚麼鬼.
但不是的,弟兄姊妹.
請你不要在回應上帝裡面的時候.
跟著你走去妄自菲薄.
上帝從來都不是要你有多大的才能.
上帝要你的是生命.
聖經裡面很多上帝對我們的呼召.
我又聽見主的聲音說.
我可以差遣誰呢.
誰肯為我去呢.
你有沒有好像以賽亞那樣.
主啊,我在這裡,請差遣我.
你有沒有.
我就對他們說.
要收個莊稼多,做工的人少.
所以你們當求莊稼的主.
他發工人出來收他的莊稼.
弟兄姊妹.
你看見莊稼嗎.
你有沒有一起去收.
所以你們要去.
使萬民作我的門徒.
奉父子聖靈的名給他們施洗.
凡我所吩咐你們的都教訓他們遵守.
我就常以你們同在,直到世界的末了.
弟兄姊妹.
上帝的旨意.
聖靈的心意.
很清楚.
問題是.
在於究竟你.
是否願意.
既看見.
我就行動.
弟兄姊妹.

$^{681}$不要小看一個的力量.
我很喜歡看傳記.
從來每一個宗派.
出現的那位領袖.
都不是想做宗派的.
他只是想好好努力在神國裡.
去做一個好的基督徒.
希望好好活他一輩子.
去做上帝的工作.
而上帝就使用他.
我之前在巡禮會.
他們的始創人是約翰.維斯尼.
最初是甚麼?.
就是一個傻瓜.
在一個很混亂的世界裡.
叫了更多的傻瓜.
一起出來的時候.
我們不要同流合污.
我們一起要過聖祭的生活.
我們一起要去禱告.
我們一起要去做神的工作.
神就將更多的就慈給他.
後來就變成了巡道會.
後來一直發展到全世界.
大家知不知道宣道會.
最初不是想做宣道會的.
大家知不知道?.
其實宣道會最初是一個宣教運動.
是我們的始創人孫信博士.
他們看見全世界有這些需要.
然後呼召了一班人.
叫了一班同道人.
我們一起去做這些工作.
然後上帝就不斷將這些東西.
不斷給他.
最後就成為了宣道會.
頂尖妹.
是這樣的.
所以我經常說.
我們不是需要.

$^{721}$一開始的時候有一個很大的宏願.
不是.
頂尖妹我從來都不會鼓勵她這樣想.
而是我只想.
我們可不可以在我們.
生命的有限之年裡.
我們能夠.
很努力地去做上帝的工作.
做好一個基督徒.
就如剛才所說.
我不知道你們有沒有感動.
如果剛才你說對於SEN那些.
你有興趣的話.
你聯絡我.
因為接下來他將會在.
元朗或其他地方.
他會開這些課程.
你只需要肯就可以了.
你願意去服侍就可以了.
會有很多的生命被震撼.
你會見到.
更多更多原來上帝將很多東西.
擺在你身邊.
原來你不是只是.
需要去過.
營營役役.
每天只是上班下班.
等放假.
等旅行的生活.
不是.
而是原來有更多的.
是上帝可以讓你去做的.
所以求神幫助我們.
好嗎.
我們一起去過一個.
順服聖靈而行的生活.
讓上帝讓我們的心靈.
能夠軟化下來.
我們能夠聽候上帝的吩咐.
當上帝你要我做事的時候.

$^{761}$我們更願意放下自己多一點.
去配合上帝的工作.
正如保羅一樣.
我們同心祈禱.
周天父我們求你幫助我們.
孩子在你面前的時候.
我經常都很清楚.
上帝你將我們每一個群體.
放在每一個區域裡面.
你就是要我們在那裡.
做這個區域裡面的工作.
很多時候我們會有很多的考慮.
我們有很多的計算.
但是神啊.
我們更加需要的.
可能我們要留意著上帝.
你告訴我們.
其實你想我們做些什麼.
主啊.
你是會告訴我們.
只要我們先求你的國和你的義.
你就會將我們所需要的一切.
你去加給我們.
主啊.
所以我們求主.
讓我們的眼光放在你的身上.
我們不是去看.
我們有錢沒有錢.
我們不是看.
我們人多人少.
而是我們知道.
只需每個弟兄姊妹.
我們都一起起來的時候.
這個地方會成為一個很美的地方.
主啊.
求主你幫助.
你叫我們更願意尋求你的心意.
我們很願意求主.
你保守住在寒心堂的童宮.
和他們的團隊.

$^{801}$是一齊的.
我們去仰望上帝在你面前.
我們去看見上帝.
你所給我們這裡的意象.
以主我們一起回應這個意象.
求你自己親自的大意.
求你親自的大力幫助我們.
求你祝福著我們.
願意我們成為.
我們每個弟兄姊妹都願意更愛你.
更願意為你服侍的一個基督徒.
我們獻上感恩.
我們多謝讚美你.
其他同學.
主耶穌得勝名至祈求.
阿們.
\newpage



\section{路加福音 19:1-10}
\label{sec:IOHR3FgR_rc}
\textbf{2022.12.25 《生命的呼喚》 路加福音 19\_1-10 (和修版) 講員 : 老耀雄牧師}
\newline
\newline
連結: \href{https://youtube.com/watch?v=IOHR3FgR-rc}{\texttt{https://youtube.com/watch?v=IOHR3FgR-rc}} ~~~~ 語音日期: 2022-12-26
\newline
\newline
\hyperref[sec:pCQk0pIrnYA]{\small{< < < PREV SERMON < < <}}
~
\hyperref[sec:index_chronic]{\small{[返順時目]}}
~
\hyperref[sec:index_scriptual]{\small{[返順卷目]}}
~
\hyperref[sec:jrHnpDp52JU]{\small{> > > NEXT SERMON > > >}}
\newline
\newline
路加福音 19:1-10
\newline
\begin{longtable}{cl}
\hline
\hline
章節 & 經文 (和合本修訂版)\\
\hline
19:1 & \begin{tabularx}{0.7\textwidth}{X} 耶穌進了耶利哥，要從那裡經過。 \end{tabularx} \\ \\ \relax
19:2 & \begin{tabularx}{0.7\textwidth}{X} 有一個人名叫撒該，作稅吏長，是個財主。 \end{tabularx} \\ \\ \relax
19:3 & \begin{tabularx}{0.7\textwidth}{X} 他要看看耶穌是怎樣的人，只因人多，他的身材又矮，所以看不見。 \end{tabularx} \\ \\ \relax
19:4 & \begin{tabularx}{0.7\textwidth}{X} 於是他跑到前頭，爬上桑樹，要看耶穌，因為耶穌要從那裡經過。 \end{tabularx} \\ \\ \relax
19:5 & \begin{tabularx}{0.7\textwidth}{X} 耶穌到了那裡，抬頭一看，對他說：「撒該，快下來！今天我必須住在你家裡。」 \end{tabularx} \\ \\ \relax
19:6 & \begin{tabularx}{0.7\textwidth}{X} 他就急忙下來，歡歡喜喜地接待耶穌。 \end{tabularx} \\ \\ \relax
19:7 & \begin{tabularx}{0.7\textwidth}{X} 眾人看見，都私下議論說：「他竟然到罪人家裡去住宿。」 \end{tabularx} \\ \\ \relax
19:8 & \begin{tabularx}{0.7\textwidth}{X} 撒該站著對主說：「主啊，我把所有的一半給窮人；我若勒索了誰，就還他四倍。」 \end{tabularx} \\ \\ \relax
19:9 & \begin{tabularx}{0.7\textwidth}{X} 耶穌對他說：「今天救恩到了這家，因為他也是亞伯拉罕的子孫。 \end{tabularx} \\ \\ \relax
19:10 & \begin{tabularx}{0.7\textwidth}{X} 人子來是要尋找和拯救失喪的人。」 \end{tabularx} \\ \\
[1ex]
\hline
\hline
\end{longtable}
$^{1}$主席.
各位弟兄姊妹平安.
今日是聖誕節.
很多朋友都知道聖誕節的來源.
就是慶祝主耶穌的降生.
聖經說.
耶穌來到這個世界的目的.
是要去拯救世人.
祂很想世人脫離一切的捆綁.
過一個豐盛的人生.
當然.
對於豐盛人生的定義.
每一個人都有不同的想法.
反而今時今日.
大部份的香港人.
其實都有很多的牽掛.
有很多事情要擔心.
好像我們的生活.
我們的工作.
我們的經濟.
家庭.
兒女.
婚姻.
移民.
很多事情都需要處理或考慮.
跟一個豐盛的人生.
似乎很遙遠.
當我們人生出現困難或逆境.
不知如何是好.
我們可以向誰呼喚呢?.
剛才主席為我們讀出了.
路加福音十九章一至十節.
在裡面有一個生命呼喚的故事.
我們一起去看看.
這個呼喚為主角帶來什麼改變.
我們一起去祈禱開始吧.
至身至善至美的主.
我們深信人類是你所創造的.
生命亦是你所賜的.
你創造我們.

$^{41}$一定有你的心意在當中.
甚願今天你親自向我們每一個說明.
你創造我們的心意.
究竟是什麼.
我們祈禱奉主耶穌基督.
明智祈求.
阿們.
要想生命改變就要鼓起勇氣.
渴求突破.
經文在第一至第四節這樣說.
耶穌進了耶利哥.
要從那裡經過.
有一個人名叫撒該.
作祟利長.
是個財主.
他要看看耶穌是一個什麼人.
只因為人多.
他的身材又矮.
所以看不到.
於是他跑到前頭.
爬上桑樹.
要看耶穌.
因為耶穌要從那裡經過.
四節的經文.
說出了故事主角撒該的身份.
身材.
和行動的目的.
他的身份是一個稅利長.
稅利長的職責是什麼呢.
是負責收稅的.
但是他收猶太人的稅.
比羅馬政府的要求多很多.
就是說.
他利用職權.
去壓榨猶太人而獲利.
以致他能夠成為一個財主.
雖然撒該很有錢.
但是沒有人緣.
是一個很討人厭的猶太人.
撒該也是個矮子.

$^{81}$所以如果有人站在他前面.
他就會被人遮住了.
什麼都看不到.
他為了要清楚看到耶穌.
於是就爬上桑樹.
他很想居高臨下.
高過下面的人群.
這樣他就可以看清楚.
四周的一切.
每一個場景.
綜合這四節經文.
給我們知道一個訊息.
一個有財有勢.
但是沒有人緣的撒該.
他很想見耶穌.
於是就想一個辦法.
去突破自己的限制.
當然不是任何東西.
都能夠驅使人去突破自己.
事實上.
如果追求的目標越大.
越吸引.
就越能夠推動人.
去超越自身的限制.
撒該的目標是什麼呢?.
他很想見耶穌.
見到耶穌.
有什麼吸引力呢?.
原來耶穌來耶利哥城之前.
其實已經有很多人提起過.
耶穌這個名字.
就說耶穌很厲害.
他治好很多人.
甚至連死人都可以治好.
又做過很多神蹟.
還可以在水面行走.
耶穌這個名字.
在猶太這個地方已經被傳開了.
有人說他是先知以利亞.
也有人說他是神的兒子.

$^{121}$究竟是什麼人呢?.
不過無論是什麼人都好.
撒該認為.
耶穌一定是一個滿有能力.
又充滿憐憫.
說話又充滿智慧的人.
耶穌的身份和能力.
對撒該構成了一個.
很大很大的吸引力.
因此撒該為了見一見耶穌.
令撒該想突破.
自己身材矮小的限制.
因此促使他要爬上桑樹.
雖然爬上桑樹看清楚.
是會方便一點.
但也有一個不好處.
它會成為全場的焦點.
一個這麼討人厭的人.
一個這麼不受歡迎的人.
一旦成為公眾注視的要點.
究竟會構成一個什麼心理壓力呢?.
對撒該來說.
他需要突破一些心理上的限制.
才能做到.
所以為了見一見.
這個被稱為神的兒子的耶穌.
撒該需要超越身材矮小的限制.
也需要超越心理上被人注視的限制.
撒該知道.
站得高就望得遠.
想提高自己的視野.
望一望耶穌.
或者望耶穌望得清楚一點.
他總要作出一些突破.
不知道你有沒有這樣的經驗.
世界盃剛剛過了.
如果你在足球場看球賽.
當球員帶著球去到禁區的時候.
大部分的人都會立刻站起來.
為什麼呢?.

$^{161}$一來是情緒因素.
另一個原因.
就是怕被前面的人擋住.
看不到精彩的入球.
如果你是第一次進球場看球賽.
你又不作出這種突破的話.
很乖的坐在球場上.
看足全場的話.
我保證你一定錯過很多精彩的鏡頭.
今天早上.
有很多朋友是第一次來到教會.
各位被人邀請來到教會的朋友.
雖然你們不是被迫.
但是按你的心意.
你們可能喜歡去喝茶.
多於來教會.
所以走進來教會的一座建築物.
對於你們來說.
是一個行為上的突破.
也突破了你們平時不會.
自己一個人走來教會的心態.
大家來教會的心態.
可能和撒哥一樣.
你充滿好奇.
邀請你們來教會的朋友或親人.
可能經常跟你說.
耶穌怎樣好.
教會怎樣好.
來聽聽吧.
給點面子吧.
大家出於禮貌.
好吧.
進來看看吧.
進來聽聽吧.
聽聽看耶穌有什麼特別.
如果你問撒哥.
你看見耶穌之後會怎樣.
我估計他未必能夠即時回答你.
不過我相信.
在撒哥的心靈深處.

$^{201}$一定有些渴求.
否則他不會突破自己.
爬上桑樹成為眾人的焦點.
因為可以和群眾一起迫.
見到耶穌便見到.
見不到便算了.
所以撒哥在心裡.
總有一種渴求.
如果自己的心態和行為都能夠突破.
那麼生命能否也能突破呢?.
能否不再那麼討人厭呢?.
能否過一個被人尊重的生活呢?.
撒哥對生命有所要求.
他充滿了好奇和期待.
他鼓起勇氣.
渴求生命上的突破.
因為他想見的是神的兒子耶穌.
弟兄姊妹.
我相信我們沒有撒哥那麼討人厭.
但我們會否也像撒哥一樣.
我們對生命有要求呢?.
我們對生命有沒有改變和突破呢?.
生命的改變關鍵.
就是回應耶穌無條件的接納.
在第五第六節.
耶穌到了那裡.
抬頭一看對他說.
撒哥快下來.
今天我必須住在你家裡.
他就匆忙下來.
歡歡喜喜地迎接耶穌.
這兩節經文.
講到耶穌到了桑樹之下.
抬頭向撒哥發出了邀請.
茫茫人海之中.
耶穌只是見到桑樹之上的撒哥.
而不是人群中的其他人.
耶穌獨獨見到撒哥.
是否因為撒哥在人群中特別顯眼.
吸引了耶穌的注意呢?.

$^{241}$抑或耶穌見到撒哥內心的世界.
知道他想突破生命的心意呢?.
我想大家在這裡想一想.
耶穌對撒哥的呼喚.
究竟是偶然還是有意的呢?.
你想好了.
記下這個答案.
最後我再和大家揭曉.
耶穌對撒哥說.
快下來.
我必須住在你家裡.
這不是一個普通的邀請.
而是一個生命的呼喚.
「我必須住在你家裡」這句話.
代表了耶穌要進入撒哥的內心世界裡.
無論撒哥以往是好還是壞.
耶穌願意走進他的世界.
也表示對撒哥以往所犯的一切過錯.
他完全的接納.
耶穌接納一個用賄賂手段去買官的撒哥.
耶穌接納一個以欺詐別人去發財的撒哥.
耶穌接納一個內心污穢到極點的撒哥.
耶穌接納一個被同胞杯葛.
生命毫無尊嚴的撒哥.
耶穌接納一個連自己都覺得生命沒有意義的撒哥.
這是生命的接納.
這是父母對兒子的接納.
這是造物主對被創造者的接納.
這是耶穌對撒該的無條件的接納.
這種生命的呼喚和接納.
能夠令人的靈魂深處得到滿足.
因為每一個曾經做錯事的人都需要被接納.
被接納就好像一個初生的嬰兒一樣.
好像被父母抱著一樣安全,溫暖和放心.
經文如何形容撒該對耶穌的回應呢?.
在行動上他急忙的下來.
現代的說法是「那那林」地走下來.
心態上他歡歡喜喜地去接待耶穌.
撒該得到耶穌無條件的接納.
他即時引起其他人的議論.

$^{281}$他們議論什麼呢?.
第七節.
眾人看見都私下議論說.
他竟然到罪人家裡去住宿.
原來猶太人不單止討厭撒該.
還認為他壓榨同胞的行為根本就是罪人.
按猶太人的觀念.
神是不喜歡污衊的.
所以屬神的人都應該遠離污衊的罪人.
因此耶穌不應該去撒該家裡.
更加不應該在家裡過夜.
不過即使撒該有罪.
得罪了所有人.
但卻得到耶穌無條件的接納之後.
生命就即時有改變.
他說:主啊!.
我把所有的一半給窮人.
我若勒索了誰就還他四倍.
首先撒該稱呼耶穌為主.
因此他已經承認耶穌就是神的兒子這個身份.
另外撒該承認自己有罪.
我若勒索了誰就還他四倍.
對於他所犯的錯.
他除了承認之外.
他還願意作出補償.
撒該起初出於好奇.
很想見耶穌.
到現在他不單單認識耶穌.
更加與耶穌建立了關係.
這個結連不是撒該做主動的.
是耶穌做主動.
如果耶穌只是經過桑樹.
與撒該擦身而過.
沒有向撒該發出這個邀請.
就這樣離開了耶利哥城.
我相信撒該的生命是不會有任何改變.
撒該現在心靈得到滿足了.
生命有所改變了.
他的價值觀也跟著改變了.
他認為人生意義.

$^{321}$以及與耶穌的關係.
比錢重要.
所以他承諾將身家的一半給窮人.
曾經勒索過的人.
就四倍償還.
撒該以往是用錢來買官的.
他重視金錢.
但今天他卻視錢財為糞土.
他能夠突破金錢對他的捆綁和限制.
各位.
如果我們願意改變.
我們也可以.
當一個人的生命被神接納.
得到赦免和寬恕.
又找到生命的意義的時候.
我們就不再需要物質上帶給我們的刺激和滿足.
而只需要屬靈上的滿足.
以及與耶穌有一個很親密的關係.
財富,名利,地位.
再不是我們終極追求的目標.
也不再是我們生命上的捆綁和束縛.
撒該的改變還影響了其他人.
因為第九節耶穌跟他說.
今天救恩到了者家.
因為他也是亞伯拉罕的子孫.
耶穌應許撒該得到救恩之外.
他整個的家庭都得到.
撒該的家族都因為他的迴轉而蒙上祝福.
一個人得福.
令到整個家族都蒙福.
經文的最後一節.
說出了耶穌進入耶利哥城的目的.
也解答了耶穌經過桑樹之下.
對撒該的呼喚究竟是偶然的.
還是有意的.
第十節.
耶穌說.
人子來是要尋找和拯救失喪的人.
耶穌要尋找人.
尋找那些生命想改變的人.

$^{361}$耶穌進入耶利哥.
就是要尋找生命想改變的人.
耶穌要拯救人.
祂要拯救那些想改變.
但是又改變不了的人.
耶穌要拯救人.
祂要拯救那些被罪捆綁.
讓他們得到釋放.
耶穌要拯救人.
祂要拯救那些承認自己不行.
但卻相信純可以將自己變為「得」的人.
耶穌對撒該的呼喚絕對不是偶然.
我自己回想自己被主耶穌呼喚的經過.
也不是偶然.
我自己出生於一個信奉民間宗教的家庭.
我住在黃大仙.
家裡對面就是黃大仙廟.
所以從小就被契給黃大仙.
家裡有一個大神台.
有祖先,地主,黃大仙.
還有很多其他的神明.
逢初一十五.
我家裡就會做節.
燒香還神.
不過我自己拜還拜.
但不相信祂可以幫到我.
認為祂只是心靈的寄託.
我自己和很多香港人一樣.
凡事都向錢看.
其實我是一個很功利的人.
出來工作幾年就申請做警察.
我不是因為有理想.
也不是想要徐步安良.
我沒有那麼偉大.
只不過是因為當時紀律部隊中.
警察的工資是最高的.
其實我是貪錢.
工資高.
一個年輕人.
二十多歲.

$^{401}$進了一個大染缸.
當時沒有一個正確的人生觀.
自然就很容易被污染.
應該這麼說.
嫖賭喝當吹.
吸煙喝酒賭錢過大海.
真的每樣都精.
每樣都好.
如是者做了警察五年.
我覺得自己都很敗壞.
又覺得自己很污穢.
有時我想.
我的本質都不是那麼壞.
我都很想脫離這種.
既荒唐又頹廢的生活.
不過很可惜.
我想做的就做不到.
我不想做的卻偏偏去做.
好像很想戒賭.
但當同事說三缺一.
自然會坐在麻將台上.
好像想戒煙.
試過十多次.
都是失敗.
甚至想戒酒.
但總是杯不離手.
後來我知道.
這種叫做罪中之律.
雖然知道不對.
但很享受過中的樂趣.
跟今天的人打機一樣.
知道不好.
但卻不可以放棄.
生命有一種無力感.
但卻無法擺脫.
有一次跟一個中學同學聊天.
說出了這種感受.
他是一個基督徒.
於是帶我回教會.
我都很抗拒.

$^{441}$一跟你說完就帶我回教會.
不過算了.
可能跟你一樣.
當作給面子派對.
心想.
回一次就算了.
當作給面子.
誰知他那次帶我.
在一個報道會.
就好像今天一樣.
他說完說.
罪可以得赦免.
生命可以重新開始.
記不記得一個戒毒廣告.
生命有第二次.
這個訊息很吸引.
生命真的可以再開始.
我可以擺脫這種.
又荒唐又頹廢的生活.
我初初以為聽完就算了.
誰知我講完在台上呼召.
他說.
如果生命想改變.
真的很想活出真我的話.
你就決志信耶穌吧.
我心想.
這麼快就叫我信耶穌.
我死都不回應.
但是神在我心裡.
一直提醒我.
他說這種活出真我的生活.
不是你想要的嗎.
你想要的.
為什麼你不回應呢.
我心裡很掙扎.
很瘀傷.
真的很瘀傷.
怎麼回應呢.
神在我心裡一直跟我說.
催逼.

$^{481}$聖靈對我催逼.
我不知道今天聖靈有沒有對你催逼.
當神第三次向我發出呼喚的時候.
我就徹底降服了.
我就舉手相信.
當然.
我信了耶穌.
不是一切都一帆風順.
我的同事不信我.
真心相信耶穌.
即是警察的同事不信.
他們認為我進教會是為了追求女孩.
你信什麼耶穌.
你追求女孩而已.
我當時謝飯祈禱.
他們把我的菜都拿走了.
他們說你信耶穌嗎.
叫耶穌把飯送給你.
不過我真正的挑戰不是被人笑.
或者我真正的挑戰不是不被人認同.
而是我是否真的可以藉著耶穌.
改變到我的心態和行為.
我記得我信了耶穌.
當時回教會.
我當時一天吃兩包煙.
如果打通宵麻雀就三包.
我回教會崇拜.
崇拜完之後.
要馬上走到後樓梯.
吃兩支煙.
真的不行.
好像煙引起.
你想想.
去完崇拜.
走到後樓梯吃兩支煙.
再回來上主日學.
所有人都知道你吃了煙.
不過沒有人揭穿我.
會友的心真好.
沒有人揭穿我.

$^{521}$但我心裡想這樣不行.
是吧.
不過我就是做不到.
就是戒不了.
你想想多愚蠢.
自己騙自己.
掩耳盜鈴.
一早回來崇拜.
然後走到後樓梯吃兩支煙.
再回主日學.
後來我祈禱.
真的祈禱.
我說耶穌你幫我吧.
我不可以這樣.
求你幫我戒煙.
幫我戒賭.
幫我戒酒.
幫我遠離這些人有.
耶穌聽禱告.
耶穌真的聽禱告.
我真的戒了煙戒了賭戒了酒.
也沒有被其他同事影響到.
遠離了這班同事.
我穩定了教會的聚會.
還有教會受洗.
生命得到改變.
這是我第一次.
倚靠耶穌.
去過我那個得勝的生活.
當然這是一個開始.
接著我的婚姻生活.
一塌糊塗.
我和我太太說.
不是我跟別人分享.
我的婚姻生活就是這樣困獸鬥.
困獸鬥.
沒得走.
如果不是耶穌介入在我的婚姻生活裡.
我和我太太早就離婚了.
沒有耶穌介入在我的婚姻生活.

$^{561}$我們不會今天走到三十年.
後來我轉工.
去做地產.
又回到那種.
升職犬馬的生活.
我太太當然很焦急.
不斷為我祈禱.
不過我是一個很頑劣的人.
很頑劣的人.
直到耶穌出手管教我.
真的將我趕到絕路.
然後我才會回轉.
信耶穌不是一帆風順.
但你向祂呼求.
祂一定會幫你.
你願意信服祂的帶領.
祂一定會幫你走一條易路.
就是這樣.
經歷過耶穌的帶領.
我的婚姻.
我的工作.
最後我太太開始讀神學.
她感染了我.
讓我可以一起進入神學院.
一起讀書.
今天一起事奉.
我不能夠說.
我現在過的是一個豐盛的人生.
我不敢說.
不知道會否像我太太說.
最終是衰收尾.
在我事奉的最終階段裡.
會軟弱跌倒.
但我仍然覺得.
只要我們跟著耶穌.
緊緊跟隨祂.
按著祂的心意而行.
這是一個很寶貴的.
生命的恩典.
各位朋友.

$^{601}$今天你們來到教會絕對不是一個偶然.
我們今天可以和一個至高至善.
至美的神相遇.
必定不是偶然.
如果我們心裡很想生命能夠.
突破.
很想突破一些困難和限制.
或者很想脫離.
一些束縛,捆綁.
或者引誘.
我們就需要接受耶穌.
這個無條件的接納.
以及耶穌的應許.
以至我們的生命.
才能夠改變.
而且必然會改變.
各位朋友.
我給大家一個安靜的時間.
你思考一下以下的問題.
你是否想.
總有一個新的生命.
一個被神去改變的生命.
不是被你自己改變.
你改變不了你自己.
你被神去改變的生命.
你是否想生命有一個突破.
讓過往對你的困擾.
束縛.
可以一掃而空.
你願不願意.
有這樣的改變.
這是你個人的決定.
是和耶穌之間的決定.
洪恩嘉的弟兄姊妹.
在這一刻.
請你為未認識耶穌的朋友.
你帶來的朋友.
切切為他代禱.
他們需要你在屬靈上的支援.
我們大家一起閉上眼.

$^{641}$安靜在座位裡.
請詩琴為我們伴奏一節.
詩歌.
詩歌:《聖經》.
我現在奉耶穌的名.
向在座未認識耶穌的朋友.
發出邀請.
這個邀請不是偶然的.
因為耶穌正在找你.
你雖然在人群當中.
但耶穌看到你.
就好像耶穌看到撒該一樣.
耶穌說:你出來吧.
今天我要進入你的生命裡.
令你得到安慰.
心靈得到滿足.
耶穌說:你出來吧.
我要讓你聚得赦免.
令你展開一個新的人生.
如果你願意回應耶穌對你的呼召.
接受耶穌無條件的接納和寬恕的話.
你願意相信耶穌.
請你在座位裡面舉一舉手.
讓耶穌知道你的回應.
你願不願意回應耶穌對你的呼喚.
請你舉一舉手.
未信的朋友.
你今天處身在教會裡面.
絕對不是偶然.
請你在心裡問你自己.
你是否想耶穌這個邀請.
是否想耶穌這個邀請.
就這樣擦身而過.
你不去回應祂呢?.
耶穌在等著你.
請你舉手回應耶穌.
撒該和耶穌相遇後.
他心裡充滿了喜樂.
如果你願意接受耶穌做你的救主.
你也會像撒該一樣.

$^{681}$心裡充滿了喜樂.
撒該渴求生命的突破.
朋友.
你也需要鼓起勇氣.
舉起你的手.
親自去回應耶穌.
耶穌再一次去呼喚你.
你願意日後不再依靠自己.
轉而依靠耶穌嗎?.
地上的父親,母親.
怎樣去愛孩子.
地上的父親願意承擔你的痛苦.
除去你一切的哀傷.
你願意進入天父的懷抱.
讓天上的父親擁抱你嗎?.
請你舉起手.
接受天父給你的祈福吧.
如果你以往曾經信過耶穌.
但生命從來沒有改變.
你想這一刻再次向耶穌說.
我願意承認你是我生命的主.
是我拯救的主.
我願意再次跟隨你.
一生的依靠你.
如果你有這個立志.
我都邀請你舉起手.
親自去回應耶穌.
縱然你以往已經缺過志了.
但我願意再次向你.
向你發出邀請.
然後再跟決志的朋友.
祝福你.
耶穌向你發出生命的呼喚.
一個有限的生命.
能夠與無限的神連結.
我們可以與創天造地的神.
有一個關係.
生命能夠回歸到創造者的手中.
讓生命可以找到真正的意義.
你想不想知道.

$^{721}$自己生命的意義是什麼?.
讓耶穌向你.
揭示生命的奧秘.
請你勇敢地舉起手.
親自回應生命的主.
向你發出的呼喚.
有些朋友很願意踏上信心的一步.
在眾人面前宣告.
他們願意相信你.
願主繼續帶領.
我們每一個在前面的路.
兼顧我們每一個信心.
請剛才有舉過手.
決志相信耶穌的朋友.
或者你心裡願意相信耶穌的朋友.
你可不可以走上台前?.
會不會有人願意相信耶穌?.
你能夠來到台前.
我們會為大家獻上祝福.
歡迎你.
除了這位朋友之外.
有沒有誰願意相信耶穌?.
歡迎你.
你相信耶穌的.
請你走上台前.
如果有人陪他.
陪他的朋友都可以陪他走上台前.
如果你認識他.
你和他一起.
給他一些鼓勵.
站在他身邊.
還有沒有朋友.
你是願意相信耶穌的?.
好.
我們一起有個祈禱.
親愛的主.
求你建立每一位.
願意決志相信你的朋友.
讓他們一次的相信.
永遠都相信.

$^{761}$求你用寶血.
潔淨他們的罪.
赦免他們以往所犯的過錯.
讓他們全然的聖潔.
他們能夠坦然無懼的來到.
施恩的寶座之前.
領受你的恩典.
讓他們的生命.
因著相信和依靠你而蒙福.
他們的人生.
因著你的應許.
得到永恆.
和豐盛的生命.
從此不再一樣.
甚願你的恩恩惠.
慈愛.
尚與他們同在.
我們一起祈禱.
奉教主耶穌基督得勝名字.
祈求阿們.
好,教會有一些陪談員.
和一些決志者做一些信仰澄清.
請陪談員帶同這些決志者.
去2037.
如果你有朋友是決志的.
你可以陪他一起去.
你都可以去2037陪著你的朋友.
又是有信仰的澄清.
好.
我們將榮耀歸一比主.
願主祝福大家.
謝謝大家.
\newpage



\section{約翰福音 1:1-13}
\label{sec:jrHnpDp52JU}
\textbf{2023.01.01 《擁抱那非一般的道》 約翰福音 1\_1-13 (和修版) 講員 : 曾敏傳道}
\newline
\newline
連結: \href{https://youtube.com/watch?v=jrHnpDp52JU}{\texttt{https://youtube.com/watch?v=jrHnpDp52JU}} ~~~~ 語音日期: 2023-01-01
\newline
\newline
\hyperref[sec:IOHR3FgR_rc]{\small{< < < PREV SERMON < < <}}
~
\hyperref[sec:index_chronic]{\small{[返順時目]}}
~
\hyperref[sec:index_scriptual]{\small{[返順卷目]}}
~
\hyperref[sec:mxoNZ6Tp6YY]{\small{> > > NEXT SERMON > > >}}
\newline
\newline
約翰福音 1:1-13
\newline
\begin{longtable}{cl}
\hline
\hline
章節 & 經文 (和合本修訂版)\\
\hline
1:1 & \begin{tabularx}{0.7\textwidth}{X} 太初有道，道與神同在，道就是神。 \end{tabularx} \\ \\ \relax
1:2 & \begin{tabularx}{0.7\textwidth}{X} 這道太初與神同在。 \end{tabularx} \\ \\ \relax
1:3 & \begin{tabularx}{0.7\textwidth}{X} 萬物都是藉著他造的，沒有一樣不是藉著他造的。凡被造的， \end{tabularx} \\ \\ \relax
1:4 & \begin{tabularx}{0.7\textwidth}{X} 在他裡面有生命，這生命就是人的光。 \end{tabularx} \\ \\ \relax
1:5 & \begin{tabularx}{0.7\textwidth}{X} 光照在黑暗裡，黑暗卻沒有勝過光。 \end{tabularx} \\ \\ \relax
1:6 & \begin{tabularx}{0.7\textwidth}{X} 有一個人，是從神那裡差來的，名叫約翰。 \end{tabularx} \\ \\ \relax
1:7 & \begin{tabularx}{0.7\textwidth}{X} 這人來是為了作見證，是為那光作見證，要使眾人藉著他而信。 \end{tabularx} \\ \\ \relax
1:8 & \begin{tabularx}{0.7\textwidth}{X} 他不是那光，而是要為那光作見證。 \end{tabularx} \\ \\ \relax
1:9 & \begin{tabularx}{0.7\textwidth}{X} 那光是真光，來到世上，照亮所有的人。 \end{tabularx} \\ \\ \relax
1:10 & \begin{tabularx}{0.7\textwidth}{X} 他在世界，世界是藉著他造的，世界卻不認識他。 \end{tabularx} \\ \\ \relax
1:11 & \begin{tabularx}{0.7\textwidth}{X} 他來到自己的地方，自己的人並不接納他。 \end{tabularx} \\ \\ \relax
1:12 & \begin{tabularx}{0.7\textwidth}{X} 凡接納他的，就是信他名的人，他就賜他們權柄作神的兒女。 \end{tabularx} \\ \\ \relax
1:13 & \begin{tabularx}{0.7\textwidth}{X} 這些人不是從血生的，不是從情慾生的，也不是從人的意願生的，而是從神生的。 \end{tabularx} \\ \\
[1ex]
\hline
\hline
\end{longtable}
$^{1}$開始的時候我們再一次一同禱告.
在新的一年仍然賜給我們這樣的福氣.
有這樣的恩典.
可以回到你的家裡去敬拜你.
很願意你對我們眾人說話.
願意聖靈充滿我們每一個的心靈.
讓我們明白神所要對我們說話.
也讓我們有很悉切的回應.
我們將這一堂道傳言交在神的手中.
願神顯出你的榮耀.
奉耶穌基督的名字祈禱.
阿們.
今天講這一堂道.
我覺得非常寶貴.
是《擁抱那非一般的道》.
大家記不記得新一年教會的主題.
我們一起讀一次.
預備起.
學務真理.
遵行主道.
今天首先要講.
《擁抱那非一般的道》究竟是一回事.
首先要解釋一下道是甚麼.
道是甚麼意思呢.
其實道的意思本來是非常豐富的.
有分開外在的意思.
有內在的意思.
內在的可以說是我們的思想.
理性.
而外在的表達.
就是說話.
是訊息.
原來這些是道.
可以很立體地看道.
不是很平面的純粹文字.
不是很小的文字.
是很立體的.
原來有思想的.
原來有說話有訊息的.
道是這麼豐富的.

$^{41}$我們現在說的是說話和訊息.
是誰的說話和訊息呢.
我們每個人都在說話.
就算現在我也在說話.
剛才細熊也在說話.
大家在台下也在說話.
有很多人的說話.
也可以發放不同的訊息.
是否一般人的訊息和說話呢.
不是.
這裡不是一般的說話和訊息.
而是神的話語.
是神的訊息.
為什麼這樣說呢.
我們看下這段經文.
我們再看一章一字的經文.
太初有道.
道與神同在.
道就是神.
太初有道.
這個道在太初的時候已經在這裡.
太初是萬有開始的時候.
是宇宙開始的時候.
原來在這個之前.
在這個之前.
道已經存在.
所以很明顯這個道不是我說的話.
是你說的話.
它是神的話語.
是神的訊息.
很久很久已經存在.
還有很特別.
世界是藉著它而來.
因為他說.
一張衣紙萬物都是藉著它而造的.
沒有一樣不是藉著它而造的.
原來在上帝創造世界的時候.
道是有份的.
世界是藉著它而來的.
原來它是一個執行者.

$^{81}$在上帝的創造工作裡.
它是執行的一個.
沒有它沒有世界.
它是世界的起源.
有它才有世界.
世界裡每一樣東西都是藉著它而有的.
所以這個道非比尋常.
這個道這個訊息.
和上帝的作為有很緊密的關係.
剛才說的是創造的大能.
一會兒我們會看到其他的大能在當中.
這個道與神的關係非常緊密.
剛才說一章一節道與神同在.
一章二節道太初與神同在.
在世界在宇宙未曾存在之前.
道已經與神一起.
重複了兩次.
是說關係緊密.
它是神的同伴.
它與神同行同伴.
已經是由世界在宇宙開始之前.
那個關係已經是這樣.
所以是非常非常緊密.
可能與神一起.
與神緊密也可以是其他東西.
經文很強調.
道就是神.
不是一個普通的人.
不是一個物件.
道自己就是神.
道就是神.
道與神同行.
它是神的同伴.
它自己就是神.
就是這麼緊密.
絕對不能當它是普通的說話.
普通的訊息.
它真的非一般.
在這個過程裡.
我們會看看這個道是什麼呢?.

$^{121}$這個道在道裡有生命.
這個生命也不是一般的生命.
是屬靈的生命.
是一個可以復活的生命.
很有權能的生命.
這個生命在人裡成為人的光.
那光又是什麼呢?.
光是在說啟示.
很特別.
這個道在人裡.
比我們經歷屬靈的生命.
是很特別的.
這個道本身有屬靈的生命.
有復活的生命.
有權能的.
這個生命進入我們裡面.
就成為我們的光.
對我們有很大的影響.
很多人對著道的光拒絕.
為什麼呢?.
這裡很特別.
我們看看光暗的分別.
大家看看經文裡.
我們回到一章四節.
一章五節.
一章四節說在他裡面有生命.
前面再說.
被造的在他裡面有生命.
這生命就是人的光.
光照在黑暗裡.
黑暗卻沒有勝過光.
光和暗有很大的分別.
暗不只是說沒有光.
黑暗不只是沒有光.
其實黑暗是在說邪惡.
活逆上帝.
原來世界是黑暗的.
因為世界充滿了罪惡.
充滿了我們人的罪惡.
世界裡的人也是活逆神.

$^{161}$所以世界是黑暗的.
很特別.
和光有很大的分別.
當光出現的時候.
它會照在黑暗裡.
黑暗不能勝過光.
但是光卻可以照在黑暗裡.
在光暗的分別裡.
人對光有不同的反應.
有很多人拒絕.
選擇是拒絕的.
為什麼呢?.
因為人知道自己的行為是惡的.
不願意讓人看到自己的行為是惡的.
所以拒絕光.
但也有人接納.
會接納這個光.
會相信這個道.
我們看回約翰福音三章十九二十節.
大家聽我讀就行了.
光來到世上.
世人因自己的行為是惡的.
不愛光 倒愛黑暗.
這就定了他們的罪.
二十節.
犯作惡的人都恨胡光.
不來接近光.
恐怕他的行為被暴露.
但實行真理的人.
就來接近光.
為要顯明他的行為是靠神而行的.
所以人對光有不同的反應.
是因為堅持在黑暗裡的人.
不想光照到自己.
是在罪惡當中.
但是接受光的人.
卻願意在光裡生活.
跟著光而行.
生命是被改變的.
是很不同的.

$^{201}$所以大家留意到.
這個道是多麼的特別.
有誰的說話.
可以在別人的生命裡成為光呢.
有誰的說話.
可以在別人的生命裡有這麼大的影響呢.
唯有上帝的說話.
唯有神自己的訊息.
所以這個不是一般的道.
不是一般的信息.
這個光對世界帶來很大的震撼.
將人分開了.
接受光的和不接受光的分開了.
這個光影響世界.
在這個道裡有生命.
這個生命絕對影響我和你的生命.
影響全世界人的生命.
也分開了接受的人.
和不接受的人.
我們的生命會有不同的結局.
這個道到今天.
在經文裡沒有跟大家一同讀出來.
是一章十四節直接揭曉.
這個道是否只是與神同在.
這個就是神.
其實這個道是什麼呢.
大家聽我讀.
一章十四節這樣說.
道成肉身 住在我們中間.
充充滿滿地 有恩典 有真理.
我們也見過他的榮光.
正是父獨一兒子的榮光.
約翰為他作見證.
哭著說 一路來.
原來這個道是誰.
大家告訴我.
耶穌基督.
去到這裡進一步揭曉了.
道成肉身.
就是耶穌基督.

$^{241}$原來耶穌基督.
就是神的話語.
就是神的訊息.
在耶穌基督裡.
有屬靈的生命.
有復活的生命.
耶穌基督進入到我們生命裡.
就會成為我們的光.
帶給我們很重要的啟示.
一章十三節這樣說.
十二節很重要.
凡接納他的.
就是信他明的人.
他就賜他們權柄.
作神的兒女.
這些人不是從血生的.
不是從情慾生的.
也不是從人的意願生的.
而是從神生的.
接納這個光的.
就可以成為神的兒女.
相信耶穌基督的.
就會成為神的兒女.
所以這個道.
不是一般的道.
這個道帶我們回到天父身邊.
讓我們和天父復和.
享有永遠的生命.
永遠和上帝在一起.
這個不是一般的道.
十四節說.
就是耶穌基督.
耶穌基督就是道.
就是上帝的訊息.
上帝的說話.
在這裡值得我們去看.
今天我和你手上.
是不是也有一本聖經.
就算有些人.
今天會在手機裡看聖經.

$^{281}$也是聖經.
有些人不看.
是聽聖經.
你聽的也是聖經.
那句聖經就是神的話語.
你有沒有想過有什麼特別.
為什麼要說道成了肉身.
太初有道說這麼多.
道就是神的話.
神的訊息.
道就是耶穌基督.
我們這本聖經就是神的話語.
我們這本聖經就是神的訊息.
所以你看聖經的時候.
看到什麼.
不敢回答我.
我剛才說.
耶穌基督就是道.
是神的話語.
是神的訊息.
不是我和你的話語.
耶穌基督是神的話語.
是神的訊息.
這本聖經就是神的話語.
是神的訊息.
所以我們看這本聖經的時候.
在看什麼.
在看耶穌基督.
你是不是在看一堆文字.
不是在看一堆過時的文字.
是很多年前寫的.
和我和你的關係.
耶穌基督的時候有手機嗎.
有電腦嗎.
有高科技嗎.
什麼都沒有.
這些東西都歪了.
和我什麼關係.
我肯定現在我們的職業.
在耶穌基督的時候是沒有的.

$^{321}$看就不要看.
很悶啊.
不知道在說什麼.
最重要是有需要的時候找人聊.
我最重要是有朋友.
有朋友什麼都可以.
這本聖經從來都沒有過時.
因為不是一堆文字.
這麼簡單.
耶穌基督自己.
在聖經裡面.
你應該會見到光.
見到生命.
聖經帶給我們的經歷.
應該可以經歷到耶穌基督的生命.
是見到光的.
行向天父.
行向永恆.
是不一樣的.
這不是一般的華語.
一般的信息.
這不是一般的信息.
這不是一般的華語.
一般的信息.
是耶穌基督自己.
你有沒有這樣想過呢.
每一天你最怕拿起這本聖經.
最喜歡打開電視機.
七八個小時.
看著看著.
看到快要收的時候.
好像電視送飯.
很閒的一天.
三餐都對著電視.
我肯定你看聖經有時間.
但電視機裡面有沒有生命呢.
你說有.
那些人活著的.
那些生命影響你的生命.
那些生命成為你的光嗎.

$^{361}$這個世代這麼黑暗.
你怎樣才能找到光呢.
就是上帝的說話.
上帝的信息在耶穌基督裡面.
你才會看得到.
才會有那道光.
你再看的時候不一樣.
但我們不願意.
我們不願意花時間看聖經.
也覺得很陌生.
我今天一定不會在這裡問問題.
誰看完這四福音書一次.
誰看完舊約聖經一次.
看完新約聖經一次.
因為我們相信.
我們實在太少看聖經了.
無論是手機版的.
還是一本本的.
硬盤的.
聽的都很少.
實在太少了.
我相信現在如果舉辦聖經問答比賽.
我們真的會很害怕.
很多名字都不認識.
不知道是誰.
什麼事件也不知道.
馬上衝出門口.
這本聖經對我們來說太陌生了.
你有沒有知道.
你看著它.
看著耶穌基督.
就是那個拯救我們生命的那個.
帶我們進入永恆的那個.
讓我們在黑暗世代看到光的那個.
上帝的話語.
剛才其實說了一樣東西.
是和上帝大能的作為有關的.
在這個世界.
在這個宇宙未曾出現之前.
道已經存在.

$^{401}$耶穌基督已經在那裡.
耶穌基督參與創造的工作.
萬物是藉著耶穌基督而來的.
今天我可以這樣說.
有任何困難.
都可以和耶穌基督說.
因為是祂創造世界.
祂知道這個世界的一切.
不要說因為這個世界已經過了很多年.
宇宙過了很多年.
變化了很多.
上帝你不認識我們了.
不是的 對不起.
上帝超越萬有.
直到今天祂都認識我們的一切.
你不找祂找誰.
這個道.
你不天天親近祂.
你親近誰.
難怪我們常常感到沮喪.
痛苦 失意 懼怕.
因為根本和道的關係非常遙遠.
上帝的話語不單止和創造有關.
耶穌基督不單止參與創造的工作.
上帝還參與什麼呢?.
拯救工作.
醫治工作.
保護工作.
耶穌基督拯救我們 醫治我們 保護我們.
有耶穌基督就是最寶貴的.
這個道原來有這麼大的功用.
顯出上帝自己的大能.
這個道的存在就是宣告上帝自己的大能.
今天我和你有沒有見識到這個道的大能呢?.
有沒有見識到?.
如果見識到 經歷到.
為什麼還是懼怕 灰心 失意 沮喪.
還是願讀.
每天糾纏在是非裡呢?.
這是值得我們問的.

$^{441}$我們講的是非還是神的話語呢?.
我記得有一次有個姐妹和我說起.
聽講道的回應.
我心裡真的很雀躍 很開心.
有人和我說起神的話語 我真的覺得很開心.
神的話語對她的影響 我真的很久沒有聽過這些分享.
通常傳道人很單向地說.
這段話是這樣 那段話是這樣的.
耶穌基督是這樣 那裡是這樣.
我很少聽回頭你怎樣經歷上帝 怎樣經歷耶穌基督.
你認識到耶穌基督的寶貴 我很少這樣聽回頭.
但其實道就是這麼寶貴 在我們生活裡 在我們生命裡.
耶穌基督就是這麼寶貴 在我們生命裡 與我們同行.
是很寶貴的 我們應該有很多經歷.
回來教會是經歷上帝.
很期望在新的一年 是洪恩嘉.
很渴望在這裡 說上帝的話 多於說是非.
是非沒什麼好說的 是不是 黑暗就是黑暗.
不用不停證明 怎樣黑暗 怎樣壞.
說完一個又說第二個 第三個第四個.
他說黑暗就是黑暗 交給上帝.
你和我一起說上帝的話 互相鼓勵 看著神.
看著洞 看著耶穌基督.
讓這個光大大的在洪恩嘉裡發出來.
這個光夠強的時候 才可以讓洪水橋區的人看到.
哇 進到教會有光 見到光 想進教會 因為有光.
不是的 如果我們裡面很黑暗 別人夠不夠膽進來.
沒有 這裡裡面都黑暗 不要進來 是不是.
讓耶穌基督在我和你的生命裡發光.
好好讀聖經 每天和道親近 和耶穌基督親近.
看著他 好好的親近.
你再去看的時候 不要再覺得一堆文字.
沒有意義的文字 過時的文字.
在這裡 我很想再一次和大家讀今年的年曆.
我希望這個不是成為我們的口號 是成為我們的行動.
帶著這份熱誠去看聖經.
如果碰了塵 好好地拿出來.
好好地拿出來看 看手機也可以 聽也可以.
有不少弟兄姊妹覺得未必懂得字 沒關係.
現在有很多方便的 有些是聽聖經的三機.

$^{481}$網上也有一些apps是聽聖經的 聽吧.
很快的 幾分鐘就有一張了.
所以我相信如果一天聽完整個約書記 一點難度都沒有.
你700小時看電視 拿一個小時也夠了.
我這樣說 不是每個人都這樣 拿700小時看電視.
我們一同禱告.
天父上帝 我們在這裡要讚美你.
讚美你把你的愛子耶穌基督賜給我們.
愛子耶穌基督是何等的寶貴.
在宇宙萬物未存在之先.
耶穌基督就在那裡和父神一同.
還有寶貴的聖靈 三一的真神.
在世界未存在之先 你已經存在.
你創造世界.
你拯救 醫治 保護 看顧 帶領你的百姓們.
給我們有數不盡的福氣 恩典.
你是配得榮耀和稱頌的.
特別要讚美耶穌基督 就是道的本身.
主耶穌在你裡面 有贖靈的生命 有復活的生命.
有權能的生命.
當你在我們的生命裡 就成為我們的光.
願主大大感動我們眾弟妹的心.
讓我們很渴望去親近你.
每天很渴望看聖經.
渴望和耶穌基督建立一個很親密的關係.
願主就以你的華語去餵養我們.
去建立我們 去牧養我們的生命.
讓我們成長成熟.
願意耶穌基督的光在我們生命裡發出來.
照亮這個黑暗的世代.
讓洪恩嘉真正在洪水橋發光.
靠主賜福保守.
奉耶穌基督的名字祈禱.
阿們.
\newpage



\section{哥林多前書 12:4-27}
\label{sec:mxoNZ6Tp6YY}
\textbf{2023.01.15 《多元中的合一》哥林多前書12\_4-27 (和修版) 趙崇明博士}
\newline
\newline
連結: \href{https://youtube.com/watch?v=mxoNZ6Tp6YY}{\texttt{https://youtube.com/watch?v=mxoNZ6Tp6YY}} ~~~~ 語音日期: 2023-01-19
\newline
\newline
\hyperref[sec:jrHnpDp52JU]{\small{< < < PREV SERMON < < <}}
~
\hyperref[sec:index_chronic]{\small{[返順時目]}}
~
\hyperref[sec:index_scriptual]{\small{[返順卷目]}}
~
\hyperref[sec:8tGomcGt7oY]{\small{> > > NEXT SERMON > > >}}
\newline
\newline
哥林多前書 12:4-27
\newline
\begin{longtable}{cl}
\hline
\hline
章節 & 經文 (和合本修訂版)\\
\hline
12:4 & \begin{tabularx}{0.7\textwidth}{X} 恩賜有許多種，卻是同一位聖靈所賜。 \end{tabularx} \\ \\ \relax
12:5 & \begin{tabularx}{0.7\textwidth}{X} 事奉有許多種，卻是事奉同一位主。 \end{tabularx} \\ \\ \relax
12:6 & \begin{tabularx}{0.7\textwidth}{X} 工作有許多種，卻是同一位神在萬人中運行萬事。 \end{tabularx} \\ \\ \relax
12:7 & \begin{tabularx}{0.7\textwidth}{X} 聖靈彰顯在各人身上，是要使人得益處。 \end{tabularx} \\ \\ \relax
12:8 & \begin{tabularx}{0.7\textwidth}{X} 有人藉著聖靈領受智慧的言語；有人也靠著同一位聖靈領受知識的言語； \end{tabularx} \\ \\ \relax
12:9 & \begin{tabularx}{0.7\textwidth}{X} 又有人由同一位聖靈領受信心；還有人由同一位聖靈領受醫病的恩賜； \end{tabularx} \\ \\ \relax
12:10 & \begin{tabularx}{0.7\textwidth}{X} 又有人能行異能，又有人能作先知，又有人能辨別諸靈，又有人能說方言，又有人能翻方言。 \end{tabularx} \\ \\ \relax
12:11 & \begin{tabularx}{0.7\textwidth}{X} 這一切都是由惟一的、同一位聖靈所運行，隨著自己的旨意分給各人的。 \end{tabularx} \\ \\ \relax
12:12 & \begin{tabularx}{0.7\textwidth}{X} 就如身體是一個，卻有許多肢體，身體的肢體雖多，仍是一個身體；基督也是這樣。 \end{tabularx} \\ \\ \relax
12:13 & \begin{tabularx}{0.7\textwidth}{X} 我們無論是猶太人是希臘人，是為奴的是自主的，都從一位聖靈受洗成了一個身體，並且共享這位聖靈。 \end{tabularx} \\ \\ \relax
12:14 & \begin{tabularx}{0.7\textwidth}{X} 身體原不只是一個肢體，而是許多肢體。 \end{tabularx} \\ \\ \relax
12:15 & \begin{tabularx}{0.7\textwidth}{X} 假如腳說：「我不是手，所以不屬於身體」，它不能因此就不屬於身體。 \end{tabularx} \\ \\ \relax
12:16 & \begin{tabularx}{0.7\textwidth}{X} 假如耳朵說：「我不是眼睛，所以不屬於身體」，它也不能因此就不屬於身體。 \end{tabularx} \\ \\ \relax
12:17 & \begin{tabularx}{0.7\textwidth}{X} 假如全身是眼睛，聽覺在哪裡呢？假如全身是耳朵，嗅覺在哪裡呢？ \end{tabularx} \\ \\ \relax
12:18 & \begin{tabularx}{0.7\textwidth}{X} 但現在神隨自己的意思把肢體一一安置在身體上了。 \end{tabularx} \\ \\ \relax
12:19 & \begin{tabularx}{0.7\textwidth}{X} 假如全都是一個肢體，身體在哪裡呢？ \end{tabularx} \\ \\ \relax
12:20 & \begin{tabularx}{0.7\textwidth}{X} 但現在肢體雖多，身體還是一個。 \end{tabularx} \\ \\ \relax
12:21 & \begin{tabularx}{0.7\textwidth}{X} 眼睛不能對手說：「我用不著你。」頭也不能對腳說：「我用不著你。」 \end{tabularx} \\ \\ \relax
12:22 & \begin{tabularx}{0.7\textwidth}{X} 不但如此，身上的肢體，人以為軟弱的，更是不可缺少的； \end{tabularx} \\ \\ \relax
12:23 & \begin{tabularx}{0.7\textwidth}{X} 身上的肢體，我們認為不體面的，越發給它加上體面；我們不雅觀的，越發裝飾得雅觀。 \end{tabularx} \\ \\ \relax
12:24 & \begin{tabularx}{0.7\textwidth}{X} 我們雅觀的肢體自然用不著裝飾；但神配搭這身子，把加倍的體面給那有缺欠的肢體， \end{tabularx} \\ \\ \relax
12:25 & \begin{tabularx}{0.7\textwidth}{X} 免得身體不協調，總要肢體彼此照顧。 \end{tabularx} \\ \\ \relax
12:26 & \begin{tabularx}{0.7\textwidth}{X} 假如一個肢體受苦，所有的肢體就一同受苦；假如一個肢體得光榮，所有的肢體就一同快樂。 \end{tabularx} \\ \\ \relax
12:27 & \begin{tabularx}{0.7\textwidth}{X} 你們是基督的身體，並且各自都是肢體。 \end{tabularx} \\ \\
[1ex]
\hline
\hline
\end{longtable}
$^{1}$(下集時間).
(主持人:我預祝大家).
(主持人:像剛才荒謬谷隊兩首詩歌).
(主持人:新春樂,樂二十年年).
(主持人:預祝大家).
(主持人:下星期就年初一了,新年快樂).
(主持人:我當然像剛才口水妹那位姐妹那樣說).
(主持人:願大家平安).
(主持人:不過我特別想說這個祝福).
(主持人:大家記得這是耶穌的祝福).
(主持人:願你們平安).
(主持人:在新年裡有很多的輝春,其中一個輝春就是出入平安).
(主持人:不過你們知道吧).
(主持人:聖經裡或耶穌的祝福,說願你們平安).
(主持人:聖經所說的平安).
(主持人:不是我們新年裡出入平安的平安那麼簡單).
(主持人:聖經裡的平安).
(主持人:其實包括另一個意思,我不知道大家明不明白).
(主持人:就是和平).
(主持人:英文PACE的字正正可以解作平安).
(主持人:又可以解作和平).
(主持人:我特別今天在開始的時候).
(主持人:將這個祝福,耶穌的安慰帶給大家).
(主持人:因為我不知道大家怎樣).
(主持人:過去一年,幾年).
(主持人:在疫情底下,很多社會的變化底下).
(主持人:我不知道你們覺得平不平安).
(主持人:我不知道你們覺得這個世界和不和平).
(主持人:俄烏的戰爭).
(主持人:今天我聽到一個消息).
(主持人:俄羅斯被導彈擊落在烏克蘭的建築物裡).
(主持人:有很多人死了).
(主持人:我們發覺這個世界很不和平,很不平安).
(主持人:頂智妹和平很重要).
(主持人:也和我們今天思想的主題息息相關).
(主持人:中國人有句說話).
(主持人:相見好同住難).
(主持人:我們神學院以前).
(主持人:我在香港神學院也事奉了十多年,二十年).
(主持人:早期我入去的時候我們有宿舍).

$^{41}$(主持人:其實在神學舊約裡宿舍很重要).
(主持人:讓一些神學生大家一起住).
(主持人:如果平時不一起住的時候).
(主持人:這些神學生通常都是好人).
(主持人:大家都是好人).
(主持人:坐在這裡).
(主持人:挺好的).
(主持人:但是一住在一起宿舍的東西就出現了).
(主持人:所以在神學舊約裡宿舍其實也挺重要).
(主持人:不過因為香港寸金隻土).
(主持人:我們終於沒有宿舍了).
(主持人:現在我們開始有新的校舍).
(主持人:裝修我們也有一點宿舍).
(主持人:相見好同住難).
(主持人:大家都知道).
(主持人:如果你拍拖的時候).
(主持人:你拍拖的時候的愛情是很浪漫的).
(主持人:一日不見就如隔三秋).
(主持人:結了婚呢).
(主持人:我經常說).
(主持人:拍拖中的愛情是浪漫).
(主持人:結婚是現實).
(主持人:在馬克上面也是一個現實).
(主持人:你一起住).
(主持人:特別是).
(主持人:我回憶我結婚很久了).
(主持人:結婚初期兩夫婦很多適應).
(主持人:因為生活習慣不同).
(主持人:性格不同).
(主持人:很多事情一起住的時候).
(主持人:又不是說拍拖的時候).
(主持人:沒有給你真面目).
(主持人:不過一起住的時候).
(主持人:更多生活習慣日常的事情).
(主持人:一些很細微的事情).
(主持人:一些性格).
(主持人:或者更加真面目).
(主持人:你就看到了).
(主持人:很多磨合).
(主持人:那是現實).

$^{81}$(主持人:家庭裡是這樣的).
(主持人:婚姻裡,愛情關係裡).
(主持人:剛才說那些).
(主持人:弟兄姊妹之間一起住).
(主持人:和睦同居是不容易的).
(主持人:我可以說大家一定同意的).
(主持人:你在職場裡工作).
(主持人:我們經常說).
(主持人:那些工作,那些事).
(主持人:很多事是做不死人的).
(主持人:最煩的).
(主持人:最難搞的就是).
(主持人:人際關係).
(主持人:這必然是的).
(主持人:我做了這麼多事).
(主持人:我做事這麼多年).
(主持人:在事奉之前).
(主持人:我都做過一段時間).
(主持人:我發覺真的很困難).
(主持人:我已經不說人性,罪性).
(主持人:有些人在辦公室裡).
(主持人:搞辦公室政治).
(主持人:為了自己上位,爭權).
(主持人:為了某些利益).
(主持人:我不說那些了).
(主持人:那些是人的罪性).
(主持人:更加麻煩).
(主持人:就算大家都是好人).
(主持人:不是有很多私心).
(主持人:不是關心自己的利益).
(主持人:但其實都是一個問題).
(主持人:每個人都是有獨特的性格).
(主持人:一起工作,同事).
(主持人:其實很難一起共事).
(主持人:我們今天說的合一).
(主持人:這是魯牧師給我的題目).
(主持人:今天跟大家說合一).
(主持人:職場裡面大家明白).
(主持人:很坦白說).
(主持人:回到教會裡面).

$^{121}$(主持人:合一是一個老生常談的題目).
(主持人:我今天跟大家說一段經文).
(主持人:也是講合一裡面其中一段經常說的經文).
(主持人:你不要期待有甚麼新的東西聽).
(主持人:老生常談).
(主持人:說就容易,但實際上是很難的).
(主持人:實踐是很困難的).
(主持人:在教會裡面不一樣).
(主持人:我們經常說我們是一個愛的群體).
(主持人:我們說我們要合一).
(主持人:但大家苦心自問).
(主持人:一起同工的時候).
(主持人:一起事奉的時候).
(主持人:很多人想法不同的).
(主持人:很難的).
(主持人:有時很多是爭執的).
(主持人:我們在職場上上班的時候).
(主持人:大家不合就不工作,就辭工).
(主持人:教會又怎樣呢).
(主持人:有時也有的).
(主持人:傳道人也是這樣,一起事奉就轉工).
(主持人:弟兄姊妹又怎樣呢).
(主持人:我又不是這份工作,不是做傳道人).
(主持人:一起去執事會也好).
(主持人:大家弟兄姊妹一起事奉也好).
(主持人:很多不和的).
(主持人:你不要告訴我你們沒有).
(主持人:我這麼大個人也不相信).
(主持人:每一道道都是一樣的).
(主持人:因為這是人性的問題).
(主持人:那怎麼辦呢).
(主持人:有些人就離開教會,去其他教會).
(主持人:不過全部教會都是這樣).
(主持人:世界上有人的地方就有人的困難).
(主持人:人際關係的問題).
(主持人:所以我經常說).
(主持人:今天這篇道理).
(主持人:說是很容易的).
(主持人:但說真的,實在是很難).
(主持人:我自己一生中).

$^{161}$(主持人:其中一個要學習的).
(主持人:一個屬靈功課就是這個功課).
(主持人:不斷要學習這個功課).
(主持人:我有些).
(主持人:現在也有陳學生在).
(主持人:我跟很多陳學生說).
(主持人:希臘文,希伯來文很難).
(主持人:我教的教義很難).
(主持人:最難學的功課).
(主持人:其實是剛才的功課).
(主持人:合一的功課).
(主持人:今天我選擇了哥林多前書).
(主持人:大家知道了).
(主持人:為什麼要哥林多前書說這些).
(主持人:因為哥林多教會本身就是).
(主持人:一個分黨分派的教會).
(主持人:哥林多教會其實就是代表了).
(主持人:我們歷史歷代教會).
(主持人:一般教會的問題).
(主持人:其實它只不過是一個反映).
(主持人:這是一個很普遍的問題).
(主持人:他們有很多分黨分派).
(主持人:不合一,很紛爭).
(主持人:這是他們教會的問題).
(主持人:當然不是只跟哥林多教會的問題).
(主持人:也是我們宗教的問題).
(主持人:甚至是社會的問題).
(主持人:所以保羅特別).
(主持人:可能問題會更嚴重).
(主持人:很具體的).
(主持人:保羅這封信).
(主持人:其中一個要處理的就是).
(主持人:這個紛爭的問題).
(主持人:我今天選擇了).
(主持人:哥林多前書十二章).
(主持人:跟大家思想的問題).
(主持人:我由第四節開始).
(主持人:其實在這段經文的時候).
(主持人:我稍為講講第一節,三節).
(主持人:很重要的).

$^{201}$(主持人:其實這段經文是講甚麼呢?).
(主持人:是講一種屬靈的恩賜).
(主持人:或者這樣說).
(主持人:我剛才說這個課題).
(主持人:你現在從一個講).
(主持人:或者一個觀念上來說).
(主持人:其實你聽過很多次).
(主持人:困難就是實踐的問題).
(主持人:實踐就牽涉到).
(主持人:所謂的那種甚麼).
(主持人:我們經常說的那種屬靈生命的操練問題).
(主持人:靈修靈修的靈不是靈魂).
(主持人:那個靈是性靈).
(主持人:我們是不是屬於性靈).
(主持人:我很坦白地跟大家說).
(主持人:一個人的屬靈生命).
(主持人:如果那個人).
(主持人:經常跟別人有很多的爭執).
(主持人:很難跟別人相處).
(主持人:其實反映他屬靈生命的狀態).
(主持人:如果將它再擴大一點).
(主持人:一間教會).
(主持人:如果你看到一間教會).
(主持人:經常在裡面有很多).
(主持人:好像那些教會).
(主持人:分黨分派).
(主持人:紛爭 撕裂).
(主持人:某程度上就反映那間教會).
(主持人:弟兄姊妹).
(主持人:那個屬靈生命的狀態的問題).
(主持人:我不知道你同不同意).
(主持人:在家庭裡兩夫婦經常吵架).
(主持人:這是什麼問題?).
(主持人:是生命的問題).
(主持人:講道理是重要的).
(主持人:但去到實踐為什麼這麼難?).
(主持人:因為去到生命就很複雜).
(主持人:這是一個實踐 一個操練的問題).
(主持人:所以第十二章第一節第三節講什麼?).
(主持人:保羅對著那些教會的弟兄說).

$^{241}$(主持人:你們以前是外邦人的時候).
(主持人:隨時被軒轅受迷惑去服侍那些偶像).
(主持人:你們知道的).
(主持人:不過現在不同了 第三節).
(主持人:我讀一本給大家聽).
(主持人:保羅說:我告訴你們).
(主持人:被神的靈 即是聖靈).
(主持人:被聖靈感動的).
(主持人:不會說耶穌是可咒詛的).
(主持人:被聖靈感動的人).
(主持人:按著聖靈的教導去宣講).
(主持人:跟別人說耶穌).
(主持人:不會說耶穌是可以被咒詛的).
(主持人:不會去奏錯耶穌).
(主持人:如果你信服聖靈的時候).
(主持人:還有下一句).
(主持人:如果不是被聖靈感動的).
(主持人:也沒有能說耶穌是主).
(主持人:聖靈 一個人按著聖靈的感動).
(主持人:他屬於聖靈的).
(主持人:於是他向別人見證耶穌).
(主持人:向別人介紹耶穌的時候).
(主持人:如果是被聖靈感動的人).
(主持人:他一定會說出來).
(主持人:耶穌是主).
(主持人:這三節經文講了兩個道理).
(主持人:第一 正在講).
(主持人:聖靈和主耶穌之間的合一).
(主持人:聖靈的工作).
(主持人:聖靈在人的心裡的工作).
(主持人:譬如我講道也好).
(主持人:你去跟別人傳福音也好).
(主持人:聖靈在裡面講工作).
(主持人:以至我來見證一些事情的時候).
(主持人:見證福音的時候).
(主持人:見證主耶穌的時候).
(主持人:如果我是信服聖靈的教導的時候).
(主持人:我傳講是按照聖靈的感動去講的時候).
(主持人:我一定是甚麼?).
(主持人:承認耶穌是主的).

$^{281}$(主持人:而聖靈就不會在我裡面講).
(主持人:你講耶穌是賊).
(主持人:耶穌是魔鬼).
(主持人:不會這樣講的).
(主持人:這裡講明聖靈和耶穌之間的那種配搭).
(主持人:那種合一).
(主持人:我經常跟弟兄姊妹和學生說).
(主持人:聖靈是很特別的).
(主持人:如果我用一個戲劇的比喻).
(主持人:主角是甚麼?是耶穌).
(主持人:聖靈就好像一些射燈).
(主持人:有一點隱藏的).
(主持人:你進去看話劇).
(主持人:你不會一進去就看著那些射燈).
(主持人:那些燈光).
(主持人:你是看著主角).
(主持人:耶穌基督來到這個世界上面).
(主持人:祂就是主角).
(主持人:而聖靈是配合祂的).
(主持人:如果你在電視台做話劇).
(主持人:沒有燈光射著主角的時候).
(主持人:你看不到主角在演戲).
(主持人:而聖靈就是一種這樣的配搭).
(主持人:跟耶穌同供).
(主持人:祂就指著).
(主持人:那個是主角,耶穌是主).
(主持人:弟兄姊妹看不看到).
(主持人:一開播十二章第一節三節).
(主持人:已經把這幅圖畫說出來了).
(主持人:祂們是同供的,合一的).
(主持人:我們上帝的生命就是這樣的生命).
(主持人:我們說是一群屬於上帝).
(主持人:我們是屬於聖靈的).
(主持人:你知道耶穌是真理,是吧?).
(主持人:耶穌說:我就是真理).
(主持人:而聖靈是甚麼?).
(主持人:真理的聖靈).
(主持人:聖靈是甚麼?).
(主持人:是幫助我們認識耶穌這個真理的靈).
(主持人:就是這段說話的意思).

$^{321}$(主持人:弟兄姊妹說到這裡的時候).
(主持人:我想說其中一件事).
(主持人:很多時候我發覺).
(主持人:我們不說外面的世界,說教會).
(主持人:教會裡面我們說都有紛爭).
(主持人:甚至乎那時候).
(主持人:例如之前我們).
(主持人:大家在社會運動期間).
(主持人:我們很多政見不同).
(主持人:我們有撕裂).
(主持人:很多的紛爭).
(主持人:其中一個問題在哪裡?).
(主持人:弟兄姊妹).
(主持人:恕我坦白說).
(主持人:大家一家人).
(主持人:我發覺有不少基督徒).
(主持人:有一個錯誤的理解).
(主持人:以為自己所說的是真理).
(主持人:然後).
(主持人:跟他們看法不同的那些).
(主持人:那個對方說的就是歪理).
(主持人:我的才是真理).
(主持人:你的才是歪理).
(主持人:以為自己擁有真理).
(主持人:弟兄姊妹,不要搞錯).
(主持人:唯有耶穌才是真理).
(主持人:而聖靈就是幫助我們認識真理的聖靈).
(主持人:你和我不是真理).
(主持人:甚至你和我不可以說我們擁有真理).
(主持人:我記得).
(主持人:好幾年前).
(主持人:我去了一間教會).
(主持人:那間教會不是很大間).
(主持人:小小一間的教會).
(主持人:差不多一年時間).
(主持人:我每一年都去港島).
(主持人:我記得有一次我去港島).
(主持人:不是第一次去港島).
(主持人:那次去港島很多年).
(主持人:我去港島講甚麼呢).

$^{361}$(主持人:當然是講聖經).
(主持人:我講關懷貧窮人).
(主持人:完了之後).
(主持人:有一位教會的領袖).
(主持人:他很客氣地跟我說).
(主持人:趙先生,你剛才講的不是福音).
(主持人:我不知道是不是有關).
(主持人:然後我就不能去港島了).
(主持人:我想舉一個例子).
(主持人:在他心目中).
(主持人:他那套福音的真理觀).
(主持人:或者他那套福音觀就是).
(主持人:沒有理由你講).
(主持人:福音不是這樣的).
(主持人:福音不是開報道會).
(主持人:然後叫人靈魂得救上天堂).
(主持人:那些不是福音嗎).
(主持人:你講關心貧窮人).
(主持人:關心弱身的需要).
(主持人:那些哪是福音).
(主持人:於是他覺得你說的不是真理).
(主持人:這個不是福音).
(主持人:在他心目中).
(主持人:他的真理就是有他的那套福音觀).
(主持人:然後他認為跟他不同的).
(主持人:你就不是真理).
(主持人:你說的那套是歪理).
(主持人:舉一個例子).
(主持人:有時候我們彼此之間).
(主持人:有些爭執).
(主持人:很多時候).
(主持人:這裡搞錯了).
(主持人:特別是基督徒).
(主持人:回到教會以來).
(主持人:你以為你自己那套東西就是真理).
(主持人:先停下來).
(主持人:唯有上帝是真理).
(主持人:唯有耶穌是真理).
(主持人:然後我們用聖靈).
(主持人:真理的聖靈).

$^{401}$(主持人:幫助我們去接近真理).
(主持人:我讀了神學這麼久).
(主持人:教了神學這麼久).
(主持人:我越讀聖經越讀神學).
(主持人:我發覺我不夠膽再說什麼是).
(主持人:我不是說沒有真理).
(主持人:我們是認識真理).
(主持人:不過我越來越發覺).
(主持人:我不夠膽說我說的就是真理).
(主持人:因為我教神學的).
(主持人:我說的都不是真理).
(主持人:不是).
(主持人:我希望大家首先放下一樣東西).
(主持人:有時候就是這些觀念).
(主持人:這些所謂的).
(主持人:我們楊木菊牧師).
(主持人:不知道大家有沒有聽過).
(主持人:他在在生的時候).
(主持人:經常有我們老到壞鬼神學的觀念).
(主持人:但我們這些壞鬼神學的觀念).
(主持人:我們以為是真理).
(主持人:那就麻煩了).
(主持人:就令到教會有很多這樣的爭執).
(主持人:我講兩個例子).
(主持人:福音和政治有沒有關係?).
(主持人:不行,教會不應該講政治的).
(主持人:福音沒理由是和政治有關).
(主持人:剛才你唱的那首歌是什麼歌?).
(主持人:有所謂萬王之王,萬主之主).
(主持人:我沒時間和大家詳細講這個問題).
(主持人:福音其實是和政治有關的).
(主持人:你還在說歪理).
(主持人:我們教會不應該講政治,聖經不應該講政治).
(主持人:我沒時間和大家詳細講).
(主持人:其實聖經從頭到尾有很多聖經講政治).
(主持人:有所謂有政治罪名被釘死的猶太之王).
(主持人:當然聖經福音的政治和我們很膚淺的理解外面的政治).
(主持人:有些東西是不同的).
(主持人:但不要這麼快,因為講政治的不是福音).
(主持人:所以要讀神學).

$^{441}$(主持人:每個人都要讀神學).
(主持人:幫助我們合一的其中一個原因是).
(主持人:你要搞清楚一些).
(主持人:譬如對聖經,如何解經).
(主持人:歷史歷代的神學家如何思想認識上帝).
(主持人:你要知道的).
(主持人:人有些觀念你卡住在那裡).
(主持人:這就是大家磨擦的其中一個原因).
(主持人:再加上很坦白地說,這個世界變得很快).
(主持人:我和我女兒其實有代溝的).
(主持人:很多代溝,我們成長不同,價值觀念不同).
(主持人:我都自問我不是一個很古板的人).
(主持人:很多東西我都嘗試從他們的角度看).
(主持人:但始終是不同的,有時就是觀念的問題).
(主持人:我經常說,特別是我自認年長的).
(主持人:我自認的,我兩元雞可以坐車,你猜到的).
(主持人:當你年紀越大的時候,人有時是越頑固的).
(主持人:有些東西是越執著的).
(主持人:這些全部都是令我們不能合一的重要原因).
(主持人:這一節三節其實說了一樣東西).
(主持人:我今天重點當然是給大家經文,大家讀那些).
(主持人:大家留意,保羅在四至六節).
(主持人:他連續用了三對平衡的句子).
(主持人:因此有很多種,卻是同一位勝利所賜).
(主持人:第五節,事奉有很多種,卻是事奉同一位主).
(主持人:第六節,工作有很多種,卻是同一位神在萬人中運行萬船).
(主持人:這是一個華修本的翻譯,如果用華本的平衡句,中文更加明顯).
(主持人:是三對平衡的句子,華本是因此有分別勝利一位,即是有分別主,卻是一位).
(主持人:公用有分別,上帝卻是一位,是很平衡的句子).
(主持人:本來在這裡要表達三節經文的意思).
(主持人:就是多,分別,分別即是多).
(主持人:我們這裡有幾十人,我不知道有多少人).
(主持人:越多人,甚麼是多?即是我們有很多分別).
(主持人:每個人都有很獨特的性格,你的成長背景,你受過教育).
(主持人:很多事情令每個人都很獨特).
(主持人:上帝創造了你,原來上帝創造了你,各從其類).
(主持人:每個人都很獨特,越多即是越多區別).
(主持人:區別的分別是多,英文是many,然後還有one,一位一位一位).
(主持人:上帝是一位一位一位).
(主持人:我這裡想強調一個觀念,分別多和一,多元和一).

$^{481}$(主持人:所以今天我這個題目是多元,多元是上帝創造的心意).
(主持人:人最尊貴是甚麼?就是你和我不同,就算是孖生的,是Trins的,都有不同的).
(主持人:這個世界你是獨一無二的,這是人最尊貴的地方,多元).
(主持人:社會越大的教會就越多,就越多不同的人).
(主持人:這個多和一,多元和一之間的關係,本來就在這裡).
(主持人:所以我今天這個題目就是多元中的一,很重要的一,我們說是合一).
(主持人:上帝從來強調的就是合一,你知道合一和統一是不同的嗎?).
(主持人:要分清楚,統一就是一把聲音,統一所有的說話,每個人都要這樣說話).
(主持人:我記得我讀大學的時候,我在中大讀書,我讀大學的年代是80年代).
(主持人:因為我們的暑假很長,我們沒有錢,不會去歐遊,我女兒現在都會,有時沒錢就回大陸).
(主持人:在早期的年代,70,80年代的時候,回到大陸的時候很特別,每個人穿的衣服都一樣,戴的眼鏡都是一樣,黑框的,而且不是那麼多花款的,都是那個框).
(主持人:說的話都一樣,你現在發覺現在還有這樣的氣質在上面,開會,官方說的話全部一樣,這個是統一).
(主持人:我沒試過去北韓,我相信北韓朝鮮是很統一的,統一有統一的好處,一把聲音就不會亂,不會有紛爭,最重要是有個人,這個就是獨裁的政權,一把聲音,每個人都這樣說話,很和平的,不過那個不知道是和平還是什麼).
(主持人:有好處就不會混亂,真的有好處的,方便管治,方便管理,但是聖經說的不是統一,是合一,合一是什麼意思呢?合一的先決條件,就是大家不同,大家不同才要學習合一,所以統一容易的).
(主持人:一個人很有權,一把聲音,如果是一個很有權威的人,我說什麼就什麼,每個人都要這樣說,可以說有某種好處,但是聖經不是這樣的,合一,所以這麼難,就是這樣,就是多元,我們的世界是這樣,上帝創造的世界也是這樣,特別現在來到的世界,有一世紀的世界).
(主持人:真是一個多元的世界,多元有麻煩的,很自由,每個人都可以說自己的話,很多聲氣,自由有自由的問題,很容易亂,吵架,但是我想說,本來這裡說的教會觀,或者我們整個信仰說的合一不是統一,背後就是信仰).
(主持人:背後就是承認多元差異,你有沒有想過為什麼我們的信仰是信三一上帝的?通常教會是不會說這些教義的東西,只學統一,很多時候三一,我信三一,但是那個太過奧秘,所以我們不要學,不需要學太過奧秘,你有沒有想過為什麼是三一?數字是不明白的,一加一加一是三,我們是不是信三個神?).
(主持人:耶穌基督是不是神?聖靈也是神,天父也是神,那不就是三個神,又不是三個神,是一個神,那是怎麼樣?一加一加一為什麼不是三等於一?).
(主持人:就是說聖父,聖子,聖靈是不同的,這裡說那麼多書信,你回去看十二,剛才說有聖靈,再讀第四節,說聖靈,然後從十二節開始,由第一節到十一節,說聖靈).
(主持人:然後十二節轉去基督,耶穌基督是不是等於聖靈?不是的,聖靈有沒有道成肉身?沒有,只有耶穌基督道成肉身,聖靈是什麼?無形無體,看不到的,它沒有肉身的,不過聖靈又住在我們裡面).
(主持人:這個世界上沒有一個宗教是這樣的,唯獨是我們的信仰是這樣,我們的上帝是這樣,這裡已經說明了,上帝觀和教會觀是相關的,為什麼我們教會講多元中的合一?因為我們上帝也很多元的,聖父,聖子,聖靈是很不同的,但是他們合一的,剛才我說一節三節就是這樣,我們的不同是我們喜歡獨立的,我們喜歡華時紀人的罪性問題,我們很難合一的).
(主持人:但是我們的上帝很特別,耶穌基督是神,他滿有榮耀的,他是萬王之王,但是他道成肉身的時候又鬼死這麼軟弱,鬼死這麼卑微,然後說他怎麼信服天父的心意,在克西瑪尼亞禱告他有他的想法,他有他的意志,但是他又放下他自己的意志,信服上帝).
(主持人:耶穌基督來到這個世界上,他透過這樣的生命展示原來三個上帝,耶穌本來是有權的,他放下,然後信服,怎樣合一一切,配搭一切工作,我們的上帝是這樣,所以上帝透過耶穌基督設立教會,他期望我們教會要學這個功課,沒錯我們很難學).
(主持人:所以你有什麼上帝觀影響你,你不是學習耶穌基督,耶穌基督也信服聖靈的,我們又怎樣信服聖靈呢?為什麼耶穌基督沒有犯罪?他進入抗日的時候試探,為什麼沒有犯罪?因為他是神,不是,這個答案錯,因為他是人,不過這個人和我們的人不同,你和我為什麼經常犯罪?因為我們不信服聖靈).
(主持人:但是耶穌基督信服聖靈,所以他性格試探,他和聖靈配搭,然後聖靈又配搭他的工作,大家看不看到?整個三個上帝是這樣的生命,那你又明白為什麼十二章之後十三章說什麼?所以我經常覺得實體聖經很重要,你看這些powerpoint又不知道下一節是什麼,除非你很熟悉聖經,你看十三章說什麼?愛的篇章).
(主持人:有人說上帝就是愛,三一的上帝就是愛,合一就是什麼?原來合一就是愛,真正的愛就是合一,愛不是統一,統一是權的問題,合一是愛的問題).
(主持人:所以保羅針對那麼多教會的問題,在十二章說這些道理,然後第十三章,為什麼能夠實踐到?我經常說這個講義,但是如何能夠實踐這個屬靈生命的粗劣問題,而當中就是生命問題,生命問題就是你和我有沒有愛?我們經常說教會是一個愛的群體,譬如過往的撕裂,大家政見上的撕裂,或者大家價值觀的不同,對一些問題的看法不同,就算是時代的問題,都可以解釋出來,但是你和我有沒有愛?你和我有沒有愛?你和我有沒有愛?).
(主持人:就算是時代的問題,價值觀的不同,就算是時代的問題,教會的資源應該用在哪裡?大家都有不同的看法,於是爭論到臉紅耳赤,我們說我們是一個愛的群體,其實教會的分裂不合一,正正是打嘴巴,愛裡面有什麼?很多的粗亂,愛恨恆久忍耐,只是這件事已經很難忍耐,要忍你?.
在你前面討厭你後面,我忍你,忍你說的話,我說的才對,我那些是真理,你那些是歪理,要怎麼忍你?你這樣的性格怎麼忍你?而且要恆久忍耐,今天我們現代人越來越缺乏耐性,忍耐是一個很難學的功課,不疾度,弟兄姊妹,按人性來說,為什麼有時候會有這麼多的紛爭?.
就是比較,競爭,我們資本主義社會的文化就是塑造我們這樣的一個人性,我們和別人比較,我們要贏,贏和輸很重要,我們很介意別人怎麼看自己,自我形象,這是人性,再加上我們的文化,我們很容易疾度,人性很容易疾度,某某人的表現開始學,每次很多人報他的班,我那班就小貓一兩隻,.
甚至開不成,就疾度他,比較,競爭,不疾度,恩慈,什麼是恩慈?上帝給我恩典我就收,我要學習恩慈,不自誇,不張狂,.
不求自己益處,我在香港出生的,整個社會,爸爸媽媽教導我,整個教育制度教導我,獅子山精神是什麼?就是求自己的益處,資本主義背後是什麼?就是講競爭,講自己的利益,你不和別人爭,你怎麼可以上位?什麼獅子山精神?不是靠勤力的,不是這麼簡單的,.
不求自己益處很難的,人性就是自我中心的,為自己的,凡事包容我,就是包容不了,包容不了就不知道什麼,包容不了就不會撕裂,和我的政見不同,見到那些同性戀者,我怎麼包容?罪就是罪,聖經說是罪,你要悔改,我包容你,什麼叫悔改?.
同性戀者除非你變成雙性戀者,不然就代表你悔改,不變成雙性戀者,不代表你不悔改,我怎麼包容你?對同性戀包容,是什麼意思?教會裡面這些不是很敏感的課題嗎?這是基督教倫理學的問題,我們贊成墮胎嗎?這是基督教倫理學有爭議的問題,什麼叫包容?.
包容背後就是接待,大家知道接待是一個很重要的功課嗎?從舊日到新日,接待什麼?什麼叫真正的接待?這段經文也有說,那些不體面的,教會裡面很弱的,用肢體身體的觀念,身體也只是一個,我們有很多肢體,你們看回時有多少肢體?.
有眼耳口鼻,我們有一條盲腸,盲腸是多餘的,割掉它,只會有盲腸炎,盲腸對我們身體沒什麼用,割掉它,手指尾腳指尾,那些不是特別嗎?.
如果你身體裡面有一樣器官,不知道為什麼,原因就是令它不美觀,不如最好弄掉它,不要理它..
那個比喻很重要的,舊日的身體,我們有很多肢體,姓名隨他自己的意思,分配給個人有不同的恩賜,有些做什麼,有些做那些,有些講語言,有些講方言,有些講什麼,有些教道,每一個不同..
第一,我們很多時候比較,為什麼上港壇當然是最重要的,為什麼我上不了,為什麼我每天回來就在下面量體溫,我回了廿年教會都是每次都在量體溫,不受重視..
這個是上帝聖靈有他的心意,他的自由,按著每個人,成就整個教會,你看看剛才的經文,上帝給不同的人有不同的恩賜,不同的職事,不同的功用,有不同的恩賜,不同的職事,發揮了功用,每個人都有功用,做到一些事情,以至整個教會,大前提是為眾人的益處,因為是common good,我們那種紛爭就是我們那時候的資本主義社會,就是我的好處..

$^{521}$我的好處,我的利益,我上不上到位,我有沒有話事權,我有沒有權力,很自我中心,越是自我中心,撕裂的分化是越大問題..
合一,為了整體的好處,為了別人的好處,第七節說的是叫人的益處,那就不只是看自己,什麼叫別人的益處呢?.
就是有一個特別弱的,但他是屬於這個肢體裡面的其中一個,你不可以說手指尾有什麼用,這條盲腸有什麼用,割了它,我們是不是有這樣的心態?.
我們的社會文化,我們的人性,我們的罪性,其實是造就我們這樣的看法,我們在社會裡面,我們很重視,我們每個人都想做強人,女強人很有社會地位,我太太是家庭主婦,沒有什麼社會地位,沒有生產能力,沒有那種生產能力就沒有什麼地位..
這是資本主義社會的其中一個問題,一個人有沒有用?是不是在社會上有地位?是視乎他有沒有用?一個沒什麼用的人,我們就會用廢言,廢青,廢柴,老人家是廢的,老的沒用的,教會不是這樣看的..
越是弱的,你就要去加級他的體面,你越是要關心他,因為從舊有開始就這樣教的,什麼叫公義?公義和愛是一樣的,同一樣的,好像一錢不是兩面,上帝教導,在舊有裡面,照顧孤兒寡婦..
所以剛才我說的不是有偏導,我是說要關心貧窮,那是福音來的,什麼叫愛?愛就是福音,福音就是愛,耶穌基督的愛,上帝的愛,三分之一的愛,神就是愛,神就是愛..
你關心一些有需要的人,福音不是一個講字的,其實就是關心他,我們曾經試過,沒有口罩,一起給予口罩,特別是沒有的人,Panadol沒有了,我家裡還有的,你沒有了,你需要的時候就來一點..
這是福音,這是接待,耶穌說愛仇敵,這是最難實踐的,我愛一些我喜歡的人就容易很多,常常為老取暖,我昨天約了一些大學的同學出來,大家對某些看法很相近,一出來就容易談很多,聊天很多,為老取暖,同聲同氣..
愛那些不難的,愛和自己很不同的人,仇敵代表什麼?和自己很不同的人,甚至害自己的人,耶穌是什麼?耶穌就是被人害到上十字架,他在十字架的時候赦免害他的人,然後教導門徒要愛仇敵,仇敵就是和自己很不同的人..
區別,差異,那些是難實踐的..
剛才說沒有體面的,弱的,我想舉一個例子,我很喜歡讀一位天主教的學者叫盧雲,我很喜歡讀盧雲的書,他講他最後十年的事奉是什麼?他最後十年的事奉就是帶領他去服侍一些嚴重傷殘的人..
有一本書講他服侍的那個人叫阿當,就像我們創世的阿當,阿當是什麼?一個嚴重弱智傷殘人,這樣形容他就是他真的好像一塊肉一樣,他不懂得說話,沒有思想,他的嚴重程度到這樣,他完全不能夠照顧自己,不得而知,他就是這樣,他的存在是什麼?.
這個沒有體面,非常之弱的肢體,他見到一個這樣的肢體的時候,他就說他見到這個好像是耶穌基督的,好像見到耶穌一樣,耶穌就說他成為一個很軟弱的人,他的軟弱是什麼?.
就是因為他不能夠照顧自己,他需要別人的照顧,於是他的存在就招聚了一班人一起去服侍,服侍這個軟弱的肢體,他的存在就好像一個和平的使者一樣,招聚了一班人,盧雲是其中一個,就一起來去服侍他,為他抹身,為他預備早餐,.
甚至他完全不能夠做任何的事,他就需要旁邊的人幫助他,一個很美麗的圖畫,一個愛的群體,原來就是一個這樣的肢體,這個肢體原來就是一個和平的使者,一個讓我們能夠學習愛,實踐愛,實踐合一的一個這樣的人,.
我們不要只尊重一些有權有勢有能力的人,我們又怎樣看一些很弱的人,弟兄姊妹我要說的都是這樣,老生常談,這裡背後有一個很重要的神學,固然是,你對上帝,一個上帝觀的問題,一個三一論的問題,然後你一個怎樣的上帝觀,你又一個怎樣的教會,.
這些觀念固然重要,我再說,這些觀念固然重要,有時候我們跟別人爭辯,因為我們沒有這些神學,但是只是這樣就不夠,怎樣實踐,這個是屬靈操練的問題,這個是你跟我一生要學的生命的問題,.
一路讀十三章,怎樣學習,一起成為一個愛的群體,剛才那些忍耐,恩慈,什麼什麼,不可以什麼什麼,凡事什麼什麼,每樣東西都很難學,你再看聖靈所記的果子,全部都是這些東西,我要說的就是這樣,.
問題就是我們怎樣實踐,怎樣學習,我們是不是成為一個屬於聖靈,屬靈人,屬於聖靈的生命,屬於聖靈的生命就是說,他引導我們回到耶穌基督的真理,耶穌基督就是一間教會的頭,是他建立一個這樣的身體,.
我經常說,頂尖會,我們再沒有撕裂的權利,我們面對的是一個很困難的情況,前面的路,我們整個社會,教會前面的路可能是比以前難走的了,.
我們同是一條船,儘管大家一定是,很多東西不同,這個是正常的,但是正正大家不同的時候,一起學習怎樣合一,互相包容,互相接待,特別接待一些軟弱的肢體,.
就是有待聖靈怎樣來到大家的生活裡面,或者在群體裡面,怎樣來操練大家,我們一起學習,對低頭祈禱..
所以我們三一的上帝,我們的信仰很特別,我們敬拜的,我們屬於的,就是想讓你是一位這樣的三一的上帝,然後聖經說,你就是愛,三一的上帝,聖父,聖子,聖靈,你們很不同,但是你們就是合一,這個就是愛..
上主我們讓你來創造,讓你來拯救,但是我們真的苦心自閉,我們在你面前再一次求你寬恕我們,這個愛的功課我們實在學得很差很差,因為我們可能有太多的罪性,太多個人的性格,很多的膚淺的看法,以前我們學不到的,上帝,求你幫助我們..
耶穌你是我們的頭,是你建立我們這個教會,這個身體,聖靈你住在我們裡面,我們整個教會,每個信徒,頂尖的屬靈生命,聖靈你是會幫助我們,求你帶領我們,讓我們能夠回轉,回到你裡面,我們一起學這個畢生裡面最難學的功課,求主帶領,使引我們這樣做..
使引我們這樣祈禱,教導感恩,奉主明求..
謝謝大家收看 再見.
\newpage



\section{約書亞記申命記 1:5}
\label{sec:8tGomcGt7oY}
\textbf{2023.01.22 《讓約書亞的忠心勒著你和我》 約書亞記 1\_5;申命記 1\_37-38;約書亞記 24\_15 (和修版) 講員 : 高淑芬傳道}
\newline
\newline
連結: \href{https://youtube.com/watch?v=8tGomcGt7oY}{\texttt{https://youtube.com/watch?v=8tGomcGt7oY}} ~~~~ 語音日期: 2023-01-26
\newline
\newline
\hyperref[sec:mxoNZ6Tp6YY]{\small{< < < PREV SERMON < < <}}
~
\hyperref[sec:index_chronic]{\small{[返順時目]}}
~
\hyperref[sec:index_scriptual]{\small{[返順卷目]}}
~
\hyperref[sec:_P22pS_FVkg]{\small{> > > NEXT SERMON > > >}}
\newline
\newline
$^{1}$(日語).
可能你今天沒有發現.
一回來教會.
門上有些泰文.
泰文就是讓我們一起來度新年快樂.
一起歡歡喜喜度過這個年.
在這個新年裡.
願神繼續祝福我們.
繼續蒙受上帝的恩典.
我們很重要的是心靈健康.
不單是身體健康.
你的心靈不健康也會影響身體.
所以三方面我們求神繼續讓我們健康.
今天很歡喜快樂.
因為大家坐在這裡都是健健康康.
最後要祝大家萬事勝意.
高傳道,你是不是寫錯字了?.
我想解釋一下.
萬事勝意的意思是.
一切的結果都會比你當初所想的還要好.
你也聽過萬事如意.
如意就低了一層.
所以接下來祝福別人.
你應該知道應該祝福什麼.
因為萬事如意.
就是你意料之中不滿足.
就是這樣.
不夠萬事勝意這麼厲害.
萬事勝意又是什麼呢?.
是不是你寫錯字了?.
我想跟大家透過約書亞的人物.
來到他特別的中心.
看完他之後.
你就會明白這個新年萬事勝意.
應該是剛才那種寫法.
約書亞的原名是荷西亞.
但是這個名字沒有人知道真名.
就好像新約裡面的巴拿巴.
也沒有人知道真名.
就一直叫他巴拿巴.

$^{41}$約書亞被摩西叫了一會.
他一直由.
甚至我們紅恩堂.
我們一起讀完約書亞.
我們已經開始進入士思記.
所以我就是欣著.
你們也讀過約書亞記.
在經文裡面有很多濃縮.
你回去看過.
你一定會得到更勝過你一籌.
你會更加明白.
他只不過是摩西的幫手.
一個助手.
但他的名字就叫約書亞.
約書亞的名字在希伯來文是約書亞.
但在希臘文裡面.
其實就是耶穌.
等於耶穌.
所以有些人就說到.
約書亞記其實有些都好像.
和啟示錄有呼召.
或者呼應的情況.
為什麼摩西叫他約書亞呢?.
我自己認為.
自己看完的時候.
我覺得亂的兒子.
他叫何西亞.
他的名字是拯救的意思.
但是摩西很想讓何西亞知道.
這個拯救是來自耶和華神.
耶和華讓我們得著拯救.
得著救恩.
所以一直叫他約書亞.
就好像巴拿巴.
巴拿巴的意思是.
他提拔一些新的人.
或者後進後輩.
叫著叫著.
我們都不知道他的真名.
就叫他巴拿巴.

$^{81}$約書亞由何西亞一直叫著叫著.
就是約書亞.
甚至我們剛剛和恩嘉一起讀的聖經.
約書亞記.
都是以約書亞記.
即約書亞這個名字來命名.
約書亞記我們看完整卷的時候.
你就會看到.
約書亞在摩西安息之後.
他就帶著以色列人去真戰的一生.
就看到約書亞都安息了.
在過程裡面.
我們的姐妹就告訴我.
名字都很難讀.
那些地方名字都很難讀.
都不知道哪裡打哪裡.
東南西北都不知道.
不要緊.
你撇開那些不明白的.
你仍然是得到益處的.
當然如果你再慢慢.
聖經是很奇妙的一本書.
你知不知道聖經這麼多年來.
仍然是最暢銷的一本書.
仍然到今時今日.
是不需要像我們教科書那樣.
年年要改版.
要符合這個世界.
是不需要的.
所以在聖經裡面.
你再慢慢看.
你可以有多一點點的得著.
但是就算你不太明白.
什麼阿摩尼人.
什麼阿納人.
一會兒又說要去迦南地.
迦南地原來這麼多.
約書亞記裡面有十個王.
什麼耶路撒冷王.
什麼王什麼王.

$^{121}$很糟糕.
但問題是.
仍然對你有幫助.
正如整條魚多不多骨.
小骨極都有骨.
吃魚有沒有益處.
有的.
小朋友.
小朋友你不要吃魚了.
很多骨.
太危險了.
會不會這樣.
不會的.
你就給他一些.
淡淡肉的地方.
給你的小朋友吃.
我們也是.
我們不太懂吃魚的時候.
就吃一些肉.
可能我們不像別人.
吃魚很厲害.
整條魚吃完之後.
是沒有什麼剩的.
面珠燈.
面珠原來也可以吃.
很厲害.
我們還沒有去到.
這麼神入化的.
吃魚方法的時候.
但是你吃到那些魚肉.
魚腩肉.
在這些肉裡面.
你也得到益處.
我們讀聖經.
再次鼓勵弟兄姊妹.
聖經有很多東西.
我們是不明白的.
未明白.
但是你去讀的時候.
對你我都有幫助.

$^{161}$例如魚的魚肉.
對我們有幫助.
在約書亞在聖經裡面.
第一次出場的時候.
我們就會看到.
約書亞他.
也是被人吩咐.
就是被摩西吩咐他.
你要起來的.
去選擇一些人.
來和亞馬力人爭戰.
而摩西在這個爭戰裡面.
你就會看到.
為什麼摩西也需要一些幫手.
幫手去托著他的手來祈禱.
那場仗結果得勝.
你不要小看.
不要覺得.
只是山上的很重要.
山下的那一群約書亞.
和他挑選的將軍去打仗.
一樣是這麼重要的.
你再去想一想.
約書亞有沒有受過軍事訓練.
一定沒有.
因為約書亞這一群人.
他們是由埃及圍奴.
做奴隸.
他不像我們泰語姐妹.
我們有些姐妹是做家務助理.
家務助理有些公眾假期.
他也會休息.
星期日也有休息.
但是奴隸有沒有休息.
沒有的.
還要是你很深印象.
當摩西要帶領他們出紅海的時候.
法老知道很生氣.
於是對奴隸說.
材料我就是這麼多給你.

$^{201}$你仍然要做這麼多份量的東西給我.
你想想就是這樣.
對奴隸是沒有人情說的.
是很刻薄的.
當然奴隸也會有一些知識訓練.
多數是工藝.
怎樣去做事.
但是軍事訓練是絕對沒有的.
你想想如果軍事訓練.
你們這群特別是埃及法老王.
他非常怕這群埃及.
埃及裡面的以色列人.
因為他們新養眾多.
越來越多人.
如果他們有軍事訓練就不得了.
但是你看到聖經第一次出場的約書亞.
沒有受過軍事訓練.
當摩西叫她去打仗.
打就打她.
你會看到中賢摩西說.
你去打.
然後就去山上祈禱.
倘若約書亞對上帝沒有那份信心.
不是忠心於上帝的拯救的時候.
給你.
你敢不敢.
我不懂打仗的.
現在還要說去打仗.
很大件事的.
但是我們看到約書亞.
在她第一個出場的時候.
她已經是願意的.
你說去就去.
就好像鐵塔尼號.
鐵塔尼號那套著名的「You jump, I jump」.
你跳下去.
跳舊生船.
當船撞的時候.
你跳我就跳.
當然我們說約書亞對上帝極有信心.

$^{241}$當以色列人走完曠野路程的時候.
預備要過約旦河.
可惜摩西在這個時候就安息了.
你想想一班以色列人過去都是看著摩西的頭.
摩西你說怎樣就怎樣.
因為實在在摩西經歷過十個災禍.
以及過紅海這些奇妙的事情.
很大的異象,恩典.
但是現在我們跟著她出埃及的領袖.
她安息了.
這班以色列人是非常絕望.
足足四十年的時間.
整個以色列的民族.
他們的生活在出埃及走.
跟著他們因為一些緣故要在曠野裡.
足足四十年裡.
他們的衣食住行.
完全都是建基在摩西一個人的身上.
他們的難處都是很容易明白的.
因為縱然後期摩西已經設立了長老去幫助他們.
但是偏偏沒有人像摩西那樣.
在這個時候.
在《約書亞記》就會看到.
第一章.
第一章裡.
第一節.
「耶和華的僕人摩西死了以後.
耶和華對摩西的助手.
亂的兒子約書亞說」.
在這段經文裡.
你會看到.
我的僕人死了.
交代了摩西安息.
吩咐他.
就是藍色這些字.
現在你要起來約旦河.
主要是要領取上帝.
其實這個英.
早早已經被阿伯拉行.
是要給你們有一個肥美的迦南地.

$^{281}$成為你們的疆土.
很簡單.
主要命令他就是這件事情.
接著你會看到.
整段裡面.
都是上帝的說話.
一直叫他要剛強壯膽.
摩西也是大大的壯膽.
還要給他很大的保障.
當你一路看書的開頭的時候.
去到第九節.
你會發現.
當剛強壯膽.
不要懼怕.
一直貫穿著.
當剛強壯膽不要懼怕.
接著綠色的字.
你會看到.
耶和華的上帝.
是必定與你同在.
必不撇棄你.
是與你同在.
我很感受到.
當一個約書亞.
一開始的時候.
他一定會想到很大的壓力.
這班以色列人.
都不是每個人都這樣.
大家都見過.
他們整天埋怨.
怨這個怨那個.
一會兒又埋怨.
這樣.
約書亞.
他跟著他的師父.
他的屬靈領袖.
摩西的時候.
他很明白.
摩西他起碼有四十年的生活.
在皇宮裡面接受訓練.

$^{321}$軍事上.
很多的知識上.
他有他的訓練.
上帝又跟著再將摩西.
在四十年在曠野生活的時候.
他去看羊.
他已經是.
帶著以色列人之前.
他已經有四十年.
在曠野裡面牧羊.
在牧羊的生涯裡面.
你都可以說.
那段時間是上帝繼續的.
不單止是知識層面.
在皇宮那種造就.
去到野地.
野地你怎樣生存呢?.
這個讓他更加能夠.
可以帶著.
這群以色列人.
結果被神罰.
他們要四十年在曠野裡面.
走的時候.
走得好一點.
因為摩西受過野地求生的.
獨特的訓練.
但是約書亞.
雖然他也有四十年的.
曠野生活的學習.
但是很多的知識.
他是沒有的.
個人的不足.
自知之明.
是沒有的.
今天竟然.
怎可以成為.
第二個摩西呢?.
但是上帝根本沒有叫他做.
第二個摩西.
只是你帶領.

$^{361}$完成第二到第四節.
這些事情.
如果你再有機會.
再看約書亞.
代入你就是約書亞.
去看的時候.
你就會很深深感到.
他應該腳震.
心震.
總之很震.
不想這件事.
最好不要淋到自己身上.
怎樣去成為.
以色列人的領袖呢?.
很多時候.
我們的人生裡面.
其實都一樣.
很多突然奇妮的事情.
出乎我們意料之外的事.
甚至大氣候.
2019年.
我們怎會想到.
戴口罩戴到今時今日.
那時候還是新年.
不要緊.
冬天當是保暖.
戴著.
夏天應該沒有了.
誰知道我們已經度過了.
三年的夏天.
還是要戴口罩.
很多事情.
突如其來.
我們可能都會顯得.
很不安.
很徬徨.
緊張.
不知怎算好.
又或者在教會裡面.
當珍姑娘.

$^{401}$魯牧師.
弟兄.
姐妹.
我接下來的三個月.
想安排你做崇拜主席.
珍姑娘.
不要吧.
找她吧.
她呢.
是不是.
你可能第一次都.
不要吧.
或者邀請你做執事.
司事.
領唱.
去報道.
你都可能有很多事情.
最好不要吧.
找她吧.
你就會有很多.
推搪.
覺得自己不行.
但是人生裡面.
都有很多角色.
你是推不來的.
在座的.
很多都結婚了.
當你第一次.
不是做人丈夫那麼簡單.
做爸爸.
你有沒有上過爸爸課.
沒有的嗎.
你都要上的了.
就是一直做爸爸.
或者做媽媽.
去學習.
我們做兒女的.
學習如何孝順父母.
或者在座的.
有些做公公爺爺.

$^{441}$或者婆婆奶奶.
沒有一本.
做婆婆奶奶的手冊.
應該是怎樣.
做媳婦.
等等的東西.
做女婿.
人生有很多角色.
沒有的.
擺到你面前.
你就要上的了.
你進取一點.
就是一直擔任著這個角色.
一直的.
好好地學習.
旁邊有個好爸爸.
這些東西我學到的.
旁邊有個好的崇拜主席.
這樣做.
這樣表達好的.
一直去吸收.
一直去學習.
在教會裡面.
我們.
當然.
我們服侍神的人.
我們有上帝的呼召.
你要放下你的工作.
去吧.
去入神學院.
做就.
然後你是全時間的事奉.
這些事情.
都是一步一步學習.
一直這樣去.
當然上帝呼召我.
全時間事奉的時候.
我都有很多藉口.
我是家裡的大女兒.
家裡的經濟之主.

$^{481}$在當時.
我怎可能去呢.
家裡.
爸爸又說到時睡街了.
很多很多的東西.
但是.
當上帝再呼召你.
一次.
透過聖詩.
透過聖經.
一直呼召你的時候.
什麼經濟之主.
什麼我不行.
你就會.
放下下來.
一步一步走上去.
然後做宣教士.
哇!不行的.
買好衣服我又害怕.
你們記不記得買好衣服是什麼.
蟑螂.
那些襪子又害怕.
蜘蛛又害怕.
蛹.
蜜蜂.
你知道等完家女.
在這個區.
我回到香港.
和爸爸在天水圍的時候.
突然間蜜蜂飛來.
你都會很害怕.
這樣怎做宣教士.
加上我以前看的宣教士傳記.
很厲害.
過海.
怎樣過海.
斬六將.
我真是這個小小的馬.
真是很馬.
不行的.

$^{521}$但上帝就說你去吧.
你就是一步.
我們是個人.
我們是有不足的.
但當上帝好像約書亞.
來鼓勵約書亞的時候.
當剛強壯膽.
的時候.
上帝說.
必與你同在.
是非常好的.
弟兄姊妹.
足足十六年.
在泰國的宣教生活.
裡面.
我是沒有幹事的.
我是沒有童工的.
一個人台前幕後.
來事奉的.
很多事情.
一開始就沒有電腦的.
後來又有電腦.
又要弄一些.
powerpoint.
又要一些詩歌.
逐步.
短宣來.
教我.
這個應該怎樣搞.
不懂的先放著.
慢慢一步一步.
那些蟑螂.
那些買好衣服.
我就是.
不要怕 只要信.
不要怕.
因為.
你知道會大叫.
那些東西來到.
今天千萬不要嚇我.

$^{561}$我一定會大叫.
我一大叫可能你不會怕買好衣服.
但你被我的大叫嚇到就不好了.
你就會看到.
很多事情.
到今天都未能克服.
但是.
一個回頭.
原來我已經在泰國.
16年了.
甚至乎回到香港.
警察會叫我做.
本地跨文化工作.
我都不懂 都不知道泰國餐廳在哪.
在童工.
什麼開飯的app.
裡面就會看到.
泰國餐廳在哪.
就跟著地圖.
三不識七泰國餐廳.
就衝進去了.
你好 請問你們是泰國人嗎.
通常門面的多數不是.
有些是好的.
就說是廚房的.
你進廚房跟我們談.
就是這樣.
不打不撞.
一段時間.
你問我有沒有經驗做.
雖然有16年在泰國宣教的經驗.
我真的未試過進去廚房.
做這個工作.
做到.
真的在那個餐廳.
因為有一間是社企的餐廳.
做wizards.
吃什麼.
因為那間餐廳.
不夠人手.

$^{601}$星期天都會有泰國的節目.
去崇拜.
在這裡.
做的那些人.
都是我們希望他們是信主的一班人.
不過現在已經.
結業了.
竟然有試過做.
鼓勵他們.
做很技巧的.
於是我.
會召喚.
弟兄姊妹過來.
一起經歷人生的新角色.
做侍應.
怎料.
到了有局時間.
還要做廚房裡.
做翻譯.
因為那些單是中文.
廚師是泰國人.
什麼菜.
我告訴你.
泰式炒麵.
海南雞飯.
Kelvin雞.
我告訴他.
讓我見識不少.
很少見到.
廚師的武功.
如何用油鹽.
炒菜.
炒了一碟東西.
原來.
有一個叫做汁碼.
他怎會記得.
那些菜的材料.
我不是很會煮.
泰美食的.
學員都要看看.

$^{641}$這個有多少斤.
都要慢慢來.
他們很快就知道.
那是泰式炒麵.
就撿起來.
一碗東西.
廚師就會按著.
來煮.
其實真的不用怕.
當人生有很多.
角色突然間.
轉到你身上的時候.
你聽我說.
是不是很精彩?.
我當時很害怕.
但問題是.
人生又經歷了.
多一點東西.
又長知識.
是一件好事.
我們看到約書.
他是靠著.
上帝的應許.
他勝過自己的不足.
約書.
我們知道他有很大的壓力.
但他忠於.
上帝的吩咐的時候.
他勝過個人的不足.
新的一年.
讓我們知道.
我們仍然是有不足的.
做人做到老.
就學到老.
一定每一天.
都有很多新事物.
讓我們去學習.
不僅是依靠自己的力量.
而是依靠.
上帝的大能.

$^{681}$還記得我怎樣形容.
上帝的大能.
是勁勃勃.
是上帝.
勁勃勃的神.
我們這些裝模作樣.
後面的支持著我們.
是勁勃勃.
在約書.
這個中心裡面.
如果他對神.
稍微移開了.
對著那種中心.
他就勝過不到.
群眾的壓力.
很恐怖.
每一個國家.
政治體系裡面.
他很害怕.
你在這裡.
看到.
當時.
很大的一個.
或者很戲劇性的.
當時.
出埃及.
不,出埃及.
不久後.
摩西就派出12個支派.
讓他們.
找你們每一個族.
最好的一個代表.
12個人.
走去迦南地裡.
來做探子.
做間諜.
去探索.
這個地是甚麼一回事.
因為他們.
由出世.

$^{721}$到現在.
他們都是在埃及.
沒有見過世面.
不知道迦南地.
甚麼流奶與物之地.
怎樣流奶與物之地.
不知道.
但是當他們40天.
在這個迦南地.
去探索的時候.
他們真的看到.
他們的報告.
出來.
都是讚嘆.
這個地.
這個迦南地.
正如上帝告訴我們.
這是流奶與物之地.
是一個好地方的時候.
他們甚至.
他們的提子.
我們就這樣稱.
哪有這麼誇張.
他們抬著提子.
聖經的描述就是這樣.
他們抬著提子.
我記得有一個牧師.
他說.
他收集油票.
他說以色列有隻油票.
都是說到.
他們的提子很大顆.
做成油票.
你會看到.
他們的出產.
是很豐盛.
非常之好.
聽到這個地這麼好.
但是很可惜.
他們說那邊的.

$^{761}$雅納人.
你不要想像到.
科幻片.
那些突然間.
或者超人片.
突然間大我們很多倍的巨人.
這個巨人的意思.
就好像姚明.
如果姚明我也很高.
姚明在這裡.
可能很高.
很高大的一個人.
在這裡.
整群的.
雅納族人.
他們真的好像巨人.
我們就像一隻蚱蜢.
蚱蜢你也知道.
一隻手掌可以承著蚱蜢.
我們真的小到不得了.
我們一定不可以和他們打.
我們一定一敗塗地.
不行的.
你看到整群的.
十二個探子裡面.
十個都是這樣說.
十個.
當時只有.
迦勒和約書亞.
不是的.
上帝應許了我們.
我們一定得到.
上帝說過.
給我們的這一片.
的好地.
你想想十把口.
兩個人的口.
當時摩西亞倫.
已經.
我不知道有沒有寫.

$^{801}$已經.
趴在地上.
根本不知道怎麼辦.
就在這裡.
如果你要再看這個經文.
你回去看.
你會看到.
摩西亞倫在那段時間裡面.
因為十個探子這樣說.
傳媒還聽到.
個個都說.
我寧願.
回去埃及.
我寧願.
繼續做奴隸.
我們現在去送死.
我寧願身為奴隸.
都不願意死在.
這些.
阿拉人的手下.
你看到.
全島.
如果我們說.
少數服從多數.
民意.
是對的.
在聖經裡面有很多篇幅.
你會看到.
不是這樣的.
十個人的.
說話已經.
引起一群群眾.
依附在那裡.
摩西亞倫在聽.
他們的報告.
都已經俯伏在地.
其中.
兩個探子.
如果你是約書亞.
你敢不敢在這個時候.

$^{841}$開聲.
你開聲前都要再三思.
我所看的是不是幻覺.
我所看的.
阿拉人是否真的這麼高大.
你都會再.
在腦海裡面再重溫一下.
四十天裡面你看的是什麼.
我相信約書亞.
和加勒.
他會重溫一下.
但是.
他們更加.
溫習得到.
這是遠古.
上帝.
應許.
阿拉罕的.
一個神.
很實在的應許.
所以.
他們兩個.
在眾人面前.
撕裂衣服.
撕裂衣服是以色列裡面的一個文.
文.
風俗.
是一個很悲哀的.
告訴他們.
是呀.
上帝.
應許我們得這一塊地.
我們去便可.
當然.
上帝就有.
榮光照耀他們.
在這裡.
有些解經學家.
說不是有道榮光照射他們.
射到他們看不到.

$^{881}$好像.
保羅遇到異象.
不是那樣.
總之上帝在這個時間裡.
做了一些工作.
以致到.
這兩個人.
雖然是少數的說話.
但是這班.
大眾當中.
是叫他們站立不住.
今天.
很多人特別在教會.
如果我們一起.
去.
馬鞍山也好.
去大帽山也好.
隨便吧.
少數服從多數沒有問題.
但是.
我們教會或團體.
的聲音.
不是多數服從少數.
是真的.
好好地在上帝裡面.
去尋求.
這個尋求.
那個根.
回到上帝要多讀聖經.
很多時候弟兄姊妹問.
我聽不到上帝.
跟我說話.
是嗎?我沒有聽過.
對不起,你是否看聖經看得少?.
多看聖經.
上帝是.
透過他的說話.
跟你說話.
所以.
我們每天.

$^{921}$一起去看聖經.
是一件好事.
你個人的靈修.
或你個人也有目標.
有些人一年要看三次聖經.
你繼續.
按照你的目標而行.
但是你是.
洪恩嘉的一份子.
一起行.
我們是每天一章.
今天已經去到.
下個世紀.
你會看到.
每天一章.
都有好處.
當然不及一年.
看三次.
去到吃魚.
吃魚的地步.
但是你起碼.
每天一章讀經.
讀經不像靈修.
不是慢慢.
沉澱,當然有時間可以這樣做.
讀經就真的.
讀了它.
如果有些.
看字有些困難.
其實網絡.
裡面也可以聽.
聽那段經文.
你可以讓速度.
慢一點再聽.
或者再聽一次.
一天,保證你已經聽完.
《約書亞記》.
一天,保證你已經聽完《時書記》.
不過我們慢慢來.
一天一章,按照進度.

$^{961}$一起去的時候.
你知道一起去.
是很感恩的.
所以.
上帝應許.
約書亞說.
我必與你同在.
這是一件好事.
少年的時候.
我們女孩一去洗手間.
就喜歡結伴同行.
去洗手間.
我不用去,我去完了.
去吧,陪我吧.
她喜歡這樣.
我以前也不覺得這樣好.
但是在泰國.
16年裡面.
我就越來越覺得有人.
與你同行,哪怕他不懂.
泰文,都是好的.
真的好很多.
例如.
總之就是好很多.
譬如你去.
一個人去.
一些地方.
很簡單.
譬如我們.
已經拿著很多東西.
你想去洗手間.
你有一個人.
可以幫你看著.
你可以.
很輕鬆地去.
過去很多的短孫.
與我一起.
有時候可能是一兩個姐妹.
他們來探望我.
就完了短孫.

$^{1001}$那七天,他們可能會.
多留十天,一起就.
很不同了.
雖然都是我開車.
他不懂,但是一起.
與他探訪.
雖然他不懂說,但是他的見證.
我就可以幫他翻譯.
又多了一個.
見證,又多了.
另外一個的經歷.
去告訴.
一些朋友.
知道,是一件好事.
所以.
你會看到.
約書.
杜杜在.
群眾的壓力.
很不容易,但是他勝過了.
今天.
這個就是當時.
在民數記的時候.
說到.
這是流奶與墨的地.
是上帝喜悅我們.
我們一定可以去領取.
這個地.
今天,世界的聲音.
就好像群眾的壓力.
很大.
但是.
今天的你和我.
求神讓我們在上帝的.
話語裡站立得穩.
你不要說.
在淑寧真賤這個層面.
很簡單,當時.
剛剛戴口罩.
你家裡明明已經有兩排.

$^{1041}$衛生紙.
廁紙.
但是新聞.
這裡搶到一片空了.
這裡的鏡頭.
你明明不用搶.
不是搶.
我們說,這麼巧.
有.
我現在不是搶買.
你都會.
我們真是一個羊群.
所以上帝形容我們是羊.
我們真是羊群.
多數人這樣做.
我們都會跟著.
我們就要留意.
求神幫助我們.
很多人會告訴你.
不要說今天是初一的主日.
每一個的主日.
我相信.
都有很大的聲音會告訴你.
你這麼難得.
才可以休息一天.
好好休息.
回來做什麼?.
現在是疫症時間.
你不就開個屏幕.
在家裡.
敬拜.
你開著電腦.
在屏幕上敬拜.
跟我們現場.
一起讚美上帝.
你會聽到.
你的聲音.
還有弟兄姊妹的聲音.
那種敬拜是很不同的.
為什麼人家要浪費錢.

$^{1081}$去買演唱會的票.
是現場氣氛.
今天.
主日裡.
我們去享受.
敬拜上帝的現場氣氛.
是貴一貴都買不到的.
是很方便.
不用周車輪動.
按一個按鈕.
是看不回來的.
再加上.
如果再早一點.
你看到領詩的形象.
唱完一唱.
看到詩琴.
他們一早回來.
準備好場地.
昨天又在插花.
IT部忙於.
做很多的事.
這些一切一切.
你在網絡裡面的.
弟兄姊妹是看不到的.
你真的要回來.
洪恩家才看到.
跟著.
很多弟兄姊妹今天回來.
恭喜恭喜.
你回來已經賺取了很多安多芬.
那種喜樂.
就是這樣.
已經有了.
今天還未到按下按鈕.
你自己不要.
自己按了自己的按鈕.
求神幫助我們.
求主給我們忠於上帝的說話.
讓我們勝過群眾的壓力.
忠於主的應許.

$^{1121}$讓約書亞勝過群眾的壓力.
求神給我們有很好的跟從神.
千萬不要跟神疏遠了.
上帝給了我們很多的應許.
我必不撇棄你.
無論在舊約.
新約都有的.
我們今天也有位泰國姊妹.
她一年到晚都沒有休息.
星期日都要工作.
今天農曆新年初一.
大部分的公司都要關門.
今天立刻回來敬拜神.
不要跟神疏遠了.
可能你因著工作的緣故.
沒有辦法現場感受崇拜.
你仍然可以讀經.
每天的靈修.
每天去親近神.
我們求神給我們忠心到底的.
去跟從神.
讓神讓我們去過每一個關.
每一個關都不容易過.
剛才我們唱的那首歌.
《數串主因》.
你看到第一第二句都很沉重.
你遇見虎狼.
好像大波浪.
當你憂愁喪膽.
幾乎好像要絕望.
這首歌應該很悲情.
為何唱到最後這麼興奮呢?.
《數串主因》.
關關不容易過.
回望過去.
你經歷過的.
你過了.
你經歷了.
上帝與你同在歷地.
經歷了.

$^{1161}$求神幫助我們.
繼續與神親密.
一生裡面都很多衝擊.
書亞亦不例外.
到了晚年.
這段話是語重心長的.
向那群以色列人說.
你們覺得事奉耶和華不好嗎?.
好吧!你們就去選擇吧!.
但如果你選擇了耶和華神.
你就好好地事奉祂.
今天初一.
這麼難得的喜慶日子.
你也回到洪恩家.
一起去敬拜的時候.
你是否已經選好了?.
選了上帝是你唯一的救主?.
如果是的.
就讓我們.
好像約書亞說.
我必定事奉耶和華.
倘若你是一家之主.
或是在你的單位裡.
更有力的話.
至於我和我的單位.
我們一定要事奉耶和華.
我相信珍姑娘,魯牧師.
站在洪恩家一定可以這樣說.
我和我的洪恩家.
我們一定要事奉耶和華.
求神幫助我們.
在約書亞.
我們看到人性有很多不好的地方.
雖然在約書亞記.
我們看到阿剛犯罪.
戰戰兢兢.
耶利哥城打得不錯.
這樣也打贏了.
但接著去艾城.
大家看到.

$^{1201}$派兩個探子去.
探子說我們少許人都可以打敗他們.
已經輕敵了.
這就是人性.
勝完一個大仗.
這些算是甚麼城.
結果他們就慘敗收場.
不是在乎人多人少.
而是因為裡面有人犯罪.
神不喜悅這件事.
接著也有埋怨.
在戰爭裡面.
約書亞也經歷了很多衝擊.
我們人生何嘗不是呢?.
最近我多陪一個姐妹.
這個姐妹在三個月前.
她的女兒.
只有一個女兒.
因為血癌很突發.
去了醫院兩個小時.
就已經安息了.
剩下一個九歲大的女兒.
和七歲的兒子.
很忙碌的丈夫.
在海關工作.
我的姐妹.
安息了一個星期.
陪她.
她說.
每個人都叫她不用傷心.
不傷心是假的.
只有我們兩個.
我心裡好像撕裂了一樣.
我睡不著.
我真的要給醫生開安眠藥.
吃一點點來幫自己睡覺.
但還要振作.
因為兩個孫子.
不能什麼都不管.
兩個孫子都要去看.

$^{1241}$但同一時間.
她本身住荃灣.
女兒的一家在樂富.
她家在荃灣.
丈夫每天要洗兩次腎.
丈夫比較不小心.
所以她兩邊走.
這樣又忙.
心裡又撕裂.
很不容易的.
她都是靠什麼呢?.
每天讀神的話.
不然不夠課.
因為過去這個姐妹.
是幫辦.
很深刻的一個姐妹.
當時我還在香港.
去募會的時候.
她的叔叔.
說:高姑娘.
我嫂子病了.
你可不可以去醫院探望她?.
我們做傳道人.
很樂意去探望她.
我記得那幅圖畫.
無論怎樣跟她說.
怎樣引她注意.
她的眼睛看都不看你.
不看著窗戶.
這樣的開始反應.
後來她信了主.
還要去泰國探望我.
還要多留一點時間.
所以她說的東西.
很道地.
每天看聖經.
不夠課來鼓勵自己.
真是.
每天看聖經.
不要覺得.

$^{1281}$看了聖經.
真姑娘就得到很多褒勵.
天上很多獎賞.
沒有的.
可能上帝都認為.
不會因為你讀聖經.
會加了我的分數.
不會的.
你讀聖經.
是加在你的份上.
你的福氣淋到你個人身上.
每天聖經的說話.
安慰.
或者聖經裡面都有見證.
就是這樣.
拉著我.
拉著我走每一步.
姐妹說不夠課.
你今天是否很夠課呢?.
要知道你不看聖經了.
求主幫助我們.
主的說話.
讓我們勝過一生的衝擊.
你會綜合去看.
你會看到約書亞的秘訣.
就是約書亞對以色列人說.
你們近前.
來到這裡.
你要聽耶和華.
你們神的說話.
今天你們和我有沒有聽神的說話呢?.
我們的見證.
有沒有很好呢?.
我們真的天天去敬畏上帝嗎?.
上帝很愛我們.
但從聖經裡面看.
上帝是有想有佛的.
聖經裡面.
你會看到.
約書亞也有說.

$^{1321}$你不好好的.
跟從神.
上帝照樣賜福.
也會照樣賜禍.
可以讓你像阿剛那樣.
速速滅亡.
不止祂一個.
還要祂一家.
今天上帝這麼憐憫我們.
讓我們得著這個福音.
求主讓我們好好地持守.
天天敬畏上帝.
大大的壯膽.
大大的壯膽的緣故就是.
上帝的說話.
放在我們心上.
讓上帝的說話.
可能你會說.
我不是很記得.
不要緊.
每天去看.
每天看到的.
就記住.
記住一小時,兩個小時.
忘記了.
有時間再看.
緊緊地跟從神.
聽從神的說話.
你也會看到.
真的聽從神的說話.
你會看到最後.
約書亞在世的日子.
他死了之後.
那些知道耶和華.
為了以色列人所做的一切的人.
他們都是很好的事奉神.
但是你會發現.
因為我們開始看《時事記》.
那一代完結之後.
信仰叫你緊緊地跟從神.

$^{1361}$是逼不了你的.
上帝的繩子是持繩外索.
不是那條俘虜的繩子拉著你.
是持繩外索.
讓你一步一步地跟從神.
你會看到《時事記》.
人的反叛就是這樣.
一會兒信.
一會兒不要.
一會兒就同流合污.
一會兒就很多很多事.
《時事記》很精彩.
看到人性.
看到我們人的軟弱.
但我們求神幫助我們.
好像約書亞那樣.
忠於上帝一切的說話.
應許和命令.
所以我仍然是.
你說我錯也好.
怎樣也好.
我覺得約書亞.
他真是萬事勝意.
所有事情.
都是上帝你的旨意.
願意成就在我身上.
所以新的一年.
求主祝福大家.
願意你成為一個萬事勝意的.
上帝的兒女.
上帝祝福大家.
\newpage



\section{哥林多後書 8:9}
\label{sec:_P22pS_FVkg}
\textbf{2023.01.29 《無比恩典》哥林多後書 8\_9 (和修版) 陳一華牧師}
\newline
\newline
連結: \href{https://youtube.com/watch?v=_P22pS_FVkg}{\texttt{https://youtube.com/watch?v=\_P22pS\_FVkg}} ~~~~ 語音日期: 2023-02-10
\newline
\newline
\hyperref[sec:8tGomcGt7oY]{\small{< < < PREV SERMON < < <}}
~
\hyperref[sec:index_chronic]{\small{[返順時目]}}
~
\hyperref[sec:index_scriptual]{\small{[返順卷目]}}
~
\hyperref[sec:CnIFsRI_OWk]{\small{> > > NEXT SERMON > > >}}
\newline
\newline
哥林多後書 8:9
\newline
\begin{longtable}{cl}
\hline
\hline
章節 & 經文 (和合本修訂版)\\
\hline
8:9 & \begin{tabularx}{0.7\textwidth}{X} 你們知道我們主耶穌基督的恩典：他本是富足，卻為你們成了貧窮，好使你們因他的貧窮而成為富足。 \end{tabularx} \\ \\
[1ex]
\hline
\hline
\end{longtable}
$^{1}$弟兄姊妹早安.
剛才聽到潤昌弟兄的分享.
今天你們回崇拜豐不豐富?.
又領聖餐又聽到見證.
還有兩樣意外的收穫.
第一就是可以看到牧師的真面目.
因為由今個星期開始.
負責宗教場所的就可以不用戴口罩.
所以三年來都沒有讓你們看到牧師的真面目.
今天你們可以看到牧師的真人面目.
另外第二樣你們的收穫.
如果你們上個星期年初一有參加新春崇拜的話.
我指的是感受堂那邊.
到今天八天內.
大家就聽了牧師兩次的講道.
這個也是意料之外.
我沒想到這麼短時間講兩次.
今天我和大家分享的題目叫做無比恩典.
剛才多謝輝英幫我們已經讀了出來.
很短的經文我們試試一起讀一次.
或者我們看看哥林多後書的第八章第九節.
我們彈一彈經文出來.
預備起 你們知道我們主耶穌基督的恩典.
祂本是富足 卻使你們成了貧窮.
好叫你們的貧窮而成為富足.
這段經文的背景就是因為保羅提醒.
哥林多的教會原本他們要去為慈惠的奉獻.
要努力的 要有一個好的捐贈.
奉獻支持其他地方的需要.
不過不知道為什麼哥林多教會很勒羅.
所以他引用了另一間的教會.
馬其頓的教會去鼓勵他們.
因為馬其頓教會在他們極困難的時候.
都願意樂意積極的去奉獻.
所以用這個例子去鼓勵他們.
不單用這個例子.
最後出了一個王牌 就是這段經文.
這段經文說什麼呢?.
它說你們知道耶穌基督的恩典.
各位聽眾.

$^{41}$耶穌基督的恩典是一個什麼樣的恩典呢?.
我們新年和拜年很多時候都會說到.
祝你福杯滿日.
主恩是什麼?.
滿滿.
或者主恩是什麼?.
滿日.
我們祝人主的恩典是豐富的.
到底主耶穌基督的恩典是什麼呢?.
今天我們要思考這個課題.
大家都知道.
耶穌基督的恩典.
這個恩典很重要.
為什麼呢?.
因為這是一個基礎性的恩典.
如果我們有這個基礎性的恩典的時候.
我們將來再得到其他的恩典.
在上面建造的時候.
這個恩典是鞏固的.
是穩陣的.
但如果我們沒有了這個基礎性的恩典.
我們有時候得到的恩典.
都可能會不穩陣.
所以這是什麼恩典呢?.
這裡說得很清楚.
這個恩典是說到耶穌原本是負足的.
但因為你和我的原本就成為了貧窮.
又因為你和我因著耶穌的貧窮.
而成為負足.
讀起來這節聖經.
有沒有一種味道.
就好像交叉異形換影一樣.
有沒有這樣的印象?.
以前有一套劇講打排球.
講到異形換影.
有沒有這樣的印象?.
好像沒有什麼人有印象.
有印象的話應該差不多一年多了.
鬼影變幻球.
記不記得?.

$^{81}$交叉異形換影的.
原來耶穌的負足.
因我們成為了貧窮.
但你要記得.
這個過程裡面.
是有四個轉變當中的.
牧師今天和大家.
很快出去看看這四個轉變是什麼轉變.
第一個轉變是什麼轉變?.
是位置的轉變.
這個位置的轉變就是說.
耶穌基督要從天上的位置.
降下來.
成為一個地上的位置.
這是很不簡單的.
換句話說.
祂從天上最高的.
降下來是什麼?.
降生在馬槽裡面.
如果做生意的人就明白.
當你做得很興旺.
賺大錢.
發達的時候.
你買一間二千多呎的豪宅.
你很開心.
住在裡面很容易適應的.
但如果因為疫情的緣故.
生意受到打擊.
不是那麼好的時候.
生意失敗了.
你就要將你的豪宅拿出來拍賣.
之後你找了一間屋.
可能是二百呎的屋.
我想問大家.
由二千呎降下來住二百呎.
辛苦嗎?.
是不是很難適應?.
這個例子也不能夠很充分地表現出.
耶穌基督為你和我的犧牲.
由最高降下來到最低的位置.

$^{121}$所以各位電視節目.
這是第一個轉變.
第二個轉變是什麼呢?.
就是身份的轉變.
身份的轉變.
這是第一個轉變.
現在我們看第二個轉變.
就是身份的轉變.
這個身份的轉變.
就是說原本我們的主耶穌.
祂是什麼呢?.
祂是有天上尊貴的身份.
但現在我們的主耶穌.
祂是因為你和我的緣故.
取了一個普通人的身份.
而且還是竇木老的兒子.
木匠的兒子.
大家有沒有留意到?.
就是說祂的身份.
是一個很極大的轉變.
你也想不到是這麼厲害的.
或者只能夠舉一個很簡單的例子.
就是我們上一任的搭手.
由去年六月三十日.
他卸任之後.
我想問大家這大半年以來.
你在電視新聞螢幕當中.
有沒有看到他的蹤跡?.
很少.
為什麼?.
當然有很多原因.
他沒有曝光在媒體當中.
但我相信其中一個原因.
他要適應.
怎樣由七百多萬人之上的頭.
一夜之間.
他要回復做回一個平民的身份.
他要去坐公共交通.
如果沒有人開車的話.
他要去街市買菜.

$^{161}$他要去超市買東西.
回復一般人的身份.
各位各位頂智慕們.
這個身份的轉變.
是不是一個很難的適應?.
所以你看到他需要時間.
調整一下自己的心態.
所以各位各位頂智慕們.
我們的耶穌基督比祂更加厲害.
祂是天上尊貴的身份.
但竟然來到這個地上.
成為一個普通人的身份.
而且是人當中最卑微的人.
第三個轉變是什麼呢?.
第三個轉變就是生活的轉變.
各位頂智慕們.
耶穌在高天上.
祂的生活是否享有一切的榮耀尊貴?.
而且又有天使天君來到.
去侍候祂.
所以我們的主耶穌.
祂的生活是一個很榮耀的生活.
但是一夜之間.
當耶穌要決定成為人.
來到這個地上.
來完成上帝的使命的時候.
祂的生活就變成一個什麼呢?.
住在罪惡的世界.
而且和很多貧苦的人.
罪人來生活.
祂又住在一個罪惡的世界裡面.
各位頂智慕們.
這種生活帶不帶改變?.
很大的改變.
寧實醫院的創辦人.
我不知道大家知不知道他的名字是哪一位.
叫做司武道教士.
有沒有聽過?.
司武道教士是挪威的一位傳教士.
他來到中國大陸宣教.

$^{201}$後來又去到香港宣教.
他來的時候.
他不是選擇一個很富庶的地方.
他去到現在.
名字很好聽.
叫做張君傲.
但以前的名字.
你知不知道叫什麼?.
調景嶺.
讀歪一點.
還要論盡.
頂智慕們.
他在那裡去做上帝的工作.
去關心醫治患了肺癆的病人.
你說當年肺癆可不可怕?.
很可怕.
死了很多人.
醫學又不昌明.
又沒有藥物.
但這位司武道教士.
他有一點點護士底.
連同其他志同道合的弟兄姊妹.
開了一個簡單的診所.
來照顧關心醫治患了肺癆的病人.
結果他的行徑激勵了很多人.
而成為了今天很出名的靈實醫院.
頂智慕們.
他原本可以在挪威過去很舒適的生活.
但他選擇什麼?.
來到香港貧民窩的地方.
做宣教的工作.
為什麼?.
是因為他學校耶穌基督的榜樣.
他將他的生活模式完全改變.
目的是什麼?.
目的是要將自己成為一個神所使用的工具.
去祝福其他的人.
令不令我們敬佩?.
各位頂智慕們.
這就是主耶穌的豐富足.

$^{241}$因你和我成為了貧窮.
第四個轉變是什麼?.
第四個轉變就是一個生命的轉變.
就是說耶穌基督知道他要來到這個世間上.
最後他要犧牲他自己寶貴的生命.
所以他已經做好一個準備.
他雖然知道這個行動是一個很大的掙扎.
很大的痛苦.
所以大家都很熟悉.
在赫西瑪利園的一幕.
耶穌在祈禱的當中.
他流出的汗好像血點.
有印象嗎?.
這就是什麼?.
表示到耶穌說.
然而不要照我的意思.
乃要什麼?.
照你的意思去做.
因為他已經知道他來到世界上.
他要面對的是要犧牲他的生命.
聖經說.
若不流血.
罪就不得赦免.
所以他要成為一個祭身獻賞.
來救贖世界上每一個犯罪的人.
難怪他出道的時候.
施洗若寒就老早告訴人們.
看吧,大家看看吧.
這個人很特別.
為什麼?.
神的羔羊.
除去世人罪孽的.
各位聽眾.
現在我們明白這四個轉變是什麼?.
就是耶穌基督的父祝.
因你和我成為了什麼?.
貧窮.
所以這是第一個部分.
是耶穌降生成人.
面對的四個轉變.

$^{281}$接下來我們看第二個部分.
就是.
聖經說.
我們因耶穌的貧窮.
而成為了什麼?.
父祝.
所以第二個部分是什麼?.
第二個部分是你和我.
需要來回轉.
悔改.
接受上帝的救恩.
以致我們能夠經歷重生得救.
各位弟兄姊妹妹.
當你和我都接受了耶穌基督.
成為上帝的兒女的時候.
重生得救之後.
上帝給我們的恩典.
是一個豐盛的恩典.
有哪些恩典呢?.
我們去看.
第一個的恩典.
就是新造的人.
這個恩典寶不寶貴?.
很寶貴.
哥林多厚書告訴我們.
若有人在基督裡.
他就是什麼?.
新造的人.
舊事已過.
都變成新的了.
我在戒毒村做導師的時候.
認識一位弟兄.
在慈雲山黑社會裡很有地位.
這位弟兄戒了毒.
信了耶穌之後.
他開始過正常的生活.
他講見證.
他說我今天非常非常感恩.
我不想到我是一個新造的人之後.
這麼多人願意去接納我.

$^{321}$這麼多人都願意去尊重我.
因為我出來之後一無所有.
但是別人看到我的生命.
是一個新造的人.
所以放膽介紹工作給我去做.
我現在有份工作去做.
很感恩.
也能夠回到教會.
有教會的生活.
有弟兄姊妹接納我.
所以他覺得.
成為新造的人.
這個身份和他以前在黑社會的身份.
完全是判若兩人.
他感不感恩?.
他非常感恩.
他覺得人生有希望.
人生有盼望.
各位弟兄姊妹們.
你和我都是.
當我們成為新造的人.
這個寶貴的恩典.
我們要牢牢的持守著.
第二個恩典是什麼恩典呢?.
第二個恩典.
就是一個尊貴的身份.
我們看到很多皇室的新聞.
很多人都很喜歡的.
尤其是一些人很蛛絲八卦.
很留意皇室的新聞.
最近皇室裡面有些人有寫書.
有很多採訪.
各位弟兄姊妹們.
一個皇室的成員.
都這麼被人關注的時候.
請你記得.
他只是一個國家裡面的王子.
或者公主.
但當你和我成為上帝的子女之後.
我們就是萬王之王.

$^{361}$萬主之主的子女.
所以我們的身份.
是不是更加的尊貴?.
更加的寶貴.
不覺得可以給我們前書告訴我們.
唯有我們被揀選的族類.
是有君尊的祭司.
有印象嗎?.
是聖潔的國度.
是屬神的子民.
所以當我們擁有這麼尊貴的身份的時候.
這個恩典.
是不是很大的恩典?.
從今之後.
我們不是外人.
而是神家裡的人.
所以我們為這個恩典獻上感恩.
第三個恩典是什麼呢?.
就是一個豐盛的人生.
當我們重新得救之後.
原來聖靈是內在我們的心裡面.
聖靈內在我們的心裡面.
一方面是成為我們的導師.
提醒我們.
好像提醒潤祥一樣.
不要衝動.
聖靈會做我們的導師提醒我們.
但另一方面.
聖靈也給我們恩賜.
讓我們得到上帝的恩賜.
就能夠事奉上帝.
我早前聽到一位姐妹的分享.
她說很感恩.
她說什麼呢?.
她發覺現在信了耶穌之後.
她的生活是很充實.
為什麼這麼充實?.
因為她覺得自己的時間很不夠用.
為什麼不夠用?.
原來她發覺上帝在她信了耶穌之後.

$^{401}$給她有新的生命.
又給她很大的恩賜.
去關心人.
去事奉上帝.
所以她發覺到.
原來她每天都很充實.
很忙碌.
她可以在家裡事奉.
她可以在教會裡事奉.
她也可以在醫院裡做義工事奉.
所以她告訴別人.
她現在信了耶穌之後.
她的生活是很充實.
活得很精彩.
為什麼?.
因為聖靈賞賜她有更多的恩賜.
事奉上帝.
各位電視節目的觀眾.
我們的人生.
我們的生命.
變得更加豐盛的生活.
所以這是上帝給我們無比的恩典.
我們再看最後一點.
上帝給我們的恩典是什麼?.
就是永生的福樂.
這一點我相信.
不需要我用太多的純涉去講.
大家都很明白.
原來我們成為上帝的子女之後.
在地上我們有一個很充實的生活.
有一個很精彩的生活.
之外.
原來我們知道.
將來離開這個地上.
上帝也為我們預備永恆的生命.
有一個永遠的家鄉為我們預備.
所以我們作為神的兒女.
我們在今生裡面.
在將來的天父的家中.
我們都是非常得到.

$^{441}$上帝的眷顧和恩典.
所以各位頂智妹妹.
我們有時候出席安息聚會.
我們心裡面.
不會好像一個絕望的人.
為什麼?.
因為我們知道.
我們對將來有盼望.
有希望.
我們和我們所至愛的人.
只不過是短暫的道別.
我們知道.
上帝賜給我們永遠的福樂.
所以各位頂智妹妹.
當你今天聽完牧師.
分享這四項恩典的時候.
我們感不感恩?.
很感恩.
現在你明白.
我們因耶穌的貧窮.
而成為什麼?.
成為富足.
所以各位頂智妹妹.
牧師今天和大家分享這個訊息.
到最後.
其實我們有兩個目的.
想你詳細記住.
這兩個目的是什麼?.
這篇道理提醒我們.
領受上帝無比恩典之後.
第一.
記得我們要常存感恩.
因為我們所得到的這些恩典.
我們其實都是很不配的.
大家認同嗎?.
我們是很不足.
我們有很多瑕疵.
很多問題.
何德何能得到這些恩典呢?.
原來完全都是什麼?.

$^{481}$是上帝的憐憫.
上帝的祝福.
臨到我們身上.
是耶穌基督的犧牲.
令我們得到這個無比的恩典.
當你常常思考的時候.
各位頂智妹妹.
無論你在順境.
在逆境當中.
你就學會凡事借恩.
常存感恩.
因為這是一個無比的恩典.
是我們意料之外.
得到的.
是我們很不配之下.
得到的.
第二個目的.
我想弟兄姊妹記得.
要廣傳福音.
為什麼?.
因為上帝給你有這個無比的恩典.
不是一個人享用的.
不是獨食.
而是將這個恩典.
去跟一些未信的朋友.
未得到的人.
去分享.
去見證.
讓他們都能夠享受到.
上帝為我們人類.
預備這一份無比的恩典.
所以我希望這兩個目的.
在我們今年的當中.
新春的裡面.
記得.
主要我們每一個人都去學習.
常存感恩.
要我們來到廣傳福音.
與人分享神的恩典.
我們低頭祈禱.

$^{521}$主耶穌我們非常多謝你.
我們禱告的時候.
要Hallelujah讚美你.
因為你的恩典.
是足夠我們使用.
你的恩典.
也是超過我們所求所想.
這一份就是無比的恩典.
求主耶穌.
讓我們能夠去思考.
這個基礎性的恩典裡面.
我們打好這個根基.
以致我們人生往後日子.
所得到的恩典.
都是一個很穩固的恩典.
多謝你.
讓我們能夠懂得常存感恩.
也能夠懂得廣傳福音.
與人分享上帝無比的恩典.
奉耶穌名求.
阿們.
\newpage



\section{士師記 19:1-30}
\label{sec:CnIFsRI_OWk}
\textbf{2023.02.05 《那時,以色列中沒有王!》士師記 19\_1-30 (和修版) 曾敏傳道}
\newline
\newline
連結: \href{https://youtube.com/watch?v=CnIFsRI_OWk}{\texttt{https://youtube.com/watch?v=CnIFsRI\_OWk}} ~~~~ 語音日期: 2023-02-10
\newline
\newline
\hyperref[sec:_P22pS_FVkg]{\small{< < < PREV SERMON < < <}}
~
\hyperref[sec:index_chronic]{\small{[返順時目]}}
~
\hyperref[sec:index_scriptual]{\small{[返順卷目]}}
~
\hyperref[sec:QG8KdqrlJlA]{\small{> > > NEXT SERMON > > >}}
\newline
\newline
士師記 19:1-30
\newline
\begin{longtable}{cl}
\hline
\hline
章節 & 經文 (和合本修訂版)\\
\hline
19:1 & \begin{tabularx}{0.7\textwidth}{X} 當以色列中沒有王的時候，有一個利未人寄居以法蓮山區的邊界，他娶了一個猶大伯利恆的女子為妾。 \end{tabularx} \\ \\ \relax
19:2 & \begin{tabularx}{0.7\textwidth}{X} 這妾對丈夫生氣，離開丈夫，回到猶大伯利恆的父家，在那裡住了四個月。 \end{tabularx} \\ \\ \relax
19:3 & \begin{tabularx}{0.7\textwidth}{X} 她的丈夫起來，帶著一個僕人、兩匹驢跟著她去，要用好話勸她回來。女子就帶丈夫進到父親家裡。女子的父親看見了他，就歡歡喜喜地迎接他。 \end{tabularx} \\ \\ \relax
19:4 & \begin{tabularx}{0.7\textwidth}{X} 這岳父，就是女子的父親，留他住了三天。他們在那裡吃喝，住宿。 \end{tabularx} \\ \\ \relax
19:5 & \begin{tabularx}{0.7\textwidth}{X} 第四日，他們清早起來，利未人起身要走，女子的父親對女婿說：「先吃點東西，加添心力，然後你們才走。」 \end{tabularx} \\ \\ \relax
19:6 & \begin{tabularx}{0.7\textwidth}{X} 於是二人坐下，一同吃喝。女子的父親對那人說：「請你答應再住一夜，使你的心舒暢。」 \end{tabularx} \\ \\ \relax
19:7 & \begin{tabularx}{0.7\textwidth}{X} 那人起身要走，他岳父挽留他，他就留下，在那裡又住了一夜。 \end{tabularx} \\ \\ \relax
19:8 & \begin{tabularx}{0.7\textwidth}{X} 第五日，他清早起來要走，女子的父親說：「來，請加添心力，留到太陽偏西吧。」於是二人一同再吃。 \end{tabularx} \\ \\ \relax
19:9 & \begin{tabularx}{0.7\textwidth}{X} 那人同他的妾和僕人起身要走，但他岳父，就是女子的父親，對他說：「看哪，太陽下山，天快晚了，你們再住一夜吧。看哪，太陽偏西了，就在這裡住宿，使你的心舒暢，明天你們一早起來上路，回你的帳棚去。」 \end{tabularx} \\ \\ \relax
19:10 & \begin{tabularx}{0.7\textwidth}{X} 那人不願再住一夜，就備上兩匹驢，帶著他的妾起身走了，來到耶布斯的對面，耶布斯就是耶路撒冷。 \end{tabularx} \\ \\ \relax
19:11 & \begin{tabularx}{0.7\textwidth}{X} 將近耶布斯的時候，太陽快下山了，僕人對主人說：「來吧，我們進這耶布斯人的城，在這裡住宿。」 \end{tabularx} \\ \\ \relax
19:12 & \begin{tabularx}{0.7\textwidth}{X} 主人對他說：「我們不可進入外邦人的城，那不是以色列人的地方，我們越過這裡到基比亞去吧。」 \end{tabularx} \\ \\ \relax
19:13 & \begin{tabularx}{0.7\textwidth}{X} 他又對僕人說：「來，讓我們到基比亞或拉瑪的一個地方住宿。」 \end{tabularx} \\ \\ \relax
19:14 & \begin{tabularx}{0.7\textwidth}{X} 於是他們越過那裡往前走，將到便雅憫的基比亞的時候，太陽已經下山了。 \end{tabularx} \\ \\ \relax
19:15 & \begin{tabularx}{0.7\textwidth}{X} 他們進入基比亞要在那裡住宿。他來坐在城裡的廣場上，但沒有人接待他們到家裡住宿。 \end{tabularx} \\ \\ \relax
19:16 & \begin{tabularx}{0.7\textwidth}{X} 看哪，晚上有一個老人從田間做工回來。他是以法蓮山區的人，寄居在基比亞；那地方的人是便雅憫人。 \end{tabularx} \\ \\ \relax
19:17 & \begin{tabularx}{0.7\textwidth}{X} 老人舉目看見那過路的人在城裡的廣場上，就說：「你從哪裡來？要到哪裡去？」 \end{tabularx} \\ \\ \relax
19:18 & \begin{tabularx}{0.7\textwidth}{X} 他對他說：「我們從猶大的伯利恆過來，要到以法蓮山區的邊界去。我是那裡的人，去了猶大的伯利恆，現在要到耶和華的家去，卻沒有人接待我到他的家。 \end{tabularx} \\ \\ \relax
19:19 & \begin{tabularx}{0.7\textwidth}{X} 其實我有飼料草料可以餵驢，我和你的使女，以及與我們在一起的僕人都有餅有酒，甚麼都不缺。」 \end{tabularx} \\ \\ \relax
19:20 & \begin{tabularx}{0.7\textwidth}{X} 老人說：「願你平安！你所需用的我都會給你們，只是不可在廣場上過夜。」 \end{tabularx} \\ \\ \relax
19:21 & \begin{tabularx}{0.7\textwidth}{X} 於是老人領他到家裡，餵上驢。他們洗了腳，就吃喝起來。 \end{tabularx} \\ \\ \relax
19:22 & \begin{tabularx}{0.7\textwidth}{X} 他們心裡歡樂的時候，看哪，城中的無賴圍住房子，連連叩門，對老人，這家的主人說：「把那進你家的人帶出來，我們要與他交合。」 \end{tabularx} \\ \\ \relax
19:23 & \begin{tabularx}{0.7\textwidth}{X} 這家的主人出來對他們說：「弟兄們，不要做這樣的惡事。這人既然進了我的家，你們就不要做這樣可恥的事。 \end{tabularx} \\ \\ \relax
19:24 & \begin{tabularx}{0.7\textwidth}{X} 看哪，我有個女兒還是處女，還有這人的妾，我把她們領出來任由你們污辱她們，就照你們看為好的對待她們吧！但對這人你們不要做這樣可恥的事。」 \end{tabularx} \\ \\ \relax
19:25 & \begin{tabularx}{0.7\textwidth}{X} 那些人卻不肯聽從他。那人抓住他的妾，把她拉出去給他們。他們強姦了她，整夜凌辱她，直到早晨，天色快亮才放她走。 \end{tabularx} \\ \\ \relax
19:26 & \begin{tabularx}{0.7\textwidth}{X} 到了早晨，婦人回來，仆倒在留她主人住宿的那人的家門前，直到天亮。 \end{tabularx} \\ \\ \relax
19:27 & \begin{tabularx}{0.7\textwidth}{X} 早晨，她的主人起來開了門，出去要上路。看哪，那婦人，他的妾倒在屋子門前，雙手搭在門檻上。 \end{tabularx} \\ \\ \relax
19:28 & \begin{tabularx}{0.7\textwidth}{X} 他對婦人說：「起來，我們走吧！」婦人卻沒有回應。那人就將她馱在驢上，起身回自己的地方去了。 \end{tabularx} \\ \\ \relax
19:29 & \begin{tabularx}{0.7\textwidth}{X} 到了家裡，他拿刀，抓住他的妾，把她的屍身切成十二塊，分送到以色列全境。 \end{tabularx} \\ \\ \relax
19:30 & \begin{tabularx}{0.7\textwidth}{X} 凡看見的人都說：「自從以色列人離開埃及地上來，直到今日，像這樣的事還沒有發生過，也沒有見過。大家應當想一想，商討一下再說。」 \end{tabularx} \\ \\
[1ex]
\hline
\hline
\end{longtable}
$^{1}$(聖母向大人致謝).
弟兄姊妹平安.
我們再一次一同去禱告.
(禱告).
天父求你今天開我們的眼睛.
開我們的耳朵.
讓我們看到你的作為.
明白你的心意.
讓我們知道神要我們思想什麼事情.
願聖靈在我們心裡動功.
帶領我們去回應神你自己的教導.
奉耶穌基督的名字祈禱.
(阿們).
我深信我們並不是常常去看這段經文.
在講壇上也不多.
今天跟大家看的時候.
真的有一些位置很想一同去思考.
今天是講「那時以色列中沒有和」.
這是《士思記》中心的一個要點.
「士思」時期所發生的事情就是.
「那時以色列中沒有和」.
《士思記》是接續《約書亞記》的一段書.
我們先看看整段經文的背景.
這個表格我發覺有些細節.
可能看不太清楚.
這個肯定是進入了迦南地.
打仗,分了地之後.
這裡有一段經文是說進入迦南地.
分地後牽涉到幾卷書都有提及.
一直到《奈及記》,《申命記》,《約書亞記》.
到之後,已經在迦南地裡.
已經分了地之後.
到《約書亞記》的時代過去後.
進入了「士思」的時期.
我們再看一看.
在這段時期,你會看到很大的以色列文的改變.
約書亞死了,和他那個時代的認識上帝作為的人都過去了.
在這個時期,新的一代,他們不認識神.
他們都沒有聽過上帝當時的作為.
來到這裡,有些地方值得我們思考.

$^{41}$為什麼到了這個時候,突然間所有東西都斷了.
沒錯,約書亞離開了.
他過世了.
他那個世代的人過去了.
但是為什麼呢?我們值得問一個問題.
就是信仰沒有傳承.
信仰沒有一代一代傳下去.
以致到那一代的人離開了之後.
他們在世世之後,新的一代就不認識神了.
沒錯,新一代是很大的問題.
同時,也在說家庭的傳承.
整個民族的傳承都出了問題.
沒有很好的一代一代.
所以今天我回看我們眾人.
無論我站在這裡說話,或者大家坐在下面聽.
其實我和你都有責任.
去將信仰一代一代的相傳.
不是普通的說.
而是很衷心的教導.
教導年輕人,教導下一代.
有子女的,教導你的子女.
育身的子女.
有屬靈的子女,教導屬靈的子女.
其實這是我們這一代的責任.
我深信無論是約書亞那一代的人.
是哪一代的人都好.
都需要將信仰傳承下去.
現在的新一代,他們不認識上帝.
接著他們開始事奉朱巴力.
他們事奉二幫的神.
大大的得罪上帝.
他們四處列國,人們拜什麼神明.
他們又去拜.
神非常的憤怒.
這個階段就是這樣.
接著神很憤怒.
又興起一些敵人去壓迫他們.
去管教這些以色列民.
是真的很嚴厲的壓迫.
是很嚴厲的管教.

$^{81}$所以以色列民在面對外邦敵人的攻擊.
那些壓迫的時候.
他們感到很辛苦,大叫救命.
向上帝祈求.
很辛苦啊,救命啊.
神在這個時候很憐憫他的百姓.
興起士師來拯救他們.
這個就是士師記的起源.
士師的興起正正是因為百姓犯罪.
得罪上帝.
上帝興起外邦人成為百姓的仇敵.
藉此要管教他們.
百姓叫救命.
神就透過士師去拯救他們.
士師有兩個很重要的角色.
第一個角色就是.
他會是一個判斷是非的人.
一個斷案的人.
譬如當中有什麼的一些.
我們社會上的一些案件.
不同的事情要判斷誰對誰不對.
通常人們去到士師那裡.
他就好像法官那樣去作出判斷.
因為其實士師的英文就是法官的意思.
所以士師是有這個角色.
但是更加重要的角色.
就是士師會帶領以色列人去打仗.
去對付那些外邦的仇敵.
所以他也有這個角色.
好了,說起.
我們當時以色列人最主要拜的.
其中一個是巴力.
另外就是艾斯塔祿.
通常他們拜這些都是跟我們.
他們肉身需要有些很實際的利益.
外邦人拜的神都是這樣的.
就是對他們肉身有些實際的好處.
他們覺得,例如巴力是掌管雨水的風暴的.
你靠著他,你討好他.
他就給你的農業,農產品都可以豐收.

$^{121}$所以當地人很喜歡拜巴力.
艾斯塔祿就跟我們的生殖歷很有關係.
是愛情和戰爭的女神.
還不靠他.
當時外邦人拜這些.
以色列人照樣去跟.
結果他們大大得罪神.
神興起的那些外邦的仇敵.
就在這幅圖裡面.
都說了有不同的外邦人.
你見到是遍滿了地圖.
遍滿了以色列的領土.
昔日,以色列人進入迦南地的時候.
神應許了,將整個地方給他們.
就算在約書亞時期.
他們打仗都還沒打完.
還有些地還沒得到.
當時上帝是答應過的.
會將這些外邦人,這些民族全部趕出去.
但重點是以色列人要跟從上帝的教導.
是忠心的跟從神,作神的指紋.
但到現在,你見到這麼多外邦的仇敵.
遍佈在這個版圖裡.
就更加證明了以色列人那種不忠心.
去背叛上帝,拜偶像.
所以上帝決定.
不會將這些仇敵從他們的應許之地挪走.
任由他們在這裡.
反而上帝會用得著這些人去管教他們的百姓.
事實時期的特色.
就好像這個循環圖一樣.
以色列人犯罪.
接著的戰罪,他們會受苦.
受外邦仇敵的逼迫,壓制.
接著他們向上帝呼求.
神仍然憐憫他們.
會興起事實帶領他們.
他們會得夢拯救.
和平一段時間之後.
他的死性不改又來犯罪.

$^{161}$於是這個循環又下降.
你看到這裡乘6.
整個《攝施記》的循環進行了6次.
就一直循環,又再循環.
這個稍稍看,有三大段.
有三大段就是一章一節到三章六節.
仍然是有講到征戰.
是怎樣為著他們的應許地去打仗.
是得著他們的應許地.
但背後亦都講到.
以色列人是速速離開他們的信仰.
背叛上帝.
令上帝非常難過,非常憤怒.
去到三章七節到十六章的三十一節.
是講不同的事司,他們的故事.
最後的,今天我們的經文是在這一段.
十七章的一節到二十一章的二十五節.
是講到事司時期,最黑暗的時期.
最黑暗的時期就是宗教人員.
都是深深的陷溺在罪惡當中.
是有令人非常憤怒的暴行.
是社會上的暴行.
甚至是由個人家庭的問題.
發展到支派之間的內戰.
這個時期是最黑暗的時期.
今天我們的經文其實是橫跨了三章.
十九二十到二十一章.
但是主要是圍繞了十九章.
大家可以去看後面的經文.
講到整個基比亞的事件.
在這個牽涉到真的令人感到很憤怒.
是不可思議的.
簡直是令人髮指的暴行.
亦因為這個暴行牽連引起內戰.
當我們去看這段經文的時候.
是很特別的.
是首尾呼應.
十九章一直到二十一章是首尾呼應.
十九章一節是這樣說.
當以色列沒有王的時候.

$^{201}$一開始是這樣說.
這個時間最大的特色就是沒有王.
沒有王又怎樣?.
「國人任意而行,我想點就點」.
其實不是沒有王.
自己就是王.
其實不是沒有王.
有人告訴你我沒有宗教信仰.
其實那個也是宗教信仰.
他相信自己.
我相信自己,我作主.
現在這個時候沒有王,誰作主?.
我作主.
我覺得是好就好.
我覺得不好就不好.
我想怎樣就怎樣.
所以我作主,我就是王.
所以那個時候以色列沒有王就很厲害.
整個周街,整個社會,每個人都是王.
想怎樣就怎樣.
去到二十一章二十五節最後那一節.
是首尾呼應.
最後呼應是這樣說.
那時以色列沒有王.
「國人照自己眼中看為對的去做」.
頭頭尾尾的句子總結了.
沒有王絕對不是一件好事.
是「國人非常自我自私」.
在一個很大很深的罪惡裡.
想怎樣就怎樣.
在這段經文,十九章的經文裡.
大家留意到一個很大的特色.
大部份的人物都是不知名的.
他提到「妾侍」,說那個「利未人」.
其實沒有說他叫什麼名字.
你看很多人物都是沒有名字的.
在聖經裡面,人物沒有名字有很多原因.
很多原因,以至到聖經不記那些人的名字.
但這裡的原因是很特別.
這裡說的每個人物都很有代表性.

$^{241}$說的那個妾侍.
不是說一個妾侍.
說的妾侍就好像說一個女性.
任何一個人在妾侍的位置.
任何一個女性在她的位置都會這樣的經歷.
說的那個利未人就不只是特別的利未人.
而是擁有這個身份的人.
都是差不多的,這樣的生命質地.
是會這樣的狀況,會經歷這些事.
而說的基比亞的暴徒都是一樣.
每個人都具有一個代表性,就是這樣.
不是這個,是那個,同樣都會這樣發生.
所以你說當時的情況是多麼的嚴峻.
不是只是一個單一的事件.
是整個社會就是這麼黑暗.
是這麼墮落,腐化.
在這個地圖上有些位置值得我們去看看.
我很快的說一說那個故事.
那個故事當時是說有一個利未人.
本身是住在以法連的山區.
雖然圖片有少許的小.
以法連山區就在你的圖片左上角的位置.
他原先是住在那裡的.
他說他娶了一個伯利恆的女子為妾.
伯利恆就在你的圖片右下角.
很下面的那個位置.
所以也有一段距離.
他娶了那個地方的一個女子為妾.
很特別,經文說什麼?.
不是為妻,就是他有太太.
他再娶為妾.
利未人是什麼人來的呢?.
事奉人員,是吧?.
就是這個聖殿裡面的事奉人員.
這麼特別.
這個應該一般人來看他.
應該是一個很認識神的人.
是事奉神的人,是吧?.
事奉神的人,這麼特別.
一開始的經文先聲奪人的不是他的事奉.

$^{281}$不是他和上帝的關係有多緊密.
他的事奉有多好.
一說就說他婚姻.
他是娶妾侍的.
不止一個女人.
事奉人員又不止一個女人.
這個立刻打個問號.
他娶了這個女子為妾的時候.
說這個妾侍對丈夫生氣.
這裡有兩個不同的翻譯.
一個是真的生他氣.
一個是說這個妾是犯奸淫.
是吧?因為這個緣故.
於是他就要麼生他氣.
要麼犯奸淫.
他就離開丈夫.
回到伯利恆他自己的爸爸的家裡.
住了四個月.
總之兩夫婦發生了一些問題.
夫婦關係裡.
大家分開了.
分開了也不是利未人馬上去找他的女人.
等了四個月.
等了四個月才去找.
找到之後.
岳父大人就好好地款待他.
一路款待他.
終於住了多久呢.
在那裡住了.
你看看.
一共住了四天.
到了第五天.
還是不行了.
岳父大人每天都熱情款待.
每天都挽留.
到了第五天.
一路去到黃昏.
太陽偏西的時候.
終於利未人覺得不行.
今天真的要走了.

$^{321}$他選擇了一個非常不好的時間離開.
一個已經太陽偏西.
快要日落的很黑的時間.
其實這個時間絕對不適合夜行.
不是今天我們現在這麼現代化的一個世界.
社會交通很發達.
周圍燈火通明.
又比較安全.
那個時間.
你要經過曠野不同的地帶.
這麼危險.
偏偏他就選擇這個時間.
晚上就出發.
他出發之後.
首先因為已經很晚了.
他就建議不如去耶布斯過夜.
但是當時經文也有說.
耶布斯這個地方是屬於外邦人的地方.
耶布斯人就不是以色列人.
是外邦人.
他們在管理.
當時這裡是耶路撒冷.
他們管理這個地方.
那個利未人就很有趣.
很堅持要住在自己同胞的地方.
覺得是很不同的.
屬於自己以色列民的地方.
是真的特別好一點的.
屬於上帝百姓的地方.
所以這些外邦人的地方不要去.
我們繞過耶布斯.
就去基比亞.
那裡一定好.
那裡是屬於上帝的地方.
正常來說真的應該會好一點.
屬於上帝的地方.
但結果是不是呢?.
大家剛才讀過經文.
去到那個地方.
沒有人接待他們一家.

$^{361}$當時那些人外出都很期望.
很需要.
人家是在以邦接待自己.
我相信當時的酒店.
都不是那麼流行.
都很靠當地人.
但是沒有人接待他們.
結果他們就要坐在城裡的廣場上.
已經很晚了.
很危險.
沒有人接待.
在這個時候就有一個老人家.
在田間做完事就回來.
原來這個老人家和利未人同鄉.
大家都是以法連山區.
這個同鄉很好.
他就接待那個利未人.
和他們的妾.
接待之下.
大家就很好.
一起吃喝聊天.
這一晚終於有個落腳位.
正正當大家聊天.
吃飯.
覺得很好的時候.
就是這個時候.
一群暴徒殺到.
這群人就在城中的霧瀨.
圍著房間.
不斷敲門.
對著老人家說.
把你家裡的人帶出來.
小心看.
把誰帶出來.
是不是把妾帶出來.
不是.
那個男人.
把利未人帶出來.
做什麼.
和他聊天嗎.

$^{401}$不是.
是做了一件惡事.
和他交合.
是在性上的交合.
這個有沒有令大家想起創世記的一段經文.
所多瑪的經歷.
天使去探羅德的時候.
去到羅德家裡.
發生什麼事.
城市裡的人走出來.
圍繞他的房間.
又是想和天使做一些事情.
行惡.
其實很相似.
這一群人.
原來他們很喜歡同性戀.
不只是喜歡同性戀.
是想強行來.
是在污辱男性的同性.
很恐怖.
可以想像那個罪惡的心.
在這個時候對老人家說.
老人家怎麼回應.
大家記不記得.
弟兄們不要做這個惡事.
老人家覺得這件事非常不好.
你們不要這樣做.
之後很特別.
他說了一些話.
你看看我有個女兒.
還是處女.
還有這個人是妾.
這個人妾.
我都可以把他們拿出來.
任你們.
怎樣對待他們.
就是對這個女人不好.
我講了什麼身份.
總之對她不好.
對那個男人不好.

$^{441}$對女人你可以這樣做.
這件事情一直發展開去.
很特別.
結果事情更加驚人的是什麼呢.
不是暴徒就這樣衝進來.
也不是屋內的人很竭力的抵抗.
而是有個人做了一件更驚人的事.
這裡我們看得不清晰.
哪人.
哪人是指誰.
是不是指老人家.
不是.
當然不是老人家.
他說完那番話.
接著有個人.
接著那些人很明顯不肯.
他們來的興趣就是一個利未人.
他們想和那個利未人的男人.
有性上的交合.
他們不會聽老人家的說話.
在這樣的形勢之下.
這個利未人.
未曾就抓住他自己的女人.
他的妾.
被那些暴徒抓了出去.
結果這一晚.
是很淒涼.
很可怕的一個晚上.
是整晚.
這一班人就這樣欺凌.
這個女性.
整段經文.
你會看到裡面有些位置很震撼.
那種安靜.
但整件事應該不安靜.
很多的暴力.
欺凌.
奸淫.
但是那個女人有沒有出聲.
沒有.

$^{481}$那個女人.
她卻是一個很大的安靜沉默.
但她卻是受一個很大的痛苦.
是整晚.
是沒有人救她.
沒有任何一個人救她.
整晚.
一直到天亮的時候.
那些人才放她走.
我深信這個位置.
是令人.
簡直是感到憤怒.
不能夠明白發生什麼事.
是利未人.
推她出去.
在家裡的男人.
都不止一個.
有沒有任何的抵抗呢?.
有沒有任何的拯救行動呢?.
沒有.
這個女人.
這個女人就是這樣在外面.
徹底地遭受到凌辱欺負.
沒有出聲.
幾乎沒有說她有聲.
直到早上的時候.
她回來.
她倒在.
留她住宿的主人的家門口.
直到天亮.
就是這樣.
整件事.
是不是一個很.
很令人難以明白.
究竟發生什麼事.
直到早上.
事件到最後的時候.
是什麼呢?.
她的主人起來.
開了門.

$^{521}$她的主人是利未人.
你會發覺到這裡.
都是很特別的.
她是如常地開門.
她昨晚應該都睡得很好.
她如常地睡覺.
她整晚都沒有想過.
出去找她的女人.
她的妾侍.
在早上開門.
她出去要上路.
如果不是這個女人.
倒在門前.
她應該不會找她.
不會理會她.
自己就上路.
經文是這樣說的.
她打算自己走.
那個妾侍很明顯.
晚上是經歷一些很痛苦.
很恐怖的事情.
但她仍然睡得著.
好像視她為不存在.
早上遇上你離開.
就是那一刻她見到你.
她的手搭在門欄上.
她怎麼說?.
她不是說.
你發生什麼情況.
我要找醫生.
我要幫你去看.
我接你去接受治療.
完全沒有問過.
她說怎麼說.
起來.
我們走.
那種冷漠.
很可怕.
真的很可怕.
我受了這麼大的煎熬痛苦.

$^{561}$而對方竟然這樣.
那個人.
婦人沒有回應.
其實因為她已經死了.
被人凌辱致死.
死了.
那個人就將她.
拖在路上.
帶回自己的地方.
整件事.
真的發展得很奇怪.
從哪裡呢?.
你再想想.
她在哪裡?.
她在基比亞.
一直上去.
上去以法連篡區.
好了.
到之後.
到這樣.
你就想.
這個女人生前.
遭受這麼大的痛苦.
是不是要好好地為她處理後事呢?.
是不是?.
我們無論什麼國家的人.
其實這件事都很體重.
讓那個人.
好像有一個很好的完結.
但是這個人.
你看她是怎麼處理後事的.
拿起一把刀.
阻住她的身體.
切成十二塊.
我在這裡.
要說.
你記住利未人.
她是事奉人員.
在聖殿裡事奉人員.
很重要的.

$^{601}$其實她應該要守一件事.
不要碰屍體.
碰屍體會令自己成為不潔.
是不是?.
她又會碰屍體.
不單是這樣.
她直接將屍體切成十二塊.
做什麼呢?.
將她分送到以色列全境.
要去傳遞一件事.
原來到了這一刻.
她的切身體成為了她說話的工具.
我將這個訊息.
發到以色列的國之牌.
告訴別人.
基比亞人做了多壞的事.
看到多壞的事.
他們怎樣對待這個女人.
你看.
看到屍體.
就看到他們的暴行.
是她說話的工具.
你說她是不是很愛這個女人呢?.
我完全看不到.
有任何一丁點的愛.
因為愛會帶來保護.
會有捨己犧牲.
她一點都沒有.
她反過來.
完全是令人為她捨己犧牲.
自私到極點.
到了這一刻.
又做這些事情.
你說她是事奉人員.
看哪一部分都不像.
整段經文沒有任何說.
她怎樣和上帝關係親密.
一丁點都沒有.
一個這樣的人.
接著.

$^{641}$大家可以回去看20到21章.
當利未人發出訊息.
之後.
國之派非常同心.
如同一人聚集.
他們說大家一起去聲討基比亞.
基比亞屬於變亞敏這個支派.
很特別.
變亞敏這個支派的人.
不是說要處理我們裡面.
基比亞人的惡行.
不是.
而是他們竟然要保護基比亞人.
所以變亞敏人就開始和其他支派的人.
進入內戰.
一直說打仗.
由個人.
本來是一個家庭的事情開始.
利未人與她的妾之間的關係.
整整下整整下.
變成支派的內戰.
死了很多人.
整件事真是.
我們不能夠明白.
為什麼會這樣呢?.
當以色列中武王.
我們來一個小小的總結.
百姓深陷罪惡之中.
基比亞人.
喜歡怎樣就怎樣.
在性方面混亂.
喜歡同性戀.
會強暴.
婚姻都不足夠.
根本不重要.
我現在喜歡什麼人.
就透過強暴去佔有.
去擁有.
這是基比亞人.
而且是大火一起來.

$^{681}$不是一個半個.
當以色列中武王的時候.
坐人者都陷於罪惡當中.
那個老人家好像是.
裡面的人物裡.
稍微一個正面一點點的人物.
稍微而已.
但你記不記得.
在事情發生的時候.
他是怎樣做的?.
是不是保護了所有的客人呢?.
不是.
你只是保護那個男人.
那些女人.
你任意待他們.
女人就不是人嗎?.
是不是?.
女人一樣是人.
一樣是這麼寶貴.
要保護.
所以這個坐人者.
你想一下他的觀念.
什麼叫幫助人?.
在過程當利未人.
把他拉出去的時候.
他也不會走出去.
勸阻這件事.
或者救那個女人回來.
不會.
因為他原先定義.
你可以傷害我的女人.
你怎樣對他們都不要緊.
所以他按自己的觀念去想.
他也一樣在罪惡裡.
所以幫助人也不是好人.
當以色列中武王的時候.
神的事奉人員.
也是深陷罪惡之中.
第一對婚姻.
在他什麼態度.

$^{721}$在自己太太之外.
還要再取妾.
你說當時的社會就是這樣.
每個人都是.
你看回時詩.
我去讀時詩記的時候.
我看回時詩.
有些真的很多太太.
所以就很多兒子.
娶這麼多女人.
但是你看回創世記的時候.
上帝本意.
神去看回.
神設立婚姻的時候.
婚姻的時候.
是一夫一妻.
一生一世.
是一夫一妻.
一生一世.
在婚姻裡是有忠誠的.
但是你看回這裡.
這個事奉人員.
對上帝的認識是怎樣呢?.
你已經打過了.
接著.
到了一個危險的情況下.
會犧牲別人.
他那種冷漠.
簡直令你鄙視他.
他沒有想過拯救他.
沒有想過保護他.
整晚還可以好好地睡覺.
第二天好像沒事發生.
我要如是者上路.
如果他不回來.
我當見不到.
算了.
這是什麼生命?.
完全沒有生命可言.
和上帝沒有任何關係.

$^{761}$沒有王就是這樣.
所以有時候.
事奉人員也可以.
只有一個名字.
沒有生命的實質.
當以色列中沒有王的時候.
衝突會升級.
從家庭問題.
還可以轉變為.
支派之間的衝突.
內戰.
可以像滾雪球一樣.
沒有王就變成這樣.
在這裡.
讓大家看看.
以色列人的重大死亡事件.
我們看看.
中間黃色的字體.
在內戰中.
變額敏人.
殺敗了十一支派.
四萬個士兵.
你可以看回《士師記》20章.
然後.
最後的時候.
十一支派最後贏了.
他們差不多剿滅.
整個變額敏的支派.
又殺了他們.
二萬五千一百個士兵.
本來是利未人和他們.
之間的事情.
最後發展死了這麼多人.
你說這個時間.
是何等的黑暗.
就是因為以色列沒有王.
他們個人生命沒有王.
整個民族沒有王.
個人任意而行.
你說關我們什麼事.

$^{801}$今天不是王的時期.
今天我和你的生命當中.
有王嗎?.
我和你的生命有沒有王?.
有.
很肯定.
你的王是.
我們的王是誰?.
耶穌基督.
是嗎?.
真的嗎?.
在哪裡看到.
我們的生命的王.
就是耶穌基督呢?.
我想問.
我們是不是真的將生命的.
每一部份都交在.
耶穌的手中呢?.
真的嗎?.
開始不敢回答我.
(笑).
要說不容易.
這些口號.
我們在教會內就會懂得叫.
什麼都交給耶穌.
真的嗎?.
真的每件事都問上帝心意嗎?.
我個人的生命.
是不是真的以主為王?.
我自己人生大大小小的決定.
家人大大小小的決定.
整個家庭.
大大小小的決定.
是不是真的在耶穌基督內?.
還是我已經決定了.
我叫他知會一下.
有時還不知會.
何必我知會.
他就知道.
我已經做了決定.

$^{841}$真的嗎?.
真的他是王嗎?.
教會一樣.
我想問.
最刺激的地方.
今天早上.
我很感恩.
有一個全教會性的祈禱會.
雖然不是每個弟兄姊妹.
都在當中參與.
可能有很多原因.
事奉.
時間.
但今早是很感動的.
今早很感動.
有一部份.
我們的禱告裡面.
我指的是.
要將我們洪恩家的主權.
要回歸給上帝.
上帝才是我們的王.
我們可以.
開很多會.
做很多決定.
可以按照經驗.
我覺得一定是這樣.
以往經驗就是這樣.
現在也是這樣.
就可以了.
可以做很多事.
我們可以開會.
做很多決定.
很多事.
但有沒有好好地.
好好地問神?.
真是負心自問.
有多少是出於神?.
多少?.
有多少是自己心意?.
我們可以騎劫教會.

$^{881}$不知不覺間.
騎劫教會.
屬於自己的.
是我們的.
你記住.
是耶穌基督的.
真的要問.
多少屬於神?.
如果一件事.
是屬於神.
是神所喜悅的.
上帝一定會大大的.
興旺這裡.
不然的話.
做幾十百件事.
都沒有意思.
我今早感動的地方.
就是我們可以.
一起為這件事禱告.
一起求神.
去憐憫我們.
我見到有很多.
弟兄姊妹.
為教會禱告.
受旺.
我很感恩.
我深信.
我們很有盼望.
真的將主權.
交出來.
洪恩階.
是屬於上帝的.
我們有王的.
是耶穌基督.
真正的王.
不是只有口說的.
我們生命.
也有王的.
你只要回去.
問一問自己.

$^{921}$是不是真的.
在人生大小決定.
都是由上帝決定.
你就知道.
是不是真的.
以主為王.
我舉我自己的例子.
我也有很多.
很多東西要學習.
怎樣真的.
以上帝為自己的主宰.
我記得去年.
和前年.
就是我爸爸媽媽.
分別安息主懷.
我記得他們.
快要安息主懷之前.
時間也很緊急.
剛好正值.
我快要講道.
其實當時心裡.
真心地想.
不如我講完道.
爸爸媽媽才安息主懷.
我真的有這樣想.
簡直是這樣渴望.
最好這樣.
我就可以很安心.
送他們最後一程.
但是很特別.
這是我想的.
我想的.
我的時間表.
上帝不是這樣想.
兩次都不是.
兩次都是.
在之前.
爸爸媽媽已經安息主懷.
爸爸那次.
不知道是之後的一兩天.

$^{961}$之後我就講道.
媽媽那次更刺激.
就是當日.
當日凌晨.
不知道幾點.
五點左右.
她安息主懷.
之後大概十一點多.
我就講道.
兩次都不是我自己心意.
很多時候我想.
最好之後某個時間.
我想定那個時間.
做某些事.
神說不是.
不是那個時間.
是有分別的.
我想講.
所以我都在學這個功課.
我在這裡彼此挑戰大家.
我們不只是說說.
而是真的邀請耶穌基督作王.
問祂.
等祂掌管.
我很有信心.
洪恩嘉也會復興.
當我們真正高舉祂為王的時候.
我們一同禱告.
天父上帝.
藉著這段經文.
是一段.
令我們感到很震撼.
很憤怒.
甚至感到很恐怖的經文.
你提醒我們.
提醒我們.
當我們生命.
沒有真正以你為王.
當教會或國家.
沒有以你為王的時候.

$^{1001}$我們就會陷入很深的罪惡.
主啊.
今天你知道我們的生命.
到了這一刻.
如果我們的生命還是那麼自我.
那麼喜歡罪惡的時候.
我求主寬恕我們.
求主憐憫我們.
讓我們能夠離開罪惡.
真正把我們生命的主權.
交出來.
主願意在我們每一個人的生命裡.
掌王權.
也在洪恩嘉掌王權.
你帶領我們.
讓我們成為你所喜悅的群體.
我深信上帝你的大能.
一定可以透過我們個人的生命.
透過教會.
去彰顯出來.
我們等候你.
奉耶穌基督的名字祈禱.
其他事項.
\newpage



\section{路加福音 10:25-37}
\label{sec:QG8KdqrlJlA}
\textbf{2023.02.12 《誰是我的鄰舍?我是誰的鄰舍?》 路加福音 10\_25-37 (和修版) 彭秀芹導師}
\newline
\newline
連結: \href{https://youtube.com/watch?v=QG8KdqrlJlA}{\texttt{https://youtube.com/watch?v=QG8KdqrlJlA}} ~~~~ 語音日期: 2023-02-13
\newline
\newline
\hyperref[sec:CnIFsRI_OWk]{\small{< < < PREV SERMON < < <}}
~
\hyperref[sec:index_chronic]{\small{[返順時目]}}
~
\hyperref[sec:index_scriptual]{\small{[返順卷目]}}
~
\hyperref[sec:_JHmy8bvPEI]{\small{> > > NEXT SERMON > > >}}
\newline
\newline
路加福音 10:25-37
\newline
\begin{longtable}{cl}
\hline
\hline
章節 & 經文 (和合本修訂版)\\
\hline
10:25 & \begin{tabularx}{0.7\textwidth}{X} 有一個律法師起來試探耶穌，說：「老師！我該做甚麼才可以承受永生？」 \end{tabularx} \\ \\ \relax
10:26 & \begin{tabularx}{0.7\textwidth}{X} 耶穌對他說：「律法上寫的是甚麼？你是怎樣念的呢？」 \end{tabularx} \\ \\ \relax
10:27 & \begin{tabularx}{0.7\textwidth}{X} 他回答說：「你要盡心、盡性、盡力、盡意愛主—你的神，又要愛鄰如己。」 \end{tabularx} \\ \\ \relax
10:28 & \begin{tabularx}{0.7\textwidth}{X} 耶穌對他說：「你回答得正確，你這樣做就會得永生。」 \end{tabularx} \\ \\ \relax
10:29 & \begin{tabularx}{0.7\textwidth}{X} 那人要證明自己有理，就對耶穌說：「誰是我的鄰舍呢？」 \end{tabularx} \\ \\ \relax
10:30 & \begin{tabularx}{0.7\textwidth}{X} 耶穌回答：「有一個人從耶路撒冷下耶利哥去，落在強盜手中。他們剝去他的衣裳，把他打個半死，丟下他走了。 \end{tabularx} \\ \\ \relax
10:31 & \begin{tabularx}{0.7\textwidth}{X} 偶然有一個祭司從那條路下來，看見他就從另一邊過去了。 \end{tabularx} \\ \\ \relax
10:32 & \begin{tabularx}{0.7\textwidth}{X} 又有一個利未人來到那裡，看見他，也照樣從另一邊過去了。 \end{tabularx} \\ \\ \relax
10:33 & \begin{tabularx}{0.7\textwidth}{X} 可是，有一個撒瑪利亞人路過那裡，看見他就動了慈心， \end{tabularx} \\ \\ \relax
10:34 & \begin{tabularx}{0.7\textwidth}{X} 上前用油和酒倒在他的傷處，包裹好了，扶他騎上自己的牲口，帶他到旅店裡去，照應他。 \end{tabularx} \\ \\ \relax
10:35 & \begin{tabularx}{0.7\textwidth}{X} 第二天，他拿出兩個銀幣來，交給店主，說：『請你照應他，額外的費用，我回來時會還你。』 \end{tabularx} \\ \\ \relax
10:36 & \begin{tabularx}{0.7\textwidth}{X} 你想，這三個人哪一個是落在強盜手中那人的鄰舍呢？」 \end{tabularx} \\ \\ \relax
10:37 & \begin{tabularx}{0.7\textwidth}{X} 他說：「是憐憫他的。」耶穌對他說：「你去，照樣做吧！」 \end{tabularx} \\ \\
[1ex]
\hline
\hline
\end{longtable}
$^{1}$各位頂上平安.
因為2107的防疫措施.
可以在教會台上不需要戴口罩.
所以今天很開心可以見到大家.
大家見到我的真樣子.
今天跟大家分享的訊息.
相信大家都很熟悉.
有兩個問題.
我想回到第一個重點.
今天想跟大家探討兩個問題.
誰是我的鄰舍.
我是誰人的鄰舍.
俗語有云.
近親不如.
對不起是遠親.
太緊張了.
遠親不如近鄰.
表達鄰舍對我來說很重要.
在聖經裡有說.
做一個好鄰舍很重要.
重要到是神的誡命.
在進入經文前.
想跟大家收看半集.
短短的短片.
大家在看的時候要想想.
在片中的主角.
他是否一個好的鄰舍.
開了.
我已經開了.
現在我們來看看.
他應該是在說泰語.
這次我們今天分享.
這個比喻.
其中一些角色人物.
想和大家.
一起.
看看.
這個比喻的開頭.
究竟這個故事是怎樣來的.
好撒瑪利亞人.

$^{41}$這個故事我想大家都.
很熟悉一點都不陌生.
這個故事的前文.
出現了一個人.
這個人就是律法師.
他已經說到.
這個律法師不單是.
問耶穌問題.
他用了一個詞形容這個人.
問的問題是試探.
他試探耶穌.
這個試探和.
耶穌受撒旦的試探.
詞是一樣的.
即是說這個律法師.
心懷不軌.
他看看有沒有捉到.
耶穌的把柄.
然後抽請他去誣蔑他.
他就問了耶穌的問題.
老師.
我們要做些什麼才能承受.
永生呢?.
這個問題很正路.
很積極.
耶穌也是高手.
耶穌不是直接回答.
耶穌以問題.
回應他的問題.
耶穌問了.
你律法書裡面.
寫了些什麼?.
你讀的又是什麼?.
因為律法師.
其實是當時.
教導人律法的老師.
所以他很熟悉律法書.
所以他就很厲害.
引用了.
申命記錄章五節.

$^{81}$要盡心盡性.
盡義盡利愛耶和華利亞神.
連同利未記.
十九章十八節.
你要愛倫如己.
這兩節的律法書.
引用拼在一起.
成為大誡明.
原來在申命記錄章五節.
是猶太人從小到大.
都要背誦的聖經.
所以他們每個人都會背.
但會背.
到底會不會做呢?.
這就是問題.
這兩段經文是大誡明.
顯示律法師很熟悉聖經.
我又想問問大家.
為什麼這麼有趣?.
耶穌覺得.
安穩.
耶穌覺得.
他答得很好.
於是耶穌就讚賞他的答案.
就說.
你答得對.
你要盡心盡性盡義盡利愛主耶和華利亞神.
愛倫如己.
你就這樣行吧.
叫律法師你就這樣行吧.
我想問問.
我們今天得永生.
是不是靠行為呢?.
為什麼耶穌叫他.
你行吧.
如果我們信的福音.
只是死後的保障.
死後的福音.
那麼我們今天的生命會怎樣呢?.
原來今天.

$^{121}$我們的福音不是靠行為.
不過.
我們的信是不可以沒有行為的.
那麼.
耶穌叫他你就這樣行吧.
原來今天我們.
決志信主.
信主.
你就得救了上天堂.
是不是這樣呢?.
我們好像留下了.
中間的重要的東西.
就是信心是有行動的.
說一個.
比喻.
如果突然間.
輝哥出來叫.
三樓火燭啊.
你們會怎樣?.
走啊!你的走代表什麼?.
你信輝哥說的話.
如果曾姑娘說.
火燭啊?你先坐一會吧.
那麼你猜曾姑娘的.
反應是什麼呢?.
信不信?.
不信啊.
今天耶穌跟他說.
你說得對,你去做吧.
即是告訴律法師.
他的信.
堅定留下嗎?.
那麼待會再探討吧.
我們的信.
是會帶出行動的.
律法師的答案.
是對的,你要盡心盡性.
盡力盡意愛主理神.
100分.
那麼.

$^{161}$我去看原文聖經.
和查我們的中文字典.
查到「盡」字.
為什麼呢?原來原文聖經.
是解「All」.
即是你有的東西.
全部拿出來愛神.
即是等於「All in」.
大家知不知道「All in」?.
所有東西都拿出來.
今天神要求我們對神愛.
就是「All in」.
那麼中文聖經.
中文字典是這麼說的.
「盡」字.
就是.
全部用出.
竭力的做到.
盡心盡力盡水.
盡職盡忠.
盡責人盡其才.
物盡其用.
很多「盡」啊.
我們的生命.
有沒有這個特質.
這樣去愛神.
「All in」去愛神.
將我們的生命「All in」在神手裡.
那麼.
這個大誡明.
有.
好像有兩個要求.
愛神和.
愛鄰舍.
我想問問你.
這個大誡明是兩個命令.
還是一個命令?.
誰覺得是兩個.
舉手?沒有人舉手.
即是每個人都是一個,你們很聰明.

$^{201}$那麼看回聖經.
中文聖經是這樣行.
那麼英文聖經就說.
「Do this」.
「This」是單數.
那麼原文聖經那個字都是單數.
就是說這個大誡明.
是捆綁式的.
你愛神,你就愛你的鄰舍.
是捆綁式的.
那麼原來神.
要我們的生命.
不只是舉手,信主上天堂.
是要我們的生命成為.
別人的祝福.
是愛我們身邊有需要的人.
那麼當然.
律法師要證明自己.
有道理.
耶穌已經解答了他的問題.
但是他覺得還不夠.
於是.
再問一個問題,那誰是我的鄰舍?.
剛才也說了.
律法師動機不良.
他很想再抽證一下.
究竟耶穌有沒有什麼把柄.
被我捉到.
證明他的不是.
律法師受過高等教育.
律法師很熟悉.
他不會不知道答案.
他一定知道.
他很想利用他的知識.
去陷害耶穌.
或者去羞辱耶穌.
原來今天.
如果我們有知識.
沒有愛神的心.
是一個很恐怖的.

$^{241}$攻擊人的武器.
也因為律法師.
這個動機.
所以就帶出.
愛倫社的主題.
就是殺馬利人的故事.
好了.
我們來看看.
以下的故事.
相信大家都很熟悉.
我短短地再.
聲演一下這個故事.
因為我很喜歡講故事.
最喜歡跟小朋友講故事.
我看到台下的小朋友.
都很留心聽我說故事.
有一天.
有一個人從.
耶路撒冷到.
耶利哥.
原來耶路撒冷位於高地.
耶利哥位於低地.
所以是下去地理位置.
這條路很險惡.
因為那裡有很多山洞.
位置會隱藏.
盜賊.
這個人很不幸運.
就是在這個時候遇到強盜.
強盜怎樣做呢?.
就把他的衣服脫掉.
沒有說他沒有財物.
可能這個人沒有財物.
原來當時的人不像現在那麼有錢.
原來衣服已經是他的財物.
脫掉他的財物.
他應該有財物.
打到他半死.
而丟下他就走了.
剛好有個祭司.

$^{281}$經過.
看到這個人.
我想這個人.
打到他半死.
心想:我又救了.
有個祭司.
祭司愛神,認識神.
一定會救我.
誰不知這個祭司在他對面.
靜悄悄地走了.
沒有施以援手.
然後不久.
就有一個尼美人.
尼美人也是.
在聖殿.
帶領人們.
去敬拜神.
這個半死人心想.
我有救了.
有位尼美人經過.
他應該會幫我.
但是一樣.
他也是從對面路.
走了.
那麼不久.
就有一個薩瑪利亞人.
經過,他就看到.
這個半死人.
他就動了憐憫心.
向前走.
處理傷口.
帶他去.
自己的生畜,女館.
照料他.
有些人說要入他數.
第三個薩瑪利亞人.
才肯出門.
幫他.
原來.
在民眾裡有三個漢奸.

$^{321}$祭司.
尼美人,漢奸.
他的反應是怎樣.
是冷漠的,視而不見.
而薩瑪利亞人的漢奸.
是動了慈心.
有憐憫的心腸.
我們探討一下.
為什麼.
祭司和尼美人.
他看到.
但是.
他為什麼會這麼冷漠.
為什麼他視而不見.
其實當你一個人.
很冷漠視而不見.
其實很可怕.
有一次.
當然我小時候.
在讀書,被熱水.
燙傷了手.
手指的位置.
小時候.
不知道死了,第二天.
去行山,我隨便.
用紗布包著傷口.
去行山,行山後.
覺得不痛,因為覺得很痛.
揭開紗布看.
我的手的位置是全都癢了.
很嚴重的發炎.
我馬上去急診室.
我還記得場面.
有一位目無表情的姑娘.
過來.
她沒有叫我過來,好像沒有跟我說話.
她拿起鉗子.
拿起棉花棒.
點了一些應該消毒的東西.
沒有說任何話.

$^{361}$很大力,我還記得.
一下子就將我.
受傷的位置.
連皮連血.
一次過掃走了.
我大叫,她仍然沒有反應.
她那種冷漠.
令我到現在.
還很記得.
原來.
我相信.
那位打得半死的人.
對於祭司和利未人的冷漠.
我想他們都會覺得.
為什麼會這樣.
其實可能有原因.
那麼.
第一個原因可能就是.
祭司和利未人.
擔心這個打得半死的人.
是假扮的.
因為當時都很流行.
假扮死,等人幫他.
人們一幫他,他就怎樣呢.
就搶他的東西.
所以他很害怕.
他們害怕,所以不幫他.
這可能.
其中一個原因.
第二就是.
他覺得有可能.
這個人受了傷.
還沒有死.
還沒有死,要醫治.
醫治要怎樣.
要錢,醫治要怎樣.
要時間.
我還有很繁忙的行程.
如果我伸出援手.
他會打.

$^{401}$擾了我接下來的計劃.
他有很多計算,計算著計算著.
沒有好處.
也搞不定,還是不幫.
第三個可能性就是.
可能受傷者.
有機會死.
因為當時.
在聖殿.
公職,他們很著重.
聖潔,不聖潔是不可以.
進去聖殿.
當職的.
他們很害怕.
如果我幫他,他死了.
那就是我接觸了屍體.
那就不聖潔,會影響我接下來.
進行一些宗教的儀式.
所以.
他想想,還是沒有幫.
其實當時.
摩西律法有規定.
接觸過屍體的人.
是屬於不潔淨的人.
所以他們.
有機會不幫他.
還有當時.
一些宗教領袖.
不單止跟摩西.
律法這個規定.
他自己再加多了.
更加苛刻的規定.
其實當時.
摩西律法是容許人.
容許人.
拯救半死的人.
和安置路邊的屍體.
但是因為.
那些宗教領袖.
顯示我自己是更聖潔.

$^{441}$所以定的要求更高.
而因為這些想法.
他拒絕.
施出援手.
這些都是我們.
推估的原因.
有可能是這樣,因為聖經沒有交代.
為什麼他不幫他,這些只是我們的估計.
好了.
原來.
今天我們幫助人也有很多計算.
很多害怕.
可能有很多條件.
底下我們才想幫不幫.
我去年.
去了一個.
機構實習.
那個機構叫「教會關注貧窮網絡」.
那個實習.
聽到一位牧者分享.
令我很深刻,這位牧者.
他服事教會在旺角.
是一間中產的教會.
在這區,旺角.
其實有很多.
基層,很多tong 房.
很多露宿者.
還有很多新移民.
其實他們有很多需要.
這位牧者很想幫助.
這些有需要的群體.
但教會裡有些.
信徒有些看法.
他覺得.
如果做基層.
如果做新移民.
其實也會帶來教會很多問題.
有什麼問題呢?.
有些為人父母覺得.
基層新移民不懂教育子女.

$^{481}$如果子女進來教會.
其實也會影響我的小朋友.
還有覺得.
他們的衛生文化.
跟香港人有點不同.
有點擔憂.
還有資源問題.
他們這麼有需要.
而他們的收入.
奉獻可能.
可能是沒有.
資源問題.
很多的量度下.
有一群中產.
覺得.
是不是真的值得做.
是不是要做.
但這位牧者就繼續.
下去接觸.
基層的市民.
還有新移民.
有一天.
崇拜完之後.
有一位很激動的姐妹.
拿著一張皺皺的紙巾.
去找那位牧者.
叫牧者看看.
有蝨子.
於是那位牧者.
就淡淡然說.
OK,有蝨子.
我們去找.
究竟這些跳蝨的源頭.
他就去找.
他發現其中一位男士.
身上.
有跳蝨.
於是.
牧者就.
建議那位姐妹.

$^{521}$不如我們進行家訪.
看看這位.
弟兄.
究竟情況是怎樣.
他去到弟兄家.
原來弟兄就住在.
樓梯底.
近垃圾站.
的一間.
單位.
那個位置很小.
他們兩個都不知道.
自己放在哪裡.
很難找到地方坐下.
因為太小.
那位牧者就揭開.
弟兄的床褥.
他也想尖叫.
當然要控制住自己.
因為他看到.
數以百萬計的跳蝨在床底.
他真的.
看到馬上就蓋起來.
之後那位牧者就.
馬上回到教會.
開緊急會議.
就滅蝨行動.
幫助教會滅蝨.
幫助弟兄家滅蝨.
幫助他賣掉那張床.
繼續.
教他清潔.
幫助他清潔.
這位原先.
很激動的姐妹.
她竟然成為.
探訪隊的核心成員.
她發現.
原來教會是可以.
幫助有需要的人.

$^{561}$令到人不只在物質上.
在生活上改善.
是那份心靈.
生命的改變.
那麼.
你們人.
就算看到半死人.
無動於衷.
其實在我們身邊.
也有很多半死的人.
其實.
回看.
政府的.
香港貧窮.
情況報告.
在2020年.
香港貧窮的情況報告.
原來有165.3萬.
的香港人.
在貧窮線.
以下.
即是說.
大約23.6\%.
即是差不多.
4個香港人中有1個貧窮.
其實有很多有需要的人.
在我們身邊.
清潔工,基層,.
tong 房戶,新移民.
他們不單止物質貧窮.
其實他們的心靈貧窮.
很多時貧窮的人不只是物質.
心靈也極度貧窮.
情感也極度貧窮.
因為他們的貧窮被困在.
自己的世界裡.
很需要有人將他們帶出來.
那麼.
當我預備這個.
政黨的時候.

$^{601}$我學會了一件事叫做堅尼系數.
堅尼系數是甚麼呢?.
堅尼系數就是.
貧富懸殊的數字.
香港的堅尼系數是0.539.
是甚麼來的?最高是10分.
香港是0.539.
那麼對於.
即是中間也不是很差.
不是,在先進國家來說是屬於很高的.
亦都是45年來.
最高的.
香港的貧富懸殊問題越來越厲害.
其實在我們身邊.
是沒有需要的人.
那麼剛才說到民眾.
在民眾裡面有3個漢劍.
第一個是特里斯,第二個漢劍是尼美人.
第三個漢劍是薩馬利亞人.
那麼薩馬利亞人.
原來就是這個故事的中心人物.
那麼為甚麼會用薩馬利亞人呢?.
我想大家都會知道.
薩馬利亞人其實是.
猶太人很憎恨他們.
見到他們會繞路走.
為甚麼呢?因為他們覺得.
這個薩馬利亞人是不潔的群體.
源於.
一個歷史.
當時北的以色列.
被亞述國.
侵吞了.
那麼他們有一個政策.
就是將很多外族人.
很多民族.
侵入當時的地方.
強迫.
以色列人.
與外族人通婚.

$^{641}$那麼純正的血統猶太人.
他們很喜歡說純正的.
覺得他們不純正.
是混血集中.
所以他們是鄙視他們.
那麼他們覺得.
薩馬利亞人是半教的.
是混血的集中.
他們猶太人覺得.
薩馬利亞人不會幫助猶太人.
不會同情猶太人.
不會施以援手.
但偏偏就是薩馬利亞人.
伸出援手.
那麼.
這裡呢.
薩馬利亞人的行動.
去說出.
甚麼是盡.
盡力愛神.
但這個題.
這個盡力愛神.
其實與愛倫社是同一個解釋.
你如何愛神.
你如何愛倫社.
薩馬利亞人的行動.
就是回應.
甚麼是盡.
盡力愛神愛人.
那麼.
剛才的經文已經讀過.
他的救助.
是很全秘的.
第一個救助.
是用油和酒倒在.
半生人的傷口上.
當時中東急救方法.
也是.
受助者.
半生人的需要.

$^{681}$按他的需要.
給他.
第二個救助.
是扶他上自己的身體.
帶他去店內照應.
療傷.
陪伴他.
給他醫療.
第三.
第二天他真的要走.
可能是一個傷人.
拿了二銀錢.
大約是.
幾天的薪金.
不少錢.
給店主照應半死傷者.
他繼續想方法.
繼續幫助.
這個人.
而且還答應店主.
如果在治療期間.
有額外的費用.
我會回來.
跟進.
他會繼續付出.
你看到他的救助不是草草了事.
很多時候我們.
做善事.
未必按別人需要.
給別人.
你收過一些你不想要的東西嗎?.
但今天.
你看到撒瑪利亞人.
他給的東西.
正正是受苦.
受難的人真正的需要.
他不是草草了事.
他是很全面.
去照顧這個有難的人.
我想大家都聽過.

$^{721}$戴德生的康康.
戴德生.
大家都很熟悉.
戴德生是英國人.
他來了中國.
傳道在華.
大約51年.
1932年在英國出生.
1905年他不是死在英國.
他死在中國.
戴德生也是一個.
很好的全面的例子.
其實當時.
他來到中國.
其實中國的政治.
很混亂.
還有很落後.
而且.
原來他的原配.
他兩任太太和他第一任太太.
和他四個小朋友.
都是來自中國.
因為在中國很多疾病.
很混亂.
如果我是戴德生.
可能我不會再回到這個地方.
但是他沒有.
打倒他.
要去幫助中國人的心智.
他繼續留在中國.
原來.
戴德生.
他在英國受教育.
受過高等教育.
還有受過醫學訓練.
其實是一個大好前途的青年.
但是他放下了這些.
他全力為自己.
為了神去到中國.
去傳福音.

$^{761}$有一次.
他不是有錢人.
他每天.
都會為大約70多個貧窮人.
預備早餐.
有一次.
他把自己所有的錢.
買夠兩天的早餐.
但他真的沒有錢了.
兩天後怎麼辦呢.
兩天後貧窮人都會來吃早餐.
但他已經沒有了.
他已經買了.
很奇妙.
就是那天他用光了錢.
他收到一張.
124元的支票.
當時是很多錢.
可以繼續供應這些窮人.
我覺得神很奇妙.
當你肯樂意.
為神去擺上.
擺上一切的時候.
神是會有不斷的供應.
他有一個很出名的名言.
相信大家都聽過.
如果我有千萬的英金.
中國人.
都可以全部的支取.
如果我有一千條的性命.
都不會留下一條.
不給中國.
不是為了中國.
是為了誰.
為了基督.
就好像我們剛才看的大解明.
你要盡心盡性盡義盡利.
我們不是為了人.
是為了神.
當你有這個愛神的心.

$^{801}$你會有愛論社的心.
我們又想看看.
究竟今天.
我們要怎樣做.
別人的論社.
我們看看下一個powerpoint.
原來今天我們是.
別人的論社.
你記不記得律法師問了一個問題.
他反過來問.
誰是我的論社.
我第二個問題是什麼.
你是誰的論社.
很不同的.
如果按律法師的邏輯.
來講.
這個問題的答案.
誰是我的論社.
答案是受打半死的那個人.
受助者的論社.
但耶穌很有趣.
他反過來說.
耶穌就說.
這三個人.
誰是.
落在.
強盜手中的論社.
耶穌反過來問.
他不是告訴他.
誰是你的論社.
而是這三個人.
誰是被打到半死.
落在強盜有需要的人的論社.
原來耶穌調整了.
律法師這個問題.
原來我們今天.
要做有需要的人的論社.
原本律法師.
很自以為是.
耶穌就將.

$^{841}$這個問題.
調轉了他看事物的角度.
律法師.
很想試圖限定.
篩選他的論社.
誰是我的論社.
正如我們剛才.
估計利未人.
和濟斯.
為什麼不肯幫助打了半死的人.
有篩選.
有計算.
有害怕.
但耶穌要我們.
像撒瑪利亞人.
有一個好榜樣.
是怎樣.
全力幫助.
自己有的東西.
幫到盡.
所以我們今天不應該問.
誰是我的論社.
我們應該問.
我是否別人的論社.
問問自己.
你是否別人的論社.
這也是我們一生要學的.
當我們肯去.
做別人的論社的時候.
你會發現你會得生命.
你會怎樣?得生命.
有一份難以形容.
從神而來的喜樂滿足.
從我們心裡.
發出來.
這個經文.
提到.
耶穌問.
哪一個.
是這個落難人的論社.

$^{881}$律法師回答得很好.
他怎樣回答?.
憐憫.
就是他的論社.
憐憫這個字.
原來在原文聖經.
憐憫這個字.
字根是指人內在的器官.
你的心肝脾肺腎.
我不知道你有沒有試過愛一個人.
愛到哪心哪肺.
我有一個很深的經歷.
我記得.
我生我的女兒的時候.
我一向都是.
做學前教育.
我做過影院.
我知道一個嬰兒.
面黃入雙照燈.
是一件很普通的事情.
當自己的女兒.
面黃入雙照燈的時候.
是絕對不普通.
那種哪心哪肺.
作為父母你會明白.
那種憐憫.
我記得我出院.
我女兒不能出院.
因為她面黃.
那一晚我不停地想.
她會怎樣呢?.
很慘被人蒙住眼睛.
很慘沒有三聚在照燈.
其實這些事我全部都知道.
但就是那一副憐憫心像.
在你心裡打動你.
旋轉你的心.
原來今天.
有需要的人.
沒有在我心裡.

$^{921}$在你心裡.
打動你的心.
打動你的心肝脾肺神.
你裡面的愛.
那個憐憫之心.
原來今天我們怎樣去幫人呢?.
我們有沒有這份愛.
有沒有這份憐憫.
這裡說到撒馬利人.
就因為有這個憐憫心像.
於是他.
打破了種族的矛盾.
其實.
按猶太人來看.
沒有可能撒馬利人.
會幫這個人.
但因為他是憐憫心像.
他打破了種族的矛盾.
亦都放下了很多的成見.
和標籤.
原來這個.
真正去.
真正的憐憫.
其實就是要看見.
有需要的人.
而不理會.
這個有需要的人有甚麼立場.
有甚麼條件.
很多時候一說到立場.
我們都會想起一些顏色.
無論是甚麼顏色.
如果我們去愛一個人.
有神從而來的心腸.
我們要放下這些東西.
而要怎樣?.
要盡力地幫他.
原來這個比喻.
就是教我們.
盡力地.
去做一個好的鄰舍.

$^{961}$而做一個好的鄰舍.
那份憐憫的心腸.
是超越了我們的積涵.
超越了我們的計算.
超越了我們的種族.
超越了我們的立場.
剛才我看了上半集.
我再看下半集.
究竟好的鄰舍.
有沒有出現呢?.
我相信大家也看到.
這個是好鄰舍.
這位好鄰舍.
就是這位露宿者.
其實他一無所有.
但他怎樣可以成為.
別人的鄰舍呢.
他也是盡他一切.
去保護一個.
守護一個不喜歡他的老闆.
不接納他的人.
他也是盡他的力量去做.
原來今天.
我們作為好鄰舍.
我們沒有什麼優越條件.
也沒有什麼好強.
我們的心也沒有盡一切去愛神.
愛人.
原來最後.
我們要怎樣做好鄰舍呢.
Do this.
就是剛才說的.
我們的大誡命.
愛神 愛人.
神就會用我們的生命.
成為別人的祝福.
我們有時間一起祈禱.
主耶穌.
多謝你.
首先憐憫了我們.

$^{1001}$多謝你.
用了你最高的愛.
來愛我們.
神啊.
多謝你親自來.
尋找 拯救我們.
你給我們的.
恩典.
讓我們白白領受.
讓我們學習像你一樣.
看見別人的需要.
看見這個世界的需要.
你幫助我們.
能夠盡我們一切去愛神.
愛人.
幫助我們成為有需要人的好鄰舍.
我們這樣禱告.
奉主名求.
阿們.
\newpage



\section{提摩太後書 2:19-22}
\label{sec:_JHmy8bvPEI}
\textbf{2023.02.19 《成為聖潔的器皿》提摩太後書 2\_19-22 (和修版) 張美薇牧師}
\newline
\newline
連結: \href{https://youtube.com/watch?v=_JHmy8bvPEI}{\texttt{https://youtube.com/watch?v=\_JHmy8bvPEI}} ~~~~ 語音日期: 2023-02-21
\newline
\newline
\hyperref[sec:QG8KdqrlJlA]{\small{< < < PREV SERMON < < <}}
~
\hyperref[sec:index_chronic]{\small{[返順時目]}}
~
\hyperref[sec:index_scriptual]{\small{[返順卷目]}}
~
\hyperref[sec:zwEIaeiy0j4]{\small{> > > NEXT SERMON > > >}}
\newline
\newline
提摩太後書 2:19-22
\newline
\begin{longtable}{cl}
\hline
\hline
章節 & 經文 (和合本修訂版)\\
\hline
2:19 & \begin{tabularx}{0.7\textwidth}{X} 然而，神堅固的根基屹立不移；上面有這印記說：「主認得他自己的人」，又說：「凡稱呼主名的人總要離開不義。」 \end{tabularx} \\ \\ \relax
2:20 & \begin{tabularx}{0.7\textwidth}{X} 大戶人家不但有金器銀器，也有木器瓦器；有作為貴重之用的，有作為卑賤之用的。 \end{tabularx} \\ \\ \relax
2:21 & \begin{tabularx}{0.7\textwidth}{X} 人若自潔，脫離卑賤的事，必成為貴重的器皿，成為聖潔，合乎主用，預備行各樣的善事。 \end{tabularx} \\ \\ \relax
2:22 & \begin{tabularx}{0.7\textwidth}{X} 你要逃避年輕人的私慾，同那以純潔的心求告主的人追求公義、信實、仁愛、和平。 \end{tabularx} \\ \\
[1ex]
\hline
\hline
\end{longtable}
$^{1}$請姐妹平安.
在此還可以趕得及和大家拜過晚年.
因為今天仍然是正月最後的一天.
恭祝大家新年進步.
主恩門日.
讓我們開始的時候.
我們一起低頭祈禱.
天父上帝願你使用孩子今天所說的.
叫孩子所說的每一句都是你要孩子所說的.
更加幫助會眾.
他們有一個能夠聽道的耳朵.
有一個能夠明白你道理的頭腦.
更加有一個心將你的道理載在心裡.
然後用他們能行道的手和腳.
將你的道理行出來.
以致他們是一個能夠聽道行道的人.
以致他們能夠土利喜悅.
願主親自使用今天的正道時間.
聽我們的祈禱.
奉耶穌基督的名求.
阿們.
今天老牧師給我一個題目.
叫我講聖潔.
什麼是聖潔呢?.
或者當我問大家.
你是否聖徒呢?.
我們很多人都會覺得.
我們比較難說自己是聖徒.
如果我問你們是否信徒呢?.
大家都會說我們信了主.
所以我們是信徒.
或者門徒都說.
這樣也好.
我也願意跟從主.
但一提起聖徒或者聖人或者聖潔.
好像有些不是很自然的感覺.
覺得我們好像做來做去都達不到標.
其實宣告會是有四重福音的.
有沒有聽過?.
如果你們聽過洗禮班或者初信班.

$^{41}$都會聽過的.
第一重福音就是.
基督是我們救贖的主.
第二就是基督是我們成聖的主.
第三就是基督是我們醫治的主.
第四就是基督是我們再來的王.
有時候某些教會比較少說主耶穌基督的再來.
我不知道洪恩有沒有多說.
耶穌是我們救贖的主.
大家都很熟悉的.
但當我們說基督是我們成聖的主的時候.
我們有時候是比較難明白.
難接受的.
當到第三點.
基督是我們醫治的主.
有時候我們都很難的.
因為很多時候就怕走了靈因.
於是我們就不說的.
但是我想我們教會是要多說的.
四重福音其實四重都要說.
不知道大家有沒有見過這本書.
《學賞聖潔》.
因為早期是宣道會出的.
就是有《學賞救贖》.
《學賞聖潔》.
然後《學賞醫治》.
然後《學賞再來》.
以前在錦繡的時候.
我們曾經是整套送給剛洗禮的弟兄姊妹.
讓大家可以靈修的時候用的.
我都鼓勵大家去問問曾姑娘.
或者問問魯牧師.
可不可以再買到.
當我們今天去讀這段經文的時候.
當我們發覺信徒能夠成為聖潔.
其實是有兩個面的.
是有兩面的.
就是雙重的根基.
因為在第十九節裡面.
聖經這樣說的.

$^{81}$「然異神堅固的根基,屹立不移,上有者印記說:主認得他自己的人.」.
又說「凡稱呼主名的人,總要離開不移.」.
所以這裡說第一面.
就是耶穌或者神或者聖靈.
是認得屬祂的人.
所以如果我們是屬祂的話.
祂是認得的.
祂會記住的.
祂會幫助的.
但是我們這樣就不能將所有的責任.
我們成聖的責任全部交給耶穌或者神.
因為我們還有第二個面.
就是說「凡稱呼主名的人,總要離開不移.」.
離開不移.
所以聖潔.
我們今天會跟大家說的.
主力就是說.
我們有什麼要做.
耶穌做了祂的部分.
神做了祂的部分.
但是我們人也有部分是我們要做的.
聖徒能夠成為聖潔.
是有這兩個面的.
所以當我們又會說.
聖潔是什麼意思呢?.
大家有沒有聽過.
「分別出來成為聖潔」.
有沒有聽過.
我們要分別為聖.
就是說是分出來的.
那些東西是歸於神用的.
所以就要分別出來.
分別出來.
分別出來成為聖潔.
所以其實聖經就將倫理行為.
和屬靈的生命.
生命是連起來的.
無論耶穌基督的登山寶訓.
普羅的聖靈果子的講論.
乃至雅各和彼得的應用性的講論.

$^{121}$都將基督徒的行為表現.
視作屬靈生命的外顯.
但是我們要小心.
不要將那些人成為一個指標.
就是我們一定要這樣這樣這樣.
這樣很多時候就會過分了.
但是我們知道.
我們也不可以忽略.
我們的生命行為.
是表現屬靈生命的外顯.
如果沒有看到我們的果子.
就不能夠聲稱有裡面的屬靈的實質.
我們憑果子就可以認出我們的樹泥.
所以其實今天我們說.
信徒是要聖潔.
或者聖徒.
我們其實每個都是聖徒的.
當我們去看普羅寫書信.
很多時候他都用聖徒.
來稱呼收信的人.
因為我們屬神的話.
我們就是聖潔.
所以當我們會發覺這樣的時候.
我們是要與眾不同.
與眾不同.
但是今天在我們的社會當中.
在我們的家庭的生活.
或者我們的婚姻的生活上.
很多時候基督徒的表現.
其實都不及猶太人的表現.
原因之一就是我們教會裡面.
沒有很整全的去教導基督徒的倫理.
當我們傳福音的時候.
我們很多時候都很著重.
我們福音要這樣.
我們要怎樣怎樣.
我們要信主.
我們要這樣那樣.
而很少是講律法的.
為什麼呢?.

$^{161}$因為我們怕別人會說我們是律法主義.
我們定了這些的規條.
但是其實當我們為什麼要聖潔呢?.
保羅很喜歡說一句.
因為神是聖潔的.
所以我們要聖潔.
所以我們的聖潔不是我們要達到了.
才可以成為信徒.
而是因為我們去做了信徒之後.
愛之後.
因為神是聖潔.
所以我們就要聖潔.
我們要跟著去聖潔.
所以不是說我們聖潔了.
我們才可以得救.
大家要很清楚這裡.
不是說我們聖潔了才可以得救.
但是當我們得救之後.
我們就要去聖潔.
因為神是聖潔.
我們通常不是這樣理解.
我們基督徒是跟一般的市民不一樣.
好像這些不是我們應該的.
所以我們很多時候就說.
我們是不是重視聖潔生命呢?.
是不是重視聖潔的品質呢?.
我們對社會上正義的衛生.
或者職場上表現誠實和表裡一致.
我們很多時候都沒有說.
這是我們不理想的地方.
所以其實應該.
應當來說.
一個信主的家庭.
它是一個以消費為導向的貪婪的社會對比之下.
顯示他們的生活是很簡單的.
是快樂和滿足的生活形態.
或者因為我們有一個穩定的家庭.
而這個家庭對於不忠.
或者對於離婚.
是不認識的,不知道的.

$^{201}$而且孩子在父母的安全環境下長成.
但是我很記得我小時候.
我媽媽教我的不是跟別人不一樣.
我媽媽教我的就是要討好所有的人.
意思是我要面面俱圓.
但是當我相信主之後.
我才發覺如果我面面俱圓的話.
對於一些不認識神的標準.
道德等方面.
我是很難討他喜悅的.
我才發覺.
原來我媽媽對我的教導是錯的.
當然我媽媽當時還未認識神.
這是她在社會裡長大.
她是一個很好人緣的人.
所以她會覺得.
我都要跟他們一樣.
她每個子女都要這樣.
是要面面俱圓.
然後我發覺原來這個是錯的.
其實宣道會.
一個屬靈的作家或偉人.
陶書有沒有聽過他的作品?.
陶書也說.
與世人無異的人.
聖靈是不會充滿他們的.
沒有被聖靈充滿的人.
是不能夠成聖.
沒有成聖的人.
生活是不能夠喜樂.
所以我們要追求.
跟未信的人不一樣.
今天的經文教導我們.
有什麼我們可以追求.
與世上其他人不一樣.
本來就舉了一個例子.
他舉了一個大戶人家的器皿的例子.
我們去看.
他第一個的說法就是.
我們要離開不二.

$^{241}$大戶人家有很多不同價值的器皿.
在第20節裡這樣說.
大戶人家不但有金器銀器.
也有木器瓦器.
有作為貴重之用的.
有作為卑賤之用的.
大戶人家有各種的金銀木瓦.
有貴重有卑賤.
這些不同價值的器皿.
有不同的用處.
當然大戶人家有金器銀器.
那些是用來服務食物的.
但這些器皿也不一樣.
就算我是小戶人家.
不是大戶人家.
我去預備菜.
我通常用梯碟來預備.
但當我要蒸魚或炒菜.
我都會放在瓦碟上.
可能我覺得瓦碟好看一點.
同時保暖會好一點.
我通常會用不同的器皿.
所以用這個例子.
不是說不同用處的東西.
而是說不同價值的東西.
我最初聽到聖結.
時時有一樣東西我不明白.
大家如果….
好像很噁心.
如果踩到地上的東西.
是否立即污染了.
即使是鞋底或是尾巴都會污染.
所以聖結時說.
當我們沾染了污穢的東西.
我們就怎樣?.
就污穢了.
但相反來說.
當你沾染了聖結的東西.
你是否會聖結?.
是嗎?.

$^{281}$不是的.
人家是聖結.
你沾染了是不會聖結的.
所以當我們要聖結的時候.
我們有些東西.
不是沾染了人就行了.
所以這點很重要.
當我們沾染了不乾淨的東西.
我們就會變為不乾淨.
所以我們要明白.
我們不要沾染不乾淨的東西.
但當我們發覺.
當我們行為端正的人.
不是生命一定是純粹的.
因為外邦人都可以有好的行為.
但當我們知道行為錯誤的人.
一定在生命出了問題.
行為好的人.
不是生命一定是好的.
但如果行為有不好的東西.
他生命一定出了問題.
所以當我們知道.
偶然被過犯所幸.
是人人都會發生的.
我們今天信了主.
為什麼我們要向神認罪?.
因為我們不多不少都會犯錯.
但我們偶然為過犯所幸應該會發生.
但我們卻不能設想一個.
享受罪終之樂的人.
我們不可以在罪裡面.
繼續繼續做下去.
我們仍然會享受.
所以當保羅提到.
大戶人家有不同價值的器命的時候.
他就會提到.
人要自潔.
脫離悲戰.
他在聖經裡說.
「人若自潔,脫離悲戰的事,必成為貴重器命」.

$^{321}$他不是完全用日常的用法.
如果是雅氣.
脫離極.
也不會變成金氣.
但我們卻知道.
在主裡面.
用這個意思是可以這樣做的.
所以聖潔就是指脫離罪惡.
專心過著像神一樣的生活.
這是伯克雷說的.
而楊少棠.
聖經學者也說.
「何時有污點,何時就少能力」.
「何時不真實,何時神就擔憂」.
「這種在基督徒的經驗中如影隨形,毫離不差」.
「越離主近,或越像主的形象中」.
「聖潔和真實是最重要的部分」.
我們怎樣去潔淨自己呢?.
這不是我們能夠做的.
我們要靠著耶穌的血才能做.
所以剛才說到是兩面的.
第一面就是神已經做了.
我們靠著耶穌的血.
已經將我們的過犯清除.
我們願意被耶穌的血洗淨.
就一定運用到我們的意志上.
願意除盡一切罪惡.
好像以色列人在除孝節的時候.
要將一切的教都除盡.
不過我聽到有些人開玩笑地說.
大家都知道.
逾越節就是除教節.
他說今時的猶太人.
有時不捨得將所有的發酵的東西都扔掉.
他們會怎樣呢?.
如果不是猶太人.
他們會將旁邊的放了七天後.
再拿回來.
明不明白?.
我們不捨得扔掉.

$^{361}$但這其實是一個極大的錯誤.
我們要自覺脫離卑賤的事.
卑賤的事比普通我們觀念當中的罪.
更加細微.
不單是一些明顯的罪惡.
更加是一些和信徒身份或體統不相稱的事.
有罪惡的成分.
不體面或不光明的事.
都是卑賤的事.
有時我們會用一些小的詭計.
暗中傷害我們不喜歡的人.
又或者佔一些人的便宜.
有些暗味的事.
不願意被人知道.
或者不敢被人知道.
我們所有的事.
可能都是在暗中進行.
但我最近聽到一個學者告訴我.
他說現在的基督徒要小心.
你可能有時以為你所做的事.
只是你和神知道.
但原來遲早外面的人都會知道.
所以是否能夠見光.
這是一個很好的標準.
我們無論怎樣做這些事.
就算是完全在沒有人知道.
或者在黑暗當中.
只有神和我們知道的話.
我們如果做錯了事.
我們都不要做.
這樣才能過得了.
這些普通人或小人所做的事.
都是卑賤的事.
所以保羅說.
人若自覺脫離卑賤的事.
就必成為貴重的器皿.
在普通人的家中.
一件金的器皿.
可能它已經是金的.
但它是要做金器的用途.

$^{401}$如果它是啞的.
它就應該做啞的用途.
但在神的家中卻相反.
如果將自己做了金器的用處.
就成為金器.
如果你自己願意做啞器的用處.
你就可以成為啞器.
所以你不脫離卑賤的事.
就是卑賤的器皿.
所以脫離卑賤的事.
而做貴重的事.
你就可以成為貴重的器皿.
保羅在此也引用一個例子.
來教導提摩泰.
因為他是提後.
告訴他.
我們要克制的.
是要逃避少年人的私慾.
在第22節上說.
你要逃避年輕人的私慾.
當時在希臘羅馬的社會.
雖然當時是羅馬政權統治.
但仍然沿用希臘的文化.
性慾和私慾是很放蕩的.
傳說當中.
希臘的女神.
叫她的女僕.
在夜裡徘徊羅馬街道上賣淫.
所以他們在敬拜中.
他們在膜拜中.
公然說出一些很卑劣的不道德的事.
有性經學者提到.
當時社會的情形是怎樣.
一個男人可以擁有情婦.
同時這個情婦.
亦是他的知識伴侶.
亦是很會思考的.
而社會的奴隸制度.
亦使當時的男人.
很輕易地擁有妻妾.

$^{441}$除了太太之外.
他們還有妾侍.
而當時的妓女.
滿足了他們萬不驚心的情慾.
妻子用來做什麼呢?.
妻子的功能就是管理家務.
是孩子的法定母親.
所以一個男人可以有妻子.
有情婦.
然後有奴隸做他的妻妾.
還可以去找外面的妓女.
我們今天可能會說.
今天香港或世界.
好像不是那麼敗壞.
是不是呢?.
我早前剛聽到紐西蘭的總理.
女總理辭職.
大家知不知道為什麼?.
因為她要回家.
照顧她初生不久的小朋友.
然後更加說.
因為在過去的COVID(新冠肺炎)當中.
她還給了原因.
所以她還未結婚.
她想著疫情過後.
她就會結婚.
但她已經是一個.
可以說是一個未婚的媽媽.
和她的男朋友同居.
但已經是總理.
所以這個社會上已經是很接納.
我又看到.
在一月尾的時候.
四川省發佈了一個.
放寬生育登記的措施.
我們知道中國一定是.
當我們發覺要有子女.
有什麼東西是要登記的.
當他們會發覺.
夫妻應當在生育前進行生育登記.

$^{481}$那一句.
改了做什麼呢?.
凡生育子女的公民.
均應辦理生育登記.
凡生育子女的公民.
均應辦理生育登記.
你有沒有聽到有什麼問題?.
即是說現在已經不再需要.
已婚的婦女.
即是已婚的丈夫.
爸爸媽媽已經結婚了.
所以才去登記.
我們大家都知道.
其實中國最近一年.
生育率是下跌的.
所以中國最近一年.
人口是退的.
是很多很多年以來都沒有的.
我想可能以前是大戰之類.
死了很多人.
所以人口是下跌的.
但這麼多年都是增加.
但最近開始減少.
所以有人在網上說.
為了生育率.
真是無所不用其極.
有人問.
這不是鼓勵未婚先人?.
不是鼓勵一些小三生蛙?.
他們說病急亂投醫.
為了解決一個問題.
而製造新的問題.
單親媽媽越來越多.
不結婚生子也越來越多.
愛譽是否更加肆無忌憚?.
我們看到.
好像很保守的中國.
也有這樣的情況.
當我們回看香港的電視劇.
我想已經很久了.

$^{521}$我小時候看的電視劇.
幾乎每一套都有人.
是很漂亮的.
可以做到愛到不能不愛的婚外情.
可能會說和原配的關係不好.
這樣那樣.
婚外情是做到很漂亮的.
但其實我們的道德標準是越來越過.
我記得我小時候.
如果是同性戀是要坐牢的.
如果被人發現.
現在是爭取他們的權益.
所以道德標準真是不斷的改變.
有時我們去看我們的子女.
男朋友和女朋友.
可不可以一起去旅行?.
有時他們一起去旅行的時候.
他們覺得兩人住.
如果參加旅行團.
兩人住一間房是否便宜一點?.
他們就會兩個人.
男和女.
就會一起去租一間房.
信徒怎樣去面對這個世界?.
保羅在這裡說.
我們要逃避年輕人的私慾.
逃避.
我們怎樣可以做貴重的器皿?.
逃避年輕人的私慾.
所謂年輕人的私慾.
或者少年人的私慾.
是特別和少年人有關的私慾.
少年人在男女關係方面.
是很容易受私慾的引誘.
慾望很高.
而節制的能力也非常小.
所以保羅提醒他們.
他不是提醒他們要忍受.
而是要逃避.
抵擋是當受敵來的時候.

$^{561}$加以反抗.
所以是被動的.
但逃避是主動的.
所以這個主動的逃避.
即是我們不去找一些試探.
來抵擋.
我們不要去接近試探.
不要去顧慮.
我們不要去接近試探.
不要去故意的去遇到試探.
然後才去抵擋試探.
逃避試探是我們要採取主動.
不要讓試探來臨.
不過在不能夠逃避.
而受到魔鬼試探的時候.
我們就要加以抵擋.
所以這裡所說的少年人的私慾.
不單單只是男女方面的情慾.
也包括驕傲.
對財富和權利的慾望.
過度的野心.
妒忌.
羨慕.
好爭辯.
同時是包括吃喝玩樂.
放縱肉體的事.
有沒有聽過.
人家說COVID三年以後.
大家就像個球一樣.
完蛋了.
看來好像是一個形容.
但其實我們沒有好好管理身體.
同時也有些奇怪的理想.
不合聖經的各種幻想和慾望.
例如好高冒遠.
喜歡出風頭.
不受勸戒.
憑血氣行事.
這些都可以說是少年人的私慾.
所以保羅說.

$^{601}$要離開不義.
要自潔.
離開不義自潔是負面方面.
對一些不乾淨的事.
我們不要踩進去.
我們不要碰到它.
我們要離開.
我們要自潔.
另一方面是怎樣呢?.
另一方面就是成為聖潔.
我們要土神的喜悅.
聖潔的希臘文.
意思就是分別出來.
基督徒在兩方面是聖潔的.
是被分別出來的.
所以消潔就是從罪裡被分別出來.
積極的就是被分別出來後.
歸給神.
我們看舊約.
會幕和聖殿中的器皿.
都是分別出來.
是讓神用.
是專門歸給神.
專門服侍神的時候用.
教會中貴重的器皿也是這樣.
其實大家有沒有想過.
我們每一個都是貴重的器皿.
而我們身為基督徒的主要目的.
一切的工作都是為了服侍神.
因此我們必須保持自己的純潔.
一個器皿不能同時用來裝垃圾.
同時用來招待客人.
擺東西上桌吃.
裝垃圾就是裝垃圾的器皿.
不會用來擺上桌.
所以我們一定要知道.
貴重的器皿一定要保持聖潔.
我們知道救贖本身就是分別為聖.
將我們分別出來歸給神.
這也是終身的過程.

$^{641}$誠信不只是既定的事實.
我們一信主就誠信.
但也是漸漸經歷的事實.
我們要一生的追求.
我們誠信做來做去都做不到和耶穌一樣.
因為我們要像主.
當我們見主面的那一天.
我相信我們沒有一個人.
可以完全和耶穌一樣.
所以誠信是兩面的.
第一面就是我們一信主就已經聖潔.
但另一面就是我們要漸漸經歷神的經歷.
一途不僅因為和神有正確的關係.
而成為聖潔.
更要成就神的心意.
過一個公義的生活.
而成為聖潔.
成為聖潔的生活.
是過清潔,聖潔,敬填的生活.
讓我們成為歸眾的器皿.
合乎主肉.
耶穌基督不僅是我們聖潔的模範.
更是聖潔的源頭.
所以第一樣就是要自潔.
第二樣就是我們也要有朋友.
我們和什麼人要一起.
保羅提醒我們不要像少年一樣有私慾.
但我們要和清心禱告的人在一起.
所以保羅說.
你要逃避年輕人的私慾.
和那以純潔的心求告主的人.
追求公義,信實,仁愛,和平.
這些屬靈品德的追求.
是要選擇那些清心禱告主的人.
人能夠清心對任何問題.
願意在神面前多多禱告.
這種人一定適宜做我們屬靈追求的同伴.
我不知道大家會否.
時時都會和一些屬靈的人一起.
有禱告的時間.

$^{681}$或者和他們一起學習.
屬靈的人可能是傳道同工.
可能是教會的領袖.
可能是你們的組長.
茶經的組長,團契的組長.
或者主日學的老師.
又或者是你們各部不同的領袖,組長.
所以要和他們一起追求成為結伴.
有些人看起來很會說屬靈的話.
但對於屬靈的問題是不肯對付的.
對罪惡的問題是含糊了事的.
我們知道他們屬靈的情況是不實在的.
所以和清心禱告的人在一起.
是個人屬靈好不好的一個很重要的標準.
約翰韋斯理也說.
人必須有朋友.
並且要和別人交朋友.
因為從來沒有人獨自一個人進入天堂.
所以一定要交朋友.
所以我鼓勵大家.
在教會裡面有團契.
有主日學.
有小組.
多參與.
多一起祈禱.
剛才的科音樂曲.
不單止台上的人唱.
我聽到台下也有人唱.
我不知道他們的要求是怎樣.
你們也可以試試報名.
除了在福音樂曲上可以進步之外.
更加可以在當中有屬靈的友伴.
可以一起成長.
然後和他們追求公義.
信實.
仁愛.
和平.
聖靈樂意將恩賜和能力賜給我們.
使我們能為主作供.
但神要求我們先有聖潔和靈的敬拜.

$^{721}$我今天很簡單地和大家說說.
公義是甚麼呢?.
當時對提摩太太有特別的意思.
因為提摩太太要設立長老和執事.
做教會的領袖.
她一定會知道對待別人.
一定要完全公平正義.
處理神的事應當不徇私.
需要代價訓練自己.
當我們要找屬靈的領袖的時候.
追求公義是很有效的事情.
而公義的標準是從何而來呢?.
就是從神的話中而來.
我們要認.
公義不是世上不同的人所說的公義.
很重要的是我們要回到神那裡.
在神聖潔的話語裡.
活出聖潔生活.
不是遵循某種程度的人類智慧.
而是忠心的追求.
遵循聖經的教導.
聖經所說的公義.
信實.
信實不是說信徒得救的信心.
而是指信徒在日常生活上的信心上的追求.
雖然已經信主得救.
但信心上仍然追求長進.
因為聖經說.
我們勝過世界的就是我們的信心.
信心長進.
就使我們過得勝的生活.
這樣我們就可以討神的喜悅.
人愛,大家都聽過了.
大家都很著重.
我們知道我們要愛人.
但你能不能像耶穌教導我們.
要有阿加皮的愛呢?.
阿加皮的愛是指那種愛.
是保羅的書信.
時時將這種愛和信心放在一起.

$^{761}$愛心能夠讓人得到人.
讓人為主發光.
不是我們私慾的愛.
也不是道德標準上的愛.
而是神的愛.
所以人愛和平也是重要的.
和平是平安的意思.
和弟兄和睦同居是屬靈的理想生活.
所以當我們要成為聖傑的時候.
都要追求和平.
第三個是合乎主用.
預備行各樣善事.
因為保羅在第21節裡說.
他說:孕育至傑.
脫離卑賤的事.
必成為貴重的器皿.
成為聖傑.
合乎主用.
預備行各樣善事.
必成為貴重的器皿.
成為聖傑.
合乎主用.
主要用的是聖傑的器皿.
他要用清傑的器皿.
所以我們要成為孕育至傑.
成為貴重的器皿.
而成為聖傑.
我們要自己看自己是貴重.
別人看我們都覺得我們是貴重.
但我們首先要自己看自己是貴重.
因為當我們不覺得自己是貴重的時候.
別人怎樣都不會看我們是貴重.
因為我們不覺得自己是貴重.
我們看自己是貴重.
很重要的是我們走的都是貴重器皿所行.
所以當人願意除盡他心裡一切的罪惡.
和一切卑賤的時候.
我們心中自然沒有什麼東西.
是難阻著聖靈的工作.
我們讓神 讓耶穌 讓聖靈充滿我們.

$^{801}$這樣我們可以成為神的合用器皿.
是合乎大戶人家所用的貴重器皿.
我們所要作的是器皿.
既貴重又合乎主用.
我們是給誰用呢?.
是給神用的.
神是否大戶人家?.
我們已經是在大戶人家裡.
所以我們要自己合乎主用.
所以保羅在信後都提到馬可.
他都用同一個字來形容馬可.
他說馬可在他的傳道時常高有益處.
所以保羅希望提摩太.
能夠合乎主耶穌所用.
就好像馬可在傳道時常高對他有益處.
保羅內心的深處就是渴望能夠合乎主用.
所以我們都要去學習成為一個貴重器皿.
然後貴重器皿是用來做什麼呢?.
保羅亦都說是預備行各樣的善事.
因為成為合用的器皿.
神才可以用我們來完成他的善事.
我們首先被主潔淨.
然後我們除去我們的罪惡悲戰的事.
然後我們繼續不斷地去除去.
我們不願意去沾染這些悲戰的事.
我們都願意去抵抗.
或者去逃避我們肉體的愛好.
這樣我們才會專一地去信服主的旨意.
行神要我們行的善事.
沒有被主潔淨而想行善.
所行的事只不過好像一個污糟的衣服一樣.
不能夠得神的喜悅.
在神家中貴重的器皿.
是要預備行好的事.
不是用來做裝飾品.
當然大戶人家都有很多裝飾品.
但是我們願意成為合用的器皿.
不是用來擺放的.
是用來用的.
不是擺放的.

$^{841}$而各樣的善事.
包括被人看來是卑微的事.
卑微的事和卑賤的事不一樣.
保羅說我們要脫去卑賤的事.
但是我們為神做的.
有的是大事.
但亦有卑微的事.
大事被人看來.
可能在外面看來是很尊榮的.
我們當然是希望可以做.
但是卑微的事都要有人做.
我們發覺原來在神家裡面.
都要有人清潔.
都要有人做周刊的事.
別人看不到不會問.
這件事是誰做的?.
不會問的.
但那些雖然是卑微的事.
但都成為大戶人家.
或者是神的合用器皿.
我們要預備好.
看神會用我們.
為他做尊榮的大事.
還是一些很微小的卑微的事.
然後第三樣我們可以看到.
其實再說.
我們只要跟隨基督.
和基督結連.
我們就可以成為聖潔合乎主用.
這是孫迅博士.
宣導會的創辦人.
他的聖潔運動的觀念.
他強調聖潔不過是順從基督.
和基督有更加深入的聯合.
我們可以用一個圖畫來解釋.
最近大家都看到馬拉松.
渣打馬拉松.
大家知不知道在跑的過程中.
有些盲眼的人也會跑.
他們怎樣跑?.

$^{881}$有一個領跑員.
手可能綁著繩子.
和參賽的黑眼的人一起跑.
領跑員有一個口號.
不用擔心目的地在哪裡.
只要跟著我.
我就將你帶到目的地.
因為他盲的.
他看不到.
他不知道.
所以他只要跟著領跑員去跑.
簡單來說.
聖潔就是活出基督的生命.
和基督有深度的關係.
我們每步跟從主.
去活出基督.
長成就好像基督一樣.
所以孫進博士教導我們不用擔心.
什麼時候可以完成.
這個遙不可及的成聖目標.
我們不需要將信仰.
變成一個又一個的外在要求.
加在我們說.
你受洗之前要戒煙.
你要戒了這個.
然後才可以做這個.
基督就在我們身邊.
基督就在我們裡面.
今天我們就可以效法.
我們可以倚靠.
我們可以踏實地去走我們的信仰.
你願意成為一個聖潔的人嗎?.
你能不能成為一個聖潔的人呢?.
保羅說能夠.
因為耶穌基督是聖潔的主.
祂是我們的聖潔.
今天只要我們跟隨著主.
和主結連.
祂就成為我們的聖潔.
我每次講道.

$^{921}$最近都喜歡用這個結束.
大家都有見過.
於我何干?.
我跟大家說了.
我忘記了多久.
可能過了一小時.
但是大家聽了之後.
你可能可以再問自己.
你願不願意成為聖潔呢?.
你能不能成為聖潔.
已經告訴你.
你能.
但是我們要願意.
神已經走了祂的那一步.
我們只要走我們那一步.
我們就可以成為聖潔.
我們知道不是馬上就達到.
我們只是百分之一百.
光環在我們頭頂.
但我們只要努力去做.
叫我們活出一個聖潔的生命.
讓人們看到.
紅銀行的弟兄姊妹是聖潔的.
是與別不同的.
在你們的裡面.
可以將基督的榮美彰顯出來.
我們低頭祈禱.
天皇上帝.
願你的話語.
成為我們的幫助.
更加願意每一位弟兄姊妹.
都願意成為聖潔.
在他們的人生的道路上.
去快跑的跟隨.
聽我們的祈禱.
奉耶穌基督的名求.
阿們.
\newpage



\section{約翰一書 4:17-19}
\label{sec:zwEIaeiy0j4}
\textbf{2023.02.26 《幸福永伴隨》 約翰一書 4\_17-19 (和修版) 講員 : 曾裔貴牧師}
\newline
\newline
連結: \href{https://youtube.com/watch?v=zwEIaeiy0j4}{\texttt{https://youtube.com/watch?v=zwEIaeiy0j4}} ~~~~ 語音日期: 2023-03-19
\newline
\newline
\hyperref[sec:_JHmy8bvPEI]{\small{< < < PREV SERMON < < <}}
~
\hyperref[sec:index_chronic]{\small{[返順時目]}}
~
\hyperref[sec:index_scriptual]{\small{[返順卷目]}}
~
\hyperref[sec:gwCa31uE_PU]{\small{> > > NEXT SERMON > > >}}
\newline
\newline
約翰一書 4:17-19
\newline
\begin{longtable}{cl}
\hline
\hline
章節 & 經文 (和合本修訂版)\\
\hline
4:17 & \begin{tabularx}{0.7\textwidth}{X} 由此，愛在我們裡面得以完滿：我們可以在審判的日子坦然無懼，因為基督如何，我們在這世上也如何。 \end{tabularx} \\ \\ \relax
4:18 & \begin{tabularx}{0.7\textwidth}{X} 在愛裡沒有懼怕；完滿的愛把懼怕驅逐出去，因為懼怕裡含著懲罰，懼怕的人在愛裡尚未得到完滿。 \end{tabularx} \\ \\ \relax
4:19 & \begin{tabularx}{0.7\textwidth}{X} 我們愛，因為神先愛我們。 \end{tabularx} \\ \\
[1ex]
\hline
\hline
\end{longtable}
$^{1}$脫下口罩可能會方便一點.
因為老牧師下星期會回來.
我的工作平日星期三來開同工會等等都已經圓滿了.
可能要隔幾個月才來一次.
看看真樣子都好一點.
感恩.
我們先做一個簡單的祈禱.
多謝主耶穌.
我們打開聖經.
我們真的很盼望.
這個永遠的福樂伴隨著我們一生.
我們明白經文.
實踐經文.
讓我們能夠享受.
你賜給我們救恩帶來的福樂.
我們這樣感恩.
奉耶穌基督的名.
阿們.
一定要定義什麼是幸福.
第十九節.
基本上婚禮一定會讀.
很多時候婚卡都是我們愛.
因為神首先愛我們.
完全沒有錯的.
因為整個經文第七節.
第十一節.
都是勉勵信徒.
在一生裡面.
你信了主耶穌.
你要實踐彼此相愛的時候.
不是給錢.
是彼此.
不是用錢去賣愛.
是在神的裡面.
你認定自己是神的子女.
你去學習愛對方.
你就會經歷一個很大的平安.
其實十七節已經說了.
還有一樣東西.
就是我們越來越像我們的耶穌.

$^{41}$因為十七節最後.
因為基督如何.
基督在世上如何.
我們都是這樣.
什麼是幸福.
你的定義是什麼是幸福.
你下個月.
三月一日.
你就飛去日本.
吃你最喜歡的魚生.
這也是幸福.
怎麼會不是呢.
你有其他的.
有些更加大的.
什麼是幸福呢.
譬如被人認出你的身份.
也是一份幸福.
剛才來的時候.
因為今天很早起床.
已經吃了早餐.
但教完主學很餓.
就趕來.
十一點.
在大發餐廳坐下.
已經點了餐.
然後有另外一個職員.
那些醫醫,嬸.
很大年紀的.
走過來拿杯水給我.
其實桌上已經有杯水.
已經剛剛下單.
他跟我說.
你去教會嗎.
我說是呀.
我吃完會去港島.
他很開心.
我不知道他是不是信主.
他認得我是.
其實我很少下來.
他認得我.

$^{81}$可能星期日穿西裝.
沒理由這麼早去飲酒.
當然是傳道.
或者牧師.
我沒有說我是牧師.
他知道我去教會.
我說是呀.
被人認出是一份幸福.
至少我們的教會.
在街坊的眼中.
原來他留意我們.
被人看到我們這一份平安.
這一份愛.
當然.
在整段聖經裡.
這裡告訴我們.
有一個永恆的福樂.
我再跟大家讀讀第十七節.
這裡告訴我們.
如果我們不去實踐.
已經在你生命裡.
由你信耶穌的一刻.
你已經有這顆愛的種子.
你如果不去.
挑灌淋水.
你讓這個種子來生長.
你就經歷不到.
這一份平安.
和永恆的福樂.
這裡告訴我們.
第十七節.
由此.
這是一個結論.
住在愛裡的人.
其實他就住在神的裡面.
然後呢.
神也是住在他裡面.
然後就由此.
由此.
如果英文就是一個結論.

$^{121}$一個結論.
因此.
可以怎樣呢.
他說愛在我們裡面.
就得以圓滿.
是有一個很整全的.
很滿足的.
這一份的愛.
我們就可以.
在審判的日子.
坦然無懼.
就是這裡.
永恆的福樂.
就是說當我們今天.
讓這個愛的種子.
來一路長大生長.
我們就能夠.
不會懼怕.
將來神審判世人的日子.
那個恐懼.
因為我們現在.
每天已經在經歷著.
神住在我們裡面.
我們又是住在神的裡面.
透過我們去學習.
用神的愛.
其實是神透過他的愛子.
將我們拯救.
生了我們.
我們成為神的子女.
然後我們就有這個.
愛的種子.
在我們的生命裡面.
當我們去愛的時候.
我們就經驗到神.
他住在我們裡面.
約翰就有一個保證.
永恆的保證.
去到將來審判的日子.
我們可以坦然無懼.

$^{161}$不需要再懼怕.
不需要擔心將來有刑罰.
弟兄姊妹.
你想想.
這是否真真正正.
永恆的福樂.
永遠的幸福呢?.
神不再審判你.
你不需要再在審判的日子.
你好像那些未信的人.
或者堅持都不肯接受耶穌的人.
面對將來的永恆沉淪.
而是我們在主耶穌裡面.
當我們去實踐.
這一份愛的時候.
我們就能夠經驗到.
神的同住.
祂的愛.
祂的平安.
在我們內心.
或者試試用另一個角度.
因為第十八節就說.
在愛裡面.
就不會有懼怕.
懼怕裡面含著的是刑罰.
我親眼.
亦都是我自己去幫忙.
去解除了這個痛苦.
就在我們自己的舞台.
有位姊妹.
她是在教會服侍的.
有一天.
很小事的.
跟我們一位同工.
傳道同工產生摩擦.
摩擦到當面在禮堂裡面說.
大家很吵.
大家的說話很重.
然後.
大家都面黑黑的.

$^{201}$然後過了一段時間.
一兩個禮拜.
大家見面都面左左.
然後.
兩個不約而同的同學跟我說.
因為我看著那件事.
很辛苦.
大家本來沒有什麼.
弟兄姊妹沒有什麼.
大家都是一起事奉的.
就因為很小的事.
產生了摩擦之後.
大家當下又不肯放下.
就有言語上的衝撞.
分別都跟我說.
很辛苦.
很簡單.
大家說對不起就沒事了.
真的大家彼此說對不起.
一個傳道同工.
向一個姊妹說對不起.
其實就是小芬姊妹.
很開心.
馬上那面牆壁沒有了.
馬上那面隔膜沒有了.
馬上大家.
當然沒有擁抱.
因為男女不會擁抱.
馬上大家有了那份的歡欣.
那份的平安.
那份的.
愛裡面的圓滿.
你想一下.
如果你的內心.
不肯去愛對方.
愛其他人.
這裡就告訴我們第十八節.
在愛裡面是沒有懼怕的.
這個懼怕.
就是有時跟人之間.

$^{241}$人與人之間的隔膜.
人與人之間的嚴規.
人與人之間的不肯寬恕.
帶來的其實就是一種.
fear 懼怕.
所以.
這是一個阻礙你自己.
去經歷.
上帝在你生命裡面的.
同住.
是最大的阻礙.
其實每一個弟兄姊妹.
由你信耶穌的一刻.
這個愛的種子.
已經在你的心靈裡面.
你需要讓這個種子.
不受阻礙的.
來生長.
所以這裡告訴我們一個反面的就是.
在愛裡面沒有懼怕.
圓滿的愛.
是會將這些懼怕驅走.
正如剛才的一個例子.
大家都願意.
互相的.
來道歉.
馬上那個懼怕.
完全驅走.
因為愛再一次.
填滿彼此的心靈.
在教會裡面.
我們就是能夠有能力.
世界我們不能夠說.
因為大家不是有共同的價值觀.
大家都不會.
懼怕敬畏.
同一位上帝.
這裡告訴我們.
上帝差遣他的愛子來到這個世界.
使得我們得到生命.

$^{281}$這個就是愛.
所以我們大家都在這個愛的源頭.
在這個共同的價值觀裡面.
來成長.
所以.
這裡告訴我們.
在愛裡面就能夠有圓滿.
因為.
懼怕的人就含著.
刑罰.
懼怕裡面含著懲罰.
懼怕的人.
在愛裡面尚未得到圓滿.
不信耶穌的人.
他沒有這個種子.
愛的種子.
在他裡面.
信耶穌的.
如果我們不讓這顆愛的種子成長.
你都會有.
隔膜.
和人的隔膜.
和神的隔膜.
你不會經驗到一份.
圓滿的愛.
是讓你整個的心靈.
可以得到滿足.
而且.
你都不敢擔保將來.
在審判的日子.
你是否能夠坦然無懼.
所以作者約翰告訴我們.
我們可以拿捏到.
我們可以抓緊的.
就是只要我們.
實踐.
要彼此相愛.
第七和第十一節.
都已經再一次呼籲我們.
第七節說親愛的.

$^{321}$我們要彼此相愛.
因為愛是從神那裡而來的.
第十一節.
親愛的.
既然神這樣愛我們.
怎樣呀?.
就是上文所說的.
不是我們首先愛神.
是神愛我們之先.
就已經猜顯他的愛子.
為我們的罪作了贖罪制.
這個就是愛了.
既然神這樣猜他的愛子.
愛我們.
所以我們都要彼此相愛.
連續說了兩次.
再推上上文.
是叫我們不要相信那些離奇古怪的邪靈.
其實我們的內心.
比一切都大.
原來神已經住在我們裡面.
藉著他的愛子.
住在我們裡面.
更加將聖靈.
都住在我們裡面.
所以其實你和我是神的子女.
當然沒有一個機器.
你一走過就知道你是甚麼身份.
但我聽一些前輩的傳道牧者說.
邪靈認識我們的身份.
知道我們屬神.
我們真的有聖靈在我們生命裡面.
比一切都大.
所以弟兄姊妹.
你有沒有開始學習.
愛弟兄姊妹.
你來到教會.
你要開始學.
讓這個種子.
一直在你的生命裡面長大.

$^{361}$這個是愛.
很大的力量.
所以當我們來到學習.
彼此相愛的時候.
神就住在你和我的裡面.
即是說我們這個群體.
有神在我們當中.
每一次你回來.
不單止你見到的老牧師.
曾姑娘.
或其他的一些堂位.
或其他的弟兄姊妹.
這個地方你會見到神.
神會住在我們當中.
在我們這個群體裡面.
有祂的同在.
有祂的保護.
有祂的賜福.
更加你會得到的.
是一份心靈的平安.
剛才17節已經說了.
我們可以坦然無懼.
在審判的日子.
不是說今天.
是將來.
神要審判全世界的人的日子.
你就站在另一個行列.
因為你一直.
來到學習.
像主耶穌那樣愛人.
所以17節.
大家有沒有留意.
最重要的就是中間17節.
當我們可以在今天.
來享用.
將來的平安.
不會有懼怕.
審判.
但是約翰居然.
說了這句話.

$^{401}$因為基督如何.
我們信徒.
在這個世界.
也如何.
你有沒有想到這一點.
約翰.
他可以這麼肯定.
基督是怎樣的.
他親眼見過耶穌.
他跟過耶穌.
三年多.
他是其中一個.
跟耶穌最貼的門徒.
他居然告訴我們.
每一個信徒.
你可以在愛的功課上.
如果學習.
基督在世界是怎樣.
你也可以是怎樣.
哇.
可以很多東西.
很宏偉的.
他可以在愛的功課上.
可以有很多東西.
好玩的.
耶穌曾經說五餅二魚.
餵飽五千人.
我待會去廣場試試.
我拿五個餅.
叫人來.
看東西.
我現在五餅二魚.
不是這樣.
神不是要我們這樣.
走奇奇怪怪的神跡.
不是這樣.
神不是要我們這樣做.
有些宣教士.
去到一些很落後的地方宣教.
有些地方很拒絕福音.

$^{441}$那些宣教士.
真是經驗神跡.
曾經在一個雲南的山區.
有一個宣教士.
叫做富能人.
他說那個地方.
因為是一些很山區的少數民族.
叫做律叔族.
整個地方的人.
不信耶穌的.
是教鬼的.
每年去到鬼節.
會弄一些刀梯爬上去.
因為當時已經有鬼符.
爬上去是不會有事的.
但後來這個富能人.
這位英國的宣教士.
去到那個地方之後.
就改變整個少數民族.
這個律叔族.
整個族的人都信了耶穌.
他就是經驗到神跡了.
因為有一個特殊的需要.
要展現給這些人看.
這位主我們信的神是真的.
你可想而知.
神可以這樣做.
但基督如何呢.
就不是說我們現在叫死人復活.
兒子不用看醫生了.
我可以叫你復活了.
你要相信我.
不是這些.
是說基督在地上.
完全沒有保留的.
愛每一個罪人.
他愛那些麻瘋病人.
親手摸那些麻瘋病人.
主耶穌接納那些碎利和娼妓.
這一份的愛是神聖的.

$^{481}$不單止不會玷污他自己.
反而這一份的愛.
去到這些人的心靈.
是會得到醫治.
是會得到真正的解脫.
脫離罪的解脫.
約翰告訴我們.
只要我們讓這顆愛的種子.
在你的內心裡.
你去愛.
你就會經驗基督是怎樣.
你這個學習愛的肢體.
你就會像耶穌.
像耶穌.
一路蒙神的愛.
一路得到別人對你的肯定.
我教主的學.
因為浸禮班.
要和大家談一下他們的得救見證.
他們不相信我.
以前未相信耶穌.
一天吃兩包煙.
飛仔的頭髮又長.
開電單車的樣子.
他們不相信.
牧師你都不是這樣的樣子.
我說就是因為福音改變我的生命.
福音令到我們每一個人.
可以得到明白神聖的愛.
不是我先愛神.
不是因為神知道我今天會成為一個牧者.
所以祂愛我.
不是的.
大家你和我那個開始.
是在神眼中一個罪人.
是祂愛我們.
譬如剛才的施班的獻唱.
平安在神裡面.
他們穿著施班的禮服.
其實是代表著天使.

$^{521}$替大家獻唱.
這個身份很神聖.
所以他們有很多時間要準備.
要去練習.
讓最好的呈獻給上帝.
所以這些身份.
代表我們在教會是與眾不同的.
弟兄姊妹.
你在世界.
你在生活.
你有沒有受到壓制呢?.
你有沒有受到別人的歧視呢?.
你有沒有在你的困難裡面.
沒有出路呢?.
教會告訴我們.
每一個都是神所愛的.
我們愛.
因為神首先愛我們每一個罪人.
所以很盼望今早.
很簡短的一個訊息.
你有沒有讓你這一顆愛的種子.
成為你永恆的幸福的伴隨呢?.
你要讓這顆種子.
一路的生長.
求主祝福大家.
盼望我們一起學習.
我們一起祈禱.
多謝主.
再一次我們在今早.
來宣認主耶穌.
你的救恩.
是將我們每一個拯救.
雖然剛才不是每一個都可以來領聖餐.
只不過是未接受水禮.
所以未得到這個機會.
但是我們坐在這裡.
大家都是很想明白.
什麼叫做救恩.
什麼叫做信主耶穌.
盼望我們在主的大愛裡面.

$^{561}$在洪恩家.
經驗到這一份的愛.
彼此切實的相愛.
還有我們可以經驗到.
上帝就住在我們的關係.
我們這個群體裡面.
求主幫助.
多謝主.
因為是你首先愛我們.
盼望我們的身份.
變得尊貴.
是因為主耶穌的愛.
改變我們.
我們這樣的感恩禱告.
奉耶穌基督的名.
阿們,多謝各位!.
拜拜.
\newpage



\section{哈該書 1:1-15}
\label{sec:gwCa31uE_PU}
\textbf{2023.03.05 《起來,建造神的家》 哈該書 1\_1-15 (和修版) 講員 : 曾敏傳道}
\newline
\newline
連結: \href{https://youtube.com/watch?v=gwCa31uE_PU}{\texttt{https://youtube.com/watch?v=gwCa31uE\_PU}} ~~~~ 語音日期: 2023-03-20
\newline
\newline
\hyperref[sec:zwEIaeiy0j4]{\small{< < < PREV SERMON < < <}}
~
\hyperref[sec:index_chronic]{\small{[返順時目]}}
~
\hyperref[sec:index_scriptual]{\small{[返順卷目]}}
~
\hyperref[sec:F6OnOYSIQDU]{\small{> > > NEXT SERMON > > >}}
\newline
\newline
哈該書 1:1-15
\newline
\begin{longtable}{cl}
\hline
\hline
章節 & 經文 (和合本修訂版)\\
\hline
1:1 & \begin{tabularx}{0.7\textwidth}{X} 大流士王第二年六月初一，耶和華的話藉哈該先知向撒拉鐵的兒子猶大省長所羅巴伯和約撒答的兒子約書亞大祭司傳講，說： \end{tabularx} \\ \\ \relax
1:2 & \begin{tabularx}{0.7\textwidth}{X} 「萬軍之耶和華如此說，這百姓說，建造耶和華殿的時候還沒有到。」 \end{tabularx} \\ \\ \relax
1:3 & \begin{tabularx}{0.7\textwidth}{X} 耶和華的話藉哈該先知傳講，說： \end{tabularx} \\ \\ \relax
1:4 & \begin{tabularx}{0.7\textwidth}{X} 「這殿荒涼，你們自己還住天花板的房屋嗎？ \end{tabularx} \\ \\ \relax
1:5 & \begin{tabularx}{0.7\textwidth}{X} 現在，萬軍之耶和華如此說，你們要省察自己的行為。 \end{tabularx} \\ \\ \relax
1:6 & \begin{tabularx}{0.7\textwidth}{X} 你們撒的種多，收的卻少；你們吃，卻不得飽；喝，卻不得足；穿衣服，卻不得暖；領工錢的，領了工錢卻裝入有破洞的袋中。 \end{tabularx} \\ \\ \relax
1:7 & \begin{tabularx}{0.7\textwidth}{X} 「萬軍之耶和華如此說，你們要省察自己的行為。 \end{tabularx} \\ \\ \relax
1:8 & \begin{tabularx}{0.7\textwidth}{X} 你們要上山取木料，建造這殿，我就因此喜樂，且得榮耀。這是耶和華說的。 \end{tabularx} \\ \\ \relax
1:9 & \begin{tabularx}{0.7\textwidth}{X} 你們盼望多得，看哪，所得的卻少；你們收到家中，我就吹去。這是為甚麼呢？因為我的殿荒涼，你們各人卻只為自己的房屋奔走。這是萬軍之耶和華說的。 \end{tabularx} \\ \\ \relax
1:10 & \begin{tabularx}{0.7\textwidth}{X} 所以，因你們的緣故，天不降甘露，地也不出土產。 \end{tabularx} \\ \\ \relax
1:11 & \begin{tabularx}{0.7\textwidth}{X} 我命令乾旱臨到土地、山岡、五穀、新酒、新油和地上的出產，也臨到人和牲畜，以及一切人手勞碌得來的。」 \end{tabularx} \\ \\ \relax
1:12 & \begin{tabularx}{0.7\textwidth}{X} 那時，撒拉鐵的兒子所羅巴伯、約撒答的兒子約書亞大祭司，和所有倖存的百姓都聽從耶和華－他們神的話，就是哈該先知奉耶和華他們神差遣所說的話；百姓在耶和華面前存敬畏的心。 \end{tabularx} \\ \\ \relax
1:13 & \begin{tabularx}{0.7\textwidth}{X} 耶和華的使者哈該奉耶和華差遣對百姓說：「我與你們同在。這是耶和華說的。」 \end{tabularx} \\ \\ \relax
1:14 & \begin{tabularx}{0.7\textwidth}{X} 耶和華激發撒拉鐵的兒子猶大省長所羅巴伯、約撒答的兒子約書亞大祭司，和所有倖存百姓的心，他們就來為萬軍之耶和華－他們神的殿做工。 \end{tabularx} \\ \\ \relax
1:15 & \begin{tabularx}{0.7\textwidth}{X} 這是在大流士王第二年六月二十四日。 \end{tabularx} \\ \\
[1ex]
\hline
\hline
\end{longtable}
$^{1}$丙姐妹平安.
來到祂面前敬拜 原來這並非必然.
讓我們一同禱告.
天父上帝.
你因為愛我們的緣故.
讓我們有生命氣息.
在我們這一生裡 讓我們有數不盡的福氣.
最大的福氣是天父你藉著你的愛子耶穌基督.
去救贖我們的生命.
今天我們成為何等的人 是因為上帝的恩而成.
天父靠你給我們感恩的心.
去過我們每一天的生活.
去敬拜 讚美.
求你現在對我們眾人說話.
讓我們的心全然向你打開.
讓我們的耳朵聽你的話.
願聖靈在我們生命裡動功.
讓我們智慧 悟性明白.
又讓我們很樂於回應.
阿們.
今天講的經文 其實心裡也很有感受.
洪恩嘉在不久前.
才過了我們十一週年的堂慶.
大家還記得當天的氣氛.
非常熱鬧 頂尖會有很多擺相.
到現在我們一直向十二週年.
我們可以為過去的事奉.
為過去洪恩嘉的發展.
作一個總結 這是很刺激的一件事.
早上我在教會的祈禱會裡.
也有說 過了十一年.
到現在我們看到什麼.
在十一年裡.
我深信由開始到現在.
我們在教會弟兄姊妹姊妹是很深刻的.
這麼多年我們經歷了很多風浪.
不只是陽光普照的日子.
也有很多暴風雨的日子.
很多很多的轉變 有很多的眼淚.
很多的困難 灰心 失意.

$^{41}$有人離開 也有人來.
非常不容易說出來.
我們自己去看的時候 今天在這裡.
我不會說結果 我們自己心裡去想.
如果這一刻叫你總結 結果是什麼.
回想我們初頭 上帝派.
李牧師帶領一班頂尖姊妹來直堂.
到現在過了這麼多年 我們聚會的人數如何.
不要只說人數 說你生命的質素.
我們和上帝的關係.
十一年是很多的時間.
一年三百六十五天.
十一年是多少的日子 和上帝的關係有沒有盡心.
還是倒退了.
為什麼今年要推動讀經運動.
我們實在十分不熟悉聖經.
我們有很多人還停留在四呼音書.
完全沒有碰過舊約.
我最近在團契裡開始講回望 講憲制.
這些我們十世都不知道是什麼.
你說不需要舊約 我們只講新約.
你不明白舊約 也不明白新約.
十一年是什麼 我們禱告期是什麼.
在這裡經歷了上帝的什麼.
有沒有被上帝得著.
講我們生命的質素 串了多少呼音.
帶了多少人信主 在拓展上帝的角度做了什麼.
我們男女的比例.
有一天有頂尖妹在後面看起來.
男女的比例很不一樣.
我們數得到有多少人是弟兄.
我肯定我們初直堂的時候不是這麼少弟兄.
初直堂的時候有很多弟兄.
弟兄們不見了 鳳爺爺的數字.
你說每次都在講數字 不是的.
你的錢在哪裡 你的心在哪裡.
講金 講心 事奉人手.
我來到紅印堂初期.
我聽到其中一樣最多的話.
那時候這個事奉有這麼多人.

$^{81}$我們經歷黃金時間有這麼多人.
現在看起來很不同了 弟兄姊妹.
我想說今天我們坐在這裡.
不要做夢.
你說十一週年.
回顧的時候我們看到一幅什麼圖畫.
認真面對自己.
對上帝說是什麼光景.
是什麼光景.
答案是什麼.
進步 退步.
當時我記得張牧師.
差派童工團隊 差派弟兄姊妹過來.
帶著一份熱誠.
看到洪水橋區有很多人的生命.
仍然需要上帝的救恩.
在這裡有一個很大的福音和唱.
在我們附近也有異端的影響.
在這裡我們有很多很大的.
服侍的需要 對教會充滿盼望.
這十一年告訴我.
做了什麼.
是不是暴風雨將我們吹來吹去.
吹到倒塌.
為什麼只是暴風雨的問題.
我需要為自己的生命負責.
為什麼今天要講這一堂.
是時候起來了 弟兄姊妹.
不要再睡覺 別再覺得自我感覺良好.
坐在這裡很舒服 有空調.
我知道你也很關心我.
真實面對我們自己.
現在教會的光景 是時候起來了.
今天這段經文.
看回哈該書的背景.
在之前.
我們講回在王國時期 以色列的歷史.
大衛可以這樣說.
王國統一了整個以色列國.
但到了他的兒子 所羅門的時候.

$^{121}$所羅門的初期是很好的 討上帝的喜悅.
但後期當他取得很多的妃嬪.
受到很多不同國家的外邦的神明影響之後.
他自己開始陷入拜偶像的罪惡的深坑.
結果他得罪上帝.
王國開始分裂 南北兩國.
北國是以色列國 南國是猶大國.
經歷了很多年之後 北國先滅亡.
接著到南國滅亡.
滅亡的時候 被巴比倫王擄走他們的國家.
以色列人感到很大的傷痛.
猶大人感到很大的傷痛.
我們竟然是亡國 失去了自己的國家.
在這個時候 他們心裡覺得很痛苦.
上帝很有憐憫.
是上帝令到亡國的事情淋到以色列國和猶大國.
是要管教他的百姓.
是要管教他們上帝的道 行在罪惡當中.
上帝要令到這些事情臨在.
但是上帝很有憐憫 亡國不是終點站.
上帝答應他們 七十年後.
要讓百姓回歸他們的故土.
的確真的是這個時間 七十年左右.
百姓回到他們的故土.
我們看看這個圖 當時經歷了.
聶布格尼撒就是巴比倫王.
當時我們大概有五萬人.
歸回耶路撒冷.
當他們可以回去的時候.
大家的心情都很激動 但同時也很複雜.
因為已經被擄之後 在當地.
巴比倫有他們的生活 生兒養女.
開始有自己的社交圈子.
有自己的工作 一切都安頓下來.
在這個時候竟然要歸回 回到耶路撒冷.
去重建 有很多事情要重建.
耶路撒冷已經變了 不是一個很繁華的地方.
不是昔日那個被上帝大大賜福.
一個很興旺 很繁榮的地方 不是.
那個地方很荒涼.

$^{161}$很多事情要等待重建.
在這裡有幾次的被擄回歸.
大家有機會可以回去慢慢看.
中間是隔了七十年的時間.
被擄也有分三次的時間.
百姓被帶到外邦的國家.
其中一個重建是重修祭壇.
這是很重要的 重建他們敬拜的生活.
重建他們整個和神之間的關係.
重修祭壇也和重建聖殿掛勾.
裡面有兩個很重要的人物.
這本書叫做《哈該書》 就是我們的先知哈該.
上帝要透過他說話.
上帝要百姓重建聖殿.
對於以色列民來說 敬拜生活是他們的核心.
我很想和你說 弟兄姊妹.
今天我和你已經信了耶穌.
敬拜生活是我和你的核心.
你不要和我說 不是 我的生活是我的工作.
我的核心才是我的工作.
反過來 先是我的家庭.
我的家庭,工作 再到回教會.
因為回教會的時間很少.
每個星期都是星期日.
頂曬佬就是星期六.
星期三晚參加一下祈禱會.
所以核心不是敬拜生活.
核心是我工作,我家庭,我個人很多的發展.
剩下的時間就給上帝.
對不起 這個不是.
因為我和你的生活的核心都是敬拜生活.
不只是回教會的時間才敬拜.
平日我們都會敬拜.
平日上帝是我們生命的主.
不只是星期日才是我們的主.
對於以色列民來說 敬拜生活是很重要的.
所以在過程中 除了上帝要他們重建他們的城牆.
重建整個耶路撒冷.
很重要的就是重建聖殿.
最初以色列民都很開心.

$^{201}$建了那個殿的根基.
但很快就受到受敵的攪擾.
接著就要停下.
這是一個很大的困局.
大家如果讀經進度 以斯拉記.
大家讀以斯拉記的時候會發覺.
有一段是說受敵的攪擾.
以至聖殿的重建停下.
一攪擾就攪了十年.
鬆跌跌.
受敵的阻撓是很厲害.
阻撓的這段時間.
除了受敵的攪擾.
還有一樣東西.
以色列民自己都出現了問題.
在哈根書這段經文裡.
看到以色列民有一個困境.
是什麼困境呢?我們去看一看.
經文這樣說.
一章二節這樣說.
萬君子耶和華如此說.
這百姓說建造耶和華殿的時候.
還沒有到.
百姓說現在不是時候建聖殿.
現在不是時候建聖殿.
為什麼呢?.
是不是因為之前受敵攪擾這麼久.
搞到地方很荒涼.
很多困難.
整個地方根本就沒有什麼人在.
只是看到困難.
只是看到荒蕪.
沒有任何的幫助.
所以百姓停下是不是因為受敵.
不是喔,不只是這樣.
你看看下面的經文這樣說.
百姓經歷了什麼.
一章六節.
你們殺的眾多,收的卻少.
你們吃卻不飽,喝卻不足.

$^{241}$穿衣服卻不暖.
領工錢卻裝入有破洞的袋中.
這班百姓在這個時候說.
建造耶和華殿的時候還沒有到.
是不是因為受敵攪擾.
還是因為現在生活有個困局.
生活的困局是什麼?殺種的多.
我付出很多.
我種大麥,種小麥.
以色列文是這樣.
猶大人種大麥,小麥.
其實這段時間說殺種多.
已經過了大麥和小麥的時間.
過了,殺的眾多,種了很多東西.
但是收的卻少,收成很少.
收成和殺種不成比例.
這個很大影響.
對於他們來說.
這些是日常的飲食.
糧食不足夠,這樣就會捱餓.
是很重要的.
你說的是糧食,吃也不夠飽.
是不是?.
接著還有渴不得足.
渴不得足就是那裡有乾旱.
水不夠.
我對於這一點,我印象.
我的感受特別深刻.
因為我曾經在港島裡說過.
我去過中東.
有跨文化的體驗.
我去過這些沙漠地帶.
中東是比較少植物.
很乾旱的地方.
在這裡沒有水是很辛苦.
我試過喝一瓶水.
很快就喝下去.
但是感覺身體好像穿了個洞.
喝完好像沒有喝.
很想繼續喝第二瓶.

$^{281}$第三瓶水.
不停地喝.
不止渴.
那個地方周圍是乾的.
曬的,熱的.
好像蒸你.
不停地乾蒸,蒸得很不舒服.
如果說這個地方.
渴不得足.
就是很乾旱.
很辛苦.
水是我們的生命資源.
沒有水是很辛苦.
穿衣服不得暖.
當然整體的經濟收入都少.
穿衣服不足夠.
還有最慘的是什麼?.
領工錢.
我辛辛苦苦勞碌.
得到的那份工資.
很快就沒有了.
領了工錢.
裝進有破洞的袋子裡.
不忘保守.
很快就沒有了.
那我怎麼辦?.
我怎麼生活?.
是我生活很困苦.
我吃喝穿都不行.
又沒有錢.
很困苦,是不是?.
是不是因為以色列民遊困苦呢?.
就是因為這樣的處境.
所以他們不建聖殿?.
不建聖殿?.
但是你記得上帝怎麼說的?.
在這麼困苦的時候.
你記得百姓在做什麼事嗎?.
百姓在做什麼?.
上帝怎麼說?.

$^{321}$一章四字.
這殿荒涼.
你們自己還住天花板的房屋嗎?.
是的,這殿很荒涼.
這個時候有很多困難.
有仇敵的搞擾.
生活很困苦.
但是這麼困苦.
人們都照著什麼?.
用盡力.
建自己的房屋.
天花板的意思是不是有遮蓋?.
你就在建房屋.
你想在裡面好好享受.
可以住,住好房屋.
上帝的殿是怎樣?.
連蓋子都沒有.
上帝的殿就荒涼.
在這裡用的字眼.
這殿荒涼.
沒有蓋子.
說什麼敬拜?.
你不理會.
不理會敬拜生活.
不理會自己和上帝的關係.
只顧著自己的房屋.
這麼困難都要顧著自己的房屋.
你說人的焦點放在哪裡?.
以色列人的困境.
後面上帝會解釋.
為什麼以色列人這麼困難?.
其實兩件事是互相有關連的.
以色列人從一開始想的焦點.
不是想要重建聖殿.
就算被人阻撓了.
我都要怎樣?.
總之有機會我就要建立.
不是這樣想.
而是我房屋比較重要.
我沒有地方住.

$^{361}$沒有地方睡.
我哪有時間理你?.
上帝現在都先理我自己.
我工作有很多需要.
先照顧我的工作.
我家庭也有很多需要.
上帝我先說說.
我先照顧好我的家庭.
當然家庭最重要.
我隔我的房子最重要.
我生活最重要.
什麼都是我自己最重要.
我有時間就來理你.
我生活中心就相反.
我先,才到你.
但你看到.
這樣的先後次序是錯的.
上帝指出這樣是錯的優先次序.
先應該求神的國和神的義.
先求上帝的電.
先求上帝的關係.
你先求自己.
以色列文反過來.
先求自己的好處.
先求自己的益處.
然後把上帝放在一邊.
結果上帝就可以說.
困境,為什麼會有困境?.
這麼淒涼.
吃又不可以吃,穿又不暖.
喝又不可以,為什麼?.
後面有說的,上帝自己解釋.
不是因為外在的什麼人造成.
上帝自己動手.
上帝會這樣嗎?.
是呀,我們的次序錯了.
我們這件事不對的時候.
上帝親自出手.
你看到一章九節說什麼?.
你們盼望多得.

$^{401}$看那所得的卻少.
你搏命眾,搏命殺種.
你想收割大麥,小麥.
搏命殺種,收割少.
你收到家中,我都會催去.
你有收穫,我都幫你催去.
為什麼?.
上帝的困境的原因是什麼?.
因為我惦芳良.
你個人就為自己的房屋奔走.
只顧著自己.
所以就是這樣.
上帝的管教臨到.
說得白一點.
一章十節,所以因為你們的緣故.
天不降甘露.
現在不是誰欺負我們的緣故.
是因為百姓優先次序錯了.
在這裡犯罪,得罪上帝.
所以天不降甘露.
天降甘露一定是出於上帝的作為.
上帝要降甘露,地才有出產.
現在就是天不降甘露.
所以地不出產.
你就沒有收成.
上帝的經文說得很清楚.
上帝命令乾旱.
臨到土地,山崗,五穀,新酒,新油.
和地上的出產.
上帝的工作.
也臨到人和生畜.
所以人和生畜都不夠水喝.
以及一切人手勞碌得來.
上帝的管教臨在.
人就缺乏.
沒有足夠的收成.
飲水不足夠,穿衣服不夠暖.
賺錢沒有用.
不務保守,全部流失了.
上帝的管教.

$^{441}$是因為我們的緣故.
這裡經文說了不止一次.
一章七節說了什麼.
你們要審查自己的行為.
不只是說.
上帝現在很慘.
有困局在.
我不明白為什麼.
我很努力的.
為我的家庭做很多事.
為我的子女做很多事.
為我的配偶做很多事.
為什麼現在的結果是這樣.
有沒有審查自己的行為.
有沒有事情是上帝不起悅我們.
我們從來都不會問的.
只會看困局.
你不要想困局的原因是什麼.
不是發生同一個結局.
所有事情都是同一個原因.
是要問,是要去審查.
這裡一章七節.
你們要審查自己的行為.
要審查.
上帝說得很明顯.
是因為我們的緣故.
我們得罪上帝的緣故.
我們得罪神.
我們有很多神的電荒涼.
我們不理會祂敬拜的生活.
在這裡,哈該先知.
就是神要去發訊息的一個出口.
上帝藉著哈該先知和以色列人說話.
但是有先後次序.
上帝先透過哈該先知向誰說話.
想問問在座各位.
先向誰說.
一次過一起說嗎.
經文有說,先向誰說.
非常好.

$^{481}$撒拉鐵的兒子.
猶大省長索羅巴伯.
他是政府官員.
是他們的領袖.
政治的人物.
政府官員.
約撒達的兒子約書亞.
這是宗教界的代表.
政治界的領袖.
宗教界的領袖.
哈該先知先向他們說.
為什麼.
有影響力.
還有要反省.
要審查自己的行為.
不只是下面的百姓審查.
領袖要先審查.
所以我經常祈禱.
要認罪.
不只是跟弟兄姊妹說.
弟兄姊妹,想想.
我們每個人都有罪.
我們同工都有罪.
教會領袖要一同認罪.
我們自己認罪悔改.
向上帝復和.
才有真正的影響力.
叫弟兄姊妹認罪悔改.
信息要臨到,要改變.
要審查自己的行為.
透過領袖.
領袖要知道做什麼.
上帝喜悅什麼,不喜悅什麼.
領袖才改變了.
才可以影響百姓.
要一同改變.
所以今天同樣地對於洪恩嘉來說.
要改變.
也要由領袖先開始.
領袖不只是同工,同位.

$^{521}$在我們當中也有弟兄姊妹.
參與領袖的事奉當中.
我們是一同有份的.
上帝會先透過我們.
傳達這個信息.
要我們改變.
然後我們才會和其他弟兄姊妹.
一同回應上帝.
走出困境的方法.
我們來看看是什麼.
走出困境的方法.
第一件事,要審查自己的行為.
知道自己哪件事.
不討上帝的喜悅.
就是要認罪悔改.
要認罪.
不要硬著心腸.
我們有很多次的祈禱會.
都有認罪.
但我總覺得那個認罪.
未是很徹底.
弟兄姊妹.
我們是否真的知道自己在那些位置.
犯罪得罪上帝.
還是覺得總是前後左右的人有問題.
不是自己有問題.
是他們有問題.
不是我們說他們有問題.
先看自己在哪裡得罪上帝.
願意聖靈光照我們.
真的認罪悔改.
這個很重要,要審查.
之所以說不要睡覺就是這樣.
迷迷糊糊回來生活.
這麼多年都是這樣.
不要迷迷糊糊.
審查,認罪,悔改.
悔改還有一件事.
一張白紙.
你要上山取木料.

$^{561}$行動.
取木料建造者殿.
你將優先次序放在正確的位置.
重建聖殿.
重建不只是用筆口說.
來,我們重建神的殿.
取木料,行動.
我很記得.
在唐慶的時候.
很多弟兄姊妹說了什麼?.
很多弟兄姊妹的心願是什麼?.
有沒有人說得出?.
有不少人說希望紅印加獨立.
希望我們這樣那樣.
對教會有很多期望.
很希望我們復興.
這是弟兄姊妹你們的心願.
那些不是口號.
我要紅印加復興.
我要紅印加很多人.
要傳很多福音.
帶很多人回來.
我要我們快點獨立.
好像口號那樣.
有沒有行動呢?.
要有行動.
什麼行動?.
我們看回.
看回那些圖畫.
我們看回在以色列文當中.
其中一個很震撼的.
建築.
這裡用一個形容詞.
第一聖殿.
第一聖殿就是那時候.
以色列國.
第三任的君王.
所羅門所建造.
第一任君王是蘇羅.
蘇羅之後就到大衛.

$^{601}$大衛傳給他的兒子.
所羅門.
神選擇了所羅門去建聖殿.
其實是很金碧輝煌.
很漂亮.
很宏偉.
當時以色列人應該引以為驕傲.
也是神為他們預備了很多材料.
很漂亮.
很高貴.
很有價值的材料.
建築得這麼漂亮.
百姓可以圍繞著這個聖殿.
去敬拜上帝.
和上帝建立關係.
這裡有很多付出.
如果聖殿這麼漂亮.
有人去獻上物資.
你想想.
金銀.
木.
還有石.
多少的物資才能建造這麼漂亮的聖殿.
除了聖殿還有人手.
建的時候牽涉了很多人.
搬搬抬抬建的.
還有很多手工的.
很漂亮的.
聖殿裡面還有婚.
有分有致聖所.
有聖所.
裡面有很多不同的部分.
你說這裡牽涉了多少人.
還有.
不只是有聖殿.
之後有敬拜生活.
有獻祭.
有祭司.
有大祭司.
有事奉人員.

$^{641}$有多少人一同擺上.
一個聖殿牽涉的人有多少.
好了.
去到.
去到什麼時間呢.
去到BC586的時候.
這個聖殿被巴比倫摧毀了.
之後第二聖殿.
就是所羅巴伯所建造的.
所羅巴伯建造的這個.
就是樣子很不一樣了.
就算是所羅巴伯建造的聖殿很不一樣.
裡面仍然牽涉很多的物資.
很多很多的資源.
很多的人手.
一樣有事奉人員的參與.
是要有行動的.
我不可以只是口說.
我現在要建聖殿了.
口號一天一天都在喊口號.
今天.
我講11週年.
就要邁向12週年的時候.
我們說要建造洪恩家.
頂梓妹起來的時候不是一個口號.
要有行動.
行動裡面包括了很多東西.
第一樣東西.
要有很充足的奉獻.
是要有很充足的奉獻.
你看到聖殿那裡.
是有很多很多人捐了很多的物資.
有很多不同的資源.
才可以建造這麼漂亮的聖殿.
今天你說洪恩家要有一個大復興.
一個復興要更多的傳播音.
容納更多人聚會.
有更多的事奉.
就需要頂梓妹的奉獻.
怎樣奉獻呢?.

$^{681}$我首先說.
最基本要做足11奉獻.
最基本.
但通常我們會做多過11.
11奉獻怎樣做?.
不是這樣的.
阿娜上帝我出糧了.
我先洗到差不多.
我去到月尾那天.
我發覺我的袋子剩下多少.
我就給你.
剩下這麼多就這麼多.
可能有時候不夠十分之一.
沒辦法.
上帝我有需要的.
我有需要.
我家裡有很多需要.
靠我養家的.
所以我給你這麼多已經很好了.
這個不是這樣的.
11奉獻不是這樣做.
你一個月開始的時候有收入.
就先拿出來.
先給上帝.
因為你去到月尾的時候.
才說看有多少.
一定不夠.
我寫包單.
是不是?.
我先放上我的11.
上帝會傾福於我.
在馬拉基書.
這個福氣多到無處可用.
我先放上我的奉獻.
上帝就會傾福給我們.
你說我也很久沒工作了.
怎麼會有收入?.
你只找一些在職人士就好了.
沒工作是不是沒收入呢?.
我想問.

$^{721}$完全沒錢花?.
我們的錢來自哪裡?.
我想說收入不是.
賺取金錢工作才有.
就算這麼久.
利是錢我都會拿奉獻出來.
有不同的方面.
上帝透過不同的人.
不同的層面給我錢.
我都會抽出十分之一.
其實更多的時候會多過.
我相信弟兄姊妹的奉獻.
會多過十分之一.
你說你對其他教會說就好了.
不是的,我想說.
有一件事我知道弟兄姊妹.
當我們心願意的時候是可以的.
弟兄姊妹的神給你是很充足的.
為什麼這樣說?.
前一陣子我們為趙姑娘.
醫療籌款.
弟兄姊妹很熱心.
因為大家都很愛趙姑娘.
但我回看的時候.
我知道不是每個人都和趙姑娘這麼熟.
不是每個人都這麼熟.
但我都願意放在心上.
因為真的愛她.
真的很願意出一分力去給她.
是不是每個人都高收入人士.
所以我給她?一定不是.
還有人根本不認識她.
不是因為我之前和她有很深厚的友誼.
所以我奉獻給她,不是.
是因為我真的很想表達一份愛.
我不認識她,我也想表達.
我也想支持她.
結果那個奉獻有一個比預期超出很多的.
一個很理想的效果.
我深信趙姑娘的家人會很感動.

$^{761}$她感受到教會的愛.
弟兄姊妹一個肢體有需要.
我們尚且會這樣.
何況整個教會是神的家.
是神的殿.
是神的國度的拓展.
我們是否應該有更多的擺上奉獻?.
是不是?.
不要輕看自己.
你的心是放在上帝那裡.
一定能奉獻.
上帝一定給我們這樣的能力.
我從趙姑娘的事件就看到.
是可以的.
很多時候我想說.
一些能夠去構置堂祀.
去重建教會的教會.
不一定是很有迷.
通常是弟兄姊妹很願意奉獻.
很痛心的教會.
不是說我很有迷.
所以很隨便.
我奉獻一部車.
一個樓的價錢.
有幾個人奉獻出來.
我們隨便吧.
下一次就去買新的堂祀.
不是.
而是那間教會很痛心.
弟兄姊妹很認定.
上帝的國度,神的家很重要.
就會放上錢.
奉獻最實際.
不要說那麼多.
我們要建什麼?.
我們要復興.
要做什麼?.
要傳很多福音.
傳福音不用錢嗎?.
你知不知道牽涉到很多東西.

$^{801}$整間教會的營運.
弟兄姊妹.
經過今天之後.
我鼓勵你重新想想.
為了上帝的國.
為了神的家.
放上多少?.
很多付出.
事奉人員也是.
事奉人員不是這樣的.
集中在一小撮.
那一小撮人同時是靈場.
同時是師班.
同時是師事招待.
同時有很多不同的事奉人員.
探訪也是他們.
傳福音也是他們.
不是一小撮人.
神的家在這裡.
我們需要彼此互相配搭.
上帝給我們的恩賜是很豐富的.
我不相信.
我不相信只有一小撮人有恩賜.
我相信在座的弟兄姊妹都有恩賜.
你願意用出來.
我保證是很豐富.
上帝很愛我們.
給我們的是比我們所求所想的多.
所以弟兄姊妹要拿出來.
無論在那裡我是搬搬抬抬的.
我在聖殿裡面.
畫一些做手工的.
雕塑也好.
刺繡也好.
我是那些做獻祭的人士.
或者是有其他的也好.
我是有份的.
是每個人都有份的.
弟兄姊妹你問問自己.
今天你在教會的事奉如何.

$^{841}$真的每個人都有事奉崗位?.
你真的將上帝給你的東西都用在這裡?.
問問自己.
你真的要起來建造神的家.
要用洪恩加福興.
是需要大家一起來.
今天早上在禱告會裡.
也有一個很漂亮的圖畫.
我很感恩很多弟兄姊妹出席祈禱會.
大家都來自不同的團契.
有不同的事奉.
有不同的成長背景.
但大家都同心.
願意為教會禱告守望.
就是要這一份同心.
弟兄姊妹這裡不單是說事奉人員.
是奉獻禱告.
是不是同心.
你說我不一定唱得歌.
好也不要緊.
神媲美的人不一樣.
我不一定做到主席.
甚至施事招待.
我可能這樣.
我有一些限制.
什麼都好.
我這些不能做那些不能做.
我告訴你禱告一定做到.
禱告也不限制幾個人禱告.
不會這樣的.
禱告的權是上帝給我們每一個.
應該全部人都可以一起禱告.
教會現在很大需要.
很大需要.
大家知道.
教會同工現在是什麼情況.
知道教會現在什麼光景.
不是靠人.
單靠人.
是需要仰望上帝的憐憫.

$^{881}$求上帝幫助我們.
所以弟兄姊妹.
回來禱告.
你說我平日不行.
星期日一個月總有一次.
第一週星期日上午十點十一點.
我們等你.
你說我不懂說.
很怕在眾人面前說什麼.
是不需要的.
不需要.
回來一起禱告.
就會學會怎樣禱告.
也求上帝幫助我們.
教我們.
不是要回來作文.
說得多華麗.
那邊土持真的.
我不知為什麼旁邊的人那麼會禱告.
他的禱告文真的很好.
很漂亮.
我一出來禱告.
我就打卡了.
我都不懂說.
很怕尷尬.
我不要回來禱告.
不需要的.
真誠地向上帝禱告.
不在乎多少.
懂不懂說.
是不是很華麗的詞語.
多一個心出來.
一起禱告.
越來越多人禱告.
教會的改變會更大.
因為上帝會在我們當中工作.
是上帝在等我和你.
等我們回來禱告.
如果可以的話.
星期三晚的Zoom祈禱會等你.

$^{921}$大家覺得.
我住得遠.
我回來又很辛苦.
很多事情.
Zoom呢.
Zoom祈禱會不行嗎.
其實更好是回來教會.
一同禱告守望.
經歷神.
禱告.
事奉.
將你整個人心交給上帝.
上帝喜悅我們.
全程投入敬拜讚美祂.
回來崇拜.
不要回來看Show.
經常說.
你不是看表演.
今天這個小組獻唱如何.
好不好聽.
好聽.
或者不是這樣.
今天主席做得如何.
先看看.
講完有沒有吹.
是不是動督笑.
我自己回來是不是感覺良好.
這樣就很久了.
這樣崇拜.
不是的.
弟兄姊妹我們所有人都是敬拜者.
只有一個觀眾是誰.
只有上帝.
祂等我們.
你坐在台下一樣獻祭給上帝.
坐台上有台上獻祭.
台下一樣.
你獻什麼給上帝.
每個星期回來崇拜.
問自己帶什麼給上帝.

$^{961}$是不是經常有人交給祂.
你如果回來看Show.
就不要回來了.
我真的這樣說.
弟兄姊妹起來.
讓我們重見神的家.
不是哭口號.
起來.
你們要上山取木料.
建造這殿.
我就因此起來.
且得榮耀.
這一句因此起來.
比較好的翻譯是.
我喜悅他.
喜悅什麼.
上帝會喜悅那個殿.
那個聖殿.
且得榮耀.
上帝會得榮耀.
今天我們復興紅印家.
上帝會喜悅我們.
上帝會因為我們的家.
得榮耀.
我們不是虧損上帝的榮耀.
是不是.
如果走出去.
別人說那間教會這樣.
爛爛的.
真的很噁心.
我們簡直虧損上帝的榮耀.
上帝會很不開心.
但不是的.
出去榮耀神.
上帝會很歡喜快樂.
很心願.
我們上帝的心.
是歡喜快樂.
心願領袖首先起來.
審查我們的行為.

$^{1001}$認罪悔改.
然後帶領弟兄姊妹一起.
歸向神.
重建神的家.
我們一同禱告.
此下這段經文要成為我們的提醒.
天父我們要向你認罪.
這麼多年.
我們都活在自己的世界.
以色列民.
先為自己的家.
去籌算.
放下神的家.
其實我們很多年都是這樣.
我們都是為自己籌算.
我們放下上帝的家.
天父寬恕我們.
寬恕我們很多的罪惡.
嘴唇的心靈的.
行事為人的.
我們有很多很多的罪惡.
上帝我知道你的心很痛.
求你寬恕我們.
寬恕我們作領袖的.
寬恕我們弟兄姊妹的.
我們有很多很多的罪.
天父求你寬恕我們.
讓耶穌基督以你的寶血潔淨.
遮蓋我們的罪污.
主啊求你光照我們生命.
讓我們知道.
我們真的這樣得罪你.
讓我們審查自己的行為.
認罪悔改.
聖靈啊.
求你充滿我們生命.
給我們源源不絕的力量.
讓我們悔改.
讓我們成長.
讓我們興建.

$^{1041}$去建造神的家的時候.
靠著你的得力.
有錢出錢.
有力出力.
將我們整個人擺上.
主啊求你重整我們整個生命.
讓我們敬拜你.
讓我們將和你的關係.
讓我們將敬拜擺在首位.
然後信多個上帝.
你就會賜福.
讓我們有這份信心.
求著私恩憐憫.
奉耶穌基督的名字.
祈禱.
阿們.
\newpage



\section{哈該書 2:1-23}
\label{sec:F6OnOYSIQDU}
\textbf{2023.03.12 《從今日起,我必賜福.》 哈該書 2\_1-23 (和修版) 講員 : 曾敏傳道}
\newline
\newline
連結: \href{https://youtube.com/watch?v=F6OnOYSIQDU}{\texttt{https://youtube.com/watch?v=F6OnOYSIQDU}} ~~~~ 語音日期: 2023-03-22
\newline
\newline
\hyperref[sec:gwCa31uE_PU]{\small{< < < PREV SERMON < < <}}
~
\hyperref[sec:index_chronic]{\small{[返順時目]}}
~
\hyperref[sec:index_scriptual]{\small{[返順卷目]}}
~
\hyperref[sec:bMJz4f2t79c]{\small{> > > NEXT SERMON > > >}}
\newline
\newline
哈該書 2:1-23
\newline
\begin{longtable}{cl}
\hline
\hline
章節 & 經文 (和合本修訂版)\\
\hline
2:1 & \begin{tabularx}{0.7\textwidth}{X} 七月二十一日，耶和華的話藉哈該先知傳講，說： \end{tabularx} \\ \\ \relax
2:2 & \begin{tabularx}{0.7\textwidth}{X} 「你要曉諭撒拉鐵的兒子猶大省長所羅巴伯、約撒答的兒子約書亞大祭司，和所有倖存的百姓，說： \end{tabularx} \\ \\ \relax
2:3 & \begin{tabularx}{0.7\textwidth}{X} 『你們中間存留的，有誰見過這殿從前的榮耀呢？現在你們看如何？在你們眼中豈不是如同無有嗎？ \end{tabularx} \\ \\ \relax
2:4 & \begin{tabularx}{0.7\textwidth}{X} 所羅巴伯啊，現在，你當剛強！這是耶和華說的。約撒答的兒子約書亞大祭司啊，你當剛強！這是耶和華說的。這地的百姓啊，你們都當剛強做工，因為我與你們同在。這是萬軍之耶和華說的。 \end{tabularx} \\ \\ \relax
2:5 & \begin{tabularx}{0.7\textwidth}{X} 這是照著你們出埃及時我與你們立約的話。我的靈仍要住在你們中間，你們不必懼怕。 \end{tabularx} \\ \\ \relax
2:6 & \begin{tabularx}{0.7\textwidth}{X} 萬軍之耶和華如此說：過些時候，我必再一次震動天地、滄海與乾地。 \end{tabularx} \\ \\ \relax
2:7 & \begin{tabularx}{0.7\textwidth}{X} 我必震動萬國，萬國的珍寶都必運來，我就使這殿充滿榮耀。這是萬軍之耶和華說的。 \end{tabularx} \\ \\ \relax
2:8 & \begin{tabularx}{0.7\textwidth}{X} 銀子是我的，金子也是我的。這是萬軍之耶和華說的。 \end{tabularx} \\ \\ \relax
2:9 & \begin{tabularx}{0.7\textwidth}{X} 這後來的殿的榮耀必大過先前的榮耀。這是萬軍之耶和華說的。在這地方我必賜平安。這是萬軍之耶和華說的。』」 \end{tabularx} \\ \\ \relax
2:10 & \begin{tabularx}{0.7\textwidth}{X} 大流士王第二年九月二十四日，耶和華的話臨到哈該先知，說： \end{tabularx} \\ \\ \relax
2:11 & \begin{tabularx}{0.7\textwidth}{X} 「萬軍之耶和華如此說，你要向祭司請教律法，說： \end{tabularx} \\ \\ \relax
2:12 & \begin{tabularx}{0.7\textwidth}{X} 『看哪，若有人用衣服的邊兜聖肉，這衣服的邊接觸了餅，或湯，或酒，或油，或別的食物，這些是否成為聖呢？』」祭司回答說：「不。」 \end{tabularx} \\ \\ \relax
2:13 & \begin{tabularx}{0.7\textwidth}{X} 哈該又說：「若有人因摸屍體染了不潔淨，然後接觸任何東西，這東西就變為不潔淨嗎？」祭司回答說：「必不潔淨。」 \end{tabularx} \\ \\ \relax
2:14 & \begin{tabularx}{0.7\textwidth}{X} 於是哈該說：「耶和華說，在我面前這民如此，這國也是如此；他們手裡的各樣工作都是如此；他們在那裡所獻的都不潔淨。」 \end{tabularx} \\ \\ \relax
2:15 & \begin{tabularx}{0.7\textwidth}{X} 「現在，你們心裡要想一想，從今日起，耶和華的殿還沒有一塊石頭放在石頭上的情況。 \end{tabularx} \\ \\ \relax
2:16 & \begin{tabularx}{0.7\textwidth}{X} 那時你們怎麼了？有人來到二十斗的穀堆那裡，卻只得了十斗；有人來到酒池那裡要取五十桶，卻只得了二十桶。 \end{tabularx} \\ \\ \relax
2:17 & \begin{tabularx}{0.7\textwidth}{X} 我以焚風、霉爛、冰雹攻擊你們，和你們手上的各樣工作，你們仍不歸向我。這是耶和華說的。 \end{tabularx} \\ \\ \relax
2:18 & \begin{tabularx}{0.7\textwidth}{X} 你們心裡要想一想，從今日起，就是從這九月二十四日起，從立耶和華殿根基的日子起，你們心裡想一想： \end{tabularx} \\ \\ \relax
2:19 & \begin{tabularx}{0.7\textwidth}{X} 倉裡還有穀種嗎？葡萄樹、無花果樹、石榴樹、橄欖樹雖沒有結果子，從今日起，我必賜福。」 \end{tabularx} \\ \\ \relax
2:20 & \begin{tabularx}{0.7\textwidth}{X} 這月二十四日，耶和華的話再次臨到哈該，說： \end{tabularx} \\ \\ \relax
2:21 & \begin{tabularx}{0.7\textwidth}{X} 「你要告訴猶大省長所羅巴伯說，我必震動天地， \end{tabularx} \\ \\ \relax
2:22 & \begin{tabularx}{0.7\textwidth}{X} 傾覆列國的寶座，除滅列邦列國的勢力，並傾覆戰車和坐在其上的。馬和騎兵都必跌倒，各人被弟兄的刀所殺。 \end{tabularx} \\ \\ \relax
2:23 & \begin{tabularx}{0.7\textwidth}{X} 萬軍之耶和華說：撒拉鐵的兒子我僕人所羅巴伯啊，這是耶和華說的，到那日，我必以你為印，因我揀選了你。這是萬軍之耶和華說的。」 \end{tabularx} \\ \\
[1ex]
\hline
\hline
\end{longtable}
$^{1}$記者:記者會平安,我記得上星期已經講了《哈該書》第一章,當時就在想,講了一半,講了我要講的內容,就停在那裡.沒想到世事都有很多變化,教會裡面也有很多變化,結果今天可以講《哈該書》的第二部分..
我覺得跟我們紅恩堂的處境也很配合,我深信上帝也有他的心意,想我們聽到他的聲音.我們一同禱告..
願你的話語臨到我們當中,願聖靈感動我們每一個,讓我們明白上帝的心意是什麼,讓我們有力量去回應,奉耶穌基督的名字祈禱,阿們..
上次我記得鼓勵弟兄姊妹,我們說起來,建造神的家,上帝的心意是要我們起來,是時候起來了,紅恩家要起來去重建..
在過程中,不是只哭喊口號,重建重建,而是我們要有實際的行動,有錢出錢,有力出力去建立這個家..
還有我們這裡每一個都有份.今天我們再接續上個星期的講道,就是講這一句,從今天起,我必賜福..
我想問問大家,剛才也有一起讀經,大家記不記得這句話是誰說的呢?.
很害怕,剛剛讀完禮,都不知道是誰說的..
沒錯,是上帝說的,從今天起,我必賜福..
這個今天是指什麼日子呢?是說建立聖殿根基的那一天,重建聖殿的那一天,上帝說我必賜福..
在這個之前,記得上次我們講哈該書的時候,那個背景,今天再去一下以色列記講多一點背景..
大家記得波斯王的時間,原來以色列民猶太人被擄70年的時間快滿了..
上帝要應驗他自己所說的話,於是感動波斯王,然後波斯王鼓勵以色列民回到耶路撒冷重建神的殿,這是之前的背景..
很特別,一個愛邦的君王卻成為了神手中的器皿,上帝透過他鼓勵百姓回到自己的地方重建聖殿,重建自己的地方..
於是以色列記一章五節的時候,就說猶太人,變亞文的一些領袖等等都被神感動,要上耶路撒冷去建造耶和華的殿..
然後就說到上次說到的,哈該書一章十四節,上帝去鼓勵去激動去激發撒拉鐵的兒子所羅巴白,和約撒達的兒子約書亞大濟斯,和剩下來的百姓去回去做神殿的工作..
上次其實不是說得很多,我只是說到所羅巴白是一個政治的領袖,其實他是君王的後人,是大衛的後人,是約瓦根的後人,所以你可以說他是有君王的血統,他這個時候是做一個政治的角色,帶領百姓..
而約書亞大濟斯很明顯是宗教上帶領百姓去重建的很重要的人,他和所羅巴白一直緊密的合作,一起的配搭..
除了這兩位領袖,也有上帝的百姓,大家一起的,這是上次我們看到的,哈該先知向以色列的百姓說話..
但是我們記得重建聖殿是不順利的,大家還記得,建了聖殿根基之後,有仇敵的搞擾,有十年多的時間他們停了聖殿重建的工作..
同時,你看到以色列民,你看到猶太人很注重重建自己的房屋,就算他們的生活條件,當時生活的環境很困難也好,你看到百姓很體重自己的家,花很多時間,很多精力,怎樣也要建好自己的家,要處理好自己的生活,卻將聖殿放在一旁..
所以上次我說到,是時候重建,上次我們回顧洪恩嘉,過去11年的時間,經歷了很多風浪,看回我們的屬靈生命,看回我們的聚會人數,我們事奉各方面,.
心上頂,只會心裡有數,是時候要重建,來到第二章的時候,我們看回的時候,神說了這一句,從這一天起,就是你建聖殿根基的這一天,我要去祝福..
百姓很聽話,在哈該先知鼓勵他們之後,他們重新去建立聖殿,上帝要鼓勵他們,既然大家是回應神的,起來建聖殿,上帝說,我要去祈福,要給福氣,但是這裡神不只是說給福氣,在裡面有很多提醒,很多鼓勵..
首先這裡有神的鼓勵,我們看回經文,神鼓勵什麼呢?忘記背後,努力面前,要重建,重建一個群體,當日重建聖殿,今天重建教會,要忘記背後,努力面前..
為什麼這樣說,要忘記背後呢?忘記什麼背後呢?這次當猶太人重建聖殿的時候,難免會有些比較,和什麼時候比較呢?和所羅門時期的聖殿比較,我們建了之後是怎麼樣的呢?.
昔日所羅門的聖殿多麼的輝煌,材料多麼的名貴,閃閃生輝,今天我們是什麼材料,建了之後是什麼質地呢?都害怕,很害怕那種比較,因為比較可能更加會疾步了,我不想向前..
我們看回這裡的圖,雖然是有些細,大家去看一看,在左邊的角落比較細的位置,就看到是所羅門的聖殿..
大家看到左邊其實是所羅門的聖殿,是在959年的時候BC966到959年的時候建的,其實是非常漂亮的,非常的輝煌..
但是去到586年的時候,BC586被毀,被巴比倫王所毀,但是至少那時候建好的時候很震撼,以色列民心心記在心中,直到現在,被擄的人歸回耶路撒冷,還記得多麼的漂亮..
所以心裡面會有這個比較,神在這裡要鼓勵他們忘記背後,為什麼?他說二章三節,你們中間存留的,就是以前看過這個聖殿的榮耀的人,見過他何等崇高的樣子..
現在你們看如何?你們現在看是怎樣?那個聖殿還在嗎?不在了,當年怎樣的輝煌都好,現在都不在了,在你們眼中豈不是如同沒有嗎?現在沒有了,我們經常說過去怎樣?經常說過去,其實那些已經成為過去,不再存在..
所以在過程裡面忘記背後,努力面前,無論過去我覺得有什麼好,有什麼做得很出色,感覺很有成就都好,過去了,努力面前..
在這裡,神鼓勵他們,忘記背後,努力面前,然後就是要剛強,好像昔日約書亞收到一個鼓勵,神鼓勵他,你當剛強壯膽,當約書亞接了摩西的棒的時候,神要叫他剛強壯膽..
今天這裡,神同樣叫人剛強壯膽,但是神這裡的鼓勵不只是給某一個人,你看看層次,二章四節,先給領袖,他在說所羅巴伯,政治領袖,你要剛強,然後大祭司,約書亞,你當剛強,不只是他們,還有百姓,這裡的百姓,你們都當剛強作供,要剛強壯膽,神鼓勵他們..
神鼓勵他們,你看看,很特別,之後所羅巴伯建造的聖殿的確在外形上,整個外觀上,可以說是不能和所羅門的時間相比,但是到之後,去到大希律的時期,擴建就很不同..
這個殿後來的榮耀的確比之前所羅門的聖殿的榮耀更大,上帝所答應的事是實現了,所以你看看圖,所羅門的聖殿在左邊,我的左邊,希律的聖殿後面是大很多的,榮耀更大..
在這裡,神要鼓勵他們剛強,不要被困難,不要被受敵嚇倒,這個鼓勵很寶貴,沒有什麼彼得上上帝的鼓勵來得更加實際有力.上帝為什麼要這樣去鼓勵他們呢?.
上帝鼓勵他們很重要,上帝記起他和以色列人所納的約,上帝鼓勵他們是因為他記起自己的約,所以二章五節這樣說,他說是照著出埃及的時候,我和你們立約的話,昔日神和以色列人說,他們是屬於上帝的,上帝是他們的神..
神希望百姓遵守他的律例,神要保護他的百姓,神和百姓立約,上帝沒有忘記這個約,直到這個時間,神仍然要提起,我和你們的祖先立約,所以照著來做..

$^{41}$為什麼可以剛強壯膽?上帝這樣鼓勵我們,為什麼實際可以做到剛強壯膽?是因為神的同在,是因為神的同在,這個很重要..
所以當神在他們當中的時候,他們什麼都不需要害怕,二章五節,我的靈要住在你們中間,即神和他們一起,你們不必懼怕..
我想說,神今天也和我們說話,不只是對昔日的以色列人說話,當我們要起來,願意回應上帝重建洪恩家的時候,神要賜福給我們,神要賜福..
神也要我們忘記背後努力面前,不要讓背後成為我們的包袱,望向前面,過去怎樣失敗,過去怎樣好,你記住已經過去了,但前面卻仍然可以努力..
上帝有一個很大的祝福,上帝的榮耀會臨在,這個榮耀要比過去更加大,當然是有前提的,這個重建不是照我們自己的心意..
你留意到領袖的團隊不只是有政治領袖,也有祭司,有宗教領袖,祭司代表百姓向上帝祈求,祭司也代表神向百姓說話,其實神要確保百姓聽到祂的聲音..
上帝要確保,當日確保聖殿按祂的心意而被重建,今天神也要確保祂的家,要按祂的心意被重建,在裡面得勝的關鍵是神的同在,神的靈與我們同在..
今天我們洪恩家看的不是看有多少年輕人,我們看經濟能力有多麼雄厚,不是看很多實際現在有什麼沒有什麼,而是上帝是不是與我們同在,上帝與我們同在就有盼望,不是按照我們的心意重建,而是按照神的心意..
在這裡,神自己親自說,祂與我們同在之餘,上帝親自工作,祂會震動萬國,萬國的增寶都必運來.上帝告訴我們,祂統管天地萬有,祂掌控王權,祂要如何去動,去改變世界,祂完全有這樣的能力,祂也掌管所有的資源..
在這裡,不需要擔心我們缺乏,神會賜給我們,然後是上帝的榮耀臨在,以至到昔日的聖殿的榮耀是很大..
以這七節來說,就使這殿充滿榮耀,是神自己讓榮耀臨在,今天神的家有榮耀,肯定是上帝的工作,上帝讓他的榮耀臨在我們當中,神繼續說,銀子是我的,金子是我的,所有的錢銀是上帝的,物質是上帝的..
現在不是在說我們要靠誰,要做什麼,而是專心靠上帝,上帝自然會把這些東西給我們,以這九節來說,這後來的殿的榮耀比大過先前的榮耀,不是建基於重建的團隊,不是當時的以色列民,猶太人有多厲害,多厲害,不是他們自己本身多有錢,是上帝的工作..
上帝自己賜下資源,賜下物質,上帝與我們同在,與當日的人同在,所以,真的是在說聖殿的榮耀更勝從前,今天神的家一樣,突然之間我們不要去看,看自己很喪膽,很害怕,總是覺得很多缺乏,怕很多不足夠,很軟弱,我們看著上帝,上帝會供應給我們,上帝會工作..
看著他,這個神的鼓勵是非常非常寶貴..
接下來有一個很大的提醒,是一段很特別的經文,突然之間插進去,你覺得很奇怪,說說一下鼓勵,說說一下神的很大的祝福,然後忽然間他就說了,二章十一節他說一段經文是說什麼?.
上帝很想當時的人去請教祭司,因為祭司是很熟悉律法的,很想下蓋先知去請教祭司,去問律法..
這裡一大段去到二章十一節,十一到十四節是說潔淨與不潔淨,怎樣去看?其實這一段是說什麼呢?是神對我們生命的要求,神希望我們聖潔..
首先我們看,如果有一件物件是很聖潔的,例如那些要獻祭的玉,獻祭的玉視為一些聖玉,是吧?很神聖的,很聖潔的這些物件..
如果放在衣服邊,我拿衣服去繞著這些玉,這些玉是聖潔的,會不會因為放在我的衣服,就令我的衣服聖潔了?於是我那件衣服又碰到其他食物,又會令那些食物聖潔呢?.
這裡的回答是不會的,原來聖潔是不可以這樣傳遞的,不可以透過物件傳遞,是吧?譬如我拿了一個物件是聖杯,或者是什麼,很聖潔的,我碰了它,我就自然聖潔了..
那物件本身是聖潔的,但如果我真的說我這個人聖潔,是要自己追求聖潔的,杯不會令我聖潔,正如這些玉不會令衣服聖潔,衣服又不會影響其他東西成為聖潔,是人本身要成為聖潔,這是其中一樣..
另一件事反過來,不好的東西就會影響,不好的東西怎麼影響呢?如果我碰到死屍,如果在舊約裡面大家看到,碰到屍體會視為不潔,碰到動物的屍體也視為不潔,更何況是碰到人的屍體更加不潔,但如果我碰到這些,我就被影響了,這些不潔的物體會影響我成為不潔,是會影響的,是吧?.
我碰到不潔的東西,我碰到的東西又會令那些東西不潔,這樣就反而負面的東西會影響,所以這個層面很有趣,他穿插在這裡說這一段是在說什麼?.
二章十四節,哈戈這樣說:「在我面前,這民如此,這國也是如此,他們手裡的各樣工作都是如此,他們在哪裡所獻的都不潔淨.」.
重點是我們自己要追求聖潔,所做的事情才能夠蒙神契約..
今天不只是重建,重建我們一起起禮,重建神的家,有什麼就拿什麼出來,我們的生命是不是聖潔的呢?.
其實今天在教會裡每一個事奉都需要有聖潔的生命,今天大家一回來,下面大堂就會看到師事招待,師事招待需要成為聖潔..
接著上來我們看到很多人,師班年過,師班需要成為聖潔,領唱是聖潔的,主席需要成為聖潔的,講完自己要聖潔的..
你說我是弟兄姊妹,坐在這裡崇拜而已,事奉人員就是,對不起,都是.其實大家坐在台下和台上事奉人員一樣,一起向上帝獻上咀唇和心靈的祭,上帝一樣在等我們,所以大家一起都需要成為聖潔..
到我們重建神的家的時候事奉更加多,教主學的,有兒童侍工的,有青少年侍工的,有IT的,影音的,有很多很多,可能我遺漏了,有探訪,有關顧的,有傳福音的,這麼多的部分,每個部分都需要成為聖潔..
上帝的要求從來不只是說某一個人,某一件事,而是整個群體是否聖潔呢?這樣去重建神的家,才可以將聖潔獻在上帝的手中..
所以弟兄姊妹在一個很重要的時候,鼓勵大家再一次認真面對自己生命.有時候我們都很輕易說,我們真的要認罪悔改,追求一個聖潔生活,好像口號一樣,太容易說了..
這就是為了什麼認罪呢?我在什麼地方得罪上帝?我的心思意念是否真的全屬上帝掌管?裡面真的在想上帝所喜悅的事情和人?還是我裡面有些意念都是上帝不喜悅的?.
或者我的態度也是,我的態度是上帝不喜悅的?我的嘴巴是做什麼的呢?除了唱詩歌,讚美上帝,敬拜神,一起讀經,學習,禱告,那把嘴巴有沒有獻給了是非呢?彼此傷害呢?.
我們今年很想挑戰弟兄姊妹,我也鼓勵自己認真對付自己的罪.上帝要的是聖潔的生命,做出來的工作上帝才會喜悅,否則所獻上的上帝必會視為不潔..
重建聖殿有一個很重要的意義,是重建我們敬拜的生活,進入這個軌道裡面.所以重建神的家,也要重建我們整個崇拜,重建整個敬拜的生命,不只是星期日..
但是重建的過程,需要我們努力追求成為聖潔,靠著聖靈的幫助,求神大大恩待我們..
我們不想每次坐下來只說自己想說的,覺得那些人很有問題,那些人很有問題之前才說自己,我很有問題,我要對付我自己..
我們要對付好,要拔掉眼中的良目,才去看,其實對方是怎樣怎樣,交給神,對付自己就好了..
今年好好對付自己,過一個聖潔的生活,盼望我們洪恩家所獻上的,是上帝喜悅的..

$^{81}$在這個之後,上帝又要提醒我們,不是只是提醒我們要聖潔,之後聖潔的生命才有聖潔的事奉,以過去為尷尬..
在這裡他在說,上一章第一章的時候就說起,以色列人回來生活很困苦,種又很少收成,吃又吃不飽,喝又是喝很乾,沒什麼滋潤,喝都不及,收入又很快沒有了,好像袋子穿了洞,生活極度極度困苦..
這是因為他們生命的優先次序調轉了,將自己放在一個很高的位置,將神的殿放到很低..
今天我們一樣,如果我們優先次序調轉了,不是先求神的國和神的義,是先求自己的享樂,享受..
那你要想一下,上帝的祝福是不是還會這樣淋到我們的生命裡面呢?值得我們想一下..
這裡就以過去為尷尬,所以他再問,你要想一下,從今天起,還沒有一塊石頭放在石頭上,那時候還沒有建的時候,你想一下我們是什麼光景?.
去到20斗的谷堆裡面,你想找多一點都不行,只有10斗,就是出產不夠.有人來到酒池,想拿50桶的酒,只能拿20桶,根本出產不夠,很多缺乏..
那時候神說,為什麼以色列會有這個缺乏?是因為上帝自己出手,攻擊他們,我以焚風,煤爛,冰雹攻擊你們,和你們手上的各樣工作,你們仍然不歸向我..
人沒有反省自己的生命,遇到這麼多的災難,只是覺得我遇到災難,沒有想過上帝要藉著災難管教我們..
所以經過管教之後,仍然不歸向上帝.弟兄姊妹,我想提醒大家,我們真的要在神面前上上,省察,要明白,上帝是不是要我們的生命回轉?.
上帝是不是要等著我們悔改?如果這是神的心意,就要悔改,不要一直等,神一直在等著我們..
接著還要說,二章十九節,倉裡面還有谷種嗎?答案是否定的,因為那時候上帝仍然在管教他的百姓,葡萄樹,無花果樹,石榴樹,橄欖樹,雖沒有結果,一直在說上帝的管教..
我們要以過去為一個尷尬,今天不要再重蹈覆轍,不要再走過去的路,今天當我重新起來,昔日當以色列人重新起來,建造神殿的時候,上帝就說這一句,從今天起,我必馳服..
上帝多愛人,給人機會改變,上帝可以不馳服,不是必然要馳服,但是上帝很愛我們,一定會馳服,今天他同樣在等著我們,我相信上帝很想馳服給我們,從今天起,上帝要馳服洪人家,要馳服我們,但這個過程需要我們的生命有轉變..
到了最後的時候,很重要,我們看回所羅巴伯的家庭圖,約瓦根的後人就是撒拉帖,撒拉帖生出了所羅巴伯,所以他是王族的後人,到了經文最後的時候,神要鼓勵他..
神要鼓勵他,他不簡單,他要帶領百姓重建聖殿,他要作為他們的政治領袖,面對很大的壓力,在這個過程,神要讓他看到他掌管萬國,他的權柄超過萬國,勝過萬國,也勝過仇敵..
所羅巴伯要將眼目放在上帝那裡,他的權柄是何等大呢?在這裡,一個很有趣的最後的總結是什麼呢?他怎樣稱呼所羅巴伯呢?.
萬君子耀號說,撒拉帖的兒子,我僕人,我僕人所羅巴伯..
哇,所羅巴伯很蒙福,上帝要大大用他,上帝要鼓勵他,我勝過萬國,很多的權柄,我勝過你的仇敵.其實根本上,你看著上帝,什麼都不需要害怕,上帝一定會為他開路,上帝可以幫助他勝過仇敵..
在這裡,但是神的稱呼就是撒拉帖的兒子,我僕人,所羅巴伯.沒錯,他是王族的後人,撒拉帖的意思就是王族,但是他也是神的僕人,所羅巴伯要謙卑在神的面前..
到哪一天,我必以你為人,因我揀選了你.上帝揀選他很榮耀,但是他也只是神的僕人,上帝才是真正的王,他要謙卑在神的面前.如果所羅巴伯要謙卑在神的面前,整個民族都一樣,謙卑在神的面前..
神是這個民族的王,神是整個重建聖殿的王,是主宰.今天神也是我們重建洪恩家的王,主宰,他帶領我們,讓我們將眼目放在他身上,我們一同禱告..
但是主人,你對我們有很多提醒,求主幫助我們追求聖潔的生活,上上審查我們的生命離開罪惡,求主幫助我們忘記背後努力面前.上上依靠的是上帝你自己,想起你與我們同在,就無所懼怕..
主啊,求你經常提醒我們,讓我們的眼睛不要看著今天我們有些什麼,沒有些什麼,不要只看到現在的光景是怎樣,而是專心看著上帝,知道你是豐富的神,全能的神,而且是滿有權能的上帝,在你的裡面一無所缺..
就讓我們知道,只要洪恩家全然以你為主為王,我們一無所缺..
在我們將來重建洪恩家的時候,滿有能力因念,這裡就成為榮耀上帝的一個燈台,能夠見證你的一個很寶貴的地方.求主帶領我們,奉耶穌基督的名字祈禱,阿們..
\newpage



\section{路加福音 19:1-10}
\label{sec:bMJz4f2t79c}
\textbf{2023.03.19 《我要這樣地委身……》 路加福音 19\_1-10 (和修版) 講員 : 張天和牧師}
\newline
\newline
連結: \href{https://youtube.com/watch?v=bMJz4f2t79c}{\texttt{https://youtube.com/watch?v=bMJz4f2t79c}} ~~~~ 語音日期: 2023-03-22
\newline
\newline
\hyperref[sec:F6OnOYSIQDU]{\small{< < < PREV SERMON < < <}}
~
\hyperref[sec:index_chronic]{\small{[返順時目]}}
~
\hyperref[sec:index_scriptual]{\small{[返順卷目]}}
~
\hyperref[sec:nJ0r0oqtEjI]{\small{> > > NEXT SERMON > > >}}
\newline
\newline
路加福音 19:1-10
\newline
\begin{longtable}{cl}
\hline
\hline
章節 & 經文 (和合本修訂版)\\
\hline
19:1 & \begin{tabularx}{0.7\textwidth}{X} 耶穌進了耶利哥，要從那裡經過。 \end{tabularx} \\ \\ \relax
19:2 & \begin{tabularx}{0.7\textwidth}{X} 有一個人名叫撒該，作稅吏長，是個財主。 \end{tabularx} \\ \\ \relax
19:3 & \begin{tabularx}{0.7\textwidth}{X} 他要看看耶穌是怎樣的人，只因人多，他的身材又矮，所以看不見。 \end{tabularx} \\ \\ \relax
19:4 & \begin{tabularx}{0.7\textwidth}{X} 於是他跑到前頭，爬上桑樹，要看耶穌，因為耶穌要從那裡經過。 \end{tabularx} \\ \\ \relax
19:5 & \begin{tabularx}{0.7\textwidth}{X} 耶穌到了那裡，抬頭一看，對他說：「撒該，快下來！今天我必須住在你家裡。」 \end{tabularx} \\ \\ \relax
19:6 & \begin{tabularx}{0.7\textwidth}{X} 他就急忙下來，歡歡喜喜地接待耶穌。 \end{tabularx} \\ \\ \relax
19:7 & \begin{tabularx}{0.7\textwidth}{X} 眾人看見，都私下議論說：「他竟然到罪人家裡去住宿。」 \end{tabularx} \\ \\ \relax
19:8 & \begin{tabularx}{0.7\textwidth}{X} 撒該站著對主說：「主啊，我把所有的一半給窮人；我若勒索了誰，就還他四倍。」 \end{tabularx} \\ \\ \relax
19:9 & \begin{tabularx}{0.7\textwidth}{X} 耶穌對他說：「今天救恩到了這家，因為他也是亞伯拉罕的子孫。 \end{tabularx} \\ \\ \relax
19:10 & \begin{tabularx}{0.7\textwidth}{X} 人子來是要尋找和拯救失喪的人。」 \end{tabularx} \\ \\
[1ex]
\hline
\hline
\end{longtable}
$^{1}$各位紅恩堂的弟兄姊妹早安.
我相信已經有超過五年的時間.
未有站在這個講台上跟大家講道.
如果你在這教會已經有一段時間.
我相信你還記得我曾經來過這教會講道.
亦都在這教會教過課程.
其實我都借過這教會的地方.
幫一些弟兄姊妹做婚前輔導.
所以對這教會的地方相當熟悉.
在3月3日.
老姚洪牧師Gary.
我一向都叫他做Gary.
Gary給了我一個短訊.
提一提我今天要在大家當中分享訊息.
以及給了我一個主題.
不知道大家有沒有留意這個主題.
在教會門口.
我留下了一些講源在這個月會講道.
在我的相片下面寫著「毀身」.
即是說今天的主題是關於毀身.
所以Gary特別強調.
你都向著這個方向來跟弟兄姊妹分享.
我相信弟兄姊妹當你回到教會一段時間後.
都發現這個毀身的題目是一個老生常談.
我們經常都會說「向神毀身」.
「向教會毀身」.
究竟我們作為基督徒.
怎樣才是毀身呢?.
或者我們怎樣去毀身呢?.
這個更加重要.
不只是我們說我們要毀身給神.
而是我們究竟是怎樣去向神毀身呢?.
今天我選用了一段大家很熟悉的經文.
一般稱為「薩戈遇見耶穌」的經文.
薩戈這個人甚至有人認為他是一個罪人.
在經文裡面看到.
當其他人看到這個人的時候.
他都沒有叫這個人的名字.
索性說他是一個罪人.
這個罪人甚至很多人會覺得.

$^{41}$他是一個不應該得到救恩的人.
但是竟然在這個人的身上.
我們看到他用了三個方式.
來向耶穌毀身.
向我們的神毀身.
我們今天就透過這段熟悉的經文.
是路加福音第十九章第一到第十節.
去學習一下今天的題目.
我們就這樣去毀身.
從薩戈的身上看看.
他怎樣用三個方式來向神毀身.
在這段經文當我們看路加記載這件事的時候.
究竟路加想他們將焦點放在什麼事情上呢?.
其實當我們讀聖經的時候.
我們要留意一下.
在這段經文裡面.
有沒有字眼是重覆重覆地出現的.
大家如果拿著程序表.
可能會容易處理一點.
因為程序表是整段經文印出來的.
如果我們看著螢光幕.
因為一節一節.
就可能較為難去追著這段經文.
怎樣去看.
大家留意到.
這段經文裡面有些字眼是重覆出現的.
譬如說第三節提到.
他要看一看耶穌.
這個漢字.
大家留意到.
然後第五節提到.
耶穌到了那裡.
抬頭一看.
就要看一看薩戈這個人.
然後第七節提到.
眾人看見.
看見耶穌說要去薩戈的家.
這裡很明顯.
作者把幾個漢字.
甚至是不同人的漢字.

$^{81}$第一個是薩戈要看耶穌.
第二個是耶穌要看薩戈.
然後第三個是眾人要看薩戈和耶穌在搞什麼.
這是路加記載這件事的時候.
一個很清楚的主題.
一個焦點.
要我們看見的.
但是在我們的和修本聖經和學本聖經裡面.
有兩個漢字是沒有譯出來的.
第一次出現.
其實在第二節.
在第二節開始的時候.
英文聖經會譯作.
Behold 請看.
有一個人叫做薩戈.
在英文聖經.
有些英文聖經也有譯出來.
但有些英文聖經也有譯出來.
這裡有一個漢字.
而另外有一個漢字.
就是到了第八節.
當薩戈站起來.
他說主啊.
請看我把所有的一半給窮人.
當然這個漢字和我們剛才所說的.
薩戈看耶穌.
耶穌看薩戈.
和眾人看耶穌和薩戈的漢字.
是兩個不同的漢字.
但是都是用漢字來講整段經文.
究竟我們讀聖經的人.
去看路加記載這件事.
有什麼焦點值得我們留意.
透過這三個漢字.
薩戈看耶穌.
耶穌看薩戈.
和眾人看耶穌和薩戈.
這三個漢字.
我們就看到薩戈對耶穌的毀身.
他用了三個方式.

$^{121}$來向耶穌毀身.
我們就進入這段經文裡面.
一段很熟悉的經文.
這段經文在開始的時候.
引出了究竟在什麼地方發生這件事.
還有這件事的主角在哪裡.
所以第一節和第二節就這樣說.
耶穌進了耶利哥.
要從那裡經過.
有一個人名叫薩戈.
作祟利長.
是個財主.
作者用了兩節經文.
很簡單帶出了.
這個地方是耶利哥.
這個地方.
當然大家都知道.
是耶利哥城.
另外也提到主角的名字.
叫做薩戈.
也提到薩戈.
他本身是做什麼職業的.
是一個祟利長.
同時間也提到薩戈.
他的身份應該也是很有錢的.
是一個財主.
短短的兩節經文.
是交代了整個故事的主角.
究竟是一個什麼人.
值得留意的.
這裡提到是耶利哥.
耶穌去到耶利哥的地方.
經過耶利哥.
耶利哥這個地方.
我相信弟兄姊妹都不會陌生.
是一個距離耶路撒冷.
不是一個很遠的地方.
耶路撒冷在山上.
而耶利哥在山下.
中間經過一條山路.

$^{161}$所以大家還記得.
好撒馬林的比喻.
也是在一個這樣的情景下發生.
那些祭司從山上的耶路撒冷宮職員.
準備去耶利哥.
這是祭司和事奉人員.
利未人聚居的地方.
然後在那裡.
他們中間見到有一個被山賊打傷了的人.
所以這個也是耶利哥的地方.
耶利哥是一個很繁盛.
是一個交通的樞紐.
尤其是要經過耶利哥.
再去到耶路撒冷的地方.
因此羅馬政府.
就在耶利哥城.
建立了一個隧關.
一個納粹的地方.
而當時羅馬政府.
他們納粹的方式.
為了更加容易收到稅款.
所以找回猶太人.
找回猶太人來幫猶太人收稅.
但是當時收稅的制度.
一般我們稱為自負盈虧.
即是說羅馬政府規定.
在這個隧關要收到多少的稅款.
然後你們總之.
到時候就將相關的稅款.
交給羅馬政府.
至於你怎樣收.
你收到多少.
有多少可以留下來.
你們可以放進口袋.
羅馬政府是不需要理會的.
因此就造成了一個做法.
就是無論稅利.
或者稅利長.
他們會將稅款可能提高一點.
因為如果提得越高的時候.

$^{201}$他的利潤就會越高.
而稅利長管住一些稅利的時候.
他的利潤相對來說.
也會越高.
再加上當時猶太人.
他們需要納聖殿稅.
因此對於這些猶太人來說.
一方面要納聖殿稅.
另一方面又要納地方的.
羅馬政府的稅.
並且稅利再提高了.
因此造成了一種張力.
他們對稅利.
很不喜歡稅利.
甚至會認為.
稅利是在幫外邦人來做事.
所以他們會看稅利.
是一些罪人.
再想想稅利長就更加重要了.
所以先讓大家了解大概的背景.
然後我就進入這段經文裡的三個漢字.
第一個漢大家留意到.
聖經這樣說.
他要看一看耶穌是怎樣的人.
只因人多.
他的身材又矮.
所以漢不見.
又出現這個漢字.
然後再留意第四節.
他說於是他跑到前頭.
爬上桑樹.
再出現一個要看耶穌.
因為耶穌要從那裡經過.
當大家看第三第四節的時候.
很明顯.
撒該當時有一個很清楚的目的.
就是要看一看耶穌是怎樣的人.
我們自然會問.
為什麼撒該這麼討人厭.
周圍的人看到這個稅利長都不喜歡.

$^{241}$但是竟然跑出來大庭廣眾.
跑出來一個很多人的地方.
就是要看一看耶穌是怎樣的人.
我們自然會問.
為什麼撒該要這樣做.
這麼犯險被其他人指指點點罵他.
結果真的有人說他是罪人.
為什麼撒該要這樣做呢.
如果你看到第十八章的時候.
這段經文的上文.
大家會留意到有些背景值得留意.
當耶穌經過耶利哥的時候.
發生了另一件事.
有一個黑眼的人.
他們聽到有很多人聽到的時候.
就問究竟發生了什麼事.
有人就說是拿撒勒人耶穌經過.
這是第十八章所記載的一件事.
然後這個人去求耶穌醫治他.
大家留意.
結果耶穌就問他.
我可以為你做什麼.
然後那個黑眼人就說.
我要能看見.
然後耶穌就醫好了他.
大家可以想像到.
同樣在耶利哥發生的一件事.
一件這麼重要的事情.
就是一個盲眼的人.
本來那些人是要打發他走.
是攔住他不想讓他接近耶穌.
但是這個人卻去到耶穌面前.
他說大衛的子孫耶穌可憐我.
大家看到人家介紹的時候是拿撒勒人耶穌.
但是這個盲眼的人.
他竟然說是大衛的子孫耶穌.
然後耶穌就醫好了這個人.
大家可以想像到.
為什麼撒該會看見耶穌是怎樣的人.
因為他可能都會聽到.

$^{281}$這個人這個耶穌.
竟然一個被人拒絕攔住他.
不讓他接近耶穌.
但是耶穌主動跟他說.
我可以為你做些什麼.
然後這個人就說.
我要能看見耶穌醫治了他.
我相信這是一個很重要的原因.
使撒該要看見耶穌是怎樣的人.
但是經文的發展就告訴我們.
因為撒該的個子很矮小.
當然聖經沒有形容到他多麼矮.
但是矮是一個相對的用詞.
譬如說一個1.8米的人.
如果站在一個籃球手面前.
你就是矮的人.
所以我們聖經沒有說到.
究竟撒該有多矮.
只是形容他的身材矮.
以至他可能被人遮擋了他的視線.
他看不到耶穌.
大家要留意.
他來是要看看耶穌是誰.
但是現在看不到.
就是說他來的目的是達不到.
但是可能有人這樣說.
矮子多計.
所以他立刻想想.
如何可以繼續看到耶穌呢.
他就跑到一棵桑樹.
然後爬上去.
然後聖經告訴我們.
他原來爬上去的目的.
都是要看一看耶穌.
因為他知道耶穌會從那個地方經過.
大家留意到.
這個撒該.
這個當時是其他人會稱呼他做罪人的瑞利長.
竟然他這麼堅心.
這麼盡力的.

$^{321}$來想見一下耶穌.
想親近耶穌.
但是姐妹們.
我不知道這件事給我們有什麼提醒.
他作為一個當時周圍的人.
周圍的人都不是這麼喜歡的人.
他說我要去看看耶穌是誰.
現在看不到.
我不知道你會怎麼選擇.
都不要繼續露面.
都不要繼續站在一個公眾的地方.
被人指指點點.
既然個子小又被人遮住了.
看不到.
算了不如回家吧.
可能我們都會這樣.
我不知道你有沒有試過準備出來.
回教會聚會的時候.
遇到一個難處.
今天還是在網上看算了.
我們看到撒該.
他雖然是冒著一些.
可能被人指指點點.
甚至被人罵的情況.
但是我們看到撒該.
他很清楚的.
無論他遇到什麼困難.
他仍然都想一個辦法.
說要去看看耶穌是怎樣的.
各位.
這個我會說是撒該對神的第一種毀身的方式.
就是一種堅定盡心的去親近神.
堅定盡心的去親近神.
似乎有些事情已經在難阻他.
但是他不怕這些難阻.
然後繼續的要去親近這位耶穌.
然後去見下這位主.
我們經過幾年的疫情之後.
有很多人都有一句說話.
回不了頭.

$^{361}$我想大家都明白.
是指我們在網上聚會.
回不了頭.
我知道剛才主席一開始提到.
網上的弟兄姊妹.
我相信在網上的弟兄姊妹.
仍然有在看同步的直播.
我說的回不了頭.
是這種回不了頭.
很多的教會.
發現原來當他們要停止直播的時候.
真的有很多弟兄姊妹.
是不回來的.
很多弟兄姊妹.
這麼多直播.
我看隔離教會的直播.
當然我們不會說.
你看隔離教會的直播.
這個不是不親近.
你也有親近主的.
但是我又記起有一次.
這是幾年前的事.
我想是18還是17年.
大家不知道有沒有印象.
在那一年裡面.
很特別的.
有三個主日.
都是打風的.
大家記不記得這件事.
十號風球,八號風球.
有三個主日.
都是排著主日打風的.
我不知道你記不記得.
在那次打風的時候.
回不了教會.
你的心情是怎樣.
或者你有沒有留意到.
其他弟兄姊妹的心情.
我當時很失落.
因為我要負責講道.

$^{401}$但是那天我沒有機會講道.
但是竟然我在教會群組裡面.
有人說.
我去找位置.
大家一起去喝茶.
好像很興奮很開心.
真是好了.
今天不用回教會.
弟兄姊妹.
當我們說我們向神毀身的時候.
我們會不會好像殺雞一樣.
堅定盡心的親近主.
縱然是遇到一些難阻.
但是我們仍然很想.
神很想去親近你.
因為你是我要毀身的對象.
這是殺雞第一個毀身的方式.
我們進入第二個看.
第二個看是第五和第六節.
第五節是這樣說.
他說耶穌到了哪裡.
就是到了殺雞爬了那棵樹上面.
然後大家留意.
聖經記載是抬頭一看.
大家可以想像到.
是耶穌刻意去找殺雞的.
所以大家看到這段經文.
最後的時候.
是說是耶穌尋找人.
所以就是這樣.
這很明顯是耶穌刻意.
大家走路的時候.
你不會這樣鱷高頭望的.
你會撞柱的.
我昨天就是一直不知道在想什麼.
去趕巴士就真的撞到柱子.
不過是矮的柱子.
所以很明顯是說.
耶穌抬頭一看.
是說他要找一找殺雞的人.

$^{441}$然後接著對他說.
殺雞快下來.
今天我必須住在你家裡.
然後第六節就是殺雞的反應.
殺雞就是他就急忙下來.
歡歡喜喜地接待耶穌.
這是第二個看.
第一個看我們看到.
是殺雞很想去看看耶穌.
然後第二個看我們看到.
就是耶穌是抬頭一看.
他是刻意的來找一找殺雞.
為什麼我們會刻意呢.
因為當他看到殺雞的時候.
他講得出他的名字.
他說殺雞快下來.
不知道大家有沒有想過.
為什麼耶穌會知道殺雞這個名字呢.
是不是殺雞臭名遠播.
所以他知道殺雞這個名字.
當然經文本身我們很難猜測到.
為什麼耶穌會知道殺雞的名字.
也不是最重要.
他為什麼會知道殺雞的名字.
但是當聖經出現.
一個人直接去呼喊對方的名字的時候.
是在表示他對這個人的認識和關懷.
雖然我再強調.
我們不知道耶穌為什麼會知道殺雞的名字.
但是經文的確記載.
耶穌是直接呼喊殺雞的名字.
殺雞快下來.
這是在表示他對殺雞的認識.
和他對殺雞的關懷.
更加我們看到.
就是跟著一句話.
今天我必須住在你家裡.
大家知道耶穌這個邀請.
或者跟殺雞說的話.
是引起其他人的非議的.

$^{481}$耶穌刻意地說要我住在你家裡.
弟兄姊妹.
耶穌抬頭一看.
他要表達對殺雞的那種關懷.
那種認識.
並且他說要住在殺雞的家裡.
然後殺雞的反應就是急忙下來.
歡歡喜喜地接待耶穌.
我會說這是殺雞的第二種方式.
來對神的毀身.
就是讓接待耶穌進入他的家裡.
所以進入他的家裡究竟是什麼意思呢?.
當然在耶穌這個時代.
進入他的家裡.
是表示他跟這個人的關係很親密.
所以大家看到在聖經記載.
耶穌通常不是進入家裡停留.
例如在馬大馬尼亞的家裡.
他經常會去.
因為是好朋友.
所以在《約翰福音》第十一章.
當馬大馬尼亞派僕人去找耶穌的時候.
他跟耶穌說你所愛的人病了.
耶穌已經知道是誰了.
所以表示是很熟悉的.
進入他家裡的時候很熟悉.
而耶穌這裡說我要去你家裡.
大家可以想一想.
耶穌去撒該家裡.
是代表著什麼呢?.
或者我這樣說.
大家可能會容易明白.
如果突然之間.
你有去過外國.
或者在外國生活過.
去別人家裡住.
其實是很普通的.
例如我要去外國領會的時候.
多數都是.
張牧師你來我家裡住吧.

$^{521}$你不用住酒店的.
通常都是這樣.
去接待你在家裡住.
但是在香港地方.
我想是少很多.
張牧師你來我家裡吃頓飯吧.
已經是很不錯了.
但是吃頓飯之前.
你也要搞一大輪的.
要收拾一下房子.
例如我每年.
學生都會來我家裡拜年.
我每次拜年之前.
都要用兩天時間收拾房子.
收拾了房子.
然後才可以接待學生.
來我家裡拜年.
所以大家明白.
原來耶穌去那個人的家.
不是只吃頓飯.
而是現在說要住在他家裡.
住在家裡是在說什麼呢.
我以下說一個故事.
大家會更加容易去明白.
曾經有一個年輕人.
他年輕的時候已經工作很順利.
所以他買了一間他自己的夢想之屋.
Dream House.
這間屋子上面有兩層.
樓上有五間房.
樓下五間房不簡單.
一間屋子有十間房.
他就說.
真是開心.
在我家裡有十間房.
而我佈置得最好的房間.
就在二樓.
然後聽到敲門聲.
他就說.
當然是我進了新居.

$^{561}$有人來探訪我了.
他一開門的時候.
發現原來耶穌來探訪他.
他非常高興.
耶穌歡迎你來我新的家.
隨便你留在這裡住.
我留下一間很好的房間.
你就住在這間房裡.
但是他卻看不到.
原來耶穌的臉色不是很開心.
然後耶穌就停留在這裡.
然後又聽到敲門聲.
聽到敲門聲的時候.
這個年輕人就想著.
又有人來探訪我.
新居入房這麼多人來探訪.
然後他一開門的時候.
突然間.
原來魔鬼馬上衝進來.
衝進來的時候.
他想推也推不走.
說不行我不歡迎你.
這個地方不讓你進來.
但是魔鬼就進來了.
所以他用盡九牛二虎之力.
很困難才勉強將魔鬼推出門外.
但是當時他想到.
耶穌住在這裡.
為什麼耶穌不出來幫我呢.
他就問耶穌.
耶穌剛才說魔鬼來過.
為什麼你不出來幫我.
耶穌就跟他說.
年輕人.
我是這間屋的客人.
你是這間屋的主人.
所以事情應該由你去應付.
這個年輕人就醒目了.
啊對.
他就跟耶穌說.

$^{601}$耶穌我願意將這間屋一半的房間給了你.
你在這裡繼續安心住.
你住哪間房就住哪間房.
然後又聽到敲門聲.
這個年輕人醒目了.
可能魔鬼又來了.
所以他開了一條窄窄的門.
誰知道魔鬼就是這麼狡猾.
這麼窄的門.
他都可以鎖進來.
同樣的畫面.
同樣的情況再一次出現.
就是那個年輕人很辛苦.
才將魔鬼推了出去.
他又再想.
他說耶穌在我房子裡面.
我已經將半間房子給了他.
但為什麼他都不出來幫忙呢.
他又問耶穌.
耶穌同樣給同樣的答案.
他說我是這間屋的客人.
你是這間屋的主人.
所以你要去處理這件事.
這個年輕人在這個時候.
他醒覺了.
他說我要將這間屋.
奉獻給神.
讓神住在這間屋裡面.
我只是住其中一間房間.
然後又聽到敲門聲.
聽到敲門聲的時候.
這個年輕人立刻想起身.
想去開門.
但他想.
我只不過是這間屋的客人.
耶穌是主人.
所以應該是耶穌處理.
結果耶穌就在樓上.
一直走下來.
他就在後面跟著.

$^{641}$看看發生什麼事.
他在後面看的時候.
他只看到耶穌開門.
看不到有人走進來.
然後聽到一把聲音.
就這樣說.
對不起.
我找錯地方了.
弟兄姊妹.
我不知道你們有沒有聽過這個故事.
當你聽到這個故事的時候.
你會明白.
當耶穌說.
我今天晚上.
我今晚.
我必須住在你的屋裡面.
是什麼意思.
其實耶穌這個邀請.
說我今天要住在你的屋裡面.
是代表著.
我要進入你的生命裡面.
每一個層面.
因為一個人來家裡吃飯.
我原來把東西放進房裡.
但是我們不讓他進房裡.
他看不到那些東西.
但是如果那個人住在我家裡的時候.
我發現原來我整間屋.
好像是透明的.
是可以讓他看見的.
所以當耶穌說.
今天我必須住在你家裡.
是在說耶穌要進入.
撒該整個生命的.
每一個層面裡面.
我們留意一下撒該的反應是怎樣.
撒該的反應就是急忙下來.
接著有一句話.
大家留意是.
歡歡喜喜的接待耶穌.

$^{681}$弟兄姊妹.
這是撒該向神毀身的第二個方式.
是歡歡喜喜的.
將自己的生命交給神.
是歡歡喜喜的.
將自己生命的每一個層面.
交給神.
這是對神毀身的第二個方式.
然後我們進入第三個罕見.
第三個罕見.
是第七和第八節.
或者可以說是整個記載這件事的高潮.
第七節說.
眾人看見都私下議論說.
他竟然到罪人家裡去住宿.
這是那些人看見耶穌和撒該.
然後就有一個議論.
竟然這個人去罪人家裡住宿.
所以大家留意到.
撒該爭奪的地步.
不叫他矮子.
不叫他有錢人.
甚至不叫他一個稅吏長.
只叫他是一個罪人.
但是竟然是耶穌進入罪人家裡住.
這裡有什麼提醒呢.
我們當然要看第八節.
第八節.
撒該站著對主說.
主啊 我把所有的一半給窮人.
我若勒索了誰.
就還他四倍.
這是撒該當時的反應.
這段經文在路加福音出現的時候特別有意思.
因為在路加福音裡面.
如果你查看路加福音.
你會發現路加福音裡面.
很常記載猶太人.
包括民事化理財人.
經常對那些他們認為是罪人.

$^{721}$很鄙視和歧視那些罪人.
包括這段經文.
都是眾人說這是罪人.
他住在罪人家裡.
所以當每一次出現這樣的情況.
譬如說在十五章記載浪子的比喻.
為什麼會出現浪子的比喻呢.
就是耶穌看到那些人.
在這時候就用了三個失的比喻來提醒.
原來是失而復得的.
所以在整個路加福音裡面.
都鋪排了很多次.
那些化理財人和民事.
都是用這個方式來看待罪人.
過往都是耶穌會出來說的.
但是在這段經文有一個特別的地方.
耶穌反而沒有站出來說.
反而是撒該自己站出來說.
在撒該站出來說的時候.
他說了兩件事情.
第一件事情就是說.
將所有的一半給窮人.
另外一件事情就是.
我若勒索了誰.
就還他四倍.
為什麼這麼特別是四倍呢.
為什麼又不是五倍.
又不是六倍.
又不是兩倍.
為什麼會四倍呢.
這個做法其實是有律法的根據的.
如果我們查考的時候.
我們發現在利未記第六章第五節那裡規定.
賠償是只需物歸原主之後.
再加百分之二十.
就是加五分之一.
這個是律法的一個要求.
而最高的要求就是給四倍.
最高的要求就是四倍.
譬如說我拿了別人的東西.

$^{761}$然後我賠償給他的時候.
我只要賠償百分之二十.
就是五分之一.
已經是滿足了律法的要求.
當然律法的要求是可以賠償到四倍.
但是這個賠償四倍.
是主要用在牛和羊的身上.
用在物件上.
不是金錢上.
但是在這件事情上.
撒該是將那些他勒索過人的.
很明顯是金錢.
然後四倍的償還.
然後另外一件事就是.
將他一半的財產分給窮人.
大家看路加福音書.
是真的傳過來很有趣的.
其實為什麼這裡分給窮人一半.
有什麼特別呢.
如果你追路加福音書前面.
有一個叫少年的官.
大家還記得嗎.
有一個少年的官然後找耶穌.
然後我做什麼可以得到永恆.
然後耶穌跟他說.
你變賣你所有的去周濟窮人.
大家還記得那件事嗎.
那個年輕人那個少年官是怎樣的.
悠悠愁愁的走了.
因為他擁有很多.
如果這樣的時候.
你做一個對比就知道.
撒該的行動是不簡單的.
因為那個是耶穌要求他變賣所有.
然後去周濟窮人.
但是他不肯.
他悠悠愁愁的走了.
但是這次耶穌沒有要求撒該.
你要做什麼就做什麼.
是撒該自己站出來說.

$^{801}$我勒索了誰.
我四倍還給他.
是超越了律法對他的要求.
然後他將他的財富的一半.
分給窮人.
是耶穌沒有要求的.
究竟這兩個動作.
是代表什麼.
這兩個動作又怎樣去反映.
撒該是用第三種方式.
來對神的毀身.
他第三種方式對神的毀身.
我會說他整個人生的價值觀.
是轉變了.
是被耶穌的福音來改變了.
大家可想而知一個稅利長.
我們剛開始的時候.
他是一個財主.
是一個稅利長.
是一個周圍的人.
不是那麼喜歡他的人.
所以其他猶太人都稱他為罪人.
為什麼仍然做這份工作呢.
我想大家很容易想到.
錢.
因為這份工作可以賺錢.
因為這份工作能夠使他成為一個財主.
如果從這個角度看的時候.
你會看到.
這個撒該應該是一個很喜歡錢的人.
但是一個喜歡錢的人.
他竟然說.
我勒索了誰.
我四倍還給他.
然後我將我財富的一半分給窮人.
是代表著什麼呢.
他的價值觀轉變了.
是代表著這些本來我牢牢守住的錢.
甚至乎被人稱為罪人的錢.
賺回來的錢.

$^{841}$但是在這個時候.
因為我要向神毀身.
因為我遇見了耶穌.
所以我願意付出.
這是一個價值觀的徹底扭轉.
弟兄姊妹.
當我們向神毀身的時候.
原來究竟我們這個人的質地.
我們的價值觀.
有沒有因著信仰的緣故.
我們作出改變.
會不會因著信主.
我們因著主去改變我們的生命.
原來我們的價值觀.
我們對一些事物的看法.
是不同了.
特別是對財富是不同了.
所以這是撒該.
他第三個向神毀身的方式.
今天我們從這三個漢字.
我們看到撒該.
他用了三種方式.
來向耶穌毀身.
這個也是今天和大家分享的題目.
就是我要這樣地毀身.
第一個就是要很堅定盡心的.
來親近我們的主.
無論有什麼難阻.
我們都要親近我們的神.
第二種毀身的方式.
就是願意將他生命的每一個層面.
來擺在神的面前.
就好像撒該那樣.
歡歡喜喜地接待耶穌.
進入他的家裡面.
這是說整個人來奉獻給神.
交託給神.
第三種毀身的方式.
就是因著與主相遇.
福音的改變.

$^{881}$以至我們的價值觀念改變了.
而這樣將自己的生命.
一個這樣被改變了的生命.
去到神的面前.
這就是我們對神毀身的第三個方式.
這段經文很熟悉.
但願一段這麼熟悉的經文.
讓我們從另一個向度.
來思考一下.
我們就是這樣去毀身給我們的神.
我們一起在這裡祈禱.
弟兄姊妹.
我也希望你們在這個時候.
可以安靜一下.
去記起這段經文.
去記起撒該.
去遇見主.
去怎樣來向我們的主毀身.
對我們有什麼提醒.
我們是不是這樣來向神毀身呢.
我想弟兄姊妹.
你們自己在位置上.
做一個很簡單的回應的禱告.
然後我會帶領大家.
做一個結束的祈禱.
親愛的主.
當我們今天來到你面前的時候.
我們讀一個很熟悉的一件事.
就是撒該和主的相遇.
但我們看到從撒該的身上.
他有三個方式來向你毀身.
主求你今天.
用這段熟悉的經文.
用這個熟悉的人物.
來提醒我們.
讓我們學習怎樣來向你毀身.
主我們知道這是我們需要學的功課.
因此我們求主你的靈.
不斷地去感動我們.
提醒我們.

$^{921}$和幫助我們.
我們這樣祈禱.
是奉靠耶穌得勝的名求.
\newpage



\section{加拉太書 3:21-29}
\label{sec:nJ0r0oqtEjI}
\textbf{2023.03.26 《律法之上》 加拉太書3\_21-29 (和修版) 陳韋安牧師}
\newline
\newline
連結: \href{https://youtube.com/watch?v=nJ0r0oqtEjI}{\texttt{https://youtube.com/watch?v=nJ0r0oqtEjI}} ~~~~ 語音日期: 2023-03-27
\newline
\newline
\hyperref[sec:bMJz4f2t79c]{\small{< < < PREV SERMON < < <}}
~
\hyperref[sec:index_chronic]{\small{[返順時目]}}
~
\hyperref[sec:index_scriptual]{\small{[返順卷目]}}
~
\hyperref[sec:7_miM8zb42g]{\small{> > > NEXT SERMON > > >}}
\newline
\newline
加拉太書 3:21-29
\newline
\begin{longtable}{cl}
\hline
\hline
章節 & 經文 (和合本修訂版)\\
\hline
3:21 & \begin{tabularx}{0.7\textwidth}{X} 這樣，律法是與神的應許對立嗎？絕對不是！如果律法的頒佈能使人得生命，義就誠然出於律法了。 \end{tabularx} \\ \\ \relax
3:22 & \begin{tabularx}{0.7\textwidth}{X} 但聖經把萬物都圈在罪裡，為要使因信耶穌基督而來的應許歸給信的人。 \end{tabularx} \\ \\ \relax
3:23 & \begin{tabularx}{0.7\textwidth}{X} 但這「信」還未來以前，我們被看守在律法之下，像被圈住，直到那將來的「信」顯明出來。 \end{tabularx} \\ \\ \relax
3:24 & \begin{tabularx}{0.7\textwidth}{X} 這樣，律法是我們的啟蒙教師，直到基督來了，好使我們因信稱義。 \end{tabularx} \\ \\ \relax
3:25 & \begin{tabularx}{0.7\textwidth}{X} 但這「信」既然來到，我們從此就不在啟蒙教師的手下了。 \end{tabularx} \\ \\ \relax
3:26 & \begin{tabularx}{0.7\textwidth}{X} 其實，你們藉著信，在基督耶穌裡都成為神的兒女。 \end{tabularx} \\ \\ \relax
3:27 & \begin{tabularx}{0.7\textwidth}{X} 你們凡受洗歸入基督的都披戴基督了： \end{tabularx} \\ \\ \relax
3:28 & \begin{tabularx}{0.7\textwidth}{X} 不再分猶太人或希臘人，不再分為奴的自主的，不再分男的女的，因為你們在基督耶穌裡都成為一了。 \end{tabularx} \\ \\ \relax
3:29 & \begin{tabularx}{0.7\textwidth}{X} 既然你們屬於基督，你們就是亞伯拉罕的子孫，是照著應許承受產業的了。 \end{tabularx} \\ \\
[1ex]
\hline
\hline
\end{longtable}
$^{1}$鄧小妹 早安.
很高興來到紅磡堂和大家一起去敬拜.
今天我們會看一段保羅在《迦太書》裡面的講道.
保羅在《迦太書》裡面.
都是講明一些很基本的信仰.
我們回到教會很多年前都認識的一些課題.
恩信 清義 福音 救恩的課題.
今天我們一起看一段經文.
重新思想律法對我們作為基督徒的意義在哪裡.
要講律法的課題之前.
我們都要講恩信 清義.
大家都很熟悉聽過的一個道理.
恩信 清義是保羅在《迦太書》裡面的教導.
在第二章十六節裡面的道理.
後來他都在《羅馬書》裡面.
特意重新很清楚地講明這個道理.
更加重要的是我們今天是基督徒.
我們是新教徒.
五百年前馬丁路德都是重新來強調.
恩信 清義的教導.
你知道對路德來說.
重新去強調信心.
不單單是我們信仰的起點.
更加是我們整個信仰的全部.
一切都是來自於五百年前.
看回五百年前路德的改教的時候.
其實路德是一個天主教會的修士.
他是一個非常認真追求自己.
要求自己很高的人.
他都有點火爆.
對於自己很認真的一個人.
所以在五百年前他在十五條裡面.
他去改教去釘上十五條的時候.
很強調這個道理.
路德其實是慢慢地講出.
他對自己對人一個最徹底的發言.
就是人是對於自己的軟弱和罪.
徹底地去發言自己的問題.
在他還未說得出恩信 清義這麼清楚之前.
其實有一個更加重要的概念.

$^{41}$就是卑微神童.
卑微神童就是說.
人對自己無能為力的那種徹底的醒悟.
面對上帝的要求.
面對上帝的審判的時候.
人是徹底地感到無能為力.
是真真正正地感覺到自己做不到.
所以卑微不是謙卑.
不是純粹一種禮貌.
或者中國人所講的一套表面的禮貌.
卑微是人對自己的罪.
徹底的認知.
我什麼都做不到.
所以唯有上帝救我.
所以人面對自己的問題.
發現自己什麼都做不到.
就學羅德發言了一個很重要的道理.
就是唯有上帝恩典.
所以羅德慢慢發現.
人在一個完全什麼都做不到的情況下.
唯有只能夠相信.
所以相信不是一個我們能夠做到的事情.
可能我都誤會.
我們相信就能夠換取救恩回來.
其實這不是羅德的意思.
信不是我們得救的交換方法.
只要我相信就能夠換取天堂的入場券.
這不是羅德的意思.
信是我們什麼都做不到.
我們什麼都做不到的時候.
只能夠靠耶穌基督的恩典來幫助我們.
所以羅德就強調這個恩信清義.
那個相信不是因為我們能夠信.
所以能夠換取這種的清義的方法.
唯獨信心這個說法其實是指向著唯獨基督.
我們唯有靠著耶穌.
單單靠著他.
信只不過是一種接受耶穌的恩典.
很重要的一種回應.
所以有些人說上帝出了雞出了豉油.

$^{81}$其實不是.
其實上帝出了豉油給我們.
所以這就是信心那種很重要的道理.
如果不是的話.
信就成為一種行為.
我只能夠信是一種行為.
所以這不是羅德的意思.
所以當我們去明白到保羅在《皆太書》裡面所說的恩信清義.
人是什麼都不能夠做到的時候.
我們就能夠明白到底律法的意義在哪裡.
當我們看回今天的經文的時候.
我們發現《皆太書》第三章說.
他說這樣.
我們讀回這個和本版本.
他說.
律法是與上帝的應許反對嗎?.
斷乎不是.
若曾傳一個人.
一個能叫人得生的律法.
亦就成而本於律法了.
但聖經把眾人都圈在罪裡.
使所應許的福.
因信耶穌基督.
歸給那信的人.
但這因信得救的理還未來已先.
我們被看守在律法之下.
直圈到那將來的真道顯明出來.
這樣律法使我們憤懵的師父.
引我們到基督那裡.
使我們因信清義.
因這因信得救的理既然來到.
我們從此就不在師父的手下了.
保羅說.
雖然因信清義在我們得救裡面.
很重要.
我們是無能為力的.
任何的律法行為.
我們行國禮.
都是在得救的時候.
是沒有功效的時候.

$^{121}$保羅就問.
這樣律法是不是沒有意義呢?.
律法是不是不重要呢?.
或者是不是律法甚至是.
違背了福音的道理呢?.
保羅就說什麼?.
他說不是的.
律法不是不重要的.
保羅說.
這樣律法是與神的應許對立嗎?.
絕對不是.
我們就看看.
律法到底今天有什麼意義和反省.
如果我們說.
清義不是靠律法的時候.
那麼律法是什麼呢?.
保羅還有一個很好的比喻.
去說明這個道理.
他說這樣律法使我們憤懵的師父.
引我們到基督那裡.
使我們因信清義.
他就像是鑊本某一種.
叫做是憤懵的師父.
雖然本來就好一點.
律法是我們的啟蒙教師.
其實原文就叫做Pedagogos.
其實就是教育或者教師的道理.
不過其實意思就不是.
解作那些很高級的教授.
或者是那些老師.
所以鑊本都還沒有好好地.
去說得出背後的意思在哪裡.
什麼叫做Pedagogos呢?.
在古代的社會裡面.
其實這個字不是解作那些.
很高級的師父或者是老師.
相反是一些陪伴小朋友讀書的奴隸.
以前在羅馬的社會裡面.
那些陪人讀書的奴隸就叫Pedagogos.
就是負責拿書,寫字.

$^{161}$如果在中國來說.
就叫什麼呢?.
以前那些叫書童.
旁邊就磨墨那些.
幫少爺拿著書架去走走.
那些人.
所以比較好的譯本就是這個.
女人中的一本.
翻譯成叫童年導師.
因為童年導師的翻譯.
不過說律法的意義就在這裡.
律法只不過是小朋友那些的童年導師.
那些書童.
那些小朋友初階班那些的BB那些的東西.
什麼意思呢?.
不過說律法是不是重要?.
律法是重要的.
不過律法都是屬於那些.
很初階小朋友BB那些的東西.
你問這些重要嗎?.
當然重要.
大家那些孫子孫女學東西.
都是要學這些的.
什麼什麼初階班,游水班,初一班那些.
都是從這些基本學起的.
小朋友學游水.
一開始總是要有個水泡仔.
學踩單車總是要裝兩個什麼.
輔助輪來學的.
BB學走路.
以前都流行學走車.
所以水泡是重要的.
那些輔助輪都重要的.
學走車都重要的.
但是等於一個人.
男人老狗都帶著水泡游水.
就很醜了.
還要踩單車.
還要輔助輪就不太好看了.
所以不太好看的醜慣問題.

$^{201}$然後你去到某個階段.
剛好相反.
這些水泡仔這些輔助.
反而會阻礙你學習去一個新階段.
你學游水總有一刻.
你要拋開水泡就要拔走.
學踩單車也是一樣.
你要拆那個輪子才能學踩單車.
這些反而是有礙你跟真正去認識學會的方法.
所以保羅就說.
律法是重要的東西.
但是律法是什麼.
就是這種階段.
一個很初階階段.
你要朝向一個更加重要的道理.
就是領我們到基督那裡去.
這就是意思.
這就是保羅所說的.
律法在我們信仰裡面的道理.
所以我們今天就要問問題.
我們的信仰.
我們回教會這麼久.
我們的信仰是不是仍然都是這些水泡.
是不是都是這些輔助輪.
仍然都是這些比較初階BB的階段.
作為一個牧師.
作為一個在網上有寫文章的人.
間中也收到很多弟兄姊妹的奇難雜症.
很多問題就問我.
通常都是這樣的問題.
問方法來問我的問題.
可不可以的問題.
我讀出來有一堆可不可以的問題.
基督徒可不可以崇拜遲到.
基督徒可不可以星期日不回教會.
基督徒可不可以不回主教學.
基督徒是否真的需要受到奉獻.
基督徒可不可以炒股票.
基督徒可不可以和微信主人結婚.
基督徒可不可以吃飯前不祈禱.

$^{241}$基督徒可不可以不靈修.
基督徒可不可以練瑜伽.
基督徒可不可以參加萬聖節.
基督徒可不可以整容.
基督徒可不可以看鬼片.
基督徒可不可以做心理測驗.
基督徒可不可以看相.
基督徒可不可以賣樂彩.
基督徒可不可以轉教會.
基督徒可不可以看中醫.
基督徒可不可以講白色的謊言.
基督徒可不可以上網崇拜.
基督徒可不可以吃早餐上網崇拜.
基督徒可不可以做面膜.
基督徒可不可以上網崇拜.
基督徒可不可以做面膜上網崇拜.
以上的問題總是有答案的.
我總可以給你一個答案.
為什麼會有這麼多的問題呢.
為什麼我們要問這麼多可不可以的問題呢.
基督徒可不可以吃早餐上網崇拜.
基督徒可不可以剪指甲上網崇拜呢.
基督徒可不可以吃早餐上網崇拜呢.
在這些可不可以的問題之前.
為什麼會有這麼多的問題呢.
很多年前我寫過一篇文章.
關於基督徒應否炒iPhone.
很多年前在iPhone 60年代.
一買一賣有什麼不可以.
神怎樣看投資.
管家比喻.
但當我寫的時候.
為什麼我們會寫這些問題出來.
為什麼我們會問這些問題.
基督徒應否炒iPhone.
我發覺這不是我們基督徒.
應該問的方法.
我們的倫理.
我們不斷發洩地問可不可以.
這樣做可不可以上天堂.

$^{281}$這樣做會不會玩聖經.
這是一個消極被動的問法.
嘗試去尋求一條緊緊合格的底線.
尋求一條不會被上帝擊殺的底線.
然後在底線上不斷徘徊.
只要我仍然可以的話.
我就做了.
這不是我們基督徒.
去問的方式.
這是一個律法的思維.
這是一個訓蒙師父.
是這些水泡.
輔助輪子的問題方法.
我們的倫理.
是一個登山寶訓的倫理.
是一個追求更加好的意義的倫理.
不是不斷去尋求最低的底線.
而是怎樣呢.
我們要不斷去竭力尋求最高的方式.
耶穌說你們要完全.
你們的天賦完全一樣.
不是問可不可以殺人.
而是問一個更加更加高的地位.
所以保祿教導我們.
我們要離開.
這個律法主義的嬰兒.
離開這個糞蒙的師父.
我們要去基督耶穌的層次.
有人問當我離開基督的嬰兒的時候.
去到耶穌的階段.
我們都以為我們恩信稱義.
基本上就是信就行了.
什麼都不用做.
我們基本上是.
基督徒.
比很多義義教徒還懶.
那些耶徵那些.
他們比我們更加勤奮.
因為我們覺得信就可以了.
也不是教行為.

$^{321}$這個當然不是那個說法.
恩信稱義不是教導我們.
做一個懶惰的人.
不是純粹尋求得救就算了.
當然.
也不是去問可不可以.
所以既然我們基督徒.
不是純粹恩信稱義.
做什麼都行.
但又不是.
尋求得不得.
那我怎麼做呢.
這個就是我們今天要去思考的問題.
其實.
耶穌有教的.
耶穌就在登山寶訓裡面.
教我們怎樣去.
去回應這個律法.
耶穌說什麼.
莫想我來要去廢掉.
律法和先知.
我來不是要廢掉.
乃是要成全.
我告訴你們.
天地都廢去.
律法的一點也不能廢去.
也要成全.
我告訴你們.
法理錯人的義.
斷不能進天國.
要求門徒的義.
勝過.
民事罰人出義的話就斷不能進天國.
是一個非常嚴厲.
認真的命令.
我們問什麼叫門徒的義.
這個義.
怎樣能勝過.
民事罰人的義.
如果不是靠行為的時候.

$^{361}$其實原文的意思.
不是解作勝過.
原文是pervasive.
意思不是贏輸的問題.
這個字可以解作多過.
大過,好過,勁過.
厲害過,超過.
都可以,總之不是勝過.
耶穌不是要你贏他們.
重點不是勝負.
不是高低之分.
這個字可以解作質素上的改變.
或者量的改變.
都可以.
這個字其實是.
豐富或者豐盛.
超越.
說得很明白.
pervasive就在五餅二魚的比喻裡面.
當耶穌行完五餅二魚之後.
食物多到爆.
多到連十二個籃子都裝不完.
那個字.
使用了pervasive.
是在說食物多到爆的情況.
食物是超過了.
超過了.
爆了燈,多到死.
這樣的狀況.
就是pervasive.
所以耶穌要求我們的義.
正正就是pervasive.
耶穌要求門徒的義是一種超越.
不斷的追求更高的境地.
如果你覺得登山寶訓很難受.
要愛仇敵.
不可以罵你的弟兄.
被人打完右邊臉,打完左邊臉.
很難.
可不可以不愛仇敵.

$^{401}$可不可以罵他幾句.
可不可以打完他幾下就走了.
這些可不可以.
仍然是一個律法問題.
所以如果你這樣想.
其實你誤會了耶穌的意思.
耶穌不是一個更加嚴苛.
更加高的律法主義者.
如果你這樣看的話.
你就不明白耶穌的意思.
耶穌不是要想做壞事.
耶穌不是說.
他們叫你不要姦淫.
我連想都不讓想.
他們叫你愛人如己.
我連傻子都要愛.
他們叫你不可殺人.
我連生氣都不讓生氣.
如果耶穌的義是一個更加高.
更加難達到.
更加嚴苛的倫理要求的話.
極其量.
這個只是更加高的.
更加嚴苛的律法主義.
更加高的更加嚴苛的律法主義.
耶穌是想我們去.
去不停留在某一個境地.
問可不可以.
可以就OK的做.
而不斷的去尋求.
一個更加高的境地.
不是一個從.
Lv5到Lv100.
不是一個從Lv5到Lv100.
不是一個低標準到很高標準.
而是你不要停留在某一個階段裡.
以為這個就叫做收工.
做足就叫做收工.
不要以為所謂一些.
尾線的事情.

$^{441}$做到某個位就叫做夠.
不要自限.
那個異議的可能.
不要局限尾線的.
那個境地.
不要以為我交了功夫.
就叫做可以做夠數的基督徒.
這個正正就是法律上的問題.
他們不是不好.
他們仍然是一些異議的東西.
但他們以為異議是一些.
做了就叫做了的東西.
主啊我一個禮拜禁食兩次.
凡我得的都是.
十分之一就是這樣.
如果你以為耶穌的異議.
是一個更加高的律法的話.
這樣我們就和法律上沒有分別.
耶穌說凡向弟兄動腦的.
難免受審判.
OK那我不生氣.
那我就暗串他.
耶穌說有人逼你走一里路.
走了兩里路.
OK我沒有走兩里路.
那我就走了.
耶穌說要為你的仇人禱告.
奉主名求眼門.
然後你就對付他.
不是這樣的.
耶穌叫我們去不斷的超越.
不斷的追求更加好的境界.
不要問可不可以.
可不可以.
可以就做,做到盡.
這個不是我們.
耶穌要求我們的意義.
耶穌要求我們.
不斷的超越.
我們的界限.

$^{481}$在我們的能力以內.
繼續去突破.
這個能力的界限.
所以大家聽過的.
免費的東西是最貴的.
耶穌基督給我們恩典.
是免費的.
是白白的.
不過他是要求我們.
要用我們的全部.
來去回應.
這個恩典.
沒錯,這個恩典是不需要你用任何東西去換的.
是無價的.
是免費的.
他同時間是枯燥的.
要去用盡一切的代價.
來去回應他.
不是一種買賣.
不是一種純粹工具.
來去換上帝恩典.
而是一種.
你甘心情願,自動自覺.
所做的事.
所以保羅在.
《哥倫比亞後書》這樣說.
原來基督的愛催逼我們.
一人忌諱眾人死.
眾人就都死了.
這個為主而活的.
那個道理是沒有人監理的.
就是你上了主之後.
你就自動自覺.
因為基督的愛的緣故.
你就自動自覺.
來做一個好的基督徒.
沒有人逼你.
沒有人催你,沒有人罰你.
都沒有人逼你回來教會.
每個禮拜你就是.

$^{521}$自動自覺.
甘心情願,起來.
來這裡敬拜他.
這些不是為了換取什麼.
而是我們.
不斷的來到.
在這個義裡.
不斷去追求更加好的地步.
不是因為律法.
而是純粹因為耶穌基督.
為我們在十字架上死.
我們就自動自覺的.
就起來回來敬拜他.
想想你們的子女.
將來如果他們和你.
喝茶,吃飯.
純粹為了每次喝茶問你.
有什麼好處的時候,就很悲哀的事情.
他們找你喝茶是因為什麼?.
因為你很愛他們.
所以他們甘心情願.
來找你喝茶.
這就是天父上帝和兒女的關係.
不是為了換取救恩.
不是問我可不可以.
給不給.
而是不斷的問.
我怎樣做得更加好.
去回應我們天父上帝.
當然這件事是很難的.
不斷追求.
耶穌登山補糞的義.
所以保羅在最後.
這樣說.
所以你們因信基督耶穌.
都是神的兒子.
你們受洗歸入基督.
的都是被帶基督了.
這個字是一個非常重要的經文.
所以,一個很重要的連接句.

$^{561}$既然你們不在律法之下.
既然你們不在.
BB初階的律法之下的時候.
保羅就勸勉等子女說.
你們要怎樣做.
你們要去被帶基督.
基督徒很難.
不在律法之內.
又不在律法之下.
要長期保持自己在律法之上.
是一件非常難度高的事情.
我們都是一些很普通的人.
不是什麼道德性人.
但是要求不斷自己要在一個更加高的境地裡.
去追求一件很不容易的事情.
親身保訓是不容易的.
因為要我們不斷地去追求更加好的.
所以保羅說.
面對著這個律法.
我們有一個方法.
就是你們要去被帶基督.
被帶基督是一個很熟悉的字眼.
一個很熟悉的字眼.
不過其實我想說.
這是一個挺搞笑的字眼.
什麼叫被帶基督.
其實是一個穿衣服的比喻.
我們要穿上基督.
一個穿衣服的比喻.
我們要被帶著基督的外衣在我們身上.
所以這個字眼其實是一個關於.
今天我想說的cosplay的東西.
就是穿著一件衣服在我們身上.
大家有沒有聽過cosplay.
就是萬聖節那些.
穿著一件衣服.
所以基督徒我們不參加萬聖節.
因為我們每天都有萬聖節.
每天都穿著基督的外衣在身上.
是很刻意的一件事.

$^{601}$你明明不是那樣東西.
但他要穿著那件外衣.
來裝扮基督的外表.
其實我小時候中五的時候.
其實做過這個行業.
就是那時候書記工就做.
穿著老虎衣服在商場裡面.
跟小朋友握手那些.
穿著那些外衣在外面做.
其實是很辛苦的.
因為那件衣服又熱又臭.
因為不是只有我穿.
書記工就好像20元一個小時.
穿著那件衣服走.
很慘的.
因為那些小朋友會攻擊我的.
在後面踢我.
歪我的尾巴.
又拉我的衣服.
後來我就想到方法去報復.
握手的時候.
聽到有人捏小朋友的手.
小朋友就馬上哭.
小朋友就問什麼事.
我就很可愛地說Hello.
那時候我還未相信主.
所以我認罪了.
後來老闆就幫我換了一件.
就是一件翠兒的衣服.
小朋友就不要碰我了.
因為是一個很奇怪的人體翠兒.
是一個人的人形.
頭髮那麼小的翠兒.
這些都不可愛的翠兒.
小朋友就不來找我了.
這個翠兒很可愛.
這個就是彼得多的意思.
你明明不是耶穌.
你卻穿著一件耶穌的外衣.
來盡力去扮成一個耶穌的樣子.

$^{641}$明明不是.
你卻用盡你的方法.
來嘗試去做到像耶穌一樣.
你沒有能力在律法之上.
不過你卻盡力地做到最好.
讓人認得出你就是耶穌基督.
所謂叫做什麼.
就是沒有材料扮四條.
我們明明都做不到的.
但都盡力地做到最好.
好讓人都能夠見到.
我自己在錦書信主.
我是中六的時候信耶穌.
我信主的時候.
就帶我兩個妹妹去教會.
我兩個妹妹就去錦書.
一家五口.
我們三兄妹去信耶穌.
我跟妹妹說.
我們三個一定要做好見證.
給我父母和我媽媽看.
因為我們中學的時候.
我中六的時候.
他們兩個就中三.
暴風少女的年紀.
回到家就都很膽小.
不做家務.
又把東西都拿出來.
又不換校服.
就在電視看.
經常被媽媽說.
我們做基督徒要做好一點.
回家就馬上換衣服.
回家就做家務.
收拾東西.
洗衣服.
幫忙做家務.
我媽媽和爸爸就立刻不同了.
你問我們是不是真的這樣.
當然不是.

$^{681}$我們一向都是比較懶.
但我們也是很願意這樣做.
明不明白.
我們明明不是我們的本性.
但我們是願意為父母和耶穌的緣故.
我們卻是盡力去做好.
這就是我們披萊基多個意思.
我們明明是做不到的.
但我們還是願意不斷去推自己多一步.
無料我們是要做到十足.
這就是我們的意思.
你說是不是很假.
這又不是真的假.
我們是願意去做.
盡力去做.
做基督徒是很難的.
做基督徒不是最本性的你.
一定是在別人面前.
盡力去做好一點.
你做回本性就很放屁.
你就會很懶.
但耶穌要求我們在這個律法之上.
不斷地去追求自己做得更加好.
這不是虛假.
這是我們盡力在軟弱的裡面.
不斷地去盡力做到最好.
我們可以的.
我們可以局部地區性.
間歇性地做到一些好事出來.
在一些親戚面前.
在一些平時和自己有嫌隙的朋友那裡.
在自己的弟兄姊妹或者父母的關係裡面.
仍然可以做到這些事.
這個不是我們本身的我.
它是被要求可以做得更加好的我.
這就是我們願意去披露的基督.
所以等下我們一起去思考.
我們有沒有一些地方.
可以順帶給我們在生命裡面.
要求我們可以取下自己.

$^{721}$面對一些很難受的人.
面對一些很不容易的關係.
我們能夠勝過自己原先的發揮.
超水準,不尋常.
這樣做得好一點.
一個律法之上.
不是問可不可以不理他.
基督徒可不可以不找他.
而是問基督徒可不可以做得更加好.
更加要求做更加好的我.
我一直祈禱.
你去教導我們.
你叫我們不要停留在一個律法之下.
不是問可不可以.
而是不斷追求更加美善.
更加合乎你心意的我們自己.
所以你知道我們是一個罪人.
我們是一個軟弱的人.
我們求主你讓我們能夠在裡面.
能夠經常超越我們自己的界限.
去愛一些我們不可愛的人.
做一些違反我們似乎是利益的事.
讓我們能夠為了你的緣故.
再一次突破我們的界限.
讓我們好好地跟隨你.
讓你知道我們骨子裡仍然是一個軟弱的人.
求主你這樣的幫助.
讓我們能夠在你裡面一次一次地突破我們的有限.
成為一個你喜悅的賢侶.
奉主名求 阿們.
謝謝大家.
\newpage



\section{約翰福音 11:28-44}
\label{sec:7_miM8zb42g}
\textbf{2023.04.02 《跟你休戚與共的耶穌》 約翰福音11\_28-44( 新漢語譯本) 彭秀芹導師}
\newline
\newline
連結: \href{https://youtube.com/watch?v=7_miM8zb42g}{\texttt{https://youtube.com/watch?v=7\_miM8zb42g}} ~~~~ 語音日期: 2023-04-04
\newline
\newline
\hyperref[sec:nJ0r0oqtEjI]{\small{< < < PREV SERMON < < <}}
~
\hyperref[sec:index_chronic]{\small{[返順時目]}}
~
\hyperref[sec:index_scriptual]{\small{[返順卷目]}}
~
\hyperref[sec:AV2ZegtPBEc]{\small{> > > NEXT SERMON > > >}}
\newline
\newline
約翰福音 11:28-44
\newline
\begin{longtable}{cl}
\hline
\hline
章節 & 經文 (和合本修訂版)\\
\hline
11:28 & \begin{tabularx}{0.7\textwidth}{X} 馬大說了這話就回去，叫她妹妹馬利亞，私下說：「老師來了，他在叫你。」 \end{tabularx} \\ \\ \relax
11:29 & \begin{tabularx}{0.7\textwidth}{X} 馬利亞聽見了，急忙起來，到耶穌那裡去。 \end{tabularx} \\ \\ \relax
11:30 & \begin{tabularx}{0.7\textwidth}{X} 那時，耶穌還沒有進村子，仍在馬大迎接他的地方。 \end{tabularx} \\ \\ \relax
11:31 & \begin{tabularx}{0.7\textwidth}{X} 那些同馬利亞在家裡安慰她的猶太人，見她急忙起來，出去，就跟著她，以為她要往墳墓那裡去哭。 \end{tabularx} \\ \\ \relax
11:32 & \begin{tabularx}{0.7\textwidth}{X} 馬利亞到了耶穌那裡，看見他，就俯伏在他腳前，對他說：「主啊，你若早在這裡，我弟弟就不會死了。」 \end{tabularx} \\ \\ \relax
11:33 & \begin{tabularx}{0.7\textwidth}{X} 耶穌看見她哭，並看見與她同來的猶太人也哭，就心裡悲嘆，又甚憂愁， \end{tabularx} \\ \\ \relax
11:34 & \begin{tabularx}{0.7\textwidth}{X} 就說：「你們把他安放在哪裡？」他們對他說：「主啊，請你來看。」 \end{tabularx} \\ \\ \relax
11:35 & \begin{tabularx}{0.7\textwidth}{X} 耶穌哭了。 \end{tabularx} \\ \\ \relax
11:36 & \begin{tabularx}{0.7\textwidth}{X} 猶太人就說：「你看，他多麼愛他！」 \end{tabularx} \\ \\ \relax
11:37 & \begin{tabularx}{0.7\textwidth}{X} 其中有人說：「他既然開了盲人的眼睛，難道不能叫這人不死嗎？」 \end{tabularx} \\ \\ \relax
11:38 & \begin{tabularx}{0.7\textwidth}{X} 耶穌又心裡悲嘆，來到墳墓前。那墳墓是個穴，有一塊石頭擋著。 \end{tabularx} \\ \\ \relax
11:39 & \begin{tabularx}{0.7\textwidth}{X} 耶穌說：「把石頭挪開！」那死者的姐姐馬大對他說：「主啊，他現在必定臭了，因為他已經死了四天了。」 \end{tabularx} \\ \\ \relax
11:40 & \begin{tabularx}{0.7\textwidth}{X} 耶穌對她說：「我不是對你說過，你若信就必看見神的榮耀嗎？」 \end{tabularx} \\ \\ \relax
11:41 & \begin{tabularx}{0.7\textwidth}{X} 於是他們把石頭挪開。耶穌舉目望天，說：「父啊，我感謝你，因為你已經聽了我。 \end{tabularx} \\ \\ \relax
11:42 & \begin{tabularx}{0.7\textwidth}{X} 我知道你常常聽我，但我說這話是為了周圍站著的眾人，要使他們信是你差了我來的。」 \end{tabularx} \\ \\ \relax
11:43 & \begin{tabularx}{0.7\textwidth}{X} 說了這些話，他大聲呼叫說：「拉撒路，出來！」 \end{tabularx} \\ \\ \relax
11:44 & \begin{tabularx}{0.7\textwidth}{X} 那死了的人就出來了，手腳都裹著布，臉上包著頭巾。耶穌對他們說：「解開他，讓他走！」 \end{tabularx} \\ \\
[1ex]
\hline
\hline
\end{longtable}
$^{1}$鼎姐妹平安.
平安.
不知道在座的鼎姐妹.
近來有沒有哭過呢?.
通常哭的時候背後有些複雜的情緒.
可能很傷心.
可能很甘心.
可能很委屈.
也可能太開心.
開心到哭.
想和你們分享一下.
近期我家裡一起哭一個事件.
我家裡人不多.
一家三口.
話說我女兒去年十月要去日本留學.
越接近女兒要離開香港.
心情就越忐忑.
因為很捨不得.
剛好星期一神學院.
會有一個心理分析小組.
可以和同學分享.
我女兒星期二飛.
星期一心理分析小組.
當然我分享這件事.
我很捨不得女兒.
我的眼淚不自然控制不到流出來.
我大約說著說著就哭了.
差不多九個小時.
因為我一個人會觸動你的眼淚.
到了星期二.
我丈夫載女兒去機場.
出車前我們一家人祈禱.
我開始祈禱.
到女兒和丈夫結束祈禱.
到我丈夫祈禱的時候.
我記得他說了一句.
天父啊!謝謝你賜給我們蜜月.
我女兒叫蜜月.
我等了很久.
為什麼沒有第二句祈禱呢?.

$^{41}$我忍不住睜大眼睛看著我老公.
我看到他淚流滿面.
他哭得像我一樣.
很捨不得.
他無法祈禱.
原來今天你所愛的人.
會觸動你的眼淚.
原來今天.
在剛才讀聖經裡.
你也看到35字說.
耶穌哭了.
耶穌為什麼會哭呢?.
今天的題目就說.
跟你憂戚與共的耶穌.
原來今天耶穌和我們是一樣的.
是有感情,有感覺,有情緒的.
原來今天.
我盼望透過拉薩路.
這個聖經.
讓我們知道.
原來耶穌和我們是憂戚與共.
在艱難裡.
在痛苦裡.
在絕望裡.
原來神要我們看見.
祂的榮耀和祂的憐憫.
大家應該很熟悉經文.
就是說到拉薩路復活.
在1-27字.
我沒有讀出來.
因為太長了.
大約就是說.
耶穌和瑪利亞.
拉薩路這家人.
因為他們是姊弟.
兄弟姊妹.
是很熟悉的.
原來在經文裡.
有介紹.
特別介紹瑪利亞.

$^{81}$瑪利亞是拉薩路的姊妹.
特別說到.
原來這位瑪利亞.
之後11章是說拉薩路復活.
12章是說.
瑪利亞獻香膏的那位.
還有用頭髮擦乾.
擦主的腳的那位.
這件事.
也告訴我們.
瑪利亞很愛主.
他刻意.
作者再次交代瑪利亞很愛主.
原來今天.
我們每一個愛主的人.
主清不清楚?.
清楚的.
主是知道的.
來紅印加.
這裡是實習.
那條樓梯.
我經常看到很美麗的背影.
看到很多弟兄姊妹.
可能腳不太聽洗.
膝蓋也痛.
但他仍然堅持.
硬上來敬拜主.
這些愛主的行為.
主一一看見.
一一記下.
主明白.
之前也說.
除了介紹瑪利亞.
也說到很有趣的.
這家人和耶穌很熟.
熟到什麼程度呢?.
他們知道耶穌的行蹤.
瑪利亞和瑪達.
打發人們.
向耶穌說.

$^{121}$阿拉什病了.
很有趣.
這個說法.
好像很隱隱約約.
說什麼呢?.
剛才沒有讀.
第三節.
主啊.
你所愛的人病了.
很有趣.
沒有求耶穌.
通常人們求耶穌醫他.
耶穌快死了.
你快點來我家醫他.
通常都是這樣.
但你看到.
他這個報信.
通知耶穌.
你所愛.
沒有說什麼名字.
也說他病了.
也沒有求耶穌去醫他.
你所愛是誰呢?.
就是說耶穌和拉撒路.
和瑪利亞.
和瑪達.
有一個很深的關係.
其實當這個報信的人.
找到耶穌.
去交代拉撒路病了.
其實沒有說.
你所愛的人病了.
應該那天.
拉撒路已經死了.
按很多醫治.
耶穌不用去現場說一句.
那個人就不用死.
但耶穌沒有這樣做.
耶穌還特意留多兩天.
報信一天.

$^{161}$再留多兩天.
然後再由他的地方.
去巴大尼.
又一天路程.
就是說耶穌總共遲了四天.
這個家庭這麼需要耶穌.
這麼絕望.
最重要的是.
耶穌很愛他們.
聖經說.
耶穌愛瑪達.
瑪利亞和拉撒路.
這麼愛他們.
這麼需要耶穌的時候.
耶穌缺席.
不在.
可能有時候.
也有這個感覺.
我這麼辛苦.
這麼多難處.
耶穌去了哪裡呢?.
耶穌好像在我人生.
很難過的關口裡.
缺席了.
就像瑪達瑪利亞一樣.
她覺得.
我明明叫人通知耶穌.
為什麼沒有反應?.
為什麼耶穌缺席了?.
可能今天你會覺得.
這個世界有很多痛苦.
很多難處.
不斷地向你湧過來.
可能你已經覺得.
我筋疲力盡.
還是沒有轉機.
可能是你的身體疾病.
可能是你的心理軟弱.
經濟困難.
或者人際關係.

$^{201}$很多很多的難處.
這些難處.
都好像在我們生命裡.
而耶穌也好像不在.
缺席了.
但是,小妹妹.
不要放棄.
原來上帝不會忘記他的兒女.
也不會.
拋棄一些.
他愛的人.
也不會向他愛的人隱藏.
在十九世紀.
有個作家.
叫做安娜華納.
原來她是一個有錢人.
的家庭出生.
她爸爸是一個很有名的律師.
但是適逢.
一件事令她家裡巨變.
就是經濟大蕭條.
1929年.
大家應該都未出生.
所以大家都很年輕.
1929年.
經濟大蕭條是什麼呢?.
按照這個讀法告訴你.
1929年華爾街股災.
席捲全球.
經濟大蕭條.
對發達國家和發展中的國家.
都帶來毀滅打擊.
人均收入.
稅收.
盈利.
價格全面下創.
國際貿易銳減50\%.
美國失業率飆升至25\%.
有些國家甚至達到33\%.
你猜到.

$^{241}$全球.
當然.
安娜在美國.
是首當其衝.
你知不知道香港失業率有多高?.
我們香港失業率大約是3.4.
1月公佈時是3.4.
但當時失業率有多高?.
25\%.
她的爸爸.
一夜之間.
蒸發了所有身家.
所以.
安娜.
原本是有錢的.
她懂得寫字.
但現在.
她一無所有.
原來她家裡都堅持.
每個星期日.
都回到教會當中服侍.
因為安娜.
她很喜歡寫作.
她也是.
主要是一個教師.
她有寫小說.
去謀生.
其中一本小說.
講一個男孩子.
死後.
男主角.
唱一首歌去安慰他.
這首歌.
大家應該都很熟悉.
我們會唱.
一起唱.
歌詞已經在PowerPoint.
(歌詞).
原來這首歌.
是上世紀.

$^{281}$我小時候.
都會唱這首歌.
回到教會.
很出名.
這首歌.
原來.
那個年代.
是第一次世界大戰.
很多軍人.
美國軍人.
唱這首歌.
這首歌.
歌詞簡單.
旋律簡單.
在他們面對恐懼.
面對死亡.
面對艱難的時候.
他們經常都唱這首歌.
原來.
在很多艱難裡.
耶穌.
是愛我們的.
我們要認識.
愛我們的主.
歌詞特別有英文.
Jesus loves me.
This I know.
即是.
我知道什麼.
知道耶穌愛我.
原來.
This I know.
是一個信仰的.
一個記號.
信仰記號就是.
都盼望成為我們信仰記號.
你知道什麼.
當你很痛苦.
很絕望時.
知道什麼.

$^{321}$知道耶穌愛我.
這句話.
也成為.
中國.
有一段時間.
遭受大迫伐的基督徒.
有些美國的宗教學士.
來到中國.
他們也遭受迫伐.
他們寫信回到自己的家鄉.
他們會怎樣寫呢.
This I know.
知道什麼.
不是知道痛苦.
不是知道眼淚.
知道耶穌愛我.
在痛苦裡.
困難裡.
在絕望裡.
他們知道什麼呢.
知道耶穌愛我.
原來.
耶穌.
是怎樣呢.
就是愛我們.
認識我們.
愛我.
認識我的耶穌.
This I know.
我都盼望.
在我們的信仰歷程裡.
都很深.
把握住這件事.
可能環境很難.
但是你知道耶穌愛你.
耶穌認識你.
其實在整件事裡.
我再說聖經.
拉撒路的事件.
你會看到耶穌刻意.

$^{361}$特意遲到.
為什麼他要特意遲到呢.
原來目的是要人.
知道他有復活的大能.
拉撒路是怎樣呢.
死透了.
四天了.
原來他要人看到.
耶穌能夠行這神蹟.
要證明.
很具體地證明.
主耶穌是掌管生命的主宰.
和滿有復活的大能.
所以在第四節之前.
我沒有讀第一節.
耶穌說過.
這病不至於死.
因為要叫神得榮耀.
叫兒子得榮耀.
看來環境很絕望.
很艱難.
但原來神要你去榮耀他.
要你經歷他復活的大能.
剛才主席已經讀了.
28節開始讀.
其實在1-27節裡.
耶穌有和瑪大對話.
我先不說.
我先看看瑪利亞和耶穌的對話.
同一個場景.
在巴大利的村莊.
耶穌還沒有進城.
瑪大看到耶穌來了.
迎接了耶穌.
回家後暗暗地對瑪利亞說.
耶穌找你.
耶穌叫你.
為什麼要暗暗地對話呢.
當時.
你看看後面.

$^{401}$猶太人已經計劃用石頭扔耶穌.
其實耶穌也冒著被石頭扔的危險.
進入村莊.
所以瑪大不想太張揚.
暗暗地對瑪利亞說.
瑪利亞就起來迎接耶穌.
因為當時喪家.
家裡有喪事.
猶太人的慣例.
親戚朋友會到他們家裡.
安慰死者的遺屬.
瑪利亞就起來迎接耶穌.
所以陪瑪利亞的人.
都跟著一起到村口找耶穌.
瑪利亞見到耶穌.
第一句說什麼呢.
和瑪大一樣.
主啊.
你若.
他說什麼呢.
主啊 要是你早在這裡.
我的兄弟就不會死.
瑪利亞說了他愛主.
但你看到他這句話.
又不太敢怪主.
但有很多不明白.
有很多遺憾.
這裡特別形容.
瑪利亞苦服在主腳前.
這幅圖畫大家很熟悉.
記不記得耶穌去他們家裡.
應約去到他們家裡.
瑪大去服侍耶穌.
瑪利亞選擇上好福分.
在耶穌的腳前聽耶穌的話.
這幅圖畫也是.
瑪利亞在耶穌的腳前.
但他很遺憾.
很悲傷.
很絕望.

$^{441}$他將他的憂傷.
帶到耶穌的腳前.
今天愛主的人.
不代表沒有眼淚.
有信心的人.
不代表沒有遺憾.
沒有悲傷.
我不知道你會不會將你的眼淚.
帶到耶穌的腳前.
耶穌沒有怪責瑪利亞.
沒有說瑪利亞你這麼不信.
他沒有.
你聽到耶穌完全明白.
接納瑪利亞當時的眼淚.
一看下去你就會知道.
為什麼我會這樣說.
耶穌是明白接納瑪利亞的眼淚.
其實這句話.
反映出瑪利亞很無奈.
如果早點在這裡.
他就不會死.
其實我們人生有很多時是這樣.
如果這樣.
主啊如果這樣就好了.
主啊如果這樣.
我就不會這樣了.
我們有很多限定對我們的主.
但主不會因為這樣.
而不去安慰瑪利亞.
或者不站在瑪利亞的旁邊.
不是將瑪利亞推開.
沒有.
耶穌讓瑪利亞繼續在他的腳前.
去表達他的眼淚.
表達他的遺憾.
那麼.
耶穌看到瑪利亞哭.
也看到和他一起來的猶太人都哭.
原來這個哭字.
我們中文就覺得.

$^{481}$流下眼淚.
不是的.
聖經說這個哭就是大哭.
就是.
怎麼說呢.
又哭又叫.
我不知道.
我們中國人有哭喪.
會專門請人哭.
當時猶太人也是.
他們有一些專業的哭喪員.
表達他們很懷念.
很懷念這位死者.
當時四周.
環境就是這樣.
人們很懷念.
大哭大叫.
那麼.
耶穌看到.
是一片愁雲慘霧.
而且死亡的罪行圍繞著這群人.
好了.
那麼瑪利亞他哭了.
我們看看第三節.
他就這樣說.
他說耶穌看到他們哭.
不好意思.
我再三十三節說.
耶穌看到他們哭.
有一句說話.
耶穌內心悲憤激動起來.
這個悲憤激動起來.
就是耶穌為什麼會激動呢.
為什麼會生氣呢.
為什麼會悲憤呢.
這個字有一個動作.
形容馬在噴氣.
不知道你們試過沒有.
有些事情令你又傷心又生氣.
你無法說話.

$^{521}$你只會這樣.
那份心情.
你猜耶穌明不明白當時.
那些人哭什麼.
他明白.
他心亂悲憤.
他悲憤什麼呢.
原來他知道死亡.
是引犯罪帶來.
他哭其實耶穌悲憤.
他不是悲憤.
搞錯了.
馬大瑪利亞和猶太人.
為什麼不明白.
我就復活主.
我一會就叫你弟弟復活.
其實在一月二十七節.
是明明說了給馬大聽.
我一會就叫你弟弟復活.
問他信不信.
耶穌不是生氣他們不信.
耶穌原來生氣什麼呢.
就是生氣死亡本身.
原來死亡帶給人.
那份絕望和破壞.
當人犯罪的時候.
就帶來死亡.
死亡就是罪生出來的結果.
其實神是很生氣這件事.
不是生氣當時的人.
例如我們看到打仗.
打仗的圖畫.
心裡是怎樣.
不是生氣那些東西都爛了.
是生氣打仗帶來的破壞.
耶穌見到人死亡.
不是生氣那人不信他死了.
而是他生氣.
罪帶來死亡那種絕望.
但原來耶穌.

$^{561}$他要率先將死亡打敗.
於是這裡說的第35節.
就是全本聖經最短的經節.
就是四個字中文.
耶穌哭了.
耶穌有眼淚.
耶穌是神.
其實他也知道.
他稍後可以叫拉撒路復活.
耶穌滿有大能.
耶穌的哭.
代表著對人的憐憫.
對人的需要.
對人的悲傷.
那份敏銳.
雖然主將會抹乾我們所有眼淚.
所以我們不會再流眼淚.
但未抹乾之前.
即是還有眼淚.
這些眼淚主是看見.
主是明白.
主是願意和我們憂戚與共.
與我們同行.
主是完完全全的一個人.
他會流眼淚.
他願意和我們一起分擔.
他願意和我們一起面對.
前年我團友的姐妹.
有一件事.
有一天我和那位姐妹.
當作是A姐妹.
參加B姐妹的爸爸喪禮.
參加完之後.
我們一起離開.
突然間A姐妹接了一通電話.
接了電話後.
神色很慌張.
又哭著想衝出馬路.
我和她一起走.
我當然拉著她.

$^{601}$究竟發生什麼事呢.
她整個人都好像失了相.
拉著她之後.
她就狂哭.
原來這通電話.
就是有人通知她.
她爸爸跳樓自殺.
她那一剎那.
其實是接受不了.
她不知道該怎麼辦.
其實她衝出馬路.
她不是想自殺.
她很想快點回家找媽媽.
所以她整個人突然失了方向.
我就陪著她回到她家.
原來去一個.
有人去世的家庭裡.
心情很沉重.
你看到她們的眼淚.
看到她們的悲傷.
看到她們很多.
不知道該怎麼辦.
因為A姐妹的爸爸.
是在二十多樓跳下來.
我心裡想.
這麼高跳下來.
當她去認屍的時候.
其實會更傷心.
因為想到場面會是怎樣.
姐妹的哥哥回到家.
因為要去現場確認.
他就說.
爸爸的屍體好像睡了.
其實當時我不太相信.
我覺得他是說謊.
騙家人.
讓他們不會那麼傷心.
到了真的要再去殮房.
認屍的時候.
她再跟我說.

$^{641}$是真的.
我爸爸好像睡了.
她會體會.
雖然很傷心.
但看到神有憐憫.
我不知道為什麼.
她不會血肉蒙混.
因為我記得我小時候.
去參加一個訓練.
我要去一層樓.
大約十步前.
有個人在我面前跳下來.
我看到他血肉蒙混.
手腳都被他咬了.
我記得.
我心裡覺得.
他爸爸的屍體一定會更傷心.
但很奇妙.
他說他好像睡了.
其實神是有憐憫.
神明白這個家庭的困難.
那個眼淚.
耶穌哭了.
代表耶穌跟你一樣.
祂有眼淚.
祂明白你.
當耶穌哭了.
有些人不是這樣想.
有些人覺得.
耶穌很愛拉薩路.
所以哭了.
但有些人覺得.
耶穌愛拉薩路.
那又怎樣?.
祂能夠醫治盲人開眼.
祂不能醫治死了的人嗎?.
其實這時候是不信的.
他們不期望耶穌有能力.
拉薩路復活.
其實這個期望.

$^{681}$這個不期望.
包括馬大,瑪利亞.
他們都不期望.
耶穌能夠讓弟弟復活.
他們的期望是.
第二次復活.
不是現在.
但耶穌不是.
不是第二次.
是現在.
耶穌的能力和幫助.
不是第二次.
是現在.
當我們看到耶穌的眼淚.
我們有沒有誤解耶穌的眼淚呢?.
覺得耶穌不明白我們.
不是的,耶穌不明白我們.
耶穌去到拉薩路墳前.
祂心裡再次悲憤.
剛才和那個字悲憤是一樣的.
祂再次悲憤.
為什麼悲憤呢?.
因為這次要正面和死亡.
這個宿敵,敵人.
正面的交鋒.
原來墳墓是代表死亡.
在38節裡.
特別說到很有趣.
在39節裡特別說到.
死人的姐姐.
特別形容是死人的姐姐.
警告耶穌.
拉薩路怎麼樣?.
已經死了四天.
其實是勸耶穌怎麼樣?.
已經沒得搞了.
有時候我們也會覺得.
沒得搞了.
神仙也難救.
因為猶太人的傳統.

$^{721}$就是死頭第四天.
頭一,二,三天死了.
可能靈魂還沒離開.
可能還在附近徘徊.
但第四天靈魂一定離開.
死了.
沒得搞了.
絕望了.
但耶穌就是要在絕望.
沒得搞的環境裡.
顯出祂的神蹟.
耶穌不單要讓拉薩路復活.
而且要叫醒死亡中的拉薩路.
死了能叫醒嗎?.
是叫不醒的.
但耶穌能叫醒.
耶穌就是復活的主.
今天可能我們的心靈也死了.
對耶穌沒有反應.
但我告訴你.
耶穌好像叫拉薩路那樣叫我們.
後面會說到.
耶穌怎麼叫拉薩路呢?.
就好像怎麼叫我們.
但在還沒叫拉薩路之前.
耶穌有一個祈禱.
祂問人.
叫人拿開石頭.
然後祂望著天祈禱.
因為你已經聽我.
在文章裡.
這個聽是過去式.
就是說耶穌已經預先為拉薩路祈禱.
今天你覺得你預先知道耶穌不知道嗎?.
神不知道?不是啊.
預先知道的.
為了你祈禱.
神已經聽了你的祈禱.
祂表明.
主耶穌和天父的關係很密切.

$^{761}$祂知道天父會聽祈禱.
其實不單止主耶穌知道天父會聽祈禱.
其實瑪利亞也知道.
馬大也知道.
門徒也知道.
祂見過耶穌.
很多次的祈禱.
很多次的醫治.
都是透過父神.
去洗很多的艱難.
或者很多的絕望改變.
這次行不行呢?.
今天很多時候.
我們知道不代表相信.
我們知道很多很多東西.
就好像瑪達瑪利亞.
其實她很認識主.
她知道很多東西.
但那一刻她信不了.
她那一刻被她眼前死了四天的弟弟.
死透了.
那個絕望被框住了.
有時候我們也會這樣.
被我們的絕望.
覺得很多的預設被框住了.
但耶穌就是生命的主.
可以將你覺得被框住的東西.
拿走.
第43,44節說到.
很大聲叫拉薩路.
聖經說到.
是很大聲呼喊的.
就好像跟一個人說話.
很大聲很大聲.
拉薩路你出來.
當拉薩路沒有死.
你不會跟死人說.
拉薩路你不會的.
你去掃墓.
你不會很大聲跟墳頭說.

$^{801}$你不會的.
你知道他聽不到.
但很厲害.
我們生命的主.
就算我們的心靈已經死寂了.
什麼都聽不到.
耶穌也可以.
令我們經過他的復活大樓.
今天耶穌可能在你心底裡.
叫我們.
彭秀琴出來.
我們要起來了.
很奇妙.
拉薩路.
聖經說.
他是死人出來.
他出來應該沒有死.
但也形容他是死人.
為什麼呢.
原來今天我們的生命.
不是在乎自己.
是在乎主.
最有趣的是.
他死人出來.
我不知道.
當時他還是包著.
手腳包著.
臉也包著.
其實看不到他的樣子.
我不知道當時拉薩路的神情是怎樣.
也可能拉薩路說不到話.
被人封住了.
我不知道瑪利亞會怎樣.
瑪大會怎樣.
我不知道有沒有人敢走前去.
幫他拔開.
解開果絲帶.
我不知道有沒有人敢.
拉薩路的復活.
嚴格來說.

$^{841}$不可以說是復活.
只是復生.
因為拉薩路生完之後.
他會再死.
這個場景很熟悉.
拉薩路的復活.
距離耶穌釘十字架.
其實不是很遠.
耶穌也會有一天.
跟拉薩路一樣.
他會被埋葬在墳墓裡.
也跟拉薩路一樣.
被人用很多塊布.
包著.
不能動.
躺在那裡.
看別人看.
他真的死了.
但他跟拉薩路不同.
耶穌復活.
不用有人幫忙拔掉帶.
不用.
因為他就是復活的主.
他一出來就永遠是復活的主.
復活之後.
他不會再死.
不像拉薩路.
拉薩路會再死.
但復活主是不會再死.
所以耶穌說什麼呢.
我是復活和生命.
耶穌說.
我是復活和生命.
相信我的人.
雖然死了.
也必活過來.
我們是會死的.
我們是有難處的.
但我們信的是誰.
是復活和生命的主.

$^{881}$This I know.
我們就要知道這些.
也要經歷這些.
原來耶穌是滿有復活大能的神.
他打敗了人類的宿敵死亡.
我們不會再受死去克制.
雖然我們真的會死.
但我們知道耶穌有復活大能.
我們有盼望.
我們能夠活出主復活的生命.
但同時雖然主是復活的主.
但他也是一個真正的人.
他會痛哭,開心,哭泣,流淚.
他和我們人生路上同行.
有一段聖經.
我以前不太明白.
他說所有活著又信我的人.
必定永遠不死.
不是的,我們每個人都會死.
為什麼會永遠不死呢.
看回英文聖經.
信我的人是怎樣.
Believe in me.
在主裡面.
是和主有聯繫.
今天我們怎樣可以體會.
經歷神的生命和復活大能呢.
你今天是否在主裡面呢.
你有沒有住在主裡面呢.
有沒有Believe in Jesus.
有沒有住在神裡面呢.
如果我們今天真的相信主.
相信主的復活大能.
我們是一個親近主的人.
我們不會少少事就斤斤計較.
我們會祈禱.
我們會去服事.
因為我們信的神是活著的.
不會少少事就斤斤計較.
所以今天我們要怎樣呢.

$^{921}$要住在主裡面.
和主有聯繫.
活出從主而來的生命.
我們今天不是靠不死的精神.
人會死的,沒有不死的精神.
我們靠什麼呢.
我們靠死而復活的主.
我們經歷到這件事.
我相信大家都看過.
魯牧師的見證.
不知道大家有沒有看過.
他發給我看.
當我預備去靜悟的時候.
他這篇見證都能夠激勵到我.
我讀一小段給大家聽.
他說.
期間,指他患病.
期間我從來沒有害怕或害怕的感覺.
一直陪著我的除了太太.
只有與神對話.
而與神對話給我一種很安穩.
很安心的感覺.
這就是神藉著我這個病.
來回應我提出的疑問.
因此我感到這次入院經歷是一個恩典.
是神讓我經歷著死亡的威脅.
卻讓我安然度過.
我心裡沒有半絲恐懼.
只有平靜安穩的心.
這都是與神同在帶來真實的感受.
這讓我從心底裡向神發出讚美和頌讚.
弟兄姊妹,信仰的寶貴.
就算無論你在低谷或面對死亡.
神都不會放棄你.
祂必定與你同在.
與你度過每一個艱難的處境.
所以大家要珍惜主耶穌.
不要覺得信仰是可有可無.
我相信魯牧師在死亡邊緣裡經歷.
是復活主的生命的同在.

$^{961}$除了你覺得魯牧師很痛很慘.
其實你有沒有看到.
神的手在魯牧師身上?.
如果不是有主.
人面對這麼大的艱難.
會有很多的埋怨或灰心失望.
但你看到他的話語裡.
又安穩又不害怕.
又滿有恩典又讚美.
弟兄姊妹,這就很奇妙.
今天愛主人不是一帆風順.
也不是沒有難處.
剛剛相反.
很多時候愛主人有很多艱難.
有很多軟弱.
但在艱難和軟弱裡.
要顯出神榮耀復活的大能.
魯牧師,今天在艱難裡.
你覺得很絕望裡.
你去經歷這位復活的救主.
他能夠成為.
不單成為你的幫助.
而他會改變你.
你能夠去榮耀主.
我們因著愛神的生命.
可以說一句.
剛才說的就是.
Shall never die.
那種不死.
那種不死不是因為我們肉體不死.
而是有主在.
我們為主是復活的主.
我們的生命能夠影響人.
最後想分享拉撒路.
這個人.
拉撒路是一個沉默的見證人.
在聖經裡沒有記載過.
拉撒路有半句台詞.
沒有.
他沒有說過話.

$^{1001}$當然我相信他有和耶穌說話.
但聖經沒有記載過拉撒路有說過話.
而且拉撒路.
他復活之後.
記載過在十二章.
和耶穌一起坐直.
好像很少提及這個人.
這個人經歷了神的厲害榮耀作為.
但他並不是光光烈烈的站出來.
好像保羅,彼得約翰.
出來正道.
有很多光輝歷史.
聖經沒有什麼光輝歷史.
但原來實際上.
他是耶穌所愛的人.
而且他是見證耶穌復活生命的人.
可能今天我們都很卑微.
人們不認識你.
但我們都是耶穌所愛的人.
是耶穌見證復活.
見證他大能的見證者.
到了十二章時.
如果你有機會可以看回楊方.
十二章說到.
很多人因為拉撒路復活而信耶穌.
今天我們的生命都是.
你是主所愛的人.
人會因著你的生命去信主.
你都是主的見證人.
原來拉撒路復活的見證.
令很多人信主.
所以祭司除了想殺耶穌.
其實連拉撒路也想殺.
拉撒路站出來作見證.
是冒著生命危險.
雖然復活了.
但也有人想弄死他.
原來今天我們去作見證.
我們生命的彰顯上的榮耀.
是有難處的.

$^{1041}$是有壓力的.
是有困難的.
但是我們要知道.
This I know.
耶穌愛我.
你是主所愛.
我們要住在主裡面.
Shall never die.
我們就會經歷神的生命.
雖然我們微不足道.
但我們是神所愛的.
我們每一個人.
都用我們的生命去見證復活的主.
也見證和我們休戚與共的耶穌.
我們有時間一起禱告.
主多謝你.
你是與我們休戚與共的耶穌.
你在裡面滿有大能.
滿有生命.
但你樂意靠近我們.
你樂意與我們同哭同笑.
多謝你今天要我們經歷你.
就是那位充滿復活生命大能的主.
主你幫助我們每一個.
每天都住在你裡面.
讓我們經歷你那份生命的大能.
讓我們的生命.
能夠見證你的榮耀.
我們這樣祈禱是奉主的命求.
阿們.
謝謝大家.
\newpage



\section{馬可福音 10:35-45}
\label{sec:AV2ZegtPBEc}
\textbf{2023.04.09 《耶穌的苦杯,我們能喝嗎?》馬可福音10\_35-45 (和修版) 曾敏傳道}
\newline
\newline
連結: \href{https://youtube.com/watch?v=AV2ZegtPBEc}{\texttt{https://youtube.com/watch?v=AV2ZegtPBEc}} ~~~~ 語音日期: 2023-04-10
\newline
\newline
\hyperref[sec:7_miM8zb42g]{\small{< < < PREV SERMON < < <}}
~
\hyperref[sec:index_chronic]{\small{[返順時目]}}
~
\hyperref[sec:index_scriptual]{\small{[返順卷目]}}
~
\hyperref[sec:b_zpDLNhNUk]{\small{> > > NEXT SERMON > > >}}
\newline
\newline
馬可福音 10:35-45
\newline
\begin{longtable}{cl}
\hline
\hline
章節 & 經文 (和合本修訂版)\\
\hline
10:35 & \begin{tabularx}{0.7\textwidth}{X} 西庇太的兒子雅各和約翰進前來，對耶穌說：「老師，我們無論求你甚麼，願你為我們做。」 \end{tabularx} \\ \\ \relax
10:36 & \begin{tabularx}{0.7\textwidth}{X} 耶穌對他們說：「要我為你們做甚麼？」 \end{tabularx} \\ \\ \relax
10:37 & \begin{tabularx}{0.7\textwidth}{X} 他們對他說：「在你的榮耀裡，請賜我們一個坐在你右邊，一個坐在你左邊。」 \end{tabularx} \\ \\ \relax
10:38 & \begin{tabularx}{0.7\textwidth}{X} 耶穌對他們說：「你們不知道所求的是甚麼。我所喝的杯，你們能喝嗎？我所受的洗，你們能受嗎？」 \end{tabularx} \\ \\ \relax
10:39 & \begin{tabularx}{0.7\textwidth}{X} 他們對他說：「我們能。」耶穌對他們說：「我所喝的杯，你們要喝；我所受的洗，你們也要受。 \end{tabularx} \\ \\ \relax
10:40 & \begin{tabularx}{0.7\textwidth}{X} 可是坐在我的左右，不是我可以賜的，而是為誰預備就賜給誰。」 \end{tabularx} \\ \\ \relax
10:41 & \begin{tabularx}{0.7\textwidth}{X} 其餘十個門徒聽見，就對雅各和約翰很生氣。 \end{tabularx} \\ \\ \relax
10:42 & \begin{tabularx}{0.7\textwidth}{X} 耶穌叫了他們來，對他們說：「你們知道，外邦人有君王作主治理他們，有大臣操權管轄他們。 \end{tabularx} \\ \\ \relax
10:43 & \begin{tabularx}{0.7\textwidth}{X} 但是在你們中間，不可這樣。你們中間誰願為大，就要作你們的用人； \end{tabularx} \\ \\ \relax
10:44 & \begin{tabularx}{0.7\textwidth}{X} 在你們中間誰願為首，就要作眾人的僕人。 \end{tabularx} \\ \\ \relax
10:45 & \begin{tabularx}{0.7\textwidth}{X} 因為人子來，並不是要受人的服事，乃是要服事人，並且要捨命作多人的贖價。」 \end{tabularx} \\ \\
[1ex]
\hline
\hline
\end{longtable}
$^{1}$弟兄姊妹平安.
亦都係度好想歡迎我地當中ge 新朋友.
如果你今日係第一次黎到我地ge 教會去聚會ge .
我想請你舉一舉手好冇呀.
我地想稍稍認識你.
我地ge 師事招待會即刻送上一份禮物比你地ge .
請你舉返手 頭先ge 新朋友 我地即刻會送上一份禮物 好冇呀.
gwo 邊都有ge  師事招待送上禮物.
好 非常歡迎你.
係 仲有前面啦.
gwo 邊都有一個ge 喎 比jor 未.
比jor 啦係咪.
好.
今日係一個好寶貴ge 時間.
我深信親友們邀請大家黎到呢.
都係覺得能夠係教會裡面認識上帝.
係一件非常非常寶貴ge 事情.
大家一睇到呢一個圖畫 諗起D 乜ye 呢.
驚ge 冠冕 係咪.
紅色ge 布.
好快會諗起壽苦節.
今日呢個題目叫做耶穌ge 苦杯.
我地能學嗎 係一個問題黎ge .
其實呢個問題呢有兩層意思ge .
第一層ge 意思呢係對住呢已經返教會.
一段時間ge 弟兄姊妹去講ge .
第二層ge 意思呢就對住我地ge 新朋友去講ge .
咁所以呢大家呢可以今日呢留意下.
呢個ge 說話呢係對住我地講ge .
有唔同ge 人有唔同意思.
講緊苦杯.
今日係咩日子呢.
仲講緊壽苦節ge 訊息.
今日好明顯係復活節.
復活節喎 嘩.
我知道周圍都有種好歡喜快樂ge 氣氛.
係咪 即係有好多ge 禮物ge 交換.
你睇D 商場ge 報紙都係好開心ge .
但係我想問大家一個問題呢.
復活節ge 意義究竟係乜ye 呢.

$^{41}$大家有冇回應呀.
復活節ge 意義係乜ye .
新朋友都可以答ge .
有冇人願意嘗試回答下.
係咪有人想回答.
復活節ge 意義係乜ye .
紀念耶穌ge 復活係非常好喎.
原來呢唔係尋找復活蛋.
我知道唔同ge 地方都有D 活動.
好開心ge  尋找復活蛋呀.
周圍走下.
其實係紀念主耶穌ge 復活.
今日呢個日子ge 意義就係呢一度.
但係我可以咁樣講.
要講到復活呢必須由.
耶穌基督ge 受苦開始講起.
冇耶穌ge 受苦.
冇耶穌係十字架上ge 犧牲.
就唔會有復活.
復活唔係就咁單單ge 存在.
係之前有個兩部曲.
受苦 犧牲.
先會黎到復活.
今日我同你好開心ge .
慶祝復活節.
歡喜快樂 好有盼望.
心情都好喜樂好雀躍ge .
要記住之前.
係經歷過兩個階段.
受苦同埋犧牲.
所以今日呢想同大家呢.
從受苦 耶穌ge 受苦同埋犧牲.
gwo 度開始.
究竟耶穌係點樣去睇自己ge .
受苦同埋犧牲ge 呢.
佢會點樣去描述呢.
睇到呢幅圖畫呢.
諗到係耶穌基督係釘係.
被釘係十字架上.
我地係稍為睇前少少經文.

$^{81}$今日我地ge 經文係馬學福音十章.
三十五到四十五節.
我地會推翻前少少睇.
睇到呢耶穌呢.
係真正開始受苦前呢.
佢係有幾次呢去預言.
佢自己即將係受苦.
係會被呢被害ge .
咁我地淨係睇下第三次.
第三次因為完jor 呢一次之後呢.
佢就開始展開佢ge 受苦之路.
第三次呢一個係十章ge .
三十二到三十四節.
當時耶穌同門徒呢.
係去耶路撒冷.
耶穌係前面行.
門徒跟住係後面.
耶穌叫jor 十二使徒黎.
將自己將要遭遇ge 事呢.
同佢地講.
佢講得好清楚.
唔係呢轉彎.
唔係拐個彎抹角咁樣.
又唔係lor 隱喻.
佢直情講呢.
la 睇下我地上耶路撒冷去喎.
人子將被交給祭司長同埋文士.
咁祭司長同埋文士呢.
就係一班宗教ge 領袖.
原來呢.
耶穌將要被捉拿.
被交比宗教領袖.
呢一班ge 宗教領袖呢.
佢地要定佢死罪.
又要交比外邦人.
跟住三十四節.
講佢地要戲弄佢.
向佢吐唾沫.
鞭打佢殺害佢.
一路由耶穌ge 受苦.

$^{121}$一路講講講.
講到佢ge 係十字架上ge 犧牲.
每一步.
耶穌係呢個位.
已經事先同門徒講.
佢講ge 時候呢.
其實呢D 事情都未發生.
未發生呀.
耶穌已經講定.
而且呢個唔係第一次.
佢係已經係第三次.
佢一路講.
越講就越詳細.
講到逐步逐步佢會經歷D 乜ye .
耶穌知道自己將要經歷ge 事.
所以呢以下將會發生ge 事情呢.
完全唔係意外.
唔係失控.
唔係呢被迫ge .
冇辦法ge 係咪.
唔知點樣呢.
耶穌比人捉拿jor 啦.
於是呢.
被人鞭打.
被人殺害.
呢一切一早.
耶穌都知道.
可以咁樣講.
係係天父ge 計劃裡面.
耶穌講比門徒聽.
要事先同佢地分享.
預備佢地ge 心.
但係門徒有咩反應呢.
我地睇下.
正常呢.
我地聽到咁樣ge 分享.
我同你會有咩反應呀.
耶穌話喎.
我即將會受害啦係咪.
被交比宗教領袖啦.

$^{161}$佢地會咁樣.
將我交比外邦人.
我會受呢D 苦.
最後我仲要被害添.
你作為耶穌門徒你聽jor 呢.
一般黎講應該有咩反應呀.
通常會有咩反應呀.
弟兄姊妹你可以回應下.
至少都講下.
冇啦.
要受咁多苦係咪呀.
可唔可以唔受咁多苦呀.
可唔可以唔好犧牲呀.
係咪呀.
唔好逃避啦.
我地而家唔好上耶路撒冷囉.
係啦走人啦.
正常反應.
我地走路啦 匿埋啦 冇比人捉住.
耶穌我唔想你.
受呢D 苦.
正常係呢個反應係咪.
但你見住呢個對話呢係直接住呢.
就係今日ge 經文喎.
係講緊乜ye 呀.
阿西比太ge 兒子雅各同約翰呢.
對耶穌講.
好似唔係回應緊耶穌講呢.
佢要受苦呀.
受害 佢講第二D ye 咁樣.
佢話老師呢無論求乜ye 呢.
我願你為我地作 求乜ye 呀.
佢地求乜ye 呀.
係你ge 榮耀裡面ge 時候呢.
請你賜我地一個坐你右邊.
一個坐你左邊.
講緊乜ye .
咩叫係你ge 榮耀裡面.
門徒.
諗緊乜ye 呢.

$^{201}$呢個係雅各.
呢個係約翰.
雅各同約翰呢同彼得呢.
一齊呢同耶穌呢.
係係山上呢見證過耶穌基督呢.
係登山.
變像.
見過耶穌基督.
同摩西同以利亞傾計.
係咪.
佢地成個過程.
佢地係係門徒當中呢.
更加親近耶穌基督.
所以係呢一刻呢.
其實佢地求一D ye .
佢覺得我同耶穌.
gwo 個關係再近D .
同佢gwo 個.
情況係唔一樣ge 比起其他ge .
門徒所以我地要求D ye .
求係呢.
耶穌ge 榮耀裡面坐左坐右.
呢個唔係講緊.
屬靈ge 層面喎.
可以.
可以咁講.
門徒完全唔明白耶穌ge 使命.
唔知道耶穌要黎呢個世界做乜.
門徒亦都唔明白.
自己使命.
咁我跟耶穌做乜ye 呢.
係咪我跟耶穌做乜ye .
但係佢有渴望喎.
佢地有渴望喎.
佢唔知自己ge 使命係乜ye .
唔知耶穌係要做乜ye .
我一陣睇下呢個渴望係乜ye .
耶穌呢.
回應佢地.
你地唔知道你求D 乜.

$^{241}$所以佢問呢個問題.
我所學ge 杯.
你地能夠喝嗎.
睇下個問題.
我想飲個杯.
你地能唔能夠飲呀.
知唔知呢個杯代表D 乜ye .
係咪.
呢個杯代表D 乜ye .
門徒以為.
耶穌將要建立地面上ge 政權.
啊.
耶穌將要復興以色列.
耶穌要作我地ge 王.
要係以色列作王.
嘩 犀利呀 點解呀.
之前我一路見到耶穌行神蹟.
一路呢.
佢醫病趕鬼 幾犀利呀.
所導之似D 人跟住耶穌.
嘩 一個咁大能ge 人.
係咪.
上帝所預備ge 拯救我地ge 人.
一定係今日係地面上.
係政權上面.
幫助我地復興以色列國ge 人.
耶穌要作王.
係政治上ge 王.
咁我就一個坐佢左邊.
一個坐佢右邊.
呢個就清晰啦.
門徒有咩渴望.
作王有呢D 好處.
有政治上ge 好處.
權柄上ge 好處.
作王有咩好處呀.
你有權柄喎.
可以管理喎.
可以掌權 可以操控.
個個聽曬你話.

$^{281}$今日如果我係特首.
嘩 幾多人呀.
要跟住我一齊 服務香港.
今日係某某國家ge 總統.
某某國家ge 總理.
嘩 呢個權大囉.
有幾多人呀跟住我呀.
真係好有權呀.
可以管理呢樣管理gwo 樣.
嘩 如果你話係個總統.
個左右手係咪好犀利呀.
佢地係咁諗.
耶穌話咩呀 你地唔知道.
你所求ge 係乜ye  乜ye 叫坐我左坐我右呀.
係咪建立地面上ge 政權呀.
梗係唔係呀 佢問.
我所飲ge 杯 你地能夠飲呀.
其實呢一個問題呢.
第一下呢一個問題呢.
其實已經預設jor 答案喎.
耶穌問呢個位.
對住門徒問ge 時候呢.
可以咁講.
耶穌要飲個杯 門徒唔能夠飲.
耶穌要飲個杯.
門徒係唔能夠飲.
因為唔係門徒諗gwo 回事.
耶穌要飲ge 係苦杯.
佢所飲ge 係苦杯.
唔係而家我有權柄啦.
我做以色列ge 王啦.
我地可以呼風喚雨啦.
我可以叫好多人跟我啦.
命令呢個命令gwo 個 我就操控哂一齊.
係咪呀 唔係呢一隻.
佢飲ge 係苦杯.
佢黎ge 世界ge 使命就要飲呢個苦杯.
係按照天父ge 旨意去飲呢個苦杯.
呢個苦杯指乜ye .
耶穌為罪人承受jor 審判ge 杯.

$^{321}$我好想問 邊個係罪人.
大聲D 講.
我地.
我地 係咪呀.
唔係一個半個.
全世界ge 人.
耶穌係為罪人承受審判ge 杯.
耶穌問ge 問題.
我飲ge  你可唔可以飲.
門徒一定唔得.
即刻我見到有人擰頭.
點解唔得呀.
耶穌你犧牲.
我都跟住你犧牲.
係咪呀.
有乜分別呢 我想問.
點解耶穌飲呢個ge .
要為眾人承受審判ge 杯呢.
佢門徒唔可以飲.
其實其他人都唔可以飲.
點解唔得呢.
其他人都係罪人.
第一個人冇犯罪.
就係耶穌基督.
得佢一個冇犯罪.
我地所有人都犯jor 罪.
呢個杯唔係亞國同約翰所嚮往ge 杯.
唔係個權力ge 杯.
呢個杯亦都唔係亞國約翰.
所可以飲 亦都唔係世界其他ge 人.
可以飲.
呢個杯只有耶穌基督.
可以飲.
耶穌基督.
為我地承受.
一切ge 審判.
仲有係審判.
帶來ge 一切.
ge 惡果.
主耶穌係邊一位呢.

$^{361}$主耶穌.
係神ge 仔.
亦都係人ge 仔.
所以佢係完全ge 神.
亦都係完全ge 人.
佢ge 身份係好特別ge .
佢完全冇犯罪.
係一粒罪都冇犯.
你話 梗係啦 佢神ge 仔.
佢點會犯罪呢.
唔係喎 頭先我介紹佢係咩.
佢亦都係人ge 仔 佢係完全ge 人.
佢曾經係以人ge 樣式.
係地面上生活過.
但係佢完全.
一粒罪都冇犯.
唔止得.
我地世人.
我地係犯罪.
得罪jor 神.
聖經講得好好 羅馬書咁講.
佢話 世人都犯jor 罪.
虧缺jor 神ge 榮耀.
世人就唔係一個半個.
係全世界.
所有ge 人.
今日係呢一度.
每一個 包括我.
企係台上ge 我.
我都犯jor 罪.
虧缺jor 神ge 榮耀.
呢間教會ge 每一個.
我地都犯罪.
得罪jor 神.
我想講呢.
現今ge 世代非常邪惡.
大家最近有冇留意新聞呀.
好震驚呀.
真係一單新聞.
未完全過去.

$^{401}$另一單又起.
首先有肢解案.
跟住呢.
有偷竊殺人案.
有姦殺案.
一路落 一路落.
gwo D 案件呢.
係非常嚴重.
人ge 罪證係越黎越深.
一個簡單ge 偷竊事件.
可以進而成為.
殺人事件.
為錢可以做到幾盡.
親情冇意思.
錢大哂.
人ge 貪婪.
人ge 兇惡.
係我地裡面煩難.
呢D 新聞都仲未完結.
其實都唔係淨係咁少.
係好多好多.
係咪.
係個過程裡面呢.
我地睇到嘩真係人ge 罪證.
係非常犀利.
你話.
我又唔係gwo D 殺人犯.
我又冇去偷冇去搶.
我冇犯法.
我有咩犯法.
我好人好者.
我記得呢.
好多年前.
應該係我讀中學ge 時候.
我都覺得.
我自己係個好人好者.
讀中學ge 時候.
我自己就係學校.
比人地ge 印象.
當然我好安靜.

$^{441}$當時唔係講好多ye .
但係係老師心目中.
我係一個好乖ge 學生.
交足曬功課.
考試成績好優異.
唔會搞破壞.
係學校冇乜是非.
總之好乖.
老師要我做乜我就做乜.
老師係好鍾意.
好鍾意我ge .
人地個個同學地睇我.
呢個係.
好好ge 學生黎ge .
有D 咩會搵我做領袖ge .
阿爸阿媽去睇我.
又有少少唔同ge .
有少少唔同.
但係屋企人睇我.
都唔係完全睇得曬.
我係個點ge 人.
但係我好想同你講.
我知道.
我自己唔係人地所睇ge 乖乖女.
我知道我唔係一個.
學生ge 典範.
唔係某某家庭ge .
女兒ge 典範.
我話返比你聽.
gwo 個真正ge 我.
非常扭曲.
係我心裡面有好深ge 憤怒.
我唔知大家.
有冇經歷過呢個階段.
係憤怒.
我睇返呢D 案件裡面.
除jor 有貪婪.
殘暴仲有憤怒.
係裡面.
我自己ge 憤怒好大.

$^{481}$憤怒到成個人係唔正常.
記得去到中學ge 時候.
當時我爸爸.
就係外面.
係有外遇.
咁佢有外遇.
返到黎一定同媽媽好多爭執.
我知道爸爸有外遇.
係心裡面.
係好傷心.
繼而變成憤怒.
我係度諗.
爸爸你點解傷害我地.
本來好好地.
一個家.
就因為你係出面有其他女人.
又唔係一個.
其實gwo D 女人.
只係貪佢ge 錢.
唔係真係咁愛佢.
但係呢D 女人影響到.
爸爸.
亦都影響我地屋企.
屋企支離破碎.
就黎要分裂.
就黎爸爸媽媽就黎要分開.
一路嘈嘈吵吵.
我發覺.
我一生人唯一覺得最重要.
要保護我個機制.
無jor 就係屋企.
我諗下.
我仲可以係邊度.
搵到保護呢.
係邊度搵到我肯定.
支持呢 無囉.
跟住我好嬲.
我嬲到.
好苦毒.
我望住我爸爸.

$^{521}$我唔想同佢傾計.
我唔望住佢.
我如果要傾計.
我講D 尖酸黑薄ge ye .
呢個就係.
學校以為ge 乖乖女.
我真係講D 尖酸黑薄ge ye .
我望住佢直情好嬲.
好嬲佢.
跟住我媽開始發覺.
哇 糟了.
我女開始唔正常.
開始gwo D 脾氣.
成個人好唔妥.
有D ye 當時屋企人都唔知.
我後黎gwo 陣時間.
差唔多接近我要會考ge 時間.
我gwo 陣時間.
差唔多接近我要會考ge 時間.
差唔多接近會考.
同學地個個好比心機.
搵書.
要做好D .
測驗.
希望準備迎接.
公開試.
大家有個期待我地考好試.
lor 好成績上大學.
跟住第二日就搵工.
過一個好ge 生活.
至少正面D 黎講.
有一個好ge 人生.
你估我點諗.
我都好比心機讀書.
好好比心機.
好瘋狂地好比心機.
你估我點諗.
係咪為jor 搵份好工.
賺大錢好好地生活.
係咪咁呢.

$^{561}$唔係.
個憤怒太大.
我當時心裡面諗.
我阿爸你睇住你.
我考上大學之後.
我會好好報復你.
你睇住你.
我會好好地考上大學.
你就知死.
咁你問我.
我當時諗住點樣報復阿爸呢.
我未諗到.
就好深ge 憤怒.
我要對付佢.
所以呢成個人係唔正常.
我唔係話為jor D .
自己將來ge 好處.
而係我更加有能力.
對付佢.
真真真真憎死佢.
受恨可以好深.
你話係呀.
個位我未犯法ge .
為jor 係地面上ge 法律黎講.
係咪冇人會拉我.
人地都唔可以話.
你犯法你有問題.
我知道自己ge 問題.
上帝都睇到.
呢份憤怒幾大.
我可以咁講.
如果冇法律呢.
我下宣佈.
香港冇法律囉.
你估我會做咩事.
好恐怖.
我裡面.
gwo 種報復念頭就出黎.
可能就變成一種暴力.
即刻出黎冇法律.

$^{601}$之前呢.
仲收係心裡面ge .
憤死佢ge .
而家就發出黎.
狠狠ge 對付佢.
係呢個階段一路.
我都係未信耶穌.
有一次有個同學.
走黎同我講.
佢話曾文.
我全名叫曾文.
你信耶穌.
如果你唔信耶穌.
你落地獄.
你知啦係咪呀.
我地唔信耶穌.
第一次聽就覺得呢D 說話好硬意.
我經歷過呢個階段.
覺得好硬意.
咩意思je .
好人好者我覺得自己唔知幾好.
有咩問題呀.
我做咩衰ye 呀.
做咩要落地獄呀.
gwo 次真係好唔舒服.
聽完之後算數.
我就唔睬佢.
後黎呢我同呢一個ge 女同學.
一齊考入中文大學.
考入中文大學之後呢.
明顯大家知道佢同我全呼音.
佢係一個基督徒.
佢又唔放棄喎.
大學一年班呢就黎結束ge 時候呢.
學校舉辦一個報道會呢.
叫Commons.
來和漢.
我同學呢有邀請我去.
佢唔放棄.
再帶我去睇.

$^{641}$就過一晚.
我真正聽到.
關於耶穌ge 事情.
耶穌係邊位呢.
阿神ge 仔黎架喎.
耶穌黎到呢個世界喎.
為jor .
我地ge 罪呢.
犧牲喎.
嘩真係.
耶穌係咁寶貴.
更寶貴呢.
係耶穌.
同我地一同生活.
係好親密.
上帝同人可以咁親密.
我唔需要自己.
係度苦苦掙扎.
自己點樣努力.
去做D 咩靠自己.
也可以靠呢個上帝.
當晚心裡面好感動.
好感動呢.
我就相信jor 耶穌.
我相信.
耶穌之後呢.
就開始祈禱.
就開始.
祈禱 我為乜ye 祈禱呢.
為我屋祈.
一路祈一路祈.
當中經歷好多艱難ge 時間.
我諗我都祈jor 成.
兩年到至少兩年到.
都試過呢.
點樣祈禱呢.
係上帝面前呢 喊住口.
講都講唔到ye 就咁喊.
一路祈禱.
神改變我地屋企.

$^{681}$首先神改變我.
大家估下呢.
我係大學有冇報復我爸爸呢.
結果係冇ge .
因為上帝已經改變jor 我.
佢寬恕jor 我ge 罪.
我亦都要去寬恕.
我爸爸.
唔可以咁樣ge .
我希望上帝寬恕jor 我啦.
係咪.
咁我又唔寬恕我爸爸.
唔可以 所以佢寬恕jor 我.
我就要寬恕我爸爸.
以前呢.
我憎死佢 我好想報復.
做D 衰ye .
但係信耶穌.
上帝改變我.
讓我愛佢.
寬恕佢.
我地呢 我想像唔到.
之前自己一路做唔到.
但係我想像唔到呢.
竟然我可以同佢又講又笑.
開心起黎.
好似朋友咁樣.
想像唔到.
跟住我見到.
上帝改變呢.
爸爸媽媽ge 關係.
爸爸媽媽.
復和.
再下一步我見到.
上帝改變.
爸爸.
佢停jor 呢D .
唔正常ge 唔好ge 關係.
佢停止.
咁樣同呢D 女人一齊.

$^{721}$去改變jor .
我信耶穌.
我返返教會.
媽媽係比較早ge 時候.
信耶穌.
但係一段時間冇返教會.
當我返教會之後.
媽媽跟住我返教會.
一步一步到家姐信主.
之後我地一齊.
好同心.
為爸爸信耶穌祈禱.
我唔係淨係成為我爸爸ge 好朋友.
我要將最好ge 福氣.
比佢.
因為我真係愛佢.
由憤怒到愛.
中間係上帝.
數不盡ge 工作.
好愛佢.
我成日為佢祈禱.
求主幫助佢 幫助佢信主.
我成日同人講為我爸爸祈禱.
我希望佢信主.
好多人為我爸爸祈禱.
上帝好多好多ge 工作.
到最後.
我爸爸真係信主.
到我爸爸離開世界前.
佢已經相信耶穌.
得夢救贖.
我地成家人.
都信jor 耶穌.
係好開心.
個關係唔再一樣.
過去有好多眼淚.
但係自從耶穌.
黎到我地屋企.
就唔同曬.
好溫暖.

$^{761}$好多個愛.
互相扶持.
我記得當我係信jor 耶穌之後.
唔單止我同爸爸ge 關係好jor .
係爸爸係越來越接納我ge 信仰.
越來越支持我ge 事奉.
記得我第一次去蒙古.
短孫ge 時候.
我父親就話.
我第一次去蒙古短孫ge 時候.
係全家人.
一齊幫我去預備.
我好感動.
爸爸gwo 時未信耶穌.
但係佢好支持我.
係一幅好靚.
好靚ge 圖畫.
係呢個過程.
就係主耶穌ge 恩典.
係我地屋企去彰顯.
如果唔係主耶穌救我地.
我.
個下場係點樣.
聖經呢度講.
人人都有一死.
我地每個人都會有.
一個死亡就係我地肉身ge 死亡.
死jor 之後.
且有審判.
我地面對ge 係面對上帝ge 審判.
呢個上帝係.
完全聖潔公義ge 上帝.
係完全唔會犯罪ge 上帝.
係完全聖潔公義.
佢會審判我地.
愛我地ge 一生.
我地一生ge 所言所行.
全部記載係行為冊上面.
有D ye 人地見到ge .
人地聽到我地講個說話.

$^{801}$睇過我地點樣行事為人.
係咪都知道ge .
但係有D ye 行為冊上記低ge ye 呢.
係人地見唔到ge .
上帝見到.
全部都記低曬.
你話得唔得人驚.
係所有.
今日有邊個知道我同你ge 所有.
心思意念裡面.
每一個意念上帝都知道曬.
係行為冊上記低ge .
如果按照行為冊咁樣去判斷我地.
我心想我地冇人可以.
逃脫.
如果唔係.
聖經都唔會話世人都犯jor 罪.
每個人都犯罪.
就因為咁仔細ge 野性.
即係係神gwo 度ge 行為冊記低曬.
我地死之後.
要面對審判.
按照咁ge 審判.
羅馬書咁講.
佢話罪ge 公間難.
係死.
但呢個死就唔係第一次ge 死.
即係肉身ge 死亡係第二次.
係永遠ge 死亡.
我永永遠遠同上帝.
分開.
同上帝分開係咩光景.
你睇呢幅圖畫.
一邊.
左上角好似金光閃閃.
好靚.
gwo 個係永生ge 路.
係上帝ge 國度.
神ge 家.
係永生ge 國度.

$^{841}$係好靚ge 金碧輝煌.
但你見到呢.
十字架下面gwo 邊係乜ye 黎架.
地獄.
罪人呢.
結局本來就係gwo 度.
gwo 度呢係同上帝分開ge 地方.
同上帝分開jor .
我地要為自己ge 罪.
承擔返ge 後果.
同上帝ge 分開係永遠ge 分開.
冇盼望冇愛.
只有痛苦.
而且係永恆ge 痛苦.
罪ge 功架就係咁.
犯罪就有後果.
原本我呢個人地體為乖乖女.
實際我裡面充滿罪惡.
我記得我信耶穌之後呢.
我睇聖經呢.
越睇我就越驚.
以前覺得自己好巴閉.
我有咩問題呀.
我唔知幾好人.
有咩問題呀.
你話我好人就真.
我越睇聖經越驚.
呢D 係罪.
gwo D 係罪.
我都有.
我有呢D 罪.
我好多罪.
越認就越多罪.
越係神面前覺得我真係唔掂.
如果耶穌唔救我.
真係係下面.
永遠永遠同上帝分開.
罪唔係淨係講緊殺人放火.
我地心思意念.
幾多ge 貪婪幾多ge 妒忌.

$^{881}$我地ge 嘴巴幾多ge 是非幾多ge 傷害.
我地好多ge 憤怒.
好多ge 憎勁.
嫉妒.
好多ge 惡.
背叛神.
背叛人.
講大話從來都唔需要人教.
講粗口浪浪上口.
即刻笑.
唔需要人教.
聽多多更加容易上口.
幾多罪.
其實我地好污糟.
邋遢.
本來我地要面對審判.
你地睇返其實左邊gwo 條路.
十字架帶去gwo 條.
gwo 條路係窄ge 路.
好細條路.
呢一邊gwo 邊ge 路係闊ge .
多人行.
十字架ge 路卻係少人行.
但十字架ge 路.
先至帶比人.
有永恆ge 盼望.
先至可以返上帝gwo 度.
返父神gwo 度.
但gwo 條闊ge 路.
卻係帶我地走向滅亡.
因為呢條闊ge 路容易行.
好多誘惑.
亦都好多人行係裡邊.
耶穌飲個苦杯係為jor 乜ye .
做乜ye .
佢要被.
佢ge 門徒出賣.
要被捉拿.
面對唔公義ge 審判.
被鞭打.

$^{921}$最後被釘係十字架上.
受咁多苦咁大ge 犧牲.
佢飲呢個苦杯.
為jor 乜ye .
就為jor 救我同你.
等我同你.
可以同天父爸爸復和.
等我同你唔想再受.
下面呢一個永恆ge 折磨.
痛苦.
煎熬.
我同你呢.
根本冇辦法處理罪ge 問題.
我盡量學乖D .
你唔識諗.
中學時候你憤怒.
你變得乖D 咪得.
就係唔得.
我真係試過係唔得.
我個人自己點乖都乖唔起.
因為我真係有罪性.
我裡面充滿jor 罪污.
係唔得就唔得.
我呢頭可以做一件好事.
就係唔得.
我gwo 陣時.
中學時候.
讀一個佛教學校.
當時教我地.
點樣去修煉自己.
達到境界.
遠離呢D .
私慾困擾.
其實佢地一樣都承認.
處理唔到.
人ge 私慾好深.
你靠自己.
永遠都唔得.
唯有靠耶穌基督.
耶穌基督為我同你犧牲.

$^{961}$為我同你受苦.
飲jor 呢個苦杯.
佢飲jor 呢個苦杯.
面對上帝.
對佢ge 審判.
佢冇犯罪.
佢代我同你受jor 呢一個.
最終極ge 審判.
同埋之後所帶黎ge 惡果.
佢就係死.
係十字架上.
佢代我同你受jor 之後.
事情就有改變.
呢個係唯一ge 方法.
好記得.
係耶穌受害ge gwo 一日.
從淨午到下午三點鐘.
遍地都黑暗jor .
好黑暗.
耶穌最後ge 時候呼喊.
以利.
拉瑪撒巴.
國大尼.
我ge 神.
點解離棄我.
大家知唔知到.
佢所受ge 苦係幾大.
全世界ge 人ge 罪.
都落係一個人身上.
唔係淨係殺人放火.
我同你ge 罪.
落哂係gwo 度.
紮jor 係耶穌身上.
全世界ge 人ge 罪.
gwo 一下.
耶穌由早上九點被釘.
到新初即係下午三點.
一共六個鐘.
係個十字架上.
前三個鐘頭就受人迫害.

$^{1001}$機潮ge 苦.
後三個鐘頭.
佢面對上帝gwo 個審判.
離棄ge 苦.
唯有佢一個可以飲呢個腐杯.
唯我同你飲ge .
單單.
唯我同你飲.
當然呢一個唔係佢ge 終局.
因為耶穌.
完全冇犯罪.
係佢被釘.
十字架之後第三日.
佢從死裡復活.
佢勝過死亡.
更加證明佢係一個義人.
佢完全冇犯罪.
當然我地知道.
係上帝令到佢復活.
你要知道一個犯罪ge 人.
死jor 就真係死.
佢唔能夠救人.
但係呢一個.
係耶穌基督佢仍然能夠復活.
除jor 話上帝叫佢復活.
亦都係因為佢係一個義人.
佢冇犯罪.
所以佢為我同你.
所成就ge ye 就有效.
佢為我同你.
飲個苦杯.
就有好大ge 意義.
可以救我同你條命.
講白D .
本來我同你就冇命.
唔係淨係簡單.
為罪受懲罰.
係我同你將會冇命.
上帝因為愛我地.
唔捨得就派佢個仔.

$^{1041}$為我地飲落呢個苦杯.
承受呢一切.
我同你先至可以.
走向永生.
復活節.
係講緊得勝.
耶穌基督ge 得勝.
上帝ge 得勝.
亦都係我同你ge 得勝.
但個關鍵.
係要係主裡面.
係要相信耶穌基督.
有兩個層面.
第一個層面.
哥林多後書.
五章十七字gwo 度咁樣講.
所以若有人在基督裡.
家就是新造的人.
舊事已過都變成新的了.
天父應承我地一件事.
只要我地願意相信耶穌基督.
相信佢.
口裡面承認.
心裡面相信.
真係真心地相信ge 時候.
神就.
抹曬我地ge 過去.
佢話舊事已過.
都變成新.
無論我舊時係一個怎樣ge 人.
都成為過去.
上帝寬恕我地一切ge 罪惡.
我地係一個新ge 人.
我地係天父ge 子女.
有耶穌基督ge 生命.
係我地裡面活出黎.
我地可以靠住主.
勝過罪惡.
靠住主過一個唔犯罪ge 生活.
新ge 生命.

$^{1081}$一路跟住耶穌基督.
走向永恆.
唔需要再害怕.
永死ge .
那個捆綁.
耶穌ge 苦杯.
只有祂一個人能夠飲.
呢一個苦杯就係為全世界ge 人.
ge 罪而飲ge 苦杯.
得祂一個可以飲.
第二個好特別ge 地方.
你又見到.
耶穌轉個頭又話.
我所飲ge 杯.
你地要飲.
你飲ge 苦杯.
我地可以飲ge 咩?.
呢個杯有D 唔同.
呢個杯唔係耶穌自己單獨飲.
gwo 個杯.
gwo 個杯得祂自己可以飲.
但呢個杯.
耶穌想佢ge 門徒都去飲.
我所受ge 洗.
你地都要受.
我所喝ge 杯.
呢個杯係咩杯呢?.
呢一個杯就係門徒.
要捨己犧牲.
謙卑服侍人ge .
以至於死ge 一個杯.
但係gwo 個.
唔係為人地.
為jor 人地ge 罪.
係咪?為jor 人地ge 審判.
要去承擔人地生命gwo 個杯.
而係一個.
跟從耶穌基督服侍主.
服侍人要謙卑捨己.
犧牲ge 杯.

$^{1121}$呢個耶穌要我地咩?.
要我地承受.
係佢ge 門徒就要去.
約翰雅各.
要飲ge 係飲呢個杯.
為jor 服侍人.
謙卑自己.
犧牲捨己.
耶穌黎到世界上ge 使命.
係呢個經文有講.
唔係要受人ge 服侍.
乃係要服侍人.
佢謙卑捨己.
係會犧牲.
佢希望佢ge 門徒都係咁樣犧牲.
係好唔同架喎.
gwo 個意義上係咪?你見到呢.
咁你究竟係邊一樣呢?係咪同一樣ye 呢?.
唔係同一樣ye 黎ge .
呢兩個講呢.
即係我要喝個杯呢,唔係同一樣ye .
好啦,咁但係呢.
至於坐左坐右呢.
上帝ge 仔呢,能夠坐佢ge .
左邊右邊係好榮耀ge .
但呢一個呢,就唔係在乎今日我做幾多ye .
唔係我地自己決定.
而係呢,即係神去決定.
邊個坐佢左坐佢右.
主耶穌呼召門徒要咁樣做.
謙卑捨己.
服侍人.
甚至係付上生命ge 代價.
之所以我地要咁樣生活係原因係乜ye ?.
你話俾我聽.
點解門徒要咁樣生活?.
因為乜ye ?.
因為耶穌為我地.
飲jor gwo 個苦杯.
佢配得我地.

$^{1161}$咁樣為佢飲.
我地gwo 個苦杯.
配得我地咁樣去服侍佢.
我地咁樣去服侍人,亦都係理所當然ge .
今日呢.
弟兄姊妹呢,係呢度呢,趁住呢個機會.
我要去問.
你係咪真係明白,耶穌幾多為我地付上jor 乜ye ?.
我想問下弟兄姊妹.
今年呢,唔知係你人生.
第幾個復活節呢?.
每一年重複又重複.
聽受苦節ge 訊息.
復活節ge 訊息,每一年呢.
我都可能覺得好感動.
我要立志,立志都唔知十年呢.
廿年呢,三十年呢.
我地做jor D 乜ye ?.
問下自己.
我係咪謙卑捨己.
服侍人.
以至於犧牲.
我都願意.
為主緣故.
我想問大家,係咪咁?.
係咪跟住?.
定係我好似門徒一樣.
我期望呢,主呀,跟你就有好處.
我明唔明白ge ?.
咩叫跟住?咩叫作佢個門徒?.
係咪真係明白?.
問下自己每個禮拜返黎係為jor 乜ye ?.
呢度有冇好處?.
有冇ye 玩?.
有冇ye 食?.
有冇D 開心ye ?.
信耶穌係咪一回咁ge 事?.
作主門徒係咪一回咁ge 事呢?.
D 門徒跟住佢三年都仲唔明白.
跟住佢三年,睇佢行神跡奇事.

$^{1201}$睇佢走遍各城.
各鄉,去傳天國.
福音,睇佢去幫咁多人.
醫病,趕鬼,服侍人.
見到耶穌咁疲倦.
咁努力擺上自己服侍.
三十幾歲人離世.
樣子似五十歲.
因為好勞苦呀耶穌.
門徒跟佢身邊咁耐,咁貼身跟佢都唔明白.
以為跟住耶穌可以作王.
真係好威呀.
有權呀,可以掌控呀.
都唔明白.
我地明唔明白?.
佢唔會咁多年,係咪知道返黎做乜?.
咩叫跟耶穌?.
聽姐妹,所以我想今日邀請大家去.
再一次諗返自己ge 信仰.
係咪真係跟耶穌?咩意思呀?.
今日會唔會有一個機會可以重新立志?.
如果你真係跟耶穌,你點回應佢?.
點跟?.
唔可以講完之後,又冇jor 件事.
下年又黎過,又講呢個訊息.
下年又重新黎過.
後年又係咁,每年都係重新黎過.
今年係咪要有行動呢?.
而對於今日arm 黎ge 新朋友呢.
卻係想問大家.
耶穌基督已經為我地上進苦杯.
受盡苦.
同埋犧牲.
大家係咪願意.
將生命交比耶穌基督?.
將自己生命交比佢.
邀請耶穌入到自己生命裡面.
成為我地ge 救主.
同埋生命ge 主.
我地係咪願意呢?.

$^{1241}$我想而家一同唱一首詩歌.
一路唱ge 時候.
我會先發出.
第一個呼召.
全新的你呢首歌.
我覺得好配合.
今日ge 天氣.
我地唱之前.
睇一睇個歌詞.
佢話陰天代表我地ge 心情.
雨天更加係我地對生命ge 反應.
今日真係好似朝頭早有D 雨.
都有D 陰天.
唔知係咪我地ge 心情.
都係咁樣有D 受到影響呢?.
會唔會每日生活都係咁?.
係咪呀?.
甚至乎有時覺得呢.
冇人了解我地.
好似好多艱難.
冇力.
我想講黎到耶穌呢一度.
耶穌可以叫一切更新.
耶穌可以改變一切.
我地先唱一詞.
第二次ge 時候.
我會發出呼召ge .
佢話陰天代表我地ge 心情.
雨天更加係我地對生命ge 反應.
佢話陰天代表我地ge 心情.
雨天更加係我地對生命ge 反應.
佢話陰天代表我地ge 心情.
雨天更加係我地對生命ge 反應.
佢話陰天代表我地ge 心情.
雨天更加係我地對生命ge 反應.
佢話陰天代表我地ge 心情.
雨天更加係我地對生命ge 反應.
佢話陰天代表我地ge 心情.
雨天更加係我地對生命ge 反應.
佢話陰天代表我地ge 心情.

$^{1281}$雨天更加係我地對生命ge 反應.
佢話陰天代表我地ge 心情.
雨天更加係我地對生命ge 反應.
佢話陰天代表我地ge 心情.
雨天更加係我地對生命ge 反應.
佢話陰天代表我地ge 心情.
雨天更加係我地對生命ge 反應.
佢話陰天代表我地ge 心情.
雨天更加係我地對生命ge 反應.
佢話陰天代表我地ge 心情.
雨天更加係我地對生命ge 反應.
佢話陰天代表我地ge 心情.
雨天更加係我地對生命ge 反應.
佢話陰天代表我地ge 心情.
雨天更加係我地對生命ge 反應.
佢話陰天代表我地ge 心情.
雨天更加係我地對生命ge 反應.
佢話陰天代表我地ge 心情.
雨天更加係我地對生命ge 反應.
佢話陰天代表我地ge 心情.
雨天更加係我地對生命ge 反應.
佢話陰天代表我地對生命ge 反應.
耶穌愛你,耶穌等你,耶穌等著我,一個全新的你.
好,再唱一次的時候呢,我想呢,邀請當中.
領唱單,我們唱的時候呢,想邀請大家.
今天第一次來到我們當中的新朋友.
在剛才呢,其實我們已經一同看到.
耶穌基督是怎樣為我們的罪去犧牲,捨忌.
是為了我和你,是喝了個苦杯.
今天呢,如果在當中有願意去相信耶穌基督的.
就說願意邀請耶穌基督進到我們生命裡面.
成為我們救主和生命的主的.
請你在台下舉一舉你的手.
如果你有願意,將你的生命交給主.
主已經為我們去犧牲.
今天只要相信祂,主就說呢,天父就說.
舊事已過都變成新的.
去歡喜我們過去的一切.
歡喜我們過去的一切.
舊事已過都變成新的,我們今天是新的人.

$^{1321}$耶穌代我和你喝了個苦杯.
我們不需要再害怕,再覺得人生是沒有盼望.
無論今天是陰天還是晴天.
人生都有盼望,是永恆的盼望.
如果有願意的,誠意邀請大家舉起你的右手.
當中有願意相信的朋友.
請你舉起你的右手.
如果有願意的朋友,誠意邀請大家舉起你的右手.
主耶穌很愛我們.
祂為我們犧牲祂的生命.
有願意的朋友,請你舉起你的右手.
期待有人能夠了解你心.
能夠愛你,持你力量的心.
如果有願意的朋友,請你舉起你的右手.
向上帝表明,主呀,我願意相信你.
請你進入到我的生命裡,成為我的救主和生命的主.
體會你的心情.
耶穌能夠改變你的曾經.
耶穌愛你,耶穌疼你.
耶穌能夠找到一個全新的你.
在這裡,我很想邀請弟兄姊妹.
如果今天我們聽到耶穌基督為我們喝了苦杯.
祂也很渴望我們跟從祂喝了祂的杯.
這個杯是一個願意謙卑服侍別人.
捨己犧牲的杯.
耶穌怎樣做,我們就跟著祂怎樣做.
如果有願意的弟兄姊妹.
我反而想邀請大家從你的座位站起來.
從你的座位站起來向上帝表明.
主耶穌,我知道今天我再次聽到.
你為我承受的苦杯是這麼大的.
主,我願意跟從你,去喝你的杯.
如果有願意的弟兄姊妹.
請你們從你們的位置站起來.
向上帝表明你們的心意.
有願意的弟兄姊妹.
主耶穌,我回應你.
我願意將生命交給你.
我願意去回應你.
我願意跟從你.

$^{1361}$你這樣去嘗苦杯,去犧牲.
雖然我不能為萬民的罪去做什麼.
但我卻願意跟從你的腳踝去服侍人.
有願意的弟兄姊妹請你起身.
請你起身.
我在這裡很想請領唱再唱最後一次.
大家先不要坐下.
最後最後一次的邀請就是.
今天新朋友你看到.
我們有很多已經跟從耶穌的朋友.
他們站了起來.
說他們要跟從上帝.
今天如果你也願意跟著這些已經信耶穌的朋友.
一同信耶穌的請你全部站起來.
請你全部站起來.
跟著這些已經信耶穌的朋友.
我們再唱一次這首詩歌.
一同唱.
你說明天代表你的心情.
明天更是你對生命的反應.
你說明天生活一樣平靜.
對於未來沒有一點信心.
親愛的朋友你是否曾經.
期待有人能夠瞭解你心.
能夠愛你賜你力量更新.
耶穌能夠讓一切都更新.
耶穌能夠體會你的心情.
耶穌能夠改變你的曾經.
耶穌愛你 耶穌疼你.
耶穌能夠一個全新的你.
我請詩琴繼續彈這首詩歌.
我看到其實新朋友都站了起來.
我不知道大家是否已經決定跟從耶穌基督.
如果有願意來到前面.
去見證上帝恩典.
想邀請大家在朋友的陪伴下.
站到台前.
我們一同向上帝述說大家的心意.
上帝我願意相信你.
如果剛才大家都跟著我們現有的基督徒.

$^{1401}$大家可以來到台前.
我們一同向上帝表明心意.
我們身邊的基督徒朋友都會陪伴你.
我們都會陪伴你.
你不會孤單.
親友們陪伴他們一起上來.
在一個很重要的時間.
我們向上帝表明心聲.
上帝我要歸向你.
你為我已經付上這麼大的代價.
我要將我的生命交給你.
從今以後只為你而活.
還有沒有親友陪伴的.
一同出來不用擔心.
這裡的基督徒都向他們表明心意.
他們要跟著主.
從今以後的路是為主而活.
你有他們陪伴身邊.
你不會孤單.
如果還有的.
想邀請你們出來.
有親友陪伴的不用擔心.
更加有耶穌基督在這裡和我們一起.
我很想帶領台前的朋友.
有一個決志的禱告.
然後我想請周牧師.
為新朋友或是站立的.
這麼多的弟兄姊妹有一個祝福.
我先想說一句.
大家說一句好不好.
我們一同有一個決志禱告.
決志禱告就是說.
我將我的心聲和上帝陳明.
從今以後我真的跟上帝.
耶穌基督是我的救主和主.
我說一句大家說一句.
親愛的天父.
猜派你的獨生兒子.
耶穌基督.
來到這個世界.

$^{1441}$為我的罪.
死在十字架上.
我承認.
我犯罪得罪了你.
我不可以靠自己.
處理罪的問題.
主啊.
我相信.
你在十架上的犧牲.
是能夠為我成就救恩.
我很相信.
耶穌基督的寶血.
可以洗淨我的罪惡.
所以這一刻.
我邀請你.
主耶穌.
進入到我生命當中.
成為我的救主.
和生命的主.
從今以後.
我只是跟從你一個.
我要切實的認罪悔改.
要按你的心意而行.
求主兼顧我.
讓我永遠都不會離開.
願你透過我的生命.
彰顯你的榮耀.
奉主名求.
阿們.
大家在眾人面前.
作了這個很美好的決志.
一個立志.
我將時間交給周牧師.
為新朋友去祝福.
亦為我們弟兄姊妹的立志去禱告.
請我們領受.
從上主而來的賜福.
願主耶穌基督的恩惠.
天父上帝的慈愛.
聖靈的感動和契通.

$^{1481}$常與大家同在.
特別今早我們能夠聽到.
主耶穌為我們世人.
為我們的罪.
改變我們的生命.
使我們已經不能改變的事實.
可以在你的生命裡.
可以來到有新的開始.
新的生命.
也站在台前的.
有願意決心跟隨你的新朋友.
成為基督徒.
但願主耶穌在他們的生命裡.
和我們每一位站立.
決心跟隨主的弟兄姊妹的面前.
我們能夠有聖靈的更生.
我們的生命能夠有新的開始.
從今時.
直到耶穌基督榮耀的再臨.
阿們.
謝謝.
\newpage



\section{路加福音 11:1-13}
\label{sec:b_zpDLNhNUk}
\textbf{2023.04.16 《宣教任務( 四) :一切從禱告開始》 路加福音 11\_1-13 ( 新漢語譯本) 楊穎兒傳道}
\newline
\newline
連結: \href{https://youtube.com/watch?v=b-zpDLNhNUk}{\texttt{https://youtube.com/watch?v=b-zpDLNhNUk}} ~~~~ 語音日期: 2023-04-19
\newline
\newline
\hyperref[sec:AV2ZegtPBEc]{\small{< < < PREV SERMON < < <}}
~
\hyperref[sec:index_chronic]{\small{[返順時目]}}
~
\hyperref[sec:index_scriptual]{\small{[返順卷目]}}
~
\hyperref[sec:c61oKrE6RXs]{\small{> > > NEXT SERMON > > >}}
\newline
\newline
路加福音 11:1-13
\newline
\begin{longtable}{cl}
\hline
\hline
章節 & 經文 (和合本修訂版)\\
\hline
11:1 & \begin{tabularx}{0.7\textwidth}{X} 耶穌在一個地方禱告。禱告完了，有個門徒對他說：「主啊，求你教導我們禱告，像約翰教導他的門徒一樣。」 \end{tabularx} \\ \\ \relax
11:2 & \begin{tabularx}{0.7\textwidth}{X} 耶穌對他們說：「你們禱告的時候，要說：『父啊，願人都尊你的名為聖；願你的國降臨； \end{tabularx} \\ \\ \relax
11:3 & \begin{tabularx}{0.7\textwidth}{X} 我們日用的飲食，天天賜給我們。 \end{tabularx} \\ \\ \relax
11:4 & \begin{tabularx}{0.7\textwidth}{X} 赦免我們的罪，因為我們也赦免凡虧欠我們的人。不叫我們陷入試探。』」 \end{tabularx} \\ \\ \relax
11:5 & \begin{tabularx}{0.7\textwidth}{X} 耶穌又對他們說：「你們中間誰有一個朋友半夜到他那裡去，對他說：『朋友！請借給我三個餅； \end{tabularx} \\ \\ \relax
11:6 & \begin{tabularx}{0.7\textwidth}{X} 因為我有一個朋友旅途中來到我這裡，我沒有東西招待他。』 \end{tabularx} \\ \\ \relax
11:7 & \begin{tabularx}{0.7\textwidth}{X} 那人在裡面回答：『不要打擾我，門已經關了，孩子們也同我在床上了，我不能起來給你。』 \end{tabularx} \\ \\ \relax
11:8 & \begin{tabularx}{0.7\textwidth}{X} 我告訴你們，雖不因他是朋友起來給他，也會因他不顧面子地直求，起來照他所需要的給他。 \end{tabularx} \\ \\ \relax
11:9 & \begin{tabularx}{0.7\textwidth}{X} 我又告訴你們，祈求，就給你們；尋找，就找到；叩門，就給你們開門。 \end{tabularx} \\ \\ \relax
11:10 & \begin{tabularx}{0.7\textwidth}{X} 因為凡祈求的，就得著；尋找的，就找到；叩門的，就給他開門。 \end{tabularx} \\ \\ \relax
11:11 & \begin{tabularx}{0.7\textwidth}{X} 你們中間作父親的，誰有兒子求魚，反拿蛇當魚給他呢？ \end{tabularx} \\ \\ \relax
11:12 & \begin{tabularx}{0.7\textwidth}{X} 求雞蛋，反給他蠍子呢？ \end{tabularx} \\ \\ \relax
11:13 & \begin{tabularx}{0.7\textwidth}{X} 你們雖然不好，尚且知道拿好東西給兒女，何況天父，他豈不更要把聖靈賜給求他的人嗎？」 \end{tabularx} \\ \\
[1ex]
\hline
\hline
\end{longtable}
$^{1}$弟兄姊妹平安.
今天他也在我們中間.
讓我今天可以在這個主日.
又回來洪恩家.
和大家一起學習神的話語.
今天的經文相信大家可能都聽過.
都是有一個主禱文的部分.
特別是耶穌教導他的門徒.
如何去祈禱.
在禱告之後.
我決定這個講道的題目.
就是一切從禱告開始.
讓我們一起在這個主日.
學習如何去禱告.
讓我們先低頭同心禱告.
親愛的天父.
求你打開我們的心.
也給我們智慧.
去明白你的話語.
透過今天的講道.
求你與我們每一位去說話.
挑戰我們每一位.
做好這個禱告的功課.
在禱告裡面參與天國的事工.
我們這樣同心合意的祈求.
阿們.
我相信大家都有禱告的習慣.
不知道大家有沒有遇過以下三個情況呢?.
第一,懷疑神有沒有聽禱告呢?.
第二,覺得自己祈來祈去.
是不是祈得不好呢?.
第三,感覺祈禱很佛力.
質疑其實祈禱有沒有用呢?.
在我服侍的群體裡.
當我關心他們幾句之後.
他們就會跟我說.
你為我祈禱吧.
通常我都會馬上跟他們做一個禱告.
我相信上帝是閱立我們每一個.
有聲無聲的禱告.

$^{41}$同時我更加相信.
一切都是從禱告開始的.
禱告是我們跟上帝同工的重要一環.
所以我們更加不要輕視禱告的功課.
好好學習禱告,然後去服侍.
最後,在今天的講道.
我會跟大家分享.
有些從禱告發展出來的事工.
根據《路加福音》的記載.
當時耶穌做了很多的服侍.
包括在人群裡.
去港島治病.
探望馬大,瑪利亞.
然後,耶穌不是回家睡覺.
不是坐著休息.
他是選擇去一個地方禱告.
我相信門徒都留意到.
耶穌禱告的時候是有點不同的.
可能門徒看到耶穌禱告.
相信也很想像耶穌一樣.
在禱告中得力.
所以他們特別要求耶穌教導他們禱告.
在今天的經文裡.
第一節.
門徒留意到耶穌禱告.
不要打擾他.
等耶穌一禱告完之後.
就提出要耶穌教導他們禱告.
就像我們在今天繁忙的生活裡.
在這個城市很忙碌.
我們禱告的時候.
真的需要靜下來.
在禱告中學習得力.
在今天的經文裡.
耶穌不單是教導禱告.
究竟要怎樣去禱告.
他還教導門徒.
你禱告的對象.
禱告的態度.
和禱告的目的.

$^{81}$我們要從這三方面學習.
怎樣禱告才能到位.
第一.
關係要到位.
耶穌教導我們禱告是要到位的.
要怎樣到位呢?.
首先要意識到.
我們不是和一個很遠很遠的神去禱告.
我們向神禱告是和我們有關係的.
為什麼這樣說呢?.
其實在今天的經文裡.
父親和兒女的關係.
其實是貫穿經文的.
可以說是很突出的圖畫.
我們一起看看第二節.
耶穌對他的門徒說.
你們禱告的時候要說父啊.
父啊.
這是我們的天父.
我們的爸爸.
然後到了第十一節.
耶穌又用了一個父親和兒女的比喻.
去表明天父的心.
耶穌基督是上帝兒子的身份.
就是讓我們得著耶穌的時候.
我們可以透過耶穌基督成為弟兄姊妹.
我們以弟兄姊妹相稱.
在上帝的家裡.
我們就成為了他的兒女.
我們都是上帝的孩子.
這個關係是很親密的.
孩子和父親聊天.
是一個很親密的時間.
就好像我回家的時候.
我的小朋友跑來跑去.
他們就會跑來抱著媽媽.
很有趣的.
跑來看著我.
然後在笑.
走過來跟我聊天.

$^{121}$抱著很親密.
我們隨著時間的流逝.
有時候我們不單止忘記了自己.
幼兒的時候是怎樣黏著父母.
是怎樣完全依賴他.
信靠我們的父母.
而且隨著時代的變遷.
我們亦都會忘記了.
可能自己也有小朋友的時候.
那種曾經擁有過的親密的關係.
到了我們成為基督徒的時候.
很多時候我們又忘記了.
我們就是天父的兒女這個身份.
我們可以走到我們的神面前.
我們叫祂爸爸.
我們叫祂天父爸爸.
一個到位的禱告.
首先就是要搞清楚.
我們是向我們的父祈禱.
所以當耶穌教導門徒禱告的時候.
都是用父親,兒女這些做例子.
我們看看第十一節.
你們中間有作父親的.
有誰會在兒子要魚的時候.
將蛇給他呢?.
或者要雞蛋的時候.
將蠍子給他呢?.
這裡是什麼意思呢?.
其實這裡有兩種對比.
一就是將魚和蛇.
另外一個就是蛋和蠍子.
其實這裡是說兩組的動物.
有一些魚看起來像蟬.
其實好像蛇.
一個不熟悉魚,蟬的人.
很容易將牠們誤會.
或者將蠍子身體的翼.
其實有些蠍子的形狀.
尤其是白蠍子.
其實會誤以為是雞蛋.

$^{161}$切開裡面其實是有些白色,黃色.
所以很相似.
說的是這兩樣東西.
其實外表看起來很相似.
所以放在一起說.
簡單來說就是地上的父親.
你怎麼不好都不會將不好的東西.
給自己的兒女.
即是魚和蛇分不開.
或者蛋和蠍子會混淆.
不會搞錯.
但我們天上的父親.
更加不會將這些東西搞錯.
所以子女向天父祈求的時候.
天父怎麼會不聽呢?.
天父怎麼會供應一些.
邪惡的東西給子女呢?.
其實早在這段經文之前.
第十章也說到.
「漢啊,我已經賜給你們權.
使你們能夠踐踏蛇和蠍子.
並且能夠勝過一切仇敵的能力.
絕對不能傷害你們」.
這裡說的是.
這些邪惡的力量.
都不會影響到天父的兒女.
當我們懷疑神有沒有聽我們禱告的時候.
或者禱告時覺得.
都沒有帶來祈求的改變.
但我們要記住.
天父其實時刻都在保護我們.
給我們的話不代表愛.
不給我們的話不代表不愛我們.
所以我們向天父禱告的時候.
我們是在信靠這個父.
能夠保護我們.
內在的光景.
第二,我們禱告的時候.
內容要到位.
耶穌跟他們說.

$^{201}$你們禱告的時候要說父.
然後就是內容.
這個內容要到位.
禱告是有目標的.
「願你的名被尊為聖」.
同時禱告是有願景的.
「願你的國來臨」.
透過禱告我們是在參與.
神的國降臨的一件事.
在上帝的事宮裡.
我們是可以透過禱告去參與.
禱告也是提醒我們.
要求去經驗神的那種供應.
求你按日次給我們.
每天所需要的食物.
我們看到我們不是為自己求的.
我們是為整個群體.
整個人類的發展去求.
而且我們不是每月頭求月尾那些.
我們是要每天去求.
每天去禱告.
其實上帝,其實天父.
是想我們跟祂有一個很親密的關係.
就是每天,每天地去禱告.
去向父祈求.
將我們的心思意念都教堂.
禱告也是提醒我們要復和.
我們要跟上帝復和.
跟人復和.
求你赦免我們的罪.
正如我們也赦免那些.
虧負我們的人.
修補我們跟人和天父的關係.
耶穌基督為我們死.
這個拯救的工作其實已經完成了.
上帝已經視我們為無罪.
但在現實的生活中.
我們的確會被很多試探.
或者陷入罪的裡面.
這是很真實的.

$^{241}$我們去求和天父建立這個關係.
我們也求這個罪.
不要使我們將我們跟人的關係.
阻隔,拉得越遠.
禱告也是提醒我們要赦免.
使我們內心沒有東西可以阻隔我們跟父的關係.
禱告也是保守我們的心.
叫我們不要陷入試探.
邪惡能夠進入我們的思想.
影響我們所想的事.
再影響我們所做的事.
所以我們要求神保護我們.
免得這些邪惡的意念.
經常出現,阻隔我們跟父的關係.
這是父願意保護我們.
但我們也要將這個心交給他.
當我們每次都這樣禱告的時候.
就會經驗到禱告的能力.
所以禱告真的需要練習.
我有一個方法幫助大家去記.
我們都有眼睛.
我們用眼睛幫助我們.
當我們望著天的時候.
我們要想起要讚頌我們的天父.
願他的名聲被稱頌.
他的國會下來.
當我們望著天的時候.
我們要讚頌他的名聲.
當我們望著地的時候.
地上仍然有很多不公不義的事情.
需要我們去禱告的紀念.
我們願上帝將這個國賜下來的時候.
將平安也賜給大地.
所以我們望著天和地.
地上有人.
我們人是會生活的.
我們每天都需要食物.
要主持我們一些供應.
屬靈的供應.
生命所需的供應.

$^{281}$所以我們為人所需祈求.
所以我們有天,地,人.
人裡面有罪惡.
我們跟隨經文的結構.
罪惡就是我們祈禱的時候.
就好像抱著天父.
我們會阻隔了天父的關係.
所以我們求天父去饒恕.
我們又去饒恕別人.
其實這是一個很難的功課.
但我們要在關係破裂.
被罪影響的時候.
我們還是交給父.
由父去修補這一切傷痕.
這一切破裂的關係.
最後經文.
我們要保守我們的心.
因為我們不想有邪惡的意念.
影響我們.
因為主知道我們心裡有什麼困難.
有什麼纏繞.
所以我們可以用這個結構.
幫助我們去想想如何祈禱.
是不是很有層次?.
所以大家要記住.
不用用紙或筆寫.
我們走出去看看.
看看天,地,人.
看看世間的罪惡.
看看自己的內心.
其實這就是禱告.
而且我們不是每天睡前才禱告.
我們走出去.
我們隨時隨刻都可以這樣禱告.
禱告的心是隨時隨地.
我們看到有什麼需要.
我們就在禱告裡交代.
第三.
態度要到位.
第九節.

$^{321}$我又告訴你們.
祈求你們就得到.
尋找你們就找著.
敲門就給你們開門.
這是一個禱告的指引.
就像一般的手冊教你如何祈禱.
第十節.
因為凡事祈求的就得到.
尋找的就找著.
敲門的就給他開門.
這就是一個禱告的應許.
就像一本手冊告訴你.
你這樣做.
然後你就會得到一個禱告的應許.
經文就是這麼簡單.
大家可以看到.
它是很相似的.
我高亮了給大家看.
這就是禱告和夢誰聽的結果.
照著做就會得著.
那為什麼不照著做呢?.
所以耶穌就講了一個比喻.
再告訴門徒.
你們的態度應該要怎樣.
第五節.
耶穌對他們說.
你們當中有朋友.
半夜去到他們那裡.
朋友說借我三個餅.
因為我另外有朋友從遠方來.
我沒有什麼招待他.
於是屋內的人就說.
你不要煩我了.
我的門都關了.
小孩都在睡覺了.
我拿不到餅給你.
但耶穌說.
我告訴你們.
即使大人不念朋友之情.
都會起來將餅交給他.

$^{361}$因為不想丟臉.
把他所需要的都拿給他.
這個比喻其實是說什麼呢?.
當時的人家裡未必有房間.
用的都是木門.
關門的時候會有聲音.
整個家裡可能都是睡在同一個房間.
如果他起來開門.
其實就會吵醒大家.
今天我們做麵包.
很可能是焗爐.
自己家裡吃多少就做多少.
但當時的社會.
他們做麵包的時候.
是一大堆婦女一起做的.
一起約一起搓麵包.
其實整個村子.
大家都知道.
大家做了多少個麵包.
意思就是.
有遠方有人來說要借麵包.
有鄰居說來要借麵包的時候.
其實大家就知道.
大家有多少個麵包.
所以知道你有才來借.
即使他不念朋友之情.
也不想丟臉.
你明知道我有的.
我為什麼不借給你呢?.
所以也會不想丟臉.
所以也會借給他.
所以整個村子一起做麵包的時候.
有朋友來.
這是整個村子的事情.
不能說.
我就不給.
你明知道他有的.
為什麼不給呢?.
所以這個其實是說.
其實是在叫門徒.

$^{401}$不怕厚臉皮.
只要拍門求就行.
意思就是你知道父有的.
你知道這間屋有麵包.
就不斷敲門叫他給.
我們同樣知道父有的.
父不會有缺乏.
萬物都是從他而來.
一無所缺.
我們需要的時候就向他求.
我們就要帶著這個態度.
我知道你有.
一直的.
去求,祈求.
你祈求的.
剛才經文怎麼說?.
你祈求的就得著了.
尋找的就找到.
敲門的就會開門.
這是一個態度.
這是我們相信.
我們的父是豐豐足足的.
我記得有一次.
我去探望一個難民的姐妹.
其實難民在香港不讓工作.
還有她的現金其實也很少.
有一次我去探望她.
帶了幾個年輕人去探望.
她準備了很多食物.
而且她買最貴的材料去做.
那時候我們也擔心.
她花了這麼多錢去做.
也送了一些超市禮券給她.
但那位姐妹跟我說.
她說她的父親是王.
她的父親是王.
她的身份你明白嗎?.
是兒女.
也是公主.
王的女兒.

$^{441}$她不擔心她沒有.
她知道爸爸有.
所以她願意去付出.
其實這就是經文引領我們.
去禱告的態度.
就是帶我們去父那裡.
我們知道我們的父.
是充滿人愛,恩典,憐憫.
我們去擁抱這個父的時候.
我們是很堅定的.
他是能夠豐豐足足賜給我們.
天父不單會拿最好的東西給兒女.
也會將我們去求他給聖靈的話.
去引導我們.
去過每一天的話.
他都會給我們.
我們去求.
就會得著.
我們去求聖靈帶領我們去很多工作.
聖靈就會作供.
就像我們今天的福音樂曲.
當我們去求.
心一意去求神去愛.
去樂意.
而甘願讓我們的生命去根深.
我們的意念的行為.
都是去根深的話.
我們就是為著我們的生命.
去彰顯主愛的時候.
主就會幫助我們.
所以我很想和大家一起分享一下.
在禱告裡面的領受.
祈求就會得到.
找著就會尋見.
敲門就會開門.
在2003年年尾.
那時候.
我本來住在將軍澳.
我搬到元朗的位置.
我記得那時候很清楚.

$^{481}$搬到來.
很想家裡能夠成為一個.
祝福其他人生命的地方.
於是我就開始買一些擺設.
耶穌是我家之主.
就放在家裡.
這裡畫一些相片.
放一些經文.
開始請朋友來.
後來因為讀神學.
就少了這些活動.
請朋友來家裡吃飯.
後來2019年.
因為又搬到附近.
開始就和大家分享過.
先生在花園每晚都開放.
有一群朋友來了.
一起分享一下生命.
先生也在當中分享信仰.
每晚都是這樣.
所以到現在.
已經四年了.
很感恩.
在這四年裡.
大概有十多位人士來過.
最多二十位.
他們是不同時期來的.
當中其實也有兩位.
在錦宣家那裡浸泥.
近來有一個.
上星期.
也是在花園那裡聽信仰.
上星期就在錦宣家浸泥.
所以我覺得很奇妙.
禱告是用在家裡.
但神也祝福了外面.
花園一個.
我都沒有擺設的地方.
神就用了那個地方.
畢業之後.

$^{521}$我去了一個堂會服侍難民.
當時我們沒有一個地方.
我們都是用一個聖堂的地方.
但教會有一個位置.
長期都是空的.
我就在想.
這個這麼漂亮的地方.
好像一個家.
有廳有桌子.
不如這裡做基地.
就是那個事工的基地.
當時教會的地方.
不是隨便說想用就用.
都是做夢.
在禱告裡教徒.
但很開心.
2021年我記得.
有這樣的想像之後.
2022年6月.
我們開始和弟兄姊妹.
我們會在這裡查經.
祈禱.
開始聚會.
定期每個星期都在這裡出現.
我們當中也有參與的同工.
也祈禱.
希望將這個地方.
能夠用食物去招待更多的群體.
於是在2023年3月.
即是上個月開始.
有附近的酒店.
開始提供食物.
讓參與這些查經.
弟兄姊妹在活動後.
可以一起進食.
很有生氣.
禱告就是.
我覺得很神奇.
願意為他們求的時候.
神就開路.

$^{561}$食物又供應.
人又供應了.
所以也很奇妙.
好像一個做夢的地方.
做夢告訴上帝.
上帝又像批准了.
讓事情發生.
這也是我們之前的堂會.
2022年2月.
也是疫情比較嚴峻的期間.
其中我們一位同工.
很想發展.
派一些水果蔬菜.
給有困難的難民弟兄姊妹.
也很開心.
有這個雜物房.
雜物房一直都是空置的.
也是透過很多禱告.
才能用到這個地方.
所以往後2022年.
2023年.
也有做食物的派送.
那個就是貨倉.
我們每星期兩次.
把食物運到貨倉.
到香港某幾個點派發.
這位同工也是非洲人.
他在2023年.
得到人道年獎.
透過這個派送食物的事工.
對他的付出有肯定.
也很神奇.
在他的分享中.
他覺得有需要去派.
然後在禱告中交託.
他也有機構贊助他.
有一個更大的地方.
剛才我們看到那個地方很小.
有機構贊助他.
接下來會有大概1000呎的地方.

$^{601}$去容納這些物資.
讓他更加方便去派給有需要的群體.
想著想著.
祈禱.
事情慢慢地出現.
記得我在2022年.
還在讀神學.
因為是另外一個課程.
那時候學到.
其實我們每個人都是.
被上帝用的雅氣.
我們的雅氣就是有裂痕的.
不完美的.
但我們裡面有耶穌基督.
我們能夠在這些裂痕裡.
綻放耶穌基督的光出來.
所以我就被感動.
很想將我服侍的弟兄姊妹.
都能夠發揮他們的才幹.
每個人在上帝裡面.
上帝也有些心意.
有些才幹因此想他們去運用出來.
所以就有這個意念.
所以在上個月.
有一個活動.
就是100個中學生.
跟一些非洲的朋友去互動.
那些非洲朋友就透過.
跟他們煮食物.
打非洲鼓.
去教中學生語言.
其實就是將大家的生命互相交流.
同時這班我服侍的非洲人.
他們也能夠用他們的恩賜.
也能夠用他們的才幹.
去找回他們自己.
因為在香港他們也不能工作.
但是給了這些機會.
有些本地的群體去欣賞.
神所給他們的才幹.

$^{641}$這也是一個很美麗的事工.
我覺得是祈禱.
為他們去求的時候.
神就慢慢開路.
這個教早前也跟大家分享過.
2014年去非洲短孫的時候.
就看到一班婦女.
這班婦女是單親的媽媽.
或者患了愛滋病的婦女.
他們來到這個中心.
去奉獻衣服.
大家互相支持.
到了2022年.
其實很感恩.
那時候就有一些醫車.
2020年11月的時候.
就有一些醫車捐贈給我們.
我們就有一個地方可以送東西.
到了2023年.
又找到一些非洲的姐妹.
真的有心去送東西.
願意每個星期去投放他們的時間.
其實就是想做一個團契.
透過送東西.
這件事很神奇.
首先就是在意念裡面.
有醫車有布工.
有線輪的工.
然後又有一些姐妹一起參加.
我們也開始去做一些衣服.
其實他們很喜歡做衣服.
媽媽能夠為女兒去做一些衣服.
運用一些創意.
去發揮他們的創意和才能.
也是一件很美麗的事.
也很感恩.
大家都是很投入.
除了做衣服.
我們也嘗試做一些圍裙.
其實我們的目的.

$^{681}$不是要訓練他們做很漂亮的東西.
而是成立一個團契.
大家互相去支持.
互相去學習.
可能他的車衣技巧比較好.
去教其他人.
其他姐妹.
大家在當中.
有時醫車也不太合作.
學習忍耐.
其實有很多信仰的質素.
可以在當中發揮出來.
然後我們也開始去嘗試畫作.
大家有沒有看到一些風景畫.
其實也是運用我們的創意.
對於這個群體也有些療癒性.
所以在2023,2024年.
也是按照這個方向去禱告祈求.
其實這個圖畫和剛才有什麼不同呢?.
大家有沒有看到有什麼不同呢?.
剛才的婦女是沒有包頭的.
現在在想像中.
也希望可以招聚一些穆斯林的婦女.
非洲籍的穆斯林婦女.
可以參加逢人的事工.
透過這個事工.
將耶穌基督的愛,信仰.
或者在耶穌基督當中.
究竟夫妻應該怎樣相處.
如何按照聖經的教導去育兒.
慢慢地把這些東西.
給予非洲的穆斯林婦女.
所以也是在想像中.
如果你問我.
其實我也不知道這些婦女在哪裡.
因為我現在認識的只有一些婦女.
但他們也不是很想去奢依.
所以也是在禱告中交託.
這件事對我來說也很大.
我也不知道這些人在香港哪裡.

$^{721}$所以我覺得越是越大.
我們就越用禱告去交託.
越是做不到的時候.
當真的做到的時候.
就是上帝的工作.
所以我們去祈求的時候.
我們是按照父的議程.
即是他的計劃去求的.
我們去見證.
父的榮耀彰顯在這個地上.
一切都是從禱告開始.
我們是否願意一起.
向上帝去發見上帝國的夢呢?.
我們可以學習.
繼續為紅人家去禱告.
紅人家有很多需要.
我們也是在禱告中去學習.
慢慢一步一步去學習.
剛才我說過.
有什麼方法去幫助大家?.
我們用我們的眼睛.
我們看天看地看人.
看地上的罪惡.
然後再紀念達國人的內心.
不要被罪去纏繞.
一起去見證禱告的力量.
就是要由心開始.
提醒我們遠離這些邪惡.
去更新我們和人的關係.
當我們去求天父去供應的時候.
我們不要只求自己.
我們乃是天天去為.
旁邊有需要的人去求.
因為如果我們記得.
其實聖經有說過.
當我們去求祂的國祂的義的時候.
上帝就會將我們所需要的.
都供應給我們.
所以我們去求.
我們去為別人去求.

$^{761}$這個大地這個情況.
有什麼需要改善改變去更新的.
或者你身邊的人.
有誰的生命.
是需要你的禱告去支持的.
就由這些事開始去學習.
如何禱告.
而讓我們一起去禱告.
這樣禱告我們就可以一起去見證.
然後為上帝的名去歡呼.
我們用我們的禱告.
爆發出來的力量.
然後我們就用眼睛一起去見證.
讓我們一起去禱告.
親愛的天父.
在今天的早上.
你教導我們如何去祈禱.
願你的名被尊榮成.
叫天上地下一切的萬物.
連石頭都要為你去歡呼.
願你的國來臨.
願這個地上有不公平不公義的地方.
你都是罕見.
你都是要去更新.
彎彎曲曲的道路你要變直.
大山小丘都要削平.
崎嶇的道路要整平.
願我們都是這樣.
用我們的眼睛去看到.
你的救恩如何降臨在這個地上.
又求你按我們每天所需賜給我們.
叫我們一同去經驗你的豐盛.
亦都是求你赦免我們的罪.
正如我們亦都是赦免虧負我們的人.
願你修補我們人與人之間的裂痕.
那種關係的破裂.
甚至是我們和你的疏離.
都願你親自去修補.
亦都求你叫我們不要陷入試探.
保守我們的心.

$^{801}$我們這樣的祈求.
全部都是為了榮耀你.
直到永遠.
阿們.
\newpage



\section{使徒行傳 5:29-42}
\label{sec:c61oKrE6RXs}
\textbf{2023.04.23 《神聖的耶穌》使徒行傳 5\_29-42 (和修版) 曾裔貴牧師}
\newline
\newline
連結: \href{https://youtube.com/watch?v=c61oKrE6RXs}{\texttt{https://youtube.com/watch?v=c61oKrE6RXs}} ~~~~ 語音日期: 2023-04-28
\newline
\newline
\hyperref[sec:b_zpDLNhNUk]{\small{< < < PREV SERMON < < <}}
~
\hyperref[sec:index_chronic]{\small{[返順時目]}}
~
\hyperref[sec:index_scriptual]{\small{[返順卷目]}}
~
\hyperref[sec:_IeOxEq_8vc]{\small{> > > NEXT SERMON > > >}}
\newline
\newline
使徒行傳 5:29-42
\newline
\begin{longtable}{cl}
\hline
\hline
章節 & 經文 (和合本修訂版)\\
\hline
5:29 & \begin{tabularx}{0.7\textwidth}{X} 彼得和眾使徒回答：「我們必須順從神，勝於順從人。 \end{tabularx} \\ \\ \relax
5:30 & \begin{tabularx}{0.7\textwidth}{X} 你們掛在木頭上殺害的耶穌，我們祖宗的神已經使他復活了。 \end{tabularx} \\ \\ \relax
5:31 & \begin{tabularx}{0.7\textwidth}{X} 神把他高舉在自己的右邊，使他作元帥，作救主，使以色列人得以悔改，並且罪得赦免。 \end{tabularx} \\ \\ \relax
5:32 & \begin{tabularx}{0.7\textwidth}{X} 我們是這些事的見證人；神賜給順從的人的聖靈也為這些事作見證。」 \end{tabularx} \\ \\ \relax
5:33 & \begin{tabularx}{0.7\textwidth}{X} 議會的人聽了極其惱怒，想要殺他們。 \end{tabularx} \\ \\ \relax
5:34 & \begin{tabularx}{0.7\textwidth}{X} 但有一個法利賽人，名叫迦瑪列，是眾百姓所敬重的律法教師，他在議會中站起來，吩咐人把使徒暫且帶到外面去， \end{tabularx} \\ \\ \relax
5:35 & \begin{tabularx}{0.7\textwidth}{X} 然後對眾人說：「以色列人哪，對於這些人，你們應當小心怎樣處理。 \end{tabularx} \\ \\ \relax
5:36 & \begin{tabularx}{0.7\textwidth}{X} 從前杜達出現，自命不凡，附從他的人數約有四百；他被殺後，附從他的人全都散了，歸於無有。 \end{tabularx} \\ \\ \relax
5:37 & \begin{tabularx}{0.7\textwidth}{X} 此後，登記戶籍的時候，又有加利利的猶大出現，引誘百姓跟從他，他也滅亡，附從他的人也都四散了。 \end{tabularx} \\ \\ \relax
5:38 & \begin{tabularx}{0.7\textwidth}{X} 現在，我勸你們不要管這些人，任憑他們吧！他們所謀所為若是出於人，必要敗壞； \end{tabularx} \\ \\ \relax
5:39 & \begin{tabularx}{0.7\textwidth}{X} 若是出於神，你們就不能敗壞他們，恐怕你們倒是攻擊神了。」議會的人被他說服了， \end{tabularx} \\ \\ \relax
5:40 & \begin{tabularx}{0.7\textwidth}{X} 就叫使徒來，把他們打了，又吩咐他們不可奉耶穌的名講道，然後把他們釋放了。 \end{tabularx} \\ \\ \relax
5:41 & \begin{tabularx}{0.7\textwidth}{X} 他們歡歡喜喜地離開議會，因他們算配為這名受辱。 \end{tabularx} \\ \\ \relax
5:42 & \begin{tabularx}{0.7\textwidth}{X} 他們就每日在聖殿裡，在家裡，不住地教導人，傳耶穌是基督的福音。 \end{tabularx} \\ \\
[1ex]
\hline
\hline
\end{longtable}
$^{1}$很感恩再一次能夠在大家當中.
一同來領受聖餐.
以及一同來分享神的話.
我們先做一個簡單的祈禱.
多謝主.
我們要打開聖經.
剛才主席所中毒的.
都讓我們看到.
我們主的聖名.
在使徒的時代.
是多麼有能力.
門徒這樣來跟隨.
宣認的時候.
會帶來這麼大的影響力.
我們很盼望.
我們每一位弟兄姊妹.
都認定主耶穌的聖名.
以及我們都能夠在我們周邊的環境.
使更多人認識我們的主.
我們在這面前等候.
禱告奉耶穌基督的名.
阿們.
今天這篇講道.
也是一個挑戰.
人怎樣看我們.
很重要的是.
我們有沒有在一個方向上.
沒有以自己做中心.
而是我們基督徒.
能夠以神.
以及神的工作.
作為我們的焦點.
以及中心.
這是一個大挑戰.
因為不容易的.
剛才是早了來.
十一點還沒到就到了.
因為教完主一學.
又去大法.
不是宣傳.

$^{41}$接著那位姨姨.
牧師又來了.
又很好招呼.
大家的睦鄰關係很好.
雖然很忙.
但是沒機會聊天.
但是感覺到他們很舒服.
看到一些熟客.
其實不是來很久.
幾個月才來一次.
但是.
今天我們看這段聖經.
這一群使徒.
被人認出他們是跟隨耶穌.
但是.
招致的是.
一個很大的攻擊.
甚至乎.
最後.
有一個法利賽人.
叫加瑪烈.
幫他們開脫.
才不用死.
但是實際上.
最後也打了他們一頓.
知道他們是跟隨耶穌.
打了他們一頓之後.
趕他們走.
接著不准他們傳耶穌的名字.
不過.
第41節這樣說.
他們離開公會之後.
對不起我用的是舊的和合本.
我帶了自己師父的聖經來.
真的盲都認出來了.
拿著這麼大的聖經去吃早餐.
一定是牧師了.
其實不是.
他認得我.
不過這樣.

$^{81}$他說他們離開公會.
心裡歡喜.
因為被算是配為這個名受辱.
配為耶穌這個名字.
來受辱罵.
所以今天這一段經文.
很值得我們反復思想一下.
一個不信神的加瑪烈.
他是法利賽人.
他是一個受百姓敬重的.
一個律法教師.
一個在社會上相當有影響力.
有地位的.
所以他講幾個案例之後.
那些百姓就接受他的勸告.
接受他的勸告就是.
你們不要動這幫人.
這幫人.
如果你搞錯了.
他們如果真的為神做事.
你們就小心一點.
大概是廣東話的翻譯.
就是說.
加瑪烈也有一個小心行事.
就是如果這幫人.
只不過是為自己.
來做一些事.
要達到自己的名聲.
始終有困難.
有阻礙.
他們就會過去的了.
但是你們要留意.
如果這幫人.
真的為神.
和有神差遣他們.
他們所做的事.
就不會敗壞的.
就是不會失敗的.
這裡有一個值得我們真的要去想.
我上個.

$^{121}$我們在武堂.
來到去在復活節期間.
壽苦節.
前後.
我有一個習慣.
就是每一個安息聚會.
我都會回家做一些記錄.
所以你看到我.
如果你有機會看我的電話.
那個行事曆的一些app.
很長.
因為我主禮過的.
那個安息禮拜.
由我自己主禮的.
我記錄了.
很長.
那是什麼意思呢?.
很多人經我們牧者.
幫他辦安息禮拜.
將來就不孤單了.
很多人等著迎接我們.
去到天堂.
很多已經離開世界的聖徒.
他會迎接我們.
我很喜歡在這些大的節期.
來去看一看.
回憶一下這些聖徒.
他們離開世界.
其中再一次.
讓我勾起一個人.
在2015年的時候.
有一位姓王.
我當時不認識他.
姓王.
他的名字很熟悉.
基學模擬的模擬.
王模擬.
不認識他.
突然有一天.
到醫院打電話來找教會.

$^{161}$同工就叫我聽.
我聽到.
什麼事呢?.
他說有一個病人.
他想洗禮.
我當然開心.
趕去醫院.
醫院去到的時候.
我不認識他.
我就表明來意.
給他一個卡片.
那個姑娘.
姑娘就說.
麻煩你現在要全套裝備.
穿保護衣.
2015年是沒有疫情的.
很自由出入.
穿了保護衣之後.
我說我又沒有病.
他說不是他感染你.
是牧師你感染的病人.
我進去之後就明白了.
原來他已經奄奄一息.
他跟姑娘說.
他說他要洗禮.
輾轉就透過老人院.
後來我問老人院.
因為是元朗一間老人院的院友.
後來送去博愛醫院.
這個病人居然要求.
有牧師幫他洗禮.
後來慢慢認識了老人院.
原來是一個已經沒有人照顧的院友.
沒有人探.
根本就自己一個人.
大概五十多歲.
但身體很虛弱.
誰人如果我們沒有保護.
是會傳染給他.
但當我問他心事的時候.

$^{201}$即是查問他的信德.
他很堅定的眼神.
和很肯定的說話.
他願意接受耶穌.
我立即幫他洗禮.
當然在護士的見證下幫他洗禮.
一個素未謀面完全不認識的人.
輾轉就回教會工作.
跟他了解.
因為問老人院拿資料.
因為要幫他做洗禮證.
沒多久之後.
一個多月後.
老人院直接打電話給我們.
告訴我們他已經在醫院離開世界.
接著老人院就請求.
牧師你是他唯一的親人.
你幫他辦安息禮拜.
我一個人加上另一位同工.
就去辦安息禮拜.
沒多久之後.
食環署就打電話給我們.
請你幫我們拿他的骨灰.
幫他去殺灰.
我一個人教會.
一手一腳幫這位黃無義姐妹.
去辦理他人生的後事.
沒親人.
一個親人都沒有.
有老人院的一位都熟的姐姐.
都來到一起.
一個這麼簡單.
簡單到只有幾個人的喪禮.
但是她的眼神.
2015年到現在已經八年了.
她的眼神.
我再去看的時候.
回憶當日.
由醫院.
接著安息禮拜.

$^{241}$接著幫他殺灰.
整整的過程.
好像在腦海裡再一次浮現.
主耶穌的名真的很偉大.
人幾無助.
得到主耶穌.
認定主耶穌是復活.
就是這麼簡單的信念.
其實我不知道是不是她小時候聽過福音.
她要找耶穌.
但是就是這麼簡單.
就感受到這個人成為神的兒女.
現在再沒有痛苦了.
在天堂上.
等待牧師將來去重聚.
去見他.
我是他的親人.
其實如果在聖經來講.
你們大家都是他的親人.
每一個在主耶穌裡面.
都是神的兒女.
不過我們回到這段聖經的時候.
我們看到.
加瑪雅烈所說的.
正正是表白使徒們他們的信心.
使徒們所做的.
就是為主耶穌所做.
你留意他們怎麼說.
第29節.
彼得和眾使徒.
就是11個使徒.
回答當時的官長.
這樣說.
他說信從神.
不信從人是應當的.
是什麼意思呢?.
上文就是說.
你們不可以再傳耶穌了.
如果再傳耶穌.
我們就要來監禁你們了.

$^{281}$這是27節這麼說.
將他們帶到公會的時候.
大祭司就問他們.
我們不是嚴嚴地禁止你們.
不可以再奉耶穌這個名字.
來教訓人.
為什麼你們現在將這個道理.
即他們認為的假道理.
傳遍了整個耶路撒冷呢?.
即現在有一個禁止.
你如果再說耶穌.
令到周圍的人都聽耶穌.
我們就會對你們不客氣了.
接著他們這樣說.
第29節.
就是這段很重要的聖經的其中的鑰節.
信從神.
不信從人.
即是意味著.
他們不是為自己去做事.
他們所做的一切是為神.
神許可他們這樣做.
神吩咐他們這樣做.
他們就繼續.
不管有多少難處.
都繼續將耶穌傳揚開去.
所以當然遭受到再一次的監禁.
再一次的審問.
甚至乎這次可能會是最後一次的.
有殺身之禍.
會被這些工會的人.
嚴刑地傷害.
所以第33節.
真的.
工會的人聽見.
就極其的惱怒.
想要殺他們.
原來我們在人生裡.
難免會有不同的挫折.
不同的遭遇和困難.

$^{321}$秘訣在這裡.
我們到底是為自己去做的事情.
還是我們為神.
以神之名去做的事情.
如果是為神.
神就會開出路.
神就會有保守.
這次神怎樣保守他們呢?.
居然在他們對方的敵對陣營裡.
在工會裡有一個最高權威的長老.
加瑪列.
來勸勉這班工會的人.
你們要小心.
他這番說話.
正是將使徒的內心.
他們的價值觀.
說了出來.
我們再看一次.
加瑪列他們怎樣說.
和這班領袖怎樣說.
和工會的人怎樣說.
留意這裡.
第34節.
但有一個法利塞人.
名叫加瑪列.
是眾百姓所敬重的教法師.
在工會中站起來.
吩咐人.
暫時將使徒帶到外面.
即是開內部會議.
不讓使徒知道.
試想想.
使徒可能會擔心.
他們在裡面談甚麼呢?.
怎樣密謀要殺我們呢?.
好像是玉錘砧板上.
跟著就說.
就和這些工會的人說.
35節.
以色列人.

$^{361}$輪到這些人.
你們應當要小心.
怎樣去辦理.
為甚麼要小心?.
權在自己手.
一切的實權.
都在工會裡面.
為甚麼要小心?.
這個加瑪列.
有一定的智慧.
有一定的謹慎.
原來在他的心底.
還害怕一件事.
就是害怕神.
猶太人有神.
雖然他們未接受耶穌.
不肯相信耶穌是他們的彌賽亞.
但是猶太人對神的觀念.
很強烈.
你看他們的文章.
他們說到那個God字.
GOD.
大家都懂得讀.
大家都知道這個字就是神的意思.
猶太人.
他們在文章裡面.
GOD這個字.
O是不可以寫出來的.
所以是沒有發音的.
發不到音.
但大家知道那個是神字.
因為他們認定了.
我們不可以忘稱神的名字.
猶太人的文章是這樣的.
他們不敢隨便稱呼神.
所以加瑪列.
他有一個謹慎.
他的謹慎就是他害怕神.
如果這班人.
真的有神.

$^{401}$和神真的叫他們這樣做.
那就很大件事了.
雖然加瑪列自己未信神.
未信耶穌.
但是.
他有一個這樣的懼怕.
然後他引了一些案例.
這些案例就是說.
其實.
由得他們吧.
如果他們不是出於神.
自然慢慢就會瓦解了.
一個群體.
如果慢慢有時候為自己個人的利益.
為個人自己的野心.
有時候你面對一些挫折.
一些說話.
你就會很容易灰心.
做得這麼辛苦.
被人這樣說.
被人這樣誤會.
不做了.
通常群體瓦解都是這樣的.
但是如果.
一心為神.
不要管閒言閒語.
也不要管一點點挫折失敗.
一心為福音為神.
慢慢慢慢神就會開出路.
會幫你澄清很多事情.
所以.
他引的案例就是說.
從前有一個叫丟大的人.
他自誇為大.
跟隨他有四百人.
但是當這個人被殺之後.
跟隨的人自然就會散開了.
很自然.
真的.
第37節.

$^{441}$然後有另外一個案例.
他說當時有一個.
人口調查的時候.
有一個加利利人叫做猶大.
他引誘一些百姓又要跟隨他.
意味著.
兩個案例.
都是針對著彼得.
這個要做大.
做大哥的領袖.
就是說現在隨他.
他有挫折.
有失敗.
如果他都這樣離開.
自然那些群眾.
自己就會瓦解了.
不用擔心.
就是說.
給一些難處給彼得.
他自然就會知難而退了.
親愛的弟兄姊妹.
我們回來教會.
跟隨耶穌.
主耶穌為我們受苦.
在十誡上犧牲.
換取的是我們得到神的赦免.
和永恆生命.
不過.
牧師勸勉大家.
我們跟隨主耶穌.
不一定有平安.
不一定有每一個人對你的讚賞.
可能你的傳道.
會帶來別人的眼紅.
老怒.
甚至要殺這些使徒.
因為.
是直接挑戰他們的信仰.
舊有的猶太教.
就是說.

$^{481}$如果你們說的耶穌.
真是彌賽亞.
那就是我們一直等候的.
錯了.
所以.
是直接針對他們整個信仰的核心.
所以這個公會.
必然會攻擊這些使徒.
一來.
大家現在是此消彼長的狀態下.
所以.
公會的人.
覺得.
有兩件事.
必須做.
第一方面就是.
都要很小心.
到底這些使徒有沒有神呢?.
有沒有神准許你這麼做呢?.
如果有的話.
就沒有人可以阻擋了.
弟兄姊妹.
我們內心必須要有一個活潑的信仰.
我們的內心一定要認定.
無論我們的遭遇環境是怎樣.
我們要認定.
主耶穌在你的生命.
是很有幫助和權威的.
因為.
主耶穌才是我們一生要追求.
要榮耀的那一位主.
所以我們跟隨的是這位耶穌.
所以也可以間接.
就是這群公會的人.
是印證彼得和他們的使徒.
他們是一群烏合之眾.
為自己的理想去做.
還是為福音呢?.
為傳揚耶穌是人類救主的福音呢?.
這是一個試金石.

$^{521}$這個環境就是出現了一個試金石.
去試驗使徒他們有沒有這種堅持.
所以公會的人聽完之後.
第四十節.
公會的人聽從了加瑪列的說話.
就叫使徒進來了.
因為他們在外面.
將他們打了一頓之後.
又吩咐他們不可以奉耶穌的名講道.
因為將他們釋放了.
這裡你們看到了.
果然他們仍然保持著自己舊有的系統.
不可以放下.
放下就什麼都沒有了.
所以既然如此.
公會只會將新的使徒逼迫.
打他們一頓之後.
就趕他們走.
叫他們以後不可以再傳.
這就印證著加瑪列所預言的.
如果打完之後.
你就四散了.
原來為耶穌傳福音這麼大代價的.
還是不要了.
頂智滿念是逢第一週下午會出去傳道.
要記住.
可能人家突然打鎚給你吐口水.
你要記住.
曾姑娘你有沒有搞錯?.
為什麼傳福音會被人打?.
不要搞了.
下次不去了.
也可能的.
但是你留意到第41節和42節.
他們離開公會.
心裡歡喜.
因被算是配為這名獸肉.
他們每天在店裡.
在家裡不住的教訓人.
傳耶穌是基督.

$^{561}$就是說已經計算了代價.
而且不會瓦解.
只會每天在不同的地方.
傳揚耶穌是基督.
所以這個試金石.
有時候困難也會使我們更加的練正.
每一個信徒的內心.
更加的去見證主耶穌.
有一位小朋友.
他十四歲.
醫生診斷說他腳有一個癌腫瘤.
需要截肢.
這個小朋友的爸爸是牧師.
他問爸爸.
我應該怎樣做.
可以在我的生活裡.
能夠人見人愛呢?.
爸爸說.
你愛人.
你愛每一個人.
這個小朋友不懂得講道.
但是他經常聽爸爸講道.
也問爸爸.
他學習去愛他身邊的人.
他在學校愛每一個同學.
有些特別性格缺陷的.
他主動去關心.
所以慢慢整間學校的學生對他都很愛.
因為他知道我愛人.
別人就會愛他.
我愛神.
神會愛我.
所以他愛人之先愛自己.
愛自己因為神愛他.
有一天.
剛才我們說了.
他的腳有癌症腫瘤.
截了肢.
有人問一些很尷尬和痛心的問題.
她姓譚的.

$^{601}$譚文君姐妹.
已經回到天家了.
你有沒有問過神.
為什麼你要有這個病.
對一個患癌症的人.
其實是很痛的問題.
這個小姐妹14歲.
她說為什麼我要問這個問題.
天父愛我.
我又愛天父.
我又根據爸爸的教導.
我愛其他人.
為什麼我要問天父.
為什麼要我有病.
她就認定自己.
一定背後有神的計劃.
她不單止不會灰心.
每次進出醫院.
醫院的病床.
成為她同學團契歡笑的地方.
令到醫院成為一個充滿愛的地方.
她的病房.
護士來到.
小聲一點.
不要太大聲.
趕走同學.
因為太吵了.
接著護士下班.
就和她一起在病房裡.
開始聊天.
談論關於神.
為什麼這個小朋友這麼喜樂.
她不是傻的.
她真的痛得很厲害.
後來過了一段時間.
醫生說不行了.
你整隻腳都要鋸.
就鋸了一隻腳.
因為她的腫瘤一直擴散.
她接著認定一件事.

$^{641}$我要在我的病患裡榮耀神.
不久之後.
就是2001年的911事件.
這位姐妹.
當時她已經是高中.
她就去災場.
有什麼可以幫忙呢.
你只是幫倒忙.
你斷了腳.
截了枝.
她說我幫不了前線.
不要緊.
她想到動員所有的同學.
教會的弟兄姊妹.
年輕人.
去紐約的逐間逐間的中餐館.
請求要製造.
免費提供那些熱食的中餐.
來給災場的救援人員.
數千份每日.
動員她.
動員所有的同學.
教會的弟兄姊妹.
這位譚文君姊妹.
她有很大的困難.
完了.
她的父親譚牧師.
是唯一一個中國人.
得到紐約市長班贈的獎狀.
其實就是她的女兒.
在這段時間.
救援的時間的動員.
這個女兒.
她去到22歲.
醫生說.
對不起.
因為你的癌腫瘤已經完全擴散.
她問了醫生一個問題.
我還有多久的壽命.
醫生說大概一個月.

$^{681}$她很積極地去繼續.
她最後的一個月.
她有一個很好的朋友結婚.
她跟爸爸說.
一定不可以告訴對方.
我有這個絕症.
只剩下一個月的壽命.
因為不希望在人家一生最歡樂的時候.
有這種遺憾和難過.
最後譚文君姊妹.
就在病床上.
跟爸爸說最後一句話.
爸爸.
我完成了天父給我的功課.
我現在要回家了.
就是天家.
她不懂得講道.
但她經常聽爸爸講道.
她知道天父愛她.
所以她學習去愛自己.
愛其他的人.
她的安息禮拜.
是數以百計的同學.
和教會的弟兄姊妹.
來出席參加她的安息禮拜.
雖然是二十二歲.
但她的人生.
不是用歲數去衡量.
是她的內心.
如何去看主耶穌這個神聖的名字.
弟兄姊妹.
復活節雖然過去了.
但復活的主耶穌.
在我們大家的心靈裡面.
繼續可以去幫助其他的人.
但願我們每一個弟兄姊妹.
都帶著這個神聖的主耶穌的名字.
在我們的生活.
在我們自己的際遇裡面.
去經歷上主.

$^{721}$我們一起祈禱.
多謝主.
你愛我們.
幫助我們能夠認識你的愛.
在每一個環境去認定.
你的愛從來都沒有改變.
我們仰望禱告.
奉主耶穌基督的名.
阿們.
多謝各位.
\newpage



\section{馬太福音 25:1-13}
\label{sec:_IeOxEq_8vc}
\textbf{2023.04.30 《有危機意識,認清問題的核心,行動起來》馬太福音 25\_1-13 (和修版) 郭鴻標牧師}
\newline
\newline
連結: \href{https://youtube.com/watch?v=-IeOxEq_8vc}{\texttt{https://youtube.com/watch?v=-IeOxEq\_8vc}} ~~~~ 語音日期: 2023-05-01
\newline
\newline
\hyperref[sec:c61oKrE6RXs]{\small{< < < PREV SERMON < < <}}
~
\hyperref[sec:index_chronic]{\small{[返順時目]}}
~
\hyperref[sec:index_scriptual]{\small{[返順卷目]}}
~
\hyperref[sec:pT_5txip_oE]{\small{> > > NEXT SERMON > > >}}
\newline
\newline
馬太福音 25:1-13
\newline
\begin{longtable}{cl}
\hline
\hline
章節 & 經文 (和合本修訂版)\\
\hline
25:1 & \begin{tabularx}{0.7\textwidth}{X} 「那時，天國好比十個童女拿著燈出去迎接新郎。 \end{tabularx} \\ \\ \relax
25:2 & \begin{tabularx}{0.7\textwidth}{X} 其中有五個是愚拙的，五個是聰明的。 \end{tabularx} \\ \\ \relax
25:3 & \begin{tabularx}{0.7\textwidth}{X} 愚拙的拿著燈，卻沒有帶油； \end{tabularx} \\ \\ \relax
25:4 & \begin{tabularx}{0.7\textwidth}{X} 聰明的拿著燈，又盛了油在器皿裡。 \end{tabularx} \\ \\ \relax
25:5 & \begin{tabularx}{0.7\textwidth}{X} 新郎遲延的時候，她們都打盹，睡著了。 \end{tabularx} \\ \\ \relax
25:6 & \begin{tabularx}{0.7\textwidth}{X} 半夜有人喊：『看，新郎來了，你們出來迎接他。』 \end{tabularx} \\ \\ \relax
25:7 & \begin{tabularx}{0.7\textwidth}{X} 那些童女就都起來挑亮她們的燈。 \end{tabularx} \\ \\ \relax
25:8 & \begin{tabularx}{0.7\textwidth}{X} 愚拙的對聰明的說：『請分點油給我們，因為我們的燈要滅了。』 \end{tabularx} \\ \\ \relax
25:9 & \begin{tabularx}{0.7\textwidth}{X} 聰明的回答：『恐怕不夠你我用的；你們還是自己到賣油的那裡去買吧。』 \end{tabularx} \\ \\ \relax
25:10 & \begin{tabularx}{0.7\textwidth}{X} 她們去買的時候，新郎到了。那預備好了的，與他進去共赴婚宴，門就關了。 \end{tabularx} \\ \\ \relax
25:11 & \begin{tabularx}{0.7\textwidth}{X} 其餘的童女隨後也來了，說：『主啊，主啊，給我們開門！』 \end{tabularx} \\ \\ \relax
25:12 & \begin{tabularx}{0.7\textwidth}{X} 他卻回答：『我實在告訴你們，我不認識你們。』 \end{tabularx} \\ \\ \relax
25:13 & \begin{tabularx}{0.7\textwidth}{X} 所以，你們要警醒，因為那日子，那時辰，你們不知道。」 \end{tabularx} \\ \\
[1ex]
\hline
\hline
\end{longtable}
$^{1}$我們一起同心祈禱.
我們親愛的天父上帝.
我們恭敬的來到您的面前.
願您藉著聖靈來讓我們明白聖經的教訓.
幫助我們在任何一個環境裡.
我們都信得過我們的上帝.
以及我們真的願意盡我們的力量來回應我們的上帝.
忠心的來做主的門徒.
我們的祈禱是奉耶穌基督的名.
阿們.
有一次我聽道的時候.
那個講員就提及.
在2015年的時候有一套美國的電影.
就不是現在是2015年.
那套電影的名字叫The Big Shot.
在台灣就被翻譯成為大賣兇.
在香港就被翻譯成為孤主一擇.
這套電影就拿到第88屆奧斯卡金像獎的最佳影片.
最佳導演等等五個獎項的提名.
就贏得最佳改編劇本獎.
這個故事其實是真人真事.
在2007-2008年環球金融危機的時候.
在美國有四位專業人士就用一些專業的手法.
來操控金融投資進行孤主一擇.
從這件事裡面他們獲得了利益.
在美國還沒有暴露次案危機的時候.
有一個基金經理叫Michael.
就通過大量的數據去分析.
就發覺地產市場由於過度的借貸.
隨時會出現爆破的危機.
由於他覺得這件事是對的.
他就不顧金主的反對.
他就孤主一擇.
就向多間的投資銀行作出了信貸違約的對賭.
結果銀行家都說你真是傻子.
有一個銀行家叫Jarrett.
他知道這個消息之後.
有時候別人都說不要執輸.
我想分一杯羹.
但是他又很想減低風險.

$^{41}$就在機緣巧合之下.
他就找到一個避險基金的經理叫Mark.
來一起合作.
另外還有兩個用車房來做投資的青年人.
叫Charlie和Jamie.
原本是想去華爾街尋求更大的投資機會.
但是在處處碰壁之下.
他們意外的發現了一份文件.
內容就是在說地產泡沫的事情.
他們兩個就覺得大有可為.
就找了一個舊的鄰居.
退休的銀行家叫Ben.
來試一下以小博大.
這套電影就是描述美國普羅大眾的投資心態.
很具體的用比喻來顯示.
不同的人都是很貪婪.
還有群眾是很愚昧的.
上至銀行家下到證券評級分析員.
和審批貸款的人.
都是向錢去看的.
大家是合力的炮製了那一次的金融危機.
影片反映當年借貸氾濫的情況.
和政界商界的人.
是合作打造一種賺緊沒虧的那種假象.
讓人民漸漸失去了危機感.
我們香港人都會說.
哪裡有這麼大的蟑螂隨街跳.
但是當人失去了危機意識的時候.
就不會這樣想的.
銀行家和廣大的投資者.
他們想了一件事.
樓市不會倒的.
實際上他們在想什麼.
誰不會去供樓.
這個是假設來的.
在香港就是在美國不一定.
其實他沒有實際的客觀分析.
只是他想是這樣的.
Mark和他的團隊聽到樓市快要倒.
這個預言他都半信半疑的.

$^{81}$他只是有這樣的想法.
為什麼沒有人發現這件事.
這麼重要的事情沒有人發現.
他就是在這裡做懷疑.
不過他懷疑歸懷疑.
他和他的團隊就去走訪各地的住客查詢.
問問借貸人或者那些證券評級員等等.
或者投資經紀等等.
他不斷去問.
他經過小心的求證.
就相信樓市會有爆破.
他就決定了買入了一些信貸違約的調期.
從中來獲得利益.
我不是跟你說經濟的東西.
而是想透過這套電影說出一件事.
當眾人以為一切如常的時候.
只有少數人才發現有危機存在.
這是歷史.
你去到哪裡都會看到.
那些能夠早些發現有危機的人.
就一定是早著先機.
現在我說這些都是馬後炮.
但是有電影出來.
真人真事的事情.
是可以令我們想很多東西的.
我就用這個例子幫我們進入馬太福音.
第25章第一到第十三節.
我們有些經文就行了.
這個故事裡面有十個同類.
我將題目就沒有那麼一般.
好像屬靈化那樣.
變成有危機意識.
有要認清問題的核心.
第三要行動起來.
我講到都是分這三個部分.
希望大家可以活化這段經文.
有危機意識.
我們看看馬太福音第25章的第三節.
第三節講到那些有五個同類.
就被經文說他愚絕.

$^{121}$他怎樣愚絕呢.
愚絕的拿著燈卻沒有帶油.
你想一下這句其實是在說什麼.
理論上他們拿著的燈是有油的.
沒有人的燈沒有油的.
不過他沒有帶油是什麼含義呢.
是他沒有另外找個容器再燃油.
為什麼說他愚絕呢.
因為他現在去迎接一個新郎.
就是從另一個地方來的.
去接人.
接人你可能就想一下我們香港的處境.
接人是很簡單的.
但是原來在中東的時候.
耶穌的時代.
去接人和我們今天接人是兩碼事.
為什麼呢.
就算你知道對方什麼時候出發.
你怎麼知道他什麼時候到.
哪天到.
你不像我們現在.
我在中環搭東涌線.
在南昌轉西鐵就來.
你大概知道幾點鐘有班車.
幾點鐘到.
但是那時候其實沒得計算.
沒得計算的時候.
所以那些童女去迎接新郎.
其實他不是拿著自己的燈就一定等到他.
道理就是這麼簡單.
只不過那五個童女覺得.
我拿著我手上的燈就行了.
另外那五個童女就不是這樣.
好了.
正經就沒有說到.
為什麼他們愚絕.
就只是用愚絕這個字.
那五個童女愚絕.
就沒有說為什麼他們愚絕.
我可以試一下將背景告訴大家.

$^{161}$其實是這樣.
好了,到第四節.
聰明的是怎樣.
聰明的就拿著燈.
又預備油再戲名裡.
就是說他和那些愚絕的一樣.
拿著盞燈,燈裡面一定有油的.
但是他們就另外找個容器.
塞著油.
就聰明一點了.
愚絕就沒有這樣做.
為什麼他們聰明呢.
正經又沒有說為什麼聰明.
但是原因,我這樣講給你聽.
他去接人的時候.
新郎其實.
大家猜他今天來.
但是不知道他幾點到.
所以聰明的就拿著那盞燈.
還有帶備額外的油在手.
就是這樣.
好了,聰明的.
就帶著另外的油.
目的是什麼.
以防對方遲到.
所謂遲不遲到都很相對的.
那個時候.
我大概,那時候就真的可以.
你看到我就看見,看不到就看不到.
這句話在耶穌時代是講得通的.
在今天我們這樣講,你明白嗎.
你的朋友就說你來不來的.
我們今天是不會的.
我們今天有手機.
你什麼時候到,你在什麼位置.
在地鐵上經常都有.
你在哪個位置.
我們不會想到問題在哪裡.
但是當我們回到耶穌的時代.
你發覺原來這個故事.

$^{201}$真的有背景的.
原來這群所謂愚絕的童女.
她是沒有一個意識.
對方可能會遲來的.
她甚至沒有這個想法.
沒有這個概念.
所以她沒有心理準備.
沒有想過要帶多些油出去.
你想一下.
所以我就說.
看完這段經文再想想這五個.
聖經說她愚絕的童女.
很明顯是沒有危機意識.
或者你再講得進一點.
根本你沒有你的常理.
那時候人人都知道.
無論是盧溪也好,駱駝也好.
你都不知道她什麼時候到.
晚上到也有.
半夜到也有可能.
你去接人.
你不是帶備好油.
其實你想都想不到.
沒有人會這麼傻的.
所以如果你這樣看的時候.
這五個童女基本上是沒有危機意識.
我們中文就說她沒有腦子.
英文就說她沒有智慧.
根本你不是做事的.
好了,為什麼我們說.
聖經說這五個童女有五個聰明呢.
就因為她有心理準備.
如果事情有變的時候.
她都可以走賺.
有個應變計劃.
在一個如果一切如常的時代.
很多人都不會覺得有問題.
那你們記不記得在舊約聖經.
就說過一個時代.
人們都是差不多的.

$^{241}$在香港有個羅亞方舟.
你們知道嗎.
羅亞的時代也是這樣的.
那些人怎麼會覺得有事呢.
唯獨羅亞與眾不同.
他走去蓋方舟.
還叫家人去蓋.
你想想周圍的人會跟他說什麼呢.
你都傻的.
這麼空閒不做事.
蓋一間方舟幹什麼呢.
你蓋一間房子就說而已.
所以很正常的.
很多時候很多人.
他是沒有了危機意識.
只是非常少數的人.
他覺得不是這樣的.
所以如果你是有這種思維.
你可能都會懷疑自己.
是不是想多了.
這幾年的變化.
就令到我們大家.
本來都不是很有危機意識.
現在都增加了危機意識.
你剛才都看到我.
我戴著口罩進來.
為什麼呢.
就因為我坐船.
現在我兩元雞.
我當然坐快船.
快船八個字.
快船是密封的.
不透風的.
旁邊那個人都沒有他扶著.
所以一下船.
我就戴著口罩.
還有坐地鐵.
你們今天有沒有坐地鐵來.
你沒有吧.
今天特別多人.

$^{281}$真的多人到不得了.
我在中環上東涌線去南昌.
那麼多.
上映多人.
在南昌又那麼多人.
是講普通話的.
今天可能有六十萬人.
來我們香港消費的.
就是說.
你戴著口罩其實沒有壞的.
我就是這樣想.
這是危機意識.
好了.
我們經過這幾年.
我們經常都有一個心理準備.
可能事情會變差的.
我們都會這樣想的.
英文又叫 Worst Scenario.
就是一個更加壞的情況出現的時候.
我們就會自然問.
有沒有 Plan B.
有沒有兩手準備.
經常這樣說的.
搞聚會也是.
如果疫情封了的時候.
還有下一個你做不做.
譬如學校的畢業禮又怎樣.
你會發覺原來我們已經是.
被迫要這樣想.
在外在環境沒有辦法控制的時候.
我們會怎樣做.
我們會擺一個防守陣容.
你明不明白.
總之就是保本止蝕.
這是首要任務.
總之就是做到最穩定的事情.
其實香港的情況不算好.
不過很多地方的情況是更加差.
這是事實.
值得我們去想的問題就是說.

$^{321}$究竟我們有沒有心理準備.
我們如果想的事情不能夠實現的時候.
或者我們根本想的事情是建立在一個錯誤的基礎上的時候.
我們應該怎樣呢.
我想我們中國人有四個字很有意思.
居安思危.
我們要想這些.
同樣的在黑暗裡面去預備黎明的出現.
經過這幾年之後你會發覺有一個道理.
就是說停擺這個名字是很普遍的.
但是如果在疫情的時候我們真的停擺.
結果你是一定會起動一萬個人的.
如果你在疫情裡面你盡量維持前進.
你還是可以走得動的.
好了 到了疫情尾的時候.
其實大家應該是想部署了.
有變化的時候你會怎樣的.
是不是這樣 一定是這樣的.
如果他在疫情尾都沒有想頭的話.
他們一定是走得比人慢的.
到現在已經不是說疫情了.
除非有大變化.
其實不是說起動.
是反彈 你有沒有聽過.
就是做生意的人.
我們怎樣彈起來.
誰彈得起來最快那個就是會贏的.
是不是這樣說.
所以在這些事情裡面.
我們可以看見的就是說.
我們要學習的就是思想.
在任何的環境下.
任何時間下.
我們都要想我們怎樣可以向前.
是不是.
可以向前的時候就向前.
不可以向前的時候.
我想一下什麼時候可以向前.
這是必然的道理.
這幾年都是這樣.

$^{361}$所以我想我們祈禱的就是.
求神幫助我們保持一種前進的心態.
我自己是很喜歡想事情的.
我是比較喜歡想方向的那些人.
有些人是很實用的.
所以我會覺得方向和方法都兩樣重要.
就是你一配合起來就行了.
我們需要去問我們往哪個方向去走.
同時要知道用什麼方法你達到那個目標.
在這裡我也很想和大家分享一下.
在這一兩個月.
建立神學院也辦了一些活動.
我覺得是神很保守的.
其中一個就是香港教會前景會.
那個反應是相當不錯的.
因為是疫情後.
似乎應該是香港教會界的一個比較大的活動.
原來很多牧者和弟兄姊妹是有一個訴求.
很想第一方面去討論那些問題.
第二方面是很想大家聚一聚.
結果我們每件事都幾百人參加.
就是一千人.
另外就是剛剛我們四月二十七日晚.
有一個聖岳培靈獻身會.
也是八九百人參加.
也有九成滿.
我們都覺得很奇妙.
大家反應是很強.
但是我回帶告訴大家.
我們計劃這些的時候.
是一年前已經有這個想法.
一年前是什麼狀態呢.
你想一想一年前都不知道什麼時候疫情停的.
所以要做任何的決定都很困難的.
問教堂借地方.
對方都不知道怎麼回答你的.
當然是說看看怎麼樣.
疫情好一點.
你明不明白關鍵就是說.
如果我們不是一年前已經決定做這件事.

$^{401}$今年的四月.
三月四月.
即是四月兩次的大聚會.
能不能搞得來呢.
你不可以搞得來.
你不可能兩個月搞得來的.
如果我們開始口罩令停了之後.
我們才說去計劃呢.
那行啦 你明年搞得來.
你明白我的意思嗎.
所以整件事就令我學習到一件事.
有些事不是我們等到環境好的時候才去想的.
而是還沒有好的時候我們已經想.
好的時候我們有什麼行動.
這樣你才有得做.
不然就一定沒得做.
大家同不同意嗎.
有人認為疫情之後教會是出現大洗牌.
這個現在是事實了.
因為有些數據出來了.
除了數據那些我不說那麼多.
我有時在不同牧者的口中知道.
有些我們叫超大型的教會.
即是叫Mega Church.
人數下跌的情況是非常非常嚴重.
非正式統計.
其中有一間是由五千人下跌到兩千人.
不要說名字.
有一間大概由一萬人下跌到四千人.
不要說名字.
去年我去過其中一間教會的成人團契的導師訓練班分享.
所以我是有真憑實據的.
我發覺除了他們的牧者之外.
有很多弟兄姊妹是做他們團契的導師.
其中一位姊妹她就厲害.
她不止做一個團契的導師.
她其實是做了三個團契的導師.
她分享當時弟兄姊妹有很多張力.
對於信仰有很多疑惑.
有些人對教會很失望.

$^{441}$有很多屬靈的衝擊.
她盡力去鼓勵那些弟兄姊妹.
我心裡想.
如果不是有這些團契裡面的導師.
其實那些只是信得久一點.
她衷心的去守護弟兄姊妹的時候.
可能更多的人冷淡.
五千人變成二千.
分分鐘沒有了這班人可能還跌得厲害.
基於危機意識.
我自己有什麼態度呢.
我就向上帝禱告.
善用時間投入事奉.
這是我自己的禱告.
上帝給我的感動就是.
讓上帝的道去牧養弟兄姊妹.
我什麼都不懂.
我只懂得講道.
寫一些靈修文章.
用神的話語去牧養弟兄姊妹.
我能夠貢獻的不是很多.
不過我只可以把握時間去寫作.
最近有一本書.
基督信仰連結整合.
這也是很感恩的.
是一本講基督信仰的書.
也很感恩.
過去我又去整理一些文稿.
今年五月初會有兩本靈修書出版.
有一本叫十二小先知書品讀.
與神同行的人生.
其實排著也還有.
為什麼我要這樣做呢.
其實我是沒有什麼特別的目標.
我是想說不要停止.
我不可以停止.
為什麼呢.
不是我亂砍不是這件事.
而是我覺得現在大家都好像有一種.
不知道前面是怎樣的感覺.

$^{481}$我自己也不太知道.
你問我我沒有答案.
不過有一件事就是說.
我可以做什麼就做什麼.
你明白嗎.
我們不要浪費時間.
我們就做可以做的事.
我問自己你做了這些事.
寫完這些書.
你可以改變現在香港教會的困局嗎.
是不是.
我問我自己.
我的答案是不一定可以.
我又問自己.
既然我不一定可以改變.
你又做來做什麼.
你這麼投入來做什麼呢.
我自己就會說.
這個是我們對上帝的呼召的一個理解.
上帝要我們去做就做了.
認認真真地去投入.
不要計算有什麼結果.
所以就是這樣.
而每個基督徒都會有他的照明.
我們在這個時候.
我們會見到很多的挑戰.
我們解決不了的.
我們每個基督徒都應該有一種態度.
就是忠於我的使命.
我去履行上帝要我做的事.
做不做得好.
我自己都不敢說自己做得好不好.
留待將來.
上帝來去評核.
現在可以做的就照做.
好過你沒得做就不用做.
信徒皆為祭司.
就是每個基督徒.
無論你是不是牧師.
是不是傳道人.

$^{521}$你都應該事奉上帝.
尤其是在現在這個環境.
更加要把握時候去事奉上帝.
現在香港教會只是轉型的時候.
教會很需要你.
你會是那五個愚拙的同類.
還是那五個聰明的同類呢.
你會不會有危機意識呢.
第二我想說的是.
認清問題的核心.
在《馬太福音》第25章第6到第9節.
就說到新郎遲遲都沒來.
那些同類就怎樣.
打頓睡覺了.
半夜有人說新郎來了.
你們出來去迎接他.
然後那些同類就起來收拾燈.
然後愚拙的就對聰明的說.
喂可不可以分些油給我們.
我們的燈就快沒油了.
聰明的就說.
那些都不夠我們用.
不如你去賣油的地方去買吧.
你可以這樣想的.
你怎麼這麼自私的.
你有就給我吧.
對不對.
可以這樣想的.
但說話又不可以這樣說.
為什麼這樣說.
你明知你不夠油.
你為什麼不早點去買油呢.
你就在那裡打頓睡覺.
你明不明.
我看這段經文都想了很久.
是不是這幾個聰明的同類.
聰明的就聰明.
但沒什麼同情心.
我會問是不是你這麼壞.
你老有機.

$^{561}$你的同事有事你不幫.
其實我想想是沒理由的.
你自己不去把機會去解決問題.
你把最後的球彈給人家.
是不行的.
如果新郎遲到.
你帶了額外油的同類.
老實說你能夠有.
就分享給人家.
不過他已經計算過.
我自己都不夠.
我怎麼分給他.
其實有個容器.
可能中間他添油.
他已經添了一點.
你還要他拿剩餘的.
他可能什麼都沒有.
如果你說那些愚蠢的同類.
他等新郎來.
自己又不夠油.
但他去了睡覺.
又不想想怎麼去解決他不夠油的時候.
不知道要等多久.
這件事你覺得他有沒有盡責任呢.
你明白嗎.
道理是這樣說的.
馬太福音第25章1-13節這段經文.
提醒我們要預備神的來臨.
同樣可以引伸到每一個事奉神的人.
我們有沒有好好預備神的來臨呢.
這個道理在這裡.
在過去的事奉經驗之中.
我發覺原來要掌握宏觀的資訊是很重要的.
包括什麼呢.
對普世教會的認識.
對有代表性的教會的認識.
以及確立我們自己對教會的理念.
什麼叫普世教會的認識呢.
其實香港都有很多宗派.
可能在在中的人.

$^{601}$都曾經來自不同的宗派.
每個宗派有他的歷史.
有他的理念.
有他的神學的傳統等等.
還有除了香港的教會.
還有海外的教會.
這麼多教會.
你要明白別人在做什麼.
對有代表性的教會的認識.
其實我們要留心.
他在不同時期.
他發展的特色.
以及他的做法的好處與壞處.
以及今天是怎樣.
所以這件事不是每個人都有這麼多資料.
但是其實找這些資料也不難的.
如果我們是牧者.
我們是需要問.
我究竟我自己對教會的理念是什麼.
做信徒領袖也是一樣要有這個理念.
我們要確定我們自己.
我所認同的教會理念是什麼.
或者我講一點點.
我剛才很感恩.
我有一個同事.
他住在錦繡.
他很早就知道我來紅印堂講道.
他一定要來我家車我.
他還豪氣到去中環碼頭車我.
我說不要客氣了.
我不是那種人.
我在錦上路.
算了錦上路吧.
他是做工程的.
在車上他告訴我.
最近有一個工程.
我們神學院的籃球場修整得很漂亮.
但是有一樣東西我不知道.
他告訴我.
他說這個籃球場四十年前.

$^{641}$應該是北宣的胡音牧師.
他是做工程的.
他帶領一些弟兄去那裡平整了籃球場.
而平整了籃球場.
有什麼東西.
你們都是第一次聽.
通常人家建一個平台.
當然會斜一點點水去別人那裡或者去外面.
道理是這樣的.
弟兄說如果我建東西我也是這樣建的.
但他說不是.
我們的籃球場不是向外斜的.
是凹下來的.
你說凹下來下雨就積水.
原來凹下來又很特別.
他又不會積很多水.
太陽出來又會很快蒸發水.
有什麼特別呢.
他說如果那時的斜水去到旁邊.
這個籃球場的地基其實不會lum 的.
水下去泥那裡是會有後果的.
四十多年來我們的籃球場的地基是沒有lum 的.
沒有什麼大問題.
其實是因為當時那個建工程的人.
他是不想影響他人.
最多雨水是流到自己那裡.
但他又不會淹到大桶水那裡.
這麼高手.
其實不只是做工程的高手.
而是心態上.
他是以不傷害人為主.
他寧願自己的地可以濕一點.
但濕得來不會有很大問題.
以致我們這次去平整地基.
就不用整個籃球場去做.
他不說我怎麼認識這些東西.
他做工程的.
所以原來道理這麼深.
當時的人如果貪方便.
地不會濕.

$^{681}$我會撤出去給人.
其實那個地基可能這幾十年.
我們要維修的時候.
是用更多的錢去做.
但重要的地方就是.
我學到一個道理.
今早坐車學的.
我們做事不是貪一時的方便.
即是說我方便不用這麼多水漬在那裡.
我做的事是要令到根基不要受損.
四十年之後我們去整.
只是要局部地整紅毛泥下去.
如果不是還要花很多錢.
任何一個想法.
他對後面的人.
是有一個長遠的考慮.
那些是上街上街的決定.
還有他的出發點.
他不是去找人笨.
他是寧願自己.
那個煩惱是給自己.
但他又會想辦法去令到煩惱.
不要搞到自己很煩惱.
那這個人是不是很厲害.
很厲害.
他不止技術厲害.
而是心態高.
所以我覺得.
我們今日有沒有辦法好像學他們一樣.
如果你們有機會去我們建築神話院.
看那個籃球場.
好像政府的運動場一樣.
我們看到他光明很漂亮那一面.
但可能沒有人記得四十年前.
去將爛地整到.
鋪平他那班人那種努力.
沒有那班人的努力就沒有今天.
就是這樣說.
所以我們要學習的就是說.
在現在這個時候.

$^{721}$我們要求變.
其實沒有一套方法永遠有效.
現在每個牧者都嘗試新的東西.
有些行有些不行.
究竟有什麼東西可以令到我們有所學習呢.
我們要認清問題的中心點在哪裡.
很重要的.
很多時候我們有很多訴求.
我們都會喜歡自己想喜歡的東西.
同時你可以想教會發展到某個地步.
就有他那種文化.
有他的組織又有他的形式.
每樣東西都定格了.
這是普遍的.
求變是很自然的.
不過應該怎樣變呢.
變的最終目標是什麼呢.
其實沒有什麼明確答案.
其實大家都不知道.
但是這段經文提醒我們一件事.
恐怕是有一些基督徒仍然是沉睡之中.
這個才是致命.
恐怕有些基督徒.
是好像那段經文所說的愚蠢的同類.
完全沒有概念沒有危機感沒有意識.
一直在他的安書區裡面.
從來都不關注教會的事或者其他教會的情況.
或者沒有其他的對不同教會的了解的經驗.
相反有些人又不能夠再等了.
就嘗試新的模式.
當然了他不斷修正一定有很開花結果的.
不過有些時候也有可能是弱石亂頭都不頂的.
在這幾年的疫情之後.
其實通很多教會的人數是下跌的.
這個是真的.
但是有個有趣的現象.
是反而200人以下的教會比100人教會那些.
那個人數下跌不多的.
很有趣的我去很多教會都是這樣的.
越是不大型的教會大家的關係仍然在那些.

$^{761}$他不是跌得很傷的.
一般來說是10\%到20\%的人數下跌.
但是有些大的教會可能跌40\%.
很恐怖的.
恐怖到不得了.
所以在這情況下我們面對的就是人數的變化.
奉獻數字的下跌.
牧者和信徒的移民.
剛才也有談到這些問題.
我們也收到消息.
譬如說英國,美國,加拿大,澳洲等等華人教會.
那些廣東話的崇拜突然間多了很多人.
其實說得不好聽.
加拿大那邊的廣東話崇拜都快要接了.
突然間不用接了.
多開心啊.
當然這個複雜了.
究竟我應該怎樣看現在的情況呢.
我就會有一個這樣的想法.
其實是本土化和國際化兩樣東西.
本土化就是說我們在香港這個地方.
我們面對著很多牧者的退休潮.
或者很多牧者的移民.
信徒領袖的流散等等這些現象.
是要重整教會.
固本培元.
這個肯定的.
一定要這樣做.
第二就是國際化.
現在香港人到處都有.
最好是在這個時候做一些傳福音建立教會的工作.
以前的情況就是說.
那些人飛到美國,英國就由他在落地生根.
但現在的情況很含糊.
有些人他在網上還在參加本地教會的團契.
或者崇拜.
那種情況是怎樣呢.
我們叫藕斷絲連.
究竟你在香港的牧者和弟兄姊妹.
關不關顧他們好呢.

$^{801}$你明不明白我說什麼.
你一是說他去了那裡不用我管.
但現實上又不是.
他經常在鏡頭那裡.
一會兒跟你WhatsApp.
好像沒有離開你那樣.
出現的問題就是說.
那就是說我要繼續關注他們.
關顧他們.
你願意不願意.
在我們的牧養和宣教上面.
其實我們是做多了很多東西.
同意嗎.
是做多了很多東西.
本地化,國際化.
都成為我們的挑戰.
似乎上帝給我們現在很重的責任.
本地都做不成還要做他們.
是不是很吃力呢.
很吃力.
但我想我們要學習的就是說.
既然我們現在面對著一個這樣的環境.
你是不能避的.
你不可以縮水地做的.
唯有我們希望就是說.
我們做得多少得多少.
大家同意嗎.
現在你看到的情況是拉扯得很厲害的.
所以在急變的環境裡面.
我們真的要看清楚.
世界的情況.
面貌是怎麼樣.
以及我們現在看清楚問題在哪裡.
去明白我們上帝的心意.
說到這裡的時候.
好像都很大挑戰.
其實我說這番話在每一間教會都適用.
隨著時代的變化.
我們自己都面對著很多生活和信仰的張力.
怎樣落實我的信仰.

$^{841}$怎樣活出信仰.
都是一個很重要的問題.
我們想想我們自己.
似乎我們自己都很多事情還沒處理.
還沒解決.
再加上聽到剛才那些東西.
其實是很沉重的.
是一種沉重的感覺.
我試過真的.
好幾年前.
我講到.
你都知道我不是很囂張那些.
我只是講一下門徒.
然後已經有一個弟兄跟我說.
我聽你說到有點壓力.
他覺得.
我給了他一點壓力.
其實我沒有.
我只是講聖經.
其實他想怎樣呢.
他想回來聽到.
開開心心聽到.
上帝愛我.
上帝祝福我.
上帝保守我.
我不是說這些是錯的.
這些是很對的.
但是.
耶穌的確呼召人去做他的門徒.
我又不可以不說的.
你明白我的情況.
所以我很明白.
很多弟兄姊妹說.
你知不知道我做事的那種吃力的情況.
我的壓力.
你明不明白.
處境就是.
我們有很多擔子.
但是外面的環境那麼多挑戰.
教會好像.

$^{881}$你不可以不接這個球.
這樣的時候怎麼辦呢.
所以我自己可以說.
我什麼都做不到.
我只是在神學院.
可以訓練神學生.
去做這些事.
但是其實我現在都開始有一個感覺.
就是說.
遠水都救不了近火.
你明不明白.
功不應求.
很多的事情是我們應付不來.
所以我們唯有祈禱.
希望我們.
看到問題的時候.
就要行動起來.
在《馬太福音》第25章第8到第13節.
如著聰明的話.
你分一點油給我們吧.
因為我們的燈就熄了.
聰明的就說.
恐怕不夠用.
不如你去賣油的地方找吧.
當他們去買的時候.
新郎就來了.
好像預場片一樣.
這麼巧合.
那些預備好的人.
就跟他一起去赴宴席.
最慘的是門就關了.
那裡才慘.
其餘的同類就來了.
主啊主啊.
你開門給我吧.
新郎就說.
我實實在在的告訴你們.
我不認識你們.
所以經文就說.
你們要警醒.

$^{921}$因為那日子那時辰.
你們不知道.
這段經文.
有時候用人的心態去讀.
用不著這麼絕啊.
新郎.
我們會有這樣的心態.
但是在信仰的層面來講.
神其實是忍耐我們.
等候我們很久.
但是我們是不是一直都沒有什麼行動呢.
總有一天.
耶穌再來的時候.
你會發現.
那時候我們已經是太慢了.
太遲了.
這段經文告訴我們.
沒有預備額外油的同類.
買油回來了.
新郎都來了.
門都關了.
遲了.
遲了大到了.
我讀經的時候.
我都會回來回去.
新郎你遲到你都有責任的.
你為什麼要這樣罰人.
你明白嗎.
我都會這樣問這個問題.
當然這是人的思維.
新郎遲到是他有責任.
但是沒有去預備額外油.
又是那五個所謂愚蠢同類的責任.
是他們有他們的問題.
如果新郎.
他會遲到.
那你自己是不是應該要想.
如果我不夠油的時候.
我是不是要多帶點油.
這是很正常的.

$^{961}$你不可以反過來.
車前馬後又說新郎你自己不對.
你明白嗎.
那個道理不可以這樣去說的.
是不是.
新郎就說.
我實實在在告訴你們.
我不認識你們.
所以這句話很小心.
我們不要去到有一天.
面對我們上帝的時候.
上帝說我都不認識你.
你是貴姓.
你是哪位.
這樣就很大件事.
所以我們今天就要警醒.
因為我們耶穌再回來的日子.
什麼時候我們是不知道的.
不過我們是會很深信耶穌會回來.
你看一下現在的世界.
你就會有一個這樣的.
你會越來越覺得.
這件事不是虛構.
我們想一下.
究竟我們是不是好像愚拙的同類.
如果將來耶穌基督對我們說.
我不認識你.
我今天可以做什麼改變.
所有現在要說的話.
是為了避免到那時候.
耶穌跟我們說這一番話.
所以我們現在要做什麼改變.
我們想到我們要改變.
我們才會行動的.
如果我覺得我不需要改變.
那我就不改變了.
我繼續是這樣的.
但如果我覺得神要我有改變的時候.
我們就應該來下定一個決心.
可能大家第一個反應就是.

$^{1001}$我老了.
就是廉頗老矣.
尚能凡否.
其實呢.
就是呢.
老不老呢.
老的字眼就沒有什麼定義的.
是不是.
這個很虛範的.
雖然現在呢.
以前有長者卡六十五歲.
現在弄個綠油卡六十歲.
你拿兩塊錢就不表示你老.
你明不明.
就是很看你怎麼看的.
是不是.
而且現在的醫學是不是發達很多.
人是長壽很多的.
八十歲九十歲.
都是很普遍.
你想想.
我現在當你六十歲.
那你可能是有八十歲命的.
你有九十歲命的.
或者九十幾的.
是不是很普遍.
不稀奇的.
真的不稀奇的.
那接下來的問題就是.
那你現在這段時間可以怎樣.
你可以怎樣.
你明不明我的意思.
你可以做很多的事情.
很多很多的事情.
說得不好聽.
現在這個時候.
有很多的.
很多人就不在香港了.
離開了香港.
但是.

$^{1041}$如果我們又沒有什麼.
太多的事情要去顧慮的.
就是子女又大了.
自己又退休了.
又.
是不是.
就是這樣的時間.
這個是不是一個好的機遇.
來去事奉神呢.
是不是.
這個是值得去想的.
你說我一向都不習慣做什麼.
其實這個不是問題.
我認識一個牧者.
他就是.
說名字也不怕了.
崇禎會的一個牧者.
他說我們崇禎會.
就是辦了一些課程.
給那些弟兄姊妹.
來去接受訓練.
接受了訓練之後.
就給他一個認可的資格.
在教會幫忙做事奉.
很威風的.
是不是.
他說第一屆也有二十多人參加.
多厲害.
我自己會想.
其實.
我就負責去呼喚人.
有一個心.
但是有不同的教會.
你去找這些人.
是由教會去認可他的.
或者給一些訓練他的.
或者一些.
你要分配的工作.
以及去督導他的.
其實.

$^{1081}$有很多人.
他本身的條件是適合的.
不過沒有人去叫他這樣做.
你明不明白我的意思.
或者.
沒有人想過他們可以做.
所以我自己會覺得.
其實方法有很多.
只要有很多人在.
他有個心願意去事奉神.
你就去尋求神的旨意.
給一些崗位給他.
給一些訓練給他.
不是你去選擇的.
當然你說.
如果選擇一些年輕又精靈的.
當然是最好.
但不是這個問題.
環境是這樣就這樣.
所以我希望行動起來.
就不只是一些年青有為的人.
而是說.
如果我不需要.
憂柴憂米的話.
我只是這樣說.
你又不用為你的事業去搏殺的話.
而且你又.
你有的是時間.
這個問題.
還有你有健康.
這個時候.
你願不願意行動起來.
你願不願意改變一下你的生活模式.
或者你想的東西.
我有些親戚是這樣的.
他們在疫情之前.
很喜歡去坐郵輪.
有些不是我的親戚.
有些朋友.
去的時候就要說.

$^{1121}$我去看北極光.
很威猛的.
但疫情的時候就不能說這些.
現在就開始有理.
現在又去通關去很多地方玩.
我自己.
我心底裡就說.
我讓你去遍全世界.
你是不是真的滿足呢.
你明白我的意思嗎.
有很多人說.
你叫我去哪裡都無所謂.
看來看去都是那些風景.
我沒什麼興趣.
我去見人而已.
會不會有些人這樣.
你可能覺得.
你是不是這麼牙刷.
你去過很多地方.
隨著人的心境改變.
你去加拿大去澳洲.
其實都是那些地方.
山山水水.
我其實沒什麼去.
我不是沒得去就說別人.
總之你會想.
有些人很想到處去.
但有些人就說.
看來看去都是那些東西.
究竟誰說得對呢.
我自己的看法就是.
其實你去得多.
你自己也會發覺差不多.
你拍很多照片回來也是高興一會.
我不會否定人去.
但我會問一個問題.
讓你去遍全世界.
你是不是很滿足呢.
是不是.
如果我們能夠看到神.

$^{1161}$在這個時代給我們一個託付.
一個挑戰的時候.
你又願不願意把時間放出來.
去為神做事呢.
不是叫你不要去旅遊.
不會這麼極端.
但不要只是想著退休之後.
我就是這種生活.
你明不明.
如果你想著生活是這樣.
很可能這個疫情.
就是一個很好的教訓.
這個疫情就是打破了我們很多的.
原有的計劃.
這個疫情告訴我們.
有很多東西不在你去控制掌握.
究竟有什麼要想下去呢.
我是逃避呢.
我負面呢.
還是說原來這是一個變幻的時候.
趁著我們能走得了.
我們就要去做全方位的工作.
剛才師母問我.
她又記得我的腳有事.
你們看我站了這麼久.
還能站得住嗎.
我就說我進步中.
你明白嗎.
原來我能夠站著唱詩.
如果行動不方便的人可以坐下.
其實我每次去教會敬拜.
我都很想坐下.
因為站著很累.
拉扯得很厲害.
能夠站著站了這麼久.
你明不明白我的心境.
能夠站著真的很感恩.
說到這麼長.
其實我早一兩年.
我的人向前了.

$^{1201}$原因我是站不直.
拉扯得很厲害.
所以我會體會到一件事.
能走能走都不必然.
我經常都想.
我太太她都是膝頭痛.
不知道各位電影節目有沒有.
去政府醫院看骨科醫生.
你知道趙亞易光.
那次去東區.
不是他那天東區有什麼問題.
而是現在普遍都是這樣.
他叫完之後.
你這個已經勞損得很厲害.
你去隔壁排隊換教.
然後就說.
你不要想著這麼快就行.
要排幾年.
你明不明白.
現在的情況就是.
現在的流程是這樣.
我就想.
你知道現在進步了.
換了教那些可以用久一點.
但問題就是.
我想.
即是.
如果不是很成功那怎麼辦.
這些想法.
有時我太太說.
那你坐輪椅吧.
很多地方輪椅不能用.
那怎麼辦.
你明不明白我意思.
趁著我們現在就好好調理自己身體.
趁著現在能走能走.
你就做多一點神的工作.
你明不明白我講什麼道理.
其實我和你一樣是病友.
那我希望透過今天的講道.

$^{1241}$讓我們去思想.
馬太福音裡面第25章.
十個同類的比喻.
其實很簡單的比喻.
但是一個很實在的問題.
我是可以歸類在哪個類別的呢.
我是應該歸類在愚拙的同類的類別.
還是我被歸類在聰明的同類的類別.
很簡單的問題.
你不需要回答我.
你是可以在祈禱裡面.
你去和上帝講.
我覺得我是屬於哪一類的.
你去聽聽上帝跟你講.
你實在是哪一類的.
願意我們是躺開我們的心.
去讓上帝跟我們說話.
各位你們對教會的承擔.
神是知道的.
我們一起祈禱.
我們親愛的天父上帝.
我們向你禱告.
因為你正如.
正是用聖經的話語來提醒我們.
我們活在一個耶穌很快回來的時代.
世界的變化.
周圍的事情是不穩定.
很多的事情我們去掌握不住.
我們再一次的去思想.
耶穌基督的福音.
是回答人心裡面很多的疑問.
但是我們自己.
是不是真的很安定呢.
我們是不是能夠有危機意識呢.
我們是不是看清楚問題的核心呢.
我們有沒有行動起來呢.
我們求主讓我們有一個信念.
我們只不過是普通的人.
我們未必做到很多事情.
我們付出的未必能夠解決問題.

$^{1281}$但是我們願意.
我們願意行動起來.
我們求你.
按你的心意來去成就人不能夠成就的事.
讓我們看見你是那位主.
我們只不過是微小的僕人.
但是你卻是願意使用我們.
當我們思想到神你是創造宇宙的神.
你說有就有 命立就立.
你根本不需要我們去為你做事.
你根本不需要.
但是你卻是願意.
讓我們參與你的聖功.
你願意我們和你一起去完成你的心意你的計劃.
這個是你對我們的肯定.
是你對我們的恩典.
就求你讓我們每一位弟兄姊妹心裡面.
都有這種渴慕的心.
來去回應我們的神.
是你自己親自發出的呼召.
阿們.
\newpage



\section{以賽亞書 1:1-31}
\label{sec:pT_5txip_oE}
\textbf{2023.05.07 《重拾對罪的警惕 重新領受神愛的管教 重燃悔改的動力》 以賽亞書 1\_1-31 (和修版) 講員 : 高淑芬傳道}
\newline
\newline
連結: \href{https://youtube.com/watch?v=pT-5txip_oE}{\texttt{https://youtube.com/watch?v=pT-5txip\_oE}} ~~~~ 語音日期: 2023-05-10
\newline
\newline
\hyperref[sec:_IeOxEq_8vc]{\small{< < < PREV SERMON < < <}}
~
\hyperref[sec:index_chronic]{\small{[返順時目]}}
~
\hyperref[sec:index_scriptual]{\small{[返順卷目]}}
~
\hyperref[sec:5NIuKMKgBWE]{\small{> > > NEXT SERMON > > >}}
\newline
\newline
以賽亞書 1:1-31
\newline
\begin{longtable}{cl}
\hline
\hline
章節 & 經文 (和合本修訂版)\\
\hline
1:1 & \begin{tabularx}{0.7\textwidth}{X} 當烏西雅、約坦、亞哈斯、希西家作猶大王的時候，亞摩斯的兒子以賽亞見異象，論到猶大和耶路撒冷。 \end{tabularx} \\ \\ \relax
1:2 & \begin{tabularx}{0.7\textwidth}{X} 天哪，要聽！地啊，側耳而聽！因為耶和華說：「我養育兒女，將他們養大，他們竟悖逆我。 \end{tabularx} \\ \\ \relax
1:3 & \begin{tabularx}{0.7\textwidth}{X} 牛認識主人，驢認識主人的槽；以色列卻不認識，我的民卻不明白。」 \end{tabularx} \\ \\ \relax
1:4 & \begin{tabularx}{0.7\textwidth}{X} 禍哉！犯罪的國民，擔著罪孽的百姓，行惡的族類，敗壞的兒女！他們離棄耶和華，藐視以色列的聖者，背向他，與他疏遠。 \end{tabularx} \\ \\ \relax
1:5 & \begin{tabularx}{0.7\textwidth}{X} 你們為甚麼屢次悖逆，繼續受責打呢？你們已經滿頭疼痛，全心發昏； \end{tabularx} \\ \\ \relax
1:6 & \begin{tabularx}{0.7\textwidth}{X} 從腳掌到頭頂，沒有一處是完好的，盡是創傷、瘀青，與流血的傷口，未曾擠淨，未曾包紮，也沒有用膏滋潤。 \end{tabularx} \\ \\ \relax
1:7 & \begin{tabularx}{0.7\textwidth}{X} 你們的土地荒蕪，城鎮被火燒燬；你們的田地在你們眼前被陌生人侵吞，既被陌生人傾覆，就成為荒蕪。 \end{tabularx} \\ \\ \relax
1:8 & \begin{tabularx}{0.7\textwidth}{X} 僅存的錫安，好似葡萄園的草棚，如瓜田中的茅屋，又如被圍困的城。 \end{tabularx} \\ \\ \relax
1:9 & \begin{tabularx}{0.7\textwidth}{X} 若不是萬軍之耶和華為我們留下一些倖存者，我們早已變成所多瑪，像蛾摩拉一樣了。 \end{tabularx} \\ \\ \relax
1:10 & \begin{tabularx}{0.7\textwidth}{X} 所多瑪的官長啊，你們要聽耶和華的言語！蛾摩拉的百姓啊，要側耳聽我們神的教誨！ \end{tabularx} \\ \\ \relax
1:11 & \begin{tabularx}{0.7\textwidth}{X} 耶和華說：「你們許多的祭物於我何益呢？公綿羊的燔祭和肥畜的油脂，我已經膩煩了；公牛、羔羊、公山羊的血，我都不喜悅。 \end{tabularx} \\ \\ \relax
1:12 & \begin{tabularx}{0.7\textwidth}{X} 「你們來朝見我，誰向你們的手要求這些，使你們踐踏我的院宇呢？ \end{tabularx} \\ \\ \relax
1:13 & \begin{tabularx}{0.7\textwidth}{X} 不要再獻無謂的供物了，香是我所憎惡的。我不能容忍行惡又守嚴肅會：初一、安息日和召集的大會。 \end{tabularx} \\ \\ \relax
1:14 & \begin{tabularx}{0.7\textwidth}{X} 你們的初一和節期，我心裡恨惡，它們成了我的重擔，擔當這些，令我厭煩。 \end{tabularx} \\ \\ \relax
1:15 & \begin{tabularx}{0.7\textwidth}{X} 你們舉手禱告，我必遮眼不看，就算你們多多祈禱，我也不聽；你們的手沾滿了血。 \end{tabularx} \\ \\ \relax
1:16 & \begin{tabularx}{0.7\textwidth}{X} 你們要洗滌、自潔，從我眼前除掉惡行；要停止作惡， \end{tabularx} \\ \\ \relax
1:17 & \begin{tabularx}{0.7\textwidth}{X} 學習行善，尋求公平，幫助受欺壓的，替孤兒伸冤，為寡婦辯護。」 \end{tabularx} \\ \\ \relax
1:18 & \begin{tabularx}{0.7\textwidth}{X} 耶和華說：「來吧，我們彼此辯論。你們的罪雖像硃紅，必變成雪白；雖紅如丹顏，必白如羊毛。 \end{tabularx} \\ \\ \relax
1:19 & \begin{tabularx}{0.7\textwidth}{X} 你們若甘心聽從，必吃地上的美物； \end{tabularx} \\ \\ \relax
1:20 & \begin{tabularx}{0.7\textwidth}{X} 若不聽從，反倒悖逆，必被刀劍吞滅；這是耶和華親口說的。」 \end{tabularx} \\ \\ \relax
1:21 & \begin{tabularx}{0.7\textwidth}{X} 忠信的城竟然變為妓女！從前充滿了公平，公義居在其中，現今卻有兇手居住。 \end{tabularx} \\ \\ \relax
1:22 & \begin{tabularx}{0.7\textwidth}{X} 你的銀子變為渣滓，你的酒用水沖淡。 \end{tabularx} \\ \\ \relax
1:23 & \begin{tabularx}{0.7\textwidth}{X} 你的官長悖逆，與盜賊為伍，全都喜愛賄賂，追求贓物；他們不為孤兒伸冤，寡婦的案件也呈不到他們面前。 \end{tabularx} \\ \\ \relax
1:24 & \begin{tabularx}{0.7\textwidth}{X} 因此，主－萬軍之耶和華、以色列的大能者說：「唉！我要向我的對頭雪恨，向我的敵人報仇。 \end{tabularx} \\ \\ \relax
1:25 & \begin{tabularx}{0.7\textwidth}{X} 我必反手對付你，如鹼煉淨你的渣滓，除盡你的雜質。 \end{tabularx} \\ \\ \relax
1:26 & \begin{tabularx}{0.7\textwidth}{X} 我必回復你的審判官，像起初一樣，回復你的謀士，如起先一般。然後，你必稱為公義之城，忠信之邑。」 \end{tabularx} \\ \\ \relax
1:27 & \begin{tabularx}{0.7\textwidth}{X} 錫安必因公平得蒙救贖，其中歸正的人必因公義得蒙救贖。 \end{tabularx} \\ \\ \relax
1:28 & \begin{tabularx}{0.7\textwidth}{X} 但悖逆的和犯罪的必一同敗亡，離棄耶和華的必致消滅。 \end{tabularx} \\ \\ \relax
1:29 & \begin{tabularx}{0.7\textwidth}{X} 那等人必因所喜愛的聖樹抱愧；你們必因所選擇的園子蒙羞， \end{tabularx} \\ \\ \relax
1:30 & \begin{tabularx}{0.7\textwidth}{X} 因為你們必如葉子枯乾的橡樹，如無水的園子。 \end{tabularx} \\ \\ \relax
1:31 & \begin{tabularx}{0.7\textwidth}{X} 有權勢的必如麻線，他的作為好像火花，都要一同焚燒，無人撲滅。 \end{tabularx} \\ \\
[1ex]
\hline
\hline
\end{longtable}
$^{1}$我們一起祈禱.
親愛的天父.
今天說的,聽的.
都是你的塞座前.
主啊,讓你的說話教導我們.
讓說的,聽的都同感一靈.
來遵從主你的教導.
以致到落實於我們行為裡面.
來調教,合乎你的心意.
主,幫助我們每一個.
我們禱告,奉主耶穌基督得勝的名字.
阿們.
剛才大家讀了整張以賽亞書.
以賽亞書第一章.
有一些講師說到.
整卷以賽亞書有66卷.
如果像看電影一樣.
要有一個預告片或是一個撮要.
就看第一章.
撮要整卷以賽亞書要我們領受的.
是甚麼?.
當然詳細請看今個月裡面.
繼續看以賽亞書.
以賽亞書是一卷.
曾經在死海裡發掘的古卷.
而整卷的書裡.
在新約聖經裡主耶穌.
都引用了很多.
以賽亞識日被上帝的靈感動的說話.
而且如果你看希伯來文.
或者不懂看的你去看色經書.
你會看到有很多的詩句.
有些比喻.
剛才你們讀的那一章.
已經充滿了很多的比喻.
可以說就算被斥責.
也很有藝術.
我覺得我們這些要去教導的人.
以賽亞書是一個很需要去學習的地方.
今日我的題目.

$^{41}$因為上星期國務司給我的靈感.
原來題目可以很長.
所以就這麼長.
今日聽完之後.
最重要是要記得這三個動作.
重拾對罪的警惕.
重新去領受神愛的管教.
重燃起我們悔改的動力.
在以賽亞書整卷裡.
我們已經打好頭號.
就是以賽亞先知所寫的.
也被譽為以賽亞書有66卷.
也是最長的.
不是啊 高傳道.
詩篇不是更長嗎?.
但是一氣阿誠整卷來看.
以賽亞書66卷是最長的一卷.
由一位先知去領受神的話語.
所以我們紅印堂不是說讀經嗎?.
你開始讀以賽亞書了嗎?.
一路去看.
你就會更加看到上帝的心腸是怎樣.
當時以賽亞寫的時代.
以色列國和猶太人.
是分成兩個大局面.
我們會稱為南國和北國.
這是在所羅門死後分開的.
而在以賽亞的時代.
從第一節可以看到.
已經經歷了四個皇朝.
其實不止的 第五個皇朝馬來西亞.
他沒有說.
因為他在馬來西亞時代.
他是不幸的歸住.
一會兒我會告訴你他怎樣死.
是很慘的.
但是以賽亞書一生的.
他的死亡也是因為忠於上帝的話.
要去教導忠言.
惹了殺身之禍.

$^{81}$他仍然要這樣去做.
在第一章我們看到.
以賽亞領受的意象.
是說到整個以色列.
不止是南國 不止是北國.
你們每一個都要聽.
因為這是從神而來的意象.
是上帝藉著先知.
要告訴他的百姓 他的子民.
今天可能會繞著手.
都是說以色列人的不是.
由創世記一直看.
我們都會搖頭.
這樣也行嗎?.
這群以色列人是上帝所愛的子民.
但他們竟然可以做出一些.
我們2023年的我們.
我們看完也會搖頭.
但請大家記住.
聖經的說話寫出來的時候.
不單是一個歷史這麼簡單.
也是讓我們從中學習.
去警惕 去悔改.
特別是他所說的這個意象.
當我們讀到第六章.
就會很詳細地看到.
他得到的意象是怎樣.
不單止是聽見.
更加看到很多圖畫.
看到這個意象是從神而來.
他要加強從上帝而來的權威性.
我們剛才讀了這段經文.
如果第二節.
如果讓在座的父母去讀.
我養育我的子女.
把他們養大.
他們竟然背叛我.
你的感受又是怎樣呢?.
我們是去教導人的.
我們很多時候都要去教導.

$^{121}$教導前其實也很不容易.
很多掙扎.
有沒有點錯了呢?.
怎樣去指導的時候.
其實他們不單止不聽.
可能還要來攻擊你.
所以我認為以賽亞寫得很好.
他竟然用牛來認識他的主人.
如果2023年.
可能貓狗也認識主人.
但當時連驢也覺得是差一級.
也不認識主人.
你們這些以色列民竟然不認得我.
不知道我是誰.
是非常違反了上帝創造的秩序.
違反了神賜予的自然天性.
我們在街上經常看到很多.
特別是年輕的夫婦.
不生了.
看到他們也會說.
你對狗和貓這麼有愛心.
不如生一個吧.
不生了,養狗養貓比吃叉燒好.
的確,狗會向主人搖頭擺尾.
向我們這些陌生人.
我們想來打他,他就吠我們.
我們是上帝,甚至主耶穌基督.
是被上帝差派下來.
去拯救我們每一個的我們又如何.
我們竟然向著上帝亂吠.
求主饒恕我們.
以色列人所犯的罪.
是一個藐視以色列的聖者.
他們離棄了耶和華.
藐視以色列的聖者.
這個藐視好像完全對你的話語是一種侮辱.
你不要以為你的說話出來的才是藐視.
你的說話出來的才是侮辱.
今天我們基督徒.
我們是神所愛的子女.

$^{161}$我們在世的時候,我們如何行事為人呢?.
我有一個姐妹.
她請了一個菲律賓的基督徒.
做她的家務助理.
每逢星期日,這個家務助理都會回教會.
但是她的工作表現如何呢?.
她主要是照顧她的媽媽.
新鮮菜自己吃.
昨天的冷飯菜汁給婆婆吃.
要出門了.
一家人出門,所以她的女兒也會出門.
女兒已經有女兒了.
又小又大.
家務助理一起出門.
有時候有些路,要用高低搭幫忙扶一扶.
家務助理完全放手.
我們說我們也要放手.
但當你是職責的時候.
不是你放手.
她就放手,之後她也沒有想再推.
輪到我的姐妹去推.
有些事情,經常向她的牧者投訴.
我的姐妹覺得她表裡不一.
就查她到底回哪家教會.
後來才知道,她掛念基督教教會.
其實是二段.
縱然我們今天回教會也是正統的教會.
但我們的行為出來的.
又是否正統教會上帝主耶穌基督的門生作為呢?.
我們求神幫助我們.
不要讓我們的言行.
都令到上帝這位神聖的主夢想羞辱.
這一群是基督徒.
我們求主讓我們女兒也是皈依.
以色列民不單一次錯,兩次錯.
身上已經有著很多很多的傷口.
是非常不健康.
你會看到,如果我們流血.
後來含糖,我媽媽早前.
因為我媽媽有糖尿病.

$^{201}$上星期不知道為什麼.
她的手突然間腫起.
腫起不久就發現有膿在裡面.
當我們要幫她塗藥的時候.
因為輕輕按摩,也有膿出來了.
都是要有流血,不乾淨的血.
或者是一些膿,也要弄乾淨.
去包紮,然後醫治.
但是以色列民卻不是這樣.
一次,兩次的管教.
他們仍然繼續犯下去.
所以我們求主幫助的時候.
我們也會經歷過一些不容易的教訓.
但是我們是否繼續這樣走下去呢?.
我們真的要去到上帝的面前.
去清理,去潔淨.
我們更加看到他們的土地變成荒廢.
聖真有被火焚燒.
我們最深刻的印象就是.
所羅門王那個時代是很厲害的.
不同的國家都要來朝拜他.
但是一代一代下去.
他們在這個時間,在北角都要慘遭.
已經被擄了一截.
再遭擄了所有的民.
現在南角的你們.
你們看到這些的時候.
你們還不警惕嗎?.
我們在2019年的時候.
我們經歷過要戴口罩被隔離的日子.
我們整個世界.
整個世界因為這個小小的病菌.
以至整個世界停了.
好像停了一樣.
大家都困在家裡.
不知道當時.
你看著彌敦道整天都這麼多車.
突然間不會再塞車.
地鐵整天都爆.
要塞進去.

$^{241}$現在不用再塞了.
因為沒什麼人在街上.
我特別在香港有沙士的時候.
我在泰國.
那一刻我已經很深刻.
這個是香港.
我看那個新聞的時候.
這個是香港來的嗎?.
那時候我覺得有慶幸.
為什麼這麼慶幸去到泰國.
因為我自己是香港人.
不可以和香港人一起面對災難.
有一種感受在當中.
香港不是這樣的.
為什麼現在百貨公司沒人走.
街上沒人.
完全很害怕.
那個時候好像20.
我們剛剛經歷過.
接近三年的時間.
以色列已經不是一次被人廢了城.
被人侵略.
他們很多很多很多次.
但是他們仍然是一而再.
再而三的犯罪.
他們覺得得罪神沒什麼所謂.
神都會救我的.
會繼續來拯救我的.
但是你知不知道.
神仍然保存著他的子民.
是因著上帝那份愛.
那份憐憫.
今天你和我在香港.
我們仍然可以自由自在地生活.
如果我們有健康.
我們還可以好好地工作.
或者來服侍神.
這些完全都是上帝給我們的憐憫.
我們不要在神的面前驕傲.
不要覺得得罪神是沒什麼問題的事情.

$^{281}$我們要對這個罪有警惕.
不要覺得是理所當然.
在一月份的時候.
在大阪裡面.
有一個70歲的男子.
他撿到一個錢包.
錢包有多少錢呢.
我先看一下.
有.
有這個8萬6.
86萬日圓.
即是香港多少錢呢.
原來大概5000元.
他撿到這個錢包.
於是他就把錢包交回警方那裡.
路不拾遺.
把錢包連同聯絡地址.
給了警方.
方便警方繼續聯絡他.
誰知道幾個小時後.
警方已經找到這個失主是誰.
原來是一位50歲的男性.
過了幾天.
這個70歲的男子撿到錢包.
為什麼一個電話都沒有來.
一點多謝都沒有.
於是他就拜託警方.
打電話給這個失主.
這個50多歲的失主說.
我很忙.
就這樣先收線.
於是警方同樣回覆.
這個70歲的路不拾遺的男士.
這個男士.
他越想越激氣.
你們會不會這樣呢.
激氣什麼呢.
不知道你們和他不同嗎.
於是他就真的要去告這個人.
因為.

$^{321}$這樣都可以告嗎.
他沒有打電話給你.
原來日本是可以告的.
為什麼呢.
因為日本有一條法例.
你撿到的財產.
你是要有5\%至20\%.
是要給那個幫你撿回來.
因為你的錢本來什麼都沒有.
你就應該給5\%至20\%的錢.
來作為一個答謝.
答謝這個幫你撿回錢包的人.
這個70歲的老先生.
他就說.
我去告他.
我要告他.
他要賠償給我20\%的答謝費.
因為.
剛才說了.
警察他們有條例.
你要這樣做.
如果這個撿到的東西.
沒有人在6個月裡面認領的時候.
這個東西可能就會被這個撿到的人.
或者被充公的.
這個就視乎警方去判決.
好了.
終於去到上課.
上去法官那裡.
這件事已經去到4月13號.
所以這個新聞都是最近我們聽到的.
他就上去.
就要控告他.
他現在應該要賠償.
他現在要給我.
我要有這個4萬4千.
是這個8萬6千的.
不是.
86萬.
他要44萬.

$^{361}$20\%的數字.
大概是5000港幣.
不是.
大概是多少呢.
你自己去看吧.
因為他給我的錢.
那些數字我就非常不懂.
那樣.
這個老先生也說.
我其實不是要那些報酬.
我是要一句多謝.
他連多謝都沒有聲音.
我就要控告他.
結果這件事庭外和解.
這位50歲的男士.
也不是給了5\%左右的錢.
給這個撿的70歲的伯伯.
在法庭裡面.
法官問你還有什麼要說.
他才說第一次的對不起.
你看到原來說對不起.
不說對不起.
都要上庭.
所以我們今天.
很多時候我們在一個很忙很忙的.
香港的節奏裡面.
我們都將這些文化帶來教會.
有時候很匆忙很匆忙.
可能突然之間很期盼.
見到你.
不只是打聲Hi這麼簡單.
很可能還想多聊幾句.
可能我們還有些事情忙.
或者不能夠跟他關注.
或者我們有時候崇拜.
我們也試過很多次.
突然間叮叮.
或者播放一些特別的音樂.
有時候我們可能會看過去.
但是有時候我們都需要.

$^{401}$當事人我們要警惕一下.
究竟我是不是按錯鍵.
或者請教一些認識的人.
怎樣調整手機的聲音.
因為當事人可能以為調整了.
其實是沒有調整.
對罪的警惕.
不是說對別人的罪.
在指.
你要記住.
當人指人的時候.
四隻手指是指向自己.
對別人的罪.
我們去體諒.
或者再勇敢一點.
去看他的罪.
為什麼會這樣.
是不是他的心靈受傷.
還沒有醫治.
似乎你這樣一句話.
可能就會像刀一樣刺著他.
所以求神幫助我們.
我們要對自己的罪.
哪怕是一個忘記說的對不起.
今天上帝對我們的恩情這麼重.
但我們對上帝有沒有感恩呢.
我們可以很僥倖地.
過了接近三年新冠疫情.
今天我們可以重見天明.
第一次不用戴口罩.
講道.
但是.
我認識的朋友.
因為打了新冠疫苗.
本來人生已經去到退休.
他的錢.
退休後可以慢慢地生活.
過活.
誰不知.
這個疫苗的針.

$^{441}$讓他內凍症.
整個人癱瘓了.
今時今日.
不單止積蓄沒有了.
因為要治療這個病.
甚至住的樓房都要賣掉.
人卻像植物人.
存留在世上.
所以我們今天.
我們有的健康.
我們有的一切.
我們更有的.
救贖的恩典.
不要輕看.
不要覺得是理所當然.
我們更加要留心.
我們有沒有得罪了上帝.
以至我們侮辱了上帝.
神同樣對我們的生命很體重.
這裡一直下去.
你會看到.
以色列人他們做了很多事.
很棒.
建制要供奉的公牛,羔羊,山羊.
什麼都齊了.
要祭司去做這些事嗎?.
打點好了.
但上帝一而再再而三.
非常之厭煩.
你們的手是沾滿了血.
就算在新約裡面.
主耶穌也有說.
你今天去祈禱.
但你心裡還很憎恨你的弟兄姊妹.
對他不原諒的時候.
你的祭.
你所獻情的.
上帝喜悅了.
所以我們求神幫助我們.
多點去檢視自己.

$^{481}$縱然你舉手的祈禱.
你懇切的祈禱.
上帝是想聽的.
但你仍然不斷地一而再再而三.
很多事情來犯罪.
上帝對罪.
不是對你厭惡.
是對罪非常之恨惡.
祂不要罪還要在你的生命裡.
所以我們求神幫助我們.
我們今天已經是上帝的兒女.
讓這些罪離開我們.
有些罪一而再再而三.
要求神用神的力量.
慢慢地去調節和改.
就好像十一週年的時候.
我們的聯合詩班.
跳來跳去.
很多個回合的練唱.
我們都未能跳得好.
我們很想跟著琴跳得好好的.
但就是不行.
當時如果你有做詩班.
心裡可能會有很大的感受.
已經還有五分鐘.
還是十分鐘就要崇拜了.
我們還唱成這樣.
自己的耳朵都過不去.
怎樣獻給上帝呢?.
當我們.
我相信在那一刻.
我們每一個的心.
都很想將最好的獻給神.
以致出乎意料地.
在集體詩班裡.
唱得最好的.
絕無僅有的.
就是那一次.
就是那一次要獻給上帝.
那一次.

$^{521}$所以你的心要著.
對這些罪.
不要說無所謂.
你要對抗那種心情.
靠著神.
不知怎樣.
有一天.
這件事會撇下.
戒毒的是這樣.
戒賭的也是這樣.
他沒有告訴你.
他無法告訴你.
我有什麼秘訣.
有什麼方法.
方法就是.
依賴神.
懇求神.
上帝要我們知罪.
是因為我們不知罪.
我們就不知道我們要乾淨.
我們更加要求神幫助我們.
在我們周圍的人.
特別是這個末世的時代.
愛心越來越少.
因為我們都領過東西.
好心幫人.
如果你在內地.
隨時就變成.
是你撞到他跌的.
或者你本來是幫他.
變成你是受害者.
是很多不容易.
今天特別喜歡.
騙騙騙的時間.
我們去幫一個人.
他可能不好.
他可能騙你.
在這裡.
很多時候我都會去幫一些.
不懂廣東話的泰國人.

$^{561}$去覆診.
或者他的孩子來了香港.
要教他們廣東話.
有些不懂聽泰文的.
他們說不要相信了.
他可能騙你的.
現在免費.
一會兒要勞你一大筆.
真的不明白.
教他的孩子說泰文.
又要幫他翻譯.
怎會不懂聽呢.
他是想告訴我.
還是告訴他呢.
不知道.
但是上帝知道.
我們是想幫他.
所以有很多事情是很不容易.
但是更不容易的.
我們洗衣服都試過.
夏天喜歡穿白色衣服.
糟糕了.
不要說是豬紅.
不要說是吃茄汁.
弄得茄汁醬難洗.
很少的可能是果汁.
橙色的.
你都覺得很難洗到白色.
何況我們的嘴.
好像豬紅.
不是吃的豬紅.
是豬紅.
紅到不得了.
單顏都是那種紅.
但是上帝有他的大能.
讓這些罪.
這麼重的罪.
都可以變成雪白.
變成樣木.
但是你想洗白白嗎.

$^{601}$上帝不會逼你的.
你可以自己去選擇.
你今天是跟從神.
聽神的說話.
還是繼續叛逆.
你可以選擇.
神的指紋.
真的需要除掉我們的惡行.
停止我們作惡的一些事情.
我們今天跟監獄裡的人.
其實是一樣的.
只不過我們幸運了一些事情.
我們的罪的意念.
沒有被人捉到.
還有我們不敢.
想了很多壞事.
有時候不敢做.
不是不想做.
有時候可能不敢做.
所以我們跟他一模一樣.
更加有一句說話.
教會其實都是開放給所有罪人.
今天站在這裡.
這個也是一個罪人.
我們都是需要上帝的管教.
上帝管教不容易.
不過還是一份福氣.
最近有個新聞.
很新鮮.
就是有一個25歲的年青年輕人.
在日本吊頸自殺.
原來他在幾個月前.
在葵芳搶劫隔離鑽石.
他不知道怎樣可以拿到隔離屋的鑰匙.
他悄悄地進屋爆竊.
誰知那個戶主回來了.
回來後很害怕.
害怕到要把人打死.
你不要以為打死的人很勇敢.
其實他就是因為害怕.

$^{641}$我們在泰國一年會有一次.
很開心一起退休.
一年才一次.
我認識很多不同的宣教士.
特別是我們單身的.
聽完他們說可能會覺得很恐怖.
我很害怕那些驗射.
一見到那些驗射會怎樣.
我會打死他.
我都很害怕.
我又不敢打死那隻蟑螂.
知道那隻蟑螂是翼蟲.
我都不敢打死.
但驗射肥嘟嘟很滋潤.
打死他.
真的很不容易.
是因為什麼原因.
你不要以為殺人的人膽大包天.
他就是害怕被殺死.
看到他做了罪行的人.
覺得沒有路可走.
不殺死他不行.
他很聰明.
他去了日本.
這個年輕人.
我看的時候覺得.
他還可以教.
他知道罪行.
有人整個集團式經營.
斬了一個人.
幾十個.
個個都不認罪.
沒有這樣的事.
有很多人犯了罪.
不知道罪行.
否決.
否認.
這個年輕人.
我相信.
第一香港沒有死刑.

$^{681}$他真的捉到回來的時候.
最多是坐很多年監獄.
但在監獄裡面.
都有監獄的生活.
因為福音傳到監獄那裡.
昔日監獄裡面.
有些人.
他們都有他們的見證.
他們信了主之後.
他們唱的歌很有感染力.
他們講的見證.
亦都很令人觸動.
大大去嘆為觀止.
上帝的作為實在是大.
連這個人.
他可以這樣改變.
所以當神還要管教我們的時候.
我們求神給我們.
警惕著我們的罪.
接受上帝的管教.
在這裡不如你讀.
21至23節.
請.
.
你會看到.
這是神的指紋.
神要控訴.
你這群不是外邦人.
你是我所愛的指紋.
不單止21至23節.
一直都有講到.
你們一直罵他們之餘.
你們真的要快點去練正.
那些渣滓.
請你除去.
要歸正.
叛逆和犯罪一定會敗亡.
上帝明明的去講.
但是那個願意去上帝的人.
他會有不同的一件事.

$^{721}$以色列國最後連南國最後的防守.
亦都是被擄.
但是他們在這個時間當中.
他們似乎還未得到教訓.
所以上帝繼續要透過不同的國家.
徹底地去管教著以色列國.
用各種的方法去練正以色列國.
要除去他們的渣滓.
神為了練正我們.
是要用很多不同的環境去管教我們.
除去我們生命那種不好.
是不得有著神所喜悅的東西.
所以我們求神幫助我們.
在過程中我們真的要不斷地尋求神.
第一,讓我們有知罪.
第二,讓我們回過頭來.
來接受上帝給我們的悔改機會.
今天還有機會.
神沒有錯的.
是有公義和慈愛.
這是千真萬確的事.
若果不是的話.
就好像羅馬書一樣.
任憑,隨便你.
在地鐵或輕鐵.
有些小朋友來到時.
會亂亂的.
不受控制.
我們在旁的小朋友.
他們還未到我們身邊.
還可以忍受一下.
但他們的父母.
就像喇叭筒.
已經把他們拉著.
或是捉著.
因為這是他們的子女.
他們不想子女一開門就衝出去.
他們要極力挽救他們.
來幫助他們的子女.
一把.

$^{761}$做媽媽的,是沒有儀態的.
大家過去經歷過的日子.
現在很有共鳴.
本來還要被丈夫不明白.
我和你拍拖的時候.
很多時候粵語片都說.
你是很溫柔的.
好像我妹妹經常投訴.
你當然要和妹夫說.
你當然要說我什麼.
但不大聲你們聽不到.
就是這樣.
男士也是.
在這個情況下.
可能下班.
下課.
穿西裝.
他都不顧情面.
因為要管教子女.
捉著跑.
捉回子女.
很多這樣的鏡頭.
我相信你們日常生活都會看到.
為什麼?.
緊張他.
我們呢?.
我們冷血嗎?.
不是,因為那不是我的子女.
我不會揭書底.
你去捉他.
教他正經的.
不會.
今日上帝.
倘若有些管教的時候.
求神幫助我們.
上帝對你還未放棄.
還很希望你.
在神的愛裡面.
默會這麼久.
做宣教士這麼久.

$^{801}$要告訴別人你做得不對.
對方未必會改.
但我們很多事.
會被人覺得我們很挑剔.
很不容易的.
做教導的人.
有一次.
我不說教導大人.
遇到的問題.
教導小朋友.
小朋友在我們的壁報.
我們都會在泰國裡面.
推一些泰國人.
有些人都去做宣教.
我們為他祈禱.
所以有他的樣子.
這個小朋友.
用按釘.
打打打打打打.
打到宣教士面目全非.
我看到.
我當然要說他.
接著被他媽媽說.
高傳道.
你說我兒子.
請你先跟我說一聲.
否則.
你就像蓋著我.
蓋著我的臉.
打打打打.
我說是這樣的嗎.
這是你們泰國文化嗎.
他說是.
我遷就這種情況.
我問了很多泰國人.
很多都不是.
各有各的文化.
可能是他的家教.
你想想.
教一個小朋友.

$^{841}$也有這樣的.
何況是一些大人.
有很多事情.
他們就會怎樣.
但是.
很多時候我都會妥協.
我不想做壞人.
一個家庭.
剛才女士們這麼有共鳴.
你們當然是做惡人.
那個好人當然是爸爸.
做惡人.
做壞人.
真的很不容易.
你罵子女.
子女最不喜歡的是你罵的.
做好人的會比較黏著爸爸.
同樣的.
我們去教導弟兄姊妹.
同樣會有這樣的情況.
上帝不要再說了.
任憑他們.
或者不是叫任憑.
上帝你透過你的說話.
你自己跟他們說吧.
不要叫我跟他們說.
我不想再做壞人.
有時候.
很多時候.
有很多攻擊.
但是.
上帝.
一次又一次說.
不是的.
你改.
叫他們改.
告訴他們.
於是我也堂服.
上帝好.
我為你願意做壞人.

$^{881}$正如在座的母親.
你也願意做壞人.
這就是上帝的心腸.
上帝的救贖.
以賽亞先知.
他的出生.
在烏西亞.
臨朝.
當時在王的時候.
他經歷了.
這麼多代的皇帝.
好的去說.
勸勉他們.
好好的去教訓他們.
但是.
有些國王.
當然對他不客氣.
特別是去到馬來西亞的時候.
因為有些.
底線.
馬來西亞.
如果你不再聽我說.
還要說上帝的話.
你就會被處分.
以賽亞.
不要受這個.
不要屈服在這個.
不要屈服在人的面前.
他只是屈服在上帝的面前.
於是馬來西亞.
這個王.
將他處死.
是怎樣的死呢.
雖然聖經沒有直接說.
是怎樣的死.
但是有些聖經學者.
很相信.
在希伯來書.
第十一章三十七節說到.
被割死的.

$^{921}$就是.
以賽亞先至.
以賽亞的成長.
是一個皇族.
是一個貴族.
他因著上帝的呼召.
他向這群.
在還的.
以色列人.
很多的勸告.
很多的責備.
為的是要他們.
回頭過來.
但是最後的時候.
以賽亞.
卻是.
這樣的.
安息住.
但是我們相信.
我相信.
他當時也會很相信.
上帝仍然.
站在他那一邊.
有著這一份.
我們求神給我們.
都有著這個悔改的動力.
讓我們真的.
我相信.
洪恩嘉這段時間.
我們很需要.
一起去黏著神.
來看神.
不同的大作為.
上帝.
不單止要管好.
我們這群人.
是讓我們這群管好的人.
成為燈.
成為.
那個炎.

$^{961}$在這個世代裡面.
成為上帝一個活生生的見證.
我們求神幫助我們.
下一次.
如果有基督教.
因為在上個星期.
兩個星期前.
也有一些基督教的.
聚會.
在大禮堂裡.
一起去.
敬拜或是去祈禱.
很多.
其實也深知這個祈禱會.
如果是報道會.
還會看兩下.
可能是新朋友.
你們有沒有經歷過.
但是那個是祈禱會.
是這樣的.
他回去.
願意不在安坐家中.
去到那麼山長水遠.
去到這個聚會裡面.
但是.
爭.
爭位坐.
爭按電梯.
爭的.
這樣.
裡面會覺得怎樣呢?.
這些都是上帝的子女.
我好像不認識他們.
有沒有這樣的感受呢?.
求神幫助我們.
洪恩嘉一站出來.
是一個好的見證.
因為我們全部都是上帝的兒女.
今天.
你可能會聽到一些.

$^{1001}$管教或是什麼.
這是上帝還是很重視你的.
我們願意.
今天在街裡面的人.
因為今天報道隊要出去.
我希望他們同樣見到你.
都是要榮耀神.
求神幫助我們.
幫助我們.
今天上帝.
還沒有放棄我們.
有幾類型的人.
知道自己有罪.
他就會遠離.
我不回教會了.
知道自己有罪.
不要重複.
撒旦的說話.
撒旦當然很開心.
你不要回教會了.
不要亂回教會.
你不會被攻擊.
安享.
安享這個位置.
不是的.
求神幫助我們.
我們越有罪.
求神叫我們.
更加親近他.
朱紅的血.
上帝可以令你的衣服染上這些顏色.
還要白如雪.
白如楊毛.
與我們言行一致.
跟從一切.
一起祈禱.
親愛的天父上帝.
多謝你.
多謝你對我們不離不棄.
你仍然願意.

$^{1041}$管教我們.
派你的獨生子為我們贖罪.
四月份.
過了復活節.
讓我們更加深思.
上帝對我們那份愛.
那份不離不棄.
我們落在被管教.
或在際遇.
不容易的時候.
主啊.
求你.
讓我們分辨.
這是我們的罪.
帶來的結果.
還是世界的罪.
影響到我呢?.
求主讓我們認清楚.
更加去審測.
我們自己.
倘若仍然有罪的時候.
求神幫助我們.
讓我們.
是呀.
我們有罪.
但我們不要像彼得那樣.
主啊 離開我 我是一個罪人.
主不想離開你.
祂只想你的罪.
離開你.
求神幫助.
讓我們.
拿走這些罪.
很難拿走的.
求神讓我們.
因著你的力量.
去改變它.
或是我們的言語.
或是我們的態度.
或是我們的行為.

$^{1081}$總要叫我們.
繼續親近你.
親近你的時候.
讓我們蒙著.
你的救贖.
已經是讓我們.
乾乾淨淨.
求主幫助我們.
叫我們.
緊緊的依靠你.
祝福洪恩嘉的每一位.
我們奉主耶穌基督的名字求.
阿們.
謝謝大家.
\newpage



\section{馬太福音 15:21-28}
\label{sec:5NIuKMKgBWE}
\textbf{2023.05.14 《一位祝福孩子的母親》馬太福音15\_21-28 (和修版) 彭秀芹導師}
\newline
\newline
連結: \href{https://youtube.com/watch?v=5NIuKMKgBWE}{\texttt{https://youtube.com/watch?v=5NIuKMKgBWE}} ~~~~ 語音日期: 2023-05-18
\newline
\newline
\hyperref[sec:pT_5txip_oE]{\small{< < < PREV SERMON < < <}}
~
\hyperref[sec:index_chronic]{\small{[返順時目]}}
~
\hyperref[sec:index_scriptual]{\small{[返順卷目]}}
~
\hyperref[sec:XkIka_Z9Y3w]{\small{> > > NEXT SERMON > > >}}
\newline
\newline
馬太福音 15:21-28
\newline
\begin{longtable}{cl}
\hline
\hline
章節 & 經文 (和合本修訂版)\\
\hline
15:21 & \begin{tabularx}{0.7\textwidth}{X} 耶穌離開那裡，退到推羅、西頓境內。 \end{tabularx} \\ \\ \relax
15:22 & \begin{tabularx}{0.7\textwidth}{X} 有一個迦南婦人從那地方出來，喊著說：「主啊，大衛之子，可憐我！我女兒被鬼纏得很苦。」 \end{tabularx} \\ \\ \relax
15:23 & \begin{tabularx}{0.7\textwidth}{X} 耶穌卻一言不答。門徒進前來，求他說：「這婦人在我們後頭喊叫，請打發她走吧。」 \end{tabularx} \\ \\ \relax
15:24 & \begin{tabularx}{0.7\textwidth}{X} 耶穌回答：「我奉差遣只到以色列家迷失的羊那裡去。」 \end{tabularx} \\ \\ \relax
15:25 & \begin{tabularx}{0.7\textwidth}{X} 那婦人來拜他，說：「主啊，幫幫我！」 \end{tabularx} \\ \\ \relax
15:26 & \begin{tabularx}{0.7\textwidth}{X} 他回答：「拿孩子的餅丟給小狗吃是不妥的。」 \end{tabularx} \\ \\ \relax
15:27 & \begin{tabularx}{0.7\textwidth}{X} 婦人說：「主啊，不錯，可是小狗也吃牠主人桌上掉下來的碎屑。」 \end{tabularx} \\ \\ \relax
15:28 & \begin{tabularx}{0.7\textwidth}{X} 於是耶穌回答她說：「婦人，你的信心很大！照你所要的成全你吧。」從那時起，她的女兒就好了。 \end{tabularx} \\ \\
[1ex]
\hline
\hline
\end{longtable}
$^{1}$鄧小平安.
今天大家都知道什麼節日.
母親節.
我祝各位母親節快樂.
今日的題目是.
一位祝福孩子的母親.
我也有一位.
祝福我的母親.
在我心目中.
我媽媽是一個很堅強的女人.
因為在我十歲的時候.
我爸爸去世了.
他獨力養大了七兄弟姐妹.
家裡小時候是耕田的.
天還沒亮就要去耕田.
八點前就要將割的菜.
拿去合作社.
小時候住在鄉村.
準時買菜回家.
照顧七兄弟姐妹.
買完菜回家.
就繼續耕田工作.
日以繼夜.
沒有大假.
沒有病假.
年終無憂.
夏天就寒冷.
就去耕田.
夏天就汗流浹背.
冬天就很記得.
我媽媽的手.
要割菜.
紙紙黑黑.
因為太冷了.
家裡不富有.
但我媽媽有個原則.
每一餐都有三餐.
有菜有肉有湯有水果.
所以令到我都很健康.
很健康地成長.

$^{41}$所以有一句說話.
有媽的孩子像個寶.
我心目中的媽媽.
是個偉大的女人.
剛才說過.
她靠自己雙手.
供書教學.
家裡清貧.
但我從來沒有因為家裡清貧.
而自卑.
因為我體會到.
有一個很愛我.
為了我很努力的媽媽.
我記得每年新年.
媽媽都會堅持.
買新衣服給我們.
新鞋新內衣.
家裡沒有錢.
但她仍然堅持.
新年要穿新衣服.
才有個祝福.
所以我媽媽都給了我.
很多童年回憶.
她賺錢不多.
但她很樂意.
將她賺的錢.
花在我們身上.
在我心目中的媽媽是個很聰明的女人.
在我小學六年級.
我媽媽和我姐姐.
回教會 信了主耶穌.
我媽媽是不懂得寫字.
一字都不懂.
我小時候跟她去銀行.
打交叉或吸印.
但她信了耶穌.
回教會.
聽到別人讀聖經.
她很羨慕別人去看聖經.
所以她就開始.

$^{81}$學看聖經.
但她一字都不懂 連自己的名字都不懂.
神給她有智慧.
我媽媽看到那邊.
是我媽媽的聖經.
看得破爛.
因為媽媽每天堅持.
每天都會看聖經.
我讀書的年代.
媽媽經常煩我.
煩我什麼呢?.
這字怎麼讀?.
有時候我們都會生氣.
青春期.
但她就說:你先教我吧!.
那種肉緊.
對於愛神.
旁邊有媽媽的見證.
我媽媽不單祝福我.
她因愛主.
認識主.
祝福很多人.
我媽媽在2020年初.
疫情開始.
被天父接了.
她安息主懷.
我媽媽不是知名人士.
只是一個村府.
不懂事.
我記得我媽媽的葬禮.
我媽媽出席.
當時疫情.
我看到神祝福我媽媽.
因為我媽媽很愛主.
我記得我媽媽.
你見她戴著厚眼鏡.
因為她有白內障.
做了手術.
她從肚子看聖經.
看到視力不好.

$^{121}$真的看不到.
因為年事過高.
她聽聖經.
從走路回教會.
坐輪椅.
她也繼續回教會.
我看到神祝福我媽媽.
因為神祝福我媽媽.
所以祝福了我.
我很感恩.
今天.
不好意思.
經文.
是講一位媽媽.
在21節.
耶穌離開.
推羅西頓境內.
遇見一位.
迦南婦人.
在這個地方.
這位婦人哭著說.
主啊大衛.
知子可憐我爸.
我的女兒被鬼纏得很苦.
這位媽媽.
似乎有點困難.
但特別有提到一個地點.
推羅西頓境.
外邦人的地方.
一個迦南婦人.
原來迦南.
和猶太人是世仇.
這段聖經交代了.
耶穌.
專門去找外邦婦人.
這位媽媽.
很有趣.
她喊主啊大衛知子可憐我爸.
她有一個請求.
但這個請求之前.

$^{161}$她認定耶穌是主.
她說可憐我爸.
這個可憐.
聖經說是憐憫.
在舊約聖經中.
用憐憫去形容.
媽媽.
懷著嬰兒的肚子.
沒有福.
聖經有說.
在舊約聖經.
以賽斯49:15.
婦人那能忘記她吃奶的孩子.
而不憐憫她.
復生的兒子.
在座當過媽媽.
當你知道自己.
有了嬰兒那一刻.
你已經和嬰兒連繫了.
是無法告訴你老公.
怎樣連繫.
總之就是連繫了.
你會感受到.
她在你肚子裡.
這個連繫包括.
對嬰兒的思念.
對嬰兒的記掛.
對嬰兒.
有很多計劃.
很多想法.
事事以嬰兒優先.
這就是一個連繫.
略甚對人的憐憫.
就好像媽媽和嬰兒的連繫.
是分不開的.
我女兒在外國讀書.
我和她.
身體上分開了.
但我常形容她的關係.
就是亦遠亦近.

$^{201}$有些事情是.
切割不到.
就算她在外國.
也切割不到.
我們每天都會通電話.
我每天都會記掛著她.
每天我都會為她祈禱.
原來.
神對人的愛.
那個憐憫.
就好像媽媽對子女一樣.
回想自己.
懷孕.
我不知道.
那個感受.
我記得我臨生前一晚.
因為我自己胎本前置.
不可以順產.
要開刀.
開刀的一晚.
我心裡也很害怕.
但心裡有一個.
壓制的興奮.
我明天見到我女兒.
很開心.
我見到她.
其實今天神對我的愛.
也是這樣.
壓制著那份興奮.
思念你.
記掛你.
很想和你有一個.
不能切割的聯繫.
有一位媽媽.
我介紹給你認識.
台灣的孩子.
叫李柏毅.
原來他.
在出生18個月.
被醫生判斷.

$^{241}$他是一個嚴重的自閉症的小孩.
他.
很難.
極難和人用語言溝通.
這個判斷.
其實對於.
這個小朋友的媽媽.
好像死刑一樣.
但因為.
媽媽.
柏毅的媽媽.
是一個基督徒.
在他成長的路上.
他見到不一樣的柏毅.
醫生說他是一個.
嚴重的自閉症小孩.
但他發現在四歲的時候.
兒子有一個.
畫畫的天份.
原來九歲.
他就鑑定為天才的畫家.
十二歲就開了他.
自己的個人畫展.
雖然柏毅.
有嚴重的用語言.
和人溝通.
但他的畫.
好像一句話.
能夠和人溝通.
他用的顏色很鮮艷.
其實他內心.
是有一份對人的熱情.
這個媽媽.
這樣說他.
自閉而不是精神病.
不是白癡,不是廢物.
需要鼓勵,需要痛惜.
需要愛.
他求上帝幫助他.
不要失望.

$^{281}$要看見不一樣的柏毅.
這個就是.
在一個無能為力的媽媽身上.
就好像嘉楠夫人.
她真的無能為力.
被鬼附著,她又趕不上鬼.
女兒又不曉好.
這個媽媽面對著一個自閉症的小孩.
他不是醫生.
就連醫生也無能為力.
但他靠著主的愛.
他看見不一樣的柏毅.
他能夠.
使這個生命.
有轉變.
因為他認定.
耶穌是主.
就好像嘉楠夫人.
主是大衛的子孫.
他知道這位就是他的主.
能夠幫助他.
在23,24節裡.
有個題目.
我加上去.
拒絕門外的媽媽.
認定耶穌是主.
如何拒絕門外?我們看看23節.
耶穌一言不發.
門徒進前來.
求她說.
這婦人在我門後頭喊叫.
請大家發她走.
這班門徒覺得這個女人.
很嘈.
叫她走.
這位媽媽.
有沒有經常被子女說.
媽媽不要問那麼多.
或者你問一句.
一道門就不想理你.

$^{321}$覺得你很煩.
其實當時門徒對這個女人.
就是這樣.
耶穌她很嘈叫她走.
我前陣子看了一套韓劇.
就說一個媽媽.
年輕媽媽.
BB出生不久.
患了一個病.
叫痴媽症.
BB痴媽症會怎樣?.
24小時都黏著.
睜眼要經常見到媽媽.
那一幕就是媽媽.
人有三急.
她要去大解.
BB又哭了.
但也沒辦法.
她就嘗試一下.
將BB放在廁所門外.
她就掩著門.
但BB哭得聲嘶力竭.
沒辦法.
於是她就要打開廁所門.
讓BB看到她.
也沒辦法.
因為她有距離.
她要再拉近BB在她身邊.
在這個尷尬的環境.
她先生下班回來.
她先生第一件事.
看到這件事有什麼反應呢?.
你有沒有衛生?.
你這麼噁心?.
你這麼骯髒?.
很多時候.
為人媽媽很多辛酸.
會覺得你怎麼教的?.
你怎麼不懂教?.
為什麼你.

$^{361}$把小朋友弄成這樣?.
會被人嫌棄.
我不知道你試過沒有.
在我自己的經歷裡.
其實.
也真的試過.
不過.
但不單止.
門徒嫌棄這個婦人.
在24,27節裡.
耶穌好像.
不太理會這個女人.
你看到23節說.
耶穌卻一言不答.
而到了.
24節.
耶穌回答她.
不是說幫她.
其實耶穌是拒絕她.
耶穌說.
其實我只是來.
以色列迷失的羊那裡.
於是耶穌就.
拒絕了她.
只是在羊那裡.
於是這個婦人.
又來拜她.
又說主啊幫幫我.
這個婦人.
其實她是恃而不捨的.
其實門徒都嫌棄她.
連耶穌都拒絕她.
但她仍然說.
主啊你幫幫我.
這就是媽媽的精神.
26節.
耶穌回答她.
孩子的餅掉級小是不妥的.
但婦人又繼續回答.
主啊不妥.

$^{401}$可恥少過也吃他主人出上吊下來的瑞士.
原來耶穌.
表面上看來.
好像拒絕這個婦人.
也沒有幫她.
而這個婦人也沒有因為耶穌的拒絕.
而收手.
或者收口.
她繼續兩次求耶穌.
求耶穌幫她.
耶穌沒有直接的即時回應.
不但如此.
還用一些羞辱她的說話.
去回應她.
她說.
稱這個婦人為.
小狗.
其實有些人.
當時猶太人看不起外邦人.
稱外邦人為狗.
這個狗就好一點.
小狗類似寵物狗.
但也不是尊重.
這個母親的意願.
但這個婦人.
她對耶穌說.
瑞渣兒的恩典已經夠了.
她認識.
這位主充滿能力.
一點瑞渣兒的恩典就夠了.
又想和你分享一個.
媽媽.
大家應該都很熟悉.
這個人叫.
力克.
他是一位戰士.
在見證裡.
我這次不是說力克.
而是說他的媽媽.
有生命的戰士.

$^{441}$當然有生命戰士的媽媽.
原來力克在澳洲出生.
他天生沒有四肢.
只有一隻小小的腳.
他有腳趾.
但不健全.
所以在醫學上.
稱為海豹.
其實當這個嬰兒出生.
力克的媽媽.
很震驚.
很絕望.
甚至拒絕.
擁抱這個嬰兒.
因為實在不知道如何面對.
力克的父母是虔誠的基督徒.
他們透過信仰.
透過祈禱.
去克服恐懼和擔憂.
盡量.
讓力克.
健康地成長.
雖然他沒有四肢.
但沒有例外.
他訓練他.
當他是一個正常的小孩.
在這個過程中.
當然有很多挫折.
因為媽媽很想.
力克讀正常的學校.
因為他沒有四肢.
但智力非常正常.
但當時澳洲的學校.
因為有法例規定.
一些雙健的小孩.
不可以讀正常學校.
但媽媽覺得.
他應該讀正常學校.
於是.
他申請了五次.

$^{481}$但五次都被拒絕.
於是他有很多挫敗.
很多拒絕.
之下.
他就積極地爭取.
修改這個法例.
在媽媽的.
不怕失敗.
繼續.
鍥而不捨的精神下.
他終於可以.
入讀.
力克.
首批可以入讀.
雙殘小孩.
正常學校的小孩.
你見到.
這位生命戰士的媽媽.
雖然.
他有一個不健全的兒子.
但他沒有例外.
他仍然當他是一個正常的小孩.
要去建立他的生命.
幫他去面對.
生命很多的挑戰.
因為媽媽見到.
他的兒子有不一樣的生命.
今天.
媽媽不能夠.
使力克.
生手生腳.
有些事情是改變不了.
是.
很沒有希望.
看下去.
但他仍然尊重.
他的兒子的生命.
信任這個生命.
雖然改變不了.
他生命有很多缺口.

$^{521}$很多不健全.
但他能夠轉化.
轉化他的生命.
成為很多人的祝福.
他將很缺陷的生命.
轉化成很有用的生命.
今天我們的主.
也是.
我們的生命有很多缺口.
有很多缺陷.
有很多不健全.
有很多軟弱.
但在神手裡.
神是生命戰士的媽媽.
也是生命戰士的主.
他能夠轉化我們的生命.
第三點.
到28節.
聖經形容.
這個媽媽很有信心.
這個信心很大的媽媽.
仍然認定耶穌是主.
當耶穌聽到.
這個婦人說.
瑞渣兒的恩典.
也夠用的時候.
耶穌回答她.
婦人你的信心很大.
只要你所要的.
成全你.
從那時起她的女兒就好了.
這裡其實在聖經裡.
在原文聖經裡.
有一個婦人.
之前有一個字.
有一個「哦」.
這個婦人.
什麼叫「哦」.
一個感嘆.
一個讚嘆.

$^{561}$就是說這個女兒厲害嗎.
厲害.
在場很多媽媽.
也覺得她很厲害.
每個人背後也有很多故事.
這個婦人如何厲害.
你要說她的信心很大.
聖經說這個很大.
英文是Mega.
通常在聖經裡.
出現了這個字.
這個很大的字.
出現了246次.
通常都是形容空間.
形容數量.
聲浪很大.
唯獨在這裡.
是用來形容這個婦人的信心.
很特別.
只有這個字.
是形容.
這個婦人的信心.
原來.
作為一個媽媽.
其實真的很厲害.
會有.
跟著愛子女的動力.
有.
是而不捨.
去幫助你的子女.
如果子女有什麼事會怎樣.
祈禱.
雖然當中.
會有很多的限制.
有很多的掙扎.
有很多的壓力.
甚至憂慮,害怕.
無能為力.
但是你看到.
這位迦南婦人的媽媽.

$^{601}$她是怎樣的.
她在很多的絕望裡.
很多的拒絕裡.
她認定主耶穌是主.
在整段聖經裡.
說了.
三次.
這個婦人呼求了三次.
主啊.
今天在我們的生命裡.
其實都有很多困難.
這些困難.
都好像.
這個迦南婦人一樣.
會有.
無能為力.
會覺得被人拒絕.
但是在很多困難裡.
這位迦南婦人她叫.
主啊,她認定耶穌是主.
今天我們有沒有認定.
耶穌是我們的主呢.
這位迦南婦人.
她認定耶穌是主.
她把她最寶貴的女兒.
把她自己的生命.
交上給主.
今天只有主.
才能夠承載我們的生命.
承載我們兒女的生命.
相信.
每一個信主的媽媽.
她都會.
有一個心願.
就是很想.
自己的生命讓主承載.
我兒女的生命.
能夠讓主承載.
因為知道有主的人生.
有主的媽媽.

$^{641}$有主的爸爸.
有主的弟兄姊妹.
是不一樣的.
我們人生不會一帆風順.
我們人生不會沒有壓力.
我們人生不會沒有難處.
但是我們人生有主.
在每個環境裡.
我們認定主就是.
我們的主人,交給他.
把我們一切的東西.
甚至最寶貴的兒女都交給主.
在主裡有最好的祝福.
我們有時間一起祈禱.
是我們人生的主.
在我們的家庭.
在我們的兒女.
在我們面對不同的環境裡.
雖然不會一帆風順.
甚至會有很多困難.
有很多的絕望.
甚至會被人拒絕.
但是主知道.
你就是我們的主.
你要陪我們去面對.
你要因著我們對你的信任.
你會祝福我們.
求你也使我們每一個人都經歷到.
你就是我們人生的主.
我們這樣祈禱.
是奉主名求.
阿們.
\newpage



\section{猶大書 1:1-4-17-25}
\label{sec:XkIka_Z9Y3w}
\textbf{2023.05.21 《聖道- 我將你的話藏在心裡》猶大書 1\_1-4,17-25 (和修版) 游志豪傳道}
\newline
\newline
連結: \href{https://youtube.com/watch?v=XkIka_Z9Y3w}{\texttt{https://youtube.com/watch?v=XkIka\_Z9Y3w}} ~~~~ 語音日期: 2023-05-21
\newline
\newline
\hyperref[sec:5NIuKMKgBWE]{\small{< < < PREV SERMON < < <}}
~
\hyperref[sec:index_chronic]{\small{[返順時目]}}
~
\hyperref[sec:index_scriptual]{\small{[返順卷目]}}
~
\hyperref[sec:33y9hhVAb_s]{\small{> > > NEXT SERMON > > >}}
\newline
\newline
猶大書 1:1-4-17-25
\newline
\begin{longtable}{cl}
\hline
\hline
章節 & 經文 (和合本修訂版)\\
\hline
1:1 & \begin{tabularx}{0.7\textwidth}{X} 耶穌基督的僕人、雅各的兄弟猶大，寫信給那些被召、在父神裡蒙愛、為耶穌基督保守的人。 \end{tabularx} \\ \\ \relax
1:2 & \begin{tabularx}{0.7\textwidth}{X} 願憐憫、平安、慈愛多多加給你們！ \end{tabularx} \\ \\ \relax
1:3 & \begin{tabularx}{0.7\textwidth}{X} 親愛的，我一直很迫切地想要寫信給你們，論到我們同享的救恩，但我覺得有必要現在就寫信勸你們，要為從前一次交付給聖徒的真道竭力奮鬥。 \end{tabularx} \\ \\ \relax
1:4 & \begin{tabularx}{0.7\textwidth}{X} 因為有些人偷偷地進來，就是早就被判定受懲罰的不虔誠的人，他們把我們神的恩典變為放縱情慾的機會，並且不認獨一的主宰—我們的主耶穌基督。 \end{tabularx} \\ \\ \relax
1:5 & \begin{tabularx}{0.7\textwidth}{X} 這一切的事，你們雖然知道，我卻仍要提醒你們：從前主只一次就救了他的百姓出埃及地，後來卻把那些不信的滅絕了。 \end{tabularx} \\ \\ \relax
1:6 & \begin{tabularx}{0.7\textwidth}{X} 至於那些不守本位、離開自己住處的天使，主用鎖鏈把他們永遠拘留在黑暗裡，等候大日子的審判。 \end{tabularx} \\ \\ \relax
1:7 & \begin{tabularx}{0.7\textwidth}{X} 同樣，所多瑪、蛾摩拉和周圍城鎮的人也跟著他們一樣犯淫亂，隨從逆性的情慾，以致遭受永不熄滅之火的懲罰，作為眾人的鑒戒。 \end{tabularx} \\ \\ \relax
1:8 & \begin{tabularx}{0.7\textwidth}{X} 照樣，這些做夢的人也污穢身體，輕慢掌權者，毀謗眾尊榮者。 \end{tabularx} \\ \\ \relax
1:9 & \begin{tabularx}{0.7\textwidth}{X} 天使長米迦勒為摩西的屍首與魔鬼爭辯的時候，尚且不敢用毀謗的話譴責他，只說：「主責備你吧！」 \end{tabularx} \\ \\ \relax
1:10 & \begin{tabularx}{0.7\textwidth}{X} 但這些人毀謗他們所不知道的。他們與那些沒有理性的牲畜一樣，只做本性所知道的事，敗壞了自己。 \end{tabularx} \\ \\ \relax
1:11 & \begin{tabularx}{0.7\textwidth}{X} 他們有禍了！因為他們走該隱的道路，又為財利往巴蘭的錯謬裡直奔，並在可拉的背叛中滅亡了。 \end{tabularx} \\ \\ \relax
1:12 & \begin{tabularx}{0.7\textwidth}{X} 這樣的人是你們愛筵上的污點；他們無所懼怕地同你們宴樂，彷彿牧人只顧餵飽自己。他們是無雨的浮雲，被風飄蕩；是秋天沒有果子的樹，死而又死，連根被拔出來； \end{tabularx} \\ \\ \relax
1:13 & \begin{tabularx}{0.7\textwidth}{X} 是海裡的狂浪，湧出自己可恥的沫子來；是流蕩的星，有漆黑的幽暗永遠為他們保留著。 \end{tabularx} \\ \\ \relax
1:14 & \begin{tabularx}{0.7\textwidth}{X} 亞當的七世孫以諾曾預言這些人說：「看哪，主帶著他的千萬聖者來臨， \end{tabularx} \\ \\ \relax
1:15 & \begin{tabularx}{0.7\textwidth}{X} 要審判眾人，證實一切不敬虔的人所妄行一切不敬虔的事，又證實不敬虔的罪人所說頂撞他的剛愎的話。」 \end{tabularx} \\ \\ \relax
1:16 & \begin{tabularx}{0.7\textwidth}{X} 這些人喜出怨言，責怪他人，隨從自己的情慾而行，口說誇大的話，為自己的利益諂媚人。 \end{tabularx} \\ \\ \relax
1:17 & \begin{tabularx}{0.7\textwidth}{X} 親愛的，至於你們，要記得我們主耶穌基督的使徒從前所說的話。 \end{tabularx} \\ \\ \relax
1:18 & \begin{tabularx}{0.7\textwidth}{X} 他們曾對你們說過，末世必有好嘲弄的人隨從自己不敬虔的私慾而行。 \end{tabularx} \\ \\ \relax
1:19 & \begin{tabularx}{0.7\textwidth}{X} 這就是那些好結黨分派、屬乎血氣、沒有聖靈的人。 \end{tabularx} \\ \\ \relax
1:20 & \begin{tabularx}{0.7\textwidth}{X} 親愛的，至於你們，要在至聖的真道上造就自己，藉著聖靈禱告， \end{tabularx} \\ \\ \relax
1:21 & \begin{tabularx}{0.7\textwidth}{X} 保守自己常在神的愛中，仰望我們主耶穌基督的憐憫，進入永生。 \end{tabularx} \\ \\ \relax
1:22 & \begin{tabularx}{0.7\textwidth}{X} 有些人心中猶疑，你們要憐憫他們； \end{tabularx} \\ \\ \relax
1:23 & \begin{tabularx}{0.7\textwidth}{X} 有些人你們要從火中搶出來，搭救他們；有些人你們要存懼怕的心憐憫他們，連那被情慾污染的衣服也要厭惡。 \end{tabularx} \\ \\ \relax
1:24 & \begin{tabularx}{0.7\textwidth}{X} 願那能保守你們不失腳，使你們無瑕無疵、歡歡喜喜站在他榮耀之前的、 \end{tabularx} \\ \\ \relax
1:25 & \begin{tabularx}{0.7\textwidth}{X} 我們的救主獨一的神，藉著我們的主耶穌基督，得享榮耀、威嚴、能力、權柄，從萬古以前，到現今，直到永永遠遠。阿們！ \end{tabularx} \\ \\
[1ex]
\hline
\hline
\end{longtable}
$^{1}$大英姐妹早安,很高興能夠在紅引堂再次分享港版神話語.
剛才的粵曲真是很好聽,而且我覺得很不容易練習.
內容也十分豐富,充滿神話語.
所以和今天講的題目也很貼切,就是聖道.
我給了他一個副題目,就是將你的話語藏在心裡.
究竟聖道又是什麼呢?我們就從一些引言開始.
我從網上找到很多學校名稱,就找了聖道.
但也找到一位古文章,他曾經提及聖道這兩個字.
他指出古文章外篇天道三節的古文中出現了聖道運而無所織,故海內服.
全部都是簡明的字,但拼在一起就不知何解.
其實他說這一段的說話是在說什麼呢?.
他說第一段是在說自然的法則是不會停留的.
所以萬物就能得以生成.
第二段是說帝王的法律也不會停滯,所以天下也能歸順.
最後他說聖人對於宇宙萬物的看法也不會停留和中斷.
所以四海之內人心也能夠去折服.
解釋就是這樣.
我們在想聖道是否在說人對宇宙萬物的看法而已.
就像這篇章裡說的,這就是聖道呢?.
究竟聖經裡有沒有提及過聖道這兩個字呢?.
又找找又找,又做一些研究.
原來在加拿大書第六章第六節,是新譯本,和合本是沒有的.
新譯本是有這樣的譯.
在聖道上受教,應該和施教的人分享自己一切的美物.
和合本修訂版就譯成「真道」,就是神的話語.
原來這裡有出現過「聖道」這個字.
可想而知,聖道和真道,和神的話語,其實有一個很大的關連.
我們繼續看看,究竟「真道」在聖經裡有沒有出現過,有什麼不同的地方呢?.
其實又發現了一些很特別的聖經裡,有提及過「真道」這件事.
這個字很小,不過我簡單的說.
原來在提多書裡,出現過三次都提及過「真道」.
不過是不同的譯本,第九節提及過要合乎教義,可靠的「真道」.
你們要用真實的道理去教導人.
所以這個字又是在說真實的道理,即是神的話的意思.
第十三節又提及到,要嚴格責備那些人,使他們在「真道」上純全無瑕疵.
在他們的信仰上是無瑕疵的,是很完善的.
最後第十四節又提及到,你們不可以聽猶太人的荒謬言語,離棄「真道」之人的誡命.
這個就是和合本說的,是在說真理的意思,即是真實的事情.
所以我們看到,就算「真道」這個字的譯法,其實都有不同的意思.
而「真道」,「聖道」裡面,其實都有不同層次的解法.

$^{41}$簡單來說,我畫了兩個圈圈.
第一個圈圈,左邊的圈圈是最窄的,是在說什麼呢?.
其實在科音書也說,「太初有道,道與神同在,道就似神」.
說到最簡單的道,耶穌基督自己就是那個道.
耶穌基督的福音,他的犧牲,他的救贖,就是那個道理的裡面.
但是最大的圈圈,就是在說「聖道」,耶穌基督的教訓,聖徒的教訓.
那個「真道」,聖經的真理的裡面,就是我們今天要看的神的話語的裡面.
所以今天大英姐妹,我們有沒有把這個「道」,「聖道」,或者這個「真道」.
不是聽了就算了,而是能不能夠放在我們心裡面.
好像剛才那首詩歌裡面,將神的話語能夠放在我們心裡面.
這樣才叫做「圓滿」,而不是聽了,水哥阿輩那樣,就叫做我們知道那個道是怎麼樣了.
我們今天看了經文一次,經文是記載在尤大書的幾節的經文.
我挑選了頭和尾的部分.
不知道大家有沒有留意這本書,尤大書裡面最重要的內容,或者他在說什麼呢.
當然第一節他說的是尤大,那尤大是誰呢?原來他自己說他是耶穌基督的僕人.
這裡當我們讀第一節的時候說,耶穌基督的僕人的裡面.
但是如果我們從另外一些的說法告訴我們,尤大和耶穌是兄弟.
有些文獻就說他是同母異父的兄弟的裡面.
但是他在書本裡面又沒有說我是耶穌的兄弟,你們要聽我的話.
他沒有這樣說的,不過他就說他是雅各,也是耶穌的兄弟.
雅各的兄弟,有一點點關係,不過他更加說他自己是僕人.
耶穌的僕人,忠於主耶穌基督,現在將他的話語跟大家去說.
他很謙卑,他說是耶穌的僕人.
那究竟尤大書寫給哪些人呢?.
尤大書所寫給的,他說是在和合本修訂版,剛才我們讀了的.
他說那些被召,在上帝裡面蒙愛,神的愛的,和耶穌基督保守的人.
就是說是上帝愛的,愛你們的,在耶穌基督的眷顧,在他的平安的裡面.
也蒙神去呼召的那幫人,寫給那些人的裡面去.
弟兄姊妹,尤大提醒我們第一件事是很重要的.
雖然我們每一個人都不同,可能大家的背景都不同.
認識神的次序,年份都不同.
但是每個人都是神所愛的,所保守的,所呼召的人.
呼召去做什麼呢?呼召就是遵行神的話,和遵行神的旨意的人.
我們每個人都是,在在你們每一位其實都是神所愛的,神所保守的.
神所呼召你們去走在神的道裡面的人.
所以在整卷書裡面,尤大書有一個很重要的目的.
如果我們剛才一直看的時候,我們就會發覺他講了一個很重要的信息給我們聽.
他就說是寫給一些人去鼓勵他們.
因為有很多假的道理,有很多似是而非的道理.
假的師父,假的教師在教會裡面,混入在教會裡面.

$^{81}$所以他提醒那些弟兄姊妹,提醒我們要怎麼樣呢?.
要堅持堅守純正的基督的道理的教訓,真實的信仰的裡面.
所以尤大他提醒,也提醒我們要守住聖經的道理當中的裡面.
這個是最重要的一句書卷,告訴我們.
如果在和合本修訂版,他就會這樣去講.
我大約折錄了一些分段出來,這個也是大家的一個參考.
但是今天我們從另一個角度去看這本書.
從一個正道,真道的角度去看這本書.
究竟他和我們在說什麼,或者對我們有什麼關係.
他說了很多東西都是當時的情況.
如果對我們來說,今時今日在在每一位究竟有什麼提醒的裡面呢?.
我會將他這樣分成不同的節數.
這個就簡略地擺出來,一會兒我們再看的時候.
就會再慢慢去看這本書究竟提醒我們什麼的當中.
他提到真道,誰是守護真道呢?.
他提到真道,聖道究竟是怎麼樣呢?.
他提到我們要怎樣去持守聖經的真道.
可以去抵抗很多不同的入侵呢?.
在這裡這樣說,第三節他說.
猶太第一提醒我們究竟今天哪一個才是真道的守護者呢?.
在第三節,我看回這裡,在第一節他這樣說.
今天我們常常以為守護著這個真道的是哪一個呢?.
是聖經的學者,他們有聰明有智慧.
神學院的教授,或者是教務同工,他們讀了神學.
所以他們就是守護這個聖經,守護這個真道的人.
要有一定的神學的水平,身份或者學識.
甚至是專業才有這個角色去維護這個真道.
不過猶太告訴我們,他這本書不只是寫給教務同工的.
他剛才已經提過,寫給我們每一個信徒.
在神的國裡面,有神的保守,神的呼召.
這些就是真道的守護者.
所以今天猶太提醒我們第一件事.
我們每一位信徒,你們信了主.
其實都有責任去守護著耶穌基督的真道.
我們大家都是這個真道的守護者.
弟兄姊妹,在這個年代,在這個年頭裡面.
你有沒有想過,其實你要做好準備.
你要準備好自己成為真道的守護者.
為主的道有時需要去爭辯,有時需要去告訴別人.
誰是真,誰是假的.

$^{121}$我們每一位都要有這樣的特性,有這樣的知道.
我每一個月都會講一些例子在裡面,我自己的故事.
我每一個月都會有幾次去到泰國人的教會.
可能大家都熟悉,九龍城宣道會有一間泰國人的教會.
我們會去到那間教會,有時會去附近的公園.
跟一些泰國的朋友講福音,派一些單張給他們.
有時會去到泰國的餐廳,我們不懂得說.
就這樣給他們一些單張,告訴他們有泰國的教會.
我有幾次都在星期日裡面,一早上.
我回教會,在公園看到有一個攤檔.
那個攤檔,我看清楚一點.
我不知道在洪水橋,還是元朗附近多不多.
原來就是夜禍見證人,所謂的摩門教徒.
他們擺了一個攤檔,有兩個人站在旁邊.
放了一些書,寫著一大張字.
宣告關於耶和華的真理.
我就說,他們在說耶和華的真理嗎?.
我們在說什麼呢?.
當然我們知道,摩門教徒站在那裡是不對的.
他們三位一體的教義,都有些錯漏在裡面.
他們所謂的聖經,都有些錯漏在裡面.
不過他們說到耶和華的真理在裡面.
有很多的誤導在裡面.
有一次,當我們去派單張的時候.
我們又看到這些人.
我就和一些泰國人,他們都會說廣東話的.
泰國的弟兄姊妹去說.
我們要學會怎麼去分.
他們寫著耶和華的真理.
但你知不知道他們在說錯誤的道.
或者是假的東西呢?.
我和他們說,我們要懂得去分辨哪些是真,哪些是假.
雖然只是一間普通的泰國人的教會.
可能只有幾層,他們可能洗碗,做一些家用.
但是我們都要懂得去分辨哪些是真,哪些是假.
因為其他人不懂得分辨哪些是真.
所以他們都要懂得去維護究竟基督教的真理.
耶穌基督是誰,他們要懂得去說.
其實二端的奇說,很多不同的道理,經常圍繞著我們.
我們要準備好,我們都要成為真理的守護者.

$^{161}$我盼望第一,大家都要記得自己是真理的守護者.
你都要告訴別人什麼是對的.
至少你可以帶他們回到紅恩堂裡面.
這個就是對的,這裡就是我們常常去的教會.
你們要認識,介紹給牧者知道.
這些就是我們所信的裡面.
第二,猶太人提醒我們,什麼是真道.
我先按一下.
剛剛提及第三節,親愛的,我一直迫切寫信給你們.
論到我們同享的救恩.
但我覺得有必要現在寫信勸你們.
要為從前一次交付及聖徒的真道竭力奮鬥.
你們知不知道猶太人一開始就說.
大家都是同享在救恩裡面.
不論是保羅傳給你,還是其他使徒傳給你.
大家都是同享救恩的裡面.
我們領受了這個真道.
我們就在基督的愛和拯救的裡面.
不過真道又是怎樣,或者聖道又是怎樣呢.
我們要留意的是從前一次交付給聖徒的真實道理.
聖徒就是指的是神的子民.
所以這裡有一個真道是指的是.
神一次而且永遠的托付給神的子民.
那些就是使徒的教訓,或者傳統的裡面.
換句話說,猶太人說了一句很簡單的東西.
告訴我們耶穌到城肉身來到這個世界裡面.
他有很多的事蹟,他跟別人說了很多話.
初期教會裡面,保羅和不同的使徒跟著他們.
他們就是耶穌基督的見證人.
他們將耶穌基督的事蹟和教訓記下來.
一次傳遞出去給我們當中.
所以其實最簡單,最簡單就是告訴我們什麼.
我們今天手上的聖經,你看的聖經.
流傳下來,特別是新約聖經裡面.
這些就是使徒的教訓.
也是耶穌基督的教訓,他的故事裡面.
這些就是最重要的信息當中.
今天我們所信的一切都是基於這本聖經.
基於我們所看到,所聽到的聖經內容裡面.
聖經裡面是神的默示,是神的說話.

$^{201}$我們今天如何去持守,看這本聖經呢?.
我們家裡可能有很多本聖經,我不知道大家家裡有多少本.
你回去數數,不同的譯本,可能我都有十本八本在裡面.
有些可能你受浸信主,又送你一本.
不同的時間又會送你一本.
我們就放了很多的聖經在書櫃裡面.
有一個這樣的故事,這個故事不是基督教的.
是說回教徒的故事.
他說在中亞洲有一個村落.
有人在河裡發現了用塑膠包好.
有一本阿拉伯語的可蘭經,就是他們的經書.
這本經書就浮浮浮,被人罵了上來.
有個婦女發現了,拉開之後,那本經書很漂亮.
罵起來的時候,看起來完好無缺.
他們就說原來那些回教徒不想那些書被毀壞.
他的經書被毀壞,就放在家裡珍藏.
當有人去搜查,可能當時有些警察會搜查,很危險.
他們就會拿出來用塑膠包好,封好,然後就扔進河裡.
那本書就飄來飄去.
故事說他們從新疆不知道什麼河,一直飄來飄去.
哈薩克斯坦,就是中亞洲的地方,可能飄了很久.
希望有人能夠將那些經書撈起來.
完好無缺,可以用.
對於回教徒來說,他們的書不單只是一本經書.
而是神聖的標記.
當我在網上看到這個故事的時候.
其實對自己也有些反思.
我們拿一本聖經,Bible,究竟是什麼呢?.
今天我們要怎麼看聖經呢?.
是一本書?是一本經書?是一本智慧的書?.
或者是一本人生道理?究竟我們要怎麼看這本書呢?.
或者你只是覺得它是放在書架的其中一本呢?.
但是主的道理是我們需要去看的.
是我們需要去閱讀的,是我們需要去尊重的神話語.
剛才我也提及過,我們可能有不同的譯本.
不同的文字在聖經裡面.
但有時候我們只是放在架子裡,不去翻閱.
究竟今天聖經這麼重要的時候.
我們要怎麼去看神的話語呢?.
第三,猶太指出.

$^{241}$那個的真道,聖道當時的處境又是怎麼樣呢?.
告訴我們為什麼要辭掉這個真實的道理.
第四節,他說因為有人偷偷進入了教會裡面.
就是詛咒被判定受刑不填成的人.
他們將我們上帝的恩典變為放縱情慾的機會.
並且不認識獨一的主宰,我們的主耶穌基督.
另一個版本,我讀一次給大家聽.
是新普及譯本,他說因為有一些不敬虔的人.
已經偷偷混入你們教會裡面.
揚言要說,有了上帝奇妙的恩典.
我們就可以放蕩地生活.
他們否認我們獨一的主宰,耶穌基督.
這種人的詛咒已經被判定在案了.
當我再讀一次的時候,不是太難.
很容易明白這句聖經的說法.
他說只是說當時有人偷偷混入了教會裡面.
不知道今天有沒有這樣的情況.
偷偷混入了教會裡面.
他們是不敬虔的人,可能很放蕩.
重要的是把上帝的恩典,神有恩典給我們.
我們就可以亂犯罪,隨意做任何事情.
最重要的是他揚言否定了耶穌基督.
否定了耶穌基督,這是很嚴重的問題.
而這些人的罪早已經被判刑.
早已經有刑罰降臨在他們的身上.
當時真道的處境是怎樣呢?.
真道的維護和延續受到很大的考驗.
因此尤大在當中很迫切和緊張.
跟每一位弟兄姊妹說.
你們要好好守住維護這個聖經的真理.
這個真道的裡面去.
當我看到的時候,對我亦有很大的提醒.
對大家的生命也有很大的提醒.
今天我們的教會,我不知道,剛才是偷偷走進來.
不知道大家有沒有認識,有誰是偷偷走進教會呢?.
你看看旁邊那個,我不認識你,你是否偷偷走進來的?.
你走進來做什麼呢?不是,大家認識的.
不過進來之後可能都是新朋友.
不過當時就有一些人困在教會裡面.
今天有沒有呢?.

$^{281}$都有這樣的情況當中.
有些是擺明的,那些叫做明刀明槍.
說不同的東西,就是那些異端的邪說.
好像前陣子說了很多報道,韓國的那個攝理教.
不知道大家有沒有看,完全是那些異端的來到我們當中.
但是有時候有一些是似是而非.
他又好像對,又好像不對.
搞到你很難去分辨的一些道理.
又會偷偷地混在教會裡面.
所以我們要很精明,很清楚究竟神的話語說的是什麼.
以致我們能夠分得到哪些是對,哪些是錯.
當然你回到教會裡面聽到牧師,聽到傳道.
又有上主一學,或者很多的教導.
我們都更加清楚究竟真的是什麼.
以致我們可以分辨到假的,似是而非的道理是什麼.
隨手在網上都找到一些資料.
他說這個在肯亞的一個教會,叫做佳音國際教會.
他的名字也很基督教,佳音國際教會.
聽起來好像沒有什麼錯誤.
但是他說和那些信徒洗腦.
他說如果你要見到耶穌,那是怎樣.
餓死你就可以見到耶穌.
他沒有關閉教會,他和會眾說.
你要做一個方法,就是要捱餓.
你要捱餓,捱餓就可以很快上天堂.
上天堂的時候就和耶穌相會.
那想不想和耶穌相會.
想,但是難道真的要捱餓嗎.
好像我們沒有這樣去說.
不過最嚴重的是當他們這樣去做的時候.
肯亞的政府當然去查他們.
後來就發現在森林裡面有很多墳墓.
他就發現很多家長餓死了小朋友.
以致就撞了在墳墓裡面.
震驚了整個地方.
發現了這個邪教的地方.
很多時候都是歪曲了聖經.
亂說一些話.
令到人犯法或者走上很多不同的路.
令到人控制了.

$^{321}$這些就是一些名刀名槍.
我們知道的時候,我們就不會聽.
也不會相信一些錯誤的道理.
但是有另外一側在香港.
不知道大家對香港的消息有沒有更加了解.
在前一兩個星期就出了新聞.
就是這個.
我想大家有看新聞可能都會知道.
因為這個新聞寫得很白.
他說在香港裡面有一對信奉.
他就這樣說,照讀.
是一對基督教靈恩派的夫婦.
他們帶著一歲的女兒去探望外婆之後.
失去了聯絡.
當時有人找他的時候.
就發現了.
消防員找到他.
就發現在村屋裡面.
他的女兒,女嬰兩歲大.
就已經去世了.
太太就抓了他們.
抓了他們就說.
他們覺得要驅魔.
就和兩夫婦在家裡做驅魔儀式.
後來也發現他們兩歲的女兒.
好像被鬼附了.
又被魔附了.
就和女兒驅魔.
驅了十一個小時的驅魔儀式.
最後他們的女兒就去世了.
不知道是不是因為這件驅魔的事件影響.
不過後來就被抓去坐牢.
前陣子在法庭裡判刑.
他們寫得很明顯.
某些報紙寫得很明顯.
就說他們是基督教靈恩派.
不過不是香港人.
好像是馬來西亞人.
他們就說他們是在香港大學的大學生畢業.
畢業後去了加拿大的維珍神學院讀完神學.

$^{361}$女孩就完成了神學.
男孩就因為抑鬱症就完成不了.
回到香港後就發生了這件事.
然後我就和其他人分享.
他們大學生又讀完.
但也會產生這類事件.
當然內裡不知道究竟是怎樣.
不過有時有些似是而非.
或者我們很容易混亂的事.
都會混入在教會當中.
所以猶太經常提醒我們.
我們要警醒.
我們要知道什麼是真什麼是假.
有時候有些很奢侈的事.
我們都要很清楚.
可能我們問一下傳道.
問一下牧者.
讓我們更加知道究竟真實的道理.
在裡面是什麼.
其實假道理很多時候都有.
但千古以來都有.
混入教會.
好像那時候的東方閃電.
到處都是.
又或者有很多同性戀的色經.
都有很多錯誤.
都好像都會混入在教會裡面.
我們聽了有時又會似是而非.
不過最重要我們要懂得去分辨.
猶太就是說.
我們要懂得去分辨.
不懂的我們就去問一下牧者.
怎樣去分辨什麼是真什麼是假.
以致我們能夠去守住.
這個真理和耶穌基督的教導裡面.
最後猶太第四樣提醒我們.
他就最嚴肅.
開宗明義就提醒我們怎樣.
要守住這個真理.
第三節他說.

$^{401}$要為從前一次交付聖徒的真道.
要怎樣.
竭力奮鬥.
竭力奮鬥是什麼呢.
就是竭力維護.
用力的.
不只去爭辯的.
並且是要付上代價.
有時要付上代價.
不是很輕鬆.
我爭辯.
我維護.
他真的要付上代價.
你要去明白.
你要去維護.
這個真實的道理裡面.
要付上代價.
今天我們有沒有維護.
竭力維護.
或者去持守.
我們自己有沒有能夠去持守.
這個真實的道理.
這個聖道呢.
說一個容易一點的情況.
我們在工作上.
我們的生活上.
我們有沒有能夠為主.
為真實的道理去作見證呢.
這個也是維護真實的道理.
我們可以想想.
我們在工作上.
我們在生活上.
怎樣能夠見證神的真道.
其實這個就是見證.
維護真實的道理.
我講一個宣教.
因為我自己也是宣教.
宣教士有一個故事.
他去到一個.
我忘記了是否澳門的地方.

$^{441}$或者是哪個地方.
他就說去吃飯.
宣教士和不同的人去吃飯的時候.
不是弟兄姊妹.
未信的.
他們吃飯的時候.
宣教士通常開始就祈禱.
我們就祈禱.
借飯祈禱.
其他人就未信的.
當他祈禱的時候.
這個故事就說.
當他一祈禱.
那些人就開始耍他.
他有一碟菜.
就收起那碟菜.
就拿走那碟菜.
就換掉那碟菜.
就夾他幾塊.
他祈禱完之後.
就打開眼睛.
咦 為什麼不同了.
為什麼雞就變成叉燒.
當然是和他玩.
因為他們不明白.
這個宣教士就會這樣想.
究竟我繼續不繼續.
去借飯祈禱.
我不知道大家有沒有借飯祈禱.
或者當是這個處境的時候.
你會不會繼續借飯祈禱.
你可能會說不要緊.
沒事的.
或者先吃完.
飯後祈禱.
但是這個宣教士就說.
他要繼續.
第二次又是被人拿走那些東西.
又是被人換走那些東西.
那他繼續吧.

$^{481}$第三次第四次.
都是繼續去做.
直到我都不記得哪一次.
直到不知道哪一次.
他那次張開眼睛.
沒事的.
原封不動.
那碟飯在那裡.
那杯飲料在那裡.
沒有加料.
沒有弄.
大家很安靜地吃.
然後他就很奇怪.
他就問.
為什麼這次沒有人換走我的菜.
當然那些人說.
都不要玩了.
玩了那麼多次.
不過他說.
那些朋友說.
我見了你那麼多次.
這樣去做.
你都繼續堅持.
他就會覺得.
你是認真的.
你的信仰是認真去做.
然後他們就說.
那我不玩了.
因為他們知道.
這是一個認真的.
對於他信仰認真.
所以那些人自此之後.
就沒有繼續去玩弄他.
我想這個很簡單.
其實對我都有一些提醒.
我們怎樣去持守.
持定很簡單的道理.
或者聖經裡面的一些提醒.
有時候我們會忽略.
不過這些很簡單的見證.

$^{521}$很簡單的持守.
這些真道.
其實就是很有力的.
持定真道的一個做法.
有時候我們未必做很多很大的判別.
但是我們生活裡面.
很需要去見證神的真道.
在猶大書裡面.
他也提到有三樣東西.
我們可以去做.
是什麼呢.
他說我們要記得.
主耶穌的使者.
從前所說的一句話.
他的提醒.
首先我們一記住.
首先剛才說.
將神的道放在我們心裡.
我們可以做到的.
記住神的說話.
第二第十節.
他說在真道上造就自己.
其實造就的意思是什麼呢.
其實我們要長進.
要讀聖經.
要學習認識神的話語.
我想這些常常都會提到.
不過猶大也是這樣提醒弟兄姊妹.
我們是需要這樣做.
記住讀聖經.
最後有些人有疑惑.
需要去練文.
需要去駁斥.
需要去忠告.
去挽回.
一些走差了的弟兄姊妹.
這些都是我們的責任.
提醒我們要記住.
我們要建立自己.
要好好讀聖經.

$^{561}$要好好去幫助身邊的人.
他們在真道上能夠去建立.
這是我們能夠.
也可以做的事情當中.
最後.
在猶大書最後那段經文.
他說願那能夠保守.
你們不失腳.
使你們無瑕無痴.
歡歡喜喜站在他榮耀之前.
我們救主獨一的上帝.
藉著我們主耶穌基督.
得享榮耀威嚴能力權柄.
從萬古以前到現今.
直到永永遠遠.
猶大提到最後誰保守呢.
其實我們遇到的衝擊可以很大.
在我們生活上遇到的困難可以很多.
但唯有上帝才是最終的保守.
願神能夠保守我們不失腳.
我們讀了這麼多聖經.
但我們求聖靈在我們當中.
讓我們明白.
讓我們認知.
讓我們能夠去分別.
讓我們能夠站在真實的道理裡.
為主去做更大的工作.
去見證上主.
所以一切都是神的憐憫和保守.
讓祂的保守保守我們每一位.
今天的衝擊很大.
有時候我們不知道很多假道理.
都會在我們身邊.
很多傳媒很多假道理.
不過讓神提醒我們.
當我們讀聖經的時候.
就告訴我們這是真的.
我們放在心裡.
切切守住.
維護.

$^{601}$也將真實的道理傳給身邊的人.
最後剛才那首詩歌.
我很想送給大家.
當中最後一個結語就是.
我將你的話語.
神的話語.
聖經的話語藏在心裡.
免得得罪神.
免得我遠離.
他說與我親近.
我愛神的聲音.
祂就是在我腳前的燈路上的光.
雖然天地都廢去.
但是神的話語.
聖經說話.
神的默示是不會過去的.
我們一起禱告.
主我們天上的父親.
你賜予揚佑我們.
我們由聖經.
甚至是中文的聖經.
讓我們能夠去看.
能夠去明白.
我們很感恩.
今日求主讓我們.
能夠多讀你的說話.
以致我們明白.
這個真實的道理.
這個是自性的道理.
將它藏在心裡.
雖然會遇到很多衝擊.
或者很多似是非的道理.
但求主你都保守教會.
保守紅印堂.
保守每一位弟兄姊妹.
他們能夠憑著信心.
能夠持定真道.
直到見主面的日子.
我們感恩.
奉主耶穌基督的名字.

$^{641}$阿們.
\newpage



\section{羅馬書 15:22}
\label{sec:33y9hhVAb_s}
\textbf{2023.05.28 《五餅二魚的團契》羅馬書 15\_22~29 (和修版) 陳麗蟬牧師}
\newline
\newline
連結: \href{https://youtube.com/watch?v=33y9hhVAb-s}{\texttt{https://youtube.com/watch?v=33y9hhVAb-s}} ~~~~ 語音日期: 2023-06-01
\newline
\newline
\hyperref[sec:XkIka_Z9Y3w]{\small{< < < PREV SERMON < < <}}
~
\hyperref[sec:index_chronic]{\small{[返順時目]}}
~
\hyperref[sec:index_scriptual]{\small{[返順卷目]}}
~
\hyperref[sec:MMKQi99_fYg]{\small{> > > NEXT SERMON > > >}}
\newline
\newline
羅馬書 15:22
\newline
\begin{longtable}{cl}
\hline
\hline
章節 & 經文 (和合本修訂版)\\
\hline
15:22 & \begin{tabularx}{0.7\textwidth}{X} 因此我多次被攔阻，不能到你們那裡去。 \end{tabularx} \\ \\
[1ex]
\hline
\hline
\end{longtable}
$^{1}$主席.
丁小姐早安.
很感恩可以來到洪恩堂跟大家一起.
用上帝的說話彼此勉勵.
今天是什麼日子?.
今天是聖靈降臨節.
也就是說耶穌基督復活之後的五十天.
就是聖靈降臨.
在二千多年前.
聖靈就像火,像風一樣.
臨到地上.
賜給信徒一個新的生命.
也賜給信徒無比的能力.
在《利未記》十二十六章是這樣說.
「必按時降雨給你們.
使地長出農作物.
田野的樹都能結出果子」.
我們知道乾旱的田地.
需要有雨水滋潤的時候.
才能結出甜美的果子.
人的心靈也一樣.
我們的心靈需要由聖靈在我們當中去囂觀.
以致我們的生命可以結出美妙的,美善的,生命的果子.
在《耶福所書》又這樣說.
「從前你們是暗昧的」.
就是說還沒信的時候.
「但現在在主的裡面是光明的」.
信了耶穌之後我們就是光明之子.
光明的子要結出良善,公義和誠實.
所以我們要經歷聖靈的風吹入我們的生命的時候.
我們就有更生.
我們就可以靠著聖靈的幫助.
我們有一個良善,公義,誠實的生命.
但是要這樣變化,變更.
除了藉著聖靈的充滿之外.
也要藉著與弟兄姊妹的相交.
我想請問Paul Paul.
一開始就去了這一頁.
我去第一頁.
對了,這裡.

$^{41}$大約在五年之前.
我就去了印尼.
去了我印尼姐姐的家鄉那裡.
去認識一下她的家人和教會.
他們印尼是這樣的.
他們的家鄉有很多的基督徒.
整條村子我想有九成多都是基督徒.
因為唯有這樣集合在一起的時候.
才不會被周圍的回教的人去攻擊.
有一天我就去她的田地那裡.
去到的時候我就聽到有人在那裡嘩嘩聲.
又看到有個人在那裡.
然後就跑去前面.
原來他們很開心.
因為那裡都沒有什麼人在外地來去.
哇,有個從香港.
都不知道香港是什麼地方.
從香港而來的人.
還要去到他們的田地.
他整天覺得他們的田地很骯髒.
也沒有椅子坐的.
所以這個很高的弟兄.
就徒手爬上一棵椰子樹那裡.
摘了幾個椰青.
他沒有手拿下來.
於是就扔了下去.
就特意扔到那些灌水渠道那裡.
然後這個矮一點的弟兄.
他就跑去那些灌水渠道.
追追追追回那幾個椰青.
那次就是第一次.
喝這個新鮮在椰子樹上摘下來的椰青.
很清甜,很好吃的.
他們的家庭很窮困.
他們的身體各自有不同的缺陷.
剛才高的弟兄是單眼的.
矮的弟兄是有癌症的.
在兩年多之前他都回到天父家了.
他們很開心有外面的人來.
他們覺得是一個榮幸,榮耀.

$^{81}$所以盡他們的所能.
就和我們分享他們覺得最好的東西.
最好的東西就是在樹上摘的椰青.
新鮮摘下來給我喝.
是一個分享,和我分享.
我們五月整個月都在講的主題就是團契.
他們同工分別一直都去.
他不太聽我的話.
對了,團契這個字.
在他寫信經過希臘文的原文裡面.
意思就是分享.
就是我剛才講的.
兩位弟兄和我分享他們覺得最好的東西.
就是這個分享這個字.
我們又看一下羅馬書的背景.
是不是向前按的呢?.
是的,特洪楷.
先分享一下羅馬書的背景.
羅馬書是在哥林多城寫成的.
當時保羅都沒有去過羅馬.
我想他死之前都沒有去過.
但從羅馬書裡面的第一章或者最後那章.
他問候了很多的人.
有些是他的親屬,有些是他的朋友.
但為什麼他和他們這麼熟悉呢?.
可能是之前保羅去過不同的地方.
去孫教,去馬其頓,阿戈亞.
去不同的地方孫教認識.
之後這班人就移民遷居到羅馬.
所以保羅寫這個羅馬書的時候.
很多人好像很熟悉.
保羅在這本書裡面表達.
剛才讀的那段經文也是.
表達了很想去到羅馬教會.
就和那班信徒在靈裡面.
一起彼此的分擔,分享.
我們就看一下.
這個團契的基礎是什麼呢?.
團契的基礎最重要的就是.
有耶穌基督在這裡.

$^{121}$在哥林多前書第一章第九節.
他說:上帝是信實的.
祂呼召你們,即是上帝呼召了我們.
與他的兒子,就是主耶穌基督.
共享團契.
上帝特別選擇了我們成為了他的子民.
目的是怎樣呢?.
與主耶穌基督連結在一起.
有一個親密,緊密的關係.
請問我按得對不對?.
現在可以了嗎?.
向這邊,不是向那邊.
OK,多謝.
第二段經文我們來看一下.
在約翰福音第十四章,耶穌說的.
到那一天,你們就會知道.
我在父裡面,你們在我裡面.
我也在你們裡面.
哇,很多個裡面,搞不清楚是怎麼回事.
我們看第一個裡面.
耶穌說,我在父裡面.
耶穌說,祂在父上帝的裡面.
那是為什麼呢?.
就是說主耶穌基督與父上帝有一個很親密的關係.
主耶穌基督在地上所做的一切.
都是按著上帝的旨意去做的.
所以當他在赫西瑪利園.
臨上十字架之前的祈禱也是這樣.
他說,如果你可以拿走我的苦杯.
不用我上十字架就好了.
不過,接著下一句就說.
不過,要按著你的旨意而行.
所以耶穌基督與父上帝的關係很緊密.
他聽上帝的話.
第二個裡面.
你們在我裡面.
你們也就是說一群信徒.
就在耶穌基督的裡面.
同樣的意思.
我們在耶穌基督裡面的時候.

$^{161}$我們與祂有很好的關係.
所以我們每天要看聖經.
要祈禱.
並且最重要的是.
我們做所有的事情.
我們的行事為人.
我們的價值觀.
都是以耶穌基督來成為準則.
這個就是我們在耶穌基督的裡面.
接著又倒過來.
我也在你們裡面.
就是說,如果當我們整個人生.
行事為人,價值觀.
都以耶穌基督為準則.
我們要緊緊地去禱告.
讀經的時候.
祂就會在我們裡面.
這個意思就是說.
在我們生命的當中.
祂看到我們沒有力氣的時候.
祂就會給力量.
祂看到我們軟弱的時候.
祂就會扶起我們.
祂看到我們危險的時候.
祂就會保護我們.
所以只要我們是.
走在耶穌基督裡面的時候.
祂都會在我們裡面.
行不行?.
這個就是一個團契的基礎.
與主耶穌基督有一個很密切的關係.
基督教是一個群體的信仰.
也是一個實踐愛心的地方.
所以我們回到這裡.
我們有團契.
我們整教會的人就是團契.
我們怎樣去實踐這個愛呢?.
去拜其他的偶像.
特別是去拜黃大仙.
我們知道.

$^{201}$可能在未信之前.
每年連三天晚上就去上頭炷香.
你可以不認識其他人.
上完炷香就回家睡覺就行了.
甚至你為了爭著去上頭炷香.
你看電視新聞.
每個人都爭著去上.
你想人家推倒你.
你可以鬆手撞人.
又可以踩人家的腳.
讓人跌倒上不了香.
你做什麼都行.
這個就是他們拜偶像的.
完全可以跟其他人沒有關係.
甚至可以彼此為仇.
但在耶穌基督的地方裡面.
在團契裡面.
這個是實踐愛心的地方.
我們看第二方面.
就是分擔和分享.
剛才我講了羅馬書的十五章.
大約保羅去了馬其頓和阿蓋亞宣教.
做的事情都差不多.
你看到這幅地圖.
中間的藍色大四方卡.
就是馬其頓和阿蓋亞.
在我的右手邊.
下面有一個紅色的橢圓形.
就是耶路撒冷.
保羅當時想著.
他在馬其頓和阿蓋亞宣教工作.
對這裡的人很好.
就收集了一些捐款.
送給耶路撒冷教會.
因為當時耶路撒冷教會很窮困.
裡面很多信徒都很艱難.
一方面因為他們在耶路撒冷的時候.
被猶太教.
猶太人迫得.
當時有饑荒.

$^{241}$所以教會和信徒都很艱難.
馬其頓和阿蓋亞教會的人.
就籌集了一些捐款.
托保羅拿回耶路撒冷.
保羅的計劃就是這樣.
把捐款拿回耶路撒冷之後.
下一步他就想去.
在左手邊的一個長方框.
西班牙.
在那時他覺得.
是一個地極.
去到最遠的地方.
所以保羅最終想去西班牙傳福音.
他計劃拿了捐款.
去了耶路撒冷之後.
接著就去羅馬.
羅馬的綠色框框.
在羅馬他會見到一些老友.
在那裡他就能夠有彼此的分享.
再在羅馬去西班牙.
請聽我讀一些經文.
都是剛才讀過的.
羅馬書15章27節.
「愛邦人既然分享了他們靈性上的好處.
就當把肉體上的需要供給他們」.
這裡說到這班愛邦人.
就是剛才馬其頓和阿蓋亞教會.
我們稱為愛邦人的教會.
他們就從保羅和宣教團隊.
認識了耶穌.
愛邦人既然分享了他們靈性上的好處.
就藉著保羅和宣教團隊得到這個福音.
就當把肉體上的需要供給他們.
他們這個愛邦人的教會就覺得.
他們現在有點困難.
我們應該去幫忙.
「供給」這個字特意用了紅色字.
告訴你.
這個字在希臘文.
聖經的新約原文裡面.

$^{281}$是團契的意思.
原來我們一班屬於上帝的人.
我們一起互相幫助.
就是一個團契.
在羅馬書另外一個地方.
原來每次都要按三次.
在羅馬書第十二章的第十三節.
「聖徒缺乏要幫補,客要一味地款待」.
我們常常都說笑.
請人吃飯請一味地送就行了.
但原來其實在這裡.
客要很專心很用心地去接待款待人.
原來這個「幫補」這個字.
應該是「幫補」這個字.
括錯了.
「幫補」這個字.
「幫助」這個字.
又是團契的意思.
團契就是弟兄姊妹一個愛心實踐的地方.
彼此的支持彼此的接待.
這個馬其頓亞哥亞教會.
是否很有錢呢?.
他們這樣去支持耶路撒冷教會.
在哥林多後書保羅又這樣說.
他們在患難終受大試煉的時候.
仍滿有喜樂.
他們都遇到逼迫的.
在極度貧窮中.
還格外顯出他們樂觀的慷慨.
原來他們自己都不夠的.
在極度貧困當中.
他們都很窮的.
我可以證明.
保羅就說這班馬其頓亞哥亞教會的人.
他們是按著能力.
他有十塊.
他就想一想.
我應該可以奉獻兩塊出來的.
但是他這樣.
而且超過了能力來捐助.

$^{321}$他本來計劃.
我應該自己家這個月要用八塊.
我就奉獻兩塊.
但是誰知道原來他們不是.
是超過了他們的能力來捐助.
結果是奉獻了四塊.
我想保羅或者那些使徒.
都知道他們艱難.
他們說不用了.
但是他們主動再三的懇求.
准他們在這供給聖徒的善上高.
有份.
原來去幫助周圍有需要的人.
有份.
他要有份參與在當中.
這個又是一個團契這個字.
所以今天原來我整篇就是講互相的幫助.
在一個很窮的山區裡面.
有一個我叫做A信徒.
他就抬了一包米.
這一包米就是牧師送給他們的.
他就抬著那包米.
就一直一直走回家.
他家就距離牧師的家很遠.
走了三個多小時.
才回到家.
他就跟家人說.
我們終於有米吃了.
於是他就在袋裡拿了一點點的米出來.
煮了一些稀粥.
他們一家十口.
就每人吃了.
覺得很滿足.
我們看完A信徒.
我們又看一看牧師那裡.
牧師給了那包米.
A信徒之後.
他家裡就剩下兩匙米.
於是叫師母就將這兩匙米煮了飯.
一家四口吃飯.

$^{361}$在吃飯的時候牧師這樣謝土.
叫到A信徒的一家和我們一家都很足夠.
這個A信徒真的很足夠.
因為整袋米.
他就拿一點點來煮稀粥.
可能這袋米就可以吃兩三個月.
但是牧師只有兩匙米.
還煮了飯.
吃完這餐之後.
就不知道下一餐有沒有東西吃.
但是他仍存著一個感恩的心.
去讚美上帝.
當我們為了我們現在所擁有的東西.
我們懂得去讚美的時候.
這是一個真真正正的信心.
但我們通常是怎樣.
通常好像A信徒一樣.
我們最好有一大包米的時候.
感恩我有大包米.
但是當我吃完這餐就沒有的時候.
這個就是一個真正的信心.
就好像我們經常都會說.
有半杯水在我們面前.
我們可以說.
慘了只剩下半杯水.
喝完就沒有了.
不知道怎麼辦.
但是我們可以說.
我們還有半杯水喝.
是吧.
就是一個感恩的信心.
那時候我們讀神學的時候.
我們也是這樣.
同學就會說.
慘了還有一個星期就要交功課了.
又要考試了.
慘了怎麼辦.
怎麼做怎麼讀.
但是有些同學就說.
還有一個星期的時間.

$^{401}$我可以做功課.
我可以去讀書.
感恩感恩感恩.
這個信心的功課.
耶和華是我的牧者.
我必不自缺乏.
我吃完這餐就沒有了.
我講不講得出這句呢.
我們繼續看羅馬教會.
剛才他說.
再讀15章27節.
這是他們樂意的.
其實也算是所欠的債.
這個就說到.
他欠的債.
就是福音的債.
他們幫我們.
好謝謝.
下一幅.
是福音的債.
就是當我們領受了福音.
我們得到耶穌基督的救恩.
我們就有責任繼續傳開它.
他就說了.
馬其頓和亞哥亞的愛邦教會.
他們得到這個福音.
他們就要繼續將這個福音傳開.
另外.
再下一幅.
另外.
我先回到上面.
可能我弄漏了一幅.
另外一個就是分享物質.
剛才那節聖經說.
因為愛邦人既然分享了他們靈性上的好處.
得到福音.
然後他們就把肉體上的衰用.
供給他們.
他們就要有責任.
見到原本傳福音的地方有困難.

$^{441}$他們就去施以援手.
雖然他們自己都不夠.
他們都很困難.
但是他們來去支持幫手.
這兩樣東西.
都是我們每個信徒都要去學習的.
第一.
我們就是這個福音的債.
我們需要將福音的恩典一路傳開去.
第二.
我們是一個團契.
不只是紅恩堂.
是我們整個教會.
甚至是外地的教會.
全世界的信徒.
我們見到哪裡有需要的時候.
我們有領受的時候.
我們就需要去施以援手.
這是上帝的心意.
好.
團契就是一個愛心實踐的地方.
所以無論我們有錢或者沒錢.
或者我們要出多一點或者出少一點.
總之在裡面我們就要學習分擔,分享.
我們不能夠攪在一起說.
我只是來到團契的時候.
我來到教會的時候.
我只是聽別人分享.
我不會說的.
或者我不會做的.
你們煮吧.
你們弄吧.
我什麼都不做的.
如果當我們抱著這樣的心態的時候.
這個不是一個真正的團契.
近日.
某一個團契.
她的那位姐妹的丈夫患了癌症.
醫生就宣佈了她已經擴散了全身的.
還有過往用的藥已經全部無效了.

$^{481}$也再也找不到新的藥可以幫助她.
很難過.
整個人的心情很沉.
那天那個印尼姐姐確診了.
現在真的很麻煩.
一確診了就不能出門.
買不到菜.
還有要跟她一家隔離.
所以什麼都好像天都塌下來一樣.
很辛苦.
她就發出呼籲.
有沒有人可以幫忙找到兼職的人來幫我.
但是有一位姐妹知道之後.
她就打電話跟她聯絡.
她說你不用找兼職的.
很難找的.
她願意每一天幫她買菜.
煮好一天三餐.
送到她家裡.
我聽到的時候很感動.
一個跟她在一起的群體裡面.
知道人家這麼困難裡面.
她就主動.
我幫你買菜.
煮好菜.
送到你家裡.
這個是耶穌非常悅耳的.
我們看看下一幅.
這個五餅二魚.
很多人都聽過這個耶穌基督的神蹟.
其實那天早上的時候.
耶穌就知道他的表哥.
屍逝若瀚.
就被斬了頭.
很難過.
他本來是想去那裡安靜一下.
但是不知道為什麼那些人.
那時候沒有手提電話.
很多人就跟著耶穌.
要聽他的講道.

$^{521}$有五千多個人.
只是男人.
未算婦女和小孩.
五千多個人一直跟著耶穌.
要聽他的講道.
但是耶穌雖然心情很煩躁.
但是也向那些人講上帝的信息.
講到黃昏的時候.
門徒就說.
講了這麼久.
可以了吧.
想叫耶穌說.
叫那些人自己回家.
快點回家吃飯休息一下.
耶穌說什麼?.
不要叫他們去找吃的.
你們給他們吧.
如果我們是當天的門徒怎麼辦?.
嚇得飛起來.
五千多人.
去哪裡找這麼多東西吃給他們?.
還有錢何來?.
但是這個時候.
一個很單純.
我們都覺得他很慷慨的小孩.
他就拿了一個藍色的食物.
去到耶穌那裡.
是五個餅.
兩條魚.
他問耶穌.
夠了沒有?.
這裡的食物夠了沒有?.
是不是很單純?.
五餅一魚.
問耶穌夠了沒有?.
肯定不夠了.
但結果大家都知道.
這五餅一魚.
可以餵飽了五千多人.
而且還剩下十二個男子的食物.

$^{561}$你和我都絕對相信.
在大群人裡面.
不可能只有這個小孩有食物.
我相信很多人都有.
不過.
那些人不肯拿出來.
只是這個很單純的小孩.
願意將他自己那一籃的食物.
拿出來給耶穌基督.
就是因為他這個單純的心.
一個這麼慷慨的心.
拿出來的時候.
耶穌基督就藉著這一籃小小的食物.
行了這個祈禱.
餵飽了五千多人.
所以,丁姊妹.
在你和我的生命裡面.
其實都有五餅一魚.
我們的能力.
我們的時間.
我們的恩賜.
我們的財綱.
我們的財物.
我們都有的.
有人多一些.
有人少一些.
但只要我們肯將他們被主使用.
藉著這些我們生命一點的.
小小的恩賜.
小小的能力.
小小的時間.
小小的財物.
我們去服侍主.
服侍人.
只要肯拿出來的時候.
上帝就可以藉著我們.
拿出來的小小.
就成為了很多人的祝福.
所以各位丁姊妹.
不要再覺得自己什麼都不懂了.

$^{601}$覺得自己微不足道.
自己是一個小小的小蝦.
不要再這樣想了.
就好像那些門徒.
拿著五餅一魚的時候.
只得五個餅.
兩條魚炸.
但是無論是多微小.
只要我們獻上出來的時候.
上帝就會使用的了.
我想跟大家唱一首.
剛才你們唱過的詩歌.
《每天每刻獻上為你》.
請梁惠玲唱.
我很喜歡後面的那句.
我要每天每刻都讚頌你是神.
你施恩典.
你手相獻.
令我們有信心.
我要每天每刻都讚頌你是神.
獻出一生教於主的手中.
丁姊妹.
無論你覺得自己是厲害還是不厲害.
或者你覺得自己是多微小都好.
只要肯交你自己出來給耶穌.
祂可以施恩典給很多很多人.
我們一起唱這首詩歌.
《每天每刻獻上為你》.
神起聖結.
當得主歸.
內心歡呼拍掌.
感我全屬你.
我的主.
神你帶領.
恩典結果.
內心感恩喝彩.
感獻上愛你.
屬於主.
我要每天每刻都讚頌你是神.
你施恩典.

$^{641}$你手相獻.
令我們有信心.
我要每天每刻都讚頌你是神.
獻出一生教於主的手中.
我們一起禱告.
因為聖靈在我們每一個相信的人的心靈裡.
給我們力量.
給我們很多的恩典.
祂真的經歷過.
在我們的生命裡.
上帝你給了很豐富很豐足我們.
多謝你.
你加入了這個大群體裡.
多謝你.
最後求你一次福給我們每一位弟兄姊妹.
不再覺得自己微小.
不再覺得自己沒用.
而在整個大群體裡.
我們怎樣呈獻給你.
主求你打開我們的心靈.
以致讓我們都經歷到.
只要我們呈獻了.
多微小都好.
主你會閱立.
並且你藉著這個微小.
可以能夠去祝福很多很多的人.
主求你就讓紅印堂裡面.
每一位弟兄姊妹.
都是你的精兵.
在全方的事情的上高.
在關顧弟兄姊妹.
在弟兄姊妹彼此的扶持.
在這個奉獻裡面.
我們大家都拜上.
盡我們的能力所有.
拜上去榮耀神你自己.
多謝你垂聽禱告.
奉主基督的聖名.
Amen.
\newpage



\section{使徒行傳 9:1-31}
\label{sec:MMKQi99_fYg}
\textbf{2023.06.04 《一個全新的你》 使徒行傳 9\_1-31 (和修版) 曾敏傳道}
\newline
\newline
連結: \href{https://youtube.com/watch?v=MMKQi99_fYg}{\texttt{https://youtube.com/watch?v=MMKQi99\_fYg}} ~~~~ 語音日期: 2023-06-07
\newline
\newline
\hyperref[sec:33y9hhVAb_s]{\small{< < < PREV SERMON < < <}}
~
\hyperref[sec:index_chronic]{\small{[返順時目]}}
~
\hyperref[sec:index_scriptual]{\small{[返順卷目]}}
~
\hyperref[sec:B5mDgas2LUk]{\small{> > > NEXT SERMON > > >}}
\newline
\newline
使徒行傳 9:1-31
\newline
\begin{longtable}{cl}
\hline
\hline
章節 & 經文 (和合本修訂版)\\
\hline
9:1 & \begin{tabularx}{0.7\textwidth}{X} 掃羅不斷用威嚇兇悍的口氣向主的門徒說話。他去見大祭司， \end{tabularx} \\ \\ \relax
9:2 & \begin{tabularx}{0.7\textwidth}{X} 要求發信給大馬士革的各會堂，若是找著信奉這道的人，無論男女，都准他捆綁帶到耶路撒冷。 \end{tabularx} \\ \\ \relax
9:3 & \begin{tabularx}{0.7\textwidth}{X} 掃羅在途中，將到大馬士革的時候，忽然有一道光從天上下來，四面照射著他， \end{tabularx} \\ \\ \relax
9:4 & \begin{tabularx}{0.7\textwidth}{X} 他就仆倒在地，聽見有聲音對他說：「掃羅！掃羅！你為甚麼迫害我？」 \end{tabularx} \\ \\ \relax
9:5 & \begin{tabularx}{0.7\textwidth}{X} 他說：「主啊！你是誰？」主說：「我就是你所迫害的耶穌。 \end{tabularx} \\ \\ \relax
9:6 & \begin{tabularx}{0.7\textwidth}{X} 起來！進城去，你應該做的事，必有人告訴你。」 \end{tabularx} \\ \\ \relax
9:7 & \begin{tabularx}{0.7\textwidth}{X} 同行的人站在那裡，說不出話來，因為他們聽見聲音，卻看不見人。 \end{tabularx} \\ \\ \relax
9:8 & \begin{tabularx}{0.7\textwidth}{X} 掃羅從地上起來，睜開眼睛，竟不能看見甚麼。有人拉他的手，領他進了大馬士革。 \end{tabularx} \\ \\ \relax
9:9 & \begin{tabularx}{0.7\textwidth}{X} 他三天甚麼都看不見，也不吃也不喝。 \end{tabularx} \\ \\ \relax
9:10 & \begin{tabularx}{0.7\textwidth}{X} 那時，在大馬士革有一個門徒，名叫亞拿尼亞。主在異象中對他說：「亞拿尼亞！」他說：「主啊，我在這裡。」 \end{tabularx} \\ \\ \relax
9:11 & \begin{tabularx}{0.7\textwidth}{X} 主對他說：「起來！往那叫直街的路去，在猶大的家裡，去找一個大數人，名叫掃羅；他正在禱告， \end{tabularx} \\ \\ \relax
9:12 & \begin{tabularx}{0.7\textwidth}{X} 在異象中看見了一個人，名叫亞拿尼亞，進來為他按手，讓他能再看得見。」 \end{tabularx} \\ \\ \relax
9:13 & \begin{tabularx}{0.7\textwidth}{X} 亞拿尼亞回答：「主啊，我聽見許多人講到這個人，說他怎樣在耶路撒冷多多苦待你的聖徒， \end{tabularx} \\ \\ \relax
9:14 & \begin{tabularx}{0.7\textwidth}{X} 並且他在這裡有從祭司長得來的權柄，要捆綁一切求告你名的人。」 \end{tabularx} \\ \\ \relax
9:15 & \begin{tabularx}{0.7\textwidth}{X} 主對他說：「你只管去。他是我所揀選的器皿，要在外邦人、君王和以色列人面前宣揚我的名。 \end{tabularx} \\ \\ \relax
9:16 & \begin{tabularx}{0.7\textwidth}{X} 我也要指示他，為我的名必須受許多的苦難。」 \end{tabularx} \\ \\ \relax
9:17 & \begin{tabularx}{0.7\textwidth}{X} 亞拿尼亞就去了，進入那家，把手按在掃羅身上，說：「掃羅弟兄，在你來的路上向你顯現的主，就是耶穌，打發我來，叫你能再看得見，又被聖靈充滿。」 \end{tabularx} \\ \\ \relax
9:18 & \begin{tabularx}{0.7\textwidth}{X} 掃羅的眼睛上立刻好像有鱗一般的東西掉下來，他就能再看得見，於是他起來，受了洗， \end{tabularx} \\ \\ \relax
9:19 & \begin{tabularx}{0.7\textwidth}{X} 吃過飯體力就恢復了。 \end{tabularx} \\ \\ \relax
9:20 & \begin{tabularx}{0.7\textwidth}{X} 立刻在各會堂裡傳揚耶穌，說他是神的兒子。 \end{tabularx} \\ \\ \relax
9:21 & \begin{tabularx}{0.7\textwidth}{X} 凡聽見的人都很驚奇，說：「在耶路撒冷殘害求告這名的不就是這個人嗎？他不是到這裡來要捆綁他們，帶到祭司長那裡去嗎？」 \end{tabularx} \\ \\ \relax
9:22 & \begin{tabularx}{0.7\textwidth}{X} 但掃羅越發有能力，駁倒住在大馬士革的猶太人，證明耶穌是基督。 \end{tabularx} \\ \\ \relax
9:23 & \begin{tabularx}{0.7\textwidth}{X} 過了好些日子，猶太人商議要殺掃羅， \end{tabularx} \\ \\ \relax
9:24 & \begin{tabularx}{0.7\textwidth}{X} 但他們的計謀被掃羅知道了。他們晝夜在城門守候著要殺他。 \end{tabularx} \\ \\ \relax
9:25 & \begin{tabularx}{0.7\textwidth}{X} 他的門徒就在夜間用筐子把他從城牆上縋了下去。 \end{tabularx} \\ \\ \relax
9:26 & \begin{tabularx}{0.7\textwidth}{X} 掃羅到了耶路撒冷，想與門徒結交，大家卻都怕他，不信他是門徒。 \end{tabularx} \\ \\ \relax
9:27 & \begin{tabularx}{0.7\textwidth}{X} 只有巴拿巴接待他，領他去見使徒，把他在路上怎麼看見主，主怎麼向他說話，他在大馬士革怎麼奉耶穌的名放膽傳道，都述說出來。 \end{tabularx} \\ \\ \relax
9:28 & \begin{tabularx}{0.7\textwidth}{X} 於是掃羅在耶路撒冷同門徒出入來往，奉主的名放膽傳道， \end{tabularx} \\ \\ \relax
9:29 & \begin{tabularx}{0.7\textwidth}{X} 並和說希臘話的猶太人講論辯駁，他們卻想法子要殺他。 \end{tabularx} \\ \\ \relax
9:30 & \begin{tabularx}{0.7\textwidth}{X} 弟兄們知道了，就帶他下凱撒利亞，送他往大數去。 \end{tabularx} \\ \\ \relax
9:31 & \begin{tabularx}{0.7\textwidth}{X} 那時，猶太、加利利、撒瑪利亞各處的教會都得平安，建立起來，凡事敬畏主，蒙聖靈的安慰，人數逐漸增多。 \end{tabularx} \\ \\
[1ex]
\hline
\hline
\end{longtable}
$^{1}$我們再一次一同禱告.
很愛我們.
你的愛總是轉化成不同的行動.
是讓我們經歷得到,感受得到.
你愛我們,就是讓我們感受得到.
你愛我們,就是渴望我們生命去改變.
願你對我們眾人說話.
讓我們明白你的心意如何.
讓我們到最後都能夠了解.
如何用我們的生命去成長和改變.
願聖靈在我們心中動功.
賜我們智慧聰明,去明白神的話語.
奉耶穌基督的名字祈禱,阿們.
今天我們要講的是一個全新的你.
我們在PowerPoint開始的時候.
不是在中間.
一個全新的你是怎樣呢?.
讓我們去思考一下.
我們都信主一段日子.
有些信主幾十年.
有些信了不久.
但我相信都是有些經歷和體會.
今天是一個時間去看看.
一個全新的你是怎樣的呢?.
我們從一個人物的生命中.
去看到全新的生命是怎樣的.
我們來看看保羅.
我想問大家一個問題.
在聖經裡有保羅書信.
保羅是作者.
一共寫了一些書信.
我想問大家知不知道.
保羅書信一共有多少卷呢?.
知道的,大聲叫出來.
保羅書信有多少卷呢?.
沒有人叫出來.
擦卷.
保羅書信有多少卷呢?.
看回聖經這個人物的見證.
是很有代表性的.

$^{41}$因為他很有影響力.
寫了很多卷書.
保羅書信一共有13卷.
所以在新約的書卷裡.
很多卷書都是保羅寫的.
而在裡面他所寫的.
他所教導的.
影響了歷世歷代很多信徒.
他不只是說.
他怎樣去看耶穌基督的舊因.
看教會.
看整個信徒群體的生命.
其實他也會以他的生命.
作為一個榜樣給你看.
信耶穌的人.
生命是會改變的.
這是最寶貴的.
講道的人不只是講道.
他會給你看一個改變的生命.
是有說服力的.
所以以後看聖經的時候.
看回保羅的教導的時候.
看回這個生命.
親自經歷的舊因.
親自被耶穌基督的舊因.
所改變的生命.
是不同的.
我們來看看保羅是怎樣改變的.
保羅的出生背景.
保羅本身是在.
大掃城.
出自一個.
是屬於變亞敏支派的.
猶太家庭.
基尼加在地圖上.
這個位置.
一個沿海的位置.
大掃是當中一個大城市.
大掃這個地方.
是一個商業中心.

$^{81}$非常繁華.
很崇尚希臘文化的思想.
但這個也是一個大學城.
什麼意思?.
是一個大學城.
即是裡面的人有什麼特色呢?.
是非常注重學習的.
很看重學習的一個地方.
一個極高文化水平的地方.
保羅就在這裡出生.
在這裡生活.
所以保羅也深受這個地方的影響.
他自己非常注重學習.
他自己也有一個很高深的教育.
當然我們一會兒看他的教育.
也包括了哪一方面呢?.
保羅的家族.
雖然處於一個希臘文化的社會.
但他爸爸是一個法利賽人.
法利賽人很熟悉神的律法.
很嚴格地遵守神的律法.
怎樣去成為一個聖傑的生命呢?.
可以遵守很多很多條例.
他就生在一個這麼虔誠的猶太家庭裡.
由他學習行路起到十二,三歲之前.
他就被送到猶太會堂的私塾裡.
那些學校裡面.
接受律法和先知的傳統教育.
在那裡他閱讀希伯來聖經.
所以他是熟悉聖經的.
也學習當時各地的通行的亞蘭語.
但很特別.
他雖然熟悉希伯來文的聖經.
卻不是真正認識上帝.
我去看聖經和真正認識神.
在寶萊來說是有一段距離.
他到了十三歲左右.
他就去到耶路撒冷和他姐姐一起住.
在那裡起先接受一般的正規教育.
他就跟隨猶太教的大師希利之孫加瑪烈.

$^{121}$加瑪烈是一個很有名的猶太人拉比.
他學習猶太的法典.
接受拉比的訓練.
大概在二十到三十歲左右.
所以可以這樣說.
你看他整個家庭的成長背景.
在猶太人的社會裡面.
他當時應該覺得自己有很多東西值得自豪.
他出生成長的城市是一個高文化教育的地方.
從小到大他在一個虔誠的猶太人家庭裡面成長.
接受一個非常豐富的.
一個很深的律法教導.
跟著一個名師加瑪烈.
跟著他學習猶太的法典.
在猶太人的社會裡面.
他是有一個很崇高的位置.
很被人尊崇.
所以可以這樣說.
其實他在他的書信裡面有時候也會提到.
說真的,以這些條件來說.
有很多東西他真的可以拿來.
我們講得俗一點叫做「曬命」.
可以炫耀的,自己的背景.
在猶太人社會裡面.
但是你看他的生命,他的轉變.
他成為拉比,他也需要學習一門手藝.
一個謀生的技能.
他選擇製造帳棚.
所以後來你會看到他一直做織帳棚的工作.
去養活自己,一直傳播音.
這個是保羅.
當時他的名字叫做「素羅」.
素羅原先是一個怎樣的人呢?.
雖然他是來自一個敬虔的猶太人家庭.
但是可以這樣說.
他不是真真正正的認識上帝.
而且他這個階段,你看這段經文.
剛才大家讀過的.
他是不斷地用威嚇,空巷的口氣.
向主的門徒說話.

$^{161}$還要發信給大馬士革的各個會堂.
如果是找著,他去見大祭司.
要求他發信.
如果是找著信奉這道的人.
無論男女都要他捆綁大道耶路撒冷.
在這個階段,他是迫不得已的基督徒.
他一直可以這樣說,他信奉猶太教.
他一直覺得有一群人出來說.
很尊崇耶穌.
又說他成就救恩,去釘十字架.
又說他復活.
一群追隨者一直在傳耶穌的事蹟.
怎樣釘十字架,成就救恩,復活.
這一群人簡直是一派胡言.
妖言惑眾.
一定要抓住這群人.
不要讓他們亂散播這些胡言亂語.
擾亂我們的百姓.
今天我們最重要的是跟從摩西的律法.
受要求的律法,很多條.
法利賽人後來訂立了很多的律例.
甚至不是上帝自己的心意.
不是上帝自己訂立的,而是自己加了很多.
為了追求聖潔的生活.
我們在聖潔的子民.
和世俗裡的骯髒的子民分開.
這是猶太教.
所以當時保羅迫使基督徒.
他不認耶穌基督是神的兒子.
也不同意基督徒所說的耶穌的事情.
他不但不同意和不認同.
還要去抓拿這些人.
這段時間他還說要去見大祭司.
請大祭司授權.
請大祭司發公銜給大馬士革的每個會堂.
其實大祭司當時活動的中心應該在耶路撒冷.
大馬士革可以說是離耶路撒冷一段距離.
大家看看有一段路,看到嗎?.
有一個很長的路.
換言之,他想大祭司的影響力擴闊一點.

$^{201}$不要只在耶路撒冷,一直到大馬士革.
誰信耶穌的就抓他.
怎樣抓呢?.
要綁住他們,拉他們回耶路撒冷.
山長水遠拉他們回來.
要他們受審判.
這裡說的是威嚇.
什麼叫威嚇?.
不要信耶穌.
聽著,你信耶穌,就抓你坐牢.
殺死你,這樣叫威嚇.
很凶悍的口氣.
他是這樣的,他真的準備去抓人.
然後還要綁住他們,帶他們去耶路撒冷.
這麼長的路綁住他們.
你看到當時的蘇羅是一個這樣的人.
一個很凶惡,迫不得已的基督徒.
不是真正認識上帝.
但他一直在為信仰大發熱心.
他覺得這樣是對的.
他原來的生命就是這樣.
但當他和主耶穌相遇之後.
生命就不再一樣.
他相遇的經歷是怎樣呢?.
大馬士革的經歷就是.
他一直由耶路撒冷前往大馬士革.
那段路很遙遠.
地圖上很遙遠,大概是六天的路程.
在這六天裡,他一直向前走.
快到的時候,忽然有一道光在天上照下來.
四面照射著他.
是一道很強的光,不是一個光柱.
是很大範圍的光環繞著他.
他感受到很重要的,超自然的.
保羅,素羅自己感受得到.
和他同行的人都感受到.
那道光是超自然的光.
不是一般的光.
這些光環繞著素羅.
以至他仆倒在地.

$^{241}$在這個超自然的情況下,仆倒在地.
素羅以前很意氣風發.
拿著大祭司的公銜.
準備捉那些基督徒,他覺得自己很有權柄.
很威風.
他就想捉那些基督徒,把他們拉回耶路撒冷.
好像不可一世一樣.
這一刻,就在這些光當中仆倒在地.
這是什麼光呢?.
是耶穌基督的光.
因為跟著耶穌基督說話.
所以是耶穌基督的光.
通常在神顯現的時候就會有光.
如果大家有機會去看經文.
大家知道,在約翰一書裡說到.
神就是光.
神是光的時候.
他的同在,他的出現就會有光.
而且是很大的一份榮光.
在神的面前,素羅.
無論多麼的厲害,仆倒在地.
要謙卑在上帝的面前.
接著聽到超自然的聲音.
有聲音跟他說:素羅,素羅,你為什麼迫害我?.
奇怪了,誰說我迫害他呢?.
他知道這一刻的經歷是超自然.
他就問:主啊,你是誰?.
主說:我就是你迫害的耶穌.
可能大家來到這裡會覺得很奇怪.
耶穌好像不是這個時候出現的.
耶穌是在之前的.
之前的時候,耶穌已經被釘死在十字架.
被埋葬,然後第三天從死裡復活.
耶穌的事已經在之前發生了.
素羅在那時候也沒有迫害過耶穌.
為什麼耶穌說:我就是你迫害的耶穌呢?.
你想想,素羅迫害了誰?.
大家回答我.
他迫害了誰?.
迫害基督徒嗎?.

$^{281}$原來迫害基督徒就是迫害耶穌.
迫害基督徒就是迫害背後的耶穌.
所以今天我們大英姐妹出去傳福音.
我們很多時候都會說:你被人拒絕.
你要知道什麼?.
別人拒絕誰?.
拒絕你背後的那位耶穌.
今天有人迫害你.
也是迫害你背後的那位耶穌.
所以今天素羅在迫害基督徒.
就是迫害他們背後的那位耶穌基督.
耶穌說:你為什麼迫害我?.
很特別的是,在這個聲音裡.
同行的人都聽到聲音.
但是不太知道裡面的內容是在說什麼.
他們只聽到聲音.
但素羅自己聽到聲音裡的內容.
所以素羅聽到耶穌問他:你為什麼迫害我?.
素羅聽到這個聲音.
他也知道是一個超自然的聲音.
而且這個聲音在說.
他自己是誰?.
他自己是耶穌.
素羅很清楚.
今天他聽到的聲音不是來自那個人.
那個耶穌基督.
他是在聽上帝.
因為是一個超自然的光.
超自然的聲音.
他知道這是神.
原來耶穌基督是真的.
原來那群基督徒是什麼是真的?.
他們真的是復活.
基督徒所傳揚的.
基督受死,埋葬,復活是真的.
耶穌基督是神.
耶穌基督真的是猶太人期望的救世主.
耶穌真是基督.
在這一刻.
素羅才知道原來自己一直以來都錯了.

$^{321}$原來以往自己迫害基督徒是不對的.
是錯了.
在這一刻.
他聽到這些說話.
他知道同行的人只是見證.
素羅經歷了一些超自然的事情.
有超自然的光,有聲音.
但他們不太知道實際發生什麼事.
在這裡.
耶穌和素羅說.
起來,進城去.
我相信這一刻.
素羅心裡一定有很多思想.
他一直大發熱心.
覺得自己是對的.
我維持我這個猶太的信仰.
在這一刻我知道.
耶穌基督是真真實實的.
是真是神的.
他真的是復活.
基督徒所傳揚的是真的,是對的.
自己錯了的時候.
你心裡有什麼想法?.
會不會很容易陷入很多的罪疚.
很多的難過.
很多很複雜的思考.
但這一刻.
耶穌立刻叫他向前走.
起來,進城去.
你應該做的事.
有人會告訴你.
但這裡很有趣.
耶穌沒有完全說明.
每一個細節要做什麼.
只是說你要進城去.
有人會將你要做的事告訴你.
在那時候.
素羅真的起來.
發生了一件很特別的事.
他發覺自己瞎了.

$^{361}$他發覺自己瞎了,看不到.
所以他不能夠自己走進大馬士革.
是有人拉著他進大馬士革.
這個和耶穌相遇的經歷.
實在是太奇妙了.
不是每個人都會有這樣的經歷.
但至少令到素羅非常深刻.
非常深刻.
一定不會忘記.
他進了大馬士革三天.
什麼都看不到.
亦不吃不喝.
這裡的不吃不喝是什麼呢?.
我們從後面的經文.
發現原來那三天裡他有禱告.
所以耶穌基督說他禱告.
耶穌和亞拿利亞說話的時候.
說素羅正在禱告.
如果從後面看.
不吃不喝,大家覺得他在做什麼?.
除了祈禱.
禁食.
他在禁食禱告.
和耶穌這個很深刻的相遇.
令到他要向上帝禱告.
生命開始改變.
他在裡面開始改變.
他的思想不再一樣.
不是像以前那樣.
我什麼時候去逼迫基督徒.
什麼時候去拘捕他們.
這些妖艷惑眾的人.
要剷除他們,不要危害我們的群體.
他開始改變,要向上帝禱告.
要向耶穌禱告.
他的禱告肯定不同了.
為什麼?他這一刻真正認識耶穌基督.
在這裡,我很想問一件事.
各位兄弟姐妹.
在教會這麼多年.

$^{401}$我們是否真的認識耶穌基督呢?.
如果我們真正認識耶穌基督.
我們對耶穌基督的感受和回應.
應該是不一樣的.
沒錯,我們不會像素羅那樣.
忽然間走著走著.
不知道去哪裡.
有大光環繞著我們.
我們的眼睛看不見.
我們看不見上帝跟我們說話.
我們不會每個人都這樣經歷.
但是我這樣寫的包單.
我們每個人都有經歷耶穌.
在這裡在座的我知道.
有很多人已經洗禮.
大家都有寫過洗禮見證,是嗎?.
是呀.
有寫過洗禮見證.
有寫過耶穌基督怎樣改變你的生命.
即是大家都有經歷耶穌.
有些人還是.
有很多神蹟奇事.
但是我想問.
大家是否真正認識耶穌基督?.
好像素羅一樣.
真真正正認識.
禱告不一樣.
禱告的對象也不同.
他不再逼迫主.
反而向主禱告.
在這個時候.
你見到素羅的生命有更多更大的轉變.
上帝有很奇妙的安排.
在大馬士革有一個門徒.
叫做亞蘭尼亞.
主在異象當中.
吩咐他去找素羅.
去到大數這個地方找素羅.
主都見到素羅正在禱告.
很特別.

$^{441}$耶穌在素羅生命裡工作.
素羅禱告的時候.
他又見到異象.
亞蘭尼亞又見到異象.
素羅又見到異象.
他見到異象的人.
叫做亞蘭尼亞.
所以耶穌同時在兩個人生命裡工作.
預備兩個人的心.
讓他們可以相遇.
素羅見到亞蘭尼亞.
進來為她按手.
讓她能夠再看見.
這是一個醫治的動作.
素羅知道有一個人.
叫做亞蘭尼亞.
會來醫治自己.
這是出於耶穌的安排.
耶穌也讓亞蘭尼亞看到.
他要去找一個人.
叫做素羅.
一說到這個名字.
亞蘭尼亞很有反應.
為什麼?.
素羅太出名了.
出名什麼?.
碧白基督徒.
很出名的.
人們一聽到他的名字都害怕.
亞蘭尼亞說.
我聽到很多人說這個人.
是不敢多苦待你的聖徒.
他是這樣的.
還有在祭司拿到權柄.
要綁那些求告你命的人.
是這樣的人.
現在耶穌要我去找他.
亞蘭尼亞心裡有掙扎.
耶穌怎樣說?.
祂說你儘管去.

$^{481}$素羅是我揀選的器皿.
要在外邦人.
外邦人這裡是說.
是列國.
原來素羅要在列國當中.
是君王在以色列人面前.
去宣揚我的名.
素羅是耶穌揀選的器皿.
耶穌要用他.
有時候上帝要用人.
不是我和你能夠明白的.
我和你去看.
總是覺得那個人.
不夠清.
那個人這麼壞.
你也會用嗎?.
上帝自有方法.
可能你曾經覺得他不夠清.
你覺得他很差的人.
他是一個很有力量.
為上帝工作的人.
只有神才知道.
上帝揀選他的僕人.
上帝自然會在那個人生命裡.
去成就他的旨意.
這裡耶穌說.
這個人是我揀選的.
他要這樣做.
這個職責.
要向外邦人傳福音.
這是其中一個使命.
我也要指示他.
為我的名必須受很多苦難.
素羅要為主受苦.
耶穌說完這番話.
亞蘭尼亞就去了.
真的去到那個家.
按手在素羅的身上.
亞蘭尼亞在這裡.
做了一件很重要的事情.

$^{521}$他不單止醫治素羅.
而是他讓素羅感受到.
教會的接納.
是怎樣呢?.
他首先去到他家.
按手在他身上.
這是肢體語言.
這個按手不單止有醫治的作用.
還有一個很重要的.
接納的表達.
他的稱呼很好.
素羅弟兄.
不是素羅先生.
素羅某某.
是素羅弟兄.
稱兄道弟.
這是主內的稱呼.
素羅那一刻真的感受到.
我被接納.
教會裡的人接納.
很好.
亞蘭尼亞很明顯知道.
素羅之前的經歷.
在他來的路上.
向顯然的主祭耶穌叫我來.
打訓我來.
讓我看見.
而且被聖靈充滿.
素羅要經歷聖靈的洗.
不單止是聖靈降臨在他生命裡.
而是聖靈充滿.
聖靈的洗禮.
這是亞蘭尼亞的使命.
要做在素羅身上.
醫治他.
素羅經歷醫治.
他經歷很大的醫治.
他的眼睛有一個像輪子般的東西.
掉下來.
他看不見.

$^{561}$他經歷醫治之後.
你看見他.
就受洗了.
他受洗.
他真真正正歸入主的名下.
他結束禁食.
他開始吃東西.
恢復了體力.
經文說.
他和大馬士革的門徒一起.
在這裡.
他有一個很寶貴的經歷.
他成為基督徒.
一個已經很大的轉變.
本身是一個迫不得已的基督徒.
他開始向耶穌基督禱告.
他成為基督徒.
他接受洗禮.
轉換身份.
由迫不得已者成為跟從者.
不單止是提到這裡.
他的生命有一個很大的改變.
蘇羅和大馬士革的門徒.
住了一段時間後.
在各個會堂串揚耶穌.
說他是神的兒子.
就立刻傳福音.
我們通常是甚麼人傳福音的?.
我們教會通常是甚麼人傳福音的?.
傳道人?.
姐妹?.
不一定是姐妹.
有弟兄.
我們最近多了弟兄.
會越來越多弟兄傳福音.
很多時候別人的觀念.
一定是信了很久.
經常都會說.
通常邀請到弟兄姊妹傳福音.
我不曉得怎樣說.

$^{601}$我先學會怎樣說.
我當然要學會.
不然我怎樣說.
我不懂得說.
我要說得很新.
不停地說.
這些信了十多年.
二十多年.
他們才傳.
排次序都這樣排.
三十多年的.
二十多年的.
十多年的.
幾年的.
等我過一段時間.
過幾年才傳.
剛剛信耶穌的人.
我們專門拉他出去傳福音.
這麼可怕.
為什麼剛剛信耶穌的人.
我們就立刻拉他出去傳福音?.
這麼恐怖.
我不回來了.
因為他那時候.
剛剛經歷耶穌新鮮熱辣的經歷.
他出去傳的時候.
再提醒自己.
我經歷耶穌基督的救恩.
他救我生命.
多寶貴.
其實這時候是最好的.
所以我們舞會的時候.
很喜歡邀請人讀三福.
讀短宣.
去傳福音.
然後信徒成長得很快.
不用經常說.
我新人你不要叫我.
我不懂得說.
兩下出去.

$^{641}$我懂得說.
我們很少人這樣說.
我來到這裡就經常聽到.
不用的.
你講自己見證就夠了.
你就講給人聽你怎麼信耶穌.
我們今天.
你去街市.
今天看到很好的魚.
很好的菜.
你需要學一下怎麼說嗎?.
讓我先回去學習.
怎樣稱讚那條魚很漂亮.
讚那些菜很漂亮很便宜.
你去七嬸八叔.
氣他們買菜.
不用的.
你轉過頭就馬上出去說.
今天街市那條魚很漂亮.
那些菜很好.
買呀.
今天去某個餐廳.
吃的東西很好.
不用別人教.
你說不行.
傳福音和買菜魚不同.
什麼不同?.
你經歷了.
你經歷了的東西可以很含糊嗎?.
不會的.
你看看.
素羅一經歷就立刻出去.
立刻很想說一個字眼.
在國會堂裡面.
不只是試試.
試試感覺如何.
在國會堂裡面.
他經歷到耶穌基督是真實的.
神的兒子這麼好.
馬上說吧.

$^{681}$救人的生命.
是不是?.
聽見的人都很驚奇.
這個在耶路撒冷殘害.
求告者明.
不是就這個人.
那時候迫百基督徒.
還說要綁著他們去哪裡.
這個人現在到處去.
去會堂就在那裡.
去傳福音.
心裡很火熱.
聽姐妹問我和你.
我們心裡有沒有這份火熱呢?.
你說這些是傳道人做的.
誰說的?.
是每個基督徒都會做的.
真正經歷上帝恩典.
耶穌基督恩典的人就會做.
因為你知道他救恩的寶貴.
你知道他的真實.
他還要和猶太人辯駁.
證明耶穌是基督.
是神的兒子.
他傳揚福音.
為主作供.
是立刻.
還有不只是傳揚福音為主作供.
過了好些日子.
他的名聲又響起來.
這次不是因為迫百基督徒.
這次是傳揚福音.
他的名氣開始響起來.
他為主發熱心.
大家都知道有個人叫素羅.
到處去傳福音.
猶太人覺得這是一個威脅.
於是就開始要殺他.
昔日他是迫百者.
今天他變成被迫百者.

$^{721}$你猜他知不知道.
自己將會落到這樣的局面?.
他肯定知道.
昔日他迫百基督徒.
迫百周圍傳揚福音的人.
他怎會不知道.
今天他出來傳福音.
別人就會迫百他.
他知道.
但他選擇要面對迫百.
由迫百人變成選擇.
願意甘心為主受迫百.
在這裡人們計劃要殺他.
還要晝夜在城門口守住.
想著什麼時候抓住他.
然後用門徒救了他.
用匡子將他從城牆除下.
救了他一命.
但在這裡你見到.
保羅為主面對迫百.
和死亡的威脅.
他轉換了位置.
他知道自己要付這個代價.
但他甘心付這個代價.
在這裡.
訴羅說他之後去耶路撒冷.
他想和門徒結交.
但耶路撒冷的門徒.
初初不敢接納他.
不敢相信他已經變了.
仍然留下印象.
是以前迫百基督徒的訴羅.
但在這裡.
你看到神有一個很好的安排.
安排巴拿巴去接待他.
如果沒有了巴拿巴.
作為中間的橋樑.
我想他很難進入耶路撒冷教會的群體.
但有巴拿巴接待他.
帶領他去見使徒.

$^{761}$將他的經歷和使徒們分享.
讓他們能夠信任他.
正正因為巴拿巴所做的事.
訴羅亦能和耶路撒冷的使徒.
教會的門徒們結連.
跟著他和他們一路出入往來.
繼續的傳揚福音.
你看到他面對的迫百越來越多.
面對的危險越來越多.
當時的猶太人.
想法子要殺死他.
他的弟兄們知道了.
又將他帶到凱撒尼亞.
然後送他往大掃去.
在這裡.
保羅他這一生為主的福音.
為傳揚福音受了很多苦.
我們看這段經文.
由23節到28節.
訴羅描述在哥林多後書11章23.
到33節.
訴羅或保羅描述自己.
一生為主受了什麼苦.
我們一起讀.
23節開始.
他們是基督的用人嗎?.
我說句狂話.
我更是.
我比他們忍受更多勞苦.
坐過更多次監牢.
受過無數次的鞭打.
上上冒死.
我被猶太人鞭打五次.
每次四十減去一下.
被棍打了三次.
被石頭打了一次.
遭海難三次.
一晝一夜在深海裡掙扎.
我又屢次行遠路.
遭江河的危險.

$^{801}$盜賊的危險.
同族人的危險.
外族人的危險.
城裡的危險.
曠野的危險.
海中的危險.
賈弟兄的危險.
我勞碌困苦.
上上失眠.
又飢又渴.
忍饑內寒.
赤身露體.
除了這些外表的事以外.
我還有為宗教會操心的事.
天天壓在我身上.
有誰軟弱我不軟弱呢.
有誰跌倒我不焦急呢.
我若必須誇後.
就誇我軟弱的事好了.
我們停在這裡.
你看到保羅的一生受了很多苦.
他甘心 甘願為主受這些苦.
原來當一個人生命轉變的時候.
當你願意為主擺上的時候.
是要付代價的.
那是生命的代價.
可以說他一生豁出去.
為主擺上一切.
為什麼這樣呢.
最後我提一件事.
為什麼願意呢.
保羅為主的命受很多苦.
為什麼呢.
因為他整個價值觀都改變了.
這個價值觀令他願意受苦.
而且令他人生方向都不同.
人生方向不同.
我想和大家一起讀.
不但如此.
我已把萬事當作是有損的.

$^{841}$因我已認識我主基督耶穌為至寶.
我為他已經凋棄萬事.
看作糞土.
為要贏得基督.
保羅由他生命改變.
到這個階段.
願意為主受苦.
願意犧牲.
為何要受那麼多苦難.
因為他真的認識.
我主基督耶穌為至寶.
我為他已經凋棄萬事.
所以我甚麼都願意受.
甚麼都願意捱.
所有事情.
他不會覺得捱.
不會覺得這麼淒涼.
所以在保羅的教導中.
其中一件事常說的是.
常常喜樂.
這句說話出於他的口中.
就有說服力.
他不是人生風調雨順.
一切都很順利.
每天都陽光普照.
沒有患難.
他的人生正正相反.
有很多患難.
很多痛苦.
坐牢.
他仍然要跟你說.
常常喜樂.
他說話就有說服力.
我們喜樂不是因為順利.
喜樂不是因為人生沒有苦難.
是因為有最寶貴的耶穌基督.
這是我們信仰很重要的核心.
你是不是視耶穌為至寶.
我問自己是否視耶穌為至寶.
你視他為至寶.

$^{881}$你甚麼都願意為他做.
你甚麼都願意為他受.
我相信如果我們每一個都這樣看.
我們走在一起就少了很多.
少了甚麼.
紛爭.
說得好.
少了很多紛爭.
少了很多埋怨.
少了很多灰心失意沮喪.
少了很多.
大家一起想.
耶穌基督甚麼都願意.
氣氛完全不同了.
這是我們信仰的核心.
是不是.
去到保羅的階段.
真的很高境界.
到最後這裡.
這是我們教會經常會金句.
金句的意思不是我跟你口啞啞.
我跟你能否做到.
我們一起讀.
預備起.
這不是說我已經得著了.
已經完全了.
這是基督耶穌所要我得著的.
弟兄們.
我不是以為自己已經得著了.
我只有一件事.
就是忘記背後努力面前的.
向著標竿直跑.
要得神在基督耶穌裡.
從上面召我來得的獎賞.
保羅.
為著耶穌基督而瘋狂.
為耶穌基督而活.
耶穌基督就是他的目標.
一生.
向著耶穌基督奔跑.

$^{921}$我和你.
是否向著耶穌基督奔跑呢?.
保羅為主而活.
要贏得基督.
最後這個見證.
很想送給大家.
大家有沒有聽過.
日本的井上勳牧師的生命見證.
原來他曾經來過香港.
有沒有聽過.
這個牧師.
他本身是一個黑社會的人.
是一個黑社會的人.
紋身.
其中一隻手指沒有了血.
一看日本人就知道.
這些是流氓.
這些是黑社會的人.
遠離他遠一點.
他說這個井上勳.
本身出生於日本北海道這個地方.
小時候家裡很貧困.
媽媽有心臟病.
經常在家裡懸懸.
因為家裡貧困.
所以井上勳小時候覺得.
有錢就有幸福.
後來他和朋友一起偷東西.
小時候就開始變成一個小混混.
一個小流氓.
曾經被判入少年院.
後來他長大後.
人生都在流氓行業裡浮浮沉沉生活.
十八歲的時候.
他開始第一段婚姻.
但他整個人狀態很差.
他經常對妻子家暴.
到二十一歲的時候.
他就已經離婚.
這幾年裡.

$^{961}$他只會鬧事打架.
回家後就對太太不好.
後來他離婚之後.
他認識了黑社會的大佬.
跟這個大佬謀生.
跟這個大佬的時候又染上毒癮.
因為這個黑社會大佬有毒癮.
他自己也染上毒癮.
在跟黑社會幾年裡.
飄飄倒吹什麼都做了.
什麼壞事都做了.
做到盡.
染上毒癮對身體造成很大的傷害.
他開始產生幻覺.
到了晚上睡不著.
耳邊傳來很可怕的幻聽.
有聲音說你死了好一點.
然後就在疑神疑鬼.
誰來找我?誰來殺我?.
他睡覺前手握著武士刀.
保護自己.
因為他有幻聽.
分不清楚什麼是真什麼是假.
只要外面有任何風吹草動.
就會拿著刀亂揮舞.
嚇壞他身邊的人.
但他心裡經常有個吶喊.
誰來救我?誰來救我?.
他很想得到拯救.
有一次他帶著刀開車到山上.
一直開車想起以前很多的事.
他竟然拿刀割自己的手腕.
鮮血不斷湧流.
後來他本來已經失去知覺.
但因為他的車擋住了後面的路.
後面有司機在按摩.
按摩的車喇叭吵醒了他.
結果他這次開車離開.
竟然又可以堅持下去.
他回想起以往的事.

$^{1001}$他覺得很唏噓.
他說自己曾經自殺三次.
但不知為何都被人救回來.
自己都不知道活著的意義.
怎樣活著去.
心裡一直在想誰來救我.
他很想去救.
他覺得沒有愛 生命沒有意義.
很絕望 不想活下去.
但後來很意外地.
他參與當地一個牙醫診所的裝修工程.
認識了一個基督徒的女子.
叫做比女子.
後來成為了他太太.
比女子在一次機會.
送了一本聖經給他.
他開始看聖經.
看聖經時給他一句說話.
很觸動他的生命.
若是有手叫你跌倒.
就撞下來丟掉.
寧可失去白體中的一體.
不叫全身下入地獄.
他望著他的手.
心裡很難過.
原來他的手曾經做了很多壞事.
偷錢 打架 用它注射毒品.
他一直想的時候一直哭.
不知過了多少個小時.
他感受到以往的污穢和眼淚.
一同被耶穌的愛沖洗了.
之後他開始很熱心讀聖經.
還要一個月內讀完整本聖經.
我們現在是不是很辛苦呢?.
讀經計劃.
讀著讀著不知去了哪裡.
但這個井上勳一個月讀完.
心裡很渴慕.
他沒有讀經計劃鼓勵他讀.
他自己就很想讀.

$^{1041}$讀完後有很多問題.
他跟基督徒的女孩子討論.
一路跟她分享.
一路去認識神.
後來被女子邀請去教會.
慢慢地他真的回到教會.
在教會裡成長.
甚至乎黑社會的大佬.
願意把他交給教會.
願意讓他脫離黑社會.
他想都沒有想過.
在日本要脫離黑社會.
小則沒有手指一撅.
大則沒有一條命.
他沒有想過黑社會的大佬.
會這樣主動讓他脫離.
讓他專心在教會裡生活.
他心裡覺得很震撼.
後來他受洗.
那天早上他去洗澡.
他哭著對耶穌說.
主耶穌 你真的要我嗎?.
我只剩下這個污穢的身體.
你真的要我嗎?.
接著有個小聲音出現.
是耶穌對他說.
我就是愛這樣的你.
你是聖潔的或污穢的都沒有關係.
我就是這樣愛你.
我愛你 我會保護你.
我會永遠和你在一起.
井上昆感到很感動.
耶穌愛我.
這麼污穢的我.
真是深深被耶穌觸動.
在洗禮當日.
井上昆看到的異象.
這個也不是每個人都有機會看到.
上帝讓他看到.
真是上帝向他表達.

$^{1081}$是何等的接納和愛他.
他看到主耶穌在天上.
在他的身邊有些孩童.
接著天使從井上昆身邊.
升上天.
很喜樂 很歡喜快樂地說.
出生了 出生了.
誰出生了?.
井上昆在神的國出生了.
在洗禮當日.
天上很大的歡呼聲.
弟兄姊妹我想說.
在你信耶穌.
在你受洗當日.
天上也有很大的歡呼聲.
天上也在歡呼.
出生了 出生了 是嬰兒.
不只是停留在那裡.
上帝 天君 天使為我們歡呼.
我們加入神的國度是何等美好.
不只是停留在那裡.
我們由嬰兒成長成為大人.
生命是180度的轉變.
蘇羅的人生轉變是否很大?.
逼迫別人.
由逼迫基督徒變成跟從耶穌基督.
甘願受逼迫.
受盡很多苦.
一生以得著耶穌基督為自保.
我和你呢?.
我和你想不想成為一個全新的人?.
信主這麼多年還在說這句.
其實這句說話很早很早很早前就應該說.
我和你要成為一個全新的人.
我們由嬰兒成長.
整個人要扭轉方向.
向著耶穌基督奔跑.
做所有事為耶穌基督.
今天大家問問你的生活.
究竟有多少是為耶穌基督?.

$^{1121}$你為誰而活?.
是否付少少代價.
二十多萬.
真是受得苦.
耶穌先給你一些好處.
這樣才行.
有什麼好處?.
你告訴我.
數幾十百有什麼好處?.
值不值得我這樣為你?.
耶穌為我們擺上他的命.
怎樣都值得.
只有我和你不配.
我們一同禱告.
在禱告的時候.
我邀請弟兄姊妹向上帝禱告.
今天的你.
是否做一個舊的你?.
是否信耶穌這麼多年都沒有改變?.
那些纏擾我們的罪.
是否還在?.
我們還是骯髒邋遢.
未信耶穌前.
為自己而活.
今天信耶穌後.
還是為自己而活.
我們和上帝斤斤計較.
還是求利益.
求好處.
是否如此?.
我想問.
我們為耶穌擺上了什麼?.
人們看到我們的生命.
是否看到主?.
請你自己也向上帝禱告.
天父上帝.
求你憐憫我們.
本來很多年前我們應該說.
一個全新的我們.
我們不再為自己而活.

$^{1161}$要為你而活.
但可能這麼多年.
我們還是一樣.
信主和未信主.
也不是很大分別.
所以身邊的人.
也不覺得我們有什麼.
很吸引他們的地方.
所以他們也未必很想認識你.
天父求你寬恕我們.
耶穌基督求你寬恕我們.
這麼輕悍.
你所為我們成就的救恩.
求你寬恕我們.
讓我們今天.
要重新起來.
悔改.
成長.
由天國裡的嬰兒.
慢慢要成長.
成為一個大人.
成為一個精兵.
主讓我們不要離開這裡之後.
就將這一切忘記得一乾二淨.
讓我們成為不一樣的人.
願意得著我們.
讓我們也得著你.
奉耶穌基督的名字祈禱.
阿們.
謝謝大家收看 再會.
\newpage



\section{馬太福音 6:19-34}
\label{sec:B5mDgas2LUk}
\textbf{2023.06.11 《今日的恩典》 馬太福音 6\_19-34 (和修版) 講員 : 曾裔貴牧師}
\newline
\newline
連結: \href{https://youtube.com/watch?v=B5mDgas2LUk}{\texttt{https://youtube.com/watch?v=B5mDgas2LUk}} ~~~~ 語音日期: 2023-06-11
\newline
\newline
\hyperref[sec:MMKQi99_fYg]{\small{< < < PREV SERMON < < <}}
~
\hyperref[sec:index_chronic]{\small{[返順時目]}}
~
\hyperref[sec:index_scriptual]{\small{[返順卷目]}}
~
\hyperref[sec:NXj4I_yTkE8]{\small{> > > NEXT SERMON > > >}}
\newline
\newline
馬太福音 6:19-34
\newline
\begin{longtable}{cl}
\hline
\hline
章節 & 經文 (和合本修訂版)\\
\hline
6:19 & \begin{tabularx}{0.7\textwidth}{X} 「不要為自己在地上積蓄財寶；地上有蟲子咬，能鏽壞，也有賊挖洞來偷。 \end{tabularx} \\ \\ \relax
6:20 & \begin{tabularx}{0.7\textwidth}{X} 要在天上積蓄財寶；天上沒有蟲子咬，不會鏽壞，也沒有賊挖洞來偷。 \end{tabularx} \\ \\ \relax
6:21 & \begin{tabularx}{0.7\textwidth}{X} 因為你的財寶在哪裡，你的心也在哪裡。」 \end{tabularx} \\ \\ \relax
6:22 & \begin{tabularx}{0.7\textwidth}{X} 「眼睛是身體的燈。你的眼睛若明亮，全身就光明； \end{tabularx} \\ \\ \relax
6:23 & \begin{tabularx}{0.7\textwidth}{X} 你的眼睛若昏花，全身就黑暗。你裡面的光若黑暗了，那黑暗是何等大呢！」 \end{tabularx} \\ \\ \relax
6:24 & \begin{tabularx}{0.7\textwidth}{X} 「一個人不能服侍兩個主；他不是恨這個愛那個，就是重這個輕那個。你們不能又服侍神，又服侍瑪門。」 \end{tabularx} \\ \\ \relax
6:25 & \begin{tabularx}{0.7\textwidth}{X} 「所以，我告訴你們，不要為你們的生命憂慮吃甚麼喝甚麼，或為你們的身體憂慮穿甚麼。生命不勝於飲食嗎？身體不勝於衣裳嗎？ \end{tabularx} \\ \\ \relax
6:26 & \begin{tabularx}{0.7\textwidth}{X} 你們看一看那天上的飛鳥，也不種也不收，也不在倉裡存糧，你們的天父尚且養活牠們。你們不比飛鳥貴重得多嗎？ \end{tabularx} \\ \\ \relax
6:27 & \begin{tabularx}{0.7\textwidth}{X} 你們哪一個能藉著憂慮使壽數多加一刻呢？ \end{tabularx} \\ \\ \relax
6:28 & \begin{tabularx}{0.7\textwidth}{X} 何必為衣裳憂慮呢？你們想一想野地裡的百合花是怎麼長起來的：它也不勞動也不紡線。 \end{tabularx} \\ \\ \relax
6:29 & \begin{tabularx}{0.7\textwidth}{X} 然而我告訴你們，就是所羅門極榮華的時候，他所穿戴的還不如這些花的一朵呢！ \end{tabularx} \\ \\ \relax
6:30 & \begin{tabularx}{0.7\textwidth}{X} 你們這小信的人哪！野地裡的草今天還在，明天就丟在爐裡，神還給它這樣的妝飾，何況你們呢？ \end{tabularx} \\ \\ \relax
6:31 & \begin{tabularx}{0.7\textwidth}{X} 所以，不要憂慮，說：『我們吃甚麼？喝甚麼？穿甚麼？』 \end{tabularx} \\ \\ \relax
6:32 & \begin{tabularx}{0.7\textwidth}{X} 這都是外邦人所求的。你們需要這一切東西，你們的天父都知道。 \end{tabularx} \\ \\ \relax
6:33 & \begin{tabularx}{0.7\textwidth}{X} 你們要先求神的國和他的義，這些東西都要加給你們了。 \end{tabularx} \\ \\ \relax
6:34 & \begin{tabularx}{0.7\textwidth}{X} 所以，不要為明天憂慮，因為明天自有明天的憂慮；一天的難處一天當就夠了。」 \end{tabularx} \\ \\
[1ex]
\hline
\hline
\end{longtable}
$^{1}$請姐妹們平安.
我們開始之前一起祈禱.
我們能夠來聆聽聖經的真理.
尤其是主耶穌在這一段教訓.
要我們先求神的國度和神的義.
求你幫助我們明白經文的含義.
也讓我們身體力行.
我們參與在神的國度的興建裡.
使主的國度早日降臨.
而我們是有份於其中的.
今天早上我們等候仰望你的恩典.
在我們中間奉耶穌基督的名.
阿們.
為什麼叫做今天的恩典這個題目呢?.
容祖兒一位很流行的歌手.
他有一首很出名的詩歌.
叫做《明日的恩典》.
取材自《奇異恩典》的故事題材.
中間一個重要的部分.
直接就是用了奇異恩典.
不知道大家有沒有聽過這首詩歌呢?.
相當流行的.
有一段時間.
今天其實這個詩歌.
是很想鼓勵聽這首歌的人.
要有盼望.
要對明天有信心.
今天的恩典.
就是透過這段聖經.
今天我們能夠由19節講到25節.
因為26到34節.
那裡是有另外一個題目.
未必夠時間講完.
在第六章.
有三件事主耶穌禁止.
跟隨祂的門徒.
要來避免的.
是關於什麼呢?.
施捨.
祈禱.

$^{41}$和禁食.
三件事都是猶太人的宗教行為.
為什麼要禁止呢?.
記住不是禁止他們不參與.
而是用什麼心態去做這些宗教行為.
留意第四,第六和第十八節.
關於施捨的時候.
心態要怎樣呢?.
要叫你施捨的事情.
行在暗中.
當你的天父.
在暗中來察看的時候.
必然會報答施捨的人.
意味著我們要做在神的裡面.
不是做給別人看.
更不是一種沽名釣譽.
第二段就是第六節.
關於禱告.
耶穌怎樣教導當時的猶太人祈禱呢?.
成為耶穌的門徒.
就不要學猶太人的拉比.
第六節.
你禱告的時候.
就要進你的內屋.
關上門.
禱告你在暗中的父.
當然不是天父躲在暗中.
而是天父無處不在.
無所不知.
天父會在暗中祈禱的人.
內心察看.
再一次出現這一節聖經.
你父在暗中察看.
必然就會報答祈禱的人.
第二段已經說了.
暗中的向天父做.
即是說所做的一切.
宗教的行為.
都是做在天父的心靈裡面.
要讓他.

$^{81}$稱他.
讓他將來賞賜這樣的人.
到第三段.
繼續是這樣說.
十八節.
不叫人看出你禁食來.
只叫你暗中的父看見.
你父在暗中察看.
必然報答.
那個暗中來禁食的人.
三個宗教行為.
包含了猶太人.
差不多是整個信仰的生活.
施捨.
對別人的慷慨.
祈禱.
為一些事情祈禱.
禁食.
有些重大的事情.
需要認真地在神面前尋求.
幾個宗教行為.
都強調一件事.
你要做在天父的眼裡.
你要做在他稱許的滿意裡面.
即是意味著什麼呢?.
我們很大部分的人.
當然包括我自己.
緊張人怎麼看我.
很少會緊張神怎麼看我.
很少會真是一心一意.
尋求天父喜悅.
就已經足夠了.
但耶穌說.
三段教導門徒.
當然就是當時跟隨他.
很想成為他的門徒的一個跟隨者.
耶穌就有這樣說的.
好像一個新的天國倫理觀.
這樣去勉勵門徒.
接著就是今天我們要說的.

$^{121}$十九到二十五節.
很特別用了三個比喻.
這三個比喻就告訴我們.
關於財寶.
關於我們生命的光.
關於你要效忠哪一個.
三個比喻都好像在在地重複.
剛才那些宗教行為.
心態要對準神.
現在就再告訴我們更正面的.
應該怎樣去很肯定.
對神國度的擺上.
和生活上的追求.
主耶穌第一樣想到的就是.
我們心靈裡面.
什麼是我們的財寶.
大家數一數.
你的小朋友小.
你就很想他入徐澤林小學.
或者是其他的名校.
上中學大學出來工作.
如果能夠入到一些比較大的機構公司.
你就好了.
這個時間是六月.
唐主任就很忙.
很多時候要寫推薦信.
推薦人入學.
就找我.
曾牧師我的小朋友要入黃塵.
幫我寫一封信.
他有沒有上兒校.
有的有的.
父母呢.
父母懂的.
洗禮的.
一般都會問這幾樣東西.
真的有穩定的上教會.
我寫的時候都不是故意作下去.
是真的有上開教會.
一路有教會事奉.

$^{161}$我們都會寫的.
這意味著我們都很緊張.
我們的小朋友的學業.
再大一點.
我們自己的財富.
愛情.
我們緊張的工作.
我們人怎樣看我們內心的尊嚴.
我們自己的形象.
有時會不順眼.
怎樣都不肯認錯.
我們都會是這樣的人.
再大一點.
可能我們很想別人讚賞我們.
再實在一點.
我們的財富可能就是現實的.
那些現金.
銀行的戶口.
那是安全感.
試想想你沒有了安全感.
而你仍然有神.
你是不是真的會很安心呢.
可能未必.
所以主耶穌在21節說了很重要.
很透視到人性裡面.
真正的問題所在.
第21節不如大家一起讀一讀.
因為你的財寶在那裡.
你的心也在那裡.
所以是心之所屬的問題.
所以不是財寶本身.
是我們對財寶執著的依戀.
那個手拿著不放手.
那件事令我們對天國.
放在一個很次要的次序.
是這個問題所在.
例如保羅怎樣將這段經文.
再一個很清楚的演繹呢.
提摩太前書第六章第十節.
但那些想要發財的人.

$^{201}$就有禍了.
他們陷在那些無知有害的迷惑裡面.
為甚麼呢.
錢財很重要的.
沒有錢財你生活也不行.
原來錢財是萬惡的嗎.
基督教是這樣的嗎.
你看看聖經.
保羅怎樣說.
他說貪財是萬惡之根.
不是錢財是萬惡之根.
是心底裡面的貪念.
所以主耶穌在這裡告訴我們.
你的財寶在那裡.
你的心就在那裡.
即是說其實你教養兒童.
你的小朋友是沒有問題的.
你追求你的事業有成.
是沒有問題的.
你追求得到盡心帶來的肯定.
是沒有問題的.
你追求的有一個安全感是沒有問題的.
不過原來最終.
耶穌就很透析地去指出.
對當時的門徒一個分面.
天國才是永恆.
天國的投資才是永遠不扭曲.
我們是為天國而生而效力.
所以你仍然可以去教養你的子女.
為他們預備最好的.
在座的工人可能已經全都離開了.
沒有年輕時那麼辛苦了.
剛才我在宣德也問.
那些人也有反應.
因為我見到很多都和我一起成長.
子女都長大了.
二十多歲現在事奉了.
真的有些事情都可以沒有以前那麼緊.
有些事情都可以慢慢放手了.
現在你追求什麼呢?.

$^{241}$耶穌問這個問題.
你為什麼而追求呢?.
耶穌問這個問題.
所以很多人只是為了.
可能他就要拼命吃最好的.
有個弟兄.
他爬到很高的山.
為了拍一張很美的照片.
但是.
是不是他的終極呢?.
他是不是人生只是為了拍很美的照片呢?.
有沒有一個為神去做的天國的價值觀呢?.
我們的事奉.
我們想尋求的是人對我們的肯定.
還是神對我們的稱許呢?.
這是很重要的內心價值.
所以耶穌一針見血.
因為你的財寶在哪裡.
你的心就在哪裡.
這是第一方面耶穌的比喻.
叫我們要小心.
我們的心如何擺放.
所以為什麼保羅也有這樣的教導.
當求在上面的事.
不要求地上的事.
不是叫我們不惜吃人間煙火.
而是我們的心注目永恆.
你想想.
如果將來永恆.
是有五百億億億億萬年.
我們都數不到了.
跟我們在地球上的一百年.
一百二十年.
你想想有多大的距離.
但往往我們人忽視了.
真是永恆的永恆.
那個永存的價值.
我們很少人真的會對永恆.
有每天的執著.
每天都想想.

$^{281}$我應該如何為神的國度.
我們去投資呢?.
耶穌說了.
要積蓄財寶在天上.
就是將我們的心力.
整個人投放在神的國度.
第二個比喻.
光明.
眼睛就是身上的燈.
你的眼睛瞭亮.
全身就光明了.
但如果你的眼睛昏花.
全身就黑暗.
而且這個黑暗何等大.
意味著如果我們的生命暗淡了.
價值觀模糊了.
我們就對永恆.
失去那種感覺和敏銳.
提到光明.
用光明.
好像身體是光明的.
驅使我們每天很想為神服侍.
我兒子很多時候回來.
他有什麼工作,家庭.
特別是工作.
會和我分享.
他說很特別.
最近他有一個祈禱.
他很想神給他一些新的事物可以經歷.
他突然想到他的部門.
他現在做了半年.
但轉變得很快.
他現在變成近乎老臣子.
他很想幫新來的實習生.
一些研究助理RA的人.
他很想讓他們融入工作.
新來很快就約團隊一起吃午飯.
給他們一個空間和接納.
好像擁抱一樣.
令團隊很開心.

$^{321}$我兒子有這樣的想法.
神給他一個機會.
有一個新同事.
很特別.
突然間他有一個朋友.
他在以前工作的朋友.
有來往的同事.
是個女生.
突然打給我兒子.
問他是不是有個新來的女生.
他說是呀.
你怎麼知道?.
說我的名字.
原來這個新同事.
我兒子很想讓這班新同事融入群體.
原來這個新同事很特別.
原來是我兒子的女性朋友.
你是女朋友.
女性朋友的中學同班同學.
還要是坐在旁邊的同學.
神奇地把這個人放在我兒子身邊.
是他的下屬.
來幫助他.
這是我兒子在這段時間的祈禱.
很希望有些新的東西可以服侍神.
這些會不會是我們大家已經模糊了.
沒有這樣的學習呢?.
你有沒有想過.
你走出去洪水橋.
你在輕鐵站.
你在亞山園.
你吃東西.
你能不能夠表達出一份善意.
你是代表著耶穌.
代表著神.
來讓人看到天國是很實在的.
在你的面容.
在你的具體行徑.
去反映出神在你的生命裡面是有價值的.
你是著緊神的國度的.

$^{361}$能不能夠我們大家一起學習呢?.
這就是生命的光.
你有沒有已經暗淡了呢?.
別說基督徒了.
那些熟悉你的人.
這個人經常說髒話的.
一離開教會就說髒話的.
那你怎樣是生命的光呢?.
大家想想.
第三個比喻再清楚一點.
一個人不可以事奉兩個主.
當時是奴隸制度.
是主人和奴隸的奴僕.
這樣的比喻.
不是愛這個就恨那個.
我們不能夠又事奉神.
又事奉馬門.
即是一心不能夠義奉義用.
和兩個主人.
就更加斬釘截鐵了.
三個比喻就是一個結論.
我們的內心.
怎樣為天國擺上呢?.
各位啊.
你會不會有很多的藉口.
哎呀.
我的子女又小.
房子又未供完.
工作又辛苦.
現在健康開始都慢慢走下坡了.
哪裡還有夢啊?.
哪裡還有事奉啊?.
不要想了.
我們是不是已經慢慢沒有了這個意象.
這個追夢的小朋友的時候.
初信就是追夢.
現在追什麼鬼啊?.
追巴士都追不了了.
我們是不是已經沒有了這種夢想呢?.
以為只是少年人的專利呢?.

$^{401}$不是的.
很快.
八月.
就是我人生一個很重要的階段.
我很想有一份生日禮物.
很想.
我事奉了三十多年了.
由我信主到現在.
好像都有點累.
但最近這段聖經提醒我.
先不要這麼快.
神的國度我們要繼續去追求.
可能那個節奏未必是這麼快.
但是那個心態不可以改變.
是一個很好的提醒.
我想用一個例子.
來彼此互勉.
有一個女孩.
今年26歲.
她是新疆一個烏魯木齊的人.
在北京讀大學.
讀完了.
如是者.
回去新疆.
有一間很大的公營機構.
做很好的工作.
的待遇和環境.
做了幾年之後.
她就覺得跟她的理想相差很遠.
她毅然.
其實上司都勸她.
你現在不要走.
很穩定的工作.
但她不理.
離開新疆.
回到北京.
在北京再找工作.
如是者.
做了幾份工作之後.
內心有個聲音.

$^{441}$都是覺得.
這些工作暫時不適合她.
又是這樣辭職.
有一天.
其實就是上個月的5月13日.
她在小紅書.
登了一個廣告.
她說她要出租自己.
一元.
出租自己人民幣一元.
她說要找到一百個人.
作為她六月生日的時候的禮物.
一百個人.
她說只要你有一個願望.
不是非法.
不是無理.
她說她都可以幫她實現.
一出了這個廣告之後.
很多人問.
真的很想了解一下.
原來有很多人有些願望.
有些心願很想達成.
其中講幾個例子.
有一個女孩.
很自卑.
對自己的身材很自卑.
她就問這個.
她叫做Peggy.
這個女孩.
26歲.
如果我用一元.
請你跟我一起.
這個女孩就是.
很自卑的女士.
她說她很想有一個心願.
就是在夏天.
穿吊帶衣服逛街.
Peggy就想一想.
這件事其實沒有什麼.
很難的.

$^{481}$一個人自己有自主權.
去穿衣服.
又不是很暴露.
為何不可以呢.
就幫她達成.
一元.
有人問為何要收一元.
一元即是每個人都付得起.
但不是代表免費.
你要付代價.
原來在這個過程.
令到那個女孩.
有勇氣面對其他的事情.
跳出自己最艱難的第一步.
來得到信心.
往後的生活.
對Peggy自己來說.
雖然她賺了一元.
但原來她自己.
體驗了不同的人生.
又可以幫到其他人.
然後另一個男孩.
他現在仍在服兵役.
他當然不是北京人.
他說我很想.
有人陪我.
參觀北京的城.
我能夠有一個.
深度的旅行.
Peggy說我本身在北京讀書.
我當然熟悉.
一元陪士兵.
當兵的.
去行北京的城.
如是者.
雖然現在.
差不多六月了.
都未夠一百個人.
但她已經幫很多很多人.
實現了那個理想.

$^{521}$那個夢想.
尤其是將那個勇氣.
一元就可以買到.
那個勇氣去做一些事情.
在我們來說.
好像天馬行空一樣.
很匪夷所思.
出租自己.
去陪一些人做一些事情.
其實是幫人達成一些願望.
親愛的弟兄姊妹.
你回教會.
你回教會是否已經沒有了你的夢想.
你是不是已經沒有了.
你為主耶穌.
為神的國度.
那個參與.
那個服事.
會不會我們再一次想一想.
無價的寶.
耶穌基督.
已經為你為我.
付出了.
祂現在邀請我們.
去到33節.
你們要先求他的國.
祂的兒.
這些東西.
都要加給你們.
去到結尾.
都呼籲門徒.
要來為神的國.
神的兒.
這樣來效力.
大家要問一問.
有什麼財寶.
是你今天都不肯放下.
有什麼生命.
來得重要的.
我們耶穌基督所說的.

$^{561}$永恆國度的財寶.
如果我們不在今天.
認真地再一次去檢視我們的內心.
而去到白白地.
糟蹋了神給我們的時間.
我們就會一直隨著我們的年日老去.
我們就會失去了對神國度的敏銳.
更重要的是.
當我們沒有了夢想.
證明我們的心.
有更多地上的財寶.
霸佔了我們的生命.
每個人都有自己心所屬的財寶.
但耶穌說.
傳到永恆才是真正的財寶.
地上的.
無論多美的相.
你的兒女.
包括你和我的健康.
包括我們的感情.
甚至物質的流域.
金錢.
我們一離開這個世界.
是什麼都帶不走的.
真的什麼都帶不走的.
我們從怎樣離開這個世界.
是完全一樣的.
是什麼都沒有.
這樣離.
這樣走.
所以耶穌說.
正正因為這樣.
我們就要明白.
永恆才是真正我們可以把握的.
很盼望今早這一段聖經.
今天的恩典.
不要等到明天.
今天.
錦宣家.
洪恩家.

$^{601}$剛才我是孫德嘉.
我們大家一起.
為神的國度.
接受耶穌的邀請.
讓我們能夠有回這份初心.
少年夢.
我們一起禱告.
多謝主.
我們在你的面前.
真的需要謙卑.
苦服.
盼望主在我們生命裡面.
有你的帶領.
感動.
讓我們一起為神的國度.
去效力.
有人有不同的恩賜.
但我們很盼望主.
真的在我們大家弟兄姊妹的內心.
有的是對主的愛.
有聖靈在我們大家的心靈裡面.
那個感動和引導.
更加有我們牧者.
在我們中間.
和我們做一個恩賜的配對.
以致我們在服事上.
更加的精準.
更加是適合主的心意.
我們這樣仰望感恩禱告.
奉耶穌基督的名.
阿們.
\newpage



\section{路得記 4:13-17}
\label{sec:NXj4I_yTkE8}
\textbf{2023.06.18 《犧牲的愛；愛的犧牲》 路得記 4\_13-17 (和修版) 講員 : 梁子勇牧師}
\newline
\newline
連結: \href{https://youtube.com/watch?v=NXj4I-yTkE8}{\texttt{https://youtube.com/watch?v=NXj4I-yTkE8}} ~~~~ 語音日期: 2023-06-18
\newline
\newline
\hyperref[sec:B5mDgas2LUk]{\small{< < < PREV SERMON < < <}}
~
\hyperref[sec:index_chronic]{\small{[返順時目]}}
~
\hyperref[sec:index_scriptual]{\small{[返順卷目]}}
~
\hyperref[sec:3dWTg4gpwpQ]{\small{> > > NEXT SERMON > > >}}
\newline
\newline
路得記 4:13-17
\newline
\begin{longtable}{cl}
\hline
\hline
章節 & 經文 (和合本修訂版)\\
\hline
4:13 & \begin{tabularx}{0.7\textwidth}{X} 於是，波阿斯娶了路得為妻，與她同房。耶和華使她懷孕生了一個兒子。 \end{tabularx} \\ \\ \relax
4:14 & \begin{tabularx}{0.7\textwidth}{X} 婦女們對拿娥米說：「耶和華是應當稱頌的！因為他今日沒有使你斷絕可以贖產業的至親。願這孩子在以色列中得名聲。 \end{tabularx} \\ \\ \relax
4:15 & \begin{tabularx}{0.7\textwidth}{X} 他必振奮你的精神，奉養你的晚年，因為他是愛慕你的媳婦所生的。有這樣的媳婦，比有七個兒子更好！」 \end{tabularx} \\ \\ \relax
4:16 & \begin{tabularx}{0.7\textwidth}{X} 拿娥米接過孩子來，抱在懷中撫養他。 \end{tabularx} \\ \\ \relax
4:17 & \begin{tabularx}{0.7\textwidth}{X} 鄰居的婦人給孩子起名，說：「拿娥米得了一個孩子了！」她們就給他起名叫俄備得。俄備得是耶西的父親，是大衛的祖父。 \end{tabularx} \\ \\
[1ex]
\hline
\hline
\end{longtable}
$^{1}$洪恩堂的弟兄姊妹大家好.
很難得和大家相隔一年左右.
我們再一次在這裡去敬拜.
上一次我和大家分享也是《路得記》的第一章.
今天我和大家看.
跳過中間的經文去聚會這段神話.
和大家在敬拜的過程中.
我覺得今天我們去看《路得記》.
也很和應洪恩堂的境況.
老人家也有,小朋友也有.
甚至今天我們看程序表才驚覺.
原來我們老牧師的身體.
去到一個這樣的情況.
看完《路得記》我們更加思想到上帝的作為.
以致我們心裡面正在燃點信心和盼望.
就如在敬拜的過程中.
你們所選取的詩歌.
以致從你們的表情流露.
我都看到上帝的恩典在洪恩堂裡面.
盼望不僅獲得鼓勵.
上帝也使用眾多弟兄姊妹在這個區份.
在上帝託付給你們這個地域裡面.
去見證,去高舉上帝.
我們一起先去祈禱.
是你光照我們.
以致我們能夠從無知,哀慼,埋怨裡面.
體會你的恩情.
遠遠深厚過我們生命裡面所歷練的.
從永遠到永遠.
你是愛.
願你光照,甦醒,引導我們.
讓我們跳出自己的眼界.
我們更加認識.
神理無邊的愛.
每一天哪怕當下.
臨到我們中間.
我們恭敬將我們個人來到教堂.
唯有你認識我們裡面.
哪怕幽暗,閉塞的地方.
求你拯救我們.

$^{41}$成為我們的幫助.
讓我們與你來到相近.
讓我們更加體會.
主啊!天下人間.
沒有別一位能夠更貼近我們的心.
願你成為我們的幫助.
誰聽我們和心中的交託和仰望.
靠耶穌基督聖靈祈求.
阿們.
剛才大家讀了後面那一段.
這是第一章的經文.
在圖片上可以看到.
那裡是描述.
這裡是拿訛米和她的丈夫離開伯利恆這個地方.
去到摩押地.
跟丈夫和兩個兒子.
兩個兒子娶了媳婦.
但是好景不常.
三個男人離開世界.
到頭來這個寡婦.
帶著兩個媳婦.
結果一個媳婦留在摩押地.
這個外邦人的地方.
另外一個媳婦叫路德.
跟著這個寡婦.
其實寡婦是沒有子女的.
是沒有子女的.
只剩下媳婦.
那種淒涼真是難以想像.
不單是今天遇到這樣的事.
當日的時候還淒涼.
沒有交通津貼.
沒有老人金.
也沒有水果金.
你知不知道?.
當拿訛米和媳婦回到伯利恆的地方的時候.
直接說,當乞丐.
也沒有拿訛米的份.
你太老了.
媳婦幫他去行乞.

$^{81}$你明白多淒涼嗎?.
後來的日子又怎樣過?.
按早不可晚.
沒有病患就好了.
沒有病患也不知道好壞.
要看你能撐多久.
有病患就快快快.
簡單簡單直接.
你接我回去吧.
就像李毓伯一樣簡單.
快點吧.
弟兄姊妹.
這是第一章的情況.
今天我們看到最後一章.
我們聽故事很喜歡看大結局.
最終結局如何?.
結果是.
我們看就以為很開心.
總算好了.
看第二章的影片.
到了第四章的時間.
我們以為都好了.
原來媳婦改嫁了有錢人家.
而且不單這樣.
生了孩子.
生孩子不僅是照顧孫子這麼簡單.
孫子就變成跟隨父親.
然後承受的遺產.
可以成為父親那一面.
好不好?.
我不知道你是撒謊還是好.
這段經文留意.
男女都沒有發聲.
後面周邊的人.
不論在拍手掌.
有這個媳婦也好.
有孫子也好.
怎樣都好.
一句話都沒有說.
那不是啊.

$^{121}$經文裡面說到後面.
到了第十七節.
論社的婦人說.
那不是得孩子了.
就給孩子.
起名叫俄比德.
最後一句.
這俄比德就是耶西的父.
耶西是大衛的父.
你和我知道大衛是誰.
還黃呢.
那是不是知道.
死了都不知道.
是不是.
死了都不知道.
還要隔代.
其實奴隸自己都不知道.
俄比德都不清楚.
耶西都未必知道.
根本不知道.
是不是.
那是什麼.
這輩子你都見不到.
你明不明白為何經文裡沒有說納和米.
好了.
將來我的後代可以出人頭地.
光宗耀祖.
很開心.
誰會說這句.
沒有人知道.
那奴隸記想說什麼呢.
奴隸記是不是想告訴你.
原來守得雲開見月明.
我終於抱到孫子.
我現在都沒有孫子抱.
我有三個女兒.
結了婚.
我不知道你期盼著什麼.
在你的人生上望著什麼.
其實.

$^{161}$就算在納和米當代.
這個所謂的俄比德.
他做皇帝.
那又如何.
你記得大衛王當年如何.
他家人一起走難.
他的子女如何.
你以為做到高處就不用面對妻離子散嗎.
其實我們都是同一樣.
在人生中遇到各樣遭遇的人.
甜酸苦辣.
你和我真的沒有經過嗎.
說去也好.
你和我都還慶幸.
前兩星期鑽石山的事有沒有嚇到你們.
有沒有看到那些錄影.
哇.
不值得看.
很可怕.
兩個女兒在街上走路.
原來不認識.
如果你在現場怎麼辦.
藍都擋不了.
沒有工具擋住.
擋一下擋不了第二下.
這幾天有沒有看新聞.
知道不是香港.
是德國發生的事.
我中間聽了未必知道.
兩個都是二十歲左右的女孩去德國旅行.
有一個美國的男孩.
說帶她去一個小路看景點更加精彩.
然後想有不懷好意.
要性侵他們.
先後把兩個女孩推了落五十米的懸崖.
我所知有一個離世.
有一個好像是重傷.
我沒有繼續跟著新聞.
是亞裔的美籍亞裔人.
是剛剛大學畢業的.

$^{201}$是五月畢業的.
我去看兩個新聞.
很相似.
人世間就是這樣發生的事.
如果那是你的子女.
你的心怎樣感受.
大學畢業.
正是青春少艾的時間.
將來的日子應該會怎樣.
為什麼會突然發生這樣的事.
我們好像不太明白.
不太理解得到.
其實當我們看路德基的時候.
我們去看納米的修明.
無論將來怎樣.
你知不知道那段日子很淒涼.
一個一個男離開.
你知不知道我怎樣從摩亞地回來伯利恆這個地方.
我面前的指日會怎樣.
就算在波亞斯的隱蔽下.
可能我內心都酸酸的.
你明白我說什麼嗎.
其實我都是寄人籬下.
我做黃就很不同了.
我就威風.
不過我是做一個老弱人家.
還是一個婦女.
只是婆婆.
我的人生是什麼.
弟兄姊妹.
求神幫助我們看到.
剛才我說的很有趣.
納米沒有說話.
但如果你看回撒母耳奇的另一位哈娜.
她也是得寺.
很不同.
她怎樣回應.
她發出禱告發出感恩.
還讚哈娜的祈禱很出名.
還將撒母耳奉獻給上帝.

$^{241}$那時候只有一個兒子.
接著有沒有生育還不知道.
這就是哈娜.
弟兄姊妹.
讓你我能夠多從路得記去看.
其實經文裡想告訴我們的是什麼.
小心啊 弟兄姊妹.
我們在整個成長的年日.
接觸聖經.
我常常以為聖經是聽故事的.
不是.
也不是告訴我們只是一個結局.
在整個過程裡.
在路得記結連著.
知道是哪一本書嗎.
是《時詩記》.
兩本書差不多是並排來看.
《時詩記》多說的是男人.
除了蒂波拉.
基本上都是說男人.
路得記說的是女人.
收到吧.
看聖經很重要的是.
不可以單單看一段.
你要整體去看.
就算不能夠整本書去掌握.
也要認識上一段是說什麼.
在路得記一章之前的一節.
就是《時詩記》最後一節.
我們一起看看.
那時以色列中沒有王.
各人任意而行.
路得記說得很清楚.
遲些孫子的孫子就做王.
大衛.
你等著就是肉眼見到的王嗎.
你搞錯了.
你沒有當上帝一直在中間.
就是在做王.
上帝好像管不了你的生活.

$^{281}$上帝好像介入不了你的家庭.
你的心.
以至你每天所思所想.
這才是聖經中所要表達和強調的.
所以其實在中間的《時詩記》是怎樣呢.
每個人都想做 只想做英雄.
誰知是什麼呢.
都是雄 都是狗雄.
每個人都像今天.
我們上網都說住紅宅.
誰知到頭來住了tong 房.
我一生真是淒涼.
真的有時tong 房都會有我的弟兄姊妹.
怎樣回望她一生呢.
我繼續住tong 房.
有時真是酸溜溜.
將來睡的地方和這個地方差不多.
多人四塊半.
是很酸溜溜的.
不過是否這樣看呢.
哪怕我們中間有上市的主席.
上市公司的主席.
又或者你真的退休.
你的人生是怎樣呢.
我們中間的姐妹.
講完《時詩記》講《路得記》.
你是貌美如花 好似路德.
無論嫁完再嫁 都可以嫁個好人家.
還是卡實卡實 你都是只會鬧米.
聽得明白我說甚麼嗎.
你叫坑吧.
你等人幫你吧.
有時我們都這樣看自己前面的人生.
哎呀 弊家父 將來怎麼辦.
我不是很想去老人院.
突然之間在我們心裡沒有那麼擔心.
突然之間不要這樣去看.
彷彿在整個《路得記》裡面神好像是真物.
他好像沒有理會.
他彷彿在等候人去作供.

$^{321}$無論在《時詩記》和《路得記》裡面.
戰亂也好 其實《時詩記》有很多戰亂.
饑荒也好.
劈頭《路得記》第一章第一節.
就已經講到他們的饑荒在中間.
又或者是腐敗也好.
這個社會怎樣腐敗也好.
但你知不知道關鍵的一件事是甚麼.
整個《時詩記》《路得記》.
它想告訴你.
是英海從未改變過.
是從未改變過.
整個《路得記》裡面很關鍵的.
我寧願跟大家說.
「伯利恆」這三個字.
原本的解釋是甚麼.
糧倉.
收到嗎.
用「糧倉」整個《路得記》也很有意思.
糧倉也飢荒的.
有沒有搞錯.
你信耶穌不是說有豐盛的生命嗎.
有沒有搞錯.
為甚麼搞成這樣.
有些時候我們信了耶穌一段日子.
我們心裡也會這樣問.
這樣就叫豐盛嗎.
我們心裡也好像迷失了.
我們困惑了.
我們是否正在經歷一個真實的信仰.
還是我們眼裡經常注目的是甚麼.
神為甚麼不給我們那些東西.
為甚麼要拿走我們的東西.
拿走我們的健康.
拿走我們的家人.
我們眼睛只看這些.
你還給我.
知道嗎.
你還給我雪糕.
你還給我.

$^{361}$哭了兩個小時也不夠.
但有些時候我們成人可能哭了幾十年.
為甚麼你不給我這些東西.
你又說聽我祈禱.
你又說愛我.
你又說是糧倉.
我餓死了.
路德基是在講一個這樣的場景.
弟兄姊妹.
其實.
當這樣看的時候.
有沒有看到其實神一直都是公認.
一直都是忍不住.
經文裡面描述.
撈米離開的時候是甚麼呢.
如果你有看.
路德基.
他自己說.
我滿滿地出去.
我的名字叫甜美.
小甜點.
不過你不要叫我小甜點.
我苦澀到難以想像.
當這樣說的時候.
是滿滿地出去.
那是在表達甚麼事.
在表達甚麼事.
你蔑視興萬gang .
你知不知道士師記和路得記在說甚麼.
上文說士師記在上的時間是因為出埃及.
經過曠野就要進去迦南.
你要進去迦南地.
我答應你要進去.
他們調頭離開.
我叫你進去伯利恆.
你滿滿地.
你為甚麼不留下來.
你為甚麼要走.
不單止這樣.
你說是饑荒.

$^{401}$你滿滿地出去.
你有沒有關心你的族人.
聽得明白嗎.
你以為只有你挨這個面前的日子嗎.
整個以色列民眾多的人留在伯利恆.
你想怎樣.
你做過甚麼榜樣.
你豈不是應該關心當地的人嗎.
所以路德基你有沒有留意.
波亞斯就是這樣.
所以波亞斯不是一個扮演出來的人物.
為甚麼你當年不是像波亞斯那樣.
兄弟姐妹.
其實你是蔑視興邁這位上帝.
內心裡面帶著不信不服的勃逆敵黨神.
我叫你在這裡.
你就走.
竟然全家千瘟萬瘟.
不聽神的吩咐和應許.
這個才是真正的悲苦.
是你的選擇.
其實哪一個真正犧牲.
哪一個真正遭受遺棄.
其實不是神遺棄人.
是撈米.
是撈米遺棄上帝和祂的族人.
他還未醒覺.
竟然當聽到這裡.
我想問你和紅恩堂有甚麼關係.
我想和大家看另外一節經文.
是主耶穌講的.
這節經文我們一起讀.
準備請.
耶穌正說話的時候.
眾人中間有一個女人大聲說.
說你胎的和乳養你的有福了.
耶穌說:是.
卻還不如聽神之道而遵守的人有福.
眾多姐妹誰想做瑪利亞.
主耶穌的母親.

$^{441}$你可要想清楚.
聖經裡面的心如刀割.
不單止這樣.
那三年不在瑪利亞身邊.
眾多姐妹.
這女人說:如果做你媽媽就好了.
很興奮很開心.
你都不知道.
不知道甚麼.
你不知道你自己有多幸福.
我有多幸福.
你給我甚麼福份.
我有甚麼值得高興.
我是女的也不是男.
整個男人社會都拿了權利利益.
有甚麼好.
看清楚下面那句.
這兩節經文我相信大家都看過.
但我們不是很認真去看.
其實每一節經文都很豐富.
你要花很多時間.
不斷去咀嚼去反思.
然後你才發現.
原來眼目看錯了方向.
我們經常看的東西.
原來不是聖經帶我們去看的東西.
難怪你越走越錯.
經文叫我們看甚麼.
他說:卻還不如聽神之道.
而遵守的人有福.
這不是我自己說的.
是主耶穌自己說的.
他愛你比愛媽媽還要愛.
聽到嗎.
你聽明白了嗎.
你不要在那裡哀嚎.
他愛你比愛媽媽還要愛.
他說你很蒙福.
我們家中誰沒有聖經.
你只有一半成功.

$^{481}$不是說笑.
以前的人沒有聖經.
回到會堂.
你揭開了嗎.
我都看不到.
你讀給我聽吧.
是這樣的.
你有聖經.
你猜你是哪裡.
我們家中誰沒有聖經.
你說:牧師,我不懂得字.
你會聽聖經的.
你知不知道.
有人讀聖經給你聽.
受訓有些程式可以幫你讀聖經.
弟兄姊妹.
你答了成功的一半.
求神幫助我們.
我們不要遺棄聖經.
容許我這樣說.
特別我們一班年長的弟兄姊妹.
是要你自己看聖經.
我負責一個長青團體.
他們抄聖經.
抄整本新書.
抄完.
可以用毛筆抄.
你不要太高難度.
用紙寫.
你知道很有幫助嗎.
可以用書法靈修.
抄聖經靈修.
重點是什麼.
原來神的話是不斷可以咀嚼的.
其實對魯德來說.
他不看重進行聖經.
以至神說了要入家門.
他不看重上教.
不看重建立神的家.
只顧自己.

$^{521}$他其實是不絕得罪神.
惹神憤怒.
糟蹋了我們的恩情.
繼續浪費那種英雄.
所以一班弟兄姊妹.
求神幫助我們.
不要再為我們的際遇.
為我們的家庭.
為周遭對待我們的人.
我們在發怨言.
我們繼續在這裡.
艾恨在這個世界上.
這節經文告訴我們.
你是蒙大恩的女子.
至於對誰說的.
瑪利亞說的.
但適用於我們中間每一位弟兄姊妹.
求神幫助我們.
其實我以後想說三點.
簡單提一提.
其實看什麼呢.
不像在羅馬看著嬰兒.
這是最後的結局.
你看看這個是誰.
不是很醜怪的出世嬰兒.
有人這樣說.
一定要很醜怪.
這樣的嬰兒.
是什麼呢.
看不到他的將來.
看什麼呢.
不過弟兄姊妹.
看看你的生命.
你的生命就已經在表明神的恩典和痛在.
你是被神所做的.
好像今天敬拜的那段.
不斷地逐場表達了.
回歸萬有的創造主.
高山.
神在這裡.

$^{561}$低谷.
神在這裡.
我應該和你們說過.
地球有多大.
太大了.
沒那麼大.
其實神在哪裡.
一直都在這裡.
你在哪裡他都知道.
所以弟兄姊妹.
我和一班長青說.
你在廚房有什麼事.
上帝在哪裡.
你一個人獨居的時間.
你在浴室發生什麼事.
誰知道.
神知道.
所以弟兄姊妹.
如果有一天你和我去到靈養院.
你知道什麼是靈養院嗎.
不要擔心沒有人探望你.
主耶穌在這裡.
所以求主.
堅固你和我.
有一天我和你去護老院.
神在這裡.
拿走心裡面的恐懼和疑慮.
不是帶著哀慨面對每一天.
以致.
這個創造主就是跨越死亡.
為什麼你還有一條命.
是因為那位勝過死亡的上帝.
所以弟兄姊妹.
我告訴我自己.
我卸任了.
所以大家在情權本沒有說.
我是屬於哪一間堂會.
不過我去看什麼.
我是黃金十年.
我說事奉了三十年.

$^{601}$不過我逐十年逐十年看.
不要太貪心.
黃金的十年.
真的.
弟兄姊妹.
你和我人生有這麼多歷練.
有這麼多豐富.
以致經歷這位上帝.
擺在你面前的日子.
上帝是一定會使用你的.
紅印堂的眾弟兄姊妹.
我看見很多年長.
在你的表達裡面.
我看到.
你的信心是堅實的.
請你活出來.
我想說一句話.
傳道書裡面有一句話.
不要說.
先前的日子.
強過如今的日子.
是什麼緣故呢.
聽得明白嗎.
不要說.
以前有.
現在就不行了.
以前我跑山.
沒問題.
現在走樓梯都不行.
你說得對.
不過跟著經文說什麼.
你這樣問.
不是出於自慰.
你夠屎就好了.
其實你說什麼.
你知不知道.
你閉家伙.
你一定走下坡了.
我經常叫長青的弟兄姊妹.
不要說這兩個字.

$^{641}$不過我心裡經常說.
說什麼.
但死.
不是.
將來的日子是更加好的.
老牧師說得對嗎.
將來的日子是更加好的.
黃金十年.
黃金三年.
黃金三個月.
為什麼.
什麼是自慰.
上帝在.
是上帝在.
他昔日能夠這樣做.
將來更加可以這樣做.
是上帝的作為.
我剛才讀的經文.
在《傳道師》第七章第十節.
你經常說以前好.
即是上帝在哪裡.
死了.
他不再做事了.
全靠你了.
求神幫助我們看到.
看一條命.
看一看你的命.
第二方面.
看嬰兒.
這個就是嬰兒.
這個嬰兒就是嬰兒.
不是他長大了.
不是他.
終於他可以.
回本了.
不是的.
弟兄姊妹.
不是這樣去看的.
不要緊.
人家對我們忘恩負義.

$^{681}$人家不喝水思源.
不要酸溜溜.
不要在這裡叫哼.
不要在這裡說.
我生錯了你.
不好了.
弟兄姊妹.
不要這樣去看.
為什麼這裡提伯尼痕.
是不是伯尼痕.
今天很重要.
我們不如去伯尼痕.
去種植.
梁滄嗎.
你搞錯了.
耶路撒冷呢.
不如去耶路撒冷.
不去也可以.
他在說的.
不是世上地域和國度的選民.
他叫你和我.
回歸萬民神的國度裡面.
所以你看看以色列.
整個歷史.
強盛的時間.
神在哪裡.
在這裡.
被擄的時候.
神在哪裡.
神都在這裡.
不過他的子民聽不聽話.
你信不信我是英雄.
弟兄姊妹.
你知不知道.
今天人世間裡面.
最重要的是什麼.
聖經.
是聖經.
一切的東西都會過去.
你身上以至所擁有的一切東西.

$^{721}$都不能夠跨世代.
錶給我兒子.
給我孫子.
地業.
留灣給子女.
什麼最好.
聖經最好.
不是因為保存.
是因為上帝說.
我乙族要成就這些東西.
為了應許.
為了聖經.
我先建立你有信心的.
弟兄姊妹.
不是因為我們一起參加師班.
一起在這裡上上有全契.
不是.
在我們周邊的人.
一個一個會離開的.
誰陪伴你最好.
不是說我們又少了一個.
你丟棄了神的話語.
是你親自看重.
神的話語.
以至去到那個地步.
晝夜思常.
我有個老朋友最近病了.
他拿出一堆藥.
哇.
紅橙黃綠青藍紫.
四種顏色.
真的.
我們長者.
可能一天.
要有一個藥盒.
等著也不記得吃哪種.
你等漏了一樣.
聖經.
不要只吃那些補充品.
不是說那些補充劑嗎.

$^{761}$聖經才是最重要.
甚至我們這班弟兄姊妹.
其實不是單單練唱的.
我們中間每一個人.
包括你包括我.
更多查聖經.
更多抄聖經.
更多物想聖經.
更多去聽錄音.
聽聖經.
讓神的話語牽動鼓勵激勵你.
你知道我為你洪水橋.
洪水橋塘.
洪恩塘.
我替你們興奮.
你知道嗎.
我上次也有說.
我們北部都會.
我以前跟他們教會在活石塘.
在粉嶺.
北部都會有很多人將會發展.
加上這邊有250萬.
很誇張.
你不要錄到後面.
下面看.
哪裡有人.
那裡還沒開發.
自這個前海.
你和我面向整個大灣區.
弟兄姊妹.
是大灣區.
不是北部都會250萬.
這麼少.
感受堂真是很有見識.
我覺得.
很看到上帝在中間給你們的感動.
能夠在洪水橋這個地圖裡面.
來建立堂會.
我這班弟兄姊妹.
求神激勵你們.

$^{801}$我想回望我信耶穌.
我家裡第一代信徒.
不是我爸爸媽媽.
甚至都不是我.
去做牧師.
現在已經死了.
我最先返教會那一位.
是我婆婆.
我還有印象.
很小的時候.
不過很多事情不記得.
只知道她帶我去我返教會.
不過在她晚年的時候.
她長年都要住院.
去到一個地步.
我都覺得很困難.
就是什麼呢?.
那時候醫學不算是很好.
要撤資.
撤資是什麼呢?.
簡單說.
你就睡在那裡.
不用想那麼多.
每一次去探望她.
我都會買燒鵝奶粉.
哄她開心.
但她的時間不多.
她回到天教.
但頂智未當的回想.
她一定有為她的家人來祈禱.
她看不到她的孫做牧師.
看不到她的孫卸任.
頂智未.
也看不到她眼前的俄彼得.
這個英女將來會怎樣.
但她就這樣渡過她的人生.
我這位頂智未.
我不知道當你和我面具整個香港的發展.
好像對一個中學生說.
你將來要做博士.

$^{841}$彭欣塘將來可以很興旺.
不是叫你現在做博士.
是真的.
有位中學同學.
我們在中學時期就叫他做博士.
結果怎樣呢?.
他真的做了博士.
是真的.
彭欣塘的頂智未.
上帝安放你在這個地圖.
不是純粹繼續維持.
我們看到你們很有宣教的眼界.
在牧者領袖帶領的底下.
一起能夠有更廣闊的眼界.
我為你們感恩.
真的.
我為你們感恩.
不要讓眼前很多叫你哀慳遮怨的東西.
錯失了看到上帝的那種恩典.
我更加想在第三點上和你們說.
不單是看生命.
不單是看這個英雄.
在路德基裡面.
最後一句話是.
你只有這個女子.
你有什麼?.
七個兒子更好.
路德基很疼愛你.
很貼近你.
陪伴你.
一直都沒有離開你.
各位頂智妹.
今天.
上帝是藉由三一上帝.
來陪伴恩愛我們.
讓你和我那派聖餐.
又或者剛才大家所提的.
是聖靈在我們整個會場裡面運行的.
誰陪伴我們?.
誰去表達愛?.

$^{881}$誰去犧牲?.
各位頂智妹.
你們知道戒命第一是什麼嗎?.
什麼是第一戒命?.
你要盡心盡聖.
盡義盡力愛主你的神.
是不是?.
我們剛剛的經文是什麼?.
聽神之道而遵守的人.
遵守什麼?.
盡心盡聖.
盡義盡力愛主你的神.
哇!上帝你這樣要求我.
過不過分一點?.
你聽不明白.
是聽不明白的.
為什麼這麼說?.
誰敢跟你說.
你愛不愛我?.
你說呢?.
陌生人跟你說.
你愛不愛我?.
誰跟你說的?.
你小時候媽媽跟你說的.
你愛不愛我?.
表達什麼?.
我都不知道多愛你.
我想聽聽你說.
爸爸媽媽今天是父親節.
說一句話他都哭了.
誰愛?.
為什麼戒命第一是你要盡心盡聖.
盡義盡力愛主你的神?.
你知道為什麼嗎?.
是因為上帝盡心盡義盡力愛你.
他才有資格這麼說.
他才不怕練武.
不怕你質詢.
不過.
我們小時候就好像畢業.

$^{921}$五億仔.
不多國.
不多資.
沒什麼回應.
當作應該.
聽不聽得明白?.
靈姐妹.
求神開我們眼睛看到.
遠遠片這個女子.
陪伴那阿米度過眾多艱難的日子.
讓你和我在艱難的日子裡看到.
不是神有沒有離開我.
是你不要離開神.
讓你和我這個生命.
你和我這個英雄.
你和我看到神萬語那種安慰.
以致我們匍匐在神的面前.
主讓我願愛你.
你讓我看到.
我就是世人眼中所看.
求你讓我們認識.
以致不單是你的生命.
你和我的成長.
你和我今天.
以致你和我將來.
是穩妥的.
我們信心能夠見到.
我們同心抵抗地去祈禱.
主讓我們知道你的話語是重要.
但我們卻是很想遺棄你的話語.
就好像撈米一樣.
明明聖經說很豐富.
很重要.
我們卻偏向記錄.
願你再一次借著.
這個不離不棄的愛.
牽動我們的心回到你的舊境.
安慰我們中間每一位的弟兄姊妹.
用你的愛榮耀懷抱.
和深深地進入我們每一個心口.

$^{961}$以致我們不懼怕.
所要面對和跨過.
知道主你必然要我們.
不單止同在.
你要將更加豐富,更加榮耀.
更加美善賜給我們.
求你燃點我們裡面那份盼望.
願你親自對我們說話.
彰顯你的榮辱.
吸引我們.
誰聽我們同心的祈禱和謝恩.
奉耶穌聖名祈求.
阿們.
\newpage



\section{路加福音 1:5-25}
\label{sec:3dWTg4gpwpQ}
\textbf{2023.06.25 《不相稱的禱告》 路加福音 1\_5-25;57-66 節 (和修版) 講員 : 陳韋安牧師}
\newline
\newline
連結: \href{https://youtube.com/watch?v=3dWTg4gpwpQ}{\texttt{https://youtube.com/watch?v=3dWTg4gpwpQ}} ~~~~ 語音日期: 2023-06-25
\newline
\newline
\hyperref[sec:NXj4I_yTkE8]{\small{< < < PREV SERMON < < <}}
~
\hyperref[sec:index_chronic]{\small{[返順時目]}}
~
\hyperref[sec:index_scriptual]{\small{[返順卷目]}}
~
\hyperref[sec:qziNPLIfqJE]{\small{> > > NEXT SERMON > > >}}
\newline
\newline
路加福音 1:5-25
\newline
\begin{longtable}{cl}
\hline
\hline
章節 & 經文 (和合本修訂版)\\
\hline
1:5 & \begin{tabularx}{0.7\textwidth}{X} 在希律作猶太王的時候，亞比雅班裡有一個祭司，名叫撒迦利亞；他妻子是亞倫的後代，名叫伊利莎白。 \end{tabularx} \\ \\ \relax
1:6 & \begin{tabularx}{0.7\textwidth}{X} 他們兩人在神面前都是義人，遵行主的一切誡命和條例，沒有可指責的。 \end{tabularx} \\ \\ \relax
1:7 & \begin{tabularx}{0.7\textwidth}{X} 只是他們沒有孩子，因為伊利莎白不生育，兩個人又年紀老邁了。 \end{tabularx} \\ \\ \relax
1:8 & \begin{tabularx}{0.7\textwidth}{X} 撒迦利亞按班次在神面前執行祭司的職務， \end{tabularx} \\ \\ \relax
1:9 & \begin{tabularx}{0.7\textwidth}{X} 照祭司的規矩抽籤，進到主的殿裡燒香。 \end{tabularx} \\ \\ \relax
1:10 & \begin{tabularx}{0.7\textwidth}{X} 燒香的時候，眾百姓在外面禱告。 \end{tabularx} \\ \\ \relax
1:11 & \begin{tabularx}{0.7\textwidth}{X} 有主的一個使者站在香壇的右邊，向他顯現。 \end{tabularx} \\ \\ \relax
1:12 & \begin{tabularx}{0.7\textwidth}{X} 撒迦利亞看見，就驚慌害怕。 \end{tabularx} \\ \\ \relax
1:13 & \begin{tabularx}{0.7\textwidth}{X} 天使對他說：「撒迦利亞，不要害怕，因為你的祈禱已經被聽見了。你的妻子伊利莎白要給你生一個兒子，你要給他起名叫約翰。 \end{tabularx} \\ \\ \relax
1:14 & \begin{tabularx}{0.7\textwidth}{X} 你必歡喜快樂；有許多人因他出世也必喜樂。 \end{tabularx} \\ \\ \relax
1:15 & \begin{tabularx}{0.7\textwidth}{X} 他在主面前將要為大，淡酒烈酒都不喝，從母腹裡就被聖靈充滿。 \end{tabularx} \\ \\ \relax
1:16 & \begin{tabularx}{0.7\textwidth}{X} 他要使許多以色列人回轉，歸於主—他們的神。 \end{tabularx} \\ \\ \relax
1:17 & \begin{tabularx}{0.7\textwidth}{X} 他將有以利亞的精神和能力，走在主的前面，叫父親的心轉向兒女，叫悖逆的人轉向義人的智慧，又為主預備迎接他的百姓。」 \end{tabularx} \\ \\ \relax
1:18 & \begin{tabularx}{0.7\textwidth}{X} 撒迦利亞對天使說：「我怎麼能知道這事呢？我已經老了，我的妻子也年紀老邁了。」 \end{tabularx} \\ \\ \relax
1:19 & \begin{tabularx}{0.7\textwidth}{X} 天使回答他說：「我是站在神面前的加百列，奉差遣來對你說話，把這好信息報給你。 \end{tabularx} \\ \\ \relax
1:20 & \begin{tabularx}{0.7\textwidth}{X} 到了時候，這些話必然應驗；只因你不信我的話，你會成為啞巴，不能說話，直到這些事實現的日子。」 \end{tabularx} \\ \\ \relax
1:21 & \begin{tabularx}{0.7\textwidth}{X} 百姓等候撒迦利亞，詫異他在聖所裡遲延那麼久。 \end{tabularx} \\ \\ \relax
1:22 & \begin{tabularx}{0.7\textwidth}{X} 到他出來，卻不能和他們說話，他們就知道他在聖所裡見了異象；他直向他們打手勢，因為他成了啞巴。 \end{tabularx} \\ \\ \relax
1:23 & \begin{tabularx}{0.7\textwidth}{X} 他供職的日子一滿，就回家去了。 \end{tabularx} \\ \\ \relax
1:24 & \begin{tabularx}{0.7\textwidth}{X} 這些日子以後，他的妻子伊利莎白就懷孕，隱藏了五個月； \end{tabularx} \\ \\ \relax
1:25 & \begin{tabularx}{0.7\textwidth}{X} 她說：「主在眷顧我的日子，這樣看顧我，要除掉我在人前的羞恥。」 \end{tabularx} \\ \\
[1ex]
\hline
\hline
\end{longtable}
$^{1}$大家好,我是羅家峰.
今天我們會討論禱告的道.
大家也很熟悉祈禱.
不過我們也一起反省一下禱告的生命.
我們藉著這段經文.
羅家峰第一章是一段聖誕節的故事.
基本上是耶穌基督出生的故事.
其實也是施西約翰出生的故事.
兩個嬰兒出生的故事.
其實也是施西約翰出生的故事.
兩個嬰兒出生的故事.
彼此交叉在故事的編排裡.
第5到25節是天使加拔列向撒加尼亞去顯現.
第5到26節是馬利亞的祈禱.
第5到27節是撒加尼亞的禱告.
羅家峰第一章是彼此平衡兩個嬰兒出生的故事.
所以羅家峰會把禱告特別強調在這本書中.
今天我們也在楊核利看一段經文.
發現很多跟禱告有關的段落.
並不是純粹馬利亞和尊主頌的祈禱.
而是整個故事的編排中有很多特別的禱告安排.
我們看看第五至十節的經文.
在希律作猶太王的時候.
阿比亞班有一個祭司名叫撒加里亞.
他妻子是亞倫的後代名叫伊莉莎白.
他們兩人在神面前都是二人.
遵行主的一切誡命和條例.
沒有可知責定.
只是他們沒有孩子.
因為伊莉莎白不生育.
兩個人又年紀老了.
撒加里亞就按班次在神面前執行祭司的職務.
按照祭司的規矩抽籤.
進到主的殿燒香.
燒香的時候眾百姓在外面禱告.
有主的使者就在香壇的右邊向他顯現.
經文記載撒加里亞是一個什麼身份.
根據福音書講他是一個祭司.
祭司是一個做什麼的.
祭司在當時即將在聖殿燒香.

$^{41}$燒香其實不是普通中國的燒香.
而是代表著百姓的祈禱.
原來禱告的希伯來文.
Tafina這字是解作如煙往上升的意思.
當百姓燒香的時候.
其實這些就代表著他們向上帝的禱告.
而祭司就負責將這些禱告呈獻給上帝.
在天上的天來垂聽禱告的上帝.
所以根據猶太人的習慣.
祭司其實有一個中介人的身份.
負責幫弟兄姊妹的禱告呈獻給神.
所以經文特別強調.
這個專門向人家幫人去禱告呈獻給神的人.
就在這個時候.
天使就向他顯現.
剛剛是什麼時候?.
就是祈禱的時候.
甚至燒香的時候眾百姓在外面禱告.
就在這個時候就有主的使者.
站在香壇的右邊向他顯現.
即是一個專門專業地.
他職責幫人祈禱的人.
就在外面禱告聲不斷的時候.
就在這個時候.
天使加勃尼就顯現.
告訴你自己的禱告就應驗了.
一個非常多禱告的時間和地點和人物.
你會發現在眾百姓禱告的時候.
一個專門幫人禱告的人.
他自己的禱告事項.
來到蒙英雲.
整段經文裡面.
其實都是很講述一個完美的時刻.
用今天教會的例子.
我們祈禱會的時間裡面.
一個專門幫人禱告的牧師.
可能是牧師或者是教會傳道人.
他自己的禱告事項應運了.
就在這個時候.
天使就跟他說.

$^{81}$你的禱告蒙英雲.
於是天使就一問一受.
就將要說的話說了.
你的妻子二十八月生一個孩子.
你要給她起名做若翰.
她必歡喜快樂.
有許多人因她出世.
也必歡喜快樂.
在主面前將要為大.
但走龍走得不渴.
從牧福裡面就被聖靈充滿.
她要使許多人以色列文回傳.
歸於上帝.
基本上是說.
你將要說若翰出世的消息.
你的禱告蒙英雲.
令人震驚的是.
似乎撒迦利亞的回應.
實在令人詫異.
十八節他說.
撒迦利亞對天使說.
我憑什麼可知道這事.
我已經老了.
我的妻子也老了.
撒迦利亞對於天使的.
禱告perfect moment的顯現.
他竟然回應什麼.
我又這樣祈禱過嗎.
他不單不敢相信.
更似乎忘記了.
自己曾經有這樣祈禱過.
撒迦利亞忘記了.
自己曾經在N年前.
這樣的一個禱告.
不過似乎也很正常.
我們一生人通常說這麼多話.
誰記得這麼多話呢.
大家也是.
說出口.
說過就不太記得.

$^{121}$聽過就記得.
無論是好說話.
還是很傷心的說話.
說過就不太記得.
聽過就很記得.
所以似乎撒迦利亞忘記了.
曾經有這樣的禱告.
而上帝又應驗了.
所以撒迦利亞似乎是這樣.
不過也不知是否正常.
你記不記得自己祈禱過.
你記不記得你今早祈禱過什麼禱告.
你記不記得你昨晚臨睡前.
祈禱過什麼禱告.
或者上星期的國主席.
你記不記得我們帶他去祈禱的時候.
祈禱過什麼禱告.
通常都不記得.
一生人祈禱這麼多.
祈禱這麼多.
所以似乎情況是這樣.
天使撒迦利亞是一個專門幫人.
將他的禱告呈獻給神.
一個專門專業地每天都這樣做的人.
他似乎自己的禱告不是很記得.
也不是很留心自己曾經這樣祈禱.
於是天使就這樣說.
到時候這些話必然應驗.
神必定會應驗.
這與你記不記得有關.
不過既然你這樣不信.
又不太記得.
你就要成為啞巴.
不能說話.
這些事實現的子.
對於天使加巴列的懲罰.
我們覺得似乎有一點點的差異.
雖然成為啞巴.
一直到嬰兒出生的時間.
都要成為啞巴.

$^{161}$十個月.
十個月要成為啞巴.
又不是很大的懲罰.
十個月不說話.
又不是短的時間.
你想想以前你做爸爸.
很多年前你兒女出生的時候.
如果你老婆出生懷孕十個月.
你不能說話.
其實都很麻煩.
不能打電話.
不能叫的士.
登記又不能說話.
其實都很麻煩.
所以為什麼天使加巴列.
要這麼嚴厲的懲罰撒迦利亞呢?.
其實信仰裡面都不少生子不信的故事.
有很多這樣的故事.
很老邁的.
生不了.
我們其中一個就是哈拿.
哈拿就是這樣.
哈拿心裡受苦.
痛痛哭泣.
土古就有話說.
萬官之說.
你若眷顧婢女的苦情.
顧念不忘婢女.
賜我一個兒子.
我必使他終身歸於王華.
不用再看他一頭.
另一段的經文都是舊約的.
就是阿巴拉罕.
阿巴拉罕和太太薩拉都是這樣.
薩拉的老邁.
月經都斷絕了.
經中說.
薩拉心裡暗笑.
我既已衰敗.
我處於也老邁.

$^{201}$豈能有這喜事呢?.
耶和華對阿巴拉罕說.
薩拉為什麼暗笑.
其實都很麻煩.
上帝看到你暗笑他.
我既已年紀老邁.
果然生養嗎?.
耶和華豈有難成的事呢?.
多了一期.
明年這個時候.
我必回到你這裡.
薩拉必生一個兒子.
薩拉就害怕不承認.
我沒有笑.
上帝就說.
不然你實在笑了.
都很麻煩.
笑上帝還要不承認.
還說你笑我.
但你會發覺.
耶和華都沒有罰薩拉.
就算薩拉不相信.
薩拉甚至暗笑都好.
但耶和華都沒有罰過薩拉.
為什麼這次.
天使加百年會這麼玻璃心.
只是不記得.
不記得自己曾祈禱.
為什麼要罰薩加尼亞呢?.
我只覺得.
這可能是對祭司的懲罰.
他是祭司.
他是一個專門.
每天都將禱告呈獻給上帝的人.
在聖殿裡事奉.
禱告是他的專業.
他每天都要做.
幫人做.
都是代求的一個人.
所以一個祭司.

$^{241}$他自己的禱告.
可能祈禱太多.
信口開河.
一時不記得.
甚至不重視.
他很賞心的祈禱.
可能有這樣的緣故.
天使特意要懲罰他.
要他十個月不能開口.
不能開口.
言下之意是不能開口去祈禱.
你說不是啊.
蜜禱就行.
不好意思.
舊約不是很流行蜜禱.
全部都是張開手.
張開眼睛大聲地祈禱.
不多蜜禱的時間.
所以對撒該來說.
這十個月.
可能是一個很特殊的經驗.
不能夠開口去祈禱.
經文最後.
剛才大家讀的經文.
說到十個月之後.
57節.
伊利沙伯的產期到了.
新一個兒子.
論社親屬聽見處.
向他大師憐憫.
就和他一同歡樂.
到了第八天.
他們來給孩子行國禮.
並要照他父親名字.
叫他撒迦利亞.
他母親回應說.
不要叫他約翰.
他們都一說.
你親屬沒有叫這個名字的.
他們就向他父親打手勢.

$^{281}$問他叫什麼名字.
約翰就說不出話來.
撒迦利亞就說不出話來.
就要求人寫一塊寫字板出來.
寫字.
他的名字是約翰.
他們就驚訝.
六十四節.
撒迦利亞的口就立刻開了.
舌頭也鬆了.
舌頭就突然說話了.
就突然開口說話.
稱頌神.
撒迦利亞一開口.
第一句說話是什麼.
就是一個Hallelujah.
讚美主的禱告.
這個Hallelujah.
不單單因為孩子出世.
也不單單因為他能說話.
而是他十個月不能說話.
想了很久.
十個月不能夠禱告.
一開口就是一個讚美神的禱告.
可能這件事對撒迦利亞來說.
是一個很特別的經驗.
曾經忘記.
然後就有一個這樣的小小的懲罰.
然後就這樣去思考自己的禱告.
一段很有趣的經文.
我們想想這段經文對我們有什麼特別的意思.
我覺得撒迦利亞最大的問題.
就是他作為一個祭司.
他太過熟悉祈禱.
可能禱告太多.
以致他不太認真看待自己的禱告.
以致他跟他的生命和禱告是相稱.
不知道大家是否這樣.
可能大家也做了很多條基督徒.
你不用別人教你怎麼祈禱.

$^{321}$甚至乎你做崇拜的事奉的時候.
崇拜主席或者團契大徒.
你經常都會隔天就熟.
我們也是.
我們在教會小組的群組裡.
我們經常都是這樣.
我們祈禱的時候.
就是按下手指出來.
然後就祈禱.
不知道是不是這樣當祈禱了.
按完那些emoji之後.
就會不會再祈禱呢?.
所以我自己也是很怕忘記了祈禱.
特別是做傳道人.
聽到頂智滿水的時候就說你祈禱.
不知道大家是不是真的有祈禱.
這樣說之後.
通常我以前也會馬上.
回家後就立即回復facebook.
求主憐憫他.
起碼先祈禱一個.
回家再提.
你重不重視你自己的禱告.
你重不重視你自己曾經說幫人祈禱.
教會有很多這些大老實行.
每個星期都有.
每個星期都有這些事情.
你有沒有真的好好地去看待這些禱告.
在團契群組裡的事情.
有沒有真的好好地去紀念它.
事實上為你祈禱.
是一句最肅靈.
而又最難度的方法.
可以跟別人說祈禱.
你真的不用去行動幫忙.
你想想教會有個弟兄姊妹.
有重病有癌症.
你會為他祈禱.
但如果你的子女有癌症呢?.
如果你子女有癌症.

$^{361}$你就不單單為他祈禱.
而是把整個心放在一起.
一起來跟他想想怎麼辦.
所以我想禱告的重點不是.
對我們來說.
不是越祈越精的事情.
我經常覺得我們祈禱太多.
以致我們的禱告是很美的.
有很多經文.
由創世記的祈禱.
到啟示錄的經文都貼在這裡.
我們祈禱很美.
但我們似乎不太記得自己祈禱了什麼.
甚至不知道是不是經驗.
在公土的時候.
有時候你會想下一句說什麼.
就會說一些話出來.
那句你不知道說什麼.
可能你還沒想下一句就說了.
問題不是我們不會祈禱.
而是我們祈禱太多.
以致我們沒有好好重視.
我們每一句的禱告.
所以我想撒嘉利亞的故事.
第一個很大的提醒.
就是你好好地來重視你的禱告.
你自己的禱告.
跟你的生命是不是相稱.
禱告裡面的你.
可以祈禱天花朗鳳.
但你是不是真的認真看待.
這些禱告.
你在群組裡按了很多手.
但是是不是真的來真正地.
把這個禱告重視呢.
這是第一點.
好好地來重視我們的禱告.
我們的生命跟禱告應該是相稱的.
第二就是珍惜每一次禱告的機會.
撒嘉利亞也是這樣.

$^{401}$一個十個月的屬靈操練.
十個月不能祈禱.
然後他一開口.
所以不知道大家有沒有試過.
被你十個月不能祈禱會怎麼樣.
你不要試了.
但是當你發現原來不能禱告的時候.
你才會去珍惜每一次可以開口的禱告.
所以好好地珍惜每一次禱告的機會.
特別是傳道人邀請你祈禱的時候.
以前我們團契是這樣的.
年輕的時候.
傳道人叫你祈禱的時候會怎樣.
就會猜輸了的棋.
或者猜贏了的棋.
感覺上好像輸了一樣.
好像有點委屈.
好像不太自願地去祈禱.
我自己也是.
我自己是傳道人.
我讀神學開始到現在.
由我拍拖到我結婚這麼多年.
每次我去外父吃飯的時候.
外父都叫我祈禱.
他叫我John仔.
John仔 祈禱.
就是大祈禱.
這二十年都是這樣.
由拍拖到現在每次都是我大祈禱.
又是我 可不可以叫別人.
但我覺得不應該這樣想.
作為傳道人或作為基督徒.
有人叫你祈禱.
是一個非常好的光榮的事.
是一件很快樂的事情.
所以我經常說.
就算任何一個人叫我祈禱.
都會立刻要去反應.
像是撞下彩一樣迎接呼召.
去大祈禱.

$^{441}$新人好好去享受.
去珍惜每一次能夠祈禱的機會.
禱告是一件我們非常重要的事情.
不要縮 不要覺得不懂.
因為禱告是一個很享受的時間.
應該是一個很享受的時間.
最後說一點.
我聽過一個最好的祈禱.
在我一生人裡.
這禱告是我媽媽的祈禱.
那時我剛回錦村不久.
我帶我爸媽回教會.
後來他在教會裡信奉耶穌.
我們住天水圍.
那時有分區祈禱會.
都有天水圍區.
那時星期三晚踢著拖鞋.
帶上了住一個月的媽媽.
回到祈禱會.
大家在家裡祈禱.
媽媽不懂祈禱.
但也為我祈禱.
那時我媽媽怎樣祈禱呢?.
她用平時跟我說話的方式祈禱.
平時跟我說什麼.
不要太晚睡.
不要傷身.
回來喝湯.
不要做得太累.
把這些說話變成了祈禱.
她不懂祈禱.
不懂得說奉主名求.
不懂得說那些.
沒有經文的東西.
只說.
神啊.
求你叫我兒子不要太晚睡.
不要太累.
不要太傷身.
這是我聽過最感動.

$^{481}$也是最好的禱告.
因為這個禱告跟她的生命是一樣的.
她自己所關心的事.
她自己所關懷的事.
每一句都是她自己.
真心誠心所願的說話.
所以我們也是一樣.
學習去珍惜.
去留心每一個祈禱.
我們不需要祈禱得太快.
祈禱得太多.
也不需要很多經文的拼貼.
單單心裡有一句.
就將那句慢慢地說出來.
無論是憤怒的.
無論是擔心的.
無論是開心的.
都將那句說出來.
這就是最好的祈禱.
有時學習慢慢祈禱.
不要太多屬靈的拼貼.
很多用來買時間的東西.
慢慢地祈禱.
將一句的禱文一句一句說出來.
我們試試一起祈禱.
我們試試在位置裡.
很簡單.
將你現在所想所求.
無論是你面對的生命.
生活,工作.
你現在擔心的事情.
甚至乎你沒有擔心的.
祈求的事情.
就將那件事求天父上帝去說.
簡單地說.
但認真用心地說.
一起祈禱.
(靜音).
響鐘.
天父,多謝你讓我們祈禱.

$^{521}$你剛才聽了我們的祈禱.
我們心裡每一位弟兄姊妹的所想所求.
求主你應驗,求主你應運我們的禱告.
讓我們每天都可以認真享受這個禱告的時間.
將我們的生命成為一個禱告,去禱告你.
為我們教會弟兄姊妹的身體去禱告你.
求主你醫治,求主你憐憫.
為我們各方面的需要禱告你.
求主你這樣的許願,所以聽我們的祈禱.
奉主名求,阿們.
(開啟字幕).
\newpage



\section{約翰福音 21:15-19}
\label{sec:qziNPLIfqJE}
\textbf{2023.07.02 《你愛我比這些更深嗎?》約翰福音 21\_15-19 (和修版) 彭秀芹導師}
\newline
\newline
連結: \href{https://youtube.com/watch?v=qziNPLIfqJE}{\texttt{https://youtube.com/watch?v=qziNPLIfqJE}} ~~~~ 語音日期: 2023-07-03
\newline
\newline
\hyperref[sec:3dWTg4gpwpQ]{\small{< < < PREV SERMON < < <}}
~
\hyperref[sec:index_chronic]{\small{[返順時目]}}
~
\hyperref[sec:index_scriptual]{\small{[返順卷目]}}
~
\hyperref[sec:5F_Xzyjh7_k]{\small{> > > NEXT SERMON > > >}}
\newline
\newline
約翰福音 21:15-19
\newline
\begin{longtable}{cl}
\hline
\hline
章節 & 經文 (和合本修訂版)\\
\hline
21:15 & \begin{tabularx}{0.7\textwidth}{X} 他們吃完了早飯，耶穌對西門‧彼得說：「約翰的兒子西門，你愛我比這些更深嗎？」彼得對他說：「主啊，是的，你知道我愛你。」耶穌說：「你餵養我的小羊。」 \end{tabularx} \\ \\ \relax
21:16 & \begin{tabularx}{0.7\textwidth}{X} 耶穌第二次又對他說：「約翰的兒子西門，你愛我嗎？」彼得對他說：「主啊，是的，你知道我愛你。」耶穌說：「你牧養我的羊。」 \end{tabularx} \\ \\ \relax
21:17 & \begin{tabularx}{0.7\textwidth}{X} 耶穌第三次對他說：「約翰的兒子西門，你愛我嗎？」彼得因為耶穌第三次對他說「你愛我嗎」，就憂愁，對耶穌說：「主啊，你無所不知，你知道我愛你。」耶穌說：「你餵養我的羊。 \end{tabularx} \\ \\ \relax
21:18 & \begin{tabularx}{0.7\textwidth}{X} 我實實在在地告訴你，你年輕的時候，自己束上帶子，隨意往來；但年老的時候，你要伸出手來，別人要把你束上，帶你到不願意去的地方。」 \end{tabularx} \\ \\ \relax
21:19 & \begin{tabularx}{0.7\textwidth}{X} 耶穌說這話，是指彼得會怎樣死來榮耀神。說了這話，耶穌對他說：「你跟從我吧！」 \end{tabularx} \\ \\
[1ex]
\hline
\hline
\end{longtable}
$^{1}$各位聽眾平安.
不知道你們有沒有被人問過.
你愛不愛我呢?.
可能是父母問子女.
你愛不愛我?.
或者是子女問父母.
你愛不愛我?.
或者是太太問老公.
你愛不愛我?.
問這個問題的人.
你猜最想聽到什麼答案呢?.
我愛你.
最討厭聽的就是我不知道.
最不喜歡聽的.
或者是有猶豫的.
你都不喜歡聽.
最想聽的是我愛你.
問這個問題.
這兩個人.
我相信關係已經很密切.
很重視對方.
你才會問.
我們分享的《約翰福音》第21章.
其實耶穌已經問了彼得這個問題.
不只問了一次.
還問了三次.
重要的事情要說多少次呢?.
三次!.
所以很重要.
在看這章經文之前.
我們先看看經文的背景.
《約翰福音》第20章.
很有趣.
好像結束了.
第十章說.
我記載了這些事.
叫你們信.
基督是神的兒子.
而且信他的名字.
生命.

$^{41}$好像結束了.
好像結論.
為什麼結束了.
還未結束.
有21章.
所以有些學者估計.
可能是約翰之後補回的.
先不說這些問題.
只是告訴大家.
第21章是一章很特別的聖經.
是一章有血有肉的聖經.
我們一起.
去到第21章的場景.
就是加利利海.
即是提比利亞海邊.
提比利亞海邊.
剛才說的又叫加利利海.
原來.
就是彼得.
約翰.
亞國安德烈.
主耶穌在那裡呼召他們.
他們當時正在打魚.
現在第21章.
他們也正在打魚.
也是這個地方.
第21章這個場景和人物.
都很熟悉.
第21章和第一章.
好像首尾呼應.
好了.
彼得從操.
僱業去打魚.
還有.
多瑪.
拿達葉.
亞國約翰.
還有兩個沒有記名的.
我們估計是安德烈和弗力.
因為通常聖經都是這個組合.

$^{81}$他們打了一整晚.
都沒有收穫.
後來聽到.
有人說.
你把魚網放在右邊.
他們聽完.
就照做.
竟然在淺水地方.
捕捉了53條大魚.
奇蹟.
主所愛門徒.
是紫若翰.
跟彼得說是主.
於是彼得就馬上跳進水裡.
游上岸.
去找耶穌.
但很有趣.
游上岸.
彼得又不敢跟主說話.
究竟他們發生什麼事呢.
我們可以一起.
看看.
究竟二十一章.
這個情節.
一群彼得.
和一群門徒.
重演了.
初初被主呼召的一幕.
這一幕讓彼得.
回想起很多事.
我們.
一起進入二十一章.
當然.
大家都知道彼得.
為什麼.
這麼沒心機.
雖然主復活了.
但他還有一些.
事.
失敗經驗或陰影.

$^{121}$未解決.
耶穌顯現.
要解決彼得的失敗.
心裡的一件事.
面對失敗迎向復興.
這就是主.
向彼得很想做的工作.
他們.
剛才說過.
開頭.
主耶穌.
去到加利利海.
找這群人.
為他們.
預備了早餐.
大家吃過BBQ早餐嗎.
煮的很特別.
BBQ為他們預備早餐.
燒魚和燒餅.
為他們預備早餐.
吃早飯.
耶穌.
吃完早飯.
跟西門彼得說.
約翰的兒子西門.
這句說話很有意思.
約翰的兒子西門.
其實在.
約翰科學第一章.
彼得.
被主呼召的時候.
主也是這樣叫他.
約翰的兒子西門.
讓彼得想起.
他怎樣.
被主呼召.
他怎樣.
放下願望.
跟從主.
他對主呼召的回應.

$^{161}$原來.
其實很刻意.
就等於我的全名.
彭秀琴.
如果有人問我彭明的女兒彭秀琴.
很刻意.
是指誰.
是我.
彭秀琴可能這個世界也有第二個.
同名的.
連爸爸也同名.
原來主是刻意.
是你彼得.
我沒有搞錯.
西門彼得.
就是約翰的兒子.
原來今天主.
是專門來找彼得.
今天主也專門來找我們每一個.
說出主對彼得的重視.
刻意的安排.
為什麼要刻意安排呢.
原來有些事情.
彼得還沒解決.
但在還沒解決之前.
記不記得.
剛才說的背景.
彼得打了一晚的雨.
也沒有收穫.
他們整晚都沒有吃東西.
在耶穌要找彼得.
聊天的時候.
我想彼得的肚子在咕咕咕地響.
那耶穌呢.
就怎樣.
原來早已預備.
早已預備.
有炭火.
有燒魚燒餅.
原來.

$^{201}$不好意思.
原來主很奇妙地.
應該說主很有溫情地.
去解決人的問題.
他知道人的軟弱.
他顧念彼得身心的需要.
知道他肚餓.
吃飽了再說.
好了.
今天我們就說到這裡.
那.
吃飽了再說.
很奇妙.
你想想當時的場景.
七個大男人在主面前.
吃早餐.
但不知為何.
沒有人敢跟主說一句話.
連平時很多事情.
說的彼得都不敢.
問你是否主.
因為他們知道.
其實.
他知道是主.
但又不敢說.
我們推測.
為何他不敢認住呢.
有些可能性.
可能是這群門徒.
知道主復活.
但對於主復活的事實.
還在消化中.
也有可能.
他沒有想過.
這位復活的主這麼親切.
會為他們預備早餐.
知道他們的軟弱.
顧念他們.
整晚打魚.
都沒有吃過.

$^{241}$他想不到復活主.
這麼親切.
明白他們.
原來主就是這樣.
這位復活主.
顧念人的.
身心需要.
體恤和溫柔.
大家記得彼得.
曾經三次不認主.
但在三次不認主之前.
彼得在公開場合.
在門徒面前.
誇口.
說:我是很靠得住的.
他對耶穌說了什麼呢.
我們一起讀.
準備.
1 2 3.
彼得對他說:主啊.
為什麼我現在不能跟你去.
我願意為你.
寫命.
耶穌回答:你願意為我寫命嗎.
我實實在在地告訴你.
雞叫而馳.
你要三次不認我.
彼得曾經在很多人面前.
誇口說:我願意為你.
寫命啊耶穌.
但耶穌怎麼回答他.
你要三次不認我.
這一次彼得.
對耶穌的承諾.
是公開的.
你想想.
你曾經在別人面前.
公開說我怎樣怎樣.
但如果你做不到.
你面對的時候.

$^{281}$你會怎樣.
會很瘀青.
無地自容.
信譽受損.
但耶穌.
他給彼得.
再一次公開.
對門徒說.
說了三次.
就是.
說:主啊.
我愛你.
給彼得一次.
公開.
讓人知道.
耶穌是接納他的.
耶穌.
確認了.
他對他的愛.
所以.
恢復彼得的.
事奉崗位.
除去彼得心中的.
很多陰影.
主還刻意.
安排了一個環境.
除了主.
初呼召彼得的環境.
還有.
BBQ的那堆火.
是否似曾相識.
對於彼得來說.
其實是似曾相識.
記不記得彼得三次.
不認住的時候.
在哪個地方不認住.
就是在.
炭火路邊.
因為當時很冷.
他在取暖.

$^{321}$有個侍女問彼得.
你不是和納撒蘭耶穌.
那一黨的人.
他就說不是.
他就是在炭火路邊.
三次不認住.
耶穌刻意安排在.
炭火路邊.
讓他說三次.
他愛主.
這個場景.
讓彼得.
回想.
他自己怎樣失敗.
今天耶穌.
不是要我們忘記自己失敗.
是忘記不了.
我們不能洗了自己的記憶.
洗了自己的腦.
忘記自己曾經不愛主.
或者對主的承諾.
我們沒有兌現.
我們忘記不了.
主也不是要彼得忘記失敗.
主是要彼得面對失敗.
在.
心理學治療裡.
有個叫記憶治療法.
就是將一個人.
他面對的痛苦.
帶回類似的場景.
讓他心裡的傷痛.
懼怕.
焦慮.
失敗的罪疚感.
去面對和處理.
彼得的失敗.
耶穌知道.
約翰知道.
彼得知道.

$^{361}$門徒也知道.
耶穌不是要抹去他們的記憶.
耶穌是要他們面對.
耶穌要將彼得的.
失敗痛苦.
帶回神面前.
得到醫治.
遙恕.
弗洛伊德.
一個很出名的心理學家.
他這樣說.
焦慮會令人產生刺痛和自責.
這個自責.
會令人覺得自己.
給自己的永遠都是負評.
我永遠都是不好.
我有一百個理由說自己不好.
然後掉入.
身心無力.
感到無助.
很多時候.
我們都會有這個狀態.
但原來主就是那位.
醫治者.
主是最明白.
願意治療我們最心底的軟弱.
和失敗.
我想和大家.
分享一個故事.
叫做黃絲帶的故事.
可能大家都聽過.
有一位.
先生結婚了.
他犯錯.
要坐牢三年.
他快要放監.
他不敢回家.
因為不知道太太會不會原諒他.
於是他寫了一封信.
和太太說如果你原諒我.

$^{401}$就在家門口.
那棵大樹上.
綁上一條黃絲帶.
我就知道你會.
原諒我和接納我.
如果我看不到黃絲帶.
在車途中.
看不到黃絲帶.
我就不會下車.
我會忘記.
我和你之間的關係.
我都不會原諒自己.
於是放監那天.
這個人很忐忑.
如果是我.
我都會很忐忑.
有黃絲帶 沒黃絲帶.
他在車上顫抖.
他就在想有沒有黃絲帶.
有沒有黃絲帶.
他很害怕.
他不知道太太會不會原諒他.
不知道太太會不會接納他.
當然他看到.
有黃絲帶.
代表什麼.
他太太願意接納他.
他當然很開心.
當他再仔細看的時候.
不單止他門前.
那棵樹上有黃絲帶.
整條村的樹都掛上黃絲帶.
表示什麼.
他很願意接納他.
他期待他回來.
就像主一樣.
他期待充滿憂傷.
充滿失敗.
充滿內疚的彼得.
再次到神面前得醫治.

$^{441}$得復興 得幫助.
除了要醫治彼得.
原來主也讓彼得再一次重新認愛.
肩負重任.
怎樣會令人信任他.
愛他.
信任他.
接納他.
他將重要的事給他做.
有些人說我相信你.
但小事都不成.
我找另一個.
那你想想那信是怎樣.
那不是真的信.
原來重任是交給你最信任的人.
你愛的人去做.
原來今天主對彼得的信任不是說說.
不是.
他將他很重要的重任.
交給彼得.
有彼得和主的對話.
有三次主問彼得的問題.
我們一起讀.
我讀主的部分.
你們讀彼得的部分.
我們準備.
若翰的兒子西門.
你愛我比這些更深嗎.
主啊 使的.
你知道我愛你.
你餵養我的小羊.
若翰的兒子西門.
你愛我嗎.
主啊 使的.
你知道我愛你.
你餵養我的小羊.
若翰的兒子西門.
你愛我嗎.
主啊 你是無數不知的.
你知道我愛你.

$^{481}$你餵養我的小羊.
彼得.
給主問三次.
你愛我嗎.
你看到第一,二,三次.
是一樣的問題.
但第一個問題有點不同.
你愛我比這些更深嗎.
這些是眾數.
是甚麼呢.
這裡沒有說.
我們猜猜.
彼得的打魚事業.
打魚工具.
彼得一班好兄弟.
其實一班門徒兄弟.
也不能說.
知道彼得失敗.
他去打魚.
這班是好同伴.
是否說這班好兄弟呢.
還是說彼得.
只彼得對著承諾沒有兌現.
那份罪疚感.
內疚感.
失敗感.
那份傷痛.
我們不知道.
這是我們的估計.
但我猜.
他曾經口哭.
沒有兌現.
人很慘的.
有內疚感.
但耶穌說.
你愛這些多過愛我嗎.
有時我們也會這樣.
這些難阻我們愛神.
彼得說.
主啊,是的,你知道我愛你.

$^{521}$彼得雖然有軟弱.
有失敗.
但他真的愛主.
主也沒有否定他.
主沒有說.
你愛他嗎.
是不是真的.
你見到他上次.
三次不認我.
我和你算舊帳.
耶穌沒有和他算舊帳.
耶穌沒有否定彼得對他的認愛.
而且每一次彼得對耶穌有愛的認信.
耶穌都將重任交給彼得.
你餵養我的小羊.
你餵養我的羊.
有些學者會說.
愛和餵養是怎樣.
在這裡我不處理.
釋經的問題.
如果有興趣.
我們可以一起討論.
因為這不是最重要.
最重要是.
究竟.
主如何回應彼得.
彼得如何回應主.
主如何醫治復興彼得.
主將重任交給彼得.
你餵養我的羊.
是我的.
不是彼得的.
是主的羊.
原來今天去服侍主.
我們去牧羊教會.
最重要的不是有多厲害.
有多厲害.
有多多風光為夕.
而是一個純粹.
沒有雜質愛神的心.

$^{561}$彼得其實也傷痕累累.
經常口渴.
有軟弱失敗.
很衝動.
但因為他真是有一個.
很真誠愛主的心.
主將重任交給他.
原來今天.
世界上沒有什麼屬靈的偉人.
他一走出來.
很有神聖.
不是.
其實我們每一個.
都是一個很軟弱.
愛神的罪人.
我們很需要神的憐憫.
真是沒有人可以誇口.
唯有倚靠主.
倚靠神.
和一個愛主的心.
和你分享一個.
一位信義代的基督徒.
他叫嘉進.
他出世就回教會.
當然他不愛主.
他又要回教會.
他十一歲開始就沒有回教會.
十三歲那年就加入黑社會.
中二就沒有讀書.
之後他經常出入青少年的監獄.
到十八歲那年.
他每天都吃六安童.
即是K仔.
最後因為藏毒吸毒和毒壓.
他被捕.
再一次坐牢.
這次十八歲要坐牢.
他都沒有想過要回頭.
他已經忘記了神.
直到他二十二歲.

$^{601}$發現人家吃K仔有很多後遺症.
他覺得自己身體極度之差.
他想找個地方休養.
重新來過.
重新來過不是轉好.
重新來過休養完.
又出來.
比WET過更差.
於是他去了西貢鄧家灣的靈愛中心.
接受科普解讀.
他不是因為信福音有大能.
而是他想找個地方休養身體.
但很奇妙地主安排了同伴和弟兄.
去跟他靈修看聖經.
他在當中重新認識上帝.
悔改.
體會上帝對他的接納.
在二十三歲那年.
他覺得讀書很重要.
於是他很勤力地讀社工副學士的課程.
他畢業後.
他一心一意很想做社工去註冊.
但因為他有太多案底.
他有太多不好的事.
社署拒絕了這次註冊.
很多時候我們犯錯.
人未必會給我們機會.
但身邊很多人鼓勵他.
當然神鼓勵他.
他繼續讀書.
他讀了社工學位課程.
在這段期間.
他經常回聚會.
他每天祈禱親近神.
他感受到神鼓勵他.
幫助他.
當他社工學位畢業.
教會這次推薦他.
他這次可以正式成為註冊社工.
而且還回到他之前的機構.

$^{641}$去做社工.
幫助吸毒的年青人.
家俊軟弱失敗.
但神不會給他機會.
但神給他機會.
他再次領受神對他的照明.
他起來服侍神.
去幫助人.
今天每一位愛主的人.
其實都有照明.
原來神今天在你身上有照明.
彭彭姑娘不要說笑.
我幾十歲.
我又不像年輕人走來走去.
我怎樣服侍主?.
不是的.
其實今天服侍主很簡單.
首先有一個愛神的心.
今天教會有很多有需要的人.
我們願不願意關心他們?.
我們願不願意每天為他們祈禱?.
我們願不願意每天親近主.
去看聖經.
去服侍有需要的人?.
主的照明不是真的難做.
我們有一個愛神的心.
願意去服侍主.
去服侍人.
彼得聽完主的重新任命後.
其實彼得不多對白.
只說主啊.
是的.
你知道我愛你.
有三次表達祂愛主的機會.
而且還是公開的.
其實讓彼得再一次受到肯定和醫治.
但原來主對彼得不單這麼簡單.
給他機會這麼簡單.
原來主想彼得想清想楚.
什麼叫愛主?.

$^{681}$什麼叫跟從主?.
想清想楚.
跟從主到死.
耶穌真的有跟彼得說.
你怎樣死.
我們一起看聖經.
我們一起讀.
準備.
(主播翻譯).
第十八節開頭.
這裡說實實在在.
這是重複公式.
耶穌要說很重要的事.
祂會說實實在在.
你看四方就是了.
耶穌說得很緊要.
我實實在在告訴你們.
這句話跟彼得說緊不緊要.
很緊要.
原來是跟彼得說.
年輕是怎樣.
老了是怎樣.
彼得年輕是怎樣.
很好.
可以蓄上戴子.
隨意往來.
這裡說到他周圍的服事主.
是自由的.
但年老的時候.
要伸出手來.
別人要將你帶去.
帶你去你不願意去的地方.
有一天有人跟你說.
你老了就淒涼.
你要去不願意去的地方.
你會怎樣.
我不知道.
真的嗎.
不是吧.
不要吧主.

$^{721}$其實主也說了後果給彼得.
說給彼得聽.
你見到彼得.
曾經沒有想清楚.
口口聲聲說.
我為你捨命都可以.
原來今天主要是讓彼得想清楚.
真的.
你真的要死.
你想清楚了嗎.
你是否願意跟我.
原來早期教會的基督徒.
會用伸出你的手來.
不是一隻手.
聖經說是兩隻手.
一對手.
代表釘十字架.
是一個暗號.
原來暗示彼得要怎樣.
死去榮耀主.
彼得曾經因為.
主被捕後.
受審受刑罰.
知道耶穌可能會死.
他害怕嗎.
他很害怕.
他害怕到不認主.
他逃避.
但這次耶穌跟他說.
你要這樣.
年輕就這樣.
老了就這樣.
但你看到經文.
彼得沒有回應.
沒有跟耶穌爭論.
可不可以不要這樣.
沒有.
也沒有說.
好,我為你死.
像年輕人一樣.

$^{761}$高言大字.
耶穌不用怕.
我一定為你死.
也沒有.
你看到彼得是無台詞.
但歷史告訴我們.
彼得是用行動回應主的呼召.
相傳彼得是倒釘十字架而死.
為何要倒釘.
因為他覺得耶穌是釘十字架.
他不配和耶穌一樣.
釘十字架.
所以倒釘十字架而死.
彼得這次不再口渴.
他用行動回應.
接受跟從主.
上清上楚.
有負代價.
甚或乎連命都沒有.
但他上清上楚.
他制.
他用行動告訴主.
大家也聽過文化大革命.
文化大革命時.
有位基督徒被紅衛兵拘捕.
紅衛兵問他.
基督教是甚麼信仰.
他想傳福音.
當然不會理會他.
於是他們問他.
你信的是甚麼.
他很大膽說.
我信的是耶穌.
他問耶穌是誰.
他說.
當時他們問.
你信嗎.
他說我信耶穌.
但他們很生氣.
很生氣地對他說.

$^{801}$你現在還信耶穌.
這樣會被定為反動分子.
於是他們憤怒地問他.
你信的是誰.
很生氣地追問他.
當時基督徒心裡也很害怕.
當時的氣氛.
比我還要害怕.
一大群人追殺.
心裡想要弄死你.
但還問你信耶穌是誰.
答不答好.
心裡也很害怕.
但他們不斷問.
快點回答.
他就回答.
耶穌是我的神.
於是全場都起哄.
有些人用狼狽的說話罵他.
有些人掌摑他.
有些人用鞭子打他.
這一刻他的懼怕消失.
換來心裡很平靜和安息.
他回想起.
其實耶穌也是被猶太人.
這樣誤會和對待他.
他決定不再解釋.
因為他知道這是神給他的十字架.
他覺得沒有需要作任何解釋.
來減輕受苦.
因為他知道主陪著他.
揹著這個十字架.
今天什麼叫跟從主呢?.
其實跟從主.
我們要想清想楚.
要有心理準備.
因為主在十字架上受苦.
付上很多代價.
今天信主不是教人.
吹吉避凶.

$^{841}$無災無病.
只有祝福的信仰.
不是.
今天我們的信仰要我們.
學習愛神愛人.
當我們愛神愛人的時候.
我們可能會面對困難.
會面對人種種的誤會.
或種種的壓迫.
你願不願意為神付上這個代價呢?.
今天聖經這樣說.
什麼叫跟從主呢?.
我們一起讀一讀.
「凡不背著自己十字架跟從我的.
也不能作我的門徒」.
大家留意這個十字架是什麼?.
是自己.
不是別人.
是自己.
我們每一個人都有十字架.
今天什麼叫愛神呢?.
當然愛神是一種情感.
也是一種態度.
以神為中心.
也是一種行動.
今天我們願不願意.
因為愛神跟從主.
休息時間少一點來崇拜呢?.
願不願意為主出隊傳福音呢?.
願不願意每天抽時間.
靈修親近主讀經呢?.
願不願意愛一些得罪你的人.
不可愛的人呢?.
原來今天愛神不是只有一個嘴巴.
或一個承諾.
愛神是一個行動.
是我們生命的行動.
願不願意在不同的環境裡.
仍然堅持跟從主呢?.
主對每一個人都有不同的處境.

$^{881}$但主對每一個人都有他的照明.
我們都回應主在我們身上的照明.
背起自己的十字架跟從他.
就像彼得一樣.
他再一次確認他對主的愛.
縱然他知道前路危險重重.
或是沒有命.
但他仍然堅持跟從主的決心和行動.
今天主給了我們一條命.
我們每個人都有一條命.
不是要我們坐著就算了.
主要我們好好愛他.
去服侍他.
怕不怕我們都拿著主給我們的命.
好好服侍主.
還有我們這位主也是充滿接納.
和意志和鼓勵.
就像彼得一樣.
可能今天有很多事情.
覺得我真的做不到.
我有很多問題在裡面.
你把你的問題帶到神面前.
神怎樣可以.
意志興起服興彼得.
他同樣都可以意志興起服興你.
因為主對我們的愛.
就是充滿接納.
充滿提率.
也充滿期待.
期待我們去回應他的召命.
我們有些時間靜一靜.
然後我們一起祈禱.
(靜音).
我們一起禱告.
對我們有最大的接納.
雖然我們有很多軟弱和失敗.
但你仍然給我們機會.
信任我們.
呼召我們來跟從你.
主你知道怎樣去意志.

$^{921}$服興使用彼得.
主求你同樣地.
賜福給紅印街的弟兄姊妹.
求你保守我們.
有一個熱愛你的心.
愛教會的心.
在跟從主的路上.
我們不怕付上代價.
以行動的回應.
神你對我們的愛.
有決心背起十字架來跟從你.
我們這樣禱告.
是奉主名求.
阿母.
\newpage



\section{尼希米記 11:1-3}
\label{sec:5F_Xzyjh7_k}
\textbf{2023.07.09 《幸福的感覺從行動開始》尼希米記 11\_1-3 (和修版) 曾裔貴牧師}
\newline
\newline
連結: \href{https://youtube.com/watch?v=5F-Xzyjh7-k}{\texttt{https://youtube.com/watch?v=5F-Xzyjh7-k}} ~~~~ 語音日期: 2023-07-11
\newline
\newline
\hyperref[sec:qziNPLIfqJE]{\small{< < < PREV SERMON < < <}}
~
\hyperref[sec:index_chronic]{\small{[返順時目]}}
~
\hyperref[sec:index_scriptual]{\small{[返順卷目]}}
~
\hyperref[sec:pH3FS1HN2vk]{\small{> > > NEXT SERMON > > >}}
\newline
\newline
尼希米記 11:1-3
\newline
\begin{longtable}{cl}
\hline
\hline
章節 & 經文 (和合本修訂版)\\
\hline
11:1 & \begin{tabularx}{0.7\textwidth}{X} 百姓的領袖住在耶路撒冷。其餘的百姓抽籤，每十人中選一人來住在聖城耶路撒冷，另外九人住在別的城鎮。 \end{tabularx} \\ \\ \relax
11:2 & \begin{tabularx}{0.7\textwidth}{X} 凡甘心樂意住在耶路撒冷的，百姓都為他們祝福。 \end{tabularx} \\ \\ \relax
11:3 & \begin{tabularx}{0.7\textwidth}{X} 以色列人、祭司、利未人、殿役和所羅門僕人的後裔都住在猶大的城鎮，各在自己城內的地業中。本省的領袖住在耶路撒冷的如下： \end{tabularx} \\ \\
[1ex]
\hline
\hline
\end{longtable}
$^{1}$很感恩再次跟大家將神的話語彼此勸勉.
剛才的聖經很簡單.
當然你要知道背後的背景.
你就會明白.
以色列人當時的內心感動和感受.
這段聖經告訴我們.
有十分之一被抽籤,抽中的人.
他們好像抽中居屋一樣.
很開心能夠住在耶路撒冷.
他們能夠以住在耶路撒冷為榮.
甚至所有的以色列人.
為這些被抽中的人祝福.
這就是這段聖經的簡單意思.
我們一起祈禱.
多謝主.
我們能夠再次用聖餐.
去紀念十價的恩情.
這個不能改變的事實.
造物主來到這個世界.
將我們拯救.
認識我們.
用永恆的愛愛我們.
多謝主.
這份愛將我們的生命改變.
弘恩加呈現在你的面前.
讓我們能夠深深地感受.
我們的幸福是來自神的.
我們這樣感恩禱告.
奉耶穌基督的名.
阿們.
在你的人生中.
有沒有曾經以居住在某一個地區為恥辱.
不敢告訴別人.
你住在那個地區.
有的可以舉手.
我是其中一個.
我在中二的時候.
因為要填手冊.
去校務處.
我們的學校的班底很低.

$^{41}$班五.
同學們都很頑皮.
去校務處交手冊.
不經意被同學看到我的地址.
笑得很離譜.
從那一刻開始.
中二.
我就很生氣.
我很怕別人知道我住在哪裡.
因為我住在九龍塘.
同學看到.
哇 九龍塘.
因為我的學校在九龍塘牛津道2A.
就是Marylock對面.
牛津道球場旁邊的學校.
現在變成幼稚園了.
那間中學.
叫做無犯中學.
班底很低.
全香港首三位.
最低最壞的.
隔著威靈頓就到無犯了.
真的.
可能在座的人會笑.
可能讀過了.
在感受有.
原來Hector和Irene是讀格字的.
不知道是格字還是威靈頓.
所以一聽就很有共鳴感.
我回頭說.
我住在九龍塘旁邊的木屋區.
叫做小西湖村.
就是現在的柔一城.
當時因為建地鐵之後.
就要整個村拆掉.
但還沒拆之前.
我就在那裡長大了.
你想想一個中二.
大概十二十三歲.
我很怕別人看到我的地址.

$^{81}$一看到他們就在那裡狂笑.
九龍塘小西湖村.
其實那區沒什麼好笑的.
那時候家家戶戶.
要不就住廉租屋.
要不就住摩區.
很少人會住豪宅.
但是有.
我班有個同學.
有一天跟我一起坐超號車.
在牛津道九龍塘會.
坐超號車.
接著去總站下車.
接著一起走.
走著走著進了那個豪宅.
回家了.
我還要走到那個村口.
一個大的隧道.
因為有個火車橋.
一個大隧道.
都還沒到.
我住在山邊.
要走半個小時.
下了總站下了車.
由那個入口.
再要走半個小時才上到山上.
所以我很怕別人知道我住在哪裡.
這是我童年的陰影回憶.
你曾經有沒有住在哪個地區.
來到覺得是一個恥辱呢?.
這一段聖經.
其實有些背景.
大家如果你有機會.
可以多讀聖經第七章.
其實整場建好了.
尼希米領導當時猶太的領袖.
在耶路撒冷的不同的領袖.
來到去建這個城牆.
始終都建好了.
城牆完全建好了.

$^{121}$城門全部安裝好了.
為著什麼呢?.
為著要讓以色滅人.
有自己的一個地方.
分別這群人.
來分別為聖.
也意味著.
外邦人以前是自出自入的.
令到那個敬拜.
令到聖殿.
令到那些外邦人來去.
示意的踐踏的時候.
就沒有了那個敬虔.
沒有了那個敬拜.
因為外邦人不會理會安息日.
安息日就特意來做買賣.
令到很多猶太人.
心裡完全不是放在神那裡.
尼希米奉召.
按神的感動.
由很遠的波斯帝國上來.
其實他是一個高官.
在亞達斯西王的身邊.
他是一個酒政.
當時是一個很高職位的官員.
但是他有感動.
他看到一個情景.
耶路撒冷破落.
百姓受盡的欺凌.
城門被火焚燒.
沒有人理會.
沒有人能夠理會.
沒有人敢去改變.
尼希米從很遠的波斯帝國上去.
上到耶路撒冷.
重建這個城門.
重建這些城牆.
第七章.
城牆修園了.
我安了門扇.

$^{161}$守門的歌唱的和利未人.
都已經派定了.
就是說每個崗位.
在聖殿裡面.
不同的崗位.
都有人來服侍了.
各按各職的服侍了.
接著第二節.
我就派我的弟兄哈拿尼.
和仍留的宰官哈拿尼雅.
管理耶路撒冷.
因為哈拿尼雅是忠信的.
又敬畏神.
過於眾人.
我吩咐他們說.
等到太陽上升.
才可以開耶路撒冷的城門.
人上看守的時候.
就要關門上山.
也當派耶路撒冷的居民.
各按班次看守自己的房屋對面之處.
第四節.
留意這裡了.
城是廣大的.
其中的民卻稀少.
房屋還沒有建造.
在這樣的情況下.
建好城牆.
聖殿.
尼熙米按照摩西的律法.
將不同的職位.
由利未人來經過訓練.
經過調查.
他真是屬於利未人的家譜.
將他們擺在他應該服侍的.
意味著什麼呢?.
首先要恢復的.
是聖殿的秩序.
首先要恢復的.
就是在聖殿裡的敬拜.

$^{201}$讓百姓慢慢恢復一個對神敬拜的習慣.
由這個習慣慢慢調整心態.
對神有一個敬畏的心靈.
這樣來守節.
這是尼熙米做好城牆之後.
原來只是一個過渡性的工程.
真正的工程.
不是硬件的城牆.
真正的工程.
是要開墾百姓.
已經沉睡很久的心靈.
對神的敬拜.
是要讓聖殿有了秩序之後.
這會帶來什麼後果呢?.
什麼正面的後果呢?.
就是當百姓開始發現.
原來你是屬神的子民.
你是神所立約的子民的時候.
外邦人在聖殿裡.
如果他們不根據那些律法來歸入猶太教.
他們是城牆以外的人.
是無份與神無份.
與神無關.
他們只能夠在外圍.
是外邦的人.
這個分別.
這樣的安排.
是至關重要的,弟兄姊妹.
因為由第一批.
稱為所羅巴伯的領袖.
和當時的大祭司約書亞.
在古列王的元年.
想重建的時候.
一路建好了聖殿.
一路建好了很多其他的設施.
但是一直以來.
那個城牆是破落的.
沒有人有那個能力.
去將所有的資源和人心.
來集中去建造.

$^{241}$帶來的就是.
與外邦人有很多瓜葛.
外邦人對猶太人的敬拜信仰.
看不起.
嘲笑.
很多時候用他們的辦法.
使猶太人難堪.
守節.
沒有那個心.
安息日.
不守安息日.
隨便買賣.
然後帶來的就是.
很多時候與外邦的人通婚.
出現的後果就是.
下一代,下一代.
來到尼希米的已經是第三代了.
過了一百二十多年了.
你可想而知.
對神更加模糊了.
對神的敬拜更加的求其了.
但是.
在這個時候.
尼希米做了一個很重要的工作.
就是首先將城牆做好.
城門也來修好.
然後有人來看守城門.
就是說不再可以隨便.
讓外邦人隨隨便便進來.
好像武演雞籠一樣.
已經不行了.
在這個時候.
要帶來一個很重要的信息.
就是什麼呢?.
就是帶一個幸福感給百姓.
每個人都知道自己的身份證就是猶太人.
誰是支派.
每個人都知道.
做了人口普查.
但是.

$^{281}$知歸知.
你能不能感受到.
這個聖城裡面有一份福氣.
有一份恩典.
你能夠感受到你是屬神的.
這一份榮耀感.
尼希米就來到這裡.
就做這個工程了.
硬件工程.
最多是三十多度.
暑熱天氣.
要休息一下而已.
都容易.
做著做著都會做成.
52天做完.
這個工程.
但是.
人心的工程.
幸福感的工程.
榮耀感的工程.
就不是簡單的.
所以你看到第七章.
建好城牆之後.
已經發現了.
耶路撒冷人稀少.
其實你一看就看得到.
每個人做完之後.
走了.
下班了.
建好了.
就回到自己的家裡.
就住在猶太國處.
所以剛才說的第四節就是說.
城是廣大的.
其中的民卻是稀少.
房屋沒有人建造.
然後第五節.
我的神就感動我的心.
去朝聚貴就官長和百姓.
照家譜來計算.

$^{321}$然後他就開始數點那個民數.
計算那個家譜.
你是屬於什麼支派.
你是屬於什麼類別的.
等等.
最終其實.
想帶出一個.
不是想讓你很煩.
這麼多的數字.
這麼多的人名.
不是想這樣.
是想告訴你.
他要找回哪些.
真的屬神的利未人.
他要找回一些.
哪些真的可以來事奉神的人.
在聖殿裡面.
來慢慢恢復那個秩序.
有回聖殿敬拜的秩序.
有回我們每一個月.
逢第二週.
我們有這樣的聖餐.
讓你慢慢慢慢深化.
你對神的愛子.
這個犧牲的愛.
你慢慢總是會得到一些.
總是會累積了.
神在你生命裡面的那個拯救.
這樣的恢復.
是利希米來做的.
所以很難怪.
我們沒時間讀的第八,第九章.
原來是一個復興會.
七月初一.
就是以色列人的摧國節.
利希米和以色列.
就在一個水門的地方.
窩吞基這個地方.
來召聚那些人.
全民聚集一起.

$^{361}$就聆聽摩西的律法書.
然後再分小組利未人的研習.
因為第七章下半段.
他已經數點了哪些具備.
這個資格的利未人.
就是說找回那些現在可以即刻事奉的人.
來參與他們這個心靈復興的工程.
這樣做了這個工作之後.
去到很戲劇性的第十章.
其實那裡是一個月的.
因為七月十五就是他們的駐彭節.
守完駐彭節之後是七天的.
去到七月二十二.
這是第九章的經文.
就有一個很奇妙戲劇性的工作.
就是全民願意簽約.
遵守摩西的律法書.
願意來到在神的面前.
照摩西的律法書的教導.
來生活.
這樣的一個奮興會帶來的.
原來是人的心.
百姓的心被挑釁.
挑釁什麼?.
以自己是神百姓.
沒錯.
不是今時今日回來這個聖殿.
是我的祖宗幾代了.
在這裡居住.
我是在耶路撒冷出生的.
有些可能還沒死的.
他們就是從老遠的波斯回來的.
有這樣的人.
也有新一代在耶路撒冷出生的人.
但是不約而同地.
他們願意簽約.
第十章.
我稍微讀給大家聽.
第一節.
簽名的哈加利亞的兒子.

$^{401}$省長尼希米和西底加.
接著一直說.
首先就是祭司簽約.
利未人簽約.
接著一直簽下去.
第三十節.
詳細讀得太多了.
二十九節.
都隨從他們貴就的弟兄法就起誓.
必遵行神.
即他僕人摩西所傳的律法.
謹守遵行耶和華.
我們主的一切誡命.
典章律例.
並不將我們的女兒.
嫁給這地的外邦居民.
也不為我們的兒子.
娶他們(外邦的女兒).
這地的居民.
如果在安息日.
或什麼聖日.
帶了貨物和糧食來賣給我們.
我們必不買.
這就像很簡單的事.
對猶太人來說.
為什麼要看得這麼重.
要簽約.
安息日不要買東西.
我不去維康.
今天上午.
不用簽約了.
散了碗就可以去.
但是當時.
他們感受到自己是神的選民的身份.
他們覺得這件事.
就是令他們過去離開了神的原因.
所以全民族.
當建好城牆之後.
外邦人再不可以隨隨便便.
進來教會.

$^{441}$進來聖殿之後.
你想想.
如果那些師班.
全部都未信耶穌.
唱完師班之後.
去茶餐廳吃東西.
然後拿煙出來.
如果教會是這樣的.
任人隨便.
因為他的歌喉好.
他唱歌好.
所以不找他不行.
這樣就麻煩了.
所以你問任何一個師班長.
首要的不是歌喉.
是敬虔的心.
從心底裡唱頌神.
引導百姓.
引導會眾.
唱頌神.
很重要的價值觀.
所以你看到.
有了這個嚴肅會的粉輕會之後.
就是我們剛才讀的經文.
哇 李兄妹真是厲害.
還有以色列.
要用一個這樣的方式.
抽籤.
抽完籤之後.
中籤之後.
好像等於中.
不是六合彩.
等於抽到居屋的心情.
百姓為這些十分之一的百姓祝福.
你就好了.
你被抽中住在耶路撒冷.
你就好了.
可以在耶路撒冷做守衛.
看守整個聖殿.
弟兄姊妹.

$^{481}$你有沒有以自己是洪恩家的弟兄姊妹.
每個星期你能夠坐在神的殿裡面.
敬拜神為榮.
上個星期我來開堂委會.
早了.
見到兩邊都有東西吃.
我就看看.
坐在太嶼那邊.
因為那邊有鮑魚.
坐下之後.
聊聊天.
不聊又折河.
那些太嶼姊妹.
其中一位.
她說她信主不是很久.
但一信主之後.
很喜歡回教會.
每個星期都很想早點回來.
很想星期天快點出現.
這樣的一個心情.
好像對神的家的羨慕.
日於言表的一份愛.
彼此的服侍.
彼此的祝福.
這就是尼希米在第十一章.
去完結了這個嚴肅會.
不如我們玩一個抽籤遊戲.
我們看看那些人.
因為現在耶路撒冷人口稀少.
房屋都沒有人建造.
那應該怎麼做呢?.
百姓的心已經是柔軟.
已經是被開懇.
感覺自己是神選民.
已經是一個很大的福氣.
福氣之中的福氣就是住在耶路撒冷.
現在尼希米說公平一點.
抽籤.
除了祭司利未人.
一定要住在耶路撒冷.

$^{521}$現在我們十二個支派.
其中十個支派就抽籤了.
逢十抽一.
那就是奉獻的意思.
十一奉獻.
第十章已經說了.
還有一個沒有讀到.
就是大家切實遵守.
我們一定會奉獻給神.
現在再進行一個奉獻.
就是人口的奉獻.
十一章.
你想想多麼美麗的一件事.
以住在耶路撒冷為榮.
百姓首領住在耶路撒冷正常.
其餘的百姓就祭籤了.
每十個人中就使一個人.
住在聖城耶路撒冷.
你想想那個抽獎儀式.
大家喝無.
抽中我就好了.
現在宣佈.
以法蓮之派.
誰是族的.
十個支派中抽一個出來.
那個人就住在耶路撒冷.
以住在耶路撒冷為榮.
現在不是那些dirty team.
進入高度感染的房間去做護士.
現在是抽到去耶路撒冷.
弟兄姊妹.
心靈需要先有一個柔軟的心.
你看到住在耶路撒冷.
聖城是榮耀的.
所以每一個崗位的弟兄姊妹.
你已能夠服侍神為榮.
這個榮耀感.
教會是要不斷提醒.
我們需要看到.
我們能夠被神選中.

$^{561}$其實每一個都是特選的.
有一本書.
在圖書館也借到.
在公立圖書館.
很好看.
作者叫劉繼榮.
是女士.
有文采的.
有一篇文章是講他的女兒.
叫家有中等生.
什麼意思呢.
他不知道為什麼.
他的女兒如何讀書.
如何勤力.
每次考試都是第23名.
全班是第50名.
他就是第23名.
每次都是.
父母當然不想.
都是有知識的份子.
當然想他的女兒讀得好一點.
可能我們沒有勸他.
幫他補習.
雞精斑等等.
他很大壓力.
讀到發高燒.
出來的成績都是23名.
接著都不行.
再推.
這次不行了.
他的成績差到33名.
父母終於放棄了.
有時會跟女兒說.
為什麼你不是神童呢.
如果你是神童就好了.
女兒是很高EQ的.
爸爸為什麼我不是神童呢.
因為你不是神父.
我當然不是神童.
令到全家都覺得很好笑.

$^{601}$但很無奈.
因為都是望子成龍的文化.
台灣人.
最慘就是過年過節.
女兒你是讀什麼科的.
將來有什麼志願.
每個表兄弟姐妹都說自己的宏願.
到女兒問他想做什麼.
我長大後想做幼稚園老師.
可以跟小朋友唱歌.
就很開心了.
第二志願呢.
第二志願就是想做家庭主婦.
我要打井井有條的家庭.
爸爸在很多親戚面前.
暗暗不敢說.
因為每個人都不想做醫生.
就只好放棄了.
有一天.
你回去看這篇文章寫得很好.
班主任召見這對父母.
很奇怪.
又不是考第一.
有什麼可以召見呢.
父母去到班主任面前.
班主任說.
我教書這麼多年.
我沒見過這樣的情況.
就是你女兒.
他們很緊張.
是不是我女兒有什麼事.
肯定成績不好.
老師第一句就說.
不是她的成績特別好.
都是23名.
不過有一件很奇怪的事發生.
就是考試老師班主任.
給了一個問題給全班同學.
就是請你選一個.
全班裡面你認為最好的學生.

$^{641}$49個人選劉繼榮的女兒.
要說原因.
樂於助人.
真誠.
各方面的待人接物的禮貌.
還有很有領導的才能.
這樣.
他的父母終於明白.
他的女兒有一個這麼獨特的親和力.
可以令到每一次她出現的地方.
帶來歡樂.
帶來的緊張都可以緩和.
親愛的弟兄姊妹.
其實你如果慢慢感受到.
自己是神兒女的幸福感.
其實你就可以製造這種歡樂.
製造這一份和睦.
製造這一份愛.
讓人感受到這個地方有神.
有神的愛.
有盼望.
有平安.
有接納.
哥林多前書第十六章.
史奴保羅在最後說了一句話.
第十五節.
大家今晚回去看.
你們知道.
斯提反一家.
是阿該亞省初結的果子.
即是初信主的媳婦.
接著那句話.
由史奴保羅的口說出來.
他們一家人.
是專以服侍聖徒為念.
專心以服侍聖徒.
作為他們信主之後的生命.
那個付出.
親愛的弟兄姊妹.
我們未必是很出類拔萃.

$^{681}$每個人都要很厲害.
但是每個弟兄姊妹.
洪恩家的弟兄姊妹.
每一個都是為神所愛.
每一個都是主所結的果子.
都是我們的牧者.
過去悉心的栽培.
來成為今天的你.
你坐在教會裡.
其實好像我們已經有一個分別為聖.
這個時刻.
這個地方就是屬神的.
你就是被神揀選的.
你的生命在神的眼中.
是寶貴的.
所以很盼望這段聖經.
第二節.
凡是甘心樂意住在耶路撒冷的.
百姓都為他們祝福.
是正面的祝福.
不是說多謝你了.
你走了我們可以住在外面.
不是的.
是為他們祝福.
弟兄姊妹感受到.
在神的家裡的.
這份屬靈的福氣.
感受到.
你是有這份的幸福榮耀感.
以致我們能夠成為別人的祝福.
我們一起祈禱.
多謝主.
多謝你藉著耶穌基督的犧牲.
無罪的為我們這些罪人犧牲.
使我們今天在神的面前.
成為神的兒.
幫助我們每一位明白所信.
又真的得著這份寶貴的救恩.
帶來的福氣.
榮耀.

$^{721}$我們成為這個社區.
洪水橋社區的祝福.
奉主耶穌基督的名.
阿們.
\newpage



\section{路加福音 8:4-15}
\label{sec:pH3FS1HN2vk}
\textbf{2023.07.16 《生命豐盛只因有......》 路加福音 8\_4-15 (和修版) 講員 : 曾敏傳道}
\newline
\newline
連結: \href{https://youtube.com/watch?v=pH3FS1HN2vk}{\texttt{https://youtube.com/watch?v=pH3FS1HN2vk}} ~~~~ 語音日期: 2023-07-17
\newline
\newline
\hyperref[sec:5F_Xzyjh7_k]{\small{< < < PREV SERMON < < <}}
~
\hyperref[sec:index_chronic]{\small{[返順時目]}}
~
\hyperref[sec:index_scriptual]{\small{[返順卷目]}}
~
\hyperref[sec:Jxn7bTnB5O8]{\small{> > > NEXT SERMON > > >}}
\newline
\newline
路加福音 8:4-15
\newline
\begin{longtable}{cl}
\hline
\hline
章節 & 經文 (和合本修訂版)\\
\hline
8:4 & \begin{tabularx}{0.7\textwidth}{X} 當一大群人聚集，又有人從各城裡出來見耶穌的時候，耶穌用比喻說： \end{tabularx} \\ \\ \relax
8:5 & \begin{tabularx}{0.7\textwidth}{X} 「有一個撒種的出去撒種。他撒的時候，有的落在路旁，被人踐踏，天上的飛鳥又來把它吃掉了。 \end{tabularx} \\ \\ \relax
8:6 & \begin{tabularx}{0.7\textwidth}{X} 有的落在磐石上，一出來就枯乾了，因為得不著滋潤。 \end{tabularx} \\ \\ \relax
8:7 & \begin{tabularx}{0.7\textwidth}{X} 有的落在荊棘裡，荊棘跟它一同生長，把它擠住了。 \end{tabularx} \\ \\ \relax
8:8 & \begin{tabularx}{0.7\textwidth}{X} 又有的落在好土裡，生長起來，結實百倍。」耶穌說完這些話，大聲說：「有耳可聽的，就應當聽！」 \end{tabularx} \\ \\ \relax
8:9 & \begin{tabularx}{0.7\textwidth}{X} 門徒問耶穌這比喻是甚麼意思。 \end{tabularx} \\ \\ \relax
8:10 & \begin{tabularx}{0.7\textwidth}{X} 他說：「神國的奧祕只讓你們知道，至於別人，就用比喻，要他們看也看不見，聽也不明白。」 \end{tabularx} \\ \\ \relax
8:11 & \begin{tabularx}{0.7\textwidth}{X} 「這比喻是這樣的：種子就是神的道。 \end{tabularx} \\ \\ \relax
8:12 & \begin{tabularx}{0.7\textwidth}{X} 那些在路旁的，就是人聽了道，隨後魔鬼來，從他們心裡把道奪去，以免他們信了得救。 \end{tabularx} \\ \\ \relax
8:13 & \begin{tabularx}{0.7\textwidth}{X} 那些在磐石上的，就是人聽道，歡喜領受，但沒有根，不過暫時相信，等到碰上試煉就退後了。 \end{tabularx} \\ \\ \relax
8:14 & \begin{tabularx}{0.7\textwidth}{X} 那落在荊棘裡的，就是人聽了道，走開以後，被今生的憂慮、錢財、宴樂擠住了，結不出成熟的子粒來。 \end{tabularx} \\ \\ \relax
8:15 & \begin{tabularx}{0.7\textwidth}{X} 那落在好土裡的，就是人聽了道，並用純真善良的心持守它，耐心等候結果實。」 \end{tabularx} \\ \\
[1ex]
\hline
\hline
\end{longtable}
$^{1}$弟兄姊妹平安.
今日在打風的日子.
仍然能夠返到神的家.
是上帝很大的恩典.
我們向神獻上感恩.
我們一同禱告.
所以你要設下你的華語.
成為我們腳前的燈 路上的光.
天父求你今天給我們每一個.
眼睛都張開一張.
耳朵張開.
要聽上帝你的聲音.
願意聖靈在我們每一個心裡工作.
是對我們每一顆心靈說話.
日後聖靈感動我們每一個.
回應上帝你的教導.
奉耶穌基督的名字祈禱.
阿們.
今日這個比喻.
我相信大家有機會在不同的場合聽過.
是很出名的殺種比喻.
當我們再聽的時候.
再看的時候.
仍然有很多值得我們思考的空間.
今日我的題目是.
生命豐盛只因有.
有什麼呢?.
神的心意是很想.
我和你有一個豐盛的生命.
這是神的心意.
不只是想某一個 某幾個.
是神渴望今天坐在台下.
每一位紅人家弟兄姊妹.
都有一個豐盛的生命.
豐盛不只是滿的簡單.
是滿到寫.
被上帝的恩典充滿.
滿有上帝賜給我們的.
各種屬靈的.
屬身體各方面的福氣.

$^{41}$今日我們問問自己.
我的生命是否豐盛的生命呢?.
如果是 因為什麼而豐盛?.
如果不是.
我仍然覺得我坐在這裡.
整個生命很貧乏 很缺乏.
是因為什麼呢?.
盼望今日這一堂.
給我們一些思考.
生命豐盛是因為什麼?.
因為有主和祂的道.
今日我們看這個殺種的比喻.
是其中一個很有特色的比喻.
這個比喻不只是耶穌說的比喻.
祂自己亦親自解釋的比喻.
所以我們會覺得不難明白.
每一個人看的時候都明白.
但明白之後.
更要進一步想想.
如何轉化成為我生命的回應.
有行動可以活出來.
這個反應是我很希望.
今天我們聽完這堂道之後.
是有這樣的一個回應.
生命豐盛是因為主和祂的道.
原來主的道可以在人的生命裡發芽生長.
是有影響力的.
是帶給我們很大的改變.
主的道可以在我和你的生命結出果實.
人們是可以透過我們的生命.
看到我們是一個豐盛的生命.
因為我們結出果實是可以看到的.
我們的生命是可以被人看到的.
我們仍然要檢視.
我們和你的生命是否有這樣的果實.
可以被人看到呢?.
但你會發現.
不是每個人都有同樣的經歷.
換言之不是每個人的生命.
被人看到是有這麼豐富的.

$^{81}$豐盛的果實.
為什麼呢?.
我們看看有幾種情況.
有三種人是沒有辦法結出果實的.
非常可惜.
這個絕對不是上帝的心意.
看回這個殺種的比喻.
這個比喻其實有另外兩本書都有提及.
今天我選擇了盧格夫音這本經文.
因為我覺得裡面有些內容.
是更加容易讓我們捕捉到.
能夠轉化成為行動.
另外馬可夫音的四章一到九節.
還有馬太福音十三章都有提及這個比喻.
有一個殺種的出去殺種.
在經文後面有提及.
這個種子是什麼呢?.
就是神的道.
這個種子一定是好的.
無論這個種子殺出去.
在什麼情況下.
這個種子都是好的.
因為是神的道.
神的道一定是好的.
神的道本來的效果是可以結出果實的.
讓人生命得到豐盛.
讓人生命得到很大很大的福氣.
這是神的道本身有的果效.
所以這個種子是沒有問題的.
我們不需要研究.
這些種子肯定是有問題.
那些種子又不行.
所以效果就不同.
絕對不是.
不是種子的問題.
神的道是沒有問題.
有問題的是什麼呢?.
為什麼有三種人不能夠結果子呢?.
第一種就是有些人落在路旁.
被人踐踏.

$^{121}$天上的飛鳥又來把他吃掉了.
為什麼呢?.
後面經文本身有解釋.
這個種子就是神的道.
落在路旁的人聽了道之後.
魔鬼就來了.
從他們心裡把道奪去.
以免他們信了得救.
這究竟是什麼情況呢?.
這在路旁的道是指什麼呢?.
這個路旁是一個什麼處境呢?.
其實原來路旁的地方是很硬的地.
重點是種子落在哪裡呢?.
落在路旁.
路旁的特色是什麼呢?.
硬的.
完全不明白神的道.
不明白之餘.
最主要是心很硬.
你跟他說什麼都好.
他都不明白.
他都不會讓你進去.
但有時你有沒有發覺呢.
其實我們聽講道.
我們每個人都有試過這樣的情況.
我們聽講道有個前設.
我來到教會裡.
我很期待.
很期待你講一些我很想聽的東西.
當你說一些我不想聽的東西.
我會自動塞耳朵.
就算你說多少都好.
我都不會聽進去.
我們有時看聖經也是這樣.
選一些.
哪些看上去對我來說是好的.
對我的心會舒服一點.
我喜歡選那些經文.
就選那些經文來聽.
來看.

$^{161}$其他的我塞住了.
完全不讓進去.
神的道完全沒辦法進去.
在這樣的情況下.
撒旦太容易下手.
馬上就把他拿走了.
其實撒旦的工作一點都不稀奇.
在創世記的時候發生什麼事.
大家記不記得.
由起初神創造天地.
然後在最後創造人類的時候.
其實你看到撒旦已經很想工作.
神很渴望和人建立一個很親密的關係.
神創造整個世界.
是為了要把整個世界給予人.
賜給人作為一份禮物.
祂想人去享受世界.
想人在裡面生活得豐盛.
和神之間有一個很甜蜜的.
親密的關係.
這是神的原意.
但從一開始的時候.
撒旦就要來破壞神的工作.
撒旦就是引誘人犯罪.
要破壞人和神的關係.
結果人抵受不住誘惑.
人真的犯罪.
犯罪之後和上帝關係破裂了.
馬上苦難與死亡進入世界.
你記得創世記的時候是很慘的.
這一幕是很悲慘的.
從一開始撒旦就要破壞上帝的作為.
所以神工作的時候.
同時撒旦也會工作.
有時候大家有沒有留意.
特別是我們出隊.
傳福音報道的時候.
總是有一種禱告.
禱告求神今天捆綁惡者.
讓他不能搞砸我們.

$^{201}$破壞我們的工作.
因為當我們出隊的時候.
直接去搶救靈魂.
這個時候你會發覺.
有很大的爭戰.
要很好的禱告.
求神工作.
因為你出去搶救靈魂的同時.
當上帝要去工作的時候.
撒旦同時間也會工作.
他不會停下來等你.
不會的.
所以這裡也是這樣的情況.
當他看到人心很硬.
根本不想聽上帝聲音的時候.
對的.
說什麼都好.
馬上拿走.
沒有機會下去.
所以他根本沒有機會.
結出一個豐盛的果子.
沒有.
我們想想.
我們自己是否這樣的情況.
是不是.
是不是這種土壤.
硬心的.
隨時被撒旦影響.
第二種.
落在盆石上.
一出來就枯乾了.
因為得不著滋潤.
這個情況.
會發一點點苗.
好像比第一種好一點.
第一種沒有機會發苗.
第二種發一點點.
但是.
他還是不能好好生存.
這個就是耶穌解釋的時候.

$^{241}$在盆石上.
讓人聽到.
很歡喜靈壽.
那一下的感受是好的.
哎呀.
你真的說得好.
小祖查經.
這段經文真的很好.
聽到覺得很好.
那一下.
但是因為沒有根.
不過暫時相信.
一碰到試煉就退後了.
這種情況.
說起盆石是怎麼樣的呢.
大家看到.
其實是有淺淺的一層泥土.
這是一層很薄的土壤.
蓋著石灰岩.
岩石是硬的.
上面只有薄薄的一層泥土.
就是土壤.
就憑著薄薄的土壤.
種子下去的時候.
就隨意地發苗.
但是.
這個地方我想說.
這個地方其實是.
熱的時候非常熱.
如果大家有機會.
看看另外一些經文.
描述是很好的.
例如我們看看馬太福音.
十三章六節這樣說.
太陽出來一曬.
因為沒有根就枯乾了.
太陽一曬.
不是普通的曬.
在這些地方.
特別是雨季過後.

$^{281}$太陽曬的時候非常熱.
你想像一下.
昨天開始嘔打風.
我聽很多人跟我說.
哎呀好熱呀.
真的熱到站在街上.
不用等一會就冒汗了.
就像乾蒸一樣.
焗呀焗呀.
就是這種感覺.
是很熱很熱.
而且是很長時間.
一路曬曬曬曬曬.
然後你發覺.
那棵東西.
那個苗.
根本不夠水分.
如果要很足夠的水分.
它需要有一個很深的根.
一直伸進去.
很深的泥土裡面.
這樣去吸取養分.
吸取水分.
才可以抵受.
這麼猛烈的太陽.
根本沒有根.
很快就會枯乾死了.
這一類人的生命.
他都叫做開放多一點.
都肯聽一下.
聽下去都覺得.
神的話語很好.
神的道不錯.
但心其實還沒有徹底開放.
還沒有徹底開放.
下面還是硬的.
也沒有紮根.
去吸取水分.
這個水是什麼水呀.
這個生命的活水是指什麼呀.

$^{321}$有沒有人能回答我.
什麼叫生命的活水呀.
沒有活水的江河在裡面.
他沒有去經歷聖靈的大能.
在我和你信耶穌之後.
神就是賜給我們有聖靈的.
其實今天和我們同住的是聖靈.
有沒有靠著聖靈.
更加深去認識真理.
有沒有紮根在神的話語裡.
經常從神的話語.
得著力量.
得著養分呢.
既有神的話語.
也有聖靈的工作.
我們有沒有這樣去得著力量呢.
很明顯沒有.
沒有的話.
我們心裡不是全然開放.
也沒有力量.
沒有養分.
反而被試煉.
嚇到了.
什麼試煉呢.
例如有信仰上的逼迫.
你說香港沒有的.
哪有逼迫呢.
現在有人拿著槍.
對著你的太陽穴.
要威脅你嗎.
今天所經歷的不是這些.
如果有家人反對.
如果有身邊很親密的人.
你的配偶反對.
有很多壓力.
你工作的環境.
有很多很多的壓力.
壓迫.
我們的回應是什麼呢.
妥協 放棄.

$^{361}$是嗎.
啊 這樣嗎.
信耶穌要面對這些壓力嗎.
是嗎.
我其實在香港聽過.
不止一次的見證.
你信耶穌.
就不要回家.
你夠膽就不要回家.
你信耶穌就大了.
有飯吃.
不行.
現在公司方面.
要求我要有這個有那個.
你信耶穌.
就不要來上班.
很多的.
那你怎樣.
很重要.
工作很重要.
家人很重要.
我現在不是說家人不重要.
但原來在這個位置.
你知道什麼才是你生命的神.
你最怕的就是你生命的神.
你如果最怕的是人.
就是人是你的神.
不是上帝.
我說真的.
如果這件事比上帝還要多.
那件事就是你的神.
不是上帝.
試煉來.
就會退後.
我曾經說過.
我去過埃及.
這個地方.
有八個星期的跨文化.
宣教的體驗.
這是第一次.

$^{401}$第二次再去埃及.
其實埃及這個地方.
有不少伊斯蘭教徒.
信耶穌.
但亦有不少.
忍受不了逼迫.
會背棄信仰.
面對生死的面前.
真的不行就不行.
我現在在這裡也不敢說自己那麼僵.
在面前.
我不怕死.
我不敢說.
但我告訴你.
生命根基不夠穩固的時候.
沒有養分的時候.
面對試煉.
就會退後.
那就不回教會.
是不是.
先放棄信仰.
人真的會這樣.
這些就不能結出果實.
第三種.
有些落在荊棘裡.
荊棘和他一同生長.
把他擠住.
荊棘是什麼.
荊棘就是那些人聽到.
走開後.
被今生的很多.
憂慮,錢財.
怨樂擠住.
看不出成術的紙粒.
這個落在荊棘裡的人.
沒有說他聽到之後很歡喜.
領受,剛才第二種也有說他們.
歡喜,也喜歡神的道.
覺得挺好,這些道不錯.
但說到荊棘裡.

$^{441}$沒有,聽了就聽了.
還要走開了.
然後生活裡.
很多東西將我擠住.
擠壓住.
有今生的憂慮.
我想說今生的憂慮.
我和你都有.
我不是要否認.
我們有憂慮.
我又不是說.
我們要去面對憂慮的事.
我們就很不屬靈.
多大得壓低自己的感受.
你不要覺得自己這樣.
不開心.
你要當憂慮不是憂慮.
這樣才可以.
我不是這樣說.
在憂慮面前就憂慮.
就不開心.
我也承認.
我很記得.
我去年.
六月的時候.
應該是五月的時候.
我開始知道媽媽有急性血癌.
當時一點心理準備都沒有.
在之前一年.
爸爸安息主懷.
沒有想到那麼快.
媽媽就有血癌.
是急性的.
生命已經去到盡頭.
當時因為沒有心理準備.
心情很難過.
我記得我回來教會.
媽媽是最看到的.
看到我在房間裡哭.
這就是我.

$^{481}$有沒有憂慮?.
有.
怎會沒有憂慮?.
憂慮很多.
師父很忙.
怎樣照顧他?.
他需要很多照顧.
一直很忙.
沒有時間陪他.
心裡有很多事情想.
我曾經有掙扎.
我想變成半職童工.
我想轉做半職.
好好照顧媽媽.
在她人生最後一程.
我有掙扎過.
有憂慮,有難過.
這些都是很真實的感受.
所以.
如果姐妹有憂慮.
我不會叫大家.
不要憂慮.
不要有很難過的感受.
我不會.
會不開心,痛苦.
會有眼淚.
因為我都經歷過.
但在那時候.
很重要的是.
在憂慮裡.
不要推開上帝.
有兩種做法.
很多時候.
我們是其中一種.
太憂慮.
我會自然地推開上帝.
我會自然地憂慮.
我會自然地想.
在痛苦的情緒裡.
無法逃脫.

$^{521}$很多時候是這一種.
不要忘記.
生命有主.
憂慮的時候.
要想著主.
我記得我媽媽那時候.
雖然我很痛苦.
有很多眼淚.
有一天.
我對著她.
哭得很厲害.
但我記得那段時間.
我和我姐夫.
很多禱告.
仍然要.
和媽媽一起.
捉著上帝.
我記著上帝.
痛苦的時候.
很難受的時候.
最後的日子.
都是和媽媽一起.
捉著上帝.
想著上帝.
在裡面.
神與我同在.
每晚和媽媽一起.
睡覺的時候.
都有上帝在.
記著上帝.
一直為她禱告.
她很不舒服.
我梳著她的時候.
不只是照顧她.
心裡默默為她禱告.
求上帝減輕她的不適.
一直神都在.
沒有離開過.
我記得有一段時間.
我真的覺得太痛苦.

$^{561}$太難受.
很沉默.
但有一天.
我聽到上帝喊我的名字.
叫我起身.
真的很寶貴.
神在這裡.
所以,弟兄姊妹.
我很想說一件事.
在艱難的時候.
在那些時候.
真的.
就算和弟兄姊妹說.
弟兄姊妹都很好.
陪我們走.
但未必最明白.
那種痛苦只有自己才知道.
但上帝是明白.
上帝是知道.
和與我們同行.
在那時候,在憂慮當中.
真的說一句.
不要離開上帝.
和祂一起.
這是我們的出路.
就算哭也和上帝一起.
捉住祂.
在上帝面前哭和不在祂面前哭.
是有分別的.
你只要自己哭.
和知道神在看著你.
哭是兩回事.
是吧.
就算這樣.
也可以結束過神.
生命有力.
我一直在過程中.
到媽媽終於.
在星期日的清晨.
去就安息主懷.

$^{601}$那時候是四點多.
五點.
我記得前一天.
我母會的牧師叫我請假.
明天不要講道.
你這樣的情況也不請假?.
我家人也跟我說.
我記得是.
我媽媽安息主懷的星期日.
她也跟我說不要回去.
要打點事情.
但是我很感恩.
上帝仍然與我同在.
就是那一刻.
我反而更加堅定地.
想到.
我要回來講道.
神和我一起.
給我有力.
你說我不難過.
傷心什麼?難過.
傷心.
神和我一起有力.
我沒有疑惑.
我便回來了.
之後整個.
媽媽喪禮籌備的過程.
非常順暢.
因為上帝從來沒有離開過我.
家裡的大小事情.
媽媽離開後.
家裡的很多身後事.
很多家裡的事情.
還要處理.
實在太多了.
上帝從來沒有離開過我一家.
與我們同在.
直到如今.
我見證神很信實.
我不會被憂慮.

$^{641}$壓制著我.
我不會被痛苦.
壓制著我.
弟兄姊妹我們也要這樣.
不要被憂慮壓制自己.
不要一直沉下去.
沉到永遠.
我們在說錢財.
每個人渴望的東西不一樣.
有些人很渴望錢財.
一生就圍著錢財轉.
我為了錢財轉.
什麼都不用想.
錢只掛帥.
錢重要.
我不會叫你.
一分錢都不要.
現在教會的運作.
也需要錢.
我不會這樣說.
錢不重要.
錢重要.
但不要讓錢做我們的主.
上帝才是我們的主.
不要被壓制著.
有錢財上的困難.
或者自己對錢財的渴望.
不要成為我們的偶像.
賢樂.
自從香港放寬.
防疫的措施之後.
社交距離.
放寬之後.
我們要上上去吃.
吃是開心的.
我也喜歡吃.
你和我說吃.
我也可以和你說很多.
我不會覺得我們不要吃.
但不要讓這件事.

$^{681}$成為我們的偶像.
我最誇張的.
我聽過一個人.
吃到什麼.
吃到破產.
你沒聽過.
最誇張的.
吃到破產.
很喜歡吃.
什麼都說吃.
到處去吃.
真的是一個很大的索菱捆綁.
吃也可以成為索菱捆綁.
很省錢.
吃看你怎樣吃.
聖經說.
不要被賢樂壓制著.
結不出成熟的紙粒.
你會發覺你整天想著.
吃啊吃啊.
去哪裡吃和誰吃.
怎樣吃.
你小心.
那些東西成為我們的偶像.
賢樂也可以.
絕對可以.
結不出豐盛的果子.
是很可惜.
可惜太輕微.
其實這段比喻.
說出來的時候.
是一個審判的訊息.
耶穌在這個比喻當中說.
有耳聽的就應該聽.
我們每個人都有耳.
有耳就要聽.
很嚴重.
其實是一個審判.
這幾種情況.
是不結果子.

$^{721}$不是隨便.
你喜歡就不結果子.
不是.
是審判.
要知道.
這樣的結果.
上帝絕對不想我們這樣.
絕對想我們悔改.
有種情況.
一定可以結果子.
就是落在好土裡.
生長起來.
結實百倍.
我問問弟兄姊妹.
剛才聽了三種情況.
是可以結果子.
馬太呼音說得很好.
十三章八節.
有的落在好土裡.
就結出果實.
有百倍的 有六十倍的 有三十倍的.
有耳的.
就應當聽.
原來最近有些研究.
說出當時.
一般農作物的收成.
農作物的收成.
大概只有四五倍.
好的收成是四五倍.
現在馬太呼音說的是.
三十倍的.
六十倍的 一百倍的.
哇 簡直是超班.
原來在神的裡面.
那個結果子.
那個豐收是怎樣.
是很超班的.
很超越的.
超過正常一般的四五倍.
可以這麼豐富.

$^{761}$這才叫豐盛.
不只是多 是豐盛.
可以結果子這麼多.
那什麼叫好土呢.
什麼叫好土.
有沒有人嘗試回應.
我想大家有些回應.
聽得到去神.
聽得到首先要.
明白.
我要去思考.
入心.
我生命是開放的.
不是聽著.
心裡很硬.
我不會讓你的壞語進來.
或者直接 我不會聽你說話.
不是這種 是開放.
聽進去 柔軟的.
在裡面.
一定可以吸收到水分.
你被根子伸進去.
很深 吸收到足夠的水分.
可以經歷聖靈的工作.
有神的話語在我們裡面工作.
是可以結出.
很豐盛的果實.
來到這裡有個重點.
我很想說 盧家呼音為什麼我要選這段經文呢.
因為它說了一個很大的重點.
結果子的生命.
一個豐盛的生命 不是一天半天的.
這裡怎麼說.
盧家呼音說.
那樂在好土裡的.
就使人聽了動.
並用純真善良的心.
持守他.
持之以行.
換言之.

$^{801}$整個結果實.
是一個過程.
不是一個短的過程.
是很一段時間的.
生命為什麼可以結出果實呢.
因為聽到.
是會回應的.
你的心是柔軟的.
開放的 被神工作的.
被神的話語在你裡面工作.
你去依靠聖靈.
你的生命是一個活的生命.
而且是持守的.
不是一天半天.
一天 兩天 三天.
一個星期 兩個星期 一年 兩年 三年.
你信主十年.
不斷的持守.
對上帝的忠心.
持守對上帝的忠心.
而且是用純真善良.
純真是什麼.
你對神的心是很純的.
很馴服的.
不是一聽到.
馬上起降.
我馬上拿刀插你.
我和你.
對著幹.
你是很馴服的.
你聽了 聽進去.
你也很有信心.
你信了神的話語.
不是只禱告.
我們腳前的燈 路上的光.
最常這樣禱告.
你真的信了 神的話語指引你一生.
神的話語帶你.
一直走向豐盛.
神的話語給你力量.

$^{841}$你信了 你也馴服.
持守.
一年 兩年 三年.
直到見主面.
還要耐心等候.
是果實.
耐心.
人生會有很多風浪.
很多苦難.
不是只有這樣.
天天都風調雨順.
天天都沒有困難.
我和你嘻嘻哈哈.
有空去德福喝茶.
去玩.
我們弟兄姊妹也有提議.
你什麼時候搞一天旅行團.
我們也很想.
和弟兄姊妹開心.
這些事是開心的.
但我們生活不是只討論這些.
我們討論生命的果實.
是耐心的.
面對很多苦難.
我反而看到.
洪恩和弟兄姊妹經歷的是什麼.
我們有很多苦難.
各式各樣的苦難.
身體的.
家庭的.
孕制的.
我們當中有很多苦難.
在這個苦難裡.
如何耐心等候.
結果實.
這就是.
我和你要去經歷的.
我們一起經歷.
我相信.
人生每隔一陣子.

$^{881}$有不同的眼淚.
但你記住.
在眼淚裡.
有更多上帝的恩典.
上帝的同在.
和保守.
這種人聽了道.
有信服.
有信靠.
有持守.
有耐心等候.
就會結果實.
生命就會豐盛.
有一天我和你要見主.
弟兄姊妹.
我們回去見主時.
是怎樣的狀態呢?.
是很貧乏的生命.
還是很豐盛?.
上帝要稱讚我和你.
嘩!厲害呀!.
良善有忠心的僕人.
你的生命活得精彩.
活得豐盛.
是哪一種?.
盼望我和你有一天.
我們都無愧於心.
在地上生活時.
事奉時.
活得豐盛.
這是非常寶貴的.
很多時候.
有很多宣教士的見證.
傳道人的見證.
他們的見證很感人.
不是因為他們實在太順利.
反而是因為他們人生很不順利.
很多苦難.
但在苦難中.
卻能夠捉住神.

$^{921}$不會離開神.
信服 信靠上帝.
讓神的話語在他們生命中工作.
又去倚靠聖靈.
他們的生命.
有很美的果子.
是這些傳道人.
這些宣教士.
他們的生命.
給你看到有很美的見證.
我很盼望有一天.
今天在座的這麼多人.
一起擁有豐盛生命的人.
我們彼此互勉.
我們一同禱告.
這個殺種的比喻.
看到.
其實上帝很渴望.
我們能夠過刀結果子.
有三十倍的.
六十倍的 一百倍的.
因為你是栽種的人.
你是殺種的人.
上帝你看著我們的時候.
等著我們結出果實.
但今天我們可能.
都有不同的情況.
我們有心硬的.
我們沒有讓上帝.
在我們裡面工作.
我們沒有真正的.
信服 信靠你.
沒有讓上帝你的話語.
在我們心中工作.
還有我們被.
世上很多的憂慮.
很多錢財.
很多的事情.
壓住了我們.
以致我們不能夠結出.

$^{961}$豐盛的果實.
天父求你憐憫我們.
你也是寬恕我們.
你幫助我們.
無論我們今天.
在我們生活環境裡.
怎樣都要抓住你.
都要依靠你.
以你為王為主.
讓你的話語在我們心中工作.
我們信服 信靠.
盼望上帝透過我們.
榮耀你的名.
主啊我深信你對我們.
洪恩家弟姐妹有你的美意.
你是希望我們每一個.
都能夠結出果實.
豐盛精彩.
求主施恩憐憫.
奉耶穌基督的名字.
祈禱 阿們.
謝謝大家.
\newpage



\section{歌羅西書 2:1-19}
\label{sec:Jxn7bTnB5O8}
\textbf{2023.07.23 《打穩根基》 歌羅西書2\_1-19 (和修版) 謝靄薇牧師}
\newline
\newline
連結: \href{https://youtube.com/watch?v=Jxn7bTnB5O8}{\texttt{https://youtube.com/watch?v=Jxn7bTnB5O8}} ~~~~ 語音日期: 2023-07-25
\newline
\newline
\hyperref[sec:pH3FS1HN2vk]{\small{< < < PREV SERMON < < <}}
~
\hyperref[sec:index_chronic]{\small{[返順時目]}}
~
\hyperref[sec:index_scriptual]{\small{[返順卷目]}}
~
\hyperref[sec:5NHMIEcWoqY]{\small{> > > NEXT SERMON > > >}}
\newline
\newline
歌羅西書 2:1-19
\newline
\begin{longtable}{cl}
\hline
\hline
章節 & 經文 (和合本修訂版)\\
\hline
2:1 & \begin{tabularx}{0.7\textwidth}{X} 我要你們知道，我為你們和老底嘉人，和所有沒有與我見過面的人，是何等地勤奮； \end{tabularx} \\ \\ \relax
2:2 & \begin{tabularx}{0.7\textwidth}{X} 為要使他們的心得安慰，因愛心互相聯絡，以致有從確實了解所產生的豐盛，好深知神的奧祕，就是基督； \end{tabularx} \\ \\ \relax
2:3 & \begin{tabularx}{0.7\textwidth}{X} 在他裡面蘊藏著一切智慧和知識。 \end{tabularx} \\ \\ \relax
2:4 & \begin{tabularx}{0.7\textwidth}{X} 我說這話，免得有人用花言巧語迷惑你們。 \end{tabularx} \\ \\ \relax
2:5 & \begin{tabularx}{0.7\textwidth}{X} 雖然我身體不在你們那裡，心卻與你們同在，很高興見你們循規蹈矩，對基督的信心也堅固。 \end{tabularx} \\ \\ \relax
2:6 & \begin{tabularx}{0.7\textwidth}{X} 既然你們接受了主基督耶穌，就要靠著他而生活， \end{tabularx} \\ \\ \relax
2:7 & \begin{tabularx}{0.7\textwidth}{X} 照著你們所領受的教導，在他裡面生根建造，信心堅固，充滿著感謝的心。 \end{tabularx} \\ \\ \relax
2:8 & \begin{tabularx}{0.7\textwidth}{X} 你們要謹慎，免得有人用他的哲學和虛空的廢話，不照著基督，而是照人間的傳統和世上粗淺的學說，把你們擄去。 \end{tabularx} \\ \\ \relax
2:9 & \begin{tabularx}{0.7\textwidth}{X} 因為神本性一切的豐盛都有形有體地居住在基督裡面； \end{tabularx} \\ \\ \relax
2:10 & \begin{tabularx}{0.7\textwidth}{X} 你們在他裡面也已經成為豐盛。他是所有執政掌權者的元首。 \end{tabularx} \\ \\ \relax
2:11 & \begin{tabularx}{0.7\textwidth}{X} 你們也在他裡面受了不是人手所行的割禮，而是使你們脫去肉體情慾的基督的割禮。 \end{tabularx} \\ \\ \relax
2:12 & \begin{tabularx}{0.7\textwidth}{X} 你們既受洗與他一同埋葬，也就在此禮上，因信那使他從死人中復活的神的作為跟他一同復活。 \end{tabularx} \\ \\ \relax
2:13 & \begin{tabularx}{0.7\textwidth}{X} 你們從前在過犯和未受割禮的肉體中死了，神卻赦免了你們一切的過犯，使你們與基督一同活過來， \end{tabularx} \\ \\ \relax
2:14 & \begin{tabularx}{0.7\textwidth}{X} 塗去了在律例上所寫、敵對我們、束縛我們的字據，把它撤去，釘在十字架上。 \end{tabularx} \\ \\ \relax
2:15 & \begin{tabularx}{0.7\textwidth}{X} 基督既將一切執政者、掌權者的權勢解除了，就在凱旋的行列中，將他們公開示眾，仗著十字架誇勝。 \end{tabularx} \\ \\ \relax
2:16 & \begin{tabularx}{0.7\textwidth}{X} 所以，不要讓任何人在飲食上，或節期、初一、安息日等事上評斷你們。 \end{tabularx} \\ \\ \relax
2:17 & \begin{tabularx}{0.7\textwidth}{X} 這些原是未來的事的影子，真體卻是屬基督的。 \end{tabularx} \\ \\ \relax
2:18 & \begin{tabularx}{0.7\textwidth}{X} 不要讓人藉著故作謙虛和敬拜天使奪去你們的獎賞。這等人拘泥在所見過的幻象，隨著自己的慾望無故地自高自大， \end{tabularx} \\ \\ \relax
2:19 & \begin{tabularx}{0.7\textwidth}{X} 不緊隨元首；其實，由於他全身藉著關節筋絡才得到滋養，互相聯絡，靠神所賜的成長而成長。 \end{tabularx} \\ \\
[1ex]
\hline
\hline
\end{longtable}
$^{1}$在小亞瑟亞爾弗所東邊的一個小鎮.
你會看到老底家.
原來老底家當時是金融中心.
很厲害的.
老底家和哥羅西距離19公里.
哥羅西教會不是保羅建立的.
不過他曾經由爾弗所派人來教會工作.
於是發現教會裡面出現假教師.
所以保羅覺得自己有責任.
應該要去教導他們.
雖然他自己不在那個地方.
書信裡面.
但他在我們這一章裡面.
講述了很多假教師所講的行為.
假教師會教導你.
原來我們要認識神.
我們要完全得到拯救.
還要加一些東西.
加些什麼呢?.
就是要崇拜一些靈界的領袖.
權威.
甚至天使.
這些不是我們聖經所教導的.
特別是要守禮儀.
例如國禮.
或者是要嚴守一些食物上的規例.
這些是人加上去的.
這些是一些異端的教訓.
進入了教會.
所以我們弟兄姊妹.
要在屬靈的生命中成長的話.
我們首先要懂得什麼呢?.
我們要懂得去分辨真理.
在第四節裡.
很清楚地告訴你.
免得有人用花言巧語.
花言巧語是什麼呢?.
有些人很會說話.
說得很動聽.
不過他會令你誤入歧途.

$^{41}$走了另一條路.
而不是走在真理裡面.
這個更厲害.
說的是用他的自學和虛空的廢話.
以前的翻譯是理學和虛空的妄言.
直接來說.
這些是虛空的.
和一些欺騙人的自學.
所以保羅叫他們要小心.
要謹慎.
世上的小學生說.
其實他們這些學說都是很膚淺的.
要小心.
否則就會被他們俘虜了.
離棄了真道.
同樣在第十六節裡提到.
在飲食上.
猶太人很注重節期.
在當中教導他們.
原來你們這些人.
已經得到耶穌基督的豐盛.
為什麼還要跟從他們的教訓.
去學習他們的學說.
或者是要守他們的節日.
千萬不要因為自己沒有守節日.
而覺得自己是低人一等.
而被人去論斷你.
或者審判你.
千萬不要小看自己.
在當時.
這個諾斯底主義.
在二,三世紀的時候很盛行.
所以當時也影響了哥羅西教會.
他們很注重強調知識.
所以他們認為.
神傳遞的豐盛.
是貫通在天地之間的.
但是要藉助天使.
或者一些諸靈.
即是星獸裡的靈.

$^{81}$去左右人的命運.
所以神並不是唯一對人命運的掌控者.
他們認為.
基督只是一個高級的受造者.
他們這樣看的.
所以人就要去敬奉一些星獸裡的諸神.
這些就是.
保羅所謂的.
這些是世上一些很膚淺的學說.
所以千萬不要上當.
人的傳統.
二端他們很喜歡加人的傳統.
但是人的傳統.
如果沒有真理.
沒有根基.
是不可靠的.
主耶穌曾經教訓這些法利塞人.
說他們死守一些傳統.
但是他們卻放棄了神的誡命.
還有呢.
這個還要拿什麼?.
叫你假裝謙虛.
其實謙虛本來是美德.
但反而是叫你假裝謙虛.
這些很糟糕.
還要叫你去敬拜天使.
我們只有這位真神可以敬拜的.
如果叫你敬拜天使.
真是違背神的旨意.
叫你假裝謙虛.
更加是不對的.
這些道理我們為什麼可以接受呢?.
他說.
免得奪去你的掌上.
神要給你一個生命的冠冕.
但是如果你因為這些學說.
離開了真神.
離開了耶穌基督的道理.
你就像他們一樣.
走向滅亡的道路.

$^{121}$所以20到21節還說.
還有一些不可拿,不可嘗,不可摸的規條.
他說.
有些人覺得.
我怎樣才可以完全呢?.
我怎樣去克制我的肉體呢?.
其實大家都很明白.
我們現在在地上還未得熟.
還未完全的.
所以我們會經歷很多試探.
來的時候.
如果不是靠著耶穌基督.
你就會跌倒.
但是如果你鞭策自己.
折磨自己,苦待自己.
難道你這樣就可以.
讓你不犯罪嗎?.
你很明白的.
我們怎樣去克制.
我們的肉體不犯罪.
我們要貪心,貪念.
我們有很多慾望無窮的.
我們要靠著耶穌基督.
我們才可以知道哪條路是對的.
哪條路是不對的.
但是我們靠著一些修行.
靠著一些苦待自己.
這些完全是不行的.
你仍然是我行我素.
仍然會貪心.
仍然會想到很多思念是為自己的.
不是為人的.
你還是得了一筋.
除了我們懂得分辨以外.
我們還要打好筋肌.
就是要生根建造.
大家都明白.
植物要向上發展.
種花也知道.
我不是很懂得種花.

$^{161}$種的東西都死了.
要除蟲.
要施肥.
有很多事情要做.
要囂觀.
要花很多心思.
生根建造.
我們在教會裡.
我們怎樣去建立信徒呢?.
首先我們要去教導真理.
就好像我們很慣常的.
有福音聚會.
有報道會.
或者做個人報道.
這些人想主.
如以訣之想主.
但下一步我們要跟進他.
要去教他.
令他明白真理.
令他明白為何我們要回來教會.
我在家裡不行嗎?.
我在家裡也可以祈禱.
我們要教他明白.
我們要敬拜我們的天父.
我們要以他為一.
我們這麼空閒.
為何不抽時間去敬拜神呢?.
真的要教導他.
當我們去教導他的時候.
他就明白了我們基督徒的責任.
所以保羅去教提摩太.
栽培他.
成為教會的領袖.
然後他說要他去教導其他人.
所以在提摩太後書.
二章十五節說.
你要竭力在神面前做一個.
經得起考驗無愧的工人.
按照正義講解真理的話.
原來我們去教導人的時候.

$^{201}$自己也要花功夫.
我們不可以抽一些東西.
照讀就行了.
不行的.
每個人都是一個活生生的人.
不可以照玩煮碗的.
他有他不同的個性.
他可能領悟快一點.
有些人領悟慢一點.
你要忍耐的.
但你自己要做些功夫的.
大家很明白大使命經文裡面.
也是講到教導.
所以教導是很重要的.
這裡說到.
「都教導他們遵守」.
即是你教導他們之後.
他們要懂得聽神的話.
懂得做出來.
直到什麼時候?.
直到世代的終結.
很長時間的.
不是說你現在教完他們就算了.
所以我們帶人來信主.
雖然你可能不是那個導師.
但你不能把責任全給導師.
你也要關心他.
你打電話問他.
你學成怎樣.
我們要關心他.
關心到他有一天.
他也可以成長到去教人.
那你就可以放心.
如果他還不懂得賺錢.
不懂得教人.
可能他還有很多東西.
還沒成長.
還是嬰兒.
那你還要跟著我講.
直到世代的終結.

$^{241}$所以我們的工作是要繼續的.
我們除了教導之外.
我們看到的是合而為一.
在教會裡.
怎樣合而為一呢?.
原來比特田書二章五節裡告訴我們.
原來你們都是活石.
原來我們每一個.
在靈宮裡.
都是一個活石.
那你要發揮到活石的功能.
要被建造成為聖年的殿.
成為聖傑的祭司.
祭司是事奉神的.
然後藉著耶穌基督獻上.
蒙神月立的屬靈的祭物.
所以原來我們要合而為一.
我們不可以各自走各自的.
好像你在房子裡.
耶穌基督.
我們就是耶穌基督耕種的田地.
我們就是祂建造的房屋.
所以你不可以藉口說.
教會裡少一個不少.
多一個不多.
多一個有什麼用呢?.
我又沒有用的.
但如果我們看你是活石的時候.
就好像這個建築物.
建築物裡面.
磚都要泥.
要黏在一起.
沒有泥的話.
它會掉下去.
大家都有功用的.
就像你建造的房子.
大家坐在一起.
你不要說少了一個不要緊.
少了一個就少了一塊東西.
那些人就凹下去就會掉下來.

$^{281}$所以原來神造的東西很有趣.
剛剛好.
就像你那塊磚.
去到祥國的時候.
它要批一批.
磨一磨它.
剛好那個位置就放進去.
其實我們每個人.
我們都有一些不多不少的.
神形我們叫過.
它要磨掉你的.
要剁掉它.
不合時的氣要剁掉它.
然後剛好凹下去.
就可以成為一個靈功.
給神使用的.
所以你會看到.
原來有破口怎麼辦?.
有破口.
那些異端就乘虛而入.
是不是?.
異端就好像.
剛才我們所說的經文.
是不是?.
異端不照耶穌基督.
所以你就會被它俘虜了.
如果你不按著神的話語去做.
你就跟了它.
因為它乘虛而入.
就進來教你一些不對的東西.
他們說.
這些敵對我們的.
律例上所寫的.
原來這些束縛我們的字句.
我們不要死守的.
我們要切走它.
釘在神事架上.
淘去它.
你不可以.
每個人這樣做.

$^{321}$我們就跟著做.
還有就是.
不要被人.
因為你沒有去守節旗.
就去評斷你.
你不要小看自己.
你要遵守真理所做的.
還有就是.
你要記得.
你不要跟著他們去敬拜天使.
假裝謙虛.
不然你就得不到神給你的掌賞.
第三方面.
就是講到我們的言行一致.
言行一致最重要是什麼呢?.
就是講到.
我們要去蒙照去事奉.
各位弟子們.
你想到一個.
言行一致的人.
有沒有想到你今天.
在座每一個都是一個屬靈人.
有沒有想到你是一個屬靈人.
講起字句好像很嚴肅.
但你的確是一個屬靈人.
那你有沒有發揮到.
神給你的生命.
怎樣去發揮成為一個屬靈人呢?.
你在世上.
有沒有發揮到這個屬靈人的功效呢?.
聖經教我們.
不可效法這個世界.
就好像一艘船.
在海裡.
你可以自由地駕馭這艘船.
很自由的.
但你會不會有時.
我要怎樣呢?.
不聽話.
我甘願墮落.

$^{361}$隨波逐流.
最後怎樣呢?.
就被海水淹死了.
是吧?.
你不去駕馭它.
你隨它怎樣走.
那你自己自甘墮落.
是吧?.
所以原來.
我們怎樣去發揮到.
言行一致呢?.
我們怎樣去站好我們的崗位呢?.
盡上我們自己的本份呢?.
就好像.
你在公司工作.
你可以盡責地為公司工作.
老闆不會看著你的.
不過你覺得你有一個責任.
因為你是一個基督徒.
你有一個責任.
不單是做好這份工作.
你也要盡責地交足一些東西.
但在小節上.
你也可以發揮到榮耀神的功用.
例如接電話.
接電話的態度可以很友善.
又可以在你心情不好時罵人.
但你想想.
你罵人的時候.
是會得罪別人的.
你自己的電話也不要緊.
但你代表公司的時候.
那就慘了.
是吧?.
那就不可以得罪別人.
說別人對公司的印象不好.
你聘請的接電話人員那麼差.
影響形象.
其實在小節上.
你也可以榮耀神的.

$^{401}$在基督徒.
如果你是基督徒.
在做好這份工作時.
罵人的時候.
你想想.
你怎樣去榮耀天父呢?.
所以我們想想.
我們蒙召去事奉的時候.
今天我們經常看聖經.
莊稼已熟.
但沒有人去收.
今天神真的呼召你.
你聽到嗎?.
神叫你去做.
在寒恩堂裡面.
有很多的崗位可以事奉.
你有沒有去參與呢?.
你有沒有留心呢?.
有沒有留心到.
在寒恩堂裡面的需要呢?.
你關不關心呢?.
你有沒有去祈禱會呢?.
你關不關心裡面的.
所發生的大小事情呢?.
今天老牧師在這裡.
我們很開心見到他.
你有沒有為他祈禱.
為他守望.
有沒有為曾姑娘.
還有為一群堂位.
在事奉.
沒有突然之間為他們禱告.
很重要的.
原來.
真是.
你過往可能沒有想過.
你有沒有想過.
要按神的旨意去服侍人.
服侍這個世代.
你有沒有想過呢?.

$^{441}$耶穌說我的食物就是.
遵行差外來者的旨意.
作成他的功.
如果你從來都.
沒有想過要服侍的話.
現在開始想吧.
有心不怕遲.
現在開始想吧.
不要想了就算了.
想了就要去做了.
做呢?.
不要緊.
不用怕.
你不會做.
有人會訓練你的.
一定是訓練了你.
才推你去做的.
不會不訓練你.
讓你出醜.
不會的.
在屬靈的事上.
我們很小心的.
一定要訓練你.
你有沒有聽過.
「書到用時方恨少」.
是嗎?.
那我們快點去做.
做了什麼時候呢?.
某個執事.
或者是曾姑娘叫你的時候.
馬上就可以接上班了.
可以做事了.
是嗎?.
快點去訓練.
因為就好像.
你打仗一樣.
養兵千日.
用了一時.
一會兒而已.
就是那一刻需要你.

$^{481}$那個時候.
需要你去跟那個人講見證.
你不要.
糟了.
我的見證要想一想.
要慢慢去做.
那就慘了.
你隨時心裡有一個稿.
已經知道我的見證是怎樣.
當然你不可以.
每年都是第一篇.
是嗎?.
你一定要有生活的見證.
你要隨時都寫下來.
隨時曾姑娘一叫你.
那個病人需要你.
跟你很相似.
你將你的見證跟她分享.
你不要到時候說.
糟了.
我要等明天才可以.
可能那個人明天就不找到了.
是嗎?.
你要隨時準備.
所以.
聖經說.
蒙召的人多.
損傷的人少.
我們不知道誰被損傷.
但你想想.
當年跟你一起洗禮的人.
二十年後.
他在哪裡.
但我們要小心.
因為將來.
我們要面對工作的審判.
所以我們真的要.
好好地在神面前祈禱.
好好地去盡責.
做好你基督徒的責任.

$^{521}$除了蒙召去示奉之外.
還要忠心.
神父召你.
接著你做的工作.
要忠心.
《阿摩西書》七章十字.
講到一個人.
叫做阿瑪. .
這個阿瑪. .
他就是.
伯達利的一個祭司.
這個祭司.
他控告阿摩西.
他控告阿瑪. .
什麼呢?.
他控告阿瑪. 什麼呢?.
他叫人告訴皇帝.
他說這個阿瑪. .
去煽動以色列人.
去背叛皇帝.
很大件事.
而這個祭司.
這個阿瑪. 的祭司.
他本身也不是很看重.
神的話語.
然後他對這個阿瑪. 說.
你不要再在這裡.
講道了.
不要再在伯達利去講語言了.
你去猶大講吧.
猶大那裡可以.
有很多時間你去講語言.
去講道吧.
那裡會供養你的.
這是什麼意思?.
他不准他在這裡講.
這番話.
對阿摩西來說.
是很大侮辱的.
因為阿摩西很清楚知道.

$^{561}$是神揀選他的.
雖然他出身是一個巫農.
但是神揀選他.
阿瑪. 是一個祭司.
他同樣.
都是在做事奉神的工作.
但是你會看到他們事奉的態度.
他們兩個之間.
是不是差距很大?.
今天你的事奉有沒有變質?.
想想我們自己的事奉.
我們剛剛信主的時候.
起初愛主的心很熱心.
有些事情我們都想去做.
但是.
過了若干年之後.
慢慢慢慢丟淡了.
可能整個世界都走到路上了.
忘記了神.
可能甚至沒有崇拜了.
大家有沒有想到.
原來做神的工作.
不是靠你的血氣可以承擔的.
阿摩斯說我原本不是先知.
我又不是先知的門徒.
我只是一個牧羊人.
我是修理桑樹的農夫.
我根本不是你講的那種.
或者錢去港渡的那種人.
他很清楚自己是.
被神揀選的.
神選了他.
他就去講寓言.
他面對攻擊的時候.
他沒有退縮.
神叫他說什麼.
他就說什麼.
他不會灰心.
他也不怕.
威嚇.

$^{601}$有人利誘他.
他都不為之所動.
但是他靠神的話.
放膽去講.
很多時候.
魔鬼會阻止你去傳講神的話.
所以你會不會.
因為.
魔鬼嚇你的時候.
你就妥協呢?.
還是你仍然.
會去忠心去服侍呢?.
在埃及有一個.
教堂叫做莫卡塔姆.
是一個洞穴的教堂.
在山洞裡面.
在埃及.
你想一下伊斯蘭教徒.
高達到百分之八十的.
這麼厲害的一個國家.
竟然擁有一間.
可以容納上萬人去聚會的教堂.
在洞穴裡面.
它的位置就在埃及的.
開羅東南部的地方.
可以說是一個最大的教堂.
每一個星期有七萬人去聚會.
在那裡敬拜,讚美.
很開心地去事奉神.
可以說是世界上.
最古老的基督教社區.
這個洞穴教會.
也被稱為聖西蒙修道院.
在那裡住了什麼人呢?.
就是以撿垃圾為生的一群人.
所以這個山.
也叫做垃圾山.
很特別.
他們.
如果他們有機會去那裡.

$^{641}$會看到那裡.
也有一些壁畫.
寫著經文.
《馬太福音》二十八章六節.
「他不在這裡,照他所說的,他已經復活了」.
即是說耶穌復活了.
據聞十世紀的時候.
伊斯蘭教徒很厲害.
他們統治那個地方.
這些基督徒遭受很大的逼迫.
甚至他們揚言要殺死他們.
他們的統治者叫做哈利法.
這個哈利法.
有一天他去圖書館.
看到一篇經文.
那篇經文說.
就是什麼呢?.
那篇經文說什麼呢?.
就是你們很熟悉的經文.
就是說.
「佳祥總」.
「佳祥總」很小的.
有這樣的信心.
真的可以移山.
他看到這樣的經文.
於是他就召見.
約基督徒的教皇.
叫他來訪談.
大家就約法三章.
他說.
如果你可以讓這個山移動的話.
我就放你們的命.
不殺你們.
如果不行的話.
我還是殺了你們.
這個教皇很害怕.
害怕的時候做什麼呢?.
祈禱.
他就向神祈禱.
求神幫助他.

$^{681}$神就派了一個獨眼俠.
一個皮匠.
跟他一起去祈禱.
他們很大聲地呼求神.
奉耶穌的名.
希望這個山可以移動.
很奇妙的.
一天到的時候.
他們繼續在祈禱.
哈利法也來了.
他們在祈禱的時候.
每一次他們去讚美敬拜的時候.
這個山好像在搖晃.
有很大的搖動.
在山下的人.
都覺得好像移動了.
看到陽光.
一個很奇妙的事情.
就是哈利法.
哈利法很害怕.
看到這個情形.
於是.
他旁邊的人都一樣.
跟他一樣很驚恐.
看到這個山搖晃.
他就跟教皇說.
神真的很偉大.
證實了你的信仰是真實的.
在埃及.
有位牧師叫約瑟夫·拿撒勒.
他說.
原來這個山大概移動了.
兩英里.
等於三公里的距離.
在2019年.
他又去台灣.
去台灣鼓勵台灣的信徒.
他說相信在3月到9月.
你們會經歷一個很大的復興.
他鼓勵他們.

$^{721}$就好像神去拯救他的子民.
就像埃及山移動的例子.
去跟他們分享鼓勵他們.
他說台灣會經歷一個.
神的應許就是那個經文.
就是一個《界賽總》的經文.
《聖經》裡面的《瑪特福音》.
17章20節.
那裡說到什麼呢?.
就是我實在告訴你們.
你們要有信心.
像一顆界賽總.
界賽總很小的.
就是對這個山說.
你從這邊移到那邊.
它也會移過去.
並且你們不會有一件.
不能做的事了.
這是對神一個倚靠的信心.
進一步.
他再次跟他們分享.
原來這個大山是有意思的.
除了一個實際的山.
另外一個意思就是.
人的心思意念.
就像一個山一樣.
它困擾著你.
他就說.
他相信台灣的基督徒.
要面對一個最大的山.
就是自己的心思意念.
每個人都需要正視自己這座山.
就是你的心思意念.
他說所有教友的成規.
或者是綑綁著你的.
你都要把它拿走.
柴火夠了.
禱告也充足了.
那些在教會裡.
已經裝備好的成熟的精兵.

$^{761}$你肯不肯勇敢地走出社會.
為耶穌做見證呢?.
當日他這樣呼籲.
這些已經裝備好了.
你要勇敢地走出去.
有時候我們在教會裡.
好像當日.
他們說我們在這裡多好.
但耶穌不是說我們在這裡多好.
是我們要走出去.
讓人可以認識神.
你不走出去.
他很怕來教會.
覺得你不知道是怎樣的.
他不會進來的.
是你要走出去的.
所以他鼓勵他們要走出去.
還有.
他說.
從前對於傳福音缺乏熱誠的基督徒.
過去可能你不是那麼注重的.
但現在你要更加用心.
當時他跟他們說.
要重燃這個福興的火焰.
他呼籲他們.
希望他們.
作為亞洲傳揚福音的起點.
主要目的是向廣大的華人世界.
傳講神的信息.
然後進而影響周邊的地域.
叫他們要把握這一波.
前所未有的福興浪潮.
其實.
同樣的.
我們今天也需要神福興我們.
在座的弟兄姊妹.
今天弘恩需要什麼呢?.
我們有沒有好好地.
去讀我們的聖經呢?.
盼望我們.

$^{801}$在這個彎曲勃謬的世代裡.
願意好好地勤讀聖經.
明白真道.
因為異端會侵入的時候.
你沒有真道是很慘的.
作神無瑕疵的兒女.
你要好好地.
去明白真道.
去實行神的真理.
當有異端進來的時候.
你懂得分辨.
你要趕走它.
不要讓它.
去孵化這間教會.
我們行事為人.
是不是跟耶穌基督的福音相稱呢?.
這也是值得我們思想的.
我們願不願意為了福音的信仰.
我們齊心努力.
不怕敵人的威嚇.
盼望我們.
好好地去謹守真道.
我們可以忠心為主.
去事奉到底.
我們一起祈禱.
愛我們的主.
愛我們的神.
主啊我們.
是不配的罪人.
但你竟然拯救我們.
主啊你愛我們.
求你幫助我們.
不要違背你的心意.
讓我們可以為主.
盡忠到底.
求你聖靈時時去鞭策我們.
何時我們走歪了.
提醒我們.
以致我們走在你旨意當中.
去滿足你的心意.

$^{841}$我們這樣祈禱.
教導下奉主耶穌基督的名教.
阿們.
\newpage



\section{馬可福音 5:1-20}
\label{sec:5NHMIEcWoqY}
\textbf{2023.07.30 《多麼大的憐憫》 馬可福音 5\_1-20 (和修版) 講員 : 黃宏達傳道}
\newline
\newline
連結: \href{https://youtube.com/watch?v=5NHMIEcWoqY}{\texttt{https://youtube.com/watch?v=5NHMIEcWoqY}} ~~~~ 語音日期: 2023-08-01
\newline
\newline
\hyperref[sec:Jxn7bTnB5O8]{\small{< < < PREV SERMON < < <}}
~
\hyperref[sec:index_chronic]{\small{[返順時目]}}
~
\hyperref[sec:index_scriptual]{\small{[返順卷目]}}
~
\hyperref[sec:cnxh3H5LvDs]{\small{> > > NEXT SERMON > > >}}
\newline
\newline
馬可福音 5:1-20
\newline
\begin{longtable}{cl}
\hline
\hline
章節 & 經文 (和合本修訂版)\\
\hline
5:1 & \begin{tabularx}{0.7\textwidth}{X} 他們渡到海的對岸，到格拉森人的地區。 \end{tabularx} \\ \\ \relax
5:2 & \begin{tabularx}{0.7\textwidth}{X} 耶穌一下船，就有一個污靈附身的人從墳墓迎著他走來。 \end{tabularx} \\ \\ \relax
5:3 & \begin{tabularx}{0.7\textwidth}{X} 那人常住在墳墓裡，沒有人能捆住他，就是用鐵鏈也不能； \end{tabularx} \\ \\ \relax
5:4 & \begin{tabularx}{0.7\textwidth}{X} 因為人屢次用腳鐐和鐵鏈捆鎖他，鐵鏈被他掙斷，腳鐐也被他弄碎了，總沒有人能制伏他。 \end{tabularx} \\ \\ \relax
5:5 & \begin{tabularx}{0.7\textwidth}{X} 他晝夜常在墳墓裡和山中喊叫，又用石頭打自己。 \end{tabularx} \\ \\ \relax
5:6 & \begin{tabularx}{0.7\textwidth}{X} 他遠遠看見耶穌，就跑過來拜他， \end{tabularx} \\ \\ \relax
5:7 & \begin{tabularx}{0.7\textwidth}{X} 大聲呼叫說：「至高神的兒子耶穌，你為甚麼干擾我？我指著神懇求你，不要叫我受苦！」 \end{tabularx} \\ \\ \relax
5:8 & \begin{tabularx}{0.7\textwidth}{X} 這是因耶穌曾吩咐他說：「污靈啊，從這人身上出來！」 \end{tabularx} \\ \\ \relax
5:9 & \begin{tabularx}{0.7\textwidth}{X} 耶穌問他：「你叫甚麼名字？」他說：「我名叫群，因為我們數目眾多。」 \end{tabularx} \\ \\ \relax
5:10 & \begin{tabularx}{0.7\textwidth}{X} 他就再三求耶穌不要叫他們離開那地方。 \end{tabularx} \\ \\ \relax
5:11 & \begin{tabularx}{0.7\textwidth}{X} 在山坡那裡，有一大群豬正在吃食； \end{tabularx} \\ \\ \relax
5:12 & \begin{tabularx}{0.7\textwidth}{X} 污靈就央求耶穌，說：「求你打發我們進入豬群，好附著牠們。」 \end{tabularx} \\ \\ \relax
5:13 & \begin{tabularx}{0.7\textwidth}{X} 耶穌准了他們，污靈就出來，進入豬裡，那群豬就闖下山崖，投進海裡，淹死了。豬的數目約有二千。 \end{tabularx} \\ \\ \relax
5:14 & \begin{tabularx}{0.7\textwidth}{X} 放豬的逃跑了，去告訴城裡和鄉下的人。眾人就來，要看發生了甚麼事。 \end{tabularx} \\ \\ \relax
5:15 & \begin{tabularx}{0.7\textwidth}{X} 他們來到耶穌那裡，看見那被鬼附的人，就是曾被群鬼所附的，坐著，穿著衣服，神智清醒，他們就害怕。 \end{tabularx} \\ \\ \relax
5:16 & \begin{tabularx}{0.7\textwidth}{X} 看見這事的人把被鬼附的人所遇見的，和那群豬的事，都告訴了眾人， \end{tabularx} \\ \\ \relax
5:17 & \begin{tabularx}{0.7\textwidth}{X} 眾人就央求耶穌離開他們的地區。 \end{tabularx} \\ \\ \relax
5:18 & \begin{tabularx}{0.7\textwidth}{X} 耶穌上船的時候，那曾被鬼附的人懇求要和耶穌在一起。 \end{tabularx} \\ \\ \relax
5:19 & \begin{tabularx}{0.7\textwidth}{X} 耶穌不許，卻對他說：「你回家去，到你的親友那裡，將主為你所做多麼大的事和他怎樣憐憫你，都告訴他們。」 \end{tabularx} \\ \\ \relax
5:20 & \begin{tabularx}{0.7\textwidth}{X} 那人就走了，開始在低加坡里傳揚耶穌為他做了多麼大的事，眾人就都驚訝。 \end{tabularx} \\ \\
[1ex]
\hline
\hline
\end{longtable}
$^{1}$(鄧家彪) 各位姐妹早安.
(各位姐妹) 早安.
(鄧家彪) 多謝鄧姑娘的邀請.
(鄧家彪) 多謝曾牧師的邀請.
(鄧家彪) 今天有機會來到這裡.
(鄧家彪) 跟大家一起分享神教話.
(鄧家彪) 很感恩.
(鄧家彪) 剛才可以跟大家一起去敬拜.
(鄧家彪) 雖然今天第一次來.
(鄧家彪) 不過來到這裡很舒服.
(鄧家彪) 原來我們是特別空曠的.
(鄧家彪) 逛市區的人.
(鄧家彪) 多點來元朗.
(鄧家彪) 前陣子跟兒子一起看Netflix.
(鄧家彪) 我介紹他看一套舊舊的電影.
(鄧家彪) 已經二十多年前了.
(鄧家彪) 可能你也看過的.
(鄧家彪) 叫做《雷霆救兵》.
(鄧家彪) 有沒有人看過?.
(鄧家彪) 看過吧?.
(鄧家彪) 都是舊舊的.
(鄧家彪) 不過太經典了.
(鄧家彪) 我兒子還沒看過.
(鄧家彪) 所以我介紹他看.
(鄧家彪) 他今年十八歲.
(鄧家彪) 剛剛考完DSE.
(鄧家彪) 如果你有印象.
(鄧家彪) 你大概知道情節了.
(鄧家彪) 開頭是一個.
(鄧家彪) 一開始就很轟烈的.
(鄧家彪) 盟軍搶灘.
(鄧家彪) 搶灘落萬幣是成功的.
(鄧家彪) 站穩陣腳之後.
(鄧家彪) 就統計一下傷亡人數.
(鄧家彪) 你也知道有多慘烈.
(鄧家彪) 死了很多人.
(鄧家彪) 那些統計資料去到總部之後.
(鄧家彪) 那個秘書打字的那個.
(鄧家彪) 他統計全球各地的美軍傷亡.
(鄧家彪) 看到落萬幣有一個陣亡的士兵.

$^{41}$(鄧家彪) 又有另一個在非洲戰場陣亡的士兵.
(鄧家彪) 又有另一個在亞洲戰場陣亡的士兵.
(鄧家彪) 原來是三兄弟.
(鄧家彪) 三兄弟.
(鄧家彪) 這個秘書就很惡心.
(鄧家彪) 就去跟上司說.
(鄧家彪) 上司又再跟將軍說.
(鄧家彪) 怎麼辦呢.
(鄧家彪) 因為還有另一個兄弟.
(鄧家彪) 就是在歐洲某個位置.
(鄧家彪) 就問將軍.
(鄧家彪) 這個媽媽她有四個兒子.
(鄧家彪) 有三個已經陣亡了.
(鄧家彪) 只剩下最後一個.
(鄧家彪) 要不要去救她呢.
(鄧家彪) 將軍想了一會.
(鄧家彪) 就決定要去救她.
(鄧家彪) 要派一支拯救的一小隊人.
(鄧家彪) 去救主角.
(鄧家彪) 叫做Ryan.
(鄧家彪) 記不記得 雷霆救兵.
(鄧家彪) 當然可以有很多思考.
(鄧家彪) 值不值得呢.
(鄧家彪) 但肯定一件事.
(鄧家彪) 其實是很高風險的.
(鄧家彪) 第一 都不知道他在哪裡.
(鄧家彪) 大概是那邊.
(鄧家彪) 但不是就這樣去的.
(鄧家彪) 中間要穿越敵人控制的地區.
(鄧家彪) 即是你預計可能會發生衝突.
(鄧家彪) 去到 找不找到有一件事.
(鄧家彪) 還有 他會不會死了呢.
(鄧家彪) 又不知道.
(鄧家彪) 即是很多不確定.
(鄧家彪) 很多的危險.
(鄧家彪) 但最後將軍都覺得.
(鄧家彪) 他很同情那個媽媽.
(鄧家彪) 那個媽媽已經沒有了三個兒子.
(鄧家彪) 這個是他最後一個.
(鄧家彪) 他很憐憫這個媽媽.

$^{81}$(鄧家彪) 雖然任務很危險.
(鄧家彪) 但他都決定要派一支拯救隊.
(鄧家彪) 去救 Ryan.
(鄧家彪) 其實我很感動這個故事.
(鄧家彪) 如果你忘記了.
(鄧家彪) 或者你未看過.
(鄧家彪) 我都鼓勵你看看.
(鄧家彪) 不過有少少血腥.
(鄧家彪) 即是頭那一幕.
(鄧家彪) 如果頂不住可以skip.
(鄧家彪) 原來在聖經裡面也有類似的故事.
(鄧家彪) 聖經記載耶穌憐憫一個沒有人重視的人.
(鄧家彪) 為他做了一件大事.
(鄧家彪) 所以我今天分享的題目是.
(鄧家彪) 多麼大的憐憫.
(鄧家彪) 今天早上讓我們和主耶穌說話.
(鄧家彪) 去看看耶穌是怎樣憐憫我們.
(鄧家彪) 特別當我們覺得沒有人重視你.
(鄧家彪) 沒有人愛惜你的時候.
(鄧家彪) 今天是經文.
(鄧家彪) 我們特別要留意.
(鄧家彪) 我們一起祈禱吧.
(鄧家彪) 親愛的天主.
(鄧家彪) 祈禱給我們.
(鄧家彪) 親愛的天父.
(鄧家彪) 我來到你的面前.
(鄧家彪) 要聆聽你的說話.
(鄧家彪) 求真理的聖靈.
(鄧家彪) 帶領我們進入你的真理當中.
(鄧家彪) 多謝你聽我們禱告.
(鄧家彪) 逢主耶穌基督明求.
(鄧家彪) 阿們.
剛才我們已經讀了這個故事.
在《馬學福音》的五章.
一至二十節.
可能你也有些印象.
其實《馬學福音》第五章.
整個第五章有三個神蹟.
一個緊接一個.
三個神蹟.

$^{121}$第一個就是剛才我們讀了的.
接著就是患血漏的婦人.
和艾路的女兒.
連同前面那一章.
最後那幾節.
即是第四章尾.
就是平靜風和浪.
大家應該也很熟悉.
所以連續四個神蹟.
放在一起.
其實就連續了前面第三章.
甚至第一章開始.
已經帶出了一個主題.
就是耶穌的權柄.
所以前面是一個大段落.
這裡第四章尾開始.
是第二個大段落.
是連續四個神蹟.
為什麼是新的段落呢?.
很簡單.
因為轉了地方.
因為耶穌在四章三十五節說.
我們到哪邊去吧.
原本是在加利利海的西岸.
即是加利利省.
大家有印象吧.
如果你有實體聖經.
其實我挺喜歡實體聖經.
眼睛舒服一點.
後面也有地圖.
你看地圖就看到了.
加利利是北部的省份.
右下角有一個.
你叫海也可以.
你叫湖也可以.
加利利海.
耶穌一出來傳道的時候.
就在加利利那裡.
特別是在加伯農.
有很多記載.

$^{161}$馬可福音頭三章.
到四章尾就說.
我們到東邊.
即是到對面海.
就是指加利利海或湖的東岸.
由西岸坐船到東岸.
所以這裡就開始新的段落了.
原本加利利就是猶太人聚居的地方.
去到東岸除了地方轉了之外.
那裡也是外邦人住的地方.
很容易的.
因為你看到有豬.
有這麼多豬.
猶太人是不吃豬的.
因為律法說是不潔的食物.
所以不吃豬.
所以當然不會養豬.
所以看到那裡這麼多豬.
你知道那裡不是猶太人住的地方.
外邦人住的地方.
所以這段經文其實很有特色.
很緊湊和很多細節.
雖然馬太和盧迦都有記載這件事.
但沒有它這麼詳細.
你看到足足有二十節.
馬太馬口很短.
只有十五章.
這麼短.
有二十節的經文.
比馬太和盧迦還要長.
所以在《講鬼》的神蹟裡.
最精彩.
最多細節的地方.
很值得我們今天仔細去看一看.
透過這段話去鼓勵我們.
這段經文一開始.
第二歲精美就說到.
去到加利利海的東邊.
那裡是外邦人居住的地方.
而格拉森是一個城市.

$^{201}$廣闊的地區.
在這段結尾.
我們看到底加玻利.
串了這個地方名.
底加玻利.
底加玻利是音譯.
這個名字是音譯.
底加即是Deca.
即是十的意思.
玻利.
Polis.
即是城市.
即是十個城市.
底加玻利這個字的意思.
即是十城.
十城是一個地區.
其中一個就是格拉森.
耶穌去到東岸.
就是去到格拉森.
底加玻利的其中一個城.
耶穌一下船.
一下船就看到一個.
從墳墓裡出來.
被污淩附身的人.
加少少想像一下.
加少少想像.
可能少一點.
你假裝上來坐車.
一下車.
你就看到一個.
由墳墓裡走出來的污淩人.
你想像一下.
這個情景都幾嚇人.
是嗎?.
那些墳墓是依山打造而成的.
和我們中國人不同.
我們古代的中國人通常是土葬.
猶太人就不是.
在巴勒斯坦地區應該這樣說.
他們是放在墳墓裡.

$^{241}$放在一個由山造的洞出來.
屍體放在那些山洞裡.
那些山洞當然比較容易造.
比一般的花崗岩更硬.
所以為什麼人會在墳墓裡走出來.
而不是在地上.
在山洞裡走出來.
是在山裡走出來.
所以如果熟悉猶太信仰的人.
很快就會聯想到什麼呢?.
律法說死人是不潔的.
住在墳墓的人.
自然是不潔.
一個不潔的人走出來.
就算不是從猶太人角度.
從外邦人的角度.
墳墓是在近岸地區.
不是在城市裡.
遠在城外.
如果說明了.
其實城市裡的人.
包括家人.
其實是不要他的.
對他很有戒心.
將他遺棄在一個很偏遠的地方.
但他真正令人害怕的.
不是因為他住在墳墓.
而是他被鬼纏身.
被魔鬼控制.
他那麼強大的力量.
令人害怕.
馬學很少用一個.
很串聯,覆聚的手法.
去強調.
好像很沉氣.
第三至第四節.
剛才大家讀過.
有些字很相似.
沒有人也不能.
大家可以回來看.

$^{281}$有些經文.
總沒有人能.
花了很多篇幅.
去描寫被鬼控制的人.
他力量太大了.
大到根本沒有人可以制服他.
就算上了鐵鐐.
都被他弄斷.
大家可以想像一下.
他的力量.
太大了.
恐怖有力.
所以.
第三章二十七節可以這樣翻譯.
總沒有人能夠.
有足夠的強壯.
去制服他.
如果你熟悉第三章.
你就會很明白.
第三章耶穌說.
要把喪事捆綁.
文字說.
靠鬼王去捆綁.
耶穌說.
是一個比喻.
人要把喪事捆綁.
才能搶他的家具.
就是喪事.
所以在呼應第三章.
馬可對這個人的自殘.
傷害自己.
他描寫得很獨特,很深刻.
馬可形容他.
是晝夜在墳墓裡哭叫.
在山中哭叫.
又用石頭捆綁自己.
大家可以想像一下.
有一個人.
晝夜在哭叫.
有沒有遇過這種情況?.

$^{321}$有時在街上.
可能在搭輕鐵.
突然有人大叫.
都會嚇到.
看看發生什麼事.
是不是兒子激怒媽媽.
還是在通電話.
很激氣之類.
他晝夜哭叫.
而且用石頭捆綁自己.
你想像一下.
我們不會這樣做的.
正常人不會這樣做.
用石頭捆綁自己.
自殘的景象.
可以大概想像他的內心.
是求生不得.
求死不能.
同一件事.
在《馬太福音平行經文》記載.
他極其凶猛.
路加沒有仔細描述.
如何傷害自己.
馬可在第二至第六節.
讓人去描述.
讓人感受到.
這個人真的很痛苦.
很痛苦.
他身上魔鬼的力量很強大.
我們跳到故事的結局.
第十五節.
還會明白多些.
第十五節他康復後.
聖經形容他是坐著.
第十五節.
坐著.
穿著衣服.
神智清醒.
換句話說.
原本他不是坐著.

$^{361}$是通山走.
沒有穿衣服.
穿上衣服.
這個人沒有穿衣服.
沒有穿衣服.
後來神智清醒.
之前神智完全不清醒.
有沒有遇過類似的情況.
可能不是很多.
除非你是精神病患的醫護.
可能會有更多機會接觸.
平時我們較少機會接觸.
所以我叫大家想像一下.
就是這樣.
赤身露體的人.
在自殘.
在哭叫.
很慘.
社會不要他.
家人也不要他.
這樣的人.
在一般人眼中.
你覺得有沒有價值呢.
有沒有價值呢.
沒有價值吧.
失敗者中的失敗者.
羅底橙中的羅底橙.
沒有人要他.
家人也不要他.
不過沒有人要他的人.
耶穌要他.
沒有人憐憫他的人.
耶穌憐憫他.
我們從第四章尾.
已經大概知道.
耶穌搭船從西岸過東岸.
剛才提過.
是否在帆風順.
不是吧.
剛才也提過.

$^{401}$平靜風和浪.
記得嗎.
那個故事你應該有印象.
每次狂風大作.
很危險.
那些漁夫已經在打魚.
已經是他熟悉的海域.
都害怕得震.
都害怕得沒命.
所以其實很危險.
那個路程還未到.
隨時有生命危險.
是不是好像雷霆救兵的情節.
真的好像.
負上生命的危險.
不單是自己去.
連同他的門徒也去.
所以耶穌是特意去找那個人.
他好端端在加利利.
為什麼突然搭船.
去外邦人的地方呢.
是不是.
耶穌是冒著生命危險.
特意去找這個沒有人憐憫的人.
差點沒命.
門徒也差點沒命.
當然沒事.
耶穌的大能.
所以弟兄姊妹.
不知道你會不會覺得.
有時自己都沒什麼價值.
我團契多數都是讀過大學的弟兄姊妹.
但都有少數沒讀過.
有時和一些姐妹聊天.
我目中的一個夫婦團契.
全部都是夫婦.
她的丈夫是大學畢業.
她自己沒讀過大學.
原來她也有些自卑.
這個團契每個人都讀過大學.

$^{441}$所謂每個人.
其實多數都不是每個人.
自己沒讀過大學.
覺得比不上別人.
有這樣的感受.
很可惜.
我從來都不會覺得.
你沒讀過大學.
你閉嘴.
少說兩句.
我從來沒這樣看她.
她不是說我這樣看她.
她有一種覺得比不上別人.
我也遇過一些姐妹.
聽她說她在九龍城寨長大的.
你知道城寨是什麼地方.
幾年前無線也有一套電影.
《黑暗的勢力》那些武打片.
幾年前的有印象.
很多罪惡的地方.
黃賭毒什麼都有.
她在那裡長大.
每次回家上學都很害怕.
當然長大後就離開了城寨.
搬了出去.
很感恩後來順了主.
她公開分享.
有時都覺得自己.
不要這樣跟別人說.
我小時候住城寨.
覺得比不上別人.
自己好像沒什麼價值.
她很勇敢.
她覺得自己也有一點自卑.
幸好順了主.
經歷上帝無條件的愛.
好多.
但她知道這是她一生.
要學習對付的地方.
所以弟兄姊妹.

$^{481}$你相不相信你在神眼中.
是很有價值的?.
無論你過去的學歷是怎樣.
無論你現在做什麼工作.
無論你的子女是怎樣.
剛剛DSE放榜.
今屆狀元歷屆最少.
有沒有人的子女剛剛DSE?.
有沒有?.
沒有?.
不要緊.
已經長大了.
已經出生了.
可能很多父母.
都是有意無意.
都會覺得慘了.
兒子考成這樣.
女兒考成這樣.
可能很失望.
可能覺得隔壁的媽媽就好了.
成績又好.
又入了神科.
讀了好大學.
她可能沒有這樣說.
但自己很自然會這樣比較.
覺得不是自己考試.
是兒子考試.
你記住.
不是你成績差.
我一定差嗎?.
是吧?.
成敗得失是他的.
不是你的.
當然我們都會著緊.
我自己剛剛考試.
我都跟他說.
一早說過.
說過很多次.
兒子無論你什麼成績.
爸爸都這麼愛你.

$^{521}$你仍然是我的寶貝.
當然我當日也有些緊張的.
放榜那天.
但嘗試不跟別人比較.
有些入了神科.
有些普普通通.
當然.
將他的成敗得失.
跟他的價值分開.
還要跟我的成敗.
我作為爸爸.
我是不是一個好爸爸.
有些人會自我質疑.
很多父母都自責.
都是自己不好.
如果我當時這樣幫他安排補習.
如果我當時這樣培育他.
他今天就不會這樣.
是不是.
很多父母的心境.
很複雜的.
其中一個不好的感受.
就是自責.
怪自己做得不好.
不必要.
當然會繼續關心.
但關心一個聾.
連自我價值.
都跟他的成敗扣上.
是很傻的.
原因是耶穌基督不是這樣看你的.
這個沒有人要的人.
耶穌都冒著生命危險去救他.
所以弟兄姊妹.
你相信嗎.
你在神眼中.
是有很寶貴很高的價值.
所以有機會.
你多點去.
這樣去到.

$^{561}$默念這個事實.
神啊.
我相信.
你是如此的看重我.
你是如此的憐憫我.
我在你的眼中.
是一個有價值的人.
試試吧.
當你有時安靜祈禱的時候.
在逛街的時候.
你這樣去到默念這個屬靈的真理.
我們繼續看下去.
相信神的憐憫.
第二讚美神的憐憫.
即使污鬼是多麼厲害.
我們上文提過.
污鬼的破壞力多大.
沒有人可以制服他.
連鐵鐐都鎖不住他.
被他弄碎.
力大無窮.
第六節.
第六節經文說到.
被鬼附的人.
一見到耶穌就拜他.
連被鬼附的人.
一見到耶穌就拜他.
他很厲害.
但是這麼厲害的人.
連污鬼都拜耶穌.
明白了嗎?.
那誰更厲害呢?.
誰更有權柄呢?.
立即去拜他.
而且留意一下.
污鬼對耶穌的稱呼.
第七節說.
大聖夫叫說.
至高神的兒子耶穌.
不是只有耶穌.

$^{601}$前面還有幾個字.
至高神的兒子耶穌.
其實第一章二十四節.
已經稱呼耶穌.
是污鬼稱呼的.
第一章已經污鬼稱呼.
是神的聖者.
三章十一節.
神的兒子.
所以魔鬼其實他認識耶穌.
在舊約裡面.
至高者是對神的稱呼.
就算不是舊約的背景.
在希羅文化裡面.
至高者.
在希臘神話裡面.
最厲害的神就叫至高者.
所以無論猶太人也好.
外邦人也好.
或者我們今天讀聖經的人也好.
我們都能夠體會到.
至高者的兒子.
其實就是神的兒子.
你看到神的兒子.
我們看到就開心.
我們看到都會拜.
但魔鬼不是.
魔鬼也有痛點.
他的痛點在哪裡呢?.
第七節他繼續說.
至高神的兒子耶穌.
你為什麼干擾我?.
我指著神懇求你.
不要叫我受苦.
受什麼苦呢?.
這裡沒有記載.
《馬太福音》就有記載.
八章二十九節.
馬太那裡.
污鬼就有提到.

$^{641}$時候還沒有到.
什麼時候呢?.
是指在末日大審判的時候.
除了未信主的人會遭受審判.
信了也會.
魔鬼也會受審判.
審完之後.
就丟入火壺裡.
這是魔鬼的結局.
在《馬太福音》二十五章十一節就說到.
魔鬼的結局.
就是被耶穌基督丟入火壺裡.
在大審判的時候發生.
所以魔鬼也知道.
其實他最終的結局.
就是丟入火壺裡.
所以他跟耶穌說.
未到時間.
時候還沒有到.
你這麼快就要我下火壺.
懇求耶穌.
不要吧.
所以他很害怕.
他看到神的兒子耶穌基督.
他很害怕.
是不是提早來罰我呢?.
請留意第九節.
第九節.
我們一起讀.
預備起.
耶穌問他他叫什麼名字.
他說我名叫群.
因為我們數目眾多.
這裡要多解釋.
第九節.
耶穌問他.
是不是.
問完之後.
他說.
是不是.

$^{681}$這兩個他.
是誰呢?.
只有兩個選擇.
要麼就是魔鬼.
要麼就是那個人.
是不是.
先說第一個他.
有多少個你覺得.
第一個他是說人的呢?.
你叫什麼名字?.
沒有.
你覺得他是魔鬼的呢?.
多數都覺得他是魔鬼.
第二個他.
我名叫群.
這個他是指一個人.
舉手.
有沒有呢?.
沒有.
魔鬼呢?.
魔鬼.
開估了.
第一個他.
不是魔鬼.
是那個人.
為什麼這樣說呢?.
因為.
我不明白為什麼.
因為大家都看中文聖經.
中文聖經是對的.
用這個他.
但是原文聖經的他.
其實是陽性.
陽性.
陽性的他.
如果是指魔鬼的他.
應該是中性.
明白嗎?.
是一個陽性的他.
所以其實耶穌是在問那個人.

$^{721}$當然那個人裡面有很多鬼.
耶穌不是在問那個鬼你叫什麼名字.
是在問那個人.
你叫什麼名字.
當然答那個呢?.
大家都是答魔鬼.
對不對?.
對的.
魔鬼又答我名叫群.
所以是魔鬼搶著答.
明白嗎?.
魔鬼搶著答.
這樣就很有意思了.
這樣搞通了這個東西.
第九節.
第一個他是陽性.
所以.
還有那個問.
原文的時態是未完成式的動詞.
中英文也沒有這個時態.
總之意思是.
保持不停去做那個動作.
就是說耶穌問那個男人.
不是問一次你叫什麼名字.
是持續地問.
持續地進行那個問的動詞.
你叫什麼名字?.
持續地問.
所以.
就是說耶穌其實.
第一他不是跟魔鬼對話.
他問那個人.
他想喚醒這個人.
喚醒這個人的意識.
明白嗎?.
他很關心這個人.
如果你第一次認識一個人.
你第一句都會說.
怎麼稱呼?.
你叫什麼名字?.

$^{761}$是嗎?.
都會這樣.
是一個尊重.
你至少知道他叫什麼名字.
就算其他東西你不知道.
尊重他.
是一個尊重.
他根本是神智不清的人.
所以想喚醒他的意識.
表達一種關愛.
當然是魔鬼在控制他.
所以魔鬼搶著回答.
我名叫群.
你想想.
如果.
你有一個名字.
譬如我當是陳大文.
我叫你.
1234.
過來吧.
你覺得怎樣?.
我喊你一個號碼.
你會不開心.
我有名字的.
雖然陳大文很普通.
但也有個名字.
如果就這樣喊你1234.
通常在監獄就是這樣.
就是這樣喊人.
但你不是監犯.
你是一個正常人.
稱呼一個正常的名字.
如果熟悉的可能會喊花名.
死肥仔.
說笑而已.
就會這樣.
喊個名字.
是對他的尊重.
和喚醒他的意識.
所以.

$^{801}$當然他被魔鬼控制.
魔鬼厲害了.
我名叫群.
OK.
你想像一下.
一個人.
你當我差不多的身形.
裡面有很多鬼.
不是一隻.
第五章之前的鬼是一隻.
這裡是群.
群即是多少?.
學者也很同意.
群其實是一個羅馬軍團.
是同一個字.
羅馬軍團的意思.
一個軍團有多少人?.
大概四千至六千個步兵.
當時第一世紀羅馬帝國.
全國有25個軍團.
OK.
有25個軍團.
一個軍團.
四千至六千個步兵.
當然.
是不是真的裡面有四千至六千隻鬼?.
不肯定.
不可以因此說有四千至六千.
但至少肯定是很多很多鬼.
對不對?.
軍團是軍隊最大的編制單位.
就是團.
很多很多隻鬼.
你當它是半價也好.
二千隻鬼.
都嚇死你了.
對不對?.
所以.
讀聖經的人.
看到羅馬軍團.

$^{841}$想到什麼?.
控制.
壓迫.
是一個踐踏人的勢力.
一隻鬼已經很可怕了.
還有這麼多隻鬼.
所以耶穌正面對一個非常強大的.
邪惡的力量.
當然.
耶穌也能驅趕.
那你知道誰才是最強最強最powerful的那個.
耶穌才是那位強者.
沒有人能夠制服的格拉森人.
耶穌就能夠.
耶穌能夠將這群污鬼驅趕進去.
豬群.
就正正證明了.
究竟誰才是強者.
就是至高神的兒子耶穌.
祂能夠將壯士捆綁.
祂能夠宣告惡勢力的終結.
所以這個人真的很慘.
群鬼控制他.
傷害他.
令他求生不得.
求死不仁.
令到社會不要他.
他的家人都不要他.
但是耶穌故意去救他.
半年前在我的團體.
發生了類似的事情.
我們牧者負責帶領新人組.
有一位未信的姐妹就來了.
我是後來才知道的.
那晚我帶茶經.
就是帶這段經文.
我也派了一些靈修材料給大家.
回家可以再靈修默想這段經文.
那位姐妹後來跟我說.
她說原來她當年被鬼壓.

$^{881}$是很辛苦的.
鬼是捏著她的頸.
所以晚上不能呼吸.
很辛苦.
她的丈夫也很相信居士.
也是未信主的.
她丈夫就帶她去看居士.
居士告訴她.
那隻鬼是男性的鬼.
當然有嘗試給她符之類的東西幫助.
但是解決不了.
總之那次她就來了我的團體.
就讀了這段經文.
晚上她回家睡覺.
她就跟鬼說.
我已經相信耶穌了.
你走吧.
她說話很斯文.
我猜她也是這樣的語氣跟鬼說.
我相信耶穌了.
你走吧.
她就感覺到一句.
啊一聲就離開了她.
她從此就沒有再有那個箍頸.
後來團體回來.
她告訴我.
她很開心地跟我說.
哇!哈利路亞.
讚美主.
最偉大,最有權能.
得勝一切污鬼的神.
所以這個姐妹其實有很多艱難.
很多不開心的.
無論是夫妻關係.
親子關係.
很多不開心.
自從她相信了耶穌.
當然不是所有問題都解決了.
但是她整個心境不同了.
她好像看到很多出路.

$^{921}$後來她也嘗試帶丈夫回來.
不過就轉了團體.
所以真的很奇妙.
不是我做了什麼.
我只是普通人.
是耶穌基督的權能.
耶穌很憐憫這個姐妹.
就好像憐憫格拉森人一樣.
我們繼續看.
這個格拉森人.
回復真相.
你又可以想像一下.
一個坐不定的人.
到處走.
到處大聲叫的人.
現在坐定了.
好像你們一樣.
坐定了.
穿回衣服了.
之前沒有穿衣服.
神智清醒了.
可以和你正常的對答.
你看到這一幕.
真的很開心.
如果我是她.
我看著耶穌.
我都想抱著祂哭.
多謝你救我.
人人都不要我.
你要我.
你救我.
你憐憫我.
你重視我.
她不是就這樣走了.
她和耶穌有一個請求.
當然群眾反應很奇怪.
群眾覺得很害怕.
看到會害怕.
經文沒有說她害怕什麼.
很多人都相信.

$^{961}$如果耶穌下次又再趕鬼.
不是又沒有多二千隻豬嗎?.
這次二千隻.
下次又不知道多少隻.
不如走吧.
很諷刺吧.
這麼善良.
這麼憐憫人的主耶穌.
那些人看重的.
可能是經濟損失.
不要耶穌.
看到那個人.
應該為他高興.
現在好了.
還算正常人.
不應該恭喜.
都沒有.
都可以這麼冷漠.
所以有時遇到人.
很冷漠.
很涼薄.
聖經時代都是這樣.
不認識神人.
但這個人.
他本人當然不是這樣.
他本人.
十七節.
不是十七.
是十八節.
「同被鬼附的人.
懇求要和耶穌在一起」.
你可以想像一下.
如果你是他.
你會不會也是這樣?.
我想你也會吧.
你連家人也不要你.
你想想.
家人如果愛他.
我會說.
耶穌你這麼愛我.

$^{1001}$反正我都沒有家了.
不如你就讓我跟著你吧.
好像你的十二個門徒一樣.
我的命是你救回來的.
從此我跟著你.
對不對?.
很正常吧?.
他很自願.
很有愛.
想報答耶穌.
是不是?.
但是耶穌的反應又很有趣.
耶穌的反應.
很有趣.
耶穌沒有答應.
沒有答應.
耶穌.
當然沒有罵他.
耶穌的反應是第十九節.
你回家去.
到你的親友那裡.
將主為你所做多麼大的事.
和他怎樣憐憫你.
都告訴他們.
你想像一下.
如果你是格拉森人.
你聽到耶穌這樣回答你.
你有什麼感受呢?.
會不會有點失望?.
會不會覺得很困惑.
不明白.
會不會生氣呢?.
這裡沒有形容這個人的心境.
我們很難肯定他在想什麼.
但是我想像.
如果我是他.
我很不明白.
第一,我很不明白.
家人都不要我了.
這裡的人又不歡迎我.

$^{1041}$剛才那群律師才剛出現.
走過來說恭喜都沒有.
是吧.
怡不得快離開這一帶.
對不對.
耶穌你這麼好.
為什麼不讓我跟著你呢?.
你又這麼愛我.
你來到這裡都沒有生命危險.
為什麼不讓我跟著你呢?.
還有.
我可能有很多掙扎.
剛才我說這群人不要我.
家人都不要我.
我想我對我的家人.
是有很多不滿.
覺得很受傷.
覺得被遺棄.
是不是.
一個連不要自己的家庭.
其實我都不想回去.
是不是.
我都不想見到我的爸媽.
兄弟姐妹.
他們都不要我.
可能當日鎖我.
他們都有份.
是不是.
你我再次面對他們.
很難.
是不是很掙扎.
是不是.
當然.
他最後有沒有做呢?.
他有.
是吧.
經文沒有描寫他有沒有掙扎.
但我想像.
如果我是他.
我都會很掙扎.

$^{1081}$是吧.
不過.
他最後順服.
是不是.
他最後有沒有掙扎都好.
他最後都順服.
他照樣去.
不單回到他家.
還在底加玻利.
剛才解釋了.
十成.
在底加玻利一帶.
就全面讓耶穌為他做了多麼大的事.
所以.
他不單面對他以前不要他的家人.
他還面對那些同鄉.
是不是.
那些遺棄他的人.
那些寧願要豬都不要他的人.
都很有勇氣.
你想想.
是不是.
他有勇氣面對那些人.
所以弟兄姊妹.
真的.
其實我很欣賞這個格拉森人.
可能他都有很多掙扎.
可能他對耶穌的吩咐都很多不明白.
但他最後順服.
他照樣去做.
我不知道你和你的家人關係如何.
我不知道你有沒有一些.
舊同事舊同學之類.
以前的同鄉之類.
關係很差.
但今天你順了.
你已經是一個新的人.
你經歷了耶穌怎樣憐憫你.
耶穌怎樣拯救你.
請你都回去和好.

$^{1121}$嘗試和他重建關係.
可能當時自己有些不對.
或者明明是他不對.
就當是他全錯.
要恕他.
原因不是因為他改變了.
可能他仍然是那個樣子.
仍然是這樣的人.
但是因為你改變了.
因為耶穌憐憫你.
以致你有勇氣面對.
可能曾經傷害過你的家人.
或者朋友.
和他重建關係.
和他分享.
他可能很好奇.
我們關係很差.
我們一直沒有聯絡.
你突然約我吃飯.
你便可以說.
是呀,因為耶穌基督憐憫你.
你是這樣經歷耶穌的憐憫.
所以你都很想和他保持聯絡.
我想這個見證真的很厲害.
可能有些難度.
你祈禱吧.
學習這個格拉森人.
修補關係.
因為耶穌憐憫你.
所以你願意和他和好.
更加願意將耶穌基督介紹給他.
可能會有掙扎.
這個格拉森人是我們很好的榜樣.
默想一下.
想像一下.
如果你是格拉森人.
你會感受到耶穌這樣憐憫你.
帶著他的信服和勇氣.
回到你以前的家.
回到以前.

$^{1161}$和你相處不好的人.
所以今天讓我們和應一下.
仕途保羅所說.
他說.
「我夢了憐憫.
好讓基督耶穌在我這罪魁身上.
顯明他完全的忍耐.
給後來信他得永生的人做榜樣.
願尊貴榮耀歸給永世的君王.
那不扭壞看不見獨一的神.
直到永永遠遠」.
讓我們一起祈禱.
好.
親愛的主耶穌.
我們讚美你.
你真是至高者神的兒子.
你來到這個世界.
你要釋放我們.
離開一切的污鬼.
離開一切罪惡的捆綁.
讓我們得以自由.
可能我們過去都經歷很多人際關係的失望.
可能我們在別人眼中.
我們覺得自己沒有價值.
你這麼偉大的神.
你都看我們如珠如寶.
求主幫助我們.
懂得這樣去自重.
也幫助我們.
可能過去有些不好的關係.
因著主的憐憫.
我們有勇氣.
我們順服你的吩咐.
去到再次找回他們.
再次將耶穌的憐憫介紹給他們.
我求主你賜福,保守.
多謝你聽我們禱告.
奉耶穌基督的聖名.
阿們.
\newpage



\section{彼得後書 1:3-11}
\label{sec:cnxh3H5LvDs}
\textbf{2023.08.06 《活出敬虔,豐富的生命》 彼得後書 1\_3-11 (和修版) 曾敏傳道}
\newline
\newline
連結: \href{https://youtube.com/watch?v=cnxh3H5LvDs}{\texttt{https://youtube.com/watch?v=cnxh3H5LvDs}} ~~~~ 語音日期: 2023-08-07
\newline
\newline
\hyperref[sec:5NHMIEcWoqY]{\small{< < < PREV SERMON < < <}}
~
\hyperref[sec:index_chronic]{\small{[返順時目]}}
~
\hyperref[sec:index_scriptual]{\small{[返順卷目]}}
~
\hyperref[sec:vjbxBlVTElQ]{\small{> > > NEXT SERMON > > >}}
\newline
\newline
彼得後書 1:3-11
\newline
\begin{longtable}{cl}
\hline
\hline
章節 & 經文 (和合本修訂版)\\
\hline
1:3 & \begin{tabularx}{0.7\textwidth}{X} 神的神能已把一切關乎生命和虔敬的事賜給我們，因我們認識那用自己榮耀和美德召我們的神。 \end{tabularx} \\ \\ \relax
1:4 & \begin{tabularx}{0.7\textwidth}{X} 因此，他已把又寶貴又極大的應許賜給我們，使我們既脫離世上從情慾來的敗壞，就得分享神的本性。 \end{tabularx} \\ \\ \relax
1:5 & \begin{tabularx}{0.7\textwidth}{X} 正因這緣故，你們要分外地努力。有了信心，又要加上德行；有了德行，又要加上知識； \end{tabularx} \\ \\ \relax
1:6 & \begin{tabularx}{0.7\textwidth}{X} 有了知識，又要加上節制；有了節制，又要加上忍耐；有了忍耐，又要加上虔敬； \end{tabularx} \\ \\ \relax
1:7 & \begin{tabularx}{0.7\textwidth}{X} 有了虔敬，又要加上愛弟兄的心；有了愛弟兄的心，又要加上愛眾人的心。 \end{tabularx} \\ \\ \relax
1:8 & \begin{tabularx}{0.7\textwidth}{X} 你們有了這幾樣，再繼續增長，就必使你們在認識我們的主耶穌基督上，不至於懶散和不結果子了。 \end{tabularx} \\ \\ \relax
1:9 & \begin{tabularx}{0.7\textwidth}{X} 沒有這幾樣的人就是瞎眼，是短視，忘了他過去的罪已經得了潔淨。 \end{tabularx} \\ \\ \relax
1:10 & \begin{tabularx}{0.7\textwidth}{X} 所以，弟兄們，要更加努力，使你們的蒙召和被選堅定不移。你們實行這幾樣，就永不失腳。 \end{tabularx} \\ \\ \relax
1:11 & \begin{tabularx}{0.7\textwidth}{X} 這樣，必叫你們豐豐富富地得以進入我們主－救主耶穌基督永遠的國度。 \end{tabularx} \\ \\
[1ex]
\hline
\hline
\end{longtable}
$^{1}$弟兄姊妹平安.
平安.
每個星期能夠回來敬拜.
親近我們的神.
都是神給我們很大的福氣.
在這裡想問大家一個問題.
我看到在座大部分弟兄姊妹.
都已經信耶穌一段時間.
可能有些是很長很長.
有些都不短.
一段日子.
想問大家一個問題.
信完耶穌之後.
神對我們的期望是什麼?.
在這個位置回答就可以了.
大聲說出來.
信耶穌之後.
神對我們的期望是什麼?.
有沒有期望的呢?.
有,這麼大聲.
大聲在這個位置說就行了.
期望是什麼呢?.
什麼?.
永生.
永生是上帝賜給我們的禮物.
信耶穌之後.
神就給我們了.
永遠和他一起.
是禮物.
很好,想起了永生.
神對我們的期望是什麼?.
有沒有人可以回答?.
聽到行到.
哇,很好.
聽神的話語.
又遵行出來.
好,非常好.
還有沒有?.
哇,很好.
信耶穌之後.

$^{41}$上帝對我們的期望.
大家應該可以回答的.
要我們得著豐盛的生命.
回應今天的講題.
非常好.
效法基督.
活得像耶穌.
原來神對我們是有期望的.
所以今天這段經文.
是一段很重要的經文.
是上帝要和我們說話.
有意聽的就應當聽.
其實這段經文.
只到十一節.
今天我想大家做一件事.
請大家翻開程序表.
一路看著這段經文.
因為裡面有些位置.
需要仔細地翻看.
看看神的心意.
如何活出一個敬虔.
很豐富的生命.
生命是豐富的.
生命不是枯乾的.
死著死著的.
也不是天天過著.
很貧窮的生命.
沒錯.
神希望我們有一個.
很渴慕祂的心.
但最終神的心意.
是要我們過得豐富.
神的心意是什麼呢?.
神的心意就是.
敬虔的生活是主所喜悅的.
而且神呼召我們.
去過一個敬虔的生活.
神很喜悅我們過得敬虔.
我們去看看吧.
在這裡有一個很重要的前設.

$^{81}$首先第一件事.
你看回經文裡.
提到一系列.
這段經文其實有結構.
有組織.
結構非常清晰.
我們跟著結構去.
神想我們過敬虔的生活.
首先第一件事.
是上帝先呼召我們.
這裡你看一章三節.
因我們認識.
那用自己榮耀和美德.
召我們的神.
神如何呼召我們呢?.
是用祂自己的榮耀大德.
是祂奇妙的作為.
上帝不是就這樣呼召我們.
是用祂的榮耀大德.
去呼召我們.
呼召我和你.
在呼召之後.
有一個箭咀說.
神的神能.
把一切關乎生命和虔敬的事.
賜給我們.
呼召我們之後.
祂又賜給我們什麼呢?.
關於敬虔生活的事情.
祂賜給我們.
在這個過程.
上帝是用祂的神聖大能.
神呼召.
神以祂的大能.
告訴我們.
如何過敬虔的生活.
上帝很想我們知道.
第三件事很重要.
上帝賜祂寶貴極大的英雄.
這裡提到一章四節.

$^{121}$因此他把寶貴又極大的英雄.
賜給我們.
這個寶貴又極大的英雄.
是在說什麼呢?.
是在說到最後.
上帝完完全全.
拯救我和你的生命.
的這個英雄.
我和你完全得蒙救贖的英雄.
上帝說要給我們.
這個他要給我們.
去呼召我們.
他叫我們過.
一個非常敬虔的生活.
整個生命的質素.
要好像上帝一樣.
然後神要給我們英雄.
是我和你得到最終極的禮物.
就是完全徹底的得勝.
去得救.
在這個大前提下.
神讓我們清楚看到一個對比.
上帝說要我們過敬虔生活.
為什麼呢?.
有一個對比.
我們要脫離情慾的敗壞.
要脫離的意思.
代表什麼呢?.
過去我們的光景.
就是在情慾的敗壞裡面.
可能你說.
不是.
我一直都沒有這些問題.
我哪有這些問題.
我不知道有多好.
在人心目中.
我是人家很乖的兒子,女兒.
我是人家很乖的老公,老婆.
是吧?.
是很好的爸爸,媽媽.

$^{161}$我去奉公守法.
哪有什麼敗壞情慾.
但是弟兄姊妹.
我們要好好想想.
如果是這樣的話.
耶穌基督的救恩.
和我們沒有關係.
是吧?.
為什麼耶穌基督要拯救全世界的人呢?.
是全世界的人.
每個人都需要耶穌基督的救恩.
就是因為我們每個人都犯罪.
得罪上帝.
這是羅馬書說的.
世人都犯罪.
虧欠了神的榮耀.
世人每一個都有犯罪.
我和你都不例外.
我和你以前就在這些情慾的敗壞裡.
今天信主之後.
上帝要我們的生活脫離.
這是一個動詞.
不是一個口號.
我們要追求聖潔的生活.
不是一個口號這麼簡單.
是一個動作要脫離.
敗壞的情慾.
然後一個轉向.
一個動詞就是.
要分享神的本性.
什麼叫分享神的本性?.
這裡是一章四節說.
要分享神的本性.
神的本性不單是說神的神性.
是說他的道德.
他的品格.
是說他的性情.
原來上帝很想我們.
在他的性情上有份.
一方面要脫離過去的所有罪惡.

$^{201}$另一方面要在上帝的性情裡有份.
就好像子女像爸爸媽媽.
神一開始的心意就是這樣.
有時候你看到不同人的子女.
總是說他的眼耳口鼻.
好像爸爸媽媽.
像爸爸像媽媽.
總是這樣說.
他的身形.
總是說話的態度很像.
性格更像.
兩父子.
母女喜歡的東西一樣.
對人的處事是這樣的.
很像.
其實今天上帝也想我們跟他相似.
特別是相似他的性情.
他怎樣對待人.
怎樣看事情.
上帝想我們相似他.
今天問問我們自己.
在上帝的本性上.
在上帝的性情上.
我們分享了多少呢?.
我們像不像上帝?.
這是很重要的.
像不像上帝呢?.
上帝是想我們像他.
正因為這個緣故.
一章五節這樣說.
正因為這個緣故.
你們要憤外的努力.
在說回應.
由一章五節開始說回應.
一章三到四節.
是說上帝的呼召.
上帝對我們的呼召.
上帝對我們的心意.
是希望我們過敬虔生活.
希望我們脫離情慾的敗壞.

$^{241}$希望我們在他的本性上有份.
我們的回應.
這裡很大段是說我和你的回應.
所以今天說完之後.
我們要在這個房間.
活出我們對上帝的回應.
回應是什麼?.
我們需要在道德本性上.
憤外的努力.
這個憤外的努力.
不是說加一兩錢肉緊.
是說盡力去做到最好.
剛才我在祈禱會.
我都有說.
最近一次培靈賢經大會.
有一次應該是晚上的.
那個講員.
粉輕會的講道.
那個講道是非常有刺激性.
但是引發人很多的心思.
那個講道是在說教會.
教會是怎樣的?.
教會不能夠吸引人回來聚會.
為什麼?.
因為我們的生命.
我們的生命的虛偽.
我們的生命犯罪.
不能夠讓人看到見證.
我們的生命不能夠吸引人回來聚會.
不是這個場地的問題.
不是什麼配套的問題.
不是我們有沒有活動的問題.
是生命有沒有吸引力.
而生命的吸引力是在說我們像不像上帝.
講員說得很有刺激性.
引發我們去思考.
是不是這樣?.
是不是人家去想的時候.
想想教會的人是怎樣的?.
見證是怎樣的?.

$^{281}$這些是令我們有很多思考.
最近我也是很刺激的一件事.
我收到一個電郵.
電郵的內容不方便詳細說.
裡面不是說我們教會的情況.
是外面有教會的情況.
是一些不理想的情況.
發電郵的人很想去伸冤.
發電郵的人很想去發出一個呼求.
很想請大家關注.
有這樣的事情發生.
我心裡很擇心也很難過.
屬上帝的人.
是不是相信了耶穌.
行了 我已經有永生的這份禮物.
好像拿了護照一樣.
我可以回到天家.
我可以在神的家裡.
很安全的生活.
沒問題的了.
行了 就做回以前的自己.
如果是這樣的話.
難怪我們不能吸引別人歸向我們.
所以這裡說要奮愛努力.
盡一切努力.
在道德的品性上要追求.
要活出來.
這裡說了一連串.
怎樣追求.
有個方向的.
有八種基督徒的品德.
今天大家數數.
看看每一個品德都為自己評分.
現在我去到哪裡為自己就好了.
是不是沒提.
我覺得你這方面.
我估計你大概三分.
我覺得前面那個人.
我看死他了 只有四分.
我覺得自己有十分.

$^{321}$不要做這些事.
你跟自己說好了.
你身邊的人交給神.
但說自己.
八種基督徒的品德有多少呢.
這裡.
它是一層一層去推進的.
但你說我們要找找邏輯性在哪裡.
為什麼從一件事.
我就跳到下一件事.
你說是不是完全可以解釋.
不一定.
但這八樣東西很重要.
排頭位的是信心.
信心 什麼是信心呢.
信心是我們所有的.
基督徒道德的.
最重要的一個基礎.
這個信心是說.
我相信上帝什麼.
相信上帝的救恩.
我相信上帝的應許.
我相信上帝的能力.
你問問自己.
對上帝有多大的相信.
什麼叫相信上帝的能力.
現在我很風平浪靜.
我講這件事很大聲.
我相信上帝的能力.
我相信上帝的救恩.
我相信祂的應許.
一道風浪來的時候.
馬上怎麼樣.
上帝去了哪裡.
上帝不見了.
上帝不幫我.
祈願禱告也不出手.
所以我有多相信.
相信上帝的能力.
很害怕吧.

$^{361}$其實紅恩堂經歷了很多風浪.
我為在座的弟兄姊妹獻上感恩.
經歷風浪之後.
大家仍然在這裡.
大家仍然心向上帝.
我為你們感恩.
這正是考驗人的時候.
經歷風浪.
是否第一時間就離開了.
又有風浪.
不行了.
怎麼辦.
我們會倒塌.
倒塌什麼.
神還在.
多相信上帝的能力.
相信祂的應許.
我將我所有的東西.
我去依靠祂.
我相信神一定會幫我.
至於祂怎麼幫.
祂的時間怎麼樣.
這就交給神.
這是第一點.
信心.
有信心的人.
也是值得信賴的人.
說真的.
我很少找到一個.
那個人沒有信心.
但不知為何.
他特別值得信賴.
我很少找到這樣的人.
通常有信心的人.
很值得信賴.
本身這些人.
也是誠信可靠的.
做事很順正.
很穩妥.
我是否一個有信心的人.

$^{401}$站在你旁邊.
是否別人感覺到.
嘩!這個人對上帝有信心.
令我也很想去依靠上帝.
還是搖來搖去.
那個神不行.
我也不想去信他的神.
是誰呢.
第二點.
有德行.
德行是指什麼呢.
在道德上表現得很優秀.
很卓越.
有很好的名聲.
這個要努力去爭取.
我們有德行.
不是單靠自己的努力.
我們的德行.
建基於耶穌基督的德行.
耶穌基督是完美的一個典範.
一個德行的典範.
也是藉著主.
靠著耶穌基督.
我們去分享上帝道德的本性.
所以和神的關係要很近.
很近很近.
但是要有這個德行.
很多時候.
在基督徒當中.
這一環.
都是會失守的一環.
就是因為沒有很好的品德.
結果在這個位置失見證.
在這個位置.
也會被人捉住大做文章.
去抨擊教會.
所以我們要非常留意.
要有德行.
我們今天在這方面.
有很好的品德.

$^{441}$之後要有知識.
要有知識.
這個知識不是說.
我讀書不是很多.
我這方面有很多缺陷.
這裡不是在說讀書的層面.
是因為耶穌基督而有的知識.
因為認識耶穌基督.
我懂得分辨.
好與壞.
是與非.
不是憑藉自己的血氣.
是憑藉自己的血氣.
憑藉自己的感覺.
我這麼多年做人.
我覺得這些是對的.
那些是不對的.
我是這樣的.
我真的覺得這樣.
不是.
而是你真的認識耶穌基督.
你真的能夠從主的角度去分辨.
什麼是真理.
什麼不是.
我懂得去遵守真理.
這個是很重要.
當異端來的時候.
我懂得分辨.
其實有一樣東西是不容忽視的.
是輕看的.
如果大家是從元朗方向.
向屯門方向乘巴士.
來洪水橋的話.
大家下了隧道.
去到洪水橋站.
然後下隧道.
走過來.
向紅印堂方向走.
其實那個隧道.
我都不止一次.

$^{481}$見到異端在那裡擺攤位.
擺攤位.
是嗎.
是有人搖頭.
他們是很積極進取的.
不停派他們的單張.
其實是很大影響的.
我都試過離開教會的時候.
向著輕鐵站方向.
見到他們.
又有穿著西裝.
大家有沒有見過.
那個牌子.
原來他們又收工了.
他們收工不是教會那邊.
我都知道他們的教會在哪邊.
距離我們不是很遠.
大概一個輕鐵站.
在那邊.
派人來到這邊.
去做.
其實我們懂不懂得分辨呢.
在人行當中.
懂不懂得分辨呢.
我記得以前最震撼的一件事.
就是有一段時間.
以前舊的名字叫東方閃電.
後來是全能神的教會.
我不知道現在叫什麼.
國道還是什麼.
我忘記了那間異端.
起初是東方閃電.
那間異端.
是很懂得潛伏在教會裡面的.
我記得以前.
我還住在上水圍的時候.
我就不知道為什麼有一次.
總之是因為要坐巴士的緣故.
認識了一個說普通話的婆婆.
我就想著.

$^{521}$很想容易一點認識人.
可以容易一點講福音.
自己就這樣想.
那個婆婆這麼好聊.
就跟她聊聊天.
一直聊聊聊.
認識了她.
都不知道她大有來頭.
後來說著說著.
有一次我跟母會那邊聯絡.
說我在上水廣場下面的巴士站.
認識了一個婆婆說普通話.
我跟她說得很好聊.
母會的牧師給我看.
是不是這張照片.
我說是呀.
東方閃電的領袖.
真的這樣.
她住在我上水圍附近.
她都知道我住在哪裡.
自此之後很刺激.
她說她很厲害.
她跟她的肢體.
潛伏在母會很長一段時間不說話.
直到某一天.
她忽然間才說.
弟兄 姐妹.
你覺得我們需要多點學習聖經嗎.
需要.
我們一起來查經好嗎.
你來我家好嗎.
好啊.
就這樣開始.
後來才發現.
原來是東方閃電的領袖.
後來更精彩的是.
她記得我住在哪裡.
有一天親自拜訪我.
這個婆婆帶著一位女士.
走來敲我的門.

$^{561}$說我要跟你說一個消息.
你知不知道將來主怎樣回來.
一個重要的消息.
一直說.
我講到一個地步.
我都在神學院接受過裝備.
我是傳道人.
她完全無視我這件事.
只有唯一一個得救的方法.
就是要相信她所說的話.
去她的教會.
當日真是多精彩.
後來我發覺.
沒有必要再繼續對話.
我相信這個對話可以持續幾日幾夜.
大家不眠不休都可以.
我打算要放下.
離開他們.
我好記得.
但我一直在想一件事.
如果今天是弟兄姊妹遇到.
或者傳道人很流行.
一直被他們迷惑.
那又怎樣呢.
所以這個知識是很重要.
要懂得分辨.
這個知識是怎樣來的.
弟兄姊妹.
你告訴我這個知識是怎樣來的.
是從神來的.
非常好.
從神來.
從聖經.
要多讀聖經.
多聽多看多學習.
這樣我才懂得分辨.
否則.
有很多人很熱心.
來我家.
我和你查經.

$^{601}$弟兄姊妹.
你真是熱心.
你不是傳道人.
你這麼好.
和我查經.
突然在我們教會裡.
有這麼好的弟兄姊妹.
我還不去你家.
這樣就糟了.
不是每個人都是這樣.
但要懂得分辨.
這些似是而非的東西.
你不知道是甚麼.
在彼得後孫這裡.
特別講到對那些假教師.
要很小心.
在今天對我們來說.
分辨異端都要很小心.
其實這些人都在我們附近出入.
我記得李牧師當初.
帶著一個團隊過來直堂的時候.
起初其中一樣東西.
他當時的心願.
都是要對付異端.
要很小心.
直到今天這件事沒有改變.
我們需要有知識.
而這個知識對我們來說.
亦是非常重要.
知識之後很重要.
有節制.
大家即時想.
可能想飲食.
飲食要有節制.
因為影響我們的健康.
我們的身形.
這裡不單是講飲食.
健康.
我講的是對付情慾.
是很多方面的.

$^{641}$有些人是飲食的.
有些人是對其他東西.
有那種慾望.
很強烈的.
慾望大過所有東西.
你問問自己.
在自己生命裡最渴望的是甚麼.
我講的是很多物質上.
很多心靈上.
一些渴望.
一些慾望.
是否已經去到一個地步.
成為我們的王.
我們的主.
我們的偶像.
節制的意思.
你要能夠控制這些慾望.
你不要讓慾望反過來控制自己.
我一直為了那些東西.
看不同情況.
每個人喜歡的東西不同.
有些人是很喜歡音響器材.
整間屋都佈滿了.
要看自己生活的情況.
有人是食物.
有人是衣服.
有人是不同的喜好.
心頭好.
不同的東西.
我不是說我沒有興趣了.
甚麼東西都不要買了.
不要用了.
不是這個意思.
而是你對那個慾望.
渴望到一個地步.
它已經成為你的王.
你的主.
你就要想想.
節制.
你要控制它.

$^{681}$不要讓它操控你.
節制之後又忍耐.
忍耐是要做甚麼呢?.
在現在這個時勢.
更加需要忍耐.
忍耐等候耶穌基督再回來的日子.
在這個之前.
有很多的艱難.
真的艱難.
其實我回看香港這幾年.
很多的艱難.
是各種不同層面的.
有政治上的.
這一年.
特別是這半年.
有一樣東西.
香港是很令人害怕的.
是甚麼呢?.
看新聞.
令人心寒.
是甚麼呢?.
斬人.
是很令人害怕的.
是嗎?.
很暴力.
好像開了個按鈕.
我看到香港人.
有些位置.
動不動就來.
是嗎?.
是很大的傷害性.
令人很恐懼.
這只是其中一點.
你看到現在整個局面.
是很混亂.
要忍耐.
是嗎?.
忍耐等候耶穌基督再回來.
還有.
不要讓人.

$^{721}$動搖到我們.
不要讓那些冷嘲熱諷者.
動搖到我們.
這個冷嘲熱諷者.
就是在彼得後書裡面.
所提到的那些人.
我們要堅持過聖潔.
敬虔的生活.
直到新天新地的來臨.
直到主耶穌基督再回來的日子.
我想說這件事並不容易.
因為整個大氣候,大環境.
都不是在說聖潔敬虔生活.
這個世界.
很多人都在說人不為己.
天誅地滅.
第一時間當然是為自己.
哪有這麼傻.
社會在提倡這些.
有一件事.
亦是很高難度.
我特別覺得在家庭裡.
弟兄姊妹會面對得比較多.
是甚麼?.
其實家人本身都有其他宗教信仰.
在生活的層面裡.
這些位置都會起衝突.
會有的.
有一些可以和平共處.
你不要教我.
我不要教你.
我尊重你.
你尊重我.
有些是這樣的.
有些是不行的.
你一定要這樣.
要跟他的那一套.
這個時候你怎樣算呢?.
這裡是需要很大的忍耐.
但亦需要我們在上帝面前有忠心.

$^{761}$這是上帝的心意.
忍耐.
不要被人動搖我們的信仰.
接著到虔敬.
虔敬是耶穌基督.
已經將關於過敬虔生活的事.
告訴我們.
我們就要活出來.
要在我們的生活上.
反映上帝的屬性.
我要說一件事.
基督徒不只是做好人.
為什麼我這樣說?.
弟兄姊妹回應我.
如果基督徒只是做好人.
有什麼問題?.
上帝的屬性.
只是一個簡單的好字嗎?.
世界其實有很多好人.
我以前去埃及.
有一個跨文化體驗的時候.
我認識了一個女士.
她信奉伊斯蘭教.
我會用好人形容她.
非常溫柔.
比我所認識的基督徒溫柔.
她的性格非常溫柔.
她是一個好基督徒.
溫柔,好人.
什麼好事都做.
夠不夠?.
所以我們經常告訴別人.
我們要做好人.
絕對不足夠.
信耶穌不是說做好人.
信耶穌是說.
有耶穌基督的救恩在我生命裡.
按照上帝的標準去生活.
活出一個聖潔的生活.
這應該要超越現在.

$^{801}$一般其他宗教信仰所說的好人.
應該要超越才是.
而且是說耶穌基督救贖的恩典.
在我生命裡.
我生命有力所不逮的時候.
我做不到的時候.
是耶穌基督的恩典.
在幫助我,奉給我.
是不一樣的.
所以這裡說的是.
敬虔活出上帝的屬性.
是什麼屬性?.
應該是超越的.
我們的層次不要只留在社會.
一般人說的好人.
如果這樣.
別人也會跟你說.
說好人?.
我很好人.
那我為什麼要信你的上帝?.
很多人都會說.
我比基督徒還要好.
那我為什麼要信你的上帝?.
做好人?.
不是的.
上帝的屬性比這一切都超越.
上帝的屬性比任何東西都超越.
我以前初期信耶穌的時候.
就覺得.
未信耶穌也有很多人是好人.
但信耶穌.
越來越看聖經的標準.
不是的.
我覺得上帝那種聖潔.
才是真正徹底的聖潔.
那個聖潔是很高的.
那個道德性是很高很高的.
遠遠超過今天一般人所說的.
所謂的好人.
現在今天我跟你.

$^{841}$生命要反映出.
上帝那種崇聖.
那種聖潔的性情.
我們就向著上帝的標準去.
說的是靠神.
靠著我們裡面的聖靈.
靠著神的能力.
去活出那種屬性.
在這個乾淨之後.
去到後面兩點.
是和愛有關.
愛弟兄的心.
什麼叫愛呢?.
我們彼此的愛顧.
我們彼此的守望.
彼此的接納.
有一樣東西很重要.
我們不要競爭.
如果說愛.
也可以很空泛.
我就算說那句話.
也可以很空泛.
我們彼此照顧.
彼此接納.
認識了.
我們照顧.
我們接納.
難在哪裡?.
如果有利益衝突的時候怎麼辦?.
有競爭的時候怎麼辦?.
我們彼此不要競爭.
是不是?.
我怎樣?.
我比你好.
你比我差.
要不你比我好.
要不我能夠怎樣.
你能不能夠怎樣.
我做什麼.
你做什麼.

$^{881}$不要競爭.
我們不是這樣.
彼此互相去消磨對方.
而是我們.
情同手足.
你知不知道?.
我們只有一個仇敵.
是誰?.
撒旦.
只有一個仇敵.
所以你旁邊的人.
我永遠都不會和他競爭.
我不會一直計算他.
我們要對付自己的妒忌.
說的是我們自己.
不要看著別人.
對付自己的妒忌.
對付自己的憎恨.
彼此相愛.
為別人祝福.
弟兄姊妹我為你感恩.
你做得好.
我支持你.
我能做也好.
不能做也好.
不要緊.
我愛弟兄.
我愛姊妹.
只有愛.
我們中間不會彼此.
一個競爭對手.
互相插對方.
弄死對方.
不會有這些.
這是很重要的.
愛弟兄的心.
愛眾人的心.
這一點到最後.
是可以解釋的.
是包括所有.

$^{921}$是愛.
今天我相信.
如果真的要吸納新朋友.
來到我們當中.
是需要有很大的愛.
我們彼此愛對方.
其實我最近在弟兄姊妹當中.
看到一份很大的愛.
我看到有弟兄要去看醫生.
我們另外有弟兄陪診.
有弟兄姊妹說.
我要買食物給你.
幫你買東西.
我要送上來.
你有什麼事.
我隨時幫助你.
如果有人有需要.
我們去探望他.
帶什麼東西.
我很感動.
其實這就是愛.
我們彼此有禱告.
這就是愛.
這就是最吸引的.
是吧?.
這就是最吸引的.
所以在這裡.
我鼓勵我們彼此之間.
彼此互相鼓勵.
無論過去的日子.
我們有什麼嚴規也好.
什麼不開心也好.
我們全部都要靠上帝對付它.
我們要復和.
我剛才說的.
不要有競爭.
不要有妒忌.
有不開心.
當面說.
面對面的責備.

$^{961}$強於背地的愛情.
是吧?.
與其背後說.
不如當面說.
有事就當面說.
處理它.
這是很重要的.
我們要有一個.
這個愛是非常重要的.
這個愛包括了所有.
我們道德的品行裡.
但要小心一件事.
愛不等於縱容.
上帝既是愛.
但上帝也是聖潔的.
你說不要緊.
有什麼犯罪也好.
不要緊.
隨它吧.
你說愛嘛.
如果單是這樣的愛.
就是縱容.
這不是上帝的愛.
要麼就沒有愛心.
插死對方.
這是嚴苛.
既是愛.
也是聖潔.
既是聖潔.
也是愛.
上帝的愛是如此完全.
這八種基督徒的道德.
大家有嗎?.
我們有多少呢?.
有兩種生命的對比.
兩種生命的對比.
我們看經文裡.
第一種生命的對比就是.
一章八字.
你們有了這幾樣.

$^{1001}$這八種基督徒的道德.
這個品行我們都有了.
作者要鼓勵我們.
要繼續增長.
這個增長不單是增加.
有更多的意思.
這個增長的意思是什麼呢?.
你要滿入.
你這種基督徒的道德.
那種品行.
要滿寫出來.
你整個生命.
被人看到有這些德行.
人家看到的.
最重要的那種增長的意思.
就必使你們在認識.
我們主耶穌基督上.
不至於懶散和不結果子.
懶散和不結果子.
是同一件事.
是沒有效益的.
沒有好處的.
當我們有這些德行.
我們有不斷的成長.
成長到一個地步.
人家自然就看到.
有好見證的.
這樣我們就一定可以.
不會不結果子.
真的會結果子.
我們的生命是豐富的.
結果子的生命.
有效益的生命.
是上帝眼中看到的效益.
不是世人眼中看到的那種.
如果有.
就會這樣.
這是第一種生命.
第二種生命.
是剛剛相反.

$^{1041}$沒有這幾樣東西的人.
沒有這些道德品行的人.
是乞眼短視的.
什麼叫乞眼短視呢?.
其實是我們不認識真理.
我們不是認識真理.
我們是盲的.
我們忘記了過去的罪.
已經得了潔淨.
這裡的意思.
令人想起洗禮.
潔淨令人想起洗禮.
人信主.
信而受洗.
歸入主的名下.
歸入教會.
這是何等的寶貴.
是神的恩典.
這些人沒有基督徒的道德.
這種品行的人.
他們是不認識真理.
忘記了自己領受了神的救恩.
在神眼中看到.
這些是嚴重的情況.
非常嚴重.
是很大的對比.
不會結果一致.
也不會有豐盛的生命.
在神眼中看到.
這些人不是敬虔的人.
這些生命.
上帝不棄穴.
所以一章十節.
那裡很重要.
作者鼓勵我們.
弟兄們要更加努力.
要加倍努力.
要實踐.
有耳聽就應當聽.
要實踐出來.

$^{1081}$你們實行這幾樣.
走出來.
不要只聽.
剛才有個姐妹說得很好.
聽道 行道.
追求這幾樣生命的特質.
使你們的夢照和被選.
堅定不移.
堅持你的信仰.
你實踐.
你堅持.
永不失腳.
你不會被打敗.
不失腳就是不會失敗.
不會被打敗.
在我們信耶穌之後.
上帝很歡喜快樂.
為了我們的生命.
得蒙救贖.
為了我和你的得救.
歡呼聲.
很開心.
為了我們得蒙救贖.
很開心.
其實我們也很開心.
我們得蒙救贖.
但你要知道.
惡者也很想打敗我們.
令我們失敗.
有很多很多生命路上的考驗.
挑戰.
要永不失腳.
就是要實踐主的道.
走出來追求這些德行.
堅持.
然後我們就永不失腳.
我們不會失敗.
我真的很渴望.
今天在地面上和大家敬拜神.
有一天在天家的時候.

$^{1121}$我也是很渴望和大家一起.
雖然我們彼此.
回到天家的時間不會一樣.
但也很渴望有一天.
在天家的時候.
我們也可以這樣去重聚.
不單止這樣.
我們和更多的基督徒.
屬上帝的人.
一路在天家敬拜讚美.
我們永不失腳.
上帝保守我們直到最後.
這裡很重要.
最後最後的一句話.
最後十一節這樣說.
「這樣必叫你們豐豐富富地.
得以進入我們救主耶穌基督.
永遠的國度」.
這個比較好的翻譯.
聽我說.
進入我們救主耶穌基督.
永世之國的福.
就會豐豐富富.
供應你們了.
這個比較好的翻譯.
在這裡說.
這一句是在說什麼呢.
說了這麼多.
作者鼓勵我們.
怎樣回應上帝的呼召.
回應神的心意.
要我們過敬虔生活.
我們要追求八種基督徒的德行.
我們要實踐.
我們要堅持.
直到最後我們不會失腳.
但最後那一句總結很重要.
是什麼呢.
今天我們能夠得到救贖的恩典.
最後最後的得勝.

$^{1161}$在於神的恩典.
是神的恩典.
沒錯.
神要我和你努力.
在德行上.
要我和你努力成長.
但是那個徹底的得救.
那一件事.
是耶穌基督為我和你成就.
是恩典.
所以他這樣說.
他說進入我們救主耶穌基督.
永世之國的福氣.
是會豐豐富富.
供應你們.
這個福氣是神給我們的.
不是我和你賺取得來.
我和你怎樣努力追求.
那個德行也好.
不是因為這樣就賺取了.
這個永生的禮物.
和賺取最後的永遠的得救.
不是.
那你說那還需要努力.
那個德行.
要努力.
神的心意想我們這樣生活.
神想我們像他.
在地面上就是要這樣生活.
是不是.
但是最後最大的福氣.
是上帝已經給了我和你.
這個是非常非常寶貴.
所以我們彼此勉勵.
很盼望今天弟兄姊妹.
我和你要去回應.
那八種基督徒的品格.
我和你今天的評分是多少.
評自己就好了.
有沒有成長的空間.

$^{1201}$經文說要用盡全力去做.
是神的心意.
有耳聽的就應當聽.
我們一起祈禱.
對我們的人生有美好的旨意.
你不是只要我們得蒙救贖.
就停在那裡.
你想我們過敬虔的生活.
你想我們過得豐富.
是你眼中的豐富.
你想我們像你.
有你的性情.
有你的品格.
你想我們有那種八種的.
基督徒的道德.
天父啊.
我求你.
我求你兼顧我們眾人的心.
讓我們今天聽了道之後.
我們感到責心.
我們走出這個門口不是忘記了.
而是真的要實踐出來.
每一天活著.
越來越像你.
別人見我們的見證.
想起上帝你.
別人見到我們整個教會的群體.
心被吸引.
要回來這裡去認識你.
天父願意你給我們有真正的愛.
是聖潔的愛.
是追求真理的.
走在上帝你道路上的愛.
讓我們無論在哪裡都好.
都顯出上帝你的榮耀.
特別在洪水橋區.
我們能夠為你作炎作光.
求主賜福我們.
大大的保守我們.
奉耶穌基督的名字祈禱.

$^{1241}$阿們.
\newpage



\section{使徒行傳 17:16-31}
\label{sec:vjbxBlVTElQ}
\textbf{2023.08.13 《進入人群的福音》 使徒行傳 17\_16-31 (和修版) 講員 : 陳麗蟬牧師}
\newline
\newline
連結: \href{https://youtube.com/watch?v=vjbxBlVTElQ}{\texttt{https://youtube.com/watch?v=vjbxBlVTElQ}} ~~~~ 語音日期: 2023-08-16
\newline
\newline
\hyperref[sec:cnxh3H5LvDs]{\small{< < < PREV SERMON < < <}}
~
\hyperref[sec:index_chronic]{\small{[返順時目]}}
~
\hyperref[sec:index_scriptual]{\small{[返順卷目]}}
~
\hyperref[sec:cIHnc4_JX1s]{\small{> > > NEXT SERMON > > >}}
\newline
\newline
使徒行傳 17:16-31
\newline
\begin{longtable}{cl}
\hline
\hline
章節 & 經文 (和合本修訂版)\\
\hline
17:16 & \begin{tabularx}{0.7\textwidth}{X} 保羅在雅典等候他們的時候，看見滿城都是偶像，就心裡非常難過。 \end{tabularx} \\ \\ \relax
17:17 & \begin{tabularx}{0.7\textwidth}{X} 於是他在會堂裡與猶太人和虔敬的人，以及每日在市場上所遇見的人辯論。 \end{tabularx} \\ \\ \relax
17:18 & \begin{tabularx}{0.7\textwidth}{X} 還有伊壁鳩魯和斯多亞兩派的哲學家也與他爭辯。有的說：「這胡言亂語的要說甚麼？」有的說：「他似乎是宣傳外邦鬼神的。」這是因保羅傳講耶穌與復活的福音。 \end{tabularx} \\ \\ \relax
17:19 & \begin{tabularx}{0.7\textwidth}{X} 他們就把他帶到亞略巴古，說：「你所講的這新學說，我們也可以知道嗎？ \end{tabularx} \\ \\ \relax
17:20 & \begin{tabularx}{0.7\textwidth}{X} 因為你有些奇怪的事傳到我們耳中，我們想知道這些事是甚麼意思。」 \end{tabularx} \\ \\ \relax
17:21 & \begin{tabularx}{0.7\textwidth}{X} 原來所有的雅典人和居住在那裡的外國人都無暇管別的事，只是談談或聽聽新聞。 \end{tabularx} \\ \\ \relax
17:22 & \begin{tabularx}{0.7\textwidth}{X} 保羅站在亞略巴古當中，說：「諸位雅典人！我看你們凡事很敬畏鬼神。 \end{tabularx} \\ \\ \relax
17:23 & \begin{tabularx}{0.7\textwidth}{X} 我到處走走的時候，仔細觀察你們所敬拜的，發現一座壇，上面寫著『獻給未識之神明』。你們所不認識而敬拜的，我現在向你們宣告： \end{tabularx} \\ \\ \relax
17:24 & \begin{tabularx}{0.7\textwidth}{X} 他是創造宇宙和其中萬物的神；他既是天地的主，就不住在人手所造的殿宇裡， \end{tabularx} \\ \\ \relax
17:25 & \begin{tabularx}{0.7\textwidth}{X} 也不用人手去服侍，好像缺少甚麼似的；自己倒將生命、氣息、萬物賜給萬人。 \end{tabularx} \\ \\ \relax
17:26 & \begin{tabularx}{0.7\textwidth}{X} 他從一人造出萬族，居住在全地面上，並且預先定準他們的年限和所住的疆界， \end{tabularx} \\ \\ \relax
17:27 & \begin{tabularx}{0.7\textwidth}{X} 為要使他們尋求神，或者可以揣摩而找到他，其實他離我們各人不遠。 \end{tabularx} \\ \\ \relax
17:28 & \begin{tabularx}{0.7\textwidth}{X} 我們生活、行動、存在都在於他。就如你們的詩人也有人說：『我們也是他所生的。』 \end{tabularx} \\ \\ \relax
17:29 & \begin{tabularx}{0.7\textwidth}{X} 既然我們是神所生的，就不應該以為神的神性像人用手藝和心思所雕刻的金、銀、石像一般。 \end{tabularx} \\ \\ \relax
17:30 & \begin{tabularx}{0.7\textwidth}{X} 世人蒙昧無知的時候，神並不追究，如今卻吩咐各處的人都要悔改。 \end{tabularx} \\ \\ \relax
17:31 & \begin{tabularx}{0.7\textwidth}{X} 因為他已經定了日子，要藉著他所設立的人按公義審判天下，並且使他從死人中復活，給萬人作可信的憑據。」 \end{tabularx} \\ \\
[1ex]
\hline
\hline
\end{longtable}
$^{1}$主席.
最緊要的是哪裡風景好,快點去搶位拍照.
還是一早上網找好地方吃好東西.
今早看新聞,上深圳的人短假期六星期一有四十萬.
上去的酒樓就滿了,反而香港的很多都沒什麼人在那裡吃東西.
然後最緊要是去購物.
但是剛才主席為我們讀出那段經文.
告訴我們保羅去到雅典旅行.
不過他去到的時候就不是來看這個很宏偉的建築物.
因為在雅典裡面,在當時公元二三百年的時候.
他們已經有很宏偉的建築物.
是不是按這個呢?.
是,有很宏偉,不過現在就剩下一些爛柱.
也看到很多很雕刻很漂亮的神像.
不過他不只是只顧著去拍照.
他在當中心裡非常的焦急.
其實這次去雅典的時候.
他去到的裡面不是去旅行,不是去參觀.
而是去旅行報道,在他的生命裡面第二次去旅行報道.
本來他和提摩太和西拉就一直去.
他首先去到這個菲臘比的地方.
在我們圖上的左上教會看到一塊紅色的.
他去到的地方有菲臘比,帖撒羅尼迦和比利亞這三個城市.
但他去到每一個城市裡面都會被人抓,鎖,打.
因為人們誤信耶穌基督復活.
覺得他在這裡搞定了.
弄得這三個城市都非常的混亂.
他們都陷在一個生命非常危險的地方.
在迫不得已之下,本來要自己一個人離開他們.
先行去到雅典.
雅典是希臘的首都.
也是希臘這麼多城市裡面最大的城市.
是一個文化的中心.
當中哲學很出名,藝術和戲劇都很出名.
他去到一個這麼有文化,藝術的地方.
為什麼他心裡這麼難過呢?.
他難過或是心急.
去到一個警方說情緒很高漲.
警方說他好像抽筋一樣.
說他很激氣,激心.

$^{41}$當他覺得在這麼有文化的地方.
但裡面整座城市都是偶像.
就是說裡面整個城市的人.
他們的盼望,思想就只是和這些偶像有關.
他很心痛.
我們的上一代,我自己的父母.
他們是先拜祖先,先拜關公,先拜王大先.
但是經過了很多很多的禱告.
也有很多的見證之後.
他們就信了耶穌,洗禮入了教會.
今年年初,我和我的家人,親戚一起去旅行了.
我發覺到現在的年輕人又怎樣呢?.
我發覺到這些三十多歲的年輕人.
他們玩這個瑜伽玩多幾年就深入了.
他們去到一個很高的境界.
他們竟然和我父母以前一樣.
從初一到十五是吃素的.
我心想,哇!又會這樣?.
另外一個年輕人,他在香港我從來沒有發覺到.
那次我們去日本旅行.
我們自己計劃去旅行.
沒有說過要進神廟.
但是什麼時候走過廟宇.
他就說你們向前走,我一會兒會追上.
他每次看到神廟的時候,原來又要進去拜一拜.
而且他不吃牛肉.
我心想,為什麼會走回他父母的年代.
不過現在就比較新潮,不再是祖先,王大仙,關公那些.
去到另一個,但是都是一個.
他自己都不知道在做什麼的境況.
當時我都很心痛.
其實這班年輕人,我以為可以脫離父母以前的境況.
讀了無論小學或中學都讀基督教的.
但是為什麼讀了這麼久.
知道了,但都不肯相信.
反而去玩一下瑜伽,玩了幾年就覺得很高深.
不知道什麼時候認識了日本的東西.
才可以信到十足呢.
當時我心靈裡面很難過.
當然沒有像保羅那樣難過.

$^{81}$心裡面很痛,好像抽筋一樣.
這個心裡面的難過.
和這個字相關的字就在出埃及記第20章第5節裡面.
當摩西帶領著以色列人出埃及.
準備進入,還沒有準備,他們要走很久.
去進入迦南地的時候,上帝已經吩咐他們.
不可跪拜那些像,也不可以事奉他.
因為迦南地有很多偶像.
耶和華說到明了.
因為耶和華說你的神是忌邪的神.
不喜歡有這些愛邦偶像的神.
這個忌邪,有些就譯為疾道.
說到他不可以有其他的偶像神明.
因為這個說的是我們和上帝的關係.
一個很親密的關係.
我們的主,我們是他的子民.
就好像一個婚姻的關係.
在這個婚姻裡面,丈夫和太太不可以有第三者.
我記得有位太太跟我說.
當她知道她的老公在外面有一個女人的時候.
她整個人的血一直飆上去頭上.
她很害怕自己會中風.
她很想衝入廚房,拿出刀來斬死對方.
這個難過,或者這個忌邪,疾道這個字.
就是說人的情緒.
在婚姻裡面,在信仰裡面.
不可以有第三者.
這個忌邪,說到很強烈的情緒.
所以,保羅給我們的榜樣是.
他很緊張一些未信的靈魂.
他很緊張他們是否得到這個救恩.
所以,丁姐妹,我們要反省一下.
我們對一些未信主的家人,親戚,朋友.
我們有沒有焦急到他們的靈魂有沒有得救呢?.
還是已經習慣了,他們信不信與我無關呢?.
保羅在這個雅典裡面.
他看到的,不是城市的繁華文藝氣息很豐富濃厚.
他看到的,是周街人的靈魂.
他們走錯了位.
所以求主幫助我們每一位.

$^{121}$我們真的想到我們的親戚,朋友.
他們的宗谷,如果沒有耶穌的時候.
他們的宗谷就會走向永死的時候.
我們由今天開始,就要緊張他們.
多為他們去禱告.
第二方面,保羅去到雅典.
他就好像去到其他地方一樣.
他先去到會堂裡面.
與猶太人自己的同鄉.
或者一些很傾慕耶和華的外邦人.
傾談信仰.
然後他每一天就在市場裡面.
與周圍的人有些辯論.
當時的事情就好像一些市集.
他們在那裡不單止有些買賣.
人們可以在市集裡面演講.
他們可以找到一個位置.
你就可以在那裡發表你自己的理論.
當你講一下,有很多粉絲認同你的.
於是人們就會推薦你去阿勒巴古.
在一個更大的廣場裡面演講.
據說,以前很出名的希臘哲學家蘇格拉底.
他最初都在市場裡面講一下.
於是就被人帶到阿勒巴古.
然後就出了名.
其實今天我們都是這樣.
有些人,你聽到有些歌星.
他們在公園裡面.
或者現在在天橋裡面.
有時我上班經過一條灣仔天橋去到會展的時候.
你會聽到有些人在彈結他.
那些歌星早上黃昏.
在那裡唱歌.
唱著唱著就參加電視台.
以前那些就是聲寶之夜.
現在就是中年好聲音.
然後唱著唱著就唱得很出名.
有一天就成為了梅艷芳.
現在的中年好聲音.
將來不知道誰會很出名.

$^{161}$原來雅典當時的文化就是這樣.
所以保羅覺得他要跟別人講話的時候.
他就在市集裡面這樣跟別人講話.
因為雅典裡面.
每個人都覺得他就是這個學者.
我有我的思想.
於是就每個人在這裡有人講.
那裡有人講.
他每次講.
那些人就無暇聽.
不是真的沒有空聽別人講話.
他們就是聽聽說說.
聽著聽著講著.
無無聊聊.
就去講著講著.
你這樣講啊.
我又有這麼多這樣的看法.
我有一位神學院的師姐.
她說她退休之後.
她就會去做管理員.
因為她發覺到.
在她住的大廈最熱鬧的地方是什麼呢?.
就是管理處.
我也發覺到.
下午管理處很熱鬧.
一方面有冷氣.
而另一方面.
退休的人士就在那裡聊天打牙交.
所以這個媳婦好不好.
然後15B的兒子有沒有工作.
然後誰又轉工.
哪裡的事事非非.
那個管理員全部都知道.
然後大家就去講一下.
又發表一下新聞的看法.
甚至有一次我聽到.
應該是很久之前.
有一個人進了銀行.
然後就下了閘.
警方就在外面猛說.

$^{201}$那些人快點出來.
在那裡說裡面有多少個賊.
人們都說是一個.
但結果出來是四個.
就是這樣.
當時那些人在那裡說.
在這個雅典裡面.
很多很多的神明.
很多的偶像.
有些人說.
這個雅典裡面的神像比人更加多.
因為當時希臘管理的地方都很多.
不過加起來.
他們所有的神明偶像.
都不及一個雅典城那麼多.
他們雅典人就覺得.
不怕萬一.
最怕一萬.
所以他們很怕拜少一個神的時候.
就怎麼辦呢.
所以他們無所不拜.
無神不拜.
最後就說我不知道那個是甚麼神.
總之就起一座壇給他.
叫做未識之神.
有拜錯.
沒有放過.
還有在當時的雅典裡面.
有兩派很出名的學派.
第一派就叫做伊比溝魯派.
是一個無神論者.
他們覺得在這個世界上.
都不知道誰是神真誰是假.
總之在這個世界上.
最重要是活得開心.
因為人生苦短.
都不知道有多久的命.
最重要是開心.
今朝聚.
明日壽來.

$^{241}$明日當.
所以他們一直主張開心.
這種哲學一直發展的時候.
從享樂主義.
貪玩貪食貪購.
然後去到放縱情慾的境況.
亂了大陸.
另外一派又和他們很極端的.
叫做斯多亞派.
因為世界上很多東西都是神明.
所以見甚麼就拜甚麼.
見到一棵樹就拜它.
於是大埔就出了一棵許願樹.
這塊石頭一塊高一塊矮.
好像一個男孩和一個女孩在拍拖.
又拜它.
於是灣仔又出了一個姻緣石.
他們凡神都要拜.
但又覺得全部都是神明的時候.
我們身體是一個臭皮囊.
我們要很謹慎對所有的神明.
所以他們覺得身體受苦是不要緊的.
於是他們覺得不吃不睡不洗澡都無所謂.
總之整個身體越辛苦越吃力的時候.
我們裡面的靈性就會升出來.
人生就會有意一點.
總之越苦人生就越有意義.
神明就會越喜歡.
這個主義在哲學一直發展的時候.
就成為禁慾主義.
有些修道士去到沙漠.
就整條柱子在那裡.
整個月不用洗澡不用剪頭髮.
這兩派的哲學一國走極尊.
時時都在罵.
在史集裡面也有人跟他們罵.
其實要稱讚一下雅典人.
因為他們很有宗教觀念.
天地間有神明.
男的神有女的神也有.

$^{281}$戰爭的神有和平的神也有.
智慧的神也有藝術的神也有.
農業的神也有.
連未識的神也要拜.
這個保羅為什麼會這麼難過呢?.
就是覺得這麼多的神他們都去拜.
為什麼沒有人知道天地間的創造者.
這個勝過死亡又復活的耶穌.
他很難過.
每樣都拜,卻沒有拜主耶穌呢?.
在過去一個星期的培靈會.
周曉輝牧師在講到.
我記得好像是最後一堂.
他就講到一個屬靈的方程式.
他說這個good減去god的時候就等於力.
這個是什麼意思呢?.
我不是講周曉輝那套.
我將他配合在雅典裡面.
他們這個多神派.
天地有很多的神.
他們有這個觀念.
天地間一定有一個神在.
這個是good的.
但當中沒有了真神.
沒有了主耶穌基督.
所以雅典人拜這麼多神的時候.
都是等於零,沒有用的.
另外一方面.
第二方面他說.
加上一個多一個零就等於good.
他的意思是.
只要我們生命裡面有主耶穌基督的時候.
其他東西有沒有都不重要.
這個和奧古斯丁的一個思想是同樣的.
奧古斯丁說我們在生命當中有很多的需要.
我們需要有衣食住.
我們很喜歡穿漂亮的衣服.
我們需要有東西吃.
我們需要有房子住.
我們也需要有金錢.

$^{321}$有被人尊重,被人肯定.
所有生命裡面有很多個洞洞.
好像蜂巢一樣.
需要這些東西填滿.
但其中在這個蜂巢的洞洞裡面.
有一個洞是屬於神的.
但你要在這個洞裡面填上這個神.
是一個真神的時候.
你的整個生命就會變得很有意義.
當你這個洞填上去是耶穌基督的時候.
就算你沒有其他.
沒有很漂亮的衣服.
沒有很多的衣服穿.
也沒有好東西吃.
人人去瘋狂旅行.
你都沒有的時候.
有時候被人冤屈你.
有時候沒有人尊重你.
不要緊,最重要有耶穌.
就像主耶穌基督有一個真神.
填滿我們生命裡面的時候.
其他的東西有沒有都不是最重要.
我們都可以有一個有意義的人生.
這個就是保羅先去了解雅典的情況.
他們有很多的神明.
很多的哲學.
唯一是沒有了耶穌基督.
所以當保羅去到雅典的時候.
他就在第23節.
他就到處走走看一看.
然後仔細地觀察他們所敬拜的.
他發現了他們有一座壇.
就是獻給未識之神.
保羅就說.
既然有一個不識之神.
我現在就向你們宣告.
我要告訴你們.
這個未識之神是誰.
這個未識之神就是一個最偉大.
最偉大,最厲害的神明.

$^{361}$是一位萬物的創造者.
第24到25節.
他是創造宇宙和其中萬物的上帝.
他既是天地的主.
就不住在人手所做的電雨裡.
也不需要人手去服侍.
好像缺少了什麼.
反而他將自己的生命,氣息,萬物都賜給人.
他說你們不要想著拜那麼多的神明.
要想著給神明什麼.
好像神明缺少了什麼.
他們真的缺少了什麼.
因為是人手做出來的.
但是這一位你們沒有拜的未識之神.
他是創造世界的.
他自己什麼都有.
相反,他可以將生命,氣息,萬物都賜給世人.
我想很久之前.
不知道有沒有在紅欣堂說過.
在禽流那一年.
所有雞都被殺了.
我在某一次坐小巴的時候.
我聽到兩位主婦在聊天.
其中一位主婦說.
真是慘了,所有雞都被殺了.
以後拜神都不知道用什麼來拜好.
很害怕,她已經吃慣了雞.
吃了千千萬萬年.
拿鴨子她喜歡吃嗎?.
拿鵝又行不行呢?.
她說都不知道以後怎麼拜神了.
另一位主婦說.
是啊,都不知道買什麼來拜好.
我就不理那麼多了.
我拿利是封放了五百元進去.
然後放在神桌面前.
我就說你自己喜歡買什麼就拿去買吧.
你喜歡吃什麼就拿去買吧.
拜完神之後就收回利是袋放進袋子裡.
這個雅典人未識的神明.

$^{401}$就是主耶穌基督.
他就是萬物的源頭.
他不需要人類給他吃雞.
我告訴你,我們的神都不需要吃羊.
不需要吃牛,不過只是看我們的心意.
他反而將萬物賞賜給世上所有的人.
另外,寶萊又說.
我們的生活,我們的行動,我們的存在.
都在於他.
他就是我們這位神.
他創造了我們,不是扔開我們就不理.
而是他一直維護著我們.
以至我們的生活,我們的動作,我們的存留.
他都掌握著我們.
在耶穌基督臨上十字架之前.
他都跟一群門徒說.
你們若常在我裡面.
我的話亦都常在你們裡面.
他就吩咐那群門徒.
一定要跟他保持一個緊密的關係.
一定要讀他的說話.
凡你們想要的祈求就給你們承傳.
意思就是門徒一定要跟他保持一個緊密的關係.
時時思想他,時時讀他的話語.
我們向他祈求,我們沒有力氣的時候.
沒有智慧的時候向他求.
他就會給我們,有什麼缺乏的時候.
他會幫我們.
最後一個,他就是那位公義的審判主.
在第三十一節說.
因為他已經定了日期.
他要藉著他所設立的人.
按公義來審判天下.
我們在地上高.
我們生活的時候.
有沒有感覺到有不公義的事情發生.
有沒有感到很不服氣.
有沒有搞錯.
我們難不成都會有這樣的經歷.
為什麼上帝沒有審判.

$^{441}$不過我們要去堅信.
惡有惡報,善有善報.
若然未報,時辰未到.
但更加堅信一件事.
這個報應是由上帝來做的.
上帝是一位公義的審判主.
他一定會為我們主持公義.
所以有什麼冤屈,有什麼事情會覺得不對.
交給上帝,他會審判.
因為雅典人是五色最偉大最聰明的耶穌基督.
也維護著我們的生命.
眷顧我們,拯救我們的上帝.
我們讀幾節詩篇.
是上帝說的.
全世界都是他創造的,屬於他的.
所以他要我們怎樣呢?.
我們一起讀,一二三.
林中的白獸是我的.
千山的牲畜也是我的.
山中的飛鳥我都知道.
田野的走獸也都屬我.
我若是飢餓不用告訴你.
因為世界和其中所充滿的都是我的.
我喜吃公牛的肉,我喜喝公山羊的血.
又要向至高者還你的怨.
並要在患難之日求告我.
我必答救你,你也要榮耀我.
我們同心禱告.
因為你讓我們先嘗到救恩的滋味.
過往我們是迷惘的.
是沒有方向的,是沒有盼望的.
讓我們因著你的拯救.
我們能夠知道我們前面的路程是怎樣.
也都叫我們在今生的裡面.
我們有力量,有方向.
主祝福我們.
讓我們嘗到救恩滋味的時候.
我們都緊張到身為未信的親戚朋友.
他們的靈魂未得救.
他們可憐,他們失望.

$^{481}$求主聖靈時時激動我們的心.
以致我們迫切地為他們去禱告.
將你都交到主你的根前.
願主來賜恩.
多謝你隨天禱告.
是奉主基督的聖名.
成名.
\newpage



\section{以弗所書 4:21-32}
\label{sec:cIHnc4_JX1s}
\textbf{2023.08.20 《脫去舊人‧ 穿上新人》以弗所書 4\_21-32 (和修版) 張美薇牧師}
\newline
\newline
連結: \href{https://youtube.com/watch?v=cIHnc4-JX1s}{\texttt{https://youtube.com/watch?v=cIHnc4-JX1s}} ~~~~ 語音日期: 2023-08-22
\newline
\newline
\hyperref[sec:vjbxBlVTElQ]{\small{< < < PREV SERMON < < <}}
~
\hyperref[sec:index_chronic]{\small{[返順時目]}}
~
\hyperref[sec:index_scriptual]{\small{[返順卷目]}}
~
\hyperref[sec:TnVaeXBM1r0]{\small{> > > NEXT SERMON > > >}}
\newline
\newline
以弗所書 4:21-32
\newline
\begin{longtable}{cl}
\hline
\hline
章節 & 經文 (和合本修訂版)\\
\hline
4:21 & \begin{tabularx}{0.7\textwidth}{X} 如果你們聽過他的道，領了他的教，因為真理就在耶穌裡， \end{tabularx} \\ \\ \relax
4:22 & \begin{tabularx}{0.7\textwidth}{X} 你們要脫去從前的行為，脫去舊我；這舊我是因私慾的迷惑而漸漸敗壞的。 \end{tabularx} \\ \\ \relax
4:23 & \begin{tabularx}{0.7\textwidth}{X} 你們要把自己的心志更新， \end{tabularx} \\ \\ \relax
4:24 & \begin{tabularx}{0.7\textwidth}{X} 並且穿上新我；這新我是照著神的形像造的，有從真理來的公義和聖潔。 \end{tabularx} \\ \\ \relax
4:25 & \begin{tabularx}{0.7\textwidth}{X} 所以，你們要棄絕謊言，每個人要與鄰舍說誠實話，因為我們是互為肢體。 \end{tabularx} \\ \\ \relax
4:26 & \begin{tabularx}{0.7\textwidth}{X} 即使生氣也不要犯罪；不可含怒到日落， \end{tabularx} \\ \\ \relax
4:27 & \begin{tabularx}{0.7\textwidth}{X} 不可給魔鬼留地步。 \end{tabularx} \\ \\ \relax
4:28 & \begin{tabularx}{0.7\textwidth}{X} 偷竊的，不要再偷；總要勤勞，親手做正當的事，這樣才可以把自己有的，分給有缺乏的人。 \end{tabularx} \\ \\ \relax
4:29 & \begin{tabularx}{0.7\textwidth}{X} 一句壞話也不可出口，只要隨著需要說造就人的好話，讓聽見的人得益處。 \end{tabularx} \\ \\ \relax
4:30 & \begin{tabularx}{0.7\textwidth}{X} 不要使神的聖靈擔憂，你們原是受了他的印記，等候得救贖的日子來到。 \end{tabularx} \\ \\ \relax
4:31 & \begin{tabularx}{0.7\textwidth}{X} 一切苦毒、憤怒、惱恨、嚷鬧、毀謗，和一切的惡毒都要從你們中間除掉。 \end{tabularx} \\ \\ \relax
4:32 & \begin{tabularx}{0.7\textwidth}{X} 要仁慈相待，存憐憫的心，彼此饒恕，正如神在基督裡饒恕了你們一樣。 \end{tabularx} \\ \\
[1ex]
\hline
\hline
\end{longtable}
$^{1}$大英姐妹平安.
平安.
很感恩今天又是八月中了.
竟然有機會在大家當中.
和大家分享神的話.
我們開始的時候.
我們一起低頭祈禱.
天父上帝你使用孩子.
幫助孩子今天所說的.
讓孩子能夠將你的說話.
說得清楚明白.
也讓會眾能夠了解.
也能夠進行.
我們今天的禱告.
是奉耶穌基督的名求.
阿們.
最近我們大家都說.
是不是要復償.
去復償的時候.
我們發覺原來開了關.
所以很多人都可以進去.
現在好像香港.
平時的人就少了.
因為那些人可能去國內.
花錢或者會出外旅遊.
或者會去探訪.
去不同的事情.
我今年都有這個機會.
六月底的時候就去了紐約.
因為我弟弟嫁女兒.
過去高興一下.
因為好像很久都沒有機會離開香港.
然後我就會帶隊.
去了南美蘇利南.
就帶了我們一行八人.
我和周牧師.
帶了六個神學生.
去蘇利南幫他們教會.
做一些海外實習的工作.
回來之後還不夠.

$^{41}$我八月初回來.
還去了黃山幾天.
因為覺得我還可以走.
最近都六十多歲了.
再不去的話.
遲早就走不了.
不如趁著走到就去.
所以其實發覺.
可能大家都是.
我發覺在網上看到.
原來Rita好像最近在英國.
我看到她的Facebook.
可能還有其他弟兄姊妹.
在不同的地方.
所以我們都是四處去.
我最近也看新聞.
因為在美國回來.
發覺原來美國那邊的.
治安或者經濟.
都不是很好.
最近看新聞.
看到有一段.
這段其實不是最近的.
不過最近看到的新聞.
最近在洛杉磯也有一件事.
原來在美國的加州.
在大概2021年.
這一段是比較久之前.
感恩節之前.
發生了一件事.
在三環市的知名的精品.
百貨公司Nordstrom.
裡面突然有八十個劫匪闖進去.
很厲害 八十個.
大家帶著雪帽.
雪帽只剩下兩隻眼和嘴.
拿著棒球棍.
然後去洗劫那間百貨公司.
那間百貨公司類似香港的.
會安 先斯那一類的百貨公司.

$^{81}$然後他們除之後.
他們就把贓物打包.
然後就拿走.
幾分鐘而已.
幾分鐘就離開那間百貨公司.
然後他們就一起走.
一起散水.
所以這個叫快閃搶劫.
但是有八十個人.
以致他們有25輛車來接他們.
你都不明白.
我也不知道他們怎樣去聯絡.
在網上聯絡.
沒有想過是這樣的.
原來不止一單.
在三環市.
其實在洛杉磯也有.
在不同的地方都有.
是加州的.
警方接報之後.
有幾十輛警車來到.
但是只抓到幾個人.
還有很出爭議的中心.
原來美國.
原來加州.
通過了一條法案.
叫SB47.
就是說如果他們去搶劫.
財物損失.
每個人搶到的東西.
是950美元以下.
就當是一個很簡單的罪惡.
不是一個大罪.
以致刑期最多只有一年.
所以有人說.
因為這樣的緣故.
所以引起他們.
可以這樣一起.
一起去搶一些名店.
去搶那些東西.

$^{121}$我們看到原來.
當一些東西的改變.
一些法例的改變.
一些法則的改變.
人的做法也會改變.
以前我沒有聽過.
這麼多人一起去搶劫.
現在發現原來可以這麼多人一起去搶劫.
其實是害怕的.
他們快要關店.
人們在走公司.
突然一堆人湧進來.
你害怕嗎.
當然害怕.
所以發現原來是很不妥.
為什麼我會提這個例子.
因為這次我提到的時候.
我問曾姑娘.
有什麼想我講.
她說講合一.
她說洪恩嘉可能很需要講合一.
我就想起了.
其實我大概一年前.
就來跟大家講過.
已忽所書第四章第一到十六節.
第四章第一節會說什麼呢.
我為主被囚的勸你們.
既然行事為人.
就要與你們所夢的呼召相呈.
我那次講了恩賜.
今天想講下半節.
下半節是講行為生活上的改變.
或者叫做倫理上的改變.
一會兒再跟大家講.
行為的改變.
為什麼會和合一有關呢.
跟大家先賣個關子.
大家會知道.
既然夢召成為合一的聖徒.
就要培養合一的美德.

$^{161}$既然夢召做一個聖潔的百姓.
我們就要有純潔的操守.
弟兄姊妹.
無論你信了主多久.
可能有些人是剛剛信主沒多久.
又或者你已經信了主很久很久.
保羅這一段其實對我們有所提醒.
他叫我們不要再像外邦人一樣.
我們會說外邦人和我們好像很遠.
但其實外邦人對今天來說.
我們用什麼來講呢.
就是非信徒.
我想大家從信主到現在.
是有了很多很多的改變.
但是大家有沒有在想.
我們其實應該是每天都在改變呢.
意思就是說.
保羅說我們不要再像外邦人.
在生活當中要有所謂典型的基督徒生活.
也有所謂非基督徒的生活.
大家也想到.
兩種的生活.
如果大家一起去比較一下.
其實是水火不相容的.
信徒的生活和非信徒的生活.
初信主的信徒.
從前是非信徒.
過的是什麼呢.
是非信徒的生活.
現在成了基督徒.
我們知道成了基督徒之後.
我們的身份馬上就改變.
但是你的行為會不會馬上就改變呢.
也未必的.
因為我以前是這樣做.
我現在也是這樣做.
特別當你是長大了才信主.
你小時候信主.
你不斷聽別人的話.
你跟著去改.

$^{201}$所以為什麼我們每天都要靈修.
我們要讀經.
我們要上主一學.
我們要聽道.
因為我們的知識上是改變了.
我們的思想上也改變了.
我們的行為也會改變.
所以為什麼會說靈修重要.
讀經重要.
聽道重要.
行道更加重要.
因為耶穌說什麼.
你們要去洗萬民作的門徒.
奉父子聖靈的名.
給他們施洗.
凡我所吩咐你們的都教訓他們.
不是吧.
不是都教訓他們就停了.
而是教訓他們什麼.
遵守.
所以我們是要不斷的去改變.
現在成為了基督徒.
不會馬上在生活上完全改變.
所以我們在意識上要有改變.
所以配合了我們改變的身份.
然後我們的行為上也要改變.
然後我們才可以過著基督徒的生命生活.
成為信徒的新身份.
就是神的新社會裡的成員.
所以我們會說我們是雄恩家園.
我們是雄恩家的一個成員.
雄恩家園當然和非雄恩家.
或者非基督徒是不一樣的.
所以我們是要有新的標準.
來做我們的參考.
就好像我們把舊的衣服脫下.
換上一件新的衣服一樣.
大家有沒有看過這張單章.
這張單章其實很典型很出名的.
是宣導出版社.

$^{241}$就是宣導書局出.
很古老很古老.
已經有接近200年的單章.
我們會看到單章的右邊的衣服是怎樣的.
好像是千瘡百孔的.
那件就是我們的舊衣服.
然後左邊的就是白色的.
所以其實有些教會的洗禮也是.
洗禮會穿一件黑的袍去洗禮.
然後上來的時候.
就應該有人會換一件白袍給他.
就好像說我洗了禮之後.
我就全部更新.
又或者我們今天就會說.
我們要脫去我們的舊日.
脫去舊我.
所以當我們去看.
已忽所書第四章第十七節開始.
就會說我們的舊我是怎樣的.
我現在就和大家去看一段聖經.
不過我就不會說很長.
因為今天我們說的經文也長.
這個脫去舊我在已忽所書第四章第十七到第十九節怎樣說呢.
他說.
所以我這樣說.
且在主裡鄭重地說.
你們行事為人.
不要再像外邦人.
存虛妄的心而活.
他們心地分昧.
因自己無知.
心裡剛硬.
以神所賜的生命隔絕了.
既然他們已經麻木.
就放縱情慾.
貪婪地行.
種種污穢的事.
這裡說舊我的影響.
其實我們對我們的舊我.
不是那麼容易覺察的.

$^{281}$為什麼呢.
因為習慣了.
我們很多時候說婚姻的主題.
主題都會說.
兩夫婦結婚.
最初最輕發生的問題是什麼.
就是擠牙膏的事.
男方擠牙膏是隨便按的.
女方是從尾部擠起.
一直捲上去.
一直都沒事.
但結婚就不同了.
結婚後大家用同一支牙膏時.
就發覺不對.
不可以這樣.
但我們是不會覺得有事的.
所以當我們信了主之後.
我們仍然用舊的身份.
穿著舊的衣服.
是不會覺得有大問題的.
所以他就說.
心裡剛硬.
都是因為自己無知.
和神所賜的生命隔絕.
但接著再說.
第20至22節說.
但你們從基督學的不是這樣.
如果你們聽過他的道.
領了他的教.
因為真理就在耶穌裡面.
你們要脫去從前的行為.
脫去救我.
這救我是因私欲的迷惑.
而漸漸敗壞的.
所以當我們信了基督.
立刻有新的身份.
剛才也說過.
我們信了基督.
立刻有新的身份.
立刻的.

$^{321}$我們有新的標準.
所以為什麼我們讀聖經.
我們會有些人說.
那些叫神學.
即是神給我們的標準.
所以我們要脫去從前的行為.
脫去舊的衣服.
是很簡單的.
大家如果.
昨天.
今天早上很大雨.
我記得我出外.
周末熙在我們後花園.
在整樹.
下大雨.
大到全身濕了.
但他在整.
他繼續整.
因為我們後面有些竹.
很長.
要斬掉之類.
然後他說.
整完之後.
立刻脫光衣服.
因為衣服很髒.
又很濕.
又大汗.
又酸又縮又臭.
他知道的.
他立刻脫光.
但脫光就算了.
還要怎樣.
洗完澡要怎樣.
要穿新衣服.
所以當我們信了基督有新的身份.
有了新的標準.
有新的神學.
我們脫去從前的行為.
我們就要.
因為我們從前的行為.

$^{361}$漸漸敗壞.
繼續敗壞.
我們不是停在那裡.
如果我們不脫去.
就會繼續敗壞.
所以我們的生活方式.
和現在世界所習慣的.
應該有截然不同的分別.
舉例來說.
現在開放文化.
很多時候原要肉體.
穿衣服是怎樣.
布料很少.
無論是衣服或褲子都是這樣.
同時很流行男教.
亦漠視精粗.
禮義廉恥等等.
去迎合人的慾念.
或者去追求名牌的東西.
去滿足虛榮感.
我們一定要抗拒.
我們要時時警覺.
所以我們知道.
我們要脫去舊人.
我不跟大家說太多舊人的事.
因為每個人來的地方都不一樣.
所以你的舊我.
每個人都不一樣.
我們要對付.
我們要脫去.
但是新我就不一樣了.
因為新我讓我們去看.
我們要去穿上新人.
在第20至21節怎樣說呢.
但你們從基督學的不是這樣.
不是你們的舊我.
如果你們聽過他的道.
領了他的教.
因為真理就在耶穌裡面.
所以我們穿上新我.

$^{401}$新我的基礎是甚麼呢.
第一我們要去學.
我們在哪裡學呢.
我們在基督那裡學.
基督的本身就是基督徒教材的素材.
我們學誰呢.
我們學基督.
所以我們叫基督徒.
我都和大家說過.
基督徒是甚麼意思.
就是基督的門徒.
即是小基督.
即是我們跟隨基督的.
基督這樣做.
我們就這樣做.
基督這樣走.
我們就這樣走.
所以報道者去傳基督.
聽眾就學基督.
然後學習關於基督的事.
關於基督的傳統.
因為基督到現在已經有二千多年了.
所以我們是要學這些的.
基督就是那位成了肉身的道.
獨一無二的神.
他成為人.
曾經受死復活.
現在去到作王.
所以我們不單止是學.
我們也應該去傳講基督的主權.
基督是我們生命的主.
傳講他所設立的公義的國度.
或者是管治權.
我們也要呼召信徒.
依從神的新標準和價值觀.
這一切都是跟未信主之前的外邦.
即是非信徒的標準是不一樣的.
所以我們是要去學.
學誰呢?.
學基督.

$^{441}$第二我們就要去聽.
去聽什麼呢?.
他說聽過他的道.
一樣是聽基督的東西.
是耶穌講的道.
今天就是在聖經裡靠著使徒記錄下來.
傳給我們.
我們今天繼續去聽.
我們不單是學習他的內容.
我們更加是要學習去聽耶穌是怎麼講.
所以我們不要輕看這個學習.
這個是誰講的?.
是耶穌講的.
是耶穌的道.
所以說基督就是在教人去認識基督.
所以我們要聽基督所講的.
我們不單是去學去聽.
我們還要受教.
領了他的教.
意思就是說.
基督除了是內容.
是教師.
他也是課室.
就是班房.
是學習的環境和氛圍.
所以你們在哪裡學的?.
在教會裡面學.
在什麼地方實行?.
在教會裡面開始實行.
當然我們不單是在教會實行.
我們在我們平時的生活上實行.
但教會是我們實行一個很開先的範圍.
我們可以說這樣的教導是屬基督的.
因為真理就在基督裡面.
所以我們可以看.
第23到第24節說.
你們要把自己的心智更新.
更新的意思是什麼?.
就是不斷的更新.
我們不單是信著那一刻.

$^{481}$有了一個新的身份.
我們要繼續去學習.
繼續的心智更新.
並且穿上的身我.
這個身我是照著神的形象做的.
有從真理來的公義和聖潔.
剛才說過.
我們要脫去舊人.
脫去舊的你.
我們每一個人都不一樣.
你有你的困難.
我有我的難處.
我有我過去不好的事.
你有你過去不對的事.
所以我們要脫去.
但我們脫去之後.
我們就要穿上新的東西.
我們要按著神的形象.
讓我們看到.
我們要穿上這件清潔的衣服.
在我們的身上.
所以在這個原因.
我們必須接受耶穌的命令.
去棄絕舊我.
以及已經被唾棄的生活方式.
新我已經穿上了.
但我們仍然要不斷的更新我們的心智.
一天一天.
或每時每刻.
經歷內在態度的更新.
所以為什麼我們說.
我們信主.
我們仍然要更新.
當然要.
你有沒有想過.
我多久沒有改變過呢?.
或者我有多久沒有想過.
我這樣做是不對的呢?.
或者我以前.
我什麼時候會有新的方向.

$^{521}$按著聖經而走的方向呢?.
其實我們信主那一刻.
是有新創造的經歷的.
那是我們的經歷.
我們要將舊我完全拋棄.
但當我們拋棄了舊我的時候.
我們仍然有一顆不斷更新的心智.
我們要滿心感恩的.
去穿上我們的新人.
很多時候.
我們信主的時候.
並未將舊我完全拋棄得一乾二淨.
所以信徒一定要留心的.
要警醒去看.
我們新的東西.
有什麼呢?.
大家想想.
如果我們去婚禮.
會穿什麼衣服呢?.
當然是穿.
結婚的當然是穿結婚的禮服.
你去參加婚禮.
都會穿一些很漂亮的衣服.
或者出席喪禮.
會穿什麼衣服呢?.
你一定不會穿婚禮的衣服去喪禮.
通常現在香港人都是習慣.
穿黑色,灰色,藍色等等的衣服.
又或者.
當你在合適的場合.
就會穿合適的衣服.
看到嗎?.
有白頭假髮.
又有一條白領.
黑袍的是什麼呢?.
是律師.
律師當你在法庭上.
就會穿律師的衣服.
又有一些.
如果我們看他下面.

$^{561}$紅色衣服的.
他有一個計時錶.
這個應該是計跑員.
旁邊那個就是醫務人員.
再旁邊那個可能是記者.
之類的.
再旁邊那個可能是空姐.
就算牧師.
都有牧師的衣服.
還有.
軍隊.
海軍,陸軍,空軍.
都有不同的衣服.
還有.
不要.
這叫什麼?.
不要理他.
犯了事.
囚犯.
都有囚醫.
當然現在囚醫不是這樣.
這是古老時候.
大家看這個就知道.
但是當我們.
改變了身份.
我們脫離了身份之後.
我們就會換不同的衣服.
假如你是空軍的.
當你.
做滿了你的.
當值.
當你.
出外的時候.
你就會脫去.
你那個空軍的制服.
穿回平民百姓的衣服.
我聽到有人說.
當你穿回平民百姓的衣服.
當你走出去.
你的辦公室的時候.

$^{601}$人不再跟你敬禮了.
你穿著制服的時候.
人們就會敬禮.
當你離開的時候.
就未必.
當然我們會知道.
那個囚犯.
囚犯當然是.
當離開了囚牢.
還回自由的時候.
就會穿新的衣服.
其實是一樣的.
我們都是一樣.
我們都是要去學習.
穿新衣服.
所以今天就和大家介紹.
五件的新衣服.
其實很多件的.
今天就在這段經文裡.
和大家介紹.
五件新衣服.
大家去留意聽.
叫做五個生活實例.
五件舊衣服.
五件新衣服.
我們先說舊衣服.
所以你們要棄絕謊言.
每個人都要和論社說誠實話.
因為我們是互為之體.
第一件舊衣服是什麼.
謊言.
說謊.
我們知道的.
是就應該說什麼.
是.
不是就說什麼.
不是.
我們應該以誠實相待.
其實嚴格來說.
在希臘文裡.

$^{641}$這個詞不是指一般的謊言.
一般的虛假.
其實是指一個很特別的謊言.
這個謊言是指拜偶像的謊言.
當人去拜偶像.
他以為拜了偶像之後.
可以得到他要有的東西的時候.
其實那是一個很大的謊言.
於是保羅在用的時候.
信徒已經離絕了最大的謊言的時候.
進一步的.
連一些較小的謊言.
都要丟棄.
只說真理.
所以第一件舊衣服.
就是不要說謊言.
沒有積極追求真.
是沒有益處的.
是跟從耶穌的人在群體裡.
應該被譽為誠實.
可靠.
賢義.
有信的人.
理由不是單單是因為.
聖經要我們去愛我們的鄰舍.
因為這裡說你們要棄絕謊言.
每個人都要跟鄰舍說誠實話.
因為我們是互為之體.
我們更加.
在這裡給了我們神學上的解釋.
我們更加是因為我們在教會裡.
我們這麼親密的關係.
我們是真誠的.
我們是彼此的信任的.
信任的基礎.
就是我們的真誠.
所以當我們彼此一起生活的時候.
我們各人就要互為鄰舍.
要說誠實的話.
我們生活在一個身體之內.

$^{681}$真理需要我們更加堅固.
所以我們要說真話.
什麼叫誠實話呢.
當我們批評別人的時候.
我們就要考慮三個問題.
當我們去想要不要說.
第一.
你說的話是不是事實.
內容是不是對的.
第二.
你是否憑愛心說的.
動機是不是對的.
第三.
是否必須在這一刻說的.
時機是不是對的.
所以當我們要說誠實話的時候.
我們要想要問這三個問題.
另外一個說誠實話的.
就是我們要為他人保守秘密.
當別人信任你.
將心中的問題困擾.
向你傾訴.
尋求你給予意見.
或者向你分享心訴的時候.
你要誠實地為別人著想.
他只是告訴你.
他沒有告訴他.
如果告訴你.
你馬上就告訴他.
那對當事人會有什麼困難呢.
所以這裡就教導我們.
我們是要說誠實的話.
我們要小心謹慎.
我們的嘴巴.
所以我們除了一件舊衣服.
舊衣服是什麼呢.
是說謊話.
說謊言.
說不真實的事.
穿一件新衣服.

$^{721}$這裡除了讓我們看到.
我們行為上要改變.
還讓我們看到神學上是怎樣.
因為我們是互惠之體.
我們是一個群體.
是一個團契.
我們是彼此.
要以真誠相待.
所以這是第一件衣服.
大家是要除掉的.
除掉謊言的衣服.
這裡也說了.
第二件衣服是什麼呢.
即使生氣也不要犯罪.
不可含怒到日落.
不可給魔鬼留地步.
所以這裡我們看到.
原來是說我們要避免發怒.
我想大家都知道.
發怒好不好.
不是不好的.
大家不要這麼快說.
他說什麼.
即使生氣.
會不會生氣呢.
是會的.
信徒不是不能生氣的.
大家要記住.
不是不能生氣的.
因為我們看過保羅也有生氣.
保羅因為哥林多教會.
有很多不道德的生活表現而生氣.
耶穌也有生氣.
耶穌去到潔淨聖殿.
他是不是生氣.
當然是.
他在邊上打人.
還不是生氣.
是生氣的.
所以我們會發覺.

$^{761}$這些是叫做疑怒.
但是保羅提醒我們.
即使生氣.
不是不能生氣.
但是生氣的人是怎樣.
容易犯罪的.
即使生氣也不要犯罪.
由氣會生什麼.
生怨.
由怨就會生怒.
由怒含在心裡.
就會生恨.
恨就會生仇.
仇恨當中.
很容易生出惡毒等等.
所以生氣的時候.
要記住我們.
我們會說錯話.
會失去我們的常態.
我們可能有時生氣.
人們就會動手.
會失手.
更加會失去我們的見證.
所以不是說我們不能生氣.
即使生氣也不要犯罪.
所以這一節告訴你.
基督徒是會發怒的.
事實上.
如果我們否認怒氣.
等於否認神.
也傷害自己.
更加縱容惡行.
特別是今天的環境.
我們更加需要.
更加多的基督徒.
去表達我們的怒氣.
如果你覺得有些事情是不公義.
有些事情是不公平的時候.
我們是要去表達的.
因為人很容易和罪妥協.

$^{801}$沒有怒氣.
當我們面對猖獗的惡事的時候.
是不對的.
我們不可以壓低.
完全不要怒氣.
但是.
當我們遇到不對的事情的時候.
我們要憤怒.
我們不要容忍.
我們要發怒.
我們不要莫言而對.
但是.
我們很重要的知道.
真正的和睦.
並不是憤色太平.
不是什麼都吞下去.
什麼都沒有.
但是這裡有三個不.
要我們留意的.
第一個是什麼?.
生氣卻不要犯罪.
很多時候.
我們生氣是因為什麼?.
是因為我的自尊心受損.
因為我覺得我被人鄙視.
我覺得別人是出於惡意.
或者別人是離間我.
或者對我報復而來.
不要犯罪.
這是第一個不.
第二個不是什麼?.
也不可含怒到日落.
不要將怒氣變成含怒.
不要讓它繼續在你裡面.
要馬上制止它.
或者你要向對方尋求和解.
或者要向對方修好.
今日是今日了.
所以不要含怒到日落.
日落之前最好處理好你的怒氣.

$^{841}$這是金方玉律.
不單是弟兄姊妹之間.
我們需要這樣.
或者同人之間也需要這樣.
我們更加是和兩夫婦之間.
都要這樣.
不要含怒到日落.
這是很好的.
第三不可以什麼?.
不要讓魔鬼留地步.
因為魔鬼深知.
異路和不異的路.
是一線之差.
要知道人很難控制我們的怒氣.
所以他最喜歡自己移動.
利用怒氣沖沖的人.
去激發仇恨.
去破壞團體.
所以我們不要讓魔鬼留地步.
戒爾文說了一句話.
他說很快地壓住你的怒火.
免得怒火因持久而硬化.
早些解決.
不要留很久.
讓他慢慢積成一個很大的怒氣.
這是第二件要脫去的衣服.
就是穿上一個很快要處理好怒氣的衣服.
為什麼?.
因為我們不要犯罪.
我們不要拖延.
不要讓魔鬼留地步.
這是神學上的解釋.
第三.
偷竊的不要再偷.
總要勤勞親手做正當的事.
這樣才可以將自己有的分給有缺乏的人.
大家如果聽到偷竊的話.
大家如果可以看回摩西的十誡.
也有說.
第八誡摩西也說不可以偷竊.

$^{881}$我們會問為什麼人要偷呢?.
為什麼?.
因為人不肯做事.
不想做事.
但他又要賺取生活費用.
所以他就用不道德或者損人的途徑.
來得到金錢來解決生活上的所需.
有時候不只是偷.
範圍內也說.
可以說什麼?.
走私漏稅.
叫國家的庫房受到影響.
有些有錢人更加少付稅.
因為他們製造更多幫自己逃稅的.
或者避稅的方向.
我們要留心.
如果大家找到又要交稅的話.
是感恩的.
因為你不單止自己找到.
更加可以叫國家.
或者叫社會或者政府.
都有幫助.
偷可能是指僱主.
剋扣工人的工資.
剝削員工.
或者僱主的角度就是剝削.
工人的角度是什麼?.
偷懶.
馬虎.
做事做得不好.
保羅重申.
他除了叫人禁戒.
不要再偷.
把偷衣服脫掉.
不要偷竊.
把過去偷了別人的東西.
都脫掉.
但更重要的是.
他叫人勤勞.
親手做正當的事.

$^{921}$他要叫人開始工作.
親手做正經的事.
自食其力.
然後怎樣?.
神學上解釋了什麼?.
然後把自己有的.
分給有缺乏的人.
自己有的當然要分給別人.
偷的時候就把別人有的拿過來.
現在反過來是怎樣?.
我有的東西.
我可以分給.
我不是分給全世界的.
因為有些很有錢的人.
都不用你派.
好像那時候別人說.
香港如果派那幾千元的消費券.
其實李嘉誠用不用?.
其實李嘉誠是不用的.
但又很難去做.
那麼多東西是阻礙的.
但如果你要去做的話.
這樣做的話.
才可以把自己有的.
分給缺乏的人.
所以這個本來不單止是偷.
現在更加是偷.
因為賊是社會的寄生蟲.
但現在你能夠回饋社會.
不單止是養活自己的一家.
更加能夠分給那些缺少的人.
所以他鼓勵信徒之間.
是要彼此幫助.
但我們看到.
他鼓勵我們信徒之間彼此幫助.
不是說你很有錢.
很多很多到不知道多少個零.
才要去幫人的.
他說是由本身經濟環境欠佳的人.
藉著努力.

$^{961}$以便為自給自足.
然後更加有餘力去資助其他人.
我們雖然不是多.
我們雖然不是大富大貴.
但我們仍然可以有餘的.
給別人去成為一個有依靠的人.
我們願意每一個信徒.
工作上都好像是為主做的.
不是為人做的.
抱著這種心態.
即使你的工作不理想.
你的老闆好像很苛刻.
對你不好.
但你都要盡心去做.
好像服侍主一樣.
願我們每一個信徒.
每天工作的時候.
都記住這番說話.
我們總要勤勞.
親手做正當的事.
在工作和職業上.
努力做出貢獻.
如果有利有不遞的時候.
除了祈禱求主加力幫助之外.
亦要找新的方法去進修.
以達到自我改進.
或自我增值.
第三件事是甚麼.
除了我們偷竊的事.
偷竊不是真的偷.
而是我們如何去貪別人的便宜.
將不應該是你的東西.
變成了你的東西.
第四是甚麼呢.
第四是二十九節.
一句壞話也不可出口.
只要隨著需要說造就人的好話.
讓聽見的人得益處.
不要使神的聖靈擔憂.
你們願受了祂的印記.

$^{1001}$等候得救贖的日子來到.
言語是神給人奇妙的禮物.
是反映神形象的能力.
我們神是說話的神.
我們也像祂一樣.
我們能夠說話.
我們的言語.
讓我們與禽獸有所分別.
保羅叮嚀地說.
他說甚麼呢.
一句壞話也不可出口.
壞話.
壞話這個字.
其實在和學本.
是污穢的言語.
是指木頭或水果腐爛.
或是質量差不合用.
或是沒有價值的東西.
即質量差不合用.
用來形容言語.
可能是奸詐,刻薄,粗俗的言語.
我估計聽到這種言語的人.
心裡會感受到很難堪.
很受傷.
所以保羅說.
一句壞話也不可出口.
反而是怎樣呢.
要隨著需要說造就人的好話.
所以除了衣服.
衣服是甚麼呢.
可能你時時都會說的壞話.
即是氣絕污言穢語.
你除了這些污言穢語.
然後你便要穿一件.
說一些造就人的說話.
箴言說.
一句話說得合宜.
就如金蘋果在銀網裡.
在人的活動背後.
存有不一樣的東西.

$^{1041}$這裡說了神學上的意思.
不要使聖靈擔憂.
最後是怎樣呢.
我們不要在我們生命上的聖靈擔憂.
因為聖靈已經在我們信主的時候.
住在我們裡面.
如果你仍然在口裡.
不造就人的話.
在你裡面的聖靈便會為你擔憂.
這裡更加說.
最初聖靈入你裡.
好像給你一些印記.
然後到最後.
當主耶穌基督再來的時候.
那些印記.
就是得贖的日子來的時候.
聖靈充滿的信號.
就會叫聖靈喜悅.
所以這提醒我們.
要棄絕我們的污言.
第五件衣服.
一切苦毒,憤怒,怒恨,讓罵,誹謗.
一切惡毒都要從你們中間除掉.
要仁慈相待,全憐憫的心.
彼此饒恕.
正如神在基督裡面.
饒恕了你們一樣.
即是叫我們戒黑薄苦毒.
言語是一樣.
另外行為也是一樣.
所以這裡說了六項.
苦毒,憤怒,怒恨,讓罵,誹謗.
和一切惡毒.
都要在我們當中離開我們.
我們應該完全摒棄.
然後怎樣呢?.
我們就會看到.
其實這些.
為什麼會有苦毒,憤怒,怒恨,讓罵.
其實是怎樣呢?.

$^{1081}$其實是因為我們生氣.
前面說了避免發怒.
這個就出來了.
但相反我們要怎樣呢?.
要仁慈相待,全憐憫的心.
彼此饒恕.
正如神在基督裡面.
饒恕了我們一樣.
因此.
憐憫.
饒恕.
我們要學習彼此饒恕.
其實這五件的衣服.
五件舊的衣服.
五件新的衣服.
其實裡面全部都和人際關係有關.
都是影響到我們之間的關係.
所以我告訴你.
我們要聖潔.
但聖潔不是孤立於人.
好像頭上只戴一個光環.
這樣那樣.
那是叫做聖潔.
我們做好的事.
不能夠在真空裡面做.
我們要在真實的人世裡面做.
所以所有的素質.
都是和教會合一.
息息相關.
如果大家想教會合一.
我們一定要在我們當中.
棄絕謊言.
避免發怒.
不要偷竊.
棄絕污言穢語.
戒黑薄.
苦毒.
我們更要穿上我們新的衣服.
我們要以誠實相對.
我們要很快地處理好我們的怒氣.

$^{1121}$我們要勞力作功.
我們更要說造就人的說話.
我們要有恩慈憐憫的說話.
每一項禁戒.
都附以一項積極相應的命令.
所以我們一定要知道.
每一個例子都告訴你.
神學上是有意思的.
是為什麼會來的.
信徒的教義和倫理當中.
信仰和行為永遠都是密切的配合.
所以弟兄姊妹.
我們一信主.
我們的身份就已經改變.
但我們不單止身份要改變.
我們還要思想上改變.
所以我們需要讀一下聖經.
讓神的話語成為我們當中.
好像很難的字.
神學.
但這就是神學.
然後我們更需要.
在我們的行為上走出來.
這樣我們就可以一起.
合一成為一個.
我們可以說是雄恩佳.
一個家是一齊去做.
彼此的包容.
彼此的饒恕.
彼此的憐憫.
彼此的恩慈相待.
我講道經常都會用這個結束.
我想大家都知道.
和你有什麼關係.
希望我們起碼找到一件衣服.
你想穿的.
找到一件新衣服.
今天找到一件新衣服.
你要穿上.
但記住.

$^{1161}$你身上有沒有一件舊衣服.
要脫下呢?.
最好已經脫光了.
但我知道可能還未.
我們繼續努力.
願我們彼此一起在主裡面長進.
我們低頭祈禱.
讓我們可以從保羅的說話裡領受.
幫助會眾.
讓他們彼此互相幫助.
讓他們有了新的身份之外.
他們更加要在信仰上.
在神學上一天的更新.
更加重要的是.
他們能夠實行出來.
走在他們的幸事為人上.
以致他們真的能夠做到彼此相助.
一起的,合一的.
成為一個你愛的奪妻.
聽我們的祈禱.
奉耶穌基督的名求.
阿們.
\newpage



\section{馬太福音 28:19}
\label{sec:TnVaeXBM1r0}
\textbf{2023.08.27 《你們要去,使萬民作我的門徒》 馬太福音 28\_19 (和修版) 講員 : 高淑芬傳道}
\newline
\newline
連結: \href{https://youtube.com/watch?v=TnVaeXBM1r0}{\texttt{https://youtube.com/watch?v=TnVaeXBM1r0}} ~~~~ 語音日期: 2023-09-01
\newline
\newline
\hyperref[sec:cIHnc4_JX1s]{\small{< < < PREV SERMON < < <}}
~
\hyperref[sec:index_chronic]{\small{[返順時目]}}
~
\hyperref[sec:index_scriptual]{\small{[返順卷目]}}
~
\hyperref[sec:iQ2CTi_2bNs]{\small{> > > NEXT SERMON > > >}}
\newline
\newline
馬太福音 28:19
\newline
\begin{longtable}{cl}
\hline
\hline
章節 & 經文 (和合本修訂版)\\
\hline
28:19 & \begin{tabularx}{0.7\textwidth}{X} 所以，你們要去，使萬民作我的門徒，奉父、子、聖靈的名給他們施洗， \end{tabularx} \\ \\
[1ex]
\hline
\hline
\end{longtable}
$^{1}$大家好.
經文只有一句.
我要講的也只有一個字.
一會兒你就知道我要講哪個字了.
如果大家看到這一句.
大家也稍為知道這是甚麼甚麼甚麼.
三個字.
是甚麼.
大使命.
我們一講到大使命.
我們不只是停在這裡.
不過今天我覺得.
我們不只是停在這一句.
還要停在一個字裡.
我們要做好一點.
有一位姐妹.
由我微信主到信主.
到今天她也很認識我.
有一天.
沒多久她問.
如果生命再給你重頭一次的時候.
你還會不會去宣教?.
這樣問我.
不知道你猜我的答案是甚麼呢?.
我先不告訴你.
這個問題我覺得也不錯.
人生我們也可以問問自己.
如果生命再給你重頭一次的時候.
你還會去做甚麼?.
例如.
經過她這樣問的時候.
我曾經問過人.
如果生命再給你重頭一次的時候.
你會不會嫁給他?.
或者你會不會娶他?.
有些人轉頭.
剛好我問的那個人又轉頭.
我心裡要默默為他祈禱.
真的.
有時候有些事情你會轉頭.

$^{41}$有時候你會繼續.
我的答案就是這個.
就是這個.
你知不知道這個甚麼字?.
這麼厲害.
我也是看Google.
我很想知道這個甲骨文是怎樣的.
剛好這個甲骨文也很有意思.
又適合我今天要說的話.
所以我就編了出來.
如果問我.
我不懂的.
但為甚麼你這麼厲害?.
這個字是去的意思.
我的答案是去.
不過.
不有不過.
不過不是說不遇見那些.
而是我要做更多好一點的裝備.
再去.
再見過這些情況.
這些事件.
再培養好一點的生命.
再去.
這個去字.
甲骨文的字型.
上面那個大字好像一個人.
下面就好像一個口.
或者一個坑.
一個洞.
一個洞口.
我覺得很有意思.
這個去字原來.
表示一個人要離開自己範疇裡面.
而出去.
我覺得很符合大使命的去字.
因為你不可以再在你的安書區裡面.
你不可以只是停在.
我是這樣.
我是很內向的.

$^{81}$你叫我去傳福音是不可以的.
或者.
我很多事都很害怕.
繼續在害怕裡面.
不是.
這個去字有個意思.
是你要離開你本身.
你可能會覺得很舒服的地方.
或者你覺得很恐懼的地方.
你都要出去.
其實大使命.
這個去字.
在原文裡面.
是放在第一個字裡面.
是很強調一個去字.
所以我今天經常有聲音在裡面.
大使命你又說不完.
還要說一句.
一句你又不知道能不能說完.
你還要只說一個字.
我心裡面.
在控訴我自己.
說得這麼少.
因為平時你看到我的講題.
很多很多.
所以我就覺得.
這麼少的東西真的很難說.
這個去字.
原來在大使命裡面.
是真的很重要的.
再說.
這個大使命.
第一個字是去字.
是你們要去.
洗萬民做我的門徒.
記住.
這個大使命不是去.
只是去叫那些未信主的人.
去相信主.
是這樣的.

$^{121}$是讓他們成為門徒.
而且不是讓未信的成為門徒.
是讓所有的人.
萬民裡面.
包括你和我嗎?.
包括的.
意思是.
你不只是想著.
怎樣搞到我旁邊的成為主的門徒.
你自己呢?.
你自己有沒有好好地跟從神成為主的門徒呢?.
這個就是大使命.
不只是向外.
同樣都要向內.
如果你記得我曾經說過大誡命.
大誡命是甚麼?.
你要愛人.
都要愛人如己.
我們經常忽略了己.
即是愛人.
你都要愛自己.
所以大使命同樣有內那邊.
亦有外那邊.
正因為去字這麼重要.
我們再看看去字.
你要離開所在的地方.
這個地方不一定是.
我要離開香港去韓國去泰國去哪裡.
不是這樣.
是一些事件.
或者怎樣去到另外一處.
總之是有移動的.
就好像你不斷地旅行.
不斷地離開.
要離開你不可以固有地站在原地裡面.
甚至可以貼切一點.
自己這方面去到別人那邊.
要去這樣.
不可以只是這樣.
自己那邊去到別人那邊.

$^{161}$例如教會的祈禱會何時開始?.
何時會有教會的祈禱會?.
下星期日.
厲害了.
下星期日每一個月的第一個主日.
這裡很歡迎.
就算你是剛剛來我們教會.
即是來了很久.
有團契也好.
沒有團契也好.
都邀請你來.
甚至我在想.
很多時候我們的IT部.
他們差不多都是會不會找一次輪流.
派一個代表去祈禱會.
因為上帝呼召我們成為門徒.
而不是呼召你.
只是星期日來守主日.
你明不明白?.
我有守主日.
意思是說我每個星期都有去教會.
不單止是這樣.
當然我們每個主日去守主日的時候.
其實已經有很多的恩典.
不過很有趣.
當你說要去教會.
去的時候.
不知道今天大家來這裡去教會.
去弘恩堂聚會的時候.
之前有沒有發生一些事件.
或者一些意外.
令你覺得早知就不要出街.
有沒有這樣?.
今天沒有?.
真是感恩.
大家都很順利.
我今早去不到電腦裡面.
我就已經非常懶惰.
早知就不會答應珍姑娘說.
太緊張了.

$^{201}$我那篇文章一直在手機裡面製作.
怎樣去到電腦裡面.
東西會走位.
又要再弄.
有時候去是這樣的.
會帶著一些不好就好了.
有時候會有這樣的心情.
但問題是在過程裡面.
當你真的完成了的時候.
有這樣的心情.
或者你見到一些卡住你的東西.
解開的時候.
你又會覺得感恩.
幸好我有去.
幸好我有這樣.
同樣我都希望你.
未來過祈禱會.
下個主日.
幾點鐘?.
本來十點鐘.
我們太與團契都會有主日學.
我們那天都不主日學.
一起祈禱.
教會裡面都有一些報道探訪.
你去不去?.
去了的就繼續去.
多走一步.
未去的就考慮一下.
會否去一次.
真的要踏出.
去到街頭又熱.
你還說去報道.
去街我都不想去.
很多人近來天氣又熱.
又經常變化.
很多人自認宅女宅男.
不想出街.
如果雪櫃裡面還有東西可以維持到下一餐.
先不要出街.
弄到沒有東西.

$^{241}$清倉了.
才去.
原來去字說是很簡單.
有時做起來都不容易.
但這裡真的鼓勵大家.
去是一個大使命.
一個大前提.
又例如教會裡面.
為何列舉這些可以繼續列下去.
淨土分煉你會否去參與?.
因為我們要成為上帝的門徒.
我們要成為門徒就更加要勤讀聖經.
這裡包括去讀主一學.
或者我們一起.
聯同弟兄姊妹.
會否有一天回來教會.
一起查一下聖經.
一起不查聖經.
我記得有一個見證.
在中國大陸裡面.
有一個姐妹只有二年級的學歷.
意思是她的字不是懂得很多.
但當她信了主之後.
當然不會一信主.
有些神仙返兵.
你就全部本聖經都可以讀.
不是.
她都不是這樣.
她就唯有身為婆婆.
問孫子,兒子,媳婦.
每天問一些字慢慢記住.
就是這樣.
於是全本聖經的字都懂得讀.
因為國內很缺乏傳道人.
信了主的人.
有些懂得字的.
他們會在一起.
每一次讀一段.
都不要解釋.
譬如《馬太號音》.

$^{281}$今天我們讀第一章.
那位姐妹就記住.
盡量記住.
不記得就問.
到有一天她竟然可以幫忙參與講道.
因為國內的信徒是這樣.
他們真的沒有辦法.
全民皆兵.
就好像我們在家.
過去我們的年代.
不是那麼流行請家務助理.
家裡的家務.
父母,大女兒.
第一,第二,第三,第四.
都有些事情要分配.
你會看到家裡髒了.
當然小時候.
我弟弟,妹妹負責.
看到髒了都不理會.
小時候還未成熟.
但長大後.
看到地面髒了.
妹妹又讀書讀得快考試了.
你會怎樣?.
你也會看到髒了就掃.
這就是那一種心情.
在教會裡面.
你會看到一些空缺.
例如招呼新朋友.
例如教會有新人來.
去跟他打聲招呼.
會否只是招待.
或者只是今天招待.
看看有沒有跟他打招呼.
不是的.
你也可以去.
因為這是屬靈的家.
我們沒有分階級.
都是一起的.
彼此配搭的事奉.

$^{321}$其實這個大使命.
不是在個人裡呈現.
是一個團體.
是一個合一的見證.
是一個team work.
是一個對工.
就算我們有牧師,有傳道.
他不是一個人包辦教會所有的事.
都是要跟教會的弟兄姊妹.
彼此事奉.
所以這就是過去的家庭.
但我們現在的教會.
很多時候變成了新的家庭.
就是現今的家庭.
看到地上一樣髒.
這是什麼責任?.
Mary Ann的責任.
就是監務助理的責任.
請她來是她的工作.
吃完飯,洗完不洗碗.
誰收拾碗?誰什麼?.
又是她的.
因為這個家庭請了監務助理去做這些.
但一個監務助理.
如果家庭有七個人.
星期六,日還要招呼親朋戚友來很多.
我們聽過的監務助理是這樣的.
中秋節,節日,農曆新年就大排筵席.
你猜一個監務助理能不能做這麼多事情?.
讓你一個人去做這麼多事情也不行.
有時候我們會站在別人的立場去看.
不是說我付了錢請了監務助理.
她就做完這些.
不是的,也要她幫忙.
所以我看到一些家庭.
雖然請了監務助理.
但那些小朋友還是要繼續看著.
繼續教.
不只是全部給監務助理去做.
今天我們屬靈的家庭.

$^{361}$不要想著教會也請了牧師.
請了傳道.
你們做吧,珍姑娘.
全部你來處理.
家裡沒有人掃地.
清潔的人不在.
你來處理.
是不是這樣?.
不是這樣的.
所以教會裡面有一些空缺.
譬如我們經常說要請什麼人?.
藍傳道.
所以在這裡的弟兄.
你也可以站出來幫幫忙.
做二十分之一的傳道好不好?.
做一百分之一的傳道好不好?.
去探訪的時候.
多去一步.
多去半步.
就是這樣.
不是說要請到人才可以去做那些事.
是現在也可以做.
我們知道在這裡.
早期有些弟兄.
他學習毛筆字.
寫大字.
也是一個平台.
就是用他的專長去多一步.
天氣這麼熱又要教.
又要想教材.
又要想怎樣教.
其實去都要有些本錢.
你都要去做.
但是都願意去籌備一切.
來做一個平台.
讓一些信的和未信的在一起.
可以透過這個平台.
讓一些不肯回教的人.
他喜歡書法.
認識上帝多一些.

$^{401}$所以其實是各盡其職.
問題是你肯不肯去多一步呢?.
教會的新人你會不會跟他們打招呼呢?.
有訓練你會不會去呢?.
去多一步.
今天有時候我們教會裡面.
我都未去過感受堂.
我現在更加不認識感受堂的弟兄姊妹.
有時候會見到一些新面孔.
我都不知道他們是新朋友.
還是感受堂的嘉賓弟兄姊妹來參與呢?.
不知道.
可能你和我一樣有這個疑問.
不知道.
但是不要緊.
我們已經很好.
我們的紅恩堂都會大家圓左崇拜.
雖然時間很緊急.
大家都會問候一下.
但是我們會不會去多一步.
每個主日找一個三面孔.
問候一下.
認識一下.
因為這是你的弟兄.
這是你的姊妹.
你都要認識一下.
求主幫助我們去多一步.
我們每個人都是不同的.
問題是我們願不願意為主去多一步呢?.
你有沒有認識一些朋友.
已經參加了旅行團.
去歐洲.
參加完報名的時候很開心.
因為他是一些小鳥價.
總之是早鳥優惠.
早點報名價錢會便宜一點.
當時報名很開心.
這個價錢可以報這個旅行團.
但是日子快到了.
快要出發的時候.

$^{441}$剛好那兩個星期.
他工作特別忙.
於是他很後悔.
竟然有後悔.
為什麼參加這個旅行團.
這個日子隔天就好了.
這麼忙.
忙到每天加班.
很忙.
收拾行李.
不知道出發前的時間夠不夠.
要帶的東西要帶.
要唱的東西要唱.
可能我們去的時候.
我們都會有些猶豫.
會有些這樣的情況.
但是終於這個朋友去完.
他說很感恩.
我當時只是收拾行李.
他說最重要是帶家裡的鑰匙.
還要帶什麼?.
護照.
還有呢?.
如果你還記得你帶身份證.
還有這些是什麼?.
帶錢.
他就是這樣.
最重要的東西.
其他真的忘記帶.
不要緊.
總之已經是記錄了.
定了這段時間去的時候就去.
他終於就坐上了.
他告訴我.
要去.
坐飛機的時候還很愁.
究竟行李裡面的東西夠不夠他用.
帶的錢夠不夠他.
還有很多東西要籌.
但他終於去完的時候.

$^{481}$很想我訂了這團.
不然那兩個星期的忙.
不能慶祝.
他覺得那幾天就好像慶祝.
我過去那麼忙.
就去了一趟.
有時候我們也是.
心裡面看到一些報告.
有時候很想參與.
但又害怕.
又有些疑慮.
所以我們就有些.
去不去好呢?.
這個他給了一筆大錢.
逼著他不去.
所以有時候我想.
教會的活動會不會.
叫大家給一筆大錢.
你不參加.
就不能退還.
或者不參加也要給一筆錢.
來支持這個事工.
這樣來做鼓勵.
做這個什麼呢?.
不過教會不會這樣.
教會很仁慈.
也有些朋友.
他在狀態當中.
朋友約了他也好.
他也很想見朋友.
約了他見面.
但三番四次都要改期.
我們就要想.
他因為有抑鬱症.
所以到差不多的時候.
他不想出門.
不想見人.
我們要留心.
如果有抑鬱症.
或者有自閉症.

$^{521}$當然我們會去治療.
但不要忽略.
屬靈,上高.
你也可能有自閉症.
你也可能有抑鬱症.
要我去做招待.
我見到人.
又說第一次見人要微笑.
我不懂得笑.
不要做了.
可能有很多猶豫.
病可以去治療.
生命可以慢慢去調教.
所以不要因為自己覺得不足.
其實今天站在這裡也有很多不足.
但上帝也猜到我去宣教.
去信主.
我覺得信主已經是很大的事.
上帝這麼寶貴.
把這個救恩給我.
今天你坐在這裡.
我相信你也經歷過上帝的救恩.
臨在你身上.
你也經歷過上帝的實在.
所以我們求神幫助我們.
我們除了守主之外.
剛才我說的訓練不想去.
什麼都不想去.
你不去的時候.
你要想想.
會不會要求神醫治我們.
求神醫治我們.
多走一步.
走出我們的區域.
擴散.
因為大使命就是要你去.
是要擴散.
你的性格內向外向.
內向要走去外向.
外向的經常不可以.

$^{561}$要走回內向.
是要走位的.
這是大使命.
在這裡.
其實我回答的時候.
我會再去的.
因為我深深地感受到.
你願意為神多走一步.
你願意要遵行神的命令.
這是因為愛上帝.
你才願意去遵行.
若不是愛.
如果上帝是用規矩.
或者法律來炮製我們.
我相信我們的頭.
斬幾十次或幾百次都不夠.
大家知道聖經裡有很多命令嗎.
今天沒有人不微笑.
聖經裡有一個命令.
要常常喜樂.
你何時不喜樂.
請你告訴我.
為何不喜樂呢?.
求神又幫你去調教.
因為太多命令在當中.
泰國不可以說黃色的不是.
很多事情要遵守.
如果你說黃色衣服很漂亮.
那也算你了.
但如果說黃色衣服有一點塵.
那也不行.
已經會被斬頭.
今天世界上的法律或制度下.
或者我們看古裝戲.
國王的命令.
你不聽會怎樣?.
你不守會怎樣?.
這樣就沒有了.
我們剛才從崇拜開始.
你已經說是認信.

$^{601}$是試圖信經.
你是相信聖父,聖子,聖靈.
我們相信上帝是我們的神.
但為何我們沒有遵行祂的命令呢?.
在《約翰福音》裡的十四章21節.
不知道有沒有人會唱這首.
(唱).
有沒有聽過?.
這首歌我也很久沒唱了.
是我那時候去短孫.
還沒去泰國孫教.
那時候去泰北的傳道人教的.
因為泰北那裡有國語.
所以這段還沒聽過.
我不知道這首歌的歌詞,歌曲.
在聖經裡說.
如果原來愛神.
今天我們口口聲聲說愛神.
但連祂的話都不遵守.
都不聽的時候.
你這個愛又是甚麼愛呢?.
如果你用中國人的角度.
有些人看那個字.
要小心說那個「愛」字.
因為也有接觸一些孫教士.
譬如在聖經翻譯.
下筆下那個「愛」字要很小心.
譬如中國古時「愛」.
譬如我愛這個榴槤.
「愛」的意思是我要吃掉它.
如果我愛珍姑娘.
我是否要吃掉它?.
有些「愛」字是這樣的.
愛是愛掉它.
直接是拿掉它.
成為我身體的一部分.
但是上帝給我們的這個愛.
是真真實實的,心心相印的那種愛.
今天我也是別人的女兒.
我爸爸媽媽要做甚麼呢?.

$^{641}$他的命令我會馬上遵從.
因為愛父母.
我們說愛這位天父.
而且這位天父是大君王的時候.
你說愛.
聖經裡面說.
愛不是吃掉它.
成為你身體的一部分.
而是遵行他的命令.
守主的道.
就像我們的主題.
所以我們求神幫助我們.
愛.
我們說「去」的時候.
因為上帝的愛.
因為願意愛.
因為我要遵行神的使命.
剛才一崇拜開始.
我們就已經一起唱歌.
因為我沒有配合過歌曲.
今天我也沒有選擇甚麼回應歌.
但是第一首歌叫甚麼名字呢?.
主使我漫遊.
求我心中的王.
究竟主耶穌.
是否你和我心中的王呢?.
求主成為我的意象.
祂的意象不是只給你安居樂業.
這一份安樂,這一份平安.
也是願意透過你和我的生命.
繼續傳遞開去.
接著你們又唱「主的愛」.
這首歌是說我們經歷過上帝的愛.
是很美好的東西.
所以今天「去」.
我們去遵行神的話語.
多走一步.
在讀聖經每日一章.
還未走的那些.
走一走吧.

$^{681}$你去聽吧.
聽也可以.
聽不入心的時候.
其實最重要是心準備好.
入到就不是只聽.
看也是.
所以求神也給你身心健壯.
因為身體不好.
讀經也會有障礙.
但問題是我們會否因著神.
搞好身體.
以至吸收聖經也會好一點呢?.
求神也給我們多走一步.
我們去另外一樣東西.
在去的過程裡.
是經歷主更深的恩典.
在這段經文.
主耶穌頒發大使命的時候.
其實是比.
其實用我們今天現身社會裡的經歷.
可以大概明白到當時的情況.
當時主耶穌受苦.
受難.
釘在十字架上.
在過程裡已經很不容易了.
更加現在.
在十字架上斷氣了.
生命完全沒有氣息了.
所以在經文.
這裡我沒有打出來.
在二十八章開頭的時候.
那些人.
你看下去的時候.
他們心情裡還會有疑惑.
究竟我一直跟隨的神.
是不是真的?.
他已經死了.
很多疑惑.
而且在那個時候.
連彼得在主耶穌還沒釘十字架的時候.

$^{721}$已經三次不認主.
在一個恐怖的氣氛下.
但是後期.
我們已經讀到「仕途行傳」.
你如何看到彼得如何願意為主去很多步.
這就是愛.
和經歷過上帝那種的心恩厚愛.
在這裡.
當然不僅是氣氛.
主耶穌釘死了.
那個時候的一片迷茫.
一到七還沒有看到主耶穌復活.
接著主耶穌復活.
顯然被那些婦女知道.
被他們看到.
然後要告訴他們.
叫門徒去加尼尼山.
我約了要跟他們見.
你會看到中元到十六到二十節.
十一個門徒.
他們去到加尼尼神約定的山裡.
去拜祂.
這群人已經很願意去.
去到加尼尼去找主耶穌.
但是他們還有人會有疑惑.
今天你和我信主了.
無論你信了長或短.
這段時間裡.
我自問也會有疑惑.
上帝,這篇道卡我不是很有信心去講.
也有疑惑.
剛才我們主一學有些姐妹.
每次求神都應該找.
但是每次都會害怕.
這就是我們人.
就是我們人的本性.
我們就是這麼軟弱.
我們真的要每時每刻去靠著主.
中元這群人已經是拜主的.
就像我們今天來到教會裡.

$^{761}$我們很拜主.
但是我們可能或多或少.
我們還有疑惑.
但是上帝主耶穌聖靈.
沒有小看我們這些脆弱.
祂仍然使用我們.
仍然將這麼寶貴的福音給我們.
因為在十八節裡.
主耶穌說.
因為他這個救贖的過程.
完全成了.
就好像十加七而最後成了.
完全是辦妥了.
所以現在天上地下.
所有的權柄都已經賜給了主耶穌.
所以.
「然後才頒布」.
這個就是完整的大使命.
「你們要去洗萬民.
做我的門徒.
奉父子聖靈的名.
為他們施洗.
凡我所吩咐你們的.
都教訓他們遵守.
我便常與你們同在.
直到世界的末了」.
這些一切一切.
如果沒有去.
你不去不動自己.
後面的事就沒有發生.
我怎麼可以和他施洗.
不要緊.
我剛才也說了.
這是一個對公.
但是傳福音的責任和見證.
在生活上的見證.
多去一步.
是我們遵行上帝大使命.
必然的去.
有姐妹說.

$^{801}$我們信了主.
經常都很吃虧.
經常要遷就家人.
我便問.
「你願不願意?」.
「願,因為我信主.
我是主的子女」.
「你願不願意為主.
去吃虧?」.
「但是你放心.
你的吃虧.
你的委屈.
你的不被人理解.
上帝與你同在.
祂知道這些事情.
祂並且要與你同在.
記住.
「以萬來彌」是好得無比的祝福.
主耶穌頒佈大使命的時候.
正如我們所說.
充滿著很多恐怖疑惑.
害怕,謊言.
為什麼有謊言呢?.
因為看守的人都見證.
都看著主耶穌的身體.
真的消失了.
在墳墓裡消失了.
但是有人以為.
錢可以做任何事.
於是用錢來收買這些冰釘.
讓他們被門徒偷取.
不是無緣無故地消失.
不是神跡地消失.
而是被人偷取.
在氣氛裡充滿著這些謊言.
今天也是.
今天這個世界.
我們有很多謊言.
有時候與舊同學.
很多都是微信主.

$^{841}$一走近.
拿著一個袋子.
這個袋子是歐洲回來的.
LV那些品牌的名字.
現在越來越多人吹捧著這些.
因為有機會去接機.
看著年輕的女生.
很年輕的.
後來聽她們聊天.
覺得是她們國內的.
於是附近說.
這些就是富二代.
拿著很多紙袋.
LV,CHANEL.
那些名牌的袋子.
我想她們只有二十多歲.
應該人生都沒有賺夠.
這麼多袋子的價錢.
她自己都肯定賺不到一個.
但她就已經買了好幾個.
很開心地炫耀.
我們有時候.
與我們的朋友.
在微信當中.
在這個世界裡.
都有很多謊言.
人生的成功.
你應該擁有.
起碼要有兩個袋子.
這樣的袋子.
或者.
你的成功應該要有這樣的車.
或者你的成功.
應該要住一間怎樣大的房子.
等等.
很多壞的謊言的價值觀.
或者.
這個社會裡我們都經歷過.
在新冠病症的時候.
我們都會恐懼.

$^{881}$這個人隨時會沒命.
在恐懼當中.
在不容易的情況.
或者現在我們香港的環境.
有很多的變化.
但是.
就算是在這樣的大氣候裡.
在這樣的香港的社會環境當中.
我們的主仍然與我們每一個人一起.
是與我們一起的上帝.
我們人往往都會有疑惑.
但是主耶穌已經承傳了整個救贖的計劃.
今天要告訴你.
這個權柄在主耶穌裡面.
這個天上地下的權柄.
都是在神的手裡.
今天我們試圖信經.
我們認信的時候.
你有沒有信呢?.
信.
有沒有信服呢?.
信服.
有沒有多走一步.
遵行呢?.
求主幫助我們.
我們求主幫助我們.
多走一步.
為了愛主.
並且這個走.
不要小看自己.
也不要小看自己太高.
但是在這個過程裡.
上帝有恩典.
你會發現很多的不足.
上帝會給我們力量.
你又再經歷.
就好像滾雪球.
你越走.
雪球不走.
永遠都是這樣.

$^{921}$但你越走.
走啊.
可能它會碰到一些雪.
你想像一下雪山.
滾著滾著.
碰到一棵樹.
可能本來那麼大.
粉碎了一點.
然後又繼續滾.
都是有收穫的.
都是有經歷的.
我們人生都是這樣.
為了緣故.
我們多走一步.
你就會更加經歷到.
上帝真的有恩典.
上帝真的有那個保守在當中.
傳福音的時候.
我想我們很多時候都很害怕.
不要說傳福音.
只是派單張.
那一刻人們都不接.
酒駕不休.
怎麼會站好.
接你.
很少數.
除非.
除非你猜什麼情況.
人們會站好呢.
不知道你會不會和我這樣想呢.
最近.
這些執法大行動.
本來紅燈.
快快就走.
誰知道.
望著在行.
每個人本來.
走得那麼急.
其實很想快點去.
他要去的地方.

$^{961}$怎麼會站好呢.
但是這些執法的人員.
一說來身份證.
你就沒辦法站好.
這個場面.
什麼時候會發生呢.
當我們去到天家的時候.
你沒得逃.
都知道.
上帝.
我們每個人都要被審.
被審的時候.
你就站好.
為什麼叫你派單張.
你又不派.
和人打招呼.
又不打.
打招呼不是聖經裡面.
聖經裡面沒有說什麼打招呼.
但是問題是表達了.
信徒.
那種.
我們是弟兄姊妹的那種相交.
見面.
送我很久了.
點頭微笑.
我們去.
今天很多的急位.
空位.
我想我們.
因為八月份.
很多家庭.
家庭真的去旅行.
去旅行.
你就會發現.
你會見到不同的人.
你如果去日本的國家.
你都會覺得.
你自己都要守規矩.
因為人家很守規矩.

$^{1001}$你不會亂來.
你就會覺得.
去到日本的國家.
它的氛圍就是.
要守規矩.
要乾淨一點.
我們洪恩家的氛圍是什麼呢?.
人家來到.
這間教會.
真是好了.
很羨慕你們.
我們求神給我們.
有一個洪恩.
屬於主的一種.
愛的關懷在當中.
一種美好的氛圍.
要知道人家來到這裡的時候.
他起碼感受到溫暖.
感受到上帝的愛在這裡.
感受到在座.
每一個不是守主人的人.
是一個遵行主道的人.
我有一個朋友.
很辛苦地帶了一個信徒.
去一間教會.
帶了一個.
都是信徒.
不過他很久沒有回教會.
所以帶他回教會.
那個信徒之後不回教會.
我已經來到這間教會兩個月了.
沒有一個人跟我聊天.
我好像一個隱形人.
我不要了.
我現在.
他在YouTube裡聽到.
但是.
我說了很多次.
在YouTube裡崇拜.
你是經歷不到同在敬拜.

$^{1041}$上帝給我們的那種應許.
不知道每個主人.
回來一起敬拜的時候.
你得不得到那種應許呢?.
你本來沒有力量.
沒有力量來敬拜.
或者沒有力量參與一些聚會.
但是來到的時候.
透過一起唱詩歌.
一起祈禱.
一起認信的時候.
你灰的心可能會紅起來.
炎亮起來.
這是一種經歷.
求神幫助我們.
我們是要在一起的.
也不要以為.
傳福音傳給神的子民.
猶太人就容易一點.
或者傳給某類人就容易一點.
我們去實行神的大使命.
無論是內向.
內向就是生命的根深.
有時候也不是容易一步升達天.
一步一步.
但是你要持續去.
同樣的去傳福音.
也是不容易的.
例如這群弟兄姊妹去了英國.
特別找一些猶太人群體.
來傳福音.
那種拒絕更加厲害.
你不要認識我們的神.
你這個.
說得不好聽的.
快點走開.
不要騷擾我們.
這樣是拒絕的.
也有.
有一些英國教會的宣教士.

$^{1081}$最近回來說.
本來我們教會只有200人.
現在有2000人.
你知不知道發生什麼事了?.
是的,有些人去了移民.
那些牧者開不開心?.
或者那些弟兄姊妹開不開心?.
每天也有一些.
開心,當然.
很多人來教會.
但是問題是.
他們過去.
不要說這是一個大浪潮.
現在突然很多人去.
過去不斷有人陸陸續續去.
過去很久.
所以這些英國的宣教士.
不約而同.
美國,澳洲那些.
我們也是差不多.
就是怎樣呢?.
很多人.
不過這一點是值得開心的.
很多人是沒有教會的人.
是很好而來教會的.
是不是?.
很多初到貴境.
都是去教會.
特別是廣東人.
有得接應.
有得問廣東人.
很多人是存在著.
怎樣在一個新的地方住屋.
在哪裡哪裡可以.
他們說不單是傳道人辛苦.
弟兄姊妹他們也很辛苦.
因為去接送.
真的用愛心去服侍他們.
是不容易的.
但是當這班新的朋友.

$^{1121}$來到這個國家.
英國也好.
美國,澳洲也好.
他們上岸之後.
可以離開了.
不理會了.
難不難?.
很難的.
當然我說有沒有成功?.
當然有.
幸好很開心.
當然有.
不然真的很灰心.
但問題是.
其實去的時候.
很多時候付出比收穫多.
很多時候真實的情況是這樣.
但上帝有個寶貴的安慰說話.
「施比受更為有福」.
當我們去施與給他的時候.
記不記得勝經裡說.
「你為我的命.
給一杯涼水給一個小子」.
上帝也記在心裡.
你這樣用愛去接待遠方.
來到你現在國家的人.
幫他打點一切.
讓他可以安定下來的時候.
你所付上的一切絕不徒然.
我們太與團契.
有些姐妹坐在這裡.
有些真的跟著.
聽我們的統籌和團長說.
真的跟了十十年.
一開始聽都不聽.
但現在已經乖乖地.
接受了成長百況.
剛剛完成了.
所以看到.
不是每個跟進的人.

$^{1161}$每個人都願意繼續成為主門徒.
在過程裡.
都會有些得到利益後.
他就覺得不要了.
他看不到更大的利益.
更大寶貴上帝的恩典.
所以彼此勉勵.
我們願意求神幫助我們.
大使命不只是宣教士.
不只是傳道人.
不只是堂委.
不只是事奉人員.
是每一個每一個都有份.
你知不知道洪水橋這個地方.
有些教會已經派了弟兄姊妹來這裡傳福音.
因為他們的堂會是在天水圍.
他們想要有個分堂.
分堂的地點在香港哪裡好呢?.
因為他們在天水圍.
就會覺得洪水橋.
因為洪水橋在未來有很大的發展.
所以他們來探測地就做探子.
這個時候是用派單張.
或者和一些社區中心去聯繫.
你知道派單張.
當傳道給上面的人看的時候.
我心裡很開心.
不是大人派.
也不只是女人派.
我們現在洪恩堂的探訪隊.
好像只是姐妹.
只是女人.
我們的弟兄們都起來了.
不只是男人,女人,大人.
小朋友也一起去派.
不知道你在洪水橋有沒有遇到他們.
他們多數派都是星期天的時間.
你看到嗎?.
他們是走一步.
求神幫助我們.

$^{1201}$要去的.
不去傳是有禍的.
聖經裡面.
哥林多前書保羅去勸勉他們.
我們傳福音沒有可誇的.
不要說我們很厲害.
我們是報道隊的.
我們是甚麼甚麼.
如果我們不傳.
我們是有禍的.
你想想.
我們是一個負心的人.
上帝給我們這麼好的心.
這麼多的恩典,這麼多的愛.
更加是最寶貴的.
就是救贖的恩典.
給我們的時候.
你竟然不繼續將這件事.
好像上帝的心意繼續傳開去.
讓更多的人得福.
你沒有去做的時候.
你說是不是上帝給了你良心.
你就當甚麼?.
狗肺.
所以求神幫助我們.
是要去的.
你還可以想到很多去的例子.
有些姐妹去過.
我不知道你們有去過.
比我還要早.
因為這個地方.
以前不是這個地址.
是樓下的.
現在搬到上面那層.
因為疫症期間.
我都沒有去過找她.
我去談的時候.
你覺得這裡像甚麼?.
你是不是去過?.
這個真的是診所.

$^{1241}$這是一個牙醫的診所.
你再看下去.
他擺了很多畫.
多謝他.
好像去了警會.
因為警會的特徵.
擺了很多不同國家的擺設.
很多不同國家的擺設.
這幅牆.
我身為泰國的宣教士.
在18年前已經接受過他的恩惠.
因為這個牙醫.
他願意幫宣教士.
免費去看牙齒.
18年前我已經去看他.
我小時候已經很怕去看牙醫.
所以我叫黎英哥一次去看牙醫.
我其實不是因為牙痛而哭.
是幫他整的時候.
我感受到他那種愛.
他幫你整牙.
但我覺得真的好.
為甚麼這麼好.
我只是做宣教士.
為甚麼他幫我看牙.
因為我小時候已經要看牙醫.
所以很怕.
一說不適合就要去剝牙.
又要整.
總之很痛苦.
但為甚麼他這麼好.
想起他這麼好.
又想起上帝的恩典這麼好.
於是眼淚就收不到了.
很感觸就流了.
結果其實很多人在等.
他竟然說.
不要緊,你們先休息一下.
全部人都走了.
連護士也走了.

$^{1281}$跟著他才回來.
我將這件事告訴一個很會畫畫的弟兄.
我說是這樣的.
不過他就將它變成卡通.
好像很害怕整牙一樣.
看不看到.
但牙醫用他的專業去服侍.
當你看到他的德行的時候.
你會建立自己.
所以我們今天.
多去一步.
剛才說我們願意為愛主.
我們受委屈,吃虧.
他也很吃虧.
他可以賺多點錢.
但他願意這個時間.
給宣教士看牙.
他不單是你們來.
他甚至去非洲.
去一些國家需要的地方.
幫當地人去看牙.
這時候你會發現.
他就是去了.
使萬民作神的門徒.
不只是傳福音.
是他的見證美麗到.
我要保持好上帝兒女的身份.
好好地成為上帝的門徒.
是這樣的.
這個去字.
我們泰國話.
先不說.
一會兒我們請泰國姐妹讀給你聽.
怎樣讀.
我很想告訴你.
有很多時候我們會覺得.
因為你留心看.
好幾個星期.
如果你看天氣預測.
我今天也看過.

$^{1321}$說9點到10點是雷暴.
還寫90\%多下雨.
結果今天好像沒有.
但問題是.
早幾個每個星期差不多.
差不多就是10點.
9點多下大雨.
所以我們回來的主學.
泰國語姐妹.
個個都濕透.
我說不要緊.
希望待會不會下雨.
等待崇拜的人.
不用濕透才來敬拜.
有時候可能因為天氣的緣故.
我們會覺得不要去教會.
不要緊.
現在流行Zoom.
在Zoom崇拜就算了.
但不是的.
我們的姐妹.
你知道嗎.
我們其中一位泰語姐妹.
在摩星嶺來崇拜.
你們有沒有遠過摩星嶺?.
我們近一點再比一比.
另外一個遠的.
也在黃大仙區.
有沒有人在黃大仙區?.
遠過黃大仙區?.
第一,第二我們也包了.
泰語團體.
他們都是來這裡敬拜.
所以我們求神幫助我們.
我們很希望.
我們紅印堂的弟兄姊妹.
固有在這裡的會友繼續在這裡.
也讓我們這群天父的兒女.
主耶穌的門徒.
好好地做好我們的本分.

$^{1361}$多去一步.
你已經多去一步了.
我看到很多姐妹很好.
不用別人叫.
我覺得就要做了.
我覺得是很好的一件事.
就好像一個家裡.
哪裡有位置.
我就互為寶珠.
就是走位.
求神幫助我們.
我們還要說一個大喜訊給你聽.
就是沒有我在這裡.
團契和統籌和我們的姐妹.
就成功地連續八課.
在星期四.
就是平日裡面.
就一起去接受訓練.
成長八課的訓練.
於是他們就一起去看這個八課.
因為裡面有一個姐妹是需要授訓的.
其他的就一起去聽.
現在暑假我們放一放假.
希望弟兄姊妹都紀念我們.
我們希望完了暑假之後.
我們繼續好不好?.
我不知道我們的泰語姐妹會怎樣.
其實我們泰語團契.
已經是每個星期都聚在一起.
因為一起的時候.
你會發現祈禱的力量是多一些.
一起查經.
就算今天.
其實你知道一早上一開始都不好.
你經常會覺得很顫抖.
今天一早開始都不好.
因為電腦裡面進不去.
所以今天早上十點.
教主就說.
你泰語怎麼都歪了音.

$^{1401}$怎麼好像深圳一樣.
他們很有趣.
他們還沒發現我深圳的時候.
他們已經一開始就不同平時.
又領詩歌又自己帶祈禱.
在接近兩年前的時候.
不是這樣.
祈禱你帶什麼什麼.
現在自發力了.
就是需要這樣.
多走一步.
在你的團體裡面.
我們知道你們有師班的組合.
有很多步.
很多探訪步.
或者師事步.
招待步.
每人多走一步.
你今天沒有步的.
多走一步.
因為如果你不去的話.
就好像我們泰文這個讀去字的音.
是怎麼讀的.
泰文字大於展會.
你們聽到什麼.
閉.
就真的閉家房.
因為我們香港人.
你不去就閉了.
我已經很快了.
\newpage



\section{詩篇 1:1-6}
\label{sec:iQ2CTi_2bNs}
\textbf{2023.09.03 《晝夜思想神的話語》 詩篇 1\_1-6 (和修版) 講員 : 老耀雄牧師}
\newline
\newline
連結: \href{https://youtube.com/watch?v=iQ2CTi_2bNs}{\texttt{https://youtube.com/watch?v=iQ2CTi\_2bNs}} ~~~~ 語音日期: 2023-09-04
\newline
\newline
\hyperref[sec:TnVaeXBM1r0]{\small{< < < PREV SERMON < < <}}
~
\hyperref[sec:index_chronic]{\small{[返順時目]}}
~
\hyperref[sec:index_scriptual]{\small{[返順卷目]}}
~
\hyperref[sec:Wzp5u5hQ940]{\small{> > > NEXT SERMON > > >}}
\newline
\newline
詩篇 1:1-6
\newline
\begin{longtable}{cl}
\hline
\hline
章節 & 經文 (和合本修訂版)\\
\hline
1:1 & \begin{tabularx}{0.7\textwidth}{X} 不從惡人的計謀，不站罪人的道路，不坐傲慢人的座位， \end{tabularx} \\ \\ \relax
1:2 & \begin{tabularx}{0.7\textwidth}{X} 惟喜愛耶和華的律法，晝夜思想他的律法；這人便為有福！ \end{tabularx} \\ \\ \relax
1:3 & \begin{tabularx}{0.7\textwidth}{X} 他要像一棵樹栽在溪水旁，按時候結果子，葉子也不枯乾。凡他所做的盡都順利。 \end{tabularx} \\ \\ \relax
1:4 & \begin{tabularx}{0.7\textwidth}{X} 惡人並不是這樣，卻像糠秕被風吹散。 \end{tabularx} \\ \\ \relax
1:5 & \begin{tabularx}{0.7\textwidth}{X} 因此，當審判的時候惡人必站立不住，罪人在義人的會眾中也是如此。 \end{tabularx} \\ \\ \relax
1:6 & \begin{tabularx}{0.7\textwidth}{X} 因為耶和華知道義人的道路，惡人的道路卻必滅亡。 \end{tabularx} \\ \\
[1ex]
\hline
\hline
\end{longtable}
$^{1}$主席.
各位觀眾姐妹平安.
對於以色列人來說.
神除了帶領他們離開埃及這個圍爐之地.
讓他們擺脫了一個奴隸的身份.
還允許他們得到迦南這塊流奶與蜜之地.
讓整個民族得到一個安身立命的土地.
不過要得到神所允許的地.
總是要付出一些代價.
約書亞在帶領以色列人進入迦南的時候.
上帝對約書亞和整個以色列民族作出了一個指示.
在約書亞記一章七至八節這樣說.
祂說只要剛強大大壯膽.
謹守遵行我僕人摩西所吩咐的一切律法.
不可偏離左右.
使你無論往哪裡去都可以順利.
這個律法書不可以離開你的口.
總要就耶思想.
好使你遵守這個書上所寫的一切話.
如此你的道路就可以亨通.
凡事順利.
上帝對以色列民族的指示很簡單.
就是遵守摩西所吩咐的一切律法.
以及律法書不可以離開你的口.
總要就耶思想.
這樣耶華所指的律法書究竟是指哪本書呢?.
舊約聖經的分類一般來說有兩個傳統.
如果以希伯來文聖經來分類的話.
就分為律法書,先知書以及聖言.
按照舊約希勒文聖經七十事譯本的分類的話.
就按體材分類.
分為摩西五經,歷史書,詩歌智慧書和先知書.
因此一般信徒都有一個錯覺.
摩西五經既然就是律法書.
就是上帝所指的律法的全部.
當然了.
五經裡不單記載了以色列人.
在西乃山上領受神所誇報的十誡.
也有利未記和申命記對於律法詳細的解釋.
不過摩西五經卻並不是律法的全部.

$^{41}$律法一詞是源自希伯來文Torah.
Torah一詞.
而Torah的原文含義卻比律法更加闊.
可以解釋為教導和教誨.
所以Torah的合義解釋可以稱為律法.
但是它講義的解釋就是神的訓誨.
而全本新舊約聖經都包含著神的訓誨.
都是需要信徒去閱讀和關注的.
因此上帝的教誨不單單來自摩西五經.
甚至詩歌智慧書以及先知書.
都同樣有著上帝的訓誨.
只不過是用不同的方法表達方式去述說上帝的聖言.
我試試舉兩個例子.
引用耶穌兩段說話給我們做一個思考.
不單單只有摩西五經才有律法的指引.
在向福音第十章.
耶穌回應那些用石頭打他的猶太人.
因為猶太人認為耶穌說了一些說述的話.
認為自己就是神.
於是耶穌在向福音第十章34節這樣說.
你們的律法書上不是寫著.
我曾說你們是諸神嗎?.
經上說是不能廢的.
如果那些領受神道的人.
神上就稱他們為諸神.
這樣父所分別為聖.
猜測世上來的那位說.
我是神的兒子.
你們還對他說.
你說是俗的話嗎?.
耶穌所指的律法書.
不是來自摩西五經.
而是來自詩篇82篇第六節.
那裡寫說.
我曾說你們是諸神.
都是至高者的兒子.
另一段經文就是在.
約翰福音十五章.
耶穌提到世人對他的憎恨.
耶穌說.

$^{81}$這也要應驗你們律法書上所寫的話.
他們無故恨我.
耶穌所指的律法書.
也不是指摩西五經.
也是引自詩篇109篇第三節.
當中這樣寫.
他們圍繞我說怨恨的話.
也無故地攻打我.
耶穌兩段對話所引用的經文.
都是將詩篇的內容.
視為律法的內容.
可見摩西五經並不是.
上帝所展示律法的全部.
或者用另一個說法.
上帝的訓誨.
貫穿整本新舊約之間.
而摩西五經內的律法.
只是上帝對世人訓誨的.
其中一個部分而已.
因此今天我想帶大家.
在詩篇的第一篇裡面.
去尋找上帝所展示的訓誨.
我們一起去祈禱.
「差因主,你話語訓誨何其寶貴.
你話如兩刃劍.
剖開我們的意念.
甦醒我們的靈魂.
使我們能夠回歸到你的身邊.
享受你賜予我們的豐盛.
逆行在你的旨意當中.
求你讓每一個渴慕你的.
都願意晝夜思想你的話語.
好叫我們能夠明白你的心意.
領受你的意象.
我們的祈禱,仰望.
是奉主耶穌基督得勝明智祈求.
阿們」.
上帝對約書亞說.
「這個律法不可以離開你的口.
總要晝夜思想.

$^{121}$好似你謹守遵行.
這書上所寫一切的話.
如此你的道路就可以亨通.
凡事順利」.
這番說話和詩篇的第一篇第二節.
「唯喜愛耶華的律法.
晝夜思想他的律法.
者人便同為有福」.
這個內容和精神是完全一致的.
詩人認為.
喜愛耶華的律法.
是上帝對人類正確生活的指示.
你喜歡什麼.
將會塑造你成為一個什麼的人.
假如你喜歡錢.
你就會敬拜馬門.
如果你喜歡權力.
你就會踐踏別人.
作為你上升的階梯.
如果你喜愛權力.
如果你喜愛自己.
你就會成為一個自以為是.
以及自私的人.
唯有你喜愛神的律法.
你才會有一個愛神愛人的人生.
亦唯有你喜愛神的律法.
你才願意遵守律法的內容.
這就是詩人對律法的一個立論.
詩人繼續論述.
怎樣是一個有福的人呢.
有福的人不單要喜愛神的律法.
還要晝夜思想.
這種晝夜思想是什麼意思呢.
不是叫我們口中經常背誦金句.
不是.
你每天一金句.
不是神的心意.
神的心意就是將他的訓誨.
內化在他的生命裡面.
生活之中每時每刻都有神的話語.

$^{161}$作為鼓勵,安慰,引導.
能夠讓信徒活出信仰.
亦能夠從信仰實踐裡面.
感到興奮和雀躍.
晝夜思想律法就是內化生命.
亦是被神塗造的一個過程.
詩人繼續論述有福的定義.
亦有別於世人對有福的看法.
世人所指的有福.
多數都是自身的利益.
例如身體健康.
有身份,有地位.
甚至有車,有樓.
生活無憂.
世人所定義的有福.
大部份都與物質有關.
但詩人所指的有福.
卻是與上帝能夠有生命的結連.
與神有份才是有福.
其他一切都不是.
亦唯有生命與神的律法內化.
才能與神有份.
否則都只是空有基督徒的名份.
卻無基督的香氣.
在身體裡流露出來.
傳道書怎樣認為呢?.
他說這種基督徒是白活一場.
只是空有基督徒的名.
卻無基督徒所踐行的信仰.
最後都是虛空.
詩人繼續論述有福的內容.
他說他要像一棵樹栽在溪水旁.
按時候結果指.
葉子也不枯乾.
凡他所做的盡都順利.
樹栽在溪水旁.
正正是樹能夠不斷在溪水中.
支取成長的養分.
生命不斷有源源不絕的溪水供應.
以致葉子是不會枯乾的.

$^{201}$這種源源不絕生生不息的供應.
令果實從幼嫩到成熟.
最終可以成為可供食用的甜美果子.
有福的人生不是只顧享受.
或將一切據為己有.
而是生命不斷被神塑造.
被滋養然後結出果子.
造福人群.
按著這種生命軌跡成長的人生.
他所做的一切盡都順利.
弟兄姊妹.
神的律法包含神的話語.
教會今年倡導讀經計劃.
正正是想肢體是在沐浴在神的話語裡.
能夠享受在溪水旁的供應和餵養.
這是教會的一番心思.
讓每一個肢體都能夠享受有福這個恩典.
你願不願意從今天起.
開始喜愛神的分魁.
開始讓神的話語內化在你的生活裡.
教會的讀經計劃你跟了多少呢?.
你會不會從今天起重拾你的讀經計劃.
讓神的話語每天都引導你呢?.
有些人說.
Gary,現在幾月了?.
九月了.
我都甩了八個月了.
九月才開始會不會遲一點呢?.
我們說他覺得很遲.
但神卻認為只要你開始.
福氣就輪不到你了.
所以不要怕遲.
總要開始.
詩篇第一篇不單單講有福的人.
詩人亦都描述那些不合神心意的人.
就是那些惡人,罪人以及傲慢人.
他們追求甚麼呢?.
就是計謀,道路和座位.
惡人所追求的是計謀.
就是為個人利益不擇手段.

$^{241}$甚至損害他人都不在乎.
惡人心裡沒有律法規範.
因為他認為自己就是律法.
罪人所追求的是道路.
亦都是一條捷徑.
或者一條容易走的路.
甚至是一條不易的路.
方向偏差了.
最終走向滅亡.
罪人心裡沒有律法.
因為他按自己的私慾而行.
傲慢人所追求的是座位.
亦都是屬世的權力.
是從驕傲所產生出來的果子.
而驕傲所帶來最大的偏差.
就是自己覺得自己就是神.
傲慢人心裡沒有律法.
因為他自視自己是神.
施人建議我們不要跟從.
不要站立.
不要坐在這裡.
就是不要跟從惡人的價值觀.
以神的心意為生命的方向.
不要站在錯誤和偏差的路口.
要走在二路走向永恆.
而不是滅亡之路.
不要貪戀在高位上的那種虛榮.
享受地上的賞賜.
要積財寶在天上.
因為施人知道這些人的後果.
施人怎說呢.
惡人像康丕被風吹散.
在詩篇三十七章.
三十七篇九至十節.
用得更加詳細.
他說因為作惡的必被剪除.
唯有等候耶和華的必承受土地.
還有遍事.
惡人要歸於無有.
你就是砌殺他的住處也不存在.

$^{281}$歸於無有就是惡人的下場.
而詩篇一百零三篇十五至十六節.
他不單單描述惡人的下場.
他更加描述所有不認識神的人的下場.
至於世人.
他的年日如草一樣.
他興旺如野地的花.
驚風一吹就歸無有.
他的原處也不再認識他.
不認識神的人最後怎樣.
歸於無有.
甚麼都沒有.
這個惡人為甚麼歸於無有呢.
施人進一步解釋.
當審判的時候.
惡人必站立不住.
其實有誰能夠在審判台前站立得住呢.
沒有.
沒有人可以在審判台前站立得住.
今日魔鬼最成功的伎倆.
就是讓人不覺得有審判的存在.
只要人認為沒有審判.
就不會再理會犯罪帶來的結果.
也不會認為神會追討所犯的罪.
以致人更加放縱地犯罪.
不斷地得罪神.
不斷地得罪人.
當施人預視了.
他惡人的道路.
卻必滅亡的時候.
我們這批屬神的子民.
會不會應該檢視一下.
我們自己的人生.
或者我們都認為.
我們都不是一個惡人.
我們又沒有作奸犯科.
我們不屬於惡人的.
牧師.
我們只是平凡的人.
難道真的平凡也有罪嗎.

$^{321}$施人一早就提出了.
他怎麼說呢.
為喜愛耶和華的律法.
晝夜思想他的律法.
這人便為有福了.
為是什麼意思呢.
為有.
為一.
為獨.
正是這句話的一個重點.
我們應該問問自己.
我們生命中的為一.
是不是只有神的話語呢.
我們生命中的為獨.
是不是只有神的心意和計劃呢.
我們生命中的為有.
是不是只有主耶穌呢.
還是我們的唯一.
仍然都是自己呢.
我們是不是在溪水旁.
接受源不絕的供應和餵養呢.
我們是不是又能夠.
定時地按時結束果子呢.
如果答案是否定的話.
我們就要回轉了.
各位姐妹.
今天講題是晝夜思想神的律法.
怎樣可以晝夜思想神的律法呢.
真的要日背金句嗎.
唯有我們多讀信經.
我們才可以將神的話語內化.
不單單是讀.
還要去研經.
在團契裡面有查經.
團契沒有查經.
小組有查經.
總是要在聖經裡面深入去認識.
唯有將信仰內化.
我們才可以被神去塗造.
唯有生命被神塗造了.

$^{361}$我們才能夠承擔神的使命和遺像.
生命中有神話語的滋潤.
專治能夠回應.
呼召活出一個豐盛的人生.
在今年年初.
一二月呢.
我就有個安息年假期.
很感恩.
能夠和家人去日本和台灣.
享受家庭樂.
在台灣我有機會探訪一個新竹的牧師.
他叫做蕭偉光牧師.
今年六十六歲了.
他和我們夫婦分享了一些生命的見證.
我也很想在這裡和大家分享一下.
蕭牧師在未進入台灣浸信會神學院.
接受裝備的時候.
他已經有個心智.
很想服侍客家人.
不過當時不知道在哪個地區服侍.
在他畢業前一年.
他領受了一個從神而來的異象.
這個異象說到神帶領他去了一個山谷.
山谷裡有一班人.
這班人很無知.
他們比賽.
誰跳進山谷裡面不死的就是英雄.
不知道他為甚麼會這樣想.
可能被人誤導.
被世俗的價值觀誤導.
最後蕭牧師見到很多人死在山谷裡面.
場面很悲慘.
於是他就呼求耶穌.
耶穌耶穌你救他們吧.
耶穌回應他.
他說你去搶救這班人.
你去搶救這班住在偏遠又失望的人.
然後牧師將異象寫成一首詩歌.
叫做《求主差派》.
後來更加製成錄音帶.

$^{401}$成為他傳福音的工具.
蕭牧師在神學院畢業之後.
領受了異象.
和太太一起去新竹橫山.
成為這個地區的一個服侍.
因為這個地區住了很多客家人.
他成立了中華托方團社區服務中心.
開方直堂.
當時山區很窮.
根本沒有人願意去服侍.
不過由於蕭牧師有異象.
因此就在橫山這個教會服侍了三十六年.
三十六年呀.
這是一個很長的日子.
三十六年同一間教會.
如果沒有異象沒有使命.
我相信這個事奉很難去維持.
為甚麼這樣說呢.
繼續聽下去.
蕭牧師提到山區教會.
山區距離教會有一段距離.
為了會友們能夠回教會.
他星期日要駕車去山區不同的地方.
去接載兒童回教會.
不是一轉而是很多轉.
主一學要等到齊人才開始上課.
探訪是蕭牧師每一天必定會做的工作.
因為沒有人會主動來教會.
這就是他去橫山教會的開方經歷.
他提到教會很窮.
但仍有宣教的活動.
他會帶會友參加馬來西亞的宣教探訪.
而會友亦都很熱切地參與宣教奉獻.
去支持不同的宣教工場.
宣教亦都是蕭牧師的另一個異象.
蕭牧師跟我分享到.
他明白新竹是一個鄉村.
所以很多年青人長大之後.
都必定會外出工作和發展.
不會有年青人願意留在一個貧窮的鄉村去為生.

$^{441}$他知道神父的使命.
只是牧養和栽種這批年青人一段時間.
他就要放手.
教會曾經有接近一百人的興旺時期.
但只有十個人這個低谷期.
不過牧師甘心樂意接受這個現實.
以及接受神父的呼召.
他提到三十六年.
曾經有十幾個童工在這裡事奉過.
但最後這些童工都選擇離開.
不願意接棒.
為什麼呢?.
因為童工的期望和童工看到的落差是有落差的.
作為一個傳道人.
當然希望教會興旺.
人數不斷增長.
但原來到某個時候又倒轉頭.
牧養了這些年青人.
長大之後就全部走了.
又再痛頭離過.
你甘心這樣去一個循環去做嗎?.
因此我自己聽完他這個分享之後.
我感受到異象和照明的重要性.
以牧人一間教會.
人數不斷進進出出.
牧者只是只問耕耘.
不問收穫的心智.
牧者需要多少堅持和忍耐呢?.
異象以及使命.
就是信徒能夠建立信仰的重要元素.
在這裡我想起洪映嘉.
洪映嘉的福音如曲咀.
正正是秋萍和惠清老師有異象和使命.
所以才有19年的歷史.
今年就踏入20週年了.
20年.
20年了.
你想想有多少堅持和忍耐才可以呢?.
我又見到簡太在洪映嘉的詩琴服侍.
簡太到今天還是有心有力.

$^{481}$我相信這個一定是簡太的異象和使命.
讓她能夠堅持到今天.
不單單是他們兩個.
我聽過很多感人的分享.
包括晚珠姐說.
我73歲了.
73歲.
我仍然可以被神使用.
可以參與崇拜主席的事奉.
很感恩啊.
我想很多教會.
73歲的會友.
今年會友.
差不多都不安排他們在教會做事奉.
休息一下吧.
不是.
晚珠姐說.
我能夠參與事奉.
很感恩.
碧霞姐說.
她和公司協調了日子.
要參加星期四晚侍班的練習.
要參與侍班的事奉.
不單這樣.
還要帶著她78歲的老公創峰一起參加.
78歲.
參與侍班事奉.
每個星期回來練歌.
不單這樣.
教會有參與家長英語班的事奉.
有Lisa.
有玉芬.
有師母.
除了平時星期二的上課之外.
三個人特別抽出星期五.
一起去跟進學員的成長.
叫他們回來教會.
一起去.
美其名就是練習一下英文.
實際上就是關心她.

$^{521}$這不是因為有意象有使命.
然後才願意付出嗎?.
我較早前聽到.
葛良參與IT的事奉.
葛良為了參加IT部早上的培訓.
因為梁權的通常他培訓是早上的.
於是葛良怎樣呢?.
原來葛良喜歡看球賽.
看球賽一般看到三,四點.
於是他就不看球賽.
不看球賽早點睡覺.
以致早上可以上梁權的培訓.
感不感動?.
我不知道當時如果是世界盃.
葛良會不會受不了.
還是要看世界盃.
不培訓了.
不知道.
星期六可能會有個動工.
不單單是這樣.
不單單我剛才所說.
我們教會有兩個姐妹.
參與了報道的課程.
很願意去毀身.
我相信在洪恩嘉年紀不是阻礙我們追求意象和實踐署名的原因.
不是的.
你現在聽到我說的話.
每個人的年紀都超級了.
不是說他們年紀的問題.
而是他們有沒有心的問題.
你願意有一個毀身的心.
你願意有一個事奉的心.
才是令我們不抗拒神對我們呼召的一個重要原因.
不要因為你的年紀而影響神對你的呼召.
當我們說教會的意象.
洪恩的意象是要成為福音的平台的時候.
洪恩嘉要成為洪水橋傳福音的平台.
我們應該問問自己.
我可以參與一個什麼的事奉.
以致洪恩嘉才能夠落實到這個福音的平台.

$^{561}$而不是一個口號.
每一週每一個星期.
那麼殷勤的回來參與崇拜.
不應該是你的使命意象.
而是你的生命意象.
而不是你的生命的平台.
而是你的生命的平台.
而不是你的生命的平台.
我們要用眼光去看我們身邊的人.
如果有看四福音的時候.
你會看到耶穌.
當祂進入人群的時候.
正經怎麼寫呢?.
祂就動了慈心.
每一次耶穌動了慈心.
祂就會施行醫治,神蹟.
我們進入人群的時候.
我們有沒有從神而來的眼光.
對著這些人動了慈心.
當是對人冷漠的時候.
當是對那些人視而不見的時候.
異象使命是不會發生的.
當是因循享受我們自己個人的信仰生活的時候.
異象使命都不會發生的.
惟有神的幻覺內化在我們生命裡面.
當我們見到有需要的時候.
那種就會動了慈心.
我上個星期有機會和華娟姐去探望甜甜.
為甜甜私餐.
因為她身體的緣故回不了教會.
回不了教會去領受聖餐.
我就去她家裡去私餐.
在聊天的過程裡面我就知道.
原來華娟姐一直陪著她去覆診.
甜甜每一次覆診華娟姐都和她一起.
沒有人知道.
沒有人知道她這樣的時逢.
她就這樣默默耕耘.
默默陪著甜甜每一次覆診.
她怕醫生和甜甜說的話她聽不清楚.

$^{601}$於是每一次覆診都陪著她.
沒有人要求她這樣做.
但我相信她做了大概三,四年.
這個就是她的使命異象.
這個使命異象燃點著她整個人生.
華娟姐都七十多歲了.
是不是?.
我們一起祈禱.
(主耶穌基督的禮物).
我們感到不悔.
你的愛是何等的財富高深.
是何等的恩典.
求你讓在座的每一個.
都不要輕看你的犧牲.
要好好珍惜.
更加不要輕視你的吩咐.
要遵守你的吩咐.
做一個有福的人.
求你恩惠常與我們同在.
讓我們的生命能夠活出信仰.
讓我們身邊的朋友.
可以從我們的生命裡面.
見到主耶穌.
讓他們都可以來到主的面前.
領受永恆的福樂.
主啊.
我們每一個都很軟弱.
求你兼顧我們的心.
讓每一個.
參與讀經計劃的每一個肢體.
都能夠享受讀經的樂趣.
享受在讀經之中.
與你相遇的甘甜.
從你的話語裡面得到滋潤.
鼓勵和安慰.
我們祈禱.
是奉主耶穌基督.
得性明智祈求.
阿們.
願主祝福大家.

$^{641}$.
謝謝大家收看 .
\newpage



\section{馬太福音 11:1-19}
\label{sec:Wzp5u5hQ940}
\textbf{2023.09.10 《不一樣的彌賽亞》 馬太福音11\_1-19 (和修版) 講員 : 謝靄薇牧師}
\newline
\newline
連結: \href{https://youtube.com/watch?v=Wzp5u5hQ940}{\texttt{https://youtube.com/watch?v=Wzp5u5hQ940}} ~~~~ 語音日期: 2023-09-17
\newline
\newline
\hyperref[sec:iQ2CTi_2bNs]{\small{< < < PREV SERMON < < <}}
~
\hyperref[sec:index_chronic]{\small{[返順時目]}}
~
\hyperref[sec:index_scriptual]{\small{[返順卷目]}}
~
\hyperref[sec:BzeYjKv9RX8]{\small{> > > NEXT SERMON > > >}}
\newline
\newline
馬太福音 11:1-19
\newline
\begin{longtable}{cl}
\hline
\hline
章節 & 經文 (和合本修訂版)\\
\hline
11:1 & \begin{tabularx}{0.7\textwidth}{X} 耶穌吩咐完了十二個門徒，就離開那裡，往各城去傳道，教導人。 \end{tabularx} \\ \\ \relax
11:2 & \begin{tabularx}{0.7\textwidth}{X} 約翰在監獄裡聽見基督所做的事，就派他的門徒去， \end{tabularx} \\ \\ \relax
11:3 & \begin{tabularx}{0.7\textwidth}{X} 問耶穌：「將要來的那位就是你嗎？還是我們要等候另一位呢？」 \end{tabularx} \\ \\ \relax
11:4 & \begin{tabularx}{0.7\textwidth}{X} 耶穌回答他們：「你們去，把所聽見、所看見的告訴約翰： \end{tabularx} \\ \\ \relax
11:5 & \begin{tabularx}{0.7\textwidth}{X} 就是盲人看見，瘸子行走，痲瘋病人得潔淨，聾子聽見，死人復活，窮人聽到福音。 \end{tabularx} \\ \\ \relax
11:6 & \begin{tabularx}{0.7\textwidth}{X} 凡不因我跌倒的有福了！」 \end{tabularx} \\ \\ \relax
11:7 & \begin{tabularx}{0.7\textwidth}{X} 他們一走，耶穌就對眾人談到約翰，說：「你們從前到曠野去，是要看甚麼呢？看風吹動的蘆葦嗎？ \end{tabularx} \\ \\ \relax
11:8 & \begin{tabularx}{0.7\textwidth}{X} 你們出去到底是要看甚麼？看穿細軟衣服的人嗎？那穿細軟衣服的人是在王宮裡。 \end{tabularx} \\ \\ \relax
11:9 & \begin{tabularx}{0.7\textwidth}{X} 你們出去究竟是要看甚麼？是先知嗎？是的，我告訴你們，他比先知大多了。 \end{tabularx} \\ \\ \relax
11:10 & \begin{tabularx}{0.7\textwidth}{X} 這個人就是經上所說的：『看哪，我要差遣我的使者在你面前，他要在你前面為你預備道路。』 \end{tabularx} \\ \\ \relax
11:11 & \begin{tabularx}{0.7\textwidth}{X} 我實在告訴你們，凡女子所生的，沒有一個比施洗約翰大；但在天國裡，最小的比他還大。 \end{tabularx} \\ \\ \relax
11:12 & \begin{tabularx}{0.7\textwidth}{X} 從施洗約翰的日子到今天，天國受到強烈的攻擊，強者奪取它。 \end{tabularx} \\ \\ \relax
11:13 & \begin{tabularx}{0.7\textwidth}{X} 眾先知和律法，直到約翰為止，都說了預言。 \end{tabularx} \\ \\ \relax
11:14 & \begin{tabularx}{0.7\textwidth}{X} 如果你們願意接受，這人就是那要來的以利亞。 \end{tabularx} \\ \\ \relax
11:15 & \begin{tabularx}{0.7\textwidth}{X} 有耳的，就應當聽！ \end{tabularx} \\ \\ \relax
11:16 & \begin{tabularx}{0.7\textwidth}{X} 「我該用甚麼來比這世代呢？這正像孩童坐在街市上向同伴呼喊： \end{tabularx} \\ \\ \relax
11:17 & \begin{tabularx}{0.7\textwidth}{X} 『我們為你們吹笛，你們不跳舞；我們唱哀歌，你們不捶胸。』 \end{tabularx} \\ \\ \relax
11:18 & \begin{tabularx}{0.7\textwidth}{X} 約翰來了，既不吃也不喝，人們就說他是被鬼附的； \end{tabularx} \\ \\ \relax
11:19 & \begin{tabularx}{0.7\textwidth}{X} 人子來了，也吃也喝，他們又說這人貪食好酒，是稅吏和罪人的朋友。而智慧是由它的果子來證實的。」 \end{tabularx} \\ \\
[1ex]
\hline
\hline
\end{longtable}
$^{1}$各位弟兄姊妹平安.
這時候我們真的很需要平安.
剛才跟姐妹們聊天的時候.
原來有些人是從黃大仙來的.
今天的情況仍然令到鞋子都濕了.
大家從不同的地方來.
相信你們都會體會到這個平安.
你們的近況有沒有令到你們有難處.
我們真的要在教會為弟兄姊妹守望.
有時真的不知道.
有時我們真的在低谷.
在我們看神的話語之前.
我們一起仰望神.
每有恩典每有慈愛主.
每當我們可以去敬拜你.
這是你給我們的福分.
你給我們的恩典.
我們向你感恩.
特別是在暴風雨之後.
你仍然讓我們聚集.
可以來到你面前.
這完全是你的恩典.
你讓我們可以守聖餐.
也可以去敬拜你.
還有求你幫助我們不忘記你的恩典.
我們時時都要數算主因.
願主這時候.
我們心中有什麼是不夢理如立的.
求聖女人當中剔除.
以致我們可以專心一意.
看到你的話語.
我們就願意去進行.
求女人當中帶領我們的心.
有種孩子心中的意念.
口中的言語是夢理如立.
奉主耶穌基督命求.
阿們.
當我們去看.
不一樣的彌賽亞的時候.
其實彌賽亞.

$^{41}$這個什麼意思呢.
這個詞.
還沒到這裡.
這個詞是說到.
基督.
彌賽亞就是基督.
也是受高者.
我們記住這個名詞.
在舊約時代.
有三種人.
是可以受高的.
有君王.
有祭司.
有先知.
他們都是受高的.
但很可惜神揀選了一些人.
來領導他們的百姓的時候.
但這些百姓.
他們繼續犯罪.
所以犯罪導致.
他們的國家怎樣呢.
被異族統治.
是很慘的警方.
他們受盡很多的痛苦.
於是先知.
就去預言.
有一位特別的彌賽亞.
或者受高者.
要來臨拯救你們.
他們就在等待.
等待一個救星.
拯救他們.
在加利利地區.
在猶大地區.
他們都在尋找彌賽亞.
每個人都有不同的動機.
去想彌賽亞是怎樣的.
而撒瑪利亞.
這些人.
正正是.

$^{81}$他們的尋找方法.
是很正確的.
耶穌很欣賞的.
所以耶穌在撒瑪利亞.
他願意去面對自己的過犯.
他說.
耶穌記得將他.
所做的一切.
都說出來了.
於是耶穌親口.
跟他承認.
他就是彌賽亞.
耶穌親自跟他說.
相比之下.
耶路撒冷.
那些宗教人士.
他們不是真正.
去尋求明白.
律法的教誨是甚麼呢.
他們不肯面對.
自己的過犯.
他們不要這位指證他們惡行的.
彌賽亞.
加利利人又怎樣呢.
他們要的是甚麼呢.
就是.
彌賽亞聽他們的話.
行他們想要的神蹟.
他們要一個這樣的彌賽亞.
只有撒瑪利亞人.
他們是尋求一個.
真正的彌賽亞.
他們正期待甚麼呢.
這些人期待.
政治的.
一個軍事領袖的彌賽亞.
好像有一個王.
以色列的王.
來請教他們.
幫他們脫離當時羅馬的暴政.

$^{121}$但他們根本沒有期待.
一個宗教的人物.
或者是祭司.
或者是先知.
他們從來沒有想過.
但耶穌是具備了這三種身份.
軍王,先知,祭司.
他是與眾不同的彌賽亞.
與眾不同的受告者.
雖然其他先知.
如以利沙.
或者是祭司亞倫.
甚至是皇帝.
大家都知道.
有索羅王,大衛王.
他們都是.
避沒高尤.
而高立的.
但他們不是救主.
大家很清楚.
而主耶穌所建立的國度.
並不是.
世人所想像的.
要基督.
他所建立的國度.
都不是現時.
在地上發生的.
他不是要建立一個成功的士工.
他不是要打造.
一個完美的教會.
而是他為了要使我們.
脫離罪惡.
而我們可以真正的愛主.
真正的愛上帝.
這是耶穌基督.
來的目的.
大家想想.
當時的門徒.
他們都不明白.
他們問.

$^{161}$耶穌在天國裡.
誰是最大的呢?.
耶穌說:凡是自己謙卑.
好像小孩子.
他在天國裡是最大的.
亞哥和約翰.
他們要求神叫他們媽媽出面.
跟耶穌說.
一個坐左邊,一個坐右邊.
做什麼呢?.
門徒聽到都不高興.
因為耶穌知道.
他們在爭誰是最大的.
但到耶穌基督.
受死.
復活.
之後向他們顯現.
這班門徒.
慢慢的明白.
神的角度是怎樣的.
外邦的角度.
主要是權勢.
有權勢.
有地位的人.
自然受到人們的敬重.
但在教會裡.
不只是敬重.
有權勢.
有地位.
而謙卑的事奉神.
是應該得到敬重的.
所以我們必須謙卑.
不為詭詐.
而且願意面對.
自己的罪惡.
這種態度.
我們去接待我們的主.
我們去迎接我們的主.
我們要跟從祂.
我們還要學孝.

$^{201}$耶穌基督的寫記.
所以一開始.
我們看到.
先知約翰.
他有疑問.
他疑問什麼呢?.
在二節.
他派門徒去問耶穌.
將要來的.
那位是不是你?.
還是我還要等候另一個呢?.
他是一個問號.
但他曾經.
為主做了美好的見證.
記不記得.
有人跟他說.
你為他做見證的那個.
他現在跟別人私下.
很多人跟了他.
他說他必興旺.
必須美.
他講了這樣的見證.
或者他說.
我先前不認識他.
可是那差我來.
用水施洗的對我說.
你看見聖靈降下來.
停留在誰的身上.
誰就是用聖靈.
施洗的.
我看見了.
所以造證這位是神的兒子.
第二天他就見到耶穌.
他就說看了.
我為耶穌的羔羊.
除去世人罪孽.
他很清楚的.
為耶穌做見證.
但他這時候.
為何忽然間有一個問號呢?.

$^{241}$為何忽然間會有疑問呢?.
李索雅是不是你的?.
為何有這樣的疑問呢?.
因為.
在《以賽亞書》.
61章的一到二節.
那裡.
主耶華的靈在我身上.
因為耶華用羔羊告我.
叫我報好訊息給貧窮的人.
差遣我.
醫好傷心的人.
報告被擄得釋放.
被捆綁得自由.
宣告耶華的恩憐.
和我們的神報仇的日子.
安慰所有悲哀的人.
原來.
這段經文.
令他想到什麼呢?.
被捆綁得自由.
我們被囚的.
就應該出監牢.
他很期望.
耶穌可以救他.
為什麼耶穌沒有救他呢?.
他開始疑惑.
他疑惑.
他是耶穌的先鋒.
他當然期望耶穌救他.
他忠於真理.
他很勇敢的去斥責那些人犯罪.
他含冤入獄.
在苦難裡.
坐牢.
他很期望耶穌救他.
耶穌既然行過神蹟.
為什麼還沒看到他行動呢?.
一個受辱的君王.
為什麼連一個.

$^{281}$做他先鋒的.
他都看不到呢?.
他都不救呢?.
我經過這麼長的時間.
我這麼痛苦.
他很困惑.
各位大兄姐們,你們有沒有試過?.
有沒有試過有些疑問?.
遇到困難的時候.
每件事都不順利.
但別人卻有升職.
我這麼用心教會.
事奉.
但我呢.
可能要被逼遷.
可能又失業了.
怎麼辦呢?.
不信的人.
好像順風順水.
更加看不開.
我有沒有信錯呢?.
我是不是跟錯了人呢?.
跟錯了對象呢?.
就好像買股票.
到底我有沒有.
投資錯誤呢?.
溫慧耀博士.
大家很熟悉的.
他發現很多次.
當他人生遭遇重大.
打擊的時候.
痛苦的時候.
他哀求神.
求神幫助他脫離苦難.
但有時神並不是按照.
他的期望去幫助他.
在漫長的忍耐.
甚至開始質疑的時候.
他發現.
很多次.

$^{321}$雖然面對.
很重大的困難.
很痛苦的時候.
竟然.
看不到出路.
看不到曙光.
看不到有改變.
但.
看他身邊.
同時又發生一些很小的事情.
很輕微的事.
譬如說.
他帶太太去醫院.
去見醫生.
但他乘車.
但這時候正在下雨.
每輛的士.
都載了客人.
心裡很焦急的時候.
突然有輛車停在他旁邊.
有人下車.
他很順利地上車.
這件事情.
他很感恩.
神就在這些小小的.
細微的生活細節裡.
去幫助他.
總是出奇意外地順利.
這些小事.
他說.
你可能覺得.
只是巧合.
但他相信神.
讓他想不到的.
神給了他方便.
而.
有些人就說.
這些只是小節.
小事情.
但我需要的.

$^{361}$不是這些小事情.
我需要的是神幫助我.
去解決我.
又大又困難.
又痛苦的事情.
他回想.
神為什麼特別在.
這些小事上.
去經驗到.
神的手.
還一直存在.
神給他經驗.
因為神知道.
那件痛苦的事.
會讓他慢慢地.
失去耐性.
他會開始對神的真實.
對神的能力.
對神的慈愛.
開始動搖.
失去信心.
就在這個時候.
這些小動作.
就在提醒他.
神在提醒他.
這些重大的事.
雖然影響很大.
但神說.
他不可以按照他的心意.
按照他的想法.
去幫他排除困難.
但是.
神卻要定義.
告訴他.
他定義要在一些你看見的.
一些你生活上.
你見到的.
一些你經驗到的.
神要介入.
神要介入在其中.

$^{401}$讓他明白.
神仍然是存在的.
神是真實的.
他仍然有能力.
可以給他.
所以.
這些小事.
繼續幫助他.
持續對神.
有不懷疑的信心.
讓他.
很多次在幽谷裡.
提醒他.
神在他身邊.
所以縱然是在一個黑暗的隧道裡.
他仍然可以靠著主.
可以走出這個黑洞.
所以.
他在當中的答案.
就是雖然.
神沒有回答他.
但是神卻在他生活裡.
大小的事情.
他透過這些細微的事情.
他就發現.
神在其中.
他不可以.
不依靠他.
他不可以不感恩.
好.
耶穌怎麼回答呢?.
當這些人將約翰的疑問.
來問他的時候.
耶穌說什麼呢?.
很簡單.
他將第四十五節.
跟他們說.
你們去吧.
將所聽見所看見的.
告訴約翰.

$^{441}$就算盲人看見.
瘸子行走.
麻風病人得意症.
聾子聽見.
死人復活.
窮人聽到福音.
凡不因我跌倒的.
有福了.
但他就.
很簡單的.
將以賽亞.
先知提過關於.
彌賽雅的寓言.
他所行的事.
就是叫他去告訴.
事實約翰.
耶穌為什麼不去建立.
約翰想要的.
彌賽雅的角度呢?.
可能包括釋放.
所有像他這樣遭遇的人.
放他監獄.
當時四周.
很多人在聽.
耶穌的答案.
所以他不會公開.
我就是彌賽雅.
你看那個.
撒瑪利亞婦人多幸福.
耶穌馬上就告訴你答案.
我就是.
但這個時候他沒有說.
你知不知道.
原來你看聖經路加福音13章.
原來有人要殺耶穌.
就是希律.
希律要殺耶穌.
所以有時他不能說得那麼清楚.
說得太清楚會落入.
不必要那麼危險.

$^{481}$所以彌賽雅的角度.
他又不是要正面.
和希律作對.
他還有事.
要去做.
他還有使命.
所以他繼續在不同的層面.
去醫治人.
但最重要就是.
凡不因我跌倒的有福.
主耶穌的意思是什麼呢.
雖然我聽別人的祈禱.
就沒有聽別人的祈禱.
我幫助別人.
就沒有幫助你.
但你千萬不要因此跌倒.
你仍然要一心的去信靠我.
那你就有真福了.
如果我們.
如果你.
沒有因著耶穌基督.
不可以滿足你的祈禱.
那你就有真福了.
所以我們.
不可以滿足你的期望.
而你沮喪.
而你失望.
如果你沒有這樣的話.
你是有福的.
主所愛的人.
有時都會遭遇患難.
雖然主是全能的.
但他暫時.
都不夠他所愛的人.
去脫離苦難.
正如《歸林多後書》.
四章十六到十八節.
那裡說到什麼呢?.
保羅說:所以我們不喪膽.
雖然我們外在的人.

$^{521}$日漸扭慰.
內在的人.
卻日日更新.
我們這短暫而輕微的苦楚.
要為我們成就極重.
無比永遠的榮耀.
因為我們不是顧念看得見的.
而是顧念看不見的.
原來看得見的.
就是暫時的.
看不見才是永久.
是永遠的.
原來.
保羅為主受了很多痛苦.
他體力都消耗了很多.
很大的耗損.
但他仍然看到.
他心靈或主的作功.
很強壯.
所以他體會到.
其實遠靠我們自己.
根本做不到事.
是神賜下恩典.
讓我們完成事工.
所以他說.
現時所受的苦楚.
只是短暫.
如果我們眼光.
放在永恆的裡面.
是多麼好.
不扭慰的.
屬天的居所是多麼好.
今年的苦楚.
這麼短暫.
跟永恆來比算得了什麼呢?.
所以我們千萬不要.
因為我們目前的事.
而對神質疑.
疑惑.
主今天愛不愛我呢?.

$^{561}$施一若不會因主.
沒有幫助他.
沒有施援手.
而跌倒.
事實上.
他被斬頭.
後來門徒.
告訴耶穌.
他信道了.
溫偉堯博士.
有一個智障女兒.
26歲離世.
他說如果過去的人生.
沒有經過這樣的苦楚.
今天的他.
不是這樣的.
他說.
正如一個運動員.
都要經驗一個無情的教練.
對他的要求.
你多麼疲乏.
撐不下去.
他說神偏偏.
揀選女兒.
進入他生命裡.
讓他學會了什麼?.
學會了在他.
讀博士學位裡.
讀不到的學問.
是多麼寶貴.
第三方面.
主稱讚施洗若寒.
在七到十節裡.
就是.
他說.
他問人們.
特別問.
跟從若寒的人.
你們從前.
去曠野看什麼?.

$^{601}$是否看風吹動的蘆葦?.
你們出去.
到底要看什麼?.
是否穿細暖衣服的人?.
穿細暖衣服的人.
在皇宮裡.
你們出去.
究竟要看什麼?.
我告訴你們.
他比先知大多了.
死的人.
就是經上所說的.
這裡說到.
施洗若寒.
比所有先知都大.
從前是指什麼?.
從前是指耶穌基督還沒出現.
施洗若寒.
他還沒坐牢.
叫從前.
蘆葦是生長在曠野近水的地方.
相信.
這群人.
他們去曠野不是要看.
穿細暖衣服的人.
這些人.
不用去那裡看.
在皇宮就能看到.
看什麼?.
當然不會去曠野.
就是為了看蘆葦.
原來.
在.
《馬頭江》第三章.
五到六節.
那裡說到.
原來.
那時耶路撒冷全猶太和.
全約旦河地區的人.
都到若寒那裡去.

$^{641}$承認他們的罪.
在約旦河裡受他們的洗.
原來他們去那裡.
為了什麼?.
原來就是要聽施洗若寒.
傳悔改的訊息.
他們去到面前.
因為施洗若寒.
要講悔改的訊息.
他們去就是為了.
這件事.
所以他們去曠野不是要看.
那些權貴階級高的人.
他們穿華麗的衣服.
不是為了看這些.
也不是為了看景.
也不是.
因為這個人是.
製作神的僕人.
神差的僕人去講訊息.
他是被差的.
耶穌告訴人.
他是很受歡迎的.
他受歡迎不是因為他的地位.
有什麼特殊.
他沒有什麼特別的.
而是神藉著他去.
傳人心靈.
所需要的訊息.
這是重點.
所以.
接著他就說.
他比先知.
更大.
他比先知更大.
就是比猶太人心目中的先知.
比他們大多.
因為怎樣呢.
因為在.
馬達福音剛才所讀的.

$^{681}$我要差已經在我使者.
在你面前.
他要在你前面為你預備道路.
其實這個經文.
是引用.
馬拉基書的三章一節.
證明到神.
透過以賽亞先知.
去講施洗約翰.
就是耶穌基督的先鋒.
透過以賽亞先知去講.
所有的先知.
都是預言.
彌賽亞要來臨.
但是施洗約翰.
他是先知預言.
要來的那位先知.
你說特別嗎.
大小之分.
是講福分.
另一方面.
眾先知都是預言.
耶穌基督.
但沒有人見過耶穌基督.
但是.
他們這些先知.
在預言的時候.
只是用信心.
遙望耶穌基督.
彌賽亞.
但是這個先鋒.
施洗約翰.
他是親自.
見到耶穌基督.
他為他施洗.
大家看到他多有福分.
如果不是他.
主的先鋒.
他看到聖靈降臨在耶穌身上.
多大的福分.

$^{721}$他親自去見證.
耶穌基督是神.
應許要背負.
世人的罪.
所以除了.
約翰和眾先知比較之外.
接著他又說.
女子所生.
經文再下去.
女子所生指什麼呢.
指肉身.
每一個都是由母親生的.
父人所生的.
天國女生的.
就是神生的.
天國女指的是神生的.
每一個人都是肉身生的時候.
這些人.
在當時很多人見到耶穌基督.
但是沒有人肯信他.
就滅亡了.
但天國裡面.
從神生的.
這些願意相信耶穌基督的.
任何一個.
由大到小.
所以他說.
最小的都比約翰大.
就是這個意思.
因為當時約翰信心軟弱.
他在苦難裡面.
他見不到神的拯救.
他看不到.
他更加不明白.
這個理所當然.
耶穌的角度.
是怎樣的呢.
他看不到.
連門徒也是.
門徒也是看到耶穌復活之後.

$^{761}$沒有什麼特別.
所以難怪.
由神生的角度裡面.
最小的都比.
施洗約翰大.
就是這個意思.
天國裡面.
最小的我們明白了.
天國是要.
努力進入的.
天國努力進入.
難道我要走很多善事.
是不是這個意思呢.
不是叫你做好事.
很多機會.
籌款.
外面很多人罵你.
是不是做這麼多就可以進入天國.
不是這個意思.
努力是表示.
全音的福音有很多阻力.
所以你要努力.
阻力令你不能前行.
阻礙你.
你要去講.
別人就潑你冷水.
所以你要很努力.
才可以進入天國.
施洗約翰.
他有以利亞心智.
他很勇敢的.
去斥責罪惡.
當時都說那些.
法利賽人.
他們是毒蛇的種類.
他很大膽的去斥責這些人犯罪.
他心裡頭的意念.
他很大膽的講.
今天我們每一個信主.
你去傳福音.

$^{801}$不是這麼簡單.
主雖然沒有.
去拯救約翰.
他被殺害了.
但他不是因為他有過失.
不是做錯什麼.
而失去聖命.
但這個世代的人.
不聽他講.
棄絕他.
但主耶穌.
卻是推舉他的工作.
推舉他所做的使命.
對他有正面的評價.
他成為.
凡不因他跌倒的人的榜樣.
這是一個.
給我們學習的榜樣.
因為這個世代的人.
為什麼要這麼努力呢?.
因為這個世代的人.
漠視神的公義.
和慈愛.
繼續下去的經文.
就被我們看到.
就說我要.
做什麼來看這個世代呢?.
這個世代.
形容孩童在街市上.
呼喊.
我們為你們吹笛.
你們不跳舞.
我們唱哀歌.
你們不垂胸.
什麼意思呢?.
吹笛,跳舞.
本來很開心的.
吹笛有喜事.
有些人還可以請樂隊.
來喜事.

$^{841}$祝賀.
這是真的.
當人有吹笛的時候.
你不是應該跳舞.
很開心跟著節奏嗎?.
表示你不重視那個場合.
有人離世了.
聖經也教我們.
與哀嚎的人同哭.
為什麼.
哀傷的時候.
你毫無表情.
為什麼你好像看不到.
為什麼.
他形容.
古時候.
那些孩童.
沒有什麼娛樂.
現在的人這麼幸運.
有個遊戲機.
沒那麼幸運.
玩什麼.
玩喪禮.
扮婚事.
很好玩.
學你吹笛.
學你一曲曲,兩曲曲.
在街上.
看到他們這樣玩.
他們玩.
就當他們玩.
不會跟著悲哀.
不會傻的.
跟著哭.
他們玩.
我就這樣走.
不當一回事.
重點是.
主要感嘆.
人們漠視.

$^{881}$他所說的話.
公義,慈愛.
人們漠視.
他宣告大喜的訊息.
或者審判來臨.
他宣告這麼嚴重的事.
人們漠視.
漠視到.
好像小孩玩遊戲.
我當他們玩.
根本不重視.
神.
在宣告這麼嚴重的事.
你要警醒.
不然就要面對審判.
但他當玩.
根本不重視.
約翰.
他穿什麼?.
駱駝毛的衣服.
吃蝗蟲野蜜.
怪怪的.
所以人們說.
他不吃不喝.
跟一般人不同.
說他像鬼婦.
耶穌基督.
他比較親民.
親近.
他比較憐憫.
他又吃又喝.
人們說他貪吃好酒.
人們.
很反常.
你吃他又說你.
不吃又說你.
他根本不接受.
施西約翰的圖.
又不接受耶穌的圖.
根本不接受他們.

$^{921}$所以這個世代.
很難服侍.
何況今天.
你要去傳福音.
有很大的難度.
我們需要很好的祈禱.
去堅持.
有一個很年輕的.
女裁縫師.
她.
很辛苦的去讀書.
她很希望.
可以考到一個時尚的設計.
將來.
可以發大財.
安樂美.
她就由佛羅里達.
來註冊.
註冊表示.
她考取了.
註冊了.
但院長不歡迎她.
院長.
說書我直言.
我不知道你是客人.
但這個時候.
這個小女孩很努力的讀書.
她就是想考上這個學院.
但這個.
院長竟然.
這樣說.
但她很低聲的跟神說.
神啊我不想離開.
求你幫助我.
她在祈禱.
但院長知道.
她不想走.
那怎麼辦呢.
她這麼堅持.
於是院長就跟她.

$^{961}$講條件.
她說你不可以跟白人.
一起上課.
你只可以打開課室門.
在旁聽.
於是.
她就這樣上課.
但她才華橫溢.
她被人提前.
六個月畢業.
她很用心讀書.
雖然她在門口聽到.
但她很留意.
去學習.
後來.
她受到.
上流社會顧客的青睞.
你一定知道.
總統.
第一夫人賈桂林.
港大夫人.
很欣賞她的設計.
她結婚的禮服.
這套婚紗是她做的.
她做了兩次.
第一次.
因為工作室.
上面水管破裂.
我看到這個.
想起我們.
禮拜五.
錦書院.
很多水流下來.
事情不難也難.
她做得這麼漂亮的婚紗.
但她沒有放棄.
要交貨.
她很努力的.
終於交了差.
她沒有放棄的精神.

$^{1001}$她就是尋求上帝幫助她.
去倚靠神.
各位弟兄姊妹.
向神禱告.
倚靠神是很重要的.
她就是這樣成功的.
所以耶穌難怪.
她做的比喻.
在路加福音說.
有一個不義的官.
和一個寡婦.
死撕著她.
很煩她.
就是這麼煩的字眼.
跟她申冤.
她經常冤著我.
耶穌要我們這麼努力.
用心的.
向神祈禱.
不要放棄.
這很重要.
因為今天人們漠視神的公義.
神的話.
當耳邊風.
甚至信徒也是.
因為我們不看聖經.
走寶.
很多神嚴厲說的話我們看不到.
於是我們犯罪.
所以我們真的要去祈禱.
我們懇求神幫助我們.
當我們願意懇求神的時候.
神就會按著.
祂的心意.
在不同的時間回應.
祈禱.
今天為了什麼事要拼命去祈禱呢?.
繼續.
神會幫助你.
會扶持你.

$^{1041}$各位弟兄姊妹.
你願不願意.
跟隨這位.
與你期望完全不同的.
截然不同的彌創下.
你願不願意跟從他呢?.
他會繼續.
對你有要求.
今天是這樣.
日後他又會更新你的看法.
不是一成不變.
你還願不願意跟從他呢?.
我們一起祈禱.
愛你的主,愛你的神.
主,我們何等不配.
你竟然連累我們.
讓我們成為你的兒女.
今天.
你就是我們的彌賽雅.
你教導我們.
要離棄罪惡.
要愛你,要愛上帝.
求你幫助我們.
不要忘記我們.
起初愛你的心.
讓我們繼續延續下去.
讓我們繼續.
為人去禱告.
為教會去禱告.
為自己去禱告.
以致我們所行的.
能夠行在你的旨意當中.
求主教導我們.
保守我們.
若然我們偏離了.
求聖靈肯當當鞭策我們.
以致我們所行.
所言的.
都能夠合乎你的旨意.
我們這樣祈禱交託.

$^{1081}$是奉主耶穌基督的命求.
\newpage



\section{哥林多後書 8:1-9}
\label{sec:BzeYjKv9RX8}
\textbf{2023.09.17 《活出一個樂捐慷慨的生命》 哥林多後書8\_1-9 (和修版) 曾敏傳道}
\newline
\newline
連結: \href{https://youtube.com/watch?v=BzeYjKv9RX8}{\texttt{https://youtube.com/watch?v=BzeYjKv9RX8}} ~~~~ 語音日期: 2023-09-18
\newline
\newline
\hyperref[sec:Wzp5u5hQ940]{\small{< < < PREV SERMON < < <}}
~
\hyperref[sec:index_chronic]{\small{[返順時目]}}
~
\hyperref[sec:index_scriptual]{\small{[返順卷目]}}
~
\hyperref[sec:1E5Bibwk4hc]{\small{> > > NEXT SERMON > > >}}
\newline
\newline
哥林多後書 8:1-9
\newline
\begin{longtable}{cl}
\hline
\hline
章節 & 經文 (和合本修訂版)\\
\hline
8:1 & \begin{tabularx}{0.7\textwidth}{X} 弟兄們，我們要把神賜給馬其頓眾教會的恩惠告訴你們： \end{tabularx} \\ \\ \relax
8:2 & \begin{tabularx}{0.7\textwidth}{X} 他們在患難中受大考驗的時候，仍然滿有喜樂，在極度貧窮中還格外顯出他們樂捐的慷慨。 \end{tabularx} \\ \\ \relax
8:3 & \begin{tabularx}{0.7\textwidth}{X} 我可以證明，他們是按著能力，而且超過了能力來捐助，主動 \end{tabularx} \\ \\ \relax
8:4 & \begin{tabularx}{0.7\textwidth}{X} 再三懇求我們，准他們在這供給聖徒的善事上有份； \end{tabularx} \\ \\ \relax
8:5 & \begin{tabularx}{0.7\textwidth}{X} 並且他們所做的，不但照我們所期望的，更照神的旨意先把自己獻給主，又給了我們。 \end{tabularx} \\ \\ \relax
8:6 & \begin{tabularx}{0.7\textwidth}{X} 因此，我們勸提多，既然在你們中間開始這慈善的事，就當把它辦成。 \end{tabularx} \\ \\ \relax
8:7 & \begin{tabularx}{0.7\textwidth}{X} 既然你們在信心、口才、知識、萬分的熱忱，以及我們對你們的愛心上，都勝人一等，那麼，當在這慈善的事上也要勝人一等。 \end{tabularx} \\ \\ \relax
8:8 & \begin{tabularx}{0.7\textwidth}{X} 我說這話，並不是命令你們，而是藉著別人的熱忱來考驗你們愛心的真誠。 \end{tabularx} \\ \\ \relax
8:9 & \begin{tabularx}{0.7\textwidth}{X} 你們知道我們主耶穌基督的恩典：他本是富足，卻為你們成了貧窮，好使你們因他的貧窮而成為富足。 \end{tabularx} \\ \\
[1ex]
\hline
\hline
\end{longtable}
$^{1}$各位姐妹平安.
今日這堂道是講奉獻的一堂道.
雖然平時我都在講道當中提及奉獻.
但並非如今日所講.
我們講奉獻的生命究竟是怎樣的呢?.
開始時我們一同禱告.
賜下你的話語.
要成為我們腳前的燈路上的光.
今日求你對我們眾人說話.
願意聖靈在我們心裡工作.
讓我們智慧悟性去明白.
我們的心真是聽到上帝你的聲音.
就去回應.
求主賜福以下的時間.
奉耶穌基督的名字祈禱.
阿們.
活出一個樂捐慷慨的生命.
我深信是上帝對我們的心意.
我們一路去看.
經文在哥林多後書八章一到九節.
哥林多這個地方究竟是怎樣呢?.
我們稍稍的認識.
原來哥林多城在希臘半島.
那裡人口眾多.
其實有點像香港.
是商業非常繁榮.
有很多發展的機會.
但是我很想這樣說.
那裡的道德是非常淪落.
道德是很差的.
亦因為這個緣故對教會都有不少影響.
所以你去看的時候.
哥林多前書裡面都有提到教會的問題.
這個地方是否很像香港呢?.
香港是很繁榮.
有很多發展的機會.
但同樣香港的罪惡一樣有很多.
對教會都有一定的影響.
我們看哥林多在哪裡呢?.
在這個圖大概左下角的位置.

$^{41}$在第二次保羅宣教之旅.
保羅在這個地方建立了教會.
一共用了一年半的時間.
時間也不短.
保羅和弟兄姊妹有很多的機會相處.
他在栽培信徒裡也傳了很多福音.
保羅和哥林多教會有不同的關係上的轉變.
在這裡,特別這一段.
應該說八章和九章都是講奉獻的.
兩章保羅都是在奉獻方面作出一些教導.
我們看看這段經文是在說甚麼呢?.
他們給我們看到一個榜樣.
是一個很樂觀很慷慨的生命榜樣.
簡單來說是一個很樂於奉獻的生命.
一個很樂於奉獻的生命究竟是怎樣的?.
為甚麼會很樂於奉獻呢?.
是否因為有人邀請?.
是否有甚麼的催逼?.
不奉獻是不可以的?.
究竟是甚麼緣故?.
我們看看這個榜樣是怎樣的.
原來是在說馬其頓宗教會的一個榜樣.
究竟這些教會是怎樣的呢?.
我們看看教會的狀況.
教會的狀況呢.
教會的狀況經文裡面是看到.
八章二節這樣說.
他們在患難中受大考驗的時候.
這個患難中受大考驗是甚麼呢?.
是指通常是因為信仰的緣故而面對迫迫.
所以那裡的信徒的生活是非常不容易.
不是每個星期日去到教會裡面聚會.
是很舒服地毫無難阻的.
沒有人去迫迫他們的.
很自由地敬拜讚美.
不是這樣的情況.
他們在信仰上面面對迫迫.
而且是大考驗.
但是教會在這樣的狀況裡面.
你見到信徒的生命是甚麼呢?.

$^{81}$滿有喜樂.
沒有被打低.
其實在這樣的情況裡.
都可以有兩種不同的表現.
面對患難 面對大考驗.
可以有甚麼不同的表現呢?.
其中一種是甚麼呢?.
離開信仰.
回教會要面對迫迫.
信耶穌要這樣.
有人拿著槍指著我.
還不離開信仰.
不要受這樣的苦.
何苦給自己一些困難.
不會的 走開吧.
今守氣 走開吧.
洩氣 灰心 埋怨 罵上帝都有.
其中一種反應是這樣的.
另一種是這種.
馬其頓宗教會是怎樣的呢?.
滿有喜樂.
不埋怨 神不灰心 不喪氣.
喜樂 是滿滿的.
這個生命是很不同的.
是很有力的生命.
接著 馬其頓宗教會的信徒是怎樣的呢?.
是否很有錢的? 中產的?.
駕駛好車 拿名牌袋的?.
穿好衣服的? 出入的?.
真的不一樣.
是否這樣? 不是的.
極度貧窮的 貧窮都不夠.
是要極度.
我們試過極度貧窮嗎?.
試過 說來看看程度有多大.
什麼叫極度?.
沒有糧食 出不了糧.
這樣叫極度.
有沒有飯吃的?.
去山灘洗衣服.

$^{121}$不是水喉水.
要用水的 但也有水的.
還有飯吃的.
什麼叫極度? 沒有飯吃的.
或者吃的東西 縮小的.
餓的.
如果對我來說 極度貧窮.
比較像非洲那些.
大家有沒有看過網上的照片.
我未曾去過那些這麼窮的國家.
如果看網上的照片.
你看到那些小朋友.
躺在地上 瘦到像棚骨一樣.
手腳像牙籤一樣.
每天都不知道能否活著.
在他身上有些蒼蠅在摟著他.
沒有水喝.
沒有衣服穿 沒有食物.
我想馬其頓宗教會可能還沒有去到極端.
但極度貧窮.
還沒有去到極端 但很窮.
如果你說奉獻.
通常人們會怎麼想.
你去找些有錢教會.
找些中產教會去講這道道.
你不要對著極度貧窮的教會講這道道.
有錢才能給得起.
這樣才像樣.
拿出來都差不多一塊.
今天保羅這裡的教導是什麼.
這個榜樣是什麼.
馬其頓宗教會是極度貧窮.
但在這樣的情況下 教會是怎樣.
格外顯出樂捐的慷慨.
是格外的.
是很願意捐獻的 很慷慨.
你想像不到.
極度貧窮是不可能的.
不是 就是極度貧窮.
還要額外樂意捐獻 慷慨.

$^{161}$剛才說的教會的狀況是這樣.
很窮 面對迫不及待 在患難中.
但教會是喜樂的.
是願意去奉獻的.
而且你看教會有行動.
我講奉獻不是說.
我現在奉獻給你.
是說這些話.
有行動的 行動是什麼.
他們真的會奉獻出來.
奉獻出來的時候是按著能力.
這裡說的是有的.
但要加上一句是什麼.
捐獻超過了能力.
我很窮 我的袋子只有兩千銀幣.
好像窮寡婦.
很窮也好.
她也是盡力.
其實這個窮寡婦應該是超過了.
窮寡婦的兩千銀幣.
是她用來養生的.
全部拿出來.
教會現在很窮.
她也是超過能力拿出來奉獻.
這樣也回答.
是呀 真是有教會這樣做.
你記住她正在面對迫不及待.
有內憂有外患.
自己的經濟有很大的缺乏.
外面面對迫不及待.
這個信仰生活一點也不容易.
簡直是艱難.
我和你試過這樣的光景嗎.
我想問.
頂枝木我想問你.
今天有沒有人拿著槍指著你.
不要回紅陰堂聚會.
你要回我就抓你坐牢.
去到沒有.
我們有沒有這樣試過.

$^{201}$現在暫時來說我們香港很平安.
沒有人這樣逼迫我們.
所以我們沒有面對她的難處.
我們是極度貧窮.
我講極度.
今時今日有沒有人還在山坑裡.
拿水洗米洗衣服.
或許今天我們停了工作.
但仍然有錢用.
仍然有飯吃.
仍然有衣服穿.
今天大家還回來教會崇拜.
我們還沒有去到.
馬其頓宗教會的這種光景.
但你看到馬其頓宗教會的生命不一樣.
他的行動是甚麼.
他有安靜能力而且超過能力.
還有教會是怎樣奉獻的.
再三去求.
去求使徒在這個.
供給聖徒的善事上有份.
其實當時出現了一個情況.
在耶路撒冷教會的信徒.
在經濟上有很大的缺乏.
在生活上有很大的艱難.
所以保羅在這裡.
反而是要幫助.
耶路撒冷教會的信徒.
去籌募捐款.
籌募奉獻.
讓這些外邦的信徒.
他們可以幫助.
幫助耶路撒冷的信徒.
在這個過程中.
你看到馬其頓宗教會.
沒有人去叫他們去捐獻.
沒有人逼他們.
沒有人請他們.
是他們自己求的.
你看到他們這麼窮.

$^{241}$你都不好意思去找他們.
去奉獻去捐獻.
但他們自己求.
不止一次.
是再三的求.
給我奉獻吧.
這麼窮還要求.
我要奉獻了.
我要供給那些聖徒.
他們有缺乏.
換言之.
這一群的信徒.
馬其頓宗教會的信徒.
是怎樣的信徒.
他們看什麼更加重要.
是看自己的需要嗎.
不是.
他們看到耶路撒冷教會的信徒.
看到其他信徒有需要.
其他人的需要更加重要.
我雖然窮.
但我都想出一分力.
這個生命是不一樣的生命.
是危險的.
不止如此.
保羅在這裡形容他們的奉獻.
他們的行動是怎樣的.
這群信徒去奉獻的時候.
是怎樣的奉獻.
他們是按照神的旨意去奉獻.
按照神的旨意是什麼.
原來神希望我們先將自己獻給主.
主很體重我們這個人.
主最體重的是我和你的生命.
最希望我們將整個人獻給祂.
羅馬書十二章一到二節是怎樣說的.
所以弟兄們.
我以神的慈悲勸你們.
將身體獻上.
當作活祭.

$^{281}$這是神所喜悅的.
將自己的生命獻上.
其實這個奉獻是一個更大的要求.
上帝的旨意是整個人要求我們奉獻給他.
很多時候我們說奉獻.
通常都會說.
這些你只對傳道人說好了.
未來傳道人說.
他將自己的生命獻給上帝.
他整個生命屬於神的.
我就是那一小撮.
我的錢包一小撮.
我的時間就是一小撮.
我什麼都是一小撮的.
我計算好那些就給上帝.
其他就是我的.
弟兄姊妹.
我想問在你生命裡.
有哪一樣是純純單單的.
純純屬於你自己的呢?.
有沒有哪一樣純純屬於你自己的.
不屬於上帝的?.
你可以在位置舉手告訴我.
沒有.
其實我們生命的一切都屬於神的.
神給我們生命.
今天才可以做這樣做那樣.
神給我們手腳能動.
眼耳口鼻能動.
才可以享受不同的東西.
還要神給我們有很多的供應.
我才可以享受到.
今天神隨便收走一樣都沒有.
是不是在過程中.
所有東西都是屬於神的?.
神今天並不缺乏.
因為所有東西都是祂的.
我們的神最富有.
我們的神也最豐富.
是祂創造世界的.

$^{321}$祂是這個世界的創造者.
祂擁有這個世界.
所有東西都是祂的.
生命的源頭在祂那裡.
今天是祂慷慨給我們.
今天祂要我們整個人是什麼?.
要我們懂得感恩.
認定祂是我們的主.
認定祂是整個世界的主宰.
創造主.
祂是我們的王.
所以要我們將整個人獻給祂.
上帝最看重的是我這個人.
祂不需要所有東西.
其實祂也不需要.
但祂看重我們這個人.
要我們去感恩.
要去敬拜祂.
尊崇祂.
王就是主.
其實這個奉獻最基本.
最原初的意義就在這裡.
這是神的旨意.
整個獻上.
所以我記得有一段時間.
香港也在推行一件事.
說基督徒是一個全人的奉獻.
不是說傳道人,神學生.
宣教士是全人奉獻.
基督徒是一半奉獻.
三分一,四分一奉獻.
不是這樣的.
其實我們每個人都是全人奉獻.
將自己先獻給主.
然後又給我們.
這群也是聽使徒的吩咐.
也會跟從使徒的吩咐.
一群這樣的人.
奉獻原來是這樣的.
而馬喬頓宗教會是這樣奉獻.

$^{361}$按神的心意去奉獻.
一個毫無保留,全然的.
將自己擺上的奉獻.
這個榜樣是不是很美呢?.
在這裡.
我們想回到八章一節.
是這樣說.
為什麼馬喬頓宗教會.
可以這樣奉獻呢?.
被逼白金窮.
都可以超過能力奉獻.
可以這麼好呢?.
很重要一句.
八章一節這樣說.
弟兄們.
我們要把神賜給馬喬頓宗教會的恩惠.
告訴你們.
原來能夠樂觀,能夠慷慨.
能夠這樣奉獻.
完全是上帝的恩典.
上帝就是這樣給了我們很豐厚的恩典.
我們才可以這樣奉獻.
能夠奉獻.
沒有什麼值得自誇.
反而見證神的恩典.
是神的恩典這樣促成.
在這裡.
保羅藉著這麼寶貴的.
這麼美好的榜樣.
對哥林多教會作出勸勉.
這裡八章八節這樣說.
「我說者說,並不是命令你們」.
我不是逼你們.
不是強勢地命令你們做.
但是保羅在這裡勸勉教會.
勸勉什麼?.
既然哥林多教會有這麼多的恩賜.
如果大家有機會去看哥林多前書.
其實教會有很多恩賜.
神給了他們很多.

$^{401}$在這裡八章七節都有說.
既然你們在信心,口才,知識.
萬分熱忱.
以及我們對你們的愛心上.
這句可以解為.
你們對我們的愛心上.
都聖人一等.
他們不單有恩賜這麼簡單.
簡直是很豐富.
豐富到比別人還多.
是聖人一等.
既然如此.
是金夢神的恩典的一群人.
更加應該在慈善的事上.
也要聖人一等.
神給他們這麼多恩賜服侍.
神給他們這麼多慈福教會.
是這麼豐富.
他們在愛人的事上.
在回應上帝的恩典上.
在奉獻上.
是否更加應該要下力呢?.
這是保羅所說的.
他們的奉獻又如何?.
現在他們奉獻的這筆錢.
是支持耶路撒冷教會.
那些很貧窮.
很有需要的信徒.
哥林多教會這群信徒.
有什麼回應呢?.
當時哥林多教會.
基本上是這樣的.
他們的聚會.
很多時候都在弟兄姊妹的家.
你要想想.
要容納這麼多人.
弟兄姊妹的家一定很大.
要大.
弟兄姊妹是很有米的.
要有一定的經濟條件.

$^{441}$你才可以有大屋.
可以容納這麼多人.
在你家裡聚會.
可以款待弟兄姊妹.
所以你看到.
哥林多教會.
反而不是極度貧窮的教會.
在能力上.
一定比馬其頓宗教會都要好.
但是在這裡.
他們反而不會主動去奉獻.
是保羅一聲說.
既然你們有這麼多恩賜神.
這麼愛你們.
你們是不是應該要多一點.
起來奉獻呢?.
保羅要考驗他們的愛心.
是不是真誠的呢?.
真心的愛人.
真心的愛上帝.
愛不是只用嘴巴說的.
弟兄姊妹問自己一個問題.
你怎樣去愛人呢?.
是不是每天都跟他說.
我好愛你.
有子女的.
你問問自己.
為子女擺上多少.
從他出生的時候.
你為他預備了多少東西.
他一路成長.
你為他擺上多少.
經濟上都數之不盡.
在學校上為他做了多少鋪排.
學校上的事.
還要想著.
他在哪裡受教育更好呢?.
移民啊.
很多事情要想.
現在這一代更加.

$^{481}$到結婚都不會離開父母.
沒有房子沒有車.
跟父母一起.
父母還要養孫子.
一路上一路下.
一生.
算錢都很難算.
因為是愛.
你不會覺得付出.
我真的很心痛.
由你出生到現在.
我花了多少錢.
很心痛.
你不會心痛.
因為你很愛他們.
不要說子女.
配偶也是.
每年.
每月每日.
為你配偶做多少事.
金錢上付出多少.
是愛就會付錢.
是愛就會付出.
你很捨得.
有很多人都很愛自己.
很捨得.
花了很多錢.
不是這樣那樣.
逐樣來.
保羅說要考驗他們的愛心.
是否真誠.
是否真愛.
愛那些有需要的弟兄姊妹.
愛那些有需要的教會.
愛上帝.
愛上帝是空口說的.
天天都這樣跟上帝說.
上帝很愛你.
是嗎?.
馬其頓宗教會和光林多教會.

$^{521}$有很大分別.
剛才有說過.
馬其頓宗教會又面對逼迫.
極度貧窮.
但主動哀求.
試圖要被他們奉獻.
哀求完一次.
又一次.
超過能力.
自己能否吃飯.
一回事.
我有沒有錢生活.
一回事.
我要去支持那些有需要的.
我要奉獻給教會.
可能這些我們很多時候說.
他們很窮.
不要動他們.
他們會主動.
反過來.
光林多教會很豐富.
神給他們的恩典很豐富.
有米的.
就不會主動奉獻.
是嗎?.
要保羅說到出聲.
保羅說不會命令他們.
因為這些要心甘情願.
弟兄姊妹.
我們問自己的問題.
我們又如何?.
我們在哪個教會?.
馬其頓宗教會?.
在光林多教會?.
在這裡.
聖經有個很寶貴的一句說話.
回應八章一節.
八章一節.
我要將神賜給馬其頓宗教會的恩惠.
告訴你們.

$^{561}$神如何恩待馬其頓宗教會?.
沒錯.
他們有很多限制和困難.
你知道神有多恩待他們?.
還可以奉獻.
弟兄姊妹.
我想說.
能夠奉獻原來是恩典.
今天記住.
能夠奉獻不是肉痛.
在錢包裡挖錢出來.
不是.
是恩典.
你能夠奉獻是恩典.
神在你生命裡有奇妙的工作.
是恩典.
為什麼我們可以奉獻.
八章九節的回應.
你知道主耶穌基督的恩典.
這個恩典是怎樣來的?.
祂本來是父族.
耶穌基督本來是父族.
祂享有全世界.
本來屬於祂的榮耀.
尊貴.
所有的父族.
是屬於耶穌基督.
我告訴你.
祂為你們成了貧窮.
祂甘心.
道成肉身.
來到這個世界.
和人一起生活.
祂甘心成了貧窮.
祂在一個木匠的家庭裡.
出生和成長.
祂出生之後.
你記得客店也沒有地方.
要睡馬槽.
祂很明白貧窮的滋味.

$^{601}$祂也和貧窮人一起.
在世33年的時間.
祂只要為人擺上自己.
為我和你最後成就救恩.
祂為我們釘在十字架.
又為我們復活.
祂犧牲捨己.
為我和你成就救恩.
因為祂這份救恩.
我們和天父才能復和.
才能離開永世的咒咒.
你說,這樣就有福了嗎?.
當然有福了.
多大的祝福啊.
萬有在主裡.
弟兄姊妹.
我剛才說.
神創造天地.
創造萬有.
其實萬有在主裡.
今天和神復和.
和統管萬有.
擁有萬有的主復和.
我們何等的豐富啊.
是何等的富足啊.
這麼大的福氣.
耶穌付出這麼大的代價.
讓我和你去得著.
好使你們因他的貧窮而成為富足.
在自己甘願貧窮.
好讓我和你成為富足.
我和你怎會沒有能力奉獻呢?.
我想問.
主耶穌基督給我們的.
是何等的豐厚啊.
這麼大的恩典.
我和你一定會有能力奉獻.
能夠奉獻是恩典.
沒有什麼值得自誇.
而且有很好的榜樣.

$^{641}$讓我和你去學校.
好.
主耶穌基督的恩典這麼寶貴.
他對我們的愛這麼深.
今天問大家的問題.
是因為他的恩典和愛.
我們才能奉獻.
頂枕外我們的回應是什麼呢?.
我們愛神嗎?.
你問自己.
我愛上帝嗎?.
你說愛.
那我們會放上什麼呢?.
我們會放上什麼呢?.
要問自己.
是否盡了能力.
我在這裡.
挑戰大家.
如果能夠像學校一樣.
超過能力.
我深信上帝的祝福會加倍的臨到.
但我相信上帝也不會這樣強求.
上帝只會教我們按著能力.
就按著你今天的能力.
你去到最盡.
可以去到多少.
為上帝擺上.
你說愛.
今天問大家.
愛不愛洪恩嘉?.
我們一定會為我們所愛.
放上金錢.
說真的.
我們會放什麼呢?.
我們愛洪恩嘉嗎?.
我們願意為洪恩嘉.
放上什麼呢?.
現在是一個禱告的時間.
領子妹.
問問自己.

$^{681}$這麼久以來.
為主放上了什麼呢?.
我們說我們愛主.
我們有沒有看到.
教會的需要.
有沒有看到.
洪水橋區的福音和長.
何等的大.
需要發展.
很需要我們多方面的支持.
經濟上絕對是其中一件事.
我們願意為教會.
放上什麼呢?.
我不需要和別人比較.
今天不是被迫的.
而是我們自己回應神.
和神說.
我們想放上多少.
願意放上多少.
是不是真的很愛主?.
在一分鐘之後.
我帶大家有一個結束的禱告.
在一分鐘之後.
我帶大家有一個結束的禱告.
在一分鐘之後.
我帶大家有一個結束的禱告.
(音量注意).
給我們很豐厚的恩典.
你是何等的愛我們每一個.
你把你的愛子耶穌基督賜給我們.
因為你愛我們.
你甘願來到這個世界.
經歷貧窮.
你捨棄一切.
為我們成為富足.
我們何等的不配.
主耶穌你配得我們去愛你.
配得我們將自己獻在你的手裡.
其實萬有本來就在你裡面.
萬有本來就在你裡面.

$^{721}$是屬於你的.
給人這麼幸福.
去享受上帝你所預備的一切.
主耶穌基督求你.
給我們.
亦都是毫無保留的奉獻給你.
毫無保留的去愛你.
將一切交在你的手裡.
讓主你來掌管.
讓主你帶領我們.
將這一切都交託給你去掌管.
主在此求你紀念我們眾弟兄姊妹.
過去的日子.
個人的奉獻如何.
主求你深深的感動我們每一個.
為你擺上.
我們能力所可以的.
盡力擺到最好.
為要回應主你的愛.
你的恩典.
要讓主你得榮耀.
要讓主你的心得到滿足.
主啊.
其實你可以不將這一切都給我們.
但你卻是愛我們.
你視我們如珠如寶.
視我們為珍貴.
所以你從來都很慷慨的和我們分享.
主求你讓我們認定這一切都來自你.
讓我們感恩.
讓我們甘心樂意擺上.
願你恩代洪恩嘉.
讓我們成為樂捐慷慨的教會.
為主你的國度.
為神你的家.
為有需要的弟兄姊妹.
擺上自己.
讓我們知道當我們願意奉獻的時候.
我們心裡會感到歡喜快樂.
因為主你的心意就是想我們這樣擺上.

$^{761}$求你帶領我們.
幫助我們.
奉耶穌基督的名字祈禱.
阿們.
\newpage



\section{列王記下 5:1-5}
\label{sec:1E5Bibwk4hc}
\textbf{2023.09.24 《神使用小人物》 列王記下5\_1-5 (和修版) 陳一華牧師}
\newline
\newline
連結: \href{https://youtube.com/watch?v=1E5Bibwk4hc}{\texttt{https://youtube.com/watch?v=1E5Bibwk4hc}} ~~~~ 語音日期: 2023-09-25
\newline
\newline
\hyperref[sec:BzeYjKv9RX8]{\small{< < < PREV SERMON < < <}}
~
\hyperref[sec:index_chronic]{\small{[返順時目]}}
~
\hyperref[sec:index_scriptual]{\small{[返順卷目]}}
~
\hyperref[sec:ILOp7icU0bo]{\small{> > > NEXT SERMON > > >}}
\newline
\newline
列王記下 5:1-5
\newline
\begin{longtable}{cl}
\hline
\hline
章節 & 經文 (和合本修訂版)\\
\hline
5:1 & \begin{tabularx}{0.7\textwidth}{X} 亞蘭王的元帥乃縵在他主人面前是一個偉大的人，得王的喜悅，因為耶和華曾藉他使亞蘭人得勝。他雖然是大能的勇士，卻染上了痲瘋。 \end{tabularx} \\ \\ \relax
5:2 & \begin{tabularx}{0.7\textwidth}{X} 亞蘭人成群出征的時候，從以色列地擄了一個小女孩，她就服事乃縵的妻子。 \end{tabularx} \\ \\ \relax
5:3 & \begin{tabularx}{0.7\textwidth}{X} 她對女主人說：「我希望主人去見撒瑪利亞的先知，他必能治好主人的痲瘋。」 \end{tabularx} \\ \\ \relax
5:4 & \begin{tabularx}{0.7\textwidth}{X} 乃縵去告訴他主人說，從以色列地來的女孩如此如此說。 \end{tabularx} \\ \\ \relax
5:5 & \begin{tabularx}{0.7\textwidth}{X} 亞蘭王說：「你可以去，我也會送信給以色列王。」於是乃縵手裡帶十他連得銀子、六千舍客勒金子和十套衣裳去了。 \end{tabularx} \\ \\
[1ex]
\hline
\hline
\end{longtable}
$^{1}$洪恩堂的弟兄姊妹 早安.
今早我們一起敬拜.
唱了很多詩歌.
也聽了很多詩歌.
心裡面暢快嗎?.
給我們這麼多詩歌可以歌頌.
也為敬拜隊 為太羽小組.
我們聽到這麼多好的音樂.
這麼多好的詩歌.
張榮耀歸比天父.
一首詩歌《人算什麼》.
掀起了我們今天要聽到的內容.
真是的 人是很微不足道的.
人是很微小的.
算不得什麼.
但你不要忘記.
原來上帝可以用微不足道的人.
甚至小人物.
上帝用不用?.
上帝也可以使用.
所以我們今早讀一段經文.
是少許陌生.
可能有些人會有印象也不一定.
預備起.
「阿南王的元帥乃萬.
在他主人面前是一個偉大的人.
得王的喜悅.
因為耶和華曾直著他使亞蘭人得勝.
他雖然是大能的勇士.
卻染上了麻風.
亞蘭人成群出征的時候.
從以色列地擄了一個小女孩.
她就服侍乃萬的妻子.
她對主人說.
我希望主人去見撒瑪利亞的先知.
她必能治好主人的麻風.
乃萬去告訴他主人說.
從以色列地來的女孩.
如此如此說.
阿南王說:你可以去.

$^{41}$我也會送信給以色列王.
於是乃萬十他年得銀子.
六千寫赫勒金紙.
和十套衣裳去了」.
今天我們看看這個小人物.
他在當時亞蘭國.
大家聽過亞蘭國沒有?.
這個國家.
如果你打開地圖.
它剛剛是在以色列國的北面.
這個國家是很能打的.
為什麼?.
它臣時某時入侵以色列.
經常造成以色列很大的煩惱.
因為這個國家有個將軍.
是很英明神武.
戰績輝煌.
讀了他的名字.
是兩個字叫做乃萬.
因為他這麼能打.
所以有一次他出去打仗的時候.
打贏了.
不但如此.
還擄了一個小女子回去.
大家有這個印象嗎?.
這個小女子為什麼是一個小人物呢?.
因為兩個原因.
第一.
聖經有沒有說她叫什麼名字?.
沒有.
對我們來說.
她無名無姓.
什麼都不知道她是誰.
她真的是一個小人物.
不知道她是哪裡來的.
也不知道她叫什麼名字.
第二個原因她是一個小人物.
是因為她被人擄走.
有沒有留意到?.
換句話來說.

$^{81}$她是很淒涼的.
打敗仗後被人擄走.
很難過的.
說得不好聽.
就是亡國奴.
這個小女子.
如果說得好聽一點.
就叫做妹仔.
叫做婢女.
說得不好聽.
就是奴隸.
大家明白這個意思嗎?.
因為她沒有身份.
沒有主權.
什麼都沒有.
就要去到外國.
去做公民.
去做服侍.
這是一個小小的背景.
一個小人物.
但是不要忘記.
雖然她是一個這樣的人.
上帝用不用她?.
上帝可以用她.
我們今年四月的時候.
大家有沒有印象.
記得參加一個很大的筵席.
去到沙田的.
有沒有印象?.
是我們的錦宣家.
洪恩.
有一個齊躍洞的.
有沒有印象?.
當中有一個老人家.
很大年紀的.
八十多歲的.
寫了一首很好的書法.
字畫.
我講講.
你記得了.

$^{121}$這個老人家是誰?.
是保叔.
沒錯.
大家不要忘記.
這個保叔原來是一個死囚.
應該死兩次.
因為他兩次有謀殺的記錄在當中.
但是他在監獄裡信了耶穌.
悔改.
整個人都改變了.
又在監獄裡練了一首好字.
也讀聖經.
勤讀聖經.
整個人改變了.
又有好行為.
後來監獄就假釋他.
一次兩次.
最後1994年的時候.
就放了他出來.
從此就可以重獲自由.
不過他也在赤柱坐了二十年監.
等於是沒有過.
一個這樣的人物.
這樣的老人家.
上帝用不用他?.
用得著他.
他可以透過抄寫聖經的金句.
來鼓勵人.
所以他去了很多的教會.
社區中心.
學校去講見證.
去分享神怎樣改變他的人生.
這個老人家.
雖然他有這樣的死囚背景.
但是神要用他的時候.
用得著他嗎?.
用得著他.
各位,真的很奇妙.
今次上帝用這個小女孩.
也不是沒有原因的.

$^{161}$因為這個小女孩的身上.
有四項的優點.
這四項的優點.
令我們明白.
上帝為何會用她.
不如我們試試看.
她有什麼優點.
第一個優點.
就是她有信心.
她跟她的主母說.
我們的主人三麻風.
在亞蘭閣見那麼多的醫生.
見那麼多的大夫.
醫治得怎麼樣?.
醫治不好.
我有一個提議.
是什麼呢?.
如果我的主人.
肯去撒瑪利亞.
即是他的鄉下.
肯去他的國家.
肯去撒瑪利亞見先知.
這樣就行了.
我們講到.
平時有一句話.
若到病除.
這次不是.
這次是人到病除.
你明白我的意思嗎?.
只要她肯去見先知.
她就有機會康復.
你說這個小女孩的信心大不大?.
非常大.
她對自己國家的先知.
有充分的信心.
相信若她肯去見先知.
謙謙卑卑去哀求先知.
先知若肯醫她.
一定能醫到.
這個信心真是一個很大的信心.

$^{201}$是否值得我們去學校?.
不單止她有信心.
你發覺她很有勇氣.
是否很有勇氣?.
因為她說得出這番話.
有沒有代價?.
冒不冒險?.
冒險的.
因為如果真是成事.
Hallelujah!.
回來後又升職.
又找個好男人給她嫁.
什麼都好.
但你不要忘記.
如果不成事.
那就大件事了.
大家明白我的意思嗎?.
即是她去騙一個大將軍.
去到撒瑪利亞.
最後破個胸.
什麼都醫不了.
回來後這個小女孩.
不知道誰說得對.
人頭落地.
大家明白我的意思嗎?.
冒不冒險?.
有冒險的.
不過這個小女孩.
因為她有信心.
所以她不怕冒險.
她都要說出來.
讓主人知道.
這是第二個優點.
第三個優點是什麼呢?.
就是她有個關心.
關心.
我們有一句話說.
事不關己.
己不勞心.
現在沙麻風是誰生的?.

$^{241}$她.
主人沙麻風.
又不是我沙麻風.
大家明白我的意思嗎?.
她可能做得壞事多.
如果想得不好.
她當然做得壞事多.
有沒有報應?.
有.
所以她可能殺了不少人.
沙麻風.
大家明白我的意思嗎?.
所以她可以想很多負面的事.
可以覺得不去關心她.
是天經地義的.
這個就是這樣的情況.
不過這個小女孩不是這樣.
她雖然做丫鬟.
做婢女.
她有沒有關心?.
她有關心.
她樂意去關心.
所以她將這個建議.
告訴主母.
這是她第三個優點.
第四個優點.
大家看到了.
有沒有愛心?.
有愛心.
為什麼呢?.
因為她想想.
這個將軍正是她的敵人.
大家有沒有留意到?.
如果不是她.
我又沒有今天.
如果不是她.
我用不著被人擄來這裡.
用不著家庭拆散.
用不著無家可歸.
大家明白這個意思嗎?.

$^{281}$如果你想起這些事的時候.
你的心裡面.
很自然地有仇恨.
其實一有恨就很笨.
很多壞事都會做出來.
很可怕.
所以不要小看人的心裡的那種恨.
那種仇恨會被魔鬼利用.
做很多壞事.
這個小女孩心裡.
發恨為愛.
她放下了這些恩恩怨怨.
你明白我的意思嗎?.
將軍怎樣對我不好都好.
今天我要以德報怨.
所以各位頂智媒媒.
她有沒有愛心?.
有愛心.
看完之後.
你現在開始明白.
上帝為什麼要用這個小女孩.
有沒有原因?.
有.
今天頂智媒媒.
如果你和我都學校這個小女孩.
擁有這四項特質的時候.
上帝用不用我們?.
上帝用我們.
所以我們今天.
真的在這個小人物的身上.
看到上帝很奇妙.
連這個小女孩都沒有想過.
她簡簡單單說了一句話.
你明白我的意思嗎?.
原來發生的事情.
就是很神奇.
想不到她一句簡單的說話.
令到整件事情.
是非常暢順.
而且是逢關過關.

$^{321}$你有沒有留意到?.
第一關是什麼?.
第一關.
大家有沒有留意到這個小女孩.
她老了做小女孩.
實際上她有很多功夫可以做.
你明白我的意思嗎?.
她可以被人派去做很多.
很多其他的打雜的工作.
又奇妙.
偏偏安排她去服侍誰?.
服侍將軍的太太.
這個真的是什麼?.
做夢都什麼?.
沒有想過.
弟兄姊妹你有沒有留意到?.
所以有時候我和一些紀律部隊.
弟兄姊妹聊天.
我說你千萬不要小看自己的身份.
雖然你今天做一個.
不是很高職級的人.
但是你留意.
上帝安排了這個崗位給你.
上帝安排了這個職份給你.
你不要小看.
分分鐘你可以接觸什麼?.
高層人士.
你明白我的意思嗎?.
分分鐘你是高層管理人員的.
左右手都未定.
大家明白我的意思嗎?.
所以各位弟兄姊妹們.
真的很奇妙.
有時候上帝安排你.
做某些的位置.
做某些的工作.
原來上帝都有計劃.
可惜時間緣故.
我不可以講這麼多例子給你聽.
否則我真的見過.

$^{361}$有些弟兄姊妹.
一個普通的職位.
但竟然是最高職位.
最高領導人的左右手.
這是什麼?.
上帝很奇妙的一個安排.
一個帶領.
還有.
為什麼奇妙呢?.
不單是這個小妹服侍乃萬的太太.
而他和乃萬的太太.
講這一番話.
這個主母的反應.
正面還是負面?.
正面.
你說奇不奇妙?.
因為這個主母一聽完她的意見.
她分分鐘可以不接納.
你明白我的意思嗎?.
分分鐘都可以罵她.
都未定.
為什麼?.
一來她的建議又怪怪的.
她叫她的主人.
不是去見撒瑪利亞的名醫.
不是去見撒瑪利亞的大夫.
你明白我的意思嗎?.
不是說撒瑪利亞有什麼出名的醫生.
顧問醫生沒見過.
她不是去見這些.
去見誰?.
去見仙智.
你說她的建議怪怪的.
是嗎?.
你都要建議什麼?.
介紹名醫給我認識.
如果這個主母聽完之後.
覺得她的提議.
是真的荒謬.
你明白嗎?.

$^{401}$她一句就可以.
埋她尾巴.
你很空閒嗎?.
沒事做嗎?.
你明白我的意思嗎?.
你沒事做不如進柴房吧.
你明白我的意思嗎?.
一句就去埋她尾巴.
多管閒事都可以.
你明白我的意思嗎?.
不是.
這個主母聽完之後.
覺得雖然怪怪的.
但未嘗不可嘗試一下.
一個機會.
這個反應差很遠.
第二關過了.
第三關.
主母聽完之後.
就把這個建議告訴她的先生.
告訴她的丈夫.
而她的丈夫聽完之後.
可否否決她?.
可以的.
為什麼?.
剛才都說了.
她這個提議怪怪的.
如果這個將軍覺得自己是大丈夫.
是大男人.
有自己的見解.
你們這些女人說這些.
這麼怪怪的.
全部都是苦如知見.
她一句都可以.
封了她.
一句埋尾巴.
但這個大將軍聽完之後.
竟然心裡有點不安.
你明白我的意思嗎?.
這個提議怪怪的.

$^{441}$但我從來未聽過.
會否是一個可能的方法呢?.
各位電視節目.
結果這個大將軍又接納這個建議.
這是不是第三關?.
好了.
現在第四關了.
見誰?.
見皇帝了.
你明白我的意思嗎?.
因為這個大將軍要去撒瑪利亞.
是否需要皇帝批准呢?.
要的.
見皇帝.
跟皇帝說.
有件這樣的事情.
有個女孩這樣這樣這樣說.
你明白我的意思嗎?.
你都知道.
皇帝想事情.
又是怪怪的.
大家有沒有留意到.
皇帝想事情.
跟你一般人想事情.
是完全可以.
截然不同的.
所以為何我們中國人有一句說話.
叫做伴君如伴虎.
皇帝想的事情.
霎時間又是甚麼呢?.
想些怪事出來.
你明白我的意思嗎?.
現在你跟皇帝說這樣的事情.
這樣的怪意見.
可能是父父得正都未定.
大家都怪.
結果皇帝接不接納呢?.
皇帝接納.
當然皇帝接納.
因為皇帝是深愛大將.

$^{481}$皇帝喜歡將軍.
因為他打得幫到國家.
對他有好感.
為了令將軍開心.
好了好了.
你說甚麼我都甚麼.
去吧去吧去吧.
你不單止去.
我還要加多封信.
你明白嗎?.
我加多封信.
寫給以色列王.
告訴他.
我的大將軍來到你那裡.
你就要甚麼?.
治好他.
這次大不大劑?.
真的大劑.
如果你再讀下去.
真的獵王技巧.
令你發覺到.
那些劇情.
比你看電視劇更精彩.
各位電視節目們.
原來到最後皇帝都甚麼?.
都接納.
都批准.
我想告訴大家.
這次整件事.
能夠這麼暢順.
逢關過關.
是不是這個小女子.
很厲害.
很能幹.
是不是?.
不是啊.
她只是一句說話.
你明白我的意思嗎?.
上帝容不容許?.
上帝容許.

$^{521}$簡單一句說話.
上帝是甚麼?.
容許.
又令到她甚麼?.
逢關過關.
所以各位電視節目們.
你要記得.
一個大將軍信耶穌.
帶來的影響大不大?.
很大.
這個大將軍乃萬.
結果被先知治好之後.
他怎麼說?.
他說我從今之後.
只知道天下間.
只有一位甚麼?.
真神.
就是你所敬拜的耶和華.
我就不再相信其他的神明.
各位電視節目們.
一個大將軍篤信耶和華.
他的影響力大不大?.
很大的影響力.
所以我相信他回到亞蘭閣.
他也會帶其他人.
認識耶和華.
歸向耶和華.
這個結果影響大不大?.
很大.
但是有第二個影響.
結果更加大.
你知不知道是甚麼?.
第二個結果.
第二個影響更加大.
是因為從此之後.
大將軍被治好.
兩個國家還打不打仗?.
還打不打仗?.
不打仗了.
休戰.

$^{561}$這個休戰對百姓來說.
是否一個喜訊?.
是否一個很大的祝福?.
電視節目們.
大將軍被治好之後.
我相信亞蘭王或者大將軍.
都要給意識別過面.
不再侵略他們.
所以各位電視節目.
你發覺這個小女孩.
是否可以做一個大故事出來?.
今天牧師想跟大家說的是.
你和我雖然微不足道.
人算甚麼.
但請你不要妄自菲薄.
為甚麼?.
只要你是神的子女.
你擁有小女孩這四項特質.
你願意聽從神的話.
就在你現在的工作崗位上.
在你的家庭裡面.
在你的教會裡面.
在你的社區裡面.
在你的社交圈子裡面.
你猜上帝用得著你嗎?.
上帝用得著你.
各位電視節目.
我們能夠見到一些很大的結論.
見到一些很美好的結果出來的時候.
不是因為我們有甚麼了不起.
是因為甚麼?.
因為上帝用得著我們.
當神用我們的時候.
整件事就改變了.
當神用我們的時候.
沒有可能的事就變成有可能.
所以各位洪恩家的弟兄姊妹.
請你記得.
我們追求被神使用.
因為神可以用得著我們.

$^{601}$在任何地方神都可以使用我們.
但我們最重要的是甚麼?.
聽神的話.
做神喜悅的事情.
有像小女孩這樣的特質.
神就能夠用我們.
我們低頭祈禱.
親愛的主,多謝你.
今天早上讓我們從一個小女孩的身上.
看到人算甚麼?.
微不足道.
一個亡國奴.
能夠被神大大的使用.
祝福多人認識上帝.
也祝福兩個國家能夠休戰.
這些重大的影響.
都全因著上帝使用一個小女孩.
你介入其中.
顯出偉大的作為.
所以求上帝都同樣使用我們.
叫我們所到的地方.
都成為天賦祝福的一個出口.
成為你自己一個流通的管子.
榮神益人.
奉耶穌名求.
阿們.
感恩.
\newpage



\section{馬太福音 6:9-15}
\label{sec:ILOp7icU0bo}
\textbf{2023.10.01 《渴望與追求》 馬太福音6\_9-15 (和修版) 黎光偉牧師}
\newline
\newline
連結: \href{https://youtube.com/watch?v=ILOp7icU0bo}{\texttt{https://youtube.com/watch?v=ILOp7icU0bo}} ~~~~ 語音日期: 2023-10-03
\newline
\newline
\hyperref[sec:1E5Bibwk4hc]{\small{< < < PREV SERMON < < <}}
~
\hyperref[sec:index_chronic]{\small{[返順時目]}}
~
\hyperref[sec:index_scriptual]{\small{[返順卷目]}}
~
\hyperref[sec:514XvjAlzD8]{\small{> > > NEXT SERMON > > >}}
\newline
\newline
馬太福音 6:9-15
\newline
\begin{longtable}{cl}
\hline
\hline
章節 & 經文 (和合本修訂版)\\
\hline
6:9 & \begin{tabularx}{0.7\textwidth}{X} 「所以，你們要這樣禱告：『我們在天上的父：願人都尊你的名為聖。 \end{tabularx} \\ \\ \relax
6:10 & \begin{tabularx}{0.7\textwidth}{X} 願你的國降臨；願你的旨意行在地上，如同行在天上。 \end{tabularx} \\ \\ \relax
6:11 & \begin{tabularx}{0.7\textwidth}{X} 我們日用的飲食，今日賜給我們。 \end{tabularx} \\ \\ \relax
6:12 & \begin{tabularx}{0.7\textwidth}{X} 免我們的債，如同我們免了人的債。 \end{tabularx} \\ \\ \relax
6:13 & \begin{tabularx}{0.7\textwidth}{X} 不叫我們陷入試探；救我們脫離那惡者。因為國度、權柄、榮耀，全是你的，直到永遠。阿們！』 \end{tabularx} \\ \\ \relax
6:14 & \begin{tabularx}{0.7\textwidth}{X} 你們若饒恕人的過犯，你們的天父也必饒恕你們； \end{tabularx} \\ \\ \relax
6:15 & \begin{tabularx}{0.7\textwidth}{X} 你們若不饒恕人，你們的天父也必不饒恕你們的過犯。」 \end{tabularx} \\ \\
[1ex]
\hline
\hline
\end{longtable}
$^{1}$各位導演姐妹早安.
終於可以不用錄影了.
我們可以實體見面.
很開心見到每一位的牧者.
很開心見到幾位牧師.
有梁傳道,有Sarah傳道.
也有很多很熟悉的面孔.
有些我們一起服侍過.
有些我們曾經在祈禱區很長時間.
我們一起祈禱,一起分享.
很記得當紅印堂剛剛建立的時候.
我還在錦川服侍.
當時我兒子才兩歲.
真的才兩歲.
女兒還沒有出生.
半年之後我就舉家回了澳門.
當時要承接澳門教會目的的工作.
轉眼間原來.
回看下一個月.
就應該是紅印堂十二週年了.
原來闊別了大家有十二年的時間.
不知道在座有多少人.
你是頭一天就已經在紅印家.
你已經在這裡了.
或者甚至你在這裡信主.
有沒有在這裡信耶穌?.
有啊.
我相信過去十二年來.
對大家來說.
有很多上帝的恩典.
對我來說也是一樣.
十二年的時間.
同樣有很多的體會.
有很多的學習.
其中一個學習.
是經歷子女的成長和轉變.
我相信在座有些人可能已經有第三代了.
是啊.
我相信你也很體會.
子女下一代.

$^{41}$他們轉眼間就有很多的不同.
我記得.
當他們剛剛升上小學的時候.
他們可以說是一張白紙.
他們對爸爸媽媽非常之依賴.
老師是他們一切答案的準則.
他們在教會學習也是一樣的.
他們喜歡聽聖經的故事.
會積極地回應主耶穌.
回應他們的祈禱.
回應主耶穌對他們的教導.
他們會回應主耶學老師的問題.
很努力地背誦.
很單純地告訴我們.
跟我們分享禱告的經歷.
不過.
不經不覺.
當他們進入了青少年的階段之後.
似乎他們自己也發現.
原來社會上.
跟隨耶穌的人並不是大多數.
面對大部分.
都沒有教會生活的老師和同學.
他們或許也會不禁地問自己.
為什麼我要花這麼多時間.
在主日的聚會上.
為什麼我要回教會呢?.
雖然爸爸媽媽從小就帶我回來.
但是長大後可能有很多的疑問.
儘管他們過去在教會.
有豐富的群體相處經驗.
但是為了融入同學的圈子.
他們需要建立另一套相處方式.
他們發覺有些事情是不可行的.
他們要自己去找出路.
而作為父母.
我們也明白當子女長大.
並不是我們這樣說.
他們就是這樣聽的.
事實上.

$^{81}$他們再不滿足在兒童主日學.
回應正確的問題.
他們會問.
到底我們這樣去遵從上帝的話.
是不是真的可能的呢?.
是不是真的做得到呢?.
對我的生活有什麼益處呢?.
我相信.
他們的成長.
跟我們的信仰歷程也很相似.
當我們剛信耶穌的時候.
我們同樣單純.
我們積極尋求信仰.
我們剛才知道有些人是.
十二年前在這裡信耶穌.
十二年前.
可能我們也是.
帶著這種單純.
也是帶著這種積極.
不過.
十二年過後.
或許我們已經熟習了教會的生活.
或許我們已經能夠.
說了很多基督徒的術語.
甚至我們參與過很多不同的事奉.
不過我們也慢慢地體會到.
要貫徹進行聖經的教導.
其實是很難的.
是不是?.
真的要按照聖經的教導去走.
其實不是那麼容易的.
我舉個例子.
先不要說要建立每天讀經祈禱的習慣.
要學習愛淪舍如同自己.
哇!.
你想想你旁邊淪舍.
你就知道有多難了.
不僅要受人的服侍.
乃至要服侍人.
耶穌說的.

$^{121}$這是YMCA的金句.
很難的.
要愛你的受敵.
為了迫害你們的人禱告.
哇!.
更加難.
這些都是在耶穌的教導裡面.
甚至我們要使不同的人成為主的門徒.
我想大部分的人.
跟我的想法都是一樣的.
能夠生活做到人畜無害.
努力做好自己的崗位.
其實已經不簡單了.
如果要進一步按照聖經的要求.
去做多一點的話.
到底可以具體怎麼做呢?.
而這樣遵行主.
是不是真的可以改變我的生活.
是不是帶來我的生活.
有一個很不一樣的幫助.
以及我可以享受到當中.
那種豐盛,那種快樂呢?.
原來在這個過程裡面.
你會發覺.
很多挫折.
不是那麼容易.
先生.
在耶穌的時代.
當時大部分的猶太人.
他們都是相信上帝的人.
當他們看到.
那些宗教領袖.
當時的法利賽人.
他們為了向人表達.
他們要過一個敬虔的生活.
他們高調地舉行一些施捨的活動.
他們在十字路口.
聖經這麼說.
大聲地來祈禱.
他們禁食也嫌不夠敬虔.

$^{161}$他們要裝扮一下.
蒙頭,臉皮地禁食的時候.
其他的猶太人.
那些信上帝的人.
他們也帶著相同的疑問.
這是不是我們遵行上帝的方式呢?.
當我信了主十二年.
甚至信了主幾十年.
當我要再做多一點的時候.
是不是就是這麼做呢?.
如果是這樣做的話.
究竟這些行動.
對我的信仰生活.
對我的人生.
有什麼幫助呢?.
這個還是只是門面功夫呢?.
還是只是做給別人看呢?.
事實上.
主耶穌並沒有反對施捨.
沒有反對禱告.
也沒有反對禁食.
這是第五章.
耶穌說的.
只是主耶穌認為.
如果信徒不能夠從心底裡認同.
履行信徒的使命和目標的話.
我們原來很容易會迷失的.
不知道大家有沒有這種感受.
有這種體會.
信主月內.
屬靈成長不是成正比的.
很多時候我們會滑了啞的.
我們有時候也會有些迷失的感覺.
好像我日復日.
周復周地回來.
究竟有些事情好像已經約定俗成.
究竟對我們有什麼幫助呢?.
我們很容易迷失.
甚至我們會擔心.
我們成為一個有名無實的信徒.

$^{201}$正正因為這樣.
主耶穌就要幫助當時的人.
也很希望今天跟大家分享這段經文.
去幫助今天的我們.
尤其是去幫助當時的門徒.
認清他們追求的方向.
以致他們用遵行主的生活態度.
來面對實際環境的挑戰.
但願我們今天都可以發現.
上帝的教導是能夠用在我們的生活上.
我知道紅人唐家.
一進來就看到我們有年提.
今年的年提是「學無真理」.
是「遵行主導」.
到底經過十二年教會的生活之後.
如何可以延續下去呢?.
如何可以決心.
正如年提所說.
我們實踐主導呢?.
今天很想跟大家分享.
剛才已經讀了我們很熟悉的經文.
不如我們再讀一次.
在螢光幕上也看得清楚.
我相信有人會背.
我們試試一起讀.
再讀一次.
預備起.
所以你們要這樣禱告.
我們在天上的父.
願人到尊你的名為聖.
願你的國降臨.
願你的子羽行在地上.
如同行在天上.
我們日用的飲食.
今日賜給我們.
免我們的債.
如同我們免了人的債.
不叫我們陷入試探.
救我們出來那惡者.
你們若饒恕人的過犯.

$^{241}$你們的天父也必饒恕你們.
你們若不要恕人.
你們的天父也必不要恕你們的過犯.
讓我們先同心禱告.
上帝我們求你親自帶領我們.
啟迪我們.
我們深信在過去.
當紅人家被建立的時候.
上帝不單建立一個屬靈的家.
提供一個地方讓我們聚會.
更加透過過去這些日子.
透過不同的牧者.
造就我們 塑造我們.
上帝我們求你親自帶領.
讓我們期待.
我們可以繼續將上帝的使命延續下去.
我們能夠知道如何去實踐你.
更加有勇氣去實踐你.
求你帶領以下的時間.
祈禱奉耶穌基督的名字.
阿們.
當我們看到主禱文的時候.
或許你會問這個問題.
如果你特別.
你可能聽過很多次主禱文.
看過很多次.
聽過很多不同的講解.
不知道你會不會問一個問題.
就是為什麼主耶穌要教他們祈禱呢?.
難道當時的門徒不懂祈禱?.
所以需要主去教他們?.
似乎不是很可能.
當時的猶太人.
不可能不懂祈禱.
他們從小就懂祈禱.
記得有一次去聖殿.
我們的導遊.
他是來自加利利的.
他就告訴我.
他不是一個很勁虔的猶太教徒.

$^{281}$不過他說他一出生.
他的家人就要教他讀經文.
要他去禱告.
五六歲的時候.
就要懂得背《摩西五經》.
每個猶太人都懂得祈禱.
為什麼耶穌要教他們祈禱呢?.
會不會是因為他們的祈禱不夠正宗.
所以耶穌要糾正他們呢?.
我相信在這裡.
主耶穌並不是要糾正他們.
祈禱不正.
而是特別針對當時的宗教領袖.
特別是一些法利賽派的猶太人.
他們對敬虔生活.
因為他們誤解.
所以主耶穌要重新讓他們的門徒.
去明白禱告的意義.
藉著這個禱告的教導去表明.
我們作為屬主的一群.
唯有我們重拾自己的身份.
重視自己的使命.
我們才可以被上帝改變.
我們才可以從裡到外.
去活出一個敬虔的生命.
這個敬虔的生命是無法偽裝的.
無法靠蒙頭救面.
在十字路口裡面大聲讀出一些熟練的經文.
可以模仿到這個敬虔.
必須要重拾自己的身份.
明白上帝在我們的心意.
我們才能做到.
那怎麼做呢?.
讓我們回到經文這裡.
事實上在這段圖文裡面.
我們大致可以分成兩個部分.
一個是渴望.
一個是追求.
就像大家的練題一樣.
讓我們逐一部分來看看.

$^{321}$第一個是我們保持對主的渴望和期望.
主耶穌一開始就要門徒稱上帝為.
「我們就是天上的父」.
說明門徒和上帝的關係.
既然祂是所有人天上的父.
所以求告祂的時候.
我們不需要像唸咒一樣說重複的話.
也不需要透過特別的儀式來向天父傾訴.
不過主耶穌卻同時要門徒明白.
這位天父是獨一的真神.
是唯一的宇宙主宰.
所以當我們要跟隨一個活在高天.
一個全能的父的時候.
我們就要保持對祂渴望的心.
我們要有這個心.
我們才可以有進一步的行動.
如果你都不渴望上帝.
你不羨慕天上的主的話.
你很難有下一步的行動.
就好像我們上班一樣.
什麼驅使我們可以戰勝每天的睡魔.
過兩天又要上班了.
可以讓我們每一天大清早就可以回到公司.
我們都需要有些渴望的.
無論你渴望發工資.
你渴望成就.
你渴望買房子.
渴望養家.
渴望貢獻社會.
渴望退休.
唯有有這種渴望.
祂才可以成為你的動力.
堅持到底.
而不同的渴望.
不同的期待.
就會帶來不同的態度.
去看待你要做的事情.
主耶穌在這裡提出.
對天父我們需要有三個渴望.
你很熟悉的.

$^{361}$「願人到尊你的名為聖.
願你的國降臨.
願你的子兒行在地上如同行在天上」.
首先我們看看這三個渴望.
都是你熟悉的經文.
都是關於你的.
都是你的.
你的命.
你的國.
你的子兒.
主耶穌告訴我們.
信徒的人生目標.
並不是一種自我實現.
我相信耶穌.
我想實現一下另類的生活.
不是.
主耶穌說.
這不是自我實現.
我們不是透過信耶穌.
實踐我理想的生活.
提升我個人的氣質.
建立我社會的地位.
當一個人.
立志跟隨耶穌的時候.
他就是要將上帝的關注.
化成我的關注.
上帝關注什麼.
我就關注什麼.
我舉一個聖經的例子.
大家很熟悉的.
是摩西的故事.
上帝的僕人摩西.
他是受委派.
帶領以色列人.
由埃及走出去.
最終要去到上帝應許之地.
這個故事.
我相信大家很熟悉.
不過.
你也知道.

$^{401}$當他們成功離開了埃及.
進入到曠野之後.
當時大部分的以色列人.
都經不起信心的考驗.
最終他們受到上帝的懲罰.
大部分人死在曠野當中.
當中.
摩西都是其中一個.
沒有辦法進入.
這個應許之地的人.
不過當摩西知道.
自己最終不能夠如願.
進入迦南美地的時候.
他怎麼辦呢?.
這才是我們值得去看的地方.
他沒有「劈炮」.
你想想.
如果是你.
上帝說.
我有一個任務交給你.
你會去到一個應許之地.
帶領著以色列人.
去到中間.
任務失敗.
上帝說.
下一代可以進入.
你不可以進入.
就算你不「劈炮」.
你應該也會隨便地做.
可能你會開始放鬆一下心情.
既然都進不了.
但是.
聖經並沒有說.
摩西放慢腳步.
根據聖經的記載.
他不單止沒有停下來.
他更加致力去栽培新一代.
整個申命記.
正正是摩西的講章.
申命記很沉氣的.

$^{441}$一次又一次.
一而再三地.
告訴下一代的人.
我們要怎樣積極地.
進入到這個應許之地裡面.
他要確保他們懷著信心.
勇敢地進入上帝賜給他們的地方.
建立使人得福的福音群體.
摩西之所以要這樣做.
是因為他一直以來.
他很清楚.
他致力實現的.
不是他個人的期望.
而是上帝的期望.
所以他並沒有因為畢生的事業.
遇到這麼大的挫敗而崩潰.
人生因為這個污點而一蹋不振.
我們在信仰的路途也會遇到很大的挫折.
甚至我們發覺.
有些失敗是很難回頭的.
很多遺憾.
可能我們在信仰上也會一蹋不振.
但是當摩西將焦點放在上帝.
他要實現的是上帝的理想的時候.
他就能夠走過這些最大的困難.
結果這位以色列歷史上.
可以說是無可比擬的領袖.
曾經是這樣形容的.
上帝以他是朋友.
不因為他的失敗而貶低了他的位置.
而他對上帝的忠誠.
反而成為他人生的標記.
成就了一生的傳奇.
今天我們稱呼摩西為神人摩西.
不是因為他特別聖潔.
他也曾經失敗.
不過他走過.
他以上帝的心意成為他畢生的目標.
弟兄姊妹.
我們是否在日常生活中.

$^{481}$仍然可以像摩西一樣.
將上帝的關注作為我們生活的優先.
我相信這是主耶穌期待我們有的一種態度.
一種渴慕的態度.
我們就要問下去.
上帝關注什麼呢?.
既然我們覺得上帝的關注這麼重要.
主耶穌就告訴我們.
上帝是天國的元首.
古代作為一個元首.
名號,國土.
他傳達的聖旨是最為重要的.
所以門徒作為天國的子民.
就應該以渴慕他的名能夠尊為聖.
他的國土可以得以擴張.
他的旨意可以得以承傳為念.
我們就要問.
人什麼時候可以尊上帝為聖呢?.
主耶穌告訴我們.
當人聽聞福音的見證.
當人接受福音的時候.
他們就能夠尊上帝為聖.
所以當我們去傳福音.
當我們使人認識上帝的時候.
我們就是在做這個行動.
人什麼時候可以見到天國呢?.
就是當信徒走在一起.
當我們彼此相愛的時候.
主耶穌說.
我們如果有彼此相愛的心.
人就認出我們是他的門徒.
所以當我們走在一起.
當我們彼此相助.
我們就能夠實現讓人見到天國.
人什麼時候可以見到上帝旨意成就呢?.
就是當門徒著實去遵行上帝旨意的時候.
人就見到上帝的旨意臨到人間.
頂尖的牧師一雙.
我們很多時候都會將信仰的焦點.
放在信耶穌這個點上.

$^{521}$沒錯的.
對基督徒來說.
信耶穌是生死的關鍵.
不過主耶穌更看重的是.
門徒信主之後的生活.
究竟當上帝有了神蹟.
去改變一個人之後.
究竟這個人往後是怎樣生活.
這個是主耶穌更加看重的.
信耶穌只是一個開始.
基督徒的屬靈生命能不能夠得到成長.
信徒能不能夠活出耶穌基督的樣式.
主同樣重視的.
事實上整本聖經絕大部分的篇幅.
並不是教人怎樣信耶穌.
你會發覺信耶穌的篇幅很少.
不過怎樣在跟隨耶穌之後.
在上帝的救恩底下.
過一個合上帝心意的生活.
你可以說整本聖經都在說這個課題.
主耶穌告訴我們.
要從心底裡流露敬虔.
生命要得到改變.
首要的條件就是這種渴慕的心.
渴慕的態度.
渴慕上帝的明媚尊崇.
渴慕上帝的國.
透過教會臨到人間.
渴慕更多人體現上帝的愛和公義.
唯有心存這種渴慕.
我們的屬靈成長才不會迷失.
好了.
我們具備了渴慕的心.
接下來怎麼辦呢.
主耶穌告訴我們.
門徒要帶著這種態度去面對社會的現實.
而當我們要面對這個現實的時候.
你就會發現我們原來很有限.
我們就要在極有限的條件底下.
靠主尋求出路.

$^{561}$所以第二個方面.
主耶穌教導我們.
我們要明白我們的需要和限制.
以致我們可以依靠主.
我們就能夠經歷主.
接著.
主耶穌就當時現實的處境.
祂提出了三個祈求.
都是大家很熟悉的.
我們一起讀好不好.
預備起.
我們日用的飲食.
今日賜給我們.
免我們的債.
如同我們免了人的債.
不叫我們陷入試探.
救我們脫離冷惡者.
相對於法利賽人.
法利賽人在第五章.
耶穌說他們以施捨.
禱告.
禁食的行動.
都是去表達他們的敬虔.
但是.
主耶穌說哪三樣東西能夠表達我們的敬虔呢?.
分別是.
飲食.
債務.
和誘惑.
你聽完之後你都覺得一點都不熟悉.
那些應該不是我們的朋友.
應該禁食.
祈禱才是.
為什麼主耶穌舉一個這麼不起眼的例子呢?.
不過.
主耶穌之所以這樣說.
其實是他沒有迴避到社會的現實.
他要門徒去學習.
在現實生活裡面.
我們才能夠表現敬虔.

$^{601}$如果純粹在禁食,在祈禱上.
那樣東西很抽離.
不在我們的生活當中.
在這裡.
是更加能夠表現出.
我們怎樣來都是去敬虔向主的.
我嘗試將這三個部分.
分成三個層面.
第一個層面.
是生存的層面.
這裡所指的飲食.
不是心靈雞湯.
不是你心靈的飲食.
耶穌說.
「我有食物是你不知道的」.
不是的.
這裡更加不是在說.
啟示錄裡面高揚的婚宴.
這裡所指的.
是真實的.
每人每一天.
賴以生存的糧食.
當時的社會.
只有極少數人比較富裕.
而跟隨主的門徒.
大部分都是來自生活貧乏的階層.
他們每一天都要為生存來掙扎.
為生存來困擾.
可能你會問.
主要門徒為飲食祈求.
是不是要鼓勵人.
望天打卦不勞而獲呢?.
我們每天.
上帝給我們飲食.
今天賜給我們.
我們就坐定一粒鹿了.
我相信不是.
事實上.
當時和今天.
有些概念是很不一樣的.

$^{641}$今天我們通常.
粗俗一點說.
我們通常是出了事才求上帝.
最後五分鐘就遲到了.
我們就開始祈禱.
求上帝給個巴士或者的士.
可以出現在我眼前.
當時的世界不一樣.
他們深深明白到世事變幻無常.
你可以想像一下.
當時其實.
特別是中東的地方.
他們沒有河流.
他們要靠雨水才能夠種植.
他們才能夠有收成.
有糧食.
所以.
他們很努力地種植.
打魚,播種,放羊.
在做這些事情的同時.
他們仍然祈求上帝保守他們.
他們不是不努力的.
他們很努力地去做.
但同時也祈求上帝保守他們.
他們每一天.
每時每刻都為了他們的生活感恩.
在大概十九世紀的時候.
如果你看一些歐洲的名畫.
因為大部分當時是信徒.
農夫晚上差不多日落的時候.
他們會在田裡祈禱.
我相信和當時的人一樣.
他們是明白到.
祈禱不是出了事才做.
時時刻刻都仰望上帝.
所以主耶穌說這句話.
是要鼓勵當時的人.
在努力尋找生活出路的同時.
我們要對上帝有信心.
要繼續堅持.

$^{681}$為每天生存去仰望上帝.
不是我能解決的事情.
就自己解決.
我解決不了的事情才找上帝.
任何事情我們都要仰望上帝.
第二方面是社交層面.
聖經裡面主耶穌提到.
「免容忘的債」.
這個「債」縱然可以指前債.
但是在當時耶穌的時代.
那些人用這個字也會表達.
他們對上帝或對人的虧欠.
這個「債」可以說是對人.
或對上帝的虧欠.
如果按照這個意思去理解的話.
求天父免我們的債.
就等於求主赦免我們對他的虧欠.
簡單來說就是求上帝赦免我們的罪.
不過這並不是經文的重點.
當時的人我相信.
大家都會做到這件事.
每一天都會求上帝免我們的債.
然後拖羊去宰殺.
去獻給上帝.
當時這一段教導的重點是.
後面那一句.
就是說當時不單止他們會求上帝赦免他們的罪.
而主耶穌也期望我們學習.
像上帝一樣.
免身邊人的債.
赦免別人對我們的虧欠.
主耶穌不單止知道當時的人.
每天為生活而擔憂.
更加知道當時的社會複雜性.
以至到在這個社會裡面有很多矛盾.
當時是有階級制度的.
我們想像一下.
羅馬人統治的世界.
所以羅馬人壓迫當時的猶太人.
猶太的平民仇視當時做抽稅的同胞.

$^{721}$那些稅吏.
當時他們很不滿意.
這些人竟然背叛自己的同胞去抽稅.
法利賽人又排擠那些外邦人.
你可以說.
仇恨是隨著彼此的衝突不斷的深化加增.
主耶穌明白現實的處境.
不過.
他也提出一個終止仇恨的選擇.
這是一個什麼選擇呢?.
就是免債.
免債並不是從此變回好朋友.
當一場誤會沒事發生.
而是為彼此的報復畫上句號.
而當時要門徒去承認.
原諒是極為困難的事.
門徒要接受自己未必有這個胸襟.
不過我們可以仰望上帝.
我們未必可以完全原諒對方.
但我們就期望上帝幫助我們.
原諒別人.
放過自己.
第三個層面是精神的層面.
主耶穌從來沒有允許信主的人不會受到誘惑.
而且性過誘惑的機率很低.
我不知道大家有沒有做過這個統計.
信主以來.
究竟我們可以性過誘惑的時間多一點.
還是失敗的時間多一點.
我自己的統計.
失敗的時間遠遠多於成功的時間.
是的.
誘惑可以說是隨時隨在.
門徒不要以為跟隨了主之後就百毒不侵.
相反.
主耶穌希望門徒更加理解自己的限制.
不要自視過高.
以至我們在生活上需要保持警惕.
求主賜給我們遠離魔鬼的試探.
尤其是當我們在失敗的時候.

$^{761}$我們就更加需要再接再厲.
明白.
我們有這麼多的限制.
我們需要重新仰望主.
來到這裡.
大家可能會問.
主耶穌要求門徒在這三種生活層面上仰望主.
到底跟之前主要門徒保持對上帝的渴慕有什麼關係呢?.
好像這個跟實踐沒有什麼關係.
為日常的生活祈求.
跟上帝的名,上帝的國有什麼關聯呢?.
第一次我們試想想.
如果我們連基本的生存能力都沒有的話.
我們拒絕適應社會的步伐.
我們試問.
又如何能夠使上帝的名受尊崇?.
如果我們不能夠為教會的合一的緣故終止仇恨.
我們如何能夠讓人看到天國在人間?.
如果我們的心智沒有完全被鍛鍊.
能夠提早避免各樣的誘惑.
上帝的愛,上帝的公義.
如何能夠透過我們實踐呢?.
人們看到的是什麼呢?.
當我們相信耶穌的時候.
我相信大部分人不是要看一個完美的人.
如果今天我們看到基督的牧師完美無瑕.
我相信他也是很完美的.
但是如果他真的完美無缺的話.
或者你今天看到我是完美無缺的話.
可能你第一件事就是說.
那我不要相信耶穌了.
因為相信耶穌要去到這麼高的程度.
我才可以相信主.
我怎麼相信呢?.
是的.
我們今天所謂的見證上帝不是完美無瑕.
而是正正是在我們掙扎的當中.
在人們看到我們都是一個很平凡的人物當中.
看到上帝的恩典.
看到我們的轉變.

$^{801}$看到我們的進步.
以致人能夠在我們裡面.
看到我們背後的上帝.
我相信這個正正就是我們.
遵行上帝說話的要訣.
正正因為我們渴望主的名.
能夠被尊為聖.
我們勇敢認識社會.
我們尋求生存的法則.
以致我們在身邊的人當中.
留下美好的見證.
基於渴慕別人看到上帝國的來臨.
我們一班弟兄姊妹願意化解一下矛盾.
保持聖明所賜給我們的合一.
最終當我們接受自己的限制.
逐漸學會對各種誘惑.
能夠提高警惕之後.
我們就能夠建立自己的誠信.
讓人透過我們看到上主的信實.
以致我們能夠體會.
上帝就在我們當中.
但願我們一起.
帶著對天父這種渴望.
努力,勇敢地在生活上實踐,學習.
並且在當中我們一同去仰望主.
靠上帝幫助我們.
做一個裡外合一,敬虔的人.
我們一起祈禱.
天父說他今天時間有限.
我們不能夠舉很多例子.
來說明敬虔的生活.
但我深信作為每一位弟兄姊妹.
特別是洪恩家的每位弟兄姊妹.
我相信他們已經認識你.
他們也熟習上帝你的說話.
但願主耶穌的教導.
成為我們的指標.
成為我們追求的目的.
但願我們能夠成為一個真正敬虔的人.
我們渴望遵行上主.

$^{841}$我們渴望去追求上帝的聖潔.
祈禱奉耶穌基督的名字.
阿們.
\newpage



\section{約翰一書 4:19-21}
\label{sec:514XvjAlzD8}
\textbf{2023.10.15 《我們愛,因為神先愛我們》 約翰一書4\_19-21 (和修版) 曾敏傳道}
\newline
\newline
連結: \href{https://youtube.com/watch?v=514XvjAlzD8}{\texttt{https://youtube.com/watch?v=514XvjAlzD8}} ~~~~ 語音日期: 2023-10-16
\newline
\newline
\hyperref[sec:ILOp7icU0bo]{\small{< < < PREV SERMON < < <}}
~
\hyperref[sec:index_chronic]{\small{[返順時目]}}
~
\hyperref[sec:index_scriptual]{\small{[返順卷目]}}
~
\hyperref[sec:EQ76QKt8XLA]{\small{> > > NEXT SERMON > > >}}
\newline
\newline
約翰一書 4:19-21
\newline
\begin{longtable}{cl}
\hline
\hline
章節 & 經文 (和合本修訂版)\\
\hline
4:19 & \begin{tabularx}{0.7\textwidth}{X} 我們愛，因為神先愛我們。 \end{tabularx} \\ \\ \relax
4:20 & \begin{tabularx}{0.7\textwidth}{X} 人若說「我愛神」，卻恨他的弟兄，就是說謊了；不愛他看得見的弟兄，就不能愛看不見的神。 \end{tabularx} \\ \\ \relax
4:21 & \begin{tabularx}{0.7\textwidth}{X} 愛神的，也要愛弟兄；這是我們從神所受的命令。 \end{tabularx} \\ \\
[1ex]
\hline
\hline
\end{longtable}
$^{1}$領子妹平安.
今日見到領子妹可以齊齊聚集在此.
很想在此和大家再讀一次經文.
這次是大家一起讀.
我們回到剛才讀經的地方.
我們看經文是約翰一書四章十九到二十一節.
雖然只是三節的經文.
但卻是很豐富的內容.
今天這一堂道旨既是對我講.
亦是對大家講.
我們一起讀.
預備起.
我們愛,因為神先愛我們.
人若說我愛神,卻恨他的弟兄,就是說謊了.
不愛他看得見的弟兄,就不能愛看不見的神.
愛神的也要愛弟兄.
這是我們從神所授的命令.
我們開始的時候再一同去禱告.
要成為我們腳前的燈,路上的光.
今日求你開我們眾人的心眼.
讓我們看到神你是怎樣的.
讓我們眾人都聽到你的聲音.
願以聖靈在我們裡面去工作.
叫我們自己去回應你.
願神在我們當中.
得著一切的榮耀和頌讚.
奉耶穌基督的名字祈禱.
阿們.
我們愛,因為神先愛我們.
我們講愛的功課.
其實教會其中一個很突出的地方就是講愛.
在約翰一書一章三節裡這樣說.
大家一邊看一邊聽我說.
這裡說.
我們將所看見所聽見的.
傳揚給你們.
為要使你們也與我們有團契.
而我們的團契是與父和他兒子.
耶穌基督所共有的.
這裡說到團契.

$^{41}$其實也說到教會一個很重要的特色.
我們回來教會是為了什麼呢?.
台下的弟兄姊妹也好.
新的朋友也好.
如果你想到的.
都大聲在位置說出來.
回來教會是為了什麼?.
非常好,敬拜神啊.
還有沒有?.
很聰明吧.
還有沒有?.
回來教會是為了什麼?.
除了敬拜.
大聲就在位置說出來.
聽神的話.
這裡除了敬拜和聽神的話.
還有團契.
很多時候我們說團契.
很自然想起我們每個星期某一天進行的團契.
我們與某某弟兄姊妹一起的.
其實團契不單是說那個團契.
是說我們整個大群體.
在弟兄姊妹裡面的.
我們一起的那個團契.
這個團契不只是有弟兄,有姊妹.
也有神.
是有我們三一的真神.
聖父,聖子,聖靈與我們一起.
所以這個團契的意義是非常非常豐富的.
教會就是要與神,與人保持團契的關係.
我們回教會這麼多年.
今天是一個時間重新的回顧.
與神,與人的團契是什麼團契呢?.
這個我們再去看.
這個約翰一書四章十六節這樣說.
我們知道並且深信神是愛我們的.
我們在教會裡說得最多的.
就是上帝很愛我們.
上帝很愛我,很愛你.
神就是愛.

$^{81}$這裡不單是說神裡面有很多愛.
不單是說.
神自己就是愛.
祂的本質就是愛.
祂是什麼?就是愛.
這個意義是很深厚的.
住在愛裡面的.
就是住在神裡面的.
神也住在祂裡面.
我們互相住在對方的裡面.
神在我們的裡面.
我們住在神的裡面.
既然神是愛.
就是我們住在愛裡面.
回教會原來是這樣的.
我們屬上帝的人是住在愛裡面的.
自從我們信耶穌之後.
我們就與主結連.
我們住在神裡面.
也住在愛裡面.
也住在我們裡面.
這是一個愛的團契.
很肯定的.
我們愛神.
我們也愛人.
愛是教會一個很重要的特色.
我們愛神和人是因為什麼?.
是因為神先愛我們.
在經文裡提到.
我們愛不是因為我們很有能力去愛.
我試過不止一次.
在教會以外與人接觸,相處.
可能當時是為人提供一些幫助.
或者是做了一些事情.
很自然的.
也通常會稱讚你.
你真的很有愛心.
但很多時候我都會補充一句.
其實我不是真的這麼有愛心.
我的上帝是很有愛.

$^{121}$我不是真的這麼有愛.
如果說到我們自己.
我們做很多事情.
我們為教會也好.
為弟兄姊妹也好.
為外面的人也好.
做很多很多事情都好.
記住一件事.
說真的.
我們這個人.
不是說我們自己.
這個人多有能力去愛.
說我這個人.
有無窮無盡的力量去愛.
不是.
因為人的愛是很有限制.
很有條件的.
今天我們說愛的時候.
不妨來看看.
我們人的愛是怎樣的?.
是不是真的無限的?.
我所有人都擁抱的.
徹底的,百分百的擁抱.
是不是?.
我看到台下搖頭的佔大多數.
不是的.
是不是?.
很多限制的.
很多條件的.
我們裡面都有很多條條框框的.
每個人都會受到我們成長的影響.
受到我們人生經歷的影響.
我們喜歡某些人多一點.
可能某些人沒有那麼多.
我們會欣賞某些人有些能力.
一說起這些人.
讚不絕口.
真的覺得他們很好.
如果交朋友和這些人就好了.
或者覺得某些人很有成就.

$^{161}$特別有欣賞他們多一點.
第四樣東西是最深刻的.
是什麼?.
都會很受別人對我們的態度影響.
如果別人對我們的態度很好.
對方通常是很欣賞自己的.
又對我很好.
很多的恩惠.
通常我們會怎樣?.
當然喜歡他們.
誰不喜歡別人欣賞自己?.
我們都很自然的.
我們都會很想別人欣賞自己.
很想別人對自己好.
如果我們接觸的人是這樣.
一定很喜歡和這些人交朋友.
我自然去付出都很甘心的.
是很甘心的.
但如果反過來.
對方是怎樣的?.
哇! 頂心慘.
我講一句.
他頂一句.
我講什麼.
他都覺得不順眼.
對了.
馬上感受就不好.
其實我這個都是對自己說.
我發覺這麼多年學愛的功課.
是很不容易的.
反而我真是越來越日子內.
看到神的愛和人的愛是很大分別.
看到自己很多限制.
愛裡面都很多限制.
很多條條框框的.
是否真的不容易?.
人的愛總帶著計算.
是真的計算的.
我付出多少?.
對方付出多少?.

$^{201}$為什麼這麼不公平?.
每次我請吃飯.
對方不可以請我吃飯.
馬上笑.
我請一餐,你請一餐.
講清楚,計算清楚.
我請的餐是什麼?.
計算你的餐是否這樣.
我通常請的都貴的.
你請的就這麼便宜的.
心不舒服?.
是吧? 很多的.
每次生日慶祝.
我總是為你.
做這麼多的安排.
你卻給我這麼多.
我心不舒服.
是計算的.
你說家人就沒有計算.
家人的愛是徹底的.
是吧?.
問一下.
家人的愛是否真的從來都沒有計算?.
也不是的.
也計算的.
而且有很多比較.
那些子女.
通常都拿很多錢回來.
那些家用.
你看我每個月收到多少.
生哪個好一點.
做哪個好一點.
我們都會的.
別說父母.
我們自己都會計算的.
父母.
給誰多一點.
給誰少一點.
怎樣關注.
我們自己都會計算.

$^{241}$不要說別人.
說我們自己.
很多人 配偶都會計算的.
一定不計算.
每樣東西都會計算.
配偶 疼娘家多一點.
呼家多一點.
對朋友好一點.
對我沒那麼好.
對同事好一點.
很多東西都會計算.
我在家裡付出多少.
都會計算.
有時對方對自己態度好.
自己容易一點.
容易受的.
我又會容易回應 多一點愛.
對方對自己不好 我計算.
我下面真的幫你.
冷戰一下.
都很多.
我從來都強調.
不只是某一個人.
是我們每一個.
都很容易計算.
有些人.
會說我會計算少一點.
或者.
我們當中計算多一點 計算少一點.
都會計算.
去到某些位置.
心裡會牢嘈.
因為我們是人.
我們都有我們的罪性.
我們有限制.
真的有限制.
人就是這樣.
但上帝的愛是怎樣的呢.
上帝的愛.
我們從幾個方面去看.

$^{281}$詩篇139篇.
13到14節.
這樣說.
我的肺腑是你所造的.
詩人這樣說.
而你是誰 是上帝.
我的肺腑是你所造的.
我在無腹中 你已偏激我.
我要稱謝你 因我受造奇妙可為.
你的作為奇妙.
這是我深深知道的.
當我們還在媽媽肚腹中.
神已經在我們生命中工作.
神正在預備我們.
這個人.
我們的性情.
我們的特質.
我們的強項.
我們的喜好.
很多東西.
上帝在我們媽媽肚腹中.
已經在預備我們這個人.
而且偏激.
這裡說的是保護.
我們在媽媽肚腹中.
神已經在保護我們.
一直保護.
祂的眼睛.
一直看著我們.
看著我們 看顧我們.
保守我們.
神一直欣賞.
祂自己的傑作.
那個傑作就是.
我和你.
我們每一個.
都是神獨立的.
單獨的做.
我和你的DNA.
在這個世界上是獨一無二的.

$^{321}$就算那些人.
是雙胞胎.
他們的DNA在這個世界上.
都是獨一無二的.
每一個都是神精心製作的.
神欣賞自己所做的.
一路在欣賞.
一路保護.
一路看顧.
一路還有更多其他的東西.
這裡139篇16節.
這樣說.
我未成形的體質.
你的眼早已看見了.
你所定的日子.
我尚未到一日.
都在你的冊子寫上了.
我的人生未過一日.
神都知道.
寫在冊子上.
祂不是冰冰冷冷地寫.
而是充滿愛地寫.
我們是祂的寶貝.
祂一路寫.
一路寫.
祂的人生在祂的手中.
祂對我們的愛.
是很長遠.
很恆久的.
甚至乎按照聖經的記載.
其實在永恆裡.
上帝已經愛我們.
不只是我們在.
媽媽肚腹裡那一天開始.
在永恆裡.
祂已經愛我們.
我們出現在媽媽肚腹裡的時候.
更加.
有很多的計劃.
這份愛已經很久了.

$^{361}$很深遠.
詩篇139.
17-18節這樣說.
神啊!你的意念向我何等寶貴.
其數何等眾多.
我若數點比海沙更多.
我睡醒的時候.
仍和你同在.
上帝對我們一生.
有數不盡的旨意.
是非常美好.
寶貴的旨意.
今日弟兄姊妹我們知道有多少.
由我們在媽媽肚腹裡.
到我們一生走到這一刻.
我們知道有多少神的旨意.
無論今天你信了耶穌沒有.
上帝對我們人生.
都有祂的旨意.
我們知道多少.
上帝原來關顧我們.
關注我們這麼久.
愛我們很久很久.
那種關心.
那種愛.
是很仔細很細膩的愛.
那種旨意.
不只是說.
祂希望我們人生怎樣.
而是上帝有很多的供應.
很多的預備.
在盧家灰音12章7節這樣說.
就算你們的頭髮也都數過了.
不要懼怕.
你們比許多的麻雀還貴重.
當時在做一個比較.
是不是一個比較呢.
和那些鳥仔比.
我們是怎樣呢.
當然我們人貴重很多.

$^{401}$神呀連我們頭髮都數過.
今日問下在座這麼多個.
在你人生裡面.
最愛你的那個人.
數過你多少條頭髮尾.
最愛你的爸爸媽媽.
子女也好.
數過你多少條頭髮尾.
還沒有啦.
你知道神細心到這樣嗎.
頭髮都數過.
不只是數頭髮.
而是我們每一樣東西.
我們的細胞所有的東西.
祂都知道了.
而且祂會供應我們.
滿足我們.
所以神想我們不要憂慮.
神知道我們的需要.
雖然在主禱文.
會向我們禱告.
用飲食求上帝賜給我們.
但其實我肯定的一句.
就是在我們禱告之先.
神已經知道.
並且預備.
在這裡你見到.
神多細心.
從小到大.
祂為我們預備了多少東西.
如果要為了這些.
去數算主因.
我相信我們是數之不盡.
我記得.
我們有一段時間.
在香港都很想去說.
一件事.
養一個小朋友有多貴.
養一個小朋友至少400萬.
到現在再說不止.

$^{441}$這個數字很小.
不止.
更多的數字.
養一個小朋友很貴.
一個家庭.
預備的就是有限數的錢.
或者有限的物質.
是不是.
但上帝為我們預備的.
何止那400萬,800萬呢.
何止那些家裡的物質.
從小到大.
所預備的何其多.
神肯定.
比我們的父母.
更愛我們.
肯定比我們的配偶更愛我們.
一定.
神對我們的愛.
是徹底的.
神認識我們.
愛我們.
樂於看顧我們.
就算眾人不喜歡我們.
如果有這樣的光景.
就算眾人不喜歡我們.
他還是能夠看到.
我們的好,欣賞和愛我們.
為什麼?.
大家告訴我.
為什麼他可以這樣做.
大聲點說.
他是神.
還有,是他做我們的.
他對我們的欣賞和接納.
是百分百的.
什麼叫百分百?.
就算我們有很多不好的東西.
神都愛和接納我們.
是真的.

$^{481}$這些可能我們以前.
理性上聽過.
但我們未能夠聽入心.
可能.
如果我們的光景.
是很多人都不喜歡我們.
如果真的有這個.
這樣看我們.
那個這樣看我們.
神看我們.
還是看到我們的好.
神還會欣賞我們.
還是很愛我們.
這和人是很大分別的.
人.
一看到一些東西.
覺得在他人身上.
是很不喜悅他.
可能整個人都會受到影響.
就不喜歡他.
就不想和他接觸.
不想和他相處.
在教會裡也是.
都是會這樣.
因為我們是人.
我們會計算.
我們會有很多限制.
但上帝不是這樣.
上帝的愛.
是很超越的愛.
上帝會徹底擁抱我們.
當我們孤單的時候.
我們想一想.
神怎樣都在我們身邊.
全然百分百的欣賞.
而且神愛我們.
到了一個地步.
不單止是這麼多年.
他賜下生命氣息.
不單止是這麼多年.

$^{521}$給我們供書教學很多的供應.
不單止是接納我這個人.
這個人的好與不好.
他愛我們到了一個地步.
是什麼?.
我已經聽到有些聲音.
為我們死.
《約翰一書》四章十節.
他這樣說.
不是我們愛神.
而是神愛我們.
差祂的兒子為我們的罪.
作了贖罪制.
這就是愛.
神很愛我們.
很愛我們.
知道我們的本相.
就是我們經常犯罪.
我們從小到大犯了多少罪.
其實很多罪都不用人教.
我們就懂得.
你想想小時候.
什麼時候開始說謊話.
我們小時候.
什麼時候會說是非.
出貓這些事.
不用人教.
口說的說話很多.
不只是說髒話.
罵人的說話.
欺騙的說話.
我們的背叛.
我們的傷害.
不單止口說.
還有很多思想言行.
很多事情就是這樣.
默默地進行.
我們可以說.
因為成長.
我們原生的家庭.

$^{561}$對我們有很多負面的影響.
我不否認.
原生的家庭對我們有影響.
但我們自己會更加變本加厲.
會繼續犯罪.
從小到大這麼多的罪.
我們帶著這些罪.
根本不可能.
原本的我們不可能和天父復和.
不可能進入神的國度.
不可能得到那份禮物.
我們帶著的那份禮物.
永生一定不可能.
等著我們的「份」在哪裡?.
我們這樣犯罪.
最終的結局應該是什麼?.
死呀.
是永死.
永死呀.
這個結局是多慘.
永遠.
我們要為我們的罪.
負上代價.
等著我們的.
原本是地獄.
這個永死的結局.
是很慘很慘.
根本我們是承受不了.
但神卻愛我們.
派祂的兒子耶穌基督.
為我們犧牲.
捨忌.
耶穌在十字架上.
為我們眾人的罪死了.
我要強調.
是眾人的罪.
為我.
為你.
為我們.
為全世界的人死了.

$^{601}$為那些眾人.
看起來只是犯輕微罪惡的人.
為那些眾人.
看起來是犯很重罪的人.
無論是犯重罪的.
輕罪的.
多罪的.
少罪的.
都死了.
這個代價是為所有人付出的.
耶穌基督不單止為我們犧牲.
捨忌.
第三天也從死裡復活.
為我們成就救恩.
好讓我和你.
得蒙救贖.
這是愛的最高境界.
耶穌不要祂自己的生命.
甘心捨棄.
祂的生命.
用我和你的生命.
是不是因為我和你很好呢?.
我想問弟兄姊妹.
今天好像重溫救恩.
回歸到我們信仰的基本.
是不是因為我和你很好.
實在太好了.
是不是?.
所以祂不救我們不行.
肯定不是.
是正正因為我們不好.
才要救我們.
就是因為.
全世界的人都犯了罪.
在這個世界.
除了耶穌基督是一個.
真正徹底的義人.
所有其他的人.
都是犯罪的.
你記得耶穌基督曾經渡成肉身來到這個世界.

$^{641}$和人一起生活過.
只有耶穌基督沒有犯過罪.
我們所有在地面上生活的人.
都是有犯罪.
無論世人都犯了罪.
規絕神的榮耀.
所有人都需要.
耶穌基督的救恩.
是因為祂的.
犧牲捨己.
因為祂為我們成就救恩.
所以在《育恆一書》四章十七節說.
愛在我們裡面.
得以圓滿.
我們信耶穌.
我們進入到神的愛裡面.
我們可以在審判的日子.
坦然無懼.
我們不用再害怕.
到時要面對.
永世的刑罰.
因為基督如何.
我們在這世上也如何.
耶穌基督.
已經拯救我們生命.
我們脫離.
永遠的死亡.
耶穌基督復活.
勝過死亡.
今天我和你.
屬耶穌的人.
我們都可以勝過死亡.
這個救恩.
是何等寶貴.
我和你.
都不用害怕.
但你記得.
耶穌付的代價很大.
不是就這樣.
我和你.

$^{681}$就可以脫離.
永世的捆綁.
是因為耶穌.
付的代價.
正正因為這樣.
大家看這圖片.
靠著十架恩典.
我們就可以.
一直直奔走向永生.
只有這條路.
是神為我們預備的.
神因為愛我們.
就預備這條路.
我和你無論做多少好事.
多少善事.
如何待人.
都不能抵消我們生命所犯的罪.
唯有耶穌基督.
救我們的生命.
我們才可以奔向永生.
這正正是.
看到上帝的愛.
多麼深厚.
在這個世界從來沒有人為我們這樣做過.
也沒有人可以.
有這樣的能力為我們這樣做.
唯有耶穌基督.
唯有上帝.
正正是因為.
神才愛我們.
所以我們有能力去愛.
我們和上帝.
結連之後.
神在我們裡面.
我們就有這個愛的能力.
去愛,不是用我和你.
本身那份愛的力量.
要記住.
我現在有多有能力.
去學愛也好,都是神在我裡面.

$^{721}$去愛,所以要將榮耀歸給神.
不是我這個人多厲害.
我們有能力.
去愛.
神希望我們彼此相愛.
在這裡說.
我們有能力去愛.
在兩個層面去說.
第一,我們有能力去愛.
那些未信耶穌的朋友.
洪水橋有很多.
街坊朋友.
未曾信耶穌.
因為有神在我們裡面.
有神的愛在我們當中.
我們有能力去愛那些街坊.
向他們傳福音.
向他們表達我們的關懷.
我們的扶持.
但在教會裡面,神希望我們彼此相愛.
這是命令.
不是說.
我想想.
如果我想做就做.
不是這樣,是神要命令我們要做.
要彼此相愛.
就算在約翰一書.
我們今天看的經文第四章裡面.
都有繼續提到彼此相愛.
約翰一書裡面提過不止一次.
彼此相愛.
要一齊.
彼此相愛的基礎.
是因為什麼?.
因為神才愛我們.
你記住不是因為.
對方很可愛.
因為對方很愛我.
對我很好.
對方是一個好人.

$^{761}$他太可愛了.
我不愛他就不行.
聖經裡面沒有說這些.
是因為神的命令.
對方可不可愛.
對方對你怎樣.
我們都彼此相愛.
你說這樣嗎?.
是的,就是這樣.
這個功課不容易學.
對我也不容易學.
但事實上.
現在我們面對彼此相愛的功課.
我們事實是什麼?.
不容易.
聖經繼續這樣說.
約翰一書四章二十節.
我愛神,卻恨他的弟兄.
就是說謊了.
不愛他看得見的弟兄.
就不能愛看不見的神.
我說我愛神.
我們說我們愛神.
但我們不愛我們的弟兄.
這是說謊.
神在這裡做了一個對比.
為什麼我們要這樣呢?.
看得見的.
對比什麼?.
是看不見的.
我們現在看見的是弟兄.
今天在禮堂裡.
坐了很多弟兄姊妹.
我們看見的就是.
弟兄姊妹.
真是有血有肉的.
活生生的.
有生命的.
在我們當中坐在這裡.
你立即和對方聊天.

$^{801}$對方有回應.
會聽到我們說話.
又會給回應.
當我們要和對方相處的時候.
我會較容易做表達.
我給他一件東西.
我為他禱告.
即時對方又有回應.
因為你看見.
交流也較容易.
但看不見的上帝.
不是說我不能對神有表達.
只不過你對比看見的人.
其實那個表達.
會再進深一層.
有些不同的.
你說現在.
我們向上帝禱告.
神呀 我很愛你.
我透過不同的方面表達.
對你的愛 例如我.
很熱心的服侍.
例如我奉獻.
或者是怎樣也好.
在過程中 你要去等.
神的回應.
應該是對比你.
對弟兄姊妹.
一定是進一層.
沒有那麼容易.
我們聽神的聲音.
一定要去等候.
但神一定會和我們說話.
一定會聽到神的聲音.
亦會有神的回應.
但對比人來說.
人會更快更直接.
有這個分別在.
但我想說.
看得見的 看不見的.

$^{841}$絕對有神的元素在裡面.
看得見的弟兄.
或者是姊妹.
都有神在裡面.
有神在裡面.
就有神的愛在裡面.
今天在聖經裡說.
如果我不愛這個.
看得見的弟兄姊妹.
其實我亦不愛神.
我騙自己說.
我是很愛神.
看得見的.
你能和他有接觸.
能夠有表達.
你不能愛他.
我們不能愛他.
那我們說什麼我們愛神呢?.
記住.
有神在裡面.
有神的愛在他當中.
我不能擁抱他的生命.
這個.
我自己都要向神交仗.
是真的很不同.
我們和弟兄姊妹一樣.
都有同一個信仰.
我們都相信耶穌基督.
在神的裡面.
在愛的裡面.
所以我們愛對方.
不是單單的.
亦是愛神.
但很多時候.
你會發覺我們有挑戰.
我們不是每個人都愛.
你會發覺.
我們不是每個人都愛.
有些我們愛.
有些我們不那麼愛.

$^{881}$我們不愛弟兄姊妹的原因是什麼?.
為什麼妹妹不愛呢?.
因為她也屬上帝.
她也與神結連.
有神在她裡面.
有神的愛在她當中.
為什麼我不能夠愛呢?.
可能我們說.
有事發生了.
在我們的相處裡.
有事發生了,我不能夠愛.
可能對方傷害了我.
或者相反.
我傷害了她.
對方不接納等等.
對方可能有不同的.
有ABCDE.
有很多不同的情況.
所以就不愛了.
我們不愛.
為什麼?.
有些地方值得我們反思.
其中一件事.
我們不太了解神的愛是怎樣.
剛才我花了很長的篇幅去講.
神的愛是怎樣的.
就算眾人不愛他.
他也會愛我們.
也會欣賞我們.
擁抱我們,接納我們.
一定看到.
人們都覺得我們有幾十八樣不好.
神也看到我這個人有好處.
從小到大.
只有祂一個.
從來都不離開我們.
在媽媽肚腹裡.
一直到如今.
都愛我們,供應我們,照顧我們.
保護我們,還要拯救我們.

$^{921}$直到如今.
神的愛是徹底的.
犧牲捨己的.
沒有保留的.
你說今天我們去愛人.
你說我們有能力.
是因為有神的愛.
神的愛在我們裡面.
也是可以發揮這個功用.
但今天很多時候.
我們用自己的血氣.
去對待人.
說到底.
這個我同樣挑戰我自己.
我是否也在用我的血氣.
對方對我好.
我就對他好.
我喜歡和他相處.
就對他好一點.
這些愛是人的愛.
是血氣的愛.
其實說真的.
未信耶穌的人也可以.
真的有不少.
未信耶穌的人.
很有愛心.
可能做得比我們更好.
你對他好,他就對你好.
他也會做很多好事.
我們和他們有什麼分別呢?.
大家也在用血氣的愛.
但現在說的是.
上帝的愛.
弟兄姊妹.
我們是否真的了解.
上帝的愛是怎樣的.
在我裡面的.
現在是上帝的愛.
不是說我自己喜歡或不喜歡.
我打算怎樣對待他.

$^{961}$我們除了不了解上帝的愛.
我們也忘記了自己的本相.
其實我們自己也是罪人.
我們忘記了.
特別是在教會越久.
我們越容易忘記.
忘記我們是罪人.
我們也有夢神的念文.
我們是得夢救贖而已.
所以我們坐在這裡.
我們是蒙恩的.
什麼?.
我們是罪人,我們有神的恩典.
神是何等的愛我們.
我們是不配的.
你要記住自己的本相是罪人.
很多時候.
信主信久了,覺得自己也很好.
是不是?.
我很好,所以看到弟兄姊妹的不好.
就很不喜歡.
非常不喜歡.
我拿著聖經.
看著他.
哪裡不好.
我要判他刑.
因為我很好.
其實我們也只是.
蒙恩的罪人.
我們不配的.
如果我們這樣捉住別人.
上帝也會捉住我們.
我覺得別人不好.
上帝也會說.
我們很好嗎?.
在這裡.
有沒有問自己.
問自己.
今天的我們是什麼情況.
在這裡.

$^{1001}$有一個挑戰.
讓我們回歸神愛的懷抱.
今天真的很想.
我們離開禮堂的時候.
我們回歸神愛的懷抱.
離開禮堂的時候.
不要忘記了.
要想起神是怎樣愛我們.
多年來.
我們怎樣經歷祂的愛.
放下自己.
特別是那份很血氣的愛.
這番說話.
我同樣對自己說.
我們一起來.
放下自己.
最愛的.
可能我們覺得可愛.
最難的是我們覺得不可愛的人.
神要我們彼此相愛.
不是強行逼出來的.
很淒涼.
夾住我的神.
我不想這樣做.
不是的.
因為上帝這樣愛我.
我很想以愛回應神.
想我愛神所愛.
神就是愛我們每一個弟兄姊妹.
上帝愛她.
因為神這樣愛我.
我愛神.
我也愛弟兄姊妹.
我回應神.
神佩德.
嘗試去發掘弟兄姊妹的好.
你看到她十幾二十樣不好.
你也去找她的好.
我很欣賞我們弟兄姊妹當中.
她們真的很好.

$^{1041}$很能夠去發掘別人的好.
我們一起來學.
學去發掘.
就算這樣也好.
但也有好處.
這麼多不好的事也有好處.
不要只看別人十幾二十樣不好.
自己也一樣.
你記住.
主禱文怎樣說.
免我們的債.
我們不是唸口黃.
免我們的債.
求神寬恕我們.
如同我們免了別人的債.
我們免了別人的債嗎.
彼此饒恕.
放下自己.
彼此饒恕.
彼此欣賞.
但在這裡有一個很重要的提醒.
最後.
無論我們自己也好.
別人也好.
有什麼不好的事.
例如我們犯罪.
生命有得罪上帝的地方.
你當看不到.
沒事.
我又沒有做這些.
彼此就算了.
沒事了.
上帝要我們彼此相愛.
你就當這些不好的事都沒有.
這樣也不是.
上帝是聖潔公義的神.
最後我一定要提醒.
上帝是聖潔公義的神.
他不會以為犯罪的.
就當沒事發生.

$^{1081}$就當沒有犯罪.
教會裡千萬不要有彼此的提醒.
也不要說.
總之什麼都好.
弟子會對我們全部都好.
我們自己全部都好.
我們從來沒有犯罪.
約翰一書不是這樣說.
我們約認自己的罪.
信實的是公義的.
必要赦免我們的罪.
洗淨我們一切的不義.
還要更好的是彼此認罪.
彼此認罪.
又不要走到另一個極端.
不面對罪的問題.
讓我們都要用愛心說誠實話.
但是要面對罪.
不要讓罪.
在我們當中發酵.
彼此用愛心提醒.
是要很敏銳.
聖靈的聲音.
我深願神對我們每一個說話.
深願聖靈對我們說話.
聽到祂的提醒.
我們就很認罪悔改.
在啟示錄.
我記得我教主學時說過.
啟示錄的頭一部分.
是說.
神審判教會.
沒錯 神是愛.
我今天說了很多愛.
有關愛的事情.
但神也是聖潔公義.
聖潔公義和愛是一切.
正正是因為神是聖潔公義.
才看到他的愛的心.
他的心喉.

$^{1121}$所以今天我們在經歷神的愛.
但神也要我們認識他的聖潔公義的性情.
離開罪惡.
在這個群體裡彼此相愛.
彼此認罪.
彼此提醒對方離開罪惡.
這樣才真正能討神的喜悅.
大概是.
幾個星期後.
就是我們.
雄恩嘉十二週年的堂慶.
那天教會也有為大家預備禮物.
我深信上帝也很想.
我們在那天領受更多祝福.
更多從神而來的禮物.
但弟兄姊妹.
你有否想過.
我們可以在堂慶那天.
領受更多祝福呢?.
弟兄姊妹你們有沒有想過.
你們想為上帝預備一份什麼禮物呢?.
我想鼓勵大家開始去想.
你們想給什麼禮物神?.
神最想看到的是什麼?.
大家回答我.
神最想看到的是什麼?.
我們彼此相愛.
這是我和你成長的方向.
讓這成為一份禮物獻給我們的神.
我們一起祈禱.
今天藉著這段經文.
去提醒我們彼此相愛.
因為上帝.
你在我們的生命裡.
靠著你這份愛.
我們就能夠去愛弟兄姊妹.
但主你知道我們軟弱.
我們都是很熟血氣的.
我們都是憑自己喜歡與不喜歡.
去愛.

$^{1161}$主啊 求你讓我們看到.
你的那份愛是不一樣的.
既然你在我們生命裡.
讓我們真的靠著你的愛.
愛我們的弟兄姊妹.
愛可愛的.
愛我們自己覺得.
我們不喜歡的.
主啊 讓我們因為這樣相愛的緣故.
就能夠為你作美好的見證.
在街外的人 在教會裡的人.
都看得到我們是熟悉你的.
我們是主你的門徒.
主啊 求你賜福給紅人家.
給我們接下來的一段時間.
我們繼續朝著這個方向去成長.
在這裡我們能夠彼此相愛.
彼此扶持.
去榮耀上帝你的名.
去滿足你的心意.
求你幫助我們.
特別是我們軟弱的時候.
求你兼顧我們的心.
與我們同在.
奉耶穌基督的名字.
祈禱.
阿們.
\newpage



\section{馬太福音 20:29-34}
\label{sec:EQ76QKt8XLA}
\textbf{2023.10.22 《因 信 奇 癒》 馬太福音20\_29-34 (和修版) 陳一華牧師}
\newline
\newline
連結: \href{https://youtube.com/watch?v=EQ76QKt8XLA}{\texttt{https://youtube.com/watch?v=EQ76QKt8XLA}} ~~~~ 語音日期: 2023-10-31
\newline
\newline
\hyperref[sec:514XvjAlzD8]{\small{< < < PREV SERMON < < <}}
~
\hyperref[sec:index_chronic]{\small{[返順時目]}}
~
\hyperref[sec:index_scriptual]{\small{[返順卷目]}}
~
\hyperref[sec:VjKjBj0FEho]{\small{> > > NEXT SERMON > > >}}
\newline
\newline
馬太福音 20:29-34
\newline
\begin{longtable}{cl}
\hline
\hline
章節 & 經文 (和合本修訂版)\\
\hline
20:29 & \begin{tabularx}{0.7\textwidth}{X} 他們出耶利哥的時候，有一大群人跟隨耶穌。 \end{tabularx} \\ \\ \relax
20:30 & \begin{tabularx}{0.7\textwidth}{X} 有兩個盲人坐在路旁，聽說是耶穌經過，就喊著說：「主啊，大衛之子，可憐我們吧！」 \end{tabularx} \\ \\ \relax
20:31 & \begin{tabularx}{0.7\textwidth}{X} 眾人責備他們，不許他們作聲，他們卻越發喊著說：「主啊，大衛之子，可憐我們吧！」 \end{tabularx} \\ \\ \relax
20:32 & \begin{tabularx}{0.7\textwidth}{X} 耶穌就站住，叫他們來，說：「你們要我為你們做甚麼？」 \end{tabularx} \\ \\ \relax
20:33 & \begin{tabularx}{0.7\textwidth}{X} 他們說：「主啊，讓我們的眼睛能看見。」 \end{tabularx} \\ \\ \relax
20:34 & \begin{tabularx}{0.7\textwidth}{X} 耶穌動了慈心，摸了他們的眼睛，他們立刻看得見，就跟從耶穌。 \end{tabularx} \\ \\
[1ex]
\hline
\hline
\end{longtable}
$^{1}$各位弟兄姊妹.
鄭華的粵劇和龍肉星的見證.
感不感人?.
很感人吧.
來將頌讚歸給天父.
也為龍肉星弟兄感恩.
因為上帝給他有很奇妙的生命改變.
不說你不知道.
原來他有陰陽眼.
不說你不知道.
他信耶穌的過程是這麼奇妙的.
信了之後又信得很堅定.
所以我們為龍肉星弟兄感恩.
也為粵劇感恩.
鄭華跟龍肉星弟兄說.
你今天暢所欲言.
你長火我就短火.
你短火我就長火.
結果我現在要短火了.
很開心.
今天牧師能夠跟大家一起.
聽到這麼豐富.
這麼精彩的見證.
心裡很舒服.
其實我跟龍肉星弟兄很相似.
他這幾天有上呼吸道的問題.
今天大家聽我的聲音.
是不是也有問題?.
是的.
原來我這幾天也是上呼吸道有問題.
所以有時候開聲不能開得太盡.
幸好大家不要要求我唱歌.
但是我這樣細聲說.
大家聽得清楚嗎?.
聽得清楚吧?.
牧師決定將今天的講道化整為零.
我很簡單.
三分鐘內我做一個總結.
然後我會邀請朋友來接受耶穌.
讓你的生命更豐盛.

$^{41}$我一個很簡單的總結.
鄭華看那部粵劇.
看到患血瘤的婦人得到醫治.
最大的秘訣是什麼?.
就是信心.
信心重不重要?.
很重要.
龍肉星弟兄.
為什麼他的生命有這麼奇妙的改變呢?.
也是因為什麼?.
信心.
他信得過的師公.
沒有不好的介紹給他.
介紹給他一定是什麼?.
好東西.
所以他自己也憑信心去接受主耶穌.
所以信心是令我們的人生有很大的改變.
因為大家知道香港這個時候.
很多人都生活得很多壓力.
不開心.
很煩惱.
很多矛盾.
或者很多重擔.
甚至有些人的情緒有很多抑鬱.
你相不相信香港人患抑鬱病的數字節節上升?.
大家相不相信?.
所以你發覺到.
是因為我們去到一個困局.
好像沒有出路.
好像進入了一條很黑的隧道裡面.
很漫長.
看不到曙光.
大家明白這個意思嗎?.
牧師今天不知道大家在當中.
你今天有遇到什麼困難.
我不清楚.
不過我相信.
你無論遇到任何困難.
都不重要.
最重要的是.

$^{81}$你是憑信心去接受耶穌的救恩.
因為耶穌來到這個世界上.
為你為我釘身十架.
流出寶血.
洗淨了我們的過犯.
所以祂來到這個世界上.
是要幫我們走出這個困局.
大家明白嗎?.
走出這個罪惡的枷鎖.
不再被我們人內心那種自私.
驕傲.
貪心.
種種不好的事情纏繞著我們.
以致我們活得不開心.
不快樂.
所以今天牧師很想.
當中每一位的朋友.
都像我們信了耶穌的基督一樣.
現在牧師有兩個呼召.
你聽清楚.
第一個呼召.
我會邀請已經信了耶穌.
你將你的人生交代給主的基督徒.
你待會站起來.
第二個呼召.
我會邀請今天在座的朋友.
牧師,我也很想相信耶穌.
好像剛才那齣粵劇.
和龍弟兄的經歷.
我想我的人生活得開心點.
快樂點.
活得更加有意義.
更加有憧憬.
我想我的人生是一個快樂的人生.
如果你願意的話.
待會第二個呼召.
你也可以站起來.
現在牧師邀請.
凡是已經相信了耶穌.
將你的生命交代給主耶穌.

$^{121}$你站起來.
現在我想邀請.
今天你都願意信耶穌的朋友.
你也站起來.
讓耶穌的生命進入你裡面.
使你活得更開心.
更快樂.
還有沒有哪一位呢?.
還有哪一位呢?.
今天你都很願意.
信靠耶穌.
好像龍弟兄.
好像粵劇當中的弟兄姊妹.
讓我們的生命.
更加精彩.
活得更加有意義.
更加有平安.
你也說很想成為上帝的子女.
你也可以站起來.
表達你的心願.
我們稍後為你祈禱.
還有哪一位呢?.
今天早上開心嗎?.
很開心,很快樂.
為什麼呢?因為我們聽到.
很奇妙的見證.
還有兩位,對不起我看不到.
還有哪一位呢?.
他們願意有.
豐盛的人生.
希望你人生.
活得不再一樣.
一個喜樂快樂的人生.
還有哪一位呢?.
好,牧師呢.
這個時候就為.
所有決志信耶穌的.
一起祈禱,好不好?.
我們低頭祈禱.
你說我不會祈禱,不要緊.

$^{161}$待會牧師說一句,你跟我說一句.
我們現在開始.
主耶穌.
多謝你.
愛我這個人.
雖然我卑微.
我不配.
你沒有嫌棄我.
你來到世界上.
釘身十架.
流出寶血.
赦免我們的罪.
讓我們人生活得快樂.
不再一樣.
多謝主耶穌.
今天我接受你.
做我救主.
進入我內心.
永不離開.
又祝福我家人.
叫他們都相信你.
得到永生的福樂.
祈禱.
奉耶穌名求.
阿們.
很開心.
歡迎今天決志信主的朋友.
就算有些人.
仍然在思考當中.
我們都歡迎他.
我們大家請坐.
\newpage



\section{約珥書 2:10-14-28-32}
\label{sec:VjKjBj0FEho}
\textbf{2023.11.05 《審判或復興》 約珥書2\_10-14,28-32 (和修版) 講員 : 老耀雄牧師}
\newline
\newline
連結: \href{https://youtube.com/watch?v=VjKjBj0FEho}{\texttt{https://youtube.com/watch?v=VjKjBj0FEho}} ~~~~ 語音日期: 2023-11-10
\newline
\newline
\hyperref[sec:EQ76QKt8XLA]{\small{< < < PREV SERMON < < <}}
~
\hyperref[sec:index_chronic]{\small{[返順時目]}}
~
\hyperref[sec:index_scriptual]{\small{[返順卷目]}}
~
\hyperref[sec:EtvLUaUHtqc]{\small{> > > NEXT SERMON > > >}}
\newline
\newline
約珥書 2:10-14-28-32
\newline
\begin{longtable}{cl}
\hline
\hline
章節 & 經文 (和合本修訂版)\\
\hline
2:10 & \begin{tabularx}{0.7\textwidth}{X} 在牠們面前，地動天搖，日月昏暗，星宿無光。 \end{tabularx} \\ \\ \relax
2:11 & \begin{tabularx}{0.7\textwidth}{X} 耶和華在他的軍旅前出聲，他的隊伍龐大，遵行他命令的強盛。耶和華的日子大而可畏，誰能當得起呢？ \end{tabularx} \\ \\ \relax
2:12 & \begin{tabularx}{0.7\textwidth}{X} 然而你們現在要禁食，哭泣，哀號，一心歸向我。這是耶和華說的。 \end{tabularx} \\ \\ \relax
2:13 & \begin{tabularx}{0.7\textwidth}{X} 你們要撕裂心腸，不要撕裂衣服。歸向耶和華－你們的神，因為他有恩惠，有憐憫，不輕易發怒，有豐盛的慈愛，並且會改變心意，不降那災難。 \end{tabularx} \\ \\ \relax
2:14 & \begin{tabularx}{0.7\textwidth}{X} 誰知道他也許會回心轉意，留下餘福，就是獻給耶和華－你們神的素祭和澆酒祭。 \end{tabularx} \\ \\ \relax
2:15 & \begin{tabularx}{0.7\textwidth}{X} 你們要在錫安吹角，使禁食的日子分別為聖，宣告嚴肅會。 \end{tabularx} \\ \\ \relax
2:16 & \begin{tabularx}{0.7\textwidth}{X} 聚集百姓，使會眾自潔；召集老年人，聚集孩童和在母懷吃奶的；使新郎出內室，新娘離開洞房。 \end{tabularx} \\ \\ \relax
2:17 & \begin{tabularx}{0.7\textwidth}{X} 事奉耶和華的祭司要在走廊和祭壇間哭泣，說：「耶和華啊，求你顧惜你的百姓，不要使你的產業受羞辱，在列國中成為笑柄。為何讓人在萬民中說『他們的神在哪裡』呢？」 \end{tabularx} \\ \\ \relax
2:18 & \begin{tabularx}{0.7\textwidth}{X} 耶和華為自己的地發熱心，憐憫他的百姓。 \end{tabularx} \\ \\ \relax
2:19 & \begin{tabularx}{0.7\textwidth}{X} 耶和華應允他的百姓說：「看哪，我要賞賜你們五穀、新酒和新的油，使你們飽足，我必不再使你們受列國的羞辱。 \end{tabularx} \\ \\ \relax
2:20 & \begin{tabularx}{0.7\textwidth}{X} 我要使北方來的隊伍遠離你們，將他們趕到乾旱荒蕪之地：前隊趕入東海，後隊趕入西海；臭氣上升，惡臭騰空。耶和華果然行了大事！ \end{tabularx} \\ \\ \relax
2:21 & \begin{tabularx}{0.7\textwidth}{X} 「土地啊，不要懼怕，要歡喜快樂，因為耶和華行了大事。 \end{tabularx} \\ \\ \relax
2:22 & \begin{tabularx}{0.7\textwidth}{X} 田野的走獸啊，不要懼怕，因為曠野的草已生長，樹木結果，無花果樹、葡萄樹也都效力。 \end{tabularx} \\ \\ \relax
2:23 & \begin{tabularx}{0.7\textwidth}{X} 「錫安的民哪，你們要歡喜，要因耶和華－你們的神快樂；因他賞賜你們合宜的秋雨，為你們降下甘霖，秋雨和春雨，和先前一樣。 \end{tabularx} \\ \\ \relax
2:24 & \begin{tabularx}{0.7\textwidth}{X} 「禾場充滿五穀，池中漫溢新酒和新的油。 \end{tabularx} \\ \\ \relax
2:25 & \begin{tabularx}{0.7\textwidth}{X} 我差遣到你們中間的大軍隊，就是蝗蟲、蝻子、螞蚱、剪蟲，那些年間所吃的，我要補還給你們。 \end{tabularx} \\ \\ \relax
2:26 & \begin{tabularx}{0.7\textwidth}{X} 「你們必吃得飽足，讚美耶和華－你們神的名，他為你們行了奇妙的事。我的百姓不致羞愧，直到永遠。 \end{tabularx} \\ \\ \relax
2:27 & \begin{tabularx}{0.7\textwidth}{X} 你們必知道我是在以色列中，又知道我是耶和華－你們的神，沒有別的。我的百姓不致羞愧，直到永遠。」 \end{tabularx} \\ \\ \relax
2:28 & \begin{tabularx}{0.7\textwidth}{X} 「以後，我要將我的靈澆灌凡有血肉之軀的。你們的兒女要說預言，你們的老人要做異夢，你們的少年要見異象。 \end{tabularx} \\ \\ \relax
2:29 & \begin{tabularx}{0.7\textwidth}{X} 在那些日子，我要將我的靈澆灌我的僕人和婢女。 \end{tabularx} \\ \\ \relax
2:30 & \begin{tabularx}{0.7\textwidth}{X} 「我要在天上地下顯出奇事，有血，有火，有煙柱。 \end{tabularx} \\ \\ \relax
2:31 & \begin{tabularx}{0.7\textwidth}{X} 太陽要變為黑暗，月亮要變為血，這都在耶和華大而可畏的日子未到以前。 \end{tabularx} \\ \\ \relax
2:32 & \begin{tabularx}{0.7\textwidth}{X} 那時，凡求告耶和華名的就必得救；因為照耶和華所說的，在錫安山，在耶路撒冷將有逃脫的人。凡耶和華所召的，都在餘民之列。」 \end{tabularx} \\ \\
[1ex]
\hline
\hline
\end{longtable}
$^{1}$提醒大家, 如果您有問題,請留言給我們,我們會盡力幫助您..
之前,我都有一個習慣,就是寫靈修分享的經文給會友,已經有三年多了..
逢星期一,三,五,我都會透過WhatsApp轉發給會友..
不過,自從三月患病以來,就停止了,因為沒有那麼多精神每次寫一千字的色經分享..
如果大家想看我的靈修分享,可以透過教會網頁,在牧者之言,就可見到..
分開新約和舊約,總共有40卷書..
雖然我沒有寫靈修分享,但沒有停止我對聖經的研究..
這段時間,我仍在讀《十二小先知書》,暫時已經完成了十卷.還有兩卷,整個《小先知書》系列已經全部完成了..
所以今天,就以《十二小仙書》的約耳書,和大家分享神的話語,我們一起同心祈禱..
讓我們透過聖經,了解你的屬性更多,更明白你的旨意..
求主讓我們能夠將你的話語藏在心裡,成為我們屬靈生命成長的土壤..
好叫我們的生命能夠茁壯地成長,成為一個合你心意的僕人,被你使用,使你的名被高舉,得到當德的榮耀..
我們祈禱,是奉主耶穌基督,聖名祈求.阿們..
我們首先了解一下約耳書的背景.從約耳書的內容,我們看不到書卷的寫作日期..
因為書卷內沒有提供任何君王的名字,也沒有指明快要壓境的大軍是哪一個軍隊..
內容只是提到猶大,不過並未提到.因此從內政就看不到約耳書的寫作日期了..
不過我們從內政可以知道,約耳書受眾受到威脅,先知約耳鼓勵他們悔改,並相信神最終會拯救他們..
約耳書一開始就描述了一個既可怕又充滿災難性的情景,一場蝗蟲的災難正臨到耶路撒冷,以及周邊地區..
世道在耶路撒冷和周邊地區的猶太人不得不懇切地請求耶和華的寬恕和拯救..
耶和華透過先知宣告,應許將神的靈傾瀉在所有生命上.約耳也預言,耶和華將審判所有國家,並建立永恆的公義..
因此約耳書的信息十分明確,唯有通過悔改才能避免災難,將審判轉化為憐憫..
約耳書有另一個主題,就是「耶和華的日子」,這個字眼在約耳書出現了五次.而在十二小先知書中的阿摩斯書,俄巴底亞書,西班牙書,撒迦利亞書都有提及..
甚至在大先知書,以賽亞書,以西結書以及耶利米,哀戈都有提及.可見「耶和華的日子」是十二小仙書書的一個重要主題..
「耶和華的日子」是好還是壞?是審判還是復興?在於他們的信仰是怎樣?對於在艱難之中,持定信仰,堅守崗位,縱有危險困惑卻不放棄的人來說,「耶和華的日子」是一個復興的日子,給予他們在困境中的盼望..
然而,倘若他們仍然醉生夢死,只顧一己的日落,「耶和華的日子」就是一個審判的日子,難以避免災禍的臨到..
我們開始進入經文,約翊書在第一章開始就描述到當地有「蝗蟲之災」,分別有「蟬蟲」,「蝗蟲」,「蚺子」,「馬乍」..
可以說是指蝗蟲生長過程的四個階段,或是不同蝗蟲的品種,每一種都比上一種更具殺傷力..
「蟲禍毀壞了當地農作物」,一章七節說,「它毀壞了我的葡萄樹,撕裂了我的無花果樹,剝乾又凋棄,似枝條露白.」.
一章十一至十二節更詳細,「農夫啊,你們要慚愧啊,收整葡萄園的,你們要哀哭啊,因為大麥,小麥與田間的莊稼全部都毀了,.
葡萄樹枯乾,無花果樹衰殘,石榴樹,棕樹,蘋果樹,田野一切樹木都乾枯,眾人的喜樂盡都消逝.」.
一片荒涼之境..
對於經文所描述「蝗蟲」,《色經學》只有三種解釋的向導:.
第一,這是真實的災禍.蝗蟲入侵是一個真實事件..
不過蝗蟲的災禍在巴勒斯坦地區是一件很經常會發生的事..
很多時蝗蟲會在春天出現,席捲北飛,進入埃及,再穿過西奈半島,吞噬以色列地區..
一般學者認為,單一蝗蟲不足以令以色列民有深刻的反思,更加不會有認罪悔改,歸回神的心思意念..
第二,這是一個隱喻.蝗蟲是代表敵人的軍隊,.
而當時是指亞述或巴比倫軍隊對猶大地進行全面攻擊..
以色列民面對敵人的入侵,當然會向耶和華呼求,.
但呼求的目的是要耶和華消滅敵人,卻不會反思自己的罪,也達不到耶和華真正的心意..

$^{41}$第三,這是指耶和華的大軍為審判以色列民而來,.
是對以色列民背叛耶和華的懲罰.在舊約,不論是摩西或所羅門的時代,.
耶和華都曾預言過這種形式的審判,只是以色列民聽得太多,無視先知的警告..
先知著耳,著著蝗蟲隱喻耶和華大軍要對猶大的背叛進行審判,.
要以色列民反思自己的罪,從而悔改,歸向耶和華..
因此,蝗禍是一種耶和華對猶大施行審判的隱喻,讓以色列民能徹底反思並認罪悔改..
先知在一章十五節向以色列人發出警告,並首次提到耶和華的日子,.
「哀哉這日子,因為耶和華的日子臨近了,好在毀滅全能者來到.」.
先知呼籲眾民,農夫,祭司都要哀號,祭司更要速上麻布痛哭,於是全以色列人向耶和華呼求..
一章十四節說:「你們要使禁食的日子分別為聖,宣告嚴肅會,召集長老和他們所有的居民來到耶和華你們神的殿,向耶和華哀求.」.
耶和華在二章一節回應:「你們要在石安,崔島,在聖山發出警報,你們所有的居民要發戰,因為耶和華的日子快到,已經臨近了.」.
這是《約珥書》第二次提到耶和華的日子,已經臨近了..
對於這個日子來臨,人的回應只有發戰,除此之外,人還能做些什麼呢?.
當先知約爾在二章十至十一節第三次提到耶和華的日子時,先知形容當時是「地動天搖,日月昏暗,星宿無光.」.
他認為耶和華的日子大而可畏,誰能擔得起呢?.
不但以色列民面對耶和華的日子不能擔當得起,其實普天下的人都不能在耶和華的日子中站立得住..
以色列民面對大而可畏的日子,應該如何自處?.
因此,對於耶和華的日子來臨,耶和華在二章十二至十三節提醒以色列:.
「你們現在要禁食,哭泣,哀號,一心歸向我,並且要撕裂心腸,不要撕裂衣服,歸向耶和華你們的神.」.
對於耶和華定義要施行審判,唯一能夠逃避的方法就是一心歸回神的身邊..
要歸回耶和華的具體行動就是認罪悔改,而禁食,哭泣,哀號,甚至撕裂衣服都只是外表的行為,真正悔改的態度是要撕裂心腸..
衣服可以撕裂,但心腸要如何撕裂呢?撕裂心腸就是要徹底否認自我,真正的悔改..
如果內心沒有真正的悔改,只是禮節性的哀悼,或者只是外在的禁食,撕裂衣服都是毫無意義的,.
因為耶和華要求的是真正和衷心的悔改,而不是禮儀性的回應..
在現實生活中,我們如何看到一個真正悔改歸向神的行為呢?.
在舊約申命記六章五節,摩西這樣說:「你要盡心盡性盡力愛耶和華你的神.」.
在新約馬可福音十二章二十九至三十一節,耶穌又這樣說:.
「第一是以色列,你要聽:主,我們的神,是獨一的主.你要盡心盡性盡義盡力愛主你的神.」.
「第二是要愛倫如己,再沒有比這兩條界面更大的了.」.
耶穌和摩西同樣要求屬神的子民不要以教條或外在的形式去回應神,而是要用心靈和誠實去回應神..
「盡心盡性盡義盡力」都不是說守律法,或者每天要祈多少節土,要奉獻多少錢,或者參與多少事奉,.
而是你是否將主耶穌放在你生命的首位,你的心思意念是否以主為先為大..
如果以色列民族能夠把先者警告和提醒,真正的認罪悔改,歸回耶和華的身邊,.
耶和華在二章十三節對他們的應許說:「因為他又因為有憐憫,不輕而發怒,有豐盛的慈愛,並且會改變心意,不降那個災難.」.
當以色列民願意認罪悔改的時候,神不單單會收回對他們的審判,.
並且更在耶和華日子未到之前,向整個民族賜下應許..
二章二十八至二十九節這樣說:「以後,我要將我的靈囂觀凡有血肉之軀的..
你們的兒女要說預言,你們的老人要做義務,你們的少年要見義將..
在那些日子,我要將我的靈囂觀我的僕人和婢女.」.
神的應許很豐富和很特別,超出了以色列民的認知..

$^{81}$首先,神的靈囂觀在凡有血肉之軀,什麼叫凡有血肉之軀?.
即是所有生命棄屍的人類,所有民族,也就是說,應許不單單輸給當時的以色列民族,.
也包括當時所有的異邦民族,包括了以色列的敵人..
第二,涉及的應許身份也很廣泛,包括兒女,老人,少年人,僕人和婢女..
兒女是指包括了男性和女性,男和女都包括在內,對於當時的女性來說是一個突破,.
因為女性在當時的猶太社群只是男性的一個財產,是附屬品,沒有地位,沒有身份..
老人固然是當時社會的領袖,少年人是社會的未來,.
但應許最劃時代的,是涉及了僕人和婢女,他們是指當時男性和女性的奴僕,.
雖然身份低微,但同樣在應許範圍裡打破了當時的社會階級之分..
這個應許就好像加拿大書三章28節那樣說:不再分猶太人或希臘人,不再分為老的,自主的,不再分男的,女的,因為你們在基督耶穌裡都成為一了..
由於這個應許範圍這麼大,涉及的人物這麼多,不禁令人提出一個疑問,究竟耶和華的日子是指什麼呢?.
是指對以色列整個民族的審判,還是指對人類的審判?.
如果是指對以色列的審判的話,那麼這個審判應該在南北角都滅亡的時候就已經結束了..
但如果是指對人類的審判的話,耶和華的日子就仍然未到了..
今天我們仍然處身在耶和華的日子來臨之前,也就是說,這個應許不單單臨到以色列,也臨到我們這批外邦人..
在《約珥書》2:10提到,耶和華大而可畏的日子臨到的時候,就會出現地動山搖,日月昏暗,星宿無光..
而在《約珥書》3:15-16再次提到,日月昏暗,星宿無光,天地震動的時候,.
就是指耶和華要審判埃及以東的時候,證明耶和華的日子並不是單一審判以色列的日子,而是審判萬國萬民的日子..
在《新約》,使徒彼得在《紹行傳》2:16-21就引用了《約珥書》2:28-32經文,.
來講明當時在五旬節聖靈降臨在耶路撒冷的門徒身上,而經歷這次聖靈降臨的人,.
包括了帕提亞人,馬代人,以蘭人和住在美索不大米亞,猶大,加帕多加,本都,亞細亞,.
佛雷加,龐菲利亞,埃及的人,並且靠近古里內的利比亞一帶地方的人,僑居的羅馬人,.
印證了《約珥書》2:28-32,聖靈囂冠凡有血肉之軀的人身上,同時也印證了2:32-34,.
那時凡求告耶和華的名就必得救,因為一切相信主耶穌基督是神的兒子的人,都必得救..
另外,史圖保羅在《羅馬書》10:13也印用了《約珥書》2:32,因為凡求告主名的就必得救,.
引證了猶太人和希臘人並沒有分別,因為人人都有同一位主,他也厚待求告他的每一個人..
所以《約珥書》宣告凡有血肉之軀的,無論是男的,女的,只要是社會一員,甚至只是奴隸身份,只要他們承認耶穌是主,就必得救..
不但得救,還會領受異象,異夢..
弟兄姊妹,從《約珥書》裡我們知道,在主刑未審判我們之前,我們需要把握回轉歸向神的機會,.
我們不單可以得到救恩,還會領受異象..
這異象就是讓世人知道,只有相信主耶穌,承認祂是生命的主,拯救的主,就必得救..
在耶和華大而可謂的審判日子來到的時候,所有相信耶穌的人都能夠永遠與祂同在,超越今生,直到永恆..
今日講題是「審判」或「復興」..
當我們仍然活在罪中,享受罪中之樂,不知悔改,或只是虛情假意的悔改,.
每個星期過的一個自我滿足的宗教活動,只是形式地撕裂衣服,沒有真正的撕裂心腸,.
縱使我們今天過得多麼滿足和開心,多麼悠然自得,.
但等待我們的卻是神那大而可謂的審判..
聖經如何形容這一批站在審判台前等著主耶穌才主的人?.
馬逝福音說:「在那裡,他要哀哭切齒了.」.
我深信這不是我們相信主耶穌的終極盼望,.

$^{121}$不過如果我們真誠地歸回神的身邊,破除自己的老我,願以順服主的帶領,.
盡心,盡性,盡義,盡力愛我們的主,又愛倫如己,.
感恩的心與生命見證在生活中彼此流露,.
我們的生命就會被復興,被聖靈囂觀,領受異象,成為主合容的僕人侍女,被主使用,.
主的聖名會因著我們的生命而被高舉,而我們的生命就能夠榮神益人..
今日是紅恩堂十二週年的堂慶日子,.
這個感恩的日子固然值得我們歡喜快樂,.
因為有很多快樂的回憶,有很多歡樂的片段,.
不過開心過後,我仍要問自己,.
在我的信仰旅程中,我有沒有被聖靈囂觀呢?.
我領受了一個什麼異象呢?.
抑或我每一天都因循例過我的信仰生活,.
對親近神毫不重視,對聖人的提醒毫不敏銳,.
甚至連神的話語都毫不追求,.
如果我們的屬靈質素只是一個吃奶的嬰孩,.
那又怎能夠領受異象,承擔使命呢?.
另外作為教會一份子,我們同為基督的身體,.
我們要問教會,教會的異象是什麼呢?.
是講傳福音?是關心區內的街坊?.
是要成為福音的平台?.
那教會有什麼具體的設施,或者具體的計劃,.
去執行神所交託給教會的異象呢?.
最後我想和大家分享一個木頭車的見證,.
在1902年,美國有一個十歲的小女孩,.
有一天她聽到父母的一個宣教士朋友講道,.
這位宣教士剛剛從中國回到美國,.
他分享道當日在中國服侍的時候,.
有一個寒冷的早上,有一個人開著一輛空的木頭車,.
趕了兩天路來到他們的報道場所,.
原來這個人從鄉下知道有些洋人帶來了什麼福音,.
他趁著春節還沒有到,還有空閒的時候,.
就遠離邀請洋人分享福音,.
他的木頭車就用來接載願意跟他回去講福音的洋人,.
可惜一班宣教士都非常忙碌,.
沒有人能夠抽空和這個中國農民到他們村莊裡面,.
這個中國人苦苦哀求,最後無奈地還是拉著木頭車離開了..
第二年初冬,這個農夫又帶著他的木頭車來了,.
很可惜,同樣得到相同的答案,.
那些宣教士實在太忙碌了,沒有人能夠跟他回鄉下,.
最後又再一次推著那個木頭車離開了..

$^{161}$當這個小女孩坐在禮拜堂聽著這個真實的故事的時候,.
忽然間她完全忘記了一切,.
還在她腦海裡就只有一幕,在中國有一條鄉村的人很想聽福音,.
但只看到一輛空的木頭車..
那個女孩心想,如果是我,我一定會去的,.
不過她想到自己只不過是一個十歲的小朋友,.
於是就向神獻上自己為祭,.
主啊,等我長大之後,我一定去..
小女孩願意回應神的呼召,.
似乎找到她生命的目標和方向了..
這個女孩漸漸長大了,.
學校生活將她帶進多姿多彩的世界,.
她開始有夢想了,就是要成為一個傑出的商人..
中國平原上那輛空的木頭車漸漸在她腦海裡消失了,.
甚至直至她在史丹佛大學讀二年級的時候,.
有一天她收到母親的信,.
她母親建議,讀萬卷書不如行萬里路,.
不如去中國走一趟,好不好?.
奇怪了,那輛早已經在腦海裡消失的木頭車,.
這一刻又再清晰地浮現在眼前了..
整個下午,這個女孩都很掙扎,.
她不能夠自制地不停地哭,.
她很後悔,她後悔小時候將自己獻上,.
她坦白地跟神說,.
主啊,雖然我答應過長大之後去中國,.
但我現在真的不想去了..
她經過痛苦的掙扎,終於做了一個決定,.
為了要證實去中國宣教,.
是自己少年時期的浪漫情懷,.
因為真是神的心意,.
她決定去一次中國,.
踏足在中國的土地,.
如果到時真的有感動的,.
就證明是神的心意,.
否則就回到美國重傷了..
在貨船駛往中國的途中,.
有一次在太平洋遇上大風浪,.
船上有人說,.
嘩,風浪很厲害,.
不會是我們中間有一個藥籃吧?.

$^{201}$這個女孩不敢插嘴,.
她感到很害怕,.
如果到時真的要拋棄貨物下海的時候,.
自己表不表態呢?.
她獨自站在船上的甲板上跟自己說,.
看看以後還敢不敢做這些傻事吧?.
然後她堅定地跟自己說,.
我要回到美國,.
做一個成功的商人..
當她這樣說的時候,.
很奇怪,.
她竟然感到這番話,.
是多麼的枯燥無味,.
更奇怪的是,.
這時候她心裡突然浮起一句話,.
為什麼不讓神為你選擇呢?.
很奇妙,.
她如釋重負,.
情不自禁地大叫,.
主啊,我願意信服,.
求你顯明你的旨意吧!.
當她說完這句話,.
一陣不能言喻的平安和喜樂,.
好像潮水般湧流在她的心頭,.
這種奇妙的平安是從何而來呢?.
就是從她對主說,.
我願意信服,.
求主向者顯明旨意的時候,.
當她願意信服的一刻,.
神的話語就進入了她的生命裡,.
這個小女孩體驗到與神相交的經歷,.
她的心靈得到滿足..
這個美國女孩原本是很勉強地去中國旅行,.
但當她能夠回應神說,.
主啊,我願意信服,.
求你顯明你的旨意吧!.
神就在旅程裡,.
她感受到上帝所愛..
這位搭上中國的女孩,.
就是宣道會希伯倫堂的創立人,.

$^{241}$何二師宣教士,.
在中國被稱為何二姑,.
她在第一次世界大戰爆發之後,.
當時她只有22歲,.
沒有讀過什麼神學,.
就隻身前往中國的官山全道,.
經歷貧窮和戰亂,.
在廣東省南海縣一個小鎮,.
建立了希伯倫教會,.
又在附近的鄉鎮建立了福音堂,.
孤兒院,婦女織造廠,.
她毫無保留地奉獻了自己,.
到80歲才退休..
我們很熟悉的李思敬博士,.
他的父親李飛吾牧師,.
就是何二師宣教士在中國收留的孤兒之一,.
所以李思敬博士說,.
如果何二姑沒有去中國,.
她的母親不肯將女兒獻上,.
今天又怎會有我呢?.
何二師宣教士為中國人的靈魂得救,.
事奉了58年,.
她一生的故事寫在一本書中,.
書名就是《誰掌管明天》,.
神秘的意象,.
令她離開本地本家,.
而她與神建立的親密關係,.
令她走得更遠,.
更持久,.
更喜樂..
甚願我們洪恩家的肢體,.
為了能夠領受神的意象,.
為主打一場美好的聖像,.
我們一起祈禱:.
「親愛的恩主,.
您又因為有憐憫,.
不輕易發牢,.
有豐盛的慈愛,.
甚願我們每一個都活在主的恩典當中,.
您求聖靈時刻提醒我們,.

$^{281}$要時刻親近您,.
要聆聽您的話語,.
要專心依靠您,.
好讓我們每時每刻都感受到.
在您的愛裡面,.
活出一個豐盛的生命.」.
我們祈禱:.
「奉教主耶穌基督,.
得勝明節祈求.」.
阿們..
願主賜福給大家..
\newpage



\section{馬太福音 11:25-30}
\label{sec:EtvLUaUHtqc}
\textbf{2023.11.12 《你們要來 , 在主裡得安息》 馬太福音11\_25-30 (和修版) 講員 : 高淑芬傳道}
\newline
\newline
連結: \href{https://youtube.com/watch?v=EtvLUaUHtqc}{\texttt{https://youtube.com/watch?v=EtvLUaUHtqc}} ~~~~ 語音日期: 2023-11-15
\newline
\newline
\hyperref[sec:VjKjBj0FEho]{\small{< < < PREV SERMON < < <}}
~
\hyperref[sec:index_chronic]{\small{[返順時目]}}
~
\hyperref[sec:index_scriptual]{\small{[返順卷目]}}
~
\hyperref[sec:xQAERVXii10]{\small{> > > NEXT SERMON > > >}}
\newline
\newline
馬太福音 11:25-30
\newline
\begin{longtable}{cl}
\hline
\hline
章節 & 經文 (和合本修訂版)\\
\hline
11:25 & \begin{tabularx}{0.7\textwidth}{X} 那時，耶穌說：「父啊，天地的主，我感謝你！因為你把這些事向聰明智慧的人隱藏起來，而向嬰孩啟示出來。 \end{tabularx} \\ \\ \relax
11:26 & \begin{tabularx}{0.7\textwidth}{X} 父啊，是的，因為你的美意本是如此。 \end{tabularx} \\ \\ \relax
11:27 & \begin{tabularx}{0.7\textwidth}{X} 一切都是我父交給我的；除了父，沒有人知道子；除了子和子所願意啟示的人，沒有人知道父。 \end{tabularx} \\ \\ \relax
11:28 & \begin{tabularx}{0.7\textwidth}{X} 凡勞苦擔重擔的人都到我這裡來，我要使你們得安息。 \end{tabularx} \\ \\ \relax
11:29 & \begin{tabularx}{0.7\textwidth}{X} 我心裡柔和謙卑，你們當負我的軛，向我學習；這樣，你們的心靈就必得安息。 \end{tabularx} \\ \\ \relax
11:30 & \begin{tabularx}{0.7\textwidth}{X} 因為我的軛是容易的，我的擔子是輕省的。」 \end{tabularx} \\ \\
[1ex]
\hline
\hline
\end{longtable}
$^{1}$(日語).
因為今天有泰國朋友來到我們當中.
這次我講的主題是.
我們要來到主廟裡得享安息.
這是一個續集.
還記得這個泰文字怎麼讀嗎?.
泰文的弟兄姊妹.
「閉」.
去就是閉.
上次結語就是說.
如果我們不去好好事奉神.
那又怎樣?中文就是閉.
所以我們一定要去.
記得要去.
但是面對教會.
或者我們自己說去洗萬民.
作為門徒.
其實這個大使命.
頒布給你和我.
你有沒有覺得上帝非常看重我們?.
自己成為門徒.
你覺得容易嗎?.
到今天我還在學.
屢拜屢戰.
屢戰又屢拜.
有時會夾雜著這樣的情況.
所以我們說還要去洗萬民.
做主的門徒.
談何容易呢?.
但是的確上帝是看重我們每一個人.
這是上帝的命令.
祂看得起我們.
祂買重了我們.
要我們去.
要完成這個大使命.
所以這個大使命.
當我們一看它的時候.
特別是在《馬太福音》裡面.
你就要記得來.
有去要有來.

$^{41}$來的泰文就是"馬".
"馬"就是來的意思.
這是當時主耶穌說.
你們繼續要做我的門徒.
秘訣就是要來.
來到我這一處.
去得到安息.
去經歷一下.
做門徒不是綁手綁腳的.
有時我們會覺得.
做了基督徒好像有很多事要做.
又好像有很多事不要做.
有些事要做又好像做不來.
有些事不要做又做了.
有時都會在這些處境裡面.
但是當我們去看這個大使命.
原來祂的秘訣是在.
十一章二十五到三十節裡面.
在當時的時候.
耶穌的事奉.
和施洗約翰的事奉.
是落在一個很大很大的低潮裡面.
如果我們看十一章.
你會看到施洗約翰.
當時已經被抓去坐牢.
他去問耶穌一句話.
你是否我要等來的主呢?.
如果你回家有機會去看.
主耶穌不是回答.
是,你不用再等第二個了.
因為施洗約翰問了他兩條問題.
兩條問題其實我們叫關閉式問題.
你是否那位主呢?.
是,主耶穌可以回答.
是的,因為祂真的是那位主.
你是不是要我.
如果不是我需要再去等第二個嗎?.
不用,因為我就是主耶穌.
主耶穌不是這樣回答他.
祂說你去看一看.

$^{81}$你去聽一聽.
你會看到很多人.
本來聾的只能聽見.
盲眼的可以見到.
死了的都可以復活.
還有窮人他們都有福音可以聽.
但是在那時候的施洗約翰.
他落在一個非常低潮的當中.
你的人生有沒有遇過低潮?.
如果只有幾位點頭.
我就有經歷過低潮.
如果你沒有經歷過低潮.
非常恭喜你.
我覺得每一個人都有高有低.
在生命裡面會遇到高的浪.
低的谷.
都有很多不容易的情況.
當時連施洗約翰.
是一位為主耶穌開路的先知.
他也將要被斬頭.
落在一個極不容易的時間.
而主耶穌回答他的時候.
說了這麼多的神蹟.
這麼多奇妙的事發生.
但是信祂的人.
真的跟從祂做門徒的人.
不是這麼多.
當你再想想.
當主耶穌被釘十架.
走那條苦路之先的時候.
祂的門徒有十一個.
一個就賣了.
剩下十一個.
那十一個怎麼辦?.
真的要走了.
去了別處.
人實在是會有軟弱.
所以今天我們自己.
我們自己也有自己本身的神學.
有時候會覺得.

$^{121}$傳道人,牧師,宣教士.
他們真的會靈明高昭.
會勁人一等.
但是我要告訴你.
其實我和你一樣都是人.
我們也會有軟弱的時間.
但是在我們健康的時候.
今天你還是聽主的話的時候.
你就真的要為自己.
多打一些定心針.
就是讓自己知道.
就算一個極大的低潮也好.
我們可以繼續去做門徒.
做門徒的意思就是說.
跟從耶穌的腳步.
學做祂.
所以就是要很好的.
來到主耶穌那裡跟著祂.
在十一章25到26節的時候.
當主耶穌.
在十一章的開頭.
當施洗約翰.
來找他的徒弟.
或者可以叫做門徒.
去問耶穌這些話的時候.
當時主耶穌.
如果你在十一章第一節.
就會看到.
主耶穌剛剛踩那十二個門徒.
出去傳福音.
去讓人知道.
主耶穌就是那位彌賽亞.
就是那位救世主.
但是施洗約翰.
竟然去問了一個.
他當時都很低落.
都很信心動搖的時候.
他很想清楚知道.
我一輩子去事奉的那位.
有沒有錯呢?.

$^{161}$他很緊張.
如果沒有錯.
我相信他不怕死.
面對著死他都不會懼怕.
在回答了施洗約翰的門徒之後.
主耶穌其實繼續.
他有說到施洗約翰.
他的事奉是蒙神喜悅的.
縱然他得到的.
將會得到的不是那麼好.
我們看聖經就知道.
他的下文是怎樣.
但是主耶穌.
仍然在讚賞施洗約翰.
接著二十五到二十六節.
是耶穌的一個祈禱.
其實當時的情況.
都不是那麼樂觀.
但是主耶穌.
他仍然回到天父的面前.
他說天地的主.
你就會看到.
他就像後期保羅書信裡.
時刻都叫我們.
我們要凡事感恩.
在失落的時候.
在低潮的裡面.
主耶穌很愛施洗約翰.
我相信他也知道.
施洗約翰.
接著下去的下場很慘.
要被斬頭.
但是他為一切的事情.
他說因為你將這些事.
向聰明通達的人藏起來.
向嬰孩顯出來.
父啊.
你的美意就是這樣.
我們讀到這裡的時候.
很多時候我們都不明白.

$^{201}$為什麼會這樣去祈禱.
你會看到當面對困難的時候.
主耶穌仍然去感恩.
去相信全福音的果效.
不是在他自己那裡.
一切是在父神的手中.
一切都是神的美意.
我們今天.
我們自己做門徒.
領別人去做門徒.
那個效果是怎樣呢?.
上個星期.
我們要加椅子.
一百多人.
今天少了一點.
這個不是我們可以控制得到的.
我們多人少人.
你會看到我們的師士.
招待.
IT部.
所有事奉人員.
我們的師琴.
我們的師班員.
領唱的.
主席.
每個人都是忠於.
那個事奉.
好好地去預備.
所以我們求神幫助我們.
當我們去為上帝.
做每一件事的時候.
結果未必是我們的意思.
但是神給我們的份.
我們就好好地做了.
因為有著上帝的美意.
而且也有人的自由選擇.
那時候耶穌說.
我想在我們來說.
我們都是要操練的.
很不容易的時候.

$^{241}$有些人會說.
「是不是阿Q精神?」.
「是不是自我安慰?」.
有時候被人說得多.
我們自己都會.
不夠信心的話.
都會怯一怯.
但是.
現在我們每一個在這裡的時候.
我們要告訴大家.
我們的主.
仍然是值得我們.
好像剛才一敬拜上帝的時候.
我們的主是值得我們去敬拜的.
敬拜的原文裡.
就是有事奉的意思.
是值得我們去事奉.
值得我們去跟從的上帝.
求主給我們.
好像耶穌一樣.
有著很重要的軸心.
去.
所以你們要去.
「使萬民做我的門徒」.
「奉父子聖靈的名」.
「為他們施洗」.
是不容易的事.
甚至我們想起來.
我們去傳福音給其他人的時候.
為什麼呢?.
他明明比我好.
讀書又比我好.
為什麼跟他傳福音傳了這麼久.
他還是不相信呢?.
有時候都會覺得.
自己比這個人不如.
但自己卻得到這麼大的恩典.
實在是很感恩的事.
所以去這個大使命.
實在是很大.

$^{281}$所以今天若我們.
不去到上帝面前的時候.
將一切.
好像二十七字一樣.
一切所有的.
都是交付給我的.
除了付沒有人知道的紙.
也都願意提示的.
沒有人知道付.
我認識一個傳道人.
他做傳道都有二十年了.
由他未做傳道的時候.
他當他蒙召的時候.
他都很希望他的家人信主.
因為家人只有他一個信主.
到他做傳道了.
現在都退休了.
家人都還未信主.
但他仍然繼續在家裡.
好好地做見證.
所以很多事我們是不知道的.
但我們都要知道.
一切的事.
有著上帝的美意的時候.
但他所託付給我們的.
我們就求神給力量我們去做.
剛才的大使命.
也是不容易的.
大使命也是當主耶穌受苦.
受死.
甚至死了之後.
所頒布給我們的大使命.
而那些門徒.
那個環境是非常惡劣.
比起施洗約翰.
將要斬頭這件事.
更加低落.
但竟然.
耶穌不是不知道.
那十一個門徒都離世了.

$^{321}$還要靠他們.
但耶穌就是這樣.
將這麼大的使命交給我們.
所以我們更加面對著.
這麼大的使命.
而上帝這麼恩寵我們.
其實我們聽.
我從信主那一刻.
已經聽很多傳道人說.
其實上帝可以不用我們.
去叫人信主.
叮一叮而已.
我們好像看卡通片.
有支魔術棒.
叮一叮.
你就信了.
但主耶穌不是這樣.
耶和華神不是這樣.
亞伯拉罕的神.
我們剛才唱.
這麼古老.
好像不可能.
這樣也可以生下兒子.
還要說你的子孫.
子孫是很多很多.
這些簡直是.
人不能預料得到.
但上帝就是這樣認去.
今天給我們的時候.
我們求神捉住上帝.
捉住主耶穌.
來認識這個身份.
我們要認定除了父之外.
耶穌一切都是以.
父得到這個肯定來滿足.
我們今天生活的時候.
我們有沒有一早的時候.
為自己起床.
可以呼吸.
可以有生命氣息.

$^{361}$還能夠走.
甚至很平安.
回到這裡.
我今早回來的時候.
地上還濕.
天是暗暗的.
是灰灰的.
想不到我坐著坐著.
回到教會窗外看.
已經是藍天白雲.
有太陽.
你有沒有為這些事情.
很平常.
你每天都是這樣.
來感恩.
求神給我們.
是因著上帝我們得到滿足.
並且當我困難來到的時候.
不是叫你笑傲江湖.
不過是合乎中道.
笑為什麼要加一個笑呢.
因為要喜樂地面對.
縱然在困難當中.
大家都有聽過健康運動.
現在叫做呵呵笑.
是做運動的時候.
要你呵呵笑.
要嘴形是笑.
原來這樣的動作.
可能當時你不是很開心.
但這個動作固定你要這樣運作的時候.
你做的時候.
原來專家說.
也有這個賀爾蒙.
幫你產生這種歡悅的感覺.
所以我們就很明白.
為什麼保羅在坐牢.
也是要我們常常喜樂.
我們不要因為人的言語.
照單傳修.

$^{401}$因為人去看人.
總是有些位置看不到.
例如我早前剛見了一位.
很久沒見的一位.
當時初信主認識的一位姐妹.
見到她.
因為過去我一直都是長頭髮.
她說你剪了頭髮.
還可以.
我說對不起.
剪頭髮是對的.
我的頭髮沒有電.
因為那天爸爸也出席.
我的頭髮就像爸爸自己捲的.
她好像不相信.
還要捉著我的頭髮.
真的自己捲.
所以這些是很小事.
被人誤會你電腦髮很小事.
但是很多時候.
我們一起去.
或者你的家庭.
家人.
甚至伴侶.
坐在你旁邊.
有些事他可能不知道.
未看到那一面.
所以當人的言語.
如果是讚賞你也好.
批評你也好.
你也不要照單傳修.
你也要在那個時間.
自省一下.
是不是真的這麼好.
是不是真的這麼壞.
明不明白.
都是要合乎中度去看.
來.
你覺不覺得.
好像人去到十字架那裡.

$^{441}$來的秘訣就是這樣.
就是來到主耶穌那一船.
當我們說要去去去.
去的時候.
請你都記得來來來.
泰文就說.
閉閉媽媽.
來來去去.
閉閉媽媽.
要有來要有去.
又不只是去.
又不只是來.
每個時刻.
你有那個來.
來到神的面前.
也會有著去的勇氣.
接著我們會很熟悉.
凡路苦單重擔的人.
可以到我這裡來.
我使你得安息.
我心裡柔和謙卑.
你當負我的壓.
學我的揚息.
這樣你的心就得向安息.
因為我的壓是容易的.
我的擔是輕散的.
我想這一段金句.
大家都可能會背.
或者我們寫卡.
過去喜歡寫卡.
都可能會贈送給很多朋友.
在剛才這裡我們看到.
因為上帝說到.
你們要來得安息.
因為我的壓是容易的.
我們要知道.
上帝幫我們.
大家的壓知道.
他是用牛去耕田.
他要有牛的壓在頸上.

$^{481}$這個壓上帝給過我們.
給過我們的壓.
不好意思.
我把人好像代為牛一樣.
牛去耕田.
他的四隻腳就是腳.
他不能用手這樣挖泥.
他那些是提.
又不像貓.
他不是這樣做的.
他需要有人做壓.
後面做爬.
他才可以在這個.
很硬的土地裡.
去耕耕他的土地.
以至農夫可以繼續種植.
這個壓要記住.
因為上帝的壓是容易的.
容易的意思是.
這個壓是為你和我.
每一個度身訂造的.
我的壓未必適合你.
你的壓也可能不適合我.
我們首先去了解壓.
這個是套著兩隻牛的壓.
你會看到.
以前不是工廠這樣做的.
所以那些壓是要木匠.
主耶穌的父親是一位木匠.
當主耶穌出來傳道的時候.
是三十歲.
因為主耶穌也是.
遵從人當時的傳統.
就是一些夫子.
三十歲之後.
才可以出來做教導的工作.
之前是怎樣呢?.
很多做夫子的.
都是跟著父親.
耕種的就做耕種.

$^{521}$種植果園的就種植果園.
或者織布.
或者做木匠.
主耶穌很多時候.
我們在聖經裡都看到.
是木匠之子.
很可能他也是一個.
跟著約瑟的養父.
他的在世爸爸.
來做木匠.
你會看到.
是一個人的.
不是,是一隻牛的鴨.
背著一個牛.
如果你要牽著東西.
就要搭車.
讓牠可以載東西.
如果是耕地的.
就要拿泥棒.
所以你看到.
牛是很需要這隻鴨的.
如果這隻鴨鬆得太多.
牠已經掉下來.
走著走著.
沒東西牽著牠.
牠就掉下來.
發揮不了作用.
牛就做不了事.
但太緊的話.
牠會動的.
你會弄傷牠的頸.
又受傷.
所以這隻鴨.
一定要很適合牛.
牠就是要這樣.
四隻腳繼續走.
繼續耕田.
你看到.
鴨是要度身去做的.
我感受到.

$^{561}$在《馬太福音》.
十一章二十八節.
所說的鴨.
不是這款兩隻.
兩個頭的那些.
因為我看過很多.
聖經老師說.
主耶穌是一個圈.
你在這裡.
你另一個圈.
跟你一起負這個鴨.
但主耶穌不是這樣說.
在《馬太福音》二十八章裡.
是你負的鴨.
是容易的.
你所負的鴨是輕散的.
是你.
如果我有時想.
這兩個圈.
哥林多書信教我們.
信與不信不可同乎一鴨.
就是這隻鴨.
就是兩個圈的那種鴨.
來去比喻.
如果你跟它一起向前走.
是很難的.
如果不是一起信的人.
或者不是一起有理念的人.
你去走是很難走的.
但是.
如果我們是一隻牛.
如果主人需要我.
自己去牽的.
就一個圈.
當主人需要我.
要跟著人一起去事奉的時候.
特別是大使命.
要幫他們施洗.
又要教導他們.
又要出去傳福音.

$^{601}$其實是很多個部門.
在教會裡面是很多個部門.
我們圈不止兩個.
可能一排很多個.
有些人不動.
只有一個人動.
就像教會.
做比喻.
僧姑娘一個圈動.
我們羊仔都不動.
你說它能不能走呢?.
走不了.
你想想圖畫是要一起的.
但是無論如何.
圖畫是一定要和牛的頸一起走.
怎樣一起拍檔呢?.
今天還未說到.
這個牛絨的圖畫是度身訂造的.
是輕散的.
不容易令你受傷的.
你要記住.
得享安息.
上帝告訴我們.
每一個人都有很多責任.
很多我們需要做的事情.
很多我們要面對的事情.
上帝給了一個很適合你的圈.
來讓我們走.
如果沒有這個圈.
我們做不到那些事情.
因為做不到爬地.
所以我們要求神.
去到上帝那裡.
去求上帝給我們一個容易的圈.
不是說容易.
是不受傷害.
例如上次說的.
祈禱會.
你是否願意去呢?.
報道去探訪.

$^{641}$現在就好了.
天氣涼快.
多去探訪兩家都不要緊.
但是熱的時候.
你只出紅恩堂門口.
已經濕透了.
還要濕漉漉去探訪.
或者去派單張.
是不容易的.
我是很被動.
我怎樣跟別人打招呼.
可能也是不容易的.
但是.
來吧.
來到上帝那裡.
去求力量.
回到上帝那裡.
更加做一個大使命.
使萬民成為主耶穌的門徒.
談何容易呢?.
讓我們可以這樣做.
這個門徒.
我再次說.
是向主耶穌學習.
主耶穌在這個不容易的時間裡.
祂還是回到耶和華神當中.
所以我們求神.
讓我們緊密地跟隨主耶穌.
接受主耶穌的管教.
遵循主耶穌所教導的真理去走.
很簡單.
也是不容易的.
我們也不能控制自己.
不要生病.
控制不了.
控制不了又多了一根皺紋.
又多了一條白頭髮.
這些事情我們都做不了.
我們唯有憑著信心.
盼望和愛心.

$^{681}$來到上帝的面前.
來到祂主的當中.
泰國有一個很有趣的傳統.
就是當他家有喜事的時候.
他就會大排筵席.
泰國是一條條村.
就算現在新榮佐那些村.
好像錦繡那樣.
錦繡花園那樣.
一間間屋都好.
他在屋的門前.
他就會大排筵席.
有很多農村.
鄉村的稻穢.
很多桌子.
類似這樣.
如果地方小.
可能會拿著.
因為鄰近的村.
可以拿著一些食物.
很願意一些喜事.
例如生日或結婚.
或小朋友很健康很平安.
一歲了.
他就喜歡請別人.
和他一起歡喜快樂.
當時我不明白.
這麼累這麼辛苦.
這一年才養大他一年.
還有一段時間要熬.
在座的父母就知道.
不是一歲就搞定.
我們還要繼續的供書教學.
還要繼續和他前行.
很不容易養一個小朋友.
但在那一年.
特別第一年.
這個小朋友很健康.
沒有疾病沒有災禍.
健健康康的.

$^{721}$他們是歡喜快樂的.
忘記了日後還要走一條漫長的路.
忘記了.
他就很開心的那一刻.
或者是滿月.
有些人剛生了小朋友.
又會請人.
今天大家就能經歷到.
我們有位泰國姊妹的孫子.
滿了一歲.
他沒有整條村.
屬靈的家就只有你們這麼多.
他已經準備了食物和飲品.
稍後你可以留在這裡.
和他一起歡喜快樂.
其實真的很不容易.
像我們這樣.
我們有時也不明白這些傳統.
例如家中可能有些特別的家規.
但無論如何.
當我們在不同的情況當中.
我們要去領人.
讓他可以相信實的時候.
我們求神給我們一個秘訣.
就是回到主耶穌的樹.
回到祂的樹來享安息.
這個享安息不是安息禮拜.
不是這種.
也不是主懷那種.
是那種.
這個安息這個字.
是去到舊約裡面.
當上帝應許那些以色列民.
你們要去到我應許給你的地.
迦南地.
那裡的字就是這個安息的字.
安息.
去到迦南地.
那些以色列人是否需要去做.
要的 要耕種的 要維持的.

$^{761}$不過是一個好的地方.
繼續去維持.
有人問我們上到天堂的時候.
我們還需要這麼辛苦招九萬五嗎?.
上到天堂的時候.
我是不是還是一個香港人.
還是很喜歡上班的那一種.
很多的《釋經》的老師.
從啟示錄去看.
我們都要做的.
起碼知道要做的一件事是什麼.
唱詩 敬拜神 讚美神.
其他的我不知道怎樣.
因為《釋經》裡面沒有說.
我們說今天去休息.
休息很多香港人.
特別是三年困在香港之後.
現在很多香港人都休息.
說我要休息去旅行.
都是做的.
不是真的預訂了一間酒店.
睡在那裡這麼多天.
都是會去做的.
但是要告訴大家.
這個休息是一個很滿足的.
這個滿足是因為我們有上帝.
因為有主耶穌和我們同行.
大使命的應許就是.
主耶穌和我們同行.
而你們來.
來得安息.
得這些勞苦.
背重擔的人來.
你們會輕散.
這個應許就是.
上帝會為你們做一個適合你們的.
適合你們的額.
讓你們去處理每一件事.
去做每一個的事奉.
甚至去完成上帝給我們的大使命.

$^{801}$求神幫助我們每一位.
在我們的事奉裡.
我們累了.
有些教會.
它就很好.
不知道你們聽起來會不會好.
它就規定.
每一個人.
每一位弟兄姊妹.
只是在教會裡.
一個服侍.
沒有事奉的.
你就要找一個事奉崗位.
有很多的.
請你選一個.
固定的事奉.
不是說你這個星期做領事.
你整年就是做這個星期的領事.
不是這樣.
是一個崗位.
意思是.
教會的理念.
就是不想弟兄姊妹.
事多.
不是買東西的事多.
是事情很多.
以致到忙亂.
事多令到你忙亂.
會的.
我曾經去參加過一些.
那些靈修大師.
給了一筆巨款.
我覺得都挺貴的.
因為他又出名.
三日兩夜.
去到第一天.
一個很簡單的講解.
每天一間房.
不准和別人聊天.
你們每個回房.

$^{841}$回房做什麼?.
睡覺,休息.
第一天.
第二天.
大家很開心.
大家學了什麼.
和教務.
第一天.
這個時間你們都是回房.
回房這次讀哪一段經文?.
讀什麼?.
不是.
繼續休息.
如果下午起不了床.
繼續睡.
晚上.
去別人的營地.
有那個時間.
六點鐘就吃晚飯.
你都還沒起床.
不要緊.
不吃四十天都不會死.
何況你的房裡已經有瓶水.
夠你今天用了.
總之你睡.
睡到今天.
不要說準時出來吃午餐晚餐.
好好地睡.
睡到自然醒.
自然醒之後.
你想讀經.
你就讀.
你想和神聊天.
你就去祈禱.
你想唱詩去歌頌神.
你呢.
可以的.
不過不要太大聲.
不要開窗.
不要吵著隔壁一個就行了.

$^{881}$就是這樣.
開頭想第一天.
第二天.
還沒睡之前.
這麼貴的靈修課程.
就叫我去睡.
我不如回家睡.
不用給營地費.
不用給導師費.
但是.
當你真的睡覺.
因為去的都是教務.
睡覺之後.
你很有力地.
去看聖經.
靈社入心.
泰文有個字叫「溝濟」.
你們經常說錯.
什麼溝濟.
這麼上下音.
不過叫「溝濟」.
意思是.
那些東西入心.
入了你的心.
很明白的意思.
就是「溝濟」.
嘩!真是.
那首詩歌.
有些人.
我們教會都會.
有些聖經.
特別那個禮拜日.
如果還要講道.
都會有些聖經.
或者準備的講章.
都會帶上.
安靜的時間.
最好準備講章.
不知道為什麼.
看的聖經.

$^{921}$很入心.
就算不是要預備講章.
平時.
拿著一首歌.
平時都是這樣.
選了幾首歌.
自己去敬拜神.
唱那些歌的時候.
很入心.
是很.
不明白的.
這就是.
你們要來.
上帝的樹.
我們知道.
紅印堂有很多弟兄姊妹.
一兼很多職.
很多,但是.
去事奉的時候.
都請你來.
回到主耶穌的樹.
去享上帝.
給你的安息.
當我們.
你不要以為.
孫教授.
不要當他們是聖人.
我在泰國的時候.
為什麼搞成.
為什麼不繼續在泰國.
回到香港.
養病.
其實有時候我們事奉.
去到一個點.
你以為靠上帝.
但有時候.
靠你的經驗.
或者你的本領.
去運作.
因為很忙.

$^{961}$很多事.
如果我們沒有將時間.
好好地.
去到上帝的樹.
安靜.
去休息.
如果身體很疲倦.
靈修大師.
讓我們學到.
睡覺都很聖潔.
很聖潔.
聖潔的意思.
是好.
是睡覺.
所以今天晚上.
或者昨天晚上.
這段時間.
有沒有好好地睡.
我那時候.
我自己都不知道.
原來自己走著走著事奉的路.
沒有少了來.
沒有好好地.
去到主的樹.
來休息.
去到.
當然每一件事.
我都祈禱.
神給我力量去做.
讓這件事.
可以很順利地去做.
有的.
但是.
人原來沒有好好地.
去到上帝的樹.
去到上帝的門口.
出去的門口.
就止瀉了.
就一路由星期日運作.
所以我們都求神給我們.

$^{1001}$知道.
去當中.
我們有個來址.
就是去到上帝的裡面.
去求安息.
今天.
我告訴你.
我曾經經歷過訓導.
如果在座的弟兄姊妹.
如果你不是偶然一兩天.
明天有件大事.
很開心著藥睡不到.
而是每一天都睡不到.
但日間你仍然要去.
運作很多事情的時候.
你真的要和你家人說.
身邊的人說.
你要好好地睡一覺.
這個睡.
特別是香港人.
現在被譽為.
香港人很喜歡上班.
特別是我們.
一定下來.
覺得好像沒事做.
其實不是.
定下來.
在神的裡面.
來到上帝的裡面.
你才會去支取這個力量.
你不要讓身體.
去到燃燒.
去到負重太多.
耗盡.
耗盡再回頭.
我都要經歷了.
五年.
我才好.
現在的我.
很不同.

$^{1041}$我昨晚還沒搞好PowerPoint.
我照樣睡.
當然我心裡.
IP部會不會罵我.
但是.
不行了 我已經支持不住了.
很安然地睡了.
我昨晚睡了.
一點多.
今天六點已經醒了.
很自然地.
我經歷了很多個星期.
自從我康復.
去事奉很多事情.
好像還沒搞定.
但是所有的事情.
真的去到神的面前.
你真的要問自己.
是否真的去到神的面前.
真的來到主耶穌的時候.
你會發現你是享安息的.
倘若我們今天.
我們每一個.
我們沒有經歷過在上帝.
在主耶穌 在聖靈裡面.
享安息.
那份滿足.
我告訴你.
你的去使萬民成為門徒.
是很艱辛的.
特別.
上一個星期12週年.
珍姑娘說.
接下來我們要.
遵行主道.
學無真理.
那個對象.
是上帝.
說容易一點.
你向上帝交往.

$^{1081}$上帝知道什麼.
這次還記不記得.
接下來這一年的.
教會主題是什麼.
上個星期.
你看回鎖匙扣.
彼此相愛.
相愛就像剛才那隻鴨.
有很多很多洞.
每個頭都放在那裡.
我愛你.
走吧.
不是走右邊.
是走左邊.
是更加的.
負挑戰性.
是更加需要我們.
來到主耶穌那裡.
因為我和神.
更容易掌控.
我如何把我的時間.
放在神的那裡.
但是你和人.
大家都經歷過.
你想想.
小學.
幸好幼稚園還沒有欺凌.
小學中學大學.
有欺凌的事.
過去三個月.
22個學生.
自殺死了.
這是看新聞.
你會看到.
很多不容易.
彼此相愛.
說就天下無敵.
做呢.
做了無能為力.
特別是那些.

$^{1121}$如何愛得下.
人和人之間很微妙.
我經歷過.
我都不認識他.
他又不認識我.
他剛來短孫才見到他.
他又見到我.
他的導師告訴他.
他很害怕你.
其實還未和他接觸.
那麼多事.
後來都要處理.
原來他過去.
有一個人對他好不好.
所以.
他一見到我.
他的恐懼又來了.
很害怕很不喜歡這個人.
就是這樣.
所以是不是很微妙.
你沒有做過什麼.
你都可能會這樣.
所以剛才我講到.
我們不要所有事.
照單傳授.
來到主耶穌的書去享安息.
更加大挑戰的時候.
我們不怕.
因為我們有上帝.
祂讓我們有應許.
來到上帝裡面.
來扶這個.
剛剛好適合你.
帶了鴨去耕田.
辛苦嗎?.
都是辛苦的.
如果沒有鴨.
他就完全什麼都做不到.
今天面對著.
整個世界.

$^{1161}$很多事.
不知道你有沒有同感.
一打開手機.
那些新聞.
多了很多罪案.
很多.
為什麼這麼多.
好像多了很多.
很多事是無能為力的.
來到.
主耶穌的書.
祂答應過.
給我們每一個.
很適合你的鴨.
是輕散的.
求主祝福每一位.
\newpage



\section{腓利門書 1:1-4-20}
\label{sec:xQAERVXii10}
\textbf{2023.10.29 42周年堂慶聯合崇拜《友伴．前行》腓利門書 1\_1,4,20 周曉暉牧師}
\newline
\newline
連結: \href{https://youtube.com/watch?v=xQAERVXii10}{\texttt{https://youtube.com/watch?v=xQAERVXii10}} ~~~~ 語音日期: 2023-11-18
\newline
\newline
\hyperref[sec:EtvLUaUHtqc]{\small{< < < PREV SERMON < < <}}
~
\hyperref[sec:index_chronic]{\small{[返順時目]}}
~
\hyperref[sec:index_scriptual]{\small{[返順卷目]}}
~
\hyperref[sec:3lGUeO3H154]{\small{> > > NEXT SERMON > > >}}
\newline
\newline
腓利門書 1:1-4-20
\newline
\begin{longtable}{cl}
\hline
\hline
章節 & 經文 (和合本修訂版)\\
\hline
1:1 & \begin{tabularx}{0.7\textwidth}{X} 為基督耶穌被囚的保羅，同弟兄提摩太，寫信給我們所親愛的同工腓利門、 \end{tabularx} \\ \\ \relax
1:2 & \begin{tabularx}{0.7\textwidth}{X} 亞腓亞姊妹，和我們的戰友亞基布，以及在你家裡的教會。 \end{tabularx} \\ \\ \relax
1:3 & \begin{tabularx}{0.7\textwidth}{X} 願恩惠、平安從我們的父神和主耶穌基督歸給你們！ \end{tabularx} \\ \\ \relax
1:4 & \begin{tabularx}{0.7\textwidth}{X} 我在禱告中記念你的時候，常為你感謝我的神， \end{tabularx} \\ \\ \relax
1:5 & \begin{tabularx}{0.7\textwidth}{X} 因聽說你對眾聖徒的愛心，和你對主耶穌的信心。 \end{tabularx} \\ \\ \relax
1:6 & \begin{tabularx}{0.7\textwidth}{X} 願你與人分享信心的時候，能產生功效，讓人知道我們所行的各樣善事都是為基督做的。 \end{tabularx} \\ \\ \relax
1:7 & \begin{tabularx}{0.7\textwidth}{X} 弟兄啊，由於你的愛心，我得到極大的快樂和安慰，因為眾聖徒的心從你得到舒暢。 \end{tabularx} \\ \\ \relax
1:8 & \begin{tabularx}{0.7\textwidth}{X} 雖然我靠著基督能放膽吩咐你做該做的事， \end{tabularx} \\ \\ \relax
1:9 & \begin{tabularx}{0.7\textwidth}{X} 可是像我這上了年紀的保羅，現在又是為基督耶穌被囚的，寧可憑著愛心求你， \end{tabularx} \\ \\ \relax
1:10 & \begin{tabularx}{0.7\textwidth}{X} 就是為我在捆鎖中所生的兒子阿尼西謀求你。 \end{tabularx} \\ \\ \relax
1:11 & \begin{tabularx}{0.7\textwidth}{X} 從前他與你沒有益處，但如今與你我都有益處。 \end{tabularx} \\ \\ \relax
1:12 & \begin{tabularx}{0.7\textwidth}{X} 我現在打發他回到你那裡去，他是我心肝。 \end{tabularx} \\ \\ \relax
1:13 & \begin{tabularx}{0.7\textwidth}{X} 我本來有意將他留下，在我為福音所受的捆鎖中替你伺候我。 \end{tabularx} \\ \\ \relax
1:14 & \begin{tabularx}{0.7\textwidth}{X} 但不知道你的意見，我不願意這樣做，好使你的善行不是出於勉強，而是出於自願。 \end{tabularx} \\ \\ \relax
1:15 & \begin{tabularx}{0.7\textwidth}{X} 他暫時離開你，也許是要讓你永遠得著他， \end{tabularx} \\ \\ \relax
1:16 & \begin{tabularx}{0.7\textwidth}{X} 不再是奴隸，而是高過奴隸，是親愛的弟兄；對我確實如此，何況對你呢！無論在肉身或在主裡更是如此。 \end{tabularx} \\ \\ \relax
1:17 & \begin{tabularx}{0.7\textwidth}{X} 所以，你若以我為同伴，就接納他，如同接納我一樣。 \end{tabularx} \\ \\ \relax
1:18 & \begin{tabularx}{0.7\textwidth}{X} 他若虧負你，或欠你甚麼，都算在我的賬上吧， \end{tabularx} \\ \\ \relax
1:19 & \begin{tabularx}{0.7\textwidth}{X} 我必償還。這是我—保羅親筆寫的。我並不用對你說，甚至你自己也虧欠我呢！ \end{tabularx} \\ \\ \relax
1:20 & \begin{tabularx}{0.7\textwidth}{X} 弟兄啊，希望你使我在主裡因你得益處，讓我的心在基督裡得到舒暢。 \end{tabularx} \\ \\
[1ex]
\hline
\hline
\end{longtable}
$^{1}$各位感宣家庭姐妹 主女平安.
很想在這個堂慶日子裡面.
去送兩句話給大家.
災於活水江河 靈命日盡.
蘊納基督盤石 誠摯宣道.
當中是蘊含著我們宣道會.
一個很強調的精神.
我們不單止自己的靈命要盡心.
我們也要繼續承擔主所託付給我們的使命.
所以在這裡也謹代表自己的教會.
北宣去恭賀感宣家42周年的堂慶夢恩.
也代表宣道會 香港區聯會去恭賀大家.
剛才主席幫我們讀了這個肥理門書.
今天和你看這段經文.
我們到底在這裡當中可以領受到一個怎樣的訊息呢.
聖經裡面又是怎樣看一個教會.
一個群體的關係是怎樣呢.
當肥理門書這樣寫的時候.
這裡提到.
我聽見你對主耶穌和眾聖徒有愛心和信心.
你看到那個愛的字.
你就知道牧師你是說彼此相愛.
不單止這樣 在這裡更加提出.
在這段經文裡面很強調一個中心思想.
就是這班彼此相愛的人是有團契關係的.
你說我知道 我也會回團契 回小組.
或者我們一起查經.
不單止是那個合意的小組群體.
而是到底我們這個家 這個整個教會.
是否有彼此團契的關係.
這個更加重要.
團契這個字的原文在新約聖經出現了有二十次.
其中有一次出現在肥理門書.
剛才我聽經文也不是很看得懂.
一會兒會告訴你在哪裡.
肥理門書是保羅在監獄裡面所寫的信件.
信當中保羅為一個由肥理門這個家裡逃走出來的奴僕.
或者奴隸歐尼西姆去求情.
歐尼西姆逃跑之後輾轉就遇到保羅.
也因此信了耶穌.

$^{41}$表面上這封信是處理一個個別的事件.
修補兩個人的關係.
但實際保羅是行文之中讓大家體會到.
他做了三個的示範.
我們走在這條天路裡面.
我們說的群體 我們說的教會.
在這條天路裡面怎樣可以得到友伴.
一個信仰群體又可以怎樣一起前行.
我們會透過這封信一起去學習一下.
我們自己的群體 我們的教會生活.
是否也可以向著這個目標.
友伴 前行 繼續邁進.
我們一起將這段聽到的時間再交代給主.
剛才我們在敬拜的當中.
我們很感受到這個家是由很多人去組成.
亦都彼此願意同心.
我們知道我們的目標是對準你.
亦領受你給我們的使命.
現在我們打開你的話語.
願聖靈指教我們當明白的真理.
多謝你.
我們的祈禱是奉主耶穌基督.
得勝名自祈求 阿們.
在剛才的四至六節那裡.
我們再一起看看.
不如我請大家一起和我讀這三節經文.
預備起.
我聽見你對主耶穌和眾聖徒.
有愛心和信心.
我每逢禱告提到你的時候.
願你與眾人分享你的信心的時候.
會產生功效.
使我們可以知道在我們中間的一切善事.
都是為基督作的.
這裡提到愛心.
你會看得很清楚.
在第四五節那裡提到.
他說:肥理門你有的愛心不是掛在口邊的.
根據聖經學者的分析.
在肥理門的年代曾經發生過一些天災.

$^{81}$極之可能的就是肥理門見到教會裡面.
因災難而生活缺乏的人.
他們缺少生活的物資 錢糧.
於是肥理門就用自己所有的去幫助這些弟兄姊妹.
其中第五節更加貼近原文的翻譯是這樣的.
因為我聽說你對眾聖徒有愛心.
對主耶穌有信心.
原來在這封信當中不斷提到的愛.
對象是弟兄姊妹.
是主內的弟兄姊妹.
是基督徒之間.
是聖徒的群體有的彼此相愛.
這個愛心不只是我們對我們見到的人.
在街外的人有連續恩賜.
而這種愛心在我們的群體都要常常有.
或者我們都稍為在想一想.
當這個時刻問你.
你在腦海裡面想起哪一兩位敢宣家的弟兄姊妹呢.
你應該想像到 馬上想到.
現在常常和你WhatsApp通訊 微信的弟兄姊妹是哪一類呢.
可能你有時候想到的人不是這些很關切的人.
可能是你要不斷學習以愛心去對待的人.
會不會有些時候自己其實對教會人的愛心.
是比對自己家裡的人多.
在我們真實的家庭裡面有時候是這樣的.
在教會的家裡面會不會都是這樣呢.
對外面的人會寬容.
但卻是對本應親近的人.
有時候會多了少少計較.
在第六節裡說.
願你與眾人分享你的信心的時候.
這個分享其實就是我一開始說的團契這個字.
在聖經的原文就是團契這個字.
當然這個團契不是限制一個小組一個群體.
而是可以說是一個教會一班人.
聚集一齊敬拜常常見面的一班人.
在這個其中 分享是有分享的.
我見到外面擺了很多食物攤位.
一會兒大家會分享食物.
不單止這樣.

$^{121}$對於肥利們 保羅是否想跟他們說.
記住了 肥利們.
你這麼有愛心 你就要有愛心到底.
剛才一開始說過這封信表面上是說兩個人要修補關係.
是不是保羅想說.
你去原諒那個得罪你的人.
不是.
保羅不是這樣去教肥利們.
或者命令肥利們.
你要這樣這樣做.
保羅是首先說我去分享我怎樣實踐.
因著信耶穌基督而對人能夠有的愛心.
對我身邊的聖徒有的愛心.
在第八第九節他這樣說.
他本應可以吩咐肥利們.
因為教會的群體裡面他的位份是很高的.
但是保羅沒有說一人少一句.
保羅也沒有計較自己是使徒地位.
你就聽我說 給我面子.
他不是這樣.
他也沒有計較自己的面子 個人榮辱.
保羅說我是憑著.
我是透過.
我是通過愛去請求你.
去求你肥利們.
我求你.
你可以像我一樣這樣去做.
我在其他教會聽過一件這樣的事.
不是錦算監 不是北宣.
有兩位弟兄.
他們一起去開教會執事會.
在過去一段很長的年日裡面.
領袖弟兄A也得到領袖弟兄B的幫助.
每一次開完會都送他回家 順道.
每次在路上開車的時候.
弟兄A和弟兄B會談剛才開會的事.
說一下教會的事情.
大家繼續分享.
有一次就是這樣開完執事會議之後.
回家的途中.

$^{161}$弟兄A說開教會的崇拜.
就和弟兄B說.
現在崇拜這樣這樣.
我覺得應該可以這樣.
於是他就把自己的意見.
這樣這樣這樣.
不斷地和弟兄B說.
在那一次車程的分享之後.
弟兄A不知為何總是覺得.
弟兄B接著就有意無意嘗試避開他.
甚至開完會之後.
總是他去過廁所.
轉頭他就開車走了.
沒有接著他一起走.
弟兄A思前想後.
是不是正因為弟兄B.
上次和他說崇拜的事.
他是有份負責的.
是不是他有點不高興呢.
但是他說大家都是領袖.
大家都是在教會裡去管理的人.
我當時已經很客氣和他.
去說這件事很中立的.
去分享一些看法.
所以我不是很明白.
為什麼他會這樣氣我.
再過多一段時間.
弟兄A再三反省.
他發現是不是我沒有換位思考.
我沒有在我這位弟兄的角度.
去看回這件事.
當我去分享我的意見.
很中肯的意見的時候.
我作為執事的資格的時候.
我用權利和地位去出發.
去反映對方B弟兄所收到的.
是一些誠實話.
B弟兄他反省.
是 我是說了誠實話.
但是我沒有對弟兄有愛心.

$^{201}$保羅為肥理門去感恩.
因為他真的可以.
展現出愛心在有災難的時候.
怎樣去幫助身邊人.
社區的人.
但是保羅他繼續去教導肥理門.
或者去求肥理門.
你可以像我這樣.
嘗試超越我們的位份.
我們的權利.
用這一份很單純的愛心.
去對待我們身邊的聖徒.
你看到嗎.
教會群體裡面.
我們因著這個大團契的關係.
我們是著重能夠.
很實在的去彼此相愛.
說的話的人.
自己實踐出來.
保羅去分享.
他也是這樣做.
這個經文對我都有很大的提醒.
回想過去這幾年.
我們所面對的人際那種.
我形容很intense.
很緊張.
很容易擦槍走火.
一句話 一句貼文.
可以變成很多爭拗的年日.
超越權利.
用愛心去對待我身邊的人.
這個是主給我的提醒.
現代人的關係的確是很麻煩.
因為我們常常被這個社會去教導.
我們要用功利 得失 效益 好處去計算.
對的 因為我們是這樣生存的.
我們是這樣生活的.
我們去賺得我們每天所需用.
是用這些方式.
所以這是很中立的事情.

$^{241}$但是如果我們所有的人際關係都是這樣的話.
就很可怕了.
例如一個人轉了工.
他自然會繼續和舊同事有一些聯繫.
不過慢慢就會越來越少.
因為大家不再圍繞同一個目標.
為了同一間公司去賺到一些利益.
以前的關係是會越來越淡的.
你轉過工 你明白的.
還有一些聯絡的 大概在工作上有需要.
又或者今天留一線.
他日還有生意可以談 還有合作的機會.
不過這些狀況如果進入到教會群體.
特別是在疫情裡面 我們會見少了.
你就會發覺有時候我們仍然是跌落這個社會.
教導我們以功利 得失 效益 好處.
去計算我們和別人的來往.
腓利門和歐尼西姆的關係 起初都是一種功利的關係.
用今天的說法 歐尼西姆是一個工具人.
為什麼呢 因為他的存在是幫腓利門去做事的奴隸.
幫他做不同的事務.
所以在十一節裡說.
保羅對腓利門說 他從前對你沒有什麼好處.
現在對你對我都有好處.
這個好處 這個致詞.
其實是用回耶穌在馬達芬講過.
五千 二千 一千的比喻裡面.
那個一千的 最後主人說他是沒用的僕人.
那個沒用的兩個字.
用就是好處或益處.
在原文是同一個字.
歐尼西姆對腓利門沒有好處.
因為他逃跑了.
不再在我身邊工作 為我做事.
所以對我來說是沒用 沒益處 沒好處.
打個比喻 今天在你公司裡面有一個超級不負責任的同事.
不盡責任 不做事.
你作為主管 你可以怎樣做.
見到他回來就是吃早餐.
很快就訂午餐.

$^{281}$然後就想喝下午茶 喝不喝珍珠奶茶.
然後等下班.
如果是這樣的話 你見到一個他都不是很productive.
很能夠為公司的好處著想的人的時候.
你會覺得你不想留他.
當然現在一個這樣的人 你都要留 因為不夠人.
但回想腓利門 歐尼西姆.
他這樣逃跑了 或者不再幫自己做事.
腓利門是真的很生氣.
當時的律例是可以抓到這個奴隸 處死他.
羅馬的法例可以這樣做.
但保羅這個時候竟然說 這些都是以前.
過去了 接著還說.
這個以前或者這個階段對你沒有用的歐尼西姆.
如今對你對我都有好處.
其實保羅在說什麼.
我們很熟悉這一節的經文.
提摩太后書.
人若自竭 脫離卑賤的事.
就必作貴重的器皿 成為聖潔 合乎主用.
你們見到的用字 這個就是好處的字.
同一個字詞.
歐尼西姆的用和今天我們一起成為聖潔 合乎主用的用是同一個字.
你想想你家裡的碗碟.
一隻碗 一隻碟 作為裝東西的器具.
如果可以用 基本上都是裝食物 或者裝湯.
最重要是什麼.
看它是金造 銀造 還是木造.
不是 我想無論怎樣用都好.
最重要是乾淨.
如果你在碗櫃裡拿碗出來 看到上面還有菜渣.
還搭著蟑螂腳.
就算相金給你幾顆閃轉 你都不會用來裝飯.
不合用.
一個器皿要給主人使用 不一定是新的.
不一定是漂亮的 不一定是名貴的.
但是這個器皿能夠給主使用 是要清清潔潔 乾乾淨淨.
洗好之後 抹乾了 放回碗櫃裡 隨時主人可以拿來用.
古羅華這個以前沒有好處 沒用的歐尼西姆現在變成怎樣.
十六節 一起讀.

$^{321}$不再是奴隸 而是高過奴隸 是親愛的弟兄.
對我固然是這樣 對你來說 不論按肉身或在主內的關係 更是這樣.
按剛才說的羅馬社會的制度 看目標功利得失效益好處的計算.
一個逃走的奴隸 沒用的 是有害的.
但是現在保羅180度的大反轉 告訴肥利門 歐尼西姆是對你對我都有好處.
不僅如此 他還將奴隸奴僕 一下提升到弟兄的地位.
保羅的思想當時非常超前.
甚至可以說是荒謬的 當時沒有人會這樣想的.
但經文講了這麼多好處和不好處 其實是在說什麼.
是提出在我們這個團契關係的群體 一個教會當中.
我們可以學的就是我們待人不求 這個世界講的好處.
各位主內的關係的確不是功利的計算.
早幾個星期 亞運 不知道你們有沒有機會了解.
香港這次很厲害 很多金牌 游水又厲害 劍擊又厲害.
很多很多項目都拿到很多獎牌.
再早十多年的時候 香港隊參加亞運比較突出的項目.
只有兩三項 其中一項就是單車 你都知道我們單車隊很厲害.
在2006年的亞運 香港單車隊拿到六面的金牌 非常厲害.
不過當時拿到金牌的人 你會想起哪些人物呢.
可能你會想起黃金寶 這個名字我們基本上香港人都認識.
但在那一年 亞運不是黃金寶拿到金牌 而是他的師弟張勁偉.
在比賽之前 其實港隊的教練已經計劃好了.
接下來的個人賽是以黃金寶為搶分的主力.
他的師弟張勁偉 這個21歲的年輕人就是做副手.
負責幫師兄在車隊前面開路擋駕.
你知道單車要贏就要破風 如果有人在前面幫你破風.
後面的人就踩得輕鬆一點 以至到時他要為搶分的時候.
他的力量可以爆發出來.
不過在比賽過程當中一定有很多變數.
令到阿寶由最初帶開的 這一組合帶開的.
去到中段開始阿寶發覺自己虛耗過量 後勁不計.
這個時候他馬上改變想法 他示意自己的師弟張勁偉.
他爬到他前面 跟他說 我跟你換.
你做主力 自己就運用經驗速度去牽制其他有威脅的對手.
去協助這個師弟衝前.
最終這個張勁偉在衝前之前 壓倒對手奪得金牌.
原本身為人人都期待他拿金牌的主將.
亦都是他應該去負責搶分拿牌的阿寶.
竟然放棄了原本應該是他的獎牌.
他以身擋敵 造就師弟一舉擲金.

$^{361}$後來有人訪問阿寶.
你這樣做很可惜 大家都期望你的實力應該是拿金牌的.
阿寶毫不介意地說 以往都是眾師弟為我開路特惠.
將機會給我 成就我拿獎牌.
今次我只是去協助自己的師弟上前.
證明他們都有奪金的實力.
而且香港單車隊不只是只有我黃金寶.
人家贏 等於自己贏.
這是一個團隊的勝利.
因為這不一定只是計算自己個人得失功利好處.
本來都是這樣想.
而且經文更加說現在的團隊是誰是主人.
誰是僕人.
肥理門以往你是主人 歐連西武是僕人.
但是在十六節的結尾說.
我們都只是在主內的關係.
恩爵信同一位主.
大家是親愛的弟兄姊妹.
所以待人不求自己的好處.
這是保羅藉著這一節的經文再一次去求肥理門.
去明白這件事.
他自己也親身示範.
在前幾節他這樣說.
歐連西武對於保羅來說是很寶貴的.
好像契仔一樣.
因為在他坐牢的時候.
歐連西武信了耶穌.
他就留在保羅身邊去服侍他.
你覺得很奇怪 為什麼坐牢有其他人進去呢.
是的 在羅馬時代坐牢是沒有食物的.
不會給你飯吃 沒有皇家飯.
你要找你外面認識的人帶食物給你去照顧你.
你有病你要外面找個人來幫你看醫生 療傷.
所有東西都是由外面供應.
甚至有時候監獄是可以容讓一些你認識的人留下來幫你 陪你一些東西.
歐連西武就留在保羅身邊幫他.
這個上了年紀的保羅.
身體不好的保羅.
但是他說我為了大家的好處.
我很想保留這個像契仔一樣的歐連西武.

$^{401}$但是保羅知道我不會求自己.
求自己的好處.
真正對歐連西武對於肥利門有好處 有用的.
不是你們抹走過去不了了之.
算了算了 我寫封信你們就這樣遙距原諒他就算了.
而是真正對歐連西武對於肥利門有好處的.
就是他們可以面對面一起去看這件事.
最終在主裡面有復和 開展.
因著同信一位主而有這個團契 彼此相愛的關係.
領事館這個道理是很簡單 我們常常聽的.
但應用在我們自己教會生活裡面.
同樣是在過去這幾年我們有大環境 小處境.
所引致各樣各樣的相處碰撞就很不容易了.
保羅示範了甚麼是不求自己的益處.
是的 他說待人就是不求自己的益處.
願意我們都一起去學習.
我們這個群體 我們這個家.
甚至我們這個宣道會 誰是主場.
是你還是他 是牧師傳道還是執事領袖.
是導師還是團友 是這個地點的人還是那個地點的人.
都不是我們所信的主.
是我們所信的主耶穌基督是主場.
所以在這樣的旗號下.
我們的確是可以學待人不求自己的益處.
繼續我們一起看17節.
我請大家一起讀這節 預備起.
所以你要是把我看作同伴 就接納他好像接納我一樣.
這個同伴在希臘文是有一個common 或公共 共同的意思.
是這個公共 共同這個字變化出來的.
所以團契相交原來都是和這個同伴這個字詞.
這個公共 這個字是相同的來源.
這樣的話帶給我們這個字詞起碼有三重的意思.
第一 同伴 中文已經告訴我們了.
是companionship 是一同 是共同 是一切.
你都明白了 一起就是這個意思 比較普通.
未講到那種很緊密的關係.
所以第二種在中間就是一種夥伴 partnership.
有些事今天我和你拍檔 今天我和你一起做這個事奉.
你明白了 一起拍著相 有多一點默契.
在這個空間時間上 我們聚在一起.

$^{441}$我們有相同的追求 相同的目標.
拍檔是甚麼呢 當然我們明白.
保羅常常提到的夥伴是在福音上一起努力的.
就好像今年的大家的年提 全陽福音.
很明顯的 很清晰的目標 一起去做這件事.
而第三種 原本原文的同伴這個字.
還有一個有福同享 有難同當的意味.
所以在最上面的字詞.
我們今天最要掌握的就是一個友伴.
一個friendship 不單止是companionship.
不單止是partnership 還是friendship.
可以彼此分擔 彼此分享.
保羅和他的同工其實就是這樣.
一起的 有一個目標 一起傳揚基督.
同時在他們的生活裡面有很多的結連.
有很多千絲萬縷相處的時刻.
一起去面對不同的事情 彼此去扶持 傾訴心事.
這個是友伴.
我們可以想一想 疫情讓我們去跟人很疏離.
今天錦宣家可以一起聚首一堂.
連同孔恩堂 連同嘉欣樓的弟兄姊妹.
連同不同語言的群體 我們有菲律賓的姊妹一起.
今天在這裡我們聚在一起.
在這麼大的群體裡面 有沒有你的同伴呢.
有啊 一起敬拜嘛 是不是.
但是否只是停留於回來教會說聲hi.
離開教會說聲bye.
有時候真的是這樣.
我有的 我有一起事奉的.
今天我們一起唱詩班.
一起在後面做影音 做招待 做預備食物.
在一起是伙伴 在事工上是伙伴.
大家是不是友伴呢.
能不能說得心事呢.
當然我們不可能跟所有人是友伴.
但是你身旁有沒有這樣的弟兄姊妹.
你可以去不單止有事奉 還有生命的交流.
還能夠做真正的對公朋友 友伴呢.
我們怎樣可以找到這樣的對公呢.
我當年自己團契選職員.

$^{481}$當時負責選舉的弟兄姊妹就公佈.
未來這一年就有這個這個.
我是其中一個名單.
我看名單一出來 哇.
未來一年怎樣事奉啊.
他都想選職員.
他這麼古怪.
這個還經常吹水 怎樣事奉啊.
哇 剛剛大家面左左.
經文都是提醒我們.
你都讓我去理解到一件事.
有時候一切的友伴事奉 伙伴沒得選.
是神放我們在這裡的.
保羅跟腓利門是團友 是友伴.
如果是這樣的話 請你也接納歐尼西姆.
歐尼西姆也是我的友伴.
也可以是你的友伴.
而17節的下半節說請你接納他.
好像接納我一樣.
這個團契關係是主安排 不是隨自己喜好.
保羅說那你接納他吧.
這個接納其實是一個很有動感的致詞.
在《少行傳》28章形容保羅一次船難之後.
他們沉船 然後游到岸邊.
有一句經文是這樣說.
當時在島上有一群土人.
看到這群船難游上來的人.
就說看待他們有非常的情份.
當時下雨天氣又冷 就生火接待我們種人.
你可以想像一下這樣的情景.
我們幻想一下.
沉船很辛苦游到岸邊.
抱著船舶 終於去到岸邊 沒氣了.
在船難之前 其實他們沒有吃到東西.
一直被風浪打擊.
上岸之後 一直趴在海灘上的時候.
看到前面有一群土人 原居民.
在生火 氣 火 氣 火.
吃人族.
他們很熱情地拉著這些在海灘趴著的人.

$^{521}$拉他們去火邊.
不是燒他們.
拉他們去火邊 讓他們取暖.
所以接納這個字詞很有動感.
不是只有一張嘴.
是歡迎你來.
這是最基本.
但不止這樣.
不只是站在這裡.
接納不只是提出一個歡迎.
而是延伸一個歡迎.
伸出你的手 拉他們來.
向大家這個群體中心多走一步.
你是主動做這件事的.
其實中文字的接納納字都很精彩.
橋詩邊加個內字.
我們可以想像.
很全神地告訴你.
你用一條帶子將那個人牽引到你這個內圈.
接納.
你這個圈外的人.
你安書區.
你朋友圈以外的其他人.
進來裡面.
可能你都會像我當年選職員的時候有抗議.
要我接納一個這樣的人不容易.
一個身家清白好人好者都沒有難度.
但要接納一個和我面左左.
和我有很多爭執的人.
還要我做主動.
不是很便宜.
他得罪我的那個.
接納這個字.
在另一邊也用過.
在羅馬書15章7節.
那裡說.
所以你們要彼此接納.
如同基督接納你們一樣.
基督要我們想起.
這個接納不單止是我們用我們的能力.

$^{561}$其實我們很多時候是沒有能力甚至不想.
我們的心不舒服.
所以我們接納的動力或者起點.
不是在我們能否解脫那些事件.
而是更重要的是.
我們是先被主接納了.
先是由耶穌來到世上.
將我們拉進祂的國裡面.
神extend a welcome.
將我們牽扯到祂的愛當中.
如果是這樣的話.
我們是信這位主.
我們是可以有能力靠著主變得慷慨.
可以接納我們從前要擋在我們圈外的弟兄姊妹.
剛才我說那個團體職員選舉的事.
好了,選了,真的選了那些人.
一年之後.
我的看法完全改變了.
我真的很感受到.
每一個人是主擺在這個團隊當中.
每一個人都是屬於主.
也可以是我的友伴.
在那一年裡.
對我最多意見的團友.
我覺得他經常不聽我的話.
經常幫我包緊的人.
當我學,我只是停下來聽多一點.
我發覺那個關係不同了.
甚至在事奉以外.
我們可以一起去街上看電影.
一起去打球.
一起玩.
不只那間由夥伴變成友伴.
神就使用他.
我發覺在一年之後.
我覺得上帝是派他來收止我的生命.
今天我們仍然是很好的團友.
友伴.
說回頭.
那些我以為相處不到的弟兄或者姊妹.

$^{601}$他們何嘗不是多走一步.
延長和歡迎.
將我接納在他們的內圈裡.
弟兄姊妹沒有主耶穌基督慷慨的接納行動.
我們也沒有這樣的.
但我們是信主的群體.
我們高舉耶穌基督的十字架.
我們這個信主的群體.
是能夠靠著主耶穌基督的愛.
去給予我們的接納.
我們也是這樣去接納.
一些和我們很不同的人.
去接納在我們看書圈以外.
主愛親愛的弟兄姊妹.
這個友伴的功課.
在服償的階段.
可以是我們最新的目標.
放下過去疫情裡.
放下過去社會很多的事情.
是我們還要去面對.
但是我們的起點.
是因為主接納了我們.
我們信主.
我們就能夠去接納身旁不同的人.
我們不需要靠自己.
是靠主.
在主裡.
你做得到.
我做得到.
保羅知道.
肥里門一定會聽他的意見.
甚至超出在信裡的要求.
所以弟兄姊妹.
今天是錦宣家四十二週年的堂慶.
十二年是一個很長的歷史.
有歷史就有故事.
有故事就有美麗的故事.
有艱難的故事.
但是無論如何.
我們在這一點可以慶祝.

$^{641}$我們依然在主這個群體裡.
我們願意走在一起.
我們就可以成為友伴.
展望前行.
疫情過後社會的變幻.
大家靠主能夠維繫這個關係.
無論如何.
經文鼓勵我們.
我們學習超越自己的權利地位.
我們用愛心去對我們家裡的弟兄姊妹.
我們去待人.
不要像社會那樣去計算功利.
我們學習不求自己的好處.
更加我們學習.
我們因為同友一位主.
同信一位主.
我們領受的慷慨.
我們都是這樣去對我們身邊.
每一個親愛的弟兄姊妹.
求主幫助我們恩著信.
不斷去建立這個真正大團契的家.
最後很快說一說.
你想不想知道肥利門他們最後怎樣.
聖經沒有寫.
歐尼西姆他們怎樣.
教會歷史是這樣說的.
後來肥利門.
看到歐尼西姆親自帶著這封保羅的書信來.
他看了之後.
他原諒了歐尼西姆.
接納了他.
拉回到他家裡.
這個教會裡.
還給他自由.
不再看他為奴隸.
看他為弟兄.
後來歐尼西姆在主後一百年左右.
成為了以弗所的主教.
這個奴隸成為了帶領教會的主教.
他還負責收集保羅的書信.

$^{681}$所以為什麼這封信.
這麼個人的信.
匯入了保羅的十三封書信.
也在我們今天新約的聖經裡.
你看看因著真正友伴的關係.
這個教會.
這個團體.
這個對公.
這兩個很不同的人.
很有過去歷史的人.
他們成為友伴.
一切前行.
為主作了美好的事.
傳揚福音.
同樣我們可以因著主的教導.
一切得著更新.
我們祈禱.
主求你大大賜福給金孫家.
當我們每一個人願意.
在你的十字架底下成為一群.
能夠彼此相愛.
互相視為友伴.
這個群體就能夠被你帶領前行.
願主繼續使這個地方.
去成為作炎作光.
一個美好的見證.
多謝你.
我們祈禱在奉主耶穌基督的聖名之祈求.
阿們.
\newpage



\section{羅馬書 12:15}
\label{sec:3lGUeO3H154}
\textbf{2023.11.19 《愛裡同行》 羅馬書12\_15 (和修版) 講員 : 曾敏傳道}
\newline
\newline
連結: \href{https://youtube.com/watch?v=3lGUeO3H154}{\texttt{https://youtube.com/watch?v=3lGUeO3H154}} ~~~~ 語音日期: 2023-11-20
\newline
\newline
\hyperref[sec:xQAERVXii10]{\small{< < < PREV SERMON < < <}}
~
\hyperref[sec:index_chronic]{\small{[返順時目]}}
~
\hyperref[sec:index_scriptual]{\small{[返順卷目]}}
~
\hyperref[sec:3jZxMYDYvEI]{\small{> > > NEXT SERMON > > >}}
\newline
\newline
羅馬書 12:15
\newline
\begin{longtable}{cl}
\hline
\hline
章節 & 經文 (和合本修訂版)\\
\hline
12:15 & \begin{tabularx}{0.7\textwidth}{X} 要與喜樂的人同樂；要與哀哭的人同哭。 \end{tabularx} \\ \\
[1ex]
\hline
\hline
\end{longtable}
$^{1}$弟兄姊妹平安.
可以在此敬拜.
亦可以彼此有相交.
實在是很多無比.
今天說愛,女同行這個題目.
是一個我和你都很需要學習的課題.
在此我們先看看.
在教會生活當中,什麼是重要的呢?.
神的話語,一定.
神的話語很重要.
彼此相愛.
因為今天要說愛.
愛很重要.
因為耶穌基督呼召我們進入群體生活.
耶穌基督本身就是愛.
呼召我們的是愛.
而祂呼召我們進入愛的群體.
在此我們彼此實踐相愛.
所以愛是非常非常重要.
我們開始時再繼續一同禱告.
此下你的話語要成為我們腳前的燈路上的光.
你的話語就是我們的指引.
就是我們的標準.
讓我們可以去跟從.
讓我們可以去學著去走.
主阿禽你今天給我們每一個都有柔軟的心.
給我們有柔軟的耳朵.
我們聽得進去.
也給我們有力量去回應.
願主帶領以下的時間.
奉耶穌基督的名字祈禱.
阿們.
在約翰福音13章說.
34和35節大家聽我讀出.
我賜給你們一條新命令.
乃是叫你們彼此相愛.
我怎樣愛你們.
你們也要怎樣彼此相愛.
很重要這是耶穌基督的命令.
為什麼我們說今年我們很想彼此相愛.

$^{41}$因為耶穌基督很想我們這樣做.
祂很想祂怎樣愛我們.
我們又怎樣彼此相愛.
我曾經在一堂道裡說過.
耶穌是怎樣愛我們的.
有沒有人記得呢.
雖然那堂道隔了一段時間.
有沒有人記得呢.
我們愛因為神先愛我們.
上次是這樣說的.
耶穌說我怎樣愛你們.
你們又要怎樣彼此相愛.
有沒有人說得少一點.
耶穌是怎樣愛我們的.
真的耶穌基督為我們捨身.
犧牲祂自己的生命.
祂是這樣愛我們的.
完全犧牲的愛.
耶穌說祂怎樣愛我們.
我們又要怎樣彼此相愛.
我們要彼此為對方什麼.
大聲一點說.
不夠膽說.
這句話很難說出口.
寫命.
聖經裡也有說.
人為朋友寫命.
沒有愛比這個更大.
不容易做.
我怎樣愛你們.
為你們寫命.
你們也要彼此相愛.
彼此寫命.
你們要彼此相愛.
眾人因此就認出.
你們是我的門徒了.
其他的人.
未信主的人.
在我們身邊生活的人.
怎樣認出我們是屬於主的呢.

$^{81}$就是彼此相愛.
原來這是標記.
這是教會的標記.
是我們一起的標記.
就是彼此相愛.
我跟你外出.
額頭不會寫著基督徒.
我是屬於耶穌的.
人們有不同途徑認出我們.
我們的言行.
我們的談吐.
怎樣為人處事.
人們會看的.
看的時候.
嘩!這一定是基督徒.
他行事為人.
處事一定與別人不同.
還有一點.
就是這個群體彼此相愛的時候.
這就是標記.
說愛要有表現.
原來愛的表現.
就是我們今天所說的.
一字.
要與喜樂的人同樂.
要與哀哭的人同哭.
在繼續說之前.
我很想問弟兄姊妹.
大家覺得這句難不難做.
要與喜樂的人同樂.
要與哀哭的人同哭.
難不難做.
有人馬上說不難.
是吧.
為何覺得不難呢.
因為愛他就會這樣做.
愛他就會這樣做.
很容易的.
愛一點都不難.
有人說難嗎.

$^{121}$難.
與別人哀傷又不一樣.
但當別人離婚.
可能又很難去妒忌.
唉!真是說得好.
有沒有同感的.
有呀!馬上有人舉手.
你與別人哀哭是容易的.
一會兒我們看看.
與別人哀哭是否真的容易.
與別人喜樂不是我們想像的容易.
要與喜樂的人同樂.
我們身邊的人.
或認識的人.
有些很值得開心的事.
這個經文很白.
與他一起開心.
是否每次都做得到呢.
是否真的這麼容易做到呢.
我記得.
以前有個青年人跟我說.
他回想起放榜那天.
放榜的成績不錯.
一拿到成績表.
那一科的成績.
五星五.
現在DSE說分數安排.
不錯.
我很開心.
我拿著成績表.
跟一個好朋友說.
你看看.
我今天放榜的成績.
是你很要好的朋友.
你覺得你的朋友會怎樣.
他會說替你開心.
考這個成績.
然後報學校.
考大學.
一會兒我們慶祝.

$^{161}$是否這樣呢.
要看看對方.
看看是誰.
對方和你差不多.
我都很好.
我給你看看.
我考了這些成績.
你看看多好.
真的要一起慶祝.
我立即打電話給誰.
你都很好.
我們一起.
要不然他更加比你好.
真不好意思.
我的成績這樣.
你都很好.
不是.
當時對方的反應是甚麼.
很不舒服.
反應很冷淡.
完全出乎他意料之外.
不是一起開心.
是很不舒服.
因為有些科目.
很明顯對方比他差一點.
是嗎.
不容易的.
同樂.
教會裡面不會的.
是嗎.
是真的嗎.
今日那個主席走出來.
真的做得很好.
很多人稱讚你.
你今天做得很好.
你跟別人分享.
我今天一路禱告.
神給我力量.
整堂誦拜在我帶領下.
很順暢.

$^{201}$上帝與我同在.
弟兄姊妹為你感恩.
你真是好.
我看到神的祝福.
是不是這樣.
都很容易想.
我那一次.
一出來吃了螺絲.
又發了台溫.
弟兄姊妹.
看到我下台時.
弟兄姊妹你很緊張.
想起我心裡很不舒服.
今日弟兄姊妹在我面前.
跟我說.
她的事奉是怎樣.
她說是不舒服.
我們會比較.
你說傳道人不會比較.
是呀,謝牧師.
今日我們在這裡.
我們傳道人不會比較.
傳道人很熟悉.
是嗎.
我特別喜歡聽某某講完講到.
他一出來講.
我覺得他的堂道.
很能夠應用在我生命裡.
一下去.
某某傳道.
你真是講得很好.
你知不知道今日你的堂道.
祝福我的生命.
我心裡很感動.
我剛才就想哭了.
這個傳道跟他的同事說.
有弟兄姊妹這樣鼓勵我.
很快樂.
你猜另一個同事.
一定是.

$^{241}$你今日的堂道真是講得很好.
同理神與你同在.
神的恩典很好.
一定.
心裡會不會想.
每次我下去都沒有人跟我說.
你真是講得很好.
你第一次去做事很悶.
我睡了.
還未睡.
是否真是.
你回到家.
真是.
你家人很親密.
我一定愛你.
你開心我開心.
你不開心我不開心.
你回到家是否沒有競爭.
沒有窒息.
是真的.
這麼多年在家裡生活.
真是這樣.
你學一下你姐姐.
是嗎.
你姐姐多會做人.
姐姐又乖.
你一坐下就這樣.
學一下你姐姐.
學一下你哥哥弟弟.
多會做人.
一出糧拿一份糧回來.
拿了多少.
今天哥哥回來.
弟弟.
我加薪.
老闆很欣賞我.
我令他開心.
自己想想.
為何我的事業發展這樣.
爸爸媽媽經常稱讚你.

$^{281}$身邊人人都欣賞你.
我呢.
是否真是.
你愛的人.
真是痛樂到.
你心裡.
真是可以.
歡樂.
你真真正正.
怎可以和喜樂的人.
真是喜樂的人痛樂.
怎可以.
你令人開心.
在你裡面.
是沒有.
競爭.
沒有窒息.
無論信主的.
未信主的.
都很受影響.
你裡面的愛有多大.
是否你真正.
可以令人開心.
我記得.
以前在舞會的時候.
弟兄姊妹都和我說.
有一團契走在一起.
那時我非常年輕.
我的弟兄姊妹已經.
大學畢業.
出來工作.
在裡面的事情.
弟兄姊妹.
各自討論的題目.
是很不同的.
弟兄走過來.
一定會說自己的發展.
自己的成就.
你會感受到.
一種無名的.

$^{321}$競爭在裡面.
一定會說.
我的工作如何.
你的工作如何.
不只是分享.
在裡面的事情.
是比較的.
姊妹會說別的.
姊妹不會說事業.
姊妹會說家庭.
說我那位怎樣.
說我們自己怎樣.
我們瘦身.
我們很多方面.
各方面很多東西.
姊妹會說.
你真的很愛你的弟兄.
姊妹,你的家人.
你身邊的朋友.
你真的可以做到不競爭,不嫉妒.
是很需要.
有神的愛在我們裡面.
我剛才開播.
耶穌的命令是.
彼此相愛.
我怎樣愛你們.
你們就要彼此相愛.
是需要有神的愛.
在裡面.
我曾經說過.
人的愛是很有限.
我們不要互相吹噓.
我這個人的愛.
有多大.
對不起,我真的告訴你.
人的愛有計較.
有限.
但上帝的愛.
卻是無限.
上帝已經為我和你犧牲了.

$^{361}$毫無保留地愛.
耶穌說要我們.
跟著祂彼此相愛.
不是說盡量吧.
你做多少就多少.
可以吧.
不是,像祂這樣.
首先我們.
要和耶穌基督有緊密的關係.
一個能夠愛的人.
和主的關係緊密.
我最近.
很喜歡.
以前有一個.
很偉大的.
報道家,布永康牧師.
他已經安息主懷.
我看回他的.
港獨帶.
心裡很感動.
我聽到.
弟兄姊妹的見證分享.
她有同一個描述.
你看到布永康牧師就像看見神.
他是一個很慷慨的人.
很有慈父的心腸.
很有愛.
在他身上你感受到那種愛.
那種愛不是人的愛.
是那種和神很近.
是神的愛在他生命裡發出來.
貼近他的時候.
你自然感受到.
不同的人都從他生命裡感受到這一點.
我深信布永康牧師.
能夠和人同喜樂.
不會計較,競爭,妒忌.
那份力量是很大.
今天我和你.
我都是和自己說些托道.

$^{401}$我不覺得我比你好.
事實上我在外的功課.
是有很多跌跌碰碰.
所以我今天都是對自己說.
我和你一同去學這個功課.
我和你都要更加親近.
和主很近.
和主得著源源不絕的愛力量.
以致和人分享那份喜樂.
今天在這裡我答應了.
一會兒再說這裡.
我們站在別人的角度.
去了解別人的快樂.
再來一次.
那個年輕人拿成績表過去.
沒錯.
雖然我自己心情不是很好.
我做另一個.
我心情不是很好.
我今天放榜.
我做得很好.
但我站在他的角度.
他真的很開心.
我替他開心.
今天弟兄姊妹做主席.
做得很好.
我為他感恩.
真是神的賜福.
就算我要想自己的事奉.
我需要成長.
也好.
試試站在別人的角度.
如果今天那個人是你.
你自己很開心.
站在別人的角度.
去想想他的快樂.
今天沒錯.
我講得不像某某牧師.
某某名傳道.
那麼容易進入人心.

$^{441}$我要去成長.
但當那個傳道人在講的時候.
我為他感恩.
神使用他.
我自己也專心聽到.
我要因為神用他.
用他的訊息.
我自己生命也要去成長.
站在別人的角度.
讓他開心.
今天要講生日見證.
稍後大家會因為.
這個生日見證.
而得到很大的福氣.
為泰語姊妹.
Arisa獻上感恩.
今天要代她講這個見證.
真的很特別.
在這個我講到的日子.
而她稍後會做一個舉動.
話說.
她在香港.
11月13日.
生日.
她生日的時候.
主人很疼愛她.
給了她一千元.
和跟她說.
你猜她說什麼?.
生日快樂.
給一千元就為了說生日快樂.
放幾日假.
也不錯.
她說.
生日快樂.
她說.
還有什麼嘗試?.
幫她買東西.
反過來.
主人幫她買東西.

$^{481}$你的神很好.
嗯.
和我們的神.
有些關係.
和教會很有關係.
她就說.
給了她一千元.
和跟她說.
你猜她說什麼?.
給了她一千元.
和跟你教會的朋友.
一起.
吃東西慶祝.
你和你教會的朋友.
拿著一千元.
吃東西慶祝.
當時Arisa不明白.
為什麼主人.
不是說.
我給你一千元.
你自己去買東西.
你喜歡買什麼.
你自己很開心的慶祝.
主人不是這樣說.
是你和你教會的朋友.
是主人這樣說.
是不是很奇妙?.
很奇妙的作為.
於是.
Arisa就很馴服.
她今天就打算.
請我們.
吃東西.
你知道她一早.
多遠?.
由摩星嶺那麼遠.
來到這裡為我們準備食物.
而且之前.
還得到主人的同意.
在家裡預備一切食物.

$^{521}$否則做不到這麼多份.
給我們這麼多食物.
是不是很厲害?.
她今天就要.
和大家一起慶祝她的生日.
讓主人給她這個錢.
我深信這個在她人生當中.
是很值得紀念的一天.
不單是她生日.
原來神.
透過主人給她的錢.
和我們一起.
慶祝.
我們有份.
和她同樂.
一起喜樂.
一起分享神的愛.
我們的聖餐很奇妙.
是很值得開心的.
不單是.
Arisa 生日.
整個教會在主的愛裡.
是有份.
在主的恩典裡有份.
一會兒我想大家做一件事.
當崇拜結束.
不單是.
很開心.
我們一起吃.
我想Arisa.
站起來讓大家認識你.
因為大家.
未必知道是誰.
不過做一件事.
祝福她.
請大家說一句祝福的說話.
祝福她.
在主裡彼此互相分享.
這是很寶貴的.
當大家.

$^{561}$祝福的時候.
想想.
上帝想她的生命.
是怎樣的.
上帝想她得到什麼福氣.
這是很重要的.
很重要.
接著.
剛才我們說.
要哀哭的人.
同哭是比較容易的.
是不是很容易.
不容易.
在小朋友哭了之後.
我們覺得不容易.
為什麼呢?.
因為她哭的時候.
我們不可以和小朋友一起哭.
是.
和哀哭的人同哭.
會遇到什麼困難呢?.
你很愛那個人.
或者很愛某個家庭.
他們可能有家人離世.
或者有家人有重病.
你很愛哀傷的人和他們一起.
和困難的人一起.
大家覺得比較容易.
但覺得不容易是因為什麼呢?.
不認識.
首先我想說.
這段時間.
我們教會有不同的喪禮.
在喪禮裡.
見證到紅人家很有愛.
我看到弟兄姊妹.
都很想透過不同的方式.
向哀傷的家庭.
表達慰問.
無論是透過花牌.

$^{601}$在喪禮裡唱詩歌.
我們教會在這方面很用心.
弟兄姊妹真的很熱心,很主動.
我回到教會.
真的很欣賞.
紅人家弟兄姊妹.
在這方面做得特別多.
我記得昨天.
最後解慰飯的時候.
只有我,簡太和慧玲.
大家還記得.
人家稱讚我們.
昨天的師班唱得很好.
我發聲說.
他們都在場.
我說.
我們的弟兄姊妹很好.
其實昨天的組合.
不只是師班成員.
也有不少不是師班成員.
在當中參加.
每一次有喪禮.
我們都公開呼籲.
弟兄姊妹有感動.
就會一起出席.
是非常好的.
在過程中.
整個小組.
師班的獻唱.
是非常好的.
聲音很好.
服侍很好.
我們除了感恩之外.
也要說.
弟兄姊妹很樂於.
在這情況下.
與哀哭的人同哭.
很願意表達那份愛意.
其實昨天在喪禮裡.
或是在之前.

$^{641}$師班奶奶的喪禮裡.
都有弟兄姊妹有不同程度.
不同層面的付出.
這裡真的看到.
弟兄姊妹的愛.
但我仍然要說一句.
與哀哭的人同哭.
不容易.
是因為什麼呢?.
我們要站在別人的角度.
去了解別人的哀傷.
這樣就不一定了解.
可能在初期.
我們可以很容易表達.
我們的陪伴.
也是容易的.
到後來.
可能我們自己裡面.
很多聲音會出來.
甚至有些一開波.
就自己裡面的聲音出來.
我們會有什麼出來呢?.
很急於教訓別人.
是吧?.
很急於教訓.
你不要這麼難過.
為什麼這麼難過呢?.
這麼沒信心.
為什麼要這麼傷痛呢?.
你應該要怎樣做呢?.
我告訴你,要怎樣做呢?.
很急於教訓別人.
你想像一下.
你站在別人的角度.
其實自己最需要的是什麼呢?.
是否ABCDE這麼多的頂音姐妹.
過來教訓自己呢?.
A教訓完,B又再教訓.
C又再教訓.
當然,頂音姐妹不是想教訓.

$^{681}$是想愛的.
但你很心急.
讓我覺得是好的方法.
我覺得是對的看法.
我要加諸於你的身上.
你聽得更好.
情況會更好.
這是否安慰呢?.
這是否愛哭呢?.
你試試站在自己的角度.
是否這樣呢?.
我記得在我媽媽安息主懷的時候.
不只一次在崇拜.
我曾經說過自己的哀傷.
其中有一個.
以前教會的資體.
是很好的.
他也安慰我.
安慰時問我.
當時的情況.
說了沒幾句,他就問.
還有什麼不開心呢?.
我心想.
這是對我人生.
一個很重要的人離開.
其實我需要一段時間.
讓我自己的情感.
去過渡.
有時我不能用理性去解釋.
為什麼我的情感是這樣.
如果可以用理性去解釋.
這些問題很容易就處理.
所有事情都列出教訓.
這樣就行了.
不是的.
其實我們的神也給我們空間.
去做愛.
有時不是要去解釋.
那一刻聽下去.
其實我自己感受不舒服.

$^{721}$還有什麼.
會不開心呢?.
是不是?.
我心想,如果反過來.
你去經歷同一件事.
我就不會這樣問你.
還有什麼不開心呢?.
因為我知道.
就是會不開心.
是不是?.
有同理心.
站在別人的角度.
說起同理心.
這件事.
我自己也有很多功課.
要去學,是不容易學的.
說起同理心.
我記得我讀神學院的時候.
第一次神很深刻.
讓我認識.
我的同理心這方面.
要增強.
當時我神學院.
有個制度.
會跟學子.
修學妹.
或是跟學兄.
修學弟.
有學兄,學弟,學子,學妹的系統.
到我修學妹的時候.
當時有兩個學妹.
其中一個行動不變.
那時候.
她的腳.
走路拐下拐下.
甚至乎有時候要坐輪椅.
她在沙士期間.
她染上沙士.
後來有骨枯.
所以腳需要.

$^{761}$很長時間的醫治.
我讀建道神學院.
在什麼地方呢?.
有沒有人知道?.
長洲.
要上山.
她的腳行動不變.
要上上落落.
很不容易的.
除了在日常生活.
我會協助她.
很多買東西.
幫她搬東西.
有時也會陪她出入.
有些事物開始累積在她心裡.
她不太舒服.
事完.
我也很操心.
有時我做事的同時.
我又會說出口.
例如.
我下斜坡.
我在前面等你.
這句說話就死了.
我走前一點.
在那裡等她.
有時也會快一點.
我在那裡等你.
多說幾次.
她終於不開心.
其實你在那裡等就好.
不用說這一句.
你知不知道.
你不停提醒我.
我是有問題的.
你在那裡等我.
不用說就等.
不然就慢一點在我旁邊.
你不要說這一句.
我沒有站在她的角度.

$^{801}$去想這件事.
其實別人不喜歡你不停提醒.
例如那些說話.
我幫你.
你拿就拿.
不用說你幫我拿.
強調我是不行的.
在這些時候,我會很操心.
沒有同理心.
去明白別人.
也很想你明白我.
相處會舒服一點.
在她說完之後.
我就會反覆思考.
你幫人也好.
也要幫到恰到好處.
而不是令人不舒服.
還有些時候會想.
很急於叫那個人.
好像在某個位置.
走的路不正確.
或者用的方法不正確.
很容易會教訓別人.
其實可能別人也知道不對.
這樣教訓又幫不到別人.
又令她覺得很不舒服.
所以,這也是我自己.
經歷過.
覺得是撞過板.
和哭的人同哭.
不容易.
真真正正.
明白別人.
還有不同的人.
在生命裡.
有些很特別的.
專屬她的經歷.
我的狀況.
和她的狀況不同.
很多時候.

$^{841}$用我的狀況去衡量.
其實你沒理由會有這樣的感受.
你沒理由這麼不開心.
你其實不應該用這個方法處理.
但其實可能對方會覺得.
你也不明白我.
你試過我這樣的經歷嗎?.
你沒這樣.
為什麼要用我.
現在這樣的情況.
去教訓她.
或者去幫別人呢?.
對方覺得.
根本不識切.
還令我很煩擾.
在過程中.
更好的就是.
安靜的陪伴.
勝過不停地說話.
勝過你做很多事情.
我想起約伯的幾位朋友.
大家有沒有看過約伯記?.
約伯真的受了很多苦.
受了很多苦.
大家回去看看.
人生真的很難忍受苦.
他幾位朋友來探望他.
初期還在陪伴.
安靜地陪伴.
陪伴了不短時間.
陪伴了不久.
陪伴了一陣子.
就說了.
你一定有什麼問題.
一定是怎樣怎樣.
我們的上帝是公義的.
不會讓這些事情.
發生在你身上.
如果沒有問題.
神不會讓你經歷.

$^{881}$開始有很多教訓.
結果越來越激起.
約伯的心.
激起他幾把火.
你不要陪伴.
還好.
你走開還好.
所以不容易.
真的不容易.
弟兄姊妹.
我們一起想想過去的日子.
有沒有試過.
我們講解他的難處.
很快我們就教訓他們.
很快用自己的方法.
去勸人.
很快就這樣下了.
試過了沒有.
如果對方是信任我們.
開放跟我們說.
還好.
我知道對方感受不舒服.
對方一直不跟你說.
不過很不舒服很久了.
這會影響關係.
彼此相愛.
有些事情是累積的.
你不知道什麼時候.
那些不好的感受.
一直累積.
累積就越長.
不是那麼容易.
剛才喪禮講.
去到喪禮裡.
其實我覺得很重要.
除了我們的付出.
還有很多的陪伴.
很重要.
禱告.
這是重點.

$^{921}$你跟喜樂的人同樂.
跟哀傷的人同哭.
禱告.
每個人很開心.
有些事情很歡喜快樂跟你分享.
你站在他的角度.
想想他為什麼那麼快樂.
替他一起快樂.
你為他禱告.
為他感恩.
主啊.
為弟兄姊妹經歷.
這麼好的事情感恩.
主啊.
我求你繼續賜福給生命.
賜福給弟兄姊妹.
一直禱告.
在喜樂裡.
與人同行.
哀哭的禱告更加重要.
這個也有人提醒我.
不要那麼快出口.
這麼快說話.
為什麼.
說多錯多.
越快越死得快.
說之前禱告.
神啊.
今天這個人的情況是怎樣.
主啊.
你想我跟他說什麼.
我試過有一次.
我去醫院.
我去看醫生.
我試過有一次.
我去醫院探望人.
一路去之前求神.
神啊.
今天我去探望弟兄姊妹.
你想我跟他說什麼.

$^{961}$不知道他的狀況是怎樣.
你想我跟他說什麼.
神真的給我一個感動.
我上去就按照上帝跟我說的.
上帝的感動.
跟對方說話.
是不同的.
如果不是.
去到就好像.
基督徒有套餐.
ABCDE餐.
這樣的情況.
我就說這些說話一定不會死.
他說B的情況我就說B的說話.
不是的.
其實每一個人都是很獨立的個體.
每個人都不一樣.
神從來不是以套餐.
來對待我們.
神對我們是每一個.
單獨的.
狀況來對待我們.
不是套餐.
不是工廠.
禱告是很重要的.
心願未來的日子.
我會和你更多在禱告中.
經歷.
神啊 你想我跟他說什麼.
他需要什麼.
怎樣安慰他是最好的.
禱告.
不要太吵.
很嚴重.
他的情況.
快點 弟兄 姐妹.
我有很多話要和你說.
可能對方覺得.
弟兄 姐妹 你們需要安靜.
你來安慰我.

$^{1001}$我需要安慰你.
你的情緒很波動.
不是的.
對方很需要你安慰時.
我們情緒很波動.
我們情緒很緊張.
反而幫助不到別人.
我們不是需要這些.
我反而覺得.
對方需要神的聲音.
我們有沒有神的聲音.
還沒聽到可以先不說.
沒有東西比神的安慰更好.
其實人的聲音實在太多了.
是嗎?.
對我來說.
所以我覺得禱告很重要.
三個很重要的事情.
同理心.
從對方想.
明白了解.
他的快樂.
他的哀傷.
有很溫柔的心.
需要溫柔.
你的心溫柔很重要.
心很硬的時候.
你了解.
你也不會很想.
按照對方的需要去回應.
需要溫柔.
禱告.
必須要禱告.
很心願.
未來的日子我們都學到這個功課.
這樣才能夠.
跟喜樂的人同樂.
跟哀傷的人同哭.
這本程序表.
由某一天開始.

$^{1041}$每星期刊登很多.
弟兄姊妹或我們同工.
患病的情況.
或生命的需要.
請大家不要.
離開這裡就扔掉.
拿回來禱告.
禱告的時候.
不要像數白籃.
跟上帝說.
上帝是否真的不知道.
我要提醒你.
弟兄姊妹或同工.
有什麼情況.
提醒你.
不是這種.
你問問神.
弟兄姊妹或同工.
需要什麼.
願上帝.
按他的心意.
同工弟兄姊妹.
每次我面對.
同工弟兄姊妹.
上帝要我說的話.
給我們的幫助是色切的.
同哭同樂.
不是自己想怎樣就怎樣.
我跟自己說.
我都是撞板多過吃飯.
這個真的要成長.
成長就很不同.
我們將來.
彼此相愛.
就會更加好.
彼此相愛就是這樣.
教會就會越來越快樂.
我們的主是滿有智慧的神.
亦是愛的本身.
他知道怎樣愛更好.

$^{1081}$心願我和你從主的身上.
學習愛的功課.
我們一同禱告.
天上上帝求你憐憫我們.
其實過往的日子.
我們都很努力.
要去愛人.
但發覺我們真的有很多限制.
我們的愛有限.
甚至我們會說錯話.
做錯事情.
我們想愛人.
反而可能大被別人傷害.
壓力.
主求你寬恕我們.
幫助我們在這個地方.
要有更美好更豐富的學習.
將來.
真正懂得怎樣去愛.
真正和人同樂同哭.
主求你給我們每一個弟兄姊妹.
每一個同工.
每一個親密的關係.
我們從你身上得著愛的力量.
願你那個犧牲的大愛.
激勵我們.
去愛我們身邊的弟兄姊妹.
愛我們的家人朋友.
主啊.
當我們彼此相愛的時候.
別人就認得我們是屬於你的.
真是讓我們紅印家弟兄姊妹.
將來可以做這個美好的見證.
別人認得我們是屬於你的.
奉耶穌基督的名字祈禱.
阿們.
\newpage



\section{路加福音 15:22-32}
\label{sec:3jZxMYDYvEI}
\textbf{2023.11.26 《天父的擁抱》路加福音15\_22-32 (和修版) 陳韋安牧師}
\newline
\newline
連結: \href{https://youtube.com/watch?v=3jZxMYDYvEI}{\texttt{https://youtube.com/watch?v=3jZxMYDYvEI}} ~~~~ 語音日期: 2023-11-27
\newline
\newline
\hyperref[sec:3lGUeO3H154]{\small{< < < PREV SERMON < < <}}
~
\hyperref[sec:index_chronic]{\small{[返順時目]}}
~
\hyperref[sec:index_scriptual]{\small{[返順卷目]}}
~
\hyperref[sec:fKZdgGDua7A]{\small{> > > NEXT SERMON > > >}}
\newline
\newline
路加福音 15:22-32
\newline
\begin{longtable}{cl}
\hline
\hline
章節 & 經文 (和合本修訂版)\\
\hline
15:22 & \begin{tabularx}{0.7\textwidth}{X} 父親卻吩咐僕人：『快把那上好的袍子拿出來給他穿，把戒指戴在他指頭上，把鞋穿在他腳上， \end{tabularx} \\ \\ \relax
15:23 & \begin{tabularx}{0.7\textwidth}{X} 把那肥牛犢牽來宰了，我們來吃喝慶祝； \end{tabularx} \\ \\ \relax
15:24 & \begin{tabularx}{0.7\textwidth}{X} 因為我這個兒子是死而復活，失而復得的。』他們就開始慶祝。 \end{tabularx} \\ \\ \relax
15:25 & \begin{tabularx}{0.7\textwidth}{X} 「那時，大兒子正在田裡。他回來，離家不遠時，聽見奏樂跳舞的聲音， \end{tabularx} \\ \\ \relax
15:26 & \begin{tabularx}{0.7\textwidth}{X} 就叫一個僮僕來，問是甚麼事。 \end{tabularx} \\ \\ \relax
15:27 & \begin{tabularx}{0.7\textwidth}{X} 僮僕對他說：『你弟弟回來了，你父親因為他無災無病地回來，把肥牛犢宰了。』 \end{tabularx} \\ \\ \relax
15:28 & \begin{tabularx}{0.7\textwidth}{X} 大兒子就生氣，不肯進去，他父親出來勸他。 \end{tabularx} \\ \\ \relax
15:29 & \begin{tabularx}{0.7\textwidth}{X} 他對父親說：『你看，我服侍你這麼多年，從來沒有違背過你的命令，而你從來沒有給我一隻小山羊，叫我和朋友們一同快樂。 \end{tabularx} \\ \\ \relax
15:30 & \begin{tabularx}{0.7\textwidth}{X} 但你這個兒子和娼妓吃光了你的財產，他一回來，你倒為他宰了肥牛犢。』 \end{tabularx} \\ \\ \relax
15:31 & \begin{tabularx}{0.7\textwidth}{X} 父親對他說：『兒啊！你常和我同在，我所有的一切都是你的； \end{tabularx} \\ \\ \relax
15:32 & \begin{tabularx}{0.7\textwidth}{X} 可是你這個弟弟是死而復活，失而復得的，所以我們理當歡喜慶祝。』」 \end{tabularx} \\ \\
[1ex]
\hline
\hline
\end{longtable}
$^{1}$(字幕提供:陳零九).
鄧子每早晨,我們來到看浪子比喻的故事.
大家都應該很熟悉這個經文.
今天你發覺讀經的時候.
你發現經文是從故事的中後段開始.
從爸爸擁抱小兒子之後.
大家慶祝的故事裡.
開始今天的經文.
實際上在路加福音裡.
耶穌在第十五章說了三個有關比喻.
三個相關的比喻.
一個是失前的比喻.
一個是失羊的比喻.
最後就是失去了小兒子的比喻.
正常來說,耶穌說到這裡就應該完結.
因為失去兒子,就回來了.
失而復得.
不過你會發現耶穌在浪子比喻裡.
其實還沒完結.
我們現在正在故事裡.
看浪子比喻後面的大兒子的故事.
實際上當你看浪子比喻的時候.
耶穌故意延續這個故事.
不單將焦點放在小兒子浪子身上.
更加將焦點放在大兒子身上.
我會說浪子比喻.
其實不是單單說一個浪子.
而且你會發覺故事裡有兩個浪子.
故事的快樂的延續.
非常特別.
是兩個浪子比喻的轉折.
父親第二次這麼說.
他就跟父親說.
將上好的袍子快點拿來.
將戒指戴在他的指頭上.
幫他將鞋穿在腳上.
將肥牛毒宰了.
我們就可以吃喝快樂.
因為我這個兒子是死而復活的.
是失而有得的.

$^{41}$這個很典型的快樂結局.
如果大家看無線劇集的話.
這個也是第20集裡.
最後那五分鐘的橋段.
都是大團圓結局.
大家一起就完結了.
不過原來在這個快樂結局的背後.
原來出現了一個.
另一個浪子比喻的故事.
故事這麼說.
25節說.
那時大兒子在田裡.
他就回來.
離家不遠聽見作樂跳舞的聲音.
然後就變叫一個僕人來.
問是什麼事.
僕人就說.
你兄弟來了.
你父親因他無災無病回來.
就將這個肥牛毒宰了.
28節大兒子就非常生氣.
甚至可以說是非常憤怒.
這個正正就是浪子比喻.
在這個快樂的現實背後.
原來是大兒子心靈的風暴.
我們聽一下大兒子的控訴.
看看這個控訴背後究竟發生什麼事.
大兒子就說.
他就對父親說.
我服侍你這麼多年.
從來沒有違背過你的命.
你並沒有給我一隻山羊羔.
叫我和朋友一同快樂.
當你這個兒子和昌記吞盡你的財產.
他來了.
你倒為他宰了肥牛毒.
故事記載.
大兒子非常的憤怒.
他完全是參與不到.
父親和小兒子的快樂.

$^{81}$甚至當大兒子聽到他們這麼開心的時候.
他就完全崩潰.
你們不要再BBQ了.
你們不要再愛回家了.
他們非常的憤怒.
很明白大兒子的感受.
當別人愈是快樂.
當你參與不到的時候.
你就是愈是更加的感到失落.
如果你上網聽YouTube.
如果你打肥牛毒詩歌.
即是有人作一首這樣的詩歌.
專門描寫這段肥牛毒.
大家慶祝的情景.
可能你們也未必這麼容易投入到.
你們唱開.
我們唱開傳統聖詩.
他們那些又跳又Disco的燈.
又在跳.
你未必會感受到那種音樂和快樂.
所以大兒子的故事.
這種的.
另一個浪子故事正正就是.
快樂的延續背後.
當中去發展.
為什麼我會說大兒子.
也是一個浪子的故事呢.
如果說浪子比喻.
是說一個兒子離家出走的故事的時候.
其實大兒子.
也可以這麼說.
他也是一個像浪子一樣離家出走的人.
不過他把這件事是離家出走的意思.
他只不過是一個離家出走的浪子.
表面上他在一個家裡.
仍然在家裡擔當一切的大兒子.
但是他的心靈裡.
卻和浪子一樣.
一早就好像不在家裡.
感覺不到自己家裡的位置.

$^{121}$我們看看大兒子這段話.
他說我服侍你子孫當年.
從來沒有違背過你的命.
你並沒有給我一隻山羊羔.
讓我朋友一場快樂.
但你這個兒子和我來他窗戶吞進你的產業.
他來你倒為他宰了你的肥牛丼.
令人注意的就是.
這段對話裡出現了兩次沒有這個字.
我沒有違背過你的命.
你並沒有給我一隻山羊羔.
中文這部分可能看得不清楚.
但在原文裡.
沒有這個否定的字是很強烈的.
不是一個普通的沒有.
而是我們英文裡.
never一個這麼強烈的字眼.
兩次的never.
我們看看這兩個的never.
第一個never.
大兒子說什麼.
他說從來沒有違背過你的命.
我沒有違反你的命令.
雖然大兒子在家裡是一個長子大兒子.
但原來他從來都不覺得自己是家裡的兒子.
他覺得自己是什麼.
是一個執行命令的人.
一個不可以違背爸爸命令的僕人.
大兒子說我服侍你這麼多年.
服侍這字其實正正就是羅馬社會裡的奴隸.
那個字.
那個工人哥哥的角色.
所以很淒涼原來大兒子.
他一直都覺得自己是一個家裡的工人.
一個執行爸爸的主人命令的奴隸.
不是爸爸的長子.
大兒子覺得自己人生裡.
只不過是一個work from home的奴隸.
不覺得自己是一個兒子孩子的角色.
第二個never.

$^{161}$大兒子說什麼.
他說你從來沒有給我一隻山羊羔.
第二句話告訴我們.
原來大兒子覺得自己是很貧窮.
他覺得自己什麼都沒有.
甚至這種一無所有.
比小兒子當他極度貧窮的時候更加貧窮.
雖然大兒子住在一個極度富有的爸爸的家.
但他不覺得他身邊的東西是他的.
他是打工的.
他是奴隸.
那些東西不是他自己的.
所以你覺得誰比較窮.
大兒子比較窮還是小兒子比較窮.
照這樣說的話.
其實小兒子雖然很窮.
但他曾經都花光了.
都曾經花了錢才窮.
大兒子他一分都沒有花過.
他覺得自己是一無所有.
覺得自己連一隻山羊羔都沒有吃過.
所以你發覺大兒子其實比小兒子更加窮.
他連坦誠反叛離家出走的勇氣都沒有.
一直都在家裡做一個work from home的浪子.
事實上最遙遠的距離.
可能就是我們每天都在天父上帝的世界那裡.
卻感覺自己一無所有.
不覺得自己是天父上帝的孩子.
所以大兒子沒有辦法去擁抱他弟弟.
不行.
試問一個更加破碎.
一個更加悲慘的人.
怎麼能夠擁抱一個.
能夠有幸被修補的人呢?.
當大兒子比小兒子更加痛苦的時候.
怎麼能夠苛求他.
去擁抱他.
替他快樂呢?.
不知道這是不是你的感受.
明明知道這是天父的世界.

$^{201}$天父在這裡.
卻有時感覺不到那種力量.
明明已經是基督徒.
卻感覺不到比其他人更加快樂.
或者更加堅強.
明明每個星期都回到教會裡.
卻未必感覺到有人去理會自己.
或者在這裡.
面對著這個世界.
特別是這個年頭.
我們怎麼去面對這個世界呢?.
這個世界越來越亂.
這個社會越來越不穩定.
甚至自己的身體.
自己面對的命運.
似乎都跟自己所想的.
未必是如期一樣.
所以當我們去看回這段經文的時候.
我們很值得去看下浪子的比喻.
這個更加棘手.
更加難解決的大兒子.
爸爸怎樣來回應呢?.
爸爸怎樣回應呢?.
爸爸就在經文裡.
在三一至三二節裡.
這樣來回答他的兒子.
爸爸說:兒啊!你尚和我同在.
我一切所有的都是你的.
只是你這個兄弟是死而復活.
失而又得的.
所以我們理當歡喜慶祝.
總的來說.
爸爸就鼓勵哥哥.
叫他擁抱他弟弟.
爸爸就給了一個理由.
鼓勵給哥哥.
說你應該欣然.
快樂地去擁抱你自己的弟弟.
基本上坦白說.
這段話其實沒什麼特別.

$^{241}$爸爸沒有說一些新的事情給他聽.
也沒有新的資訊.
也沒有新的秘密告訴他.
爸爸大概說了一些.
誰不知道爸爸是男人的事情.
一些他也知道的事情.
兒啊!你尚和我同在.
我一切的都是你的.
父親稱呼他為兒子.
是一個很好的肯定.
不過他說這些都是你的.
我也是尚和你同在.
大姐似乎也不是不知道的.
這麼多年我都在家裡.
我一切的都是你的.
我也知道是.
但是我感受不到這些是我的.
所以不知道.
當我不斷思考這段經文的時候.
作為一個牧者.
如果突然之間聽到.
天父上帝這樣跟他說.
這句話是不是真的能夠幫助我們.
去面對一個這麼艱難的世界.
或者這麼艱難的自己的命運呢?.
這些道理大家都知道.
都是正確,很熟悉的.
每一句話都是真的.
我們沒有什麼反對的理由.
但是對於大兒子來說.
他究竟怎樣可以真正去理解這番話.
不過他發現耶穌就停在那裡.
整段浪子的比喻就停在這段經文裡面.
耶穌就再沒有去交代大兒子的回應.
經文就停在那裡.
究竟大兒子是否明白爸爸的心意.
或者能不能夠明白這句話背後的意思呢?.
我們不知道.
不過我們作為看聖經的人.
我們就需要更加明白.

$^{281}$這句父親的話背後是什麼意思呢?.
爸爸說:.
「兒啊,你尚和我同在,我一切所有的都是你的」.
在浪子的比喻裡面.
爸爸就是代表著我們的天父上帝.
他說:.
「兒啊,你尚和我同在,我一切所有的都是你的」.
這句話其實不是普普通通一句.
很表面的安慰的說話.
就是跟大兒子說:.
「你看看你多好,你這麼多年都在家裡享受冷氣」.
「不用在外面工作,你多好」.
「你看看你弟弟多慘」.
「你快點去接受他,去擁抱他」.
不是,爸爸不是要他為一個哥哥去接受這段關係.
爸爸是很想知道他的擁抱和快樂.
究竟是從哪裡來的.
或者應該怎樣來參與這份擁抱和快樂.
爸爸說:「你尚和我同在,我一切所有的都是你的」.
爸爸說一個更深層次的意思.
當你面對著你的遭遇,面對著這個世界的時候.
你只能夠參與我的一切來擁抱這個世界.
你只能夠去擁抱我一直擁抱的世界.
你只能夠跟我一起去快樂.
在我們的聖經裡經常有這個模式.
我們所做的其實是天父上帝一直都在做的.
我們才去做的.
這個模式.
譬如我們愛,因為神先愛我們.
你們要去,因為我已經為你們預備了地方.
你們要背十字架,因為耶穌早就背了十字架.
我們信,因為耶穌基督的信.
我一切所有的都是你的.
今天天父上帝和弟子說的不是叫你自己去勉強去擁抱一件事.
而是我們去擁抱一些天父上帝去擁抱的東西.
我們去快樂不是叫你很悲慘都要你去快樂.
而是我們去嘗試去參與天父上帝的快樂.
所以,我們理當歡喜快樂.
這種快樂是參與我們天父上帝的快樂.
我一切所有的都是你的.

$^{321}$這個正正就是我們今天要思想的主題.
擁抱一切.
擁抱當然是一個很重要的事.
面對著這個世界的時候.
我們怎樣去擁抱這個世界.
甚至擁抱我們的遭遇.
我們身體的疾病.
我們家人的疾病.
我們所面對的命運和際遇.
怎樣能夠擁抱呢?.
不知道大家有沒有聽過尼采.
一個德國哲學家尼采.
可能大家聽他比較多的都是一些比較負面的.
因為他是一個反對基督教的人.
提出上帝已死的哲學家.
他是一個虛無主義者.
即是說他覺得生命是沒有意義的.
沒有上帝,沒有美善,沒有價值,沒有真理.
一切都是很虛無的.
沒有所謂絕對的東西.
不過雖然他是一個虛無主義者.
但是他對於生命或命運絕對不是悲觀消極.
他提出一個很特別的說法.
叫做熱愛命運.
他這樣寫到.
熱愛命運.
讓這句說話成為我一生的所愛.
我不會向醜陋宣戰.
我不會指控醜陋.
我也不會指控他的人.
無視他們是我唯一的否定.
總而言之我要成為生命的肯定者.
提出這個說法是什麼意思呢?.
既然他覺得活著是虛無的.
生命是沒有意義的.
那麼對他來說唯一拿著的是什麼呢?.
就是他自己的生命.
雖然是虛無的.
但是他唯一知道要堅持擁有的就是生命.
所以他知道要熱愛命運.

$^{361}$無論他的生命遭遇到的多麼的醜陋.
多麼的痛苦.
多麼的絕望.
他都要成為他生命裡擁抱的對象.
簡單來說就是我們以前Nokia手機的貪吃蛇的概念.
一看到蛇就不斷吃那些豆豆.
越吃越長越吃越大.
總之生命裡遇到什麼.
我就擁抱它.
我就吞噬它.
成為我生命裡的強大力量.
所以彌采有句很有霸氣的話.
可能大家聽過.
凡殺不死我的都必使我更加強大.
當然我們不認同彌采的看法.
我們是一個信上帝的人.
不過我想如果一個不信神的人.
他都可以毅然去擁抱他自己的遭遇.
我們作為一群基督徒.
我們作為一群相信天父的愛.
和天父上帝擁抱的愛的時候.
我們豈不會更加有力量和勇氣.
去擁抱我們所面對的命運呢?.
當我們面對我們人生裡的遭遇的時候.
你可以去逃避它.
可以去不開心的去苦苦哀求.
也可以用很多不同的方法去面對.
但是擁抱是一種很強大的霸氣.
我們廣東話有一句話.
我擁你.
當你做事的時候.
你上師說我擁你.
或者你作為上師說我擁你的時候.
一句非常有霸氣的話.
什麼意思呢?.
無論有什麼事我都會承受的.
我都會傳言去抵擋的.
這種接納.
這種完全的擁抱投入.
是一個很大的霸氣.

$^{401}$因此面對著我們這個世界.
我們一群基督徒.
我們願意去擁抱我們所面對的遭遇.
我們未必認同.
更加不是傻乎乎的去認同這個世界裡不好的事.
但是我們去擁抱我們所面對的遭遇.
更加用一份勇氣和信心去面對.
一個非常重要的事情.
特別是你知道擁抱其實不容易.
甚至乎是有些危險的.
我們看80年代的港產片.
你看周潤發擁抱的人通常會怎樣.
有六成機會都是會怎樣做.
人家就會在肚子裡開兩槍.
插兩刀.
擁抱是一件非常不容易的事情.
因為你有信心去躺開你的手.
完全去接受.
甚至乎去擁抱.
但是因為天父上帝的雙臂.
緊緊擁抱這個世界.
你知道這個世界.
正正是天父上帝的世界.
更加是去擁抱這個世界的時候.
我們一班天父上帝的兒女.
正正有這份本錢.
同樣去擁抱這個世界.
我一切所有的都是你的.
因此面對著這個命運.
我們都去學習去面對這個命運.
我爸爸在暑假的時候.
六月尾發現胃有血.
以為是胃潰瘍.
發現原來有些癌細胞在上面.
一開始以為只是切一點點.
後來因為在上面的位置不能切一點點.
整個位置都要切除.
很快六月尾知道.
七月檢查.
七月尾就要去動手術.

$^{441}$很記得和他吃最後晚餐.
他有胃之前的最後一餐.
在汀水圍吃一人火鍋.
一起吃最後一餐.
切完之後很快.
就要過一個無謂的生活.
由食道直接到小腸.
每一口吃的東西.
營養奶或其他食物.
一吃下去都很不舒服.
身體很痛.
腸胃很痛.
爸爸最喜歡吃東西.
突然間不能吃東西.
人生沒有了90\%的樂趣.
當我去暑假和他去看病的時候.
香港政府的醫療系統都很物流式.
基本上大家人就拿著皺紋.
在那裡分流看.
營養師看外科.
每人看五分鐘就走出來.
再去看別的地方.
當你在排隊看皺紋的時候.
你發現原來在香港.
和你家裡差不多租慾的人很多.
多到整個街市都一樣.
皺紋就很熱鬧地坐著等待看醫生.
基本上我們就是這樣去面對.
沒有胃就沒有胃.
這件事就沒有了.
人生之後的日子.
就是這樣慢慢去學習.
適應這樣的生活.
你可以覺得很悲慘.
但同時間你知道.
你只能在天賦的懷抱裡.
去擁抱你所面臨的遭遇.
我想大家都面對很多類似的情景.
面對我們人生裡.
或者家人裡的一些遭遇.

$^{481}$我們似乎不能改變.
我們只能憑著信心.
帶著一點快樂.
去擁抱我們面前的遭遇.
不過大家聽過這個故事.
一位稱為包山王的黎智偉的故事.
黎智偉在十幾年前.
是一個人稱為包山王的香港運動員.
最高曾經是全球排名第八名的攀石運動員.
2011年有一次他在屯門公路開電單車的時候.
突然有一輛超速的私家車撞向他.
然後他撞向隔離線.
然後突然有另一輛車再駛過去.
令到他整個人的脊椎.
從此下半身癱瘓.
他也想不到自己當年28歲.
從此終結了他運動員的生涯.
下半世都要在一個冷冰冰的輪椅裡度過.
當然有很多無數的眼淚.
很多的困苦.
很多的痛苦.
幾年之後我們就知道黎智偉決定擁抱他的命運.
甚至乎擁抱他一直要坐著冷冰冰的輪椅.
今天沒有帶這張照片出來.
大家可以看看.
自己搜索一下.
可能大家見過.
就是他抱著他坐的輪椅.
再一次來攀石.
一個很震撼的圖畫.
雖然他要坐輪椅.
但他願意來擁抱他的命運.
繼續向前走下去.
所以我想我們都是這樣.
面對著我們人生裡的困難.
這個世界,這個社會.
我們自己很多不同的遭遇.
我們嘗試去學習去擁抱.
不是因為你擁抱到.
我們去擁抱的是天父一直擁抱的世界.

$^{521}$我們的命運就交託在上帝的手裡.
我們一起祈禱.
因為你是擁抱著你所愛的世界.
你去快樂地去面對這個世界.
因為知道這個世界無論多麼黑暗,猖狂.
但都仍然在你這份懷抱裡.
所以我們作為你的兒女.
我們求主你幫助我們.
能夠可以同樣地學習.
去擁抱你所擁抱的東西.
我們都將我們手裡.
或者眼前見到的很多問題.
我們解決不到的問題.
我們教會裡的弟兄姊妹.
老牧師等等.
我們都交託給你.
我們能夠可以一起去擁抱.
去憑信心去面對我們所面對的.
我們都知道我們無論遭遇什麼.
都仍然在你的手裡.
求主你讓我們能夠快樂.
能夠堅強.
能夠有信心去發現你給我們的恩典.
奉主名求.
阿們.
好了.
謝謝大家.
\newpage



\section{詩篇 126:5-6}
\label{sec:fKZdgGDua7A}
\textbf{2023.12.03 《誰是收割者?》 詩篇 126\_5-6 (和修版) 謝靄薇牧師}
\newline
\newline
連結: \href{https://youtube.com/watch?v=fKZdgGDua7A}{\texttt{https://youtube.com/watch?v=fKZdgGDua7A}} ~~~~ 語音日期: 2023-12-07
\newline
\newline
\hyperref[sec:3jZxMYDYvEI]{\small{< < < PREV SERMON < < <}}
~
\hyperref[sec:index_chronic]{\small{[返順時目]}}
~
\hyperref[sec:index_scriptual]{\small{[返順卷目]}}
~
\hyperref[sec:7y90EbIBR_w]{\small{> > > NEXT SERMON > > >}}
\newline
\newline
詩篇 126:5-6
\newline
\begin{longtable}{cl}
\hline
\hline
章節 & 經文 (和合本修訂版)\\
\hline
126:5 & \begin{tabularx}{0.7\textwidth}{X} 流淚撒種的，必歡呼收割！ \end{tabularx} \\ \\ \relax
126:6 & \begin{tabularx}{0.7\textwidth}{X} 那帶種流淚出去的，必歡呼地帶禾捆回來！ \end{tabularx} \\ \\
[1ex]
\hline
\hline
\end{longtable}
$^{1}$各位弟兄姊妹平安.
很開心見到你們.
我們好像經常見面.
在我們看神話之前.
我們一起祈禱.
愛你主 愛你神 我們每一生都歡喜快樂.
因為今天是你的日子.
我們要歡喜快樂地.
去到你面前敬拜你.
願你在當中帶領我們的心.
讓我們被聖靈感動.
讓我們在看你的話語的時候.
懂得摸著你的心意.
願孩子心中的意念.
口中的言語是夢裡如立.
奉主耶穌基督的名求.
阿們.
再過三個星期就是聖誕節.
這個好像屬於基督教的節日.
我們會叫做張臨期.
但是全世界都知道.
當我們看到聖誕節.
就知道是聖誕節.
我們看到聖誕節就對的時候.
做生意的人.
就很聰明.
他們想到很多方法.
讓他們可以賺一筆.
其實耶穌基督.
真正出生的日子.
是沒有人知道的.
因為當時.
人們看祂是一個木匠的兒子.
根本沒有人去留意祂.
直到一些學術家.
就說.
祂是一個木匠的兒子.
直到一些學術家.
或者一些.
帝王的時期.

$^{41}$開始注目的時候.
不過已經是.
沒有辦法去.
去稽考.
祂到底是哪個日子.
因為沒有記錄下來.
所以.
我們想.
到底神叫耶穌基督.
由天來到地上.
道成肉身.
到底祂的目的是什麼呢?.
其實很清楚.
就是祂是我們的救主.
祂要讓.
獨生兒子來到地上.
要傳揚天國的訊息.
所以我們看到.
原來.
耶穌基督成為救主.
在一般的人.
全世界的人都知道.
不用說.
但真正的日子不知道.
但我們很清楚.
當我們看到這個圖畫的時候.
世人都會想到是聖誕節.
想到一幕.
但其實.
如果再留意經文的時候.
是有些不同了.
既然神要派獨生兒子.
來到地上成為救主的時候.
原來這是.
神的應許.
在《耳熟眼書》61章.
1到2節.
就清楚知道耶穌來有一個使命.
就是叫盲的看見.
叫.

$^{81}$被老的得到釋放.
軟弱的得到強壯.
我們看到.
神有祂的心意.
甚至看到耶穌基督.
是怎樣為我們犧牲在十字架上.
流出寶血.
為什麼?.
祂就是令我們這班.
沒辦法得救的人.
耶穌要救我們.
令我們不要沉淪.
所以祂要來地上.
甚至為我們寫命.
當祂要離開的時候.
要升天的時候.
祂就祝福他們.
特別是那些人會很害怕.
特別是那些門徒.
我跟了耶穌這幾年的光景.
祂忽然間.
走了.
我還沒有學到東西.
怎麼辦?.
我們要休守.
祂是為我們預備地方.
然後祂要他們.
大使命.
要去.
我怎樣去?.
師父走了.
通常你們學習的時候.
起碼老師都在旁邊幫你.
怎麼辦?.
在這個經文就很重要.
但聖靈降臨在你們身上.
你們就必得著能力.
不用怕.
聖靈降臨的時候.
你們就可以剛強壯膽.

$^{121}$因為祂給了你們能力.
要在耶路撒冷.
是他們的中心點.
然後要發散出去.
猶太全地,撒瑪利亞.
甚至是直到地極.
要做什麼?.
做耶穌基督在你身上的見證.
要繼續傳講.
當我看到這個經文的時候.
心靈可以定一定.
你可以有能力.
是很大的.
祝福我們.
神給你這個能力的時候.
今天我們在座.
你已經信了耶穌.
撒瑪利亞.
你今天.
問問你.
問問你自己.
耶穌基督給我們一個使命.
你做了多少?.
你開始去了沒有?.
這個.
大使命裡面要我們去.
去第一步,你去了沒有?.
去之前有些事要做.
做什麼?.
要裝備.
你裝備了沒有?.
不是出去就啞口無言.
不知道說什麼.
你去了沒有?.
耶穌叫我們要去.
使人做祂的門徒.
你只要帶一個人信主.
等很久.
你聽見見證.
甚至你信了主.

$^{161}$叫你洗禮.
你都要等很久.
你不會這麼快.
有些人馬上.
不聽主的吩咐就去.
但有些人不是.
有些人等到他.
回到天間的時候.
他想起有件事沒做.
我們不要等那麼久.
我們不知道有多少壽命.
如果要等.
等到那一刻.
你會覺得錯失了很多機會.
可以服事主.
起落弟兄姊妹都知道.
服事主很開心.
但我們不要等,我們要去.
我們信了主,你起步了沒有?.
耶穌叫我們要遵守.
祂教我們的.
起碼第一步要去.
去之前要好好準備.
然後你就可以很開心快樂的.
將見證告訴別人.
所以我們知道.
今天經文很短.
相信你會背.
希望你走的時候記得.
第五節第六節說什麼.
「流淚撒種,必歡呼收割」.
這個很開心的經文.
在《詩篇》126篇.
所以我們首先去看看.
誰是收割者呢?.
當我們讀這個.
這麼開心的經文的時候.
誰是收割者呢?.
這個收割者.
讓我看到.

$^{201}$原來在經文裡.
他告訴我們.
你們不是說還有四個月嗎?.
但耶穌說什麼呢?.
收割的人.
已經得功錢.
為永生儲存五穀.
使到撒種的.
和收割的.
一同快樂.
人家看到是四個月.
但耶穌要他看到熟成的眼光.
原來已經是熟透了.
好像你家裡種的東西.
熟透了.
怎樣呢?.
爛了.
其實還沒爛之前.
小鳥都先吃了一口.
他提醒你可以吃了.
你要快點收割.
要搶著收.
但神給你機會去做這個事奉.
你要去收割.
要搶著去收.
不收就爛了.
所以這個時間要把握.
你才可以做到.
所以現在是一個.
收割的時候.
你準備好了嗎?.
當耶穌基督.
要差派你.
做這個工人.
要去收割的時候.
原來.
做這個工夫很辛勞.
甚至我們好像.
在享受.
別人辛勞的成果.

$^{241}$好像坐享其成.
但沒辦法的.
有時候可以生下一個時代.
前人種樹.
後人就成樑.
我後人出世.
前人做了很多成就.
我們就可以享受.
香港每年都有不同的建設.
年輕人.
就要享福.
很多科技.
現在的手機多方便.
以前的時代.
所以現在老人家要學.
科技.
以前哪有這麼方便.
一按就有你心目中所要的.
你想看哪一條片都很方便.
但以前沒有.
他們很努力去耕耘.
種東西.
種東西很辛苦.
有些人在家裡.
習慣冷氣.
你叫他出去曬曬太陽.
他不願意做這個工作.
所以做農夫的工作.
是很辛勞的.
當我們看到做了收割的人.
我們又不想.
佔別人便宜.
我們又不想坐享其成.
那誰是殺種的呢?.
誰是收割的呢?.
可能左手是做殺種的.
右手是做收割的.
可能反過來.
右手是做殺種的.
左手是做收割的.

$^{281}$都可以.
同時期.
你也是殺種的,也是收割的.
但是.
其實紅恩堂.
我不是這麼熟悉的.
我想問問你們.
你們剛剛過了生日.
你們.
幾久的日子?.
12還是11?.
12年.
是12年,沒錯.
12年了.
那誰是殺種的呢?.
誰是收割的呢?.
我不知道的,但是有一個知道.
這個創造我們的知.
是他知道.
知道你心中的意念.
甚至因為.
他是賜給我們生命的.
他也是收回我們生命的.
所以他賜給你生命的時候.
他也會傳福音給你.
要你做見證.
你說我不懂說.
但起碼你有見證.
你不會沒有見證的.
你不會失見證的.
因為你是.
屬神的兒女.
起碼你的故事.
你的生命故事可以跟別人分享.
這是第一步.
你懂得說你的生命故事.
已經在傳福音.
讓別人看到你為什麼不一樣.
為什麼你會坐在這裡.
為什麼我每次叫你的時候.

$^{321}$你都說我要做禮拜.
當別人問你.
你去哪間教會的時候.
已經開始了.
或者你在茶樓吃東西的時候.
你在祈禱已經在做事了.
這班基督徒是不同的.
其實你在這裡.
開始為主工作.
在做見證.
所以當我看到.
其實你們做的功夫.
應該比我多一點.
你們傳福音應該比我多一點.
你去到人群當中.
那些主婦就最厲害.
看到很多婆婆.
有很多事情可以說.
那時候你就可以做見證.
你認識的朋友比我多.
你每個人.
帶十個.
是不是比我厲害.
我一個人.
我一個人做不了多少.
我最厲害就做十個.
但你一個人做十個.
你再做也不止十個人.
你加起來.
也很多.
所以你要去.
不要嫌棄自己什麼都不會做.
起碼你有口.
你的口就可以說見證.
你可以為人祈禱.
也可以做.
我們來聽聽.
這個殺種的.
這個收割的.
我差你們去.

$^{361}$收你們所沒有的辛勞.
你們享受他們辛勞的成果.
今天我們可以在這麼乾淨的地方.
我們可以享受.
我們可以聽到.
但我們不要單單聽到.
我們要出去.
雖然前人很辛勞.
他們做了的我們要繼續.
他做了的我們不要停在那裡.
停在那裡就沒有了.
雖然以前很辛苦的.
我們要繼續.
做得更加完善.
流淚殺種的人.
必歡呼收割.
所以我們看看.
收割者的代價是什麼呢.
就是忍耐.
就是要等人.
等一分鐘.
兩分鐘也不要緊.
可能你等你女朋友.
等超過半個小時.
你就會開始發牢騷.
但今天我們看看.
這個種東西.
可能一年半分鐘.
可能一年裡面的心血.
可能一風一泥.
可能做生意的.
我的心血沒有了.
我的錢賺不了.
我放了這麼多本.
要忍耐的.
要去忍耐做一些事情.
你知道我們《聖經》裡面.
有講到殺種的比喻.
又在好土的.
又在荊棘的.

$^{401}$又在慈善的.
又在荊棘的.
又在錢的石頭的.
又被鳥擔了.
殺種也不容易.
你又不可以說.
我只要在這個地方.
你一殺就這樣去.
你怎麼知道你去到哪個角落.
但是這個殺種裡面.
你也要花點心機去做.
這個東西.
你殺出去了.
但是風一吹.
可能就沒有了.
鳥擔了又沒有了種子.
成功率就是一把風.
你就計算不了.
流淚殺種.
我想問你.
流過多少眼淚.
為自己.
自己的失敗.
可能會懺悔的.
但是為了神的事工.
我們流了多少眼淚呢.
我們有沒有去想一下.
哪個弟兄姊妹的需要.
我們有沒有去為她祈禱.
可能我們.
去數一下的時候.
可能流淚殺種的.
我們可能還沒有到.
耶穌的標準.
給我們看到還有一個距離.
但我們可能有為了.
自己的需要.
或者我們遇到一個挫敗.
我們遇到痛苦.
或者我們遇到患難當中.

$^{441}$或者我們身體的需要.
我們都是一根刺.
就是你會痛苦.
痛苦的時候怎麼辦呢.
原來我們痛苦的時候.
我們可以祈禱.
有沒有想到.
我們的幫助是從哪裡來的呢.
是從創造天地的神而來.
你要想到你要祈禱.
雖然你收割的過程很辛苦.
你要忍受.
可能人家會嘲笑你.
像羅亞的時代.
整個方舟天氣這麼好.
但我們聽神的話的時候.
我們就要去做殺種.
我們就要按時按後.
去做應該做的事.
令到果實或者你所種的東西.
可以有收成.
有忍耐.
如是者每天都要去做.
去祈禱.
你在這個患難的裡面.
你在逼迫的裡面.
特別在疫情裡面.
遇到很多不同的事.
遇到很多不同的遭遇.
你自己都好像中了招.
那時候沒有現在這麼好.
現在中招還可以外出.
以前是不行的.
可能封住你的樓層.
人們不讓你外出.
很多的艱難.
有些心裡想做的事.
或者你和外面的人.
很難人們做生意.
要和外面有交流.

$^{481}$飛機不飛.
很多事情都做不到.
很多事情停頓了下來.
要忍耐.
很長的日子要忍耐.
到現在疫情還在.
還有不同的事要發生.
我們不知道的.
但我們還要忍耐.
因為我們的目標還未到.
但我們做不到的時候.
我們很艱難的時候.
就是祈禱.
祈禱是很重要的.
很多時候我們知道.
祈禱是很重要.
但祈禱會有多少人?.
當然我們有教會的事工.
很熱心去做的人很多.
或者我們可以奉獻.
甚至有些人.
我不出名.
我都是默默的去奉獻.
很多人.
但是祈禱有多少人呢?.
有多少人真的願意.
放一些時間.
去做守望的功夫呢?.
守望的功夫是最重要的.
有沒有看到?.
原來神給我們看到.
要復興是要祈禱的.
甚至我們看一些比喻.
原來祈禱.
好像發電機一樣.
很多人在聚會之前.
默默的去祈禱.
去守望.
為了別人的需要去禱告.
然後神就開始動功.

$^{521}$然後他們就崇拜.
他們崇拜.
人們就很熱心.
不單是聽完禱告.
他們還很火熱.
想著他們去到人群當中.
他們要去跟人傳福音.
為什麼呢?.
因為他們有人會去祈禱.
甚至自己都去祈禱.
他們是一個火熱的心.
於是神就去復興教會.
我們紅印堂可以很興旺.
祈禱很重要.
我們去祈禱.
我們要知道要怎麼做.
所以我們看到耶穌.
在聖經裡面.
他不是去強調講道.
他強調什麼呢?.
他每一次都安靜的.
退下就去祈禱.
耶穌時時都這樣做.
耶穌這樣做.
我們是凡人.
我們是不是更應該退下.
我們要去祈禱.
向主求能力.
所以祈禱很重要.
保羅也教我們.
在提摩太后書裡說.
要為萬人代求.
今天我們看到.
新聞裡面發生很多的事.
你有沒有為這個世界去祈禱呢?.
你有沒有為在上掌權去祈禱呢?.
下個星期又選舉了.
你有沒有為這些事去祈禱呢?.
你有沒有為你心目中想要選擇的人去祈禱呢?.
會不會現在呢?.

$^{561}$是講一個版本.
上了台呢?.
我們不知道的.
我們是不是為這些事去祈禱呢?.
為一些很小事.
很靠近的事.
你有沒有去祈禱呢?.
我們要求神幫助我們.
才破碎我們自己的生命.
才可以祝福別人.
好像耶穌基督.
祂自己也是破碎了自己.
然後去成就這麼拯救世人的工作.
我們要破碎自己.
所以呢.
除了忍耐,祈禱.
還有什麼呢?.
裝備.
裝備很重要的.
每一個修過的人.
要裝備的.
聽到裝備就很害怕.
聽到培訓.
要去訓練.
那些人不喜歡上課的.
你告訴我吧.
說個答案給我聽吧.
不想自己去花點時間.
裝備很重要的.
你學每一樣東西.
都要花時間.
你想想你小時候.
做學生都要花時間.
今天我們可能教你傳呼音.
不用這麼多年.
是不是?.
相信你們在珍古樑.
你們有些教你們傳呼音的方法.
不用讀幾年書.
是不是?.

$^{601}$很快的.
可能都不用兩個月.
就完成了.
但完成之後要做什麼呢?.
這件事要繼續做.
你不做.
丟下了.
白費了又要重來.
因為要操練.
所以我們知道.
在醫院裡.
做醫務人員的.
隨時都有人去急症室.
急症室.
急症室做什麼呢?.
當然有需要去.
有時有些人.
覺得你沒死就坐久一點.
繼續等.
生命危險的.
會先救他.
急症室.
很多時候有些車禍鬧事的人.
會去當中.
可能這些人會多一點.
會搞傷.
所以醫務人員很忙.
可能有些小小的刺傷.
他不太理你.
因為你沒死.
先救多一點危險的.
醫生去診斷的時候.
他來到你面前.
他會不會.
找我書的時候.
他會不會這樣.
很有信心.
如果那個人來到你面前.
找書的時候你對他沒信心.
你快點走.

$^{641}$他不會慢慢研究.
這個病症.
我應該看哪條方.
其實他心目中已經知道.
感冒就應該給什麼藥.
他已經很清楚.
除了你那些.
癌症.
他要特別找一班人去研究.
但通常去到急症.
那些不是那麼嚴重.
應該是他應付得到.
醫生不會慢慢.
再看書本去研究.
所以我們看到.
裝備.
你平時的裝備.
到時急到來你面前.
可能你以前是做小兒科的.
但那晚你當值.
你什麼科.
來到你面前你都應付.
所以有些基本的東西.
其實醫生可以做到.
他就會去應付這班人.
不知道他是來自什麼病症.
他都要看.
所以你看到這些就是.
平時的操練.
平時讀那麼多書.
到時可能就是這個時候.
就需要了.
我們讀聖經也是.
我們不會.
到時去探望病人.
讀哪個經文.
其實你平時靈修.
隨時都有很多經文.
你便寫下它.
祝福人.

$^{681}$到時便可以和人分享.
或者你研經.
這些都是時間.
都是慢慢浸出來的.
培訓都是你受過訓練之後.
你就要去用.
你需要時間需要操練.
這些就是裝備.
你開始裝備了嗎?.
或者你已經裝備了.
你有沒有操練呢?.
曾經在打仗的時候.
有個實習醫生.
他和人群一樣走難.
在急難的時候.
有一個孕婦肚痛.
那些人就知道.
他是實習醫生就叫他.
你知道他平時已經操練了.
只不過他還沒有拿牌.
他在急的時間.
可以幫助孕婦.
幫她去接生.
所以這些就是平時的操練.
你一定要裝備.
你不裝備的時候.
隨時來到你身邊的時候.
其實神就要用你.
如果你平時都不去準備.
你就會很害怕.
但你平時都去操練的時候.
你不用害怕.
神就要祝福你.
與你同工.
你就可以大道一緊順.
可能祂在憂愁當中.
需要你為祂祈禱.
如果你平時都沒有祈禱的時候.
你就會很害怕.
要說什麼呢?.

$^{721}$你就會很緊張.
但你平時都祈禱的時候.
你想到很多神的話語.
你就可以祝福祂.
其實這個是要操練.
要裝備的.
所以這個操練是平時慢慢積累的.
不是我現在要去.
今天我們要去.
哪個醫院.
我們要去派單張.
我現在才開始去準備.
等等我見到那個人.
我要怎麼說呢?.
不是這樣做的.
平時你就要去神面前安靜.
要去靈修.
這是很重要的.
修過者.
祂有什麼喜悅呢?.
當我們今天看經文的時候.
祂有很多喜悅.
首先.
在喜悅看到祂們的希望.
譬如在第126篇.
第一節裡面.
當耶和華在錫安被擄的人.
歸回的時候.
我們很像做夢的人.
那時我們滿口喜笑.
滿舌歡呼.
那時列國中就有人說.
耶和華為他們行了大事.
行了大事.
我們就歡喜.
在這篇看到什麼呢?.
我們被擄.
真的想著被人俘虜.
你想想猶太人.
很長的時間.

$^{761}$第二年的時候已經被人俘虜了.
俘虜了這麼久的時間.
俘虜不是好玩的.
你不是在自己的國家.
沒有享受.
做奴隸.
你的地位永遠都是被人低一級.
低很多級.
奴隸.
是洗碗的.
在這樣的光景下.
看到他們.
現在回來了.
所以第一節裡面看到他們.
好像做夢.
我們有機會回到自己的家園嗎?.
帶著希望.
可以回去了.
我們被人俘虜的一群.
現在皇帝肯放我們.
讓我們回去.
還.
你見到.
充滿著希望.
他們回到裡面.
看到他們很荒涼.
因為很久沒有運作了.
你見到.
他們的希望.
是多麼的大.
以前.
是要被人洗碗的.
縱使.
你的房子是很好的.
但你也不舒服,因為那不是你自己的地方.
但當你回去的時候.
原來帶著希望.
這些回歸的人.
他們好像.
蒙恩.

$^{801}$我們得救蒙恩.
不知道你信了主之後.
得救蒙恩的時候.
你開心嗎?.
你是一個蒙恩得救的人.
你是屬於主耶穌基督的.
你很開心.
他們很開心.
開心的警方.
他們有很多的困難.
但今天竟然可以自由了.
帶來希望了.
不會再被人抓回去了.
回到自己的地方.
回到自己的國家.
自己的家園.
神為他們行了大事.
他們好像做夢一樣.
你看看.
以色列和巴勒斯坦打仗.
放了人質.
是不是很開心.
有希望了,不用死了.
有希望了.
你可以感受到.
那種希望.
甚至你看耶穌基督死而復活.
當門徒見到耶穌的時候.
多麼開心.
又有希望了.
原來那夫子沒有死.
帶來盼望了.
因為他是死而復活的主.
所以你看到那種開心的心情.
你看到他們.
當耶穌又離開.
叫他們不要憂愁.
但今天.
我們有了耶穌基督的時候.
祂叫我們做事.

$^{841}$就算修國也是神的事工.
我們要憑信心.
我們不要看到風.
像小朋友說.
下雨了,不上學了.
看到颱風了,不工作了.
但原來做神的事工.
不是看風的,不是看天氣的.
什麼時候都要做.
好像醫院那些.
多緊急都要做.
都要救人.
所以原來.
是不能休息的.
雖然天氣不好.
我們仍然要守望.
你可以祈禱.
所以提摩太后說了什麼?.
保羅教我們.
是.
無要全道.
無論得時不得時.
所以無論是.
好的氣氛.
還是不好的氣氛.
你仍然要去全道.
要告訴別人.
這麼開心的事情在我身上發生.
我也要告訴你.
所以要去殺種.
殺種的功夫.
有時說.
我去殺種.
等待別人修國.
可能我等不到.
可能我死了之後.
那時候信主.
我看不到他.
你就是殺了種子.
在他心裡慢慢萌芽生長.

$^{881}$有一個人問他.
他就馬上決志.
有時很妒忌.
我經常跟他說.
他也不肯相信.
那個阿B一來.
他跟他說兩句.
竟然肯舉手決志.
我以前叫他這麼久.
經常問他.
他都說等一會.
他就信主.
很妒忌.
不要緊的.
神知道你辛苦.
神會紀念你的勞苦.
忍耐.
你見到他.
信主你很開心.
你替他開心.
他終於和舊恩有關.
願意信主.
蒙福的人.
所以我們不要眼淺.
看到別人做的不開心.
有時可能.
你是殺種的.
都是收割的.
你很開心.
所以我們看不見不要緊.
還有收割者的喜悅.
是怎樣呢.
是帶來歡樂的.
流淚殺種的.
必歡呼收割.
很開心的.
你帶了信主很開心.
種植無論多麼努力.
更運.
就沒有了.

$^{921}$你很失望.
但我們的努力好像是.
前功盡棄.
很捨棄.
天災人禍.
但沒辦法控制.
疫情一來.
所有事都停頓.
全世界都受影響.
全世界的人都病了.
這些你沒辦法避免.
忽然失業.
也沒辦法避免.
但我們做了一整天的工夫.
我們看到.
雖然是疫情.
雖然有些人面對打仗.
或者是面對生意失敗.
或者面對失業.
但我們仍然要憑信心去種.
去殺種.
去收割.
因為叫他成長的.
是誰呢?是天父.
就算你種得多漂亮.
他都會死的.
是吧.
我不會種東西.
我種東西都沒收成.
但我們知道原來那樣東西.
叫他成長的是神.
所以我們看別人做得那麼好.
其實不是.
當然他都很努力.
但不是因為人的關係.
他會成長得那麼好.
是天父叫他成長的.
叫他成長得那麼好.
所以我們可以做的就是.
盡心去做好你的那份.

$^{961}$神讓你殺種.
你就去殺種.
讓你去收割.
你就要好好栽培他.
因為有些人不是那麼快就信主.
可能你教完他主一學.
教完他所有的東西.
上了洗禮班.
接著他又想.
過了一年.
你要繼續打電話.
去關心他.
你不要放棄他.
你要關心他.
直至他有一天都像你一樣.
可以教人.
你才可以鬆一口氣.
他沒教人可能隨時都會離開教會.
所以你要繼續守望.
繼續為他祈禱.
繼續關心他.
這就是我們今天.
收割者為什麼會喜悅.
因為前面流了很多眼淚.
流淚殺種.
必歡呼收割.
你會有收成.
你收成的時候.
神將責任交給你的時候.
你可以很開心地將成果獻給神.
然後你就會很滿足.
因為你所做的事.
神會叫你福杯滿人.
你就會得到滿足的喜樂.
所以聖經後面有說到.
好在南地的河水伏流.
滋潤貧瘠的土地.
你要知道.
本來南地是很乾旱的.
可以等到河水伏流.

$^{1001}$河水伏流.
真是興旺.
其實也有另一個意思.
當時回來的人不是很多.
人家肯讓你回來.
但是回來的人不是很多.
但是你要興建.
要很多人.
河水伏流.
很多人回來.
會繁榮.
城市會興旺.
人多了.
新聞說很多人放假.
就走上去.
這裡好像很蕭條.
但是我們都要為城市祈禱.
為城市的人.
我們不只是注重物質.
他們的心靈很饑渴.
我們要為他們祈禱.
希望他們都可以有福音的種子.
給他們.
所以盼望在座的.
你願不願意起來.
做一個忠心的神.
做一個忠心的守望者.
忠心的收割者.
為主去收割.
要忍耐的.
要有神給你的信心.
要歡喜快樂的祈禱.
為你的事工祈禱.
你就會看到成果.
你願不願意呢.
我們一起去祈禱.
愛你的主 愛你的神.
祝我們今天這麼太平.
這麼舒服.
這麼好的環境.

$^{1041}$你幫助我們不是坐在這裡.
你幫助我們.
我們要起來.
我們要去聽你的吩咐.
我們要去使萬民做你的門徒.
我們知道要付很大的代價.
但是主要求你幫助我們.
有一個堅定的信心.
我們靠著耶穌基督.
我們就可以有能力.
去跟人分享.
求你使用我們所說的.
讓人都願意.
歸順你.
讓教會可以繼續.
去做傳福音的工作.
帶更多人順著.
因為我們知道主要求我們回來.
回來之後就要審判世人.
很多人就會被失落.
就會去到地獄的地方.
主要我們都不希望.
這麼多人.
能夠離開你.
我們不希望.
你幫助我們.
努力的.
照著你的吩咐.
我們好好去做好你給我們的使命.
我們要去做你忠心的管家.
求你使用我們每一個.
我們這樣祈禱交託.
是奉主耶穌基督的命求.
\newpage



\section{馬太福音 6:9-13}
\label{sec:7y90EbIBR_w}
\textbf{2023.12.10 《主禱文》 馬太福音6\_9-13 (和修版) 講員 : 陳秀賢牧師}
\newline
\newline
連結: \href{https://youtube.com/watch?v=7y90EbIBR_w}{\texttt{https://youtube.com/watch?v=7y90EbIBR\_w}} ~~~~ 語音日期: 2023-12-10
\newline
\newline
\hyperref[sec:fKZdgGDua7A]{\small{< < < PREV SERMON < < <}}
~
\hyperref[sec:index_chronic]{\small{[返順時目]}}
~
\hyperref[sec:index_scriptual]{\small{[返順卷目]}}
~
\hyperref[sec:drs0Z9UDTE0]{\small{> > > NEXT SERMON > > >}}
\newline
\newline
馬太福音 6:9-13
\newline
\begin{longtable}{cl}
\hline
\hline
章節 & 經文 (和合本修訂版)\\
\hline
6:9 & \begin{tabularx}{0.7\textwidth}{X} 「所以，你們要這樣禱告：『我們在天上的父：願人都尊你的名為聖。 \end{tabularx} \\ \\ \relax
6:10 & \begin{tabularx}{0.7\textwidth}{X} 願你的國降臨；願你的旨意行在地上，如同行在天上。 \end{tabularx} \\ \\ \relax
6:11 & \begin{tabularx}{0.7\textwidth}{X} 我們日用的飲食，今日賜給我們。 \end{tabularx} \\ \\ \relax
6:12 & \begin{tabularx}{0.7\textwidth}{X} 免我們的債，如同我們免了人的債。 \end{tabularx} \\ \\ \relax
6:13 & \begin{tabularx}{0.7\textwidth}{X} 不叫我們陷入試探；救我們脫離那惡者。因為國度、權柄、榮耀，全是你的，直到永遠。阿們！』 \end{tabularx} \\ \\
[1ex]
\hline
\hline
\end{longtable}
$^{1}$各位弟兄姊妹早安.
很感恩和感動和大家一起敬拜.
可以在聖餐出前紀念主耶穌基督.
我今天舌頭長了很多痱子.
所以我怕走音.
求主祝福我說話不要走音.
我想問問弟兄姊妹.
誰會背主禱文?.
舉手給我看看.
很多人會背.
我們很多人都會背主禱文.
我們今天很想看看聖經主禱文是怎樣的.
剛才主席已經幫我們讀一次.
我逐句會說.
我們在天上的符.
紅色那裡.
耶穌基督立下一個很好的榜樣.
耶穌基督自己也這樣祈禱.
我們可不可以讀一下馬考弗14:36這段經文呢?.
預備請.
耶穌說:阿爸父啊,在你凡事都能,求你將這杯撤去.
然而不要照我所願的,而是照你所願的.
耶穌基督其實他向天父祈禱的時候.
他說天父,在你凡事都能,求你將這杯撤去.
不要照我的意思,乃是照你的意思.
在這裡大家有沒有印象這段經文在哪裡祈禱?.
在赫失馬利園.
對,赫失馬利園,大家都很熟悉聖經.
在赫失馬利園,耶穌基督面對很大的苦難.
祂只要自己釘十字架的時候.
是很悲嘮.
所以在赫失馬利園,祂說天父爸爸.
如果你將這杯撤去.
你是能夠的,你凡事都能.
不過我都照你的意思.
所以耶穌基督都是向天父祈禱.
各位等於姐妹.
其實當初信耶穌的時候.
向天父祈禱,難不難呢?.
如果說天父.

$^{41}$我記得我初信耶穌的時候.
回教會,回宣道會信律課.
我見到一個年輕人.
他是做社工的.
他經常祈禱天父爸爸.
我真的毛骨悚然.
當時我初信天父爸爸.
我真的覺得毛骨悚然.
這麼親切.
大家知不知道.
有些人是不能叫天父爸爸的.
叫不出來的.
因為有些是單身家庭.
我今早很早來.
我十點二來到這附近.
我走過這附近.
我想起洪福村.
我其試探過一個人.
有一個SEM.
讀寫障礙的小朋友.
和媽媽照顧他.
在香港很多人是單身家庭.
沒有爸爸.
可能爸爸離婚.
所以叫天父爸爸.
是否很艱難呢?.
是有點艱難的.
聖經告訴我們.
我們所領受的不是奴僕的靈.
營救害怕.
而是我們所領受的.
是兒子的靈.
因此我們呼叫阿爸父.
聖靈和我們的靈同症.
我們是上帝的子女.
因為我們信了耶穌基督.
聖靈在我們裡面.
而我們是神的兒子.
上帝是我們的神.
所以我們可以很親切地.

$^{81}$叫阿爸父.
在《迦勒西書》四章六至七節.
因為你們是神的兒子.
上帝就猜他的兒子的靈.
進入我們的心.
呼叫阿爸父.
可見我們不是奴隸.
而是神的兒子.
各位弟兄姊妹.
以前我們未信耶穌的時候.
我們是罪的奴隸.
我們是撒旦魔鬼的奴隸.
以前我們未信耶穌.
我們說謊話.
是不用眨眼的.
我們說一個謊話.
有四千個謊話蓋著他.
都不用眨眼.
當我們信耶穌的時候.
聖靈經常提醒我們.
你剛才說話得罪了人.
你剛才說謊.
你說話又不真實.
聖靈經常提醒我們.
為什麼呢?.
因為我們已經成為神的兒子.
所以神愛我們.
神會教我們.
我們祈禱的時候.
我們可以叫阿爸父.
主童文物「我們在天上的父」.
我們在天上的父.
我們可以叫親愛的天父.
我們這樣祈禱.
但「我們在天上的父」.
主童文教我們非常好.
是「我們」.
紅恩堂是隸屬宣導會.
你們的舞堂是感受堂.
所以感受堂.

$^{121}$是差遣很優秀的師班長.
來教我們音樂.
這是宣導會.
如果你在香港.
有機會在這個職場.
或者在不同的地方.
你會見到宣導會的弟兄姊妹.
你會見到浸信會的弟兄姊妹.
你會見到中華基督教會的弟兄姊妹.
你會見到聖公會的弟兄姊妹.
你會見到不同宗派的弟兄姊妹.
我們懷著什麼心呢?.
可以稱兄道弟.
雖然他們是不同宗派.
我們仍然可以跟他們很友好.
因為他們都是天父爸爸的子女.
剛才我坐在這裡的時候.
有位獅班員.
獅班員是天父的兒子.
獅班員走過來跟我說.
跟我打招呼.
他說:陳牧師.
以前我們在屯門.
我見到你帶著不同宗派的弟兄姊妹.
一起在屯門的街市.
幾百人傳福音.
唱歌,敬拜,讚美.
走過來跟我說這件往事.
其實我們可以跟不同宗派.
在嘉年華會或街市傳福音.
因為我們跟他們沒有競爭.
在24日你們報街音.
我相信不只是宣道會.
紅春堂報街音.
還有其他教會等於報街音.
當你看到他們的時候.
Hello 祝福你.
彼此祝福.
為什麼你可以彼此祝福?.
因為我們在天上的父.

$^{161}$所以各位等於姐妹.
這種合一心.
因為我們所有人.
雖然來自不同宗派.
有不同的讀聖經的方法.
祈禱的方法.
聖餐的方法.
但我們都可以宣告.
我們是天父爸爸.
我們可以有朋友.
第二個主土文.
我們讀讀這個主土文.
紅色字.
願人都尊你的名為聖.
我們很希望.
每一個人都尊神為聖.
在以西結書36章23節說到.
神說.
我在他們眼前.
在你們身上顯為聖的時候.
他們就知道我是耶和華.
上帝自己都會.
表揚祂自己的名字.
神在很多的先知書裡.
說出自己.
是聖潔的.
神是公義的.
神是慈愛的.
神是無瑕疵的.
神是分別為聖的.
所以祂會表述祂自己的屬性.
行不行.
我們做人.
我們都會榮耀我們的主耶穌基督.
在馬太福音5章16節.
都讓祂的百姓.
表揚祂自己.
我們讀讀聖經.
你們的光也要這樣.
照在人前.

$^{201}$叫他們看見你們的好行為.
就將榮耀歸給他們.
天上的父.
我看到.
可能和我年紀差不多.
你們有沒有經歷過.
以前內地經濟不太好.
我們要擔很多油.
很多米.
很多衣服.
過羅湖的深圳河.
就拿去內地.
有沒有這個經驗.
我們這樣做的時候.
真的將經濟.
將我們所有的物質.
送給我們內地的親戚朋友.
所以中國人有句說話.
就是江中要什麼.
祖.
我們要江中要祖.
很開心地將很多物資.
江中要祖.
但我們信了耶穌之後.
我們江中要什麼.
神.
江中要主.
江中要神.
我們要榮耀神.
這兩三年退休.
去了不同的教會.
不同的宗派去港島.
認識了很多大英姐妹.
我.
你知道.
很多大英姐妹認得我.
有一次我在大興花園的港島.
我走進大興堂.
有一個姐妹.
迎面來說.

$^{241}$陳牧師.
我不認得她.
我馬上說.
我是屯門堂的大英姐妹.
她說是嗎.
我馬上想一想.
我住在迎風園.
屯門那邊.
經常在街上買菜.
在超市買菜.
真的要小心一點.
因為在買菜.
你跟別人很沒禮貌的時候.
旁邊的姐妹跟你說.
如果你跟買菜的人.
很沒禮貌的時候.
旁邊的姐妹跟你說.
陳牧師早安.
要煮.
是不是不能說要煮.
大家要聰明一點.
剛才有主席敬拜.
唱得很好.
你們唱歌都穿了斯班圖.
其實你認不出會眾.
如果有個新朋友來了.
紅欣堂今早唇拜.
他認得到你.
如果你去餐廳吃東西.
很粗魯地跟別人說話.
你在江巢?我這麼久都沒來.
新朋友就認得到你.
你是主席.
斯班圖指揮他認得到你.
你認不出他.
這樣就會失手.
所以我們要小心.
我們的行為和表現.
我們要光宗耀祖.
我們要榮耀神.

$^{281}$願人都尊你名為聖.
願每一個人都honor your name.
願每一個人.
藉著詩歌.
藉著我們的行為.
要榮耀主.
要尊主為聖.
繼續讀主禱文.
預備請.
願你的國光臨.
願你的旨意行在地上.
如同行在天上.
我們主禱文.
很熟悉.
但你留意.
其實我們背主禱文.
我們很熟悉.
但有時會口渴.
什麼叫口渴.
我們會這樣說.
奇怪了.
但心是否這樣.
願你的國光臨.
願你的旨意行在地上.
如同行在天上.
那種一統化.
願神你的旨意.
行在地上.
就像在天上一樣.
那種一致性.
神的旨意是甚麼.
神的旨意.
神的價值觀.
願你的想法.
願你的旨意行在地上.
如同行在天上.
這是我們背主禱文的心願.
但往往.
現實世界.
真實又未必這樣.

$^{321}$其實我們會一國兩制.
我們實踐.
一國兩制.
神你的國就是你的國.
我的國就是我的國.
即是一國兩制.
神你的價值觀.
願你的旨意.
有時我們是不能實行的.
我記得今年.
三月份.
我要去緬甸一次.
因為幾年內戰.
不能乘飛機.
去不了.
又COVID.
七月我要去泰國一個月.
去宣教支援工場和宣教士.
我退休之前.
我有一個計劃.
因為退休後沒有工資.
只有出沒入.
所以計劃一年去一次宣教.
一年去一次.
誰知今年.
要去兩次.
計算下.
還是去一次緬甸.
去一次緬甸.
因為機票很貴.
現在一萬五千元去一次緬甸.
以前七千元.
包括旅費.
現在是雙倍.
所以我計算.
去一次緬甸或泰國.
我就想.
不去緬甸去泰國.
因為泰國去一個月.
洗費更大.

$^{361}$但我心裡知道.
抗疫中止.
因為神都感動我去緬甸.
見到神學院的同工.
神學生.
因為他們內戰.
神都感動我.
去緬甸見弟兄姊妹.
今年三月停戰.
所以可以去.
但我抗止下.
就很混亂.
我現在決定.
去緬甸.
去泰國.
兩邊都去.
當然輸了.
香港警察會.
免費讓我去泰國.
緬甸我自己付錢.
原來做牧師都會抗止.
我經常抗止.
但我經常要學習.
要遵行神的旨意.
各位弟兄姊妹.
願你的旨意.
行在地上如同行在天上.
《瑪撻科》七章有一節.
不是每一個人.
稱呼為主啊.
就可以進天國.
唯有什麼人呢?.
就是遵行我父旨意的人.
才能夠進天國.
假如你的子女跟你說.
媽媽.
我很孝順你.
我很有愛心.
你跟兒子說.
是嗎?.

$^{401}$你很有愛心嗎?.
你一年都沒有和我去飲茶.
只有「港」字.
其實.
媽媽我很有愛心.
其實沒有實踐.
如果我們面對神的旨意.
要說主啊主啊.
主啊主啊.
其實我們沒有遵行神的旨意.
這樣是不行的.
所以耶穌基督.
教我們.
先求神的國.
一切東西都要交給你.
學習先求神的國.
學習先求神的義.
神一切東西.
都會交給你們.
現在香港教會.
面對一個很大的困局.
香港經濟.
有一點倒下.
香港很多教會.
經濟都下滑.
急遽.
首當其衝.
有些教會會戒掉.
宣教.
戒宣教.
宣教的支出.
我知道宣教會有難處.
但我們香港教會.
要一起反省.
我們支持一個宣教士.
他可能去某個國家.
他在牧養.
50萬他們區域的人.
他們的.
微德斯文.

$^{441}$可能他擁有50萬人.
我們只是在香港.
我們香港基督徒.
我們仍然節省一點.
節省一點錢.
節儉一點.
減少生活.
讓這些錢.
十一奉獻.
十二奉獻.
仍然奉獻給差權.
讓神的家有餘.
不會缺乏.
這才是先求神的國.
和神的義.
一切東西都會交給我們.
主禱文是耶穌基督教我們祈禱.
我們一起讀.
預備請.
我們日用日的飲食.
今日賜給我們.
我們日用日的飲食.
今日賜給我們.
可能早餐午餐晚餐.
我們都會用這節經文.
祈禱.
豐富的食物.
我們日用日的飲食.
今日賜給我們.
我們日用日的飲食.
明天賜給我們.
我們日用日的飲食.
明天賜給我們.
今日賜給我們.
我們想起.
以色列人在抗疫的時候.
他們吃什麼大的.
馬拉大.
馬拉是每天供應給他們.
但不能儲存到日後.

$^{481}$如果儲存到日後會發生什麼事.
會發霉會變.
所以馬拉是每天賜給以色列人.
我們日用日的飲食.
今日賜給我們.
我們日用日的飲食.
明天賜給我們.
我們日用日的飲食.
明天賜給我們.
我們日用日的飲食.
明天賜給我們.
我們日用日的飲食.
明天賜給我們.
我們日用日的飲食.
明天賜給我們.
我們日用日的飲食.
明天賜給我們.
我們日用日的飲食.
明天賜給我們.
我們日用日的飲食.
明天賜給我們.
我們日用日的飲食.
明天賜給我們.
我們日用日的飲食.
明天賜給我們.
我們日用日的飲食.
明天賜給我們.
我們日用日的飲食.
明天賜給我們.
我們日用日的飲食.
明天賜給我們.
我們日用日的飲食.
明天賜給我們.
我們日用日的飲食.
明天賜給我們.
我們日用日的飲食.
明天賜給我們.
我們日用日的飲食.
明天賜給我們.
我們日用日的飲食.

$^{521}$明天賜給我們.
我們日用日的飲食.
明天賜給我們.
我們日用日的飲食.
明天賜給我們.
我們日用日的飲食.
明天賜給我們.
我們日用日的飲食.
明天賜給我們.
我們日用日的飲食.
明天賜給我們.
我們日用日的飲食.
明天賜給我們.
我們日用日的飲食.
明天賜給我們.
我們日用日的飲食.
明天賜給我們.
我們日用日的飲食.
明天賜給我們.
我們日用日的飲食.
明天賜給我們.
我們日用日的飲食.
明天賜給我們.
我們日用日的飲食.
明天賜給我們.
我們日用日的飲食.
明天賜給我們.
我們日用日的飲食.
明天賜給我們.
我們日用日的飲食.
明天賜給我們.
我們日用日的飲食.
明天賜給我們.
我們日用日的飲食.
明天賜給我們.
我們日用日的飲食.
明天賜給我們.
我們日用日的飲食.
明天賜給我們.
我們日用日的飲食.

$^{561}$明天賜給我們.
我們日用日的飲食.
明天賜給我們.
我們日用日的飲食.
明天賜給我們.
我們日用日的飲食.
明天賜給我們.
我們日用日的飲食.
明天賜給我們.
在主禱文繼續說.
我們一起讀.
預備請.
免我們的債.
如同我們免了別人的債.
免我們的債.
如同我們免了別人的債.
各位靚姐妹.
下一次你背主禱文小心一點.
因為每一次背主禱文.
你是提醒上帝.
你首先看看.
如同我們免了別人的債.
邏輯是說天父爸爸.
你是不用免我們的債.
如果我們不免別人的債.
這是邏輯.
陳牧師你有沒有解錯正經.
有沒有讀錯神學.
其實神在十八七言.
祂經常赦免我們.
耶穌基督愛我們.
祂長闊高深.
祂一定會赦免我們的罪.
等姐妹們弄清楚.
如果我們不赦免別人的債.
神也不會赦免我們.
你記不記得.
在《馬太空間十八章》.
有一個很熟悉的比喻.
耶穌基督見了彼得.

$^{601}$彼得說.
夫子.
我搖搖搖搖了人多少次.
我有些得罪我的人.
有些得罪我的人.
搖搖搖搖了七次.
夠不夠.
彼得已經覺得很醒目.
七次.
但耶穌說不是七次.
是多少次.
是七十個七次.
完完全全的寬恕.
耶穌說完這句話之後.
他說了一個比喻.
耶穌說有一個很有錢的人.
有一次他借錢給一個A君.
A君欠他債.
這個有錢人說.
你欠我債.
你賣了你的老婆.
你賣了你的子女.
你值得還.
這個A君說我實在是沒有錢.
給我一點時間.
我會還錢給你的.
香港.
如果欠人債.
很嚴重.
會被人淋紅油.
貼街招.
其實是很嚴重的事.
但當他去.
呼求主人的時候.
這個主人.
說不用他還.
但這個人很有趣.
他看到B君.
B君又欠他一筆錢.
大概一百元.

$^{641}$很小的數字.
他就說.
你欠我錢.
你快點還.
B君說無法償還.
於是他就被抓去坐牢.
這件事實在是.
很多人看不過眼.
在看不過眼之下.
你知不知道.
你原諒那個A君.
A君竟然打B君.
被抓去坐牢.
所以這個主人.
很不留情.
就說了這句話.
第34節.
我們一起讀.
主人就大勞.
把他交給私營的.
直到他還清了所欠的債.
你們各人.
若不從心裡描述你的弟兄.
我天父也要這樣.
待你們.
各位兄弟姐妹原來是這樣的.
我們是蒙赦免.
我們是蒙恩.
我們是恩信稱義.
但當我們信耶穌之後.
你讀過雅各書.
雅各書就是你有行為.
你有信心就有行為.
你有愛心.
你就會孝順父母.
有信心就有行為.
如果你是經歷過耶穌基督的救恩.
這個恩典的時候.
你是可以用神這個.
莫大愛去饒恕別人.

$^{681}$主人就說如果你個人心裡.
你不饒恕你的弟兄.
我也不會饒恕你.
所以在.
主禱文的.
14至15節.
14至15節是這樣說的.
你們若.
饒恕你的過犯.
你們的天父也必.
饒恕你們.
你們若不饒恕人.
你們的天父也必.
不饒恕你們的過犯.
各位靈姐妹.
你們心裡有沒有.
仇家.
那些仇家.
是否誤會了你.
或者你誤會了仇家.
你有沒有被冒犯過呢.
你有沒有被得罪過呢.
今天.
我很想和你們講主禱文.
我們選擇.
去饒恕.
在情感上你受傷害.
別人冒犯你.
在情感上.
你想想你自己也不開心.
這是情感上.
但在理智上和真理上.
我們今天要選擇.
饒恕.
因為當我們不饒恕別人.
我們會影響我們的情緒.
影響我們睡眠.
影響你.
事奉神的功效.
和果效.

$^{721}$你得罪了別人.
你如何高歌讚美主呢.
唱詩班呢.
所以.
會影響我們服侍的果效.
影響我們和天父.
父親的關係.
所以我們快快求主幫我們.
選擇.
我們饒恕心裡的.
曾經冒犯過我們.
得罪過我們.
不友好我們的弟兄姊妹.
或朋友.
我們今天就選擇.
因為我們.
如何對待別人.
影響上帝如何對待我們.
所以你下一次.
背主禱文的時候.
主禱文很有意思.
提醒我們要饒恕別人.
因為當.
我們饒恕別人的時候.
天父才會饒恕我們.
不叫我們遇見斯坦.
救我們脫離凶惡.
其實.
我有一個朋友他做地盤工作.
他每天早上都要.
祈禱不叫我們遇見斯坦.
救我們脫離凶惡.
救我們脫離斯坦.
救我們脫離凶惡.
我們作基督徒.
每天一個人.
早上起床已經.
面對很多很多.
試探.
星期日起床.

$^{761}$你已經有一個掙扎.
回來實體還是直播.
你一早起床.
一開機就看youtube的短片.
還是靈修.
這是試探.
很多的試探.
叫我們離開上帝.
親近撒旦.
所以我們在《轉圖文》裡.
我們呼求神天父爸爸.
求你叫我們.
遠離這些試探.
救我們脫離這些凶惡.
我年輕的時候.
我的牧師就跟我說.
教我們如何逃避.
試探.
他說假如你知道附近有一個.
報紙攤.
有很多色情的報紙和雜誌.
你就繞路走過.
繞路走.
這就叫逃避試探.
各位靈異姐妹.
有時活在香港逃避是很好.
但另一方面我們有時避不過.
很多事.
很多事是湧過來.
我記得.
幾個月前我聽過.
神學院有一個講座.
是講IT.
現在的人工智能.
對我們的試探是很大的.
我們很明白.
各位靈異姐妹.
每個人都有手機.
智能手機.
非常快.

$^{801}$5G.
你發現一件事.
這個手機給我們做了什麼試探.
我們用很多時間看短片.
一個短片可能十分鐘.
十分鐘十分鐘.
十分鐘二十分鐘.
我們用了很多時間.
我們沒有時間靈修.
沒有時間讀聖經.
退休了.
我們都要過一個很紀律的生活.
除了服侍.
還有讀聖經.
現在12月.
我將會把新舊約都讀一次.
各位靈異姐妹.
那個智能手機.
是令我們很大.
很大的試探.
我們小心一點.
今年3月我去了緬甸.
在緬甸這十幾天.
發現緬甸的詐騙.
很嚴重.
很多人投入詐騙的行業.
很多人用電腦.
騙很多人.
很多年青人不相信耶穌.
去了詐騙公司賺大錢.
我思想.
他們會國語.
和電腦.
這個世界告訴我.
他騙什麼人.
就是騙14億同胞.
和東南亞人.
因為是國語.
今天我們所面對的試探.
是很靈禮.

$^{841}$最近我看了一套電影.
叫《孤豬一擲》.
這套電影.
是講詐騙公司的手段.
其實.
為什麼很多人受騙.
東張西望.
經常把例子講出來.
但也有很多人受騙.
為什麼.
因為貪心的試探.
貪心就是.
投資這些東西.
要少本.
就賺很多錢.
貪心就是.
我還未失業.
現在有份工作.
三萬元不用上班.
在家中工作就能賺到三萬元.
你又有應機.
貪心.
不甘心.
你輸了錢.
就被人騙.
我們基督徒是沒有問題的.
各位近年居民.
我們是沒有免死金牌.
我們真的沒有免死金牌.
我認識一個弟兄.
他平日不奉獻.
經常說好像沒有錢.
裝窮.
不肯十一奉獻.
但你知道嗎.
他被一個朋友.
騙了.
投資.
他整副家產.
原來我們真的沒有免死金牌.

$^{881}$我有一個朋友.
他讀博士.
很厲害.
但他很寂寞.
他上網站.
他認識了一個朋友半年.
後來他警覺.
這個人應該詐騙他.
各位姊妹.
我們要小心.
現在的WhatsApp會突然彈出來.
WhatsApp會問你.
吃飯了嗎?.
你會怎樣?.
你會怎樣回答?.
你會問他吃飯了嗎?.
不會的 會立即舉報他.
這是詐騙.
有一次WhatsApp.
問我.
很久沒有跟你行過山.
我心想.
我不認識你.
又舉報他.
早前天氣下降.
在香港冷風.
冷風冷.
你穿回衣服.
我不認識你.
又舉報他.
各位弟妹.
有人真的被騙了.
一副退休金.
就是那種孤單.
這些事態.
隨時出現.
在WhatsApp.
你沒有定力.
好心回覆.
你就被騙了.

$^{921}$各位弟妹.
我們面對的.
很多試探.
你有經驗嗎?.
如果你血壓高.
你上YouTube.
那兩三天會有.
高血壓的資訊.
我這樣想.
就是因為IT.
大量收我們的數據.
知道你的興趣.
它會彈一些資料給你.
糖尿病.
它會彈一些糖尿病資料給你.
但這個科技很危險.
如果你不小心彈了色情畫像.
這兩天.
會鋪天蓋地彈給你.
弟兄也好.
姐妹也不要嘴笑.
我們都會遇到困難.
所以如果你真的這樣.
你只有呼求神.
我不知道你的組長.
會否用你的手提.
我不知道你的傳道會否用你的手提.
假如你不小心.
這樣入手機.
真的看了色情畫像.
你就說求主.
救我脫離兇惡.
在心坑.
在這個沼澤快點拉我上來.
因為在色情上.
眼目的情慾.
今生的驕傲.
眼目的情慾會破壞我們和神的關係.
破壞我們和配偶的關係.
破壞我們和妻子的關係.

$^{961}$破壞和丈夫的關係.
破壞我們的生命.
令我們心敗名裂.
所以試探是很凌厲.
求主幫我們.
救我們脫離.
這些坑.
因為國道權並榮耀.
全是你的直到永遠.
阿們.
主禱文教我們.
如何做一個禱告的人.
禱告人很重視.
和天父爸爸的關係.
所以我們在天上的父.
禱告人不是口哼哼地祈禱.
而是願你的旨意.
行在地上.
如同行在天上.
禱告人是深深相信.
天父是公認我們.
禱告人是願意學習.
被寬恕.
而且去寬恕人.
而禱告人.
是要靠主勝過.
一切一切的試探.
我們一起祈禱.
弟兄姊妹.
剛才我提到試探.
你現在不小心陷入試探.
你在罪裡.
你現在就在聖殿裡.
你呼求神.
求神救你脫離這個試探.
剛才我剛才說到要饒恕人.
你心裡想起仇家.
你很憤怒他.
我邀請你現在選擇.
寬恕他.

$^{1001}$原諒他.
假若你常常都很憂慮.
耶穌現在跟你說.
不要憂慮.
不要受煩.
不要憂慮.
假若你今天聽到.
你立定志向.
去遵行神的旨意.
讓神的國度.
真的完全降臨在你心靈裡.
你遵行神的旨意.
我們一起.
祈禱.
如果我問你的問題.
你很想求主幫你.
你把你的右手.
放在.
你的心口前.
作一個祈禱.
你需要神.
主耶穌.
你在主禱文裡教我們.
我們紅恩堂弟兄姊妹.
心裡呼求你.
口裡呼求你.
手所擺的姿勢.
你一一都看到.
我們在你的桌前.
很認真的禱告.
求你幫助我們.
按我們不同每一個人的需要.
你幫助我們.
憐憫我們.
主我們願意.
願你的國降臨.
願你的旨意.
行在我們心靈當中.
如同行在天上.
多謝耶穌.

$^{1041}$聽我們祈禱.
\newpage



\section{創世記 32:22-32}
\label{sec:drs0Z9UDTE0}
\textbf{2023.12.17 《真正的得勝》 創世記32\_22-32 (和修版) 楊小玲傳道}
\newline
\newline
連結: \href{https://youtube.com/watch?v=drs0Z9UDTE0}{\texttt{https://youtube.com/watch?v=drs0Z9UDTE0}} ~~~~ 語音日期: 2023-12-20
\newline
\newline
\hyperref[sec:7y90EbIBR_w]{\small{< < < PREV SERMON < < <}}
~
\hyperref[sec:index_chronic]{\small{[返順時目]}}
~
\hyperref[sec:index_scriptual]{\small{[返順卷目]}}
~
\hyperref[sec:ZLhqPiukKhM]{\small{> > > NEXT SERMON > > >}}
\newline
\newline
創世記 32:22-32
\newline
\begin{longtable}{cl}
\hline
\hline
章節 & 經文 (和合本修訂版)\\
\hline
32:22 & \begin{tabularx}{0.7\textwidth}{X} 他夜間起來，帶著兩個妻子，兩個婢女和十一個孩子，過了雅博渡口。 \end{tabularx} \\ \\ \relax
32:23 & \begin{tabularx}{0.7\textwidth}{X} 他帶著他們，送他們過河，他所有的一切也都過去， \end{tabularx} \\ \\ \relax
32:24 & \begin{tabularx}{0.7\textwidth}{X} 只剩下雅各一人。有一個人來和他摔跤，直到黎明。 \end{tabularx} \\ \\ \relax
32:25 & \begin{tabularx}{0.7\textwidth}{X} 那人見自己勝不過他，就摸了他的大腿窩一下。雅各的大腿窩就在和那人摔跤的時候扭了。 \end{tabularx} \\ \\ \relax
32:26 & \begin{tabularx}{0.7\textwidth}{X} 那人說：「天快亮了，讓我走吧！」雅各說：「你不給我祝福，我就不讓你走。」 \end{tabularx} \\ \\ \relax
32:27 & \begin{tabularx}{0.7\textwidth}{X} 那人說：「你叫甚麼名字？」他說：「雅各。」 \end{tabularx} \\ \\ \relax
32:28 & \begin{tabularx}{0.7\textwidth}{X} 那人說：「你的名字不要再叫雅各，要叫以色列，因為你與神和人較力，都得勝了。」 \end{tabularx} \\ \\ \relax
32:29 & \begin{tabularx}{0.7\textwidth}{X} 雅各問他說：「請告訴我你的名字。」那人說：「何必問我的名字呢？」於是他在那裡為雅各祝福。 \end{tabularx} \\ \\ \relax
32:30 & \begin{tabularx}{0.7\textwidth}{X} 雅各就給那地方起名叫毗努伊勒 ，說：「我面對面見了神，我的性命仍得保全。」 \end{tabularx} \\ \\ \relax
32:31 & \begin{tabularx}{0.7\textwidth}{X} 太陽剛出來的時候，雅各經過毗努伊勒，他的大腿就瘸了。 \end{tabularx} \\ \\ \relax
32:32 & \begin{tabularx}{0.7\textwidth}{X} 因此，以色列人不吃大腿窩的筋，直到今日，因為那人摸了雅各大腿窩的筋。 \end{tabularx} \\ \\
[1ex]
\hline
\hline
\end{longtable}
$^{1}$各位姐妹平安.
不好意思.
大家見我有少少失常.
平時上到台應該準備好.
但今日有點手震.
有點無所適躁.
但相信大家會明白.
因為正如.
今日要分享的主題.
是真正的得勝.
相信我們.
會在耶穌的眷顧下.
每一個人都過一個得勝的生命.
今日.
自從我服侍以來.
一個最特別的.
正道.
第一,今日我原本.
在宣道會恩橋堂.
正道.
但幾日前因為.
Gary牧師的情況.
轉差,所以恩橋堂的.
同工很好,就說.
你不如留在家陪牧師.
照顧牧師.
當時我還堅持.
但對方又堅持作出.
一些調動.
以致我今日.
不用去那裡.
第二,就是.
平時我都會想.
如果家中有人過身.
有甚麼重大事,那天又要講道.
那我去不去呢?.
我的決定是.
去的,但沒有想過.
今日是成真.
第三,雖然.

$^{41}$我比Gary早.
出來服侍.
但我從來.
都沒有他的福氣.
可以在母會.
或洪恩堂正道.
我記得Gary.
在三月病發前.
有一天,他下班.
他神色.
很凝重地對我說.
細節我忘記了.
但要不這樣.
他說我決定了.
我決定了,我要邀請你.
來洪恩堂講道.
因為我相信你的道.
對回眾有益.
但沒有想到.
他做了這個決定後.
很快就病了.
在病中.
他說.
不會再有可能.
讓你來洪恩堂講道.
當時我也是這樣想.
唐董.
一群弟兄姊妹.
唐偉.
能夠讓我們夫婦.
今天可以達成這個心願.
這個心願就是.
願意Gary.
過去的服侍.
我未來的服侍.
我們同樣.
成為大家的.
祝福.
更願主是德榮耀.
我們一起禱告.

$^{81}$親愛主.
我們讚美你.
因為你掌管生命.
所以每一個人.
是生是死.
都是在你的美意當中.
只要在你的美意當中.
此生無憾.
亦無懼.
前面是否波濤洶湧.
掌管生命.
坐著為王.
這一刻我們知道.
你就是在這裡.
我們同心仰望你.
我們要敬拜的是你.
我們要聽的訓誨.
是從你而來.
願講的聽的.
都能夠得著.
神你的月亮.
願講的聽的.
都懷著謙卑的心.
以致我們一生.
都能夠.
倚靠神.
過得勝的生活.
願你.
坐著為王.
願你得榮耀.
我們祈禱交托.
豈奉靠我主耶穌基督.
得勝的名字祈求.
我今天要分享的題目.
是真正的得勝.
得勝就是得勝.
如果我們要真正.
把兩個字放進去.
可能就是暗示.
這個得勝.

$^{121}$是有待商榷的.
不知道大家有沒有想過.
在自己的生命裡.
怎樣才是真正的得勝呢?.
今天我們.
分享的經文是.
《創世記》三十二章.
二十二至三十二節.
但其實我們不只是.
了解這幾節的經文.
而是我們需要知道經文的背景.
我簡單講一下背景.
耳塑和雅各是親兄弟.
大家都知道耳塑是長子.
可以得到父親.
很特別很寶貴的祝福.
但雅各.
就串通媽媽.
去騙祝福.
令到耳塑很生氣.
所以耳塑就決定.
等爸爸過世之後.
我就要將雅各殺了.
於是雅各就逃跑.
逃出了大陸.
雅各就逃跑.
逃到舅父的拉班.
在舅父的家裡.
他很喜歡他的表妹拉潔.
於是他為了娶拉潔.
就為舅父服侍了二十年.
他很本事.
他在這二十年裡.
贏得了老婆.
子女,瀑婢.
和生畜.
財產豐盛了很多.
又一無所有.
到現在成為一個大財主.
其實應該叫得勝.

$^{161}$但很可惜.
因為他和拉班有利益衝突.
於是他又要帶著.
他的妻兒,財產.
再一次.
要踏上逃亡的路.
拉班當然不肯放過他.
用了七天.
終於追到他.
但很感恩.
上帝阻止拉班.
加害雅各.
經文去到.
三十二章的時候.
正正是說到.
拉班放過雅各.
於是大家和平.
道別.
然後雅各又帶著家人.
和財產繼續.
上路.
我們回望雅各的過去.
他似乎一直.
都活在一個真境裡.
而且他的堅持力.
是非常強的.
他是一次又一次地贏.
現在他又.
可以.
拉班放過他.
他可以保留他的家人.
他的財產.
他又再次得勝.
但他的心情.
在這一刻沒有得勝的喜悅.
除了回到老家.
他再沒有其他的選擇.
但前路是生是死.
我相信.
他一點把握都沒有.

$^{201}$因為他知道他當年.
欺騙了耳掃.
拿了應該.
屬於耳掃的祝福.
他不知道耳掃.
會不會放過自己.
其實過去.
雅各曾經用一些計謀.
是一次又一次地得勝.
所以在這一刻.
他的老外又出來.
他又要施計.
首先他用什麼計呢?.
他打發人去耳東地.
去探查.
究竟耳掃的態度是怎樣.
等著等著終於探子.
回來了,就報告說.
原來耳掃帶了四百人.
迎著雅各而來.
雅各很擔憂.
嚇死.
耳掃是不是要來殺自己呢?.
於是他又.
繼續用計.
他就將那些瀑被,牲畜.
分成兩隊.
為什麼?萬一耳掃.
擊殺其中一隊.
起碼另一隊都還有機會.
逃跑.
保留那些資產.
同時他很會說話.
他很聰明.
在三十二章九節.
他向上帝禱告.
他怎樣禱告呢?他說:耶和華.
我祖亞伯拉罕的神.
我父親以撒的神.
你曾對我說.

$^{241}$回你本地本族去.
我要厚待你.
雅各很聰明.
他刻意提醒上帝.
上帝啊!你曾經賜下應許的.
當然.
他在十一節也很坦白.
跟上帝說.
他很害怕,他求上帝.
求你救我脫離我哥哥.
耳掃的手.
因為我怕他來殺我.
連妻子,大兒女.
一同殺了.
他好像很勁虔,可惜他求神.
不過他祈禱完之後.
他還是不放心.
他還是要自己用計.
我們會不會是這樣的?.
這次祈禱完,那次用自己的方法.
在三十二章十四至十六節記載.
雅各很匆忙.
預備了要送給耳掃的生畜.
每樣分一堆.
每一堆.
都有一個距離.
為什麼?一群一群.
出現在耳掃面前.
第一群耳掃看到可能更生氣.
第二群可能更生氣.
第三群又更生氣.
可能最後就全部被拒絕.
因為看到那麼多禮物.
萬一發生不利的情況.
還是生氣的話,那怎麼辦?.
一群一群去,雅各看不對路.
就可以逃亡.
按道理來說.
雅各應該勝算在握.
因為他過去沒有輸過.

$^{281}$他也是贏的.
首先他騙取了長子的祝福.
然後他在寇府家贏得豐厚的家財.
最後寇府又放過他.
不追討他.
但是在二十二至二十三節看到.
雅各似乎一點得勝把握都沒有.
他連夜帶著兩個妻子.
兩個婢女和十一個兒子.
來到雅各渡口.
為什麼呢?.
原來雅各想要將他們打發才過河.
不過自己就獨自留下.
有些人就說雅各愛惜自己的生命過於一切.
所以兒女先去.
如果你們讓二嫂去對付的話.
其實我還在這邊.
我還可以走.
不知道是不是他這樣想.
真的不知道.
要殺就殺他們吧.
我留到最後.
要走的時候方便一點.
但肯定的是.
一個人面對生死未卜的威脅.
我相信絕對不好受.
就在這個孤單迷茫的時候.
在第二十四節就說.
「只剩下雅各一人.
有一個人來與他率交直到黎明」.
咦?.
有一個人來與他率交.
這個人是誰呢?.
我們其實看聖經要通的.
要看其他的.
我在《河西亞書》十二章找到線索.
那裡說.
「他在福中抓住哥哥的腳筋.
壯年的時候與神交力.
與天使交力並且得勝」.

$^{321}$這裡是在說雅各.
他說雅各壯年的時候與神交力.
與天使交力也得勝.
這句經文讓我們想到兩個問題.
第一個問題.
究竟雅各與神交力還是與天使交力呢?.
為什麼兩個都說呢?.
沒有換人.
經文沒有說打打打換別人.
有解經家認為.
人是不能與神面對面的.
所以應該是一個代表神的使者.
但即使不是上帝自己出現.
使者的出現是甚麼?.
都是上帝猜派的.
使者是代表上帝的.
所以其實都是與上帝面對面.
所以我覺得我們不用理會.
究竟這是上帝還是天使.
總之我們知道.
上帝是要藉著摔跤.
讓雅各去經歷.
第二個問題就是.
雅各在這場交力裡是否真的得勝呢?.
這裡我要賣一個關子.
我講道很喜歡這樣.
你們要專心聽.
大家可以慢慢判決.
我們回到二十四節.
當雅各孤單一個人的時候.
上帝使者出現.
其實這一幕對雅各來說是似曾相識.
因為在二十年前.
雅各為了逃避以掃帚報復.
他孤單一個人去到伯特利.
不知道大家是否熟悉這個經文.
在伯特利上帝出現.
是透過夷夢.
雅各在夢裡與上帝相遇.
現在雅各又逃亡.

$^{361}$同樣地孤單一個人.
上次一個人都沒有.
這次有老婆子女.
但是都不在一起.
上帝又再一次讓他去經歷.
經歷上帝的同在.
弟兄姊妹.
我相信我們大部分人.
都會經歷過孤單和困惑.
不會說洛天派到一個地步.
從來都不覺得難過.
在孤單裡的時候.
你有沒有經歷過上帝與你一起呢?.
如果有的話.
我想問.
這些經歷對你的生命有什麼作用呢?.
有時我覺得很可悲的.
就是基督徒的屬靈經歷.
竟然成為他日後自誇的網絡.
我那次那次怎樣經歷上帝.
上帝怎樣對我說話.
他說來說去都是那個見證.
他很得戚.
他沒有想過.
他已經是佔奪了上帝的榮耀.
還有你去經歷上帝.
經歷上帝的恩典.
是不是讓我們覺得.
都很幸運.
上帝放過了我.
上帝都看著我.
以致日後我們更加放肆.
其實每一次經歷上帝.
就好像儲糧一樣.
面對饑荒的時候.
我們就不怕.
因為我們已經儲了不良.
我記得那時候.
Gary牧師在服侍的時候.
他跟很多的NGO都有合作.

$^{401}$幫紅印堂申請了很多很多.
柴米油鹽.
禮物,口罩堆滿了.
那時候多開心.
在疫症之下.
又可以照顧周圍的街坊.
又可以派給弟兄姊妹.
多溫暖那種感覺.
所以其實經歷上帝.
是要讓我們好像儲糧一樣.
有那種滿足.
那種安全.
意思是如果下次再有同樣的孤單.
難受的時候.
其實我們不怕.
因為我們的上帝不是喜怒無常的神.
他過去怎樣看顧我們.
他今天仍然要看顧我們.
他將來仍然要看顧我們.
現在雅各就是這樣.
他很孤單.
上帝又來關照他.
這一次用了一條很新穎的方法.
竟然跟他摔跤.
雅各依然貫徹他的宗旨.
他很有毅力的.
來吧!打吧!.
他竟然打了一整晚.
打到黎明.
在二十五節經文這樣說.
乃人見自己勝不過他.
就將他的大腿窩摸了一把.
雅各的大腿窩正在摔跤的時候就扭了.
大家留意經文這句說話.
乃人見自己勝不過他.
大家認為這個人.
這個可能是神.
這個可能是神的使者.
是否真的勝不過雅各呢?.
上帝有沒有機會勝不過我們呢?.

$^{441}$其實上帝給予我們自由權.
很多時候我們堅持用自己的方法.
堅持要走自己的路.
上帝好像真的勝不過我們.
但是上帝沒有任由我們.
明不明白我的意思?.
他不是獨裁者.
他會尊重.
應該說他也任由我們.
但是他不會任由我們.
到一個地步.
就是讓我們自生自滅.
大家記不記得.
當雅各第一次逃亡.
就是剛才我說他騙了阿爾素的祝福.
然後逃亡的時候.
上帝讓他在伯特利.
在夢裡面見到異象.
是記載在《創世記》28章.
在13節上帝還自我介紹.
自我介紹.
我是耶和華.
你祖先亞伯拉罕的神.
也是以撒的神.
然後上帝就向雅各賜下應許.
在14至15節.
他說:你的後裔會像地上的塵沙一樣.
要向東西南北開展.
地上萬族因你和你的後裔得福.
我也與你同在.
你無論往哪裡去.
我必保佑你.
領你歸回這地.
總不離棄你.
直到我成全了.
向你所應許的.
上帝很人性化.
他很明白雅各在那一刻的無助和孤單.
上帝也很知道我們的壞事.
就是我們無力了.

$^{481}$我們就會仰望他.
而上帝不介意我們到那一刻才找他.
上帝讓雅各在伯特利經歷異象.
為的是提醒他什麼.
就是你的身份是何等的寶貴.
你有上帝做你家族的神.
你有上帝的應許作為保障.
弟兄姊妹你們明白嗎?.
對於一生拼搏.
只顧著贏贏贏的人來說.
其實這是很好的提醒.
就是你要繼續用自己的方法.
還是你放手交託.
這位已經向你納下保證的上帝.
他照著你.
那你還要靠自己去拼搏.
還是你相信.
你雙手去接他要給你的應許呢?.
當然人是環顧無知的.
一生人當中.
如果上帝只是提醒你三,四次.
你就已經醒了.
已經很有福了.
有些人提醒你幾十次.
上百次也說不定.
一些很熟悉Gary和我的弟兄姊妹.
應該也聽過.
在他很聲色犬馬的時候.
我曾經祈禱.
Joe最近也提起.
祈禱什麼呢?.
求神將他趕去絕路.
不認識,未聽過的弟兄姊妹.
你不要以為我是一個惡毒的妻子.
其實在這裡先聲明.
我不是教妻子這樣跟老公去怠妒.
而是當時我是懷著一個破釜沉舟的心.
因為我一心.
我決心要將我的丈夫挽回.
純粹是經過不斷向神呼求之後.

$^{521}$一種個人的領受.
所以你不需要學我的.
我做好心理準備.
如果神真的將Gary趕去絕路.
其實是什麼呢?.
也是我的絕路.
嫁雞除雞,嫁狗除狗.
還是我一對子女的絕路.
他們現在坐在這裡.
也沒有去絕路.
果然在名利相失的時候.
走肉朋友全部離開了他這條絕路.
他當時做了很多錯誤的選擇.
將來有興趣可以來問我.
Gary開始了倦之還.
他尋求神的憐憫.
我還記得他邀請我一起跪下.
當時我們還住在錦繡.
我們在睡房裡跪下.
我們一起向神認罪.
我們修復婚姻.
其實也是修復和上帝的關係.
在那一刻.
我們真的切實感受到.
上帝是何等的慈愛.
是可敬可畏.
雅各也是.
當他走投無路的時候.
上帝親自來到安慰和兼顧他.
他在夜夢中醒來.
他敬畏之情就油然而生.
於是他將要做枕頭的石.
拿來立碑.
在上面囂遊獻給上帝.
就叫那個地方做伯特利.
即是上帝的家.
又向上帝許願十一奉獻.
但是弟兄姊妹.
雅各做了這麼神聖的儀式.
是不是從此脫胎換骨.

$^{561}$立地成聖呢?.
不是的,大家都知道.
整個聖經也知道.
雅各在日後的人生路上.
依然喜歡用自己的方法來做人.
剛才我分享.
Gary被神趕去絕路.
他重新回歸神.
也不是一步成聖.
他依然有很多的老鵝.
好像以色列人一樣.
起起跌跌.
每次神又把他拉回來.
他好像雅各一樣.
他喜歡靠自己.
大家是否相信我這麼說.
他過去是很靠自己.
但是每一次跟神相遇的經歷.
其實是什麼呢?.
其實好像我們建房子,打樁.
打樁不可能打一次就堅固.
如果你附近正在建房子.
你就知道有多痛苦.
一條一條地打.
打完一次又一次.
根基才會穩固.
我們的屬靈生命也是這樣.
不可能一步登天.
生命的建立就是一次又一次的打樁.
又喧鬧又痛苦.
是不是?.
但是一次又一次.
直到人生的某一刻.
我們忽然醒覺.
原來上帝一直在這裡.
原來祂早就似下應許.
我想到那個時候.
我們就懂得真正的放手.
讓祂作生命的主.
Gary就是這樣.

$^{601}$雅各又如何呢?.
他來到人生最疲憊的一刻.
前路茫茫.
打了一整晚樁.
心心靈都很疲乏.
這個時候.
天使只不過在他的大腿窩裡摸一下.
他就扭了.
這一刻好像真的什麼都沒有.
連摔樁都好像不行了.
因為跛了.
這一刻不知道前塵往事.
是否立刻浮上雅各的心.
你想想年輕的時候.
用一碗紅豆湯換到長子的名分.
種下的是兄弟之間的仇恨.
當爸爸已殺.
誰誰老去.
雅各又騙了應該屬於已數的祝福.
換來的是長達二十年的離鄉別井.
沒有根.
雖然在舅父拉班取得心愛的妻子.
用計致富.
最後又要帶著所有的資產逃亡.
你覺得他有什麼感受呢?.
在第26節.
這裡說那人說.
不知是上帝還是使者說.
天來明了,容我去吧,你放我走吧.
雅各說:你不給我祝福.
我就不容你去.
事實上雅各真的這樣.
他在被扭傷的腳之前.
他是很強壯的.
否則他怎會整晚摔跤.
摔角手也不一定可以.
如果天使不摸他的大腿窩.
我想他應該會一直打.
過去在雅各的每一項計劃裡.
他從來沒有被上帝做過主導.

$^{641}$他做完義夢.
獻了做了祭禮.
很神聖的儀式.
然後他又用自己的方法.
但現在他最後一個把握.
就是他的腿扭了.
他還有什麼把握呢?.
他這一刻唯一能夠依靠.
無非就是每一次在他軟弱的時候.
幫他的神.
曾經次下應許的神.
都保證了.
於是他想到了.
他拼命纏著這個人要祝福.
因為雅各很清楚.
祝福是來自上帝.
大家有沒有看過新約希伯來書.
七章七節怎麼說.
他說:向來是位份大的.
比位份少的祝福.
這是無可爭議的.
就是說雅各終於怎樣.
終於承認我是渺小的.
我是不行的.
整晚我和你打架.
我要比你厲害.
但是這一刻我承認了.
我承認你的位份比我大.
他默認了唯有從上帝而來的祝福.
才是最重要的.
而不是靠自己的聰明才智.
在二十七節.
這個人就說了.
你叫做什麼名字?.
雅各說:我叫做雅各.
剛才我說了.
這個人是上帝或上帝的使者.
所以不可能他不知道雅各的名字.
我經常說.
上帝的發問不是上帝不知道.

$^{681}$我經常用這個例子.
因為在《創世記》裡.
阿當,喀哇犯罪了.
上帝問你們在哪裡.
難道上帝真的不知道嗎?.
被你們騙了嗎?.
上帝問是想人自己去想.
想一想.
是呀,我在哪裡呢?.
我是誰呢?.
是很真實的.
我記得我和Gary那時.
婚姻出現了問題.
我從神而來.
我領受到我先生對我不忠.
我當時是很想我先生.
能夠向我坦白.
以致他將來不用再被魔鬼控訴.
但我高估了我屬靈的能力.
我以為我能夠撐得住.
誰知道他向我坦白的時候.
我哭了一場.
很怕他將來再背叛我.
於是我和神對話.
神啊,當你知道Gary犯罪的時候.
你是怎樣的?.
神說我很傷心.
我比你更傷心.
神還說我到處去找他出來.
那時我不是很熟悉聖經.
我想不到創世記.
我事後才知道.
原來人犯罪.
上帝是會到處找他出來.
很真實.
現在上帝問雅各.
你叫什麼名字.
雅各也要思考.
因為這個名字代表什麼?.
因為雅各這個名字代表.

$^{721}$抓住腳跟.
欺騙者或取代者.
當你報上名來.
你說我是一個欺騙者.
我是一個抓住腳跟的人.
其實很醜怪.
所以當他報上名字的時候.
他一定會不期然地想起.
我這一生是怎樣拼搏.
我又要抓住這個,又要抓住那個.
我都是用一些狡猾的手段.
在過去又一次一次地抓住.
他是不是真的抓到呢?.
如果抓到.
為什麼在這一刻他會有這樣的下場呢?.
他壓上了妻兒,他壓上了一切.
他卻連保住自己的命的把握都沒有.
上帝只是要雅各看清楚認清楚這一切.
在28節.
乃人說:你的名不再叫雅各.
要叫以色列.
因為你與神,與人教力都得了聖.
大家知不知道聖經裡面的改名是有特定的意思的.
它可以標記某一種事件的起源.
也可以代表某一個人身份的轉變.
例如在《創世記》17章.
我們都記得上帝叫亞伯蘭改名.
給了一個名字,他叫做亞伯拉罕.
就是標記了上帝與人立下應許之約.
又標記了從今以後.
亞伯蘭就成為多國之父亞伯拉罕.
同樣地,當雅各要改名.
改名做以色列的時候.
是標記了上帝揀選的民族以色列.
在歷史上.
往往在宗主國要立繁殊國的王的時候.
有時都會幫繁殊國的王改名.
為什麼呢?.
代表宗主國的權勢凌駕於繁殊國之上.
而雅各被改名.

$^{761}$也代表上帝在行使他在雅各之上的權柄.
是上帝讓雅各成為上帝子民領袖的身份.
上帝給雅各這個新名以色列.
意思是上帝勝過.
上帝要在雅各的生命中得勝.
上帝要讓雅各由一個喜歡用妓女去考取豪奪的人.
和一個經常做壞事躲避別人報復而逃亡的人.
脫胎換骨.
成為上帝和人訂立文約的代表繼承人.
因為雅各的祖先已經和上帝立約.
雅各現在是一個代表繼承人.
這才是真正的得勝.
因為上帝在人的生命中得勝.
在三十節裡.
雅各給地方起名叫比魯伊勒.
意思是我面對面見到神.
我的生命得著保存.
其實舊約講得很清楚.
上帝的榮耀是人不能夠承擔的.
人見到上帝的榮耀一定死的.
但是上帝卻是很慷慨.
他願意以不同程度的榮耀.
向不同的人去顯現.
好像摩西.
上帝讓他見到他的背面.
這是上帝給摩西獨有的尊榮.
也表明了摩西和上帝的關係.
是何等親密.
現在雅各就是這樣.
雅各經歷了上帝的臨在.
他感覺到他已經和上帝面對面.
有了真正的關係和接觸.
在第三十一節,三十二節.
日頭剛出來.
雅各經過比魯伊勒.
他的大腿變得很強.
然後知道以色列人不吃大腿窩的筋.
直到今天.
為何?因為要紀念雅各.
祖先.

$^{801}$大腿窩的筋已經沒有了.
扭傷了.
上帝透過使者在雅各的大腿窩.
摸了一把.
讓他成為強子.
你想想在現實生活裡.
其實是一個徹底的失敗.
是否一個徹底的失敗?.
因為你已經成為殘廢.
並且你身體留下了一個很醜陋的印記.
一個慘敗的終身印記.
現在連腿都變得強了.
雅各又不能走.
前路茫茫,生死未卜.
在人們來說,雅各已經輸了.
但是,大英姐妹,你有沒有想過.
其實這個強腿的印記.
其實正是他生命裡一個最寶貴的印記.
甚至是整個以色列民族的印記.
因為雅各的強腿記號背後.
是上帝對以色列民族的祝福.
上帝揀選以色列成為他的選民.
所以,親愛的弟兄姊妹.
我很想大家去想想.
在你的生命裡.
你有沒有像雅各那樣的醜陋印記呢?.
我們叫他做屬靈的強腿印記.
沒有人敢回答我有.
我告訴你,我和Gary有.
而且不止一個.
我們都有很多.
例如,屬我們的弟兄姊妹都知道.
我們曾經先後破產.
我們幾乎婚姻破裂.
但是,在每一個人體內不光彩的印記.
可以在神的恩典裡.
成為一個榮神益人的見證.
所以,Gary從來都不介意.
我講道時提起他的壞事.
以前,我初初實習.

$^{841}$Gary還未讀神學.
我去講道時.
他會載我和兒子和女兒.
當然是和我.
我在台上講道.
他和兒子和女兒在台下聽我講道.
不知為何,我記得的位置就是今天.
笑紅和艷瓊夫婦的方向.
他們是我們認識了四十多年的弟兄姊妹.
今天特意從希伯倫堂過來.
每一次我講道.
兒子和女兒可能都會打瞌睡.
因為年紀太小.
但是Gary.
眼泛淚光.
很感動地看著我講道.
所以我知道今天我在這裡講道的時候.
Gary如果知道.
他一樣會眼泛淚光.
很開心,老懷安慰.
姓老的,真的老懷安慰.
自從今年三月.
Gary患上重病.
身體一日一日衰殘.
但是他對生命的反省.
和對神的信服.
我看到比過去每一刻都來得深.
來得真,來得堅定.
主的僕人Gary已經息怒歸主了.
特別多人稱讚他,不知為何.
是否見人身體有事,哄哄他.
不是的,其實我都為他驕傲.
這些說話都是由心而出.
不過我們不會因為他離世.
就刻意美化他,是嗎?.
他和我們一樣有翹腿的印記.
看起來不光彩.
但是充滿上帝的賜福.
試問一個完美的人.
還需要完美的神嗎?.

$^{881}$不可能.
正正是這些印記.
讓我們知道我們要謙卑,要依靠神.
弟兄姊妹,今天Gary不在了.
但將來我們會在天家相遇.
我們走的路不一樣.
但目標是一樣的,是嗎?.
大家還記不記得Gary經常說.
他說要盡心盡意盡力.
盡性盡意盡力.
我不聽話了,留下了一個盡性.
去愛神.
盡心盡性盡意盡力愛神.
很盼望我們在神的塗灶和管教.
不要怕管教,上帝愛他才管教.
我經常求上帝管教我.
管教我老公,我的兒子,我的女兒.
在他的塗灶管教之後.
願我們都能夠活出得性的生命.
囉嗦一點,再說一次.
上帝在我們生命裡得性.
才是我們真正的得性.
弟兄姊妹,願不願意上帝在你生命裡得性?.
不如我們一起站立.
我為大家禱告,祝福大家好嗎?.
我們一起禱告.
親愛的主,你看見現在站立的弟兄姊妹.
不是因為我叫他們客氣,給面子.
是因為他們真的被你的靈感動.
他們說願意.
他們願意上帝在他們生命裡得性.
上帝,我們人是善忘的.
但是你記得的.
你就牢牢記住他們說.
他們願意上帝在他們生命裡得性.
以致他們能夠得性的許願.
願上帝成就.
願上帝堅固.
願上帝在每一個人的生命裡動工.
願上帝除去我們生命裡所有的癡慾邪情.

$^{921}$縱然我們起起跌跌.
但是靠著你每一次再起來.
每一次走得更穩固.
直走到你為我們立下的標竿.
主啊,願你得榮耀.
我們拭目以待.
要看我們生命裡如何因著你.
發出光彩,以致神你得榮耀.
我們祈禱交托,奉靠我主耶穌基督.
得性的明智祈求.
阿們.
請坐.
\newpage



\section{路加福音 15:11-32}
\label{sec:ZLhqPiukKhM}
\textbf{2023.12.24 《回家》 路加福音15\_11-32 (和修版) 曾敏傳道}
\newline
\newline
連結: \href{https://youtube.com/watch?v=ZLhqPiukKhM}{\texttt{https://youtube.com/watch?v=ZLhqPiukKhM}} ~~~~ 語音日期: 2023-12-27
\newline
\newline
\hyperref[sec:drs0Z9UDTE0]{\small{< < < PREV SERMON < < <}}
~
\hyperref[sec:index_chronic]{\small{[返順時目]}}
~
\hyperref[sec:index_scriptual]{\small{[返順卷目]}}
~
\hyperref[sec:6_IYA7gmRpA]{\small{> > > NEXT SERMON > > >}}
\newline
\newline
路加福音 15:11-32
\newline
\begin{longtable}{cl}
\hline
\hline
章節 & 經文 (和合本修訂版)\\
\hline
15:11 & \begin{tabularx}{0.7\textwidth}{X} 耶穌又說：「一個人有兩個兒子。 \end{tabularx} \\ \\ \relax
15:12 & \begin{tabularx}{0.7\textwidth}{X} 小兒子對父親說：『父親，請你把我應得的家業分給我。』他父親就把財產分給他們。 \end{tabularx} \\ \\ \relax
15:13 & \begin{tabularx}{0.7\textwidth}{X} 過了不多幾天，小兒子把他一切所有的都收拾起來，往遠方去了。在那裡，他任意放蕩，浪費錢財。 \end{tabularx} \\ \\ \relax
15:14 & \begin{tabularx}{0.7\textwidth}{X} 他耗盡了一切所有的，又恰逢那地方有大饑荒，就窮困起來。 \end{tabularx} \\ \\ \relax
15:15 & \begin{tabularx}{0.7\textwidth}{X} 於是他去投靠當地的一個居民，那人打發他到田裡去放豬。 \end{tabularx} \\ \\ \relax
15:16 & \begin{tabularx}{0.7\textwidth}{X} 他恨不得拿豬所吃的豆莢充飢，也沒有人給他甚麼吃的。 \end{tabularx} \\ \\ \relax
15:17 & \begin{tabularx}{0.7\textwidth}{X} 他醒悟過來，就說：『我父親有多少雇工，糧食有餘，我倒在這裡餓死嗎？ \end{tabularx} \\ \\ \relax
15:18 & \begin{tabularx}{0.7\textwidth}{X} 我要起來，到我父親那裡去，對他說：父親！我得罪了天，又得罪了你， \end{tabularx} \\ \\ \relax
15:19 & \begin{tabularx}{0.7\textwidth}{X} 從今以後，我不配稱為你的兒子，把我當作一個雇工吧。』 \end{tabularx} \\ \\ \relax
15:20 & \begin{tabularx}{0.7\textwidth}{X} 於是他起來，往他父親那裡去。相離還遠，他父親看見，就動了慈心，跑去擁抱著他，連連親他。 \end{tabularx} \\ \\ \relax
15:21 & \begin{tabularx}{0.7\textwidth}{X} 兒子對他說：『父親！我得罪了天，又得罪了你，從今以後，我不配稱為你的兒子。』 \end{tabularx} \\ \\ \relax
15:22 & \begin{tabularx}{0.7\textwidth}{X} 父親卻吩咐僕人：『快把那上好的袍子拿出來給他穿，把戒指戴在他指頭上，把鞋穿在他腳上， \end{tabularx} \\ \\ \relax
15:23 & \begin{tabularx}{0.7\textwidth}{X} 把那肥牛犢牽來宰了，我們來吃喝慶祝； \end{tabularx} \\ \\ \relax
15:24 & \begin{tabularx}{0.7\textwidth}{X} 因為我這個兒子是死而復活，失而復得的。』他們就開始慶祝。 \end{tabularx} \\ \\ \relax
15:25 & \begin{tabularx}{0.7\textwidth}{X} 「那時，大兒子正在田裡。他回來，離家不遠時，聽見奏樂跳舞的聲音， \end{tabularx} \\ \\ \relax
15:26 & \begin{tabularx}{0.7\textwidth}{X} 就叫一個僮僕來，問是甚麼事。 \end{tabularx} \\ \\ \relax
15:27 & \begin{tabularx}{0.7\textwidth}{X} 僮僕對他說：『你弟弟回來了，你父親因為他無災無病地回來，把肥牛犢宰了。』 \end{tabularx} \\ \\ \relax
15:28 & \begin{tabularx}{0.7\textwidth}{X} 大兒子就生氣，不肯進去，他父親出來勸他。 \end{tabularx} \\ \\ \relax
15:29 & \begin{tabularx}{0.7\textwidth}{X} 他對父親說：『你看，我服侍你這麼多年，從來沒有違背過你的命令，而你從來沒有給我一隻小山羊，叫我和朋友們一同快樂。 \end{tabularx} \\ \\ \relax
15:30 & \begin{tabularx}{0.7\textwidth}{X} 但你這個兒子和娼妓吃光了你的財產，他一回來，你倒為他宰了肥牛犢。』 \end{tabularx} \\ \\ \relax
15:31 & \begin{tabularx}{0.7\textwidth}{X} 父親對他說：『兒啊！你常和我同在，我所有的一切都是你的； \end{tabularx} \\ \\ \relax
15:32 & \begin{tabularx}{0.7\textwidth}{X} 可是你這個弟弟是死而復活，失而復得的，所以我們理當歡喜慶祝。』」 \end{tabularx} \\ \\
[1ex]
\hline
\hline
\end{longtable}
$^{1}$(自動翻譯).
各位聖誕快樂.
感恩在這個普天同慶的日子.
我們回到神的家裡.
一同去慶祝.
家對我們每個人來說是很重要.
每個人都需要回到家裡.
我們需要有自己的家.
剛才一路聽梁傳道講的時候.
心裡也很感動.
我認識梁傳道這麼久.
今天才真正聽回.
以前梁傳道作為爸爸的時候是怎樣.
我認識他的時候.
他已經是一個廿四孝的爸爸.
這個假期梁傳道也是拿來荷蘭探望女兒.
他經常是這樣說.
我的女兒怎樣?我的孫子怎樣?.
我很快樂.
整個口都說他的女兒和孫子.
我知道他很疼愛他的女兒.
我相信作為他的女兒是很幸福的.
直到今天才聽回.
原來他過去也有這些經歷.
在裡面看到耶穌基督對他有很大的改變.
我很同意梁傳道所說的.
今天各位朋友.
你能夠坐在這裡是一個很大的福氣.
那個帶你來的朋友.
他很關心你的生命.
因為今天我們要說一件事.
對你的生命是很大的幫助.
我們開始之前有一個禱告.
很愛我們每一個.
今天你召集我們回來.
是要去聽你的聲音.
去聽你的話語.
求你讓我們每一個的心都開放.
我們的耳朵很柔軟.
願意聆聽.

$^{41}$願意讓我們聽入心裡去.
讓我們明白.
我們就回應你.
奉耶穌基督的名字祈禱.
阿們.
今天大家讀的經文.
也說到一個爸爸.
梁傳道是一個爸爸.
聖經裡有一個爸爸.
這個爸爸有什麼特色呢?.
我們看整段經文.
知道這個爸爸非常非常有錢.
很富有.
他不單有兩個兒子.
可以給他們很多財產.
他還有很多的故宮.
在這個爸爸的家裡.
從來不缺錢.
可以說是生活無憂.
非常豐富.
做他的兒子很有福.
不用工作.
就是這樣享受.
這個爸爸應該很愛他們兩個兒子.
但是他說到.
有一天.
這個兒子就跑去跟爸爸說.
爸爸.
把我的財產給我.
哇 你的語氣這麼兇.
就是這樣說.
爸爸還沒死.
我們什麼時候分家產的?.
很年輕的時候.
爸爸就跑去跟他說.
每個人都是這樣.
爸爸三十多四十歲就已經說.
快點把你的財產拿過來.
讓我享受.
我要去享受人生環遊世界.

$^{81}$我不用工作了.
會不會這樣呢?.
一定不會.
是爸爸快要死了.
快要死的時候才跟他說.
甚至有些已經死了.
遺產.
我才去分他的遺產.
現在爸爸不是快要死.
沒有死.
兒子就跑去跟他說.
把我的財產給我.
是不是給了爸爸.
我們去環遊世界.
享受人生.
我們一起去.
不是.
他說這裡是在收拾.
收拾好.
變賣了.
然後拿去自己走.
我不要這個家.
我不要跟爸爸一起生活.
我要跟哥哥一起生活.
我要走.
我遠走高飛.
我有我的人生.
我有我的夢想.
去到有多遠.
去到多遠.
我想怎樣就怎樣.
我的人生.
我作主.
剛才說.
是不是.
在山上.
有山靠山.
有什麼靠什麼.
是不是.
我這些都是不靠的.

$^{121}$靠自己.
有錢就行.
最重要是有錢.
所以爸爸現在拿過來.
作為爸爸.
你作為他爸爸有什麼感受.
我很愛我的兒子.
想跟他共享人生.
我將我最好的.
所有的東西.
我都想留給他.
我想我兒子不用捱苦.
但現在這一刻.
我還沒死.
很久都沒有.
我兒子就要我分家產.
還說不想跟我們一起生活.
我要遠走高飛.
我可以說什麼.
聖經這裡沒有說任何話.
是沉默的.
沒有說爸爸有說任何話.
但我們可以想像一下.
爸爸這一刻的心情.
是不是很好受.
一點都不好受.
過了沒幾天.
兒子就真的這樣.
整齊了.
可能就是變賣了.
變成現錢.
這樣才可以走得遠.
拿著牛.
拿著羊.
走得這麼遠.
變賣了.
拿著錢.
走到遠方去.
在那裡.
什麼都做完.

$^{161}$爸爸不在身邊.
家人不在身邊.
鄉親父老不在身邊.
我想怎樣過生活.
是我的事.
姓氏犬馬.
嫖飲賭取.
什麼都做完.
所以都要說放蕩.
我有的是錢.
拿出來.
我怎樣花.
我作主.
開心.
現在我作主.
沒有人管我.
沒有人在身邊說我是怎樣.
花著花著花著錢.
剛開始還有很多很多很多.
不用怕.
繼續花.
花完一天.
兩天.
三天.
花完花完.
沒有了.
竟然沒有了.
竟然我花光了錢.
這樣都可以.
現在怎麼辦.
找份工作.
真的要想些方法.
現在我手上都沒有東西.
沒有東西拿了.
真的最糟糕.
最糟糕的時候.
我這個地方.
有饑荒.
到處都沒有東西吃.
食物多缺乏.

$^{201}$我現在要投靠誰.
沒有東西吃.
我又沒有錢.
你想有錢去買食物.
都沒有錢買.
到處都這麼缺乏.
糧食比什麼都貴.
糧食很貴.
比黃金還貴.
那我怎麼辦.
在這個時候.
身邊什麼人都沒有.
什麼親人都沒有.
我很窮.
這個時候小兒子真的很窮.
這麼窮的時候.
可以怎麼做.
都要找人.
要投靠一些人.
都要給我食物.
不然我怎麼可以挨下去.
我已經去到窮途末路.
沒有路可走.
個個都投靠不了.
有個養豬的.
我都投靠他.
今天我們可能覺得.
養豬的都不錯.
香港一頭豬有多貴.
有沒有計算過.
曾經我們粗略計算過.
一頭豬800元比較便宜.
要千多元.
一頭豬都很貴.
如果全部養十幾頭豬.
或者幾十頭豬.
很厲害.
那是一筆很可觀的費用.
這麼興旺養豬的人.
當時的猶太人覺得.

$^{241}$養豬是一個非常卑賤的職業.
猶太人自己本身.
一定不會養豬.
骯髒邋遢.
在信仰上.
禮儀上都是不潔的.
哪有人從事養豬的.
不會吃豬肉.
也不會養豬.
外邦人.
不屬上帝的民族.
才會去養豬.
你呀你呀你呀.
千萬不要養豬.
我看不起你.
比地底泥還要差.
現在沒有錢.
沒有錢.
如果鄉親父老知道我投靠養豬.
真的被人看得我多低.
看得我多差.
比地底泥還要差.
我真的沒得生活.
我真的捱不下去.
我需要.
我需要投靠別人.
找份工作.
養豬也要.
就去養豬.
去養豬的時候.
我希望別人給我食物.
大饑荒期間.
養豬.
我都沒有一口好吃.
真的沒有一口好吃.
我這麼慘.
錢收入微薄.
沒有錢.
我見到豬.
我已經流口水.

$^{281}$見到豬吃豆鴿.
我都很想.
有沒有人分些豆鴿給我吃.
我想和豬爭吃那些豆鴿.
因為我好餓.
我不想死.
豆鴿都沒得我吃.
真的這麼慘.
怎麼辦.
再這樣下去.
會死在那裡.
真的會死.
餓都會餓死.
我不能捱很久.
人生怎會去到這樣.
去到窮途末路.
連豬的豆鴿都吃不到.
在這個時候.
小儀姐真的想起.
我爸爸.
我還有個爸爸.
很遠很遠的地方.
我爸爸在那裡.
我爸爸就有錢.
是不是只能分財產.
給我和哥哥.
說真的.
哥哥一定分得多些.
爸爸那份財產應該分三份.
哥哥拿兩份.
大兒子就多些.
我做小兒子一份都很多.
撇除這些.
我爸爸還有錢養很多故宮.
那些宮人肯定有得吃.
有瓦遮頭.
有工作.
他們一定不能死.
我爸爸那麼有錢.
不好了.

$^{321}$快去找爸爸.
投靠爸爸.
我相信我不會餓死.
真的不會餓死.
是不是比在這裡這樣好.
連豆鴿都沒得吃.
餓死嗎.
我現在就回去.
我和爸爸說.
我真是個壞兒子.
我得罪了天.
又得罪你.
我求你寬恕我.
我不配稱為你的兒子.
從今以後你把我當故宮.
我知道我離開家的那一刻.
多傷你心.
我真是個壞兒子.
你現在接納我.
做故宮都好.
我幫你打工.
至少有飯吃.
我不會餓死.
就這樣.
回去吧.
我要回去找我爸爸.
於是小兒子就回去.
起來回來.
去找他爸爸.
這一刻.
大家想想小兒子的心情.
想著回去投靠爸爸.
但心裡會不會有些擔心.
爸爸會不會接納我.
我真是個壞兒子.
你看隔離鄰舍.
怎樣看我爸爸.
他的兒子真是壞到這樣.
未死就希望他快點死.
現在就拿著財產.

$^{361}$走去家裡.
不要和家人一起生活.
這個壞兒子回來.
可能鄰居到時知道我回去.
都會笑我.
爸爸可能會生我的氣.
不接納我都不奇怪.
但我都要回去認帝位.
要認罪.
他一路回去.
大家猜猜.
看到些什麼.
大家猜猜.
看到些什麼.
他一路回去.
可能心情都有些忐忑.
不知道爸爸會不會接受我.
會不會把我趕出家.
不讓我進去.
誰知道.
原來他不用擔心.
相隔很遠很遠.
他看到誰.
很遠.
是很遠.
看到爸爸.
爸爸你竟然在那裡等我.
真的爸爸在那裡等他.
真巧合.
我那天回去.
他就在那裡等我.
只是那一天.
不止.
爸爸應該等了很多天.
每天都在那裡.
想著兒子會回來.
等他.
我想他回來.
一路等.
一日兩日三日.

$^{401}$等等等.
終於就那一天.
兒子決定回去.
爸爸等到了.
所以他遙遠看到他.
爸爸看到的時候.
心就動了.
動了一次.
心就跑過去抱著他.
被別人.
很多人可能被我.
記得在那裡罵一下.
臭小子.
知錯了嗎.
走得那麼爽快.
你有沒有想過我的感受.
你們這麼多年.
你們怎樣對我.
搞到我這樣.
可能我罵他一下.
因為爸爸沒有.
抱著他.
很疼他.
還要親他.
好像中間.
沒有發生過什麼事.
很寶貝.
抱著他.
這個時候.
兒子認錯.
爸爸我得罪了天.
得罪了你.
我不配稱為你的兒子.
爸爸卻吩咐.
把這個上好的袍子.
給他.
還要把戒指.
戴在他手上.
給他穿鞋.
一連串的動作.

$^{441}$爸爸要告訴所有人.
我要恢復他.
是兒子的身份.
兒子有沒有說.
爸爸你當作是.
故宮就行了.
我不配做你的兒子.
你當我故宮就行了.
有沒有這樣說.
兒子去到這個階段.
末端.
未,只是說我不配做你的兒子.
未說到.
我現在為你打工.
爸爸已經要還回他.
抱著他.
而且.
恢復他做兒子的身份.
兒子和故宮.
是很不同的.
身份不同.
故宮只是服侍.
主人怎樣吩咐.
工人怎樣回應.
怎樣做.
是沒有決定權的.
而且.
主人再怎樣要求.
都好.
故宮不能說.
這樣不合理.
這樣太辛苦了.
工時太長了.
工資太少了.
我要有我的權益.
在這裡如果真的要做故宮.
是沒有這個權益的.
地位亦都很低微.
但兒子不同.
兒子.

$^{481}$是要去享福.
去享受.
父親給他的種種福氣.
在家裡的位置.
非常尊貴.
寶貴,被愛護.
被愛惜.
父親不會把他當工人去看.
就是他的寶貝兒子.
恢復他的身份.
還要搞一個慶祝宴會.
把一隻肥牛毒.
叫人分款拿過來.
燙了他.
一同慶祝.
還要奏響樂曲.
一同宣佈.
我兒子回來了.
簡直是死而復活.
是否失而復得.
很開心,一起慶祝.
很開心.
爸爸很開心.
在這時候.
是否所有人都很開心.
大家覺得.
如果我作為家庭一份子.
我們應該要替他感到開心.
爸爸得到失去的兒子.
弟弟又回來了.
可惜.
家裡的哥哥.
不開心.
他的反應.
完全超出我們想像.
超出我們想像.
當時.
慶祝的宴會.
開始的時候.
哥哥在田裡工作.

$^{521}$他回來.
離家不遠.
就聽到樂器的聲音.
跳舞的聲音.
問他一個同伴.
發生什麼事.
為何這麼熱鬧.
好像很開心.
同伴就告訴他.
弟弟回來了.
他多開心.
看到他沒事.
人多好.
很開心就慶祝.
哥哥本來應該很開心.
這麼好,弟弟回來了.
不是.
馬上就生氣.
生氣得不肯進去.
要爸爸出來.
哄他進去.
接著.
哥哥.
積了這麼久的氣.
就出來了.
忍了很多年.
你看.
我服侍了你很多年.
我沒有違背你的命令.
我很聽話.
我很乖.
你叫我做什麼我就做什麼.
我多勞苦.
你連生養都沒有.
你連令我和朋友開心都沒有.
好像爸爸很嚴苛.
爸爸只要諸多要求.
只要勞役他.
一點好處都沒有.
你看.

$^{561}$那個臭小子.
你弟弟一回來.
你知道他在外面.
花天酒地.
不斷放蕩.
花光你的錢.
和這麼骯髒的人一起.
花光你的錢.
回來.
你竟然.
還在那裡宰肥牛毒.
生養一隻都沒有.
你就在那裡宰肥牛毒.
你說你是不是大小超.
你偏心.
很生氣.
其實你也很妒忌.
為什麼爸爸對弟弟這麼好.
為什麼你不是對我好.
他一點都沒有.
因為弟弟.
失而復得.
死而復活.
而快樂.
為了這個家有這麼大的消息而快樂.
沒有.
只是這樣.
很難過.
很不開心.
爸爸怎樣回應他.
爸爸說.
你經常和我一起.
我的東西.
就是你的.
那裡有很多.
我們現在去看.
大兒子.
剩下的財產.
要給他是兩份.
是小兒子的雙倍.

$^{601}$很多.
其實在爸爸身邊.
我相信不只是那兩份財產.
所有爸爸的.
都是他.
你開口一句.
今天我請朋友來.
可不可以宰一隻肥牛.
宰幾隻羊.
不要說一隻羊.
爸爸都給.
我相信爸爸真的很慷慨.
一點都不會吝嗇.
為了那隻羊.
那隻肥牛毒.
還要計算你做了多少.
值不值得那隻羊.
值不值得那隻肥牛毒.
不會的.
爸爸是很愛我們.
但很明顯.
大兒子.
不明白爸爸的愛.
所以心裡有很多埋怨.
很多苦毒.
很多憤怒.
在這裡.
我們一路去看的時候.
好像流浪的.
在外漂流的浪子.
就是小兒子.
其實大兒子何嘗不是呢.
人就在爸爸身邊.
但其實.
整個生命.
整個心都在外漂流.
他一樣沒有享受到.
家的溫暖.
享受到家的好處.
爸爸說完那番說話.

$^{641}$就是你的弟弟.
是死而復活.
失而復得.
所以我們應當開心之後.
沒有說到哥哥的回應.
就停在這裡.
我們不知道哥哥.
最後是怎樣看.
但值得我們去想一想.
今天我們看回的時候.
這個比喻是.
耶穌說的.
耶穌今天也要.
和我們說這個比喻.
今天究竟在主耶穌的眼中.
誰是小兒子呢.
誰是小兒子呢?.
大家有沒有想過.
我們是小兒子還是大兒子呢?.
我又不是小兒子.
我又不是大兒子.
小兒子在經文裡.
是說一些稅利.
說一些罪人.
統稱來說.
是犯罪的人.
罪惡.
令我們和上帝.
分開了.
隔得很遠.
罪惡.
深深的.
是綑綁著我們的生命.
你說我哪有犯罪.
殺人放火.
才叫罪.
去偷去搶.
就叫罪.
奸淫擄掠就叫罪.
你看看亡民.

$^{681}$我們全場個個都不知道記號.
哪有罪?.
不是的.
剛才梁全說得很好.
我們為了自己.
為了自己的利益.
我們會犯罪.
不用別人教.
小時候.
很多事都不用別人教.
我們都會做.
而是我們已經犯罪.
得罪上帝.
這麼多年.
我們人都是自私.
會為自己的利益.
自私.
犧牲別人在所不惜.
將最好的東西.
給自己,令別人一定有損失.
可能別人不察覺.
周邊的人不察覺.
但自己知道.
為了爭取利益做了些什麼.
我們的心思意念.
很多東西都是上帝不起源.
我們裡面都會有貪婪.
我們有很多憤怒.
這些憤怒.
傷害自己,傷害別人.
裡面的憤怒.
可以轉化成.
我們語言方面的暴力.
不只是打別人一拳.
打腳踢.
語言也可以暴力.
我們自己被人傷害過.
其實我們自己也傷害過別人.
從小到大我們從來都沒有吵架.
說話的時候.

$^{721}$可以毫不留情.
語言很傷害.
群體裡面可以說很多是非.
透過是非.
去傷害排擠別人.
這麼多年.
有多少東西.
還有很多.
心裡面的污穢的說話.
意念.
口裡的污穢的說話,意念.
不只是青少年.
會說很多髒話.
成年人也有很多.
說真的青少年的髒話從哪裡來.
就是在成年人身上學的.
一回到家裡,關上門.
沒有人聽到.
憤怒的時候就來了.
很多很多.
你問問自己.
如果今天.
上帝說將我和你的一生.
不只是外在的東西.
是所有我們說過的話.
想過的意念.
和人的相處.
每一樣在明處暗處做過的事.
全部上帝將我們拍成一部電影.
今天很簡單.
就在紅恩堂播.
我想問在在.
有沒有誰會說.
好啊,播我的一生出來.
播我的一生吧.
沒問題的.
我覺得自己所說的,所想的.
所做的,完全經得起考驗.
沒問題.
一點點的罪也沒有.

$^{761}$我相信我們不夠膽.
至少我不夠膽.
因為我知道我自己也有犯罪.
所以聖經說得很好.
世人都犯了罪.
規絕神的榮耀.
我們每個人都有犯罪,得罪上帝.
犯罪令我們和上帝分隔開來.
而最後帶領我們面對的是永遠的死亡.
所以梁存道一直說的是生命.
大家剛才有沒有問一個問題.
關生命什麼事.
是不是我健康有什麼問題.
是不是我健康有問題.
所以我就要離開世界.
不是說健康的問題.
是永遠的死亡.
永遠和上帝分開.
我還要為我的罪.
承受永遠的懲罰.
是永遠的.
今天我想和大家說.
我們承受不了.
我和你都面對不了.
這是永遠的死亡.
這是骨頭旱.
沒路可走.
所以我們一定要回家.
回到天父的家.
我們不能繼續流浪.
走向永遠的死亡.
我們沒路可走.
小兒子一定要回家.
你是不是那個小兒子.
小兒子需要回家.
天父上帝何等愛我們.
藉著他的愛子耶穌基督.
為我和你的罪.
已經死在十字架上.
今天我和你很肯定.

$^{801}$我們沒辦法處理我們的罪的問題.
你說我犯了十樣罪惡.
我做十樣好事.
就抵銷.
對不起,抵銷不了.
罪就是犯罪.
這個紀錄就在這裡.
而我要為這個罪付上代價.
所以今天天父藉著他的愛子耶穌基督.
為我和你的罪.
付上十字架的代價.
他死在十字架上.
代我和你的罪已經死了.
今天我們願意相信耶穌基督的時候.
天父就會寬恕我們的罪惡.
好事已過,都變成新的.
所以我們要回家.
回到神的家.
天父很愛我和你.
至於這裡的大兒子指的是誰呢?.
大兒子指的是誰呢?.
剛才有人可能心裡想.
我又不是大兒子,又不是小兒子.
大兒子其實一直都在爸爸身邊.
大兒子一直都在工作.
在服侍爸爸.
那就是誰呢?.
有些人想說,大聲一點說.
大兒子指的是誰呢?.
信徒.
屬上帝的人.
我們已經進入到神的家裡去生活.
在神的身邊.
但我們一直都不明白.
神對我們的恩典是何等的大.
不明白神的愛是何等的寶貴.
總覺得神都沒有給我什麼.
神給了我什麼呢?.
沒有.
只是在這裡做,做,做.

$^{841}$服侍就做很多.
做完這樣做那樣.
你有什麼好處給過我呢?.
我想要這些,要那些.
我看到別人眼都紅了.
你給那個這麼好.
你給那個這麼好,這個這麼好.
那我呢?我也想要.
神給很多恩典給其他人的時候.
我們也很不滿意.
為什麼神可以恩待別人呢?.
我們有沒有為了神.
真的恩待不同的人.
我們要獻上感恩.
弟兄,姐妹為你感恩.
上帝很疼愛你,真是很好.
我和你一起歡喜快樂.
同時間,不是覺得自己很慘.
很淒涼.
上帝對我很嚴苛.
對別人特別好.
不是的,神對我們每一個都很好.
那份愛心,我們是不明白的.
弟兄,姐妹.
我們是不是大兒子.
雖然在神的家生活.
但我們不明白他的愛.
我們的心在流浪.
離開家很遠.
天父很遠.
你有沒有想過天父也很想念我們.
人就坐在這裡.
但心卻從來不在.
你知道天父也很愛我們.
很想念我們,很難過.
有沒有想過.
你是不是大兒子.
天父也等待大兒子回家.
回歸他愛的懷抱裡.
大兒子.

$^{881}$你心裡是怎樣想的.
我很想在這個時候.
和領唱們一起唱兩首詩歌.
領唱一同上台.
今天這首詩歌是唱給大兒子聽.
也是唱給小兒子聽.
很想小兒子和大兒子.
都一同有個回應.
親愛的耶穌.
給我盼望.
給我一個家.
我希望我能夠在這裡.
跟你相遇.
我希望我能夠在這世界.
跟你相遇.
我希望我能夠在這世界.
跟你相遇.
我希望我能夠在這世界.
跟你相遇.
我希望我能夠在這世界.
和你相遇.
我希望我能夠在這世界.
跟你相遇.
我希望我能夠在這世界.
跟你相遇.
我希望我能夠在這世界.
跟你相遇.
我希望我能夠在這世界.
跟你相遇.
我希望我能夠在這世界.
跟你相遇.
我希望我能夠在這世界.
跟你相遇.
我希望我能夠在這世界.
跟你相遇.
我希望我能夠在這世界.
跟你相遇.
我希望我能夠在這世界.
跟你相遇.
我希望我能夠在這世界.

$^{921}$跟你相遇.
我希望我能夠在這世界.
跟你相遇.
我希望我能夠在這世界.
跟你相遇.
我希望我能夠在這世界.
跟你相遇.
我希望我能夠在這世界.
跟你相遇.
我希望我能夠在這世界.
跟你相遇.
我希望我能夠在這世界.
跟你相遇.
我希望我能夠在這世界.
跟你相遇.
我希望我能夠在這世界.
跟你相遇.
我希望我能夠在這世界.
跟你相遇.
我希望我能夠在這世界.
跟你相遇.
我希望我能夠在這世界.
跟你相遇.
我希望我能夠在這世界.
跟你相遇.
我希望我能夠在這世界.
跟你相遇.
我希望我能夠在這世界.
跟你相遇.
我希望我能夠在這世界.
跟你相遇.
我希望我能夠在這世界.
跟你相遇.
我希望我能夠在這世界.
跟你相遇.
我希望我能夠在這世界.
跟你相遇.
我希望我能夠在這世界.
跟你相遇.
我希望我能夠在這世界.

$^{961}$跟你相遇.
我希望我能夠在這世界.
跟你相遇.
我希望我能夠在這世界.
跟你相遇.
我希望我能夠在這世界.
跟你相遇.
我希望我能夠在這世界.
跟你相遇.
我希望我能夠在這世界.
跟你相遇.
祂想我們回到神的家裡.
去生活.
今天如果您.
願意相信耶穌基督.
願意回到神的家.
請您舉起您的右手.
舉高一點.
舉起您的右手.
天父正在等您.
舉高一點您的右手.
舉高您的右手.
天父很想邀請您.
回到祂的家裡.
請您舉高您的右手.
祂正在等您.
正在等您.
將您的生命交給祂.
在神的家裡.
很豐富.
我們不需要再流浪.
我們再唱一次.
親愛耶穌.
給我盼望.
給我一個永恆的家.
親愛天父.
請賞我.
回到祂面前.
不再流浪.
我看到家的光.

$^{1001}$回家回家.
回到永恆愛的家.
喜樂充滿我的心.
活不出讚美.
回家回家.
回到永恆愛的家.
天父張開愛的雙臂.
我一生屬於祂.
回家回家.
回到永恆愛的家.
喜樂充滿我的心.
活不出讚美.
回家回家.
回到永恆愛的家.
天父張開愛的雙臂.
我一生屬於祂.
回家回家.
回到永恆愛的家.
天父張開我的心.
活不出讚美.
回家回家.
回到永恆愛的家.
天父張開愛的雙臂.
我一生屬於祂.
天父張開愛的雙臂.
等著你回到神的家裡.
可能你說其實我不是小兒子.
我沒有啊,我很積極的生活.
我每一天都很努力的工作.
照顧好我的家人.
我常常做好事.
去祝福人群,祝福社會.
我不是小兒子.
我亦都不是大兒子.
但我想和你說.
朋友,我們一樣都是犯罪.
得罪了上帝.
而這個罪同樣.
讓我們和上帝的關係隔得很遠.
這個罪亦都是.

$^{1041}$令我們走向永遠的死亡.
所以在這裡.
我都是邀請你.
即使你不是.
那個虛到盡的小兒子.
我們都是犯罪.
得罪神,走向滅亡.
走向滅亡.
今天很想邀請你同樣的.
舉起你的手.
就是邀請耶穌基督進入你的生命裡.
你回到神的家裡生活.
如果是這樣的話.
我都誠意邀請你.
舉起你的右手.
我們不需要去到人生盡頭.
才能回到.
神的家裡.
去生活.
今天就好,回到這裡.
你說我願意.
相信耶穌基督.
讓祂成為我的救主和生命的主.
你願意的時候,請你舉起你的右手.
將你的生命交給祂.
有沒有願意的朋友,請你舉起你的右手.
有願意的朋友,請你舉起你的右手.
你未曾相信耶穌基督.
今天你聽到.
天父很愛你.
很願意你回到祂的身邊.
如果有的.
有這樣的朋友,請你舉起你的右手.
讓我們知道.
讓天父都知道.
你願意回來了.
有沒有這樣的朋友,請你舉起你的右手.
不要等了.
你不知道什麼時候有機會.
今天就是一個機會.

$^{1081}$天父想你回來,請你舉起你的右手.
這樣的愛.
這樣高深的愛.
願至今古不變.
直到未來.
降臣人世內.
彰顯神的愛.
接近人群.
抹去病菌的愛.
這樣的愛.
甘願犧牲的愛.
這愛傾出寶血.
送我回來.
這愛傾出寶血.
送我回來.
這神除障礙.
血盡流的愛.
世上無人.
如同這份愛.
世上無人.
如同這份愛.
我的主恩心.
比海更深.
這份愛.
解我心中罪困.
淒清眾伴我.
同行孤單中.
伴我前行.
這份情仍是最真.
我的主恩.
高比天更高.
這份愛.
今我竟可遇到.
漆黑中.
沒有迷途.
有谷中.
是我路途.
這份情難以盡數.
我心中有數.
我想邀請.

$^{1121}$如果剛才.
你不敢舉高手.
你心裡也決定.
邀請耶穌基督進入你生命裡.
成為你的救主和生命的主.
你要回到神的家裡.
去生活的朋友.
這一刻.
在你的朋友的陪伴之下.
能夠來到台前.
我們很想為你禱告.
很想為你去祝福.
今天是一個很寶貴的時機.
歸向耶穌基督.
回到天父的身邊.
天父很愛我們.
很愛我們.
不願意我們走向永遠的死亡.
不願意我們再繼續被罪惡捆綁著我們.
天父在等我們.
我們很願意.
如果我們願意.
向天父表明我們的心意.
我回來了.
天父我知道你愛我.
我知道耶穌基督很愛我.
為我犧牲.
如果你願意的話.
邀請你站在台前.
很想為你禱告.
為你祝福.
我們一同在神面前.
作出一個很重要的決定.
如果有願意的.
邀請你現在把握這個機會.
來到台前.
小兒子.
是時候要出來了.
在你的朋友的陪伴之下.
一同出來.

$^{1161}$我們要為你禱告.
邀請大家一同出來.
不需要害怕.
有朋友陪伴你.
你的朋友昔日回歸.
天父懷抱.
就是這樣回應.
在神的愛中.
經歷了很大的恩典.
從今以後我們人生.
不是自己一個.
是天父和我們一起.
如果有願意的朋友.
一同出來.
我們要向上帝.
作出決志的禱告.
我們繼續等.
我知道天父在等待我們當中的人.
我們繼續等.
天父在等待我們.
是何等寶貴.
在等待的時候.
我們一路默默的禱告.
是何等寶貴.
這一份愛.
如果願意的時候.
邀請大家站出來.
我們為你們禱告.
將我們的心意向神說.
從今以後.
我們不需要走向永遠的死亡.
是何等盼望.
是何等愛.
還有沒有朋友願意.
在眾人的鼓勵之下.
在眾人面前.
我們做一個很重要的決定.
將心交給主.
在眾人的鼓勵之下.
在眾人面前.

$^{1201}$我們做一個很重要的決定.
將心交給主.
在眾人的鼓勵之下.
在眾人面前.
我們做一個很重要的決定.
將心交給主.
在這個時候.
我想做另外一個呼召.
我想呼召大兒子.
你知道的.
你知道自己是不是大兒子.
雖然人在教會當中.
但是心都在外面流浪.
一路不明白.
神是很愛我們.
神一直都很愛我們.
很願意將很多福氣賜給我們.
其實神已經賜了這些福氣給我們.
但是我們不明白.
一路在外面流浪.
一路我們的心很難過.
但是神的心更加難過.
今天是時候大兒子要回家.
大兒子要回家.
我們真正明白神的愛.
回應祂的愛.
一生不要再為自己而活.
為上帝而活.
如果有這樣的弟兄姊妹.
請你在座位裡站起來.
去回應上帝的呼召.
你說:神,我今天不再流浪.
我不要做大兒子.
我今天答應你.
我要回到你的身邊.
我要好好地在神的家裡生活.
經歷你的愛,回應你的愛.
如果今天有這樣的弟兄姊妹.
願意委身的弟兄姊妹.
請你在你的位置站起來.

$^{1241}$這是一個很重要的表達.
我們去回應上帝的心意.
我們知道.
自己離開神有多遠.
我們知道.
今天如果我願意.
回到神的身邊.
沒錯,我是大兒子.
我已經離開神很久的時間.
我的心裡有很多傷痕.
很難過,很久.
今天是一個時間.
回到神的身邊.
如果有願意的弟兄姊妹.
我從今以後要委身.
在上帝的愛裡.
一生跟從祂.
真正活出上帝那份愛.
回應神,愛神,愛人.
有這樣願意的弟兄姊妹.
請你在位置裡站起來.
請你站起來.
上帝也想看到我們.
向祂表明.
我們的心意.
有這樣的弟兄姊妹請你在位置裡站起來.
神看到的.
我們的心不要再流浪.
我想邀請大兒子.
剛才我們說的弟兄姊妹.
也請他們站起來.
我們也有一個禱告.
請剛才站在位置的弟兄姊妹也出來.
因為接下來.
請剛才站在位置的弟兄姊妹也出來.
我們也想為你們禱告.
我同樣邀請所有在座的弟兄姊妹站立.
無論今天你帶著什麼心情去參與這個報道會.
我很想你自己也去回應神.
今天對你說的話.

$^{1281}$我們不單止為台前的弟兄姊妹去禱告.
請你也為自己的生命去禱告.
上帝很愛我們.
我們是否真的明白祂的愛.
真的能夠回應祂.
是愛神愛人.
請我們一同開聲禱告.
之後我帶領大家禱告.
祝福你們的生命有美麗.
吃飽了.
( 感 competitions of kindness and kindness of Brother Chan. ).
今天,您藉著浪子的比喻.
您呼喚我們回到您的家裡生活.
天父,我知道今天在新朋友當中.
有感到很感動.
天父,他們未必會出來回應您.
但您知道他們的心想回應您.
天父,求您待會繼續帶領他們.
真正將他們的心門向您打開.
他們就會相信您.
天父,我知道在天上有天君天使等待.
很想為了這個寶貴的時間歡呼拍掌.
因為它是何等寶貴.
從今以後,這些生命是屬於您的.
這些生命不需要走向永遠的死亡.
他們永遠在神的愛裡,在神的家裡生活.
我求您祝福和賜福這些新朋友.
讓他們得著主你要給他們的.
主要是為了帶兒子去禱告.
主要是我們每一個都是帶兒子.
我們人在教會,但我們的心卻很遠離.
我們不知道您這麼愛我們.
您的救恩是如此寶貴.
您給我們的東西遠遠超過我們今天所知道的,所了解的.
天父,求您去寬恕我們.
我們實在得罪您.
求您寬恕我們如此軟弱.
寬恕我們一切的罪惡.
我們要回到您的身邊.
天父,求您兼顧我們的心.

$^{1321}$讓我們的心回到您的身邊.
讓我們的心在您的愛裡繼續走下去.
讓我們為了愛您,愛人,用盡我們的一生.
在台前有一群弟兄姊妹站出來立志.
他們要回到神的家裡.
回到您的懷抱裡生活.
我求您給他們大大的賜福.
我為他們的生命獻上感恩.
因為他們願意在您面前立志為生.
主求您賜福,讓他們不會退後.
讓他們將生命交在您的手裡.
從今以後是很有力的活著.
所靠的是聖靈的大能.
願主您的話語成為他們腳前的燈,路上的光.
主在來的時間,更願主您的愛每一天伴隨著他們.
充滿他們的生命,讓他們感到溫暖有力.
我們有很多弟兄姊妹沒有站出來.
主您也知道他們的需要.
求主同樣要兼顧他們的生命.
要讓他們離開一切他們不討您喜悅的罪惡.
要過聖潔的生活,為您而活.
為上帝的國度而打拼.
主啊,您說過我們先求您的國和您的兒.
其他一切都要加給我們.
所有弟兄姊妹都朝著這個方向努力.
我們唱這首詩歌的第二節.
這樣的愛,這份一生的愛.
這愛永擇不退,但很精彩.
您正在我內,空虛而不再.
晝夜同行,從來沒有分開.
這樣的愛,傾盡一生的愛.
縱有困苦厄運也不更改.
世間無替代,這情和這愛.
我們傾出一生,報答這份愛.
我們的主,一心比海更深.
這份愛,解我心中罪困.
淒清灶,伴我同行.
孤單中,伴我前行.
這份情,仍是最真.
我們的主,仍高比天更高.

$^{1361}$這份愛,今我竟可遇到.
漆黑中,沒有迷途.
憂谷中,是我路途.
這份情,難以盡收.
漆黑中,沒有迷途.
憂谷中,是我路途.
這份情,難以盡收.
請坐,可以回到座位.
謝謝大家.
\newpage



\section{詩篇 100:1-5}
\label{sec:6_IYA7gmRpA}
\textbf{2023.12.31 《謝主洪恩》 詩篇 100\_1-5 (和修版) 講員 : 高淑芬傳道}
\newline
\newline
連結: \href{https://youtube.com/watch?v=6_IYA7gmRpA}{\texttt{https://youtube.com/watch?v=6\_IYA7gmRpA}} ~~~~ 語音日期: 2023-12-31
\newline
\newline
\hyperref[sec:ZLhqPiukKhM]{\small{< < < PREV SERMON < < <}}
~
\hyperref[sec:index_chronic]{\small{[返順時目]}}
~
\hyperref[sec:index_scriptual]{\small{[返順卷目]}}
~
\hyperref[sec:r5CRKC0_eBw]{\small{> > > NEXT SERMON > > >}}
\newline
\newline
詩篇 100:1-5
\newline
\begin{longtable}{cl}
\hline
\hline
章節 & 經文 (和合本修訂版)\\
\hline
100:1 & \begin{tabularx}{0.7\textwidth}{X} 普天下當向耶和華歡呼！ \end{tabularx} \\ \\ \relax
100:2 & \begin{tabularx}{0.7\textwidth}{X} 當樂意事奉耶和華，當歡唱來到他面前！ \end{tabularx} \\ \\ \relax
100:3 & \begin{tabularx}{0.7\textwidth}{X} 當認識耶和華是神！我們是他造的，也是屬他的；我們是他的民，是他草場的羊。 \end{tabularx} \\ \\ \relax
100:4 & \begin{tabularx}{0.7\textwidth}{X} 當稱謝進入他的門，當讚美進入他的院。當感謝他，稱頌他的名！ \end{tabularx} \\ \\ \relax
100:5 & \begin{tabularx}{0.7\textwidth}{X} 因為耶和華本為善；他的慈愛存到永遠，他的信實直到萬代。 \end{tabularx} \\ \\
[1ex]
\hline
\hline
\end{longtable}
$^{1}$(手挽手當炒).
因為今天穿了太裝不講泰文好像有點差.
今天的主題是「謝主雄恩」.
其實當我接收到雄恩堂邀請我今天去講的時候.
我已經覺得真的要謝主隆恩.
為什麼呢?還記不記得萬事勝義?.
我在穆會七年,宣教二十年裡.
從來沒有試過在一個年度裡.
從頭可以講到尾.
這個在事奉裡我覺得是一個很榮幸.
很慎很恩代的感覺.
但是在這段時間教會發生的事情都不是很難過.
我自己家裡也發生了一些事情.
我就在想,上帝我還是不要講這個題目了.
我轉大,但是還是來不及.
到了星期四,問我講題是什麼?信息在哪裡?.
我再問神,我能不能講出我能不能做到?.
我不是只能講出我能不能做到.
其實一直想著轉大.
想著轉大,轉到.
如果這段時間講一些憂傷的題目可能會比較容易.
但是上帝讓我定經.
我喜歡自己說,定了經文就不能走.
告訴Joann,就是這個經文題目的時候.
我求神給我力量,繼續講好這件事.
並且做得好.
今天領唱的大英姐妹選的歌曲.
真的很應景.
我們值得拿著程序和歌詞.
在今天最後的晚上.
我們可以好好思想到上帝在我們人生裡有很多很多的恩典.
不單是2023年.
是多到數不盡的恩典.
我一直凝注紅水反藍.
紅字不是這樣寫的.
是另外一個字.
但是我查字典,借助紅因,真的用了紅字.
不知道福曲部是否對中文字有更多的了解.
我又感受到上帝很恩代我們紅因堂.
我們正正是紅因.

$^{41}$紅的意思是大過大,乖乖的.
是非常大的恩典.
我們今天坐在紅因堂裡敬拜神.
這是非常感恩的事情.
今天我講的這一篇100篇.
其實由第93篇到100篇.
是一個群組的詩篇.
今天我們很喜歡有什麼群組.
這是一個群組.
總共有8篇詩篇.
前面的7篇是講.
這位耶和華的意思是什麼.
他願意跟我們立約的神.
他願意跟我們做朋友.
願意,其實上帝是高高在上.
但是耶和華的意思是.
他願意跟我們好像平等一樣.
如果我是一個主.
如果這裡還是一個僕人的話.
我不用跟你簽立約.
因為我比祂高.
我的命令就是代表一切.
但是上帝在這一篇.
他講到他是耶和華的王.
他是願意跟我們一起.
其實他的位分比我們高很多.
但是他願意跟我們立約.
納的約是很多.
在聖經裡納了很多很多.
有恩典的約,有應許的約.
到前面的7篇.
7字大家知道在聖經裡代表什麼嗎?.
很完全.
上帝是一個很完全很完全的王.
所以詩人在100篇裡.
再做一個總結.
總結前面的7篇.
來撰寫這100篇.
這一篇是最簡單的.
五節的經文.

$^{81}$要分兩個大段.
就是一到三就是一個大段.
四到五又是一個大段.
兩個的結構都是很相似的.
去頌讚神.
為什麼呢?.
接下來就告訴你為什麼.
兩段的結構都是這樣.
在這100篇裡.
我們也叫做「清謝詩」.
也是頌讚的詩.
所以怪不得保羅在監獄裡.
其實是不容易的.
但是在腓立比書中.
我們看到他要叫我們.
不是呼籲.
是命令我們凡事謝恩.
在詩篇100篇裡.
雖然沒有說到凡事.
但是裡面的意思是.
當我們一息尚存的時候.
其實感恩,頌讚上帝.
是我們離不開的.
我們唯一可以回應給神的回應.
這首歌更加特別.
你知道詩篇全部都是歌.
只差我們現在唱不出.
這首詩歌是以色列人.
當他進入聖殿的時候就去唱.
但是我發現.
這五節裡.
非常有前瞻性.
特別是第一節.
普天下當向耶和華歡呼.
我們先說普天下.
為什麼?.
這首詩歌明明說到.
以色列人去哪裡?.
去哪裡的時候要唱?.
進入聖殿的時候唱.

$^{121}$應該說是我們.
我們應該向耶和華歡呼.
但這句詩為什麼不是用我們?.
是用普天下.
這不單是中文.
中文看到這個字.
看到這個原文.
是用了普天下.
普天下這個.
我們剛剛過了聖誕節.
聖誕節有首歌叫做什麼?.
普天同慶.
當我們聖誕節.
你知道整個世界.
普天就是整個世界.
全部的人.
不管你信了耶穌還是不信耶穌.
你認識上帝有多深.
全部的人.
我們中文最好.
全部人.
全部人應該向耶和華歡呼.
普天下就是全天下的詩.
包羅一切.
很豐富.
沒有分泌到你是泰國.
我是香港.
我是英國.
我是美國.
沒有分泌到你是黑,白,紅,黃.
沒有分泌到你的皮膚的膚色.
也沒有分泌到你的地域.
也沒有分泌到你的來自.
就算泰國也好.
我們也會分泌到你是南部的.
你是北部的.
你是東北部的.
我們也會分泌到.
但是去到這裡的時候.
上帝物是這位詩人.

$^{161}$當他寫的時候.
他完全沒有分泌.
你知道猶太人對這件事很著重.
去到新約.
我們還看到這些傳統.
或者這種模式.
那些是外邦人.
在以色列人看我們.
我們是外邦人.
如果在聖經裡.
我們是野橄欖.
根本是不入流的.
但是竟然在這裡.
完全不分你和我.
是普天下.
看到嗎?.
當我們是一個外邦人.
這個人都可以來向神歡呼的時候.
這已經是恩典.
就好像現在連演唱會去撲飛.
透過演唱會.
你可以買到很貴的票.
有黃牛票.
或者你被人欺騙.
為什麼?.
你看現在什麼都可以WhatsApp.
什麼都可以電話.
為什麼不在電話裡看呢?.
因為去到那個場地.
去看到你所愛的在唱歌.
那個靈性是很不同的.
我現在沒有去看了.
以前我也很喜歡葉德嫻.
我去聽她的演唱會.
不知道你們喜不喜歡她.
我去聽演唱會.
真的很開心.
今天來到上帝的面前.
我們是一起去的.
是那種歡喜快樂的.

$^{201}$是普天下的有我份.
這個有我份.
是一個很深很深的.
不單只是表面字面的意義.
是一個屬靈深處裡.
來了一份很大很大的意義.
特別是今天.
我們再去看以色列的時候.
我告訴你.
連以色列的總理.
他都還沒有相信耶穌.
他都不可能忘記.
他要靠著耶穌得救.
他只知道他是神所揀選的子民.
所以你看不看到.
我們這個外邦人.
比起現在有很多很多.
神所揀選的子民.
是來得更加有為份.
所以我們覺得.
詩人寫到這句的時候.
是一早這麼多年前.
都預了你和我.
是非常有前瞻性的.
大家要留意.
接著又有一個七.
哪個字出現了?.
七次.
泰文的聖經就是「仲」.
「仲」的面很接近.
「當」.
「當當當當當當當」.
七個音.
七個音很完全.
因為音樂的表達.
也是一個完整的音.
接著你會加高低.
發展出很多很多美妙的.
敬拜師很多的音樂.
在這一書.

$^{241}$你不要小看這個「當」.
是一個命令的動詞.
不是靜止的一個名詞.
看到就算了.
是要做的.
你做不到.
今天學一點.
明天再學一點.
一路去學.
是終身學習的.
怎樣繼續再做好一點.
雖然我的名字叫高勳.
但去到高勳還不夠.
要去到上帝給我們的分數.
去到上帝給我們那種.
合乎神旨意的分數.
第一個環節在第一節裡.
普天下當向耶和華歡呼.
這句說到「當歡呼」.
你知不知道什麼是歡呼?.
你可不可以示範一次?.
我說一二三.
你們就示範歡呼.
一二三.
(歡呼).
你們真的很有前執.
因為我們紅恩堂真的不同.
你們知道什麼是歡呼.
有一個師兄.
他試過叫弟兄姊妹們歡呼.
他們就大聲歡呼.
我不知道是不是開玩笑.
他就說你們再來一次歡呼.
你們太斯文了.
再歡呼.
你們剛才是讀歡呼.
已經提示了他們.
他們仍然大聲一點.
歡呼.
歡呼是由心裡面研發的.

$^{281}$你不會再理會環境.
我是傳道.
我是一盞.
我現在看足球.
安靜一點.
我其實沒有看足球.
但我在泰國學泰文的時候.
我的同學是很多國家都有的.
其中.
你想想已經是二十年前.
其中有一年.
在韓國做世界盃.
總之不知道什麼盃.
韓國的同學非常.
「來啊,去看球」.
我人生沒看過球.
我不懂得看.
「你來吧」.
「這次是在韓國打」.
「現在是韓國隊出場」.
跟著他去.
才知道原來泰國裡面.
除了有華人街外.
也有韓國街.
有一個商場裡面.
一二三四五.
全部店舖都是賣韓國的東西.
連韓國的石碗.
金子.
也有很多碗.
原來在那個位置.
你就可以買到很多韓國的東西.
那個商場裡面.
就弄了一個很大的銀幕.
大家一起看.
我如果沒有印象錯.
他不知道是不是叫我們穿紅色衣服.
還是全場都是紅色.
總之.
他一進球.

$^{321}$如果我們弟兄看球.
或者姐妹看球.
你就知道會怎樣.
我也不懂.
「嘩」.
真的很開心.
我想問.
他們這麼開心.
他們有沒有用金收.
沒有的.
但他們就算韓國隊贏了.
政府會不會派給他們.
一人一萬元.
一萬韓元.
也沒有的.
什麼也沒有的.
但他們就是很開心.
自己的國家隊.
踢了一球進去.
他們就叫了很開心.
又打擊一下股.
總之很開心.
不過我看到電視裡面的人.
也很開心.
我們拍到那些啦啦隊.
你知道嗎?.
那種歡呼就是這樣.
如果相比我們.
平時我們說.
我們去敬拜向神歡呼.
其實我們剛才的詩歌.
也去向神歡呼.
我們就很斯文.
但我們心裡面很澎湃.
當我們見到上帝的恩典.
為什麼剛才我說.
你再看看這些詩歌.
再回去.
手手的詩歌也說到.
上帝在過去.

$^{361}$給了我們很多很多的恩典.
我想你在心裡面.
向神發出這個.
嘩!上帝真是勁XX.
記不記得什麼叫勁XX.
嘩!很厲害的神.
用你的語言去歡呼.
這個歡呼這個字.
再加上普天下.
那個氣勢.
是非常配合.
而且是非常強大的.
這個歡呼這個字.
在舊約裡面.
我們說要用吹號角.
來歌頌神.
要吹大聲.
那個字是一模一樣的.
就是那種向神發出的歡呼.
甚至乎.
這種歡呼是喜樂的.
是得勝的.
就好像約書亞帶領著那班以色列人.
去攻陷那班敵人的時候.
特別是耶利哥城.
城都未塌下來.
最後一下.
他們走了多少個圈?.
走了多少個圈?.
沒有人回答.
只是做手勢.
你回去看吧.
他們走了七個圈.
第七個圈.
他們一叫就塌了.
那種歡呼其實還未得勝.
但他們知道上帝應他們得勝.
一歡呼就塌下來.
你看到嗎?.
這個讓我看到.

$^{401}$當我們覺得現在很不容易的時候.
我們真的要看著.
這些聖經裡面這麼多人.
這麼多的事.
來到學校.
雖然我媽媽進了醫院.
三個星期了.
還未痊癒.
但我相信上帝一定有祂的美意.
這個要學習的.
第二個當.
當樂而時逢耶和華.
中文就是這樣.
可能中文的翻譯.
怕我們知道.
時逢根本不知道搞什麼鬼.
是否要負責任.
是否死死地氣.
是否很膚淺.
是否因為我是這個角色.
我要好好去做.
不是這麼簡單.
是傳樂意的心.
時逢.
當我們說今天.
如果我們主日崇拜.
英文是Sunday Service.
Service就很好.
外國人就很好.
經常提醒.
敬拜其實就是時逢.
敬拜和時逢.
在原文裡是一個字.
今天我們有招待.
我們有師事.
我們有師班員.
我們有師琴.
我們有IT部.
我們有主席.
每一個都有崗位的時逢.

$^{441}$或者之前我們有主日學的學生.
我們有主日學的老師.
他們在做崗位.
這個是時逢.
但我告訴你.
你沒有崗位.
你坐在這裡.
誠心地來敬拜神.
同樣是時逢神.
這個時逢.
不知道你是否甘心樂意.
今天還未到九度.
很溫暖.
是否難阻了你.
很感恩你已經回來了.
這個時逢.
時逢就算是耕田,耕種.
都是這個字.
時逢就是做一些事.
就像管理園裡的一切.
管理,做修剪的事.
這個都叫做時逢.
所以這個時逢的態度.
是甘心樂意的.
並且是自己願意的.
是活祭的.
《新約聖經》裡的羅馬書.
說到我們要像活祭.
活祭的意思是.
聖經裡沒有特別說要綁你.
因為死祭都不用綁.
活祭不會綁.
你隨時都可以跳下來.
我們人就是這麼壞.
有時候都可以跳下來.
但是很合乎上帝的心意.
獻上是一個活祭.
你願意去時逢.
有時候我們去時逢那些朋友.
真的很不容易.

$^{481}$有時候你去時逢群體,組員.
你真的覺得很不容易.
時逢就沒有了「憑」字.
就是時逢那些「友」.
真的不容易.
但是為什麼我要時逢?.
為什麼這麼不容易?.
這麼受氣.
都要去做.
我想我們當中有份去幫忙時逢.
大部分都是無身的.
都是樂意去時逢.
但是又受氣.
但又不容易.
那群「友」又很難搞.
為什麼呢?.
是因為我們信服神.
我們願意去時逢上帝.
我們知道今天是.
表面上是去時逢這群「友」.
但是裡面的核心.
我是在時逢上帝.
所以今天每一個時逢的人.
不要看這群「友」.
當上帝擺在你的崗位時.
求主給我們願意的心.
因為我信服上帝.
我願意起來.
為上帝做一些事.
你不做不要緊.
上帝還可以派其他人.
沒有一個崗位.
不會說沒有了你不行.
還會有很多很多的人.
但是要留意.
我們有份可以去時逢神.
並且樂意去時逢耶和華時.
這正是我們每一個感恩時.
感恩而帶出來的一件事.
第三個「當」.

$^{521}$就是來自第二節.
「當歡暢來到神的面前」.
我講過經文.
「來」.
還記得「來」的泰文是什麼嗎?.
泰語姐妹.
「麻」.
「麻」是很重要的.
你是願意來到上帝的面前.
每一個時逢.
你都要去到上帝的面前.
我的時逢不是站在那裡表演.
站在那裡炫耀.
不是的.
要誇的就是誇我的上帝.
我的上帝有.
我才可以去時逢神.
是要「來」.
來到上帝的面前.
當然「歡暢」.
在原文裡的「來」.
你來得就是自願.
在原文裡的「來」.
已經包含了這件事.
但是在中文裡.
因為有「歡暢」的意思.
所以在中文譯本裡.
「來」不是逼你來.
是你歡喜快樂.
來到上帝的面前.
來到祂的面前.
來到祂的時逢.
去讚美祂.
去感恩.
第四個「當」就是「曉得」.
越是去稱頌上帝.
越是去事奉上帝的時候.
其實在原文裡.
就是「曉得」這個字.
我們中文聖經或泰文聖經.

$^{561}$都很好.
七個「頌」字裡都有說到.
去認識上帝.
認識就是「曉得」.
你「曉唔曉」.
中文很好.
你「曉唔曉」.
不單是認識這麼簡單.
我當然是認識.
這張叫「紙」.
這支叫「筆」.
但問題是你「曉唔曉得」寫.
寫之中.
你「曉唔曉得」寫主禱文這三個字.
你「曉唔曉」.
是你要從經驗裡體會.
體會上帝是真真實實的上帝.
當「曉得」這個字.
在原文裡.
你會發現.
當夏娃和阿當去同房的時候.
他們都是用這個字.
是彼此很深很深的認識.
我們對上帝有沒有這份認識呢?.
如果有的.
繼續.
因為這個認識.
當我們還有一息尚存的時候.
仍然是未停止的.
然後進入.
進入他的門.
進入.
即是去.
去的進入.
去又記不記得了?.
去.
閉.
泰文就是閉.
這樣.
進入是你要走去.

$^{601}$出到你家.
由你的安書區進入.
今天要呼籲.
我們仍然是在網上敬拜神.
真的要呼籲你.
去進入上帝的殿.
今天的你.
你知不知道.
在屬靈層面裡.
你是進入了上帝的門.
而去敬拜上帝.
在屬靈的意義是深一層的.
你就好像去到演唱會.
在現場.
現場的氣氛.
在現場你感受到.
聖靈的工作.
你會感受到.
當然我們現在看電話.
在電腦裡.
我們看到很多的敬拜.
我們可以參與很多不同的崇拜.
我們可以聽到很多的好道.
這個對你也是有幫助.
但是.
每一個星期.
回到你屬靈的家.
進入上帝的殿.
這個對你.
在靈靈裡是有祝福的.
所以這份祝福.
我們經歷過那幾段時間.
我們不可以回來一起敬拜神.
你的感覺是怎樣?.
會比現在好嗎?.
你自己評分吧.
我喜歡現場的敬拜神.
進入上帝的門.
而且這個門.
是很有意思的.

$^{641}$上帝是開了這個門.
你知道嗎?.
你和我都是一個罪人.
我們根本就進不了這個門檻.
我們的家長不要輸在起跑線.
不要輸在起跑線.
每一個家長.
帶著肚子也好.
已經去問門檻了.
只會問住院的門檻.
怎樣進去?.
有沒有門進去?.
有沒有門路?.
看到嗎?.
上帝開了這個門.
給我們每一個人.
所以今天教會的門是開的.
來的人.
是可以普天之下的人.
問題是我們今天已經坐在這裡.
洪恩家裡的人.
我們會否接納到這個人來了?.
這個人又來了.
求神也撫慰我們的心.
讓我們進入上帝的聖殿.
來知道這個門是寶貴的.
上帝已經為你而開了.
當札美進入他的院.
不是你可以去看聖經參考的.
是你要發自你的內心.
發自你的內心之餘.
你是會越來越深的.
去深藏著上帝給你的恩典.
這個在我們每一個人.
我們說基督徒最重要.
你有一本聖經.
和一本人的書.
這個人的書.
你這一本書.
寫得怎麼樣呢?.

$^{681}$裡面多不多著.
寫到很多上帝對你的恩典呢?.
由救贖之恩.
開始的時候.
是上帝怎樣來呼召你呢?.
我前幾天有一個弟兄.
也是我母會的執事.
以前我進去教會的時候.
他是團契的團長.
突然間問我一句.
你記得誰帶領你順主嗎?.
當然不是他.
我說記得是誰誰.
他說不是.
不是?是他們.
沒有錯的.
原來他最近跟了我一個小學同學.
一年級的小學同學.
他還沒有揭曉的時候.
不是他揭曉了.
我記得了.
三個字.
三個字說了出來.
哦,是呀.
他是第一個帶我回教會的.
當時是一年級.
你想不想到一年級的同學.
竟然可以邀請另一個同學回教會.
當時我住在荃灣.
現在是地鐵.
以前我們住在石圍角.
不是現在的石圍角.
在老圍村再上一點.
再嚴格一點再細緻一點.
就是石圍角二波鎮.
我這個同學住在山下.
我們叫山中間.
如果老圍村是山上的話.
我們叫中間.
他星期六已經告訴我.

$^{721}$明天要回教會.
好呀好呀.
誰知我那個星期日.
不知道為什麼.
這個一年級的事情.
到現在都不知道是多少年了.
但我記得.
那個一年級的事情.
就是那天我仍然像一般的小學同學.
星期日.
做什麼?.
早上睡覺.
還在睡覺.
直到這個小學同學來到家.
有人來找高淑芬.
喂!有人找你呀.
你快點起來.
我才起來.
這個小學同學看著我梳洗完.
穿上外出的衣服.
和她去教會.
很感恩.
這個小學同學牽起了我.
由小學一年級.
到六年級.
我都有信義會天恩堂.
我們讀信義小學.
現在又搬到象山村.
是很感恩的.
她一說起.
燈燈燈燈.
我便說了這個故事.
我問她為什麼跟著她.
因為這個家庭.
當我當時還是傳道人.
我見到伯母.
因為有一年.
嘉利大廈大火.
燒死了很多人.
因為通道走火.

$^{761}$走不了.
已經有二十多年前了.
我的同學妹妹.
在那場火災裡燒死了.
這個家庭.
不知道是誰介紹.
因為後來.
那裡被拆了.
現在荃灣地鐵.
大家都各散東西.
我都沒有見到她.
那一次見到.
伯母這麼面熟.
是否有什麼態度.
她說是.
我說我是阿公淑芬.
見到我同學.
見到她弟弟.
原來她弟弟和伯母.
後來都有回教會.
唯獨是我的同學.
只顧著賺錢.
到現在.
我經常陪伴著.
所以我現在的執事.
跟著她.
所以我覺得.
很感恩.
一個傳給別人的人.
自己失落了.
所以今天我們坐在這裡.
求神祝福你.
上星期你帶了很多朋友來.
你自己千萬不要失落.
你還是要繼續帶朋友來.
但自己要站立得穩.
怎樣站立得穩.
人生有很多不容易的事.
夾在你心裡.
夾用了泰文.

$^{801}$收藏在你心裡的夾萬.
一定要收起來.
因為我們很容易會忘記.
你再重新看的時候.
你會發現神的恩典一直都在.
最後就是當稱頌.
稱頌你.
很多時候我們在聚會祈禱裡.
分得很細節.
要去讚美神.
在讚美的環節.
很多人只說兩三句就停了.
有些背經文.
有些又多幾句.
你可以這樣學習讚美神.
但是稱頌神.
是你一件一件的恩典.
經歷你自己已經經歷.
你就是過來人.
你就是本體經歷了這些恩典的時候.
你越經歷上帝很豐富的恩情的時候.
你可以用你的言語去稱讚神.
例如上帝真的是一個.
說來就來的神.
說得就得的神.
不成功就成功的神.
關關很難過.
但是神就給你一個關去過的神.
你會更加更加懂得.
頌讚這位上帝.
正如我剛才所說.
這篇詩篇七個當都說了.
我們有沒有聽到神的命令.
來履行這麼完全的感恩.
頌讚讚美神.
求主幫助我們.
第一段和第二段.
祂呼籲你們.
就是剛才那七個當.
來去頌讚神的時候.

$^{841}$祂說了很多很多理由.
例如第三節.
我們是祂做的.
在聖經裡面.
我們是祂做的.
因為祂是上帝.
我們是上帝所做的.
自己做的事是很寶貴的.
所以我會很珍惜.
有個姐妹送了一條頸巾.
所以我就抱著她.
因為覺得親手做的很珍貴.
好像我們上次收到.
親手做的曲奇餅.
寬鬆的.
那件餅如果不吃.
就發霉.
所以不捨得.
也要拿來吃.
是親手做的.
還有上星期的愛言頌巴閉.
是怎樣的.
弟兄姊妹親手做的.
親手做的東西.
和外面茶餐廳的東西.
是不同的.
所以如果你請我吃東西.
最好請我去你家吃家常小菜.
不用去酒樓.
因為這些是值得更加欣賞.
我自己去理解.
我自己做的事.
我都會珍惜.
亂寫就慘了.
這件事是自己做的.
那件事是自己寫的.
但是一本本.
不扔掉.
就沒有了位置.
扔掉又很可惜.

$^{881}$於是下面的雜誌.
雖然如果再重覆.
幾年.
不知道會不會升值.
有些雜誌會升值.
但想想.
既然不是我做.
現在還未升值.
你會先刪走那些.
你那些肯定100年都不會升值.
但你寧願留住.
你自己手上所做的事.
你想想.
你啊我啊.
是神所做的.
是做你的神.
在第三席裡.
我們是神做的.
還要屬於祂.
我做了之後.
我送給梁傳道.
就屬於梁傳道.
他不喜歡扔垃圾桶.
我無權過問.
因為屬於梁傳道.
但屬於我的.
就很不同.
你想想.
你做的事.
這些事還屬於我.
我會珍愛有加.
上帝對你.
就是這樣.
珍愛有加.
更加厲害.
第三節.
我們是.
祂的民.
現在的難民.
大家知道.

$^{921}$他去到法國.
能否入境.
對不起.
貨物是入境.
能否入境.
我去到土耳其.
能否入境.
等於以前越南難民.
來到香港能否入境.
香港境.
在這裡.
我們是屬於.
耶和華神.
耶和華的國在哪裡.
叫甚麼國.
天國.
這個天國.
的皇帝.
我們是他的子民.
大家知道.
現在有多少人.
去拿英國護照.
不是英國的.
是加拿大的居留證.
都要幾年.
在這裡.
或者在台灣.
有規條在這裡.
在這裡.
我們是他的民.
是好心好心的意思.
我們是.
有訂酒的.
我們去到最終的一天.
我們是去到.
上帝的國裡面.
這句更厲害.
是他草場的羊.
一片的草場.
我們是羊.

$^{961}$這片草場的草.
Oh, you can eat.
像不像吃自助餐.
你喜歡.
整包.
好像自助餐.
不單是.
吃.
全部是軟綿綿的草地.
你可以怎樣.
在這裡.
休息.
睡覺.
歇一會.
你看到.
我們的皇帝.
我們的主.
對我們好不好.
Di ma ma.
ma ma 當然是非常.
Di 就是好.
非常好.
上帝是非常好的民.
第三節解釋到這件事.
到最後的時候.
他就說到.
第五節.
他的慈愛.
是傳到永遠.
這份慈愛.
你看看這篇詩篇.
不是早於主耶穌降世的時候.
當我們唱聖誕節.
當時還沒有聖誕歌.
還沒有聖誕節的時候.
上帝已經.
給我們這麼大的應許.
每一個.
慈愛由那一刻.
上帝創造我們.

$^{1001}$到現在是傳到永遠.
並且他的信實.
直到萬代.
今天.
我們香港的.
銀行體制.
都信得過.
否則我們全錢.
傳進去.
會不會一下子.
沒有了.
哥倫比亞或阿根廷.
有些國家.
本來是中產.
階層的人民.
一夜之間.
他在銀行的錢.
完全歸零.
全部零擔.
沒有了價值.
一夜之間.
政權一改完全沒有了價值.
這些國家.
這樣的情況有.
或者我們上一輩.
他們說擔住的錢.
馬上變了.
以前一張張.
現在一擔擔都不值錢.
那些糧票.
以前很值錢.
去到中國.
生活的時候.
糧票很重要.
但是上帝的信實.
直到萬代.
意思是.
這是最信心的.
保障.
所以我對自己說.

$^{1041}$神是很看顧我的.
我媽媽也是屬於.
祂的.
上帝會繼續帶領.
其餘一個人.
我擔心的就是我妹妹.
因為我妹妹.
媽媽還未出事.
於是要看醫生.
很多事情.
發現腦內有個水瘤.
但未去是良性的.
很多事情她已經睡不著.
我問媽媽現在入院.
她試過呼吸不周.
壓力這麼大.
她也有看中醫.
看醫生.
她這樣.
每天探訪媽媽.
我經常叫她.
要不要回去休息.
她說明天一定不能來.
今天我一定不能去.
探望我媽媽.
她說明天休息.
她也來了.
我們三個.
畢竟在床位.
三個星期.
我看到一幅.
很美麗的圖畫.
因為在醫院.
我只會.
動來動去.
所以我看到人們.
坐在床上.
我心想為何我們這麼忙.
不是摸媽媽的肩膀.
而是摸她的腳.

$^{1081}$做茶的事情.
清潔的事情.
所以我沒有跟別人.
有兩句.
我專注在媽媽身上.
我們一開始.
也不懂得幫媽媽.
因為她要活插為喉.
不准飲食.
九天內.
不知如何處理.
有床隔壁.
教我們可以用這些.
來清潔她的口.
其實我沒有問她.
因為我媽媽的位置.
有近洗手的地方.
她在洗手.
我看著她.
她說不是用這些來教我.
我說謝謝.
我只告訴我妹妹.
那個床隔很好.
她教我這個那個.
我妹妹.
跟別人聊天.
在那三個星期.
也會聊.
她主要是到處聊.
其實平時是我.
但在這個場景.
變了.
不幸地.
教我比較多的床位.
那位伯母.
過身.
我妹妹竟然問我.
你會不會去跟別人祈禱.
我說.
我還未做完.

$^{1121}$我媽媽.
我沒有兩句.
你跟別人聊的句數.
比我多.
你去吧.
這個也不是很好的見證.
高傳道.
去到病房.
我只尊顧我媽.
沒有其他.
我妹妹.
已經呼吸不了.
她竟然.
自發地.
我看在眼內.
她竟然自發地.
這裡要清潔.
封了.
家人在外.
她出去跟家人祈禱.
我祝福她.
你去吧.
我問她.
信主還未信.
她說她願意跟我一起祈禱.
我覺得是一幅.
很美麗的圖畫.
高傳道.
沒有傳道.
是我妹妹傳.
所以.
當我看到這些美麗的圖畫.
的時候.
我覺得上帝是看著的.
而且.
我又要很理智地.
告訴自己.
新老病死.
我們每一個都會經歷.
不是說我們信了耶穌.

$^{1161}$不會經歷.
我記得當時.
探癌症醫院.
在泰國.
六歲.
那星期去.
很開心地跟她玩.
鼓勵她.
因為全院最小的六歲.
沒多久.
那星期去.
哇!這麼小.
這麼小.
這麼小就走了.
希望她聽到.
真的上天.
但我們看到.
很多是突然去.
有些是病慢慢去.
很多種.
總之生離死別.
亦都是人生.
這就是人生.
我好像傳道書裡.
生有時死有時.
在這個刻.
就算我們有很多不容易的事.
哇!倘若沒有神.
我想會更慘.
更不知如何是好.
所以.
神是信實.
直到萬代.
所以我們願意我們洪恩家.
不單止萬事.
是照著.
上帝的意旨.
而且是.
時刻的.
謝主洪恩.

$^{1201}$願神祝福各位.
\newpage



\section{撒母耳記上 15:17-23}
\label{sec:r5CRKC0_eBw}
\textbf{2024.01.07 《作主聽命的兒女》 撒母耳記上 15\_17-23 (和修版) 曾敏傳道}
\newline
\newline
連結: \href{https://youtube.com/watch?v=r5CRKC0_eBw}{\texttt{https://youtube.com/watch?v=r5CRKC0\_eBw}} ~~~~ 語音日期: 2024-01-11
\newline
\newline
\hyperref[sec:6_IYA7gmRpA]{\small{< < < PREV SERMON < < <}}
~
\hyperref[sec:index_chronic]{\small{[返順時目]}}
~
\hyperref[sec:index_scriptual]{\small{[返順卷目]}}
~
\hyperref[sec:wCnaB7fTUZk]{\small{> > > NEXT SERMON > > >}}
\newline
\newline
撒母耳記上 15:17-23
\newline
\begin{longtable}{cl}
\hline
\hline
章節 & 經文 (和合本修訂版)\\
\hline
15:17 & \begin{tabularx}{0.7\textwidth}{X} 撒母耳說：「你雖然看自己為小，你豈不是作了以色列諸支派的元首嗎？耶和華膏你作了以色列的王。 \end{tabularx} \\ \\ \relax
15:18 & \begin{tabularx}{0.7\textwidth}{X} 耶和華差遣你，吩咐你說：『你去除滅那些犯罪的亞瑪力人，攻打他們，直到把他們完全滅盡。』 \end{tabularx} \\ \\ \relax
15:19 & \begin{tabularx}{0.7\textwidth}{X} 你為何沒有聽從耶和華的話呢？你為何急著撲向掠物，行耶和華眼中看為惡的事呢？」 \end{tabularx} \\ \\ \relax
15:20 & \begin{tabularx}{0.7\textwidth}{X} 掃羅對撒母耳說：「我聽從了耶和華的話，行了耶和華派我行的路，擒了亞瑪力王亞甲來，滅盡了亞瑪力人。 \end{tabularx} \\ \\ \relax
15:21 & \begin{tabularx}{0.7\textwidth}{X} 百姓卻從掠物中取了牛羊，是當滅之物中最好的，要在吉甲獻給耶和華－你的神。」 \end{tabularx} \\ \\ \relax
15:22 & \begin{tabularx}{0.7\textwidth}{X} 撒母耳說：「耶和華喜愛燔祭和祭物，豈如喜愛人聽從他的話呢？看哪，聽命勝於獻祭，順從勝於公羊的脂肪。 \end{tabularx} \\ \\ \relax
15:23 & \begin{tabularx}{0.7\textwidth}{X} 悖逆與占卜的罪相等，頑梗與拜偶像的罪孽相同。因為你厭棄耶和華的命令，耶和華也厭棄你作王。」 \end{tabularx} \\ \\
[1ex]
\hline
\hline
\end{longtable}
$^{1}$丁姐妹平安.
我們在新的一年.
仍然有福氣回到神的家.
敬拜和聆聽祂的話.
我們再一同禱告.
讓我們有何等的福氣.
我們在這一生可以認識你.
藉著你的愛子耶穌基督.
救贖我們的生命.
讓我們離開永遠的死亡.
今天我們在屬於上帝的群體裡生活.
是何等的有福.
是何等的榮耀.
主啊,求你讓我們感恩的心.
去過每天的生活.
在這一刻,願你讓我們每一個的心靈.
完全的向你打開.
讓我們聽到主你自己的聲音.
讓我們的心被觸動.
就要去回應你.
願以聖靈對我們眾人說話.
奉耶穌基督的名字祈禱.
阿們.
其實今天我要說的這段經文.
遠遠是多過剛才大家所讀的那段經文.
是非常豐富的.
因為我們都需要上文下理.
知道之前發生什麼事.
為什麼去到剛才那一道呢.
或者之後的後續是怎樣.
待會有些經文是我們剛才沒有讀過.
大家聽我讀就行了.
今天很想說.
作主聽命的兒女.
在一年新的開始的時候.
很想成為我和你的一個很大的鼓勵.
一個祝福.
希望我們都是這樣.
新一年的願望就是作主聽命的兒女.
這段經文我們去看的時候.

$^{41}$就會看到訴羅這個人物.
我們看看.
在這裡有一個非常重要的命令.
在剛才那段經文之前.
其實是一個很重要的命令.
是什麼呢?.
撒母耳對訴羅說話.
訴羅是以色列第一任的君王.
上帝揀選他.
揀選他做以色列人的君王.
在過程中.
臣透過撒母耳.
讓訴羅知道一件事是很重要的.
這件事是所有事情之先的.
第一件事不是要吩咐他.
第一件事是要管理以色列民.
在管理上要做ABCDE.
然後哪件事是怎樣的安排.
不是說這些安排.
也不是說那些事情實際是怎樣.
最重要的.
今天一定要離開這裡.
大家要記住這句話.
就是要聽從耶和華的話.
這是撒母耳對訴羅說的.
很重要.
你什麼都不記得.
你離開的時候就要記得.
聽從耶和華的話.
是所有事情之先.
聽從耶和華的話.
第一個很重要的命令是.
原來撒母耳要對訴羅說.
你要去攻打亞馬利人.
要去打掌神的百姓.
我們的上帝不是很渴望和平嗎.
是很溫柔的.
喜歡和平的.
第一件事就是要去打仗.
我有沒有聽錯?.

$^{81}$沒有聽錯.
是要去打仗.
要攻打亞馬利人.
但是這裡很清晰.
很肯定一件事.
為什麼要去打這場仗.
這是來自上帝的審判.
是上帝對亞馬利人的審判.
在後面.
其實剛才讀經文的時候.
都看到.
在後面.
神提到亞馬利人是犯罪.
在這個民族.
他們是犯罪的亞馬利人.
神提他們是犯罪的.
是得罪自己的.
不是犯一個小的罪.
是一個半個罪.
算了.
你放過他.
不是的.
這個民族一直一直犯罪.
而且他們還做了什麼事呢?.
我們來認識一下.
首先我們認識這個民族是怎樣來的.
原來他們是以素的後裔.
大家還記不記得以素是誰呢?.
有人可以告訴我.
以素是誰嗎?.
大聲說出來.
亞國的什麼?.
哥哥.
他們的爸爸.
以素亞國的爸爸是以撒.
以撒旗下有兩個兒子.
以素是哥哥.
亞國是弟弟.
原來是哥哥的後裔.
以素有一個兒子叫做以利發.

$^{121}$他娶了一個妾叫做婷納.
婷納就生了亞馬利.
亞馬利人是這樣來的.
所以可以這樣說.
亞馬利人和以色列人是什麼關係呢?.
以色列是亞國的後裔.
亞國和神摔跤之後.
神幫他改名叫以色列.
以色列人就是亞國的後裔.
以色列人和亞馬利人的關係是什麼呢?.
大家直接說是什麼關係.
其實是有血緣關係的.
是親屬來的.
親屬都這樣.
是喔.
以素的後人和亞國的後人.
在這個衝突裡面.
其實他們不是完全不相關的.
是有關係的.
是有親屬的關係.
但神卻要素羅帶軍隊去打仗.
為什麼呢?.
在《申命記》25章這樣說.
你要記得你們出埃及的時候.
亞馬利在路上是怎樣對待你的.
在路上迎擊你.
在你被罰昆厥的時候.
擊殺所有在你後面軟弱的人.
並不敬畏神.
大家是兄弟的後人.
是親屬.
有沒有那種親情呢?.
沒有.
完全沒有親情可言.
不單沒有親情.
還要落井下石.
趁你病 趁你命.
趁著你在這種困境裡.
還要再害幾下.
在路上是怎樣呢?.

$^{161}$已經很疲乏.
出埃及的時候.
你知道他們大大小小.
拿著所有家當.
你要知道後面是埃及的追兵在追他們.
所以以色列人有多大壓力.
多疲累.
多害怕 恐懼.
很多種種的.
他們屬於親屬的人.
不但沒有幫助他們.
扶持他們 保護他們.
還要再在你後面軟弱的人.
擊殺.
這些絕對是罪人的所為.
神看著的時候.
到了時候.
上帝會審判.
上帝也會施行公義.
神不會看人犯罪.
當那人沒事.
沒事就算了.
都過去了.
就當沒事發生.
不是的.
去到時候.
神會審判.
神會施行公義.
就是現在這個時候.
上帝要掃羅領兵去打仗.
現在要去攻打亞馬利人.
滅盡他們所有.
這場仗.
神的吩咐是很重要.
我們要聽清楚.
是不是就這樣打仗.
你打贏就可以了.
就放過他們.
他們輸了.
就知道自己得罪上帝了.

$^{201}$不是的.
要留意神的命令是很清晰的.
要滅盡他們所有.
不可以連累他們.
將男女孩童吃奶的.
以及牛羊駱駝和驢.
全部殺死.
所有人連同牲畜.
一個不留.
這是命令的所有.
什麼叫聽從耶和華的命令.
不是說我聽一點不聽一點.
我聽你的話.
我就去打仗.
我叫聽你的話.
你也要稱讚我.
乖乖.
不是這樣的.
整個命令要聽清楚.
整個命令是滅盡所有的.
不可以連累他們.
將男女孩童吃奶的.
以及所有牲畜.
全部殺死.
一個都不能夠整.
這是命令的全部.
蘇洛是要聽全部的.
這樣叫聽命令.
今天我們也問問.
我們人生過去的日子.
我相信上帝對我們人生有很多旨意.
很多神的心意.
神很想跟我們說.
整本聖經.
裡面有很多神的心意.
整本聖經.
我們看了多少聖經.
這本聖經看了多少.
不只是選一句.
選一段.

$^{241}$聽了.
哪一段最重要?.
哪一句最重要?.
只選那一句聽了.
不足夠.
整本.
我聽了多少?.
我知道神對我的旨意有多少?.
你聽了那一句就是那一句.
還要看自己有沒有回應.
聽的時候覺得.
真的很好.
轉眼間忘記了.
聽完等於沒有聽過.
我們一年有很多堂道.
我們總會有些時候.
記得這些,不記得那些.
不要緊,每一堂道你記得一句都很好.
你聽了之後有回應都很好.
一年就很多了.
問問自己一年聽了多少堂道.
有多少實行出來.
問問自己.
這件事很重要.
這個命令之後.
我們看看蘇羅的回應.
蘇羅的回應是怎樣呢?.
我們看看這個圖.
小許細.
蘇羅有聽命令.
他說他去攻打亞馬利人.
從哈菲拉直到.
埃及東邊的樞爾.
一直去打.
他們捉了亞馬利的王雅.
用刀殺盡亞馬利的眾百姓.
我們看看.
左邊的圖是埃及.
中間的是樞爾的曠野.
亞馬利的圖不是十分清晰.

$^{281}$是稍為偏右一點.
蘇羅有聽命令.
做了一跌.
做了什麼?.
去攻打.
一直打,打成仗.
但我們要留意.
這場打成仗是因為什麼?.
蘇羅太好打了.
一個勇猛的戰士.
一個勇猛的帶領者.
他的軍隊太好了.
所以就贏了.
是不是這樣呢?.
並不是.
亞馬利人很能打.
大家記得出埃及記.
有一場仗是怎樣打的.
那場仗.
摩西在山上禱告.
約書亞在地面上帶領部隊.
他所揀選的人.
出去和亞馬利爭戰.
摩西何時舉手.
那場仗是怎樣呢?.
就贏了.
摩西何時舉手禱告.
就贏了.
當上帝痛哉.
上帝祝福的時候.
何時累了.
手一垂低就輸了.
所以亞馬利人不是那麼弱.
不是那麼容易打的.
真的靠上帝去打贏仗.
之前試過.
現在是神將亞馬利人交在.
以色列人的手裡.
所以那場仗很順利.
打贏了.

$^{321}$他們捉了王.
也滅盡了百姓.
但是很快地看到經文.
蘇羅的回應有些什麼.
蘇羅和百姓很憐惜阿鴿.
很憐惜那個王.
我們不知道.
亞馬利人的王.
有什麼令人很憐惜他.
他是不是有什麼魅力.
是不是有什麼迷惑.
令蘇羅和百姓覺得很可惜.
不忍心下手殺你.
不知原因.
聖經沒有說.
總之就留下了阿鴿.
還有些事情又不捨得.
那些好的牛,羊,牛獨,羔羊.
一切的美麗.
太好了.
好像很可惜.
上帝要我們殺光.
很浪費.
這麼好的東西.
留下吧.
怎會這樣就殺光.
只是殺光那些看不上眼的.
沒有價值的.
那些就殺光.
去到這裡大家怎麼看.
蘇羅是否遵守命令全部.
不是.
一看就知道不是.
留有餘地.
黃又留下了.
那些牛,羊又留下了.
好的那些.
做完這些事之後.
神有沒有回應呢.
是吧.

$^{361}$算了.
都打贏仗了.
那就算了.
都滅了那些百姓.
放過他吧.
差一點而已.
不要緊.
做人不要這麼執著.
是不是這樣呢.
不是.
神的命令是全部都要聽的.
回應是百分之一百.
神一直在看蘇羅怎樣回應他.
正如神今天也在看我們.
好了.
神看到了.
這次到神去審判蘇羅.
因為很明顯.
大家都知道.
他沒有聽上帝的命令.
神在的.
神昔日一直看著蘇羅.
看著他的軍隊.
看著整個以色列民.
他當然也看著阿瑪利人.
看著所有的民族.
每一個人都不能逃過上帝的眼光.
今天神同樣看著我們.
神同樣看著我們.
今天我們在崇拜當中.
神也看著我們.
看著我們每一個人.
講道的.
在台下聽道的.
服侍的.
上帝看著我們.
看著我們外面的表現.
看著我們的內心.
上帝是有回應的.
上帝去審判了.

$^{401}$耶和華的說話.
就臨到撒母耳說.
在這裡我納蘇羅為王.
我感到遺憾.
在這裡表達的一個很大的失望.
神對蘇羅是有期望的.
他揀選他為王.
這並不是一件小事.
經文後面有提到.
不是一件小事.
不是每個人都會選擇他為領袖.
特別是這個位置.
是上帝的選民百姓.
是多麼的尊明的群體.
上帝特別去選擇那個領袖.
上帝對那個領袖是有期望的.
但這一刻.
神已經看到.
他不是徹底聽上帝命令.
聽一點不聽一點.
做一點不做一點.
因為他轉去不跟從我.
聽從是甚麼?.
聽和甚麼?.
跟從.
遵從.
走出來.
是不是?.
蘇羅已經轉去了.
不跟從我.
他不遵守我的命令.
神感到很失望.
撒慕爾也很生氣.
很生氣.
他也有份去高辣蘇羅.
他有份去扶持他有多少個禱告.
有多少個扶持.
但這個王這麼快就已經背離神了.
這麼快就已經不遵守神的命令.
所以神是很難過.

$^{441}$神感到遺憾和失望.
然後神下一步就是嫌棄.
他不要他作王.
在經文裡說到.
神已經選擇了一個更好的人.
可以代替蘇羅去作王.
當然我們知道.
神是很有憐憫的.
也不是即時將蘇羅從他的位置挪開.
不讓他作王.
一段時間也很長.
讓蘇羅有機會悔改.
神還是很憐憫人.
但神已經很嫌棄他.
神興起一個更好的人.
那個更好的人大家知不知道是誰?.
大衛.
以色列的第二任君王就是大衛.
神要興起一個更好的王.
去接觸他.
去跟從上帝.
而且神為自己的決定絕不後悔.
立定心.
興起一個新的人.
不要他.
人嫌棄我們.
我們感受上會有難受.
我不會說沒有感受.
但是那個嫌棄是很短暫.
不會影響到那麼深遠.
但我告訴你神的嫌棄是很深遠.
神的嫌棄是很深遠.
是很大.
神的嫌棄才是真真正正對我們來說.
是值得我們注意.
是我們害怕的.
人要我們不要害怕他.
除非是我自己犯罪.
不單得罪人還得罪上帝.
那我要害怕.

$^{481}$因為我得罪的是上帝.
如果不是的話不用害怕.
但如果今天是神的嫌棄.
是一件很大的事情.
值得我們反思.
也要立即認罪悔改.
而且神一嫌棄.
去到這個地步就不會後悔.
不會說我又回頭想.
我又改變了.
人就會改變.
神不會.
所以你看神有回應.
神看著我們也有回應.
譬如今天.
耶穌基督如果今天在這裡.
讓我們見到祂和我們一起.
你覺得耶穌基督今天會對我們說什麼話?.
不會.
祂只會對蘇羅感到遺憾.
只會嫌棄祂.
對我不會.
肯定?.
我的生命有沒有夜神不喜悅的?.
有沒有?.
這一刻.
主真的在這裡親自跟我們說話的時候.
你覺得主會跟我們說什麼?.
很欣賞我們?.
還是怎樣?.
值得我們思考.
我們是否很在乎主的心意?.
蘇羅有好幾樣東西他不討主喜悅.
我們看看他不討主的喜悅.
第一.
愛自己的榮耀.
多於愛上帝的榮耀.
這不只是蘇羅獨有.
我們從蘇羅身上看人性.
看人性的意思是.

$^{521}$我們也有機會是這樣.
我們會喜歡自己的榮耀多於上帝的榮耀.
哪裡看到呢?.
這裡說.
撒母耳清早起來見蘇羅.
有人告訴他.
蘇羅去到加密這個地方.
看他在那裡為自己做了什麼事.
拿了一個紀念碑.
再轉身去到吉格.
在打勝仗之後.
蘇羅為自己拿紀念碑是為了什麼?.
為自己.
為了什麼?.
榮耀自己.
我真是一個很有能力的人.
有能力的.
一個帶領軍隊的元帥.
一個元首.
你看我帶領軍隊得勝.
我真是很能打.
是我鎮住了亞馬利亞人.
亞馬利亞人不為患.
為那一區帶來多少患難.
我一起來我就贏了.
在這裡他忘記了是誰令他得勝.
是神自己.
是神自己令他得勝.
不是他自己.
他忘記了.
我們看回加密在右邊上方的地方.
中間是舒爾的抗議.
蘇羅在那裡為自己立了一個紀念碑.
紀念自己的豐功偉績.
在過程中他忘記了.
要將榮耀歸還給神.
讚美我們的上帝.
他帶領以色列軍隊得勝.
榮耀是屬於上帝的.
我們就是應該這樣.

$^{561}$其實我們在我們的事奉中.
我們都應該要這樣.
我們事奉有一個很好的.
譬如在人看來有一個很好的成就.
或者有一個很好的效果的時候.
我們不是在說.
我們那個效果有多好.
我這個人做得多好.
而是我們的上帝.
他能力我們保守整件事.
神親自工作.
令到我們的事奉的果效.
何等的好呢?.
我們要將榮耀歸還給神.
但我們問自己.
我們有什麼時候真的歸還給神?.
我們真的不愛自己的榮耀?.
真的專愛主的榮耀嗎?.
回顧我們之前的服事.
有多少時候真真正正將榮耀歸還給神?.
問問自己.
求神寬恕我們.
如果我們真的經常偷取上帝的榮耀.
經常渴望別人稱讚自己.
肯定自己的時候.
如果真的這樣.
而不是將榮耀歸還給神.
我們要小心.
我們渴望別人的肯定.
都是很自然的.
如果完全別人不肯定我們自己.
就很淒涼.
就好像服侍完全沒有作用.
肯定歸肯定.
仍然要將榮耀歸還給神.
如果不是神保守.
可以任何效果都沒有.
結果可以很差.
所以要記住這一點.
蘇狼這個位置是一個很大的軟弱和失敗.

$^{601}$他自我感覺良好.
看不到自己的罪惡.
這是第二個問題.
為什麼我們發現他自我感覺良好呢?.
撒母耳一到蘇羅.
蘇羅未曾說.
我已經遵守了王的命令.
我已經遵守了命令.
撒母耳覺得他自知之明.
明明命令很清晰.
要做到什麼程度.
他沒有做到.
他自己稱讚自己.
我已經遵守了神的命令.
撒母耳說.
我耳中聽到的羊叫牛鳴的聲音是什麼?.
究竟是什麼聲音?.
很想蘇羅自己可以認罪.
蘇羅有沒有認罪呢?.
沒有.
沒有,他覺得自己也很好.
蘇羅還說.
這些百姓從亞馬利亞得來的.
最重要的是什麼?.
是上好的.
是很好的.
要獻給神.
是獻祭的.
我為了上帝.
我嘗試解釋.
神一定喜歡的.
在舊約.
上帝教導他的百姓去獻祭.
透過獻祭.
學會如何和上帝相交.
在獻祭裡.
我們有贖罪制.
我們向上帝認罪.
或者我們完全奉獻.
我如何將自己的生命.

$^{641}$完全奉獻.
等等.
學會去感恩.
學會如何和上帝建立和平的關係.
等等.
神透過獻祭教導我們和他相處.
神也很體重獻祭.
是很重要的.
但是在這裡.
當蘇羅拿這件事作為藉口的時候.
神如何去看待.
一會兒有一個很重要的回應.
蘇羅說:我拿來獻祭.
我是要做給神你的.
有什麼不好.
還要解釋.
這裡薩姆爾忍不住.
要叫他閉嘴.
閉嘴.
讓我把上帝的話跟你說.
蘇羅經常不知道自己有什麼問題.
不意識到自己有罪.
我們的生命又如何.
我很多服事.
我做很多很多事情.
我參與很多事情.
我剛來紅印加的時候.
我真的很佩服弟兄姊妹.
心裡也很欣賞.
因為當時我們的事奉人手不多.
所以每一個弟兄姊妹都做很多.
同步做很多事奉的項目.
我心裡欣賞的是弟兄姊妹很願意為神.
很願意擺上.
也很愛神的家.
知道神的家有很大的需要.
後來慢慢發展的時候.
我就覺得弟兄姊妹可以去分工了.
多一點分工.
彼此配搭.

$^{681}$也是另一個很美好的事情.
在這裡我想說.
是否值得問一個問題.
神要我們只是參與了很多事奉.
做很多事情.
我們在這裡很忙碌.
就已經很好了.
就已經足夠了.
神看的是什麼.
是否就是這些憲制.
是這些很多的.
神都看重.
神看重我們服事.
從來不是不看重.
只是看重這些.
是否這些就是唯一的呢.
為什麼撒母耳要素羅閉嘴呢.
他講得自己這麼好.
有什麼問題呢.
值得我們看.
是否只是外在的事奉工作呢.
值得我們留意.
一會兒這裡會有回應.
第三件事.
素羅也很推卸責任.
他很推卸責任.
其實說真的.
是我的問題嗎.
是那些白色的.
是那些白色的收了牛和羊.
他們說要獻給上帝.
所以你要怪罪.
至少是誰起的.
是白色的問題.
他忘記了自己作王.
他有份的.
經文是說他和白色去連繫.
去連繫上好的牛羊.
他有份的.
白色這樣做.

$^{721}$他不可以純粹是他們問題.
自己沒有份.
他有份.
但他推卸責任.
純粹是他們的問題.
還有呢.
之後在經文裡看到他.
是在說什麼呢.
他懼怕白色.
在經文後面有說.
他很懼怕白色.
大家聽我讀那段經文後面.
後面當他終於認罪.
這次真的認罪.
不認自己犯罪都不行.
在這時候.
他說什麼.
在十五章很後.
在二十四節.
素羅對撒母耳說.
我有罪了.
我違背了耶和華的指示.
你的命令.
因為我懼怕白色.
我懼怕白色.
我聽從他們的話.
我怕他們.
他們說很想留下牛羊獻給上帝.
群體壓力.
群眾壓力.
大家都覺得對的.
我怎麼可以說不對呢.
我怕他們.
所以我就聽從他們的話.
是不是大部分人覺得對的.
覺得對的就一定對呢.
這樣很危險.
如果大部分人說的是犯罪的話.
大部分人覺得對.
這樣很危險.

$^{761}$我跟著一起.
我跟他們一起得罪上帝.
要想一想.
有沒有試過這樣的場景.
大火而一切都覺得那件事是對的.
但其實那件事是不對的.
但我因為群體壓力.
我不跟別人這樣.
我會被杯葛的.
在學校裡.
這些事情很容易明白.
同學們一起.
非理性的去欺負一個同學.
我們不由自.
其實都不會分辨是非.
我們要加入一起去欺負同學.
我不去欺負.
那些人不讓我靠近.
他們又杯葛我.
接著到他們欺負我.
我怎麼知道這個同學發生什麼事呢.
總之我知道大火而就是這樣.
我也要一起.
我怕別人不接納我.
你說學生是這樣.
大人不會.
大人也會的.
其實有一樣東西.
說起來大家可能不太留意.
很多年了.
在香港有一個同志運動.
這是世界性的.
很厲害的.
同志運動在推動同性戀.
一直不斷地推動在不同的國家.
其實有不少有心智的基督徒.
一直努力在同志文化抗衡.
但你會發覺.
其實有不少基督徒都被影響.
抗衡.

$^{801}$你知不知道.
其實現在整個社會.
很接受同志文化.
首先我必須強調.
我們不是要非常憎恨同性戀者.
不是.
上帝也愛世人.
我們也一樣愛世人.
就算是同性戀的人.
我們也一樣愛.
但我們不能認同罪.
這要分得很清楚.
不能認同這些同志運動.
在過程中.
我發覺很多基督徒受影響.
世界也是這樣.
我們要跟在一起.
為何要執著呢?.
態度好像很殘酷.
沒有愛心.
要小心.
愛不是對著罪惡.
小心分辨.
這裡素羅真的大聲說.
我只怕百姓聽從他們.
都是他們的問題.
我怕.
在這裡我想說.
素羅問題是怕人多於怕上帝.
人只是說.
這一生這幾十年影響我們.
說真的.
有沒有幾十年影響呢?.
是否真的可以每天這麼大影響力.
影響我們呢?.
我不相信.
但上帝是從過去到現在.
直到永遠都影響我們.
我的生死在上帝的手.
不是在人的手.

$^{841}$人可以奪取我的性命.
但不可以奪取我的靈魂.
能夠奪取我的性命.
也是神允許才可以這樣做.
但神影響我一生直到永遠.
素羅怕人卻不是怕上帝.
怕錯了.
你說這是他.
不是我.
我們很多時候都怕人.
很體重人的情面.
所以我最近在一些分享中.
我都會說我們會討好人.
我們希望人喜歡自己.
希望人喜歡自己也很自然.
但你喜歡到一個地步.
你渴望到一個地步.
是多於你渴望神喜歡自己.
這是一個問題.
經常都體重人.
是否喜歡我.
那些是否喜歡我.
是否喜歡我很重要.
他們不喜歡我就不開心.
你有沒有問過上帝喜不喜歡自己.
應該是求這件事第一.
應該敬畏上帝多於怕人.
我們很多時候就是怕人.
很多事情擺位就錯了.
素羅也是這樣.
所以失敗.
神很厭棄他作王.
我不願意神厭棄我們.
厭棄我們服侍他.
厭棄我們為他擺上什麼.
神厭棄我們才害怕.
在這裡.
經文的重點是什麼.
搏抑和占卜的罪相等.
還梗和拜偶像的罪孽相等.

$^{881}$素羅因為厭棄耶和華的命令.
耶和華也厭棄他作王.
他不聽上帝的命令.
上帝看他很大件事.
好像占卜的罪.
拜偶像的罪.
是這麼嚴重的.
神很憎恨這些事情.
拿回那件事來.
之後我們再說.
不聽從神的命令.
是最大最大的問題.
我想說這件事大到什麼.
在聖經裡說.
剛才也有說過.
耶和華喜愛繁祭和祭物.
喜如喜愛人聽從他的話.
上帝喜悅我們獻祭.
但上帝更喜歡我們聽命令.
聽命令啊弟兄姊妹.
是吧.
看啊聽命聖於獻祭.
信從聖於公羊的知方.
聽祂命令是神最喜悅.
我今天可以很多服侍.
參與很多活動.
不過我不聽祂命令.
我的生命和上帝的教導.
是背道而馳的.
我是空蕩蕩的整個人.
只有一個殼.
只有一些很外在的形式活動.
沒有實質的.
我仍然在罪惡裡.
我違背上帝的命令.
對不起神不會喜悅.
神不會喜悅.
是吧.
你看啊.
神會厭棄我們.

$^{921}$厭棄我們的服侍.
厭棄我們的一切.
上帝要我們整個人去歸向祂.
這是非常重要的.
喜愛繁祭和祭物.
喜如喜愛人聽從他的話.
看啊聽命聖於獻祭.
信從聖於公羊的知方.
我怎知道上帝對我的人生.
這些命令是什麼呢?.
所以回歸聖經裡.
我的聖經已經封塵.
我都不知道上帝對我說什麼.
那你又何從來有什麼命令呢?.
禱告也不多幾次.
我也沒有打算聽上帝的聲音.
我已經人生幾十年經驗.
我知道怎樣做是好.
我覺得上帝一定會喜歡.
我就這樣去做.
就這樣在教會生活幾十年.
我相信上帝沒有什麼異議.
這樣就不是神的心意.
神是活人的神.
神是有回應的神.
神的回應.
有些時候你感覺好像很久.
但有些時候很即時.
神在看著我們.
在等著我們.
隨時會回應我們.
但有沒有開放你的生命.
去聽神的聲音.
還有沒有一個跟從的生命.
你將你的生命放在神的手裡.
不是自己決定.
是憑經驗.
憑自己喜悅.
你說我喜悅的上帝一定喜悅.
是不是?.

$^{961}$那你可以做神了.
那一定不是.
我們的神是很奇妙的神.
祂有祂的旨意.
心願在新的一年.
我們彼此去鼓勵.
這是很重要的.
聽從神的命令是第一.
聽啊!弟兄姊妹.
聽神命令.
因為素羅不聽神命令.
神厭棄他.
不要他作王.
到最後神將那個國.
給了大衛和大衛的後裔.
素羅聽到撒母耳的責備之後.
他也沒有徹底悔改.
他沒有徹底悔改.
什麼叫徹底悔改?.
即刻將牛羊滅了.
將瓦格抓出來殺了他.
這樣叫徹底.
不單只是認罪.
實際行動要做一些事情.
命令所要求的.
他也要徹底地做.
但他沒有這樣做.
最後誰替他埋沒?.
是撒母耳替他埋沒.
撒母耳替他埋沒.
就是找人叫瓦格出來.
要殺死瓦格.
是撒母耳最後埋沒.
不是素羅.
素羅聽了神的責備之後.
也沒有動手.
還是不肯悔改.
所以難怪神最後徹底地厭棄他.
神不後悔.
是因為神太認識他.

$^{1001}$知道他是一個不會悔改的人.
我們今年新的一年開始.
其實神最喜悅我們是聽他的命令.
要聽他的命令.
要給時間神.
我們給了很多時間.
我們的家人.
我們所愛的配偶.
兒女.
我們的朋友.
所有這些人都是重要的.
但我們更要給時間上帝.
讓他跟我們說話.
告訴我們人生怎樣走下去.
大小決定.
他對我們人生.
我們這個人.
怎樣活著.
才是他喜悅的.
讓我們過一年.
是一個順從上帝命令的一年.
我們整個教會也一樣.
教會也要順從上帝的命令.
走在神的心意當中.
我相信神必定會祝福.
我們一同禱告.
天父上帝.
其實我們並不比素羅好.
但上帝你的憐憫.
就淋到我們身上.
讓我們看到這個版.
我們可以引以為鑒.
我們還有機會去悔改.
在新的年.
新的開始.
讓我們很愛慕你的話語.
很想聽到你的聲音.
明白你對我們生命的旨意.
知道上帝你的命令是甚麼.
願你讓我們去回應你.

$^{1041}$主啊 縱然我們有軟弱.
但我們不用怕.
因為上帝你的靈與我們同在.
有聖靈在我們生命裡.
一定會幫助我們.
幫助我們.
可以行出上帝你的旨意.
在新的年.
將曾睿貴牧師帶到我們當中.
帶領教會的發展.
主啊 求你.
在當中賜曾睿貴牧師力量.
給他力上加力 恩上加恩.
祈你藉著他帶領我們眾弟兄姊妹.
要去聽上帝你的聲音.
帶領我們去愛慕主.
帶領我們這一年.
只要去服從上帝你的命令.
成就你的心意.
好讓洪恩家大大的復興.
我們就是這樣.
將整間教會交在你的手裡.
求你賜福.
奉耶穌基督的名字祈禱.
阿們.
\newpage



\section{馬太福音 7:1-6}
\label{sec:wCnaB7fTUZk}
\textbf{2024.01.14 《以愛結連》 馬太福音 7\_1-6 (和修版) 曾裔貴牧師}
\newline
\newline
連結: \href{https://youtube.com/watch?v=wCnaB7fTUZk}{\texttt{https://youtube.com/watch?v=wCnaB7fTUZk}} ~~~~ 語音日期: 2024-01-16
\newline
\newline
\hyperref[sec:r5CRKC0_eBw]{\small{< < < PREV SERMON < < <}}
~
\hyperref[sec:index_chronic]{\small{[返順時目]}}
~
\hyperref[sec:index_scriptual]{\small{[返順卷目]}}
~
\hyperref[sec:sZSvUc8jKA0]{\small{> > > NEXT SERMON > > >}}
\newline
\newline
馬太福音 7:1-6
\newline
\begin{longtable}{cl}
\hline
\hline
章節 & 經文 (和合本修訂版)\\
\hline
7:1 & \begin{tabularx}{0.7\textwidth}{X} 「你們不要評斷別人，免得你們被審判。 \end{tabularx} \\ \\ \relax
7:2 & \begin{tabularx}{0.7\textwidth}{X} 因為你們怎樣評斷別人，也必怎樣被審判；你們用甚麼量器量給人，也必用甚麼量器量給你們。 \end{tabularx} \\ \\ \relax
7:3 & \begin{tabularx}{0.7\textwidth}{X} 為甚麼看見你弟兄眼中有刺，卻不想自己眼中有梁木呢？ \end{tabularx} \\ \\ \relax
7:4 & \begin{tabularx}{0.7\textwidth}{X} 你自己眼中有梁木，怎能對你弟兄說『讓我去掉你眼中的刺』呢？ \end{tabularx} \\ \\ \relax
7:5 & \begin{tabularx}{0.7\textwidth}{X} 你這假冒為善的人！先去掉自己眼中的梁木，然後才能看得清楚，好去掉你弟兄眼中的刺。 \end{tabularx} \\ \\ \relax
7:6 & \begin{tabularx}{0.7\textwidth}{X} 不要把聖物給狗，也不要把你們的珍珠丟在豬面前，恐怕牠們踐踏了珍珠，轉過來咬你們。」 \end{tabularx} \\ \\
[1ex]
\hline
\hline
\end{longtable}
$^{1}$弟兄姊妹早安.
再一次我們祈禱.
感恩我們再一次能夠聆聽上主的真道.
求你藉這一段聖經.
讓我們能夠學習彼此在主裡面相愛.
我們仰望等候你的教導.
奉耶穌基督的名.
阿們.
其實非常感恩的.
因為我已經病了一個星期.
看了兩次醫生.
昨晚回到家裡.
咳嗽變得更加頻密.
整晚都睡不著.
終於在床上.
因為我沒有看時間.
不知道幾點.
睡在床上.
知道自己是醒了.
不是做夢.
我說我明天要施餐.
要為你講道.
讓我由起床時不要咳嗽.
到現在.
已經沒有咳嗽了.
我也不知道為什麼.
感恩.
這個訊息.
剛才主席已經讀完了.
如果你很心水清.
主耶穌教我們去挑刺.
你有沒有試過被刺刺到.
很痛的.
眼睛也刺刺.
怎麼挑刺.
是大手術.
是不是.
我們人有很多的習慣.
都不是一下子可以改變.
我這個星期.

$^{41}$其實不是這個星期.
是兩個星期.
我要從家裡元朗.
過來雄恩.
洪水橋.
我連續兩次.
開車回到感受.
改變還沒轉.
腦袋告訴自己.
我正在雄恩.
但那個習慣.
原來我已經走去感受那條線.
然後就要很刻意地轉.
我們有很多這些的習慣.
當然這些沒什麼所謂.
繞一圈.
在湖泉處轉一轉.
那些很小事.
但是我們人都有很多不同的習慣.
不同的記憶.
剛才和弟兄吃早餐的時候.
他留意到.
很細心留意到有一樣東西.
他說曾牧師.
原來你是左手的.
是啊.
小時候一出生就已經習慣用左手.
左手的人特別聰明.
他這樣說.
我沒理由說不是.
當然是稱讚自己聰明.
其實就不是.
我們很多這樣的約定俗成.
之前冬至.
你記不記得很冷.
有些專家就說.
其實就是那些.
媽媽說.
冬至冷.
過年一定暖.

$^{81}$個個都點頭.
原來我們有這麼多這些習慣的記憶.
什麼十個光頭.
九個就會富.
又是這樣說的.
但其實沒有根據.
很多.
很多這些的習慣.
很多這些約定俗成.
有些是無雙大雅的.
但是.
主耶穌在這裡告訴我們.
我們不可以評斷別人.
首先就是這段聖經肯定告訴我們.
不是不可以做判斷.
因為第六節.
耶穌就已經教門徒.
另外一個角度.
看事物.
看人.
有些人原來是.
抗拒.
刻意敵視真理福音.
耶穌看這些人是豬是狗.
當然是用比喻的形式.
這樣來說.
如果你沒有這個判斷力.
你怎麼能夠告訴我們.
這些人是一直抗拒真理.
攔阻人信耶穌.
例如最近在香港.
有一個叫新天地的教派.
到處去教會.
找教務拍照.
拍完之後.
到處告訴人.
我和這個教務拍過照.
他認同我們的工作.
有一個大會在香港舉行.
如果你沒有做判斷.

$^{121}$你不知道他是異端的.
所以這段聖經肯定告訴我們.
不是不可以判斷.
而且人類與生俱來.
神給我們判斷的能力.
動物是沒有的.
只有人類才有的.
不過問題就出在這裡.
耶穌說.
我們如果用一個沒有救恩的生命.
過去救人的那種價值觀.
去看其他的人.
我們內心的無論自我.
無論是自異.
無論是我們內心那種.
不肯解除冤仇的那種價值觀.
就會對人有一個很強烈的先入為主.
我們就沒有做一種.
好像主耶穌所說的.
你要怎樣看自己呢?.
這段耶穌說得很重.
不過如果你看回上文第六章.
主耶穌要求我們要完全.
好像天父完全一樣.
要我們能夠在這個天國裡面.
要學孝臣.
上文一直在說的門徒.
聽了很多主耶穌的教導.
知道那些是很重要的真理.
耶穌甚至重新去解釋十誡.
那個背後的精髓.
你可以想像得到多麼高要求.
門徒可以聽到這麼好的真理教導.
但慢慢地.
有時將自己那個不知不覺的.
以為做到的價值觀.
套用在別人身上的時候.
就會有這個危險.
就會有這個警告.
耶穌就要提醒他的門徒.

$^{161}$不可以用你自己那個狹隘的心靈.
去判斷別人.
這個也是相當重要的.
各位弟兄姊妹.
這個告訴我們.
因為你怎樣去判斷別人的時候.
也必怎樣被.
我的舊的和合本就是論斷.
和判斷一樣的.
都是有一個判斷的.
這種這樣的判斷.
有時不知不覺我們就成為了審判官.
耶穌告誡我們.
不是不可以看一些事物.
有一個分析的能力.
應該要有分析的能力.
但是不是以審判官的姿態.
去做這件事.
是要讓人看到我們在救恩底下.
大家是神的子女.
我們大家是以憐憫.
恩慈.
彼此的寬恕.
彼此的接納.
這樣來看待一些事情的發生.
這時候.
又會換取到一份彼此的愛.
又說剛才的例子.
我吃早餐和弟兄一起.
我叫那個包變成多士.
加一元無所謂.
因為多士比白的麵包.
那個升糖指數沒那麼高.
叫多士.
很快就整碟.
應該是油尖多.
送到我的桌上.
弟兄就說多士來了.
我馬上拿起吃.
接著那個服務員就來到.

$^{201}$不好意思原來是隔壁桌的.
我吃了怎麼辦.
他就很不好意思.
隔壁桌就問.
為什麼我吃多士.
怎麼辦.
那個服務員就很不好意思.
就問我.
問我們兩位.
如果入你們的帳.
好不好.
沒問題.
我吃了當然好.
不然你很難做.
那個服務員.
千多謝萬多謝.
本來我可以說不關我的事.
是你給我的.
我吃了你自己搞定.
你自己跟上司說.
但不應該.
因為我們多走一步.
就建立到一份彼此的寬恕.
彼此的一個接納.
何樂而不為呢.
臨走的時候還要跟他說聲多謝.
我們走了.
回教會了.
那個服務員就感受到.
為什麼這個人會這樣.
給多點錢.
都不關你的事.
吃了他那碟多士.
親愛弟兄姊妹.
這裡告訴我們.
我們要有很強的判斷力.
但對人.
要有很強的愛.
因為主耶穌這樣說的時候.
祂有一個很強的信念.

$^{241}$就是原來我們自己.
在上帝的眼中.
其實有很大很大很粗的良目.
耶穌自己這樣用很誇張的比喻.
相信大家都肯定知道.
可能會有一塊木在你的眼裡.
第三節.
為什麼看見你弟兄眼中有刺.
相信這個是我們每個人.
其實不斷不斷都會看到別人有刺.
但是很少人.
能夠苦心自問.
能夠每天的來到去看自己.
虧欠上主的恩典.
而想到自己其實是有良目的人.
而神不會挑我們的良目.
反而會寬恕白白的赦免我們.
但是這一點.
已經足夠我們大家.
來到去可以在神的恩典裡面.
去寬恕別人.
去接納別人.
使大家在一個和諧裡面.
多開心呢.
我下次再去那間大法.
新的大法不是舊的那間.
新的大法.
再見那個職員.
跟他說我是牧師.
會很不同.
如果你據理力爭.
其實沒有錯的.
我沒有做錯事.
是你放在那裡.
我當然會吃.
但不是.
我們應該比其他人多走一步.
那你就會見到.
神有很奇妙的恩典和工作.
你自己眼中有良目.

$^{281}$怎麼能夠對你的弟兄說.
用我去掉你眼中的刺呢.
這個又是另外一個問題了.
看不到有良目.
然後不斷以自己為審判官.
因為認為自己做得好了.
認為自己看事物的判斷很準確.
來到對別人.
不是對自己.
是對別人.
是有一種這樣的.
而且不是一種善意的包容提醒.
很開心.
榮翔上一個星期跟他吃飯.
他說曾牧師.
有一樣東西我很認真問你.
如果我聽到一些事是跟你有關的.
我能不能夠直接問你.
我說當然好.
應該.
正正就是.
如果我有什麼做得不對的.
你就要提點我.
我們就是弟兄姊妹.
我是崗位上.
我是你的牧師.
但我是你的弟兄.
同不同意嗎.
正正就是這樣.
但我們就很容易.
不單不理自己有良目.
一來到就看到別人的刺.
就來到對別人.
有很多很多的評斷.
甚至超過神的寬恕恩典的一種責備.
這樣是不行的.
耶穌說.
所以耶穌刻意跟門徒說.
一定要這樣來處理.
這個心態價值的問題.

$^{321}$4月1號.
弟兄姊妹.
個個都很快想扔掉家裡的垃圾.
第二個星期.
1月.
還有很多時間.
有一條看不到的良目.
在我們大家的價值觀裡面.
我們能不能夠在主的恩典裡面.
再看自己.
其實救恩令我們成為今天的人.
我嘗試用一個例子來說.
這個成見.
如果能夠消除.
是帶來一個很大的福樂.
真的是一個群體的平安和喜樂.
有一個年輕的.
在台灣的.
真人真事.
在台灣的一個年輕的女孩.
她自小就很想做記者.
她在新聞系如願的畢業.
真的去做一個報館.
做記者.
跑前線的新聞.
很開心.
能夠她的心願達成.
但不知道什麼原因.
有一次的意外.
她的聽覺嚴重的受損.
她不能夠再好像平常那樣.
來做訪問.
報館就安排她做後勤資料的搜集.
和同事一起.
有一次很大的突發新聞.
同事要趕著跑前線.
回來要準備材料.
要馬上出街報到.
這個同事始終她的手腳慢.
大家溝通上是有一定的困難.

$^{361}$同事因為太過緊急.
太過快.
就好像對她站在一旁.
我們很趕.
令她很受傷.
令她覺得她現在這個意外.
她自己現在這個困難.
已經被自己的同事.
都好像有一些對她的不好.
她一路自己思前想後.
一路自己這樣去想.
越想就越不開心.
其中有一次.
令到她決心要離職.
每個月都會有大家一起.
同事之間的一個.
在一些房間裡面.
開開心心的吃東西聊天.
每個月都有一次的歡聚.
居然不再叫她了.
很激氣.
寫了封信辭職信.
想著進去敲門.
想著就走了.
一路去的時候.
越走就越激氣.
房裡面小聲說大聲笑.
激到她沒辦法.
真的決心打開門.
丟下封信.
怎知打開門之後.
眼淚就這樣來了.
原來一班同事.
在學手語.
希望可以和這位同事.
往後的工作上協調.
小聲說大聲笑.
就是那些手勢錯了.
很搞笑.
不是大家玩耍.

$^{401}$是樂她.
這個女同事很受感動.
她自己親自去報館投稿.
告訴別人.
這些誤會可以解除.
當然人與人之間.
未必真的這麼戲劇性.
很多時候.
隱藏在心底的誤會.
一直都解除不了的時候.
就會帶來人與人之間.
有很多傷痕.
傷痕的累積.
我們記住我們是人.
我們很容易就會慢慢.
來到只會繼續.
在我們內心積累一些苦毒.
越來越我們就更加沒有出路.
所以主耶穌很清楚告訴我們.
教我們要撩開那些刺.
在祂的眼中.
我們只不過是刺.
就是說可以拿走.
如果不是.
我們如果扮演著.
主耶穌的神的身份的角色.
來審判別人.
我們就會變成日積月累的良木.
所以你可以想像到.
第十三節.
十二節.
所以無論何時.
你們願意人怎樣待你們.
你們也要怎樣待人.
因為這就是律法和先知.
神的真理.
一個信了主耶穌的人.
首先要學習的是.
你想人怎樣對待你.
自己就首先走出那一步.

$^{441}$每個人走出一步.
我在廣場慣了.
看到影子們回來很開心.
就會打招呼.
每個婆婆都說.
我看到牧師很開心.
現在第二個星期.
你看到牧師在門口.
雖然剛剛太陽升.
很曬.
但其實很開心.
看到影子們慢慢走回來.
其實我們大家能夠每個星期聚會.
是一件很開心快樂的事.
影子們你要穿最漂亮的.
因為你要見神.
端莊.
整潔.
最好和你平時沒有關係的衣服.
更加好.
因為我們來朝見真神.
帶著最好的敬拜神.
我們錦繡堂有菲律賓語.
有印尼語的弟兄姊妹.
他們每個星期穿著自己的制服.
來獻唱.
是很開心的事.
平日他們是一個工人.
但來到星期天.
他們就是神的兒女.
洪恩家的弟兄姊妹.
你們是神的兒女.
你們要學習.
如何得到回報.
就是你想人怎樣對待你.
你就要怎樣對待人.
記不記得中間夾雜著第七節.
到第十一節.
教祈禱.
所以這段聖經是一個首尾的呼應.

$^{481}$頭尾有關連.
中間三文字的餡就是祈禱.
所以這個祈禱.
是祈求神讓你看到自己的內心.
不是求那些飛機大炮.
不是求那些六合彩的號碼.
不是這些.
是求你和人的關係.
是可以你的人生與人為善.
神一定聽的.
因為這個正正就是先知的道理.
律法的道理.
就是主耶穌來這個世界.
建立的就是和平之主的福音.
所以你想想.
你走出來.
見到每個人都和你是很開心.
是沒有任何的嫌隙的.
這豈不是救恩給我們最大的福樂嗎?.
與神已經和好.
與人更加的和好.
我們一起祈禱.
多謝主.
讓我們能夠短短的經文.
已經有很深的教訓.
我們真的要在4月1日之前.
將我們自己的.
無論是良目.
或者我們都認為可能是次.
要拿走.
要來到去在神的恩典裡面.
學習我們的價值觀.
是要追求愛神.
更加是要愛人.
更加是要與人和睦.
來到去在主的恩典裡面.
弘恩加成為一間.
真的能夠叫其他人羨慕的教會.
這樣感恩禱告.
奉耶穌基督的名.

$^{521}$阿們.
真的很感恩.
到現在聲音還是這樣出來.
是完全沒有cut.
謝謝大家.
\newpage



\section{約書亞記 1:1-9}
\label{sec:sZSvUc8jKA0}
\textbf{2024.01.21 《使命的承傳》 約書亞記 1\_1-9 (和修版) 郭文池院長}
\newline
\newline
連結: \href{https://youtube.com/watch?v=sZSvUc8jKA0}{\texttt{https://youtube.com/watch?v=sZSvUc8jKA0}} ~~~~ 語音日期: 2024-01-23
\newline
\newline
\hyperref[sec:wCnaB7fTUZk]{\small{< < < PREV SERMON < < <}}
~
\hyperref[sec:index_chronic]{\small{[返順時目]}}
~
\hyperref[sec:index_scriptual]{\small{[返順卷目]}}
~
\hyperref[sec:dSTGrtAycig]{\small{> > > NEXT SERMON > > >}}
\newline
\newline
約書亞記 1:1-9
\newline
\begin{longtable}{cl}
\hline
\hline
章節 & 經文 (和合本修訂版)\\
\hline
1:1 & \begin{tabularx}{0.7\textwidth}{X} 耶和華的僕人摩西死了以後，耶和華對摩西的助手嫩的兒子約書亞說： \end{tabularx} \\ \\ \relax
1:2 & \begin{tabularx}{0.7\textwidth}{X} 「我的僕人摩西死了。現在你要起來，和眾百姓過這約旦河，往我所要賜給以色列人的地去。 \end{tabularx} \\ \\ \relax
1:3 & \begin{tabularx}{0.7\textwidth}{X} 凡你們腳掌所踏之地，我都照我所應許摩西的話賜給你們了。 \end{tabularx} \\ \\ \relax
1:4 & \begin{tabularx}{0.7\textwidth}{X} 從曠野和這黎巴嫩，直到大河，就是幼發拉底河，赫人的全地，又到大海日落的方向，都要作你們的疆土。 \end{tabularx} \\ \\ \relax
1:5 & \begin{tabularx}{0.7\textwidth}{X} 你一生的日子，必無人能在你面前站立得住。我怎樣與摩西同在，也必照樣與你同在；我必不撇下你，也不丟棄你。 \end{tabularx} \\ \\ \relax
1:6 & \begin{tabularx}{0.7\textwidth}{X} 你當剛強壯膽，因為你必使這百姓承受那地為業，就是我向他們列祖起誓要給他們的地。 \end{tabularx} \\ \\ \relax
1:7 & \begin{tabularx}{0.7\textwidth}{X} 只要剛強，大大壯膽，謹守遵行我僕人摩西所吩咐你的一切律法，不可偏離左右，使你無論往哪裡去，都可以順利。 \end{tabularx} \\ \\ \relax
1:8 & \begin{tabularx}{0.7\textwidth}{X} 這律法書不可離開你的口，總要晝夜思想，好使你謹守遵行這書上所寫的一切話。如此，你的道路就可以亨通，凡事順利。 \end{tabularx} \\ \\ \relax
1:9 & \begin{tabularx}{0.7\textwidth}{X} 我豈沒有吩咐你嗎？你當剛強壯膽，不要懼怕，也不要驚惶，因為你無論往哪裡去，耶和華你的神必與你同在。」 \end{tabularx} \\ \\
[1ex]
\hline
\hline
\end{longtable}
$^{1}$各位弟兄姊妹早晨.
主今次是我第一次來到紅屋堂的地方跟大家分享神的話.
多謝曾牧師和教會邀請給我這個機會來到大家當中.
時間是比較倉促的.
教會也能夠遷就時間讓我今天跟大家分享.
今天我想跟大家思想的題目就是「使命的承傳」.
基督徒也知道我們每一個人.
你和我都是應該有耶穌基督的使命的.
我希望今天能夠將這個題目講得清楚.
我首先可以代表博路神學院向大家問安.
博路神學院過去幾十年一直在九龍塘的地方作為我們的校舍.
由2019年開始我們就搬到美孚.
更加近這裡.
也因為這個緣故我們有這個聚會也想請大家考慮參加.
如果當中有人在考慮會否一生會回應上帝給你的感動,呼召.
作傳道人的話.
請你可以考慮參加一個叫做「時代的呼聲」.
一個全職事奉探討的聚會.
3月2日星期六.
詳細情況請你可以留意我們的網頁.
我們其實神學院也有一些課程給信徒的.
是一些證書課程.
即是說一些普及課程.
是不需要做功課,不需要考試,不需要入學等等的資格的.
你願意讀就可以了.
是一些普及課程,好像講座那樣.
當然不是學位.
有一個日間的金靈課程.
如果你是星期二,星期四上午可以去到美孚上課的話.
歡迎你考慮參加,現在也在招生中.
還有一個就關於我自己了.
也要講一講,不過就很快很快.
因為這個星期三就要開學了.
這個星期三晚.
如果你真的很有興趣的話.
請你盡快這一兩天請牧者簽名報名.
就是神學專題探討.
不知道大家信了主這麼久會不會覺得.
很難理解,怎樣解釋三維一體呢?.
有沒有聽過聖經無誤的教義呢?.

$^{41}$靈因運動怎樣影響我們教會的生活呢?.
基督,又是人又是神,怎樣去理解呢?.
比較常被問的問題.
我就會在星期三晚有一個神學專題的探討.
也是普及課程來的.
如果你有興趣的話,有六堂的時間.
可以請你考慮參加.
好了,就這麼多廣告時間了.
現在就是和大家思想,使命的承傳.
我想到教會,想到你們的情況.
想到神給我的感動.
我就選了這段聖經.
約書下記一章,一到九節.
我相信對大部份的弟兄姊妹來說都不陌生.
這段經文就好像開始一個新的時代.
摩西一個很偉大的領袖.
雖然歷盡艱辛.
但真是一個在以色列人歷史裡面.
可能是最,如果不是最.
都是其中一個最偉大的領袖.
到現在,這個好像成為了他們歷史上的代表人物.
不過進入到他的時代結束.
進入到約書也,他幫手的時代.
開始了約書下記.
面前有很大的挑戰,有很大的使命.
但是領袖的改替.
卻是帶給這段經文.
可以說對我,對你.
可能都是一個很大的參考作用.
聖經開始這樣說.
耶和華的僕人摩西死了.
然後這段經文最後是怎樣呼召.
或是再次呼召約書也回應使命.
請你留意經文.
這裡說上帝說摩西死了.
不過上帝的作為,上帝的使命並未完成.
要繼續的,這個說得非常清楚.
他說現在你要起來帶領百姓.
來到我所應許給你們的地方.
這個地方是非常非常大的.

$^{81}$這本和學本就寫著.
是從曠野和者尼巴嫩直到伯拉大河.
即是這個又法拉底河.
黑人的傳遞又到大海日落之處.
大海對以色列人來說即是地中海.
即是他們地圖的西邊.
伯拉大河是地圖的東邊.
東到西,然後北邊尼巴嫩.
然後南就是這個曠野.
可以說這個整個地方.
以色列從來未試過成為它的國家.
太大了,大成怎樣呢.
你看看地圖.
如果你看看地圖.
現在今天打仗的地方是靠著地中海.
細細很窄的地方,是加薩的地方打仗.
其實以色列的地方大概是皇上的左邊.
近海綠色長長的那一截就是它的地方.
但其實上帝若用約書客的經文說.
整個綠色的地方都是他們的.
很大的.
然後如果放大看.
你會看到地中海裡面的地圖更加清楚.
你會看到原來上帝答應給以色列人.
是很大的地方.
不過種種原因.
特別是以色列人.
他們沒有履行使命.
所以最終.
以我所知道歷史上.
他從未試過這麼大的地方.
還有這段經文.
不單止是地方的問題.
你還看經文多次提到一件事.
我的僕人沒勢死了.
你要起來,你要得這個地方.
但是得這個地方談何容易呢.
怎麼能夠一班從埃及走出來.
抗日飄來幾十年的一班以色列人.
他現在能夠去到這麼大的地方.

$^{121}$這麼大的挑戰呢.
所以我說當以色列人往外看.
看看周圍環境的時間.
是困難重重的困難.
踏入2024年.
我們想2023年.
或者2019年到2023年疫情過了.
應該經濟好一點.
民生好一點.
然後整個環境好一點.
大家知道2024年開始的時候.
一點都不是我們想像.
疫情未完結.
戰爭越來越多.
也沒有好像完結的地方.
似乎中東.
就是以巴的戰爭.
還蔓延到整個中東.
甚至是歐洲的地方.
經濟,香港也好.
中國國內也好.
或者是這個世界的情況也好.
似乎都不是太過樂觀.
教會也是.
我不知道你們有沒有受移民潮的影響.
我們教會在新蒲崗的地方.
大概,最低限度.
大大小小的.
少了40\%移民了.
我去英國和他們坐下來吃飯的時候.
要三四圍桌.
他們專門見我,走在一起.
我們教會不是很大.
教會只有大約100多200人.
走了差不多40\%.
包括六個傳代裡面.
四個都走了.
連同唐主任.
所以,當你在這種環境裡面的時候.
是很捨棄的.

$^{161}$有些甚至是長者.
子女都走了.
可能你也祝福他們,為他們的好處.
為他們孫的好處.
你祝福他們.
但是流下來的那種感覺.
那種孤單.
那種無助.
還有不知道前面會怎樣.
我快要退休了.
開始想一下.
特別是見到老人家退休生活.
開始投入去想一下.
年輕不會想的.
但是到了你這麼上一年級想.
你就見到那種都很辛苦.
將來要不要進老人院呢.
一人照顧怎樣呢.
這些都是當你往外看的時間.
重重困難.
你同不同意嗎.
在2024年好像到了新的年度.
疫情過了以為有一個好的開始.
但是人因為罪的緣故.
直接間接不停受到影響.
約書亞他帶領這班的以色列人.
要完成上帝給他的使命.
差不多是一句說話.
用英文說.
Mission Impossible.
不可能的.
不可能完成的.
不單止這樣.
你再看經文怎樣說.
三節聖經.
第一節.
第二節.
第三節.
開始三節.
三節都提到摩西死了.

$^{201}$提到這個偉大的領袖.
第一節說.
而我僕人摩西死了以後.
第二節.
我這僕人摩西死了.
我知道他死了不用說兩次.
但都要繼續說.
摩西死了.
很強調這件事.
第三節說.
我會照著摩西的話賜福給你.
這個是稱呼這個.
曾經在這個時代帶領以色列人的偉大領袖.
三次提及到他.
兩次提到他死.
第三節說.
我答應的話都會成就.
現在將這個使命交給約書亞.
約書亞怎樣稱呼.
大家有沒有留意到.
稱呼摩西叫做耶和華的僕人.
第二節.
我的僕人.
第三節.
我所應許摩西的.
全部神和摩西好像很密切的關係.
叫做耶和華的僕人.
但約書亞呢.
摩西的幫手.
約書亞好像沒有自己的身份.
要說摩西才知道誰是約書亞.
就好像師母.
大家知不知道曾牧師太太叫什麼名字.
大家知不知道.
沒有自己的身份.
要說曾師母才知道她是誰.
大家明白嗎.
我太太也是.
我問了一些回族.
很多都不知道她姓什麼都不知道.

$^{241}$只知道她是郭師母.
這個同樣都是摩西的幫手.
要這樣介紹她.
還要介紹她是爸爸.
亂的兒子約書亞.
從來不需要介紹摩西的爸爸是誰.
大家都知道.
她是一個偉大的人物.
但約書亞呢.
可以說是一個無名小卒.
她的身份只是協助這個偉大的僕人.
然後聖經還這樣說.
他說我怎樣與摩西同在.
也必照樣與你同在.
我必不撇下你.
也不丟棄你.
然後三次說.
你當剛強壯膽.
不用怕.
但你有沒有試過看醫生呢.
如果進到醫生房.
他第一句說話是.
我們做了化驗.
不過你不用怕.
那是什麼意思呢.
就是你應該很害怕.
當醫生跟你說不用怕的時候.
就是你很害怕.
很久之前.
我剛在神樂院畢業的時候.
吃錯東西.
進了醫院.
很嚴重的腸胃炎.
醫生小心一點.
幫我做很多檢查.
檢查了第三天.
也找不到什麼原因.
然後醫生跟我說.
你不用怕.
我們再幫你做多一個實驗.

$^{281}$再做多一個化驗.
我們懷疑你是血癌.
你說你不用怕嗎.
不用怕嗎.
通常當你特別難的時候.
才跟你說不用怕.
我和你一起的.
我不會丟下你的.
這件事也很害怕.
這段經文不斷重複第六第七第九節.
你當剛強壯膽.
這件事是非常難的.
事實上當我們看回背景.
約書亞圍爐的埃及地長大.
但摩西呢.
次郎傳說.
他學了埃及人一切學問.
他是一個非常學問的領袖.
但約書亞呢.
大概他沒有受過教育.
只能在家庭受希伯來人家庭的教育.
是一個不能夠和摩西比較的教育背景.
然後他帶領的人.
是以色列人當中逃難的烏合之眾.
但迦南人堅誠善戰.
並且有很大的後援的力量.
然後更難的就是.
這個偉大的使命.
就是這個偉大的摩西也不能夠完成.
現在交給摩西的僕人.
亂的兒子約書亞.
叫他不用怕完成.
是不是已經很怕呢.
所以不單外面困難重重.
你還往內看.
看看自己.
看看自己所有的.
你覺得自己很渺小.
弟兄姊妹.
我發覺教會裡面.

$^{321}$有自大的人.
但是比起自卑的人呢.
我覺得在我觀察裡面.
自卑的人多很多.
我不行的 你不用找我了.
特別到選舉的時候.
大家最自卑的.
我不行的 你不用找我了.
特別當中的長者.
我不知道你們有沒有這樣的情況呢.
我們教會裡面長者.
我發覺那段時間他很喜歡說一句話.
我年紀老了.
不懂得做了.
我左而入 右而出.
我現在也有這樣的感覺.
有些人是過目不忘的.
我現在差不多過目即忘了.
然後就說.
不用找我事奉了.
找年輕的那些.
這些說話好像傳染病.
這些本來七八十歲說的人.
慢慢五六十歲就說了.
我快要退休了.
教會做了二十多三十年.
教給年輕的 他們多一點活力.
三四十歲就跟著說.
我做了團結職員很多年了.
上有高堂 下有細蚊仔.
所以不要了 找年輕人多一點活力.
十多歲做好 交給他們吧.
好像傳染病.
傳染到整個教會都覺得交給其他人吧.
然後就跟長者說.
你說我年紀大了 沒用了.
讀書又不記得了.
就不用事奉了.
哪裡聖經教你的.
當然沒有 聖經哪有這樣說.

$^{361}$聖經沒有教 為什麼你當真理傳的.
傳福音就不見你這麼熱心.
傳這些很熱心的.
弟兄姊妹.
會不會你都是往來漢.
覺得自己很難承接使命呢.
說到使命.
都好像跟你沒有關係.
我年紀大了.
我不識字.
或者是我身體不好.
不關我的事.
我請我們重新再想想.
上帝喜歡我們怎樣生活.
過去這個星期.
我機會.
雖然在外地 剛從外地回來.
我機會都在外地的時候.
抓緊的機會 早上六點多.
都跟在美國一位退休的宣教士.
我的老師去事奉.
他今年九十多歲.
身體很差.
然後經常跌倒.
給我看他的照片.
和跟他事奉的時候.
頭撞到腫了.
瘀了 眼睛瘀了.
很辛苦.
他曾經有一段時間說.
我想快點回天堂.
我很辛苦的.
但後來有很多人鼓勵他.
他重拾.
他過往.
這位老師叫許得利姑娘.
是我認識的所有人當中.
最熱心傳福音的人.
任何情況都是傳福音的.
我們叫做神奇女俠.

$^{401}$任何情況都是傳福音的.
他在香港生活了很多年.
一口流利的地道的粵語.
和街邊的掃地都是傳福音.
買菜都是傳福音.
有一次他更神奇.
他去廣州 大陸.
他去到那些吃雲吞麵的地方.
吃魚蛋粉的地方.
他坐下來很擠迫地吃.
他是一個老外.
即是一個外國人.
很高大的女士 西人.
所以很多人看著他很有趣.
有時候有些人說.
這個鬼婆也來這裡吃的.
然後他就說 你是不是在說我.
原來這個人很聰明.
有一次.
他在廣州吃魚蛋粉.
他有感動地.
和對面不認識的人傳福音.
但怎樣傳呢.
他就說 這些東西好不好吃.
那個人說 好吃啊.
他會說廣東話.
然後說 小心吃啊.
然後就啃死了.
很奇怪.
然後就問他死了在哪裡.
然後就開始說福音.
結果神使用他.
那個人信了主.
如果你很想學.
我都祝福你可以試試.
一會吃飯的時候這樣做.
不過我不擔保結果會是這樣.
他一直都很熱心傳福音.
但到他有困難的時候.
他有一段時間想回天堂.

$^{441}$結果有人鼓勵他.
他重拾說.
我繼續在這個病房.
繼續在病房裡面.
醫生 護士 病人.
我作見證.
和他們傳福音.
不會有一個人到一個地步.
是沒有使命不用給神用的.
你今天無論在你所住的街坊.
在你親人當中.
都可以有神的使命.
讓他認識神.
在教會裡面.
你出席聚會 走上這裡的樓梯.
你本身就是一個見證.
這裡周圍有很多街坊未曾認識主.
無論買菜還是喝茶.
都可以成為一個見證.
最大的問題是往內看.
當我們在自己的時間.
我不行了 什麼都不行了.
你這裡已經將它關上之後.
就以後聽不進去.
你看我們的樂曲隊.
很厲害.
將他們所學的 所明白的救恩.
和他們的中國文化連在一起.
來作見證.
但是現實一點說.
使命這麼大.
我只是一秒.
其實我可以怎樣做呢.
不如你看看神怎樣建立約書.
他怎樣能夠繼續完成使命.
用神的方法跟從他.
我現在有個錦囊和大家分享.
以下經文我發覺.
當我再讀的時候.
原來很寶貴.

$^{481}$我們平時未必有留意到.
第五節到第九節.
我不知道你讀的時候.
或現在讀的時候.
有沒有感覺就是很多重複的說話.
有沒有感覺.
然後你覺得好像說來說去.
都差不多.
如果這樣想.
我想請你跟我一起.
去欣賞這段經文.
一個非常美麗的結構.
這段經文有兩處地方.
開始的時候第五節.
我怎樣與摩西同在.
亦都與你同在.
結束亦都是.
而會說第九節.
而會說你的神必與你同在.
一頭一尾.
開始結束都在說.
神會與你同在.
然後再請你再看下.
接下一個因素是.
第六節 你當剛強壯膽.
後面再數上來.
第二個因素都是.
你當剛強壯膽.
跟著.
你再看下.
數來第三個因素.
到了.
這個第七節.
可以凡事都可以順利.
然後第八節.
你凡事都順利.
中間就放了一句.
中間就放了第八節.
請一起讀出第八節.
這律法書.

$^{521}$不可離開你的口.
總要晝夜思想.
好使你謹慎.
遵行這書上的.
所寫的一切.
或者我將它.
將它.
這樣排列方法可能更容易.
看得到.
神比約書亞能夠完成.
這個使命的錦囊.
透過精密的結構.
要凸顯中間最重要的一件事.
開始.
神應許約書亞同在.
結束.
神應許約書亞同在.
然後鼓勵他因為神的同在.
所以他當剛強壯膽.
不是自己夠大膽.
是因為神同在所以大膽.
同樣地結束的時間.
神鼓勵約書亞剛強壯膽.
不要看自己渺小.
也不要看困難很大.
因為神比一切都大.
然後當你有這樣的態度的時候.
神說你會凡事順利.
凡事順利.
但是有一個條件就是中間那裡.
這個律法書.
不可以離開你的口.
你要常脈想.
就是說.
外面你要熟讀.
裡面要思想.
更加重要的.
聖經不只是用來.
研究和讀.
聖經是進行的.

$^{561}$就好像禮堂這兩個對聯.
是.
真理進行主的讀.
神真的要實踐.
當你這樣做的時候.
這樣結構就表達了.
要凸顯中間那部分.
錦囊就是如果我能夠.
紀念上帝的說話.
脈想上帝的說話.
遵行上帝的說話.
我就能夠完成這個使命.
摩西的偉大不是因為他學了埃及人一切學問.
也不是他很能打.
結果埃及人差點死了.
摩西的偉大是因為.
和約書亞的要求是完全一樣的.
有神的真理.
回應神的真理.
和遵行神的真理.
所以如果我們能夠像摩西那樣完成使命.
不是因為我要學一些東西.
或者我什麼恩賜.
或者問題大小我能夠處理.
如果我照神的方法.
能夠完成神給我的使命的話.
你和我和摩西和約書亞一樣.
就是聽神的話.
也要遵行神的話.
反過來說.
為什麼我不能夠完成了.
上帝給我的使命呢.
為什麼我生活裡面.
不能夠成為一個認真的基督徒呢.
不是因為我學歷的問題.
年紀的問題.
環境的問題.
是因為我對神的話的回應.
所以你看見第八節.
他這個法書不可以離開你的口.

$^{601}$當時不是像我們手拿聖經那樣.
很難才有神的話.
是靠一代一代.
一句一句.
甚至一卷一卷的書來背誦.
不可以離開你的口.
當然對約書亞來說.
就是摩西念經.
不單止像蓮藕說.
還要思想.
將神的話成為生命的一部分.
是透過一個默想的過程.
這句話.
這個靈修的時候聽到.
然後我將它變成.
對我現在.
對我現在實際的生活.
對我的性格.
對我處事為人的方法.
神有什麼話跟我說呢.
然後默想之後.
不是停在想那裡.
我接著要走的.
我嘗試去做.
雖然周圍環境挑戰很大.
我的性格跟神的吩咐.
好像格格不入.
但我願意去做.
我學習去做.
譬如我很生氣一個人.
不知道你有沒有呢.
有沒有人你很生氣呢.
你跟他談的.
幾十年都記住在心裡的.
老死不相往來的.
有沒有啊.
不用舉手.
有沒有啊.
有沒有這樣的人啊.
但如果你聽到神說.

$^{641}$你一定知道是不對神的意思.
是,我很生氣他.
但神叫我這樣做.
我想一下,我願意做.
接著我向神禱告.
神知道我不行.
你幫我做.
行不行.
這樣就將神的話變成了我的思想.
也變成了我的行為.
但如果說不行.
不用說了.
那你就停在那裡.
十一奉獻.
很敏感這個題目.
大家知不知道是聖經的教導.
但會不會很難.
一講十一奉獻.
大家差不多按著錢包.
千萬不要說.
這樣就是我們拒絕神的話.
無論我是一個.
低收入的人.
或者我是一個百萬富翁.
我都願意因為神的話.
所以我嘗試去想.
什麼困難呢.
有什麼不捨得呢.
然後我嘗試去做.
你看到這樣的基督徒的生活.
就很有意思.
是不停不停有神的話.
然後我經過思想之後.
我就去實踐出來.
你說多好呢.
我越來越像上帝的要求.
傳福音也是.
我知道神叫我傳.
但我很害怕,我不知道怎麼傳.
我又是經過上帝的說話.

$^{681}$這樣說.
為什麼我會這麼害怕呢.
然後我嘗試去做.
求神幫我.
我求神祈禱.
叫祂幫我.
這樣慢慢地神的話成為你生命的一部分.
但是不單止看見困難很大.
不單止看見自己渺小.
還有一樣很重要的.
往上看.
看上帝自己的話.
我們就會天天有精彩的人生.
最後.
你記不記得剛才怎樣稱呼摩西.
耶和華的僕人.
約書亞呢.
摩西的助手.
亂的兒子約書亞.
這裡呢.
就是剛才我看的經文.
約書亞的一章一節.
耶和華的僕人摩西死了.
約書亞呢.
就是摩西的助手.
但你知不知道經過了.
他遵行神的話.
他以照著神的說法.
去背誦.
去.
默想.
去遵行.
到了約書亞結束的時間.
約書亞的二十四章.
到了約書亞的二十四章的時候.
你猜猜怎樣稱呼約書亞呢.
二十四章九字.
我們一起讀出.
一起讀出.
這些事以後.

$^{721}$耶和華的僕人亂的兒子.
約書亞.
正一百一十歲就死了.
是不是很美麗.
怎樣稱呼他.
耶和華的僕人約書亞.
這個已經是.
跟摩西完全一樣.
我再說.
不是因為他的背景.
歷史是不會有改變的.
也可能是他恩賜不及摩西.
他的膽識不及摩西.
他的領導的才華.
不及摩西.
他的命都沒有那麼長.
比摩西還少一點.
但是.
上帝揀選人.
不是因為剛才那些.
上帝用不同的崗位.
有不同的背景去完成.
不過能夠被神使用.
一個條件.
遵行神的說話.
所以.
我希望摩西承接了.
這個使命.
在神的計劃裡面.
他走了上帝心意的其中一步.
今天我問你.
上帝有沒有給你使命.
有沒有感動你.
在洪福堂.
上帝有沒有使命給你.
不是.
我說錯了.
洪恩堂.
我看到大家的臉色有點不同.
對不起.

$^{761}$在洪福村的洪恩堂.
我應該說得清楚一點.
很明顯我對教會不認識.
在洪恩堂.
有沒有使命給你.
那個條件是.
遵行神的說話.
然後.
我一樣.
希望我們的盼望是.
將來在天家.
耶和華的僕人.
上帝所喜悅的人.
我們一起祈禱.
但願主.
你的靈在我們當中.
你所聽到的.
叫我們能夠應用在.
我們的生命當中.
我將洪恩堂的弟兄姊妹.
恭敬交在神的手裡.
可能當我們想到.
有牧者的過去.
我們心裡面很難過.
可能當我們想到.
周圍的環境.
我們沒有盼望.
甚至現在.
我們弟兄姊妹正面對很大的難處.
但願聖靈對我們說.
說話.
摩西死了.
若是後來興起.
同樣地.
求你叫我們繼續.
能夠傳承.
神給洪恩堂的使命.
在這裡.
我們能夠參與.
可以將這個教會.

$^{801}$成為這個社區.
成為我們親朋好友的一個明亮的燈台.
求主使用我們.
這裡當中的每一個.
奉耶穌基督的聖名祈禱.
阿們.
\newpage



\section{路加福音 23:40}
\label{sec:dSTGrtAycig}
\textbf{2024.01.28 《談敬畏》 路加福音 23\_40 (和修版) 陳韋安牧師}
\newline
\newline
連結: \href{https://youtube.com/watch?v=dSTGrtAycig}{\texttt{https://youtube.com/watch?v=dSTGrtAycig}} ~~~~ 語音日期: 2024-01-30
\newline
\newline
\hyperref[sec:sZSvUc8jKA0]{\small{< < < PREV SERMON < < <}}
~
\hyperref[sec:index_chronic]{\small{[返順時目]}}
~
\hyperref[sec:index_scriptual]{\small{[返順卷目]}}
~
\hyperref[sec:HRSr7qaJnjw]{\small{> > > NEXT SERMON > > >}}
\newline
\newline
路加福音 23:40
\newline
\begin{longtable}{cl}
\hline
\hline
章節 & 經文 (和合本修訂版)\\
\hline
23:40 & \begin{tabularx}{0.7\textwidth}{X} 另一個就應聲責備他，說：「你是一樣受刑的，還不怕神嗎？ \end{tabularx} \\ \\
[1ex]
\hline
\hline
\end{longtable}
$^{1}$(陳志全) 定智媒,早安.
今天我們談敬畏.
先說一個片段.
有一次跟一位屬靈教會的前輩吃飯.
大家聊聊天.
聊聊天的時候,他突然問我.
你覺得崇拜是甚麼.
突然問了這個問題,也知道有些事.
崇拜是甚麼.
原來他有一次看了我禮拜六.
開一間比較年輕的教會流堂.
其中一個崇拜.
那次我們有定智媒一起.
那次是開學崇拜.
我們定智媒一起穿著以前的校服回來.
一起崇拜.
大家大部分穿著以前的校服回來.
一起崇拜.
主題有點意思.
不試探討崇拜學的禮儀問題.
不過那位前輩有一個很好的提醒給我.
我覺得很深刻,也覺得很有得著.
他說:陳慧安,你要學習敬畏.
只要你敬畏神,你怎樣走都不會錯.
這句說話我覺得是非常好的一句說話.
我不是一個很傳統的牧師.
平時我裝扮,在流堂也不是穿西裝,戴領帶.
所以今天跟大家一起思考這個題目.
有關敬畏的題目.
不知道你有沒有這樣的感覺.
當你聽到敬畏這個題目的時候.
其實這個題目既熟悉又陌生.
敬畏這個字,英文是fear of God.
一直都知道這個字眼.
但不知道如何去理解這句說話.
如果你在網上搜索fear of God.
其實第一個看到的是什麼?.
其實是一個時裝品牌.
你會看到很多名貴的衣服.
一萬元一件衣服的名貴品牌.

$^{41}$其實這個fear of God的品牌的創文.
本身都是懷著信仰的理由.
都想讓人可以去敬畏神.
誰知道就成為一個越來越貴的品牌.
人們就追名牌,就追了這個.
敬畏在這個年代裡面.
都很容易被人忽略.
因為我們都不是經常提及.
特別現在是一個趨勢.
我們教會裡面都很講安慰.
我成為一個安慰的宗教.
上帝憐憫,恩典拯救.
恩典太美麗,耶穌很愛你.
當我們知道耶穌上帝很深深地愛我們的時候.
敬畏這個字我們反而越來越少提及.
如果你去書店買一些屬靈書籍的話.
你都很少看到有關敬畏的書.
相反就叫你不要害怕.
不要懼怕,愛裡面沒有懼怕.
要當強壯但不要懼怕.
所以面對海量的恩典念文.
愛的時候.
我們都不是很拿捏到.
究竟我們如何去敬畏.
fear我們的上帝.
所以今天剛才讀的那段經文.
其中一段.
在十字架的兩旁裡面.
其中一個強盜就去問一個強盜.
他說你還不怕神嗎?.
一個臨死之前的一個凡人.
一個值得我們去思考很嚴肅的題目.
這個題目都成為我們今天的思考.
今天不是一個經文講道.
今天我們會去看看整本聖經裡面.
如何去談及敬畏這個題目.
敬畏這個字在整本聖經裡面.
如果是中文的話就出現了191次.
舊約裡面的敬畏這個字.
就叫做Yarat Elohim.

$^{81}$這個就是以色列人一個很基本的信仰基礎.
申命記這樣提.
我要叫他們聽見我的話.
使他們存活在世的日子.
可以學習敬畏我.
又可以教訓兒女將行.
學習敬畏上帝.
是以色列民信仰教育一個很重要的傳承.
不過我們問以色列人如何去敬畏呢?.
如何算敬畏神呢?.
所謂敬畏神對他們來說.
其實就是一個人不去任意妄為.
而是無形中有一個更加高的上帝在他的生命裡面.
所以去到兩個時期.
這種敬畏神是慢慢消失我們中國人所謂敬天的概念.
雖然你未必能夠見到那位的神.
但是在你的生命裡面.
你在頭頂上就好像有位無形中的上帝.
在影響著我們的行為.
所以我們中文也這麼說.
論語裡面說君子有三畏.
畏天命,畏大人,畏聖人之言.
當中那個畏天其實正正都是一種這樣的態度.
所以你發覺因為這種態度.
在經文裡面,特別是舊約裡面.
你會見到一些很特別的故事.
《出入》的第一章是這麼說的.
提到有兩個修身婆的故事.
他們聲音說去修身婆去敬畏神.
不按照埃及王的吩咐行.
竟然就存留男孩的性命.
聲稱這兩位修身婆正正因為是敬畏神的緣故.
她就冒著危險.
甚至乎是不去遵行王的吩咐.
即是公然抗命.
這樣來違反掌權者的要求.
十八節提到你們為什麼要行此一事.
存留男孩的性命呢?.
明顯地他們去存留男孩的性命.
就是因為他們敬畏神.

$^{121}$所以敬畏神成為了一個很巨大的約束.
叫人有勇氣去堅決不做一些事情.
這是其中一個舊約裡面的例子.
另一個相反的例子正正就是亞伯拉罕的例子.
亞伯拉罕勸以撒的故事可能大家都熟悉.
不過可能大家都覺得這是一個信心的故事.
不過如果看回經文的話.
當耶和華去考驗亞伯拉罕之後.
他這樣對亞伯拉罕說.
他說你不可再藉同以身上下手.
一點不可害他.
現在我知道你是敬畏神的了.
所以原來亞伯拉罕勸以撒的故事.
除了是信心的故事之外.
其實同時間也是一個敬畏的故事.
因為亞伯拉罕他敬畏神.
他就堅決去做一些上帝要求他做的事情.
所以發覺修身婆的故事.
或者亞伯拉罕的故事.
正正是說了我們敬畏的兩面.
一個是你因為敬畏的緣故.
堅決不做一些上帝不喜悅的事情.
同時間因為敬畏的緣故.
我們都會堅決去做一些上帝要我們做的事情.
不過我還沒有回到問題裡面.
對於猶太人來說.
這個敬畏的畏是畏什麼呢?.
究竟他要怕什麼?.
他要恐懼什麼?.
究竟這兩個修身婆.
他們究竟是怕上帝什麼?.
或者亞伯拉罕是怕上帝什麼?.
難道這兩個修身婆.
是因為恐懼上帝會擊殺他們?.
相比起法老的刑罰.
他更加怕上帝.
所以兩害取其輕.
寧願得罪法老都不想得罪神?.
是這樣的緣故嗎?.
難道亞伯拉罕都是出於恐懼.

$^{161}$純粹因為怕上帝會擊殺他.
所以寧願自己的兒子死.
都不想自己死.
沒理由的.
這個不是聖經所說的那種敬畏.
以色列人所說的敬畏.
我告訴你.
我們fear of God.
這種敬畏.
不等同於fear of punishment.
不是純粹因為怕落地獄.
純粹因為怕上帝會罰.
所以才來敬畏他.
大家都是.
大家不是純粹因為怕上帝會罰你.
所以今天上崇拜.
怕上帝會落地獄.
所以就做基督徒.
不是這樣的.
我們純粹是因為什麼?.
因為上帝的救恩.
上帝的愛去吸引我們.
往往都是一種pull factor.
多於一種push factor.
不是出於懲罰.
而是上帝的愛.
可能大家不知道有沒有試過.
我自己初信主的時候都這樣想過.
有時候自己肚子痛的時候.
在洗手間蹲下來都會想.
是不是自己做錯了?.
是不是上帝會這樣罰我?.
是不是因為我最近大病.
因為上帝所以要警戒我?.
可能大家都這樣想過.
不過我想長遠來說.
大家都不能夠靠這種想法來做基督徒.
純粹因為擔驚受怕.
怕自己被上帝罰.
怕自己將來會覺得不好.

$^{201}$所以今天才來做一個好基督徒.
所以不是出於懲罰.
也不是出於一種功利.
而是另一種原因.
所以這種消極的恐懼.
不是我們所說的敬畏.
正如這樣的敬畏.
對於我們以前一位古代的神學家教父.
Huston所說的.
這種敬畏,害怕敬畏.
是叫做奴僕的恐懼.
是我們基督徒一開始早期的看法.
純粹因為怕上帝會罰.
因為律法,懲罰.
我們才願意做一些事情.
但Huston認為當我們基督徒.
越來越經驗上帝的愛的時候.
這種敬畏就會慢慢被取代.
正如聖經所說.
愛裡沒有懼怕.
愛的完全就怕懼怕除去.
出而代之的就會出現一種昇華了的敬畏.
這個敬畏不再是純粹因為懲罰.
而是同時間包括了上帝的愛在我們的生命裡.
這種敬畏.
因為上帝的愛的敬畏.
正正是我們所說的那種敬畏.
事實上.
愛上帝和敬畏上帝是不可分割的事情.
Love of God和Fear of God.
其實是源於同一件事情.
在撒相辯中.
中文就叫做令人可畏而又吸引的奧秘.
拉丁文我就不說了.
因為叫做令人可畏而又吸引的奧秘.
就是說我們正正是所說的話.
因為上帝那種偉大.
上帝的慈愛,上帝的恩典,上帝的憐憫.
是完全傳來震懾我們.
所以令我們同時間產生一種既是羨慕.

$^{241}$同時間又是震驚的那種感受.
這個敬畏正正是超越我們人世間一種情緒.
就是恐懼和驚喜.
是同時間湧現一種很複雜的情感.
就像我們唐伯虎點秋香所說的.
我的心又喜,我的心又慌.
就是喜和慌是同時都會一起出現的.
當你真的去感受到上帝那種高度.
祂的尊榮,祂的崇高.
上帝不需要穿任何東西.
祂的權威,那種莊嚴和肅立.
都令我們很震懾.
但同時間又不是那種怕佛的震懾.
而是祂的慈愛和偉大是同時間的出現.
所以這種恐懼和驚喜的交集.
就像詩篇42篇所說.
詩人說:你的瀑布發聲,心冤就與心冤響應.
你的波浪洪濤漫過我身.
八州耶和華必向我詩詞哀.
黑夜我要歌頌禱告賜我生命的神.
上帝的瀑布強大地告訴我們祂的力量.
上帝的心冤去淹沒詩人十萬八千尺裡面.
上帝的波浪就像巨型海嘯一樣掩蓋了整個詩人.
上帝的氣場過盛.
我們被上帝的慈愛淹沒.
我們被上帝的同在完全包圍.
這種複雜的偉大又慈愛的感受.
正是我們面對上帝的時候.
那種叫人去既震驚又羨慕的傲備.
所以我們作為基督徒.
不能夠失去這種戰略和震驚.
作為新教徒.
你不會出於恐懼.
不會怕將來會落地獄.
不怕將來會受到什麼刑罰.
因為相信我們恩信稱義.
我們因為信的緣故.
上帝已經拯救我們.
但我們又不會單單過份去強調.
上帝呵你,疼你,哄你.

$^{281}$而忘記了上帝和我們那種距離.
和那種肅立,起勁的態度.
所以當我們面對上帝的時候.
我們仍然要保留最基本而不可缺少的第一步.
正如以下書所說.
你要以敬畏而無華為之寶.
這是我們世上最基本的一點.
也是很重要的第一點.
你永遠都不能夠丟棄最基本的第一步.
大家有沒有聽過《奇恩典》這首詩歌.
這首詩歌的背景大家知不知道.
其實是1748年.
大概幾百年前.
這首詩歌的詩人John Newton.
面對著一場在海裡面非常大的暴風雨.
幾百年前一艘船去大西洋.
突然面對著一個非常可怕的暴風雨的時候.
John Newton就向上帝去呼求.
後來他就將這個經歷寫成《奇恩典》這首詩歌.
如果你看回這首詩歌的歌詞第二節的時候.
一個非常有意思的一段英文.
他說.
This is what grace that taught my heart to fear.
And grace my tears relieved.
How precious that the grace appeared.
The hour I first believed.
當中提到這個恩典.
是叫我們心裡面去敬畏.
The grace that taught my heart to fear.
當然我們知道.
恩典永遠都是我們信仰裡面的主線.
恩典是最重要的.
不過這個恩典卻同時間不是和這種敬畏分開的.
沒有了敬畏就沒有真正的恩典.
這種恩典永遠都是這種敬畏戰略裡面.
作為我們的背景.
沒有了這份敬畏的高度.
任何很親密的恩典就失去了它最重要的意義.
所以這都是猶太人信仰的價值.
一切耶和華的信仰.

$^{321}$都必須建基於這個敬畏的起點裡面.
所以猶太人給了一句很有智慧的話.
如果我們沒有了對於上帝的敬畏的話.
就好像你得到一顆寶藏裡面的寶箱.
卻沒有了一條鎖匙一樣.
一切東西都有了.
今天你能夠有很多聖經知識.
有很多基督徒應有的東西.
但如果沒有了這份最基本的敬畏的時候.
其他東西等同於沒有一樣.
所以這種敬畏對我們來說.
今天有什麼東西值得我們去思考呢?.
最基本的就是敬畏去教導我們.
我們其實要知道我們在世界裡面.
是要怕什麼?不怕什麼?.
正如老方所說.
耶穌說:那殺身體以後不能再做什麼的.
不要怕他們.
我要指示你們當怕的是誰?.
當怕那殺以後又有權柄丟在地獄裡的.
我實在告訴你們.
正要怕他.
五個麻將不是賣義勳銀子嗎?.
但在神面前一個也不忘記.
就是你們的頭髮也被數算過.
不要懼怕.
你們比許多麻將還貴重.
耶路經文裡所說的.
不單單是上帝像看顧麻將般看顧我們.
更加是什麼前提呢?.
正是叫你不要害怕其他的東西.
面對今天令我們很害怕的東西.
很多的病毒.
很多的擔心.
甚至乎很害怕自己踩界.
很多不同的令我們很擔心的事情.
但聖經說.
當你要知道你真正要怕的是什麼的時候.
其他的東西.
你就發覺你不需要去怕.

$^{361}$不知道大家怕不怕鬼?.
很多基督徒都很怕看鬼片.
或者來到去很怕鬼.
你只要懷著你怕神的話.
其實就應該不會怕鬼.
因為你知道.
當你更加要怕的是一個愛你的神的時候.
你就懷著這份敬畏.
其他鬼都不敢來教育你.
因為你更加敬畏的是神.
面對這個社會也一樣.
令我們今天有很多令我們很懼怕.
很害怕.
很擔心的事情.
無論是我們的前途.
或者是這個社會也一樣.
當我們真正捉到我們要敬畏的.
我們反而有勇氣.
面對著我們不同的情況.
當然面對著社會世界的議題.
更加令我們重要的.
就是提醒我們.
敬畏叫我們回到敬畏的初心裡面.
正如我所說.
敬畏是我們信仰的第一步.
它是關鍵的第一步.
當我們經歷了這一點之後.
當你爬上去之後.
你不能爬上去之後.
它是丟棄這個第一步.
不知道你最敬畏神的時候是什麼時間.
是你剛剛信耶穌的時候.
還是什麼時候呢?.
很多人都是剛剛信主的時候最敬畏神.
拿著本身買聖經.
很敬而有之地逐頁逐頁去看和揭.
教會裡面的聚會.
每一件事都是非常認真嚴謹地去面對.
可能那個時間.
雖然你很多事情都不懂.

$^{401}$但它是你最敬畏神的時間.
受孝院也是一樣.
我自己發覺.
我在受孝院裡面觀察.
受孝山也是這樣.
通常最敬畏是哪些同學呢?.
一年級的同學.
剛剛來到受孝院讀書.
就很敬畏神.
恭恭敬敬.
跟隨受孝院的規矩.
然後就很好的去讀書學習.
他們的敬畏程度就會逐年減少.
跟他們自家的功課.
開始不守規矩.
所以我們很重要.
懷著一種我稱之為菜鳥的敬畏.
回到我們信仰開始認識神的時候.
那種既是發現上帝很愛你.
同時間也是很認真尊重這份信仰.
上帝叫你做什麼.
我們就敬畏地去做.
所以敬畏是一件很重要的事.
可能你今天已經長大了.
甚至年紀也很成熟.
很多時候我們都會越來越少了這份敬畏.
因為我們覺得自己做什麼都可以.
自己一個人這麼大.
沒有人管你.
也沒有人可以自己做很多事情.
但我想仍然的來到去.
帶著這份敬畏.
我們基督教沒有任何神聖的物品.
沒有聖物.
也沒有聖像.
沒有聖神聖的地方.
但我想我們要好好地來到去.
在生命裡面保留著一塊.
我稱之為神聖的石頭.
這塊神聖的石頭正正在你的心裡面.

$^{441}$你要好好保護這塊石頭.
在你的工作裡面好好地帶著它.
在你的家裡面好好地保護著它.
無論是很多人還是自己一個人.
仍然在心裡面存留著一塊神聖的石頭.
在你的生命裡面.
當你拿著這塊石頭的時候.
正正好像那位前輩跟我說.
只要你敬畏神.
你怎樣走都不會走錯.
你無論去到哪裡做什麼.
帶著這塊神聖的石頭的時候.
我們就能夠很自由地.
懷著勇敢地去做很多不同的事情.
敬畏耶和華是智慧的開端.
正正大概就是這個意思.
我一起祈禱.
因為你給我們一份很大的厚愛.
但我們知道你是那位偉大的神.
一位超越我們崇高致敬的上帝.
你是跟我們很近.
他同時間是天上的父神.
我們求主你幫我們學習去敬畏你.
將我們生命裡面每一件的事情.
我們都知道你在我們的生命裡面.
有我們這份敬畏去管理我們的行為.
有我們去堅決不去做那些你不要我們做的事情.
去堅決去做那些你吩咐我們做的事情.
有我們去敬畏你.
有我們這個.
本堂這個教會都是一個敬畏的群體.
我們一起彼此去學習敬畏.
實踐你給我們這份信仰裡面的功課.
讓我們能夠建基於這裡.
成為一個你所希願的教會.
奉主名求.
阿們.
\newpage



\section{出埃及記 33:12-17}
\label{sec:HRSr7qaJnjw}
\textbf{2024.02.04 《認識神的摩西》 出埃及記33\_12-17 (和修版) 樓恩德牧師}
\newline
\newline
連結: \href{https://youtube.com/watch?v=HRSr7qaJnjw}{\texttt{https://youtube.com/watch?v=HRSr7qaJnjw}} ~~~~ 語音日期: 2024-02-04
\newline
\newline
\hyperref[sec:dSTGrtAycig]{\small{< < < PREV SERMON < < <}}
~
\hyperref[sec:index_chronic]{\small{[返順時目]}}
~
\hyperref[sec:index_scriptual]{\small{[返順卷目]}}
~
\hyperref[sec:IFtn6wInBFk]{\small{> > > NEXT SERMON > > >}}
\newline
\newline
出埃及記 33:12-17
\newline
\begin{longtable}{cl}
\hline
\hline
章節 & 經文 (和合本修訂版)\\
\hline
33:12 & \begin{tabularx}{0.7\textwidth}{X} 摩西對耶和華說：「看，你曾對我說：『將這百姓領上去』；卻沒有讓我知道你要差派誰與我同去。你還說：『我按你的名認識你，你也在我眼前蒙了恩。』 \end{tabularx} \\ \\ \relax
33:13 & \begin{tabularx}{0.7\textwidth}{X} 我如今若在你眼前蒙恩，求你將你的道指示我，使我可以認識你，並在你眼前蒙恩。求你顧念這國是你的子民。」 \end{tabularx} \\ \\ \relax
33:14 & \begin{tabularx}{0.7\textwidth}{X} 耶和華說：「我必親自去，讓你安心。」 \end{tabularx} \\ \\ \relax
33:15 & \begin{tabularx}{0.7\textwidth}{X} 摩西說：「你若不親自去，就不要把我們從這裡領上去。 \end{tabularx} \\ \\ \relax
33:16 & \begin{tabularx}{0.7\textwidth}{X} 現在，人如何得知我和你的百姓在你眼前蒙恩呢？豈不是因為你與我們同去，使我和你的百姓與地面上的萬民有分別嗎？」 \end{tabularx} \\ \\ \relax
33:17 & \begin{tabularx}{0.7\textwidth}{X} 耶和華對摩西說：「你所說的這件事，我也會去做，因為你在我眼前蒙了恩，並且我按你的名認識你。」 \end{tabularx} \\ \\
[1ex]
\hline
\hline
\end{longtable}
$^{1}$吾皇萬歲萬歲萬萬歲.
主席.
信耶穌.
我相信這三個字.
全世界人都應該在做的.
又不是,很多人還沒信.
沒信他就不會放聖誕假了,是嗎?.
給神的愛激勵感動.
讓我們每次在聖餐的時候.
有反省的機會.
神的愛能讓我們感受到.
神愛世人.
但是我們到底知道有耶穌的降生.
我們可以說有些經歷,他愛我們.
但是我們有多少人相信耶穌已經復活.
留心,《仕途信經》最後一句是.
我信身體復活.
誰的身體?不是耶穌的身體,是誰的身體?.
是你的身體.
這點我相信很多弟兄姊妹沒想過.
我怎麼知道我身體會復活?.
很簡單,如果耶穌沒復活.
你和我都不用知道會復活.
耶穌在歷史上已經復活了.
那你復活了會看到他身體嗎?.
他在哪裡?在太空?不用穿太空衣嗎?.
他的身體已經是超自然,超時空.
但又實體真正的存在.
這顯然不僅僅是他復活之後.
當他登山變貿的時候,變象的時候.
已經展示了一次給誰看?.
給德約翰雅圖看.
當時和他一同出現的還有誰?.
還有摩西,爾利亞.
問題是很多人不信耶穌復活.
所以你和我都會怕死.
但記住一個怕死的人,一定會死.
不怕死的人都會死.
不過死了之後會不會復活?.
就是另外一件事.

$^{41}$詩篇23篇第4節.
他使我的靈魂蘇醒.
為自己的名引導我走異路.
我雖然行過死刃的幽谷.
沒有一個人不會死.
但遭不遭害?一定遭害.
但下面他說也不怕遭害.
因為你與我同在.
你有耶穌同在就不會遭害.
沒有耶穌同在就一定死.
死到永遠,滅亡.
清楚了嗎?.
所以基督耶穌給我們的信息告訴我們.
天國不是死後想靈魂搬家搬到那裡永恆.
不是.
基督的降生,基督的受死,基督的復活.
已經發生了.
是事實,是歷史,是無法改變的.
還有一件事,教會很少提及.
我想你幾年都沒聽過.
基督的升天.
你何時聽我講基督升天的道理給你聽.
很少講.
他不升天他怎麼會回來.
這是一個邏輯.
他還要再來.
所以耶穌升天,耶穌還會回來.
這是我們信仰福音的傳辟福音.
那些房子應該有五支裝.
如果裝了兩支半.
你說黑色暴雨的時候.
衝了山泥.
那些房子危不危險.
所以很多信徒遇到小小事件.
自然災害,人為的問題.
就怕得七彩一樣.
怕什麼.
你走過死人的幽谷.
一定會遭害.
不過詩人已經很清楚地說.

$^{81}$也不怕遭害.
下面怎麼說.
說啊,你不說就糟了.
下面說下去.
因為.
誰是你啊.
專科醫生啊.
你的親人啊.
不是,是耶穌.
是這位基督.
不要只說神的字.
拜神都叫神.
是不是.
伊斯蘭人說他是神,真主.
天主教又說他是神.
個個都說他是神.
我神很偉大很偉大.
但是偉大的神你不相信他就是基督.
就是道成肉身的那位耶穌.
猶太人今天的慘況就是這樣.
他到今天都不相信.
還有一段時間的.
只要他不承認耶穌是基督.
大災難一定會淋到他們身上.
真的說的.
今天和大家看這個題目.
認識神的摩西.
很簡單.
我知道我時間有限.
所以30分鐘其實不夠我講.
你知道我沈氣.
再說.
不過你約了人吃午餐怎麼辦.
不過我講件事實給你聽.
我今天只講一個主題.
你認不認識這位神,這位耶穌.
最重要的一件事剛才讀經.
不是你認不認識他.
不是你知不知道他.
你認識他,他認不認識你.

$^{121}$今天你都不知道他認不認識你.
將來最怕一件事就是你以為.
罕然無據去見他.
耶穌怎麼說.
我不認識你.
不是面不面蒙的問題.
是永恆的問題,弟兄姊妹.
我舉一個例子.
這個例子不一定是金科玉律.
如果這本聖經全是講耶穌的.
你都沒有看過.
你認識多少人.
我是說你識字的人.
你可以看很多手機資訊的人.
你財經界是專科.
醫學界是專科.
你識很多很多東西.
但你不認識耶穌.
你可以很多方面成功.
但我恐怕一件事.
今生你活得精彩.
不過永生如果悲慘可憐的話.
人生會是什麼.
信耶穌得什麼.
不是人生嗎.
耶穌說的.
我來了是要叫人得生命.
活得豐盛.
只要我這生命活得光光彩彩.
這樣就行了.
光彩不光彩都走過一條路.
最後床位都是兩尺半.
三六尺長.
動一動手就不是了.
動一動手就化灰了.
一尺多二尺高.
我不是說.
新年快到了你說什麼就什麼.
如果言語不好.
我送幾個灰塵給你.

$^{161}$很棒的.
弟兄姊妹.
我講的是真實.
一到我這個年紀.
我真的想這些問題.
我曾經和我的孩子商量.
綠化了就算了.
不用這麼笨.
睡在這裡不睡在那裡.
將來你哥哥不在了.
還找人來搞事.
扔了我下海.
以後你游泳就記得我了.
是不是.
這個是題外話.
今天和大家看這個話題.
認識神的摩西.
我想透過歷史上發生的事件.
以色列人本來要大遭毀滅.
如果上帝只要發脾氣.
不用上帝出手.
以色列人全部死在曠野.
但摩西一個土狗.
要說什麼.
你認識我.
我都認識你.
原來摩西給上帝的認識.
是透過祂的祈禱.
你認識神.
都是透過祈禱.
你經歷過神的大能作為.
你有信心.
相信神的話.
你禱告上帝應允.
你做見證.
你沒有經歷過.
你經常覺得這個可以.
這個可以.
為什麼我不可以.
因為你的禱告.

$^{201}$未達到上帝的標準.
你為你人生禱告.
你很少為你永生禱告.
一個永生禱告的人.
是一個榮耀神的人.
一個為人生禱告的人.
只是滿足自己的私慾.
年三十晚有很多地方.
有人排隊爭頭位.
他求什麼.
我想教會都可以這樣教.
整個年三十晚.
誰進來第一個祈禱的就.
可以的.
不過最重要的是.
你心裡不是這樣想.
耶和華神知道摩西的禱告.
認識他.
為什麼呢.
和大家看今天的經文裡面.
我剛才有講過.
或者講一些背景.
我今天用這段經文.
是摩西在他第一次上山.
上帝給他.
上帝親自做了兩個石板.
寫了字給他.
但是當摩西經過四十天之後.
從山上拿著石板下來的時候.
發生了什麼事.
上帝寫字說.
除我以外不可有別的神.
摩西一到山谷中間.
下面菲利普拉拉吵得很大聲.
若孫雅娜說下面有天神.
不是啊.
很歡樂啊.
真的有神.
在拜神.
在拜什麼神.

$^{241}$你又知道.
有看聖經的.
我不會拜金牛的.
你不會拜金牛.
那只牛可能不是金的.
是銅做石做的.
在中環廣場那只牛在那裡.
你知道我說什麼.
不知道就算了.
不是中環廣場.
就是中環的兩層大廈.
中間有一隻牛在那裡.
你會不會拜那些.
心中那些.
當時神在山上頒布命令.
告訴你.
你要這樣這樣這樣守.
不可這樣不可那樣.
第一就說不可以有別的神.
不可以為自己雕刻偶像.
他們就這樣在做.
這幫人不是上帝作難他們的.
在埃及430年.
400年的時間.
做奴隸這麼生活.
還有.
當時希伯來人在埃及地方.
幾乎是種族滅絕的.
為什麼.
凡是女人生小孩.
一生了男孩就要弄死他.
所以摩西尖被他姐姐扔到水里.
這個真的是種族滅絕的.
那怎麼辦.
於是在那個地方叫苦連天.
求高爺說個名.
一求高爺說個名的時候.
上帝興起摩西.
不是就這樣.
找運輸機搬他出去.

$^{281}$行了十個神跡.
將埃及人打得幾乎慘了.
過紅海.
到曠野地方.
沒東西吃.
草原巴閉.
上帝給他吃.
晚霞.
吃銀春.
還有.
沒水喝.
給他喝水.
不是一律.
不是同軍去旅行.
有百萬大軍在那裡.
這些是經歷神跡.
大家知道嗎.
在你過去日子中間.
有沒有經歷過上帝的供應.
有沒有體會過上帝的同在.
不是信心大小問題.
就是你信不信得那些神器.
我們日常飲食.
今日似給我們.
以前中國人過年.
我聽人說.
那些米缸要寫滿字在那裡.
沒得吃.
現在要寫滿字嗎.
不用寫了.
米都不吃了.
你看看.
所以主要是問你背後.
背後說話.
詩圖信經也是這樣背下算數.
所以我背到最後一節.
我信復活.
復活.
誰復活了.
牧師復活了.

$^{321}$你會復活了.
今日頂智回.
當摩西下到山上.
見到這樣的時候.
摩西是一個.
很火爆的人.
年輕的時候在埃及.
習武習文什麼都會.
一見到不平.
馬上打兩拳.
兩個人死了.
一個人死了.
是埃及人.
第二次死了.
自己人.
於是.
他要逃命.
上帝等他曠野.
四十年的時間磨練他.
不知道弟兄姊妹.
你過去日子中間.
有沒有被上帝磨練你四十年.
四十日你都已經叫苦連天了.
四十年磨到火氣都沒有了.
真的火氣都沒有了.
我做了二十五年校長.
十五年在牧會.
加起來四十年.
我現在火氣真的少了很多.
真的.
不是沒有.
就是不敢發.
因為血壓會高.
但是摩西都忍不住.
拿著上帝給他的那份東西.
人家送他很珍貴的禮物.
讓你親手寫給你.
你會不會在那裡.
打爛他.
當他打爛了之後.

$^{361}$上帝的憤怒都臨到.
當時已經擊殺了很多人.
我們剛才讀經文就是.
你再上來.
我要告訴你.
接著說.
因為上帝應許了.
要將希伯來人.
猶太人.
或者伊斯蘭人.
離帶他們離開.
進入迦南.
上帝應許的話.
一定要成就.
你用信心去接受.
你就經歷了神跡.
你不信.
你就遭受上帝的審判.
當信主耶穌.
你和你的家都必得救.
我經歷了很多次.
我看過很多見證.
只要你肯堅定相信.
到最後那一刻.
上帝都可以救他.
除非你半信半疑.
自己擔心擔愧.
我不要對他寫福音.
他不要受洗.
他又沒有受洗這件事情.
但他最後會不會呼求上帝的名.
和你的信心很有關係.
我曾經在學校工作的時候.
和工友都有定期的福音工作.
每個星期左右.
我都有和工友一起查經.
特意批準時間.
他不用工作.
他們和我查經的時候.
他們都很客氣.

$^{401}$但他信不信耶穌.
我就不逼他.
有一個工友有一天早上.
來不到學校上班.
就收到家裡人電話.
說他已經去急診室.
因為他離開家的時候.
走過路被車撞到.
撞到之後立刻送去急診室.
很嚴重.
我去醫院探他.
我立刻去醫院看他.
當我看他的時候.
他已經昏迷不醒.
那怎麼辦.
我最後在他耳邊和他說話.
我說你要信耶穌.
很奇妙的.
一個不省人事的人.
會流眼淚.
不是流淚.
是真的流眼淚.
我為他禱告.
弟兄姊妹.
人在最後那一刻間.
能不能夠接受救恩.
你不知道的.
只要他心裡願意.
認罪悔改.
只要肯接受.
肯求告主名.
他就可以得救.
所以我自己很相信.
神的應許是真實的.
我在經歷很多事情上.
沒辦法的時候.
不知道怎麼處理的時候.
我會有禱告.
當我禱告的時候.
用信心交托依賴.

$^{441}$神就一件一件事情.
帶領我經過.
你都可以經歷到.
摩西當他在這個經歷中間的時候.
和上帝對話.
對話的時候.
曾經有一段時間.
有一段經文.
聖經是這樣說的.
我剛才讀的.
初一級第33章1到6節.
我和大家看看.
經文打出來.
耶和華吩咐摩西說.
去離開這地.
你和你從埃及地領出來的百姓.
耶和神已經很生氣了.
你呀.
帶到你那些百姓.
從埃及出來.
要上到我興許亞伯拉罕.
耶和華吩咐.
我興許亞伯拉罕.
以撒下面.
和瓦割之地去.
我曾對他們說.
我要將這些地施給你的後裔.
上帝說的話一定要做到.
下去.
有沒有.
我要差遣使者在你面前.
將迦南人.
比利亞摩利亞人.
黑人.
比利時人.
耶和華人.
趕出去.
第三節.
去到羅馬爾等之地.
但我不與你們上去.

$^{481}$因為你們是硬著頸項的百姓.
免得我路上板不滅絕.
你這麼硬頸.
整天頂頸.
整天拼嘴.
是不是.
做個老師都知道這些事.
做個爸爸媽媽都知道.
說真的都頂.
一言九鼎.
你氣不氣.
下面第四節.
百姓聽完你話消息.
就悲哀.
沒有人佩戴首飾了.
第五節.
耶和華對摩西說.
你對以色列人說.
你們是硬著頸項的百姓.
我若在你們門中間一起上去.
只一瞬間就必把你們滅絕.
滅絕.
我上山上面經過這麼多日子.
幾乎幾個月的時間.
我行了這麼多神跡.
你們跟著第一件.
我說不要入別的神.
你馬上拜別的神.
你搞什麼.
很難怪.
他四百年在埃及地.
以魚目掩.
聽來聽去.
每個新年都說恭喜發財.
是不是.
你聽得多.
你都想著發財是好事.
所以吃.
找味菜都要找點發財來弄.
那有什麼好事.

$^{521}$好事發財.
意頭.
心中有偶像.
下面那句.
第六節.
他們手下留情.
我好知道怎麼處置你們.
以色列人離開河內山之後.
將他們知死.
但已經做錯.
都知死.
那怎麼辦.
下面第七節.
不要了.
第六節不要了.
所以以色列人知道闖禍.
摩西那時發了鬧之後.
那怎麼辦.
我們看採取第十三章.
上帝怎麼跟他們說話.
法老放他們走之後.
神卻不令他們走那些路.
因為神說.
恐怕白色余季.
戰爭就後來.
上帝帶領他們走的路.
都安排好了.
令他們不用一出來就遇到戰爭.
第十九節.
摩西將約瑟的骸骨帶走.
約瑟很有信心.
知道自己將來一定能走.
骸骨要搬走.
第二十節.
時間不多.
從他們起營到伊藤安營.
二十一節.
耶和華走到他們面前.
每天走路.
日間用雲柱引導.

$^{561}$夜間用火柱照亮.
日夜都照顧.
弟兄姊妹.
你有沒有時間覺得.
耶和華好像不願意同在.
他們經歷過的.
到二十二節.
夜間日間.
總不離開白色的面前.
上帝的同在是這麼寶貴.
問題是.
你看不到的.
不是.
失憶的.
忘記了.
所以弟兄姊妹.
CPU二十三篇第四節.
走過死人沒有谷.
不用擔心.
你只要常常有主願你同行同在.
你不會遭害.
但是你說.
我不用你.
我不聽你的.
我喜歡自己行事.
那你就自行己路.
說得好聽就自行己路.
說得不好聽就自行什麼路.
我沒說過你說的.
猜測記第三十二章.
我們看.
摩西對百姓說.
你們犯了大罪.
我如今要上耶和華那裡去.
或許可以為你們贖罪.
三十一節.
摩西回到耶和華那裡說.
這百姓犯了大罪.
為自己做了金的神明.
都是神.

$^{601}$不過是金做的神.
弟兄姊妹.
是你的理想.
是你的心願.
你想達到.
但上帝說不行就不行.
除我以外沒有別的神.
凡以別的代替耶和華的人.
他們的心裡仇苦就不斷增加.
你有任何東西代替了永生的神.
賜福給你的神.
拯救你的主耶穌.
我告訴你.
你不會有平安喜樂這四個字.
之後什麼.
勞苦受煩.
是不是.
用不著靠安眠藥吃睡覺.
用不著吃開心果來開心.
越吃越糟.
第三十二節.
然而靠你赦免他們的罪.
這句話很重要.
不然就把我從你寫的冊上除名.
摩西圖古到這步就是.
你赦免他們.
如果不是.
我頂替他們.
你將我除名.
這個圖古只有耶穌基督才圖古過.
摩西在那時候情遲逼切.
他真的站在神的立場.
看神的眼光去看事物的時候.
他說算了.
你赦免他們.
如果不是.
我頂替了.
這個是誰說的.
神的心意說的.
今天我們禱告就是.

$^{641}$你都不照我的心意做.
你都不是神來的.
你聽過這些見證沒有.
我聽過.
真的這樣說那些人.
我少年來的時候.
我曾經坐在那裡聽講道的時候.
我聽過一個旁邊的叔叔伯伯伯.
那時候不是六合彩那些.
馬彪.
拿著馬彪祈禱.
上帝求你賜福.
我心想你搞錯了.
可能他和上帝準備對婚.
我給你一半.
拿著馬彪祈禱.
我那時候都不懂事.
怎麼這麼有趣.
求神嘛.
等下.
不然就把我從你所寫的名冊上除名.
耶穌基督在約翰福音第十四章.
講過兩個土文.
父啊求你救我脫離這個時刻.
下面跟著說不是.
跟著說什麼.
我知道我是為這個時刻而來的.
跟著下一句.
父啊求你永耀你的名.
完了這兩句土文之後.
跟著天上有聲音.
那些人聽不到.
那些人說打雷了.
有天使說話了.
不是.
然後上帝說.
這句話是上帝回應給耶穌基督.
我要永耀我的名.
基督徒在一生人.
你的祈禱只有一件事乾.

$^{681}$是要永耀神的名.
神的旨意在你身上成就的時候.
他就得永耀.
神的旨意在你身上不能夠成就的時候.
你以為自己很有face.
我昨天下午在一個團體講話那裡.
很簡單的經歷.
我做了二十五年校長.
二十四年的時候.
我以為2003年.
我講完很快就算了.
你不想聽.
SARS停課一個月.
我的英文科主任跟我說.
校長.
停課那些成績不能夠補習.
不如我們介紹一些網上給他學習.
我有網上學習嗎.
他說有.
有個課程叫global English.
我說要不要請學家要錢.
我說多少錢.
一百塊美金.
大概.
對我那間學校.
我們在石泥區來說.
那些子弟家裡不是很有錢.
於是我就請一個推銷員來.
簡介之後.
我請一個英文科老師聽完.
都說好.
就想想.
不是了.
錢不是我給的.
不是學校給的.
是家長給的.
我請家長代表來聽.
都好.
兩個人都好.
我就出封信.

$^{721}$跟學生說.
跟家長說.
你們如果想冒免參加.
這個課程.
九個月.
一百塊美金.
我不是跟他們做宣傳.
現在有沒有我不知道.
可以從一個初學者.
beginner.
去到考Tofu.
這麼厲害.
那我沒問題.
我問了.
打聽了.
其他學校.
一間學校都是幾十人參加.
我寫一封信給家長.
就說如果你們想參加.
就自己拿入數紙.
入數.
入他的戶口.
於是他們成交成.
不關我的事.
我不可以做商人推介.
出了之後.
不到一個星期.
我問同事.
我說.
多少人參加會議.
你猜多少人會議.
我說一千二百人.
你猜多少人會議.
七百多個人會議.
我算一算.
七百多人.
大概七百美金一個人.
收了多少錢.
你心算得多聰明.
我心想.

$^{761}$哇 這麼大單.
我就隨口說.
那間店鋪在哪裡.
他說在旺角.
我說你們有沒有人去看過.
他說每次他來.
我立刻找了兩個同事.
我說你們上去看看.
他們兩個人不去.
都沒有那麼害怕.
一去完之後.
打電話問校長.
水牌沒有公司名字.
接著.
我說糟了.
那怎麼辦.
我說水牌沒有.
上去看看.
上去一看.
重大一件事.
在裝修.
美容院來的.
我好像被滾水從頭滾到尾.
死了怎麼辦.
那怎麼辦.
我第一時間回到辦公室.
坐下.
我說神啊.
這件事這麼嚴重.
我打算明年我就退休了.
2003年來到林美香.
我見一見報都死人.
真的很害怕.
你知不知道當時我得到什麼提醒.
我第一次不去怪責人.
我自己有責任.
我也有責任.
我就唯有認罪悔改.
認什麼罪.
你們不知道.

$^{801}$我做了24年校長.
學校出了很多事乾.
每次出事都解決了.
學校從班5一直升到班1.
算是這樣了.
我真的跟自己說.
我不敢跟別人說.
我怕說出你驕傲.
跟自己說.
你都算是這樣了.
當我在那時候一看到這個光景.
真是盛名光照.
你都算是這樣了.
臨了一步.
你都猜不到會出事的.
我跪下祈禱.
認罪悔改.
求神給我智慧怎麼處理.
真是智慧.
求主憐憫.
其中一個事乾.
我打電話給我校董會的律師.
律師說.
那你想怎麼辦.
因為他星期四收了錢.
星期一就把密碼要給我們學生.
他說你現在是星期幾.
我現在是星期四.
他說沒辦法了.
那怎麼辦.
他說你等到星期一再算吧.
哇.
星期五 星期六 星期日.
這三天夠熬的.
真是受難.
我打電話給一個學生在警界做.
做得挺高級的.
我問他.
他說這樣啊.
你現在報警都不受理.

$^{841}$因為星期一才.
他能不能給你貨.
你報警沒用的.
我和兩三個副校長和科長聊天.
他們打聽.
他打電話問總代理.
總代理說是有這間公司的.
不過.
他的牌到期了.
我們會不跟他作牌.
那怎麼辦.
等到星期五怎麼辦.
自己按摩包填了他.
不行.
結果.
三天時間裡面.
我前兩晚都好好睡.
星期一晚上睡覺.
到星期一早上起床的時候.
我不想起床.
我不想上班.
真的.
但奇妙的事.
星期四.
星期五我請他上來辦公室.
我和他聊天的時候.
他一坐下.
他對我說得很客氣.
校長多謝你這次.
我心想你當然多謝我.
我不知道什麼事過.
於是我和他聊天.
我說你是不是星期一.
可以派阿叔給我們學生.
他說行沒問題.
我說不過聽說你的牌照到期了.
他臉色一轉.
你怎麼會知道.
我見他這種情況.
我就單刀直入.

$^{881}$上帝讓我信他.
我說可不可以這樣.
你星期一拿到密碼之後.
他說我海東南亞拿貨給你.
我說你拿完貨之後.
沒拿完密碼之前.
你把這些錢退回給我學校的戶口.
我說的不是學校戶口.
行不行.
你以為他怎麼回答我.
好啊.
一個好之後.
我想那不行.
那張支票空頭怎麼辦.
我跟著說我想要現金.
你以為他怎麼答應我.
好啊.
這兩個好之後我沒話好說.
我一日土古代之後.
我把前件事交給上帝.
我說一定是你自己帶領我.
所以到星期一早上.
我很不想上學.
這個人性是這樣.
但叫他出現.
他出現的時候.
我工友告訴他.
我叫他回學校.
他說要回學校.
坐下把長桌放在那裡.
他拿著袋子拉著拉鍊.
捧著那些現金.
多少錢你知不知道.
五百塊錢一張.
二十五萬.
推到我面前.
校長你怎麼收這些.
二十五萬.
我數都不敢數.
看都不敢看.

$^{921}$心裡想二十五萬.
他說我還有些數.
今天下午明天早上.
他把那兩條數說了給我聽.
我心算一加.
都不是齊頭數.
那怎麼辦.
收不收.
你們這麼聰明.
損失少一點.
我立刻找兩個身體健碩的男同事.
開車立刻把錢先還給銀行.
先入數.
有學校沒有.
有些損失怎麼辦.
叫我真是感動他.
下午第二天他都入數.
不過不平.
都給我九成.
那怎麼辦.
上帝聽到告.
結果我們和總代理打電話聯絡.
總代理肯給我們這批貨.
給個折扣是十分之一.
事情就這樣解決了.
解決之後我不敢出聲.
因為我不知道那些密碼做什麼.
到那些發情差不多.
三四月發生的事件.
到七月的時候.
我開最後一次的.
會員會議.
我把事情都告訴同事.
同事說哇這樣都行.
我說不是我都行.
是上帝行.
我告訴他們.
我說我沒跟人說.
你們明年還準備和我榮優.
我不想榮.

$^{961}$我虧缺了上帝的榮耀.
我就罪人了.
你讀回詩.
羅馬書吧.
所有都是恩典.
所以我後來給了十個字.
他們給我寫了個牌.
一切都是恩典.
榮耀傳歸上主.
不是你讀書厲害.
不是你身體健碩.
不是你有經驗.
一次意外足以致命.
我看中間很多都上了年紀.
是不是.
所以給我一個很經歷的學習.
神的名得榮耀.
是我們人生裡面.
最主要最主要做的一件事.
他不得榮耀.
你很有非事.
So what.
是不是.
所以耶穌基督在地上.
約翰福音17章.
講了兩句話很有意思.
我.
他向父神禱告.
他說他自己來到世界上.
讓人認識耶穌神.
就讀一個真神.
認識你讀一個真神.
得到永生.
跟著下面說.
我已經榮耀你.
因為我已經完成了.
你交托給我要做你的工作.
每個人神都有工作給你的.
你只要好好地完成那件工作.
你就完成這位.

$^{1001}$就是永耀神.
不是你賺很多錢就永耀神.
不是你選到.
總統你就做永耀神.
你完成你的工作.
摩世.
他很辛苦帶領的猜測.
到這個地步.
看到爆火.
扔了上石板.
後來自己要作兩塊石板.
跟上帝交數.
後來他自己作.
還有一樣東西.
當他跟上帝說話的時候.
他說我帶領你去.
不過我不會去.
我找個侍者跟你去.
我差一工人跟你去.
他說不行.
你一定要跟我們去.
如果你不跟我們去.
那我想是.
我今天題目裡面.
第二個.
如果你走人生路.
你沒有上帝.
你說行不行.
都不麻煩了.
如果你走永生路.
沒有耶穌同行.
誰跟你走.
一定死.
所以人生中間要與主同行.
有神的話.
要在我腳前的燈路上的光.
上上給聖靈感動.
在呼應基礎上.
記住.
你的使命.

$^{1041}$是要完成神交托給你的工作.
如果你能夠完成的話.
你就是我的榮耀神.
你不能夠完成的話.
你不能夠是榮耀神.
對不起.
我時間真的過了.
不過我希望.
這些個人的經歷.
所以當時的血壓高高低低.
是真的令我.
以後做事.
人家每次說.
你怎麼怎麼.
不要說那麼多.
為什麼.
一切榮耀都屬於神.
摩西在禱告事情上.
讓上帝看到他的心意.
所以耶和華神說.
我認識你.
今天你的生活.
如果不能夠榮耀神.
你的生活不能夠見證到神.
哪怕你做得多麼的輝煌.
就算今生.
你大名大陸.
但是你如果不能夠榮耀神.
頂著門.
你知道發生了什麼事嗎.
只有今生.
但永生呢.
你怎麼都要揹著它.
求你與我同行.
你那天.
我不知道你小時候怕不怕老師.
怕不怕你爸爸媽媽.
我有一次在培訓會聚會出來.
有個女士急急忙忙追問我.
牧師牧師問這個問題.

$^{1081}$我說什麼問題.
她說我們.
你經常說.
她說得很白.
她說經常被神盯著.
不是太好.
我說被神盯著不是太好.
我看她.
善艾.
我想不想你兒子.
經常盯著兒子.
我說她也不想.
她做好事想不想你盯著.
她做壞事不想.
我說上帝的看顧.
不是在你患難時候看顧你.
在你心高氣傲的時候都看著你.
是不是.
所以摩西禱告.
知道這幫人壞了.
一定闖禍了.
那怎麼辦.
上帝.
有耶穌這個心腸.
念在他們身上.
將愛命除下.
讓他們做.
這個禱告滿足上帝的心意.
因為他認識這位是.
信實慈愛憐憫的神.
任何人只要肯.
回轉悔改去到神面前.
神一定隨聽他的禱告.
結果神聽不聽摩西禱告.
聽.
有沒有同行.
有同行.
所以不用死光.
死那些不信的人.
願主賜福.

$^{1121}$和大家分享一起讀禱告.
讓我們今天早上一同.
透過摩西的禱告.
看到他怎麼認識你.
你怎麼認識他.
求主幫助我們上上.
警醒禱告.
好好在話語上學習.
使我們不單止認識你.
使我們的禱告.
能夠蒙你圓立.
都成為你所認識的人.
求你隨聽我們的禱告.
奉耶穌名字祈求.
好.
\newpage



\section{路加福音 12:33-40}
\label{sec:IFtn6wInBFk}
\textbf{2024.02.11 《龍年祈福》 路加福音12\_33-40 (和修版) 曾裔貴牧師}
\newline
\newline
連結: \href{https://youtube.com/watch?v=IFtn6wInBFk}{\texttt{https://youtube.com/watch?v=IFtn6wInBFk}} ~~~~ 語音日期: 2024-02-13
\newline
\newline
\hyperref[sec:HRSr7qaJnjw]{\small{< < < PREV SERMON < < <}}
~
\hyperref[sec:index_chronic]{\small{[返順時目]}}
~
\hyperref[sec:index_scriptual]{\small{[返順卷目]}}
~
\hyperref[sec:Yl0EevfFnJU]{\small{> > > NEXT SERMON > > >}}
\newline
\newline
路加福音 12:33-40
\newline
\begin{longtable}{cl}
\hline
\hline
章節 & 經文 (和合本修訂版)\\
\hline
12:33 & \begin{tabularx}{0.7\textwidth}{X} 你們要變賣財產賙濟人，為自己預備永不壞的錢囊和用不盡的財寶在天上，就是賊不能近，蟲不能蛀的地方。 \end{tabularx} \\ \\ \relax
12:34 & \begin{tabularx}{0.7\textwidth}{X} 因為你們的財寶在哪裡，你們的心也在哪裡。」 \end{tabularx} \\ \\ \relax
12:35 & \begin{tabularx}{0.7\textwidth}{X} 「你們要束緊腰帶，燈也要點著， \end{tabularx} \\ \\ \relax
12:36 & \begin{tabularx}{0.7\textwidth}{X} 好像僕人等候自己的主人從婚宴上回來。他來叩門，就立刻給他開門。 \end{tabularx} \\ \\ \relax
12:37 & \begin{tabularx}{0.7\textwidth}{X} 主人來了，看見僕人警醒，那些僕人就有福了。我實在告訴你們，主人會叫他們坐席，自己束上腰帶，前來伺候他們。 \end{tabularx} \\ \\ \relax
12:38 & \begin{tabularx}{0.7\textwidth}{X} 他或是半夜來，或是天亮之前來，看見僕人這樣，那些僕人就有福了。 \end{tabularx} \\ \\ \relax
12:39 & \begin{tabularx}{0.7\textwidth}{X} 你們要知道，一家的主人若知道賊甚麼時候來，就不容賊挖穿房屋。 \end{tabularx} \\ \\ \relax
12:40 & \begin{tabularx}{0.7\textwidth}{X} 你們也要預備，因為在你們想不到的時候，人子就來了。」 \end{tabularx} \\ \\
[1ex]
\hline
\hline
\end{longtable}
$^{1}$請主使我們在初二的日子裡敬拜神, 大家這麼早回來, 先敬拜神, 開年跟家人相聚. 和大家先簡單禱告, 多謝主, 我們要將這段聖經講解..
很盼望主的聖靈開啟我們大家的心靈, 讓我們看到永恆的國度才是真實, 而且那是永遠都不會改變的, 幫助我們能夠在我們有限的人生, 為天國, 為神的家,.
為我們將來要豐富進入神的家, 神的國度去準備, 幫助我們每一位今早能夠在聖靈的感動下, 一同在敬拜裡面, 遇見我們的主, 得到你的啟迪, 你的安慰, 和你的教誨..
我們這樣等候奉耶穌基督的名, 阿們. 在牧堂每年的過年, 我們都會有初一的崇拜, 好像有一個非成文的規例, 每年因為我們中國人, 我們會有十二生肖,.
所以每到初一, 港道的同工, 以往的安排, 通常是新同工要講的, 特別是剛剛來, 最早的, 不是太久的, 安排他講道, 又很特別的就是每次講道都講那個生肖, 這個是不成文的, 沒有規定的, 仍然都是按同工自己的感動去講的,.
今年農年, 很多教會的牧者講龍, 聖經裡面的龍, 你想到有多少龍呢? 如果最多人講的就是啟示錄十二章, 講那個大龍, 就是古蛇, 是魔鬼, 講他的結局, 講到主耶穌怎樣來去勝過他,.
信主的人怎樣能夠抵擋這個魔鬼, 就是靠著那個的認信耶穌和他的寶血, 很多人講的, 我上網查過, 以往當然是十二年前, 2001年的時間, 講這一段經文, 很多教會講的,.
我今天和大家預備的, 我就不是講龍了, 中國人相信你和我都一定認定我們是龍的傳人, 雖然沒有根據的, 我們都不知道, 因為從來沒有人見過真正的龍, 是不會出現的, 其他的生肖的動物就有的, 仍然有的, 繼續都會有的, 但是龍就沒有的,.
很特別的地方就是, 中國人和龍的關係, 又是很密切, 很有關係的, 譬如講到龍的祝福語, 很多的, 龍飛鳳舞, 龍躍, 突然間想不起, 很多祝福語的,.
很讓人感受到那種活力澎湃的, 而且大家都知道, 中國人一去到過年, 就憤外很想親近神, 很想親近他認為的神, 而得到祂的祝福,.
你可以想像得到, 年三十晚爭著去上頭炷香的, 往往就是夏惠姨, 還有就是最近她開始回教會, 這樣又有些奇怪的, 到底你想要的是神的, 真神的祝福, 還是入鄉隨俗, 你要去做那個, 讓人看到你仍然緊張地上的,.
你個人的, 自身的福氣, 這個都是我們值得去想, 很多不信耶穌的人, 一去到過年, 就很緊張自己有沒有福, 很緊張自己要去祈福, 要得到福樂,.
那就要問一個問題, 到底福樂從何而來呢? 如何得到福樂呢? 如何能夠每年都有福樂呢? 其實剛才福音樂曲的弟兄姊妹, 已經給了我們答案, 相信今天大家弟兄姊妹來到教會, 你都知道真正的福樂平安, 是在神, 真神的裡面,.
在主耶穌的救恩裡面, 所以今天早上我選取這段經文, 其實很寶貴, 很重要, 告訴我們, 我們的神樂而將他的國度給過我們, 我們是一班在世上很渺小的人,.
但是因為你和我認識了真神之後, 我們就開啟了一個無窮的豐富, 就是神永恆的國度是屬於我們的, 現在就要問這個問題, 也是今天早上和大家在這個講道時間, 一同要去思想, 一同在新的一年, 我們要去追求, 我們要得著的,.
就是到底我們怎樣為永恆準備呢? 怎樣將我們自己的日常生活放在永恆裡面, 來思想和做積極的準備呢? 這段聖經就給了我們兩個很重要的方向,.
首先第一個就是, 原來你要得到神的福樂, 首先原來你需要看到生命真正的價值意義是在哪裡, 這裡告訴我們, 生命不是在乎你擁有多少的東西,.
不是在乎在你的生活裡面, 你有多少可以在今生享用的資源物質上, 真正的豐富是在於我們的生命和主耶穌和神有這個認識, 所以生命的價值不是在於我們的財產擁有多少,.
生命的價值在於我們認識這位生命的主宰, 就是這位神, 因為耶穌基督在他講的這段上文, 是講了一個很有智慧的有錢的財主,.
很勤奮的工作去積聚他在地上的財富, 勤奮和積蓄很豐富, 但是唯一他在耶穌的眼中是無知, 是什麼無知呢?.
他沒有想到原來靈魂不是用物質就可以填滿的, 靈魂需要回到次生命的源頭, 耶穌那裡, 他從來沒有為自己的靈魂去準備, 所以他在耶穌的眼裡面是一個無知的財主,.
還有就是沒有人知道自己什麼時候會被神收取生命, 這兩方面就帶進耶穌基督對門徒的勸勉教導, 就是我們要怎樣能夠為永恆做準備呢?.
第一方面就是我們必須要認定生命的真正價值在於對這位生命的主宰神有認識, 這個很重要的, 今天很多人, 譬如明天的初三, 人們就去車公廟去做那些儀式, 敲那個鼓來祈福,.
轉那個風車, 然後賣那些風車, 完全跟他們所信的所謂車公是完全沒有關係的, 只是做了一個動作, 就覺得那一年就會有順風順水的了, 但是基督徒不是這樣的, 基督徒最重要的生命價值就是你要對主宰, 生命的主宰有認識,.
當你認識了這個生命的主宰, 你就能夠在你的生活裡面有意義地, 有目的地這樣去生活了, 所以第三十三節就立刻告訴我們一班有神國度的人, 他已經看到生命的意義不是在於不斷的擁有, 記住在聖經裡面金錢財富是忠誠的, 絕對不是邪惡的,.
提摩一帶前書也說過, 貪財才是萬惡之根, 不是錢財是萬惡之根, 不要搞錯, 貪戀錢財才是萬惡之根, 是心底裡面好像不能夠滿足的那個黑洞, 那個貪,.
所以這些告訴我們真正對生命有意義認識的人, 首先第一個行動就是三十三節, 主耶穌這樣說, 變賣所有周濟人為自己預備永遠不會壞的前浪, 用不盡的財寶在天上, 是沒有東西可以奪去的財寶,.
因為三十四節主耶穌就說出一個很重要的內心價值, 到底我們的內心注目的是永恆不會改變的, 還是短暫的地上的物質的個人的無盡的自己的私慾呢?.
是兩者完全不同的, 所以有了三十四節, 你明白原來我們的心, 我們的財寶是在哪裡? 當然很清楚是要在天國, 為著將來的這個我們要進入的天國做我們的準備,.
你就會有一個很大的轉向, 就是你不會再介意分享帶來的喜樂, 那份平安, 那份做在神裡面對別人的周濟和幫助, 這個是一個很重要不同的心態看永恆的福樂,.
如果不是我們只會不斷祈福, 祈我們自己的福, 我們變成越來越有一種自我中心, 所以基督教在神國裡面, 你會見到往往是基督徒願意付出, 多為別人行一步,.
有一個例子, 我很久之前說過, 應該不是在這裡說, 在武堂說過, 有一位中文大學總腦科腦外科教授, 他叫做王海東, 他有一天放大假, 正常的日子, 他會看一些關於腦外科的一些網站,.
上網去瀏覽, 可能增進知識, 可能有時候有些case他可以看看別人怎麼做, 一般他看完就馬上delete, 就會刪除, 因為有時候都多, 他剛剛那天放假, 無意中見到一個女孩, 15歲, 是青海人, 她的名字叫做糾格,.
他見到這個網址的時候, 就吸引他的注意力, 因為那個標題說這個女孩只剩下四個月的壽命, 裙衣束手無策,.
他看到這個標題就很有興趣, 因為和他的本科有關, 因為這個女孩腦裡面有一個很大的瘤, 好像一個哥爾夫球那麼大, 醫生說沒得救了, 會診過, 當然是在青海三間醫院的醫生,.
這位黃教授就說為什麼這麼快就放棄, 他就問了那裡的醫院, 聯繫了, 就拿他的一些斷層的片拿來香港去看, 他看完之後發現原來不是絕症, 還有六七成的機會可以做手術,.
就馬上找了中文大學威爾斯的醫生, 腦腫瘤科的醫生幾個去會診看這件事, 大家都認同真的有六七成的機會, 還不是絕症, 就馬上安排, 原來這個發出網址的是基督教例行會的總幹事張紅秀美這位女士, 這位姐妹,.
她因為知道她在她的兒童福利院裡面, 在青海的福利院裡面有一個女孩, 本來這個女孩很乖巧的, 但是無奈突然之間發現有腫瘤, 而且越來越大, 影響著她有時會突然間暈倒, 有時會有很大的問題出現,.
有些甚至乎大小便失禁, 影響著她的腦神經, 但是這個總幹事覺得大家都沒有辦法, 不如我嘗試一下把她的情況放上網址, 如果有心人能夠看到, 又有能力改變的, 就希望可以幫到這個女孩,.
果然是一個神蹟, 一位姐妹只是希望有人能夠看到這個故事, 這個情況就來幫助她, 剛好黃教授剛剛放大假, 看到這個事件, 他就想多走一步, 不如拿這些片下來看一看, 看完之後果然真的有機會可以改變,.

$^{41}$當然他們要籌錢, 因為她不是香港的公民, 看醫生手術一定要立即做, 很快就籌了13萬5000元, 就是這個手術費, 幫這個15歲的青海女孩在香港威爾斯做了手術,.
親愛的弟兄姊妹, 主耶穌的教導, 原來神國度裡面的福樂, 是認定我們在這位生命的主宰, 這位做物主的神面前, 發揮我們在地上的作用, 就是多些施捨, 幫助有需要的人,.
在我們不同的崗位, 不同的領域, 你怎樣來運用你的人生, 耶穌說要變賣的, 可以是實際的金錢, 可以是你的才幹, 重點是我們要來有這份施捨的心靈, 施比受更為有福,.
我們將我們自己的生命, 在神的手中被運用, 這個是一個很奇妙的例子, 讓我們看到原來神國度裡面, 不是不斷的祈福, 是獲得從神那裡去幫助人的福樂,.
第二方面很重要, 主耶穌說神國度裡面的人, 要怎樣來為永恆準備呢? 相信第35到40節, 耶穌用了一個很簡單的比喻, 來講到一個等待永恆降臨的人, 應該怎樣來準備自己, 就是等待, 怎樣的等待呢?.
大家喜歡釣魚嗎? 我絕對不喜歡, 因為最難受的就是放一支竿在那裡, 等待, 等待完之後, 有時水動了, 你以為有魚了, 拿上來, 魚餌沒有了, 仍然是這樣, 所以我很不喜歡釣魚,.
我只知道張美眉牧師的丈夫曾錫華牧師, 他在神學院時常常叫同學一起去釣魚, 去捉墨魚, 同學很喜歡, 但我很不喜歡, 因為我沒有忍耐, 不喜歡等待, 但這件事已經教我們, 要等待,.
但不是純粹忍耐等待, 而是要很積極的準備自己, 你看看主耶穌怎樣來勉勵他的門徒, 就是說你們腰裡要束上帶, 燈都要點著, 自己在好帳篷人等候主人, 從婚姻的筵席上回來, 留意, 他來到叩門, 就立刻給他開門,.
主人來到, 看到這個僕人警醒, 那個僕人就有福了, 我實在告訴你們, 主人必叫他們坐直, 自己束上帶, 進前侍候他們, 調轉了身份, 主人變了僕人, 在古代似乎是不可能的,.
現代也不會有, 你不會叫飛鴻, 今天年三十晚, 團年飯我煮, 你坐在那裡吃, 不會的, 很少的, 我昨晚回到媽媽那裡和我兒子煎年糕, 剛剛媽媽的飛鴻回來了, 去完街八點多回來, 跟著就立刻換衣服, 就進廚房和我姐姐一起煎年糕,.
其實他放假的, 但他也很自然的進去, 看到有媽媽的兒子回來, 很難得, 他就煎年糕了, 很少人會這樣, Mary你坐下, 我們煮給你吃, 但這裡耶穌這樣說, 到底這個是將來永恆的延籍, 還是在這個比喻裡面, 對一個警醒的獎賞呢?.
積極的等待其實就是警醒, 準備自己讓那個降臨, 你是有福和有份, 這個是很難的, 長時間的警醒是最難的, 無盡的等待最難的就是考驗我們的忍耐,.
所以不是只是一種被動的坐下等待, 是積極的準備自己的警醒, 大家可能都認識鄭丹瑞, 通常叫阿旦, 現在他有一個叫健康蛋, 沒錯了,.
他是一個YouTube的頻道, 不斷推廣健康的生活等等, 有很多不同的醫學知識, 我想說一下就是偶然間聽到他一個分享, 說他很多年前出道的時候,.
他由商台被香港電台重金禮聘, 他由商業電台轉去香港電台工作, 他主持的當時叫做音樂蛋, 是黃金時間早上九點到十一點的, 港台一個重要的節目,.
他從商台過來, 被人重金禮聘, 在這個時段他就一直做, 感覺自己很得心應手, 也有很多人聽, 接著突然間有一天看報紙, 看到娛樂版,.
就說港台要精簡人手, 要cut預算, 高層說到明開刀, 要將音樂蛋這個主持cut, 是他看報紙看到自己會被cut,.
就很激氣, 很生氣, 為什麼在報紙娛樂版說我要被裁員呢? 馬上和他相熟的高層吳嘉年談, 高層說是真的事實, 電台是沒有資源, 現在要cut, 要找你開刀, 是黃金時間,.
他後來和高層爭取, 可不可以不cut我呢? 他說不行, 另外我幫你這樣, 給你一個時段, 你去做吧, 是什麼時段呢? 不是黃金時間, 是深夜的兩點到六點, 可想而知那些節目哪會有人聽呢?.
但沒辦法, 至少不要裁員, 就在那裡做, 你可以想像得到他的心情是怎樣的, 那個叫輕彈淺唱不夜天, 有人聽過嗎? 會嗎? 你們晚上不用睡嗎? 兩點到六點, 你可以想像他怎樣做, 敷衍著做, 那時候沒有敷衍這個字,.
很生氣, 完全沒有mood, 是發脾氣地做, 做到一個地步是很古怪的, 播完梁祝的粵曲, 那一邊播hotel california, 完全不拉奸的, 因為他根本是很生氣, 也沒有高層會聽你的節目, 任他做也沒有人會知道的,.
因為是深夜的節目, 沒多久, 他們港台有一個工友, 每個同事都會派信, 因為他們會有聽眾來信, 阿旦, 有封信是你的, 打開封信, 沒有署名的,.
我直接讀給大家聽, 如果你不喜歡, 你就不要做, 我每晚四點鐘, 我需要有音樂陪伴我工作, 這句話好像當頭棒喝, 原來有人留意他, 不過是你敷衍著做, 你這麼不喜歡就不要做, 我需要有音樂陪伴我,.
這樣令他很感到羞恥, 因為原來有人聽這個節目, 原來有人留意他做得好不好, 所以由那封信之後, 他很認真去準備, 很認真去服侍, 慢慢開始收到不同的信,.
但是阿旦心底裡面, 他原來在等一封信, 一封可以被肯定的信, 再過了一段時間, 接著那個工友又再來告訴阿旦, 你有信了, 他一看筆跡, 他就知道上一次很不禮貌的聽眾來信, 這次是什麼呢?.
這次是有署名的, 阿旦, 我想四點鐘播這首歌, 聽到這個聲音, 這個說話, 這個肯定, 他整個人充滿了那種力量, 親愛的弟兄姊妹,.
主耶穌說, 第37節, 主人來了, 看見僕人警醒, 那個僕人就真的有福了, 朋友, 兄弟姐妹, 我們在新的一年, 我們有一個積極的等待, 就是讓我們的心靈準備好, 去進入神永恆的國度,.
在這個準備過程, 我們是要積極的等待, 不單純是忍耐, 而是我們能夠讓我們自己的生命成為別人的幫助和祝福, 好像剛才黃海東教授無意之間看到這個網誌, 就救了一個生命, 就是這位張洪秀美, 這位總幹事,.
她不甘心, 希望這個事件不應該只是一個人體的絕症, 給這位15歲的女孩一個機會糾葛, 也讓我們想到阿旦, 她因為人對她的責備, 成為她生命的動力, 她自己能夠有一個生命的改變,.
很盼望新的一年, 我們能夠真正的福, 在神裡面的福, 就是我們能夠將我們自己的生命放在永恆裡面, 去看我們自己, 我們一起禱告, 感恩天父給我們新的一年的開始, 我們不單止有元旦的1月1號, 更加有我們中國人很體重, 很寶貴的農曆新年,.
我們可以再一次在新的年度裡面去立志, 幫助我們能夠對於永恆很有信心把握, 更加我們能夠每一位弟兄姊妹, 我們為主的國度, 一同去服事, 去效力, 讓更多人因為我們的見證, 認識我們的主, 我們的真神,.
我們這樣感恩奉耶穌基督的名, 阿們!.
\newpage



\section{創世記 11:1-9}
\label{sec:Yl0EevfFnJU}
\textbf{2024.02.18 《宣教使命:從分散一族到召聚萬國》 創世記11\_1-9 (和修版) 楊穎兒宣教士}
\newline
\newline
連結: \href{https://youtube.com/watch?v=Yl0EevfFnJU}{\texttt{https://youtube.com/watch?v=Yl0EevfFnJU}} ~~~~ 語音日期: 2024-02-20
\newline
\newline
\hyperref[sec:IFtn6wInBFk]{\small{< < < PREV SERMON < < <}}
~
\hyperref[sec:index_chronic]{\small{[返順時目]}}
~
\hyperref[sec:index_scriptual]{\small{[返順卷目]}}
~
\hyperref[sec:lg5Gw9lhP5s]{\small{> > > NEXT SERMON > > >}}
\newline
\newline
創世記 11:1-9
\newline
\begin{longtable}{cl}
\hline
\hline
章節 & 經文 (和合本修訂版)\\
\hline
11:1 & \begin{tabularx}{0.7\textwidth}{X} 那時，全地只有一種語言，都說一樣的話。 \end{tabularx} \\ \\ \relax
11:2 & \begin{tabularx}{0.7\textwidth}{X} 他們向東遷移的時候，在示拿地找到一片平原，就住在那裡。 \end{tabularx} \\ \\ \relax
11:3 & \begin{tabularx}{0.7\textwidth}{X} 他們彼此商量說：「來，讓我們來做磚，把磚燒透了。」他們就拿磚當石頭，又拿柏油當泥漿。 \end{tabularx} \\ \\ \relax
11:4 & \begin{tabularx}{0.7\textwidth}{X} 他們說：「來，讓我們建造一座城和一座塔，塔頂通天。我們要為自己立名，免得我們分散在全地面上。」 \end{tabularx} \\ \\ \relax
11:5 & \begin{tabularx}{0.7\textwidth}{X} 耶和華降臨，要看世人所建造的城和塔。 \end{tabularx} \\ \\ \relax
11:6 & \begin{tabularx}{0.7\textwidth}{X} 耶和華說：「看哪，他們成了同一個民族，都有一樣的語言。這只是他們開始做的事，現在他們想要做的任何事，就沒有甚麼可攔阻他們了。 \end{tabularx} \\ \\ \relax
11:7 & \begin{tabularx}{0.7\textwidth}{X} 來，我們下去，在那裡變亂他們的語言，使他們彼此語言不通。」 \end{tabularx} \\ \\ \relax
11:8 & \begin{tabularx}{0.7\textwidth}{X} 於是耶和華使他們從那裡分散在全地面上；他們就停止建造那城了。 \end{tabularx} \\ \\ \relax
11:9 & \begin{tabularx}{0.7\textwidth}{X} 因為耶和華在那裡變亂了全地的語言，把人從那裡分散在全地面上，所以那城名叫巴別。 \end{tabularx} \\ \\
[1ex]
\hline
\hline
\end{longtable}
$^{1}$弟兄姊妹早晨,新年夢幻,很開心能夠回來紅印家見到大家.
剛才美娟也提到,我在2020年實習的時候,在紅印堂做神學生的時候.
我還記得再次領受這個呼召,那時候是和曾姑娘魯牧師一起在祈禱裡面去服侍非洲人.
所以我畢業之後在堂會服侍非洲人,現在轉去警察局繼續專注服侍非洲人.
所以也很開心,紅印家也見證著我服侍的成長,所以在宣教裡面大家都有份.
我想這樣來介紹自己也是比較親切的,我們一起來祈禱.
慈愛的天父,我們今天聚集在你的面前,都是尋求你的同在和引導.
求你用你的智慧去充滿我們,讓我們明白你在巴別塔的故事裡面的教訓.
讓今天的講道從孩子的口中講得清楚,滿有你的智慧.
我們奉主耶穌基督的名而求,阿們.
今天簡選的經文其實對我們肩負耶穌給我們的大使命是很有關係的.
大使命這個課題是很廣闊,從個人的信仰,教會做地區的猜拳.
以至教會在世上如何參與宣教等等,都是和大使命有關.
所以今天很想和大家從上帝的創造心意裡面,我們一起看看如何在大使命裡面有份.
可能大家已經想到,今天講大使命,為什麼不是講馬太福音第28章19節.
它是一講大使命我們就會提到的經文,經文是這樣說的.
你們要去,使萬民作我的門徒,奉父子聖靈的名給他們施洗.
為什麼今天我們會選擇巴別塔的經文去講呢?.
如果像我們的講題那樣說,我們是從一族要去到萬族.
我們今天就回到創世記去看看神在大地創造的心意,它是如何連接到大使命.
可能你會想,創世記裡面有這麼多經文,有亞伯拉罕,有約瑟,為什麼偏偏要選擇巴別塔這個片段呢?.
就好像神懲罰人類的片段去再講大使命,巴別塔和大使命有什麼關係呢?.
連主人學的小朋友都會知道,巴別塔的故事都是讓我們引以為鑒.
我們慣常都會知道,人類的高傲自大就是建造了那座城,那座塔,而且塔也是通天.
就代表著人以為自己的能力可以和耶和華相比,去藐視了神在大地的權能.
引起耶和華極大的不滿,於是懲罰人類分散他們的語言,使他們語言不通,不能夠建好這座塔.
我們就要問,巴別塔的重點是不是講懲罰人類的驕傲呢?.
在今天早上,我們很想進入經文的世界,從神的創造心意去想想我們如何承擔宣教使命.
第一,當時神的性情,我們要想一下,是不是特別喜歡審判人呢?.
其實如果要我們描述神是一個審判的神,其實是沒有錯的.
因為由第一章至十一章,其實都記錄了人不斷墮落,離棄了神的教導,以致神要赤下洪水,消滅整個人類.
到十二章開始,就是亞伯拉罕從一族往外看,是一個走信心之路的故事.
神作為審判的神,從第十二章開始,這個角色就漸漸式微.
下一次要和整個子民去審判他們的時候,可能是要去到被擄的時間.
我們所相信的神是不變的,如果我們說舊約的神,是經常審判散播人類的,又消滅人類的.
現在我們想想,我們經歷的那個上帝,是否一做錯事就立刻懲罰我們呢?.
所以我們去想巴別塔是否真的懲罰人類呢?我們去看下去.
我們繼續看的時候,我們看看當時人的特徵,是否特別邪惡,要引致上帝要一個極大的審判呢?.
當時傳遞了一種語言,說一樣的話,而且大家彼此相諒同心,要一起拿磚頭去造成那座塔,那座城.
同時大家亦成為了一個民族,去說一個語言,很同心,很團結.

$^{41}$第五節,耶和華去看到世人所造的城和塔,一看之下,看到他們塔頂通天,很有計劃,他們造磚的技術很成熟.
而且他們彼此相諒,要將這些磚一起消透,將這些磚拿來造石頭,又放一些泥漿進去.
我們可以看到人類怎樣的同心合意,而且是很有計劃的.
他們亦遷移到一個士拿的位置,士拿的位置是指昔日的巴比倫,即是今日伊拉克的地方.
在士拿這個地方,很有趣,是一個河流交匯,其實島地是很肥沃的.
某程度上,這些經文的背景,是指當時的人類,文化,建築,科技,其實都是很興旺的.
大家同心一致,團結一致,城市是非常繁榮,而且我們不要小看,同心合意的力量是可以很厲害的.
因為當神見到人同心的時候,第六節這樣說,他們現在做什麼事,都沒有什麼可以難了他們.
似乎人類用同心去建設,去將帶到人類,整個城市進入文明呢?.
是不是上帝要懲罰人類的原因呢?其實我們現在生活的都是一樣的,一路有科技,一路發展,有AI,都是一個過程來的.
我們在家庭,我們在職場,我們在教會,我們都是需要同心的.
好像中國人有這樣說的,團結就是力量,夫妻同心,其利斷金.
如果不同心,我們就會有很多爭拗,很多不團結,事情就會有拖延.
很明顯上帝懲罰人類,不是因為他們同心,同心有很多好處的.
在以往的解經,很多時候說,也是人類一起同心,聯繫上驕傲自大.
但奇怪的是,當我們去看經文原文的時候,其實沒有一個字是和驕傲有關的.
但我們可以肯定,一定是人做了一些事情是神喜悅的.
我們看看第三節,第四節,他們為自己立名.
我們會想到,他們為自己立名是不正確的,因為他們只是傳自己的名,沒有傳揚神的名.
但我們今天要進入創世的故事去看.
創世的故事裡面,其實由第一章至十一章,都沒有提到神要人去傳揚神的名.
其實在創世以來,神的尾線已經在大地顯現出來.
那種完整,完美,萬物是按其線,按其規律去生長的時候.
其實已經是無言無語地在傳揚神.
而且當神做人的時候,是有人的形象.
人的形象是會反映著上帝那種線,那種光.
而當人沒有遵守神的說話,要墮落的時候,就不能夠反映神的形象.
不能夠顯出神的那種榮美,不能夠再顯出神的那種名的尊貴.
所以我們再去想想,為自己立名又是不是這麼邪惡呢?.
當我們去看「立名」這個字的時候,其實是跟工作有關的.
我們可以在我們現在的生活中去想想.
我們附近有些中學叫做「陳XX紀念中學」.
其實這些中學都是在做一些教育的工作,用了一個名字去有一個意義.
有些公司也會用後面那兩個名字去做公司的名字.
所以這些「立名」其實不是說我們要傳揚自己的名字.
都是跟一些工作有關.
而且我們去看回頭十一章的時候,我們看到第四章,第五章,第十章.
其實說完巴別塔之後,都是列了很多的名字.
如果名字不是重要的話,頭十一章不用花那麼多篇幅去說一些家譜,一些名字.
再說下去,我們在十二章第二節也看到耶和華神親自主動向亞伯拉罕賜下一個祝福.

$^{81}$而其中一個祝福就是使他的名為大.
所以有些學者就說其實「立名」也不是一個彌天半役的大罪.
我們就在想巴別塔的故事在創世記裡面.
重點未必是說上帝懲罰人類,也不是人類的驕傲,也不是人沒有傳揚神的名.
其實巴別塔的重點是什麼呢?.
經文很有趣,當經文有重點要說的時候,是會重覆地說的.
第四節,他們說「來,我們做一座城,一座塔,塔頂通天,我們要為自己立名,免得我們分散在地面上」.
第七節,耶和華變亂他們的語言,使他們語言不通.
第八節,於是耶和華神使他們從那地分散到全地面上,他們就停止做那座城了.
第九節,因為耶和華變亂了全地的語言,就把人從那裡分散到地面上,所以那座城就叫做巴別.
大家看到嗎?重要的事情是要重覆地說三次.
剛才在經文看到耶和華在他們當中要分散他們.
他們知不知道耶和華神要分散他們?.
他們知不知道這是神的心意?.
他們當然知道,不知道是不會說的,他們說什麼?.
「免得我們分散在地面上」.
如果他們不知道的話,他們是不會說的.
我們看到他們知道,但他們不想做.
知而不做,為什麼?.
為什麼知道也不做?.
原來有學者發現當時有個傳統,是有紀念死人的傳統.
紀念死人是要靠後世的人去紀念死去的人.
不會是死人去紀念死人.
所以當人聚居生活的時候,後世的人才可以對死去的人進行紀念.
但問題來了,當一分散到全地的時候,那時候沒有我們交通現在這麼方便.
問題是誰還去紀念這些死人呢?.
在我們今天的生活中,其實我們也有不同的信仰,也有類似的制治.
在我們的傳統中我們也認識很多.
所以說的也是人關心的問題,就是死後去哪裡?.
死後的生活呢?.
如果這樣說,我相信我們也很能夠理解.
所以我們看到,從創天造地的時候,耶和華神不是要人聚在一起.
不是要人去紀念死去的人.
不是聚焦在死後的世界.
而是要人活在當下,要去分散到全地.
所以我們看到了,問題不是在於塔有多高.
不是在於那時候的人類發展有多文明,或者文明帶來的自大.
而是在於建築這座城,這座塔的背後的目的不是和神的心意合而為一.
他們是想免得我們分散在地面上.
我們又想想,神是何時說過要人類分散全地呢?.

$^{121}$我們去看創世記第一章28節的時候.
就這樣說道,神賜福給他們,即是人類.
神對他們說,要生養眾多,遍滿全地.
而到了羅亞洪水之後,神也賜福給羅亞和他的兒子.
要他們生養眾多,遍滿全地.
可能你會說,生養眾多真是難道.
這是中文的譯本,但我們看原文解釋的話.
其實生養眾多是和結果子的概念有關.
我們住在接近郊外的地方,也很自然會接觸到大自然.
結果子對我們不陌生.
因為有大的果子,有小的果子.
果子也會有些結得漂亮,有些結得不太漂亮.
或者看它有些光澤.
果子要營養充足,要沒有蟲豬.
要按著季節去結果子.
而且果子有什麼特點呢?.
它是一棵植物生長出來的,不只是自己一顆.
一棵植物會有很多果子結出來.
我們去到新約的時候,聖靈的果子也有很多特質.
聖靈的果子有人愛喜樂和平忍耐.
恩慈良善信實溫柔節制.
果子結得漂亮,不只是漂亮,而是多的.
所以這個生養眾多的概念.
其實是要由一變到多,要從一而出.
這是神在大地的心意.
人類是從亞當一人而出.
神賜福給人類,要他們生養眾多.
但因為人的墮落,神賜下洪水.
洪水過後,神又興起羅亞一家.
要他們生養眾多,遍滿大地.
到了亞伯拉罕的時候.
神要這一族成為大國.
使萬族必因他而得福.
當我們看到整個人類的救贖歷史的時候.
神也猜派自己的愛子成為人.
從一人的死而令全地的人得救.
大家看到嗎?這些都是通的,從一而出.
今天救恩已經在耶穌基督完成了.
但仍然有很多人仍未相信耶穌.
所以教會是有責任的.

$^{161}$繼續秉承著大使命.
向外伸展,向萬族奉父子聖靈的名.
與人施洗,遍滿全地.
今天我們能否將從一而出的概念.
擁抱在我們生活的每一個環節呢?.
從我們的家庭,我們的工作,我們的教會的事奉.
其實我們面對著很多挑戰.
如果我們從一而出.
特別是年長的時候.
子女都長大了,他們都會離我們而去.
甚至會去到外國的地方.
我們能否願意他們拿著神的祝福.
去到遍滿全地的地方,繼續祝福不同的人,不同的社區.
宣教其實是每一間教會都要承擔的使命.
而且我們宣道會的教會.
也帶著宣道的基因,宣道的精神.
我們會在地區做猜拳的工作.
去接觸這個社區.
我們也會有宣教士在前方.
去到神國福音最需要的地方.
去面對那些屬靈的爭戰.
去將福音帶給美德的群體.
而教會就在後方為他們禱告,守望,奉獻支持.
這是一個神創造心意的宣教模式.
大家一起去遍滿全地.
去看到大地上神的作為.
而且這個使命是人人都可以參與的.
不要以為宣教只是宣教士做.
我們就是從我們平日的生命已經擁抱著這個概念.
要生養眾多,要結好的果子,遍滿全地.
明明知道要遍滿全地,卻去做一座城一座塔.
不去遍滿的話.
是這一點引起耶和華神的不滿.
耶和華神要求人們向橫向多元的發展.
但他們偏偏選擇成為一個內聚的群體.
當他們不向外面發展.
神創造的大地還有誰來關心和管理呢?.
十多年前,洪恩堂在洪水橋這個地區開始服事.
正正就是由錦繡花園,錦宣家,錦繡堂,錦繡地區.
去分散出去,擴展出去.

$^{201}$同樣,今天在洪水橋的社區.
我這次來到,看到很多不同的大廈不斷建造.
福音的禍長很大.
在洪水橋做地區工作.
不是要我們去內聚.
而是要我們散出去.
去做好這個社區的猜拳工作.
知道大家早前有一個派輝春的活動.
吸引了社區很多人去認識這個教會.
所以當人交了權力給人去做的時候.
人沒有好好按著神的心意去做.
是人不馴服.
無論有多同心,都是枉然.
神要藉著一個教會,一個生命.
去祝福大地更多的人.
當人不馴服的時候.
對神來說就是一個危機.
因為直接影響著神賜福給大地的心意.
我問,如果是這樣的話,神要不要出手?.
要啊.
第七節.
剛才說到神是要變亂他們的語言.
要他們語言不通.
然後要分散他們.
使他們停止建造那座城.
亦是變亂他們的言語.
去分散他們在地面上.
其實是重複說同一個概念.
耶和華說創天造地的時候.
我們要細心去想.
我們做植物的時候,不是只做一棵.
做樹葉的時候,不是只做一款.
做花朵也不是只做一種顏色.
而動物也不是只做一種.
是很多種類的.
在座的弟兄姊妹.
你看看身邊的弟兄姊妹.
我們大家有不同的五官和樣子.
但我們在神裡面.
都是得著神的恩典.

$^{241}$而生命也是可以發出上帝的光.
但當我們去看這個世界.
是有不同的語言.
不同的人.
不同的膚色.
不同的種族.
不同的信仰.
這也是神的心意.
神的心意是蘊藏著很多元的元素.
我們再想想巴別塔這個故事.
最終變亂全地的語言.
由一個語言變成更多的語言.
是不是神真的在說要懲罰人類呢?.
還是說祂要將這個祝福繼續延伸到大地呢?.
很有趣.
又有學者去看「變亂」這個字.
原來「變亂」這個字.
可以解釋為「混合」.
即是混合在一起.
或者有一個「調和」的意思.
而且在聖經裡出現過43次.
都是說要將他們混合去調和.
就好像我們做餅.
只放麵粉是不行的.
我們要放些油去搓得它軟軟的.
才可以做到一些餅.
所以人馴服的時候.
神就要出手了.
要調和.
就從這個語言開始.
使人類能夠分散各地.
以致他們可以生養眾多片滿全地.
我們今天在我們的社會裡.
都看到.
其實香港也是一個很獨特的社區.
我們身邊有很多的照顧者.
我們有很多的家庭傭工.
他們是來自菲律賓,印尼.
或者其他的東南亞國家.
他們都是來到我們中間.

$^{281}$都是在我們單一的華人社區裡調和.
就好像我們教會.
我們也有泰語的士工.
我們在錦宣家那邊.
在佳印堂那邊.
也有一些服侍傭工的士工.
所以我們看到教會.
其實不是只有一個模式.
教會就是信徒.
教會就是要分散.
從一族到朝聚萬國.
特別在宣教的層面.
我也很希望可以和洪恩嘉的弟兄姊妹.
一直看得遠一點.
在基督宗教之下.
我們相信主耶穌基督的拯救.
我們會包括有天主教.
東正教.
新教.
新教裡面也有宣道會.
有播道會.
有浸信會.
其實我們都是同屬於基督的肢體.
在基督宗教裡面.
不只是這個宗教.
我們也有其他的宗教.
我想給大家看一張照片.
我現在在一個警察會服侍.
警察會叫做非洲內地會.
專注服侍非洲人.
在非洲的非洲人.
也有一些散居在外地的非洲人.
這張就是我在警察會.
在外國的一個圖片.
他們舉辦了一個活動.
請來很多北非的群體.
北非不只是非洲人.
有摩洛哥.
利比亞.
埃及.

$^{321}$他們也不是黑色皮膚.
但他們有一個特點.
他們一出生就信奉另一個宗教.
所以在他們的生活上.
他們會包著頭.
這個警察會舉辦了一個活動.
有一個聚餐活動.
請來很多穆斯林的群體.
同台一起分享生命.
分享信仰.
很美麗.
很興旺.
能夠招聚到這麼多穆斯林群體.
另外一張圖片.
我在網上找到.
原來在非洲.
有一些婦女.
她們是來自不同的宗教.
例如左邊那位.
大家會看到她是信奉基督教.
而白色衣服的.
其實她是一位天主教的修女.
而兩位穿黑色衣服的.
她們是信奉穆斯林的女士.
大家在建立關係的時候.
原來她們會先交換經書.
大家交換本《可蘭經》和《聖經》.
以表示首先要尊重大家和大家的不同.
我覺得這幅圖畫很特別.
所以我想跟大家分享.
很多時候我們會想到.
我們不和其他宗教的人來往.
或者我們會很容易有一種驕傲.
去指責其他宗教的邪惡.
其實我們是很難去分享愛的.
第一步去建立關係都這麼難.
怎樣可以繼續說下去呢?.
我現在還會在神學院修讀一些跨文化的課程.
其實我們學習的不是宗教的技巧.
不是三福四律.

$^{361}$反而是學別人的宗教.
他們究竟在說什麼呢?.
例如我們會去考察其他穆斯林的寺廟.
去看看他們的書籍.
很多時候我們會覺得我們不能看.
這是我作為宗教學者要做的工作.
但對我來說.
如果我不去看他們說什麼.
我怎知道他們在學什麼宗教信息呢?.
就好像我們去打仗.
我們都要知道對方是拿刀還是拿槍.
這樣我們才可以去爭戰.
這不是說我們放棄宗教.
不是放棄我們相信耶穌基督為神的兒子.
不是放棄用祂的寶血拯救我們每一位.
而是我們嘗試收起我們的成見.
收起手指.
然後跟別人同座.
去聽別人說的話.
或者他們的掙扎和恐懼.
在我服侍本地非洲的穆斯林群體的時候.
我發現他們經常被一些做夢的情況去擔心.
疑生疑鬼.
家中牆上可能有濕了水的污漬.
就覺得有邪靈在家中.
其實這些就是他們的宗教給他們的一個綑綁.
所以在我的工作中.
我都要去接觸他們學的是什麼.
在座的未必有膚照去做宣教士.
我們今天是怎樣去承擔這個大使命呢?.
其實很簡單.
就是從生活開始.
很好的靈修是很重要.
但更大的考驗就是我們在群體生活的時候.
我們怎樣將神的形象去顯示出來.
由起床的一刻到晚上睡覺的一刻.
我們都能不能夠去表現那種屬靈的果子的特質.
我們見到我們的家人.
我們的家庭傭工.
我們父母子女.

$^{401}$走在街上見到的鄰居.
我們職場的同事.
我們都是否有這樣的樣式去感染他們呢?.
看到耶穌的拯救在我們生命的改變.
我想有一幅圖畫給大家看看.
大家認不認得這個地方在哪裡?.
對呀!很熟悉呀!洪水橋.
我最喜歡在哪裡呀?.
前面呀!.
對呀!前面呀!旁邊呀!.
換個角度給大家看看.
其實大家都可能經過這個地方.
為甚麼我會拍這個地方呢?.
因為自從我在這裡做神學生的時候.
我已經在這間店舖買Pizza.
這間Pizza是由一個巴基斯坦印度籍的夫婦去經營的.
他們很有善.
而且他們懂得說中文.
有一次我去買Pizza的時候.
聽到原來她播放著一個收音機.
收音機是在播放一些關於信仰的訊息.
但因為拿Pizza的時間很快.
我還未聽得清楚她在聽哪一個信仰在說神.
很有趣呀!我們在這麼近已經有一個不同國籍.
不同膚色.
不同語言.
不過也有學我們語言的人在我們附近.
所以弟兄姊妹.
我們去想一下.
甚麼是宣教呢?.
宣教其實不是像我這樣去到前方去看別人的經書.
但是是在說我們今天有某一個概念.
就是要從一而出.
要擁抱著這個神.
就是用我們每一個人.
在座每一個弟兄姊妹的生命.
去參與宣教.
就像我開始這樣說的時候.
當時我是領袖著去服侍非洲人的呼召.
更加確定的時候.

$^{441}$是在這個地方.
所以我很感動.
是在座各位一起去見證的.
所以當我們去了解宣教的時候.
我們不要想又是這個宣教士回來述職了.
關我甚麼事呀.
但我們可以去聽一下他們的事供.
又或者不要想.
經常是宣教士又來問教會要錢.
其實不是的.
做宣教的工作.
神是會有充足的供應.
是不用擔心的.
但更加要明白的是.
神要我們去分散.
我們要去看到神的國.
如何在遠方那裡彰顯.
去將萬族歸回祂名下去敬拜.
然後我們在世界各地.
我們用不同的語言.
不同的敬拜方式.
在世界上不同的地方去敬拜.
其實這個就好像天國的景象.
上帝關心的事就是我們關心的事.
所以用一幅圖畫.
用不同種族的人.
去總結今天的講道.
其實上帝是透過猜拳宣教.
帶我們去到每一個角落.
我們去見到人生命那種困苦流離.
那種流離失所.
然後叫我們謙卑.
我們去學習用耶穌的心腸去憐憫人.
這是我們參與宣教的其中一個重要目的.
就是叫我們的生命有更新有改變.
叫我們不再關顧自己的事.
而是關心別人的事.
關心上帝的使命.
上帝是願意使用遠方的生命.
只要我們教會願意加入.

$^{481}$去承擔宣教的使命.
上帝就是用遠方的生命去祝福我們每一位.
所以我都很想在這個時間.
再講多一點.
我們怎樣可以改變我們去擁抱這種概念.
其實我們可以做一些自我的審查.
我們可以從成長的過程中.
我們會不會對某些種族.
某些語言文化有些愛恨.
會不會有一些偏見.
會不會對某些群體有些定型.
他們一定是這樣.
或者很難接受他們來到我們中間.
在我們的社交圈子裡.
我們會不會去開放認識他們的文化.
我們不是單方面的接納.
我們要在這個和而不同的裡面.
去堅持去維護我們的信仰.
我們要站立得穩.
然後我們會不會由家庭做起.
會不會家庭裡面已經有一個不同我們同種族的人.
或者會不會讓我們的小朋友再去到遠方.
這些都是我們心中有些掙扎.
教會的使命.
教會是否可以有不同的種族.
不同的言語的人進入教會這個地方.
這個都是很值得我們繼續思考.
當教會是同心.
而且較正目標.
注目點是要分散傳遞.
要伸兩眾多的時候.
我們可以回到今天的經文.
只要做任何事.
神都不會阻止.
沒有什麼可以阻攔他們.
願神都在這個當中.
繼續賜福洪恩堂.
我們同心的祈禱.
親愛的神.
你創造的心意是滿有智慧.

$^{521}$這個智慧是要我們緊緊的去跟從你.
去信靠你.
去依靠你的時間去學習.
你帶給我們一個新的觀點.
去看看今天的教會是可以怎樣.
你從中謙卑我們.
去接納不同的人.
不同的事.
求你繼續引領我們.
帶領我們在前面.
叫我們追隨著你的事工.
你的發展.
你的擴展.
我們看到你的賜福.
怎樣賜福這個大地.
求你這樣教導我們.
如果我們當中今天聽到一些訊息.
是需要你的力量.
讓我們去改變.
求神你真的去改變我們.
也寬恕我們從前.
讓你喜悅的事.
讓你聖靈的微星.
經常提醒我們.
要我們都行在你的正道裡.
願你鼓勵我們.
激勵我們.
幫助我們每一位.
不是頭腦上.
而是我們可以心體力行.
每一天每一個時刻.
都是活出你的那種豐盛.
我們亦都求你.
像是這個聖靈.
驅趕一些阻攔我們去親近你的事.
我們即將繼續每一天的生活.
求你賜福給我們每一位.
我們這樣奉主明求.
阿們.
\newpage



\section{提摩太後書 1:1-7}
\label{sec:lg5Gw9lhP5s}
\textbf{2024.02.25 《最有影響力的栽培員》 提摩太後書1\_1-7 (和修版) 講員 : 曾敏姑娘}
\newline
\newline
連結: \href{https://youtube.com/watch?v=lg5Gw9lhP5s}{\texttt{https://youtube.com/watch?v=lg5Gw9lhP5s}} ~~~~ 語音日期: 2024-02-28
\newline
\newline
\hyperref[sec:Yl0EevfFnJU]{\small{< < < PREV SERMON < < <}}
~
\hyperref[sec:index_chronic]{\small{[返順時目]}}
~
\hyperref[sec:index_scriptual]{\small{[返順卷目]}}
~
\hyperref[sec:DSsURlo4WPg]{\small{> > > NEXT SERMON > > >}}
\newline
\newline
提摩太後書 1:1-7
\newline
\begin{longtable}{cl}
\hline
\hline
章節 & 經文 (和合本修訂版)\\
\hline
1:1 & \begin{tabularx}{0.7\textwidth}{X} 奉神旨意，按照基督耶穌裡所應許的生命，作基督耶穌使徒的保羅， \end{tabularx} \\ \\ \relax
1:2 & \begin{tabularx}{0.7\textwidth}{X} 寫信給我親愛的兒子提摩太。願恩惠、憐憫、平安從父神和我們的主基督耶穌歸給你！ \end{tabularx} \\ \\ \relax
1:3 & \begin{tabularx}{0.7\textwidth}{X} 我感謝神，就是我接續祖先用純潔的良心所事奉的神，在祈禱中晝夜不停地想念你。 \end{tabularx} \\ \\ \relax
1:4 & \begin{tabularx}{0.7\textwidth}{X} 我一想起你的眼淚，就急切想見你，好讓我滿心快樂。 \end{tabularx} \\ \\ \relax
1:5 & \begin{tabularx}{0.7\textwidth}{X} 我記得你無偽的信心，這信心先存在你外祖母羅以和你母親友妮基的心裡，我深信也存在你的心裡。 \end{tabularx} \\ \\ \relax
1:6 & \begin{tabularx}{0.7\textwidth}{X} 為這緣故，我提醒你要把神藉著我按手所給你的恩賜再如火挑旺起來。 \end{tabularx} \\ \\ \relax
1:7 & \begin{tabularx}{0.7\textwidth}{X} 因為神賜給我們的不是膽怯的心，而是剛強、仁愛、自制的心。 \end{tabularx} \\ \\
[1ex]
\hline
\hline
\end{longtable}
$^{1}$.
我們一同禱告.
天父上帝是何等寶貴.
你讓我們有福氣成為你的兒女.
而且你上上設下你的話語.
說要成為我們腳前的燈路上的光.
當我們在此聚集的時候.
就求你對我們眾人說話.
願聖靈在我們心中動功.
賜我們智慧悟性去明白.
心願神親自對我們說.
又給我們力量去回應神的教導.
奉耶穌基督的名字祈禱.
阿們.
昨晚在這個地方.
我們有青少年的活動.
去到很晚.
昨晚一共有14位青少年.
當中大部分都是中學生.
有一位是大學生.
昨晚我們印象非常深刻.
因為大家在玩遊戲.
我們在唱詩歌.
我們的事奉人員昨晚去到很晚.
今天仍然在當中.
為大家的生命獻上感恩.
心裡也很歡喜雀躍.
我們有一段時間有斷層.
現在慢慢地可以開始這個工作.
除了興奮之外.
其實心裡也覺得有很大的責任.
因為不只是和他們玩遊戲.
認識建立關係.
是要帶他們信耶穌.
去栽培他們的生命.
我們感恩已經有導師在當中.
帶了幾位成員去信耶穌.
接下來還有一些中學生,大學生.
還沒有信耶穌.
我們還要繼續工作.

$^{41}$信耶穌之後還要栽培.
其實栽培是更加長遠.
所以今天想起栽培員.
想起這個題目的時候.
覺得是非常非常重要.
我們不單是需要栽培年輕人.
成年人需要栽培.
我們年紀很成熟的弟兄姊妹.
也需要栽培.
我們每個人都需要栽培.
究竟一個怎樣的栽培員.
是很有影響力的呢?.
今天很想和大家說.
究竟誰是栽培員呢?.
一定是傳道人.
習慕旁貸.
一定是成熟的弟兄姊妹.
或者說明是栽培員.
是不是呢?.
我希望今天完結這一堂課後.
大家心目中有一個答案.
誰是栽培員呢?.
耶穌!非常好!.
這個一定是最終極的栽培員.
耶穌基督會派誰去栽培.
我們身邊不同的生命呢?.
什麼?.
有心的弟兄姊妹.
剛才有人說什麼?.
聖靈!聖靈一定會栽培我們生命.
神一定會的.
神會在我們生命工作.
剛才我聽了一個很好的答案.
每一個人都可以.
非常好!.
我們去看看最有影響力的栽培員.
是誰?.
他有什麼質素?.
不是一般的影響力.
是很有影響力.

$^{81}$是誰?有什麼質素呢?.
今天我們透過一個人物的生命.
看一個樣板.
一個示範.
究竟生命的成長是怎樣的.
受到誰的影響呢?.
如果生命能夠茁壯成長.
在耶穌基督裡面.
有什麼人會影響他呢?.
提摩太!.
大家聽過提摩太的名字了嗎?.
聽過吧!.
一聽提摩太.
很自然就會想起誰呢?.
很自然就會想起保羅.
他的屬靈師父的屬靈爸爸.
在這裡稍稍說說提摩太的背景.
提摩太的爸爸是一個希臘人.
所以他爸爸不是猶太人.
媽媽反而是猶太人.
而且是一個猶太的基督徒.
提摩太的家鄉是在路斯德.
在羅馬帝國加拉太省境內.
其實是今天土耳其的地方.
原來提摩太的家鄉是在那裡.
提摩太本身的性格是怎樣的呢?.
他是比較溫馴的.
但同時也有些怯懦.
所以他一開始的時候.
不是一個剛強壯膽.
非常勇敢的人.
走到最前線.
一個勇猛的戰士.
他不是這些.
他有些怯懦.
而且身體又不是很強壯.
有少少腸胃的問題.
所以你看到.
他的情況就是這樣.
但你看到在聖經裡面.

$^{121}$一路看著看著.
看到他生命的成長.
看到他最後承載著.
上帝所給他的使命.
他成為一個很剛強壯膽.
一個能夠完成神的託付.
一個這樣的人.
一個很忠心惠報的人.
神很喜悅的人.
但你想起他起初.
他有些條件的限制.
有他的軟弱.
在整個過程.
生命的成長究竟是怎樣的一回事呢?.
我們又看看.
第一個人.
這段經文裡面一定有提到.
一定是保羅自己.
他屬靈的爸爸.
保羅我們都知道.
在聖經裡面.
他寫了多少卷書畫?.
沒有人夠膽說.
保羅書信有多少卷畫?.
27?.
13?.
定下來沒有?.
27,13.
保羅書信有13卷.
在聖經裡面有很多.
一個人.
寫了這麼多卷書.
在歷世歷代影響了多少.
弟兄姊妹的生命.
影響了多少人的生命.
這是上帝非常重用的一位僕人.
保羅.
保羅和提摩泰的關係是怎樣的呢?.
在提摩泰的前書裡這樣說.
寫信給大恩信主.

$^{161}$作為真兒子的提摩泰.
保羅和提摩泰有同一個信仰.
大家都是信耶穌的.
因著這個同樣的相信.
保羅和提摩泰可以成為父子的關係.
這是何等的寶貴.
不單是因為信仰的結連.
這裡也是在說那種信任.
那種親密度.
不是只是弟兄弟兄之間.
而是父親與兒子.
何等的親密.
提摩泰的前書是這樣.
提摩泰的後書更加厲害.
寫信給親愛的兒子.
親愛的.
保羅很愛提摩泰.
很喜歡他.
很信任他.
很欣賞他.
爸爸是怎樣看待兒子的.
不是普通弟兄之間的那種欣賞.
是一種很深很深的關係.
原來保羅與提摩泰之間的關係是如此親密.
我親愛的兒子.
保羅不是對每個信主都這樣說.
保羅也有稱提多為他的兒子.
但好像提摩泰那樣.
我親愛的兒子並不是每個人.
所以他與他的關係是更加親密.
保羅很愛他.
在經文中也看到.
他在這裡說.
就是接觸祖先.
用純潔的良心所事奉的神.
在祈禱中晝夜不停地想念你.
他很愛他的屬靈兒子.
愛到怎樣?.
晝夜不停地想念.
而且是禱告.

$^{201}$我經常為他禱告.
晝夜想念他.
一想起他的眼淚.
在這裡經文沒有實際說明什麼事件.
為什麼提摩泰會哭呢?.
我們不知道.
但肯定保羅知道.
保羅一想起.
他心裡的愛意又被激發起來.
我急切地很想見你.
是愛的表達.
很讓我滿心快樂.
他很掛念他.
真的好像育身的爸爸.
很掛念他育身的子女.
一直在說他多麼愛他.
多麼肯定他.
這位屬靈的爸爸和兒子.
那段關係何等親密.
引著信仰.
兩個人可以同心.
兩個人在靈裡有很多交流.
很多的結連.
就是這段關係.
在裡面.
你看到保羅對提摩泰有很多的栽培.
很多的栽培.
他有很多的鼓勵.
很多的勸勉.
就看回剛才的經文.
大家看回程序表也有這樣說.
保羅鼓勵他什麼呢?.
他鼓勵他一章六節.
為這個緣故提醒你.
提醒你要把神藉著我.
按手所給你的恩賜.
災如火挑黃起來.
保羅很想提摩泰挑黃他的恩賜.
很鼓勵他要更好的.
更好過用他的恩賜去服侍.

$^{241}$因為神給我們的不是膽怯的心.
是剛強人愛自制的心.
保羅在鼓勵他.
剛強壯膽人愛自制.
是何等的寶貴呢?.
屬靈的爸爸鼓勵他的兒子.
怎樣去事奉.
怎樣去盡用他的恩賜.
還有提摩泰後書三章十到十二節.
這裡大家看到.
跟我一起讀好不好?.
好,預備起.
但你已經追隨了我的教導.
行為,志向,信心,寬容,愛心,忍耐.
以及我在安提我.
以歌念路斯德所遭遇的迫害和苦難.
我忍受了何等的迫害.
但從這一切苦難中.
主都把我救了出來.
其實凡立志在基督耶穌裡.
敬虔道逸的也都將受迫害.
這段經文說.
保羅栽培提摩泰的生命.
他不單用言語去鼓勵.
其實他也以他的行為.
直接讓他看到自己的生命是怎樣.
是一個生命的板.
所以他說你已經追隨了我的教導.
那是可以是言語的教導.
也是生命的教導.
我的行為是走出來的.
保羅教不單是說.
你今天要有信心.
你今天要剛強壯膽.
你今天要對人有愛心.
不是這些,不單是這些.
是走出來.
他告訴你他的志向是怎樣.
他看到他的志向.
看到他的信心.

$^{281}$他寬容人.
真的讓他看到他的胸襟.
他的愛心.
他對人是怎樣的忍耐.
看到的.
走出來的生命.
所以他才說提摩泰是追隨了.
提摩泰是跟著他的行為.
是言行.
也是新教.
不單是說.
很多時候當我們去教人的時候.
別人就會看著你.
你有沒有活出你所教的呢?.
你的行為和你的說話.
是否相配合呢?.
人很自然看到的.
在這裡保羅是成功了.
他是這樣教.
他也是這樣走出來.
提摩泰是跟著他的.
不單是跟著他.
這些很好的生命的質素.
連同保羅自己所遭遇的迫害.
苦難.
忍受的迫害.
其實提摩泰也看到.
提摩泰也是跟著他去.
這是何等的寶貴.
在過程中主有拯救保羅.
深信主同樣帶領提摩泰.
保羅說得很好.
納至在基督耶穌裡敬前度日.
也會受到迫害.
他跟提摩泰說這番話.
是想鼓勵他.
想勸勉他.
你跟著我一起去敬前度日.
你也會受迫害.
我會這樣經歷.

$^{321}$你也會這樣經歷.
那要做什麼?.
我忍耐.
我堅持到底.
你也要忍耐.
你也要堅持到底.
我就是一個活生生的板.
所以保羅的生命是非常有力.
作為一個栽培員.
是很有力.
也很有影響力.
提摩泰看著他的生命.
他這樣實際走出來.
實際這樣經歷.
他就跟著他去.
一點也不抽象.
不只是舌頭上的說話.
所以這是很好的.
現教新教.
今天也想問我們在座的每一位.
在我們身邊.
我們有沒有提摩泰呢?.
或者我們自己是提摩泰.
我們有屬靈的保羅.
多好啊.
我們可以跟著屬靈的保羅.
去學他生命的功課.
如果我們自己是那位保羅.
我們栽培提摩泰的時候.
我們自己的生命有沒有活出來.
你說:我就是覺得我今天的生命.
還沒有成功.
我不夠膽去做這個保羅.
弟兄姊妹鼓勵你.
鼓勵你挑戰自己.
要做屬靈上的保羅.
當你願意這樣做的時候.
你的生命會成長得很快.
你不只是用口去教人.
你想想你的生命要怎樣活出.

$^{361}$你所學到的一切.
從哪裡開始學?.
每個星期你回來聽到.
哪怕你只聽到一點.
你那一點都要活出來.
不要只是跟講員說.
你真的講得好.
下次要問他哪裡講得好.
你能不能夠用得出來.
活出來.
用得出來就是你的.
神的話語.
當我們這樣吸收的時候.
成為你生命的幫助.
你就可以去影響別人.
當你去教人的時候.
別人就會看著你.
生命是如何.
我相信將來會有越來越多的人.
會起來影響別人的生命.
想想自己今天去到哪一個階段.
我的生命要怎樣去成長呢?.
好了,大家可能會說.
今天說了最有影響力的栽培員.
原來就是像保羅那樣.
可能是教會裡的屬靈領袖.
栽培員,或者傳道童工等等.
不是的.
這段經文並未停在這裡.
有兩個人物出現了.
是提摩太太的外祖母.
叫什麼名字呢?.
剛才我剛剛讀完.
叫羅爾.
外祖母羅爾.
和媽媽有尼基.
這兩個人物是非常非常重要.
原來他們是影響提摩太太的生命.
外祖母羅爾,媽媽有尼基.
他們在信仰的層面.

$^{401}$很深的去影響提摩太太.
提摩太太後書一章五節這樣說.
我記得你無偽的信心.
你的信心先傳在你外祖母羅爾.
和你母親有尼基的心裡.
我深信也存在你的心裡.
這裡兩次提到提摩太太的信心.
兩次.
第一次提到你無偽的信心.
之後又說我深信也存在你的心裡.
說的都是無偽的信心.
保羅很肯定很欣賞提摩太太.
那個信心是無偽.
說什麼信心?.
信耶穌基督.
提摩太太對耶穌基督那份信.
是無偽的.
有虛偽的假信心?.
有.
不是真的.
兩下就被打倒.
不是真真正正堅持地相信耶穌基督.
那些就是假的.
不是真的信心.
但這裡保羅欣賞她無偽的.
徹徹底底真實的.
而且是怎樣?.
這些字眼的描述很好.
這個信心是先存在.
而不是說我發現你愛祖母羅爾有這個信心.
是存在在她裡面的.
在她的心裡面代表什麼?.
是恆常地常常都有信心.
是她生命的一部分.
是常常都有信心.
很真實的.
這些叫做無偽的信心.
首先是在愛祖母羅爾那裡.
接著在她媽媽那裡.
一代傳一代.

$^{441}$又去到提摩太太那裡.
聖經裡面提得很好.
羅爾會影響有尼基的生命.
兩位又會影響提摩太太的生命.
提摩太太小時候已經在看希伯來的聖經.
在看舊約.
並且知道你從小明白聖經.
這聖經能使你印在基督耶穌裡的信有得救的智慧.
提摩太太小時候就學習聖經.
當然那不是整本的.
今天我們看的新舊約.
是舊約的聖經.
她那時候已經在看.
這麼小的時候看是誰教她的?.
她爸爸是希臘人.
是誰教她的?.
愛祖母,媽媽.
從小就教.
一直傳在她的心裡面.
成為她生命的一部分.
她很堅定地,持久地,恆久地去信耶穌基督.
是無偽的信心.
這個無偽的信心.
不是保羅放在她裡面.
是她家人栽培她的.
有血脈的家人栽培她的.
這是很寶貴的.
在《箴言》裡說.
教養孩童走當行的道.
就算到老他也不偏離.
是誰教養的?.
一定是教會傳道人.
教聖經的人.
他們教養孩童走當行的道.
是不是只有這樣呢?.
其實這裡是在說父母.
父母啊.
教養孩童走當行的道.
就算到老他也不偏離.
家是何等重要?.

$^{481}$是一個很重要的基督教教育場所.
今天說到教育.
我想說.
有很多家庭想將教育的責任.
第一,推給學校.
我將我的子女送到學校.
學校就要教好我的子女.
我子女回到家裡.
她就要成為一個.
不單止是知書識禮.
她還要懂得.
人如何行事為人.
每樣東西.
都是很優秀的.
很有品格的.
一個這樣的人.
如果我子女經過.
每個星期幾天.
在學校生活回來之後.
她還變差了.
她說什麼都好.
我覺得我要投訴學校.
是不是這樣?.
現在很多都是這樣.
有不少家長都是這樣.
覺得學校很大的責任.
沒有教好我子女.
第二個責任.
將她推給教會.
很重要.
我星期天將我子女.
放到你這裡.
你就要教好我子女.
星期天她回教會.
是回幾多個小時?.
大家跟我算一下.
一個多小時吧.
我們上面現在進行.
兒童團契.
崇拜.

$^{521}$很厲害的話.
兩個多小時.
早點回來.
坐在裡面.
有聖經的教導.
之後他們有些很開心的.
活動時間.
一個多小時.
兩個小時.
每個星期只有兩個小時.
她回到學校.
學校時間一定比教會長.
但最長的時間在哪裡?.
在家裡.
每晚.
家長都會見到子女.
我看著在座的弟兄姊妹.
有些是自己.
現在還在帶著.
小朋友的.
有子女.
有些已經是做爺爺奶奶.
有孫子的.
家裡一定有子女.
有小朋友.
你對著她的時間.
肯定比學校多.
對著她.
也比教會多.
對著她.
在那些時間.
就是你的時間.
是要教養孩童.
走當行的道.
聖經教的.
不行.
周牧師,你快點來我家.
我每一刻的時間.
全部給你.
我子女就是要周牧師.

$^{561}$教好她為止.
周牧師就想.
我可以去多少個家呢?.
每晚都要去教她的子女.
講一堂道.
她去不了多少個家庭.
真的去不了.
但是父母們.
爺爺奶奶們.
你卻每一天與你的子女同在.
你的孫兒孫女同在.
你就在她身邊.
你就是那個栽培員.
你說不要緊.
我每晚的時間.
全部都不會用.
我會與她玩.
我會建立親子關係.
我就等她.
星期日回到教會.
然後傳達人就教她.
或者那個道師就教她.
已經錯過了黃金時間.
教養孩童.
走當行的道.
就是道路也不偏離.
你看看.
提摩太的信心.
是怎樣開始建立的.
不是小時候.
很久不知在哪裡見到她.
然後捉住她來栽培她.
是她的外祖母.
她的媽媽才栽培她的.
她的根基是這樣回來的.
小時候.
不只是教她人生的道理.
教她聖經.
弟兄姊妹.
你有沒有教你的子女聖經?.

$^{601}$我想問你們有沒有教你們的子女聖經?.
怎樣教?.
羅爾和尤尼基怎樣教?.
活出來.
我一張嘴巴跟他們說.
那個聖經是怎樣怎樣.
怎樣愛.
怎樣去敬畏上帝.
另一方面.
我要活出來.
我去愛我自己.
去敬畏上帝.
我去愛自己.
我去愛我自己.
去敬畏上帝.
我的子女看到.
我的孫兒看到.
他們就懂得去敬畏上帝.
他們才懂得去愛上帝.
去愛人.
活出來.
同樣都是言教身教.
這是很重要的.
甚麼叫無畏的信心?.
我們再看看.
《雅各書》這樣說.
我的弟兄們.
若有人說自己有信心.
卻沒有行為.
有甚麼益處呢?.
我常常叫你.
有信心嗎?.
但我沒有行為.
你不會再聽我說話.
因為你覺得我的說話沒有份量.
信心也是一樣.
若沒有行為是死的.
只有一張嘴巴是死的.
可見信心是與他的行為.
相輔並行.

$^{641}$而且信心是因著行為.
才得以成全的.
真正無畏的信心.
你是可以活出來的.
人們透過你生命的果子.
透過你生命.
看到信心.
這份信心.
在提摩太的外祖母.
羅爾的裡面.
在她媽媽有尼基的裡面.
亦都進入到她自己的生命裡面.
她也有無畏的信心.
人們是透過她的生命看到的.
弟兄姊妹.
今天我們的生命如何?.
如果今天說要做一個.
最有影響力的栽培員.
我可以這樣說.
在座的每一位都可以做.
不單是有心.
難道沒有心的就不可以做?.
你打算將你的子女.
你的孫兒孫女.
全部扔給別人去教?.
你打算交給別人去教.
自己不去教?.
但在聖經裡說.
我們每個人都有責任.
你說我沒有育身的子女.
你有屬靈的兒女.
是嗎?.
我們永遠都不能這樣.
我永遠都做屬靈的兒女.
我不會做別人的.
屬靈的媽媽,父親.
這樣我們就不會成長.
如果要成長的時候.
一樣都是賢教,新教.
每一個都有份做栽培員.

$^{681}$每一個.
今天在座的每一位.
你們是栽培員.
如果你要成為一個有影響力的.
而且是很有影響力.
甚至是最有影響力的.
我們每一個需要神的話語.
看聖經,學聖經.
我們今年暫時沒有機會去推讀經.
去年有推讀經.
我們讀了多少聖經?.
就算今年.
我們的主力是說彼此相愛.
都需要聖經.
自己去學.
參加主教學.
活出來.
成為你生命的一部份.
預備好自己.
去影響別人生命.
所以成為一個栽培員.
不是不停上很多課.
我先接受裝備.
栽培員第一步是甚麼?.
第二步是甚麼?.
第三步是甚麼?.
第四步是甚麼?.
重點是你的生命.
今年現在是二月份的時間.
還在年頭的時間.
我挑戰大家.
鼓勵大家.
追求生命更大的成長,成熟.
你看看自己生命的範圍.
有哪些人會受到你的影響?.
你是他生命的保羅.
或者你是他肉身的爸爸媽媽.
肉身的爺爺媽媽.
請你要起來.
做一個很有影響力的栽培員.

$^{721}$栽培他們生命.
感受他們.
你們的子女.
兒孫們.
一路走在主的道路上.
一生不偏離.
這是一個最大最好的福氣.
我舉一些例子.
我曾經看過一本書.
其中一部份.
那本書叫做.
《孩子反叛有理》.
有聽過這本書嗎?.
這本書叫做.
《孩子反叛有理》.
是講述基督徒家庭的子女.
為何會反叛.
你說基督徒家庭的子女會反叛.
是的.
為何會反叛呢?.
有很多原因.
其中一個原因.
子女會看著父母的生命.
我父母很喜歡這樣.
星期日就說.
兒子,女兒,回教會了.
我最初也跟著他們回教會.
接著我發現.
其實我父母.
只把信耶穌這件事.
或者回教會.
當作宗教活動.
知不知道甚麼是宗教活動?.
活動跟真正的關係.
有甚麼分別?.
活動是有很大分別的.
可參加或不可參加.
對我的生命不一定有甚麼益處.
這是父母喜歡的活動.
跟父母和耶穌基督有一個.

$^{761}$真實的,親密的關係.
是兩回事.
在父母有困難的時候.
他看不到那份信心.
看不到那份倚靠.
看不到父母如何.
靠著耶穌基督.
去走過艱難.
艱難.
還有.
在聖經裡的話語.
如何影響父母的生命.
令父母的生命成長.
是不同的.
我看不到父母和外面的人.
生命質素有甚麼不同.
這樣.
作為父母.
很難影響子女的.
淑女生命.
是很難的.
因為你不只是現教.
是需要新教.
我們教會將會大大進行.
兒少青的事工.
沒錯.
教會一定會舉辦很多不同的活動.
會大大栽培兒童的生命.
栽培青少年的生命.
但是第一個戰場.
不是在這裡.
在哪裡?.
在大家家裡.
在你們家裡.
很重要.
子女兒孫在看著你們的生命.
通常我們帶人去教會.
是最容易的.
平時的生命.
能夠彰顯耶穌基督.

$^{801}$你去邀請人去教會.
是比較容易的.
因為別人覺得很有說服力.
你和外面的母信主的人.
是很不同的.
如果這樣邀請家人去教會.
是容易很多的.
同樣子女跟著父母也是.
我們的生命是很重要.
所以今天我很想.
透過提摩太的生命.
讓大家看到.
我和你都有份.
可能是栽培育生子女的生命.
可能是栽培.
屬靈晚輩的生命.
我們每個人都有份.
心願主使用我們每一個.
今天如果看回這段經文.
最有影響力的栽培員是保羅.
是外祖母的母親.
我也不敢說.
神也用他們幾個.
去栽培提摩太的生命.
你想想自己的生命.
是誰栽培我們的?.
我相信神也放了不同的人.
在我們身邊栽培我們.
我相信神將來會用我們每一個.
現在可能已經在用.
在座的每一個.
去栽培其他人的生命.
心願神保守我們.
建立我們.
兼顧我們.
讓我們成為最有影響力的栽培員.
重點是甚麼?.
言教.
心教.
活出信仰.

$^{841}$我們一同禱告.
天父上帝.
是何等寶貴.
從提摩太身上看到.
上帝你何等愛我們.
你為了讓我們生命成長.
總會將不同的人放在我們身邊.
在提摩太的生命裡.
我們至少看到.
由保羅成為他屬靈的父親.
一直帶領他栽培他.
由羅爾.
由有尼基.
去栽培提摩太的生命.
由他從小就認識聖經.
從小就認識上帝你.
主下你一直栽培提摩太的生命.
以至他能夠由一個.
有限制的軟弱的生命.
成為一個很忠心.
上帝你很喜悅的.
大大重用的僕人.
他有份於保羅的宣教.
他可以去牧養教會.
去教導.
做很多很多上帝你.
託付給他的工作.
這顆生命活著的時候.
就是永遠要神你自己.
主下今天我們看回.
你在我們的生命裡.
都行奇妙的事.
你猜探不同的人.
去栽培我們的生命.
以至到我們今天.
我們的生命.
都可以被你用.
我深信神要使用我們每一個.
去栽培別人.
在座的每一個弟兄姊妹.

$^{881}$主你都可以用.
所以求主叫我們今天開始.
更加的去體重我們和你的關係.
要我們更多的去閱讀聖經.
更多的思考.
更多的去活出上帝你的話語.
更多的是靠著聖靈去成長.
讓我們的生命成為別人的祝福.
主呀求你幫助我們.
幫助我們成為很優秀的栽培員.
幫助這裡這麼多個家庭.
父母們.
爺爺媽媽們.
都可以好好栽培他們的子女.
二孫的生命.
好讓他們一生.
走在主你的道路當中.
主呀我們心願未來的日子.
我們整個教會.
一同去經歷那個復興.
整個教會的弟兄姊妹.
我們一同去經歷生命的成長.
願你愛我們的心.
得到很大的滿足.
願你因為看到我們生命的成長.
而你得榮耀.
而你感到喜樂.
奉耶穌基督的名字祈禱.
阿們.
大門.
\newpage



\section{約翰福音 4:31-38}
\label{sec:DSsURlo4WPg}
\textbf{2024.03.03 《不一樣的看見》 約翰福音 4\_31-38 (和修版) 講員 : 梁友東牧師}
\newline
\newline
連結: \href{https://youtube.com/watch?v=DSsURlo4WPg}{\texttt{https://youtube.com/watch?v=DSsURlo4WPg}} ~~~~ 語音日期: 2024-03-08
\newline
\newline
\hyperref[sec:lg5Gw9lhP5s]{\small{< < < PREV SERMON < < <}}
~
\hyperref[sec:index_chronic]{\small{[返順時目]}}
~
\hyperref[sec:index_scriptual]{\small{[返順卷目]}}
~
\hyperref[sec:sd0jGt0zjkQ]{\small{> > > NEXT SERMON > > >}}
\newline
\newline
約翰福音 4:31-38
\newline
\begin{longtable}{cl}
\hline
\hline
章節 & 經文 (和合本修訂版)\\
\hline
4:31 & \begin{tabularx}{0.7\textwidth}{X} 就在這個時候，門徒求耶穌說：「拉比，請吃吧。」 \end{tabularx} \\ \\ \relax
4:32 & \begin{tabularx}{0.7\textwidth}{X} 耶穌對他們說：「我有食物吃，是你們不知道的。」 \end{tabularx} \\ \\ \relax
4:33 & \begin{tabularx}{0.7\textwidth}{X} 門徒就彼此說：「難道有人拿甚麼給他吃了嗎？」 \end{tabularx} \\ \\ \relax
4:34 & \begin{tabularx}{0.7\textwidth}{X} 耶穌對他們說：「我的食物就是要遵行差我來那位的旨意，完成他的工作。 \end{tabularx} \\ \\ \relax
4:35 & \begin{tabularx}{0.7\textwidth}{X} 你們不是說『到收割的時候還有四個月』嗎？我告訴你們，舉目向田觀看，莊稼熟了，可以收割了。 \end{tabularx} \\ \\ \relax
4:36 & \begin{tabularx}{0.7\textwidth}{X} 收割的人已經得工錢，為永生儲存五穀，使撒種的和收割的一同快樂。 \end{tabularx} \\ \\ \relax
4:37 & \begin{tabularx}{0.7\textwidth}{X} 『那人撒種，這人收割』，這話可見是真的。 \end{tabularx} \\ \\ \relax
4:38 & \begin{tabularx}{0.7\textwidth}{X} 我差你們去收你們所沒有辛勞的；別人辛勞，你們享受他們辛勞的成果。」 \end{tabularx} \\ \\
[1ex]
\hline
\hline
\end{longtable}
$^{1}$好 大家平安.
我高興今天跟大家分享神的話.
今天來到西貢漢羅營.
一個郊外的地方.
今天編導我也是很切合.
是郊外來講.
因為今天講的經文.
處境是在一個郊外.
剛才說有田有野.
還有一個井.
一會兒我會講.
不知道在人生中.
你覺得什麼是很重要的呢.
很多人都覺得錢很重要.
還有很多東西很重要.
今天想跟大家分享.
什麼是很重要.
主題叫不一樣的漢奸.
剛才說在約翰福音四章.
耶穌和門徒經過薩馬利亞的地方.
當時是正午.
正午的時候很熱.
他們走路困倦.
門徒去買食物.
耶穌坐在井旁.
那時有一個薩馬利亞婦人.
中午來打水.
耶穌問這個婦人拿水喝.
這個婦人說你是猶太人.
我是薩馬利亞人.
猶太人和薩馬利亞人沒有來往.
好像以色列和巴勒斯坦.
都有敵對關係.
門徒回來也看到.
耶穌和這個婦人聊天.
覺得很奇怪.
打破了當時的傳統.
男人和女人要保持距離.
耶穌問她拿水喝.
這個婦人說井很深.

$^{41}$你有沒有打水.
為什麼要拿水喝.
耶穌說.
如果你知道神的恩賜.
你一早求祂.
祂就會賜給你活水.
喝了永遠不渴.
你喝這些水.
仍然會渴.
但喝了神賜的活水.
就永遠不渴.
這個婦人就問耶穌.
讓我喝吧.
耶穌就問她.
叫你丈夫來吧.
這個婦人如何回答耶穌.
她立即防衛機制.
我沒有丈夫.
耶穌說.
你之前有五個丈夫.
現在這個丈夫也不是你的丈夫.
這個婦人就很驚訝.
你是不是先知.
你知道我的事.
你們禮拜在這個山上.
在耶路撒冷.
我們禮拜就在這個山上.
耶穌說.
遲到了.
今天拜父.
不是在這裡.
今天拜父就像大家剛才敬拜.
我們用心靈和誠實.
來敬拜.
這個婦人就說.
莫非你就是.
要來的救主.
來賽亞.
耶穌就說.
換圖回來的時候.

$^{81}$看到耶穌和這個婦人.
聊天也很奇怪.
然後.
換圖去買食物.
買食物回來就給食物耶穌.
耶穌說.
我有食物.
是你們不知道.
剛才說耶穌有吃過東西嗎.
沒有吃東西.
耶穌說.
我的食物就是進行猜我來者的旨意.
作成他的功.
耶穌接著說.
看周圍.
莊稼已經熟了.
可以收割了.
但當時的門徒說.
不是啊.
我看不到.
還有四個月才收莊稼.
耶穌說不是啊.
事後到莊稼已經熟了.
他說收割的人.
即從五穀到永生.
叫撒種的.
日同快樂.
一會兒我講的是.
關乎到剛才.
耶穌如何與.
撒瑪利亞婦人相遇.
撒瑪利亞婦人不單這樣.
當她知道.
耶穌是一位要來的救主.
她就放下水桶.
走去城裡.
和城裡的人說.
撒瑪利亞城裡的人說.
這個.
將我過往的事.

$^{121}$說出來.
其實過往的事.
是甚麼事.
說出來.
莫非他就是那位要來的救主.
人們都奇怪.
救主?我也要認識他.
走去城裡.
找耶穌.
耶穌就留在那裡.
和城裡的人說.
說甚麼?當然是說福音.
說救恩.
結果.
整個城裡的人很多人都相信.
耶穌就是要來的救主.
這位救主.
今天我就開始和大家說.
耶穌所看見的.
和我們今天所看見.
所看重的有甚麼不同呢?.
第一個的看見.
就是超越身份的看見.
剛才說猶太人和撒瑪利亞人.
從來都沒有來往.
甚至是一種敵對.
耶穌是怎樣去看人呢?.
耶穌看人不同.
很多人說.
先見來醫.
就是後見人.
有沒有地位.
有沒有身份,有沒有錢等等.
耶穌不是.
耶穌去看人.
在神的眼中.
每一個人都是神所愛,神所寶貴.
超越種族,階級.
性別甚至仇恨.
如果你看福音書.

$^{161}$你看耶穌最喜歡和甚麼人在一起呢?.
罪人.
大患麻瘋病的人.
傳染病.
他都會親近他們.
接近他們.
所以神看我們每一個人.
在神的眼中.
都是寶貴,都是神所愛.
你覺得你自己寶不寶貴?.
你和隔壁的人說.
你很寶貴.
真的.
我們每個人的生命都是很寶貴.
神看我們.
甚麼階級,甚麼性別.
神看我們都很寶貴.
甚至耶穌是怎樣?.
愛了那些罪人.
當時的瑞利人.
人人都很憎恨.
他是罪人.
反而耶穌是怎樣?和罪人在一起.
當罪人親近耶穌的時候.
他們的生命就被改變.
生命就被改變.
所以今天.
我們怎樣去看人呢?.
今天你怎樣去看自己呢?.
我告訴你.
耶穌很愛你.
神很愛你,神看重你.
等等.
耶穌是怎樣?.
超越了受助者和幫助者的關係.
我們不是由上而下.
我們幫你.
耶穌問那個婦人.
拿水喝.
不是一種.

$^{201}$高高在上.
我來幫你,施捨你.
不是.
人與人之間.
彼此之間的分享.
生命就是一種分享.
很重要的.
今天在我們的福事裡.
都同樣很重要.
這段經文.
講到水和食物.
我講水的故事給大家聽.
好嗎?.
有一次我去到.
香港有些偏僻的鄉村.
有一次我去鄉村暴動.
去到一個寮屋.
住了一個獨居長者.
這個獨居長者.
有一點殘疾.
咳嗽.
我們去探望他.
探望他的時候.
他就接待我們.
進入了那間寮屋裡.
就倒水給我們喝.
那杯茶壺有很多汁.
那杯茶壺也有很多汁.
茶有很多汁.
倒水出來.
是冷的.
不知道放了多久.
讓你喝嗎?.
考慮一下.
當我喝了.
當我拿起杯子.
喝下去的時候.
伯伯就說.
他們不喝我的水.
原來這個伯伯.

$^{241}$可能生命有很多坎坷.
如果不是獨居在.
鄉郊的寮屋裡.
原來這個伯伯.
以前是做.
做三行的.
做建築.
開始去聽聽.
聽聽他的故事.
這條街我建過.
這棟大廈我建過.
伯伯你好厲害.
你今天香港有這樣的.
繁華.
建設.
其實我們上一代所付出的.
血汗,勞力.
以至今天我們得到的成果.
伯伯就很開心.
有人肯定他.
剛才說伯伯.
他不喝你的水是甚麼.
原來伯伯年紀大.
沒有工作能力.
自己很坎坷.
要靠種援金.
生育過活.
有時保障部.
有些人去探望他.
幫他搞一些東西.
給他喝水.
他們喝水他總是不會喝.
令到伯伯有一種.
好像.
我現在沒有用了.
人生年紀大了.
有一種乞憐的感覺.
以後人就去周濟了.
其實.
當我和伯伯說.

$^{281}$肯定他.
其實不是.
今天我們的繁華繁盛.
都是你所付出的努力.
你應該被肯定.
其實值不值得肯定.
當然值得.
這杯水.
這杯茶.
就進入了我們彼此間的生命.
我喝了的時候.
開始.
伯伯說完他的故事.
就說我們的故事.
我們的故事當然是說耶穌的故事.
很感恩我們帶了伯伯去.
信耶穌.
很多人去探望他.
他都很開心.
有一次他進了靈石醫院.
我們去探望他.
伯伯說.
他能夠和我們說.
耶穌愛我.
我愛耶穌.
同樣今天我們怎樣去看人.
教會服事當中.
我們是不是用一種.
心態來服事.
我們的靈石.
或者服事有需要的人.
如果我們往往.
只是一種很慈慰式.
很高高在上.
其實不舒服的.
中國人有句說話.
得人恩果.
千年報說他不能夠去.
參與投入的時候.
關係都是有一種距離.

$^{321}$如果假如.
我們能夠人與人之間.
彼此能夠分享.
就像耶穌一樣.
向婦人拿水喝.
耶穌進入罪人家中.
被他款待被他接待的時候.
我們的關係就很不同.
是不是.
如果有一天你在服事你的靈石.
靈石請你去家裡吃飯.
你吃不吃.
當然要吃.
我們都做很多新來港的工作.
除了我們.
有時候都會服事他們.
都給他們一些回饋.
在國內帶一些手信回來.
可能對你來說.
都不是很需要.
你接不接納呢.
我們一定要接納.
一定要接納.
接了之後.
彼此之間.
就有一種分享.
所以今天在我們服事靈石.
裡面.
我們會看中每個人.
在耶穌裡面都是寶貴.
我們每個人都是.
生命裡面彼此的分享.
當我們的靈石.
進入到我們的教會當中.
當他能夠找到自己的角色.
找到自己的身份的時候.
能夠自己有貢獻的時候.
他就能夠建立一種歸屬感.
他來的就不像.
我來是拿東西.

$^{361}$我來只是被服事.
所以有時候.
他們不同鄉下.
來的有不同的小吃.
不如你來做小吃.
一起吃好不好.
一起有機會去參與的時候.
他們就會很開心.
找到自己的價值.
自己也能夠有貢獻.
他就能夠建立一種歸屬感.
歸屬感.
前兩天.
我們有一個.
籌款晚會.
有一間教會的牧者來分享.
他就說是了.
接觸到一些街坊.
包餃子.
一起去吃.
他就很開心.
感覺到這是一個家.
我們很欣賞.
宣道威錦繡堂.
他們的分堂在哪裡.
洪水橋.
是一個分堂.
還有一個是加州花園.
你知道加州花園做什麼嗎.
我們有些服事.
加州花園的人.
菲律賓.
印尼.
早前我去港島.
他們感覺到.
跟一個家一樣.
他們有一個廚房.
他們經常在廚房.
一起聊天.
煮東西.

$^{401}$就像一個淑齡的家.
洪恩堂也是.
第一個超越的.
體見就是超越身份.
的體見.
超越的體見是什麼呢.
物質的體見.
在剛才所讀的經文中.
耶穌和撒瑪利亞婦人.
談論到.
在亞國井裡喝的水.
門徒又帶食物回來.
又提到食物.
我想問.
食物和水重要嗎.
當然重要.
是維生的.
沒有水沒有食物我們會死.
食物和水也帶出.
貧窮的問題.
香港有沒有貧窮呢.
有.
香港有百四十多萬人.
在貧窮線以下.
我們的機構也做很多扶貧的工作.
我們有食物銀行.
食物援助.
緊急援助基金.
物資的賄送.
幫助很多貧窮家庭.
各方面的工作.
貧窮是不是單單物質呢.
不是的.
雖然.
水和食物都是重要.
但耶穌所看到的.
不是單單這樣.
當耶穌和撒瑪利亞婦人.
談論到.
喝的水.

$^{441}$然後耶穌說.
如果你喝了我賜給你的活水.
在你生命中永不渴.
永不渴.
耶穌所看到的.
這婦人不是單單喝的水.
而是她生命中的飢渴.
是什麼呢.
耶穌問她.
叫你丈夫來.
她說我沒有丈夫.
之前我有五個.
現在這個也不是.
意思是撒瑪利亞婦人.
在她生命中.
非常破碎.
她看望的是很多感情關係.
她看望可能有人醫護.
有人保護.
能夠帶給她安全.
但可能一段感情.
一段關係.
帶給她很多生命.
反而是破碎.
傷害.
人生命中.
都有不同的飢渴.
這些飢渴.
如果不是耶穌的活水.
永留在她生命中.
她仍然會渴.
你喝多少都會渴.
就像人喝海水.
越渴.
但很多人的飢渴.
有不同層面.
有些人要賭錢.
有些人要感情關係.
有些人要權力.
有些人要吸毒.

$^{481}$等等.
很多東西都被這些東西.
捆綁住.
這個婦人.
在感情上.
受到種種的創傷.
所以耶穌看到她的需要.
她的渴.
不是感情滿足.
她的渴.
是需要耶穌的活水.
永留在她生命中.
叫她不再渴.
其實今天我們每一個人.
心靈中都有一些不同的渴.
不同的貧窮.
我們說說貧窮.
有什麼貧窮.
這個婦人的貧窮.
不單是物質的貧窮.
這個婦人也有什麼.
關係的貧窮.
為什麼她中午出來取水.
你會不會中午出來取水.
一定不會.
巴勒斯坦地區.
中午是很曬很熱.
你不會中午出來取水.
因為她想避開人.
人們通常.
都是早上去取水.
太陽不猛不曬的時候.
去避開人.
可能她會.
她覺得人會.
很多閒言閒語.
這個不知道是什麼女人.
是不是.
她感覺自己.
生命中很感到.

$^{521}$一種羞恥感.
所以這個羞恥感.
令到她和人之間的關係.
是怎樣的 破裂的.
你看看.
耶穌的活水湧流在她的生命中.
她不再渴的時候.
你看看這個婦人的改變是怎樣.
她扔下桶子.
回到城裡.
告訴城裡的人.
是怎樣的.
她將過往不見光的事.
不想讓人知道的事都說出來.
莫非她就是那位救主.
她再能夠面對自己.
她不再羞恥.
她從那種.
釋放出來.
她和人之間的關係.
可以建立一個和好的關係.
她不再.
看人的眼光.
來看我.
其實人很多時候會看別人的.
社會的價值觀.
來看自己.
但剛才說的.
如果你知道耶穌很愛你.
不論你過往做過什麼.
神都願意愛你.
神都願意來拯救你.
就像薩瑪利亞婦人一樣.
她生命的改變.
她將她帶回群體裡.
我自己是做一些新移民的工作.
你知道新移民來到香港.
人生路不熟.
來到異議的文化當中.
很多的適應.

$^{561}$工作的適應.
兒女的適應.
家庭的適應.
我們的福祉就是.
能夠將新移民帶到哪裡呢?.
神的家.
屬靈的家裡.
讓他們能夠教會.
成為他們的陰秘地方.
因為教會就是一個家.
教會是一個.
對人能夠有愛.
能夠接納.
有包容.
因為有基督在我們當中.
有神在我們裡面.
所以.
耶穌解決了婦人的關係的貧窮.
第三個貧窮是甚麼?.
動機的貧窮.
中國人有句說話.
人窮.
接著是甚麼?.
人窮志短.
有人說積極一點.
人窮志必窮.
但事實上不是.
人窮志短.
有一次.
在青少年.
香港剛才說有貧窮.
在天水圍住.
他們家裡人都是窮的.
拿縱門.
問小朋友.
你未來的心志是甚麼?.
你猜是甚麼?.
拿縱門.
真的.
問過.

$^{601}$在貧窮家庭裡面.
成長裡面.
他們真的失去那種志氣.
原因是甚麼?.
有錢的人.
可以怎樣?.
去補習.
學一樣.
現在入學校.
都很看你的資料.
你學過甚麼?.
學過甚麼很貴的.
我們說.
跨代貧窮.
就是這樣.
因為他們在家庭的經濟.
原因下.
沒有機會向上流動.
有人做過調查.
這些在貧窮家庭裡面.
成長的小朋友.
他們上大學的機會.
是比較弱的.
所以在我們的服務裡面.
怎樣能夠解決這種.
動機的貧窮或跨代的貧窮.
就要有幫助.
甚麼?下一代.
今早.
和曾牧師談過.
教會.
紅印堂.
我們會開展青少年的崇拜.
我們開始去服務.
貧窮家庭的青少年.
幫他們怎樣?.
補習.
提升他們的學習動機.
動力.
剛剛來了一個中五的學生.

$^{641}$可能要考DSE.
外面補習很貴.
教會願意幫他補習.
他們回到團契.
決志了順主.
是不是?.
我自己教會也是.
我記得在十多二十年前.
我們開展了新來港的工作.
和我們教會去tong 房.
探望一些青少年.
很小的少年人住tong 房.
很困難.
我們就幫他們補習.
陪伴他們一起成長.
十幾二十年後.
他們已經成長.
已經做老師.
做專業人士.
現在他們回饋.
去服侍教會其他基層的人.
我幫他們做.
婚禮的時候.
他們很感恩.
感恩什麼?.
不單是上帝的恩典.
他們感恩的是教會.
能夠陪伴他們一起同行.
以至他們能夠經歷上帝的恩典.
生命能夠不一樣的改變.
關於動機的貧窮.
但最重要的貧窮是什麼?.
心靈的貧窮.
這個婦人.
撒瑪雅婦人.
她就是心靈的貧窮.
是不是?.
耶穌說:唯有耶穌所賜的.
活水湧流在她的生命裡.
叫她永遠不渴.

$^{681}$結果這個婦人.
真是被耶穌的活水湧流.
她不再渴.
因為耶穌進入她的生命.
她認識耶穌.
就是那位.
要來的基督,那位的救主.
在我們福事裡.
同樣都是.
不是單單一個社福機構.
一個社福機構提供服務.
完了就完了.
但我們所看到的是.
很多人的生命需要耶穌.
需要福音,需要耶穌的活水.
能夠湧流在我們生命中.
叫我們不再渴.
兄弟姐妹朋友.
今天你心靈裡有沒有饑渴呢?.
你有哪些不滿足呢?.
我相信.
唯有耶穌的活水湧流.
在你生命裡.
你經歷耶穌的愛.
他就不會再渴.
就好像撒瑪雅婦人一樣.
好了.
第三個.
超越的體見.
超越傳統的漢奸.
當撒瑪雅婦人和耶穌.
話題一轉.
轉到你們禮拜在.
耶路撒冷,我們禮拜.
在山上.
耶穌就和她說.
「夫人,你當信我」.
「時候到了」.
「你們拜父不在山上」.
「也不在耶路撒冷」.

$^{721}$「真正拜父的」.
「要用什麼?一起讀」.
「因為父要這樣的人害怕」.
今天有很多傳統.
「你禮拜一定要在耶路撒冷」.
「我們禮拜在山上」.
耶穌不是這樣看.
不是那種形式主義.
「我每個禮拜」.
「這段時間就回來」.
不是的.
耶穌超越了那種時間.
地域.
所有種種的限制.
是不是?.
今天同樣是.
時代在變遷.
如果教會不變.
我們就會.
走在時代的後面.
有一本書想大家認識.
相信大家認識.
藤斤威牧師.
他應該是五十年前寫的.
一本書叫「時代的挑戰」.
在第一章說到.
時代的夾縫和紅衣.
他說到.
如果教會沒有聽到時代的呼聲.
沒有回應時代的需要.
他就會跌在時代的夾縫裡.
他提到.
在四五十年前.
怎樣做學生工作.
那時提到突破.
蔡元文醫生.
蘇銀泰姐妹.
突破時刻.
突破雜誌.
那時提到學生分團契.

$^{761}$FES.
他提到藤斤威牧師怎樣做學生工作.
他說你沒有聽到時代的呼聲.
沒有回應時代的需要.
你就會跌在時代的尾巴後面.
教會就會窒息.
今日香港都是教會.
面對一個很大的時代的變遷.
教會你願不願意.
有一種更生和改變.
他提到當時不單止突破.
年青人的工作是怎樣.
我覺得藤斤威牧師.
是一個很有先知的體傑.
他提到當時做學生的工作.
是甚麼.
是開咖啡,茶座.
五十年後.
我們才知道要開茶座.
所以在.
在嘉欣堂又做過.
現在很多家教會.
有咖啡室.
這個就是.
有咖啡室.
有氣泡.
你為何要叫年青人喜歡去哪裡.
去starbucks.
去吃東西.
去甚麼.
藤斤威牧師已經看見.
是不是.
藤斤威牧師做基層工作要怎樣做.
他提到.
工業呼音團契.
我老師劉達芳博士.
四五十年前他就在功福.
去哪裡做.
那時是有工廠.
進入工廠裡面.

$^{801}$做女工.
進入到裡面.
很枯燥.
做電子.
一度到一度敬拜.
嘗試過當中去接觸.
工廠女工.
跟她們傳呼音.
我想問.
今日香港還有沒有工廠.
工業呼音團契存不存在.
有.
只不過他們.
焦點不在工廠.
焦點在貧窮裡面.
焦點在香港做了很多賭博的問題.
做了很多.
街賭的工作.
機構都不斷的.
去更新去改變.
去回應我世代的需要.
新福都是.
我們九七年成立.
當時.
我們已故的總幹事李健華牧師.
去四十日禁食禱告裡面.
領受關心從國內來的新移民.
特別是.
貧窮基層的新移民.
這二十多年裡面.
我看見上帝的恩典同在.
他的預備和供應.
去回應在世代裡面.
那個的需要.
今日香港都在變遷.
如果我們.
教會不去改變.
我就會走在時代的尾巴後面.
譚牧師是這樣說.
跌在時代的夾縫裡面.

$^{841}$超越了耶穌說.
不是在這裡在那裡.
用心靈誠實.
來敬拜.
怎樣讓人能夠與上帝相遇.
製造這種環境空間.
是很重要.
有一次我講兩個宣道會的.
見證給大家聽.
大家都是宣道會.
有一次我去大埔的宣道會.
去港道.
港道的時間不是星期一.
港道的時間是星期四上午.
來的一班人都是一班婦女.
為何你星期四上午崇拜.
不是星期日.
星期六崇拜.
原因是.
他們有一班信了主的婦女.
但她丈夫沒有信.
如果星期日回到教會.
丈夫不跟你回去.
不信的時候.
那怎麼辦.
回到教會會令到家庭.
要家庭聚會.
要家庭的親子.
家庭丈夫的家庭.
不如這樣.
就在星期四.
上午崇拜.
那裡有幾十人.
大部分是婦女.
她的丈夫很多都未信主.
他們就可以去敬拜上帝.
可以親近上帝.
是不是超過了那種時間.
傳統的限制.
接著.

$^{881}$另一間教會.
可能大家都認識.
在屯門.
屯門宣導會.
有一次我去教會.
星期四.
我們有一個波亞斯的計劃.
他們有一個姊妹.
新內港的姊妹.
信了主.
很追求參加波亞斯婦女.
領袖培訓計劃.
因為我也要見一下牧者.
去到我覺得很奇怪.
他們有很多房間.
每間房間都有小組.
晚一點就有派飯.
派飯完之後.
有聚會.
給予很多長者.
我們不只是滿足.
物質層面的需要.
其實需要心靈.
怎樣能夠創造一個空間.
讓人與上帝.
能夠相遇.
今天我們說無牆的教會.
是不是.
好像.
淑娟牧師有來過港島嗎.
淑娟牧師.
村道堂.
不知道有沒有說過.
教會有六千多呎.
聚會的人只有幾十人.
曾經新連崇拜.
只剩十幾人.
淑娟牧師來到.
開放.
來找我.

$^{921}$問可不可以和新方合作.
當然可以.
怎樣合作.
在社區裡.
看到媽媽和小朋友.
在玩.
應該上學前班.
為何不上.
因為沒有錢.
新移民又不懂說話.
早上開學前班.
做媽媽的工作.
中午吃飯.
大角咀區有很多基層.
獨居長者.
清潔工人.
等基層人士.
很需要吃飯.
開了飯堂.
百多人來吃飯.
下午看到.
區裡有很多學生需要.
如果有錢.
我介紹你去補習.
如果沒錢就來我們補習.
特別關心.
學習需要.
有學習需要.
基層家庭.
都不知道怎樣教導子女.
他們去服事.
和我們合作很多婦女的工作.
那時疫情前.
婦女來煮飯.
星期六就預備.
客人要一味款待.
預備星期日.
二百多人吃飯.
整個教會.
短短一段時間.

$^{961}$被復興起來.
連廚房.
都做了煮日學.
一日在教會.
的人流是百多人.
教會.
我們願不願意.
將新酒.
放於新皮袋.
有這種思維.
甘願改變.
如果你沒有回應時代的需要.
就像譚牧師說.
你會跌在時代的夾縫裡.
時代的尾巴.
求主給我們.
一種願意更新改變的.
一種精神.
一種態度.
好,接著我們說.
超越責任的體見.
門徒買食物.
回來的時候.
給耶穌食物嗎?.
耶穌說.
我已經有食物了.
是你們不知道.
門徒真的很奇怪.
誰給食物給耶穌呢?.
我的食物就是.
進行猜我來者的旨意.
作成他的功.
兄弟姐妹.
我們每個人都有工作.
是不是?.
你不是只打好這份工呢?.
只是說.
當然啦.
做事一定要盡上工作的責任.
但在你的工作裡.

$^{1001}$在你的人生裡.
你看不看到.
是不是在行神的旨意.
在作神的功.
在你人生裡.
你有沒有看到人生的照明.
是甚麼呢?.
很多人.
一生去追求的是甚麼呢?.
金錢.
你看看這幾年.
這兩三年裡.
世界的改變.
曾經是.
首富的.
恆大.
首富等等.
一失那間.
都會怎樣呢?.
過去.
我們所追求的.
是不是一桶金兩桶金.
人生是不是.
金錢為了很多東西.
你看重的是甚麼呢?.
耶穌提醒門徒.
他所看重的是甚麼呢?.
我們今天所做的.
我們所看重的.
就是很多的甚麼呢?.
莊稼熟了.
可以收割了.
莊稼是指甚麼呢?.
生命.
在馬太福音九章.
耶穌說.
就是要收的莊稼.
作工的人是怎樣呢?.
少啊.
耶穌看見很多失喪的靈魂.

$^{1041}$今天我們有沒有.
看到耶穌所看見.
當耶穌走遍各城各鄉.
去宣講去服事.
他就有一種看見.
很多人困苦流離.
像羊毛牧人一樣.
他就憐憫.
同樣今天教會需要一種憐憫的看見.
不是我們只是一種責任交代.
有時教會.
都會面對一樣甚麼呢?.
說了就等於做了.
聽了就等於做了.
會不會啊?.
有時我都會參與.
每年教會都有不同的主日.
猜拳主日.
請宣教士來講.
聽完教會就做了宣教.
奉獻一下就做了宣教.
我們報道主日.
街頭報道.
就去街頭報道.
一日.
但如果你看重生命.
生命不是一種責任交代.
不是一種工作.
不是一種應付.
而是你有沒有將生命.
放在心裡.
很緊張.
你不收割就會怎樣?.
就會發白,就會殘.
我們粉嶺有一個農莊.
有時我都很心急.
種了的東西.
你不收割,它會怎樣?.
纖維化,會爛.
你的心血甚麼都沒有.

$^{1081}$今天我們.
所積蓄的是甚麼?.
積蓄五穀道永生.
那些不衰殘,不扭壞,不沾污的.
存留在永恆的基業裡.
在我自己的服事裡.
詞窮裡.
三四十年.
應該有.
我會看到.
我的滿足.
不是我賺到多少.
我的滿足是.
我看到上帝的工作.
各音的大能.
怎樣將人生命轉化.
我怎樣看到神的榮耀.
在我們教會.
在我們當中彰顯的時候.
你覺得那種喜樂.
那種滿足.
不是你在地上.
所能夠賺取的東西.
所能夠滿足.
我們每個童工.
都能夠帶著一份使命.
說回屬於牧師教會.
早期的執事.
都說.
好好預備編導給我們.
餵養我們.
但今天他們教會改變了.
每個人都找到.
人生的照明.
每個人都有他的位份.
甚至他的影響.
影響到在澳門的.
宣道堂.
他們的長輩.
坐船過來.

$^{1121}$去宣道堂.
做服事.
那種使命.
今天在教會.
每個人都找到人生弟兄姊妹.
找到人生的照明.
無論你是長者也好.
無論是什麼也好.
神給你的恩賜.
神給你的才幹.
你都可以來發揮.
你都可以來.
去做服事.
去建立.
盼望我們的教會能夠成為一個.
什麼?神喜悅的教會.
盼望我們在裝甲.
已經熟悉.
我們都成為收裝甲的工人.
看完.
還有沒有.
很重要的一件事.
超越個人的成就的體建.
無論是殺還是收.
是不是.
有些人都很功利.
很功利.
你做新移民.
在舊區特別.
做了.
有沒有收穫.
教會會不會增長.
教會不增長.
那怎麼辦.
舊區住tong 房.
三五七年會搬去新市鎮.
是不是.
那他就回不了教會.
那你做不做.
確實有些教會.

$^{1161}$選擇不做.
沒有再生能力.
做青少年好一點.
為什麼要做新移民.
為什麼要做長者.
如果我們用一種很個人的.
個人成就.
個人的.
沒有看到神覺的胸襟.
神覺的體建.
殺和收.
你享受人奴婦所得的.
你看不到神的心意.
剛才曾牧師和我聊天.
說到在.
洪福村.
有一間教會.
在香港的.
生命堂.
每個星期四.
有二十幾個弟兄姊妹.
來到洪福村.
關心.
探望家庭.
很想去服侍.
借了一個.
餐廳.
我們也有和他們聯繫.
接觸服侍.
他們很想去服侍家庭.
來和.
曾牧師聊天.
大家聊得很開心.
很開心.
看見神的覺.
無論我做了多少好.
或者將來.
歸入朝人生命堂.
開放.
總是.

$^{1201}$不是競爭.
為了神的覺.
我們彼此都能夠.
貢獻自己.
彼此我們為了.
神的旨意成就.
神的功德而達成.
這是我們很需要的.
一種胸襟.
我們看最後一個.
超越困難的體見.
門徒說.
你們豈不說.
到修葛的時候.
我告訴大家.
舉目向田官.
莊家所可以修葛.
今天我們所看見.
在我們社會的變遷.
人只會看見困難.
但你知不知道.
困難背後.
帶有很多需要.
需要我們回應的時候.
正正是我們服侍的契機.
正正如新福所說.
在97年回歸.
經歷金融風暴.
釋法風波.
種種很多需要下.
反而我們回應上帝.
給我們呼召.
你就會經歷上帝的困難.
人在很多困難當中.
正正是需要服任.
在安逸裡.
不會帶來復興.
在很多困難裡.
人不能再靠自己.
需要靠上帝.

$^{1241}$今天香港面對變遷.
教會面對很多困難.
很多人移民.
我們下一代.
好像在斷層一樣.
很多需要在這裡.
這些需要正正是我們服侍的契機.
需要求神給我們一種信心的體見.
而這種契機是一種迫切性.
有一種迫切性在裡面.
耶穌說裝嫁淑女.
你不收就會怎樣.
就沒有機會.
所以今天我們要把握這個機會.
洪福村開了多少年.
大家在.
幾年了.
八年十年了.
都是一個新的時辰.
在我這裡.
接下來都會繼續有.
有些新的發展.
是一個新的機遇.
你來到的時候.
你去關心他.
他來到面對很多適應.
服侍他的時候.
他只會心容而開.
如果你不服侍他.
三五七年他已經融入了.
適應了.
你再跟他說.
在香港今天你移民了去外國.
在香港沒有信主.
去到外國.
我當那裡有教會來關心他們的時候.
信主的機會都比較大.
高.
今天我們有沒有把握住.
這個機遇.

$^{1281}$機遇過去了就不會再有.
求神給我們一種體見.
有一種體見.
最後我想說的見證.
在我的同工當中.
有些以前是受助者.
現在成為我們的同工.
關心的.
有一個來自天水圍.
因為元朗.
天水圍的一個姐妹.
前兩天.
我們在一個感恩會.
一個籌款晚會去分享見證.
她來自東北.
在家裡曾經四年.
四年不出來.
自己有抗躁性.
自己有抑鬱.
經常打兒子打女兒.
兒女形容媽媽是恐龍.
是獅子.
她曾經沒有出門.
四年.
想想那個人是怎樣.
打開窗戶想跳下去.
但原來上面一個.
比她早跳.
跳下去.
人們說很不吉利.
弄髒了.
死都被人罵.
我不死了.
神的恩典念文.
臨到她.
她有機會順著主.
認識了神父.
她做手工藝做得很出色.
所以我們有一個社企.
叫做手情家.

$^{1321}$做一些手工藝.
培訓一些婦女.
做手工藝 做導師 教班.
找到自己生命的意義和價值.
將福音傳給她.
這個姐妹.
她叫阿靖.
她信了主後.
生命有很大的改變.
兒女作文.
不再形容媽媽是恐龍.
形容媽媽是袋鼠媽媽.
是不是很抄原來.
袋鼠媽媽.
很保護.
最近她的兒女都來幫忙.
幫媽媽.
在社企裡面.
她不單是做手工藝做得很出色.
有一次我們同工查經.
牧師.
雖然我做手工藝.
做得出色.
但我現在不是做手工藝.
我現在是做什麼.
做人.
她看到.
她所接觸的很多婦女.
很多人.
很大的需要.
現在我們成家有二百多個.
被培訓的婦女.
她們來幫她們.
不單是教她們做手工藝.
透過這個平台.
帶她們來信主.
看見很多生命的改變.
她現在負擔是做人.
雖然神給她的恩賜.
做了一些事.

$^{1361}$她的見證也感動了很多人.
感動了很多人.
在疫情底下.
帶著一班婦女.
進入社區.
關心長者.
派飯給長者.
計劃完結了.
牧師有沒有資源給我.
繼續派.
因為我的長者和我的婦女.
建立了很深厚的關係.
計劃完結了.
就不要就這樣完結了.
你看到她.
對人的那種關心.
那種體諒.
生命的改變.
今天我看見很多.
就像薩瑪利亞婦人一樣.
生命過往是很破碎.
但當耶穌的活水湧流.
到生命當中.
她的生命就不一樣.
可以將整個城的人.
帶到耶穌面前.
讓他們知道耶穌就是真,及全.
我希望.
彭欣彤我們的教會.
都有這五方面的體諒.
超越身份的體諒.
體諒神愛我們每一個人.
不論我們.
什麼背景.
也讓我們體諒.
不只是物質層面.
體諒人心靈裡有很多不同的需要.
不同的貧窮.
唯有耶穌的活水能夠滿足我們.
神的救恩.

$^{1401}$才將我們請教.
讓我們體諒.
教會需要改變.
需要更新.
超越傳統的規範.
對應時代的需要.
來作出伏線.
超越的是什麼?.
我們不是責任的體見.
而是對生命.
對永恆的價值.
我們渴望,追求.
就是生命.
很多失望的生命.
最後我們不只是體諒困難.
是體諒困難背後.
因為全能的神.
他的應許與我們同在.
當我們憑著信的時候.
你能夠經歷.
神的榮耀,神的作為.
在我們當中彰顯.
我們一起祈禱.
天主上,多謝你今天.
讓我們在這個戶外的地方.
來崇拜.
正是讓我們有一種體見.
讓我們敢於超越環境,空間.
傳統的體見.
因為耶穌的體見.
很多失望的人.
他們困苦流離.
像羊毛,牧人一樣.
求你差遣我們.
成為你收莊稼的工人.
讓我們每一個弟兄姊妹.
都能找到我們人生的照明.
當我們願意將自己擺上的時候.
我們就能夠體現.
神的作為.

$^{1441}$神的榮耀就在我們當中.
願神祝福我們的教會.
祝福紅恩堂.
祝福我們每一個弟兄姊妹.
讓我們都能夠帶著基督的愛.
來服侍見證.
願神讓我們的教會.
能夠成為.
洪水橋,洪卜村.
這一帶裡面的一個明亮的燈台.
能夠見證榮耀.
多謝主,誰聽我們的禱告.
共耶穌名下,阿們.
\newpage



\section{腓立比書 1:1-11}
\label{sec:sd0jGt0zjkQ}
\textbf{2024.03.10 《健康的關係 ‧ 喜樂的泉源》 腓立比書 1\_1-11 (和修版) 張美薇牧師}
\newline
\newline
連結: \href{https://youtube.com/watch?v=sd0jGt0zjkQ}{\texttt{https://youtube.com/watch?v=sd0jGt0zjkQ}} ~~~~ 語音日期: 2024-03-12
\newline
\newline
\hyperref[sec:DSsURlo4WPg]{\small{< < < PREV SERMON < < <}}
~
\hyperref[sec:index_chronic]{\small{[返順時目]}}
~
\hyperref[sec:index_scriptual]{\small{[返順卷目]}}
~
\hyperref[sec:wuCk59M9Tmo]{\small{> > > NEXT SERMON > > >}}
\newline
\newline
腓立比書 1:1-11
\newline
\begin{longtable}{cl}
\hline
\hline
章節 & 經文 (和合本修訂版)\\
\hline
1:1 & \begin{tabularx}{0.7\textwidth}{X} 基督耶穌的僕人保羅和提摩太寫信給住腓立比、在基督耶穌裡的眾聖徒，以及諸位監督和執事。 \end{tabularx} \\ \\ \relax
1:2 & \begin{tabularx}{0.7\textwidth}{X} 願恩惠、平安從我們的父神和主耶穌基督歸給你們！ \end{tabularx} \\ \\ \relax
1:3 & \begin{tabularx}{0.7\textwidth}{X} 我每逢想念你們，就感謝我的神， \end{tabularx} \\ \\ \relax
1:4 & \begin{tabularx}{0.7\textwidth}{X} 每逢為你們眾人祈求的時候，總是歡歡喜喜地祈求， \end{tabularx} \\ \\ \relax
1:5 & \begin{tabularx}{0.7\textwidth}{X} 因為從第一天直到如今，你們都同心合意興旺福音。 \end{tabularx} \\ \\ \relax
1:6 & \begin{tabularx}{0.7\textwidth}{X} 我深信，那在你們心裡動了美好工作的，到了耶穌基督的日子必完成這工作。 \end{tabularx} \\ \\ \relax
1:7 & \begin{tabularx}{0.7\textwidth}{X} 我為你們眾人有這樣的想法原是應當的，因為你們常在我心裡；無論我是在捆鎖中，在辯明並證實福音的時候，你們都與我一同蒙恩。 \end{tabularx} \\ \\ \relax
1:8 & \begin{tabularx}{0.7\textwidth}{X} 我以基督耶穌的心腸切切想念你們眾人，這是神可以為我作證的。 \end{tabularx} \\ \\ \relax
1:9 & \begin{tabularx}{0.7\textwidth}{X} 我所禱告的就是：要你們的愛心，在知識和各樣見識上，不斷增長， \end{tabularx} \\ \\ \relax
1:10 & \begin{tabularx}{0.7\textwidth}{X} 使你們能分辨是非，在基督的日子作真誠無可指責的人， \end{tabularx} \\ \\ \relax
1:11 & \begin{tabularx}{0.7\textwidth}{X} 更靠著耶穌基督結滿仁義的果子，歸榮耀稱讚給神。 \end{tabularx} \\ \\
[1ex]
\hline
\hline
\end{longtable}
$^{1}$紅人家的聖徒們平安.
我們一起用祈禱來開始今天的靜悼時間.
我們低頭祈禱.
天主上願你使用孩子.
叫孩子今天所說的.
都是你要孩子說的.
更加幫助會眾.
他們有一個能聽到的耳朵.
有一個能明白你話語的頭腦.
更加有一個心.
將你的話語藏在他們的心裡.
更加保守他們.
有一個能行道的手和腳.
將他們所聽到的.
行在他們的生活當中.
聽我們的祈禱.
奉耶穌基督的名求.
阿們.
各位弟兄姊妹平安.
幾個月沒見了.
通常半年來一次到.
我想大家會覺得.
最近香港人都很抑鬱.
可以說全城都抑鬱.
全城之前抗疫.
但最近都多了.
我發覺不同的人都開始染病.
可能是COVID.
亦都可能是流感等等.
全城都因為各種政治和經濟因素抑鬱.
好像每天聽到的.
都是不好的消息.
其實沒有一個人想抑鬱.
每一個人都在尋求快樂的門.
過去很多人都在追求名譽.
追求地位.
追求學問.
追求錢財.
追求物質等等.
得到之後.

$^{41}$很快人就會明白.
這些都不能夠帶來永久的快樂.
沒有地位的人憂愁.
甚至自殺.
有地位的人都一樣.
沒有學問的人憂愁.
甚至自殺.
有學問的人都一樣.
沒有錢的人憂愁自殺.
有錢的人都一樣.
因為我們知道快樂不在外面.
快樂在我們內裡.
很多時候不單止一般的世人抑鬱.
其實基督徒有沒有抑鬱?.
有!.
好像沒有例外.
當我們要尋找快樂.
在哪裡找呢?.
信徒都知道.
其實我們在尋找.
我們很多時候都不叫快樂.
叫什麼?.
喜樂!.
所以我今天做的題目是.
喜樂.
我們想尋找.
什麼是喜樂呢?.
無論快樂還是喜樂.
並不是從外面找.
其實應該是從內裡尋找.
我們今天嘗試從.
腓立比書的作者.
知道是誰寫的嗎?.
是保羅寫的.
我們在當中.
尋找到我們怎樣.
可以有喜樂的門徑.
很簡單.
其實腓立比書是.
保羅一封很短的書信.

$^{81}$只有四章.
但書中卻充滿了.
喜樂,快樂等等的言詞.
它有喜樂,歡樂,歡喜,喜悅.
這些詞語.
總共出現了17次.
四章書中出現了17次.
但知不知道.
保羅是什麼時候寫的?.
他在羅馬的監獄裡寫的.
其實一封這樣的書信.
這麼充滿喜樂的書信.
作者卻在沒有自由.
受限制.
甚至物質很缺乏.
溫飽也不能擔保.
前路也不明朗的情況下寫的.
這裡更加可以顯示.
其實喜樂真的不在外面.
喜樂在我們內裡.
所以今天想跟大家說.
其實喜樂是一種習慣.
喜樂也是一種選擇.
我們可以選擇.
是你和我的選擇.
這個習慣是要去建立的.
讓我們今天去看看.
這個喜樂的習慣.
就是去建立一個健康的人際關係.
今天的題目是.
健康的關係 喜樂的傳閱.
好像不能下文.
可不可以幫我下文.
我下下一張.
好.
為我們讀出了整篇經文.
第一章第一到第十一節.
我今天不會全部說.
因為時間關係.
我最初預備的時候.

$^{121}$我不知道今天是你們的首勝餐.
所以我要把兩個字退出來.
今天說兩點.
第三點就不說了.
因為我通常都喜歡首時.
當我們要尋找快樂.
我們說過在內裡.
保羅告訴我們.
很重要的一樣.
就是人際關係.
當我們和我們的家人.
當我們和太太或者丈夫.
和我們的子女有爭吵的時候.
你猜你可不可以快樂.
我記得我聽孫先生說的見證.
他那時候和孫太.
很早期的時候.
吵架的時候.
他們來到教會門口.
他們就閉上嘴.
然後大家不再吵架.
但崇拜完又怎樣.
繼續吵架.
繼續板著臉.
因為其實當我們不舒服.
關係不好的時候.
是不可能很快樂的.
我們看看保羅.
在開頭的內容.
怎樣去說他的快樂.
保羅說我每逢想念你們.
當他想起人的時候.
我們看看他寫給什麼人.
我們看第一節.
基督耶穌的僕人保羅和提摩泰.
寫信給泛著菲勒比.
在基督耶穌裡面的眾聖徒.
諸位監督 諸位執事.
他這裡寫給聖徒 監督和執事.
其實對我們一般人來說.

$^{161}$通常不是這樣擺的次序.
你會寫給什麼.
你會寫給監督 執事和信徒.
明不明白.
我們通常會寫給位高的.
這裡最高的位是監督.
第二就是執事.
第三就是信徒.
但保羅不是.
他首先寫給聖徒.
他寫給信徒.
這個好像跟例子不太合.
但我們知道.
其實一個人怎樣可以.
從信徒變成執事.
他要努力.
他要在主裡面繼續成長.
然後到最後他才可以去到監督.
所以保羅是寫給.
全教會的所有信徒.
眾信徒的一個詞.
以後時時會出現.
我們可以說.
他在第一章四節.
就是說為眾人祈求.
他一章七節.
就是關心他們眾人.
然後一章八節.
就是想念他們眾人.
然後一章二十五節.
就是希望他們眾人和他們同住.
然後在第四章二十一節.
也在耶穌基督裡.
去問眾信徒的安.
是眾信徒.
但你有沒有發覺.
我好像有個字.
跟保羅的字不一樣.
他是眾什麼.
聖徒.

$^{201}$你今天是否流行用.
聖徒來形容自己.
好像不是很流行.
好像覺得勝是怎樣.
好像有光環.
是很厲害的.
我現在才介紹.
保羅是會說眾聖徒.
所以你見不見.
我剛才和大家打招呼是什麼.
是紅恩堂的眾聖徒.
怎樣可以成聖呢.
有人說是否教會.
好像天主教會.
捉了聖才成聖呢.
這裡不是這樣說的.
又是不是要被封的才會成聖.
不是.
這裡是說.
我們在耶穌基督裡的聖徒.
所以無論你.
無論我.
我們在耶穌基督裡.
就已經是聖徒.
所以重點是什麼.
我們在基督裡.
在基督裡.
我們在基督裡.
是從神的恩典.
好像今天首聖餐.
我們做一次紀念.
我們為什麼會得到我們現在這個身份.
就因為耶穌基督.
他在十字架上.
為我們眾人的罪而死.
而我們願意向主認罪.
認罪的意思是怎樣.
我們離開我們的罪行.
我們願意在主裡成為一個信徒.
我們就有這個身份.

$^{241}$所以我們教會每個月都會紀念一次.
讓我們不會忘記.
我們知道這是神給我們的恩典.
在基督裡是很重要的.
我們要在基督裡.
就好像我們知道.
小鳥在天空.
如果牠離開了空氣.
牠會怎樣.
沒有空氣會怎樣.
即刻窒息.
魚就在水裡.
魚沒有水又怎樣.
又是即刻死.
植物一樣.
植物雖然根在地裡.
但其實牠一樣是靠著空氣.
靠著根去取得東西.
我們知道.
我們任何一件事.
都不可以離開我們適應的環境而生存.
所以這裡講得很清楚.
一個的信徒要在哪裡.
要在耶穌基督裡.
知道你們最近都有一些主一學.
學怎樣去做信徒.
這是很重要的.
所以很鼓勵大家繼續去學習.
學我們怎樣可以繼續在基督裡成聖.
如果我們時時住在基督裡.
我們就蒙主保守成為聖傑.
但如果我們離開了主又怎樣.
就好像葡萄樹的枝子被斬掉.
丟掉乾了又怎樣.
死了.
枝子沒有用.
葡萄樹枝很幼很硬.
就算用來燒也沒有用.
所以要在主基督裡才可以成為聖徒.
但這裡不只是講基督裡.

$^{281}$因為耶穌基督的僕人保羅和提摩泰.
寫信給什麼.
凡住在菲律賓基督耶穌裡的眾聖徒.
他們除了在基督耶穌裡.
他們還住在菲律賓.
因為我們有兩個身份.
一方面我們要在基督裡.
另一方面我們要怎樣.
我們要活在我們今天的地方.
當時來說.
保羅寫給菲律賓.
他們做菲律賓的聖徒.
不只是他們去崇拜周日.
才去做菲律賓的聖徒.
他們不是每個星期一次.
出席一些團契的活動.
或者在教會其他聚會的時候.
才做聖徒.
他們在菲律賓裡.
他們都是聖徒.
今天我們就知道.
我們不在菲律賓.
菲律賓不是很多人去過.
可能有些人去過.
但也不是很多人去過.
我自己也沒有去過.
但今天我們在哪裡.
我們在洪水橋.
我們在元朗.
我們在香港.
其實當我們發覺.
當我們今天.
當日的菲律賓.
今日的洪水橋.
元朗或者香港.
其實都不是一個沒有罪惡.
好像天堂一樣的地方.
其實四周都有罪惡.
我們最近才看到新聞.
昨天還是前天.

$^{321}$在家裡發現.
兩個嬰兒的屍體.
看到我也害怕.
其實真的不知道怎樣.
有些人說竟然是這樣.
其實四周都是罪惡.
但是保羅在說.
我們活在世上.
但我們卻不和世人同流合污.
我們要過聖潔的生活.
因為我們被召去過聖潔的生活.
我們要被召.
將我們從耶穌女.
我們在基督女.
聖徒的身份.
活在今日的菲律賓.
今日的洪水橋.
今日的元朗.
今日的香港.
我們有了聖徒.
這麼尊榮的地位之後.
我們就要在世上的生活裡.
活出我們的使命.
在地上活出聖潔.
在地上活出美好的見證.
所以保羅就說.
你們從頭一天直到現在.
都是同心合意的興旺福音.
所以興旺福音很重要.
在菲律賓的時候.
就會興旺福音.
他寫給聖徒之後.
就寫給監督和執事.
執事就是教會的一些位份.
監督就是一些管理的工作.
保羅要多謝教會的監督.
執事他們發起.
讓有弟兄姊妹.
或者是爾巴忽提.
送禮物給他們.

$^{361}$同時也直指提醒他們.
要留意保羅在信裡所講及的事情.
請這些監督.
這些執事.
或者這些教會的領袖.
是教導弟兄姊妹的.
因為要不斷地教導.
所以這些信.
當時來說.
就是在崇拜的時候讀出來的.
讓弟兄姊妹同樣都可以得到幫助.
保羅在菲律賓.
是否真的很開心呢?.
大家如果想看的話.
你們可以回到《仕途行傳》第十六章.
第五節至第四十節.
其實是講到保羅.
最初去到菲律賓的時候.
發生了什麼事?.
如果大家有聽過.
就是馬其頓的異象.
他好像在未到達歐洲.
未到達菲律賓之前.
他就有一個異象.
有異象的意思.
就是那裡有人去求他.
去傳福音給他們聽.
那你會期望什麼呢?.
你期望在菲律賓.
就有很多人聽到.
有很多人信主.
但大家如果看《仕途行傳》第十六章.
會看到什麼呢?.
不是的.
保羅去到菲律賓.
在城裡住了幾天.
他找不到猶太人的會堂.
只是說他找不到聚集的猶太人.
或者根本找不到機會.
和猶太人的男人去溝通.

$^{401}$那時候男女是分得很清楚的.
男人是一家之主.
不是特別有社會上的地位.
他找不到.
結果他只在河邊.
和幾個女的.
一直在分享.
當中一個是女底亞.
後來他在街上看到一個侍女.
那個侍女有鬼在她身上.
結果保羅不耐煩.
因為侍女跟著保羅.
於是保羅不耐煩.
就趕了那隻鬼出去.
結果怎樣呢?.
結果侍女的主人竟然抓了他.
然後關在監獄裡.
半夜還打到他受傷.
然後就被關在監獄裡.
嘩!開心嗎?.
當然不開心.
被抓了.
然後被打.
所以如果保羅停留在菲律賓.
不開心的事上.
保羅可以很沮喪.
但為什麼他可以在腓立比書裡.
寫出這麼開心的事呢?.
因為他想起他們的好處.
他說從頭一天直到如今.
你們是同心合意的.
慶旺福音.
即是第五節裡說.
他可能想起女底亞在河邊.
信主之後.
他就開放家裡讓保羅回家住.
然後菲律賓教會.
就在女底亞的家裡成立.
他為菲律賓教會有一個.
很愛主.

$^{441}$願意擺上的姊妹.
雖然他被關在監獄裡.
結果怎樣呢?.
禁足全家信主.
很難得.
他也想起菲律賓教會.
一次兩次的打發人.
給他的需用.
所以保羅就在第三節裡說.
願因為平安.
從神我們的父並主耶穌基督歸於你們.
然後我每逢想念你們.
所以想和大家分享的第一點.
建立健康的人際關係.
是怎樣呢?.
想起別人的時候.
不要想壞事.
不要想不開心的事.
想什麼呢?.
因為怎樣呢?.
剛剛說了.
第五節說.
從頭一天直到如今.
你們是同心合意的興旺福音.
所以保羅選擇.
剛才說過快樂是選擇.
喜樂是選擇.
保羅選擇紀念別人的好處.
別人討神喜悅的事.
別人的好處的時候.
今天其實我們和別人的關係都一樣.
比如說.
你和你的配偶.
有沒有好的關係呢?.
大家在別人面前一定說.
一定好.
其實配偶和你做了很多好事.
比如說.
我星期四晚上教書.
我穿不夠.

$^{481}$我帶不夠衣服.
因為我一感受回去.
晚上就要回長洲.
很冷.
我就WhatsApp曾牧師.
我今天不夠衣服.
我沒有頸巾和帽子.
曾牧師又怎樣呢?.
他又來接船.
我坐10點船.
他10點7就到碼頭.
拿了帽子和頸巾給我.
是不是做得好?.
好.
但我又不好.
後來到他講道的時候.
才知道原來他也介意這些事.
因為我們覺得做得好應該是應份的.
做得好應份的.
但是我們將焦點放在哪裡?.
他做不好的事.
他有什麼做得不好.
我們就說.
哎呀 為什麼他會這樣?.
我們時時紀念雙方.
相處不好的事.
我們知道.
其實有些事.
不好的關係是沒有問題的.
譬如你在街上.
見到一個你從來沒見過的人.
你可能不會再見到他.
他走路的時候撞了你一大下.
你就很痛.
你看看他.
誰知道他怎樣.
他掌摑著你.
你想想.
我當然要罵他.
但你想想.

$^{521}$算了.
我以後都未必再見到他.
那又怎樣?.
一聲把氣吞下去.
那就算了.
因為那個人你不認識.
你都未必再有機會再見到他.
但是如果那個人是你的兒子.
那又怎樣?.
為什麼你走路不打眼睛?.
為什麼你踢到我?.
為什麼?.
很多時候就是這樣.
不理想的關係.
和熟悉的人.
是怎樣?.
是揮之不去的.
越熟就會越生氣.
越熟就會越不開心.
好像和你的配偶.
你的兒女.
你的父母長輩.
你的同事.
你的鄰居等等.
我們要保持關係是一生的.
是要一生一世.
我們不單只要保持.
我們更加希望關係可以改善.
我們知道兩個人的關係.
都要改善.
我有事要改.
對方呢?.
對方也有事要改.
我們很多時候都希望對方知道要怎樣做.
他是女兒.
他應該怎樣?.
他是女兒.
他應該怎樣?.
他是我先生.
他應該怎樣?.

$^{561}$他是我鄰居.
他應該怎樣?.
我們都知道.
其實你能改變別人嗎?.
你改不了別人.
改變別人其實是神的工作.
但是我們怎樣?.
我們可以改變自己.
所以保羅說什麼?.
我每逢想念你們.
所以保羅就是改變自己.
我們要好像保羅一樣.
管理我們的記憶.
選擇去紀念別人的好處.
忘掉一些不好的東西.
打個比喻.
假如你心目中有敵人.
你心目中有沒有敵人?.
通常都有.
我們當然說沒有.
但很多時候都有.
假如你心目中有敵人.
你將他困在你心中的監倉.
你想著想著.
想著想著.
不太對勁.
於是將他拿出來.
打打打打打.
打了他一輪.
懲罰了他.
然後將他放回監倉.
你猜他會不會痛?.
他會不會知道?.
他不會.
但怎樣?.
你就不是了.
你的情緒就會受到波動.
會不開心.
所以你去以為懲罰了別人.
其實是怎樣?.

$^{601}$其實是懲罰了自己.
所以管理我們的記憶.
是將對方的事.
所以這裡會說.
我們建立健康的人際關係.
第一個法門是怎樣?.
紀念別人的好處.
其他的都忘掉.
其實這個是可以學的.
有沒有學過正向思想?.
這就是正向思想.
正向思想.
我們只想正面的事.
正面的事是很好的.
這是第一樣.
是正面的思想.
但只有正向的思想.
夠不夠呢?.
是不夠的.
我們是什麼?.
我們是信徒.
我們還有第二個利器.
第二個利器是什麼呢?.
第二個利器.
就是為別人歡喜的祈求.
就是正向的禱告.
歡歡喜喜的為別人祈求.
保羅要為跟他有關係的人.
我們現在不是說街上不認識的人.
是跟我們有關係.
但可能關係不是很好.
你心目中通常都會有不同的人.
但最好不要對號入座.
要為他歡喜的祈求.
你有沒有想過.
有人為你祈禱.
你會怎樣?.
是很開心的.
我和大家分享一些我的經歷.
我1998年從美國回來香港.

$^{641}$我回來香港的時候.
開始在錦繡堂募會.
那時候我有些過敏.
我還沒有回來募會之前.
每次放假回來.
通常回來兩三個星期.
我都會聲沙啞才回到美國.
但我過來之後.
我發覺我一年都沙啞.
沙啞的情況.
到了晚上的時候.
有些音是發不出來的.
我不知道你們是否明白.
有些音是發不出來.
但早上更容易發.
但晚上就發不出來.
如果要說到就很擔心.
但當我發覺.
那時候教會有一個.
白天的祈禱會.
有幾個老人家.
每次都參與祈禱會.
星期二,三,四.
就會一起祈禱.
我和謝牧師.
那時候我們兩個都是.
謝姑娘和張姑娘.
我們兩個就在這裡.
開一個祈禱會輪流帶.
有一個是簡老太.
簡明光.
你媽媽在這裡.
簡老太每次祈禱.
就口氣數全道人.
宣教士.
你.
你太太.
還有你的親戚朋友.
不斷地數.
是很暢順的.

$^{681}$你會發覺他們.
時時祈禱都暢順.
家人不….
我們會明白.
但原來教會的全道人.
和宣教士他們都是.
很難得.
時時為他們祈禱.
另外一個就是朱太的爸爸陳伯.
陳伯也是時時為我們祈禱.
他經常說.
他說張姑娘.
不用怕.
我為你祈禱.
你講道的時候.
我就為你祈禱.
真的.
我每次講道.
我都講到出來.
我不會突然講到.
沒有了那個音.
我不知道大家是否明白.
如果你試過聲沙.
就會發覺.
去到講到某個位置.
是沒有了那個音.
但我一直靠他們.
為我們祈禱.
然後就….
有人為你祈禱.
是很開心的.
你知不知道陳伯.
是怎樣離世的.
陳伯是在一個.
祈禱會裡離世的.
那時候教會是有.
福音主日.
星期六晚的時候.
朱太是老師.
校長.

$^{721}$於是他請了一班同事.
在他家燒烤.
吃完之後.
就帶回教會.
然後就聽福音主日.
陳伯就在.
那個輔堂那裡.
幫忙祈禱.
他在祈禱的時候暈倒.
然後第二天離世.
我當時是很難過.
但想想.
其實可以在服事當中離世.
都是一個很難得的事.
所以我們發覺.
為教會禱告.
為傳道禱告.
為執事和領袖們禱告.
為幹事禱告.
為教會其他弟兄姊妹禱告.
為宣教士禱告.
為福音能夠在世界廣傳禱告.
為福音能夠在中國.
傳開歡喜的祈求.
是很難得的.
所以每逢為你們中人祈求的時候.
時時都是歡歡喜喜的祈求.
所以當有人得罪你的時候.
你和他的理論有幫助嗎?.
沒有幫助.
你向別人投訴.
期望別人能夠幫助你解決關係.
有沒有幫助?.
都是沒有幫助.
但是你卻可以為他祈禱.
你歡歡喜喜的為他祈禱.
我們不能夠改變別人.
但是我們可以為他祈禱.
讓神出手.
神要改變他就改變他.

$^{761}$神要改變你就改變你.
神要改變環境.
就會改變環境.
所以第二個竅門是什麼?.
開始為別人禱告.
是最快能夠改變一個不好的關係.
這是正向的祈禱.
有正向的思想是不足夠的.
要有正向的祈禱.
就好像在第四節.
每逢為你們中人祈求的時候.
時時都是歡歡喜喜的祈求.
所以保羅不會想到.
他在菲律賓找不到男士.
不會想到他在菲律賓被打.
反而會想到.
教會在女帝雅家開展.
為他們祈禱.
因為福音從第一天開始.
就在他們那裡出來.
那如何為他們祈禱呢?.
很快地.
我用五分鐘跟大家說完.
我就結束了.
為別人祈禱的方向.
我們看第九到十一節.
我屬禱告的.
就是要你們的愛心.
在知識和各樣見識上.
多而又多.
所以你們能分別是非.
將誠實無過的人.
直到基督的日子.
並靠著耶穌基督.
結滿了人意的果子.
叫榮耀稱讚歸給神.
所以我們要為別人祈禱.
如何祈禱呢?.
大家可以回去慢慢查經.
查這三節.

$^{801}$第一.
讓他們愛心增長.
你為他們祈禱.
他們可能得罪了你.
跟你吵架.
當然你要為自己愛心增長祈禱.
但還要為對方的愛心增長祈禱.
更重要的.
不只是愛心增長.
他們還要有知識和見識.
不是男愛或有愛無慮.
人家做什麼都愛.
不是.
是要有知識.
希望他們在知識和見識上.
愛心都能夠增長.
另外.
是能夠分別是非.
當然你會以為別人錯.
當你跟別人吵架.
最好就是求神幫助.
讓他們能夠分別是非.
我當然覺得自己對他們錯.
我都求神幫助.
幫助他們知道自己錯.
但其實我祈禱的時候會發現.
原來我錯他們是對的.
當我為別人祈禱的時候.
我都為自己祈禱.
知識和見識.
常常多而又多.
使我們能夠分別是非.
另外.
做誠實無過的人.
要很真誠.
要很誠實.
不會有過錯.
或過失.
我們要誠實.
但也要在對的時候.

$^{841}$說合理的話.
例如.
當你發現一個人.
他的手上有一個很大的瘤.
長出來.
你當然說.
你去看醫生吧.
但你不會說.
讓我來幫你去割.
如果你是金峰醫生.
當然可以.
但如果你不是.
你幫人去割.
會亂了.
所以你要誠實.
要知道自己做到什麼.
做不到什麼.
然後按著需要去做好.
所以這是很重要.
然後就期望.
他靠著基督.
去結滿人意的過節.
為別人祈禱.
就是為這四樣東西祈禱.
因為在保羅的《斐納比書》.
第一章第九到十一節.
就說了出來.
我還有第三個.
以基督的心牆愛別人.
也一樣有經文.
但最重要的一句.
抱著基督的心牆去為別人.
第二章是為人代禱.
剛才就說了.
保羅在這裡也說了.
那四個代禱的方向.
然後第六到八節.
我們就看到.
原來我們要用基督的心牆.
去愛別人.

$^{881}$心牆是什麼呢.
其實是那兩點.
所以我通常說到結束.
我都會這樣結束.
於我何干.
大家聽了半個小時.
我嘗試了.
七個字.
不好意思 遲了.
大家聽了之後.
和你有什麼關係呢.
想想有什麼人.
是你很親的.
關係理論上很好.
但好像有點不理想.
你能不能想起他的時候.
能不能為他祈禱.
成為你們之間的關係.
讓你們的關係能夠重新建立.
因為健康的人際關係.
是我們喜樂的泉源.
剛才說過喜樂是怎樣.
是一個習慣.
亦是一個選擇.
我們要選擇.
你今天會否做這個選擇呢.
讓我們一起低頭祈禱吧.
天主上願你使用孩子今天所說的.
願會眾對你的話語有所回應.
聽我們的祈禱.
奉耶穌基督的名求.
阿們.
\newpage



\section{路加福音 13:22-30}
\label{sec:wuCk59M9Tmo}
\textbf{2024.03.17 《當努力進窄門》 路加福音13\_22-30 (和修版) 講員 : 曾敏姑娘}
\newline
\newline
連結: \href{https://youtube.com/watch?v=wuCk59M9Tmo}{\texttt{https://youtube.com/watch?v=wuCk59M9Tmo}} ~~~~ 語音日期: 2024-03-18
\newline
\newline
\hyperref[sec:sd0jGt0zjkQ]{\small{< < < PREV SERMON < < <}}
~
\hyperref[sec:index_chronic]{\small{[返順時目]}}
~
\hyperref[sec:index_scriptual]{\small{[返順卷目]}}
~
\hyperref[sec:hGBqBs0DEW4]{\small{> > > NEXT SERMON > > >}}
\newline
\newline
路加福音 13:22-30
\newline
\begin{longtable}{cl}
\hline
\hline
章節 & 經文 (和合本修訂版)\\
\hline
13:22 & \begin{tabularx}{0.7\textwidth}{X} 耶穌往耶路撒冷去，在所經過的各城各鄉教導人。 \end{tabularx} \\ \\ \relax
13:23 & \begin{tabularx}{0.7\textwidth}{X} 有一個人問他：「主啊，得救的人很少吧？」 耶穌對眾人說： \end{tabularx} \\ \\ \relax
13:24 & \begin{tabularx}{0.7\textwidth}{X} 「你們要努力進窄門。我告訴你們，將來有許多人想要進去，卻不能。 \end{tabularx} \\ \\ \relax
13:25 & \begin{tabularx}{0.7\textwidth}{X} 等到一家之主起來關了門，你們才站在外面敲門，說：『主啊，給我們開門！』他要回答你們說：『我不認識你們，不知道你們是哪裡來的。』 \end{tabularx} \\ \\ \relax
13:26 & \begin{tabularx}{0.7\textwidth}{X} 那時，你們要說：『我們在你面前吃過喝過，你也在我們的街上教導過人。』 \end{tabularx} \\ \\ \relax
13:27 & \begin{tabularx}{0.7\textwidth}{X} 他要對你們說：『我告訴你們，我不知道你們是哪裡來的。你們這一切不義的人，給我走開！』 \end{tabularx} \\ \\ \relax
13:28 & \begin{tabularx}{0.7\textwidth}{X} 你們要看見亞伯拉罕、以撒、雅各和眾先知都在神的國裡，你們卻被趕到外面，在那裡要哀哭切齒了。 \end{tabularx} \\ \\ \relax
13:29 & \begin{tabularx}{0.7\textwidth}{X} 從東從西，從南從北，將有人來，在神的國裡坐席。 \end{tabularx} \\ \\ \relax
13:30 & \begin{tabularx}{0.7\textwidth}{X} 看吧，在後的，將要在前；在前的，將要在後。」 \end{tabularx} \\ \\
[1ex]
\hline
\hline
\end{longtable}
$^{1}$弟兄姊妹平安.
每一次能夠回到神的家裡敬拜神.
都是我們的福氣.
神仍然將祂的話語賜給我們.
也是我們的福氣.
我們開始的時候再一同禱告.
讓我們又再可以聆聽你的說話.
今天求你讓我們眾人有一個開放的心.
開放的耳朵.
讓我們聽 明白 又去回應.
心願聖靈親自對我們眾人說話.
奉耶穌基督的名字祈禱 阿們.
很快的三月尾我們有一個大假期.
是什麼假期呢?.
是呀 受苦節 復活節.
那個假期大家有沒有計算過?.
有沒有計算過?.
上班的人可能會計算一下.
只要請幾天假.
請幾天假將會放一個非常長的假期.
哇 真的很舒服.
整個假期休息 去旅行.
很多人就在計劃.
去哪裡旅行 去玩什麼 吃大餐.
有很多禮物.
也可以去探訪很多人.
做很多自己很想做的事情.
就是歡喜快樂.
就像這幅圖一樣.
復活節是不是令大家想起長假期?.
開心 是不是這樣一回事呢?.
耶穌基督希望我們怎樣去過這個假期.
受苦節 復活節.
耶穌對我們的心意究竟是什麼?.
這個反而值得我們去問一問.
今天當我聽福音樂曲隊去唱福曲的時候.
心裡覺得非常配合.
其實我們是沒有配合過的.
但原來大家很配合.
那兩首福曲都表達到主的心意.

$^{41}$這個長假期究竟是為了什麼?.
主想我們怎樣呢?.
剛才那兩首福曲令我們想起什麼?.
主的受苦 主的犧牲.
祂是釘死在十字架上.
也想起在第三天主復活.
原來受苦節 復活節主最希望我們想起的.
是祂的受苦 祂的受死 祂的犧牲 祂的復活.
為什麼呢?.
因為主的受苦 被釘十字架 犧牲 祂的復活.
正正就是為我們成就了舊恩.
為什麼我們今天可以坐在紅印堂敬拜我們的神?.
為什麼我們的生命有一個全新的開始?.
我們不需要再懼怕永遠的死亡.
我們人生這麼有盼望.
其實信耶穌的人是最應該去慶祝復活節的.
但不是世人那種慶祝的方式.
是那種對主的敬拜讚美.
為什麼呢?.
我們得到一份很寶貴的禮物.
是主拯救我們的生命 離開永遠的死亡.
讓我們得著永生 是永遠的和祂在一起.
是何等的美好.
原來主對我們的心意是要這樣紀念祂的受苦 祂的復活.
我們的回應是什麼呢?.
我們對主的舊恩回應是什麼呢?.
我們的回應決定了我們的將來究竟是永生還是永死.
你說我已經坐在這裡了.
坐在這裡還說這個題目.
這些像是對未信主的人說的 是不是?.
我現在應該立即走出去找個地方對未信主的朋友說.
這個不應該對我們說的.
一定是永生的 是不是?.
是不是呢?.
為什麼我打個問號呢?.
我們的回應是怎麼樣呢?.
今天要去提一提.
主耶穌在這裡說了一件很重要的事.
是祂說的 祂說你們要努力進閘門.
閘門是指什麼呢?.

$^{81}$在這裡耶穌用一個很好的比喻去形容.
閘門是指通往永生.
一個很重要的一件事.
透過閘門我們進入永生.
耶穌說要努力進閘門.
是祂鼓勵我們 是祂勸勉我們 吩咐我們 甚至是命令我們.
要努力進閘門.
在這裡有三個問號 為什麼呢?.
為什麼?我要問 為什麼要努力呢?.
要很努力去得永生?.
這個觀念好像很奇怪.
我們記得約翰福音十四章六節是這樣說的.
耶穌說 我就是道路真理生命.
若不直著我 沒有人能到乎哪裡去.
什麼叫我要努力進閘門.
是不是現在說的要靠自己的努力 自己的行為.
我要有好行為 我要好像世俗所說的 直得行善.
那我就可以得到永生.
是不是說這些呢?絕對不是.
聖經說得很清楚 耶穌說祂是道路真理生命.
要直著祂才能去到乎哪裡去.
所以一定要靠著耶穌基督 直著耶穌基督.
直著祂所成就的救恩 我們才能得到永生.
這是很肯定的.
直著我們那份信 耶穌基督所成就的救恩.
天父上帝的恩典 我們才能得到永生.
這是肯定的 不容置疑.
但是耶穌又叫我們要努力進閘門.
是甚麼一回事? 即是要努力甚麼?.
所以我們要問 甚麼是努力?.
努力是指甚麼? 有兩方面.
第一方面是我們對福音要有真實的 熱切的回應.
沒錯 我們是要靠著救恩 靠著耶穌基督.
靠著那份信 我們得到永生.
但是主期望我們聽了福音之後.
我們明白救恩 需要有一個真正的回應 是真實的.
不是吹出來的 不是假的.
是虛浮的 是熱切的回應 甚麼是努力?.
是指我們的回應 是主耶穌很想要有的.
再說得白一點 是甚麼? 生命要悔改.

$^{121}$我聽了福音 我聽了救恩 主為我犧牲 主為我成就救恩.
我邀請耶穌基督進入到我的生命裡.
不是這樣畫上一個句號 而是生命要改變 悔改.
以前我是個怎樣的人?.
那些不好的事 我所犯的罪 今天我不要再犯了.
我要離開 以前說謊話 現在就不要說謊了.
以前我是貪心的 現在不要再貪心了.
以前我推卸責任的 我經常都A字簿的.
現在不要A字簿了 我要去承擔 做出來.
別人是看到你的生命的 是不同的 你是另一個人.
是按主的心意而行 以前我是按自己心意而行的.
我話勢是 我是我生命的主人 信主之後不是 是按主的心意而行.
而這一點我發覺對於很多基督徒來說是很難做的.
在甚麼時候看到呢? 很多時候做人生決定的時候.
很快就會自己拿主意了 我讀哪間學校?.
我做甚麼職業? 選甚麼男朋友? 女朋友?.
我人生怎樣決定? 移不移民?.
我自己怎樣? 我的子女怎樣? 我父母怎樣?.
全部都是自己決定的.
我做好決定之後 我就通知上帝一聲.
上帝 我做好決定了 你覺得怎樣?.
這樣是否按主心意而行呢?.
有多少人願意付上禱告的代價 問神.
神你覺得接下來是怎樣?.
在神那裡得著答案 按神的帶領去行呢?.
當我邀請耶穌進入我的生命那一刻開始.
我的生命至少有這幾樣東西 是不同了.
我會遵守主導 聖經裡怎樣教? 我會怎樣去做?.
我看事情的觀點角度是不一樣的.
我未信主前去看 是怎樣去看一件事?.
信主之後去看 是怎樣去看?.
看開心的事 看不開心的事 特別是不開心的事情.
信主之後 我們的觀點角度是應該有不同的.
我不會和以前未信一樣 都是這樣去看事情.
你去看 你會不會問神?.
神啊 讓我明白這件事 我應該怎樣去看?.
有沒有我要學的功課? 有沒有我要明白的道理?.
有沒有一些事情你想我回應?.
或者你希望別人回應 我可以為別人禱告.
信主的人不一樣.

$^{161}$所以要努力進去的時候 這一切都是有不同的.
生命要改變.
第二件事 努力進閘門 要努力是指什麼?.
我們要為主受苦 付代價.
這件事很多人都覺得不中聽.
我比較反對傳福音時 經常說一個很廉價的救恩.
經常說到一個樂園天地的信仰.
你來信耶穌 信耶穌有很多福氣 很多好處.
肉身的福氣 拿著護照上天堂 真是好 很多無比.
讓別人聽起來感覺上 信耶穌之後身體不會病.
事業順利 讀書一定好 我人又漂亮.
戀愛婚姻順利 事業也好 什麼都好.
我不會有苦難 信耶穌總之每樣都好.
信耶穌之後感覺上被人騙了 騙了他.
我好多苦難 身邊很多人不開心的事.
我要和上帝理論 上帝是你欠我的.
上帝從來都沒有欠我們.
我們不應該有廉價的信仰.
傳福音也不應該有廉價的福音.
是很重價的 拿耶穌基督的生命換回來的.
要為主受苦 信耶穌那一刻就開始了.
不只是傳道人宣教時才會的.
是每一個跟從主的人就要有心理準備.
信主那一刻就開始了 所以不簡單的.
努力進閘門 努力的意思是很重要的.
為主付代價 認不認識他.
認識他嗎?看上去他像什麼職業?.
明星?他很有名氣.
報導這個新聞的時候 我還在想他純粹是一個日本人.
在那個事件中被人斬首 我不知道這麼多.
後來上帝提醒我 要留意這個人.
原來他是一個基督徒 而且是一個非常虔誠的基督徒.
當然他不是宣教士 也不是傳道人.
就是一個很虔誠的基督徒.
後藤建義是1997年受洗成為基督徒.
他當時加入日本基督教團和東京都大田區的田園調報教區.
去回教會 虔誠.
在網上你可以看到有關他的文章.
他強調的不只是基督徒的身份 而是他的生命.
他活出他的信仰 活出他所相信的福音.

$^{201}$後藤建義是特別選擇了一些題材去報導的.
作為記者 他是一個戰地記者.
所以經常出入一些很危險的地方.
他報導過伊拉克 索馬利亞和敘利亞等等戰亂衝突地區的一些新聞.
特別他的焦點是去報導受苦的群眾.
例如兒童 婦女 貧窮的百姓等等.
目的不是為了出圍 目的是要將他們的苦況拍攝出來.
讓世人留意他們 為他們的苦難提供一個出口.
讓世人留意 有不同的人為他們禱告 關注他們.
他們願意透過自己的職業去為人伸張正義的那種人.
其實常常處於很危險的狀況.
你說他有沒有害怕過呢? 一定有.
他也是一個人 換作我們處於戰況當中 我們也會害怕.
但他心裡很清楚 如果有一天離開了這個世界的時候.
他不害怕 為什麼呢? 天父會為他有最好的安排.
天父會帶他離開 這一份相信成為了他很大的幫助.
每一次他出去採訪 他都會帶著聖經去.
為什麼呢? 讓上帝的話成為了他的帶領.
所以不只是他很忙 衝出去做很多事情.
是有上帝的話成為了他的幫助.
在過程中 他的心常常與上帝結連.
上帝在聖經裡所教導的 他也會活出來.
所以在網上給他一個很高的評價.
後來他的好友被伊斯蘭國脅持 被綁架.
為了去營救他的好友 他不顧自己的安危.
當時是他的第二個女兒 出生不久.
但他要放下家人去救他的好朋友.
結果他在拯救的過程中 自己也被伊斯蘭國綁架 脅持.
經過了好一些日子的協商.
日本政府也沒辦法與伊斯蘭國達成協議.
結果他被斬頭.
當時日本人很震驚.
竟然這樣傷害我們的記者.
家人也很難過 但家人很有信心.
他的家人也是信耶穌 以他為自豪.
後人一直說起的時候 想起的都是他對上帝堅定的信.
一生去活出他的信仰.
你說這樣就很不值.
他的人生有一個很美好的前程.
他在做他很喜歡的記者職業.

$^{241}$他有美好的家庭 很愛他的太太 有子女.
這麼幸福 你就安安定定.
其實做記者有很多題材可以採訪.
一樣可以做得很出色.
不用特意去戰區 每天都冒著生命危險.
你不需要 還有你的好友發生什麼事 你也不需要去營救.
永遠都想著怎樣救自己的生命.
他就沒有這樣的選擇.
他活出他一生所相信的生命.
是甘願受苦.
如果大家看聖經裡其中一個說受苦很突出的是什麼人物.
再大聲一點 約伯.
這是舊約裡的一個人物.
他經歷了很多苦難.
新約裡的保羅 大家有機會回去聖經看.
他受了很多苦.
上上都在危難裡 上上都在死亡的邊緣.
為了服侍主 他甘心.
其實主呼召他那一天 在大馬士革跟保羅相遇的那一天.
主不單去拯救保羅的生命.
其實也是給他一個很重要的使命.
就是要向外邦人傳福音.
主已經定義保羅的一生要為主受很多苦.
結果保羅真的受了很多苦.
你說那是保羅 這是後藤健兒 不是我.
對不起 主也呼召我們.
為什麼?.
主耶穌自己為了拯救我和你的生命 祂就受了很多苦.
昨晚我很感恩.
當我參與青少年的工作的時候.
和江澤夫夫一起 有機會幫一班DSE的考生.
和他們分享考試的事情 前路等等.
在裡面有沒有想到上帝很奇妙.
為我們開路可以和其中兩位學生去傳福音.
一路去傳福音的時候.
再一次有機會去說耶穌基督為我們受的苦.
祂被釘在十字架的苦.
我問那位學生.
你知不知道人們怎樣形容.
當釘釘到耶穌身體的時候 痛苦是多大.

$^{281}$有沒有人聽過這個形容.
當釘釘到祂身體的時候 痛苦是多大.
有沒有人可以去描述.
我曾經….
十級痛.
我聽說是生孩子.
生孩子是最痛 是不是這樣.
是很痛的.
龐靈珠說是很痛 很痛.
是痛入心.
相信生孩子也會痛入心.
但有人形容過那種痛是好像在你身體裡放了一百斤炸藥.
一聲炸 那種痛是很厲害的.
耶穌基督為我們被釘在十字架上.
當年這是一個很殘酷的死刑.
耶穌並沒有犯任何罪.
祂為我和你付上了代價.
很多時候我們說的時候 變成了好像吃生菜一樣.
很隨便 很隨意 已經淡忘.
也不是很深入 很能夠明白.
什麼叫為我和你受苦呢.
是為我和你付上生命代價.
多痛啊 是不是.
祂為了救我和你 受了很多很多的苦.
祂已經做了受苦的先鋒.
叫一切相信祂的人 踏上這條路的時候都會受苦.
如果我們立定志向要敬虔道日 聖經有說過.
如果我立定志向要敬虔道日 會怎樣呢.
一定會多受苦難.
除非我今天決定 就赤過這一生.
但是你這樣就枉然了.
你沒有真真正正的跟從主.
沒有真正的邀請祂進入你的生命.
也不明白祂為我們犧牲 拯救是什麼一回事.
主已經為我們犧牲 跟從祂的人也會走這條路去受苦.
事實上有很多見證人如同雲彩一樣在我們身邊包圍著我們.
他們做了一個很好的榜樣.
他們沒有吹噓自己的生命有多好 只是默默地做.
在今天 耶穌要我們把握機會進閘門.
把握機會 因為時間有限.

$^{321}$經文裡這樣說 等到一家之主起來關門的時候.
你才會突然間想起 啊 我想進門了.
我想進入永生裡 我這個時候要開門了.
主啊 讓我們開門吧 這個時候已經來不及了.
祂要回答你們說 我不認識你們 我不知道你們從哪裡來的.
時間非常有限.
你說 行了 我不知道主什麼時候再回來.
可能很長很長很久才回來 我有的是時間.
但對不起 時間從來都不是這樣掌握在我們手裡.
我們知道我們人生什麼時候是最後一天.
現在就算在禮堂發生什麼事我也不知道.
出去教會門外發生什麼事我也不知道.
明天如何我也不知道 是不是.
時間有限 要把握機會 要進閘門.
否則 當這個時間已經到了.
主再回來的時候 我們就沒有這個機會.
這是耶穌說的 一家之主 是不是.
這就是掌管著永生的一家之主.
我們的神 我們的耶穌基督.
會有關門的一天.
不要覺得那道門永遠為我打開 永遠等著我.
會有關門的一天 把握時間.
可能有些人說 不要緊 我和主有些關係.
我和主吃喝過 他認得我.
到時一定開門給我 是不是.
和主吃喝過 有些關係 他記得的.
吃喝 通常是接納的人才和他一起吃.
就是飲飲食食 不認識不認識怎會一起飲食.
主認識我 記得的.
還有 我聽過主教導.
我聽過很多堂導.
是不是聽過主教導就可以得救呢?.
夠了沒有?.
還未夠 經文說還未夠.
主耶穌要的是我們和他有真實的關係.
真正的相信.
我聽了他的導 我就真正相信 我接受.
我會去聽從 我會走出來.
我會跟著他 真真正正擁抱耶穌.
聽也可以 不入耳.

$^{361}$所以 由耳出 聽完好像沒聽過.
我最怕下星期再回來問你.
還記不記得上星期說過什麼.
我已經不記得了.
我怕我自己也是這樣.
所以要彼此提醒.
真的要活出來 要有真實的關係.
你真的要接受.
有些人聽完聽了很多.
他也可以站上台說.
可能他說得比我們好.
但不代表他相信 不代表他會聽從.
這是真的.
我曾經說過一件事 令人很害怕.
我跟別人說 你相不相信.
可能這個世界有些人 最好的講員.
有機會不相信耶穌.
有人立即走進去.
最好的講員不相信耶穌.
他真的很懂得怎麼說.
可能會震動你的心靈 但他不相信.
這樣沒意思.
他也沒有進閘門.
他也不會得著永生.
只是說而已.
有很多人口才很厲害.
懂得說話的方式.
不代表是生命 不代表是關係.
主要是關係 真實的關係.
真正的相信.
將來在上帝的國度裡有很多人.
來自世界各地的人.
那些是真真正正相信耶穌基督的人.
聽了主的話.
會真的去擁抱.
會去聽從 會去信服 跟從.
走出來的人.
有世界各地很多的人.
就會在那裡相聚.
一同敬拜 讚美我們的神.

$^{401}$心願我和你進閘門.
心願我和你進閘門 得到永生.
在上帝的國裡是一份子.
在這裡我和你是一份子.
今天問自己兩個問題.
第一 你聽了這麼多.
你是否真正相信.
剛才說要努力進閘門.
這麼久以來你很熟悉舊音.
很清楚福音是說什麼.
你真的回應了.
你真的回應了 你真的明白.
第二 主問我們要不要付代價受苦.
你的回應是什麼.
剛才那兩首歌曲是很寶貴.
很詳盡說的是舊音.
主已經付代價了.
今天跟著他走的人也會這樣.
我想問大家願意為主付什麼代價.
我想在這個受苦節 復活節.
請大家思考這個問題.
不要只說我為你付代價.
一邊要用時間 主要是用錢的時候.
主要是用我的恩賜的時候.
主啊 你不要過來.
還有很多東西 要為主付代價.
我也不願意 是不是這樣呢.
我們一起思考.
藉著你的話語提醒我們.
主耶穌說要努力進閘門.
這是對我們的勸勉.
對我們的提醒 是要努力.
我們不能只在嘴巴上說願意 相信.
要把整個生命放在上帝的手裡.
真實的回應你.
在主你的教導上 我們切實地擁抱 跟從.
在受苦的路上 我們有份.
主啊 我們這樣去禱告其實並不容易.
求主幫助我們 這麼多人去思考.
究竟我們願意為主你擺上什麼.

$^{441}$受苦節 復活節 對我們來說.
不只是一個長的假期.
而是我們要回應主你的日子.
所以求你深深地感動我們 幫助我們.
心願我們每個人都活出保羅的生命.
侯騰建義的生命.
更活出耶穌基督的生命.
我們仰望你 奉耶穌基督的名字祈禱.
阿們.
\newpage



\section{羅馬書 8:31-34}
\label{sec:hGBqBs0DEW4}
\textbf{2024.03.31 《一切從復活開始》 羅馬書 8\_31-34 (和修版) 曾裔貴牧師}
\newline
\newline
連結: \href{https://youtube.com/watch?v=hGBqBs0DEW4}{\texttt{https://youtube.com/watch?v=hGBqBs0DEW4}} ~~~~ 語音日期: 2024-04-01
\newline
\newline
\hyperref[sec:wuCk59M9Tmo]{\small{< < < PREV SERMON < < <}}
~
\hyperref[sec:index_chronic]{\small{[返順時目]}}
~
\hyperref[sec:index_scriptual]{\small{[返順卷目]}}
~
\hyperref[sec:Ux8UdMSUx6c]{\small{> > > NEXT SERMON > > >}}
\newline
\newline
羅馬書 8:31-34
\newline
\begin{longtable}{cl}
\hline
\hline
章節 & 經文 (和合本修訂版)\\
\hline
8:31 & \begin{tabularx}{0.7\textwidth}{X} 既是這樣，我們對這些事還要怎麼說呢？神若幫助我們，誰能抵擋我們呢？ \end{tabularx} \\ \\ \relax
8:32 & \begin{tabularx}{0.7\textwidth}{X} 神既不顧惜自己的兒子，為我們眾人捨了他，豈不也把萬物和他一同白白地賜給我們嗎？ \end{tabularx} \\ \\ \relax
8:33 & \begin{tabularx}{0.7\textwidth}{X} 誰能控告神所揀選的人呢？有神稱他們為義了。 \end{tabularx} \\ \\ \relax
8:34 & \begin{tabularx}{0.7\textwidth}{X} 誰能定他們的罪呢？有基督耶穌 已經死了，而且復活了，現今在神的右邊，也替我們祈求。 \end{tabularx} \\ \\
[1ex]
\hline
\hline
\end{longtable}
$^{1}$各位弟兄姊妹早安.
大家見到禮堂的擺設.
有弟兄姊妹愛心的.
在昨天用了很多時間來佈置.
很想告訴我們復活節是值得慶賀的.
因為如果我們真的認識我們的救贖主耶穌.
我們的生命就會有改變.
正如剛才弟兄在領唱時所說.
不單止是耶穌的身體有復活改變.
其實在每一個信徒的心靈裡.
這個復活生命是已經開始的.
所以今早盼望藉著這段聖經.
跟大家一齊思想.
一齊從復活的主的生命進入我們裡.
來有這個改變.
有這個生命的變化.
跟大家有一個簡單禱告.
再一次我們將這段時間.
神的話語被打開.
在很多弟兄姊妹的心靈裡.
再一次被提醒.
聖靈的工作在我們大家的心靈裡.
來喚醒我們.
主耶穌已經復活.
我們等待的是再來的君王.
是這位將來要改變整個世界宇宙的救贖主.
這位彌賽亞.
所以我們今早在慶賀的過程.
我們的內心帶著盼望.
奉耶穌基督的名.
阿們.
選取這段聖經.
其實三十四節你就會知道.
誰還可以定我們基督徒的罪呢?.
因為有基督耶穌已經死了.
保羅特別強調.
而且是從死裡復活.
現在是坐在神的寶座右邊為我們代求.
所以我們要看到這個景象.
雖然物質看不到.

$^{41}$我們不會高一點.
四樓還是五樓.
就看到耶穌在那裡祈禱.
就不會的.
因為是天堂靈界.
我們要有一個按真理來看到.
這個永恆已經在發生.
同步發生.
主耶穌現在在天上為我們祈禱.
我們看到的是什麼呢?.
很多的生命的故事.
很多人的生命改變.
三十一節如果翻譯得好一點.
就是神幫助.
其實你看英文都已經知道.
如果是God for me.
for us.
如果是站在我們這邊.
那就沒有人可以抵擋我們.
神是不是站在你那邊呢?.
我記得我中一的時候.
由一間學校.
都已經是差的學校.
踢出校.
再去讀另一間.
要重讀.
那間學校就更加差.
全香港三間最差的其中一間.
進去之後.
你看到高聯班那些.
個個都凶神惡煞.
因為真的是Band 5的學校.
很害怕.
後來不用怕了.
為什麼不用怕呢?.
因為我的哥哥.
對上的哥哥.
讀高聯班的.
還是校隊的.
足球校隊.

$^{81}$那一群他的同學.
走進來我的班房.
那時候我是方斌仔.
這個是我弟弟.
那就是說.
其他人就不敢欺負我了.
因為高聯班的.
還是足球隊的校隊成員.
是我哥哥.
那一群他的同班同學.
踢球的那些.
就走進來我的小息.
走進來我的班.
你想想從高班.
高班是高一點層數的.
從六樓走進來二樓.
就說這個是我弟弟.
那就是說誰都不要碰他.
不過保羅這裡不是說黑社會.
保羅說的是.
神如果站在我們這一邊.
沒有人可以抵擋我們.
不過我們不要想錯了.
因為保羅很快就告訴我們.
一路去數後面那些.
35節一路去到38節.
是數很多很多的患難.
人生的不同的際遇.
大部分是不好的.
他說難道.
誰使我們可以跟基督的愛隔絕呢?.
難道是患難?.
難道是困苦?.
難道是逼迫?.
難道是飢餓?.
難道是赤身怒體?.
難道是危險?是刀劍?.
你可以想像得到.
保羅他列舉這些在全道的過程.
所面對的人生.

$^{121}$就是說原來.
我們不要有一個錯誤的想法.
耶穌復活了.
整個世界宇宙改變了.
任何的人都要對基督徒.
給幾分薄面我們.
沒有人會來騷擾我們.
令我們感覺在信仰上難堪的.
其實不是的.
保羅說世界沒有改變.
改變的是我們對神的愛.
對主耶穌的愛的認識和經歷.
這就會令我們仍然在這個墮落的世界.
宇宙裡面有力量面對我們的人生.
因為基督的復活.
使我們有力量面對人生.
就是說就算不好的時代.
你想想我們樓下的地鋪.
已經關閉了幾年.
都還沒有人來承租.
你可以想像得到.
如果現在你在元朗.
在洪水橋.
你想做食肆.
到星期六你就很擔心.
那些人會不會全部走到大陸.
就不來吃東西呢?.
其實這兩天已經在發生了.
昨天和我兒子出去旺角.
不用排隊.
泊車免費.
你可以想像得到.
時代可以是很特別的轉變.
突然間的改變會帶來很多我們措手不及的.
所以信仰告訴我們.
不是這個世界改變了.
是我們因為認識了主耶穌.
而生命改變了.
我們有力量面對明天.
所以不好的環境.

$^{161}$不會令我們.
如果有信仰.
有真理.
我們會有意志消沉.
沒錯.
我們會有困難.
和世人一樣面對困難.
但是我們有智慧.
我們有神.
我嘗試用一個例子.
大家可能聽過.
有一個叫王永慶.
台灣從前的首富.
經營塑膠業.
但原來他早期不是做這些的.
一個很有智慧的少年人.
16歲.
和爸爸借錢.
借了很多錢.
去到台灣.
不是台北.
去到中部嘉義.
一個地方.
很窮困的地方.
開米舖.
但那個地方.
其實有很多米舖.
他呢.
真是有智慧的人.
是很不同的.
16歲.
和兩個兄弟一起.
三個人.
就開米舖.
用爸爸給他們的資金.
有什麼不同呢.
你可以想像得到.
當時呢.
那些稻米呢.
因為那個技術.

$^{201}$不是那麼先進.
很多很多的米呢.
夾雜著很多沙石.
王永慶這一位的年輕人.
他看到.
我怎樣可以和其他.
米舖的人不同呢.
幾個兄弟.
選那些沙.
選了出來.
變了他的米.
大部分真是米.
不需要再去.
回到家裡.
要自己選.
然後.
還是第一個.
到回服務.
送貨服務.
送到你家裡.
這樣就很方便.
對於那些購買的人.
都不單止這樣.
去到之後.
將那個米缸.
舊的米.
先倒出來.
抹乾淨那個米缸.
然後將新的米倒下去.
再將舊的米放在最頂.
變成舊的米.
就不會不斷放在沉底.
這些是家籍的服務.
很快很快.
那些家家戶戶.
都看到這個年輕人.
他的經營.
真是帶來他們很多的方便.
你想想一個不好的時代.
很窮困的時代.

$^{241}$當時他的家境.
很窮困的.
但是.
他有這個思維的想法.
所以很快.
他的米舖.
就已經脫穎而出.
成為在那條街裡.
最出名的一個米舖.
很快就已經有他第一桶的金.
後來再轉型.
其他的事業.
各位弟兄姊妹.
保羅告訴我們.
我們是無懼那些困難的.
我們是因為神揀選.
這是八章的28節.
26節開始就說到.
聖靈在我們心裡的幫助.
令我們能夠有力量.
面對不同的事情.
為什麼呢?.
因為28節.
我們曉得萬事是互相的效力.
叫愛神的人得到益處.
就是按他旨意被召的每一個人.
按神的旨意被召的每一個人.
就是說原來這一份的恩典.
改變我們生命.
面對不同困難.
困境的那個力量.
是給每一個被召的人.
每一個都可以的.
不是那些特別高級的.
特別精英的.
才可以經歷這種能力.
不是.
每一個基督徒.
那哪些人是呢?.
被按他的旨意被召的人.

$^{281}$第29節.
因為他預先所知道的人.
神的預知.
如果你和我是被神揀選.
在神的愛裡面.
今天我們能夠成為基督徒.
你就是可以經驗上主.
復活大能的那個基督徒了.
這裡告訴我們.
而且給我們一個很美麗的人生.
就是預先定下.
效法他的兒子的模樣.
有一個模特兒.
就是耶穌基督.
他的美德.
他的美麗.
他的永恆的屬性.
當然不是神聖.
是這位主.
他的愛.
他的真理.
他的品格.
在我們每一個基督徒上.
神已經給了我們一個模樣.
所以基督徒原來是一條朝聖的人生的路.
很多弟兄姊妹不相信.
我年輕的時候.
二十多歲的時候.
是一天吃兩包煙.
滿口都是粗口的一個少年人.
每個星期.
最多站在主任室門口的就是我.
很多弟兄姊妹不相信.
我未信主之前.
開電單車.
穿飛仔衣服.
牛仔外套.
可惜那時候的數碼技術不夠.
很多照片都留不回.
如果不是就應該播出來給大家看一看.

$^{321}$這個就是神將我的生命改造了.
每個人的生命都不同.
但這裡告訴我們.
每個人都是預先定下.
效法他兒子耶穌基督的模樣.
使他的兒子在很多弟兄姊妹當中.
作長子.
預先定下的人.
就召他們來.
在萬世以前召我們回來.
所召來的人又稱他們為義.
就是神不再追究我們這個.
生命的舊人的罪.
稱為義的人.
我們可以得到榮耀.
神的榮耀.
各位弟兄姊妹.
這就是復活的意義.
一切必須是由復活的主.
他的復活開始.
在我們大家的生命裡.
有一位小朋友十二歲.
是外國人.
南非人.
故事就是這樣.
他是被收養的.
故事就在1997年11月18日.
在台北.
有一個逃犯叫陳俊興.
十惡不赦的.
殺了很多人.
不斷被追捕.
不斷逃走.
在那晚.
七點多晚上.
他就走去台北南非的一個大使館.
走進去之後.
挾持裡面的人.
第一個挾持的就是這個女孩.
十二歲的.

$^{361}$她叫Christine.
挾持了這個小朋友.
然後就叫他.
抓住他上樓.
因為他的父母在上面.
是收養的父母.
這個南非的一個將軍.
是跟台灣有邦交的.
很特別的地方就是.
台灣跟南非差不多要斷交了.
所以那個大使館.
一個月之後就要close(關閉).
他們都準備全家要回南非.
突然之間出現了這個陳俊興.
闖入他們這個大使館.
就挾持住人質.
然後他要公諸於世.
就叫南非的大使.
來通知全球的最重要的媒體.
CNN和其他的媒體.
就是要全世界看著他現在.
挾持人質的事件.
於是警察不讓步.
對戰 槍戰.
就射傷了大使的太太.
大使的女兒.
後來陳俊興走火的那支槍.
射傷了他的大使.
這位中文的翻譯.
就是 忘記了名字了.
兩個人傷了.
警察在外面隨時要衝進來.
但這個兇徒就說.
如果你們衝進來.
我就要殺光那些人質.
於是談判 輾轉談判.
後來終於肯放兩個傷者離開.
然後僵持住.
一直到深夜.
這個十二歲的女孩.

$^{401}$是一個很好的基督徒.
我們沒辦法想像.
面對一個非常兇惡的陳俊興.
他就拿一支筆在畫畫.
畫了什麼呢?.
畫了一個十字架.
然後有一個大的愛心.
包圍著十字架.
畫完之後.
他內心有感動.
他說他畫的時候.
其實是很害怕.
他祈禱神幫他.
讓他不要害怕.
但沒多久之後.
他突然間想到.
其實這個陳俊興.
也需要神的愛.
他將他畫好的這幅圖畫.
送給了陳俊興.
這個陳俊興突然間.
看到這幅畫之後.
是很愕然的.
為什麼這個女孩.
她會不害怕.
送這幅畫給我呢?.
然後他有很大的感動.
從來沒有人敢夠膽.
在他面前這樣.
大家四目交投.
每次見到他.
都要不被他殺.
要不被他凌辱.
之前97年的時候.
這個陳俊興還犯下了很多罪案.
有十幾個婦女被他強姦.
這幅圖畫.
令到陳俊興軟化下來.
然後願意談判.
來妥協.

$^{441}$來投降.
於是警察進來.
就抓了陳俊興.
這幅圖畫.
改變了陳俊興的內心.
更加令我們沒法想像.
一個月之後.
因為當時是1997年11月18日.
這件事發生.
在12月的聖誕節.
因為南非的大使.
很快就要回南非.
臨走之前.
這對夫婦.
去監獄探望陳俊興.
說我們不會怪責你.
已經饒恕了你.
那個大使被他射傷了膝蓋.
他的大女兒中槍危殆.
這樣一個用愛去寬恕.
這個十惡不赦的囚犯.
這個舉措.
令到陳俊興.
真的決心信主.
由一個12歲的小女孩.
一幅圖畫.
一個很勇敢的心.
可以改變一個.
十惡不赦的大罪犯.
所以陳俊興洗禮那天.
很多人很不開心.
因為沒可能的.
基督的愛可以寬恕這個罪人.
不可能的.
可能他是一個扮演的.
很兩極的.
在台灣的輿論.
當時是1998年的時候.
1999年陳俊興.
說完最後的見證.

$^{481}$就來執行死刑.
他回到天家.
不過.
你可以想像得到.
一個小女孩.
裡面有神的愛.
令她鎮定.
裡面有神的愛.
令到她能夠立即.
可以用這幅圖畫.
給這個罪犯.
改變這個罪犯的生命.
弟兄姊妹.
我們往往不希望有不好的事情.
發生在我們身上.
但是原來.
聖經告訴我們.
萬事是有神掌管.
互相的效力.
使愛神的人會得到益處.
我們是屬主的兒女.
我們盼望的是永恆的生命.
世界沒有改變.
不會因為今天是復活節.
很多人會有生命的轉變.
如果我們沒有主在我們內心.
復活節只是多了一天假期.
讓你可以多住兩天.
在深圳.
或者去日本.
去放鬆.
去休息.
但是有復活主的生命在我們裡面.
我們就可以面對不同的人生困難.
再說一個例子.
最近我也不斷地說.
這位姊妹是很奇妙的一件事.
我逢星期四晚.
在錦繡堂有查經班.
有一天.

$^{521}$有一位我不認識的姊妹.
就來找我.
跟我說.
你能不能把你的WhatsApp電話.
給我的妹妹呢?.
因為她很想問我一些.
聖經神學的問題.
我一點都不介意.
牧者當然是很高興.
她問准我之後.
把我的電話給了她的妹妹.
她的妹妹是住在三藩市的.
這位姊妹.
去三藩市.
在香港住的姊妹.
住在沙田.
姓陳.
去三藩市探望這位妹妹.
她的妹妹原來在2018年的時候.
開車上班.
差不多還有兩個街口就到達公司.
被一個大學生的車輕微碰撞.
但可能她不知道是什麼原因.
不知道開車還是怎樣.
很害怕還是怎樣.
碰撞之後.
把她的車推出去公路.
這位妹妹第二次被高速公路的車.
衝過來撞.
她的車完全損毀.
由那個時候開始.
她在頸以下.
沒有能力活動.
每天98\%.
需要有專人照顧.
你可以想像到那種困難,痛苦.
這樣的人生.
而且會繼續.
到離開世界才能改變.
說回剛才的故事.

$^{561}$姊姊為什麼要拿我的WhatsApp呢?.
原來姊姊去探望她的時候.
在早期去探望她的時候.
姊姊在感受堂聽了很多講道.
又聽到我的茶經班.
覺得挺好,挺舒服.
就拿這條連結給她的妹妹.
在三藩市.
妹妹就不知道什麼原因.
一直聽了一年.
她還沒相信耶穌.
她第一次跟我說話.
其實挺生氣的.
她說:曾牧師.
你說的道理我完全聽不懂.
不知道你說什麼.
但不知道為什麼.
我聽完你講了這一年.
我的心裡很平安.
我內心裡有一個盼望.
我不懂得說.
我不知道為什麼.
我就很簡單地跟她在WhatsApp裡對話.
邀請她祈禱.
去倚靠神.
過了一段時間.
她跟我說.
她說她現在沒有那個迷惘.
上個星期.
她說她堅持要奉獻給紅恩堂.
在三藩市的姊妹.
她問了一個不知如何回答的問題.
其實很簡單.
她說:我回不了教會.
神會不會不喜歡我?.
我說:你不是回不了教會.
你是暫時不能回教會.
你的行動是需要你的先生.
你的家人幫你.
你才可以去教會.

$^{601}$你不可以自己去教會.
她會問這樣的問題.
一個完全不知道什麼叫基督教的人.
她問我什麼叫三一論.
因為我的查經班有講.
她很想了解一下什麼叫三一論.
她很想了解一下.
怎樣的一個永恆生命.
她已經信了主.
在主裡面的弟兄姊妹.
她現在為她的先生每天祈禱.
希望先生的工作順利.
因為COVID(中共病毒)倒閉,失業.
要轉行.
要在唐人街開車開很多個小時.
很多的轉變.
但是弟兄姊妹.
基督的復活生命進入了她的心.
有時她會分享:我很無助.
我要等丈夫起身.
才可以照顧我.
每天就是在等.
你試想一下.
從現在開始.
你坐到下午.
其實我們都不能夠接受.
何況一個人已經繼續.
只能夠在一個這樣的輪椅裡面.
等待別人的幫助.
但是她裡面的內心有了盼望.
裡面的心靈有了耶穌.
復活的主在她的生命裡面.
今天是復活節.
復活節告訴我們.
不是所有的事物改變.
其實世界沒有改變.
剛才主席也說了.
俄烏的戰爭.
以巴的戰爭.
無情的轟炸.

$^{641}$還有一些更加慘絕人寰的.
因為我們有時也要追蹤一些新聞.
東歐的娼妓的問題的嚴重.
很多少女在羅馬尼亞.
在那些很窮的東歐國家.
被欺騙去了西歐.
大部分去了英國.
就被鎖在一些地下室裡面.
暗無天日.
任由那些嫖客來虐待.
不受煩的就毒打至死.
大有人在的.
在今天這樣的現代的21世紀.
你可以想像到.
罪惡從來沒有消失.
不會因為耶穌的復活.
不過弟兄姊妹.
我們經歷了主耶穌的復活.
所以一切的生命改變.
是由我們內心裡面的主開始的.
保羅沒有告訴我們.
一個美麗的圖畫.
這個世界因為信了耶穌.
就沒有任何的痛苦困難.
保羅告訴我們.
他為這個福音.
他願意忍受各種的困難.
因為他真的親眼見過復活的主.
向他這個人顯現.
他的生命現在成為了寶貴.
他不會覺得對外邦人傳道.
對他來說是一種侮辱.
因為他是一個猶太的拉比.
他很不喜歡外邦人.
但是主耶穌改變了他的心靈.
他看外邦人是主的兒女.
改變了他的光彩.
是他的榮耀.
他要將人帶進那個榮耀裡面.
所以弟兄姊妹.

$^{681}$盼望我們今天走出我們的會堂.
我們帶著的是基督復活的榮耀.
我們自己也都經驗上主的能力.
在我們的生命裡面.
我們一起祈禱.
感恩.
多謝主.
我們可以慶賀基督的復活.
主在我們內心成為榮耀的盼望.
願你復活的大能.
在我們每一位弟兄姊妹.
每一個處境.
每一個困難裡面.
有些可能都不足為愛人道的困難.
我們都可以向主來交託.
更加我們有力量去面對人生.
即使人生未必如我們所願的美麗.
但我們知道只要我們有主.
那個地方.
那個情景.
那個人生.
都會變得很不同.
因為主是完全的給我們的.
願神得到一切的榮耀.
奉主耶穌基督的名.
阿們.
\newpage



\section{ }
\label{sec:Ux8UdMSUx6c}
\textbf{2024.03.31 復活節佈道會《活在每一天》 李炳光牧師}
\newline
\newline
連結: \href{https://youtube.com/watch?v=Ux8UdMSUx6c}{\texttt{https://youtube.com/watch?v=Ux8UdMSUx6c}} ~~~~ 語音日期: 2024-04-02
\newline
\newline
\hyperref[sec:hGBqBs0DEW4]{\small{< < < PREV SERMON < < <}}
~
\hyperref[sec:index_chronic]{\small{[返順時目]}}
~
\hyperref[sec:index_scriptual]{\small{[返順卷目]}}
~
\hyperref[sec:EYrQ_D9F5Wo]{\small{> > > NEXT SERMON > > >}}
\newline
\newline
$^{1}$其實這齣戲是由我們自己編寫的.
好像是開玩笑一樣.
好像是自己編寫.
但是這齣戲背後的意思是什麼呢?.
就是說一個很有問題的家庭.
這個家庭已如我詐.
大家彼此欺騙過了一種不正當的生活.
你看見他們的情況就知道了.
但是當他們發現到.
原來人生是這麼短暫的.
隨時可能是沒有機會的.
他們才發現到.
原來我們這樣對自己的配偶.
對自己的爸爸.
他們就良心發現.
當人發現到最後的時刻.
就會覺得自己真的要把握機會.
其實在我們平靜安穩的時候.
我們不會想這個問題.
但當你發現到自己有問題的時候.
就發現到我們應該要把握機會.
來做正 說正.
這個故事的結束似乎是令人失望的.
因為以為沒事.
誰知道原來真的有事.
結果怎樣呢?.
我不回答了.
因為你自己想了.
但是這個故事讓我們看見.
當我們需要到一個問題發生的時候.
我們立刻就發現到.
原來人生是非常短暫的.
家庭的夫婦 妻子兒女的相處也是短暫的.
如果我們不能真誠相處的時候.
我們也是活在虛偽的當中.
所以各位朋友.
今天這齣話劇是提醒我們.
其實我們應該活在每一天.
因為每一天都有事情發生.
我們不知道將會發生什麼事.

$^{41}$這件事讓我們想到.
我們怎樣過著一個有意的人生.
每一天都過著有意的人生.
我想這個題目的時候.
我立刻想到一件事.
我有一個教友.
他是一個很有名的醫生.
但突然間他很熱心地回來我茶經班.
奇怪 他平常都很忙的.
為什麼來參加我的茶經班呢?.
原來他告訴我.
他快要說再見了.
原來他患了白血球過多症.
他很害怕.
他叫我為他祈禱.
一邊說一邊哭.
他又不敢告訴太太.
和他媽媽和爸爸.
在這個時候我只是跟他祈禱.
我又不是醫生.
他在醫院進行的時候.
原來他自己因為是醫生.
他可以想到一個出路.
白血球過多只是換骨髓.
但是他自己製造骨髓.
也就是用幹細胞製造.
大家會覺得真假.
結果他真的成功了.
在多個範例當中.
他是成功的一個.
他告訴我.
當他在醫院掙扎的時候.
痛苦的時候.
他用幹細胞來克服疾病的時候.
實在非常痛苦.
有一天他睡覺的時候.
有個護士進來.
拿他的手來打針.
誰知護士看著他的手.
他說你的手已經打了針.

$^{81}$不能再打下去了.
他自己是醫生.
他拿來看一看.
原來這些針打一次的地方.
第二次不能再打了.
我也不懂.
是不是這樣我不知道.
結果他說.
明天我就再見了.
因為我不能再打針了.
這種針原來是這麼重要的.
不能再打了.
他自己交給上帝了.
他是一個很難得的.
因為他患了腫瘤的時候.
他發掘到自己有一個寫作的恩賜.
他就寫了一個故事出來.
我讀了之後很感動.
原來他第一天起床的時候.
有個男護士.
年輕人來到他前面.
醫生今天要打針了.
他說還打什麼.
我沒得打了.
但這個男人不理會他說什麼.
轉動他的手.
這裡還有一針.
轉到他臉上.
原來下面還有一個地方可以打.
他看到原來這個男護士.
這個年輕人竟然找到一個針口.
讓他可以打下去.
結果他打了這針.
真的勝過這個患疾.
結果他痊癒了.
現在他在澳洲居住.
非常精彩.
有空畫畫.
有空寫文章.
很精彩的文章.

$^{121}$你猜他的文章說什麼呢.
他說我們不知道將來會怎樣.
但只有活在每一天.
如果這個家庭.
剛才演戲的家庭.
能夠真真正正活在每一天的時候.
家庭不會變成這樣.
就算大亞灣電廠爆炸.
那又怎樣.
大家一起去祭拜.
人生就是這樣.
原來有很多事情.
我們不止預測.
活在每一天.
這位醫生寫了這個題目.
活在每一天.
我想到這個題目的時候.
在復活節下午.
跟大家分享一個很重要的訊息.
各位.
有位老人家.
他很大年紀.
但由於家庭經濟有問題.
他每天都推車到外面做小販.
但有一天他遇到交通意外.
有輛車撞到他的小販車.
結果他自己受了傷.
小販車當然是全散了.
他太太趕到醫院.
很久後醫生出來.
在手術室說.
出來的時候問.
阿婆我想問問你先生今年幾歲.
哇 八十幾歲了.
什麼八十幾歲.
為什麼他的身體健康不像八十幾歲.
什麼醫生你聽說他不像八十幾歲.
是啊 我驗他的身體的時候.
我跟他做檢查的時候.
我發現他的身體很多方面都很健全.

$^{161}$不像八十幾歲.
為什麼.
為什麼你這麼大年紀還出來做小販.
沒辦法 他喜歡做的.
當然推翻他.
但這個醫生說.
他真的不像這麼大年紀.
你告訴我他有什麼秘訣.
怎麼能夠活得這麼好.
這位阿婆就說.
醫生我告訴你.
我老公有個習慣.
我時常叫他小心.
但他每次出去的時候.
他說雖然我也有很多不妥的地方.
但我相信明天會好一點.
明天會好一點.
他就活在每一天.
盡他的力量活得精彩 活得活潑.
各位朋友.
活在每一天其實是一個秘訣.
其實我和你都不知道將來會怎樣.
但如果我能夠真正把握每一個機會.
能夠盡量活得好的時候.
將來發生什麼事都不要緊.
最重要的是我們活得精彩.
最近有人叫我錄影做港島.
有個題目叫做.
「我知誰掌管明天」.
我就引用這個故事.
真是 各位.
原來今天我們不知道將來會怎樣.
大家知不知道有一句說話.
「明天會更好」.
這句說話是怎樣寫的.
原來這句說話是江澤民寫給我們的特首.
我曾經有機會去過特首的辦事處.
幸好我做教會以外的一些服務工作.
我有機會見到她.
原來她寫著貼著江澤民的手筆.

$^{201}$「明天會更好」.
這個作者是誰呢.
這個作者原來是台灣台北衛理堂的師班指揮.
叫做詹宏達先生.
他在囚過.
有一天他來找我.
他說牧師我聽到你在港島提過我的名字.
所以今天想見你.
我說甚麼事.
我跟他說了一會兒.
就說到這首「明天會更好」.
我就貪得意跟他說.
我說詹先生你不要介意.
我想改一下你的歌.
「明天會更好」.
沒錯.
但是我想加一句.
「今日都唔錯」.
他看著我.
睜著眼睛看.
甚麼今日都唔錯.
我說真的.
我說詹先生.
今日如果不好.
你想明日會好嗎.
是不是.
我們今日都唔錯.
所以我希望大家.
今日都唔錯.
明天會更好.
好不好.
各位.
所以復活節.
其實是讓我們知道.
我們復活.
不是單說到耶穌基督死了.
他就復活.
其實復活節的故事.
是繼續發生.
我曾經有一個題目說到.

$^{241}$就是復活的故事.
還未說完.
為甚麼.
因為每天都在發生.
為甚麼.
我們說復活.
不是單單肉體的復活.
是心靈的復活.
今日我們很多人.
活在死人的環境當中.
每天都被罪惡.
被困擾.
被不安.
來困擾我們.
我們沒有平安.
我們沒有喜樂.
各位.
我想要你的同意.
其實我們都生活在得失之間.
其實我們有很多事情.
是遭遇到的.
甚麼是得失.
生死.
成敗.
愛恨.
貧富.
凶吉.
樂憂.
疾病和健康.
信逆.
安危.
其實這些都是我們要面對的困難.
贏輸.
高低前後.
進退榮辱.
生死.
都是影響了我們的情緒.
我們不平衡.
各位.
我做牧師已經做了很多年.

$^{281}$我告訴大家的時候.
你可能會想我幾多歲.
我去年因為慶祝85歲生日.
嘩.
像不像.
我去年.
我去年因為慶祝85歲生日.
我告訴大家.
我做了60年牧師.
他們說真的嗎.
其實我見過很多問題.
今天很多人.
你見過好好的.
其實他生活在很多問題當中.
今天我剛見到一個精神病醫生.
我跟他說.
我剛有個教友.
他的女兒突然情緒不安寧.
每天都不理自己的女兒.
牧師怎麼辦.
我去見見醫生好不好.
他不肯去.
今天我跟醫生說.
你叫他來吧.
你吩咐我做就行.
因為他是我的教友.
我覺得今天很多人.
都活在情緒當中.
各位知道張國榮是怎樣死的嗎.
其實他不是那麼大膽.
那麼高樓跳下去的.
其實那些死是有聲音的.
這個聲音是什麼.
是情緒的問題.
精神的問題.
聽到這些聲音.
跳下去是沒有問題的.
有個女醫生因為面對考試.
突然間打電話給我.
問我什麼事.

$^{321}$她說有聲音跟我說.
叫我吃藥.
這些藥我知道.
吃了一定死.
牧師怎麼辦.
我問他為何突然這樣.
他說不知道為何.
我經常聽到聲音.
各位如果聽到聲音.
要小心一點.
不要以為這些聲音是真實的.
其實是你自己的幻覺.
你聽到這種聲音提醒你.
叫你不要害怕.
跳下去吧 吃了它.
其實今天我們很多人.
都因為這些不平衡的環境.
很多高低起落.
生死考試等等的事.
影響了我們寧靜的內心.
各位你有沒有試過睡不著覺.
其實睡不著覺.
這是很正常的.
其實不正常.
其實要小心.
原來有很多精神病醫生告訴我.
一個失眠的人是很危險的.
為什麼.
因為早上天空亮的時候.
是最危險的時候.
原來我曾經查過.
林黛 樂蒂 還有幾個明星.
都是自殺死了.
他們都是早上吃藥.
原因是什麼.
想不通.
最重要的是什麼.
這個女醫生很聰明.
打電話給我陪著她.
馬上打電話給她的父母.

$^{361}$她說父母去旅行.
快點回來陪著她.
現在女醫生在屯門醫院做得很成功.
各位這些都是真實的事實.
各位朋友.
復活節的故事仍然是說下去的.
耶穌基督沒錯.
在死裡復活.
帶給我們希望.
人在絕望中.
困難中.
死亡中.
我們看到耶穌基督復活.
讓我們知道我們有機會.
有盼望.
將來不是走一條骨頭的道路.
今天我告訴大家.
不是等到要死的時候.
才要走一條骨頭的道路.
今天很多人的心靈沒有平安.
平安很重要.
平安就是沒有病沒有痛.
就是平安.
其實今天有病有痛.
當然是沒有平安.
但是沒有病沒有痛.
都沒有平安.
原因是什麼呢.
就是平安不是遭受什麼困難.
而是我們內心裡.
不得寧靜.
各位.
我引用一句說話.
我希望你們能夠聽得進去.
我說.
今天我看到很多人.
帶著昨天的倦態.
明日的仇容.
過現在的日子.
對不對.

$^{401}$為什麼我會這樣說呢.
有一次我在學校教書.
我教聖經.
我看到同學.
好像很憂愁.
問我做什麼事.
牧師明天要考試.
考試就是考試.
那怎麼辦呢.
不是啊.
我昨晚沒有睡覺.
昨晚沒有睡覺.
所以我今天很累.
但是昨天沒有睡覺.
為什麼要這麼不精神.
剛才不是說了.
我明天要考試.
帶著昨晚沒有睡覺.
明天考試的心情.
過現在的日子.
各位.
你們不用考試.
但是我告訴大家.
有很多時候.
猶憂千年無米煮.
猶憂什麼.
無命享千年.
其實我們每天都是憂.
又不知怎樣.
每天都過著困擾不安的心境.
各位朋友.
今天我要說復活節的訊息.
不是說復活節.
而是因為耶穌在死裡復活的時候.
讓我們知道.
人生是有盼望的.
人生是有機會的.
而且不要被這些事情困擾我們.
各位.
復活是表示什麼呢.

$^{441}$一個人在罪惡當中.
其實是一種死亡.
一個人如果有心靈的不平安.
是另一種死亡.
如果有心靈的不平安.
是另一種死亡.
其實.
今天我們活在不安的狀態當中.
是另一種死亡的表現.
復活是讓我們知道.
這種可以擺脫出來的.
所以我今天願望大家.
能夠有復活的生命.
什麼叫復活的生命.
不是叫你死了又復活.
而是你有復活的生命.
即是什麼呢.
就是你活在不平安的狀態當中.
復活的生命.
即是什麼呢.
就是你每天都活潑潑的生活.
過著喜樂平安的生活.
好不好.
今天很重要的.
就是希望大家能夠.
過著喜樂平安的生活.
怎樣過著喜樂平安的生活呢.
我希望大家能夠.
真正找到一個出路.
這個出路.
就是會讓你們知道.
原來神帶領著我們.
為我們開一條出路.
我在做牧師這麼多年.
我見到很多復活的生命.
我退休之後.
我一直在組織報道團.
有人問.
牧師你不悶嗎.
經常說復活.

$^{481}$經常說耶穌.
經常說是信仰.
喂 你不悶嗎.
我告訴你.
當我見到大家.
決志信耶穌的時候.
我開心的程度.
是比較中馬錶.
現在不流行買馬錶.
或者買六合彩.
更加開心.
是的.
因為一個人舉手.
接受過著新的生命.
這個就是復活的生命.
是不是.
我希望大家.
在復活節的時候.
真正享受復活的生命.
我要告訴大家.
一些很真實的見證.
我見到有很多人.
走到絕路的時候.
突然間改變了.
今天有幾個.
很天真很可愛的小朋友.
我為他們感恩.
有一間學校.
是我們教的學校.
專門收藏.
有問題的小學生.
這間學校.
有很多學生是自閉症.
有很多學生是精神不全.
有一個校長叫我去報道.
我說糟糕了.
這次去廣道的時候.
怎樣能夠令他們信耶穌呢.
他們聽不明白.
但是很奇怪.

$^{521}$他說我去廣道的時候.
他們的反應很好.
很留心聽.
我說給你聽.
很多人害怕.
原來那些自閉症的學校.
學生帶著鋼盔.
來聽道.
為什麼.
原來我過了時間.
自閉症的小朋友.
撞頭撞牆.
帶著鋼盔.
那些老師跟我說.
牧師你不要害怕.
原來自閉症的小朋友.
過了時間很敏感.
過了時間就大叫.
或者發脾氣撞頭.
我心理準備.
沒有過時.
沒有撞頭撞牆.
但是他們竟然接受福音.
我叫他們相信耶穌.
他們舉手.
我叫他們走出來的時候.
他們舉手.
我叫他們跟我一起祈禱.
我看到他們流眼淚.
過了一段時間.
有一個家長.
牽著小朋友.
來到校長的辦公室.
一進來的時候.
校長還以為是甚麼事.
突然間校長.
我想問問.
誰叫我兒子信耶穌.
嘩.
害怕.

$^{561}$不是我.
是李牧師.
誰知這個家長說.
是李牧師.
我怎麼找到他.
他不是回學校.
不如我和你約他.
不要了.
現在我也不方便.
你和我多謝李牧師.
原來有一件事發生.
這件事發生是.
這個家長說了一件很真實的事.
有一個賭徒.
把他兒子丟到街上.
然後自殺.
聽到後很傷心.
原來這件事發生在家長身上.
原來這個家長.
他是做股票行的.
他犯了法.
結果.
可能大小報紙可以看到.
結果他本來.
很舒服的.
很無雜.
又很舒服的生活.
雖然兒子不是很正常.
但因為這個緣故.
崩潰.
失去工作之外.
還要破產.
搬到沙劇灣的tong 房.
這個家長.
老婆又走了.
有一天.
家破人亡.
他做家長的.
他告訴校長.
校長有一天想不通.

$^{601}$自己哭得很厲害.
對不起我的太太.
對不起我的小孩.
現在只剩下我一個.
我怎麼能養大他.
於是他牽著兒子.
跟兒子說.
兒子 爸爸沒用.
爸爸不能夠帶大你.
你跟我一起上天台.
我和你去一個地方.
我和你一起.
你不用掛心 很快的.
為什麼要去一個地方.
你先跟我上來.
他上天台的時候.
家長說.
你跟我一起跳下去.
很快的結束我們的生命.
不在這個可怕的世界生活.
他一直說到哭.
誰知道.
他不停地說.
誰知道.
每個家長.
在這個最艱難的時刻.
每個家長都有問題.
兒子竟然說這樣的話.
他說.
爸爸 你千萬不要.
耶穌愛你.
誰說的.
牧師說的.
耶穌愛我.
耶穌怎麼會愛我.
是啊 耶穌愛你.
你不要這樣做.
結果.
這個傷心的父親.
抱著兒子.

$^{641}$跪下在地上.
兩父子大哭一場.
結果有沒有死.
沒有.
他下來.
放棄了這個可怕的念頭.
現在他們重獲新生.
他們的心.
已經重獲新生.
這個就是復活.
我覺得這個就是復活的故事.
不斷地發生.
今天可以在我們的日常生活當中發生.
你覺得失望.
我告訴你.
有希望.
你覺得痛苦不安.
我告訴你.
上帝會為你開一條出路.
如果你覺得有很多不平安的心境.
耶穌基督說.
在我裡面有平安.
你們心裡不用憂愁.
各位朋友.
復活節的訊息.
最感人的地方.
不只是說回.
耶穌基督以往.
死而復活的故事.
而是今天仍然發生在我們心中.
其實我今天沒時間說這麼多.
因為我可以說很多這些故事.
大家聽到.
但是最重要的就是.
我自己感受到復活的真實.
我希望各位朋友.
都能夠感受到復活的真實.
在你的家庭當中.
在你的心靈裡面.
在你的朋友當中.

$^{681}$在你的家裡的人當中.
繼續發生.
你想不想.
耶穌基督讓你能夠.
享受一個復活的生命呢.
各位.
我希望各位朋友.
各位.
我盼望大家.
都能夠在復活節的時候.
過一個不一樣的復活節.
這個不一樣的復活節.
就是說有復活的主.
在你的生命當中.
帶領你前面的道路.
有主引領你.
一些覺得困擾不安.
不明的將來.
但是.
聖經說.
我知誰掌管明天.
有神掌管明天.
你知道復活節的故事.
繼續的發生.
在復活節之後.
耶穌顯現.
每一個故事.
都是人家最困難的時候.
最絕望的時候.
耶穌就在他旁邊.
今天耶穌仍然在你旁邊.
在你的生活當中引導你.
使你渡過這個難關.
你想不想依靠他呢.
你想不想伸出你的手.
讓他拖帶你呢.
如果你能夠的話.
今天你來這裡.
沒有枉費.
我盼望你能夠過一個.

$^{721}$有意義的復活節.
現在.
我希望大家做個決定.
如果.
你想接受這個復活的主.
在你的生命當中.
能夠真正的.
引領你的前途.
帶領你的將來.
帶領你的家庭.
帶領你個人的生活.
活在每一天.
好不好.
如果你想這樣做.
想我為你祈禱.
請你舉高手來.
祖母 舉高手來.
不要害羞.
祖母 舉高手來.
祖母.
祖母 請你舉高手來.
今天第一次來的朋友.
看到了 看到了.
好難得.
還有哪一位.
剛才舉手的.
請你站起來.
站起來 我為你祈禱.
還有哪一位.
站起來 我為你祈禱.
還有哪一位.
沒有人相信耶穌.
你的心裡沒有平安.
如果你想得到平安的心境.
你想有復活的生命的話.
請你站起來 好不好.
還有哪一位.
我想你鼓勵.
我想在座的基督徒.
基督徒.

$^{761}$請你一起支持他們.
支持他們 好不好.
你們站起來 基督徒.
如果你發現到.
有些還沒有站起來.
你鼓勵他們幫忙.
鼓勵他們一起.
好.
現在差不多.
全部站起來了.
我特別有一個要求.
如果你第一次站起來.
第一次舉手.
或者你是基督徒.
你想今年這個復活節.
那個有意義的復活節.
讓神帶領你.
享受這個復活的生命.
請你走出來.
我一起跟你祈禱.
好嗎 走出來.
還有哪一位.
請你走到前面來.
不要錯過這個機會.
還有哪一位.
請你跟我一起祈禱.
還有哪一位.
聖靈感動你.
願意是你的.
願意是接受這個挑戰.
聖靈感動你.
使你能夠過著.
活潑的生命.
復活的生命.
是不是很難得.
請大家在後面鼓掌.
好.
今天.
我引領大家.
做第一次的祈禱.

$^{801}$決志的祈禱.
或者重新的.
或者重新建立信心的祈禱.
好嗎.
我講一句 你跟我講一句.
在座當中 你願意參與的.
一起好嗎 我們同心祈禱.
親愛的主耶穌.
多謝你.
因為你復活.
帶給我們有希望.
帶給我們有出路.
讓我們知道.
人生有困難.
有很多我們不歡喜的事.
但是主沒有離開我們.
主答應我們同在.
與我們同行.
今天就是復活節.
我們願意跟隨你.
一起享受復活的生命.
天父啊.
幫助我們.
活在每一天.
每一天都有你在我們當中.
每一天都有新的機會.
每一天都有新的人生.
多謝天父.
願你繼續與我們同行.
今天在你面前.
我們願意決志.
跟隨耶穌.
你赦免我們的罪.
從新開始.
成為你的兒女.
求你接受我們.
我們就不一樣的人生.
我們享受你新的生命.
多謝天父.
我們懇求祂禱告.

$^{841}$奉主耶穌的聖名.
我們懇求祂禱告.
奉主耶穌的聖名.
阿們.
Hallelujah.
\newpage



\section{雅各書 1:1-8}
\label{sec:EYrQ_D9F5Wo}
\textbf{2024.04.07 《作個成熟的基督徒》 雅各書1\_1-8 (和修版) 講員 : 曾敏姑娘}
\newline
\newline
連結: \href{https://youtube.com/watch?v=EYrQ-D9F5Wo}{\texttt{https://youtube.com/watch?v=EYrQ-D9F5Wo}} ~~~~ 語音日期: 2024-04-08
\newline
\newline
\hyperref[sec:Ux8UdMSUx6c]{\small{< < < PREV SERMON < < <}}
~
\hyperref[sec:index_chronic]{\small{[返順時目]}}
~
\hyperref[sec:index_scriptual]{\small{[返順卷目]}}
~
\hyperref[sec:i_n9uN87S9M]{\small{> > > NEXT SERMON > > >}}
\newline
\newline
雅各書 1:1-8
\newline
\begin{longtable}{cl}
\hline
\hline
章節 & 經文 (和合本修訂版)\\
\hline
1:1 & \begin{tabularx}{0.7\textwidth}{X} 神和主耶穌基督的僕人雅各問候散居在各處的十二個支派的人。 \end{tabularx} \\ \\ \relax
1:2 & \begin{tabularx}{0.7\textwidth}{X} 我的弟兄們，你們遭受各種試煉時，都要認為是大喜樂， \end{tabularx} \\ \\ \relax
1:3 & \begin{tabularx}{0.7\textwidth}{X} 因為知道你們的信心經過考驗，就生忍耐。 \end{tabularx} \\ \\ \relax
1:4 & \begin{tabularx}{0.7\textwidth}{X} 但要讓忍耐發揮完全的功用，使你們能又完全又完整，一無所缺。 \end{tabularx} \\ \\ \relax
1:5 & \begin{tabularx}{0.7\textwidth}{X} 你們中間若有缺少智慧的，該求那厚賜與眾人又不斥責人的神，神必賜給他。 \end{tabularx} \\ \\ \relax
1:6 & \begin{tabularx}{0.7\textwidth}{X} 只要憑著信心求，一點也不疑惑；因為那疑惑的人，就像海中的波浪被風吹動翻騰。 \end{tabularx} \\ \\ \relax
1:7 & \begin{tabularx}{0.7\textwidth}{X} 這樣的人不要想從主那裡得到甚麼。 \end{tabularx} \\ \\ \relax
1:8 & \begin{tabularx}{0.7\textwidth}{X} 三心二意的人，在他一切所行的路上都搖擺不定。 \end{tabularx} \\ \\
[1ex]
\hline
\hline
\end{longtable}
$^{1}$(自從2014年7月1日開始,我們開始了一同禱告).
天父上帝,我們讚美你,因為你統管萬有.
所以今天我們才可以在上帝的國度裡有份.
在神的家裡有份.
今天心願你對我們眾人說話.
願意聖靈在我們心中去工作.
讓我們切實地聽到上帝你的聲音.
又很樂意地去回應神的教導.
阿們.
最近教會開始很多更新.
很多優化的工程,很多改變.
在過程中,我自己去想.
除了在硬件上去更新,去優化.
更重要的是我們生命的成長和成熟.
大家想像一下,將來有很多新朋友進入我們當中.
和我們一起聚會.
在這個時候,我和你更加重要.
我和你要做什麼呢?.
禱告.
新朋友來,我們就開始禱告.
禱告是很重要的.
預備自己.
預備我們的生命.
將他們擁抱,帶他們去耶穌基督的面前.
所以我們這個人是很重要的.
今天很想和大家說說.
作一個成熟的基督徒是怎樣的呢?.
我相信今天的內容不會停在這個題目.
只是一個開始.
整本聖經有很多位置帶領我們去學習.
怎樣去成為一個成熟的基督徒.
第一個成熟的特徵是謙卑.
我們去看自己,不要去看別人.
看自己是否成熟.
第一,我是否謙卑的人.
要做到謙卑,是否容易呢?.
(我看到有些頭在動,不容易).
為什麼不容易呢?.
要謙卑,要牽涉很多東西.
謙卑和自卑是不一樣的.

$^{41}$通常自卑是很難謙卑的.
自卑很容易就會驕傲.
因為自卑需要有些東西填滿我的心靈.
所以很容易,有些東西我要告訴你.
我是多麼的厲害.
真的很不容易,所以要謙卑下來.
謙卑是那個人真真正正對自己.
對上帝的信是很大的.
能夠從上帝的角度去看自己.
這是不容易做到的.
直到今天我還在學這個功課.
怎樣在神的眼中看我自己.
看得合乎中度.
一自卑就很容易驕傲.
成熟特徵是謙卑.
為什麼我一開始就說這個呢?.
《雅各書》的作者是誰呢?.
短短的一句話.
我們看到他的生命.
他的氣質是很不同的.
一個人怎樣介紹自己是很不同的.
一開始,我可以拿很多頂的高帽來跟你說.
我在哪裡畢業的?.
我之前做過什麼?.
12345678910.
我告訴你我有什麼資歷.
我在哪裡工作的?.
我是CEO,我有很多東西.
一路堆上去.
我介紹自己的時候.
一路說說說.
總之是世界上的人很欣賞的名號.
我都告訴你.
我是很有耐力的.
可以這樣介紹自己.
但《雅各書》的作者.
他怎樣介紹自己呢?.
其實他是一個很謙卑的雅各.
聖經裡有不只一位雅各.
這位雅各是誰呢?.

$^{81}$這位雅各是耶穌的兄弟.
是肉身的兄弟.
如果要說關係.
在神的家.
在神的國度裡說關係.
說我是耶穌肉身的兄弟.
這個厲害嗎?.
別人一聽就覺得很有份量.
我介紹的時候可能就這樣介紹.
我是他肉身的兄弟.
瑪利亞的兒子.
還有.
在地面上我是耶路撒冷教會的領袖.
一出來就跟你說我是領袖.
有誰不認識我呢?.
當時很有影響力的.
每個人一提我名字就知道.
你問吧,這裡沒有人不認識我的.
一說就這樣說.
他可以這樣介紹.
但你看到雅各都沒有這樣說.
還有他可以這樣說.
我見過復活的基督.
我可以跟你說.
耶穌基督有很多奇妙的事.
但見過復活的基督.
可以作為一個見證.
反而不是一個最大的問題.
但如果我那張.
我所擁有的.
覺得對我有利的.
能夠抬高我的位置的.
那些描述都說出來.
那沒有人不認識我了.
你也應該要認識我.
但雅各不是.
雅各一開波是怎樣說.
是說神和主耶穌基督的僕人.
僕人這裡令我們想起當時的奴隸.
是一個位置非常謙卑.

$^{121}$奴隸的位置很謙卑.
但在裡面不自卑.
為甚麼.
很榮耀的.
是神和主耶穌基督的僕人.
是很榮耀的.
卻並不驕傲.
是恰到好處的介紹.
雅各自己寫經文的時候.
他以生命.
以他的談吐告訴你.
我們為人可以這樣謙卑.
不妨大家看看.
就一個短短的位置.
想想今天的自己.
是一個怎樣的人.
你出來會怎樣介紹自己.
人們怎樣認識你.
我們會否每一個紅印堂的電影節目出來介紹.
我們是神和主耶穌基督的僕人.
我們是哪位哪位陳大文.
等等等等.
我只是舉一個例子.
這個問題是值得去思考的.
在這裡我們去了解一下.
雅各書的背景是怎樣的.
雅各書本身是要寫給一群.
分散在各地的猶太基督徒.
例如分散在耶路撒冷附近的地方.
這一群基督徒.
當然他們肯定是決志信主的一群人.
亦有機會曾經被雅各目擁過.
剛才我們說雅各是教會很重要的領袖.
很重要的領袖.
做很多目擁的工作.
可能這一群人曾經被他目擁.
但在這段時間他們因為信仰的緣故.
要面對迫不及待.
很多的艱難.
很多的挑戰,苦難.

$^{161}$甚至乎是四處的漂流,流蕩.
生活非常艱難.
可以這樣說.
就像後面的經文所說的.
落在各種的試煉裡面.
是各種的.
你數得出的.
身心靈的試煉裡面.
是很多.
而雅各就是要這樣.
真的去寫這本書.
要去鼓勵他們.
要去安慰他們.
究竟是怎樣的情況呢?.
成熟的第二個特徵.
是能夠以一個很超越的眼光.
去看困難的處境.
說真的.
我們每一個人都會面對困難.
或大或小.
其實這裡說的試煉.
不只是困難.
是很多艱難.
人生總有艱難.
中國人很喜歡說.
人生不如意事.
是十常八九.
是不如意事.
是十常八九.
如意事是一二左右.
我和你很自然.
常常都會有機會去面對.
去面對的時候.
既然大家都不能避開.
大家都有機會去經歷.
去見過.
不同之處.
成熟的生命特徵是什麼呢?.
我能夠以一個很超越的眼光.
去看這個處境.

$^{201}$當面對各種試煉的時候.
我們一般人很容易有的反應是什麼呢?.
大家想想.
我們很容易有什麼反應呢?.
埋怨.
灰心.
很多灰心.
為什麼我常常都遇到這些事呢?.
一而再再而三.
這裡說各種試煉是很多的.
試煉有些人想嘗試分.
會不會是從撒旦那邊來的試探呢?.
真的很多.
很想令我們都落在罪惡當中.
很想令我們徹底失敗.
又或者來自上帝那邊陣營的.
要考驗我們的生命.
那個堅韌度是多少.
究竟是試驗嗎?.
是試探嗎?.
我們不知道.
是很豐富的.
各種試煉是你想得到的.
都是包羅萬有的.
現在你想想我們生活.
試奉.
你能數得出來的.
你的健康.
經濟.
家庭.
家庭關係.
工作上的人際關係.
你的前途.
教會試奉各方面.
很多.
我們會面對各種試煉.
而當時這群信徒.
直接地面對迫伯.
流離失所.
我們面對試煉.

$^{241}$很容易灰心,沮喪.
我們懼怕.
我們埋怨.
埋怨人,埋怨神.
特別會失去對神的信心.
在平順的日子.
我們說信心是沒有意思的.
為什麼?.
當時很平安.
我很有信心.
但在平安的日子.
看不到生命的堅韌度在哪裡.
看不到成熟度在哪裡.
要看到成熟度.
就在試煉當中.
在困難,很大的困境裡.
才能看到.
所以平時一般人.
包括我們在內的反應.
很多是很負面的.
大部分都是很負面的.
瓦割要我們什麼?.
要我們面對各種試煉.
要認為是大喜樂.
要超越.
要超越它.
不要被它困住.
綁住你.
要覺得它是大喜樂.
不只是喜樂.
是很大的喜樂.
不是壞事.
試煉.
不要看得這麼負面.
不要灰心.
很快放棄.
埋怨這個,埋怨那個.
是好事.
大喜樂.
為什麼?.

$^{281}$待會兒會有解釋.
後面有解釋.
但我們能夠這樣看的時候.
真的是一個生命成熟度的指標.
不容易做到.
今天我講的堂導.
也是要鼓勵自己.
不要被困境困住自己.
在面對困難的時候.
要看到那種喜樂.
一個超越的眼光去看.
為什麼我能夠有超越的眼光呢?.
繼續吧.
成熟的特徵.
為什麼可以超越呢?.
信心要經過考驗.
成熟的特徵也是.
信心經得起考驗.
信心要經得起考驗.
雅各一路描述的時候.
怎樣經得起考驗呢?.
在種種的難處裡.
種種的困難,低谷裡.
你想想.
我們的信心要開始經歷考驗.
那個考驗可能不是一分鐘,兩分鐘.
一個小時,兩個小時.
不是一天,半天.
可能是以年計.
一年,兩年,十年,二十年.
一直下去,持續很久.
在過程中.
我們要堅持一件事.
信靠神.
信心不是對自己自我欺騙.
我就是要對自己有信心.
我要信有一個對象.
信靠我們的神.
就是那個拯救我們的耶穌基督.
那個成為我們的救主.

$^{321}$也是我們生命主人的耶穌基督.
祂超越萬有.
我要信祂.
這份信經不經得起考驗?.
成熟的人會經得起考驗.
你想知道自己是否成熟.
不要在現在平安的時候.
說自己很有信心.
而在患難裡.
夠膽讓別人看到自己.
有對神,信靠神的信心.
這樣就行了.
就是要這樣.
或者不是讓別人看.
重點是.
自己知道那份信心不被動搖.
這裡說到信心有多耐.
信心一直下.
一直堅持.
就會去到忍耐.
忍耐是什麼意思?.
是恆久.
很持久的信心.
不是一次性的信心.
是一直忍耐.
一天,兩天,三天,十年,二十年.
一直堅持.
由忍耐一直堅持到底.
這裡有個經文說.
又去到又完全.
又全被,又完整.
它說的是要去到底.
徹底.
今天我和你.
我們對上帝的信靠.
要堅持到什麼時候?.
問問大家.
這個考驗.
兩天而已.
所以我們堅持兩天就夠了.

$^{361}$之後就放鬆了.
我們要堅持什麼時候?.
對上帝的信靠.
直到永遠.
直到見主面.
至少直到永遠.
未見主面.
都不要放鬆.
堅持到底.
在《雅各書》裡.
最後說完全的字眼.
是在說成熟.
信心去到一個地步.
我們整個人成熟.
信心成熟.
原來生命的成熟.
是要這樣.
對上帝的信靠.
一直堅持.
橫行不打沉.
不拖低.
一直成熟.
甚至這份成熟.
是毫無瑕疵的地步.
是一份這樣的信.
我鼓勵大家.
在完結這堂道之後.
回去量度一下自己.
如果我們當100分.
就是我的信心徹底成熟.
今天給自己的評分是多少.
不要評別人.
我就知道.
那個人是沒有信心的.
我就知道.
前面後面誰就不及我這麼有信心.
我們不需要看這些.
看自己.
我去到哪裡了.
是不是.

$^{401}$你越多患難.
就越多時候看我去到哪裡.
我經常說.
不要在平安的日子.
說自己很有信心.
因為這是未經過考驗.
要去到這樣為止.
一個成熟的生命.
成熟的特徵.
另一樣東西就是.
懂得向神求智慧.
這個智慧在亞覺書這一段.
是說什麼呢.
是不是我有智慧.
去做不同的工作.
我現在上一份新工作.
我需要有智慧.
於是我求上帝.
其實是好的.
我們凡事求神.
但這個智慧有一點不同.
成熟的特徵是什麼呢.
當我面對試煉的時候.
面對艱難患難的時候.
我需要有智慧去明白.
其實我應該怎樣去看這個患難呢.
我應該怎樣去看這個患難呢.
上帝要透過這個患難.
要我學什麼功課呢.
上帝要怎樣透過這個患難.
要令我生命成長.
去成熟呢.
我需要明白.
很多時候患難出現的時候.
我們的第一反應就是.
為什麼輪到我.
接著下一次就對上帝懷怨.
下一次就罵神.
你有沒有想過.
其實我們需要有智慧去明白.

$^{441}$神啊.
在這件事裡.
你要我明白什麼.
怎樣去看.
不要這麼快.
按照我的人生經驗.
假定這件事發生.
一定是因為ABCD.
我用套餐一粒粒的.
就這樣解釋.
為什麼不問神呢.
直到問到為止.
不只是問.
那個智慧是什麼智慧.
在這麼多的患難裡.
在這些試煉裡.
怎樣走出去.
上帝下一步我怎樣做.
你帶我.
神一定會很清楚.
列一張清單給你.
你帶我走.
可以了.
上帝我知道你.
很了解我的人生.
知道我這一刻的經歷患難.
是為了什麼.
我求你帶領我.
讓我智慧明白.
讓我智慧走出去.
我需要這個智慧.
成熟的特徵不是.
我已經什麼都明白了.
什麼都知道了.
而是我懂得在這時候.
求神.
不是求人.
在患難的時候.
不是找十個人.
重複地說我人生悲慘的故事.

$^{481}$接著那十個人.
還要跟你說十個很不同的版本.
可能令我人生更加混亂.
但大家並不是惡意的.
大家都是很關心我的.
很想幫助我.
但那十個都是人的聲音.
成熟的.
真正成熟的基督徒會問上帝.
神你告訴我.
要等神.
求智慧向神求.
要等.
我們在這方面聽過.
不同的見證是很美好的.
有些人移民到海外.
或是在香港.
都跟我說過他們等候上帝.
很寶貴的見證分享.
真的等神.
要等.
要等一段時間.
直到神跟你說.
神因為有主權.
決定什麼時候.
怎樣回應我們.
但等到之後.
真的很不同.
令自己的信心大增.
因為是神的建議.
是神的帶領.
跟人的自以為的聲音.
是很大分別.
所以我想說.
我們要懂得向神求智慧.
心願我和你一起學這個功課.
等候上帝.
求問神.
當然神亦為我們預備弟兄姊妹.
預備好好的親友在身邊.

$^{521}$鼓勵我們.
與我們同行.
但我相信.
成熟的弟兄姊妹.
不是急於給他建議.
我和你說.
以我幾十年的人生經驗.
我告訴你怎樣走下去.
不是應該這樣.
而是你應該說.
弟兄姊妹我為你禱告.
和他同心禱告.
陪他走.
這是最好的.
心願我們不單是成熟的人.
做一個成熟的陪伴者.
是很重要的.
在這裡.
剛才說.
如果我欠缺智慧.
我不一定明白.
為何這件事發生.
或我怎樣走出去.
求主給我智慧.
求智慧的時候.
都要有信心.
等神的時候.
要有信心.
不要疑惑.
祈願土的時候.
神好像沒有回答我.
神不回答我.
我試過問不同的人.
你問了神沒有?.
你覺得神怎樣回答你?.
沒有啊.
神沒有跟我說過話.
神會怎樣跟我們說話.
我都不知道.
我想神不會回應我.

$^{561}$我試過.
不是的.
神一定會回應我們.
亞國教我們.
憑著信心.
不要求.
求神.
就要對他有信心.
不要疑惑.
不要懷疑.
神很不喜歡我們懷疑.
在這裡用了兩個例子.
說疑惑的人.
懷疑的人是怎樣.
可能這樣.
神不會回答我.
或是神不會幫我.
神沒有辦法.
這是更大的問題.
什麼叫神沒有辦法?.
神不懂電子.
神不懂這裡那裡.
現在是現代社會.
好像我感覺上.
聖經的神.
舊約很久之前的神.
新約的年代.
跟我們都很遠.
他不明白我們現代的事情.
他不知道.
這些簡直是.
可以這樣說.
何止為神呢?.
之所以為神.
就是色在,金在,永在的上帝.
無論什麼時間.
什麼科技.
什麼社會.
什麼環境.
上帝完全了解明白.

$^{601}$今天我們覺得.
那天我們同工一起走去.
逛商場的時候.
看手機.
覺得最新的.
哪個型號的手機.
我心想.
上帝根本看這些手機不上眼.
上帝用什麼手機?.
不用的.
上帝跟我們說話.
用什麼手機?.
上帝的東西才是高科技.
我和你都做不到.
怎麼會不明白呢?.
不單上帝不明白.
而且上帝是有能力的.
但他回答我們的時間.
不是這樣的.
我今天馬上.
一撥通電話.
就要他馬上回答我.
馬上按我的方法解決我的問題.
不是這樣的.
神有神的時間.
神的方法.
祂有主權.
我們要尊崇祂為上帝.
但要相信祂有力量.
一點有什麼疑惑呢?.
這裡有兩個例子.
第一件事.
他說疑惑的人就像什麼呢?.
像波浪一樣.
波濤洶湧.
你看看.
在翻滾.
求上帝.
但心裡好像忐忑不安.
覺得神不幫我.

$^{641}$又不回應我.
你想想那些波浪.
那種感覺很不穩定.
剛才讀經文說.
三心兩意的人是怎樣的呢?.
三心兩意的人.
可以得著什麼呢?.
可以從神得著什麼呢?.
可以得著什麼呢?.
一無所得.
什麼都得不著.
之後.
我們自己的生命.
陷入了一個很慘的漩渦.
死路一條.
這是非常可惜的.
上帝要我們不要疑惑.
你求得祂.
就要對祂有信心.
一心一意.
一心一意尋求祂.
一心一意相信祂.
一心一意等候祂.
神很想賜福給我們.
很想賜福給每一個尋求祂.
專心等候祂的人.
不要疑惑.
今天在這裡.
最後很想說一個見證.
我很欣賞.
這一位.
有人知道他叫什麼名字嗎?.
非洲之父.
我很欣賞的宣教士.
有很多宣教士.
我都很欣賞.
其實大家知不知道.
耶穌基督.
是一位很偉大的宣教士.
耶穌基督的宣教才厲害.

$^{681}$真的很厲害.
耶穌基督的宣教.
由神進入到人類的世界.
那個跨越是多大.
這個非洲之父.
叫做李文斯登.
我很喜歡.
很欣賞他.
本身他的父母.
是非常虔誠的基督徒.
他是一個蘇格蘭人.
他生在113年.
小時候家裡非常貧窮.
但貧窮對他生命的裝備.
我用這個字眼.
裝備是很重要.
因為貧窮.
所以他不可以像其他人一樣.
讀書時間就上學.
到中學的時候.
他日間就要打工.
晚上才可以讀書.
就算是之前的日子.
都是這樣.
很長的.
小時候已經出來上班.
但因為父母是虔誠的基督徒.
他小時候就學習聖經.
九歲.
你猜他可以背什麼經文呢?.
我們覺得.
我們這麼大都背不出經文.
現在笑了.
我們都很難背.
他九歲.
九歲可以讀到.
詩篇119篇.
小時候他很喜歡聖經.
父母的栽培很好.
小時候已經很熟悉.

$^{721}$這成為他一生的補複.
一生的力量.
你知道他中學畢業是多少歲呢?.
23歲.
這個年紀.
中學畢業.
因為沒有辦法.
家裡貧窮.
一直要上班.
後來很感恩.
神給他充滿了聰明智慧.
在這麼遲的時間.
一直讀書的時候.
神仍然都恩待他.
裝備他.
讓他考上大學.
在大學裡學了很多知識.
有醫術方面的.
有天文地理各方面的.
這些都很方便他將來.
去宣教.
他讀書的時候.
並沒有想太多.
只是讀書的時候.
有這些方面的裝備.
但他曾經想過.
會不會我學了醫.
我就可以透過我的醫學.
去宣教呢?.
起初.
他對我們中國很有負擔.
他有想過來中國宣教.
後來神對他.
有特別的帶領.
帶他去非洲.
他就去非洲宣教.
李文斯頓的特色是什麼呢?.
他是很體貼人.
但他很愛神.
有些人體貼人.

$^{761}$但卻不愛上帝.
愛人愛得過頭.
有些人真的很厲害.
捉住上帝的規條.
對人少了一份愛.
但他剛剛好.
他真的很適合去非洲宣教.
當時非洲很多土著.
世人來看.
很野蠻.
如果西方國家去.
可能很喜歡拿刀,槍.
嚇住人.
不要來搞我.
很想對著當地的土著.
但他不會.
他很明白人為何會這樣.
人因為有罪.
不認識神.
所以才會用這麼野蠻的方式待人.
所以他對土著很好.
有一次他透過一個方法.
他不是故意的.
但那件事之後.
土著真的知道.
他很愛他們.
為什麼呢?.
話說他去宣教的地方.
有獅子.
因為非洲獅子很厲害.
有一段時間.
他在當地成立了一個報道所.
當地有人誣告他們.
你違背我們祖宗的宗教.
竟然在這裡設立報道所.
他們多陰險.
他們會去打發獅子.
當地人很懂得打發猛獸.
他們會攻擊土著的牛羊.
獅子出沒.

$^{801}$不要說牛羊.
人也很危險.
當時李文斯敦決定.
跟一個不怕死的土著.
一起拿槍對付獅子.
結果他們雙方受了重傷.
獅子在他身上咬了11口.
特別影響他左手.
在這意外之後.
他終生左手不能舉重物.
不可以拿重物.
他那次流了很多血.
但透過這件事.
他是為了當地人.
為了保護當地人對付獅子.
而甘願自己受傷.
當地人知道他很愛他們.
對他們非常非常尊重.
是很不同.
如果說是試煉.
我們是否憑著一腔熱誠去服事.
就足夠呢?.
其實我們的生命也需要經歷考驗.
要經歷考驗.
如果我期望每天都是平順的日子.
我們的生命很難成長成熟.
但在考驗中要堅持到底.
他堅持什麼呢?.
其實當地的生活環境非常惡劣.
每隔一陣子他就被人誣告.
被人對付.
有獅子的問題.
當地生活的艱難等等.
他一直從小到大.
對於復活基督的信心.
在他內裡.
對基督的信.
信靠 就是這份信靠.
一直堅持在試煉中.
他不會停止.

$^{841}$可能對很多人來說.
有獅子的問題.
我不去了.
或者被人事奉得很好.
又被人誣告.
我會很灰心.
我為神而已.
為什麼要誣告我.
為什麼要迫我 傷害我.
又在我背後說我是非.
我不做了.
我放棄了.
很多人都是這樣.
有沒有留意.
教會是普遍性的教會.
是非真的會令人離開.
是非令人灰心 放棄.
困難也是.
我們常常灰心.
很多不信 對上帝的不信.
很懷疑.
沒有理由的.
我是愛你 服侍你.
為什麼不能順利一點.
但他一生的服侍.
就是被人看到很多艱難.
艱難是恆常的.
但他沒有想過放棄.
在他的事奉生涯裡.
他有幾次不同宣教的旅程.
每一次都有很好的貢獻.
其中有一段時間.
他找到販賣奴隸的道路.
是一個很大的貢獻.
對後來推翻商業活動.
造成一個很大的貢獻.
他找到那條路.
後面的人要推翻.
就容易很多了.
所以在過程裡.

$^{881}$他一生不是為自己而存活.
而是為上帝而擺上.
後來即使經歷了什麼.
他也沒有想過退休.
我不事奉了.
享受一下人生好不好.
我回去西方國家很舒服.
我回去蘇格蘭.
舒服一下 享受一下人生好不好.
他沒有這樣.
在1873年的時候.
他患上非常嚴重的痢疾.
也沒有回到蘇格蘭.
他就病死在非洲.
他被稱為非洲之父.
這個聲譽是名不虛傳.
別人問他.
李博出位嗎?.
有些人覺得.
你特意去這麼艱難的地方.
你知道有很多是非.
有時見你這樣事奉艱難.
你也想博出位.
想別人稱讚你.
為上帝拋頭顱 灑熱血.
不是的.
他問我這麼多年得到什麼.
清潔的良心.
如果說真的.
要榮譽的話.
我相信上帝會給他.
他渴望的從來不是人的清譽.
是生命的成熟與成長.
我相信經過這麼多年的風浪和試煉.
那份信心一定是經得起考驗.
你叫他出來講唐道.
生命成熟的特徵.
信心經得起考驗.
他一定可以講到.
他親身經歷過.

$^{921}$是為了服侍神堅持到底.
不是說些吹噓的話.
在這個世界只有一個李文斯登.
但我和你是屬於上帝.
我們可以信靠神.
今天神有給我和你的功課.
每個人都不同.
有給我和你的個人困難.
可能是患難.
各位姐妹今天我想問問大家.
想不想去迎接這些挑戰?.
給自己的信心經歷考驗.
給自己的生命成長.
夠膽有一天經歷所有事.
出來跟他們說.
在種種試煉中我們要認為是大喜樂.
心願我和你可以說得出.
我們一同禱告.
藉著亞哥的話.
成為我們很大的鼓勵.
神你知道.
經前到日一定會多經歷患難.
除非今天我們不經前.
天父靠你給我們經得起風浪.
信心經得起考驗.
在患難中.
我們的眼光要超越.
看到上帝給我們經歷的美意.
讓我們每一個生命都成熟.
讓我們每一個都謙卑.
都成為忠心的僕人.
有一天見到你的時候.
得到最寶貴的一件事.
是主耶穌基督稱讚我們.
有良善有忠心的僕人.
心願我們每一個都得到這個獎賞.
奉耶穌基督的名字祈禱.
阿們.
\newpage



\section{哥林多前書 4:1-6}
\label{sec:i_n9uN87S9M}
\textbf{2024.04.14 《克服身份危機》 哥林多前書4\_1-6 (和修版) 講員 : 曾裔貴牧師}
\newline
\newline
連結: \href{https://youtube.com/watch?v=i-n9uN87S9M}{\texttt{https://youtube.com/watch?v=i-n9uN87S9M}} ~~~~ 語音日期: 2024-04-14
\newline
\newline
\hyperref[sec:EYrQ_D9F5Wo]{\small{< < < PREV SERMON < < <}}
~
\hyperref[sec:index_chronic]{\small{[返順時目]}}
~
\hyperref[sec:index_scriptual]{\small{[返順卷目]}}
~
\hyperref[sec:jSl7w85PRnU]{\small{> > > NEXT SERMON > > >}}
\newline
\newline
哥林多前書 4:1-6
\newline
\begin{longtable}{cl}
\hline
\hline
章節 & 經文 (和合本修訂版)\\
\hline
4:1 & \begin{tabularx}{0.7\textwidth}{X} 人應該把我們看為基督的執事，為神的奧祕的管家。 \end{tabularx} \\ \\ \relax
4:2 & \begin{tabularx}{0.7\textwidth}{X} 所求於管家的，是要他忠心。 \end{tabularx} \\ \\ \relax
4:3 & \begin{tabularx}{0.7\textwidth}{X} 我被你們評斷，或被別人評斷，我都以為是極小的事；連我自己也不評斷自己。 \end{tabularx} \\ \\ \relax
4:4 & \begin{tabularx}{0.7\textwidth}{X} 雖然我不覺得自己有錯，卻也不能因此判為無罪；審斷我的是主。 \end{tabularx} \\ \\ \relax
4:5 & \begin{tabularx}{0.7\textwidth}{X} 所以，時候未到，在主來以前甚麼都不要評斷，他要照出暗中的隱情，揭發人的動機。那時，各人要從神那裡得著稱讚。 \end{tabularx} \\ \\ \relax
4:6 & \begin{tabularx}{0.7\textwidth}{X} 弟兄們，為你們的緣故，我拿這些事應用到我自己和亞波羅身上，讓你們從我們學到「不可過於聖經所記」這話的意思，免得你們自高自大，看重這個，看輕那個。 \end{tabularx} \\ \\
[1ex]
\hline
\hline
\end{longtable}
$^{1}$我們一起祈禱.
感恩再一次我們聚集.
這段聖經都相當清楚告訴我們.
在教會裡面.
我們要學校仕途保羅.
他對自己的身份.
神賦予他的人生的職份.
是恰如其分的.
這樣能夠去自己.
和他的同工阿波羅.
都需要去認定.
甚至乎可以成為教訓.
叫其他哥林多教會裡面的弟兄姊妹.
去學習怎樣去欣賞.
和對同工的遺訓認定.
我們很盼望主.
讓我們能夠一起在洪恩家.
或者在我們人生.
我們總會有身份的模糊的時候.
幫助我們能夠真的可以透過主的真理.
去認識我們自己.
在神眼中的我們.
願主大令.
感恩禱告.
奉耶穌基督的名.
阿們.
大家如果有留意到.
今天有少部分的姐妹.
穿得特別漂亮.
因為是泰國的新年.
我們的泰語姐妹.
其實由十二年前.
在我們中間一直下來的席堂.
沒有間斷地在我們當中聚會.
不過在室內不方便潑水.
有連續三天的.
如果待會散會.
你需要潑水.
待會你就問一問.
這個題目很重要.

$^{41}$每個人都可能會有身份危機.
我講一下我的姐妹.
姐妹的女兒.
大女兒.
從小就在加拿大出生.
加拿大的中國人.
嚴格來說是香港人.
從小開始上學.
到星期六.
她和弟弟.
必須在星期六的下午.
要去唐人街學中文.
小時候沒問題的.
就去吧.
媽媽載她去.
爸爸載她去.
到十二,三歲.
開始問一個問題.
我周圍的同學.
來自不同國籍的.
都是講英文.
純正的英文.
為什麼我要學中文.
又不懂.
又難寫.
又難讀.
開始抗拒.
還未能夠獨立之前.
一定是要求.
逼迫.
威嚇.
一定要去.
不去不行.
於是就在這些拉鋸.
完成了中學畢業.
她有一個心願.
讀書的時候.
很喜歡做航空公司的空中小姐.
到了高中.
她自己去應徵.

$^{81}$加行.
一次就成功.
為什麼呢.
原來加行說.
你會講廣東話.
你又這麼純正的英文.
又會法文.
所以她的優勢就是.
她可以講廣東話.
看就不懂.
寫更加不懂.
所以她開始對自己的身份危機.
慢慢過渡了.
我相信無論是美國.
或者是加拿大.
很多在那邊出生的小朋友.
都或多或少會有身份危機.
在2008年的時候.
有一個新聞都很有趣.
有一位中大碩士畢業的.
統計學的學生畢業.
他的GPA是3.6分.
相當高.
但蘋果日報就追訪這個人.
為什麼呢.
他叫梁建鵬.
因為他畢業之後.
只做了兩份工作.
每份工作都做不長一個月.
就被人解僱.
然後他努力每天去應徵.
在兩年裡面.
應徵了200份工作.
兩年200份.
差不多每天都要去應徵.
所以他每天出門就穿著同一套西裝.
只有一套.
因為他家裡沒有錢.
後來輾轉都不行.
每次見工都沒有人請他.

$^{121}$最終要怎樣呢.
一個碩士畢業生.
就要拿縱緩1830元.
當時2008年.
接著都不夠.
上網去控訴教育制度.
為什麼香港的教育制度.
訓練到我變成一個這樣的人.
200份工作都沒有人請我.
他說我都已經苦咒了.
我去麥當勞應徵了.
麥當勞的經理就說.
你不要玩了.
你是碩士來的.
你來做麥當勞.
他控訴.
後來蘋果日報就去追訪他.
就幫他.
找他.
帶他去見獵頭公司.
教他怎樣去應徵的技巧.
最終發現.
他的智力有問題.
更加不是他不承認自己是高分低能.
而是說.
原來他的社交問題.
他不和別人相處.
他陷入自己身份的危機裡.
當然現在已經2024年了.
梁先生現在情況怎樣呢.
就沒有人再跟進了.
我也不知道.
原來人生我們都會有不同階段的身份危機.
剛好我預備這個講章.
其實幾個星期已經有這個題目想講了.
上個星期.
樹人大學做了一個報告.
做了一個問卷調查.
訪問中學生.
高中和大專.

$^{161}$原來有八成的高中生.
是對自己的身份.
模糊和有危機.
不知道自己想怎樣和做什麼.
原來有六成的大專還沒有畢業.
都有身份模糊.
身份的危機.
好像還沒有找到自己的人生方向.
這個很快去到職場.
就要面對很多困難.
所以這個題目原來都頗切身.
頗靠近身.
剛才這段聖經.
其中第五節.
尤其是第六節就更加寶貴.
因為第六節保羅告訴我們.
所以我要將這件事轉給我和阿波羅.
第六節弟兄們.
我為你們的緣故.
這些事情就是上文所說.
在教會裡面.
弟兄姊妹千萬不要看牧者.
或者當時的仕途過高和過低.
這裡是延續一到第三章.
哥林多教會有分門結黨.
我是屬於彼得的.
我是屬於保羅的.
我是屬於耶穌基督的.
這樣的分黨.
分門結黨.
所以來到第四章.
保羅給予一個很重要的詮釋.
到底一個服事主的人.
應該如何恰如其分地看待自己.
和弟兄姊妹如何看待牧者.
這裡有很詳細的.
至少有兩大方面.
我們今天一起思想.
第六節我還未讀完.
這些情況.

$^{201}$轉給自己和阿波羅.
叫你們效法我們.
千萬不要過於聖經所記.
免得你們自高自大.
貴重輕看其中一個.
這樣.
不要比較.
每個牧者.
神安插他在不同的地方.
有不同的恩賜.
有不同的安排.
和服事.
是有他得定的目的.
所以我們必須一起學習.
第五節.
更值得我們思想.
所以時候未到.
什麼都不要論斷.
只等主來.
他要照出暗中的隱情.
顯明人心的意念.
暗中隱情.
沒有人知道的情況.
暗中的.
還有人心的意念.
在裡面.
心裡面的意念.
神在將來.
藉著耶穌基督的顯現.
要怎樣呢?.
是會顯明出來.
顯明意念.
顯明暗中的隱情.
我們要想.
各人在主面前.
受到的是什麼?.
一定是指責.
暗中.
誰人的暗中.
是完全沒有不好的意念.

$^{241}$但是留意保羅這裡.
他有沒有用錯了形容詞.
和那個情況呢?.
原來.
使徒保羅的內心.
這樣看十個救恩.
帶給我們的生命改變.
原來.
更加不應該.
很快就對別人下一個死刑的判斷.
看死你.
不應該.
因為連自己都不好看.
這樣看死自己.
因為這裡告訴我們.
有一個很寶貴的將來.
會這樣來看我們.
就是主耶穌怎樣看我們呢?.
各人要從神那裡得著的.
不是指責.
是稱讚.
你想想.
你心裡面的意念.
暗中的忍情.
是要被稱讚的.
這不是救恩生命的改變是什麼呢?.
一個人不信耶穌.
他的暗中忍情.
他心裡的意念.
沒有經過主耶穌的愛.
十個救恩的改變.
人的內心.
只會是非常詭詐恐怖.
怎麼可能將來會有稱讚呢?.
但是原來在主裡面.
我們看我們自己原來有一個將來的稱讚.
是這麼崇高這麼寶貴.
而且是這位主.
直接對我們的稱讚.
所以,弟兄姊妹.

$^{281}$一定要追求內心.
讓主來潔淨.
以至我們的生命力能夠服侍主.
這是第五節.
重要的是保羅對自己的身份.
有一個很肯定的認識.
以至他能夠在神.
在哥林多教會.
大家帶著不同的眼光去看牧者.
看阿波羅.
看基法.
看保羅的時候.
保羅很清楚自己有的位份.
職份.
以至他能夠用他的能力.
用神給他的恩典.
和恩賜.
來服侍哥林多教會.
相信如果我們有的能力.
不是用在神安排的人生.
有時會浪費了.
試著說一個例子.
最近有一套電影.
在內地.
香港不流行.
叫做.
一套電影.
高處除三害.
有沒有聽過這套電影.
小學已經讀過了.
周柱.
不好意思,周柱.
進朝的時候.
當然這套電影不是講進朝.
是講現代版的周柱.
這個典故.
這個故事告訴我們.
周柱在進朝的時候.
在一個地方.
他是一個很有神力.

$^{321}$很高大的一個人.
很有力量.
不過.
很多村民.
不太喜歡他.
告訴他.
我們這條村.
有兩個.
影響我們.
令我們很懼怕的.
在北邊有一個老虎.
在南邊有一個蛟龍.
在海裡.
這個周柱就說.
我有辦法.
去把他們搞定.
他真的走到山上.
打死老虎.
還有其中一個害.
就是蛟龍.
南邊的蛟龍.
周柱就跳到海裡.
跟他搏鬥.
三天三夜都沒有回來.
村民在那裡手掌慶祝.
他不知道還沒有回來.
回來之後.
發現村民不再慶祝.
又很害怕.
後來他知道.
原來除了兩害之後.
還有一害還沒有除.
原來就是他平日的作風.
平日對村民的威嚇.
令到村民對他懼怕.
比起老虎和蛟龍.
更加的懼怕.
這個周柱就有很大的醒悟.
原來他的出現.
不是村民開心.

$^{361}$而是懼怕.
最終他把自己.
完全徹底的改變.
似乎是真人真事.
在晉朝有記載.
世說新語的記載.
但是各位弟兄姊妹.
能力越大.
如果我們的心不是很認定.
在神裡面.
我們自己所獲得的身份.
我們所得到的主的憐憫和恩典.
來服侍的時候.
我們會不會對別人.
帶來一個影響呢?.
哥林多教會就是這樣.
幾位牧者的出現.
先後的出現.
保羅帶哥林多教會信主.
後來輾轉他離開.
一年半之後他離開.
去到歐洲其他地方馬其頓.
阿波羅曾經是阿波羅幫助過他.
他來到這個地方.
那些人開始有比較.
開始去看彼得,阿波羅,保羅.
誰比較強?.
誰能幫助我們多一些?.
這樣互相較量的時候.
就出現了一種不成熟的互相較量.
不是牧者之間較量.
是弟兄姊妹之間.
有這樣的問題.
出現在哥林多教會裡面.
就開始有一種分門結黨的情況.
在哥林多教會出現.
所以保羅在這裡.
來勉勵他自己.
透過這樣.
來向哥林多教會的弟兄姊妹.

$^{401}$來看清楚.
千萬不要效法我們.
過了聖經所記.
即是說.
很清楚第一節就告訴我們.
我們只不過是基督的僕人.
這個執事是跟第二節有關的.
是神奧秘時的執事的一個管家.
是受託付.
去將這個奧秘講明白給人知道.
解清楚,解明白.
給每一個的回眾知道.
這是我們一個尊貴的職份.
是神給教會,給保羅的.
使得保羅和他的同工.
是能夠將神的奧秘解釋清楚,明白.
所以只有一個要求.
對於牧者神安排的人生的職份.
一個要求.
所求於管家.
是要他有忠心.
對神的事工.
要有這種忠於所託.
很重要的內心.
例如我嘗試說說舊約的人物.
舊約神興起很多先知.
其中一個很特別的.
叫做以尼亞.
一出現他就跟阿哈和說.
如果我不祈禱就不下雨了.
果然開始旱災.
接著就饑荒.
三年六個月沒有下雨.
輾轉那些溪水就沒有了.
他就要去到別的地方.
接著差不多時間他就祈禱了.
真的傾盆大雨.
但是當時阿哈王的太太耶洗別.
這位皇后.
拜引渡城的國家.

$^{441}$拜偶像,拜巴黎,拜亞瑟拉.
拜得很厲害.
所以耶洗別就要追殺.
這個以尼亞先知.
這個以尼亞先知突然之間.
發現自己所做的好像一無是處.
國民都不幫他.
他就不斷地走.
由北面走到南面的曠野.
到了一個地步.
仗著飲食的力量.
四十晝夜走了.
到了一個羅藤樹下.
他要來求死.
接著有天使就拍他.
為什麼你要這樣做?.
你做什麼?.
他說了兩次.
我為神大發熱心.
整個以色列拜偶像.
現在只有我為神大發熱心.
連續說了兩次.
各位,似乎是以尼亞在他.
做了那麼多為神清算的假先知.
在加密山上.
八百五十個先知.
都求雨,求不到.
他一祈禱,天就降雨.
燒光那些祭物.
降火下來,燒光那些祭物.
很明顯他的祈禱是神閱立的.
那些假先知祈禱.
用很多方式都不行.
但是事與願違的就是.
夜洗別無悔改.
阿蝦也不理他.
甚至下一個追殺令.
要來抓這個以尼亞.
他就害怕.
他就開始感懷身世.

$^{481}$身份出現了危機.
他就逃走.
因為只有我一個先知為神大發熱心.
只有我這個人大發熱心.
其他人都不行.
是我仍然為神的真理.
在以色列國站著.
是不是想多了.
是不是開始思覺失調呢?.
以尼亞.
開始有一個問題.
開始為自己的內心.
那種自我中心來辯護.
神好像不理他.
繼續進入一個山洞裡.
然後神向他顯現.
不過很特別的.
很微小的聲音呼喚以尼亞.
以尼亞你在哪裡?.
然後有些山的地震等等.
神不在那裡.
後來神用微聲跟以尼亞說.
然後神跟以尼亞說.
現在你回去.
回去之後你要告立幾個人.
來接觸神的工作.
首先就是告立這個靈士的孫.
這個耶護.
然後還要告立一個叫做.
這個以尼沙.
做先知接觸你.
即是說.
原來神不單止沒有收回.
以尼亞先知的那個位份.
還給他一個很大的尊榮.
他還知道這個國家.
在他不在之後.
這個國家繼續存在.
繼續有神先知的密使.
還有那個教導.

$^{521}$還有那個工作.
即是說.
以尼亞其實你不要搞錯了.
我對你的事奉是肯定的.
是你自己面對一個身份危機.
所以現在給回這個.
告立的這個職份.
回到上去.
親自去告立幾個王.
還有告立這個靈士.
這個以尼沙的時候.
對不起.
沙伐的兒子以尼沙.
來到這樣去接觸的時候.
以尼亞其實一直回歸.
回歸他自己的身份.
回歸他自己.
原來在神的眼中.
他是完成了神對他的託付.
所求於管家的.
是要有忠心.
親愛的弟兄姊妹.
你找到你自己在教會.
在社會.
在人生.
在洪恩家的那個位份了沒有?.
你找到你自己能夠為神.
可以大發熱心的那個職份了沒有?.
你已經領受了神給你的恩賜了沒有?.
所求於管家.
只需要有忠心.
人應當看我們是基督的執事.
是神奧庇時的管家.
所以弟兄姊妹.
你一樣有你那個管家的.
守住某一點.
某些的奧庇.
去幫你的朋友.
幫你的家人.
帶他們認識主耶穌.

$^{561}$這個永恆救恩.
用你的脈絡.
用你自己的context(場面).
去幫你的家人.
幫你的朋友.
教會的牧者.
就更加有這個位份.
要解明福音的奧庇.
要將這個執事.
完整地展現在神的弟兄姊妹當中.
所以這段聖經真的很寶貴.
告訴我們.
我們需要去尋索.
自己在神面前的身份.
每個人都會有危機.
身份的危機.
就是過多看自己.
過低看自己.
都會有不正常的心理問題.
但是這是短暫的.
神會照出其中的忍情.
顯明你和我將來心裡的意念.
記住.
救恩會使我們將來得到稱讚.
不是責備.
所以如果我們不追求.
不去尋索自己的身份.
你繼續讓你的忍情.
讓你的心裡的意念.
是骯髒的.
將來就是指責.
所以保羅真的很寶貴.
他看同宮是從欣賞的角度.
不是從一種責備.
挑剔.
判斷.
批判的角度.
所以我和阿波羅這樣說.
我們是神中心的僕人.
所以你們看我們.

$^{601}$一定要俠如其分.
盼望這段聖經成為我們弟兄姊妹.
一起學習.
我們要尋索神在我們生命裡的照明.
那個身份.
使我們大家一起.
不需要比較.
大家盡我們人生的力量去服侍主.
記得周處未悔改之前.
他是帶來很多人的影響.
但是當他真的明白自己.
原來他的力量是用錯方向的時候.
他就帶來很多對人的傷害.
但是當他改變.
成為了很多人的祝福.
盼望我們的心讓神去管理.
我們一起感恩禱告.
多謝主今早一個很簡短的訊息.
但是很重要.
因為我們真的需要面對神.
在我們內心的審查.
求主幫助我們.
以致我們的位份.
照明.
是在主的面前獲得.
以致我們能夠被肯定.
不需要別人對我們的肯定.
而是主對我們的肯定.
使我們能夠忠於所託.
在主給我們的恩賜.
我們的照明.
一起去服侍.
使主在洪恩家.
在洪水橋.
在香港得到榮耀.
我們這樣感恩禱告.
奉耶穌基督的名.
阿們.
\newpage



\section{歷代志上彼得前書 16:23-24}
\label{sec:jSl7w85PRnU}
\textbf{2024.04.21 《宣教由 「信」做起》 歷代志上16\_23-24;彼得前書1\_17節 (和修版) 講員 : 高淑芬姑娘}
\newline
\newline
連結: \href{https://youtube.com/watch?v=jSl7w85PRnU}{\texttt{https://youtube.com/watch?v=jSl7w85PRnU}} ~~~~ 語音日期: 2024-04-23
\newline
\newline
\hyperref[sec:i_n9uN87S9M]{\small{< < < PREV SERMON < < <}}
~
\hyperref[sec:index_chronic]{\small{[返順時目]}}
~
\hyperref[sec:index_scriptual]{\small{[返順卷目]}}
~
\hyperref[sec:bW5i5_FWoTI]{\small{> > > NEXT SERMON > > >}}
\newline
\newline
$^{1}$Sawasdee ka.
剛才我們泰語詩歌都跟大家說了.
Sawasdee ka.
就是很歡迎大家的意思.
最後我們就說.
Khob khun ka.
知道是甚麼嗎?.
多謝你們.
Sawasdee ka是怎麼寫的呢?.
現在你看到的泰文就是Sawasdee ka.
你和我都一樣.
在我早期還沒有學泰文的時候.
我看到很多蝌蚪.
還有一個燙豆.
你見不見到.
在地上有一個燙豆.
泰國文還分了有蒸氣燙豆.
燙豆還會有些東西在.
我其實也不太相信.
我竟然會學到泰文.
因為當時學泰文.
真的學到哭了.
但是上帝猜我去泰國宣教.
所以我都繼續憑著上帝的恩典.
繼續來走.
我不知道按哪個按鈕才上.
不如你幫我控制.
下一幅.
是不是下箭咀.
但是沒有反應.
是一個信字.
大家都知道.
我們中國有一句話.
官字兩個口就好做事.
宣教一個口開始.
這個信很重要.
剛才我們一起讀《仕途信經》.
你會發現很多個口.
每個字都有個口.
這個信無論是宣教.

$^{41}$或者教會有很多關顧.
我們的報道.
遲些我們的咖啡室.
其實我們都要由信做開始.
這個信是怎樣的呢.
首先我們很想看看.
雖然在聖經裡面.
我們看不到宣教這個字.
但是宣教的意識.
在整本聖經裡面是有的.
甚至乎在舊約裡面.
我們翻譯為mission.
sending這個字.
都是由古音.
希伯來話也好.
新約的希臘文.
裡面都是有講.
雖然沒有直譯叫做猜.
或者是宣教.
但是就好像在歷代至上.
這裡十六章23到24節.
請大家再讀一次好嗎.
(全地都要響耀國歌).
(以祈全人他的救恩).
(在列邦中述說他的榮耀).
(在萬民中述說他的旗幟).
在這裡你都看到全地.
然後是萬民.
列邦所有的幫派.
雖然當時在舊約裡面.
以色列文就是以色列文.
然後有米甲人.
或者很多族的不同的人.
亞東人等等.
但是上帝的訊息.
這句歷代至上.
其實是這段裡面.
是在講述大衛王的經歷.
在前面剛才是十六章.
是他所寫的詩篇之一.

$^{81}$在歷代至上.
前面發生了甚麼事呢.
那件事就是這樣的.
因為約櫃終於非常艱辛地.
終於能夠接到耶路撒冷.
耶路撒冷就是他們很重要的一個城都.
他們被譽為聖城.
他們可以把耶和華的約櫃.
來到他們自己的地方.
回到自己的家裡.
甚至這個聖城才是名符其實.
因為連約櫃都在其中.
大衛當時已經是一位國王.
他手舞足蹈地.
真的很開心.
看得很帥.
你還記得他哪位老婆在那裡紮馬呢.
米鴿.
他的太太就在那裡說.
你真是丟臉.
竟然身為皇帝.
竟然是這樣的.
沒有儀態.
但無論如何.
大衛王真的歡喜快樂地.
歡迎上帝的約櫃來到耶路撒冷這個地方.
甚至用很大的儀式.
來護送這個約櫃.
我們都知道.
在聖經裡面.
大衛是稱為合神心意的人.
大衛非常體寵神和神的同在.
弟兄姊妹要留心.
體寵神和體寵神的同在是兩回事.
兩回事我們都要體寵.
我們很多人.
我們有種牧師林群體.
牧師林群體去體寵神.
你想想他們一日五次去敬拜神.
如果你坐輕鐵.

$^{121}$我突然忘記去哪個站.
有一個站還未去到元朗.
你會看到中東的店鋪.
那裡平時.
如果不是他們敬拜的時間.
你會吃到中東的薄餅.
他們的咖喱餐.
如果你是星期五12點.
你坐車又經過.
我先查一下那個站.
我忘記那個站的名字.
你們知道就好了.
如果你們有機會.
星期五12點左右.
因為星期五的回教旨是最齊人的.
那個時刻他們覺得.
最重要一定要來敬拜神.
就像我們今天說的.
主日我們是要來敬拜神.
很可惜.
2019,2020發生了甚麼事.
新冠疫症蔓延了世界.
新冠疫症還未完.
到今時今日還有.
不過全世界的國界已經開放.
於是就算香港本地都開放.
你不用以馬上進去吃飯.
又可以出去其他國家.
於是香港就出現了很多的報復式.
報復式的去吃.
報復式的去遊玩.
但很奇怪.
教會從未出現過報復式崇拜.
不單止我們香港.
不單止我們洪水橋.
全世界都是這樣.
有些教會.
因為弟兄姊妹.
依靠Zoom來敬拜.
Zoom的弟兄姊妹.

$^{161}$我不是說你.
我都知道你有些苦衷.
來不了.
不過問題是.
當你有機會回來的時候.
現場的崇拜是很不同的.
當中有神與我們每一個人同在.
所以認定神是很重要.
亦都看重神很重要.
同樣看重神與我們的同在.
同樣是很重要的.
神的同在有四個字代表.
是甚麼.
你們真的不知道.
你們不要這麼謙卑.
合乎中道就行了.
「以馬來利」就是神同在.
在新約裡面.
主耶穌基督降生來到這個世界上.
要為祂起名.
叫做「以馬來利」.
意思是.
神是上帝應許與我們同在.
昔日主耶穌親身來到現場.
來到這個世界上的現場.
成為一個宣教士.
祂就是我們的榜樣.
從天家.
我們現在多渴望去天家.
天家不會有眼淚.
不會有近視.
不會有頭痛.
不會有腳痛.
腰骨痛等等.
不會有衰老.
不知道神會定型我們.
在十多歲還是二十多歲.
怎樣也好.
是沒有衰老的.
或者一百歲也好.

$^{201}$但每個人都很健康.
怎樣也好.
神的同在.
當主耶穌基督復活後.
被接生到天空的時候.
上帝同樣應許.
聖靈與我們每一個同在.
弟兄姊妹.
今天你有沒有歡迎.
聖靈入住你的心裡?.
上帝創造這個世界.
由一開始已經給人自由意志.
所以我們剛才試圖.
在信經裡.
我們信三位一體的神.
當中有上帝是我們的父.
有主耶穌基督成為我們的朋友.
也有聖靈是我們的保衛師.
與我們一起的.
今天你有沒有邀請到.
聖靈在你的心靈裡面?.
所以大衛已經懂得體重神.
以及體重神與他同在.
所以為何約季去耶路撒冷.
是一件如此重要的事.
令大衛蒙形.
歡天喜地讚美神.
他全心全意敬拜神.
讚頌神.
並且在當時歌頌神.
我們在上星期.
因為是泰國新年.
我們有載歌載舞.
但我們沒有給你看.
有興趣可以問泰國姐妹.
看她們願不願意給你看片段.
我們同樣有祝福.
這個在潑水節.
如果不特別告訴你.
可能你會以為是九龍城那種潑水.

$^{241}$或是在泰國電視新聞潑水.
其實最重要的意義.
是互相的祝福.
我們泰式的祝福.
就是有些水,花來祝福.
上帝的祝福同樣是很重要的.
你想想,未信的人.
一年一度的祝福都很重要.
農曆新年我們都很重要的來祝福.
上帝的祝福是天天祝福給我們的.
你信不信?.
不信,所以什麼都要從信開始.
我們亦看到.
剛才說到這段經文前面的背景.
當十六章整個儀式.
好像在講述給我們今天看聖經的人.
中人把神的約櫃放進去.
但當時未有聖殿.
只是放進了帳幕.
或者我們可以稱為「會幕」.
因為帳幕是特別敬拜神的.
大衛在當時獻上梵祭,平安祭.
來奉耶和華的名.
來祝福所有的以色列人.
在這個梵祭是完全的擺上.
平安祭亦是一個完全.
所有的都是一個完全.
我今天很感恩.
詩歌團隊你們所選的歌.
或者崇拜聚會的詩篇的宣講宣召裡面.
其實若果你看到.
其實都是在說我們要向萬民傳遞.
來宣講神的榮美.
祂行的神蹟騎士.
亦都在獻唱裡.
一切獻在壇上的時候.
其中領師的姐妹說到.
羅馬書將身體獻上當作活祭.
是我們每一個人.
昔日我們都不太懂甚麼是梵祭,平安祭.

$^{281}$總之你很簡單.
記得是完全的擺上.
是完全的感恩.
今天你是一個祭.
你這個祭是否完全的呢?.
還是你的軀殼在壇上.
你的心靈已經離開了呢?.
不知道.
但是我們懇求神幫助我們.
讓我們既然是相信神的時候.
求主給我們同樣獻情給上帝的時候.
這個活祭.
我們每一個人的活祭.
都是完完全全的獻情給我們的上帝.
在歷代至上.
講述這件事.
然後到這裡.
大衛在當時.
來讚美上帝的說話.
全地都要向耶和華歌唱.
雖然他在他們以色列人的儀式當中.
祝福以色列的子民.
但是如果你記得很久很久在創世記.
有一群外邦人.
其實在約書亞.
早期我講約書亞時代.
他們要每一個的迦南美地.
每一個都要征服.
因為是神賜給他們的迦南美地.
其中有一個群族.
他裝模作樣.
準備了薄餅.
大餅也好.
是發霉的.
走去以色列人那裡.
我們很遠來的.
我們很想投靠你的神.
誰不知其實是很近的一個大族.
欲知詳情請你看約書亞記.
因為我今天只有半小時的時間.

$^{321}$所以你會看到.
其實這個以色列民當中.
其實早就已經參集了一些外邦人.
這群外邦人雖然他們有國王.
有很高的領袖.
他們都願意來服侍這群以色列人.
為了什麼?.
他們認定他們所信的上帝是真實.
所以他們不惜輸諸降貴.
都來投靠上帝.
甘願為這群神子民的子民.
去做勞瀑.
來服侍他們.
以前為什麼大衛可以.
有這樣的感受去寫了一篇這麼好的詩歌呢?.
因為大家都知道.
自從薩姆爾高立他為王的時候.
當時還有一個王還在生.
叫做什麼王?.
蘇羅王.
是蘇羅.
我怕你們答的不對.
大家答的很對.
是我沒有信心給你們.
當他去做皇帝的時候.
其實不是那麼容易.
繼位登位就登基.
不是的.
是真的蘇羅當時都要追殺他.
都要打了很多場仗.
甚至有一次蘇羅.
有一支箭差一點.
神保守.
他經歷過很多不容易的事情.
在當中體驗到上帝是真實的.
現在他成為以色列王.
他深知一切都是上帝的帶領和祝福.
以至到他今天可以成為以色列王.
所以他急不及待地.
很想把這個藥櫃帶到耶路撒冷.

$^{361}$今天來到.
他就在讚美.
其實我們詩篇也好.
剛才那段經文也好.
都是他發自內心向上帝的讚美.
他呼籲以色列人一定要歌頌神.
很奇妙的作為.
要尋求耶和華的面.
今天我們同樣.
你和我.
今天不知道有沒有旺雨.
我沒有看新聞.
但很大很大雨.
面對接近十一點半的弟兄姊妹.
我看到你們身體都濕了.
希望你們不要著涼.
不要著涼.
求主保護.
你們為何風雨之下都願意來呢?.
因為一齊敬拜.
現在我們還可以有這麼自由的敬拜時間.
我們真的要好好珍惜.
一齊敬拜上帝.
尋求上帝的面是好得無比的事.
我們信主到了今天.
大家經歷過.
意指到你是很願意.
星期日來敬拜上帝.
所以我們都要一起彼此的勉勵.
我們一起歌頌讚美上帝奇妙的作為.
時時刻刻都要尋求耶和華的面.
尋求耶和華的人.
應該有些特徵.
你覺得那些特徵是甚麼?.
你不用告訴我.
你做給我看.
你應該是很歡悅的面容.
心裡歡喜.
要知道你的面容也會歡喜.
我讀神學的時候.

$^{401}$最後的論文寫了.
我發現你那一刻.
我們僅有一天24小時.
僅有那些時刻.
我們真的笑不下.
真的笑不下.
那一刻更加要回到神的面前.
因為我們何時沒有這份喜樂的心.
原來我們信不過神.
信不過神即是得罪了神.
你明不明.
你今天跟大佬.
你不信大佬.
你怎樣威風.
你是否得罪了大佬.
何況我們得罪了天上的父神.
所以我們的情緒幫助我們.
我們有憂慮.
我們回到神的面前.
你就會有這份力量.
所以第一個口是信.
第二個口是回.
回到神的面前.
得取這份力量.
我們求主給我們.
時時刻刻懷著感恩的心.
不斷向神呼求.
有很多時候我們會靜下來祈禱.
有很多時候我們一直做我們這個人.
一直去求神.
你時時刻刻都可以祈禱.
不用有些危難.
你等一會.
當然你可以有這樣的時刻.
但很多時候.
你在走路的時候.
求神給我有平安.
或者你有些事很煩的時候.
都可以在這裡祈禱.
我很想告訴你.

$^{441}$我們可以時時刻刻.
尋求神的面.
你會看到.
如果是說十六章.
這個約季去到聖殿的時候.
嚴格來說.
這篇詩篇.
大衛去唱那首讚美歌.
不是我這麼短.
我只有23,24.
他再多十節.
再多一點.
如果你慢慢回去.
拿回聖經.
打開列王記上十六章.
這裡去看的時候.
你看到23到34節.
是不斷重複.
傳遞列邦,萬民.
民眾的萬族.
我們香港人已經很多族.
有南亞裔.
有些泰裔的姐妹.
已經是香港的市民.
有印尼族.
有菲律賓.
有很多.
最近我才知道.
我開始時不是很有知識.
我以為香港沒有難民.
自從越南難民走了之後.
原來香港還有難民.
因為去探望一位泰國朋友.
一位女士.
她已經和香港人結婚.
生了三個孩子.
但她仍然是難民身份.
原來她每個星期.
還要去屯門移民局簽到.
在她的口中說.

$^{481}$原來真是萬族萬民.
很多不同國家.
非洲有什麼.
原來我們香港.
只是難民.
也有很多不同的群體.
正式成為我們.
所以為什麼是萬民中.
一開始我以為萬民就行了.
為什麼還要是民族.
原來這個民.
就像香港為例.
是香港的.
不過裡面有很多不同的族裔.
天,地,海,田中所有的樹木.
甚至上帝所創造的.
都要去讚美神.
不知道你有沒有經歷過.
當我們說安靜在神面前的時候.
除了唱詩歌,讀聖經,祈禱之外.
我們去看一看上帝所創造的.
特別是花草樹木,天.
每天看天都很不同.
雲,不要說顏色.
有黑色,有紅色.
黃昏的時候紅色.
不同的藍色.
或者今天黑色.
這樣也好.
你也會看到不同的變化.
其實你安靜在窗外看一看的時候.
或者對著盆栽也好.
或者你養的金魚也好.
都可以做這個靜觀.
你看看上帝所創造的.
多複雜.
所以聖經裡面有說.
連花的一朵.
多謝大家.
我也為了太魚姐妹.

$^{521}$多謝大家今天讚美我們.
其實花的一朵.
比我們今天穿的衣服更美.
這些是上帝所創造的.
可以在這一刻.
連昔日大衛王也要呼籲.
要一起去讚美人.
讚美神.
述說祂的作為.
述說上帝的作為.
就是說好消息.
說好消息就是傳福音.
我們現在說的都是上帝創造天地.
主耶穌到新約.
主耶穌來洗濁我們的罪.
所以我們求主幫助我們.
讓我們時刻.
宣教不再是宣教時才有的.
特別是今天你進了宣道會.
今天要考慮一下.
有沒有進錯會.
宣道會就擺明了.
信徒和宣教士一起去宣教.
一起不是叫報道會.
也不是叫宣教會.
你要去問ABC Simpson.
孫信博士.
他整個人就是宣講.
上帝的榮美.
你今天坐在這裡.
在座的年長的.
是否是四十九歲的那位.
即是九十四歲的那位.
應該是最年長的.
你知不知道我們這麼年長的.
雖然他自稱四十九歲.
但他九十四歲.
他這麼年長的子會也好.
他坐在這裡已經是活活的見證.
你當天不上班.

$^{561}$你會否覺得慚愧呢?.
Zoom的弟兄姊妹.
如果你今天是為了下雨而不上班.
你應該向我們九十四歲的姊妹.
鞠躬,真是慚愧.
這麼大雨,他也願意來.
濕透了.
搭著車也濕透了.
也來了.
所以你坐在這裡敬拜神.
這就是一個活活的敬拜.
就是一個見證.
看到這裡著重全地.
列邦萬民.
當上帝呼召我的時候.
我以為要去泰國.
去泰國就傳福音給泰國人.
那些是佛教徒.
誰知在我最後的宣教.
二十年最後的三年裡.
三,四年裡.
上帝讓我認識了一群穆斯林群體.
或者簡稱為回教徒.
你會看到.
因為我後期認識他們.
我沒有這方面特別的裝備.
我上學的時候.
他們會覺得男和男.
女和女.
但你看看這裡.
為了主要在男的當中.
他們全部是長老.
他們的婦女準備了一些東西.
我在泰國裡面.
上帝讓我開了門.
讓我接觸那些穆斯林.
我真的對著那些長老.
今年七月份.
我有機會回到泰國.
我們會去泥府一個地方.

$^{601}$去到泰國之後.
還會再搭內陸機去到泥府那邊.
因為泥府.
我們會去清康縣這個地方.
不是清內,不是清邁.
那裡已經有很多宣教士.
我們會去清康縣的地方裡面.
來做很多小朋友.
因為那些小朋友的家長.
父母會輪流去坐牢.
為甚麼要坐牢?.
不是殺人放火.
而是吸毒.
這樣的情況.
所以當時上帝.
我的宣教觀是很闊的.
所以上帝開了這個門.
我就做了.
也是由信.
信得過上帝.
把我送到泰國.
我現在在泰國裡面遇到的事.
信得過上帝.
仍然和我同在.
其實當我們去宣教.
傳福音.
關顧別人也好.
其實這些事.
都是我們每一個.
上帝頒佈給我們.
由你一信.
上帝已經告訴你.
是一個大使命.
今天我就不再說大使命.
總之.
信仰和生活是連結在一起的.
他們是拍拖的.
你不可以丟棄一個.
他們是一起的.
所以你的生命.

$^{641}$就是.
你的生命.
就是活著.
是生活著.
這個就是使命.
你要過得好.
我們的信徒.
以前我們比較.
我早期.
很早期.
的時候.
我的觀念也是.
我雖然不至於.
好像菲律賓每年復活節.
壽難節.
他們又釘自己.
再一次十字架.
那一種的受苦.
但我覺得.
自己捱得就捱吧.
拼命地服侍神.
結果拼出燒身.
耗盡.
但是.
記得我和你說道.
我說過道.
你要愛自己.
神給我們的誡命.
但要愛神.
其次要愛自己.
然後愛其他人.
我不看了.
不知道有沒有得看.
到時去看吧.
總之.
宣教其實是一種生活的方式.
當我去到泰國.
其實有很多事發生.
不過今天先不說了.
總之你的生活.

$^{681}$就是一個很有威力的見證.
當時和這個長老比較.
因為他家算是較近教會.
我由教會去.
或由泰國去教會.
不,是去他的家.
都較近.
近之中.
都在我這裡.
其實開放了.
自從新官開放之後.
我都有報復式.
報復式回泰國.
因為我自己沒有車.
沒有車我就不能開到他家.
這麼遠.
連的士都不肯載我.
因為太遠了.
之後他沒有車回頭.
在我們巴吞塔尼的地方.
如果你在曼谷.
你不會明白.
什麼Uber.
叫的士.
在我們那裡沒有.
連叫的士都可以不來.
所以我就沒有辦法去那裡.
但在未來的七月.
就有機會去他們那裡.
因為我們是包車去他們那裡.
這個家庭因為比較近.
所以和他們接觸得多.
想不到有一次我回來.
那次是我最後回來香港.
他整家人.
你看不看到有多少個子女.
現在你看到抱著的那個.
因為黑漆漆的.
所以你看不到.
這裡已經有六個了.

$^{721}$到後期再聯絡.
原來已經有七個了.
最後終於是兒子.
抱著的那個都是女.
就是這樣.
所以繼續提醒大家.
一個口真的要相信.
要去相信.
要去述說.
要向著上帝唱讚美歌.
不單是向上帝.
是向萬族萬民.
就好像今天這樣.
我們當然是一路祈禱.
希望今天步行籌款.
好天氣.
在這裡想.
如果真的.
因為這幾天都會下雨.
新聞報告這樣說.
我昨天就想.
就算下雨.
就風雨同路.
是不是.
抱著雨傘.
可能很多年來的.
我知道籌款.
步行籌款已經很多年.
可能今年你會很特別.
抱著雨傘.
大家濕了身.
不要緊.
回去一樣是洗澡.
無論是汗水.
我們都要洗.
何況是雨水.
一樣洗.
所以.
信得過上帝的大靈.
就好像我媽媽.

$^{761}$早幾個月前她入院.
因為一入院.
她的狀態突然很低很低.
連醫生都說.
要不要急救.
急救會傷脈骨.
你做不做.
叫我們不要做.
好吧不做.
又要氧氣.
總之.
當時我就在想.
如果上帝是接她的.
就等她不要那麼辛苦的時候.
接她回天家.
我媽媽一開始探望她的時候.
就看到很多漂亮的花.
有人穿白色衣服.
是不是耶穌.
她又沒有說.
頭兩天都是這樣.
很漂亮.
很漂亮的花.
如果上帝接.
就趁她那麼開心.
就接她.
如果留.
如果上帝讓她繼續留.
就留.
我們人做什麼.
她來到我們.
我們只有兩個小時探訪.
就很忙.
幫她做這個做那個.
總之.
無論我們遇到什麼事情.
我們相信上帝是和我們一起的.
我們都凡事交給祂.
剛才我不是說到回.
我們凡事都要回到上帝的根前.

$^{801}$而且上帝在《仕途信經》也有說.
我們不單回到上帝的面前.
Zoom的朋友.
我知道你們是回到上帝的面前.
但是不要忘記.
上帝也要叫你們.
回到你的群體當中.
屬靈的家.
我不是很好見證.
媽媽進了醫院.
一開始我沒有告訴弟兄姊妹.
有說.
不過沒有呼籲.
沒有告訴你多少樓.
多少醫院.
因為心裡想.
不要煩著大家.
而且那兩個半小時只有兩個人.
我和我妹妹都非常忙.
幫她擦身.
幫她轉身.
不要搞砸弟兄姊妹.
但是有肢體就真的問.
回答她.
我真的來嗎?.
這樣很不同.
真的很不同.
就是和媽媽.
叫媽媽.
媽媽我來了.
她回應一下.
又昏迷了.
但是我的弟兄姊妹來.
跟她一說.
起碼說了三四句.
她才昏迷.
你說是不是真的.
回到神的面前之餘.
你將你的困難.
因為聖經應許.

$^{841}$不是你一個人祈禱.
當然你一個人祈禱很重要.
聖經應許.
兩三個人禱告.
這個上帝實在的聽.
是聽得更實在一點.
所以為什麼教會說.
你一定要回團契.
回團契不是綁著你.
是讓你認識一些肢體.
和他有團契的生活.
有什麼可以說.
這樣來到一起.
回到上帝的面前.
接著三個口.
你厲害了吧.
我姓高也是兩個口.
我感受到上帝不只是.
要我們一個口信.
要我們回到上帝.
還要我們的品.
我們的品是怎樣.
我們的Being是怎樣.
我們的行為是怎樣.
以前我都很天真.
那個人說是基督徒.
我就很相信.
我爸爸更天真.
我爸爸相信了之後.
我爸爸開的士.
我爸爸喜歡聊天.
我是基督徒.
他就不收人家的錢.
本來我也不知道.
剛好有一次有個姐妹.
坐過他的車.
剛好我爸爸說.
他認得我爸爸.
那天沒有收我錢.
我才知道.

$^{881}$這個秘密.
意思不是叫你不收錢.
因為基督徒也有開餐廳.
開很多公司.
不是叫你不收錢.
不過你的行事作為.
有沒有稱得起是基督徒呢.
今天我反而.
很多時候我在墓會.
甚至在泰國裡.
我都要跟新人說.
新朋友來到.
我不想上帝的名被侮辱.
新朋友來到.
特別是女生.
不知道我講過沒有.
有些弟兄是基督徒.
後來讀神學.
不知為何不反對他.
後來去了另一間教會.
他經常去追女生.
所以有些新的女生來到.
我就告訴朋友.
你要小心.
他很喜歡追女生.
雖然我是那間教會的主持人.
但都要先說出來.
免得你看不住.
發生事.
所以我們的品.
你的品.
厲害的兩把口.
官字有兩個口.
品三個口.
你所作所為.
你不用說.
別人都看到.
雖然我媽媽在病床裡.
我們姐妹來唱歌祈禱.
大家都知道我們是基督徒.

$^{921}$想不到有一天.
有一個年輕人.
我不太相信來自年輕人的口.
你真的很神心.
神心甚麼呢.
他不是說我神心去看媽媽.
祈禱.
是神心.
所以這個神心.
歸於你們關顧部.
你們關顧部不單讓我有神心的稱號.
在病房裡.
還有我們早前有報道會.
其中一個來的人.
是我媽媽旁邊的床.
他出去了.
我們有繼續聯絡.
不單我聯絡他.
我們關顧部認識了他.
又和他有聯絡.
又帶了他回來.
不過都要繼續為他祈禱.
你看到保羅.
就算坐在監獄裡.
你看到他的書信不是埋怨.
埋怨這裡很冷.
埋怨這裡濕濕淋淋.
埋怨不夠衣服穿.
不夠食物.
不是.
你要喜樂.
在腓立比書.
亦是監獄書信.
他仍然有喜樂的心.
他仍然在監獄裡.
監獄才是國際化.
如果你有機會去監獄探訪.
有很多香港監獄也好.
泰國監獄也好.
來自世界各地的人.

$^{961}$在香港做完世界.
被人捉.
不幸被捉了.
所以就變成國際化.
我相信當時羅馬的監獄都是這樣.
所以保羅在當時.
他不單是品質好.
連監獄長都信了.
你會看到.
無論你的際遇是怎樣.
你仍然可以宣揚上帝的訊息.
在這裡.
沒有機會再繼續長長地說.
彼得前書一章十七節.
不過是非常令人明白的.
這裡說.
你們既然稱這位不偏待人.
個人行為審判人的主為父.
不單稱這位為上主.
我們稱他為阿爸符帖.
請大家要全敬畏的心.
到你們寄居的日子.
正如剛才泰語詩歌去唱的那首歌.
我們這條信主的路.
在這個世界上.
是像一個閘門.
是一個閘路.
有時是很不容易的.
但是.
上帝要回來的日子.
同樣是不知道何時.
所以今天就是宣教的日子.
今天你已經為神做好見證.
我承認在教會裡.
我們回來見證是很容易的.
但是最難度高的.
就是回到家裡.
好好見證.
這是上帝同樣祝福你的地方.
所以我們做什麼最辛苦呢?.

$^{1001}$你想過還做嗎?.
做基督徒也挺辛苦的.
跟從上帝的人.
24小時除了睡覺.
你也不在狀態裡.
沒有辦法警醒.
但是你放心.
主耶穌說.
上帝是不打頓的神.
祂還會看著你的.
求神祝福我們.
就算是難度高的家裡.
或者向世界出發也好.
我們整個人.
就是一個很美很美的活祭.
獻給上帝.
宣揚上帝拯救我們這個大好福音.
願上帝祝福大家.
\newpage



\section{腓立比書 2:17-30}
\label{sec:bW5i5_FWoTI}
\textbf{2024.04.28 《同心事主的知己》 腓立比書2\_17-30 (和修版) 講員 : 馮耀榮牧師}
\newline
\newline
連結: \href{https://youtube.com/watch?v=bW5i5-FWoTI}{\texttt{https://youtube.com/watch?v=bW5i5-FWoTI}} ~~~~ 語音日期: 2024-05-01
\newline
\newline
\hyperref[sec:jSl7w85PRnU]{\small{< < < PREV SERMON < < <}}
~
\hyperref[sec:index_chronic]{\small{[返順時目]}}
~
\hyperref[sec:index_scriptual]{\small{[返順卷目]}}
~
\hyperref[sec:W1Okagljeug]{\small{> > > NEXT SERMON > > >}}
\newline
\newline
腓立比書 2:17-30
\newline
\begin{longtable}{cl}
\hline
\hline
章節 & 經文 (和合本修訂版)\\
\hline
2:17 & \begin{tabularx}{0.7\textwidth}{X} 我以你們的信心為供獻的祭物，我若被澆獻在其上也是喜樂，並且與你們眾人一同喜樂。 \end{tabularx} \\ \\ \relax
2:18 & \begin{tabularx}{0.7\textwidth}{X} 你們也要照樣喜樂，並且與我一同喜樂。 \end{tabularx} \\ \\ \relax
2:19 & \begin{tabularx}{0.7\textwidth}{X} 我靠主耶穌希望很快能差提摩太去見你們，好讓我知道你們的事而心裡得著安慰。 \end{tabularx} \\ \\ \relax
2:20 & \begin{tabularx}{0.7\textwidth}{X} 因為我沒有別人與我同心，真正關懷你們的事。 \end{tabularx} \\ \\ \relax
2:21 & \begin{tabularx}{0.7\textwidth}{X} 其他的人都求自己的事，並不求耶穌基督的事。 \end{tabularx} \\ \\ \relax
2:22 & \begin{tabularx}{0.7\textwidth}{X} 但你們知道提摩太是經得起考驗的，他與我為了福音一同服侍，待我像兒子待父親一樣。 \end{tabularx} \\ \\ \relax
2:23 & \begin{tabularx}{0.7\textwidth}{X} 所以，我一看出我的事怎樣了結，我希望立刻差他去， \end{tabularx} \\ \\ \relax
2:24 & \begin{tabularx}{0.7\textwidth}{X} 但我靠著主自信我不久也會去。 \end{tabularx} \\ \\ \relax
2:25 & \begin{tabularx}{0.7\textwidth}{X} 然而，我想必須差以巴弗提到你們那裡去。他是我的弟兄、同工和戰友，是你們差遣來供應我需要的。 \end{tabularx} \\ \\ \relax
2:26 & \begin{tabularx}{0.7\textwidth}{X} 他很想念你們眾人，並且極其難過，因為你們聽見他病了。 \end{tabularx} \\ \\ \relax
2:27 & \begin{tabularx}{0.7\textwidth}{X} 他真的生病了，幾乎要死。然而神憐憫他，不但憐憫他，也憐憫我，免得我憂上加憂。 \end{tabularx} \\ \\ \relax
2:28 & \begin{tabularx}{0.7\textwidth}{X} 所以，我更要盡快送他回去，好讓你們再見到他而喜樂，我也可以減少憂愁。 \end{tabularx} \\ \\ \relax
2:29 & \begin{tabularx}{0.7\textwidth}{X} 故此，你們要在主裡歡歡喜喜地接待他，而且要尊重這樣的人， \end{tabularx} \\ \\ \relax
2:30 & \begin{tabularx}{0.7\textwidth}{X} 因他為做基督的工作不顧性命，幾乎至死，為要補足你們供應我不夠的地方。 \end{tabularx} \\ \\
[1ex]
\hline
\hline
\end{longtable}
$^{1}$今天跟大家分享的是一個比較個人的見證.
或者心路歷程.
剛才我們看到唱粵曲的弟兄姊妹.
都是女生多於男生.
教會在年輕的時候團契.
男生女生都差不多.
因為我們帶小朋友回來.
我們生的小朋友應該是男女差不多的數字.
到青少年的時候.
教會的男生都有一定的人數.
在學校裡傳福音.
其實也有很多男生順序進教會.
但到了長者就開始出現女生多於男生.
主要原因是以前男生是早離世.
60歲或70歲就離世.
教會的女生90歲也有很多.
以前的比例少了.
男生也可以理解.
但現在醫學發達.
男生像我岳父是90多歲.
現在只要他不去看醫生.
男生也會比較長命.
教會裡的男生仍然是雙歲少.
什麼原因呢?.
我剛剛退休.
所以我開始進入退休的男性.
會是什麼心情或表達.
今天在座的姐妹.
你起碼理解一下你的丈夫.
他的心情.
或者我們年輕一點的.
你理解你爸爸的心情.
我也應該進入這個狀態.
所以我今天和大家選一段經文.
一個很出色的人物保羅的個性.
他怎樣看他和他周圍的人的關係.
他的側重點.
和他身為一個男性.
他會有什麼獨特的地方.
這段經文的程序表是有的.

$^{41}$所以大家可以去看一看.
我沒有把PowerPoint做得好.
所以字比較小.
但本身的秩序裡有一段經文.
這段經文簡單來說.
有一節經文是20節.
我們一起讀一讀這一節.
我看到你的程序表.
20節我們一起讀和合本.
因為我沒有別人與我同心.
真正關懷你們的事.
這裡是和合本的翻譯.
沒有人與我同心.
去關懷菲律賓教會的人的事情.
這裡是呂鎮忠譯本.
我們教會是用和修本.
和合本是很久以前的.
經過現代化後.
和修本我也有負責.
其中一兩卷書叫做審核.
在大約一千多年的時候.
我們開始將和本比較古舊的文字修改.
所以在香港負責或有新譯本.
都是找一些聖經科老師做審閱.
我們不是做翻譯的.
呂鎮忠這個人其實翻譯聖經很早.
我想五,六十年之前他已經開始翻譯.
呂鎮忠的舊約到新約已經出版了很久.
在很多和合本的修改本或新譯本之前.
已經有呂鎮忠譯本.
他的譯本在這裡.
我也是最近準備講章的時候發現的.
他將「與我同心」這樣翻譯.
因為在這裡我沒有別的知己.
去關心教會的人.
所以和合本的譯法或和修本的譯法.
和大部分中文聖經的譯法.
都是將到保羅的心境說.
沒有人與我同心.
只有兩個人.

$^{81}$這段經文其實我們不會太多講這兩個人.
我們只會講保羅的心境.
他提到提摩泰.
接著提到爾巴忽提.
他說這兩個人是唯一和我同心的人.
呂鎮忠譯本說.
他是我這兩個唯一的知己.
爾巴忽提後來病了.
他生病的時候.
保羅就叫他回到菲律賓.
本來這個爾巴忽提是菲律賓教會.
差去關心保羅.
當時保羅在坐牢.
或者年紀大的時候去照顧.
保羅就送這個爾巴忽提回到菲律賓教會.
聖經就講到這段經文.
他應該是退休年齡.
開始離世.
在腓立比書中說.
我人生已經到消電的時間.
我整個人生奉獻給神.
走到盡頭.
他用一個獻祭.
一個灑消電.
好像一些水.
消電在祭壇.
或者一些酒消電.
所以用消電的祭去表達.
在他人生中.
他說其實沒有什麼人和我同心.
我沒有什麼知己.
我太太前一兩個月問我.
我有沒有知己.
其實我朋友也不多.
男生想清楚怎樣定義知己.
我一定沒有太太有這麼多知己.
因為我太太的WhatsApp是很多的.
經常都嘟嘟嘟.
我睡覺都聽到她嘟嘟嘟.
我將WhatsApp.

$^{121}$有WhatsApp的人下載.
我關掉了.
我不會讓他嘟我.
我WhatsApp是有限的.
我的群組就是我負責的小組.
其他全部都沒有.
所以我太太和我比較.
我太太也是一個傳道人.
她的工作是做院務.
我在神學院教書.
兩個都是全職事奉.
男女是有別的.
平均來說都是這樣.
男生是沒有什麼朋友.
不要說知己.
朋友也不多.
他有一個同班同學畢業了二十多年.
他在讀書的時候.
每一個星期一起祈禱.
一個小時一起祈禱.
聊天和分享.
畢業後沒有經常見面.
但大家都結婚了.
他一年有兩三次約出來.
吃飯,聊天和祈禱.
到現在都有.
他同班同學都退休了.
去了英國.
他視像和她見面.
每年有一次祈禱.
我沒有這樣的人.
我有同班同學.
我記得有一個同班同學.
是韓國來的華僑.
我年輕的時候也有朋友.
他會帶我去參觀教會.
所以年輕人也有朋友.
年輕人也有知己.
但大家忙.
有自己的工作職責.

$^{161}$所以這個同學畢業後回到韓國.
十多年回來香港事奉.
我們在香港幾年有一次機會見面.
到他和我退休後.
他打電話和我聊天.
問我何時有空去中環.
他住在馬灣.
我住在長洲.
所以我們多數在中環約見.
都是他約我的.
兩三年一次.
出來吃下午茶.
所以我和太太.
都是他約我的.
只是出這個.
其他都.
其實我作為一個教務同工.
其實也有很多應酬.
我們學生有關心小組.
每一個月兩個月會見面.
我有很多同事.
下午基本上都會去飯堂吃飯.
都要在不同的台.
和不同的人聊天.
所以我不是那麼孤獨.
在未退休之前.
我退休後.
這些事情可以不用做.
可以不用應酬.
我大部份時間.
在辦公室和宿舍.
我今天出來港島.
是很少的.
我很少出來香港.
雖然兩元也好.
很少出來.
沒什麼人.
因為兩元便宜一點.
會出來香港逛逛.
在長洲也可以逛逛.

$^{201}$為什麼要出香港.
我到現在還沒去過國內.
為什麼要去大陸吃東西.
我在長洲吃個麵包就行了.
所以我很心虛.
退休後不用應酬.
以前說一個月.
都要出來港島一次.
現在可以不用.
間中都出來港島.
所以男生.
會比較.
孤獨一點.
是一個普遍的現象.
就算我們是一個.
服侍人的職業.
即是我們.
職責上.
要應酬.
要見人.
好像崇拜完之後.
要跟頂尖的打招呼.
職責上都會的.
都可以談一些事情.
但如果可以不用.
退休後不用自己職責.
躲在家裡坐坐.
輕鬆很多.
不知道大家是否明白.
究竟我們需不需要知己.
從一個男性的角度.
需不需要朋友.
或者需不需要很多朋友.
特別在我們這個年齡.
需不需要重新.
去認識朋友.
多數男生.
覺得這個不是.
他的首選.
不是他的首選.

$^{241}$所以當我們.
教會或者我們.
去向這一批.
年長的男性.
退休的時候.
現在他還有二三十年命.
現在的男生.
還有二三十年.
還很長.
怎樣應對.
所以我自己就想一下.
其實我的個性都是很.
孤獨.
所以.
以前的人六十歲.
都已經老了.
現在六十歲.
開始一個新的人生.
起碼有三十年命.
如果正常.
我們現在平均香港年齡的男生.
都八十多歲.
怎樣可以過得這麼久.
這麼孤獨.
連太太都開始有距離.
我和我太太.
大家工作忙的時候.
大家有自己的工作.
是不覺得距離這麼遠.
因為你沒有機會.
有太多時間坐在一起聊天.
因為大家都忙.
各自的工作.
到我開始退休.
她快要退休的時候.
想一下將來我和她一起.
廿四小時有什麼可以說.
已經說了.
其實都有點擔心.
有什麼可以說.

$^{281}$這個現象.
暫時沒有太大的衝突.
因為我太太仍然是全時間工作.
所以她和我溝通的過程.
以前都沒有什麼問題.
因為大家.
以前我因為有工作做.
我是一個牧師.
我是一個神院牧師.
所以我都要說很多話.
所以說.
我最近和我太太.
她的事奉是院牧.
所以她.
就好像上兩天.
她開會.
她負責籌錢.
牧師或掌執.
那邊灣仔.
她在律敦自做醫院.
所以幾間教會支持她.
她就要去.
和他們談一談明年的預算.
怎樣做.
她回來就和我說.
她之前和我說過.
她會做這件事情.
她和我說.
首先開始說.
幾點鐘約他們.
下午的時候.
我和牧師.
牧師遲了來.
跟著她告訴我.
其實她由開始說到最後.
其實我最有興趣.
究竟有沒有人支持你.
其實她應該很早說.
沒問題.
現在教會有個特別奉獻.

$^{321}$可以支持二十萬.
但我要聽久一點.
才知道結果是怎樣.
其實有時她說了之後.
沒有說結果出來.
遇到一個人.
發生什麼事情.
告訴我.
現象是女生說話很多.
男生呢.
以前我做教師.
說話很多.
但不用說的時候.
就不說了.
我們的距離就越來越遠.
有些時間.
我們有些衝突.
以前我叫她早一點說.
說完結果.
她就會不開心.
我見到她不開心.
下次也不要了.
忍一下.
無謂不開心.
但我忍.
我的臉色可能在忍.
她也開始感覺到.
所以溝通不是.
我們都是牧師.
都是傳道人.
所以我們不會吵架.
不過說話多少.
距離是很遠.
我想她很想找一個.
知己能溝通.
我沒有興趣.
要找一個很有知己.
能溝通.
我其實沒有興趣.
沒有興趣說話.

$^{361}$如果不是職責上.
我不是唯一這樣的人.
我觀察教會很多弟兄.
都是這樣.
起碼我發現我爸爸也是這樣.
我爸爸小時候.
他沒有說話.
他看報紙.
很早退休.
在家裡看完報紙.
有一兩次他帶我去喝茶.
在喝茶的時候.
他也沒有說話.
我沒有說話.
因為我小學六年級.
所以不知道要跟爸爸說什麼.
所以我不知道要跟爸爸說什麼.
爸爸也不打算跟我說什麼.
所以我見他的人生.
很孤獨.
因為我回教會.
所以我職責上需要聊天.
要教書.
男生在教會長大.
我年輕的時候在教會長大.
如果你不喜歡聊天.
很難回教會團契.
很多不喜歡聊天的男生.
留不住在教會做團契.
因為團契就要聊天.
要分享.
你勉強分享一兩次.
但你不喜歡分享.
你不能生存在團契裡.
所以不喜歡分享的男生.
很少留在教會.
所以到我們這個年齡.
如果要向年長的人.
人生比較傾向.
比較內斂的.

$^{401}$比較寡言的男性.
接觸他們.
帶他們順處.
一定不能用傳統的方式.
舉辦福音樂曲.
不是我不贊成.
對女性可能比較容易.
對男性比較難.
還有.
我有印象.
我太太也有做福音工作.
做讚美操.
讚美操裡面.
男女比例更加激烈.
30人裡面.
只有一個男性.
我忘記了.
導師也是女性.
所以如果靠福音樂曲.
或者讚美操.
帶人順處.
多數姐妹都可以.
既然有這麼多姐妹.
更加男性不敢進去.
我年輕的時候.
是跳舞的.
但到了我這個年齡.
要在別人面前跳舞.
其實是比較麻煩的.
比較不習慣.
所以.
很多長者的活動.
要請男性.
是比較難的.
最簡單的.
就是進老人院.
這個年齡.
都要準備一下心境.
有些人就說.
這個老人院好.

$^{441}$但有很多活動.
很多活動.
很麻煩.
其實.
進老人院.
對男性來說.
還是不要去.
在家裡好.
很多人都喜歡回家.
已經不想去老人院.
老人院還有很多活動.
通常我們的想法.
老人家就要有群.
要有多些活動.
但對男性來說.
很多活動其實更慘.
我們的心境.
不是從當事人來說.
當然有少部分的男性.
社交的能力.
或者興趣.
會強一些.
但大部分的男性.
都是一個孤獨的.
孤獨的甚麼.
生物.
孤獨的人.
這個比較普遍的.
現象.
不知是甚麼現象.
我想西方也有這個現象.
你想想.
我興趣做甚麼.
我去旅行的時候.
有兩個地方會去.
第一個就是杭州.
我很年輕的時候.
去過一次杭州.
那時候是80年代.
剛開始開放的時候.

$^{481}$到現在.
我間中就想去杭州走走.
在湖邊走一圈.
走一圈.
是沒甚麼事做的.
不用做甚麼.
就是走一圈.
你就會有.
一千年.
蘇杭是千年的歷史.
你就好像.
蘇東坡見過面一樣.
你走在蘇堤.
你就超越了時空.
這是我第一個地方會去.
第二個地方就是京都.
京都在鴨川走.
由下面一直走上去.
這是我很有興趣的地方.
所以我太太.
我在安息年有一年.
在京都住了一年.
在一間大學做交換的老師.
本來她可以來多一點時間.
最後她還是來了三個星期.
她放下了我.
有一年在那裡.
其實我的宿舍是很大的.
我住在那裡.
其實我的宿舍是很大的.
但她工作也是很有.
她很想工作.
她不想放假.
沒甚麼事做.
但我想最關鍵是她不想那麼悶.
一個人在那裡.
沒甚麼事做坐在那裡.
坐.
不要說坐一年.
我太太是一個.

$^{521}$很工作型的人.
她很忙.
她在家裡坐一會.
十分鐘不夠.
就開始掃地.
做很多事.
煮飯.
她是一個工作型的人.
她是不能坐下的.
她不想浪費時間.
所以去一個地方.
沒甚麼事做.
坐在那裡.
看著河邊.
人生的時間.
對她來說是浪費了.
我比較極端的.
我是一個.
可以坐在那裡.
思考的人.
很長時間.
沒甚麼興趣.
沒甚麼特別想做甚麼.
要去那裡.
其中一個旅行後.
我太太.
最後一兩天就會買手信.
買手信給誰呢?.
她的同事.
全部日本的手信.
都是在香港買到的.
我去旅行.
還要想買甚麼手信給朋友.
我是不會的.
如果我一個人去.
一定會.
就是應酬才做.
所以男性.
我是比較極端.
但我想不是一個極端的現象.

$^{561}$女性是關心周圍的人多一點.
所以她會.
用自己的時間.
自己的旅行.
都會想到.
買些東西給朋友.
那些東西.
吃不吃不重要.
所以我們.
有沒有朋友呢?.
應該不多.
比較孤獨.
做了基督徒.
如果你看保羅.
其實這句說話.
裡面都挺獨特的.
他說.
除了這兩個人.
提摩太和.
以巴弗提.
其實我身邊都沒有人.
沒有朋友.
沒有知己.
所以我有點安心.
我像不像一個牧師.
要有.
群體的生活.
原來保羅其實都很孤獨.
他只有兩個朋友.
兩個知己.
我還差一個.
神學院的畢業生.
可以說是知己.
有約會.
其他我可以約會.
但沒有動力.
原來保羅其實在人生裡.
只有兩個朋友.
兩個知己.
由朋友進入知己.

$^{601}$都可以理解.
為何男生會.
特別是.
保羅這些人.
他的興趣,他的負擔.
不是買甚麼食物.
甚至他的重點.
都不在乎朋友對他的想法.
如果我不買他會怎樣想.
其實我都不會想這些.
都不會想別人怎樣想我.
因為我一個人.
已經是存在的價值.
不會因為別人.
怎樣想我去存在.
所以在這點裡.
保羅.
他不太在乎別人.
如果你想想.
腓立比書講了一個很重要的金句.
或是一個總體的訊息.
就是我的人生.
只是為了神.
只有基督.
是我生命的至寶.
其他一切都不重要.
既然是這樣.
其實他都不需要太多朋友.
不知道大家明白嗎.
因為有耶穌基督在我心.
就是我最終極的.
所以.
這個其實有些難處.
難處的意思是.
我今天希望不會.
給大家誤會.
不是叫弟兄不參加.
教會聚會.
或者不是請一些.
弟兄去.

$^{641}$或者請一些男性回來教會聚會.
不過男性.
是會躲起來多一點.
這個都是一個現象.
我們要理解一下他們的心境.
他願意接受應酬.
其實不是他本能的.
一個現象.
保羅的人生.
就是說.
我的生命只有基督.
耶穌基督的心境.
我不在乎別人.
怎樣看我.
所以他所要討的喜悅.
就在基督裡面.
相對來說.
他不是太多人際的牽掛.
或者牽引.
最重要的是.
屬世的人還要結婚.
還要遷就太太.
或者配偶的想法.
保羅的人生.
說.
我為了基督.
連婚姻都不結.
我的心境.
是要討耶穌基督的喜悅.
所以他可以結婚.
他結婚的人應該彼此相愛.
丈夫愛妻子 妻子愛丈夫.
哥林多前書後書全部都是保羅寫的.
但當他選擇的時候.
他將整個人.
投身在基督身上.
所以保羅需不需要朋友呢?.
他需不需要知己呢?.
其實只有耶穌.
才明白我們.

$^{681}$所以從我的角度.
有些人的世界.
是會比較單向.
比較.
沒有那麼社交性.
沒有那麼.
在意別人.
在意別人的看法.
男性.
特別是有一定事業的男性.
他會傾向這樣.
他不在乎太多別人的看法.
所以在教會裡.
有些男生在工作上.
我們出來工作.
無論是哪個工作.
都要應酬.
所以我們都會正常的有社交.
但到了一個退休下來.
我們會有些人.
會有些人.
但到了一個退休下來.
應酬不是我們的工作.
是.
我現在的人生.
有自己的時間.
當然這個自己的時間.
不一定是.
不可以和其他人接觸.
但會想想我可以.
還有人生裡面.
我面對自己的時候.
我怎樣改人生.
不需要再.
為了責任.
向別人交代.
這是普遍的男性.
的做法.
當然.
這樣是要改的.

$^{721}$太過孤獨.
不是太好.
但你要他.
繼續.
去用職責.
和其他人應酬.
接觸我們不算很理解.
他的本性.
我最近.
剛剛在教會裡面.
有一個.
同工的哥哥.
或姐姐.
有一個親人.
有病,他們兩個還未結婚.
所以不可以叫他們結婚.
他有病.
都是嚴重的病.
他願意信耶穌.
所以我作為牧師.
到他家裡面.
探訪他.
確定他.
他就是一個很典型的.
年長的人.
其實他信不信.
我們現在這個年齡的人.
其實在小時候.
都接觸過基督教.
不是完全沒有接觸.
所以到和他談的時候.
其實他會信的.
他清楚.
當然開始的時候.
他先說說他的情況.
他很準確地說出.
自己的病的病情.
很準確地說出自己的病情.
他是屬於.
顫動人.

$^{761}$這個病.
這是一個通俗的說法.
他大概知道.
什麼時候說話會不清楚.
什麼時候會吞噬.
開始出現問題.
所以他上網.
看到很多資訊.
他也有上網看.
一些類似病發的人.
他們之間有一個小組.
所以他在家裡面對這三四年.
他是很獨立能力.
去清楚交代.
亦都很清楚地說出他信著.
接著我幫他受洗.
所以年長的人.
他比較.
看回他之前.
問一問他親人.
他之前.
其實他都不多回教會.
其實他都是比較清楚.
比較內斂的人.
但信主仍有機會.
仍有機會.
他們會清楚知道.
自己想要做什麼.
但一個人信主.
特別是我們這個.
上了年紀的.
香港的男孩.
他不容易.
帶他回教會.
因為他要重新適應.
一個社群.
他是不是很有動機.
可以嘗試.
但大部份情況.
其實我們可以單刀直入.

$^{801}$處理他的信任.
他們都肯信主.
所以.
福音的工作.
對我們一些退休了.
還有十年,二十年.
或者就在病患裡面.
好像我太太做阮木.
她帶信主的人裡面.
男孩和女孩一半半.
因為在教會.
男孩會少.
但在醫院裡面的男孩.
一樣會多.
因為現在男孩的壽命.
高了.
所以他一個星期,兩個星期.
就有人帶信主.
他又開始洗禮.
都是男孩.
這些男孩講一下他的人生過去.
又是很孤獨的.
大部份都很孤獨.
沒什麼人講.
家人不太喜歡他.
甚至有時都沒什麼人探訪他.
因為他不像女性.
有多些家人.
親人之外的朋友.
男生.
你病了.
有朋友來探訪你的機會是少的.
但女性病了.
女性的朋友.
探訪的又高了一點.
其中一個行業.
就是做老師.
很多女老師病了.
都有人去探訪她.
但都沒人探訪她.

$^{841}$(笑).
女性比較有人情.
所以當老師.
病了,你的學生可能已經.
十多二十年了.
我們有的印象.
多數是女老師.
男老師都會罵我們.
他們不會很親近.
所以.
如果我有父母.
男孩女孩都要探訪.
父母都要.
其中一個影響我們.
最深的.
就是小學,中學.
甚至在神樂院的老師.
都會出現這樣的現象.
在神樂院我是老師.
如果是女的.
我以前有一個老師.
叫陳淳華老師.
她是媽媽.
她是全界學生.
東北是媽媽.
男生女生.
陳淳華老師.
她是媽媽的角色.
是一個很人際關係.
很好的關係.
所以我和太太結婚的時候.
她都請了我們一起去吃飯.
所以.
到她去世的時候.
很多人接觸她.
但我們有一些.
很出名的男老師.
他在醫院裡.
通常都是.
後來的老師.

$^{881}$或後來的主任牧師去探望她.
回眾探望她的機會很少.
因為她不是一個.
或者不像曾牧師.
有情色一點的.
男性的領袖.
傾向職業上.
或職責上的指導.
人與人之間的關係.
是相對會淡薄一點.
人與人之間的關係.
是相對會淡薄一點.
人與人之間的關係.
是相對會淡薄一點.
這是不是對錯的問題.
這是現象.
如果你看看你的家.
你的媽媽和爸爸.
是誰比較內向.
是很難請他們回來教會.
我媽媽很早就請他們來.
去到長洲.
她開放日都會來.
我爸爸完全不會帶她去長洲.
他都關心我的.
我去外國讀書的時候.
我還有一張照片.
我和媽媽拍照送我去.
去外國的時候.
80年代.
我記得我爸爸是關心我的.
他送了一件大衣給我.
但他沒有來送機.
那件大衣.
是很舊的.
很舊款的.
我都不好意思穿出來.
但我會保留著.
所以爸爸不會來送機.
但不是表示他沒有感情.

$^{921}$但他的表達手法.
是有一點距離的.
我們接受.
我們的親人.
是有不同的.
男性的感情.
是比較隱藏一點.
比較含蓄一點.
比較不是用言語去表達.
現在我們開始.
教年輕的男孩子.
說多一點.
以前上一代的父母.
說我愛你.
現在男孩子都可以說.
所以我們新一代可能有一點不同.
但起碼我這個.
六十多七十歲的男性.
要說到我愛你這個詞.
都是挺大的勇氣.
當然在教會裡面唱詩.
一路說我愛你就很簡單.
所以表達感情.
男孩子是會遲一點的.
不是綁心舞的.
所以保羅.
在這裡.
他說我沒有什麼朋友.
沒有什麼知己.
其實都可以理解.
因為他是一個成功的人士.
他是一個獨立的人.
他不需要在感情上.
和周圍的人有太多的關係.
在人類社會裡面.
男性相對來說是比較獨立.
所以我小時候.
已經不在家裡住.
我家裡是水上人.
所以有十兄弟姐妹.

$^{961}$如果我要在家裡住.
就要在船上幫忙.
我就讀不到書.
所以我的哥哥姐姐都沒有機會讀書.
當然我是第八.
所以我就不用幫忙家裡.
就可以去讀書.
我去讀書.
以前水上沒有機會.
可以固定學校.
我就在.
一個親人的家裡.
他本身也有七八兄弟姐妹.
我和他長大了.
但由於我不是在家裡.
所以我好像有很多朋友在旁邊.
很多表兄弟姐妹.
但又好像我不是在家裡.
所以我是一個很孤獨的人.
很喜歡孤獨的人.
所以我在長洲住了幾十年.
都不覺得悶的原因.
就是因為這個.
所以我剛才說的.
不算是很正常的成長年齡.
由小時候已經不太多溝通.
不是到處都沒有人.
只是由於.
小時候比較.
好像失去了家庭.
就有一種.
內向.
但內向之餘.
由於信仰改變了.
所以我教會有很多弟兄姊妹.
然後我有很多朋友.
我教會有很多朋友.
很多弟兄姊妹.
我有很多同學.
有很多校友.

$^{1001}$所以正常的接觸是有的.
但保羅.
他也有很多教會.
寫過很多信.
他也有很多弟兄姊妹.
但他最後離開世界之前.
其實他有一個心境.
真正和我同心的人.
都不多.
只有提摩太和爾巴弗提.
所以我有些安樂.
有些安慰.
不是我有太大的問題.
我一定有少少問題.
因為自小成長比較孤獨.
到現在都喜歡孤獨.
這個也不是完全沒有.
我教會以前在長洲一間教會.
有一個執事.
他是牙醫.
他太太.
剛和我太太一起吃飯.
這位醫生.
他一年多.
過世了.
他過世之前.
他出了香港.
整家家庭搬了.
我們沒有見面很長時間.
他太太.
在建道.
和我做過同事一段時間.
所以我見他先生離開了.
我太太.
最近碰到他.
知道他先生離開了.
他本身.
有些抑鬱.
因為先生離開.
先生離開之前有病.

$^{1041}$醫生多數都不看醫生.
他知道自己應該很嚴重.
所以他都不打算要醫.
所以做太太的.
就推他去醫.
但到最後.
他都沒有醫.
最後發現有症狀.
已經去到後期.
短時間內.
一年多就離開.
所以.
他太太有些抑鬱.
最近.
走出來的心境.
碰到我太太.
在醫院幫忙做義工.
所以我們一起吃飯.
這就很典型.
他做醫生的時候.
以前做教會執事.
他都會一起團契.
他是一個可以應酬.
團契的人.
但當他一個人出去.
他不用.
退休後.
他喜歡聽音樂.
一個人聽音樂.
不需要去團契.
因為他說.
以前做醫生的時候.
轉牙的聲音.
令他很煩躁.
他能夠解決的問題.
不是去做團契.
他能夠聽音樂.
喜歡聽音樂的人.
應該是比較孤獨的人.
比較孤獨的人.

$^{1081}$他喜歡聽音樂多於.
聽別人的說話.
其實都挺多人.
男性.
他不是沒有專注的東西.
但他相對來說.
職業上.
要應酬或回教會.
聚會之外.
他可以選擇.
都是比較孤獨的.
這個不說對錯.
這是一個社會現象.
香港的男性.
年紀大的男性.
你看看在老人院.
外面樓梯.
坐下的那些人.
女性會聊天.
男性都是坐在那裡.
呆地坐在那裡.
偶爾會有些維園阿發.
偶爾有些阿發都會說很多話.
但大部份的男性.
都是.
沒甚麼話可說.
我可以想像一下.
我已經是屬於維園阿發一類.
我都是屬於坐在那裡.
想想事情.
所以為何我會喜歡呂振中.
為何會特別對這段經文.
有興趣.
其實.
沒甚麼人有興趣呂振中.
因為這個人本身.
比我還孤獨.
他起碼差不多.
不會說到.
他是一個事奉者.

$^{1121}$他的人生.
從創世記到啟示錄.
把希伯來文和英文.
和所有英文翻譯.
中文翻譯.
做一個思考的人.
他一定要捱得住這種孤獨.
否則他不會翻譯到這本書.
所以我可以想像.
呂振中是一個.
聖經學者.
是一個.
很內向的人.
他已經不需要.
他有職責上.
和別人溝通.
他人生唯一的興趣.
就是翻譯聖經.
他把他的人生.
整個人生燃燒在這裡.
現在他有的成果.
其實都不是很多人.
能夠欣賞到.
因為我們沒有訓練.
把中英文.
和希伯來文.
希臘文.
去看究竟這個翻譯.
是否合適的翻譯.
最後.
我很想.
將來我們有機會在天家再見.
問一問.
其實他把這個翻譯成知己.
可能是錯的.
因為大部份的翻譯.
都叫做同心.
我都覺得.
和本的翻譯應該是好一點.
為何他會.

$^{1161}$翻譯成知己.
其他大部份的翻譯都是很好.
這一部的翻譯.
他突然中了寶來.
說沒有別的知己.
其實我印象很可能就是.
呂振中在連繞的時候.
他其實都想有知己.
只是.
他流露了.
在他的世界裡.
把寶來的心境的另一種表達.
表達出來.
其實寶來.
不太在乎有沒有知己.
在這段經文裡.
用和學本的翻譯.
他說與我同心.
剛才所說的.
不是對錯的問題.
不過他有一個盲點.
就是我們想找朋友或知己的時候.
都是一個比較自我的世界.
知己這個字眼.
就是一個自我的世界.
就是要對方明白我.
我是很喜歡.
研究聖經的.
我太太.
都讀了很久神學的.
不過她有沒有興趣知道我這個翻譯是哪一個呢.
我想她都沒有興趣.
因為.
她的職責是做院務的.
她探病的時候.
病人不想知道翻譯是怎樣翻譯的.
她的病人需要知道.
怎樣得救.
她的病好不好翻譯.
醫生有沒有錯.

$^{1201}$我是不處理這些問題的.
醫生有沒有錯.
應該吃什麼.
所以我去探病.
我是不會插手.
吃什麼好一點.
因為我又不是醫生.
叫他吃什麼好一點.
吃錯了東西.
我有問題.
我沒有興趣.
所以我去探病大概十分鐘就好.
讀一段詩篇.
鼓勵一下.
我太太去探病一個小時都不出來.
因為她說什麼.
吃了什麼好.
你兒子沒有來探病是什麼原因.
她說了很多跟病沒有什麼關係的東西.
我是忍受不到.
無謂的東西說出來.
但她是容易.
單人順處.
因為大部分人.
你跟她聊天.
單刀直入說信耶穌她多數都不相信.
但如果你關心她.
你跟她說了很久.
你講什麼她都相信.
所以我.
我多數有病人.
都是拉我太太一起去探.
如果不是.
我忍很久都不得救.
這是本性.
甚至是職業病.
她的工作.
跟我的工作.
是不同的.
我的職責是處理聖經的解讀.

$^{1241}$當然我是牧師.
我都要關心.
傳福音的事情.
但我太太.
比我方面好很多.
所以傳福音這方面.
需要有社交能力.
需要有朋友.
需要有多些願意做知己.
所以還好我們結婚了.
我沒有的東西她有.
她沒有的東西都不重要.
因為她對李振聰.
是否認識都不重要.
她其實不需要李振聰.
她問我.
其實最簡單.
究竟是什麼原因告訴我.
然後我去講道.
所以她不需要去做研究的過程.
她需要花時間.
跟當事人.
建立朋友的關係.
所以我不是說沒有朋友.
不是說做朋友不重要.
我知己.
只是男生傾向這方面.
是很缺乏的.
男性傾向這方面很缺乏.
當然有少部分男性.
是不同的.
亦有少部分女性.
都是很孤獨的.
這是.
平均的問題.
所以我們面對.
不同性別或不同的個性.
我們怎樣處理呢?.
這段經文.
給我有一個這樣的想法.

$^{1281}$現在我不是用它.
去證明自己的對或錯.
我們看看保羅的心境.
他怎樣跳出.
這個.
他的個性和其他人.
之間的關係.
保羅沒有找一個人.
比他更好的.
跟他做朋友.
他找的是提摩泰.
是他下一代.
而提摩泰不是很出色.
因為提摩泰沒有寫過書.
記住提摩泰書是保羅寫的.
通常你找的朋友.
都是跟你工作能力.
或恩賜.
差不多.
如果他跟你.
可以一樣的時候.
就變成知己.
我找的朋友.
或知己.
他是做中國侍工.
他是做廣播.
所以在韓國時.
有電台廣播.
回來香港的電台廣播.
所以國內很多人聽過他廣播.
而我本身神學教育.
在我最後這十多年.
我多數入國內.
做神學培訓.
我們找知己.
其實我們同樣有負擔.
因為我們在社會中的傳呼音.
和廣播有關係.
所以在現實世界.
我們有朋友.

$^{1321}$和知己.
男性多數和他的興趣.
恩賜.
或才港.
差不多.
所以.
不是很多人有這個負擔.
或興趣.
所以能夠有的朋友就少一點.
我們不是特別有興趣.
去哪裡旅行.
雖然也會去.
但不是特別有興趣.
所以保羅不需要有很多朋友.
因為他比較事業型.
所謂事業型.
就是想達到做那件事.
所以.
這種情況.
他不算和同輩的人太多.
但他也看.
提摩泰.
和爾巴忽提.
這兩個人.
基本上在才幹.
在成就.
他的成就連提摩泰也不如.
因為.
沒有人叫做爾巴忽提團契.
所以.
提摩泰團契.
也因為保羅的原因.
提摩泰在歷史上有一個名稱.
但保羅他選擇的人.
這兩個人.
他不是特別強的恩賜.
好像保羅.
或他的成就也像保羅.
保羅選擇他的時候.
就是我們剛才說的.

$^{1361}$最後第二節的部分.
其他人都求自己的事情.
不求耶穌基督的事情.
我今天和大家說.
最關鍵的部分.
知己本身.
我覺得不是一個很好的翻譯.
因為我們也很想.
別人理解自己.
所以我會找一些.
和我相似的興趣.
能力和背景的人.
所以做護士的.
做醫護的.
回來教會有醫護的人.
比較容易一起.
或者喜歡唱詩的.
就會和唱詩班一起.
做生意的.
會多些和做生意的一起.
或者做社工的.
或者做公務員的.
會有一個群體.
這是知己.
知道自己和自己一起.
因為我們都是一起信耶穌.
所以都是以自我為中心.
這是正常的.
社群的.
社團之間的關係.
這不是錯的.
很正常的.
我看到神學院的同學或是校友.
都會多說兩句.
因為大家有同樣的背景興趣.
或者做法.
但是在這個基礎之外.
保羅他認為提摩泰.
和以法不提.
是他的朋友.

$^{1401}$或者知己.
因為這兩個人.
都是不求自己的事情.
是求耶穌基督的事情.
所以就是說.
只要我們有人.
願意放下自己.
去傳福音.
去求耶穌的事情.
這些人都可以是我們的知己.
不知道大家明白嗎?.
有很多神學院的畢業生.
來到教會.
都是很孤獨的.
因為他們作為傳道人.
已經有少許不同.
所以他們不敢說一些.
剛才我說的話.
我所說的.
說我是有些自廢.
你做牧師傳達.
不要說這麼快.
說我有些自廢.
他不知道找你聊天好.
不找你聊天好.
如果找你聊天好.
好像煩著你.
不受歡迎.
所以我這個年齡退休.
可以這樣說.
你沒有職責.
我也沒有職責找你聊天.
所以大部份傳道人.
一開始不可以說得太快.
所以他們在教會裡.
是很孤單的.
一旦很孤單.
他就會容易.
失去朋友.
我通常這樣鼓勵.

$^{1441}$我們過性.
是可以過孤獨.
但我們好像保羅.
他對提摩泰.
對伊巴菲提.
他當他是知己.
因為他們本身都是為了福音的緣故.
不求自己的事情.
願意去做.
所以我們有人.
唱福音的詩歌.
有人去跳舞.
健康的舞.
帶人去信主.
其實都是為了福音的緣故.
這些人就是我們的知己.
這樣我們就不孤單.
所以就算做傳道人.
雖然職責上有些事情.
不適合說.
但我知道我們教會的弟兄姊妹.
都是為了福音.
願意去參與.
這些就是我們的知己.
所以不一定要.
跟我一樣興趣.
去杭州.
或者喜歡做什麼.
在我們的興趣上.
可以很不同.
但在我們同一個心智.
同心去求耶穌的事情的人.
都可以是我們的知己.
所以我最近.
教會的一個牧師.
他也很孤單.
因為剛剛來我教會.
我另一個牧師就走了.
所以他有點孤單.
找了一個很好的傳福音的人.

$^{1481}$他帶我們教會的弟兄姊妹去傳福音.
我說其實.
任何一個人.
當你覺得孤單.
都是正常的.
我們不是世上興趣.
恩賜和很多人是一樣的.
有時我們.
特別是男性都不在乎這些.
但我們需不需要有知己呢?.
需不需要有人.
和我同心呢?.
保羅都需要.
這個人已經快離世.
在離世之前.
我已經囂墊了自己的生命.
這個人仍然在在世的時候.
他能夠說出.
雖然很多人沒有和我同心.
但有兩個提摩太.
以巴弗提.
他都是不求自己的事情.
求耶穌基督的事情.
所以他是我知己.
所以我會和太太說.
他的問題.
他問我.
問我有沒有知己.
這個問題.
其實是問另一個問題.
就是他問.
他當不當他是我知己.
如果是說.
旅行的興趣.
選擇.
我們下一次的旅行.
就是我們應該結婚.
修年紀念.
我還在決定.
究竟跟他還是他跟我.

$^{1521}$這方面很難做知己.
但我會用普羅的想法.
我太太是一個很全科音的人.
而我也很支持這方面的工作.
所以我剛才說的.
能夠有機會帶人信主.
特別是比較孤獨的男性.
我已經進入了這個狀態.
所以我不是太期望男生.
是來跳舞.
來團契.
說說自己的人生.
其實他多數都不會在這方面.
去做.
他會在這方面.
去做人生.
他多數都不會在公開講.
他多數都是個別講.
他都想有人跟他聽.
如果有病患.
有探病.
有機會.
特別現在我們的兒女走了.
很多人有兒女離開.
留在香港的父母有多少.
也留在香港的男性有很多.
其實我們有機會.
去聽他們講一講.
但又不要叫他.
講太多他自己.
一開始不要叫他太多團契.
否則他會覺得很煩.
是事實.
我們.
可能.
我也進入了這個狀態.
我們更加.
孤獨的人需要.
人進入他的世界.
我們有職責.

$^{1561}$去進入一個比較.
內斂的世界的男性的心.
這是我們的祈禱.
而.
我印象.
我的人.
很孤獨.
也很想有朋友和知己.
所以到保羅.
雖然他得到神是我的滿足.
其實我也是這樣.
我很早已經準備一個心境.
無論我的兒女或我太太.
我都以耶穌基督為我.
自保.
有一天我要做決定的時候.
我會不捨得對世上的人.
或是市民.
但我一定要跳越.
人之間的糾纏.
進入天國.
進入和神絕對的關係.
對我們.
對人的關係也會好很多.
我們不捨得世上的東西.
就不能夠絕對求到基督.
但我們絕對求到基督.
就好像保羅.
也會有知己.
所以.
我們回教會一段年.
弟兄姊妹.
是一個很強烈的張力.
我們如果想尋求.
弟兄姊妹的溫暖.
會走火入魔.
我們會覺得.
現在很多人都不回教會.
回到教會的人都沒有什麼.
可以聊天的.

$^{1601}$沒有人可以理解我.
可能回教會就更加孤單.
因為我們和神的關係.
不夠深切.
我們不在乎.
人世間的分離.
有些人離開了.
有些人留在這裡.
有些人留在這裡又不回教會.
其實在這個世界的情感裡.
如果糾纏要.
令到大家關係很好.
就解決不了.
現在很多上一代在香港.
下一代走了.
其實人世間有很多的分離.
情感的事情.
我們作為基督徒應該有怎樣自處.
我們和神有絕對的求.
耶穌基督的事情.
絕對可以滿足到.
我們的心境裡的友誼.
不過我們和神好.
亦都在世上會有.
一些同心的.
為基督耶穌的人.
和我們需要的友誼.
我們可以施出.
所以保羅是一個很獨立的人.
但是保羅亦都是一個.
很有情感的人.
保羅以耶穌基督為.
他的自保.
不過他對一個.
比較陌生的提摩太太.
其實提摩太太是他的乾兒子.
叫做比較熟悉.
但他說提摩太太比提摩太太多.
表示其實他的朋友.
沒有分資深一點的.

$^{1641}$長一點的時間.
只要對方有一個心.
在基督裡面.
為基督來釋放的人.
他都是耳包不提.
所以我覺得保羅的視覺.
是不會因為年齡.
因為日子的長久.
他會選擇做自己.
或者朋友的人.
他不會脫離了和基督的聯繫.
所以基督是我們的自保.
我們都求.
我們能夠成為.
求基督耶穌的人.
我們一起祈禱.
七日天父.
多謝你讓我們在座.
不同的男性或女性.
我們都有不同的.
一些本能.
或者一些不同的.
這個品性.
多謝你讓我們.
能夠有機會看到保羅.
他都很想.
有世上能夠.
可以理解他自己.
一個願意將人生.
囂墊在十字架上.
我們將生命擺上.
以致我們世上有更多人.
成為我們的知己.
我們將弟兄姊妹.
特別是我們有很多長者.
男性的長者.
我們有機會接觸他們的時候.
我們能夠理解.
用他們所能夠.
適應或者能夠.

$^{1681}$創開自己新的心境.
去帶領他們信主.
我們將自己交託.
奉耶穌聖名結求.
謝謝大家收看 拜拜.
\newpage



\section{雅各書 1:12-18}
\label{sec:W1Okagljeug}
\textbf{2024.05.05 《過得勝的生活》 雅各書1\_12-18 (和修版) 曾敏姑娘}
\newline
\newline
連結: \href{https://youtube.com/watch?v=W1Okagljeug}{\texttt{https://youtube.com/watch?v=W1Okagljeug}} ~~~~ 語音日期: 2024-05-06
\newline
\newline
\hyperref[sec:bW5i5_FWoTI]{\small{< < < PREV SERMON < < <}}
~
\hyperref[sec:index_chronic]{\small{[返順時目]}}
~
\hyperref[sec:index_scriptual]{\small{[返順卷目]}}
~
\hyperref[sec:Q94i4NHF3WY]{\small{> > > NEXT SERMON > > >}}
\newline
\newline
雅各書 1:12-18
\newline
\begin{longtable}{cl}
\hline
\hline
章節 & 經文 (和合本修訂版)\\
\hline
1:12 & \begin{tabularx}{0.7\textwidth}{X} 忍受試煉的人有福了，因為他經過考驗以後必得生命的冠冕，這是主應許給愛他之人的。 \end{tabularx} \\ \\ \relax
1:13 & \begin{tabularx}{0.7\textwidth}{X} 人被誘惑，不可說：「我是被神誘惑」；因為神是不被惡誘惑的，他也不誘惑人。 \end{tabularx} \\ \\ \relax
1:14 & \begin{tabularx}{0.7\textwidth}{X} 但每一個人被誘惑是因自己的私慾牽引而被誘惑的。 \end{tabularx} \\ \\ \relax
1:15 & \begin{tabularx}{0.7\textwidth}{X} 私慾既懷了胎，就生出罪來；罪既長成，就生出死來。 \end{tabularx} \\ \\ \relax
1:16 & \begin{tabularx}{0.7\textwidth}{X} 我親愛的弟兄們，不要被欺騙了。 \end{tabularx} \\ \\ \relax
1:17 & \begin{tabularx}{0.7\textwidth}{X} 各樣美善的恩澤和各樣完美的賞賜都是從上頭來的，從眾光之父那裡降下來的；在他並沒有改變，也沒有轉動的影兒。 \end{tabularx} \\ \\ \relax
1:18 & \begin{tabularx}{0.7\textwidth}{X} 他按自己的旨意，用真理的道生了我們，使我們在他所造的萬物中成為初熟的果子。 \end{tabularx} \\ \\
[1ex]
\hline
\hline
\end{longtable}
$^{1}$弟兄姊妹平安.
我們一同禱告.
天父上帝是何等的美好.
我們可以在神的家裡一同去敬拜讚美你.
這是我們不配得的.
是你百百次給我們的恩典.
要賜下你的話語.
成為我們腳前的燈和路上的光.
心願今天你對我們每一個說話.
也都很願意聖靈在我們心裡動功.
讓我們明白又去回應.
奉耶穌基督的名字祈禱.
阿們.
其實我們回到教會裡.
不單要我們去敬拜.
我們有不同的聚會.
我們很希望我們的生命是一個得勝的.
一個這樣的生命.
因為這是神的心意.
在這裡要去得勝.
去榮耀我們的神.
究竟怎樣去過一個得勝的生活呢.
今天很想透過這段經文.
和大家一同去學習.
亞國書這段經文很寶貴.
我們回到PowPoint裡去看看.
得勝生活的秘訣究竟是什麼呢.
首先這裡這樣說.
他說得勝的人有什麼特色呢.
就是能夠忍受試煉.
這裡說忍受試煉的人有福.
其實這些人就是得勝的人.
不是說我遇到試煉.
我就是一個有福的人.
而是你遇到我們遇到.
在裡面我們要去勝過.
我們是要經得起考驗.
這樣我們合格了.
才算是有福的人.
可以說我們是得勝的人.

$^{41}$在這裡我們去看看.
這裡說的試煉.
究竟是在說試煉還是試探呢.
兩個是有一點點的分別.
試煉可以這樣說.
是來自於神.
給我們一些考驗.
試探有些人嘗試分開.
就是失彈.
給我們一些攻擊.
一些的搞擾.
在這裡忍受試煉.
是綜合了的.
神給我們的考驗.
也可以是失彈那邊.
給我們一些攻擊.
我們經過了考驗合格.
我們有得勝的時候.
我們就是一個有福的人.
怎樣去忍受試煉呢.
首先忍受試煉這四個字.
很值得我們去研究一下.
什麼叫忍受呢.
是有忍耐的意思.
是忍耐.
不是說我現在經歷一個困難.
一個患難.
一個考驗.
我馬上就希望神.
就是馬上.
馬上這一刻你就挪走.
你馬上就幫我解決它.
改變它.
就很好了.
我需要的是馬上.
我不需要任何半刻的忍耐.
這個就不是忍受試煉.
忍受試煉強調的是忍耐.
這段時間我一直都在思考.
忍耐這件事.

$^{81}$原來我發覺是很不容易做的.
譬如對於我自己來說.
就很有難度了.
是忍耐.
我的性格是怎樣呢.
就可以跟大家分享一下.
例如如果我有機會去看韓劇.
我去追劇的話.
我怎樣看呢.
看了頭一段.
都挺好看的.
你知道韓劇很長的.
一百集.
看了頭一段之後.
說到男女主角.
說到當中的人物經歷很多風浪.
很多風波.
就很想追了.
通常有些人就很有耐性的.
一直下.
逐集逐集追.
一百集都可以追到後面.
我就已經沒有耐性了.
馬上就拿大結局來看.
我由一開始追到大結局.
快點看下結局是怎樣.
有時候我就想.
可能我就是那種人.
一開始可能自己都是這樣.
經歷風浪.
很想在神那裡馬上就知道.
怎樣解決了它.
怎樣可以處理到.
最好你馬上.
好像看電視劇一樣.
會的.
但神要我們忍耐.
我又看回.
我自己的性格.
真的不容易.

$^{121}$這個對我來說是挑戰.
可能不是每個弟兄姊妹都知道.
其實現在我還在看《夕衣》.
因為大概是.
差不多一年前.
有人就介紹我.
不如你去看下《棚骨》.
我經常都背重物件的.
還有一些不好的動作.
一些壞習慣.
所以《棚骨》其實是有不同的問題.
有些歪了.
這裡那裡.
自己都感受到不舒服.
於是有人推介我去看《夕衣》.
《夕衣》初期會看得比較密.
做到差不多.
各方面好像好一點.
可以放鬆一點.
不用看得太密.
譬如由逐個星期看.
逐個星期.
一個星期看兩次.
到一個星期看一次.
到一個月看一次.
一直看.
我的性格就沒什麼忍耐.
然後我就想.
哎呀.
看到什麼時候呢?.
我很想快點再拉開一點距離.
很想快點.
最好不用看了.
其實就由這些細節看.
我自己是一個不容易忍耐的人.
但是在這裡.
得性的人就要忍耐.
不能快.
兩下就看大結局.
兩下就好了.

$^{161}$最好馬上就好了.
每一天都行神蹟.
我從來不否認神有能力行神蹟.
但神有時就是要我們去經歷.
要忍耐.
忍耐是很重要的一件事.
是對神有信心.
你信得過.
無論神說馬上幫你.
一天幫你.
一個月幫你.
一年才幫你.
十年才幫你.
祂都是愛我們.
祂都有祂的能力.
祂有祂的方法.
祂的時間表.
祂的安排是最好的.
神是不斷的都要跟我說這件事.
神都是要跟你說這件事.
要對祂有信心.
這樣才能夠忍耐.
而經文裡說.
是主應許給愛祂之人.
對神有信心.
亦都是愛上帝的表現.
我信上帝.
每天我們都讀《使徒信經》.
信這樣信那樣.
是否真心地信.
信得來背後是否有一份對上帝的愛.
愛祂.
愛祂就不一樣.
我經常都挑戰弟兄姊妹.
一樣東西.
你真正去愛一個人.
你對他的回應是不一樣的.
你可以很理性去分析很多事情.
你很理性去說一些話.
背後可以不帶著愛的.

$^{201}$但你愛的時候.
那份信任是百分百.
真正的愛是這樣.
真正的信心亦都是這樣.
何況這個上帝是創天造地的上帝.
今天的詩歌是很多敬拜讚美的詩歌.
在裡面我一路去唱的時候.
一路看到.
詩歌裡面不停地讓我們去.
看回上帝的奇妙.
弟兄姊妹.
那些詩歌是否都是你的心聲.
我們都是真正地信上帝.
真正地愛祂呢?.
我現在很想說一個人物.
蘇欽培.
我想問在座有多少人聽過他的名字?.
舉手.
都有一些弟兄姊妹.
證明我們是同一個年代成長的.
我們蘇欽培.
你看回那個日子.
其實他在世的日子是不長的.
是不是?.
大家有數得計.
他在世的日子是四十多年.
他在生命裡面很重要的一個生命表現.
他的見證就是.
只有祝福.
沒有咒助.
蘇欽培的生命是否很強壯呢?.
不得了.
是不是?.
不是.
正正相反.
他的身體不強壯.
但是他留下了一個名言.
就是與其咒助黑暗.
不如燃燒自己.
這個名言.

$^{241}$而燃燒自己.
被人看到的是.
只有祝福.
沒有咒助.
在四十多年他短暫的人生裡.
他在很多人的生命裡面.
留下了一個很深刻的印象.
他在香港大概三十年代末的時候.
那個時間是出世的.
中學的時候是讀英語.
然後在中學的時候.
他信主.
在信主的同時.
你不要看他那時候很年輕.
他已經決定奉獻他的一生給主耶穌.
給主用.
後來他放棄香港大學.
而入讀師範.
其實以他的資歷.
他是可以入讀香港大學.
為什麼他不讀呢?.
在當時的年代.
他是有一番夢想.
他很想去教育.
從事教育的工作.
他也曾經當了六年小學的教師.
他去教書.
那個想法不是.
我有穩定的工作.
我有很好的一份糧等等.
而他當時對一班貧窮的孩子.
很有負擔.
所以他教書的時候.
也是很想幫助那些貧窮的孩子.
而回看他的一生.
他所追求的.
只是對上帝傳言的信服和委身.
記住這一句話.
他一生所追求的.
不是名不是利.

$^{281}$不是自己的名聲.
不是自己的成就.
只是對上帝傳言的信服和委身.
這是後人對他一個很大的肯定和評價.
他在26歲的時候.
被診斷有甲狀腺癌.
26歲.
非常年輕.
無論什麼年紀患上癌症.
我們都很擔心.
都很自然.
但他在年輕的時候.
26歲就被診斷有甲狀腺癌.
當時他不知情.
因為家人初期都不想告訴他.
後來去到七年之後.
他的癌症又復發.
這個時候他就知道了.
他覺得很晴天霹靂.
原來我有甲狀腺癌.
但晴天霹靂只是感到驚訝.
他沒有怨天尤人.
他沒有埋怨.
他也沒有感到灰心,沮喪.
也沒有覺得人生這樣就算了吧.
或者反過來.
為什麼我這麼年輕就這樣呢?.
上帝你對我真是不公平了.
又不見你對其他人這樣.
我人生這麼年輕.
有很多美好前途.
他都沒有這樣想.
他繼續信服神.
繼續委身在神裡面.
他放下教學的工作.
他去進修.
其中的進修.
他去了無敵聖經學院.
學習一些聖經方面的.
神學方面的東西.

$^{321}$接著在好一些時間裡.
他展開了在學生當中的工作.
也開創了好幾本雜誌刊物.
他對年輕人,少年人.
是有很大的負擔.
因為他看到當時的少年人.
陷於罪惡當中.
受到世俗很大很大的影響.
所以他很想為這群年輕人.
去做一些事情.
他就將自己擺上.
曾經創辦了不同的刊物.
也是突破這個機構.
其中的創辦人之一.
但又不要看他.
身子好像看上去很軟弱.
很多的長年受到疾病的纏繞.
看上去好像一個弱質女子.
但生命卻是閃閃發光.
是為上帝而發光.
可以為神做很多的工作.
神透過他.
不單是創辦那些刊物.
也在青少年的運動裡.
帶來很深遠的影響.
你真的沒想過.
神可以這樣去用一個人.
去到1970年的時候.
他就積怒成疾.
他的身體崩潰了.
太累了.
也是因為疾病的緣故.
他要放下他的工作.
曾經有一段時間.
他跟死亡很近.
即使是這樣.
仍然沒有埋怨.
他的生命經常說的是.
只有祝福.
沒有咒詛.

$^{361}$他的口是這麼說.
他的生命也是這樣表現.
他認定自己的生命是這樣.
所以他寫了一篇文章.
叫做只有祝福.
帶來很大的迴響.
是祝福很多人的生命.
這個不容易說.
今天我們在平順的日子.
我們可以很容易去祝福人的生命.
是說你要有信心.
我們有信心.
我們要克服困難.
但在癌症當中.
在長期病患當中.
可以這樣說.
不是每個人都做得到.
他就做得到.
他身體稍有改善.
他又做事了.
為誰做呢?.
為神而做.
將生命交在神的手裡.
他的一生就是這樣.
是非常奇妙.
對於他來說.
我相信這個長期的患病.
就是一個試煉.
一個很大的考驗.
他不可以像其他人一樣健健康康.
很強壯的.
在這樣的狀態下服侍神.
就在這麼軟弱的裡面.
就是這麼多限制裡面.
很辛苦的去服侍神.
但在試煉裡面.
他是怎樣的?.
他是得勝.
他是得勝的.
他沒有被拖垮的.

$^{401}$我們看看一個試探的例子.
大家猜猜這幅畫是說誰呢?.
沒錯,是說藥室.
藥室在試探裡面.
是一個很大的得勝.
大家記得.
他被哥哥們賣掉了.
賣了之後.
人們輾轉將他帶到埃及這個地方.
去到波提佛的家.
主母天天對著他.
結果就引誘他.
很想要他犯罪.
在這個時候.
你嘗試代入一下.
藥室的生命裡面.
他深深的受傷.
受家人的傷害.
家人賣掉他.
家人那些哥哥們很妒忌他.
對他這麼大的傷害.
他要離鄉別井.
離開父母.
去到一個陌生的地方.
而且要圍爐.
在這裡又面對一個這麼大的試探.
在過程中.
換作有些人.
很容易就會掉進陷阱犯罪.
因為主母不是一次性的生他氣.
你跟我一起吧.
跟我一起不是一次.
說完一次又一次.
又找不同的機會.
甚至有一次.
還拉著他的衣服.
在這個時候.
在試探當前.
一個年輕人.
不容易.

$^{441}$但是你看到.
藥室.
藥室得勝.
勝過試探.
在裡面.
一會兒我們看到.
有些元素.
藥室的得勝出於什麼.
為什麼可以在這樣的狀況.
我可以拒絕.
不犯罪呢.
敬畏神.
信服神.
將神擺在首位.
在這個時候.
拒絕那惡者的引誘和試探.
所以他得勝.
他的人生.
都是一個得勝的人生.
無論是試煉的.
無論是試探的.
都能夠得勝.
是何等的寶貴呢.
在這裡經文有說到.
忍受試煉的獎賞.
原來我們過得勝的生活.
神會獎賞我們.
那個獎賞是什麼呢.
在十二節這樣說.
因為他經過考驗以後.
必得生命的冠冕.
神會將這個永生的生命.
賞賜給得勝的人.
是何等的寶貴.
這是主應許給愛他之人.
有獎賞的.
弟兄姊妹.
我們是否渴望得到神的獎賞呢.
我最近多了和小朋友.
一起在主教學中.

$^{481}$和他們一起生活.
一起敬拜.
我看到小朋友很單純的心.
他們去聽聖經故事.
去聽分享.
他們是很用心的.
心也很單純.
他們很渴望得到禮物.
這是好的.
我經常覺得神很愛他們.
事實上聖經也有這樣說.
神很愛這些小孩子.
神也很樂意去獎賞他們.
小朋友很渴望得到獎賞.
無論是零食.
是一些圖畫.
是什麼特別的禮物.
其實我相信我們也渴望得到禮物.
我相信我們渴望得到獎賞.
更何況這個獎賞來自神.
來自神的獎賞和人是不一樣的.
地面上人的所有賞賜.
其實也會過去.
不會去到永恆.
但上帝的獎賞是不一樣的.
是傳到永恆.
這個價值是何等的大.
在這裡說我要得到神的獎賞.
是多麼寶貴.
神給的獎賞也是永恆的.
是生命的冠冕.
是永生的冠冕.
是很寶貴的.
但要得勝才可以得著.
是要得勝.
這是非常重要的.
在這裡有些問題很重要.
我們需要處理.
過得勝的生活.
有些事情我們需要認識清楚.

$^{521}$究竟是怎樣造成的.
得勝在我們面前擺著一個很大的難阻.
是罪惡.
罪惡令我們失敗.
是罪惡的問題.
我們要問一個問題.
我們被罪惡誘惑的原因是什麼.
需要一同了解.
因為很多時候在試探.
在試煉裡.
很多時候我們會將原因放在神身上.
例如我身體長期不適.
或者我經歷很大的患難.
第一時間就說.
神是你.
在背後做了什麼.
為什麼令我受到這麼大的痛苦.
為什麼你不幫我.
是不是你在背後做了什麼.
令我跌倒.
我們經常很容易怪罪於上帝.
在這裡我們要了解清楚.
究竟被罪惡誘惑的原因是什麼.
這個你見到的人頭上有什麼.
有一個烏鴉.
這個字比較小.
這個烏鴉原來是在說試探.
在經文裡很詳盡的.
很真實的說.
為什麼會被罪惡誘惑.
這個罪惡影響我們.
不能過一個得性的生活.
為什麼.
13節這樣說.
人被誘惑.
不可說我是被神誘惑.
因為神是不被惡誘惑的.
他也不誘惑人.
我們經常怪是神在背後做了什麼.
令我跌倒.

$^{561}$我的人生就這麼坎坷.
我們每件事都不好.
在這裡其實要說.
你說到罪惡的層面.
弟兄姊妹我們要小心.
我們要看清楚一件事.
其實我們陷入罪惡裡面.
不是因為神的緣故.
第一 神自己不被罪惡誘惑.
我們的上帝是聖潔的神.
和罪惡劃清界線.
神很憎惡罪惡.
所以神自己一定不會被罪惡誘惑.
他也不會拿這件事誘惑我們.
以致我們犯罪.
所以我們在罪惡當中.
更加不應該將那個球拋給上帝.
上帝是你造成的.
不是 不是這件事.
在這裡你看到進程.
是怎麼起的.
為什麼我們開始犯罪.
首先是因為我們裡面有私欲.
我們有私欲.
我們有慾望.
有不好的念頭.
在我們裡面.
先在我們裡面.
慢慢慢慢.
你看看經文怎麼說.
私欲既懷了胎.
它一路成長.
一路累積.
一路成長.
好像嬰兒在媽媽的肚腹裡面.
起初慢慢結合一個胚胎.
它一路成長.
一路成長.
越來越大.
成了一個胎.

$^{601}$接著一路下去.
其實私欲也是這樣的.
好像嬰兒成長的過程.
私欲懷了胎.
它長大了.
它就生出罪來了.
是這樣開始的.
有時候你想想.
人們怎樣走上.
去偷東西.
去搶東西的地步.
可能有些人.
因為他們生活的貧窮.
有需要.
也不是每個人因為生活貧窮.
有需要會去偷去搶.
有些人純粹是因為貪.
有一次.
有個人說了一個例子.
我覺得值得我們.
真的引以為鑑.
有一次有間酒樓.
有一段時間發覺.
為什麼整整一整杯子都沒有了.
整整一整杯子都沒有了.
奇怪.
那些茶樓的杯子.
為什麼會不見了.
誰拿了去.
慢慢去調查.
發覺是一個婆婆.
一個婆婆長期習慣性去偷杯子.
拿回去了解一下.
跟婆婆聊天.
甚至約婆婆的家人聊天.
看看家裡是不是有很大的經濟需要.
為什麼要偷茶樓的杯子.
結果你猜發現什麼.
不是.
婆婆家裡經濟挺富裕.

$^{641}$是什麼問題呢.
為什麼要去茶樓拿杯子呢.
說真的.
這麼多人用過.
用完又用.
很損耗.
又不是說圖案特別漂亮.
造型美.
為什麼呢.
原來就是裡面那份貪.
裡面有一份私慾.
總之我要拿一些東西.
拿完一個又拿一些東西.
總之裡面是一個無底洞.
要不停地拿.
拿了就很滿足.
所以大家想了很久.
都不明白為什麼.
就是這個原因.
這樣去偷.
就是這個東西.
在裡面的概念.
慢慢地懷胎.
成為一個實際.
行動上的罪.
說真的.
首先發生的罪.
在這裡發生.
在我們的腦海裡.
先有個貪念.
但這個貪念是怎麼來的呢.
是不是.
真的值得再去.
再跟這個婆婆談多一點.
就是她裡面這件事.
會生出一個罪來.
當這個罪一直長盛.
出現了.
原來罪帶來的後果是什麼.
生出死.

$^{681}$後果非常嚴重.
就好像這幅圖畫.
起初有一個屍肉在我們裡面.
試探出現.
好像一個烏鴉.
在我們的頭上.
接著就越來越大.
越來越大.
成為罪惡.
一直腦海裡的罪惡.
成為實際的行動出來.
就好像這幅圖一樣.
上面就開始抓個巢.
罪既長盛.
就生出死來.
巢之後還要變成什麼呢.
這個.
是不是.
罪既長盛.
接下來的就是死亡.
是何等的可怕呢.
這個死.
說的是什麼.
是永遠的死.
永永遠遠.
和我們的神分開.
在神裡面我們沒有份的.
我們在救恩裡面.
我們在神裡面沒有份.
我們在神的國裡面沒有份.
是何等的可怕呢.
在永恆的黑暗裡面.
是何等的可怕.
在這個過程.
作者要提醒我們.
提醒我們.
有些人就說.
我犯罪.
為什麼犯罪.
一定是和神有關的.

$^{721}$一定是神將這些東西.
放在身邊.
引誘我的.
為什麼我會妒忌呢.
因為隔壁的女士.
神為什麼派這麼漂亮的女士.
在我身邊.
相比之下.
我顯得沒有這麼漂亮.
就是神因為你放她在我身邊.
所以我就妒忌.
為什麼我會出貓.
考試.
就是因為神給這麼厲害的學生.
在我旁邊.
不出貓就對不起自己.
出言是不信的.
說話也是大聲的.
總之是可以賴就賴.
都是神不好.
你不放這些東西在我身邊.
你不做這些.
我就不會這樣.
為什麼你現在要做這些東西.
這麼好.
我沒有份.
我也很想要.
所以我就犯罪.
去偷去搶.
去不同的東西.
將這些東西賴給神的時候.
這對神是極大的羞辱.
極大的不尊重.
經文這樣說.
神是第一.
自己不會被惡誘惑.
第二.
他一定不會這樣誘惑人犯罪.
他絕對絕對不會這樣做.
我們要認清楚.

$^{761}$我們的神是一個怎樣的神呢.
在這裡這樣說.
他說各樣.
17節這裡這樣說.
各樣美善的恩澤.
和各樣完美的賞賜.
都是從上頭來的.
從重光之父那裡降下來的.
說到這個很美善的恩澤.
是真的很善良的.
說這些好的地方.
是很善良.
是很好的東西.
是非常好.
對人有益的.
神給我們的東西是對我們有益的.
是良善的.
是好的.
而且這裡說到什麼呢.
是各樣完美的賞賜.
不是一般的賞賜.
是完美的.
一點瑕疵都沒有.
神的賞賜是這麼好的.
神給我們的東西.
都是這麼好的.
我們的神是怎樣的神呢.
是重光之父.
是光明的.
神在這個世界.
做了不同的光體.
有日月星神.
是成為我們的光.
去照著我們.
神有位處置那些光.
讓那些光繼續照著我們.
但是神自己就是我們的光.
其實神自己是我們的光.
說到重光之父.
由神那裡來的.

$^{801}$這個賞賜是何等的寶貴.
弟兄姊妹我想說.
神不單止不會引誘我們犯罪.
神是將美善的.
好的完美的東西給我們.
神讓祂的光去光照我們的生命.
是何等的好呢.
世界會改變.
但是神裡面沒有改變.
那些光體呢.
你看看太陽每一天.
都會照著我們.
照著我們的時候.
那些影子都有轉變.
在不同的時刻.
那些太陽照射的影子.
都有改變.
但是神裡面沒有改變.
祂由起初就是良善.
就是完美.
就是聖潔.
就是光.
祂以前是這樣.
祂現在是這樣.
祂永遠都是這樣.
祂都不會改變.
祂都會將好的東西給我們.
這個才是神.
那我們自己為什麼會犯罪呢.
說到尾.
我們犯罪是因為我們自己的私慾.
我們要承擔的責任.
是我們的私慾.
我們的貪心.
我們的自私.
我們裡面很多的比較.
自卑.
很多的嫉妒.
是我們裡面自己的生命的情況.
影響到我們犯罪.

$^{841}$去得罪神.
是我們自己.
我們不妨回想一下.
這一切是怎麼來的呢.
既然我們自己才是犯罪的真正原因.
在過程裡面可以怎麼做.
怎麼去勝過試探.
怎麼去過一個得勝的生活呢.
這裡是很重要的.
有兩段經文想跟大家分享.
所以要信服神.
要抵擋魔鬼.
魔鬼就必逃避你們.
要親近神.
神就必親近你們.
沒錯.
烏鴉開始在我頭上想搗亂.
我裡面的私慾對我有影響.
私貪來到了.
撒旦也很會找位置.
怎麼會無緣無故就這樣去試探呢.
一定是看到我們生命裡面.
有些破口.
有些位置.
有私慾.
他就試探一下.
撩一下你 動一下你.
看你會不會承認.
如果我裡面看中一些東西.
很貪婪.
撒旦試探.
拿吧 拿吧.
用什麼方法拿 偷搶都好.
用不同的方法.
我很喜歡講一個例子.
就是你不要說.
中學生都學會一些方法.
滿足自己的慾望.
我曾經跟一些人分享.
譬如中午回學校.

$^{881}$怎麼可以不花錢都吃得到東西.
每天中午.
回到學校怎麼可以.
自己不花錢都可以吃得到東西.
很簡單.
Liz 你今天帶飯.
帶什麼.
魚蛋很好 拿一顆來.
拿來吃一顆.
今天你家.
徐太.
做了些炒米粉.
拿些炒米粉過來.
兩口 兩口拿一些過來.
於是那餐飯.
這裡拔一點 那裡拔一點.
有個學生真的學會.
每餐拔別人的飯糊過來.
拿一些 結果他.
不用帶錢包出街.
我不用回學校買東西.
我拿同學的東西吃就行了.
這個後來當然.
我也有教導他.
其實這個過程就是.
在我裡面的私慾.
想一些方法 是不是.
其實我們不是這樣生活.
你說是不是他完全沒有錢.
所以我就想這個方法不是.
這個方法令自己很滿足 很快樂.
我不需要自己出錢.
我可以滿足到我的慾望.
但這個是非常自私.
我們也是犯罪 得罪上帝.
小時候已經找到方法這樣做.
你長大了更加不得了.
在過程中.
開始有些怎樣.
怎樣成為親近神.

$^{921}$中學生需要上帝.
青年人需要神.
成年人需要神 長者需要神.
兒童需要神 每一個都需要.
每一個我們都有我們的慾望.
有我們的私慾.
在裡面是要信服神.
信服神 神教我們的是甚麼.
神教我們的不是自私.
從別人身上得到自己所滿足的.
神教我們的是愛倫如己.
你愛別人好像愛自己一樣.
施比受更為有福.
你付出吧.
付出的人是有福的.
在過程中 神教我 我信服他.
撒旦你退去.
我今天要得到這些東西.
我有正當的渠道.
正當的方法可以得到.
我不需要透過犯罪的方法去得到.
現在好像別人比我漂亮.
別人比我厲害.
別人有甚麼好.
撒旦你退去.
我不會去嫉妒.
上帝愛我.
我怎樣 上帝都愛我.
我不需要做這個世界.
全世界第一的人.
上帝才愛我.
我是感神就愛我.
我就滿足了.
不需要和別人比較.
不需要妒忌.
這是很重要的.
抵擋魔鬼.
跟我說你退去吧.
你不要以為你這樣可以搞到我.
我們要跟自己說.

$^{961}$魔鬼就會逃避我們.
還要親近神 神就會親近我們.
長期不親近神.
撒旦看到.
來來來 我和你做朋友.
散播負面的思想給你.
讓你非常不開心.
很多比較 很多不開心.
讓你經常都想犯罪.
等你做我的徒弟.
我們親近神.
撒旦會退去.
神會親近我們.
神會幫助我們.
在這裡是何等寶貴.
另一段經文也是很寶貴.
我說你們要顧著聖靈而行.
絕不可滿足肉體的情慾.
因為肉體的情慾和聖靈相爭.
聖靈和肉體相爭.
這兩個是彼此敵對.
使你們不能做所願意做.
在我們信耶穌之後.
神讓聖靈住在我們當中.
讓我們信服聖靈的大靈.
聖靈會用微弱的聲音.
對我們每一個說話.
我們聽了要去回應.
不要去壓低聖靈的聲音.
也不要去被自己的心梗.
不去聽聖靈的提醒.
這是很重要的.
我們需要住的地方.
當我們能夠這樣做的時候.
這個烏鴉.
這個試探的烏鴉.
沒有位置可走.
神自然會幫助我們.
這一個得性的渠道.
是很重要的.

$^{1001}$在這裡很想鼓勵弟兄姊妹.
鼓勵弟兄姊妹.
一起過得性的生活.
一起在這裡.
不是說我信了耶穌就夠了.
什麼是真正的信.
真正的信.
是我們整個生命.
要屬於耶穌基督.
耶穌基督是得性的.
我和你都要過得性的生活.
心願我們聽神的話.
聽聖靈的提醒.
順服他們而行.
其實蘇欽培的一生.
是一個很美好的見證.
他比我和你更加艱難.
在一個這麼大的疾病裡.
他仍然選擇順服神.
委身於神.
他沒有埋怨.
他的生命只有祝福.
沒有就坐.
他都是這樣鼓勵人.
弟兄姊妹.
可能我們的生命裡.
都有很多困難.
有很多挑戰.
我們不否認.
我們有困難.
有挑戰.
但我和你都可以得性.
不只是蘇欽培的人生可以得性.
我和你都可以.
心願我們的生命.
能夠榮耀上帝.
有一天見到神的時候.
我們是無愧於心.
而且我和你一起.
可以得神的掌上.

$^{1041}$我們一同去禱告.
藉著這段經文.
給我們提醒.
神我很相信.
在我們裡面.
我們每個人都有我們的私慾.
可能是不一樣的.
但在當中求你給我們.
能夠勝過試探.
讓我們更多的親近你.
更多的去信服你的教導.
信服聖靈的提醒.
我們要上上宣告.
我們的生命是屬於耶穌基督的.
要讓惡者退去.
祂不能去試探我們.
讓我們的生命在人群當中.
能夠作美好的見證.
就好像蘇欽培姐妹.
所做的見證一樣.
是那麼美好.
又或者好像昔日的約瑟一樣.
能夠這樣去勝過一切的試探.
主求你幫助我們.
奉耶穌基督的名字祈禱.
阿們.
\newpage



\section{箴言 31:1-9}
\label{sec:Q94i4NHF3WY}
\textbf{2024.05.12 《利慕伊勒王的真言》 箴言31\_1-9 (和修版) 講員 : 謝靄薇牧師}
\newline
\newline
連結: \href{https://youtube.com/watch?v=Q94i4NHF3WY}{\texttt{https://youtube.com/watch?v=Q94i4NHF3WY}} ~~~~ 語音日期: 2024-05-13
\newline
\newline
\hyperref[sec:W1Okagljeug]{\small{< < < PREV SERMON < < <}}
~
\hyperref[sec:index_chronic]{\small{[返順時目]}}
~
\hyperref[sec:index_scriptual]{\small{[返順卷目]}}
~
\hyperref[sec:mgPJ9gnhbtQ]{\small{> > > NEXT SERMON > > >}}
\newline
\newline
箴言 31:1-9
\newline
\begin{longtable}{cl}
\hline
\hline
章節 & 經文 (和合本修訂版)\\
\hline
31:1 & \begin{tabularx}{0.7\textwidth}{X} 瑪撒王利慕伊勒的言語，就是他母親教導他的。 \end{tabularx} \\ \\ \relax
31:2 & \begin{tabularx}{0.7\textwidth}{X} 我兒，怎麼了？我腹中生的兒，怎麼了？我許願而得的兒，怎麼了？ \end{tabularx} \\ \\ \relax
31:3 & \begin{tabularx}{0.7\textwidth}{X} 不要將你的精力給婦女，也不要有敗壞君王的行為。 \end{tabularx} \\ \\ \relax
31:4 & \begin{tabularx}{0.7\textwidth}{X} 利慕伊勒啊，君王不宜，君王不宜喝酒，王子尋找烈酒也不相宜； \end{tabularx} \\ \\ \relax
31:5 & \begin{tabularx}{0.7\textwidth}{X} 恐怕喝了就忘記所頒的法令，顛倒所有困苦人的是非。 \end{tabularx} \\ \\ \relax
31:6 & \begin{tabularx}{0.7\textwidth}{X} 可以把烈酒給將亡的人喝，把酒給心裡愁苦的人喝， \end{tabularx} \\ \\ \relax
31:7 & \begin{tabularx}{0.7\textwidth}{X} 讓他喝了，就忘記他的貧窮，不再記得他的苦楚。 \end{tabularx} \\ \\ \relax
31:8 & \begin{tabularx}{0.7\textwidth}{X} 你當為不能自辯的人開口，為所有孤獨無助者伸冤。 \end{tabularx} \\ \\ \relax
31:9 & \begin{tabularx}{0.7\textwidth}{X} 你當開口按公義判斷，當為困苦和貧窮的人辯護。 \end{tabularx} \\ \\
[1ex]
\hline
\hline
\end{longtable}
$^{1}$各位大英小英朋友,很久不見了.
今天是母親節,祝各位在座的母親節快樂,身體健康,能夠永生一人.
我們在看神話之前,一起去祈禱.
每有恩典,每有慈愛的主,當我們能夠來到你面前,何等的福分.
求主在這時候去潔淨我們的心思意念.
有什麼是不蒙你喜悅的,求聖靈眼當中去剔除.
以致我們能夠專心一意地進入神話裡.
求主親自教導我們,讓我們能夠所做的合乎你的樣式.
又願孩子心中的意念,口中的言語,蒙你如納,奉主耶穌基督命求,阿們.
我們剛才讀了利牧伊拉王的真言.
到底利牧伊拉王是誰呢?.
我們對他知道的不是很多.
只不過我們知道他是一個君王.
而且他從母親那裡學到一些東西,得到智慧.
他的名字叫做奉獻給神的意思.
有些人認為他和前面那張阿古爾.
可能都是同一個地方,在阿拉伯北部的馬薩國.
都是同類的人.
到底母親對兒子有什麼期望呢?.
當你看經文的時候,你一開始會看到.
剛才讀的聖經裡,說到關於君王的職責.
他告訴你,是由他母親那裡學到的.
所以你想到一個母親,特別是國王的母親.
在皇宮裡一定有他的影響力.
所以他對兒子有很大的期望.
你看到他充滿智慧地教導兒子.
在前面兩節裡,他不斷地說.
我兒啊,怎麼了?.
我福中生的兒,怎麼了?.
我許願而得的兒,怎麼了?.
你會看到他在當中對兒子很有期望.
表明他是愛兒心切.
所以他教導他的時候.
你會看到當母親的,他很有權力和責任去教導兒子.
而兒子就希望兒子聽他的話.
兒子也有責任和義務去聽從.
因為他母親懷胎十個月,這麼辛苦生出來.
他在這麼艱難的情況下.
他母親這麼愛他,是沒有假的.
就算他罵他,也不會有惡意.

$^{41}$他母親這麼愛他,並不是有目的.
不求回報,在家的母親.
你知道他愛兒子,他生病的時候最辛苦.
你沒想到,日後要好好孝順他,對他好一點.
不會的,很辛苦的,多麼辛苦的.
要養大子女.
而你看這段經文的時候.
繼續看三四五六.
你會看到我畫的那些叫你不要做的事.
他教導他的兒子不要做一些敗為君王的行為.
還有叫他不要喝酒.
其實君王擁有三宮六院是很平常的事.
但你會看到聖經裡.
這個所羅門王因為妃嬪太多.
他隨心所欲.
所以你會看到他跟著妃嬪去拜他們的偶像.
反而忘記了這位向他顯現的耶和華神.
他忘記了這就是他一生裡最失敗的地方.
喝酒,其實君王怎會不喝酒.
但這個很明智的母親.
教導他喝酒會使人亂性.
顛倒是非,失去公義.
所以由此我們想到.
其實神也給我們人有自由意志.
但叫我們不要把這個自由.
當作放縱情慾的機會.
所以這個自由在神的角度裡.
他怎樣教我們的呢.
當你看神的話語的時候.
你就會明白神叫我們做一件事.
叫我們自由.
今天每個人來到神面前.
心靈都有自由.
但當你有自由去做很多事的時候.
叫你選擇的時候.
你會不會有障礙呢.
很多東西讓你選擇的時候.
但如果你是一位熟悉神的人.
你一定會按照神的標準.
這個世界雖然價值觀有很多不同的尺度.

$^{81}$但當你是熟悉神的人.
你要選擇自由的時候.
你也不是隨便去選擇的.
你看我們在聖經裡.
一些好的君王.
行神,看為正的人.
他背後應該有一個好的母親去教導他.
母親在家庭裡扮演的角色是很重要的.
因為她和子女相處的時間很長.
比父親還要長.
所以你看到他.
媽媽做什麼他都會學他.
所以他很像他的母親.
所以媽媽的一舉一動影響了他一生.
那怎樣做一個稱職的母親呢.
相信在座的為人母親的.
你們很有資格很有經歷.
不過我以一個子女的身份.
去看一看一個做一個母親的.
其實他應該有一個好的榜樣.
新教還要現教.
你的一舉一動.
隨時你的子女都在模仿.
你忽然間會聽到一些小朋友出口驚人.
他們說的話其實是由大人去學的.
你有沒有覺得.
所以我今天回到五方面去想一想.
既然要做榜樣.
祈禱我讀到.
加菲德美國的第20任總統.
他的童年意識裡.
感受到母愛真的很偉大.
他自幼就失去父親.
他媽媽帶著一群子女去養大他們.
他心裡很希望他們能夠受到高等教育.
所以小朋友從小就擔起養家的責任.
而加菲德這個總統.
他很喜歡航海.
夢想航海.
所以他找了一份船員的工作.

$^{121}$有一天晚上他就掉進了混河裡.
他不懂得游泳.
他喊救命.
沒有人聽到.
怎麼辦呢.
他很辛苦地掙扎.
終於可以爬回船上.
他坐在甲板上.
他就想起一個很嚴重的問題.
他想到.
今天神救了我.
神救了我.
一定要我對我的母親負責任.
同時他想到一件事.
比航海更重要的使命.
等著他去做.
於是他決心要發奮圖強.
努力讀書.
他就放棄了船上的工作.
回家了.
回到家的時候.
他在窗外看著.
看到他母親跪在那裡為他祈禱.
他很感動.
他說今天神救了我.
就是因為母親為他祈禱.
所以他看到.
母親為子女祈禱有多重要.
原來他的人生觀也會有很大的改變.
關於信仰的問題.
不知你覺得你的小孩.
何時帶他回教會好呢?.
有時小孩都很吵.
小孩那麼吵.
不要帶他回家.
我都不知道坐在那裡.
每個人都看著我那麼吵.
所以有些人就想著.
他長大了再帶他去.
但等他長大了.

$^{161}$他就有自由意志.
他有很多東西玩.
他沒興趣跟你回教會.
所以看到有些母親都很後悔.
或者他們流淚祈禱.
不知怎樣帶兒子回教會.
我們要趁他沒有毛有翅的時候.
快點帶他回來.
在教會裡浸浴.
都有教會的味道.
所以一定要帶他回來.
特別是在紅印堂.
都有兒童崇拜.
不要怕煩.
老師不怕煩.
你帶他來就好.
神自然會教導他.
所以盼望你趁他小時候帶他回來.
其實耶穌基督都叫母親.
把小孩子帶到他面前.
所以在小學的時候.
我父親在寮國工作.
他放暑假.
短短的時間.
才跟我們相聚.
看聖經.
明白神的話.
平時就是我外婆和我媽媽.
帶我們回教會.
上主日學.
在那裡我們學習認識神.
去敬拜神.
這是很好的學習.
就像提摩太的信心.
聖經說是由外婆.
萊爾和母親得到的.
由那裡來的.
所以從小就要讓他去敬拜神.
其實這個功課.
我們幾兄弟妹.

$^{201}$都是受用不淺的.
整個話語.
幫助我們在什麼環境下.
我們都懂得交託給神.
將我們的事情.
有時候.
你一舉一動.
其實弟弟是看著你的.
他在觀察你的.
他的小眼睛.
他會聽你的話.
他也會聽神的話.
所以家庭.
如果家庭有家庭崇拜.
更加好.
你可以去書室.
他可以提供資料給你.
其實家庭.
無論父親帶你去聚會.
或者母親.
都不重要.
最重要是.
讓他一齊去敬拜神.
家庭崇拜.
這都是神所喜悅的.
任性.
現在時代進步了.
物質充裕.
生活富庶.
養成小朋友.
嬌生慣養.
你見到小朋友.
在街上鬧彆扭.
鬧你買東西.
你不買他怎麼辦?.
他會睡著.
不起來.
哭著.
要買到為止.
其實我自小.

$^{241}$家庭很單純.
我媽媽認識的朋友.
都是舊的長輩.
同輩.
其實我們生活很簡單.
平時.
功課讀得好.
有獎勵.
當然.
最多獎勵的不是我.
但我們會珍惜父母給我們的獎勵.
我們不會羨慕.
別人拿到好的物質.
我們不會有這樣的心情.
但我們見到.
小朋友.
他們往往.
帶人送禮物給他.
很快就拆了紙.
然後去玩.
玩了一下就丟在一旁.
或者.
不久後.
他又去找他媽媽再買.
但他好像一點都不會.
覺得.
父母這麼辛苦.
賺錢去買給他.
他從來都沒有這樣想.
自然就好像.
養成他們.
沒有.
我女兒就有了.
很自然.
搞成很多浪費.
難怪現在要環保.
女兒就有了.
這些是不好的習慣.
所以做父母.
做事要有原則.

$^{281}$難怪我們.
小時候看書.
盲目了三千.
無論現在.
居住什麼環境.
你要去觀察.
他交什麼朋友.
他在看什麼書.
可能你要幫他選擇.
因為他還小.
所以不要讓他任意而行.
因為他還不懂分辨.
是非.
今天教育很難嗎?.
今天小朋友教育很難嗎?.
有些做老師的都很害怕.
教育的那些.
都是不聽話的.
但現在有新式的教育.
是愛教育.
愛教育又怎樣?.
不主張去打他.
罵他.
但箴言多次說.
要以仗打.
希伯來斯十二章六節八節.
因為主所愛的.
他必管教.
他又鞭打.
他所接納的每個孩子.
為了管教.
你們要忍受.
這是神的話.
神待你們如同待兒女.
哪有兒女不被.
父親管教的呢?.
管教原是眾兒女.
共同所領受的.
你們若不受管教.
就是私生子.

$^{321}$不是兒女了.
愛者,我們曾有肉身之父.
管教我們,我們尚且敬重他.
何況靈性之父.
我們豈不更.
當信服他.
而得生命嗎?.
其實你短短看經文.
就有.
因為管教不是.
任由他亂來的.
好像有時.
你會發脾氣.
會有憤怒.
可以帶弟弟來.
糟蹋他.
發洩自己的脾氣.
好像《耳法所書》.
教我們做父親的.
不要激怒兒女.
原來的.
話本就是惹怒兒女的氣.
要我們按照.
主的話語去教導他們.
原來.
不管他怎樣.
或者.
你要去管他,要打他.
當然.
愛的教育就是愛.
但要教他.
但愛的成份要在.
你要罵他.
要有愛心在.
為他好,你才罵他.
去責備他.
愛是要平時去培養.
不是你打罵他.
那時你已經很生氣.
怎會有愛呢?.

$^{361}$愛是平時去培養的.
所以你要去跟他聊天.
甚至.
有時他不說話.
你要想他發生什麼事.
其實現在就要想這麼多.
所以做父母是不簡單的.
不是說.
你代溝.
所以要多聊天.
父母又要.
說民主.
但小孩子的階段.
他不知道.
什麼是民主.
你以為拿旗子抗議就行嗎?.
你要教他.
所以孩子的階段.
都要給他一些規矩.
要教導他.
嚴格一點是沒問題.
小時候不教他,長大後就很麻煩.
自找麻煩.
你不知道怎樣教他.
因為他已經有毛有翼.
已經長成了.
到他懂得分辨的時候.
才說民主.
教育.
在家庭裡.
你都希望.
他受高等教育.
要求成績高.
要好一點.
是無可厚非的.
誰不想自己的兒子出人頭地.
但你追求體面的時候.
你不要太過要求.
否則.
有悲劇出現.

$^{401}$人會向高處爬.
你都要看他的能力.
如果他像猴子.
懂得爬樹.
你不懂爬樹.
是否就.
不可以做人呢?.
就像兔子跳.
你不可以叫他爬樹.
每個人的才幹都不同.
每個人的個性都不同.
你不可以.
要求他每個人都是.
這個樣子的.
神造人都有不同的樣子.
所以每個人的個性都不同.
千萬不要和.
隔壁的比.
或者和哥哥姐姐比.
姐姐讀書很厲害.
你不要和她比.
否則她就會很自卑.
我都抬不起頭來.
或者和隔壁的比都不對.
因為是別人的子女.
你不可以和她比.
她做十次對一次.
你都要稱讚她.
只有一次對.
有些負面的說話.
她就沒有人稱讚你.
不是說.
你要怎樣捧她.
但你都要鼓勵她.
鼓勵一下.
幫助她建立一個正確的人生觀.
很重要的.
你不要常常說她蠢.
有時有些中國人.
好像不知為何.

$^{441}$都說自己的子女.
是賤.
有什麼就說什麼.
不要太謙卑.
反而造成.
子女不知怎樣看.
常常說她蠢.
她可能以為.
自己真的蠢.
其實她可能不蠢.
所以我們不妨試試.
去鼓勵一下她.
讓她正面去看事情.
所以.
當我們去看.
尼姆伊達王.
他的言語.
都激勵我們要向神求智慧.
像他媽媽一樣有智慧的教子.
所以.
令他成為.
一個俠義的君王.
我有一位.
很尊敬的長者.
她就是信我愛.
兒童之家的.
王守信師母.
平日人們稱她為王奶奶.
她是.
孤兒院的院長.
她在上海的時候.
創立了一個.
互助家庭.
後來在上海失手.
她去了香港.
等候神帶領她下一步.
最後就去了台灣.
她去過很多地方.
嘉義,高雄,屏東.
她每一個地方都和傳教士.

$^{481}$一起工作,傳福音.
很熱心的.
那些地方都陸陸續續設立了教會.
她又辦托兒所.
後來終於成立了.
信我愛兒童之家.
她離世的時候.
兒童之家收容了120多個.
弟弟.
有30個是嬰孩和幼兒.
其他已經是讀小學.
中學或高中.
她自己很注重.
兒童的教育.
鼓勵他們向上.
去讀書.
她早晚都會靈修.
任何一個長者.
早晚都會靈修.
所以養成小孩.
小孩就懂得向神祈禱.
尋求神.
小孩就學會了.
倚靠天父.
她很重視兒童的飲食.
營養,體格的發展.
為了充裕.
兒童之家的經費.
她養豬,養牛,養羊.
找專人去打理.
補充.
兒童的經費.
雖然有人救濟.
但不是全部靠人.
她希望.
自己有辦法.
去耕耘.
讓小孩有書讀.
生活好,有好營養.
她收容了.

$^{521}$孤兒寡母.
貧病無依的人.
不計其數.
很多.
她不單養也教他們.
她自己是一個模範.
所以由她兒子的孝宗說.
她好像一棵樹.
樹向上的主幹.
迎接著神給他們福氣.
然後.
樹葉.
就好像隱蔽著.
他們後輩的人.
無依無靠的人.
這棵樹.
很正直.
對神敬愛.
對人.
仁慈.
她一生中.
無數人都經歷過她的幫助.
當時她去到美濃鎮.
一個客家莊.
講客家話.
這些人不認識.
什麼是宗教.
當時她的存亡很艱難.
人們冷言冷語對待她.
後來她真的.
做到將神的愛.
彰顯在社會.
讓人們.
有機會.
信主.
所以她老人家.
全心愛主.
一直默默耕耘.
她一生中.
神沒有偏待她.

$^{561}$雖然我們.
曾經有一句說話.
天下無不是的父母.
但現代思潮下.
已經推翻了.
她創立兒童之家.
卻是天下.
無不是的孩子.
她說看你怎樣愛他們.
所以神沒有偏待她.
原來她那麼辛苦.
去栽培他們.
當中有些淑齡的兒女.
長成了.
他們還願意.
奉獻自己給神.
去做傳道人.
兒童之家.
由無到有.
你會看到.
她本來很簡單.
但她在默默耕耘中.
都是血汗.
所以令到社會.
人們.
懂得支持.
都要效法她的精神.
默默地做.
今天兒女.
就是神賜給你們的產業.
我們.
看到她.
用心地.
照顧.
無依無靠的.
孤兒.
今天我們.
成為神的兒女.
神賜給你們做管家.
我們是否願意.

$^{601}$做主中心的管家.
為子女祈禱呢?.
我們有沒有遵守.
神賜給我們的吩咐.
帶他們.
去信主呢?.
當他們偏行己路的時候.
我們會否幫助他們呢?.
免得他們.
誤入歧途呢?.
我們有沒有按真理.
去教導子女.
讓他們懂得用神的價值觀.
去待人處事呢?.
我們一起祈禱.
愛你的主.
愛你的神.
現代做.
父母的很艱難.
我們也感恩.
因為祂曾經讓父母.
辛苦地養大我們.
栽培我們.
當我們今天看到.
一代一代的出來.
看到父母的艱難.
與主一樣.
是艱難賜禮給他們.
讓他們在社會裡.
能夠滿有智慧.
去教導下一代.
教導他們的子女.
使他們成長成熟.
能夠合乎你的樣式.
能夠在人前.
有基督的樣式.
做好榜樣.
求主祝福每一個.
母親節的.
我們慶祝.

$^{641}$為他們感恩.
父親節也一樣.
做父母不容易.
求主給他們.
加倍的力量,加倍的智慧.
讓他們去教導子女.
讓他們願意.
認識神.
以致他們將來.
為主所用.
我們這樣祈禱交代.
\newpage



\section{馬可福音 4:35-41}
\label{sec:mgPJ9gnhbtQ}
\textbf{2024.05.19 《有主在船中》 馬可福音4\_35-41 (和修版) 講員 : 陳一華牧師}
\newline
\newline
連結: \href{https://youtube.com/watch?v=mgPJ9gnhbtQ}{\texttt{https://youtube.com/watch?v=mgPJ9gnhbtQ}} ~~~~ 語音日期: 2024-05-21
\newline
\newline
\hyperref[sec:Q94i4NHF3WY]{\small{< < < PREV SERMON < < <}}
~
\hyperref[sec:index_chronic]{\small{[返順時目]}}
~
\hyperref[sec:index_scriptual]{\small{[返順卷目]}}
~
\hyperref[sec:JJ0c8EbTgFo]{\small{> > > NEXT SERMON > > >}}
\newline
\newline
馬可福音 4:35-41
\newline
\begin{longtable}{cl}
\hline
\hline
章節 & 經文 (和合本修訂版)\\
\hline
4:35 & \begin{tabularx}{0.7\textwidth}{X} 那天晚上，耶穌對門徒說：「我們渡到對岸去吧。」 \end{tabularx} \\ \\ \relax
4:36 & \begin{tabularx}{0.7\textwidth}{X} 門徒離開眾人，耶穌已在船上，他們就請他一同去；也有別的船和他同行。 \end{tabularx} \\ \\ \relax
4:37 & \begin{tabularx}{0.7\textwidth}{X} 忽然狂風大作，波浪打入船內，以致船灌滿了水。 \end{tabularx} \\ \\ \relax
4:38 & \begin{tabularx}{0.7\textwidth}{X} 耶穌在船尾上，枕著枕頭睡覺。門徒叫醒他，說：「老師！我們快沒命了，你不管嗎？」 \end{tabularx} \\ \\ \relax
4:39 & \begin{tabularx}{0.7\textwidth}{X} 耶穌醒了，斥責那風，向海說：「住了吧！靜了吧！」風就止住，大大平靜了。 \end{tabularx} \\ \\ \relax
4:40 & \begin{tabularx}{0.7\textwidth}{X} 耶穌對他們說：「為甚麼膽怯？你們還沒有信心嗎？」 \end{tabularx} \\ \\ \relax
4:41 & \begin{tabularx}{0.7\textwidth}{X} 他們就非常懼怕，彼此說：「這到底是誰？連風和海都聽從他。」 \end{tabularx} \\ \\
[1ex]
\hline
\hline
\end{longtable}
$^{1}$陳浩思.
你們兩個都姓陳的.
謝謝.
我也要扮陳友站在這裡.
陳友唱歌大家有沒有共鳴.
很有共鳴.
他的見證是不是很感人.
很感動.
我們真的發覺.
陳友和他太太的經歷.
他不是時必要告訴我們.
你明白我的意思.
但是他真的很真誠地.
將他的經驗.
經歷一五一十.
告訴我們.
我們真的心很感動.
所以今天來這個聚會.
是不是很感恩.
將榮耀歸給天父.
特別我們今天很開心見到.
不單有陳友和他太太Pauline.
兩位一起.
在我們當中.
特別Pauline是陳友背後.
一個很重要的女人.
因為神做男做女.
是很奇妙的.
我剛開始的時候已經跟陳友說.
我說你放心.
我會說到.
長短火的.
你暢所欲言.
你如果有感動.
你就盡量說.
你長火我就短火.
你短火我就長火.
結果大家知道我今天選擇了什麼.
短火.
我就將我的講章.

$^{41}$cut了三分之二.
我就講三分一.
其實今天的見證.
跟今天的經文很吻合.
為什麼呢.
因為今天我想說.
有一次耶穌和門徒在船上.
突然之間就起了什麼.
風浪.
實際上那個地方不應該有風浪.
是罕有的.
為什麼呢.
如果你讀過聖經你會明白.
耶穌和門徒在那個海.
那個海叫加里尼海.
又叫格尼撒勒湖.
這個海.
實際上是內海.
是內湖.
為什麼呢.
因為這個海是眾山環抱.
大家明白我的意思嗎.
很難有什麼.
有這麼大的風會.
吹到下去.
所以這一幕的情景.
就讓我想起.
是我們人生的寫照.
大家不知道你認不認同.
有一句說話.
天有不測的什麼.
風雲.
人有十時的和福.
你認不認同.
這個風浪.
這樣起.
真的令到門徒.
很驚慌.
臉色很純白.
自己就快死了.

$^{81}$但是.
有一個很大的對比.
竟然.
在同一條船.
同一個處境.
耶穌.
竟然心平氣靜.
就睡在.
船尾的地方.
很安心.
睡覺.
為什麼有這麼大的分別.
關鍵就是信心的問題.
大家認不認同.
信心重要嗎.
很重要.
所以在路加福音.
在聖經記載的時候.
多說一句.
耶穌問門徒.
你們的信心.
去了哪裡.
這句說話.
當頭棒喝.
很值得我們去反思.
各位.
人生風暴.
好說不好聽.
真是避無可避.
你不想要的困境.
你不想要的疫情.
你不想要的逆境.
到我們避嗎.
避無可避.
與其逃避.
我們就去正視它.
各位.
我相信.
可能你的人生經驗未必有.
但我相信.

$^{121}$你在人生歷程上.
你不多不少都會遇過.
有風暴.
可能是什麼.
你的人生風暴.
可能是失利.
可能是生意失敗.
可能是.
經濟不太好.
老闆要裁員.
這麼不幸.
裁到自己.
你明白我的意思嗎.
可能你的人生風暴是什麼.
是自己身體健康.
出了問題.
可能是什麼.
可能是.
你與家人的關係很惡劣.
或與同事的關係很惡劣.
大家有心病.
有嚴規.
你又常常.
盯著我.
我又常常防備你.
你明白我的意思嗎.
各位牧師.
今早不知道大家.
你的人生經歷遭遇什麼.
但我總相信.
在我們人生經歷裡.
一定有起有跌.
有順有逆.
像剛才陳友華.
有一首詩歌.
上帝未曾應許.
天色常藍.
原來.
信了耶穌和未信耶穌.
我們都沒有分別.

$^{161}$我們都是人.
上帝都很公平.
你信了耶穌也好.
未信也好.
原來我們都會遇到人生風暴.
你相不相信.
不過最大的分別.
就是當你遇到人生.
有困難有逆境的時候.
你的信心.
在哪裡.
你的信心.
放在自己身上.
或者你的信心.
放在別人身上.
或者你的信心.
放在環境上.
還是你的信心.
放在耶穌身上.
就像陳友華.
當他走到人生.
窮途末路的時候.
當他人生.
跌到一個低谷的時候.
他很混亂的時候.
突然間他會遇到.
一位牧師.
突然間會提醒他要去祈禱.
突然間.
他自己去祈禱之後.
又有很特別的經歷.
各位.
這就是他的信心.
放在主耶穌的身上.
當我們發覺.
我們遇到人生風暴.
原來我們不是絕路的.
有沒有出路.
有出路.
出路是什麼呢.

$^{201}$出路是你的信心有沒有放在對的位置.
有沒有放在耶穌的身上.
為什麼呢.
因為主耶穌是大能的神.
祂是掌管萬有的.
祂也是掌管大自然的.
祂是很明白你和我的需要.
當我們遇到困難.
有逆境的時候.
我們將信心放在耶穌的身上.
祂就能夠.
第一祂會陪著我們.
祂不會離開我們.
第二祂會幫我們解決這個問題.
所以今天在陳友遞兄的身上.
你看不看到.
他想著太太.
真的說得不好聽.
好像沒有機會了.
是不是.
但是.
柳暗花明.
又一串.
上帝透過奇妙的醫治.
能夠令到陳友的太太Pauline.
可以再一次.
重新去做人.
你說是不是很奇妙.
所以各位朋友.
各位弟兄姊妹.
牧師要和大家做一個呼召.
很想鼓勵今天的新朋友.
牧師我也很想有一個.
決志的祈禱.
我想將我的信心.
放在對的位置.
將我的信心放在耶穌的身上.
我將我一生都交託給主耶穌的話.
待會祈禱的時候.
你跟我一起祈禱好不好.

$^{241}$你說我不會祈禱也不要緊.
牧師跟你一起祈禱.
我說一句你跟我說一句.
我將我的心志告訴天父.
在座當中的基督徒弟兄姊妹.
你們會不會都.
陪新朋友一起決志祈禱.
都會.
我們大家一起來向天父.
來去禱告.
我們低頭同心祈禱.
主耶穌.
多謝你.
愛我這個人.
雖然我微小.
我不配.
我甚至忘記了你.
偏行幾路.
但你愛我.
來到世界上.
被釘十字架.
流出寶血.
洗淨我罪.
今早我願意相信你.
接受耶穌.
做我的救主.
我將信心.
放在你身上.
求你帶領我人生.
祝福我的人生.
也祝福我的家人.
多謝耶穌.
聽我祈禱.
奉耶穌名求.
阿們.
我想問一下.
在座當中有幾多位朋友.
有跟我一起祈禱.
你舉手.
我看到很多.

$^{281}$有祈禱.
跟著牧師做決志祈禱.
不如我想請.
剛才跟我一起做決志祈禱的朋友.
你都站起來.
還有誰.
你繼續站起來.
不要坐下.
我現在想邀請.
決志信耶穌的朋友.
來到台前.
牧師跟你祈禱.
祝福你.
歡迎.
剛才有.
你決志信耶穌.
恭喜.
陳友遞兄.
也恭喜你們.
祝福你們.
恭喜你們.
還有誰.
新朋友.
今早跟著牧師做了決志祈禱.
現在來到台前.
我看到了.
還有誰.
我現在要交託在.
上帝的手中.
我的一生就能夠.
行走安穩.
因為耶穌是大能的神.
我看到還有人.
走出來.
還有誰.
不要錯過這個機會.
人生就像剛才.
陳友遞所說.
我們要把握現在的機會.
來到耶穌的面前.

$^{321}$把自己的生命.
交託給主耶穌.
還有誰.
你還未做.
你未決志的.
現在來到台前.
牧師為你祈禱.
不要讓他停.
還有誰.
人生就交託在耶穌的手中.
我不靠自己.
我不靠主耶穌.
來到耶穌的面前.
我們看到有長者.
恭喜.
願意決志信耶穌的.
每一位朋友.
陳友遞.
和他太太的人生.
是我們一個很好的人版.
他不是靠自己.
因為他知道人的力量.
是有限的.
我們要靠.
創天造地的上帝.
還有誰.
我們再一次.
為新朋友們祈禱.
我們一起祈禱.
親愛的天父.
我們今早很興奮.
很開心.
很雀躍.
因為有很多新朋友來到你的面前.
願意將他們的一生.
交在耶穌的手中.
願意信靠耶穌.
跟隨耶穌.
以致他們的人生.
是一個豐盛的人生.

$^{361}$他們的生命.
不再一樣.
我求祝福今天來到台前.
祝福他們的家人.
一人信主.
全家得救.
多謝耶穌聽我們的祈禱.
奉耶穌名求.
阿們.
很開心.
有沒有辦法.
和真朋友拍張照.
我們真的看到.
我們拍張照.
我們.
陳友遞兄你過來.
你上來.
陳友遞兄你上來.
我們望向回眾.
望向回眾拍張照.
拍張照.
我們今天決志信耶穌.
我們一起拍張照.
留念.
Pauline你也來.
Pauline.
我們歡迎Pauline.
Pauline.
我們和心慧.
拍張照.
我們用V字.
靠耶穌得勝.
V字.
多謝主.
靠耶穌得勝.
非常好.
因為我們地方.
淺窄.
現在我們好不好.
陪談完.

$^{401}$各位決志的朋友.
去到.
後面.
305的地方.
我們有個陪談.
各位.
決志的朋友.
跟著.
陪談完.
去到後面.
有個房間.
305的房間.
我們有聊天.
做一些陪談工作.
有個聊天的時刻.
好.
大家.
就.
你有大心朋友.
今天決志.
你也陪他去305的房間.
我們一起.
有個聊天的時刻.
因為這個地方淺窄.
就比較.
沒那麼寬鬆.
大家可以去到.
去到後面.
我將以下時間.
交給主席.
大家坐好就OK.
坐好.
OK.
\newpage



\section{詩篇 46:1-11}
\label{sec:JJ0c8EbTgFo}
\textbf{2024.05.26 《上帝是我們患難中隨時的幫助》 詩篇46\_1-11 (和修版) 鄧泰康傳道}
\newline
\newline
連結: \href{https://youtube.com/watch?v=JJ0c8EbTgFo}{\texttt{https://youtube.com/watch?v=JJ0c8EbTgFo}} ~~~~ 語音日期: 2024-05-28
\newline
\newline
\hyperref[sec:mgPJ9gnhbtQ]{\small{< < < PREV SERMON < < <}}
~
\hyperref[sec:index_chronic]{\small{[返順時目]}}
~
\hyperref[sec:index_scriptual]{\small{[返順卷目]}}
~
\hyperref[sec:QSB7W2xKaig]{\small{> > > NEXT SERMON > > >}}
\newline
\newline
詩篇 46:1-11
\newline
\begin{longtable}{cl}
\hline
\hline
章節 & 經文 (和合本修訂版)\\
\hline
46:1 & \begin{tabularx}{0.7\textwidth}{X} 神是我們的避難所，是我們的力量，是我們在患難中隨時的幫助。 \end{tabularx} \\ \\ \relax
46:2 & \begin{tabularx}{0.7\textwidth}{X} 所以，地雖改變，山雖搖動到海心， \end{tabularx} \\ \\ \relax
46:3 & \begin{tabularx}{0.7\textwidth}{X} 其中的水雖澎湃翻騰，山雖因海漲而戰抖，我們也不害怕。（細拉） \end{tabularx} \\ \\ \relax
46:4 & \begin{tabularx}{0.7\textwidth}{X} 有一道河，這河的分汊使神的城歡喜，這城就是至高者居住的聖所。 \end{tabularx} \\ \\ \relax
46:5 & \begin{tabularx}{0.7\textwidth}{X} 神在其中，城必不動搖；到天一亮，神必幫助這城。 \end{tabularx} \\ \\ \relax
46:6 & \begin{tabularx}{0.7\textwidth}{X} 萬邦喧嚷，國度動搖；神出聲，地就熔化。 \end{tabularx} \\ \\ \relax
46:7 & \begin{tabularx}{0.7\textwidth}{X} 萬軍之耶和華與我們同在，雅各的神是我們的避難所！（細拉） \end{tabularx} \\ \\ \relax
46:8 & \begin{tabularx}{0.7\textwidth}{X} 你們來看耶和華的作為，看他使地怎樣荒涼。 \end{tabularx} \\ \\ \relax
46:9 & \begin{tabularx}{0.7\textwidth}{X} 他止息戰爭，直到地極；他折弓、斷槍，把戰車焚燒在火中。 \end{tabularx} \\ \\ \relax
46:10 & \begin{tabularx}{0.7\textwidth}{X} 你們要休息，要知道我是神！我必在列國中受尊崇，在全地也受尊崇。 \end{tabularx} \\ \\ \relax
46:11 & \begin{tabularx}{0.7\textwidth}{X} 萬軍之耶和華與我們同在，雅各的神是我們的避難所！ \end{tabularx} \\ \\
[1ex]
\hline
\hline
\end{longtable}
$^{1}$早安,有一段時間沒有見大家了.
今天我一進來的時候.
施主編很好,派了一張大版的程序給我.
隨著我年齡的增長,我從你們這裡由細版變成大版.
我估計我應該都可以開始做長者的工作了.
今天跟大家分享的時候,我想跟大家分享一個很貼題的情況.
我想先問大家,大家昨晚睡得好嗎?.
真的假的?你們睡得那麼好?.
通常這個情況我就會立刻想一下,可能我比較陰謀論.
通常你問鄧恕露好不好,我說OK,其實我OK.
男人通常都是這樣的,你問我我就回答你.
最近睡得好不好?這一年睡得好不好?.
如果大家真的睡得好,我真的想說為他瞭瞭眼.
為什麼呢?中國人有句說話說得很好.
人無遠慮,就必有近憂.
如果你是沒有憂慮的話,其實這個才是不正常.
我想說這個是人性,如果撇開了信仰.
同樣都是,我都是很軟弱,其實從這裡就看到人是很軟弱.
在最近,整個世界,先說香港,大家最近有沒有不上消費?.
應該有些人都參與在這個行列.
你會發覺一件事,你做老闆的,打工的,你就看到整個香港.
平時排隊那麼辛苦,不出去,現在不用排了.
那你出去嗎?你還是不出去,一滿的,特別是做生意的.
特別是打工的,慘了,遲些的時間工作都不知怎樣.
好了,子女教育問題,我兒子以前讀一間小學,我不說名字.
我最近有機會見到他的班主任老師.
他說鄧全道,現在小朋友讀一年級的東西,是你兒子以前讀三年級的東西.
你就可以想像到現在小朋友讀書的壓力,再加上家長的壓力.
現在的家長一見到他就會說,晚上我和他溫習都未溫習好.
明天就考試了,都未溫習好.
但我回想,我兒子以前讀書好像不是這樣的.
我晚上只是和他玩,我沒有和他溫習.
我想說整個社會推著你,好像要向前走.
你不走,你會覺得有問題.
既然香港這麼大壓力,有空放假,不如去旅行.
大家最近旅行有沒有綁好安全帶?.
一見到新聞的時候,哇,害怕.
坐飛機都很大壓力.
去到其他地方,大家有沒有帶好防災設備?.
這裡有地震,那裡有水災,那裡有雪災.

$^{41}$整個社會,整個環境,令到我們大家都會覺得很有壓力.
所以我覺得今天和大家分享這段經文是很適切的.
剛才我提過,人的憂慮是正常的.
你沒有憂慮才是正常的.
但問題是,不是停在這裡.
我們作為一群信主的人,或者我們尋求上帝的人.
我們必須知道,接受自己是事實,我們是會有憂慮.
但是,最重要的是你有沒有出路.
信耶穌的人是必須要有出路.
因為我們的上帝在聖經裡面,他給了我們很多寶貴的說話.
而這些說話,我想告訴各位,你不要小看.
因為今天的救恩都是從這些說話而來.
是上帝對大家的應許而來.
是上帝由最初對以色列人的應許,直到今天對我們的應許.
從來沒有改變而來.
所以如果你不相信上帝的話,我可以告訴你們.
你們就要繼續憂慮下去.
你們就等於信耶穌,等於沒有信耶穌.
因為其實我們和外面的人,最大的分別就是我們有上帝.
今天和大家分享這篇詩篇的時候,是一篇很經典的詩篇.
大家看詩篇的時候,一定要看著腳.
一開始的時候,他講這篇詩篇是可拉的後裔.
他將這篇詩篇交給司令長,即是司班長.
亦都寫明了要用女高音去唱.
有目的的,大家不要不看這些.
看聖經一定要看到頭一天.
看這篇詩篇的時候,他講的是女高音.
其實你就要想像,我去到一個歌劇院.
一個女高音站在中間的時候,和你去唱一首歌.
整個歌曲,整個環境給你的,是一個什麼樣的氣氛.
是很澎湃的,你記住,是很澎湃的.
和很情緒高昂的.
你一定要拿著這樣的氣氛去看這段聖經.
你才會有感覺.
這段聖經開始的時候,其實他用了兩幅很對比的圖畫.
一個圖畫是外面非常混亂.
外面非常混亂,但是在上帝的裡面很安靜.
兩個圖畫是很不同的,所以你就可以想像.
當這個女高音唱的時候,唱到外面有地帶震動.
咆哮的時間,是令你震撼的.

$^{81}$但是接著他一唱,唱到我們也不懼怕的時候.
他就突然間收了聲.
他唱到我們在神面前要安靜的時候,他就收了聲.
這個就是音樂的能力.
這個作者要帶領你進入詩歌當中.
要讓你在你最不容易,最困難的時候.
讓你有共鳴,帶你回到一個安靜的環境.
這段經文是這樣的.
所以地需改變,其實是在說大地震動.
山水搖動到海深,崩塌入海.
你會看到一個場景,是震到一個地步的時間.
就好像之前花蓮地震,大家記不記得太老閣的片段.
是整個飛沙走石,衝下海的環境.
然後你留意地震完之後衝下海.
海受到石頭撞下去的時候,就變成了翻騰.
其實是一連串的.
因為給你的感覺是沒有喘息的機會.
然後你就看到海,因為海漲而戰鬥.
因為衝下去,海一升起的時候,海浪就在轉動.
你想像一下,你身處在當時.
很描繪當日地震的時候,在太老閣的朋友.
在山上所經歷的事情.
我回看那些短片.
其實真的很不容易.
如果我身處其中,我想我心裡真的很慌.
在上一次台灣花蓮地震的時候.
曾經有一個很出名的,我忘記了是統帥酒店還是總統酒店.
在上一次地震的時候,地震塌了.
整層壓下來.
就在塌之前的兩個月.
我去了台北,就是在花蓮,因為我大伯住在花蓮.
他就是請我去那間酒店吃自助餐.
所以當他塌的時候,我心裡很有那種感受.
我會立刻想像當時的環境.
你會發覺整個外面的世界.
沒有一個地方可以讓你安靜和安心.
等到今天你去面對著,剛才我問大家睡得好不好的時候.
我不知道你腦裡會衝出來一些什麼形象.
有很多這樣的環境會令我們無法安睡.
因為好像沒有蠟燭之地,好像沒有其他東西.

$^{121}$但是這個作者卻是告訴你,我們也不害怕.
儘管是這麼厲害的情況之下.
但是我們不害怕.
但這不是我們很想要的嗎?.
然後他就告訴你,有一道河是分叉的.
使得神的城歡喜.
意思就是將這些歡樂帶到上帝的城當中.
而這道河是說的是從上帝而來的祝福.
在啟示錄曾經有一段經文.
一道生命的河,光亮如水晶.
從神和羊的寶座流出來.
你會發覺,這道河是在說上帝的恩典.
其實是在你最困難的時候,他會流到這裡.
這是一個很重要的地方,是上帝的城.
所以你立刻要問,你今天是否在上帝的城裡.
聖經裡面四福音,其中一個很大的觀念.
耶穌給你的觀念,叫做圈內和圈外.
今天我們很多人,你可以說自己信了耶穌.
但是你可以是一個圈外人.
但究竟關鍵是你今天是否在神的城當中.
而這個城是至高者的聖所.
你不要立刻想到,就說是教會.
有什麼事就衝過來教會.
不是的.
記不記得勝經裡面曾經說過,你是神的殿.
神的殿裡面最關鍵,你有沒有神的靈.
聖靈有沒有內在.
聖靈有沒有充滿你.
你是不是屬於聖靈.
我們剛剛在五月感受,那裡有一個金句.
神的靈在哪裡,哪裡就有自由.
弟兄姊妹,你今天是否在神的城當中.
神在其中,神成必不動搖.
到一天亮的時候,神就幫助者成.
我們在神的城當中,儘管外面有很多震撼.
你也不需要擔心,因為等到天亮的時候,神就會做事.
這裡說的是上帝的時間.
是上帝的時間,不是你認為的時間.
今天很多環境我們都不知道為什麼.
以前的時間,我想說我們人最糟糕的地方.

$^{161}$就是跟著整個環境走,不是跟著上帝走.
香港很風光,如果大家還記得.
我都很感恩,我在那段受惠的時間成長.
80年代,我經常跟弟兄姊妹說,我不會跟你比.
因為你現在的情況跟我那時候是沒得比的.
那時候我又不是讀書,又不懂讀書,成績又不好.
出道之後,我都找到一個文員做.
跟著我每天回到公司的時間,很準時.
可以朝九晚五,其實是朝九晚四點半洗手.
跟同事打下牙,就可以下班.
很輕鬆地工作,一個部門的時候.
我記得老闆跟我說,我自己在裡面工作.
明明知道部門所做的事情,不需要有十個人.
但老闆說不可以不請十個人,知道為什麼嗎.
因為你今年停了,老闆就不會再給你.
所以就算沒事情做,你都要請十個人.
跟著時間,我們原本是這麼少的事情,分十個人.
大家想一下大家在做什麼.
所以我很記得,現在大家還有沒有三點三下午茶.
還有沒有?應該沒有吧.
現在做到都不夠時間,是不是想死.
我們不是的,一到三點鐘的時間,就接到秘書打電話問老闆.
老闆,三點三嗎?.
我們有大食大,當然老闆就給錢.
所以當時我很記得,如果大家懂得喝這個東西,就像我一樣老.
我記得下午同事問我喝什麼,我回答,滾水蛋.
很經典的,大家未喝過.
在這個情況下,你會發覺,我們在那個快樂的年代.
我記得曾經在九十年代衝起來之後,大家如廁怒放.
連我媽媽,一個小學畢業,對股票什麼都不懂的人.
之後就去買股票,我還記得那時義氣風發.
我未見過媽媽跟我說話那麼義氣風發.
一回到家,就跟我說,兒子,我今天又賺了多少錢?買一點.
你知道我們這些是孝子,媽媽叫你給錢她去買股票.
你怎麼可以不給呢?她說得很清楚,不是給我.
是你給錢我,我幫你投資,我多厲害,我幫你賺錢.
我聽她說,就給她幾千元,她不多,給她,然後她就幫我買.
每次我回家吃飯,她說,我今天又幫你賺了多少錢?.
這些錢我從來沒有見過,她只是告訴我賺了多少錢.
她說,我繼續幫你買,繼續幫你買,大家還記得吧.

$^{201}$股票,股災,金融海嘯,大家記得吧.
我記得那段時間,我回去就特意問我媽媽,媽媽,股票怎麼樣?.
然後我媽媽就說,你不要問了.
我很記得,我們在整個世界很好的時候,我們認為很發展的時候,我們就很快樂.
我們就跟著這個世界一起快樂,很安心,什麼都不用怕,很闊綽,如廁撈飯.
到了這個世界搞不定的時候,我們就一起向下.
頂智媒是不是要這樣?我們是有上帝的,頂智媒.
我們從來都不是跟著這個世界外面的環境而活的.
我們是跟著上帝的應許而活的.
儘管在現在這麼困難的時候,請你記得,上帝有祂的時間.
因為祂從來沒有變過,在祂的時間到的時候,神的幫助就會臨到當中.
因為上帝在這裡說,外邦軒然,外邦一片混亂,列國動搖,那些地方地震,神發聲,地變融化.
上帝的說話,你要記住,是明立就立,說有就有.
你不要小看我們聖經,頂智媒.
為什麼我們做傳道人,說到吐血,都要提醒大家要看聖經,要學聖經,要靈修.
你沒有上帝的應許的時候,你怎樣去面對這個世界呢?.
你怎樣去面對這個混亂不安的地方呢?.
你有什麼材料去對抗呢?.
所以我們必須要回到上帝的這個應許當中,你要相信祂.
然後這個作者很有這種志氣,他說,你來看夜話的作為,看他使地怎樣荒涼.
他指識刀兵直到地極,得節斷弓斷槍,把戰車焚燒在火中.
當時的混亂,其實就是在說戰爭,如果大家很清楚,馮檢基的詩歌出現的時候.
就是以色列和其他地方,很多地方在打仗.
他在這裡說的時候,其實看他使地怎樣荒涼,其他譯本是譯到看他如何摧毀世界.
但是我們一看,當然不是,這個世界上帝那麼愛,怎會摧毀呢?.
所以你弄一個易城是什麼呢?你看看上帝怎樣去處理這件事.
你看看上帝怎樣去摧毀這些不合他心意的事情.
他的方法是什麼呢?他就是要結束這些戰爭,將所有的武器銷毀.
從來他處理戰爭的方法都不是戰爭,你不是越打越贏就越搞定,從來都不是.
從來都是在上帝的愛當中,從來都是在上帝自己的心意當中.
上帝不是用我們人認為好的方法去處理的.
我們曾經讀過一些經文,一些根據,在上帝的裡面,我們不是靠兵,我們不是靠馬.
我們是靠耶和華,方能成事.
但是今天我們在教會裡面,我們是用什麼智慧去面對教會,去面對我們所有的生命.
我們是用我們的方法,還是我們是用上帝裡面所說的方法.
沒有上帝的地方,那個地方不會有自由,那個地方不會有安靜.
所以我說當你自己抓不緊這個說話的時候,你就很容易陷在這些混亂教會,我照樣,但是其實我心裡面是空蕩蕩,完全沒有上帝.
所以在最後那段經文,是我們最熟悉的.
你們要休息,要知道你是神,我必在外邦中被尊崇,在遍地上也被尊崇.
所以在患難裡面,在外面很混亂的時候,我們最需要的是安靜,而去認清祂才是上帝.

$^{241}$甚至乎你要去敬拜上帝.
當你今天有很多困難的時候,你不是像,中國有一個形容詞,叫熱窩上的螞蟻.
大家很急的時間就上網找,我舉個例子,我鼓勵大家千萬不要這樣做.
一見到上帝,醫生就跟你說,照到陰影的時候,你上網立刻怎樣.
身體左邊陰影,這些關鍵字,我就列了一大堆嚇死你的東西.
可能是癌症,可能是什麼什麼.
我們人很喜歡用這些字,你以為我知道的越多,我就越能夠去處理這件事情,你錯了,弟兄姊妹.
之前我剛教完傳道書,傳道書的作者告訴你,你知道的越多,其實你就越多煩惱.
我不是反智,你弄清楚,我不是反智.
但我想說的時間,如果你把這些知識取代了上帝的話,你就糟糕了.
我們人要努力,但基於有上帝的裡面.
我們有知識,儘管我知道我是癌症,那上帝讓不讓我死呢?.
弟兄姊妹每一個都跟我說,鄧傳道很危險的,想清楚,不要做.
然後我回應的只有一個答案,上帝要我死,我就死了.
上帝不讓我死,你想死都死不了,我的做法就是有病你就醫.
我在前面有路我就走,而其他的都是上帝的.
我們今天面對著整個世界,我經常說你很憂慮,剛才說大家分不分得準.
我很憂慮,我很多事情想,你憂慮完之後可以處理這件事嗎?.
其實一點幫助都沒有的.
其實我們要在神面前學習一件事,剛才我也說了,我們有沒有憂慮?有,我們有憂慮.
請你好好,好像聖經所教,安靜回到神的面前,認清祂才是神.
是祂去決定所有事情的結果,弟兄姊妹.
你唯一需要的是什麼呢?就是去敬拜和尊崇祂.
所以當大家最憂慮的時候,我真的鼓勵大家,唱歌,唱詩歌.
繼續去讚頌上帝,因為詩歌能夠提醒大家,你能夠在安靜當中.
有些很基本的,都是說到嘴裡都裂開的靈修,你要透過安靜靜觀.
你要去尋找你自己的內心,你在裡面的時候去找回上帝的作為.
在混亂的裡面不是做更多的事情,而是在裡面的時候安靜下來去找回上帝,這才是方法.
所以有一個很重要的信念,就是屬於神的人,我們一樣會面對著患難.
我們不會讓其他人面對患難,我們一樣有病,跟他們一樣,別人水浸,我不會不水浸.
我們一樣會面對很多這種情況,但是唯一的分別就是我們仍然要確定.
讓我們能夠得到保護安穩有力量的是上帝.
為什麼呢?剛才那段經文當中,如果大家有留意,第一節.
神是我們的避難所,其實避難所這個字是堡壘.
神是我們的堡壘,是我們的力量,是我們患難中隨時的幫助.
第七節,萬君之的話與我們同在,雅各的神是我們的避難所.
第十一節,萬君之的話與我們同在,雅各的神是我們的避難所.
作者不斷一次又一次再一次地提醒大家,祂才是你的堡壘.
弟兄姊妹,今天我們回來教會也有一段日子.
你對神的信心有多少?.

$^{281}$你是不是真的信得過神?.
我們當日決志信耶穌是你講的,沒有人拿著耳機逼你.
我們願意將我們的生命完完全全交給神.
你才能做基督徒,大家知道的,沒有人騙你回來.
我們每個人都很清楚地告訴你,你必須將你整個生命傳言交給神.
是你答應了的,但我們今天有沒有像一個圖畫.
當我們拿著很多重擔來到神的面前,向祂哭訴和祈禱.
上帝,我的重擔很多,上帝邀請你,你將你的重擔放在我私隱座前.
但很多時候我們都不願意,我們喜歡自己死攬的時間,然後離開.
知道為什麼嗎?因為你信不過神,你信自己.
你覺得你不去做,你不去處理,事情是解決不了的.
所以最後我用一個例子和大家做結束.
我一直有做一些婚前輔導,婚姻輔導.
有一次有一個人轉介,我的婚前輔導和婚姻輔導不一定是信徒.
我覺得如果對方肯聽我講,我也照做.
因為我不是信徒的男生,兩個都不是信耶穌.
他找我,其實不是婚前輔導,因為女生很想結婚,男生很不想結婚.
他來找我談,女生試圖找我勸男生和她結婚.
男生試圖找我勸女生不要和他結婚.
大家各自表述,大家明白嗎?.
搞了一段時間,男生是一個聰明的男生.
但他沒有告訴我一些特別的原因.
總之他就說了很多結婚的不好處,怎樣怎樣.
總之他一直推.
女生最後的殺手簡是甚麼?.
你在何時何時之前你不和我結婚,我和你分手.
最後就是這招,否則沒有其他可以做.
後來我發覺我都搞不定,我用了很多方法幫他都搞不定.
有一段時間他找我,我就停了,我開始停了.
反過來我覺得暫時要停一停,不可以再和他談.
我主要是為他祈禱,我用了一些時間為他祈禱.
我祈禱,我說上帝你搞定吧,我這麼大了,我真的搞不定.
我用盡了我的輔導技巧,我都搞不定.
然後我就忙其他事情,男生也沒有再找我.
我記得我最後的時候,我和他一起祈禱.
我說雖然你不信耶穌,但你信我,我相信上帝可以幫你.
我說上帝,他在裡面發生甚麼事我不知道.
我不知道他為何這麼不想結婚,我不知道他甚麼情況.
不過上帝來幫他解虛,問題不是結婚和不結婚.
而是他被一些東西卡住他不想結婚.

$^{321}$我大概是這樣祈禱,跟他祈禱完就再見.
好了,隔了半年,有一天我突然收到這個男生的電話.
他很開心地打給我,說等著瞧,很久沒有找我了.
其實我心裡很怯,我說你又找我,我搞不定的.
然後他打給我,說等著瞧,我通知你我結婚了.
我說這麼奇怪?然後我很想找你聊天.
我說你肯結婚,當然找我聊天,我也想知道發生甚麼事.
他們找我聊天,一坐下來,他哭著對我說.
他不是因為結婚而感動,他是說這段日子發生甚麼事.
因為這個男生,他從小父母離婚.
其實他很想自己跟媽媽住,但他跟他父親沒有聯繫.
他心裡這麼多年,從小到大都有個問號.
是不是我爸爸疼我?是不是我爸爸不要我?.
後來他就發生了一件事,因為他突然醫生說他有個病.
而這個病他不方便跟其他人說.
他一直想,想不到誰可以說.
然後他竟然想起他爸爸,他打電話給他爸爸.
跟他爸爸說他現在有這個病,然後他爸爸就安慰他.
他爸爸說了一句話,兒子你放心.
你要處理,你要看醫生,我一定支持你.
我傾家蕩產,我都會幫你.
他甚至甚麼都說不出口.
他在那一刻突然間發現,他終於知道為何他這麼怕結婚.
因為他很怕自己好像他爸爸.
他自己都不確定他爸爸是怎樣.
但現在上帝親自幫他醫治了他,透過這個病醫治了他.
讓他解答了這個答案,原來我爸爸一直很愛我.
原來雖然我爸爸媽媽離婚,但我爸爸對我的愛從來沒有變過.
所以他就發覺他有這份勇氣.
他經歷這件事之後,他就跟他女朋友說.
我願意結婚.
所以我想問給各位聽,很多時候當我們面對一些情況的時候.
不是用我們人所想的方法,我們可以盡力.
但上帝是怎樣去做,是上帝的方法.
所以今天我都鼓勵,我經常說這是一個很重要的平衡.
當今天我們面對這個環境,特別是現在整個香港的環境.
其實我們可以做的事很少,但我還是那句.
有甚麼是按照上帝的說話,我們做出來的時候.
是可以祝福到我們身邊的人,包括你的家人.
包括你身邊的朋友,包括你身邊的同事,我們盡力去做.

$^{361}$而最後將所有的事做完後,交給上帝.
我夠膽絕對告訴大家,你一定會經歷到上帝在你前面.
開一條一條你覺得不可能的路.
你一定會看見上帝的作為,你一定會看見上帝的同在.
你就會有像作者所說的信心.
儘管外面是怎樣也好,但我不懼怕.
因為我的上帝在當中.
我們聽者會很盼望,我們相信耶穌.
相信到他那麼驕傲,相信到他那麼有信心.
他夠膽告訴其他人,我的上帝很棒,我的上帝很厲害.
希望我們都成為一個有生命的基督徒.
我們一起去成為神的見證.
能夠活得更有盼望,我們一起同心祈禱.
秋天我們求主,你自己親自來幫助我們.
我們在這段日子裡,我們或多或少每一個人都在我們生命裡.
有我們的懼怕,有我們的擔憂.
主,我求主,你親自來保守帶領我們.
我們很盼望,像詩歌裡的詩人一樣.
對你的信心,能夠帶到一個地步.
儘管在外面的環境我們面對不到.
我們的能力沒辦法處理的時候.
我們仍然相信上帝在裡面.
你會用你的方法,你會在你的時間裡面.
主,你會去處理.
最後我想求主,你讓我們在你面前的時候.
成為一個寄獻碑,我們知道我們其實很軟弱,我們很微小.
但同樣的,我們都帶著上帝的大能.
去成為其他人祝福的一個生命.
主,我們就將整個紅恩堂的每一個弟兄姊妹.
他們的領袖,牧者,我們交在神的面前.
求主,你的恩典臨到他們身上.
希望他們能夠從你的裡面得到這一份肯定.
得到這一份盼望.
最後願意他們能夠成為一個有靈的活人.
成為一個有信心的信徒.
將一切在所有的家庭面前祈禱,禱告.
奉主耶穌基督名字祈求.
\newpage



\section{希伯來書 11:1-3}
\label{sec:QSB7W2xKaig}
\textbf{2024.06.02 《真實的信心》 希伯來書11\_1-3 (和修版) 曾敏姑娘}
\newline
\newline
連結: \href{https://youtube.com/watch?v=QSB7W2xKaig}{\texttt{https://youtube.com/watch?v=QSB7W2xKaig}} ~~~~ 語音日期: 2024-06-03
\newline
\newline
\hyperref[sec:JJ0c8EbTgFo]{\small{< < < PREV SERMON < < <}}
~
\hyperref[sec:index_chronic]{\small{[返順時目]}}
~
\hyperref[sec:index_scriptual]{\small{[返順卷目]}}
~
\hyperref[sec:EkAEiOY56rU]{\small{> > > NEXT SERMON > > >}}
\newline
\newline
希伯來書 11:1-3
\newline
\begin{longtable}{cl}
\hline
\hline
章節 & 經文 (和合本修訂版)\\
\hline
11:1 & \begin{tabularx}{0.7\textwidth}{X} 信就是對所盼望之事有把握，對未見之事有確據。 \end{tabularx} \\ \\ \relax
11:2 & \begin{tabularx}{0.7\textwidth}{X} 古人因著這信獲得了讚許。 \end{tabularx} \\ \\ \relax
11:3 & \begin{tabularx}{0.7\textwidth}{X} 因著信，我們知道這宇宙是藉神的話造成的。這樣，看得見的是從看不見的造出來的。 \end{tabularx} \\ \\
[1ex]
\hline
\hline
\end{longtable}
$^{1}$(自動翻譯).
讓我們可以回到你的家裡.
一同去敬拜,讚美你.
我們是何等的不配.
因為你很愛我們.
要藉著你的話語.
成為我們腳前的燈,路上的光.
求你幫助我們眾人.
有一個徹底開放的心靈.
去聽上帝你自己的聲音.
願聖靈在我們裡面.
去幫助,教導我們.
去明白神的話語.
我們又能夠作出悉切的回應.
我們這樣的禱告.
是奉耶穌基督的名字.
阿們.
剛才主席說愛護的.
我想起去年.
對於我來說是非常特別的一年.
我經歷了弟兄姊妹的愛護.
去年我們有同工的患病.
周牧師也過來,有很多協助.
同時弟兄姊妹也對我很多幫助.
很愛惜我,知道我很孤單.
去年也是一年.
我在屬靈上有很多學習的一年.
我相信跟我比較熟的弟兄姊妹.
也有聽我分享過.
去年是一月份的時候.
是2023年.
一月份的時候.
在我完全毫無心理準備的情況下.
上帝就將一個炸彈.
放在我手裡.
我為人是這樣的.
我很看重家庭.
我很需要家庭生活.
那時候爸爸媽媽仍然在世的時候.
我真的很依賴他們的家庭.

$^{41}$我以他們為我的根基.
跟他們一起生活.
總之他們在哪裡生活.
我一定會跟著他們一起住.
後來爸爸媽媽相繼安息主懷.
只剩下我自己一個.
剩下什麼呢?.
剩下那間屋.
那間屋某個程度上.
就是我的一個安定的地方.
每天我期望自己.
完成了教會的服事.
就可以去一個固定的地點生活.
所以那間屋是很重要的.
如果連那間屋都沒有了.
對我來說是有影響的.
我記得一月快要接近農曆新年的時間.
竟然有一天.
有一晚我洗完澡之後.
忽然間我發覺完全沒有水.
是一滴水都沒有.
然後我就想.
是不是水務處有什麼事情.
可能很快會收到他們的通告.
應該明天去到差不多的時間.
會收復了.
去到第二天.
仍然是一滴水都沒有.
真的很刺激.
那怎麼辦呢?.
我的基本生活.
大家都想像得到.
很多東西是要用水的.
而且我那裡.
因為是鄉村地方.
我們是接駁不到鹹水的.
是沖廁所.
我們是需要吃水.
沒有水就代表.
每樣東西都會影響.

$^{81}$洗手間,洗澡.
日用飲食很多東西.
我就開始很焦急.
很焦急.
最基本還要處理.
洗澡的問題.
很好.
當天我記得有弟兄姊妹都說.
你可以因為住得比較近.
你可以過來我那裡洗澡.
我姐姐在上水市中心.
她都說可以過去洗澡.
不過我想想.
我洗完澡.
那裡的走路都要十幾分鐘.
洗完澡回來.
全身都是汗.
還要加上地上的塵.
而且沒有水.
真的有很大影響.
我隔壁有個公廁.
很不方便.
那段時間.
心情一直很徬徨.
可以這麼說.
很焦急.
很想快點弄好水.
後來問問水務處.
最後發現.
原來不是水務處出了問題.
最後發現.
就是我那一座斷水管.
或者漏水.
整條村都沒有人發生這樣的情況.
只有我那一座.
我和樓下有一戶人家.
我們兩戶人面對這個情況.
整條村都沒有.
只有我們那一座水管.
這是神為我度身訂造的功課.

$^{121}$是信心的功課.
那我可以住在哪裡呢?.
你說一天半天都有方法處理.
但如果斷了水管.
或者搞不定.
都不知道什麼時候.
很快我們就嘗試找師傅過來.
看看水管的水管口在哪裡.
原來住鄉村.
真的是很困難的事情.
經過這件事我學會了.
花了很長時間.
才大概了解到水管的路徑.
原來水管經過私家樓.
經過私家地.
我們的師傅也不太懂得處理.
就在年三十晚.
就在私家地作開.
就想幫我們看看.
是不是斷了水管.
結果惹來很大的不滿.
又說要報警.
於是就掀起一場又一場的風浪.
然後我們又要蓋起那些掘起的地方.
然後村民紛紛拒絕.
拒絕我們去掘地.
了解水管在哪裡.
就算知道水管是不是在他們附近斷開也好.
他們拒絕讓我們掘開的地去修理.
我們就站在那裡.
我站在那裡完全沒有辦法.
在這個過程中.
你會感受到上帝完全關了門.
是沒有辦法.
我姐姐直接說.
你要住酒店了.
住酒店住多久.
很不方便.
我的家當都搬了出去.
想像一下住天水圍的酒店.

$^{161}$每天來回洪水橋.
上教會.
但是真的很不方便.
很多的陪徒都不在.
我記得那段時間.
上帝要拆毀我一些東西.
我經常要抓住一些東西.
不是抓住人.
是抓住房子.
抓住固定的東西.
現有的東西.
抓住一個很穩定的生活.
抓住規律.
什麼都想抓住.
現在上帝一下子就打倒它.
你沒辦法抓住.
我連房子都不讓你住.
但是上帝的過程.
要我學信心的功課.
你沒東西可以抓住.
你想投靠家人.
家人叫你去住酒店.
你可以怎樣呢.
你也不可能每天走去弟兄姊妹家.
我想過來你那裡洗澡.
洗了很多次.
都不方便.
在這個過程上.
上帝要我學信心的功課.
你現在沒東西可以依靠.
你只有我.
你只有我一個.
你沒東西可以抓住.
你是不是傳道人也好.
都要學信心的功課.
這個當時我很深刻.
當然我來到現在.
我今天願意在講壇講整件事的時候.
我已經經過了.
裡面有很多神的恩典.

$^{201}$到最後結束的時候.
我講當中怎樣經過.
但我想說在過程中.
神要我一步一步走.
放手.
你就是不可以靠這個.
不可以靠那個.
全都關了路.
當時村民對我們的態度非常惡.
完全一個幫助都沒有.
無論怎樣我都搞不定他們.
但這個過程就是上帝為我開籠.
上帝要我學.
抓住我吧.
你不要抓住地面上的人和事.
抓住我.
就是這樣.
由2023年的年頭.
一直要學抓住上帝.
有時候要疑惑.
有時候又要掙扎.
在過程中一直去到.
大概2023年.
好像差不多年尾.
整件事就解決了.
整個信心的功課.
都要經過一段時間去學.
我就發現到其實信心是很重要.
沒有信心走不下去.
但我相信不只是我需要信心.
而是在在每一個.
每一個屬上帝的人都需要信心.
我們在講壇裡聽過不同的講員講信心.
也聽過一些人講憲政.
但我覺得最重要是今天我和你去實踐信心.
是需要的.
沒有信心我們很難走人生這條路.
很難走向永生.
分分鐘半路就跌倒了.
不知去到哪裡.

$^{241}$所以今天我要講真實的信心是很重要.
沒有信心的人是怎樣的呢.
這個不一定是講別人.
可能是講我們自己.
沒有信心的人是怎樣的呢.
大家想一想.
很多疑惑.
上帝是否真的在這裡.
上帝你是否真的有能力幫我.
你是否真的會幫我.
我覺得上帝你應該不會幫我.
我覺得很沮喪.
很不開心.
今天我和A君講了一大輪.
和B君講了一大輪.
都是不太行.
我明天和C君又講一次我人生的經歷.
明天我又和第四個人講我人生的經歷.
從頭到尾不停述說.
不停不停講我人生很多的沮喪.
很多的困難.
我講完又講.
我講完之後的定律都是覺得.
上帝你不是很行.
只有你不幫我.
我人生是很慘的.
你知道我很慘嗎.
這個就是沒有信心的人.
有時候我們不要只看隔壁.
前後左右.
我們要看我們自己.
我們是否真的自己.
沒有信心就是這樣.
我覺得上帝不會幫我們.
有時候沒有信心到一個地步.
聽別人講見證都懷疑.
別人出來說我病得醫治.
聽說病得醫治這麼神奇.
你們想想.
是不是真的.

$^{281}$這麼神奇.
耶穌可以醫治你.
你有沒有想過這個法治基督徒內心的疑惑.
我沒有誇張.
真實發生在我們當中的事情.
我怎會懷疑.
聽別人講見證.
多神奇.
上帝的大能彰顯.
我們不是讚美上帝.
而是.
是嗎.
真的.
耶穌是行這神蹟.
不是騙人.
我心想.
這個真的沒有信心.
到一個地步很疑惑.
你除了懷疑別人講見證.
還懷疑什麼.
懷疑上帝.
為什麼懷疑上帝.
因為我自己經歷不到.
所以上帝不是真的.
上帝不是幫我的.
其實這樣很危險.
是不是.
大家看不看到這幅圖畫.
沒有信心的人是什麼.
疑惑的人是什麼.
鎖鏈.
鎖著腳.
看不看到.
大重錘.
寫著英文.
SIN 罪.
不信.
不信是罪.
其實不是困難.
是罪拿命.

$^{321}$我患病也不是罪拿命.
我沒錢也不是罪拿命.
我沒水也不是罪拿命.
我沒房子也不是罪拿命.
沒得洗澡也不是罪拿命.
罪拿命是不信.
我懷疑上帝.
我懷疑的時候.
這個就是鎖鏈.
其實是我的不信罪鎖著我.
我走也走不到.
我會灰心沮喪.
撒旦的頭這樣打敗我們.
這是很慘的.
最慘的是什麼.
如果我是一個不信耶穌的人.
人家會說.
靠自己力量.
信耶穌的人竟然沒有這份信.
我疑惑.
我懷疑上帝是不行的.
我懷疑上帝是沒能力的.
我祈禱又不見你馬上回應.
我祈禱你馬上修好水管.
我明天就想有房子.
可以嗎.
上帝就停在那裡.
祂的時間表和我的時間表不同.
但上帝的能力從來都沒有改變.
但我們會疑惑祂.
就是這樣.
這不是誇張.
是真實的.
還有我們像什麼.
就像船在波浪中飄來飄去.
我真的要倒了.
船要入水了.
我會淹死的.
你想想2023年全年到現在.
我們生命的光景是怎樣.

$^{361}$是充滿信心迎接挑戰.
還是每天都死死走.
別人跟你講見證.
我還要懷疑.
假的.
是真的.
裝模作樣.
上帝要透過他們見證.
自己的榮耀.
我懷疑他們.
這個不信去到這麼大.
如果是這樣.
我們的生活會很慘.
每天都打敗仗.
這樣的生活是怎樣的.
在這裡.
第一件事很重要.
我們需要放下.
不信的惡心.
弟兄姊妹我想說.
上帝很愛我們.
祂的能力.
不容置疑.
要行什麼神跡.
神跡對神來說從來都沒有難處.
要扭轉人生命從來沒有難處.
難的就是我和你沒有信心.
只是這裡.
所以第一件事要我們放下不信的惡心.
希伯來書三章十二節.
這樣說.
弟兄們你們要謹慎.
免得你們中間有人存著邪惡不信的心.
離棄了永生的神.
你不要以為.
不信就是信心軟弱.
你體諒一下我.
我經歷少.
我不明白.
但我想說神體是一件大事.

$^{401}$是罪.
弟兄姊妹.
我對自己說也對大家說.
免得你們中間有人存著邪惡不信的心.
離棄永生的神.
所以希伯來書三章十三節這樣說.
總要趁著還有今天.
天天彼此相勸.
我們彼此勸對方.
很重要.
免得你們中間有人被罪迷說心腸剛硬了.
什麼剛硬.
不信的剛硬.
神很憎人不信.
希伯來書是大悲觀.
大悲觀.
希伯來書特別強調.
希伯來書三章十四節這樣說.
只要我們將起初確實的信心.
堅持到底.
起初我們信主的時候.
我們是否很火熱.
心裡很雀躍.
我今天信耶穌.
我很願意.
生命改變.
我就禱告.
上帝就改變我.
你記不記得.
我們當初是這樣的.
信著信著就開始疑惑.
當然了.
苦難多了.
懷疑了.
苦難時間長了我就會懷疑.
但神在此鼓勵我們.
起來我們互相鼓勵.
你要堅持.
堅持我們要抗拒.
下一次在我們心裡.

$^{441}$聽別人說見證的時候.
如果我們再聽到聲音.
跟我們說是不是.
說說而已.
這個見證是假的.
哪有這麼神奇.
上帝不是真的.
如果我們心裡聽到這個聲音.
要奉耶穌基督的名抵抗它.
我相信冷漠者.
要誘惑我們.
要我們懷疑上帝的作為.
要奉耶穌基督的名抵抗它.
我們頂著頂著記著.
不要被它影響我們.
我們堅定相信上帝就是有能力.
神的喜悅我們有信心.
希伯來說.
很喜悅.
好像這幅圖一樣.
你張開你的手.
向上帝禱告.
你向祂呼喊的時候.
心裡就相信.
要做的事神要去成就.
但不要去強求.
上帝用你的方法.
你的時間表去成就.
上帝自有.
祂的方法和時間表.
交託給祂.
信得過.
希伯來書11章6節這樣說.
沒有信就不能討神的喜悅.
今天討神的喜悅.
不是做十幾二十幾.
三十幾樣事.
然後向上帝說.
上帝你看我做了多少事.
不是.

$^{481}$上帝喜悅我們有信.
你看看整本聖經.
貫穿在整本聖經裡.
上帝很喜悅我們有信.
沒有信就不能討神的喜悅.
今天每個禮拜.
回來崇拜.
如果我完了之後.
上帝的話都不能成為我們的幫助.
我們的生命就落入一個很可憐的光景.
每個禮拜崇拜之後.
不能令我的心更加貼近神.
更加能夠依靠神.
你出去紅欣堂之後.
又充滿沮喪灰心的時候.
你覺不覺得這樣很可憐.
上帝不想我們這樣生活.
上帝想我們活得有力.
所以要我們去崇拜.
沒有信就不能討祂的喜悅.
因為到神面前來的人.
必須信有神.
什麼叫信有神.
想問問大家.
就是神.
人們常常跟我說耶穌.
釘十字架.
什麼叫神.
超越萬有的上帝.
創造了我們.
我們的生命.
超越萬有的上帝.
創天造地的上帝.
我昨天栽培一班中學生的時候.
我跟他們說.
上帝多厲害.
連我們頭髮多少條都數過.
你媽媽有沒有數過你頭髮多少條.
不會吧 都數不到.
你覺得醫生會不會數你頭髮多少條.

$^{521}$不會吧 都數不到.
沒有人數過你的頭髮.
上帝數過.
數你的頭髮.
認識你整個人.
不是認識一個.
整間教會那麼多個.
他完全認識.
今天在看直播的.
弟兄姊妹上帝都認識.
全世界的人神都認識.
我們的上帝是一個.
超越的上帝.
我們信有神就不是.
普通的說話.
信一會就信.
你信一位怎樣的神.
是很超越的.
而且你信他.
會賞賜尋求他的人.
上帝會賞賜我們.
當我們尋求他的時候.
帶著這份盼望.
要帶著這份盼望.
要尋求他.
上帝會賞賜我們.
你有沒有期盼.
我真的很希望今天崇拜完之後.
我很有期盼.
這個星期很想更多的尋求上帝.
神好喜悅.
我們對他有信心.
我們要離開.
不信的惡心.
真實的信心是怎樣的呢.
希伯來書十一章一節.
剛才大家讀過.
信就是對所盼望之事有把握.
這是修訂版的翻譯.
和合本的翻譯.

$^{561}$在這個情況下.
信可能比較貼近.
信是什麼呢.
是所盼之事的實體.
實質.
盼望的事情.
是將來的事.
未實現的.
我現在未經歷到.
但我有盼望.
他成為一個實體.
是實際的.
就是因為這份信.
你有信心的時候.
就算我未經歷.
我未見到.
也好.
就是這份信心.
令我知道他真是真實.
是會實現的.
在希伯來書裡.
所說的盼望.
是什麼類型的盼望呢.
有幾個類型的盼望.
大家聽我稍稍讀一讀就行.
希伯來書九章十五節.
所說的盼望.
是所應許的永遠的產業.
是屬天的.
神要賜給我們永遠的產業.
神真是.
我們天上的爸爸很愛我們.
真的會賜給我們產業.
大家有沒有想過.
另外.
九章二十八節是說.
為了拯救.
熱切等候他的人.
就是完全的救恩.
今天我們已經得蒙救贖.

$^{601}$但救贖的恩典還未完全.
要到耶穌基督.
回來那一天.
這一切就徹底的完成.
這是我們盼望.
還有盼望什麼呢.
十章三十四節.
就是你們知道.
自己有更美好.
更長存的家業.
這是屬天的家業.
我們盼望這些.
十三章十四節就是說.
我們本沒有永存的城.
城市的城.
以至再尋求那將要來的城.
我們盼望那將要來的.
永恆的.
不會扭壞的.
那個屬上帝的城.
這些我們對永恆的盼望.
我們還沒有.
真的經歷到.
還沒有經驗到.
但因為那份信.
我們知道.
這個是會實現的.
在我們的.
屬靈的經歷裡我們知道是實際的.
不是騙人的東西.
是實際的.
有這樣的信.
就算還沒有經歷.
我都知道.
是真真實實的.
對未見之事有確據.
未見之事.
不只是將來的事.
現在也會有.
例如什麼呢.

$^{641}$大家有什麼問題嗎.
我不敢說完全沒有.
大家見過神了嗎.
我不敢說.
當中可能有人見過神.
大家見過神了嗎.
說真的.
沒有吧.
大家見過神了嗎.
大家見過上帝創造天地的過程嗎.
沒有吧.
是呀.
上帝都答應我們很多事.
是呀.
上帝答應我們很多事.
說真的我們未見過.
但不是將來的事.
是現在的.
是呀.
上帝會賞賜.
上帝會審判.
上帝有一個很寶貴的身份.
祂是信實的.
還有呢.
大家有沒有見到耶穌基督.
在天父的右邊.
為我們代求.
看到祂怎樣為全世界各地的人代求.
大家有看到嗎.
大家看不到吧.
但聖經告訴我們.
耶穌基督在天父的右邊.
長期為我們祈禱.
大家看不到.
但大家聽過.
這些都是未見之事.
還有聖經提到天上的耶路撒冷.
其中的居民.
這些我們都未見過.
未見過.

$^{681}$還有不被震動的事.
不能震動的國.
很多事是現在已經存在.
我和你們都看不到.
是真的看不到.
在信耶穌之後.
上帝讓他的聖靈.
住在我們生命裡.
大家有看到嗎.
我們會經歷.
我們聽到聖靈的聲音.
但是.
如果說肉眼能夠看到的東西.
其實很多我們都看不到.
很多屬靈的事.
但那些東西已經現存.
但那份信心.
對於這些未見之事.
我都相信是真確的.
是真實的.
但我不會懷疑.
信心就是這樣一回事.
如果我只是.
每次都要看到那些東西.
我才相信.
那些就不是信心.
是吧.
所以.
是吧.
說真的.
就算不要說屬靈的事.
你們看不看到空氣.
看不看到空氣.
看不看到氧氣.
看不到.
但你們知道它存在.
是吧.
但我說的神.
我們真的看不到.
信心.

$^{721}$讓我們知道.
上帝是真實的.
信心就是這樣.
我們需要有這份真實的信心.
你看不到.
知道上帝是真實.
你看不到.
但你知道上帝會賞賜我們.
你看不到.
但你知道神很想我們有信心.
你害怕的時候.
你看不到神.
但你知道神會保護我們.
如果說保護我.
這麼多年都可以說很多.
保護我的經歷.
我通常都看不到上帝.
但我真的經歷到.
上帝保護我.
是否在過程中.
沒有人可以做你.
24小時的保鏢.
有些人常常陪伴你.
常常陪伴你.
也不代表他可以常常保護你.
但我們的神就可以.
有信心.
信得過.
他會這樣做.
信得過他在.
那間屋.
在斷水管後.
很長時間.
上帝都沒有告訴我.
幾年幾月幾日.
我就會修好水管.
你的屋就會修好.
沒有說.
他也沒有說.
之後就修好.

$^{761}$透過什麼人.
什麼事.
他都不會說.
但我就禱告.
在過程中.
我的信心也有起伏.
但我記得一直堅持.
我和我姐姐.
一直堅持禱告.
經歷很多風浪.
經歷惡人.
在過程中很多難阻.
一直禱告.
到後來我宣告.
我的屋是屬於上帝.
我將它交給你.
你作主.
上帝你作主.
你掌權.
如果不是出於你允許.
任何人都不想得到那間屋.
因為後來我們發覺.
原來有些村民想逼我們.
減便宜賣出去.
真的去到.
一個這樣的地步.
但在過程中.
我們堅持禱告上帝.
我見不到神.
我不知道神的安排.
但我相信.
上帝是我們一家之主.
我相信上帝會為我們掌權.
會為我們伸冤.
我相信上帝.
也會眼目看顧這間屋.
所以一直禱告.
我和姐姐也是這樣.
一直為這間屋禱告.
堅定的相信.

$^{801}$神會工作.
我鼓勵弟兄姊妹.
也是這樣.
在這裡.
聖經給了很好的例子.
希伯來書十一章二字.
古人因著這信獲得了讚許.
一份對上帝的信.
得到上帝的稱讚.
神是很喜歡我們有信心.
很喜歡.
在過程中.
古人做得很好.
做了榜樣.
因為信上帝而得到稱讚.
給大家看一些例子.
這裡是講羅亞的例子.
因著信.
羅亞夢神指示他.
要做方舟.
當時神說會命令洪水毀滅這個世界.
有誰會相信呢.
我相信世界上的人.
也不會相信.
怎麼會相信呢.
這麼無稽的事也有.
你怎樣呢.
洪水毀滅的世界這麼大.
你能浸得完嗎.
然後神說叫羅亞建方舟.
這件事是前面沒有人做過的事.
沒有例子讓你跟進.
做方舟這麼大.
還要將動物帶入方舟.
傳流品種.
你想想對於羅亞多震撼.
由建方舟到後面.
要動物帶入.
這件事是否需要很大的信心.
如果前面有前人走過.

$^{841}$做過那件事.
你跟他走.
就說有板給你看.
完全沒有板.
就是上帝這樣吩咐.
你就相信他.
你怎知道之後的洪水.
是怎樣的情況.
很多細節神不會跟你說.
你就要相信他.
在過程中.
羅亞有敬畏的心.
信心背後.
也能看到人是否敬畏上帝.
懷疑上帝的人.
對上帝.
也沒有敬畏的心.
所以為何神這麼不喜悅.
說是不信的惡心.
因為過程中.
沒有敬畏的心.
但羅亞是敬畏上帝.
令他全家得救.
所以希伯來說.
直至他定了世代的罪.
世代還是不信的.
羅亞建了這麼久的方舟.
人們明明看到.
也沒有人說.
羅亞我們一家都要進方舟.
跟你一起逃避洪水.
沒有.
可能大家都笑他.
那個傻子.
建方舟.
那些動物要一起進去.
傻傻的.
精神不正常.
不信.
所以羅亞的信.

$^{881}$反而定了世代的罪.
羅亞自己一家.
承受了重新而來的疑.
羅亞承受了這個疑.
這是一個很大的福氣.
全家得夢拯救.
上帝真的會想賜人.
這是什麼例子呢?.
騎駱駝的什麼人呢?.
這是說亞伯拉罕.
因為信亞伯拉罕蒙召的時候.
就遵命出去.
往將來要承受為基業的地方去.
他出去的時候.
都不知道往哪裡去.
我們人很多時候都要這樣.
先講清楚.
你去哪裡.
怎麼去.
要用多少天.
你講清楚我就信你.
但上帝就不會講清楚.
想你一步一步跟著他.
在那個過程.
亞伯拉罕出去的時候.
真的不能詳細知道地點在哪裡.
但不簡單.
你看那幅圖.
為什麼.
一家大小.
財產又大又小.
全部浩浩蕩蕩出發.
你又不知道去哪裡.
是不是很要命.
真的很要命.
聖經稱他為信心之父.
不簡單.
不要這樣說.
他移民而已.
所有家當.

$^{921}$你又不知道去哪裡.
你移民去英國.
我不知道.
不知道下一步神要我去哪裡.
法國 加拿大.
我不知道.
家當又大又小.
又大又小.
就是信心.
上帝答應他.
呼召他出來.
有一個很寶貴的英雄.
所以這個很重要.
很重要.
在過程不簡單.
那份信.
所以古人在信上做了一個很美好的見證.
得到上帝的稱讚.
在你人生的經歷裡.
是上帝為你度身訂做.
正如他為我度身訂做.
一定有信心考驗.
弟兄姊妹.
我們不要想避開.
也不要想著.
上帝最好快點挪走我的困難.
挪走我的痛苦.
不用質疑.
好好學習這個功課.
要信他.
看不到也要信他.
你不知道他詳細安排也要信他.
因為有見證人.
見證人.
到最後.
三節說.
因著信我們知道.
這個宇宙是直神的華造成的.
拿一個例子.
什麼叫未見之事.

$^{961}$的確具.
你沒有見到這些東西.
但你知道它是真的.
給你一個例子.
你因著信.
你知道宇宙是直神的華造成的.
剛才我問你.
你有沒有見過.
上帝創造宇宙的過程.
他用他的說話.
去造成的.
或者透過他的說話.
帶來整個宇宙的秩序.
整個過程其實沒有人見過.
透過聖經.
你知道.
是神創造天地.
神透過他的說話.
去整理整個大地.
上帝的說話.
在這裡工作.
你透過聖經知道.
但我沒有見過.
但因為這份信.
我信.
神就是這樣造成.
整個宇宙整理大地.
我信上帝.
就是這樣行不同的神蹟.
就是這份信.
是何等的重要.
在起初神創造天地.
還有很多奧秘.
根本沒有人.
看到過.
沒有人.
上帝只是.
自己去創造.
天地萬物.
經歷一些描述.

$^{1001}$啟示地面上.
熟悉他的人.
讓熟悉他的人.
寫出來.
我們看到那些文字.
憑著信去相信.
上帝就是這樣造天地萬物.
上帝就是這樣行事奇妙.
但這份信是很重要.
是很重要.
我和你需要.
有這份信.
這樣看得見的.
是從看不見的.
造出來的.
看得見的是整個宇宙.
宇宙萬物.
造出來的時候.
是從看不見的.
造出來的.
神的創造.
是很奇妙.
因為信.
我們相信是這樣.
我都說.
如果每樣東西.
我都要看到才相信.
我們真的沒有了那份福氣.
上帝想我們有份單純的信.
今天我們需要回轉向小孩子.
小朋友.
你和他說的時候.
你有沒有發覺小朋友的心很單純.
他就相信.
耶穌行神蹟.
很厲害.
他禱告的時候.
他會很有力.
求耶穌醫治你.
求耶穌幫助你.

$^{1041}$相信耶穌可以幫助我們.
小朋友的信很單純.
我和你.
為什麼沒有了那份單純的心.
日子內.
心很容易蒙了油.
是不是開始我們很靠自己.
這麼久做人.
都是靠自己的聰明智慧.
靠自己能力.
每樣東西都看得住.
所以靠自己.
讓我們回轉向小孩子.
一個單純的信.
真實的信.
看不到,但我相信.
我相信上帝在.
我相信上帝愛我們.
上帝愛我們.
剛才我在祈禱會說.
上帝愛我們.
不需要我們不斷地做很多事情.
證明我是怎樣怎樣好.
不需要.
上帝很愛我們.
很想我們知道他很愛我們.
心願我們都相信.
他愛我們.
讓我們回應他的愛.
我相信這樣.
人生是很不一樣.
我記得去年.
一直在禱告.
禱告上帝.
我都不知道上帝什麼時候會.
幫我們處理水管的問題.
一直有很多反對的意見.
很多的難阻.
我們都做不到水管.
在過程中.

$^{1081}$就禱告禱告禱告.
到一個時間.
很奇妙,神的時間到了.
上帝就開路了.
為我們預備一些人.
讓我們可以見到一些人.
可以和他們談.
事情就很奇妙地.
一步一步地.
開路.
水管可以在某個位置.
接駁到.
水務處為我們預備的位置.
甚至乎.
水務處真的很好.
在整件事上我很感恩.
但是神的恩典.
水務處對我們有很大的幫助.
在整件事上.
再加上.
在村民裡.
有些人的協助.
結果可以搭一條新的路.
水管.
就是這樣落成.
到去年年尾的時候.
整座的水管.
重建好了.
如果接下來再有.
斷水管的時候.
是很容易修理.
不需要再經過.
某某私家地.
沒有人再難阻我們.
還有我記得.
有一天很奇妙.
上帝在我.
洗澡的時候.
不知道為什麼.
神的聲音就臨到.

$^{1121}$忽然間就說.
其實不需要很多錢.
就會修理好.
就這樣說了.
到後來.
我們回想.
是真的便宜了很多.
修理水管.
是可以很貴.
非常貴.
但在過程上.
神的神奇.
幫助了我們很多.
所以可以.
以一個非常合理的價錢.
修理好.
而且最奇妙的一件事.
我記得.
在一月份.
斷了水管的時候.
我不知道何去何從.
忽然間.
無家可歸.
神就很奇妙地為我預備了一個家庭.
我相信.
我人生下半場.
都會和這個家庭一起生活.
這個家庭帶給我很多照顧.
我很感恩.
上帝透過這個家.
給我很多很多的愛.
很多的呵護.
在過程中.
對我自己的屬靈生命.
有很大的幫助.
見到上帝很疼愛我.
很多事.
我不知道.
很多事我看不到.
但神已經預備.

$^{1161}$而且神讓我看.
祂想我搬去元朗住.
如果不是去年.
全年的時候.
我們有同工.
在很大的病患當中.
我怎可以承擔事奉呢.
真的很遠.
所以我很明白.
住上水粉嶺區的弟兄姊妹.
是很不簡單.
他們很努力.
過來教會聚會.
其實交通時間很長.
在忙碌當中.
絕對是很大的壓力.
但去年這樣全年.
在元朗居住.
省下很多事奉的時間.
我很感恩.
神讓我看.
祂很奇妙.
因為上帝固然會為整件事開路.
上帝要我看.
要我看.
祂的工作.
祂在我身上做什麼.
我相信今天.
神也要大家看.
要你們放手.
相信祂會幫助你們.
相信祂會賜福氣給你們.
你覺得很沮喪.
很灰心的地方.
很想大家相信祂.
而且要給大家看.
祂有很奇妙的安排.
上帝對我們每一個人.
都有很奇妙的計劃.
所以我很盼望.

$^{1201}$今年大家有一個很豐富的一年.
很深刻的一年.
要有真實的信心.
我們一起來祈禱.
天父上帝.
真的要放下我們.
不信的惡心.
因為你喜悅我們.
有信心.
你很想賜福氣給我們.
你愛我們.
你很想我們經歷你.
很想我們知道.
神對我們每個人的人生.
有很奇妙的安排.
你的道路.
高過我們的道路.
你的意念高過我們的意念.
你需要我們.
不是我們拼命在做什麼.
你需要我們.
只是要對我們有信心.
求神賜給我們弟兄姊妹.
我們眾同工.
有一顆真實的信心.
今年活得不一樣.
求主帶領我們.
奉耶穌基督的名字.
祈禱.
阿們.
謝謝大家收看 再見.
\newpage



\section{詩篇 118:1-13}
\label{sec:EkAEiOY56rU}
\textbf{2024.06.09 《與神同行的人生》 詩篇118\_1-13 (和修版) 曾裔貴牧師}
\newline
\newline
連結: \href{https://youtube.com/watch?v=EkAEiOY56rU}{\texttt{https://youtube.com/watch?v=EkAEiOY56rU}} ~~~~ 語音日期: 2024-06-13
\newline
\newline
\hyperref[sec:QSB7W2xKaig]{\small{< < < PREV SERMON < < <}}
~
\hyperref[sec:index_chronic]{\small{[返順時目]}}
~
\hyperref[sec:index_scriptual]{\small{[返順卷目]}}
~
\hyperref[sec:WqOKT5MNC_k]{\small{> > > NEXT SERMON > > >}}
\newline
\newline
詩篇 118:1-13
\newline
\begin{longtable}{cl}
\hline
\hline
章節 & 經文 (和合本修訂版)\\
\hline
118:1 & \begin{tabularx}{0.7\textwidth}{X} 你們要稱謝耶和華，因他本為善；他的慈愛永遠長存！ \end{tabularx} \\ \\ \relax
118:2 & \begin{tabularx}{0.7\textwidth}{X} 願以色列說：「他的慈愛永遠長存！」 \end{tabularx} \\ \\ \relax
118:3 & \begin{tabularx}{0.7\textwidth}{X} 願亞倫家說：「他的慈愛永遠長存！」 \end{tabularx} \\ \\ \relax
118:4 & \begin{tabularx}{0.7\textwidth}{X} 願敬畏耶和華的人說：「他的慈愛永遠長存！」 \end{tabularx} \\ \\ \relax
118:5 & \begin{tabularx}{0.7\textwidth}{X} 我在急難中求告耶和華，耶和華就應允我，把我安置在寬闊之地。 \end{tabularx} \\ \\ \relax
118:6 & \begin{tabularx}{0.7\textwidth}{X} 耶和華在我這邊，我必不懼怕，人能把我怎麼樣呢？ \end{tabularx} \\ \\ \relax
118:7 & \begin{tabularx}{0.7\textwidth}{X} 在那幫助我的人中，有耶和華幫助我，所以我要看見那些恨我的人遭報。 \end{tabularx} \\ \\ \relax
118:8 & \begin{tabularx}{0.7\textwidth}{X} 投靠耶和華，強似倚賴人； \end{tabularx} \\ \\ \relax
118:9 & \begin{tabularx}{0.7\textwidth}{X} 投靠耶和華，強似倚賴權貴。 \end{tabularx} \\ \\ \relax
118:10 & \begin{tabularx}{0.7\textwidth}{X} 列邦圍繞我，我靠耶和華的名必剿滅他們。 \end{tabularx} \\ \\ \relax
118:11 & \begin{tabularx}{0.7\textwidth}{X} 他們圍繞我，圍困我，我靠耶和華的名必剿滅他們。 \end{tabularx} \\ \\ \relax
118:12 & \begin{tabularx}{0.7\textwidth}{X} 他們如同蜜蜂一般地圍繞我，他們熄滅，好像燒荊棘的火；我靠耶和華的名，必剿滅他們。 \end{tabularx} \\ \\ \relax
118:13 & \begin{tabularx}{0.7\textwidth}{X} 你用力推我，要叫我跌倒，但耶和華幫助了我。 \end{tabularx} \\ \\
[1ex]
\hline
\hline
\end{longtable}
$^{1}$(主席向觀眾展示禮品).
領子們一起同心祈禱.
我們可以來到有聖餐的主日.
我們剛才來到以聖餐的餅和杯.
紀念主耶穌十字架上的大愛.
和這份恩典.
我們雖然在天雨下了很久.
好像沒有停過.
但是我們能夠憤外感受到.
主耶穌在我們的生命裡.
無論是天晴或雨天.
或是我們生活順暢的時候.
或是有不如意的時候.
我們的主這份救贖恩典.
這份情愛沒有改變.
幫助我們能夠透過信仰.
使我們的生活能夠有與神同行的人生.
我們今早等候在你的面前.
也有頂子妹在我們中間一同虔誠敬拜.
很盼望這個屬靈的家.
一起在這個地方為主的福音.
來慶旺我們這樣感恩禱告.
奉耶穌基督的名.
阿們.
很久沒有講詩篇.
這篇詩篇其實相當重要.
有極多的解經家告訴我們.
《馬太福音》的二十六章.
主耶穌唱了詩之後.
在最後的晚餐.
祂就出去.
唱《哈利路亞詩》.
就是這一篇118篇.
其中因為詩篇很長.
這個段落.
22,23,24節.
今天就不夠時間去講.
很重要的匠人所棄的石頭.
已經成為了房閣的頭塊石頭.
是誰做的呢?.

$^{41}$在我們眼中看為稀奇.
這是耶和華所定的日子.
我們在其中高興歡喜.
如果去參加婚禮.
第24節就很受歡迎.
很多人就用那節聖經.
作為他們婚禮的重要開始.
不過這節聖經其實是在講.
那些完全不顧耶穌的人.
對祂的生命的蹂躪.
耶穌唱這首詩.
接著走出去準備接受.
在赫西瑪利園猶大的出賣.
在這首詩歌中所說的.
祂就是被棄的石頭.
但是令我們稀奇的就是.
這塊被棄的石頭.
成為了房閣的頭塊石頭.
所以這篇詩篇.
我一直去查一些資料.
原來很重要.
在500年前的馬丁路德.
他說這篇詩篇.
是他整個人生.
他唯一喜歡的詩篇.
這篇詩篇盛載著的.
當然極多的解經家.
都認為是大胃王.
他經過了很多危難之後.
寫的這篇詩.
其中有第五到十三節.
講述他的生命故事.
所以差不多馬丁路德.
可以對號入座.
他在面對宗教改革.
他寫了的九十五條之後.
輾轉受到教廷.
他所愛他所尊敬的領袖.
當時稱為教皇.
或者是地方的.

$^{81}$在威登堡教堂的領導人.
向他發出最後的通牒.
要麼收回對教廷的指控.
要麼被趕逐離開教堂.
馬丁路德最有名的言論.
除非你在聖經裡.
找到我有什麼錯誤.
否則願主幫助我.
就是這篇詩的靈感.
很寶貴.
所以我們今天.
跟大家去看看一到十三節.
將來有時間.
我們再看下一個半的一段.
教會.
如果大衛王.
他到了一個有人生經歷.
很多生命故事的人生階段.
他緊張國家什麼呢?.
一到第四節.
已經充分反映大衛王的內心.
要人學習.
學習神是怎樣的神.
你們要稱謝耶和華.
因為祂的良善.
因為祂的慈愛.
永存的.
我們就這樣唱詩歌.
一點都不難.
但是要在生活.
在你實際的處境.
去體驗耶和華本為良善.
祂的慈愛永遠常存.
不只是大衛.
他說願以色列所有神的恩若.
摩西之若的子民.
即猶太人.
都要像他這樣.
能夠這樣說.
神的慈愛永遠常存.

$^{121}$不單止.
收窄範圍.
你是以色列人.
一個星期才上去聖殿.
但是在聖殿裡工作的人.
他就說願亞倫加.
都要這樣認定.
祂的慈愛永遠常存.
不單止.
還有所有敬畏神的人.
都要這樣宣告.
在禮拜堂.
在聖殿.
當時還沒有聖殿.
因為他的兒子.
所羅門起聖殿.
在會幕.
當時的侍羅會幕.
但是要每一個人.
在生活裡.
在人生裡.
都經歷.
都經驗.
都能夠宣告.
祂的慈愛永遠常存.
這個就要很多時間.
每個人有不同的故事經歷.
如果大家有機會.
看新聞.
你們會知道.
前幾天.
有一宗轟動的新聞.
有個媽媽.
在保安道圖書館.
教女兒.
做功課.
我忍了你十分鐘了.
你還是不懂.
然後馬上捏她的耳朵.
都不夠了.

$^{161}$那個火還沒降.
就趁她後樓梯.
不斷打.
然後不知道為什麼會有人拍到.
就抓了.
她很後悔.
為什麼我會這麼衝動呢.
大家有沒有經歷過呢.
搭地鐵.
有些人因為搶錯位置.
在那裡吵架.
吵架然後就說髒話.
整個車就很麻煩.
有沒有經驗過這些.
有個神是良善.
祂有永遠慈愛.
最近有一位弟兄.
給我看了一條片.
是巴勒斯坦哈馬斯.
底下的那些.
沒有死的小朋友.
他們那些驚慌的樣子.
很多個在那裡不停地震.
因為那些炸彈的轟炸.
你想想如果你經歷到這些難民.
你見到這些事情.
你能不能夠說.
祂是本為善呢.
祂的慈愛永遠長存呢.
如果你親身經歷這樣的事情.
我在我姐夫的安息禮拜整個過程.
我的外甥.
她是一個藥劑師.
在病床.
我去到和我的姐夫.
輕輕地跟他說.
他馬上眨眼.
我的外甥就很激動地跟我說.
因為我躺在姐夫的床邊.
去到靠近她的身邊.

$^{201}$我的外甥就站在那裡.
看到.
她很激動.
為什麼爸爸會醒呢.
原來我去之前已經昏迷了三天.
我下飛機之前.
所以她很激動地看到.
爸爸有那個反應.
我再呼喚她.
跟她決志.
她又眨眼一次.
很清楚地聽到我想跟她說.
帶她信主.
我想這個經歷對我的外甥.
應該有很大的意義幫助.
她也看到整個安息禮拜的安排.
已經完結了.
她開車載我去一個大商場買東西.
有機會只有兩個人.
姐姐不在.
因為姐姐在那裡說話.
我們沒有機會說話.
外甥開車.
我說.
姐夫信得很好.
真的看到神的恩典.
看到她幾次眨眼.
我的外甥沒有否認.
接著我問她.
不如你也有機會這個星期回教會吧.
她很清楚.
當然她很小聲地跟我說.
我是不會回家的.
我說為什麼你不回家呢.
她說九年前她的姐姐.
也是這樣肺癌死了.
現在我的爸爸剛剛退休兩年.
也是這樣不明不白地.
突然間肺癌死了.
她說.

$^{241}$如果你說的神是真的.
為什麼會這樣呢.
她說我不能夠接受.
一位這樣沒有慈愛的神.
給我有很大的反思.
我當然不能代表神回答她.
我絕對不能夠說.
其實不是的.
神是慈愛的.
我不能夠這樣回答她.
但是在她的親身經歷.
她說你的神我不會相信.
因為我感受不到.
祂是有慈愛.
是有良善.
就是這樣的一個.
很清楚對神的否定.
我的外甥.
弟兄姊妹.
當然這是很正面地告訴我們.
作為一個國家的領袖.
作為一個領導的人.
是要呼喚所有人.
無論你的經歷.
你現在可以解釋.
不可以解釋.
每個人一定是每天.
發生著一些事情.
這些事情慢慢.
會在你的生命裡面.
成為你自己的故事.
每個人都有自己的故事.
所以為什麼第五到十三節.
大衛王.
作詩歌的人.
就講他的故事.
但是在講之前.
他要讓國民.
來看到.
到底你認不認同.

$^{281}$我們這一位做物主.
他是本為良善的.
還有他的慈愛.
一直到永恆.
要認同.
要在生活裡面.
無論你看到的東西.
是多麼不如意.
是多麼不符合你的想望.
但是你要學習.
認定這個信仰的衝擊.
往往在青年的階段.
就最大了.
很多一大學畢業出來工作.
那個衝擊就很大了.
使得他很難去接受.
小時候回教會.
完全沒有問題.
完全不會有任何的質疑.
但是到他真的可以去思想.
開始看到人生的時候.
例如很多的不公平.
例如上司也是基督徒.
那就更加慘了.
如果上司也是跟你回一間教會.
就更慘了.
耶和華本為善.
唱歌的時候.
突然間想起他.
他還是幫手事奉的.
我不是太想講太多的東西.
因為我去到士嘉堡的華人宣導會聚會.
有一個弟兄.
就是我的外甥女的團契.
他提到他正正就是面對.
一些他覺得不可能的事情.
太白不可以了.
因為好像有些具體就不好了.
不過真的.
我們每個人都必須在信仰裡面.

$^{321}$要學習認定.
我們的神.
祂本為良善.
祂的慈愛永遠長存.
你可以想像得到.
有什麼證據可以來看到.
神本為善呢?.
從第五節開始.
因為一到第四節.
是要讓不同的持份者.
你只是恩若的子民.
例如你是猶大之派.
你不是利美之派.
你是不可以參與聖殿的事奉.
你是一個只不過我每天日出而作.
日入而息.
對你什麼的獻祭我完全不懂.
我只懂得帶羊去獻祭.
其他一切一概我不懂.
但是我知道耶和華本為善.
我要定期帶我的小朋友.
家庭和家族去聖殿守節.
一個以色列人就是這麼多.
因為知道自己是神的選民.
但是亞倫加就不止了.
每個星期參與事奉.
做錯了一些事.
可能就會招致神的審判.
這就是亞倫加的生命.
他盛載的耶和華本為善.
是不同的.
有事奉的弟兄姊妹你會明白.
參與事奉的那一刻.
你要憤愛.
你要審查自己.
讓神在很多人的生命裡得到建立.
願敬畏耶和華的.
都來感受到神的慈愛.
你可以想像到大衛王.
他是多麼的體重.

$^{361}$不是他個人.
是整個國家.
這個國度.
神的管轄的地方.
所有的人.
都要對神的屬性.
有掌握,認識和宣告.
弟兄姊妹你回來教會.
你就是要學這些真理.
令你心悅誠服.
對神在你的生命裡的工作.
你只能夠宣告.
祂真是本為善.
祂的慈愛永遠的長存.
又不是無的放矢.
不是只唱歌.
因為有實際的生命故事.
第五到十三節.
出現了很多次.
我求告耶和華幫助我.
耶和華幫助我.
在投靠的人當中.
唯有耶和華幫助我.
這樣來到.
好像呼之欲出地.
他找過不同的幫助.
經歷過不同的難處.
但是唯有耶和華.
才是真正唯一可靠的幫助.
這樣的意思.
重複地說.
那個倚靠,那個幫助.
說了很多很多次.
在第五到十三節那段.
說了很多次.
提到急難.
有沒有試過急難?.
急到你沒辦法.
你只能夠很快祈禱.
疫情前.

$^{401}$星期一.
跟媽媽和飛鴻.
出了九龍喝茶.
喝完茶.
很大雨.
差不多是黑雨.
開車回元朗.
進長青隧道.
很自然地將水撥調低一點.
因為準備出隧道之後.
再開回水撥.
準備出隧道.
長青隧道進元朗.
準備出隧道.
在走快線.
開回水撥.
完全失靈了.
是直接壞了.
你可以想像得到.
在快線,公路,水撥.
沒有了.
傾盆大雨.
完全看不到面前的路.
你是不可能再開車了.
停在快速公路.
黑雨的時候.
一直不敢走.
停下了不敢走.
後面的車子如果看不到就撞過來.
一秒時間.
馬上祈禱.
簡直是害怕.
因為我自己一個.
再加上媽媽,飛鴻.
害怕到一個地步.
你沒有想像的那種.
是裡面震出來的害怕.
一秒而已.
突然間999.
馬上打999.

$^{441}$都不用五分鐘.
完全不用五分鐘.
那個對方的接線員.
安撫我.
我很害怕.
現在我在快速公路的快線.
我的水撥壞了.
他說你不用怕.
馬上有人來了.
都不需要五分鐘.
就馬上有工程車.
不知道從哪裡來的.
就一直來到.
就引我.
開了那些死火燈.
帶我去工作的地方.
就解決了這個急難.
我小時候.
我的兒子很小的時候.
三歲還是四歲.
小朋友.
完了崇拜.
準備讀一學.
走來走去.
因為我們只有一層.
我以前的舞會.
只有一層.
幾個小朋友.
不知道怎樣走來走去之後.
玩耍.
突然間有個小朋友.
衝出來叫我.
就說曾牧師.
那時候我不是牧師.
就叫我快點進去房間.
原來我的兒子.
手指.
不知道為什麼.
夾在門的轎子上.
如果有人一拉.

$^{481}$就斷了.
我見到之後.
我害怕到一個地步.
很害怕.
不知道怎麼辦.
當然當時沒有祈禱.
馬上就想辦法.
叫人不要動門.
然後慢慢安撫我的兒子.
叫他把手慢慢伸出來.
如是者就解決了那個急難.
親愛的弟兄姊妹.
大衛說.
我在急難當中.
求告耶和華.
祂就應允我.
把我安置在寬闊之地.
原來人生的苦難.
是讓我們的生命更加經歷.
祂的幫助.
帶來我們的生命更加寬闊.
原來基督徒.
我們不要逃避苦難.
不要逃避一些神要我們學習的人生的功課.
當然.
生老病死.
已經是人類的命運定律.
必定會有的.
好像.
你可以想像得到.
上星期.
有新聞報告.
在菲律賓.
有警察打劫.
穿著制服.
開著警車.
截停了私家車.
有三個中國人.
一個馬來西亞人.
叫他們把證件拿出來.

$^{521}$接著有另一個貨車衝過來.
那些蒙面人就抓他們.
把他們綁架.
幸好有兩個中國人逃脫了.
報警.
終於發現.
而且還是很高級的警察.
菲律賓.
所以.
一個女孩千萬不要去.
真的不要去.
不要開玩笑.
所以.
如果你在KK園看到就慘了.
是很恐怖的事.
但現在我不是想講政治.
不是想講罪惡.
我是想講人生.
其實.
不可能沒有苦難的.
但是屬神的兒女.
就有一個很重要的生命功課學習.
第五到第十三節.
不斷講.
神怎樣幫助.
神怎樣帶領大衛王.
他經過不同的水火.
不同的仇敵.
例如.
下面這樣說.
有些人是想他.
推他跌倒.
但是耶和華幫助我.
所以我看見那些恨我的人.
遭到報應.
蘇羅.
就是這樣死了.
亞沙隆.
他的親生兒子.
第三個兒子.

$^{561}$就是這樣遭遇神的審判.
報應.
他誇耀自己的頭髮.
每年才剪一次頭髮.
從宜遷到舍克勒.
這個亞沙隆.
你這麼誇耀自己的才華.
你的樣貌.
比爸爸更加的俊美.
他怎樣死的?.
騎著小驢.
經過那些密林.
撬著頭髮.
然後被虐壓.
用槍刺死.
這樣在仇殺中死了.
因為亞沙隆.
要反叛自己的爸爸.
要殺自己的爸爸.
你可以想像到.
大衛經歷了很多恩怨情仇.
大衛感受到的就是.
他所投靠的耶和華.
是讓他的生命去到寬闊之地.
他沒有錯.
他自己也有罪.
所以下面你會看到第十八節.
耶和華雖嚴嚴懲治我.
卻未曾將我交於死亡.
大衛可以有這樣的生命經歷.
他可以帶領國民.
給我闖開二門.
進去稱謝耶和華.
這就是一個有經歷的人.
不是沒有罪.
而是在神裡面.
罪得到赦免的一個經驗.
他仍然可以進入二門.
打開耶路撒冷的門.
當然他不是祭司.

$^{601}$他不能進入回望獻祭.
但是他是君王.
他可以在神赦免的公義裡面.
他能夠看到這是二門.
他可以進入.
所以一個有生命經歷的人.
即是有故事的人.
不是沒有罪.
而是在神裡面.
罪得到赦免.
而是能夠在救恩裡面.
看到神在我們生命裡面的作為.
一個領袖應該是能夠講故事的人.
不是動聽的故事.
是講生命成長的故事.
來了第六個月了.
今天已經是九號.
一月七號過來.
我們述志之後.
勵善牧師願意留守在母堂.
去年的九月.
內心一直擔心洪恩堂.
老牧師的身體越來越差.
我和幾位領袖分享.
我說曾姑娘自己一個人.
真的很辛苦.
因為我隔個星期就要來開同工會.
很多事情根本沒有人去幫她.
我有越來越重的負擔.
很想在洪恩裡面幫忙.
於是大家很認真祈禱.
我去完土耳其回來.
他們就約我.
因為他們給了我一個訊息.
如果你過來洪恩.
就會有一個很大的改變.
他們可以勸服勵善牧師.
我都可以加上勸服勵善牧師.
願意留守在母堂.
使我沒有後顧之憂.

$^{641}$我就可以在洪恩和大英節目一同去傳福音.
我在這幾個月裡面.
很多大英節目問我怎樣.
我說我從未試過這麼開心.
在每一天回來洪恩.
現在我能夠做到初中的同學.
和他們聊天.
和他們一起透過補英文.
接觸到他們.
有很多在這裡補習的年輕家長.
可以和他們坐下來聊天.
我們快要開始開講座.
讓他們可以探索福音.
看到街坊們.
看到我又叫我牧師.
很多人看到我們的改變.
我真的看到神在過去的感動.
一直讓洪恩家的弟兄姊妹一同去看到.
而一起去成長.
看到這就是神的一個很清楚的帶領.
我從來沒有試過.
這麼開心地做到一些中學生的工作.
因為大的母堂.
根本就已經有不同的同工.
做那個年齡層.
我們的中學生很多都是在幼稚園那邊聚會.
後來我離開之前.
堅持一定要有青少年的崇拜.
在自己的禮堂.
怎麼會禮堂沒有青少年崇拜呢?.
而用來做婚禮呢?.
不應該這樣的.
青少年一定要有自己的地方.
在自己的母堂的禮堂崇拜.
這件事發生了.
現在越來越多人青少年的崇拜.
是很感恩的.
所以真的看到.
現在我能夠在洪恩.
可以接觸到很多不同的弟兄姊妹.

$^{681}$不同的街坊.
簡直是一件很開心.
好像是最後的五年.
一個新的工場.
是我從來沒有試過.
這麼近身可以做到這些模樣.
以及這些全副工作.
是很開心的.
還有看到弟兄姊妹的生命的改變.
弟兄姊妹.
剛才說的急難第五節.
但是你留意第八,九節.
「投靠耶和華,強恥倚賴人」.
我這本舊的和合本更加精簡.
不僅是權貴.
更加強恥倚賴王子.
就是說大衛的人生.
都有很多不同的資源.
他可以來到去倚賴.
但是一個事奉者.
一個跟隨神的大衛.
他最終發現.
一切人.
一切的資源.
如果不是在神裡面.
其實都是不可靠的.
唯一可靠的.
就是這位耶和華.
好像一個對比的.
希伯來的詩歌.
來做一個這樣的對比.
讓你更加凸顯.
唯有神才是真正可靠.
真正我們要倚靠的這位耶和華神.
這就是大衛的一個經歷.
強恥倚靠人.
強恥倚靠王子.
就是說人和王子.
權貴都是不可靠的.
唯有耶和華才是真正的可靠.

$^{721}$所以在座如果你有身體的困難.
有生活的困難.
人際關係的困難.
你去倚靠神.
將你的情況交託給神.
讓神介入.
改變.
所以一個有故事的人.
一個去到有經驗過神的人.
其實從他的眼神已經觀察到.
因為眼是靈魂之窗.
耶穌說的也對.
眼睛就是身上的燈.
你的眼睛瞭亮.
全身就光明了.
你的眼睛如果昏暗.
全身就黑暗了.
那個黑暗何等大.
所以其實眼神反映到.
我們和神的關係.
眼神反映到.
我們是否真的有在神的倚靠裡.
而看到有很多出路.
很多幫助.
所以你可以想像到.
一個如果內心有很多的冤屈.
很多苦楚的人.
眼神都充滿了苦澀.
看到你都會覺得要走遠一點.
因為他的眼神很苦.
其實每個人都有不同的故事.
一定有的.
但是有沒有將你的故事.
帶到神那裡去倚靠他呢?.
在他那裡讓你有釋放.
慢慢地.
你心靈有釋放.
眼神自然就會明亮.
這是很重要的一個真理的教導.
最後.

$^{761}$第十四節和第二十九節.
就是反映了大衛王.
他內心這些故事.
如何化為詩歌.
耶和華就是我的力量.
是我的詩歌.
他也成了我的拯救.
第二十九節.
再一次來總結.
他想教導所有的回眾.
國民要留心的.
就是你們要稱謝耶和華.
因祂本為善.
祂的慈愛永遠長存.
各位這一首主題曲.
這一首耶穌基督.
在《馬太福音》最後晚餐.
來唱頌的詩篇118篇.
希望是你的生命詩歌.
這一首詩歌帶領馬丁路德.
經過很多的艱難.
很多當時的德國的一些政治領袖.
宗教的領袖.
對他的背叛.
對他的攻擊.
和追殺.
馬丁路德靠著這一首詩歌.
尤其是第八節.
成為他一生生命裡的力量.
投靠耶和華.
強恃倚賴人.
強恃倚賴王子.
當時他真的倚靠德國一個君王.
一個城邦的君王.
但是最終都幫不到他.
但是耶和華就是他的幫助.
很盼望我們作為教會的領袖.
和大家一起去學習.
我們的主是本為善.
他的慈愛永遠的長存.

$^{801}$我們一起祈禱.
感恩天父讓我們能夠.
可以安靜的去聆聽你自己的真理.
多謝你藉著大衛王.
來去譜寫他的生命的故事.
來去成為很多國民的學習.
不是他沒有做過任何錯事.
而且試過是錯得很嚴重.
但是神啊.
你仍然給他闖開二門.
因為十價的救恩.
主的寶血.
是洗乾淨他的罪.
藉著他獻祭.
向先知拿丹的認罪.
來去反映他內心敬畏神.
是願意歸向神.
無條件地去承認自己需要救贖.
多謝天父讓我們看到.
無論人間有多少的罪惡,痛苦.
但是你的本性.
你的屬性的良善和慈愛.
是永遠不會改變.
但願我們每一個弟兄姊妹.
我們都是屬神的群體.
我們都是敬畏主的人.
我們去認識你的慈愛永遠長存.
在我們的生活,在我們的生命裡.
以致將來我們都說我們自己的故事.
去勉勵其他人在主的愛裡成長.
感恩禱告奉耶穌基督的名.
阿們.
\newpage



\section{使徒行傳 6:1-7}
\label{sec:WqOKT5MNC_k}
\textbf{2024.06.16 《lor 上帝做藉口之:當有投訴時》 使徒行傳6\_1-7 (和修版) 張雲開老師}
\newline
\newline
連結: \href{https://youtube.com/watch?v=WqOKT5MNC_k}{\texttt{https://youtube.com/watch?v=WqOKT5MNC\_k}} ~~~~ 語音日期: 2024-06-17
\newline
\newline
\hyperref[sec:EkAEiOY56rU]{\small{< < < PREV SERMON < < <}}
~
\hyperref[sec:index_chronic]{\small{[返順時目]}}
~
\hyperref[sec:index_scriptual]{\small{[返順卷目]}}
~
\hyperref[sec:RaRAzKcNgEg]{\small{> > > NEXT SERMON > > >}}
\newline
\newline
使徒行傳 6:1-7
\newline
\begin{longtable}{cl}
\hline
\hline
章節 & 經文 (和合本修訂版)\\
\hline
6:1 & \begin{tabularx}{0.7\textwidth}{X} 那些日子，門徒增多，有說希臘話的猶太人向希伯來人發怨言，因為在日常的供給上忽略了他們的寡婦。 \end{tabularx} \\ \\ \relax
6:2 & \begin{tabularx}{0.7\textwidth}{X} 十二使徒叫眾門徒來，說：「我們撇下神的道去管理飯食，是不合宜的。 \end{tabularx} \\ \\ \relax
6:3 & \begin{tabularx}{0.7\textwidth}{X} 所以弟兄們，當從你們中間選出七個有好名聲、滿有聖靈和智慧，我們派他們管理這事。 \end{tabularx} \\ \\ \relax
6:4 & \begin{tabularx}{0.7\textwidth}{X} 至於我們，我們要專注於祈禱和傳道的事奉。」 \end{tabularx} \\ \\ \relax
6:5 & \begin{tabularx}{0.7\textwidth}{X} 這話使全會眾都喜悅，就揀選了司提反—他是一個滿有信心和聖靈的人；他們又揀選了腓利、伯羅哥羅、尼迦挪、提門、巴米拿，並皈依猶太教的安提阿人尼哥拉， \end{tabularx} \\ \\ \relax
6:6 & \begin{tabularx}{0.7\textwidth}{X} 叫他們站在使徒面前，使徒禱告後，就為他們按手。 \end{tabularx} \\ \\ \relax
6:7 & \begin{tabularx}{0.7\textwidth}{X} 神的道興旺起來；在耶路撒冷門徒數目增加得很多，也有許多祭司聽從了這信仰。 \end{tabularx} \\ \\
[1ex]
\hline
\hline
\end{longtable}
$^{1}$各位節目早安.
 很多謝洪恩唐邀請 曾牧師.
 邀請讓我這次.
 六十幾歲的人都會做一些以前沒做過的事情.
 有些人說.
 臨死之前要跳一次降落傘.
 搭飛機上去 單引擎飛機上去.
 跳下去 之類之類.
 今天就是我第一次.
 第一次進入洪水橋這個位置.
 未進過元朗屯門 天水圍都進過.
 但去到洪水橋這個位置就是第一次.
 我不知道大家.
 其實我香港很多地方都沒去過.
 雖然那幾年的疫情.
 香港人本來就喜歡出去.
 一向都喜歡出去.
 因為香港覺得自己是很小的地方.
 最重要是出去 飛出去 到處玩.
 我們去我們的後花園.
 台灣 日本那些是我們後花園.
 隨時去的.
 我就比較 我又不是很宅的.
 但我比較是 即使疫情期間沒得出去.
 走香港走 走到盡頭.
 我發覺長洲都多很多人.
 疫情期間開始多的.
 繼續持續都多的.
 每一條路 每一塊基石.
 那些叫什麼石頭.
 又不是路標石.
 機子 電機的介石.
 那些叫介石 沒錯.
 介石的人都擺上去.
 去過哪裡拍一張介石的照片.
 小紅書就更多這些東西.
 我就仍然是在疫情期間.
 走來走去都是長洲.
 所以很多地方都沒走過.
 今天有機會來到洪水橋這裡.

$^{41}$ 跟大家一起聚會.
 多謝曾牧師給我這樣的機會.
 能夠參與.
 跟大家一起敬拜我們的上帝.
 我的聲音有點沙啞.
 因為昨天去了一個婚禮.
 我做了訓練.
 跟著我們校友.
 我以前關心小組.
 整班同學晚上就喝酒.
 很開心 喝酒的地方.
 是很吵很吵.
 整個桌子都要很大聲很大聲說話.
 所以昨晚基本上叫了一整晚.
 訓練的時候就沒有叫.
 因為有咪高峰.
 在教堂裡面就沒事.
 晚上喝酒的時候就叫了一整晚.
 聲音沙啞.
 我們先再一次低頭祈禱.
 殿上我們將時間交託給你.
 今天是我們父親節.
 我們知道你是我們在天上的父.
 你也在地上讓我們有原生的家庭.
 有親生的父母.
 在我們成長的期間.
 我們得到我們的長輩.
 我們父母的照顧.
 我們今天來到你面前敬拜你.
 也是紀念我們父親在我們生命裡面.
 恩情的時候.
 他們的辛勞的時候.
 你與我們同在.
 叫我們讀聖經.
 我們思考.
 我們都能夠心被恩感.
 也救贖我們的上帝.
 真的有很多扶植過我們的人.
 尤其是辛勞在公書教學帶領我們長大.
 我們的父親 我們的母親.

$^{81}$ 求天父誰聽我們禱告.
 奉主明求.
 其實如果你說父親節母親節.
 一向節日都比父親節大.
 熱鬧過父親節.
 也比早過.
 應該沒有搞錯.
 我這件事應該沒有fact check過.
 但感覺上應該比父親節早.
 既然有母親節沒理由沒有父親節.
 今天跟大家不是思考單一一位男士.
 他在家庭中所成就的貢獻.
 不是.
 就是一班男人.
 一班男人在一起會做些什麼呢.
 碰到一個問題的時候.
 會如何處理問題以至得到解決呢.
 我們剛才所讀的.
 《士林傳》六章一至七節.
 其實一個典型.
 這個故事有好結局.
 你知道一班男人在一起做事情.
 未必一定有好結果.
 教母會不會打起來呢.
 不知道.
 但這個故事是由一個問題開始.
 但有一個好結局.
 所以我們可以拿來做一個範例.
 讓我們去看.
 要得到一個好結局.
 就沒有保證的.
 真的沒有保證的.
 你可以每件事情.
 每個步驟都做對了.
 但仍然有些不受控的因素.
 那你的結局都會變差.
 都不是如你預期.
 想像中想得到的.
 所以又不是是非必然的一個好結局.
 但無論結局好與壞.

$^{121}$ 其實我們都要知道.
 有沒有正確的行路方法.
 行走解決問題.
 一起共事.
 你不只是一個人工作.
 即使家庭裡都不是一個人工作.
 父母要一起合作工作.
 如果父母兩個人這樣的話.
 這個家庭都很難健康.
 甚至很難持續下去.
 一起工作的時候又是一回什麼事情呢?.
 我們先看看問題在哪裡.
 我們先看看經文的第一節.
 問題很簡單.
 有些人被忽略了.
 這個就是最早期的教會.
 現在是《史壇傳》第六章.
 所以五旬節過了之後沒多久.
 十二使徒就繼續在聖殿.
 他仍然在耶路撒冷.
 在聖殿附近宣講.
 沒多久之前就被釘起.
 釘起十字架上.
 卻是三日後復活而升天的.
 為以色列人所選擇.
 要拯救他們的基督彌賽亞.
 宣講這位復活的主.
 我們的盼望不再是虛幻的一件事情.
 是真真實實.
 但我們的主已經復活.
 死亡的勢力即使是現存.
 但事實中死亡的勢力已經過去.
 因為有更大的力量已經戰勝死亡.
 你和我的年齡都慢慢長大.
 都在看地上生活的終點.
 但你和我都要知道.
 因著主的關係.
 地上的終點卻不是我們生命的終點.
 宣揚基督.
 教會一向都有一個.

$^{161}$ 由最早期教會都有一個.
 慈慰的基因 教會的DNA.
 教會的DNA當然是敬拜.
 斷溪 傳福音 禱告服事上帝.
 亦都是彼此的服事.
 彼此的服事當中.
 又不只是教會裡面弟兄姊妹彼此服事.
 教會是外向型的.
 由第一天開始都是外向型的.
 所以他們會做慈慰工作.
 慈慰工作一方面當然是對人家好.
 令人家眼中有個印象.
 將來你和人分享一些東西的時候.
 人家都會比較容易聽到.
 這是人之常情.
 所以合理的 是合理的做法.
 但其實有更加深的一層.
 就是基督徒都認識到.
 即使我們是信主的.
 外面的人是不信主的.
 但他們都是上帝所創造.
 上帝的心意就是要讓所有人都歸向他.
 萬人得救 不願意 一個人沉淪.
 而我們對他人好.
 是建基於他們都是上帝所創造的前提上.
 不單是希望有個務實的動機.
 希望我們將來向他們說傳福音.
 他們會聽.
 而事實上他們和我們在上帝面前.
 都是上帝所愛的人.
 而且這個世界上永遠都有人有需要.
 所以就持畏基因就在這裡出來.
 所以早期教會人數三千五千人信主.
 當中有很多窮人 以前世界上有很多窮人.
 現在世界上當然也有 現代世界.
 但是剛才曾牧師很辛苦.
 主席學完之後就去車站載我過來.
 他當然要介紹 開車一路走介紹.
 低密度住宅 兩等是低密度住宅.
 中產或中上 中間夾住那個.

$^{201}$ 現在叫做簡約公屋.
 兩等中間夾住簡約公屋.
 所以我們中間又有需要人.
 以前世界更加是.
 這個不需要我跟大家解釋.
 所以早期教會做最簡單的事.
 就是一日三餐解決.
 一個人要溫 要飽.
 溫未必一定提供到那麼多地方.
 住一下 只住溫.
 但是飽很重要 要飽.
 如果吃不飽的話.
 連自己有份工作給你做.
 你都沒有力量 你會爭輸.
 被人搶了份工作做.
 越來越每況愈下.
 變成沒得救了.
 所以最重要是飽.
 你有份工作做.
 都可以做到那份工作能夠幫到自己.
 所以就做食物銀行.
 不是食物銀行.
 開條食物線.
 是飯食.
 不過呢.
 很壞不壞呢.
 就是被人投訴.
 就說忽略了.
 現在就叫歧視某一類的人.
 耶路撒冷 猶大地.
 以色列人說的以色列話.
 以色列話我們說的所謂的希伯來話.
 現在的以色列的摩登以色列話.
 是由以前的希伯來話變過來.
 但是跟以前很不同的了.
 其實.
 就等於我們現在的白話.
 廣東話都好.
 都跟以前是會有文言文.
 就算是說白話.

$^{241}$ 以前的白話都跟我們現在有些區別.
 但是猶太人.
 大家讀舊約就知道.
 因為以色列犯罪.
 上帝是讓外幫人擄了.
 滅了北面以色列.
 南面猶大國.
 結果猶太人就有一部.
 很長時間就四散.
 後來就有一部份就回到以色列.
 就造成在耶穌時代.
 猶大地有以色列人.
 但是也有很多呢.
 分散到外地避老.
 他們不是去移民.
 我知道新界有很多.
 去英國很早期.
 去英國開餐廳.
 移民去.
 不是移民而是避老.
 打仗的時候被人俘虜.
 被人帶到異地外地很慘的.
 但是經過十代.
 二十代之後.
 就算有得回來也好都不回來.
 所以周圍都有很多猶太人.
 但是他們仍然是猶太人.
 所以他們仍然知道敬拜的中心在耶路撒冷.
 所以他們有時會過節.
 會回去某些大節期.
 所以看到五旬節.
 就很多周圍外地的猶太人都會回去耶路撒冷過節.
 那些人回去耶路撒冷過節.
 那些人就不是說當地話.
 就是說.
 說雞腸.
 就是我們這些小朋友出去讀書.
 讀了十多年回來差點連廣東話都不懂.
 就是說英文.
 說希臘話.

$^{281}$ 有些人回來之後就滯留了.
 這裡是說那些寡婦.
 可能是住了一段時間.
 先生過世了.
 又或者是有其他的因素.
 令到他們那些說希臘話的猶太人.
 寡婦滯留在耶路撒冷.
 那你知道.
 寡婦在現在世界都很難.
 我的爸爸.
 我爸爸是我中學畢業那年.
 我中五那年過世的.
 家裡有五個小朋友.
 我最大的.
 最小的弟弟那時候在讀幼稚園.
 我媽媽是一手.
 搞定五個人的.
 寡婦.
 我讀中學畢業那年.
 媽媽就變成寡婦了.
 所以這些事情我是很深刻的感受.
 以前的寡婦更加慘.
 是你沒有經濟能力.
 完全沒有經濟能力.
 我媽媽因為住公屋.
 都還好一點.
 起碼那個灣那裡有一個地方.
 租金很便宜.
 所以在這些情況下.
 那些滯留寡婦.
 是需要幫助的.
 不過早期教會.
 沒有想到一件事.
 大家四處圍在一起.
 四一人個個講四一話.
 圍在一起 台山人個個講台山話.
 照顧台山人.
 我們在香港講廣東話.
 內地有講普通話的.
 他們又不懂聽你講廣東話.

$^{321}$ 就算有一個地方有飯吃.
 都不為意.
 不知道沒有人叫他.
 所以就被忽略.
 忽略下來.
 就有人為他們發聲.
 就投訴你們在每天的公認上.
 忽略了那些講外地話.
 其實都是猶太人.
 猶太人來的 猶太人寡婦.
 有極大需要 忽略了他們.
 所以.
 我分享的題目就是.
 當有投訴事.
 不過這個是.
 有投訴.
 香港多不多人投訴.
 很多人投訴.
 我要問問 建州.
 弟兄就知道.
 他以前在警隊服務就知道.
 我打電話要向警察投訴.
 哪裡哪裡出現這些事.
 某某人 怎樣怎樣.
 香港都很盛行一種投訴文化的.
 其實是.
 有投訴 你怎樣.
 如果你是接受投訴的.
 或者你是被投訴的一方.
 你是被投訴的一方.
 你那晚在家裡.
 開小早 突然很開心很嘈.
 隔離 嘎嘎嘎嘎 搞到十點鐘.
 十點鐘很早而已 在香港.
 你不要那麼嘈.
 被投訴的一方.
 那你會怎樣做.
 你說 嘈什麼 嘈 這麼早.
 還沒睡覺 哪有人這麼早睡覺.
 十點鐘我們會閉嘴的 你回去吧.

$^{361}$ 你這樣.
 這樣做就升級了.
 我們說將事情慢慢升級.
 被人投訴的時候.
 要解決這個問題.
 你知道.
 飯食.
 我們不知道聖經沒有說.
 是早期教會.
 每天要餵多少人.
 你當他餵.
 例如我們坐在這裡100人.
 每天三餐要餵100人.
 也不簡單.
 你要煮100人.
 沒有現代廚房.
 以前的燒柴.
 要餵100人.
 那些錢從哪裡來.
 之前有記載過.
 有些人變賣家業.
 販賣公用.
 有些錢就放在.
 使徒那12個的腳前.
 讓他們分配.
 要做很多功夫.
 錢是因為那些人相信他們.
 他們是領袖.
 將錢給過他們.
 所以要分配.
 馬首是在那個階段.
 他們自己都要做很多買菜.
 設置.
 每天都要.
 崇拜那天都要.
 他們就很忙.
 但無論如何.
 持有的就是那12個.
 所以那些人投訴.
 那12個使徒.

$^{401}$ 你們做事.
 為什麼要這樣做事.
 你們當我們不是以色列人嗎.
 怎樣怎樣.
 你想像你自己是那12個男人.
 男人來的.
 我煮飯給你吃.
 還被人嘈.
 好像你父母買菜回來.
 煮飯給子女吃.
 又吃這些東西.
 以色列人向上帝這樣說.
 每天都吃這些東西.
 通常媽媽.
 尤其是媽媽買菜又煮飯.
 那些家庭.
 不喜歡吃嗎.
 你自己去買吧.
 不要說煮 你自己去買吧.
 去到街市不知怎樣買菜.
 只會去超市買東西.
 你想像那12個男人.
 他們會怎樣做事.
 這很重要 一路下去.
 他們的處理手法.
 其實要看到.
 他們第一步做甚麼.
 經文去做下去.
 第二節也可以.
 回到第二節.
 12個仕途叫門徒來說.
 說甚麼呢.
 其實我們看聖經.
 是要仔細讀的.
 每一句都不需要看完整節才去想.
 一句都已經停在這.
 12仕途叫眾門徒來.
 叫弟兄姊妹聚集在一起.
 他們解決問題的第一個方法是甚麼.
 叫所有人來.

$^{441}$ 很有趣 叫所有人來做甚麼.
 叫做甚麼.
 叫做開會.
 大家喜不喜歡開會 舉手.
 有沒有人喜歡開會.
 沒有人喜歡開會.
 開會是否浪費時間呢.
 又要開會.
 執事會又要開會.
 團契導師又要開會.
 堂委又要開會.
 業委 業管又要開會.
 公司又要開會.
 開會又不關.
 如果你要開會.
 那你就壓力很大.
 不用開會 坐在這.
 按下手機.
 都不想開會 很浪費時間.
 但是聖經在這裡給我們例子.
 聖經在這裡沒有教導.
 沒有說大家.
 不是不可停止聚會.
 不可停止開會.
 聚會也是開會的中氣.
 實在有說.
 不可停止開會.
 聖經不是吩咐.
 這是故事.
 我們在看故事的時候.
 望聞新意.
 開會其實很重要.
 開會在這裡.
 重要在哪裡.
 如果別人向你投訴.
 那你就找所有人.
 所有人.
 眾門徒.
 眾門徒在這裡沒有分哪類人.
 眾門徒.

$^{481}$ 因為你飯吃的就算是做好事.
 也有些持份者.
 你幫哪些人受到照顧.
 現在有一部分.
 照理應該受到照顧.
 但你忽略了他.
 所以你叫所有持份者.
 來談對策.
 因為不是只有你一個人.
 有些事情是你自己能處理的.
 別人也有份的.
 如果是一件事情.
 要解決的.
 那件事情只關你自己的事.
 別人沒有份的.
 那你就不用開會.
 你問自己 你肚餓.
 你肚餓就只關你自己的事.
 你的肚子向你投訴.
 你就不需要.
 老婆.
 我和老婆開會.
 其實你只是吩咐她去工作.
 你就這樣.
 我肚餓了 我自己找東西吃.
 我自己想想.
 去吃什麼 去吃什麼.
 自己的腦和自己開會.
 跟著去找東西吃就搞定了.
 因為只關你自己的事.
 不關別人的事 你肚餓.
 但在這裡關別人的事.
 就一定要開會.
 所以開會最基本條件就是.
 你要先問這件事關誰的事.
 關別人的事.
 你就不要自己一個人獨斷獨行.
 一個人搞定了.
 其實十二個人自己都要先開一次會.
 聖經沒有說.

$^{521}$ 但第二節說十二仕徒叫眾門徒來.
 即是他們已經同意了.
 自己都談過 我們應該叫眾門徒一起來再談.
 我們十二個人都沒有辦法解決這件問題.
 就算我們有同一個看法.
 都是不好 叫所有人再諮詢一下 談一下.
 所以十二仕徒之前已經開過會.
 起碼是第二個會 不是第一個會.
 起碼是第二個會.
 就將持份者帶在一起.
 所以不可停止開會.
 就真的是一件很重要的事情.
 開會要有東西談 開會不要浪費時間.
 開會談什麼呢.
 這個就是第二節的繼續下去.
 一開會 他們就劈頭說.
 弟兄們 我們都要回到第二節.
 還沒讀完 麻煩你.
 我們撇下上帝的道去管理這件事 是不合宜的.
 麻煩你 可以去第三節了.
 哇 他們一開會.
 十二個 應該不是十二個一起說話.
 應該是代表 代表應該是彼得.
 不要說了 應該是彼得.
 找代表說.
 我們劈下上帝的道去管理飯食 本來就不應該的.
 哇 你聽聽 開會第一件事.
 主持會議 召集會議那班人.
 或者被投訴那班人 一開始他們說什麼.
 什麼叫做 我們劈下上帝的道去管理飯食 本來就不應該的.
 你是怎麼解釋這句話 他們在說什麼.
 當你投訴我吧.
 然後我一出來我就說.
 弟兄姊妹 其實我劈下的道去管理飯食 其實是不應該的.
 那是什麼呢.
 真是卸膊.
 現在有問題就是說 你本來就不應該做這些事.
 本來就不應該是你們做的.
 如果不是你們做 誰來做呢.
 十二時一開始就卸膊 卸膊 真是很嚴重的.

$^{561}$ 被投訴的人說 與我們無關 我們也不想做 我們不應該做的.
 不知道為什麼要我們做.
 這樣就很嚴重 如果你的老闆是這樣的老闆 你就慘了.
 很難做的嘛 不過我們要小心一點 小心一點看.
 剛才那句話 其實也比第三句說得快了 回到第二節.
 小心一點 逐句逐句看 我們聖經是要逐句逐句看的.
 我們撇開上帝的道去管理飯食是不應該的.
 的確說我們不應該做這件事 這件事與我們無關.
 他沒有給理由 為什麼不應該做這件事.
 有比 當然還有很多其他理由的 比如說我們不會煮飯.
 我們不會買菜 我們不會組織 都是理由.
 我做得不好 總之你怎麼找我做呢 找一些厲害的去做吧.
 我都不會彈琴 你找我彈琴做什麼呢.
 合理的 但他們這裡沒有說這件事.
 他們說 我撇開上帝的道去管理飯食 什麼意思呢.
 這個就是先後次序了 即是說十二使徒表示.
 其實他們有其他事情要做的 而因為管理飯食.
 就令到他們撇下了另外一件事做不到.
 因為三餐飯食是全時間的 大家知道.
 所以叫做家庭主婦嘛 是不是.
 一做家庭裡面的老頭 你就是做主婦了.
 不是主人的主啦 當然.
 很難的 主人的主 煮飯的主 主婦 是不是.
 那你的時間都去了 你買菜要量度 買什麼.
 接著去買 接著抬回來 接著要切 要割 要煮 要分配 還要洗.
 哇 接著第二餐就來了.
 都不用休息了 對了 都不用休息了.
 為什麼今天是父親節呢 應該是母親節才對.
 無論如何 所以我們又明白 他們一煮飯吃.
 他們就做不到宣講上帝話語的事 很明顯的.
 但是他們這樣說是什麼意思呢 其實他們自己也有抱怨.
 他們這樣說是表達了一個優先次序 他們自己也向會眾說.
 不一定是抱怨 只是將事實告訴你 其實我們這樣做 令我們放下了重要的東西.
 重要的東西就是宣講上帝話語 事實上第一章小人傳第一章的大使命.
 你們耶路撒冷 猶大 全地 三不一 直到地極 要去做我的見證.
 主耶穌向他們說要為我做見證.
 哪些人才能夠為基督做見證呢 所有信徒都可以.
 但不是在當時見證 當時見證著.
 所以十二個使徒當猶大死了 少了一個的時候.
 一定要補選一個出來 而補選的條件就是補選的人的資格.

$^{601}$ 必須要由開始一直跟其他十一個人 跟基督站在一起.
 而且見證耶穌基督復活和升天.
 你見證見證見證見證 你都沒見過 做什麼呢.
 你是真的經歷過 意思就是說那十二個使徒是很重要的.
 有些事情是他們做到 後來信主那些是沒辦法做的見證.
 一定要他們繼續做宣揚基督 是他們才可以這樣宣揚基督.
 我們各自可以宣揚 見證基督在我們身上的工作.
 但我們沒見過基督釘十字架 沒見過基督重焚無復活 沒見過基督升天.
 我們做不到那部分的 我們相信那部分.
 但那些見證人一直流傳下來的見證.
 第一世紀的見證就是他們目擊證人的見證.
 在法庭裡面目擊證人的見證.
 所以是他們才做到的 現在因為管飯吃 更加重要的事情就做不到.
 所以開會 我自己領受就是 開會第一件事可能要做的.
 如果發生的問題在那裡來就需要澄清的.
 第一件事情就是一定要界定幼稚.
 你遇到這個問題真的跟幼稚有關係的.
 你為了一件次要的事情 花了功夫在次要的事情上.
 結果第一件事 最重要的事情做不到.
 所以是A字博 是卸博.
 卸博不是什麼都不做 我們人老了 無能為也.
 你們做了不想做 翹起雙手 你們做什麼.
 我做得這麼辛苦 你們還投訴我.
 他們不是這麼說 而是他們說其實我們有一件事更加重要的.
 本來是我們去做的 其他人都做不到.
 現在因為管飯吃的緣故 放下了就做不到.
 那就變成一個合理的理由.
 所以這個交代不是A字博 是解釋 讓所有持份者明白到底發生什麼事.
 到底12時內他們心裡在想什麼 溝通中.
 他們的心裡在想什麼.
 所以開會的時候一定要說 心裡在想什麼 到底什麼重要 什麼沒那麼重要.
 否則就花很多時間在不重要的事情上.
 還有要把心裡在想什麼說出來.
 你不要依賴別人猜你的心裡在想什麼.
 否則你肚子裡的蟲 我怎麼知道你的心裡在想什麼.
 你說就說這件事 原來你想就想另一件事 沒那麼困難.
 所以四人團就把這件事說得清楚.
 其實很重要的 男人溝通不是最擅長的 女性溝通好一點的.
 這個很多科學證明 不知道怎樣科學證明 證明到的.
 好像是這麼說的.

$^{641}$ 溝通 這裡十二個男人都懂得溝通.
 不過有一點點愚蠢的 我們就這樣看.
 我們放下上帝的道德觀反思原本是不應該的 有一點點愚蠢.
 但是就把他心裡的說話說出來.
 但重點就是說分優次 然後表達溝通好這件事情.
 好了 終於來到第三節了.
 這樣要解決問題嘛.
 那你家裡有些更重要的事情 你回去做重要的事情了.
 這些難事怎麼辦呢.
 要找人來做 真是要找人來做了.
 所以弟兄們 由我們中間 要說選出.
 選出有七名有名的明星 還有姓名和智慧.
 我們就派他們來管理這件事情.
 又是隻隻 字字珠璣 隻隻字都重要的.
 第一 就是在你們當中選出.
 大家都同意找代表.
 為什麼七個我就不知道了.
 總之七個吧 不會多過十二個.
 七個 但是就要set條件了.
 你要找人來做 因為他們做不到.
 那就要找人來做 找人來做.
 就不可以媽媽 狗都跑來做某些事情.
 你說管飯吃而已嘛 那就是買菜 煮菜.
 分飯 洗碗 誰做不到啊.
 不好意思喔 早期教會連這些事情都很認真.
 你想做就讓你做.
 而是要有好明星 哈哈 真是名氣好.
 滿有姓名 所謂好明星.
 明星是時間累積的 人家對你的觀感是怎樣.
 就是說 你有人知道你做事做得來.
 你有好明星 信得過 所謂好明星就是信得過你.
 就是你每一次做一件事情都炒火了.
 每一次都賣不完 掃走大便 讓人來撿爛攤子的話.
 那你不可能有好明星 做得到事情的人才是好明星.
 滿有姓名就是基督徒 領受過姓名的是他們當中自己人.
 還有有智慧 為什麼要有智慧.
 十二仕徒就是正正面對別人的投訴.
 一面對別人投訴 你就需要找個智慧人出來.
 因為你不能罵回頭 投投投投 投什麼數啊.
 這樣就沒有智慧了.

$^{681}$ 要找個有智慧 懂得解決問題.
 因為十二仕徒真的很有智慧.
 他知道這只是第一個問題 以後還有其他問題出現.
 所以一定要找一個有智慧的人.
 煮飯 分飯 都要有智慧才行.
 如果不是又被人投訴 或者食物中毒怎麼辦.
 以前食物中毒沒辦法投訴.
 你沒有洗手.
 要有智慧 所以這條件看起來很籠統.
 其實是很重要的條件.
 你要找人做 你一定要設定條件.
 我們說門檻 一定要有門檻.
 如果不是 你找人做都未必做得成那件事.
 你想做好那件事 有個門檻 訂合適的條件.
 然後就去找.
 然後找怎樣 你看看.
 我們派他們管理這事.
 來到這裡你會發覺.
 12個仕途是不是A字幕.
 不是的.
 他說我們派他們管理這件事.
 他們背上身的.
 那12個是這班所有的持份者當中選出來.
 有門檻 有條件.
 不是媽媽舊都找出來.
 但是到了給權柄他們去做.
 是誰給的 是12個仕途給的.
 他們背上身的.
 因為這7個人你們選出來的.
 但這7個人出事了都是我們負責的.
 是我們派他們去的.
 其實沒有寫上去的.
 只不過我們不做那件實質的事.
 但我們仍然.
 英文叫做The Bug Stop See.
 有投訴都會回到12仕途那裡.
 仍然要背上身的.
 即是說這個方法解決得不好.
 我們再想第二個方法.
 他們都要背上身的.

$^{721}$ 其實有一個很重要的原因.
 因為錢的人.
 就算有7個執事選出來管飯吃.
 錢的人是不是放在7個人的腳前的.
 也不會的.
 那些錢仍然去到哪個位置.
 仍然回到12仕途的位置.
 所以他們不能甩掉這個責任.
 所以這7個人都是他們派出去.
 所以沒有甩掉.
 A字都A不完的.
 派他們去管理這件事情.
 要說回他們的.
 我們現在會問.
 叫人做你們做什麼.
 耗時懶飛.
 至於我們呢.
 我們是專注在祈禱和傳道的職事.
 本業.
 做回本業.
 教會的本業.
 我們要知道.
 傳道人的本業.
 老師的本業.
 父母的本業.
 我們以前聽新聞都聽過.
 這是六七年前.
 有對夫婦在過大海回來的時候.
 在信德中.
 還是在澳門.
 在信德中.
 上岸之前.
 就被警方拘捕.
 接著就說.
 獨留兒童在家中.
 自己年輕夫婦.
 就過大海.
 很開心.
 將兩三歲小孩留在家裡.
 晚上想睡覺.

$^{761}$ 過大海他們很威風.
 第二天早上.
 因為小孩醒了.
 便坐下來.
 隔壁鄰舍發現不對勁.
 沒大人在 敲門沒人回應.
 於是報警.
 結果兩夫婦.
 回來的時候就被捉拿.
 就算被捉拿.
 也對小孩不好.
 因為.
 誰照顧呢.
 交給社署也不行.
 所以這些事情很難.
 我也沒有跟進.
 但有些父母.
 本業 他們忘記了本業.
 父母要犧牲很大.
 想去威風.
 也威風不了.
 為了小孩 也要犧牲.
 遲一點吧.
 等他們20歲 遲20年後才去威風.
 之類.
 本業.
 所以他們.
 回歸本業.
 所以開會重要嗎.
 到現在.
 開會是澄清很多事情.
 最怕開會.
 整天在聊天.
 不澄清一些事情.
 結果開完會也等同沒有開會.
 那就不好了 開會要知道要做什麼.
 開完會之後.
 回歸正業.
 祈禱傳道為事.
 我們可以看.

$^{801}$ 下面那段 經第五第六節.
 大家都很開心.
 他們也明白.
 這些是合理的.
 然後就揀選出來.
 一揀選出來.
 七個人.
 我們平時查經.
 人名不多.
 這些不是張雲開 曾銳季.
 歐戰超 這些名字.
 這些名字是什麼 都不認識.
 外國人名.
 但是原來聖經裡面.
 連名字都很重要.
 這些名字讀起來很繞口.
 肥里 伯羅 哥羅 尼加諾 提門 巴米那.
 還有一些詳細的解釋.
 關於猶太教的安提亞人.
 尼哥拉.
 還有斯提反.
 這裡還沒完.
 斯提反在前面.
 斯提反是排第一.
 我真的忘記了.
 這些名字.
 學者 一揀選.
 學者嘛 學者.
 看學者 留意學者 我自己一開始也沒有留意.
 然後一揀選的時候.
 聖經書裡面就有些學者.
 不是每個學者都提.
 有些學者就提 這些人是什麼人.
 這些全部都是希臘人.
 因為這些全部都是希臘名字.
 Sorry 不是希臘人.
 全部都是希臘名字.
 是猶太人.
 因為他們最早期的信徒.
 全部都是猶太人.

$^{841}$ 不過全部都是希臘名字.
 沒有一個是猶太人名字.
 猶太人名字 約瑟.
 約拿丹 約翰.
 約什麼的那些.
 都是猶太人名字.
 沒有一個是猶太人名字.
 全部都是希臘名字.
 什麼意思.
 你記住問題是哪裡來的.
 問題是講希臘話的猶太人寡婦.
 在每天供應的膳食上被忽略.
 沒有人幫他們.
 所以講希臘話的猶太人.
 投訴講本地話的猶太人.
 信徒就說.
 你選班人出來解決這件事.
 擺平這件事.
 應該有個比例.
 猶太人當中有講希臘話的.
 有講本地話的.
 中華人當中有講普通話的.
 有講廣東話的 廣東話是本地人.
 一件事情是大家都有份.
 你選個委員會出來管這件事.
 普通話的佔一個比例.
 廣東話的佔一個比例.
 4 3 行不行.
 我們4 你們3.
 那就不好了 你們4 我們3.
 都有份 用什麼逆向歧視.
 倒過來我們本地寡婦被忽略怎麼辦.
 所以大家一定有比例.
 誰知道選出來後很奇怪.
 7個名字都是外邦人名字.
 即是希臘名字 猶太人來的.
 還有一個簡直是外邦人.
 並歸於猶太教的安提亞人尼戈拉.
 尼戈拉是外邦人.
 歸化了猶太教.

$^{881}$ 變成了恆國禮守律法.
 變成猶太人 但其實血統是外邦人.
 這個就是第二點.
 這個就不關開會的事.
 這個就是你明白早期教會的亮度.
 他們說 我們這件事.
 他們承認這的確是一個問題.
 要解決這個問題.
 選吧 一選.
 通通給你們投訴.
 選出來的7個 全部你們當中選出來.
 一般不會這樣做.
 一般就是大家一半一半.
 但他們說既然是這樣.
 他們真的不擔心逆向歧視.
 這個就是早期教會的彼此信任.
 不會說你們比例都不要那麼高.
 連比例4352還是16都爭拗.
 零7 我們零得了你們7.
 這個不是大方.
 大方是什麼意思呢.
 大方其實是沒有意思的.
 其實是信任.
 我信任 全都給你們做.
 你們不會翻臉.
 不會倒過來.
 我記得你以前怎麼做.
 以其人之道 還是其人.
 不會的 他們信得過這班說外地話.
 外地過來的猶太人.
 是不會這樣做的.
 然後讓他們最明白.
 說希臘尼亞話 寡婦需要.
 可能連吃的東西也有點不一樣.
 讓他們搞.
 所以這件事就變成了一種有容乃大的精神.
 是他們有這樣的胸襟.
 而且選出來的條件.
 是我們剛才看過.
 是能力和人格並重的.

$^{921}$ 他們不只是說你做得好就行了.
 人格很重要 有好明星.
 但也不是說有好明星就行了.
 也要有能力.
 大家知道 我很尊重.
 是曾牧師.
 曾牧師明星 性靈充滿.
 怎樣的.
 但是 找他做什麼.
 找他做控制怪獸.
 要地基我們要動土禮.
 要挖地基我們要建堂.
 找他搞這件事.
 因為這不是他的專長.
 他有好明星.
 被人尊重 性靈充滿.
 過去做事都辦得好.
 但他沒做過的事不代表他會做得好.
 所以在這些位置.
 人格和能力都並重.
 他們需要處理錢的問題.
 所以好明星也很重要.
 因為錢是給仕途的.
 但買菜就是那七個去買.
 仕途要分錢給那七個人去買.
 大家哪些節目家裡有外傭?.
 外傭當然有.
 沒人家裡有外傭.
 姐姐啊.
 姐姐我們不叫外傭.
 我們叫姐姐.
 你要付錢去買菜.
 我也聽過一些故事.
 我自己家裡沒試過有外傭.
 但有時上車都知道.
 很多時候你會發覺.
 早期現在不是了.
 因為那些僱主不讓八達通去.
 所以每次坐巴士都是付錢.
 數數錢.

$^{961}$ 家裡沒有八達通機.
 查看你到底有多少錢剩.
 所以他要付錢.
 其實他們要管錢.
 拿錢去買菜.
 剩下多少錢買什麼菜.
 會不會有人打虎頭呢.
 所以好明星是重要的.
 當中的人是重要的.
 他們自己信任的那個是重要的.
 所有的事情.
 不單只是人格.
 而且他們都要處理好.
 能夠有能力處理辦好那件事.
 所以從路加的敘事來看.
 能力和人格並重.
 而他們一直都是建立信任.
 而不是製造差異.
 他們的七個選出來.
 就去做了.
 我們看看最後的兩節.
 結果如何.
 他們過來.
 他們站在使徒面前.
 使徒背起飛機.
 他們的七個是那十二個所猜險的.
 所以最終使徒背起飛機.
 結果就是第七節了.
 上帝的道興旺起來.
 耶路撒冷門徒的數目增加很多.
 很多祭司聽從了.
 甚至 應該加上「甚至」在翻譯.
 也有很多祭司聽從了這個信仰.
 這個其實.
 覺得使徒做得很好.
 那七位也做得很好.
 但是你想清楚一點.
 這個結論是很古怪的.
 完了之後這個故事就完了.
 第七節就完了 以後都不再提這件事了.

$^{1001}$ 但是第七節的結論是極其古怪的.
 因為投訴是什麼呢.
 投訴是說.
 你們忽略了我們說希臘話的猶太人寡婦.
 那麼就嘗試解決問題了.
 問題解決了.
 所以本來的故事結論應該是怎樣呢.
 從此 說希臘話的猶太人寡婦的日常膳食.
 有需要的寡婦 寡婦一定有需要的.
 膳食需要都解決了.
 如果你跟我寫故事 一定會這樣寫的.
 這是叫做首尾.
 複認嘛 你守了 美不認 怎麼行呢.
 所以你一定要表明問題解決了.
 大家都覺得故事很完整.
 看完了 很開心.
 她沒有提 泛釋問題解決了.
 反過來提哪件事呢.
 上帝的道興旺起來.
 由這一點 現在是路加在說故事.
 所以你明白路加對整件事情的看法.
 路加為什麼要提泛釋問題這件投訴事件呢.
 投訴值得提的 教會怎麼解決投訴問題.
 都是值得提的.
 但在路加生活中 教會的主業仍然是傳上帝的道.
 最終都是傳上帝的道.
 所以原來泛釋事情做得不好.
 帶出來的一件事.
 其實是十二仕徒沒有做好他們的主業.
 他們回去做他們的主業之後.
 接著上帝的道興.
 之前一段時間很多人信主.
 接著開始做泛釋 信主人數少.
 不是因為做慈惠的事情 令信主人數少.
 而是那十二仕徒去做慈惠的事情 結果人數少了.
 那十二仕徒回去做宣講 做見證的事情.
 人數又多了 慈惠的事情繼續做.
 但從來沒有說不做.
 那十二仕徒說不如不要做了 這樣就糟糕了.
 他們不是這樣的A字幕.

$^{1041}$ 而不是我們不應該做這件事.
 我們應該做那件事.
 所以找其他人做這件事 只是這個意思.
 所以當他們做回他們正職.
 宣講為基督做見證的時候 人數就增加了.
 所以這個就是路加的眼光.
 他明白早期教會存在的目的.
 最終都是為基督做見證的.
 但為基督做見證的同時.
 我們是有DNA 也有驅使.
 你要對人好嘛.
 你不可以為基督做見證.
 轉頭過來就見死不救 這樣也不行.
 這個是同一件事來的.
 主耶穌也說過很撒馬利人的比喻.
 要實踐的誡命 你就照著去做.
 你做別人的鄰舍.
 所以這是一件事來的.
 但12個誤導有他自己更加重要 首要工作.
 而這件事做不到的話.
 就算其他事都是對的事 應該做的事.
 但他們做不到這件事的話.
 很多事都是白費.
 教會就變成一個慈善機構.
 洪恩棠是不是一個慈善機構?.
 是不是?.
 喂喂喂 先轉頭過來.
 洪恩棠是不是一個慈善機構?.
 你們又錯了.
 建政府註冊呢?.
 是的 哈哈哈哈.
 88 啊 88.
 建政府註冊法例上.
 沒有所謂法定宗教機構這件事.
 因為香港沒有宗教法.
 所以沒有法定宗教機構.
 所有宗教機構 尤其是有錢過手的.
 都會列入做慈善機構.
 但在我們心目中我們是一個什麼樣的慈善機構?.
 當然不是 我們是傳福音 敬拜上帝.

$^{1081}$ 互相團契 建立上帝國度.
 信了主之後都要互相建立的群體.
 是天國子民.
 宣揚那位.
 照顧我們出黑暗 入奇妙 光明者的美德 宣揚上帝的機構.
 所以一旦本末 本和次 擺好了.
 所有其他東西都會一起來.
 你的本擺得不對 連次要的東西都做得不好.
 所以路加就讓我們看到.
 是早期教會的智慧.
 我想那些男人一起開會有什麼好結果.
 吃喝喝 談了很多.
 就算不是言不及義也好.
 浪費時間 結果也做不到什麼.
 但不是 而是早期教會的智慧上帝.
 所以聖靈充滿有智慧的意思.
 是讓他們懂得解決問題.
 讓他們懂得知道什麼是重要 什麼是次要.
 展示了早期教會彼此的信任.
 我放心你們去做 全都給你們去做.
 因為我們怕我們再做.
 四三也好 五二也好.
 我們也有些想法未必是不盡完美.
 這次就完全選出你們當中的人出來做代表去做.
 這裡沒有說過界別.
 會不會一年一屆還是幾年一屆 不知道.
 但如果再出事的時候 再開會再談.
 開會就沒有一次開會這回事.
 就有第二次第三次 繼續.
 所以給大家的說法就是.
 盧家最終解決問題.
 他認為最終解決問題就是大眾的參與.
 持份者一定要有份.
 公開的說明.
 你不要只是你自己跟自己說.
 放在心裡不說.
 第三 明白分數.
 投訴那邊要說清楚.
 你聽那邊也要解釋清楚自己的立場.
 透明的安排 合理的賦權.

$^{1121}$ 選出來有門檻有條件的.
 不要亂找一個人.
 責任的承擔 他們願意承擔責任.
 頭頭也要承擔責任.
 這樣就是將榮耀歸給上帝.
 依賴上帝給我們的所有男人.
 所有女人 弟兄姊妹.
 各自有智慧各自有恩賜.
 來做好我們教會應該做的事情.
 所以今天在父親節.
 跟大家看這段經文.
 作為一種勉勵 我們一起祈禱.
 天父我們感恩.
 你是我們的上帝.
 我們每天都像櫻桃小丸子說的.
 每天問題都多 每天都是問題.
 但你也給我們有頭腦.
 頭腦懂得解決問題.
 有些是我們自己一個人.
 只能夠一個人面對的問題.
 有些是我們一家人面對的問題.
 有些是我們兩夫婦面對的問題.
 有些是整間教會面對的問題.
 求天父幫助我們.
 無論我們面對什麼問題.
 我們都能夠藉著你給我們的恩賜和智慧.
 跨越問題 將問題轉變成一個機會.
 將問題轉變成一個彼此建立.
 讓你自己的恩典能夠彰顯的一個場合.
 你施恩給我們.
 叫我們這班蒙你恩典的人.
 懂得怎樣接受恩典.
 亦懂得怎樣互相分享恩典.
 誰聽我們禱告 奉主名求.
 阿們.
 奉獻是基督徒的本份.
 若在座當中不明白奉獻的意義.
 可以自由參與.
 我們現在為奉獻祈禱.
 又賜各人有不同的恩賜.

$^{1161}$ 讓我們有能力賺取日用所需.
 更教導我們不要單顧自己的需要.
 更要顧念別人屬靈的需要.
 願祝我們獻上當立的金錢.
 求主帶領教會的教目.
 領導洪恩嘉成為洪水橋明亮的燈台.
 照亮每個黑暗的角落.
 榮耀歸主名.
 以上合一的禱告.
 乃奉教主耶穌基督得勝的明志.
 阿們.
 現在開始向神奉獻.
(主席奉獻).
(主席奉獻) 都怪耶穌 都怪耶穌.
(主席奉獻) 我願全生都怪主.
(主席奉獻) 願我言語行為應驗.
(主席奉獻) 眾生歲月都怪主.
(主席奉獻) 都怪耶穌 都怪耶穌.
(主席奉獻) 都怪耶穌 我求主.
(主席奉獻) 都怪耶穌 都怪耶穌.
(主席奉獻) 都怪耶穌 我求主.
(主席奉獻) 願我雙手為主作歌.
(主席奉獻) 願我雙足跟主走.
(主席奉獻) 願我上目只見耶穌.
(主席奉獻) 願我口是眾主名.
(主席奉獻) 都怪耶穌 都怪耶穌.
(主席奉獻) 都怪耶穌 我求主.
(主席奉獻) 都怪耶穌 都怪耶穌.
(主席奉獻) 都怪耶穌 我求主.
願我們都能領受神的真光.
帶領我們知道祂的旨意.
盡心盡性盡義盡力愛我們的主上帝.
主愛光輝正照耀人間.
照念康榮的風象明燈.
主祈禱帶真光照亮萬國.
賜與憎美事釋放你罪所.
照耀我靈照亮我心.
真光普照充滿天地.
來勝負永遠可驗勝利.
此次火獸主愛擁抱.

$^{1201}$使人顛覆不夠送眾心.
請至真理將主光輝照耀.
皇帝正前仰關我廣廢.
讓我彰顯主聖凡應聚.
從未了老應要不斷感濕.
在我心上還應有主神的照耀.
照耀我靈照亮我心.
真光普照充滿天地.
來勝負永遠可驗勝利.
此次火獸主愛擁抱.
使人顛覆不夠送眾心.
請至真理將主光輝照耀.
會眾請坐.
主日送拜結束.
會眾安靜脈套後可以自由交通.
願上帝賜福給各位.
《天地之祈禱》 詞:陳家麗 曲:陳家麗.
天地之祈禱 詞:陳家麗 曲:陳家麗.
\newpage



\section{創世記 18:16-33}
\label{sec:RaRAzKcNgEg}
\textbf{2024.06.23 《與神同行》創世記 18\_16-33 (和修版) 許淑芬牧師}
\newline
\newline
連結: \href{https://youtube.com/watch?v=RaRAzKcNgEg}{\texttt{https://youtube.com/watch?v=RaRAzKcNgEg}} ~~~~ 語音日期: 2024-06-24
\newline
\newline
\hyperref[sec:WqOKT5MNC_k]{\small{< < < PREV SERMON < < <}}
~
\hyperref[sec:index_chronic]{\small{[返順時目]}}
~
\hyperref[sec:index_scriptual]{\small{[返順卷目]}}
~
\hyperref[sec:6OJe2UUeX7c]{\small{> > > NEXT SERMON > > >}}
\newline
\newline
創世記 18:16-33
\newline
\begin{longtable}{cl}
\hline
\hline
章節 & 經文 (和合本修訂版)\\
\hline
18:16 & \begin{tabularx}{0.7\textwidth}{X} 三人從那裡起程，面向所多瑪觀望，亞伯拉罕與他們同行，要送他們一程。 \end{tabularx} \\ \\ \relax
18:17 & \begin{tabularx}{0.7\textwidth}{X} 耶和華說：「我所要做的事豈可瞞著亞伯拉罕呢？ \end{tabularx} \\ \\ \relax
18:18 & \begin{tabularx}{0.7\textwidth}{X} 亞伯拉罕必要成為強大的國；地上的萬國都必因他得福。 \end{tabularx} \\ \\ \relax
18:19 & \begin{tabularx}{0.7\textwidth}{X} 我揀選他，為要叫他命令他的子孫和後代家屬遵行耶和華的道，秉公行義，使耶和華所應許亞伯拉罕的話都實現了。」 \end{tabularx} \\ \\ \relax
18:20 & \begin{tabularx}{0.7\textwidth}{X} 耶和華說：「所多瑪和蛾摩拉罪惡極其嚴重，控告他們的聲音很大。 \end{tabularx} \\ \\ \relax
18:21 & \begin{tabularx}{0.7\textwidth}{X} 我要下去察看他們所做的，是否真的像那達到我這裡的聲音一樣；如果不是，我也要知道。」 \end{tabularx} \\ \\ \relax
18:22 & \begin{tabularx}{0.7\textwidth}{X} 二人轉身離開那裡，往所多瑪去；但亞伯拉罕仍然站在耶和華面前。 \end{tabularx} \\ \\ \relax
18:23 & \begin{tabularx}{0.7\textwidth}{X} 亞伯拉罕近前來，說：「你真的要把義人和惡人一同剿滅嗎？ \end{tabularx} \\ \\ \relax
18:24 & \begin{tabularx}{0.7\textwidth}{X} 假若那城裡有五十個義人，你真的還要剿滅，不因城裡這五十個義人饒了那地方嗎？ \end{tabularx} \\ \\ \relax
18:25 & \begin{tabularx}{0.7\textwidth}{X} 你絕不會做這樣的事，把義人與惡人一同殺了，使義人與惡人一樣。你絕不會這樣！審判全地的主豈不做公平的事嗎？」 \end{tabularx} \\ \\ \relax
18:26 & \begin{tabularx}{0.7\textwidth}{X} 耶和華說：「我若在所多瑪城裡找到五十個義人，我就為他們的緣故饒恕那整個地方。」 \end{tabularx} \\ \\ \relax
18:27 & \begin{tabularx}{0.7\textwidth}{X} 亞伯拉罕回答說：「看哪，我雖只是塵土灰燼，還敢向主說話。 \end{tabularx} \\ \\ \relax
18:28 & \begin{tabularx}{0.7\textwidth}{X} 假若這五十個義人少了五個，你就因為少了五個而毀滅全城嗎？」他說：「我在那裡若找到四十五個，就不毀滅。」 \end{tabularx} \\ \\ \relax
18:29 & \begin{tabularx}{0.7\textwidth}{X} 亞伯拉罕又對他說：「假若在那裡找到四十個呢？」他說：「為這四十個的緣故，我也不做。」 \end{tabularx} \\ \\ \relax
18:30 & \begin{tabularx}{0.7\textwidth}{X} 亞伯拉罕說：「求主不要生氣，容我說，假若在那裡找到三十個呢？」他說：「我在那裡若找到三十個，我也不做。」 \end{tabularx} \\ \\ \relax
18:31 & \begin{tabularx}{0.7\textwidth}{X} 亞伯拉罕說：「看哪，我還敢向主說，假若在那裡找到二十個呢？」他說：「為這二十個的緣故，我也不毀滅。」 \end{tabularx} \\ \\ \relax
18:32 & \begin{tabularx}{0.7\textwidth}{X} 亞伯拉罕說：「求主不要生氣，我再說一次，假若在那裡找到十個呢？」他說：「為這十個的緣故，我也不毀滅。」 \end{tabularx} \\ \\ \relax
18:33 & \begin{tabularx}{0.7\textwidth}{X} 耶和華與亞伯拉罕說完了話就走了；亞伯拉罕也回到自己的地方去了。 \end{tabularx} \\ \\
[1ex]
\hline
\hline
\end{longtable}
$^{1}$相信大家都有去過旅行,旅行最重要不是去哪個地方,無論本地遊或外地遊,最重要是看看跟誰去,誰去才重要,誰去誰有話題..
正如莎士比亞所說,如果在人生兩半同行,路遙不覺得遠,這個我相信大家都有感受..
在人生旅途上,大家回顧一下,有沒有一些常常或不同階段都會有的同行者呢?我們相信朋友未必跟你一生一世,身邊只有一個朋友,跟你一生一世走,可能這是你的配偶也說不定..
但如果在人生中,總會有不同時段有不同的朋友,我記得我的弟兄姊妹,因為我弟兄姊妹每隔一段時間都會去釣一釣鰻頭,去釣一釣鰻頭的時候,大家一起吃飯,去farewell等等..
其中有些朋友都會分享,總是在這些時候很不捨得,覺得又過了,又要認識另一班朋友或戰友..
但在我們人生路途上,我們知道生離死別對我們來說是人間最難的地方,但我們很感恩,上帝給我們有最好的朋友,耶穌就是我們最好的朋友,讓我們能夠體會這位朋友永不分離,與我們同在..
在聖經裡,舊約裡說過很多與神同行的人,像以洛,除了阿伯以洛是最短命的,不過因為與神同行,上帝就接走他..
另外就是羅亞,羅亞也被稱為與神同行,神吩咐他去方舟,他就去做了,救贖了他一家,我相信羅亞本身也不只是想救贖一家,不過結果也是救贖了一家..
今天我想和大家分享亞伯拉罕,亞伯拉罕也被上帝稱為與神同行,也是以色列人一個很重要的人物,是以色列人的先祖..
還記得神親自與他同行的一段經歷,多謝主席剛才幫我們一口氣讀完整段與神同行的過程..
亞伯拉罕也因為與神同行的經歷,他的生命出現了一個很大的改變..
所以剛才我讀了一大段經文,我想用三方面和大家一起思想,亞伯拉罕如何與神同行,他的生命因著,為什麼與神同行,他的生命有著改變..
第一件事我們要知道的是,亞伯拉罕能與神同行,他最大的得著,就是因為他認識神自己,認識神的名字是很重要的..
我們知道摩西最初帶領以色列人出埃及的時候,亞伯拉罕之後的一個人物,他帶領以色列人出埃及的時候,.
他對上帝說:上帝,我不知道你是誰,這個名字我不認識,我怎麼介紹給你聽呢?.
於是上帝就說了一個字,不過在舊約,我們知道的創世記就告訴他,這個字叫做自由永有..
但其實這個字本身就叫做I am,其實他回答摩西的時候,不知道摩西怎麼聽的呢?.
回答摩西之後,他就說我就是我,他就是這樣跟摩西聊天,所以摩西就介紹給以色列人聽,誰介紹你來帶領以色列人出埃及,就是我就是我,就是那個人物..
這個名字我們叫做耶和華,我們不詳細說這個字,耶和華本身就是我,就是我的意思..
在聖經,在新約,約翰福音有一個很有趣,或者給我們一個很大的安慰,或者呼應的一段經文,.
就是使徒約翰用了七個的「我是」來表達耶穌的自己,意思就是他在呼應舊約的這位耶和華,其實就是耶穌基督「我是」的自己..
所以約翰用了七個「我是」,我是好牧人,我是世界的光,我是道路真理生命,這些全部都是舊約或者以色列人很熟知的名字..
在舊約常常都會看到很多不同的名字,很多不同的名字,當你打仗的時候,遇見敵人的時候,哇,神就是我的避難所,.
或者當你要去打仗找個頭出來的時候,哇,神就是我的新旗,其實在舊約去說耶和華的名字是很普遍的,耶和華就是我的什麼..
我記得我小女兒小時候有一次去旅行,我和展秋的子女差不多大,一起去戶外露營,小時候帶她去行山樂,.
我們團契每個月都一次行山樂,每天一次行山樂,到她大了一點我就不去,行山樂對我來說很辛苦,這麼熱,.
但是小孩子我就想帶她多一點去行山樂,有一次我就和她說,明天就去行山樂,看你們去不去..
我女兒當時,我兒子小一點,我女兒就說很想去,我就告訴她,我明天很忙,我沒空又叫你又叫你起來,.
我說你們自己祈禱吧,求耶穌給你,叫醒你,她真的走去窗口那裡祈禱,在祈禱之後,第二天早上她很開心告訴我,.
天父叫她起來,我說為什麼?早上六點,正正六點鐘,弟弟就喊叫,喊醒了她,她就很開心和我一起去行山樂,.
其實我女兒小時候就經歷,耶和華去哪裡?就是流忠,我記得我初信主的時候,祈禱很厲害,祈禱好像法術一樣,.
沒有車,耶和華就是巴士,最快的那架,最貴的那架,馬上就走到,或者沒錢的時候,耶和華就是我的銀行,耶和華就是我的朋友,.
當我們去經歷我們這位神,耶穌說將祂的名字給我們的時候,我們教徒很厲害,我們大家都有法術,不是法術,我們有耶穌術,.
我們可以藉著耶穌的名字去經歷我們這位神是又真又活的,所以認識神是要我們更深的體會神,我不是叫大家亂說神,.
我只是叫大家去做一些壞事,警隊和黑幫都同拜一位偶像,這位神如果真的是神的話,就有點怪怪的,黑幫又拜,正義的警察又拜,那怎麼辦呢?.
好像不太能接受,但這位上帝是正直信神,是公義的神,無論你是怎樣都好,我給大家看一幅圖,大概十幾年前,也有二十年前,是2010年,十幾年前,.
我們紅黑黃報道會,在過去這麼多年,辦過兩屆,二十年一屆,是一個很大的盛事,如果我們看到紅黑黃這三件衣服,黑色代表戶外,可能是一群戒毒,一群黑幫的人士,.
或者紅色代表藝人,很紅的,大家看到他,黃色代表黃色事業,是黃家警察,以前香港還沒回歸的時候是黃家警察,無論你的背景是什麼都好,信愛這位神是最大的不同,.
每個人都要歸向神,無論你是怎樣的人,背景是怎樣都好,我們大家都要跟著聖經的指引,這才是我們體會的上帝,也是我們所共信的同一位神,所以基督教有一個很獨特的特徵,我們同信這位神,我們去到世界各地也好,我們所同信的神,我們大家都有些樣子,都像樣子,.
因為耶穌給我們兩個命令,一個是彼此相愛,一個是傳呼音,也是愛人如己,愛人靈魂,我們去跟別人傳呼音,當我們有這兩個特徵的時候,我們就說看看,都像基督徒,這是我們世界共有的一個特徵,所以亞伯拉罕不單止與神同行,去認識上帝,去認識上帝的名字之外,.

$^{41}$第二方面我們要思想的是,亞伯拉罕要明白上帝的計劃是什麼,很多時候我相信我們每一個基督徒都很渴望去明白,究竟神在我身上有什麼計劃呢?但是我要問多一句,當我們說要明白上帝在我們心意的時候,究竟是你的心意,還是你真的很渴望是上帝的心意呢?.
這一點我們基督徒一定要懂得去反問,懂得去辨別,當我們反問的時候,很多人可能都會啞口無言,我不知道是我喜歡還是上帝喜歡呢?我不知道怎樣,但是我們看這段經文其實很特別,上帝是自言自語,來講出亞伯拉罕,說上帝在亞伯拉罕身上的計劃,.
他自言自語的說,我所做的事豈可瞞著亞伯拉罕呢?如果你是亞伯拉罕,你會感動嗎?上帝說不會瞞著你,他做什麼都會跟你說,多麼感動啊這件事,可能我們看福音電影的時候,我們都會覺得這些偉大的傑作,就是人人都知道他人生的目標是什麼,.
如果大家看過荷蘭的姐妹,福音的歌手,她爸爸是牧師,她出生天生沒有手,腳長短腳,跟阿叻四肢都沒有了,她就是長短腳,有人訪問她的時候,問她會不會因為沒有手很難過,她就很坦誠的說,她從來沒有為過沒有手難過,.
她可能為了其他事情很難過,但是沒有手對她來說一點都沒有問題,所以你看到她的,如果你沒見過她,有空去看看她,她的腳很厲害,她的腳可以開罐頭,她的腳可以彈琴,她的腳可以開車,長短腳都可以做的,所以當大家心情不好的時候,你看一看她,你沒理由心情不好的,.
她的人生使命是什麼呢?她的人生使命擺在那裡就是激勵人的,她不用說話,很多人因為她的緣故,本來有點不順氣的,都會復興的,我們可能會看這些,在她身上,我們覺得大家都很明白她的人生計劃是怎麼樣,但是我們不妨留意一下,上帝對阿叻的計劃,大家要站好,千萬不要錯過,.
我們每一個都是上帝認定我們的,或者亞伯拉罕的計劃,也都在我們身上要成就的,上帝對亞伯拉罕的計劃有三大,如果在創世記的十一章,我們可以看見,上帝要亞伯拉罕成為大國,他的名字是大的,.
他成為萬人的祝福,他的國是大,他的名字是大,他的福是大,今天我們教徒不要忘記,上帝在你的生命裡面同樣有這幾個目的,而這幾個目的,我們今天活的時候,求主讓我們活得出來,.
或者我們先講這裡,新約,或者我們去體會,新約也是同樣講一樣東西,正如我剛才說的,上帝給我們人生計劃,給門徒的吩咐只有兩樣,第一樣就是我們綜合了,不用看十誡,我們看新約那兩條,就是彼此相愛,對著我們弟兄姊妹,.
不是基督徒,我們就需要以大使命為我們的目標,其實這個也是亞伯拉罕的使命,我們要去建立神的國,這是天國最大的,所以如果在2024年4月29日剛出來的數字,他講到全球的人口突破81億,很厲害吧,.
如果在2023年那個計算,就是全球的基督徒人口大概是24億,意思就是說比例大概是30\%基督徒,其實地上沒有一個國家是有這麼多人救基督徒的,我們大家都是天國的子民,我們在地上是30\%個佔,.
其實很簡單,如果我們的數字厲害一點的話,我們大家都知道,如果我們一個人帶3-4個,就怎樣,全世界都信耶穌了,如果我們大家努力一點,一人帶3個,現在就用你的本份,一人帶3個,全球就已經信耶穌了,其實如果我們用這個角度來看,.
上帝要藉著他的子民,這個全球至少20多億或30億的人口,我們應該怎樣去履行上帝的吩咐,上帝要我們的國家,要我們的國民,是怎樣的呢,就是要遵守神的道,負上帝旨意,是要不斷傳揚上帝的話語,傳揚上帝的福音,以致人也都因為我們的生命受感動,來歸向上帝,.
所以今天我們在神面前的領受,我們怎樣傳福音,我們怎樣為主,來事奉主,這個是我們今天應該要有的規劃,我們應該想怎樣呢,你儘管去想,最近有個姐妹跟我說,有個異象,因為她的小朋友去到teenage的階段,她就跟我說,很想見到很多的新生子,當他少年之後,就會很多出走,離開神的,.
她很想去建立一個祈禱,招聚一些父母,願意為少年人,為兒童祈禱的家長,或者弟兄姊妹,一起去為我們少年的一代禱告,她有個這樣的計劃,當然是拍手掌支持她,好啊,你去做,大家去支持你,.
真的,剛剛開始,前幾天就開啟了這個,呼籲弟兄姊妹,如果你都願意為少年人祈禱的話,就加入這個群組,然後我就每天發一些圖文出來,來禱告,其實如果我們每個人都發起一些,不一定你要大批全部人都要跟你,不是的,有感動的就去做,.
我們最近有咖啡因,大家都知道,展秋一力承擔,其實展秋真的很厲害,我覺得真的要稱讚一下,因為在去年的時候,他跟我的同工,一直推動,我們有個使命空間,邀請一些人來跟我們分享,其實我們同工都是50/50,因為我們已經夠忙了,忙到出煙的時候,咦,又搞咖啡因,怎麼搞,不是我們熟悉的範疇,.
但是一力承擔,不單止是去建立咖啡因,然後我們有個小小的空間,不要理這麼大,就真的擺了幾張沙發,幾張圓桌,我們就開始這個咖啡的地方,買了一個咖啡機,買了一些奶茶,然後沖,一是不搞,又自可,一搞,原來不得了,原來真的很開心,.
突然之間,坐在那裡,拿著咖啡杯底,然後就聊天,在我們特別的中間,我們有個EBS,就是我們有個熱絡的Bible Study,就是我們周六崇拜的中間,就在那段時間,就喝咖啡聊天,哇,熱烈很多,人數都增長,我想至少增長了十分之一,.
這個都是我們看到,原來有心搞,肯去搞,原來就會有一些不同的特色,其實在整個咖啡因裡面,我真的出過一點力都沒有,但是上帝就會差演人,來到他出力,所以如果突然之間,我相信大家都是臥虎藏龍,待會搞的一定是臥虎藏龍,曾慕斯,個個都有他的恩賜才幹,大家就發揮,我覺得樂曲隊挺好的,.
我們發揮得到,我看到長者,很願意去事奉神,我很感動,我都希望我媽媽,不過她九十幾歲了,帶她出來唱一下樂曲,我覺得很感動,你們不用說那麼多,見到你們,就覺得信耶穌是好的,或者如果我們有班年青人,你們去跟人家打一下球,可能就會很有吸引力,.
所以當我們願意去為神的計劃,我們怎樣去參與,我們去認識上帝在我們身上的作為,有人問我,經常有人問我,究竟什麼叫恩賜和才幹,其實我簡單跟他說,你願意將你的才幹奉獻給神,就是你的恩賜,但是如果你將你的才幹收起來,不用的話,這個只是你個人的才幹,對身邊人一點作用都沒有,.
但是你將它奉獻給神的時候,你就可以成為上帝使用的恩賜,我們願不願意讓上帝所賜給你的才幹,去好好地發揮,成為神給你的恩賜,我什麼都不會做,有個肢體跟我說,有個老人家剛剛進了醫院,他那段時間進了醫院,是個嬌徒,睡在床上,什麼都做不了,什麼都不做,還想死,.
我就告訴他,你要跟他說,你想死那麼容易,你的最重大的使命是什麼?你就是睡在那裡也要為人祈禱,祈禱的價值挺有價值的,你要鼓勵他,不要讓他抑鬱到想死,覺得什麼都沒有用,什麼都沒有用,你的腦還在運作的時候就有用,你的腦不運作的時候也不用怕,因為不用你運用,因為我們已經是安然享天家就可以了,.
但是我們今天如果上帝給我們有一天的,我們可以活的時候,我們就好好去配合,或者去為上帝計劃而活,第三樣,我留了一張,這個是全球規劃信徒的人口,經常都說穆斯林就快超越我們了,好像還沒見,其實都超越了,好像都在超越了,我們都要努力,.
不過也不用擔心,因為基督徒的信主跟穆斯林的信仰是很不同的,角度是很不同的,第三樣,我們需要配合上帝的計劃,願意告訴亞伯拉罕他的計劃,而亞伯拉罕明白上帝計劃的同時,他有什麼行動呢?這個才是最重要的,.
我們大家都知道上帝無端端自言自語,說的時候又讓亞伯拉罕知道,如果上帝做的事是這麼秘密的話,根本就不會讓亞伯拉罕知道,就是因為要讓亞伯拉罕知道,所以就透露了一件事,這個索多瑪,摩拉的地方都是罪惡滔天的了,如果不是,我也早就知道了,所以現在要看看,我要去毀滅這兩個城,.
如果是一個秘密行動的話,亞伯拉罕就不會知道,但是亞伯拉罕知道的時候,他有什麼行動呢?其實我們很多時候都聽到小道消息,為什麼上帝讓我知道呢?所以你不要當不知道,不要假裝不知道,當上帝讓你知道的時候,可能就要你行動,不要說這是秘密,不要讓人知道,不是的,這些秘密如果上帝讓你知道的時候,你會怎麼樣呢?.
亞伯拉罕就是知道這個秘密,這個大秘密,上帝要審判這兩個城,於是亞伯拉罕就定下來,然後就用和上帝同行的一個行動,就去向上帝禱告,去問上帝,所以原來當上帝透露給亞伯拉罕聽的時候,是沒有令上帝失望的,上帝就是要他這樣,.
上帝就是要他去回應,借助這個知道,以致他借助向上帝禱告,而這個禱告我們也知道,其實都是幫助亞伯拉罕去更新,去明白,究竟他對上帝認識,他那個神學,他過去那套神學,可能是有問題的,所以他就要向上帝了講,.
去和上帝了講,他說難道惡人和義人你都滅,這個不是你來的,這個就是他認信的神學,但是在實際的生活上面,他就是看到上帝說要滅他們,如果滅他們,萬一那裡有50個義人,那怎麼辦呢?.
上帝就解答亞伯拉罕這個很重要的問題,所以當亞伯拉罕由祈禱50人到40人到30人到10個人的時候,我們可能問,亞伯拉罕為什麼10個就停了呢?其實可能都禱告到不好意思,50個上帝說不滅,45個不滅,30個不滅,10個不滅,他再祈禱難道一個呢?.
好像不是很好意思,如果一個都沒有,真是心想都該死了,這兩個城都該滅,兩個城當時可能有2萬人,這樣去滅法,真是都該死了,而且罪惡是很大,當他這樣禱告之後,亞伯拉罕就明白了,原來這兩個城果然是惡人的城,所以他就沒有辦法再去禱告,.
這是一個很重要的基礎,最近有一個肢體,你有提到,有一個同事,平時脾氣很大,還要時時都破壞公物,在警署破壞公物,今天被同事當場抓到,刑事毀壞,偷笑,毀壞這麼久才抓到,.
但這個肢體竟然這樣說,很難過,我要為他祈禱,希望他沒事,我就問他,他常常都是破壞公物,那你為他祈禱什麼?希望他沒事,我再問大家,當我們為一個人這樣祈禱的時候,是不是這樣呢?.
如果破壞公物,以前找不到兇手,現在找到兇手,我們是不是應該感恩?我們應該誠之於法,應該感恩,但是我們生活上,其實我們也要問一下,我們的禱告會不會濫用?我們的禱告會不會是在求上帝不公義的事情,求上帝變做成憐憫的事情,這個我們也要很小心,.
所以我們整個禱告裡面,我們看到阿巴拉禱告很有原則的,很有原則的,他相信上帝是公義的神,正直的神,他又相信上帝是有原則的神,他用這個做一個前設,很重要,既然上帝是公義正直的神,他就不是求殺人放火,奸淫擄掠那些,上帝你叫他無罪釋放,不會的,因為他知道上帝不是這樣的屬性,.
所以當我們今天祈禱的時候,我們就求上帝幫助我們有新的眼光,去看什麼是合他心意,因為就算我和你們之間,現在我們不是一個同行的關係,因為我看到的東西和你們看到的東西完全不同的,我都看不到港台,如果我不知道的話,都沒有港台,這樣也可以的,.
但是你所看到的東西,我們怎樣才能夠同行呢?就是我們大家站在一邊,我站在你的角度,我看的東西和你一樣,我和上帝同行都一樣,我們不要和上帝看到不同一樣的東西,當我們看到天災人禍的時候,我們就說上帝為什麼這麼可憐,那個腰折的,那個小孩為什麼這樣都要死的,.
我們就需要去轉一轉,謙卑,去同神並肩同行,讓我們的禱告,就好像亞伯拉罕那樣,一邊禱告的時候,他的生命就開始改變,然後上帝一一答應的過程,他就更加經歷,原來上帝是那位憐憫人的上帝,.

$^{81}$所以當他停下來之後,第二天,我們知道,當亞伯拉罕遠遠去望的時候,就看到這兩個城,慢慢無煙,因為很遠,慢慢無煙,他就知道上帝的審判臨在,他那時候都沒有再祈禱,因為他已經找到答案,他知道原來那裡連十個贏都沒有,但是即是說他沒有祈禱,上帝都因為他懸固,紀念羅德和他一家,於是上帝就救了他,.
救了三個人,所以我們今天頂尖會當我們要履行上帝計劃,我們要配合上帝的心意的時候,去求主幫助我們願意信服的心,而信服的過程,我們都要相信,禱告不是改變人類命運的一個行動,.
但是我們相信,禱告,直接禱告,是上帝可以改變人類命運的,這件事我們千萬不要走在一起,免得大家高舉,就是祈禱,我們祈禱就可以了,不是,我們祈禱背後,是祈禱邀請這位上帝去介入,去參與才重要,.
曾經有一個姐妹,她為了一個家庭很苦惱,她說這個家庭是她的至親,是因為虐待的問題,又面對離婚,但是她很不開心,她就說為她祈禱,但是上帝好像沒有做事,做不到功,我問她,你想求神什麼?.
她說求神不要打架,希望他們好好的,基本上我們都會這樣祈禱,結果是什麼呢?他們打架,然後警察上門,然後就分開他們,小朋友就去了福利署,她說結果為什麼會這樣?我說上帝不是聽了我們祈禱嗎?至少上帝聽了我們祈禱,分開了之後是不是沒有打架?.
小朋友是不是沒有被虐待?可能不切,不是你的預期,但是你要知道觀察上帝的作為,你的祈禱上帝是聽的,你的憐憫心腸上帝比你更憐憫,所以我們要知道,要相信上帝會明白我們的心意,所以我們儘管去禱告,所以當我們的方向正確的時候,方向比努力更加重要,.
當我們方向正確的時候,我們去履行上帝的心意的時候,我們的努力才有價值,我們配合上帝的工作才有價值,不然我們一路祈錯禱,我們努力所做的事情都不是上帝要我們所做的事情,我們就會前功盡廢,所以求主要幫助我們,今天我們要講到與神同行,我們要重申,我們首先要去明白,究竟我們會是上帝是誰?.
我們首先要認識這位上帝,他既然是那位大能者,上帝在你生命的計劃裡面,大家千萬不要只是做六天的基督徒,我們一至五天都是基督徒,我想不是那首歌,就是儀哥六天披著羊皮,其他日子披著羊皮,其他日子披著狼皮,.
所以我們六天才做羊,一至五天就披著狼皮,其實我們不是,我們不應該這樣,我們一至日整個生命都是屬於神的,我們沒有請假的,基督徒沒有請假的,我們每天都是基督徒,而且上帝在我們生命裡面,當我們和祂同行的時候,.
當我常常想到上帝原來要在我身上建立一個大的國度,我覺得很偉大,上帝現在也是在做的功夫,我們人的功夫很小,但是當我們願意做一點點功夫的時候,配合整個上帝的國度,這個上帝的國度是在行的,原來上帝要我們的名字是大名,我們的大名是什麼?我們的名字是什麼?我們的名字是香還是臭?.
這個我們也需要去分辨,願神要幫助我們,我們知道很多前人建立基督徒的名字是非常好的,教徒有愛心一點的,教徒不卸膊的,教徒有責任感的,教徒不能不畫的,很幫忙的,很關顧人的,這些的名字是非常好的,我們希望我們繼續去承接,承接前人給我們這個名字這麼好的,然後原來上帝很重要,.
我們的目的,我們生存的目的就是祝福別人,讓我們的人生成為更多人的祝福,願上帝幫助我們,今天我們基督徒使命感如果強的時候,我們每人其實問一下ABC就可以了,我們找一下ABC有三個人,我們就可以在全球領人歸屬上有份的了,願神祝福我們,我們今天沒有什麼藉口,我們不與神同行,.
希望我們這個與神同行,就是讓我們在神面前去禱告,讓我們真的去交代給神,先讓我們經歷,我們知道神的名字這麼大,然後上帝要用你這個名字,教徒的名字,我們求主幫助我們,如果我們想我們的名字是大的話,我們連一點點脾氣都搞不定的話,我們怎樣去見證上帝呢?.
我們面對一些困難的時候,我們都不能靠上帝經歷得到,我們怎樣有能力告訴別人上帝可以的,我們可以靠上帝,所以願神幫助我們,我們一起去祈禱,讓神的生命繼續在我們生命裡面建立..
差害天父上帝,我們要多謝你,我們要多謝你為你揀選我們,在這個時代,我們仍然可以成為一個大國,我們看到天國這個國度是沒有收縮的,反而是不斷地將見擴大,我們願意在全球有三分一的人口是屬於你的國度,我們求你祝福全球的人口,特別是教徒人口,都去努力的,.
要將福音傳給一些未信主的人,但是我們求你讓我們的生命就是將見福音,就是將見這個好消息,讓人看到原來我們的生命是有上帝的大能去彰顯,以致我們的生命是像基督的生命,我們求你幫助我們,.
今天如果我們生活上遇到神學仍有困難,困擾問題的時候,我們求你教導我們謙卑的,與你並肩同行,我們不要和你對著幹,我們不要去問為什麼上帝要你這樣做,而是我們要並肩和你一起走,求上帝幫助我們,能夠更深的明白你的心意,令我們能夠更大的,以更大的行動來配合你的計劃,主耶利幫助我們,你知道人的信心都是很少,.
但是我們願意,我們這個信心是讓你建立,讓我們弟兄姊妹互相建立,激勵,讓這個紅恩堂,這間教會能夠被你使用,在這個地方的教會裡面,你要去彰顯,高舉你的名字,讓更多人,特別是紅水橋地方的人,願意的,你要前來,來得著這個福音的好處,願你祝福,多謝你聽我們祈禱,奉耶穌基督的名,阿們..
拜拜!.
\newpage



\section{歌羅西書 1:4-6-13-23}
\label{sec:6OJe2UUeX7c}
\textbf{2024.06.30 《頂級耶穌》 歌羅西書1\_4-6,13-23 (和修版) 曾立華牧師}
\newline
\newline
連結: \href{https://youtube.com/watch?v=6OJe2UUeX7c}{\texttt{https://youtube.com/watch?v=6OJe2UUeX7c}} ~~~~ 語音日期: 2024-07-01
\newline
\newline
\hyperref[sec:RaRAzKcNgEg]{\small{< < < PREV SERMON < < <}}
~
\hyperref[sec:index_chronic]{\small{[返順時目]}}
~
\hyperref[sec:index_scriptual]{\small{[返順卷目]}}
~
\hyperref[sec:W9_hqGLMr2Q]{\small{> > > NEXT SERMON > > >}}
\newline
\newline
歌羅西書 1:4-6-13-23
\newline
\begin{longtable}{cl}
\hline
\hline
章節 & 經文 (和合本修訂版)\\
\hline
1:4 & \begin{tabularx}{0.7\textwidth}{X} 因為聽見你們對基督耶穌的信心，並對眾聖徒有的愛心。 \end{tabularx} \\ \\ \relax
1:5 & \begin{tabularx}{0.7\textwidth}{X} 這都是因著那給你們存在天上的盼望，它就是你們從前所聽見真理的道，就是福音； \end{tabularx} \\ \\ \relax
1:6 & \begin{tabularx}{0.7\textwidth}{X} 這福音傳到你們那裡，也傳到普天下，並且繼續增長，不斷結果，正如自從你們聽見福音，真正知道神恩惠的日子起，在你們中間也是這樣。 \end{tabularx} \\ \\ \relax
1:7 & \begin{tabularx}{0.7\textwidth}{X} 這福音是你們從我們所親愛、一同作僕人的以巴弗學到的。他為我們作了基督的忠心僕役， \end{tabularx} \\ \\ \relax
1:8 & \begin{tabularx}{0.7\textwidth}{X} 也把聖靈賜給你們的愛告訴我們。 \end{tabularx} \\ \\ \relax
1:9 & \begin{tabularx}{0.7\textwidth}{X} 因此，我們自從聽見的日子就不住地為你們禱告和祈求，願你們滿有一切屬靈的智慧和悟性，真正知道神的旨意， \end{tabularx} \\ \\ \relax
1:10 & \begin{tabularx}{0.7\textwidth}{X} 好使你們行事為人對得起主，凡事蒙他喜悅，在一切善事上結果子，對神的認識更有長進。 \end{tabularx} \\ \\ \relax
1:11 & \begin{tabularx}{0.7\textwidth}{X} 願你們從他榮耀的權能中，得以在一切事上力上加力，好使你們凡事歡歡喜喜地忍耐寬容， \end{tabularx} \\ \\ \relax
1:12 & \begin{tabularx}{0.7\textwidth}{X} 又感謝父，使你們配與眾聖徒在光明中分享基業。 \end{tabularx} \\ \\ \relax
1:13 & \begin{tabularx}{0.7\textwidth}{X} 他救了我們脫離黑暗的權勢，遷移到他愛子的國度裡。 \end{tabularx} \\ \\ \relax
1:14 & \begin{tabularx}{0.7\textwidth}{X} 藉著他的愛子，我們得蒙救贖，罪得赦免。 \end{tabularx} \\ \\ \relax
1:15 & \begin{tabularx}{0.7\textwidth}{X} 愛子是那看不見的神之像，是首生的，在一切被造的以先。 \end{tabularx} \\ \\ \relax
1:16 & \begin{tabularx}{0.7\textwidth}{X} 因為萬有都是在他裡面造的，無論是天上的、地上的，能看見的、不能看見的，或是有權位的、統治的，或是執政的、掌權的，一概都是藉著他為著他造的。 \end{tabularx} \\ \\ \relax
1:17 & \begin{tabularx}{0.7\textwidth}{X} 他在萬有之先；萬有也靠他而存在。 \end{tabularx} \\ \\ \relax
1:18 & \begin{tabularx}{0.7\textwidth}{X} 他是身體（教會）的頭；他是元始，是從死人中復活的首生者，好讓他在萬有中居首位。 \end{tabularx} \\ \\ \relax
1:19 & \begin{tabularx}{0.7\textwidth}{X} 因為神喜歡使一切的豐盛在他裡面居住， \end{tabularx} \\ \\ \relax
1:20 & \begin{tabularx}{0.7\textwidth}{X} 藉著他，神使萬有與自己和好，無論是地上的、天上的，都藉著他在十字架上所流的血促成了和平。 \end{tabularx} \\ \\ \relax
1:21 & \begin{tabularx}{0.7\textwidth}{X} 從前你們與神隔絕，心思上與他為敵，行為邪惡； \end{tabularx} \\ \\ \relax
1:22 & \begin{tabularx}{0.7\textwidth}{X} 但如今，他藉著他兒子肉身的死，已經使你們與他自己和好了，把你們獻在他的面前，成為聖潔，沒有瑕疵，無可指責。 \end{tabularx} \\ \\ \relax
1:23 & \begin{tabularx}{0.7\textwidth}{X} 只要你們持守信仰，根基穩固，堅定不移，不致動搖，離開了你們從前所聽見的福音的盼望；這福音也是傳給天下一切被造之物的，我—保羅作了這福音的僕役。 \end{tabularx} \\ \\
[1ex]
\hline
\hline
\end{longtable}
$^{1}$其實我不需要繼續講下去,因為葉劍霞姐妹的親身經歷已經很感動我們人心.
這就是基督教的福音信仰的精神親持示範.
福音是什麼呢?福音就是十字架.
十字架已經放在全世界的教堂上,很多西方教會都將它放在頂部.
高舉耶穌基督為人犧牲寫命的勇氣而放出來的寫氣生命.
為了救贖我們這些有罪的罪人.
剛才葉姐妹已經演繹了,經期本身有三種缺陷.
常人都覺得不可能做這麼冒險的手術,用命來碰?.
你命很大嗎?但現在你真的命很大.
坐在這裡告訴我們,活生生的經歷,如果不是耶穌基督的十字架大愛感化你,你肯做嗎?.
本來你都要退縮,最好是那個先生代替我,我不用做.
父母不用擔心,連乾爹國民前院長都不用擔心.
他也會質疑,做神學院院長也會質疑,為什麼你可以這樣做.
不過你就是這樣,本著福音生命在你裡面的那種作為.
你已經做出了非常偉大犧牲寫命的故事放在我們面前.
所以今天我想從保羅所寫的一封所謂坐牢的書信,叫做《歌羅西書》.
去年我九月到十一月有機會去意大利羅馬教書的時候.
我周末有很多時間去尋找保羅的蹤跡.
在羅馬二千多年前他留下來傳福音的蹤跡.
他寫這封書信的時候,他坐在羅馬政權附近的一個地區的監獄裡.
我親身進去看過,他就在那裡寫了幾卷所謂監獄詩信.
而《歌羅西書》就是他高舉耶穌福音的其中一卷書信.
剛才你看到,因為我聽見你們對耶穌基督的信心.
並對眾聖徒的愛心,這都是因著耶穌給你們存在天上的盼望.
祂就是你們從前所聽見的真理的道.
這就是福音,福音指向誰呢?指向這位頂級的耶穌.
保羅本來是迫害基督徒的兇手.
但是上帝就在這個人的身上動功,在這次行禮裡.
上帝的愛感召他,他就在大馬骰路上信了耶穌之後.
他生命就完全翻轉了,他從前是迫害耶穌的.
現在變成一個非常熱心傳耶穌基督為人寫命的福音信息的一位使徒.
傳了三十多年,到他人生最後,他還是犧牲了自己的生命.
在《釋量傳》第28章裡最後一句,他在坐牢裡.
有監察著他,但是羅馬政府沒有禁止他不傳福音.
所以每一天都有很多人去探監的時候,他都努力向他們傳福音.
為什麼會有一個這樣的人,生命是不是一個絕緣的,翻轉之後,扭轉之後的.
不可能他自己做,肯犧牲自己的.
保羅就是為了這個福音而寫命的人.
後來我去了他另外一個監獄,最後他人生的監獄更加小.
然後他走過一條他自己要殉道的路,我搭在他路上.

$^{41}$他被圍住了,但你可以跳過去,走到他被斬頭的地方.
這個就是今天葉子妹所帶出來的福音信息.
一個寫命,將他生命原意,將他的肝捐出來.
我問他先生,你們有沒有子女,他們沒有.
希望上帝賜福你們,這麼年輕.
但他們做了這個手術,可能各種原因,不可以再懷孕.
這個就是福音,什麼叫做福.
就是你自己做不到的,你自己本來沒有的.
我們世界上所有的人都是罪人.
從歷史的記錄來看,人性基本上沒有改變過.
兩千年來的社會形態雖然不一樣.
但人的需要,人的軟弱,人的罪性,人的恐懼.
人的目標,人的抑鬱,仍然在打擊著沒有盼望的人類.
千百年來不同的國家,人民,階層.
但正因為耶穌基督去到他生命裡面,觸碰他,他的生命就改變了.
你告訴我,這個世界任何的學說,任何的理論,任何的哲學.
任何的東西,能否改變一個人生命的本質和罪性呢?.
到現在為止,大家都很清楚.
我們裡面都有黑暗陰暗的一面.
我們做牧師,我們自己知道,一樣面對人性的試探.
會跌倒,會失敗,因為我們的罪性.
就像亞當夏娃所犯的罪,罪性在裡面,一直傳承下來到今天.
沒辦法,藉著教育,藉著哲學,藉著政治,藉著各種經濟.
更加令更多人因為沒錢犯罪.
用各種欺騙方法來詐取我們,敲詐我們的罪.
對不對?人性真是醜惡到極.
所以保羅說,只有福音的真道,福音已經傳到普天之下.
這是真實的,全世界120多個國家裡,信耶穌的人佔大多數.
72億人,裡面有30多億人,未過半,但都是信耶穌的.
他們信這位頂級的耶穌,他們從來沒聽過福音.
但當他們聽到耶穌的福音之後,憑著簡單的信心,接受福音.
接受耶穌在十字架上釘十字架所流出的寶血,來洗淨我們人類的罪.
還有一個很特別的儀式,每個月都會有一次聖餐.
紀念耶穌為我們捨身流血,就是十字架本身的重要真理.
我們每個基督徒都知道,每個月都要定期洗一次罪.
清洗我們裡面的污穢敗壞.
因為沒有辦法藉著人類的一切行為,甚至所謂的行為治療,可以改變人性.
到底保羅如何述說耶穌基督呢?.
很簡單,因為周末斯提醒我,要簡短地說.
在第十五節,保羅已經在第十五節說清楚了.

$^{81}$「愛子是那看不見的上帝之像,是一切被造物的守生者」.
即是說耶穌是永恆神的反應,是具有超越性的.
我們現在說超越性,甚至香港年青人都說超越性.
耶穌是最超越性的,因為祂本身就是神.
但祂鍍成了肉身,降卑成為了一個嬰孩.
在馬槽裡成長過來,成為窮的人.
但祂的生命顯露了一種很不同的氣質.
所以耶穌基督是太初有道的道.
沒有見過神,今天我們都沒有見過.
基督徒經常說,信耶穌是神.
「將耶穌拿給我看」.
耶穌已經復活了,祂今天在天上.
祂是最勁抽的君王,是永恆的神.
是創造宇宙萬物的神.
因為聖經裡很清楚說,只有腐壞裡的獨生子子,耶穌基督將祂表明了出來.
保羅要證明這個愛子是人體不見的神的像,真像.
就是耶穌等同神.
所以耶穌基督具有永恆不變的超越本質和最高的權威.
你說是不是頂級?.
希伯來書第一章三字說,神是神榮耀的光輝.
是神本性最超越的位置的神子.
比任何高等的文物都超其無比.
是無法比擬的.
用廣東話來說,無法頂級.
是無法頂級的.
這是香港人的清譯,無法頂級的.
保羅高舉耶穌基督.
耶穌基督就是為我們捨身流血的位置.
生存了世界三十多年.
傳了三年的福音,教導了他的門徒.
這群人全都是不肖子,不肖男子.
但他們每個人都脫胎換骨,生命變化.
耶穌復活升天之後,差派他們到世界各地.
在意大利我們見到很多講十二使徒的事蹟.
他們全都是垃圾.
收稅佬,打魚佬,賣鹹魚腩的.
做這樣做那樣的.
不是正統的,受過教育.
官府的人說,這群人是沒有受過教育的人.
但在耶穌基督復活升天之後,回到天上.

$^{121}$不需要一個世紀,這十二個門徒.
已經翻轉了整個歐洲.
全片歐洲,然後去到美洲,然後來到亞洲.
直到今天,耶穌的福音.
藉著教會的傳道人.
藉著我們這群牧師傳道人.
將聖經的真理講解給你聽.
他是復活的主.
他所成就十字架救贖.
因為他生命中的聖靈能力.
改變了人類很多的生命不同階層.
幾年前,我和一位在大陸的高層人士.
我接觸他之後,他最初不太接受我.
因為他知道我是牧師.
我跟他說福音,他不聽不聽.
最後私人飛機飛到浸會醫院做手術.
醫生說,你只有兩三天的命.
我當時向他在牆邊傳福音,他相信了.
喪禮就在灣仔長途會的一間很漂亮的禮拜堂.
裡面有很多高官,很多黨員來參加.
在吃解慰飯的時候,我和幾位同學聽我在安息禮拜講了一篇福音的信息.
他們傳口跟我說,牧師你剛才所講的很有道理.
已經結了果子,有朋友告訴我,那些人信了耶穌.
五六年前,我帶領一個查經小組,十二又是十二門徒.
那位你們都知道,是趙世曾的第一個女人,姚慧.
這個姚慧在我查經班裡,我帶她三年.
她當時已經信了耶穌,她信耶穌就是迷迷糊糊.
有人帶她去聖地看看,她當時都膽大麻雀,天分地暗那種人.
有人帶她去在約旦河耶穌洗禮的地方.
當時她感觸到耶穌,你可以在YouTube上找她見證.
我今天就不跟你講了.
她在約旦河全身受到聖靈的感動,觸動了她.
她覺得自己已經脫胎換骨,變了.
因為福音的耶穌在她的生命裡的改變.
她回到香港之後,人們罵她打麻將,她非常厭惡.
就離開了,以後她改變了.
她就走來屯門做很多服務的工作.
她整個人都變了.
什麼階層的人我都看到他們的生命的轉化.
問題是她願不願意謙卑下來.

$^{161}$成為耶穌頂級的主神.
第二方面,在第十六十七節.
耶穌是宇宙的自尊創造者.
保羅繼續說,耶穌基督是萬有的,也是教會的元首.
他在教會生活中,是宇宙最高的創造者.
我去過美國侯斯頓的太空科學館參觀的時候.
我進入了正門,如果你有去過,你都知道我不是說謊.
他在正門上寫了一句英文說話.
這代表了美國最尖端的太空科學家.
探索太空的尖端人士,他們的信念.
他們的那句說話是.
Man's flight through life is sustained by God's knowledge.
是什麼意思?.
人類要探索太空,其實我們知道是很有限.
但是我們知道一件事.
就是這位神耶穌基督.
他的智慧讓我們能夠探測到一些東西.
但是太空太宏大,太寬宏.
我們這三十年來的太空科學.
已經脫離了教科書,已經改版了.
無數的黑洞在太陽系之中.
無數的大洋系都在存在.
這個宇宙的所謂穩定,到今天都沒有變亂.
就是因為這位神的創造者.
他自己在維持著.
我們在中國讀的什麼法律,什麼定律.
就是神自己鋪排在那裡.
非常穩定地運轉.
沒有碰到錯亂,人才能夠生存得這麼安全.
你看很多災難片,虛構出來的.
多可怕,多恐怖.
一顆小小的行星衝入地球.
整個地球就可以毀滅.
是誰做的?我不想再說了.
當然你可以反駁我,說什麼理論.
我今天不是跟你辯駁這件事.
我只要證明給你聽.
這個宇宙這麼穩定.
幾十億年來這麼穩定.
就是這位神在托著這個宇宙.

$^{201}$最後,保羅說他是教會的元首.
為什麼?.
他坐在教會裡.
我們今天敬拜一位看不見的.
但復活在天上的耶穌.
世界上哪裡有宗教領袖.
死了之後能夠翻生的?.
No! Not at all!.
沒有一個,只有耶穌基督.
這位神的救主,世人的救主.
是坐在教會裡,繼續統治教會.
所以教會裡的人.
因為生命改變了之後聚集在一起.
我們有愛心,有捨棄之心.
我相信如果葉子妹你不信耶穌.
你不會這麼寬宏大量.
把自己這麼寶貴的肝臟捐給人吧.
人是自私的.
但因為有了主在心中.
你感動了,你連自己的肝也願意捐給人.
這是愛心的流露.
就是耶穌基督的愛.
傳流在教會之中.
我告訴你,在我自己.
事奉了牧羊教的幾十年.
最安全分享你自己生命軟弱的地方.
你不會在寫字樓.
一出地鐵你就知道了.
在地鐵裡聽到很多人說寫字樓的事是非非.
大家都是鼎鼎壯壯.
你敢不敢在寫字樓裡跟人說你的心事.
說你的煩擾呢?.
沒有的.
但教會的信徒可以安全地在耶穌的愛裡.
可以開闖了我們的心.
願意跟人分享我們內心的秘密.
我們的軟弱,我們的失敗,我們的罪惡.
尋求弟兄姊妹的祈禱.
代求,接納你和我.
父母子女的割裂關係.

$^{241}$大家彼此之間的愛.
可以接納在一起.
在教會的群體裡.
牧羊的生命當中才能顯露出來.
牧師不要只看我們在港島的那一刻.
幾十分鐘.
其實我們在背後做很多事你們是不知道的.
我們探訪每個家庭.
觸碰你們生命裡各樣的事.
你們需要為你們祈禱.
耶穌基督臨到你們的家.
在這裡解決你們的問題.
如果耶穌不在教會裡.
我們可以這樣做嗎?.
我們就彼此鬥爭.
這個世界就是鬥爭的世界.
你看見香港正在鬥爭.
哪有人性的?.
睜大眼睛說謊.
但我們會不會呢?.
我們教會的信徒多數都比較真誠.
光明磊落.
跟大家分享.
我已經信耶穌七十多年.
我今年八十歲.
下個月十八日就是我八十大壽.
你說呢?.
這樣的歲數的人.
都這麼有魄力去繼續傳講這個福音.
是什麼?.
不是我.
是耶穌灌滿我的生命.
我繼續努力.
我退而不休地繼續去傳講這個福音.
福音很簡單.
就是承認自己有罪.
你有很多軟弱.
你有很多毛病.
你有很多事情自己做不到.
你根本無能為力去改善自己的人生.

$^{281}$學校的公書教學.
幫不到你.
你自己沉溺在網絡世界裡.
那些罪,色情的網絡.
令很多年青人沉溺死亡.
你父母都幫不到你.
父母只能罵你,沒有用的.
你要解脫這種罪性的纏繞.
只要你把這個罪交給耶穌.
求耶穌的罪洗淨你就可以了.
你就有能力.
不是你,是聖靈.
三位一體的聖靈在你裡面.
幫你約束自己.
不會再觸碰壞的訊息.
成年人一樣.
你生命裡面,你夫妻裡面的關係.
幾十年來已經有些糾纏的事情.
要恢復和好關係.
因恩愛愛只有耶穌在你心裡.
幫助你解決家庭問題.
因此,憑藉耶穌基督的權柄.
希望你已經知道.
人類得悲釋放的唯一之路和方法.
就是你肯謙卑下來.
把你人生的主權交給他.
幾十年來,我把人生的主權.
交給神,由神來掌管.
他就留眷顧我.
剛才很欣賞施班.
唱這首中國本色化的詩歌.
開初那幾句.
林中陰很有悲情地唱出他的音調.
但一轉,就轉成喜樂.
很厲害的情緒上的變化.
就是福音改變的能力.
就是有這種果效.
就是因為這位頂級的人的耶穌.
在各階層的人士觸碰他們的生命之後.
你能獲得尊貴的生命.

$^{321}$尊貴的自尊.
尊貴的人性.
你會在很多失意的人生中獲得鼓勵.
在受苦中受安慰.
剛才那首詩歌背後不是「受安慰」嗎?.
是不是?.
所以今天又是福音主日.
又是培靈會.
多謝施班.
你們真的唱得非常好.
多謝施班指揮和音樂專家指揮.
坐在這裡打鼓.
很好聽.
再唱一次給我們聽.
在唱之前請大家低頭祈禱.
頂級耶穌.
我們今天高舉你的名.
高舉你自己十字架的拯救大能.
高舉你自己為人捨命的大愛.
真是人感覺我們何德何能.
可以這樣接受恩典的大愛.
我們坐在這裡的基督徒.
我們自己都知道.
心濃大愛.
我們才能繼續在生命的路上.
走得安康.
走得愉快.
因此我願意憑耶穌基督.
聖靈的大名.
向會眾當中一些.
心裡聽聞了耶穌.
你願意接受這位耶穌基督.
做你生命的救主.
今天你就是福音改變人生的一個起點.
你願意這樣勇敢的表達.
你要相信.
如果你願意.
我憑耶穌基督邀請你.
舉起你的手我為你祈禱.
有沒有這樣的人.

$^{361}$或者大家相信了耶穌.
或者已經在尋道.
但你還未肯定你的內心.
耶穌是否真的這麼厲害.
你不需要再和我爭辯.
也不需要辯駁.
你先進入耶穌的生命.
你就會完全明白.
福音在你生命裡帶來人生大轉化的真正意義.
你才會真正覺得寶貴.
你那時就會感恩.
祈禱感恩.
有沒有這樣的人.
主耶穌我們知道.
你的名被傳開.
福音是真直接為我們擺開.
保羅這位福音使者.
他曾經迫害你.
但他現在勇敢傳揚福音.
將你生命的道.
將你的超越性放在我們的根前.
讓我們明白.
知道你是維護我們人生道路真理生命的主.
求主讓這位主進入.
願意開拓他們的內心.
願意謙卑承認自己是罪人.
這種誠意.
以致他們在生命裡有大轉化.
好像歷來世代不同的基督徒.
所經歷的舊屬福音的寶貴經歷.
求主知道我們這個信仰.
是道理,是生命的源頭.
是我們盼望的基礎.
願主耶穌基督,福穆的大能.
繼續在我們生命裡使用我們.
改變我們的生命.
祈禱交託,奉耶穌寶貴名字祈求.
\newpage



\section{哥林多後書 3:4-6}
\label{sec:W9_hqGLMr2Q}
\textbf{2024.07.07 《我要感謝你》 哥林多後書3\_4-6 (和修版) 曾裔貴牧師}
\newline
\newline
連結: \href{https://youtube.com/watch?v=W9_hqGLMr2Q}{\texttt{https://youtube.com/watch?v=W9\_hqGLMr2Q}} ~~~~ 語音日期: 2024-07-08
\newline
\newline
\hyperref[sec:6OJe2UUeX7c]{\small{< < < PREV SERMON < < <}}
~
\hyperref[sec:index_chronic]{\small{[返順時目]}}
~
\hyperref[sec:index_scriptual]{\small{[返順卷目]}}
~
\hyperref[sec:CrvvxFK_CfQ]{\small{> > > NEXT SERMON > > >}}
\newline
\newline
哥林多後書 3:4-6
\newline
\begin{longtable}{cl}
\hline
\hline
章節 & 經文 (和合本修訂版)\\
\hline
3:4 & \begin{tabularx}{0.7\textwidth}{X} 我們藉著基督才對神有這樣的信心。 \end{tabularx} \\ \\ \relax
3:5 & \begin{tabularx}{0.7\textwidth}{X} 並不是我們憑自己配做甚麼事，我們之所以配做是出於神； \end{tabularx} \\ \\ \relax
3:6 & \begin{tabularx}{0.7\textwidth}{X} 他使我們能配作新約的執事，不是文字上的約，而是聖靈的約；因為文字使人死，聖靈能使人活。 \end{tabularx} \\ \\
[1ex]
\hline
\hline
\end{longtable}
$^{1}$各位姐妹早安.
我們一起同心祈禱.
多謝天父今早我們再一次聚集.
我們很盼望在我們大家的心靈裡面.
因為主耶穌的救恩.
我們能夠常常想起感恩.
成為一個願意學習感恩的人.
以及在我們的心靈裡面不再有任何的仇恨敵人.
而是不斷有很多的恩人環繞著我們的心靈.
以致我們每天都是學習在我們見主之前.
常常感恩.
為到就算很微小的事情.
更何況親自幫過我們的人.
我們要憤外擺在心靈裡面.
所以今早的訊息.
成為我們每一位弟兄姊妹的學習.
提醒我們這樣祈禱.
奉耶穌基督的名.
阿們.
嚴格來說今天的講道.
不算是很懂得經的講道.
當然上下文是講兩封箭訊.
保羅說你們是我們的箭訊.
我們能夠被周圍的.
馬其頓一帶地方的人認識.
是因為哥林多教會你們的生命.
接著保羅再強調.
其實你們是基督的信.
藉著我們這些傳道者去修成的.
不是用墨寫.
是用永生神的靈去寫的.
即是說保羅不斷為他們的生命.
為他們的成長去禱告.
去流淚.
去擔心.
使得哥林多教會.
能夠成為神的兒女的一個見證.
在暑假.
這個星期.
上個星期.

$^{41}$都出席了兩個畢業禮.
七月八月.
特別是七月.
「憤愛亡」的牧者.
因為幼稚園.
小學.
有很多畢業禮我們都要出席.
充滿了感恩的文化.
基督教的辦學.
很著重這一點.
今天很開心.
林校長.
我們營辦的錦瑟幼稚園.
他今天在我們中間.
我剛才跟他說過.
我會借用你們的小朋友.
你的學生.
你的子女.
他們的見證.
成為一個例子.
近兩年.
幼稚園有一個很大的轉變.
就是.
邀請每一個畢業的家長上台.
親自幫自己的小朋友.
戴四方帽.
然後.
詩儀就來叫.
每個小朋友.
現在媽媽,爸爸幫你戴了四方帽.
你就抱媽媽吃一口.
然後拍照.
這兩年成為一個很好的氣氛.
是很開心的.
你想想.
如果你的小朋友在讀幼稚園.
而你能夠上台.
幫他戴四方帽.
我感覺興奮開心.
前天.

$^{81}$又出席了徐澤林小學的畢業典禮.
很感恩被邀請.
跟麗善牧師一起.
我們在台上頒獎.
當然家長不可以上台.
因為這麼多人.
還有小朋友六年級.
已經要抱著媽媽.
開始有些難度了.
他們可能已經抗拒了.
中學就更加嚴重.
但不是的.
原來徐小.
他們又教小朋友.
我聽到校長這麼說.
原來他們教小朋友.
你要學習每天感恩.
為一些事情.
為一些人.
你要學感恩.
你在街上看到一些朋友.
你也要展現一種笑容的文化.
這樣的教育.
從小就這樣訓練.
真的帶給我們.
基督徒.
基督教.
是不是一間感恩的教會呢?.
保羅.
他說他自己能夠成為新約的執事.
他看這件事.
因為上下文是說.
他跟摩西的十誡做比較.
摩西差不多是以色列人裡面.
除了亞伯拉罕.
是最偉大的人.
摩西領受十誡法板.
現在保羅說他是新約的執事.
就是摩西之外就是他.
當然不算主耶穌.

$^{121}$他就是新約的一個好像電機人一樣.
他是一個猶太的拉比.
其實是很仇恨外邦人.
說透過耶穌.
就可以得到天堂.
等於我們猶太人的永生.
他覺得是一件很難想像的事情.
所以他逼迫基督徒.
但是.
他現在跟哥林多教會說.
你們就是我的見證人.
你們的生命改變.
是藉著我們用聖靈去挑剔的.
你帶給外邦人信主.
其中有一點.
必須.
如果保羅沒有這個特質.
他不可能心悅誠服地.
被救恩改變.
被神的愛感動.
就一定沒有了那種.
因為錯誤的神學.
錯誤的價值觀而來的仇恨.
過去就是因為他是典型的猶太教.
教法師.
法利賽人.
他熟悉律法.
精通舊約的每一個文句.
他知道.
什麼是猶太教典型的神異類.
神的選民.
除了猶太人.
沒有人是神的選民.
你想想這個包袱大到一個地步.
怎麼可能透過一個耶穌.
一個木匠的兒子.
就說可以.
你相信他.
你就可以得到永生呢?.
不可能的.

$^{161}$但現在保羅因為主耶穌的愛改變他.
主耶穌感化他.
他整個人的生命改變.
他說他現在是心悅的執事.
整個新時代開始.
就藉著這幾個使徒了.
所以.
如果我們現在能夠解剖使徒的內心.
唯一一件事.
就是再沒有對任何的民族.
有偏見.
有仇恨.
只有感恩.
能夠領受這個職份.
盡力帶領這些外邦的弟兄.
成為神的兒女.
這就是這段聖經.
大概解經上下文的意思了.
心悅的執事.
那麼.
我現在還有些燒香的味道.
不是我今天拜完神.
而是我今天在下村.
圍村.
祠堂.
做一個安息禮拜.
有一個很德高望重的長輩.
在下村裡面的幾高位份的.
信了主耶穌.
生命改變.
成為了基督徒.
我們能夠在下村的祠堂.
幫他辦安息禮拜.
當然你可以想像到.
那些汗是不停流的.
然後那些村的兄弟都不等你了.
繼續在那裡燒香.
他在那裡燒香.
你在那裡講道.
唱聖詩.

$^{201}$這樣的一個.
我第一次出席.
其中我想說的就是.
鄧小白的女兒.
雖然她還沒有信耶穌.
但是她對爸爸的愛.
那份感恩.
那份關心.
是日於言表.
感動我.
我每次去到醫院見到她.
她就在爸爸的床邊.
不斷地幫他抹汗.
幫他移除床的位置.
希望他睡得舒服一點.
這個女兒很有孝心.
感動我自己.
我也嘗試在忙裡.
我找媽媽喝茶.
晚上約定媽媽.
明天早上和她喝茶.
星期一.
媽媽說好.
八點鐘元朗.
那沒問題了.
八點鐘元朗OK了.
誰知道我去了平時她去的那間.
上去之後.
慘了沒有人.
打電話給她.
我已經不在那間很久了.
在另外元朗的再大一點的.
原來她懂得性價比.
她說那間舊酒樓的職員現在不好.
而且茶又不美.
現在在元朗那間.
原來我一直以為在家附近那間.
接著我就.
摸門釘.
再坐車下去她那裡.

$^{241}$就變成了.
我這個感恩.
不能夠隨著她的轉變.
就是說我原來很久沒有關心媽媽.
很久沒有時間和她聊天.
很久沒有和她一起坐下來.
這是一個要學的功課.
一位未信的姐妹.
子見弟兄的姐姐.
能夠這麼有對爸爸的愛.
這種愛是很難想像的.
今早和大家看的這段經文.
保羅告訴我們.
其實是說有一些人.
混入了哥林多教會.
用他們自己出色的口才.
用他們自己某些.
譁眾取寵的演講技巧.
來擄獲哥林多教會人的心靈.
使得他們開始對保羅.
尤其是保羅所傳遞的神學.
都開始有些懷疑.
保羅說.
我們豈會像那些人.
來推舉自己.
強調自己有什麼強調的神學.
或者自己口才的銳利等等.
保羅說我們不強調這些東西.
我們不會像這些人那樣.
用其他人的戰訓來給你們.
或者用你們的戰訓來給其他人.
在那裡誇耀誰的出色.
誰來強勁.
保羅說我們不要強調這些東西.
不誇耀這些東西.
你們就是我們的戰訓.
而且是寫在我們使徒的心靈裡面.
這是很難想像的.
弟兄姊妹在傳道者的內心.
你的開心.

$^{281}$就是我們開心的原因.
你的離開.
你不再在我們中間一起出入.
一起事奉.
或者你最近已經不再熱心事奉.
是我們所擔心的.
童宮我們有時候祈禱.
主要為弟兄姊妹的情況祈禱.
這就是保羅的.
上心於弟兄姊妹的心靈需要.
一個牧者.
不只是關心整頓教會.
整頓教會.
其實是讓大家更加安靜的敬拜神.
更加好的同心事奉.
所以這是一個很重要的牧羊情懷.
使徒就是這樣.
他說寫在我們的心靈裡面.
被眾人所知道所念頌.
就是說保羅他們的內心.
是人都知道的.
他就是這樣來愛每一個.
歐洲當時有很多地方的城邦.
帖撒羅尼迦.
有菲勒比.
有保羅去過的比利亞.
有去過希臘的雅典.
接著在哥林多.
就是現在這個地方.
接著輾轉回到亞洲的有爾弗所.
還有一些保羅未必有去過的.
譬如提多在那個地方.
堅格利的地方.
他未必有去過.
羅馬他都未去過.
後來他差不多死了才在羅馬.
但是眾多的歐洲不同的民族的弟兄姊妹.
都是在他的心靈裡面.
現在他就要寫這一封信.
就是很想告訴哥林多教會的人知道.

$^{321}$我們不是好像這樣.
要來嘩眾取寵.
要來藉著什麼名牌.
來推薦自己.
你們就是我們的戰訓了.
而且是寫在我們大家的.
牧者的心靈裡面的.
弟兄姊妹.
洪恩家的弟兄姊妹.
我們的服事.
希望的是你們的生命成長.
希望是你們能夠越來越認識主耶穌.
越來越你們能夠.
真的可以為主而活的獨當一面.
寫在我們的心裡.
個個都知道大法是我的飯堂.
那些職工一見到我.
就問牧師今天吃什麼.
今天有好介紹.
他會介紹我吃什麼.
我不想想.
希望能夠建立一個.
大家彼此很好的關係.
讓更多的街坊去認識.
譬如說單車店.
Pizza.
另外那個賣酒的男士.
都叫我牧師了.
雖然我沒有幫他買酒買煙.
但是他都走過了.
牧師早晨.
很開心的一個牧鄰關係.
現在我們已經擴展到Leo.
我認識了他的主席.
希望能夠大家.
你們是我們的戰訊.
如果有一天.
大法的職員和我說.
誰誰你們的教友.
說髒話的.

$^{361}$那就不是一個戰訊了.
是一個尷尬了.
所以真的弟兄姊妹.
你們就是我們的戰訊.
寫在我們的心靈裡面.
所以我們每一次都希望.
用心靈去大家服侍.
很重要就是同心.
很重要就是大家一起.
在主裡面尋求.
主耶穌得到榮耀.
很重要就是我們能夠.
不要有太多我們個人的.
很強的主觀.
大家願意一起放下.
以致尋求最好的方案.
這樣去服侍.
這就是神賞教會的成長.
不過要記得.
你們明顯是基督的信.
所以不是保羅的戰訊.
是基督的信.
是用我們透過聖靈來寫的.
不是寫在石板.
十誡.
石板就是十誡.
是寫在你們心靈裡面的心板.
所以每一個基督徒.
每一個弟兄姊妹.
你要記得.
必須記得.
你是基督的信.
你們才是基督的信.
是藉著我們牧者去修成.
所以牧者只不過是.
在你的生命裡面.
留下一些痕跡.
我們不是主角.
主耶穌在你的生命裡面.
才是主角.

$^{401}$剛才我開始已經說了.
保羅為什麼能夠看自己是.
新約的執事.
而且他還說.
不是我們憑自己能夠承擔什麼.
我們能夠承擔的是出於神.
這份對自己的評價.
對自己謙卑的認識.
很大的原因就是.
他有一個很正確的神學.
就是他再沒有仇恨.
再沒有偏見.
再沒有自己很主觀的那種有色眼鏡.
否則的話.
你根本不可能說自己是新約的執事.
新約的執事就是說.
舊約已經過去了.
一個猶太人夠膽說舊約過去.
肯定就是他有更大的發現.
就是新約耶穌的救恩.
取締了舊約.
舊約曾經出現過.
但是主耶穌已經成就了.
人類唯一的盼望.
就是神兒子降世.
所以開展的是一個新的開始.
所以保羅是介乎舊約和新約.
這麼重要的人物.
他一定是有一個很大的.
他自己生命的突破.
再沒有偏見對外邦人.
再沒有自己的優越猶太教.
有的就是神的僕人.
所以他是新約的執事.
我舉個例子.
今天可能很短.
因為我們希望多點時間.
大家準備下午開幕感恩.
一個呈獻給神的地方.
大家可能有去過Sara買衣服.

$^{441}$Sara的一個宗旨就是.
最平價的衣服.
但是最時尚的.
堅持了由1985年到現在.
就是每一個人.
就算你不是很高收入.
你都可以穿得起時尚的衣服.
是他們的宗旨.
所以你敢走進Sara.
因為他們的衣服便宜.
二三百元就可以有交易.
而且他們的款式是很時尚的.
是不會過時的.
而且很緊貼潮流.
這是他們的宗旨.
但是原來.
我看回他們的歷史.
奧爾特加西班牙人.
這位老闆.
原來他在西班牙一個很小的村莊.
一個漁村長大的.
他在那時候受盡屈辱.
十二歲.
家裡實在太沒錢了.
太窮困了.
他就要去做學徒.
開始輾轉做過不同工作之後.
後來在一間服裝店做學徒.
在那裡開始.
是一些高級的服裝店.
輾轉很勤奮.
老闆看得起他.
就讓他開始在店門.
做一些銷售等等.
老闆的女兒很喜歡他.
他也很喜歡這個女孩.
於是就出入.
後來有一個很有錢的顧客.
來到去買衣服.
慣常都是熟客.

$^{481}$就問老闆.
你的兒子在哪裡?.
他以為這個學徒就是他的兒子.
後來原來不是.
原來是一個很窮困的學徒.
那個顧客.
那個很有錢的女士.
對這個學徒.
奧爾特加.
沙瓦的創辦人.
非常之白眼.
給盡面色去看.
對他尖酸刻薄.
其實是一個顧客.
不是老闆.
這個奧爾特加.
當時是一個年輕人.
是受盡這些屈辱.
但是他內心裡面.
不知為何有這麼高的情操.
他內心裡面.
很艱苦,很勤奮去工作.
輾轉後來就和他太太.
開了第一間在西班牙的店.
他就是要賣什麼呢?.
他就開始發現.
我要決心.
將來我長大.
如果我有能力.
我要讓西班牙人.
即是他的同胞.
要有尊嚴地.
穿一些有體面的衣服.
但不是昂貴的.
這樣的一個信念.
成為一個企業.
後來有人訪問他.
他說他的成長是很刻苦.
受盡屈辱.
但是他說.

$^{521}$我從來沒有任何一次怨恨人.
沒有一個人是他仇恨.
這句話.
打動了我的心靈.
我不知道他有沒有信仰.
但他也成立了一個很大的基金.
來支持很多西班牙的福利事業.
他是現在全世界.
如果是服裝界.
第三名全世界有錢的人.
如果是全世界.
他是第八個.
大家可能沒有數得這麼遠.
的富豪.
他的富豪榜.
起碼有九十多億的美金.
資產等等.
這個就是奧爾特加.
他說我從來.
不因為我的困苦.
不因為我的成長受人的屈辱.
我來去仇恨人.
內心是這麼坦然.
內心是這樣來去.
執著於任何的人.
都可以有尊嚴地.
穿起時尚的衣服.
他說穿得起時尚的衣服.
就是莎拉的創辦宗旨.
親愛的電影姐妹.
今天一篇不算很悉經的講道.
分享更加多.
我們有沒有在我們的內心.
我記得我一直寫這個講道的時候.
2005年9月.
神學院開學禮.
就是我入學.
梁家倫院長也是剛好就職典禮.
他那晚講了一篇信息.
只是跟我們校內的同學說.

$^{561}$我是新入學的新丁.
什麼都不懂.
學兄就帶我去看看自己的房間.
自己的地方.
看圖書館.
教我很多東西.
那晚的真的好像一個陪靈會.
梁院長只是說一個信息.
你能夠坐在那裡讀神學.
是有很多人為你做了很多事.
我還沒完聚會.
眼淚就已經流了.
就是很多人令我能夠入讀神學.
我有15封推薦信.
讓我能夠進入.
見到神學院讀神學.
全部是個人推薦信.
沒有教會推薦我.
所以神學院也要破格.
來開會.
討論一下這個學生.
應不應該進來.
接著由我四年的神學.
學院也是破格.
為我開了一個基金.
是有弟兄姊妹奉獻在基金裡.
讓我讀完四年神學為止.
從來沒試過.
以後也沒有這麼特別的讀神學.
所以梁院長應該不知道這麼多.
或者可能知道.
我不知道.
他第一篇的信息.
已經令我很感動.
在座的每一個神學生.
你可以坐在那裡讀神學.
是有很多人在你的生命裡.
做過很多事.
你要帶著感恩.
去完成每一個學科.

$^{601}$我記得四年畢業之後.
因為我們要有散學禮.
接著會上台拿獎.
梁院長在我的耳邊.
跟我說了一句話.
我上去領獎.
「瑞貴,我真的很以你為榮」.
因為他看到我四年的困難.
成長.
由一個福音堂的教會.
轉向一個福音派的群體.
這樣來完成神學的畢業.
所以這篇信息更多是我個人的.
藉著剛剛這兩次的畢業禮.
的一個分享.
一個互動.
也藉著我們下午的開幕.
一個新的地方.
我們更加可以被神用.
一個分享.
很想邀請弟兄姊妹.
大家真的心裡要懷著感恩.
要時時有感恩的人.
在你的生命裡.
時時想起.
讓你能夠衷心地.
透過感恩.
驅除內心的偏見.
內心的仇恨.
使我們能夠一生.
活在一份愛裡.
一起禱告.
多謝主.
多謝你讓我們身邊.
願你得到一切榮耀.
在我們大家的心靈裡.
都來繼續工作.
奉主耶穌基督的名.
阿們.
\newpage



\section{約翰福音希伯來書 13:15}
\label{sec:CrvvxFK_CfQ}
\textbf{2024.07.14 《敬拜為主.為主敬拜》 希伯來書13\_15;約翰福音4\_24 (和修版) 吉中鳴牧師}
\newline
\newline
連結: \href{https://youtube.com/watch?v=CrvvxFK-CfQ}{\texttt{https://youtube.com/watch?v=CrvvxFK-CfQ}} ~~~~ 語音日期: 2024-07-15
\newline
\newline
\hyperref[sec:W9_hqGLMr2Q]{\small{< < < PREV SERMON < < <}}
~
\hyperref[sec:index_chronic]{\small{[返順時目]}}
~
\hyperref[sec:index_scriptual]{\small{[返順卷目]}}
~
\hyperref[sec:in2H2R6OX8g]{\small{> > > NEXT SERMON > > >}}
\newline
\newline
$^{1}$主席.
今天有首詩歌想送給大家.
牧師,還有沒有時間?.
我想送這首詩歌給你們.
《牧師》.
牧師.
每個人都能夠對主耶穌說.
主耶穌說我要原諒,請你原諒我的警察.
牧師.
祝福大家.
\newpage



\section{啟示錄詩篇 24:7-10-100:5}
\label{sec:in2H2R6OX8g}
\textbf{2024.07.21 《Church of Turkey and Smyrna》 詩篇 24\_7-10,100\_5;啟示錄 2\_8-11( NIV)  Pastor Ozan}
\newline
\newline
連結: \href{https://youtube.com/watch?v=in2H2R6OX8g}{\texttt{https://youtube.com/watch?v=in2H2R6OX8g}} ~~~~ 語音日期: 2024-07-21
\newline
\newline
\hyperref[sec:CrvvxFK_CfQ]{\small{< < < PREV SERMON < < <}}
~
\hyperref[sec:index_chronic]{\small{[返順時目]}}
~
\hyperref[sec:index_scriptual]{\small{[返順卷目]}}
~
\hyperref[sec:1eBxYFSJlEI]{\small{> > > NEXT SERMON > > >}}
\newline
\newline
$^{1}$[Pause].
Hello. Okay. Good morning, brothers and sisters..
[Cantonese].
I'm very excited. I'm very happy to be here..
[Cantonese].
Also, my wife and my son, they are very happy to be here among you..
[Cantonese].
This morning, when I came to the church, this church, also when I entered the door,.
I felt like I'm in my family..
[Cantonese].
I think this is very important because we are far away from our country, our home,.
maybe 10,000 kilometers, but still we are feeling like in the same family..
[Pause].
Yes, thank you..
[Cantonese].
Yeah. Sorry. Which?.
[Pause].
Okay. Okay. Yeah, I am Ozan..
My family and my wife, Esen, my son is Yalan..
We are coming from Turkey. Also, we are coming from Muslim background..
[Cantonese].
I came to faith 17 years ago..
[Cantonese].
Also, my wife and I, we are joining the same church for 17 years..
[Cantonese].
Our church is the local church, the biggest and also the first local church in Turkey..
[Cantonese].
We are serving in our church. Also, I am serving one Christian organization as a full-time worker, as a staff..
[Cantonese].
We are trying to reach people to share the gospel and share the God's love to people,.
especially young people in Turkey..
[Cantonese].
It's not easy, I know, because Turkey population is 85 million people..
[Cantonese].
But just 20,000 Christians, 20,000 believers in Turkey..
[Cantonese].
That's why we have a lot of work. We have a lot of work to reach people, to share the gospel, to make a discipleship in Turkey..
[Cantonese].
As I said, I feel very happy to be here, to be among you, because you know thousands of churches in the world,.
but at the same time, there is one church in the world..

$^{41}$[Cantonese].
We are members of that church, that one church..
[Cantonese].
Also, we worship the same God. Also, we say as father, the same God, as the same father..
[Cantonese].
That's why, thank you so much. You make us feel we are in the same family..
[Cantonese].
This is my third time in Hong Kong. I always told my wife and my son, I love Hong Kong, I love Hong Kong brothers and sisters..
Their faith are very strong, they are powerful in the faith..
[Cantonese].
They always say, "Oh, okay, okay, why Ozan said like this, why Ozan said that?".
But that time, we came as a family to Hong Kong, and they felt like me, the same thing, same feelings..
[Cantonese].
For example, today, it was maybe 10 minutes, but this 10 minutes is very important for us,.
because they felt like we are in family, same family, even my son..
[Cantonese].
Yeah, the Turkey is 85 million people live in Turkey, 99\% Muslim people..
[Cantonese].
25\%?.
99\%..
99\%, oh, 99\% are Muslim..
You can, yes, you can see..
Yeah, but there were many Christians in this land, in Turkey..
[Cantonese].
Many Christians, you know, from the screen, 19, 20, 27, no, 19, 14, 25\% Christians in Turkey before..
[Cantonese].
But now, it's too less, too few, just 20,000 people..
[Cantonese].
But we want to revival in Turkey, in our country..
[Cantonese].
We want to revival come to Turkey again..
[Cantonese].
Also, we want to rebuild destroyed walls in Turkey..
[Cantonese].
We want to rebuild destroyed walls Church of Turkey again..
[Cantonese].
As you know from the Nehemiah..
[Cantonese].
Nehemiah wanted to rebuild destroyed walls of Jerusalem..
[Cantonese].

$^{81}$Also, we want this same thing in Turkey, Church of Turkey..
We want to rebuild destroyed walls Church of Turkey..
[Cantonese].
You know, this one?.
Okay, as you know from the Bible, Turkey has a very important place for history of the Christianity..
[Cantonese].
Also, you know from Abraham, city of Haran, can you see?.
[Cantonese].
Haran, Antioch, Cappadocia..
We know from Bible, this place. Also, Abraham and his family lived in Haran..
[Cantonese].
Also, mountain Ararat is very important place in the Bible. Noah's Ark landed in mountain Ararat..
[Cantonese].
Also, Hittite civilization, they are very dominant civilization in Turkey in this land before..
[Cantonese].
Also, you know from the New Testament, most of Jesus disciples spent their gospel journey in Turkey, in this land..
[Cantonese].
When we look at the Acts in the Bible, many Christian communities planted this place, this area..
[Cantonese].
Also Ephesus, Colossae and Galatia are towns and regions in Turkey..
[Cantonese].
You know Apostle Paul, he was born in Tarsus..
[Cantonese].
Most of his ministry was in Tarsus, in this land..
[Cantonese].
Also, the seven churches is very important in the Revelation..
You can see from the map, Pergamon, Smyrna, Ephesus, Thyatira, Sardis, Philadelphia and Laodicea..
[Cantonese].
Lots of place from the Bible..
[Cantonese].
That's why we want to come to revival in Turkey again, like before..
[Cantonese].
Also, the other important place is Cappadocia. You know Cappadocia from the Bible? This is very, very important place..
[Cantonese].
After persecution, the Christian people went to escape to Cappadocia place, this place you can see..
[Cantonese].
If you go to Cappadocia, you can see, you can visit many underground cities in Cappadocia..
[Cantonese].
These cities built by the first believers, first Christians..
[Cantonese].

$^{121}$Also, you can see many Bible motifs on the walls in cave..
[Cantonese].
Also, Antioch. Antioch is very important in Christianity, for Christianity. Also, Turkey, maybe you heard before, there was an earthquake, very big earthquake in Antioch, in Turkey..
[Cantonese].
Modern name is Antakya..
[Cantonese].
Ancient church destroyed, but spiritual church will be stronger than before. Kingdom of God will come again to Antioch..
[Cantonese].
Many buildings collapsed, destroyed, many people escaped from Antioch..
[Cantonese].
But we want to rebuild again in Antioch, but not just building, we want to rebuild spiritually..
[Cantonese].
As I said, we are coming from Izmir, Smyrna..
[Cantonese].
The Bible says about Smyrna church in the Revelation book called Faithful Church. Smyrna is Faithful Church..
[Cantonese].
I can say this, Smyrna Church is still Faithful..
[Cantonese].
Smyrna Church is reflecting God's love to people in Turkey, in Izmir..
[Cantonese].
You know from the Revelation 2, 8, 11..
[Cantonese].
Brothers and sisters, we want to raise the church up with your prayers..
[Cantonese].
I want to encourage you to pray for Turkey, pray for Middle East..
[Cantonese].
Because it's hurt my heart because there were many Christians in my country, in my land..
But now, just 20,000 people..
[Cantonese].
God's unique love must reach again in Turkey..
[Cantonese].
May the seven churches arise again in Turkey..
[Cantonese].
May the seven churches rise again in order to Kingdom of God come to Turkey again..
[Cantonese].
And may the church in this land be stronger and greater than before..
[Cantonese].
Maybe I can read Psalms 24, 7, 10..
[Cantonese].
This time I can read, English?.

$^{161}$Yeah..
I don't know Cantonese..
I will translate to Chinese..
Lift up your heads, you gates, be lift up, you ancient doors, that the King of Glory may come in..
Who is this King of Glory?.
The Lord is strong and mighty..
The Lord is mighty in battle..
Lift up your heads, you gates, lift them up, you ancient doors, that the King of Glory may come in..
Who is He?.
This King of Glory..
The Lord Almighty, He is the King of Glory..
[Cantonese].
In Turkey, to be Christian, it's not easy. It's difficult..
[Cantonese].
There are social pressure from your families, from relatives, from job, friends, from everywhere..
[Cantonese].
You have to pay many costs from your life..
[Cantonese].
If you come to Jesus in Turkey, you may lost everything in your life..
[Cantonese].
Sometimes it's very difficult to walk with God..
[Cantonese].
Even you may lost your identity..
[Cantonese].
Because if you say to people, I was Muslim before, but now I am a Christian, they are looking you very weird, very strange..
[Cantonese].
As I said, I am coming from a Muslim background. My whole family, my whole relatives are Muslim. Just I am a Christian in my family..
[Cantonese].
I came to Jesus through two dreams..
[Cantonese].
The dreams are too long, but long story of short, I can say, maybe you heard before, in the Middle East, God speak people through dreams..
Like me..
[Cantonese].
My first dream, God touch my back and I got healed and I asked to sky, who is this hand?.
[Cantonese].
Sorry..
Who is this hand?.
[Cantonese].
And the majestic voice said, this hand coming from David's descendant..
David's descendant..

$^{201}$[Cantonese].
But,.
[Cantonese].
But when I woke up, I was shocked because I didn't know who is David..
[Cantonese].
I didn't know Jesus descended. Who is Jesus, where is coming from Jesus?.
[Cantonese].
In my second dream, that time we were just friend, my wife in university and we were walking on the flat area..
When we were walking, the big mountain appeared in front of us..
[Cantonese].
The big mountain..
[Cantonese].
And she said, we have to climb this mountain. I said, no we cannot do that because no one has been able to climb this mountain..
[Cantonese].
When we were discussing, the man who is coming from the sky, he said, Ozan and Esen, I will suffer for you and I will reach top of the mountain for you..
[Cantonese].
And after that, Jesus climbed the mountain. There were many details but I am passing details..
When Jesus reached top of the mountain, the mountain suddenly disappeared..
[Cantonese].
In my dream, I was shocked because no one has been able to climb this mountain but Jesus did. And I said, how could Jesus do that?.
[Cantonese].
And after that, I woke up, three days later, middle of night, I said, okay Jesus, I accept you, I believe you, you are my savior..
[Cantonese].
And then I came to Jesus, 20 years old and I graduated. But that time, I thought that, oh, my life will be very nice, very good, I will be a powerful Christian..
[Cantonese].
I graduated university, I went back to my parents' house..
[Cantonese].
But every Sunday, I tried to go to church..
[Cantonese].
And one day, I went to church and come back to home and my father asked me, where did you go today?.
[Cantonese].
I said, I went to, I met with my friends, just friends..
[Cantonese].
But he asked again. Normally, my father never asks two times. But that time, he asked me, where did you go today?.
[Cantonese].
And this time, I said, I went to church..
[Cantonese].
He said, asked me, why did you go to church?.
[Cantonese].
I said, church is very beautiful place, I love Christian people, they are very friendly, they love me..

$^{241}$[Cantonese].
And after that, he started to say bad things to Christian people, and church, and Christianity..
[Cantonese].
I said, why do you say bad things to Christian? Because I know them, I know church, I know Christianity, I know brothers and sisters, I felt love from them..
[Cantonese].
But my father was getting very angry..
[Cantonese].
And I realized that, when my father said bad things to Christianity, I realized that I defend my brothers and sisters, also church..
[Cantonese].
But I was just six months believer, I didn't know many things about Christianity, many things about Jesus, but I know Jesus changed my life..
[Cantonese].
In my first dream, Jesus touched me, God touched me, and I got healed physically, but I know Jesus changed my life, and Jesus gave me spiritual healing..
[Cantonese].
We were discussing, it's very, very strong, my father was getting angry, and my mother, my sister, they were crying, they were saying bad things..
[Cantonese].
And after three hours, my father, he came to my, near my knee, and he said like this, "You have to choose right now, me or Jesus.".
[Cantonese].
It was very difficult, because as I said, I didn't know many things, but also, I had a very nice childhood, I felt love from my family, I love them so much..
[Cantonese].
But that time, it was very difficult, when I remember this time, I also feel like very sad. But in this time, I said to my father,.
"I don't want to choose right now, because I love you, and I know you love me.".
[Cantonese].
But he said, "You have to choose right now.".
[Cantonese].
And I said, "Okay, father, okay, I choose Jesus.".
[Cantonese].
The most difficult decision for me..
[Cantonese].
Because I love, you love family, right?.
[Cantonese].
Family, your family loves you, you know, you feel..
[Cantonese].
But that time, I believed that I have to choose Jesus..
[Cantonese].
And my father looked my eyes and he said, "Okay, you are not my son, you have to go from my home, this is my home.".
[Cantonese].
I was crying and said, "How can I go? Where can I go? I had no money, zero, I had no job.".
[Cantonese].
But he said, "You have to go from my home, because this is my home, you cannot stay here.".
[Cantonese].

$^{281}$My mom and my sister were crying, and I took my bag and I went out from the home..
[Cantonese].
I was crying, I had no money because the time is 11pm, 11-12pm..
[Cantonese].
I was very angry to God because I said, "I choose Jesus you, but I lost my family.".
[Cantonese].
I saw many miracles in my life from God..
[Cantonese].
Some miracles I cannot remember, but there is one miracle I cannot forget..
[Cantonese].
First miracle from God in my life, first miracle..
[Cantonese].
I took my bag and I went out from home, I was crying because I had no money, nothing..
When I was praying and also when I was crying, I opened my eyes and I found very, very small Turkish tailor on the ground..
[Cantonese].
This detail is very important for me, also for God I know..
Just this small tailor, just was enough to use one-way train..
[Cantonese].
Not for return ticket, just one way..
[Cantonese].
I put my card and I went to my Christian brother's house..
[Cantonese].
I stayed with them three months..
[Cantonese].
Also I joined the church, I joined youth group, my brothers and sisters take care of me..
[Cantonese].
My parents didn't talk with me five years, they didn't connect with me, also they didn't come to our wedding..
[Cantonese].
I tried to speak with them but they don't want to do connection with me..
[Cantonese].
But when I started to speak, I said to you, if you come to Jesus, you may lose many things from your life..
[Cantonese].
But at the same time, when you come to Jesus in Turkey, God gives to you new brothers and new sisters, also new family..
[Cantonese].
Like you, even you are far away from your country, thousands of kilometers, you may feel like in same family, your brothers and sisters like you..
[Cantonese].
That's why I said you, I'm feeling like in my family with you, among you..
[Cantonese].
Also Jesus gives to you new identity..
[Cantonese].

$^{321}$Also this identity cannot be changed, cannot be lost because this is coming from the heaven, coming from the God..
[Cantonese].
As I said you, as I told you, we are coming from Muslim background, but also Esen and I, my wife and I are attending same church for 17 years..
[Cantonese].
Our church has 250 members, local Turkish..
[Cantonese].
I encourage you, if you come to Turkey, if you come to our church, when you speak people, when you talk with people, you can hear lots of miracles from God..
[Cantonese].
I'm just one of them, just me..
[Cantonese].
But if you talk with other members of church, you will find more miracles happening to you..
[Cantonese].
God is doing many, many miracles in Turkey, in Izmir..
[Cantonese].
Our church planted more than 40 years ago..
[Cantonese].
Old building was under the apartment..
[Cantonese].
Has been subjected to many pressures from, in the past, from the government, from the different faith groups..
[Cantonese].
Even from my father..
[Cantonese].
After my family reject me, my father came to church..
[Cantonese].
And middle of the sermon, he stopped the sermon, and he said bad things..
[Cantonese].
And he said, you steal my son, I will complain you to government, complain to police..
[Cantonese].
Despite all things, our church stay faithful..
[Cantonese].
It was closed many times in the past, many times oppressed by the police, and slandered by the people..
[Cantonese].
Like in the revolution, still is faithful..
[Cantonese].
25 years, our church met in this building, under the apartment..
[Cantonese].
25 years, but it's not too big, it was not too big..
[Cantonese].
But our numbers, the people's come to faith, and one time we couldn't fit this church..
[Cantonese].

$^{361}$And we pray for new place..
[Cantonese].
10 years later, God provides new place for us, for our church..
[Cantonese].
And we moved new church, and now our church is growing..
When you come to Turkey, when you join our church, you can see our church is growing because people coming to faith..
[Cantonese].
Again, I encourage you, brothers and sisters, come to Turkey, come to Izmir, see our church, the other churches, because there are a lot of churches in Izmir..
[Cantonese].
The harvest is ready in Turkey, but workers are very few, that's why we need to sometimes, we need to help from brothers and sisters like you..
[Cantonese].
Because as far as I see, I understand you are very faithful and strong in the faith..
[Cantonese].
Asian church is growing fast, so fast..
[Cantonese].
And I think God has a purpose for your nations..
[Cantonese].
To go to other nations, to bless the other nations..
[Cantonese].
To reach the other nations..
[Cantonese].
To bless the other nations, be guide to other nations..
[Cantonese].
That's why I encourage you, come to Turkey, come to Izmir, pray for us and hug brothers and sisters..
Your brothers and sisters, ask them what do you need, how can we pray for you. This is very important..
[Cantonese].
Harvest is ready. Do you know the olive tree is very famous in Turkey?.
[Cantonese].
Olive is very delicious, olive oil actually very delicious, very valuable in Turkey, also everywhere I know..
[Cantonese].
But the olive harvest is, to reap the harvest is very difficult..
[Cantonese].
You cannot do that alone, you need the people..
[Cantonese].
You cannot collect alone, you need the people to collect the olive..
[Cantonese].
Because the harvest is, of the olive tree is abundant..
[Cantonese].
That's why God has a purpose, I think your nations, because in order to bless other nations, to reach other nations, to go to other nations..
[Cantonese].

$^{401}$God wants to use you brothers and sisters, for the all nations..
[Cantonese].
I want to pray for you brothers and sisters, thank you so much for the invitation, good opportunity..
But I will pray Turkish, ok?.
[Cantonese].
I mean, thank you..
Thank you..
\newpage



\section{馬太福音 9:35-38}
\label{sec:1eBxYFSJlEI}
\textbf{2024.08.04 《求主差遣》 馬太福音 9\_35-38 (和修版) 譚立晶姑娘}
\newline
\newline
連結: \href{https://youtube.com/watch?v=1eBxYFSJlEI}{\texttt{https://youtube.com/watch?v=1eBxYFSJlEI}} ~~~~ 語音日期: 2024-08-06
\newline
\newline
\hyperref[sec:in2H2R6OX8g]{\small{< < < PREV SERMON < < <}}
~
\hyperref[sec:index_chronic]{\small{[返順時目]}}
~
\hyperref[sec:index_scriptual]{\small{[返順卷目]}}
~
\hyperref[sec:sujYDAwGYXA]{\small{> > > NEXT SERMON > > >}}
\newline
\newline
馬太福音 9:35-38
\newline
\begin{longtable}{cl}
\hline
\hline
章節 & 經文 (和合本修訂版)\\
\hline
9:35 & \begin{tabularx}{0.7\textwidth}{X} 耶穌走遍各城各鄉，在他們的會堂裡教導人，宣講天國的福音，又醫治各樣的病症。 \end{tabularx} \\ \\ \relax
9:36 & \begin{tabularx}{0.7\textwidth}{X} 他看見一大群人，就憐憫他們；因為他們困苦無助，如同羊沒有牧人一樣。 \end{tabularx} \\ \\ \relax
9:37 & \begin{tabularx}{0.7\textwidth}{X} 於是他對門徒說：「要收的莊稼多，做工的人少。 \end{tabularx} \\ \\ \relax
9:38 & \begin{tabularx}{0.7\textwidth}{X} 所以，你們要求莊稼的主差遣做工的人出去收他的莊稼。」 \end{tabularx} \\ \\
[1ex]
\hline
\hline
\end{longtable}
$^{1}$(自動翻譯).
我們一起去敏銳神的心意,好嗎?.
我們一起先祈禱.
主耶穌,多謝你.
多謝你昔日願意走遍各城各鄉.
主耶穌,多謝你.
因為在我們生命可能曾經歷過漂泊迷茫.
有時我們或許是經歷孤立被邊緣.
又或者我們走迷了路.
你會差派一些很愛我們的門徒,基督徒圍繞在我們身邊.
以致我們可以經歷到上帝的真實和愛.
主啊,求你幫助我們在座的每一個.
因為我們每一個都是你所愛的.
亦都是你所揀選的.
願意我們被你揀選的時候.
我們知道這是一個很大的救贖.
亦都是一個很大的榮譽.
如果我們能夠一生事奉你.
是真的好得無比.
我們所賺取的不是世上短暫的一些財寶虛榮.
而是可以賺取永恆.
和你的關係.
能夠經歷你自己寶貴的愛.
亦都可以經歷原來大人信主那份無比的喜樂.
主啊,求你使用我們每一個.
願我們所說的,所聽的都同感日靈.
願耶穌,聖靈,天父親自對我們說話.
我們誠心祈禱,交託.
奉主名求,阿們.
聖經也有說過.
學生不能大過先生.
即是說耶穌所做的事.
我們可否不跟著做.
其實不可以的.
我們很多事都不會被豁免.
耶穌是怎樣做.
我們就跟著怎樣做.
在35節的經文.
「耶穌走遍各城各鄉.
在會堂裡教導人.

$^{41}$宣講天國的福音.
又醫治各樣的病症.
耶穌說.
走遍.
其實是願意去走.
走去哪裡.
原來原本的翻譯是這樣的.
不是怎樣到怎樣.
其實是圍繞一個地方.
即是圍繞圍繞去走.
所以昔日耶穌去傳道的時候.
他是願意去不同的地方.
去找不同的人.
如果你們記得.
新約其實是.
耶穌是到處去遇見人的.
找人的.
有些是特意避開人群的.
但是耶穌就找到那些人.
然後可以跟他們.
談一些心事.
然後又將他們帶到神面前去經歷.
在我們去澳門的日子.
我們就學了.
大家知不知道澳門.
乘車經常遇到一個問題.
是什麼問題.
多不多人去澳門.
塞車.
有時可能你走路.
比乘車還要快.
即是當交通擠塞的時候.
我們就學了.
耶穌去走遍各城各鄉.
難得我們一家人.
可以去到澳門這個地方.
說要認識當地的文化.
為什麼我們不跟耶穌一樣.
就一直走走看看呢.
我不是典型的香港人.

$^{81}$我不知道你是否典型的香港人.
原來當我去到澳門的時候.
澳門人對我們香港人.
有不同的評價或者感覺.
有些人很喜歡我們香港人.
為什麼呢.
好像給錢他們賺刺激經濟.
但有些人又不太喜歡.
他們會說每逢假日.
你們快要踩倒我們的島了.
我不是典型的香港人.
因為以前很多香港人.
真的會去澳門.
但我不是這樣去的.
我真的去到有機會.
做一個跨文化體驗.
我老公之前在.
宣導會員機堂服侍.
我在宣導會的福音堂服侍.
因為他的安息年.
可以申請一個跨文化體驗.
我們可以去九個月.
去看看上帝.
給我們下一里路怎樣走.
所以我們就去體驗去經歷神.
我們就跟著去走.
我為什麼對澳門沒有特別的感覺呢.
就是因為我爸爸以前經常去澳門.
我爸爸去澳門做什麼.
大家猜到嗎.
你們很聰明.
所以我對澳門沒有什麼好感.
我爸爸就是經常來澳門.
然後回來.
看他風光是怎樣.
一條一條鏈加上去.
銀鏈金鏈手鏈.
越來越粗.
風光.
但什麼時候看到他倒霉呢.

$^{121}$你們這麼開心.
一條一條少一條.
拿去當.
以前很流行用當舖.
所以我們家裡.
很多不愉快的經歷.
因為會倒錢的人.
會有迷信.
信不信扯.
我們小時候沒有街去.
我最深刻是.
我星期天不能看卡通片.
只可以看一件事.
賽馬.
我不知道這個是不是你的童年.
我的童年就是.
星期天沒有街去.
因為爸爸要專心做功課.
然後努力功課完.
看看有沒有成績.
我有一次很深刻的經歷.
就是到今天還記得.
有些人說原來在童年.
那幾年十年.
如果家裡的關係不好.
其實是影響他一生.
你不想記得長大也會記得.
我記得那天小朋友的時候.
看馬.
但你知道小朋友的世界.
很繽紛純真.
沒東西玩看馬會怎樣.
也會用心去看.
大家知不知道賽馬之前.
會做什麼.
繞圈.
你有看.
多謝你給我反應.
我真的看馬繞圈.
然後我看中了一隻.

$^{161}$我真的覺得牠好像很輕盈.
毛身很亮.
很有活力.
我就記住牠.
十號.
為什麼到今天還記得呢.
你一會兒就知道了.
到開跑了.
小朋友會怎樣.
十號 十號 十號 十號.
誰知道.
我那隻跑中了.
但糟糕了.
我爸爸不是買那隻.
我馬上被他罵到死.
然後就.
雞毛掃.
以前最流行.
所以很痛苦的經歷.
他不會稱讚你.
有眼光.
他只會說你害我輸了.
原來在一個不信主的家庭裡長大.
被周圍的風俗影響.
有時真的會破壞親子關係.
所以那時我對澳門.
只有一個負面的感覺.
我覺得澳門和香港很近.
如果我要去宣教.
我想去遠一點.
更缺少資源的地方.
所以我沒有想澳門.
誰知道很奇妙.
那年有機會去澳門的時候.
神開了我的眼界.
不知道大家有沒有去過澳門旅遊區.
有嗎?有吧?.
路很闊.
每間酒店都很香.
刺激你的嗅覺.

$^{201}$很榮華富貴.
但我正正被安排住在舊區.
舊區40年的唐樓.
每天上街都看到活生生的蟑螂.
手指頭那麼大.
有時會看到老鼠.
很平常.
讓我看到原來可以有很大的對比.
以為澳門很繁華的一面.
原來只是一部分.
真正本土的澳門.
其實他們的生活不是我們想像中那麼理想.
我記得當時去到澳門.
我在新橋塘有機會去參與服務.
如果大家在宣道差會的網頁.
找澳門工場和新橋塘.
你就知道是那間教會.
因為我是在9月去到.
去到另一年8月回來.
差不多臨近中秋.
當地教會就想.
不如我們也做一些街坊的報道.
接觸街坊的活動.
我們要親自去接觸街坊.
才知道街坊的狀態.
所以我覺得很好.
剛才主席說流淚與默的地方.
希望能夠吸引街坊一起來.
我老公的經歷.
也是在街上走著走著.
被教會做外展的時候.
被拉到旁邊去決志祈禱.
所以擺街站很重要.
讓街坊看到我們.
我們去接觸街坊很重要.
我們不是主要的童工.
所以童工說要做什麼.
我們支持就可以了.
因為團隊合一.
做什麼都是夢主如立.

$^{241}$神都會使用.
童工就說有感動.
不如我們擺街站.
既然臨近中秋.
不如我們玩猜燈謎.
在香港還會不會玩猜燈謎呢?.
有部手機搜尋.
那句是什麼.
都出了答案.
何必猜呢?.
但我們覺得.
跟隨.
既然上帝給你有感動.
你在這裡久一點.
跟隨.
好.
總之我們大家配合.
很簡單.
拿幾塊剪板下去.
放十幾條燈謎在上面.
然後圍著環境.
邀請街坊進去猜燈謎.
猜完後.
我們教會又有份心意禮物給你們.
我們都在想.
究竟上帝有什麼作為呢?.
如果在香港就應該不可行.
但在澳門會怎樣呢?.
怎料令我很讚嘆.
當日接觸了二百多個街坊.
有二百多人願意進去玩猜燈謎.
發生什麼事呢?.
原來澳門人根本很純樸.
本土澳門人很純樸.
很簡單.
有東西玩就去玩.
真的可以放下手機.
不看.
進去做.
拿份禮物很開心.

$^{281}$二百多人.
到聖誕節的時候.
我們想做多一點.
因為有二百多人.
可能第一次試.
不敢接觸街坊.
輕盈一點.
到第二次.
不如和街坊一起做手工.
你猜在香港有沒有人陪你做手工呢?.
去到那裡.
做聖誕節的手工.
我們有小朋友部分.
有中小學生.
也有長者做手工.
也是二百多人.
二百多人願意進來做手工.
藉著做手工我們就說耶穌的誕生.
說耶穌的救恩.
我們每個人都覺得神很奇妙.
讓我們看到原來澳門還有得做.
就像香港八九十年代那種很純樸的氣氛.
一起去做.
一起去參與.
令我很讚嘆神的作為.
所以我們在澳門.
通常都是用腳遊走.
可以走的我們都走.
當然我們沒有試過跨橋走.
原來去到不同的小街.
小巷.
你會看到一些民生的生活.
很多的不簡單.
那時候.
如果你有印象.
澳門那時候一開始的防疫不是做得很好嗎?.
他們經常很讚嘆.
他們是零個案的.
怎樣都是零.
我們外面的人不清楚.

$^{321}$就覺得真的很厲害.
原來進去的時候.
了解他們的生活.
了解他們的醫療.
終於明白多了一點.
我不知道你們有沒有聽過.
他們說當地的醫療.
有兩個體制.
有私營有公營.
他們說.
一間是謀財.
另一間是甚麼?.
你們都知道了.
我說這麼誇張?.
怎會這樣?.
因為當時還是疫情.
所以我們要住14天酒店.
所以我開了澳門的新聞.
澳門的介紹.
賣私營醫院.
很多設備.
很漂亮.
很專業.
一聽到弟兄姊妹說醫療的黑歷史.
我說這樣的嗎?.
然後我說不是.
但我看你們的器材設施很好.
對呀.
所有東西都很好.
但沒有人.
你就明白了.
原來他們缺乏很多專業人士.
後來慢慢了解.
原來沒有疫情.
港澳來往很方便.
原來很多澳門人.
寧願來香港看醫生.
他們不會留在澳門看醫生.
如果可以.
因為他們覺得.

$^{361}$我付一筆錢.
我在香港看醫生.
可以應付.
可能我付同樣的錢.
在澳門本身是小事.
但還處理不到.
然後還要搞到很大件事.
然後還是要回香港醫.
會這樣的.
令我覺得.
我估不到.
我想像不到.
一河之隔.
諸葛口之隔.
有這麼大的差別嗎?.
所以我很感恩.
原來耶穌走遍各城各鄉.
按著上帝給我們的感動.
上帝要我們怎樣走.
在哪裡走.
我們要看到什麼.
我們要做什麼.
盼望我們跟著耶穌的腳踏.
耶穌做了什麼.
在35節裡.
其實是說耶穌傳道的方法.
在會堂裡教導.
也宣講天國的福音.
也醫治各種病症.
所以耶穌最主要做的是什麼.
祂會說聖經的真理.
也會說福音的重要.
因為福音是永恆的.
救贖的.
有很多難關.
不容易過.
但原來對天國有盼望.
人就會容易過.
然後也醫治各種病症.
因為昔日有些.

$^{401}$無論是大麻瘋.
不潔症的病.
就會被人鞭完.
沒有人理會.
好像等死一樣.
但耶穌也沒有忘記這些人.
所以耶穌做什麼.
我們就去學.
怎樣去做.
36節其實是說耶穌傳道的心腸.
其實我們每個人都是耶穌所拯救的.
我們有沒有耶穌的眼光呢.
我們有沒有耶穌的心腸呢.
祂說.
祂看到一大群人.
就憐憫他們.
其實最重要的是憐憫.
其實我們每個人都有眼睛.
我們都可以看到很多東西.
但有沒有動到你的心呢.
憐憫就是我們.
好像耶穌說.
很多時候大家聽過的描述.
就是動了慈心.
如果是其他人.
麻瘋病人當然會走.
但耶穌動了慈心會摸.
所以當我們看到人.
或者我們坐立在洪水橋.
我們看到街坊.
我們有沒有動憐憫的心呢.
如果他們沒有信主.
究竟他們將來會怎樣呢.
我們會否為他們願意去做一些事.
或者看到他們的屬靈光景的時候.
我們願意去為他們多做一步呢.
那裡說.
36節說.
因為他們困苦無助.
像羊沒有牧人一般.

$^{441}$即是說原來沒有信仰.
沒有耶穌的人.
可能他們會走迷了.
他們也不知道自己在做什麼.
如果在困苦無助那裡.
原本的翻譯是說什麼呢.
其實是說他們被掠奪.
即是好像搶東西.
被剝皮.
而且那個無助.
簡直是形容到.
跌倒了.
躺在地上.
即是沒有人扶起的狀態.
所以我們看不看到.
雖然每個人好像正正常常.
走來走去.
但我們看不看到他們的屬靈光景呢.
其實他們很需要住呢.
我印象很深刻的是.
爸爸這樣去帶動我們家.
其實我們家裡有很多眼淚.
我們小時候被訓練.
哭著來吃飯.
但不可以說話.
明白嗎.
所以飯是什麼味道呢.
是鹹的.
爸爸因為很暴躁.
有時他好像控制不了自己.
好像一爆發就收不到了.
爸爸一生氣.
一不高興.
就會進廚房拿刀.
所以我們小時候這樣的成長.
很不容易 很艱難.
但我自己經歷過.
以前華人社會都說.
家醜不出外傳.
所以有苦就自己知.

$^{481}$而且很明白.
對人歡笑就背人仇.
我在街上走路.
沒有人知道我有什麼事.
但自己眼淚在心裡流.
什麼都知道.
所以我很明白.
就算我看到很多人.
很正常地在街上走路.
我會知道.
每個人都可以有故事.
所以你要肯去觸碰.
你要肯去跟他有關係.
他才會告訴你一些事情.
不然那些人就很慘了.
因為他也要每天過一天.
我記得青少年期.
每一天都很痛苦.
因為每一天都是上學.
上學是最好的.
因為可以離開家.
但放學又要回家.
我記得經歷了很多年.
鎖匙到了門口.
要掙扎開不開門.
一進去就要對著一個.
今天不知是否正常的爸爸.
我很羨慕別人.
人家在街上.
兩父女可以牽手.
可以翹手.
有話題.
我和爸爸是沒有話題的.
他叫我.
只有幾件事.
當我是小妹妹.
拿水來.
拿什麼來.
小妹妹這樣做.
要不叫我來就是發洩.

$^{521}$我下面還有兩個孿生弟弟.
那個年代還有那些重男輕女.
所以怎樣.
都會投訴我.
那些很喜歡挑撥離間的弟弟.
很不容易.
總之覺得自己很卑微.
很渺小.
但沒有辦法.
所以.
生得我這麼小.
是因為怎樣.
又不開心.
又沒有一口好吃.
又經常哭.
怎可以長高.
長大.
所以我特別對小小的.
有感動.
很喜歡捉小小的一起聊天.
看看這麼小的是不是有其他原因.
關心一下.
所以.
我很希望去到澳門的社區.
我很願意和不同的澳門人聊天.
和他們了解澳門的情況.
了解他們的情況.
有沒有事情可以做到.
以致他們沒有覺得無助.
原來澳門真的很小.
其實澳門的面積.
當然現在不斷在填海.
但如果你用腳步走.
其實好像元朗的大小.
人口只有幾十萬.
其實我們香港是他們的十倍.
所以他們有很多事情.
不是很專業.
所以就.
困在裡面.

$^{561}$所以有很多艱難.
進去住.
進去生活.
進去和他們聊天.
就明白了.
去到澳門的時候.
另一個令我驚訝的地方是什麼呢.
很多.
補習社.
很多託管.
後來這些弟兄姊妹說.
你知不知道在澳門除了旅遊業賺錢之外.
還有一樣必定賺的.
我們就說什麼什麼.
就是補習.
為什麼.
因為原來有賭業,旅遊業.
父母進去上班.
有沒有人理小朋友.
小朋友不如就拿去補習.
拿去託管.
很誇張的.
就是將小朋友從科學去那裡.
吃飯,工作.
做練習,溫習.
到晚上八,九點.
最厲害的.
才爺爺,媽媽.
爸爸,媽媽來接.
你想想.
如果親子關係不好.
到青春期,暴風期的時候會怎樣.
其實衍生很多情況.
所以有一些.
比較久的澳門人.
有兩邊不同的看法.
有些覺得澳門多了一個地方.
可以去旅遊.
但對一些本土的澳門人來說.
他們覺得.

$^{601}$澳門怎樣了.
變了.
以前澳門不是這樣的.
應該很魚香,很樸素.
現在混雜了很多人.
完全不同了.
這些就是上帝開了我們的眼界視.
所以那年我去完之後.
我們一家人都受感動.
明白上帝的呼召.
所以我們計劃之後會再回澳門.
現在在這段短時期.
可能即日來回.
或短期去幫忙.
很多事情不簡單.
但要去到才知道.
原來澳門是一個溫水煮蛙的狀況.
其實是放了很多代價.
當看到這些處境的時候.
耶穌怎樣說呢.
在37,38節.
其實是在說耶穌傳道的吩咐.
祂怎樣吩咐我們呢.
我不知道上帝讓你看到的是洪水橋的需要.
讓你看到的是外地的需要.
是哪裡的需要.
祂怎樣跟門徒說呢.
37到38節.
不如我們一起讀吧.
預備起.
「於是他對門徒說:要收的莊稼多,做工的人少.
所以你們要求莊稼的主,差遣做工的人出去收他的莊稼」.
你看到耶穌不是馬上叫你們去做.
不是的.
耶穌也有祂吩咐門徒的次序.
祂先說了.
根本現在的處境是不合比例的.
收的莊稼很多.
但做工的人很少.
然後祂說:所以怎樣?.

$^{641}$不是叫你馬上去.
我們要留神.
祂說:所以你們當求莊稼的主.
到差遣做工的人出去.
收他的莊稼.
你們記住.
所有莊稼都是誰的?.
都是主的.
所以是按主的心意去領受差派去做他的工作.
有時候我們說在傳福音.
做神的工作中.
有時候我們界線拿捏得不好.
就會製造很多不咬牙的事件.
為什麼?.
有時候會以為莊稼是我的.
不是的.
莊稼是主的.
是否我喜歡收那些莊稼就收那些?.
我喜歡收那些莊稼就收那些?.
也不是的.
原來主要我們收哪些.
我們就去收哪些.
莊稼是否我們選擇喜歡收哪些就收哪些?.
不是的.
所以有時候我們說.
每一個宣教士都很奇妙.
上帝要猜那個去某個國家.
就是收那些.
那個就是收那些.
不能說對比.
神啊為什麼不猜我去那裡.
為什麼要猜我去這裡.
我也覺得很奇妙.
因為有一次我去泰國短孫.
我和一對泰國的家庭夫婦去說.
原來他們是澳門人.
他們的經歷是什麼呢?.
原來他們家人都是賭錢的.
但那時候有宣教士去家訪他們.
很順利.

$^{681}$那次家訪他們很聽宣教士的話.
然後父親就不賭錢了.
真的很乖地重量改過.
然後他們就信主了.
他們都問.
我們都覺得很特別.
神又不是要我們留在澳門做宣教.
要我們去泰國.
所以有時候我們要敏銳.
上帝叫我們做什麼.
原來我們知道的時候.
我們就不會有那麼多埋怨.
那麼多生氣.
因為是神叫的.
那神會負責.
如果只是說那些是我決定的.
我想怎樣就怎樣.
有時候可能會有一些跌宕.
會有紛爭.
因為你當了是我.
但不是的.
耶穌說.
要求莊家的主.
所以莊家是主的.
但法宮人猜誰去哪裡.
大家去回應就行了.
有時候去某些地方可以收很多莊家.
但你知不知道去某些地方.
可能一生一個年日裡.
收很少莊家.
但都要去做.
就是忠心的擺相.
所以求主幫助我們.
明白多些神的話語.
原來上帝要我們去收莊家.
我們要尋問.
上帝你想我們收哪些莊家.
你讓我們看見什麼.
求你讓我們有憐憫的心腸.
以致我們做得來.

$^{721}$不需要和其他人.
其他教會去比較.
其實大中小型教會.
大家全部都是主內一家.
所以我們收莊家.
不是說競爭.
你收多些你收少些.
不是.
原來收主的莊家.
忠心做在上帝叫你做.
你做了.
上帝就.
如何?.
閱立.
上帝就知道.
所以我也很感恩.
神留下一些聖經很好.
他說流淚殺腫的.
都歡呼收割.
那個歡呼收割不一定是.
每個人都收.
可以是其他人收.
但一起去歡喜快樂.
我很想和大家看一些.
我們的照片.
因為是上年開始通關.
但通關的時候.
我小朋友上學了.
不能這樣走來走去.
所以我們上年到暑假的時候.
我帶我女兒一起去體驗.
看看還有沒有東西可以做.
支援.
上帝想我們怎樣做.
藍色衣服的就是我女兒.
我和她親子裝.
戴眼鏡.
我們去哪裡做呢?.
在公園.
看看能否接觸到小朋友.

$^{761}$你去到澳門一些公屋村居屋村.
他們叫醒屋和驚屋.
他們的特色是.
他們依然會有小朋友.
自己下來玩.
他們不擔心.
為什麼呢?.
因為澳門有一個很好的設施.
稱之為天眼.
所以其實大家都很放心.
我初時去到不明白.
為什麼小朋友這麼放心.
這麼敢下來玩.
不怕有風化案.
不怕被人拐帶.
不怕那些事嗎?.
怕什麼?.
到處都覆蓋著.
但原來還可以做.
像八九十年代.
很多小朋友都會外出玩.
父母沒有空理會.
我們就很有感動.
想接觸這些小朋友.
和他們聊天和玩.
他們未必有父母可以陪伴他們.
我們去陪伴.
我都很感恩.
雷雷是我一個很好的精兵和助手.
她專心和其他小朋友玩.
玩完之後.
她會叫小朋友過去.
找我們大人聊天.
你一定是有關係.
然後就可以靠近.
我們是很好的合作.
然後到了十二月.
我們邀請街坊一起參與.
你看到有小朋友和大人.
我們一起參與.

$^{801}$這次我丈夫也有一起.
我們借用了一間餐廳.
在那裡可以做聖餐主日.
特別的地方是.
因為是餐廳.
聚會後.
大家拿一份茶餐回家吃.
很有人情味.
不同地方有不同的文化.
這是茶餐.
很開心.
有很多.
我女兒也不害羞.
到處看到街坊就派.
祝福你.
聖誕節快樂.
耶穌愛你.
這樣說.
所以很感恩我們一起配搭.
然後到今年.
繼續陪伴他們.
這些小朋友都是宣教時.
一直在街頭接觸.
玩耍.
所以可以叫他們一起回教會.
一起參與.
到了新年.
我們又再聚集多些人.
有時我們也會覺得.
好像小貓幾隻小貓幾隻這樣做.
也不肯定上帝有什麼心意給我們.
但原來做澳門人.
我不知道你們知不知道.
澳門很流行一句話.
叫「識人好過識字」.
澳門人是很講求關係的.
如果你跟他沒有上下關係.
他不會理你.
不會理睬你.
澳門人也不容易被你打入他的圈子.

$^{841}$所以到了這次.
做了一輪聚會的時候.
小朋友或家長終於說了.
「下次是什麼時候?」.
你就知道原來有些東西真的建立了.
有個關係可以後續.
所以我們說做科工工作.
千萬不要看短時間.
原來上星期我們在社區裡.
我們在裡面生活.
其實也不是短時間.
我們用每個星期.
每個月 每年.
這樣建立建立建立.
然後他真的可以歸屬.
跟我們一起了.
特別這個家庭很想提一提.
因為很想大家一起祈禱.
大家有沒有看到中間那幅.
藍色衣服的是我的女兒.
旁邊有四個小朋友.
然後有個格子衣服.
另外一個家庭的小朋友.
其實中間那四個小朋友.
是來自同一個家庭.
有大女兒媽生.
然後有第三個女兒.
然後有第四個兒子.
那次聚會的時候.
他們有小朋友.
他們自己製作.
我是帶家長的.
聊天.
聊天的時候.
四個孩子的媽媽哭了.
我說「你有什麼心事嗎?」.
「有什麼可以跟大家分享嗎?」.
誰知道她就哭了.
然後她說原來.
她先生發現末期.

$^{881}$肝癌.
她就想「怎麼辦?」.
四個小朋友.
怎麼辦?.
她也很徬徨.
所以聽完她說.
我們已經手上沒有錢了.
所以你看到我們下面.
就跟他們去吃東西.
我說「我跟你們一起吃東西吧」.
她也表示她覺得.
做媽媽很愧疚.
丈夫有病.
又沒有錢.
媽媽也要照顧小朋友.
又沒有錢.
她說已經很久沒有跟他們.
去快餐店吃東西了.
很久沒有買東西給他們了.
我自己也做媽媽.
也很明白.
我說「不如我們一起去吃東西吧」.
「一起去試試吧」.
其實我昨天也去完澳門.
昨天去了.
再了解那個消息情況.
那個丈夫已經離世了.
他們收到一些消息.
可能要他們搬遷.
撤離.
但很難過.
怎麼可以這麼快做這些事呢?.
所以也有.
社工局.
社工正在幫他們.
希望可以扯求情信.
幫他們可以留下.
原來上帝很奇妙地.
讓我們接觸到他們.
知道了他們的情況.

$^{921}$然後陪他們過.
希望讓他們點滴經歷到.
原來神真的會找他們.
神真的會幫他們.
上帝會出手.
他們不會孤單.
其實也有很多不同的教會.
持續到澳門服事.
有時我們去服事.
會碰面,會遇到.
颳颱風我們也會去.
風雨同路.
原來.
農夫會否颳颱風就不用做了?.
颳颱風就更快收.
否則損失慘重.
原來我們做科學工作.
風雨不改.
求主給我們.
這個信念.
這個要求.
這些就是.
近期的.
上個月.
我們接觸到的.
所以盼望在靈裡.
給予刺激.
一起做科學工作.
神給我的份和給大家的份不同.
給我的份.
每逢週六.
我便會去澳門.
每次我都要用.
五,六小時來回.
但我可能只做幾個小時.
我也有問過上帝.
這樣會否.
香港很講效能.
是否很適合.
這個比率呢?.

$^{961}$但上帝只告訴我.
收他的倉價.
從來不是講效能.
是講忠心.
我也會想.
其實我可以來回也很好.
上帝讓我想起以前的.
先賢.
他們是有去.
沒有回頭.
我可以來回.
還有命回來.
已經很蒙恩.
之前來傳福音的人.
他們就是預了.
一來就不能回.
所以求主幫助我們.
一起敏銳上帝的心意.
求主去差遣我們.
我們也求問神.
神要我們做什麼呢?.
我們一起為主去多作主公.
我們一起同心同道.
賜給人類的天賦.
因為你愛我們每一個.
你揀選我們每一個.
你用你的愛和真理.
栽培我們每一個.
幫助我們知道.
我們的生命.
有你自己奇妙的作為.
你也要我們成為.
其他人的祝福.
主求你幫助我們.
去敏銳你的心意.
求問你的心意.
讓我們的餘生好好被你使用.
好好做好你要我們做的.
忠心去做.
以致將來我們見主面的時候.

$^{1001}$我們不枉此生.
我們無悔.
能夠果子累累地去見你.
主啊 求你.
去閱立我們所獻情的一切.
我們誠心禱告交託.
奉主明求 阿們.
\newpage



\section{以弗所書 4:20-24}
\label{sec:sujYDAwGYXA}
\textbf{2024.08.11 《帶著耶穌》以弗所書4\_20-24 (和修版) 曾裔貴牧師}
\newline
\newline
連結: \href{https://youtube.com/watch?v=sujYDAwGYXA}{\texttt{https://youtube.com/watch?v=sujYDAwGYXA}} ~~~~ 語音日期: 2024-08-12
\newline
\newline
\hyperref[sec:1eBxYFSJlEI]{\small{< < < PREV SERMON < < <}}
~
\hyperref[sec:index_chronic]{\small{[返順時目]}}
~
\hyperref[sec:index_scriptual]{\small{[返順卷目]}}
~
\hyperref[sec:_zoyMoXK3Ek]{\small{> > > NEXT SERMON > > >}}
\newline
\newline
以弗所書 4:20-24
\newline
\begin{longtable}{cl}
\hline
\hline
章節 & 經文 (和合本修訂版)\\
\hline
4:20 & \begin{tabularx}{0.7\textwidth}{X} 但你們從基督學的不是這樣。 \end{tabularx} \\ \\ \relax
4:21 & \begin{tabularx}{0.7\textwidth}{X} 如果你們聽過他的道，領了他的教，因為真理就在耶穌裡， \end{tabularx} \\ \\ \relax
4:22 & \begin{tabularx}{0.7\textwidth}{X} 你們要脫去從前的行為，脫去舊我；這舊我是因私慾的迷惑而漸漸敗壞的。 \end{tabularx} \\ \\ \relax
4:23 & \begin{tabularx}{0.7\textwidth}{X} 你們要把自己的心志更新， \end{tabularx} \\ \\ \relax
4:24 & \begin{tabularx}{0.7\textwidth}{X} 並且穿上新我；這新我是照著神的形像造的，有從真理來的公義和聖潔。 \end{tabularx} \\ \\
[1ex]
\hline
\hline
\end{longtable}
$^{1}$我們一起同心祈禱.
再一次感恩天父讓我們聚集.
我們有敬拜的時間.
亦有剛才我們聖餐紀念主十家大愛的時間.
有我們在這段時間聆聽主的真道.
多謝主讓我們有脫去舊人.
穿上新人的教導.
我們很盼望在香港.
在這個時代.
我們來學習.
試圖補羅他的教訓.
二千年前對已不所成的人.
基督徒的教導和勸勉.
使我們能夠在香港.
成為讓人看到.
我們是屬主的兒女.
新生命,生活的流露.
我們一起等候在你的面前.
奉耶穌基督的命.
阿們.
很感恩.
因為我們作為牧者.
我們不單教初信班.
帶領到信主.
我們在準備洗禮之前.
我們要來跟每一位準備接受洗禮的.
我們要查問信德.
了解他們的內心信.
在上個星期.
跟唐慧,曾姑娘.
一起去接見每一位準備洗禮的.
其中一位弟兄.
他在座一定知道.
他說我信了耶穌.
我不知道應該怎樣在生活裡.
作見證.
他說很容易就會發脾氣.
很容易舊有的性情就會出現.
可能是跟工作.
跟工作地方的人士的關係.

$^{41}$他說不知怎麼辦.
突然間其中一位提到.
你帶著耶穌上班.
這就想起.
對耶穌和我一起.
我有力量.
就想起.
這段聖經很重要.
跟大家一起彼此勉勵.
相信我們一點也不陌生.
已不鎖書有三章聖經.
頭三章聖經.
講到關於信了主耶穌的神學.
在這個救恩裡.
得到天上屬靈的福氣.
是很詳盡地講述.
但第四章.
講到我們的行事為人.
又相當詳盡仔細.
教我們怎樣做基督徒.
又教得很仔細.
雖然我們今天選取的是四章二十至二十四節.
但其實延伸至五章第一段.
都是講到我們怎樣脫去舊人.
首先怎樣脫去呢.
就是你一定要意識到.
那些是舊的.
是神所不喜歡的.
是要脫開的.
你要有自己那種意識.
看到這些是神不喜歡的.
例如二十五節我們沒有引述的.
要氣絕謊言.
氣絕.
不是慢慢講少一點.
某些場合應該講.
有些場合不要講.
白色的謊言就不怕講.
沒有這件事.
是氣絕.

$^{81}$氣絕謊言.
還有各人語淪陷.
是要講真實的說話.
謊言和真實的說話.
是一對孿生兄弟.
很肯定.
信了主耶穌.
你要認識神一定不喜歡講謊言.
你要脫去.
但想想我們有很多機會.
要我們講謊言的場合.
做了基督徒.
你要意識到這樣的危機.
神是不喜悅的.
你可以想像到.
還有其他很實際的教導.
當時的已乏所人.
講謊言當吃生菜.
很多人.
大家理如我詐.
大家不誠實.
保羅這樣教導.
不是應該講一些很崇高.
很深奧的神學.
才是聖經.
但保羅不是.
他看到基督徒.
很實在地活在地上.
地上的生活.
正正就是你能不能活出.
耶穌基督的新生命.
你能不能夠在你的生活裡.
在你的說話裡.
讓人看到新生命.
所以每一個教導背後.
都是保羅意識到.
這些舊有的性情.
舊有的舊人.
是會阻礙我們穿上新人.
來成為在地上.

$^{121}$神的兒女的代表.
剛剛上星期一.
我有機會和同學.
上去深圳.
要買一些零件.
去一間餐廳吃東西.
那間餐廳我也算很熟悉.
因為經常去.
職員見到我.
今天這麼早上來.
很開心地打招呼.
吃東西之後.
突然有一個伯伯.
在香港上來.
他拿了一個外賣的盤子.
那個盤子應該是餐廳的盤子.
很生氣地質問那些職員.
為什麼上次我買外賣.
買牛腩.
一塊牛腩都沒有.
他很生氣.
我在香港上來.
我已經八十歲了.
不斷地罵.
不單止這樣罵.
你們大陸人.
做的全部都是假的.
地溝油.
把很多的觀念都帶上去.
但我作為那個餐廳的熟客.
不是的.
那些職員很好相遇.
很好的.
接著他再說.
我們大家都忍不住笑.
那些職員都忍不住笑.
本來已經很生氣.
大家在吵架.
他說明明你那張相.
宣傳的那張相.

$^{161}$整個盤子是充滿牛腩的.
為什麼沒有牛腩.
原來是那張相.
我就忍不住笑了.
我真的要出去解圍了.
我說爸爸.
我也是從香港上來的.
我跟這些職員很熟.
其實人家不是這樣的.
我說很小事.
入我帳.
叫職員給他一盤.
入我帳.
在我的帳戶扣錢.
因為我有帳戶在餐廳.
那個爸爸不肯.
小事.
一定要換一盤牛腩給我.
我說不關你的事.
這樣.
我就想到一件事.
香港人在這些職員面前.
其實很誤會.
那個爸爸是看著那張宣傳相.
拍得很漂亮.
滿滿牛腩.
他要有這樣的陣容.
才叫做要對板.
原來是一個這樣的誤會.
而他帶著這麼舊有的執著.
令到餐廳很尷尬.
開始慢慢大家越來越火滾.
大家彼此在叫.
對方又罵香港人.
很多誤會.
帶來.
我那時候內心很感動.
真是內心有感動.
要洗人和睦.
正正是餐廳的職員和我很熟.

$^{201}$我說算了.
爸爸這麼大年紀.
他可能真的誤會.
其實我們大家忍不住笑.
不如算了.
入我帳.
職員又不肯.
不關你的事.
最終因為真的吵得很厲害.
那些職員很晦氣.
舀了很多牛腩給他.
送走了.
這樣的圓滿的解決.
兄姐妹.
在我們的生活裡.
其實充斥著這些人和人之間.
少少的嫌隙.
這些可以引發.
可以令大家很不開心.
本來是很開心的一件事.
你這麼遠上到深圳.
吃東西.
你喜歡吃什麼就吃什麼.
為什麼為了這麼少的東西.
雖然兩塊錢坐車.
但也要一個小時.
你住得這麼遠.
八十歲.
你也要為了一盤牛腩.
罵得很厲害.
其實真的不關職員的事.
兄姐妹.
我們有沒有帶著耶穌.
在生活裡.
來流露出我們那份.
你留意一下.
保羅他怎樣形容這個新人.
穿上得好呢.
第24節.
他說穿上了新人之後.

$^{241}$這個新人是照著神的形象去做的.
有真理的仁義和聖潔.
穿上了這個新的形象.
這個新的形象.
就是說.
神是照著自己的形象.
樣式做人.
就是《創世記》第一章的24節.
神就做阿當.
成為人類的第一個人.
第一個人在犯罪之前.
是有神的形象的.
原來《創世記》第五章.
在講洪水之前.
神透過摩西再寫加普的時候.
阿當是照著神的形象去做.
第五章第一節.
阿當的後代記載下面.
當神做人的日子.
留意.
是照著自己的樣式去做的.
並且做男做女.
在他們被做的日子.
神賜福給他們.
稱他們為人.
就是說與神做的所有的萬物.
都是不同的.
唯有人有神的形象樣式.
未墮落之前.
你可以想像到.
這是第五章告訴我們.
重述創世的故事.
因為加普是後來再詳盡.
記載到第七代.
特別要記載的就是以洛.
你可以想像到.
還未完.
第三節.
當他活到130歲.
生了一個兒子.

$^{281}$然後聖經如何描述呢?.
形象樣式與自己相似.
就幫他改名為「失特」.
形象樣式出現了三次.
第四次就是保羅說.
保羅在神的啟示下.
他知道任何一個基督徒.
就是新的阿當.
你就有神的形象樣式.
你就能夠穿上神的性情.
神的形象.
當然不是一個站著的人.
這位造物主.
不是一個有身體的物體.
他是靈體.
但是神的性情.
我們一信了主耶穌.
就擁有了真理的人意和聖潔.
這就是穿上的重要.
所以保羅不斷教導我們.
我們要穿上之後.
更重要的是要帶著耶穌.
在我們的生活裡.
你看看五章已不鎖書提到.
或者我們不看五章.
四章這裡.
第三十節開始.
他說不要叫神的聖靈去擔憂.
你們原是受了祂的孕寄.
等候得贖的日子來到.
這裡告訴我們.
由信了主耶穌之後.
我們就有了聖靈.
在我們的心靈裡成為了一個孕寄.
聖靈就在我們裡面.
這個孕寄就是第一章提及過的.
就好像一個字.
是一個訂金.
是神.
在我們信主的時候.

$^{321}$給了我們訂金.
就是神已經在地上.
知道誰是贖祂的兒女.
已經買贖了.
所以我們有了這個孕寄.
就是這個訂金.
我們將來就可以得著救贖之後.
就有永恆的基業.
這是一章的經文教導.
現在第四章再一次提及.
有了這個孕寄之後.
會有一個聖靈的工作在我們裡面.
就是第30節說.
不要叫神的聖靈去擔憂.
你們原原本本是受了祂的孕寄.
這個孕寄使得我們將來得贖.
這個是憑據.
這個憑據就是我們有聖靈.
所以這裡告訴我們.
等候得贖的日子.
就是神再回來.
和我們完成地上的責任.
離開這個世界.
這兩個的日子.
到底是主耶穌回來.
還是我們離開這個世界.
就是得贖的日子的意思.
就是說還沒有得贖之前.
我們就要面對地上人生.
每一個生活裡面.
我們需要去小心.
因為第31節再一次教我們.
聖靈的擔憂.
多直接和我們的言語有關係.
一切的苦毒,怒恨,憤怒,洋流,毀謗.
並一切的惡毒.
都當從你們中間除掉.
全部是負面的.
全部是由內心的苦毒.
導致言談之間.

$^{361}$對人有一種很大的刻薄.
甚至是毀謗.
不存在的都要去詆毀它.
要置它於死地的這樣的一個心態.
你可以想像到.
聖靈的擔憂.
就在我們的言語之間.
在內心出現.
所以第32節.
當我們意識到.
我們不應該再這樣了.
第32節.
剛剛這個對比.
一個有聖靈來充滿的人.
並要以恩慈相待.
傳一個憐憫的心.
彼此的饒恕.
因為人與人之間總會有誤會.
總會有少許的一些嫌隙.
彼此的饒恕.
正如神在基督裡面.
饒恕我們一樣.
要饒恕人.
要處處從一個同理心的對別人.
看別人.
總是可以接納的.
對自己當然要認真.
對自己要律己以嚴.
所以這裡就是一個.
保羅很細緻的教導.
如何去除舊.
如何去穿上.
以及如何可以帶領著主耶穌.
在我們實際的生活裡面.
其實每個人都知道.
如果我們不去面對.
我們只能夠讓聖靈擔憂.
我們不能夠反映神的形象.
最近一直找一個例子.
我覺得很詫異.

$^{401}$有一個.
在網上如果你打這個詞彙.
4美元.
即是4元美金的教授.
就出現了很多關於他的事跡.
這個原來是一個哈佛大學的教授.
有4個學位.
其中兩個是博士學位.
美國的哈佛大學.
法律學博士.
經濟學博士.
4個學位.
而且是立即成為哈佛的教授.
當然是副教授.
由助理教授慢慢變成副教授.
副教授之後.
如果你的教學.
你的研究.
足夠成為終身教授.
就非常顯赫了.
為什麼他叫4美元教授呢?.
如果你很快地搜尋.
立即就找到他了.
原來在2014年12月5日.
在美國的麻省.
這個教授叫外賣.
在一間四川的餐廳.
叫了外賣之後.
在網上訂購.
然後就去到餐廳拿外賣.
他就發現.
因為他是法律出身.
看事物看得很精準.
為什麼會多計算了4美元呢?.
回到家之後覺得不服氣.
他就打了一封很長的電郵.
給那個老闆.
為什麼你多收了4美元呢?.
明明網上的價錢.
和實際見到的帳單.

$^{441}$就是去到現場拿外賣.
為什麼會多了4美元呢?.
質問老闆.
老闆整天很忙.
到晚上看到一封電郵.
發現有些投訴.
立即回答解釋.
原來你不明白了.
因為我們人手不足.
網上的資料還未更新.
所以我們在現場加價.
每一份餐點.
他點了4個菜.
每一份菜加了1美元.
其實是網上的資料還未更新.
這個教授好像發現了一個金礦.
立即運用他的法律知識.
再寫一封電郵.
根據馬賽諸夏州的法例.
犯了欺詐消費者的罪.
是嚴重的指控.
現在有辦法.
你要賠3倍給消費者.
即是他.
4×3=12.
還要扣我4元.
前後要給我16元.
老闆解釋.
其實我們一直收緊這個價錢.
真是不好意思.
是誤會.
其實是未更新網上的資料.
他完全不擔心.
他找到一個機會.
可以來據理力爭.
他的理由是.
我是網絡的警察.
我是為民除害.
你不知道已經騙過多少人.
這樣.

$^{481}$記住在哈佛大學.
當時只是26歲的教授.
哈佛大學是很顯赫的大學.
在世界的排名很高.
一直不肯放過餐廳老闆.
還要去政府的當局投訴這件事.
老闆說既然你已經投訴.
就等政府決定.
如果要我賠多少就賠多少.
就這樣交著了.
這個教授很意氣風發.
將這個案例.
成為一個給學生做的功課.
你看看.
我運用法律的知識.
據理力爭.
令這些小商人.
不可以再影響消費者.
叫同學對這個案例.
如勢加珍的描述之後.
叫學生來做功課.
得到的結論.
你知不知道是怎樣.
那些哈佛大學的學生.
沒有一個人看這件事是這樣.
覺得是為何一個教授.
會咄咄逼人對一個小商戶.
作出這樣無情的控訴.
然後那些學生.
你將電郵那些東西.
給了我們作為教材.
那些學生將這一切.
全部放上網.
然後波士頓日報.
將這件事拿出來.
這個教授叫做四美元教授.
在2014年的12月.
是全網上最討厭的人.
是真有其事.
還未止.

$^{521}$去到2015年.
哈佛大學要決定.
因為這個教授叫艾德蒙.
他的中文翻譯.
就申請終身教席.
哈佛大學破天方破例.
要觀察兩年.
兩年後2017年.
因為從來未試過.
要由所有教授投票.
他可不可以成為終身教席.
相信大家都已經有答案.
有三分之二覺得這個教授.
因為他的學生.
不斷去搜尋網上.
原來他不是第一次.
已經找到不同的案例.
去攻擊那些小商店.
可想而知一個.
學觀這麼豐富的人.
他的知識.
不單止不是去幫助人伸張正義.
反而是要令自己.
得到最好的好處.
所以很多學生.
發起一個運動.
四美元運動.
大家奉獻四美金.
幫助食物銀行.
去使那些有經濟困難的.
低下階層.
美國麻省裡面的.
可以得到救濟.
很多很多的錢.
來去收到.
因為哈佛大學的學生.
覺得這個教授.
影衰了我們.
我們要去跟他劃清界線.
只是因為四美金.

$^{561}$親愛的弟兄姊妹.
我們有沒有帶著耶穌.
在我們的生活裡.
我們有沒有在我們的生活裡.
別人看到我們是基督徒.
我們是有神的形象.
樣式.
記住.
阿當真是神做的唯一一個.
但他的兒子 瑟特.
是有父親的形象 樣式.
即是說原來這件事.
就算墮落了.
也沒有神收回.
現在到了新約.
保羅還可以很清楚地告訴我們.
每一個基督徒.
都有神的形象 樣式.
而且有聖靈在我們大家的內心.
來帶領我們.
所以我們千萬不要讓聖靈擔憂.
今早.
因為下午我們有七位弟兄姊妹.
其中有三位是泰國的弟兄姊妹.
我們待會在舞堂.
也有七位弟兄姊妹.
十四位弟兄姊妹.
成為新人.
下水.
埋葬.
上水.
就是新做的人.
就是有神的形象 樣式.
很盼望.
今天我們大家去觀禮.
去為他們喝彩 感恩.
更加要記住.
我們每一個弟兄姊妹.
都有神的形象 樣式.
我們一起禱告.

$^{601}$多謝主.
我們真的看到.
我們是屬於你.
是多麼寶貴榮耀的事.
本來我們是屬於舊有的世界.
舊有的人.
我們沾染的是很多.
不同世界上的人的思想.
價值觀.
都是神所不喜歡的.
我們很盼望.
我們能夠在主的愛.
聖靈的幫助.
感動和啟迪.
我們可以來到去.
穿上新人.
而且是不斷的.
在我們的生活.
帶著主耶穌.
在生活裡.
以致我們有力量.
我們有一個新的價值觀.
去看事物.
奉主耶穌基督的名.
阿們.
\newpage



\section{約翰福音 1:14-18}
\label{sec:_zoyMoXK3Ek}
\textbf{2024.08.18 《上帝的說話》約翰福音1\_14-18 (和修版) 陳韋安牧師}
\newline
\newline
連結: \href{https://youtube.com/watch?v=-zoyMoXK3Ek}{\texttt{https://youtube.com/watch?v=-zoyMoXK3Ek}} ~~~~ 語音日期: 2024-08-19
\newline
\newline
\hyperref[sec:sujYDAwGYXA]{\small{< < < PREV SERMON < < <}}
~
\hyperref[sec:index_chronic]{\small{[返順時目]}}
~
\hyperref[sec:index_scriptual]{\small{[返順卷目]}}
~
\hyperref[sec:QmBSplE5ueA]{\small{> > > NEXT SERMON > > >}}
\newline
\newline
約翰福音 1:14-18
\newline
\begin{longtable}{cl}
\hline
\hline
章節 & 經文 (和合本修訂版)\\
\hline
1:14 & \begin{tabularx}{0.7\textwidth}{X} 道成了肉身，住在我們中間，充充滿滿地有恩典有真理，我們也見過他的榮光，正是父獨一兒子的榮光。 \end{tabularx} \\ \\ \relax
1:15 & \begin{tabularx}{0.7\textwidth}{X} 約翰為他作見證，喊著說：「這就是我曾說：『那在我以後來的先於我，因為在我以前，他已經存在。』」 \end{tabularx} \\ \\ \relax
1:16 & \begin{tabularx}{0.7\textwidth}{X} 從他的豐富裡，我們都領受了恩典，而且恩上加恩。 \end{tabularx} \\ \\ \relax
1:17 & \begin{tabularx}{0.7\textwidth}{X} 律法是藉著摩西頒佈的；恩典和真理卻是由耶穌基督來的。 \end{tabularx} \\ \\ \relax
1:18 & \begin{tabularx}{0.7\textwidth}{X} 從來沒有人見過神，只有在父懷裡獨一的兒子將他表明出來。 \end{tabularx} \\ \\
[1ex]
\hline
\hline
\end{longtable}
$^{1}$各位兄弟姐妹,午安!.
我們一起來聽上帝的話.
今天的講題是上帝的話.
我們一起來認識和思想上帝的話.
講道本身就是上帝的話.
上帝透過講道在聖經裡對我們每位弟兄姊妹來說話.
不知道大家喜不喜歡說話.
有些弟兄姊妹喜歡說話.
我自己就一般般喜歡說話.
我也是比較內向性格的人.
我用MBTI性格測試.
我叫INTJ.
比較喜歡聆聽和沉默.
說話都是偏向簡潔.
一針見血準確.
說完就完.
所以講道對我來說是一個很享受的事情.
因為能夠慢慢精心炮製一段說話.
和大家分享.
當然作為一個INTJ.
我想大家也一樣.
作為一個成年人.
大家長大後.
縱然不是什麼說話高手.
大家大概都掌握了說話的重點.
學會何時打開話題.
何時說幾句簡單的chitchat.
何時不適合說什麼話題.
大家也知道.
在這個社會裡.
說話其實很重要.
也是無可避免.
我們長大後.
說話的機會隨之而來也越來越多.
說話不單單是一個很純粹的事情.
說話成為了一種工作.
功能,禮貌,要求,方式,職業.
所以我想大家在社會裡都知道.
說話基本上是一件很必然的事情.
我自己到了中年的時候.

$^{41}$也發覺到時不想說也不說.
所以也明白公園伯伯的心情.
基本上沉默不說話就不說話.
當然說話也不一定是說話.
有時候我們文字上的說話.
也是一種方式.
大家用手機和別人溝通的時候.
很多時候都不是說出來.
而是在手機打出來.
現在不同了.
20年前我們想打電話給別人.
就會打電話給別人.
現在想打電話給別人會怎樣.
會先發訊息給別人.
約時間問對方是否接受才打給他.
所以無論是文字,打電話,留言.
其實都是一個目的.
去傳遞我們本身想說的訊息.
說話是一種思想的表達.
一種交流,訊息的傳遞.
不過不知道大家有沒有同感.
在這個年頭.
說話和語言是越來越卑賤.
不知如何說.
特別是我們現在看到.
這兩年裡面出現了AI的時候.
基本上你一按一個按鈕.
字和文字就不斷地出現.
所以發覺在一個這麼訊息澎湃.
文字廉價,說話失真的年代裡.
我們越來越多快餐式的交流.
或者商業的說話套路.
或者是誇大其詞,言過其實的說話.
我們有很多聽到的說話.
但真正能夠聽到的可能是很少.
今天我們回到最基本.
我們思想上帝,我們的神.
祂自己的說話.
我們透過聖經來認識上帝的說話.
究竟我們如何和上帝的說話相處.

$^{81}$去回應上帝的說話.
我祈禱.
天父多謝你幫我們透過崇拜來敬拜.
我們更加在聖道的時間裡.
去恭敬,聆聽你的說話.
求主你親自對我們說話.
透過耶穌基督的彰顯.
聖靈當中的傳遞.
我們能夠真正地認識.
用我們的生命來回應你的說話.
求主你來帶領我們.
奉主名求.
阿們.
究竟什麼是上帝的說話呢?.
不知道大家有沒有聽過神跟你說話?.
有沒有?.
應該有的吧?.
讀了那麼久應該都有的吧?.
你覺得怎樣說?.
在什麼情境裡聽到神跟你說話?.
神是說廣東話的嗎?.
是不是?.
也是的吧?.
還是一個黑人穿西裝.
說英文的樣子.
還是一陣風吹過.
突然有聲音出現.
然後突然醒神的樣子.
大家應該有試過吧?.
有沒有試過很迷茫的時候.
突然想問神的心意.
翻一本聖經.
中了一頁就看看那一頁.
有點像問米.
打開看看神跟我們說什麼.
所以我們很多時候都會嘗試.
聽或體會神跟我們說話.
今天看經文正正是說這段經文.
是說上帝的說話是一回事.
可以讀到嗎?.

$^{121}$這段是《讓福音》第一章.
十四到十八節經文.
我們一起讀吧.
準備,一二三.
道成肉肉生,住在我們中間.
充充滿滿地有恩典有真理.
我們也見過他的榮光.
正是父獨身兒子的螢光.
約翰為他作見證,喊著說.
這就是我曾說.
那在我以後來的先於我.
因為在我以前他已經存在.
從他豐富裡我們都領受了恩典.
並且欣賞加恩.
律法是藉著兩西斑部的.
恩典和真理卻是由耶穌基督來的.
從來沒有人見過神.
只有在父懷裡獨身兒子將他表明出來.
在《讓福音》第一章裡.
一開始的時候稱之為神學引言.
簡單說說這段經文.
如果你熟悉的話.
你會發現《讓福音》跟其他三段福音書有點不同.
《讓福音》用了一個很有趣.
很獨特的神學論述.
來介紹耶穌基督的起初.
太初有道,對神同在,多謝神.
即是太初與神同在.
其他的福音書並非如此.
其他的福音書.
例如馬太老家是用聖誕節.
即是馬槽降生來介紹耶穌基督.
馬可是以耶穌的出道來介紹耶穌基督.
《讓福音》很特別.
是用道,一個很哲學.
很深奧的道理.
道來介紹耶穌基督的起初.
什麼叫道呢?.
道的希臘文就是Logos.
他說耶穌基督就是上帝的道.

$^{161}$道與神同在,道就是上帝.
Logos可以有很多不同的意思.
可以解作語言,說話,理性.
所以是Word of God,上帝自己的說話.
上帝今天不說,但總之就是說.
道成肉身,住在我們中間.
充充滿滿地有恩典有真理.
這是對耶穌基督的來臨.
一個很有趣精彩的描述.
不是聖誕節馬槽降生的故事.
而是充滿了哲學味道的道成肉身的道理.
什麼叫道成肉身呢?.
可能大家也聽過,也經常說.
道成肉身就是道成肉身.
但我們知道道成肉身並不是一個動詞.
不是一個人落手落腳就道成肉身.
其實這個字是一個動詞.
是說上帝的道成為肉身.
所以道字不是一個動詞.
不是多人向善,而是道字.
而是說上帝的說話成為了身體.
所以道成肉身的意思是說話成為了身體.
上帝的說話成為了一個具體的人.
一個具體的肉體.
所以道成肉身住在我們中間.
充充滿滿地有恩典有真理.
上帝的說話成為了人.
這個說話就是耶穌基督自己.
這是永恆信仰最重要的信息.
上帝的說話成為了一個人.
我們可以去認識和思考這個課題.
耶穌基督是上帝對我們的說話.
這是我們很重要的事情.
上帝有說話對我們說.
並且藉著耶穌基督來對我們說.
祂自己就是上帝說話的本身.
所以當我們說的時候.
這個說話就推翻了我們對於說話的理解.
你會發覺原來上帝的說話是另一種層次.
舉個例子.

$^{201}$我家裡有隻貓.
我家裡的貓養了兩年.
基本上牠這兩年都沒有笑過.
都是很生氣的.
扁鼻貓很肥胖的.
很生氣的.
整天看著窗口.
牠偶爾都會生我的氣.
偶爾會咪一聲.
早上想吃東西的時候.
就會靠近我咪一聲.
有時候想玩的時候會咪一聲.
有時候想吃零食的時候.
就會過來咪我一聲.
咪一聲好像有很多意思.
咪一聲可以解釋很多東西.
解釋吃東西 玩耍 很多不同的意思.
牠都是咪一聲就說出來了.
其實我也不明白牠明不明白我說什麼.
牠的名字叫做喵仔.
有時候叫喵仔牠也不太理會我.
我也聽不懂牠明不明白我說什麼.
這兩年來我也不知道大家明不明白對方說什麼.
牠也不知道我說什麼.
不知道喵仔說什麼.
不過其實貓的語言應該和人的語言不同.
貓的語言應該很簡單.
是不同的當次.
人的語言應該是更加深層次 更加複雜.
所以同樣道理.
上帝的說話應該不是和我們人的說話是同一個當次.
當然神可以藉著廣東話跟你說話.
或者用福建話跟你說話.
但是人的說話是透過空氣的震盪.
聽到聲音.
但是你知道神說話其實不需要這樣.
不需要透過人的方式來說話.
當然上帝可以留一段文字給你.
或者說一段說話給你聽.
但是最重要最首要的是什麼.

$^{241}$上帝的說話成為了肉身.
上帝要和我們說話.
並且透過一個很具體很立體的方式.
去和我們說話.
這句話就是耶穌基督自己.
所以這正正就是道綽心的意思.
上帝要對我們說話.
這句話成為了肉身.
處在我們中間.
有恩典 有真理.
所以等於如果耶穌基督自己就是上帝的說話的時候.
我們對上帝的說話有完全不同的理解.
上帝要對我們說話不單單是透過一段文字.
或者是留言 或者是一段聲音.
上帝要我們上帝的說話成為了一個生命的時候.
我們要用我們的生命來好好地回應這句話.
祂將祂的生命放在我們心中.
我們要全程用整個的生命來聆聽.
去回應上帝的說話.
這正正就是很重要的意思.
尤其是今天我們出現了很多AI的時候.
要整一段文字出來不容易.
整一段有意思的也不難.
但是真真正正能夠去改變生命.
它是一個不同的檔次的東西.
我們要用我們的生命來回應上帝的說話.
神也是用祂的生命來對我們說話.
所以我們要用我們的生命來捉緊祂.
去回應祂 去聽從祂.
所以經文裡面最後一句話是這麼說.
從來沒有人見過神 只有在父華裡的多生子 將他顯明出來.
這個正正就是耶穌基督自己.
唯有耶穌是真真正正上帝的說話.
我們要用整個生命和耶穌建立關係.
來聆聽上帝的說話.
既然我們知道耶穌是上帝的說話的時候.
我們怎樣來跟上帝的說話相處呢.
怎樣去處理上帝的說話呢.
我們看看下面的經文.
第十六和十七節.

$^{281}$經文說.
從他豐滿的恩典裡 我們都領受了.
並且因上加恩.
律法本是接著摩西傳的.
恩典和真理都是由耶穌基督來的.
所以說恩典和真理都是由耶穌基督來的.
經文說.
耶穌基督作為上帝的說話有兩個很不同的特點.
第一個是上帝的說話是真理.
第二是上帝的說話是恩典.
所以上帝的說話有這兩個很重要的特性.
今天我們不說真理了.
沒有時間.
我單單來說恩典這個課題.
上帝的說話是恩典.
這是我們今天一起來思想的.
很有趣的.
我告訴大家一個聖經的冷知識.
可能你會覺得很震驚.
你會發現原來整卷的樣方裡面.
恩典這個字是很稀有的字.
有沒有發現.
整卷的樣方出現了兩次恩典這個字.
就在這裡.
恩典是一個很普遍的字.
但整卷的樣方裡面.
只有這裡出現了恩典這個字.
對於樣方來說.
恩典這個字不是一個隨隨便便說出來的字.
不是一個慣常.
好東西.
正的意思.
他說什麼.
從他豐滿的恩典裡我們都領受了.
並且因上加恩.
律法本是至多不失傳遞.
恩典和真理都是由基督耶穌那裡來的.
用風說.
上帝的說話就是恩典.
這個恩典不是一件隨便偶然的事情.

$^{321}$你不可以take it for granted.
不可以隨隨便便當是必然的事情.
不知道大家是不是這樣想.
我們太多人將上帝的說話take it for granted.
覺得這件事反而是有的.
手機裡面有一個app.
是不是一個.
不知道大家有沒有發現.
我們越來越方便用手機的時候.
越來越隨便.
手機有聖經的時候.
會不會越來越少看.
以前我們背著聖經在背包裡.
整本大短書.
今天不需要.
手機裡面有一個app叫聖經.
但是越方便大家就越少看聖經.
打開openrise多一點.
打開聖經多一點.
基本上聖經這個app是一個工具.
有時候要查.
要找就用.
但是拿來看反而很少.
我們太過將上帝的說話看成是必然的事情.
反正都是一個app.
download來十秒鐘就download完.
但我們越來越忽視.
去珍而重之的.
將上帝的說話好好重視.
反正是一個其中一個app.
會是.
楊豐說上帝的說話是恩典.
楊豐強調.
這個恩典更加是什麼.
是恩上加恩.
可以這麼說嗎.
是一個恩上加恩的恩典.
什麼是恩上加恩呢.
意思是.
我們可能認為恩上加恩.

$^{361}$都是很多很多的恩典的意思.
有點像李答李.
上節.
或者喜上加喜的朋友.
有點像是很多很多的恩典的意思.
但其實.
原文裡面的意思是很特別的.
恩上加恩的原文.
如果翻譯成英文.
字面上的意思.
是Grace Anti-Grace.
用了anti這個字.
Grace Anti-Grace.
意思是.
Anti 意思是很多的意思.
是取代.
或者代替的意思.
所以這個恩典是一個.
Grace Anti-Grace.
Grace Instead of Grace.
恩典取代恩典.
什麼是恩典取代恩典呢.
我告訴你.
上帝的恩典永遠都是新的.
永遠都是不斷的更新.
去代替.
去替換 不斷的更新.
就像你上網看影片.
現在是串流的.
不斷的有一個streaming.
你不停的上網.
不斷的更新.
去更新.
這個就是這個方式.
你不能夠下載.
你不能夠offline去看.
你只能夠不斷的去欣賞佳音.
不斷的去更新.
新的恩典會取代舊的恩典.
現在的恩典會取代過去的恩典.

$^{401}$將來的恩典又會成為新的恩典.
我們每天清晨.
都在神裡面得著新的恩典.
你不能夠保留舊的恩典.
只能夠不斷的和神裡面的關係.
在神裡面去明白.
和得到更多的恩典.
任何人要抓住舊的恩典去奔放.
他就變成了拒絕了基督耶穌的豐盛.
所以這個欣賞佳音.
告訴我們恩典的本質.
恩典不是一些你能夠累積下來的東西.
不是你儲積分.
這樣儲起來的.
而是不斷在基督裡面.
去支取 去取得上帝新的恩典.
這個正是欣賞佳音的意思.
不是加很多.
而是不斷的更新 更換.
不斷的去代替.
這是恩典的本質.
因此 既然這是恩典的本質的時候.
我就明白上帝的說話是恩典的意思.
上帝的說話是恩典.
所以上帝的說話是不停的和你更新.
在這個關係裡面.
你用你的生命不斷的和上帝接觸.
上帝不斷的和你說話.
你不斷的去聆聽.
不斷的去更新 不斷的去連接.
然後不斷的再次和上帝交流 去認識.
然後被上帝恩典去豐盛 去壓倒.
這個正是恩典的特性.
所以最後用摩西的律法來做對比.
他說 律法本是藉著摩西傳的.
恩典和真理都是從耶穌基督來的.
沒錯的 寫在法版上的律法是容易的.
因為是有板有眼.
白紙一紙寫清楚 是不是.
但是上帝的說話是在這個之上.

$^{441}$他不是要推翻律法.
不過說律法以上.
更加要不斷的在關係裡面.
去聆聽 去聽到上帝的說話.
上帝的說話永遠都是不停的和你有關係.
你只能不斷的聆聽 留心去查看上帝的說話.
簡單的做個總結.
我們說幾點.
第一個是什麼.
耶穌基督就是上帝的說話.
這個是用分辨最重要的信息告訴你.
這個說話就是耶穌自己.
第二 上帝的說話是真理.
今天我們沒有說的.
不過大家都明白的.
第三 上帝的說話是恩典.
所以我們不能夠take it for granted.
要好好地去珍惜它.
好好地承載它.
第四 恩上加恩.
上帝的說話是要求我們不斷在基督裡面.
去根深蒂固 用我們的生命來承載它.
我們簡單的來說一些應用.
既然上帝的說話是恩典的時候.
有幾點可以跟大家想一想.
第一 我們說話的本質本來就是恩典.
如果上帝的說話是恩典的時候.
我們作為基督徒.
無論你是外向也好 內向也好.
我們願意和別人說話.
本來就是一件恩典的事情.
我們因為愛的緣故.
因為神的緣故.
我們願意開口和別人說話.
你每一個願意打開口和別人說話的對象.
無論是7-11的職員.
還是你的父母 還是你的兒子都好.
當我們願意和別人開口說話的時候.
這都是帶著上帝的恩典來和他們說話.
因為你願意愛他們.

$^{481}$願意和他們結合關係.
我們就這樣和他們說話.
無論你多麼內向也好.
我們都願意作為基督徒.
都願意以說話來成為和別人結合關係的方法.
當然留心我們的說話.
因為這是一個恩典的出口.
也是一個愛的表達.
所以我覺得沒有所謂說話的藝術.
任何一個說話的藝術.
單單有技巧而沒有那種愛的時候.
這只是表面的功夫.
無論你的說話多麼口賤拔拙.
多麼不會說話也好.
帶著這份恩典和別人說話.
是我們基督徒說話的一個很重要的本意.
好好想想你自己的說話.
我們每一次和別人說話的時候.
會不會有太多套路 公式化的東西.
還是真真正正和別人有著恩典的關係呢.
第二 我們要想的是聖經作為一個上帝的說話.
耶穌基督是上帝說話的本身.
聖經正是很重要被寫成的上帝說話.
去見證著耶穌基督自己.
不知道你們怎樣理解聖經.
我們從宣道會 由錦書到現在.
聖經都可以的.
但不及曾牧師那麼好.
因為他的宗派非常厲害.
背誦聖經是很厲害的.
背誦聖經和金句是非常厲害的.
重點不是要大家背誦熟多少.
或者能夠很熟地轉回來背誦.
而是好好地將這句說話放在你心中.
我想沒有將你心中結連的一句說話.
哪怕成為了很熟悉的金句也好.
這句也是一個沒有意思的說話.
唯有你自己的心命和它有交流.
用你的心命去實踐出來的神話.
才是一個有意思的說話.

$^{521}$真日你放在手機裡不打開來看.
基本上和其他APP沒有分別.
都是一個普通的手機APP.
當你用你的生命 真義眾知來打開來讀.
哪怕只是一節經文.
當你用你的生命投入進去.
這篇經文就正正是上帝的說話.
一個有意義有力量的上帝說話.
一本放在書架裡的聖經是沒有意思的.
一本放在手機裡的APP也沒有意思的.
唯有你自己用生命投入的說話才有意思.
聖經裡有很多不同的理解.
例如放在酒店裡的一本書.
當我們基督徒打開來閱讀的時候.
用生命投入去回應的時候.
才是一個有意義的上帝說話.
所謂學武神的說話.
未必一定要你熟悉到經文裡的說話.
懂得背誦或讀某些神學課程.
而是你把上帝的說話看為上帝的說話.
認為能夠跟你生命有關係的神話才是重要的.
真人 靈修也一樣.
如果沒有用生命去連結的話.
只不過是讀理解.
一篇文章 讀小學的時候也讀理解.
看完之後就做了.
靈修的文字其實也是一些文字.
唯有你用生命投入的時候.
將神的話語珍惜地回應的時候.
這才是真真正正上帝的說話.
最後一點就是生命的說話.
唯有將我們的生命參與的時候.
我們才能夠來到最正確的回應神的說話.
用你的生命很好的方法來投入.
透過這句話來聆聽 敬拜 求敏.
我自己在生命中有幾個階段.
很明顯知道神跟我說話.
例如我夢照的時候.
我讀書的時候.
或者是開新的教會的時候.

$^{561}$都很明顯知道神跟我說話.
但神怎麼跟我說話呢?.
沒有一段經文叫你開教會.
沒有一段經文叫你去德國讀書.
但當你用你的生命去參與的時候.
你透過你自己的投入的時候.
神透過不同的方式跟你說話.
可能是一段經文.
可能是弟兄姊妹的聊天.
可能是牧者的提醒.
當我們願意去真而重之.
將上帝的說話看為上帝的說話的時候.
神就肯定會跟我們說話.
今天我們教會.
試試在這個星期裡.
大家試試.
認真地將上帝的說話看為寶貴.
這句話可能很簡單.
但我可以用粵語來形容.
將上帝的說話.
看為跟你生命有關係的東西.
不是純粹課程內容.
或者是諸侯的教導.
而是跟你生命有關係的時候.
當你願意這樣去理解的時候.
神必定會對你說話.
因為上帝仍然在說話.
縱然今天我們不容易見到神的時候.
但上帝仍然在說話.
只要我們願意將上帝的說話看為寶貴的時候.
這是非常重要的.
不妨大家試試.
在這個星期裡.
你面對一些情況的時候.
當你願意相信神跟你說話的時候.
上帝必定會對你說話.
我們一起祈禱.
耶穌基督.
我們知道你是那位.
上帝到誠肉身的上帝.

$^{601}$我們能夠透過你認識上帝.
能夠聆聽上帝對我們生命的命令.
我們求你讓我們再次將你的說話看為寶貴.
我們不單要背多少金句.
或者讀多少次聖經.
這些都重要.
但最重要的是我們將你的說話.
當作上帝的說話.
一個跟我們生命有關係的說話.
求你這樣幫助我們.
也是鼓勵我們.
在這個星期裡.
我們能夠聽見你自己的說話.
你的說話我們願意去跟從.
願意去實踐.
我們求你在我們生命的環境裡.
無論是我們的工作.
我們的生活處境.
我們生活的難題.
將你的說話成為我們很重要的一部分.
我們在這個情境裡.
來聆聽你的說話.
求主你這樣對我們說話.
求你這樣幫助我們.
我們每一個基督徒.
我們都珍惜來聆聽.
因為你的說話是恩上加恩.
不斷地來更新.
不斷地來跟我們說話.
求主你這樣的幫助.
逢主明求.
阿們.
\newpage



\section{馬可福音 2:1-12}
\label{sec:QmBSplE5ueA}
\textbf{2024.09.01 《行出看得見的信心》 馬可福音2\_1-12 (和修版) 講員 : 陳燕雯傳道}
\newline
\newline
連結: \href{https://youtube.com/watch?v=QmBSplE5ueA}{\texttt{https://youtube.com/watch?v=QmBSplE5ueA}} ~~~~ 語音日期: 2024-09-02
\newline
\newline
\hyperref[sec:_zoyMoXK3Ek]{\small{< < < PREV SERMON < < <}}
~
\hyperref[sec:index_chronic]{\small{[返順時目]}}
~
\hyperref[sec:index_scriptual]{\small{[返順卷目]}}
~
\hyperref[sec:XLqWIi0boGI]{\small{> > > NEXT SERMON > > >}}
\newline
\newline
馬可福音 2:1-12
\newline
\begin{longtable}{cl}
\hline
\hline
章節 & 經文 (和合本修訂版)\\
\hline
2:1 & \begin{tabularx}{0.7\textwidth}{X} 過了些日子，耶穌又進了迦百農。人聽說他在屋裡， \end{tabularx} \\ \\ \relax
2:2 & \begin{tabularx}{0.7\textwidth}{X} 於是許多人聚集，甚至連門前都沒有空地；耶穌就對他們講道。 \end{tabularx} \\ \\ \relax
2:3 & \begin{tabularx}{0.7\textwidth}{X} 有人帶著一個癱子來見耶穌，是由四個人抬來的； \end{tabularx} \\ \\ \relax
2:4 & \begin{tabularx}{0.7\textwidth}{X} 因為人多，無法抬到耶穌跟前，就把他所在那房子的屋頂拆了，既拆通了，就把癱子連所躺臥的褥子都縋下去。 \end{tabularx} \\ \\ \relax
2:5 & \begin{tabularx}{0.7\textwidth}{X} 耶穌見他們的信心，就對癱子說：「孩子，你的罪赦了。」 \end{tabularx} \\ \\ \relax
2:6 & \begin{tabularx}{0.7\textwidth}{X} 有幾個文士坐在那裡，心裡議論，說： \end{tabularx} \\ \\ \relax
2:7 & \begin{tabularx}{0.7\textwidth}{X} 「這個人為甚麼這樣說呢？他說褻瀆的話了。除了神一位之外，誰能赦罪呢？」 \end{tabularx} \\ \\ \relax
2:8 & \begin{tabularx}{0.7\textwidth}{X} 耶穌心中立刻知道他們心裡這樣議論，就說：「你們心裡為甚麼這樣議論呢？ \end{tabularx} \\ \\ \relax
2:9 & \begin{tabularx}{0.7\textwidth}{X} 對癱子說『你的罪赦了』，或說『起來！拿你的褥子行走』，哪一樣容易呢？ \end{tabularx} \\ \\ \relax
2:10 & \begin{tabularx}{0.7\textwidth}{X} 但要讓你們知道，人子在地上有赦罪的權柄。」就對癱子說： \end{tabularx} \\ \\ \relax
2:11 & \begin{tabularx}{0.7\textwidth}{X} 「我吩咐你，起來！拿你的褥子回家去吧。」 \end{tabularx} \\ \\ \relax
2:12 & \begin{tabularx}{0.7\textwidth}{X} 那人就起來，立刻拿著褥子，當著眾人面前出去了，以致眾人都驚奇，歸榮耀給神，說：「我們從來沒有見過這樣的事！」 \end{tabularx} \\ \\
[1ex]
\hline
\hline
\end{longtable}
$^{1}$(字幕提供:MG).
(字幕提供:鄭穎宵).
多謝鄭穎宵邀請我到鴻恩堂.
其實我也是元朗區的街坊.
所以我今早乘輕鐵就可以過來.
很開心來到當中.
原來我也有認識一些弟兄姊妹.
所以在這裡見到 譬如高全道.
感覺很親切.
在幾個月前有一個廣產片叫做.
《九龍城寨之圍城》.
因為這樣 九龍城寨就掀起熱潮.
成為城中熱話.
連在五十多六十年前.
來到九龍城寨傳福音的Jackie Pullinger.
人們稱她為潘小姐的宣教士.
都被人翻查她的資料.
潘小姐是英國人.
她在二十二歲的時候.
神感動她.
她有一個心很想去宣教.
她就祈禱 踏出她信心的一步.
結果藉著神的帶領來到香港.
我們今天的講題就是.
行出看得見的信心.
我們一起禱告.
主佑多謝你 多謝你帶領我們.
你和我們同行 你沒有撇棄我們.
求主你的靈和我們會眾同在.
幫助說的能夠說得清楚明白.
亦幫助聽的能夠不單聽到.
亦能行到.
求你當中看顧保守.
你得著一切榮耀頌讚.
我們這樣祈禱感恩.
是奉主耶穌基督的名祈求 阿們.
講題是行出看得見的信心.
經文是《馬學福音》一至十二節.
《馬學福音》是四本福音書中.
最短的一本.

$^{41}$《馬學福音》的寫作風格很快.
很多時會一浪接一浪.
一件事件緊接另一件事件記載.
所以在福音書裡.
弟兄姊妹會看到.
立刻 除夕 這些字眼.
在第一節已經開宗明義地說.
神的兒子耶穌基督 福音的起頭.
在第一章時.
不單記載了施洗約翰.
預備主的道路.
耶穌基督專講謝罪的福音.
亦記載了耶穌基督在加利利.
到會堂教導.
甚至做了很多神蹟.
醫治了很多不同的人.
他這樣做 令到他的明星.
在當時加利利一帶聲名大噪.
傳到加利利周圍的人都聽到.
來到第二章.
馬可交代了背景.
第二章 經文裡說.
過了些日子 耶穌又進了加拔隆.
又進了加拔隆.
即是說他之前去過.
其實是 第一章第21節提及到.
耶穌和門徒到加拔隆.
耶穌在安息日的時候.
去到會堂教訓人.
大家對他的教訓都感到驚訝.
因為覺得他的教訓是有權柄的.
不像解經學家那樣.
除此之外.
那天耶穌在會堂.
命令污鬼離開.
那些人看到耶穌的作為.
感到驚訝.
他們就因為這樣.
將耶穌的明星傳播到加拉利周圍.
所以當耶穌再次來到加拔隆的時候.

$^{81}$那些人聽到耶穌在屋內.
馬上有很多人聚集.
甚至連門前的空地都沒有.
耶穌就跟他們講道.
有四個人.
抬著個毯子來見耶穌.
但由於人多.
他們沒有辦法將毯子抬到耶穌面前.
他們怎麼做呢.
他們居然在耶穌頭頂上的屋頂上.
拆了屋頂.
將毯子連人帶的床褥.
聚到耶穌面前.
聖經說.
耶穌看見他們的信心.
就對毯子說.
孩子,你的罪赦了.
經文說.
耶穌看見他們的信心.
原來信心是可以被看見的嗎.
信心不是虛無飄渺的東西.
難怪雅各在雅各書中這樣說.
我藉著我的行為.
將我的信心給你看.
信心和行為是相輔並行的.
而且信心是因著行為才得以成全的.
如果信心沒有行為.
就是死的.
耶穌看見他們的信心.
這四個朋友從哪裡來的信心.
大到可以連屋頂都拆得開.
他們的信心又是基於什麼呢.
剛才提到耶穌去過加伯隆.
我猜那四個朋友.
一定在加伯隆聽過耶穌的道.
不單這樣.
他們可能親眼聽過.
耶穌所行的神蹟奇事.
治好不同的人.
這樣我就相信.

$^{121}$他們有這樣的信心.
可以做出這樣的行為.
是因為他們知道耶穌是誰.
他們相信耶穌有能力治好他們的朋友.
所以他們做出這樣的行動.
但我們要小心.
我們的信心不要放在錯的地方.
如果信心投放在錯的地方.
可以很危險.
甚至會沒命.
不知道大家看見這條片.
是否記得這段新聞呢.
這是去年六月的時候.
有五個人.
他們將信心投放在泰坦號潛水器裡.
這個去觀光鐵達尼號的潛水艇.
結果離難失去生命.
弟兄姊妹.
我們的信心是投放在哪裡呢.
我們的信心又是基於什麼呢.
在希伯來書十一章裡.
提到很多不同的信心偉人.
其中提到拉俠的信心.
他說:寄女拉俠恩著信.
接待探子.
就沒有與那些不順從的人一同滅亡.
拉俠的信心是基於什麼呢.
不如我們看看約書亞記的第二章.
那裡記載了約書亞暗中派了兩個探子.
去窺探那個地方.
特別是去看耶利哥城.
那兩個探子去到拉俠家裡的時候.
拉俠的弟弟劉叔.
耶利哥王就派人去到拉俠那裡.
要他交出那兩個探子.
拉俠把那兩個人藏在他的屋頂.
為什麼拉俠這麼有信心.
藏這兩個探子在他的屋頂呢.
難道他吃了豹子膽.
他不怕耶利哥王對付他嗎.

$^{161}$我們聽聽拉俠對那兩個探子的說話.
在約書亞記的第二章9-11節.
拉俠對那兩個探子說.
我知道耶和華已經把這地刺給你們了.
你們使我們十分害怕.
這地所有的居民都在你們面前膽戰心驚.
因為我們聽見你們在埃及出來的時候.
耶和華怎樣使紅海的水在你們面前乾了.
以及你們怎樣對待約旦河東阿摩利人的兩個王.
西弘和鄂.
把他們完全消滅.
我們聽見了就都心裡驚怕.
沒有一個人有勇氣在你們面前站立.
因為耶和華你們的神是天上也是地上的神.
原來拉俠不怕耶利哥王對付他的信心.
是基於他認識耶和華.
他知道耶和華把這個流奶與蜜之地.
刺給以色列人.
他有聽見過耶和華怎樣使紅海分開.
叫埃及人在前面走過乾地.
拉俠認識耶和華.
他甚至說耶和華是你們的神.
是一個天上地上的神.
拉俠不怕耶利哥王對付.
是他對耶和華有信心.
他認識耶和華是一個怎樣的神.
剛才提到傑克布林吉.
潘小姐知道神要她出去.
雖然她不知道去哪裡.
但因為她認識她所信的神.
對她有信心.
知道主耶穌愛她.
而且主耶穌是一個全知全能的神.
主耶穌不會點一條黑路給她走.
所以她很放心和平安踏上信心的一步.
我相信潘小姐未來到香港前.
她不知道香港在哪裡.
但她知道她所信的是誰.
所以她就上船來到香港.
弟兄姊妹.

$^{201}$我們的信心是基於甚麼呢.
是我們的知識,才幹,人脈,經驗.
還是我們對主耶穌的認識呢.
行出漢德建的信心.
我們要認識我們的主耶穌.
我們要認識我們的信心對象.
你所認識的神是怎樣的神呢.
我們對主耶穌的認識有多少呢.
你對主耶穌的認識能否讓你有信心踏出你的安殊區.
你對主耶穌的認識能否順靠祂帶領你走前面的路.
耶穌能否告訴你我漢德建你的信心呢.
行出漢德建的信心.
行出信心祝福他人.
馬可沒有記載那四個朋友的名字.
對我們來說他們都是寂寂無名的.
可能你在教會裡也覺得自己是個小鼠.
但不要緊.
我可以肯定知道的是.
你微小的信心擺上.
你信心的行動.
耶和華主耶穌是會看得到.
就正如聖經所說.
記載耶穌是看到那四個朋友的信心.
耶穌看見他們的信心對坦子說.
孩子你的罪赦免了.
因為這四個朋友的信心行動.
這個坦子就可以被耶穌赦罪醫治.
我們看到耶穌和人之間的互動.
這個坦子很有福氣.
他有四個幫助他的朋友.
其實同樣這四個朋友也很有福.
因為他藉著這樣的做法.
他可以親眼看到.
見證到耶穌的大能和榮耀.
原來當我們有信心的行動出來的時候.
是可以祝福身邊其他的人.
而耶穌也可以用我們的信心行動.
去祝福身邊的人.
我想和大家分享一個真人真事.
這個故事發生在19世紀的時候的美國費城.

$^{241}$主角是一個小女孩.
名叫凱蒂·華爾特.
Hattie Mae Wyatt.
她去到一間教會叫Grace Bethtips Church.
她想回主約河.
但那個教會的地方很小.
她不能進去.
所以她就在教會的門外等.
當教會牧師看到凱蒂的時候.
凱蒂就和牧師哭訴.
「我很想回主約河.
但我進不去,很多人,沒有位子,很難過」.
牧師就親自帶她去教堂.
牧師就和她說.
「希望有一天我們的教堂可以大一點.
可以容納多一點的小朋友來這裡.
回主約河認識耶穌基督」.
她聚會後回到家.
她想到今天回到主約河很開心.
很感恩.
但同時她又想到.
很多小朋友都沒有這個機會去到教堂.
來認識主耶穌,敬拜神.
所以她心裡有個想法.
大約過了兩年之後.
凱蒂因為病離世.
之後發現凱蒂裡有個小錢包.
小錢包裡有些錢和一張字條.
那張字條寫著.
「這是獻給神,要把我們小小的教會建得大一點.
讓更多的小孩子能夠上主約河」.
錢包裡有五十七仙.
五號七仙.
很多教會的弟兄姊妹.
甚至其他人聽到凱蒂這個信心擺上的行動.
深受感動.
他們作出回應去奉獻.
最後不單止建立了更大的教會.
還有主約河的教室.
有學校和醫院.

$^{281}$而這間教會今天叫做.
天普浸信會.
The Baptist Temple.
凱蒂的心願其實是希望.
有更多人能夠回到教會.
認識主耶穌.
她憑著信心去存錢.
一心想奉獻給教會.
五十七仙.
神就用這五十七仙.
將她倍增去祝福她.
就好像那個小孩拿著五個餅.
兩條魚.
來到主耶穌面前.
主耶穌祝福.
倍增.
就能夠祝福當時聽耶穌說的.
單是男丁也有五千個.
我相信加上婦女和小孩.
是會過萬的.
洪水橋是香港現在.
其中一個重點的發展社區.
在2015年的時候.
洪水橋第一間公共屋村.
洪福村入伙.
而洪水橋的屯馬線站.
也在興建當中.
預計在2030年建成.
大概還有五六年左右.
到時我相信會有更多的人和家庭.
搬進洪水橋這個地區.
你們洪恩堂在2011年的時候.
在這裡成立.
來服侍社區.
作炎作光.
我相信神會繼續使用你們每一個.
盼望神就使用你們信心的行動.
無論是微小的團契小組裡.
或是社區服務中.
祝福弟兄姊妹和街坊.

$^{321}$神可以用我們信心的行動.
幫助,鼓勵和祝福弟兄姊妹和鄰舍.
除此之外.
其實在宣教層面.
每一個門徒都可以進行大使命.
神母呼召到每一個人都成為宣教士.
好像Jackie Pudding 潘小姐那樣.
但我們每一個信徒都能夠做宣教的事.
我教會有一個長者.
我非常敬佩他.
雖然他不能去什麼短宣因為年紀大.
但他對猜拳,對宣教很有熱誠.
非常關心.
他會仔細閱讀宣教士寄來的代禱信.
他會點名為宣教士祈禱.
也會為那些未得之盟來祈禱.
他不單在禱告擺上.
他也在金錢裡奉獻在宣教上的事工.
他是我一個很好的學習榜樣.
弟兄姊妹.
你有否將你的五十七仙或五餅二魚.
交給神來獻上,給神用呢?.
九龍城寨的潘小姐.
因為她信心的行動祝福了很多人.
數以千百計的人因為潘小姐的行動.
生命得到改變.
他們有這個救恩得到永生.
我們的信心行動是可以祝福身邊其他人.
耶穌基督可以使用我們當中在座的每一個.
不單這樣.
同時我們也蒙受神的祝福.
看到神的作為和見證祂的能力.
行出看得見的信心.
行出信心榮耀神.
馬可福音二章六至八節說.
當時有幾個文士.
即抄寫聖經的人坐在那裡.
心裡有異論說.
這個人為何這樣說呢?.
他說了些褻瀆的話.

$^{361}$除了神一個之外.
還有誰能夠赦罪?.
耶穌知道他們心裡在想甚麼.
就對他們說.
你們心裡為何這樣議論呢?.
對個癱子說你的罪赦免容易.
還是對他說拿起你的玉子行走容易呢?.
耶穌這樣問他們.
那些文士怎樣回答?.
他們沒有作任何回答.
當然耶穌也不期待他們回答誰容易.
因為無論是赦免癱子的罪.
還是醫治癱子.
叫癱子站起來走出去.
都是從神而來的力量和權柄.
耶穌基督的目的是甚麼呢?.
耶穌問他們的目的.
就是要他們知道人子.
在地上有赦罪的權柄.
《馬可福音》第一章.
已經記載了耶穌宣講赦罪的福音.
到第二章.
耶穌不單宣講赦罪的福音.
更加證明自己有赦罪權柄.
證明自己是神的兒子.
就像《馬可福音》第一章.
第一節已經開宗明義說.
神的兒子耶穌基督福音的起頭.
在十一節的經文.
耶穌用了三個動詞.
這三個動詞都是用命令的語氣.
就是「起來」「拿」和「去」.
耶穌吩咐癱子用命令語氣說.
叫他起來.
拿起你的玉子.
去 回家吧.
癱子聽到耶穌這樣說.
他又如何回應呢?.
他立即起來.
拿著箭玉.

$^{401}$在眾人面前走出去.
他相信耶穌的話.
立即行動.
癱子也有行出看得見的信心.
大團圓結局.
那四個朋友憑著信心拆了屋頂.
抬起癱子在耶穌面前.
期望耶穌治好他.
現在他也被治好了.
還在眾人面前離開.
故事可以在此完結.
不 馬可沒有在此完結.
他還記載了.
眾人都驚奇榮耀歸於神.
眾人說.
我從來沒有見過這樣的事.
行出信心榮耀神.
想不到癱子行出來的信心.
能夠榮耀神.
叫眾人讚美耶和華.
開始的時候.
耶穌看見那四個朋友的信心.
結尾的時候.
眾人看到耶穌的榮耀.
癱子可能沒有想到.
他這個信心行動帶來這麼大的震撼.
能夠讓人看到.
耶穌的能力和榮耀.
將榮耀歸於神.
在2015年2月的時候.
伊斯蘭國ISIS發佈了一個視頻.
顯示了21個被綁架的.
埃及基督徒集體被處決.
大家可能都記得新聞上.
看到這個景象.
但大家可能沒有看到的.
就是他們的信心帶來的影響.
令到多人將榮耀歸於神.
在這21個基督徒當中.
有兩個是親兄弟.

$^{441}$即是當日他們的媽媽.
同時失去了兩個兒子.
接下來和大家分享一個視頻.
視頻記載了媽媽的一個訪問.
她說到她很難過.
但她雖然難過.
她知道她的兩個兒子為主殉道.
她覺得很安心 很安慰.
他們的哥哥甚至說.
伊斯蘭國留給他們的.
實在過於他們所想.
他們留下這段視頻.
讓人看到那些殉難者.
在被斬頭的時候.
仍然供奉主耶穌的名.
令到我們的信心增加.
他們的媽媽訪問時被問到.
如果伊斯蘭國的人來到你家.
你會怎樣?.
她說:我會祈禱.
求神打開他們的心.
去認識主耶穌的信仰.
主持人問:你不怕.
他們將你殺了嗎?.
她說:我不怕.
我的兒子也不怕.
我還怕什麼?.
伊斯蘭國留給他們的.
實在過於他們所想.
他們的媽媽訪問時被問到.
如果伊斯蘭國的人來到你家.
你會怎樣?.
他們說:我會祈禱.
求神打開他們的心.
去認識主耶穌的信仰.
主持人問:你不怕.
他們說:我會祈禱.
我還怕什麼?.
他們說:我會祈禱.
我還怕什麼?.

$^{481}$他們說:我會祈禱.
我還怕什麼?.
他們說:我會祈禱.
我還怕什麼?.
他們說:我會祈禱.
我還怕什麼?.
他們說:我會祈禱.
我還怕什麼?.
他們說:我會祈禱.
我還怕什麼?.
他們說:我會祈禱.
我還怕什麼?.
他們說:我會祈禱.
我還怕什麼?.
他們說:我會祈禱.
我還怕什麼?.
他們說:我會祈禱.
我還怕什麼?.
他們說:我會祈禱.
我還怕什麼?.
他們說:我會祈禱.
我還怕什麼?.
他們說:我會祈禱.
我還怕什麼?.
他們說:我會祈禱.
我還怕什麼?.
他們說:我會祈禱.
我還怕什麼?.
他們說:我會祈禱.
我還怕什麼?.
他們說:我會祈禱.
我還怕什麼?.
他們說:我會祈禱.
我還怕什麼?.
他們說:我會祈禱.
我還怕什麼?.
他們說:我會祈禱.
我還怕什麼?.
他們說:我會祈禱.
我還怕什麼?.

$^{521}$他們說:我會祈禱.
我還怕什麼?.
他們說:我會祈禱.
我還怕什麼?.
他們說:我會祈禱.
我還怕什麼?.
他們說:我會祈禱.
我還怕什麼?.
他們說:我會祈禱.
我還怕什麼?.
他們說:我會祈禱.
我還怕什麼?.
他們說:我會祈禱.
我還怕什麼?.
他們說:我會祈禱.
我還怕什麼?.
他們說:我會祈禱.
我還怕什麼?.
他們說:我會祈禱.
我還怕什麼?.
他們說:我會祈禱.
我還怕什麼?.
他們說:我會祈禱.
我還怕什麼?.
他們說:我會祈禱.
我還怕什麼?.
他們說:我會祈禱.
我還怕什麼?.
他們說:我會祈禱.
我還怕什麼?.
或者可能.
走出去.
去到你的鄰舍.
去到你的親人.
去到你的朋友當中.
甚至去到你家裡的工人姐姐那裡.
又或者對一些人來說.
可能就是走出去.
走出你的安書區.
去到可能一個新的事奉崗位.

$^{561}$你們的領袖是什麼我不知道.
但我可以知道的就是.
走出來的信心.
能夠榮耀到神.
能夠祝福到人.
就讓我們去認識.
這個我們所信靠主耶穌基督.
這個信心的對象.
正如你們這裡說.
我們去認識我們的主耶穌.
使到主愛結連.
行出看得見的信心.
信行合一.
我們就有一段時間.
給大影子有一點時間.
去和神有一個回應.
一段時間的密討.
之後我會祈禱結束.
(音樂結束).
親愛的天父.
我們知道你是那位信實慈愛的主.
神讓我們祈禱.
求你帶領我們.
讓我們更加去認識.
我們這個生命的主.
以致到我們能夠將我們的生命.
擺上當作活祭.
呈獻給神.
因為這個是聖潔.
是你所喜悅.
我們這樣事奉是你所當然的.
使用我們成為別人的祝福.
使用我們叫到人去認識主耶穌基督.
將榮耀歸給主.
我們這樣祈禱感恩交託.
是奉主耶穌基督的名字祈求.
阿們.
\newpage



\section{傳道書 6:1-12}
\label{sec:XLqWIi0boGI}
\textbf{2024.08.25 《你想人生》 傳道書6\_1-12 (和修版) 講員 : 周曉暉牧師}
\newline
\newline
連結: \href{https://youtube.com/watch?v=XLqWIi0boGI}{\texttt{https://youtube.com/watch?v=XLqWIi0boGI}} ~~~~ 語音日期: 2024-09-04
\newline
\newline
\hyperref[sec:QmBSplE5ueA]{\small{< < < PREV SERMON < < <}}
~
\hyperref[sec:index_chronic]{\small{[返順時目]}}
~
\hyperref[sec:index_scriptual]{\small{[返順卷目]}}
~
\hyperref[sec:6BPEnUSuXaU]{\small{> > > NEXT SERMON > > >}}
\newline
\newline
傳道書 6:1-12
\newline
\begin{longtable}{cl}
\hline
\hline
章節 & 經文 (和合本修訂版)\\
\hline
6:1 & \begin{tabularx}{0.7\textwidth}{X} 我見日光之下有一件禍患重壓在人身上， \end{tabularx} \\ \\ \relax
6:2 & \begin{tabularx}{0.7\textwidth}{X} 就是人蒙神賜他財富、資產和尊榮，以致他心裡所願的一樣都不缺，只是神使他不能享用，反被外人享用。這是虛空，也是禍患。 \end{tabularx} \\ \\ \relax
6:3 & \begin{tabularx}{0.7\textwidth}{X} 人若生一百個兒子，活許多歲數；他即使壽命很長，心裡卻不因福樂而滿足，又不得埋葬；我說，那流掉的胎比他倒好。 \end{tabularx} \\ \\ \relax
6:4 & \begin{tabularx}{0.7\textwidth}{X} 因為這胎虛虛而來，暗暗而去，名字被黑暗遮蔽， \end{tabularx} \\ \\ \relax
6:5 & \begin{tabularx}{0.7\textwidth}{X} 而且沒有見過天日，甚麼都不知道，這胎比那人倒享安息。 \end{tabularx} \\ \\ \relax
6:6 & \begin{tabularx}{0.7\textwidth}{X} 那人雖然活千年，再活千年，卻不能享福；眾人豈不都歸同一個地方去嗎？ \end{tabularx} \\ \\ \relax
6:7 & \begin{tabularx}{0.7\textwidth}{X} 人的勞碌都為口腹，心裡卻不知足。 \end{tabularx} \\ \\ \relax
6:8 & \begin{tabularx}{0.7\textwidth}{X} 智慧人比愚昧人有甚麼益處呢？困苦人在眾人面前知道如何行，有甚麼益處呢？ \end{tabularx} \\ \\ \relax
6:9 & \begin{tabularx}{0.7\textwidth}{X} 眼睛所看的比心裡妄想的倒好。這也是虛空，也是捕風。 \end{tabularx} \\ \\ \relax
6:10 & \begin{tabularx}{0.7\textwidth}{X} 先前所有的，早已起了名，人早知道人是如何的，不能與比自己強壯的相爭。 \end{tabularx} \\ \\ \relax
6:11 & \begin{tabularx}{0.7\textwidth}{X} 話語多，虛空也增多，這對人有甚麼益處呢？ \end{tabularx} \\ \\ \relax
6:12 & \begin{tabularx}{0.7\textwidth}{X} 人一生虛度的日子，如影兒經過，誰知道甚麼才是對他有益呢？誰能告訴他身後在日光之下會發生甚麼事呢？ \end{tabularx} \\ \\
[1ex]
\hline
\hline
\end{longtable}
$^{1}$如果願望可以成真.
你想你的人生要些什麼呢?.
剛才很感恩我們聽到Mary姊妹的見證.
可能在她人生不同的年歲.
想要的東西都會不同.
可能在我們中間.
我們都會有很多東西可以想.
我給你十秒想想你想要些什麼.
這些?.
其實能夠心想事成,誰不想呢?.
剛才主席幫我們讀了聖經的一本書.
Mary姊妹就有美貌與智慧.
我只講智慧.
這本書是講智慧的.
傳道書.
傳道書講的智慧就是.
如果我們可以心想事成的人生.
是不是最好呢?.
在經文第九節說.
眼睛所看的比心裡妄想的都好.
這也是虛空,也是暴風.
聖經直接說出.
人心想事成的真相其實反轉就是心裡妄想.
如果真的看著看著就滿足就好了.
看著看著就好了.
但我們人偏偏多數不是這樣的.
有沒有試過去吃自助餐.
自助餐是怎樣吃的?.
排隊進去已經想好.
我一會兒要先去拿長腳蟹壽司.
吃長腳蟹壽司的時候.
就想著何時去拿燒牛肉.
你一邊吃那樣東西.
一邊想下一樣東西.
可能我們吃的時候.
沒有好好面對那件食品.
要學學麗兒.
怎樣尊重食品.
我們人心就是這一類.
想,我們想很多東西.

$^{41}$於是實際我們可能都覺得OK.
不過我們是不滿足.
我們很希望心想事成.
這個在我們新年也會互相祝賀.
心想事成不是不好.
不過是否就得到你想要的人生呢?.
傳道書就給了幾個關於心想事成的例子.
二至三至七至九節都有心裡這隻字.
重複的.
它告訴你.
有這一切你心裡想要的東西.
到底會有什麼情況出現.
其中第二節我們先看.
那裡說.
就是人夢神似他財富資產和尊榮.
以至他.
怎樣呢?大字一起讀.
心裡所願的.
一樣都不缺.
你要的東西都有.
生命裡很豐富.
從小到大.
你都是有滿足的.
聖經說有.
都要用得到才行.
有都要用得到才行.
你同不同意嗎?.
可能你也聽過這樣的例子.
有人在供一個全維多利亞海警的大單位.
有這麼大的房子.
自己清潔也很麻煩.
就請傭人姐姐幫忙.
不過因為這個人工作時間也很長.
基本上一個星期.
只有一天半天真的在家.
平日早上出晚歸工作.
有很多其他活動.
你猜誰每天都在享用全維多利亞海警?.
可能是姐姐吧.
收集回來給人用.

$^{81}$聖經很早就指出這個智慧.
所以從結果而言.
心想事成.
其實可以是一種.
不是真是福氣.
我想起我自己的祖父.
上世紀的五十年代.
帶著一堆金條來香港.
聽聞家裡有一塊土地.
是今天的跑馬地.
當時他六個小孩很小.
如果他能好好運用金條.
很容易養活一家八口.
綽綽有餘.
但他被朋友帶去.
想想這些.
試試這些.
玩玩這些.
搞搞那些.
結果全部家產漸漸被洗掉.
被人騙走.
賭錢.
他們怎樣賭啤牌呢?.
拿著啤牌.
攤得很開.
周圍的人.
前後都有人圍著他.
其實很容易打到龍通.
其他人就很容易騙走他的錢.
這些東西一樣都沒有留下.
本來看著那堆金條是很安心的.
但結果他是不是真的能用到.
到他很快的年日.
中年的時候.
心裡裝著很多東西.
變成完全空空的痛.
傳道書說第二個心想.
不是說錢.
說多少歲.
這裡說人若生一百個兒子.

$^{121}$活許多歲數.
他即使壽命很長.
心裡卻不因福樂而滿足.
其實我們中國人.
常常都渴望長壽.
一般人都認為長壽是福氣.
當然是好的.
問問大家.
世上哪個地區的人.
平均年齡是最高的?.
有些人知道.
有些人猜到.
以前有印象是日本的.
因為我們經常看到.
世界之最.
說日本有很多人歲.
一百一十幾一百二十歲.
的確.
他們人歲比較多.
但平均.
壽命最長的地方.
不是日本.
是香港.
是我們香港.
2021年發表全球調查.
發現香港當時是連續.
七年成為世界上.
人均年齡最高的地區.
當時的數字.
香港女性平均是89.6歲.
很長.
你看看你到90歲了沒有.
有很長時間.
弟兄不用擔心.
我們都有84.65歲.
是全世界所有國家地區裡.
最高的.
長壽.
開心.
現在數數.

$^{161}$差多少年了.
但你想想.
為什麼會是香港呢.
天水圍這裡不錯.
地方闊落.
我在港島.
住的地方環境狹窄.
一出街空氣污濁.
生活匆忙.
工作壓力很大.
香港人吃東西.
肉多菜少.
你去街不會點菜.
都是點肉.
所以我們都是吃三高食品.
雖然標榜三低.
但都是三高食品.
煎炸油炸.
濃味的食物.
我們竟然可以是.
世界人均年齡最高的地區.
如果對於生活富足.
沒有太多病痛.
能夠養天年.
退休.
可以這樣活這麼高壽.
其實都是一個好消息.
但長壽是否等於滿足呢.
不少上了年紀的朋友跟我說.
有時候都想不要太長命.
有時候長命大.
多病痛.
就比上班更忙.
一三五就去醫院覆診.
二四六就做物理治療.
你回到醫院.
診所時間都很多.
出去排一排.
如果是這樣的話.
我要長壽.

$^{201}$是否真的要覆診呢.
這裡有經文說.
那個人有100個兒子.
其實是說他有很多代.
表達他的生命可以延續很長.
才可以生這麼多兒子.
是否等於快樂呢.
有一個大的家庭.
是否就是快樂呢.
世界上被談論得最多的家庭.
你有沒有想起哪個家庭呢.
英女皇的家庭.
間中有新聞.
最近有兩兄弟的新聞.
特別是這五年.
風風雨雨.
兒孫輩都很多事.
前年九月.
女皇九十六歲高齡離世.
她離世前幾年.
她的孫子都很多新聞.
哈里王子夫婦.
威廉王子夫婦.
兩兄弟夫婦不和.
女皇的丈夫再早一年.
即2021年已經離世.
當這個菲納親王離世的時候.
傳媒就將他一生說過的.
有智慧的說話.
刊登了一些出來.
留意到其中一句.
她自己並不是很明白.
什麼叫做獨立自主.
因為她自己是全世界.
被管治得最多的人.
女皇身邊又管治她.
不止.
她說皇室有很多規條.
全部都是別人決定.
就連她的兒子改什麼名字.

$^{241}$她都沒有權.
要跟很多規則改.
這個丈夫離世之後.
皇室就送了幾隻新的小狗給英女皇陪她.
英女皇聽聞很喜歡一隻小狗.
叫做哥基狗.
很短腳的那些.
很希望陪她.
但怎知沒有幾個月.
其中一隻小狗就病死了.
很慘.
女皇接連受到這些打擊.
還有在那一兩年.
有些調查報道發現.
找出當年.
在十多二十年前.
其中BBC的記者做了一幾則假新聞.
令到當時她的前媳婦.
戴安娜王妃.
因為那則新聞變得疑神疑鬼.
間接造成.
其中可能走向離婚.
走向後來很多很多悲劇的原因.
女皇什麼都有.
龐大的皇室家庭.
但是不是等於滿足呢?.
有很多兒孫.
是不是等於滿足呢?.
聖經說.
再活千年.
卻不能夠享福.
這樣是不是好呢?.
牧師我做人.
都不求這些.
我對得起天地良心.
所以我很腳踏實地.
很平實生活.
傳道書再說第三個心想.
人一切勞碌都為口福.
可能你現在回到崇拜.

$^{281}$你正在想吃什麼.
我們都是為了吃.
我們的心始終都填滿.
這裡也說有自慰人.
有貧窮人.
但這裡說窮人.
如何行事為人.
的確貧窮不是一個最大的難處.
而是那個人如果是窮人.
行為端正.
可能生活上是好的.
香港上世紀六七十年代.
不少戰後一代其實很清楚.
正如Mary的見證一樣.
禮儀舉止重要.
有些家庭都教自己的孩子.
雖然我們沒有錢.
雖然窮.
但你都要穿戴整齊乾乾淨淨.
你看得起自己.
不要被人看扁.
不過人始終都會比較.
不知道大家有沒有一些重聚的聚會.
和小學同學中學同學.
或者什麼時候的舊同事.
每次出來.
總會聽到對方有些不同的發展.
消息.
這個升了什麼位.
那個又做了什麼.
他又生了多少個孫.
其實大家就開始很多的比較.
很多人在心裡面.
其實都會很執著.
都會去排.
都會不滿足.
聖經是不是叫我們安分守己.
放下執著.
其實其他哲學宗教都會這樣說.
但聖經還揭開一種實況.

$^{321}$就是我們人心裡面的想法.
這些慾望.
是讓我們看不見我們眼前已經有的.
所以你怎樣想你的人生都好.
最終你要留意一個智慧.
就是我們的心怎樣填都填不滿.
怎樣填都填不滿.
所以你現在在追逐什麼呢.
無論我們信主或不信主.
我們都可以想.
我們追逐的所求是什麼呢.
可能這一刻你說.
我只要有這個就行了.
我只要找到這個.
我去到那個位置.
我就快樂了.
到做到的時候.
你就發覺.
是嗎.
我只要去到那個地方.
我就平安.
我就安全.
我就沒事了.
是嗎.
其實這一生我們所追的.
是不是真的能夠讓我們滿足呢.
或許有些人說.
我的確需要眼前的一些成功.
牧師當然說我們.
不要這樣去追逐事業.
難道有人喜歡失敗嗎.
我當然不是這個意思.
其實聖經傳道書.
在前一兩章都有說到.
人生有四種成功的程度.
最低那裡.
排在最高.
在第四章有說到.
由最底層的被壓的一群人.
全沒有機會的受壓者.

$^{361}$去到上一點點.
能夠去努力一下.
可以上到一點點的位置.
然後有些真的能夠去.
得到財富.
得到名譽.
他不斷的勞碌是有成果的.
也都對最高級別.
當然他是有權有勢.
什麼都有.
呼風喚雨的領袖首富.
在這四種裡面.
大家自己排一下你是哪一種.
聖經反問人.
你排到最高.
成為最多.
最大.
最厲害.
是不是最好呢.
當然不是說成功不好.
是好的.
不過是不是對你一生.
這件事是最好呢.
我第一份工作是在政府裡面工作.
五月畢業.
六月馬上上班.
那個時候還是叫做鐵飯碗.
在那個年代.
政府工作都是一種體重年資的工作.
即是你安分守己.
好好地做.
基本上排著排著都可以升級的.
所以絕大部分不是說很成功.
不過總能夠慢慢爬上去.
最起碼比起同期的同學.
一開始就已經是物質充裕的一群.
我想想當年.
都快三十年了.
當年一畢業出來拿個人工.
比起現在大學生的人工.

$^{401}$竟然還是多他們一大截.
當然這一代有這一代的情況.
比起同期的同學.
好像好一點.
但這樣代表什麼呢.
說到最被羞欺壓的那一群.
其實人類的制度.
常常都有腐敗的.
結果呢.
最貧窮的人.
利益是得不到保障的.
因為的確是一層吃一層的.
不同的社會制度都是這樣的情況.
追求幾均分公平.
都仍然會出現的.
最高的當權者.
到最被羞欺壓的社會底層.
每一層其實貫穿的都是物質.
而那些物質其中一個來源.
傳道書說是地.
地就代表利益.
第五章其中一節這樣說.
況且地的益處歸眾人.
就是君王也受田地的供應.
其實那裡是說.
就算說到你最高級別的成功者.
君王領袖.
他都是被那塊土地困住.
為什麼這樣呢.
古時田地的出產代表財富.
種到多少東西或有多少收益.
但到今天仍然是這樣.
土地可以是核製人.
所以要明白為什麼.
在世上不同的領袖.
當權者.
會那麼介意土地的大小.
一分一吋都很緊張.
不斷思想如何令自己的土地.
有最大的收益.

$^{441}$有些拼命製造圍牆.
封住土地.
有些拼命建造軍隊.
防住敵人.
說穿了.
一直努力努力的.
都是混在一起的.
這些物質.
好像沒有帶給他們滿足.
反而他們還不斷為這些物質而活.
聖經提出了三個比較.
沿著剛才對田地的理解.
你就明白.
一個可能我們近一點.
我們不是做管治的人.
說普通我們人人都會有的.
我們有一些財產.
喜愛錢財的不得安睡.
哀財卻是不安.
擁有更多是不能夠滿足.
近年.
這一類的投資.
有沒有見過.
加密貨幣.
早一兩年.
曾經有個階段.
世界首富.
即是賣電動車那位.
他在網上發了一則推文.
我現在要買某某的加密貨幣.
於是就馬上一群人衝著跟著他.
這樣真的升了.
大家一起見到.
真的升得很厲害.
跟他買的人會怎樣.
買完之後會怎樣.
就天天看著.
不知道你有沒有投資的經驗.
你放五千元.
一萬元.

$^{481}$五萬元.
十萬元.
你放了去投資.
你會怎樣.
就天天看著.
最初一天看一次.
然後半天看一次.
然後每個小時看.
因為上落得很快.
在這個首富這樣說完之後.
很多人跟風之後.
突然間過了幾個月之後.
他又在上面發了一則貼文.
他說我賣了我所有的加密貨幣.
結果那個加密貨幣突然間大跌.
還引起當時一個小的.
他們不是叫股災.
叫幣災.
幣災.
如果你是跟著買的.
你還能睡得著嗎.
連你睡覺都拿走.
富有的人積聚財富.
同樣你賺了錢財去換物質.
這些物質不會帶給你滿足.
香港我們自己也有首富.
大家都認識他.
李生.
當然他投資很成功.
最近大家看他在海外怎樣投資.
有些說他將資金搬回大中華.
曾經何時想香港發達的人.
都會參考他的成功經驗.
會看看他怎樣工作或投資.
但同時我們有沒有參考他的人生.
他家裡也曾經出現一個大事.
就是他的大兒子曾經被綁架.
非常驚嚇.
對於這個成年人.
即是掌管他爸爸的企業的人.

$^{521}$對他一生有怎樣的影響.
其實聖經說得很對.
當積聚財富的時候.
很多時候可以是同時積聚災禍.
最近網上已經不流行看IG.
看叫threads.
其實是一些文字的小故事.
其中有一則小故事.
有一個香港工作人員.
說他最近怎樣接待一些.
從外國來的富有的人.
入住一個酒店.
那個人帶著四五個保鑣.
一個化妝師和廚師.
去住酒店.
然後怎樣招呼他.
他帶著幾個保鑣.
出入其實就在他身邊.
越多財產可能也會令你越驚恐.
所以就算你成為最成功的人又怎樣.
其實你有豐衣足食.
但你這樣的心想事成.
並沒有帶給他真正的滿足.
他是煩惱.
病痛怒氣可能依然在他生命當中不斷出現.
所以各位我們.
我們得到我們所想的.
會怎樣呢.
我想起自己的一個經歷.
瑪麗姐妹是最在行的.
拍照.
我曾經有一個階段.
就是攝影.
年輕的時候.
喜愛上攝影.
當時沒有現在的電話那麼方便.
拍張照片.
拍一些好的照片.
是要用那些單鏡反光機去拍的.
一喜歡了那個嗜好.

$^{561}$就經常想著要換相機.
有新相機就想買新相機.
新鏡頭就想買新鏡頭.
有了長焦的鏡頭.
又要找一個人像的鏡頭.
見到別人拍得很漂亮.
又要找一個微距的鏡頭.
很多東西玩的.
新婚那兩年.
我就和太太去旅行.
就會去一個地方.
知道那裡夏天的煙花很漂亮.
晚上放煙花.
滿天煙花.
當時我馬上上網.
去搜尋一下.
怎樣拍煙花的呢.
其中在討論區.
那個年代還沒有Facebook.
有人就說.
如果你要去拍煙花.
還要定的.
不過你的手很難定的.
一定要有一支腳架才行.
腳架還是不要普普通通的.
要很穩固的.
有機有鏡.
就要買腳架.
出發去機場那天.
我才發現.
為什麼有這麼多東西.
一大個相機背包.
還要一大把相機腳架.
很重.
那好吧.
滿心期待就去吧.
終於到了那一晚.
要去河邊看煙花.
去已經人多了.
背著一袋二袋又熱.

$^{601}$夏天.
勉強找到一個好的位置.
就要將背上的相機袋放下.
開始砌東西.
設好腳架.
然後整個相機放上去.
又試一下調光圈.
調一下角度.
整晚就在這裡拍煙花.
拍拍拍.
結果整晚煙花.
我都是透過相機的洞看.
我沒有真的用眼睛看.
拍完就要拆.
拆完又放回倉袋.
倉袋又要背回酒店.
最糟糕的是什麼.
從頭到尾都沒有聽到我提醒太太.
當然太太不會開心.
我完全被這個物質.
佔據了整個心思.
合計了.
一點都沒有享受到旅遊的樂趣.
拍照出來.
現在我都不敢拿出來看.
原本應該屬於兩夫婦很快樂的旅遊時光.
就沒有了.
很久很久之後.
有一次我又帶.
那時候三位孩子.
我又帶孩子一起.
再和太太去那個地方旅行.
又看煙花.
這次我聰明了.
我拋下相機.
全程用眼睛看.
各位聖經告訴我們.
我們心想很多東西.
以為我們有很多.
放在眼前.

$^{641}$背上身.
我們就快樂.
其實不是福.
我們得到你想想的意料.
其實是困住我們的人生.
你始終都說.
無論如何.
我都想試試心想事成.
各位我請你想想.
你認識的人.
如果他是一個能夠心想事成.
即是他想怎樣就怎樣.
想要做到就達到.
你覺得他會變好還是變壞.
世界是否會因為他而變得更好.
大概多數人會答變壞.
你答得很對.
聖經就是在裡面不斷告訴我們.
由第一本書《創世記》已經說.
人在地上的罪惡是罪大惡極.
心裡面終日所想的都是壞事.
新約也常常說.
我們人就是很容易被我們眼睛看到的東西.
以致我們放鬆自己.
放縱自己.
做很多其他事.
你會說我不是那些放縱私慾.
貪行污穢的人.
你想想.
我們現在在一個暢通無阻的網絡世界.
一個匿名隱蔽的環境.
沒有其他人知道你做甚麼看甚麼的時候.
你是怎樣.
人家不知道.
你知道.
聖經不斷提醒我們.
這就是世人的寫照.
告訴我們每個人都有罪過.
人總有做錯的時候.
我們既然得罪人.

$^{681}$亦可能得罪神.
更加我們是得罪自己.
所以傳道書常常說虛空.
虛空的意思其實是一陣風一口氣.
人生到頭可能只是這一口歎氣.
令到本來的日子好像虛度.
回望我們有多少人生真的過得好呢.
罪是會壓著我們的心.
沉重到無法釋懷.
以至你覺得.
以為心想事成.
隨心所欲就可以擺脫這些感受嗎.
是不可以的.
聖經這裡是說了.
你以為嗎?不是的.
一下過的時候.
其實你甚麼都不知道.
誰才知道真是好呢.
這裡說了.
神告訴你.
他就知道了.
正正就是因為我們人.
剛才說我們是有罪的.
罪就是令我們和永遠.
永遠和創造天地的上帝.
賜福下的上帝.
隔開了.
分隔開了.
完全隔離.
你記得COVID的日子嗎.
知道甚麼是隔離嗎.
一點都不見得.
罪就是令我們無法靠近這個.
永生的上帝.
或者生命的源頭.
而且這個鴻溝我們自己解決不了.
任何的財物都填不了.
任何的名成利就都無法幫助我們.
填這個缺口.
那可以怎樣呢.

$^{721}$上帝不是讓我們.
個個都成為好人.
聰明人.
能夠自我完善才幫我們.
他在我們還是一個.
都爛爛的人.
或者我們自己知道自己有多好.
他已經幫我們.
給了一個唯一的解決辦法.
他就差派他的兒子耶穌基督.
來幫我們從這個漩渦當中.
走出來.
我們人生其實很多時候都兜兜轉轉.
不是這個漩渦.
就是跳落另一個漩渦.
我們心所想的只是另一個漩渦.
但是神卻是.
派派了他的兒子耶穌基督來.
讓我們去面對這個人裡面.
很難面對的自己內心的罪.
他差派了他的兒子.
寫在十字架上面.
所以你去教會.
今天可能你第一次.
或者你去教會幾次.
你一定會看到這個十字架的標誌.
代表神給了我們救恩.
讓我們能夠填補這個.
跟他分隔的空洞.
讓我們跟他重新接連.
有永遠的生命.
各位可能你說.
我是一個好人.
我講什麼罪呢.
但是一個好人都需要耶穌.
可能你聽過這個人.
田家丙.
有沒有聽過.
他是香港的大慈善家.
建了一千間學校.

$^{761}$香港有十一間中小學.
幼稚園在內地更多.
助人無數.
八十歲的時候.
他還把自己九龍塘的房子賣掉.
那時候賣了五千六百萬.
全部用來助學.
自己就租房子住.
他曾經說過.
很多人都不明白.
為什麼他會在九十五歲.
這麼高齡的時候.
信耶穌.
他這樣說.
坦白說.
我平生都是憑著良心做人.
我會毫不面紅地說.
我是一個好人.
根本不需要耶穌.
但他後來坦白地說.
但在人眼中的好人.
都有七情六慾.
一樣會犯錯.
所以人眼中的好人.
都是罪人.
的確人人都需要神的救恩.
能夠脫離罪惡.
田家丙也明白這個道理.
人人都需要上帝給我們的恩典.
你說只有你的上帝嗎?.
是的.
聖經裡說.
創造天地的上帝.
就是給了你這種解決方法.
如果你知道.
普天下當時有疫症的時候.
只有一種可以解決疫症的藥物.
你會不會再吃一些沒用的藥.
或者只有一種解決方法.
一種方式可以救你.

$^{801}$你會不會做其他無謂的事.
不會的.
所以聖經清楚說.
人和神分隔開罪的問題.
唯一解救出路就是耶穌.
你仍然想留在黑暗裡.
等待心上事成?.
還是願意去找這位耶穌呢?.
或者耶穌來找你.
你願意接受祂呢?.
\newpage



\section{以弗所書 6:18-20}
\label{sec:6BPEnUSuXaU}
\textbf{2024.09.08 《凡人的祈禱》 以弗所書6\_18-20 (和修版) 講員 : 曾裔貴牧師}
\newline
\newline
連結: \href{https://youtube.com/watch?v=6BPEnUSuXaU}{\texttt{https://youtube.com/watch?v=6BPEnUSuXaU}} ~~~~ 語音日期: 2024-09-08
\newline
\newline
\hyperref[sec:XLqWIi0boGI]{\small{< < < PREV SERMON < < <}}
~
\hyperref[sec:index_chronic]{\small{[返順時目]}}
~
\hyperref[sec:index_scriptual]{\small{[返順卷目]}}
~
\hyperref[sec:w9qOKksHuCE]{\small{> > > NEXT SERMON > > >}}
\newline
\newline
以弗所書 6:18-20
\newline
\begin{longtable}{cl}
\hline
\hline
章節 & 經文 (和合本修訂版)\\
\hline
6:18 & \begin{tabularx}{0.7\textwidth}{X} 要靠著聖靈，隨時多方禱告祈求，並要為此警醒不倦，為眾聖徒祈求。 \end{tabularx} \\ \\ \relax
6:19 & \begin{tabularx}{0.7\textwidth}{X} 也要為我祈求，讓我有口才，能放膽開口講明福音的奧祕， \end{tabularx} \\ \\ \relax
6:20 & \begin{tabularx}{0.7\textwidth}{X} 我為這福音的奧祕作了帶鐵鏈的使者，讓我能照著當盡的本分放膽宣講。 \end{tabularx} \\ \\
[1ex]
\hline
\hline
\end{longtable}
$^{1}$請各位主人平安.
先和大家做一個簡單的祈禱.
在上主的面前.
求您將您的話語教導我們.
我們知道我們要祈禱.
但求您讓我們認識我們是什麼人.
向哪一位禱告.
向哪一位要求我們為他祈禱.
讓我們深深地認識.
使我們寶貴我們這麼重要的身份.
和禱告守望的職事.
使地上因為我們有祈禱的人.
神的國度可以成就彰顯.
而且是超過我們所求所想.
今早我們一同等候在您的面前.
奉耶穌基督的名.
阿們.
相信答案是顯而易見的.
簡單.
這段聖經已經點出了.
哪一位要求我們祈禱呢?.
戴著鎖鏈的保羅.
是這位使徒.
為著神的福音.
要來坐牢.
在當時羅馬尼祿皇帝的時候.
一而再說了幾次.
你們要為我祈禱.
為什麼要祈禱呢?.
就是要講明福音的奧秘.
保羅不是在《新經八戰》.
已經是很有經驗的使徒嗎?.
為什麼要叫門徒.
在外面有自由的.
已不鎖教會的弟兄姊妹.
要為他祈禱呢?.
什麼口才呢?.
他需要廣東話?.
英文?.
相信不是語言的問題吧.

$^{41}$他認識希臘文.
他是希伯來人受訓練的法利賽人.
一個拉比.
但重點是.
他要有膽量.
一個這麼有經驗.
在君王面前.
在這麼重要的場合.
加瑪列門下.
很多的訓練.
他需要口才.
需要膽量.
來講明福音的奧秘.
可想而知.
保羅不是隨口要求.
他真的很認真.
如果他要放監.
他就要講明福音的奧秘.
我們要思想一下.
今天我們要去看.
為什麼保羅要我們祈禱?.
為什麼他需要教會為他祈禱?.
弟兄姊妹.
有沒有想過.
你可以為這麼重要的人去祈禱?.
你有沒有看重你個人的祈禱?.
剛才在祈禱的時候.
有一樣東西其實很重要.
認不認識我們自己的身份?.
我用的題目是凡人.
即是我們不懂得飛天遁地.
我們祈完禱.
不會像家裡的小狗一樣說話.
我們感受到母堂有位姊妹.
她家裡有一段時間.
在她的祈禱區.
她說她的小狗會祈禱.
張開雙手.
當然不是講人話.
是講狗話.

$^{81}$叫兩聲.
不知道是對號入座.
還是她說會祈禱.
動作好像是一個條件反射.
但是弟兄姊妹.
我們是凡人.
不過.
已不所書.
保羅.
很想我們認識.
我們為什麼可以為保羅去祈禱.
我們有什麼身份.
可以為這位使徒.
帶著鐵鏈鎖鏈的使徒.
在內監裡面.
來到去都不知道明天.
我們可以為他祈禱.
祈禱不是他不再貪生怕死.
而是有膽量的.
放出去之後.
他能夠講明.
福音的奧秘.
將奧秘講明白給人聽.
這多麼偉大.
就是每一個外面的凡人.
可以為保羅祈禱.
所以你會看到.
重點是講明福音的奧秘.
我們繼續思考這段經文.
很重要的情況就是.
真的要講明.
怎樣去講明呢.
各位.
我們真是凡人.
我們是誰呢.
不過.
來到教堂.
來到教會.
我們要在神的眼中.
去看你和我的新身份.

$^{121}$不是你再是你自己.
你要在主的救恩眼中.
去看我們自己的新身份.
這是很重要的.
我們告訴我們.
由第六章的第一到第十九節.
有一段經文很重要.
告訴我們.
我們在家庭倫理裡面.
這裡提到.
當然五章不夠時間說.
六章的一到第九節.
有兩個大段落.
就是當時羅馬帝國之下.
信主耶穌的家庭生活.
是應該怎樣呢.
保羅很期望弟兄姊妹.
能夠有父母.
有兒女.
一起去聽他的教導.
讓那個家庭.
在救恩.
在神的真理下.
成長.
剛才領聖餐的時候.
泰國的姊妹.
安娜帕桑.
抱著孩子出來領聖餐.
這樣的場景是很美的.
保羅就是期望.
每一個家庭的成員.
在教會裡得到神的真理.
作兒女的應該怎樣呢.
是理所當然.
理所當然的.
聽從父母.
作父母的.
又要用主的教訓和警戒.
去養育自己的兒女.
只是這兩點.

$^{161}$在羅馬帝國當時已經是很顛覆性的.
已經告訴我們.
原來神的期待的家庭的倫理.
是要這麼美和和諧的.
興有弟恭.
父母子女是很融合的.
大家的相處.
很開心的生活.
有談有講.
但是有規範.
有教育.
甚至有警戒.
你可以想像到.
保羅期望每一個家庭.
都是這樣來到.
在這樣的愛裡成長.
然後.
做奴隸的.
要用誠實的心.
去聽從你肉身的主人.
好像聽從基督一樣.
一個奴隸信了耶穌.
繼續做奴隸.
但是現在有一位主在天上.
不再看我是奴役.
是奴隸.
看我是他的兒女.
但是在地上.
仍然要用地上的責任身份.
來信服誠實的心.
來聽從肉身的主人.
好像聽從基督一樣.
做主人的.
你就要來到.
好像你自己服侍基督一樣.
不要虐待.
你想想那個倫理.
是多麼和諧合一.
保羅眼中的教會弟兄姊妹.
是如此的平凡.

$^{201}$是如此的服侍.
你留意一下圖片裡.
提了多少次在主類.
提了多少次照著主的教訓.
好像聽從基督.
要像基督的僕人.
從心靈裡進行神的旨意.
必按所行的.
得到主的賞賜.
主耶穌是不偏待人的.
這樣來建構一個新的群體.
這個群體.
就是可以為保羅.
守望.
禱告.
使他能夠更加領受.
神的福音奧秘.
講明白.
講明白很重要.
人家聽不明白你說什麼.
或者人家看見你的樣子.
也不想聽你說.
那就不是講明白了.
沒有溝通的機會.
看見你的樣子已經想走開了.
那不是真正的溝通.
保羅說他要有口才.
他要在不同的人當中.
講明白這個福音的奧秘.
講明聽得明.
自然就會相信.
因為聖靈可以藉著保羅的講解.
對方就會打開心靈.
相信主.
你說多麼重要.
但是弟兄姊妹.
保羅先告訴我們.
我們每一天的生活.
都是在這些範疇裡.
家裡如果是融洽相處.

$^{241}$家裡如果是父母子女.
是非常有愛.
有話題.
有溝通.
有相愛.
可以一起去吃自助餐.
可以一起去看電視.
可以一起在教會裡敬拜神.
昨晚我在晚堂講道.
兒子就在那裡領唱.
不是我編的.
應該是師班的編.
不關我的事.
兒子就說.
為什麼這麼巧又跟你一起同台呢.
事奉呢.
是很美好的一件事.
大家一起服侍神.
大家一起很開心地.
在不同的場景.
弟兄姊妹.
你有沒有這樣的家庭工作倫理呢.
我們是不是有這樣的新群體.
來承載神的福音呢.
不要小看這樣的家庭倫理.
因為第二段聖經.
突然帶我們去看一個.
轟天動地的屬靈征戰.
去到一個有規劃.
有秩序的.
屬天的黑暗勢力.
要我們去面對.
我們從來沒有想過.
信了耶穌.
我們要面對這些魔鬼的戰爭.
保羅就帶我們去看.
這樣的天上的戰爭.
是阻礙我們建構福音.
和家庭生活.
我們可以去看一看後面的經文.

$^{281}$如果你想想.
一個家庭不和睦.
作為父母.
如果用當時羅馬帝國的.
作為父母的權柄.
那些無理.
虐待自己的兒女.
就可能會有這樣的例子.
大家知不知道這位叫遼國豪.
台灣一個很轟動的人.
在2010年的時候.
他公開向傳媒自首.
他要找當時的台灣議員.
一個立委.
指定要他來.
因為我只信任你.
要他去到一個地方.
來自首.
由他陪同.
去到台中的警署.
原來這個遼國豪.
是當時一個通緝犯.
是一個職業殺手.
完全殺人不眨眼.
但是他只是一個年輕人.
他第一句說話就是說.
台灣的教育害死我.
因為他的成長.
爸爸黑社會.
坐牢後來很早就死了.
自小就在這樣的破碎家庭長大.
他說他很深刻.
跟這個立委說.
我的小學.
很早就已經被老師放棄.
完全把他當透明看不到.
有什麼事情就賴他.
他就在這樣的成長裡面.
沒有父母.
沒有父親.

$^{321}$沒有父親的教導影子.
完全不知道什麼對錯.
很自然就被收納.
成為一個少年人.
16歲.
拿槍.
很鎮定的對一個黑社會的頭目.
近距離射殺.
職業殺手.
你回去查一查他的資料.
沒有了這些家庭規範.
一個很好的人.
可以變成一個壞人.
變成一個危害社會的人.
李家同教授.
台灣清華大學的教授.
很出名的.
評論這個廖國豪.
台灣的教育要思想.
我們拼命把人放進名校.
最好的學校.
拔尖.
對這些永遠在邊緣.
在家庭已經有困難的.
我們不能夠教育.
伸出援手.
對這些問題學生.
加以輔導.
用教育.
用警戒.
使他們能夠成長.
這是其中一個很典型的例子.
不過是由他的口.
去控訴台灣的教育.
這件事已經是2010年的時候.
這個廖國豪.
就是非常出名的.
其實他的樣子是很斯文的.
一點都不是殺人不眨眼的.
你可以想像.

$^{361}$他的心靈受到的損害.
感恩.
在普羅的眼中.
已不鎖的地方.
信了主耶穌是一個新群體.
開始要想想.
作為兒女的.
你要怎樣在主裡面聽從父母.
是大應許的誡命.
聽從.
信服.
原來在神的眼裡.
是蒙福的.
作為父母的.
又不可以濫用你的權.
不可以喜樂無常.
這樣的成長.
造就的是一個家庭.
一個家庭.
在教會裡面.
就會有很大的力量.
就會放重心思.
在神的國度的奧秘裡面.
我們進入了那個屬靈戰爭.
第二個段落.
告訴我們.
我們有一個很大的戰爭要面對.
中馬田牧師.
大家可能不太認識.
上個世紀.
50年代.
60年代.
英國.
威斯敏斯特大教堂的主任牧師.
在英國來說.
是最重要的一個教會.
他講了一句話.
教會非常值得憂慮.
我們關注很多事情.
唯一.

$^{401}$我們沒有關注的.
撒旦的存在.
撒旦對教會.
對人的影響.
告訴我們.
我們需要去裝備自己成為剛強的人.
來到這個磨難的日子.
我們要抵擋仇敵.
而且抵擋完之後.
保羅告訴我們.
我們要站穩.
穩定是很重要的.
教會的弟兄姊妹.
信主月來的.
你的屬靈生命是要穩定.
並且要站立得穩.
經過一切的爭戰.
不同的風波.
你也能夠站立得穩.
是保羅的要求.
難怪.
如果我們不認定魔鬼.
對神的工作.
攻擊的時候.
我們就很難.
在這些我們覺得.
初信班就應該學的東西.
去追求.
信主月來的就不需要了.
保羅提出了六件事.
去裝備我們.
面對這場屬靈戰爭.
而且是有規模.
有秩序的.
那個詭計.
所以真理.
公義.
福音.
信心.
救恩的頭盔.

$^{441}$神的道.
那種游刃有餘的運用.
是很重要的.
基督徒要裝備這些.
保羅看到一個在他身邊.
根據一些歷史說.
那些囚犯是要跟那個士兵.
綁在一起的.
不讓他走.
香港或者看電影都有.
手銬扣住.
不讓走.
一起回警署.
這個也不奇怪.
但在坐牢的時候.
那個囚犯跟那個士兵是綁在一起的.
所以保羅看到的士兵.
就令到他有一個屬靈的聯想.
在外面的平凡的弟兄姊妹.
要面對屬靈的戰爭.
保羅不是瘋癲.
是他突然間看到.
一個屬靈的很重要的戰場.
所以弟兄姊妹.
信了主耶穌的弟兄姊妹.
一個家庭一個家庭的信主.
在主的教育底下.
來去有好的家庭關係.
但不只是停留在家庭的日常生活裡面.
而是我們看到.
原來我們信了主.
就加入保羅給我們看到的屬靈戰爭的陣營.
所以要裝備.
如果不裝備.
就不知不覺間.
就被魔鬼帶我們離開了那個真道.
不來參與為神國奧秘的福音禱告.
所以你會看到這六樣東西.
其實都是告訴我們.
救恩令到我們有屬靈的保護.

$^{481}$這個屬靈的保護.
使我們對永恆更加的認識.
使我們知道.
我們雖然生活在地上.
但是我們看到屬天的真實.
其中不要忘記有魔鬼存在.
如果保羅說家庭.
他只需要拆毀家庭的倫理.
他只需要令到做兒女的.
越來越暴戾.
越來越反叛.
明明是對的都不肯聽.
為反對而反對.
作父母的控制不了自己的情緒.
無故無理的找自己的兒女出氣.
拆毀這些關係.
在工作地方反對.
藐視.
不尊重.
欺壓.
各方面.
由你睡醒到你睡覺.
你都是支離破碎.
任何關係都是不開心的.
你何來感受到屬天的能力.
屬天的平安.
在你的生命裡流露出來呢.
自己每天都面對很多憤怒.
很多不和.
很多疑神疑鬼.
各樣的猜測.
你何來有平安的福音.
跟別人分享呢.
一出電梯.
見到一個看監.
你已經想發脾氣了.
你何來能夠有這份平安呢.
何來你能夠再加入保羅這個祈禱的陣營呢.
親愛的兄弟姐妹.
今早最重要.

$^{521}$不可以忘記魔鬼的存在.
但不是叫你看多些恐怖片.
看看魔鬼是怎樣.
整段聖經保羅完全沒有說他的來源.
沒有說他是甚麼樣子.
但只告訴我們.
我們要站穩在真理.
所以要裝備這些真理.
成為剛強的人.
所以要對福音有一個全面投切的認識.
和常常去用這個福音.
你可以用你的口.
唱頌福音的聖詩.
用你的音域.
引領人去到救恩的泉源.
你可以用不同的方法.
首先就是你一定要你自己的內心.
感受到你是屬主的.
你要站穩在神的真理.
穿戴神所賜的全副的軍裝.
神已經給了我們.
主耶穌是全被的救主.
我們需要對救恩有全面的認識.
以致我們能夠在危難的時間.
我們站穩在神的真理.
弟兄姊妹,這是一個很重要的課題.
如果我們不去裝備對神的認識.
我們只會慢慢停留在自己的感覺.
我們會面對困難的時候.
魔鬼就不知不覺.
令我們失去信心和信仰.
這位大法官約翰·羅拔.
他的兒子哈佛大學畢業.
他剛好被邀請去作訓練.
第一句說話.
各位同學,你們都是社會的精英.
我不會像其他畢業演講一樣.
來恭喜你們,祝福你們.
他說,我希望你們在畢業後離開校門.
經常遇見.

$^{561}$是哪五樣東西呢?.
不是祝福.
我希望你們經常得到不公平的對待.
這樣你們就會知道什麼是公義.
什麼是公平公正.
非常有智慧.
但真的很殘忍.
去到哪裡都不公平對待.
在麥當勞排隊,人家排隊打你尖.
職員就說應該打你尖.
沒理由這樣的.
當然,社會不會這麼簡單地不公平對待.
一定會有些技術性.
有些集團式的不公平對待.
如果一個精英不斷有不公平對待.
而他慢慢能夠認識.
什麼是真正的公平.
他渴望的公平.
那就很不同了.
就是說他將來有權的時候.
他就會讓不公平盡量消失.
第二方面.
希望你能夠感覺孤單.
讓你感受到.
友誼原來不是必然的.
真正的友誼.
是你應該要經營和珍惜的.
很有見地的大法官.
再下去.
再恐怖一點.
被背叛.
經常走霉運.
這樣的糞便.
如果有機會.
今晚你去搜索一下.
這位大法官的訓詞.
英文的.
很有智慧.
叫大學生不要囚著滿志.
出去社會.

$^{601}$慢慢建立你的權力.
成為將來危害的人.
相反.
帶著謙卑的心靈.
在各樣的困境裡.
被背叛.
就感受到公正的重要.
遭遇厄運.
原來成功不是必然的.
而且你要更加珍惜.
有同理心.
有些人沒有那麼好運.
原來自己的不幸.
可以建立一份同理心.
對人會有一份憐憫.
最後.
就是被人忽視.
你就會感受到.
你應該學習聆聽.
學習對別人有一種體諒.
他不是要無的放矢.
不是要去喜樂.
包括他自己的兒子畢業.
他真的由衷地將他自己的人生經驗.
去講出來.
楊希望這一班.
快要走出社會的大學精英.
哈佛大學裡面的這一班學生.
能夠成為真正.
可以承載別人的故事的人.
基督徒的我們.
我們在我們的社群裡面.
我們是不是站穩呢.
我們是不是有真理在我們的內心裡面.
我們去分享神的福音呢.
我們對鞋子是不是平安的福音.
我們的頭盔是不是有救恩呢.
我們是不是有公義.
良善.
正直.

$^{641}$去護衛我們的心胸呢.
因為我們隨時都想整人.
這個是保羅期望教會的弟兄姊妹.
現在教會很有自由在外面.
是有保羅和他的同工坐牢.
包括這個以巴弗.
還有這個推基古等等.
但是保羅要求的.
正正就是我們能夠在家庭社會的生活工作的倫理中間.
在神安排我們的人生裡面.
我們怎樣去自處.
怎樣去活呢.
這個是至關重要.
我們需要在我們這些倫理裡面.
看到我們需要站穩.
站穩之後.
不只是家庭.
不只是工作.
不只是這個地球物質的世界.
我們看到的是屬靈的世界.
是有這個戰爭.
所以我們要認定這個福音的大能.
和認定基督是已經得勝.
我們就將人帶進救恩裡面.
這條領帶.
是從美國寄回來的.
我提過有一位的姐妹.
我完全不認識她.
2018年的時候.
她在美國開車.
走出去.
被一個大學生開車.
將她的車推到外面.
推出公路之後.
有一輛高速的車迎面衝過來.
從那一刻開始.
這位姐妹就變成植物人.
到今天.
為什麼我會認識她呢.
我不認識她的.

$^{681}$她的姐姐去美國探望她.
遇到這個特大的車禍.
她的姐姐輾轉.
在沒去之前.
居然不知道是誰告訴她.
錦繡堂有那些叫聖師.
我忘記了名字了.
勵士牧師錄製的.
很短的聖師和聖師的故事的講解.
這個姐姐聽完之後.
姓陳的姐妹.
聽完之後很感動.
去到美國之後.
將錦繡堂的教會網頁的連結.
給她的妹妹.
她妹妹是不信耶穌的.
聽完這些聖師.
太多時間了.
每天就要在床上.
去教會的網頁聽道.
聽著聽著.
聽到我講道.
就是星期四晚的茶經班.
每一堂是一個半小時.
她聽完之後.
就跟姐姐說.
姐姐你能不能拿到曾牧師的電話和WhatsApp.
我想和他談一些事情.
姐姐就問香港.
她不是回我們教會的.
輾轉就問了.
找了我的電話.
告訴我她妹妹的故事.
我說當然可以.
輾轉就來對話了.
她妹妹第一句說話.
她說你講的道我完全聽不明白.
不知道在說什麼.
因為她真的不信耶穌.
她說不知道為什麼.

$^{721}$我聽完你的道.
我的心靈裡面有很大的平安.
我不知道為什麼.
我就說.
我邀請她信耶穌.
她說我還不行.
我還不能夠接受我自己現在的狀況.
我不可以相信.
但繼續.
每個星期聽我的講道.
在美國三藩市.
網上真的這麼好.
聽完之後.
那段時間沒有再找我.
很久了.
她說我已經信了耶穌.
我自己決志.
這樣信耶穌.
她說她現在的痛.
是你們想像不到的.
動不了.
但痛得很厲害.
還要面對高昂的醫療費.
丈夫還不相信耶穌.
不明白他.
在疫情期間失業.
面對很多困難.
我安慰他.
當然是WhatsApp.
他說你不用擔心我.
我有了這個信仰.
因為我有時候講道會分享自己的近況.
我就說我是八月生日.
他記得.
叫姐姐去G2000買一條胎給我.
這樣一份珍貴的禮物.
是救恩帶來的生命改變.
保羅說.
使我有口才.
開口放膽.

$^{761}$講明福音的奧秘.
改變人的生命.
每個場合.
對著一班士兵要怎樣講道.
對著一班法利賽人要怎樣講道.
是需要口才.
是需要膽量.
是需要恩典.
是需要整個生命的改變.
講明福音的奧秘.
弟兄姊妹.
保羅邀請我們去祈禱.
我們其實是平凡人.
但在主耶穌的救恩裡.
我們是新的群體.
以弗所書很強調.
我們是在主裡合一的群體.
我們可以發很大的力量.
使福音被人聽明白.
所以很盼望神的恩典.
在我們洪恩家.
在不同的教會.
我們一起努力.
將神的福音傳揚.
願主祝福大家.
我們一起祈禱.
感恩天父讓我們今早.
可以一同領受聖餐.
更加我們能夠在神的恩典裡.
去服侍.
多謝主.
但願我們能夠在你的救贖裡.
更加明白我們身份的尊貴.
一起去服侍你.
奉耶穌基督的名.
阿們.
\newpage



\section{約翰福音 12:20-26}
\label{sec:w9qOKksHuCE}
\textbf{2024.09.15 《作一粒落在地裡的麥子》約翰福音12\_20-26 (和修版) 曾敏姑娘}
\newline
\newline
連結: \href{https://youtube.com/watch?v=w9qOKksHuCE}{\texttt{https://youtube.com/watch?v=w9qOKksHuCE}} ~~~~ 語音日期: 2024-09-16
\newline
\newline
\hyperref[sec:6BPEnUSuXaU]{\small{< < < PREV SERMON < < <}}
~
\hyperref[sec:index_chronic]{\small{[返順時目]}}
~
\hyperref[sec:index_scriptual]{\small{[返順卷目]}}
~
\hyperref[sec:upkPX3gy4Ao]{\small{> > > NEXT SERMON > > >}}
\newline
\newline
約翰福音 12:20-26
\newline
\begin{longtable}{cl}
\hline
\hline
章節 & 經文 (和合本修訂版)\\
\hline
12:20 & \begin{tabularx}{0.7\textwidth}{X} 那時，上來過節禮拜的人中，有幾個希臘人。 \end{tabularx} \\ \\ \relax
12:21 & \begin{tabularx}{0.7\textwidth}{X} 他們來見加利利的伯賽大人腓力，請求他說：「先生，我們想見耶穌。」 \end{tabularx} \\ \\ \relax
12:22 & \begin{tabularx}{0.7\textwidth}{X} 腓力去告訴安得烈，然後安得烈同腓力去告訴耶穌。 \end{tabularx} \\ \\ \relax
12:23 & \begin{tabularx}{0.7\textwidth}{X} 耶穌回答他們說：「人子得榮耀的時候到了。 \end{tabularx} \\ \\ \relax
12:24 & \begin{tabularx}{0.7\textwidth}{X} 我實實在在地告訴你們，一粒麥子不落在地裡死了，仍舊是一粒；若是死了，就結出許多子粒來。 \end{tabularx} \\ \\ \relax
12:25 & \begin{tabularx}{0.7\textwidth}{X} 愛惜自己性命的，就喪失性命；那恨惡自己在這世上的性命的，要保全性命到永生。 \end{tabularx} \\ \\ \relax
12:26 & \begin{tabularx}{0.7\textwidth}{X} 若有人服事我，就當跟從我；我在哪裡，服事我的人也要在哪裡；若有人服事我，我父必尊重他。」 \end{tabularx} \\ \\
[1ex]
\hline
\hline
\end{longtable}
$^{1}$弟兄姊妹平安.
開始時再有禱告.
天父上帝 求你讓孩子的生命被聖靈充滿.
讓孩子的生命說出上帝要求的訊息.
也深願神讓聖靈充滿每一個弟兄姊妹的生命.
讓弟兄姊妹的心傳言向上帝的道開放.
讓我們聽了 願意回應 生命就成長成熟.
奉耶穌基督的名字祈禱.
阿們.
這個功課非常寶貴.
墨寺的功課很值得我們思考.
我們今天再仔細看這段經文.
作一粒落在地上的墨紙.
這裡說了一粒墨紙的比喻.
注意是一粒 不是一堆墨紙.
不是一大群墨紙 是一粒.
一粒墨紙和一大群一堆有什麼不同呢.
一粒其實說真的.
它的影響力有多大.
如果在地面上的話.
不是在地底的話.
是沒有影響力的.
是一粒的 是不是.
有原因的 說個比喻.
一堆又不同.
一堆一群的時候.
一大堆就有很多食物.
你直接想到這一點.
是很有影響力.
一粒在地面上不是地底.
好像沒有什麼意思.
但是今天耶穌要我們.
透過這一粒墨紙.
去思考我們生命的功課.
看看背景是怎樣的.
當時有幾個希臘人.
去求見耶穌.
安德烈和肥力就告訴耶穌.
耶穌的回應很有趣.
沒有直接說什麼.

$^{41}$直接說一個比喻.
就是這一粒墨紙的比喻.
在說之前 他說.
孕子得榮耀的時候到了.
原來這個時段.
是耶穌基督已經進入耶路撒冷.
將會面對後面的受苦之路.
在這裡他就說.
孕子得榮耀的時候到了.
通常我們想得榮耀是什麼.
這個人被高升 被高舉.
得到很大的冠冕.
有很多很大的稱讚.
我們想地面上的時候.
總會想到一些權力.
很重要的人物.
什麼的家面.
但是耶穌在這裡說.
孕子得榮耀的時候.
是不是在說這些呢?.
原來不是.
是在說耶穌基督.
即將要踏上受苦之路.
是即將要經歷.
被他所默養幾年的門徒.
去出賣.
他要經歷被人出賣.
他要經歷不公義的審訊.
被人奚落嘲笑.
經歷很殘酷的鞭打.
最後要被釘死在十字架上.
他要擺上他的姓名.
十字架在古羅馬的時間是死刑.
誰被釘在上面.
其實說真的要經歷的是什麼.
是羞辱.
這簡直是我們人生裡.
最低谷,最羞辱,最痛苦的時間.
但耶穌基督說.
孕子得榮耀的時候到了.

$^{81}$為什麼呢?.
他在十字架上.
那個犧牲是為了什麼呢?.
是我們非常熟悉的.
我們的福音的核心內容.
就是為了我和你的罪.
為全世界的人的罪死在上面.
是為了我們去負上生命的代價.
去甘願受羞辱,受痛苦.
去到人生最低點.
就是這裡.
卻是得榮耀的時候.
原來世人看的和天父去看的是不同的.
世人看這個很羞辱.
一個很低微的位置.
但天父去看他的兒子的時候.
這個卻是他得到榮耀的時候.
最高峰的時候.
這個是必經之路.
受苦之路.
就在這個大前提之下.
開始說一粒墨紙.
說這一粒墨紙是怎樣的呢?.
大家看回經文.
一粒墨紙不落在地裡.
死了,仍舊是一粒.
剛才我說的.
是一粒.
那一粒就放在地面上.
總之它打死都不落在地裡.
不被埋在地裡.
它永遠就是那一粒.
到時候其實沒有用了.
就成為垃圾.
拿去燒.
完全沒有價值可言.
所以這裡說.
一粒墨紙不落在地裡.
死了就是一粒.
但如果它落在地裡.

$^{121}$就結出很多紙粒.
什麼叫落在地裡呢?.
我就稍稍了解了一下.
原來落在地裡.
至少要有一個深度.
不是很薄很薄很薄的一層泥.
而是要埋在泥土下面.
三厘米以下.
被泥完全覆蓋著.
在這個時候.
墨紙經歷的是什麼呢?.
完全和外界隔絕.
它面對很大的壓力.
是泥土的壓力.
在這個時候.
包圍著它的是什麼呢?.
是泥土裡面的濕度.
有溫度,壓力,熱力.
還有落肥料的話.
有些臭味.
是不是?.
這就是一粒墨紙.
在泥裡面經歷的.
如果這粒墨紙.
願意這樣被埋在裡面.
它會面對這一切.
接著.
這一切外在的東西.
濕度,溫度,熱力,壓力.
肥料等等.
就會令到這粒墨紙.
有一個很堅硬的外殼.
就會裂開.
令到它裂開.
這個時候.
我們會形容這個墨紙.
死了.
墨紙原來就是這樣.
這個時候.
是一個屬於死亡的階段.

$^{161}$死了之後.
它反而會分解.
它的杯牙.
漸漸地成長起來.
接著就紮根.
在泥土裡面.
吸取養分.
慢慢地.
就發出很漂亮的幼嫩的牙苗.
然後大概.
經過四個月的時間.
新一株的苗.
就會長成.
就會結出.
很多很多的紫粒.
很厲害.
你想像不到.
這個紫粒的影響力.
是很多很多很多粒.
是造福很多的人.
成為不同的人的食物.
這個就是一個埋在地裡.
死亡.
然後再.
生出來的一個過程.
所以這裡說.
當它在地裡.
死了.
就結出很多紫粒.
經文下面是說到.
下面有一個平衡的句子.
這樣才說.
耶穌基督在說什麼呢.
下面這裡說.
愛惜自己的性命.
就喪失性命.
其實這個世界.
很多人.
非常多人很愛惜自己的性命.
整個生命的重心.

$^{201}$是為自己而活的.
我整個方向.
是為自己而活的.
我所有的東西.
所有的東西都是希望自己.
我的慾望.
得到最大的滿足.
我所有的東西都是自己喜歡的.
我想怎樣就怎樣.
我自己決定.
我拿著權.
全世界圍著我轉.
為什麼這麼多憤怒.
為什麼這麼多不開心.
因為我期望別人都為我而轉.
所有人都為我而滿足.
但偏偏這個世界.
就不是這樣的.
因為很多人都是為自己而轉.
為自己而生活.
大家都感到很不滿足.
別人都不能滿足我.
達不到我的期望.
好像我的慾望.
是一個無底洞.
很深很深.
上上上著別人虧欠我的.
得罪我的.
我們愛惜自己的性命.
就是這樣.
一生都是為自己而活.
在這裡.
經文很直接說.
就喪失性命.
這麼愛自己.
反而失去性命.
很直接.
這裡沒有說.
很愛自己性命.
你就得到滿足.

$^{241}$你的生命就非常美好.
很圓滿.
不是.
相反.
你越愛惜就死得越快.
死得越徹底.
反過來.
恨惡自己在世上的性命.
反而.
要保存性命.
去到永生.
這裡的意思不是.
我由現在開始就傷害自己.
我讓自己沒一口好食.
傷害我的肉體.
傷害我的精神.
各方面苦練.
教導自己整個人很痛苦.
煎熬自己.
要保持良好.
要肯定自己的性命.
所以我現在就來了.
不是說這一種.
而是你真的要將自己放下.
以後人生的方向.
不要為了滿足自己的慾望.
我拿著權柄.
我希望怎樣就怎樣.
我全世界都希望為自己.
不是這樣.
放下這一切.
反而.
以上帝為我們生命的主.
以上帝為我們的主.
跟祂走.
為了滿足神.
而去生活.
真的要放下自己.
要放下自己.
恨惡自己.

$^{281}$你不是想.
我這個肉體.
怎樣的慾望得到滿足.
你是想神的心意.
怎樣得到滿足.
上帝喜悅我怎樣.
我跟祂走.
我想上帝很滿足.
很快樂.
上帝希望我在這個世界.
做不同的事.
是為祂緣故.
那我去做.
跟著上帝的心意.
不一定要有東西令你很快樂.
可能神要你服事.
很難服事的人.
很難忍受的人.
做一些你不喜歡做的事.
什麼叫恨惡自己.
就是放下自己.
不是放1\%.
2\%.
是全放下.
這叫恨惡自己生命.
耶穌基督.
做了一個很好的榜樣.
祂說.
自己的人子.
在榮耀的時候就到了.
耶穌基督.
作為神的兒子.
享有一切的豐盛.
是在經歷完全滿足的愛.
他根本不需要.
來到這個世界.
不需要讓自己.
經歷這麼多的限制.
之後經歷這麼多的痛苦.
為了什麼.

$^{321}$犯罪的人類.
不是非常可愛的.
很聖潔的人類.
不是.
是為了犯罪的人類.
你想想自己的生命.
想想整個世界.
發生什麼事.
多少個罪惡.
我們自己多少個罪惡.
世界多少個罪惡.
耶穌基督就是要來.
為了這麼不可愛的人.
犧牲自己的生命.
全然的.
放下自己.
祂就是那顆墨紙.
落在地裡.
去謙卑.
去捨己.
整個生命.
拿出來.
祂死了之後.
被三天從死裡復活.
成就救恩.
拯救.
數不盡人的生命.
千千萬萬人的生命.
是因為耶穌基督.
得夢拯救.
離開永遠的死亡.
哇!.
耶穌基督的影響力.
多麼的厲害.
祂就是這顆墨紙.
死了之後.
結出許多紙粒來.
去放下祂的生命.
而為世界.
帶來許多新的生命.

$^{361}$是屬於耶穌基督的.
新的生命.
是屬於上帝的生命.
祂做了一個榜樣.
祂說這個比喻.
不只是說.
自己的生命.
祂說這個比喻做什麼?.
就是希望他的跟從者.
都是這樣做.
希望我和你.
都是這樣做.
我和你可以選擇.
是不是?.
選擇就是要留住自己.
我就是要做這顆墨紙.
我自己是主人.
我要滿足我的慾望.
這一生都是為自己而活.
如果是這樣的話.
永遠就是那一顆墨紙.
就在地面上.
不會為別人帶來祝福.
不會帶來更多的生命.
最後.
都是失去.
聖命.
被銷毀.
還是我跟著耶穌基督.
擺上我的聖命.
是不是?.
我含糊在這世上的聖命.
我放下它.
我的聖命就是屬於耶穌基督.
上帝要怎麼用.
就怎麼用.
我交出來.
這個交出.
一點都不簡單.
當日.

$^{401}$其實有.
12個門徒.
跟著耶穌基督.
有一個是.
背叛耶穌的,那位叫猶大.
我們不說他.
另外的11位.
大家知不知道那11位門徒.
後來的結局是怎樣的?.
除了約翰.
大部分都殉道.
彼得是.
十字架.
死得很殘忍,很痛苦.
他就擺上自己的聖命.
每個人都跟著.
耶穌基督,做那顆.
落在地裡的墨紙.
將生命交上.
上帝要我做什麼就做什麼.
上帝要我傳福音就傳福音.
在那時面對很多逼迫.
我願意面對逼迫.
那時要我殉道,我殉道.
你想想,彼得.
之前是怎樣的?.
三次不認主.
主耶穌被捉拿的那一夜.
彼得三次不認主.
很害怕.
但到了耶穌基督復活.
向他們顯現之後.
我們去到使徒行傳的時候.
看到他們一路.
經歷聖靈的洗的時候.
你看到門徒.
多有力量.
為主大放光芒.
傳揚福音.
講道,帶領多人信主.

$^{441}$到最後.
殉道,也在所不惜.
這就是.
跟從耶穌基督的墨紙.
你說最多的.
都是門徒.
這裡.
大家有沒有聽過這本書.
回首百年.
殉道血.
有沒有聽過這本書.
這一本.
簡直是經典.
裡面記載了.
在1900年.
義和團事件裡.
殉道宣教士的.
事情.
1900年很久了.
義和團事件.
裡面有大批的基督徒.
被殺害.
也有大批的宣教士.
被殺害.
這些基督徒.
這些宣教士.
都是為了福音的緣故.
而在中國.
這個地方殉道.
在殉道者裡.
總共有多少人呢.
有189人.
189人裡.
其中有.
52位是.
宣教士的子女.
哇!小時候就做墨紙.
小時候就跟著.
父母做墨紙.
自己的生命完全在.

$^{481}$上帝的手裡.
52位是宣教士的子女.
佔了總人數的四分之一.
當中52位裡.
有8個人是未滿周歲.
有29個人.
是一到五歲.
這麼小的年紀.
就跟著父母.
到中國宣教.
在這個義和田的事件裡.
就被狂徒殺戮.
我們也看過.
有這麼多殉道的人當中.
分別是屬於哪些的.
差會呢.
原來最多的是屬於兩個差會.
分別是什麼呢.
內地會.
就是大德生所建立的.
在裡面有很多很多.
宣教士.
大德生自己非常愛中國人.
他說有多少錢.
有多少性命.
他不會留下.
任何一點不給中國人.
他創立了這個.
差會也感動了.
很多很多的宣教士.
來中國傳福音.
當時很有趣的.
那些宣教士去到中國.
總是選擇一些比較開發的地方.
很多人都堆在那些地方.
去傳福音.
但是內地會的宣教士就特意要選擇一些比較艱苦的.
是福音還沒到的地方.
去傳福音.
就是這樣.

$^{521}$而另一個.
大家有沒有想到哪個差會.
多人去殉道呢.
就是我們的宣道會.
紅印堂是屬於宣道會的.
紅印堂屬於宣道會.
宣道會是非常重視宣教.
原來我們的宣道會當時.
在內也是其中一個.
很多宣教士殉道的差會.
宣道會的宣教士.
也是一樣.
不想堆在很多宣教士集合的地方.
比較開發的地方.
特意就走去偏遠的.
偏僻的地方.
去那裡宣教.
所以在這個過程裡.
兩個差會旗下的宣教士.
是很多人殉道.
這一本書.
是記下他們的事跡.
全部是為福音緣故.
我相信裡面有很多人才.
如果要做一顆墨紙.
為自己而活.
他們可以這樣做.
一生享福.
有很多的才幹.
能力.
有的是青春.
做一顆為自己而活的墨紙就好了.
開心一世.
想享盡人生的各種福浪.
不好嗎?.
但這一群人.
就是為了福音.
為了我們中國的同胞.
甘願將所有東西撇下.
連命都豁出去.

$^{561}$那條命交給主.
就去傳福音.
發生這樣的事情.
殉道在所不惜.
子女跟著一起殉道.
今天.
特別是在香港這個時代.
這些小朋友.
真的被父母寵愛到不得了.
真的寵愛到不得了.
要什麼就什麼.
一點艱難都沒有.
一聽到.
去哪裡?.
或者辛苦一點的行程.
太辛苦了.
兒子,女兒,你不要這樣.
我真的覺得很痛.
一點點都不能承受.
真的很難想像.
這麼多的小朋友.
在這個事件裡殉道.
很難想像.
今天我們的大人也好.
小朋友也好.
在這樣的世界裡.
就是一粒墨紙.
我要在地面上的墨紙.
我不要在地底下的.
但是你看看.
這麼多宣教士殉道.
讓中國的福音事供.
遍地開花.
一直越來越多人信耶穌.
他們帶來一個.
很美好的榜樣.
很美好的見證.
讓人認識原來神是這麼愛我們的.
中國同胞神是這麼愛我們中國人的.
派這麼多人來殉道.

$^{601}$為我們犧牲.
就是這樣.
教會被建立起來.
生命成長改變.
宣道會.
也是這樣的.
裡面的宣教士是這樣擺上犧牲.
香港的教會.
德蒙祝福.
都是因為宣教士的捨己.
犧牲.
是不是一粒墨紙.
落在地面上.
帶來多大的奇妙的影響.
但是這個代價很大.
耶穌基督.
祂做了第一個這樣做的人.
祂的門徒跟著.
歷世歷代.
有不同的人跟著.
當然不單止這些宣教士.
還有世界不同地方的.
基督徒也是.
有不少.
在伊斯蘭國家的.
裡面殉道.
有很多.
為福音,為信仰緣故.
擺上自己的命.
這個很寶貴.
今天我問弟兄姊妹.
無論你信主幾個月.
一年,十年.
幾十年.
你有沒有問過自己.
這粒墨紙.
在地面還是地底.
你這粒墨紙.
在地面還是地底.
如果在地面.

$^{641}$你有沒有想過.
我挑戰你.
把你的性命交出來.
放在神的手裡.
埋在地裡.
讓祂經歷死亡.
讓上帝用你的生命.
祝福多人的生命.
無論代價有多大.
你有沒有想過.
這個除了.
犧牲自己的生命.
我把我的命都擺上了.
為上帝的國度.
為上帝的福音.
為拯救人的靈魂.
還有一件事.
這粒墨紙在地面.
第二件事很重要.
放下救我.
或者致死救我.
致死救我.
剛才我說.
墨紙在地面會死.
要經歷死亡.
才會有轉變.
會發芽生長.
結出不同的子粒.
祝福多人的生命.
要致死自己.
我們這麼久以來.
我想問問大家.
我們的生命轉變了多少.
信耶穌第一年.
心裡很火熱.
有很多東西很渴望.
很想追求.
第二年 第三年.
好像沒有了那種動力.
感覺上也沒有分別.

$^{681}$我和信耶穌之前.
一樣都是這樣生活.
我以前就是這樣.
去惹人生氣.
我現在一樣.
以前是貪心.
我現在一樣貪心.
以前說是非 現在一樣.
以前很喜歡.
生命的主權在我手裡.
現在我一樣.
那顆墨紙的殼.
不知道有多硬.
完全沒有改變.
一樣地面上的生命.
我不會埋在地底.
讓自己的生命去改變.
弟兄姊妹.
我們的生命是不是這樣.
這樣的生命不能夠祝福別人.
沒有意思.
在聖經裡說.
其實是很嚴厲.
不埋在地底死就是一顆.
一顆.
我這麼愛惜自己的生命.
就會喪失生命.
那一顆墨紙.
不願意.
自死自己的那顆墨紙.
最終變成垃圾.
我們的生命.
是不是為了變成垃圾.
我們不願意.
我們的生命.
是不是為了變成垃圾.
這不是天父的心意.
天父的生命是要我們改變的.
大大的改變.
今天說一些唐道的時候.

$^{721}$我也很想說一件事.
我們的感恩.
青少年的團契成立了.
應該是上個禮拜六的時間.
開始了我們第一次.
正式的團契的活動.
很感恩.
當天也有.
十位年輕人.
參加我們的團契.
我們仍然在當中經歷.
在過程裡.
我想說我們的團契叫什麼名字呢.
叫墨紙團.
這個名字.
是神秘的.
神要我們的團.
做墨紙團.
為什麼呢?.
因為要讓我們的生命成長.
改變.
真的做落在地裡的墨紙.
致死自己.
信耶穌之後真的改變.
將生命主權交出來.
徹底的認罪悔改.
讓舊人.
真真正正和耶穌基督一起釘十字架.
這些生命要成長.
要結出子粒來.
祝福多人的生命.
因為這些生命.
神也很想他們學習謙卑.
以致他們可以成就.
上帝很奇妙的作為.
神想他們做墨紙.
但是落在地裡的墨紙.
這個過程心裡也很感動.
也很想藉這個機會.
和大家分享過去的一些經歷.

$^{761}$這些相片.
是我們過去.
為了一路凝聚這群年輕人.
我們有不同的活動.
打羽毛球.
去過長洲.
我也帶他們去過.
周一,日讀書的地方.
神學院.
很多的生日周.
不同的活動遊戲.
一群年輕人的關係.
慢慢慢慢建立起來.
在裡面想說說.
自雖然小.
相信大家看到.
我們在2023年10月成立以來.
神很奇妙.
透過英文時宮.
在裡面吸引了來自團圓天.
不同背景.
不同程度年級的.
一些中學生來到我們的教會.
讓他們在這裡.
去認識教會.
認識信仰.
直到如今我們做了一個統計.
原來透過這個英文班.
我們接觸了.
接近40位的年輕人.
他們會去上英文課.
當中有8位.
已經是決志信主.
也有接近10位.
之前參加.
我們的信仰探索班.
現在他們.
參與我們的團契.
也有接近20位.
在我們每一次的活動裡.

$^{801}$他們會穩定地去參與.
我感恩.
在整個過程.
神的工作.
神親自預備了.
服務的導師.
整個團隊.
神親自帶領.
這群年輕人在教會裡.
參與,是很奇妙的.
整個事宮的需要.
很大,趁著這個機會.
我也很想邀請弟兄姊妹.
為我們禱告.
他們固然覺得學英文.
不容易.
在裡面有很多挑戰.
還有其實每個年輕人.
都面對著他們家庭不同的問題.
他們自己的情緒.
學習上的需要.
人際關係裡的挑戰.
等等.
在裡面,我也很同意.
導師們所說的.
他們最需要的是耶穌基督的愛.
是耶穌基督救贖的恩典.
所以請大家.
切切為這群年輕人.
禱告,為我們的脈子團.
禱告,讓一群中學生.
他們的生命在裡面.
繼續地成長.
為裡面服侍的導師們.
禱告,求主.
加他們的力量,好嗎?.
每一次你想起,就為他們禱告.
因為我們教會,實在是很需要.
有新一代.
透過曾牧師給我很大的支持.

$^{841}$很想我專心地做.
兒童青少年的工作.
在裡面,我很經歷.
上帝祂的大能的供應.
祂奇妙的作為.
真是看到神.
很愛年輕人.
非常地愛他們.
在這裡,也很想說.
神對他們的心意.
就是要他們做那個落在地裡的脈子.
要他們成長.
要他們成熟.
要祝福多人的生命.
最後我想說,神對我們每一個的心意.
都是,想我們落在地裡.
弟兄姊妹.
現在我想大家.
一同進入禱告裡面.
如果一直以來.
我們只想做地面上的脈子.
不要落在地裡.
今天求天父真的寬恕我們.
寬恕我們,讓我們悔改.
甘心將生命的主權交出來.
致死的老鵝.
讓上帝帶領我們人生.
做一顆結子.
一顆這樣的脈子.
我們一起來祈禱.
天父上帝.
我們在這裡獻上感恩.
很愛我們每一個.
神啊,縱然我們生命.
真的很軟弱.
我們信耶穌之後.
仍然會繼續犯罪,得罪你.
還帶著老鵝一起生活.
天父啊,我們.
是何等的虧損你的榮耀呢?.

$^{881}$可能我們有不少人.
仍然堅持地面上的脈子.
不要被埋在地裡.
天父求你寬恕我們.
寬恕我們每一個.
寬恕我們自我.
自私,寬恕我們很多的罪.
我們還和舊人一起.
心願耶穌基督.
以你的寶血,傳染潔淨.
遮蓋我們的罪惡.
聖靈求你充滿我們的生命.
你給我們力量.
致死老鵝.
給我們將生命的主權.
交出來.
我們做地裡的脈子.
我們願意死去.
我們願意這一生.
只為耶穌基督而活.
我們的生命.
是被主所用.
無論上帝要怎樣用我們.
你讓我們甘心樂意.
然後,主啊,心願你可以.
透過我們.
祝福多人的生命.
讓越來越多的人.
因為我們.
認識你.
得蒙救贖.
我們更願意我們的年輕人.
我們的中學生.
都是做這一粒.
埋在地裡的脈子.
讓他們生命成長.
他們成為教會新一代的接班人.
你們祝福多人的生命.
永遠要上帝的名.
向你們紀念脈子團.

$^{921}$紀念我們每一個.
奉耶穌基督的名字.
祈禱.
阿們.
\newpage



\section{申命記 18:15-22}
\label{sec:upkPX3gy4Ao}
\textbf{2024.09.22 《時代先知的工作》 申命記18\_15-22 (和修版) 講員 : 黃永強牧師}
\newline
\newline
連結: \href{https://youtube.com/watch?v=upkPX3gy4Ao}{\texttt{https://youtube.com/watch?v=upkPX3gy4Ao}} ~~~~ 語音日期: 2024-09-24
\newline
\newline
\hyperref[sec:w9qOKksHuCE]{\small{< < < PREV SERMON < < <}}
~
\hyperref[sec:index_chronic]{\small{[返順時目]}}
~
\hyperref[sec:index_scriptual]{\small{[返順卷目]}}
~
\hyperref[sec:dufzJl9BQpc]{\small{> > > NEXT SERMON > > >}}
\newline
\newline
申命記 18:15-22
\newline
\begin{longtable}{cl}
\hline
\hline
章節 & 經文 (和合本修訂版)\\
\hline
18:15 & \begin{tabularx}{0.7\textwidth}{X} 「耶和華－你的神要從你弟兄中給你興起一位先知像我，你們要聽他。 \end{tabularx} \\ \\ \relax
18:16 & \begin{tabularx}{0.7\textwidth}{X} 這正如你在何烈山大會的那日向耶和華－你的神所求的一切，說：『求你不要再叫我聽見耶和華－我神的聲音，也不要再叫我看見這大火，免得我死亡。』 \end{tabularx} \\ \\ \relax
18:17 & \begin{tabularx}{0.7\textwidth}{X} 耶和華對我說：『他們說得對。 \end{tabularx} \\ \\ \relax
18:18 & \begin{tabularx}{0.7\textwidth}{X} 我必在他們弟兄中給他們興起一位先知像你。我要將當說的話放在他口裡；他要將我一切所吩咐的都告訴他們。 \end{tabularx} \\ \\ \relax
18:19 & \begin{tabularx}{0.7\textwidth}{X} 誰不聽從他奉我名所說的話，我必親自向他追究。 \end{tabularx} \\ \\ \relax
18:20 & \begin{tabularx}{0.7\textwidth}{X} 若有先知擅自奉我的名說了我未曾吩咐他說的話，或是奉別神的名說話，那先知就必處死。』 \end{tabularx} \\ \\ \relax
18:21 & \begin{tabularx}{0.7\textwidth}{X} 你心裡若說：『我們怎能知道那話是耶和華未曾吩咐的呢？』 \end{tabularx} \\ \\ \relax
18:22 & \begin{tabularx}{0.7\textwidth}{X} 先知奉耶和華的名說話，所說的若沒有實現，或不應驗，這話就是耶和華未曾吩咐的，而是那先知擅自說的，你不必怕他。」 \end{tabularx} \\ \\
[1ex]
\hline
\hline
\end{longtable}
$^{1}$洪永堂的電影節目平安.
很開心曾牧師邀請我來到大家當中.
宣講成日華語.
首先跟大家說聲抱歉.
大家看到牧師好像匆匆忙忙地走進來.
我真的沒想過.
我以為很近的.
園基堂在Daiso樓上.
十二蚊店.
有沒有電影節目去過Daiso買東西.
有是吧 很方便.
我想20分鐘一定來到.
原來我等車比坐車的時間多.
有些事我們遇不到.
這個時代有些事我們遇不到.
但神在這個時代作事是奇妙可畏.
今天跟大家分享時代先知的工作.
先知我相信大家回教會一段日子.
一定聽過這個名詞.
也知道先知在裡面代表什麼人物.
在舊約裡面談論.
神揀選很多的先知.
但黃牧師在先知前面加了一個時代兩個字.
不是在說歷史之前的事.
是和我們現在這個時代有關.
和你和我有關.
最後先知做了什麼.
今天我們這個時代.
我們向上帝拯救.
跟從耶穌的基督徒.
我們又做了什麼呢.
大家想不想知道.
我們一起去祈禱.
願保衛師聖靈.
也是真理的聖靈.
來自每一個信耶穌基督的基督徒生命裡的真理聖靈.
我們呼求你.
呼求你這一刻運行在我們當中.
幫助我們每一位頂尖的姐妹.
能夠聽得明白上帝的話.

$^{41}$有被上帝的話的感動.
也在當中憑信心去回應上帝的吩咐.
在這個時代去活出真理裡.
要我們所見證的這位榮耀的上帝.
願主你閱臨我們.
將耶誓教堂禱告奉主耶穌基督的聖名而求.
阿們.
我想大家一邊聽一邊擺一條問題.
在你的心裡.
在你的腦海裡.
這個問題就是今天.
是教會實踐大使命的時代.
OK.
舊約先知的身份.
角色和工作.
給我們有什麼行動力呢.
我相信大家回到教會.
大家坐在這裡.
你都曾經經歷過大使命的運動.
簡單來說.
有人給你傳了福音.
以致你有機會聽.
然後知道耶穌是真神.
以致我承認自己是罪人.
我決定去跟從耶穌.
所以今天你坐在這裡敬拜上帝.
是不是.
可能今天我不知道當中有沒有新朋友.
如果有新朋友.
我為你感恩.
因為你竟然不選擇其他東西.
你不去飲茶.
不去威風.
不去玩.
你今天來到教會.
跟其他基督徒一起敬拜上帝.
你好好的來上天聽.
耶穌基督今天要祝福你.
所以.
當我們去想回舊約的先知.

$^{81}$他的身份.
他的角色和工作的時候.
對應我們今天.
我們繼續在大使命的運動裡.
我們怎樣做一個時代的先知呢.
我們進去吧.
首先主席讀過這段經文.
我不再重複.
希望大家仍然記在心裡.
但我想給一個背景.
一個背景.
這段經文在一個什麼背景下.
申命記的作者去記錄.
上帝和摩西在當中對話.
在當中要聽這段說話的受眾.
和我們今天要聽這段說話的聽眾.
有什麼背景我們要知道呢.
申命記基本上一開始已經告訴我們.
其實上帝揀選了摩西.
摩西可以說是舊約聖經裡的一個先知的典範.
一個先知的模範.
上帝在他的時間.
救贖的歷史裡.
揀選摩西.
拯救以色列人出埃及.
接著他們離開埃及之後.
申命記一開始第一章告訴我們.
他們種種原因.
我不再重複.
不然我30分鐘都說不完.
不過我估計30分鐘都說不夠.
但我盡力.
要關咪就要關咪.
他們因為種種原因漂流了曠野40年.
OK.
現在漂流到時間.
神所定的時間.
夠了.
他們去到約大河東.
摩西帶領第一代的以色列人出埃及.

$^{121}$去到現在第二代.
因為第一代已經過去了.
第二代的以色列人.
他們將要進入上帝應許賜給他們的.
流奶與麥芝地.
過約大河之後去到西面.
就是那個迦南地.
就是一個這樣的背景.
但對於摩西來說.
這段說話.
對於他有很複雜的情感.
情感是什麼呢.
他現在做先知.
但上帝好像跟他說.
耶和華.
你的神要從你弟兄中.
給你興起一位先知像我.
你們要聽他.
這個我就是指摩西.
但大家想.
摩西在那個背景.
摩西那時候年紀大概是120歲.
現在120歲了.
似乎聽這句說話.
聽到什麼呢.
是不是我沒得做先知.
好像要找第二個先知來接棒.
是不是有這個意思.
所以我想大家代入多一點經文.
如果你是摩西.
你是什麼心情去迎接.
神在當中興起新的先知呢.
但首先要讓大家知道.
先知不是你做先知就做先知.
誰有膽量說我要做耶和華先知.
你要小心.
因為先知是神自己親自去選擇.
他親自興起.
是神的主權.
不是你做就做.

$^{161}$你要記住一件事.
第二.
神所興起的先知.
有兩個原則.
在經文中.
必須在哪裡興起.
在哪裡去選擇.
在摩西你的兄弟之中.
簡單來說.
在以色列文當中.
再對應今天.
在神的子民當中興起.
不是在國徽信徒.
當然神可以透過一些非信徒.
去說一些神要說的話.
這是神絕對有的主權.
聖經裡也有記載.
但這段經文.
想告訴大家.
神要興起一個像摩西.
要在以色列文當中.
對應我們今天在教會裡.
每一個都是神的子民.
我們都有機會.
被選擇.
被興起.
去做先知的職份和工作.
第二個原則.
就像摩西.
大家有沒有想過.
像摩西是一件什麼事.
如果就像黃博士說.
摩西是先知的模範.
做得模範做得典範.
就沒有什麼別的了.
所以這個字.
留意這個字是像.
像即是有機會不是百分之一百.
一模一樣.
但這是標準.

$^{201}$是標準.
我們提一提.
摩西這個先知的模範.
其實在申命記最後的時候.
申命記的作者給摩西.
做了一個評價.
做了一個描述.
其實在這裡摩西已經離開了人世.
在申命記最後的時候.
她已經完成了她的先知的使命.
申命記的作者給摩西一個什麼評價呢.
不如我們一起讀.
大家看不看得到.
小了一點.
盡力而為.
預備起.
以後以色列人中再沒有興起一位先知像摩西的.
他是耶和華面對面所認識的.
耶和華差派他在埃及地.
向法老和他的一切神僕.
以及他的傳遞.
行了各樣神蹟奇事.
又在以色列眾人眼前.
顯出大能的手.
行了一切大而可為的事.
從這段經文我們看到摩西這個先知的典範.
她是怎樣的呢.
她是面對面和上帝認識的.
摩西是面對面和上帝見面.
和上帝聊天.
而摩西其中一個工作.
這個經文沒有提及.
但其實經文很清楚.
摩西經常站在上帝和以色列民眾中間.
做一個很重要的工作.
就是祈禱.
摩西經常為那群正在吵架.
還半億的以色列民眾去祈禱.
求上帝憐憫.
求上帝審判不要臨到他們當中.

$^{241}$摩西是面對面神所認識的.
神認識摩西.
摩西又很認識神的屬性.
所以摩西又很清楚.
怎樣和上帝幫助以色列人代禱.
然後呢.
這裡提到.
她在埃及地.
神差遣她在埃及地.
做了什麼.
她行了很多神蹟奇事.
在這圖畫中.
神呼籲摩西的時候.
都是單對單.
十災.
神透過摩西.
施行十災.
那些全部都是神蹟奇事.
為的要讓埃及人知道.
哪一個才是上帝.
要放神的指紋.
去敬拜他們.
然後提到最後.
摩西這個先知典範.
在神的指紋當中.
在以色列人眼前.
同樣顯出很多大能.
大而可畏的事情.
大家想起嗎.
摩西舉起了像.
舉起了手.
紅海又分開了.
我們去想.
這個真是.
應該沒有什麼機會再發生.
我看過一齣電影.
開玩笑地說.
那碗羅宋湯.
分開了兩邊.
大家有沒有看過這齣電影.

$^{281}$搞笑的.
但前無古人後無來者.
神真的透過摩西.
施行了很多大而可畏的事情.
沒水喝.
苦水變甜.
沒水喝.
盆石可以出水.
大家有什麼缺點.
就說 曾老師你做個神蹟.
如果神真的叫曾老師這樣做.
絕對可以做.
我真的這樣相信.
在神裡.
豈有難成的事.
是不是這樣相信.
是不是真的這樣相信.
有些事情.
教會面對著一些什麼挑戰.
面對著一些什麼困難.
我們是怎樣去呼求上帝.
但同樣主權交給上帝.
OK.
一個這樣的典範.
我們知道.
剛才略略提到.
先知的工作.
在這段經文裡.
略略提到.
有一段對話.
這段對話是說昔日.
神和以色列人立約的時候.
在荷列山立約的時候.
一個情景.
當時上帝的榮耀顯現.
是轟天動地的.
在火焰當中.
整個說話出來.
地震山崩.
是很震撼的.

$^{321}$當時的以色列民.
是以色列民的領袖.
其實是很正常的.
明不明白.
可能我們還沒見過.
但如果我們遇到地震.
你都知道那種震撼.
遇到颱風.
那種風力.
大家最近有沒有.
深存一個很大的感恩.
那隻摩羯.
距離香港三百公里.
其實香港一樣會受災.
有些事.
我們.
可以怎樣去回應.
裡面那份害怕.
但對於上帝從信仰的角度.
看那份害怕.
是導致我們.
有一份對上帝的.
敬畏.
就很不同.
其實舊約原文.
敬畏那字其實是fear.
同樣都是害怕的意思.
有時我經常說笑.
我們怕老鼠怕蟑螂都不怕上帝.
是不是.
我不知道大家是不是這樣.
但在裡面我們看看.
這段經文說回那件事.
跟著他們的對話.
以色列民的領袖就跟摩羯說.
不如這樣我們真的很害怕.
我們很怕被燒死.
不如你在我們當中做先知.
做那個中介者.
在神和我們中間.

$^{361}$你幫我們傳話.
確實上帝都認同.
這個安排是好的.
是好的.
所以在裡面.
我們很清楚看到一個中介的角色.
先知.
在這段經文裡面再重新.
讓我們知道先知.
有一個中介的角色在當中.
做一個中介者的工作.
而先知.
他不是說他自己想說的東西.
先知是說.
上帝說了給你聽.
你聽一些什麼.
律例典章或者一些什麼說話.
你接收了之後.
你就要忠於地.
不加鹽巴.
加醋地去傳遞.
上帝要你傳遞的.
那個手續.
就是一個這樣的意思.
而聽的那群以色列民.
他們是要.
順服.
神跡先知所吩咐的說話.
大概這段.
中間這幾個經文.
讓我們很清楚知道先知的角色.
先知不是說.
他自己的東西.
先知是要忠於.
宣講上帝.
吩咐給他的說話.
上帝.
放了一些什麼說話.
在先知的口裡.
他就照樣這樣去吩咐.

$^{401}$告訴神.
要告訴哪些人聽.
就哪些人聽.
要聽從他奉我的命.
我必親自.
向他追究.
OK.
但是可能.
頂尖我有時都會問.
會不會有些人亂說話.
真的.
確實在我們現在這個時代.
確實有些人亂奉.
上帝的命去說一些.
不是聖經裡面所說的.
在我們這個時代裡.
真的有很多異端.
甚至邪教.
有很多假的先知站出來.
有一段時間很流行.
大家一定都聽過.
在東方會閃電的教派.
東方閃電.
我教會都曾經被滲透過.
好在我們後來發現.
拍了照片.
通知大家這班人要小心.
他真的有很多方法.
這段經文說的一件事.
若有先知擅自奉我命說.
我未曾吩咐他講的話.
或者奉別臣的命說.
才知就必處死.
這些人.
他不知道在上帝裡面.
是有嚴重的審判.
但是我們可能.
聽的時候我們心裡就問了一個問題.
我相信昔日會這樣問.
今日我們都會這樣問.

$^{441}$我們怎樣知道.
那些話是未曾吩咐他的呢.
當然我剛才說了一些.
我們可以.
看回聖經我們就知道.
聖經沒有說.
你憑什麼說.
他可能有一天上帝親自告訴我們.
童生你聽這些.
你要豎起耳朵.
是真不是假.
就算真的跟你.
我們認識了幾十年.
我們都知道大家.
來龍去脈.
特別一些陌生人.
突然間很熱心.
要豎起耳朵.
我們這麼近這麼方便.
門是打開的.
但這段經文有一個提示.
特別在先知的工作.
其中一個範疇.
特別講一些預言的東西.
就有一個提示.
提示是什麼呢.
先知奉耶和華明說話所說的.
若不是實若若沒有實言.
或不應驗.
這句話就是耶和華.
未曾吩咐的.
但先知擅自說的.
你必不怕他.
我曾經聽過.
一些頂尖會.
真的被人哄.
你不要亂說話.
你不要把什麼.
告訴教會的傳道人.
每逢你聽到這些.

$^{481}$你更加要找曾牧師.
找其他傳道人說.
有一個人好像很勤力在教會出現.
很熱心事奉.
但他這樣說話.
好像有點不對勁.
找一些團契小組.
職員去找.
當場去找牧者.
這裡說了一樣很重要的東西.
真正的先知.
時間可以證明一切.
而且假施之的下場.
上帝會審判.
其實在舊約聖經裡.
都有很具體的例子.
我提到其中一個例子.
在耶利米書.
大家知道耶利米一定是真先知.
聖經很清楚告訴我們.
但在他去作先知工作.
去宣講神的審判訊息.
宣講神安慰訊息的時候.
曾經有一些假先知.
其實耶利米很早已經說.
其實因為以色列人.
不聽神話.
結果最後審判.
就是一個.
聯合王國要分裂南北國.
後來北國被亞述所滅.
後來南國被巴比倫所滅.
這些我相信大家應該知道.
要回到一段日子.
大家知道.
我問大家.
神說.
南國被巴比倫所滅.
他們要服侍巴比倫.
多少年才有機會.

$^{521}$被老回歸.
我聽到了.
70年.
曾牧師流眼淚.
很開心.
突然有個人.
叫哈娜尼亞走出來.
說安慰訊息.
大家放心.
兩年內巴比倫會被滅.
所有聖殿被老到.
巴比倫的聖物會回歸聖殿.
我們可以在聖殿裡.
重新敬拜上帝.
你聽到的時候第一個反應是.
哈娜尼亞也是先知.
大家覺得會怎樣.
哈利路亞.
兩年內.
之前有人說過70年.
誰是真誰是假.
耶利米怎麼回應.
當時耶利米也在.
哈娜尼亞.
她說情願你所說的話是.
真的.
但其實不是.
耶利米就.
然後神就跟耶利米說.
你去跟她說.
她擅自.
說不是上帝所吩咐的說話.
大家知道.
哈娜尼亞最後的結果是怎樣.
她還沒等到.
兩年她自己死了.
還沒等到她所預言的.
兩年的回歸.
她自己已經被上帝處死了.
她的生命到了終結.

$^{561}$當我們剛才聽到.
舊約先知.
那個.
神怎樣出現.
神怎樣揀選.
她的身份角色.
她的工作.
我們有什麼反思呢.
我有兩個反思.
第一個反思.
在我們這個時代.
我們要好好地講好神的故事.
好好地講好神的話.
然後可能你會說.
黃母子.
我是不是要等上帝.
你手上有沒有聖經.
有沒有.
有沒有電話.
手上有沒有電話.
手機裡面有沒有聖經的app.
我們作為基督徒.
我真的很相信.
我們不可能沒有聖經在手.
不過看不看.
所以問題.
在我們作為一個時代的先知.
我們要好好地講好神的說話.
當然.
當我們要好好地講好神的話的時候.
當然我們要好好地.
讀好聖經.
教會有提醒大家.
頂尖會有互相激勵.
我們要靈修.
我們要讀經.
有主學要把握機會上.
有講座要把握去聽.
但可能黃母子.
我做完這些事.

$^{601}$其他事不用做.
總之按你的情況去選擇.
我們生命.
是需要在上帝裡面.
浸的.
其實我在這裡說.
宣講神的話語.
我有時都會問.
其實我的信仰生命.
是否真的靠那四年的神學.
其實那四年的神學只是拿個牌.
其實我過去的信仰歷程.
我很多事其實是由第一天.
一路一路累積.
對上帝的認識.
對神話語的認識.
生命當聽到神話有吩咐.
我生命去選擇這樣去跟.
這樣去改變.
是要用力的.
但有時力自己不夠.
但你相信聖靈在你裡面.
加力給你.
有時真的聖靈幫我們.
心意更生.
生命變化.
然後走到今天.
我.
八四年開始回教會.
我不知道何時信主.
因為我讀基督教學校.
舉過很多次手.
我都不知道哪一刻是真正.
但我知道我回教會.
一段時間我很確信.
我是跟耶穌.
由我回那間教會.
到如今.
我半個人生都是在那間教會.
由開始被栽培.

$^{641}$然後去事奉.
最後去做執事.
然後去奉獻讀神學.
然後回到自己的母會事奉.
直到如今.
四十年.
大家不要猜我幾歲.
但有句說話.
有些金句.
提示.
聖經全都是神所物事的.
我這齊讀.
我讀.
聖經全都是神所物事的.
對於教導,責備,糾正,訓練人.
行義都是有益的.
讓熟成的人得到充分的裝備.
可以行各樣的善事.
弟兄姊妹.
你信不信.
當然這裡所說的聖經.
很大部分是舊約聖經.
因為聖經有些部分.
都未成為正典.
但我們知道.
我們是很蒙福的一群.
時代的先知.
因為神完整的吩咐.
律例典章已經在整本的正典.
舊約和新約聖經裡面.
問題是.
願不願意講好神的故事.
還有怎樣講.
我分享我最近一個經歷的體驗.
我很喜歡年輕人.
我今天見到大家當中有年輕人.
很開心那邊有些年輕人.
其實我很喜歡香港的年輕人.
年輕人有.
創意,有活力.

$^{681}$有他們的看法.
但我年長的.
我經常鼓勵我們要支持他們.
我有三個小朋友.
都是很另類的.
有時會激動我噴火.
但我是很支持他們.
走他們要走的路.
我早前帶了幾個年輕人.
都是年輕人.
因為我帶了一個.
門訓的小組.
我都很久沒有帶了.
不過教會有需要.
近年都需要門訓.
我帶了八個頂尖妹.
三十多歲.
你們願不願意承諾一年.
每個月都要開早.
每個月都要見我一次.
要見一個小時.
我陪你們去找上帝.
要在裡面有很多要求.
最後一定有一個.
很重要的要預留時間.
大家的年假是怎樣安排的.
年假當然是香港人.
年假放假當然是去日本.
去泰國東南亞.
我說今次一年後.
你們要預留年假.
不過不是去玩那麼簡單.
是要有生命的挑戰.
結果我們在海米.
打大風.
打到幾百公里.
在台灣那天.
我們就在香港飛過去.
我預了炒魷魚.
我預了執線問責.

$^{721}$你這樣都去.
在過程中我祈禱上帝跟我說.
這段時間你不可以主導一切.
你讓我在頂尖妹.
生活裡面主導他們.
你陪他們.
他們說去.
你就陪他們去.
就是這樣.
結果我們在八號風球.
安全登陸.
經歷了五天很特別的事情.
其中我們又遇上了六位年輕人.
他們用了一種新的時代方式.
去講好聖經的故事.
他們每一個都是專業的人.
他們有些是原住民.
有些不是原住民.
但他們回到一個原住民的生活.
鄉下.
用他們的專業.
他們開了一個工作室.
在裡面.
他們有.
鹽和江工作室.
不好意思.
位置走了.
他們在工作室裡面.
有咖啡廳.
有背包客棧.
他們的咖啡廳背後有工作室.
在裡面他們有不同的主修.
英文音樂.
藝術.
在裡面搞觀光.
他們又設計了一些.
現在很流行的卡機遊戲.
桌上遊戲.
用很多不同的方法.
希望接觸回那班年輕人.

$^{761}$特別是在鄉下裡面.
因為他們的感情.
在他們的鄉下裡面.
大家有沒有聽過一個地方叫南莊.
就是我們教會.
差演出來的地方.
叫小學.
就是我們教會差演出來的地方.
宣教士他們搬去那裡.
做聖經翻譯.
搞一個共享空間.
希望去接觸那個鄉下裡.
那裡其實不是他們.
首先有其他人.
都在那裡工作.
這個建設我想說什麼呢.
上帝在裡面在不同的時代裡面.
在摩西的時代.
神有那種的作事方法.
神在這個時代有他.
另外去用我們的生命.
所作的奇事.
那些神蹟奇事.
為的是要讓人認識.
這位上帝是何等的奇妙可畏.
當你聽到.
年輕人這樣的時候.
你有沒有想過上帝.
放一些什麼在裡面.
用你的專業用你的生命.
去講好上帝聖經的故事呢.
創業很難的.
在分享當中.
他們說第二年就想結業.
第二年就已經想結業.
很困難.
在教會以外去做這些事情.
不容易.
但是我們去聽他們分享的時候.
他們已經八年了.

$^{801}$去見證.
曾經有一個報道裡面.
這個總監說了一句話.
他說工作是上帝的神蹟.
他沒有讓我們缺乏過.
讓我們有動力.
除了對南方的歸屬感.
還有對這地的使命感.
上帝把我們派在這裡.
就是要把.
這邊的事情傳承下去.
就是上帝給他們的.
對那個地方.
那裡的人.
那份感情.
那份歸屬感.
他們對上帝的使命感召.
他們被差派.
其實在城市他們可以有很好的發展.
最無大學生.
有尊嚴.
但是神就要差派他們.
去承傳.
這位陳總監.
陳姐妹.
她爸爸.
當時是這個莊裡面的一個牧師.
牧師是另類的牧師.
他做了另類的事情.
他就是透過音樂.
透過師班.
去建立那間教會.
去傳揚福音.
當然他爸爸離開的時候.
有些事情就轉變.
但他們承傳.
這班年輕人.
他們用了另一種方法去承傳.
而承傳的那件事.
交給他.

$^{841}$那個大使命.
都是傳福音.
如果教會不傳福音不如結束它.
如果我們的生命不去見證上帝的榮耀.
不如你接手它.
年輕人可以有這樣的勇氣.
如果你聽到一些年輕人.
這樣分享的時候.
我想特別和在座.
連接你們有能力的.
你們可不可以支持他們.
無論實質的支持.
禱告的支持.
信任的支持.
你們可不可以支持他們.
第二個.
我們才知道.
當神揀選的時候.
他會見到很多神作的奇事.
當見到的時候.
不可以只放在自己心裡.
是要見證出來.
所以第二個.
給我的反思.
我都希望和紅印堂的弟兄們一起去反思.
我們要好好地去見證.
我們所看見神的作為.
大使命要我們出去.
當我們出去的時候.
我們會見到你.
我們不要只在教會的封牆.
我為了紅印堂感恩你們有地舖.
這段時間我知道牧師一直在用很多方法.
善用地方.
去接觸街坊.
中秋節剛好有個活動.
我也想過來見識一下.
實在是耽誤.
但我知道.
曾牧師經常去分享.

$^{881}$去激勵屯源天的童工.
我也為你們去感恩.
善用神給你們的土地.
善用神給你們的資源.
為的是要見證.
耶和華是上帝.
而在裡面.
所見到什麼呢.
我今天見到紅印堂.
都幾平均.
也有些練長.
在這次裡面.
我見到.
有一班長者.
好型.
他們穿的衣服.
有些logo.
有些徽章.
他們叫基督徒騎士團.
那個是一個.
什麼情景之下.
我們那天去一個山上.
的教會去崇拜.
參加一個原住民的教會去崇拜.
那間教會.
兩年前只有五個人在崇拜.
那個傳道.
江地頭的傳道.
叫油幹傳道.
他在那天的崇拜裡面去分享.
感覺他被拆敗的尼因教會.
是一種被鞭的感覺.
我很欣賞.
他這麼真實地說出來.
但他這種感覺.
他有沒有放棄神.
選擇他在山上原住民教會裡面.
作先知的角色.
作牧者的角色.
作牧羊的角色.

$^{921}$他有孤單的時候.
他那天唱了一首歌.
就是他最孤單最困乏最無奈的時候.
上帝給了他一首詩歌.
安慰他.
給他力量.
在這兩年你知道有什麼變化.
他現在有十幾人聚會.
但當日我們崇拜的時候.
有百多人在山頭裡面.
露天的教會.
踏上去在山頭裡面.
百多人在那裡敬拜.
你想像一下是很震撼的.
在神所創造的大自然裡面.
敬拜的聲音響起.
禱告的聲音響起.
周邊的人有機會.
去聽到.
在裡面敬拜的人.
好像和上帝很近的一個情景.
因為他領受了一件事.
他說這個地方這麼特別.
要讓更多更多的弟兄姊妹.
可以一起在這裡敬拜上帝.
結果在這兩年裡面.
有超過二千人.
曾經在這裡崇拜.
是不是很奇妙.
沒有想過.
但來敬拜的人是做什麼的.
那天我看到有兩個群體.
一個群體就是.
他們崇拜裡面的一間教會.
一群弟兄姊妹來.
大敬拜.
好像剛才C班那樣.
他們有一個全隊.
來大敬拜.
另外有一個派對.

$^{961}$就是基督徒騎士團.
全部都是退休人士.
但他們非常有型.
駕駛機車.
台灣的重型電單車.
駕駛幾十公里甚至過百公里.
有些在基隆泥有些在高雄泥.
未天亮就開車.
我們上山時發現.
那山發生什麼事.
為什麼幾十輛電單車.
十幾輛私家車.
最後發現還有一輛.
三噸半的貨車.
就是這班退休的弟兄姊妹.
他們帶上所有音響器材.
所有混音器材.
所有麥克風.
用三噸半的貨車運來.
打好排行.
預備好一個敬拜上帝的地方.
他們獻唱又唱詩歌.
我裡面.
我最近越來越喜歡.
除了我喜歡年輕人.
我還很喜歡將會退休或退休的弟兄姊妹.
其實你有沒有想過.
不要想著退休就退.
我發現原來.
上帝要成全的要揀選的.
不是揀選年輕人.
為什麼你們沒有聽過一個觀念.
我們經常說成全就是年長就交給年輕人.
你們是不是這樣聽的.
好像社會都是這樣說.
但我最近在說不是這樣.
成全就是.
將使命.
去成全.
上帝感召的年齡.

$^{1001}$不論大小.
總之他回應上帝的感召.
回應上帝的使命.
管你多少歲.
我作為牧仔一定支持.
我相信曾牧師一定會支持大家.
教會一定會支持大家.
有沒有想過.
香港沒有很多重機車.
但我知道.
黃絲踩單車的.
我有班非基督徒的單車朋友.
我在踩單車的過程中.
我見到很多退休人士踩單車.
有沒有想過我們組織一個單車隊.
基督徒騎單車團.
在元朗.
新界.
周圍.
遇上這些人.
現在有一班.
跟我踩單車的人不是基督徒.
但他們每次跟我踩單車.
一定要給我祝福.
為什麼呢.
每次吃飯時我都會說.
讓我祝福一下大家.
他們很樂意的.
現在他們還在啞門.
我沒有教他們.
總之我一吹哨子.
馬上回應牧師何時進去.
我剛剛.
我剛剛吹哨子.
明年.
我去台灣環島.
大家有沒有聽過台灣有一個叫.
不老騎士.
有吧.
其實.

$^{1041}$我自己從熟人的角度.
什麼年紀我們都可以回應.
大使命的召命.
你看到一些什麼見證呢.
你看到一些神秘的作為呢.
記得要分享.
分享時會激勵人心.
我這樣分享會不會激勵到你.
馬上想一些事.
神作事是可以奇妙的.
不說那麼多了.
但見到的東西.
記得回來告訴鄧英文.
結果那班悶訓的組員回來.
在自己團體分享了三個小時.
都說不完.
因為那五天都差點死.
暴風雨都死不了.
結果去到那些地方很多經歷.
整條鄉都知道.
有十幾個香港人.
不知道發生什麼事.
暴風雨又水浸又山崩.
走來我們這條鄉.
因為我們有兩個朋友.
香港人會住在你們當中.
然後他們告訴我們一些秘密.
其實我們最近.
縣政府在附近會建一個樓盤.
你們不如去買.
那裡是秘密消息.
會告訴你一些這樣的事.
已經做好朋友.
他叫不要告訴別人.
(笑).
OK.
舊了先知.
神在那個時代.
他所揀選的先知.
在那個時候作成他要作的事情.

$^{1081}$在今天我們這個時代.
神同樣揀選.
你不要用舊了先知來引用我們.
我們講好聖經的故事.
我們見證好上帝的作為.
就已經很.
很神奇.
我下台了.
謝謝大家.
\newpage



\section{尼希米記 8:1-12}
\label{sec:dufzJl9BQpc}
\textbf{2024.09.29 《建立讀經群體》 尼希米記8\_1-12 (和修版) 講員 : 蔣文忠傳道}
\newline
\newline
連結: \href{https://youtube.com/watch?v=dufzJl9BQpc}{\texttt{https://youtube.com/watch?v=dufzJl9BQpc}} ~~~~ 語音日期: 2024-10-01
\newline
\newline
\hyperref[sec:upkPX3gy4Ao]{\small{< < < PREV SERMON < < <}}
~
\hyperref[sec:index_chronic]{\small{[返順時目]}}
~
\hyperref[sec:index_scriptual]{\small{[返順卷目]}}
~
\hyperref[sec:8mvS2zOfCvg]{\small{> > > NEXT SERMON > > >}}
\newline
\newline
尼希米記 8:1-12
\newline
\begin{longtable}{cl}
\hline
\hline
章節 & 經文 (和合本修訂版)\\
\hline
8:1 & \begin{tabularx}{0.7\textwidth}{X} 那時，眾百姓如同一人聚集在水門前的廣場，請以斯拉文士將耶和華吩咐以色列的摩西的律法書帶來。 \end{tabularx} \\ \\ \relax
8:2 & \begin{tabularx}{0.7\textwidth}{X} 七月初一，以斯拉祭司將律法書帶到聽了能明白的男女會眾面前。 \end{tabularx} \\ \\ \relax
8:3 & \begin{tabularx}{0.7\textwidth}{X} 他在水門前的廣場，從清早到中午，在男女和能明白的人面前讀這律法書，眾百姓都側耳而聽。 \end{tabularx} \\ \\ \relax
8:4 & \begin{tabularx}{0.7\textwidth}{X} 以斯拉文士站在為這事特製的木臺上。站在他旁邊的有瑪他提雅、示瑪、亞奈雅、烏利亞和希勒家；站在他右邊的有瑪西雅；站在他左邊的有毗大雅、米沙利、瑪基雅、哈順、哈拔大拿、撒迦利亞和米書蘭。 \end{tabularx} \\ \\ \relax
8:5 & \begin{tabularx}{0.7\textwidth}{X} 以斯拉站在上面，在眾百姓眼前展開這書。他一展開，眾百姓都站起來。 \end{tabularx} \\ \\ \relax
8:6 & \begin{tabularx}{0.7\textwidth}{X} 以斯拉稱頌耶和華至大的神，眾百姓都舉手應聲說：「阿們！阿們！」他們低頭，俯伏在地，敬拜耶和華。 \end{tabularx} \\ \\ \relax
8:7 & \begin{tabularx}{0.7\textwidth}{X} 耶書亞、巴尼、示利比、雅憫、亞谷、沙比太、荷第雅、瑪西雅、基利他、亞撒利雅、約撒拔、哈難、毗萊雅和利未人使百姓明白律法；百姓都站在自己的地方。 \end{tabularx} \\ \\ \relax
8:8 & \begin{tabularx}{0.7\textwidth}{X} 他們清清楚楚地念神的律法書，講明意思，使百姓明白所念的。 \end{tabularx} \\ \\ \relax
8:9 & \begin{tabularx}{0.7\textwidth}{X} 尼希米省長、以斯拉祭司文士，和教導百姓的利未人對眾百姓說：「今日是耶和華－你們神的聖日，不要悲哀，也不要哭泣。」這是因為眾百姓聽見律法書上的話都哭了。 \end{tabularx} \\ \\ \relax
8:10 & \begin{tabularx}{0.7\textwidth}{X} 尼希米對他們說：「你們去吃肥美的，喝甘甜的，有不能預備的就分給他，因為今日是我們主的聖日。你們不要憂愁，因靠耶和華而得的喜樂是你們的力量。」 \end{tabularx} \\ \\ \relax
8:11 & \begin{tabularx}{0.7\textwidth}{X} 於是利未人叫眾百姓安靜，說：「安靜，因今日是聖日，不要憂愁。」 \end{tabularx} \\ \\ \relax
8:12 & \begin{tabularx}{0.7\textwidth}{X} 眾百姓去吃喝，也分給別人，都大大喜樂，因為他們明白所教導他們的話。 \end{tabularx} \\ \\
[1ex]
\hline
\hline
\end{longtable}
$^{1}$各位兄弟姐妹平安.
主席的恩典也讓我能夠和大家一起.
今天一起敬拜.
一起分享上帝的華語.
主席剛才和我們讀出了一段尼希米記第八章的經文.
可能我們未必很熟悉第八章的經文.
或者我們會很熟悉尼希米前面.
帶領著一些百姓回歸到耶路撒冷重建城牆.
如果大家更熟悉以色列的歷史.
也知道以色列尼希米記.
其實就像上一集一樣.
先有索羅巴伯和以色列帶領著第一批早期歸回的以色列人回來.
先重建聖殿.
讓這些百姓回到耶路撒冷的時候.
可以有像大家一樣舒服的環境.
有電.
去敬拜.
獻祭等等.
到第二個階段.
大家讀過書也聽過.
是一個工程師帶領各個工匠.
各個製造業的人.
弟兄姊妹.
一起重建城牆.
重建城牆重要的是讓人們在當中安居樂業.
因為城牆抵禦外敵.
讓他們在城內安穩生活.
為什麼會發生這些事呢?.
大家都知道歷史.
因為以色列人被上帝審判.
被捕七十年.
回到巴比倫其他異族的地方.
七十年.
然後上帝讓他們一批一批回到自己的地方.
回到巴比倫時我們要問.
他們要建立什麼呢?.
像我今天說的題目.
建立讀經群體.
他們不會用這些字眼.
我之前也有推動一些教會同行的時候.

$^{41}$也有說教會的核心是什麼呢?.
可能就是用讀經群體.
一群人一起都很願意讀聖經.
去思考教會的定位.
最重要的價值等等.
但當時不會這樣說.
但他們也會想.
這群人回來了.
回到一個荒涼之地的時候.
一個他們故鄉的時候.
最重要的是做什麼呢?.
他們的硬件都做好了.
又有聖殿.
又有城牆.
到了尾八章重要的是.
一群整個群體.
如果剛才讀經的時候.
眾百姓如同一人.
在七月初一的時候聚集.
在耶路撒冷的水門前.
然後才讀了很多名單.
讀名單大家可能會覺得很膽怯.
每個名字都要讀.
但我們沒有這麼深入地跟他們的關係.
但如果我們想深一層.
這些名單其實都是很感動的.
就是說每一個部分.
每一個事奉人員.
我們有時候教會.
年中年初的時候.
都會列出一些事奉名單.
其實那個名單是什麼意思呢?.
先有民事以色列的同伴.
那些名單站在民事以色列左右.
站台.
好像跟他們有一種.
有時候你看到很多人.
在這裡選舉.
都會有些人站台.
表示一個加持.

$^{81}$一個支持.
或者是一種同伴.
那是一班民事的團隊.
有時候大學就說.
Teaching Faculty.
一個講師.
一個教學的團隊.
就是第一次讀的名單.
再讀下去的時候.
另一個名單在第七節的時候.
那個是利未人的名單.
是一班助手.
一班教師的助手.
去到百姓中間.
做分組討論.
如果大家剛才讀的時候.
你很留心時序.
其實剛才讀的經文.
不知道教會有沒有辦一些讀經日營.
有沒有?.
我們沒有會很流行的.
就是公教就會搞讀經日營.
但是他們想象.
就是一個讀經日營的流程.
經文在第三節說.
從清早到中午.
就是早上的Session.
早上就聽到.
就是民事師奶就說到.
下午的Session.
下午就是剛才那班利未人.
就帶分組查經.
你可以想象.
在第七節開始.
去到百姓中間.
到剛才第九節開始.
黎謙蕊出場.
就應該肚子餓了.
過了一整天就要吃飯了.
所以就沒東西吃的就分給她吃.

$^{121}$但是可能他們都沒有預料過.
這麼長都不等.
不知道時間安排.
然後就一起哀怨的時間.
是一個一整天的流程.
為了建立這班歸回的子民.
在上帝的話當中.
重新讓他們可以有宣講.
有聆聽.
有教導.
有分享.
重建一份和上帝的話.
很重要的關係.
所以就和我今天這個講題.
有一個Echo的地方.
對於我們來講.
其實我們教會.
有沒有這麼重視上帝的話呢?.
求主幫助.
我們一起低頭禱告.
然後就進入今天的經文信息當中.
主多謝我們今天.
能夠在這個聖日當中.
就好像剛才讀經的裡面.
黎謙蕊提醒百姓.
今天是主的日子.
是主的聖日.
我們分別出來歸向給你.
敬拜你.
讀你的話語.
專注質疑.
聽你的話.
以致聽了我們能夠明白.
以致能夠進行.
在生命當中.
我們有走錯路的.
求你施恩憐憫寬恕我們.
以致我們歸回主你的道路當中.
讓你繼續牽引我們.
無論我們整個教會群體.

$^{161}$我們每一個弟兄姊妹的生命.
都在你的手中.
願你保守我們以下的時間.
同心祈禱奉主耶穌得勝名字.
很快說一些前言.
我們剛才主席介紹了FES.
我們是一個做學生機構的組織.
學生科學團體裡面.
我們在六月的時候.
搞了一個青少年牧羊研討日.
是一個端午節日.
是星期一.
都大掙扎的.
因為全段放假.
會不會要放假就不參加呢.
我還要整天像剛才那樣.
整個日程.
早上下午都有環節的.
很感恩我們有超過一百十幾位.
一些傳道.
主要是清晨的導師一傳道.
來參加.
感酸也有弟兄姊妹一起參與.
我們在問的時候.
其實我們除了傾清青少年牧羊上面.
今天的何去何從之後.
我們也問在這個時候.
這個時勢裡面.
我們在香港繼續牧羊教會.
究竟是為了什麼而做呢.
我們可能突然間面對很多的衝擊.
很多的人才資源的流失.
但是這些外在的改變.
對於我們內心.
對於我們教會上帝面前的領受.
有沒有不同呢.
我們就問了.
什麼是不能夠缺少的呢.
對教會.
是不是地方呢.

$^{201}$是不是聚會呢.
還是一群人呢.
一群人裡面的關係.
他們走在信心之中.
才是最重要.
我們走一層.
好像一路拆解那樣.
去問自己.
問屬於教會群體當中.
究竟我們那群人.
在做什麼呢.
追求什麼呢.
一個很重要的自我檢視.
我們再問.
我們教會要做什麼服事呢.
我們多搞聚會.
剛才說了.
有團契.
我們有不同的部.
團契.
就是分享相交.
我們固然也有詩歌敬拜.
這些都很重要.
但是原來我們有一個很核心.
就是一個叫聖道的即時.
The Ministry of the World.
這個字.
這個稱謂.
其實由我們中國革命運動開始.
由馬丁路德加爾文.
當時十六世紀.
和當時所謂大公教會.
天主教會.
有一個分道揚鑣.
一個最核心的問題.
他也問這個問題.
究竟做什麼服事.
教會.
真正的教會在哪裡.
他們就推動一件事情.

$^{241}$其實是一個.
回到上帝的道當中.
他們當然說.
聖道和聖禮.
都要被宣講.
聖道.
聖禮就要很合宜地執行.
那裡就是教會.
就是一個以上帝被高舉.
他話語被宣講.
流通成為弟妹之間的那種分享.
和交流的一個內容.
這個地方那裡的人.
才是教會.
這個很重要.
帶回我們常常想起.
我們服事裡面.
我們做很多的活動.
做很多的群體聚會.
很多不同的分享裡面.
我們有沒有忽略了.
回到道.
或者基礎裡面.
是要在道當中.
建立我們一切的服事呢?.
我和大家去.
講今天這段經文.
主要我不會整段都很詳細地講.
我會抽取幾個重點.
和大家去分享.
第一個就是.
關乎到.
剛才說了.
這批歸回的人.
特別紅色.
和大家都有提示.
聽了能明白的會眾.
他怎麼形容這批百姓.
這批站在水門面前的.
這批人呢?.

$^{281}$他說他們是一群聽了.
會明白的人.
然後去到再講.
第三節那裡.
剛才說了清晨到早午.
有一個講道.
再形容一群男女.
男女和能明白的人.
去讀這個律法書.
然後再多一句.
怎麼再講聆聽的東西呢?.
就側耳而聽.
今天大家坐得很公正.
他突然側身.
耳朵聽.
我也有點害怕.
我們不習慣這樣做.
但當時那群人是.
你記住他們70年沒聽過上帝的話語.
是不是?.
搞完一大頓飯.
又弄電又弄.
現在終於有個地方了.
聚集到.
真的好好的.
有一天分別出來.
聽上帝的話.
所以再讀下去.
你會看到後面的訛言時間.
很激動.
又哭又笑.
我第三節會講一點.
但我第一節跟大家分享的是.
這群人是很願意去聽.
側身去聽.
很渴望.
很想明白上帝的話語.
這是第一點.
聆聽以致能夠明白.
大家有沒有看過聆聽.

$^{321}$今天我們經常都會說.
聆聽很重要.
因為其實我們用耳朵聽.
但大家常常都聽到一些人.
譬如一些男女爭吵.
都說你不要聽我說話.
你左耳入右耳出.
然後突然間聽.
可以是一個很重要的話.
留在心裡面也可以.
可以是一些音波.
震動你的耳膜.
然後走了也可以.
你可以很多角度來.
科學角度也可以.
有一個很真誠的話的內容也可以.
你可以有很多原因.
很多藉口去講聽話.
但是.
聆聽是很重要.
聆聽是我們要聽了.
先去回應.
先去批評.
那些是後話.
所以我們有時候聽港獨學.
所以我們學港獨.
我老師第一個就說.
是聆聽先於宣講.
問港獨者.
首先有沒有聽到上帝的話語.
還是你急急腳就講你自己想講的話.
這是對於港獨者一個很大的操練.
原來我們台上台下的人.
是一起在聽上帝的話.
不是我聽了.
然後我消化完.
有些好東西就給你.
不是這個意思.
而是一起在那個港獨的時間.
一起在聽上帝要講話.

$^{361}$那你再問.
不是啊.
都是我說的.
是不是.
那怎樣能夠對於我們的.
每個弟兄姊妹.
都能夠有一個對應呢.
那這個就去到一個三一上帝的工作.
我們要相信聖靈.
在我們每個人的生命工作.
事實上.
就算牧師上來港獨.
去熟習大家.
都未必會知道每個禮拜.
大家發生很仔細的事情.
究竟這個禮拜我講什麼道.
對應到每個弟兄姊妹的生命.
是很難的.
沒什麼可能.
但我們相信.
就是講道的時候.
或者聽道的時候.
道要被高舉的時候.
聖靈在我們每個人的生命工作.
以致那個道安置在我們內心.
成為我們每個弟兄姊妹個人.
我們是一個個人.
我們每個人自己要聽的東西.
有一個客觀的信息.
轉化為一個你主觀的領受.
這是一個很重要的.
你聽完.
接著就明白了.
然後就是要去那個禮拜當中.
切實去進行.
當行的道路.
所以其實很簡單.
是不是.
我們平常都會找一些很基本的教理.
從耳朵聽.

$^{401}$用心去明白.
不一定要用腳的.
我覺得整個生命都很重要.
有時候我們言行舉止.
價值觀是決定我們的行動.
那裡有沒有衝擊我們整個人的優先次序.
價值觀.
我覺得有時候是更加寶貴.
更加重要.
所以我就問大家.
我們有沒有如此三步曲.
進入上帝的話.
這樣學習.
進行他的指示的道路.
還是我們停在第一步.
是不是.
《新老經》也說.
耶穌說.
道路落在路邊.
轉頭被烏鴉帶走了.
轉頭落在樓梯被美食帶走了.
顧著吃飯.
都不記得今天聽了什麼.
只停在聽那裡.
聽很多.
聽完就以為自己做了.
其實不足夠.
還是停在頭兩個步驟都不行.
你只是聽和不明白.
可能你上了主一學.
你是主一學出席最高拿勤道獎的學生.
你看很多書.
你頭腦裝了很多知識.
但是你只是停在一個知識的層面.
理性的層面.
你說了很多道理出來.
但卻沒有讓道挑戰我們生命最核心的地方.
所以是要三步.
聆聽.
明白和跟從.

$^{441}$才完整地讓上帝道得以被高舉.
上帝道的工作才叫做完滿地實現.
求主憐憫幫助.
如果我們停在第一階段.
求主幫助我們.
起碼去到第二階段.
如果只是第一第二階段.
求主給我們更多信心和勇氣.
更多的歸回他的心.
去到完整地走出來將他的說話.
去到第二個的分享.
除了聆聽.
明白.
跟從.
這是當時的會眾重新學習上帝的話語.
第二點想和大家分享的.
除了早上是一個講道的環節.
下午其實是一個分組茶經.
不知道大家是否流行分組茶經.
一會兒我會說一些大學的工作.
經常都會分組茶經.
基本上是FES一個核心部門.
沒有茶經就沒有學生工作.
很多地方.
無論香港或其他世界各地.
學生團契的興起.
都是有老師或師兄師姐.
帶著師弟同學茶經開始.
才發展成為很多團契.
很多校園,傳福音,宣教工作等等.
核心是從凝聚在上帝道面前開始.
剛才說了這裡有第二個名單.
是一班利未人.
這班利未人是祭司的助手.
你記住.
以色列是兩頂帽子.
又是文士又是祭司.
好像一會兒會有出場.
兩頂帽子一起.
當他為著百姓宣講教導的時候.

$^{481}$他就是文士.
但他是獻祭的鐘堡.
讓上帝的教導提示.
傳遞給以色列人.
或者代表著以色列人.
將有需要的獻祭感恩代求.
贖罪等等獻給上帝.
他也是中間的角色.
這是祭司要做的職責.
利未人就是幫祭司的助手.
當他教導元律化的時候.
這班利未人也會出場.
但他們想像這個安排.
不是隨心而發地就這樣做.
但他們想像文士.
以色列是有準備的.
可能早就跟他們講解了.
下午你們要將我早上講的道.
再跟他們分享討論.
然後就說誰去哪裡.
百姓站在自己地方.
第八節.
第七節都站在.
不動百姓.
因為很多人.
但起碼有幾千上萬人.
是利未人很有秩序.
知道怎樣安排之下.
去到他們中間.
可能一個對幾十個.
來分組交流上帝的話.
他們任務做什麼呢.
紅色的代碼.
使百姓明白律法.
然後再說.
他怎樣做呢.
他再講解明白意思.
使百姓明白所念的律法.
是不是以色列有問題呢.
講得不清楚.

$^{521}$要再講明呢.
其實不是的.
這個是一個很重要的任務.
如果今天那些主學老師.
小組組長.
去輔助牧者上來港島的任務.
這個講明意思是什麼意思呢.
我跟大家分享一下.
這個講明這個字.
在原文來說.
最核心的意義是translate.
就是翻譯.
可能由外語翻譯到希伯來文.
他們回到猶太不需要翻譯.
我們不需要.
英文譯成廣東話.
廣東話普通話譯成廣東話等等.
這些過程.
基本上全部都是自己人.
第二層的意思.
就是那個explain.
或者power phrase.
是一個更加技術性的字.
power phrase就是將一些很複雜的語言.
再用你的語言.
去講給一些不懂專業的人.
聽得明白.
講什麼.
你講得到.
你就power phrase了.
其實是另一種的恩賜.
是一種教學.
傳遞.
信息.
講解.
溝通.
的恩賜.
所以做組長是很難的.
做查經組長是更加大的挑戰.
但是很重要.

$^{561}$這件事很重要.
也在回應.
剛才說路德加爾文在十六世紀的時候.
他除了高舉.
聖道即是之外.
都要work(行動)的.
都要做的.
他另外再說路德.
大家熟悉的信徒皆濟私.
那為什麼呢?.
不是每個人都在殿裡面獻祭嗎?.
大家要明白.
當時已經是十六世紀.
不是留在舊約時代的那種環境.
其實意思是說.
你再讀多一點路德的時候.
你就會明白.
原來是叫他們每個信徒.
都可以參與在聖道即是當中.
突然間想一想.
你明白了.
十六世紀是一個什麼世代呢?.
印刷術已經興起了.
拉丁文聖經.
已經譯作特文.
透過馬丁路德.
他們可以拿著一個.
好像我們那間珍珠和合本和修本.
你可以想像這樣的情況.
有自己語言的聖經.
而且已經是mass production(全面製作).
每個人手上都有一本.
和以前的十幾個世紀裡面.
大家都看尼沙.
只看著神父讀拉丁文.
完全不知道他說什麼的時候.
另一個世界.
所以路德在這樣的時代之中說.
信徒皆濟私.
每個信徒都可以拿著一本聖經.

$^{601}$一起分享教導.
所以到今天.
我們都在這樣的後繼之下.
繼續推動教會.
我們一起推動群體當中去讀經.
我們不需要手上有聖經.
我成長的時候.
我仍然是拿著硬碟.
家裡還有五六本聖經在手.
好像我這個年紀.
或者比我年紀更大.
才會知道家裡有多少本聖經.
你明白我的意思嗎?.
今天過渣,APP那些不用.
但是其實是一樣的.
就是說你已經是隨手可得.
大家知道什麼叫隨手可得.
隨意拿著一本聖經.
問題是我們怎麼讀.
我們有沒有珍惜去讀.
資源都在身邊.
我講聖經.
當然還沒算那些講座.
識經書.
所有幫助你讀聖經的人.
物件,書本等等.
你都是隨手可得的.
你有心去.
就不會找不到.
問題是我們有沒有一個意象.
有沒有一個負擔去推動這件事.
我和大家分享一點我的經歷.
這個就是.
剛剛上學期.
上年還沒完之前.
大概三四月左右.
在大學裡面一個團契的周會.
他們每年都一定有.
一起查經的周會.
我做什麼呢?.

$^{641}$通常在查經周會裡面.
我很少出場.
我通常都是最後那十分鐘.
才出來做一點總結.
保證大家查了一個多小時.
都沒有失誤.
都不知道去了哪裡.
或者離經半導.
那些而已.
保證大家在那個框裡面.
繼續讀上天的話語.
我就完了.
但是更多的功夫在哪裡呢?.
在這個周之前.
晚上和一群組長.
三更半夜他們pre-study.
通常我就在zoom裡.
和他們在一起.
一個多兩個小時.
去做一次那個過程.
一起體驗上帝的話.
其實是一起去經驗上帝的話.
認識上帝.
和分享生命的過程.
然後他們就用他們的語言.
化為一些問題.
就帶回一個晚上的實體的聚會當中.
帶同輩一起去經驗.
在那個預查pre-study經驗的過程.
通常他們查的時候.
他們年輕人會怎麼說呢?.
我就會問.
他會突然間說.
我get了.
我就會問.
你get了什麼呢?.
就會問.
然後就會說.
說著說著.
說出自己get的東西出來.

$^{681}$然後就說.
好啊.
你這樣get的時候.
你可以再想多一點.
就是一個和他.
將他想到.
就是我傳遞給他的時候.
他想到想到的東西.
說出來.
然後再確認他.
想到的東西.
是不是真的合乎聖經裡面.
或者是查考過程裡面當中.
那個訊息的意義.
是很重要的.
是一群人一起.
伏在上帝的道的中間.
權柄裡面.
所以我經常跟同學說.
如果你肯做查經組長.
你是最有利的.
因為你學的東西最多.
你查完之後.
get了一些東西.
剛才get了嘛.
然後再帶的時候.
你透過同學再查.
你會再get多一點.
是一個不住地認識上帝說的話的過程.
是一個不住地盡心.
跟上帝的話建立關係的過程.
話說回來.
我經常用我那本我信讀經的書.
整個training manual.
給他們拿著手一邊學.
其實一些對聖經和對於讀經的神學反省.
我就會用我的書.
來跟他們做一些材料.
引導他們.
不會說.

$^{721}$五個同學都在不同教會背景.
都不知道教會裡面教了什麼.
我就有時候用一些比較.
我自己做了的一些整理.
跟他們分享.
讓他們都能夠同步地去學習.
對上帝話語的即時.
這個就是教會.
這個就是一間.
剛剛幫他們帶來的讀經日營.
教會跟大家分享的經驗都挺特別.
大概兩年前左右.
教會找我聊天.
牧師.
看了我那本書.
然後很想推動建立讀經群體.
在教會當中.
但他說很傳統.
也不知道怎麼開始.
他有一個.
就想膽小膽大地搞一個營會.
那次的營會我就幫他多一點.
他叫了一班組長出來.
我就用幾晚時間跟他們.
首先跟大學生一樣.
跟他們pre-study.
十個八個左右.
其實教會不是很多人.
八十人左右.
去camp大概五十人左右.
其實都算挺好的.
大概有六七個組長.
每組大概七八個人.
就跟他們聊天.
分享.
也用我教材的材料給他們.
明白多一點讀經.
或者是其實要追求什麼目標.
很感恩就是.
到了今年.

$^{761}$這個是今年一張照片.
五月的時候.
我比較忙碌.
我不有空跟他們.
再這麼深度的幫助.
我只能夠來一天.
他就說好啊.
那你最後一天.
你過來跟我們做個總結.
沒問題.
接著功夫怎麼做呢.
牧師經過一年多兩年.
不如我跟他pre-study.
好啊.
牧師你也願意做多一點.
有信心做多一點.
於是.
本來我跟他做的.
很多時間的部分.
他們的牧者已經做到了.
然後我回來.
最後一天就跟他們.
再交流一下.
整個營會當中.
他們讀上帝話語的過程.
我感恩的地方是什麼呢.
我發覺他們.
過了兩年之後.
那些小組討論的內容.
是不同了.
怎麼不同呢.
其實我覺得他們.
好像我這麼說.
其實他們能夠.
學習實踐上帝的道.
在生活之中.
然後就會.
分享裡面就是說.
我得有什麼掙扎.
還是我做到.

$^{801}$然後感恩的分享.
上帝的話真的成為了他們.
整個群體當中.
無論是一個.
分享裡面的一個核心.
一個重心.
還有每個弟兄姊妹.
都是學習得到.
是哦.
原來讀經.
群體讀經裡面.
最後是要將上帝的話.
應用在生活當中.
或者回應生活裡面.
一些我自己的習慣.
有沒有一些.
其實我們是不願意.
跟隨上帝的話.
或者是我們仍然很想.
留戀自己的方式.
然後將這些經歷和掙扎.
去分享.
我說很不同.
這次你們已經不是.
只是說茶餘飯後的事情.
只是說大家什麼時候去旅行.
什麼時候在這裡做什麼.
打一下golf.
或者其他遊戲.
他多說了自己.
生命當中一些核心的東西.
然後就說上帝的話語.
給了什麼提醒.
我為地教會感恩.
逐漸向一個.
真實的讀經群體.
一步一步地向前走下去.
求主幫助.
祝福.
紀念我們的工作.

$^{841}$這個也是很重要.
無論是在大學生的群體當中.
在教會群體當中.
如何可以不斷地更新他們.
對學習上帝的話語的文化.
和那種習慣.
這個也是不容易的.
但是有弟兄姊妹.
有牧者.
或者有些領袖.
帶起的時候.
就逐漸地整個群體.
這樣去推動.
這樣去改變.
去到最後的第三點.
剛才說了.
這裡就是到了黃昏時間.
是不是?.
尼熙是工程師.
所以剛才讀經的都沒什麼出場.
但是到了最後的時候.
尼熙米拉.
這麼重要的領袖.
以斯拉文士祭司.
剛才說了兩句.
都說完了.
還有剛才那些利未人.
那些組長.
都出場了.
就是跟他們說.
今天是耶和華你們上帝的聖日.
不要悲哀.
也不要哭泣.
是不是?.
這是因為眾百姓.
聽見律法書都哭泣.
有個原因的.
他們哭泣.
聖經說得很清楚.
因為眾百姓聽了律法書而哭泣.

$^{881}$然後利美出場了.
你們去吃肥美的和甘甜的.
沒有預備的就給他們.
因為今天是主的聖日.
不要憂愁.
因靠著耶和華而得的喜樂.
都是你們的力量.
不過我們一起讀最後兩節.
十一節開始.
十一,十二節.
預備.
一,二,三.
於是利未人叫眾百姓安靜說.
安靜.
因今天是聖日.
不要憂愁.
眾百姓去吃喝.
也分給別人.
都大大喜樂.
因為他們明白所教導他們的話.
為什麼要讀最後呢?.
因為他們大大喜樂.
是不是因為有東西吃?.
不是.
不是因為有東西吃就大大喜樂.
是因為最後一句.
他們明白了.
他們所讀.
到現在才明白.
為什麼不是一早明白了?.
你想想.
明白上是一路一路.
你再明白多一點.
再明白多一點.
再個人一點.
再跟你建立關係多一點.
最後牽動你的情感.
你看到兩個的反應.
我特別跟大家分享.
他們哭泣.

$^{921}$是因為聽到上帝的話而哭泣.
他們喜樂.
是因為明白了.
他們所聽的教導而喜樂.
回到一個更大小的處境.
剛才說了.
他們是被擄歸回的一代.
被擄是一個很沉重.
悲痛的經歷.
某程度上.
是離開上帝的後果.
是不是?.
如果再具體一點.
是離開上帝的話語的後果.
你記得.
被擄之前.
上帝是三番四次.
藉著他的先知.
你也聽過很多先知.
以賽亞.
以西結.
尼欽.
耶利米.
再加一些小先知.
很多.
一次又一次提醒他們.
要依靠上帝的話語.
不要被外敵.
很多那種假先知.
假教師的話.
聽他們的話.
以為對於應許之地.
對於這個殿.
有些迷思.
簡單來說.
如果特別看耶利米書的時候.
大家可以回去慢慢詳細看.
但是他們沒有聽上帝自己的先知的話.
去痛改前非.
去靠著上帝.

$^{961}$就算下了一個很危急的處境.
他仍然可以臨危立馬.
但是沒有這樣做.
結果上帝.
用一個七十年的被擄.
換來一個信仰上.
生活上.
一個很沉重的打擊.
所以去到異地的時候.
其實分種已經沒有了.
上帝的話語的手.
讀的都是巴比倫的語言.
你記得《但義理書》.
他們被人恥笑.
被人奚落.
各樣東西.
他們想念耶路撒冷.
想念昔日的日子.
好像《詩篇42篇》那些說法.
他沒有想過.
都不知道上帝這個懲罰.
是會到什麼時候.
所以到了他們能夠回來.
你想象.
你算一算.
七十年之後.
可能老人家都過世了.
在異地.
你明白我的意思嗎?.
回到來的時候.
就算有第一身的經驗.
那時候可能都是小朋友.
回到來都是老人家.
帶著下一代.
回到自己的地方.
然後竟然呢.
做了很多工程之後.
建完殿.
建完牆之後.
可以聚集在一起.

$^{1001}$有上帝的一排的.
原來上帝的文士還在.
有很多人.
祭司還在.
利民還在.
所以有一個這麼莊重的聚集.
然後這樣去高舉上帝的話語.
所以他們一聽到.
其實是一種很責心的感受.
可能他們想起自己.
避老前.
或者上一代.
避老前那種.
藐視了上帝的話語.
但是上帝竟然再給一個恩典.
再給他們一個機會.
再次聚集在一起.
重新讀上帝的話語.
這是一個對他們來說.
很責心.
一個很痛悲的經歷.
怎麼樣呢?.
上帝再給我們機會的時候.
我們會不會珍惜.
重視上帝的話語呢?.
這是他們那種情緒激動的原因.
又珍惜.
又想起以往那種沉重的打擊.
然後今天好像又重新開始.
我們怎麼開頭呢?.
怎麼去應接呢?.
是一種很大的.
心裡面很大的情緒反應.
然後當尼熙米.
當這些利民人出來鼓勵他們.
安慰他們的時候.
他們為什麼會有大大喜樂呢?.
我想當經過了早上的講道.
下午的討論.
再去到一個黃昏時間的.

$^{1041}$一種這麼個人性.
這麼情感性的.
真情的對話的時候.
其實他才在當中.
更加接觸得上帝的話語對他的意思.
對生命那種帶來.
好像尼熙米那樣說.
帶來那種喜樂和力量.
你再想下去其實是不容易的.
因為被老歸回.
百廢待興.
沒有人會看得起他們.
或者會重視他們.
可以將整個耶路撒冷復興.
但是他們就是來到的時候.
要擔起往後重建耶路撒冷的工作.
往後的日子怎麼走呢?.
需要一些力量的支持.
最重要的是.
上帝的話原來在當中.
成為一個支撐著他們.
往後可以繼續向前走的力量.
今天我們都是面對這個情況.
香港裡面我經常說.
其實很多事情都要重建.
教會要重建.
社會要重建.
人的關係要重建.
很多事情我們都好像很荒涼.
不知道大家有沒有一種荒涼感.
有時候在街上走過.
你會感覺到一種比較荒涼的感覺.
香港好像不是這樣的.
我們又會回想過去的日子.
想起今天的日子.
但是我們作為上帝的子民的時候.
我們的力量的根源在哪裡呢?.
我們的方向感在哪裡呢?.
我們的喜樂的源頭在哪裡呢?.
我們對未來的盼望在哪裡呢?.

$^{1081}$不是在地上的事情.
不是在經濟裡面.
也不是在社會的建設裡面.
也不是在什麼繽紛裡面.
在我們自己作為上帝的子民.
在當中靠著上帝的話語.
一步一步地建立一個群體.
成為這個地區這個社區的崩台見證.
這個才是我們真正的存留的目標.
所以到最後的時候.
我的總結.
甚至可以說這個.
建立一個讀經的群體.
是我們教會很重要的.
存在的價值.
存留的目標.
昔日的子民去被擄歸回.
首要.
其實我覺得不是他的硬件.
雖然次序上是先建了電線.
然後建了牆.
再建立信仰.
但是倒過來價值上.
其實第八章反而是最重要的.
甚至可以說比你建牆還重要也說不定.
大家可以討論一下這個問題.
是很重要的.
是整群人一起再次推翻信仰.
再次聽回上帝的話語.
這個是建立信仰生命群體是最重要的.
所以到我們今天來說一樣.
我們每一個生命都要上帝說的餵養.
我們活著不是單靠食物.
是上帝說給我們作為生命的糧.
給我們每天的指示.
當行的路.
當做的事.
所以要以道去連繫我們每一個人的生命.
無論是個人群體都要有道去建立.
所以盼望我們能夠讓道.

$^{1121}$成為整個群體裡面的關係的根基和養分.
和我們那個內容.
以致我們彼此勉勵.
彼此同行.
彼此督促的時候.
都是有上帝的話語的同在.
這個就是我們能夠切實的建立成為獨經群體之後.
那個群體生活當中那種關係的美好.
而且我們帶著使命.
再去成為別人的祝福.
向前再走下去.
在一個很不明朗很黑暗的世代當中.
我們真的可以為世人燃點一點光芒.
讓世人都能夠進入到一個.
以耶穌基督為中心的一種生命的內容.
盼望我們能夠在當中領受上帝的話去回應他.
我們最後一起在這裡禱告.
天父多謝你給我們今天一起.
從《聖經》當中.
昔日的子民.
去歸回.
很迷失.
其實很多事情都好像迷失的狀態當中.
你的話再次去到中間.
好像一個定海神針那樣.
安定他們子民的心靈.
其實我們都需要一樣.
你今天可能在香港處境.
何況都不是一種像荒涼的狀況.
都需要你的話成為生命中的一種穩妥.
一種掌舵.
這樣的力量.
給我們盼望.
給我們經過教會建立一個成為愛慕你.
重視你懷舊的群體.
我們的群體就是一個見證.
將你的愛.
將你的真理和公義.
向身邊的人宣講.
願你親自建立.

$^{1161}$祝福我們每一個.
聽我們同心禱告.
奉主耶穌得勝.
領治而救.
阿們.
\newpage



\section{使徒行傳 16:6-10}
\label{sec:8mvS2zOfCvg}
\textbf{2024.10.06 《來自遠方的呼喚,你們聽見嗎?》 使徒行傳16\_6-10 (和修版) 講員 : 楊穎兒姑娘}
\newline
\newline
連結: \href{https://youtube.com/watch?v=8mvS2zOfCvg}{\texttt{https://youtube.com/watch?v=8mvS2zOfCvg}} ~~~~ 語音日期: 2024-10-08
\newline
\newline
\hyperref[sec:dufzJl9BQpc]{\small{< < < PREV SERMON < < <}}
~
\hyperref[sec:index_chronic]{\small{[返順時目]}}
~
\hyperref[sec:index_scriptual]{\small{[返順卷目]}}
~
\hyperref[sec:jVaUl1bB0os]{\small{> > > NEXT SERMON > > >}}
\newline
\newline
使徒行傳 16:6-10
\newline
\begin{longtable}{cl}
\hline
\hline
章節 & 經文 (和合本修訂版)\\
\hline
16:6 & \begin{tabularx}{0.7\textwidth}{X} 因為聖靈禁止他們在亞細亞講道，他們就經過弗呂家、加拉太一帶地方。 \end{tabularx} \\ \\ \relax
16:7 & \begin{tabularx}{0.7\textwidth}{X} 到了每西亞的邊界，他們想要往庇推尼去，耶穌的靈卻不許。 \end{tabularx} \\ \\ \relax
16:8 & \begin{tabularx}{0.7\textwidth}{X} 他們就越過每西亞，下特羅亞去。 \end{tabularx} \\ \\ \relax
16:9 & \begin{tabularx}{0.7\textwidth}{X} 夜間，有異象向保羅顯現。有一個馬其頓人站著求他說：「請你過來，到馬其頓來幫助我們！」 \end{tabularx} \\ \\ \relax
16:10 & \begin{tabularx}{0.7\textwidth}{X} 保羅既看見這異象，我們就立即設法往馬其頓去，認為神呼召我們傳福音給那裡的人。 \end{tabularx} \\ \\
[1ex]
\hline
\hline
\end{longtable}
$^{1}$弟兄姊妹平安.
很開心見到大家.
很多熟悉的面孔.
我見到你們的時候.
很高興今天可以再回來.
紅印堂分享一些.
尊教猜拳的訊息.
我有這個感覺.
好像經常回到娘家.
是不是.
好像大家可能都被我年長.
看著我長大.
所以很感恩.
今天我們一起.
去看看保羅宣教的工作.
我們在進入經文之前.
我們一起低頭祈禱.
讚美你.
是你讓我們聚集在一起.
是你讓我們有機會去敬拜你.
你亦給了我們一個職份.
去傳揚福音.
願今天的訊息更加堅固.
我們教會.
我們眾弟兄姊妹的心.
我們這樣祈禱.
奉主名求.
阿們.
仕途保羅在.
仕途行傳裡面.
他進入一個.
要往.
外邦地區宣教工作的階段.
那時候在耶路撒冷.
開了一個大會.
決定要猜派保羅.
去做一些宣教的工作.
大家很同心.
隨著時間的轉移.
他們去完第一次的宣教旅程.

$^{41}$回來了.
很興旺.
他們去了小亞細亞的地方.
就是圈住的地方.
圈住的地方.
小亞細亞的地方.
傳講神的話語.
仕途很有信心.
保羅做得好.
在外邦地區擴展福音.
將福音帶到.
不單止是猶太人.
也帶到.
外邦地區.
也在那些地方接觸.
多元文化.
和不同文化背景的人.
去交流.
也在當中幫助.
初期教會的仕途成長.
去堅固他們的信心.
告訴他們.
受的迫迫將來會得到賞賜.
他也用了一些.
很靈活的宣教技巧.
和一些同工一起去宣教.
當他們去到.
佛裡加.
加拉泰這一帶的地方.
聖靈禁止他們在.
亞細亞講道.
亞細亞就是紅色圈住的地方.
綠色就是他們.
需要走過的道路.
我們看地圖.
按道理.
經過當然是順道.
傳道.
沒理由去到那裡又回來.
但這次是聖靈禁止他們.

$^{81}$不單止這樣.
他們去到.
美西亞邊界.
想被推離的時候.
耶穌的靈再次不讓他們再去.
我們看地圖.
我們看到.
美西亞.
又過了一些地方.
他們想經過美西亞再上去.
被推離.
去俄羅斯再往中國那條路.
但聖靈說不讓.
耶穌的靈不讓他們去.
那時候保羅和他們的同伴.
很渴望進入.
去到的地方都希望福音可以去到.
他們那種激情.
偏偏遇到聖靈的阻攔.
我們看經文.
可能是一兩句已經跳到這個位置.
但在他們的宣教旅程.
其實是很漫長的.
他們在過程中會沮喪.
他們走每一步路.
都受到阻礙.
就好像我們看箴言.
有一段經文.
人心是有很多計謀的.
我們有很多計劃.
就算我們在每個人生中.
我們都會看到很多希望.
無論是工作.
無論是夢想.
無論是學業.
我們都很想在環境中.
我們可以控制.
我們會有計劃.
我們會有目標.
我們知道什麼是最好的.

$^{121}$我們向前衝.
我們去祈禱.
我們去策劃.
我們去走自己的道路.
但當前面突然之間.
沒有路的時候.
東又走不到.
西又走不到.
還要是經過一段很漫長的時間.
那種感覺.
就好像被拒絕.
被離棄.
那種希望.
被擊碎.
又或者我們會很想責怪身邊的事情.
阻礙著.
都是那樣東西阻礙著.
唯有耶和華的.
酬酸才能立定.
我們會不會.
在下一次有機會再遇到.
這種無路可走的時候.
我們可以向神.
去求智慧.
去求智慧去辨識.
並且求.
聖靈帶領我們.
去停止現在做的事.
而能夠重新去評估.
一個新的方向.
就在這個時候.
他們就越過了.
有條路能夠走得了.
然後去了特羅亞.
又過了一點.
去了特羅亞這個地方.
夜間.
有異象向保羅顯現.
有一個馬其頓人.
站著和他求.

$^{161}$你過來吧.
到馬其頓這裡幫我們吧.
我們不知道這個現象.
是以什麼形式.
向保羅顯現.
也不知道他要幫.
馬其頓的哪些人.
但我們知道這個現象.
是很針對保羅當時的情況.
去說.
前面沒有路了.
現在有一個異象告訴他.
在別人的需要身上.
保羅.
有一個新的任務.
你有沒有試過這樣呢.
看到別人的需要.
而你找到自己的價值.
在這個關鍵的時刻.
其實就是.
要我們放下原有的計劃.
而去跟隨.
一個更高更遠.
更大的旨意.
在做宣教的工作裡.
我發現.
我負責的很多角色都是開荒的角色.
開荒就是本來水不到的地方.
你就要開到一條路.
於是有水可以到.
不知道大家有沒有耕田的經驗呢.
我玩過耕田.
發現耕田的時候.
我不是只坐那段路.
因為你以為水會灌溉到.
旁邊都是種植的地方.
你也要爬鬆旁邊的土壤.
所以爬鬆的時候.
也會做旁邊的其他東西.
就好像有時計劃了一些東西.

$^{201}$也覺得有可能.
但原來不是.
又要回去那條路繼續爬.
所以我們最初有些計劃.
我們去推進.
但有時推進不了.
我們就需要暫停.
去聆聽.
去重新調整.
以致能夠做得更好.
更遠.
這個停頓其實是將.
我們原有的計劃放下.
我們去學習謙卑.
我們去學習信任上帝的那種大靈.
意味著我們知道.
我們自己的限制.
而我們看見上帝.
更遠的地方.
人會覺得.
我要計劃的.
如果不是我走錯.
但其實.
我們今天在保羅這裡看到.
他們其實是要.
有信心去認定.
那條是上帝.
讓他們走的路.
我們看看下一段.
保羅就機看見.
這個看見者現象.
保羅看到而已.
不是保羅和隊友.
接著他們就立即設法去.
馬其頓去.
當然保羅會在當中和隊友分享.
看到這個現象.
於是隊友一起.
設法.
都要設法去到馬其頓.

$^{241}$並且認為.
神是這樣呼召他們.
去傳福音.
我們看到.
保羅的計劃被聖靈打亂.
不單止要面對很多障礙.
現在還要.
去一個沒有想過去的地方.
而我們也看到.
保羅意識到有一種迫切性.
馬其頓那裡.
會有人.
是一個異象而已.
有人需要.
福音.
然後就去認定.
認定是需要信心的.
你也不知道.
是否這樣的時候.
去認定.
接著去踏上.
保羅的經歷.
讓我們看到.
保羅的使命.
其實是很敏銳.
在上帝的大靈裡.
可能我們當中.
也有面對這樣的情況.
可能在我們的工作.
或是在找學校.
或是在找人生的下一站.
我們也要很敏銳地.
去看看神的大靈是怎樣.
我們在當中.
可能會考慮到很多自己的需要.
但我們會否停一停.
我們想想.
踏出去的那一刻.
是需要上帝去幫助你的.
試試.

$^{281}$試試走那一步.
走了那一步.
看看上帝的大靈.
是怎樣奇妙的.
這個社區大家很熟悉.
洪水橋社區.
洪水橋這個社區是很獨特的.
我記得我在這裡做神學生的時候.
魯牧師帶我走出去.
大馬路.
你看看.
這邊就是某某村.
這邊是曠野.
我們就在曠野的邊邊.
很明顯.
洪恩堂的教會.
是有自己的獨特性.
我相信在我們.
做一些地區性差傳的事工.
也不乏要聽聽.
周邊的鄰舍.
的聲音.
或是他們的需要.
以致我們更加可以.
幫助鄰舍的福音.
在世界上.
有一個叫.
洛桑世界福音宣教大會.
這個可能.
大家很新.
未聽過.
這個福音大會.
其實是在1966年.
由葛培里報道團.
葛培里聽過了.
葛培里報道團.
是贊助和組織的.
這個組織其實是想.
集合.
宗教會.

$^{321}$跨宗派宗教會.
有志宣教的人.
來到這個大會.
去討論下一步.
是如何一起去宣教.
在當年第一次的宣教大會.
因為開第一次的時候.
沒有想到會有第二三四次.
第一次開的時候.
當中有牧者提到.
有一個宣教的故事.
就是有五個年青人.
可能大家聽過.
五個年青人去了厄瓜多爾.
一些土著宣教.
一去到被殺.
死了.
年青人很有為.
被殺了.
很害怕.
在大會上分享.
原來這五個宣教士.
他們死了.
傳不到福音.
但他們的太太.
知道他們的丈夫死了.
繼續回到那條村.
去服侍.
這群村民.
村民問他們.
為何還要來.
他們已經殺了丈夫.
但這群妻子說.
因為耶穌的愛.
教我們去寬恕.
所以我們再來.
將耶穌的愛帶給你們.
後來這件事.
是真人真事.
拍了一部電影.

$^{361}$宣教大會裡說.
宗教會是很火熱.
以致.
這幾位就是那幾位宣教士.
年青.
有為 有學識.
但一去到就死了.
在宣教大會裡.
第一次.
主題是眾人都要聽到.
宣教會的聲音.
宗教會於是在開會後.
就向著目標進發.
在七,八十年代.
教會就要傳揚福音.
到八,九年代.
在宣教樂喪大會裡.
第二次.
這次就提出了.
有一個.
十,四十的窗口.
甚麼意思呢.
就是在北緯十度.
和四十度中間的居民.
想不想到有甚麼地方.
有中國.
有中東地區.
有中美洲地區.
中間的地方.
他們生活很貧困.
生活質素很低落.
他們有95\%.
信奉伊斯蘭教.
印度教.
佛教.
無神論.
這些是福音去不到的地方.
就是世界的中間.
於是就興起了很多宣教士.
我們就是要去.

$^{401}$十,四十窗口的地方.
去宣教.
如果我們有.
收過這些宣教的訊息.
就知道.
現在非洲,印度.
有很多宣教士.
都曾經去過那裡.
而且建立了教會.
然後去到2010年的時候.
又有一個.
第三次的宣教大會.
就說神在基督裡.
是要叫世人和自己和好.
於是又帶領著宗教會.
向著這個目標.
去進發.
這個和好是人與人之間.
與上帝之間.
與大地之間的和好.
到了2024年.
其實是上個月.
這個消息是新鮮滾熱辣.
上個月有第四次的.
洛桑宣教大會.
就是要宣告.
要向全民宣告.
讓上帝去顯現.
怎樣去讓上帝顯現呢?.
就是我們.
每個信主的基督徒.
我們都可以帶著.
我們的屬靈生命.
我們的過往生命.
我們怎樣認識上帝.
去到我們的工作場所.
去到我們坐巴士的地方.
去到我們的教會.
去到任何一個角落.
去到一個會受迫迫的地方.

$^{441}$就是我們全民皆兵.
成為一個福音的出口.
讓天國的福音.
繼續傳下去.
這些會很興奮的.
怎樣形容呢?.
很多集合了.
全部對宣教有熱誠的人.
一起去共同.
去得到一個這樣的領袖.
這些會最近這一次.
2024年這一次.
是有五千人參加.
世界各地.
而我們今天看的經文很有趣.
其實我們今天看十六章.
其實在十五章的時候.
當時後都開了一個大會.
教會開了一個大會.
他們就說了.
我們同心定義.
揀選幾個人.
猜他們和我們親愛的巴拿巴.
和保羅.
去你們那裡.
他們就開了一個會.
有一個共識.
要一起去猜.
一些有屬靈.
生命.
比較有屬靈帶領的.
領袖.
去傳揚這個福音.
就是這樣.
在這個會裡.
保羅得到眾教會的肯定.
肯定他的工作.
我們去.
反思保羅在.
仕途行轉的那種經歷.

$^{481}$我們其實都可以.
帶回我們自己人生的旅程.
我們教會的旅程.
我們其實有沒有.
一起去.
擁抱一些使命呢.
我們是一起.
透過教會的聚會.
一起認識.
一起禱告.
一起成為神的耳朵.
我們去.
接收一些.
遠方的需要.
就好像在.
仕途保羅.
這次的行程裡面.
他有一個意象就是.
去到馬其頓那裡.
傳福音.
他不知道的.
在他去接受這個意象的時候.
他都不知道那些人是誰.
但當我們繼續看十六章.
十七章的時候.
我們會發現保羅.
去到菲勒比.
有女弟子.
去歸順.
去耶穌那裡.
她是一個買賣的女人.
她聽了福音之後.
她全家都接受了.
並且洗禮.
當中有一些奴隸的女子.
沾了邪靈.
而保羅需要.
將這個靈驅走.
到菲勒比的時候.
他被收監.

$^{521}$但同時.
他們透過唱詩歌的時候.
監獄的門就打開.
他們又得到自由.
在監獄的玉卒.
他們看到這些事情的時候.
又一家全家.
接受耶穌受洗.
然後他們去到.
帖撒羅尼迦的時候.
又和很多猶太人辯論.
跟著跟他們說了很多.
關於他們的神.
希臘神.
不是正統的信仰.
你們要信耶穌基督.
然後又遇到很多猶太人.
去反對他們.
很多很多這些故事.
在他接受這個異象的時候.
我一定相信.
他不知道有這麼多事.
要發生.
唯有他認定.
憑信心去認定的時候.
才會有.
這些事件去發生.
有這些人去得到福音.
是要用信心去踏上.
保羅就是.
這樣認定了.
在剛才我提到的.
洛桑福音世界大會裡面.
我看到我沒有去.
覺得很可惜.
但在當中.
有一個伊朗人.
這個伊朗人.
他大概五十六十歲.
他的故事.

$^{561}$其實他在一個穆斯林的家庭.
成長.
媽媽甚至在.
一個廟裡面.
教.
可蘭經.
你看到他的家庭.
成長的伊斯蘭.
是很重的.
後來1996年4月.
他悔改.
他跟從耶穌.
他得到的快樂.
是不能形容的.
他離開教會的門口.
他就說我接受了耶穌.
很開心.
我可以離開了.
他在1996年信耶穌.
到了2010年.
過了十幾年.
當中在伊朗的地方.
有很多教會建立.
到了2010年.
12月26日.
剛剛復活節之後.
有一大隊的士兵.
來抓了他.
因為他犯了國安法.
因為在他們的國家裡面.
他們是.
全陽福音的話就是犯了國安法.
於是抓了他.
就關了他在一個.
兩米乘兩米的房間.
沒有窗的,黑漆漆的.
他過了兩年的時間.
而到了.
2015年.
他在被收監.

$^{601}$1820日之後.
他被釋放.
為他禱告.
原來在他出來的時候.
他才知道.
伊朗的教會.
突然蓬勃得很厲害.
在20多個城市.
有50個教會.
已經誕生了.
我們看到.
讓他困在一個.
很困苦.
一個受迫迫.
一個黑漆漆的.
但是.
被困住.
教會得到興旺.
迫迫是不能夠.
去拆毀.
我們的教會.
只會令到福音.
更加廣傳.
而且這個迫迫.
在神的故事裡面.
是一個很重要的部分.
如果我們看回聖經,很多人都受過苦.
我們是不是很想.
福音可以去到很遠的地方.
是不是.
是.
我聽到點頭.
我們有迫迫.
迫迫.
迫迫靠近我們.
很人性.
我們會覺得.
我們很想傳福音.
但是這個迫迫.
不要了.

$^{641}$不是我.
在我服侍的群體裡.
有一部分的人.
是穆斯林的轉信者.
他們面對很大的迫迫.
因為他們生出來是穆斯林.
那張身份證就寫著.
他們的信仰.
是還沒出生已經有這個身份.
當他們歸信的時候.
不是說.
你進來教會參加崇拜.
他們就可以有基督徒的身份.
他們是要去拆毀.
他們穆斯林的名字.
他們要離開他們的家庭.
因為他們的家庭覺得.
他們是羞恥.
離開了他們的信仰.
要政府通緝他們.
他們的家人要殺害他們.
所以穆斯林的歸信者.
是面對極大的迫迫.
但是這個迫迫.
其實我們很難去理解.
是不是.
我們不會知道.
為什麼要走.
或者我們不會知道.
被親人眾叛親離.
因為信耶穌的緣故.
而要被殺害.
我們不會明白.
我服侍的群體.
當他們來到香港的時候.
他們發現.
都安全了.
香港都可以在這裡尋求庇護.
但他們這次不是面對迫迫.
他們面對的是.

$^{681}$很多人很冷漠.
很不明白.
很不了解.
跟別人說的時候.
甚至會有歧視.
甚至會有問題.
或者聊天.
會侵犯到他們.
侵犯到他們.
他們不明白.
他們繼續在受苦的裡面.
只不過形式轉了.
有時候我們在社區.
實在太安恕了.
我們聽不到.
這個世界的苦難.
甚至我們.
自稱為基督徒.
但我們可能聽到這些苦難的時候.
我們不懂得怎樣面對.
我們不懂得.
活出那種愛.
我們覺得不關我們事.
不過.
很好啊,我聽了你的故事.
但其實這些不是真正的愛.
我們教會.
今天香港的教會.
我們夠不夠火熱呢.
我們為福音的擴展.
更加火熱.
在我們監堂.
我是監堂負責一個英文侍工.
我們來了一個穆斯林的歸信者.
他被迫走到.
我們當中.
跟我們一起聚會.
我看到他受了很多苦.
在他們的國家裡面.
從小就被基督徒欺負.

$^{721}$女性.
要做一個國泥.
基督徒.
因為是伊斯蘭國家.
女性.
基督徒會移去.
垃圾山的地方.
你們住在那裡.
你們吃豬的.
我們國家倒垃圾的時候.
倒到那裡.
所以我覺得地區很臭.
當然.
他轉信的話.
就是.
你要撕掉自己的東西.
而帶他轉信的人.
其實都是要走的.
因為都是犯了國安法.
你不可以讓他本身是伊斯蘭教的.
你去將他轉信.
他來到我們教會.
因為我們中間.
服侍一些在旅程裡面.
受苦的人.
在過程裡面.
他終於來到香港.
而有另外一間神學院.
去支持他.
去讀神學.
他現在是在信修神學.
我看到很大的機會.
讓這個弟兄.
他知道那裡的穆斯林文化.
也都是要接受神學裝備.
去拆毀他原有的.
他自己.
然後神會怎樣利用.
他的生命.
去傳揚福音.

$^{761}$很感動.
在我服侍的時候.
我都是盡量幫他.
在我們有的平台.
去幫助這個人.
將來能夠成為一個.
繼續去牧養教會的一位弟兄.
神就是要我們有.
這些耳朵.
我們去聽到周邊的聲音.
這些聲音未必是.
我們的真相.
可能是周邊的人說話.
可能是你看到的一些東西.
我們去反思的時候.
我們有沒有這個耳朵呢.
如果我們在過往的日子.
沒有不要緊.
我們求就有.
我們的主耶穌是.
應許過我們.
我們求我們就有.
神在這個世界上.
所做的事很多.
這些都是神關心的事情.
神關心的事情.
是不是我們的事.
神關心的事情.
當然是我們的事.
我們一起去參與.
不是說神你在那裡做.
我在這裡參加敬拜.
不是.
如果我們可以.
在祈禱上.
在支持上.
在關心上.
在行動上.
我們可以盡分力去做.
猜拳.

$^{801}$這個有點朦朦地.
因為我在YouTube拿下來.
我們的猜拳模式是怎樣呢.
這個課程.
是我剛剛在錦繡堂教完的.
是一個猜拳士工聯會.
制定的一個課程.
教會是怎樣去.
做一個宣教的教會.
其中有一個我覺得挺好的.
所以抽了出來.
跟大家分享.
我們今天教會的模式.
宣教是一個.
士工的一小部分.
是好像其他的士工.
敬拜.
報道 團契 門訓的一部分.
還是我們將宣教.
放在一個核心.
以致我們在我們的報道.
我們敬拜.
我們團契 我們門訓.
裡面都是有宣教的火.
而去推動.
我們所有的小組.
將我們.
全部人的生命一起去成長.
好像透爐.
宣教就像透爐.
大家有燒過東西吃.
你要起好爐.
火起得好.
燒出來的東西才好吃.
火經常都是小小小小.
雞翼中間.
燒得不熟.
其實道理是一樣的.
所以他擺這個圖也挺好.
讓我想起一個郊區的燒烤爐.

$^{841}$今天我們在地區做猜拳工作.
我們要有這團火.
今天我們做宣教工作.
我們去看遠一點的地方的時候.
我們也要有這團火.
無論你在哪個崗位.
你都可以參與.
猜拳的服侍.
在地區性的猜拳服侍.
我們去接觸這一區的居民.
宣教的服侍.
我們看遠一點.
我們去聽聽他們的故事.
我們去感受他們.
我們去跟他們聊天.
我們去支持土共.
不要想.
很多事情要做.
猜拳很多事情要做.
我做不來.
不要這樣想.
你不會做不來的.
沒有你.
這些事工.
也都可以繼續進行.
只不過當我們.
不去參與的時候.
損失的是我們.
我們看不到神在做什麼.
祂在做的事與我們無關.
神在遠方.
給大家有的經驗.
我們沒有.
你不參與就沒有.
剛才主席也介紹了.
我們在奸堂有一個.
負責英語的士工.
其實是.
今年3月才開始.
在過往的日子.

$^{881}$去年的時候.
曾牧師在場.
他可以做見證.
去年的時候.
我跟曾牧師說.
3月我要去工場.
今年3月偏偏停下來.
要開這個士工.
我那時候.
覺得我都不是牧會的.
我的呼召一向都是宣教的.
但現在要停在牧會.
什麼都不會做.
嘗試去踏上.
踏上的時候.
最困難的第一件事就是找敬拜隊.
因為我們英語的士工.
我們有人帶敬拜.
沒有敬拜隊.
找了一陣子都沒有人.
最後我找了我的先生.
沒有人.
不好意思.
要你一起.
他說好 試試吧.
最初他領唱.
之後他就開始試彈琴.
他猜不到.
他之前是懂得彈琴.
但他沒有想過現在什麼歌都懂得彈.
而且他一聽就懂得彈.
是一個.
技能的解鎖.
我們沒有想過.
我跟曾牧師去年.
去談的時候完全沒有說過這件事.
現在繼續.
做這件事的時候.
發現原來我和我先生.
可以這樣拍檔.

$^{921}$一起去.
在不同的地區.
或是世界各地.
我們都可以以這樣的模式去建立教會.
我覺得這是很大的祝福.
不單是我們的家庭.
但是在說.
教會或是一些遠方.
準備要有教會的人.
我不知道上帝會帶我去哪裡.
但如果有需要去建立教會的時候.
原來我和我先生是可以做得到的.
我很驚訝.
是不知道的.
就是認定.
認定的聲音就去走.
在我夢召的時候.
我很少和別人說.
在我夢召的時候.
我經常做一個夢.
晚上做一個夢.
這個夢就是.
很多人.
來到我面前.
他又不說話.
但嘴巴就這樣.
不開心.
然後就默默有淚流下來.
這個夢有好幾次出現.
我覺得很害怕.
小時候覺得是不是邪靈.
很害怕.
於是有一次我決定.
又來我面前哭的話.
我就和他說.
你不要來找我了.
你去找耶穌吧.
真的有一次.
最後做一次夢的時候.
我和他說你不要來找我了.

$^{961}$你去找耶穌吧.
這個夢我一直記住.
我發現原來被帶到.
尋求被迫.
這些群體裡面的時候.
我見到他們了.
真的有這些人.
他們見到我的時候.
他們未必有眼淚流出來.
但那些眼淚是在心裡面流出來.
讓我有機會去接觸到這些人.
這些就是在社會裡面.
最沒有人喜歡的人.
因為他們就像香港尋求庇護.
就是拿香港的資源.
但上帝帶我們去.
和這些人一起.
去哭泣.
和一起經過這些階段.
然後我就開始去服侍.
這些被很多迫迫的歸主者.
就好像剛才說的.
我有機會.
好像栽培一位.
歸信的.
他日他有一日.
可以成為牧者.
願意踏上的時候.
有很多故事可以聽.
大家有機會去看《仕途保羅》.
很多故事很奇妙.
上星期有機會和曾牧師.
和曾姑娘去吃飯.
我們同工會吃飯.
有機會吃飯的時候.
都說到我們會去一個工場.
在一個小鎮.
小鎮很好,交通很方便.
方便我們做事工.
不用浪費時間.

$^{1001}$在旅程中.
很有趣.
這個星期初.
見到警察會的總幹事.
他來到.
他告訴我.
我們有一對夫婦在那個省.
但不是那個鎮.
但他們準備退休.
沒有人接.
後來我看了一些代討信.
原來這對要退休的宣教士夫婦.
已經有一些.
記憶不太好的症狀.
而且在過往的一年.
都是這樣.
禱告需要人去接事工.
我和先生分享的時候.
我們覺得.
不如我們也試試.
我們不懂做.
但我們去試試.
先生說好.
我們去試試.
在那個鎮裡住了一堆.
北非.
一些很藍.
他們的民族是99\%.
不信耶穌的群體.
很有趣.
就是要有這樣的靈活性.
去分享的時候.
我不怕去分享.
因為分享的時候.
曾牧師就有份去經歷.
我所經歷的那種過程.
那種宣教的過程.
所以不要怕和別人分享.
我一領的時候.
怕錯了.

$^{1041}$不要緊 錯就錯.
神會關門.
在我們生命的路上.
多美好.
神的關門也是他的動作.
值得我們去讚美.
而且我們可以.
招聚到大家一起去見證.
所以各位.
在座的各位.
我今天很想邀請你們.
我們一起去立志.
我們為宣教.
為猜拳去立志.
我們去想想.
在未來的日子裡.
在崗位裡去服侍.
不是每個人.
都要用行動去踏出去.
去愛地.
但我們可以開我們的耳.
我們可以開我們的眼睛.
我們去看片段.
我們可以去聽他們的故事.
我們可以去接觸他們.
如果我們有機會的話.
我們接觸他們.
我們再聽多一點.
我們用心去聽.
我們用禱告一起去結合.
很想挑戰大家.
每一位.
有機會的話.
在你的團契.
在你的事奉崗位.
大家一起去討論.
去找這個崗位.
你究竟可以在宣教猜拳裡.
哪裡有份.
我剛才進來的時候.

$^{1081}$很開心.
又來到這麼美的咖啡室.
相信在建這個咖啡室的時候.
都是認定了一個目標.
在這個咖啡室.
去吸引地區的居民.
我們可以怎樣去.
繼續承接著.
上帝的宣教工作呢.
是我們在座每一個都有份的.
認定你是有份的.
上帝就會帶領.
我們一起祈禱.
秋月天父.
你是那個.
瑞領教會的主.
你是那個瑞領教會的主.
你是不斷在爭戰.
你是不斷興起很多.
受苦的群體.
去跟著你的名利.
去受苦.
我們求你繼續去堅固他們.
你去猜派使者.
去幫助這些人.
建立教會.
主啊.
當我們看到你.
做這麼多的事的時候.
我們不得不來到你的面前.
我們去求你的.
寬恕你的憐憫.
因為我們當中可能.
不懂得.
去聽.
去看.
去愛.
去關心.
我們這樣去求的時候.
因為你.

$^{1121}$應許過.
當我們向你求的時候.
你就會開門.
我們求的時候你就會得著.
我們要求的就是教會能夠興旺.
如果宣教能夠.
教會能夠興旺的時候.
我就求你.
將宣教猜拳的心.
給我們.
洪恩家.
我們錦宣家.
我們是在地區火熱.
我們在世界火熱.
我們這樣求.
我們認定.
我們是得著的.
主啊.
就將這個.
很好的種子.
放在我們心裡.
去成長.
以致有一天我們回來的時候.
我們看到.
主啊.
主席.
阿們.
\newpage



\section{歌羅西書 4:7-15}
\label{sec:jVaUl1bB0os}
\textbf{2024.10.13 《來自保羅團隊的問安》 歌羅西書4\_7-15 (和修版) 講員 : 曾裔貴牧師}
\newline
\newline
連結: \href{https://youtube.com/watch?v=jVaUl1bB0os}{\texttt{https://youtube.com/watch?v=jVaUl1bB0os}} ~~~~ 語音日期: 2024-10-14
\newline
\newline
\hyperref[sec:8mvS2zOfCvg]{\small{< < < PREV SERMON < < <}}
~
\hyperref[sec:index_chronic]{\small{[返順時目]}}
~
\hyperref[sec:index_scriptual]{\small{[返順卷目]}}
~
\hyperref[sec:dXTahKH7Xis]{\small{> > > NEXT SERMON > > >}}
\newline
\newline
歌羅西書 4:7-15
\newline
\begin{longtable}{cl}
\hline
\hline
章節 & 經文 (和合本修訂版)\\
\hline
4:7 & \begin{tabularx}{0.7\textwidth}{X} 推基古是我親愛的弟兄，忠心的僕役，和我一同作主的僕人；他要把我一切的事都告訴你們。 \end{tabularx} \\ \\ \relax
4:8 & \begin{tabularx}{0.7\textwidth}{X} 我特意打發他到你們那裡去，好讓你們知道我們的情況，又讓他安慰你們的心。 \end{tabularx} \\ \\ \relax
4:9 & \begin{tabularx}{0.7\textwidth}{X} 我又打發一位親愛忠心的弟兄阿尼西謀同去；他也是你們那裡的人。他們會把這裡一切的事都告訴你們。 \end{tabularx} \\ \\ \relax
4:10 & \begin{tabularx}{0.7\textwidth}{X} 與我一同坐牢的亞里達古問候你們。巴拿巴的表弟馬可也問候你們。關於他，你們已經得到指示；他若到你們那裡，你們要接待他。 \end{tabularx} \\ \\ \relax
4:11 & \begin{tabularx}{0.7\textwidth}{X} 稱為猶士都的耶數也問候你們。奉割禮的人中，只有這三個人是為神的國與我作同工的，也是使我心裡得安慰的。 \end{tabularx} \\ \\ \relax
4:12 & \begin{tabularx}{0.7\textwidth}{X} 有一位你們那裡的人，作基督耶穌僕人的以巴弗問候你們。他禱告的時候常為你們竭力祈求，願你們能站穩而成熟，充分確信神一切的旨意。 \end{tabularx} \\ \\ \relax
4:13 & \begin{tabularx}{0.7\textwidth}{X} 他為你們、老底嘉和希拉坡里的弟兄多多勞苦，這是我可以為他作見證的。 \end{tabularx} \\ \\ \relax
4:14 & \begin{tabularx}{0.7\textwidth}{X} 親愛的醫生路加和底馬問候你們。 \end{tabularx} \\ \\ \relax
4:15 & \begin{tabularx}{0.7\textwidth}{X} 請問候老底嘉的弟兄以及寧法，和她家裡的教會。 \end{tabularx} \\ \\
[1ex]
\hline
\hline
\end{longtable}
$^{1}$我們一起同心祈禱.
很感恩天父讓我們能夠由剛才的聖餐.
我們可以省察,我們可以思想主救恩和我們生命關係.
多謝你讓我們能夠上上,可以有機會去牽手聖餐.
和弟兄姊妹一同去確認和去認識主長闊高深的愛.
今早我們藉著這一段最後的文案.
通常在保羅書信裡面都會有的文案.
裡面的描述是非常有感情.
和真的有一番給我們去研讀.
很多當時的童工團隊他們生命的表現.
特別是為福音那份擺上.
是很值得我們後世的教會一同去學習.
我們今早在你面前等候.
奉耶穌基督的名,阿們.
團隊的文案意味著.
不是保羅自己一個人.
是一班童工跟隨保羅在不同地方傳道.
小徒在哥羅西羅馬坐第一次監.
他來寫這封歌羅西書.
嚴格來說,哥羅西教會不是保羅建立的教會.
但是他了解了似乎是以巴弗和阿尼西姆在羅馬認識保羅.
而知道哥羅西教會的情況.
面對著一些屬靈的問題.
所以保羅就寫這封信.
而且是透過推基古將這封信帶回哥羅西教會.
讓他們去了解.
在行文裡面,去到第四章.
保羅就藉著很多童工.
保羅描述他們是怎樣的童工.
是怎樣的為人.
對福音是怎樣的人品.
這樣來問安.
激勵哥羅西教會.
收到這封信後.
他們會感受到.
到底原來保羅的事奉和他們的童工的事奉.
是怎樣興旺建立教會.
所以我跟大家一同在這幾組裡.
去看到這些童工是怎樣的為福音擺上.
我們一同進入文安裡面.

$^{41}$抱歉,第十五節提到的女門徒靈法.
不夠時間講,所以我沒有準備.
靈法是一位當時很有錢的希臘婦女.
家裡有教會.
當時有這樣的習慣.
我們看,這裡有大約四組.
我只能選取頭三組.
推基古,阿尼西姆.
或者現在的新翻譯叫阿尼西謀.
以巴弗.
第二組就是阿里達古,馬可.
馬可就是寫馬可福音的馬可.
耶穌.
其實他叫尤士都.
第三組就是路加和底瑪.
路加醫生就是路加福音.
和寫道行傳的新約的作者.
新約寫最多書卷的.
按章數來說.
路加醫生是寫得最多的.
是一個被神用的學者.
當然最後就是靈法.
我們今晚回去看看靈法.
保羅如何稱讚他.
我們由第一二三組去看文案.
保羅的報道團裡面.
團隊成員的描述.
我們先看第一組,推基古,阿尼西謀,以巴弗.
他們的表現如何.
有簡單的經文.
我高亮了紅色的.
保羅想強調.
推基古是我親愛的弟兄.
在主耶穌的救恩中.
忠心的僕役,即是神的僕人.
忠心的.
保羅稱讚年輕的推基古.
王一同作主的僕人.
都是神呼召的僕人.
他要把我一切的事都告訴你們.

$^{81}$幫大家做一點功課.
因為在不同地方描述推基古.
不斷說他在保羅身邊.
把保羅的消息帶給不同教會.
他是一個傳遞訊息的同工.
忠心.
忠於保羅現在的狀況.
去告訴教會.
那時候沒有WhatsApp.
一個訊息可能要半年.
還要坐船很危險.
才可以傳到對方的地方.
所以要說明白.
遭遇現在的需要.
其實很重要.
所以這些傳遞訊息的工人.
其實他們的責任都很大.
又很能夠做一個串連.
相信大家都知道.
馬拉松的故事.
跑到馬拉松後.
把一個打勝仗的訊息說完就死了.
紀念這個人成為了.
日後馬拉松的一個運動.
可以想像推基古是一個怎樣的人.
我特意打發他到你們那裡去.
好讓你們知道.
我們的情況.
又讓他安慰你們的心.
推基古扮演一個這樣的.
傳遞訊息.
好訊息的一個工作.
而且是保羅信任的同工.
下面還有提摩特後書.
去到很後了.
保羅要離世了.
要往已不所去.
當時的小亞細亞的地方.
傳遞工作是這樣流通.
他們去不同地方傳遞.

$^{121}$保羅關心他的足跡.
他的心靈.
不只是停留在某個教會.
而是整個亞洲.
是保羅自己的牧區工場.
推基古也是被打發的年輕同工.
去到保羅要他傳遞訊息的地方.
推基古終於他所託.
去到這些不同的地方.
提摩特後書第四章.
其中一個重點.
就是已經有很多人變節.
底瑪已經貪愛現今的世界.
就離棄我.
往帖撒羅尼迦加去.
有些人去了不同的地方.
但推基古始終在保羅身邊.
是保羅信任.
是保羅打發.
是保羅要他將重要的訊息.
講解給當地教會.
這些串連和聯繫是很重要的.
親愛的弟兄姊妹.
我們的訊息.
是在傳遞什麼訊息呢.
我們在傳遞一個.
很好讓人有成長的訊息.
折毀人心的訊息呢.
我們在某程度上扮演武堂.
子堂和嘉欣堂的角色.
很多時候動員弟兄姊妹.
看到情況很好感恩.
將很多我們的需要帶回去.
下午我們開堂董會.
很重要的.
眼光,視野.
前瞻性的抱負.
有沒有大比教會一起去想.
一起去望,一起去發展呢.
因為你已經局限了你自己.

$^{161}$那種心態呢.
有一個訊息很重要.
要有屬靈的眼光.
推基古肯定有保羅的世界觀.
那個教會觀.
那個擁抱周圍都是牧場,牧區的心境.
很重要的.
否則我們就會偏狹於我們自己.
所以現在很多人說的圍爐取暖.
其實是不健康的.
只是說你自己安好就算了.
你不會想到其他地方,其他人的需要.
基督的教會,基督的工人就不應該這樣了.
推基古有這樣的特質.
值得我們去想保羅如何來稱讚他.
亞歷西茂.
亞歷西茂的舊翻譯.
你看這個腓利門書很清楚.
哥羅西人.
亞歷西茂其實是一個奴隸.
這個奴隸走到羅馬.
輾轉要坐牢認識保羅.
保羅知道.
帶著亞歷西茂信了耶穌.
一個奴隸成為了基督徒.
保羅就很清楚當時奴隸制度的弊端.
奴隸如果現在脫離了主人.
一方面就是一個亡命之徒.
但是保羅基於基督的愛.
就打發亞歷西茂回去哥羅西那裡.
去找他的主人.
他的主人就叫做腓利門.
腓利門書是同時間.
保羅寫的在監獄裡面的書信.
所以亞歷西茂就要帶著一個道歉.
回去他的主人的家.
到底是吉還是空是沒有人知道的.
腓利門最慘就是傳道人.
傳道人如何理解一個奴隸偷了你的東西.
又跑去羅馬.

$^{201}$現在保羅叫他回去.
不過在這裡.
保羅如何稱讚亞歷西茂.
他沒有做什麼,他是一個奴隸.
剛剛信了耶穌的小弟兄.
他不是推基督的一個傳道.
我又打發一位親愛忠心的弟兄亞歷西茂和他.
他們也是你們那類的人.
所以第一組的推基督.
亞歷西茂和伊巴忽.
全部三個都是哥羅西教會的人.
對哥羅西教會有很大的一份感動.
這幾個人.
有一個奴隸偷了東西.
被人抓到.
在監獄裡信了主.
生命有改變.
福音如何能夠覆蓋這個改變呢?.
有沒有一個生命的故事.
可以讓大家感受到這份基督的接納.
寬恕和愛呢?.
當然要看腓利門.
他如何去看他自己這一份.
會不會是一口氣.
會不會是他一直很生氣.
自己的面子都沒有了.
在自己的同伴面前.
這樣被亞歷西茂奚落.
這些全部是應該要問的問題.
不過保羅這樣稱讚.
在哥羅西教會面前.
這樣稱讚亞歷西茂.
他是你們那類的人.
他們這三個人.
要把這類一切的事.
都告訴你們教會知道.
保羅說亞歷西茂是親愛忠心的弟兄.
我們不是太多知道他日常的生命改變.
沒有機會教他初信班,朕黎班.
也沒有聽過他的見證基督教故事.

$^{241}$但是保羅很快就看到他的生命改變.
他是親愛忠心的弟兄亞歷西茂.
所以救恩可以帶來這麼戲劇性的生命改變.
第三個,有一位你們那類的人.
再一次提到這三個人都是哥羅西人.
他們是如何來慰祖.
保羅認識以巴弗.
對他的為人,其中一件事.
就是他是一個很祈禱的人.
在教會裡面有一個很祈禱的人.
他的祈禱保羅形容.
作基督耶穌僕人的以巴弗問候你們.
同工在哥羅西教會.
輾轉又去到羅馬要坐牢.
他禱告的時候常為你們竭力.
竭力可能根本就是用心用力用時間.
相信不是五分鐘.
相信是用很多時間為哥羅西教會的弟兄姊妹祈禱.
開聲提名的祈禱.
祈什麼呢?.
願你們哥羅西人.
站穩在真理.
而且是成熟的基督徒.
成熟就是充分確信神一切的旨意.
就是你祈禱的,你要確信和明白.
而在明白裡面去求神的旨意.
成就在你或教會的事工上.
這樣的祈禱一句的描述.
就可以想像到保羅對以巴弗的形容.
他是一個幾個成熟的祈禱人.
這個時代有沒有祈禱人呢?.
有時還要再問一個更尷尬的問題.
這個時代有沒有祈禱的人呢?.
祈禱人是一個已經進入生命成熟的事奉崗位.
看祈禱是看得很緊要.
但今天祈禱的人.
這樣吧,我們星期三晚上.
現在我們透過科技的在線.
弟兄姊妹準時上線.
不是在乎人多.

$^{281}$在乎我們看到禱告的寶貴.
禱告的重要.
一起為神家,神國守望.
所以不是在乎人數.
在乎我們有沒有識見.
看到原來祈禱人是這個時代最大的欠缺.
很需要祈禱人.
但需要由祈禱的人變成祈禱人.
你都不祈禱.
你說自己是祈禱人.
其實是不可能的.
以巴弗能夠願賀羅西教會.
站穩和成熟.
還要充分的確信神一切的旨意.
否則祈禱來做什麼呢?.
祈禱就是要在神的旨意裡面去尋求.
以致能夠認識.
所以這道求主讓我們能夠看到.
大家看看這三個人.
是哥羅西教會的弟兄.
兩個是傳道人.
一個是奴隸.
但已經由奴隸.
當然固然繼續是奴隸.
由奴隸成為神的兒女.
在主裡面的弟兄.
所以.
如何看待我們這個身份呢?.
歌羅西書三章.
保羅自己寫的教義.
就是這樣很具體的有人去落實.
在文安裡面我們已經看到.
保羅真的這樣理解這些真理.
在這世上.
並不再分希臘人和猶太人.
受國禮的,未受國禮的.
未開化的那些當時的巴比利人.
那些叫西古提人.
遺奴的自主的奴隸.
一些主人.

$^{321}$自由人.
不過記住.
你有你這個地上的身份.
你有你這個地上的人生際遇.
不過唯獨基督是一切.
又在一切之內.
很清楚告訴我們.
在救恩裡面.
我們是成為一.
我們有我們地上的身份.
我們有我們地上的不同的崗位.
但是在基督裡面.
保羅理解.
每個人都是神所愛的兒女.
所以這裡很重要.
如何看待我們的身份.
我們在基督裡面的新身份.
這是很重要的.
如何去斷定我們自己.
屬於誰.
我們怎樣的.
誰的性質.
這是很重要的.
在基督裡面.
身份的尋索.
弟兄姊妹.
對自己的身份.
有沒有一個遺憾.
有沒有一個惋惜.
有沒有因為成長.
而帶給你很多.
對家庭.
對自己.
有很多的憤怒.
我很多時候在初信班有說.
我信主一出生在九龍塘.
不過隔離火車橋的母區.
一出生就有十幾人在我家.
全部是我的哥哥姐姐.
你怎樣成長.

$^{361}$家裡有兩個媽媽.
兩個.
你怎樣成長.
成長就有很多很多的問題.
很多很多的掙扎.
很多很多的哀殺.
這樣去塑造我這個人.
所以成長階段.
是有很多憤世嫉俗.
所以這個我是誰.
是我一直都問自己.
但是自從信了主.
整個生命改變.
給大家一個例子.
這位還是現任的.
韋爾比主教.
聖公會最大的主教.
有什麼特別呢.
原來他在成長之後.
他以前是一個石油公司的高級管理人員.
後來就蒙召做傳道.
做主教.
後來去到.
成為了聖公會的總主教.
有人告訴他.
咦 你可能不是你爸爸所生.
怎會呢 難道你媽媽是甚麼.
原來他真的聽聞很多謠言.
說他現在的親生爸爸.
不是他爸爸.
他就去做DNA的檢查.
他想證明自己的身份.
怎知.
終於真相揭露.
原來是邱吉爾的私人秘書.
就是他爸爸.
當年他的媽媽.
就跟這個秘書有婚外情.
生了他.
當然周圍的人就晴天霹靂.

$^{401}$這個韋爾比主教.
他說 沒錯.
物質的DNA的確不是現在爸爸所生.
但我的身份.
我的存在.
不是用物質的DNA.
我在主耶穌裡面.
我接受我自己是一個私生子.
上一代的錯誤.
我不需要付那個代價.
我仍然是我母親所生的一個韋爾比.
這個人.
但我在主耶穌裡面.
我有一個新的身份.
很不容易的.
因為你是一個名人.
很難接受到自己這樣的身份.
被揭露出來.
所以一句名言.
不是DNA決定我的身份.
是在主的救恩裡面.
決定我的身份.
所以剛才我們看的第一組的.
推基古 亞歷西茂 以巴弗.
三個哥羅西的人.
就讓我們看到保羅如何推廣.
救恩的普遍.
每個人都在主的救恩裡面.
成為神的兒女.
成為一.
不會再有分別.
記住.
一至六是有分別的.
因為一至六可能是穿西裝 打領帶.
有些人是送外賣的.
他可能很辛苦 出汗.
沒錯.
一至六真的有不同崗位的.
但記得來到教會.
我們就是神的兒女.

$^{441}$這個信息是要很緊要牢固記住的.
所以你是CEO嗎.
一樣你可以做師事.
派程序表.
我們在感受堂真的有CEO派程序表.
很謙卑地你不知道他是CEO.
因為他來到星期天.
他是神的兒女派程序表.
是其中一個師事.
這樣的學習真的很美.
每個人來到神的面前.
就是神的兒女.
不再分誰是誰.
不再分猶太人 希利尼人等等.
我們看第二組.
阿彌達古 馬可 耶穌.
令我一路研究的時候更加感動.
因為保羅怎樣稱呼他們三個呢.
我們看下去.
讓我一同坐牢的阿彌達古問候你們.
很正常.
不過留意.
他是帖撒羅尼迦人.
有一隻亞大米田的船.
這是小行傳的二十七章.
很後期了.
開往亞細亞沿海一帶去的時候.
我們上了船.
當然這裡數了一大堆人.
我只是引述阿彌達古.
有馬其頓的帖撒羅尼迦人.
阿彌達古和我們同去.
因為你要留意.
這一組人跟保羅一樣.
是奉猶太教的.
稱為猶士都的耶穌.
也問候你們.
馬可我就沒有引述了.
奉國禮的人當中.
只有這三個人.

$^{481}$是為神的國與我作同工的.
也是使我的心得到安慰.
只有這三個人.
肯定不是.
當時的猶太人傳道者.
約翰還在.
你當我死了嗎.
他們是猶太人.
彼得也是猶太人.
為什麼要稱呼阿彌達古.
馬可和耶穌.
這三位年輕人.
奉國禮.
不是猶太人.
不過奉國禮.
不是天生的猶太人.
他們奉國禮.
也很嚮往猶太教.
但是.
記住.
他們不是守猶太教的規例.
而是他們信了主耶穌.
很清楚自己.
是神的兒女.
但是他們又願意.
跟隨保羅.
來奉國禮.
這樣來認同使徒保羅.
保羅肯定是猶太人.
肯定是第八天.
他自己說的.
第八天行國禮.
他是汴雅敏茲派.
即是皇室的後裔.
汴雅敏茲派的.
一個猶太人.
他在神的律法裡.
他是一個法利賽人.
保羅這樣稱呼自己.
不過後來當他遇見主耶穌.

$^{521}$他要去傳講一個.
外邦人得救的福音.
他就要告訴.
很多外邦人.
你可以因信就清二了.
但是這樣的訊息.
就帶來猶太人.
到處迫害保羅.
你是我們同胞猶太人.
你說外邦人.
可以無條件因信就得救.
這個訊息是猶太人.
不能夠接受的.
到處去攻擊迫害.
和要控告他.
所以哲行主是控告.
保羅去到羅馬.
要他去坐牢.
要上到該殺的面前.
受審判.
所以猶太人.
這樣對保羅.
保羅是很明白.
自己因為有一個.
這麼神聖的訊息.
表面看就是便捷.
就是不再叫人信猶太教.
不再叫猶太人.
信猶太教.
更加讓外邦人.
白白就可以得救.
只是你口裡承認耶穌.
心裡相信他復活.
你就可以得救了.
這樣的一個這麼簡單的訊息.
但是.
當保羅在寫文安的時候.
在他身邊.
有三個.
奉割禮的人.

$^{561}$不是猶太人.
但是奉割禮.
心明保羅的困難.
那個.
內心的傷痛.
因為猶太人.
不明白自己.
以為保羅是便捷.
背棄了猶太教.
所以這三個人.
是能夠讓保羅.
心裡得到安慰.
就是說他現在.
傳說的這個.
恩信稱義的道.
居然有人明白.
而且明白得.
跟他同宮.
奉割禮.
馬可更加是巴拿巴的表弟.
耶穌.
猶時都本身就是奉割禮.
耶穌本身.
就是希臘文的耶穌.
耶穌的名字.
舊約就是約書亞.
這樣的三個人.
在神的國裡.
跟保羅同宮.
只有這三個人.
當然在他身邊.
只有這三個人.
或者相對於.
下面說的那組.
第三組的外邦人.
這三個奉割禮.
更加令到保羅的內心.
有一個共鳴.
有一個精神的支持.
又想起一個例子.

$^{601}$很值得和大家想一想.
這三個人奉割禮.
跟保羅很親的背景關係.
就有這樣的共鳴感.
大家看到海特王妃.
今年患癌.
剛剛完結了化療.
她擁抱的是一個16歲的少女.
本來想做導演拍攝.
而且非常有拍攝天份.
她的作品很好.
不過今年.
這個16歲少女.
今年1月.
2024年1月.
確診了一個.
我不知道名字怎麼說.
名字很長.
一個癌症.
絕症.
醫生說她只有半年至兩年的壽命.
這個少女很積極.
列了一個清單.
叫做維園清單.
她媽媽很心痛.
就把這份清單.
和女兒的故事擠上網.
海特王妃無意中.
看到這份清單.
很感動.
就邀請她去.
溫莎堡.
去接見她.
就是這個擁抱.
幫她達成.
其中一個心願.
因為他們有一個拍攝.
一些運動員.
很出色的英國運動員.
受婚儀式.

$^{641}$得到一個婚銜.
邀請這個16歲的少女.
拍攝.
拍攝完之後.
邀請這個少女.
去她家內室.
家庭的私人聚會.
和她分享.
海特王妃自己的.
抗癌的心路歷程.
讓這個16歲的少女.
有一份.
很大的感動.
一個絕症的少女.
可能她的才華.
不能夠再展露.
但她在這個維園清單裡.
又達成了一項.
海特王妃幫她達成.
親愛的弟兄姊妹.
在我們的事奉路上.
有沒有感覺孤單.
感覺一些困難.
或者我們做基督徒.
我們自己也有很多.
感覺到的迷惘.
保羅說這三個.
奉國禮的人.
使我的心得暢快.
與神同工.
和我一樣.
為神的國同工.
保羅不覺得自己.
原來走錯路.
其實他不是走錯路.
是猶太人不明白他.
他一直叫人不要再信猶太教.
不需要守國禮.
不需要安息日.
只要信耶穌就行了.

$^{681}$很難想像猶太人.
有多憤怒保羅.
但有三個奉國禮的.
不是猶太人.
奉國禮的猶太人.
愛保羅.
信任保羅.
跟隨他.
成為他心裡得到暢快.
神有時會給我們.
有不同弟兄姊妹.
出現在我們身邊.
給我們心裡有安慰.
我們不要拒絕這些安慰.
不要自視自己是強人.
不需要別人的安慰.
千萬不要這樣想.
不要這樣做.
一定要與人為善.
讓不同的同工.
在我們身邊出現.
下星期.
有一位很高的西人.
他是丹麥.
瑞典.
瑞典前財務司長.
財務司長.
但後來做了傳道.
不再做財政司的.
師父的司長.
做了傳道.
他最近回到香港.
一定要見我探望大家.
下星期他就會出現.
他身高約六尺.
是瑞典人.
是我們一直有支持的.
Z7的總幹事.
歐洲區的總幹事.
下星期會來.

$^{721}$因為他知道我去了紅欣服侍.
本來他要在錦獸.
但他說不要.
把時間給我們.
下星期.
就會和我們一起分享Z7的發展.
下星期再說.
不要拒絕.
同工同路人的安慰.
保羅.
被這三位弟兄.
感動安慰.
最後這一組.
我們看看.
很簡單的一句文案.
親愛的醫生路加.
和底馬文侯.
一句.
親愛的醫生.
弟兄姊妹.
我們需要看到.
在主耶穌裡面.
希臘人.
猶太人.
沒有再分這個身份.
路加醫生.
沒有說他是哪一個民族.
小腿行主.
十六章.
是不是那個.
馬其頓的異象.
就是路加醫生.
很多的神學家.
這樣推論.
但聖經從來沒有說.
讓我們知道的就是.
優美的希臘文.
路加福音和小腿行主.
是最多寫新約的.
雖然是兩卷.

$^{761}$但最多就是路加醫生.
由十六章開始.
一直跟隨保羅.
到提摩大學後書的四章.
都是跟隨保羅.
默默的.
記錄著保羅傳道的事情.
小腿行主十六章開始.
一個很明顯.
很大的轉向.
就是不再說.
他們.
用第一身說我們的傳道.
一直到二十八章.
這是小腿行主的特色.
也是路加醫生.
對保羅的品格.
對他的傳道的高尚.
和他為主耶穌奉獻的.
單純的心.
深深的折服.
成為了保羅的跟隨者.
去記錄保羅.
傳道.
中心地面對什麼事情.
他的逆境和順境.
路加醫生.
如數無遺地記錄.
這是路加醫生.
親愛的.
醫生路加.
一句.
到提摩大學後書四章.
剛才看的第十一節.
獨有路加在我這裡.
一句.
自稱路加就是一句.
但是.
我們的傳道.
我們的人品.

$^{801}$我們的志向.
有沒有激發一些有心人.
有能力的人.
有智慧的人.
投放他的人生.
為你所做的傳道職業.
來擺上.
有沒有感動到這樣的人.
成為你的跟隨者.
同路人.
這是一個指標.
對一個傳道牧者.
或者你已經是一個很有經驗.
卓有所成的一個領導層.
的一個.
見證指標.
如果你中伴.
是沒有證明你的失敗.
你多成功都是失敗.
因為.
真的人家.
放心下去.
觀察了你很多年.
很久.
你一如既往的初心.
才能夠.
接服這樣的人.
跟隨你.
這是很重要的.
所以我們傳道的人生.
其實應該建立不同的梯隊.
建立不同的下一個梯隊.
建立下一個接班人.
任何一個大機構的主管領導.
都第一時間要想.
我要怎樣去栽培.
更多的同路人.
所親愛的.
醫生,路家.
問你們安.

$^{841}$分享一個.
我自己很受感動的故事.
我有一位老師.
那時候我們在讀四年班.
這是神學系的老師.
因為他已經離開了.
回到美國.
我不說他的名字了.
他講他自己的夢照和讀神學.
讀博士學位.
他提到他有一次.
終於要去美國.
讀博士學位.
那個教授.
看到他的成績.
其實不足夠.
收他成為.
博士生.
但被他的誠意.
各方面感動.
就破格.
收了這個學生.
即是我的老師.
成為博士研究.
過了很多年了.
四五年了.
終於完成了他的博士學位.
也真的順利畢業.
也有聘約.
回到建道事奉.
準備回香港.
我的老師.
他說很簡單.
中國人當然是尊師重道.
他就去教授那裡敲門.
say goodbye.
敲門.
教授坐在那裡.
還沒坐下.
就給他一個信封.

$^{881}$叫他打開.
本來就是多謝教授.
當年.
幾年的指導.
現在順利畢業.
回香港.
準備在神學舊約裡事奉.
他打開信封.
教授叫他打開.
是一張Pedicash支票.
一些現金.
可以兌錢.
他當然是愕然.
我來say thank you.
應該是我送禮物給你.
為何你會有錢給我.
當然教授用英文說.
他很愕然.
不知如何回應.
教授就說.
不知如何回應.
say thank you.
and serve the law.
我的老師很詫異.
他為他預備少少金錢.
在他旅途上.
擔心他會被騙.
他為他預備少少金錢.
在他旅途上.
擔心他會被騙.
給少少金錢.
不是多的.
是少的.
但這個舉措.
令到教授.
令到我的老師.
放在心裡.
成為他事奉的動力.
他很多年都告訴我們.
他的人生.

$^{921}$好像一個伯樂.
讓他在下一個階段的事奉.
他有一個.
好像Mentor.
給他有很大的靈魂啟發.
不只是書本的知識.
是生命的一個知識.
這個老師.
這樣的少少舉措.
我們都要學習.
預備這些信奉.
有些人很失落的時候.
或者他可能有一個新的挑戰.
的時候.
我們可能一杯涼水.
給他喝.
這對他很大的鼓勵.
放在心裡.
成為他事奉的動力.
親愛的醫生.
路家.
問你們安.
親愛的醫生.
路家和底馬.
問你們安.
親愛的弟兄姊妹.
初期教會不是保羅.
他有神蹟.
他有三層天的經歷.
他講方言.
他又可以其他的東西.
那些是的.
但是更實在的.
他如何提拔後輩.
多少人願意.
見到苦難都跟隨保羅.
成為他的同路人.
其實應該說.
成為主耶穌神國度的工人.
一起服侍.

$^{961}$紅恩家的弟兄姊妹.
我們一同透過這個問安.
我們收了他.
然後我們一起去學習.
在神感動你的不同範疇.
一齊的.
同心的服侍.
我們一起祈禱.
感恩天父讓我們能夠看到.
聖經這麼寶貴的.
這些童工們.
他們那份對福音的熱誠.
在保羅的號召下.
其實應該是在主耶穌.
的選召.
和召命下.
大家在不同的崗位.
在不同的呼召的場景.
來服侍主.
但願主耶穌繼續幫助.
感動我們.
讓紅恩家的弟兄姊妹.
我們在整個洪水橋區.
在新界西.
在我們一堂三點.
我們有錦宣.
我們有佳恩堂.
我們在整個三角形的地方.
讓福音能夠在新界西區.
有人聽明白.
明白.
和在生命裡成長.
加入事奉的行列.
求主幫助.
的分享.
傳遞的訊息.
奉主耶穌基督的名.
阿們.
\newpage



\section{列王記上 17:8-16}
\label{sec:dXTahKH7Xis}
\textbf{2024.10.20 《順從主的話,上帝必供應》 列王記上17\_8-16 (和修版) 講員 : 譚立晶傳道}
\newline
\newline
連結: \href{https://youtube.com/watch?v=dXTahKH7Xis}{\texttt{https://youtube.com/watch?v=dXTahKH7Xis}} ~~~~ 語音日期: 2024-10-21
\newline
\newline
\hyperref[sec:jVaUl1bB0os]{\small{< < < PREV SERMON < < <}}
~
\hyperref[sec:index_chronic]{\small{[返順時目]}}
~
\hyperref[sec:index_scriptual]{\small{[返順卷目]}}
~
\hyperref[sec:xtW6V_2_PsE]{\small{> > > NEXT SERMON > > >}}
\newline
\newline
列王記上 17:8-16
\newline
\begin{longtable}{cl}
\hline
\hline
章節 & 經文 (和合本修訂版)\\
\hline
17:8 & \begin{tabularx}{0.7\textwidth}{X} 耶和華的話臨到他，說： \end{tabularx} \\ \\ \relax
17:9 & \begin{tabularx}{0.7\textwidth}{X} 「你起身到西頓的撒勒法去，住在那裡，看哪，我已吩咐那裡的一個寡婦供養你。」 \end{tabularx} \\ \\ \relax
17:10 & \begin{tabularx}{0.7\textwidth}{X} 以利亞就起身往撒勒法去。他到了城門，看哪，有一個寡婦在那裡撿柴。以利亞呼喚她說：「請你用器皿取點水來給我喝。」 \end{tabularx} \\ \\ \relax
17:11 & \begin{tabularx}{0.7\textwidth}{X} 她去取水的時候，以利亞又呼喚她說：「請你手裡也拿點餅來給我。」 \end{tabularx} \\ \\ \relax
17:12 & \begin{tabularx}{0.7\textwidth}{X} 她說：「我指著永生的耶和華－你的神起誓，我沒有餅，罈內只有一把麵，瓶裡只有一點油。看哪，我去找兩根柴，帶回家為我和我兒子做餅。我們吃了，就等死吧！」 \end{tabularx} \\ \\ \relax
17:13 & \begin{tabularx}{0.7\textwidth}{X} 以利亞對她說：「不要怕！你去照你所說的做吧！只要先為我做一個小餅，拿來給我，然後為你和你的兒子做餅； \end{tabularx} \\ \\ \relax
17:14 & \begin{tabularx}{0.7\textwidth}{X} 因為耶和華－以色列的神如此說：『罈內的麵必不用盡，瓶裡的油必不短缺，直到耶和華使雨降在地上的日子。』」 \end{tabularx} \\ \\ \relax
17:15 & \begin{tabularx}{0.7\textwidth}{X} 婦人就照以利亞的話去做。她和以利亞，以及她家中的人，吃了許多日子。 \end{tabularx} \\ \\ \relax
17:16 & \begin{tabularx}{0.7\textwidth}{X} 罈內的麵果然沒有用盡，瓶裡的油也不短缺，正如耶和華藉以利亞所說的話。 \end{tabularx} \\ \\
[1ex]
\hline
\hline
\end{longtable}
$^{1}$(自我介紹).
今天很想和大家分享主的說話.
我們藉著昔日先知去順從神的話.
而神也藉著先知去將神的話傳遞.
然後其他聽的人都去順從的時候.
他們怎樣經歷上帝的公認.
我們一起去看看.
因為今天很緊湊.
可能我講話的語調速度會快一點.
大家就順從耳朵盡力去聽.
我們先看看上文.
剛才主席幫我們讀了那段經文.
我們說要理解聖經認識聖經的時候.
我們要看看今天所講的經文.
和整個背景發生了什麼事.
我們先看看上文.
為什麼會發生今天這段經文的事呢?.
為什麼以利亞要求神去一個寡婦的地方呢?.
原來他們正正經歷著寒災.
但為什麼會有寒災呢?.
原來他們得罪了神.
但為什麼關不下雨?.
從哪裡開始說呢?.
我們看聖經第八章.
那時候他們正正建好了聖殿.
我們來看看聖經背景.
是誰為耶和華建聖殿的呢?.
是不是大衛?.
不是,是誰呢?.
是他的兒子,所羅門.
因為神說不要大衛來建殿.
因為他要出征,打仗.
神說不要你建殿,但我要你的兒子.
就是要任意的.
所以就找他的兒子所羅門去建殿.
這段經文出自於約季進了聖殿.
所羅門就說了一番話.
接著就是所羅門的祈禱.
其實建殿的原因是什麼呢?.
就是要將上帝的名去高舉.

$^{41}$要大家來這個殿去投靠神.
要去敬拜神.
但所羅門這個皇帝是怎樣呢?.
世上最有智慧的皇帝.
其實他懂得一切.
他懂不懂得人性呢?.
人會不會離開神呢?.
人會不會叛逆呢?.
會的.
所以所羅門王昔日就有這個祈禱.
不如我們一起讀.
「你的民因得罪你.
你懲罰他們使天閉塞不下雨.
他們若向此處禱告.
承認你的名離開他們的罪.
求你在天上垂聽.
赦免你僕人以色列民的罪.
將當行的善道指教他們.
且降雨在地的地.
就是你次級離民遺孽之地」.
所以原來今天所說的經文.
是早早開初.
建元殿所羅門王已經.
一個很重要的祈禱.
已經祈禱了.
原來人會跟隨上帝離開.
為何今次他們會離開呢?.
我們就看看一個比較近的上文.
正正是同一個17章.
那裡說.
「住在激烈的提斯比人爾利亞.
對阿蝦說」.
所以是神要出動.
派他的先知和代言人.
要跟當時的皇帝阿蝦王說.
如果大家熟悉舊約.
阿蝦王已經是很壞的.
但更厲害的是.
背後的女人是誰呢?.
我們說是惡毒的女皇后.

$^{81}$耶洗別.
耶洗別就要特意帶其他巴力先知的敬拜.
進入以色列這個地方.
要其他百姓不要去聖殿敬拜.
太遠了,你們不要過去.
其實這是關乎當時的政權.
剛才我一直聽福音樂曲的獻唱很好.
不知道你們有沒有留意歌詞.
有暴政的.
有些人是驕驚的.
就是在說當時的阿蝦王.
就是這個處境.
他不要上帝.
不認上帝.
還要跟他的老婆.
耶洗別聯手一起.
把巴力先知的敬拜.
帶入整個百姓的會中群眾當中.
所以神就要出手了.
神就派了以利亞跟阿蝦王說.
我指著所事奉的永生的耶和華以色列的神.
這幾年若不禱告.
就不會降露水.
也不會下雨.
這是一個什麼呢?.
審判.
懲治.
教訓.
神要告訴.
他的以色列民.
也要告訴這個惡王.
誰才是神.
不是巴力是神.
所以當時.
神就跟以利亞說.
你離開這裡.
往東去.
躲在約旦河東邊的基納溪旁.
為什麼要離開?.
為什麼要躲?.

$^{121}$我們就看看原文的意思.
躲的動詞.
就是躲藏.
不要讓人看見知道.
但沒有逃避的意思.
離開.
神要透過先知的離去.
象徵神的同在和祝福.
暫時離去.
先知代表神.
即是說.
上帝不在了.
因為你們不聽話.
你們不聽上帝的話.
所以神的福也不在了.
所以不降雨了.
正如神所說.
不會使雨露落在以色列地上.
所以就是審判的來臨.
所以就來到今天這段經文.
我把一些螢光的字高光了.
你們可以深刻地去看.
這個故事的場景.
這個經歷裡有什麼人呢?.
有什麼經歷呢?.
耶和華說.
臨到以利亞.
你起身到西頓的撒拿法去.
住在那裡.
我已經吩咐那裡有一個寡婦供養你.
剛才說的是她去基納溪旁.
基納溪旁.
如果大家有聖經看上文.
那時候神派烏鴉去供養她.
我不知道大家有沒有試過餵鳥.
現在是不准餵的.
但以前我們的年代.
很開心地餵鳥餵魚.
很治癒的.
現在不准餵的.

$^{161}$但你有沒有逆地而處.
你有沒有想過是小鳥餵你.
小鳥餵食物給你.
食物剛好在你面前.
你吃不吃?.
哎呀,好骯髒啊.
但這個神跡正正是在.
以利亞的身上這樣經歷.
但為什麼現在在溪邊要走呢?.
為什麼?.
沒有下雨.
溪邊還有水嗎?.
沒有了.
乾塘了.
沒有了.
所以神要派以利亞去找一個寡婦.
然後神就派寡婦去供養她.
當我一直看的時候.
其實有些不明白的地方.
因為我們說我們要研讀聖經.
看清楚裡面說什麼.
有時候我們要問一下問題.
神為什麼你這麼特別?.
為什麼你不找一個有錢人去供養以利亞?.
為什麼你要找一個寡婦去供養以利亞?.
我一直看到這裡的時候.
我就想起一句.
我想大家都聽過的俗語.
在乞丐的碗裡拿飯吃?.
神為什麼?.
為什麼你有很多東西?.
為什麼你不找一個富翁?.
不找一個財主?.
為什麼你要找一個寡婦去供養她?.
究竟上帝你有什麼作為呢?.
但是以利亞很順從.
以利亞就起來了.
就去了.
她在溪邊.
其實她要走好一段路.

$^{201}$才去到撒勒法.
所以她去到城門口.
她說那裡有一個寡婦.
在那裡執柴.
以利亞就呼喚她.
我又有問題要問.
為什麼以利亞一去到城門口.
就知道那個是寡婦?.
你會不會一出門就知道那個是寡婦?.
問問題.
為什麼會這樣?.
原來這個和昔日的文化背景有關.
以前的女人通常在哪裡?.
在家裡.
其實不常外出.
就算外出是怎樣?.
三五成群.
一起外出是他們的風俗.
如果大家看過新約聖經.
通常女人一起出去會做什麼?.
打水.
很厲害.
很熟悉聖經.
但是現在無端端就看到只有一個女人.
還要做執柴粗重的事.
執柴應該由誰來做?.
男人.
就是因為她家裡沒有男人.
但是家裡也要顧.
她就要一個人出去執柴.
所以是一個很不容易的處境.
又沒有雨.
沒有雨代表什麼?.
農產品失收.
沒有糧食.
以前有個很好的制度.
叫人們種田不要收掉.
要留一點.
讓其他孤兒寡婦可以去執.
但是現在農作物又失收.

$^{241}$大家又不夠食.
又沒有水.
怎麼辦?.
但是神依然叫以利亞去找寡婦供養.
很匪夷所思.
我們繼續看.
以利亞.
你要想像.
走到身水身汗.
已經不是溫柔.
是呼喚.
呼喚是怎樣?.
快點拿個器皿.
拿點水來給我喝.
寡婦就預備去.
去的時候又再叫她.
你手裡有什麼呢?.
拿些餅來給我吃.
那個叫做拿些餅來給我吃.
是在說她身上有沒有現成做好的餅.
如果我現在問.
我今天早上還沒吃早餐.
大家有沒有現成的東西可以給我吃?.
有沒有?.
有,你們可能有.
有帶在身上.
但是呢.
誰知道寡婦怎樣說?.
她說.
我指著永生的油.
說你的神起誓.
我沒有餅.
我沒有現成的餅.
壇裡只有一把麵.
瓶裡只有一些油.
你看.
我去拿兩斤柴.
帶回家.
是為我和我的兒子去做餅的.
我們吃了之後.

$^{281}$我們就等死了.
這樣.
其實寡婦都是很老實.
將她的情況說出來.
那個一把麵.
那個一些油是在說剩下多少呢?.
可能就是在說那一頓.
一把麵.
就是在說那些麵粉.
可以做餅.
然後有些油.
這裡呢.
有沒有人會做餅的?.
會做蛋糕的?.
會烘焙的?.
我就不懂的.
因為我覺得工作很繁複.
我就查了一下.
為什麼餅和油有什麼化學作用呢?.
原來有的.
如果我說錯了.
你們就補充我.
原來麵粉裡加上油.
會令到麵粉團.
會更加脹大.
更加飽肚.
更加柔軟.
可以充飢多一點.
是這樣的.
所以他們是不可或缺的.
寡婦就將她的苦情說了出來.
以利亞就跟她說.
其實以利亞都很淋漓盡致.
很想吃東西.
不用怕.
你照你所說的去做吧.
只要先為我做一個小餅.
還是先做給我.
拿給我吧.
然後你再為你的兒子做餅.

$^{321}$我們有沒有覺得他很咄咄逼人.
嘩 上帝 為什麼你這樣.
以利亞 為什麼你這樣.
究竟你想怎樣.
但我想到這個時候.
我覺得很奇妙.
神又令我想起一些.
有關係的經歷.
我問問你們.
你們覺得像不像.
你們有沒有經歷過.
上帝要你傳福音給.
比較貧困基層的人.
你會鼓勵他做一些突破.
你會鼓勵他去經歷神.
要他放下一些東西.
有沒有試過.
但可能那些微信者覺得.
嘩 你為什麼還要.
拿我的時間 拿我的東西.
你是不是幫我.
但又很奇妙.
當那個人願意相信.
倚靠的時候.
他又會經歷到.
原來他願意經歷神的時候.
他又會經歷到上帝.
無微不至的看顧.
我們去做傳道人.
我們去做宣教士.
我們有一個責任義務.
和權利是什麼呢.
就是將工場裡的需要.
帶回來給各位.
但有時候我們都尷尬.
因為什麼呢.
有時候我們都知道.
可能我們在座.
有些弟兄姊妹都艱難.
可能有時候我們都知道.

$^{361}$在座的弟兄姊妹.
都不是很富裕.
但我們都要.
上帝給了感動我們.
我們就要說出來.
然後呼喚大家.
為我們一起祈禱.
你可以經濟上支持我們.
就支持我們吧.
還可以出人力.
一起去短孫.
一起去探訪那裡的人.
就更加好了.
是不是.
但有時候你說.
我們這個身份都有點尷尬.
好像以利亞一樣.
咀嚼咀嚼.
但上帝要我們經歷什麼呢.
為什麼以利亞可以這麼盡這麼咀嚼呢.
他背後有一個很重要的信念.
他是不是真的可以逢生呢.
有危是不是真的有機呢.
原來神的話很重要.
以利亞說.
因為耶和華以色列的神這樣說.
潭內的面必不容盡.
貧裡的油必不短缺.
直到耶和華洗雨降在地上的日子.
原來以利亞是因為背後這個重要的信念.
他相信神的供應.
相信神的話.
所以他就這樣去做.
他也將這個話傳遞給婦人聽.
十五節就看到婦人也是這樣.
憑信心信從.
婦人就照以利亞的話去做.
她和以利亞以及她家中的人.
是要許多日子.
潭內的面果然沒有用盡.

$^{401}$貧裡的油也不短缺.
正如耶和華說.
直以利亞所說的話.
神奇嗎.
很神奇.
所以有時我看到這些經文的時候.
想起有時我們也不容易走前線.
但原來我們也會經歷到.
將這個消息傳遞出去.
呼喚給其他人聽.
而其他人一起去聽.
一起去回應的時候.
就經歷到上帝的供應是很微妙.
是真的不會盡.
真的不會短缺.
我們看看經文裡有什麼不同的角色和元素.
有耶和華神.
我們在我們生命裡有沒有將神的話放在心裡.
有以利亞這個角色.
被神差派.
神去供養他.
用一些特別的方法.
而神的話也會臨到我們每一個信徒.
不只是昔日的先知.
神也讓我們人有不同的經歷.
有寡婦.
有家人.
兒子.
也有其他元素.
有柴米油鹽.
柴面油.
有家人.
起初我們見寡婦只提到她的兒子.
沒有留意到她有其他家人要供養.
但是因著以利亞去到他們的家.
原來上帝就是要藉以利亞去祝福寡婦的家庭.
為什麼要選擇寡婦的家庭呢?.
我去問這個問題.
為什麼要選擇寡婦的家庭呢?.
其實上帝是知道哪些人是敬虔而靠他.

$^{441}$所以上帝知道這個婦人的敬虔.
所以他知道這個婦人會聽上帝的話.
所以上帝也知道寡婦的苦情.
所以他差派以利亞去供養寡婦.
看似是寡婦供養以利亞.
但實際上是上帝供養以利亞,寡婦,兒子和家庭的人.
很奇妙,整件事逆轉過來.
所以有時候我們要求主給我們莫大的信心.
有時候我們看不清這個環境.
看不清這個經濟.
看不清各種政權的轉變.
但是我們要認清我們的身份.
我們是神的子民.
我們是屬於神.
我們要敬虔,我們要順從.
原來我們的順從是會幫助到我們家庭裡的其他人.
都可以經歷上帝.
求主去幫助我們.
我自己說一點見證分享.
昔日我都覺得我是被神差派的人.
差派到我家裡.
我媽媽雖然有一個在生的丈夫.
但都如同一個寡婦.
為什麼呢?.
因為我之前說過.
我爸爸都很喜歡去澳門.
但那時候他不是去澳門宣教.
去澳門做什麼?賭錢.
我有爸爸就好像沒有爸爸.
我媽媽有丈夫都好像沒有丈夫.
但神有什麼能夠臨到我呢?.
他呼召我進入一個全福音訓練的中心.
受訓了兩年.
受訓了兩年是說什麼意思呢?.
就是說我是家裡的大家姐.
我去了奉獻,做裝備.
我有沒有錢拿回去給家裡?沒有.
但憑信心去做.
我媽媽也讓我去做.
但那時候我媽媽是未信主的.

$^{481}$到我讀完短孫的時候.
我就收到上帝另一個說臨到我.
你去讀神學.
哇,讀神學是說什麼?四年.
那怎麼辦呢?.
媽媽讓不讓我又再四年.
沒有經濟的支持.
你知道以前的大家姐要養家的時候.
是不能做自己想做的事.
就是怎樣?出去工作.
然後賺錢拿錢回去養其他家人.
我就祈禱上帝如果你真的感動我去讀神學.
你感動我媽媽,雖然她未信主.
我很深刻記得.
我有感動就帶我媽媽去參加一個大型的報道會.
去到大型報道會之後一起聽.
其實很感動.
我一直為我身邊的媽媽祈禱.
神啊,你真的讓我媽媽可以信主.
讓我媽媽也經歷到你的真實.
我很希望我媽媽也經歷到.
我記得那兩年受訓去傳福音.
我經常把我的經歷告訴媽媽.
讓她知道人信了主之後會怎樣.
會怎樣經歷到改變.
很奇妙的,那天報道會.
當講完呼召,有沒有人願意信主呢?.
我媽媽舉手.
哇,好感恩啊.
我說媽媽你走不走出去?.
不用走,我坐下來跟著講就行了.
我說你認真心跟著講吧.
很開心,信了主.
但沒想到隨即講完.
馬上做了全時間的呼召.
通常報道會很少.
你說陪靈粉兄會.
就說會呼喚人去奉獻.
誰知在報道會裡.
隨即帶完人信主.

$^{521}$隨即又呼喚有沒有人願意為神奉獻.
哇,那一刻我就七上八落.
媽媽剛剛在旁邊信主.
突然我就說神要你女兒再去讀四年神學.
哇,好大挑戰.
但我很想起來,坐在椅子上.
很想起來.
但我又怕會刺激到我媽媽.
我祈禱,祈禱.
誰知我媽媽就在旁邊看著我.
看著我.
起床吧.
起床?.
是啊,我知道你想的.
你是不是想?.
我說我想啊,但可以嗎?.
要讀神學讀四年.
可以了,起床吧.
哇,那我就起床了.
哇,一起床.
心就不停地跳.
我感受到聖靈的感動和催迫.
哇,神真的很奇妙.
幫了我這個做女兒的.
幫了我這個寡婦的媽媽.
但我們一家人就是經歷了上帝.
由兩年的短孫終生訓練.
到四年神學訓練.
然後又再有些突發事發生.
我再讀教務輔導.
讀完道石才離開.
但我們一家人無穿無濫.
是甚麼作為?.
是上帝的作為.
求主幫助我們.
順從祂的話.
上帝必然供應.
求主讓我們打開耳朵.
聆聽上帝藉著不同的僕人.
去呼喚我們做的每一件事.

$^{561}$我們竭力去做.
上帝會供應我們,會幫助我們.
我們一起同心同意.
慈悲仁愛的天父.
因為你是我們的神.
你是我們的主.
你是我們的王.
求主讓我們可以真真正正地.
將你放在我們生命當中的首位.
相信你的同在.
相信你的供應.
相信你讓我們去做的每一件事.
一定是有你的心意.
不要讓困難.
不要讓逆境.
不要讓挑戰.
難阻我們向前為你奔跑.
天父求你幫助我們.
在這個世界,在這個世代.
特別是現在香港經濟越來越差的環境裡.
幫助我們去回應你的呼召.
你感動我們去關心其他基層的人.
其他街坊.
或者其他家裡有需要的.
天父求你使用我們每一個.
我們誠心將我們心裡面.
大大小小的憂慮.
我們會怕的東西.
我們都向你陳明.
就好像寡婦一樣.
她就是向你陳明她自己的苦況.
是的,大家都有難念的經.
但我們相信上帝.
你容得著我們每一個的時候.
我們其實是很榮幸.
很榮譽的可以與神同工.
我們誠心去獻上禱告感恩.
求你給我們有力量.
給我們有信心為主多作供.
奉主名求,阿們.

\newpage



\section{路加福音 8:1-3}
\label{sec:xtW6V_2_PsE}
\textbf{2024.11.03 《跟隨基督的腳蹤》 路加福音8\_1-3 (和修版) 講員 : 馮浩鎏醫生}
\newline
\newline
連結: \href{https://youtube.com/watch?v=xtW6V_2_PsE}{\texttt{https://youtube.com/watch?v=xtW6V\_2\_PsE}} ~~~~ 語音日期: 2024-11-04
\newline
\newline
\hyperref[sec:dXTahKH7Xis]{\small{< < < PREV SERMON < < <}}
~
\hyperref[sec:index_chronic]{\small{[返順時目]}}
~
\hyperref[sec:index_scriptual]{\small{[返順卷目]}}
~
\hyperref[sec:lQqzMqR_pN8]{\small{> > > NEXT SERMON > > >}}
\newline
\newline
路加福音 8:1-3
\newline
\begin{longtable}{cl}
\hline
\hline
章節 & 經文 (和合本修訂版)\\
\hline
8:1 & \begin{tabularx}{0.7\textwidth}{X} 過了不久，耶穌周遊各城各鄉傳道，宣講神國的福音。和他同去的有十二個使徒， \end{tabularx} \\ \\ \relax
8:2 & \begin{tabularx}{0.7\textwidth}{X} 還有曾被邪靈所附，被疾病所纏，而已經治好的幾個婦女，其中有稱為抹大拉的馬利亞，曾有七個鬼從她身上趕出來， \end{tabularx} \\ \\ \relax
8:3 & \begin{tabularx}{0.7\textwidth}{X} 又有希律的管家苦撒的妻子約亞拿，和蘇撒拿以及好些別的婦女，她們都是用自己的財物供給耶穌和使徒。 \end{tabularx} \\ \\
[1ex]
\hline
\hline
\end{longtable}
$^{1}$聽節目早晨,很開心,也跟一些長輩說聲早晨,我見到很多長輩和有些在這裡獻詩,心裡充滿感恩,今天很特別的日子,也是洪恩家,洪恩堂十三週年的堂慶,我覺得很榮幸今天表演能夠分享整個話語..
剛才大家所聽到的經文是在路加福音第八章,很多時候我們去到路加福音第八章的時候,就會留意那個題目叫做殺種的比喻,很多時候我們就會很快跳到殺種的比喻這個故事裡面,但是我很喜歡路加福音,有路加醫生,路加醫生和我是同行,特別喜歡他..
除了這個之外,路加的特別是因為路加很多時候留意一些小人物,是其他福音書沒有記載的,所以今天我們跳到殺種的比喻之前,我想今天特別留意我所聽的經文,是第一節到第三節,這三節的經文..
當時路加記載的處境是耶穌周遊各城各鄉傳道,當然他說有十二個門徒,很出名的,可能是在前頭大練的,有十二個門徒和耶穌一起..
但是路加特別說除了這十二個門徒之外,還有一些小人物,教會也有很多小人物,但是今天路加特別說我們要留意這三個小人物,當然不止這三個,還有其他婦女,有三個人名特別提及的..
第一個就是瑪利亞,我不知道大家中間有沒有Mary瑪利亞,第二個叫約阿娜,約阿娜中文可能不太明白,Joanna就是英文..
第三個就是蘇薩娜,就是Susanna,他說這三個人物都是跟隨基督,服侍基督,和他周遊各城各鄉..
我想大家先用想像力,不是寫記載,是想像一下,你在他們當中,你剛剛在一個城鎮裡面聽完耶穌講道,耶穌準備離開了,去另一個城鎮,還說耶穌不要走,用現代版,我們先拍一張合體照,Selfie,.
如果真的要合照的話,大家可以想像站在中間的是誰呢?應該就是耶穌基督,耶穌站在那裡拍照,我想猜猜站在耶穌的左邊和右邊又是誰呢?當然不知道,我沒有說,只是猜想..
大家記得馬可福音第十章裡面有兩個門徒,一個叫約翰,一個叫雅各,當耶穌說我要快受死,我會準備釘十字架的時候,當時門徒所問及的問題不是關於耶穌受死是什麼,當你載來得榮耀的時候,讓我一個坐在你的左邊,一個坐在你的右邊,其他十個門徒覺得很開心,不是,很生氣,你以為你是誰,一個要坐左邊,一個要坐右邊..
門徒跟隨了耶穌三年,他們最關心的不是基督的受死,最關心的是他們的排名先後,他們的明星有多少,所關心..
所以今天我很開心來到路加坤第八章裡面,路加說你要留意這三個小人物,怎樣跟隨基督,也給我自己心裡面有大啟迪,從這三個人物裡面有什麼學習..
第八章第一節,耶穌周圍各城各鄉傳道的時候,有一個瑪利亞,這個瑪利亞是甚麼呢?他說有稱為勃大拉的瑪利亞,甚麼叫勃大拉的瑪利亞呢?.
對不起,這個螢幕不是現在拍的,可以放了,因為我今天沒有powerpoint,看經文就可以了..
甚麼叫勃大拉的瑪利亞呢?大家有留意嗎?譬如我的名字叫馮浩樓,我告訴你是soldier Bazaar的馮浩樓,你可能聽不懂我說甚麼,甚麼叫soldier Bazaar的馮浩樓呢?.
我和太太Jenny在巴基斯坦哈拉奇一個大城市裡面,一個鎮的地方裡面行醫,做醫生,醫療專家,那個城鎮就叫做soldier Bazaar,那個地方不是有名的,你講洪水橋我也知道洪水橋在哪裡,但是勃大拉是不知道在哪裡來的,我的地方soldier Bazaar沒有甚麼出名,最出名的是樂陀,因為我家門口有個樂陀站,香港沒有樂陀的..
樂陀不是給人拍照下來拍照遊客的,樂陀是當時在巴基斯坦運送貨物,很重的貨物,所以門口的樂陀來休息然後再上路..
soldier Bazaar不是一個有名的地方,勃大拉也不是一個有名的地方,盧家驊說這個是一個小人物,當時我去以色列探望的時候,勃大拉是一個很小的城鎮,當時我坐巴士,導遊說我們來到勃大拉了,還不到30秒鐘,我們離開勃大拉了,這樣你一度已經過了,因為是一個不出名的地方..
盧家驊說這個人不是來自名門望族,不是來自一個大教會,不是來自一個很出名的教會,是一個很小的地方,沒有人留意的,除了之外,他也說這個人的背景,他說曾有七個鬼從他身上趕出來,我們想像這個人的背景是很痛苦的..
如果今天勃大拉說申請參加警察會的話,我們也會想,他是不是正常呢?以前有被鬼附,他行不行呢?他可不可以參與宣教呢?可不可以傳福音呢?可能對勃大拉有很多很多的懷疑..
但我可以想像,可能他不是和十六門一樣在前面帶領,但是我很想像他在人群當中,當他周遊各城各鄉的時候,他跟別人說,你有沒有時間,我跟你說耶穌基督怎樣改變我的生命..
很多時候我們想,資格服侍神,我的才幹,我的學位,我的經歷,甚至我的財富,有沒有奉獻的能力.但今天神提醒我,第一個服侍神的資格,不是才幹,不是能力,不是學位,是一份感恩的心..
瑪利亞因為可以告訴別人,基督改變我的生命,以前我是被黑暗,被魔鬼所附,今天他經歷了整個改變,他可以服侍神.大家記得有一首古老的詩歌叫做,我有一個故事傳給萬邦,大家有沒有聽過這首詩歌,很好聽的..
每個人都有一個故事傳給萬邦,你今天坐在這裡,因為你經歷了基督的改變的時候,你可以告訴別人,你不需要在台前彈琴,就像瑪利亞一樣,他可以告訴別人,基督已經改變我的生命,一份感恩的心,才是真正第一個服侍神的資格..
我參加過一個叫OMF的警察會,以前是前身中國內地會,中國內地會的創辦人是達德生先生,1865年.當時由英國坐船去中國宣教,不是現在這麼簡單,是不容易的.因為在1865年坐船由英國去中國,大家知不知道要多久的時間,是要五個月的時間..
所以當時假裝是短孫,最好多五個月,現在只能坐船五個月.所以有人申請是很不簡單的,所以大家知道如果有人想去中國,已經是很開心..
當時有一個年輕人,現在中文名叫曹亞直George Stott,很可惜他申請了三次去中國,他不成功,別人不讓他去,他最後問,你給我多一次機會,我求達德生讓我去中國..
他說OK,我讓你去見達德生先生.當他和達德生先生見面的時候,他走進來,眾人都跟他說,你還是留在英國為中國祈禱,你不可以去.達德生先生也是這樣說,弟兄你還是留在英國..
為什麼呢?原來當他走進來的時候,人們看到他拿著一隻拐杖,他只有一隻腳,因為他左腳在年輕的時候發炎,醫生要把左腳切割,就是說他走進來只有一隻腳,人們說你去中國宣教已經不容易,你一隻腳根本不可以在中國生存..
當時曹亞直和達德生先生說了一句話,他說英國有兩隻腳的人很多,但他們不是在中國,你給我一隻腳去中國吧.當他說了這句話的時候,達德生先生很感動,OK,他說我讓你去..
這裡是當時溫州的地方,大家知道溫州是沿岸的地方,中國南部.當時1865年的溫州是沒有教會,很少基督徒,今天的溫州是700萬人口,10\%是基督徒,而且是全中國其中所有大城市最多基督徒的一個城市..
我相信曹亞直是一個從人看來沒有資格去中國的人,他沒有想過在1865年之後,神怎樣興起溫州的教會.我想他從來沒有想過,一個從人看來沒有資格的人,為什麼他要去呢?一份感恩的心..
今天當我們在寒冷堂慶祝十三週年的時候,老道上神提醒我們,你不要說我沒有資格,我沒有才幹,我又不是像他那麼厲害,不是的,神說你經歷了基督的改變,你是有資格,你有一份感恩的心去服侍神,能夠跟隨基督的腳蹤..
我們今天看第二個人物,約阿娜,Joanna,約阿娜和瑪麗亞有什麼分別呢?其實差不多是天淵之別,為什麼我會這樣說呢?瑪麗亞是活在痛苦黑暗裡,被鬼折磨,但是約阿娜她說她的丈夫是撫薩,撫薩是跟誰打工呢?是跟希律打工的..
為什麼這麼重要呢?希律是當時加利利一帶地方的市長,香港沒有市長,香港只有特首,比如說多倫多市長,或者紐約市長,如果你跟市長打工,你的生活應該是不錯的,不是說一定富裕,但是你應該是無憂無慮,你住在一個不錯的地方..
但是今天神提醒我說約阿娜願意離開她的安樂窩,願意跟隨基督,周遊各城各鄉,當然她去了多久我們不知道,但是她願意離開她的安樂窩,如果我說第一個是服侍神的資格是一個感恩的心,第二個是一個願意的心,不是你厲害不厲害,不是你有沒有財運,而是你願不願意..
當我讀醫學院的時候,我其實入醫學院第一年信耶穌,直至現在,但是我記得我讀醫學院第三年的時候,神給了我一句聖經的金句,直至今天,都是對我來說很重要,是在箴言的23章26節,四個字而已,很簡單,直至今天我還記得,箴言23章26節,將心歸我..
將心歸我是什麼意思?神說我不是要你醫學院的財國,不是要你醫學院的學位,那個都是次要,不是要你有多厲害,而是說將心歸我.直至今天神提醒我說,今天我們事奉神除了感恩的心之外,是一個願意的心,是你願不願意很重要,不是說神說我不行,不是,是神說你肯不肯最重要..
我剛才說過我參加的參會叫OMF前身是中國內地會,當時申請去中國的人,我就說坐船很不容易,要四個月五個月的時間,1876年有一位牛津大學的畢業生,醫學院,一位很優秀的醫生在牛津大學畢業,當時大學很喜歡他,就說邀請他來這裡牛津大學做研究,.
你將來一定成為一個教授,因為他很優秀,去醫院不是神要我去中國事奉,所以他就申請中國內地會,想去中國.很奇妙的在記錄裡面,他申請又不成功,我說中國內地會真麻煩,你一隻腳不能去,我明白,但是這個是牛津大學優秀畢業生不讓他去,為什麼呢?.

$^{41}$我在當時英國大學的文獻府找到一個記錄,其實這個文獻府的記錄和耶路大學的記錄一樣,我找到兩個記錄,為什麼不讓他去?一個原因,因為他單身未結婚,為什麼不讓他去呢?.
因為當時在面試委員會裡面,覺得當時的英國弟兄很不懂得照顧自己,他說你生存能力不是那麼高,你去中國死定了,所以他說你回去祈禱,神給你太太,你再回來第二次面試,我不知道現在的警察會有沒有這樣的要求,.
但是很奇妙,當你翻查記錄的時候,四個月之後真的有第二次面試,而第二次面試的記錄裡面,第一句說話很精彩,他說Scopio醫生現在已婚,真的神在四個月之後,給他太太去,皆大歡喜,派他去山西太原的地方行醫,因為他很優秀的醫生,學中文傳福音,做醫生,醫療專教,.
1880年去到中國,1883年8月他照顧一個病人,當然現在回頭看我們知道發生什麼事,因為這個病人有感染病,當他看完這個病人之後,一個星期之後,這個Scopio醫生過世,享年只有33歲,我看這個記錄的時候,我心裡面很大聲說,神啊,為什麼你不讓他活長一點呢?.
更好的事奉,更多人信耶穌,為什麼你只讓他活到33歲這麼短的時間,還有一個這麼優秀的醫生,是不是很浪費呢?我心裡面也有掙扎,但是我想起耶穌那一句話,他說,一粒麥子不落在地裡死了,仍然是一粒麥子,但是落在地裡就結出許多的子粒來,當我們翻查記錄的時候,看到他的生命影響了當時很多英國年輕人的生命,.
奉獻給中國,然後之後,一浪子一浪的很多人去中國事奉,因為Scopio醫生的生命的影響,大家知道墓碑,墓碑通常是一個名字,有兩個日期的,一個是什麼?出生日期,一個死亡日期,中間是有什麼在這裡的?.
一條線,連接兩個日期,你有沒有留意過墓碑上面的線,長命的人,那條線是不是長一點呢?短命的人,那條線是不是短一點呢?如果你不相信,你自己去墓碑看看吧.
那條線長短,其實沒什麼分別的,長命短什麼,但是神提醒我說,你投資什麼在這條線上面才是最重要的,那條生命線上面,你放什麼在上面?.
今天約阿娜,一份願意的心,Scopio醫生,一份願意的心,他只是可能33歲,但是因為他願意,神用他,好,我們去到第三個人物,第三個人物是蘇薩娜,蘇薩娜我就覺得很困難,為什麼呢?.
你看聖經裡面,中文的聖經說,並蘇薩娜,原文是和,蘇薩娜,我對盧家醫生有些投訴,盧家醫生你只說蘇薩娜這個名字,他做了什麼呢?他有什麼好呢?他因此在哪裡來的呢?.
他是不是款待人很厲害呢?他究竟為耶穌做了什麼,為什麼你什麼都沒告訴我呢?只是給了個名字蘇薩娜,我想了很久,我說神啊,這個人名對我有什麼啟迪呢?.
我想了很久,我說神啊,原來你教了一個很重要的功課,很多時候當我們事奉神的時候,我們很希望別人記得什麼,記得我做了什麼,重心是我,不是基督,我做了什麼,你在教會事奉,也希望牧師記得我做了什麼,不是基督做了什麼,如果牧師多謝各位,提不同的人名,漏了提你的名,你心裡也不開心,是不是?.
為什麼牧師你沒提我名,我有幫忙的,但是今天在這裡提醒我們說,蘇薩娜這個人名裡面,正是因為他沒有提過他做過任何事情,對我有一個很大的啟迪,原來神提醒我,重點不是別人記得我做過什麼,是基督在我身上做了什麼,如果我說一份感恩的心,一份願意的心,第三個是一份謙卑的心,.
蘇薩娜正是因為他什麼都沒有提過他做過什麼,神提醒我一份謙卑的心,當時我提過內地會有很多宣教士,包括在Scofield醫生去世之後,很多年輕人奉獻,包括當時我們說的劍橋七傑,不是牛津,這次是七個從劍橋年輕有為的奉獻,.
去中國,坐同一條船,一起185年去中國事奉,有位人物叫做何詩德,何詩德他有很特別的性格,來自一個非常有錢富有家庭,貴族家庭,爸爸阿爺是很有影響力地位的人物,來自一個這樣的背景,.
他奉獻的時候去到中國什麼地方呢?去到中國最窮的山西洪洞的地方,你可以看到很不容易,一個很有錢富有貴族家庭,來自一個中國最窮的地方,他做什麼呢?他當時幫助內地會差不多幫助當時吸鴉片的人去戒毒,香港有戒毒中心,陳牧師也很熟悉戒毒的工作,當時他就做戒毒工作,.
當時大家知道英國賣了很多鴉片給中國人,而中國因為這樣很多人受害,所以中國內地會決定要參與戒毒工作,當時他參與戒毒工作的時候,何詩德要跟隨一位中國牧師叫做直勝摩,直勝摩當時辦一個叫天招局,幫助吸鴉片的人戒毒,.
何詩德寫信給他媽媽說,當時我跟這班人一起生活戒毒,他很不容易,為什麼呢?他說我跟150個以前吸毒現在戒毒的人一起生活,我沒有自己的空間,沒有自己的房間,150個人同一個房間一起生活,起居飲食,日夜對住,.
他說我來中國之前,我以為我將所有東西奉獻放在十字架面前,但今天神說你還有一件事未做,還要放在十字架面前,是什麼呢?他說你將個人的空間放在十字架面前,個人的空間,.
神挑戰他,他就這樣奉獻,跟這班人一起生活,何詩德來中國的時候25歲,何詩德離開中國的時候85歲,你不可以想像,其實他一生人就在中國裡面,25到85歲,60年在中國服侍,其實他一回英國,85歲,未嫁一年就去世了,.
當他去世的時候眾人很懷念他,不約而同,鄧炳芝妹說不如我們寫些東西紀念他,為中國人做過些什麼呢?不約而同,人們寫不同的東西,當鄧炳芝妹寫的時候,有一句話不謀而合,同一句話出現,是什麼出現呢?.
何詩德一生人的目的,就是要叫別人忘記他,以致別人記得主耶穌基督,一份謙卑的心,我自己心裡很感動,一個60年在中國服侍的人,他一生人的目的不是叫別人記得他做過什麼,而是叫別人記得耶穌基督做過什麼,一份謙卑的心,.
所以今天看這三個人物,瑪利亞一份感恩的心,約瓦那一份願意的心,蘇珊娜一份謙卑的心,也從三個中國內地人物裡面去看,今天神給我們挑戰,.
我們結束的時候,我仍然引用了《箴言》23章26節,四句話來結束,神給我們今天的挑戰,讓我們面對也感恩十三週年唐慶的時候,神說的四個字,將心歸我,是神對我們的勉勵,將心歸我,讓我們的心能夠歸向神,.
主耶穌基督,讓神祝福,洪恩唐,多謝大家..
多謝大家..
(拍攝).
弟兄姊妹平安,現在可以播放那張照片了..
1870年,在英國一個小地方,有一個十七歲的年輕人,他俯伏在神面前,他信了主只有幾個月,他在神面前表達他心裡的感恩,他一而再地向神說,求你給我一點點的工作,.
無論是多艱難,或者是多渺小,只是讓我去表達我對你的感恩和愛.四年之後,他踏足中國,大得生二十一歲,去到中國上海,他懷著一份感恩願意的心去到那裡,他要將福音傳遍中國每一個角落,但是當時沒有人信他可以做得到..
但是五十四年之後,他在湖南長沙離開世界,回天閣,而湖南是他所創立的中國內地會所進入的最後一個省份,神用一個平凡的年輕人的生命去完成了從人眼中是不可能的事..
可不可以幫我前進到下一張雲燈片?.
今天同樣有很多年輕人願意超越自己,用他的生命去回應神的愛,去將福音,將神的作為傳遍天下萬民..
我們很想他們都知道,神願意用平凡人成就非凡事,以致他們可以和大得生一樣,全然的為生信靠,去將榮耀歸給神..
我們很盼望,對不起,下一張slide倒轉了,我們很盼望用一套電影,不是說大得生怎樣好,而是說神怎樣去用一個生命去影響更多的人..
我們很想用一個愛情故事去描繪神和大得生瑪利亞,和他們和中國人之間的愛,去講述一個神的作為..
我們盼望這套電影可以向不同的人都有一個訊息,我們覺得一套電影你看完無論多感動,很快就會過去,所以我們想有一個跟進的行動,讓那些人一邊看一邊可以去問,今天神要我做什麼?.
我們盼望未信主的可以去尋求,去探索信仰,可以找到生命的主,信了主的可以去進一步進深成長,成為門徒,去尋找合適他們的事奉機會,去找到裝備事奉的途徑..
我們盼望以一個電影去成為多人的祝福,拍電影不便宜,按照現在的劇本我們需要七百萬美金,七百萬美金不是一個小數目,但神已經供應了三百六十萬,.
我們盼望在2024年剩下的兩個月,我們可以籌足七百萬去完成這個電影的製作,我們感恩有電影姐妹奉獻一百萬,五十萬,但是我們更加感恩是一些小額的奉獻,我自己最深感動的是兩年前認識一對韓國的夫婦,.

$^{81}$他們跟我說他們很想參與這個項目,但是他們沒有固定收入,他們說我們有工作的時候就會奉獻,之後我每隔幾個星期,有時候幾個月,我就收到他們奉獻的十元,二十元,他們衷心的為這個項目去拜上,而我深信他們每一次奉獻的時候,他們都會為這個項目去祈禱..
而代禱就是這個項目最最最需要的一件事,也是我相信我邀請你和我們一起去做的一件事,在這個二維碼你可以收到最新的資訊,我們也希望可以留下一些代禱卡,讓你拿回去,繼續為我們去禱告,也分享給其他人,更多人可以一起去做..
可以一起參與,一起蒙受神的福氣.多謝大家..
\newpage



\section{歷代志下 1:7-12}
\label{sec:lQqzMqR_pN8}
\textbf{2024.11.10 《蒙神喜悅的禱告》 歷代志下1\_7-12 (和修版) 陳一華牧師}
\newline
\newline
連結: \href{https://youtube.com/watch?v=lQqzMqR-pN8}{\texttt{https://youtube.com/watch?v=lQqzMqR-pN8}} ~~~~ 語音日期: 2024-11-11
\newline
\newline
\hyperref[sec:xtW6V_2_PsE]{\small{< < < PREV SERMON < < <}}
~
\hyperref[sec:index_chronic]{\small{[返順時目]}}
~
\hyperref[sec:index_scriptual]{\small{[返順卷目]}}
~
\hyperref[sec:Y3F7I4ibIlg]{\small{> > > NEXT SERMON > > >}}
\newline
\newline
歷代志下 1:7-12
\newline
\begin{longtable}{cl}
\hline
\hline
章節 & 經文 (和合本修訂版)\\
\hline
1:7 & \begin{tabularx}{0.7\textwidth}{X} 當夜，神向所羅門顯現，對他說：「你願我賜你甚麼，你可以求。」 \end{tabularx} \\ \\ \relax
1:8 & \begin{tabularx}{0.7\textwidth}{X} 所羅門對神說：「你曾向我父親大衛大施慈愛，使我接續他作王。 \end{tabularx} \\ \\ \relax
1:9 & \begin{tabularx}{0.7\textwidth}{X} 耶和華神啊，現在求你實現向我父親大衛所應許的話；因你立我作這百姓的王，他們如同地上的塵沙那樣多。 \end{tabularx} \\ \\ \relax
1:10 & \begin{tabularx}{0.7\textwidth}{X} 現在，求你賜我智慧聰明，好在這百姓面前出入；不然，誰能判斷你這麼多的百姓呢？」 \end{tabularx} \\ \\ \relax
1:11 & \begin{tabularx}{0.7\textwidth}{X} 神對所羅門說：「你有這心意，不求資財、豐富、尊榮，也不求滅絕恨你之人的性命，又不求長壽；我既立你作我百姓的王，你只求智慧聰明，好審判我的百姓， \end{tabularx} \\ \\ \relax
1:12 & \begin{tabularx}{0.7\textwidth}{X} 我必賜你智慧聰明，也必賜你資財、豐富、尊榮，在你以前的列王未曾有過，在你以後也不會再有。」 \end{tabularx} \\ \\
[1ex]
\hline
\hline
\end{longtable}
$^{1}$上星期是我們的堂慶,有多少位有出席堂慶?.
今天有朕禮,我剛才問曾牧師有多少位洗禮?他說有四位洗禮.
你說天父給我們洪恩家是否有很多恩典?.
今早牧師和大家分享的題目已是耳熟能詳,和祈禱有關.
牧師為什麼要分享這個題目?.
這個題目由我未信耶穌到信於耶穌經常做.
我未信耶穌之前,朋友勸我祈禱.
那時我不懂祈禱,只閉上眼,人們教我祈禱.
到我信了耶穌後,我就經常祈禱.
所以這個題目對我們來說,一點也不陌生.
頂智妹妹,我們很多時候都有祈禱.
但有時我們要問,我們這麼多次祈禱.
到底有多少次我們的禱告是蒙神喜悅?.
這就不同了層次,如果我們的祈禱是蒙神喜悅.
我們的禱告就會與眾不同.
今天我們學習一個人物,所羅門.
剛才大家也讀過那段經文.
學習所羅門如何向上帝禱告.
以至上帝喜悅他的祈禱.
話說當所羅門造王登基不久.
有一次他帶著文武百官到一個地方祈禱.
那個地方的名字比較陌生,叫做基片.
為什麼會帶這麼多人去獻祭呢?.
因為大家都知道那時候聖殿還沒有建好.
以色列人如何向上帝獻祭呢?.
原來在基片的地方記載了一個很大的休壇.
所以他們就去那裡向上帝獻祭,向上帝祈求.
我相信因為所羅門做得很好.
那一晚就特別了.
聖經記載當晚,即是他去獻祭的那一天.
上帝向所羅門顯現.
上帝對所羅門說.
無論你向我求什麼,我都會給你.
看見沒有?.
很棒吧?.
是不是有點像阿拉丁神燈一樣.
擦兩擦,要什麼就有什麼.
是不是很棒?.
但你留意,所羅門的回應.
他向上帝的禱文是很特別的.

$^{41}$以至他的祈禱是神所喜悅的.
有四個原因在當中.
我今天顧著說都忘了給大家看泡泡.
證明我不是常用泡泡的人.
所羅門的禱文有四個重點.
我們來看看.
第一個重點,原來他的祈禱.
是建基於神的應許上.
他這樣向上帝禱告.
神啊,今天你立我為以色列人的王.
是為什麼?.
是你要實現和我父親大衛所應許的話.
有沒有印象?.
原來他的禱文是說.
今天我做王不是要我威風百面.
不是要我高高在上.
也不是要我君臨天下.
是要我什麼?.
是要實現你和我父親大衛所應許的話.
換句話說,頂智穆門.
所羅門的禱文是建基於神的應許上.
什麼是神的應許?.
神的應許就是神的話語.
就是神的道,就是聖經.
所以我們現在明白.
為什麼我們每次回到教會崇拜.
會讀很多聖經.
甚至有時候教會叫我們背讀經文.
還讀經文.
不單讀還背.
原來我們每一次朗讀經文.
或者我們每一次背誦經文.
或者我們每一次自己靈修看聖經.
或者我們參加查經班去研讀經文.
各位頂智穆門.
原來神的話語不知不覺之間.
融化在我們心中.
所以我們的禱告若果能夠建基於神的話語.
神的應許上就很不同了.
譬如舉個例子來說.

$^{81}$有一段時間可能自己發覺.
好像近來做什麼都不順利.
你明白我的意思嗎?.
每件事都好像碰著黑一樣.
人都在這裡有少許十五,十六.
到底上帝是否仍然聽我的祈禱呢?.
人都很掙扎.
是否應該去祈禱呢?.
有時候我們都會陷入這些低谷.
頂智穆都試過了.
但當你在掙扎的時候.
忽然間你想起一節經文.
那節經文是耶穌說.
你們祈求就給你們.
尋找就尋見.
溝門就是為我們開門.
當你想起上帝有這樣的應許的時候.
你就更加放心.
向上帝祈禱.
為什麼呢?.
因為你建基於神的話語上.
或者弟兄姊妹當你自己信了耶穌.
都很愛主很熱心.
我見有不少弟兄姊妹.
都會在自己的職場.
或者在自己的公司.
或者自己工作的地方.
去開一些小組.
我以前在警察團體.
在不同的警署.
都有開一些午間聚會.
利用食午飯的時間.
來開聚會.
有時候你會發覺.
上帝啊.
我開了這個小組之後.
好像人數都沒有增長.
你明白我的意思吧?.
來來去去都是只有兩三個人.
出席這個小組聚會.

$^{121}$有時候心裡有一點什麼呢?.
有一點像灰心一樣.
為什麼人數不多呢?.
不過當你有這樣的感覺的時候.
忽然間你又想起.
耶穌說的那句話.
耶穌說在你們中間.
只要有兩三個人.
奉我的名來聚會的時候.
我就在你們中間了.
你想起來.
雖然我的小組.
只有三四隻小貓.
但原來有主耶穌的應許.
主會與我們同在.
你的信心就大大的振奮了.
原來神與我們在一起.
所以各位觀眾.
記得原來當我們的禱告.
能夠建基在神的話語.
神的應許上的時候.
你這樣去禱告.
一定是很穩陣.
不會出錯的.
為什麼呢?.
因為你將神的話語.
成為你禱告的土文.
是不是很寶貴呢?.
這是第一個原因.
為什麼所羅門的禱告.
是神所喜悅的.
我現在去看第二個原因.
第二個原因是因為.
所羅門的禱告.
他向神禱告.
祈求神的旨意.
來成全.
就是說.
神今天你納我為以色列人的王.
是要我來建立大衛的國.

$^{161}$鞏固大衛的家.
來完成上帝你救贖的計劃.
換句話來說.
他清楚知道.
這次上帝納他為王.
是在神的計劃當中.
原來神有祂的旨意.
因為上帝興起大衛家.
要讓大衛家成為上帝救恩的渠道.
有時候你讀經文.
讀到舊約.
有兩隻字比較陌生.
但是很有意思的.
不知道你有沒有印象.
讀到上帝說.
我要興起大衛家的苗裔.
苗裔這兩個字.
大家有沒有印象?.
沒有什麼反應.
沒有什麼反應就是沒有印象.
苗裔這兩個字是什麼意思呢?.
是特定指到.
在大衛的家譜裡.
上帝要興起一位與眾不同.
就是世間做救贖主的耶穌基督.
所以他是大衛家的苗裔.
他出現的時候.
要成為上帝的彌賽亞.
拯救地上所有的人.
各位觀眾.
現在你明白.
為什麼有時候讀《馬太福音》.
讀《盧家福音》.
一讀第一章就開始頭痛.
為什麼?.
為什麼要這麼多名字?.
誰生誰?.
你以為好像很多餘.
實證有沒有意思呢?.
有的.

$^{201}$為什麼?.
因為要記載耶穌基督的家譜.
原來就是源自於大衛家.
一直延續下來.
所以現在大家明白.
所羅蒙他知道上帝要他做王.
不是要他高高在上.
而是什麼?.
要完成上帝救贖的計劃.
要讓神的旨意去成就.
祂不過是工具.
所以祂也知道我是王.
但是我都要謙卑.
因為我不過是神要用的器皿.
直至什麼?.
神的救恩.
神的救贖.
能夠延伸下去.
所以各位聽眾.
如果你明白這個意思.
你看下面這幅圖.
就更加有意思了.
這幅圖.
耶穌在哪裡祈禱?.
在喀西瑪利園.
大家都很熟悉.
在喀西瑪利園祈禱的耶穌.
那一刻是非常沉重.
也非常痛苦.
也有很大的壓力.
祂在神面前禱告的時候.
是非常掙扎.
因為祂知道.
祂面前十字架的凝聚.
是一個很殘酷的凝聚.
是一個很痛苦的凝聚.
到底我要不要上十字架的凝聚呢?.
各位頂尖的朋友.
那時候的耶穌.
如果身邊有多些人支持祂就好了.

$^{241}$但很可惜.
事與願違.
身邊的門徒.
每個人都睡著了.
只有耶穌獨自面對一個這麼沉痛的時刻.
各位頂尖的朋友.
雖然祂的祈禱很掙扎.
聖經記載.
祂的汗點滴在地上.
好像血一樣.
換句話來說.
這個十字架的凝聚.
我要不要去上呢?.
所以祂向上帝禱告.
祂說:神啊!如果可以的話.
求你將這個苦杯移走.
然而不要照我的意思.
而是照你的意思去行.
所以各位頂尖的朋友.
耶穌的禱文是神所喜悅的.
因為祂的禱告是要讓神的旨意去成就.
如果你明白耶穌的禱告的時候.
你會發覺所羅門的禱文和耶穌的禱告.
有種前後呼應的作用.
因為老早之前所羅門的禱告已經展現出.
神啊!願你的旨意成就.
現在耶穌也說:神啊!讓你的旨意去成就.
當我們這樣去禱告的時候.
頂尖的朋友.
你會發覺我們的禱告是什麼?.
是神所喜悅的.
所以在這裡也要提醒頂尖的朋友.
有時候我們犯了一個毛病.
我們好像自己說了算.
有時候會否像我這樣向神禱告.
神就要聽我的禱告.
神就要給我禱告.
你不聽不給我就不理你.
有時候會否有這樣的情況.
各位頂尖的朋友.

$^{281}$如果我們的態度是以自我為中心.
現在我們要將它更改.
要什麼?要讓神的旨意成就.
因為神的計劃是最好的.
有時候我們不明白為什麼.
神啊!我身上這個病已經兩三年了.
但好像還沒有徹底地斷尾.
還是有這個病.
為什麼呢?.
有時候也會有心裡很不舒服的感覺.
各位頂尖的朋友.
無論在任何的情況下.
記得有時候上帝做每一件事.
都有祂很奇妙的安排.
奇妙的計劃.
祂的旨意在我們身上.
所以我們的禱告是什麼?.
是要讓神的旨意成就.
我們再看第三個原因.
為什麼所羅門的禱文是神所喜悅的呢?.
因為祂的祈禱內容是不畏己求.
或者是小畏己求.
因為祂向神禱告.
祂求什麼?.
神啊!求你讓我有聰明智慧.
為什麼祂求上帝讓我有聰明智慧呢?.
一來祂是初出茅廬.
剛剛做王.
沒有什麼經驗.
另一方面最主要是.
神啊!你派我管理的這群百姓人數.
多到好像地上的塵土.
剛才不是這樣記載嗎?.
弟兄姊妹們我們能不能數得出來?.
是多不勝數的.
上帝要我管理這麼多人.
我怎麼能夠管理得到呢?.
弟兄姊妹們.
我們不要說管理這麼多人.
如果你在家裡有兩隻猴子.

$^{321}$你也很難管理.
會不會呢?.
這兩隻猴子專門跟你作對.
你叫牠去東牠就去西.
你叫牠坐下牠就跳高跳低.
每樣都跟你作對.
有沒有這樣的小朋友?.
管理得到嗎?.
頭都痛了.
何況現在要管理這麼多人.
我怎麼能夠令到我的國家.
能夠國泰民安呢?.
我怎麼能夠令到我國家的百姓.
都敬畏上帝呢?.
弟兄姊妹們.
我怎麼能夠在眾多百姓當中.
施行公義的審判呢?.
容易嗎?.
一點都不容易.
我怎麼能夠在百姓中自由出入.
聽受他們的訴訟呢?.
弟兄姊妹們.
容易嗎?.
簡直是一個很艱鉅很大的任務.
所以即使如此.
所羅門向上帝禱告.
他求上帝給他有聰明智慧.
也不是為自己的私利去祈禱.
他的禱告是.
我祈求聰明智慧.
其實我的目的.
都是要完成上帝給我的託付.
完成上帝給我的任務.
做得好.
管理得好.
弟兄姊妹們.
如果這樣看法.
就很不一樣了.
原來他的禱告.
都不是為自己的利益.

$^{361}$而是為整個百姓.
整個國家.
能夠合神的心意.
完成這個任務.
來去禱告.
所以這裡又給我們一個好好學習.
各位弟兄姊妹們.
你們有沒有覺得.
我們就算信了耶穌一段時間.
可能我們的祈禱.
來來去去.
都是這樣的.
神啊.
求你讓我發達吧.
神啊 求你讓我小屋換大屋吧.
最好抽到一間大屋吧.
神啊 求你讓我舊車換新車吧.
神啊 求你讓我一年去三次旅行吧.
去不到也去兩次吧.
我們有很多時候.
來來去去都什麼.
好像祈這些的.
你們有沒有留意到.
弟兄姊妹們.
如果我們的禱告.
來來去去.
是為個人利益去禱告的時候.
你說上帝是不是也很失望呢.
是不是.
為什麼.
你信了主這麼久.
為什麼我們的禱文.
來來去去都是什麼.
環繞著自己的利益呢.
忘記了什麼.
忘記了去求上帝.
讓我有能力.
去完成你給我的任務.
是不是.
完成你給我的託付.

$^{401}$因為你在我身上.
神你有你的計劃.
是不是.
所以我希望能夠達成.
你給我的任務.
弟兄姊妹們.
當我們這樣禱告的時候.
就是一個蒙神喜悅的禱告.
是不是.
好 我們再去看最後.
最後.
所羅門的禱告.
是一個很單純的祈求.
為什麼呢.
因為他整個禱文裡.
只是單一求一件事情.
就是求神給他有聰明智慧.
你會發覺他的禱告.
一點都不複雜.
他有沒有求上帝給他榮華富貴.
有沒有.
有沒有求上帝給他長命百歲.
有沒有.
沒有.
有沒有求上帝給他戰勝敵人.
打贏仗.
有沒有.
沒有.
他沒有求這些.
他實正是可以求的.
你明白我的意思嗎.
因為神都應許你求什麼.
我都給你.
但他祈得好.
他沒有求那些.
他只是求什麼.
求一件事.
就是上帝給我有聰明智慧.
一個很單純赤子之心的禱告.
難怪神喜不喜悅.

$^{441}$神喜悅.
然後神就這樣說.
神說.
你沒有向我求的.
我都什麼.
我都給你.
很棒吧.
你沒有求的.
我都什麼.
給你.
而且我要給你的祝福.
可以用四個字去形容.
哪四個字呢.
就是空前絕後.
大家認不認同.
上帝對所羅門說.
我給你的福氣.
舊陣子的皇帝.
有沒有品嘗過.
沒有.
將來出現的皇帝.
可不可以得到.
不可以.
唯獨是你.
能夠擁有我給你的福氣.
各位兄弟姐妹們.
你和我羨不羨慕.
都很想做所羅門2.0.
如果上帝這樣祝福我們就好了.
空前絕後.
厲害吧.
很加勢.
各位兄弟姐妹們.
雖然你和我.
不是所羅門.
雖然想做2.0.
但是你和我都有這樣的福氣.
為什麼呢.
在我結束我的講道的時候.
我就要送一節經文給大家.

$^{481}$也是主耶穌的應許.
只要我們的禱告.
在夢神喜悅的時候.
我們同樣都得到神的恩典和祝福.
是一節什麼金句呢.
我們一起來讀這節金句.
你們要先求他的國和他的兒.
這些東西都要加給你們了.
這就是在新約.
耶穌親口的應許.
你們要先求什麼.
他的國.
什麼是神的國呢.
神的國就是神的家.
神的家就是教會.
教會就是神在地上計劃的一個表徵.
所以弟子們.
當我們今天很有福氣.
信了耶穌.
又能夠在神的家.
在神的教會裡.
一起來敬拜.
一起生活的時候.
這是很大的祝福.
但你們有沒有想過.
當上帝的家.
當神的教會.
很需要多些人起來事奉的時候.
你們有沒有挺身而出呢.
各位弟子們.
很感恩我知道.
洪恩堂是很夢神的恩典.
我們有很多服侍的機會.
弟子們.
當教會缺乏人手的時候.
你們有沒有主動自動去承認呢.
或者當洪恩家.
需要在經濟上支援的時候.
我們有沒有很樂意的去奉獻.
支持神的聖功呢.

$^{521}$又或者.
當我們發覺到.
原來神的家.
是很需要.
有一個傳福音的行動的時候.
有報道聚會的時候.
我們有沒有很積極的.
去帶領人到神的面前呢.
各位弟子們.
如果你和我經常記掛著.
神的家.
神的國的時候.
我們的禱告就很容易想起.
教會的需要.
弟兄姊妹的需要.
為他們禱告.
牧者的需要.
為牧者禱告.
社區的需要.
我們紀念社區的需要.
弟兄姊妹們.
我們就先求什麼呢.
神的國.
什麼是神的義呢.
神的義就是神所喜悅的事情.
神所喜悅的事情.
莫過於我們去遵行祂的大使命.
就是將福音廣傳開去.
正如耶穌在《馬太福音》28章.
提醒我們.
要往普天下去傳福音.
所以弟兄姊妹們.
當我們積極去傳福音的時候.
你就會發覺.
天國的國度就會不斷地擴展開去.
在疫情過後.
復常的這兩年.
我也觀察到.
香港的教會有一個現象.
就是不少教會.

$^{561}$就再次起來.
很主動積極地去傳福音.
開很多報道會.
福音餐會.
個人報道.
甚至最近有一個.
不算新興.
但也有一個很難得的現象.
就是做街頭報道.
街報的工作.
弟兄姊妹你們逛街坐地鐵.
有沒有見到有些人.
在地鐵站附近.
派單張.
唱詩歌.
講訊息.
有沒有遇過.
這個就是什麼.
就是街頭報道.
使用一位很年長的牧師.
柳振平牧師.
他八十多歲.
但身先士卒.
老當益壯.
天天都街報.
厲不厲害.
他天天都街報.
我們年輕人不敢躲起來.
都要和柳牧師一起出去街報.
現在香港十八區裡.
差不多每一區都有街報.
每一區的教會都有響應.
都有弟兄姊妹.
有牧者.
一起在街上報道.
你說這是不是很大的祝福.
這是很大的恩典.
整個中國只有兩個地方.
可以有自由報道.
有街報的.

$^{601}$一個就是香港.
一個就是澳門.
大家明白嗎.
我們應該要把握很多很好的機會.
廣傳福音.
神給我們有智慧.
將福音傳開.
以致我們會為這些事祈禱.
願意先求神的國和神的義.
聖經說.
這些東西都要什麼.
家給你們要.
我們同心低頭祈禱.
親愛的主.
今天我們見到你祝福所羅門.
是因為祂的禱告是你所喜悅的.
但願我們今天都學校所羅門的禱告.
能夠在我們的禱告當中.
你們去追求神你所喜悅的事情.
以致我們的禱文是蒙神閱立.
也是蒙神賜福.
求你祝福眾弟兄姊妹.
弘恩家我們繼續努力.
先求你的國和你的義.
好讓我們領受上帝豐盛的恩惠.
祈禱奉耶穌名求.
阿們.
\newpage



\section{創世記 16:7-14}
\label{sec:Y3F7I4ibIlg}
\textbf{2024.11.17 《神與你同行》 創世記16\_7-14 (和修版) 講員 : 許淑芬牧師}
\newline
\newline
連結: \href{https://youtube.com/watch?v=Y3F7I4ibIlg}{\texttt{https://youtube.com/watch?v=Y3F7I4ibIlg}} ~~~~ 語音日期: 2024-11-17
\newline
\newline
\hyperref[sec:lQqzMqR_pN8]{\small{< < < PREV SERMON < < <}}
~
\hyperref[sec:index_chronic]{\small{[返順時目]}}
~
\hyperref[sec:index_scriptual]{\small{[返順卷目]}}
~
\hyperref[sec:lRt1SGTCNSE]{\small{> > > NEXT SERMON > > >}}
\newline
\newline
創世記 16:7-14
\newline
\begin{longtable}{cl}
\hline
\hline
章節 & 經文 (和合本修訂版)\\
\hline
16:7 & \begin{tabularx}{0.7\textwidth}{X} 耶和華的使者在曠野的水泉旁，在書珥路上的水泉旁遇見夏甲， \end{tabularx} \\ \\ \relax
16:8 & \begin{tabularx}{0.7\textwidth}{X} 對她說：「撒萊的婢女夏甲，你從哪裡來？要到哪裡去？」她說：「我從我的女主人撒萊面前逃出來。」 \end{tabularx} \\ \\ \relax
16:9 & \begin{tabularx}{0.7\textwidth}{X} 耶和華的使者對她說：「你要回到你的女主人那裡，屈服在她手下。」 \end{tabularx} \\ \\ \relax
16:10 & \begin{tabularx}{0.7\textwidth}{X} 耶和華的使者對她說：「我必使你的後裔極其繁多，多到不可勝數。」 \end{tabularx} \\ \\ \relax
16:11 & \begin{tabularx}{0.7\textwidth}{X} 耶和華的使者又對她說：「看哪，你已懷孕，要生一個兒子。你要給他起名叫以實瑪利，因為耶和華聽見了你的苦楚。 \end{tabularx} \\ \\ \relax
16:12 & \begin{tabularx}{0.7\textwidth}{X} 他為人必像野驢。他的手要攻打人，人的手也要攻打他；他必常與他的眾弟兄作對。」 \end{tabularx} \\ \\ \relax
16:13 & \begin{tabularx}{0.7\textwidth}{X} 夏甲就稱那向她說話的耶和華為「你是看見的神」，因為她說：「他看見了我之後，我還能在這裡看見他嗎？」 \end{tabularx} \\ \\ \relax
16:14 & \begin{tabularx}{0.7\textwidth}{X} 所以這井名叫庇耳‧拉海‧萊，看哪，它位於加低斯和巴列的中間。 \end{tabularx} \\ \\
[1ex]
\hline
\hline
\end{longtable}
$^{1}$頂智慧平安.
我感恩有機會再次跟大家分享上帝的話語.
轉眼間 原來已經過了幾個月.
幾個月前 六月的時候.
跟大家分享一個主題.
用亞伯拉罕作為核心.
我們與神同行.
今天我想用另一個角度去看.
我們要看神與我們同行.
神很樂意與我們每一個人同行.
只是今天我們願不願意與神同行呢?.
曾經有一個姐妹分享她的見證.
讓我看到神在她生命中的經歷.
如何與神同行的經歷.
話說這個姐妹很年輕.
她從小到大都很幻想.
有一天她很想結婚.
她也很心急結婚.
她二十歲時已經找到一個對象.
她想著嫁了就對了.
於是就結婚了.
但結婚後不久就生了一個小朋友.
小朋友還未到一歲的時候.
老公跟她說.
其實我還未準備好結婚.
那怎麼辦?.
要抱著小朋友.
誰知還要加上一首歌.
「一走了之」.
「一走了之」還有一句俗語.
就是一屁股債都給了她.
意思是她要抱著小朋友.
又要爭人家的周身債.
在她很絕望的時候.
她抱著小朋友.
本來想著死了之.
她分享她記得勝經裡有一個人物.
叫做夏革.
她問上帝如何去看顧夏革.
她記得夏革是單親.

$^{41}$當然她是走的.
她單親.
她說上帝也紀念她.
她也覺得上帝應該可以紀念她.
幫助她.
於是她暫時不死.
於是她繼續重新回到神那裡.
回到教會.
再重新經歷神.
在這個階段她學到一件事.
在教會聽到一個十一奉獻.
她聽到之後又爭人家的錢.
又要養B.
還要奉獻.
但她心裡覺得應該要聽從神的話語.
於是她堅定聽從神的話語.
不過在我們外人看來.
屋漏兼逢夜雨.
不單她的經濟有困難.
她還在公司準備裁員.
一裁就裁了她.
照理說很困難.
要面對生活.
但很奇妙的是.
當她找到一份工作.
裁員介紹工作給她.
她就找到了.
不久前老闆就叫她.
不如你回來做兼職.
於是她有一份半工.
即是兩份工資.
她幫她還債.
還可以繼續好好生活.
又可以回到教會.
又可以養小朋友.
這個姐妹的經歷.
我們看到她祈禱之後.
她的生活有沒有完全改善.
會否因為她祈禱之後.
不用再工作.

$^{81}$靠上帝的手.
再調錢給她.
又不會.
她的身體和肉體的辛苦.
比平時更辛苦.
還要做多份工作.
睡覺都不夠時間.
但她因為經歷上帝的緣故.
心靈裡很平安.
她相信一切都不是偶然.
不是巧合.
就是因為這樣.
她就繼續堅定依靠上帝.
她經歷上帝.
就是因為這個人物.
叫做夏甲.
我也很想和大家透過夏甲.
這個姐妹.
去經歷上帝對夏甲的愛.
同樣對我們每一個人的愛.
所以為了大家不要睡著.
你和旁邊的人說.
神很愛你.
神很認識你.
神很認識我們每一個人.
就讓我們在這個時刻.
一起了思想.
夏甲心靈裡有一個很重要的問題.
就是問.
因為你的版本和我看的版本.
有少許不同.
他問這個問題.
用我這個版本.
在我這個階段.
我有沒有真的遇見這位看顧我的神呢.
上帝就在這段經文裡解答他.
一個很重要的問題.
他能不能找到這位看顧他的神呢.
我透過三方面和大家一起了思想.
第一件事就是.

$^{121}$我們要知道上帝原來對我們每一個人都這麼認識.
他知道你是誰.
神聽見夏甲的苦情.
他根本就完全明白夏甲是誰.
所以在經文裡.
你告訴我們.
就是使者遇見他的時候.
就說薩萊的侍女夏甲.
你從哪裡來.
你要去哪裡.
夏甲就說.
我是從主母薩萊逃出來的.
其實如果我們這樣看的時候.
好像很多餘.
明明你知道.
薩萊的侍女夏甲.
你明明都認識她.
為什麼要問一遍呢.
你問一問她.
其實可能我們說.
上帝沒有錯認識你.
他怕你不認識自己.
所以多問你一遍.
不過我們要知道.
聖經的故事都很特別.
因為當摩西寮寫創世記的時候.
可能他都要用一些.
我們以前知道是用口傳的方法.
記錄這些事跡.
其實夏甲本身的名義.
就叫做逃跑.
可能以前就是.
拿下來的時候.
就是要讓人更加深刻.
意思就是說.
逃跑啊逃跑啊.
你去哪裡啊.
就是這樣的意思.
但是我們知道.
夏甲這個人.

$^{161}$是真有其人.
以實瑪利.
他生的兒子.
以實瑪利都是真有其事的.
不過在聖經記錄的人物裡面.
可能未必一定是真的姓名.
因為告訴你.
我們都不認識.
阿拉伯文.
或者是希伯來文.
我們都不認識.
但是上帝就藉著.
這個撒萊的侍女夏甲.
讓我們看到.
上帝知道他的苦情.
他怎樣苦呢.
如果我們看.
本來他就是一個埃及的.
被賣的奴隸.
當時阿巴拉罕.
因為饑荒的時候.
落到埃及.
可能就是因為這樣.
他就得著了夏甲.
做撒萊.
服侍撒萊的婢女.
服侍撒萊.
成為撒萊的婢女.
他原本想著.
這一生都是奴隸.
但是當撒萊看上了他.
不如這樣吧.
你做我丈夫的妾.
我們俗語說.
當然有反應.
於是他不單成為阿巴拉的妾.
還要新.
真的有新人.
於是他就開始對主母.
有少許的面色.

$^{201}$撒萊因為面色受不了.
就要趕他走.
所以在這個苦情裡.
夏甲本來也得了少許.
希望人生的盼望.
但是突然間.
撒萊對他的虐待.
所以苦從中來.
不如逃走.
這一切上帝都知道.
上帝知道他所有的經歷.
但是上帝在他知道.
他最徬徨的時候.
他仍然要問他.
兩邊都要離開.
如果我們問上帝.
不如明知道他這麼徬徨.
馬上把他的應許告訴他.
讓他大快活.
大家開心.
但不是的.
上帝問完他之後.
然後指引他一些方向.
所以我們看.
為何上帝要問呢.
其實在輔導學上.
有一個叫做敘事治療.
敘事治療的意思是.
原來人有很多過去發生的事.
讓我們自己都偏激自己的故事.
而這些故事可能有偏差.
可能對自我的觀點上.
有模糊和不認識.
所以這個敘事治療.
是幫助一個人.
重新認識自己.
或者說一個很簡單的例子.
如果一個受虐待的人.
他的認知自己是什麼呢.
就是他自己不可愛.

$^{241}$他沒有價值.
被人拋棄.
他對自我的價值很低.
但當他用敘事治療的時候.
來幫助他重新思考.
究竟他是誰.
虐待的人既然不對.
就不應該因為虐待者.
而錯誤了待自己.
把他的故事扭轉.
讓他自己重新編寫故事.
或者說另一個例子.
我記得曾經做過婚前輔導的時候.
幫一對準備結婚的人.
來談談他們的情況.
其中一個姐妹.
她是未婚妻.
她就很情緒低落.
要和我們分享.
她說剛剛去見南方的時候.
覺得南方有些親戚.
看不起她.
說話很多骨.
她說你們吃這些東西.
帶著父母來.
大家很尷尬.
於是她一而再.
禁止父母說話.
在這頓飯裡.
她不斷禁止.
不要讓父母說話.
我問她為什麼要禁止.
父母說話呢.
她說怕要保護父母.
怕南方的親友.
看輕自己的父母.
我問她.
自己有沒有輕看父母呢.
她說沒有.
沒有輕視父母.

$^{281}$但因為她父母是農民出身.
所謂的蠻子都不認識.
比較粗人.
被輕視她也不察覺.
只有她自己察覺.
但當她受輕視的時候.
她心裡很不舒服.
我便和她談談.
輕視一個人.
其實是誰對呢.
在我們的判斷上.
究竟那個人是怎樣呢.
如果輕視人.
是因為他沒有修養的緣故.
因為他不能夠馬上體恤別人.
因為他不懂得尊重別人.
如果因為他的輕視.
而你和他埋單.
因為他輕視而你不舒服.
然後去保護父母.
帶來自己的情緒困擾.
這樣是否值得呢.
她想的時候.
為何要為別人的錯誤而埋單呢.
可能在我們今天的生活中.
常常為別人的錯誤而埋單.
但當她重新看自己的故事的時候.
她的生命開始有些改變.
她開始釋懷了.
她開始看別人怎樣看父母.
有甚麼重要呢.
最重要是我尊重父母.
最重要是我怎樣面對原生家庭.
我能夠自在地面對原生家庭.
這也是我們今天要注意的地方.
其實這個所謂的敘事治療.
在上帝也不陌生.
在聖經裡可以看到.
記不記得當夏娃犯罪的時候.
當夏娃躲起來.

$^{321}$上帝就問她.
你為甚麼躲起來.
她說她看到自己赤身露體.
上帝問她.
你為甚麼知道自己赤身露體.
如果是我們.
我們早就已經指出了問題.
當然是你吃了樹的果子.
我們一定會這樣指出.
但上帝就不問.
任由她說.
再讓她重新說這個故事.
這個故事是怎樣呢.
其實也很扭曲.
最後也好像上帝說錯了.
你做了這個女人.
做錯了.
但上帝任由她們說.
原來當一個人.
你要將她的故事說一說的時候.
她的生命會有改變.
所以我們也求主幫助我們.
其實很多時候.
是因為我們太不認知自己.
我們對自己不認識.
所以上帝問我們是誰的時候.
我們也不知道是誰.
我們是誰呢.
我是誰呢.
但我們很渴望.
上帝讓我們有弟兄姊妹.
讓我們心中有弟兄姊妹.
我們就去學習.
怎樣去分享自己的故事.
讓我們學習.
怎樣去聽別人的故事.
原來當我們聽一遍的時候.
或者在過去的時侯.
當聽別人說一遍的時候.
其實我也沒有很多力氣.

$^{361}$她自己就通了.
聽完一遍之後.
她就學會原來自己是這樣的.
然後她的生命故事就開始改寫.
上帝讓我們今天.
我們信了主之後.
不是立刻抽離我們.
不是立刻接走我們.
其中一個很重要的原因就是.
祂要讓我們更加認識自己.
我們是誰.
而我們越認識自己.
我們才能夠越認識這位上帝.
當然也可以反過來.
當我們越認識上帝的時候.
我們就越體現我們自己是誰.
越體現自己是誰的時候.
我們的生命才能夠有所改變.
才能朝著上帝的方向去走.
如果我們不認識自己的話.
第二點對我們一點意思都沒有.
但當我們認識自己的時候.
上帝就幫助下甲走第二點的方向.
就是當他認識自己是誰.
我就是這個婢女.
我就是撒萊的一個僕人.
一個奴隸.
然後上帝就指引他的方向.
他知道下甲.
然後他指引下甲應有的方向.
其實上帝很明白下甲的苦情.
但上帝沒有馬上將下甲抽走.
既然你逃走了.
好吧,我讓你逃走了.
離開這個這麼痛苦的路.
不是的.
上帝指引下甲做什麼呢.
就是你回到撒萊的身邊.
來去信服他,服侍他.
其實對下甲來說.

$^{401}$是一個痛苦的經歷.
他剛剛才逃走.
你就叫我回去.
他怎樣去面對呢.
但我們不知道.
不知道為什麼上帝不馬上讓他逃離.
其實當下甲十五年之後.
他一定都同樣要離開阿巴拉罕和撒拉.
但是十五年歸十五年.
上帝仍然要他在這十五年裡.
回到撒萊的身邊.
我們不明白為什麼.
但這個偏偏是要他學習的地方.
我相信上帝做的事.
很多時候都是一箭雙鵰.
他不單止要幫助下甲.
怎樣面對他不喜歡的女主人.
或者他要面對困苦的經歷.
他要釋放她.
我們不知道.
也可能上帝都同樣要讓撒萊去學習.
怎樣因為阿巴拉罕的緣故.
因為這個福中的寶寶的緣故.
都要忍受下甲.
或者為阿巴拉罕的緣故.
讓他不要過一個這麼內疚的生活.
如果下甲一走了之.
可能阿巴拉罕一生都會很內疚.
或者心裡的懸念都很掛念.
所以如果我們去看阿巴拉罕的生平.
阿巴拉罕其實是一個很軟弱的人.
我們可以看到.
如果在歷史上讓我們看.
當阿巴拉罕送走下甲二十馬來之後.
其實阿巴拉罕從此就住在別士巴.
為什麼要住在那裡呢?.
其實如果我們看.
他會常常和下甲有通訊.
或者和二十馬來有通訊.
當阿巴拉罕死的時候.

$^{441}$以撒和二十馬來是同時安葬阿巴拉罕.
所以我們看到.
這個阿巴拉罕同樣是一個.
如果我們正面說.
是一個有情有義的人.
他都要幫助阿巴拉罕在這件事上.
怎樣去面對.
但是我們不知道.
上帝做很多事都不用告訴你.
上帝做很多事需不需要告訴我們呢?.
我們都需要靠著主去學習.
上帝要我們去信服的功課.
我記得我年輕的時候.
信主的時候.
我們那時候是牧者.
你要教導我們.
如果你在工作裡面不如意.
人際關係不好.
你千萬不要辭職.
為什麼?.
你要搞好人際關係才辭職.
或者我們和現有的警隊同事都是.
很多時候.
萬一人際關係不好就打紙走.
我們都教他們不要打紙走.
先搞好人際關係.
因為原來跟人.
人際關係是跟人的.
你這個人際關係不好.
去到那邊都不好.
所以我們就要去學習.
所以我們那時候有主學.
我們一起和幾個姐妹一起.
除了上主學之外.
我們有自己的茶經.
有幾個好姐妹.
來跟我們說.
我現在可以辭職了.
為什麼?.
我們就跟著拍手.

$^{481}$為什麼?.
我們都明白.
因為他搞好人際關係.
他心裡可以坦蕩蕩辭職.
不是因為他要逃離老闆才辭職.
而是因為他知道他已經越過了.
這個都是我們要學習的地方.
我不知道.
但是我們就要學習.
當下甲回到撒萊身邊的時候.
他的生命就開始不同了.
這個上帝對他的引導.
而帶來給他生命的改變.
所以今天我們被上帝引導.
我們都同樣扮演一個很重要的角色.
我們都要去引導別人.
或者在座中間可能有父母.
或者有兒孫.
當我們要面對.
引導我們的兒女的時候.
我們怎樣去引導呢?.
我記得我兩個小朋友都大了.
小的都大學畢業了.
我記得我第一次懷孕的時候.
很開心的.
也很緊張.
我就說.
順便懷孕怎麼辦呢?.
怎樣去照顧小朋友呢?.
我想做姐妹的可能都會有這個心情.
怎樣去做一個好的父母呢?.
當時我自己有個原則.
讓我看見.
我問自己.
我對兒女有什麼期望呢?.
我就有一個答案給我自己.
我很渴望我的兒女能夠自負責任.
這個自負責任.
也是我照顧兩個小朋友.
一個很重要的原則.

$^{521}$所以當小朋友能夠自己吃飯的時候.
我從來都不追著他們.
我希望大家做媽媽,爺爺,婆婆.
不要追著小朋友吃飯.
因為他要自負責任.
或者當他長大了.
他自己洗澡的時候.
他就要自己洗澡.
雖然我都試過.
我兒子洗澡的時候.
我有一個星期.
看見他頭髮一團團的.
原來他不太懂得洗澡.
但是我都要他學習.
當他頭髮一團團的時候.
我當然幫他處理了.
但是他跟著下去就懂了.
這個自負責任是他要做的.
或者當他學寫字的時候.
他不會太心急.
因為寫得不好.
會被老師勸導.
所以他就拿著手.
其實不是他寫.
是父母寫的.
或者當他要畫畫的時候.
他很想子女貼堂.
於是我看見很多畫.
肯定不是三歲小朋友畫的.
就是父母畫的.
或者是婆婆,爺爺,奶奶畫的.
所以當我們以為.
他貼了堂.
幼稚園,小學貼堂.
以為他很厲害.
但是其實是你做的.
他都還沒有自負責任.
所以其實我很坦言地說.
我從來沒有捉著我兒子的手.
寫過一隻字.

$^{561}$或者幫他畫一幅畫.
全部都是他自己得來的.
所以沒有貼堂.
但是這個最重要的.
就是在他成長裡面.
我的原則是什麼.
我怎樣引導他.
正如這個原則.
我覺得上帝都是這樣引導我.
上帝沒有常常告訴我.
指著我的鼻子.
要拉著我走.
但是他就要我去學.
怎樣去自負責任.
所以我們今天教徒都是一樣.
在這個罪惡的世界裡面.
其實上帝是要你我.
都要自負責任.
怎樣過一個誠信的生活.
你不能夠說.
上帝啊 請你把罪拿開.
那我就不會犯罪了.
好像犯罪是跟上帝有關.
但是上帝就要我們.
怎樣自負責任.
幫我們怎樣去抉擇.
我們的人生.
而這個抉擇.
我們不是行為主義.
而是一個誠信的生活.
我們恩著耶穌都可以.
但是我們都需要靠著.
耶穌的信服聖靈的生活.
以致我們的生活.
能夠過一個合神心的生活.
其實昔日教會的二元論.
今天都同樣影響教會.
要麼就是法理錯人主義.
要麼就是律法主義.
就常常覺得這個不對那個不對.

$^{601}$去指責別人.
要麼就是放縱主義.
神都很愛我.
我犯了一點錯也不要緊.
但是我們求主幫助我們.
這個二元論.
這兩個極端的東西.
都不是我們要的.
我們要的是什麼呢.
我們要的是此時此刻.
我們要如何被聖靈光照我們.
我們要如何過一個信服的生活.
所以其實今天靈修是必須要的.
靈修不是一個律法主義.
靈修是我們一個必然的.
跟上帝的關係.
其實每天的靈修.
就好像跟上帝的價值觀.
一樣.
上帝的價值觀.
擺在這裡就是這樣.
但是我們今天怎麼想呢.
上帝告訴你要愛仇敵.
不是吧.
上帝說那人怎麼愛.
我還要為他祈禱.
這樣我受不了.
那怎麼行呢.
覺得上帝是離地的.
我們不聽上帝的話.
當我們以為是離地的時候.
我們的生命是十年如一.
不改變的.
所以我們求主幫助我們.
聖聖的生活.
是我們必須要學習.
信服 遵從.
當我們肯信服 遵從的時候.
我們不是說.
基督徒聖聖的生活.

$^{641}$不是說蓮枝.
而是說你對上帝的遵從.
有多少.
我們即是說信主.
我們剛剛信主.
但是我們就學習.
如何去遵從上帝的生活.
你的靈性的進步.
是比一個十年都不去.
遵從神的話語的人.
可能是更加成熟的.
所以我們求主幫助我們.
當然我們不要看別人.
我們看到上帝很愛我們.
但我們需要學習什麼呢.
我們需要看清楚.
上帝給我們的方向.
還有我們需要信服主.
又說大家不要睡覺.
我們又跟隔壁的人說.
上帝很愛你.
不過要你遵從.
信服祂的話語.
所以很渴望大姐姐妹.
去分析上帝給我們手上的聖經.
我們這一代的基督徒.
我們這個新的基督徒.
我們是一個很幸福的基督徒.
因為我們人手都拿著聖經.
其實現在不是了.
現代我們人手拿著手機.
手機可能有十個聖經的版本.
其實以前列王記的時候.
是皇帝才有聖經.
是皇帝才有機會抄一本聖經.
或者才有一本聖經給他看.
我們今天大家都做皇帝.
我們都是公主,皇子.
我們在神裡面是很有榮耀的身份.
而這本聖經就是成為我們的一個指引.

$^{681}$願神幫助我們.
我們熟讀聖經.
熟讀神的話語.
我們生命才有改變.
上帝不單認識我們.
上帝很渴望我們遵從祂的話語.
第三方面是什麼呢.
上帝應許我們未來.
同樣上帝將下甲鳥帶去指引他.
你回到主母身邊.
但同時給他一個最重要的應許.
給他一個應許是什麼呢.
就是他和撒拉平分抗禮.
下甲將來會生十二個族長.
而撒拉同樣是十二個族長.
以色列有十二個支派.
原來上帝祝福下甲同樣是有份量的.
不過有不同的地方.
今天很多巴勒斯坦地區的人.
我知道看錯了.
他們死爭爛爭.
上帝不是要應許這些.
原來上帝要藉著以色列人.
成為萬人的祝福.
但上帝應許下甲.
祝福下甲同樣是安慰他.
指出他人生路向.
有上帝同在.
有上帝看顧.
有上帝同樣將這個應許給他.
撒拉因著有上帝的應許.
他的生命有一個重大的改變.
即是他回到撒拉面前.
他很恐懼.
他可能很害怕撒拉來虐待他.
但因為有這個永恆的應許.
帶來給他能夠衝破所有的困難.
所以我們知道今天.
其實我們今天面對人生的際遇.
是順境逆境其實不重要的.

$^{721}$最重要的是你怎樣過.
你的態度是怎樣面對.
原來當我們清楚上帝.
給我們一個美好的應許的時候.
我們就有能力去面對我們的人生.
所以很多時候.
我都會和大家分享.
上帝其實沒有為我們做很多事.
上帝最重要做一樣很重要的事.
就是為我們.
藉著耶穌基督釘在十字架上.
就是為要將我們的罪了交換.
賜給我們一個新的生命.
而這個新的生命.
帶來給我們一個新的眼光.
今天我們有永恆的盼望.
我們不只是有今生的祝福.
當然今生的祝福.
上帝是樂意讓今生的祝福伴隨我們.
但如果我們今天只是為了今生的祝福.
我們是得不到永恆的盼望.
所以我們看四福音的時候.
大家要看清楚.
四福音沒錯.
看到耶穌的憐憫.
對人的憐憫.
你因著信就可以得到耶穌的醫治.
但耶穌的醫治也好.
趕鬼也好.
教導人也好.
最重要的是.
當你面對上帝對你的醫治.
教導趕鬼.
得到自由之後.
你是怎樣.
你究竟是否選擇.
一路去跟從耶穌.
榮耀歸給耶穌.
我們看到在四福音裡面.
耶穌醫治的人.

$^{761}$不是所有人都跟著耶穌.
沒錯有很多人得到耶穌的醫治.
但這個身體得救.
不等於他永恆的得救.
他永恆的得救是什麼呢.
就是要他甘心樂意的了根重柱.
可能在人生裡面有很多風浪.
但這些風浪如果你知道.
你走得準.
你看得準.
我們跟著耶穌.
就得到盼望.
今天我們就會輕視這些苦難.
正如我們知道亞伯拉罕.
亞伯拉罕的人生.
我們知道他有175歲.
聖經形容他75歲.
由喀拿再次被上帝呼召出去.
我們不知道他離開迦勒底的吾爾是多少歲.
但他本來是離開迦勒底的吾爾.
帶著他拉.
即是他爸爸.
他所有的家產.
到喀拿之後.
直到他爸爸離開世界之後.
他再次蒙召.
75歲才到迦南地.
但如果我們知道.
當他生以實瑪利的時候.
是86歲.
意思是又過了11年.
但當他生以實的時候.
他幾歲?.
99歲.
意思是在人生裡面.
有很多時候.
至少有四分一世紀.
25年都看不到上帝英渠.
但上帝說.
你的子孫像海沙一樣.

$^{801}$好像天上的星一樣.
都看不到.
我們今天可能在聖經的英渠.
或上帝的英渠裡面.
你是不是?.
當我們在雨中.
40年一雨.
50年一雨.
天災人禍的時候.
我們就說.
上帝是不是快來了?.
耶穌是不是快來了?.
我們很多時候都會問一遍.
沒錯我們可以問.
但我們今天最重要.
我們有一個.
相信上帝英渠是永不落空.
這樣的心態.
所以其實.
耶穌什麼時候來又如何呢?.
耶穌什麼時候來.
我們有沒有影響我們的生活呢?.
所以我們求主給我們一個眼光.
我們驕途相信上帝的大能.
福音的大能.
是幫助我們.
要搞一切相信的.
但我們不要只是頭腦的相信.
而是我們整個生命.
怎樣來相信耶穌拯救我們.
而且不單是今生的福祿.
還有永恆的盼望.
我信主.
我18歲信主的時候.
好像很早.
但我18歲信耶穌的時候.
我覺得很遲.
為什麼這麼遲.
為什麼這麼遲信耶穌.
早一點就好了.

$^{841}$但當我信耶穌的時候.
我的教會有一個字句.
因為我們是二樓的教會.
幫我們很多人.
我們幾百人的教會.
在二樓寫著.
我來了.
小結糧得生命.
並且得的更豐盛.
我當時回教會.
因為剛出來工作.
有一個同事邀請我回教會.
其實他邀請了我很多次.
其實我因為不好意思.
不好意思再推他.
好吧隨便吧.
邀請你一次之後.
邀請我一次.
我大概坐在第三第四行.
邀請我一次之後.
以後踢我也不走.
因為當我那天回到教會之後.
我整個人完全的改變.
我這個改變是什麼呢.
我本來在少年的時候.
在我讀書的時候.
我人生覺得很迷茫.
我覺得人生出來.
面對死亡.
生老病死.
是一件我覺得很荒謬的事.
為什麼人生出來要面對死亡.
而且我以前的耳朵很靈.
我聽著新聞的時候很害怕.
新聞不是家人不努力.
而是殺人放火.
真的出街很危險.
那時候我覺得這個世界很不平安.
還要加上我眼見很多人.
生活很痛苦.

$^{881}$樣子又不開心.
面對地盤工人.
這一生就是這樣.
人生哪有開心的地方呢.
但當我回教會之後.
我相信了主.
我看到.
當我第一天回教會.
應該說可以回轉.
我看到神原來很愛人.
而這個很愛.
就是沒錯人生出來面對死亡.
但原來上帝要我們.
藉著人生.
經歷更豐盛的生命.
所以我由那時候開始.
我頭幾個月.
我真的很看見.
什麼叫豐盛生命.
這個豐盛生命是什麼呢.
就是當我和人聊天的時候.
我發現我對人的認識了解.
突然間多了.
我突然間用了.
我用了另一個眼光.
去看身邊的人.
甚至我覺得.
我用一些說話.
用聖經金句也好.
用神的話也好.
用我自己認知的說話也好.
我和人聊天的時候.
我發現原來帶給人生命改變.
這個也是給我有一個很初嘗.
什麼叫豐盛生命.
所以我整個人扭轉了.
所以我不知道大家.
怎樣去經歷這個信仰.
我越覺得.
我越初信.

$^{921}$一個初信主的時候.
你怎樣信耶穌呢.
其實對我們的生命也很重要.
所以我們不要隨便叫人.
決志信耶穌.
我們需要讓那個人.
知道他真的有罪.
認罪悔改.
他的生命才有改變.
我記得我有時帶一些.
陪談訓練的時候.
到今天我帶很多次陪談訓練.
很多大型節目都說.
他很難得決志信主.
不如我們先不要說罪.
覺得說罪有點困難.
不如說一下.
耶穌怎樣祝福你.
愛護你.
以後一生跟你同行.
如果你陪談的時候.
不說罪說什麼.
我們如果讓那個人不清楚.
什麼叫信主.
信主就是將我們的罪了.
交給耶穌.
本來我們罪該萬死.
但是我們靠著耶穌.
有一個新生命.
如果我們不是這樣說的話.
我們怎樣能夠得著.
讓人得著真正的福音.
讓人看到自己有罪.
你會去決志歸向主呢.
所以有些人.
我去訪問的時候也有很多.
有很多人都還沒得著.
都還沒看到自己有罪.
你先不要幫他決志.
通常那些我不幫他決志.

$^{961}$讓他慢慢等.
有些人等了十年.
有些人十年.
突然間光照他.
我因為兒子的緣故決志.
我也有看到.
但是這個決志就很重要.
看到自己的罪.
是一件很重要的事.
所以願神幫助我們.
耶穌願意和我們一起同行.
我們的神是藉著聖靈.
當我們今天決志信主之後.
是聖靈和我們同行.
靠著聖靈我們有一個新的生活.
而靠著聖靈.
上帝就光照我們.
你究竟是誰.
當我們認識自己是誰的時候.
耶穌就指引我們.
或者聖靈就指引我們.
走當行的路.
當我們走當行的路的時候.
耶穌也讓我們看到.
去永恆的盼望.
在人生的路上.
即使有多大的困苦.
我因著耶穌的緣故.
我有力量去面對.
所以最後我想和大家說.
或者大家可以問自己.
不用問隔壁的那個.
其實你想去哪裡呢?.
你現在在哪裡?.
你想去哪裡呢?.
你願不願意.
神和你一起.
去走你人生的路呢?.
如果你願意.
神和你一起.

$^{1001}$走人生的路.
你今天應該怎樣.
抉擇你的人生呢?.
願神祝福我們.
先有一個祈禱.
大家先來祈禱.
你要告訴耶穌.
告訴耶穌你想去哪裡.
然後你願不願意.
耶穌和你一起去呢?.
(靜靜地看著自己的影子).
七位主人.
我們都謝謝你.
我們謝謝你.
為你甘心地死在十字架上.
這個重要的標誌.
你要告訴我們.
原來當我們繁仰望你的.
就必得著這個救恩.
但當我們仰望你的時候.
我們都需要被聖靈光照.
讓我們看到.
我們真是一個罪人.
我們本來就要死.
甚至我們本來是一個.
永恆的死亡.
但我們為著.
耶穌你甘心地為我們死.
讓我們今天.
回看我們過去人生的時候.
求你教導我們.
怎樣自己說自己的人生.
讓我們對自己人生的.
一些錯誤的看法.
一些錯誤的價值觀念.
求你先改變我們.
讓我們用聖經.
一個最正確的話語.
來帶領我們.
引導我們.

$^{1041}$讓我們走正路.
讓我們走誠誠的路.
求你教導我們.
向遠方望.
讓我們看到.
你給我們永恆的盼望.
因著你的永恆.
讓我們看到今天.
我們即將面對人生的痛苦.
我們的身體可能會漸漸衰敗.
我們生活可能會有很多意外.
但我們能夠靠著耶穌.
過一個充滿信心.
充滿喜樂.
充滿平安.
充滿盼望的生活.
我們求耶穌.
你將這個盼望.
放在我們心裡.
也讓我們今天.
面對我們生活際遇的時候.
求你指引我們.
讓我們知道.
因為有耶穌.
讓我們同行.
我們就無所懼怕.
多謝你聽我們祈禱.
我們祈禱過.
仰望你.
祝耶穌基督的名.
\newpage



\section{提多書 1:10-2:1}
\label{sec:lRt1SGTCNSE}
\textbf{2024.11.24 《講真話的教會》 提多書1\_10-2\_1 (和修版) 曾裔貴牧師}
\newline
\newline
連結: \href{https://youtube.com/watch?v=lRt1SGTCNSE}{\texttt{https://youtube.com/watch?v=lRt1SGTCNSE}} ~~~~ 語音日期: 2024-11-25
\newline
\newline
\hyperref[sec:Y3F7I4ibIlg]{\small{< < < PREV SERMON < < <}}
~
\hyperref[sec:index_chronic]{\small{[返順時目]}}
~
\hyperref[sec:index_scriptual]{\small{[返順卷目]}}
~
\hyperref[sec:KpxypRPehLQ]{\small{> > > NEXT SERMON > > >}}
\newline
\newline
提多書 1:10-2:1
\newline
\begin{longtable}{cl}
\hline
\hline
章節 & 經文 (和合本修訂版)\\
\hline
1:10 & \begin{tabularx}{0.7\textwidth}{X} 因為也有許多人不受約束，說空話欺哄人，尤其是那些奉割禮的人。 \end{tabularx} \\ \\ \relax
1:11 & \begin{tabularx}{0.7\textwidth}{X} 這些人的口必須堵住，因為他們貪不義之財，將不該教導的事教導人，敗壞人的全家。 \end{tabularx} \\ \\ \relax
1:12 & \begin{tabularx}{0.7\textwidth}{X} 克里特人中有一個本地的先知說：「克里特人常說謊話，是惡獸，貪吃懶做。」 \end{tabularx} \\ \\ \relax
1:13 & \begin{tabularx}{0.7\textwidth}{X} 這個見證是真的。為這緣故，你要嚴厲地責備他們，使他們在信仰上健全。 \end{tabularx} \\ \\ \relax
1:14 & \begin{tabularx}{0.7\textwidth}{X} 不要聽猶太人無稽的傳說和背棄真理之人的命令。 \end{tabularx} \\ \\ \relax
1:15 & \begin{tabularx}{0.7\textwidth}{X} 在潔淨的人，凡物都潔淨；在污穢不信的人，甚麼都不潔淨，連心地和天良也都污穢了。 \end{tabularx} \\ \\ \relax
1:16 & \begin{tabularx}{0.7\textwidth}{X} 他們宣稱認識神，卻在行為上否認他；他們是可憎惡的，是悖逆的，不配做任何好事。 \end{tabularx} \\ \\ \relax
2:1 & \begin{tabularx}{0.7\textwidth}{X} 至於你，你所講的總要合乎那健全的教導。 \end{tabularx} \\ \\
[1ex]
\hline
\hline
\end{longtable}
$^{1}$我們一起同心祈禱.
可以在我們洪恩家.
看到有很多的弟兄姊妹.
慢慢認識真理.
信主.
得救.
連真姐妹.
她在晚年.
經歷了這麼多的困難.
但她經過教會.
主的恩典.
主的聖靈.
感動她.
跟朋友進來.
就是這樣開始.
他們在我們洪恩家.
過去大半年.
信主的成長.
雖然她不太懂得詳細說明.
但她的內心.
之前的痛.
和現在在主耶穌的恩典裡面.
弟兄姊妹的關心.
在代禱裡面.
她的改變.
我們看到神的真理.
真的在人的心裡面.
會有增長.
會有改變.
求你使我們洪恩家.
每一個弟兄姊妹.
都能夠對真理有認識.
剛才主席讀的經文.
已經在在告訴我們.
原來方話.
只能騙到很少部分的人.
而且會讓人最終發現.
但其實更加重要.
騙不到我們這位.
無方言的真神.

$^{41}$所以我們懇求主.
讓教會是一個.
講真道的教會.
是一個講真話的教會.
我們一起在神的面前.
去讓真理在教會裡滋生.
去成長.
我們仰望主.
今早的聚會.
崇拜的講道.
求你聖靈的工作.
在我們每一位弟兄姊妹的心靈裡.
去產生思想.
一起去學習.
奉耶穌基督的名.
阿們.
很感恩有姐妹提及.
我們過去在舞堂.
感宣.
因為我們浸禮聚會.
我們每個人都要講見證.
但我們很少回來.
我們很多弟兄姊妹沒有參加.
但又聽不到弟兄姊妹的見證.
我就想.
為什麼我們不在崇拜裡.
讓弟兄姊妹可以逐一分享她的見證.
讓大家一起.
更加彼此有共鳴.
一起去學習.
大家看到.
很美麗的花.
姐妹們為昨天.
二十一週年的福音樂曲小組.
的服侍.
擺上.
又看到.
弟兄姊妹.
昨天我們去了.
我和曾姑娘去了大堂.

$^{81}$不是去燒烤這麼簡單.
而是有兩間教會的合作.
一起在洪福村做一些福音工作.
差不多有六十人的家庭.
這樣來到一起.
讓教會可以更加認識到.
這些基層的家庭.
其實是一個很大的祝福和和祥.
所以.
教會.
我很想將這個題目.
分兩次.
因為下個星期都會是我講.
分兩次我們去看提多書.
很重要.
保羅在這個時候.
他已經坐完第一次牢.
他要來吩咐他的年輕傳道者.
去到不同的地方.
提多.
就被托付在克里特島.
在地中海一個算大的島嶼.
相當美麗的地方.
不過.
在當時.
保羅派提多去傳道的時候.
那些島嶼上的人.
是什麼情況呢?.
他們的特質是什麼呢?.
在智慧行間已經浮現了.
還要保羅引述.
一個克里特本地的先知.
似乎也是信主的先知.
說克里特人有一些特徵.
他們常常說謊話.
而且.
很貪吃.
但是又懶.
那怎麼維持這幾樣東西呢?.
我們就要看看.

$^{121}$保羅怎麼來提多.
有一個機會.
去在那裡事奉.
改變很多人的心靈呢?.
有沒有吃過這條香蕉?.
6,520萬人民幣.
其實值兩個半而已.
只是貼在牆上.
卡特蘭這個意大利藝術家.
他給了一個名字.
那個膠布.
貼在牆上.
他的名字叫做喜劇演員.
他說他刻意做出來.
是要諷刺人們.
亂來的.
這樣去投標的藝術品.
他覺得這是一個價值連城的藝術品.
你懂得欣賞的話.
你就會起標的了.
起標價是500萬.
居然美金620萬成交.
加上蘇富比的拍賣佣金.
以及行政費用.
是6,520萬人民幣成交.
那個買家就說他承諾.
三天內他會吃那條香蕉.
有說明書.
有出售紙.
教你怎麼做.
其實在座的每個人都懂得做.
世界充滿了很多我們覺得很荒謬的事情.
連這個拍賣官都認為.
他這樣叫出來.
他自己覺得很荒謬.
從來沒想過貼一條香蕉在這裡.
就這樣來拍賣成交.
他都覺得好像很假.
世界充滿著這些荒謬.
當然能夠給錢的人.

$^{161}$一定是賺錢很容易.
就是加密貨幣的一個有錢人.
太容易了.
又藉這些機會.
推高他自己的加密貨幣的價值.
這樣的一個社會.
這樣的世界.
對於我們在在來說.
我們笑完就沒事了.
因為我們根本不會有這樣的機會.
去做這些荒謬的事情.
不過.
荒謬.
謊言.
會不會充斥著在我們周邊的地方.
周圍的環境呢?.
我們要小心.
因為這一段經文告訴我們.
我們信奉的是一位不說謊話.
沒有謊言的真神.
我們一起來看看.
我們今天要跟大家說的經文.
大家見過的這一段經文.
這一節就叫做.
「真的假不了」.
我引述陳志雲所說的.
因為當時他被無線去告他.
後來呢.
最終他的申辯.
宣判無罪.
他一句話.
真的假不了.
我們一起進入經文裡面.
再加一個故事給大家去了解一下.
很恐怖的.
原來我們人生.
如果有一個我們內心裡面.
一個習慣的.
錯誤的價值觀.
可以帶來很大的傷害.

$^{201}$有一個亞洲的年輕人.
相當努力讀書.
最終畢業.
去應徵.
據聞是在法國的一個地方.
不約而同.
他的申請信.
很出色的畢業的成績.
不被錄用.
他就很奇怪.
問了一間一間公司都不成功.
沒有人請他.
他終於想通了.
因為我是亞洲人.
受到歧視.
就向僱主查詢和質問.
為什麼我的履歷這麼好.
會不被取錄呢?.
也符合你們公司入職要求的文化等等.
僱主終於跟他說.
我們每次接見一個求職者.
我們會審查你的背景.
發現你有三次坐火車沒有買票.
這個年輕人就抱怨.
有什麼大不了的事呢?.
我三次沒有買票而已.
全部是情有可原的.
僱主就說.
是的.
第一次你沒有買票的時候.
是你剛剛來到法國一個月.
當時你的辯解是.
你不知道這個國家的火車系統.
所以你忘記了買車費.
但是之後的兩次.
你是有心真正的騙我們國家的系統.
國家是用一個信任制.
不會有人查票.
也很容易.
你跳過去.

$^{241}$就可以過了坐車.
果然.
這個僱主告訴他.
我們很著重的是你的誠信.
以及你對制度的尊重.
你沒有這麼做.
我們的公司沒有能力.
也不可能不斷監察著你有沒有誠信.
所以我們是不會僱用你.
也說不上是對你的歧視.
如果這件事是真的.
我們就要不斷問自己.
我們在人生裡面.
我們有沒有我們自己的誠信系統呢?.
沒有人監察.
我們仍然要有誠實呢?.
所以我們最近有時跟堂委.
跟公說.
參與事奉的弟兄姊妹.
如果被發現.
原來是很虛荒的.
我們立即不會再用.
因為你的態度根本不適合事奉神.
我們事奉的是一位沒有謊言的神.
這段聖經其實挺重的.
提到要責備那些克里特人.
我們再進入經文裡面.
我們一起來看.
神的僕人耶穌基督的使徒保羅.
為了使神的選民.
順從和認識.
什麼是敬虔的真理呢?.
就是這個真理.
就是我們有盼望.
這一位沒有謊言的神.
一開始第一節.
第二節.
已經將神這一方面的屬性.
神有很多很多的屬性.
愛,公義等等.

$^{281}$全能,全知的屬性.
特別強調.
在這個教會現在的處境.
神的無謊言.
在萬古之先.
已經應許的永生.
所以保羅在這裡這樣的描述.
這樣的開場白.
就是想說什麼呢?.
就是神是一位沒有謊言的神.
就是說他已經用預言.
啟示將永恆不改變的真理.
來講述出來.
人才會發現.
發現之後.
神就設立全道的人.
去講明神的敬虔的真理.
所以神的道是不變.
是永恆.
是真實.
但是.
全神道的人.
就需要有這個配合.
你留意一下這裡下面就說了.
可能頂住了.
先說自己一點的例子.
初信的時候.
我的教會,母會很嚴格.
那個全道人也很嚴格.
剛初信.
每個星期回團祭.
星期六.
「銳貴,今天有沒有靈修?.
有沒有讀經?」.
每天都問.
那慘了.
真的有時候沒有.
那怎麼逃避呢?.
有一段時間不敢回教會.
因為怕說謊.

$^{321}$就逃避.
但後來想想也不是辦法.
沒理由這樣逃避.
只能夠.
應該怎麼說呢?.
試過真的有一次.
真的說謊.
在全道人面前說謊.
有啊.
每天都有.
很好.
其實就是沒有.
後來真的不行了.
就要坦白地承認.
我們還沒到這種生命的成長.
原來坦白.
比假裝還要重要.
坦白就能夠很坦然地面對自己.
面對自己現在的階段.
慢慢地拾級而上.
基督徒,基督教.
我們是不是只停留在一種.
說完就當做了呢?.
其實我們大家一起要反思.
很久之前聽一個童工說.
現在祈禱手太容易了.
按一按WhatsApp就已經祈禱了.
但按完之後就忘記了祈禱.
當然這絕對不算是說謊.
但是我們應該要有一個具體的行動.
真的要為我們所關心的事情.
感受祈禱.
這就是我自己起初的時候.
初信階段.
因為被傳道人迫得很緊.
所以在這樣的緊張之下.
試過去否認,說謊.
試過假裝好像很好.
但其實還沒有到那個階段.
還沒有到生命的成長.

$^{361}$盼望這成為我們大家.
互相提醒,彼此勉勵.
我們繼續下去的時候.
我們留意一下.
經文繼續這樣告訴我們.
到了適當的時機.
就是神的時機.
藉著傳揚福音.
把祂的道顯明.
保羅很清楚認定.
神的工作是怎樣的運作.
神的啟示.
神的預言.
甚至已經透過人寫在聖經裡面.
是否就這樣完結呢?.
不是的.
原來神的啟示是完備了.
但是需要有不同的人.
被神興起,感動的人.
來負起傳揚的責任.
傳揚的責任.
就托付在歷世歷代.
每一個時代的工人.
有主學老師.
有講福音見證的.
有為神做其他的工作的.
都是一個見證.
就是這個責任.
被保羅托付.
按照著救主神的命令.
這樣交托了.
所以一個事奉的人.
你要知道.
你要忠於神的托付.
來做好你要做的工作.
這是其中一件事.
所以保羅看.
這個預言,啟示.
去到成為一本書本.
但是需要提多去傳揚.

$^{401}$提多要負起這個責任.
要傳揚.
那怎麼傳揚呢?.
下面就繼續說.
第五節我就沒有引述.
第六,第七和第十,十一節.
就可圈可點了.
有一個教會已經成立了.
這個教會的人.
在信耶穌之前.
有三個重要的特徵.
性格.
習慣.
說謊.
貪吃.
懶做.
是不是元朗的特質呢?.
是不是元朗的特產呢?.
是不是洪水橋的人有呢?.
大家不需要回答.
因為似乎全香港.
我想每個人都或多或少會有.
但是當時就有這樣的教會.
在這個克里特島裡面.
那是不是遺棄這個島.
任由它自生自滅呢?.
不是的.
保羅知道可以改變人心的.
就是第五到第十一節.
在克里特島裡面.
去找一些神改變的弟兄姊妹.
去做長老,做執事,做教會領袖.
很特別.
越是那個島的國民.
那些人是這麼負面的性格.
那種壞習慣.
那種內心.
我更加要在那個教會裡面.
提醒大家.
要找到這樣的人去設立.

$^{441}$就是說.
如果有這樣的弟兄姊妹.
那個教會,那個社會.
那個海島的人就有生命了.
若有無可指責的人.
只作一個婦人的丈夫.
兒女也是信主的.
家庭是經過了牧者的牧養,教導和觀察.
看到他們家庭的信仰是很真實的.
沒有人告他們放蕩.
他們可能是指他們的子女.
甚至是整個家庭.
不受約束.
放蕩不受約束.
就是說喜歡做什麼就做什麼.
我的家庭的事是我們家庭的事.
不關牧者的事.
我們自己的事.
我離開教會之後做什麼事.
關你什麼事.
這樣就不是不受約束.
這樣就變成有可指責.
保羅勸勉提多.
要這樣留意,觀察和牧養這樣的家庭.
這樣的家主就可堪被設立.
成為教會的領袖.
一般人看到都會很信服.
這樣的領袖在教會做治理,帶領.
你可以想像得到.
這樣的無可指責在克里特島.
相對於在教會以外.
這麼混亂,這麼虛荒,這麼謊話.
這麼貪吃,這麼懶做事.
這麼好吃懶飛.
教會就成為了明亮的燈台.
所以教會對領袖的牧養,揀選,嚴格是必須的.
這樣就會令到神的福音真理.
被人尊重,被人高舉.
你可以想像得到.
如果教會的人比外面的人更虛荒.

$^{481}$更貪吃,更懶做.
更只會指點而什麼都不做.
那怎麼會帶來其他人對你產生尊敬呢?.
所以是一個相輔相成的影響.
也不要變成一個惡性循環的壞榜樣.
再下去就是第七節.
又是講監督.
監督既然是神的管家.
你想想不是管理你的家.
是管理神的家.
必須要公平,公義,有慈愛,做事有原則.
更重要的是不可以虛荒.
必須無可指責.
已經是第二次講了.
不要自負,不要暴躁.
那個性情.
別人對你說的一些話.
未必是合你的內心.
立即發脾氣.
立即有自己的位份就拍桌子.
這樣就絕對不適合做了.
因為你控制不了自己的情緒.
不於酒,不好鬥,不貪財.
當時的歐洲人.
大家看過《戰狼三百》.
那個斯巴達人很兇狠.
很暴戾地成長成為戰士.
相信在古代都有這些性情的人.
這些海島的居民.
可能以前的世代是否做海盜.
這個就不太清楚.
但你可以想像到.
在這個島的弟兄姊妹.
他們信了主.
要將他們提升被揀選成為神的管家.
所以有這些條件.
其實大家看上去很合理.
很應份,很應該.
應該是這樣.
弟兄姊妹你也應該這樣看.

$^{521}$我們的領袖應該是這樣.
如果不是這樣.
那就真是教會的失職.
因為聖經寫得很清楚.
時代可以不同.
但人心的形態其實差不多.
所以找一個很暴戾的人做領袖.
就預計每次開會都會拍桌子.
因為他很容易就暴戾.
應該是很溫柔.
很有原則地溫柔.
對人很客觀,持平.
不會很主觀,偏見.
這樣才讓教會做事.
在原則和法理下.
法理下又有愛,包容.
教會就會很健康地成長.
不要貪財.
領袖收藏很多禮物.
你不給我禮物.
我就看你是一個二等的信徒.
請你坐在一旁.
教會就慘了.
原來爭相巴結不同的領袖.
那就很大件事.
一視同仁的眼光.
每一個都是神的兒女.
每一個都是神愛的兒女.
不貪財背後的價值觀.
丙姐妹我們一起去看到.
保羅對於教會.
如何面對在克里特島這麼多困難.
原來是要找對的人.
去在治理之中.
即是說外面風風雨雨.
不是太大問題.
因為神的道屹立不變.
因為做物主是沒有謊言的.
他會運用他的恩典.
去改變人心.

$^{561}$所以如果教會裡有這樣的弟兄姊妹.
外面的人只會羨慕.
只會很願意成為.
克里特教會的門徒.
第十節之後我不再引述了.
因為也有許多人不受約束.
說空話欺控人.
尤其是受奉國禮的猶太人.
或者進猶太教的外邦人.
這些人的口必須堵住.
要封實這些人的口.
如何封呢?.
可以用膠紙最快.
但是那個心沒有被封住.
改變不了他的心有什麼用呢?.
如何改變人的心呢?.
教會設立監督和長老.
教會設立被神默許的人.
就是經過觀察.
目樣是可堪能夠具備成為神管家的人.
這樣的教會就有規矩.
有制度.
有大家認同的那種心靈的.
那種阿們的教會.
這是我們大家一定要追求的方向.
和為教會的領導層祈禱的.
很重要的方向.
這裡說.
這些人的口必須堵住.
因為這樣說話的人.
背後一定有他自己的私人的.
不好的動機.
這樣說話.
說些虛空的話.
又欺控別人的東西.
欺控即是會令人害怕.
令人怕你.
其實不是.
動機就是貪不義之財.
將不該教導的事教導人.

$^{601}$就敗壞人的全家.
在前後的大段經文裡.
由第五到第十一節.
中間夾雜的就是.
一個好的監督.
一個好的長老.
一個被驗證.
可以堪稱為一個好的領袖的長老.
就是一個教會的重要基石.
就是教會能夠成長.
必不可少的元素.
當然教會可以有很多東西是重要的.
硬件是重要的.
各樣的東西都要重要.
但是保羅緊張體重的.
很多.
雖然你是一個年輕的傳道者.
你未有我的使徒的位份.
但是你已經負起了.
這個全陽真道的責任.
全陽真道第一件事.
你要找到合適的人作為領袖.
弟兄姊妹大家一起監察.
大家一起去為教會祈禱.
需要這樣的人作領袖.
剛才讀了.
克里達人的特徵.
本地先知這樣說.
克里達人常說方話.
是惡獸.
貪吃.
懶做.
原來再加多一件事.
惡獸.
似乎很暴戾.
第十三節保羅說.
這個見證是真的.
這樣寫聖經跟提多說.
提多也很害怕.
我要服侍要牧羊的島裡面的人.

$^{641}$原來未信主之前是這樣的性情.
那怎麼去做呢?.
不過保羅說.
為者緣故你要嚴厲責備他們.
使他們能夠在信仰上健全.
仍然說的是信仰的範疇.
不是干涉誰人的私事.
不是干涉你任何的事.
不是進入僭越你的個人空間.
而是關乎到信仰.
成為神的兒女.
弟兄姊妹.
我們就要追求神的性情.
我們就要在神的真理下.
互相服侍.
所以提多是碰這個課題.
所以不是說什麼私隱的問題.
是提多要牧羊教會.
要有智慧地選擇領袖.
保羅已經說了他的標準.
哪些人才能合資格成為領袖呢?.
就是剛才信民說的那些.
經過長時間的觀察.
牧羊能夠深交.
能夠看到他內心的動機.
一心為神.
不為什麼.
不為錢.
不為名.
這樣來服侍.
就真是一個可以被神用的教會.
堵住是否可以是一個好辦法呢?.
相信不是.
相信必須是我們能夠在信仰上.
健全的弟兄姊妹.
每個弟兄姊妹.
神都給你均等的機會.
神的道去到你的內心.
你怎樣去承接神的道.
怎樣去追求.

$^{681}$怎樣去活出在家庭裡面.
作美好的見證.
怎樣去在生活裡面.
讓人感受到你內心.
那個對神的愛.
那個主的敬畏呢?.
在網上找到關於「誠信」.
待人真誠.
不說謊.
不歪曲事實.
信.
就是信守信用.
一落千金.
就是說了就兌現承諾.
除非不承諾.
履行自己應該承擔的義務.
這就是「誠信」.
今天我們這個講道裡面.
很緊張.
很重要的是.
教會是否一個講真話的教會呢?.
我們弟兄姊妹是否一群誠信的人呢?.
教會絕對需要拒絕.
講是非的人.
教會最大的自命.
是很多講是非的人.
講一些可能根本不真實的事.
來影響人心.
勉勵弟兄姊妹.
我們知道有些事有疑惑.
直接找當事人.
去尋求真相.
當然不是大庭廣眾地說.
是私下去尋問他.
你這樣提出.
是不是一個這樣的問題.
你可以反映你的角度.
反映你聽到的.
不是對質.
是彼此勉勵分享.

$^{721}$當事人說是.
就處理了誤會.
否則根本不需要.
在當事人耳外傳.
這絕對是教會的傷害.
克里特這個地方.
本身已經是一個講謊的小島.
保羅很清楚.
那個島的人.
不是這樣的人不可以信耶穌.
更加要信耶穌.
不過信了耶穌之後.
裡面的弟兄姊妹.
就要拒絕謊言.
因為神禁止.
因為神的聖情.
是誠信真實的.
所以保羅就要叫提摩太多.
在教會裡面.
必須要選擇這樣的人.
做領袖做監督.
教會就帶來真理的模樣.
所以求神幫助.
我們洪恩家.
能夠在誠信.
在真實裡面.
大家有看這套在印度.
算是最賣座的.
全球的這套電影.
很勵志的講教育.
探討教育.
特別是印度的處境.
是一個很高的學府.
等於是印度的MIT.
理工大學裡面.
其中的第三個.
我右手邊的第三個.
我們的翻譯叫拉朱.
這位年輕人.
他很害怕.

$^{761}$因為他家出身很貧窮.
貧窮得很厲害.
他讀書很厲害.
父母很想他將來成為.
出色的工程師.
他就不斷求神拜佛.
不斷在他讀書期間.
什麼神都拜.
希望保佑他考試成績好.
最終他的成績都不行.
後來出貓.
後來那個主角.
就來幫他.
那個老師院長.
就逼他.
你要向我承認.
是那個主角教你.
這樣出貓.
他兩為其難.
不知道怎麼辦.
最終他選擇在三樓.
跳樓下去希望自殺.
跌斷了雙腳.
後來因為他那間.
是很出色的理工大學.
就有些印度的大公司.
畢業之前會認證招聘.
一間大公司就約見他.
聽他的故事.
問他為什麼成績那麼差.
為什麼會跌斷雙腳.
他終於說.
因為他已經昏迷了兩個月.
在那兩個月昏迷裡.
他終於醒了.
那些假神幫不了他.
那一切都令他更加懼怕.
更加害怕.
更加不可以有好的表現.
他坦白向幾位面試官說.

$^{801}$他過去的遭遇.
還有為什麼他要跳樓.
他不想陷害他的好朋友.
這個主角.
那個面試官.
很巧妙地這樣說.
其實你不需要那麼坦白.
我們的公司需要一個.
很有擅長外交手腕的人.
你不需要那麼坦白跟我們說.
你那麼坦白.
我們很難會聘請你的.
當然這是故事.
是電影的橋段.
這個拉朱他說了一句話.
他說對.
以前我是這樣的.
但經過這兩個月.
我的很深的反思.
我終於發現.
我斷了兩條腿.
但我真正能夠起來做人.
我可以對我自己的良心.
對我的好朋友出賣.
我坦白.
你們不聘請我不要緊.
你們收回這個offer.
這個履歷我送給你們.
但我可以站起來.
頂天立地做人.
我走出去.
我會找我另外的工作.
跟著那幾位面試官.
感受到這個年輕人.
在自己的利益最大前提下.
居然不隱瞞自己的軟弱.
過失.
坦然的承認.
這個面試官他們談完之後.
立即叫他回來.

$^{841}$我們正正需要你這個人才.
因為我做面試做了幾十年.
未見過一個人在利益當前.
會毫無保留的.
為自己的信念去捍衛.
而不是為了一份工作作假.
承信啊!弟兄姊妹.
這是電影.
電影也很想告訴我們.
原來承信最終可以有收穫.
承信永遠是最重要的.
所以這裡告訴我們.
我們的弟兄姊妹.
我們在神的家裡.
我們學習的神的真理.
我們就需要在我們的生活裡運作.
讓人感受到我們是洪恩家.
不僅是硬件有優化.
是我們這群弟兄姊妹有成長.
因為有神的真道.
我們每個星期一起學習.
一起在愛裡成長.
最後.
大家都應該認識.
星期六科學技大學頒發一個博士學位給他.
榮譽博士學位.
還有梁朝偉先生.
沒有人不認識他.
在科技界.
現在是全球三間最大的公司.
他已經是最大.
就是NVIDIA.
輝達.
是AI的電腦晶片的運算.
是他來公司發明.
而且潛力是繼續無限擴大.
所以是全球最大的市值公司.
但是.
現在成功當然是各樣東西都理所當然.
每個人如果有他的股票.

$^{881}$他的升值是倍升不得已.
但是他訪問的時候.
他說我曾經做過洗廁所.
他讀大學的時候家裡很窮.
台灣的移民.
去到美國的波特蘭大使.
他說我讀書的時候做過洗廁所的工作.
但他有一個信念就是.
我要做一個真的我.
即便我一生都是做洗廁所.
我要做到最出色的最好的.
這些信念令他不斷尋求成長.
突破.
當然後來開了這間公司.
和幾個好朋友.
一路發展.
但原來曾經.
如果你是對股票很認識.
曾經2019年的時候.
他公司的股值.
由100\%回到20\%.
每個人都拋棄他的股票.
因為當時未有人知道.
AI的重要.
影像晶片的運算的重要.
當時還是微軟.
當時硬件的晶片.
主導的市場.
可以想像當時.
他是一個很大的低潮.
不過他和他的同事說.
那些同事沒有一個人離職.
那些股票跌得這麼厲害.
跌了八成.
他說我每天仍然想要做什麼.
我的家人仍然愛我.
信任我.
我公司的同事仍然信任我.
我仍然有這份初心要去做.
就成為了今天.

$^{921}$全球市值最高的輝達公司.
前天就拿了科技大學頒發的榮譽博士學位.
就是這位黃仁勳先生.
現在你看他.
我看他.
看他的報道.
理所當然.
這麼厲害.
這麼出色.
但是人家在低潮的時候.
在人家內心的初心.
不斷堅持裡面.
就成為了今天的成功人士.
所以教會沒有任何的捷徑成長.
是我們每一個弟兄姊妹.
都尋求真實的神.
都在我們的生命裡面.
讓神的道改變我們.
為大家做一個願意開放自己.
與人相交.
用愛與人交往的弟兄姊妹.
教會就會成長.
一起禱告.
感恩.
天父讓我們看到.
雖然他過去的經歷是很困難.
但是主耶穌的愛.
成為他最好的保護和安慰.
我們很盼望他的兒女.
都慢慢可以順主.
多謝你們讓我們成為一間.
在紅水橋的教會.
弟兄姊妹也是來自不同背景.
我們甚至有泰國的華僑.
很盼望我們大家一起.
在誠信.
在真理上去追求.
成為一間合主心意的教會.
在我們的生活.
尤其是在我們低潮.

$^{961}$陷在個人利益的時候.
更加我們要揀選這位真實的神.
以致我們能夠有美好的見證.
為主耶穌作見證.
感恩禱告.
奉耶穌基督的名.
阿們.
\newpage



\section{提多書 2:11-15}
\label{sec:KpxypRPehLQ}
\textbf{2024.12.01 《你仲返教會?點解ge ?》 提多書2\_11-15 (和修版) 講員 : 曾裔貴牧師}
\newline
\newline
連結: \href{https://youtube.com/watch?v=KpxypRPehLQ}{\texttt{https://youtube.com/watch?v=KpxypRPehLQ}} ~~~~ 語音日期: 2024-12-01
\newline
\newline
\hyperref[sec:lRt1SGTCNSE]{\small{< < < PREV SERMON < < <}}
~
\hyperref[sec:index_chronic]{\small{[返順時目]}}
~
\hyperref[sec:index_scriptual]{\small{[返順卷目]}}
~
\hyperref[sec:3EdQX7vZPlY]{\small{> > > NEXT SERMON > > >}}
\newline
\newline
提多書 2:11-15
\newline
\begin{longtable}{cl}
\hline
\hline
章節 & 經文 (和合本修訂版)\\
\hline
2:11 & \begin{tabularx}{0.7\textwidth}{X} 因為，神救眾人的恩典已經顯明出來， \end{tabularx} \\ \\ \relax
2:12 & \begin{tabularx}{0.7\textwidth}{X} 訓練我們除去不敬虔的心和世俗的情慾，在今世過克己、正直、敬虔的生活， \end{tabularx} \\ \\ \relax
2:13 & \begin{tabularx}{0.7\textwidth}{X} 等候福樂的盼望，並等候至大的神和我們的救主耶穌基督的榮耀顯現。 \end{tabularx} \\ \\ \relax
2:14 & \begin{tabularx}{0.7\textwidth}{X} 他為我們的緣故捨己，為了要贖我們脫離一切罪惡，又潔淨我們作他自己的子民，熱心為善。 \end{tabularx} \\ \\ \relax
2:15 & \begin{tabularx}{0.7\textwidth}{X} 這些事你要講明，要充分運用你的職權勸勉人，責備人。不要讓任何人輕看你。 \end{tabularx} \\ \\
[1ex]
\hline
\hline
\end{longtable}
$^{1}$我們一起同心祈禱.
感恩天父讓我們能夠在今早.
我們有歌唱敬拜你的時間.
更加有嘉豪弟兄的得救見證.
看到他的內心.
由軟弱.
在主裡面剛強.
能夠開始有信心面對星學.
人生.
逐步逐步走出.
他自己的一些陰暗作弊.
都看到神的救恩在他的生命裡面.
今早我們真的很想思考這個題目.
為什麼我們在星期日.
我們要回教會.
而不是去喝茶.
當然喝茶絕對不是醉.
但是我們好像將上午分別為聖.
來敬拜我們造物的主.
如果有人問我們為什麼.
這真是一個很好的回答機會.
天父讓我們能夠有這份肯定.
從心裡面認定我們是屬於你的.
以致我們能夠成為一個.
很榮耀見證主的群體.
今早我們打開聖經.
提多書這一段經文.
盼望主的榮耀顯現.
成為我們大家的盼望.
和努力面前的動力.
感恩禱告.
奉耶穌基督的名.
阿們.
剛剛收到網絡斷了.
可以了是吧?.
Sorry.
有沒有人問過這個問題?.
為什麼你要回教會?.
我在人生裡面.
有三個人問過我這個問題.

$^{41}$剛剛結婚.
1990年.
住在元朗.
有一次那個劉宇.
因為租地方剛剛結婚.
租地方.
業主住在麗港城.
年輕.
因為是他爸爸的物業.
在麗港城和她的男朋友進來視察.
因為好像有些工程他要幫忙做.
就上來我家.
就問了一個很有趣的問題.
為什麼你家裡沒有電視機?.
他覺得很好奇.
為什麼.
你也是年輕人.
那時候我二十多歲.
為什麼你家裡沒有電視機?.
我就解釋給他聽.
太太就希望小朋友出生之後.
就不要看電視.
就專心來讀書.
還有不要受很多電視劇集的影響.
他覺得很詫異.
為什麼呢?.
不看電視怎麼生存呢?.
這是當時我還記得的一個片段.
這是為什麼呢?.
有沒有人問過你呢?.
所以你會知道.
我兒子從來沒有看過TVB.
所以他有時候去朋友家裡.
他很好奇.
原來TVB是不好看的.
這是其中一件事.
接著要問一個問題.
有人問我.
為什麼你要做牧師呢?.
這是很多年前了.

$^{81}$有一個上司.
那時候我在政府裡面工作.
我決心讀神學.
他說為什麼你要做傳道.
那時候他不知道什麼叫牧師.
為什麼你要做牧師呢?.
我就解釋給他聽.
我的信主.
我的改變和我的呼召.
他覺得很匪夷所思.
政府這麼穩定.
你去做傳道.
那時候我應該是剛剛兩年的批准過了.
就是去到永續的終身制.
你為什麼不做呢?.
這麼好.
他覺得很詫異.
最終還是真的離開了.
第二個問題問我為什麼呢?.
還有.
這次是我媽媽問我.
她叫我貴仔.
為什麼你要進神學院呢?.
那裡是做什麼的呢?.
我媽媽不識字.
她知道我要去長洲.
知道我連工作都不做.
知道我要進長洲神學院.
就要去住宿.
已經結婚了.
已經有孩子了.
為什麼你要這樣進去長洲呢?.
她就很奇怪,很詫異.
為什麼你又政府工作不做呢?.
她叫我貴仔.
為什麼呢?.
她完全不懂.
那時候她還沒有信耶穌.
三個問題問為什麼呢?.
今天這段經文很想跟大家.

$^{121}$當然我是從牧者的角度.
來跟大家分享.
為什麼你還在教會呢?.
大家真的要回去問一問.
為什麼你仍然是回教會呢?.
我們開始今天的經文.
我們一起來看看這裡.
在十一到十五節.
不知道你有沒有留意到.
在幾節經文裡面.
保羅說了兩次主耶穌的顯現.
第一次固然就是.
耶穌基督的降生那一次.
不過沒有明顯說祂的降生的顯現.
而是說一份恩典的顯明.
這是十一節的經文.
接著到了第十三節.
說到主耶穌一定會再回來的榮耀顯現.
那就是祂的肉身的顯現了.
所以在十一到十三節.
說了兩次主耶穌的顯現.
這裡很值得我們去看看.
經文的情況.
成為我們大家彼此的鼓勵.
為什麼呢?.
當人問你為什麼你還回教會.
為什麼你還信耶穌成為基督徒呢?.
十一節因為神拯救萬人的恩典.
已經顯明出來了.
其實恩典是看不到的.
恩典只會在信耶穌弟兄姊妹的生命裡面.
流露才會看到.
所以這裡保羅比較隱晦地.
其實是說到救主的顯現.
降世.
祂成為了人.
將人帶進救恩裡面.
這個恩典就顯明了.
其實十四節.
待會我們讀的時候.

$^{161}$你就會看到.
再一次說到主耶穌的降生.
所以這裡說的顯明.
其實就是主的降生.
祂成為人這件事.
十二節.
「這恩典管教我們.
要我們棄絕不敬虔的生活.
和屬世的私慾.
在今世過自律.
正直敬虔的生活」.
這個你可以想像得到.
救恩帶給我們.
作為傳道牧者.
有一個託付.
有一個責任.
就是要將人.
管教.
帶進去永恆裡面.
在今世.
就是有限地上的人生.
我們可以去改變人心.
改變了人心.
他就會過一個自律.
正直敬虔的生活.
所以傳道牧者.
是不會變成夕陽工業的.
永遠有人類.
永遠有市場.
永遠要將人改變.
這裡聖經就用管教.
其實是practice(練習).
是要操練的.
跑長跑就知道了.
如果你一個月都不跑.
你的肥肉就出來了.
你就很難再要鍛鍊.
要經常運動訓練.
管教的字就是訓練的意思.
這裡告訴我們.

$^{201}$主的降生.
這樣的顯明.
原來是恩典的顯明.
使我們能夠讓人.
透過認識恩典.
認識耶穌.
就可以過一個.
在屬世的生活裡面.
跟第一章說的那些.
格理底人.
有懶惰.
有惡.
有貪吃.
那些人是不同的.
在教會裡面.
原來是有另外一批人.
同樣都是面對當時的社會的勢情.
但是過的生活.
是一個自律.
正直.
敬虔的生活.
十三節.
期待那福樂的盼望.
就是我們偉大的神.
和救主耶穌基督榮耀的顯現.
我們留意到.
兩節聖經裡面.
說了兩次主耶穌的降生.
和祂的將來的再回來.
兩個時空裡面.
就是夾雜著十二節.
要訓練很多的.
信主的格理底的弟兄姊妹.
格理底島上的人.
來去有改變.
當然不只是那個島.
是全世界聽聞福音的人.
都可以在今世.
來過一個自律.
正直敬虔的生活.

$^{241}$留意十四節.
耶穌基督為我們獻上自己.
是為了救贖我們.
脫離一切不法的事.
其實只是補充十一節那段經文.
就是主的降生.
恩典的顯明.
這一件事.
所以.
你可以讓我們看到.
並為自己潔淨屬於他的子民.
就是那熱心行善的人.
接著第十五節.
保羅再一次.
來講完那個.
恩典的顯明之後.
就提多.
他有什麼責任要做.
這些事你要講明.
要運用一切的權柄.
勸勉人.
指正人.
不要讓人輕看你.
這個就是.
提多他在格里底這個地方.
他需要去做的工作.
就是作為一個勸導人的人.
他自己要講明.
這些真理.
叫人明白.
在這個勸導過程.
就會將人帶進那個.
救恩的改變的行列裡面.
就是說.
我們透過每一次的接觸.
講解.
分享.
和勸導.
是會將人改變的.
我們又再看看.

$^{281}$這段經文的結構.
剛才已經提過了.
兩次的主降臨.
一次就是恩典的顯明.
保羅提到主的降生.
是神恩典的顯明.
顯明出來.
就是讓人可以透過認識恩典.
就可以得著生命的成長.
這個成長.
是要透過管教.
操練.
和訓練的.
所以.
不訓練的基督徒.
你的成長是慢的.
訓練的基督徒.
願意被訓練的基督徒.
那個成長是顯著的.
是的.
第十三節就是榮耀的再臨.
就是耶穌基督.
將來肉身會再回來.
所以兩個顯然的中間.
就是十二節了.
就是牧者要訓練每一個信徒.
信徒可以領受一個新生命.
但是這個新生命是要不斷經過訓練的.
要過一個正直.
以及能夠在神的面前.
敬虔的生活.
弟兄姊妹我們要留意一下.
牧者.
很特別.
第一節.
其實上一個星期我們已經講了第一節.
和第十五節.
兩者好像有一個括弧.
先講一下這個例子.
剛剛29日.

$^{321}$不知道大家有沒有留意新聞.
有兩位拿了樹仁大學的榮譽博士.
一位就是蕭永青先生.
他是中原地產的創辦人.
拿了當時工商管理的榮譽博士.
他說他覺得受之有愧.
他自己也沒有讀過大學.
說不上什麼管理人.
但其實他太謙虛了.
他的成就對香港也有很大的貢獻.
另外一位就是這位.
稱為港版的德蘭修女.
有沒有人認識她呢.
潘寧卓.
這位在城寨做了超過40年的宣教士.
這位女士.
她1966年.
當時你可以想像得到.
香港是暴亂的.
其實香港當時很亂.
她當時只不過是22歲女.
由外國來到香港.
只身一個人.
只剩下100元港幣.
不知道怎麼辦.
有另一位她認識的宣教士.
帶她去城寨.
從此在城寨結合了一份感動.
在城寨裡面一直做到她退休為止.
超過40年了.
他說上一年.
2023年.
當年她的父親是一個引君子.
一個導遊.
說你收留我的父親.
讓他有改變.
當然了.
她說不介意.
因為她的心.
根本上就是被這份愛感動.

$^{361}$她很想將裡面的引君子.
黃賭毒.
黑社會.
帶她去信耶穌.
你可以想像得到那個不容易.
她開始的時候在裡面租了一個地方.
開辦青年會的一個活動中心.
開始希望多些人來參加這些活動.
打康樂棋.
其他的活動.
希望可以打發時間.
不要去做犯罪的事情.
但是開始的時候.
那些黑社會來恐嚇她.
站在門口.
說我幫你守護著.
我做你的保鑣.
她說我不需要.
接著就讓人來作弄她.
那個宿舍裡面.
全部找一些糞便.
塗滿了窗戶.
塗滿了地面.
希望她知難而退要走.
接著就自己一個人清潔.
一步一腳人.
慢慢去感化那些人.
有兩個.
一個叫做Winson.
一個叫做阿賢.
兩個是黑社會的打手.
看到這個女士.
為什麼會有這樣的生命力呢?.
由反對.
由嘲笑.
由攻擊.
因為是黑社會的大佬.
叫她去搗亂.
騷擾他們.
到後來.

$^{401}$看到這位潘寧卓.
那份生命的愛.
無私的愛他們香港人.
慢慢改變.
真的成為了基督徒.
由吸毒.
變成開始不反對.
慢慢站在門口.
進去裡面.
信主.
後來要帶她的朋友.
那個阿賢就是Winson帶信主.
你可以想像得到.
其實一點都不容易.
要管教.
要訓練.
是經年累月的時間.
慢慢勸導一個人.
看到基督福音的真光.
是需要很多的心力.
很多的時間.
去做這些灌溉的工作.
弟兄姊妹.
你和我都是一個心田.
我們用什麼去灌溉我們的心靈呢?.
有沒有一些傳道牧者.
一些我們熟悉的前輩.
在你的生命裡面做灌溉.
使你能夠有這份生命力的成長.
到今天你現在能夠在教會裡事奉.
來榮耀神呢?.
這是我們大家一起要面對.
要回答的問題.
當如果再有你微信的朋友.
很好奇為什麼你堅持要回教會.
堅持要帶你的子女回教會.
帶你的兒孫回教會.
你要告訴他.
這裡是有生命的地方.
這裡能夠讓你的心田成長.

$^{441}$能夠遇見神的地方.
我們回到經文.
一和十五節.
很清楚提多.
去到格里底.
他很快就發現.
那個地方一點都不簡單.
裡面的人.
保羅都說得很清楚.
是很麻煩的.
又懶惰.
又貪吃.
又惡.
好像惡獸.
保羅這樣說.
保羅說.
這個見證是真的.
你就不要走.
你就要嚴嚴地教導他們.
傳道人就是要在一些.
沒有人想去的地方.
你就要幫助裡面的人.
不要看死那些人.
你要幫助裡面的人.
提多收到這封信.
看到和他看到的是一樣.
這是天和地連成一線的呼召和感動.
呼召你可以不回應.
但如果你一回應.
那個後果.
就要面對很多的困難.
至於你.
這是二章一節.
你卻要按著健全的教義說話.
健全的.
是要很全面的.
不要主觀的.
以及要客觀地將神全面的真理說清楚.
例如.
今天有一些福音派是靈恩的.

$^{481}$很強調神蹟的.
什麼都強調神蹟的.
你不講方言就不夠屬靈的.
甚至第四波已經到了.
在現在多倫多的Toronto Blessing.
已經到了很誇張地.
要恢復先知.
恢復使徒的職份.
但是.
我們會不會過度了.
變成過度理性呢?.
我們什麼都用我們的理性去分析.
就好像不能夠容許.
超自然的任何事情發生.
兩者都是極端的.
所以一個健全的教義.
牧者在教會裡面的治理.
需要有健全的教義.
全面不要有偏離一些自己的喜好.
這就是一個很重要的.
保羅的教導.
接著到了十五節.
剛才我們讀過了.
這些事你要講明.
要運用一切的權柄.
勸勉人.
指正人.
不要讓人輕看你.
相信是提多.
要看自己是神的僕人.
有神的呼召.
他不是保羅的跟班.
不是保羅的紅人.
不是保羅來去欽點.
不是.
是神的僕人.
那個呼召.
所以所謂的權柄.
是忠於將神的道講明和勸勉人.
這樣的牧者提多.

$^{521}$要在這個格里底這個地方.
這個海島裡面.
來做這些工作.
我們又繼續下去看看.
由第二節一直到十四節.
很特別的出現了.
不斷提多.
要在純正的教義上來說話.
第二 第三 第四.
一直到第六 第七 第九 第十.
為什麼沒有了第八節呢?.
因為待會兒會再講.
這裡提到有家庭的倫理.
在社會上的倫理.
就是說.
只不過是舉例的.
涵蓋的.
就是大部分的人.
提多就要在教會去教導.
去管教這樣的人.
四次提「勸」.
勸年長的男士.
勸年長的婦女.
勸年輕的人.
勸僕人.
就讓我們看到.
提多就是做這個工作.
第八節很特別.
你要以身作則.
樹立好行為的榜樣.
在教導上要正直.
要莊重.
言語要純全.
無可指責.
這樣反對的人.
找不到什麼可說我們的壞話.
就自覺羞愧.
這裡提了好幾樣言語.
你自己的言語.
要成為榜樣.

$^{561}$你自己的為人.
要成為別人的榜樣.
這是其中一件事.
說的話要純全.
純全就是不要有你自己的偏見.
有你自己的限制.
盡量來保持客觀.
另外.
這裡也告訴我們.
反對的人會找.
牧者傳道的把柄.
來說壞話.
找不到也說.
何況找到.
更加說.
所以就要留意.
保羅說.
你不要去找其他人說什麼壞話.
你做好你自己作為牧者的呼召.
既然神安排你在格里底島上.
幫助那些格里底人.
他們說明是這樣的.
未信耶穌.
那個生命就是這樣的.
好像惡獸.
他們懶惰.
又餵食.
這樣的一群人.
但是原來是可以改變的.
原來是基督耶穌的恩典.
淋到他們身上.
你講明的福音.
和講明的教義.
他們是會改變的.
人心.
原來你不要以為不能改變.
人心有神的道.
就可以改變了.
所以.
我們不是要和反對的人對抗.

$^{601}$而是要告訴他們.
其實不要浪費時間去敵對.
大家一起為那些未信的人.
去傳福音.
使更多人.
得到神的救恩.
這樣互相的勉勵和互相的激勵.
就是勸勉的很重要的出發點.
這是第八節.
夾雜在第二第三節.
老年人.
老年婦女.
老年男人.
勸年青人.
勸年青婦女.
還有勸地上的奴隸.
即是包含主人.
不過時間不夠.
本來不講這麼詳細.
就是說每一個階層.
都涵蓋的.
牧者有一個天職就是.
每一個都是你的羊.
你都要去教導牠.
老人家.
你都要勸他.
老年女士.
你又要勸她.
年青人.
和你一樣年紀的.
你更加要勸導他.
僕人.
你又要勸導他.
那個奴隸我勸他.
我應該虐待他才對.
不是.
你當他是一個人.
勸導他.
主人.
一樣要勸導他.

$^{641}$有錢的.
沒有學識的.
都要勸導他.
這就是神安排教會有牧者的原因.
所以第八節.
真是一個很好的提醒.
和一個勉勵.
要以身作則.
來成為別人的榜樣.
這裡就是提多書.
簡單的一個分段.
你就可以想象.
在教會裡面.
有多少事情我們需要做.
有多少奇妙的事情.
教會可以帶來生命的改變.
我想起這個例子.
很值得和大家分享.
這位是印度的約翰·卡米根.
這位印度的.
其實他的身世是很坎坷的.
因為他有嚴重的麻風病.
麻風病就是侵蝕神經系統.
使你不覺得自己會痛.
就是沒有痛的感覺.
所以.
你走路不小心踩到.
你是不會避開的.
所以很多麻風病人.
為什麼手腳會骨折呢?.
因為麻風是很恐怖的.
那時候是不能醫治的.
或者不懂得醫治.
所以就會把他放在一些地方.
自生自滅.
沒有人理會.
所以生活環境是很困難的.
在印度就更加了.
看你為很低等的人.
自生自滅.

$^{681}$所以有時候那些地方有很多老鼠.
老鼠咬手腳.
咬鼻子.
咬耳朵.
很多人就這樣咬到體無完膚.
流血不覺得痛.
咬完也不覺得痛.
因為他已經沒有痛覺神經.
這是麻風最恐怖的地方.
卡米根就是一個這樣的人.
他的臉都扭曲了.
血肉模糊.
輾轉.
他在那個地方經常搞事.
因為沒有人接受他.
沒有人接納他.
覺得他很恐怖.
好像一個怪人.
所以越是這樣的時候.
他就越是用一種對抗.
來回應身邊的人.
很想製造麻煩.
引人注意.
後來作者的媽媽.
就是一個去印度的宣教士.
非常有愛心的一位宣教士.
看到卡米根.
就走過去.
用很多時間安慰他.
幫助他.
跟他有很多的禱告.
最終卡米根.
真的決心信主.
信主.
但是仍然面對他自己的身體.
他的面容.
仍然擔心.
人們會不會接納他.
有一次他信了主之後.
洗禮.

$^{721}$洗禮就是權宜的洗禮.
不是去大堂洗禮.
他提出一個挑戰.
我現在是不是神的兒女.
宣教士說當然是.
我想回教會.
怎麼辦呢.
這個朋友信了主洗禮.
他想回教會.
宣教士的媽媽.
她的兒子其實是一個醫生.
在英國去印度傳教.
他在印度出生的一個醫生.
很厲害的.
保羅班德醫生.
面對一個很兩難.
要不要接納卡米根呢.
如果不接納.
就會令他更加挑戰基督教.
原來根本不是講愛的.
保羅班德醫生.
之前去了當地的教會.
跟當地的牧者說了.
他的病是不會傳染的.
但他的樣子很醜.
很恐怖.
你們接納他吧.
是不會感染的.
他們想了一段時間後.
就讓他來.
到了當天.
連卡米根自己也覺得很詫異.
有人肯接受我來崇拜.
他就帶他們去.
和宣教士的媽媽一起去.
去到教堂.
大家都很凝重.
人們會怎麼看他呢.
會不會彈開呢.
過了不久.

$^{761}$有一個印度的弟兄.
遠距離看到卡米根.
就握手叫他過來坐.
叫他坐在他旁邊.
一起敬拜.
若無其事.
當他是一個普通人.
對卡米根從來沒有被人接納.
人見到他都這麼恐怖.
都會走開.
第一次感受到.
原來在基督裡面.
是會被接納的.
改變了他對人.
對教會的看法.
慢慢地過了一段時間.
一段很長的時間.
保羅班德醫生再去印度這個地方.
去看當地的復修.
一些地方.
幫助麻風病人.
做一些復健等等的地方.
還有一些基本的手工.
他們會幫忙做.
可以來賺錢.
他要介紹一個人給你認識.
這個人是我們全院社裡面.
最優質做小器材零件的人.
在我們印度這個社區.
他的工作是得獎的.
就是約翰卡米根弟兄.
他見到保羅班德.
很多年沒見的醫生的時候.
他的笑容.
你可以想像得到.
半個臉都沒有了.
但是張開嘴巴在笑的那種感情.
那種美麗.
是一種扭曲的美麗.
基督教往往就是這樣.

$^{801}$往往就是在一些不可能的地方.
有宣教士堅持.
有不同的牧者.
默默地去做.
去勸勉.
去帶領.
去讓人看到基督.
看到神的改變能力.
弟兄姊妹.
一起在洪恩家.
我們來做福音的工作.
改變不同的人.
大家可能來自不同的背景.
我都說了.
我小時候家裡十三.
已經開始家裡又可以拿到煙牌了.
在家裡吸煙沒有人再打我.
罵我了.
因為已經不能打了.
已經罵不了那麼多了.
喜歡的話還收起來.
回家就收起來.
但是收起來不了就是那股味道.
被媽媽姐姐罵.
中三都不管了.
由得你了.
一起吸煙了.
跟爸爸哥哥那些.
變成是這樣的了.
是這樣成長的了.
但是.
恩典可以改變.
我信耶穌的一刻.
是突然間早上起來.
以後不再吸煙.
到今天都沒有再吸.
聞到煙味都會感覺到厭惡.
吸兩包煙一天的一個年輕人.
是突然間的改變.
你可以看到舊恩.

$^{841}$恩典的改變能力和奇妙.
小時候八歲九歲.
已經每天去偷家裡的錢.
去吃東西.
已經是比家豪還要慣犯.
去做這件事.
後來信主真的慢慢知道.
不應該做.
就沒有再做了.
當然十幾歲已經不做了.
因為被家人發現了.
不斷打到跳起來就沒有做了.
弟兄姊妹.
我們這裡是在說.
為什麼你還要成為基督徒.
當很多人覺得信耶穌已經是很outday的時候.
信耶穌幫不了他什麼的時候.
信耶穌只不過是一種自我安慰的時候.
其實不是.
恩典改變人心.
改變你和我的心.
我們要讓恩典在教會裡滋生.
使更多人得到神國的福音.
譬如上兩個月.
在下村.
一位很熟悉弟兄的爸爸.
是很原居民.
很迷信的.
但最終晚年的時候.
要趕去醫院跟他說.
第一次跟他說福音.
兒子叫我來為他祈禱.
他說可以.
祈禱沒問題.
但你不要叫我信耶穌.
因為他真的很辛苦.
祈禱沒問題.
這次不說耶穌.
只祈禱.
他說可以.

$^{881}$再過了兩個星期.
真的不行了.
真的來不及了.
真的很嚴重了.
兒子再叫我去.
他是很迷信的原居民.
他要做很多儀式.
這次去為他祈禱.
他看到我來.
我為你祈禱,他說好.
這次不如信耶穌吧.
他就點點頭.
我說我看不清楚.
因為他的手上的膠布.
纏著儀器.
夾著的.
手指綁得很厚.
手指頭那麼大.
我說這樣吧.
如果你真的信耶穌.
你舉起手指頭吧.
他慢慢地舉起.
手指頭很大.
我說清楚了.
看到了.
跟他一起決志.
信主.
馬上洗禮.
那天晚上離開世界.
你可以看到.
鄧紫堅的父親.
很難信主.
他能夠信主.
也是在座彭先生的好朋友.
可以想像到.
救恩可以改變的能力.
是很大的.
所以我們傳道牧者.
今天我用我的角度去看.
還回教會?.

$^{921}$為什麼?.
我是傳道牧者.
我們接觸這些生命.
有時候一兩句說話.
改變一個人的生命.
是很感恩的.
那種快樂是很大的.
在在未必是前線.
去遺留的地方.
去其他地方探訪.
但是.
其實你自己的生命.
都是經過主耶穌的改變.
而帶來的生命.
很盼望大家一起努力.
在神的家裡.
我們去將恩典.
藉著我們自己的生命.
去改變其他人.
一起感恩祈禱.
多謝天父.
你的愛就是這樣.
藉著耶穌基督帶領我們.
兩次的降臨.
主耶穌在十誡上的犧牲.
你降生成為人.
將我們拯救.
用十誡的大愛.
用你的身體.
用你的寶血.
但是第二次的降臨.
就是關乎到我們每一個人.
信和不信.
都要面見你.
多謝你讓我們有生之年.
成為基督徒.
成為你的兒女.
成為可以改變別人生命的使者.
多謝主奉耶穌基督的名.
阿們.

$^{961}$謝謝大家.
\newpage



\section{哥林多後書 4:5-7}
\label{sec:3EdQX7vZPlY}
\textbf{2024.12.08 《光明之子》 哥林多後書4\_5-7 (和修版) 陳韋安牧師}
\newline
\newline
連結: \href{https://youtube.com/watch?v=3EdQX7vZPlY}{\texttt{https://youtube.com/watch?v=3EdQX7vZPlY}} ~~~~ 語音日期: 2024-12-09
\newline
\newline
\hyperref[sec:KpxypRPehLQ]{\small{< < < PREV SERMON < < <}}
~
\hyperref[sec:index_chronic]{\small{[返順時目]}}
~
\hyperref[sec:index_scriptual]{\small{[返順卷目]}}
~
\hyperref[sec:rX2pkPqwylA]{\small{> > > NEXT SERMON > > >}}
\newline
\newline
哥林多後書 4:5-7
\newline
\begin{longtable}{cl}
\hline
\hline
章節 & 經文 (和合本修訂版)\\
\hline
4:5 & \begin{tabularx}{0.7\textwidth}{X} 我們不是傳自己，而是傳耶穌基督為主，並且自己因耶穌作你們的僕人。 \end{tabularx} \\ \\ \relax
4:6 & \begin{tabularx}{0.7\textwidth}{X} 那吩咐光從黑暗裡照出來的神已經照在我們心裡，使我們知道神榮耀的光顯在耶穌基督的臉上。 \end{tabularx} \\ \\ \relax
4:7 & \begin{tabularx}{0.7\textwidth}{X} 我們有這寶貝放在瓦器裡，為要顯明這莫大的能力是出於神，不是出於我們。 \end{tabularx} \\ \\
[1ex]
\hline
\hline
\end{longtable}
$^{1}$丁師傅早安.
我們今天會看《哥林多後書》第四章.
《哥林多後書》是一段保羅所寫的書信.
一封其實充滿著很多掙扎的書信.
保羅在當中面對著哥林多教會很多的指控.
每一句說話都是他自己表白.
說出自己的掙扎.
所以我們可以這樣留意.
在這段書上我們每次看到「我們」字頭.
都可以用「我」字代入.
保羅自己的獨白和心聲.
這是第一點.
第二點我要補充的是.
當你打開《哥林多後書》第三四章的時候.
其中一個很重要的字眼就是「榮耀」這個字.
今天沒有時間跟大家看《長勝經》.
但如果我們大概去看第三四章的話.
你會發覺書中充滿著「光」,「榮耀」這個字.
你知道書中「榮耀」這個字.
「Dark Sir」這個字其實和光有很大關係.
所以你會發覺書中經常看到「榮光」.
或者「光輝」這樣的字眼.
所以今天我們不詳細說這裡.
不過你會發覺「光」這個字.
在經文中不斷出現.
有這樣的背景我就明白今天的經文.
經文說.
我們原不是傳自己,乃是傳基督耶穌為主,並且自己因耶穌作你們的僕人.
那分輻光從黑暗裡出來的神,已經照在我們心裡.
叫我們得知上帝榮耀的光,顯在耶穌基督的面上.
今天特別要處理第六節.
第六節是有點難明白的.
如果你去看第六節的時候.
你會發覺這是一個很難翻出來的經文.
當中中文聖經的譯本差別之大.
是令你大吃一驚的.
不同的譯本有不同的翻譯.
例如:那分輻光從黑暗裡出來的神,已經照在我們心裡.
使我們得知上帝榮耀的光,顯在耶穌基督的面上.
這個也比較難明白.

$^{41}$新學譯本的翻譯.
那位說光要從黑暗裡出來的神,已經照在我們心裡.
使我們因認識耶穌基督的面上那神的榮光而得到亮光.
是有一點不同的意思.
新學譯本的翻譯.
那說要有光從黑暗裡出來的神,已經照在我們心裡.
要我們把神的榮光照出去.
就是使人可以認識那在基督面上的榮光.
完全不同意思.
是照給別人而不是照給我們自己.
另一個譯本,女人中的譯本.
那分輻光當從黑暗裡出來的神,照徹了我們心裡.
使我們在耶穌基督面上,接受了上帝的榮耀之光.
是很難解釋的.
不過這是最接近原本的解釋.
我們問到底發生什麼事呢?.
當上帝的光照我們的時候,有什麼後果呢?.
你會發覺,剛才的翻譯有不同的答案.
第一個就是,當你被光照的時候.
你就會知道耶穌基督面上有榮光.
這是我們看到法國的那一頁.
新的譯本是這樣說的.
當你被光照之後,你就會因為認識耶穌面上的榮光得到亮光.
是有一點不同的原理.
你看到耶穌的臉上就有光.
新的譯本就說,當你被光照之後.
你就叫其他人認識基督耶穌的榮光.
是完全不同的意思.
兩種是最難明白的.
當你被光照之後,你就會接受了上帝的榮耀之光.
這是最接近原文的翻譯.
翻譯裡面主要是沒有動詞.
原文的意思是.
Enlightenment of knowing of glory in the face of Jesus Christ.
大概是你能夠得到啟蒙,能夠知道耶穌面上的榮耀.
所以無論如何,重點在於知道.
當你被光照之後,你能夠知道一些東西.
能夠被啟蒙一些東西.
不過,究竟是哪些人知道呢?.
是你知道,還是其他人知道呢?.

$^{81}$其實其他人知道是比較合理的.
當你被光照之後,你就成為了別人的光.
也能夠讓人認識耶穌的光輝.
這也是聖經其他意思.
光照在你面前,讓人認識你的開明遺物.
不過,原文裡面沒有其他人的字眼.
所以要加上字眼才能夠給予合理的解釋.
這可能就是說,不是其他人知道,而是你自己知道.
當你被上帝的光光照之後.
你就會知道基督耶穌的光光照你.
所以有點無聊,說完這句話,我也不知道該說什麼.
所以是有點不容易理解的.
我的看法是,兩個角色同時都是同一個意思.
剛才唱歌的時候,我們會唱一首歌.
《真光普照》,英文詩歌的名字叫《Shine, Just Shine》.
這首是我以前在教會裡經常唱的詩歌.
說的是什麼呢?.
耶穌基督是世界的真光,我們被光照.
Shine, Just Shine, Shine on me.
基督耶穌的真光光照了我們每一個人.
當我們被耶穌的光光照的時候.
才會給我們一個很美麗的稱呼.
叫做光明之子.
這是一個非常美麗的稱呼.
這也是聖經說的.
你們當趁著有光,順從者光,使你們成為光明之子.
我們是光明之子.
耶穌說你們是世上的光,你們是光明之子.
我剛剛信主的時候,覺得這個字是挺有型的.
我覺得有點中二病.
有點像光明之子,暗如光劍,光明騎士,LX那些.
我覺得是一個比較漂亮的名字.
不過很吊詭,很諷刺,很陰默的事情是什麼呢?.
跟雞尾包裡面沒有雞尾,中環廣場不是中環一樣.
光明之子其實一點都不光明.
其實他不懂得發光,光明之子.
唯有耶穌才是世上的真光.
所以你們知道,我們基督徒不是真正的光.
我們只不過是什麼?.
我們只不過是光輝的反照.

$^{121}$耶穌基督是世上的真光.
我們是這個真光首先照亮我們.
我們成為基督徒,然後我們就用這個真光照出來.
就好像玻璃一樣.
玻璃本身是沒有光的.
但它很多時候比其他光源更加發亮,更加光輝.
近代基本上每一個建築大師都喜歡用玻璃來做建築材料.
最著名是什麼?就是貝勒銘.
1989年羅浮宮,大家可能都去過.
一個非常漂亮的玻璃金字塔.
為什麼人們會越來越喜歡用玻璃來做建築材料呢?.
第一件事就是因為工程的進步.
現在就不需要用外牆來承擔功能.
所以用一些比較薄的東西都能夠做到.
第二就是玻璃比水泥更加輕.
減少整個建築的重量.
最重要的是什麼?當然就是玻璃有採光的作用.
它能夠讓外面的陽光照進來,令到有很多天然的光線.
所以頂尖門我們是一塊玻璃.
我們世上的光其實是一個玻璃的原理.
不是因為有什麼光在裡面.
而是我們是讓裡面的真光照出來.
所以這個道理正正是耶穌所說的.
耶穌說當你的眼睛光明,全身就光明了.
當你的眼睛昏花,全身就黑暗了.
所謂的明亮,光明其實是一種透徹,簡單,通透的意思.
就是說如果你的人是通透的話,你就全身就光明了.
如果你的眼睛是封閉的話,黑暗就會出現.
所以重點不是你有多大的發亮,發光.
而是你的人有多大的透徹,能夠讓裡面的光輝照出來.
所以回到保羅經文,我就明白了保羅的意思.
保羅說那分副光從黑暗裡照出來的神,已經照在我們心裡.
叫我們得知上帝榮耀的光,顯在耶穌基督的面上.
保羅說當你被上帝的光光照,你要好好地做一件事.
你要知道,你要好好地記住,你要搞清楚.
這個光其實是耶穌基督那裡而來的.
不是你身上而來的.
你唯一要做的就是要知道.
你能夠讓其他人真真正正看到這個光輝是來自耶穌基督自己.
所以就說One of both.

$^{161}$你要知道耶穌基督的光輝不是你的光輝的時候.
你才能夠讓人在你身上看到耶穌基督的光輝.
所以保羅說第五節說什麼.
我們原不是傳自己,乃是傳耶穌基督為主.
當顛覆的光照耀你的時候,這個光就是耶穌基督.
然後你就成為了光明之子.
但是作為光明之子,你只需要做一件事情.
你要深深地知道,認識明白,這個光是來自耶穌基督自己.
你作為光明之子,是這樣的來去治癒其他人.
所以整件事就像周星馳所說的,太陽能電筒一樣.
唯有耶穌基督是光輝光照的時候你才會亮.
如果沒有光,絕對不會亮.
唯有耶穌基督光照遍你的時候你才會有光輝在裡面.
所以我們所謂的作鹽作光,其實你沒有辦法作鹽作光.
你只能夠是鹽,你只能夠是上面的光.
你沒有辦法作鹽,你不能夠裝鹹.
你只能夠是鹽,或者不是鹽.
光也是一樣,你不能作光,你明明不是光.
你不能裝光,你走到人身邊去裝光,這是什麼?.
這是不好的事,你作不了光,因為你沒有光輝在裡面.
你唯一可以做的就是讓耶穌基督的真光從你身上照出來.
所以我們作為光明之子,我們其實要好好地去做一塊玻璃.
你的光不在乎你多努力去亮麗,多努力去發光去耀眼.
而是去努力保持自己生命的通透.
保持一塊玻璃一樣的單純,透明,簡單.
讓你真真正正裡面的光輝,沒有添加地照出來.
這就是我們要認真去做的事情.
信仰裡面的第一步就是將你裡面的善良和價值.
真真正正建基於你裡面的耶穌基督去表明出來.
不是我們去特意去營造人造的光輝.
簡單去說一些形容.
正如耶穌說你們是世上的光,你們的光也當借盡人前.
叫他們看見你們的行為,變相榮而歸給正聖的父.
我們很多時候以為我們基督徒要做的事情.
好像要不斷逼出好行為,不斷去發揮好行為.
明明本身沒有的,都要用法術去擠出來.
我們經常說基督徒在別人面前做好見證.
這句話其實有點奇怪,什麼叫做好見證?.
這句話其實有點像做好見證,有點像開工.
明明沒有的,就馬上做出來給別人看.

$^{201}$有點像打卡做單,做好見證出來.
不知道大家有沒有去過四五點鐘下場的餐廳吃飯.
那間餐廳基本上沒有開冷氣,只會關燈.
那些員工就坐在座位上看手機,睡覺.
當你走到餐廳裡面,突然間就坐在那裡,沒怎麼開工.
突然間那本書就出來,做好見證出來了.
所以整件事有點像明明沒有,但要強行做出來的感覺.
做好見證其實不是一件我們要特意做出來給別人看.
或者是你特意要擠出來的東西.
耶穌要求我們要好行為,但是一個真真正正.
出於你內心做出來的好行為.
正如我們所說,在幾百年前,楊惠詩莉的故事.
楊惠詩莉當時在牛津大學讀書.
她當時有一個club,叫做Holy Club.
Holy Club是一個很特別的地方.
首先你讀牛津,是這樣的Club.
每個星期大家一起聚會,一群牛津大學的大學生聚在一起.
做事前有一個很嚴重的問題,彼此問對方.
問題就是說,在過去一個星期裡面.
你有沒有有意無意地讓人覺得你比真正的你更好.
我再說一次,你有沒有有意無意地讓人以為你比真正的你更好.
這是一個很嚴重的問題.
我們有沒有有意無意做出來的更好的自己呢?.
我經常問我們怎樣來做基督徒.
基本上我們這麼多星期回來教會,基本上都是很友善的.
大家都是很好的人.
不知道在教會裡面的你比較好,還是對著媽媽的你比較好呢?.
在家裡的你比較好,還是在教會裡面的你比較好呢?.
很多時候都覺得做基督徒總是有善笑容.
我們做傳道人也是.
有一個詞語叫做會有笑.
看到頂字背後的笑容是特別好看的.
我覺得都是想衷心地真誠地對待別人.
怎樣可以真心地用真真正正的耶穌基督給我的善良和光輝去變出來.
學習去嘗試去表裡如一.
我想做基督徒很簡單的.
大概有一個樣子是可以學的.
回來之後有些是友善的,有些是友善的.
大概是可以學的.
但真真正正難的是什麼呢?.

$^{241}$是真真正正讓耶穌基督改變你的生命.
表裡如一的做出來.
這是一個很漫長的旅程.
我們可以很簡單地去學做基督徒.
但真真正正成為基督徒.
成為基督耶穌的真光給你的光輝.
是一個很漫長的過程.
嘗試去學習做一個我稱之為本真的人.
一個真正的人.
一個真真正正的你.
不要去學做基督徒.
而是真真正正成為一個光明之子.
不是叫你隨便發脾氣.
把最差的表現出來.
而是真真正正的讓你裡面去改變.
我們不是花功夫去學做哪一層.
而是花時間在禱告裡面.
不斷修練自己裡面的部分.
這個反而是難的.
反而是需要我們很漫長的時間.
當你習慣了這樣的時候.
當你習慣了讓你的生命在裡面發放出來的時候.
當人處事 做人 做基督徒.
都是從裡面基督耶穌的真光發出來的時候.
當你不斷地做的時候.
你就會發現裡面有一個很驚人的東西.
這個東西將會帶領你生命很多很多不同的階段.
你就會明白保羅這句話.
第七節說什麼.
保羅說:我們有這寶貝放在瓦器裡.
要顯明這莫大能力是出於神.
不是出於我們.
在我們裡面的正正是一個寶貝在瓦器裡.
最好的東西在我們生命裡面.
我們正正是將這個最好的東西.
去簡單地表現出來.
當你這樣做的時候.
你就會發現原來我們可以很坦蕩蕩地.
將裡面耶穌的美好變出來.
學習去做一個這樣的基督徒.

$^{281}$我想當我們習慣了嘗試去表如一的時候.
這個改變是更加長遠的.
更加讓我們可以走得更加遠.
我們教會裡面也是.
每個星期回來的時候.
嘗試去保持真誠.
回到家的時候.
保持我們真正的善良.
在裡面嘗試去表如一.
我們一起祈禱.
主你求求你幫助我們.
我們能夠在你的光輝裡面.
真正被你照遍.
將我們黑暗.
將我們當中很多生命裡面的短處.
慢慢被你修補.
讓我們去學習.
不是去遮蓋或掩飾我們的短處.
而是我們在裡面嘗試去做一個真真正正的基督徒.
讓你的真光照遍我們的生命.
讓我們每分每刻都能夠表如一來對待人.
讓我們真真正正成為合你使用的器皿.
能夠承載你寶貝的雅氣.
這個寶貝的雅氣.
因為好的緣故.
不是因為我們自己有多好.
而是因為你裡面的生命有多好.
求主你照遍我們.
讓我們在堂頂之內都能夠成為一個美麗的教會.
彼此用真誠.
彼此用真正的光輝.
彼此去相愛.
求你幫助我們.
奉主名求.
阿們.
\newpage



\section{創世記 39:2}
\label{sec:rX2pkPqwylA}
\textbf{2024.12.15 《活出神的同在》 創世記39\_2 (和修版) 盧思源姑娘}
\newline
\newline
連結: \href{https://youtube.com/watch?v=rX2pkPqwylA}{\texttt{https://youtube.com/watch?v=rX2pkPqwylA}} ~~~~ 語音日期: 2024-12-17
\newline
\newline
\hyperref[sec:3EdQX7vZPlY]{\small{< < < PREV SERMON < < <}}
~
\hyperref[sec:index_chronic]{\small{[返順時目]}}
~
\hyperref[sec:index_scriptual]{\small{[返順卷目]}}
~
\hyperref[sec:j8cE70aJucM]{\small{> > > NEXT SERMON > > >}}
\newline
\newline
創世記 39:2
\newline
\begin{longtable}{cl}
\hline
\hline
章節 & 經文 (和合本修訂版)\\
\hline
39:2 & \begin{tabularx}{0.7\textwidth}{X} 約瑟在他埃及主人的家中，耶和華與他同在，他是一個通達的人。 \end{tabularx} \\ \\
[1ex]
\hline
\hline
\end{longtable}
$^{1}$弟兄姊妹平安.
很感恩今天第一次來到紅引堂當中.
跟大家分享上帝的話語.
願意我們開始之前我們都有個祈禱.
我們回到你的殿當中去敬拜.
一同去聽你的話語.
一失間孩子要去分享你自己的信息的時候.
求真理的聖靈與我們無論說的.
我們聽的同在.
願意我們進入你自己的真理當中.
進入主你自己的話語當中.
我們被主你自己的話語去得著.
我們恭敬的將以下的時間去交賞.
奉主耶穌基督明治祈求.
阿們.
弟兄姊妹.
不知道大家有沒有看過.
Star Wars星球大戰.
這套戲已經很多年了.
其中有一句對白.
May the force be with you.
中文就是願原力與你同在.
在當中在星球大戰中.
一個主角絕對沒有使用你來祝福對方.
很多時候在道別的時間.
說這句話就好像取代了說再見.
而這句話如果在今時今日.
對一些星戰迷.
很多時候就像祝福的句子.
通常在離別的時候.
祝福對方好運的意思.
希望他們將來如果面對什麼困難.
好像絕地無事一樣.
都有一些原力與他們同在.
如果大家有留意.
如果是星戰看開的話.
因為May the force.
與May the 4th很相似.
所以如果在5月4日.
甚至是星球大戰日.

$^{41}$即是大家要來紀念星球大戰.
但為什麼會說這句話呢.
這句話與基督徒有句話很相似.
May God be with you.
願神與你同在.
有時候我們說這句話的時候.
可能都是想與對方祝福.
希望對方可以有神同在.
希望可以有神的保守保佑我們.
其實都幾像我們中國人的民間信仰.
有時候求神拜佛.
我們都希望保佑我們順順利利.
平平安安.
這也沒錯.
因為我們知道在聖經當中.
我們的神是一位保守我們的神.
也是一位賜下平安給我們的神.
但我們的神是否只是這樣呢.
如果只是這樣的時候.
與其他宗教有什麼分別呢.
又或者當我們遇到一些不順利的時候.
是否代表神不可信呢.
難道祂沒有與我們同在嗎.
今天我們看的經文.
雖然剛才我只用了一字的經文.
但今天我們也會看很多的聖經.
也會看39章和其他經文.
所以鼓勵大家如果有聖經的話.
也可以翻開.
雖然我Power Point也有些聖經.
也有經文.
今天這段經文提到約瑟.
約瑟大家可能記得的話.
他是雅各的兒子.
他與拉傑在年老時所生.
如果聖經提到約瑟的時候.
約瑟特別受到雅各的寵愛.
而且為約瑟做了一件彩衣.
就是因為這件彩衣.
令約瑟的十個哥哥妒忌他.

$^{81}$而且約瑟也曾經做夢.
哥哥要向他下拜.
所以令約瑟的哥哥.
也不太喜歡約瑟.
而如果我們剛才說的39章.
前面的經文37章.
裡面提到.
一天雅各叫約瑟去看他哥哥.
去放羊.
誰知道那天他哥哥就動了殺機.
將約瑟丟到坑裡.
將他賣給以實瑪利亞.
之後就帶到埃及.
這個就是剛才我們說.
看第39章當中的背景.
約瑟帶到埃及的時候.
他被法老的護衛長波提弗.
買了作為綠布.
而剛才我們看約瑟.
在39章裡面.
今天我沒有把所有的句子都說出來.
但如果我們看39章.
原來有四次的經文提到.
耶和華與約瑟同在.
分別就是剛才我們讀的第二節.
和第三節當中.
那裡提到約瑟在他埃及主人的家中.
耶和華與他同在.
他是一個通達的.
他主人見耶和華與他同在.
又見耶和華使他手裡所辦的事都順利.
這裡是頭兩次.
而之後如果我們看後一點.
21,23節同樣都記載到約瑟在監獄裡.
上帝與他同在.
如果我們看新譯本.
剛才我們讀那段經文的第二節.
原來新譯本那裡這樣說.
耶和華與約瑟同在.
約瑟就事事順利.

$^{121}$住在他主人埃及人的家裡.
他的主人見耶和華與他同在.
又見耶和華使他手裡所作的盡都順利.
在這裡我們見到.
約瑟凡事都順利.
如果我們看英文.
其實這個字是叫做successful.
或者succeed.
即是做什麼都很成功.
大家可以想像一下.
如果我們認識一個人.
如果他讀書的話.
讀書的科目都很厲害.
或者如果他出來工作.
做什麼都很順利.
上班又升職又加薪.
工作都很醒目.
頭頭是道.
你會怎樣看待這個人.
你可能都會對他.
都很與別不同.
好像都有點特別.
沒錯.
今天我們在看聖經裡面.
約瑟都是很與別不同.
令波提弗對他有另眼相看.
而且聖經裡面提到.
波提弗都見到約瑟的神與他同在.
大家要知道.
波提弗是法老的內神.
他不是信耶和華.
但是當他見到約瑟所做的時候.
他見到神與約瑟同在.
所以如果我們剛才說.
讀聖經39章.
我們見到約瑟之後.
原本約瑟是一個奴僕.
把他的身份升為管家.
讓約瑟管理家裡一切大小事.
聖經裡面提到.

$^{161}$除了他自己吃的東西.
波提弗是什麼都不知道.
如果大家可以想像.
波提弗對約瑟非常信任.
什麼都交給他.
第五節提到.
耶和華因約瑟的緣故賜福.
給埃及人的家.
凡家女和田間一切所做的.
都蒙耶和華賜福.
我們見到.
因著神與約瑟同在.
約瑟事事順利.
波提弗作為一個愛邦人.
他也見到神與約瑟同在.
甚至因著約瑟的關係.
連同波提弗一家也蒙神賜福.
弟兄姊妹.
當我們今天說.
神與我們同在的時候.
不只是單單一種感覺.
或單單一個願望.
我們很想上帝我們同在.
我們在聖經裡看到.
在約瑟身上.
原來可以活出神的同在.
原來可以被人看到.
我想一個見證.
昔日我在神學院讀神學的時候.
有一位神學院的老師.
他是一位宣教士.
他去了一個地方住了兩年.
但那個地方是一個創啟的地方.
即是不可以公開地傳講福音.
傳講耶穌.
在那兩年裡.
他沒有特別跟周圍的人說.
我是信耶穌的.
他沒有特別去說.
不過在兩年的生活裡.

$^{201}$他一直關心周圍的靈社.
關心周圍的人.
在他身上看到一些很特別的特質.
有一天他問.
有一個人住了兩年.
有一天有一個靈社轉頭問他.
跟你一起生活了兩年.
你好像有點特別.
你相信甚麼.
覺得很不同.
最後這位宣教士.
他可以跟這個靈社.
因為不是他自己去傳講.
純粹分享自己的信仰.
分享耶穌在他生命裡的經歷.
給這個靈社聽.
弟兄姊妹.
今天在我們的生命裡.
有沒有讓人看到神的同在.
特別如果我們身邊.
有些未信主的人.
他們藉著我們見到神.
昔日波提佛見到神和藥室同在.
不單單是藥室做事時事順利.
聖經描述.
除了吃東西.
他甚麼都不知道.
剛才第23節提到.
當藥室收監的時候.
監獄長把所有東西交給藥室.
監獄長一概不察.
要怎樣才可以不理會很多事.
一定要對那個人非常信任.
因為首先.
就算你多聰明多聰明.
如果信不過.
其實你也不會把所有東西交給他.
回到我們今天的時候.
弟兄姊妹.
今天我們身上.

$^{241}$我們有沒有一些特質.
讓人看到上帝和我們同在.
舉個例子.
剛才說的宣教士.
我們有沒有關心身邊的人.
對人有沒有憐憫的心.
還是當我們跟別人說.
我們是基督徒的時候.
大家會反應.
你是基督徒嗎.
你不像.
還是怎樣.
盼望我們的生命.
好像藥室一樣.
我們可以活出神的同在.
讓人看得見.
讓人藉著我們認識神.
接著.
如果我們繼續回到聖經的時候.
剛才說.
神的同在.
是否只停留在一切很順利.
當我每一次讀這段經文.
讀藥室這段經文的時候.
剛才說到這段有四次提到神的同在.
藥室做事都很順利.
但如果我們去聖經.
看中間七至二十節的時候.
我們看到藥室.
他被波提弗的老婆冤枉.
之後被人收監.
如果我們再看第四十章.
我們今天沒有讀.
當中提到.
當藥室被收監的時候.
有個酒精和善長為他們解夢.
善長在夢裡說.
善長會被處決.
酒精會復職.
最後真的這樣發生.

$^{281}$藥室就叫酒精.
如果真的當他復職的時候.
叫他和法老說好話.
希望可以放出來.
因為我們知道藥室沒有做過什麼錯事.
但如果我們一路看的時候.
我們看到當酒精復職的時候.
他忘記了藥室.
是兩年後才記得.
如果我們以一個普通人去想.
會想.
這樣都算順利?.
擺明一點都不順利.
又被人冤枉.
之後還要幫了人.
人家又不記得.
但我們看到聖經裡說.
藥室雖然遇到這些處境.
但神仍然與他同在.
讓我們去思想.
究竟神的同在.
是否真的在於凡事順利?.
當遇到逆境的時候.
神是否缺席?.
有時可能我們作為基督徒.
我們很難明白.
為什麼神容許一些逆境.
或者一些苦難.
去發生在我們身上.
是不是上帝故意.
弄出來要我們去承受?.
但我們在聖經裡知道.
並不是這樣.
聖經讓我們知道.
我們所信的神.
是一位聖潔美善的神.
在祂裡毫無黑暗.
祂賜給我們兒女是最好的.
我們有時會覺得.
神為什麼會這樣?.

$^{321}$我們要知道.
我們的神不是變態的神.
祂不是刻意糟蹋我們.
然後要我們去聽祂.
神不是這樣的神.
聖經裡有很多經文讓我們知道.
我們信主的人.
神會保護我們.
保守我們.
但不代表我們不會經歷苦難.
不代表我們不會經歷逆境.
作為基督徒.
不代表我們一定會一帆風順.
上帝沒有允許天色常難.
不過祂允許.
祂在任何時間.
任何處境.
祂和我們同在.
所以在耶穌在世.
在《約翰福音十六章》裡也說.
在世上你們有苦難.
但你們放心.
我已經勝過世界.
所以逆境苦難是有的.
但回到今天的經文.
我們看到約瑟所經歷.
雖然他在逆境中.
在困苦中.
但他仍然可以活出神的同在.
在過去三個星期前.
在錦宣家有一個報道會.
分享的見證是我們的弟兄.
Ken徐正寧弟兄.
我想有部分弟兄姊妹可能都認識他.
他是我們教會的會友.
在十多年前.
他32歲很年輕的時候中風.
大家可以想像.
如果發生在自己身上.
這麼年輕的時候中風.

$^{361}$其實真的不容易去面對.
不過在當晚的報道會裡.
記得Ken分享.
當他中風的時候.
那時候他還沒有真正認識神.
不過因著這件事的發生.
他有機會接觸福音.
有機會認識耶穌.
之後還回教會成為上帝的兒女.
縱然今天其實如果大家看到他.
他的行動仍然有不便.
信了耶穌並沒有令到他身體.
回復未中風之前的狀態.
但在他的分享裡.
他說他回望.
他看到神在他生命裡的恩典.
對於Ken來說.
環境沒有改變.
因為他仍然要面對行動的不便.
面對生活上的一些困難.
不過我們看到.
在他生命裡.
神與他同在.
因為上帝透過他昔日的經歷.
今天使用他成為別人的見證.
成為別人的祝福.
弟兄姊妹說.
今天我們可能會面對逆境.
或者將來我們有機會面對一些.
我們未知或未必所想發生的事情.
有時我們可能不明白.
但我們要確信.
無論在任何處境中.
上帝仍然與我們同在.
神的同在並不是我們所想.
快點把困難拿走.
把迫害拿走.
而是神的同在是在困難裡.
上帝應許他對我們不離不棄.
他應許和我們一起.

$^{401}$他應許在困難中.
我們可以靠著主得勝.
接著回到經脈.
我們繼續看今天39章.
第六節當中提到.
約瑟其實很帥.
聖經形容他容貌盡美.
而波提弗的太太.
見到他時經常盯著他.
甚至想和他一起上床.
而約瑟的反應是怎樣.
約瑟的反應.
在第八第九節說.
約瑟拒絕.
他對主人的妻子說.
看啊 一切家務我主人一概不知.
他把所有的都交在我手裡.
在這家裡沒有人比我更大.
除你意外.
他也沒有留下一仰不及.
交給我.
因為你是他的妻子.
我怎能恨這麼大的惡得罪神呢.
在這裡聖經說.
雖然波提弗的太太.
經常叫約瑟和他上床.
但約瑟不聽.
避開他.
不和他一起.
大家知不知道約瑟當時是甚麼年紀.
聖經提到他被賣的時候是十七歲.
大家想像一下.
一個十幾二十歲的男子.
一個血氣方剛的男子.
有一個女人.
天天都和他說.
誘惑他.
其實也挺不容易抵擋的.
特別如果我們回到現在這個世代.
很多人都會覺得.

$^{441}$哇 這麼近.
有人送到你身邊.
你不是不接受吧.
可能現在的處境.
有些世代是這樣的觀念.
不過我們看到約瑟的反應.
就是他堅持不會做這些事.
甚至聖經裡提到.
有一天整個家裡都沒有人.
這個女人再次誘惑他.
但約瑟說他知道.
他不可以做這件事去得罪神.
雖然沒有人在.
但神知道.
約瑟知道神和他同在.
所以他不會做一些得罪神的事.
回到我們今天.
有時我們有沒有意識到.
其實神和我們同在呢.
記得很久之前.
曾經和一個弟兄姊妹去聊天.
不是我們教會.
他有記憶圖.
不過聊天裡提到.
有時不想神和我們同在.
為什麼呢.
因為如果我做了不好的事.
其實我也不想神知道.
也不想被看到.
之前我有個朋友.
他在社交平台上.
發了幾歲的女兒.
去到外婆家裡.
忽然不見了.
她躲在房間裡偷偷地吃巧克力.
吃到滿嘴都是.
然後她媽媽就拍下了.
我想這個小朋友.
平時應該沒什麼機會吃巧克力.
所以去到外婆家裡.

$^{481}$看到沒有人.
沒有人又有巧克力.
就偷偷地躲在房間裡吃.
以為沒有人發現.
誰知道被他媽媽發現.
還拍了照片.
放在臉書上.
糟糕了.
回到我們的信仰.
剛才我們一方面.
我們很想上帝和我們同在.
很想神保守我們.
事事都順利.
但會不會有時我們做了一些事.
其實我們不想上帝知道.
特別如果那些是神不喜悅的時候.
但我們知道.
神是無所不在的.
我們有沒有像弱勢一樣.
對神存有那份敬畏的心呢.
因為我們知道.
祂常與我們同在的時候.
以致當我們面對世界的誘惑.
或者引誘的時候.
我們記得神在這裡.
以致我們可以靠著.
神給我們的能力.
去勝過面前那份試探.
或者那份引誘.
接著下去.
剛才說弱勢面對主母引誘.
祂抵擋了.
不過換來的結果.
剛才說被人誣陷.
明明女人自己拉開衣服.
接著說弱勢偷了她的衣服.
然後誣告弱勢想跟她睡.
還要留下衣服成為證物.
你想想其實水洗都不清.
這件事都是一宗冤案.

$^{521}$但我們最後看到.
弱勢因為這樣.
雖然被波提弗收監.
但聖經提到.
弱勢沒有特別為自己辯護.
我們看到弱勢.
雖然最後被收監的時候.
面對這麼大的冤屈.
在人的心裡.
我們大可以埋怨.
為什麼會這樣.
神你又說與我同在.
你又說和我一起.
但我這麼大的冤屈.
被人委屈.
但聖經沒有提到她埋怨.
相反聖經提到.
耶和華與弱勢同在.
向她施恩.
使她在監獄長眼前蒙恩.
剛才說弱勢最後在監獄裡.
連屍肉都信任她.
把囚犯交給她管理.
而且她做的都盡順利.
弟兄姊妹.
在這裡我們要想一想.
究竟什麼是神的同在.
正如之前我們所問.
神的同在.
是否代表一切事情一定順利.
不過我們在弱勢身上.
我們看到神的同在.
不一定凡事順利.
相反雖然在逆境中.
但神的同在.
讓人可以跨過處境.
如果我們看弱勢整個人生.
由她哥哥被賣給20瑪利人.
成為埃及人的奴僕.
又被人冤枉坐牢.

$^{561}$在這些事上一切都不如人意.
但我們看到在弱勢身上.
神與她同在.
神在弱勢逆境中仍然成就上帝的美事.
如果我們繼續看創世記.
我覺得很好看.
我們知道最後她幫他解夢的酒精.
剛才說記得弱勢.
之後她在法老面前解夢.
弱勢被法老封為宰相.
以色列有機方.
她的哥哥要去埃及借糧.
重遇弱勢.
那一幕很好看.
當他們重遇那一幕在45章.
當時弱勢可以說出來.
叫她的哥哥不要因為賣她來埃及而自責.
因為在弱勢的心裡.
雖然她經歷了這麼多不順利的事.
但她仍然確信.
這一切是神為了保護她的性命而來到埃及.
她相信上帝有她的心意.
但當然我們不是說.
神是否要把弱勢賣來埃及.
不是這個意思.
而是神可以在人的錯誤或失敗當中.
成就她要成就的事.
所以當弱勢遇到這麼多事情.
她不會埋怨神.
因為在弱勢的心裡.
她相信神一直與她同在.
弟兄姊妹.
願神幫助我們.
兼顧我們的信心.
不要因著環境有時的不信.
或一些困難去打擊.
而令我們懷疑神的同在.
令我們懷疑神好像不愛我們.
盼望弱勢的見證.
告訴我們.

$^{601}$神無論任何時間.
無論在我們的逆境.
在我們的順境.
神都與我們同在.
因著神與我們同在的時候.
我們不單讓人在我們身上.
讓人見到神.
而且因著神與我們同在的時候.
我們可以勝過一些試探.
勝過引誘.
以致我們可以靠著主.
在無論任何環境當中.
靠著主跨過.
活出神的同在.
我們一同祈禱.
藉著昔日弱勢的經歷.
你讓我們見到.
上帝你與弱勢同在.
今天你也與我們同在.
求你幫助我們.
我們活出上帝的同在.
我們知道上帝的同在.
你不單單在順境當中.
而是當我們面對人生.
我們的高低起跌.
神也與我們同在.
也求你聖靈上上提醒我們.
因著神你與我們同在.
以致我們在行事為人中.
我們敬畏神.
我們真的當面對試探.
面對引誘的時候.
我們靠主去勝過誘惑.
以致我們活出上帝你所喜悅的樣式.
願意上帝你自己的話語.
成為我們心中的幫助提醒.
願意我們在你自己的愛.
在你的話語裡面.
繼續去成長 扎根.
天父多謝你.

$^{641}$聽我們的祈禱.
奉主耶穌基督明智祈求.
請自期待 阿們.
\newpage



\section{創世記 4:1-16}
\label{sec:j8cE70aJucM}
\textbf{2024.12.22 《愛子》 創世記4\_1-16 (和修版) 陳晟希傳道}
\newline
\newline
連結: \href{https://youtube.com/watch?v=j8cE70aJucM}{\texttt{https://youtube.com/watch?v=j8cE70aJucM}} ~~~~ 語音日期: 2024-12-23
\newline
\newline
\hyperref[sec:rX2pkPqwylA]{\small{< < < PREV SERMON < < <}}
~
\hyperref[sec:index_chronic]{\small{[返順時目]}}
~
\hyperref[sec:index_scriptual]{\small{[返順卷目]}}
~
\hyperref[sec:tR6kQF63JQE]{\small{> > > NEXT SERMON > > >}}
\newline
\newline
創世記 4:1-16
\newline
\begin{longtable}{cl}
\hline
\hline
章節 & 經文 (和合本修訂版)\\
\hline
4:1 & \begin{tabularx}{0.7\textwidth}{X} 那人和他妻子夏娃同房，夏娃就懷孕，生了該隱 ，她說：「我靠耶和華得了一個男的。」 \end{tabularx} \\ \\ \relax
4:2 & \begin{tabularx}{0.7\textwidth}{X} 她又生了該隱的弟弟亞伯。亞伯是牧羊的；該隱是耕地的。 \end{tabularx} \\ \\ \relax
4:3 & \begin{tabularx}{0.7\textwidth}{X} 過了一些日子，該隱拿地裡的出產為供物獻給耶和華； \end{tabularx} \\ \\ \relax
4:4 & \begin{tabularx}{0.7\textwidth}{X} 亞伯也把他羊群中頭生的和羊的脂肪獻上。耶和華看中了亞伯和他的供物， \end{tabularx} \\ \\ \relax
4:5 & \begin{tabularx}{0.7\textwidth}{X} 卻看不中該隱和他的供物。該隱就非常生氣，沉下臉來。 \end{tabularx} \\ \\ \relax
4:6 & \begin{tabularx}{0.7\textwidth}{X} 耶和華對該隱說：「你為甚麼生氣呢？你為甚麼沉下臉來呢？ \end{tabularx} \\ \\ \relax
4:7 & \begin{tabularx}{0.7\textwidth}{X} 你若做得對，豈不仰起頭來嗎？你若做得不對，罪就伏在門前。它想要控制你，你卻要制伏它。」 \end{tabularx} \\ \\ \relax
4:8 & \begin{tabularx}{0.7\textwidth}{X} 該隱與他弟弟亞伯說話 。 二人正在田間時，該隱起來攻擊他弟弟亞伯，把他殺了。 \end{tabularx} \\ \\ \relax
4:9 & \begin{tabularx}{0.7\textwidth}{X} 耶和華對該隱說：「你弟弟亞伯在哪裡？」他說：「我不知道！我豈是看守我弟弟的嗎？」 \end{tabularx} \\ \\ \relax
4:10 & \begin{tabularx}{0.7\textwidth}{X} 耶和華說：「你做了甚麼事呢？你弟弟血的聲音從地裡向我哀號。 \end{tabularx} \\ \\ \relax
4:11 & \begin{tabularx}{0.7\textwidth}{X} 現在你必從這地受詛咒，這地開了口，從你手裡接受你弟弟的血。 \end{tabularx} \\ \\ \relax
4:12 & \begin{tabularx}{0.7\textwidth}{X} 你耕種土地，它不再給你效力；你必流離飄蕩在地上。」 \end{tabularx} \\ \\ \relax
4:13 & \begin{tabularx}{0.7\textwidth}{X} 該隱對耶和華說：「我的懲罰太重，過於我所能承當的。 \end{tabularx} \\ \\ \relax
4:14 & \begin{tabularx}{0.7\textwidth}{X} 看哪，今日你趕我離開這塊土地，不能見你的面；我必流離飄蕩在地上，凡遇見我的必殺我。」 \end{tabularx} \\ \\ \relax
4:15 & \begin{tabularx}{0.7\textwidth}{X} 耶和華對他說：「既然如此 ，凡殺該隱的，必遭報七倍。」耶和華就給該隱立一個記號，免得人遇見他就殺他。 \end{tabularx} \\ \\ \relax
4:16 & \begin{tabularx}{0.7\textwidth}{X} 於是該隱離開了耶和華的面，去住在伊甸東邊挪得之地。 \end{tabularx} \\ \\
[1ex]
\hline
\hline
\end{longtable}
$^{1}$洪恩堂的弟兄姊妹大家好,早安.
我想大家第一次見我,其實我自己也是錦繡堂的弟兄.
我好像也是第一次來洪恩堂,已經十年以上了.
但我原來是第一次來,所以也很感恩.
還有好像也有些人是見過的,以前在錦宣那裡.
所以我們也是很有淵源.
還有謝謝周牧師邀請,給我機會過來服侍.
還有幾天就是聖誕節.
其實每一年在教會過聖誕,我不知道大家是第幾年了.
我自己也十幾年了,每一年我們都會去紀念耶穌.
對我們的意義.
不知道大家每一年對於聖誕節的感受是怎樣呢.
我想今年跟大家說一個跟聖誕無關的經文.
通常都會說一些聖誕經文.
不過我想應該之後聖誕崇拜的周牧師會再說.
我想跟大家由一個不開心的故事去想聖誕節.
雖然剛才說的那段經文好像跟聖誕沒什麼關係.
但其實也有些東西跟聖誕很相似.
大家一起聽聽吧.
我想哥仁和阿伯大家都認識吧,這兩個人.
有沒有人第一次聽這兩個人的名字.
我自己比較少聽傳道人說這兩個人的名字.
但其實這兩個人的名字是很有意義的.
但不是一些好的意義.
他們是人類歷史上第一宗謀殺案.
裡面的兩個主角.
也是亞當夏娃第一對子女.
聖經裡面說的.
先看看.
對了.
剛才經文有介紹到.
這兩兄弟都是農夫.
阿伯是牧羊.
哥仁是種地.
他們有一天拿他們自己養殖的東西.
獻給耶和華上帝.
獻祭大家不知道大家是否了解.
獻祭其實是一個感恩的舉動.
就是我送禮物給上帝.
聖誕節我們互相交換禮物.

$^{41}$或者你送禮物給誰.
你也收一下禮物.
獻祭就是人類想獻禮物給上帝.
多謝他.
哥仁種地.
種農作物.
阿伯養羊.
所以順理成章他們獻自己種的東西.
在經文裡面說.
哥仁獻地類的出產.
即是水果,菜等等.
而阿伯獻自己羊群裡頭生的.
經文特別說到他獻了頭生的羊和豬油.
頭生的意思是很棒的.
還有很特別的意義.
因為頭生在聖經裡面也是很重要的.
兩兄弟獻祭之後.
劇情下去就是.
上帝喜歡阿伯的公物.
喜歡羊.
然後對哥仁和他獻的東西就不接受了.
哥仁就很生氣.
因為其實我們不知道他們是怎樣互動的.
你有沒有試過獻一些東西給上帝不要的.
你奉獻了我一百塊.
然後突然彈回你的銀行.
不會這樣的吧.
其實我也不太想像到他們是怎樣互動的.
但總之哥仁就很收到.
上帝不喜歡我獻的東西.
就生氣了.
變了面色.
我不知道你會不會問一個問題.
我自己看的時候就問.
為什麼上帝不喜歡我獻的東西呢.
我讓你選擇吧.
如果你是上帝的話.
左邊那個就是阿伯獻的.
右邊那個就是哥仁獻的.
你是上帝.

$^{81}$你要選擇哪一個呢.
如果你問我.
是不是我都一定選擇左邊的.
有沒有人喜歡吃菜多過吃肉的.
有啊,我看到有人點頭.
那哥仁就很開心了.
如果你是上帝的話.
如果給我.
我猜應該很多人都不喜歡吃扒.
喜歡吃菜多過吃肉的人.
可能在少林寺會多些找到.
那我就想.
會不會上帝跟我們一樣呢.
都是食肉獸.
喜歡吃香味的東西.
吃扒.
很容易理解.
如果是這樣的話,那就是神偏食.
是不是呢.
其實可能哥仁也會這樣想.
我也會代入他的角色去想.
咦,上帝為什麼不喜歡我獻的菜呢.
是不是嫌我的菜太清淡呢.
難道他看不起我嗎.
可能他也會覺得.
有沒有搞錯,上帝偏心.
我跟我弟弟兩兄弟.
大家都是各盡其職.
為什麼你選他不選我.
實在太不公平了.
偏心.
如果你想想.
發生在一個現代的家.
有兩兄弟.
哥哥有這樣的反應是不是很正常.
是啊.
有什麼理由,哥哥跟弟弟兩個.
都是獻一些東西出來.
為什麼爸媽選弟弟那些不選我呢.
都很不開心.

$^{121}$如果你是哥哥.
就是很痛苦.
好像覺得自己是次等一點.
那.
所以如果這樣看.
似乎神是有些問題.
神偏心.
為什麼會有一個這麼不懂.
做父母的神呢.
究竟神是不是真的偏心呢.
我自己也聽過有人這樣說.
他覺得神就是會偏心的.
他是神.
他喜歡偏心誰就偏心誰.
你管他吧.
但我自己覺得應該有點誤會了.
我不覺得神會像人一樣.
會有這麼狹窄的心胸.
有什麼可以幫我們理解.
究竟發生什麼事呢.
其實在經文沒有直接說.
但有些線索我們可以試試去推理.
首先我們看看.
神怎麼回應該人呢.
他就說你為什麼發怒.
你為什麼要變了臉色.
你如果行得好.
你豈不蒙越立.
你如果行得不好.
罪就伏在門前.
神的回應.
他就說那個人似乎有些問題.
暗示了他本身已經有些不對勁.
以至神拒絕他的獻祭.
如果他有問題.
那問題是什麼呢.
第二條線索我們再看看.
究竟作者如何描寫.
他們兩兄弟獻的東西呢.
有些學者.

$^{161}$因為經文有些隱晦.
有些學者就覺得.
你看看.
該人就這樣寫.
該人是他自己地裡的出產.
而亞伯就是他投生和資油.
似乎有些對比.
亞伯獻就是他特意選擇最好的.
而該人就隨便吧.
不理當不當做.
總之湊夠數就將菜扔上去給神.
這樣.
所以換句話說.
似乎該人的動機和亞伯的動機.
那個心是有點不同的.
似乎他有些對比.
似乎該人的問題.
有些暗示.
他好像交行貨.
其實一般獻祭.
獻完之後.
都期望神會有些回應.
所以.
可能.
暗示了亞伯和該人的分別.
就是亞伯.
他真是很誠心誠意.
很想神都開心.
他就獻最好的給他.
但該人呢.
他似乎.
好像想.
做了形式就算了.
總之獻了東西.
交了貨就算了.
我不理了.
我做了我的部分.
祝福給我.
如果這樣看.
又好像說得通.

$^{201}$為什麼該人不開心呢.
可能他覺得神偏心.
但也可能有些失望.
有沒有搞錯.
我做了我要做的事.
我竟然換不到.
我想要的祝福.
有些挫敗.
其實也挺有意思.
我看短短的經文.
有點像.
人與人互動.
好像.
子女與父親.
或者下屬與上司.
有些日常的感覺.
我們再試試代入.
該人很生氣.
父親應該很開心.
神有什麼感受呢.
其實.
我自己看.
神是否刻意羞辱該人.
就是說.
我就是不選你.
讓你覺得沒面子.
其實神.
應該不是想這樣做.
雖然這個拒絕.
令該人覺得好像.
很不被接納.
但神說的話.
重點不是羞辱.
而是想讓他知道.
你有些問題.
你要注意.
而事實上.
如果神不重視該人.
其實他不太疼愛.
這個兒子.

$^{241}$他自己創造出來的人類.
他根本不用像媽媽.
在含恨他.
如果他走得不好.
又怎會這樣呢.
當一個人不被重視.
他就真的沒有人提醒他.
無論他做得怎樣.
不懂做人.
不懂禮貌.
都不關我的事.
我為什麼要提醒他呢.
反過來.
就是因為我很緊張他.
所以我才浪費時間.
去吻他.
想他知道自己.
要怎樣改善.
所以神很關心該人的態度.
出了問題.
他不只是看.
你冒險獻最好的菜給我.
該人可能就以為.
你看不起我的菜.
但其實神從來都不太.
介意獻祭的東西.
如果神真的.
那麼介意食物.
這個神其實都很膚淺.
所以這裡看到神和人.
有些分別.
該人將神看成一種.
祝福交易的對象.
可能我們今天.
要做很多事奉.
但為了得到讚賞.
和我做很多事奉.
我就可以肯定神一定會.
很喜悅我.
我做得多一點.

$^{281}$神就祝福我多一點.
但在這個經文看.
其實神完全不是這樣的.
神由始至終都沒有說過.
他很介意他們兩個獻什麼.
而是關心他們生命的狀況.
雖然.
反過來經文沒有說.
很多阿伯的事.
阿伯下一段再出現就死了.
但你可以推斷.
阿伯獻祭的心態.
其實帶著一種.
愛和感情在.
就好像你將禮物送給.
你愛的人.
你想選擇好的.
其實不是因為你覺得對方會不喜歡.
而是你真的很想.
將最好的東西給他.
你這樣做的時候你自己也開心.
所以他和神之間.
有真正的感情.
而反過來.
該人就.
有點像走肉朋友.
總之不是說.
完全是最真誠.
有利可圖.
他就走過去.
但沒有回報的時候.
就生氣失望了.
所以其實.
再說得準確一點.
神很看重的.
我相信是他和該人之間的感情.
不只是說.
該人這個人是否乖.
懂得獻祭.
而是他不想該人.

$^{321}$你只是來找我換祝福.
我其實.
想和你建立一段.
真正有愛的關係.
所以神一直.
其實在過程中想提醒該人.
其實我是.
愛你的.
我和你之間不用這麼功利.
不過該人可能.
他沒有這種心.
所以提醒.
姐妹其實我們也可以.
去想想其實究竟我們在.
和神的關係裡面.
我們是怎樣的呢.
我和神之間.
其實感情是如何的呢.
很可惜.
經文說.
該人被神提醒之後.
他沒有理會上帝.
他就很生氣.
神叫他控制自己的罪.
不要被罪帶來的衝動.
控制你.
所以其實.
我不知道有沒有遇過這樣的情況.
總之有些不如意的時候.
有時候可能會生神氣.
或者總之.
有些做得不好的地方.
神也說.
神沒有什麼特別的力量.
去令你不犯罪.
而是我會將選擇帶到你面前.
神提醒.
該人.
你可以選擇.
你永遠都有機會選擇.

$^{361}$但你看到神.
他將選擇給該人.
他要自己肯去選擇.
不去繼續沉淪下去.
大家最後很可惜.
該人選擇了做一些.
很不對的事情.
創了歷史第一宗殺人案.
如果說該人這個人.
其實他某程度上也很慘.
他很妒忌他弟弟.
妒忌也很辛苦.
你面對一種很強烈的情緒.
有時候真的.
很難控制.
如果你有試過很妒忌一個人.
可能你也會有些體會.
我們很熟悉的.
三國演義.
也有很多例子.
例如周瑜妒忌諸葛亮.
他最後妒忌到死了.
如果你有聽過這個故事.
因為妒忌.
令他情緒有問題.
身體很差.
這些情緒.
這些罪其實可以影響一個人.
很多事情.
他不開心.
身體很差.
命也短短十年.
雖然這段經文不是說.
我們如何處理妒忌.
但是你會看到.
其實神的角色是怎樣.
他不是追究.
該人說.
你搞什麼鬼.
你不應該這樣.

$^{401}$麻煩你自己好好處理.
其實我更加相信.
神在提醒他的時候.
其實神是預備好.
幫該人克服.
他內心的問題.
如果神是一個.
扔下包袱.
讓你自己去處理.
自己的嫉妒.
他也不用緊張.
他也不是.
該人的天賦.
至少我相信.
神在說的時候.
他已經預備好.
只要願意.
我會跟你一起面對你的問題.
並且我會讓你看到.
我拒絕你這次的憲制.
不是為了懲罰你.
而是我覺得.
你這一刻需要.
跟我一起面對你內心的問題.
如果該人.
他也願意.
去選擇.
去到神的面前.
你不接納我的憲制.
我是很不開心的.
但是你這樣說的時候.
我都發覺我態度可能有些問題.
對不起.
你可不可以幫助我去改變.
如果該人又這樣.
跟神祈禱.
我相信阿伯就不會死.
故事的發展也很不一樣.
經文其實我自己看到.
神作為一個天賦.

$^{441}$和一個牧者.
他的心腸.
其實他一直在對話中.
看到神很想幫助該人成長.
他也有罵他.
但罵他不是為了拼他.
其實是想幫他站起來.
從錯誤中悔改.
而即使該人.
他的態度很惡劣.
其實這個人也很惡.
但是神和他依然很近.
其實.
會跟他說話.
神和這個人疏遠.
你犯罪了.
我疏遠你.
該人是不會聽到神和他說話.
該人就算在這樣的狀態.
他也聽到神和他說話.
其實神.
特別在這個狀態.
在罪裡面.
掙扎很辛苦的時候.
神和他特別近.
因為他很需要聽到神和他說話.
而這種很近的關係.
在該人殺了他弟弟之後.
有些改變.
但依然沒有影響到神對該人的關心.
我們看下去就會看到.
該人殺了他弟弟就很嚴重.
經文說神也問他.
你做了什麼.
我聽到你兄弟的血.
在地裡出來和我哀告.
神.
其實那時候.
也知道.
回不了頭.

$^{481}$人也死了.
他就要宣判有些後果.
他說什麼呢.
都很嚴重.
第一就是他種地.
地不再幫他效力.
他會流離飄蕩.
另外就是他會受詛咒.
該人也知道自己很嚴重.
首先他種地不再幫他效力.
你知道他種地.
斷了他的路.
如果今天我解僱你.
你不用再指望入這一行.
這樣也很嚴重.
我做什麼好呢.
失業了.
難道我只靠綜援過日子.
還要流離飄蕩.
失業也算了.
現在連房子也沒有了.
不給你房屋住.
那怎麼辦.
所以該人也知道.
但他立刻說.
刑罰很重.
他還說.
他說.
你趕著我離開這裡.
我見不到你.
他和神的關係也破損了.
我流離飄蕩.
人們見到我一定會殺我.
其實補充一點.
因為那時候.
那年代.
有殺人填命的觀念.
總之你殺了一個人.
其他人就可以報仇.
類似這樣.

$^{521}$所以他也很嚴重.
你想想.
他們應該全部都是親戚.
那時候阿當夏娃.
全部都是他們生出來的.
所以可能他的姨媽姑姐.
表弟那些.
你殺了阿伯姐.
你殺了表哥.
全部都走來.
總之簡單來說.
很嚴重.
他整個人生.
基本上是受到破壞.
你甚至可以說.
他玩完也不過分.
一個人失業.
沒有房屋住.
被人殺.
還有什麼情況比這個更差呢.
所以你看到.
他做錯事之後.
那個嚴重性.
是真的去到這樣.
難怪上帝.
這樣去告誡他.
那個人.
他為了自己的利益.
於是乎.
他錯失了所有.
最重要的東西.
但是.
在神裡面.
究竟是不是這樣就等於玩完了呢.
今天其實.
我也很少見.
有些人犯了罪之後.
我們每個人都會犯不同的罪.
但是你有沒有聽過.
有這樣的例子呢?沒有的.

$^{561}$在現代來說.
是匪夷所思的,神很有恩典.
是不是舊約的神.
比較殘忍呢?.
其實也不是.
再看下去,發覺其實神.
依然繼續有給恩典的個印.
例如剛才.
剛才那個人就說.
那些人殺我,怎麼辦呢?.
於是乎神就立刻.
跟他說.
我給你一個記號,用來保護你.
總之殺你的.
就租報七倍.
就是說,我確保你.
不會被人殺,那些人動不了你.
那就給他一個記號.
不知道是什麼記號.
可能印在額頭上.
說不准殺我.
不知道,總之.
就保護了該人.
其實.
你想想,神是沒有義務.
開恩的,法律上.
你犯了法,那你要坐.
這麼多年牢,就要坐足他.
本來.
該人接受懲罰是公義的.
但是神仍然.
願意保護他.
該人我相信他那一刻.
就很後悔,但是.
他在這個後悔裡面,我也會想.
其實他有沒有發現.
神仍然都很愛他.
當你做了這麼嚴重的事.
其實.
我想起浪子的比喻.

$^{601}$那個浪子其實也不好意思.
跟他老爸說,老爸.
你像以前那樣,給我住漂亮的房子.
每天給我吃好東西.
他也覺得.
如果你肯收留我.
我已經偷笑了.
所以.
我自己猜的.
其實不要說.
其實可能對一個犯了這麼嚴重的錯的人來說.
他覺得.
我以後都不會再.
神不會再愛我.
我人生已經玩完了,你給我一條生路.
我以後都不再期待什麼.
但其實.
不是這樣.
神繼續都很愛他.
即使他犯了這麼重罪.
連他的姨媽姑姐.
都說要追殺他.
神仍然很重視這個人的生命.
他說我不想死.
神就說好,我保護你.
你一定不會死.
而且.
神是沒有停止供應恩典給戈仁.
經文就停在16節.
但其實16節之後.
經文繼續說下去.
戈仁也要離開那裡.
他也被人趕走了.
但他又找到老婆.
之後還生了兒子.
生了一個叫以樂的兒子.
還建了一座城.
你都不知道他.
為什麼可以這麼厲害.
被人趕走了之後.

$^{641}$自己找到地建房子.
可能是做了房地產.
我猜他都做得很棒.
他從人生是零.
去到之後.
他好像重新.
又再踏上另一個人生高峰.
所以.
我相信.
如果沒有神的保守和祝福.
他不可能這樣.
一個已經去到.
人生絕路的人.
這樣都可以翻身.
其實.
我們看到神一直都沒有放棄戈仁.
之後發生什麼事我們不知道.
但可能戈仁他都知錯了.
神就繼續賜福給他.
他就再站起來.
看到神不願意.
戈仁這個人.
他的人生就這樣毀了.
他仍然很希望戈仁的生命有改變.
其實.
整件事裡面.
我相信神自己都很受傷.
見到自己家裡.
都是他的兒子.
阿伯戈仁.
兩兄弟.
哥哥竟然殺死了弟弟.
對亞當夏威這個家庭.
都是很大打擊.
神一定都很傷心.
但就算很傷心.
神都沒有放棄過.
戈仁這個兒子.
只不過.
我不知道戈仁.

$^{681}$他看不看到這一切.
至少在這段經文裡面.
戈仁他不明白.
原來他一直都是神的愛子.
他不明白.
這個身份.
他是神所愛的.
所以.
當一個人不明白.
自己真的很被愛的時候.
他就會用很多手段.
想去肯定自己.
是可以得到祝福的.
但如果他明白的時候.
他就知道.
無論是什麼時候.
他只要肯回到天父那裡.
那些祝福永遠都是屬於他.
弟兄姊妹.
這段經文雖然和聖誕節沒有關係.
但其實我看到.
舊約這一位神.
和聖誕節降生的那位神.
神的愛子耶穌.
都是一樣.
他從來都沒有改變過.
同樣很相似.
其實在聖誕節.
你看回福音書.
其實那些人在耶穌.
降生的時候.
都不太知道祂來做什麼.
祂只知道這是彌賽亞.
還有知道祂會救.
以色列這個民族.
但祂究竟實際上想怎樣救他們.
其實那些人.
對耶穌有很多誤解.
就好像那位神.
其實不太體會到天父的心腸.

$^{721}$那些以色列人.
都不太明白耶穌的心腸.
很多人都拒絕耶穌.
不相信祂是神的兒子.
甚至有些門徒.
覺得耶穌應該是.
來搞一些軍事革命.
搞政變.
推翻當時的政府.
然後自己做皇帝.
後來那些人.
才慢慢開始明白.
其實耶穌是要來跟他們重建關係.
更新他們的生命.
我看到.
該人的上帝.
和新約降生的.
耶穌.
其實從來都是一樣.
祂來跟我們重建關係.
建立有感情有愛的關係.
並且去更新我們的生命.
弟兄姊妹.
我們盼望.
在聖誕節.
我們繼續.
去認識更多.
這位神.
我們的主對我們的心意.
特別.
不知道你會不會.
對該人有些共鳴.
如果有些弟兄姊妹.
覺得自己正在經歷神的拒絕.
管教.
或者你覺得很不如意.
神對你很不公平.
會有些憤怒.
該人的經歷告訴我們.
在最差的情況.

$^{761}$神的愛.
依然一樣.
不會改變.
而且你願意.
轉向他的時候.
他就會跟你一起面對.
神從來都不是要去懲罰.
不會去羞辱我們.
神永遠都很想建立我們.
盼望.
大家更加認識.
我們的主.
在聖誕節.
其實不一定只是聖誕節.
聖誕節是一個.
紀念的日子.
去認識神.
我們每天都會繼續.
最後不如我們一起去祈禱.
為大家去祈禱.
天父我們都.
獻上感恩.
因為你對我們.
忠誠.
是信實的上帝.
無論我們怎樣.
你對我們的心.
都是一樣.
就算是對最差勁的人.
你一樣.
都想去挽回.
想去建立我們.
盼望你幫助我們.
更加去體會到.
我們每一個都是.
所愛的兒女.
就好像.
你喜悅主耶穌.
我們都是一樣.
神.

$^{801}$我們這個時間.
都是紀念主.
降生的日子.
天父願你幫助我們.
我們在聖誕節裡面.
能夠去更新.
對你自己的認識.
幫助我們.
去明白.
你為我們所做的.
是怎樣的事.
幫助我們更加信任你.
幫助我們更加.
在你裡面有那份安穩.
願你叫我們.
聖誕節能夠.
為我們這個年度.
很好的總結.
讓我們展望將來.
更加叫我們.
不單只是聖誕節.
每一天我們都同樣去紀念你.
賜福紅恩堂的弟兄姊妹.
我們禱告教託.
奉主耶穌基督的名求.
阿們.
\newpage



\section{使徒行傳 26:1-23}
\label{sec:tR6kQF63JQE}
\textbf{2024.12.29 《沒有違背那從天上來的異象》 使徒行傳26\_1-23 (和修版) 講員 : 李富成牧師}
\newline
\newline
連結: \href{https://youtube.com/watch?v=tR6kQF63JQE}{\texttt{https://youtube.com/watch?v=tR6kQF63JQE}} ~~~~ 語音日期: 2024-12-29
\newline
\newline
\hyperref[sec:j8cE70aJucM]{\small{< < < PREV SERMON < < <}}
~
\hyperref[sec:index_chronic]{\small{[返順時目]}}
~
\hyperref[sec:index_scriptual]{\small{[返順卷目]}}
~
\hyperref[sec:0jw5FB1CQ6s]{\small{> > > NEXT SERMON > > >}}
\newline
\newline
使徒行傳 26:1-23
\newline
\begin{longtable}{cl}
\hline
\hline
章節 & 經文 (和合本修訂版)\\
\hline
26:1 & \begin{tabularx}{0.7\textwidth}{X} 亞基帕對保羅說：「准你為自己申訴。」於是保羅伸手辯護說： \end{tabularx} \\ \\ \relax
26:2 & \begin{tabularx}{0.7\textwidth}{X} 「亞基帕王啊，猶太人所控告我的一切事，今日得以在你面前辯護，實為萬幸。 \end{tabularx} \\ \\ \relax
26:3 & \begin{tabularx}{0.7\textwidth}{X} 更慶幸的是你熟悉猶太人的規矩和他們的爭論；所以，求你耐心聽我。 \end{tabularx} \\ \\ \relax
26:4 & \begin{tabularx}{0.7\textwidth}{X} 「我自幼為人如何，從起初在本國的同胞中，以及在耶路撒冷，所有的猶太人都知道。 \end{tabularx} \\ \\ \relax
26:5 & \begin{tabularx}{0.7\textwidth}{X} 他們若肯作見證，就知道我從起初是按著我們教中最嚴緊的教門作了法利賽人。 \end{tabularx} \\ \\ \relax
26:6 & \begin{tabularx}{0.7\textwidth}{X} 現在我站在這裡受審，是為了對神向我們祖宗的應許存著盼望。 \end{tabularx} \\ \\ \relax
26:7 & \begin{tabularx}{0.7\textwidth}{X} 這應許，我們十二個支派，晝夜切切地事奉神，都指望得著。王啊，我正是因這指望被猶太人控告。 \end{tabularx} \\ \\ \relax
26:8 & \begin{tabularx}{0.7\textwidth}{X} 神使死人復活，你們為甚麼判斷為不可信呢？ \end{tabularx} \\ \\ \relax
26:9 & \begin{tabularx}{0.7\textwidth}{X} 「從前我自己認為必須竭力反對拿撒勒人耶穌的名， \end{tabularx} \\ \\ \relax
26:10 & \begin{tabularx}{0.7\textwidth}{X} 我在耶路撒冷也曾這樣做過；我不但從祭司長得了權柄，把許多聖徒收在監裡，而且他們被殺，我也表示贊成。 \end{tabularx} \\ \\ \relax
26:11 & \begin{tabularx}{0.7\textwidth}{X} 在各會堂，我屢次用刑強迫他們說褻瀆的話，我非常厭惡他們，甚至追逼他們，直到外邦的城鎮。」 \end{tabularx} \\ \\ \relax
26:12 & \begin{tabularx}{0.7\textwidth}{X} 「那時，我帶著祭司長的權柄和命令往大馬士革去。 \end{tabularx} \\ \\ \relax
26:13 & \begin{tabularx}{0.7\textwidth}{X} 王啊！我在路上，中午的時候，看見從天上有一道光，比太陽還亮，四面照射著我和跟我同行的人。 \end{tabularx} \\ \\ \relax
26:14 & \begin{tabularx}{0.7\textwidth}{X} 我們都仆倒在地，我就聽見有聲音用希伯來話對我說：『掃羅！掃羅！你為甚麼迫害我？你用腳踢刺棒是自找苦吃的！』 \end{tabularx} \\ \\ \relax
26:15 & \begin{tabularx}{0.7\textwidth}{X} 我說：『主啊，你是誰？』主說：『我就是你所迫害的耶穌。 \end{tabularx} \\ \\ \relax
26:16 & \begin{tabularx}{0.7\textwidth}{X} 起來，站著，我向你顯現的目的是要派你作僕役，為你所看見我的事，和我將要指示你的事作見證人。 \end{tabularx} \\ \\ \relax
26:17 & \begin{tabularx}{0.7\textwidth}{X} 我也要救你脫離百姓和外邦人的手。我差你到他們那裡去， \end{tabularx} \\ \\ \relax
26:18 & \begin{tabularx}{0.7\textwidth}{X} 要開他們的眼睛，使他們從黑暗中轉向光明，從撒但權下歸向神；使他們因信我而得蒙赦罪，和一切成聖的人同得基業。』」 \end{tabularx} \\ \\ \relax
26:19 & \begin{tabularx}{0.7\textwidth}{X} 「因此，亞基帕王啊！我沒有違背那從天上來的異象； \end{tabularx} \\ \\ \relax
26:20 & \begin{tabularx}{0.7\textwidth}{X} 我先在大馬士革，後在耶路撒冷和猶太全地，以及外邦，勸勉他們應當悔改歸向神，行事與悔改的心相稱。 \end{tabularx} \\ \\ \relax
26:21 & \begin{tabularx}{0.7\textwidth}{X} 為這緣故，猶太人在聖殿裡拿住我，想要殺我。 \end{tabularx} \\ \\ \relax
26:22 & \begin{tabularx}{0.7\textwidth}{X} 然而，我蒙神的幫助，直到今日還站立得穩，向尊貴的和卑微的作見證。我所講的，並不外乎眾先知和摩西所說將來必成的事， \end{tabularx} \\ \\ \relax
26:23 & \begin{tabularx}{0.7\textwidth}{X} 就是基督必須受害，並且首先從死人中復活，把亮光傳給猶太人和外邦人。」 \end{tabularx} \\ \\
[1ex]
\hline
\hline
\end{longtable}
$^{1}$各位弟兄姊妹平安.
很高興能夠在今年的年底.
這個主日和大家在紅欣家一起敬拜事奉.
一起來分享上帝要向我們傳遞的訊息.
保羅在面對殉道的關頭.
在他回顧為主奔跑一生的時候.
他寫信給提摩太.
他說了一段可以說是豪情壯語.
這段說話大家可能都很熟悉.
他說美好的將我已經打過了.
當跑的路我已經跑盡了.
所信的道我已經守住了.
從此以後有公義的冠冕為我傳流.
就是按照公義審判的主.
到了那天要賜給我的.
為什麼保羅如此肯定.
甚至大家看他有沒有說大了.
他有沒有大言不慚.
他說主一定會賜他公義的冠冕.
剛才我們一起中讀的史徒行傳第26章1至23節.
就為我們提供了很清楚的答案.
原來史徒保羅如此肯定自己能獲得冠冕.
是因為他沒有違背從天上來的意象.
並且他終其一生以信服這個意象為他的職志.
赴湯蹈火 萬死不辭.
正如史徒保羅和已忽所的長老.
在離別時曾經說過一段字八.
我們一起讀.
我卻不以性命為念 也不看為寶貴.
只要行完我的路程.
成就我從主耶穌所領受的職事.
證明神恩惠的福音.
今天早上讓我們透過這段保羅向雅基拍王自辯的內容.
去了解保羅沒有違背從天上來的意象的原因.
帶來什麼結果.
以及他要為此付出什麼代價.
作為我們弟兄姊妹去迎接你們25年的主題.
懷抱報道初心 福音遍及社群.
這個主題的借鑒與激勵.
在這段經文中我們看到第一方面.

$^{41}$史徒保羅信服意象的原因是什麼呢.
在第四第十六節中我們看到.
是因為救贖的恩情帶來他生命的轉化.
大家留意在第十九節保羅這樣說.
我固此沒有違背那從天上來的意象.
請留意固此這個字眼.
這個字眼的重要是告訴我們.
保羅信服這個意象沒有違背這個意象是有原因的.
他並不是傻傻的一時衝動心血來潮.
所以他走上這條事奉的道路.
原因是在第四到第十六節.
保羅和雅基拍王的自辯裡表達出來.
跟大家溫習一下保羅基本的生平.
保羅原名叫做.
未信主之前.
這條難回答.
他使成一個猶太的經學大師.
他的師父叫什麼名字.
有人答到了.
加瑪列.
從這個角度來說.
他是系出名門.
用今天說的名字來形容.
即是名校出身.
他是希伯來人所生的希伯來人.
可以說是根正苗紅.
他是嚴守律法的法利賽教門的青年才俊.
可以說是前途似錦.
甚至他是為了消滅當時剛剛崛起的耶穌教派的人.
他是不擇手段去逼迫基督徒.
對他來說他是大發熱心.
雖然他滿手基督徒鮮血.
這個和耶穌完全敵對的猶太教狂熱分子.
他是前往大馬士革去捉拿基督徒的途中.
主親自向他顯現.
這個突如其來的經歷.
令到保羅永遠都不會和滅.
所以在《使徒行傳》第22章.
保羅向群情洶湧的猶太人自辯的時候.
他亦都曾經憶述這段被主光照的經歷.

$^{81}$在第22章第10節.
保羅說被主光照的時候.
他的回應是什麼呢.
他說主啊 我當作什麼.
大家留意對猶太人來說.
當他稱呼對象是主的時候.
他是承認猶太人引頸以待的那位救主.
這個稱呼包含他認信耶穌就是那位救主.
保羅遇見主蒙恩得救生命轉化.
他本來人生的職志是做什麼的.
本來是打壓耶穌的.
現在一個轉變.
現在他去傳揚耶穌.
再直白一點說.
他願意為耶穌賣命.
正如保羅在《菲勒比書》第3章7至8節說.
他說 只是我先前以為與我有益的.
我現在因基督耶穌的緣故都不是這樣看.
我當作有損.
不但如此 我也將萬事當作有損.
因我已認識我主基督耶穌為至寶.
我為他已經凋棄萬事.
看作糞土 為要得著基督.
保羅在哥林多前書第15章.
也很感恩地說.
我們一起讀.
我原是使徒中最少的.
不配稱為使徒.
因為我從前迫不得神的教會.
然而我今日成了何等人.
是夢.
為何保羅能夠信服異象呢?.
是因為他深深體會耶穌基督的救贖恩情.
所以他的生命取了一個180度的變化.
可以說.
耶穌的愛改編了保羅人生的惡章.
我今日第一次和大家見面.
當然在座有幾位我是認識的.
所以待會認識我這幾位是不能回答的.
你猜我未做牧師之前我做哪一行呢?.

$^{121}$給你三個選擇.
第一是殺手.
第二是歌星.
第三是校長.
有樣子在看嗎?校長.
其實我真是殺手.
因為我84年開始.
在宣道會陳元梨小學教書.
在沙田.
我是訓導主任.
大家心目中的訓導主任是空神惡煞嗎?.
當然在我年代已經不是了.
不過在84年我教書.
校長說「富城你高高大大不如你管學生」.
我做了訓導主任.
但當時做訓導主任很厲害.
好像警察拖槍一樣.
我還可以拖著藤條.
當時的制度還可以校長授權一個老師.
用藤條輕打同學.
記住是男同學.
女同學不能碰.
男同學手板幾下.
所以我做訓導主任那八年.
我已經打過無數男同學的手板仔.
所以那幾十年他們除了長大.
請我吃飯.
在食飯局的房間.
他們就玩我.
伸手出來說「李主任很久沒打過了」.
我說「你是不是很想打你兒子?」.
他們說「是呀」.
我說「你小心,現在不能打」.
歌星我真的開演唱會.
不過是自己搞好.
然後派票給別人聽.
差點要送叉燒飯.
我真的是校長.
我在2013年9月提早退休.
展開了我全時間事奉的人生下半場.

$^{161}$到今年已經是第11年.
我提早退休之前.
我做了29年的小學教育工作.
我頭八年就是做訓導主任.
後那21年就在人際醫院.
在大埔新辦的一間小學.
叫做人際醫院賽影圖小學.
擔任一位創校校長.
就好像我們教會開荒一樣.
我這29年教育工作裡面.
那些同事,家長,學生,校友.
甚至我的校董,人際的老闆.
都覺得這個校長有點特別.
有什麼特別呢?.
我特別疼愛那些.
聽清楚.
我特別疼愛那些叫學習能力薄弱的後進生.
這個名字好聽嗎?.
學習能力薄弱的後進生.
如果有上一年多.
就即是說沒用.
不要再讀書了.
罰他們出門在家.
即是一班有特殊學習需要的同學.
另外就是一些.
他沒有特殊學習需要的.
不過是專門賭博.
專門跟老師過不去的那些所謂的頑皮.
為什麼我對這些所謂一般我們覺得不合模的學生.
我特別疼愛他們呢?.
很簡單的.
想到嗎?.
你很有愛心.
因為我是過來人.
我小學的成績和我的操行.
比我過去的29年接觸過的所謂頑皮.
或者跟不上的學生.
他們還差我很遠.
意思就是.
他們所謂的頑皮和我小時候的頑皮.

$^{201}$還差很遠.
知道這條什麼村嗎?.
厲害吧.
現在是這樣.
以前是這樣的.
彩虹村是1960年開始入伙.
分幾期.
是一條很大的屋村.
1961年我就入住這條屋村.
我媽媽跟我說.
我是第三兒子.
我說我腳頭好.
因為我出生.
我在港華園出生.
我就入住彩虹村.
因為之前住旺角那些叫頭房,尾房,中房.
類似現在的tong 房那些.
我出生就全家搬去彩虹村.
知道彩虹村最偉大的地方是什麼嗎?.
是因為有獨立的廁所和獨立的廚房.
在60年代的公屋.
很多還是洗字區那些.
七層用公廁.
廚房在自己門口.
有個整個.
就是這樣.
我們那裡已經是獨立廚房和廁所.
所以是很偉大.
很有人性化的屋村.
但是就住了我這個嘍囉.
這張照片是我小六升中.
當時要考一種市分配學位叫升中市.
歐sir都看起來是考升中市.
誰考升中市的舉手.
哇這麼多.
不明白升中市的要解釋一下.
這個見證弟兄一定不明白.
升中市是什麼呢.
79年之後轉了學能測驗.
之前升中市就是.

$^{241}$李富城去考中英數三科.
你考到的成績就是你的了.
就像考DSE一樣.
我當時考到一個很亮麗的成績.
叫445.
反過來就叫做叻.
111就是以利沙白.
黃仁等著你.
445就基本上沒有書讀了.
其實.
就是沒有津貼學位了.
不過剛剛74年.
政府就開始幫一些私校買位.
我僅僅不是555.
555就是地底.
445就是還在邊上.
就派到去我認識曾牧師的一個機會.
就是在伯特利中學.
當時還在九龍迦南邊道的中文部.
有一個津貼的上午校的學位可以讀.
我媽已經開心到瘋了.
因為我兩個哥哥已經跟著我爸爸出去學師了.
當時.
是第三個兒子.
讀得上中學.
讀得上中學已經在我那層樓很厲害了.
還有中學讀嘛.
其實我小時候是草根階層.
我那層樓的地方叫波底.
知道什麼叫波底嗎.
誰知道波底的翻譯應該叫什麼.
球場.
球場現在是上網訂的.
有界線有龍門.
有人來爭場.
你可以叫管理員趕他們走.
管理員都不聽.
就叫警察出來.
不能爭的.
但是當時.

$^{281}$我想我大概四五歲已經流連街頭.
流連球場.
現在拍照這麼漂亮的地方.
那個籃球場以前是平地.
因為要弄時租停車場所以升高了.
當時那個大平地就是我們叫波底.
整個彩虹村的小朋友.
包括牛池灣和釜山道大漢村那些.
全部都是來這個最大的波底爭場踢球的.
大家知不知道爭場踢球用什麼來分配.
兩樣東西.
首先就是粗口.
罵走他.
你夠粗口罵走他會害怕.
不怕.
胸口回來.
第二就是拳頭.
所以波底就叫英雄地.
所以波底就是很容易學壞的地方.
媽媽就最擔心我們爭球地.
但是當時上下午班的制度.
我們除了上學就是爭球地.
所以波底裡面所有的壞習慣.
我都浪浪上口.
包括粗口當然.
打架.
偷東西.
現在叫店舖盜竊.
挺好聽的名字.
偷東西.
看色情書.
全部在球場集團的壞行為.
我都是一個當時的一個球場的小混混.
在黑社會群的同黨的邊緣裡面擦身.
為什麼我有轉變呢.
是和這個女士很有關係的.
你猜一下她是我誰.
我問過很多次.
九成都猜是我太太.
她是我在小五的時候遇到一個代課的老師.

$^{321}$她叫做劉月琴.
英文名叫Queenie.
秘密是這張照片是去年拍的.
去年我62歲.
她72歲.
所以女士的樣子是不能從外貌看的.
女士的年紀不能從外貌估計的.
她長大了十年.
這個秘密就不是秘密了.
她做我的代課老師難道她長大了兩年嗎.
我當時是小五十歲.
我當時在班房裡面都是一個壞的小領袖.
原任老師我們是不敢去.
因為原任老師當時是可以看法學生.
所以是絕對不敢和原任老師抬槓的.
但是代課老師就慘了.
因為代課老師不敢亂來.
所以我們就專門找代課老師來發洩.
對了班主任進入大修.
來一個代課老師.
大家就看著進來的是什麼樣子.
想著作弄她.
一進來.
咦這麼漂亮的女生.
又年輕.
你想想當時我十歲她二十歲.
現在她六歲.
她七十二歲都這樣.
你想想她二十歲的時候是什麼樣子.
我們一群小五的男生已經青春期.
嘩這麼漂亮.
於是我這個壞的小領袖就叫兄弟們先不要騷擾她.
先看她怎樣.
她教得很好.
很有心機.
雖然沒有受過師資訓練.
但是她很有心機教.
加上很有活力.
我們很喜歡上她課.
李富城你回答.

$^{361}$嘩整個倒下了.
但是快樂的時光會過的.
原任老師會回來的.
因為原任老師還是穿長衣的.
因為落差太大.
我們很失望.
她要走.
她說當時她暴露身份.
她說我是基督徒.
我們雖然在讀聖公會的小學.
都聽過耶穌.
但是都不太知道是什麼.
她說基督徒.
我星期日上教會的.
你們哪個小朋友想跟我上教會的.
舉手.
那你猜我們班的同學舉不舉手.
舉手了 旁邊的聽到都舉手了.
接著我們就很失望.
因為她告訴我們.
我們的教會和我們住的彩虹村很遠.
我住彩虹村.
但是當時的教會還是在.
這個深水Po 的舊樓.
這幅是在我母會的舊樓的照片.
怎麼去啊老師.
她說不用怕.
你們夠多人去.
我可以叫教會安排一部旅遊車.
1972年.
逢星期日來彩虹村接你們去.
還有我假假地都是老師.
原來她也住彩虹村.
我假假地都是老師.
聽她跟你媽說.
有老師帶旅遊車坐.
我相信你家長會讓你跟我去.
我真的這樣回去跟我媽說.
我媽第一句就說.
臭小子 跟你老師去吧.

$^{401}$比拉波地好.
是不是.
我就是這樣踏足我的母會.
回到兩個月後.
母會的姐姐就說.
傅仔 我們搬了.
搬去哪裡呢.
就搬到何文田的常和街.
以前的何文田村旁邊.
的一個新的座堂教會.
這個座堂大概有52年了.
我回到我成長的母會.
已經超過半個世紀了.
我回到教會之後.
很快有人跟我說.
有下令會去.
知道什麼是下令會嗎.
就是營會.
還是五天四夜.
當時很喜歡小朋友.
不喜歡在家睡.
喜歡在街上睡.
五天四夜拿正牌.
去一個叫宣道園.
宣道園現在都暫停了.
因為疫情的緣故.
還很記得.
當時的聚會.
不會分青年人成年人.
什麼都混在一起.
聽牧師說.
說了一小時.
我們都流口水了.
是睡著的流口水.
但我們識字.
我們看到時間表.
熬完一小時.
就看著宣道園外面的大球場.
我們是踢球出身的.
有草場踢.

$^{441}$就帶著球走.
一起來.
有個哥哥走過來.
小褲子先不要吹.
我說什麼事呢.
跟你聊天.
知道聊什麼嗎.
知道什麼叫個人談道.
我告訴大家.
第一句.
今天大家個人談道.
不能用.
第一句就問我.
你是不是醉人.
如果你今天跟人個人談道.
第一句問人.
你是不是醉人.
我想人家要麼起身走.
要麼摑你一巴.
第一句就問我.
是不是醉人.
他當然知道我的背景.
因為我們進營會.
還在說髒話.
還被哥哥姐姐罰.
他問我是不是醉人.
我當然說是.
很簡單的一個帶領.
一個很簡單的信心.
凝訊開始進入上帝的家.
在基督耶穌的愛裡面.
慢慢來學習成長.
正如哥倫西書三章九至十節話.
我們一起讀.
因你們已經脫去舊人和舊人的行為.
穿上了新人.
這新人在知識上漸漸更新.
正如做他主的形象.
我這個邊緣街童.
因為耶穌的愛改編了我人生上半場的樂章.

$^{481}$我一邊去做教育工作.
一邊在教會有很積極參與的事奉.
去做很多建立生命的工作.
也是由於耶穌基督這份深深的愛.
他改編了我人生下半場的樂章.
呼召我用我的人生下半場.
來全時間服侍他.
我們每一個人蒙恩基督教的切入點.
都不同.
剛才這位小弟兄分享.
他有兩個姑姐.
又遇到一位陳校長.
可能我們每次洗禮聽這位弟兄姊妹講他的見證.
你會發覺神真的很奇妙.
他用不同的方法.
有些人陪女朋友回教會.
陪自己信了.
有些人因為自己的子侄回教會.
老人家覺得子侄回教會很好.
來教會聚餐聚餐聚餐聚餐.
就信耶穌.
我爸媽跟我回我母會聖誕節平安夜聚餐.
起碼超過四分一個世紀.
因為我們每年平安夜晚都會有一次聚餐.
聚餐完有一個方聚會.
不是吃完酒.
一定會跟著我上去.
他也很自然會跟著我上去.
難道先離開嗎?因為要等我載他走.
其實他聽了二十多年的方法.
有一年很特別.
突然跟我說.
我們一起住.
阿父明天我跟你回教會.
在我爸爸媽媽七十之靈.
他跟我回母會受災陪受浸.
有十多年的時間.
他們才安息主懷.
每個人的切入點歷程不盡相同.
不過我相信我們有相同的經驗.

$^{521}$就是我們有被主光照我們犯罪本上的經驗.
正如保羅所說.
我今天成了何等人.
是蒙神的恩才.
頂止每樣我們一同審查.
基督的救贖恩情.
是否就是激勵你願意順服上帝旨意.
實踐去得人使命的原動力和原因.
或者我直接問.
耶穌的愛在今天.
是否仍然改編你人生的樂章.
我在教會一段很長的時間.
有很多信徒.
我們常說笑.
他叫老信徒.
什麼老呢.
就是信徒之後原地踏步.
他不動.
耶穌說不可愛.
但每次說耶穌愛他的經驗.
都是說他幾十年前的事.
其實耶穌和我們愛的關係.
就像夫婦一樣.
是歷久常新.
如果我們用一個.
以前耶穌不知道在什麼機會拯救我.
如果我還在說小午回教會的經驗.
我這五十年和主的愛關係.
沒有更新的話.
我相信我沒有足夠的動力.
去回應主託付給我的使命.
弟兄姊妹.
我仍然要問一次.
耶穌的愛在今天.
是否仍然改編你人生的樂章.
第二是這段經文.
我給你看到.
很信服異象的結果是什麼呢.
對保羅來說.
就是他終生傳揚福音真道.

$^{561}$保羅明白信服異象的原因.
我們就看看他帶來什麼結果.
原來保羅所說的異象.
是包含除了有蒙主拯救的部分.
還有被主差派的部分.
在第十六節我們一起讀.
你起來站著.
我特意向你顯現.
要派你作執事.
作見證.
將你所看見的事.
和我將要指示你的事.
證明出來.
然後去到第十八節.
主又這樣和保羅說.
我差你到他們那裡.
要叫他們的眼睛得開.
從黑暗裡歸向光明.
在撒旦的權下歸向神.
由因信我得蒙赦罪.
和一切成聖的人同德紀念.
請大家留意.
保羅所領受的救恩.
除了是他個人蒙恩得救的部分之外.
其實還有被主差派.
將這個救恩繼續傳揚的部分.
保羅的見證和宣講.
引導了很多人在撒旦的權勢下歸向神.
保羅領受從主差派傳福音的意象很清晰.
所以在以弗所書第三章.
保羅描述自己是福音的執事.
他就一生勞苦盡心竭力.
用各種智慧勸誡人.
教導各人.
也福音隨著保羅宣教的足跡廣傳.
不同城市的教會就一個一個被建立.
例如和保羅非常密切關係的肥納比教會.
我小時候就不懂肥什麼肉教會.
經常會幫聖經改化名.
肥納比教會是在說什麼呢.

$^{601}$是在說保羅第一個踏足歐洲宣教的橋頭堡.
大家留意當時的歐洲不是你現在心目中想去歐洲旅遊的那種歐洲.
當時的歐洲是落後的.
是貧瘠的.
例如你今天想去哈薩克斯坦宣教.
保羅的宣教旅程.
幫這個宣教的歷史揭開了輝煌的序幕.
不知道大家有沒有認識或聽過.
這位前中神院長李思敬博士.
李思敬牧師的講座或看過他的書.
李思敬牧師的父親叫李飛吾牧師.
和大家的宗派有關係.
是宣道會希伯倫堂服侍了相當長日子的一位忠心的上帝使用的僕人.
李氏父子兩人蒙恩得救.
廣傳福音造就無數的信徒.
但大家知不知道李氏父子是受一個人影響.
被匿稱為何二姑的一位宣教士.
叫做何二思教士.
李思敬牧師稱她為何教士婆婆.
這位小女子有什麼把炮可以影響李氏父子.
間接影響了我們環教會的力量.
話說在二十世紀初.
美國加州有一個基督徒家庭.
他們很喜歡接待宣教士回來述職請他們回家吃飯.
美國人的文化吃完飯就會去客廳喝咖啡聊天.
在這個時間聽很多宣教士.
在相對落後的地方聽很多宣教故事.
他們在分享的過程中.
有位何牧師在中國某個村落.
有一個村民推著一輛木車來到他們的報道所.
問他們的洋人有什麼耶穌.
可不可以派人到他們的村落說說.
如果你是宣教士你已經發達了.
因為他們開門想找人去.
何牧師說不知道為什麼那次沒有人回應.
於是這個村落的村民.
推著木車在黑夜裡長長的身影消失了.
在這個一起聽故事的家庭當中.
有位小女孩只有十歲.
她彈起身說了一句.

$^{641}$在中國有條鄉村想聽福音.
但看到的只有一輛空的木車.
如果是我我就一定去.
如果大人聊天時有個小朋友彈起身說出這樣的豪情壯語.
你們會猜拳還是壓下去.
她當時被人勸小孩坐下.
她這個回應.
她這個立志藏在她心裡.
但隨著她升上史丹福大學.
她已經忘記了.
這個木車的意象.
她已經忘記了.
不過她當時立志要做松商.
在她升大學二年級.
她媽媽寫封信給她.
不知為何在信裡提到中國.
這個木車的意象再一次在這個女士.
小女孩的心裡走出來.
終於經過一番掙扎.
這個小女孩信服上帝給她的意象和猜拳.
在1913年11月.
這個21歲的小女子遠竊重陽來到中國.
輾轉間來到廣東省南海縣西朝的觀山傳教.
建立禮拜堂孤兒院.
稱為李飛吾牧師的原名叫李真徒.
是一個孤兒.
被人送到這個孤兒院裡被接濟.
當然也因為這個緣故聽聞福音.
以致她獲得照顧聽聞福音.
這個小女子就是希伯倫會的創辦人.
何二師教士.
何二師教士21歲踏足中國.
直至1971年78歲才退休.
她在華人裡服侍了超過60年.
因為信服木車的意象.
她拋下了加州的陽光.
她捨棄了高校的學位.
埋葬了她從商的壯志.
她離別了她溫馨的家庭.
留意她是千里迢迢跑到一個陌生的國度.

$^{681}$她經歷過的是什麼.
中日戰爭.
太平洋戰爭.
國共內戰.
種種的困境和危險.
但她始至終生傳揚.
福音真道.
這本是大家在坊間書局買到的.
叫《誰掌管明天》.
就是何二師教士的傳記.
李子英博士在一個序言裡說.
如果何教士沒有來到中國.
今天不會有我.
弟兄姊妹在回應主的意象當中.
有一首詩歌的副歌.
我覺得很窄心來挑戰大家.
不知道大家看過這首歌.
有沒有印象.
唱過.
在《青年聖歌》裡有.
叫做《萬世戰爭》.
副歌裡問了四個問題.
是否你見此意象.
是否你心激動.
是否你應主呼召.
是否你願聽從.
真的唱得到.
懂得唱的同學唱一次.
是否你見此意象.
是否你心激動.
是否你應主呼召.
是否你願聽從.
第三.
信奉意象是要付代價的.
保羅付什麼代價?.
他是獻奉所有他擺上一生.
史陶保羅提到他信奉意象之後.
剛才我們讀到第20,21節.
就說他去了大馬士革.
耶路撒冷.

$^{721}$猶太傳遞.
和外邦傳揚悔改歸主的好訊息.
不過.
他是因此觸怒了猶太人.
因為他原本是他們的陣營.
現在竟然調轉槍頭去傳揚耶路撒冷.
所以有一群狂熱的猶太人.
要追著保羅.
如果你讀史陶行傳.
你會留意到保羅去到哪裡.
都有一群人追著他.
不是粉絲.
是殺手.
為什麼說是殺手?.
保羅的命其實是命是一線.
在史陶行傳第23章12節第15節.
提到他們想迫害保羅.
但當時的駐防官覺得.
第一保羅有羅馬公民.
不知道發生什麼事.
一群猶太狂熱分子要迫害一個人.
不知道發生什麼事.
就將保羅囚禁起來.
實際就變相保護了他.
但那群猶太狂熱分子就不打算這樣.
他們想怎樣呢?.
他們一起同謀起誓.
什麼叫同謀起誓?.
如果要畫面.
就是一起喝完酒.
丟在地上.
不殺他又怎樣?.
若不先殺保羅.
就不吃不喝.
意思是什麼?.
不是他死.
是我死.
我將命放在身上.
一定要將保羅置諸死地.
還不是殺手?.

$^{761}$他們預備好.
要準備去伏擊保羅.
不過這個消息外洩.
所以這個駐防官不敢怠慢.
認為自己那裡的保安系數不夠高.
所以千呼掌就將素羅要移送去.
另一個菲力斯巡撫官邸.
即是加強保安.
但大家請留意.
移送隊伍.
壓界隊伍.
龐大的保安隊.
不如先問你.
你看到紅色字有多少人?.
編制.
470人去護送一個人.
習主席出巡.
我不知道保安是怎樣.
但表面看都沒有470人.
可能背後有.
470人全部是荷槍實彈.
全力武裝裝備.
步兵200 馬兵70 長槍200.
是一個什麼陣勢.
其實你要留意千呼掌是職業軍人.
他評估風險才會派多大的保安隊.
他不是貪玩.
推470人出去玩.
不是的.
這個職業軍人其實是評估.
如果素羅在途中.
不夠足夠的震懾力.
猶太的狂熱殺手就會出現.
所以這個評估其實是反映保羅的命是怎樣.
比水還要冷.
保羅亦曾經在哥林多後書第十一章.
列舉過一張受苦的清單.
我小時候讀這段的時候.
有時覺得保羅是在曬命.
想說他為主怎樣辛苦.

$^{801}$想抬高自己.
其實不是.
當我細心閱讀的時候.
保羅其實是想和教會說清楚.
當我們要跟隨主的路上.
去順服主的意象.
你是要付代價.
他只是將自己付的代價.
如實地向哥林多教會的信徒表達出來.
保羅為了宣揚真道.
他作為基督僕人所受的患難.
我們是很明白.
順服意象是要付代價.
而且付這個代價是有切腹之痛的.
什麼叫切腹之痛呢?.
這四個字太文雅了.
有沒有人可以用廣東話直譯給我聽.
切腹之痛是甚麼?.
廣東話常常說的這兩個字.
你有甚麼消息.
你會切腹之痛?.
例如我有兩個孫子.
我小時候的男外孫.
大約三個月.
如果我女兒突然傳訊息給我.
爸爸她撞破頭了.
我會有切腹之痛嗎?.
當然有.
如果我不用切腹之痛.
你會用廣東話的兩個字來形容嗎?.
這兩個字.
今天整堂課你記得這兩個字都可以了.
肉.
差一個字.
肉赤.
有甚麼令你肉赤呢?.
你今天進來.
十一奉獻你有沒有肉赤?.
今天一早回來.
預備大靈師的敬拜.

$^{841}$這麼冷起來有沒有肉赤呢?.
原來去信奉.
意象到一個境地.
你要有肉赤的味道.
你就知道甚麼是付代價.
剛才我說我很喜歡唱歌.
如果你過個星期.
周牧師說你這麼喜歡唱歌.
過個星期你也領師吧.
我肉都不用赤.
我多開心.
吃我那頓飯.
過個星期都像自己唱卡拉OK.
最好有全組.
甚麼是肉赤?.
當你信服的時候.
甚麼是信服?.
你要做一些你不喜歡做.
但上帝要你做的事才叫信服.
如果上帝叫你做一些你很開心做的事.
你都不知道.
你何時有掙扎呢?.
就是當你.
很清楚主呼召你回應.
他的心意做了一些事.
但那不是你的議程.
不是你的心水.
不過你要調整.
去滿足上帝的心意.
這才叫信服.
這才叫沒有違背.
從上帝而來的.
那個意象.
對我來說.
大家知道.
我剛才說我真正的身份是做校長.
我要離開.
我校長的.
角色的時候.
有沒有掙扎呢?當然有.

$^{881}$但當我.
其實蒙召.
40歲左右已經.
有這種經歷.
我自小回教會.
甚麼會?甚麼奮興會?.
下令會?總之那些講完.
一生事奉主,全職事奉神.
那些我全部不會回應.
因為我覺得那些不是我的.
那些是隔離的.
但很奇妙到我40歲.
當了10年校長,上帝真的透過很多.
不同的印證.
來挑戰我.
不過我拖了很久.
剛才你聽我說.
我52歲才全職事奉.
我拖了12年.
不敢謝主.
在2009年.
我出席戴紹增牧師.
的一個追思會.
戴紹增牧師是甚麼邊位呢?.
是戴德生牧師第四代孫.
他的下一代.
也是在華人地區事奉的第五代.
戴繼宗牧師.
他1929年.
因為宣教的緣故.
他的父祖輩在中國.
所以在中國出生.
他的普通話比我們任何人流利.
包括就算你在內地.
來香港都不夠他厲害.
他的流利的普通話.
令到他有一段時間.
在國內做扶貧工作的時候.
內地的官員很敬重他.
因為他根本是中國人.

$^{921}$中國出生的.
他曾經.
坐過日本人的.
集中營.
十多歲的時候.
他立志一定不會像他爸爸.
爺爺那樣.
做宣教士.
就算做也不在中國做.
不過很奇妙神又帶領他.
後來釋放接受了神學裝備之後.
他又回到我們華人地區.
做了很多.
服侍華人的工作.
直至在香港離開世界.
我們那次安星.
那次追思會.
其實我是特意.
沒有叫其他弟兄姊妹.
我只是想看看這位淑玲的.
先鋒前輩的生命.
很奇妙神透過那次的聚會.
肯定了我的夢照.
我是因為.
我從小回教會.
聽很多很感動的會.
我都不會哭的.
但是聽很多人.
上去說大少周牧師的見證.
我不斷流淚.
說到盡都說不完.
裡面的催逼.
李富城你拖了很久.
你看人家.
你還拖什麼.
所以那次2009的聚會.
我就立志.
當然要安頓好學校的工作.
所以我2013.
踏上.

$^{961}$大家想我要回應這個呼召.
有什麼肉痛.
你們一定想.
當然是薪金.
大家知不知道.
小學校長薪金有多少錢.
現在算吧.
不要說.
目前一間.
標準夠24班或以上的小學校長薪金.
是12萬8千元.
還有那筆公積金.
承諾的紅利.
你又想.
馬上質問我.
牧師你想浪費命嗎.
我想告訴你.
這份薪金和公積金.
絕對不是我當時覺得肉痛的東西.
因為我常常覺得.
金錢物質.
夠就行了.
我又不相信做傳達人會餓死.
因為我肚餓就會.
找你吃飯.
你當然請我吃.
我又找秋先生.
喝杯茶.
秋先生當然請我喝.
這是華人文化敬重目的.
所以我當時覺得.
錢絕對不是我考慮的原因.
雖然我的行家.
個個都很緊張.
你有沒有想清楚.
做校長教會事奉就好.
為什麼要做傳達人.
他們全部知道.
掛勾於收入.
是他們關心的.

$^{1001}$價值觀也不同.
我今天穿這件西裝.
我穿了20年.
原來一件深藍色的西裝.
對男士來說.
這是我做校長的工作服.
我可以去婚禮.
可以去喪禮.
肥一點穿也可以.
你看到很大件的.
我瘦一點穿也可以.
我懷疑這件衣服穿到天父.
見天父也可以.
我不用買很多件西裝.
去弄一個很大的衣櫃.
每天早上起床就有選擇困難症.
今天選哪件.
所以我相信有很多蒙召者.
他們在物質生活上.
他們都很容易能夠跨越.
對我來說.
最難的是什麼.
我想和你分享.
我要離開一個安書區.
這是更難.
什麼叫安書區.
我做了21年校長.
當時我在教育界已經有些名氣.
如果我繼續選.
如果有選立法會的工人界別.
還有的話.
我多數會選到.
因為我當時也是教協的熱心分子.
因為我年輕時做校長.
我30歲已經做.
所以我已經在業界.
有一些名氣.
我在大埔區做了21年校長.
大埔區裡面.
官政商民警都很通的.

$^{1041}$所以你在大埔區打橫行.
不是叫你亂來.
意思是你在裡面有很多連繫.
笑話就是.
我有一次去酒樓.
當時酒樓還是很旺市.
想去喝茶.
看到一條龍很長沒有位置.
想調頭走.
有個知客就叫.
校長.
校長想喝茶.
沒有位置.
排隊很久.
不是 你當然有位置.
有些排隊的.
有沒有搞錯.
那個知客就說.
傻的嗎?那個是我的校長.
在華人文化.
雖然不是很好.
但也很快樂.
我會離開一個.
我已經很有安全感的地方.
我有秘書.
我有整隊隊伍.
我的電話號碼做得好嗎?.
我告訴你.
我是離開了校長之後才一筆一手學的.
以前是怎樣.
修文進來幫我做powerpoint.
幫我做好.
你幫我做好.
以前是這樣.
還有一個校長室.
有時累到可以打瞌睡.
我離開一個我熟悉的行業.
現在是我喜歡的.
我現在做很多校董.
我很喜歡做教育界.

$^{1081}$我要離開一個.
我熟悉的環境.
我做傳代人.
我要重新學記名.
我重新要學沒有秘書.
重新要學沒有辦公室.
沒有獨立辦公室.
其實對一個已經習慣了的人來說.
是很難的.
弟兄姊妹.
不過我說這麼多.
我就想和你分享.
相對於初代教會.
或者教會歷史上.
為主擺上命的那些人.
我們這些算是什麼.
算是什麼.
有什麼這麼肉赤.
當然我現在和你這樣說.
我已經跨越了這個安書區.
你現在叫我不做牧師.
我又要離開另一個安書區.
因為我已經明白了.
上帝要使用我在這個目的.
去服侍他.
剛才我讀到第23節.
其實經文有一節.
有一節很有趣的.
是第24節.
當時.
陪聽保羅說.
阿基帕王其實是八卦的.
提醒保羅出來.
讓他說一下他做街下囚的事情.
陪聽的有一個巡撫.
叫做菲斯都.
菲斯都是外邦人.
他不知道.
保羅前面的23節說的話.
他根本不想說.

$^{1121}$於是他就回應了.
他說保羅你癲狂.
你的學問太大.
叫你癲狂.
頂主妹.
有時候會覺得有些微信的朋友.
親人會覺得你做基督徒癲癲的.
星期日好好的.
不去喝茶回教會.
你傻的嗎.
什麼十一奉獻.
自己都還沒夠小.
你還奉獻你癲的嗎.
小朋友讀書這麼忙.
現在起跑線很重要.
明天考試.
你今天早上還叫他回兒童組.
你有沒有幫小朋友想過.
其實你在職場裡.
乾乾淨淨.
君君珍珍的做過職場人.
你的同事都會暗地裡.
傻的嗎.
其實基督徒.
真的要癲癲的才行.
如果有些人跟你說.
你信耶穌.
信到癲癲的.
我為你感恩.
因為你在做的事情.
跟他不同.
這是從信仰而來的.
不過我告訴你.
那個笑你癲癲的人.
他遲些就會羨慕你.
他會敬重你.
他開始會欣賞你.
甚至到他.
聖靈感動他.
他會拉著你三尾.

$^{1161}$今年聖誕節你教會有什麼聚會.
我都想去聽聽.
我鼓勵靈姐妹.
在新的一年.
在教會的主題帶領之下.
我們在這個社區.
在你個人所遇到的人際環境.
你的家族.
你的職場.
你退休人士的茶局.
你北上消費的那班老友.
如果在讀書的你的校園.
上帝都仍然在問我們.
好不好像保羅一樣宣告.
我們一起讀紅色的紙.
我固此沒有違背那從天上來的意象.
有一首詩歌是非常好聽.
不過我預告給你聽.
是極之難唱.
不過你跟著我唱你會唱到.
是一首很美麗的旋律.
譜了保羅的生平的一首詩歌.
我一路唱完三節.
你何時想加入都歡迎.
我們一起.
戰下天起保羅之光.
透過這首詩歌.
來回應我們剛才所聽的道.
大麻色圖上主基督顯現.
榮耀聖光忽照在人面前.
沒法解釋其已刻骨的經驗.
但卻深知一生將要改變.
從前拼命迫伯世上殉途.
還為自己當天所作自豪.
然而耶穌的天然光照路途.
十家捨身的哀感撐上渡.
自那天起全力奔跑崎嶇窄路.
無懼困逼不霸浪淘淘淘.
夜冷風高仍未想過實屬腳步.
生中主所託眾路.

$^{1201}$條條萬里即使翻過暴雨空途.
茫茫夜晚怎守至處全無.
然而神恩典依舊鋪滿路途.
讓我窮盡一生將主愛回報.
為了基督無論經過多少試驗.
名利負眾一概拋卻身邊.
病痛較戰仍未改心底的信念.
直至那天必達不留觀念.
全力進醫!全力進醫!將福音繼續篇傳.
無問代價只想天國萬言.
全憑神恩典相伴這數十年.
就算由天傾出身他也無怨.
讓我們同心禱告.
多謝主你今天透過使徒保羅在雅基柏王面前的自白.
讓我們明白我們每一個跟隨你的人.
我們要信服你給我們從天上來的異象.
我們一定要深深再一次察慮我們與你的關係.
被你的愛所激勵.
以至我們願意跟人分享這份寶貴的救恩.
又因此求你讓我們明白我們要回應你的愛.
我們不單只是在唱詩.
在口頭.
是落實在我們的福事上.
今天可能在當中不同的弟兄姊妹有不同的掙扎.
有不同在新的一年正在計劃他們怎樣去過他們的日子.
求主就將你的國優先的考慮放在我們的計劃當中.
讓我們去回應主祈義的恩典.
祈禱奉耶穌基督的聖名.
阿們.
\newpage



\section{瑪拉基書 3:6-12}
\label{sec:0jw5FB1CQ6s}
\textbf{2025.01.05 《我們一起來試驗神》 瑪拉基書3\_6-12 (和修版) 講員 : 曾裔貴牧師}
\newline
\newline
連結: \href{https://youtube.com/watch?v=0jw5FB1CQ6s}{\texttt{https://youtube.com/watch?v=0jw5FB1CQ6s}} ~~~~ 語音日期: 2025-01-05
\newline
\newline
\hyperref[sec:tR6kQF63JQE]{\small{< < < PREV SERMON < < <}}
~
\hyperref[sec:index_chronic]{\small{[返順時目]}}
~
\hyperref[sec:index_scriptual]{\small{[返順卷目]}}
~
\hyperref[sec:gjk_0mIGHss]{\small{> > > NEXT SERMON > > >}}
\newline
\newline
瑪拉基書 3:6-12
\newline
\begin{longtable}{cl}
\hline
\hline
章節 & 經文 (和合本修訂版)\\
\hline
3:6 & \begin{tabularx}{0.7\textwidth}{X} 「我－耶和華是不改變的；所以，雅各的子孫啊，你們不致滅亡。 \end{tabularx} \\ \\ \relax
3:7 & \begin{tabularx}{0.7\textwidth}{X} 從你們祖先的日子以來，你們就偏離我的律例而不遵守。現在你們要轉向我，我就轉向你們。這是萬軍之耶和華說的。你們卻說：『我們如何轉向呢？』 \end{tabularx} \\ \\ \relax
3:8 & \begin{tabularx}{0.7\textwidth}{X} 人豈可搶奪神呢？你們竟搶奪我！你們卻說：『我們在何事上搶奪你呢？』其實就是在你們當納的十分之一奉獻和當獻的供物上。 \end{tabularx} \\ \\ \relax
3:9 & \begin{tabularx}{0.7\textwidth}{X} 因你們全國上下都搶奪我的供物，詛咒就臨到你們身上。 \end{tabularx} \\ \\ \relax
3:10 & \begin{tabularx}{0.7\textwidth}{X} 你們要將當納的十分之一全然送入倉庫，使我家有糧，以此試試我，是否為你們敞開天上的窗戶，傾福與你們，甚至無處可容。這是萬軍之耶和華說的。 \end{tabularx} \\ \\ \relax
3:11 & \begin{tabularx}{0.7\textwidth}{X} 我必為你們斥責蝗蟲，不容牠毀壞你們的土產。你們田間的葡萄樹，果實未熟以先也不會掉落。這是萬軍之耶和華說的。 \end{tabularx} \\ \\ \relax
3:12 & \begin{tabularx}{0.7\textwidth}{X} 萬國必稱你們為有福的，因你們必成為喜樂之地。這是萬軍之耶和華說的。」 \end{tabularx} \\ \\
[1ex]
\hline
\hline
\end{longtable}
$^{1}$剛才主席已經和大家有新年的文案.
也是新一年的開始.
祝福大家在靈命上.
身體上.
以及在教會裡的合一.
感受到這是神的家.
一起在主裡彼此同心的服事.
我們在開始講道之前.
一起同心禱告.
再一次多謝主.
我們要開始來打開聖經.
很盼望剛才這一段聖經.
我們能夠掌握明白.
以及到了新約已經是二千多年後.
我們應該如何來落實.
形容在我們每一個地方的教會.
求主幫助.
奉耶穌基督的名.
阿們.
試驗.
開了不好意思.
大家都知道了.
這一段經文.
是神藉著馬拉基先知.
來邀請.
不是一種命令.
因為說明要試驗神.
是一個邀請.
如何來試驗神的信實.
如何試驗呢.
很清楚第八第九節.
就是你們已經虧欠了我.
搶奪了屬於我的財物.
你們沒有履行.
全國上下沒有去做.
當時馬拉基的時代.
是主前大概四百二十五年.
是第二聖殿已經建好之後.
一段長時間百多年了.
因為五三六年的時候.

$^{41}$這個第二聖殿重建.
所羅巴巴已經建好了.
就是說猶太人.
已經在外邦的波斯地.
回到重建聖殿.
在那裡慢慢安居樂業.
有輾轉的.
有所羅巴巴的第一代.
也有第二代的以色列.
和第三代的尼希米.
他們已經回到重建.
有自己的國家.
有自己的城牆.
有自己的敬拜.
有自己國民的身份.
但是到了這個時候.
主耶穌還沒有來這個世界之前的主前.
四百多年.
神就藉著馬拉基先知.
來到有不同的.
對當時的以色列人.
有那個指控.
就是提出他們現在.
屬靈的整體的那個光景.
我們收集的範圍就是第三章.
提出了一個.
好像說你們有這樣的欠缺.
一個對話的形式.
那些國民就說.
我們哪一方面來藐視你.
我們哪一方面做得不夠好.
我們怎樣來搶奪.
提出一個反問的問題.
問回馬拉基先知.
先知很清楚告訴他們.
也給了一個方向.
要他們來有一個迴轉.
所以今天很重要的題目.
就是迴轉向神.
還有在我們今年建構一個.

$^{81}$個人的平台.
教會的平台.
去測試這一位不改變的神.
第六節他是不改變.
所以亞國之子沒有因為犯罪而滅亡.
因為三章五節我就沒有引述.
三章五節列舉他們種種的罪惡.
在社會上.
哈哈,爸爸,有錢人,欺壓窮人.
各種罪惡.
神列舉出來.
藉著先知.
但是因為神是不改變.
他是信實對這個民族.
有那份牧羊永恆不死的愛.
所以讓這個民族仍然.
可以在神的面前.
有這樣的方向.
改變,迴轉.
所以今天藉著這個迴轉.
我們要來到一齊.
紅恩間的弟兄姊妹.
一起去學習這個課題.
很清楚,很簡單的.
就是什麼迴轉呢?.
就是拜託你們以色列人.
在神的聖殿裡面.
你們做好你們自己的十一奉獻.
這個就是當時的經文的背景.
但是怎樣延伸到我們今天.
新約的信徒,新約的教會呢?.
這裡有一個很重要的進程.
我們一起去學習.
首先就是試驗.
我有一個很特別的經歷.
當我預備的時候.
就立刻浮現了一個.
我很多年前初信的時候.
的一個試驗.
那時候剛剛在醫院做.

$^{121}$剛剛信耶穌.
就很熱心.
一放工就上去醫院.
去那些病床.
因為當時沒有疫情.
也不會怕什麼人白撞.
你去哪裡人家都不會阻撓你的.
跟那些護士站說一聲.
就可以.
因為更加我是在醫院做.
更加身份容易.
五點鐘放工就上去.
派單張給那些病人.
慢慢醫院很多的同事.
都知道我是基督徒.
還有我又很熱心跟那些同事講福音.
有一次.
有一個同事就說.
我們現在就去吃午餐.
因為我們每個月.
會有一些假期.
是政府准許的.
可能部門的上司准許.
就是那個小時吃飯.
可以延伸到一個半小時.
就一群人一起出去吃飯.
去到紅橋吃飯.
大家一起吃飯.
後來那個上司就告訴我.
我們不只是去餐廳吃飯.
我們現在要去唱卡拉OK.
在當時我信主初期.
教會教導我.
當時去卡拉OK是不對的.
我們是不應該去的.
當時不是說現在.
當然我也不喜歡唱歌.
所以很少去.
我兒子就會去.
我聽到之後.

$^{161}$我有一個決定.
我可以跟你一起吃飯.
但去這些地方我就不去了.
其實是午餐.
其實是唱卡拉OK.
不是唱什麼流行曲.
但因為當時我初信.
教會的教導是不要去這些地方.
這些地方未必那麼屬靈.
不要理這個神學.
當時我真的這樣.
回應我的同事.
其實很兇.
兇得很厲害.
好像沒有彎轉.
接著我堅持不去.
原來我的上司是特意針對我.
我現在就試一下他.
看看他是不是真的.
說得自己那麼勁乾.
那麼屬靈.
那麼什麼.
原來他們跟那些同事打賭.
是我不知道的.
但當時我就堅持.
當然那個對錯.
今天一定不是應該這樣看.
但當時我的神學.
我的信仰.
就說我不應該去.
我告訴他們.
我不是跟他們有什麼過節.
我不跟他們吃飯.
而是要去那個地方我就不去了.
這個試驗.
原來不是普通大家去吃飯.
原來我的上司是一直針對我.
特意耍我.
特意說他一定會去.
你放心.

$^{201}$那些同事說不會.
他說玉貴不會去.
後來真的.
原來是一個打賭.
這個上司就輸了.
這樣的一件事.
是我之前不知道的.
試驗可以是這樣的.
不過弟兄姊妹.
今天我們一起.
每一個弟兄姊妹.
聽完這些堂道之後.
我們要有一個行動.
試驗神.
是神邀請我們去試驗他.
就是說給了一個方向我們.
我們真的這樣做的時候.
他真的會傾福給我們.
這個是一個很大的考驗.
盼望我們聽完這些堂道之後.
我們大家一起.
去進入對神的試驗.
在我們今年特別是年初的時候.
甚至盼望我們大家.
有一個立願.
真的願意這樣去做.
原來是一個試探.
對當時一個.
剛剛初訓的弟兄的一個試探.
後來輾轉還有更大的試探.
就是我的上司.
當時我們開始.
因為屯門醫院就越開越大.
新院就沒那麼多工作.
後來開始不同的部門要輪班.
怎麼輪班呢.
上司就定了.
因為星期日和公眾假期.
就需要有兩組人.
那怎麼輪班呢.

$^{241}$接著他就說我編了.
不可以再改了.
逢星期日就A組.
公眾假期就B組.
那你想到我是拍哪一組了.
全年A組就是你.
B組就是另一位同事.
再加上另一位下屬.
B組是一個主婦.
她就說每次都是公眾假期.
我和我子女去街上玩都不行.
可不可以和我調呢.
我說當然好.
我當然想.
我逢星期日都要回教會.
那時候我是很熱心的.
上司說不改.
不准改.
我已經改了.
不准改.
你可想像得到.
連那些不信的同事都知道.
原來就是這樣繼續玩我.
繼續的瀏覽我.
讓我逢星期日回不了教會.
起碼有一段五個月的時間.
每個星期都回不了教會.
多麼不開心.
多麼激氣.
這樣的情況.
回來又不是太多事做.
就這樣坐在那裡.
突然之間聽到一堂道.
他說你回不了教會.
不是你不回教會.
回不了教會和不回教會是兩回事.
你回不了教會.
你有很多事可以做.
譬如你坐在那裡.
你睜大眼睛看著那些人來光顧.

$^{281}$因為我是做會計的.
很多病人來交院費.
我們就要幫他們結帳.
就要做那些工作.
但是星期日不是很多人出院.
星期是不會有門診的.
所以時間是很多的.
但需要有一個人在工作.
那個時候.
我就想起.
這個時候就是崇拜.
我就拿一支筆拿一張紙.
為崇拜的牧者祈禱.
接著我就為弟兄姊妹祈禱.
原來我坐在那裡.
心靈是在參與現場的崇拜.
當然沒有zoom.
更加不可以這樣做.
但是那支筆那張紙和心靈.
進入了那個敬拜裡面.
維持了五個月.
接著到了星期三.
要補假的.
政府是不補錢的.
補假.
星期三我就可以放下午.
因為下午我就可以回教會.
找牧者聊天.
幫教會做的時候是用錄音帶.
用港島的錄音帶.
那段時間那半年.
讓我有很大的生命建立.
反而應該是.
多謝這位上司針對我.
針對一個基督徒不准回教會.
所以這個試驗.
原來帶來一個很大的.
生命的改變.
盼望弟兄姊妹.
我們進入這一堂道的時候.

$^{321}$你想一想.
你自己今年的困難.
你面對神.
你有什麼可以來試驗他呢?.
你有什麼不妨向這位.
信實不改變的神.
來禱告.
尋求他的帶領呢?.
當然這是一個很簡單的例子.
第二個的聖殿.
神不能夠從敬拜的.
以色列人當中.
得到那個榮耀.
因為他們帶著的.
是很多很多的怨言.
很多很多的問題.
所以你看到.
一章第七第八節.
已經很快就責備以色列人.
你們將核眼的獻為祭物.
這不算為惡嗎?.
將瘸腿有病的獻上.
這不算為惡嗎?.
哇!很直接.
你想想我們回到聖殿.
假如你是一位的祭司.
根據利未忌百姓來獻祭.
首先你要問他是獻贖罪祭.
個人的贖罪祭.
還是群體的贖罪祭.
還是你獻平安祭.
還是你獻繁祭.
是有分別的.
獻的方式是不同的.
但是第一樣.
離遠你見他拖著羊拖著牛來.
或者斑溝或者搓甲來.
第一樣你要問的就是.
有沒有殘疾.
已經看到了.

$^{361}$核眼的.
瘸腿的.
有病的.
這樣來獻上.
問題就是.
祭司照單全收.
你可以想像得到.
當時聖殿的混亂.
完全沒有規矩.
聖殿的敬拜者.
你要不要不要我就拖走.
或者不拖走.
你要就要.
不要就丟下.
祭司因為個人利益的關係.
病的.
瘸腿的.
盲的.
照收照獻.
但是獻給誰.
是獻給這位造物主真神.
耶和華神.
神造.
才來譴責他們的情況.
然後告訴我們.
你們將不潔淨的食物.
獻在我的祭壇上.
卻說我們在何事上.
使你不潔淨呢.
你們說.
耶和華的公祥.
是可藐視的.
可藐視.
其實就是在百姓的心靈裡面.
敬拜可以求求其其.
你不要那麼多嘰噤.
祭司你不要那麼多吵吵.
我們來給你面子.
我們來給神面子.
但是神說.

$^{401}$不可以.
不真正被神閱立的敬拜.
因為你們藐視.
耶和華的公祥.
就是那個祭壇.
就算多麼香火鼎盛.
全部是一些瑕疵.
而且瑕疵背後的心態.
才是真正的問題所在.
根本不是藐視祭司.
是藐視.
聖殿事奉的神.
這位真神.
所以不是代表祭司沒有權威.
沒有那個吉.
是代表百姓整體.
在那個時代.
對神的一種藐視.
這是一個.
先知已經在《至女行間》.
很清楚提出.
亦都很想挽回百姓的心靈.
接著我們再看下去的時候.
萬軍之耶和華說.
還是第一章的經文.
從日出之地到日落之處.
我的名在列國中必尊為大.
在各處人必奉我的名.
笑向獻潔淨的貢物.
因為我的名在列國中必尊為大.
已經提了兩次.
神的心意.
祂的標準.
神希望祂的選民.
成為列國的榮耀.
成為列國不信神的人.
未認識耶和華神的人.
神要被尊為大.
這是神對祂的子民的期望.
神已經很清楚地說了.

$^{441}$所以要在神的裡面.
來去奉祂的名.
笑向要獻潔淨的貢物.
潔淨的貢物.
表示你是真誠真心實意的.
來去敬拜.
造物主整個宇宙萬物都是他創造.
在詩篇也說過.
種生的牛羊都是他所做的.
他根本不稀罕地上的人.
給他有多少的東西.
但是他寶貴他的選民.
在這個聖殿裡面的獻祭.
背後是蘊含著救恩的元素.
蘊含著這個民族.
一直等待真正的彌賽亞.
神的羔羊降生.
所以神自己有他的計劃.
有他的心意是不會改變的.
所以你會看到以色列民.
他們在一個時代一個時代.
雖然回歸了.
雖然已經是第三代的選民.
但是可能礙於當時的社會.
可能待於當時波斯帝國的管治.
他們仍然是一個亡國.
沒有自己地方的民.
但是慢慢忘記了自己真正的國民身份.
就是要有一個敬虔的對神的敬拜.
弟兄姊妹我們一起下去看這個時候.
神要讓自己的名得到榮耀.
那我們敬拜的人可以得到什麼呢?.
我們如何可以在神的敬拜裡面.
得到我們自己的榮耀.
和我們那份平安福樂呢?.
和大家一起去看的時候.
我們回到三章第七到第十二節.
有幾樣東西我們必須要學習.
我們要轉向神.
我們要試驗神.

$^{481}$我們就會獲得神所應許的.
我們成為一群有福的人.
最後十二節說.
這個地方以色列耶路撒冷.
成為一個喜樂之地.
充滿喜樂.
充滿平安.
充滿神同在敬拜帶來的那份心靈的喜樂.
我們開始在新一年的時候.
從你們祖先的日子以來.
你們就偏離了我的律例而不遵守.
現在神就透過先知來勉勵他們.
這是一個很溫柔的勉勵.
不是一個命令的警告.
現在你們要轉向我.
我就轉向你們.
這是萬君之耶和華說的.
轉向.
第十節.
你們要將當立的十分之一全然送入倉庫.
使我家有糧.
以此試試我.
是否為你們闖開天上的窗戶.
傾福於你們.
甚至無處可容.
這是萬君之耶和華說的.
為什麼神要追究他們的奉獻十分之一呢?.
為什麼要這樣試神呢?.
為什麼轉向不是其他東西.
為什麼一定要講金錢物質的奉獻呢?.
為什麼在當時代要這樣追討呢?.
我們又要問一個問題.
新約還有沒有十一奉獻呢?.
新約有沒有提過必須要做十一奉獻呢?.
這是我們信徒必須要問的問題.
不信的人一開始就說.
教會要奉獻十分之一嗎?.
那我不信耶穌了.
但是信主的弟兄姊妹.
我們要問這個問題.

$^{521}$這些經文.
主前的四百多年到今天.
是否還有有效性的延伸呢?.
或者怎樣的延伸呢?.
我們就要再看這個課題了.
仍然是第十二節.
如果我們去試驗神.
如果我們履行了祂對我們的要求.
而我們真心實意去做的時候.
萬國必稱你們這班選民.
在聖殿裡敬拜的這班人.
或者在耶路撒冷稱為神明下的這班子民.
萬國就稱你們是有福了.
因為你們必成為喜樂之地.
這是神的應許.
好像有一個印.
就是萬君之耶和華說.
這裡在整個《馬拉經書》出現很多很多次.
萬君之耶和華說.
就是說不是先知說的.
是神透過先知去說的應許.
所以這是一個帶著很具體的能力的宣告.
好了.
我們要去思想這個.
很重要和我們有關連的.
在奉獻的事情上.
有兩點透過這段經文和新約的提醒.
我們就要很小心了.
律法主義的自以為是.
自以為宜.
或者是另一個心靈的問題就是.
自我的放縱.
有經文支持的.
耶穌基督在路加福音第十一章.
說你們這些法利塞人有禍了.
因為你們將薄荷,芸香和各樣的蔬菜獻上十分之一.
但是有些更重要的事情你們疏忽了.
就是公義愛神的事.
那只是一個類別.
公義可以衍生就是你這個人的公義.

$^{561}$你如何塑造社會的公平.
公義.
譬如在你的家庭.
你如何讓你的家人或者你的傭工.
能夠感受到這份公平,公義呢?.
或者你在人生的處事.
你如何讓社會在你的職場.
來衍生這個公義呢?.
耶穌基督叫我們去.
舉一反三.
而且做了一個比較.
公義和愛神的事.
耶穌基督是當面責備法利塞人.
他們的旨意就是說.
凡是我所有的已經奉獻十分之一.
但是主耶穌在這裡這樣責備.
這裡說這原是你們該做的.
至於其他都不可以忽略.
這裡就是暗示著什麼呢?.
耶穌基督是舊約到新約的人.
聖經的舊約他是相當熟悉.
就是說他一生下來這個世界.
他就知道自己是一個敬虔的猶太人.
他就知道摩西的律法.
那個教導.
他自己一生三十多歲.
他就要遵守十一對神的敬拜奉獻.
一定是金錢物質上的服事.
對神會帶到聖殿.
因為是經由祭司去收取的獻金.
來到聖殿運用.
但是他比我們有更大的靈感.
不只是物質上做了應該做的.
原來神體更重要更大的.
是我們的生命是否符合神的公義.
愛神的事.
他說這些是該做的.
不可以疏忽.
但其他那些相對成為次要.
不是不做十分一.

$^{601}$而是相對公義愛神的事.
那些反而是次要.
花地主人就剛剛本末倒置了.
做給人看.
十分一我做了.
我做到足夠.
薄荷雲香很少的他都計算.
但是更大的心靈的.
和人際的.
在聖殿裡面的.
有身份有地位的.
居然不做這些公義愛神.
招致的是做物主.
耶穌基督對他們公開的譴責.
甚至是一個定罪的審判.
不改變.
花地菜人是要落火.
這樣.
就是說耶穌看到他們的十分一.
不會因為這樣而做足夠.
很讚賞你們.
不是這樣.
原來相對於物質.
心靈裡面對神的敬拜.
神看我們的內心更加重要.
你可以想像到.
主耶穌的眼裡面的十分之一.
不只是物質這麼簡單.
是你整個人對神的認識和敬拜.
生命的舊因大比例的改變.
那個的服侍.
這一點是大家一定要很關注的.
因為是由主耶穌親口說出來的.
各位.
我們用一個小小的例子.
我不知道是真實還是假的.
但是在網上就有這個故事.
有一個年輕人在他公司工作.
美資公司.
聖誕節會有大抽獎.

$^{641}$每人把十元美金放進去.
抽獎箱.
當然不是十元.
是他的名字.
大家寫一個名字.
每人拿十元美金出來.
有三百個僱員.
每人寫一個名字.
一抽到一個.
一個而已.
只有頭獎.
一個人拿了三千元美金.
因為每人十元.
三千元美金.
是很大的獎品.
每個同事都寫自己的名字.
這個年輕人知道.
最近他的一個同事.
是做清潔的事務員.
兒子患病要做手術.
但是沒有錢.
他的醫醫很擔心.
分享.
這個年輕人.
因為每個人要放一個名字.
在抽獎箱.
這個年輕人.
放棄自己的權利.
寫了醫醫的名字.
放在抽獎箱.
他希望三百分之一.
加上醫醫的名字.
即是三百分之二.
有兩個機會.
希望抽到他.
他祈禱有神蹟.
到解餃豬.
一抽.
真的這個醫醫中獎.
開心到三千元美金.

$^{681}$可以解決.
做手術的燃眉之急.
這個年輕人很開心.
因為他有份在三百分之二.
去做這個.
很開心.
開始準備吃飯.
大家聊天玩遊戲等等.
這個年輕人不經意.
沒有人.
走過去挨一下抽獎箱.
試試抽出來.
為什麼我抽出來.
為什麼有這個名字.
原來有很多人.
都寫醫醫的名字.
這是人間的創造神蹟.
已經很奇妙.
如果人肯放棄自己的權利.
但你想想教會.
剛才Cindy和和哥的見證.
當然我們不是.
伴隨著他們成長.
但偶然聽和哥.
其實太太已經回教會很多年.
她每次都是送到門口.
就不回教會.
真的今年明顯看到改變.
所以教會是真正創造神蹟的地方.
弟兄姊妹.
我們來自不同的階層.
我們來自不同的背景.
教會真的讓我們.
在這個大的生命平台裡面.
去經驗到神的奇妙.
神的神蹟.
我們不是刻意營造神蹟.
很自然地當教會尊重神.
我們將最好的獻給神.
神就會在教會裡面.

$^{721}$施行祂自己的作為.
這就是創造奇蹟的地方.
弟兄姊妹.
我們一起要去學習.
所以關乎奉獻.
新約原來保羅有很多教導.
這個教導豐富了我們對所謂十一奉獻的了解.
甚至原來是超越了所謂十一的律法的規定.
我給大家一些經文你就會知道.
在這三方面.
保羅就帶出了.
哥林多前書十六章的第二節.
哥林多前書其中第九章的第七節.
就告訴我們.
每個人都要甘心樂意去做.
神會閱納每一個人.
但是第九章其實更大的下文.
保羅說.
我們作為教導者.
將福音傳給你們哥林多教會.
我們收取養生之物.
根本就是必須.
應該是神安排的.
牛在牆上踩谷.
不可以擋住牠的嘴.
《心靈記》二十五章六節.
保羅用這節聖經講了兩次.
在哥林多前書九章第七節.
和提摩太前書的第六章又再講.
不過這裡告訴我們.
很特別的.
每逢七天的第一天.
即是主日.
每人要照自己的收入.
抽出藥乾.
保留起來.
免得我來的時候現臭.
保羅在建構一個教會常規的運作.
讓這個運作一路一路順暢.
不過你要記住.

$^{761}$教會外面是尼祿做王.
教會從來不受當時羅馬帝國的歡迎.
從來不稀罕你們這些宗教事業.
因為羅馬皇帝是敬拜假神.
他們不會容許基督教的存在.
不過這幾段聖經.
最寶貴的地方就是.
外面多反對基督教.
保羅在建構教會傳福音.
他要建構一個穩定的.
恆常的.
可以有生命建立的平台.
很清楚的教導.
在真道上受教的就是弟兄姊妹.
要把一切美好的東西.
與施教的人分享.
加勒太書六章六節.
什麼意思呢.
教會必須要有一個穩定的講台.
傳遞神的福音.
真理.
信息.
弟兄姊妹要成長.
保羅過去了.
下一代又再出現.
譬如提摩太,提多.
或者其他的人.
可以繼續傳道.
直至主耶穌再回來.
要建構一個穩定的.
可以傳遞真理的地方.
弟兄姊妹回來.
可以聽到神的真道.
被神的真道感動.
清楚自己神給他的召命.
這樣去成長.
這裡告訴我們.
保羅要建構一個穩定的.
不是一步十行的機制.
要讓這個機制.

$^{801}$一直穩定成為教會的常規.
記住外面還是尼祿作王.
很快到提摩太後書.
保羅要殉道.
因為羅馬的皇帝尼祿.
燒了舊的羅馬城.
被當時的元老院來指責.
他嫁禍給基督徒.
說這群人不跟我們敬拜.
我們敬拜的是耶穌.
嫁禍給基督徒.
大肆地去捕獲和捉拿.
就是這樣.
但很奇妙地.
教會保羅試圖清楚知道一件事.
在他們有生之年.
盡快要寫聖經.
盡快讓教會常規化的.
穩定崇拜聚會.
讓弟兄姊妹可以去到一個固定的地方.
敬拜真神.
一個固定的系統.
大家來到慢慢地.
不是潛移默化.
是慢慢地讓神在聚會當中.
不斷地建立不同的信徒.
小的,年長的.
在一個平台裡面.
被神的愛建立.
這裡就是告訴我們.
就如摩西的律法記著.
牛在牆上喘哭的時候.
不可攏住他的嘴.
保羅解釋這一節聖經.
25章6節.
《申命記》.
難道神所掛念的是牛?.
他不全是為我們說的.
的確是為我們說的.
這是哥林多前書.

$^{841}$剛才所引述的經文.
9章9至10節.
但你讀下面更加精彩.
保羅說我不用這樣的權柄.
因為你們哥林多人.
誤解了我們這個使徒的身份.
所以在哥林多教會.
是不要弟兄姊妹的奉獻.
其他教會全部都要.
你可以想像到.
哥林多教會內部混淆了.
教會的存在的目的.
其實是要傳遞神的福音真理.
讓每一個人得到神的救恩.
再一次我們看到.
因為耕種的要傳著子望去耕種.
收割的也要傳著分享穀物的子望.
去收割.
然後是提摩太前書.
5章16至17節.
再一次告訴我們.
教會存在.
外面的社會未做.
教會已經首先去做.
因為這些是公義愛神的事.
信主的婦女.
如果有親戚是寡婦.
就要自己去救濟他們.
不可以拖累教會.
好使教會救濟.
真是無依無靠無助的寡婦.
保羅在提摩太前書.
其實剛剛放監.
他就去到馬其頓.
留下提摩太在爾弗所教會.
一個很大的教會.
保羅就說.
提摩太你要在教會內.
設立一個恆常安排的機制.
就是辭位的機制.

$^{881}$要恆常這樣做.
要有標準.
有些人原來寡婦有家人.
她很慘.
但她的慘都不夠真正無依無靠的寡婦的慘.
因為當時的社會.
是不會有社會福利署的.
是不會有眾援的.
是不會有政府會理會這些窮苦人的.
教會已經第一時間去做了.
教會已經去幫助這些信主的寡婦.
教會必須要成為常規化的.
神的愛公義的平台.
原來教會是公義和愛神的平台.
弟兄姊妹.
我們這樣去學習.
你可以想像到.
教會為什麼要有這樣的分別.
原因就是因為教會是有限的資源.
幫助真正有需要的人.
辦事的人要有智慧去分辨真正無助的寡婦.
你可以想像到.
善於督導教會的長老.
尤其是勤勞講道教導人.
應該得到加倍的敬奉.
因為經常說牛在踩谷的時候.
不可攏住他的嘴.
又說工人得到工資是應當的.
再一次勉勵提摩太.
在教會中心做你自己神給你的職份.
去建立一個完善的.
恆常的一個可以貢獻.
輸送生命真理的平台.
神自然會祝福事奉的人.
這樣建構一個穩定的平台.
教會外面仍然是非常敵視基督教.
保羅現在面對這個年青的傳道者.
其實很害怕的.
在一個大教會裡.
在一個大的都市裡.

$^{921}$有時已不為所人.
是拜那個女神亞底米的.
是很迷信.
是很攻擊基督教的.
但是保羅說.
提摩太你不要擔心.
善於督導教會的長老.
尤其是勤勞教導的人.
應該得到加倍的敬奉.
這是保羅勉勵提摩太.
所以兄弟姐妹.
學習慷慨,甘心,樂義.
還有關乎猜拳宣教.
因為馬其頓和阿該亞人.
樂義奉獻一些捐項.
給耶路撒冷的聖徒中的窮人.
這固然是他們樂義.
其實也算是所欠的債.
因為外邦人既然分享了.
他們靈性上的好處.
當把肉體上的衰庸.
即是錢.
供給他們.
現在耶路撒冷有一個饑荒.
神就透過保羅.
寫信給羅馬書教會的弟兄姊妹.
他說他現在要上耶路撒冷.
他帶著的是什麼呢?.
馬其頓宗教會的一些捐項.
帶上耶路撒冷.
去奉獻給耶路撒冷的猶太人.
不過他加了一件事.
這是基本的責任.
因為外邦人領受了.
本來不屬於外邦人的那份恩典.
那份救恩.
本來神和猶太人立約.
但現在主耶穌成為萬國的救主.
所以馬其頓人就很明白這件事.
他們可能不是錢多錢少的問題.

$^{961}$是意義的問題.
外邦人蒙了神的救恩.
現在聽到耶路撒冷的聖徒有缺陷.
他們願意集齊一些物資金錢.
透過保羅帶上耶路撒冷.
去供應裡面的聖徒的缺乏.
這告訴我們.
既然分享了靈性上的好處.
就應該將肉體的物質的好處.
分享給猶太的聖徒.
咦?.
教會就是一個可以讓福音.
傳給不同民族的平台.
教會就是要這樣.
我們今日香港人.
我們現在很富足.
我們現在有福音.
我們有一位信徒孫教士.
余建峰弟兄.
剛剛在星期日過來分享.
我們今年的猜拳的事工.
我們都會支持他們的事工.
在禱告上支持.
他現在正在中東.
做這些難民工作.
做土耳其一些教會的重建.
和分發聖經的工作.
現在我們將福音.
帶回這些土耳其的地方.
土耳其就是已不鎖的地方.
一個這樣的回饋機制.
不能讓我們有份參與.
所以現在外邦人將這些獻金.
透過保羅.
埃保羅帶上耶路撒冷.
說是我們外邦人奉獻的.
對一個猶太人.
從來看外邦人.
與神的救恩是沒有份的.
但現在不是.

$^{1001}$因為主耶穌的救恩.
大家原來都是神的兒女.
你想想背後的意義是多麼大.
教會需要有這樣的平台.
最後,弟兄姊妹.
你知不知道這段經文.
寶貴之處.
最大的啟發.
已不鎖書.
教導一至三章.
講很多真理.
教義第四章開始.
講生活倫理.
簡單來說就是.
在阻礙主耶穌之後.
舊時的行為.
心思意念.
要改換一生.
要穿上一個新人.
其中有很多例子.
從前偷東西.
不可以再偷.
總要勤勞.
親手做正經事.
正當的事.
這樣才可以把自己有的.
寶來仍然看到一件.
很重要的回饋機制.
原來個人要搭建.
自己的生命平台.
你說很簡單.
我以前偷東西.
經常大洗.
現在不偷.
現在勤奮做事.
不懶惰.
越來越開始有積蓄.
保羅說有積蓄未得.
有積蓄不去幫有需要的人.
不去行公義愛神的事.

$^{1041}$你的累積機制停在那裡.
某程度上這是罪.
偷竊是罪.
肯定是.
偷竊完之後勤力做事也是罪.
不是吧.
勤力做事之後.
原來整個心態.
知道自己是神的兒女.
開始要幫助有需要的人.
這次聖經很特別.
我們從來很少這樣想.
勤力做事.
不要再懶惰.
不要偷東西.
我們不會再教人.
有了積蓄之後.
神給你有一個好的生命.
願意勤奮.
有力量勤奮.
有力量改善自己的生活.
就停在那裡.
那些東西全部是我的.
我的勤奮做回來的.
我已經不再做慣犯偷東西.
我現在很勤力做事.
我現在很勤力儲蓄.
不浪費.
就停在那裡.
原來弟兄姊妹.
在保羅的眼中.
他看十一奉獻.
和公義愛神的事.
是完全拉上一個關係.
所以基督徒今天.
我們要試驗神.
我們必須搭建.
在教會裡一個.
恆常的.
穩定的.

$^{1081}$可以傳道傳福音的一個機制.
我們有敬拜.
我們有講道.
我們要一起去傳福音.
我們有不同的聚會.
讓每一個弟兄姊妹.
來得到神的牧羊.
要越穩定.
越恆常.
越好的.
你就能夠不斷的.
生生不息的.
互相提醒.
互相建立.
但原來不單止這樣.
我們要試驗神.
我們自己就要搭建這個平台.
有沒有看到.
我們的一切是從神而來.
能不能夠從心裡這樣說.
我從神裡得到的祝福.
我現在要想想.
怎樣做公義愛神的事.
就是分級有缺乏的人.
盼望弟兄姊妹不要停留在這裡.
不再偷東西.
勤力做事.
勤力儲蓄.
但是忘記了.
幫助有缺乏的人.
所以教會就是一個.
神安排在地上的一個機制.
我們有詞彙.
我們有猜拳.
我們弟兄姊妹一起去參與.
一些的服事.
譬如今天.
待會兒會宣佈.
最近有弟兄姊妹奉獻了.
一千支牙刷.

$^{1121}$因為現在酒店不准再派.
很多獨立包裝的.
未用過新的牙刷牙膏.
所以農曆年你要去過年.
回去玩.
我們每人可以拿十支.
因為有一千支.
這些機制在教會裡發生.
又分享一點.
大概六月七月左右.
有個弟兄.
我不能夠說名字.
曾牧師.
這裡有一張支票.
你看看教會有什麼需要你們要用.
二十萬.
完全不是我叫他.
你可以想像得到.
弟兄姊妹看到.
他們知道.
教會需要一個穩定的.
傳遞恩典的機制.
公義愛神的事.
主耶穌說.
比十一奉獻更加重要.
但是有了公義和愛神的事.
弟兄姊妹自然會知道.
一切都是從上主那裡而來.
我們按照著.
神給我們的感動去服事.
教會自然會幫助更多更多的人.
我們一起感恩祈禱.
多謝天父給我們剛才.
把新約保羅的教導.
主耶穌的教導.
和在馬拉基書的應許.
仍然到今天都兌現.
我們知道萬物都是從上主那裡而來.
我們能夠生命改變.
我們能夠珍惜自己的工作能力.

$^{1161}$我們可以運用自己的金錢.
一切都是恩典.
一切都是神給我們的生命氣息.
但願在我們新年的開始.
這不是一個勸捐.
而是我們每一個弟兄姊妹都看到.
原來我們可以用金錢去服事神.
去建立一個可以建立生命的平台.
就是教會.
就是能夠讓更多人來到神的家裡.
得到牧羊.
得到生命的建立.
奉耶穌基督的名.
阿們.
\newpage



\section{哥林多前書 10:23-11:1}
\label{sec:gjk_0mIGHss}
\textbf{2025.01.12 《忠於使命,效法基督》 哥林多前書10\_23-11\_1 (和修版) 講員 : 曾敏姑娘}
\newline
\newline
連結: \href{https://youtube.com/watch?v=gjk-0mIGHss}{\texttt{https://youtube.com/watch?v=gjk-0mIGHss}} ~~~~ 語音日期: 2025-01-14
\newline
\newline
\hyperref[sec:0jw5FB1CQ6s]{\small{< < < PREV SERMON < < <}}
~
\hyperref[sec:index_chronic]{\small{[返順時目]}}
~
\hyperref[sec:index_scriptual]{\small{[返順卷目]}}
~
\hyperref[sec:15mYsSTo_Pc]{\small{> > > NEXT SERMON > > >}}
\newline
\newline
哥林多前書 10:23-11:1
\newline
\begin{longtable}{cl}
\hline
\hline
章節 & 經文 (和合本修訂版)\\
\hline
10:23 & \begin{tabularx}{0.7\textwidth}{X} 「凡事都可行」，但不都有益處；「凡事都可行」，但不都造就人。 \end{tabularx} \\ \\ \relax
10:24 & \begin{tabularx}{0.7\textwidth}{X} 無論甚麼人，不要求自己的益處，而要求別人的益處。 \end{tabularx} \\ \\ \relax
10:25 & \begin{tabularx}{0.7\textwidth}{X} 凡市場上所賣的，你們只管吃，不要為良心的緣故問甚麼， \end{tabularx} \\ \\ \relax
10:26 & \begin{tabularx}{0.7\textwidth}{X} 「因為地和其中所充滿的都屬於主」。 \end{tabularx} \\ \\ \relax
10:27 & \begin{tabularx}{0.7\textwidth}{X} 倘若有一個不信的人請你們吃飯，而你們也願意去，凡擺在你們面前的，只管吃，不要為良心的緣故問甚麼。 \end{tabularx} \\ \\ \relax
10:28 & \begin{tabularx}{0.7\textwidth}{X} 若有人對你們說：「這是獻過祭的物」，那麼為了那告訴你們的人，並為了良心的緣故就不吃。 \end{tabularx} \\ \\ \relax
10:29 & \begin{tabularx}{0.7\textwidth}{X} 我說的良心不是你自己的，而是他的。我的自由為甚麼被別人的良心評斷呢？ \end{tabularx} \\ \\ \relax
10:30 & \begin{tabularx}{0.7\textwidth}{X} 我若謝恩而吃，為甚麼因我謝恩的物被人毀謗呢？ \end{tabularx} \\ \\ \relax
10:31 & \begin{tabularx}{0.7\textwidth}{X} 所以，你們或吃或喝，無論做甚麼，都要為榮耀神而做。 \end{tabularx} \\ \\ \relax
10:32 & \begin{tabularx}{0.7\textwidth}{X} 你們不要使猶太人、希臘人，或神教會中的人跌倒； \end{tabularx} \\ \\ \relax
10:33 & \begin{tabularx}{0.7\textwidth}{X} 但要像我一樣，凡事都使眾人喜歡，不求自己的益處，只求眾人的益處，使他們得救。 \end{tabularx} \\ \\ \relax
11:1 & \begin{tabularx}{0.7\textwidth}{X} 你們該效法我，像我效法基督一樣。 \end{tabularx} \\ \\
[1ex]
\hline
\hline
\end{longtable}
$^{1}$頂小妹平安.
在我們一年開始的第二個星期.
我們可以繼續透過神的話語.
去思想神自己的教導.
今天這句話比較像一個什麼.
忠於使命 效法基督.
像什麼.
一個立志.
好像一個立志.
一年開始的時候.
讓我們一同在神面前立志.
是重要的.
我們一同禱告.
要此下你的話語.
成為我們腳前的燈.
路上的光.
由年頭一直去到年尾.
你的話語就成為我們的指引.
求你賜福給我們每一個.
讓我們的心開放.
我們的耳朵很專注地聽.
我們的心回應你.
願寶貴的聖靈親自對我們說話.
讓我們每一個都明白神的心意.
奉耶穌基督的名字祈禱.
阿們.
今天選擇的這段經文.
其實是神秘的一段經文.
裡面有些很寶貴的學習.
特別是農曆新年快到了.
很想藉著這段經文.
彼此有一個勉勵.
首先說回.
這段經文來自哥林多前書.
十章二十三到十一章一節的部分.
哥林多是一個什麼地方呢?.
哥林多其實真的有點像香港.
或者香港像哥林多.
為什麼呢?.
是非常繁華.

$^{41}$很繁榮.
商業的發展很頻密.
但是裡面有些東西是不好的.
就是拜偶像也是很盛行的.
在神廟裡.
除了拜祭之外.
也會有行淫的情況.
所以神廟裡有很多廟祭.
同時神廟裡也會設立燃席.
拜完偶像的人.
也會在裡面一同去吃燃席.
燃席不只是吃一頓飯那麼簡單.
燃席本身跟拜偶像有緊密的聯繫.
這就是哥林多的情況.
哥林多教會裡也有一部分信徒.
是來自原先不是相信耶穌基督的人.
是一些外邦,異邦的人.
本身是拜偶像的.
以前也會到神廟裡拜祭.
參加燃席.
裡面有這樣的信徒.
我們看看這間教會有什麼特色.
看看在哪裡.
哥林多教會是什麼時候建立的呢?.
是保羅在第二次宣教之旅建立的.
這個地方逗留了很長時間.
我們看看那條深色的線.
就是第二次報道之旅.
保羅去的地方.
也去了很多很多地方.
當時交通不是很方便.
但保羅周車勞頓.
為了福音的緣故.
一直去不同的地方.
哥林多在哪裡呢?.
就是在這個地方.
去了很遠.
在這個地方逗留了多久呢?.
一年六個月.
是相當長的時間.

$^{81}$他做了很多福音工作.
他自己一負心血.
在裡面建立教會.
所以可以這樣說.
保羅對於哥林多教會是什麼呢?.
好像媽媽對子女一樣.
那份心是很熱切的.
很關注他們的屬靈生命的情況.
哥林多教會的特色.
是他們有很多恩賜.
如果大家詳細看經文的時候.
會發覺他們的恩賜是不缺乏的.
是非常豐富的.
神給他們的是很好的.
如果說繼續去拓展上帝的國度.
去發展很多的事工.
是綽綽有餘的.
同時你會看到這間教會有很多問題.
真的非常多問題.
所以保羅在哥林多的前書裡.
是逐一逐一地說教會的問題在哪裡.
是要教會去悔改.
問題是很多.
其中一個問題是什麼呢?.
就是有一群人.
他未信主之前會拜偶像.
他會去神廟.
不只是一般的拜拜.
一定會進去參與的筵席.
信了耶穌之後.
是不是就完全悔改了呢?.
又不是.
還有人繼續去神廟.
還會去參與筵席.
你說信了耶穌吧.
是吧?.
就是凡物潔淨.
吃這些東西有什麼問題?.
不是的.
你看看我們今天這段經文.

$^{121}$之前那段經文.
保羅就責備進入神廟裡.
去吃這些筵席的人.
為什麼呢?.
首先這些筵席不是一般的筵席.
它和拜偶像是有關係的.
是上帝非常喜悅的.
而邪靈,鬼就會在這些筵席裡工作.
它會透過偶像.
透過敬拜偶像的過程.
整個筵席裡迷惑人心.
叫人和他們去結連.
所以保羅的指責是非常嚴厲的.
他說這些筵席.
不是去獻祭給上帝.
而是獻祭給鬼.
你去吃的筵席.
是什麼呢?.
和鬼一起有筵席.
信耶穌的人.
我們的生命.
有聖靈住在我們裡面.
我們和誰有筵席呢?.
和上帝有筵席.
我們剛剛進行了筵席.
是什麼呢?.
聖餐.
聖餐是主的筵席.
我們的生命.
是分別為聖.
是屬於耶穌基督的.
我們和主有筵席.
所以經文裡說到.
神是很憎惡我們.
一方面說屬於我們.
想和我們有筵席.
另一方面就大祭這些鬼.
和鬼有筵席.
邪靈的鬼.
這些骯髒的.

$^{161}$上帝是聖潔的.
最憎恨人是骯髒的.
所以神是非常不喜悅的.
這一群信徒.
信了耶穌.
還要回去.
不改變.
還要去吃他的筵席.
神真的非常不喜悅.
我今天很想說一件事.
我都知道.
香港是一個很豐富的城市.
有各種不同的飲食.
除了一大類.
和肉食有關.
有不同的肉食的筵席之外.
有一類是素食的.
吃素食的有不同的.
你去買回來吃.
有些是怎樣的.
直接去哪裡吃.
進廟裡.
進佛寺裡吃.
你說沒問題.
我信耶穌的.
百毒不侵.
有什麼問題.
進去吃.
我想和你說.
不要進去吃.
第一.
這些地方是很多邪靈鬼的地方.
你將自己.
處於一個什麼樣的地步.
放在一個試探當中.
在污穢當中.
上帝一定不喜悅.
是上帝一定不喜悅.
你說我不是拜.
但那些賢侄.

$^{201}$擺在裡面的賢侄.
就是為了這樣的目的.
你有沒有問過神.
神怎麼看你.
你帶著什麼身份進去.
那一刻我就不做基督徒了.
有這樣的嗎.
信了耶穌就完全屬於上帝.
上帝是用他愛子耶穌基督的生命.
買熟我們.
不能那一刻就說.
忽然間我不做基督徒了.
我不信耶穌.
那一刻我先吃完.
吃完再信.
沒有這樣的.
信了那一刻.
你完全屬於耶穌基督.
是聖結的.
你需要和這些分割出來.
不要讓自己陷入污穢當中.
所以我會反對.
我真的反對人.
進入寺廟裡.
進入佛寺裡.
我們吃吃齋.
真的很好吃.
身體好.
那些齋真的做得比較完全.
有不同的款式.
有什麼款式都是假的.
撒旦也會炮製很豐富的賢侄.
你分清楚誰炮製的賢侄.
那些賢侄多豐富都不去.
你知道吃了會有什麼後果.
為上帝懸故.
一定不去.
在我們今天這段經文之前.
就說到這一點.
在過程中.

$^{241}$我們看到.
其實人是很喜歡自由的.
不只是當時哥林多的人崇尚自由.
其實我們香港人都很強調自由.
我們要強調.
我們什麼都是自主.
我們要自己作主.
我們想怎樣就怎樣.
我們事實上都是很自由的.
所以經文裡說.
凡事都可行.
你想做什麼都好.
其實都沒有人攔阻你.
但是.
但是寶萊這裡的教導是什麼.
寶萊直接下了一個命令.
是一個命令.
不是一個勸勉這麼簡單.
不要求自己的益處.
而是要求別人的益處.
為什麼.
弟兄姊妹.
今天我們不是單獨地生活.
你看看前後左右.
前後左右是誰來的.
你是不是單獨地在這個世上活著.
你不是.
每個星期日回來這裡聚會的時候.
我們公開做一個見證.
整群信耶穌的人做一個見證.
我們是屬於上帝的.
不是單獨地做一個見證.
是一齊.
整個群體.
你有弟兄姊妹.
你是互為之體的.
彼此之間.
所以我們會影響別人.
別人都會影響我們.
這個也不是一個普通的群體.

$^{281}$是屬上帝的群體.
而且是一個愛的群體.
為愛的緣故.
不要求自己的益處.
而是要求別人的益處.
這裡說什麼呢.
接觸上一段經文.
說不要進廟裡吃那些筵席.
但是哥林多整個城市.
因為拜偶像太厲害了.
事實上他們預備的食物.
就算是街邊的食物.
很多食物都是拜過偶像的.
那你是不是真的不吃呢.
保羅就要回應這個問題.
吃不吃呢.
是不是所有東西都不吃.
那就不用吃了.
每樣都問過.
這個有沒有拜過偶像.
拜過我不買了.
這個有沒有拜過偶像.
拜過偶像我就不吃了.
那不用吃了.
保羅回應了.
在不影響別人的情況下.
不帶來負面影響的情況下.
是可以吃的.
例如有些情況是說.
市場上所賣的.
其實不用問那麼多.
市場上所賣的.
不需要再問.
那你是不是有拜過偶像.
不用的.
也都甚至有不信的人.
請你去吃飯.
那些食物都很大機會.
是拜過偶像的.
他請你去吃飯.

$^{321}$那你又先問了他.
先吃或者不吃.
不會 照樣去吃.
為什麼呢.
有一個原則.
在這些情況下.
原則是什麼呢.
剛才有一首詩歌.
是什麼呢.
這是天賦世界.
這是天賦的世界.
所有東西都是天賦的.
詩篇二十四篇一節.
都是這樣說.
地和其中所充滿的都屬於我.
一般來說.
那些食物.
你都不需要問什麼.
你就可以去吃.
是不是因為所有食物.
都是屬於神的.
食物本身是沒有什麼.
那你就放膽去吃.
不會因為這樣.
又戰戰兢兢.
有些東西又怕.
那些又怕.
去別人家裡又想來想去.
吃不吃.
在街上買又想吃不吃.
不需要擔心.
就去吃.
因為這些是屬於上帝的.
就去吃就行了.
但是有一種情況.
擺明了.
有人跟你說.
這個是獻過祭的.
那個人為什麼要這樣說.
試探你.

$^{361}$很多時候.
如果有人這樣說.
這個是獻過祭.
對你來說.
對方知道你是信耶穌的人.
對方知道你是信耶穌.
有些人.
他真的信心軟弱.
他知道你是信耶穌.
他覺得.
他好心要提一提你.
我知道你們信耶穌的人.
是不會拜偶像的.
我提一提你.
這些食物是獻過祭的.
在他心目中.
他軟弱.
他覺得獻過祭的食物.
你吃了就好像拜偶像一樣.
所以我提一提你們.
信耶穌的.
這些就不要去吃了.
甚至有些軟弱的信徒.
真的這樣想.
這些食物比偶像還要大.
不好吃.
好像一番好意.
要幫你.
甚至好像救你.
如果有些人這樣說.
說到明.
那吃不吃.
就不吃.
不是食物本身是怎樣.
而是那個人的良心.
那個人的良心.
他提你.
他覺得基督徒.
是不會拜偶像的.
我提一提你.

$^{401}$你吃這些.
就好像拜偶像一樣.
所以提一提你.
食物有這樣的情況.
在這個時候.
為那個人.
為他良心的緣故.
就不要吃.
免得他跌倒.
免得他覺得.
基督徒都拜偶像.
在他心目中.
可能他都覺得.
吃這些偶像之物.
就是拜偶像.
基督徒吃.
你也是拜偶像.
那我都可以.
軟弱的基督徒會這樣想.
未信主的人.
將來都會這樣想.
他跌倒了.
那就不好.
為他緣故不要吃.
所以一個行動.
你選擇進行一個行動前.
你要想一想.
不是我自己喜不喜歡.
我有沒有自由.
你一定有這個自由.
或者你都很喜歡做.
但你想一想.
會不會令人跌倒.
如果令人跌倒.
就不要做.
在這裡.
來到一個很重要的位置.
去說.
保羅說.
所以你們或吃或喝.

$^{441}$無論做什麼.
都要為榮耀神而做.
不只是吃喝.
是你做什麼都是.
問問自己.
這件事榮不榮耀神.
神是否得著讚美.
人是否尊崇上帝.
如果我做完這件事.
人們會覺得.
那些信耶穌的人都這樣做.
都不要跟他的神.
他的神真的不太好.
不太行動.
不要跟他.
還說帶我回教會.
我都不想跟他回.
這樣就不好.
這樣就真的不好.
這樣就不要做.
因為不榮耀神.
今天有一首詩歌.
真的很配合.
我記得Don給我的詩歌.
我沒有想過自己去說的時候.
我覺得很配合.
活著為耶穌.
活著不只是為自己.
是為耶穌.
單單為主活.
為主耶穌而活.
我們要為榮耀神而活.
為榮耀神而做每一件事都是.
吃喝是.
其他事都是.
你問問你自己.
問問我們自己.
今天生活裡.
是否每一件事都榮耀神.
我們買東西.

$^{481}$是否榮耀神.
大食會.
農曆新年有很多機會吃.
喝喝吃吃.
去哪裡玩.
去哪裡參觀.
看什麼.
是否榮耀神.
我今天坐電梯.
我住那裡.
我住樓宇.
我坐電梯到樓下的時候.
我看清楚.
我們住的地方.
他們都有組團.
農曆新年去玩.
通常農曆新年.
團都是喝喝吃吃.
我看見很多齋菜.
每次去都是廟.
我去到就顯明.
信心的剛強.
我的堅定.
齋菜那麼多款式.
但你想想你的見證.
屬於耶穌的人.
農曆新年全部去廟參觀.
逐個廟去.
為了什麼.
有些人說.
當然了.
不去是為了認識他們的情況.
你認識完之後.
是否打算傳福音.
將整個廟裡.
參拜的人.
全部救他們出火海.
如果是我為你感恩.
為你禱告.
但說真的.

$^{521}$大部分人都不是.
想吃而已.
就去參觀.
但當別人都知道你這樣的時候.
你說的是什麼見證呢.
是否很榮耀神呢.
我講得極端一點.
譬如說我自己.
有一天.
曾姑娘.
知不知道曾姑娘去哪了.
她參加他們住宅旅行團.
逢廟必入.
哪裡都去看.
拍照都在廟上.
在偶像面前.
你覺得曾姑娘怎樣.
你印象是否很榮耀上帝.
真的很聖潔.
是不是.
我見到大家在笑.
都知道答案了.
當然不要.
但這些不榮耀神的事.
我們不要去做.
真的要想榮耀神.
才去做.
所以是選擇.
要選擇過.
不是橫衝直撞去做.
我們去挖穴.
挖磕.
做任何事.
不是為了滿足自己的慾望.
哇.
人的慾望是無窮的.
中午吃一個很大型的自助餐.
晚上再來一個大型的餐.
第二天早上再吃巨型的餐.
一天三餐都吃滿.

$^{561}$真的很快樂.
很開心.
很榮耀.
真的很快樂.
但真的要問自己.
這個吃穴榮耀什麼.
我做的事是否榮耀.
還有另一種情況.
我不是滿足慾望.
有一種是.
榮耀自己.
我做很多很多事.
我告訴你.
我真的很好.
數百男.
一二三四五.
八成數.
真的很榮耀.
是否為自己.
那首詩歌.
是說什麼.
活著為耶穌.
經文是說什麼.
為了榮耀神.
不是為了榮耀自己.
做這麼多事都是榮耀神.
讓我們的上帝得榮耀.
我們的上帝真的配得.
有誰像我們的上帝.
創天造地.
以大能統管萬有.
有誰像我們的上帝.
這麼愛萬物.
思憐民.
有誰像我們的上帝.
派他的兒子來救贖我們.
沒有.
所以上帝配得.
一切榮耀頌讚.
所以我們活著做任何事.

$^{601}$包括我們的事奉.
都是為了榮耀神.
有時候我們很有趣.
我們嘴裡說榮耀神.
心裡總是有動機.
總是很想.
你有沒有留意我.
你可不可以讚我.
總是很想這樣.
求神光照我們.
讓我們看到我們的動機.
我們的動機.
我們整個人都是為了榮耀神.
不是為了榮耀自己.
保羅在這裡說.
凡事.
什麼都要使眾人喜歡.
不求自己的益處.
只求眾人的益處.
使他們得救.
這就是使命.
得救.
得著福音的好處.
保羅就是一個板.
剛才我這麼簡單.
很快就給大家聽完.
我再說一次.
保羅是一個板.
我再簡單.
很快就給大家聽完.
第二次宣教之旅.
我沒有說完整個行程.
宣教之旅就是從這裡到那裡.
做福音的工作.
當中經歷苦難.
被人逼白.
傷害.
被石頭打.
試過差不多死.
付了很大的代價.

$^{641}$他就是這樣不求自己的好處.
什麼是求自己的好處.
來來來.
不用這麼拼.
我們在一個地方.
串一會福音.
到處觀光.
吃吃有什麼好吃的.
玩完一大輪住下好的酒店.
開心.
大家彼此聯誼.
真的很舒服.
住夠了.
去下一個城市串兩下.
做做.
如是者飲飲食食十幾個城市.
真的很快樂.
那些都有串福音.
但也有很多求自己的益處.
保羅做到很盡.
真的不求自己的益處.
不浪費時間.
為福音緣故.
就是這樣做.
令他們得救.
這是很重要的.
在這裡保羅說.
你們該效法我.
效法基督一樣.
保羅.
你也很驕傲.
你是不是很驕傲.
叫人效法你.
你覺得自己很好.
不是.
效法我好像什麼.
為什麼要效法我.
因為我效法基督.
最終是要效法基督.
我想給大家看看.

$^{681}$保羅.
為福音的緣故.
如何放下自己的自由.
為福音的緣故.
為別人的好處.
大家聽我讀這段經文.
在哥林多前書.
九章十九到二十三節.
保羅這樣說.
我雖然是自由的.
不受人管轄的.
但我甘心作了眾人的僕人.
是僕人.
服侍人.
卑微的謙卑的.
為贏得更多的人.
對猶太人.
我作猶太人.
為贏得猶太人.
對律法以下的人.
我雖不在律法以下.
還是作律法以下的人.
為贏得律法以下的人.
對沒有律法的人.
我作沒有律法的人.
為贏得沒有律法的人.
其實我在神面前不是沒有律法.
對軟弱的人.
我作軟弱的人.
他很苦就卑微的人.
為贏得軟弱的人.
一直說.
無論如何.
我總要救一些人.
整天想著.
有靈魂未得救.
我就去傳福音.
弟兄姊妹.
問問大家.
由這個年頭到年尾.

$^{721}$我們傳了多少次福音.
其實上一次.
我參加教會的.
感恩分享會.
我很感動.
聽阿謙的分享.
他這幾個月.
傳福音.
我心裡很感動.
我在他身上看到上帝的工作.
很奇妙.
我很感動.
神會感動我們.
神會在我們裡面工作.
叫我們愛靈魂的得救.
不是每個人都放下工作.
去做傳道人.
在你的職業裡可以服侍上帝.
在你的生活裡.
可以服侍上帝.
每個人家裡.
都有不同的家人.
有不同的親戚.
鄰居 朋友 同學.
我們每個人都很成熟.
我所說的年紀.
我們身邊一定有很多人.
有多少人信了耶穌.
神讓你在圈子裡.
信耶穌的人.
你有多想念他們.
有多用力.
去帶他們信耶穌.
這是很重要.
你有沒有想著他們.
上帝很愛他們每一個.
你有多少付上代價.
從年頭到年中.
做了多少福音的工作.
這裡每個人都可以成為上帝的精兵.

$^{761}$每個人都可以成為上帝的精兵.
每個人都可以傳福音.
坐在這裡都可以傳福音.
我經常說.
紅印堂的年紀比較成熟.
年長一點.
就算公公婆婆.
去社區中心做運動.
坐在這裡都可以傳福音.
拉著東西做運動.
隔壁有阿嬸.
有阿婆.
都可以說.
有沒有去教會.
都可以說.
很開心.
不用去街上.
拿個單張.
傳福音.
買菜都可以.
你會接觸到人.
其實我知道.
我以前帶開恩光團的時候.
當中.
都有長者.
很受歡迎.
認識很多街坊朋友.
我都試過跟他們.
在村裡跟街坊聊天.
我心想你認識很多人.
不用的.
認識人就可以說.
認識人就可以跟他們傳.
你為他們的好處.
是為他們的好處.
不是只是為自己.
你看保羅就算怎樣.
總要救一些人.
凡我所做的.
都是為福音的緣故.

$^{801}$為要與人共享福音的好處.
為了讓人得到福音.
我們自己的見證.
很重要.
是要造就別人.
如果今天我的見證.
是失見證.
你不用再說了.
是不是.
買菜.
我都要跟他吵幾場架.
然後.
中間可能爆一些坐雨刺.
對方一看.
這是來自哪裡.
來自教會.
教會的人是這樣的.
態度是這樣的.
很難再傳福音.
為他的好處.
為他們的好處.
做好見證.
所以做任何事.
生活的形態.
我是不是對別人有益處.
建立別人的生命.
深深的自己.
真是為福音緣故.
這是很好的.
保羅在這裡說.
你們不知道在運動場上賽跑.
大家都跑.
但得獎賞的只有一人.
你們都是要這樣跑.
很盡全力地跑.
好使你們得獎賞.
凡參加競賽.
在各方面都要有節制.
我有自由.
但我不會亂用.

$^{841}$我的自由.
我有節制.
是用得恰到好處.
那些人在地面上參加比賽.
是要得會扭壞的冠冕.
但我們是要得不會扭壞的冠冕.
所以我們奔跑.
不像沒有目標.
有目標的.
為了上帝的國度.
為神.
也為了榮耀神.
凡事做什麼都是為了榮耀神.
要讓人得救.
榮耀神.
這是我們的使命.
在過程中.
保羅說我克制己身.
使他完全信服.
免得全福音級別人.
自己反而被淘汰.
小心.
我在濫用自由.
不要以為純粹影響全福音.
我自己沒問題.
不是.
不要自己成為被淘汰的那一個.
在此保羅很重要.
他說.
效法我.
就好像我效法基督一樣.
最重要.
是要效法基督.
基督我們.
實在有太多經文.
很寶貴.
但我看回這一段.
很想和大家看.
在《路加福音》22章42節.
是一個什麼經文呢.

$^{881}$《路加福音》22章42節.
說.
耶穌說.
父啊!你若願意.
求你將這杯撤去.
然而不是照我的意願.
而是要成全你的旨意.
什麼時候說的呢?.
赫西瑪利元.
他真是很大的掙扎.
快要去經歷.
被出賣.
被捉拿.
整個受苦的過程.
在赫西瑪利元這裡.
他幾次去禱告天父.
在這裡.
他說.
真是一個很大的苦杯.
只有他一個可以喝.
他說如果你願意.
就將這杯撤去.
但不是照我的意願.
而是要成全你的旨意.
耶穌基督選擇.
順服天父旨意.
喝了這個苦杯.
接著的後果.
一定是.
就是經歷被出賣.
被捉拿審判.
被鞭傷.
釘上十字架.
受很大的苦.
是他選擇的.
他是這樣用他的自由.
去選擇.
去受苦.
為成全別人.
為了全世界.

$^{921}$人的好處.
如果你說.
為自己體貼自己的肉體.
我就不會去犧牲.
我不會去犧牲.
受很多苦.
不會去犧牲.
不會被人出賣.
我不要去.
受鞭傷.
被人嘲笑.
我都不會去受.
為自己肉體的緣故.
為自己的好處.
但如果這樣為自己.
全世界的人.
就在永遠的死亡裡.
就是耶穌基督看到.
他很愛我們.
為了拯救我們生命.
他選擇犧牲.
這一個.
就是保羅效法的榜樣.
保羅叫我們效法他.
由他效法基督.
今天來到這裡.
想問問大家.
我們的使命是甚麼.
何欣嘉在這裡.
剛剛過了.
十三週年的時間.
在過去雖然有很多風浪.
但仍然在傳福音.
我相信將來.
我們會進行很多活動.
我們會有很多活動.
或有很多安排.
都是為了福音的緣故.
為人得好處.
不單止這樣.

$^{961}$我們在座的弟兄姊妹.
今天我們活著.
每個人活著.
都應該為了福音的緣故.
為別人得好處.
因為我們自己.
就是這樣得到好處.
我們亦需要.
讓別人得好處.
我們跟隨的是.
跟隨耶穌基督.
所以是效法基督.
我們的使命是甚麼.
傳福音.
榮耀神.
今年.
我很希望成為我們每一個.
其中的立志.
是嗎.
是要忠於使命.
不是三分鐘熱度.
是呀.
年初的時候說說.
傳福音到年尾的時候.
我都不知道去了哪裡.
沒有了.
是否效法基督.
不是了.
不知道效法了誰.
不是.
如果真的立志的時候.
我由年初到年尾.
我立志就是這樣.
我忠於使命.
由年初到年尾.
一直到年尾.
都是這樣.
跟隨耶穌基督.
效法他.
不要看著別處.

$^{1001}$農曆新年有很多假期.
期盼大家.
在神面前.
向上帝傾心.
禱意.
求上帝幫助我們.
我們一同禱告.
天父上帝.
我們獻上感恩.
你很愛我們.
此下你的話語.
要成為我們的幫助.
其實我們身邊真的有很多見證人.
好像雲彩一樣圍著我們.
這些見證人為了福音的緣故.
為了榮耀上帝你的緣故.
擺上自己的生命.
擺上自己的一生在所不惜.
而主耶穌基督.
你更是這一切榜樣中的榜樣.
你為了愛我們.
拯救我們生命.
你付上了生命的代價.
在你第一次來到世界的時候.
你都不是為自己著想.
你從來都是為我們.
主耶穌你佩得我們的愛.
求你讓我們跟你的腳踏中行.
讓我們回應你.
讓我們為福音的緣故.
擺上自己的一切.
讓我們忠心地跟從你.
奉耶穌基督的名字祈禱.
阿們.
謝謝大家.
\newpage



\section{約書亞記 1:12-18}
\label{sec:15mYsSTo_Pc}
\textbf{2025.01.19 《不做局外人》 約書亞記1\_12-18 (和修版) 講員 : 鄧泰康傳道}
\newline
\newline
連結: \href{https://youtube.com/watch?v=15mYsSTo-Pc}{\texttt{https://youtube.com/watch?v=15mYsSTo-Pc}} ~~~~ 語音日期: 2025-01-19
\newline
\newline
\hyperref[sec:gjk_0mIGHss]{\small{< < < PREV SERMON < < <}}
~
\hyperref[sec:index_chronic]{\small{[返順時目]}}
~
\hyperref[sec:index_scriptual]{\small{[返順卷目]}}
~
\hyperref[sec:8UN_bXXTUSo]{\small{> > > NEXT SERMON > > >}}
\newline
\newline
約書亞記 1:12-18
\newline
\begin{longtable}{cl}
\hline
\hline
章節 & 經文 (和合本修訂版)\\
\hline
1:12 & \begin{tabularx}{0.7\textwidth}{X} 約書亞對呂便人、迦得人和瑪拿西半支派的人說： \end{tabularx} \\ \\ \relax
1:13 & \begin{tabularx}{0.7\textwidth}{X} 「你們要記得耶和華的僕人摩西所吩咐你們的話說：『耶和華－你們的神使你們得享安寧，必將這地賜給你們。』 \end{tabularx} \\ \\ \relax
1:14 & \begin{tabularx}{0.7\textwidth}{X} 你們的妻子、孩子和牲畜可以留在約旦河東、摩西所給你們的地。但你們中間所有大能的勇士都要帶著兵器，在你們的弟兄前面過去，你們要幫助他們。 \end{tabularx} \\ \\ \relax
1:15 & \begin{tabularx}{0.7\textwidth}{X} 等到耶和華使你們的弟兄和你們一樣得享平靜，並且得著耶和華－你們神所賜他們為業之地的時候，你們才可以回到你們所得之地，承受為業，就是耶和華的僕人摩西在約旦河東、向日出的方向所給你們的地。」 \end{tabularx} \\ \\ \relax
1:16 & \begin{tabularx}{0.7\textwidth}{X} 他們回答約書亞說：「凡你吩咐我們的，我們都必做；凡你差我們去的地方，我們都必去。 \end{tabularx} \\ \\ \relax
1:17 & \begin{tabularx}{0.7\textwidth}{X} 我們在一切事上怎樣聽從摩西，也必照樣聽從你。惟願耶和華－你的神與你同在，像與摩西同在一樣。 \end{tabularx} \\ \\ \relax
1:18 & \begin{tabularx}{0.7\textwidth}{X} 無論甚麼人違背你的命令，不聽從你所吩咐他的一切話，就必處死。你只要剛強壯膽！」 \end{tabularx} \\ \\
[1ex]
\hline
\hline
\end{longtable}
$^{1}$我想問一下剛才Ken的太太是誰?.
哦 難怪你今天唱歌唱得這麼開心.
應該最開心的是你太太.
因為看到身邊的人的轉變其實是一種很大的興奮.
剛才Ken分享的時候令我想起一個片段.
就是我媽媽 為什麼呢?.
因為我家裡也是這樣 從小到大他們都去拜偶像.
他們家人每年都接醫治 做很多事情.
到了年老之後 偶像就沒有機會再做了 為什麼呢?.
因為我們家人一個一個信著.
以前由我們開始 誰誰誰負責葬香 你們明白的.
信了耶穌之後 誰誰誰做 阿妹做.
全家都信了主之後 沒人再去葬了.
我媽媽又老了 她又不能再去葬了.
所以就放在神桌上.
很多年之後我們問她.
媽媽 說真的 你又不再去葬香了 你放在那裡做什麼?.
她說 我不能去葬 放在那裡吧 別理她吧.
我想說我媽媽的感覺是什麼呢?.
她不敢去 她又不相信.
最奇妙的就是.
我說不如我們祈禱吧 如果有機會我們就去吧.
有一次我媽媽回鄉的時候.
突然有一天神桌整個斜了下來.
像是在叫什麼 比薩斜塔一樣.
我弟弟馬上通知我妹妹.
去和我媽媽在鄉下.
神桌斜了下來 我媽媽就扔了它.
所以我想說很多時候是一種心結.
所以上帝會做工作的 很奇妙的.
大家都可以一起去為家裡有需要的禱告.
今天和大家分享這段經文.
也是一本我很喜歡的書.
我被蒙召去做傳道人.
也是這本書裡有一節很有名的金句.
約書亞記第一章第九節.
你們不要懼怕 也不要驚慌.
因為無論往哪裡去 耶和華的神都必與你同在.
整個約書亞記其實是整本書一本很特別的書.
為什麼呢 因為整本書裡只有約書亞記是一路勝利的.

$^{41}$基本上在裡面 約書亞記只曾經有一次很細微的錯誤.
而令他們輸給艾成.
因為他們沒有依靠上帝 他們靠自己.
所以後來輸了 而之後所有的戰役全部都是贏的.
所以這本書特別的地方是什麼呢.
它是奠定了一些屬於上帝的軍隊.
我形容為屬於上帝的軍隊.
當你要走上帝的路的時候.
一些很重要的基礎 特別第一章.
如果大家有機會可以回去看.
第一章 由第一節到第九節.
其實當時是在說上帝再一次呼召約書亞.
呼召約書亞接替摩西成為大領者.
但問題是你會發覺在第一節到第九節有一個特色.
就是在第六 第七 第九節的經文裡.
不斷重複著一個主題.
就是你不要懼怕 你不要驚慌.
不斷說 其實我就是想說.
什麼人才會不斷跟他說你不要怕.
就是你很怕.
其實我想說約書亞在他的背景裡.
沒錯 我很相信上帝.
我很願意跟從上帝.
但不是代表他不怕.
問題是他以前一直是一個副手.
他不是大領的那個.
如果大家留意約書亞的特色是什麼呢.
他很馴服.
他馴服是什麼呢.
他馴服就是經過摩西馴服.
就是說摩西叫我做什麼我就做什麼.
他習慣了做這個角色.
也做得非常好.
如果大家記得他們那次打亞馬利人的時候.
記得摩西舉高之上舉到手軟.
其中一個就是約書亞.
你過來幫我托住手.
他就是很忠心地完成了摩西叫他做的事之後.
你會發覺那件事就成就.
透過摩西去看迦南的時候.

$^{81}$他也是其中一個很相信上帝的人.
我想強調他的特性就是.
他很忠於上帝所給他的託付.
他對上帝很有信心.
但他很怕.
所以第一至第九節的時候不斷重覆.
上帝必與你同在.
如果你是上帝的軍隊的話.
你必須要首先確定.
上帝一定與你同在.
今天我們做基督徒.
其中一個很重要的地方就是.
你要記住你要做得很有信心.
你要做得非常有勇氣.
這樣才能夠告訴別人.
我誇張一點.
你不要怪鄧全.
我以前是做話劇的.
所以有時候我講的會比較誇張一點.
譬如很簡單.
你找到一顆糖.
很好吃.
然後你很想介紹給你的朋友吃.
你的表情是怎樣的.
很好吃啊.
很好吃.
吃到別人流口水都很想吃那顆糖.
然後你回味了.
很好吃啊.
那顆糖很好吃.
你試一下.
你吃不吃啊.
你不會的.
我想強調的是.
我們做基督徒.
首先我們要很清楚.
我們的上帝你要確信.
上帝是你唯一的救贖.
上帝是你唯一的出路.
你要有這種勇氣.

$^{121}$才能夠介紹給別人聽.
這個上帝是值得去信的.
所以上帝在這裡跟約書亞說.
我一定與你同在.
所以在這方面他是沒有問題的.
你會發現很多人忽略了下面.
剛才大家讀的那段經文.
其實很重要的.
為什麼呢.
第一章的時間.
這裡才是一個整體.
有上帝的同在.
但同樣需要人的合一和同心.
今天在教會裡面的時間.
我可以完全夠膽告訴大家.
教會裡面從來出的問題.
都不是上帝不在.
從來都不是上帝沒有能力.
從來都是人的問題.
聽清楚了嗎.
所以你會發現.
剛才這章的經文.
你記住他們預備打仗.
他們要挑戰他們從前很害怕的迦南人.
接著這班人一起打仗的時候.
必須要整個群體是一條心.
所以你會發現他們首先.
約書亞他自己意識到一件很重要的事情.
就是這兩個半支派.
其實在約旦河東已經得到土地.
說真的.
我一直按照上帝的應許.
其實得到什麼呢.
就是得到土地.
上帝賜給我可以安心的地方.
而這兩個半支派已經得到.
按邏輯的話.
我可以收工.
明白我意思嗎.
在跟隨上帝的裡面.

$^{161}$我要得到的我都得到了.
我就收工.
但你會發現約書亞很意識到一件事.
你記住永遠以色列12支派都是一體的.
而最重要的就是當日.
這個土地不是靠這兩個半支派打回來的.
是整個12支派一起幫忙打回來的.
而只不過後來上帝按照他的心意.
把約旦河東的地方賜給這兩個半支派.
就是說這群人拿了這個福分.
是因為有其他一些.
他不介意.
他也很忠心.
跟你一起去努力的人.
你才能夠得到這個福分.
所以為什麼他要他們一起在當中.
就是因為你們要記住.
你們要記住是有人幫過你們的.
現在你們還有些兄弟是還沒得到的.
所以你們不要停下來.
你們要繼續為他們的緣故努力.
你們要繼續一起打.
裡面出現什麼問題.
我想大家想想.
人性裡面很簡單.
我以前跟你們一起去打.
現在你們拿到那個土地.
你不理我嗎.
你明白我意思嗎.
裡面就出現了這樣的情況.
會很不舒服.
會帶來未來很多隱患.
所以你會發覺.
約書亞在剛才那段經文裡面.
很清楚跟他們說.
你們的家人.
你們的生畜.
是可以留在約旦河東.
沒問題的.
不過你們所有能夠打仗的大能勇士.

$^{201}$請你跟他們一起去打.
直到所有的人都得到平安為止.
對不對.
這個很重要.
弟兄姊妹.
過去你們能夠得到這個福音.
同樣也有一群人.
他們曾經努力過.
所以你們才得到這個福音.
有些是信義代.
是因為爸爸媽媽.
他們從小到大帶你們回來教會.
讓你們去親近上帝.
去認識上帝.
你們才能夠得到.
有些人很直接.
有人傳福音給你們.
就像剛才Ken所說.
有人為你們努力.
曾經跟你一起研究.
跟你去做初訓班.
跟你去做浸禮班.
有人幫你的.
有些人很厲害.
不是的.
我自己中進來的.
我也告訴弟兄姊妹.
這裡是爆出來的嗎.
不是的.
這裡也是有一群人很努力.
去建立了這個教會.
開始了這個地方.
你才能夠有教會來到當中.
每一個人其實在過去.
都有一群人很努力.
不計較.
我不會知道這個地方.
能夠祝福什麼人.
然後他們就將這個福音留給你.
而你得到.

$^{241}$同樣道理的.
就是今天其實同樣.
你得到這個祝福.
你得到祝福之後.
洗完禮了.
然後怎樣.
搞定.
收工.
不是的.
你要記住原來我們今天.
在我們這個福音裡面.
是永遠都要有一個.
生生不息的循環.
在感受堂.
大家如果還記得.
因為銀銀堂是感受堂分出來的.
感受堂在最初.
大家知不知道是怎樣開始的.
有沒有聽過.
其實是由西教士.
帶著一些弟兄姊妹.
他們在家裡有祈禱會開始.
而你記住西教士.
是原本在當時比較富庶的國家.
他們沒有安於他們的舒服.
而他們願意毀身來到這個地方.
去為這個地方.
相對當時落後的地方.
而去做這份工作.
所以就建立了感受堂.
你就會發覺.
同樣道理.
就是後來有一群弟兄姊妹一起起來.
然後就建立了這個教會.
我想告訴弟兄姊妹.
如果你只需要在教會當中.
有一刻突然有一種很安恕的感覺.
就是我得到永生就行了.
其他事情都不關我的事.
我想告訴弟兄姊妹.

$^{281}$你們這個教會.
就一定會停下來.
其實我們每一個人都需要.
為什麼我叫生生不息的循環.
就是當你想要耶穌的那段日子.
你就會成為別人的老師.
或者榜樣.
大家有沒有試過和小朋友去賣旗.
特別是小朋友.
幼稚園小學那班賣旗.
大家有沒有這些經驗.
你有沒有試過.
你和他們賣旗會不會這樣.
你去賣旗吧.
不會的.
大家會做什麼.
當時我兒子還很小.
我就會和他說.
一會兒我們賣旗.
你就試試和叔叔說.
叔叔可不可以幫我買支旗.
你是要教他怎麼做.
你要教他說什麼.
更加你要做榜樣.
你要在他面前示範給他看.
先生你可不可以幫我買支旗.
是你的示範令到小朋友會發覺.
原來賣旗是要這樣的.
是要說這些.
原來不是這麼難.
他才會買.
同樣的.
每個弟兄姊妹.
坦白說信了耶穌之後.
我是不是就認識了基督徒.
當然不是.
是我們這些信了耶穌的人.
我們要做榜樣.
坦白說.
如果你問我的話.

$^{321}$我覺得時代進步了.
但我們的傳播心是退步了.
在我剛剛信了耶穌.
三十多四十年前.
當時我們的社會沒有這麼複雜.
複雜的意思是什麼呢.
沒有那麼多困難.
現在連公屋下面也要保安.
也要有身份證.
以前我們不是這樣的.
我們當時在荃灣.
荃灣在屯門也好.
你們的屋村全部下面是沒有保安的.
是自由往來的.
我們每一個星期.
是每一個星期.
我們信了之後也會洗房子.
大家知道什麼叫洗房子嗎.
我們會大家衝上去.
每一個家族都會拍門.
所以我裡面的記憶是什麼呢.
被人扔花生殼.
被人三字經問候.
被狗吠.
被人潑水.
什麼都試過.
但就是因為這個緣故.
坦白說.
我這些人.
其實我很怕面對這些事情.
但我看到那些哥哥.
那些姐姐.
他們那麼勇敢.
他們衝出去.
我在後面衝不衝.
也要衝的.
然後就經歷了很多上帝的奇妙.
原來這樣也可以傳福音.
明明看到那裡這樣的情況下.
我就知道原來這樣叫洗房子.

$^{361}$這樣叫傳福音.
我們才能夠跨出一步.
坦白說.
如果今天我們在座的弟兄姊妹.
你都坐在這裡的話.
我可以保證告訴大家.
你整個教會一定坐在這裡.
你不會有增長.
甚至乎.
之前有一句話很喜歡說.
佛系.
有就有,沒有就沒有.
然後我們就美其名.
推給上帝.
上帝心意.
你搞清楚.
你講上帝心意之前.
你已經盡力了嗎.
等於這群弟兄姊妹一樣.
當然他們一樣可以回覆約書.
看上帝心意.
他想拿到的時候.
上帝給就給,不給就不給.
你就講得輕鬆.
但你坐在這裡不動.
所以其實你會發覺.
首先我們今天要去做上帝的工作的時候.
為什麼我今天沒有一個叫局外人的原因.
就是你今天來到這個紅恩堂的地方.
就是上帝給你一個機會來到這裡.
我不知道什麼原因.
我不知道你是不是路過.
我不知道你是不是最近搬來.
我不知道.
但你要去問的就是.
上帝為什麼你讓我來到這個地方.
接觸到這個地方.
有沒有一些是你的心意.
而講求的就是.
大家你是否願意將你最好的.

$^{401}$我去奉獻出來.
而去幫助這些弟兄姊妹能夠成長.
另外一件事你會發覺.
這段經文令我感動.
就是第十六節至第十七節.
其實這個圖畫是很漂亮的.
為什麼呢.
你會發覺.
為什麼約書會害怕.
因為他過去見證著.
那些人怎樣去攻擊摩西.
大家記得嗎.
多次.
所以如果不是這樣的話.
在早一代的那群人.
就不會全部死在抗爭.
所以他害怕.
但現在這個圖畫是什麼.
竟然是你放心吧.
你叫我怎樣做我就怎樣做.
你叫我去哪裡我就去哪裡.
和以前摩西是一樣.
這給了約書一個很大的定心丸.
弟兄姊妹.
什麼叫做同心.
我想告訴弟兄姊妹.
你今天回來教會.
你不是去跟某一個領袖.
或者傳道人或者牧師.
不是的.
你弄清楚.
你今天是跟隨上帝.
今天我們為什麼要去馴服.
或者我們去跟隨領袖.
是因為他是上帝選擇.
就是這麼簡單.
很多時候我們的人性裡面.
我們很喜歡去說.
我喜歡和我不喜歡.
我喜歡你.

$^{441}$我固然的就願意跟隨你.
但我不喜歡你又怎樣.
我懷之你.
我都不喜歡你.
但我想告訴弟兄姊妹.
每一個領袖每一個牧者.
其實他們都有不同的處事風格.
和不同的恩賜.
你教清楚.
大家有沒有去玩過一些性格測驗.
有一個我最近不斷在研究.
因為我自己接下來.
我想是不是開一些個人成長.
其中有一個很好用的工具.
就叫做 MBTI.
他講了就是.
性格裡面大家很不同.
你要記住.
不同性格的人走在一起.
是未必合的.
但問題不是合與不合.
你要去認住的不是這個傳道人.
或者這個領袖.
你要去認住的是上帝.
你會發覺在昔日裡面.
我聽過有一個教會.
很出名.
他們也做很多工作.
也祝福了很多人.
有一位他們很資深的牧者.
是帶領著整個堂會.
甚至整個宗派.
他們發展得非常好.
大家很敬重這位牧者.
但這位牧者有一天也要退休.
他就很好.
他在退休之前.
他就祈禱之後.
他就去揀選了一位牧者.
接替他的位置.

$^{481}$然後那個牧者也很謙卑.
然後就去到教會.
退休牧師就叫他過來.
這間房間將會是他的.
將會是他自己的房間.
他就坐在這裡.
他也很謙卑.
我不敢坐.
這間房間是你的.
我說叫你坐你就坐.
然後他就坐在這裡.
在他上任的第二天.
有一個資深的信徒.
走進他的房間.
很兇的跟他說.
你沒有資格坐在.
XX牧師的這張椅子.
因為他說那張椅子是我買給他的.
我很感恩這位牧師.
他做分享的時候.
他很感恩.
他在神的裡面很平安.
他沒有很生氣.
很不開心.
他就很平安.
他說其實不要緊.
他說如果你覺得我沒有資格坐的話.
我坐第二張椅子.
很小事.
但你為我祈禱.
然後他就繼續沒有理會這件事.
繼續很努力.
繼續做事奉.
隔了幾個月之後.
這位弟兄姊妹.
走回他房間.
就跟他說.
他說XX牧師.
我要向你道歉.
因為實在我太無禮了.

$^{521}$但在這幾個月裡.
我終於明白.
為何牧師會揀選你去代表他.
他說我很樂意去跟隨你.
我很願意跟你在一起.
所以我想說.
很多時候人的裡面有很多盲點.
你為我們自己的一套.
但最重要的是.
我們在神的裡面.
究竟我們是否真的去跟隨上帝.
大家在聖經裡面.
還記不記得有一個叫可拉黨.
在民數記裡面.
就有六章裡面曾經說過可拉黨.
可拉黨的經文.
曾經這樣說.
你們的曾孫.
哥核的孫子.
以斯哈的兒子可拉.
和裡面子孫中.
以利亞的兒子大丹.
亞比蘭與比勒的兒子安.
並以色列會眾的250個首領.
就是有名望選入會眾的人.
在摩西面前一同起來.
聚集攻擊摩西亞論說.
你們擅自尊權.
傳會眾個個既是聖傑.
搖滑也在他們中間.
你們為何自告超過搖滑的會眾呢.
這些人.
可拉他們是因為是祭司的後代.
他們一直都會覺得.
我也是神教揀選.
神也在我們裡面.
為何你要告我呢.
他們就是看不起摩西.
他們認為自己比摩西厲害.
這件事的故事.

$^{561}$最後的結果是怎樣呢.
摩西亦都很鎮定.
他說不要緊.
我們每一個人.
這250人包括我和阿倫.
我們一同拿著上帝的香壇.
每一個人有一個金灰籠.
我們去到神面前.
我們去呈獻.
我們看看上帝的答案.
他說如果你們每一個人.
自然地死去.
那就是我有問題.
就代表上帝不揀選我.
但如果你們按我所說的.
上帝懲罰你們.
就代表你們的呈獻上帝不願立.
然後他即時在會眾面前說.
他說上帝.
你會看到這群人是怎樣的.
他說上帝如果是的話.
你不接受他們的話.
請你將地裂開.
讓這些人即時墮入深淵.
下去陰間.
所以當時的聖經裡面.
地即時裂開.
這250人全部掉進坑裡.
然後聖經裡面地就合起來.
所以弟兄姊妹.
我想給大家看到.
你不去馴服於上帝的揀選.
其實是很大的一件事.
是整個群體和上帝.
是會很反感的.
所以問題是.
我想大家留意.
剛才我說.
我們今天馴服的是.
上帝所選召的僕人.

$^{601}$但是你要記住.
這些的領袖.
這些的僕人.
絕對不是完美.
一定不是完美.
所以在這樣的情況下.
大家請記住這個特點之後.
你要留意.
如果今天我們領袖的能力.
和恩賜是有限的話.
我們需要的是.
大家的補足和同心合力.
剛才我提過.
最近我自己研究 MBTI.
聽你明白就聽著.
如果不明白.
等下完結之後再問我.
我是一個I人.
I人是什麼呢.
就是比較內向.
即是內向傾向.
那你怎麼搞的.
鄧傳道你在台上講導引.
你又跟人探訪人.
怎麼會是內向.
上帝就說我的特性是怎樣的呢.
我的特性就是.
如果我有一個很清楚的工作崗位.
我有一件事.
我很清楚要怎樣做的話.
我是沒有問題的.
我是絕對應付得到的.
我最不濟是什麼時候呢.
大家千萬不要找我吹牛.
我是完全不會找話題的.
我是沒有辦法在一個.
我沒有工作崗位.
或者角色的底下.
我沒有辦法自處的.
在這樣的情況下.

$^{641}$我舉個例子.
如果突然之間吃飯的時候.
應該這樣說.
你有事找我聊天.
我很樂意跟你聊三個小時都沒有問題.
但你說鄧傳道.
我只是想跟你一個人吃飯.
聊聊閒聊.
那我就死了.
除非你是一個很懂得聊天的人.
不然我坐在這裡.
我不知道跟你聊什麼.
其實某程度上.
我是很努力的去找話題的.
我想說我的特性就是這樣.
所以我剛出來事奉弟兄姊妹.
大家猜猜你是誰.
人們猜我做什麼.
你猜人們以前多數猜我做什麼.
在我沒有說過的情況下.
人們猜我做醫生.
他說很像.
平常我看你不是怎麼說話.
但反過來.
有診症.
我來找你的時候.
你又說OK.
然後診症完結了.
就完結了.
我又不懂繼續建立.
我想說這是我做傳道人裡面.
其中一個很大的缺陷.
坦白說.
所以坦白說我很羨慕.
一華牧師.
我很羨慕他.
因為他很能夠和大家交流.
他很有那種.
向外的那種.
我是很羨慕的.

$^{681}$所以其實我有缺陷.
我需要弟兄姊妹的醫人.
外向型的人幫我.
其實我很需要這樣的組合.
所以我以前有一次很成功.
就是我在之前的事奉教會.
我們比較基層一點.
我就說一起去探訪.
但我一定要有一班姊妹和我一起去.
他們就幫我約一些人.
聊聊天.
他們聊聊天的時候我插嘴.
好了,接近主題要講福音.
我來吧.
就拍得非常好.
所以那段時間大家都明白這個道理之後.
很多人進耶穌那段日子.
但我也曾經去過一段時間.
就是怎樣呢.
陳鎮成你是傳道人.
誰誰需要你找他.
我自己一個去就死了.
問題就是我連話題都想不起來.
他怎樣跟你講福音呢.
我想說.
我講這番話.
不是說有些你不要預我.
不是,而是反過來你了解我的缺陷.
我一些能力上的不足.
所以需要弟兄姊妹一起起來.
配搭,這個才是教會.
所以反過來你不需要討厭我的缺陷.
因為我想說這個也是上帝做我的一些特點.
但又倒過來.
所以我就需要大家一起起來一起配合.
當你的領袖對人和對事出了問題的時候.
他需要你的什麼呢.
是善意的提醒.
我想說的是你直接的詢問.
直接的分享.

$^{721}$其實是可以的.
我相信如果在神裡面一個謙卑的僕人.
他絕對願意聽大家的意見.
但最忌諱或最不好的是什麼呢.
就是在背後的批評和流傳.
我不喜歡某某誰.
然後我不跟他說.
但我在背後說很多.
然後你說著說著之後.
但是弟兄姊妹中文有句話非常好.
叫做.
我突然飛了.
就是言語的那種破壞力.
所以你就會發覺.
當弟兄姊妹你跟很多人說的時候.
並不是每一個人都能夠有能力去分辨.
真還是假.
所以當你說著說著之後.
那些人就當真了.
而我想告訴弟兄姊妹.
其實在一個群體當中.
你特別對於領袖.
對於牧者.
甚至其他弟兄姊妹.
你破壞了他.
你拆毀了他.
對教會有什麼好處呢.
原來耶穌在他的教導裡面.
他很清楚地跟我們有一個教導.
他說倘若你的弟兄得罪你.
你就趁著只有他和你在一起的時候.
指出他的錯來.
他若是聽你.
你便得了謝你的弟兄.
他若不聽.
你就另外再帶一兩個人去.
要憑兩三個人的口作見證.
句句都可定準.
若是不聽他們.
就告訴教會.

$^{761}$若是不聽教會.
就看他像愛邦人和瑞利一樣.
耶穌在這些教導裡面很清楚.
最初就是你和他.
還要在你們兩個.
不要在其他人面前.
為什麼呢 弟兄姊妹.
我們做家長講座.
教家長集你們很明白.
讚你的子女.
在哪裡讚.
在人面前.
中國人倒轉.
在人面前.
這個是壞人.
生了叉燒的那個生了他.
全部都在人面前講.
有多壞有多不好.
人家讚你真好.
那裡有女孩.
你女兒好美.
哪裡美 黑衣女.
我們看起來好像很謙虛.
但其實不是.
你知不知道其實在拆毀很多人的心靈.
但當你要教你的小朋友的時候.
閉嘴.
這個其實我們都在講對這個人的心靈的尊重.
其實我們都在講.
整個都是神的群體.
我們在講愛.
坦白講.
剛才我們說我們一定有瑕疵.
包括自己都有瑕疵.
我們想提出來的原因.
並不是要去拆毀他.
我們反過來是要建立他.
等於姐妹.
當我們講他肢體關係.
記不記得.

$^{801}$我們身體裡面有些肢體.
他們不體面的時候.
就怎樣.
要幫他越發加上體面.
所以你所做的一切.
是為了要建立這個人.
你要讓他在神的面前.
看見自己生命當中.
有一些問題.
然後我們一起努力.
所以等於姐妹.
你會發覺.
如果我們的教會裡面.
真的可以這樣去做的話.
我想告訴弟子們.
教會裡面才能夠一起向前行.
就不會我們一路向前行.
後面一路有人拖回去.
整個教會.
我們才能夠真真正正.
走在神的道路上.
所以我再講.
我們沒有局外人.
我們每一個人都有這個義務.
但請你記得剛才我說的話.
請你善意提醒他.
而所有的領袖.
所有的教會人.
都需要你在禱告裡面.
去守望他.
無論他是好.
或者他是不足.
都將他交在上帝的手中.
因為他是屬於上帝.
昨晚我們的團契.
我們才不斷討論同一個問題.
談了很久.
究竟我們可不可以做審判官.
其實我們沒有一個人.
有資格做審判官.

$^{841}$如果大家還記得勝經裡面.
耶穌講的例子.
教會裡面.
或者在這個世界當中.
一定會有墨子有拜子.
我們很容易就自己分.
那個墨子那個是拜子.
你千萬不要這樣做.
聖經說什麼.
是到末後的時候.
由上帝所差派的使者.
天使.
他們去分辨的.
而他講了一個原理.
當他很擔心.
當你去拔這些的拜子的時候.
連墨子都會拔掉.
你有沒有看到.
這個圖畫裡面最重要的是什麼.
是整個群體.
我們今天所做的任何事情.
有沒有以整個群體去做出發.
我們今天所做的事情.
是不是祝福整個教會.
我們所做的事情.
只是一時自己之氣.
一時自己的喜好.
喜歡或不喜歡.
還是我們真的去祝福整個教會.
你想教會一起起來的時候.
我們就需要有這份同心.
所以剛才那份就是.
你吩咐我去哪裡.
我就去哪裡.
你叫我怎樣做.
我怎樣做.
以前我怎樣聽從摩西.
我現在就怎樣聽從你.
背後就是因為我聽從上帝.
所以沒有一個方法叫絕對好的.

$^{881}$我很肯定.
我都說了.
就算我今天都算是一個帶領者.
我都很清楚.
我所講出來的.
不一定是最好的.
但是我在這個崗位上面.
所以我記得有一次.
有一年我在上一個教會事房.
我們教會都有十幾間堂會.
雖然全部都是小堂會.
有一次我們舉行一個.
很大型的全港報道會.
同一個機構.
就在一個很大型的.
類似大專會堂.
就是浸會的地方.
當時我負責是整個中派裡面的傳道部.
我是傳道部部長.
大家十幾個堂會就一起去.
我不知道大家是否明白.
十幾個堂會一起去做的.
最難的是什麼.
大家知不知道.
就是大家都很多意見.
因為現在大家都很愛主.
又很有經驗.
但是我想告訴你.
你會知道的.
大家都給意見.
大家都這麼愛主.
大家都這麼有經驗.
結果是什麼.
就是亂七八糟.
所以那次我自己很大膽.
做了一個決定.
我就是當大家一起來到.
當日參加的時候.
我裡面做了一個簡報.
開始的時候我做了一個簡報.

$^{921}$我今天以傳道部部長的身份.
想邀請大家一件事.
我說今天這個報道會是屬於上帝.
而上帝選了我.
今天負責這個報道會.
所以今天在這個報道會當中.
請大家暫時收起你們的意見.
只需要做一件事.
順服.
聽從.
我們猜拳你什麼崗位.
你就什麼崗位.
我們想你怎樣做.
你就怎樣做.
當時我聽到一些涉事私語.
有些人說.
這樣就有很多聲音.
但我沒有理會.
我們繼續去.
有些電影節目我記得圖畫是什麼.
我們差不多在不同的地方.
有些電影節目很順服.
有些電影節目有些樣子.
扭扭捏捏.
後來去到最後的環節.
去決戰.
後面有些電影節目.
當時我們叫他們站在那裡.
當時他們的意見是什麼.
為什麼不叫我們什麼什麼.
我說不是.
終於站在那裡.
後來我們見到.
有些人你都知道.
有些人.
決戰是喜歡玩彈弓手.
大家知不知道.
不是這樣的.
是這樣的.
是不是.

$^{961}$我們旁邊的就標記著這些人.
在彈弓手那些.
我們就叫旁邊的電影節目.
那個彈弓手.
站在他旁邊.
在這個過程裡面.
全部分佈齊了.
我們決戰的時候.
很感恩.
當晚發生什麼事.
其實當天是爭霸號風球.
我們在舉行.
沒有現在的天文台那麼好.
當時去到最後一刻.
他們都不說.
然後我們在最後.
都不知道是不是爆號風球.
還是雷電交加.
大家都很沒信心.
究竟那些人來不來.
那晚我們一起去做這些決戰.
完了之後.
所有的事情都完了.
才爆號風球.
然後電影節目在最後.
我們做回預測的時候.
其實很感恩.
電影節目說.
最初是.
最初我不開心.
我這樣這樣的時候.
其實我心裡很混亂.
為什麼不叫我做這些.
但後來他說.
你猜到我去做的時候.
我經歷了一些.
我沒有經歷過的事情.
而是上帝讓我看見他.
是異口同聲去說.
是上帝看見.

$^{1001}$上帝在當中的工作.
上帝的奇妙.
所以在那次B片裡面的時候.
大家都經歷到上帝的同在.
所以弟兄姊妹.
我想跟大家說的就是.
整個群體裡面.
最需要的.
剛才我說.
在最初上帝和他們有了應許.
你不要懼怕.
不要驚惶.
我必與你同在.
這是上帝的保證.
而第二件事.
是需要弟兄姊妹之間的.
互相支持和肯定.
同心一起去做.
所以我很盼望弟兄姊妹.
如果我們想在神的角度裡.
我們得勝的話.
請你記住.
若書也記.
這兩個重點是從哪裡開始的.
亦都互相提醒.
坦白說.
很多時候人有軟弱.
我對某個人很憤怒的時候.
我會跟他說.
Frankie 神.
哎呀 怎樣怎樣.
另外一個就要醒覺.
你就要提醒他.
如果有問題的話.
我們會不會一起找牧者傾傾.
一起找誰誰傾傾.
將你的感受告訴他.
我們會不會趁著還未很嚴重的時候.
我們就去處理問題.
令到大家在神的角度裡.

$^{1041}$可以同心一起向前.
所以坦白說.
相反.
如果群體內耗.
自己打自己.
坦白說.
不需要魔鬼出手.
我們自然都會一盤散沙.
而這是魔鬼最樂見的.
所以求上帝幫助我們.
我經常說.
我們的群體大不大 多不多.
不重要.
真的完全不重要.
我想問大家.
哪間教會 哪個宗派.
不是從小開始的.
很多時候我們只會看到.
我們的環境怎樣怎樣.
很難做.
不關事的.
我想告訴大家.
上帝在的時候.
你只會更加經歷到上帝的奇妙.
你只會經歷到上帝更多的準備.
最重要的是.
你的心裡面.
整個群體當中.
是否一切都很清楚.
我們一切同心跟隨上帝.
向著一個方向向前.
所以求上帝.
將這些祝福.
臨到給紅恩堂.
亦都臨到給錦繡堂 佳恩堂.
求上帝幫助我們.
一起同心祈禱.
在我們很需要在神的面前的時候.
我們的合一和同心.
我們去認定.

$^{1081}$我們去清楚.
我們人是有缺陷的.
但是因為上帝在你當中.
所以我們不懼怕.
最後就求主你幫助.
求主你給我們.
能夠有從神裡面的勇氣.
我們能夠去令到整個群體.
我們一起朝向你.
我們因著過去裡面.
其他人對我們的祝福.
我們今天去祝福其他人.
我們去帶領著一些新生.
初信的弟兄姊妹.
我們和他們一起起來.
能夠可以去得到更大的祝福.
而去令到將來.
更多人可以信耶穌.
同樣的我們也求主你自己幫助.
叫我們在我們人與人之間的關係裡面.
在我們帶領當中的時候.
求主你讓我們放下我們個人的一些想法.
而我們真的能夠可以.
循循善勸.
我們在過程裡面.
將我們的心可以和彼此的交流.
與上帝你自己幫我們拿開一切.
阻礙上帝你國度擴展的事.
將一切交在神的面前.
將所有的我們交在你當中.
其他奉教主耶穌.
得逞名字祈求.
阿們.
\newpage



\section{路加福音 6:20}
\label{sec:8UN_bXXTUSo}
\textbf{2025.01.26 《崇拜與群體德行陶塑》 路加福音6\_20 (和修版) 講員 : 何世傑牧師}
\newline
\newline
連結: \href{https://youtube.com/watch?v=8UN_bXXTUSo}{\texttt{https://youtube.com/watch?v=8UN\_bXXTUSo}} ~~~~ 語音日期: 2025-02-01
\newline
\newline
\hyperref[sec:15mYsSTo_Pc]{\small{< < < PREV SERMON < < <}}
~
\hyperref[sec:index_chronic]{\small{[返順時目]}}
~
\hyperref[sec:index_scriptual]{\small{[返順卷目]}}
~
\hyperref[sec:9D6qqGrY1os]{\small{> > > NEXT SERMON > > >}}
\newline
\newline
路加福音 6:20
\newline
\begin{longtable}{cl}
\hline
\hline
章節 & 經文 (和合本修訂版)\\
\hline
6:20 & \begin{tabularx}{0.7\textwidth}{X} 耶穌舉目看著門徒，說：「貧窮的人有福了！因為神的國是你們的。 \end{tabularx} \\ \\
[1ex]
\hline
\hline
\end{longtable}
$^{1}$我是不是要低一點?.
我比較小顆.
我們禱告.
願主幫助我們,恩著你的話語,叫我們生命有著改變.
叫我們每天有更新的時候,願主幫助.
禱告奉主明求,阿們.
很多謝曾牧師邀請我來這裡講道.
我們是同班同學.
他是我神學院時期的祈禱的土伴.
所以他有什麼事我都為他祈禱.
很開心可以跟大家宣講一個課題.
崇拜與群體德行的圖索.
頭部份會比較…不是聽得明白,但不要緊,聽聽吧.
但後面的部分會比較活潑.
在當代的語境中,我們經常用社群或共同體來描述群體.
很簡單地說,一群人成為一個團體.
這個團體包含了結構性的關係.
首先,我們可以看到教會組織中.
結構性透過教會體制得以彰顯.
展現秩序及層次的美感.
第二,我們成員之間有一個關係.
彼此奉獻,有互相的愛心流露出來.
好像一首交響曲,什麼都有,很開心.
剛才也是,很溫暖.
第三,生活中的共同體關係超越了心靈及思想的契合.
在日常生活中,我們亦會彼此聯繫,共同生活.
形成一種溫暖,又有真實的相依之情.
這個群體不單是形式的結合.
更是一種關係,一種很有生命的圖畫.
其實是一幅美麗的圖畫,展現在上帝的眼前.
展現彼此扶持共同生活的豐富意義.
耶穌基督以人的身份降臨在世間.
與人同住同行.
祂的生活方式及教導我們.
祂的教導成為了我們的典範.
啟迪我們這個群體,我們要學習,要實踐.
祂的教導不單充滿智慧.
更以生命的仁技,彰顯了天國的真理.
例如天國的八福.
向我們揭示了群體應該追求屬靈的品格和祝福.

$^{41}$我們聖經也很熟悉,葡萄樹的比喻.
約翰福音十五章.
也提醒我們,唯有連在這棵耶穌基督的樹裡.
才能結出美好的果子.
你離開了,就沒有果子.
脫離基督,我們的生命就會失去豐盛的源頭.
信仰群體唯有持續和耶穌基督保持一個靈性的連結.
才能慢慢塑造我們的好的行為,好的品行.
散發出屬天的光芒.
就是你的奶茶.
是好東西.
是見證.
如果你的奶茶是黑面的.
我就不回來了.
如果是笑容,早安.
我比較開心.
我在長洲.
我在金巴倫長老會長洲堂裡募會的.
我跟真牧師畢業後就出去了.
畢業後就留到現在.
我在長洲當牧師.
長洲比較簡單,快樂一點.
可能你見到我都會笑一笑.
沒那麼大壓力,很開心.
剛才說到,你遞出去的時候.
你服侍人,你由心裡出發.
是很開心的.
下一次,人們就會回來.
我們基督徒的生命是連結的.
不單止是你開心.
你們也會感染到其他人.
他們也會很開心.
這也是主給我們的恩典.
我們每個人的生命目的.
和群體的存在意義.
其實息息相關.
這麼複雜.
其實很簡單.
我們究竟是甚麼人呢?.
我最近跟一個八十歲的長者.

$^{81}$他在服侍教會已經很厲害.
已經很久了.
他在二十多歲的時候.
已經買了太古城的房子.
那時候已經很厲害.
又有老婆.
基本上每天都很開心.
但他說晚上的時候.
晚上寧靜的時候.
就會思考.
他做生意.
做他的工作.
他說每天已經安排好所有工作.
一放工就去吃飯喝酒.
最喜歡就是多喝兩杯啤酒.
然後就醉了.
醉了就可以亂說話.
可以亂說話.
我真的挺有趣的.
我沒有醉過酒.
我都不是很喝.
亂說話.
不斷說話.
但周圍的人.
他們不會給你一個心.
在亂說話.
星期六,日就約了親戚朋友.
打麻將.
八點打到晚上十二點.
再打一次.
每個星期都很豐盛的人生.
但他說一離開吃飯的地方.
自己一個人上的士.
靜了.
他覺得很空虛.
他說每天都是這樣?.
不是吧.
他有錢的.
有朋友的.
有娛樂的.

$^{121}$他說究竟我還欠什麼呢?.
最後他上了一個茶經班.
之後他就信了主.
他說耶穌基督就是填補我這件事.
他是不是做了傳道人呢?.
沒有.
他仍然去上班.
但是他每個星期日一放工.
昨天每個星期日就在教會裡服事.
很豐富.
很開心.
現在八十歲了.
繼續做.
我們生命的目的.
其實只有教會給我們生命的目的.
教會作為基督的身體.
存在的目的是身體的使命.
不會分開的.
因此每個個體.
都應該竭力地.
讓群體可以彰顯耶穌基督的形象.
我們要努力的.
不只是門口的每個人努力.
是所有人都要一同努力.
顯示耶穌基督的形象.
成為世人的光和炎.
當群體逐漸反映基督的形象的時候.
彼此關懷的裡面.
實踐愛的誡命的裡面.
就會自然流露出道德的向導.
會越來越好的.
你有沒有發覺.
你上了很多教會.
眼睛是會有些變化的.
要越來越好的.
道德的向導提醒我們.
所有的成員.
我們都要彼此相依.
彼此相依.
不可分離.

$^{161}$而教會更以神的真理.
作為道德的基石.
為社會提供無可替代的價值.
一個真正的基督徒的群體.
會帶來信仰.
行為和言語的轉變.
這正是品格塑造的過程.
基督的群體.
不單止展示了基督徒應有的品性.
更加見證了神的本質.
在生命共同體裡面.
我們一同生活.
彼此扶持.
甘心服侍他人.
譬如今早有人載我過來.
正在下雨.
我看到一輛車的時候.
我心裡面多開心.
曾牧師.
甘心樂意服侍.
成為神恩典活的見證.
崇拜不單止在日常生活的延續.
更加是和神.
深深相遇的時刻.
而這個相遇更加反映過來.
塑造我們日常不同的生活.
當信徒.
以有節奏.
循環的方式.
每天與神同行的時候.
他們的生命和品行.
將會被日漸更新.
這裡有一個節奏.
第一個節奏是.
TAP TAP TAP TAP TAP TAP TAP TAP TAP TAP TAP TAP TAP TAP TAP TAP TAP TAP TAP TAP TAP TAP TAP TAP TAP TAP.
大家有沒有手?有!.
我們試一下拍不拍到.
一二三.
厲害呀!.
又有節奏.

$^{201}$節奏是什麼呢?.
教會有沒有問過一個問題?.
如果有留教會的弟兄姊妹.
可能會問.
牧師傳道為什麼教會有這麼多活動呢?.
星期一,二,三,四,五,六,日都有.
總是有些活動.
為什麼呢?.
很想幫助我們.
跟隨耶穌基督.
但很多.
沒問題的.
希望大家跟隨教會的節奏.
來生活.
不要被社會的節奏.
影響你的節奏.
教會每一個活動.
其實每一個都是節奏.
譬如星期一崇拜.
每個星期一回來.
有一個節奏在裡面.
我們在節奏裡面會慢慢成長.
越來越像耶穌基督.
如果沒有了.
你就會.
去了第二個節奏.
信徒的生命.
常常受到世界的挑戰.
因為世俗的生活模式.
跟基督的生命力.
本質上.
有些不同.
這個世界跟教會所說的.
是有些不同.
我們的崇拜.
不單止要顯示.
這種差異.
更為了信徒.
可以開啟天國的門.
我們可以進入去.

$^{241}$每週崇拜的裡面.
信徒領受屬天的意象.
體驗主的美善.
和恩典.
這個不單止是一個.
屬靈的盛宴.
更加是一個生命的更新.
我們每一次翻崇拜之後.
我們有一點點的更新.
教會作為屬靈的燈塔.
引領群體.
站在真理的裡面.
幫助會眾.
跟世俗誘惑.
我們要選擇正路.
選擇對的路.
崇拜塑造群體的習慣.
我們每個星期日.
都會回來.
慢慢幫助我們.
越來越像耶穌.
逐步形成獨特的屬靈文化.
每一間教會的說法.
都有些不同.
去不同的教會靜悟的時候.
都發現每一個群體都不同.
我在這裡.
基本上是第一次來.
在其他日子來.
這裡的崇拜群體.
很活潑的.
很有喜樂在裡面.
每一個群體都有點不同.
崇拜不單止是.
一時的屬靈的經歷.
更加成為.
信徒生命的節奏.
持續地帶領我們.
進入真理的裡面.
成為世上的炎和光.

$^{281}$在世界上見證天國的榮耀.
崇拜怎樣塑造群體的德行.
如何塑造群體的德行呢.
後面有一個表.
這本是我的書.
已經很舊了.
已經出版了.
有四年了.
曾牧師有的,你問他.
我借來看,一定會送本給他.
那是在說什麼呢?.
我們如何在崇拜的裡面.
如何可以有一個.
群體的德行的改變呢?.
第一步.
群體首先要認識.
耶穌基督為主.
無論我們如何.
認同基督教的價值觀.
如果那個人.
未信耶穌基督.
未成為他生命的主的時候.
教會的模樣.
對他來說其實也沒什麼用.
因為他聽到.
教會學校很好.
團契很好.
人們都比較乖.
兒女們當然是回教會比較好.
不會學壞.
但他不是信耶穌.
還是差一點.
唯有讓基督成為.
我們生命的中心.
才能夠真誠地.
敬拜他.
沒有救恩的話,我們如何能夠.
真正地敬拜和讚美主呢?.
在崇拜的裡面.
我們群體一同承認.

$^{321}$耶穌基督是我們的救贖主.
以心靈回應他的大愛.
第二步.
群體會全心全意地.
專注敬拜神.
更願意跪下讚美敬拜神.
你要將人願意交給主.
才不會翹腳.
是很難的.
因為椅子比較矮.
我去過一間.
不同教會的教會.
坐在那裡很帥氣.
翹腳在那裡享受.
看何博士講什麼道.
上帝在那裡.
你不害怕嗎?.
你批評.
上帝的僕人.
看我有什麼說錯.
你不是吧.
不可以的.
我們要.
願意全心全意.
交給上帝.
知道上帝在那裡宣講.
上帝的話語.
你的習慣.
是慢慢會改變的.
我以前也很難的.
我也是會翹腳的人.
但翹的時候.
我心想不是在這裡敬拜.
還是不要翹腳.
第三步.
群體在崇拜的裡面.
有一個悟性的塑造.
認識到屬靈被.
塑造.
在崇拜的裡面.

$^{361}$我們深深體會到神的愛和美善.
這一份體悟.
進入我們生活的裡面.
成為力量,傳言.
不單止更新我們的心靈.
更加塑造了我們的品格.
第四步.
群體會參與崇拜.
裡面的聖樂福瀉.
因為我是聖樂.
雖然我比較活潑.
我也很嚴謹.
有些是很嚴謹的.
但我比較活潑一點.
可能是性格問題.
我們會在聖樂裡面服事.
生命會被塑造.
聖樂在崇拜裡面.
好像天國的旋律一樣.
激盪我們的心靈.
好的詩歌.
真的好的詩歌.
是會令我們有些不同.
透過音樂和讚美.
信徒的靈明得以滋潤.
心靈被滿足.
在五惡音的裡面.
我們經常可以與神同在.
靈魂因此被高舉.
整個群體的信仰.
氣氛更加濃烈.
什麼音樂.
是好的呢?.
大家有沒有唱過.
生日歌會唱到哭?.
在聖誕節的裡面.
我也有做兒童活動.
所以我有些兒童的心.
兒童的詩歌裡.
有一首.

$^{401}$聖誕節的詩歌.
最後會唱.
Happy Birthday to you.
Happy Birthday to Jesus.
Happy Birthday to you.
幼稚園生.
在唱.
他們在唱的時候.
我會哭.
小朋友.
還有多少時間.
有多少小朋友.
還有機會可以.
回到教會呢?.
我們知道有些地方.
小朋友不可以回教會.
我看見這些幼稚園的小朋友.
我一邊唱一邊拍手.
心裡流淚.
生日歌.
也可以流淚.
也可以感動.
什麼是好的聖樂呢?.
慢慢要上課.
我是一個.
建道神學院教崇拜學的.
網上的老師.
實體看不到我.
因為要上班.
我是教會教堂的.
牧師.
網上會看見我.
還有.
第五步.
群體直著身體的.
每個部份.
感受神的同在.
我們有.
眼耳口鼻和.
皮膚.

$^{441}$也是一個器官.
眼睛看見上帝.
看見恩典.
看見人.
可以經歷上主的同在.
耳朵聽到.
眼耳口.
口就唱歌.
鼻.
鼻怎樣可以感受到神的同在呢?.
我也有做.
長洲什麼都要做.
生育.
死葬都與我無關.
長者.
我們要推輪椅.
帶他們上崇拜.
長者一到那地方.
一聞到教會.
就算看不到.
在那裡聞到一陣氣味.
教會有陣氣味.
不是人情味.
是有陣氣味.
我母會是生命堂.
我太久沒有回去.
回到教堂裡.
崇拜場地裡.
一走樓梯.
走回崇拜場地.
因為他又請回港督.
一聞到那陣地氈味.
我以前做私琴.
做了十幾二十年.
走到鋼琴台.
那裡也是冷冷的.
不滑的位置.
每間教會都有陣氣味.
我們知道.
這裡就是上帝的地方.

$^{481}$還有什麼?.
感覺.
皮膚.
皮膚怎樣感覺到上帝呢?.
你旁邊有朋友.
你和他摸摸.
摸摸他的真肉.
你不要摸他的衣服.
摸摸他的真肉.
溫暖的.
和我們一起崇拜.
一起去敬拜.
旁邊的弟兄姊妹.
彰顯上帝的榮耀.
我們一起感受到.
主的同在.
主的臨在.
第六,崇拜在.
這個秘書做的過程裡.
經驗與神對話.
即是與神交心.
我們在崇拜裡.
與神相遇這段聖潔的時間.
我們可以在.
他的同在裡.
傾心吐怡.
你裡面.
可以和主說話.
不要怕.
有時機會過了.
就沒有了.
如果有機會信耶穌.
早些去決志.
如果還未洗禮.
快些進入.
上帝的恩典禮儀裡.
我們可以聆聽他的聲音.
我們在崇拜裡.
可以一同唱詩.
可以回應主的呼召.

$^{521}$第七步.
群體一同.
享受聖餐.
回應主的愛.
得到力量.
聖餐是一場與基督相聚.
神聖的筵席.
藉此我們可以追憶耶穌.
為我們的愛.
和他的犧牲.
這個儀式不單止.
呼喚我們的感恩之心.
更加設下屬天的力量.
激勵我們.
最後一步.
群體有能力在真實的社群裡.
活出基督的形象.
基督的教導.
不單止停留在.
崇拜的時間裡.
而是延伸到.
我們生活的每一個地方.
我們在崇拜裡所領受的.
信仰和德行.
應該在群體生活裡.
得到實踐.
成為耶穌基督的真實見證人.
我們不單止能夠.
反映神的榮耀.
更加將他的愛和真理.
傳遍世界.
讓世人可以看到.
基督的美善和大能.
這些步驟.
都是.
讓我們可以轉化.
的步驟.
我們可以成為.
榮耀基督的群體.
讓他的生命.

$^{561}$藉著我們在這個世界可以彰顯.
崇拜的意義遠超我們的想像.
遠超我們的想像.
每個人都應該.
擁有參與.
每個人都應該有機會.
參與的.
無論你是身處何種處境.
境況.
無論你是貧窮.
富足.
每一個靈魂都很渴望.
與神相遇.
都需要教會的擁抱.
讓他們可以與上主同在.
用一個真實的例子.
展現這份普世性.
精神復康的群體.
我們在裡面.
有服侍.
他們可否參與崇拜呢?.
他們有些精神.
有些問題.
可否參與崇拜呢?.
可以.
真的可以嗎?.
我和一個很好的同學.
不是曾牧師.
我們在灣仔裡.
啟動過一個精神復康的崇拜.
有十年的時間.
每三個星期.
一兩次.
由社工幫忙.
聯絡會員.
參加崇拜.
藉著崇拜.
我們越來越明白.
其實我和患病後.
康復的.

$^{601}$他們都有些問題.
我和他們差不多.
同樣需要耶穌的愛和救贖.
他們很渴望.
他們的信仰.
他們的渴望和信仰.
叫我見證崇拜.
其實對每一顆心靈.
都有意義.
無論是什麼角色.
有什麼疾病.
都有意義.
有無窮的力量.
這段經歷提醒我們.
崇拜不是一場儀式.
更加是一個橋樑.
叫人的心.
和神的心.
可以緊緊相連.
每個生命都可以在主的光裡.
得到醫治.
和找到盼望.
裡面有四位成員.
第一個叫阿操.
他經常沉溺.
在不愉快的記憶裡.
這份記憶.
主要針對外面的人.
但隨著這個群體.
他慢慢有信仰.
我們信任.
他經歷了主和他說話.
他最近開始分享.
他自己的經歷.
他很多時候都不敢分享.
怕別人標籤他.
他分享了小學時.
不愉快的經歷.
他已經50歲.
他發病時只有十多歲.

$^{641}$三十多年都在疾病裡.
他分享了小學時.
不愉快的經歷.
他分享了小學時.
不愉快的經歷.
三年級時.
別人說他是白痴.
一句說話.
記憶了幾十年.
他說他都不想.
他控制不到自己.
另一位.
阿仇先生.
他常常帶憂愁.
但他很有個性.
他很有型.
他懂彈結他.
他有他的才華.
因此崇拜時.
擔任結他手.
他有自己的風格.
何牧師可不可以.
我有自己的風格.
可以彈前奏.
你彈前奏.
我們唱.
他彈了很久.
我想像得到嗎.
他很有風格.
很有性格.
另一位.
阿珍.
他常常坐立不安.
會無端端笑.
而且坐立不安.
他會說不行.
就會去洗手間.
到處走.
幾個月.
因為天氣的原因.

$^{681}$我們改為Zoom崇拜.
精神科也有Zoom崇拜.
他說.
在Zoom會議.
會分享.
他說.
他說了很多次.
很氣憤的言論.
什麼事.
你這麼生氣.
因為旁邊有社工.
我怕他生氣得拍桌子.
怕旁邊的社工被他打.
他說.
何牧師住長洲.
在Zoom崇拜.
因為雨太大.
另一位朋友.
他住地鐵站.
你可以坐地鐵回來.
我們這麼辛苦.
才有一次聚會.
為什麼你不回來.
很生氣.
幸好我住長洲.
OK.
一個會員.
他對崇拜的.
執著.
那種難得.
他也有對上帝的尊敬.
還有一位叫阿樂.
他是一位單純而且快樂的人.
看見其他人有需要.
毫不猶豫地出手幫人.
他總是感到.
他到教會.
似乎人們對他不友善.
對他的要求很多.
他曾經提到.

$^{721}$有人在崇拜後.
可能有些教會的文化.
崇拜後.
分為小組祈禱.
他聽到.
旁邊的小組祈禱.
他便祈求.
叫旁邊的姊妹.
下星期回另一間教會.
也有.
害不害怕.
是他的經歷.
我不知道他有甚麼行為.
四位.
精神復康群體.
其實只是被忽略的群體.
其中一個小群體.
其實我們周圍.
還有很多默默存在的群體.
譬如跨文化的人.
我們會否接納.
南亞裔人士進入崇拜.
印尼人.
可否進入崇拜.
一起敬拜.
他們和我們同樣需要教會的溫暖.
和他們要去崇拜.
知道崇拜的美善.
如果我們已經體會到.
回教會有益處.
教會是很美好的.
是否可以更加主動地.
讓他們融入上帝的節奏.
陪伴他們.
不要只是周末.
中環不是很多人.
這就是他們的節奏.
和朋友從早上到晚.
這就是他們的節奏.
可否幫助他們.

$^{761}$進入上帝的節奏.
可否為他們提供.
與神相遇的空間.
我知道.
教會是非常忙碌的.
星期日.
當然沒有時間幫助他們.
可否在其他時段.
有這些可以幫助他們.
主的愛.
無處不在.
他們的角度.
可以讓我們看到.
很多榮耀.
最後,創建和接納.
新崇拜的群體.
敬拜.
群體同來.
跪拜.
讚美上主.
群體心被恩感.
歌頌上主.
感恩.
群體可以向神認罪.
回歸上主.
群體可以呈現.
基督的德行.
活出基督.
今天原本是邀請我.
講聖樂主日.
主要是講崇拜.
但程序沒有講聖樂主日.
所以我比較安樂.
我帶了一支手風琴.
希望可以.
拉手風琴.
我做了.
傳統牧師.
很少拉手風琴.
因此我的功力.

$^{801}$只剩下幾成.
以前也有拉過.
小時候.
拉手風琴.
稍後會演奏.
一首歌.
火車名笛楊主恩.
我記得這首歌.
以前在幫教會籌款.
去拉這首歌.
或者去.
華人落後的地方.
因為手風琴.
很接納.
拉完手風琴後.
可以派大白兔.
小朋友可以.
一起聽聽耶穌的故事.
火車名笛楊主恩.
這首歌是改編自.
火車飛過苗山寨.
也是我老師.
所寫的作品.
已經在天堂.
這首作品.
他很厲害的.
拿了世界冠軍.
手風琴世界大賽冠軍.
有他一成.
也很開心.
這首歌原曲描寫.
村民.
首次目見.
火車興奮的時候.
在我的.
改編中.
幫他加了一些.
屬靈的意義.
有兩次模仿.
火車進行.

$^{841}$和名笛的聲音.
象徵上主的呼喚.
喚醒子民.
要認罪悔改.
回歸在他的懷抱.
最後.
名笛聲更加.
迫切地猜險上主的工人.
將福音可以傳遍.
地上每一個角落.
拉完這首歌後.
我們會唱一首.
回應詩歌.
奇異恩典.
以奇異恩典.
來祝揚上主.
我用手風琴伴奏.
你們可以一起唱.
有四節.
在群體裡.
我們一同用這首詩歌.
回應主的恩典.
這首詩歌提醒我們.
無論貧窮或富足.
神的國.
都為我們躺開.
我們歸向祂.
我們的生命就會得到豐盛.
你們貧窮的人有福了.
因為.
神的國是你們的.
崇拜不單止.
帶領我們親近神.
更加塑造我們的生命.
塑造群體的德行.
叫我們成為.
活在基督故事裡的見證者.
叫我們在崇拜裡.
得到更生.
榮耀上主.

$^{881}$更加可以將福音.
傳到地域.
我去準備一點.
好,那我呢.
(....聽不明教練說了多少次ayana是英文不懂).
(教練說了多少次).
我們同心禱告.
你賜下音樂給我們.
讓我們可以藉著音樂來敬拜上主.
上主你也幫助我們.
可以在音樂裡.
在崇拜裡.
可以經歷到上主的恩典.
求主.
求你繼續賜福教會.
讓教會成為一個.
更加美麗的群體.
禱告奉主明求.
阿們.
(教練說了多少次).
\newpage



\section{箴言 29:18}
\label{sec:9D6qqGrY1os}
\textbf{2025.02.02 《洪恩家齊躍動》 箴言29\_18 (和修版) 陳麗蟬牧師}
\newline
\newline
連結: \href{https://youtube.com/watch?v=9D6qqGrY1os}{\texttt{https://youtube.com/watch?v=9D6qqGrY1os}} ~~~~ 語音日期: 2025-02-03
\newline
\newline
\hyperref[sec:8UN_bXXTUSo]{\small{< < < PREV SERMON < < <}}
~
\hyperref[sec:index_chronic]{\small{[返順時目]}}
~
\hyperref[sec:index_scriptual]{\small{[返順卷目]}}
~
\hyperref[sec:Ho71_1KRGq8]{\small{> > > NEXT SERMON > > >}}
\newline
\newline
箴言 29:18
\newline
\begin{longtable}{cl}
\hline
\hline
章節 & 經文 (和合本修訂版)\\
\hline
29:18 & \begin{tabularx}{0.7\textwidth}{X} 沒有異象，民就放肆；惟遵守律法的，便為有福。 \end{tabularx} \\ \\
[1ex]
\hline
\hline
\end{longtable}
$^{1}$丁字妹新年快樂,主因滿滿.
很感恩,不過我講到要用我的powerpoint的remote暫時未有.
所以沒有powerpoint都可以講到.
很感恩,來到這裡大家非常開心.
來到這裡大家招呼我喝奶茶,吃賀年食品.
我原本應該在去年六月要退休.
在六月底要退休.
在2023年11月左右,曾牧師跟我說.
要請我退休後繼續在錦繡半職事奉.
當時很感恩,能夠退休後繼續事奉.
不過我有很多時間開始計劃.
計劃在退休後不會立即回來.
有人邀請我去美國.
有人邀請我去新疆,南疆,北疆旅行.
有人邀請我去創啟地方做全道訓練領袖.
我答應了.
報名參加退休營報了四個.
接近一萬元都交了錢.
但在12月初有三位唐董找我談話.
我心想唐董可能覺得已經夠人了.
你都不用做了.
可能我又要再多作安排.
誰知唐董告訴我.
因著童工有些調動.
曾牧師會來紅恩堂和大家一起.
繼續發展福音.
我會留在錦繡堂推動方工作.
很感恩,所以一年沒有來港.
今天看到很多弟兄姊妹.
舊的在這裡.
亦有很多面孔我不認識.
今天只講一節聖經.
就是《箴言》29章18節.
不如我們一起讀.
一二三.
沒有異象,民就放肆.
為遵守律法的便為有福.
我應該按哪個按鈕.
太陽的按鈕.
向下那個.

$^{41}$好,謝謝.
我們看看異象.
異象這個字不是說晚上發夢.
更加不是我們白天發白日夢.
異象在聖經裡的意思.
就是上帝說話的一些默示.
讓我們知道人生裡的一些方向.
給我們一個指引.
讓我們看到他的心意.
當人有異象的時候.
我們就會有目標,有方向.
有目標,有方向的時候.
我們自然會帶來動力.
我們看看這個小朋友.
他的姐姐跟他說.
「快要考試了,快點讀書吧」.
弟弟說「不用了,我已經足夠了」.
「我學了很多東西」.
姐姐說「你現在只是在讀小學」.
「你打算將來可以做什麼呢?」.
這個弟弟有他的想法.
他說「我比BB班的人懂得多一點」.
「所以我可以去教幼稚園」.
這個究竟是不是一個異象.
是不是一個使命呢?.
這是一個難度.
我記得在2023年底的時候.
魯木斯離開了我們.
我們就去籌備安息禮拜和追思會.
當中有很多瑣碎的事情.
很趕著去做.
裡面有些場景令我很感動.
弟兄姊妹說.
他們很想為魯木斯做一件.
在他最後為他做一件事.
要做到最好.
所以在他們同工的時間.
我看到他們真的盡力了.
其中一幕就是他們要找吉兒的一元.
他們估計要找三百個.

$^{81}$後來還是不夠的.
要三百個.
不知道為什麼銀行找都找不到.
只找到一百五十個.
弟兄姊妹竟然回家的床底.
籠子裡的東西.
湊起來就找到三百個.
我看到他們這麼小的事情.
但他們希望都能夠做到最好.
我想他們一開始的大前提就是.
他們要為他們最愛,最尊敬的魯木斯的安息禮拜.
追思禮拜.
能夠做到榮神益人.
周圍的弟兄姊妹都可以得到安慰.
因為這個愛的意象.
他們學習不求自己的益處.
所以整個跟他們一起的事奉過程.
他們要誰載他們就載誰.
就算沒有空.
他們都會努力想到辦法.
就是他們的這個意象.
成為了當時的使命.
也成為了他們整個事奉的動力.
既然我們在紅恩堂.
我們有我們的意象.
有我們的使命.
不過我們就要說一下我們的行動.
我們要一起去起行.
在新的一年.
2025年我們要一起去起行.
我問一下弟兄姊妹.
新的一年有沒有立了今年在屬靈裡面的目標和方向?.
有沒有?.
有些人又up 肚了.
有些人就沒有動了.
不過這麼老套的事情.
很久沒有做了.
但記住.
意象,使命是叫人有方向.
有這個方向的時候.

$^{121}$就帶來這個動力.
也帶來我們的勤奮.
昨天有位姐妹跟我聊天.
她說陳牧師.
將來我們某一個組.
我們要方便溝通.
去進行一些的事奉.
我們就要推出一個電腦的新方案.
我就告訴她.
我是老人家.
可能我要學很久很久才會.
她就這樣說.
這樣啊?.
算了算了.
用回舊方法.
不要這麼煩.
我說不行啊不行.
沒試過.
我怎麼知道我一定不行.
所以先學.
學了或許可以.
有方法有新方法.
如果學了不行的時候.
才用回舊的方法.
她就立即拍手掌.
陳牧師.
欣賞你啊.
這個就是我們做人的態度.
有些事情出來.
先試了.
不試的時候.
才用回舊的方法.
這個就是我們推動.
有一個使命.
有一個意象.
就是希望能夠有更好的溝通.
2023年的時候.
我又和一班姐妹.
大家都說靈命很枯乾.
不如我們一起努力.

$^{161}$希望能夠靈命高升.
我們說不如.
大家彼此勉勵.
一起多讀聖經.
我們七個人.
大家懂的文字不同.
讀書的程度也不同.
但不要緊.
我們一起立下目標.
一起讀經.
其中有五個人.
說要全年讀完整本聖經一次.
包括我在內.
另外有兩個人.
說自己對文字比較弱.
不行.
他們說喜歡讀甚麼呢.
他們說全年希望讀完四卷福音書.
好.
我們在這年裡一起努力.
到2023年底.
結果是怎樣呢.
有兩個人.
曾經說要讀完整本聖經.
他們說真的.
全年讀完新舊約一次.
但沒有包括我在內.
另外那兩位姐妹.
他們說全年讀完四卷福音書.
他們很厲害.
兩個人都讀完四福音.
並且開始了少許的仕途行傳.
他們也完成了他們的目標.
但我和另外三位完成不了.
全本聖經完成不了.
不過我們也讀了大半本.
雖然完成不了.
但也有努力過.
我們有使命.
到上年2024年.

$^{201}$我覺得經過一年.
今年大家自己來吧.
到2024年底.
我問他們.
大家上年做了多少.
七個人.
全部的答案都是一樣.
沒有做過.
所以人就是這樣.
當我們有異常.
知道自己不行.
我們要拼命高升.
我們就要有使命.
有使命的時候.
我們就需要有行動.
否則一件事都做不成.
所以盼望在新的一年裡.
我們要為上帝做更多工作的時候.
我們就先需要在靈命上高.
我們一起努力讀經和祈禱.
立這個目標.
立這個方向.
第二方面.
我們如何一起更衷心的去事奉呢?.
有一次有位弟兄問我.
應該是2023年的時候.
應該是2024年初.
有位弟兄問我.
究竟有多少位同工.
曾牧師就來洪恩堂了.
我們感受那邊剩下多少.
就去數.
他說:鄭牧師只有半職.
但他在嘉欣那裡.
星期天的嘉欣.
有一天探訪.
感受這邊就剩下鄧全道.
楊姑娘,施元姑娘.
當時姑娘還在.
還有我.

$^{241}$我數完之後去定一定.
數完了.
就是這麼多了.
她說:哦.
定一定之後.
她就說:不要緊.
我們錦宣家有很多弟兄姊妹.
可以動員.
她倒轉做一個很積極的層面.
就來去看.
同工就是這麼少.
不過弟兄姊妹.
我可以告訴大家.
香港所有的教會.
裡面的同工都不多.
洪恩家有兩個很幸福的.
我們感恩.
感恩我們在感受堂.
我們有二,四,五個.
我們也很感恩.
不過這個同工不是很夠.
這個不要緊.
最重要是全部的弟兄姊妹.
都可以齊齊起來.
一起事奉上帝.
洪恩家本來以前是小小的港台.
現在就非常廣闊.
在下面.
剛才來到真的很開心.
大家弟兄姊妹熱烈的歡迎.
令到每一個回來的弟兄姊妹.
或者新來的朋友.
都感到被歡迎的感覺.
我知道在這麼多裝修工程裡面.
感受家,洪恩家和嘉恩家.
都有來幫忙.
你們有沒有來幫忙過?.
做這個裝修.
有沒有來幫忙做過這個清潔.
這個是我們的家.

$^{281}$所以理應我們每一個.
都要一起來下手,下腳去做.
有時感受弟兄姊妹跟我說.
陳牧師.
我這個禮拜或幾個禮拜後.
我要去洪恩堂做主席.
我說好.
在洪恩暫時人手缺乏的時候.
要下去幫忙.
大家彼此支援.
不過要謹記一件事.
不可以全力去做.
因為全力去做的時候.
洪恩家的就不能去發揮.
每一個地方,每一個教會.
我們都要各人按著恩賜去事奉上帝.
我說是要的.
大家彼此的支援.
不過一定要留一些機會.
給洪恩家在裡面去發揮.
我很感恩.
在上一次的禮儀中.
我見到洪恩裡面有年青人洗禮.
出來講見證的時候.
很感動.
知道在洪恩裡面的兒童工作.
學生工作.
青年工作一直發展得很好.
我知道曾姑娘.
還有江澤夫婦.
還有很多兒童工作者的弟兄姊妹.
我不認識.
她們很用心地一直在這裡.
為洪恩家下一代去事奉.
她們很用心.
一路數一路數.
很多弟兄姊妹.
我又問一下.
弟兄姊妹.
在這個淑齡的家.

$^{321}$在洪恩家的擺上有多少呢?.
可不可以再多一點呢?.
除了要為這個淑齡的家.
多一點去祈禱之外.
去想想自己有什麼恩賜.
可以做主席.
可以去訓練小朋友.
可以參與兒童的事工.
還有青年工作.
需要很多人和他們一起去同行.
所以想想.
想想有什麼清潔.
這些不行.
下屬幫忙沖奶茶.
幫忙預備很多食物.
幫忙接待.
加入這個師班.
各方面.
這個淑齡的家.
很想洪恩裡面的每一位弟兄姊妹.
都齊齊地去參與.
所以當你找到自己有什麼恩賜的時候.
或者找不到.
你問上帝.
問曾牧師,曾姑娘.
他們告訴你.
然後請他們為你們安排合適的崗位.
讓我們整個洪恩家齊齊地躍動.
同心地去事奉榮耀上帝.
短暫性在外來有不同的支援.
我們知道小玲或者我們師班.
有不同的弟兄姊妹,主要姐妹.
下來幫忙訓練師班,領唱.
很感恩.
這是很需要的.
但在長遠計算的裡面.
我們必須要洪恩裡面的弟兄姊妹.
自己起來承擔在這裡的地方.
福音的使命.
這才是真正的成長.

$^{361}$真正的增長.
我們一起去三間堂.
感受嘉恩和洪恩.
我們在不同的地方.
齊齊地努力.
第三方面.
我們齊躍動要多禱告.
我們看回教會歷史裡面.
很多很重要的事件.
都是從禱告開始.
當耶穌基督接受了水禮之後.
祂就去禱告.
因為祂接受了水禮之後.
祂就要踏上前面事奉的路程.
當祂在禱告的時候.
上帝就來跟祂說.
「你是我的愛子」.
「是我所喜悅的」.
耶穌基督將要踏上事奉的路.
然後跟著上十架的時候.
祂在禱告.
上帝就認定祂是神兒子的身份.
加祂的力量.
當耶穌基督要上十字架路.
祂就帶著彼得,約翰,雅各.
一起上山去禱告.
又是禱告的時候.
祂的面容變得很光芒.
登山變相.
因為當祂禱告完之後.
祂要下山上十架路.
又是上這個艱難的路之前.
祂又是禱告.
踏入這個使徒行傳.
彼得有一次在這個流房裡祈禱的時候.
他就看到一個大布.
在大布裡有不同的動物.
在禱告的時候.
上帝的靈臨到他.
他就呼召彼得.

$^{401}$他要去做這個愛邦人的工作.
又是禱告.
到巴拿巴和保羅.
他們要起行宣教.
第一次宣教旅程的時候.
他們就禁食禱告.
原來所有教會歷史裡面.
一些很重要的事情.
開展的時候.
都由禱告開始.
所以洪恩家的弟兄姊妹.
我們要在上帝裡面發展.
上帝會怎樣洗禦洪恩家.
我知道曾牧師已經在這裡兩間中學聯絡.
嶺南大學也在聯絡.
我們怎樣去發展青少年的事工.
在這些聯絡裡面.
我們必須要先從禱告開始.
我又感恩.
知道你們一星期有兩次早上祈禱的聚會.
一會崇拜完之後.
我就知道會有十幾二十位的弟兄姊妹.
自然進入這個房間裡.
齊齊地去禱告.
有這麼多人願意加入禱告的行列.
可能在下一年.
或者不知道什麼時候.
我再來港島的時候.
已經全部坐滿了.
然後就要開房間.
所以弟兄姊妹.
我們在洪恩家齊躍動的時候.
我們必須要從禱告開始.
今年的教會年提是什麼呢?.
不如你們讀給我聽吧.
懷抱報道初心.
福音獻給社群.
這個年提就是大家訂出來的.
同工和一些管委一起訂出來.
不過不夠.

$^{441}$就要所有的肢體.
一起來起行的時候.
才能夠達成今年的年提.
否則就只是一個口號.
盼望每位弟兄姊妹.
我們都用實際的行動.
我們去到這個社區.
去關顧文后.
新出主耶穌基督愛的援手.
很記得是一套電影.
不過是好幾年之前的.
《降巨靈》.
我讀不出這三個字.
這套是真人真事改編的電影.
很感人的.
是講到在第二次大戰的時候.
一位美國的軍醫.
是一個基督徒.
隨軍隊去到日本沖繩島打仗.
他當時很堅持.
堅持上帝所頒布的十誡.
不可殺人的教訓.
所以他堅持不向敵人開槍.
並且向上帝作出了一個毀身的使命和祈禱.
他說救得一個就一個.
所以在整個戰爭裡面.
他救了很多受重傷或者行動不便的同胞.
雖然也有很多人很憤怒他.
為什麼你不開槍?.
為什麼你不開槍去殺人?.
但是在第二次大戰結束之後.
美國總統就頒了一個勳章給他.
表揚他救了很多同胞的偉大精神.
所以在搶救靈魂的事情裡面.
我們都需要有這個精神.
救得一個就一個.
所以我們要努力救得一個就一個.
沒有異象,民就放肆.
放肆是什麼意思?.
其實不用說也知道.

$^{481}$罵小孩的時候不要放肆.
就是說不要輕率自把自為.
所以今天在這個真言裡面.
他說的那兩句就是沒有異象,民就放肆.
下面那兩句就是為遵守律法的便為有福.
所以在我們的家庭裡面.
或者在我們的社區裡面.
當所有的人都能夠遵守上帝所定的律法.
遵守上帝給我們的教導.
我們接受來走出來.
我們的生活和人的關係.
自然能夠有一個快樂,幸福和諧的生活.
所以這個錦宣家我忘記了.
改了他做洪恩家.
在上年我跟他們說的時候.
我們需要在這個全部的屬靈的弟兄姊妹.
一起來去慶行的時候.
我們需要有這些認定.
我們一起來去讀好不好?.
1,2,3.
我生命最高目標是為了榮耀上帝.
並頌讚他的美善.
我願意多操練上帝給的恩賜.
並且運用他來榮耀上帝.
所以我們記得一件事.
我們所作的一切要榮耀上帝.
並且我們每一個人都要按著自己的恩賜.
我們來去發揮,事奉上帝.
我們一起讀第三,第四.
1,2,3.
我會努力建立與弟兄姊妹之間的友情.
彼此相愛.
我願意勤讀聖經,祈禱.
讓上帝改變我的生命.
所以弟兄姊妹.
我們要去跟弟兄姊妹有一個和諧的關係.
然後我們活出這個愛.
但是我們自己的屬靈生命如何被提升呢?.
多禱告,多讀聖經.
唯有聖經的話語改變我們的生命.

$^{521}$這句是在網上抄出來的.
很好.
有意象我們才有使命.
有使命我們就知道我們要去傳遞.
如果不是有意象在我們裡面.
我們沒有事做是沒有用的.
但當我們去傳遞意象的時候.
我們才有事做.
但如果我們每個人都不做事的話.
我們就只能坐在這裡.
但當我們肯為上帝工作的時候.
就經歷到上帝的能力.
所以你經常覺得我沒有什麼恩賜.
我的恩賜只有很少.
只要肯踏出去做的時候.
會經歷到上帝奇妙的能力.
經歷到上帝奇妙的能力.
我們就可以出來講見證了.
所以弟兄姊妹.
我們經常說我們沒有什麼話說.
沒有什麼經歷.
就是因為我們沒有向前踏出一步.
所以最後結束的時候.
我們就唱一首.
簡太告訴我們.
大家都很熟悉的.
《要忠心》這首副歌.
要忠心 忠心.
切忘我 要我是否.
忠心 要忠心.
忠於你的永要我.
要忠心 忠心.
忠別人 都我尊重.
無能可敬方.
必要著自薄.
願至雙健於終生.
我想請問有多少人是不認識這首詩歌的呢?.
一,二.
只有兩個人不認識.
三個人不認識.

$^{561}$四個人不認識.
五個,六個,七個.
沒理由是這樣的聲量.
是不是?.
我們要為上帝去啟行.
事奉的時候.
可不可以用盡全力.
再來快一點.
請大家等一等.
一,二,三.
要忠心 忠心.
切忘我 要我是否.
忠心 要忠心.
忠於你的永要我.
要忠心 忠心.
忠別人 都我尊重.
無能可敬方.
必要著自薄.
願至雙健於終生.
我們一起同心圖局.
看到洪恩家他們一路的增長.
弟兄姊妹一路的熱心事奉.
他們和諧的相處.
最後願意在新的一年裡.
給每一位弟兄姊妹.
都很認定.
這就是我們屬靈的家.
我們要一起在當中參與.
我們來事奉.
我們有多的做多.
有少的做少.
但我們每個人都可以做的就是禱告.
主就求您.
求您來幫助.
幫助我們眾弟兄姊妹.
惠著主您的國度.
獻上更多的禱告和獨經.
求主幫助.
幫助讓我們能夠同心合意的.
慶旺福音.

$^{601}$求主使用.
多謝您垂聽禱告.
是奉靠主耶穌基督得勝的名字而求.
Amen.
\newpage



\section{約翰一書 5:18-21}
\label{sec:Ho71_1KRGq8}
\textbf{2025.02.09 《信仰要尋真》 約翰一書5\_18-21 (和修版) 曾裔貴牧師}
\newline
\newline
連結: \href{https://youtube.com/watch?v=Ho71-1KRGq8}{\texttt{https://youtube.com/watch?v=Ho71-1KRGq8}} ~~~~ 語音日期: 2025-02-10
\newline
\newline
\hyperref[sec:9D6qqGrY1os]{\small{< < < PREV SERMON < < <}}
~
\hyperref[sec:index_chronic]{\small{[返順時目]}}
~
\hyperref[sec:index_scriptual]{\small{[返順卷目]}}
~
\hyperref[sec:ulbWKYhlm8s]{\small{> > > NEXT SERMON > > >}}
\newline
\newline
約翰一書 5:18-21
\newline
\begin{longtable}{cl}
\hline
\hline
章節 & 經文 (和合本修訂版)\\
\hline
5:18 & \begin{tabularx}{0.7\textwidth}{X} 我們知道，凡從神生的，必不犯罪；從神生的那一位，必保守他，那邪惡者無法加害於他。 \end{tabularx} \\ \\ \relax
5:19 & \begin{tabularx}{0.7\textwidth}{X} 我們知道，我們是屬神的，而全世界都伏在那邪惡者的權勢之下。 \end{tabularx} \\ \\ \relax
5:20 & \begin{tabularx}{0.7\textwidth}{X} 我們知道，神的兒子已經來到，並且將悟性賜給我們，使我們認識那位真實者，我們也在那位真實者裡面，就是在他兒子耶穌基督裡面。這是真神，也是永生。 \end{tabularx} \\ \\ \relax
5:21 & \begin{tabularx}{0.7\textwidth}{X} 孩子們哪，你們要遠避偶像。 \end{tabularx} \\ \\
[1ex]
\hline
\hline
\end{longtable}
$^{1}$我們一起祈禱.
再一次感恩.
我們可以透過聖餐.
紀念我們的主.
來得到你永恆生命的模樣.
使我們能夠生命得到飽足.
今早我們再一次聚集.
聆聽上主的教導.
尤其是我們在過年.
我們見到有很多很多的香港人.
他們都是這樣來到.
由初一一直到今天.
都是這樣來到去拜.
很想尋求真神的祝福.
但我們確實透過聖經知道.
這一切都是徒然無用的.
我們需要逃避偶像對我們的影響.
今早我們盼望透過使徒約翰.
給我們的教導.
讓我們能夠認識.
我們的信仰的真確性.
以及我們已經進入尋真的過程.
今早我們等候在神面前.
求主的聖靈開啟.
引導我們大家的心靈.
奉耶穌基督的名.
阿們.
隨便可以看到很多的相片.
你見到年三十晚要裝頭炷香的.
特別醒目的大家都知道.
每年都穿著賀年的每年生肖.
夏惠姨或者其他的人.
爭相頭炷香.
在在反映人心很想得到真理.
很想得到從自己以外的祝福.
不知道大家有沒有看過這套電影.
來自外太空的PK.
PK不是髒話.
印度話的傻瓜.
相當有深度的一套電影.

$^{41}$當然當時引起印度很廣泛的討論.
以及認為阿米爾罕想嘲諷印度的宗教.
列舉不同的宗教.
這個外星人他失去了遙控器.
回不了他自己星球的聯繫.
他就要去問哪裡可以找到呢.
有人就告訴他.
你要找你的遙控器.
一定是找神.
神就會啟示你.
他真的不斷去找印度教裡面.
大大小小的神.
當然是一套笑話的電影.
但裡面的深度就是.
最終有一樣東西.
是很值得我們去想想.
他說原來他一直的尋找過程.
很多人.
很多宗教的代理人.
來告訴他.
這些的代理人.
不同宗教的代理人.
來告訴每一個信.
例如你是信婆羅教.
或者是其他的.
例如濕婆等等的印度宗教.
你就要用什麼方式去做.
但他全都做了.
真誠地做了.
都不能夠讓他找到他自己的遙控器.
但他最終真的發現了.
他說真神只有一位.
在電影裡面.
真神只有一位.
他不需要有代理人.
來跟他做一個代理.
真神是會向人啟示的.
很特別的地方就是.
這一套電影.
反映了一個很重要的真理.

$^{81}$雖然這套電影從來沒有想過.
要告訴我們基督教才是真理.
因為他也去找基督教.
主要找天主教.
但令他不能夠得到答案.
弟兄姊妹.
我們今天也是敬拜神的群體.
我們來到神的面前.
我們如何尋找這一位真神.
而獲得我們自己的真理呢?.
在《約翰一書》很特別的地方.
最後這一段經文.
到五章的二十一節.
那一節聖經與整個上下文.
是沒有任何直接的關聯的.
就是小子們.
我們要自首.
要很警醒自己.
要遠避偶像.
當時很明顯的.
在爾弗所的地方.
因為仕途約翰年紀大了.
他牧羊的教會.
就在爾弗所的大城.
爾弗所的大城是一個守護女神.
就是亞底米女神.
Diana.
現在成為一個很重要的古蹟.
任何旅行團都會去的.
就是一條柱的位置.
是一個很多羽房的位置的女神.
可想而知.
當時爾弗所是非常浩大的敬拜假神的活動.
但是約翰是否只是限於.
你千萬不要去亞底米女神廟拜.
是否只是限於這一點呢?.
而在整個《約翰一書》裡面.
有幾件事是很重要的.
約翰是重複地說的.
居然在最後這三節聖經.

$^{121}$五章的十八.
一直到二十節.
出現了不斷說的主題就是.
我們知道十八節.
第十九節.
我們知道第二十節.
我們又知道.
三個知道的.
一個基督徒尋找神.
你就應該會得到神給你的教導.
十八節.
我們知道凡是從神生的.
必不犯罪.
從神生的那一位.
是說那個人.
必保守他.
表面上因為新的和合本.
有些難解釋的.
但是舊的和合本就很簡單.
或者漢語聖經也有簡單說.
主體是那個人.
就是凡從神生的那個人.
他必定不會犯罪的.
從神生的那個.
會保守他自己的.
這是舊的和合本.
接著留意這裡.
那邪惡者無法加害於他.
這裡有兩個問題.
我們可以思考.
誰是邪惡者呢?.
上下文就應該知道了.
因為下文就說魔鬼了.
另外一個問題就是.
從神而生的基督徒.
信了主耶穌.
有神生命的基督徒.
不再犯罪了.
這是一個很大的信仰難題.
任何一個信了耶穌的人不再犯罪.

$^{161}$你說.
未信耶穌之前.
犯罪是很合理.
先信耶穌.
才需要令我們信耶穌.
但是信了耶穌.
從神而生的.
就不再犯罪了.
這就衍生了一個很大的信仰困難問題.
沒有人敢說自己是從神而生的.
因為很難我們在信了耶穌之後.
不再犯罪了.
這是什麼意思呢?.
當然後面的重點.
第十九二十節.
都會繼續說的.
邪惡者不能夠再加害於.
從神而生的一個聖徒.
相信大部分的解經家都認為.
不是從此.
就像你的罪根被拔掉.
不再犯罪.
不會犯罪.
不能夠犯罪.
相信不是這樣.
相信是不會再習慣犯罪.
例如.
上文.
你看看有很特別的地方.
十六至十七節已經說了.
人如果看見弟兄犯了不至於死的罪.
就要為他祈求.
神必定會將生命.
好像再一次賜給犯罪的人.
使他不至於.
好像跟神斷絕了生命的關係.
這樣理解十六,十七節的時候.
回頭看十八節.
我們就不難去理解.
在約翰的認知裡面.

$^{201}$絕對不是說.
從神而生之後.
不能夠或者不會再犯罪.
絕對沒有.
上下文從來都不是這麼說.
反過來就是.
從神而生的.
我們有一個教導學習.
就是什麼呢?.
就是我們看見弟兄犯罪.
不至於死的罪.
我們可以在神的感動.
以及在聖靈的能力下.
我們去為這樣的人祈禱.
神就會聽祈禱的人的呼求.
會將軟弱跌倒犯罪.
這裡特別強調是不至於死的罪.
到底哪些是至於死呢?.
因為十七節我就沒有引述了.
十七節就說.
不至於死的罪.
約翰就說.
我不會認為應該要這樣祈禱.
這裡有很大的神學難題.
超越了今天我們能夠說的時間.
不至於死的罪.
就是說我們肯定了一件事.
信了主耶穌.
有神生命的人.
應該不會再習慣犯罪.
犯罪變成我們會有一種敏銳.
不輕易隨隨便便.
再活在罪惡之中.
而肆無忌憚地生活.
這是約翰告訴我們的.
這就意味著有一件事突顯了出來.
我們有了神生命.
我們要知道一個很重要的真理.
就是基督徒.
原來我們有了神的永恆生命在我們裡面.

$^{241}$我們就開始跟罪有一個很重要的了解.
和厭惡.
和習慣學習不再犯罪.
這是很重要的.
跟一個未有神生命的人.
從來不會有這樣的想法.
這就比我們去莊主香.
得到所謂的任避.
去觀音開庫.
給一千萬那張紙.
你又回來還願.
似乎不是這些.
似乎是我們裡面有一個新生命的發現.
和由信那一刻.
開啟了跟神的生命有份的人生旅程.
這豈不是更加重要?.
你想想任何人照X光,照MRI.
都不會照到自己裡面有罪.
照到可能有其他異常的細胞.
或者器官有其他問題.
但沒有人會照到罪.
唯有從神而生的人.
你就開始對罪有一個敏感的厭惡.
和你不習慣再犯罪.
接著我們下去下面經文.
所以這裡第十六,十七節已經印證了.
不是信了耶穌從神而生而不犯罪.
而是不會再輕言蜕語.
好像未信之前的犯罪.
我們講一些例子.
原來在日常的生活裡面.
我們都可以經驗,認識和遇見神.
我借用一位如果大家認識的.
孫寶齡博士孫寶齡牧師.
曾經是浸上神學院的教授.
後來去了新加坡神學院做院長.
輾轉退休去了台灣的神學院任教.
後來有幾年的時間做教授.
他退休了.
講他的成長,講他的人生.

$^{281}$講他剛剛讀神學的時候.
第一年級.
進入不同教會的同學一起.
聚在一起.
在台灣讀書.
他說頭半年.
有一個神學院的同學.
很生氣一個外籍的教授.
生氣到一個地步.
處處令這個教授很難堪.
孫博士當時是一個神學生.
同班同學.
很不喜歡這個同學.
不知道什麼原因有什麼瓜葛.
後來過了不久.
這個孫博士的同學.
因為某些原因被學校勒令停學.
孫博士當然很開心.
因為他不喜歡.
而且對自己的教授這麼不敬.
他就想著可以在教授面前.
為他出一口氣.
這個外籍教授.
跟孫博士說.
他說這個同學.
一定是在他入學之後.
家庭或者個人有很大的困難.
所以導致他要這樣離開停學.
他說他衷心希望.
也為他祈禱.
希望他快點解決了他生命的問題.
他個人可能不能跟別人說的困難.
以致他能夠有神的感動再入學.
其實很多年.
退休的時候.
孫博士回想回神學的時間.
他說為什麼這位外籍教授.
會這樣對這個學生呢.
不僅不會生他的氣.
反而為他祈禱.

$^{321}$希望他能夠有神的感動.
去走回這個蒙召,事奉的路.
孫博士回想回.
其實他第一年畢業.
他說他因為剛出來.
什麼都不懂.
性格有很多的靈覺.
已經跟別人吵架.
被教會第一年未過試用期.
就解僱.
他想想我的人生就這樣完結了.
我還怎麼可以事奉呢.
神學院畢業出來.
第一年就已經被解聘.
沒有續約了.
但是神感動他.
一路慢慢地成長.
成為今天在華人社會.
非常在解經上讀到的一位牧者.
在我的辦公室裡.
我有他參加全套的書.
很得到那個激勵.
但是神學院一年級的時候.
已經有這樣的生命.
畢業第一年.
他已經被解僱.
你可想而知.
原來我們人生.
我們一路去到成長階段.
其實是不斷對自己的發現.
每一個信耶穌的弟兄姊妹.
其實我們可以在神面前.
不需要隱瞞我們自己的本相.
那個本質.
信仰告訴我們.
我們可以知道.
凡是從神而生的.
必不犯罪.
這個必不犯罪不是不准.
不是不需要.

$^{361}$是可以向罪.
慢慢地我們在神的救恩裡面.
更加認識自己.
更加讓自己的靈命成長.
這是其中一個很重要的知道.
所以弟兄姊妹.
你回教會.
得到的就是對自己的認識.
也不是靠自己.
逼自己不要犯罪.
而是在恩典裡面.
我們已經在神的生命裡面.
那個惡者就不會加害於我們了.
就是說.
我們已經有力量.
面對世界不同的事情.
我們可以有一個抗衡.
我們初信的時候.
因為是福音堂很嚴格的教會.
規定.
走過戲院.
拜託你繞另一條路.
不准看到戲院.
因為當時很多的艷情片.
初信一定要繞過另一條街.
這個教導本身是好的.
不過.
再信久一點你就發現.
其實這樣只是斬腳趾鼻沙蟲.
當然不是說.
那我更加進去看.
不是這個意思.
而是.
我們不是用這種方式逃避的主意.
而是我們要看到.
我們要遇見神.
凡從神生的.
就必不犯罪.
這個屬靈生命的不斷成長.
第二個十九節了.

$^{401}$所以.
保守自己.
是脫離了對罪.
對自己的恐懼.
開放自己讓神在我們的生命裡面.
這是我們必須要在第十八節裡面學習的.
十九節.
我們知道.
第二個.
所以你看到那個修辭.
都已經是張希利文這個知道.
不斷地讓.
修信的讀者.
以弗所教會的弟兄姊妹.
要認識.
我們是屬神的.
而全世界都服在那邪惡者的權勢之下.
剛才十八節就已經在說.
那位邪惡者是不會加害於.
凡從神生的人的了.
第十九節就說.
這裡強調的那個歸屬.
我們在神的陣營裡面.
世界世人.
沒有神的.
就服在邪惡者的權勢之下.
好像有一個隔離主義教導我們.
如果我們學得不好.
我們就姓俗而分.
福音堂就有這個困難.
姓和俗一定要分得很清楚.
不要沾染.
所以我們是不准派利是的.
到今天都是.
我的姨仔和我的妹夫.
我的姨仔的先生.
堪弟.
他們是不准女兒收利是的.
我們派呢.
不行的.

$^{441}$不准收利是.
因為利是就是迷信.
是這樣的.
這樣姓俗而分.
是不是這裡的意思呢?.
會不會帶來一個問題就是.
好像美國.
曾經的隔離政策.
隔離主義.
黑人.
是不能夠和白人坐公共汽車.
讀書是不可以進白人的學校.
原來當年.
馬丁路德金的年代.
五六幾年的時候.
這些路牌是隨處可見的.
這區是白人區.
這一區才是黑人區.
這裡是不是想告訴我們.
我們要有一個這樣的隔離主義呢?.
我們是屬神的.
那些不是屬神的.
也對.
他們是落地獄的了.
我們這一堆就上天堂了.
是不是這樣的意思呢?.
似乎.
約翰不是想灌輸一種這樣的排擠.
這樣的自我的優越.
不是這樣的意思.
說的是歸屬於神.
我們屬於神的人.
留意一下這裡.
我引述的《爾夫所書》四章二十四節.
屬神.
是有這個很重要的新生命.
你們要把自己的心智更新.
並且穿上身我.
這身我是照著神的形象做的.
有從真理來的公義聖潔.

$^{481}$這是神的永恆生命.
在我們裡面.
我們所謂的屬神.
所謂的不再是在世界魔鬼.
邪惡者的那個.
核制,控制,操控下.
是這樣的意思.
是我們來到這裡有一個很主動的意識.
既然我們從神而生.
我們就有這個新的心智了.
舊我的心智.
是要改換一生的.
這是四章二十一到二十三節.
不斷列舉舊的人.
有什麼罪惡等等.
現在四章二十四節就很清楚了.
神的形象讓你想起誰呢?.
是阿當.
一章二十六節.
我們要照著我們的形象.
樣式去做人.
於是就吵了阿當.
所以保羅告訴我們.
原來我們在主耶穌的救恩裡面.
我們就是新的阿當.
新的創造.
這個新的創造不是神再用泥土做.
是用聖靈將我們生入神的家裡面.
我們就是神的子女了.
所以這個新我.
就是要照著神的形象去做.
神的形象就是真理而來的公義.
真理而來的聖潔.
在神的生命裡面.
所以剛才所說的隔離.
不是一種對立性的隔離.
是我們一群在地球上居住的弟兄姊妹.
主耶穌還沒回來之前.
我們都可以孕育神的生命在地上.
藉著我們的見證.

$^{521}$孕育主的愛在不同的地區.
在人還沒見到真神之前.
他們見到我們.
因為我們是照著神的形象做.
就會見到神了.
你想想我們每一個人.
原本我們是不屬於神的.
現在我們可以去代表神.
在這個地上有神的形象,樣式.
所以弟兄姊妹我們或多或少.
在不同的成長過程.
我們都彰顯著神的形象.
所以信仰群體每一個星期回來.
我們就是要教導大家.
我們在一個很恆常的穩定的機制下.
讓大家發現自己是有神的形象,樣式.
去讓更多的人認識神的救恩.
譚文君.
2007年已經離開世界.
他從十四歲開始.
就已經有天生的骨癌.
但他的故事.
在美國長大的.
他爸爸是牧師.
譚進德牧師.
其中出名的地方.
就是他在911紐約世貿大廈.
恐怖襲擊之前.
他已經讀美國的長春藤大學.
但因為他的病多次需要休學.
他的骨癌到了一個地步.
要截肢.
即是說他是一個拿著拐杖上學的.
一個充滿陽光的女孩.
在他的內心.
你見不到任何的埋怨.
爸爸是牧師.
為什麼我會有這樣的重病.
沒有.
完全沒有.

$^{561}$認識他的教授,同學.
都感受到他內心的盼望.
對神的信靠的陽光.
是一位很陽光氣的姐妹.
911出現了.
晚上爸爸沒事回家.
但兩個主學老師在世貿中死了.
他很生氣.
躲在衣櫃裡哭了一整晚.
哭完之後.
很憎恨這些恐怖襲擊.
和這些恐怖分子.
怎麼辦呢?.
他想到了.
我們要去支持現在救援的人員.
爸爸問他.
你怎樣支援?.
你行動不方便.
他突然想起了.
我可以回到少年團.
叫青年人一起參加.
我發起的向唐人街.
去買熱食物.
給消防員去做救援.
六個月.
每天沒有停過.
他自己親自開車.
去送熱食物.
去給這些救援人員.
連續六個月的時間.
911這個女孩.
譚文君.
這種展現出來的生命力.
那種從神而生.
而有神的形象,樣式的人.
基督徒就是這個示範.
讓人看到.
在多艱難,困苦裡.
仍然有神的形象,樣式的人.
讓人不信的人.

$^{601}$沒錯.
全世界真是伏在惡者的底下.
但原來神的能力.
已經在他的兒女.
信耶穌的基督徒的內心.
展現出來.
今早弟兄開車來.
做主席.
是在嘉欣堂的弟兄.
他昨晚問我.
開車來是不是要問.
要說一個密碼.
14座地下.
就可以泊車了.
我說不用了.
因為.
去年.
近年底的時候.
柏宇花園的主席.
業主立案法團的主席.
通知管理處.
告訴我們以後教會來泊車.
就可以時租了.
不用用暗號了.
可以名正言順地泊車了.
還有.
業主主席很快就申請.
我們三樓遲些開業主大會.
很開心的一個.
好像過年的一個.
恩典的祝福和禮物.
讓人看到.
我們一群.
熟神的紅印堂弟兄姊妹.
很多管理員.
大讚我們的弟兄姊妹泊車.
有時失驚無神就給一份禮物.
給一點糖.
給一包餅.
給一瓶水.

$^{641}$這樣當然不是賄賂.
而是一份木輪和諧.
是我們真的感受到.
我們很想將這份真理的仁義聖潔.
讓人看到神的形象.
弟兄姊妹.
我們是熟神的.
全世界都臥在惡者手下.
絕對不是叫我們很負面的.
做一個避難主義.
我們只需要回來這個地方.
感受一下靈氣就行了.
外面好像一個殺戮戰場.
不是這樣.
是將我們的真理.
在不同的地方.
讓人看到基督.
讓人感受到神的愛.
神的真理.
最後第二十節.
我們知道神的兒子已經來到.
並且留意強調這個悟性.
這個知識.
將悟性賜給我們.
使我們認識那位真實者.
我們也在那位真實者裡面.
就是在他兒子耶穌基督裡面.
這是真神也是永生.
相信約翰一路寫的時候.
肯定他對於主耶穌.
在第十七章的約翰福音第三節的祈禱.
認識你獨一的真神.
並且認識你所差來的耶穌基督.
就是永生了.
很特別.
約翰再一次重覆這個語調.
在他裡面認識一件事.
就是在耶穌裡面.
我們有這種永恆.
屬天的知識.

$^{681}$知道我們是神兒女.
神所生和永生.
這是很不簡單的一件事.
當然如果你對整個約翰一書有認識.
當時有一些近乎諾斯底主義的.
否定耶穌基督的神聖的異端.
開始在教會裡面衍生.
所以約翰需要再一次在他有生之年.
親手摸過耶穌的這一位年老的約翰.
來講述.
所以不是告訴你.
我從來都不再洗手.
因為我摸過耶穌摸過神.
不需要這些東西.
不是這些東西.
是我們裡面的那個悟性.
是我們裡面的真知識.
就是對真理的認順的知識.
使我們有這個力量.
知道自己是有永生的人.
還有聖靈已經在我們裡面.
賜給我們這個神的形象.
新的創造.
所以很盼望我們能夠在新的年.
我們一起去學習.
如何將我們所擁有的.
這份悟性,這份認識.
在這位真實者裡面.
來生活.
來讓人看到神的大愛.
很想用這位很重要的改革家.
他是蘇格蘭改革的先鋒.
洛克斯.
被加爾文的影響更加大.
蘇格蘭,英國都受他的影響.
他有一句臨終之前很重要的一句話.
如果一個人.
你認識了神.
你在神的裡面.
你就是大多數.

$^{721}$記住當時的場景就是.
他受著天主教和當時的政權對他的迫害.
基本上是很慘的.
但是他看到這個真理.
唯有因順清義的這個救恩.
叫人可以白白得到神的救贖.
所以他帶來的就是整個蘇格蘭的改革.
他受到當時有三個瑪麗皇后.
其中一個可能大家聽過的.
叫做Buddy Mary.
血腥瑪麗.
英格蘭的一個女王.
對他的迫害很嚴重.
但是在他的內心.
原來我站在神這邊.
我為真理去捍衛.
我才是大多數.
我不是一個孤單的少數.
同學們.
我們去理解這些先賢.
和剛才聖經給我們的教導.
我們要記住.
我們是在神這邊.
我們是屬於神的.
我們當然不是用刀槍對付其他人.
我們是用神給我們的真理.
仁義和聖潔.
彰顯的是神的永恆生命.
願主祝福大家.
我們一起祈禱.
感恩天父讓我們能夠在今早.
我們真的有一個很深的渴慕.
就是要尋找信仰.
這個真實.
求主幫助我們能夠透過.
死徒約翰的三個在他劫獄的時候的知道.
我們來認識我們和神的關係.
還有透過主耶穌的救恩.
我們遠離了邪惡魔鬼的侵害.
我們已經在神的永恆真理的行列.

$^{761}$求主幫助.
使我們內心能夠.
越來越在真理的悟性底下.
有這一份真正的知識.
使我們不會再重複一些迷信的事.
或者追求一些在神不喜悅的事情裡面.
求主幫助.
也讓我們能夠看到.
我們真是一群可以不再犯罪的人.
不是我們有能力不再犯罪.
是主的聖靈在我們內心.
使我們能夠可以向罪說不.
願主幫助.
祝福我們每一位.
奉主耶穌基督的名.
阿們.
\newpage



\section{腓立比書 }
\label{sec:ulbWKYhlm8s}
\textbf{2025.02.16 《做一個內心強大的基督徒》 《腓立比書》4-10-15 (和修版) 李文耀牧師}
\newline
\newline
連結: \href{https://youtube.com/watch?v=ulbWKYhlm8s}{\texttt{https://youtube.com/watch?v=ulbWKYhlm8s}} ~~~~ 語音日期: 2025-02-17
\newline
\newline
\hyperref[sec:Ho71_1KRGq8]{\small{< < < PREV SERMON < < <}}
~
\hyperref[sec:index_chronic]{\small{[返順時目]}}
~
\hyperref[sec:index_scriptual]{\small{[返順卷目]}}
~
\hyperref[sec:TScbhjJAAD8]{\small{> > > NEXT SERMON > > >}}
\newline
\newline
$^{1}$大家好,我是大英哲學,很感恩能夠在節目中分享上帝的話語.
這是我長這麼大第一次來洪水橋.
事奉神有好處,可以去一些從來未去過的地方.
今天的題目是因為有一天在社交平台上看到一則貼文.
這則貼文上寫著做一個內心強大的人.
下面放了苗僑偉一張很有型的照片,我們忍不住點進去看看.
裡面是這樣寫的.
如果別人說你兩句你就受不了.
被兩句話干擾得吃不好睡不好.
你得有多脆弱你要明白.
能干擾你的往往是自己太在意.
能傷害你的往往是自己想不開.
你若和平無人能擾.
所以不要過於敏感.
內心要不斷地強大.
別被別人左右你的情緒.
有心者有所慮.
無心者無所謂.
做一個內心強大的人.
走自己的路.
讓別人去說罷.
看完之後我就想到這個講道的題目和內容.
一般人都知道內心強大是很重要的.
信主的人更加要不斷提醒自己.
我們要做一個內心強大的基督徒.
能不能做到呢?.
或者怎樣做呢?.
我們要先問一問內心脆弱有什麼問題呢?.
我們每個人都多不少會受外在事物影響.
例如別人說一兩句話或做一兩個動作.
都會影響到我們的情緒.
內心脆弱的人更容易受到外在環境的刺激.
很容易憤怒.
睡不著.
有時有些有興趣的事.
都變成沒興趣沒動力去做.
嚴重些可以引發一些極端的行為反應.
要將困擾自己的人或事徹底剷除.
要麼剷除他者.
要麼消滅自己.

$^{41}$剷除他者方面.
我想到在聖經裡面有一個故事.
我想大家都熟悉的.
就是該隱的故事.
記載在創世記第四章.
那裡說到該隱因為上帝看中了他弟弟阿伯的供物.
看不中自己的供物就怎樣呢?.
非常生氣.
到一個地步.
有一天在田間就將他弟弟阿伯殺死.
為何會這麼生氣呢?.
生氣到一個地步.
要將他弟弟阿伯殺死.
背後可能夾雜著其他因素.
例如他自我形象低落.
他內心缺乏安全感.
或長期得不到他父母的肯定.
有很多背後心理因素在影響著.
如何解決他內心的困擾呢?.
該隱採取了一個極端的消滅他者的方法.
只要導致困擾情緒的人在我眼前消失.
一切就會回復正常和平靜.
其實我們回看整個故事.
上帝從來沒有說過不接納該隱.
只不過在那次不知什麼原因.
他看中了阿伯的供物.
內心脆弱的人有一個顯著的特徵.
就是很容易將一件事放大.
例如我這個人.
有沒有遇過這樣的事?.
我明明是說這件事.
他就會覺得你只是說他這個人不好.
那裡不行.
因為一次看不中阿伯的供物.
該隱就覺得他被上帝針對.
是上帝偏心.
是上帝針對我.
是上帝不喜歡我.
其實該隱比我們更加緊張.
在意上帝是否看重他.

$^{81}$如果一個人不看重上帝.
其實是無所謂的.
你沒看中我弟弟的供物.
無所謂.
正正因為很在乎上帝怎樣看待他.
於是上帝因為這一次看不中他的供物.
他就非常生氣.
該隱的問題不是他不看重上帝.
他很看重上帝.
他的問題是他內心不夠強大.
去克制他內心的嫉妒和憤怒的情緒.
最後用了一個非常極端的消滅他者的方法.
來使自己內心平復下來.
另一個極端的行為就是.
消滅自己的方法.
他者消滅不到.
最後就是以消滅自己的方法.
使內心平靜下來.
聖經有沒有一些例子說到這方面呢?.
我想到的是.
以利亞在羅藤樹下求死的故事.
我想大家也聽過.
記載在聖經上的十九章一到八節.
那裡說到.
以利亞在加密山上.
他一個人抵擋450個巴力的先知.
後來將他們通通殺死了.
很厲害吧.
一點畏懼一點恐懼也沒有.
一個對著450個.
不是一個打10個.
一個對著450個.
一點恐懼畏懼也沒有.
很厲害.
不過這件事之後.
當耶駒別.
大家知道他是誰嗎?.
當時亞哈王的妻子.
他知道了之後就派人去追殺以利亞.
以利亞突然間很害怕.

$^{121}$然後他就逃跑.
去了曠野.
然後在一棵羅藤樹下求死.
是不是很大的差別?.
一個對著450人也不害怕.
現在只不過有一個風聲.
皇后耶駒別要追殺他.
有多久也沒有殺到馬利生.
但是他突然間很恐懼.
他起來逃命並且要求死.
每一個人無論有多厲害有多厲害也好.
也會有脆弱的時候.
越是成功的人.
就越難去面對人生的挫敗挫折.
於是最後他會想到死亡.
死亡就是一個最後的解決.
解脫自我肯定自我的最後手段.
困擾自己的他者解決不了.
最後就唯有解決自己.
這是他最後肯定自己的一個手段.
其實自殺的行為也是這樣.
自殺就是最後一條命也是我的.
是我去決定我自己的命運.
這兩個例子就告訴我們.
內心脆弱的人.
其實可以有很多激進的行動產生出來.
當然我們日常不會有這麼激進的行動.
但內心脆弱也會帶來很多心理各方面的人際關係的影響.
於是如何成為一個內心強大的人.
真的是一個很重要的問題.
外面的人一般的建議就是.
叫我們看開一點.
有心者有所慮.
無心者就無所謂.
基督徒有什麼方法.
基督徒我們說要做一個內心強大的人.
我們有什麼方法.
我們所說的和外面不相信的人.
有沒有什麼不同呢.
有沒有什麼分別呢.

$^{161}$我們是不是只說看開一點呢.
是不是只有這樣呢.
聖經裡面其實有很多內心強大的人物.
例如大衛摩西.
今天時間關係.
我就特別選擇了保羅和大家來思考一下.
保羅的內心如何強大呢.
為什麼可以這麼強大呢.
今天就選擇了.
菲勒比書四章十到十五節.
這段經文和大家來看一看.
剛才主席已經讀了.
我就不再重複和大家讀一次.
保羅是在什麼處境下.
寫下這段文字呢.
我們大概知道.
菲勒比書的開頭就說到.
他是為基督的緣故受困所.
有預認的全軍在看守.
就是說他是在坐牢.
當時保羅在哪裡坐牢呢.
學者估計有幾個可能性.
一是羅馬.
一是蓋薩尼亞.
一是以弗所.
一是哥林多.
不過無論什麼地方都好.
坐牢絕對不是一件開心的事.
坐牢始終是一個叫人痛苦的經驗.
譬如沒有了活動的自由.
沒有了正常社交的活動.
不是說自己喜歡做什麼就做什麼.
睡也睡不好.
吃也吃不好.
還有很多監倉裡面的潛規則.
你要遵守的.
旁邊的囚犯.
可能脾氣不是很好.
所以在監獄裡面生活.
其實一點都不容易.

$^{201}$很痛苦的.
我自己就沒坐過牢.
不過我看監獄風雲看了很多次.
我不知道當時保羅的情況是好一點還是差一點.
不過坐牢是痛苦的.
是不容易的.
不單止他要面對坐牢的痛苦.
書信的開頭說到.
他當時還要面對一些說話的攻擊.
外面有很多閒言閒語.
有很多壞話在四處散播.
保羅說.
有些是出於質度.
有些是出於自私動機不順.
企圖要加增他受困所的苦楚.
保羅當時要面對很多造謠生事.
很多壞話.
很多捏造的.
不真實的說話在攻擊他.
痛不痛苦啊?.
頂尖有沒有試過被人說壞話?.
有沒有試過在外面被人說一些說話.
其實那些是不真實的.
但是不斷在說你.
有沒有試過?.
很痛苦啊.
很多年前我去德國讀書的時候.
曾經幫過一個學生.
因為他要找一個地方住.
當時大學已經沒有位置了.
我們那個住所有一間房.
可以暫時借給他住.
那間房平時我是拿來去溫習.
去做研究去寫作.
不過見他找不到地方.
他又是住在內地的人.
我們就說讓他在那裡住一段時間.
我們說不收費的.
等他找到一個地方.
他就搬出去了.

$^{241}$大概住了一個月之後.
他就搬出去了.
後來沒多久.
我從一些朋友聽到.
不知道為什麼這位.
我們曾經幫過的朋友.
主內的學生.
弟兄姊妹.
竟然在外面.
用一些不太好的說話.
去說我們.
其實我們都沒有對他做什麼.
他說很多說話不真實.
我們聽完之後.
很痛苦.
我們是無償地幫他.
為什麼他出去.
會這樣去說一些.
重傷我們的說話.
我們做錯了什麼.
被人說一些壞話.
是很難受的.
保羅他為基督的緣故.
去傳福音是一件好事.
為什麼他坐牢的時候.
會有這麼多說話在攻擊他.
其實是很難受的一件事.
內心脆弱一點.
應該都崩潰了.
撐不住了.
保羅當時他怎樣去回應.
去面對這個處境.
在四章的第十節.
首先保羅說.
我靠主大大喜樂.
在原文裡面.
那個說話應該是說.
我在主裡面大大喜樂.
I rejoice in the Lord greatly.
在主裡面大大喜樂.

$^{281}$他當時是一個.
極度艱難困苦的裡面.
但是他靠主.
在主裡面仍然有喜樂.
不是小小的喜樂.
是大大的喜樂.
面對煩惱.
面對挫折的時候.
外面的人一般的提議.
就是我們要看開一點.
不要想一二邊.
但我們都知道.
其實能不能做到.
做不到吧.
因為我們就是.
不能夠看開一點.
大家有沒有試過失眠.
有沒有試過睡不著覺.
如果身邊的人跟你說.
不要想那麼多.
安靜就能睡了.
有沒有用.
沒有用吧.
內心平平靜下來.
那個是情感的問題.
不是思想的問題.
不是說我的思想.
我的意志說要靜下來.
不要想那麼多就行了.
情感的問題只能夠通過.
情感的方法去處理.
得不到尊重的人.
你要讓他感受到尊重.
得不到安全感的人.
你要讓他感受到安全感.
一個長期得不到肯定的人.
你要讓他經歷到愛和接納.
情感的問題只能夠通過.
情感的方法去處理.
為什麼保羅被人抓去坐牢.

$^{321}$在外面有很多說話攻擊他的時候.
他仍然可以大大喜樂.
因為他在主裡面.
經歷到那種平安.
是大大蓋過他的苦楚.
丙姐妹.
喜樂是一個經歷的問題.
喜樂是主賜下的一個經歷.
不是頭腦和理智上的問題.
大家吃完飯回家.
試試跟自己說三次.
我要喜樂.
看看能不能立刻喜樂.
是不行的.
喜樂是要在主裡面.
是主賜下來的一個經驗.
主耶穌和門徒說過.
在《讓福音》十五章.
耶穌說我是葡萄樹.
你們是枝子.
常在我裡面.
我也常在他裡面.
你了我.
你們不能作什麼.
為什麼有些事情我們做不到.
是不是因為能力不夠.
是不是因為資源不夠.
因為我們人脈不多.
這些都是原因.
不過更多時候.
我們之所以有些事情做不到.
因為我們內心出現了問題.
有焦慮和恐懼.
本來那件事情很簡單.
很容易處理.
因為我們內心有恐懼和焦慮.
那件事情就變得很困難.
做不到.
為什麼以色列人的先祖.
要在曠野漂流四十年.

$^{361}$大家都知道吧.
因為有十二探子去了迦南.
有十個回來回報的時候.
那些人很大隻.
很厲害.
我們好像蚱蜢一樣.
本來以色列人要得到迦南.
其實很容易.
因為上帝應許了.
耶利哥的城是很難攻陷的.
走七個圈已經搞定了.
其實很容易的.
隨手可得.
為什麼最後他們要漂流四十年.
因為內心恐懼.
內心焦慮.
他們忘記了上帝是他們最終的依靠.
為什麼人會失敗.
因為我們內心出現了問題.
我們有恐懼.
我們會焦慮.
我們覺得很無助.
甚至我們覺得絕望.
前面沒有路可以再走下去.
怎樣去面對這些內心的情緒.
當然.
見父道是一個方法.
但更加重要的.
就是在主裡面.
通過與主結連.
好像葡萄樹與枝子一樣.
在主裡面經歷祂所賜的平安和愛.
在那裡克服我們內心的恐懼和焦慮.
耶穌曾經應許過我們.
在《約翰福音》14章27節.
耶穌說我要將我的平安賜給你們.
我所賜的平安.
不像是世人所賜的.
你們心裡不要憂愁.
也不要膽怯.

$^{401}$我們每個人都經歷過膽怯的時候.
當憂愁或膽怯來到的時候.
主耶穌應許過.
祂說祂所賜的平安喜樂.
就會給我們力量去面對去克服.
耶穌可以通過很多的渠道.
給我們平安和喜樂.
可以通過一首詩歌.
通過一篇訊息.
通過一個崇拜.
可以通過很多不同的渠道.
祂會賜下平安和喜樂.
在德國進修那幾年.
曾經也試過有迷茫和焦慮的時候.
特別是頭一年是最艱難的.
人生路不熟.
語言又不通.
不容易.
曾經想過撿包袱回香港.
不過又想起很多靈姐妹.
背後有支持和祈禱.
這樣回頭好像不知如何交代.
很多內心的掙扎.
很多困難和困擾.
試過很多次.
人生的低谷.
很想放棄的時候.
耶穌就通過一首詩歌.
一個敬拜.
讓我重新得力.
在一個很沮喪的情緒裡.
重新靠主堅持下去.
去面對這些困難.
喜樂是一個經驗的問題.
情感的問題只能通過情感的方法去處理.
於是當我們覺得很累.
覺得很沒力量.
覺得很焦慮的時候.
我們應該更加要回教會.
在教會繼續和耶穌有一個連結.

$^{441}$然後等候主耶穌基督.
祂親自來賜下平安.
愛和喜樂.
這是第一個重點.
靠主大大喜樂.
到了一個程度.
我們各樣的痛苦都可以被蓋過.
被克服.
在主裡面靠著主.
第二方面.
為何保羅內心這麼強大.
可以面對各樣的攻擊和挑戰.
因為他看到身邊有一群支持他.
明白他.
能夠和他分擔憂患的弟兄姊妹.
這個體見也是觸發保羅在主裡面.
大大喜樂的一個原因.
剛才說到我們經歷主親自賜下的平安.
會給我們有喜樂.
這是真的.
不過更多時候.
主耶穌是通過弟兄姊妹的關懷和同行.
讓我們更加具體去經歷主的同在和看顧.
四章十節的下半節.
我靠主大大喜樂.
因為你們關懷我的心.
如今又表現了出來.
其實你們一直都關懷我.
只是沒有機會罷了.
十四節又再說一次.
然而你們能夠和我分擔憂患是一件好事.
他看到身邊有弟兄姊妹和他能夠分擔憂患.
肥仔比教會是保羅在第二次宣教旅程所建立出來的教會.
根據《小童行傳》十六章十四節記載.
第一個信徒應該是一個姐妹.
她叫做呂底亞.
大家有印象吧.
呂底亞應該是一個商人.
一個有錢人.
因為聖經記載她是一個賣紫色布的婦人.

$^{481}$紫色在當時的社會文化裡.
是屬於君王貴就所穿的顏色.
就等於以前黃色和金色是屬於什麼人去穿的.
皇帝皇族的人去穿的.
她賣的衣服只有有錢有地位的人.
皇族的人才可以買.
所以她應該是一個有錢的商人.
除了呂底亞之外.
當時記載河邊還有一群婦女聽保羅的講道.
我們估計肥仔比教會應該大部分都是姐妹.
看上去就像現在這樣的情況.
不過通常都是香港教會.
香港教會的真牧師河主要是姐妹多一點.
男性弟兄好像比較少.
當時是這樣的.
肥仔比書四章二節記載有兩個姐妹.
有阿諜和秦都基.
保羅在那裡勸勉他們要在主裡同心同腹一握.
大家看到那裡應該都意會到.
這兩個姐妹應該有些爭執.
如果沒有爭執保羅也不用特意勸勉.
這兩個姐妹有阿諜和秦都基.
要在主裡同心.
應該有些東西出現了.
這兩個姐妹的爭執應該會帶來教會有一定的影響.
以至他們不能夠同心見證基督傳福音.
我真的沒有性別歧視.
不過大家知道姐妹吵起來都挺糟糕的.
所以保羅要特別寫信勸勉他們要在主裡同心.
二章二節提醒他們.
這群人的信徒要義志相同.
愛心相同.
有一致的心思.
有一致的想法.
無論當時教會面對什麼情況.
這兩個姐妹的爭執如何影響其他弟兄姊妹.
但在某一件事上他們是非常同心的.
什麼事呢?.
就是幫助保羅同行.
跟他分擔憂患這件事上.

$^{521}$他們是非常同心的.
即使兩姐妹有爭執.
即使弟兄姊妹在教會有很多意見不和.
在這件事上他們教會是非常同心的.
弟兄姊妹如何跟他們分擔憂患.
四章十五節提到在收支的事上.
除了他們之外並沒有別的教會跟他們分擔.
這裡說的很明顯是什麼?.
是一些很物質上的支持.
應該是金錢上的支持.
菲律賓教會是支持他們去不同的地方傳福音.
在這件事上.
保羅看到他們的支持.
鼓勵和同行.
四章十七節提到.
菲律賓教會對保羅的關懷.
如今又表現出來.
即是說有一段時間沒有了.
不過現在他在監獄落難的時候.
保羅再一次收到菲律賓信徒的支持.
問安和祝福.
即使在監獄裡面要面對很多困難.
即使在外面有人抹黑他,詆毀他.
保羅仍然在主裡面大大喜樂.
除了他經歷到主給他的平安和愛之外.
他看到身邊還有一群弟兄姊妹.
來跟他同行支持他,明白他.
以利亞在羅藤樹下求死.
因為他覺得他自己一個人去面對這個處境.
只有他一個人為上帝爭戰.
保羅在這個處境裡仍然大大喜樂.
因為他在困苦裡看到還有人明白他.
他不是孤身一個人.
還有人支持他,跟他同行.
弟兄姊妹,有人曾經說過.
只要身邊有人仍然信任我們.
仍然肯定我們,對我們有期待.
我們就沒有放棄自己的理由和藉口.
咬緊牙關,一步一步走下去.
總會走出那個難關.

$^{561}$不要氣餒,不要放棄.
身邊還有弟兄姊妹,還有人明白我們.
我們不是一個人,我們不是孤單.
我們不是獨力去承受這一切的痛苦.
我們要看到還有人跟我們同行.
有同行的人,我們在主裡面.
我們會經歷到平安和喜樂.
這是基督徒沒那麼容易被打敗的原因.
我們不是孤單,在教裡面有同行者.
有明白我們的人.
最後,保羅之所以內心強大.
因為他領悟了一個生活的秘訣.
四章十二到十三節.
我知道怎樣處悲戰.
也知道怎樣處豐富.
或飽足或飢餓,或有餘或缺乏.
任何事情,任何境況,我都得了秘訣.
我靠著那家給我力量的,凡事都能做.
他找到一個秘訣.
他領悟了一個秘訣.
就是靠著那家給我力量的,凡事都能做.
這節經文對我影響很大.
在二十多年前,當我在建道畢業的時候.
去到一間教會面試.
有一個執事問了我一個問題.
他問我,你覺得自己有什麼條件.
去升任做這個傳道牧者的工作?.
很難回答.
如果你說,我什麼都可以,我可以的.
就很驕傲.
如果說,我不行的,我水皮很差.
別人就不會聘請你.
很難回答這些問題.
當時就靈機一觸,想到這節經文.
靠著那家給我力量的,凡事都能做.
然後就聘請了我.
現在我還在這間教會幫忙,教導.
凡事都能做.
靠著那家給我力量的,凡事都能做.
這句話的意思不是說什麼都做得到.

$^{601}$什麼都做得出色.
這裡不是說成功神學.
不是說我們信耶穌,什麼都可以.
什麼都順利,心想事成.
當然,新年期間我們有時候也說這些話.
但我們知道不是的.
信耶穌不是凡事順利,心想事成.
如果是這樣,保羅就不會為了主的緣故被抓去坐牢.
這節經文也不是說上帝會幫我們達成心中的夢想.
凡事都能做.
不是說上帝幫我們實現我們的夢想.
今天很多人都說夢想.
有時候走過地鐵,或者看到標語.
都是說要追求我們的夢想.
叫我們一生無憾.
夢想現在不只是年輕人,中年人也行.
大家有沒有看過中年好聲音3.
有一次我太太師母跟我說.
你也唱了幾句,不如你去參加.
我就說不要玩了.
有沒有上過牧師,參加中年好聲音.
雖然她可以賦予我見證神,但也不喜歡.
這裡不是說上帝幫我們實現自己的夢想.
保羅從來都不是追求自己的夢想.
你看回《一章20節》.
保羅說凡事坦然無懼.
無論是生是死,總要讓基督在我身上照常顯大.
保羅做任何事情,不是為了自己的夢想.
他不是要上帝滿足他的期待.
他做的一切都是在傳揚基督這個大前提下.
然後他說凡事靠著上帝都能應付,克服.
凡事都能做,這裡不是說凡事順利.
不是說成功神學.
這裡也不是說上帝要達成我們個人的夢想.
而是我們為主所做的一切.
無論遇到什麼情況.
上帝都會給予我們力量應付,面對.
保羅有這個領悟,並不是一朝一夕的事.
他是經過長時間,經過很多的經歷.
才有這個肯定和確信.

$^{641}$在整個事奉生涯裡.
保羅告訴我們,他曾經被抓去坐牢.
他曾經受過無數次鞭打.
他曾經經歷過海難,盜賊的危險.
他曾經忍受的勞苦比我們要多.
他經常睡不著覺,又飢又渴,赤身怒體.
他經歷了很多事情.
他經歷了很多事情之後.
然後得出了人生的一個秘訣.
靠著那加給我力量的.
凡事都能做.
我在1994年夢召,我感到很感動.
回應這個呼召,踏上事奉的人生.
1994年到現在都差不多30年了.
多不多?長不長?.
OK啦,肯定沒有曾牧師那麼長.
但這三十多年,我都經歷了不少的事情.
今天沒有時間和大家分享.
在二十多年前,在應徵做傳奴的時候.
我講這句話和現在我再講這句話已經不同了.
這是我的真實的體驗,真實的經歷.
弟兄姊妹,信仰是要用經驗去沉澱的.
信仰是要用經歷去累積的.
沒有經歷,沒有沉澱,沒有實踐的信仰.
耶穌用過一個比喻去講.
就好像一個無知的人將房子蓋在沙土上.
風吹雨打,水沖撞擊,就會怎樣?.
全部倒塌下來.
信仰是要用人生的經歷去累積出來.
人生有起起跌跌.
不只是風乾的時候去經歷上帝.
在我們最失意的時候都可以經歷上帝.
經歷了人生幾十年之後.
回頭再看,我們就知道上帝的恩典處處.
什麼可以讓我們內心強大.
面對任何的衝擊,我們都能夠屹立不倒.
是因為我們相信上帝的保護.
我們相信靠著哪一家給我們力量的凡事都能做.
是真的.
過去神學院有幾年.

$^{681}$我不知道大家有沒有聽到我們一些故事.
特別是一九年之後都經歷過很多的挑戰.
有社會運動,有新冠疫情.
還有我們學院裡面的人事風波.
那幾年真的有很多的問題.
剛好當時我接手做了教務長的位置.
真的痛苦到不得了.
天天要開會,天天都要危機應對.
因為沒有試過這樣的情況.
有老師跟我這樣說.
他想安慰我,其實也不知道是不是安慰.
他說見到幾十年都沒有試過這樣的情況.
不過靠著上帝給我們的智慧能力.
總算渡過了.
昨天我們見到有個開放日.
我們紀念一百二十六年的歷史.
見到經歷了一百二十六年很多的戰爭,很多的風波.
到今天仍然屹立不倒.
不是我們有什麼本事,而是上帝的本事.
保羅之所以能夠面對逆境.
無論在什麼情況下都可以知足.
可以內心強大.
這段經文提醒我們有幾點.
第一,要在主裡面經歷喜樂.
是主賜下來的平安和喜樂.
我們要看到除了我們去承擔這些困難之外.
還有一群同行者明白我們.
他們會陪我們一起去渡過.
他們會為我們分擔各樣的痛苦.
還有最重要的是我們要知道.
上帝一直看顧著我們.
靠著他我們什麼事情都可以克服和面對.
遇到挫折,有人一敗塗地.
有人逆境自強,人生起起跌跌.
有順境的時候,也有逆境的時候.
外在的環境我們控制不了.
不過我們可以叫內心強大一點.
帶著喜樂,感恩,使命去面對我們的人生.
不信上帝的人都可以內心強大.
基督徒就更加可以內心強大.

$^{721}$無論得時不得時,無要全道.
無論得時不得時是什麼?.
無論順境和逆境,無要全道.
無論得時不得時什麼處境都可以去全道.
因為我們的內心夠強大.
求主幫助我們,我們一起祈禱.
天主上我們在人生裡的確會面對很多風波.
很多挫折挑戰.
有很多事情不是我們想的.
也有很多事情不是我們自己造成.
有些是我們的責任.
但有很多都不是我們可以控制得到.
求主幫助我們,無論我們遇到什麼情況.
我們都靠主來有喜樂平安的心去面對.
幫助我們做一個內心強大的基督徒.
無論在什麼處境當中.
我們都願意來向人分享主的恩典和愛.
因為在外面實在有太多的人在流離,困苦的裡邊.
他們需要福音,需要主耶穌.
求主幫助我們.
我們這樣簡單地祈禱.
奉主耶穌之名.
阿們.
\newpage



\section{哥林多後書 8:1-15}
\label{sec:TScbhjJAAD8}
\textbf{2025.02.23 《信徒愛心的實在》 哥林多後書8\_1-15 (和修版) 講員 : 陳淑娟牧師}
\newline
\newline
連結: \href{https://youtube.com/watch?v=TScbhjJAAD8}{\texttt{https://youtube.com/watch?v=TScbhjJAAD8}} ~~~~ 語音日期: 2025-02-25
\newline
\newline
\hyperref[sec:ulbWKYhlm8s]{\small{< < < PREV SERMON < < <}}
~
\hyperref[sec:index_chronic]{\small{[返順時目]}}
~
\hyperref[sec:index_scriptual]{\small{[返順卷目]}}
~
\hyperref[sec:QcQ_nHz9bjU]{\small{> > > NEXT SERMON > > >}}
\newline
\newline
哥林多後書 8:1-15
\newline
\begin{longtable}{cl}
\hline
\hline
章節 & 經文 (和合本修訂版)\\
\hline
8:1 & \begin{tabularx}{0.7\textwidth}{X} 弟兄們，我們要把神賜給馬其頓眾教會的恩惠告訴你們： \end{tabularx} \\ \\ \relax
8:2 & \begin{tabularx}{0.7\textwidth}{X} 他們在患難中受大考驗的時候，仍然滿有喜樂，在極度貧窮中還格外顯出他們樂捐的慷慨。 \end{tabularx} \\ \\ \relax
8:3 & \begin{tabularx}{0.7\textwidth}{X} 我可以證明，他們是按著能力，而且超過了能力來捐助，主動 \end{tabularx} \\ \\ \relax
8:4 & \begin{tabularx}{0.7\textwidth}{X} 再三懇求我們，准他們在這供給聖徒的善事上有份； \end{tabularx} \\ \\ \relax
8:5 & \begin{tabularx}{0.7\textwidth}{X} 並且他們所做的，不但照我們所期望的，更照神的旨意先把自己獻給主，又給了我們。 \end{tabularx} \\ \\ \relax
8:6 & \begin{tabularx}{0.7\textwidth}{X} 因此，我們勸提多，既然在你們中間開始這慈善的事，就當把它辦成。 \end{tabularx} \\ \\ \relax
8:7 & \begin{tabularx}{0.7\textwidth}{X} 既然你們在信心、口才、知識、萬分的熱忱，以及我們對你們的愛心上，都勝人一等，那麼，當在這慈善的事上也要勝人一等。 \end{tabularx} \\ \\ \relax
8:8 & \begin{tabularx}{0.7\textwidth}{X} 我說這話，並不是命令你們，而是藉著別人的熱忱來考驗你們愛心的真誠。 \end{tabularx} \\ \\ \relax
8:9 & \begin{tabularx}{0.7\textwidth}{X} 你們知道我們主耶穌基督的恩典：他本是富足，卻為你們成了貧窮，好使你們因他的貧窮而成為富足。 \end{tabularx} \\ \\ \relax
8:10 & \begin{tabularx}{0.7\textwidth}{X} 我在這事上把我的意見告訴你們，是對你們有益，因為你們開始辦這事，而且起此心意已經有一年了。 \end{tabularx} \\ \\ \relax
8:11 & \begin{tabularx}{0.7\textwidth}{X} 如今就當辦成這事，既然有願做的心，也當照你們所有的去辦成。 \end{tabularx} \\ \\ \relax
8:12 & \begin{tabularx}{0.7\textwidth}{X} 因為人只要有願做的心，必照他所有的蒙悅納，並不是照他所沒有的。 \end{tabularx} \\ \\ \relax
8:13 & \begin{tabularx}{0.7\textwidth}{X} 我不是要別人輕鬆，你們受累，而是要均勻： \end{tabularx} \\ \\ \relax
8:14 & \begin{tabularx}{0.7\textwidth}{X} 就是要你們現在的富餘補他們的不足，使他們的富餘將來也可以補你們的不足，這就均勻了。 \end{tabularx} \\ \\ \relax
8:15 & \begin{tabularx}{0.7\textwidth}{X} 如經上所記：多收的沒有餘，少收的也沒有缺。 \end{tabularx} \\ \\
[1ex]
\hline
\hline
\end{longtable}
$^{1}$紅映堂的弟兄姊妹 早安.
主內平安.
感恩牧師邀請我們.
來到這裡分享.
其實不是很遠.
我發覺從銅鑼灣過來.
一個小時就來到.
我以前住紅水橋.
所以我都熟悉.
一同聽主的說話.
我們先一同禱告.
天上的八父.
願你的聖靈充滿著我們.
願今天我們能夠開我們的眼睛.
願我們看得到.
上帝要我們所作的.
讓我開我們的耳朵.
聽得見在水深火熱的人.
我們如何與他們同行.
開我們的心.
讓我們聽道能夠行道.
讓我們所作的都是蒙你的喜悅.
我們將以下時間.
恭敬交託 以上不完全的禱告.
奉耶穌的名祈禱 阿們.
我知道紅映堂.
在關懷這個社區裡.
不為餘力.
所以今天我希望藉著這篇經文.
作為彼此的鼓勵提醒.
去看看.
甚麼是真正的愛倫如己.
這篇經文是在.
哥林多後書.
八章一到十五節.
先讓我介紹一下.
哥林多其實八章到第九章.
保羅一直討論.
一個重要的事項.
就是周濟耶路撒冷貧窮的聖徒.

$^{41}$所以很多的.
牧者都會.
用這段經文討論.
主日奉獻的事.
其實保羅所討論的.
有更深的神學意義.
他就從.
一個行動的角度.
去看愛心這件事.
他說愛心.
是甘心樂意.
不被迫不勉強.
最重要保羅認為捐獻這件事.
不是單單善舉.
那麼簡單.
從神學意義來看.
其實就是在討論.
在成聖的過程中.
如何與主同工.
領受上主的應許和信實.
今天我就從這幾方面.
和大家一同分享.
我可以自己按.
那個powerpoint.
第一個分題.
叫做愛心的原則.
我們從.
馬其頓宗教會裡面.
我們發現他捐獻的時候.
有一些屬靈的愛心原則.
我們先看看這些屬靈的愛心原則.
我邀請大家.
我們一起讀.
第一到第五節.
讓牧師節省一點氣.
第一到第五節.
預備起.
弟兄們.
我們要把神給馬其頓.
的恩惠.

$^{81}$告訴你們.
他們在患難中.
受大考驗的時候.
仍然滿有喜樂.
在極窮貧窮中.
急.
(念經).
.
好 先來到這裡.
在這短短的五節裡.
我們發現了愛心屬靈的原則.
第一點 就是超越限制.
剛才讀的時候.
不知道大家有沒有覺得很奇怪.
他們正在患難中受到大的考驗.
又在極度貧窮中.
就是說馬其頓教會正處於什麼狀況呢.
就是又窮 又有逼迫.
但是他們的快樂是開心的.
因為他們額外樂意的捐獻.
這件事很奇怪.
跟我們正常人所想的很不同.
首先馬其頓宗教會.
我先介紹一下.
其實他們是.
肥納比教會.
帖撒羅尼迦教會.
和比利亞教會.
如果根據帖撒羅尼迦前書.
和肥納比書.
你會知道這間教會.
不是富有平穩的教會.
如果用我們現在香港.
或者華人教會來看.
不是那些中產教會.
不是那些巨型教會.
甚至他們人數不多.
不但如此.
他們當時的環境正在受到迫迫.
他們是極之窮困的.

$^{121}$這跟我們想的邏輯有些不同.
當我們很窮困的時候.
我們內部都還沒處理好.
又窮人又不多.
又受迫迫.
通常我們第一時間想什麼.
先處理好自己.
我去很多教會.
做教育的工作.
同行的工作.
過往我和大約110間教會同行過.
在上一年我和27間教會同行.
最初和他們聊天的時候.
大部分教會都說.
牧師我們現在怎樣幫人.
怎樣和人同行.
第一移民潮.
我走了三分一.
第二現在嬰兒潮.
每一間教會都老齡化.
很難見到一間教會.
大部分都是年青人.
我們都要面對現實.
第三我們內傷.
我們經歷社會事件.
我們有內傷.
自己都還沒處理好.
不要搞那麼多.
我都聽過很多這樣的狀況.
這都是我們要面對的事實.
我沒有否定這些事實.
事實上我們是要面對的.
不過我想從馬克龍宗教會的例子.
我們看到.
縱然他們當時在患難裡面.
縱然他們很窮困.
但是他們依然超越自己的界限.
不單止盡自己的能力.
因為你看看第三節.
他說.

$^{161}$我可以證明他們是按著能力.
而且超過.
甚至超過自己的能力.
去和別人.
耶路撒冷的貧窮信徒同行.
我們看到他們的超越界限.
就是他們沒有計算自己的能力.
也沒有看當時的環境.
他們看什麼.
如果你繼續看下去.
你知道他們只是看著神的恩典.
他靠的是神和他們同工的恩典.
他沒有計較任何的環境.
包括他們當時的貧窮和自己的不足.
他們都樂於捐助別人.
其實很多時候.
當我們越沒有能力的時候.
我們越願意去做的時候.
往往顯出救恩的寶貴和能力.
當然這裡也直接告訴我們.
愛心的捐獻和錢有沒有關係.
這是你願不願意的問題.
他們如何超越界限.
其實他們超越了兩個界限.
第一個界限.
馬其頓宗教會是外邦教會.
耶路撒冷是十萬九千七丈遠.
不像現在我們.
如果找牧師來幫忙.
也算是很近.
就算去銅鑼灣幫忙也很近.
因為現在只坐車.
去美國幫忙也很近.
但當時沒有電郵.
沒有長途.
沒有IDD.
你幫了他.
其實你這輩子也不會看到耶路撒冷的信徒.
第一他超越了一個地域的界限.
一個陌生人.

$^{201}$也不打算有回報.
只是一個陌生人.
第二就是超越了自己能力的界限.
這個值得我反思.
正如我剛才說.
我去很多教會分享.
關於如何進入社區.
與有需要的人同行的時候.
大家第一個迷思就是.
我們都沒有資源.
我們又沒有人.
我們又沒有很多的錢.
我們如何投放這麼多人.
甚至不要說教會.
弟兄姊妹自己也會這樣想.
我有份工作我很忙.
我退休了.
我又要照顧孫子.
我已經做了一輩子.
你又叫我做就不要了.
我們都會想這些.
我們事實上有很多要估計.
我經常說.
如果教會要進入社區.
我不反對要做估計.
甚至要做研究.
上我的課要做研究.
不過這樣.
我們估計的時候.
你不要拿計算機出來計算.
拿計算機來計算就糟糕了.
因為你用一個商業的角度.
去營運教會.
這樣不對嗎?.
不對的.
因為計算機裡面.
加減乘除什麼制都有.
不過沒有上帝的恩典.
應許和祝福.
對嗎?弟兄姊妹.

$^{241}$我記得我的教會最初進入社區.
我的教會以前是一間中產教會.
在大角咀區.
很小的.
地方很大.
七千呎.
但會友只有五六十人.
連同小朋友.
但我們覺得很好.
很自貼.
我們每層都買.
你們都租.
我們每層都買.
覺得又可以.
又有奉獻.
不過當一進入社區的時候.
我的執事會就跟我說.
牧師說要.
那時候我只是傳道.
要計算一下.
我計算什麼數.
他們說要計算一下.
我不反對有些初步的評估.
但我真的很反對.
拿著計算機去計算.
因為那時候我很記得.
我跟我的執事會說.
是不是要計算KPI.
大家知不知道什麼是KPI.
做商界你就知道.
我更厲害.
因為我做了三十年商界.
我年年都要做KPI.
我是一間上市公司的行政總監.
要計算P and L.
我更厲害.
不過我們不能夠用這個模式去計算.
因為我們知道.
我們所做的.
是靠著神的恩典.

$^{281}$舉兩個小的例子給大家聽.
我記得在疫情.
第五波的時候.
很混亂.
我們已經做了八九年的社區服務.
當時那麼混亂的時候.
其實最大問題就是.
真的沒有什麼物資.
沒有測試劑.
沒有藥物.
當時因為社區很混亂.
教會有599G.
大家記得了.
就不能開門.
但我知道很多街坊都很困境.
當時我就跟我的同工說.
不如開一條24小時熱線.
如果油尖旺中了就打來.
我們送物資給他.
送什麼呢.
七天的藥物.
七天的食物.
我的同工說.
牧師你傻傻的.
他說你看看我們的倉有多少.
看完我們估計.
如果一天有100人打電話來.
我們可以送兩天.
我就跟他們說.
我就送兩天.
中了多少就多少.
接著我的同工說.
牧師那誰聽電話.
24小時.
我說我聽.
不過你們送.
因為我要留一條命.
我當時開玩笑.
因為我們有9個同工.
我們分開3組.

$^{321}$我們有3個地點.
每3個人為1組.
第一天接了多少電話呢.
300多個.
我們都不知道怎麼辦.
不過經篩選之後.
我們做了最急的200個.
就這樣送了.
很奇妙的.
第一天之後就沒有藥物派.
但我們繼續做.
我們都看Facebook post.
未必有藥物派.
測試劑還有一點.
但我們照樣買菜.
因為他不能上街.
我們就送7天的食物給他.
都很厲害.
有雞有魚有菜有肉.
我們就繼續做.
最初的同工說誰送.
我們很快.
半天就找到100個義工.
在區內很多弟兄姊妹.
很多未信主的人.
都來幫忙.
煮到第三天.
不知道熬了多久.
突然有個電話打來.
我知道你們送藥物給街坊.
我說是呀.
他說你需要多少.
我記得當時的對話.
我問他.
你有多少.
我心想你都滿足不到我.
其實頭幾天真的很好.
有些婆婆.
即是我們的街坊.
一代二代的便利貨.

$^{361}$原來他們平時存了很多.
我不知道你有沒有存.
如果有存.
你都差不多上一年多.
他們看完公立醫院.
每次都給便利貨.
每次都給感冒藥.
一兩天就這樣捱過.
因為到處都沒有便利貨.
到處黑通都沒有.
然後那個人就跟我說.
你要多少.
我就問你有多少.
猜了兩眼之後.
他說其實我們是藥廠.
你想要多少.
我就說我需要給他聽.
很感恩.
第二天他就送了.
又不是很多.
是一輛貨車這麼多.
後期他差不多兩三個星期.
又送一次給我們.
其實我們COVID完之後.
沒有之後.
即是沒有了之後.
我們整個倉都還是藥物.
我們要分發給弟兄姊妹.
我想告訴你.
很多時候我們沒辦法這樣計算.
因為上帝的恩典事實上很多.
剛才說過我教會很多教會.
在上一年我記得有一間教會.
他們接近倒閉.
他們很想做社區工作.
但他們面對一個很大問題.
第一沒有錢.
因為他們經過殖堂和移民潮之後.
教會只有50多人.
平均年紀是65歲以上.

$^{401}$沒有年輕人.
只有一兩個.
如果那兩個不在.
我懷疑拉到70歲.
最老的長者90多歲.
不是一位.
是好幾位.
當時他來跟我談的時候.
他說牧師我們都沒什麼希望.
他還想找郭文慈牧師.
說想花錢結業.
真的.
然後郭牧師來找我.
他說這間教會有辦法嗎.
我說我去看看.
我看到覺得很寶貴.
這麼多退休長者.
真是發達.
不是開玩笑.
不是吧.
是的.
長話短說.
到最後他上了我的課程.
開始做社區.
我給了他一些方法.
他在短短的九個月裡.
多了一百個年輕人.
現在他們教會的崇拜是百多人.
我最近找鄧睡文去跟他們做報道會.
坐滿了信主的幾十人.
他都說牧師我來不及.
但搞不定都是沒錢.
現在越來越多人越來越沒錢.
怎麼辦呢.
我就介紹了基金給他.
他真是好笑.
隔了不久我就問他.
怎樣能不能申請.
他說我打算申請五萬元.
我心想這麼少.

$^{441}$怎料基金批了六十萬給我.
五萬元也能申請.
幸好我熟悉基金.
因為人家說百多個年輕人.
四十多五十個新移民婦女.
為什麼這麼差.
我想告訴你基金不是基督徒.
但他很願意讓我們教會進行社區工作.
我看到上帝的恩典.
是我們人猜不到的.
但問題是.
我想告訴大家.
當我們沒有信心踏前一步的時候.
我又怎樣看到上帝的五餅二魚神蹟呢.
所以馬其頓宗教會這個超越界限.
讓我們提醒我們中賢有很多困難.
例如這班退休長者.
我上去看這班教會.
有一個九十多歲的長者.
原來他最厲害的是畫畫.
他有很多年輕人的粉絲.
他現在每個星期回來教三堂畫畫.
有一班婦女.
都是年輕人 六七十歲.
都是婆婆.
每天輪流回來.
煮愛心午餐給中學生吃.
中學生休息的時候.
下午來這裡吃飯.
給十元.
來這間教會吃飯.
結果晚晚又做了一批中學生.
他和隔壁教會合作.
所以現在年輕人多到不得了.
最近這一屆很多年輕人洗禮.
一間以為要倒閉的教會.
當他把最有限的放上去.
生命就不一樣.
第二 甘心樂意.
在第三到第四節.

$^{481}$愛心的屬靈原則.
除了超越界限.
就是要甘心樂意.
在第三到第四節說的是.
他們主動再三的請求.
目的是想在神的裡面有份.
那份恩情有份.
不知道你在事奉裡面.
是否甘心樂意.
真正的奉獻是要甘心樂意.
不是求.
所以我經常說.
如果有一天有弟兄姊妹.
來求牧師就好了.
這是求的.
我們通常是牧者來求.
你有沒有空?.
幫幫忙教什麼?.
幫幫忙在下面沖奶茶.
這是求的.
有探訪活動.
是求的.
幫幫忙.
通常是這樣.
求的時候.
年青人就說.
不要啦 我好忙.
我正在大學.
沒空.
你叫年長一點的.
年長一點的.
那些年長一點的.
還沒退休的.
又跟你說.
我好忙.
工作快要退休.
你不要搞我.
然後你再找年長一點的.
再年長一點的.
又跟你說.

$^{521}$我做了一世.
又是我.
通常是我們求的.
但根據聖經的原則.
你看看馬其頓宗教會.
主動再三請求.
准他們在供給聖徒的善事上有份.
今天我們有沒有求你的牧者.
在社區服務.
在福音工作上我有份呢?.
是你主動的.
還是求你.
你都要說.
我先祈禱.
看看有沒有感動.
這篇道我經常說.
如果第一次去教會.
我多數都說這篇道.
當然我的道很多.
但通常第一篇都說這篇道.
因為有很多例子.
兩年前我去一間中學講道.
那天剛好他們連同五,六年級的小朋友.
因為他們預備進入大堂.
坐在後面聽道.
我只說這篇道.
就是說求.
完了之後.
有一個十歲的小朋友.
五年級.
牽著爸爸衝來前面.
牧師牧師.
求你給我工作.
十歲.
我就說好啊.
因為他就在我們教會的隔壁街.
我們後面在講Band One的中學.
我就說好啊.
我接下來就要去暑期班.
我們教會的暑期班.

$^{561}$由第一天開到最後一天.
早上八點就開到晚上八點.
任何時間都可以來幫忙.
我記得牽著他的爸爸.
還很好笑.
他爸爸看不起他的樣子.
他就說你又說暑假是旅行.
接著小朋友就說我不去了.
爸爸就說聽著吧.
這樣就完了.
當然當時我的心也是想.
聽著吧.
這麼認真幹什麼.
隔了一個多月.
到我們暑期班.
有一天我在我們教會的走廊.
看到一個小朋友跑出來.
我告訴你.
我認人很弱的.
我一點都不認得.
他跑出來.
我問牧師.
我以為他上了暑期班.
小朋友就問你哪一班.
他說我幼稚園班.
我心想這麼大個幼稚園班.
我問他為什麼是幼稚園班.
他說你不認得我了.
所以大家千萬不要問我認不認得你.
真的不認得.
剛才有位弟兄要脫了口罩才認得.
牧師比較難認人.
你最重要是告訴我.
你哪個堂會可能我認得.
他問我是哪個學校.
上幾個月跟你說過.
他說在大龍.
我終於記得了.
我還很壞.
我說你真的認得我嗎.

$^{601}$他說是的.
我問他你做什麼.
他說很重要.
我現在在幼稚園班.
負責帶小朋友上廁所.
問他重要嗎.
他說重要的.
如果你做過暑期班.
或者主日學班.
你就知道了.
這個崗位很重要.
我很欣賞這個小朋友.
一個十歲的小朋友.
他聽道就會行道.
弟兄姊妹.
如果你信了主十年.
你聽過五百二十篇道.
請問你有沒有行道呢.
我們要甘心樂意在主的道裡.
讓我們跟他同工恩情有份.
第三.
為什麼他們可以這樣做.
是因為第五節.
他們奉獻給主.
他們之所以可以超越界限.
之所以甘心樂意.
不是因為別人叫他們.
是因為他們知道.
他們的生命早已奉獻給了神.
然後才用這個生命.
去幫助有需要的人.
所以我經常說.
如果你要進入社區.
去關懷弱勢社群.
照顧孤兒寡婦.
甚至去幫助整個社區.
一起認識主.
這是不能選擇做或不做的.
很多時候我們說要祈禱.
有沒有感動.

$^{641}$弟兄姊妹.
我們知道基督徒.
有兩件事一定要做.
第一.
叫做大使命.
你們要去使萬民作我的門徒.
第二.
叫做大誡命.
我們要盡心盡性愛我的神.
並且愛倫如己.
沒得選擇的.
今早你來崇拜.
你不會說.
今早起床.
祈禱看看有沒有感動.
有感動就回家.
沒有感動就不回家.
你不會的.
我不反對.
祈禱你做什麼群體.
你怎樣參與.
是對的.
但是教會和你的生命的使命.
大使命和大誡命.
這是必然的.
弟兄姊妹.
我們經常唱詩歌.
都說要將身體獻上.
當作活祭.
今天你的生命是不是這樣.
如果今天主在這裡問你.
我可以差遣誰的時候.
請問你會怎樣回答.
請差遣他.
他又懂打鼓.
剛才那位說多厲害.
他在社區打鼓.
又擅長唱詩.
唱詩的也很厲害.
我什麼都不懂.

$^{681}$所以請差遣他.
然後你還要加一句.
我為他祈禱.
基督徒.
弟兄姊妹.
我們是不是這樣呢.
我想大家去思想.
第二.
第六到第八節.
有些照片不讓大家看了.
我想大家幫忙.
讀第六到第八節.
哥林多後書八章.
第六到第八節.
預備起.
因此.
我們勸提多.
既然在你們中間.
開始這善事的事.
都把它辦成.
既然你們在信心.
口才.
知識.
萬分熱誠.
亦有在我們對你們的愛心上.
都情人一等.
難忘當在這次善的事上.
也要情人一等.
我遂處這罪.
並一切事明明理會.
以示直接對你們的熱誠.
代表你們愛心的真誠.
來到這裡.
第二點.
愛心的真誠.
哥林多教會大家都知道.
其實他們非常多恩慈.
什麼都可以.
不過很可惜.
哥林多教會.

$^{721}$你讀得很通透.
你就知道.
有一個諺語叫.
問題多多.
哥林多.
而他們最大的問題是什麼呢.
縱然他們有恩慈.
但是他們沒有愛.
與他們無益.
沒有愛.
保羅用馬其頓宗教會作榜樣.
鼓勵.
當時哥林多教會.
他說我不是命令你們.
我是用其他教會.
去鼓勵你們.
要有一個愛心的真誠.
或者真實.
不同的版本.
不同的譯法.
簡單來說.
就是要試試你們的愛心真假.
是不是只有個「講」字.
愛心是必須有實際的行動.
真正的愛心.
是要有行動力.
愛神愛人.
不是一句口號.
弟兄姊妹.
我們很多時候.
基督徒都說.
信忘愛.
甚至我們的愛.
什麼時候說得最多.
你知不知道.
唱詩的時候.
你看我們的詩歌.
手手都有愛.
當然是好的.
因為提醒我們.

$^{761}$但問題是.
如果我們的愛.
只留於詩歌裡.
這個愛就變得不實在.
在《約翰一書》三章十八節說.
小子啊.
我們相愛不只在言語和舌頭上.
總要在行為和誠實上.
保羅提醒慈惠之士.
是愛心的真誠.
因為真正考驗愛心的.
是我們願不願意走出去.
超越界限.
接著主動的再三的請求.
希望去做.
而這個做的.
去幫助別人與人同行.
發展愛心.
不是要讓人得著回報.
更不要說.
我認識他就有愛.
有時都要檢討一下.
我們信徒愛心都比較偏窄.
因為我們愛弟兄姊妹.
不是那麼難的.
因為你跟他熟.
你愛不愛旁邊那個.
不愛都不要說.
都愛的.
你弟兄姊妹.
熟的.
十幾年了.
我在這裡.
但我要愛一個陌生人.
容易嗎.
其實不容易.
但聖經的教導.
不只是叫我們.
閉上門.
Hallelujah.

$^{801}$愛自己友.
他要我們這個愛.
是彰顯出去.
因為我們願意去愛.
所以別人看到上帝的Being.
上帝就是愛.
要在我們作為信徒裡.
活現出來.
有一次.
很記得.
我去一間Mega Church 港道.
大家都聽到很興奮.
很流行.
大家完結後就去酒樓吃飯.
我剛好去那間酒樓吃飯.
不過我約了其他牧者聊天.
去到的時候.
在我前面有一個婆婆.
走在我前面.
不遠.
她去到酒樓沒有位置.
因為團契很早就爆了.
每張都有兩三個人.
有位置.
有一個婆婆在我前面.
街坊走進去.
逐張問有沒有位置.
兩夫妻都要.
沒有.
有.
我又看到.
張張桌的人的樣子.
就是剛才桌下的樣子.
當我走進去的時候.
距離不夠二十步.
當我走進去的時候.
張張桌.
牧師這裡有位置.
弟兄姊妹.
愛應該是怎樣.

$^{841}$我想大家去思想.
甚麼叫愛心的珍惜.
不是只有「港」字.
我們教會有沒有真正的.
走出社區去實行愛心.
第三.
第九節.
這個很重要.
是這一節的經文約時.
我請大家一起讀.
第九節 預備起.
你們知道.
我們主耶穌基督的恩典.
他本是富足.
卻為我們成了貧窮.
好聖.
這句是整段經文的約時.
信徒的愛心在哪裡來.
為甚麼我們有能力去愛.
愛一些我們完全不認識的人.
甚至去愛一些.
我們覺得一點都不可愛的人.
弟兄姊妹.
其實不只是基督徒會說愛心.
你覺得不容易.
我們基督徒很多人.
現在做飯堂去派飯.
我們教會也有飯堂.
做了十年.
逢星期一至五.
星期六日就吃飯盒.
堂食的.
做五晚餐.
現在有120個長者和露宿者.
每日都這樣吃飯.
有些已經回到天家.
又有新一批.
總共11年了.
但不只是基督徒會派飯.
道教派不派飯.

$^{881}$佛教更多.
沒有信仰.
排名都沒有信仰.
都派飯.
但問題是我們基督徒的愛心.
和沒有信仰或異教的愛心有甚麼不同.
我們要弄清楚.
分別是我們所認識的愛.
是從神而來.
不是我們自己裡面.
我想做善事.
我愛人.
不是.
是因為神先愛我們.
讓我們將神的愛植根在心裡.
我們就有上帝那種形象.
去學習怎樣愛人.
很不同的.
因為如果我們的愛.
是留於我們自己.
當我們遇到不可愛的事.
不可愛的人.
你就會覺得.
他很煩.
我為何要服侍你.
我就會這樣想.
如果你的愛是紮根於上帝裡.
我們的愛心是有根有機的.
聖經說.
基督因你們的信.
住在你們心裡.
叫你們的愛心有根有機.
所以真正的愛是從上主而來.
這段第九節經文.
正正告訴我們.
主耶穌的恩典是甚麼.
祂本來富足.
在天上擁有一切榮耀尊貴.
祂卻因為我們成為貧窮.
來到世界盜成肉身.

$^{921}$因著祂的貧窮.
使我們有救恩.
生命得著富足.
寄予是因為神愛我們.
這是恩典最大的特點.
馬其頓宗教會為何可以超越界限.
為何可以願意在基督上.
背後最大的動力.
就是來自這份救恩.
弟兄姊妹.
如果用我們的愛.
我們很多時候很薄弱.
例如我們剛才說的教會飯堂.
最初弟兄姊妹來幫忙派飯.
有些婆婆覺得你要服侍我.
當然有些婆婆最初的時候.
要飯心的飯.
不要側邊的飯.
菜太黃不要給我.
太油格油.
有時服侍的姐妹覺得.
這麼麻煩.
但當我告訴姐妹們.
不要緊.
你還用神愛你的愛心愛他.
他們問我甚麼是神的愛心.
你麻不麻煩.
她想想我真的挺麻煩.
上帝都愛你.
你是否罪人.
上帝都愛你.
你只要用愛心去愛他.
結果我們的飯堂.
現在有七成的長者.
是洗禮信了主.
我記得有一次看見一個見證.
那個婆婆真的這樣寫.
竟然有個地方不嫌她煩.
不嫌她身體一股味.
依然這樣服侍她.

$^{961}$這樣和她同行.
所以她看到基督.
她在信徒身上看到基督.
弟兄姊妹.
今天未信的人.
在我們身上看到甚麼.
所以我們要學習的.
不是憑我們那種.
那麼膚淺的愛人的愛.
我真的很膚淺.
和你友善一點.
就是愛你.
你真是刺眼刺鼻.
不想理你.
很膚淺.
但基督的愛.
只要你將它長居於心裡.
你就有能力去愛.
愛一些不喜歡的人.
愛一些不可愛的人.
甚至當你很累的時候.
你都有能力去愛.
舉一個例子給大家聽.
我其中一個同工.
一個傳道人.
是做露宿者的.
做街友的.
他每一個露宿者.
都跟很多年.
因為不是一兩年.
就可以突然間.
一聲變好.
而且他跟的很多露宿者.
都是吸毒的.
我記得有一次.
晚上他打電話來.
他哭得很厲害.
其實他很少哭.
因為他很大隻.
他整個變形合一.

$^{1001}$很大隻.
這麼大隻.
有些人哭.
我忍不住笑.
但又不行.
我問他.
你做什麼.
回到寫字樓.
他說我很生氣.
我說發生什麼事.
因為他有一個弟兄.
跟了十多年.
是街友.
由他吸毒.
在街上.
慢慢幫他陪他.
尾沙塘五道五道地戒.
戒了之後.
幫他找光房.
上了樓.
又幫他找工作.
陪他在地盤門口.
上班一個月.
因為怕他走.
每天陪他上班.
好了.
接著好一點.
又跌低.
好一點又跌低.
來來回回十多年.
最近這兩年.
真的變乖了很多.
結果他被放進了光房.
但光房也有規矩.
但有一次.
他可能在工作中.
遇到不愉快.
他就去喝酒.
喝了酒.
又回去找他的朋友.

$^{1041}$在橋底.
又再上了一次電.
回到光房的時候.
情緒很差.
因為吃了毒品.
又喝了酒.
打架.
我們叫做.
寫監.
叫傳道人馬上去.
見到他很失控.
沒辦法.
唯有做一件事.
報警.
報警的時候.
警察來到.
上無價值的媽葉的時候.
這個人竟然很兇狠地.
看著我的傳道人.
他說.
竟然是你親手報警.
我很恨你.
我出來一定報仇.
這個傳道人.
傷心到.
他不是怕他報仇.
而是他這樣和他同行十多年.
最後竟然在這個時刻.
他這樣和他說.
所以他回來就和我說.
我以後不會理他了.
很傷心.
他說我這麼多人做.
我終於做了.
以後不理他.
我沒什麼可以安慰他.
我和他一起祈禱.
祈禱之後.
我還記得我和他說.
明天你不要回家.

$^{1081}$你休息一下.
你都弄到這麼晚了.
不要上班了.
他說不行.
明天一早要去荔枝角.
去荔枝角做什麼.
接著我就這樣.
我說去荔枝角吃早餐.
他說.
你知道我.
說說而已.
其實到最後.
到現在這個弟兄.
當然他已經盡了主很多年.
現在又回到一個宿舍.
我們都繼續為他祈禱.
其實你想想.
跟一個人十幾年.
他最後當你是仇人.
其實你是理所當然放下.
但這個傳道人為什麼沒有這樣做.
是因為他的愛是從神而來.
最近他都做過一個案.
這個家友.
不理他八九年了.
但很奇妙.
最近這個家友的腳要鋸.
因為糖尿.
爛得很厲害.
他就夾進了醫院.
醫生說要鋸.
傳道人就力求.
找醫務社工各樣.
我們覺得如果非必要.
不要鋸.
搞了一大輪.
爭取了一大輪.
結果用了一隻新約.
現在不用鋸.
很感恩.

$^{1121}$做完手術之後.
這個盧小姐.
從來都不理我這個傳道人.
做完手術.
清醒.
她第一句就說.
可不可以幫我安排戒毒.
我想福音戒毒.
你可不可以安排我回教會.
怕不怕我很臭.
其實每一次可能要等很長時間.
但是.
如果我們的愛心不是從我們來的.
我們一定不會灰心.
我們一定有力量.
最後.
送這句給大家.
這是我自己的座右銘.
我覺得耶穌基督是我們一切愛心的根基.
他的恩典.
是信徒愛心的動力.
最後.
以天國的價值觀實踐真正的愛心.
這段經文比較長.
10到15節.
有君弘這個字.
讀一讀.
10到15節.
請大家讀.
我沒氣了.
預備起.
(念經).
我想大家留意這段.
由一開始到現在.
保羅整段經文從來沒有提過.
提哥林多教會要奉獻多少錢.
其實多少不是問題.
在這裡再一次強調.
有沒有願做的心.
願做的心.

$^{1161}$不是做多少的問題.
是你願不願意做.
做是最實際的行動.
其實這個背景.
哥林多在捐獻事上已經搞了一年.
超過一年.
但不知為何.
到現在還未搞成.
所以保羅再次.
寫這段經文.
勸勉他們盡心完成.
所以不是做多少.
而是你願不願意踏出.
但很奇怪.
去到第13第15節.
突然說君弘和合本叫做君平.
是甚麼.
我覺得君弘也更加好.
君弘.
甚至他用了嗎哪作比喻.
多收的沒有如少收的沒有缺.
是甚麼意思.
其實君弘或君平.
在舊約和新約.
一直是天國的價值觀.
其實所有世界的東西.
都屬於上主.
而今天神分給我們的時候.
我們已經學習怎樣做君弘.
即是說.
今天我們去幫助別人.
其實我們不是施予者.
他們也不是受助者.
只是今天我們付足.
將一部分君弘給他們.
將來到他們付足的時候.
可以君弘一部分給我們.
這是人人平等.
所以在第15節.
他用了嗎哪.

$^{1201}$如果大家記得當時.
收取嗎哪.
有能力的那些.
是否幫助那些完全行不通的.
收起來.
這就是比喻.
其實教會也應該是這樣.
今天我們教會有地方.
我們就應該分享給別人用.
好像你們樓下的奶茶地方.
非常好.
咖啡雖然剛才還沒趕得及喝.
待會再喝一杯.
非常好.
是一個開放的.
給街坊使用的地方.
我有一本書.
如果大家第一本書.
看過的話.
你會知道.
我有一個地產朋友說.
全香港最浪費地方的.
就是教會.
一至五都沒有人.
我以前教會也是.
我的教會七千呎.
一至五都沒有人.
我記得我最初做傳代人的時候.
真的是抑鬱症.
頭兩個星期.
電話又不響.
我們做商界.
電話又響.
電郵又響.
手機又響.
全日電話都不響.
一響就很興奮.
不過原來是打錯的.
電話不響.
然後沒有人進來.

$^{1241}$無論如何都沒有人進來.
除了外賣那個.
多恐怖.
七千呎.
星期六日就有人.
現在我的教會.
星期一至星期五.
星期六日當然多人.
每日有三四百人.
你隨時上來.
沒有人理會你.
你坐下來睡覺都可以.
沒有人理會你.
不會有人問你的.
很多我的學生試過.
他說牧師真的沒有人理會我.
真的沒有人理會你.
你拿著菜進來煮飯.
都沒有人理會你.
我們有個社區廚房.
我們現在甚至要.
多買一個地方.
和多租四個單位.
才夠用一至五.
弟兄姊妹.
我很想告訴你們.
什麼叫做均平.
均平就是我們願意.
分享無論空間.
金錢.
財幹.
時間.
我們都要服侍.
你說不是.
這裡說到我們不足的時候.
人家會補我們的富足.
會不會呢.
當然會.
我舉一個例子.
我們最初最主要服侍.

$^{1281}$就是新移民群體.
他們有很多不足.
所以我們會幫他們.
做英文班.
補習班.
適應課.
學前班等等.
很多人說.
你哪有這麼多人想服侍.
沒有的.
不過現在我的教會.
最多服侍的群體.
就是新移民群體.
譬如我們有一個長者計劃.
我們有百多個長者.
剛才說了.
長者又要覆診.
又要陪診.
又要家居清潔.
又要幫他們做身體檢查.
誰來做.
我們所有新移民婦女.
我們有40個新移民婦女.
加上年青人.
他們做什麼呢.
他們做了一個義工隊.
就是六個人為一隊.
照顧20個長者.
那怎麼照顧呢.
1,3,5就去我們的中心.
做運動.
看著他們.
20分鐘吸引.
吸完引.
夠引就去夾公仔.
我們真的有一部夾公仔機.
很開心的.
那些長者夾餅乾.
接著有剪頭髮服務.
有量血糖量血壓.

$^{1321}$每個星期都做.
有執藥.
因為我們發覺那些長者亂吃藥.
喜歡當差不多顏色就吃.
幫他們執藥.
陪診.
很多長者去覆診很淒涼.
還有幫那些獨居長者入院.
我們幾乎全道人.
不用探訪都可以.
因為那些弟兄姊妹.
一有探訪時間.
這些新移民婦女就輪流去探訪.
她在補我們的不足.
甚至我們的長者啟發課程.
是新移民幫忙開的.
我們的飯堂.
我們現在11個同工.
沒有人要進飯堂.
每天.
因為全部都是那些婦女做的.
我們一進廚房就被人罵.
進來做什麼.
倒杯水.
回去坐.
坐.
他們會說.
要煉水還是燉水.
幫我們倒.
是非常好的.
無論清潔.
我們教會沒有請清潔工.
7000呎都不用請清潔工.
全部.
信主的街坊做.
當日我們覺得.
我們很怕沒有人手去服侍.
今天反而.
我們更多的.
社區服務的人手.

$^{1361}$當我們不足.
他們補足.
均平.
因為每個人都有自身的資源.
可以分享在社區裡面.
譬如這個就是我們的義工隊.
我們甚至有社區大使.
逐家敲門去洗樓.
找一些不懂得找資源的人.
我們甚至會派他們去教會.
開啟發課程.
教他們去教會做婦女查競班.
這些是我們的服侍.
甚至我們有一個.
我很感恩.
有一個社企.
我們沒有想過會開一個社企.
我們有一個叫心更世作的社企.
因為我發覺那些婦女.
實在太厲害了.
煮東西吃.
不知道大家有沒有吃過酸菜魚.
知道有多少款嗎.
原來酸菜魚也有過十款.
不同的地區.
有不同的酸菜魚.
所以我弄到這樣.
我未做傳統人之前是八一磅.
現在是八四磅.
多大件事.
因為一天七餐.
八餐都是他們煮的.
煮同工的那份.
由於我發覺他們煮東西吃很厲害.
我們希望再幫助他們.
轉化他們的生命.
我們在教會弄了一個社企.
幫助一些SEN的家長.
一會兒他會介紹.
做了一個社區的服務.

$^{1401}$我們很感恩.
他在幫助我們的宣教基金.
這個社企三年都賺錢.
三年我問他們.
究竟用來做什麼的錢.
他們說放在教會做傳福音工作.
所以每年我們都有幾萬元進帳.
在宣教基金裡面.
今天我請了這位.
其中一個做心更世作的婦女.
跟大家做一個見證.
她叫李玉紅姊妹.
她說普通話的.
大家聽得懂嗎.
聽不懂我重要的就翻譯.
我的普通話是她那裡學的.
玉紅.
不如你介紹一下自己.
怎樣來宣導堂教會的.
大家好.
我當時去教會的時候.
是因為孩子申請奶粉.
然後孩子是S英小朋友.
當時申請奶粉的時候.
正在那裡填表的時候.
突然間孩子不見了.
其實我一直都是講普通話.
到了這邊就是.
孩子不見了就很緊張.
然後到處找.
她從禮堂.
因為我當時是在婦堂.
但是她是從禮堂那邊.
跑出來.
她說媽媽那邊有歌.
我們去聽歌好不好.
你緊張不緊張.
她指揮我要怎樣去做.
我說聽歌你不要到處亂跑.
一會兒被人罵.

$^{1441}$她第一次來教會是申請奶粉.
申請奶粉的時候.
突然間孩子不見了.
因為我們有兩邊門口.
原來她的孩子就跑去禮堂那裡.
她的孩子很好笑.
說那裡唱歌的那個很好聽.
這個一定要講的.
因為有人在唱福音樂曲.
她的孩子真的很特別.
特殊學習障礙.
她覺得樂曲很好聽.
而且因為唱的那個是我.
真的是我.
所以一定要講.
我教長姐唱樂曲.
她覺得好聽.
最後媽媽找到她.
罵她一頓.
因為她不知道教會是甚麼.
甚至她很害怕.
來到香港她不認識人.
她經常在街市被人罵.
因為她講普通話.
所以她找了一個不懂廣東話的童工.
去接待她.
叫May姐.
當時May姐接待我的時候.
我跟她講.
她對我講.
我們兩個誰也聽不懂.
誰在講甚麼.
講了半天以後.
她說你搞滅人走了.
我發脾氣了走了.
我就想拖著孩子跑.
當時接她進去一會兒.
分街出來.
我問她說我可以幫到你.
當時她一說完以後.

$^{1481}$因為她是用普通話講的.
我不知道為甚麼眼淚嘩啦一下就流下來.
我就覺得好親切好親切的.
那個感覺一出來以後.
我就問她這是哪裡.
然後她就.
跟我講這裡是教會.
我說孩子想聽歌.
當時就是在問.
她怎麼樣可以來.
她說這邊有孩子的補習班.
然後周末的時候又可以來.
我聽不明白.
補習班的話我家孩子年齡小.
他是不收的.
但是她就對我講了一句.
我幫你申請一下.
然後說完了之後.
就留了個電話.
第二天給我打電話的時候.
就說我幫你申請了.
可以帶孩子過來.
但是那一天是下著雨.
她又很溫柔的對我講.
外邊下雨不急的.
等到時候雨不下了你再來.
那種感覺很溫馨.
我就覺得那是家.
我一定要來.
芬姐其實是我們十年前.
接觸的一個SEN家長.
她小時候來的時候.
應該只有一歲.
後期她順了主.
在我們洗禮.
然後進了神學院.
她是我們一個牧養婦女的童工.
現在她一個人牧養200個婦女.
所以當天就由她.
接待她.

$^{1521}$她很感動.
開始帶小朋友來暑期班.
慢慢接觸教會.
一段時間就洗禮了.
孩子就在這裡長大.
你小時候來的時候是多大.
三歲多.
現在多大.
十幾了.
十一,十二歲了.
中一.
我記得我也有幫她的兒子補習.
好了.
你一直在這裡成長.
成為門訓的門徒.
我們教會開了心耕細作.
不如說說你們做心耕細作的過程或感受.
其實我們每天在家帶孩子.
煮飯帶孩子.
天天煮.
老公不會覺得你煮東西好吃.
孩子也不會覺得你煮東西好吃.
其實每天都說.
你在家做什麼.
那種感覺被打擊得什麼都不是.
但是我們在教會的一幫婦女.
在那裡我們煮東西.
大家都是那種.
你煮家鄉的東西.
我也煮家鄉的東西.
大家互相欣賞.
那種感覺完全不同.
就很開心.
自從開始做了心耕細作以後.
做出來的東西很多人欣賞.
我記得特別有印象的.
就是那個時候.
第一次是去那個波道神學院.
波道神學院那個時候.
我們堂到那邊.

$^{1561}$基本上是秒沒.
因為我們有很多時候在一起.
大家都是分享.
煮飯煮完了以後.
孩子不吃老公不吃.
外邊東西好吃那種感覺.
但是做完了以後.
出去了以後.
我們那一天出去.
秒沒的時候.
我就打電話回去給他們.
當時我就很開心.
那邊一點聲音都沒有.
接著就很難過.
他說就是說被人認可的那種感覺.
很開心.
就覺得終於有人認可了.
其實在這個過程當中的話.
我們做了這個以後.
現在就是說是我們回去煮飯.
老公和孩子都有說.
什麼時候煮一餐飯.
我們以前是煮飯.
他們就說出去吃.
但是現在說的就是.
什麼時候可以煮飯給我們吃.
那種感覺完全不同了.
他就覺得我被需要了.
我們也覺得你不需要我.
別人自然有需要的.
那種感覺我們的自信就有了.
所以大家不要嫌棄你媽媽煮的飯不好吃.
原來很大打擊的.
很感恩看到他們的成長.
我想問你們做新聞細作.
有沒有什麼目的.
有什麼目標.
有什麼意象.
其實我們做這個.
也有很多和我們一樣.

$^{1601}$就是新來的.
到了這邊各種不適應的.
然後我們是希望通過這個.
讓更多的人認識到教會.
可以找到家.
可以認識到.
我們回到這個大家庭.
多謝姚紅.
我幫你介紹一下.
今天我們也帶了他們一個產品.
叫做雪花素.
其實他們只有一個產品.
這個產品很特別.
第一就是他們自己重研發.
雪花素不是他們發明的.
但是他們改良了.
他們用的所有產品.
我記得最初.
因為我們是怎樣做.
我們開一個社企.
將他們個人的恩賜放在這裡.
找我們弟兄姊妹的恩賜配搭.
例如我們有弟兄是做食物研發的.
我們找他.
有人是做廚藝.
我們找他.
有人做酒樓.
我們找他入貨等等.
有人做CS的.
我現在做CS的姊妹.
原本兩年前是電腦都不會開的.
現在她是完全用電腦下單.
寄貨什麼的都是他們.
在做這個過程.
我自己很深刻的一件事.
就是他們很嚴謹.
他們覺得這個是代表基督去做.
第一他們所用的都是有機的產品.
我記得綠茶粉.
他們試了十幾次.

$^{1641}$因為他們不要加牛油的糖.
所以其實只有乾果的甜味.
為什麼綠茶粉要試十幾次呢.
因為他們覺得.
每一種都沒有抹茶味.
結果他們買了一種.
日本入口的綠茶粉.
然後被我罵.
我做了虧死你們.
他們說不要緊的.
只要有人欣賞.
我們要拿回那種尊重和尊嚴.
後期我終於明白他們.
所以在這個雪花show裡面.
他們每一次出去做攤位都叫我提醒.
他們做了試食版.
請大家不要因為他們是新移民.
不要因為他們是SCA的家長.
而光顧這盒東西.
千萬不要.
請你一會兒下去.
你們有奶茶嗎.
反正送奶茶的價位一流.
請你下去試食.
你覺得好吃.
我歡迎你買.
很多人買來送給朋友.
我們最好生意就是聖誕,復活節,新年.
因為新年買一盒餅都只需一百元.
很多人買來送禮.
我都鼓勵你送給未信的人.
或者拿回家吃.
糖尿病是可以吃的.
不過不要吃太多.
因為它都有乾果的甜味.
雖然沒有額外加糖.
每一件都是手製的.
但不是用機制的.
都是用心的.
裡面都有福音的訊息.

$^{1681}$鼓勵大家去試食.
我們不是拿了很多貨來.
買貨不是拿了很多.
反而是試食拿了很多.
所以鼓勵大家.
如果一會兒你吃得到.
跟他們說聲欣賞.
因為我下面還有一個姐妹幫忙賣的.
說聲欣賞.
另外賣完最後的廣告.
原來沒有了.
我帶了我的新書來.
舊的那本不能帶來.
因為已經斷版了.
只有這本新書.
《行在社區的福音》.
記載了甚麼呢.
第一記載了整個入社區的理念.
記載了上年跟我一起.
我coaching的27間教會.
其中13間的個案.
剛才說的「執泥執笠」的個案.
有在裡面的.
有三個院長所寫的文章.
最重要是裡面有一個research.
這個research給我很大的反省.
我們找了一個research center.
做了一個research.
這個research出來.
這27間教會.
參加了《行在社區的福音》.
即是進入社區之後.
他們的信徒的靈性.
是大大的增長.
這個我覺得正如.
聽道行道才是生命的實踐.
這本書我買了多少錢.
128元我也不知道.
因為這本書雖然由我出.
但我會存回教會的木外中心.

$^{1721}$做慈善用途.
希望大家藉著這本書.
祝福你們的教會.
祝福大家好嗎.
不好意思.
我知道今天overrun了一點.
願主祝福大家.
拜拜!.
\newpage



\section{使徒行傳 9:10-20}
\label{sec:QcQ_nHz9bjU}
\textbf{2025.03.02 《善待不可愛的人》 使徒行傳9\_10-20 (和修版) 講員 : 李富成牧師}
\newline
\newline
連結: \href{https://youtube.com/watch?v=QcQ_nHz9bjU}{\texttt{https://youtube.com/watch?v=QcQ\_nHz9bjU}} ~~~~ 語音日期: 2025-03-02
\newline
\newline
\hyperref[sec:TScbhjJAAD8]{\small{< < < PREV SERMON < < <}}
~
\hyperref[sec:index_chronic]{\small{[返順時目]}}
~
\hyperref[sec:index_scriptual]{\small{[返順卷目]}}
~
\hyperref[sec:oUBZwJOpFGQ]{\small{> > > NEXT SERMON > > >}}
\newline
\newline
使徒行傳 9:10-20
\newline
\begin{longtable}{cl}
\hline
\hline
章節 & 經文 (和合本修訂版)\\
\hline
9:10 & \begin{tabularx}{0.7\textwidth}{X} 那時，在大馬士革有一個門徒，名叫亞拿尼亞。主在異象中對他說：「亞拿尼亞！」他說：「主啊，我在這裡。」 \end{tabularx} \\ \\ \relax
9:11 & \begin{tabularx}{0.7\textwidth}{X} 主對他說：「起來！往那叫直街的路去，在猶大的家裡，去找一個大數人，名叫掃羅；他正在禱告， \end{tabularx} \\ \\ \relax
9:12 & \begin{tabularx}{0.7\textwidth}{X} 在異象中看見了一個人，名叫亞拿尼亞，進來為他按手，讓他能再看得見。」 \end{tabularx} \\ \\ \relax
9:13 & \begin{tabularx}{0.7\textwidth}{X} 亞拿尼亞回答：「主啊，我聽見許多人講到這個人，說他怎樣在耶路撒冷多多苦待你的聖徒， \end{tabularx} \\ \\ \relax
9:14 & \begin{tabularx}{0.7\textwidth}{X} 並且他在這裡有從祭司長得來的權柄，要捆綁一切求告你名的人。」 \end{tabularx} \\ \\ \relax
9:15 & \begin{tabularx}{0.7\textwidth}{X} 主對他說：「你只管去。他是我所揀選的器皿，要在外邦人、君王和以色列人面前宣揚我的名。 \end{tabularx} \\ \\ \relax
9:16 & \begin{tabularx}{0.7\textwidth}{X} 我也要指示他，為我的名必須受許多的苦難。」 \end{tabularx} \\ \\ \relax
9:17 & \begin{tabularx}{0.7\textwidth}{X} 亞拿尼亞就去了，進入那家，把手按在掃羅身上，說：「掃羅弟兄，在你來的路上向你顯現的主，就是耶穌，打發我來，叫你能再看得見，又被聖靈充滿。」 \end{tabularx} \\ \\ \relax
9:18 & \begin{tabularx}{0.7\textwidth}{X} 掃羅的眼睛上立刻好像有鱗一般的東西掉下來，他就能再看得見，於是他起來，受了洗， \end{tabularx} \\ \\ \relax
9:19 & \begin{tabularx}{0.7\textwidth}{X} 吃過飯體力就恢復了。 \end{tabularx} \\ \\ \relax
9:20 & \begin{tabularx}{0.7\textwidth}{X} 立刻在各會堂裡傳揚耶穌，說他是神的兒子。 \end{tabularx} \\ \\
[1ex]
\hline
\hline
\end{longtable}
$^{1}$各位靜姐妹平安.
我今次是第幾次來貴價會廣道.
即是來過第一次.
我第一次有介紹過我背景.
我未做這行之前我是做哪一行.
於是我們就玩一次.
訓導主任叻.
訓導主任之後跟住做校長.
在我做校長工作的時候.
做老師工作的時候.
很喜歡同同學玩一個小遊戲.
叫做老師畫.
我現在就同你玩.
怎樣玩呢.
就是如果我待會請你給一個指令.
前面有老師說你先好跟住做.
如果沒有老師說.
就這樣給你一個指令.
你不要跟住做.
如果你跟住做.
就不聽老師的話.
那你就輸了.
老師說舉起你的右手.
放下.
輸了一半了.
就是這麼簡單.
大家都不聽老師的話.
再來一次.
老師說舉起左手.
舉起右手.
站起來.
厲害了.
沒有人站起來了.
請放下.
這個小遊戲.
在我做校長工作的時候.
經常跟小朋友玩.
因為有很多寓意.
可以用這個小遊戲.
跟同學來分享.

$^{41}$可以轉做老師畫.
校長畫.
媽畫.
爸爸畫.
四大長老畫.
等等都可以.
弟兄姊妹我們記憶圖.
我們聽誰話.
對啦.
聽耶穌畫.
耶穌的說話.
有些很合聽的.
有些我們覺得.
真的聽多些都不要緊.
譬如你對人好.
人對你好.
耶穌說.
你賜我們平安.
這些很合聽.
但是大家有沒有想過.
耶穌有些說.
你聽起來.
你覺得有點不滿意.
你覺得做不到的.
你會覺得.
理想來的.
不現實的.
例如這一段.
耶穌在路加福音第六章.
32-35節.
有一段這樣的說話.
你讀過沒有.
我們重溫一下這段的聖經.
我們請弟兄讀弟兄那些.
講完即是我.
就讀紅色那些.
姐妹就讀紫色那些.
我們合作讀好這段.
很工整的耶穌的教導.
弟兄開始.

$^{81}$(你若只愛愛你不必).
(有是罪人 假愛罵 愛他悶的人).
(你若善待 那善待你悶的人).
(有是罪人 那善者入座).
(你若借給人 希望從他收回).
(有是罪人 那善者入座).
(你若借給人 希望從他收回).
(有是罪人 那善者入座).
(你若借給人 希望從他收回).
你們倒要愛仇敵.
要善待他們.
基督教講愛.
我想大家不會陌生.
愛到什麼級數呢?.
愛自己喜愛的人.
我現在最愛當然是我兩個孫子.
我女兒已經差了一點.
愛自己喜歡.
同聲同氣.
愛自己喜愛的人.
愛自己喜歡.
同聲同氣的人.
一齊喜歡唱歌的人.
一齊喜歡.
同一個興趣的人.
性格相近的人.
暗key的人.
血緣相親的人.
我都認識有句說話叫.
什麼有別話?.
親疏有別的嘛.
教會的弟兄姊妹.
都愛的.
不過耶穌在這裡.
正在講一個饒恕的問題.
她說你們的愛.
要去到一個能人所不能.
基本上不是人做到的東西.
叫做.
愛你的仇敵.

$^{121}$要善待他們.
你讀這段經文的時候.
你會覺得理想.
做不到的.
不過向著這個目標.
所以如果我做不到.
你千萬不要拿這段經文來量度我.
因為我是一個有血有肉的人.
怎麼能夠做到.
愛仇敵.
愛那些不可愛的人.
善待那些不可愛的人.
剛才主席帶領我們讀的這段經文.
有一個主角.
他叫做阿拿尼亞.
主政政就是要.
呼召阿拿尼亞.
去幫一個.
絕不可愛的數落.
主政政就是要.
呼召阿拿尼亞.
去幫一個絕不可愛的數落.
要去善待他.
究竟.
阿拿尼亞.
又怎樣去回應.
耶穌這個吩咐.
如果你有留意.
剛才讀這段經文.
是一段很公整的.
故事.
起成轉.
合是齊的.
開始的時候.
這位大馬士革.
的門徒.
阿拿尼亞.
有些聖經研究學者.
認為他應該是當時這個.
在大馬士革的.

$^{161}$一班猶太銳為主的.
耶穌基督的門徒.
的領袖.
為什麼呢.
都值得我們接受這個.
的推斷.
因為他和主的關係是非常密切的.
他可以.
在意象中和主對話.
主就.
叫他.
去幫一個叫做.
數落的人.
為他來去.
按守.
大家覺得.
耶穌這個命令這個任務.
有什麼難度呢.
有沒有難度.
沒有啊.
為什麼會有難度.
你是耶穌的門徒.
是教會的領袖.
耶穌說你去那裡找一個人幫他.
如果今天有一位新朋友.
來了來了幾次聚會.
是一位男士.
和阿富城差不多背景的.
那祖師就說.
有個新朋友來了.
和你差不多背景你可不可以代表.
教會去和他聊聊天.
好像他決了志.
你可不可以和他做初信栽培.
那我會不會托.
祖師手肘看來都不會.
雖然我都會覺得.
早半個小時回來.
或者遲半個小時走.
要和他做初信栽培.

$^{201}$我都要付一些代價.
不過沒什麼難度.
應該的.
阿拉利亞對於主.
這個的吩咐.
確實.
有一個這樣的回應.
我們一起讀引號裡面的字.
開始.
主啊,我聽見許多人.
講到這個人.
這樣在耶路撒冷.
多多輔代你的聖徒.
並且他在這裡.
有從祭司長得來的.
權兵.
要捆綁一切求告.
你命的人.
對阿拉利亞來說.
在當時.
他的認知.
蘇羅是哪一位.
蘇羅是殺星.
蘇羅.
已經臭名遠播.
因為他.
已經在耶路撒冷.
作為猶太教的狂熱分子.
他得到.
祭司長的仇意.
可以任意捉拿.
殘害基督徒.
這個消息.
我們在大馬士革的信徒.
大家都知道.
也聽到有一個情報.
這個蘇羅正是.
正在來大馬士革.
捉拿在大馬士革.
的猶太裔的基督徒.

$^{241}$因為祭司長的權兵.
即是這個拘捕令.
是對猶太人才有效.
然後捉到他們.
之後就戒他.
回到耶路撒冷的宗教法庭.
接受一個審訊.
你要信耶穌.
就虐待你.
你要信耶穌就殺你.
所以對當時的阿拉利亞來說.
蘇羅是殺星.
可不可愛.
不可愛.
耶穌就叫我去找他.
那不就是宋陽嗎.
想不想得明白這件事.
宋陽入虎口.
耶穌為什麼會叫我去宋陽入虎口.
所以.
阿拉利亞這個回覆.
對主的差派.
作出了一個很大的疑問.
他其實是質詢主.
你為什麼叫我去找這個人.
這個人已經是.
罪惡超章了.
可能我們讀聖經的人.
就會想.
會不會阿拉利亞收錯情報.
蘇羅不是那麼壞.
不是.
在路加寫.
《史圖行傳》到後期.
第26章.
當然我們大家都知道後來.
蘇羅變成了史圖保羅.
史圖保羅.
在面對他.
自己曾經.

$^{281}$同夥的陣營的人.
對他的追逼的時候.
他曾經面對.
群情洶湧的一班.
猶太的狂熱分子.
他曾經做過一段.
這樣的自白.
他說我從前是怎樣的.
我從前是竭力反對拿撒拉人耶穌的名字.
我在耶路撒冷都是這樣做過.
我不單止從祭司長那裡.
拿到權兵.
將很多的聖徒收在監裡.
而且他們被殺.
我也表示贊成.
簡單來說.
不是他下刀都好.
他有沒有份決定.
有.
就是鄒牧師要信耶穌.
殺了他好不好.
從這個角度去講.
當時的蘇羅.
自己去紙白的時候.
他未歸主之前.
他是滿手基督徒鮮血.
他在各會堂.
我屢次用刑強逼他們.
說褻族的話.
如果這個要信耶穌的.
就盡情的羞辱他.
虐待他.
我很憎恨他們.
很厭惡他們.
追逼他們直到外邦的成真.
這個外邦的成真.
應該就是指導在第九章講到的.
去大馬士革捉拿.
基督徒的這段故事.
大家留意到.

$^{321}$路加記載這一段.
張力就很明顯.
耶穌叫他的門徒去幫人.
順理成章.
但是這個.
阿拉利亞就說.
怎麼幫這個人.
這個人幫不下去.
我去到是自頭羅網的.
簡單來說.
就是死硬的.
怎麼去.
幫不過.
我們同不同情阿拉利亞.
為什麼阿拉利亞.
對幫蘇羅.
有很大的猶豫.
因為蘇羅絕不可愛.
不單止不可愛.
不是性格不同.
不是樣貌可憎.
是這個英文字.
這個是什麼英文字.
敵人.
什麼敵法.
不是普通的敵人.
還是死敵.
什麼叫死敵.
不是他死就是我死.
去到這樣的張力下.
這個人的家.
去幫這個.
要來捉我和我班弟兄姊妹的.
蘇羅.
好.
我們疑似創作.
我現在就是阿拉利亞.
我正是跌在這個兩難.
不明白的情況.
剛剛今天就大馬士革的.

$^{361}$主日崇拜了.
我就帶著這個疑問.
來到大家的當中.
我就想問一下鼎姊妹.
你們全部現在扮演.
是大馬士革的信徒.
我就問大家.
我有一個這樣的情況.
哪位鼎姊妹贊成我.
阿拉利亞去幫蘇羅的.
有沒有.
舉手.
沒有嗎.
有.
有人贊成我去的.
贊成我去的那些.
看來我都是和你是敵人.
你贊成我去.
不就是等我去死.
你可能會說.
不是吧,夫人.
因為主吩咐你去.
你講得很響.
主不吩咐你去.
你贊成我去.
你都知道.
在當時大家的認知裡面.
我去.
是自己去送死.
為甚麼你還贊成我去.
沒辦法.
你又沒有叫我.
有沒有人說.
反對我去的.
有沒有人反對我去的.
只有一個.
你都夠大膽.
我明明是主吩咐叫我去的.
你都叫我不要做.
大家明明跌在一個.

$^{401}$去又死.
你都會說下一句.
再來一次,去又死.
那怎樣.
死就死,有些人叫我去死就死.
去吧.
不過.
阿拉利亞這個疑問.
這個兩難.
你留意這段經文.
在轉那個位.
主就很明白.
所以主給了答案.
主對阿拉利亞說.
你只管去.
它是我所選的器皿.
要在外邦人和君王並以色列人面前.
宣揚我的名.
我亦都會指示他.
為我的名必須受.
很多的苦難.
大家留意你今日.
你和我今日都知道後來的素羅.
真是這樣的.
後來的素羅.
真是這樣的,記不記得我第一次.
講到12月最後的主人.
就是講沒有違背那從天上.
來的意象,他正是在阿基帕王.
面前受審.
正是一路傳福音.
一路被猶太的狂乏.
狂熱分子追逼.
素羅曾經.
後來的保羅曾經描述過.
他遇過的各種.
患難.
他開列過他受苦的清單.
我們大家都知道.
不過你設身處地.

$^{441}$如果你是當時的阿拉利亞.
你知不知道?不知道.
那你說不是?那耶穌就跟你說了.
他未來是這樣.
你的理性.
接不接受到?今日這個是.
滿手基督徒鮮血的人.
他日會為了耶穌基督.
被人迫害?.
你會不會理性接受到.
這個今日是極端反對耶穌的人.
他日會傳揚耶穌的名字?.
我講了個答案給你聽.
不等於你可以理解.
不等於你可以.
耶穌已經將未來的水晶球.
都講了給你聽.
但是如果我們.
設身處地我們是阿拉利亞.
我們當時的理性.
和情緒.
我們聽不明白.
耶穌講.
這個素羅的.
未來的日子.
是這樣.
不過.
有四個字一定聽得明白.
哪四個字?.
哪四個字.
一定聽得明白.
一起講一次.
這句明白.
素羅將來會為我的名字受害.
傳揚我的名字.
全部都不理解.
因為根本沒有可能發生.
在當時的阿拉利亞心目中.
不過耶穌說.
你只管去這四個字就聽得很明白.

$^{481}$就是說.
你明白也好.
你不明白也好.
都要去.
我們大家留意.
整段經文最後還有一個環節.
叫起承轉.
合個合個位呢.
是大團圓結局的.
阿拉利亞最終是有去幫.
這個不可愛的素羅.
那我們就問.
為什麼.
他肯去.
那你說是啦有人答到啦.
因為聽從主耶穌的話.
不過這個答案.
正確了九成九.
差一點點.
剛才我一開始講.
耶穌的話.
正確聽的我當然聽.
但有些不正確聽的.
那不正確聽的那些都做.
叫什麼啊.
不正確聽的那些做叫什麼啊.
正確聽的要做.
不正確聽的都要做那些叫信服.
什麼叫信服.
我很喜歡唱歌的.
如果曾穆斯說夫城.
每個禮拜你都做帶領詩歌的workship leader.
我說多謝多謝多謝.
非常樂意很正確聽.
最好full band.
每次都好像我開演唱會.
如果現在耶穌說.
夫城你放下你校長的工作.
去做傳道人.
這個不正確聽.

$^{521}$我喜歡做校長多一點.
什麼叫信服.
信服啊.
你對的也好.
你不對的也好.
你都以耶穌基督的.
旨意為優先的.
的抉擇.
這個才叫信服.
我小時候.
開始有電視.
當時.
現在我們很怕兒女.
兒孫他們受電子奶嘴的影響.
他們沉迷在這個.
虛擬世界我們很擔心.
我們小時候.
我們的父母即是六十年代.
當開始家裡有電視的時候.
我們的父母就很擔心我們.
只顧著看電視就怎樣.
不讀書嘛.
大家當時的電視大家有沒有留意.
我不知道你們有沒有試過.
有些電視是有百頁門的.
拉起來可以鎖住它.
媽媽解鎖你才可以看.
不解鎖不可以看.
當時我很喜歡看一套卡通片.
鐵甲萬能俠.
有一天我家.
對著電視.
看到入迷的時候.
我媽媽在廚房我住的彩虹村.
記住彩虹村是一個非常.
有人性化的公共屋村.
因為它有獨立的廚房.
我媽媽從獨立的廚房.
伸頭出來.
阿夫煮一煮飯.

$^{561}$沒有豉油下去士多買一支豉油上來.
我想問.
我想不想去.
為什麼不想.
因為在看電視.
我跟一般小朋友的反應一樣.
就叫做.
可以了可以了.
現在叫做氣我媽.
我媽媽也聰明的.
過了十分鐘.
都不見鐵閘開.
有聲的.
我們有鐵閘的.
現在就伸第二個頭出來.
不像這次的面容和聲調不同了.
是凶惡了.
憎憐了很多.
我知道如果我不起身.
出行.
我當時所受的後果是什麼.
當時住屋村.
每天都鬼哭神號.
打掣的聲音.
不是這邊打那邊打.
不過不要打得太誇張.
如果打得誇張對面的黃太太會說.
你看打一下好了.
現在就不是了.
現在打到都報警了.
我知道如果我不起行.
我會有什麼後果.
於是.
這次就這樣玩了.
因為我也很想追.
我用最快的九秒九的速度.
跑下去士多.
買一支豉油上來.
一堆人到廚房.
豉油.

$^{601}$然後又得到電視看電視.
我想問.
我想去嗎.
不想.
最後有沒有去.
那就順服.
當然了.
媽媽那麼辛苦在廚房為整個家庭.
預備一個晚餐.
你這個貪看電視的兒子.
竟然托手肘.
是不是值得.
打呢.
頂紙媒.
阿拉利亞在這個.
那麼難的抉擇裡面.
他最後選擇了.
接受.
耶穌這個看起來.
根本完全不合理的.
吩咐.
因為順服的緣故.
他真的去.
找這個素羅.
於是最後這個.
段落就是講到.
素羅得幫助.
通常我們讀.
《史上行傳》第九章.
我們會將焦點放在.
史圖保羅當時的素羅.
在大馬士革的途中.
主向他光照.
以至他.
跌倒在地.
看不到東西.
有人領了他入城.
三天他是跌在人生一個.
最大的黑洞裡面.
在這個時候.

$^{641}$在上帝永恆的慈愛裡面.
就引領.
阿拉利亞.
介入在這個保羅.
的黑暗.
人生的當中.
阿拉利亞去到那個人家.
跟他說.
素羅弟兄.
然後就告訴他.
誰叫我來找你的.
就是那個向你顯現的主.
然後.
讓他能夠看回東西.
又被聖靈充滿.
然後.
受了洗.
誰幫.
素羅洗洗.
雖然這裡沒有主語.
說是阿拉利亞為他洗洗.
不過從經文的上下文.
我們完全可以推斷.
一定是阿拉利亞為他洗洗.
我們也.
有理由相信阿拉利亞.
是一個人去找他的.
因為他都想著.
有機會多或少.
他怎會找一群弟兄姊妹去.
所以他幫助.
這位素羅先生.
這個殺星的.
阿拉利亞.
經歷了一個他理性根本想像不到的事情.
他見到.
這個本來凶神惡煞.
威風凜凜的素羅是一個.
軟弱無力.
看不到東西沒有吃飯.

$^{681}$好像一塊人肉那樣放在那裡.
在這個時候.
耶穌基督就使用他.
成為素羅的.
福音確認的.
一位牧者.
也可以說是他的.
初信哉配樂.
然後大家留意聖經說.
當素羅恢復.
體力受洗.
素羅就和大馬士革.
的門徒一起.
住了一些日子.
然後就開始在各會堂.
即是猶太人聚集的地方.
去傳揚耶穌.
說他是上帝的兒子.
即是反轉窗口.
大家留意我紅色的字.
素羅和大馬士革的門徒.
一起住了一些日子.
有什麼特別.
大馬士革的信徒的認知裡面.
和阿拉利亞一樣.
素羅是哪位.
是我們的人.
是殺星.
是我們要逃避的追緝我們的人.
我們正是怎樣防他.
為什麼會住在同一些日子.
因為他信得住.
我想問他信主受洗.
這位殺星信主受洗.
這個戲劇性的轉變.
有直播的嗎.
有社交媒體報道嗎.
誰知道.
只有誰知道.
阿拉利亞知道.

$^{721}$阿拉利亞親自經歷.
這個不可能的轉變.
這個已經從本來是敵對的死敵.
現在變成弟兄了.
很正常.
阿拉利亞就在某一個主人.
帶著素羅進入到信徒群裡當中.
我們扮演一次.
你們記住.
你們還是記住素羅是來捉你的.
跟著見到我出場.
我是阿拉利亞.
丁姐妹早安.
今天介紹一個新朋友給你們認識.
哪位.
請出來請出來.
這位就是大名鼎鼎的素羅先生.
你有什麼反應.
傻了嗎.
阿拉利亞.
你帶他來.
快點派人出去看看有沒有兵馬.
圍住我們.
不是.
他信了耶穌.
傻了嗎你沒看過無間道嗎.
可能他用虎肉計.
然後來到我們信徒群裡當中.
將我們一網成禽.
阿拉利亞就跟丁姐妹說.
不用怕不用怕.
我親自.
用我去保證.
我經歷.
這位本來是來捉拿我們的素羅.
今天.
他已經被主光照.
跟我們一樣.
成為主合用的器皿.
跟著你的反應就從本來.

$^{761}$拖了嘴.
慢慢因為阿拉利亞.
那個領袖的.
將他自己壓了上去.
成為那個的保證.
你開始從張開嘴巴.
轉回就叫做.
way and see.
看看怎樣.
可能看看怎樣.
慢慢完了聖會.
就有人上來.
跟阿蘇羅握手.
然後蘇羅.
真的好像不是他們想像中.
那個來捉拿他們.
滿手基督徒鮮血的.
猶太狂熱分子.
然後.
接納他.
在信徒群裡當中.
蘇羅.
這個初信者被.
接納在信徒群裡當中.
誰壓了自己上去.
阿拉利亞.
幫到底.
阿拉利亞.
現代蘇羅帶來什麼結果.
我想今天我們很容易回答.
在教會的生活裡面.
特別講到新約聖經.
除了講耶穌.
我想講得最多一定是士奴保羅.
保羅被稱為.
外邦人的使徒.
他有三次的宣教旅程.
在建部,史瓦西亞,希臘,意大利各地.
從這個角度講.
他是一個偉大的宣教士.

$^{801}$下午你們要.
按守那位弟兄.
都是去到這一大地方.
做宣教的工作.
我們今天很多教會.
在猜傳宣教的.
聖經基礎裡面.
都從士奴保羅的歷程裡面.
來去結取.
非常好的信仰的素養.
他不單止是宣教士.
保羅並建立了.
很多間教會.
而且是選立領袖去傳承.
從這個角度講.
他是一個偉大的牧師.
雖然他不是一個住堂很久的牧師.
他也是一位牧師.
你留意肥納比書一開始.
肥納比是.
士奴保羅搭載歐洲宣教第一站.
所建立的教會.
他一開始就說.
諸位監督.
諸位執事.
意思就是這個大概他建立的十年的教會.
經過十年.
這個教會已經有規模.
有他們自己的領袖.
去帶領.
而第三個.
更加重要的.
或者說我體會.
是對整個教育運動.
這二千多年來.
教會上非常重要的.
一個奠基者.
新約聖經有27卷.
有13卷是.
保羅書述.

$^{841}$從這個角度講.
他是一個偉大的神學家.
如果大家有涉獵.
你就算讀初信栽培也好.
你在崇拜聽道也好.
你有機會讀一些相對.
有系統的神學課程也好.
在整個教義的建立裡面.
保羅書述.
是其中一個非常重要的素材.
就以我們馬丁路德.
所謂更正教這五百多年.
馬丁路德是看.
保羅書述裡面的.
加拉太書和羅馬書.
來到再一次.
去提倡.
什麼稱義.
這個我們.
大家都擁有的.
這個教義.
恩信稱義.
不只是在保羅.
特別在加拉太書和羅馬書裡面.
當然二佛所書都有講.
一個很清楚的.
信仰的根基嗎.
阿拿利亞是建立了.
一個可以說是千稀人物.
一個偉大的宣教士.
一個偉大的牧師.
一個偉大的神學家.
不過阿拿利亞去幫他的時候.
知不知道自己有這些結果.
不知道.
他只知道四個字.
他不知道他信佛主.
的吩咐會為.
教會這二千年.
帶來有這麼大的影響.

$^{881}$如果說.
保羅是一個偉大的教會領袖.
他這個初信.
輩員都非常之光榮.
但是他去的時候.
是不知道的.
他只是信佛主.
的吩咐.
在我事奉的歷程裡面.
我曾經都接待一個.
非常之不可的年青人.
他叫Jason 就是這個肥仔.
我已經在弘恩事奉.
我在弘恩之前.
在另外一個.
機構.
YMCA當時成立了.
在油麻地的YMCA成立了一個.
小型的.
神學訓練的基地.
我們英文簡稱叫ICM.
我們當時的.
目標是什麼呢.
傳統神學院不會收的那些.
就來我們讀的.
就是說他的學歷又不夠.
教會都不太想推薦.
不過他們.
又想接受一些訓練.
去服侍主.
在教會界裡面有這樣的.
弟兄姊妹.
有一位牧師.
認識我.
為我介紹了這位同學來.
我當時所事奉的ICM.
這位牧師已經跟我說明.
這個同學.
會考.
沒有分.

$^{921}$這位同學來面試的時候.
大概是20歲左右.
當然面試.
都不會一個人的.
我們有兩個同工一起面試.
完了大家都想.
如果收他.
有機會讀不完.
有機會.
實習唾衰我們的名.
有機會帶壞.
同學.
從面試所收到的.
信息.
他不是很準備.
不是一塊很適合的素材.
來讀神學.
不過.
很奇怪他又說到自己.
很想事奉主的熱切.
又真的打動我們.
當時.
由於我做小學校長的時候.
都接觸很多看起來.
不起眼的學生.
大家聽過我分享.
所以帶著這個心腸.
我就說給他一個機會讀.
再加上我們又不夠學生.
他上學.
他上學最大的特色就是.
趴在這睡覺.
有時候我們自己老師.
教還好.
有時候有些客席的老師.
我們兩個字形容叫什麼.
失禮.
有些牧長請他來教.
有塊肥豬肉在你面前.
趴在這.

$^{961}$牧長完了都會說.
這個什麼同學.
是不是我很催眠他.
即是把我們推回頭來.
我們很尷尬.
慢慢的同學都說.
Jason又睡覺了.
我們就去了解一下.
再深入了解.
原來他是一個.
很悲慘的家庭出生.
父母很早離異.
很小的時候他也患病.
經常出入醫院.
自然小學的成績.
一落千丈.
再加上不知道是不是病的.
緣故影響他肥仔.
什麼.
大家就知道很有趣的.
在人際關係裡面肥仔.
要麼就欺負人家.
要麼就經常被欺負.
大家知道有個廣東話的族語.
什麼肥仔.
見稱肥仔就想撕他.
他小學已經被人欺凌.
但是由於.
派位制度他小學成績.
很高漲.
他自然也派到地區裡面.
一些所謂龍蛇混雜的中學.
他這樣的背景.
去到中學.
更加強他被人欺凌.
的經歷.
由於那家學校是教會辦.
介紹他來讀書的那位牧師.
又入學校教聖經.
這位黃牧師見到Jason.

$^{1001}$憐憫他.
好似爸爸.
這樣跪跪他.
雖然他學業成績.
沒辦法追得上.
但是竟然去到中四中五的時候.
他因為參加基督少年軍.
他開始對信仰.
對服侍主.
有一種很有趣的領袖.
他就跟黃牧師說.
我中學畢業無論考成怎樣.
我都想繼續去接受裝備.
去服侍主.
黃牧師看著他都.
不知道哪一間會接受你.
後來知道我們出現.
就介紹給我們.
為什麼他會上課躺著睡呢.
是因為他已經要.
自尋自食.
自食其力.
為什麼他會躺著睡呢.
因為他要做餐廳的侍應.
做到.
當時沒有疫情.
他要做到一點多鐘.
再搞定所有東西.
三四點.
你要到九點回來上課.
再加上他有一個身體上的需要.
他自然就會躺著睡.
我們了解了他的背景.
我們提議他.
Jason不如這樣.
你不要做.
夜晚的兼職行不行.
或者做短一點.
你其實想讀完.
這個課程去服侍主.

$^{1041}$你做兼職.
你做到你維生到就OK.
好不好.
他聽說.
他知道我們疼愛他.
於是他真的做短了.
而那個的收入.
只是緊緊救他生活.
當然他住沒問題.
因為他爸爸有間屋.
在荃灣他住沒問題.
但食是要自己搞定.
生活要自己搞定.
他聽我們說.
於是你看到這張相影得很好.
在ICM事奉的最後一年.
當時有位牧長.
我們請他教我們一課.
他就說我這麼巧在暑假.
在馬來西亞暫訊會神學院.
都是開這課.
不如合起來一起上.
你跟同學不如組織他們來.
馬來西亞好像遊學團.
我又覺得幾有人.
於是就問下同學有沒有興趣.
真的有五六個同學.
有興趣.
我們組織了一個神學院級的遊學團.
在暑假去馬來西亞暫訊會神學院.
跟內地去的馬來西亞華人.
一起上課程.
同時是我們的退休.
這是張耀畢業那一天.
讀完後你覺得他是不是很正常.
我第一次見他打一條胎.
跟我一起在馬浸訊拍了這張相.
我很喜歡.
他最後畢業.
很多時候有些教會邀請我講青少年報道會.

$^{1081}$大家知道我這種的樣子去講青少年報道會.
青少年都有點抗拒.
所以我通常.
我遇到Jason的時間.
我就會跟Jason講.
你可不可以在這個報道會.
我講到之前你分享見證.
因為他的見證.
就對著青少年報道會的見證.
因為他的見證對著青年人來說.
就比較容易入口.
而他的見證也比較過山車.
他這樣的背景.
仍然蒙主的恩典和使用.
這次我在馬鞍山一個教會講報道會.
青少年報道會之前.
他跟我在前面拍檔.
他講他生命的見證.
最近我亦都介紹了他接受《恩與之聲》的訪問.
他現在在做什麼呢.
他現在除了在一間教會裡面.
有部分時間的工作.
那裡是相對穩定一點的收入之外.
他參與了一個.
專做中學生科研工作的小型機構.
叫做小火箭.
有一次在他的Facebook.
我見到的是什麼呢.
他很注重帶領.
就是要讓禱告去轉發.
校園的屬靈氛圍.
他竟然能夠去動員那些中學生.
7點半回到學校神討會.
那如果那些學生要7點半回來神討會.
那他可不可以遲到.
不可以.
我們見到他找到上帝給他生命的照明.
不可能他會7點半起來的.
更不要說7點半去到別人的中學.
但是事實就是這樣發生.

$^{1121}$開始的時候我們讀這段經文.
我故意是斬斷了的.
我們讀到這裡.
你們賭要愛仇敵,要善待他們.
那下面是說什麼.
我們一起讀下去.
並要借給人,不只望常還.
你們的賞賜就很多了.
你們必作至高者的兒子.
因為他恩待那忘恩的和作惡的.
你們要仁慈,將你們的賦是仁慈的.
教會善待李富城.
帶來了一個果子.
這位姐姐是我.
72年10月踏足我母會的第一個主學老師.
她叫做少霞姐.
當我在博神畢業的時候.
這位我聖經的啟蒙老師.
我當然會邀請她出席.
她也非常興奮.
因為她做夢都猜不到.
在72年10月她曾經教過的一個蛙鬼.
竟然完成了一個神學訓練.
她現在還在我母會聚會.
其實她不是帶我很多.
她只是帶我十多年.
她仍然在教會中心的參與事奉.
她在72年教主學的時候.
她是不是看到這個人將來有水晶球.
說這個李富城會做校長.
這個李富城會做牧師.
知不知道.
不知道.
當時當時當時透過教會給她一個任務.
當時有個代播老師帶了一班蛙鬼去教會.
即時要分很多班.
要動員很多弟姐妹去分班教學.
剛好她就教中我.
她不知道在她手下這個小朋友.
在幾十年之後.

$^{1161}$有一個這麼大的戲劇性的轉變.
就是這種我曾經被人善待.
被人接納.
被人幫助的經歷.
以致無論在我以前的教學工作.
我現在遇到的神學生.
就算我現在在弘恩.
都有同學是永遠deadline都不交功課的同學.
我都用無比的慈愛去逼迫她.
聽清楚.
用無比的慈愛逼迫她.
但是我不會採取一個你死就死.
我知道她有她的困難.
特別有很多是有在職的工作.
她又要讀書那些.
我有特別的恩典.
因為我都捱過兼讀的催涼.
在我做校長工作的時候.
我發現哪些老師最接納不到那些.
叫弱勢的學生呢.
就是那些老師.
他自己成長的時候是含金鎖時代的那些.
他就發現.
交功課是應該的嘛.
媽媽簽齊的手冊是應該的.
他很多應該的.
他不明白這個小朋友的不應該.
有他自身的困難.
老師需要屈膝.
放下身段.
去輔助他的不應該.
當然.
不是要認同他的不應該.
是要去輔助他.
然後幫他的應該慢慢走上來.
媽媽.
丁姐妹.
我們是因為我們都曾經被人善待.
我們才懂得用這一份.
的愛去善待其他人.

$^{1201}$耶穌說你們要仁慈.
好像你們的父是仁慈一樣.
有時候教會久了.
我們會有一個感覺.
我們都是一群好人.
但我們忘記了我們蒙恩的罪人.
所以我們慢慢會構成一種.
抗外的文化.
這個回來教會穿得不整齊.
那個小朋友回來就串串共.
你猜不到那個回來串串共的小朋友.
將來會是你的堂主任.
丁姐妹.
在家庭.
在校園.
在職場.
在社會.
在族群.
甚至在教會裡面.
有沒有人你覺得你是不能夠善待的呢.
你今天來崇拜之前.
你還想著那個人.
那個明天上班又會見到他.
最好不要見到他.
如果我今天這篇道給你有激勵的.
我鼓勵你.
我不會認為你明天會那麼戲劇性.
明天你見到他會抱著他.
最低限度.
他在這邊進.
你不要在這邊走.
他在這邊進你迎面過去.
帶著主角愛去跟他說早安.
他不理你嗎.
明天再喊.
喊了五天之後.
我相信下星期一他怎樣都.
點一點頭.
你的主動.
你的善待.

$^{1241}$會帶來一個不一樣.
的變化.
如果有機會.
我們玩一個這樣的遊戲.
當我這裡前面有一個人.
叫做曾銳貴.
曾銳貴.
你估你很可愛.
不過因為耶穌愛你.
所以我都善待你.
那麼曾銳貴就會跟我說.
李富城.
你估你好可愛.
不過耶穌愛你.
所以我都善待你.
只有一個人.
我都善待你.
只有那些認為自己好可愛的人.
才不懂得去善待其他人.
如果你發現你自己.
都不可愛.
你接受別人對你的恩典.
你就用這份愛.
去回應.
去愛其他的人.
說到最後.
能夠善待不可愛的人.
是因為信服主的吩咐.
按照人的常態.
自己的兒孫.
當然是愛.
同聲同氣的當然是愛.
智趣相投的當然是愛.
合理的當然是愛.
善待不可愛的人.
只有信服主的吩咐.
才會創奇.
做出一些.
別人會覺得.
為什麼你們這群基屋徒.

$^{1281}$真的有點不同.
如果我們願意這樣做的話.
我相信貴教會.
不單止在內部.
會有更加的友愛.
而這份愛是會.
滿到流出去.
寫出去.
讓我們帶來的親朋戚友.
社區的人士.
和我們一樣去經歷.
這份在耶穌裡面的恩典.
所以最後跟大家唱這首短詩.
叫在耶穌裡.
我們會唱三次.
第一次唱好朋友.
第二次唱好弟兄.
自然就是弟兄唱.
用普通話唱的.
大家一齊來參與.
在耶穌裡面.
在耶穌裡.
我們是好朋友.
在耶穌裡.
我們是好朋友.
在耶穌裡.
我們是好朋友.
從今時一直到永久.
在耶穌裡.
我們是好朋友弟兄.
我們是好弟兄.
在耶穌裡.
我們是好弟兄.
在耶穌裡.
我們是好弟兄.
從今時一直到永久.
在耶穌裡.
我們是好弟兄姊妹.
在耶穌裡.
我們是好朋友.

$^{1321}$在耶穌裡.
好姊妹.
在耶穌裡.
好姊妹.
從今時一直到永久.
在耶穌裡.
好姊妹.
我們再唱最後一句.
不過下一句會唱高八度.
你不明白我說什麼.
你跟著我就行了.
在耶穌裡.
我們是好朋友.
好.
謝謝祈禱.
我們同心一起禱告.
當你主.
你透過昔日阿拉利亞的榜樣.
還有一次激勵我們.
原來愛去到仇敵.
仍然是切實可行.
不過我們要學習對你.
旨意馴服的操練.
又讓我們更多去體恤別人的.
我們認為不足之處.
我們自己都有很多軟弱.
需要人善待的地方.
就讓我們彼此的包容.
彼此的饒恕.
彼此以恩慈相待.
去建立一個滿有愛的教會.
又讓我們這份愛能夠揚日流露.
去到我們的親友.
我們的社區求主幫助.
祈禱奉耶穌基督的聖名.
\newpage



\section{使徒行傳 8:4-8}
\label{sec:oUBZwJOpFGQ}
\textbf{2025.03.09 《聖靈引領的傳道》 使徒行傳8\_4-8 (和修版) 講員 : 曾裔貴牧師}
\newline
\newline
連結: \href{https://youtube.com/watch?v=oUBZwJOpFGQ}{\texttt{https://youtube.com/watch?v=oUBZwJOpFGQ}} ~~~~ 語音日期: 2025-03-09
\newline
\newline
\hyperref[sec:QcQ_nHz9bjU]{\small{< < < PREV SERMON < < <}}
~
\hyperref[sec:index_chronic]{\small{[返順時目]}}
~
\hyperref[sec:index_scriptual]{\small{[返順卷目]}}
~
\hyperref[sec:td1EIi6hWco]{\small{> > > NEXT SERMON > > >}}
\newline
\newline
使徒行傳 8:4-8
\newline
\begin{longtable}{cl}
\hline
\hline
章節 & 經文 (和合本修訂版)\\
\hline
8:4 & \begin{tabularx}{0.7\textwidth}{X} 那些分散的人往各地去傳福音的道。 \end{tabularx} \\ \\ \relax
8:5 & \begin{tabularx}{0.7\textwidth}{X} 腓利下撒瑪利亞城去，向當地人宣講基督。 \end{tabularx} \\ \\ \relax
8:6 & \begin{tabularx}{0.7\textwidth}{X} 眾人都聚精會神，同心合意地聽腓利所說的話，一邊聽他的話，一邊看他所行的神蹟。 \end{tabularx} \\ \\ \relax
8:7 & \begin{tabularx}{0.7\textwidth}{X} 因為有許多人被污靈附著，那些污靈大聲呼叫，從他們身上出來；還有許多癱瘓的、瘸腿的都得了醫治。 \end{tabularx} \\ \\ \relax
8:8 & \begin{tabularx}{0.7\textwidth}{X} 那城裡，有極大的喜樂。 \end{tabularx} \\ \\
[1ex]
\hline
\hline
\end{longtable}
$^{1}$各位姐妹平安.
我們一起有祈禱.
感恩天父給我們今早一同聚集.
敬拜你.
多謝你每一次都閱立.
每一個地方奉你名的聚會和敬拜.
我們能夠成為神的兒女.
是一件何等偉大的事.
我們能夠信到這位救贖主.
是你的啟示.
是你的恩典.
是你的揀選.
其實更加是一份愛的恩約.
恩典的美好的約.
我們很盼望.
我們珍惜.
我們把握.
我們更加要裝備自己.
向不同的人去傳福音.
今早的經文都真的告訴我們.
傳道的重要性.
我們一起仰望.
等候在你的面前.
求你藉著這段聖經.
讓我們能夠看到.
神的作為何等奇妙.
我們這樣仰望禱告.
奉耶穌基督的名.
阿們.
《使徒行傳》是一本歷史的書卷.
勞艱醫生是保羅的跟隨者.
也是同心的傳道者.
很後期信主.
根據記載教會的歷史告訴我們.
《使徒行傳》的馬其頓異象.
保羅祈禱的時候.
有一個人.
馬其頓人.
很特別.
他連國籍都能看到.

$^{41}$向他呼求.
你可不可以到我們東歐的地方.
跟我們講福音.
就是這位勞艱醫生.
一直由十六章開始.
到二十八章.
都是勞艱醫生.
用我們來記載.
一起報道宣教的旅程.
勞艱醫生.
記載了他目擊和知道的歷史的時候.
有幾段值得我們一起去想想.
第八章.
很重要的一段經文.
就是《肥利》.
《肥利》是哪一個呢?.
就是當時耶路撒冷教會.
七位的執事.
其中一位猶太人.
被選立出來.
要處理當時的不公.
猶太人和希臘人.
因為很多人聚會.
有不公平的分配.
可能是言語上.
可能是大家有一些矛盾.
所以導致教會內部.
已經出現了這樣的紛爭.
但是.
選立這七位執事之後.
立即那件事情.
就得到妥善的解決.
原因.
作者有說.
因為選立出來.
其中一個很重要的.
就是聖靈.
所以大家可能要留意一下.
在第六到第八章.
關乎施提反和肥利這兩位執事.

$^{81}$一而再提到他們有聖靈.
聖靈在他們的生命裡.
在他們的事奉裡.
主宰的作用.
所以很重要.
跟大家讀讀第六章.
這一段經文.
第三節.
使徒們認為要找人出來.
處理寡婦的分配.
每日的飯食問題.
第三節有這樣的描述.
要選出來的人的資格.
第三節.
所以弟兄們.
即是全教會耶路撒冷的弟兄姊妹.
當從你們中間.
選出七個有好的明星.
被聖靈充滿.
智慧充足的人.
我們會派遣他們去管理這件事.
第五節.
大眾都喜悅.
即是全教會都認為.
使徒們這樣列舉.
未選之前的條件.
是非常平衡.
非常客觀.
真的可以讓這件事.
得到一個妥善的處理.
有好的明星.
有聖靈充滿.
聖靈充滿是怎樣看的呢?.
聖經沒有記載.
聖靈充滿是怎樣看的.
聖靈在裡面.
心靈裡面.
好的明星.
我們日子內就會知道.
這個人真的有好的明星.

$^{121}$這個人的明星一般般.
我們會知道.
但是.
有智慧管理.
也是需要在工作的時候.
策略.
等等的分配.
是能夠很有條理.
做得好.
是很自然的.
但是聖靈充滿.
是怎樣的一回事呢?.
不過第六到第八章.
一而再出現這兩個人.
是有聖靈充滿的.
所以我們真的要看一看.
這段經文.
為什麼路加醫生.
他會這麼著重.
他們有這樣的表現.
和有這樣的內在特質呢?.
第五節.
這是第六章的第五節.
大眾都喜悅.
使徒們的這番說話.
就開始選人了.
選了斯提反.
乃是大有信心.
本來是沒有要求有信心的.
不過再加上.
他真的很有信心.
有信心真的很重要.
對一些事情的看法.
眼光.
眼界.
如果本身有信心.
是很不一樣的.
這裡說大有信心.
聖靈充滿的人.
然後再選其他六個.

$^{161}$又選了肥利.
第二個就提他了.
白羅哥羅.
尼加羅.
提門.
巴米娜.
和進猶太教的安提阿人.
尼戈拉.
因為有些是希臘的婦女.
所以有猶太人.
但是七個裡面.
都有一個進猶太教的安提阿人.
就是土耳其那個地方.
你可以想像到.
都很有代表性.
不只是清一色七個.
都是猶太人.
進猶太教信主的.
都有聖靈充滿.
不是猶太人專利.
所以路加醫生很仔細記載.
他真的從歷史的切入點.
去記載這些事情.
路加的《士徒行傳》是很好看的.
因為他有很多個人的.
做一些注釋.
就讓我們看到.
接著第七節.
七個人被揀選.
七個人有這幾樣的恩賜.
神給的生命的成熟.
很好的屬靈基督徒領袖.
就帶來教會有什麼後果呢?.
你想想.
第一節已經說有法願言了.
第七節.
經過了這個行政.
不是干預.
是行政的管理.
妥善的管理之後.

$^{201}$第七節.
路加給我們一個很振奮的消息.
原來教會是可以改變的.
願言,不公.
還有很多的誤會.
只要有聖靈.
只要有管理者.
有聖靈.
教會就可以改變.
大家一起讀第七節.
大家不能一起讀.
因為你們的無聖經在手上.
有神的道興旺起來.
在耶路撒冷門徒數目加增的甚多.
也有許多祭司.
信從了這個道.
你可以想像到路加醫生.
記載這件事的時候.
是多麼的振奮.
當然他不是在場.
因為他還沒信耶穌.
他還在很遠的地方居住.
在馬其頓那一帶的地方.
當他知道這些事情.
他再記錄的時候.
他有一個很肯定的注釋.
教會經過了神的工作.
有人代表神.
在教會裡面工作.
就真的帶來生命改變.
不單止裡面似乎一定不會再有願言.
神的道興旺.
歡欣,平安,喜樂.
一定是隨之而來.
這不在話下.
教會成長.
一定是每個人的內心有喜樂.
有平安,有合一.
充滿了一份對神的感恩.
對其他還沒信耶穌的人的靈魂的熱愛.

$^{241}$所以耶路撒冷的門徒.
數目加增的甚多.
還有許多的祭司.
都順從了這些道.
《小道行傳》每一次的事件.
都會有一些注釋.
路加寫作的作風就是這樣.
因為他真的以一個歷史的學者.
來記載這些事情.
記載這些事情有什麼發生呢?.
很快就出現了一個一個的爭端.
這次就不是較內的小小的誤會.
而是較外有非常大的基於對道的誤解.
甚至乎是猶太人自己的偏見.
首當其衝對施提反有一個很大的攻擊.
第八節作者很快將鏡頭轉一轉.
記載了一個獨特的事件.
就是耶路撒冷的教會.
為什麼在這麼興旺的情況下要分散呢?.
作者拿捏到一個原因.
就是施提反受到嚴重的迫害.
教會就需要分散在不同的地方.
大家留意第六章的第七節.
尤其是第十節就更加要留意.
施提反滿德因為能力.
在民間行大其事和神蹟.
大家要留意原來施提反可以行神蹟.
待會第八章肥利又行神蹟.
值得我們再想一想.
今天我們傳道和道本身帶來的生命影響.
我們一路讀下去就會看到端倪和作者想要帶進的真理.
留意第九節.
當時有清利伯地拿會堂的幾個人.
猶太人有很多會堂.
其中一個似乎很顯著的會堂.
其中的人和古利奈,亞歷山泰,基利加,亞西雅各處會堂的幾個人.
起來和施提反辯論.
開始要挑戰這些執事在教會裡面.
為什麼會宣講耶穌.
為什麼會對會堂慢慢.

$^{281}$有些人因為很多祭司都信得住.
那些祭司就去聚會.
可能會堂的人數開始遞減.
所以出現了一個他們很想了解.
但挑戰的成份居多.
再一次留意第十節.
作者路加再一次告訴我們.
施提反為人和內心.
這次不用處理飯食.
因為已經解決了.
教會內部已經非常平安.
大家彼此同心去傳福音.
第十節.
施提反是以智慧和聖靈說話.
眾人敵黨不住.
只能夠用旁門左道.
買出人來說.
我們聽見他說.
旁讀摩西和神的話.
然後煽動會堂的人.
民間的百姓,文士,長老.
來捉拿他們帶去公會.
當時還有這些公會的人.
因為羅馬的皇帝.
採取了宗教懷柔政策.
所以容許猶太教.
你們有自己的會堂.
關於你們律法的審理權.
所以耶穌被解到.
祭司長彼拉多和戈埃亞法.
所以是兩邊.
當然猶太教不可以有殺人的權柄.
只可以治他宗教的罪.
所以耶穌最後被釘十字架.
被彼拉多釘十字架.
所以這個時候.
民間的宗教法庭.
即是會公會.
還有這個權.
相當大的罪名.

$^{321}$就是旁讀摩西和神的話.
再加上第十三節.
設下了假的見證.
說這個人斯提反.
他說話不斷糟踐聖所和律法.
即是踩低聖所.
藐視律法.
他又不是一個祭司.
他又不是一個拉比.
他又不是一個教法師.
他只不過是教會裡信主的弟兄.
當然沒有說他之前做過什麼工作.
但肯定在教會裡是一個執事.
但是居然那個罪名說他侮辱踩低聖所和神的律法.
摩西五經.
這些罪名對猶太人來說.
是相當嚴重.
等於如果你是原居民.
你連夜去祠堂.
把燭泡燒毀.
把要拜祭的東西弄髒.
這樣就很嚴重.
真的非常嚴重.
你可以想像到.
這個罪名絕對不輕.
所以這個時候.
聖經又再一次很有趣地.
路加醫生捕捉那些場景.
他說私提犯被人捉去公會.
這樣的罪名.
他的面目是很虛怯.
是很害怕.
還是很生氣.
還是怎麼樣呢.
留意第十五節.
六章十五節.
在公會裡坐著的人.
都定睛看這個私提犯.
看到他的面貌.
好像天使的面貌.

$^{361}$有沒有人見過天使?.
我的舊教會有位姊妹.
我不知道真假.
她說她在聚會的時候.
見到天使.
她很正常.
她是做護士的.
她說她真的見到天使.
我也想啊.
我不會這樣祈禱的.
如果神要我見到就見吧.
但是當時是眾人見到天使.
天使的面貌.
天使的面貌是怎樣呢?.
磁場?.
經常笑?.
還是怎麼樣呢?.
但是見到.
相信是我們難以形容的.
出奇的平靜.
內心充滿了一份對神.
對自己的信心.
因為這是一個很嚴重的指控.
被人這樣誣告.
但仍然保持自己那份.
對神,對人的誠信.
那份穩定的平安.
不容易的弟兄姊妹.
這個罪名是相當嚴重的.
所以第七章.
大祭司就容許斯堤反.
來到他自己.
不會請律師.
自己去辯護一下.
大祭司就說.
這些事情.
是不是果然有啊?.
然後呢.
斯堤反就說了一篇很長的道.
來將由亞伯拉罕被神.

$^{401}$去拯救.
帶進迦南地.
一直說到大衛建造聖殿.
我們留意第七章到第四十六節.
這一段.
大衛在神面前蒙恩.
祈求為雅各的神預備居所.
卻是所羅門為神造聖殿羽.
其實至高者.
並不著人手所做的.
就如先知所言主說.
天是我的座位.
地是我的腳凳.
你們要為我造何等的殿羽.
那裡是我安息的地方.
這一切.
不都是我手所做的嗎?.
引述的是.
二賽亞書六十六章的經文.
斯堤反.
不過第五十一節的結論.
就出事了.
一直說.
好評鋪直敘說.
我們的民族.
我們的祖先.
更重要的焦點是聖所.
是如何來的呢?.
是如何確立的呢?.
是大衛的兒子所羅門建造.
來讓神有居所.
不過留意第五十一節.
這個結論.
由沖兒發出的.
對當時的群眾.
對當時的祭司.
利未人.
這些拉比.
發出一個控訴.
第五十一節.

$^{441}$你們這硬著脖子看.
心與耳未受割禮的人.
嘗試抗拒聖靈.
你們的祖宗怎樣.
你們也怎樣.
哪一個才知.
不是你們祖宗逼迫呢?.
他們也把預先傳說.
那二者要來的人殺了.
如今你們又把那二者.
即是耶穌.
賣了.
殺了.
你們受了天使所傳的律法.
即是摩世悟經的十誡.
竟然不遵守.
作出一個這樣的.
在聖靈的感動下的指控.
當時的人.
你可以想像.
本來已經有敵意.
本來已經對你有針對性的.
旁獨.
借題發揮.
想攻擊斯提反.
但斯提反還在說.
這麼嚴重的對他們的指控.
當然一定是事實.
當然是這群人.
他們不認識救主.
十誡的救恩.
當然不知道.
耶穌就是神的兒子.
而把祂賣給羅馬.
而把祂殺害.
把這位二者耶穌殺了.
他們其實一直從來都沒有遵守律法.
你們說受了天使所傳的律法.
其實從來都沒有遵守.
相信大家都不用想太多.

$^{481}$就知道後果的下場.
眾人聽見者說.
極其的惱怒.
向斯提反咬牙切齒.
但斯提反.
再一次提.
只是這兩章已經不住提一件事.
斯提反一直被聖靈充滿.
第55節.
斯提反被聖靈充滿.
定睛望天.
看見神的榮耀.
又看見耶穌站在神的右邊.
全本新約.
只有這地方.
說到主耶穌站在神的右邊.
所以很多教經家都認為.
在這個時候的這個姿勢.
其實是站在.
準備迎接.
快要死的殉道的斯提反.
因為斯提反說完道之後.
那些人咬牙切齒.
多謝你.
牧師說得很好.
下星期再來.
幾千塊石頭扔死斯提反.
所以主耶穌站在迎接.
殉道.
快要殉道的斯提反.
不過我們要留意.
今天的主題.
聖靈的很重要主宰.
聖靈的介入.
在這些事情上.
所以.
盧家醫生作為一個歷史學家.
如果按當時記載歷史的希臘手續.
是會去神跡化的.
不會記載神跡.

$^{521}$不會將神跡放在記錄裡面.
但是盧家的《少道行傳》.
一個本來是醫生.
應該是很科學的頭腦.
一而再一而再說.
很多事情是有聖靈.
有神跡在其中來管理介入的.
所以很特別的一個《少道行傳》的記載.
你可以想像到.
盧家自己本來是深信不疑.
他更加希望有更多的人.
認識神的作為.
所以斯提反這個殉道.
是多麼美麗的一件事.
當然一點都不美麗.
面容模糊.
你想像一下整個人.
慘死的當場.
但是他看到的是神的兒子.
他看到的是天開了.
他看到的是神的榮耀.
他看到的是主耶穌站起來迎接他.
當然我們不是鼓勵大家去殉道.
我們最近為一位在土耳其被監禁的伊朗人.
Mustafa.
我再問清楚Chris弟兄.
因為他們有直接的之前的對話.
這位弟兄因為他是伊朗人.
在土耳其被政府抓了.
因為他信了耶穌.
所以要遞誡回伊朗.
未遞誡之前就在土耳其的監獄服刑.
其實等待那個什麼時候批就將他遞誡.
同一時間.
在加拿大的教會知道這件事.
就為他申請難民資格.
希望他能夠去到加拿大.
免去這個災難.
但是不知道的.
一直的移民審批都被拒絕.

$^{561}$很快的如果再沒有上訴.
他就要遞誡回伊朗.
不過想說的就是.
Chris弟兄告訴我們.
這個伊朗人在監獄裡.
居然得到土耳其監獄裡的獄卒的喜愛.
來看到他的生命.
然後將囚犯交給他.
因為他之前是做一些營商.
幫人剪頭髮.
類似美容.
獄卒就叫Mustafa幫其他囚犯去做理髮.
在監獄裡.
給他一定的自由.
有一次他居然帶了一個土耳其的囚犯.
信了主耶穌.
在監獄裡本來自己差不多要死的.
回伊朗.
他帶了這個人信耶穌.
又讓教會其他人知道這件事.
好奇妙地看到神的聖靈工作.
所以這裡告訴我們.
弟兄姊妹.
我們今天透過這段經文.
很快我們就要去第八章.
去到《肥尼》.
作者先交代了.
為什麼這麼多的信徒要急切離開耶路撒冷.
就是因為斯提反.
首當其衝的.
受到一個嚴重的.
公開的.
正式的迫害.
所以他們就來將這個人打死之後.
這裡也值得我們再讀後面那段經文.
第57節.
七章的57節.
眾人大聲喊叫.
掩著耳朵.
齊心地湧上前去.

$^{601}$把他推到城外.
用石頭打他.
作見證的人.
把衣裳放在一個少年人.
名叫素羅的腳前.
他們正用石頭打的時候.
斯提反呼籲主說.
當然是很痛的.
被石頭.
你有沒有試過打過腳趾.
小小的石頭.
有意識地自己打到自己也很痛.
也不知道怎麼辦.
很多石頭同一時間掉下來.
他呼籲的.
不是自己的痛.
他呼籲的是.
求主耶穌接收我的靈魂.
再一次呼籲.
就是很想學校耶穌.
不要將這些罪歸給他們.
這樣的心腸.
被人這樣打到死.
所以有聖靈充滿.
你可以看到斯提反.
他的心.
說了這番話之後.
就睡了.
就是死了.
素羅都喜悅他被害.
所以第六章到第九章.
作者是按部就班地記載.
三個重要的聖經.
初期教會的人物.
斯提反,肥利.
第九章就是素羅歸主.
很有意識地.
聖靈如何在這三個人身上.
帶來教會的向外擴展的一個很大的工作.
今天我們要講的這段經文.

$^{641}$就是第八章.
我只是選取第四到第八節.
其實下面還有更多.
值得大家去看.
聖靈如何引導肥利.
去帶領一個太監信主.
第四到第八節.
我稍為再快一點讀給大家聽.
相信最重要的鑰節.
就是宣講基督帶來的生命平安.
第八節特別強調.
那個城裡面.
那個是什麼城?.
撒瑪利亞城.
一個對猶太教有敵意的地方.
其實是猶太人對他們有敵意.
不是撒瑪利亞人對猶太教有敵意.
這樣一個很長的時間.
一種宗教的偏見.
所以這個城從來都沒有試過這麼大的歡喜.
就是這個城裡的人.
相信基督帶來的歡喜.
相對第七章.
作者很想做一個對比.
司提反的正道.
向一班自己的同胞猶太人.
講一個大家都耳熟能詳的一個.
自己的民族的歷史.
講自己的民族的這個聖殿.
都是神藉著所羅門興建.
來到今天都有這個敬拜.
但是一個敬拜神的耶路撒冷.
宗教的中心.
敬拜真神的地方.
充滿的只是因為對真理的誤解.
而來的偏見.
來的怨氣.
而來的憤怒.
對神的僕人司提反.
來攻擊.

$^{681}$來殺害.
但是第八章.
只不過宣講基督.
整個城.
撒瑪利亞城.
原先是猶太人所鄙視的地方的人.
整個城裡面.
得著基督.
得著那個歡喜.
得著醫治.
得著的是再無歧視.
至少他們去到撒瑪利亞.
從來沒有一個對別人.
對這個民族的低等的歧視.
這是過去猶太人的歷史.
如果大家熟悉.
約翰福第四章.
你和一個撒瑪利亞人.
來到港島.
猶太人和撒瑪利亞人是沒有來往的.
他們是繞道不去撒瑪利亞的.
但是.
你看看.
初期教會的領袖.
在聖靈的充滿裡面.
連帶這些難以梳理的種族問題.
一下子就沒有了.
因為在主耶穌裡面.
所有的民族都是合一.
歸於一.
這裡告訴我們.
肥利的宣講.
第七章就沒有了.
第八章就有了.
那些人看著.
留意聖經的描述.
第六節.
眾人聽見又看見肥利所行的神蹟.
就同心合意聽從他的說話.
再加上原因就是.

$^{721}$因為有許多人當場被污鬼附著.
那些鬼大聲的呼叫.
從他們身上出來.
還有許多癱瘓的.
瘸腿的.
都得到醫治.
所以在城裡面.
如果是一次性的報道會.
就真的要賣新聞.
上東張宣告了.
但不是的.
這個城裡面大有歡喜.
聖靈的工作.
不過留意了.
斯提反在港獨裡面.
完全沒有任何的神蹟.
因為原來.
保羅在哥林多前書一章告訴我們.
原來猶太人有一個.
他們認為先知.
認為神的代言人.
其中的條件.
猶太人是要求要有神蹟.
希利尼人.
就是希臘人.
就要你講話.
你要有智慧.
你要有哲學.
我就會聽了.
如果你說話沒有實力.
又沒有什麼嘔吐文字.
我們就不聽了.
保羅是很了解.
那個時代的特性.
所以猶太人.
聽著斯提反的講道.
他們看不到斯提反有什麼神蹟.
走出來.
所以他們聽完.
再加上斯提反的總結.

$^{761}$對那些回眾的對號入座.
你們就是殺害先知的子孫.
你們拒絕律法.
把矛頭指向回眾.
回眾當然是憤慨.
施提反於死地.
不過在這段經文.
非常寶貴.
肥膩.
充滿了聖靈.
和智慧.
來到這個薩馬利亞.
很多時候第一到第三節.
都是很多時候.
在近年的社會運動.
很多人喜歡引用的經文.
教會受迫迫.
所以我們要走是很合理的.
完全沒有問題的.
每個人有自己的決定.
有自己的想法.
希望為下一代著想.
是完全沒有問題的.
但這段經文不要被濫用.
這段經文我們要留意.
由第六到第九章.
作者是很有意識地記載聖靈.
在幾位主要領袖的心靈裡的工作.
和他們帶動的教會.
是有意識地向外.
去履行大使命的福音.
所以如果你們沒有.
不知道大家是否懂得背.
《小徒行傳》第一章.
主耶穌的吩咐第八第九節.
你們要在耶路撒冷等待.
直到什麼時候.
聖靈降臨在你們身上.
你們就必得著能力.
並要怎樣呢?.

$^{801}$在耶路撒冷.
猶太傳地.
撒瑪利亞直到地極.
作我的見證.
所以原來.
作者是很想告訴我們.
他這樣根據耶穌的吩咐.
現在到第八章.
就是這樣進行.
是怎樣的時態背景呢?.
這個不是人可以操控.
因為的而且確.
神可以用耶路撒冷誤會基督教的人.
來作為一個踏腳石.
好像一個令事情白熱化的原因.
就是當時整個教會受到迫害.
大遭逼迫.
第八章第一節.
由那天開始.
司提反遜道那天開始.
耶路撒冷的教會大遭逼迫.
除了使徒以外.
門徒也分散在猶太和撒瑪利雅各處.
所以不要那麼容易將這些歷史的經文.
很快就對號入座.
每個人都可以離開.
一定要離開.
全教會都要離開.
使徒不准離開.
這樣就有時過了火位.
每個人有自己的決定是完全沒有問題的.
而且我們明顯看到.
現在在英國.
在加拿大.
那個復興是明顯看到的.
真的看到神的工作.
所以我們在當下.
未必馬上可以分析到.
馬上可以看到端倪.
但是每個人按照著神的感動.

$^{841}$是沒有問題的.
不過經文告訴我們.
在整個耶路撒冷第二節.
就已經出現了信徒的隨空大哭.
那些得勝的會堂.
那些猶太教的人.
充滿了憤怒的勝利.
還不死?.
你這麼厲害.
這麼有勝利充滿.
就先殺了你.
得勝了.
但是神的工作.
怎麼會這樣被憤怒.
被敵視影響呢?.
神的工人.
在這樣的情勢下.
不是四散逃跑.
是分散在各處傳道.
大家有沒有留意這些描述.
是帶著意識.
帶著感動.
帶著使命.
這樣去傳道.
是聖靈在他們生命裡的感動.
所以要走的人.
真的要記住.
沒有問題的.
不過真的要有這樣的考慮.
就是要去到各處.
為主耶穌作見證,傳道.
所以我們在撒瑪利亞城.
感受到的就是.
宣講基督.
整個城.
經歷一個大的復興喜樂.
就是救恩臨到這些被藐視的猶太人.
藐視的撒瑪利亞人.
整個城大有歡喜.
因為宣講基督.

$^{881}$弟兄姊妹.
我們真的必須要有這個.
生命裡面的內容.
就是基督在我們的生命裡面.
和在別人的生命裡面的重要.
很深刻的例子.
一直預備的時候浮現.
有一位很熟悉的弟兄.
他是第二代的基督徒.
在原浸的弟兄.
在讀神學的時候.
在他的團契.
過了一段時間.
現在我們每個月都有聚會.
都有弟兄的茶經.
他們想我帶茶經.
沒所謂的,當大家吃東西.
或者一晚.
總之每個月一次.
一個茶經班.
他講一個見證.
他有一個好朋友.
知道他離流.
去醫院盡量探望他.
和他聊天.
他有一個內地結婚的太太.
不信耶穌.
他和這個朋友.
這個朋友當然不信耶穌.
和這個朋友在醫院.
簡單地說.
分享福音.
不行了,最終都離世.
都不信.
過了一年.
這個在深圳住的太太.
其實不認識,不太熟悉.
就打來給這個弟兄.
你那次和我先生講道.
他不信.

$^{921}$但我信.
他回到深圳之後.
他自己不知為何去找教會.
信.
再過了兩年.
就邀請弟兄上深圳教會洗禮.
見證他的洗禮.
好奇妙的一個工作.
弟兄說我從來沒有想過.
和他講話.
我只是安慰一下那個朋友.
那個朋友也不肯信耶穌.
但是居然.
誤中附居.
他的太太.
聽完之後回到深圳.
其實沒有人再跟進,沒有人目仰.
他自己去深圳找教會.
宣講基督.
就帶來生命的改變.
其實很多這些例子.
我們每一個弟兄姊妹.
我們都可以成為一個宣講基督的人.
大比別人有生命的改變.
大比別人有工作.
不過記住.
第六到第九章.
這三個重要的人物.
都有一個共同的特徵.
都真是願意讓聖靈管理人生.
讓聖靈在生命裡面.
來作帶領.
你想想.
很快去到第26節.
又記載了另外一件.
肥利的事件.
有主的使者.
對肥利說.
起來向南走.
往那從耶路撒冷.

$^{961}$去加薩的路上.
那條路是曠野.
不知道是不是那條King's Highway.
就是去埃及的.
以前的古道.
就是從以色列的南面.
最遠去到別士巴.
下去到加薩的走廊.
似乎都是一些王道.
因為這個人是一個顯赫的.
總管銀庫的.
好像財政司的官位.
他就上耶路撒冷去守節.
即是說他入猶太教.
不過留意聖經.
作者又再一次告訴我們.
聖靈在這件事上.
扮演的角色.
就是29節.
聖靈對肥利說.
你去貼近那輛車.
肥利就跑去太監那裡.
聽見他在唸先知以賽亞的書.
他有一卷卷軸.
以賽亞書是66張.
以賽亞書的卷軸是兩卷.
如果在死海古卷發現的.
是兩卷.
因為不夠寫.
他們是古卷來的.
卷軸來的.
真跡的卷軸.
不是我們現在的印刷本.
是抄本來的.
你可以想像得到.
這個人一定很有錢.
能夠擁有這些卷軸.
而且他一直這樣回去的時候.
居然在唸先知以賽亞的書.
而且很特別.

$^{1001}$是53章第6節.
他像羊羔被牽到宰殺之地.
又像羊羔被剪毛的人手下無聲.
他也是這樣不開口.
他卑微的時候.
就去到11節了.
以賽亞書的53章.
不按公義審判他.
誰能夠述說他的後代.
因為他的生命從地上奪去.
太監就見到肥利.
這個陌生人這樣走過來.
有聖靈叫他.
就問這位先知.
請問先知說這句話.
是指著誰說的?.
是指著自己.
是指著別人.
肥利就開口.
從聖經開始.
對他傳講耶穌.
聖靈和傳道者.
和傳道者講的信息.
這三樣東西結合.
帶來的就是一個人.
無論他多有社會地位都好.
都在主的大愛.
救恩底下來信.
來得到生命改變.
所以埃塞俄比亞據聞就是這個人.
成為了宣教士.
回去帶來很多人信主.
親愛的弟兄姊妹.
今天我們很想將這幾章聖經.
有一個很重要的主題.
聖靈主宰著我們的人生.
聖靈引導我們內心.
我們在神的旨意下生活.
宣講.
求主讓我們常常有一個很柔軟的心靈.

$^{1041}$讓神帶領著我們.
你看看第七章.
人認為這些道.
只會帶來很多猶太教的影響.
所以他們用他們人很有限的力量.
攻擊神的僕人斯蒂反.
以為可以這樣置他於死地.
這個教就會完結.
但是換來的就是.
信徒當然很傷心.
一個這麼敬愛的斯蒂反.
忠心的服侍.
換來的就是這樣慘死.
為他大哭大哭.
整個城,耶路撒冷城.
本來是喜樂之都.
以上的說話.
但是在這個時候.
他們不肯歸向基督耶穌.
只會帶來很多很多.
整個城裡的那種嫉妒,憤怒.
但是你看看第八章.
作者很想特別形容.
這個本來被人藐視的撒瑪利亞城.
這個時候得到基督帶來的.
就是充滿了歡喜,平安和永恆的福樂.
洪水橋越來越多居民入住.
最近這兩個星期.
有家庭來找教會.
接下來這個星期四.
他們會帶著他們的奶奶來團契.
長者團契.
越來越多人很想認識福音.
當然很想先幫小朋友找社區中心.
一些班祖.
那是一個切入點.
最近都有好幾間教會.
很想了解一下我們的茶座.
為什麼可以有這樣的想法.
有這樣的事供.

$^{1081}$很多的交流.
我們能不能夠讓洪水橋成為一個歡樂之城呢?.
我們能不能夠將基督有力量的宣講呢?.
透過詩歌.
透過我們大家一起不同的事奉.
有人可能是師事.
有人可能是插花.
有可能是其他的侍工.
有詩禽,簡太默默地每個星期.
為我們準備.
當然有古王.
為我們每個星期來服侍神.
我們每一個人能不能夠找到自己的崗位.
很重要在聖靈的充滿下.
來一齊同心的服侍呢?.
很想用這段經文.
藉著這幾個星期所經歷的.
上一個星期我們在嶺南大學.
原來現在有很多學生回應了那些問卷.
所以現在我們有四場茶座的沖奶茶.
我們就會開始.
所以盼望我們能夠藉著宣講.
讓人得到神的救恩.
而帶來心靈充滿了平安和歡喜.
我們一起禱告.
感恩天父讓我們能夠有份於.
神國度的福音的擴展.
多謝你我們能夠得到神的救恩.
真的有很多的前輩.
很多愛我們的人為我們祈禱.
使我們最後都信了主.
我們很盼望我們能夠裝備我們自己.
讓我們的心靈.
很重要讓聖靈管理.
帶來感動.
使我們能夠將神的愛.
有意識地向每一位我們所認識的朋友去宣講.
使他們得到神的救恩.
再一次多謝主給我們今早聖餐的主人.
我們一同來敬拜你.

$^{1121}$我們感恩禱告奉耶穌基督的名.
阿們.
\newpage



\section{路加福音 22:39-46}
\label{sec:td1EIi6hWco}
\textbf{2025.03.16 《願你旨意成就》路加福音22\_39-46 (和修版) 講員 : 曾敏姑娘}
\newline
\newline
連結: \href{https://youtube.com/watch?v=td1EIi6hWco}{\texttt{https://youtube.com/watch?v=td1EIi6hWco}} ~~~~ 語音日期: 2025-03-16
\newline
\newline
\hyperref[sec:oUBZwJOpFGQ]{\small{< < < PREV SERMON < < <}}
~
\hyperref[sec:index_chronic]{\small{[返順時目]}}
~
\hyperref[sec:index_scriptual]{\small{[返順卷目]}}
~
\hyperref[sec:wH4miLAekc8]{\small{> > > NEXT SERMON > > >}}
\newline
\newline
路加福音 22:39-46
\newline
\begin{longtable}{cl}
\hline
\hline
章節 & 經文 (和合本修訂版)\\
\hline
22:39 & \begin{tabularx}{0.7\textwidth}{X} 耶穌出來，照常往橄欖山去，門徒也跟隨他。 \end{tabularx} \\ \\ \relax
22:40 & \begin{tabularx}{0.7\textwidth}{X} 到了那地方，他就對他們說：「你們要禱告，免得陷入試探。」 \end{tabularx} \\ \\ \relax
22:41 & \begin{tabularx}{0.7\textwidth}{X} 於是他離開他們約有一塊石頭扔出去那麼遠，跪下禱告， \end{tabularx} \\ \\ \relax
22:42 & \begin{tabularx}{0.7\textwidth}{X} 說：「父啊！你若願意，求你將這杯撤去；然而，不是照我的意願，而是要成全你的旨意。」 \end{tabularx} \\ \\ \relax
22:43 & \begin{tabularx}{0.7\textwidth}{X} 〔 有一位天使從天上顯現，加添他的力量。 \end{tabularx} \\ \\ \relax
22:44 & \begin{tabularx}{0.7\textwidth}{X} 耶穌非常痛苦焦慮，禱告更加懇切，汗如大血點滴在地上。〕 \end{tabularx} \\ \\ \relax
22:45 & \begin{tabularx}{0.7\textwidth}{X} 禱告完了，他起來，到門徒那裡，見他們因為憂愁都睡著了， \end{tabularx} \\ \\ \relax
22:46 & \begin{tabularx}{0.7\textwidth}{X} 就對他們說：「你們為甚麼睡覺呢？起來禱告，免得陷入試探！」 \end{tabularx} \\ \\
[1ex]
\hline
\hline
\end{longtable}
$^{1}$各位弟兄姊妹平安.
感激神讓我們每個星期日都有福氣.
回到神的家去聆聽祂的話.
我們開始的時候.
再有一個禱告.
讓我們仍然活著.
我們仍然可以去聆聽你的話.
去回應你.
甚至為了我們生命的罪去認罪悔改.
主在這個世界上沒有東西彼得上.
我們和你的關係這麼重要.
所以每一次崇拜的時候.
我都求你對我們說話.
今日都求你對我們眾人說話.
讓我們每個人都聽得清楚明白.
願意聖靈在我們心裡動功.
讓我們明白神你自己的心意.
更願意聖靈親自對我們說話.
奉耶穌基督的名字祈禱.
阿們.
其實大概二二年的時候.
我都曾經和大家看過這段經文.
願你旨意成就.
我當時用了耶穌基督的圖文.
願你旨意成全.
去說那篇信息.
今日神感動我.
從一個很不同的角度.
去說這段經文.
盼望成為我們的提醒和幫助.
願你旨意成就.
這是出自耶穌基督的禱告.
如果大家記得.
耶穌基督本身是最體重.
父神的心意.
是最體重父神的旨意.
我們從哪裡看到呢?.
有幾個地方.
我們一同去看看.
第一.

$^{41}$原來祂很小的時候.
大家記不記得.
祂十二歲的時候發生什麼事?.
十二歲祂做什麼?.
大家都有印象.
我聽到有些聲音.
原來祂十二歲的時候.
大家記得.
祂和父母上耶路撒冷去守節.
守完節之後.
父母已經回到家鄉.
但祂仍然留在耶路撒冷.
留在耶路撒冷什麼地方呢?.
在聖殿那裡.
一路和人分享和傾談.
然後祂的父母.
過了三天才在聖殿找祂.
非常焦急.
他們見到兒子.
當然很重要.
於是當時.
祂的母親怎樣跟他們說?.
兒子,為什麼你這樣對我們?.
你看看你爸爸和我.
很焦急.
到處找你.
耶穌的回應是什麼?.
為什麼要找我呢?.
難道你們不知道.
我應當在我父的家裡嗎?.
誰最重要?.
天父最重要.
天父的家最重要.
其實關乎天父的旨意.
也是最重要的.
雖然地面上.
父母是重要的.
我們不會說.
地面上父母不重要.
但天父更加重要.

$^{81}$他關心的是天父的事.
是天父家的事.
所以你看看他.
焦點是很不同的.
想想大部分小朋友.
12歲在做什麼?.
玩耍.
對人生有很多期待.
在品嘗父母給他的美食和好處.
但耶穌基督12歲.
就在掛念.
很注重天父家.
下一個就是主禱文的教導.
大家記不記得主禱文的教導.
是怎樣說的?.
頭幾句.
接著.
大家今天真的很好.
非常精神.
記得清楚主禱文所說.
願你的旨意.
行在地上.
如同行在天上.
主耶穌教導門徒禱告.
首先不是說.
我知道大家身體有很多不適.
首先求耶穌基督醫治你.
或者我們有很多生活,工作,家庭的需要.
你就不停將這些需要交到上帝的手裡.
這些我們都可以向神禱告.
是可以的.
但首先要關注的不是這些.
關注的是什麼?.
是願人尊上帝的名為性.
是願祂的國降臨.
是祂的國.
然後再說祂的旨意.
是祂的旨意行在地上.
如同行在天上.
這才是重要的.

$^{121}$比我們肉身的需要重要.
比我們地面上的一切.
肉身的很多關注重要.
耶穌基督是這樣教導門徒禱告.
你記住主禱文不是說.
傳道人的禱告.
也不是宣教士的禱告.
是給每一個真正跟從耶穌基督的人.
跟從上帝的人的禱告.
換言之今天這個主禱文.
是對我們每一個都有關的.
而上帝就是要我們關注祂的旨意.
這是耶穌基督的教導.
祂的教導是這樣.
然後去到第三個位.
那個高峰就是橄欖山上的禱告.
在赫西瑪利園的禱告.
更加去到一個極點.
祂不是教人祈禱.
祂自己也是這樣走出來的.
在這個赫西瑪利園裡.
在橄欖山上.
這個位置是一個很艱難的位置.
耶穌基督即將踏上受苦之路.
即將要面對被出賣被捉拿.
一連串的受苦.
最後要釘十字架.
在這個禱告裡.
你會看到祂在說什麼.
我們回到經文裡.
祂這樣說42節.
「父啊!你若願意.
求你將這杯撤去.
然而不是照我的意願.
已是要成全你的旨意」.
祂最後的這個禱告也是說.
成全你的旨意.
是願你旨意成就.
去到這裡.
就要踏上受苦.

$^{161}$被害.
犧牲的這一步.
祂的禱告也是願你旨意成就.
是要成全你的旨意.
祂以祂自己的生命.
告訴我們什麼是最重要.
祂由始至終.
都是堅持這一點.
祂的旨意是最重要的.
那麼父神的旨意又是什麼呢?.
我們分開兩個層面.
第一個.
父神對於我們.
整個宇宙萬物.
我們人類.
有一個很終極的旨意.
是什麼呢?.
約翰福音三章十六節.
是在說什麼呢?.
有沒有人背得到?.
背得到就大聲說.
應該是我們的金句.
是不是?.
我記得周牧師.
是要我們背金句的.
說什麼呢?.
真是太好了.
周牧師在這裡就很開心.
在這裡.
和修版是這樣說.
「甚愛世人.
甚至將他獨一的兒子.
賜給他們.
叫一切信他的人不自滅亡.
反得永生」.
我們天父上帝的旨意.
就是要藉著他的愛子.
耶穌基督成就舊恩.
拯救萬民的生命.
讓萬民都進入.

$^{201}$他愛子的國度裡.
可以永遠和上帝一起生活.
這是他最終極的旨意.
也是最重要的.
上帝的國.
福音工作.
這是我們都有份的.
我不斷強調.
聖經裡說的.
不是對著傳道人說的.
而是對著宣教士說的.
今天在座每一個.
聽過這堂道的人.
我們都有份的.
沒有人是例外的.
我想說.
這個父神的旨意.
是對著我們每一個人說的.
父神的旨意是這樣.
而耶穌基督很重要的.
就是要去成就.
父神的旨意.
就是要去成就舊恩.
這是他當時.
渡成肉身來到這個世界.
很重要的.
很重要的去成就.
當然父神的旨意.
對於整個人類.
是一個救贖的計劃.
但是對於我們每一個人.
神都有他的旨意.
是每一個人.
上帝對你和我.
有旨意的.
不會只是對著.
某一個人才有的.
從一開始.
我們在媽媽肚腹的時候.
我們都不認識她.

$^{241}$到上帝來到這個世界.
父神對我和你就有旨意的了.
來到這個世界之後.
一路經歷人生的高高低低.
之後我們有機會.
信耶穌了.
我們進入神的家.
一樣上帝的旨意.
都沒有離開過我們.
不過重點是.
我和你.
知不知道神的旨意.
我問問在座這麼多人.
問問在座這麼多人.
回教會這麼多年.
大家是否真的知道.
神對你個人的旨意.
是甚麼.
我們是否真的.
行在其中.
是不是.
神你有你的旨意.
我有我的旨意.
永遠都是分開的.
你和你的 我和我的.
神對每一間教會都有祂的旨意.
我們全世界這麼多地方.
每個地方各自都有教會.
大部分的地方.
可能有少部分地方.
福音未曾編輯.
這不奇怪.
但大部分地方都有教會.
神對每一間教會都有祂的旨意.
是不一樣的.
每一間教會.
是否真的明白神對我們的旨意.
是甚麼.
是否真的知道.
還是上帝.

$^{281}$你有你的旨意.
我們教會喜歡怎樣.
我們就有我們的旨意.
這是兩回事.
上帝是有祂的旨意.
教會有沒有尋求.
神的旨意.
這一定要問.
教會不是純粹人的教會.
是耶穌基督的教會.
你記住.
我們說.
耶穌基督是作王.
在作主.
祂是我們的頭.
但我們是否真的.
問耶穌基督.
你對我們教會的旨意是甚麼.
有多少間教會真的會問.
這是很重要的.
值得我們留意.
耶穌基督很重要.
就是真理.
道路 生命.
祂就是那條橋.
要透過耶穌基督.
我們才可以走向永生.
在神的國度裡有份.
真的需要耶穌基督這條橋.
這裡問.
我們最看重的是甚麼.
耶穌基督最看重.
父神的旨意.
我們最看重的是甚麼.
這裡有沒有我們.
所看重的呢.
金錢.
地位.
你說不會的.
信耶穌的人.

$^{321}$也不是的.
我們的老外沒有被致死.
我們轉了場景.
進入教會.
在教會裡的渴望.
得到很多人的讚賞.
肯定 名譽.
一樣有地位.
我們追求甚麼.
甚麼是最關注的.
我們是看人.
最重要的是看人的關係.
最重要的是我配偶.
最重要的是我子女.
甚麼都是最重要.
是我身邊最關心的人.
他就會影響.
我的一切.
他就是我的一切.
一動了他.
我就跟你拼過.
我不放過你.
甚麼是你最關注的.
其實自己心裡知道.
我們每個人都很想有健康的身體.
是好的.
神都想我們管理好我們的身體.
管理好我們的健康.
但如果去到一個地步.
去到一個極點.
已經主宰了.
已經超越了我們對上帝的渴望.
成為所有東西當中.
最頂級的.
對我們渴望的東西.
但又要小心.
是甚麼.
在地面上.
我們看重甚麼.
問自己.

$^{361}$如果做一個比較.
我們在座的這麼多人.
包括我自己.
我是否最看重.
父神的旨意.
還是看重.
我自己看重的.
我們最看重的.
是否上帝.
是否成就我心中的渴望.
這樣有分別的.
我心中渴望的東西.
不一定是上帝的旨意.
我一直渴望.
上帝快點.
按照我渴望的.
你成就吧.
我日夜祈禱就是這樣.
我告訴你.
我最渴望的東西.
你為我祈禱吧.
你有沒有祈禱呢.
我想明白.
上帝對我的旨意是甚麼.
我很想.
上帝的旨意成就我身上.
如果上帝的旨意.
和我的心意不同呢.
那怎麼辦.
打算怎麼辦.
不行.
我要成就我的心意.
我要成就上帝的心意.
我不同的.
他的心意和我的心意不同.
是不是呢.
記住.
在耶穌基督拯救我們那一天.
在我們相信耶穌的那一天.
我們的生命是屬於誰的.

$^{401}$不是再屬於自己的了.
耶穌基督用他的生命.
拯救我們生命的那一天.
記住.
我也跟自己說.
我要記住.
我們的生命是屬於耶穌基督的.
是不是.
這幅圖在座的.
大家是很用心傳福音的.
弟兄姊妹.
一定見過的.
這幅圖.
我們經常跟人說.
你看圓圈.
代表你的生命.
在圓圈當中有張椅子.
代表你的生命的核心.
你的中心.
誰坐在椅子上.
代表你的人生.
是誰在掌管的.
在未信耶穌前.
就我.
坐上那張椅子.
十字架代表耶穌基督.
原來他在圈外.
他和我的生命沒有關係.
我在管理我的人生.
我們是這樣.
跟未信主的朋友說的.
於是就邀請他.
你決志信主.
讓十字架.
坐在寶座上.
你坐在旁邊.
等神帶領你.
我想問.
到現在當我們信耶穌之後.
這幅圖有沒有改變呢.

$^{441}$我們誠實面對自己.
真的讓十字架坐在椅子上嗎.
我們真的.
真的讓十字架.
坐在椅子上嗎.
真的.
凡事以耶穌基督的心意.
為最重要.
以父神的旨意為尾.
真的嗎.
有多少.
真的交出來.
給主耶穌去管理.
有多少我仍然留給自己.
其實這幅圖不只是對於.
未信耶穌的朋友說.
會不會都是我們的光景.
我們仍然坐在上面.
不過耶穌的位置好一點.
十字架應該是在圈內.
但他不在椅子上.
大部分時間都是我作主.
你也為我祈禱.
有些事我很想達成.
你也為我祈禱.
但還是坐在上面.
不要真正將你生命的主權.
交出來.
如果願父神的旨意成就.
就真的要知道.
他對我們每一個人.
的旨意是甚麼.
對教會的旨意是甚麼.
明白上帝對萬民的.
旨意是甚麼.
我們就投身在當中.
主耶穌要去成就.
父神旨意.
是有掙扎的.
不是沒有掙扎.

$^{481}$你看看經文.
首先在禱告裡面.
他真的有掙扎.
所以剛才才說.
你若願意.
就將這個杯切去.
我就不需要去喝.
這個杯.
然而不是照我的意願.
而是要去成全你的旨意.
耶穌基督知道.
要去成全父神的旨意.
但也有求.
父神將這個杯切去.
因為祂知道這個杯.
不是一個很簡單的杯.
這個杯只有祂一個.
可以去承受.
只有祂一個人可以去喝.
但卻是非常艱難.
在經文裡面.
你看到有位天使.
天上顯現.
加添祂的力量.
這一刻是一個很激烈的.
拉鋸.
自己的心靈的掙扎.
耶穌非常痛苦.
焦慮.
禱告更加懇切.
是一個掙扎.
汗如大血點滴在地上.
很厲害.
汗是大滴滴下來.
掙扎很激烈.
這不是一件普通的事.
在這個世界只有祂一個人可以做.
我和你做不到.
我和你真的做不到.
沒有人做到.

$^{521}$只有祂一個人.
整件事.
祂是否要去信服.
父神的旨意.
是一個很深的掙扎.
結果主耶穌選擇了.
信服父神的旨意.
如果我們看.
在聖經裡面.
如果我們看.
馬太福音.
記載.
赫西瑪利元禱告的時候.
在馬太福音的記載裡面.
直接說到.
三次的禱告.
三次裡面.
都是有說祂的掙扎.
第一次說.
我父啊,如果可能,求你使這杯離開我.
然而不是照我所願.
而是照你所願.
其實這個.
所願的.
願意和旨意.
是屬於同一個詞組.
在原文裡面.
所以耶穌基督說.
你願不願意.
不是你喜不喜歡.
你想不想.
而是說這件事.
天父,如果這是你旨意的一部份的話.
如果是你旨意的一部份.
你就叫這杯離開我.
離開我.
祂還是說.
是父神的旨意.
祂的心很維繫在這件事裡.
第二次.

$^{561}$祂又去禱告.
我父啊,這杯若不能離開我,必須我喝.
就願你的旨意成全.
還在求.
還有沒有機會呢?.
這杯.
真的可以離開祂呢?.
會不會在父神的旨意裡面.
也有這一部份呢?.
但祂仍然最後還是說.
願你的旨意成全.
還在禱告.
最後在禱告過後.
祂信服.
選擇信服的代價.
很大.
今天我就要和大家說.
很多時候.
我們覺得我們選擇信服.
神的旨意.
神就開了路.
非常順暢.
毫無難阻.
所有事情都順利.
我就可以.
去到終點站.
這就是父神的旨意.
是不是呢?.
不是啊.
是付上很大的代價.
首先.
在選擇信服的時候.
都要面對.
屬靈的爭戰.
我們回看經文.
在路加福音.
說了.
耶穌說.
39節說:出來照常往.
40節說:到了那地方.

$^{601}$他就對他們說.
是說門徒.
你們要禱告.
免得陷入試探.
這裡有個爭戰.
撒旦從來都沒有缺席.
耶穌基督很重要.
祂要成就父神的旨意.
這個旨意很重要.
關乎萬民的得救.
撒旦總是要.
迷惑人心.
要人心軟弱.
不要跟從父神的旨意.
是不是?.
所以耶穌基督.
不單止自己禱告.
祂想門徒和祂一起禱告.
也想門徒在禱告裡.
要站立得穩.
不要陷入試探.
陷入惡者的詭計裡.
惡者從來都很希望.
耶穌的跟從者會失敗.
從來都沒有放棄過.
要做這件事.
今天當我們說.
要去服事神.
我們要希望.
為主做更大的工作的時候.
惡者都沒有忘記.
也要跟我們來爭戰.
所以在這裡要禱告.
是要面對爭戰.
是要禱告.
這是非常非常重要.
在《馬太福音》裡.
也看到耶穌基督.
很希望門徒.
和祂一起.

$^{641}$在禱告裡.
經歷很重要的時間.
但是門徒.
心靈固然懸疑.
肉體卻軟弱了.
在那時候就睡了.
他們不知道.
這時候的爭戰是何等激烈.
在這裡.
更大的是什麼呢?.
在這裡.
經文說.
父啊,你若願意求.
你將這杯切開.
這杯是指什麼呢?.
這杯是指.
神的烈怒.
和耶穌基督.
將要面對的苦難.
神憤怒.
因為什麼而憤怒?.
因為我們世人犯罪.
虧欠上帝的榮耀.
哇,說起人的罪就多了.
不只是現在外面.
人們犯的罪.
我們也有犯罪的.
教會圈子裡.
教會外面.
歷世歷代.
多少人犯罪.
得罪上帝,多少罪.
有顯現的,有隱而未然的惡.
世人的罪.
我們也在裡面.
何其多.
神在何等憤怒.
本來是我們全世界的人.
要承擔,要面對.
上帝的憤怒.

$^{681}$但是現在就在這個杯裡.
耶穌基督所承受的.
裡面.
包括了上帝的憤怒.
上帝對全世界人的憤怒.
落在耶穌基督.
一個人的身上.
是何等大呢?.
這是何等難的承擔呢?.
接著就說苦難.
苦難的,大家都知道.
之後一直下去.
是被他的門徒出賣.
被捉拿.
面對不公義的審判.
戲弄,鞭打.
最後被釘死在十字架上.
為的是什麼?.
就是全世界的人的罪.
只有他一個可以做.
何其多.
我們也有犯罪的.
只有他一個可以做.
我和你.
都不能去做,都背不住.
都不能去喝這個杯.
就是耶穌基督去喝的.
嘩,這個.
是多難去承擔呢?.
主耶穌就是一個.
去承擔這一切.
這個代價.
是非常非常大.
是因為付上了這個代價.
今天我和你才可以坐在教會的門內.
在這裡崇拜.
生命改變.
從今以後好有盼望.
由現在.
直行到永生,永永遠遠.

$^{721}$是何等寶貴.
但這個代價是非常非常大.
耶穌基督所要承受的.
是要殉道艱難.
他的痛苦,他的死亡.
也是獨一無二的.
因為耶穌基督所作的.
是為全世界的人所做的.
信服是要付代價.
今天很想分享.
幾個不同的人的經歷.
說說付的代價是怎樣的呢?.
這個有沒有人知道他是哪一位呢?.
好像有些年代.
原來他是一位宣教士.
一位宣教士曾經.
想去中國宣教.
結果他在北朝鮮.
這個地方殉道.
他叫做湯瑪士.
很多年前.
他27歲的時候.
想去中國宣教.
但他未曾去到中國.
在坐船去到韓國的時候.
在大同江邊.
被人捉了.
當時他帶著聖經.
被捉後被斬首處死.
死的時候.
27歲.
非常年輕.
死之前.
他兩手高舉的聖經.
大聲喊著.
耶穌.
在這個過程裡.
有人一直目睹整個過程.
大家記得.
他有聖經.

$^{761}$有一個11歲的少年.
叫做崔志良.
把他那本聖經.
拿回家.
但他也知道.
當時這本聖經是禁書.
他不能夠就這樣留下.
於是他就把他給了.
當時在平壤.
作為監聽警備的.
一個叫朴永植的人.
把這本聖經給了他.
想讓他處理.
很有趣.
朴永植沒有把他.
銷毀掉.
朴永植就做了一件.
很奇怪的事.
他回家把聖經的內頁.
全部撕下來.
然後把聖經當成牆紙.
貼在牆上.
但朴永植不是信耶穌.
不敬畏上帝.
當時其實可以這樣說.
福音根本就沒有遍傳在.
北朝鮮.
是一個福音非常封閉的國家.
傳福音是很艱難的.
但聖經就這樣貼在牆上.
發生什麼事呢?.
後來神的話語很有能力.
非常有權能.
朴永植就因為牆上的經文.
久而久之又看.
又看.
就入了他的心.
因為經文接受了耶穌基督.
神多厲害.
聖靈親自的工作.

$^{801}$藉著自己的說話.
要他信耶穌.
後來.
交聖經給崔志良.
因為他常常.
去朴永植的家.
受到朴永植的影響.
他又信了耶穌.
後來這個崔志良.
成為了平壤的.
愛野村教會的長老.
帶領教會.
多奇妙.
而朴永植的房間.
是貼滿聖經的房間.
後來成為了平壤的第一間教會.
叫做張台彥教會.
哇.
這間教會神大大的使用.
就是1907年.
平壤大復興的發源地.
整件事.
是神的心意.
神的旨意.
神的旨意是要福音傳到北朝鮮.
要福音一直在人心裡.
工作.
這就是教會.
你想都沒想過.
如果讓你想.
最好是他的國家很開放.
你自由的進去傳福音.
然後很舒暢地.
跟教會和信耶穌的人.
做栽培.
光明正大地.
跟著跟他們祈禱.
栽培他們.
我們正常都會這樣想.
但神不是用這個方法.

$^{841}$是有人信道.
他付上了生命的代價.
付上了生命的代價.
我想.
這個湯瑪士離世的時候.
他根本看不到.
之後發生的事.
他只是付上了.
自己的生命代價.
但上帝旨意成就.
上帝只是.
過了一下.
讓他的生命順服.
你聽了.
那你想想.
付代價.
你看看後來.
很多人進入教會門內.
是一個大復興.
湯瑪士的代價.
是很值得.
是生命的代價.
但要這樣付上.
大家有沒有想過.
我們聖經裡的保羅.
保羅的呼音工作.
是做了很好的呼音工作.
也帶領很多人.
歸向他.
上帝對他的旨意.
一開始就是要他.
向外邦人傳呼音.
他為此付上很大的代價.
他受了很多苦.
而他去到最後的時候.
大家記得.
在《使徒行傳》20章的時候.
22節大家聽我說.
他說現在我被聖靈催迫.
要往耶路撒冷去.

$^{881}$我知道在那裡會遭遇什麼事.
但知道聖靈在.
覺醒裡向我指正.
說有困所與患難等著我.
我卻不以.
性命為念.
只要走完我的路程.
完成我從主耶穌所領受的職份.
父神的旨意.
就是要他向外邦人.
傳呼音.
現在.
聖靈催迫他.
去耶路撒冷.
但在那裡患難等著他.
他要付代價.
他想不想呢?.
他想!.
他不以性命為念.
他會付這個代價.
就要去完成.
從主耶穌所領受的職份.
今天我想說.
弟兄姊妹你們準備好了嗎?.
其實我們是要付代價的.
每個人付的代價.
都不一樣.
神也有給我們.
教會的旨意.
我們願不願意.
順服主的旨意.
成就在我們身上呢?.
我們要去知道.
神對我們每個人的旨意.
是什麼.
我們願不願意為此付上代價呢?.
你一定要問自己.
不要就這樣活著.
這樣活著.
差不多離開世界的時候.

$^{921}$感覺到.
好像沒有為主做什麼.
似乎做了很多事.
但都不是主要我們做的.
是不是?.
這是很不值得.
活著,現在我和你活著.
不應該再為自己而活.
乃是為耶穌基督而活著.
就願主的旨意成就在我們生命裡.
並且要準備.
為此付上代價.
我記得我最近也有和電影姐妹.
分享自己的夢照.
我記得我讀書的時候.
心裡.
總是很想.
將來有一天畢業.
去找一份工作.
或者是做生意.
賺大錢.
過舒適的生活.
好像很多人一樣,組織家庭.
總之是人生的享受.
這是我.
個人喜好.
所喜歡的.
但父神對我的旨意.
不是這樣.
父神呼召我.
去做一個福音戰士.
要傳福音.
在這個過程裡.
我有我的心意.
父神有父神的心意.
在裡面有一個很大的拉力.
拉鋸在裡面.
經歷了很多年的時間.
最後我終於.
聽到父神的旨意.

$^{961}$我發覺.
這是無比的.
因為我想說.
我所得著的.
遠遠勝過今天在地面上.
所有的金銀財寶.
所有的生活的享受.
吃喝玩樂.
主耶穌基督比萬有都寶貴.
能夠得著.
就是我這一生.
最大的福氣.
我自問自己不配.
各位弟兄姊妹.
不是每一個人.
做傳道人宣教士才會這樣想.
是每一個人.
甚願大家以得著耶穌基督.
為自保.
我們一同禱告.
主耶穌基督在這裡.
以你自己的生命.
去見證.
如何看重父神的旨意.
願父神的旨意.
成就在我們的生命裡.
成就在教會裡.
成就在整個宇宙裡.
主耶穌求你幫助我們每一個.
放下自己.
其實那幅圖畫.
不只是給未信主的人.
可能今天就是我們的光景.
我們仍然坐在.
我們自己生命的座位上.
沒有真正將主權交出來.
主耶穌今天求你.
真的去寬恕我們.
寬恕我們那種自我.
寬恕我們仍然.

$^{1001}$是這樣的.
是真的違背神的心意.
求你讓我們.
真正.
讓你坐在.
我們生命的寶座上.
掌管我們一生.
讓我們的生命活著.
只是願意.
天父的旨意成就.
願我們的生命榮耀你.
願教會榮耀你.
整個世界裡.
萬事萬聞.
所有事證都要榮耀.
我們的父上帝.
奉耶穌基督的名字祈禱.
阿們.
\newpage



\section{哥林多前書 15:55-58}
\label{sec:wH4miLAekc8}
\textbf{2025.03.23《絕不徒然的事奉》 哥林多前書15\_55-58 (和修版) 講員 : 林誠信牧師}
\newline
\newline
連結: \href{https://youtube.com/watch?v=wH4miLAekc8}{\texttt{https://youtube.com/watch?v=wH4miLAekc8}} ~~~~ 語音日期: 2025-03-25
\newline
\newline
\hyperref[sec:td1EIi6hWco]{\small{< < < PREV SERMON < < <}}
~
\hyperref[sec:index_chronic]{\small{[返順時目]}}
~
\hyperref[sec:index_scriptual]{\small{[返順卷目]}}
~
\hyperref[sec:gznE_6kHeYc]{\small{> > > NEXT SERMON > > >}}
\newline
\newline
哥林多前書 15:55-58
\newline
\begin{longtable}{cl}
\hline
\hline
章節 & 經文 (和合本修訂版)\\
\hline
15:55 & \begin{tabularx}{0.7\textwidth}{X} 「死亡啊！你得勝的權勢在哪裡？死亡啊！你的毒刺在哪裡？」 \end{tabularx} \\ \\ \relax
15:56 & \begin{tabularx}{0.7\textwidth}{X} 死亡的毒刺就是罪，罪的權勢就是律法。 \end{tabularx} \\ \\ \relax
15:57 & \begin{tabularx}{0.7\textwidth}{X} 感謝神，他使我們藉著我們的主耶穌基督得勝。 \end{tabularx} \\ \\ \relax
15:58 & \begin{tabularx}{0.7\textwidth}{X} 所以，我親愛的弟兄們，你們務要堅固，不可動搖，常常竭力多做主工，因為你們知道，你們在主裡的勞苦不是徒然的。 \end{tabularx} \\ \\
[1ex]
\hline
\hline
\end{longtable}
$^{1}$領子妹平安.
多謝曾牧師的邀請.
我今天很開心能夠來到宣道會紅恩堂.
和大家在主裡面彼此交通.
和將神的道互相渲染.
今天早了一點來.
我是很impressed(印象深刻).
我一進入你們教會的門口.
就已經覺得非常滿有喜樂.
和很有那種溫度.
不是因為外面天氣熱.
是因為領子妹那種熱情.
剛才德蒙Joseph弟兄.
他招呼我.
也已經享受了你們的咖啡.
不是,我喝了奶茶.
喝了咖啡,我也知道應該是相當好的.
也看到有慕道的親友來了.
你們也有空間接待他們.
我更加impressed(印象深刻).
看到你們的社區.
你們已經在做得很好.
又服侍到泰語的群體.
和有弟兄姊妹.
Sawasdee krap.
其實我和師母.
我們很喜歡去曼谷.
休息一下.
因為我們也很熟悉.
不過我們熟悉的意思是.
我們去那裡也是那幾樣東西.
吃東西,休息,按摩.
所以我們很感恩.
每一次我自己體會.
原來我們的人生.
其實上帝給我們有很多恩典.
但我們的人生也不是永遠的.
上帝給我們原來是有一個時限.
今天我和大家來看看.
剛才主席姐妹已經幫我們看了這段經文.

$^{41}$我嘗試和大家多看一點.
在這段經文之前.
其實在第五十節開始.
是有這個.
對不起,我應該調轉.
不如我們一起讀完這幾節.
在第五十節到第五十四節.
我們一起了讀.
弟兄們,我要告訴你們的是.
血肉之軀不能承受神的國.
必扭壞的也不能承受不扭壞的.
我如今擺一件奧秘的事告訴你們.
我們不是都要睡覺都要改變.
就在一剎那,眨眼之間.
浩同活似吹響的時候.
因浩同要吹響.
死人要復活,成為不扭壞的.
我們也要改變.
這會扭壞的必須變成不扭壞的.
這會死的總要變成不死的.
當這會扭壞的變成不扭壞的.
這會死的變成不會死的.
那時經上所記.
死亡已被勝利吞滅了的話.
就應驗了.
我們在這裡看到.
其實保羅為何這樣提醒他們呢?.
其實因為前文保羅用了相當長的篇幅.
就是講到我們基督的信仰.
一個很重要的關鍵.
就是我們會有復活的.
他說如果我們沒有復活的話.
我們是沒有指望的.
跟其他人有什麼分別呢?.
但今天我們看到我們的人生雖然是短暫.
但我們有一個永恆的家鄉等待我們.
而且我們會在那裡永遠的安息.
要進入那裡的地方.
他說要進入神的國.
是用我們現在這個血肉之軀.

$^{81}$是不能夠承受的.
因為我們這個血肉之軀.
是因為犯了罪的緣故.
我們死亡會臨到.
我們的身體會隨著年紀.
一直都會有變化.
會有扭壞的.
但是在這裡講了有幾個事實.
一個就是必扭壞的.
是不能夠承受不扭壞的.
而且另一件就是說必須要改變.
今天有很多人都不喜歡改變.
他好像想我們一生一世都是無風無浪.
什麼都不改變.
其實不可能不改變.
你看看你身處的社區有沒有改變.
元朗這個社區變了多少.
你看看你旁邊的那位弟兄姊妹.
她有沒有改變.
我的意思是.
她和你幾年前見她的樣子.
會不會有些改變.
都會有的.
我們真的每個人都一直在改變.
但是一個最重要的改變.
就是在耶穌回來的時候.
天使想吹起這個號角.
神的呼喊.
我們見到所有的死人.
要從死裡復活.
如果仍在生的基督徒.
會立時被提.
去到空中與主相遇.
然後我們的身體.
在一秒鐘都不夠.
就會變化.
而那個變化.
就變成一個永遠不再扭壞.
是永遠是強壯.
而且是榮耀的一個復活的身體.

$^{121}$這個變成一個不扭壞.
不會死的事實.
這兩樣東西變成.
我們期待耶穌.
有一天祂一定會回來.
不過今天仍然有很多人取笑我們.
你們基督徒說耶穌回來.
說了多久.
超過二千年.
有些人說.
你們說而已.
耶穌哪裡回來.
我們說不是.
你看看我們今天的聖經裡面.
說到主在來這些預言.
超過百分之九十九.
已經是應驗了.
現在唯有說.
如果大家都不醒覺的話.
連神的選民.
以色列國.
1948年復國.
這個是不爭的事實.
但這個事實.
在歷史學家的研究裡面.
是一個很重要的關鍵.
因為你翻查.
世界所有的歷史.
列國的歷史.
沒有一個國家可以亡國超過五百年.
都可以翻身.
有沒有?.
你試試找找.
真的沒有.
一個也沒有.
但是以色列國亡國了多久?.
亡國了多久?.
亡國超過二千五百年.
到1948年才復國.
他們在公元前722年.

$^{161}$北角以色列十個支派被亞述國滅.
到公元前586年.
南角僅餘的兩個支派.
猶大和汴南文.
已經被巴比倫摧毀.
連聖殿都被毀.
什麼都好像沒有了.
國破家亡.
他們所有的這群神的選民.
都很悲傷.
神你去了哪裡?.
神好像是否丟棄了我們?.
我們今天的人生.
有很多人都可能會這樣想.
我們面對.
會衰殘.
有一天都會面對死亡.
死亡之後我們去哪裡?.
我們有時候回看我們的人生.
很多的際遇.
為什麼會這樣發生?.
為什麼有些悲劇會這樣立時呈現?.
我們就算信耶穌的人.
有時候我們的信心.
如果你想得不通達的話.
你也有時候會質疑.
上帝的信實.
或者上帝你去了哪裡?.
是不是神缺席呢?.
其實神永遠都沒有缺席.
神永遠在我們最憂傷.
就是落在一個最悲痛.
我們很困惑的境況裡面.
我們由舊約的歷史看到.
伯伯局善治怎樣向神呼喊.
我們看到耶利米善治怎樣流淚地.
去等候上帝應驗.
七十年期滿.
我們選民這群被擄的.
可以回到聖地.

$^{201}$我們看到但以利.
他在但以利書第九章.
他就看到耶利米書這樣記載.
他就向神迫切祈禱.
上帝讓我們真的可以回轉.
這一切莫非是因為什麼呢?.
就是他們相信神的話.
神的話不會失信.
人的承諾有時候我們是兌現不了的.
特別今天你有時候聽一些國家元首說話.
你會質疑而聽.
聽清楚.
有一個比較離譜的總統.
經常在上台前說要怎樣怎樣.
一上台就有九十多條行政指令.
這裡說那裡就改變了.
二十四小時不夠他就變了.
為什麼會這樣呢?.
因為人根本就是不可信.
我們今天信心的真正根基.
就在神的話裡面.
我們很感恩.
神讓我們可以看見這個事實.
所以剛才主席姐妹帶領我們讀的這段經文.
很有意思.
死亡我們每個人都要面對.
但是死是否可怕呢?.
原本是可怕.
如果我們沒有耶穌就真的可怕.
為什麼呢?.
因為死亡有一個毒刺.
就是我們的罪.
罪的權勢就是律法.
律法向我們每一個犯罪的人指控我們.
我們去到神的審判台前.
我們死定了.
所以聖經在羅馬書三章二十三節怎樣說呢?.
因為世人都犯了罪.
規缺了神的榮耀.
沒有一個人是例外.

$^{241}$所以今天包括在聖殿裡面.
我們神的子民.
我們沒有一個人是完美的.
但都是恩典.
所以我很喜歡你們教會的名字.
洪恩堂.
是因為基督這個偉大的救恩.
洪大的救贖功勞.
這個代價是主耶穌為我們賦上.
所以我們今天我們看到.
撒旦在神的審判台前想指控我們.
叫神的百姓好像沒有盼望.
神是透過耶穌基督得勝.
原因何在?.
原因是耶穌代替我們.
在十字架已經死了.
而且他的死和一般人的死.
是很大分別.
我們的死就是我們該死.
我們罪的功架就是死.
耶穌的死.
是他不應該死.
因為他一生沒有犯過半點的罪.
所以他的死是不正常的.
他是不應該要面對死亡的.
但是他為何要忍辱負重.
甚至被人取笑.
你說你是神的兒子.
在十字架下來了.
耶穌可以這樣做.
但是他選擇留在十字架.
就是要擔當我們世人的罪.
成為神的羔羊.
代屬我們每一個人所犯的罪.
罪債得以還清.
是因為基督為我們傾盡他的生命.
他愛世間屬於他的人.
是愛到底.
毫無保留.
毫無條件的.

$^{281}$去為我們贖罪.
這件事我相信我們一生.
都不可以做一個忘恩負義的人.
今天我們是被主耶穌的寶血所賣贖.
是因為這樣的緣故.
也是因為我們有這樣的期望.
所以我親愛的兄弟們.
你們務要堅固.
不可動搖.
你們知道死亡之後.
有一天耶穌要回來接我們回永恆的天家.
而這個榮耀的再來.
加上我們榮耀的復活.
這件事足以叫我們好好選擇.
在我們一個短暫的人生裡面.
今天你做什麼投資呢?.
你做什麼選擇.
去用你的時間.
用你的光陰.
我們要竭力去多做主公.
因為知道在主裡面的奴婦.
是不會徒然的.
因為將來在審判台前.
上帝已經知道我們每一個人.
在基督耶穌裡面所作的.
有很多事未必別人會知道.
甚至別人看不到.
但是天父記著.
所以你只要至死忠心.
你不是事奉為了要做在人的眼前.
當然我不是這樣說.
叫你眼前的事奉.
你就不敢去做.
好像我們唱敬拜隊的.
師班帶我們唱詩的.
在人的眼前.
在人的背後.
我們仍然忠心事奉.
有很多別人看到.
看不到.

$^{321}$但是主全部看到.
我們用真心誠意.
我們去好好事奉神.
這件事.
弟兄姊妹.
記住.
絕不徒然.
真的.
我今天嘗試和大家做一個總結.
血肉之體不能夠進入神的國.
既然是這樣.
我們期待這個不會扭壞的身體.
一個復活的盼望要出現.
在耶穌回來的時候.
死人要復活.
要帶來一個很清楚的事實.
就是我們一定要改變.
不死不扭壞.
不死不扭壞是進度永恆.
我們也無懼於死亡.
因為這是罪的結果和律法的指控.
我們也無懼於審判.
因為死後.
我們要面對永生神的審判.
原本我們就好像死囚.
死定了.
被定罪.
沒有盼望.
但因為耶穌的緣故.
祂的待贖.
我們今天的罪債已經被還清.
我在這裡和大家做兩個見證.
都是在這一兩年間發生的.
大家都知道我們經過三年的疫情.
在那時候.
其實剛剛在二,三年的時候.
疫情剛剛復常.
有一次我在開執事會的時候.
我收到一個訊息.
說可不可以趕去醫院.

$^{361}$為一個病危的病人祈禱.
這個病人其實是我認識的.
不過我不是和他交往很多次.
在疫情前.
原來張瑪麗姐妹是我們教會的.
她就將這位舞蹈者.
是她的老友記.
認識了很多年.
她就帶了來我們教會.
在一個福音主日的崇拜裡面.
有機會聽福音.
很好.
她聽完這個舞蹈的朋友.
都非常之心裡很歡喜.
她雖然那天沒有舉手決志.
但是她內心就已經預備好.
準備繼續會回來.
可惜跟著之後就有疫情.
疫情之後她就一直都沒有什麼機會回來.
但是張瑪麗姐妹就和我說.
她私下和她做祈禱.
和她決志了.
感恩.
她也送了一些詩歌帶.
見證帶給她繼續聽.
我那天收到這個訊息.
我就正在開執事會.
走不掉.
其實我跟著.
我就.
其實下午還有一個喪禮要去火葬.
我就唯有和那個訊息的發佈人說.
我必須要去完那個喪禮之後.
我盡快來醫院.
來完醫院之後.
我晚上還有一個聚會.
所以我整天.
我基本上只有那個空間.
原來神為我預備了那個空間.
我去到醫院.

$^{401}$我見到這一位臨終的病人.
原來他的徒弟都離開了病床.
這些徒弟全部都很孝順他.
當他是他的爸爸.
你知不知道這個是誰呢?.
這個師父是誰呢?.
原來他不是打功夫的.
他是美食的一個專家.
他是專做鮑魚的.
就是.
是呀.
楊師父.
他就是一鮑魚.
當時我自己和他的女兒.
還有我們有一個女執事.
是他們很熟的朋友.
我在他身邊.
他知道我們在這裡.
雖然他開不了眼.
但是他聽到.
我為他祈禱.
為他感恩.
也即時幫他做了洗禮.
在這個後來幫他做了安息禮拜.
他們家人得到很大的安慰.
更加奇妙的是.
我在這裡遇上了一個以前的街坊.
當時我還是小朋友.
我現在是認不出他.
但原來這個街坊.
是一個很出名的藝術家.
畫畫的 畫油畫的.
和我們國家領導人畫的人像.
那時候我才知道.
這麼厲害.
他在那裡和我相認.
他說林牧師你好.
我說你好.
請問你是哪位?.
原來是我媽媽幾十年前.

$^{441}$是帶他去掘字的.
你說兜兜轉轉.
原來我們今天在WhatsApp.
做了好朋友.
所以我看見神的國很奇妙.
無論我們住在哪個城市.
哪個境域.
但我們彼此都是.
神的國 基督的徵兵.
我們都有同一個使命.
就是要大人信主.
原來這位藝術家.
也是楊師傅多年的朋友.
一直都勸他要信主.
另一個見證是甚麼呢?.
都是在二,三年的時候發生的.
有一天我收到一個電話.
有人跟我說.
牧師你可不可以去醫院.
幫一個臨終的病人.
他說他癌症末期.
醫生說不能逆轉.
他太太是教徒.
原來他太太多年參加.
Bible Study Fellowship.
大家聽過這個BSF.
很多年研經.
他的子女都信主.
但唯獨這個先生還未信.
但這個先生一直不會阻礙他們回教會.
但陪他們回到教會門口.
他就走了.
載他們去教會.
他就走了.
當我接到這個邀請的時候.
當時我正在去醫院.
我只知道名字是誰.
我嘗試上Google搜尋一下.
這位是甚麼人.
我一看的時候.

$^{481}$原來這個人很厲害.
他是國際知名的環境保護專家.
他做了很多的研究.
為了氣候.
在聯合國發表了一些文章.
得到很多國家認同.
一起簽署協議.
減少碳排放.
他是一個很出色的科學家.
但後來他為了爸爸的緣故.
他回到香港.
他父親是做珠寶生意.
爸爸去世後.
沒有人接手.
他一定要頂上.
所以他回到香港.
已經有一段時間.
我去到醫院.
我見到躺在病床的人.
他當時是很清醒.
如果不告訴你.
他是末期.
你以為他出院了.
原來不是.
他真是末期.
當時我跟他說.
我去到床邊.
我肯定他一生所作.
我自己是科學出身.
我是工程師出身.
我很欣賞.
你為我們整個地球.
所付出的努力和貢獻.
這個大地都是神的創造.
你已經盡了你的努力和責任.
為了我們的地球作出貢獻.
我肯定他所做的一件事.
但我亦告訴他.
我們人生最重要的一件事.
不可以缺失的.

$^{521}$就是你一定要有主耶穌基督.
你就算在人面前.
是一個多好的丈夫.
好的父親.
好的人.
這件事不能夠叫我們去天國.
你需要是耶穌基督的恩典.
這個寶貴十字架的救贖恩典.
我講到最後的時候.
他的眼睛眼泛淚光.
然後我勸他.
在他太太和兒女的側巾.
就和我一起做這個決志祈禱.
一做了這個祈禱之後.
我就為他做了施洗.
解釋了洗禮的意義.
問他願不願意.
他很願意.
結果.
我將這兩個見證告訴大家.
我很想大家知道.
我們今天無論遇到任何人.
都一樣.
我自己一樣.
我是接觸很多基層的人.
我們將會在6月1日.
如果上帝許可的話.
我們就在香港仔的牌坊後面.
我們找到一個地方.
有二千多呎.
我們在那裡租了下來.
我們開展一個廣福關愛中心.
就是接觸基層的家庭.
幫助那些SEN的小孩子.
是因為基督的愛激勵了我們.
我們必須去還福音的債.
我們是去關心別人.
所以我們講傳福音.
不只是說教.
傳福音是要用我們身體力行去活出來.

$^{561}$給我們有很多這方面的恩典.
我們今天需要與神與人.
我們去同心.
我們去履行這個大使命.
而教會需要得著福氣.
彼此相愛.
我們傳承的信仰.
是要給下一代.
今天我是很多不同教會.
有時接觸港獨的時候.
我心裡頭.
都是有一個負擔.
我就要鼓勵.
其實包括我們自己的廣福堂一樣.
我已經五,六年前.
跟自己的教會說.
我們必須要防止教會老化.
我們要年青化.
我這樣說是有原因的.
我這樣說不是說我們年長的弟兄姊妹.
你就站在一邊.
你不用事奉了.
不是.
我們要跨代同行.
所以去年我們是四十週年.
我用這個做我們教會主題.
跨代同行.
傳揚主恩.
但今天有時我們越愛教會.
我們很多時候.
我自己的舞會是九龍城朝的浸信會.
我到現在跟他們還有結連.
我們有一班叫兒童師班的老友.
我以前是小朋友.
一直唱師班.
為何我唱到現在.
我回到自己的舞會.
原來你們還在唱.
他們有些都出來做了指揮.
教下一代兒童師班.

$^{601}$不過他們那種忠心和擁抱.
我說不要緊.
你越愛那個服侍.
沒有問題的.
但同時間我們要提攜年青人.
新生代.
是你了.
你們要起來了.
家長和我們做牧者也要留意.
原來手沖咖啡.
現在年青人很喜歡.
我們大人不要只沖.
慢慢讓他們一起沖.
唱歌也是.
可以跨代同行去唱.
所以今年我們教會有一個轉變.
師班邀請一些年青人.
和他們一起唱.
邀請年青人進我們的師班.
他們不敢來.
他說那些全部都是父母.
他們不敢來.
不敢來我們就轉變方式.
我們主動出擊.
邀請他們.
我們小組一起上台唱.
加上一些孩子.
一些八歲到十二歲的小孩子.
原來他們彈小提琴.
奢羅.
吹那些旋律.
我們雖然搞不成一個交響樂團.
不需要.
但他們有這樣的恩賜.
我們給他們機會.
鼓勵他們.
全家出來事奉.
那些家長都不知道多興奮.
神的家是需要這一種.
我們一起得著鼓勵.

$^{641}$是整個家庭一起去事奉神.
我們要將信仰傳承給下一代.
教會在社區裡.
要樹立這種金燈台的見證.
愛神愛國愛家愛人.
我看見你們.
我為你們很感恩.
我由我第一刻踏入你們的門口.
我已經真的.
剛才我說過.
你們的共享空間做得很好.
做得非常好.
這個就讓路過的街坊都可以進來.
領人歸主.
客不容緩.
因為沒有人知道.
人的壽數幾何.
在各個專業的領域.
我們都用生命.
毀身的職場見證.
貢獻社會和我們的國家.
我在疫情的時候.
我給這幾張相片大家看看.
我和師母.
我們是帶著夫婦區.
整個團契.
我們去到這個.
在屯門.
一個屋苑裡面.
我們和屯門的播道會.
一起有一個合作.
就是洗樓.
幫那些社區.
我們將那些物資.
將那些口罩.
我們帶去送給那些街坊.
這些是約好了.
是星期日下午.
我們見證完了之後.
我們就去.

$^{681}$左手邊上面.
左手邊那個.
大家看看.
這個就是一個老人家.
另外我們看.
大家看不看到.
看不到.
不知道為什麼照不到.
這幾個是同一個人.
這是第二個家庭.
第一個家庭.
是一個婆婆.
我們也有和她傳福音.
不過她完了之後.
她很客氣.
我們就趕著去樓下.
探望另外一個街坊.
但這個街坊很特別.
粉紅色衣服的那位老人家.
他是星期日.
原來那天他要上班.
他專登告假等我們來.
大家猜猜他幾歲.
八十幾九十歲.
他七十五歲.
但他是做清潔工.
他是政府僱用的清潔工.
但他專登那天知道我們要來.
他告假等我們.
然後我們去到他的樓層.
都不用去到他門口.
他就在樓梯等我們來探望他.
他已經很準備好.
結果我和他聊了兩句.
我聽到他的口音.
我說你是不是潮州人?.
他說是的.
我說你家鄉在哪裡?.
原來他是潮洋.
我說我也是潮洋.

$^{721}$加傑林.
然後我和他用潮州話.
就含冤了幾句.
我就問你什麼時候來香港?.
原來他十八歲就來香港.
嫁給了人.
他老公也是潮州人.
他們全家.
原來他這麼多年.
都已經將他的生命奉獻給了他家.
無悔無憾為了兒女.
為了丈夫.
他用他的生命.
就是這樣為家庭貢獻.
不過因為經濟緣故.
到現在七十五歲.
都未可以退休.
都是為了家庭.
我當然很感動.
然後我說.
你有沒有接觸過基督教?.
他說他十八歲的時候.
在潮州去過教堂.
進去教堂聽過道.
幾十年前.
我聽完之後.
我自己說.
上帝今天一定是差遣了我們這對天使.
要和你再傳福音.
然後我很簡單和他再講福音一次.
然後講完之後.
我說你願不願意信耶穌?.
他原來已經預備好了.
他說願意.
我們一起幫他決志.
旁邊最左邊的姐妹.
就是在屯門的.
我們播道會的教會的姐妹.
他就帶她回教會.
我們有跟進的.

$^{761}$我們很感恩.
我們教會願意和其他社區.
其他的教會一起合作.
結出那紙.
因為天國的緣故.
不一定要回到自己的教會.
你就回到你自己社區的教會.
我們就是這樣彼此服事.
彼此祝福.
很感恩.
這裡可能字比較小.
其實我很喜歡這張海報.
這張海報.
他說時間是自由的.
但它是無價的.
你不可以擁有它.
但你可以使用它.
你不能夠保存它.
但你可以花掉它.
一旦你失去.
你就永遠都拿不回來.
真的.
我們今天在主裡面.
要把握這些機會.
這些時間.
所以今天你要明白.
在主裡面的勞苦.
是永不徒然.
但你怎樣去用你有限的人生呢?.
我們一起讀出來.
這一節聖經.
已不鎖書.
五章十五到十八節.
你們要謹慎行事.
不要像無知的人.
要像智慧的人.
要把握時機.
因為現今的世代邪惡.
不要作糊塗人.
要明白主的旨意如何.

$^{801}$不要醉酒.
酒能使人放蕩.
要被聖靈充滿.
在這裡保羅說了.
我們要在這個邪惡的世代.
如果說是二千年前的爾忽所的城市.
為何說它是邪惡呢?.
因為他們滿個城市都是偶像.
人的心是被迷惑.
所以他說我們要小心.
其實每一個時代都有這些危機.
我們要謹慎行事.
不要做無知的人.
不要做糊塗人.
我們必須要明白神的心意的話.
要做智慧者.
必須要把握時機.
原本的黃河本說要愛惜剛蔭.
其實原文希臘文有一個更好的意思.
愛惜剛蔭.
剛蔭是甚麼呢?.
剛蔭不只是流水作業.
原來他在說的是一個奇異物.
一個機遇.
但機遇是要把握的.
機遇不把握的時候.
就會輕輕被它流失.
我在這裡想和大家講一個見證.
你的人生裡有沒有一些你很清楚.
是神給你的機遇.
錯過了回不去.
你試過有一些遺憾.
明明是同學來的.
大學一起讀書幾年.
我以為還有很多時間可以和他傳福音.
怎知走了.
有沒有試過類似這樣的經驗.
或是你的同事或親友.
如果有的話請你記住.
你今天要反思.

$^{841}$原來你現在還有神給你的機遇.
很多時候我們人生不是你自己去.
你填好日程表.
特別我們做牧師傳道.
我經常和神學生和自己的同工說.
你不要以為這個星期.
你已經好好計劃好你的行程.
你預備好星期天要講的講章.
你以為安枕無憂.
忽然間突然收到電話.
忽然間一個WhatsApp來.
原來我們做神的僕人.
就算突然姐妹一樣.
你可能隨時都有一些突發的需要.
有一些事情.
我們怎樣去回應呢.
這就是神給我們的Divine Appointment.
是一個神聖的約會.
我昨天在自己的總會開了一整個早上.
然後我要趕回下午要去教會又要見人.
我打算在茶餐廳吃個快餐.
怎知整個餐廳很滿人.
原來很多國內同胞來做.
因為那個地方在尖沙咀.
結果我走進去餐廳.
他說沒有位子了.
要搭桌子.
我說好啊 搭桌子也好.
我坐下去.
真是Divine Appointment.
原來茶餐廳搭桌子.
很好傳福音.
我一坐下去.
在我前面.
一個很年輕的太太.
原來兒子十五歲了.
我是用普通話跟他們交談.
起初他又不知道我是牧師.
我說歡迎你們來到香港.
你們覺得這裡怎麼樣.

$^{881}$很好吃啊.
他要吃些菠蘿油.
那些最邪惡的.
又要喝奶茶.
我跟他們行圈幾句之後.
原來跟他們講福音.
原來他們從來沒有機會接觸過.
這些他們自小無神論主義.
一直就是這樣.
打開福音畫鴿子.
然後我就不是拿刀出來.
拿些卡片和福音邀請卡.
我告訴他們.
原來有兩個星期.
媽媽陪兒子來放假.
我說對了.
這兩個星期你來我們.
我們有普通話崇拜.
你來中環.
你順便來我們教會.
中午十二點你就來.
繼續為他們祈禱.
求主也是這樣開他們的心.
他們因為早坐下.
他們吃完.
但我還沒喝完.
我的那杯凍奶茶.
所以他們走了.
拜拜之後.
有兩個年輕的女孩.
也是國內人.
坐過來對著我.
她又不知道我是牧師.
我就跟她們.
由普通話.
「歡迎你們」.
我好像做了香港的接待大師.
我們很開心.
我們聊天.
原來今天你的人生.

$^{921}$是應該這樣.
你不只是福音主人.
你也不只是在教會裡.
你才可以跟別人講福音.
你在茶餐廳.
就可以跟別人講福音.
所以她說著說著.
後來這兩個女孩.
這位人士也很特別.
後來我給她一張卡片.
她說你是念牧師.
牧師跟神父有什麼分別?.
她說很好.
這些問題就是好.
結果我又邀請她來普通話崇拜.
我講一些見證給大家知道.
我1922年做了通波仔手術.
是神給我一條命.
我當時真的很感恩.
長話短說.
因為時間不夠.
我在1922年9月.
我做了通波仔.
我起初醫生說要盡快做.
我說醫生我已經訂了機票.
我要去加拿大探望一些親人.
和答應了幾個教會港督.
等我先回來.
她馬上滿面愁容.
你先回去考慮一下.
她叫我回去考慮.
她又不可以逼我.
結果我原來回去也不能考慮.
第二天我就來了曾牧師.
那時你記得去了錦繡堂.
就是那天去錦繡堂港督.
講完堂.
曾牧師很客氣.
請我和師母一起去喝茶.
多找幾個弟兄姊妹.

$^{961}$我不知道你記不記得那次.
那次原來是有一個故事.
原來大家都脫了口罩.
因為那時是1922年.
還是疫情.
一脫了口罩.
吃完那頓飯就再見了.
有個弟兄是我自己教會.
Mary Cheung的老公.
載我回去香港.
我在車上已經開始不舒服.
到那天晚上我累得不得了.
我人生中未試過這麼累.
我和師母說我要早點休息.
那時是九點多.
師母說是不是你今天出去港督累了.
我說沒有理由.
我平時在自己教會講四堂.
我都是這樣.
我不知道為什麼.
不過我想我睡完覺.
希望明天早上好一點.
怎知道明天早上累得起不了身.
整個人都軟了.
然後去看醫生.
醫生說你九成是COVID.
給我一個小瓶去驗.
一驗結果是.
原來神出手.
神要我乖乖地.
我差點就捉了去高灣.
幸好家裡有多個洗手間.
所以不用去隔離.
就在家裡隔離.
結果在那七天裡.
神很清楚跟我說.
你留在香港.
哪裡都不用去.
教會都去不了.
然後神跟我說.

$^{1001}$神那我留在香港做什麼.
還可以做什麼.
你說你可以做什麼.
做手術啊.
做完手術.
醫生跟師母跟我說.
他說你三條血管都要放支架.
已經塞了超過70\%.
他說其中一條已經有機會爆.
如果爆裂的時候.
我去加拿大.
在機上下了機.
出現這個情況.
我應該都不能站在這裡.
我今天已經不在這裡.
其實我回看22年9月.
原來神給我.
生命可以延續.
是神的恩典.
所以今天我一息尚存.
我仍然想跟大家說.
你要把握你的時機.
有另一個見證.
我講不了.
沒有時間了.
但是有很多見證.
我很想告訴你.
原來我們在神的國裡.
你願意在主裡.
你所做一切的勞苦.
絕不徒然.
而且神在他將來.
永恆的角度裡.
你要得你的賞賜.
求主幫助我們把握時機.
遵守神的旨意.
而且對準彪剛.
竭力奔跑.
我們一起讀這節聖經.
弟兄們.

$^{1041}$我不是以為已經得著了.
我只有一件事.
就是忘記背後.
努力面前的.
向著彪剛直跑.
要得神在基督耶穌裡.
從上面.
召我來得的獎賞.
弟兄姊妹.
不要將來有任何的遺憾.
你就要把握機遇.
忘記背後.
努力面前.
向著彪剛.
永恆的目的地.
我們進發.
我今天沒有機會.
講這個見證給你聽.
不過這是一個四十歲.
不幸肺癌的去世.
新加坡一個億萬富豪.
這個張慶祥.
二十多歲他去過教會.
在當地.
但他那時不可一世.
他以為自己做了美容的醫生.
賺了很多錢.
但他想不到.
三十九歲的時候.
他就患了肺癌.
一患了的時候.
已經是4B.
原來是最近的那種肺癌.
他後來就很後悔.
他在神的面前懺悔.
我給最後這個圖片.
就結束了.
這個圖片是我們教會.
一個長者團契.
在疫情的時候.

$^{1081}$這班長者都願意出去.
派物資給有需要的街坊.
我想大家看.
這個圖畫裡面.
是唯一一個男生.
這個男生.
就是這位.
見不見.
你猜他幾歲?.
是九十歲.
你猜對.
九十歲.
他在疫情前很有氣.
吹薩斯風.
很好聽.
但他當時.
他其實那個身體.
都是比較虛弱了.
但是他烏底的身.
他都願意出去.
派街坊.
幫助有需要的人.
我經常用這幾張圖片.
我去挑戰.
鼓勵教會.
我說:靖姐妹.
如果一個九十歲的長者.
都可以出去服侍社區的人.
我們今天如果健健康康.
我們走得走得.
無論你是哪個年紀.
我們有什麼藉口.
我們要挽手.
我們剝花生.
在這裡只是看別人.
去服侍.
求主給我們就有這個看見.
是給主的愛.
去激勵我們每一個.
我們真是很好的.

$^{1121}$絕不徒然的.
我們在基督耶穌裡面.
每一個你所花的.
每一個的時間.
時刻和金錢.
原來做在主的身上.
你給一杯涼水.
一個小紙.
都是做在主的身上.
這些經文提醒我們.
你不要錯過.
今天你搭地鐵.
或者好像我剛才說的.
在茶餐廳.
三不識七.
我簡直覺得.
昨天是上帝為我預備的.
其實就好像福音茶座.
原來都不用叫人來.
就坐下來吃飯.
是自動送上門的.
大家就.
祂又不能走.
叫了食物.
這樣就多好.
下次原來去喝茶.
看有沒有搭台.
都好的.
我想說這些例子.
不要怕去接觸陌生的人.
原來這一切都是神給我們的恩典.
我們一起祈禱.
我們的人生雖然是短暫.
但在永恆裡面.
我們知道我們去哪裡.
主耶穌給我們這個.
復活永存的盼望.
將來他們在天家.
我們要晝夜和二十四位長老.
我們不斷去敬拜.

$^{1161}$頌讚神.
何等的榮耀.
何等美麗的一幅圖畫.
我們在那裡再沒有任何的.
真的疾病.
疼痛.
哀哭.
來纏繞著我們.
我們得到完全的自由釋放.
因為今生我們可以無悔.
無憾地去活一次.
是為著主.
你呼召我們.
這個冥定.
我們一起願意獻祭給主.
原來在我們頌鼎姐妹的生命裡.
成就你永恆計劃的美意.
主你賜福曾牧師.
和整個團隊教會的領袖.
一起去服侍這個社區.
祝福更多更多的人.
我們看到將來北部大都會.
在元朗區成為一個很大的機遇.
成為一個福音的大和場.
求主感動更多更多弟兄姊妹.
被激勵.
一起我們還福音的債.
我們一起被建立成為主喜悅的門徒.
同心合二禱告.
奉主耶穌基督的聖名求.
阿們.
(影片完結).
\newpage



\section{哥林多後書 4:16-5:10}
\label{sec:gznE_6kHeYc}
\textbf{2025.04.06《活潑的盼望》 哥林多後書4\_16-5\_10 講員 : 曾敏姑娘}
\newline
\newline
連結: \href{https://youtube.com/watch?v=gznE-6kHeYc}{\texttt{https://youtube.com/watch?v=gznE-6kHeYc}} ~~~~ 語音日期: 2025-04-06
\newline
\newline
\hyperref[sec:wH4miLAekc8]{\small{< < < PREV SERMON < < <}}
~
\hyperref[sec:index_chronic]{\small{[返順時目]}}
~
\hyperref[sec:index_scriptual]{\small{[返順卷目]}}
~
\hyperref[sec:AR6jAGDUTDM]{\small{> > > NEXT SERMON > > >}}
\newline
\newline
哥林多後書 4:16-5:10
\newline
\begin{longtable}{cl}
\hline
\hline
章節 & 經文 (和合本修訂版)\\
\hline
4:16 & \begin{tabularx}{0.7\textwidth}{X} 所以，我們不喪膽。雖然我們外在的人日漸朽壞，內在的人卻日日更新。 \end{tabularx} \\ \\ \relax
4:17 & \begin{tabularx}{0.7\textwidth}{X} 我們這短暫而輕微的苦楚要為我們成就極重、無比、永遠的榮耀。 \end{tabularx} \\ \\ \relax
4:18 & \begin{tabularx}{0.7\textwidth}{X} 因為我們不是顧念看得見的，而是顧念看不見的；原來看得見的是暫時的，看不見的才是永遠的。 \end{tabularx} \\ \\ \relax
5:1 & \begin{tabularx}{0.7\textwidth}{X} 因為我們知道，我們這地上的帳篷若拆毀了，我們將有神所造的居所，不是人手所造的，而是在天上永存的。 \end{tabularx} \\ \\ \relax
5:2 & \begin{tabularx}{0.7\textwidth}{X} 我們在這帳篷裡嘆息，渴望得到那從天上來的居所，好像穿上衣服； \end{tabularx} \\ \\ \relax
5:3 & \begin{tabularx}{0.7\textwidth}{X} 倘若脫下也不至於赤身了。 \end{tabularx} \\ \\ \relax
5:4 & \begin{tabularx}{0.7\textwidth}{X} 其實，我們在這帳篷裡的人勞苦嘆息，並不是願意脫下地上的帳篷，而是願意穿上天上的居所，好使這必死的被生命吞滅了。 \end{tabularx} \\ \\ \relax
5:5 & \begin{tabularx}{0.7\textwidth}{X} 那為我們安排這事的是神，他賜給我們聖靈作憑據。 \end{tabularx} \\ \\ \relax
5:6 & \begin{tabularx}{0.7\textwidth}{X} 所以，我們總是勇敢的，並且知道，只要我們住在這身體內就是離開了主。 \end{tabularx} \\ \\ \relax
5:7 & \begin{tabularx}{0.7\textwidth}{X} 因為我們行事為人是憑著信心，不是憑著眼見。 \end{tabularx} \\ \\ \relax
5:8 & \begin{tabularx}{0.7\textwidth}{X} 我們勇敢，更情願離開身體，與主同住。 \end{tabularx} \\ \\ \relax
5:9 & \begin{tabularx}{0.7\textwidth}{X} 所以，無論是住在身內或住在身外，我們都立了志向要得主的喜悅。 \end{tabularx} \\ \\ \relax
5:10 & \begin{tabularx}{0.7\textwidth}{X} 因為我們眾人必須站在基督審判臺前受審，為使各人按著本身所行的，或善或惡受報。 \end{tabularx} \\ \\
[1ex]
\hline
\hline
\end{longtable}
$^{1}$願至聖聖三一保佑.
能有福氣回到神的家裡.
去敬拜祂.
我們開始的時候一同禱告.
此下你的話語.
要成為我們腳前的燈.
路上的光.
我們生命裡總是會面對很多挑戰.
但你總是會透過你的話.
去鼓勵和帶領我們.
今天求你讓我們每一個人.
都有一個很開放的心靈.
願以聖靈親自對我們每一個說話.
讓我們聽到心裡的去.
也去回應神的教導.
求你賜福以下的時間.
奉耶穌基督的名字祈禱.
阿們.
最近有很多事情發生.
世界的.
香港的.
我們自己的.
在過程中.
可能我們心裡也會覺得很亂.
很多很多的思考.
也有很多的討論.
是不是?.
在這裡我們真的要問一個問題.
我們的盼望在哪裡?.
每個星期日我回來聽到.
是不是一個很分割的事情?.
星期日聽到.
回來敬拜神是一件事.
我一回到我的生活.
我就是面對林林總總的.
數不盡的挑戰.
我也不知道怎麼辦.
是不是分割的呢?.
一定不是分割的.
是很有關係的.

$^{41}$所以很盼望今天.
這一堂道可以鼓勵我們的心.
我們不單有盼望.
還是很活潑的一個盼望.
在這裡.
剛才說到.
現在最近我們的世界很亂.
生命中有很多事情.
讓我們感到懼怕.
例如.
討論得相當多的.
緬甸的地震.
非常多人離開了世界.
我們其實都不知道.
有多少離開世界的人.
在當中已經相信了耶穌.
如果沒有相信耶穌.
他們之後去了哪裡呢?.
還有.
之後可能很大機會是.
爆發疫情.
緬甸的地震又影響到泰國.
最近有很多的討論.
人們一直討論.
可能台灣什麼時候地震.
日本什麼時候地震.
所以旅行要趁著多久之前.
就要快點去.
不停地說地震.
這裡地震那裡地震.
很多的恐懼.
接著.
有關稅戰.
接著有很多討論.
經濟會受到影響.
我們香港怎麼樣呢?.
因為我們真的很關心我們的生活.
我們的工作怎麼樣.
我們的收入怎麼樣.
還有.

$^{81}$現在不只是成年人會擔心經濟.
學生也會擔心.
我現在生活最基本的.
我的需要.
供應在哪裡.
我也有不足夠的地方.
學生有學生的擔心.
成年人有成年人的擔心.
有今天的擔心.
有將來的擔心.
將來我不夠錢又怎麼樣呢?.
誰來照顧我呢?.
我們比較成熟的弟兄姊妹.
是佔多數.
我們越來越成熟.
越來越年長的時候.
我們的健康.
可能越來越多出現問題.
吃藥也越來越多.
我們不只是吃藥.
我們擔心我們的健康怎麼樣.
我發覺弟兄姊妹聚在一起.
討論很多健康的事情.
吃什麼.
怎麼吃.
什麼時候.
做什麼.
曬太陽.
怎麼樣健康.
不停討論這些問題.
聚在一起就說.
因為我們真的很掛心.
很擔心的事情.
就會在口裡經常說.
我們除了健康不好.
還有擔心.
誰來照顧我們呢?.
我們年紀大了.
如果子女在身邊.
就說子女照顧.

$^{121}$還要子女孝順.
如果不在身邊.
我們又怎麼辦呢?.
這些很實際.
還有.
我們的人際關係.
婚姻關係.
關上道門.
兩個人所經歷的事情.
別人都不知道.
裡面有很多難處.
都不可以跟別人說.
我們人際關係.
同事關係.
弟兄姊妹關係.
那些困難.
好像經常解決不了.
一個是非.
完結後又一個是非.
一說到是非.
我們有沒有這個懼怕?.
我們每一個人.
有沒有這個懼怕?.
別人會不會在背後說我們呢?.
他這個是不是在說我?.
是不是?.
他在說我是非.
或者他會怎麼看我?.
他們在討論的.
是不是在討論我呢?.
我們都很懼怕的.
很多這些.
天天前前後後.
不停地糾纏.
你會發覺.
人真的很難.
好像很不容易有平安.
今天有一首詩歌說平安.
你會發覺我們的平安.
常常都被影響了.

$^{161}$坐在這裡崇拜.
可能東想西想.
這件事很不開心.
那件事也很不開心.
那怎麼樣呢?.
有沒有盼望呢?.
信耶穌的人.
我們是不是和外面的人.
是一樣的.
沒有盼望呢?.
這個值得我們去問一問自己.
這個問題.
在這裡保羅自己.
做了一個很美的.
一個生命的見證.
非常好的見證.
我們聽我讀就行了.
我們看回哥林多後書.
四章.
其實還有其他的經文.
說到保羅受苦.
說保羅受苦的經文.
也不只是一處.
他受的苦是非常多.
我敢說.
我們在座的.
沒有人受的苦.
是比保羅多的.
沒有.
他所受的苦.
有肉體的苦.
有心靈的.
各方面的.
我們沒有人敢說自己.
比他受的苦多.
但是你看保羅受那麼多苦.
生命是有一個很活潑的盼望.
在哥林多後書四章這樣說.
他說我們處處受困.
卻不被困住.

$^{201}$內心困擾卻沒有絕望.
遭受迫害卻不被撇棄.
擊倒在地卻不自滅亡.
我們身上尚帶著耶穌的死.
使耶穌的生也在我們身上顯明.
因為我們活著的人.
尚為耶穌被置於死地.
這個是最盡的了.
為福音的緣故.
為福事的緣故.
被置於死地.
是最盡的了.
但是卻在裡面.
多有盼望.
使耶穌的生命.
在我們這必死的人身上顯明出來.
是有生命的.
不被打倒的生命.
所以保羅才可以這樣說.
叫我們什麼?.
上上喜樂.
那些笑容不是擠出來的.
是上上喜樂.
坐牢也喜樂的.
被人逼迫也喜樂的.
保羅自己是這樣經歷的.
他自己是這樣經歷的.
他就可以在裡面.
去作這個見證.
而且他會教我們.
生命為什麼有盼望.
我們有盼望是因為什麼?.
弟兄姊妹.
我想說這個道理.
一點也不高深.
我們有盼望是因為.
我們有寶貴的救恩.
信耶穌.
得夢拯救之後.
你馬上看事情.

$^{241}$你的眼光不應該再是一樣.
不單是眼光不一樣.
你整個生命的經歷.
事實上都是不一樣的.
就因為這份救恩.
我和你都有這份救恩.
我們不是讀很高深的神學.
一路考究.
逐個字逐個例研究.
搞了一大輪.
寫了很多論文.
才知道.
這個這麼高深的理論就是這樣.
不是的.
就這麼簡單.
我和你信耶穌得救.
就有盼望.
在這本經文.
我們逐步看下去.
做了一些不同的比較.
這個盼望在哪裡看得出來呢?.
四章十六節.
那裡這樣說.
所以我們不喪膽.
很多困難.
很多的沮喪.
很多的迫害.
但是我們不喪膽.
雖然我們外在的人日漸老懷.
內在的人卻日日更新.
這裡有個比較.
外在的人和內在的人.
外在的人在說什麼?.
我的身體.
我的身體在衰殘.
如果大家有機會看回.
我們年輕時的照片.
大家會不會有些感慨.
我年輕時是這樣的.
我試過有一次.

$^{281}$分享給別人聽.
我那時候讀書的時候.
有這些經歷.
之後是這樣.
別人就說.
那個是你嗎?.
都不像你現在.
我感受到這個.
外在的人日漸老懷的真理.
一直我們的身體在衰殘.
以前跑跑跳跳.
跳跳跳.
現在不行了.
走路走久一點.
都喘口氣.
感到很辛苦.
但是記住.
這個身體的衰殘.
內在的人卻日日更新.
內在的是指什麼?.
不是說另一個人.
是說我們的情感.
我們的意識.
我們對上帝的認識.
我們對上帝的回應.
很豐富的.
內在的這個人.
我們心裡的力量.
卻是日日更新.
我們是不受外在身體的衰殘.
內在仍然有力量.
我們對神的認識.
是越來越多的.
很有力量.
所以大家才可以聽到.
在不同的報道會當中.
聽到不同的弟兄姊妹講見證.
他身體可能經歷很多疾病.
很多患難,艱難.
但內在他是活潑的.

$^{321}$有力量的.
他經歷了神.
他會跟你說.
他如何有魄謀還可以服侍神.
還可以做多少不同的事.
為神作見證.
裡面是另一個世界.
不是自欺欺人.
不是自我告訴自己.
我一定可以的.
我們真的可以的.
真的可以的.
就是這樣日日更新.
還有.
保羅做了一個比較.
就是現在這一生所經歷的.
和永恆的事情做一個比較.
我們現在這個短暫而輕微的苦楚.
今生的.
很短暫.
幾十年就過去了.
那個苦楚.
在保羅看來.
都是很輕微.
保羅經歷的苦楚就不輕微.
被石頭打.
經歷海難.
差不多死.
又很飢餓.
又睡不著.
那些苦楚怎會輕微呢?.
不輕微.
但在他看來.
很短暫又輕微.
你知道這一切.
要為我們成就達至什麼嗎?.
你看看將來.
永遠的那一件事.
價值多重量.
是極重的.

$^{361}$是無比的.
榮耀是永遠的.
現在算不得什麼.
哇!弟兄姊妹.
信耶穌之後.
有救恩在我們生命裡面.
原來我們就應該這樣看.
是不同的.
不是看得.
現在的苦楚.
真的難當.
現在的日子.
真難捱餓.
這麼難受.
比死難受.
很多時候.
我們實際的感受是這樣.
但保羅這裡說.
不是的.
短暫的.
很輕微的.
你想想將來.
極重無比.
永遠的.
是榮耀.
今天是不是很有力量?.
我面對苦難.
很有力量.
為什麼?.
我們不是顧念看得見.
顧念的意思.
就是專注度.
我不是經常看著.
那些看得見的.
哇!看得見.
我的身體衰殘.
年紀越來越大.
不是看得見.
現在電視上.
有很多很多東西.

$^{401}$只是看著這些東西.
而是顧念.
專注看不見.
看不見永恆的生命.
永生.
看不見的是將來.
上帝賜給我們.
很榮耀的.
一個新的身體.
看不見很多.
更加好的東西.
人的眼目.
不是就是這樣.
經常看著地面.
我很喜歡用一個例子.
就是雞和鷹的分別.
雞和鷹.
雞和鷹的分別.
看著地面的東西.
不斷地.
永遠看著地面.
鷹可以飛得很高.
哇!真是.
看著整個大地萬物.
我們就是應該.
有鷹一般的眼光.
保羅教我們就是這樣.
你不要只是.
看著現在看得見.
他要看著看不見的.
原來看得見的是.
暫時的.
你一個衰殘的身體.
你的日漸衰殘的生命.
很短暫.
我們的困難都很短暫.
地面上的苦難都很短暫.
但是看不見的是永遠的.
上帝賜給我們.
數不盡的祝福.

$^{441}$永生.
一切榮耀的新身體.
一切所有的.
那些是永遠.
永遠的先有價值.
昨天我們青少年區.
有一個活動.
在感受進行.
我們去思想.
耶穌基督的受苦.
思想死亡.
對我們的意義是什麼.
其中有一個.
少年人的分享.
我覺得是非常好.
就是說.
因為之後.
我都有再問他們.
他們拿出來的物資.
就是他們身上的財產.
當時身上所有的東西.
包括手機.
身份證.
所有東西.
我說要將他們奉獻給教會.
問他們願不願意.
其中有一個少年人分享.
因為他在體驗死亡.
他說我一想.
其實我都已經死了.
是不是這一切的財產.
有什麼這麼重要呢.
他就願意獻出來.
是的.
你正在經歷死亡的時候.
你就知道.
其實那些東西.
地面上的東西.
真的沒有這麼重要.
價值不是真的這麼大.

$^{481}$你很自然就放手了.
你這些地面上的東西.
你不會進入到永生裡面.
它不會有永恆價值.
所以現在放你.
你都不會覺得可惜.
就是這樣.
所以現在看得見.
暫時看不見的.
才是永遠的.
在這裡.
接著保羅繼續做一些對比.
他說因為我們知道.
我們這個去到第五章了.
我們這個地上的帳篷.
如果拆毀了.
我們將有神所做的居所.
不是人手所做的.
而是在天上永存的.
這個帳篷是指什麼呢.
也是指我們地面上.
我們的身體.
身體是衰殘了.
我們生命離開的時候.
我們不要擔心.
其實是更好.
有神所做的居所.
不是人手所做的.
而是在天上永存的.
是上帝賜給我們.
一個榮耀的新身體.
是直到永遠.
那個就不會再有疾病.
不會再經歷痛苦.
現在我們地面上的身體.
就算怎樣也會衰殘.
總有一天.
就要放下它.
你要想下一個對比.
地上的帳篷.

$^{521}$和天上永存的.
差多遠.
那個是神所做的.
我知道有一天.
我會有一個新的榮耀的身體.
我們在這個帳篷裡.
嘆息渴望得到.
那從天上來的居所.
好像穿上衣服一樣.
這個嘆息.
這個位置不是說.
哎呀 真是絕望了.
我真是不開心了.
不是說這裡.
是說我們在這個帳篷裡.
我很渴望.
我很渴望得到.
從天上來的居所.
我希望有這個榮耀的新身體.
第三節.
倘若脫下也不至於赤身了.
比較好的翻譯應該這樣說.
倘若脫下應該說.
倘若著想.
如果我著想.
如果我可以擁有.
這個新的很榮耀的身體.
我不會感到赤身.
不會感到羞恥.
那會覺得很好.
其實我們在這個帳篷裡的人.
勞苦嘆息.
並不是願意脫下地上的帳篷.
而是願意著上天上的居所.
好使你必死的生命.
被生命吞沒了.
這裡說的是.
今天我們都會勞苦嘆息.
這個真是嘆息.
在地面上有很多壓力.

$^{561}$有逼迫.
有困難.
有沮喪.
有傷害.
有很多東西.
我們真的會嘆息.
在這裡.
我們在這個身體裡.
所經歷的就是很多的無奈.
所以我們會嘆息.
我們其實也不是很想脫去.
不要這個身體.
但是我們有一樣東西更加渴望.
有一個比較.
就是我很願意著上天上的居所.
我很想有一個榮耀的新身體.
我很想.
我們到時就不用再嘆息了.
不用再勞苦了.
很自由.
很舒暢.
而且是榮耀.
是有對比的.
弟兄姊妹.
我們的眼睛是看著哪裡呢?.
是不是每天都看著自己呢?.
我相信.
是有很多事情不開心的.
有沒有看著.
永恆裡的新身體.
那種榮耀.
有一種渴望的心呢?.
我們真的希望.
這個必死的.
就是現在地面上.
這個會衰殘的身體.
有一天就會被榮耀的新生命.
那個身體取締了.
我們很渴望.
很渴望這件事.

$^{601}$而那為我們安排這事的是神.
是神成就整件事.
祂賜給我們聖靈作為憑據.
是聖靈給我們信心.
信心.
這裡說到一件很重要的事.
我們有舊恩.
所以我們有盼望.
而且為什麼會有盼望.
因為聖靈賜給我們信心.
信耶穌之後.
神就讓聖靈住在我們的生命裡.
聖靈會給我們信心.
我們有信心.
知道神會為我們安排.
很奇妙的.
很美好的事情.
所以生命就有盼望.
我看這個世界.
看我所遭遇的事.
不一樣.
盼望有多大.
後面再看.
所以我們總是勇敢的.
是勇敢.
這個勇敢出現了不只是一次.
第八節又說我們勇敢的.
我們看事物有點不同.
今天.
這裡說我們住在這個身體裡.
住在這個身體裡.
我們還在世界上生活.
原來我們和天家是分開的.
我們在地面上生活.
我們當然不是在天家生活.
神即是主.
祂就在天家裡.
在那裡坐著為王.
我在地面上生活.
就離開了主.

$^{641}$不是和主一起生活.
應該這樣說.
在短暫的時間裡.
等你有一天.
我離開這個身體.
我回到天家的時候.
就不一樣了.
我和主一起生活.
而且有一個榮耀的新身體.
是很多無比的.
所以我們勇敢.
更加情願要離開身體.
我不怕死亡.
不怕身體衰殘.
因為我可以與主同住.
很多無比.
這是保羅很有信心的地方.
所以.
這個很重要.
無論住在身內.
或者住在身外.
我們都立了志向.
要得主的喜願.
今天知不知道.
什麼叫有盼望.
我面對現實生活.
我很有力量.
不是磨爛.
磨爛.
天天呼天搶地.
叫生叫死.
很淒涼.
我現在活著.
我很慘.
不是這些.
我很有力量.
很有盼望.
我知道我已經得著.
因為救恩.
而神給我得著的一切.

$^{681}$永恆的.
將來有榮耀的新身體.
可以永遠與主一起.
有永遠的.
很有重量的.
無比的榮耀.
很多無比的福氣在等著我.
今天地面上的一切都很短暫.
算不得什麼.
而且有主同在.
我一定可以勝過.
有這個盼望.
我們就勇敢.
勇敢.
而且勇敢到一個地步.
無論住在身內.
或住在身外.
我們立了志向要得主的喜悅.
活著有方向.
有盼望的人活得有方向.
沒有盼望的人.
活得很淒涼.
信主比不信主的人更淒涼.
我說真的.
但信主有盼望的時候.
活得很有力量.
方向就是要討主的喜悅.
這個身外是說.
我仍然有今生的現在的生命.
或我生命結束了.
在身外.
我進入新的階段的時候.
我們都立了志向要得主的喜悅.
活著不是為了討自己喜悅.
弟兄姊妹.
活著是要討主的喜悅.
我經常聽很多的分享.
我計劃了.
我打算今年這樣.
我打算明年這樣.

$^{721}$我想做這些.
我想得到那些.
我想怎樣怎樣.
我想和誰一起.
我.
我想做什麼.
經常都說.
好像整個「我」字很大.
活著就是為了誰.
我.
活著是為了我自己.
你記住.
在我們信主的那一刻.
主耶穌基督用祂的寶血.
埋屬我們生命的那一刻.
我們的生命不再屬於自己.
我都不止一次在講壇上這樣說.
不是屬於自己了.
弟兄姊妹.
不要再為自己而活著.
乃是為主而活著.
要討主的喜悅.
那主喜悅我們什麼.
大家知道嗎.
我們知道嗎.
有多少人.
我們真的有好好讀聖經呢.
我不讀聖經.
怎知道主耶穌喜悅我做什麼.
或者.
我們可能.
譬如我們看文字有困難.
不要緊.
現在是很好的.
可以聽的.
那就聽聖經吧.
我們可能提倡.
你今天聽了聖經尾.
也可以的.
聽完一次.

$^{761}$或者入不到心.
不要緊.
就聽第二次.
第三次.
第四次.
你聽多一點.
才知道.
原來主耶穌喜悅我這樣.
那我就要這樣生活了.
是立志要得主的喜悅.
為什麼呢.
因為知道有一樣東西.
我們眾人必須站在.
基督審判台前受審.
為使各人按照本身所行的.
或善或惡受報.
小心這裡不是說.
那些未信耶穌的人.
他們會受審判.
不是的.
這裡是說我們.
弟兄姊妹.
我們有一天要面對神.
而且聖經說得很清楚.
審判先從神的家開始.
是從神的家開始.
從我們開始.
我們要面對神.
我們一直活著.
是為什麼而活著呢.
做什麼呢.
按照本身所行的.
或善或惡要受報.
我們要面對神.
你準備好了嗎.
我們不知道主耶穌何時回來.
但知道祂回來那天.
就會審判全世界.
審判先從神的家開始.
我們準備好了嗎.

$^{801}$在這裡我鼓勵弟兄姊妹.
接下來我們的生命有一個大更新.
我今天所說的.
並不是高深的理論.
一個相信耶穌基督有救恩的人.
生命就應該不再一樣.
是活得很有盼望很有力量.
面對困難的時候.
不會被打倒.
或許有情感上有情緒上的影響.
但不應該被打倒.
不會受困擾.
不會失敗沮喪.
每天都像自己的生命失敗一樣.
不會的.
有救恩的人就不再一樣.
聖靈會給我們有信心.
而且有一句話.
我們看第四章說得很好.
第四章第七節是這樣說.
我們有者寶貝放在雅氣裡.
說要顯明者莫大的能力是出於神.
有耶穌基督在我和你的生命裡.
上帝會透過我們彰顯祂的大能.
神的大能要在我們生命裡彰顯出來.
所以基督徒才有那股力量.
在癌病裡不懼怕.
在生意失敗裡站起來.
婚姻失敗裡站起來.
還可以在眾人面前作見證.
這是很實在的.
在我們生命裡能夠走出來.
聖靈會給我們信心.
但今天有件事要問問大家.
我們是否真的明白救恩呢?.
明白啊.
坐在這裡十多年還不明白?.
那為什麼我們沒有盼望呢?.
我們是否真的明白救恩呢?.
如果真的不明白的時候.

$^{841}$我誠意邀請你找屬靈的友伴.
教會的領袖和公.
好好地談一談.
我們身邊需要有弟兄姊妹與我們同行.
為我們禱告守望.
有救恩就有盼望.
有聖靈在我們裡面就有信心.
還有,我怎知道這麼多呢?.
有沒有看聖經呢?.
有沒有聽聖經呢?.
我又不聽聖經,又不看聖經.
我都不知道神的心意是什麼.
我又不知道這麼多.
不知道在說什麼.
這樣就不行了.
第一年是這樣,第二年是這樣.
我們今年十三年也是這樣.
二十年,三十年都還是這樣.
沒有盼望,這樣活著.
是很辛苦的.
我們教會其實沒有耶穌基督的身體.
但我們一起都是一個軍隊.
軍隊打仗,士兵們就要很有勇氣.
很有力.
如果活著,沒有盼望,就一點力都沒有.
走去跟別人傳福音,可能都會很怯.
你要有盼望嗎?.
朋友心裡想,其實我也覺得很沒有盼望.
我也不知道怎樣跟他說要有盼望.
你跟別人說要有信心.
其實我自己也很害怕,很多的沮喪.
怎樣跟別人說要有信心呢?.
我們現在的時候是一個.
我相信主耶穌基督再回來的日子.
近了,要更加多的傳福音.
是每個弟兄姊妹都需要起來傳福音.
但我們現在這樣,如果沒有盼望的生命.
是一點力都沒有的.
弟兄姊妹是時候要起來了.
我們一同禱告.

$^{881}$要去鼓勵我們.
要有盼望.
當我們的眼睛看著地面上所發生的事情.
看著我們衰殘的身體.
看著我們經歷的種種困難的時候.
求主耶穌幫助我們.
看到,因為你為我們成就的救恩.
我們充滿了盼望.
我們將來可以有榮耀的身體.
有榮耀等著我們.
有主耶穌基督賜給我們的永生.
我們什麼都不怕.
今天地面上一切困難.
我們都可以靠著你勝過.
有聖靈賜給我們信心.
我們不會動搖.
主啊,在這裡求你給我們眾弟兄姊妹.
生命有一個很大的跨越.
讓我們要好好地看聖經.
讀聖經,聽聖經.
我們要跟著上帝的教導而行.
從今以後,活得不一樣.
主啊,當洪恩家要邁向未來的日子.
我們的生命需要起來.
求主給我們真的有很大的成長.
奉耶穌基督的名字祈禱.
阿們.
(影片完結).
\newpage



\section{使徒行傳 1:6}
\label{sec:AR6jAGDUTDM}
\textbf{2025.04.13《直到地極》 使徒行傳1\_6―8 講員 : 李思敬牧師}
\newline
\newline
連結: \href{https://youtube.com/watch?v=AR6jAGDUTDM}{\texttt{https://youtube.com/watch?v=AR6jAGDUTDM}} ~~~~ 語音日期: 2025-04-15
\newline
\newline
\hyperref[sec:gznE_6kHeYc]{\small{< < < PREV SERMON < < <}}
~
\hyperref[sec:index_chronic]{\small{[返順時目]}}
~
\hyperref[sec:index_scriptual]{\small{[返順卷目]}}
~
\hyperref[sec:95mRci2_UGA]{\small{> > > NEXT SERMON > > >}}
\newline
\newline
使徒行傳 1:6
\newline
\begin{longtable}{cl}
\hline
\hline
章節 & 經文 (和合本修訂版)\\
\hline
1:6 & \begin{tabularx}{0.7\textwidth}{X} 他們聚集的時候，問耶穌：「主啊，你就要在這時候復興以色列國嗎？」 \end{tabularx} \\ \\
[1ex]
\hline
\hline
\end{longtable}
$^{1}$主席.
如果我沒有記錯.
宣道會很少聖餐是要走出來領的.
你有沒有印象?.
很重要的.
如果你去聖公會.
還要是牧師一個人派給你.
然後他派給你的時候.
他講你的名字.
是啊.
很感動的.
那很浪費時間.
不是的,那段時間在神的面前.
難得專心安靜.
你們今年的主題.
記得吧.
已經過了三個半月了.
懷抱報道初心.
接著呢.
福音遍及社群.
這裡有寫著的.
哈哈哈.
使徒行傳的金句.
我相信宣道會的弟兄姊妹.
一定滾瓜難熟.
聖靈降臨.
得著能力.
從耶路撒冷直到地極.
作主的見證.
其實這句話不是耶穌本來準備說的.
所以我們從第六節開始讀起.
就明白那個「但」的連接詞開始.
門徒問耶穌.
你復興以色列國就在這時候嗎?.
今天也是的.
總之世界有什麼變動.
加關稅.
有人出來說主在內.
聖經的預言應驗.
不是的.

$^{41}$為什麼你會知道.
耶穌說的.
耶穌說什麼.
我們只記得第八節聖靈降臨.
我們不記得耶穌一開始說了什麼.
「乎憑著自己的權柄所定的時候日期」.
多謝PowerPoint帶出來了.
一起讀下面那句.
「不是你們可以知道的」.
任何一個講完站出來說我知道.
他比主耶穌厲害.
不可能的.
我們需要知道.
可以知道.
應該知道.
所以鮑威昆牧師有一次陪靈會說得好.
他說我們是主的僕人.
主人什麼時候回來.
僕人不會問的.
你家的飛龍,太龍,印龍.
每次出門都問你.
什麼時候回來的.
上次遲了.
他不是僕人.
反過來.
他是主人.
所以耶穌提醒我們.
不是問這個問題.
不過.
但.
要問什麼.
聖靈降臨.
得著能力.
福音是怎樣從耶路撒冷直到地極的呢.
這個要問.
所以整卷的使徒行傳.
其實路加醫生就是寫下.
福音是怎樣從耶路撒冷直到地極的經過.
你說沒有啊.
沒有去到南極.

$^{81}$沒有去到北極.
不是.
我們不是在說南極的小島都要徵稅.
我們不是在說地理上的地極.
也沒有來到中國.
使徒行傳沒有.
古代四大文明古國.
哪有啊.
是.
但是如果你回去讀.
以賽亞書49章第6節.
那裡說什麼.
那裡上帝差遣他的百姓.
使你們作外邦人的光.
並且施行我的救恩直到地極.
光就是代表救恩.
地極呢.
外邦.
未認識神的外邦人.
就是聖經裡面所說的地極.
所以使徒行傳28章停在哪裡.
保羅去到羅馬.
地中海的世界.
羅馬帝國條條大道.
通羅馬.
羅馬就是外邦世界的一個代表.
福音去到未認識神的外邦人中間.
使徒行傳由耶路撒冷主復活升天.
一直到地極.
現在我們懂得解釋了.
外邦羅馬.
福音在耶路撒冷很興旺.
你記不記得.
給你溫習一下這些數據.
耶穌升天之後.
留下在耶路撒冷同心祈禱的.
等候聖靈的有多少人.
120個.
但是很快的五旬節來了.
聖靈降臨.

$^{121}$一次報道會.
請問多少人決志.
3000.
第二次再開報道會.
5000.
你會加數的.
3000加5000.
8000.
還有開頭那120個呢.
不止.
路加還記載主將得救的人數.
什麼呢.
天天加給他們.
教會就是在一個這樣的恩典當中.
從120人.
去到一萬人.
天天增加.
少數怕長計.
香港回歸每日150個單程證.
25年之後多少人來了.
你算一算.
真的你走去中環置地廣場.
站在那裡聽.
全部都是普通話.
所以耶路撒冷教會很忙碌.
但是很快樂.
因為每日都見證到福音的大能.
拯救失喪的人.
很多人來信耶穌.
不止他們凡物公庸.
記得.
出過事的.
但是很快就解決了.
有怨言的.
說希臘話的寡婦.
在每日的照顧上.
好像漏了他們.
其實可能都不是漏了.
不過派到他們的時候.
那碟飯就凍了一點.

$^{161}$是啊.
是不是歧視我們.
很多事情說的.
試圖沒有責備.
試圖站出來說.
是啊我們真的不夠人手.
第二代的弟兄姊妹興起.
所以十二試圖之外.
有七個執事.
太快樂了.
耶路撒冷教會.
成為了Mega Church.
超過一萬人.
要買一間大戲院.
是嗎?是啊.
耶路撒冷教會.
舞會.
座堂.
以後來朝聖的基督徒.
都不可以在二樓.
明白嗎?.
是一個很快樂的教會.
不過我稱之為.
快樂的樽頸.
一個Happy Bottleneck.
因為耶穌的吩咐是什麼呢?.
耶路撒冷.
猶太傳遞.
撒瑪利亞.
還有直到外邦地境.
耶路撒冷教會再大.
再多人.
都是耶路撒冷.
聽得明白嗎?.
但是試圖已經盡力了.
每晚開會開到半夜三更.
是不是?就是搞不定.
你試試給一萬人來.
紅恩堂.
你去哪裡聚會?.

$^{201}$你去單車徑吧.
被人撞.
沒辦法解決的問題.
但是就是要面對.
很快樂,很感恩.
大家祈禱的時候滿懷高興.
今天有多少人信耶穌.
聖靈怎麼工作?.
你看到嗎?.
聖靈沒有跟教會開會.
商量宣教計劃.
為什麼停滯不前.
聖靈一個不該.
讓執事會主席斯蒂凡.
信道.
記得嗎?第七章.
哇!大家都呆了.
所以八章第一節.
《士道行傳》記載.
碧白臨到教會.
都不是羅馬帝國的碧白.
其實是猶太會堂的碧白.
但是除了使徒以外.
門徒分神.
去哪裡?.
在猶太和撒瑪利雅各處.
還沒醒來吧?.
路加畫龍點睛.
不是,畫公仔,畫出腸.
八章四節.
那些分散的人.
往各處去傳道.
一夜之間碧白臨到.
教會一萬二千人.
除了十二使徒以外.
全部都分散.
教會人數暴跌.
網上有人這樣說.
是好事.
怎麼會是好事?.

$^{241}$我們奉獻不足.
是不是要炒魷魚?.
還是要動身?.
還是要減薪?.
不是問這些.
聖靈將耶路撒冷教會的.
每一個弟兄姊妹.
都成為宣教士.
都成為福音的見證.
他們很自然分散出去.
去哪裡?.
從耶路撒冷.
去到猶太傳遞和薩瑪利亞.
福音在聖靈的帶領之下.
是這樣離開耶路撒冷的.
這是《使徒行傳》的頭三分之一.
一至八章.
接著呢.
福音是怎樣從亞洲.
今天我們說土耳其.
過到希臘.
馬其頓.
歐洲.
十五章的結束.
保羅和巴拿巴.
吵架.
他們是誰?.
我們今天經常說門徒訓練.
就說保羅和提摩泰.
從來都不提.
其實保羅有個師父.
帶他出生傳道的是巴拿巴.
巴拿巴是個好人.
讀了副道.
MCS畢業.
他是勸惠子.
有沒有看到《使徒行傳》介紹他.
巴拿巴叫保羅來安提柯教會.
大家一起同心服事.
聖靈感動.

$^{281}$分派他們出去傳道.
於是第一次.
很成功帶了外邦人信主.
回來.
舞會.
質詢.
很難想像.
這麼多人信主.
舞會不是.
慶祝.
不是慰問.
不是恭賀.
是質詢.
《使徒行傳》十五章.
很精彩.
讀這些全部是歷史的經過事實.
聖經是這樣記載下來.
結果他們回到安提柯教會.
又要出去.
但這次竟然兩個人.
聖經記載得路加醫生很客氣.
二人起了爭論.
甚至彼此分開.
巴拿巴留在錦繡.
保羅就來了加州.
你明白我的意思嗎?.
分開.
為什麼事分開?.
約翰·馬可.
在第一次的傳道旅程.
約翰·馬可不是一個實習神學生.
他是舞會耶路撒冷派來的代表.
很有身份.
所以出去的時候三個同工.
中間馬可離開.
聖經沒有說什麼原因.
但現在巴拿巴說我們再出去.
我們找馬可一起去.
保羅認為不可以.
巴拿巴是一個好人.

$^{321}$幫表弟.
你知道馬可是他的表弟.
讀過羅西書.
幫人幫到出面.
搞了一輪又賺錢.
不是的.
巴拿巴不是特朗普.
巴拿巴是給保羅第二次機會的師父.
我們什麼時候會忘記.
保羅是受過巴拿巴栽培的恩惠門徒.
他應該很明白師父的心腸.
給多一次機會.
為什麼這麼難?.
保羅心裡認為神的工作不是這樣的.
神的工作很嚴肅.
不可以手扶著泥向後看.
一開始爭論.
我們就馬上問.
我們不是問誰對.
我們是問你幫誰.
你幫誰?.
路加醫生告訴我們.
安提柯教會並沒有因此分裂.
看到嗎?.
沒有因為兩個牧師吵架.
不止.
今天我們香港會說什麼?.
撕裂.
於是就有聖保羅堂.
聖巴拿巴會.
從來都沒有.
仍然是一個教會.
安提柯教會.
但是他們做了一件很重要的決定.
本來一對宣教隊伍出去.
現在雙倍了他們的任務預算.
他們的宣教預算沒有錢.
現在分開了.
大家都出去.
巴拿巴大馬可.

$^{361}$保羅大西拉.
兩對宣教隊伍出去.
福音就是這樣.
更加雙倍的努力.
不是幫誰.
沒有這回事.
不需要考慮這件事.
不需要這樣做.
教會是不需要因為爭吵而分裂.
可以分開.
都是服侍同一位主.
巴拿巴和保羅留下一個.
未必是他們自己決定的.
但是整個教會很成熟.
能夠有這樣的決定.
於是你就會發覺.
在十五章記載下去.
保羅去到哪裡.
耶穌的靈又禁止.
聖靈又不許.
莫名其妙.
保羅不是去犯罪.
保羅不是去度假享福.
保羅去傳福音.
為什麼聖靈會禁止.
耶穌的靈不許.
你讀下去.
就在這樣的莫名其妙當中.
保羅夜間就見到.
馬其頓的異象.
路加立刻記載.
我們隨即上往馬其頓去.
看到聖靈如何預備好.
這個團隊的心.
透過困難.
透過難阻.
透過分開.
透過吵架.
撕裂.
奇怪.

$^{401}$從耶路撒冷出來.
是逼迫.
從亞洲去到歐洲.
原來是巴拿巴和保羅的爭論.
事到行轉的說法.
他們有沒有和好.
我留給你自己查經.
今天時間不夠.
你看保羅的書信.
留意他後來怎樣說馬可.
有沒有留意.
文案的說話裡.
很精彩.
福音是怎樣去到羅馬.
再遠一點.
你今天去地中海.
坐郵輪.
郵園希臘愛琴海.
要過去意大利.
都要走兩天.
是很遠的.
事到行轉.
二十一章.
告訴我們.
保羅趕回耶路撒冷.
為了要遵守聖經的教導.
守五旬節.
不是慶祝耶路撒冷教會周年會慶.
是要守五旬節.
回到耶路撒冷.
事到行轉記載.
他去見雅各.
不是兩個人沒話說嗎.
只是寫迦拉泰書和雅各書.
隔空.
北戰.
因信稱義.
還是因行為稱義.
你又搞錯了.
不要聽這些fake news.

$^{441}$看回聖經的記載.
保羅很尊敬雅各.
所以他回耶路撒冷第一件事.
他去見雅各.
雅各亦很關心保羅.
當然.
雅各會跟保羅說.
耶路撒冷教會.
很多弟兄姊妹都很熱心.
律法成千上萬.
沒有問那句.
你在外邦建立的教會每間多少人?.
坐滿一個地舖已經不得了.
哪有成千上萬?.
那不是重點.
雅各的重點是勸保羅一件事.
說什麼?.
雅各跟保羅說.
有人說.
你教訓猶太人離棄摩西的律法.
不如這樣.
你用行動證明他們這個看法.
是以訛傳訛的.
好不好?什麼行動?.
你帶四個人進去聖殿.
行潔淨禮儀.
眾人就可以知道.
先前聽見你的事.
都是虛的.
假的.
錯的.
下面那句才慘.
就知道你自己為人循規蹈矩.
進行律法.
保羅還需要傳恩信稱義的福音嗎?.
透過這個行動.
讓人知道原來他是循規蹈矩.
進行律法的.
讓你做保羅.
你聽不聽?.

$^{481}$保羅聽.
所以你認識保羅.
你認識他的為人.
他不是那個不斷跟人吵架.
又跟巴拿巴吵架.
又跟彼得吵架.
又跟雅各吵架.
總之所有人都跟他不對.
不是的.
保羅很謙卑.
接受雅各的提議.
就真的帶著四個還願的弟兄.
翌日.
第二天就去聖殿.
結果怎麼樣?.
你讀聖經.
你應該可以上來講的了.
結果暴動.
蘇聯.
猶太人起哄.
聖殿要關門.
不僅如此.
他們要打保羅.
不僅如此.
羅馬的軍隊出動.
介入.
將保羅帶回駐軍的營樓.
在那裡問他.
耶路撒冷都搞不定.
結果羅馬的軍隊.
要將保羅押到蓋薩尼亞.
在軍營還押候審.
我們很熟悉這兩個字.
還沒審的.
都不是定他有罪沒罪.
總之就關著他.
先解決了蘇聯的問題.
羅馬的巡撫.
探戶.
希望教會會送禮.

$^{521}$贖回大使圖.
教會當然不會這樣做.
然後他又要討好猶太人.
所以想將保羅送回耶路撒冷.
一隻眼開一隻眼閉.
猶太人也有陰謀.
要在路上解決保羅.
你讀《士徒行傳》.
好像看韓劇.
其實你不看的.
歷史現實就是這麼貼地.
不是在講神學的理論.
結果保羅就拿出他的BNO.
上訴蓋薩.
公堂送他從巴勒斯坦去到羅馬.
原來《士徒行傳》告訴我們.
聖靈的大能.
不是醫病趕鬼.
講謊言.
行神蹟奇事.
講到很有能力.
這些事情都有發生.
但是福音如何從耶路撒冷直到地極.
路加分開三個階段.
而關鍵和焦點.
我們現在懂得掌握了嗎?.
逼白.
撕裂.
掃亂.
三樣都是我們不喜歡的.
主啊!.
萬萬不要讓這件事淋到我們香港.
不要淋到我們紅印堂.
讓這些惡事遠離我們.
但是聖靈的工作.
《士徒行傳》在第一世紀.
頭一次.
福音如何傳到外邦.
這個事實的記錄.
給我們一個典範,一個基礎.

$^{561}$今天你再看到這些事情.
我們不需要喊生喊死.
我們知道.
聖靈的工作.
初心不變.
聽得明白嗎?.
聖靈的工作不是要保守教會.
沒有逼白臨到.
教會和睦同居.
不會發生爭吵.
然後出到哪裡.
人見人愛.
車見車載.
不是啊.
你讀聖經.
《士徒行傳》告訴我們.
透過逼白教會從耶路撒冷.
就出到撒瑪利亞猶太傳地.
透過撕裂教會.
就從亞洲小亞細亞.
Asia Minor.
去到馬其頓,希臘,歐洲.
透過蘇聯,保羅,上訴,蓋薩.
福音就是這樣.
去到羅馬.
外邦世界的核心.
原來不是我們的計劃.
是上帝的掌管.
上帝做光又做暗.
祂施平安又降災禍.
有沒有讀過這節金句.
不讀的.
我們就是要救苦救難的.
那個是觀世音.
那個不是主耶穌.
分得出嗎?.
耶穌呼召我們.
天天背起十字架.
彼得提醒我們.
你們蒙召是為了什麼?.

$^{601}$因為基督受苦.
留下這個受苦的腳踵.
跟從主不是為了自己的福樂.
跟從主是見證主的大能.
120個人怎麼可能將福音傳到地極.
讀28章的《使徒行傳》.
你就發覺不是靠他們的能力.
不是靠他們的本事.
不是靠他們的熱心.
不是靠他們為主忠心.
迪伯林道沒有人想的.
殉道誰想做撕裂犯?.
我叫Steven.
我也想改名字.
撕裂大家都呆了.
不要發生在我們中間.
上帝可以化腐朽為神奇.
掃斷.
暴動.
我們不喜歡的.
但是原來這是神的大能.
保守他的使徒.
平安去到羅馬.
有士兵.
那個叫做天使的保護.
我們想像不到.
以前的世界一出門.
離開洪水橋區.
你就可以被人搶.
想不到嗎?.
有沒有看《水滸傳》?.
每間客棧都是黑店.
翻開床.
下面就有人伏在那裡拿著刀.
你不讀.
那是宋朝太平盛世的事實.
羅馬帝國.
條條大道都平安.
是因為帝國的軍隊.
還有兩個保住.

$^{641}$跟著保羅去到羅馬.
神的旨意.
真是高過人的旨意.
今天我們說要事奉他.
福音直到地極.
我們不過是沒用的器皿.
放在主的手中.
五個餅.
兩條魚.
可以叫五千人得飽足.
昔日聖經的記載是歷史事實.
今天我們都可以同樣經歷.
這個出人意料的神蹟奇事.
福音如何可以.
不單止在香港.
九龍新界.
現在你們成為焦點的焦點.
核心的核心.
大灣區接駁的.
就是洪水橋.
就是北區.
上帝要做什麼.
上帝要用教會.
來使大灣區蒙福.
聽聞福音.
沒可能的.
我們人數這麼少.
上帝可以逼迫的.
不如趁逼迫之前.
自己先出去.
不用祂動手.
撕裂.
不撕裂都可以去.
不是要靠這些.
要先弄一場暴動.
不用的.
一樣可以去.
神的恩典.
神的能力.
從來都沒有改變.

$^{681}$指導行傳.
就是讓我們見到.
一個這樣的事實.
今天繼續要發生.
在我們中間.
我們同心低頭祈禱.
親愛的主.
我們來到受苦節的前夕.
我們一同參加聖餐.
紀念主為我們釘十字架.
為我們流出你的血.
成就莫大的救恩.
其實我們每一個人.
相信耶穌.
都是一個神蹟.
沒可能的事.
我們的家族.
能夠認識神.
棄假歸真.
都是一件想像不到的恩典.
今天你要繼續透過教會.
透過弟兄姊妹.
是,我們軟弱.
我們很有閒.
我們沒辦法想像.
為什麼你要揀選.
一群軟弱的人.
去完成你的救恩.
直到地極.
但這就是你的心意.
也是你聖靈大能的工作.
在我們的眼前.
向我們顯明.
多謝你.
我們恭敬將自己奉獻.
交託在恩主的手中.
不是祈求出入平安.
凡事順利.
而是祈求.
你用得著我們每一個.

$^{721}$成為你的見證.
成為光.
成為別人蒙恩的管道.
求你聽我們不配的祈禱.
奉耶穌基督的名祈求.
阿們.
\newpage



\section{羅馬書 8:28}
\label{sec:95mRci2_UGA}
\textbf{2025.04.20《哭泣的祝福》 羅馬書8\_28 講員 : 曾裔貴牧師}
\newline
\newline
連結: \href{https://youtube.com/watch?v=95mRci2-UGA}{\texttt{https://youtube.com/watch?v=95mRci2-UGA}} ~~~~ 語音日期: 2025-04-23
\newline
\newline
\hyperref[sec:AR6jAGDUTDM]{\small{< < < PREV SERMON < < <}}
~
\hyperref[sec:index_chronic]{\small{[返順時目]}}
~
\hyperref[sec:index_scriptual]{\small{[返順卷目]}}
~
\hyperref[sec:o8E8B2LqAe4]{\small{> > > NEXT SERMON > > >}}
\newline
\newline
羅馬書 8:28
\newline
\begin{longtable}{cl}
\hline
\hline
章節 & 經文 (和合本修訂版)\\
\hline
8:28 & \begin{tabularx}{0.7\textwidth}{X} 我們知道，萬事都互相效力，叫愛神的人得益處，就是按他旨意被召的人。 \end{tabularx} \\ \\
[1ex]
\hline
\hline
\end{longtable}
$^{1}$今天在座也有一些新朋友第一次來參加.
也代表教會歡迎大家.
很盼望今早我們的訊息能夠.
剛才永文弟兄他的成長.
我們看到如果沒有神在他的生命.
在他的家庭.
藉著一個基督徒.
一群愛他的朋友.
他不會有今天能夠在基督教的機構裡.
服侍很多人.
用他的熱心.
在傳媒界不斷製作不同的影音節目.
可以幫助很多人.
所以我們很想今天能夠藉著這個.
非常特別的題目.
哭泣何來會祝福呢?.
不過有主耶穌.
有神的恩典.
由這節聖經所說的.
好,讓我來.
謝謝.
我們讀一節這節聖經.
很簡單的.
我們知道萬事都互相效力.
叫愛神的人得益處.
就是按他旨意被召的人.
這裡要留意.
這節聖經有好幾個重點.
如果我們打壓不了.
不明白.
就會失之交臂.
我們就會錯失了裡面的寶藏.
萬事可以互相效力.
萬事其實不會自己互相效力.
相信一定是我們信仰最重要的.
就是神使到萬事互相效力.
大家同不同意嗎?.
神使到萬事互相效力.
剛才你聽到美文弟兄所說.
由中學開始.

$^{41}$已經是一條不容易.
不是我們正常的少年人成長的路.
如果沒有神在他的生命裡面.
其實學壞的機會是很大.
怎會有今天.
可以在我們大家當中.
來分享他的人生故事.
怎會夠膽可以很坦誠.
和我們分享人生故事呢?.
神使到每件事.
不好的事.
都能夠效力.
為誰效力?.
就是為神所呼召揀選的人效力.
所以今早我們這個福音聚會.
我們很想告訴你.
到底信耶穌的人.
是一個怎樣的人呢?.
剛才我說有三個很重要的重點.
在這節聖經裡面.
我們說了萬事是因為神.
使到萬事效力.
為誰人效力呢?.
效力又有什麼帶來的好處呢?.
我們可以再看看.
這裡告訴我們.
萬事原來可以是人生裡面.
大小的事.
好的事.
不好的事.
就是萬事.
我們未信耶穌的人.
我們希望的就是.
每天都平安.
每天都不會有任何的不好的事.
不好的意外.
不好的身體的狀況.
我們希望是完全沒有.
但是.
大家在在都知道.

$^{81}$你每天面對的.
有些可能是你預計到.
更多的是我們預計不到的事情.
預計不到的事情.
我們自己都不能夠去掌握.
其實很多這些事情.
聖經告訴我們.
原來信耶穌的人.
不是寄望萬事都要如意.
萬事都要和好.
萬事都要平安.
原來不行的.
原來神使萬事去效力.
最重要的就是.
那個益處.
不是萬事都必須要平安.
得到什麼益處呢?.
這就是這節聖經.
和今早.
無論剛才弟兄的見證.
和牧師今天很簡單的訊息.
很想大家一起來思想.
到底一個信耶穌的人.
到底得到什麼益處呢?.
什麼才是我們生命裡.
最重要的益處呢?.
原來不是祈求萬事順利.
祈求神要使每一樣東西都如我的心意.
一定不是.
我們自己都證實不是.
但是.
這告訴我們.
原來就是第29節.
因為他預先所知道的人.
就預先定下.
效法他兒子.
就是耶穌的模樣.
使他的兒子耶穌.
在許多弟兄當中作長子.
原來這節聖經回應那個益處.

$^{121}$原來不是我們的物質.
我們每天的生活都要平安無事.
原來不是.
原來神使萬事效力.
為著的是我們的生命成長的益處.
為著我們能夠始主耶穌.
為了我們能夠有神的愛的生命.
為了我們內心能夠有平安.
有永恆盼望.
這個益處.
這個就一定是信耶穌的人才有.
不信耶穌.
未認識神.
不肯相信神的人.
你就沒有這個益處了.
當然.
信和不信.
我們每天有時會很平安.
有時我們會很富足.
但是更多的時間.
我們面對的是很多不同的艱難.
譬如現在的貿易.
現在的關稅.
或者世界的疾病.
我們逃不過的.
我們都要面對的.
但是信耶穌的人有獨特的福氣.
就是神為我們的生命效力.
最終是怎樣?.
令我們得到益處.
這個益處就是我們的心靈富足.
我們有神的愛在我們的生命裡面.
嘗試用一個例子來想一想.
我們怎樣去面對.
人生是很無常的.
沒有東西是絕對你可以擁有的.
在台灣.
有一個年輕人.
一次不小心的意外.
一隻手指斷了.

$^{161}$他就終日.
在一種痛苦.
埋怨.
關在房門裡不下樓.
媽媽很擔心.
不知怎辦.
不上學.
活在他自己的絕望痛苦裡面.
家人都要陪伴他面對這個痛苦.
媽媽打電話去台灣.
一間殘疾的護養院.
希望有義工來探望他.
安慰他.
開解他.
聯絡了義工.
到那天.
還未來之前.
媽媽就跟他的兒子說.
今天會有義工的姐姐來探望你.
得到的答覆就是.
我不要人幫助.
你不要找人來.
媽媽很不好意思.
打電話去叫他們不要來了.
義工的姐姐說.
不要緊.
反正都來了.
我繼續來.
來到的時候.
媽媽見到這位義工.
都不好意思.
原來是雙腳有殘疾.
他的腳是一個很嚴重的殘障.
要穿著一些很重的鐵鞋.
穿著這樣兩個拐杖.
他的小朋友住在二樓.
怎辦呢.
不如你不要去了.
這麼麻煩.
義工的姐姐說都來了.

$^{201}$不要緊.
你介不介意將我其中一個拐杖.
先拿上二樓.
我扶著樓梯這樣上去.
於是這個義工姐姐.
就這樣上去.
剛才都說了.
他兩隻腳.
是穿著很重的鐵鞋.
來支撐著他的腳的力.
所以每一步上樓梯.
這樣走上去.
很重的力.
單手扶著那個樓梯.
一步一步這樣上去.
那個兒子.
在房裡聽到.
為何這麼嘈的聲音.
自己打開門出來.
聽聽是甚麼事.
就見到這個姐姐.
這樣扶上來.
她一見到之後.
她整個的心都好像融化了.
原來她自己的手指.
沒有了.
是多麼的微不足道.
就令到她整個人的生命.
又重新燃點一種希望.
弟兄姊妹各位朋友.
我們的人生.
真的不到我們掌握的.
我們會遇到大大小小的意外.
問題就是.
甚麼令到你.
能夠在你不開心.
或者不順利的環境.
有一個重新的力量.
去解釋和去面對.
使到你有力量呢.

$^{241}$剛才這個年青人.
他感受到原來有人比我更加不幸.
有人比我更加嚴重.
他就有力量面對他自己微小的困難.
信仰我們所信的耶穌.
就是很想我們得到真正的益處.
不是米缸要不斷滿.
我們的銀行戶口要不斷都要這麼多.
我每個月都可以去旅行.
不是這些.
這些是物質上面的.
有些人有.
不代表他一定會很開心.
但是信仰是告訴我們.
我們的心靈像耶穌.
這是神一定要在我們生命裡面建立和做的.
每一個基督徒都可以享受與神的和好.
帶來的平安和那個福樂.
弟兄姊妹.
這是一個很重要我們要發現的.
我們需要認識的就是.
原來人生是神使萬事互相效力.
所以這個世界沒有神.
就很混亂.
因為太多的事情我們掌握不到.
太多我們自己個人的事情.
家庭的事情我們掌握不到.
萬事不是益處本身.
所以你會見到.
有時就算是很出名的人.
很有錢的人.
很成功的人.
他都未必是開心的.
當然窮的人不等於他會開心.
每個人都有他的困難.
哪些人才真正得到益處呢?.
得到益處的人只有一類的人.
就是認識耶穌.
持主耶穌的人.
他就是因為神在他生命裡面去效力.

$^{281}$讓每一件事情.
好像永雲大兄所說.
其實不是他想這樣.
但是姐姐的遭遇.
帶來的是一個家庭.
一個很大突然的變故.
但是神叫萬事互相效力.
他們的家庭就得到益處了.
生命建立.
有了這份和諧和滿足.
最後我想用一個例子.
結束這個講道.
這位叫張煒醫生.
在2002年的時候.
他在香港的浸信會洗禮.
其實他不是香港人.
他是在遼寧的一位專科醫生.
在1997年5月12日.
他是當時的骨科的專科醫生.
響應國家的調配.
去了也門.
來到一個門外的醫生.
有一天休假去游泳.
在一個泳池跳水.
誰知道原來那個泳池的水不夠深.
他一跳下去.
整個人下半身完全不能夠再動.
到今天他還是一個嚴重殘障的人.
你可以想像一下.
一個年輕有為的專科醫生.
突然之間變成一個嚴重殘廢的人.
回到內地.
每天躺在床上.
動不了.
痛苦到一個地步.
每天要媽媽幫他轉身.
去清潔.
整個人就是活在這樣的狀態.
他有一次跟媽媽說.
因為每天都推他出去曬太陽.

$^{321}$吸新鮮空氣.
媽媽你推我去懸崖.
推我下去吧.
我不想再生存.
一個二十多歲的年輕人.
一個專科醫生.
一心想著加入政府.
加入黨.
去做一個對社會有貢獻的人.
在他成長的時候.
不相信有神.
相信個人的能力可以決定一切.
而且讀書又很驕傲.
很自恥.
但是這樣的打擊.
令他整個人落入絕望.
算了.
你幫我解決我的人生吧.
突然有一天.
在醫院做復健治療.
有一個女士.
看到這個年輕人.
為什麼這麼憂愁呢?.
為什麼會這麼痛苦呢?.
本來已經走了.
走遠了.
回頭跟她聊天.
安慰她.
這位是一位日本的交流教授.
叫做矢谷靈子.
是日本人.
但是這位姐妹.
是信主的.
安慰她.
你不要這麼憂愁.
來到她為她祈禱.
接著.
回到日本之後.
送了一本書給她.
就是後來.

$^{361}$這位張醫生.
來做翻譯的書.
這本書的原作者.
是一位美國的姐妹.
也是跟張醫生完全一樣的遭遇.
也是一次的跳水.
就這樣令自己整個人.
因為神經中斷.
就殘缺了.
但是這個鍾離的女士.
美國的女士.
她展現的是一個很大的生命力.
很有那種的愛和平安.
所以對張醫生來說.
她收到這本英文的書.
覺得很奇怪.
因為你想想.
張醫生由她這樣遇到這個意外之後.
每天差不多是24小時.
都是要在床上.
所以她太悶了.
終於有一天叫媽媽揭那本書給她看.
因為她的手不能動.
她的腦是完全清醒.
完全健康.
但是除了頭以外以下.
全部都不能動.
叫媽媽揭那本書給她看.
越揭越感覺到一種力量.
為什麼這個人跟我完全一樣的遭遇.
她會有這麼大的平安喜樂.
她就根據那本書.
向神祈禱.
她自己又感動.
就將這本書翻譯成中文.
希望幫助更多跟她遭遇的人.
就是這樣.
張醫生開始接觸基督教.
起初這本書裡面有很多經文.
她完全看不明白.

$^{401}$因為她是一個不信主的人.
看不明白.
但是後來她一直看這些聖經.
她居然在翻譯過程.
跟聖經說的完全一樣.
她就因為這本書.
自己信了耶穌.
向神禱告.
信了耶穌.
又感覺自己突然間有一個很大的生命力.
接著輾轉有她認識的基督徒.
再跟她很清楚的決志.
她一直回教會.
後來終於在2002年.
就被邀請來到香港的其中一間浸信會.
好像是北角的浸信會.
接受洗禮.
跟她媽媽一起.
信主耶穌.
見證神的恩典.
她現在在遼寧省的安山市.
開了一間福建中心.
曾經在中央電視台訪問她.
為什麼妳會有這樣的生命力.
去介紹耶穌.
大家想想.
一個二十多歲的這麼厲害的人.
一下子掉進一個低谷.
一個絕望裡面.
怎樣能夠展現這種生命力呢?.
原來就是神使我們認為不好的事.
來效力使得她可以得到這樣的生命的轉變.
是神在人的內心.
各位朋友.
各位弟兄姊妹.
今天很想藉著永文弟兄剛才很真誠.
還有一些絕密的相.
在感受我都沒有看過.
但是很坦誠地分享我們的過去.
我們不是超人.

$^{441}$不是因為我們特別厲害.
所以營造了一個好的家庭.
好的父親兒子的關係.
而是因為耶穌.
而是因為神使得萬事.
我們認為不開心不好的事.
來效力最終使得那個人.
他所愛的人得到益處.
親愛的朋友.
有些人第一次今天有家人邀請你們參加我們的報道會.
這個福音主日.
其實我們很想介紹的就是.
耶穌你相信他.
你不會突然之間有什麼改變.
你不會突然之間你的家庭.
你的自己的人生會有馬上的改變.
但是相信了耶穌.
他會使得萬事互相效力.
他會使得你的心靈得到益處.
他會使得你整個家庭因為有耶穌的同在.
會帶來和睦.
會帶來你內心的平安.
會帶來你對永恆有盼望.
你不會再懼怕明天.
你不會再在一些你駕馭不了的事情上.
你會失控.
你會無助.
你會跌入絕望裡面.
親愛的朋友.
張醫生.
永文弟兄.
其實我自己也是.
未信耶穌之前.
我就真的自己離家出走.
最高紀錄就是兩個星期不回家.
因為極限了.
因為沒錢了.
要回家了.
每個人都有我們的過去.
但是有耶穌.

$^{481}$會使得我們的生活.
使得我們的所謂命運.
有一個真正永恆的意義.
在耶穌裡面.
所以今天很想.
在這個時候牧師代表著神.
去邀請大家.
你的生命.
如果不在耶穌裡面.
我們自己駕馭不了的.
如果沒有耶穌.
我們自己不能夠面對困難的.
如果沒有耶穌.
我們是沒有力量.
跟人有好的和平的交往.
所以信耶穌.
不是你的生命馬上改變.
但是信耶穌.
你會得到神的同在.
帶給你整個的生命.
慢慢就有益處了.
各位.
這個益處.
是耶穌已經為你預備了.
他已經釘在十字架上.
每一個信的人.
就是他呼召他所愛的人.
有哪一位朋友.
勇敢的.
你接受這個邀請.
你要跟耶穌建立這個生命關係.
那個益處就會臨到你的生命.
臨到你的家庭.
臨到你身上.
有沒有朋友.
第一次.
尤其是男士.
勇敢的走出來.
接受這個邀請.
得著神為你萬事效力.

$^{521}$有沒有這樣的朋友.
你願意來接受的.
你可以舉手.
甚至乎.
如果你夠膽的.
你可以走出來.
勇敢的承認.
我願意跟隨我的主耶穌.
接受他成為你的生命的救主.
同行人生的這位救贖主.
有沒有哪一位朋友.
你願意接受耶穌呢.
勇敢的.
因為耶穌已經為我們犧牲在十字架上.
毫無保留地來寫去他的聖名.
只要我們有耶穌.
我們的生命就有真正的益處.
否則我們只能夠自己面對.
有沒有這樣的朋友.
你願意接受主耶穌呢.
得著這個永恆的救恩.
有的請你舉手.
牧師稍後會為你祈禱.
來讓你成為神的兒女.
有神在你的生命裡.
在你的生活裡來效力.
有沒有這樣的朋友呢.
你不需要理會其他人.
因為這個時候是神在你的內心的感動.
的一個帶領的一個催促.
有沒有這樣的朋友.
願意相信耶穌基督呢.
接受耶穌成為你個人的救主呢.
如果有的請你舉手.
不要緊.
我們很歡迎大家今天來參加我們的福音主日.
也能夠在今早我們很感恩.
永文弟兄給我們有很好的信息.
他的生命的成長.
和神在他的生命裡的見證.

$^{561}$更加今天很想用這一節聖經.
羅馬書的八章.
來和大家彼此的勉勵.
我們有這個很重要的效力.
就是神在我們生命裡的效力.
我們能夠得著這個最重要的.
就是我們生命裡的益處.
所以很盼望這節聖經送給大家.
在今年的復活節.
成為你生命裡不一樣的一個日子.
不如我們一起讀一次這一節聖經.
我們盼望大家帶著這一份的祝福.
這樣回家.
再來思想主耶穌和你的生命的關係.
我們一起讀一次.
我們知道萬事都互相效力.
叫愛神的人得益處.
就是按他旨意被召的人.
羅馬書八章二十八節.
這節聖經是很寶貴送給大家.
\newpage



\section{路加福音 1:39-45}
\label{sec:o8E8B2LqAe4}
\textbf{2025.05.04《載道,信道,行道》 路加福音1\_39-45 講員 : 蔣文忠傳道}
\newline
\newline
連結: \href{https://youtube.com/watch?v=o8E8B2LqAe4}{\texttt{https://youtube.com/watch?v=o8E8B2LqAe4}} ~~~~ 語音日期: 2025-05-04
\newline
\newline
\hyperref[sec:95mRci2_UGA]{\small{< < < PREV SERMON < < <}}
~
\hyperref[sec:index_chronic]{\small{[返順時目]}}
~
\hyperref[sec:index_scriptual]{\small{[返順卷目]}}
~
\hyperref[sec:vW2xmW8knK8]{\small{> > > NEXT SERMON > > >}}
\newline
\newline
路加福音 1:39-45
\newline
\begin{longtable}{cl}
\hline
\hline
章節 & 經文 (和合本修訂版)\\
\hline
1:39 & \begin{tabularx}{0.7\textwidth}{X} 在那些日子，馬利亞起身，急忙前往山區，來到猶大的一座城， \end{tabularx} \\ \\ \relax
1:40 & \begin{tabularx}{0.7\textwidth}{X} 進了撒迦利亞的家，向伊利莎白問安。 \end{tabularx} \\ \\ \relax
1:41 & \begin{tabularx}{0.7\textwidth}{X} 伊利莎白一聽到馬利亞問安，所懷的胎就在腹裡跳動。伊利莎白被聖靈充滿， \end{tabularx} \\ \\ \relax
1:42 & \begin{tabularx}{0.7\textwidth}{X} 高聲喊著說：「你在婦女中是有福的！你所懷的胎也是有福的！ \end{tabularx} \\ \\ \relax
1:43 & \begin{tabularx}{0.7\textwidth}{X} 我主的母親到我這裡來，為何這事臨到我呢？ \end{tabularx} \\ \\ \relax
1:44 & \begin{tabularx}{0.7\textwidth}{X} 因為你問安的聲音一入我耳，我腹裡的胎就歡喜跳動。 \end{tabularx} \\ \\ \relax
1:45 & \begin{tabularx}{0.7\textwidth}{X} 這相信的女子是有福的！因為主對她所說的話都要應驗。」 \end{tabularx} \\ \\
[1ex]
\hline
\hline
\end{longtable}
$^{1}$(主持人:很感恩在此大家中間一齊去敬拜,一齊去聽上帝的話語,剛才讀的經文可能大家有時會在聖誕節,或者有時會輕讀過,未必會很留心這段經文,一會兒我們會詳細看一下這段經文當中,特別是兩個姊妹,兩個主角,.
在與道,與聖道要降臨,與道成肉身的奧秘發生的時候,他們的回應,他們的關係是怎樣?.
大家可能都知道我們忙碌了復活節和壽難節,其實在教會的節期,這段時間是復活節期,有七個主日,直到所謂的聖靈降臨節,小人傳是第二章,聖靈降臨在使徒的生命中間..
耶穌其實之前陪伴門徒多少天?小人的第一章,陪伴了四十天,然後升天,然後再過了一天,聖靈降臨在使徒的身上..
在復活節期裡面,其實特別是拉丁教父和特土良都說,這是一個最喜樂的空間,這段五十天,這七個主日,其實是繼續讓我們在教會當中回味和感受復活帶來的喜樂..
食活的能力帶給我們生命的改變,整個人世間的局面,其實不是讓我們世界當權的人帶動,而是讓世界掌權的主繼續與我們同在,引領著我們人類得到生命真正的拯救..
所以我們可以想想,我們不要停下來,我們懷著什麼心情去期待聖靈,即復活的能力給我們的改變呢?我們是不是一直在經歷復活帶來給我們這種新的生命,新的機遇,新的盼望呢?.
這種更新,這種新機,就是在復活節當中賜給我們.當然,其實我們慶祝的節期,或者我們能夠對生命有一種新的改變,有新的盼望,那個根基,那個力量,其實是源自於,像今天所說的上帝的道..
在上帝的道,與人同在,上帝的道,復活帶給人一種力量的改變,成為我們一個生命可以持續往前走,就算遇上不同的艱難,不同的挑戰,其實我們都不需要懼怕,不需要憂慮..
所以今天我們看這個經文,其實是想跟大家一起去思考,究竟我們弟兄姊妹,我們教會,其實我們與道之間有一種怎樣的關係呢?.
這個我今天的題目是講再道,進道,行道,好像有三種不同的樣式,模式,同道的關係,但其實我最後都講了同一件事情,希望大家留意,希望跟大家一起分享的時候,你會明白得到..
我們從剛才讀的經文,路加封第1章39到45節,這裡是一個短短的記載,但它是一個很重要,也都是在表示瑪利亞和以利沙白的時候都忽略了的老姐妹,當時來說,她與道之間的關係..
我們通常會將重點放在瑪利亞,但有時候因為我們不想太高舉瑪利亞,所以我們不會看很多關於瑪利亞在過程中的生命經歷..
這本書很簡單,前面就講到天使加巴利與瑪利亞報信,大家應該很熟悉,你會懷胎聖靈感恩,伸出我們的救主.瑪利亞就是一個心懷相信,回應願主的話,成就在她生命當中..
她第一個反應就是在這裡,第39節,就是去找以利沙白,跟她分享這個情況,是不是很特別?其實是一個很獨特的,在這段時間裡面,瑪利亞與以利沙白的一個很獨特的團契..
瑪利亞沒有什麼人明白她,她還沒有出嫁,大家明白,天使告訴她,聖靈要讓你福中懷胎,然後要繼續長大成人,成為一個救主..
這就是所謂的道成肉身,一種很立體,一種我們很細心去想的時候,是一個我們想的時候,瑪利亞的角色,做母親的角色是非常重要的..
所以她一直相信,一直思考,究竟是怎麼發生的,其實還沒有發生的,剛剛聖天使說完之後,究竟其後的九個月的懷胎日子,然後胎兒出生,長大成人的過程是怎麼樣呢?其實瑪利亞也沒有路燭,沒有一個男徒知道怎麼做,只是一步一步學習怎麼去回應..
所以她找了以利沙白去分享她的一個特別的經歷,可能有一點忐忑,但是憑信繼續往前走,因為天使跟她說,你的親戚都懷胎,當時已經懷了六個月了,她是一個很獨特的兒子,都即將透過以利沙白,一會兒我會講多一點她的故事而出世..
所以當她相遇的時候,就有一個經文,以利沙白跟瑪利亞說,你福中的胎兒是有福的..
瑪利亞有一個懷度的經歷,是在預備道成肉身,有沒有想過,有時候我們一看耶穌的生命,就想到祂,最多可能是我們去受洗,受試探,然後出來傳道..
有沒有細心思考過,原來我們的救主,他定義道成肉身的時候,其實是完完全全地取了一個人的樣子,他完全地不是一出來就成人,他完全地就是一個母親的福中的胎兒出世,然後學習走路,這樣成長的過程..
其實他取了一個最脆弱,最不可避免的形式,大家可能都聽過很多悲劇,或者一些很不開心的事情,有時候很多母親懷胎,但很不想要的小朋友,她就把他放下,或者出世之後,將他棄掉等等..
我們也聽過很多時候,在社會當中,很多故事,很多我們很不如意的事情.孩子福中的胎兒是不能夠作主的,他沒有能力的,他完全依靠著他母親的母體,給他什麼,給他什麼營養,怎樣照料他,他就怎樣成長..
大家在座的姐妹,有小朋友,你見過,有過孫子等等,就明白這個過程,是很不容易,但也是一個互相影響的過程..
瑪利亞就要承擔著,為著他的孩子成為救主,擔起母親的角色,是一個很獨特的遺訓,要在瑪利亞福中去承認..
甚至是一見面的時候,以利沙白土宗的施西漢一見到救主,救主不一定是我們那時候的畫像,一定是長頭髮,有鬍鬚的耶穌,很慈悲,慈愛的耶穌.有沒有想過一個胎兒裡面,都能夠回應一個胎兒的耶穌?.
有時候我們母親知道懷孕的時候,她會踢我,沒有什麼可以做的,她最多是跳躍,在那裡活動..
現在在以利沙白的母宗的施西漢,六個月的施西漢,他就在那裡興奮跳動,以致令到以利沙白,原來我面對著的瑪利亞福中,就是他自己生命的救主..
對著胎兒,他可以宣認他是他生命的主..
是很特別的,正正是這段時間才會發生的..
他成為了一個人,成為了一個呼召門徒,已經是另一個故事..
到我們今天,我們常常去,剛才說,我們從聖誕節到壽難節到父節,我們都是在想一個成人形象的耶穌,很多時候..
沒想到這個道成有生,原來是一種這麼徹底,這麼將自己放在人的手中,讓自己的生命經歷一個家庭當中的成長..
大家可以繼續去思想,去默想耶穌這種道成有生的奧秘,那種徹底的謙卑,去到什麼程度?.
我們特別要看馬尼亞這個生命的樣式的時候,他說一個在道的生命是怎樣呈現給我們看到的呢?.
第一步就是,剛才說的經文,39節,馬尼亞起身急忙前往山區..
為什麼這麼急呢?很快,立即,他沒有猶豫..
他這個承載著,懷著道的生命,他沒有為自己很多憂心,很多考慮,他即時回應著..
雖然他在途中,他即時用行動來回應這個救主..

$^{41}$這個即時的,是表示一種馴服的回應,是全然接受一個新的身份..
他也不知道會怎樣,要成為救主的母親..
大家有沒有留意,其實這個馬尼亞的角色,在初代教會是一個很爭議性的討論..
當然這個討論不是因為馬尼亞,而是因為耶穌基督的神人異性..
引致他母親牽連在內,其實是這個意思..
究竟馬尼亞是所謂人之母,還是神之母,就是這個討論..
當然應該是神之母,因為他母腹中的那個是完全的神..
所以就爭論這些問題..
待會兒,你看到伊爾沙伯的宣講,你就知道,宣認他..
馬尼亞是即時,殷然,不是殷然,戰戰兢兢可能都是,接受一個新的身份..
沒有人能夠定義的,是神正在定義他這個身份..
新的挑戰,新的方向..
大家有沒有想過,上帝呼召你,你正在做的事情是全新的,你不是在重複前人做的事情..
因為都是屬於你,給你一個獨特的使命和任務..
這是很神奇,這是很奇妙的事情..
在上帝眼中,每一個人的生命都是很獨特..
在上帝國當中,扮演的角色都是很有獨特的一個角色,一個恩賜..
沒有了你是不行的..
你當中發揮的時候,就能夠祝福到其他人..
可能聽了這麼久,你會說,馬尼亞又是神之母,又道成肉身,.
成為祂身體中間的一個胎兒..
我們不是,我們不是馬尼亞,那跟祂有什麼關係呢?.
當我們思考這個位置的時候,這個再度的生命的時候,其實跟我們息息相關..
大家記得耶穌有殺種的比喻,.
有落在淺土上面,轉頭被烏鴉淡走了,.
我們被很多世務纏繞著,忘記上帝的道..
也有落在好土裡面,總之能夠健康地紮根,一直茁壯成長,結出很多的果子..
這個土壤就是我們的生命,其實是一樣的..
上帝的道,好像今天孫子說的,我們聽道,我們看聖經,.
落在我們心靈的土壤裡面..
我們是一個什麼樣的土壤呢?.
質量,質地是怎麼樣的呢?.
我們有很多的淺土,很多的荊棘,不能夠讓道在當中佔一個位置,影響我們整個生命..
還是我們預備好自己,是一個很接受的,很願意領受上帝話語的土壤呢?.
以至上帝話落在我們生命當中,我們是能夠欣然接受,然後讓它可以發芽生長呢?.
所以讓道在生命中的成長,是我們每個弟兄姊妹都能夠學習得到瑪利亞這個再度的生命的樣式..
我接著說杜伊尼撒白,我想用更多的時間說他..
因為我們每個人都信道,再度好像很難理解,怎樣承載著呢?.
我們有福音,有上帝話,我們相信,我們大家都會說,.
很容易,我應該可以在當中更多的學習..

$^{81}$而撒白是一個很好的樣式,讓你去思考,怎樣為之信道的生命..
大家聽聞當瑪利亞急忙到山區找他探望的時候,他的反應就是,四十三夕,我主的母親竟然在新日本..
其實不緊要,大家和日本的和諧差不多,不過我覺得日本的神髓更加到位..
你竟然來到我這裡?這件事怎會來到我的生命當中?.
很有趣吧?你們記住,伊尼撒白當時是年紀老邁,我想像起碼五六十歲,懷孕,很神奇,是一個神蹟的事..
瑪利亞還未成年,未出嫁,十幾二十歲,相差一截,年紀上..
伊尼撒白是他的長輩,想像一下,大家在拜年的時候,你的長輩去後輩的家,還是後輩去長輩的家?.
理所當然的,我們有大食大,中強的文化,一定是去最年長的家,我們全部子子孫孫都去那裡,一定是這樣,.
一定是拜年後再去,所以,瑪利亞向一個懷胎六個月的前輩,一個親戚,去拜會他,送燕窩,送什麼給他吃,理所當然的,.
不需要這麼大反應,但是伊尼撒白倒轉了整個人倫的關係,他見了一個十幾歲的妹妹,.
他說,我主的母親竟然來到我姊女,這件事怎麼會發生?正常人不會這樣說話的,就算你見到後輩來到,你很尊重他,你也不會這樣說話的..
這是一個順道讓他見到在面前的瑪利亞,我主的母親,已經回答了神之母的問題,其實伊尼撒白一句已經回答了,KO了問題..
我主當然是神,不是母親厲害,是他肚中的我主,使這個人叫我主的母親,大家明不明白?不是我主的母親,是我主才是最首要,伊尼撒白要回應的對象,不過對著一個腹中的胎兒這樣說..
為什麼會來到我這裡?為什麼會由我主來拜會我?所以他覺得是很不配,很受寵若驚的反應,是不是信心很大?.
他見到整個在神國當中的秩序,雖然我年紀大一點,但我是一個排在幕後的角色,對著瑪利亞,對著她母親腹中的胎兒,對他來說是一個更加尊貴的角色,所以他覺得怎麼會發生這件事?萬萬不能,很不配,你來到我家中..
更加反映他早已預備好主動迎接道的降臨..
我們問,如果今天好像葛國考,或者像啟示所說的,耶穌在我們心門外等候我們的時候,我們又會作出什麼反應,回應上主的道降臨呢?.
這是很值得去想的,其實常常都在我們耳邊,常常都在我們心門外,祂想進入我們的生命當中,進入我們的空間裡面,我們是不是都好像以利沙白那樣預備好,欣然地一說就說中一個應該說的話,還是祂還猶豫?.
我們做了很多的盤算,很多的考慮,都還沒有認出主在我們面前,你見到很多,剛才說復活的時候,大家都認不出主,女士也認不出主,門徒也認不出主,直到耶穌自己顯現的時候,說了祂的話的時候,才認出主來..
我們不能夠一見到主就認出主,祂可能不是用一個我們常慣的面孔見面,腹中的胎兒一定不是常見的,沒有想過..
以利沙白是全然相信道成的新告庇,她才見到馬來的母親,她的胎兒,可能剛剛開始就說懷胎,不過司祭還提醒她,她腹中的胎兒跳躍,讓她知道這個母親,這個眼前的親戚是一個主的母親..
她即時就作一個反應,作一個相信,作一個信服,我們有時考慮很多,盤算很多,猶豫很多,但你見到,無論是馬來的即時去到山區,或者是一開門,以利沙白一見到馬來就說,我主的母親竟然來到我家中,這些反應,證明他們信心很大..
證明他們早就預備好自己,去迎接主降臨,即使他們不是用一個我們很想見到,或者很慣常的形式去出現在我們面前..
有時我們的主很謙卑,謙卑到一個程度我們未必能夠接受,我們太習慣一種有我們的生活模式,我們待人接物的禮節,沒有想過我們的主會打破這些我們已經習以為成的習俗傳統,能夠用一個全新的形式進入心中,為的是看我們的心是否全然的相信他..
再說的時候,一記經文很有趣,到最後一節,即是45節,他說:「這相信女子是有福的,因為主對她所說的話都要應驗.」.
我用我自己的中文對著希臘文重新翻譯一次給大家聽,因為通常「有福的」是擺頭的,「Blessed is who」,「有福的人是」,「有福的人是」,「有福的人是」,是誰呢?是這麼說的..
但通常中文是放在最後,所以我將「是有福的」放在最後,但中間那句是很陋的,就好像這樣說:「這個相信主對她所說的話必定會實現的女子是有福的.」.
整句意思就是這樣.首先,這是女子,這是一個很大的提示.考考大家,這個女子是誰呢?瑪麗亞,有沒有人說以利沙白呢?有的舉手.沒有,我就告訴你,其實很大機會是以利沙白.為什麼呢?我先解釋給你聽..
首先,你看著我們的華本,一定是瑪麗亞的,因為她是生括胡的.我們新標點華本,你看到42節「高聲喊著說」,這是以利沙白說的話,又42節「講講講」,說了45字.這是新標點華本的先賢,他為這個經文決定了是瑪麗亞..
但大家知道,如果讀過一些…可能曾牧師教過大家一些新的概論,你都明白,當時希臘文一排字沒有的,你不要說標點符號,你連剪輯那些字,你都要靠你很熟悉希臘文,古典希臘文那種知識,你才可以剪輯得準確.所以是有很多變數在當中的..
如果這個關括胡關在44節,而不是關在45節的話,這個話就不是以利沙白說的了.這個話就是盧家作為作者評論這個事情.大家明白我的意思嗎?.
盧家在直接跟著這個話說的時候,以利沙白的反應就是盧家說以利沙白這個女子,相信主的話對他說過必定實現的女子以利沙白是有福的..
她有多大福呢?那就要追溯一點,如果大家有聖經,可以回去慢慢看,整個盧家峰第一章.你記得整個故事其實源於第五節,開始她和她的先生薩加利亞,兩個都是祭司家庭,忠心事奉主..
你記得嗎?很有趣,他們兩個一二人,然後在第六節尾,第七節,指示沒有孩子,因為以利沙白不生育..
有些想像,兩個幼小讀神學,服侍主,想主有很大的祝福,想建立家庭,有小朋友,特別當時猶太人,其實很重要,母親,然後有孩子,這是整個家庭的一種很重要的尊貴價值在當中..
但是祈禱內討,也沒有,那怎麼辦呢?很多人事奉的弟兄姊妹,她仍然忠心屬事,所以一生投入在聖殿當中服侍..
突然間,同一個天使加巴列,就向薩加利亞說,跟著祝福完之後,就告訴她不要害怕第十三節,你的祈禱已經蒙英雲了,可能已經放下了,也沒有可能了..
按照我們的計劃,當然希望大家年輕力壯,有小朋友,有力氣,照顧他長大,然後就建立家庭..
上帝沒有開路,其實也接受了,然後慢慢就已經不想這件事了..
沒有想過年紀老邁才生孩子的,大家有沒有想過年紀老邁才生孩子的?不會的,是吧?.

$^{121}$有時候也會想,不知道有沒有體力來陪伴他,懷胎的時候,母親又擔心很多危險的變數,自己的身體機能也未必很好,很多擔心這些事情..
但是,上帝的角度不是這樣.上帝天使對他說,不要害怕,你的禱告已經蒙英雲了,要生一個兒子,然後就說施摘約翰的角色..
然後撒吉利亞的反應是不信,我和我太太年紀老邁豈能成事呢?於是就啞巴,大家知道,出來的時候說不出話,直到施摘約翰出生的時候,他才說話..
男人通常都是信心小一點..
撒吉利亞的反應是一直都不知道的,直到24-25節,這四天以後,他的妻子以利沙白懷孕了,就隱藏了五個月,說主在眷顧我的日子,這樣看待我,要把我在人間的羞恥除掉..
他也知道自己年老,也有很多不穩定的事情,可能等到自己的胎穩定了,我們通常是三個月,四個月,五個月就更加穩定,現肚子了,然後才出來見鄉親父老,告訴他們我懷孕了..
鄉親父老見到他的時候,哇,你真是好了,好像一籠糯米一樣,沒有兒子,很悲哀,其實那是一個祝福..
然後撒吉利亞的反應就是主眷顧他,看待他,要把他在鄉親父老,親戚面前的那種蒙羞除掉..
他一直都相信,他不是撒吉利亞那種一下子的反應就不能成事,他一直都是主說的話,他就自己預備自己,轉變心態,等待著這件事情發生..
我就說,信心不是自己推理出來的,相信主要是在最好的時間實現,這個就最難了..
你能不能夠撒吉利亞,二二三八,說這個時間是最好的呢?按我們人的推算,人的計劃,三十歲,轉到四十歲,生小朋友,你不會年紀再大,很多事情要計劃,又要計錢,又要計什麼,很多事情..
但如果從上帝的角度來說,不是這樣計的,我們能不能夠真正全心地相信呢?想想這件事情,他出生的孩子,施慈若瀚,要為耶穌開路,他不能夠比耶穌大很多,一出生就大一點點,你記得施慈若瀚成長的時候就在約旦河向猶太人施行悔改的洗,然後耶穌出來..
但問題是兩個母親的年紀差很遠,那你怎麼辦呢?難道要耶穌遷就你嗎?當然不會的,你當然是施上遷就耶穌的,你作為他的開路先鋒..
施慈若瀚要懷胎的時候,只能夠在這個年紀懷胎,成為耶穌基督的開路先鋒.他見到在上帝的角度當中,他生出來的孩子好得無比,因為他能夠成為耶穌開路,向世人宣講悔改的福音,讓耶穌出來的時候,更多人預備好自己,去聽聞福音,去接受耶穌..
在上帝的角度當中是最好的時間,不是按人的想法,而是按人的計劃當中是最好的時間..
這是最難的,我相信大家,我們很想神配合我們的時間,我們沒有想過我們要傳言要配合上帝的時間,但是你看到一個相信主的話必定要實現的女子,她傳言相信在配合上帝在生命當中的祝福大能的時間,她是真正有福的..
她有福不只是一些生命,很多的享受,她的祝福是在上帝的角度當中,她能夠有份參與,能夠有份在當中成為上帝角度的一個角色..
她看到這個屬靈的祝福是最為至寶,她能夠感受到生命的滿足,她是一個一直都服侍主的人,以致她能夠在這個時候,遇到一個人間當中不是好的時間,但她覺得這個真的是上帝最好的,她相信,她沒有計自己的年紀,沒有計自己將來孩子長大的時候,她和她先生有多少體力,陪伴孩子成長..
她知道,她就盡力去做,回應就可以..
沒有因環境困難,計算代價,反而願意去傳言的相信和信服..
這個就是一個信道的生命..
說到最後,行道的一方面的時候,我想起一節經文,這個經文,《雅各書》一章二十五節,這個這麼說的時候,都是一個有福的人,是不是?.
唯有詳細察看那使人自由的傳避的律法,並且在時常遵守的人,他不是聽了見,聽了就忘記,而是實前出來,因為他所作的是蒙福..
行道的,令他蒙福..
這個行道的生命,不是純粹聽了上帝的話,然後想想怎樣去應用,怎樣去在什麼事情上,去參加一些計劃,自己有一點點的付出,這麼簡單..
其實是啟動我們整個人,去產生了一個新的力量,去改變我們,帶來我們的根深,而我們的根深也在改變我們身處的地方的一些根深..
特別大家看到藍色那裡,用在我下面的英文的譯法更加貼近..
他這裡說的一個不是聽了就忘記,而是實人出來,其實如果看真一點,那個是說那個人自己,不是說做了什麼..
我們不要成為一個善忘的聆聽者,一個人,但是一個active doer,一個主動,一個積極,不是說我們要做很多事情,不是這樣,都要排到,由早七晚十一,整個時間變得密集,做到自己崩潰了,自己很枯竭那樣東西..
有一種好像瑪利亞和伊斯拉伯那樣,是一種隨時準備好那種active的狀態,active mode,我們有時候電話,電腦,我們不是sleep了,那種active mode的意思..
這種active mode是在背後有什麼指示給我的時候,上帝點擊哪個apps我們要做的,就去行那個..
我們不會手機停了,不會等了很久,不會sleep了,不會醒,是那個active的意思,那種積極,那種預備好自己,是這個意思..
一個事後順著後命的行動者,我們主一到一到,我們就會行動,就會回應..
好像瑪利亞那樣,隨即去山區找伊斯拉伯,伊斯拉伯一開門,隨即就見到我主的母親,喜約吾教來到我生命裡面,是很配合上帝的道,來到的時候,作出一個應該有的回應,說出應該說的話,是一種模式..
不是怎樣應用,不是怎樣想一個說話,是怎樣一個原則,而是整個人要投入,被根深改變..
剛才說到耶穌在瑪利亞的福中的胎兒成長,當然耶穌的童年,不多記載,大家記得嗎?只有一段,就是路加福第十二章最後,五十一至五十一節那段經文..
大家記得嗎?十二歲的耶穌,你也背了給我聽,在這裡陪家人過節,然後全部回到那裡,上了那裡,但走了三天,突然間不見了耶穌,然後就回去找,直至找到為止..
然後,瑪利亞和父親約瑟見到他的時候,就很生氣,為什麼要我們生氣,要找你呢?.
然後耶穌很厲害,一句話,回到瑪利亞,你豈不知道我以我父家的事為念魔?.
然後,這個就是接著的反應,他母親把一切的事都存在心裡,很有意思..

$^{161}$如果你的兒子十幾歲,或者小一點,你母親就算是自己之前的經歷,她頑皮,她反叛,你的反應如何?.
就算在別人面前把自己克制住,你內心也很顧慮,忍不住,你又火得來,那種反應..
很不容易的,因為胎兒的耶穌,你耶穌就動不動跳,到了成長的耶穌,特別到了十二歲青春期,反叛的耶穌,有自我建立的耶穌,他講的話,講出了自己想怎樣,當然他想怎樣的東西,是他跟著上帝的旨意而行..
但作為母親,看著他長大的時候,聽到一些不順耳的話,有什麼反應呢?.
你很難馬上覺得兒子是對的,是吧?沒有什麼可能的,很少的..
這個母親就是在不同的階段,對著一個道成肉身的救主,她要繼續學習怎樣應對,一直這樣應對..
第一次被兒子這樣得罪,她就將這些話反覆思量,留在心裡面,想想究竟是什麼意思..
她需要她自己很謙卑,不是按自己的身份地位權柄去為重要,而是留意道的說話,她的工作,看看自己怎樣去配合..
這個樣式就是一種行道生命的樣式,不是要自己做什麼,不是要取代道,不是要替天行道,是吧?我們經常這樣說..
或者是用我的意思,代替了道,是吧?.
現在耶穌還沒長大,還沒出來教導,他講的說話不用聽了,你還是小朋友,可以這樣說的..
沒有啊,因為他知道那個不是一個常人,是一個孩童的耶穌,都是上帝的道..
福中的道,都是上帝的道..
十二歲的耶穌,都是上帝的道..
釘十字架,都是上帝的道..
我們面對著同一個救主,我們的人的身份地位,就很考驗我們能不能夠讓一個這麼謙卑的救主形式在我們面前的時候,我們懂得去回應呢?.
有時候我自己都做了十幾年學生的工作,基本上每年都跟學生查經,或者跟他們分享..
當然有時候,都會有很多的牧者被我提醒,但是,偶爾,一些年輕的人,講的話,其實是在對我提醒..
有時候我自己都要好像瑪麗亞一樣,見著一個比我年輕,好像我提攜他,我教導他的人,他可以再講一句責心,挑戰著我的說話..
我聽到的時候,我有什麼反應呢?.
我立刻抹殺他,我立刻迴避他,立刻擠壓了他,還是能夠像瑪麗亞一樣,將他反覆思量放在心裡呢?.
這個很不容易..
就是一樣,道的採取一個樣式,不一定是我們期望一些權威人士,一定是一個很有經驗,一定是一個你覺得他一定是,你很期望已經在教導你聖經的話的人..
可以透過一些不經意的人物,一些小人物,一些在你之下的人物,來提醒你,上帝的道在當中,要提示我們..
所以行道是一種不住反思的過程,你才會行,就不會一味地衝,一味地聽完那句話,就在做,就做得到..
所以到最後總結的時候,我們怎麼回應呢?.
怎麼傳人,投入,這種由在道,信道,到行道,其實是一整體..
我們要行道的時候,我們一定要有上帝的道在心裡面,一定要相信它work,你才要行出來..
所以三者是不能夠分割的,你在著它的時候,你必然是相信,否則你做不到它,而且必然是要令你行出來..
你信的時候,你信的是什麼呢?不是信一些離開你的生命的道理,而是信一些在你裡面不住提醒你,在的道理..
而且你信的時候,全心相信,你必然影響你的生命,整個生命活現,行出來,這個上帝的話語..
所以三者是分割不了的,所以我的結論就是講這件事..
我們怎麼去回應呢?大概有三個程度..
可能按照弟兄姊妹,我們對上帝的話的熟絡,我們的關係是多麼緊密,或者我們是不是初信,還是信了很久..
可能我們初信的時候,我們上初信班,洗禮班,我們聽很多上帝的話,都似乎是一些個人的生活..
我們就學習,活出一個應有的行為樣式,這個聖經想教我們,我們要脫去舊我,要離開黑暗,進入光明..
有時候我們以前的習慣,舊思想,我們讓上帝的話常常在我們的生命當中指引我們,改變我們的行為習慣,這是第一個程度,第一個level..
但當我們生命已經過了初信,我們不能一生都停在初信的階段,其實我們要繼續往上行,我們的生命要提升自己,.
就到了第二個層次,就是人生的照明..
剛才說了,馬尼亞伊斯拉伯有一個獨特的生命的路要走,剛才我講過,我們每一個人都是一樣,這個就是calling的意思..

$^{201}$上帝呼召你的,上帝講的話,就像剛才一樣,天使向著馬尼亞講,向著撒加拉亞講,向著伊斯拉伯講,是獨特的,是關乎他的意思..
所以上帝的話語是很針對性的,也很明白我們的需要,很了解我們能夠怎樣回應他,補足我們的信心..
所以這就是人生的照明,問題是我們是不是好像伊斯拉伯和馬尼亞一樣,不懼困難呢?.
兩個都很大困難的,一個不夠青,要生小朋友,一個已經很老,生小朋友,都是很有困難的,整件事基本上都是很卡的..
但他們很多環境,很多人要交代,講多久都不明白,但為了上帝要他,他自己才繼續走下去,一步一步走下去,然後就逐步顯明,別人就接受到這件事情..
很多的環境因素,很多我們的擔憂,恐懼,照明就是要不讓這些東西佔據我們的內心,準備自己回應道的引領和指引..
最後一個層面就是最難的,不過也要講,就是不只是我們個人,而是整個社群性的制度,向導,關係都要被導去更新..
有時候我們容易留下一個個人的信仰,教會一個群體,有一個制度在這裡..
我們身處的工作環境,你的角色,你不知道是不是一個經理,一個主管,一個下屬也好,你的社群當中上上下下,有一個制度,有當中的價值觀,當中的文化,當中的人際關係..
有沒有想過上帝的導,透過你,都要改變這一切呢?.
這個就是最後想和大家分享,這個也是最重要,在我們生命當中,能夠讓到顯大的時候,能夠為上帝的國當中更切實地降臨在我們人世間中間..
這個最後,我自己的事工機構,都記念我們機構,教會的事工,中學生的事工,大學生的工作,文字的工作,和一些畢業生的職場信徒的工作..
都盼望教會和上帝的國度當中,都彼此同行,彼此紀念,一起在上帝的手中,有給我們的生命更多的偉仰,更多的明白上帝很大的工作..
在教會,在上帝的國度當中,透過不同的人,在做什麼呢?.
我們心靈都向上帝去慷放,更見到上帝更偉大的作為,我們一起得有禱告..
我見到兩個姐妹,她們都有她們很大的挑戰,很大的經歷,一個很大的使命,但我更加見到是一個再度順道行道的心靈..
她們沒有放大到她們面對的困難,也沒有給自己的情感,恐懼,擔心,完全的主導著她們生命前的決定..
她們將自己開放給你,全然讓道當中顯大,讓道當中成就你的心意,實現你的國度,透過你的生命..
她們有福的人,她們所經歷的福是很奇妙的,在上帝的國度當中,她們看到自己的角色的時候,她感到很大的福樂,很大的滿足..
拍賣中頂的姐妹們,都是學習成為一個再度行道,順道的生命,同樣嚮往著這份屬靈的福氣,.
能夠為著上帝的國度當中,我們個人生命的照明回應,教會,以至我們所處於的工作環境,社會,能夠徹底被你的道去盛化,去更新..
盼望你祝福使用弟兄姊妹,成為你國度的出口..
多謝您,聽我們同心禱告,奉主明求..
阿們..
(影片完結).
拜拜!.
\newpage



\section{路加福音 1:39-45}
\label{sec:vW2xmW8knK8}
\textbf{2025.05.04《載道,信道,行道》 路加福音1\_39-45 講員 : 蔣文忠傳道}
\newline
\newline
連結: \href{https://youtube.com/watch?v=vW2xmW8knK8}{\texttt{https://youtube.com/watch?v=vW2xmW8knK8}} ~~~~ 語音日期: 2025-05-13
\newline
\newline
\hyperref[sec:o8E8B2LqAe4]{\small{< < < PREV SERMON < < <}}
~
\hyperref[sec:index_chronic]{\small{[返順時目]}}
~
\hyperref[sec:index_scriptual]{\small{[返順卷目]}}
~
\hyperref[sec:KCtIY5_ZgcE]{\small{> > > NEXT SERMON > > >}}
\newline
\newline
路加福音 1:39-45
\newline
\begin{longtable}{cl}
\hline
\hline
章節 & 經文 (和合本修訂版)\\
\hline
1:39 & \begin{tabularx}{0.7\textwidth}{X} 在那些日子，馬利亞起身，急忙前往山區，來到猶大的一座城， \end{tabularx} \\ \\ \relax
1:40 & \begin{tabularx}{0.7\textwidth}{X} 進了撒迦利亞的家，向伊利莎白問安。 \end{tabularx} \\ \\ \relax
1:41 & \begin{tabularx}{0.7\textwidth}{X} 伊利莎白一聽到馬利亞問安，所懷的胎就在腹裡跳動。伊利莎白被聖靈充滿， \end{tabularx} \\ \\ \relax
1:42 & \begin{tabularx}{0.7\textwidth}{X} 高聲喊著說：「你在婦女中是有福的！你所懷的胎也是有福的！ \end{tabularx} \\ \\ \relax
1:43 & \begin{tabularx}{0.7\textwidth}{X} 我主的母親到我這裡來，為何這事臨到我呢？ \end{tabularx} \\ \\ \relax
1:44 & \begin{tabularx}{0.7\textwidth}{X} 因為你問安的聲音一入我耳，我腹裡的胎就歡喜跳動。 \end{tabularx} \\ \\ \relax
1:45 & \begin{tabularx}{0.7\textwidth}{X} 這相信的女子是有福的！因為主對她所說的話都要應驗。」 \end{tabularx} \\ \\
[1ex]
\hline
\hline
\end{longtable}
$^{1}$(主持人:很感恩在此大家中間一齊去敬拜,一齊去聽上帝的話語,剛才讀的經文可能大家有時會在聖誕節,或者有時會輕讀過,未必會很留心這段經文,一會兒我們會詳細看一下這段經文當中,特別是兩個姐妹,兩個主角,.
在與道,與聖道要降臨,與道成肉身的奧秘發生的時候,他們的回應,他們的關係是怎樣?.
大家可能都知道我們忙碌了復活節和壽難節,其實在教會的節期,這段時間是復活節期,有七個主日,直到所謂的聖靈降臨節,小人傳是第二章,.
當中聖靈降臨在使徒的生命中間,耶穌其實之前陪伴門徒多少天?.
小人傳第一章,盤踴四十天,然後升天,再過一天,聖靈降臨在使徒的身上,在復活節期裡面,特別是拉丁教父和特土良人都說,這是一個最喜樂的空間,.
這段五十天,這七個主日,其實是繼續讓我們在教會中回味和感受復活帶來的喜樂,.
讓我們學會了食活的能力,帶來我們生命的改變,整個人世間的局面,其實不是讓我們世界當權的人帶動,而是讓我們世界掌權的主繼續同在,引領著我們人類得到生命真正的拯救,.
所以我們可以想想,我們不要停下來,我們懷著什麼心情去期待聖靈,即復活的能力,給我們的改變呢?.
我們是不是一直在經歷復活帶來給我們這種新的生命,新的機遇,新的盼望呢?.
這種更新,這種新機,就是在復活節當中賜給我們,當然我們慶祝的節期,或者我們能夠對生命有一種新的改變,有新的盼望,.
那個根基,那個力量,其實是源於,好像今天所說的上帝的道,在上帝的道與人同在,上帝的道復活帶給人一種力量的改變,成為我們一個生命可以持續往前走,就算遇上不同的艱難,不同的挑戰,其實我們都不需要懼怕,不需要憂慮..
所以今天我們看這個經文,其實是想跟大家一起去思考,究竟我們弟兄姊妹,我們教會,其實我們與道之間有一種怎樣的關係呢?.
我今天的題目是講再道,進道,和行道,好像有三種不同的樣式,模式,和道的關係,但是我最後都講了其實是同一樣的東西,希望大家一起分享的時候,你會明白得到..
我們從剛才讀的經文,路加封第1章39到45節,這裡是一個短短的記載,但是它是一個很重要,也是在表示瑪利亞和以利沙白的時候都忽略了的老姐妹,當時來說,就是她與道之間的關係..
我們通常會將重點放在瑪利亞那裡,但是有時候因為我們不想太高舉瑪利亞,所以我們不會看很多關於瑪利亞在過程中的生命經歷是怎樣..
這故事很簡單,前面就講到天使加巴利與瑪利亞報信,大家應該很熟悉,你會懷胎,聖靈感染,生出我們的救主.瑪利亞就是一個心懷相信,回應願主的話,成就在她生命當中..
她第一個反應就是在這裡,39節,就是去找以利沙白,跟她分享這個情況,是不是很特別?其實是一個很獨特的,在這段時間裡面,瑪利亞與以利沙白的一個很獨特的團契..
瑪利亞沒有什麼人明白她,她還沒有出嫁,大家明白,天使告訴她,聖靈要讓你福中懷胎,然後要繼續長大成人,成為一個救主..
這就是所謂的道成肉身,一種很立體,一種我們很細心去想的時候,是一個我們想的時候,瑪利亞的角色,做母親的角色是非常重要的..
所以她一直相信,一直思考,究竟是怎麼發生的,其實還沒有發生的,剛剛天使說完之後,究竟其後九個月的懷胎日子,然後胎兒出生,長大成人的過程是怎麼樣呢?其實瑪利亞也沒有路燭,沒有一個男徒知道怎麼做,只是一步一步學習怎麼去回應..
所以她找了以利沙白去分享她的一個特別的經歷,可能有一點忐忑,但是憑信繼續往前走,因為天使跟她說,你的親戚都懷胎,當時已經懷了六個月,她是一個很獨特的兒子,都即將透過以利沙白,一會兒我會講多一點她的故事而出世..
所以當她相遇的時候,就有一個經文,以利沙白跟瑪利亞說,你福中的胎兒是有福的..
瑪利亞懷著一個懷道的經歷,正在預備道成肉身,有沒有想過,有時候我們一看耶穌的生命,就想到祂,最多可能是受洗,受試探,然後出來傳道..
有沒有細心思考過,原來我們的救主定義道成肉身的時候,其實是完完全全地取了一個人的樣子,不是一出來就成人,完全就是一個母親的福中胎兒出世,然後學習走路,這樣成長的過程..
其實,她取了一個最脆弱,最不可避免的形式,大家可能都聽過很多悲劇,或者一些很不開心的事情,有時候很多母親懷胎,但很不想要的小朋友,她就把他放下,或者出世之後,將他棄掉等等..
我們也聽過很多時候,在社會當中,很多故事,很多我們很不如意的事情.孩子福中的胎兒是無法作主的,她是沒有能力的,她完全依靠著她母親的母體,給她什麼營養,怎樣照料她,她就怎樣成長..
大家在座的姐妹,有小朋友,你見過,有過孫子等等,就明白這個過程,是很不容易,但也是一個互相影響的過程..
瑪利亞就要承擔著,為著她的孩子成為救主,擔起母親的角色,是一個很獨特的遺訓,要在瑪利亞福中去承認..
甚至是一見面的時候,以利沙白土宗的施西漢一見到救主,救主不一定是我們那時候的畫像,一定是長頭髮,有鬍鬚的耶穌,很慈悲,慈愛的耶穌..
有沒有想過一個胎兒裡面,都能夠回應一個胎兒的耶穌?有時候我們母親知道,如果懷孕的時候,她踢我,沒有什麼可以做的,最多是跳躍,在那裡活動..
現在在以利沙白母親的施西漢,六個月的施西漢,她就在那裡興奮跳動,以致令以利沙白,原來我面對著的瑪利亞福中,就是她自己生命的救主,對著胎兒,她可以宣認她是她生命的主..
是很特別的,正正是這段時間才會發生的,她成為了一個人,成為了一個呼召門徒,已經是另一個故事..
到我們今天,我們常常去,剛才說,我們從聖誕節到壽難節到父節,我們都是在想一個成人形象的耶穌,很多時候..
有沒有想過,這個道成肉身,原來是一種這麼徹底,這麼將自己放在人的手中,讓自己的生命經歷一個家庭當中的成長..
大家可以繼續去思想,去默想,耶穌這種道成肉身的奧秘,那種徹底的謙卑,去到什麼程度?.
我們特別要看瑪利亞這個生命的樣式的時候,她說一個再度的生命,是怎樣去呈現給我們看到呢?.
第一步就是,剛才說的經文,39節,瑪利亞起身急忙前往山區,為什麼這麼急呢?很快,立即,她沒有猶豫..
她這個承載著,懷著道的生命,她沒有為自己很多憂心,很多考慮,她即時就回應著,雖然她在途中,她即時就用行動來回應這個救主..
這個即時的,是表示一種信服的回應,是全然接受一個新的身份,她都不知道怎樣,要成為救主的母親..
大家有沒有留意,其實瑪利亞的角色,在初代教會是一個很爭議性的討論..

$^{41}$當然,這個討論不是因為瑪利亞,而是因為耶穌基督的神人異性,引致到她母親都牽連在內,其實是這個意思..
究竟瑪利亞是所謂人之母,還是神之母,就是這個討論..
當然,應該是神之母,因為她母腹中的是完全的神,所以就爭論這些問題..
待會兒,你看到以利沙白的宣講,你就知道,宣認她..
瑪利亞是即時,殷然,不是殷然,戰戰兢兢可能都是,接受一個新的身份,沒有人能夠定義的,是神正在定義她這個身份,新的挑戰,新的方向..
大家有沒有想過,上帝呼召你,你做的事情是全新的,你不是在重複前人做的事情,因為都是屬於你,給你一個獨特的使命和任務..
這是很神奇,這是很奇妙的事情..
在上帝眼中,每一個人的生命都是很獨特,在上帝國當中,扮演的角色都是很有獨特的一個角色,一個恩賜..
沒有了你是不行的,你當中發揮的時候,就能夠祝福其他人..
可能聽了這麼久,你會說,瑪利亞又是神之母,又道成肉身,成為她身體中間的一個胎兒..
我們不是,我們不是瑪利亞,那跟我們有什麼關係呢?.
當我們思考這個位置的時候,這個再度的生命的時候,其實是跟我們息息相關的..
大家記得耶穌有殺種的比喻,有落在淺土上面,一會兒被烏鴉淡走了,我們被很多世務纏繞著,忘記上帝的道..
也有落在好土裡面,總之能夠健康地紮根,一直茁壯成長,結出很多的果子..
這個土壤就是我們的生命,其實是一樣的..
上帝的道,好像今天孫子說的,我們聽道,我們看聖經,落在我們心靈的土壤裡面..
我們是一個什麼樣的土壤呢?質量,質地是怎樣的呢?.
我們有很多的淺土,很多的荊棘,不能夠讓道在當中佔一個位置,影響著我們整個生命..
還是我們預備好自己,是一個很接受的,很願意領受上帝話語的土壤呢?.
以至上帝說,落在我們生命當中,我們是能夠欣然接受,然後讓它可以發芽生長呢?.
所以讓道在生命中的成長,是我們每個弟兄姊妹都能夠學習得到瑪利亞這個再度的生命的樣式..
我接著說杜伊尼撒白,我想用更多的時間說他..
因為我們每個人都信道,再度好像很難理解,怎樣承載著呢?.
但是我們有福音,有上帝話語,我們相信,我們大家都會說,很容易的,我應該可以在當中更多的學習..
而杜伊尼撒白是一個很好的樣式,讓你去思考,怎樣來自於這個信道的生命..
大家經常聽到,當瑪利亞急忙到山區找他探望的時候,他的反應就是,四十三夕,我主的母親竟然在新日本..
其實不緊要,大家和日本的和諧差不多,不過我覺得日本的神髓更加到位..
你竟然來到我這裡?這件事怎會來到我的生命當中?.
很有趣吧?你們記住,伊尼撒白是年紀老邁,當時,我想起碼五六十歲,懷孕,很神奇,是一個神蹟的事..
瑪利亞還未成年,未出嫁,十幾二十歲,相差一截,在年紀上..
伊尼撒白是長輩,想想,大家在拜年的時候,你的長輩去後輩的家,還是後輩去長輩的家?理所當然的,我們是有大食大,中強的文化..
一定是去最年長的家,我們全部的子子孫孫都去那裡,一定是這樣,一定是拜年後再去..
所以,伊瑪利亞向一個懷胎六個月的前輩,一個親戚,去拜會他,送一些燕窩,送一些什麼給他吃,理所當然的..
不需要這麼大反應的,但是伊尼撒白是倒轉了整個人倫的關係..
他見了一個十幾歲的妹妹,他說,我主任母親,竟然來到我姐女,這件事怎會發生呢?.
正常人不會這樣說話的,就算你見到後輩來到,你很尊重他,你也不會這樣說話的..
這是一個信道,讓他見到在面前的瑪利亞,我主任母親,已經回答了神之母的問題,其實伊尼撒白一句已經回答了,KO了問題..
我主任當然是神,不是母親厲害,而是他肚中的我主,使這個人叫我主任母親..
大家明白嗎?不是我主任母親,而是我主任才是最首要,伊尼撒白要回應的對象,不過只是對著一個腹中的胎兒這樣說而已..
為什麼會來到我這裡?為什麼會由我主來拜會我?所以他覺得很不配,很受寵若驚的反應..

$^{81}$是不是信心很大?他見到整個在神國當中的秩序,雖然我年紀大一點,但我是一個排在幕後的角色,.
對著瑪利亞,對著她母親腹中的胎兒,對他來說是一個更加尊貴的角色..
所以他覺得怎麼會發生這樣的事情?萬萬不能,很不配,你來到我家中..
更加反映他早已預備好,主動迎接到的降臨..
我們要問,如果今天好像各國敲門,或者像啟示所說的,耶穌在我們心門外等候我們的時候,我們又會作出什麼反應,回應上主的導降臨呢?.
這是很值得去想的.其實常常都在我們耳邊,常常都在我們心門外,祂想進入我們的生命當中,進入我們的空間裡面..
我們是不是都像以利沙白那樣預備好,欣然地一說就說中一個應該說的話,還是祂還猶豫?.
我們做了很多的盤算,很多的考慮,都還沒有認出主在我們面前..
你見到很多,剛才說復活的時候,每個人都認不出主,記不記得?.
女士又認不出主,門徒也認不出主,直到耶穌自己顯現的時候,說了祂的話的時候,才認出主來..
我們不能夠一見到主就認出主,可能不是用一個我們常慣的面孔見面..
福中的胎兒一定不是常見的,沒有想過.以利沙白是全然相信,已經導成了新告庇..
她才見到瑪麗的母親,她的胎主,可能剛剛開始就懷胎,不過師姐提醒了她,福中的胎兒跳躍,讓她知道這個母親,這個眼前的親戚是一個主的母親..
她即時就作出一個反應,作出一個相信,作出一個信服..
我們有時考慮很多,盤算很多,猶豫很多,但是你見到,無論是瑪麗的即時去到山區,或者是一開門,以利沙白一見到瑪麗就說,我主的母親竟然來到我家中..
這些反應證明了他們信心很大,證明他們一早預備好自己去迎接主降臨,即使他們不是用一個我們很想見到,或者很慣常的形式去出現在我們面前..
有時我們的主很謙卑,謙卑到一個程度我們未必能夠接受,反而我們太習慣一種,有我們的生活模式,我們待人接物的一種禮節..
沒有想過主會打破這些很多,這些我們已經習以為成的習俗傳統,能夠用一個全新的形式進入心中,為的是看我們的心是不是全然的相信他..
再說的時候,一句經文很有趣,到最後一節,45節,.
「這相信女子是有福的,因為主對她所說的話都要應驗了.」.
我用我自己的中文對著希臘文重新翻譯一次給大家聽,因為通常「有福的」是擺頭的,就是「Blessed is who」,通常是這樣說話的..
「有福的人是有福的」是誰呢?就是這樣說的,但通常中文是放在最後,所以我就將「是有福的」放在最後..
但中間那句就很漏漏,就好像這樣說,「這個相信主對她所說的話必定會實現的女子.」.
大家明白嗎?這個相信主對她所說的話必定會實現的女子是有福的,整句意思是這樣..
首先,這是女子,這是一個很大的提示,考考大家,這個女子是誰呢?.
瑪麗亞,有沒有人說以利沙白呢?有的舉手.沒有,我就告訴你,其實很大機會是以利沙白,為什麼呢?先和你解釋一下..
首先,你看著我們的華本,就一定是瑪麗亞的,因為她是生括胡的,我們新標點華本,你看到42節「高聲喊著說」,以利沙白說的話,又42節「講講講」,說了45字..
這個是新標點華本的先賢,他為這個經文決定了是瑪麗亞..
但大家知道,如果讀過一些…可能曾牧師教過大家一些新的概論,你都明白,當時希臘文一排字沒有的,你不要說標點符號,你連剪輯那些字,你都要靠你很熟悉希臘文,古典希臘文那種知識,你才可以剪輯得準確,所以是很多變數在當中的..
如果這個關括胡關在44節,不是關在45節的話,這個話就不是以利沙白說的了,這個話就是盧家作為作者評論這個事情,大家明白我的意思嗎?.
盧家在直接跟著這個話說的時候,現在以利沙白的反應,盧家就說這個以利沙白,這個女子,相信主的話對他說過必定實現的女子以利沙白是有福的..
有多大福呢?那就要追溯一點,如果大家有聖經,可以回去慢慢看,整個盧家峰第一章,你記得整個故事,其實源於第五節,開始她和她的先生薩加利亞,兩個都是祭司家庭,忠心事奉主,你記得,很有趣,他們兩個二人,然後第六節尾,第七節,指示沒有孩子,因為以利沙白不生育..
有些想像,兩個幼小讀神學,服侍主,想主有很大的祝福,想建立家庭,有小朋友,特別當時猶太人,其實很重要,記得,母親,然後有孩子,那個是整個家庭的一種很重要的尊貴價值在當中,但是其中的內討都沒有,那怎麼辦呢?.
很多事奉的弟妹姊妹,仍然忠心屬事,所以一生投入在聖殿當中服侍,突然間同一個天使加把年,向薩加利亞說,跟著祝福完之後,就告訴她不要害怕第十三節,你的祈禱已經蒙英雲了,可能已經放下了,也都沒有可能了..
放下了計劃,當然希望大家年輕力壯,有小朋友,有力氣,照顧他長大,然後就建立家庭,上帝沒有開路,其實都接受,然後慢慢就已經不想那件事了,沒有想過年紀老邁才生小孩,大家有沒有想過年紀老邁才生小孩?不會的,是吧?.
有時都會覺得,不知道有沒有體力去陪伴他,懷胎的時候母親又擔心很多危險的變數,自己的身體機能未必很好,很多擔心這些事情,但是上帝的角度不是這樣,上帝天使對他說,不要害怕,你的禱告已經蒙英雲了,要生一個兒子,然後就說施摘約翰的角色,然後薩加利亞的反應是不信的,是吧?.
我和我太太年紀老邁,豈能成事呢?於是就喇叭,大家知道,出來的時候說不出話,直到施摘約翰出生的時候,他才說話,男人通常都是信心小一點,是吧?.
薩加利亞的反應是一直都不知道的,直到24-25節,這四天以後,他的妻子以利沙白懷孕了,就隱藏了五個月,說主在眷顧我的日子,這樣看待我,要把我在人間的羞恥除掉,是吧?.
他也知道自己年老的,是吧?也有很多不穩定的事情,可能等到自己的胎穩定了,我們通常是三個月,四個月,是吧?五個月就穩定一點,懷孕了,然後才出來見鄉親父老,告訴他們我懷孕了,這樣..
那些鄉親父老見到他的時候,哇,你真是好了,好像一籠糯米一樣,沒有兒子,很悲哀,其實那是一個祝福,接著薩加利亞的反應就是主眷顧他,看待他,除掉他在這些鄉親父老,親戚面前的那種蒙羞..

$^{121}$他一直都相信,他不是薩加利亞那種一下子的反應就不能成事,他一直都是,可能是主說的話,他就自己預備自己,轉變心態,等待著這件事情發生..
我就說,信心不是自己推理出來的,相信主要是在最好的時間實現,這個就最難了,是不是?你不夠薩加利亞,二二三百萬,這個時間是最好的,是吧?.
按我們人的推算,人的計劃,三十歲,最多四十歲,生小朋友,你不會年紀再大,很多東西要計劃,又要計錢,又要計什麼計什麼,很多東西..
但如果從上帝的角度來說,不是這樣計的,我們能不能夠真正全心地相信呢?想想這件事情,他出生的孩子,施慈若瀚,要為耶穌開路,他不能夠比耶穌大很多,一出生就大一點點,你記得施慈若瀚成長的時候就在約旦河向猶太人施行悔改的洗,然後耶穌出來..
但問題是兩個母親的年紀差很遠,是不是?那你怎麼辦?難道要耶穌遷就你嗎?不會的,當然是你施上遷就耶穌,你作為他的開路先鋒..
施慈若瀚要懷胎的時候,只能夠在這個年紀懷胎,成為耶穌基督的開路先鋒.他見到在上帝的角度當中,他生出來的孩子,好得無比,因為他能夠成為耶穌開路,向世人宣講悔改的福音,讓耶穌出來的時候,更多人預備好自己,去聽聞福音,去接受耶穌..
在上帝的角度當中是最好的時間,不是按人的想法,而是按人的計劃當中是最好的時間..
這是最難的,我相信大家,我們很想神配合我們的時間,我們沒有想過我們要傳言要配合上帝的時間,但是你看到一個相信主的話必定要實現的女子,她傳言相信在配合上帝在生命當中的祝福大能的時間,她是真正有福的..
她有福不只是一些生命,很多的享受,她的祝福是在上帝的角度當中,她能夠有份參與,能夠有份在當中成為上帝角度的一個角色..
她看到這個屬靈的祝福是最為至寶,她能夠感受到生命的滿足,她是一個一直都服侍主的人,以致她能夠在這個時候,遇到一個人間當中不是好的時間,但她覺得這個真的是上帝最好的,她相信,她沒有計自己的年紀,沒有計自己將來孩子長大的時候,她和她先生有多少體力,陪伴孩子成長..
她知道,她就盡力去做,回應就可以..
沒有因環境困難,計算代價,反而願意去傳言的相信和信服..
這個就是一個信道的生命..
說到最後,行道的一方面的時候,我想起一節經文,這個經文,《雅各書》一章二十五節,這個這麼說的時候,都是一個有福的人,是不是?.
唯有詳細察看那使人自由的傳避的律法,並且在時常遵守的人,他不是聽了見,聽了就忘記,而是實前出來,因為他所作的是蒙福..
行道的,令他蒙福..
這個行道的生命,不是純粹聽了上帝的話,然後想想怎樣去應用,怎樣去在什麼事情上,去參加一些計劃,自己有一點點的付出,這麼簡單..
其實是啟動我們整個人,去產生了一個新的力量,去改變我們,帶來我們的根深,而我們的根深也在改變我們身處的地方的一些根深..
特別大家看到藍色那裡,用在我下面的英文的譯法更加貼近..
他這裡說的一個不是聽了就忘記,而是實人出來,其實如果看真一點,那個是說那個人自己,不是說做了什麼..
我們不要成為一個善忘的聆聽者,一個人,但是一個active doer,一個主動,一個積極,不是說我們要做很多事情,不是這樣,都要排到,由早七晚十一,整個時間變得密集,做到自己崩潰了,自己很枯竭那樣東西..
有一種好像瑪利亞和伊斯拉伯那樣,是一種隨時準備好那種active的狀態,active mode,我們有時候電話,電腦,我們不是sleep了,那種active mode的意思..
這種active mode是在背後有什麼指示給我的時候,上帝點擊哪個apps我們要做的,就去行那個..
我們不會手機停了,不會等了很久,不會sleep了,不會醒,是那個active的意思,那種積極,那種預備好自己,是這個意思..
一個事後順著後命的行動者,我們主一到一到,我們就會行動,就會回應..
好像瑪利亞那樣,隨即去山區找伊斯拉伯,伊斯拉伯一開門,隨即就見到我主的母親,喜約吾教來到我生命裡面,是很配合上帝的道,來到的時候,作出一個應該有的回應,說出應該說的話,是一種模式..
不是怎樣應用,不是怎樣想一個說話,是怎樣一個原則,而是整個人要投入,被根深改變..
剛才說到耶穌在瑪利亞的福中的胎兒成長,當然耶穌的童年,不多記載,大家記得嗎?只有一段,就是路加福第十二章最後,五十一至五十一節那段經文..
大家記得嗎?十二歲的耶穌,你也背了給我聽,在這裡陪家人過節,然後全部回到那裡,上了那裡,但走了三天,突然間不見了耶穌,然後就回去找,直至找到為止..
然後,瑪利亞和父親約瑟見到他的時候,就很生氣,為什麼要我們生氣,要找你呢?.
然後耶穌很厲害,一句話,回到瑪利亞,你豈不知道我以我父家的事為念魔?.
然後,這個就是接著的反應,他母親把一切的事都存在心裡,很有意思..
如果你的兒子十幾歲,或者小一點,你母親就算是自己之前的經歷,她頑皮,她反叛,你的反應如何?.
就算在別人面前把自己克制住,你內心也很顧慮,忍不住,你又火得來,那種反應..
很不容易的,因為胎兒的耶穌,你耶穌就動不動跳,到了成長的耶穌,特別到了十二歲青春期,反叛的耶穌,有自我建立的耶穌,他講的話,講出了自己想怎樣,當然他想怎樣的東西,是他跟著上帝的旨意而行..
但作為母親,看著他長大的時候,聽到一些不順耳的話,有什麼反應呢?.
你很難馬上覺得兒子是對的,是吧?沒有什麼可能的,很少的..
這個母親就是在不同的階段,對著一個道成肉身的救主,她要繼續學習怎樣應對,一直這樣應對..
第一次被兒子這樣得罪,她就將這些話反覆思量,留在心裡面,想想究竟是什麼意思..
她需要她自己很謙卑,不是按自己的身份地位權柄去為重要,而是留意道的說話,她的工作,看看自己怎樣去配合..

$^{161}$這個樣式就是一種行道生命的樣式,不是要自己做什麼,不是要取代道,不是要替天行道,是吧?我們經常這樣說..
或者是用我的意思,代替了道,是吧?.
現在耶穌還沒長大,還沒出來教導,他講的說話不用聽了,你還是小朋友,可以這樣說的..
沒有啊,因為他知道那個不是一個常人,是一個孩童的耶穌,都是上帝的道..
福中的道,都是上帝的道..
十二歲的耶穌,都是上帝的道..
釘十字架,都是上帝的道..
我們面對著同一個救主,我們的人的身份地位,就很考驗我們能不能夠讓一個這麼謙卑的救主形式在我們面前的時候,我們懂得去回應呢?.
有時候我自己都做了十幾年學生的工作,基本上每年都跟學生查經,或者跟他們分享..
當然有時候,都會有很多的牧者被我提醒,但是,偶爾,一些年輕的人,講的話,其實是在對我提醒..
有時候我自己都要好像瑪麗亞一樣,見著一個比我年輕,好像我提攜他,我教導他的人,他可以再講一句責心,挑戰著我的說話..
我聽到的時候,我有什麼反應呢?.
我立刻抹殺他,我立刻迴避他,立刻擠壓了他,還是能夠像瑪麗亞一樣,將他反覆思量放在心裡呢?.
這個很不容易..
就是一樣,道的採取一個樣式,不一定是我們期望一些權威人士,一定是一個很有經驗,一定是一個你覺得他一定是,你很期望已經在教導你聖經的話的人..
可以透過一些不經意的人物,一些小人物,一些在你之下的人物,來提醒你,上帝的道在當中,要提示我們..
所以行道是一種不住反思的過程,你才會行,就不會一味地衝,一味地聽完那句話,就在做,就做得到..
所以到最後總結的時候,我們怎麼回應呢?.
怎麼傳人,投入,這種由在道,信道,到行道,其實是一整體..
我們要行道的時候,我們一定要有上帝的道在心裡面,一定要相信它work,你才要行出來..
所以三者是不能夠分割的,你在著它的時候,你必然是相信,否則你做不到它,而且必然是要令你行出來..
你信的時候,你信的是什麼呢?不是信一些離開你的生命的道理,而是信一些在你裡面不住提醒你,在的道理..
而且你信的時候,全心相信,你必然影響你的生命,整個生命活現,行出來,這個上帝的話語..
所以三者是分割不了的,所以我的結論就是講這件事..
我們怎麼去回應呢?大概有三個程度..
可能按照弟兄姊妹,我們對上帝的話的熟絡,我們的關係是多麼緊密,或者我們是不是初信,還是信了很久..
可能我們初信的時候,我們上初信班,洗禮班,我們聽很多上帝的話,都似乎是一些個人的生活..
我們就學習,活出一個應有的行為樣式,這個聖經想教我們,我們要脫去舊我,要離開黑暗,進入光明..
有時候我們以前的習慣,舊思想,我們讓上帝的話常常在我們的生命當中指引我們,改變我們的行為習慣,這是第一個程度,第一個level..
但當我們生命已經過了初信,我們不能一生都停在初信的階段,其實我們要繼續往上行,我們的生命要提升自己,.
就到了第二個層次,就是人生的照明..
剛才說了,馬尼亞伊斯拉伯有一個獨特的生命的路要走,剛才我講過,我們每一個人都是一樣,這個就是calling的意思..
上帝呼召你的,上帝講的話,就像剛才一樣,天使向著馬尼亞講,向著撒加拉亞講,向著伊斯拉伯講,是獨特的,是關乎他的意思..
所以上帝的話語是很針對性的,也很明白我們的需要,很了解我們能夠怎樣回應他,補足我們的信心..
所以這就是人生的照明,問題是我們是不是好像伊斯拉伯和馬尼亞一樣,不懼困難呢?.
兩個都很大困難的,一個不夠青,要生小朋友,一個已經很老,生小朋友,都是很有困難的,整件事基本上都是很卡的..
但他們很多環境,很多人要交代,講多久都不明白,但為了上帝要他,他自己才繼續走下去,一步一步走下去,然後就逐步顯明,別人就接受到這件事情..
很多的環境因素,很多我們的擔憂,恐懼,照明就是要不讓這些東西佔據我們的內心,準備自己回應道的引領和指引..
最後一個層面就是最難的,不過也要講,就是不只是我們個人,而是整個社群性的制度,向導,關係都要被導去更新..
有時候我們容易留下一個個人的信仰,教會一個群體,有一個制度在這裡..

$^{201}$我們身處的工作環境,你的角色,你不知道是不是一個經理,一個主管,一個下屬也好,你的社群當中上上下下,有一個制度,有當中的價值觀,當中的文化,當中的人際關係..
有沒有想過上帝的導,透過你,都要改變這一切呢?.
這個就是最後想和大家分享,這個也是最重要,在我們生命當中,能夠讓到顯大的時候,能夠為上帝的國當中更切實地降臨在我們人世間中間..
這個最後,我自己的事工機構,都記念我們機構,教會的事工,中學生的事工,大學生的工作,文字的工作,和一些畢業生的職場信徒的工作..
都盼望教會和上帝的國度當中,都彼此同行,彼此紀念,一起在上帝的手中,有給我們的生命更多的偉仰,更多的明白上帝很大的工作..
在教會,在上帝的國度當中,透過不同的人,在做什麼呢?.
我們心靈都向上帝去慷放,更見到上帝更偉大的作為,我們一起得有禱告..
我見到兩個的姐妹,她們都有她們很大的挑戰,很大的經歷,一個很大的使命..
但我更加見到是一個再度信道行道的生命,她們沒有放大到一些她們面對的困難,也沒有給自己的情感,恐懼,擔心,完全的主導著她們生命前的決定..
她將自己開放給你,傳言讓道當中顯大,讓道在當中成就你的心意,實現你的國度,透過你的生命..
她們有福的人,她們所經歷的福是很奇妙的,在上帝的國度當中,她見到自己的角色的時候,她感到很大的福樂,很大的滿足..
帕瓦頌頂的姐妹們都是學習成為一個再度行道,信道的生命,同樣嚮往著這份屬靈的福氣,.
能夠為著上帝的國度當中,我們個人生命的照明回應,教會,以至我們所處於的工作環境社會,能夠徹底被你的道去聖化,去更新..
盼望你祝福使用弟兄姊妹,成為你國度的出口..
多謝您,聽我們同心禱告,奉主明求..
(影片完結).
拜拜!.
\newpage



\section{希伯來書 13:7-13}
\label{sec:KCtIY5_ZgcE}
\textbf{2025.05.11《祂的故事,我的故事》 希伯來書 13\_7-13 講員 : 曾裔貴牧師}
\newline
\newline
連結: \href{https://youtube.com/watch?v=KCtIY5_ZgcE}{\texttt{https://youtube.com/watch?v=KCtIY5\_ZgcE}} ~~~~ 語音日期: 2025-05-13
\newline
\newline
\hyperref[sec:vW2xmW8knK8]{\small{< < < PREV SERMON < < <}}
~
\hyperref[sec:index_chronic]{\small{[返順時目]}}
~
\hyperref[sec:index_scriptual]{\small{[返順卷目]}}
~
\hyperref[sec:_TgXxd1W4v8]{\small{> > > NEXT SERMON > > >}}
\newline
\newline
希伯來書 13:7-13
\newline
\begin{longtable}{cl}
\hline
\hline
章節 & 經文 (和合本修訂版)\\
\hline
13:7 & \begin{tabularx}{0.7\textwidth}{X} 從前引導你們、傳神的道給你們的人，你們要記念他們，效法他們的信心，回顧他們為人的結局。 \end{tabularx} \\ \\ \relax
13:8 & \begin{tabularx}{0.7\textwidth}{X} 耶穌基督昨日、今日，一直到永遠，是一樣的。 \end{tabularx} \\ \\ \relax
13:9 & \begin{tabularx}{0.7\textwidth}{X} 你們不要被種種怪異的教訓勾引了去，因為人的心靠恩典得堅固才是好的，並不是靠飲食。那在飲食上用心的，從來沒有得到益處。 \end{tabularx} \\ \\ \relax
13:10 & \begin{tabularx}{0.7\textwidth}{X} 我們有一祭壇，上面的祭物是那些在會幕中供職的人無權可吃的。 \end{tabularx} \\ \\ \relax
13:11 & \begin{tabularx}{0.7\textwidth}{X} 因為牲畜的血被大祭司帶入至聖所作贖罪祭，牲畜的體卻在營外燒掉。 \end{tabularx} \\ \\ \relax
13:12 & \begin{tabularx}{0.7\textwidth}{X} 所以，耶穌也在城門外受苦，為要用自己的血使百姓成聖。 \end{tabularx} \\ \\ \relax
13:13 & \begin{tabularx}{0.7\textwidth}{X} 這樣，我們也當走出營外，到他那裡去，忍受他所受的凌辱。 \end{tabularx} \\ \\
[1ex]
\hline
\hline
\end{longtable}
$^{1}$謝謝.
我們一起同心祈禱.
天父我們將今早我們能夠敬拜.
能夠聽到宣教士的召命.
他們那份內心的單純.
對福音能夠傳給未得之民.
尤其是在一些不同國家的難民.
我們憤外感動.
很盼望主在我們今早聚集的時候.
你的聖靈感動我們.
藉著這一段經文.
主耶穌怎樣來到這個世界.
為我們成就救恩.
成為我們人生生命的故事.
一路我們去編寫.
直至我們見主面前.
我們是一班走出營外的基督徒.
我們這樣禱告.
奉耶穌基督的名.
阿們.
再一次在母親節都祝願各位.
尤其是節期.
憤外想起不同的感人故事.
每逢母親節.
憤外這兩天大家留意那些新聞.
特別是報紙網上的新聞.
就很多感人故事陸續出台.
父親節就少一點.
對的.
父親節比較含蓄一點.
很想今天用故事.
來帶進我們今天.
怎樣來看我們的人生.
其實我們大家是在編寫.
將來的人看我們.
是怎樣的一個基督徒的故事.
當然現在.
我們在看前人.
他怎樣走他的人生路.
成為我們的激勵.

$^{41}$成為我們大家的支持.
今天其實很想就著.
阿時司的分享.
因為都看著他們的夢照.
一直Albert信主.
一直他們經歷到這個.
照明.
我算是跟他很緊密地.
看著他的發展.
由他來做神學生.
過來弘恩.
都是由我去編排等等.
看著他們的家庭.
很清楚神一直帶領.
很想就著這個故事.
去看看我們自己的故事.
跟大家選讀的經文.
就是希伯來書十三章.
剛才第七節已經很重要.
作者去到最後的時候.
祝福信徒.
要有這樣的眼光.
去編寫你的故事.
你需要有材料的.
人的故事就是最重要的材料.
他說從前引導你們的.
傳神的道給你們的人.
當然就是傳道者.
或者哥哥姐姐.
一些屬靈的領袖.
你要紀念他們.
效法他們的信心.
以及要回顧他們為人的結局.
那個為人.
就是他們的人生是怎麼走的呢?.
他們內心的價值觀又是怎樣的呢?.
這些是很重要的.
要紀念,要效法,要回顧.
就是說你的信仰原來是要自己參與下去的.
但是你不會自己閉門就有一個故事的.

$^{81}$一定是別人在你的生命裡面.
帶給你的衝擊.
帶給你的影響.
而且那個還要是好的影響.
當然了.
更多的可能是不好的影響.
我們都很難去分辨.
有時候一些不好的領袖.
會影響我們信仰的.
會帶給我們有時候會有挫敗感.
當然這個是正面的.
所以我們要學習.
如何去學西選.
大家有沒有去謝霆鋒的四場演唱會?.
五四中十八萬多人.
是如癡如醉的.
因為二十五年已經沒有開演唱會了.
那個公仔遮著頭.
就是一個當年.
謝霆鋒二十歲的時候.
開第一場的Viva演唱會.
然後隔了二十五年了.
剛剛開完的演唱會.
那麼這位女士今年四十歲了.
她去完兩場演唱會.
買最貴的票.
可想而知她受著謝霆鋒的影響有多深.
她一直說她十五歲那年.
周圍的人都不理她.
為什麼她喜歡這個年輕人.
這麼討人厭.
為什麼你會喜歡她.
她爸爸都不准她跟這個偶像.
她不理那麼多.
十五歲那年.
不斷去追星.
去到不同的地方.
這張照片.
就是終於她追星.
謝霆鋒跟她拍的了.

$^{121}$因為她在看完這場演唱會之後.
在網上講述她二十五年的人生.
由她十五歲的時候.
瘋狂的去追謝霆鋒.
在每一個場合.
每一次能夠聽到他在公眾地方來見面.
受著謝霆鋒的影響.
她就說.
不無悔她現在這二十五年.
所以她知道二十五年之後.
這個偶像.
唯一的偶像再開演唱會.
她說.
你再不開.
我就不可以再叫.
不可以再跳.
因為我已經四十歲了.
親愛的影子妹.
一個明星,歌星.
可以影響著一個人.
他的人生.
來到現在.
她感覺無悔了.
因為我能夠二十五年後.
再有機會在最貴的門票.
很近的在謝霆鋒面前.
又唱又叫.
好像已經完成了她人生的一個夢.
相信我們作為基督徒.
我們多不多回來聚會.
見到一些不同的領袖.
屬靈的長輩.
傳道給我們的人.
他們的為人是怎樣?.
他們的心是怎樣跟隨耶穌?.
有沒有嘗試去靠近.
聽聽他們的故事?.
有沒有嘗試去看看他們的日程表是怎樣做?.
怎樣去服侍他們?.
二十五年不是短.

$^{161}$二十五年可以帶來很多很多的.
生命的影響.
不用回答我.
有沒有人現在還在影響你的生命?.
很重要.
有沒有人在你的內心.
不知不覺跟你一起編寫著你自己在主面前的故事?.
如果沒有.
這是問題所在.
因為你對信仰根本就不是太熱心擺上.
不過仍然都是.
人始終是人.
有他的叛逆.
有他的不羈.
有他自己的生活各方面.
始終是人.
如果我們將我們的心靈.
放在一個人上.
始終有時會令我們很大的失望.
但是第八節很特別.
耶穌基督是永遠都不改變.
我們看到一個人.
越留意.
英雄見慣都變成常人.
但是如果一個人.
他能夠這樣來到我們走的時候.
今天母親節都要講一下女性.
在我的生命裡面.
這兩位女士.
是影響得我很深.
有一位大家可能不太認識.
就是剛剛退休的建築神學院神學系的副教授.
黃玉明博士.
上過他一年的課程.
他教教會歷史.
第二年班開始要上他的課程.
第二年雖然已經過了第一年班.
最困難的希臘文的三課.
全年三課.
不用測驗.

$^{201}$不用再碰那些外語.
但是開始要很深學習寫論文.
有一天.
二年班剛剛開學沒多久.
如常進去長洲.
來回要兩個小時.
一個單程要兩個小時.
當時住在元朗.
疊翠峰.
可想而知兩個小時才回到家.
四點三一落頭要趕下去.
衝下去坐船.
坐慢船.
回到家差不多到元朗.
拖著行李箱一路走回家的時候.
突然收到黃老師的電話.
很奇怪很特別.
普通一個同學.
全班有十八個人.
再加上其他班級一起上.
基本上不會認識我.
剛剛開課沒多久.
差不多六點的時候.
回到家差不多.
他打電話給我.
我聽黃老師.
睿貴.
我感覺到你很不開心.
很大壓力.
什麼事啊.
我簡直呆了.
他怎會知道我最近這麼不開心.
因為開始的功課壓力.
走來走去.
走讀生.
兩個小時坐船趕回家.
回到家六點已經很累.
家人又不明白你為什麼會這麼累.
讀書不是應該很開心嗎.
不是應該很空閒嗎.

$^{241}$很多壓力.
當時學校和同班同學.
大家來自不同的中派.
開始有不同的碰撞.
有時候在組會等等.
開班會的時候.
都不是這麼同心這麼開心.
他感受到.
他知道.
其實我比他大一點.
黃老師.
他很年輕.
他感受到我不開心.
感受到我有很大的壓力.
他說你有什麼要跟我分享.
你隨時來我的辦公室.
來找我.
整個人.
一個大男士.
被這位博士.
感化,感動.
到我一路畢業.
對這位老師.
越了解他之後.
原來他對學生的愛.
無微不至到一個地步.
他能夠陪學生突然患病.
急病.
或者有些會被蛇咬到.
他就要立刻搭飛機.
過去香港.
他要陪著學生.
直到家人來到.
他才會離開.
很多時候他的模樣.
比教書更加得到學生的愛戴.
認同.
這位是我所敬重的黃玉明老師.
學習這些人.
原來他們做神學.

$^{281}$是如此落地的.
原來關心那些同學.
是如此無微不至的.
還有當然就是中間張美眉牧師.
剛好也是建道神學院.
也是我前上師.
這位牧者是我真的非常敬重.
對他的為人.
因為我在2004年.
對不起.
我是2009年畢業.
應該是2010年的時間.
我畢業之後第二年.
我就進去錦繡堂.
去探望張牧師.
因為當年我進神學院的時候.
如果不是他為我填寫推薦信.
我是進不了神學院的.
他是相信我.
信任我.
認為我是有神的呼喚.
他願意個人推薦我進建道神學院.
所以我這麼多年.
過了畢業之後.
我很想去探望他.
很老土地將我四年的成績表.
和第三年,第四年拿到的獎狀.
告訴他.
很想說一聲多謝.
如果沒有你當年這麼決心.
願意幫我寫推薦信.
是不會有我今天這條事奉的路.
可以再走的.
我衷心地去多謝他.
說他對我的恩情.
跟著和他分享.
講我畢業之後.
在一間小堂會裡面的福事.
和他分享的時候.
我問他一些意見.

$^{321}$他第一句說話.
「瑞貴,不如來感受吧」.
於是在當年的一零年.
我就在錦繡堂開始那個福事.
很感恩.
在我未按目之前的五年.
由一零年到一五年的按目.
張牧師很經常.
應該叫做請我入房.
來講很多牧羊的生命.
和教我怎樣做一個傳道.
和他的為人.
相信在錦繡堂認識張牧師的人都會知道.
他的為人.
那種坦蕩,正直.
和沒有任何的第二層層.
或者隱藏層層是沒有的.
他講就是這樣.
不是有任何的東西.
我講了這樣.
其實不是這樣的.
不是的.
他整個人就是被人感覺到那種坦蕩.
他那份對別人的欣賞,信任.
是造就今天的我.
所以對牧者這一份的愛.
這一份提攜的恩典.
我是銘記在心靈裡面.
也成為我自己的重要的座右銘.
也看著他.
上個星期他剛好在錦繡堂講道.
明仔回來跟我說.
張牧師說他為什麼退休.
其實我知道的.
我和勵善牧師都知道.
因為他六十歲那年.
驗出輕微有一點點腦的瘀失.
很輕微而已.
不嚴重的.
但他已經知道.

$^{361}$他不應該再阻礙錦繡堂的發展.
他需要多一點的休息.
和要轉另一個跑道.
他很想在神學教育上來服侍.
就是這樣毅然放下.
牧堂很繁重的牧養工作.
這樣放下.
他有一個很值得我.
效法和敬重的.
他說我離開就是離開.
我不會干涉牧堂的發展.
他覺得不應該有尾巴.
他也很想和牧堂有好的關係繼續.
他真的講得出做得到.
他在新的工場.
在見道神學院裡面.
有一個教席來服侍.
他很開心能夠帶領神學生.
在在學期間學習.
這位是我很敬重的張牧師.
鼎姐妹.
有沒有人一步一步陪伴你.
編寫你自己的故事.
你事奉的故事.
有沒有一些你很敬重的人.
你衷心對他敬佩.
不是巴結不是奉承.
是敬重.
從他身上那種屬靈的氣質.
成為你的成長養分.
來編寫你自己的故事.
如果沒有的話.
有首詩歌.
這首是很出名的歌.
時光可變世界可變.
人情亦許多都變遷.
有共情不變.
那種爭找不到缺點.
你我再次相見.
其中一節很感動我自己.

$^{401}$真的.
人一定會有轉變.
但是耶穌基督.
這就是十三章第八節.
當我們去看這些帶領我們.
傳道給我們的人.
看得多.
我們要看回耶穌.
真正永不改變的是耶穌.
我們的信仰的堅固.
是耶穌基督.
所以耶穌基督昨日.
今日一直到永遠是一樣的.
人可以不一樣.
人的事奉的心可以不一樣.
隨著年日.
隨著環境.
可以有很多很多的變遷.
但是我們信仰的核心.
是在耶穌那裡.
所以耶穌才是永遠不改變.
我們要跟隨.
我們要效法.
所以他的故事.
就應該成為你和我.
編寫的人生事奉的故事.
我們繼續下去第九節.
突然間說.
有些似乎當時.
真的再次回到猶太教.
有些可能全假道的人.
你們不要被那些種種怪異的教訓.
勾引了他.
因為人的心.
是要靠恩典得堅固.
才是好的.
不是靠飲食.
所以飲食不是普通的飲食.
不是吃喝喝的飲食.
是說那些祭祀的禮儀.

$^{441}$宗教的條例.
可能當時有很多.
眾說紛紜的.
說要這樣的禮儀.
你才能夠蒙神的.
祝福恩典的.
作者告訴我們.
不是的.
一方面也告訴我們.
傳道者是要傳講.
是要將恩典林黎晉寺這樣講解.
叫人不是靠誰.
我們以前的舞堂.
我自己的成長教會.
是很傳統很飢餓的福音堂教會.
我們會以福音堂為榮.
如果學得不是最好的時候.
福音堂以外就好像沒有救恩一樣.
是很恐怖的.
不知不覺.
講得多.
我們要很純正.
我們要很堅持.
我們要赤裸裸跟隨神.
福音堂以外的人.
你不是這樣.
你不是這樣的時候.
你就是次等了.
原來追求得好.
認真.
原來是等同於看不起其他人.
所以絕對不出奇.
有人會傳這些似是而非的道.
所以作者說.
不要被這些種種怪異的教訓.
是教訓.
就是說傳給你.
你也不知道.
似是而非.
所以傳道者的其中一個重要的生命.

$^{481}$就是要講明耶穌的愛和救恩.
不是其他你自己的傳統.
你自己的執著.
你自己的看法.
是很重要講恩典.
是恩典將人拯救.
不是靠這些飲食.
從來不是這樣.
然後這裡告訴我們第十一節.
我們有一個祭壇.
上面的祭物.
是那些在會幕裡面公職的人.
是沒有權可以吃的.
就是說那是全年最重要的贖罪日.
七月初十的贖罪日.
那些血要到致勝所去彈.
那些肉就要燒了.
是不可以來祭司吃的.
因為生畜的血被大祭司帶入致勝所的時候.
作贖罪祭.
生畜的身體.
就是獻祭的身體.
就是那些祭生.
就要在營外來燒掉.
是不可以吃的.
大祭司只可以吃平安祭的祭肉.
突然之間講這個獻祭.
突然間講這個.
似乎就有一個相對.
這裡告訴我們.
所以耶穌也在城門外.
耶路撒冷.
各國他這個俘虜頭地.
這個山上面.
來受苦.
為要用自己的血.
使百姓能夠永遠的成聖.
就是說第十到十二節.
就明顯是一個對比.
就是說舊約的舊有的獻祭.

$^{521}$已經失效了.
已經沒有用了.
唯有耶穌才能夠帶給我們真正的成聖.
而且他還要在城門外來受苦.
各位.
這裡告訴我們.
外面的營.
十三節就是這一段聖經.
是七節到十三節的一個最重要的總結句.
也都是.
真的很想回應時詩.
他們走出營外.
離開香港.
去到宣教的工場.
當然絕對不是告訴我們.
我們每個人都要開始辦移民.
我們去宣教了.
不是這個意思.
我們的安書區.
我們自己以為可以來學校耶穌.
作者在這裡告訴我們.
一個事奉者.
是要這樣來成為別人的祝福.
這樣我們也當走出營外.
到他那裡去.
忍受他所受的凌辱.
我相信任何一個事奉者.
任何一個願意為主.
來擺上我們人生的故事.
如果沒有了那個凌辱.
沒有了那個出到營外.
是很難說得通.
你說你是有一心.
要為神擺上.
連我們的救主.
都是這樣走出營外.
成就這個永恆救恩.
最後想用一個例子.
去結束今天的講道.
大家有沒有看過麥迪文的電影.

$^{561}$我相當喜歡這個明星.
他有一次被邀請.
去到MIT麻省理工.
這麼高的學府.
來作一個畢業禮的致詞嘉賓.
他和這班大學生.
準備要出社會了.
這班高材生來勉勵.
他說其實他自己.
是在波士頓南部成長的.
他自己也曾經有機會.
在哈佛大學裡面.
不過讀不成.
因為很快.
他在在學期間.
就出去拍電影.
後來就不再讀書.
成為了一個舉世的明星.
他勉勵這些同學.
他說在MIT或者其他的大學.
特別是在牛津大學.
有一位教授.
曾經有一個這樣的推論.
這是一個AI.
很重要的一個電機教授.
他說有可能.
我們現在存在的星球.
即是這個地球.
是有外星高等的生物.
為我們設計的一個模型.
即是說我們的存在.
只不過是一個模擬機器.
是有外星的很高智慧的生物.
設計我們存在的.
我相信這個人.
這個教授一定是未信主.
他有這樣的想法.
麥迪文就用了這個例子.
來勉勵在座的大學生.
當然他否定了這件事.

$^{601}$他說如何能夠證明.
我們不是一個模擬機器.
被人操控.
隨時按個按鈕.
關掉了.
這個地球就玩完了.
就是這個了.
他說轉身面對你看見的問題.
走出去.
看到這個社會.
人類有什麼問題.
你不要視而不見.
轉身面對這個問題.
用你的辦法.
用你的人生的激情.
去改變那個環境.
麥迪文現在有很多的基金.
其中一個叫water.org.
這個水基金.
就是他創辦的.
為什麼呢.
因為他去過危地馬拉.
去過非洲那些落後的地方.
是完全沒有水.
基本的.
在一個物質富裕的西方.
歐美.
水根本不是一件什麼事.
但居然在這些地方.
在一個絕對貧窮的地方.
這樣的一個問題.
他覺得沒有辦法.
他不需要去批判.
而是他立即找到有心人.
創辦基金.
去改變那個社會的問題.
丁姐妹.
你的故事是怎樣呢.
你的人生.
人家將來怎樣看你呢.

$^{641}$我都期望將來的牧者.
後輩.
看我這個人.
曾牧師.
他已經盡他的能力.
走完他人生的路了.
你的路是怎樣走呢.
首先就是你必須要看到前人.
在你生命裡面的.
捉著你的手教你初信班.
慢慢看到你成長.
在教會事奉.
不斷讓你看到.
人生的不同的信心.
到你自己.
能夠獨當一面的來服侍.
耶穌基督在你的生命裡面.
這一份恩情.
這一份恩典.
是怎樣不斷令你能夠有力量.
這樣擺上呢.
轉身面對那個問題.
神的國度有很多的需要.
不要再浪費時間.
在一些無謂的事情.
集中心靈精力.
去做一些神喜悅我們做的事情.
我們一起同心抱告.
感恩.
我們能夠聽到孫教士.
他們知道困難.
但是他們願意面對.
在這些難民群體裡面的需要.
就是靈魂的得救的需要.
我們求主.
祝福穎兒姐妹一家人.
也都求主.
在母親節的今早.
來感動我們.
當我們編寫我們人生.

$^{681}$自己的故事的時候.
我們能夠讓後來者看到.
我們就是這樣.
走完我們人生的路.
由這一位永遠不改變.
昨日.
今日.
一直到永遠都不改變的主.
成為我們的目標.
成為我們追求.
事奉的力量.
我們這樣禱告.
奉耶穌基督的名.
阿們.
\newpage



\section{路加福音 7:36-50}
\label{sec:_TgXxd1W4v8}
\textbf{2025.05.18《看見尊貴》 路加福音7\_36-50 講員 : 黃國維牧師}
\newline
\newline
連結: \href{https://youtube.com/watch?v=-TgXxd1W4v8}{\texttt{https://youtube.com/watch?v=-TgXxd1W4v8}} ~~~~ 語音日期: 2025-05-20
\newline
\newline
\hyperref[sec:KCtIY5_ZgcE]{\small{< < < PREV SERMON < < <}}
~
\hyperref[sec:index_chronic]{\small{[返順時目]}}
~
\hyperref[sec:index_scriptual]{\small{[返順卷目]}}
~
\hyperref[sec:xtvbpi6Ud4M]{\small{> > > NEXT SERMON > > >}}
\newline
\newline
路加福音 7:36-50
\newline
\begin{longtable}{cl}
\hline
\hline
章節 & 經文 (和合本修訂版)\\
\hline
7:36 & \begin{tabularx}{0.7\textwidth}{X} 有一個法利賽人請耶穌和他吃飯，耶穌就到那法利賽人家裡去坐席。 \end{tabularx} \\ \\ \relax
7:37 & \begin{tabularx}{0.7\textwidth}{X} 那城裡有一個女人，是個罪人，知道耶穌在法利賽人家裡坐席，就拿著盛滿香膏的玉瓶， \end{tabularx} \\ \\ \relax
7:38 & \begin{tabularx}{0.7\textwidth}{X} 站在耶穌背後，挨著他的腳哭，眼淚滴濕了耶穌的腳，就用自己的頭髮擦乾，又用嘴連連親他的腳，把香膏抹上。 \end{tabularx} \\ \\ \relax
7:39 & \begin{tabularx}{0.7\textwidth}{X} 請耶穌的法利賽人看見這事，心裡說：「這人若是先知，一定知道摸他的是誰，是個怎樣的女人；她是個罪人哪！」 \end{tabularx} \\ \\ \relax
7:40 & \begin{tabularx}{0.7\textwidth}{X} 耶穌回應他說：「西門，我有話要對你說。」西門說：「老師，請說。」 \end{tabularx} \\ \\ \relax
7:41 & \begin{tabularx}{0.7\textwidth}{X} 耶穌說：「有兩個人欠了某一個債主的錢，一個欠五百個銀幣，一個欠五十個銀幣。 \end{tabularx} \\ \\ \relax
7:42 & \begin{tabularx}{0.7\textwidth}{X} 因為他們無力償還，債主就開恩赦免了他們兩個人的債。那麼，這兩個人哪一個更愛他呢？」 \end{tabularx} \\ \\ \relax
7:43 & \begin{tabularx}{0.7\textwidth}{X} 西門回答：「我想是那多得赦免的人。」耶穌對他說：「你的判斷不錯。」 \end{tabularx} \\ \\ \relax
7:44 & \begin{tabularx}{0.7\textwidth}{X} 於是他轉過來向著那女人，對西門說：「你看見這女人嗎？我進了你的家，你沒有給我水洗腳，但這女人用眼淚滴濕了我的腳，又用頭髮擦乾。 \end{tabularx} \\ \\ \relax
7:45 & \begin{tabularx}{0.7\textwidth}{X} 你沒有親我，但這女人從我進來就不住地親我的腳。 \end{tabularx} \\ \\ \relax
7:46 & \begin{tabularx}{0.7\textwidth}{X} 你沒有用油抹我的頭，但這女人用香膏抹我的腳。 \end{tabularx} \\ \\ \relax
7:47 & \begin{tabularx}{0.7\textwidth}{X} 所以我告訴你，她許多的罪都赦免了，因為她愛的多；而那少得赦免的，愛的就少。」 \end{tabularx} \\ \\ \relax
7:48 & \begin{tabularx}{0.7\textwidth}{X} 於是耶穌對那女人說：「你的罪都赦免了。」 \end{tabularx} \\ \\ \relax
7:49 & \begin{tabularx}{0.7\textwidth}{X} 同席的人心裡說：「這是甚麼人，竟赦免人的罪呢？」 \end{tabularx} \\ \\ \relax
7:50 & \begin{tabularx}{0.7\textwidth}{X} 耶穌對那女人說：「你的信救了你，平安地回去吧！」 \end{tabularx} \\ \\
[1ex]
\hline
\hline
\end{longtable}
$^{1}$第一天早安.
今天很感恩能夠來到大家一起去到崇拜.
我也是宣道會的.
我的母會是北宣.
我自己是住在九龍東那邊.
所以我平常是在遊堂宣道會裡面.
不過來到看到宣道會的標誌也很熟悉.
今天是遠了一點.
不過很感恩周牧師邀請.
也感受到.
大家很樂意去服侍.
臨死的地方.
見證上帝.
所以今天非常感恩能夠來到.
不如我們一起做一個祈禱.
因為你曾經渡盛肉身來到我們當中.
你去傳講天國的道理.
你在地上.
做了我們的工作.
因為我們的聖經再一次.
讓我們看到你怎樣去工作.
我們就求你今天臨在我們當中.
讓我們親眼可以見到你.
求你的說話.
更新我們的心靈.
讓我們更加做你.
合用的工人.
聽我們祈禱祈禱告奉耶穌的名.
阿們.
今天的經文在路加福音第七章.
裡面.
剛才主持幫我們讀.
我們發覺這段經文裡面.
有幾個人.
一開始的時候36節和37節.
就介紹了三個人物.
第36節就說有一個.
法利賽人.
法利賽人就請耶穌.
去吃飯.

$^{41}$然後耶穌就去了和法利賽人.
在家裡坐直吃飯.
然後就介紹.
城裡面有一個女人.
是一個罪人.
在這段經文裡面就有耶穌.
有一個女人和法利賽人.
問的問題就是誰是主角呢.
究竟這段經文想說誰呢.
聰明的弟兄姊妹就一定.
跟我說耶穌一定是主角.
就是永遠都是耶穌的主角.
不過究竟那兩個人.
女人和法利賽人.
究竟這段經文.
想說些什麼.
注意力在誰身上呢.
我們不如就.
逐個看一看.
我們首先看一下女人.
聖經第37節開始就已經介紹.
那個女人沒有名字.
不過說她是什麼呢.
是一個罪人.
整段的聖經沒有否定她是一個罪人.
甚至是眾所周知.
她是一個罪人.
這個女人做了什麼呢.
她知道耶穌.
在法利賽人家裡坐直.
她知道.
法利賽人請耶穌.
去她家裡吃飯.
她就聽到.
她就想去找耶穌.
就拿著聖滿香膏的玉瓶.
去到那時候就站在耶穌背後.
靠著他的腳哭.
眼淚滴濕了耶穌的腳.
就用自己的頭髮抹乾.

$^{81}$又用嘴連連親他的腳.
把香膏抹上.
這個罪人.
女人是沒有被邀請的.
很明顯她之前.
是遇見了耶穌.
聖經沒有說什麼.
不過肯定是遇見.
她就很想去找耶穌.
做剛才所說的那些事.
就.
有一點不請自來.
去到那裡.
大家就坐下來吃飯.
想像一下.
那個女人做的事都很尷尬.
很露骨的.
拿著聖滿香膏的玉瓶.
當然大家知道.
這個香膏是很昂貴的.
之後就拿著香膏去摳耶穌的腳.
玉瓶的香膏.
可以說是幾個月的工資.
是很貴的東西.
來到耶穌背後.
就靠著他的腳哭.
眼淚滴濕了耶穌的腳.
又用自己的頭髮抹.
還要用嘴連連親他的腳.
不要說以前的世界比較保守.
今天這樣.
男女授受不親.
其實都很露骨的.
所以我相信.
那時候的氣氛都相當尷尬.
大家都知道.
為什麼這個女人做這些事呢.
後面的經文其實有解釋.
我們跳一跳.
到後來耶穌說了一個比喻.

$^{121}$去解釋這個女人.
為什麼這樣做.
到第四一節耶穌就說了.
說了一個比喻.
兩個人都欠了某個債主錢.
一個欠了五百個銀幣.
一個欠了五十個銀幣.
一個多一個少.
一個少一點.
兩個人因為他們無力償還.
債主就開恩赦免了他們兩個人的債.
兩個人都已經還不了錢.
債主就兩個都免了.
問題是.
那麼這兩個人哪一個更愛他呢.
原來那時候.
耶穌和一個法利賽人.
那個法利賽人叫西門.
他們有一個對答.
那麼耶穌就透過這個比喻.
去解釋那個女人為什麼這樣做.
西門回答.
我想是那個多德赦免的人.
耶穌對他說.
你的判斷不錯.
那麼在耶穌這個比喻裡面.
他其實是在說.
那個女人那個罪人.
是欠了債主.
欠了債主.
欠了耶穌 欠了上帝的錢.
那麼那個女人是誰呢.
欠了五百個銀幣.
還是欠了五十個呢.
肯定是五百個 因為每個人都知道她是罪人.
所以這個.
也說到原來之前發生過什麼事呢.
一定是耶穌.
遇見了這個罪人女人.
而赦免了那個女人.

$^{161}$所以那個女人對耶穌.
有一個很大的感激之情.
那麼就透過剛才那些.
露骨在人面前.
很尷尬的行動.
去到回應.
或者感激.
耶穌 特別的地方.
是什麼呢 在比喻裡面.
耶穌就問.
法利賽人.
究竟那兩個人 誰更加愛他呢.
西門.
法利賽人回答得對.
我想是那個多德赦免的人.
耶穌對他說.
你的判斷不錯.
這個意思是什麼呢.
這個意思是那個女人即使做的事很尷尬.
即使很不舒服.
其實她做的事很合理.
連西門也說.
這個女人其實她因為被赦免得多.
其實她所做的事.
是相當之合理的.
即使.
是很尷尬的.
所以.
耶穌就.
接著是肯定了.
這件事.
然後他就說 轉過來向著那女人.
對西門說 你看見這女人嗎.
即是西門.
不明白那個欠錢的女人.
現在耶穌.
劃出場跟你說.
為什麼那個女人做事那麼尷尬.
在你家裡做那麼奇怪的事.
因為.

$^{201}$你看看那個女人.
她就是表達.
對我赦罪的那種.
感激 所以她就這樣做.
接著他就說 你看見這女人嗎.
我進你家裡.
你都沒有洗腳 這個女人就用眼淚.
來幫我洗腳.
又用頭髮抹乾.
你沒有親我 這個女人.
就不住親我的腳.
你沒有用油膏抹我的頭.
這個女人用香膏.
抹我的腳.
解釋那個女人 其實很合理的.
她做的事.
最後一句 所以我告訴你.
她的罪都赦免了.
因為她愛的多.
少得赦免的 愛的就少.
第47節解釋.
這個女人.
罪赦免 因為她愛得多.
不過我看那個翻譯.
我比較喜歡.
新漢語譯本 比較清楚.
新漢語譯本是怎樣說的.
因此我告訴你.
她愛得多.
是證明了.
她很多罪都赦免.
她愛得少.
就證明她赦免的少.
這個意思是甚麼.
不是說我們先愛耶穌.
所以他才赦免我們.
我們知道不是這樣的.
我們知道是甚麼.
是我們犯了很多罪.
是耶穌先赦免我們.

$^{241}$無論多少多少.
五百個銀子 五十個.
都是被耶穌赦免.
既然赦免得多.
那個人的愛.
愛得多證明了.
她很多罪都赦免了.
所以我覺得很特別的地方.
原來我們.
對耶穌那種愛的表達.
好像那個女人.
有甚麼.
香膏 眼淚.
親耶穌的腳.
這些表達是證明了.
我們原來被赦免.
我想這對我們有提醒.
做基督徒.
如果我問大家 你愛不愛主.
你一定會說愛.
為甚麼 因為我們說.
神愛我們 我們願意愛上帝.
甚至耶穌吩咐我們.
要愛神愛人.
我們每個基督徒都會說.
我們愛 我們願意學習.
愛上帝 愛耶穌.
這句說話就問.
那你怎樣去證明你的愛.
那個女人用甚麼證明.
就用她的香膏.
用她的眼淚.
用她的眼淚 愛耶穌的腳.
用她在眾人面前.
做一些很尷尬的事.
去證明.
這個女人.
都給了我們一些提醒.
我們知道我們每個人.
多多少少.

$^{281}$最多是給上帝赦免.
我們都說我們愛耶穌.
不過我們用甚麼去證明.
我們的愛呢.
聖經說我們的愛不是只用口說的.
我們的愛帶來行動.
我們怎樣證明.
我們的愛主.
這個女人告訴我們她的愛主是獻上香膏.
香膏可以是.
很昂貴的東西.
幾個月的工資.
我願不願意透過.
我付出的金錢.
我付出的資源.
去愛耶穌.
又或者.
這個女人用她的情感.
很感激耶穌.
我們是不是.
用我們的情感.
用我們的眼淚.
去愛耶穌呢.
又或者這個女人是很願意.
在眾人面前.
做一些很尷尬.
令人覺得這樣的事情.
去表達對於耶穌的愛.
我們又樂不樂意呢.
有時候我們覺得.
今天要傳福音.
有時候都很尷尬.
今日的社會會覺得.
這些信仰的東西很私人.
公司裡面最好不要提.
你在別人面前見證耶穌.
你好像很愁.
今日的社會不是很接受.
我們願不願意.
因為愛耶穌的緣故.

$^{321}$放下自己.
願意在別人面前.
做一些別人都覺得尷尬的事情.
去表達對耶穌的愛呢.
他說.
愛原來是透過.
付出很多的錢.
做一些尷尬的事情.
其實今天也有.
不過很多時候未必是基督徒.
是什麼呢.
我常常覺得.
有些人想求婚的時候.
都會做這些事情.
即是.
男生想結婚.
向女生求婚.
當眾跪下.
送戒指.
幾個月工資.
用這些去證明他的愛.
大家都知道.
這樣也可以.
但這個女人.
就是用這些去證明.
對耶穌的愛.
提醒我們.
我們是否也願意.
去回應耶穌的愛呢.
這是女人.
給我們的提醒.
接著我們看看.
法利賽人.
這裡就是.
一開始.
這段經文就介紹.
法利賽人請耶穌和他吃飯.
耶穌就應約.
去到法利賽人的家.
去坐席.

$^{361}$猶太人和中國人.
都有一個相近的地方.
請人吃飯.
其實是一種社交的.
大家都知道.
中國人吃飯.
不只是吃飽飯.
很多意思.
我們會知道.
尤其是法利賽人.
法利賽人是猶太人.
那時候的宗教領袖.
特別是會教人.
律法的.
教聖經.
像上學院的院長.
今天會教人.
大家都很尊敬的人.
所以.
這個看來.
是一個官方的.
隱喻.
為什麼法利賽人會這樣呢.
後面有些線索.
當那個女人.
做了剛才露骨的事.
之後.
聖經記載39節.
請耶穌的法利賽人.
看見這事.
這件事是什麼.
就是那個女人做露骨的事.
心裡面就說.
這人 耶穌.
若是先知.
一定知道我他的是誰.
是個怎樣的女人 是罪人.
這裡.
法利賽人認為.
耶穌是誰 是個先知.

$^{401}$所以.
為什麼法利賽人想請耶穌吃飯.
耶穌那時候.
已經出來傳道.
做上帝工作 行了很多神蹟.
出名的 所以很多人就說.
耶穌是從神而來的.
身為宗教領袖.
當然要看看他怎樣.
所以很合理的.
請你吃飯.
這是其中一個原因.
又或者 哇 這麼出名的人.
我有面子跟他吃飯.
跟他打卡.
其實都挺提升身價的.
可能是這樣的.
所以就請耶穌吃飯.
不過.
吃飯.
就有個罪人衝進來.
耶穌竟然.
讓他做這樣的事.
所以他的心裡面就說.
如果真的是先知.
一定知道那個女人是罪人.
為什麼會讓他這樣做.
所以我們會看到這個法利賽人.
其實是因為.
知道或者認為耶穌是從神而來的先知.
誰知道.
在他家裡面.
坐直發生的事.
叫到這個法利賽人西門.
失望 怎麼會這樣.
所以我們看到.
那個法利賽人其實為什麼請耶穌.
可能.
都想認識一下耶穌.
打卡 搞一下關係.

$^{441}$之類的.
接著到了第40節.
耶穌就回應.
這個西門.
其實他心裡面說話.
不過耶穌是看得出來的.
大家看得出來的.
所以耶穌就回應他.
就說 西門.
我有話要對你說.
西門老師請說 其實很特別的.
第40節.
可以不出現.
你在家裡請人吃飯.
你用不用跟人說.
我有話要對你說.
人家說 老師請說.
其實不用的.
為什麼出現這件事呢.
我們可以相信.
在那時候.
耶穌在那個坐直裡面.
都沒有說過話.
或者沒有機會說話.
如果不是 為什麼.
我想說一句話要跟人說.
所以我們從這裡都可以看到.
其實坐直裡面.
可能發生的事.
就是西門跟他的朋友.
因為正經有記載 整個圍人.
跟他的朋友.
一句兩句.
耶穌都沒有怎麼說話.
或者他們都不是請耶穌來.
虛心去聽耶穌教導.
耶穌要插嘴.
所以我們看到.
這個西門 原來.
他這樣的動機.

$^{481}$他看到有個從神而來.
先知的人不如請你吃飯.
抬高自己的身價.
做一些交際的活動.
其實這個西門.
在這裡都對我們有些提醒.
我們都可能信主.
年日都不短.
有時候習慣了教會.
有時候我們自己都要問一下自己.
究竟信仰對我來說.
是什麼呢.
是不是都 你明不明白.
就是有些都.
給我一些好處.
我都被人在人面前覺得你是基督徒.
都幾好啊.
有時候回教會 為什麼回教會 其實我不想回教會.
不過今天不回 被牧師見到.
問我為什麼不回教會.
我就不開心 你明不明白.
我們可能很多時候 信仰對我們來說.
都慢慢慢慢變得.
我考慮一下去哪些.
做些什麼 提高自己的身價.
我覺得這可能是很自然的反應.
不過.
我想今天這兩個人.
給我們一些提醒.
我們可以.
信很久其實沒問題.
法律出人可以都沒問題.
不過問題就在於我們慢慢慢慢.
信得久 會不會信仰都僵化了呢.
我問一下自己.
今天我回了很多年教會.
我做了很多年基督徒.
我的生命是不是仍然.
渴望聽到.
耶穌講話.

$^{521}$而不是其實耶穌坐在桌子上.
我都不願意聽他講話.
我想這都成為.
我們一些的提醒.
接著.
耶穌因為西門的反應.
哇.
這個人是誰.
才會這樣.
耶穌就說不如我有話跟你說.
西門就說請說.
耶穌就講了那個比喻.
剛才已經講過那個比喻.
那個比喻是講.
有兩個人欠了錢.
不過耶穌.
真的很厲害.
他那個比喻裡面.
就是講兩個人.
所以他一個比喻.
是講了或是回應了那兩個人.
一個是那個罪人女人.
是不是.
回到那個比喻裡面.
兩個都欠錢.
除了那五百個銀幣.
那個罪人女人之外.
耶穌說.
還有另一個人.
是欠五十個銀幣.
兩個人都在.
接著我們看到耶穌去解釋那個比喻.
44節.
耶穌轉過來對著那個女人.
就跟西門快利卓人說.
你見不見到那個女人.
接著他就說.
我進你家.
你沒有洗腳.
這個女人用眼淚.

$^{561}$洗腳.
你沒有做這些.
這個比喻裡面.
講了兩個人.
一個欠債很多.
一個欠少很多.
一個是女人.
另一個是誰.
另一個是西門.
西門忘記了.
原來他也是欠債的人.
他就是.
很覺得自己沒有欠債.
沒有欠債.
我生命很圓滿.
我沒有問題.
我做到一個這麼高位.
我教人律法.
宗教的領袖.
每個人都尊敬我.
我不需要.
但耶穌提醒他.
西門其實你也欠債.
正正是因為西門.
不覺得自己欠債.
所以他沒有愛耶穌.
沒有愛上帝.
那種表達.
我想這個也給我們提醒.
其實是甚麼呢.
到底是甚麼呢.
我們基督徒的生命.
上帝提醒我們.
我們仍然.
是欠上帝的債.
我們的生命仍然.
是需要成長.
即使那時候.
的法利塞人.
其實.

$^{601}$每個人都是欠上帝的債.
生命裡面.
仍然需要成長.
我在神樂院教書.
做院長.
可能對很多弟兄姊妹來說.
神樂院裡面.
很多人都很認真信仰.
上帝呼召願意.
放下高薪厚職.
裝備自己做傳道人.
很尊敬的人.
其實是的.
其實我裡面很多同學.
其實都很感動.
他們的故事.
不過你知道.
我們發覺我們怎樣栽培.
一個上帝的僕人.
就是在說.
原來我們都承認.
我們的生命很多軟弱.
我們的生命可以有很多的.
眼淚很多的創傷.
我們要不斷地去.
更新自己的生命.
即是在神樂院裡面.
做些甚麼的.
可能沒有讀過.
可能有很多神秘.
但其實我們每個學期一開始的時候.
我們有一個小組.
是做甚麼呢.
是跟同學一起去做一件.
叫做生命線的事.
其實是一個信任的環境.
大家回望.
上帝在我們生命裡面的工作.
往往不是很光彩的.
往往都是我發覺原來.

$^{641}$自己很多的軟弱.
很多的眼淚.
很多的創傷.
我們在恩主的面前.
去看清楚自己的生命.
就讓到.
讓到聖靈在我們的心裡面.
去到工作.
一路一路去到醫治.
我想這些不只是神樂生.
因為我在神樂院裡面.
有些真是校友.
出來教會去服侍.
有時候教會服侍都經歷很多.
不開心.
有時候都回來跟我們聊天.
生命不斷.
要去到成長.
西門就是這樣去被.
耶穌去到提醒.
任何人.
都軟弱.
包括我想我們.
每一個在內.
我們都渴望自己去到被.
這個耶穌去到提醒.
問題就是.
這個西門.
不是很願意.
這樣開放自己.
在聖經裡面很特別.
他的記載裡面.
是在說西門.
西門和他同籍的法利賽人.
無論遇到不同的事情.
都很喜歡心裡說.
很有趣的.
他們在桌上可能很多東西說.
不過真那句在哪裡說呢.
在心裡面說.

$^{681}$兩次.
第一原來上面那裡.
其實是最後的.
耶穌對著那個女人最後說.
你的罪赦免了.
法利賽人在第49節的回應.
又不說出聲.
同籍的人.
那班法利賽人心裡說.
這個耶穌就像什麼人.
竟然赦免人的罪.
你知道對他們來說.
對猶太人來說.
神才是赦免人的.
這個耶穌為什麼可以赦免人的罪.
因為他們不承認耶穌是神.
心裡說.
但早一點點也是.
39節是說那個女人.
做了露骨的事情之後.
請耶穌的法利賽人.
看到那個女人做了這些事情.
耶穌也不阻止.
心裡又說多一次.
這個耶穌如果是孫子.
就一定知道這個女人是誰.
所以這些法利賽人和西門.
很喜歡心裡說.
為什麼心裡說我們很明白.
也不錯啊.
給你面子.
人家覺得你做了一些不合眼的事情.
你就會覺得很有禮貌.
心裡說.
不過耶穌就不容許.
耶穌就說.
你心裡有什麼.
其實你要表達.
不過心裡說.
我覺得是發生什麼事.

$^{721}$其實是他心裡用了一把尺.
他用了一把尺去量度別人.
他兩次都量度耶穌.
第一次是量度耶穌.
為什麼讓罪人女人做這些事情.
第二次又用心裡的一尺去量度耶穌.
沒有人能夠赦免其他人.
為什麼一個人會用自己心裡的一尺量度別人.
因為覺得自己是最對的.
或者其實某程度上.
你會發現.
他們那一把尺不是完全沒有道理的.
你想想.
聖經說我們要聖潔.
不過問題是.
他把聖潔對自己的要求放在別人身上.
不是完全沒有道理的.
在教導的過程中.
他也會覺得.
其實是神才赦免人.
但問題是.
他們到了某個位置.
覺得因為他們懂得真理.
因為他們從小到大讀聖經.
他們一定是最對的.
所以每次見到任何人.
他們都用心裡的一尺量度別人.
或者另一句說話是.
見到任何人.
他們都有一個標準去看待別人.
用自己的一尺量度別人.
意思是任何人在他面前出現.
其實他看不到別人.
他只有一把尺.
見到任何人.
他都有一個真理在他身上.
他有一個準則.
他就用這個準則去量度別人.
這個又不好,那個又不行.
其實我們很明白.

$^{761}$其實這件事對基督徒來說.
真的有時候會有些引誘.
為什麼呢?.
因為我們基督徒.
說我們從小到大學聖經的真理.
第一,是好的.
不過.
不理睬人提醒我們要怎樣.
我們可能到了某個地步.
覺得我們真的.
好像上帝那麼厲害.
就只是用一把尺.
去量度身邊很多很多的人.
第一,我現在不是說我們不需要有真理.
聖經在這地方.
沒有說那個女人不是罪人.
不過問題是.
我們怎樣去抬高自己.
認為.
我就是真理.
所以不斷地.
用心裡的一把尺.
去量度別人.
看不到別人.
只看到一把尺.
你明白嗎?.
是這樣的.
我覺得整段經文.
最感動我的是第44節.
44節.
耶穌和西門.
去到聊天的時候.
講完比喻.
耶穌做了一個動作,44節.
於是耶穌.
轉過來.
向著那女人.
對西門說,你看見.
這女人嗎?.
為什麼要轉過來呢?.

$^{801}$因為原來早一點.
那女人衝進屋裡的時候.
在吃飯.
大家圍著圈.
那女人當然不能坐下.
那時候猶太人都是.
靠著地面坐下.
所以唯有那女人.
站在耶穌的背後.
靠著他的腳.
哭著,親他的腳.
做那些事情.
所以整個環境就是.
西門和耶穌坐在圈裡.
大家相對.
那女人在耶穌後面做的事情.
一直去討論.
直到44節,耶穌怎樣呢?.
轉身過來.
對著那女人.
不過就跟西門說話.
問她.
你見不見到她?.
其實在.
耶穌轉身過來之前.
他們是在討論.
那個女人.
英文就叫做Talk about her.
我是在討論她.
她是罪人.
為什麼要做這些事情.
但耶穌說完比喻之後.
他轉過來.
看著那個女人.
問她,你真的.
見不見到她?.
是見到她.
還是你見到你心裡的那把尺子呢?.
我們問.
見到什麼呢?.

$^{841}$我們看回.
那個女人被耶穌看見.
耶穌看見她.
耶穌看到她.
整個的生命.
耶穌明白她.
理解她.
和饒恕她.
接納她.
所以因為耶穌的接納和看見.
這個女人.
很感覺到.
被耶穌的接納和.
饒恕.
跟著的回應就是.
在她面前哭.
用眼淚滴濕她的腳.
頭髮抹乾,連連親她的腳.
用香膏.
抹上.
為什麼這個女人願意將自己最寶貴的東西.
獻給耶穌呢?.
因為耶穌.
真真正正.
看見她.
耶穌看到她.
是的,她犯了罪.
但其實她是一個.
很寶貴,很尊貴.
上帝所創造.
有神形象的人.
耶穌看到她.
一生她是犯罪.
但她是被罪纏繞.
耶穌看穿.
她的罪的表面.
看到她生命.
的需要.
在聖經裡面.
這段經文沒有清楚講.

$^{881}$不過新約聖經講到.
罪人還是女人.
通常可能是犯一些.
姦淫之類.
或者男女關係.
大家知道其實在那個處境裡.
可能為什麼有些人要做妓女.
未必是因為他自己想.
做這行的.
可能他被丈夫拋棄.
生活不到.
唯有是這樣.
她是犯了罪.
但耶穌.
看穿她.
犯罪的.
表面的行動.
而看得到.
她的尊貴.
她願意去接納.
但我想.
對於法理錯人來講.
剛才有講.
看見人都是用自己的尺.
看到任何人都用.
你和我不行.
我和你畫一條界.
你因為不聖潔.
我聖潔.
你不懂這些.
你沒有真理,我有真理.
我只是畫一條線在大家當中.
但耶穌就突破.
這個界線.
他看到真是尊貴的人.
你知道嗎?.
我們生命都要成長.
今天就經文.
女人和法理錯人.
他們哪一個人.

$^{921}$生命裡面有成長.
很明顯是一個女人.
那個女人.
其實是一個.
模範的門徒.
為什麼這樣說.
願意悔罪.
願意在人面前見證耶穌.
願意將最寶貴的.
金錢奉獻.
你說是不是一個最模範的門徒.
為什麼.
就是因為.
她經歷了.
主耶穌給她的接納.
我們很多時候都想.
讓人生命成長.
想讓自己生命成長.
或者我們很多時候會用.
法理錯人的方法.
你做得不好.
你犯罪.
做好一點.
批判你.
我們認為批評得多.
人們就會進步.
但上帝的工作原來不是這樣.
耶穌的工作不是這樣.
當他真正見到人的尊貴.
完全明白一個人的過去.
接納一個人的軟弱時.
人就給予一份恩典.
去改變你.
反而將自己最寶貴的.
金錢奉獻.
去改變你.
反而將自己最寶貴的東西.
去獻給上帝.
耶穌就是一位這樣的主.
所以.

$^{961}$女人被感動.
願意將自己的生命.
將自己最寶貴的東西.
去獻上.
所以聽到這段聖經.
沒有跟我們說.
我們不需要處理罪.
沒有這樣說.
聖經也沒有說.
罪人沒有犯罪.
不是這樣的.
上帝很渴望我們脫離罪.
很渴望我們身邊的人.
都脫離罪.
不過問題是.
耶穌來到.
他說.
上帝是施恩的上帝.
恩典的上帝.
他不選擇做法理裁人.
指責,抗拒.
而是選擇去接納.
回到我最初的問題.
究竟這段經文.
誰是主角呢?.
是女人還是法理裁人?.
我想三個都是主角.
其實我們可以想想.
我們可能是.
女人需要去.
透過我們的行動.
去表達對上帝的愛.
但聖經裡面.
出現法理裁人.
我經常覺得不是被我們罵.
是被我們提醒自己.
其實我們都很知道.
聖經裡面真的很有.
法理裁人的傾向.
我們都很喜歡.

$^{1001}$我對你錯,我有真理,你沒有.
不是要放下真理.
不過就是.
當我們面對身邊的人.
我們遇見人的時候.
我們可不可以.
好像下法主耶穌那樣.
真正看到別人.
看見別人.
進入別人的生命裡面.
而將從耶穌而來的恩典.
和別人分享.
我想這就是耶穌基督.
給我們那種無比的福音.
由我們經歷這份恩典.
生命被改變.
亦都將到同樣的恩典.
可以分享給身邊的人.
由他們給自己.
我們一起去祈禱.
因為你是那位恩典的上帝.
你來到工作.
你發覺.
你真的進入.
在當代裡面.
最有需要的人.
那些被社會.
去到拒絕的人.
你看得到.
他每一個的生命.
求你幫助我們.
其實你都首先看見.
我們每一個人.
其實我們每一個.
在認識你悔改之前.
我們都心硬.
我們都犯罪.
我們領受你.
無比的恩典.
即使我們爭取了你.

$^{1041}$多少都好.
你都赦免我們.
我們很感激因為你將我們的生命還回.
同時你.
呼籲我們.
你想我們效法你一樣.
對其他人都一樣.
能夠將到你以來的恩典.
和人分享.
所以主耶穌求你去幫助我們.
我們效法你.
去拆毀我們和人之間的界線.
更新我們.
將我們的恣意.
都拿開.
讓我們效法你.
讓我們帶著上帝的恩典.
亦都將到這份恩典.
和身邊的人去分享.
求你這樣幫助我們.
聽我們祈禱禱告耶穌基督的名.
阿們.
\newpage



\section{撒母耳記上 1:1-28}
\label{sec:xtvbpi6Ud4M}
\textbf{2025.05.25《借錢梗要還》 經文:撒母耳記上1\_1-28 講員 : 張雲開老師}
\newline
\newline
連結: \href{https://youtube.com/watch?v=xtvbpi6Ud4M}{\texttt{https://youtube.com/watch?v=xtvbpi6Ud4M}} ~~~~ 語音日期: 2025-05-27
\newline
\newline
\hyperref[sec:_TgXxd1W4v8]{\small{< < < PREV SERMON < < <}}
~
\hyperref[sec:index_chronic]{\small{[返順時目]}}
~
\hyperref[sec:index_scriptual]{\small{[返順卷目]}}
~
\hyperref[sec:HMy_7Yhf_bE]{\small{> > > NEXT SERMON > > >}}
\newline
\newline
撒母耳記上 1:1-28
\newline
\begin{longtable}{cl}
\hline
\hline
章節 & 經文 (和合本修訂版)\\
\hline
1:1 & \begin{tabularx}{0.7\textwidth}{X} 以法蓮山區有一個拉瑪的瑣非人，名叫以利加拿，他是蘇弗的玄孫，託戶的曾孫，以利戶的孫子，耶羅罕的兒子，是以法蓮人。 \end{tabularx} \\ \\ \relax
1:2 & \begin{tabularx}{0.7\textwidth}{X} 他有兩個妻子：一個名叫哈拿，另一個名叫毗尼拿。毗尼拿有孩子，哈拿卻沒有孩子。 \end{tabularx} \\ \\ \relax
1:3 & \begin{tabularx}{0.7\textwidth}{X} 這人每年從本城上到示羅，敬拜萬軍之耶和華，向他獻祭。在那裡有以利的兩個兒子何弗尼和非尼哈當耶和華的祭司。 \end{tabularx} \\ \\ \relax
1:4 & \begin{tabularx}{0.7\textwidth}{X} 每逢獻祭的日子，以利加拿把祭肉分給他的妻子毗尼拿和毗尼拿所生的兒女。 \end{tabularx} \\ \\ \relax
1:5 & \begin{tabularx}{0.7\textwidth}{X} 他給哈拿的卻是雙分，因為他愛哈拿。耶和華卻不使哈拿生育。 \end{tabularx} \\ \\ \relax
1:6 & \begin{tabularx}{0.7\textwidth}{X} 她的對頭毗尼拿因耶和華不使哈拿生育，就常常惹她發怒，要使她生氣。 \end{tabularx} \\ \\ \relax
1:7 & \begin{tabularx}{0.7\textwidth}{X} 年年都是如此。每當她上到耶和華殿的時候，毗尼拿就這樣惹她發怒，以致她哭泣不吃飯。 \end{tabularx} \\ \\ \relax
1:8 & \begin{tabularx}{0.7\textwidth}{X} 她丈夫以利加拿對她說：「哈拿，你為何哭泣？為何不吃飯？為何傷心難過呢？有我不比有十個兒子更好嗎？」 \end{tabularx} \\ \\ \relax
1:9 & \begin{tabularx}{0.7\textwidth}{X} 他們在示羅吃喝完了，哈拿就站起來。祭司以利坐在耶和華殿門框旁邊的位子上。 \end{tabularx} \\ \\ \relax
1:10 & \begin{tabularx}{0.7\textwidth}{X} 哈拿心裡愁苦，就痛痛哭泣，向耶和華祈禱。 \end{tabularx} \\ \\ \relax
1:11 & \begin{tabularx}{0.7\textwidth}{X} 她許願說：「萬軍之耶和華啊，你若垂顧你使女的苦情，眷念不忘你的使女，賜你的使女一個子嗣，我必使他終生歸給耶和華，不用剃刀剃他的頭。」 \end{tabularx} \\ \\ \relax
1:12 & \begin{tabularx}{0.7\textwidth}{X} 哈拿在耶和華面前不住地祈禱，以利注意她的嘴。 \end{tabularx} \\ \\ \relax
1:13 & \begin{tabularx}{0.7\textwidth}{X} 哈拿心中默禱，只動嘴唇，聽不到她的聲音，因此以利以為她喝醉了。 \end{tabularx} \\ \\ \relax
1:14 & \begin{tabularx}{0.7\textwidth}{X} 以利對她說：「你要醉到幾時呢？不要再喝酒了！」 \end{tabularx} \\ \\ \relax
1:15 & \begin{tabularx}{0.7\textwidth}{X} 哈拿回答說：「我主啊，不是這樣。我是心裡愁苦的婦人，清酒烈酒都沒有喝，只在耶和華面前傾心吐意。 \end{tabularx} \\ \\ \relax
1:16 & \begin{tabularx}{0.7\textwidth}{X} 不要將你的使女看作不正經的女子。我因極其難過和生氣，所以一直禱告到如今。」 \end{tabularx} \\ \\ \relax
1:17 & \begin{tabularx}{0.7\textwidth}{X} 以利回答說：「平平安安地回去吧。願以色列的神允准你向他所求的！」 \end{tabularx} \\ \\ \relax
1:18 & \begin{tabularx}{0.7\textwidth}{X} 哈拿說：「願你的婢女在你眼前蒙恩。」於是婦人上路，去吃飯，臉上不再帶愁容了。 \end{tabularx} \\ \\ \relax
1:19 & \begin{tabularx}{0.7\textwidth}{X} 他們清早起來，在耶和華面前敬拜，就回去，往拉瑪自己的家裡。以利加拿和妻子哈拿同房，耶和華顧念哈拿。 \end{tabularx} \\ \\ \relax
1:20 & \begin{tabularx}{0.7\textwidth}{X} 時候到了，哈拿懷孕生了一個兒子，給他起名叫撒母耳，說：「這是我從耶和華那裡求來的。」 \end{tabularx} \\ \\ \relax
1:21 & \begin{tabularx}{0.7\textwidth}{X} 以利加拿和他全家都上去，要向耶和華獻年祭和還願祭。 \end{tabularx} \\ \\ \relax
1:22 & \begin{tabularx}{0.7\textwidth}{X} 哈拿卻沒有上去，因為她對丈夫說：「等孩子斷了奶，我就帶他上去朝見耶和華，讓他永遠住在那裡。」 \end{tabularx} \\ \\ \relax
1:23 & \begin{tabularx}{0.7\textwidth}{X} 她丈夫以利加拿對她說：「就照你看為好的去做吧！可以留到兒子斷了奶，願耶和華應驗他的話。」於是婦人留在家裡乳養兒子，直到他斷了奶。 \end{tabularx} \\ \\ \relax
1:24 & \begin{tabularx}{0.7\textwidth}{X} 斷奶之後，她就帶著孩子，連同一頭三歲的公牛，一伊法細麵，一皮袋酒，上示羅耶和華的殿去。那時，孩子還小。 \end{tabularx} \\ \\ \relax
1:25 & \begin{tabularx}{0.7\textwidth}{X} 他們宰了公牛，就領孩子到以利面前。 \end{tabularx} \\ \\ \relax
1:26 & \begin{tabularx}{0.7\textwidth}{X} 婦人說：「我主啊，請容許我說，我向你，我的主起誓，從前在你這裡站著祈求耶和華的那婦人就是我。 \end{tabularx} \\ \\ \relax
1:27 & \begin{tabularx}{0.7\textwidth}{X} 我祈求為要得這孩子，耶和華已將我向他所求的賜給我了。 \end{tabularx} \\ \\ \relax
1:28 & \begin{tabularx}{0.7\textwidth}{X} 所以，我將這孩子獻給耶和華，使他終生歸給耶和華。」他就在那裡敬拜耶和華。 \end{tabularx} \\ \\
[1ex]
\hline
\hline
\end{longtable}
$^{1}$Angie 早安.
現在港督比以前複雜很多.
以前拿著聖經,最多一張稿就可以了.
我不知道平時港員過來領會或主日崇拜港督會不會用powerpoint.
我其實比較少用powerpoint,大部份時間都沒有用powerpoint.
不過今天我覺得有個powerpoint好一點.
結果做遲了,才給後面的同工.
但用powerpoint就要控制住.
這個,未開.
開了.
有了.
要用按鍵按,要拿著聖經.
手機要用來做什麼呢?.
萬一你知道所有東西都可能出事.
要有後備,我都開了自己的powerpoint在手機.
看多兩張,知道後面說什麼,一會兒忘記了.
就比較頭痕.
所以科技帶來,人們說科技可以空出更多有限的時間.
但其實我們看過往的日子,科技是令我們節省了某些東西的時間.
但人又會將時間填滿,結果有科技做多了東西.
但更忙.
現在有三隻手可能比較好,兩隻手可能不夠.
很高興,這是我第二次過來.
第一次還需要勞煩曾牧師開車接我.
這次我問了曾牧師,查一下地圖.
曾牧師都告訴我,這條水路轉輕鐵,很方便.
我問他走到門口要走多久,因為我上次崇拜回家的時候.
都是坐輕鐵走的,發現很短,來都是那段路.
剛才就自己過來,不需要勞煩曾牧師.
曾牧師今天出了去港島,也不需要安排其他弟兄姊妹節.
發現真的很方便,很快就到了.
今天跟大家分享的經文,剛才是王大兄帶我們讀了.
我給他一個題目,借錢當然要還.
剛才將PowerPoint下載的時候,已經有了.
下一句是甚麼?.
「別給錢中介」.
其實這個故事都會講中介的.
但怎樣才叫做「給錢中介」呢?.
所以就算了,借錢當然要還.
這是政府的宣傳片.

$^{41}$不要叫我們看著還款能力來借錢.
如果不是的話,就有很大問題.
會爆煲,被人叫洞,很多事情都出問題.
這個故事跟借有關係,但不是借錢.
是Hannah,是求子.
這裡有沒有姐妹叫Hannah?.
Hannah不是一個特別普通的外國女士的名字.
但我也認識一兩個Hannah.
Hannah是Grace,你叫Grace就等於希伯來文的Hannah.
上帝的恩典的意思.
求子,我自己是大兒子.
當我讀這段經文的時候,都會想起我自己的媽媽.
這篇不是母親節的訊息,我母親節就不是講這篇訊息.
不過這次想講這篇訊息.
我作為大兒子,媽媽在我出生之前已經是信主.
但她很久都沒有告訴我,我的由來.
當然是媽媽和爸爸生了我.
但問題是由於她在我出生之前已經是信主.
她有沒有在上帝面前做了手腳呢?我不知道.
但後來我讀了神學之後,她才告訴我.
她說,雲開,早就已經奉獻給上帝了.
你早點告訴我,不要走那麼多冤枉路.
十八歲就去讀神學了,不要等到我三十四歲才讀.
不要等那麼久.
這對我來說都是一個很特別的經驗.
當媽媽這樣跟我說的時候.
我家裡有兩個傳道人,我和小妹都是傳道人.
我五兄弟姐妹,跟著我那個也是傳道人.
不過她現在不是,她現在是做輔導的工作.
她自己在香港讀大學,做了幾年工作.
然後去讀神學,在自己的舞會服侍了好幾年.
所以這個故事是說我們向神求.
我們看看哈娜向神求子的經歷.
然後也反思上帝和我們之間的關係.
關係是很複雜的,關係有很多種方式.
其中一種關係裡面表達方式就是.
受受的關係.
受就是有剔手邊,沒有剔手邊.
有剔手邊的就是給出去.
沒有剔手邊的就是接收,接受的那個受.

$^{81}$受受是人與人之間的關係裡面必然會發生的.
尤其是跟朋友,如果你稱呼他為朋友.
總是會有受受之間的關係.
你喜歡為朋友做事,送禮物給朋友.
他作為你的朋友也會有不同的對你好.
可能會送禮物給你,可能會為你做一些事情.
就算一句知道你不舒服,一句問安的話.
也是跟受受有關係的.
你知道他紀念你,他讓你心裡面舒暢.
這些就是往往人與人之間的關係會出現的.
其中受受是一個很重要的層面.
也讓那個關係繼續能夠維持住.
如果是很久,或者你跟某人認識他.
但是從來都沒有受受的關係.
你就很難說你跟他很熟的.
我發覺我跟熟的朋友全部都有比較認真一點的.
經常或者很多次的,甚至是一些大的事情.
受受的關係.
你可以認識一個人二三十年從來都沒有受受.
你幫我,我幫你的那些關係.
關係一般來說很難說是深入的.
往往有受受的關係.
慢慢在接受別人的好處.
你給別人好處的過程裡面.
建立比較深刻的友誼.
人就是這樣.
你不要說是貪,你又想別人給你什麼.
你又沒有吝嗇,不想別人給你什麼.
但是你又不要把這個人與人之間的關係看成.
這樣啊,英文說很transactional.
很功利主義.
一來一往,我們中國人往往叫來往.
受受其實來往是另外一個說法.
一來就一往,是不是太過計數呢?.
可以變成很計數的.
可能你跟我認識的一些朋友.
其實我們平時也會有的.
有些在習俗上真的有一點點計數的味道.
但是如果真的是熟的朋友.
或者建立關係深厚的朋友都不會計較的.

$^{121}$但是疏的也會計數的.
最簡單的例子,雖然現在也不經常做.
就是過年的時候.
過年去拜年的時候派利是.
拜年我們也說會回禮的.
所有這些禮就是習俗.
不一定與關係深厚不深厚有關的.
因為可能親戚朋友也是一年見一次.
未必一定很深.
但是要回禮啊.
你收別人的利是,你一定帶著利是去別人家裡.
也要派的.
你不能拿別人的利是.
別人有小孩你也不派利是.
是沒有這件事的.
以後別人就不跟你拜年了.
所以這是人之常情.
在一些比較習俗的場合.
真的是一來一反來往是有預期的.
你不能說是功利,不能說是計較.
本來就是這樣的.
但是在真的朋友關係裡面.
一來就必有一往,不一定是.
因為在不同的時刻.
有時候是有需要的時刻.
對方會想起你,會陪著你,會為你做些事.
在另外一個時刻可能對方有需要你又做些事.
你又不會想著其實是還債.
而是單純因為你關心他,你想對他好.
所以你做些事.
所以不一定是這麼重的計數.
這樣的觀念.
我們跟神之間的關係又如何呢.
今天的故事就是你問神拿了一件事.
你怎樣看這件事呢.
這就是我們要看的.
經文裡面要思考的東西.
我們再一次祈禱.
我們在崇拜,來到你面前聚集的時候.
我們願意我們的心思意念.

$^{161}$我們口所說的,我們所唱的詩歌.
還讚我們的祈禱.
我們所擺上的奉獻.
都蒙你自己,願納讓你自己的名德榮耀.
我們知道在這裡敬拜.
不單是我們,也有很多其他屬於你.
是你自己的兒女在不同的.
不一定是香港,其他世界,其他時區.
都一樣聚集敬拜的.
都求天賦月.
是所有你自己的兒女在你面前.
誰聽我們禱告,逢主命求.
阿們.
好了,先問三個問題.
不用舉手了.
大家有沒有向上帝求過什麼.
應該都有.
我們向上帝祈求.
是很自然的事.
等於小朋友問父母要什麼.
是自然的.
某程度上是天經地義.
因為是他長大了.
小朋友自己做不到,拿不到東西.
在這個世界上,人都很厲害.
萬物之靈.
但有太多東西我們想要,得不到.
要兒子是其中一件事.
所以我懷疑大部份人都會舉手.
問題是,求有沒有得到過.
求會不會是.
求居然上帝就給你.
求未必一定會有.
又或者他給了你,你不一定覺察到.
但有沒有得到過呢.
如果有得到過,就第三條問題.
得到過,在上帝面前求.
他又按照我們所求.
有什麼回應呢.
這就是問題.

$^{201}$這就是借錢還的其中一個.
有沒有回應呢.
我們想的時候,會不會考慮.
上帝有沒有預期我們有回應呢.
如果上帝預期.
他私恩典,或者按照我們所求.
給我們東西的話.
他又預期有回應.
他會不會太計較呢.
就是我剛才所說的計數.
上帝是不是這麼計較呢.
會不會太計較我們的回應.
如果他真的計較的話.
他又想見到什麼回應呢.
這些可能都沒有什麼太清晰的問題.
今天的經文只是一個例子.
我們也不需要把它變成一個通例.
是一個求子的例子.
我們在這個世界求是很多東西.
有些你求了上帝給了你.
你都不知道怎麼回應.
譬如說我記得我中學的時候.
媽媽住屋村.
家境比較媽媽夫婦.
爸爸是做海員.
長年累月都不在家.
一年回來一次或兩次.
每次最多就是一個星期.
所以媽媽都經常不見老公.
不至於是守山寡.
我見過爸爸數數賣賣.
因為我中學畢業那年.
爸爸就去世了.
都是幾十次.
有記憶的話就沒有幾十次這麼多.
有一次爸爸在行船.
不在家.
媽媽突然肚子痛到.
你沒見過有人痛到.
在床上扭來扭去.

$^{241}$嚇死我.
我那時候讀中學.
讀中學也是長大了.
小不更事.
還有小妹四個.
小妹四個在那裡.
真的不知道怎麼做.
嚇到媽媽就說.
你為我祈禱吧.
把你床上的東西捲起來.
捆起來.
為媽媽祈禱.
求神醫治.
又不知道什麼事.
打999叫救護車.
媽媽又不肯.
沒得打999.
只能夠祈禱.
當然求上帝醫治媽媽.
無論是什麼問題.
希望媽媽能夠恢復.
最重要是不痛.
但起碼在短時期裡.
上帝沒有應允.
還是有點痛.
結果還是要進醫院.
進醫院原來是膽石.
膽石割了膽.
回來復康就沒事了.
我小時候.
一段時間已經在團契.
但就不明白.
我覺得小朋友向上帝祈禱.
上帝也不聽.
小朋友很天真.
十多歲也不太天真.
但也沒有惡念為媽媽祈禱.
為媽媽祈禱也是天經地義.
為什麼上帝不聽呢.
後來久了之後.

$^{281}$這件事一直在我心裡.
有一根東西在那裡.
為什麼上帝不聽祈禱呢.
第一我假設上帝一定會聽.
尤其是小朋友.
更加一定會聽.
第三就是我假設上帝不聽.
過了很長一段時間.
這件事一直在我心裡.
有一根東西在那裡.
後來我一直看聖經.
成長了.
讀了聖經又多了.
然後又讀神學.
發覺其實我自己很傻.
上帝有聽祈禱就醫好了.
問題解決了.
這件事跟漢娜後來有個兒子一模一樣.
漢娜可以向上帝求兒子.
那上帝怎樣給兒子呢.
就好像耶穌一樣.
童貞女神節.
不需要老公.
我給你兒子.
當然不會.
所以你會發覺.
這是我們剛才所讀的.
這是第十九節.
我用舊和合本.
剛才大家讀的是和修版.
第十九節說.
次日早他們起來在耶華面前敬拜.
漢娜那次去跟伊利加那她的老公.
上去聖殿一年一次敬拜之後.
她在上帝的殿門口站在那裡.
求上帝.
祭司伊利還以為她喝醉酒.
叫她不要喝那麼多酒.
這裡一路都很暗.
聽不到說什麼.

$^{321}$又不走又暗暗沉沉.
伊利就以為她喝醉酒.
後來漢娜就跟她解釋.
我是心裡受苦.
求耶和我說.
你不要當我是一個不正經的女人.
之類之類.
她說我只是在耶華面前求.
你平安走吧.
你所願耶華賜你.
將你所祈求的給過你.
她就回去了.
這就是第十九節.
次日清早他們起來在耶華面前敬拜.
就回到喇嘛.
到了嘉義後伊利加那和妻子漢娜同房.
她要生兒子始終都要同房.
夫妻行周公之禮.
我祈禱我是要什麼.
我是預期上帝什麼.
我祈禱.
媽媽就不痛了.
我是要求.
祈禱是求一個結果.
但我一直耿耿於懷的原因.
就是我祈禱沒有即時的作用.
而即時作用就不是醫好.
就不是普通醫治.
看醫生醫好.
就是要一個神蹟.
我十幾歲的心裡面.
我祈禱沒有作用.
就是因為我看不到神蹟.
媽媽要找醫生醫好.
找醫生醫好就不是神蹟.
要一下子康復才是神蹟.
就是說漢娜求子.
她不需要和老公行周公之禮.
生了個兒子就算是神蹟.
她和伊利加那同房.

$^{361}$然後後來懷孕生了個兒子就不算神蹟.
這個就是一個很過分.
這個我就當是我小時候.
做小孩子笑不更時.
就是說我不單求上帝醫治我媽媽.
我還求上帝要這樣醫治我媽媽.
連方法都列出來.
你不是這樣醫治我媽媽.
你還沒有醫治我媽媽.
醫生醫好就不是你醫的.
後來我就在上帝面前悔改.
但是這個思考思想是留了太久.
二十年了.
所以我經常都要矯正自己.
你向上帝祈求是一件事.
你不要要求上帝一定要照著你的方式.
將那件事實現給你.
你現在大了嗎.
我不單要你給我一份禮物.
我還要那份禮物一定要這樣包裝.
一定要那時候那個地方.
在那間公司買.
用那間公司來快遞給我.
什麼都要求.
所以在這些時候我們就要問.
其實我們要矯正.
很多時候上帝真的給了我們一份禮物.
按照我們要求.
按照我們請求.
給了我們一份禮物.
我們自己不察覺.
因為我們總覺得上帝沒有給我們.
因為我們求上帝的時候.
心目中連方法都告訴上帝.
你要透過這個方法給我.
而上帝繞個圈給你.
你已經不察覺.
冤枉上帝說上帝沒有給你.
其實上帝是治好了我媽媽的.
藉著去醫院.

$^{401}$醫生的手治好我媽媽.
所以這件事是我讀哈娜的故事的時候.
有一次勾起我過去的記憶.
求上帝要小心求.
要醒定一點求.
你是求.
不是要.
我要.
我要你這樣這樣.
是求.
求人的時候就要小心求.
但是我們回到《三寶寓記》上.
《三寶寓記》記載了以色列歷史的轉捩點.
讀《三寶寓記》上很多都是故事.
記載著以色列烈王時期的開始.
但是烈王時期的開始都是出於一個祈求.
今天全天都是講祈求.
所以紅字版求求.
《三寶寓記》的敘事本身不單止是以色列歷史裡面.
以色列人記載了第八章.
求要一個皇帝.
拉開了以色列王朝的序幕.
《三寶寓記》本身都是用一個求的故事開始的.
所以這個很特別.
以後你想《三寶寓記》上下.
猶太人是不分上下的.
猶太人的希伯來文聖經裡面.
《三寶寓記》上是一卷書.
所以就只有《三寶寓記》.
後來翻譯成希臘文的時候.
就分為上下.
所以以後教會裡面新約.
基督徒教會裡面對《三寶寓記》才分為上下.
但是我以後讀《三寶寓記》的時候.
就想著《三寶寓記》的故事是跟求有關的.
哈那求上帝.
以色列人求上帝.
各有不同的結果.
追得到也未必好.
這個要很久.

$^{441}$這些年紀大才知道的.
現在就一定要追追追.
最好是.
現在怕你生得少.
都不生了.
以前呢.
幾多個夠數啊.
兩個夠數.
怕你生太多.
我家我媽媽五顆.
我們這裡有沒有六顆七顆的.
我讀在美國讀書的時候.
有家美國人.
他們自己有農場.
生了多少顆啊.
十二顆啊.
哇.
真是我說.
十二顆是什麼觀念呢.
我都不知道.
可能爸爸媽媽叫子女的名字都混淆了.
無論如何.
在這些位置.
就一切都是跟祈求有關的.
追一下.
我們看剛才一張的故事.
就有幾個角色的.
有些生卡拉.
那些叫卡拉菲亞不是生卡拉.
卡拉菲亞我們就不理了.
主要角色.
伊利嘉娜.
老公.
哈娜就是女主角.
比尼娜就是對手.
敵對者.
然後伊利就是烏哈烏哈那個.
上帝僕人.
希望不是所有上帝僕人都是烏哈烏哈.
然後當然有上帝自己.

$^{481}$耶和華自己.
其實如果你看名字.
深水清還有兩個名字.
我沒有加上.
其實有很多名字.
因為講到伊利嘉娜的族譜的時候.
講她的.
曾祖父之上那個叫什麼.
太太.
她的元孫.
一直數數來.
那幾個根本就沒有角色.
只是族譜名我們就不理了.
伊利有兩個兒子.
何忽尼和菲尼哈.
都提了伊利當時做士師.
或者當祭師.
他有兩個兒子.
那兩個兒子是.
壞人.
往後就會知道.
在這裡只是提他的名字.
但是在這個故事裡面他們沒有角色.
所以我也不將他放在這裡.
先看伊利嘉娜.
伊利嘉娜喜歡愛不愛他老婆.
愛.
但是愛情.
不可以當什麼喝.
不可以當水喝.
愛情只是愛情不夠的.
就伊利嘉娜有錢.
你會發覺.
有錢人到今天也是.
有錢人.
娶了媳婦.
或者有錢人娶了老婆.
最重要是什麼.
你跟我生嘛.
是不是.

$^{521}$不然我的錢給誰.
傳不下去.
然後流入外人的手.
窮人也無所謂.
窮人可以生很多.
不知道為什麼.
非洲的生育率很高.
相對世界上窮的國家.
但是西方發達國家.
東方發達國家.
日本韓國香港台灣.
不生.
已經是沒有辦法.
接續維持人口的地步.
每個國家都要輸入人口.
所以伊利嘉娜當然要愛情.
但是也要子子.
兩邊都要.
所以有兩個老婆.
一個老婆生孩子.
一個老婆疼愛.
只有有錢人才做到.
其實開始的時候.
未必是這樣.
他怎麼知道.
誰有誰沒有.
如果哈娜是先的.
哈娜一直都沒有孩子.
他就娶第二個.
這個可能性比較高.
剛才我們讀聖經排位.
最先出現的都是哈娜的名字.
但是他又要愛情又要子子.
剛好兩個人.
很難擺平這兩方面的需要.
鍾無艷和夏迎春.
怎麼辦呢.
各有各的好處.
所以家庭是注定不和的.
而在不和當中.

$^{561}$誰會吃虧一點.
不是生了孩子的人會吃虧一點.
有愛情卻沒有子子.
是一定吃虧的.
可能今天世界也是這樣.
圍村的世界.
我都寫了一份報告.
如果今天還可以一夫多妻的話.
都是有孩子的和沒有孩子的.
但是老公愛沒有孩子的.
一定是沒有孩子的吃虧.
平時生活.
難道沒有孩子.
得到寵愛的那個.
跟對方說.
老公今天不配你.
都會有壓力的.
對方有孩子.
最終老公都會.
所有的東西都會給孩子的家.
所以一房二房都是看你生孩子.
我們在香港見到大有錢人.
都是一房二房三房在罵很多東西.
其實都是這樣.
所以家庭一定不和的.
注定不和.
而在這個情況下.
Hana注定.
尤其是古時代.
一定吃虧.
沒有孩子就一定吃虧.
而且伊利加那有些盟室.
盟室的盟字是這樣寫的.
我查過字典才寫的.
打下去.
他明白Hana的處境.
但是又明白她的心思.
她是一個幾無能的老公.
很被動.
她自己又要這樣又要那樣.

$^{601}$不專一.
她有子子的需要.
她又有愛情的需要.
她又疼這個.
但是疼的來.
她又不明白Hana的心.
Hana每次哭.
應該不是每年一次去上帝的殿獻祭.
吃燒肉.
燒肉不是豬肉.
他們不獻豬的.
牛或者羊.
才會哭.
平時都會哭.
背人垂淚.
不過伊利加那誰知道.
她哭.
不吃飯.
猶太人從來都不吃飯.
所以吃飯是中國人的說法.
翻譯的時候的說法.
不盡善.
我們正確一點的翻譯應該說.
一直哭.
第七節那裡.
哭泣不盡善.
第八節聖經就說.
她丈夫伊利加那對她說.
Hana你為什麼哭又不吃飯.
為什麼心裡受悶.
有我不是比有十個兒子好嗎.
這些就是很蒙塞的男人的說法.
你跟他說.
你有我比十個兒子好.
你又提醒他.
他一個兒子都沒有.
他只是沒有兒子.
我們姐妹.
丈夫應該怎麼說才會稍微好一點.
有我比有十個兒子好.

$^{641}$這樣說就很明顯不行.
因為每次都提醒他沒有兒子.
應該說.
我是愛你.
你是我心中唯一的一個.
從來都不提醒兒子.
不能提醒子女.
你看看這個世界上我最愛的就是你.
你還想怎樣.
這樣都好一點.
千萬不要提醒兒子.
但你不能說.
有我一個頂得住十個兒子.
這樣還可以嗎.
每個人都可以無恥.
有愛情就兒子都不用了.
你又過去二房那邊.
做什麼呢.
所以在這些位置.
他都是很蒙塞的.
這個就是伊利加拿大.
你會發覺在故事裡面.
伊利加拿大其實沒有什麼角色.
總之做了丈夫.
他又幫不到他所愛的哈娜.
幫不到.
他完全幫不到忙.
他說話又說得不好.
又令到哈娜媽媽傷心.
他自己兩家在這裡.
他在那裡一三五二四六.
就是這樣.
伊利加拿大.
另外一個男人.
伊利.
濟斯.
偏偏他兩個兒子.
何芬妮和菲尼哈.
現在不知道的.
往後知道他們兩個廢仔的父親.

$^{681}$兩個兒子都很廢.
廢我括住了.
這個是聖經裡面的說法.
這個聖經裡面的說法.
廢這個字.
在我們剛才讀的故事出現過.
不過就不是用來形容.
伊利那兩個兒子.
何芬妮和菲尼哈.
是哈娜說.
哈娜祈禱的時候.
伊利暗暗沉沉的站在那裡.
又不開眼.
總之暗暗沉沉一段時間.
伊利說他喝醉酒.
伊利說你不要再喝酒了.
夫人.
不是犯婦人.
夫人你不要再喝酒了.
哈娜就睜開眼.
我主啊.
其實不是.
我是心裡面仇婦的夫人.
清酒濃酒都沒有喝.
新和修版翻譯了清酒烈酒.
其實不是烈酒.
以前釀不到烈酒.
以前喝酒.
我們知道希臘羅馬人一直過去.
都是這樣.
就是你釀了出來.
一定混水的.
混到很稀的.
一定混水.
那些叫清酒.
不混水那些叫濃酒.
所以就不是白蘭地威士忌的烈酒.
啤酒就差一點.
就不是這樣分化.
所以仍然維持舊和學本翻譯好很多.

$^{721}$清酒濃酒都沒有喝.
在上帝面前傾心吐意.
不要將婢女漢做不正經.
和修版也翻譯了.
不要將婢女漢做不正經的女子.
不正經.
在二章十二節.
同樣用來描繪伊利那兩個兒子.
但是二章十二節如果你看.
翻譯成惡人.
他兩個兒子是惡人.
這裡不正經過頭惡人.
兩個都不是很好翻譯.
廢是最好翻譯.
其實字面的意思就是沒用.
useless.
在以前的世界.
老公娶了老婆.
老婆不生孩子.
基本上就沒用.
某個程度上.
這樣說很難聽.
我們現在男女平等.
你當我是生孩子機器.
不是這個意思.
但是那時候真的這樣.
家婆就會踢.
踢白燕.
踢櫃子.
家嫂.
家嫂當然生孩子.
很多孩子.
送很多豬耳繩給她.
結婚要送一大堆.
托著的豬耳繩.
好像豬一樣.
生十幾二十隻.
生出來.
趴在那裡.
喝奶就最好了.

$^{761}$都是這個意思.
一腳就被人踢出去.
所以沒用.
這裡翻譯做沒用是合理的.
但是翻譯做廢.
廢就真的沒用.
所以同樣以利那兩個兒子.
後來就被視為廢.
做不到事.
幫不到以利手.
做不到祭司.
Hannah對於利說.
我是一個仇婦的僕人.
不要把婢女看成廢女人.
以利就是兩個廢仔的爸爸.
沒教好那兩個兒子.
判斷能力又差.
人家明明在祈禱.
說人家喝醉酒.
在祈禱.
曾牧師說你不要喝那麼多酒.
曾牧師.
你需要換眼鏡嗎.
判斷能力成疑.
現階段大概只是一個廟祝.
只是看著檔口.
做不到什麼.
帶領以色列人都出問題.
以利家還會持續一段.
會等撒母耳.
Hannah生了一個兒子叫撒母耳.
長大之後.
整個家族都會被廢.
被上帝廢了.
接著是女人.
女人厲害了.
比尼娜.
比尼娜有子並非萬事足.
因為老公不愛她.
有兒子沒愛.

$^{801}$都是很頭痛的事情.
所以她也有切身之痛.
切膚之痛.
也不能夠太怪她.
不過.
雖然她是.
老公不愛她.
她知道老公愛.
另外一個女人叫Hannah.
可能我們當Hannah是阿弟.
老公愛阿弟.
但我和她生小孩.
她偏偏老是在阿弟那邊.
過到來好像心不甘情不願.
她也是受害者.
我們也要明白.
女人在家庭裡沒有老公的愛.
其實也很慘.
當她是一部機器.
你有貢獻.
你有沒有聽過老公跟你說.
你在家庭裡很有貢獻.
什麼貢獻.
你跟著問他什麼貢獻.
生小孩.
你一巴掌打過去.
什麼貢獻.
貢獻你當我是什麼.
她是典型的受害者.
但她又變成加害者.
她本來是受害者.
但她看不順眼就加害Hannah身上.
打她.
你知道生不出小孩的痛也是很痛.
打她一輩子也很慘.
受害者變成加害者也很慘.
可能大家也覺得自己人生活了幾十歲.
某段時間被人害過.
做過受害者.
但我們要想.

$^{841}$做受害者的時候有沒有變成加害者.
受害者受害到一個地步.
變成欺負其他人.
被人欺負我們也會變成欺負其他人.
要小心.
千萬要小心.
別人欺負你誰壞.
他壞.
別人欺負你然後你再欺負別人.
扭曲了我們的心理.
變態了.
然後再欺負別人.
誰壞.
他壞.
你也壞.
所以要小心.
被人欺負就變成這樣.
Hana.
愛有什麼用.
愛可以拿來吃飯嗎.
沒用的.
愛不能替補遺憾.
愛也不能停止傷害.
她被二婆欺負.
老公對她的愛也不能停止.
制止不了.
沒辦法制止.
因為一見到對方的小孩.
對方不需要說話.
對方很開心.
跟小孩玩耍.
有空過來.
今天天氣很好.
抱著小孩過來.
已經刺激到你.
七曉新煙.
你過來過來.
抱著小孩過來.
幹什麼.
給我看.

$^{881}$簡直是眼淵.
所以停止不了傷害.
問題不一定是比利娜身上.
是她自己身上.
她想有卻沒有.
所以丈夫愛也停止不了.
傷害不一定是比利娜傷害.
傷害好像丈夫呼呼.
十個兒子.
丈夫也會被傷害.
她自己也會傷害自己.
想一邊.
執一邊.
我們想一邊就自己傷害自己.
所以當我們很灰心.
或者很不開心.
很糾結.
很低落的時候.
就不要想一邊.
想一邊.
出不來.
因為其實這樣是傷害自己.
有時要聽其他人的話.
她沒辦法依賴丈夫.
Hannah的突破是怎樣.
當然Hannah比較多.
Hannah的突破是她沒有始寵新嬌.
或者她無法.
她也沒有Jackson始寵新嬌.
但是她決定不做加害者.
她沒有做加害者.
你發覺整個故事裡面.
就算Hannah後來生了孩子.
她也沒有傳給你.
只有你有嗎.
我也有.
她也沒有.
你記得雅各的兩個老婆.
兩個姐妹.
你和拉傑就是這樣鬥生孩子.

$^{921}$我生得不夠.
這陣子沒有.
我找個妹仔.
我的妹仔給老公.
等妹仔生孩子.
那些兒子是我的.
兩邊都轉轉貢.
Hannah沒有.
她不相信命運.
只相信夜話.
她將事情.
這是她的突破.
她身邊的男人都幫不到她.
她認為是.
以妮加拿幫不到她.
現在生不到孩子.
那時候生不到孩子一定是賴誰.
賴女人.
尤其是以妮加拿和比尼拿生不到孩子.
當然是賴女人生不到.
男人沒事當然是你.
現在男人都會去做檢查.
不會只是賴女人.
所以現在公平很多.
是你精子數目不夠.
補一下.
兩邊都會.
但是以前一定是賴女人.
但是她發覺.
賴女人.
男人有沒有幫忙.
幫不到忙.
以妮加拿幫不到忙.
以妮的祭司.
他去主殿祈禱.
以妮也在呼呼.
也幫不到忙.
她只能夠直接將她的個案.
直接帶到上帝面前.
這就是她威風的地方.

$^{961}$以前世界是富裕社會.
男人最威風.
在地上男人就是最威風的.
尤其是這兩個男人更加威風.
老公和祭司.
兩個都不找.
直接就找上帝.
月級就找上帝.
最大的那個.
這就是我們今天.
也要在哈拿身上學習的.
跳過那些都做不了的.
直接找老闆.
大老闆就是上帝.
不求問中介人士.
借錢的中介找中介.
給錢的中介就是這個意思.
借錢的中介說不是給錢的中介.
直接去求問耶和華.
哈拿在這個位置就極端懂事.
祈求的往來本質我們不理.
哈拿許的願在哪裡呢.
一章十一節.
我們讀出來.
一章十一節.
她所許的願其實在這裡.
這裡我已經將經文列了出來.
她許願就是說.
十一節是她在上帝面前.
哭泣祈禱.
直接記載她所說的話.
這是以利.
以利在那裡含血.
你在那裡喝那麼多酒幹什麼.
借錢的含血.
所以這就是她許的願.
前面還有很多東西.
求上帝眷顧婢女的苦情.
眷顧不忘婢女.
然後下面這兩句.

$^{1001}$耶和華若給婢女一個兒子.
我就將她給耶和華.
一個我括弧裡有個刺.
當然你求上帝給其實是刺給.
但是你將兒子給上帝.
就是還給.
所以刺和還.
但願問是同一個字.
給.
求上帝如果給我一個兒子.
給小女子一個兒子.
我就將他給耶和華.
很有趣.
這就是一來一往.
馬上回應.
這不是禮.
而是說我就是要有個兒子.
有了兒子之後我不要也可以.
這也是很古怪的想法.
當然有了兒子就不會不要.
所以我們讀到後面的時候.
又是一年又一年又要上去獻祭的時候.
Hannah怎麼說.
跟老公說你上去吧.
我在這裡照顧兒子.
還沒戒奶.
先不上去.
她真的想儘量留兒子留久一點.
以前戒奶可能去到四五歲.
現在四五歲.
留到四五歲小小的.
可以走可以走的時候.
才將他帶到侍羅那裡.
給過兒女.
在上帝面前服侍.
所以那時候她應許上帝許願.
你求上帝給你這麼重要的東西.
上帝說.
這樣嗎.
你有什麼東西給我.

$^{1041}$她覺得上帝會這樣想.
所以她立即向他求.
表示出她的真誠.
我將它給你.
這是她許的願.
到了還願.
就是我們剛才所讀的最後兩節.
27 28 兩節.
這個是戴了奶.
然後再上去的時候.
就帶小十母耳上去.
又去到耳裡面前.
我讀回經文.
26節.
哈娜聖經說.
那個婦人.
當時連名字都不用.
哈娜說主啊.
我敢在你面前起誓.
從前在這裡站著.
祈求耶和華的女人.
就是我.
不知道耳裡記不記得.
接著27節.
我祈求為得這個兒子耶和華.
已經將我所求的賜給我了.
所以我將這個孩子歸於耶和華.
使他終身歸於耶和華.
但是.
中文很通順的.
很理解的.
但是那個求字.
其實譯得不好.
我在這裡再用叻吻一點.
直接將希伯來文的味道帶出來.
哈娜說.
我就是那個女人了.
27節說.
我為這個孩子祈禱.
耶和華賜我所請求他的那個請求.

$^{1081}$就是耶和華賜給我.
我求他的那個求.
就是很強調的求.
那樣東西.
耶和華就給了我求的那樣東西給我了.
所以.
我也按照他所求給我的借給耶和華.
很有趣的.
他這裡用了.
那個字可以就這樣給.
但是在這裡反而借比較好一點.
他按照所求的借給耶和華.
而他終身的年日.
就是被借給耶和華.
或者就這樣給耶和華都可以.
所以在這些位置上.
我們發覺.
哈娜對上帝還願的時候.
她也有一點不情願.
借給耶和華.
希望將來可能會不會拿回來.
但是她明白一件事.
就是她開始許願的時候.
已經說了.
說了.
是要還的.
所以這個就是借錢當然要還.
因為.
她其實是問耶和華借回來.
差不多.
因為她要還.
而她記得.
她就要還.
所以還願的時候.
就將兒子.
她也盡量拖延時間.
借奶遲一點.
所以這個就是.
我在這裡標題.
極端懂事的哈娜.

$^{1121}$她完全知道她要做什麼.
她不賴貓.
不賴帳.
她向上帝求.
上帝給她.
她當時求.
也說我將兒子會給你.
給回你的.
終身服侍.
她就在這裡.
將上帝給她的東西.
又給回上帝.
所以一切都是和給這個字有關.
和求有關.
我求上帝的東西.
求在希伯來文是很有趣的.
求如果做被動的語態.
又可以變成.
不是求.
倒過來是被求.
借給對方.
而後面就會提到.
以色列求皇帝.
求皇帝.
以色列人唾棄上帝.
所以結果就很壞.
同樣是求.
上帝同樣都給他.
給以色列人.
但結果就很壞.
在這裡.
哈娜求子.
上帝給了撒母耳給她.
結果是好的.
但是去到以色列人求皇帝的時候.
上帝給了皇帝.
都是同樣的字.
求.
結果就很壞.
而且這裡和後面有關.

$^{1161}$我終生的年日.
他終生的年日要被借給.
被借給.
這三個字.
希伯來文是.
Shah-u.
讀法是Shah-u.
Shah-u.
會令到你想起.
一個什麼人的名字.
其實就是蘇羅.
蘇羅.
英文是sor.
其實蘇羅那個名.
希伯來文.
就是exactly就是Shah-u.
一點都沒有變.
他終生的年日是.
要被耶和華.
做Shah-u.
後來第一個以色列皇帝.
叫做Shah-u.
Shah-u就是求回來的意思.
那爸爸改兒子的名字叫求回來.
就是說蘇羅的爸爸都求兒子.
而且求到將來高大衛猛的兒子.
他也要求回來.
他也是以色列人向上帝求回來的.
但是兩家對求的東西.
沒有悉切的回應.
所以結果問題就很大.
當然上帝也是.
那個角色.
最重要的角色.
在故事裡面我們看到上帝掌管一切.
生孩子不生孩子都是上帝掌管的.
爾尼加拿和哈娜同房.
不知道多少次了.
每次都沒有.
就是那一次上帝顧念哈娜.

$^{1201}$她懷孕生子.
我們學院有個老師.
以前也是我們學院的同學.
兩夫妻一起進來讀書.
我還記得我教她的時候.
兩夫妻一起讀書.
但是讀了一年.
太太同班.
讀了一年太太就不讀了.
我問她為什麼呢.
後來才發現原來她是小產.
他們想生孩子生了很久.
第一個第二個都是流產的.
所以你想要又沒有.
很慘.
結果他們就在上帝面前求.
求到了.
他們兩夫妻就寫了一本書.
這本書還可以賣的.
你去基督還有得賣的.
是台灣校園出的.
叫做給我一個Baby.
這本書名叫做給我一個Baby.
他說上帝給.
還有一個女兒.
一男一女.
上帝給他們這樣求.
一直都生不了.
中間就沒有了.
身體很慘.
媽媽的身體很差.
如果這樣教的話.
會很虛的.
最終上帝給了他們.
這簡直就是一個神蹟.
我不知道他有沒有跟Hannah說.
你給我我就給你.
我就不知道.
我沒有問他.
這次預備講章之後.

$^{1241}$我都覺得我回去應該問他.
你有沒有說把兒子給上帝.
但他就寫了一本書.
他講明在書裡面.
這本身對他們兩夫妻來說都是一個神蹟.
上帝掌管的.
生孩子用現代科技.
他用了現代科技的.
因為生不了用了現代科技.
都是上帝掌管.
上帝顧念.
讓人能夠得到他們所求的.
而在第二章一至十節.
很重要.
記載了Hannah的.
作了一首歌.
歌頌上帝.
我們就不看了.
其中一句就是.
上帝由沙塵塵土中抬舉人.
在糞堆.
其他的詩歌.
帝玻拉作曲.
米利安作曲.
摩西作曲.
大衛也作詩歌.
都提到.
是上帝抬舉.
是窮乏人.
有需要的人.
在糞堆裡面.
將人抬舉出來.
清理乾淨.
將你提起.
這裡就講到塵埃裡面.
上帝抬舉窮乏人.
上帝給予.
我們就需要還.
尤其是許願.
很多時候求.

$^{1281}$我求上帝醫好媽媽.
我都不知道還什麼給她.
我當時當然沒有.
起碼我記憶當中.
沒有跟上帝說.
你醫好我媽媽.
我就怎樣怎樣.
我就去讀神學.
肯定沒有這樣說.
你說是否求上帝.
都要有些東西還呢.
我覺得不是.
因為我們求的五花八門.
有些東西真的不能還.
而且是求上帝幫你媽媽.
難道你叫媽媽去還嗎.
那是怎樣呢.
所以還不還.
往往是看你求的時候.
你怎樣求.
你說那我知道了.
哈娜的故事.
教訓是什麼.
千萬不要答應上帝.
你會給他兒子.
不要許願.
求還求.
還還還.
千萬不要求到一個好東西.
還給上帝.
你不要這樣說.
所有事情都是出自內心.
一會兒我們就會奉獻.
一定是主席都會跟你說.
是甘心樂意.
所有事情都是來自上帝.
但是我們一定要從哈娜的故事中.
看懂一件事.
在上帝面前.
我們就需要學做.

$^{1321}$好像哈娜一個極端.
我叫做極端懂事的女人.
極端懂事.
極端到連兒子都給上帝.
這個其實是極端了一點.
我做拿西爾人也是一段時間.
他說一輩子.
一輩子都在上帝面前服侍.
是極端了一點.
但是問題是她懂事.
她知道上帝給她兒子.
是她本身沒有的.
也沒有說我配有.
誰說你一定是配生兒子或女兒.
既然她求上帝.
上帝給她.
她求的時候也說明.
我會將兒子歸你.
給你.
她就這樣做.
所以在這些時候.
哈娜可以說是極端懂事.
以後我們求上帝什麼都好.
起碼得到上帝英雲的時候.
我們也需要存一個感恩的心.
最起碼就是一個感激的心.
我們說多謝.
你試過有些人你幫了他一事.
或者你給了一點點好處.
他連多謝都沒有.
你試過碰過這些人嗎.
有啊.
可能有些甚至是朋友.
一般來說通常都比較三寶.
你會感覺不良好.
因為你預了.
我不是叫你.
我給你就是給你.
不需要你還.
但是有些表示.

$^{1361}$人家跟你開門.
你也說一聲謝謝.
沒有碰過.
這是基本禮儀.
是潤滑劑.
讓人與人之間的關係.
就算與陌生人.
不認識的.
三不認識七的.
關係都能夠好的.
當然你說這樣很搞.
去到日本.
但是有禮貌.
也是日本人的好處.
我們去日本舒服.
就是覺得別人對我們有禮貌.
你真的舒服.
你不要去日本.
你去日本做什麼.
你去日本就是因為舒服.
他心裡怎麼想你都不管.
但起碼他服務是舒服.
禮貌禮儀是舒服.
去到地鐵上是靜悄悄.
這些是好處.
當然去到極端也很虛假.
有外面沒有裡面.
裡面不得不捏死你.
外面就有個公.
這個當然是病態.
如果去到極端.
但是仍然禮.
中國人的禮是重要的.
尤其是在一來一回.
一來一往的事情上.
在求的事情上.
在受的事情上.
而這件事我們可以向Hana學習.
所以今天就跟大家分享.
如何面對上帝將各樣好處賜予我們.

$^{1401}$我們要想想.
我們一起祈禱.
天父我們感恩.
你在我們出生的時候.
甚至可能是還沒出生之前.
你已經應允了.
如果是我們父母有求的時候.
已經應允了他們.
我們今天能夠成為何等樣的人.
都是你自己在我們生命中提攜的帶領.
尤其是我們在生命中居然能夠認識你.
成為你自己的兒女.
能夠更多的與你親近.
能夠更多的經歷你給我們各式各樣的好處.
我們感恩.
為了你在主耶穌基督裡面給我們最大的好處.
我們特別深被恩感.
求天父都幫助我們.
在生命中成為懂得如何做感恩的人.
不單止是與你相處.
與你關係.
即使是人與人之間的關係.
願我們同在.
誰聽我們禱告.
奉主明求.
\newpage



\section{約書亞記 14:6-15}
\label{sec:HMy_7Yhf_bE}
\textbf{2025.06.01 《憑信得著應許的勇者――迦勒》 經文:約書亞記14\_6-15 講員 : 鄧泰康傳道}
\newline
\newline
連結: \href{https://youtube.com/watch?v=HMy-7Yhf-bE}{\texttt{https://youtube.com/watch?v=HMy-7Yhf-bE}} ~~~~ 語音日期: 2025-06-06
\newline
\newline
\hyperref[sec:xtvbpi6Ud4M]{\small{< < < PREV SERMON < < <}}
~
\hyperref[sec:index_chronic]{\small{[返順時目]}}
~
\hyperref[sec:index_scriptual]{\small{[返順卷目]}}
~
\hyperref[sec:iSHaOgcIOTI]{\small{> > > NEXT SERMON > > >}}
\newline
\newline
約書亞記 14:6-15
\newline
\begin{longtable}{cl}
\hline
\hline
章節 & 經文 (和合本修訂版)\\
\hline
14:6 & \begin{tabularx}{0.7\textwidth}{X} 猶大人來到吉甲，約書亞那裡，基尼洗族耶孚尼的兒子迦勒對約書亞說：「耶和華在加低斯‧巴尼亞指著我和你對神人摩西所說的話，你都知道。 \end{tabularx} \\ \\ \relax
14:7 & \begin{tabularx}{0.7\textwidth}{X} 耶和華的僕人摩西從加低斯‧巴尼亞差派我窺探這地的時候，我剛四十歲。我把心裡的話向他報告。 \end{tabularx} \\ \\ \relax
14:8 & \begin{tabularx}{0.7\textwidth}{X} 雖然同我上去的眾弟兄使百姓膽戰心驚，我仍然專心跟從耶和華－我的神。 \end{tabularx} \\ \\ \relax
14:9 & \begin{tabularx}{0.7\textwidth}{X} 那日，摩西起誓說：『你腳所踏之地必要歸你和你的子孫永遠為業，因為你專心跟從耶和華－我的神。』 \end{tabularx} \\ \\ \relax
14:10 & \begin{tabularx}{0.7\textwidth}{X} 現在，看哪，耶和華照他所說的使我活了這四十五年。當以色列人在曠野飄流的時候，耶和華曾對摩西說了這話。現在，看哪，我已經八十五歲了。 \end{tabularx} \\ \\ \relax
14:11 & \begin{tabularx}{0.7\textwidth}{X} 現今我還很健壯，像摩西差派我去的那天一樣；無論是戰爭，是出入，我現在的力量和那時的力量一樣。 \end{tabularx} \\ \\ \relax
14:12 & \begin{tabularx}{0.7\textwidth}{X} 請你將耶和華那日所說的這山區給我。那日你也曾聽說，這裡有亞衲族人，以及寬大堅固的城，或許耶和華會照他所說的與我同在，我就把他們趕出去。」 \end{tabularx} \\ \\ \relax
14:13 & \begin{tabularx}{0.7\textwidth}{X} 於是約書亞為耶孚尼的兒子迦勒祝福，把希伯崙給他為業。 \end{tabularx} \\ \\ \relax
14:14 & \begin{tabularx}{0.7\textwidth}{X} 所以希伯崙成了基尼洗族耶孚尼的兒子迦勒的產業，直到今日，因為他專心跟從耶和華－以色列的神。 \end{tabularx} \\ \\ \relax
14:15 & \begin{tabularx}{0.7\textwidth}{X} 希伯崙從前名叫基列‧亞巴；亞巴是亞衲族最尊貴的人。於是國中太平，沒有戰爭了。 \end{tabularx} \\ \\
[1ex]
\hline
\hline
\end{longtable}
$^{1}$電影節目早安.
今天很感恩來到大家當中跟大家分享.
嘉勒這個角色其實很特別.
在我看來是一個很成功.
但我要形容他是一個很低調的人.
其實你可以看在他的所謂成就.
特別是他當時應該和約書亞平起平坐.
但你會發覺很有趣.
在後段好像嘉勒已經消失了.
一直是約書亞帶領.
所以我說這個人很特別.
其實如果大家有留意聖經.
你會發覺這個人在神的眼中.
可能比約書亞更重要.
其實嘉勒的第一次出現是在民數記.
大家都很清楚.
當上帝差派十二個探子去到迦南地的時候.
那些人去了四十天後.
歡天喜地拿著土產回來.
很豐富很漂亮.
那些人應該很興奮.
上帝賜給我們這個地方原來這麼豐富.
不過發生了一件事.
裡面有十個探子.
他們回來後結果是怎樣呢.
我讀給大家聽.
在民數記十三章.
我又告訴摩西說.
我們到了你所打發我們去的那地.
果然是流那雨蜜之地.
這就是那地的果子.
應該很興奮.
然而往那地的民強壯.
誠邑也堅固寬大.
所以我在這裡提醒大家.
做基督徒最忌的就是什麼呢.
是我知道.
不過最忌這件事.
我要告訴大家.
之後他們是怎樣看待這件事呢.

$^{41}$但那些與他同去的人說.
我們不能想去攻打那民.
因為他們被我們強壯.
探子中有人論道所窺探之地.
向以色列人報惡信說.
我們所窺探經過之地.
是吞吃居民之地.
我們在那裡所看見的人民都身量高大.
我們在那裡看見阿納族人就是偉人.
他們是偉人的後裔.
據我們看就如渣蠻一樣.
據他們看我們也是如此.
當時派了十二個探子出去.
其中有十個探子.
他們就有這個反應.
唯一的就是你會發覺.
迦勒他不是.
迦勒和約書也不是.
聖經裡面怎樣說呢.
窺探那地的人中.
亂的兒子約書也和耶夫尼的兒子迦勒.
撕裂衣服.
對以色列傳媒眾說.
我們所窺探經過之地是極美之地.
耶和華若喜悅我們.
就必將我們領進那地.
把地賜給我們.
那地原是流奶與麥之地.
但你們不可背叛耶和華.
也不要怕那地的居民.
因為他們是我們的食物.
並且任被他們的已經離開他們.
有耶和華與我們同在.
不要怕他們.
我不知道大家有沒有看到對比.
在這十個人裡面的時候.
他們眼中就是這些困難.
這些人是又高又大.
是大到我們高不可攀的.
他們是偉人.

$^{81}$我們是甚麼?.
砸猛.
留意對比.
但在剛才那段經文是甚麼?.
他們是我們的食物.
用今天很貼地的話.
我吃硬他.
你明白嗎?.
他們不是看著這些所謂偉大.
他們看著的是甚麼?.
因為有耶和華與我們同在.
我想等下會開始記得這些很重要的片段.
但你會發覺為何說迦勒這麼特別呢?.
其實第一個發聲反對的不是約書尤.
那是迦勒.
迦勒在摩西面前安撫百姓說.
我們立刻上去德拉地壩.
我們足以能成.
是迦勒發聲.
而你會發覺是甚麼呢?.
在上帝眼中.
當他稱讚的時候.
亦都是只有迦勒.
唯獨我的僕人迦勒因他另有一個心志.
專一跟從我.
我就把他領進他所去過的那地.
他的後裔也必得那地為業.
今次讀這段經文.
整件事的發聲是由這裡來的.
所以我剛才說.
我很肯定的.
在上帝眼中.
迦勒有一個很特別的位置.
而這個人不簡單的地方是甚麼呢?.
你會發覺其實他很忠於上帝.
很信靠上帝.
所以在他心目中只有上帝.
就因為這種信靠和肯定.
所以他得到了上帝所給他的回應.
而他就能得到今天所帶來的這麼大的祝福.

$^{121}$這個因果關係.
大家作為基督徒一定要記得.
我經常說.
如果你回到教會.
我們沒有了上帝.
大家不是確信上帝.
我要很坦白地說.
我不知道你來做甚麼.
如果你說我純粹住在附近.
我來到教會是找朋友.
我認識一些朋友.
這個在最初你來教會的時候.
動機是沒有問題的.
最初是沒有問題的.
但問題是.
如果到最後.
你回到教會一段日子.
你仍然只是因為認識Thomas.
和Thomas很要好.
所以每次回來都會找Thomas.
所以有些人出現了一個情況.
有些弟兄姊妹不在就不在.
他回來我就回來.
他不回來我就不回來.
你就會在這裡見到.
但我想告訴弟兄姊妹.
如果你真的這樣想的話.
你回到教會.
你極其量只得到一個甚麼.
友誼.
我可以夠膽告訴大家.
如果你說好玩.
教會絕對比不上北上班的人.
是不是.
你要好玩.
你要好吃.
你要很開心.
我可以說.
在教會絕對沒有他們那麼開心.
但你來到教會的時候.

$^{161}$你不是要找這些友誼.
你是要找上帝.
現在很清楚的.
讓他得到45年後的祝福.
就是因為他確信上帝的每一句說話.
你來到教會.
你要得到真正的祝福.
你唯有走這條路.
你才能夠得到這種祝福.
所以大家回來的時候.
千萬不要搞錯.
千萬不要搞錯了焦點.
我經常說.
你回來崇拜的時候.
同樣我都說.
其實我今天雖然跟大家講到.
但你要去聽的是甚麼.
是後面十字架的那位耶穌.
是上帝.
你回來崇拜他.
我自己分享經文的時候.
我都要很小心.
我講個例子.
最近發生了一件事.
我不講背景.
我知道發生了.
有一間教會.
來了一個講員的時候.
上了台的時候.
他分享了45分鐘.
全部都是講他自己的事情.
是完全沒有講聖經的.
我叫做.
按照程序.
擺了一段經文出來.
但45分鐘我從來沒有講過聖經.
接著就出現一個情況.
那些心水清的弟兄姊妹.
他們就會走來跟牧師講.
阿... 牧師.

$^{201}$今天這個講員.
他完全沒有講聖經的.
下次還請不請他.
但很有趣.
後來有些弟兄姊妹.
走來跟牧師講.
今天這個講員很棒.
知道為什麼嗎.
很多笑話.
你看了45分鐘.
洞篤笑.
你明白我的意思嗎.
也值得的.
你不用錢給了45分鐘.
黃子華也要收不知多少錢.
這種情況下.
有些人很開心.
但你會見到兩個很不同的向度.
你今天來聽哪個講什麼.
當然我也提醒自己.
我也不希望我講到的時候.
下面的弟兄姊妹都睡了.
所以我盡量講到大家不要睡.
這是我每次祈禱的禱告.
所以你會發覺.
迦勒得到什麼祝福.
其實很不簡單.
第一 應許之地.
這就是當日他們要進去的地方.
想得到的地方.
而就是因為這份信的緣故.
他們特別得到.
你要留意.
其實他們那群人.
最初也是想得到這個地方.
但因為你不信上帝.
你看到身邊環境的緣故.
他們就要後來走了40年.
而按聖經的記載.
除了迦勒和約書之外.

$^{241}$這群男人全部死在曠野裡.
即是你不信上帝.
你是沒有資格進去.
這個應許之地的.
所以你一定要相信.
而他得到一個很特別的地方.
就是希伯倫.
待會我才告訴大家.
希伯倫有什麼重要的地位.
而他得到這個地方.
為什麼我說他比約書更加厲害.
就是因為全個族裡.
唯一得到他.
是以一個家族個人名義.
得到上帝的分地.
其他全部都是以一個族.
大家有什麼留意.
後面就是猶大族.
猶大之派.
約瑟之派.
是這樣分的.
但唯一得到他.
所以這是很特別的.
第二樣東西是很多人都很渴望的.
就是身體健康.
看啊我現今85歲了.
我還是強壯.
像摩西打發我去的那一天.
無論是征戰的出入.
力量那時如何.
現在還是如何.
在座的都有很多長者.
長者很喜歡說什麼.
我老了.
我沒用了.
你們做吧.
我不知道大家想不想得到約書.
我40歲的時候是這麼健康.
我今天85歲也是這麼健康.
大家應該要有這種魄力.

$^{281}$我相信這是大家很想得到的.
我媽媽也說的.
我媽媽今年80多歲了.
我每次和她的時候.
都是一個典型的長者.
一問她的時候.
我沒用了.
因為我媽媽每天只做幾件事.
就是睡覺.
起床.
吃東西.
看電視.
然後去廁所.
然後睡覺.
再吃東西.
然後再看電視.
睡覺.
她說我就是這樣.
所以你再說得差一點.
我等死.
我媽媽經常這樣說.
但我經常都和她說.
我媽媽說.
除非你給我健康.
我多少歲都沒問題.
可以走路.
90歲都可以.
但弟兄姊妹.
我想告訴大家.
為什麼嘉勒能夠有這樣的祝福.
因為他一直為上帝爭戰.
由40歲.
他一直和他們打.
打到現在85歲.
他沒有停過.
大家還要留意後面的經文.
他和阿約書亞怎麼說.
如果是上帝容許的話.
就讓我繼續打下去.
所以後面的經文.

$^{321}$是他帶頭85歲的伯.
帶頭整個群族走去.
打了這個地方回來.
所以我想告訴大家.
你想健康.
為上帝爭戰.
什麼叫為上帝爭戰.
不是叫你打架.
是你的心裡.
一直都渴望著.
上帝我要為你做多一點.
大家有沒有聽過聖經裡有句話.
凡扣門的.
我就給你開門.
這個扣門.
不是你扣什麼都可以.
不然就變成了黃大仙.
這個扣門是什麼.
你扣門.
你想為我工作.
你想做我的事情.
我一定開門給你.
領子妹今天你去禱告.
你多數求什麼.
除了三餐裡求上帝潔淨之外.
除了在你病的時候.
求上帝醫治之外.
除了平時你為家裡的家人去禱告之外.
你還為上帝.
你還祈禱什麼.
我不是說剛才那些不對.
那些是對的.
但是我想說是要更好.
你禱告當中.
為什麼你不求上帝.
上帝你就賜給我健康.
我能夠繼續為你爭戰.
所以我可以告訴頂子妹.
什麼人最有魄力.
眼中心中只有上帝的人最有魄力.

$^{361}$我可以告訴你.
我們教會有一個長者是寶.
為什麼呢.
我們叫二妹.
在教會內你都認識.
二妹以前很厲害的一個姐妹.
後來中風.
現在要工人姐姐推她回來.
但是她每個星期都回來.
她不是住在錦繡.
她住在元朗.
她每個星期都要坐那班車.
因為那班車是方便上輪椅.
她每個星期都準時要上那班車.
回來的時候.
崇拜完還要上主教學.
她是上完主教學之後才走的.
而還有一件事.
她最厲害的是什麼.
就是每天在群眾發emoji.
她會在群眾發一則關於基督教的emoji.
上帝祝福你.
上帝愛你.
我經常跟她說.
繼續發.
發到你有一天安息主懷.
真的.
每次你看到她的時候.
其實她仍然是一個很有魄力的婆婆.
雖然已經要坐輪椅.
已經要被人照顧.
但她從來都沒有說要躺在家裡.
我相信很多人今天都沒有動力出街.
行動不方便.
每次出來都會有壓力.
但她不是這樣.
如果你想得到上帝給你的身體健康祝福.
請你一直為上帝爭戰.
另外一件事.
是整個族裡面最大的尊榮.

$^{401}$為什麼我剛才再強調.
在神的眼中.
其實約書亞不及她.
原因就是.
我不知道大家有沒有留意剛才那段經文.
讀到最後的時候.
突然間有答了一段.
希伯倫從前名叫基列阿巴.
阿巴是阿納族中最尊大的人.
有沒有留意.
而接著到了他.
拿了這個地方之後.
他再說阿巴是阿納族的始祖.
為什麼呢.
剛才我說希伯倫是有一個很特別的意義.
原來就是整個迦南族.
他們的發源地.
而他們最尊貴的人.
原來就在這裡.
我用今天的例子.
為什麼人們這麼喜歡住在北京.
中國的首都北京.
因為那裡有什麼.
我假設那些人是跟毛主席的.
那裡就是毛主席的發源地.
他就在那裡立國.
甚至現在人們去大西北.
很喜歡去那裡看.
我以前又怎樣怎樣.
我想大家留意.
意思就是整個族裡面.
他們都很尊重的一個人.
其實上帝特意將這個地方給他.
這個地方就代表了.
你將會是我們整個族裡面最尊大的人.
他得到這些祝福.
不只是得到分地這麼簡單.
而是他得到了上帝完完全全的確定.
這個是他個人能夠得到的.
第二件事就是他對整個群族的祝福.

$^{441}$大家留意.
當他猜拼12個族長的時候.
其實12個是族長.
12個探子不是二打六.
是族長.
是每一個支派的族長.
即是頭兒.
其實這個頭兒是很影響整個支派.
所以你就會發覺.
其實他的信心.
你看到後來是很不同.
因為他成為了猶大支派的一個榜樣.
所以你會發覺.
在整個南北分地.
在分了地之後.
南邊最大的是誰?.
猶大.
你會發覺整個群族裡面.
某程度上帶頭的就是猶大.
在之後的時事記當中.
當他們要去得到上帝所應許之地.
上帝說了之後.
只有猶大支派是完成了所有分地的佔據.
其他支派都沒有.
有些就沒有趕出.
有些就趕不出.
有些就特意留下別人的.
用來做你的奴隸.
你會發覺後面的全部是完成不了的.
但是只有猶大支派是完成了整個佔據的分地.
而接著猶大也是上帝揀選彌賽亞出現的支派.
還有一件事.
你會發覺原來這個祝福是一直臨到現在的.
當耶穌出現之後.
其實他透過這個支派所給他們整個族群.
給全世界的祝福就是從這裡而來的.
所以我想告訴大家.
你對上帝的信心其實是可以祝福你整個家族的.
其實有些人都很好奇.
我有一個兒子.

$^{481}$他很早之前.
20年前我剛讀完神學.
他跟我一起到錦繡去實習.
所以他也是錦繡幼稚園的畢業生.
後來他在外面孫德直堂成長.
後來我出去事奉.
一直在教會成長.
這個小朋友我形容在教會裡面是放養的.
問題就是每次接他回來之後.
他就會說再見.
然後他就會找自己的朋友.
整個崇拜.
因為我事奉.
由早上一直到下午五點.
我是看不到他的.
然後他很自動自覺.
四點多就會走回來辦公室等我.
我們一起走.
到他再大一點.
我們教會沒有那麼多這些.
我們那時候教會小.
沒有什麼青少年崇拜.
我們那時候十三歲就要參加成人崇拜.
所以我們那些年輕人是習慣了.
沒有那麼多古怪的事情.
然後有些人就很好奇.
他就問我兒子.
他說為什麼崇拜你不尷尬嗎?.
你爸爸在上面講道.
下面的時候他講一些事情.
他經常講家裡的事情.
你不覺得做兒子尷尬嗎?.
這個是很多傳代人的通病.
但是我兒子怎麼回答呢?.
我得意耳聽到.
我兒子說有什麼尷尬?.
我爸爸在台上只是講一些好東西.
他從來都沒有罵過我.
他從來都沒有說過我.
他也不會去將家裡.

$^{521}$當然爸爸和兒女之間有沒有一些作假.
當然有.
但是你會不會搬上去港台講?.
我不會的.
原因不是搞.
我想講是因為這個是聖經的教導.
我想告訴大家.
所以我兒子在他一路的過程裡面.
其實很坦白講.
他是看著你這個爸爸.
還是傳代人的爸爸.
你在家裡和教會是不是同一個樣子?.
這個是我們今天做基督徒.
最容易出現的問題.
為什麼信義代這麼容易流失?.
原因就是因為我們第一代那一班.
在教會裡面是一個樣子.
回到家裡是另一個樣子.
然後這一班是最容易看到什麼叫做虛假.
所以在這樣的情況下.
其實我想講.
作為傳道牧者.
其實是我們的考牌.
我想講不是裝.
問題是你在台上講得這麼鏗鏘有力.
你回家是不是做的?.
你是不是真的走的?.
所以我覺得這個生命的影響性是很重要的.
你對上帝有信心.
他自然對上帝也有信心.
今天我兒子已經在台灣生活了.
他也在台灣發展.
但其實在他臨去台灣讀書之前.
他跟我們講他16歲.
他說我想爸爸.
我想在香港受洗.
他主動提出的.
然後他也跟我們分享.
他說其中一個原因就是.
因為當時媽媽病.

$^{561}$他想在媽媽離開之前.
他想能夠見證著他能夠有這個的浸泥.
然後他到了台灣.
說真的我離他這麼遠.
我不會再見到他.
很在乎的是你跟隨上帝是怎樣.
但是正因為他有這個舉動.
他願意受洗.
去見證著上帝.
今天我去台灣.
我完全不需要擔心.
今天我們做爸爸媽媽.
最怕我經常說.
我不是要你很有成就.
不過我不想你學壞.
但反過來對於這個兒子.
我是完全不擔心的.
因為什麼呢?.
不是因為他乖.
是因為他有上帝.
這個是直接影響了你自己的家族.
你不要小看這個影響性.
所以我經常鼓勵弟兄姊妹.
在教會裡面.
你一定要做到一件事.
你去學的.
你要帶回家.
在坊間也有一個說話.
為什麼你要將你最好的脾氣.
給了陌生人.
但將你最差的脾氣給了家人.
其實我想說的整件事.
是很不合邏輯的.
所以這件事是.
上帝怎樣教我們的?.
我們的子女是.
上帝賜給我們的箭.
是我們作為神的勇士.
我們的箭.
你射出去的箭有沒有力量?.

$^{601}$你的箭是否可以繼續為上帝工作?.
你一定要考慮這個問題.
如果當你心裡有這件事的時候.
你的小朋友自然就在上帝裡面成長.
是一定的.
所以你會發覺.
其實我們今天面對我們的環境.
我們有沒有因為環境的膽怯.
而失去了上帝的應許.
給我們的祝福呢?.
大家想一想.
大家用一點時間想一想.
在這一刻.
在最近的裡面.
你有沒有一些困難.
好像剛才那些人所說.
好像偉人那麼大.
那個困難大到我都快要炸了.
有沒有這些?.
如果當你有的時候.
請你記得迦勒所說的.
我們有耶和華與我們同在.
我們一定得勝.
所以任何的困難.
都不應該成為你的困難.
為什麼我這樣說呢?.
耶穌說.
在馬太福音六章二十七節.
你們哪一個能用思慮使受歲多加一克呢?.
沒錯.
我們都控制不了你所想的.
我經常跟別人說.
你去關心別人不要說廢話.
什麼叫廢話?.
他今天很擔心.
不要想那麼多.
這些是廢話.
你以為他想這樣想嗎?.
你明白我意思嗎?.
所以你要學.

$^{641}$我經常鼓勵弟兄姊妹.
你有關心別人的心.
你去學怎樣關心別人.
有些事情是不需要說的.
很簡單.
一個人已經用了很多.
很加油,很努力.
都還是在困難當中.
接著你繼續跟他說什麼?.
加油!.
他真的想打你.
你明白我意思嗎?.
我已經很加油了.
你還叫我加油?.
不是的,弟兄姊妹.
我們要這樣去學.
弟兄姊妹不是他自己想想.
人是不會不想的.
你不想的時候.
你已經魂歸天國了.
每一個人每一天都在想事情.
但你要去學的.
幫他想什麼.
所以請大家多學聖經.
熟習多些聖經.
當弟兄姊妹有困難的時候.
請你用聖經幫助他.
當在困難當中的時候.
你告訴他.
不用擔心.
神會跟你在一起.
我也會跟你在一起.
當他很努力的時候.
他還未OK的時候.
你告訴他.
上帝知道了.
上帝知道你盡力了.
你放心吧.
上帝會看著你的.
我們也肯定了你的努力.

$^{681}$好不好?.
所以我們這些是需要去學的.
弟兄姊妹.
去學成為一個做盡的人.
所以我很盼望我們今天不要走了冤枉路.
其實祝福一早就有.
大家有沒有聽過一個說話.
叫做轉念之間.
知道什麼叫轉念之間嗎?.
就是當你今天的眼目.
明明看著這個世界裡面的困難.
但是你將你的眼目.
轉去看著上帝.
轉念之間.
你的問題就可以解決了.
因為我想說.
在那個困難當中.
這些問題是常見的.
其實我經常說.
我最近在教會教的那個叫價值觀重整之類.
我教弟兄姊妹一樣東西.
最簡單的.
大家都認識的.
你拿著一杯半杯水.
究竟你是去問上帝.
上帝怎麼搞的.
我只有半杯水.
誰誰有一杯水.
為什麼你只給我半杯水.
還是你向上帝祈禱感恩.
幸好我還有半杯水.
這個就是意念之間.
你在困難當中.
全體你自己心裡面.
你怎樣去看這些事情.
所以求上帝幫助我們.
在他當中.
為什麼迦勒和約書.
他能夠在一個這麼大的困難.
這麼高這麼難的情況之下.

$^{721}$他能夠可以超脫.
就是因為這個意念.
所以我弟兄姊妹我經常說.
我們作為基督徒.
請你有這一份在神面前的底氣.
我們什麼都沒有.
我們有什麼.
有神.
是不是.
所以聖經說.
不是靠車.
不是靠馬.
我們單靠什麼.
而是我們有神.
當你懂得這樣想的時候.
教會裡面所有的問題.
都不再是問題.
另外神的信實.
是展現在迦勒的身上.
你就會發覺.
隔了45年.
其實上帝從來對他的應許.
都沒有落空.
而你也會發覺.
其實在這45年裡面.
迦勒是默默等待的.
我不知道大家有沒有發覺.
在後來的時間.
當上帝揀選了約書亞的話.
其實說真的.
迦勒是絕對有資格去投訴上帝的.
是不是.
我和他一樣這麼有信心.
上帝為什麼你揀選他.
為什麼不揀選我.
但你發覺他沒有.
他沒有.
他完全尊重上帝他的主權.
他完全尊重上帝.
他自己擺我在什麼位置.

$^{761}$我人生在什麼位置.
他從來不會自己走去埋怨上帝.
但是有一樣東西.
他是不會放下的.
就是上帝你答應過我的應許.
我是不會忘記的.
為什麼呢.
你發覺剛才那段經文當中.
和民數記那段經文是一樣的.
唯獨我的僕人迦勒.
因他另有一個心智尊一跟從我.
我就把他領到他所去過的那地.
他的後裔也必得那地為業.
所以去到剛才讀的那段經文.
就是然而我和上去的洞弟兄.
使百姓的心消化.
但我專心跟隨耀華我的神.
當日摩西起誓說.
你腳所達之地.
定要歸你和你的子孫永遠為業.
因為你專心跟隨耀華我的神.
差不多一樣.
這句話上帝給他的應許.
45年都沒有忘記.
弟兄姊妹.
你對於上帝所給你的應許.
你相信多少.
我經常鼓勵弟兄姊妹讀經.
不是因為我是傳道人.
經常說讀經吧.
我經常說好像做銷售員.
好像賣聖經的那個.
但我想告訴弟兄姊妹.
我們今天所有人生的答案.
在聖經那裡.
你連聖經都不熟悉.
你怎會找到上帝的應許呢.
你沒有讀過上帝的應許.
你怎會想要上帝的應許呢.
你明白我的意思嗎.

$^{801}$所以你一定要掌握住.
這個上帝所給你的應許.
你要確信你不要放手.
你要絕對相信上帝.
你講出你做得到.
大家有一段最挑起大家的挑戰.
應許是什麼.
若你有這麼大的信心.
你都能夠移山.
如果你們不相信.
哪有移過山呢.
有個小朋友很小的時候.
你騙人的.
我問他為什麼騙人呢.
聖經裡面說我有信心的時候.
我就可以移山.
我昨天和媽媽去行山的時候.
我就跟上帝說.
我移走這座山,他都不動.
我不知道我們今天是不是有這種信心.
但不是的.
我經常問你一件事.
上帝真正的主權.
為什麼要由你決定呢.
為什麼不是上帝決定呢.
上帝可不可以移那座山呢.
可以,絕對可以.
但就是因為我要給小朋友證明.
告訴他們我講這個說話.
然後在他們面前不敢移.
如果上帝是這樣,就很大件事了.
有些小朋友很調皮.
然後他對著你的屋.
然後上帝說.
你幫我剷平他的屋.
然後上帝剷平他.
上帝有他的主權.
我們真正相信上帝的信心.
是因為我們相信上帝的主權.
你所做的每一件事.

$^{841}$都不會沒有意思.
都不是落空的.
這才是真正的信心.
所以今天大家.
你去面對人生當中很多的挑戰.
上帝沒有即時跟你解釋.
為什麼會發生這件事.
但如果你對上帝有信心.
這個困難不會成為你的困難.
他最多成為了你一個小小的考驗.
我形容.
然後你可以在困難當中再成長.
所以我很盼望我們能夠有.
加勒這份對上帝的堅信.
所以你會發覺到最後的時候.
加勒這個事件表達了什麼呢.
就是他這份勇氣和堅持的重要.
你會發覺他四十歲為上帝奮戰.
八十五歲他仍然一樣.
他沒有看到自己的限制.
他一直只看到上帝.
他一直堅持下去.
我們面對人生的挑戰其實都一樣.
其實我想說現在這個時代.
特別現在很多的傳媒.
很多的不斷告訴你.
經濟學者告訴你.
經濟有多差.
有些人在這段日子裡.
最多的問題是.
為什麼特朗普會當選.
明知不應該這樣.
搞到關稅.
搞到經濟差了.
我回答不了你.
因為不是上帝.
我不知道上帝為什麼讓他當選.
有些人喜歡他.
有些人不喜歡他.
沒錯.

$^{881}$因為這些緣故.
也令到很多真實的事情發生.
我們教會裡很多頂尖的會做生意.
他說送去的貨櫃裡.
有兩個已經要徵收百分之一百二十.
接下來都不知道怎麼辦.
但他說很奇妙.
船只能行到一半的時候.
幸好我用快船.
我用慢船.
他說行到一半的時候.
突然間特朗普改口.
突然間變回一百零幾.
你說他人生是不是又高又低.
但我想告訴你.
正好就是告訴你你是誰.
你可以掌管什麼.
你可以去做的就是.
信靠上帝做好自己.
我再說一點觸碰性的.
可能今天說完之後.
明天上帝就回來.
明天回來的時候大家在做什麼呢.
我想說我們每一天的時候.
我們所做的任何事情.
其實我們不需要去關注一些.
超過你能力所求的事情.
其實我們今天做人.
為什麼我們這麼不開心.
就是有很多事情根本控制不了.
你都要抓緊它不放.
人與人之間的關係.
你自己的能力.
各樣東西你都要抓住.
你都要不放.
所以最近的時間.
我看了一套戲.
我都看了幾次.
因為最近他都上了迪士尼.
名字我就不再說了.

$^{921}$但他這套戲最後的一句說話.
我覺得是很有意思.
不斷提醒我的.
他說其實我們今天能夠出世.
是一個好運.
現在是一個運氣.
而我們與其經常擔心.
究竟我們什麼時候下車.
我們不如去欣賞一下.
沿途誰和我們在車上.
和我們沿途的風景.
好好享受你的人生.
無論世界變成怎樣.
你身邊總是會有朋友的.
你身邊總是會有些令你感恩的事情.
所以最後我問大家一個問題.
今天大家歲數都不同.
老實回答自己.
不用回答我.
你從出世到現在.
困難多還是恩典多.
其實如果你懂得真實回答自己.
一定是恩典多.
為什麼我們要為了.
我們所面對的小小困難.
我們要將它放到無限大.
變成一個人.
為什麼你不看著上帝.
在這麼多年.
他完全沒有問你.
拿過任何回報的情況下.
給了你這麼多恩典.
為什麼你不相信上帝.
所以求上帝幫助我們.
讓我們和迦勒一樣.
憑著信心.
能夠得著應許.
我們每個人都能成為上帝的勇者.
我們一直為上帝爭戰.
直至到你光榮.

$^{961}$回到上帝面前.
我們求上帝幫助.
我們一齊動心祈禱.
秋天我們求你幫助我們.
讓我們都很盼望迦勒這個例子.
能夠成為我們的祝福.
因為這個僕人很謙卑.
他一路會跟隨你.
他一路很清楚.
上帝你才是上帝.
所以他不會用他自己的方法.
用他自己的想法.
去爭取自己想要的東西.
但他卻一路記得.
上帝你所給他的應許.
所以他夠膽問你.
但如果我們剛才記得的時候.
他到最後都說.
如果上帝願意的話.
我們就可以得到這個地方.
他到最後一刻仍然很謙卑.
完全看上帝你自己是如何.
我求主讓我們每個弟兄姊妹今天都一樣.
我們能夠不斷去看上帝你的心意.
我們不斷願意跟隨你.
但我們確信上帝你的恩典.
當我們有困難的時候.
我們不要過於去問.
為何上帝你這樣做.
而在過程中我們相信上帝的同在.
就如你將這一切我們交託你.
祈禱奉靠主耶穌.
得成名字其及.
(吃完).
\newpage



\section{路加福音 15:11-24}
\label{sec:iSHaOgcIOTI}
\textbf{2025.06.15 《筵席的再思》 經文:路加福音15\_11-24 講員 : 曾裔貴牧師}
\newline
\newline
連結: \href{https://youtube.com/watch?v=iSHaOgcIOTI}{\texttt{https://youtube.com/watch?v=iSHaOgcIOTI}} ~~~~ 語音日期: 2025-06-17
\newline
\newline
\hyperref[sec:HMy_7Yhf_bE]{\small{< < < PREV SERMON < < <}}
~
\hyperref[sec:index_chronic]{\small{[返順時目]}}
~
\hyperref[sec:index_scriptual]{\small{[返順卷目]}}
~
\hyperref[sec:XbnJouzOuIw]{\small{> > > NEXT SERMON > > >}}
\newline
\newline
路加福音 15:11-24
\newline
\begin{longtable}{cl}
\hline
\hline
章節 & 經文 (和合本修訂版)\\
\hline
15:11 & \begin{tabularx}{0.7\textwidth}{X} 耶穌又說：「一個人有兩個兒子。 \end{tabularx} \\ \\ \relax
15:12 & \begin{tabularx}{0.7\textwidth}{X} 小兒子對父親說：『父親，請你把我應得的家業分給我。』他父親就把財產分給他們。 \end{tabularx} \\ \\ \relax
15:13 & \begin{tabularx}{0.7\textwidth}{X} 過了不多幾天，小兒子把他一切所有的都收拾起來，往遠方去了。在那裡，他任意放蕩，浪費錢財。 \end{tabularx} \\ \\ \relax
15:14 & \begin{tabularx}{0.7\textwidth}{X} 他耗盡了一切所有的，又恰逢那地方有大饑荒，就窮困起來。 \end{tabularx} \\ \\ \relax
15:15 & \begin{tabularx}{0.7\textwidth}{X} 於是他去投靠當地的一個居民，那人打發他到田裡去放豬。 \end{tabularx} \\ \\ \relax
15:16 & \begin{tabularx}{0.7\textwidth}{X} 他恨不得拿豬所吃的豆莢充飢，也沒有人給他甚麼吃的。 \end{tabularx} \\ \\ \relax
15:17 & \begin{tabularx}{0.7\textwidth}{X} 他醒悟過來，就說：『我父親有多少雇工，糧食有餘，我倒在這裡餓死嗎？ \end{tabularx} \\ \\ \relax
15:18 & \begin{tabularx}{0.7\textwidth}{X} 我要起來，到我父親那裡去，對他說：父親！我得罪了天，又得罪了你， \end{tabularx} \\ \\ \relax
15:19 & \begin{tabularx}{0.7\textwidth}{X} 從今以後，我不配稱為你的兒子，把我當作一個雇工吧。』 \end{tabularx} \\ \\ \relax
15:20 & \begin{tabularx}{0.7\textwidth}{X} 於是他起來，往他父親那裡去。相離還遠，他父親看見，就動了慈心，跑去擁抱著他，連連親他。 \end{tabularx} \\ \\ \relax
15:21 & \begin{tabularx}{0.7\textwidth}{X} 兒子對他說：『父親！我得罪了天，又得罪了你，從今以後，我不配稱為你的兒子。』 \end{tabularx} \\ \\ \relax
15:22 & \begin{tabularx}{0.7\textwidth}{X} 父親卻吩咐僕人：『快把那上好的袍子拿出來給他穿，把戒指戴在他指頭上，把鞋穿在他腳上， \end{tabularx} \\ \\ \relax
15:23 & \begin{tabularx}{0.7\textwidth}{X} 把那肥牛犢牽來宰了，我們來吃喝慶祝； \end{tabularx} \\ \\ \relax
15:24 & \begin{tabularx}{0.7\textwidth}{X} 因為我這個兒子是死而復活，失而復得的。』他們就開始慶祝。 \end{tabularx} \\ \\
[1ex]
\hline
\hline
\end{longtable}
$^{1}$姐妹主女便平安.
父親節快樂.
相信在街上.
大家今天不會收到花的.
父親.
香港人也不會這麼浪漫.
很感恩.
我們今天特別安排小師班.
也是第三周.
福曲的獻唱.
所以是一個雙重的.
敬拜和享受.
我們父親節特別感受到.
這份窩心.
這首歌的內容.
小朋友可能會唱.
但未必深究內容.
內心的含義.
對父母的.
奎勞,他們一生的擺上.
相信父親節.
父親未必想要.
有美味的食物.
有固然好.
今晚有沒有一些的筵席.
是你已經預備好的.
你的子女為你預備好的.
相信父親.
人生裡.
應該更深感受到.
在我們的生命裡.
我們的成長.
我們對下一代的了解.
關注.
使他們能夠.
可以在成長的過程.
越來越看懂.
人生.
相信如果一個父親能夠去到.
這樣的人生階段.

$^{41}$有這樣的心境.
就真是一件很大很美好的事.
今天這一堂的道.
很想透過路加福音第十五章.
我們大家一起去想一想.
主耶穌所形容的天父.
是怎樣的心腸呢?.
今天有好幾個言籍.
在福音書.
耶穌時時透過一些.
很簡單的活動.
帶出很深很深的.
天國的那個意義.
和傲備.
我們一起.
透過這一大段聖經.
我們去思想一下.
我們的主耶穌.
祂的心意.
謝謝.
言籍的再思.
我們一起去看看.
在十四.
十五章.
總共有四個言籍.
十四章有一個.
很特別的言籍.
就是有一個叫大的言籍.
耶穌就講述.
這個言籍.
家主興高采烈.
不知為了什麼原因.
要去鋪排一個很大的言籍.
就叫那些.
已經請的賓客.
去參加言籍.
但那個比喻裡面.
很特別.
所有的賓客.
收到請帖的.

$^{81}$一概不來.
你想想.
如果你是那個家主.
你會怎樣呢.
聖經的描述就是.
那個家主是極其的老奴.
然後怎樣呢.
因為你擺滿了.
那些言籍.
那些賓客一個都不來.
那段經文告訴我們.
耶穌說那個家主.
雖然很老奴.
很憤怒.
但找了那些在街外.
貧窮的殘疾的.
低下的階層.
來出席這個言籍.
這就說.
這個言籍的賓客.
座上客.
全部都是後備.
全部都是次選.
不是首選.
怎會有這樣的言籍呢.
為什麼耶穌會說這樣的言籍的故事呢.
想帶出什麼心意呢.
帶出.
家主的言籍不受歡迎.
因為怎樣呢.
然後第十五章.
很特別.
有一點點的背景.
因為當時耶穌.
和一些.
罪人,瑞利.
來坐席吃飯.
就招致了.
有一些人.
就是法利賽人.

$^{121}$看到耶穌和這些人.
來吃飯.
就開始議論.
為什麼這些人.
明知道那些人是罪人,瑞利.
娼妓.
為什麼要和這些人吃飯呢.
就是說.
是耶穌的身份不配.
是那些人不應該.
和耶穌一起.
還是這些法利賽人.
根本看這些人不上眼呢.
這是怎樣的一個場景呢.
大家又可以想一下.
耶穌和.
那些.
當時社會上.
很邊緣化的人.
去坐席.
吃飯.
而那些人是聽耶穌講道的.
那又是怎樣的一個.
場景故事呢.
耶穌怎樣去回應.
這些議論呢.
相信第十五章.
就有三個比喻.
也都是三個很特別的言直.
第一個.
就是一個叫尋找.
失楊的一個比喻.
有一個人.
有一百隻羊.
走失了一隻.
九十九隻都在羊圈.
這個人呢.
就去到.
外面.
尋找那隻羊.

$^{161}$這是第四到第七節.
第一個比喻裡面的故事.
很特別的地方就是.
三個比喻.
都帶進一個.
的言直.
三個比喻.
和十四章之前講的那個比喻.
都是一個言直.
我們可以看看.
這三個言直.
帶給我們.
主耶穌怎樣帶出天父的心.
天父的心腸.
天父的心意呢.
這就告訴我們.
這個牧羊人因為一個罪人悔改.
牧羊人就大肆的慶祝.
他找回他自己的羊.
但耶穌呢.
就不是想講這隻羊.
耶穌在這裡的時候.
就和那些法利賽人.
和當時仍然在現場的.
瑞利和那些罪人.
倡計.
來講這個比喻的目的.
目的就是.
一個罪人悔改.
在天上的喜樂.
要給已經有.
九十九個.
無需要悔改的義人.
那個喜樂.
更加的大.
來得更加大.
耶穌是直接和.
這些法利賽人和.
在現場的那些罪人講.
就是迷失了的羊.

$^{201}$被找回.
耶穌說.
天上的喜樂.
也就是說天父的喜樂.
也就是說.
我們雖然看不到.
天上有一個言直.
言直帶來的就是.
很大的一個歡樂.
就是迷失的.
被找回.
這個就是第一個比喻的.
那個目的.
那些法利賽人.
在現場的.
那些罪人.
明不明白.
耶穌是不是若有所指呢.
耶穌講這個比喻.
解釋那個目的.
大家有沒有能夠明白呢.
如果第一個的言直.
他們不能夠感受到.
那種的歡樂.
那有第二個的.
這裡有一個例子.
是發生在我自己身上.
早兩個禮拜.
出了一部新電話.
我打算買定.
明仔的那部機很舊了.
希望如果它壞了.
我就可以馬上給他.
就去我的電訊公司.
打算出一部新機.
出一部新機.
他就按電腦.
曾先生你是我們的特選客戶.
你買這部機.
可以回贈三千元.

$^{241}$原來那部機已經減了一半以上了.
都未止.
原來新機可以去那個機的.
那個牌子的機.
去換禮物.
又送一個平板電腦.
又一千多元.
那就是說我現在的新機.
完全新的.
只需要三百多元.
但是年價是五千多元.
剛剛我們有個弟兄.
也是跟我一樣.
出了同一部機.
他就說要原本價錢.
就給.
差不多五千元.
你可以想像得到.
特選客戶.
當然這是開玩笑的.
迷失一隻羊.
天父.
透過耶穌跟我們說.
一隻羊.
迷失了.
他要尋找.
尋找到了.
在天上有大的喜樂.
如果這個都不夠清楚.
第二個比喻.
就是尋找.
那個失落的銀幣.
有多少個銀幣.
剛才是一百隻羊.
現在是十個銀幣.
這個婦人.
怎樣來做呢.
聖經告訴我們.
耶穌好像不斷增加.
讓人感受到.

$^{281}$你應該歡樂的.
如果你找到.
第二個比喻.
這個婦人.
點起燈來.
打掃整間屋.
可能要找掃地機械人.
去鑽下洞裡面.
找到就找.
現代的就用掃地機械人.
古代就沒有.
可能就要拿開桌椅.
去找.
仍然是那個目的.
找到了.
找到之後.
很特別.
居然他找到這個銀幣.
邀請他的朋友鄰居.
來一起慶祝.
一個銀幣.
和一個筵席.
你想想.
很懸殊的比喻.
落差.
慶祝找到一個銀幣.
不是一個金幣.
是一個很普通的銀錢.
你可以想像得到.
耶穌特顯的.
那個比喻.
是找到之後的開心.
滿足延籍的歡樂.
你想一想.
你怎會為了一個.
錢幣.
帶親朋好友鄰居.
來慶祝.
第二個比喻.
就特顯.

$^{321}$找到之後的快樂.
要有延籍的歡樂.
第二個比喻.
第三個呢.
剛才我們讀的經文.
就告訴我們.
有一個浪子.
他有一個行動.
什麼行動呢.
所以其實要問.
這一段經文.
11-24節.
是浪子的比喻.
還是父親的故事呢.
我們看看經文.
某人有兩個兒子.
小兒子對父親說.
爸爸.
請你現在就把我應得的產業分級.
是另一個翻譯.
更加.
特顯了.
這個小兒子.
在古代.
現代應該不會.
古代就更加大逆不道.
沒有一個人會說.
這個小兒子應該這樣做.
你想一想.
爸爸現在請你.
現在分家業給我.
是沒法想像的.
五逆.
是沒法想像的.
一種敗壞的行動.
現在請你.
現在分我.
應該得到的家業.
給我.
這個小兒子得到之後.

$^{361}$大家看到耶穌說的比喻裡面.
他怎樣來到.
立即的將他有的.
帶著離家.
接著就放任.
接著就繼續.
因為有大饑荒.
他就用了.
他所有的.
變成一個在養豬的地方.
要吃豬的豆鴿.
來衝擊.
這樣的落魄.
的人生.
然後就記述他要回去找爸爸.
這裡告訴我們.
沒有人會對這個浪子.
有同情的.
沒有人會覺得他所做的事.
是對的.
因為他已經在一個不合適的時間.
在一個不好的.
一個語調.
來跟爸爸.
這樣說這些話.
相信大家都看到耶穌說的這個比喻.
這個小兒子.
根本是罪無可恕.
在當時的社會.
是一件很難想像.
耶穌說出這樣一個.
的人.
好像很形象化.
一個罪人.
來告訴我們.
我們又讀一讀這段經文.
小兒子離家還遠.
父親忘見了他.
就充滿了愛戀.
奔向前去.

$^{401}$緊緊抱著他.
不停的親吻.
這個小兒子.
這就是父親.
一直的等待.
找回這個兒子.
不是父親.
到處尋找.
不是父親點上燈.
似乎一個人的生命.
跟一隻羊和錢.
是不同的.
一個人會抗拒.
他還沒有認真悔改.
就算父親找他.
也不會聽.
父親的召喚.
父親是等待.
去尋找他.
所以這個等待的尋找.
是很特別的一種動作.
父親每天.
就看著.
什麼時候這個小兒子.
會回來,會醒悟.
大家留意一下這段時間.
小兒子回來之後.
有什麼動作.
就是爸爸和小兒子.
和吩咐他的僕人.
準備好.
那些筵席.
原來那個筵席.
一早每天都準備好.
要用的材料.
全部準備好.
所以在這個時候.
是一個很歡樂的筵席.
就是小兒子.
悔改回來.

$^{441}$爸爸和他大排筵席.
不過在這個時候.
有一個.
從來都沒有浪子的浪子.
來到家門外.
在野外.
種田.
這個大兒子.
他也分了爸爸的產業.
因為有幸於小兒子說.
麻煩你分兩份.
一份是我的,一份是哥哥的.
哥哥拿了那份產業之後.
繼續很忠誠地.
在家幫爸爸做事.
不過.
留意一下這裡的描述.
這個沒有離家的浪子.
在不適合的時間.
什麼不適合的時間呢.
所有人.
看到小兒子回來.
無論家主.
爸爸和所有的僕人.
和小兒子.
在房間裡面.
大肆慶祝.
高興歡樂.
但是唯有一個.
在外面.
怎樣都不肯進去.
在不適合的時間.
所有人慶祝的時候.
他不慶祝.
他很生氣小兒子回來.
他很生氣爸爸.
為什麼為這個小兒子.
一個敗家的兒子.
來慶祝.
為什麼我的忠誠.

$^{481}$我的每天的生活.
耕種.
為什麼我不得到.
爸爸的喜愛呢.
他當時的內心.
就是在這樣的狀態.
而且他有一種.
很不正確的態度.
去藐視自己的父親.
他要爸爸.
出來勸他.
僕人出來勸他.
他都不聽.
他都不進去.
這個小兒子.
是死而復活.
失而又得的.
所以我們理當的.
歡喜快樂.
爸爸很語重心長地.
勸這個大兒子.
一個從來.
好像沒有任何的行動.
要做浪子的大兒子.
各位.
整個十五章.
三個比喻.
和十四章.
那個大言直的比喻.
都是講言直.
講歡樂.
都是講一件事.
就是我們明不明白.
天父的心腸呢.
十四章那裡.
似乎所有的人.
突然被呼召.
在這個大言直裡面.
是一個後備.
是一個次選.

$^{521}$但是.
如果我們這樣想的時候.
我們就會錯過了.
因為剛才十五章.
三個連續的比喻.
正正就是要回應.
十四章那個大言直裡面.
告訴我們.
原來每一個.
在天父的眼裡.
都是極其的寶貴.
失去一隻羊.
是要細意的尋找.
找著的.
天上有大歡樂的言直.
這個婦人.
她在地上這樣擺言直.
其實耶穌也是這樣說.
在天上.
得到這個小錢.
這個銀幣.
天上同樣都有大的歡樂.
再加上.
耶穌說的這個父親.
日夕盼望.
這個小兒子的歸來.
又是有一個大的言直.
來到所有人.
來到歡慶.
不過.
這裡就要開始問一個.
很重要的神學問題.
我們.
敢不感受到.
我們在.
耶穌的比喻.
教導帶出的.
那個勸勉.
在場的罪人.
在場的瑞尼.

$^{561}$在場的娼妓.
感受到自己.
好像和神的國度.
比起那些法利賽人.
真是一個後備.
真是可有可無的.
但是如果聽完這三個比喻.
耶穌說.
不是啊.
每一個在神的眼中.
都是這麼有價值.
這麼寶貴.
一個都不可以少.
一個都不可以失落.
九十九隻.
羊圈裡的安妥的.
穩妥的羊.
但是有一隻在野外.
耶穌說要細意的去找.
一個的銀錢.
耶穌說.
就算是十個.
都要找那個失落的錢.
小兒子.
更加是父親的.
兒子.
不是二分之一這麼簡單.
所以三個比喻.
很想告訴那些罪人.
因為當時.
真是法利賽人對他們.
白眼.
對他們有一個批判.
不應該.
靠近我們的耶穌.
不是他們很尊重耶穌.
剛剛相反.
是很不尊重耶穌.
他們認為耶穌.
不應該和這些人來往.

$^{601}$親愛的弟兄姊妹.
來到神的家裡.
有沒有感受到.
我們不是一個後備.
我們不是在街上.
被人邀請來.
你坐滿我們的筵席.
我們是神.
在冥冥當中.
在愛裡面揀選的.
我們每一個弟兄姊妹.
在主的眼中.
是寶貴.
父親節.
我們有沒有想到.
天父這一份心腸.
我們越來越信主內.
越來越體會.
天父的心呢.
我們見到有很多人.
不知怎樣的原因.
他最終能夠信主.
得到神救恩.
成為神的兒女.
他能夠在神的.
恩典裡.
我們有沒有為此而感覺.
一份很大的歡樂.
因為他能夠被邀請.
去到神.
神國度的大賢直.
有沒有這一份心腸.
我們能夠.
這樣看到.
另一個角度.
耶穌又給一個.
很開放式的.
一個提問.
叫這個大兒子.
你的弟弟.

$^{641}$是死而復活.
失而又得.
所以我們是理當.
大樂的.
聖經就停在那裡.
我們不知道大兒子.
最終怎樣想.
最終有沒有死死地氣.
進去賢直裡享受.
那個歡樂.
還是仍然顧我的.
在外面發脾氣.
不肯進去.
這是一個很僵的局面.
裡面大歡樂.
外面一個大的憤怒.
大兒子.
擔心爸爸被弟弟.
騙光錢.
也都敗光了.
但是他仍然覺得.
爸爸不應該接納他.
他回來可能有動機.
想再騙你的錢.
這是比喻裡面.
不過爸爸的心腸是怎樣呢.
爸爸似乎不是.
看那些物質.
是看到一個.
心靈的悔改改變.
爸爸看到的就是.
大兒子.
如果都進來.
在這個歡樂的賢直裡.
可能會有很大的改觀.
看到這個弟弟.
真是一個這樣回頭的生命.
相信.
整個十五章和十四章的比喻.
就很清楚.

$^{681}$那個受眾.
對象聽耶穌說的那段訊息的.
大兒子肯定就是.
罪人.
瑞利.
不要感覺自己是.
無份無關.
可有可無.
每一個.
無論過去多大的罪.
在神的愛裡.
都是神特選的兒女.
經過神聖靈工作.
假以時日.
生命可以很大的蛻變.
可以很大的.
為主發光.
在人生裡.
罪人可以得到很大的生命建立.
這個就是小兒子.
大兒子呢.
相信.
就是在現場.
議論耶穌的法利賽人.
他們有很多神學.
有很多他們自己個人的見解.
認為耶穌應該貼近法利賽人.
和我們討論不同的經課.
不同的舊約聖經.
不同的禮儀.
不應該和這些罪人來往.
他們是罪人,是罪有應得的.
不應該和他們來往.
不過.
耶穌不是這樣看的.
天父不是這樣看的.
原來.
從來不離開的浪子.
就是大兒子都有份.
不只是小兒子.

$^{721}$小兒子是名副其實的浪子.
因為真的走了.
真的敗光了家產.
當然他認為是屬於他自己的.
因為分了.
不過小兒子回來.
又凸顯了大兒子.
從來沒有離開爸爸.
從來不敢離開爸爸.
不過原來從此.
由始至終.
他有浪子的心境在內心.
他的內心裡.
有一個浪子.
彭黃弟兄姊妹能夠看到.
今天有多少人還沒回教會.
離開了神之後.
不再回教會.
很盼望有一天他們走投無路的時候.
他們思考.
在我的父家裡.
口糧有餘.
為什麼我要在這裡.
可能是馬場.
可能是酒吧.
可能是其他不好的地方.
為什麼我要在那裡.
飢餓.
墮落.
為什麼我不決心.
回到爸爸的父家.
爸爸是不是口糧有餘的地方呢?.
教會是不是一個.
全真道福音的地方呢?.
教會是不是一個.
透過無論是詩歌.
講道.
聚會去接納罪人的地方呢?.
這個地方.
是一個浪子.

$^{761}$他感覺到是口糧有餘的地方呢?.
另外一個問題.
大兒子.
他沒有離開父親.
不過大兒子的內心.
就是一個浪子.
大兒子看不到弟弟.
那個悔改的寶貴.
看不到一個人的生命轉變.
經過聖靈工作.
改變的新生命的美麗.
那個歡欣.
那個應該慶祝的歡樂.
我們有多久沒有在教會裡.
感覺歡樂.
還是感覺混亂.
還是感覺很多人.
很多問題.
我們仍然.
忘記了.
教會是神.
上帝歡慶的地方.
因為有很多的罪人.
經過主的救恩.
而得到改變.
朕的聚會.
分別感受到那種歡樂.
很多弟兄姊妹.
他準備洗禮講他的見證.
說全場一來坐滿了人.
很開心.
感受到神救恩的改變.
但其實每個星期.
我們來到神家.
其實就是一個大筵席.
我們應該感受到.
有神的同在的歡樂.
見到不同弟兄姊妹.
他們生命的改變.
我們在教會裡.

$^{801}$應該就是一個大筵席的歡樂.
如果我們的內心.
真的有像大兒子一樣的心.
求主幫助我們.
可能你沒有看錯.
小兒子真的是這樣的.
沒有錯的.
完全是對的.
不過能不能夠再有一個生命.
看到神的心腸呢?.
看到神愛每一個人.
神知道每一個人的時間.
以致能夠.
那個等候.
終於有收成呢?.
盼望今天這一段經文.
父親節送給大家.
今晚.
如果你有筵席.
就想想這四個筵席.
其實是一個訊息.
你當在救恩裡.
歡喜快樂.
盼望我們每一個.
真的越來越感受到.
天父在我們不同的生命裡的工作.
我們一起感恩禱告.
多謝主.
讓今天成為美好.
每一次的聚集.
我們看到弟兄姊妹的面容.
不同的事奉.
如果不是救恩.
如果不是主呼召我們.
我們根本就不會在一個地方.
一同的服侍.
敬拜我們這一位真神.
多謝你閱立我們.
無論有聲無聲.
有崗位或沒有崗位的.

$^{841}$我們來到聚會的.
每一位弟兄姊妹.
甚至乎在座有些是未信主.
但是主都是這樣愛和閱立.
幫助我們明白.
在整個大言寺的比喻裡.
我們感受到.
我們的主.
就是這樣愛每一個人.
你知道小兒子的問題.
你知道大兒子.
他的軟弱.
因為你仍然是將.
那個真理.
那個心靈裡.
一個很深度的.
那種關懷.
來讓每一個.
無論是當時的法利賽人.
和聽這個故事的.
罪人瑞利.
來陳述.
多謝主.
因為救恩從來不是去強迫.
而是讓人.
選擇.
揀選我們這一位主.
求主幫助我們.
讓我們今天在主面前.
我們有我們自己那份的.
福樂,平安,安寧.
在主的恩典裡.
我們去成長.
這樣感恩禱告.
奉耶穌基督的名.
阿們.
祝福大家.
\newpage



\section{創世記 5:21-24}
\label{sec:XbnJouzOuIw}
\textbf{2025.06.22 《與神同行人生路》 經文:創世記5\_21-24 講員 : 許淑芬牧師}
\newline
\newline
連結: \href{https://youtube.com/watch?v=XbnJouzOuIw}{\texttt{https://youtube.com/watch?v=XbnJouzOuIw}} ~~~~ 語音日期: 2025-06-24
\newline
\newline
\hyperref[sec:iSHaOgcIOTI]{\small{< < < PREV SERMON < < <}}
~
\hyperref[sec:index_chronic]{\small{[返順時目]}}
~
\hyperref[sec:index_scriptual]{\small{[返順卷目]}}
~
\hyperref[sec:jY3KiVb_LXo]{\small{> > > NEXT SERMON > > >}}
\newline
\newline
創世記 5:21-24
\newline
\begin{longtable}{cl}
\hline
\hline
章節 & 經文 (和合本修訂版)\\
\hline
5:21 & \begin{tabularx}{0.7\textwidth}{X} 以諾活到六十五歲，生了瑪土撒拉。 \end{tabularx} \\ \\ \relax
5:22 & \begin{tabularx}{0.7\textwidth}{X} 以諾生瑪土撒拉之後，與神同行三百年，並且生兒育女。 \end{tabularx} \\ \\ \relax
5:23 & \begin{tabularx}{0.7\textwidth}{X} 以諾共活了三百六十五年。 \end{tabularx} \\ \\ \relax
5:24 & \begin{tabularx}{0.7\textwidth}{X} 以諾與神同行，神把他接去，他就不在了。 \end{tabularx} \\ \\
[1ex]
\hline
\hline
\end{longtable}
$^{1}$請姐妹平安.
雨神同行是一段很奇妙的旅程.
首先要多謝曾牧師邀請我來這裡.
有機會跟大家分享訊息.
今天應該是我第三次來這裡跟大家分享.
過去兩次我都用雨神同行成為主題.
今天我想繼續用雨神同行的角度.
思想上帝的說話.
今天我想跟大家一起思想.
以洛在人生的路上如何與神同行.
其實聖經有幾位人物都被稱為與神同行.
但當然其中最具代表就是以洛.
因為以洛的生平是最短命.
但如果他生了兒子.
就是最長命的那個.
所以長一段給我們看到一個很大的對比.
也是聖經唯一記載一個與神同行.
記載了他的年份的人.
神將他接去他就不在世.
特別記載他有三百年的日子.
神直接將他接走.
有聖經學者說到.
啟示錄有兩個人.
其中一個就是以洛.
因為以洛未死.
有一天啟示錄時.
傳講福音的其中一個就是他.
我們拭目以待.
但今天我們一起思想以洛.
首先我們問.
究竟以洛發生了甚麼事.
令他選擇與神同行的生活.
聖經特別記載一段很短的經文.
特別記載為何他會與神同行.
原來就是在他生了兒子後.
他生了兒子後.
不知道是否有一個很大的頓悟.
就是因為這樣.
他就改兒子的名字.
叫做馬土撒拉.

$^{41}$這也是他生命中一個很大的轉捩點.
可能大家回顧自己的人生.
也有些轉捩點是很重要的地方.
我自己信主是一個很大的轉捩點.
夢照是一個很大的轉捩點.
同樣生了神父後.
也是我生命中一個很大的轉捩點.
可能大家也有不同的轉捩點.
我們看到以洛.
因為他生了兒子後.
他生命有一個很大的改變.
於是他就思想一個很重要的問題.
就是在人生路上他應該怎樣走.
他就決定了與神同行.
或者聖經來觀議一個與神同行的名字.
究竟以洛與神同行的生活.
會是一個怎樣的生活呢.
是不是不用我按.
我想和大家一起思想.
以洛與神同行的三方面.
我們一起思想.
他怎樣來決定.
或者為什麼會決定與神同行呢.
當他決定與神同行之後.
他究竟在人生發生什麼事呢.
三方面和大家一起思想.
第一方面就是.
當他決定與神同行之後.
其實是一個不容易的決定.
也就是說他決定.
與這個世界有一個很大的分別.
他要抵抗這個俗世洪流.
所以聖經就說.
他就決定了抵抗.
不是抵抗一天半天.
就是抵抗了三百年的日子.
到底以洛他生養的那段時間.
是一個怎樣的世代呢.
如果從聖經裡面.
我們看到至少有三件事.

$^{81}$是當時的世代發生的.
就是當時的道德敗壞.
我們知道當時.
阿當生了該隱.
然後又生了阿伯.
兩兄弟也開始了一個謀殺案.
第一宗的由家庭謀殺案開始.
就是該隱殺死阿伯.
當殺死或有人殺.
如果大家要注意.
當以洛出生的經歷裡面.
他到了三百多歲.
阿當才死.
而他當時是沒有人死過.
只有阿伯死了.
所以當時仇恨就已經有些蔓延.
當時該隱有個後裔.
就叫做拉麥.
拉麥就去公然.
他說殺該隱的.
遭報七倍.
而殺拉麥就遭報七十個七倍.
可想而知當時有一批人.
崇尚強權.
用自己的強權去踐踏.
或者欺凌其他弱勢的人.
這也是當時普遍的現況.
也是我們今天看到的現況.
第二件事至少當時是放縱情慾.
本來是設立一夫一妻制.
但當時我們看到聖經形容.
神的兒子們去娶人的女子們.
因為看見人女子們的美貌.
我們就不用學者說一些很離奇的東西.
我們就說是指敬虔的人.
和一些未信的人的結婚.
這個婚姻的狀態.
隨意的挑選.
以至在家庭生活裡面.
開始出現混擇.

$^{121}$甚至多娶妻子.
所以聖經裡面我們看到.
有很多的妻子帶來互相的紛爭.
家庭撕裂.
就是當時的社會的常態.
也都同樣聖經裡面記載.
當時在創世第六章裡面記載.
他說上帝見到當時的人是怎樣.
就是人在地上的罪惡很大.
人終日所思想的事惡.
他們的思想或人的思想.
就是充滿敗壞.
充滿令人痛苦的事.
或充滿自我中心的事.
讓自我充滿.
以至這個世界敗壞.
目中無人也都目中無神.
以下就身處這樣的世代.
一個所謂黑暗的世代.
但是以下去選擇.
去決定要活出一個與人不同的生命.
所以當以下來改馬杜撒拉.
叫馬杜撒拉並不是隨便一個名字.
馬杜撒拉的名義就叫.
拿著矛槍的人.
拿著矛槍的人是什麼意思呢.
就是上帝的審判即臨.
上帝的審判來臨.
所以事實當馬杜撒拉死了的時候.
就是洪水即將來到.
如果大家看這個表就有少少的理解.
當時人是未死的.
很多人都未死.
除了阿伯之外.
阿當都未死.
阿當都是在以洛的308歲的時候才死.
但是意思是說.
以洛在65歲的時候.
已經是人生一個很大的頓悟.
他知道上帝一定會審判.

$^{161}$上帝的審判亦都令到地上的人有一個結局.
就是一個死亡的結局.
所以由大書來到一個很正面的了講.
以洛是先知.
以洛來到宣告.
上帝有千萬的使者來到.
來到審判這個世界.
所以馬杜撒拉就是一個.
上帝即將要審判這樣的意思.
所以當馬杜撒拉死的時候.
我們都知道馬杜撒拉最長命.
最長命亦都給我們一個很大的體會.
就是上帝對人的憐憫.
去到最盡.
忍耐人去到最盡.
但是他的說話是必定會應驗.
所以當以洛.
我們說他要選擇與神同行.
他就是要選擇那個不隨波逐流的生活.
他要選擇當時的價值觀的人.
所思想的盡都是惡.
盡都是利己.
盡都是為自己的生活方式.
而以洛就要堅決去站在上帝那邊.
他要確定他怎樣與神同行.
所以今天我們去看.
很多人改小朋友的名字.
都有自己的意願在.
我們有個弟兄改了兒子的名字.
叫君林.
名字很好的.
就是相信上帝即將的來臨.
原來這些名字.
正正是讓我們活生生去見證.
你的信仰的機會.
如果我們看以洛的生活.
這個社會的混亂.
道德的敗壞.
和我們今天這個時代都很相似.
我們會不會都像以洛一樣.

$^{201}$要選擇那個對抗.
抵抗逐勢洪流的生活.
我們再說是一件不容易的事.
亦都不能夠靠自己去面對.
如果我們大家都知道.
1974年成立了ICAC之後.
香港有一個很大的轉變.
我有個弟兄就是在當時.
是基督徒加入警隊的一個見證.
在不說這個弟兄之前.
李君孝先生可能大家都知道.
第一個的華人處長.
我在很多年前.
都有30年前.
看他的一篇見證.
在一個公開的文章裡面.
他來分享.
他說當時的黑暗的時代是怎樣呢.
他打開抽屜.
就會見到有錢.
有錢在裡面.
如果你看到有錢.
人人都會自動放進自己的袋子.
你會怎樣抉擇呢.
當時他的文章.
我很深刻記得.
他就說他自己關上抽屜.
我不知道結果會是怎樣.
當人們看到他沒有收到的時候.
可能立場會有些不同.
也不一定.
亦都是一個對抗.
這個俗世洪流的一個很重要的方法.
剛才我也說.
有個弟兄是我們團契.
每個人都尊重.
他也是我們開山的始祖.
45年前.
開始成立的時候.
一個弟兄.

$^{241}$這個弟兄.
俗語所說.
就是擔善不偷食.
當然我們不太明白這個說話.
為什麼會擔善不偷食.
但是他剛剛對抗.
當時他守油麻地.
或者油尖旺的時候.
遇到一個當時黑幫的大佬.
亦都是之後戶外的.
現在我們說戶外前的負責人.
當時他們兩個相遇.
當時我們這個弟兄就信了主.
這個弟兄就未信主.
以前警察貪污.
個個都知道.
亦都很容易辦.
經過就給錢就行.
但當他見到這個弟兄的時候.
他就要拔腿走.
為什麼呢.
因為他不貪.
他就會被捕.
所以他就用這樣的行動.
見到他就要離開.
所以這個弟兄就成為當時.
收黑錢歪風的一個無聲的抗議.
這個時代的一個不同的反應.
我記得很久之前.
在未接受接受接受的時候.
剛出來工作.
我老闆是很粗口的.
但是當我和他聊天的時候.
他就從來都不說粗口.
不過有時候漏嘴的時候.
就會和我道歉.
原來當你不說粗口的時候.
原來是抵抗世俗潮流的.
一個很重要的方法.
我現在見到今天的粗口.

$^{281}$是這個時代的.
好像是甚至納入香港的文化.
我千萬不要.
祈禱千萬不要納入香港文化.
香港不要這些文化.
但是我見到無論穿著校服也好.
所有的工種也好.
都是一樣.
都會去說粗口.
我記得有一個姐妹.
她教書的.
當時她教書.
我不是抹黑老師.
但是她自己這樣說.
不關我的事.
她自己分享的時候.
當時她同時.
是說三級的笑話.
她心想.
真是很好笑.
她也很想大笑一頓.
不過她就去想.
如果我笑了.
不就和你一樣.
和你一樣.
我就要忍住不笑.
這個肢體真是.
有一個很好的見證.
是當時的.
她自己一個教學的群體裡面.
保持不笑.
所以抵抗俗世洪流.
或者我們自己團契也是.
我們同時都有粗口文化.
但是我們有很多弟兄姊妹.
你會說.
如果你不說粗口.
你身邊的人就會自然不說粗口.
如果你整輛車.
都是說粗口的時候.

$^{321}$他就不說.
同事就自然正常.
這個是很奇妙.
但是如果你都說.
你以為不說不好意思.
我都和他說了一份.
他就會和你大講特講.
所以今天我們怎樣抵抗.
這個俗世洪流.
或者我們今天.
我們怎樣教導我們下一代.
我們兒女也好.
我們現在年少的弟兄姊妹也好.
我們怎樣幫助他們.
去對抗這個俗世洪流.
以為人人都是這樣.
他都要是這樣.
還是有一個榜樣.
可以勇敢為主作見證.
所以我們都求主.
給我們有這個勇氣.
我們不要這個文化.
我們要抵擋這個文化.
所以和你旁邊的人講一下.
或者提醒他.
我們要靠著主.
勇敢抵抗這個俗世洪流.
和你旁邊的人提醒他.
提醒他怎樣去教導下一代.
你要勇敢抵抗這個俗世洪流.
免得和這個世界一樣.
這是第一方面.
第二方面我們要思想.
以諾不單止抵抗這個俗世洪流.
第二方面就是要忍耐平淡的生活.
聖經說以諾三百年與神同行.
意思不是一天.
不是一個月.
不是一年.
是三百年.

$^{361}$天天都忠心堅持.
這樣來維持與神同行的生活.
但聖經形容以諾是一個很簡單.
非常平凡的一句說話.
就是以諾生兒養女.
這個簡單的記載.
就蘊含一個很重要的功課.
就是今天你在這個生活上.
你怎樣過與神同行的生活.
我要記住.
以諾並不是在修道院生活.
以諾就是在這個思想.
或者所有人.
這個世人思想盡都是惡的世界裡面.
或者在這個敗壞的世代裡面.
去過一個普通人.
去過一個生兒養女.
建立家庭.
承擔父母責任.
這樣的生活方式.
這個給了我一個很大的提醒.
當我們今天.
如果大家鼓勵.
彼此鼓勵.
我們過一個與神同行的生活.
我們不是大家鼓勵.
我們過一個轟轟烈烈的生活.
或者我們過一個驚天動地的生活.
不是的.
有更多時候.
我們要過一個平常不過.
或者是承擔著.
我們日常應該要有的責任.
這樣的生活.
讓人要看到.
你怎樣與神同行.
不只是見到你高高在上.
這樣的生活.
以諾不因為這個世代敗壞.
而拒絕承擔責任.

$^{401}$這個是一個很重要的訊息.
我們今天來到看到.
今天很多人怕生兒養女.
我們警隊同事都一樣.
他怕生了一個.
將來要抓的人.
這麼負面的.
這個就令他們不生兒養女.
很多人這樣想.
所以寧願養隻狗仔貓仔.
養隻寵物.
都不去生兒女.
因為實在覺得生兒養女.
負責任.
或者負這個責任太大.
但是以諾他的選擇.
是敬畏上帝.
為什麼我們不更正面去說.
我們要生一個敬虔的後裔.
我們要敬畏神的生活.
而以諾他的後裔.
就是一個敬虔.
上帝揀選的後裔.
讓我們在那個世代裡面.
來見到.
我們知道家庭是上帝建立的.
一個很重要的.
基本的單位.
家庭健康.
其實社會就有出路.
但是家庭破壞.
我們見到很多的罪行.
其實背後.
就是因為家庭破裂的原因.
每一個問題的少年.
或者問題的人.
背後還有一個問題的家庭.
所以當我們今天.
我們要選擇.
在這個世代裡面.

$^{441}$我們怎樣去建立.
我們家庭的生活.
如果上帝給你有婚姻的話.
我們怎樣去建立一個神聖的家庭單位.
我們要持守一個愛的責任.
我們每一個人都要過一個.
健康的生活.
因為都是履行上帝.
給我們一個重要的社會責任.
所以當我做婚前輔導的時候.
我總是第一句.
恭喜我們大影姐妹.
恭喜你肯結婚.
這是第一件事.
如果你生了十個小孩.
恭喜你為社會.
你又生了一個社會棟樑.
我們值得恭喜.
如果你看到有大影姐妹.
又結婚生了十個小孩.
因為這個責任.
不是每一個人都願意承擔.
但是我都很鼓勵大影姐妹.
真的要生.
至少要生一個.
當我生第一個小孩的時候.
我發現我整個生命有改變.
這個改變就是更加明白.
什麼叫無條件的愛.
我對我先生是很有條件的愛.
但是對兒女.
你真的無條件.
想將最好的東西給他.
這個就是當你做了父母之後.
你會完全願意付出的.
一個很重要的元素.
所以當我們今天.
我們說以樂去忍耐平凡的生活.
我們今天要怎樣過呢.
所以我們今天要排好次序.

$^{481}$當我做門訓的時候.
我就會問一下大影姐妹.
你怎樣去排這個次序.
工作先?事奉先?還是家庭先?.
大影姐妹看一看我.
可能為了討好我.
將事奉放在首位.
或者有些人很誠實.
就是工作放在首位.
但是哪樣先?.
大家真的非常好.
沒錯就是這樣.
其實上帝的心意就是.
要將家庭放在首位.
家庭放在首位.
但實際上你生活.
你八個小時睡覺.
八個小時工作.
可能你根本沒有八個小時.
要對著你的家庭.
但其實在家庭裡面.
或者我們說擺放時間.
並不是在說時間的長短.
其實是在說生活的.
或者時間的質素是怎樣.
所以當我兩個小朋友.
小時候的時候.
很小的時候.
我有一個立志.
當我生了十個小朋友的時候.
我的立志就是.
要幫助他們自負責任.
就是我對他們的目標.
第二個當我慢慢.
工作又忙.
其實當我們在最中年.
或者最年輕的日子.
亦都是我們最繁忙的日子.
亦都是我們付出的代價最大.
那你怎樣去分時間呢?.

$^{521}$那時候我自己立志.
每天抽三十分鐘給十個小朋友.
嘩這麼少的.
好像很少.
但曾經突破了一個研究.
父親給兒女.
每天平均是七分鐘.
所以我這個很豪華.
三十分鐘.
三十分鐘是在說.
有質素的陪伴.
千萬不要和他們吵架.
三十分鐘.
說說讀書怎樣.
怎樣去收拾床.
怎樣去有規矩.
不說這些.
我們建立關係.
三十分鐘去建立.
我發現其實是很有幫助的.
當你三十分鐘去陪伴的時候.
至少你要知道.
家裡也有傭人姐姐.
這個傭人姐姐不是幫你照顧兒女.
她只是幫你照顧兒女的起居.
但實際她的生活.
她的道德標準.
她的承擔責任感.
她的生存的日後技能.
全部都是靠你.
靠你怎樣去灌輸她的價值信念.
所以當我們今天.
我們說要建立一個家庭的生活.
其實就是一個很平淡的生活.
但在這個平淡的裡面.
她是教我們.
告訴我們這個世界怎樣幫助我們.
來抵抗我們面對這個壓力的時候.
我們用家庭怎樣幫助我們.
抵抗這個俗世洪流.

$^{561}$所以以諾同樣.
她因信了與神同行.
去抵抗俗世洪流.
她都忍耐怎樣過一個平淡的生活.
所以我們都溫馨提示旁邊那個.
我們要靠著主建立溫暖的家庭.
讓我們面對這個俗世洪流才有力量.
我們溫暖一點.
微聲地向旁邊那個呼喚.
我們要靠著主.
建立溫暖的家庭生活.
以諾就是這樣.
她在平平淡淡的生活裡面.
她就建立敬虔的後裔.
這就成為我們今天所見到的榜樣.
第二方面我們要看.
以諾不單止抵抗這個俗世洪流.
她不單止忍耐平淡的生活.
第三方面我們要思想.
她就要專注這個目標.
她有一個永恆的目標.
以諾的目光超越了地上的短暫歲月.
她能夠看到.
定睛在上帝身上看到.
上帝是那位掌權者.
上帝必定會審判.
因為這樣.
所以她敬畏神.
所以她與神同行.
所以她就專注在她人生裡面.
怎樣專注神.
當她專注神的時候.
她過著與神同行的生活.
是意味著什麼呢?.
就是意味著她能夠專一的.
去過一個怎樣選擇與神同行.
以致面對這個世界的時候.
她怎樣抵抗這個世界的不同.
而她這個信心.
是一個特徵讓我們看到的.

$^{601}$這個特徵就是讓我們看到.
她用一種堅持持久的方式.
來活出她這個信心.
這個信心千萬不要說是過去式.
不要說過去很熱心.
不關事.
不要想著將來很熱心.
很多大影節目說將來會奉獻.
你現在就奉獻.
你不要等退休.
你現在全部都用現在式.
不要讓未來的空想.
以為我過了這個階段就行了.
我這些子女長大了就行了.
我父母安全了就行了.
我退休了就行了.
不是.
我們要說今天我怎樣堅持與神同行.
為什麼專注有這麼大的困難呢?.
就是因為這個世界實在有市場.
這個世界是不容易對抗.
這個世界的潮流是不容易對抗的.
我們如果活在這個罪惡的世界.
我們很正常跟著這個世界走.
就像史蒂夫·約翰·紐所說.
這個世界是充滿眼目的情慾.
肉體的情慾.
和今生的驕傲.
所以我們要面對.
其實我們相信有神.
我們也要相信有魔鬼的存在.
由古至今我們知道.
魔鬼是用盡所有的方法.
試圖讓我們將我們的注意力轉移.
離開上帝.
讓我們去看看這個花花碌碌的世界.
讓我們去看看這個世界.
是多麼能夠滿足你的需要.
所以魔鬼用盡這個方法.
我們也要用盡方法去專注神.

$^{641}$專注神是一件不容易的事.
我以前教會有一個姐妹.
她的女兒學芭蕾舞.
我很好奇問她.
學芭蕾舞不會暈眩嗎?.
不會轉到暈眩嗎?.
她就教了我們一個竅門.
她說不會暈眩.
她看著一個目標就行了.
如果她看著螢光幕.
她轉一圈看著它.
這樣就不會暈眩了.
大家不妨試一下.
我們看看是否暈眩.
她就是這樣不斷地去專注.
這個世界可能有很多吸引的地方.
但她仍然要認定目標在哪裡.
目標在哪裡很重要.
因為當我們不知道目標的時候.
就中了魔鬼的詭計.
因為魔鬼很有目標.
他的目標就是要吃掉你.
他要遍地遊行.
尋找那些沒有目標的人.
尋找那些流離朗蕩的人.
尋找那些迷失方向的人.
但當我們目標清晰的時候.
我們的方向就不會被這個世界.
也不會被魔鬼來干擾.
所以我們需要定下.
你人生目標到底是什麼.
當然我們門徒訓練.
我們做栽培也好.
做洗禮班也好.
我們去建立弟兄姊妹.
都要常常問他.
你有什麼人生目標.
最近我問一個小弟兄.
他順著沒多久.
就問他.

$^{681}$你有什麼人生目標.
他就有點不好意思地告訴我.
我問他職場有什麼目標.
他說希望能夠平平安安退休.
我說這個好目標.
這個真的是好目標.
上個月我們剛剛有個弟兄.
準備退休.
他退休之前就搞了一個小型報道會.
我第一句就要恭喜他.
恭喜你可以平安退休.
其實警察退休是一件要恭喜的事.
因為我們常常說.
警察一隻腳踏在監獄的門口.
一有衣服.
可能一時行差踏錯.
就會有問題.
但是能夠平平安安退休.
我就跟這個弟兄說.
非常好的目標.
所以當你有這個清晰的目標的時候.
你就要想一下.
我怎樣才能夠過一個平平安安退休的生活.
你今天的生活狀態是怎樣.
你怎樣跟你的老闆合作.
你怎樣跟你同事合作.
你怎樣做好你的工作.
你有你很好的專業的態度.
這個你就可以平平安安的退休了.
所以你每一個目標都是很好.
那我就問他.
你家庭有什麼目標.
他又不是很好意思地跟我說.
他說我希望女兒健健康康成長.
我說這個好目標.
好目標來的.
你能夠看著你的兒女健康成長.
是不是一個好目標.
那你就要去想.
究竟你跟你的太太怎樣保持良好的關係.

$^{721}$你跟太太怎樣合作.
教導兒女.
現在兒女只有幾歲大.
你怎樣去教導他能夠成長.
健康的成長.
你給他什麼價值觀.
你怎樣帶領他引導他.
這個是不是一個好的目標.
所以什麼目標都不要緊.
最重要是我們有專注.
我們的專注在哪裡.
我們怎樣將這個目標呈現給上帝.
每一個目標都是好的.
所以有些弟兄姊妹退休很開心.
我們的目標是什麼呢.
吃喝玩樂.
這個目標好不好.
當然好.
但你要幫助自己能夠吃喝玩樂.
不持久.
你就怎樣吃喝玩樂呢.
你吃得健不健康呢.
你玩得正不正確呢.
你會不會帶自己跌倒呢.
這個都是一個好的目標.
所以凡是所有目標.
我們都要呈現給上帝.
當我們能夠奉獻這個目標給上帝的時候.
我們就將它成為一個屬靈的目標.
成為一個上帝都喜悅的目標.
以致我們與神同行.
我們同走這條路的時候.
我們就能夠體會我們自己的生命的不同.
所以同樣鼓勵大家.
必須要去讀上帝的話語.
當我們去專注的時候.
我們有什麼成為我們的準繩.
有什麼成為我們的支持點.
有什麼成為我們比較的轉捩點呢.
就是上帝的話語.

$^{761}$上帝的話語就能夠幫助我們.
看見前面的路對還是不對.
所以大家做了這麼多年教徒.
有時候會聽弟兄姊妹講幾句話.
你大概會知道他的靈情在哪裡.
這個大家心裡有數.
不要告訴別人.
我覺得你是BB班.
不用跟他說.
但是為什麼會知道呢.
就是我們與弟兄姊妹聊天的時候.
就會大概知道.
他對上帝的話語熟不熟悉.
他對上帝的話語究竟熟悉到.
有沒有遵行到神的話語呢.
這個是我們心裡有數.
但我們求主讓我們.
不是只用來量度別人.
我們量度自己.
我們有沒有心裡有數.
我們的生命是不是不斷的要更新.
我們是不是建立到.
我們一個屬靈的自信.
我們要知道.
我所做的事是合乎上帝的心意.
很多弟兄姊妹在這裡很忐忑.
他們說我不知道對不對.
有一個姊妹告訴我.
她不知道今天吃了一頓飯.
都需要去祈禱.
我說不用吧.
因為上帝給我們自由意志.
我們不是缺乏.
經常吃肥豬肉就行了.
我們缺乏的都是合乎上帝的心意.
我為什麼不可以有自信.
去缺乏一些這麼簡單的事情呢.
或者當我們缺乏一條路.
有十條路都很好.
譬如我們的文化.

$^{801}$我們警隊的文化.
就是過幾年就會打紙酒.
打紙酒一樣很好.
有什麼不好呢.
你拍門好不好.
一樣很好.
但當我拍門的時候.
我會不會先拍上帝的門.
然後拍完之後.
你就會拍其他的門.
讓我們去看看上帝的心意是怎樣.
以洛紀一生去專注.
在這個背逆的世代.
怎樣專注同神同行.
他就成為在這個黑暗的裡面.
一盞燈.
一盞亮光.
或者一個Walking Bible.
讓人看到你.
就看到耶穌的生命是怎樣.
我有沒有這樣的信心.
你相信你所行的.
就好像活動的聖經一樣.
讓人看到神的話語.
讓人看到.
我今天應該向哪一方面走呢.
所以我們求主幫助我們.
我又再大聲提醒旁邊的.
要與神同行.
讓我們的人生充滿上帝給我們的激情.
讓我們受的激勵.
以致我們的生命.
在這個荒謬的世代裡面.
我們過一個.
可能是一個沙漠的生活.
可能是一個我們覺得很平淡的生活.
但其實不重要.
最重要是有沒有神在當中.
我們去哪裡都不重要.
最重要是你與神同行.

$^{841}$當我們願意與神同行的時候.
就像以洛一樣.
去面對的就是一個生命很重要的見證.
去抵抗.
去自己成為一個說話的一個Icon.
一個見證.
讓這個瀰漫著惡的.
瀰漫著自我中心的生活.
今天的自我中心的生活是越來越離譜的.
但是我們今天如何選擇與神同行.
同行的意思是.
我們是與神站在一邊.
我們不是與神對立.
我們不是面向神.
有時候我們都看不到上帝的角度.
我們是與神同行的時候.
我們就看到前面.
上帝看到什麼路.
我們都同樣看到什麼路.
上帝要讓我們在生活中有一個漫長的歲月.
我們的人生沒有辦法知道自己有多長.
我們都沒有辦法選擇.
我們的人生是否健康.
但是我們可以選擇.
我們是不是與神同行.
在一個平淡的日子中.
我們過一個負起敬畏神.
負起一個建立敬虔家庭.
或者負起一個我們能夠見證.
基督在我們人生中的狀態.
願上帝幫助我們.
我們千萬不要幻想.
或者我們千萬不要天真地.
以為這個世界會因你而改變.
但是我們會相信.
上帝會恩著我們.
這個光可以照亮我們旁邊的那幾個.
當他願意照亮旁邊的那幾個.
旁邊的那幾個又照亮旁邊的那幾個.
我們的生命就會有力量.

$^{881}$願神改變我們幫助我們.
我們站在上帝那裡.
我們願意立志過一個與神同行的生活.
願上帝的福音大能幫助我們.
我們都能夠成為每一個時代.
或者我們今天這個時代的見證人.
讓聖靈不斷地根深我們的生命.
也讓我們活在這個時代.
我們不斷去抵抗這個世俗洪流.
我們能夠忍耐這些所謂平淡的生活.
但是我們不斷地專注上帝給我們的目標.
以致我們能夠去到見主面的時候.
我們真的過一個無愧.
也都是滿有豐盛的生活來見主.
願上帝祝福我們每一個人一起祈禱.
你讓我們知道我們不是靠著我們自己.
就可以過一個與人不同的生活.
而是我們靠著去明白你的話語.
明白你的心意.
以致我們選擇過一個與你同行.
主要我們知道我們今天要選擇與你同行.
一點都不容易.
因為這個世界會帶著我們走.
這個世界的一些價值觀.
我們以為是對的.
我們有時都水過鴨背.
我們沒有想清楚的.
但是我們求你幫助我們.
正如以洛一樣.
我們求你開我們眼睛.
你要讓他看到他當時所看到的.
也都開我們今天的眼睛.
讓我們看到我們今天這個世代所看到的.
你放我們在這個世代.
一定有你的心意.
所以我們求你幫助我們.
使用我們每一個都成為這個世代的見證人.
讓人們看到我們.
羨慕我們過一個與神同行的生活.
讓我們有一個與神同行的美好見證.

$^{921}$願你祝福我們每一個.
多謝你.
我們願你祝福紅印堂每一個弟兄姊妹.
讓我們在這個社區裡面.
繼續建立教會.
也都會在這個教會裡面.
成為燈台.
讓更多人得著.
你的福音和盼望.
願你祝福.
多謝您聽我們祈禱.
奉耶穌基督的名.
阿們.
\newpage



\section{哥林多後書 4:16-18}
\label{sec:jY3KiVb_LXo}
\textbf{2025.06.29 《如果我們的信仰是一次變身》 經文:哥林多後書4\_16-18 講員︰陳韋安牧師}
\newline
\newline
連結: \href{https://youtube.com/watch?v=jY3KiVb_LXo}{\texttt{https://youtube.com/watch?v=jY3KiVb\_LXo}} ~~~~ 語音日期: 2025-07-08
\newline
\newline
\hyperref[sec:XbnJouzOuIw]{\small{< < < PREV SERMON < < <}}
~
\hyperref[sec:index_chronic]{\small{[返順時目]}}
~
\hyperref[sec:index_scriptual]{\small{[返順卷目]}}
~
\hyperref[sec:m1thLVxgZE8]{\small{> > > NEXT SERMON > > >}}
\newline
\newline
哥林多後書 4:16-18
\newline
\begin{longtable}{cl}
\hline
\hline
章節 & 經文 (和合本修訂版)\\
\hline
4:16 & \begin{tabularx}{0.7\textwidth}{X} 所以，我們不喪膽。雖然我們外在的人日漸朽壞，內在的人卻日日更新。 \end{tabularx} \\ \\ \relax
4:17 & \begin{tabularx}{0.7\textwidth}{X} 我們這短暫而輕微的苦楚要為我們成就極重、無比、永遠的榮耀。 \end{tabularx} \\ \\ \relax
4:18 & \begin{tabularx}{0.7\textwidth}{X} 因為我們不是顧念看得見的，而是顧念看不見的；原來看得見的是暫時的，看不見的才是永遠的。 \end{tabularx} \\ \\
[1ex]
\hline
\hline
\end{longtable}
$^{1}$鄧小姐,早安.
我們今天會討論改變的題目.
不知道大家覺得改變難不難.
有些人覺得三歲還是八十.
改變似乎是一件很困難的事情.
我們小時候看卡通片.
很期望改變.
我們小時候看卡通片.
一時變身就能改變.
如果跟我差不多年紀.
或是弟兄.
相信也曾經被一幕震撼.
在《龍珠》裡.
孫悟空變成超級殺人那一刻.
震撼了全世界的小朋友.
因為突然覺得原來人可以.
一下子變得很厲害.
你也應該試過吧.
小時候洗澡時用洗髮水.
把頭髮弄髒.
然後很憤怒地覺得自己變厲害了.
或者你是姐妹.
小時候也很渴望撿到一些東西.
撿到一些鏡盒.
撿到一些神仙棒.
就能夠變成一些很漂亮的東西.
很厲害的東西.
不過對於變身.
我有一個比較哲學的問題想問.
究竟一個人變身.
是因為他變得厲害了.
所以才能變身.
還是因為他變身就變厲害了.
有沒有想過這個問題.
一個人變身.
是因為他自己已經夠厲害了.
所以才能變身.
還是因為他變身.
所以才變厲害了.
其實卡通片是有兩種的.

$^{41}$有些人是因為不斷鍛鍊.
最後能夠變身.
有些人純粹是因為.
得到一些物件.
所以才能夠變身.
所以今天我們來到思想.
其實我們回看我們的信仰生命.
其實我們的信仰生命.
也可以說是一次變身.
特別最直接的經文是什麼.
就是哥林多後書第五章.
大哥叔經文.
有人在基督裡.
他就是新造的人.
舊事已過.
都變成新的了.
經文說任何人在基督耶穌裡面.
他都可以變身.
當然所說的變身.
那個新字是新舊的新.
同時間是包括新舊的新.
所以其實這個變身.
不單單是一個轉變.
更加是一個新的創造.
一個新創造.
上帝應許過我們.
任何人信耶穌之後.
他都有一個新的開始.
他有一個不同的生命.
更加豐盛的生命.
這個正正是一個變身.
雖然我們的頭髮顏色沒有轉變.
也沒有一套很華麗的制服.
但是我們作為基督徒.
其實本來就是一次變身.
至少我自己是這樣.
我在錦繡堂信主.
我是十七歲的時候信耶穌.
那時候我記得在錦繡堂.
星期天上教會就決志信耶穌.

$^{81}$那時候我是住元朗的.
我就坐錦繡巴在十惠那站下車.
一下車我就覺得很開心.
因為是一個風和日麗.
藍天白雲的日子.
我就跑跳回家.
有沒有試過開心到跑跳回家.
好像男精靈那樣.
就跑跳回家.
上帝給我一個很新的生命.
感覺自己是一個.
完全一個全新的生命.
上一世前一個很低落迷茫而滑.
十七歲的人.
但是信耶穌之後.
我知道耶穌愛我.
耶穌也給我有快樂.
黑暗不再成為一個光明之子.
這是一個不折不扣的變身.
我們變成一個新造的人.
所以我們以前也經常唱那首詩歌.
讚美之泉那首.
全新的你.
它這樣說.
耶穌能夠一切都更新.
耶穌能夠改變你的曾經.
耶穌愛你.
耶穌痛你.
祂能夠做一個全新的你.
成為基督徒.
你感覺是一次的變身.
舊事已過.
都變成新的.
其實這也是很多年前的事.
可能也是很多年舊事的事.
當你做了那麼多年基督徒之後.
不知道你有沒有發現.
你自己有沒有改變.
二十年前的新造的人.
還是新的.

$^{121}$舊了的新造的人.
究竟是一個什麼樣的狀態.
一個全新的你.
似乎也不是那麼新.
這個變身還是不是一個變身呢.
所以我們今天就一起去思想.
究竟一個基督徒.
怎樣來轉變.
或者有什麼方法.
其實是讓我們能夠去真的轉變.
保羅說.
若有人做基督徒,他就是新造的人.
究竟這句話是什麼意思.
今天我們看的經文是哥林多後書第四章.
正如剛才的經文.
《信不信你》的經文的之前一章.
第四章的經文.
實際上保羅在哥林多後書裡面.
是一封他自己很個人經驗的書信.
雖然當中有很多聖經的教導.
但是你不要忘記.
這本書是當時保羅和哥林多教會.
很多個人的真論的書信.
哥林多後書是一封保羅來爭辯的書信.
因為當中有很多教會裡面的敵人.
很多的攻擊.
很多的誣告誤解.
令到保羅不能夠不寫這本書信.
來分享他們自己的經驗.
老實說.
當你讀這本書的時候.
一方面是聖經真理的教導.
更加是保羅自己的信仰經驗.
所以哥林多後書第五章十一節.
「洋著基督女,她就是新族的人」.
這句不是單單一個神學理論或者得救方法.
而是保羅自己的信仰經歷.
一個新族的人.
正正是保羅自己的經驗.
什麼時候?.

$^{161}$這件事是很多年前在大馬士革路上的轉變.
但原來剛剛已經二十年.
即昔日保羅在大馬士革路上遇見主耶穌的轉變.
和今天哥林多後書的寫作.
正正是相隔二十年之久.
二十年.
保羅昔日在大馬士革路上遇見主耶穌.
耶穌還問當時的靖年.
叫做掃羅的人.
掃羅 掃羅 你為什麼逼迫我?.
然後他的眼睛就被耶穌基督的榮光光照著.
這也是在第三章裡面說的.
哥林多口信三章說.
那分副光從黑暗出來的神.
已經照在我心裡.
讓我得知上帝榮耀的光顯在耶穌的面上.
這就是保羅昔日在大馬士革的經歷.
保羅看不到東西.
三個看不到東西.
然後這個叫掃羅的人就悔改.
回傳.
舊事已過都變成新的.
掃羅就變身做保羅.
就成為一個基督徒.
二十年之後.
當這個新做的人保羅再寫一封書信.
叫哥林多口書的時候.
他就去描述他自己當下.
即是今天的狀態.
正是今天所說的經文.
所以我們可以用這段經文去看.
我經常說.
當保羅說到我們的時候.
我們可以更加去譯做我這個字.
因為是說自己的經驗.
不只是我們,是我的經驗.
所以這樣說.
所以我不喪膽.
外體雖然毀壞.
內心卻一天卻生似一天.

$^{201}$我這自輕自暫的苦楚.
要為我成就極重無比永遠的榮耀.
原來我不是顧念所見的.
乃是顧念所看不見的.
因為所見的是暫時的.
所不見的是永遠的.
這正是保羅自己的信仰經歷.
保羅二十年前.
變身從新族人成為一個新族的人.
經歷了二十年的風風雨雨.
很多不同的遭遇.
教會受到迫害.
教會裡面有很多不同的鬥爭.
保羅被教育哥林多教會誤解.
否定,攻擊,誣告.
教會艱難,世道艱難.
面對著很多不同人和事.
保羅再次分享自己這二十年.
耶穌基督如何改變他今天的樣子.
不知道當中有沒有跟保羅差不多時間.
有二十年信主還是信仰主嗎?.
大家想想也有吧?.
信仰主有沒有人是二十年以上的?.
大部分都是二十年以上.
可能在二十多年前.
當你還是很年輕的時候.
在某個中學福音營舉手.
某個報道會走出來.
決志,你便變身成為基督徒.
我自己也二十多年了.
你重新做人.
有一個新的改變生命.
你經歷了很多在教會裡面的第一次.
很多第一次都是在教會裡面發生的.
第一次去祈禱.
第一次嘗試去學習祈禱.
第一次去洗禮.
第一次洗完禮後去吃聖餐.
還記得嗎?.
那是第一次吃聖餐的感覺.

$^{241}$第一次去事奉.
第一次去參加教會斷宣隊.
第一次去翻崇拜.
第一次沒有去翻崇拜.
第一次為信仰感動.
第一次慢慢失去信仰的感動.
第一次信了主後再犯罪.
第一次開始不想回教會.
第一次,有很多的第一次.
然後這麼多年來.
在教會裡面問.
究竟我有什麼的轉變呢?.
作為一個老信徒.
一個有經驗豐富的基督徒.
你是一個怎樣的狀態呢?.
其實我在教會裡面.
和一個初信者做栽培.
很多年沒有做栽培.
在哨院教書較多.
都是教哨學生.
很少和初信者做栽培.
我基本上都是將哨學的事告訴他.
所有哨院的事都一次過精華告訴他.
不過初信者有一個很特別的情況.
他很容易醒問我.
信仰主之後要做什麼.
他就急不及待地問我要做什麼.
我跟他說你有三個身份.
你要知道你是天父上帝的兒女.
你是耶穌基督的門徒.
你是一個見證耶穌的基督徒.
所以你是天父的兒女.
你要和神有好的關係.
你是門徒.
你要學習主耶穌基督的榜樣.
你是基督徒去見證神.
你要去將耶穌見證給人聽.
他就很熱誠地抄下來.
這是一個很特別的經驗.
我發覺初信者是一個很特別的動物.

$^{281}$他是很有信仰的熱誠.
馬上就想知道有什麼可以做.
反而我們信仰主很久的人.
就不是這樣想的.
很多信主都會問.
教會有什麼福利.
教會有什麼活動讓我參加.
教會有什麼活動讓我投票.
其實我發覺很特別.
其實做基督徒就是要做這些.
做基督徒就是要做這些很簡單的事.
所以我發覺原來在我們信心病裡.
最爆炸的階段.
可能就是初信主的階段.
所以我經常和自己說.
快點把握這個機會和他做栽培.
有什麼就快點叫他做.
慢慢他就開始不做了.
所以初信主是一個很特別的時間.
實際上在教會中牧養這麼多年.
其實牧養是一件很不容易的事情.
因為牧養就是希望弟子姐妹能夠.
生命有改變.
但是牧養叫人生命有改變.
是談何容易的事情.
我自己港島數數下手指都說了20年.
我24歲港島到現在都有20年.
那一年裡面都或多或少有一兩篇.
是有些挺好的港島會.
但是我經常問.
這些港島有沒有真的能夠幫助弟兄姊妹.
或者幫到多久呢.
看完篇道之後.
那天可能有些火熱有些感動.
可能一個星期一個月.
之後又會怎樣呢.
究竟什麼東西是叫我們真真正正.
生命可以得著改變呢.
是一件很不容易的事情.
所以我經常覺得牧養就好像種花一樣.

$^{321}$開花是很漂亮的.
但是開花是一件不容易的事情.
你種了花你就知道.
一年裡面一個植物開花.
是一個偶爾發生的事情.
開完花都會凋謝.
所以牧養是一件很不容易的事情.
所以我們一起去聽一下.
保羅這段話.
保羅說.
所以我不喪膽.
外體雖然毀壞.
內心卻一天生似一天.
保羅說內心是一天生似一天.
這裡其實原文裡面不是用一個新這個字.
原文不是用一個新這個字.
就是canol這個字.
而是一個很特別的字叫renew.
就是anarchanormal.
renew這個字是一個很特別的字.
審訊中出現了兩次.
一次是哥林多後書裡.
一次是歌羅西書經文裡.
這個就是我不喪膽.
外體雖然毀壞.
內心卻一天生似一天.
就是用了renew這個字.
更新這個字.
另一段就是歌羅西書經文.
穿上了新人.
新人和知識就漸漸更新.
這個就是renew這個字.
renew這個字是一個很少的字眼.
因為很少機會用得著.
為什麼?.
因為renew這個字正正是只能夠針對new.
renew是專為new而設的.
很特別的定律.
有三個不同的定律.
世界上任何東西都會怎樣?.

$^{361}$會從new變old?.
明明是一個新的東西.
就變成舊的.
你買一雙球鞋回來.
穿著就會變成舊的球鞋.
你一棟新樓.
慢慢就變成一棟舊樓.
你結婚.
就是新婚夫婦.
慢慢就變成老夫老妻.
老婆老公.
所以這個是世界定律.
從new變成old.
但是上帝可以相反.
任何old的東西.
一些舊的東西.
都可以變成新的.
這就是保羅海頓經文.
都變成新的了.
任何舊的東西都可以變成新的.
但原來還有第三個定律.
任何新的東西.
在上帝裡面.
都要變成renew.
需要更新.
renew是專為新而設的.
任何一個新做的人.
都需要更加變成更新.
更新是必須的.
renew是必須的.
因為這個專為每一個新做的人.
而設的改變.
經歷了20多年.
風風雨雨的保羅.
他告訴我們.
我們的身體雖然毀壞.
內心卻一天生似一天.
所說的外體.
不是單單所說的身體.
而是整個世界的環境.

$^{401}$整個人.
所以不是說你現在老了.
頭髮少了.
肚子大了.
風濕多了.
而是你整個人的外面的改變.
所以保羅不斷去強調.
這個對比經文裡面.
外體內心.
所見的.
所不見的.
暫時的.
永遠的.
保羅說外體雖然毀壞.
但他裡面的人.
他仍然能夠是一個新做的人.
所以並不存在.
所謂舊了的新做的人.
一個新做的人.
他需要不斷在上帝裡面.
不斷的去renew.
今天講題叫.
如果我們的信仰是一次變身.
其實我們的信仰從來都不是一次的變身.
我們的信仰是一次又一次.
每天慢慢不斷的.
少少的改變.
我們經常都想.
我們的信仰就是.
在二十年前那一次就改變了.
某一次我回某一個培靈會.
我很感動就改變了.
某一次我信仰裡面.
很特別的經驗就改變了.
我們經常都想.
以前看卡通片一樣.
一下子就變得很厲害.
但保羅告訴我們.
這些信仰裡面.
是一個慢慢每天少少的改變.

$^{441}$不斷的持續的改變.
我們經常都覺得改變是什麼呢?.
就好像你看電視劇一樣.
通常電視劇去到最後那一集.
最後那十分鐘.
有些問題解決不了.
主教會怎樣做呢?.
通常會坐飛機離開.
坐飛機離開.
寫著五年後.
五年之後會怎樣呢?.
突然發現主角在巴黎的咖啡店.
寄了一張明信片回來.
你覺得五年之後.
人們就會突然的豁然開朗.
人們就會看開了.
放下了感情.
什麼都解決了.
我們很期望我們的生命可以這樣.
突然間可以這樣.
什麼都不一樣了.
什麼都改變了.
事實上我們作為一個新做的人.
正正是一個不是一次過突如其來的轉變.
而是仍然去堅持.
堅持我們生命裡面的更新.
所以保羅一開始就開宗明義.
來勸勉弟兄姊妹.
他說不要喪膽.
其實和本的翻譯不是最準確.
不要喪膽的意思其實和膽子沒有什麼關係.
不是大膽或小膽的問題.
而是在乎於你的心的問題.
不要喪膽的意思是.
英文可以說是Don't lose your heart.
不要失去刁鋸你的心.
不要放棄不要氣餒.
保羅勸勉弟兄姊妹跟他自己說.
不要失去你的心.
不要丟棄願意改變的心.

$^{481}$好好保守著它.
縱然這個世界不斷告訴你.
事情是改變不了的.
縱然你的經驗告訴你.
你很想轉變都好像無法轉變.
但是仍然要保持願意改變的那份心.
縱然覺得好像很難.
但仍然知道我們基督耶穌裡面仍然可以轉變.
我們的生命仍然是一個不斷改變的途徑.
上帝仍然不斷每天更新我們.
改變裡面最大的問題.
其實不是改變不了.
而是我連願意改變的心都沒有了.
回到離教這麼多年.
已經沒有什麼期望.
只不過很多責任.
當你沒有了這份信仰的期望.
當你沒有了這份初信者的熱切期待的時候.
其實這正正就是我們走錯了方向.
實際上聖經裡面「變」這個字.
是一個很普通到不得了的字眼.
希臘文叫做「基督埋」.
「道成肉身」的字眼.
就是改變的意思.
所以是一個很普通的字眼.
聖經裡面出現了二千多次這個字.
基督埋,成為,改變.
上帝其實可以叫任何東西變成任何東西.
他可以將水變成酒.
可以叫摩西的手杖變成蛇.
可以叫風浪平靜.
可以叫石頭變成水,出水.
但是聖經裡面.
上帝要轉變一個人.
一個普普通通的人.
是一件最高難度的事情.
不是上帝做不到.
而是當一個人的轉變.
上帝期望他自己都一齊來參與.
上帝不會將你變成一個沒有感情.

$^{521}$只是做好事的好人.
讓你一齊來參與.
所以願意去改變.
回到那麼多年.
到今天仍然願意渴望改變.
正正是我們很期待的事情.
願意在今天早上.
我們今天離開教會的時候.
擁有這顆初信者的心.
大家可以跑跳回家.
傳授這份信仰的快樂.
這份新的生命.
我們一起去面對我們的人生.
我們祈禱.
天父多謝你讓我們在基督裡面成為新造的人.
你讓我們的生命.
可以在你裡面不斷地更新.
求主你叫我們不要放棄.
不要刁談,不要失望.
我們仍然對這份信仰.
一份我們很熟悉.
一份我們仍然有很多期待的信仰.
我們仍然有著這份盼望.
讓我們仍然去思考我們怎樣去更新自己.
對於一些改變不了的事情.
我們仍然樂於去看見這個改變.
求主你的聖靈去更新我們.
讓你的聖靈在當中.
叫我們仍然在你裡面成為一個新造的人.
讓你來光照.
奉主名求.
阿們.
\newpage



\section{猶大書 1:20-21}
\label{sec:m1thLVxgZE8}
\textbf{2025.07.06 《保守自己常在神的愛中》 經文:猶大書1\_20-21 楊明麗姑娘}
\newline
\newline
連結: \href{https://youtube.com/watch?v=m1thLVxgZE8}{\texttt{https://youtube.com/watch?v=m1thLVxgZE8}} ~~~~ 語音日期: 2025-07-08
\newline
\newline
\hyperref[sec:jY3KiVb_LXo]{\small{< < < PREV SERMON < < <}}
~
\hyperref[sec:index_chronic]{\small{[返順時目]}}
~
\hyperref[sec:index_scriptual]{\small{[返順卷目]}}
~
\hyperref[sec:V1IfIE0yrWg]{\small{> > > NEXT SERMON > > >}}
\newline
\newline
猶大書 1:20-21
\newline
\begin{longtable}{cl}
\hline
\hline
章節 & 經文 (和合本修訂版)\\
\hline
1:20 & \begin{tabularx}{0.7\textwidth}{X} 親愛的，至於你們，要在至聖的真道上造就自己，藉著聖靈禱告， \end{tabularx} \\ \\ \relax
1:21 & \begin{tabularx}{0.7\textwidth}{X} 保守自己常在神的愛中，仰望我們主耶穌基督的憐憫，進入永生。 \end{tabularx} \\ \\
[1ex]
\hline
\hline
\end{longtable}
$^{1}$我們看兩節經文.
或許有些人很熟悉.
以前背過的金句也不一定.
但我們也嘗試一起看看.
一起認識一下.
我應該向著.
啊 Sorry.
嗯?.
好.
不好意思.
我們這次看的這本書.
叫做《猶大書》.
《猶大書》是誰寫的呢?.
是猶大,大家認不認識猶大?.
這個不是大家很熟悉.
那個出賣耶穌的猶大.
不是的,他已經去世了.
這是誰的猶大呢?.
這是耶穌基督的弟弟.
耶穌基督的弟弟不是反對耶穌傳福音嗎?.
他不是不相信嗎?.
原來耶穌基督.
死後復活.
有一段時間跟他們一起.
耶穌基督復活後,過了40天才升天.
那段時間.
他寄了弟弟,他應該有四個弟弟.
那四個弟弟和妹妹.
都認識了這個真系.
哥哥真的是救世主.
他們真的相信了主.
不過他們很特別.
他們不敢稱呼他為哥哥.
他們稱呼他為耶穌基督.
我是他僕人.
猶大這樣介紹自己.
我是耶穌基督的僕人.
他是一種謙卑.
來稱呼耶穌.
他說他是雅各的兄弟.

$^{41}$雅各應該是大哥哥.
就是耶穌基督之後的弟弟.
寫雅各書的雅各.
這位猶大.
也是耶穌基督的弟弟.
也就是說.
他對耶穌基督的生平很熟悉.
甚至他未出來傳道的時候.
他的生平也很熟悉.
知道他的敬虔.
知道他怎樣成長.
仰望神.
敬拜神.
這是猶大.
猶大寫信當中.
提了一群人.
提了一群.
這個世界有些人是不敬拜神.
有些人是不敬拜神.
這個世界有些人是不敬虔.
你不認識一些人.
是不敬虔的.
他們不懂得敬拜.
真正的神.
或者他們說信主.
但不懂得.
時時都和神傾談.
問問.
說信主,不過和神沒有什麼聯繫.
另外有些人.
喜歡分黨.
喝茶.
隔壁桌.
有些人喜歡.
有些人就說.
是不是?.
有什麼事.
一會兒外面再談.
不是的,信主的時候.
你就會改變.

$^{81}$但有一群人.
猶大形容他們.
沒有聖靈.
心裡沒有聖靈和他們溝通.
心裡沒有想.
神是想他們怎樣.
原來有這樣的人.
是不是呢?.
猶大寫給他們呢?.
我們看看猶大寫給.
哪些人.
猶大寫給這群人.
他們是在父上帝裡.
蒙愛的.
是神所愛的.
我們唱的詩歌.
是給父神愛的.
神時時賜恩典給我們.
這群收信的人.
都是在神面前被愛的.
保守的.
保守這個字.
照顧著,看著他們.
耶穌基督原來保守這群人.
天父有愛他們.
這群人還有一個特色.
他們熟悉聖經.
當時還沒有新約聖經.
新約的書信.
還在開始寫.
我們現在有的新約聖經.
他們還沒有.
所以他們讀舊約的聖經.
但他們對舊約聖經很熟悉.
為什麼這樣說呢.
因為猶大寫這封信的時候.
又說以洛.
又說摩西.
又說索多瑪,又說俄摩拉.
好像他們每個人都知道這些事.

$^{121}$這群人是熟悉聖經的.
還有一些當代的文獻.
裡面有些說的事.
就說天使去抓摩西的遺體.
這些我們不熟悉的.
但他們是懂的.
這些是什麼呢?.
是一些當代的文獻.
當代還有一些.
我們稱之為.
龐經,次經.
現在要找,網上都找到.
不一定要買,網上都找到.
但我們沒有列入聖經.
不是我們教會不列入聖經.
是當時.
編輯的時候,列出哪些為正典的時候.
很多當代的神學家.
教父.
他們覺得.
那些不是正典.
可以做參考.
大家可以讀.
不過正典的聖經.
就是舊約三十九卷,新約二十七卷.
所以我們就.
經常讀聖經,就少說.
叫大家讀一些當代的文獻.
但這群人.
他們當時都熟讀文獻.
等於姐妹,你有沒有熟讀一些呢?.
有沒有讀一些屬靈書呢?.
我相信這裡有很多書.
有些人除了看聖經.
他也看一些參考書.
看一些屬靈書籍.
這些人他們看的.
還有.
猶太人稱呼他們.
親愛的.

$^{161}$猶太人是親愛的人.
我聽姐妹.
哪些是你親愛的人.
和你親愛的人.
你給什麼他.
最重要的一定跟他說.
說多幾次.
外面沒有人覺得你暗沉.
但你親愛的人就覺得你暗沉.
因為是親愛的.
怕他得少了利益.
這群猶太的.
修信人.
猶太書的修信人.
就是猶太怕他們.
吃虧了.
所以一定將最好的東西.
一而再地跟他們說.
他跟他們說了幾件事.
雖然說了很多.
現在又有假教師.
有這個有那個.
最重要.
你要在至聖的真道上造就自己.
藉著聖靈祈禱.
藉著聖靈禱告.
保守自己.
上載上帝的愛宗.
仰望我們主耶穌基督的憐憫.
進入永生.
他說這幾樣是很重要.
我們又看看這幾樣.
很重要的是什麼.
第一樣就是.
要在至聖的真道上造就自己.
你多不多.
看一看聖經讀一讀聖經.
現在的聖經很好.
以前的聖經只有.
直行本.

$^{201}$後來就有雙行本.
但有的都是細字本.
是不是.
我覺得.
這個屏幕對我很好.
如果屏幕放在後面.
我就要看看.
看不看得到.
大家看聖經的時候.
現在有大字版的聖經.
還有什麼版.
是方便你看的.
是的,後面有機器.
讓我看看,有電子版.
電子版.
你想多大字都可以.
想多小字都可以.
電子版輕很多.
以前我們上教會.
我想曾牧師也是.
拿著一本大大又厚的.
又捨不得不帶那本串珠.
串珠就更大本更厚.
是不是.
但是一定帶著聖經.
回去崇拜.
回去崇拜帶著聖經.
有什麼牧師說完.
就可以在旁邊寫下來.
很重要的.
有沒有帶聖經.
有沒有看聖經.
原來第一樣.
猶大和信徒說.
第一樣不是要防住假教師.
第一樣是要搞定自己.
讀好本聖經.
原來要搞定自己.
讀好本聖經.
你會說我怎麼讀都不明白.

$^{241}$你現在有電子版.
電子版一開著.
它就讀給你聽.
你讀不明白.
我讀給你聽.
你喜歡男聲還是女聲.
你喜歡快還是慢.
讀給你聽.
真的讀給你聽.
那我們怎樣跟神說.
天父我看不到聖經.
今天不讀.
不可以吧.
今天只是不舒服.
你很累也好.
我們一起看聖經.
讀聖經.
我相信曾牧師有計劃給你們.
怎樣讀聖經.
還有一樣.
背一下.
背一點.
我聽過一些牧師.
他們是背聖經.
不是一節.
背聖經的意思不是一節.
一張.
有些地方是背一本.
我們以往.
北上傳福音.
現在人家是.
南下.
讀神學.
時代的轉變.
當時北上傳福音是怎樣呢.
他們是沒有聖經.
於是他們是上本背.
還有很多地方是沒有聖經的.
現在中國很進步.
是印聖經的.

$^{281}$我去參觀過一間.
在南京的.
印製聖經的.
印刷廠.
人家說幾畝地的.
那間廠.
要開車去那幾畝地的廠.
去不同的房間參觀.
參觀他們的紙張.
他們怎樣印刷.
出來之後怎樣釘裝.
然後是.
金邊皮面.
或者紅邊皮面.
或者不同的青年版的皮面.
很漂亮的.
運送出去.
外國給人家可以讀.
甚至運送出去.
一些他們不能自己印聖經的地方.
我們這麼方便.
在處也有.
去哪裡也有.
我們要熟讀我們的聖經.
聖經不是要放在這裡打開它.
聖經是要打開它.
我們就讀.
讓我們能夠熟讀我們的聖經.
有金句的時候我們記住它.
記得記住地址.
記不到地址不能探訪.
所以它是哪一章哪一節.
記得.
有沒有注意學.
我其實多遺憾.
中學是基督教學校.
我讀書的時候.
有團契的.
學校本身是教會學校.
教會就在這裡.

$^{321}$但是有一樣東西缺乏.
就是缺乏主日學.
到現在我都覺得是一種遺憾.
因為能夠從最小的時候.
認識聖經.
認識它正確的解經.
其實是很好的.
如果有主日學.
你參加吧.
你不用擔心我聽不聽得明白.
牧師不會問你.
今天幾分.
聽到幾成.
不會的.
你聽到十成當然很好.
聽到九成都賺了.
記得.
主日學查經班.
是為幫我們.
不是為講完而設.
是為幫我們.
讓我們能夠在至聖真道上.
做受到自己.
當我們學了之後.
還有一樣東西要做.
就是用.
你用了就會記得.
就好像.
我剛才享受了很好的咖啡.
但我不懂得製造.
但製造的.
他是懂得製造.
他用著懂得製造的方法.
他一直繼續製造.
就會越來越有進步.
我們讀了聖經.
就算背了它.
都好像唐詩一樣.
很熟悉.
但月光.

$^{361}$宋詞的月光.
都不是和我們有很大關係.
關係是我們的聖經.
我們用了它.
是很親切的.
我要向高山舉目.
我的幫助從何而來呢.
你知道祈禱.
就知道你的幫助.
從何而來.
求主幫助我們.
我們接下來要做的.
就是要祈禱.
原來我們信了主之後.
有一個權利.
由你決志信主那一刻開始.
你就有一個權利可以禱告.
那怎樣祈禱呢.
信主那一刻就有人帶著你祈禱.
你回來教會有人帶你祈禱.
那我在家裡怎麼辦呢.
都有很多書講怎樣祈禱.
不過最重要.
靜心下來.
將你要的事.
你要講的事.
跟天父講.
就算你很疲累.
你告訴天父我今天很疲累.
告訴他.
你覺得有些事很難.
我不知道怎樣做.
天父這件事很難.
我不知道怎樣做.
最好多加一句.
你幫我.
如果他是天父.
他時時都愛我們.
我們告訴他.
我家的WIFI很慢.

$^{401}$轉不到公司.
原因是每間公司都說.
你們那間.
屋苑我們不做.
於是我只能一個.
每間公司做.
價錢很便宜.
但不是便宜貴.
我想用的時候有時會輕.
上去zoom的時候.
我不知道去那裡zoom.
我們.
祈禱.
跟神說話的時候.
何時傳遞到上去.
你試一下.
我試過一些時候.
天父說.
無情百事下大雨.
我沒有戴傘.
感受總站一下來.
就沒有戴傘.
是沒有戴傘的.
每條街都沒有戴傘.
怎麼辦?.
天父.
現在都要下車.
來不及停雨.
誰知他就是停雨.
原來這樣都叫祈禱.
這樣神都收貨.
馬上停雨.
回到教會.
進到教會.
外面下大雨.
是不是?.
有時我常覺得神幫我撐傘.
是不是?.
祈禱.
神的Wi-Fi很強.

$^{441}$誰的Wi-Fi都能收到.
神的容量.
電腦的容量.
很強.
可以收到所有人的禱告.
不會延遲回覆.
遲答覆你那些.
他做得很好.
正在做預備功夫.
冷水咖啡好像要先泡一晚.
沒錯.
神.
他聽我們禱告.
不過.
不要只自己祈禱.
自己祈禱.
我們看到的只是自己的事.
未必看得很闊.
就算有.
我們拿著程序表.
多祈幾樣.
但回來,弟兄姊妹的分享.
你可以彼此代禱.
還有做多了一件事.
關心.
我們都很關心.
回來的時候.
牧師告訴你.
我們為了哪個國家禱告.
你可以祈禱得更闊.
在你的心靈裡.
做了宣教的工作.
多划算.
不用去到那麼遙遠.
你的心靈開始宣教.
求主幫助我們.
原來.
祈禱不是我想怎樣說.
就怎樣說.
祈禱我們可以.

$^{481}$問天父你教我禱告.
今天我祈禱甚麼.
跟神溝通.
不只是我們說.
我們也聽聖經.
祂引導我們.
在聖靈裡我們去禱告.
讓聖靈帶著我們去祈禱.
我中學.
教會學校.
中學會考之後要做甚麼.
就要來DSE放榜.
要做加分還是減分.
我們需要找找.
去哪間學校讀甚麼科.
這個比較重要.
去哪間學校.
選甚麼科.
這件事.
對一個十多歲的年青人來說.
是很陌生的.
亦很複雜的.
我怎樣做呢.
我將它放在祈禱裡.
放在祈禱裡.
除了會考之外.
還要考考入學試.
找找學校.
考考入學試.
我給了三份入學試.
其實中學.
入學試.
要家長簽名.
是嗎.
原來入學試要家長簽名.
我三份的.
投考表格.
媽媽只肯簽兩份.
第三份也不肯簽.
她說.

$^{521}$那間學校.
媽媽太太告訴我.
不要讓女兒去讀.
媽媽不肯簽.
我沒得去.
但另外兩間.
投考表格.
我自己決定.
交了.
其實兩間我簽不同的表格.
我打算.
我想讀這幾科.
隨便一科.
因為中學生很懂.
看看哪間收到我哪一科.
起碼有一間收我一科.
最後.
我發覺媽媽不讓我讀那間.
是挺好的.
因為山長水遠 坐車也很遠.
另外兩間.
一間中途距離.
一間近距離.
他們收我的竟然是同一科.
我去到哪裡都是讀那一科.
很確定.
神祈求我讀這一科.
後來我蒙召做傳道人.
我發覺.
那一科對我很重要.
為什麼呢.
那一科.
很多時候我要和陌生人見面.
做面試.
寫報告.
其實我之前在中學.
我是不說話的.
我完全不說話.
我可以一個星期除了上學.
我閉上嘴巴.

$^{561}$沒必要不說話.
原來現在可以這麼長氣.
其實那幾年大學.
對我很好幫助.
神知道我需要什麼.
所以祂就給我什麼.
神都知道你什麼.
知道你是怎樣的人.
知道你需要什麼.
還有知道你前面的路要怎麼走.
祂為你預備的.
多過你看得到.
祂為你準備的.
多過你看得到.
祂為你預備的.
讓你可以預早在起跑線.
祂帶你走前面的路向.
這是神幫我們的.
我們一定要在禱告裡面.
做一件事.
去聽神說什麼.
按著神的心意去祈禱.
我鼓勵你做一件事.
寫下祈禱的結果.
如果不是.
神回答了你你都不記得.
沒有回答你.
或未回答你的事.
你都要記住.
這樣很不公平.
不是對神不公平.
是對自己不公平.
忘記了神幫了我.
忘記了下次可以找神.
這句話我覺得很特別.
保守自己常在神的愛中.
不是說神已經愛我們嗎.
不是說耶穌保守我們嗎.
為什麼我們還要.
保守自己常在神的愛中呢.

$^{601}$各位爸爸媽媽.
你都知道.
各位爸爸媽媽.
你都很愛你的兒女.
但有時候總覺得.
他們乖的時候很可愛.
但有時候總覺得.
他們很激動.
不知道怎麼回事.
是不是.
神愛我們.
我們都要看著自己.
讓神愛得樂得.
是不是.
求主幫助我們.
第一件事.
他們.
這班人.
修信的人.
是夢神所愛的.
神愛他們.
你是不是.
你是不是夢神所愛的.
是吧.
我們大力一點.
不會掉下來的.
我們是夢神所愛的.
我們不是那種人.
我們不是好嘲弄的.
不是不敬虔的.
不是只有放縱私慾的人.
是不是.
我們不做這種人.
我們要看著自己.
神看著我們.
我們也要看著自己.
看著自己什麼.
不做就不做.
不做好嘲弄的.
不做不敬虔的.

$^{641}$看著自己走一條正路.
正的那條路.
就算是一條小路.
那條是正路就要過去.
還有.
看著自己健康.
健健康康.
你就能回來.
是不是.
健健康康你又能幫助別人.
看著自己.
健健康康可以事奉得更久.
更長久.
都是一個很好的見證.
是.
彼此鼓勵.
在主日學也好.
在祈禱會也好.
回來崇拜也好.
在咖啡茶座也好.
我們可以彼此互相鼓勵.
不是.
每次都一定是你鼓勵我.
或者是.
每次一定是我鼓勵你.
是我們彼此都會有時候有需要的.
讓別人.
別人又成為我們的幫助.
我們又可以成為別人的幫助.
這個.
不是代表我們很卑微.
能夠被幫助.
都是蒙恩.
能夠被幫助都是一種自信.
只有自信的人.
才能接受別人的幫助.
求主紀念我們.
讓我們能夠做好一個準備.
可以進入永生.
沒錯.

$^{681}$得救了.
但有些事要時時做好準備.
記不記得有十個童女.
那十個童女.
正在等新郎來.
怎知只有五個人.
可以迎接新郎.
為甚麼?有五個沒有做準備.
主何時要來?.
我怎知道?.
你知不知道?.
你都不知道,你怎知道?.
我們沒有人知道.
我們不知道,所以我們要做一件事.
時時都做好準備.
做好準備.
神何時要來.
何時召喚我們.
我們都做好準備.
時時存在盼望.
喜樂盼望地面對.
應該是復活節後.
我們收到這個消息.
我們已經是第一對夫婦.
收到這個消息.
我先收到,然後我才殺手.
我經常這樣說.
這是一個.
不是那麼好的消息.
我先生的其中一個哥哥.
對我來說是大伯.
但可能因為我是傳道人.
我們認識的時候我已經是傳道人.
所以幾位大伯和嫂子.
他們都有一個情況.
有甚麼事,是先找我.
找我,我會為他們祈禱.
還有他們會後面加一句.
可以讓我先生知道.
明白嗎?.

$^{721}$雖然大家是一個家族.
但他們一定會加這一句.
因為他們跟我說的時候.
沒有說我先生可以知道.
所以我先知道.
後來我先生知道.
是甚麼呢?就是我們七哥.
我先生排第八.
我們七哥應該是.
復活節的時候.
開始不舒服.
我先生說得很輕鬆.
我是自己走進急症室.
意思是.
我不是叫白車進去.
他告訴我們的時候.
他已經在醫院住了幾天.
我又要等復活節完結之後.
復活節的日日.
我才去探望他.
去探望他的時候.
我已經覺得打了問號.
為甚麼你每件事都說得這麼輕鬆.
你以為公立醫院.
床位多得太多嗎?.
讓你住幾天,住一個星期.
他的樣子.
真的沒有病.
他的樣子真的不覺得他有病.
但我真的要相信.
公立醫院留他在這裡.
做很多不同的診斷.
人家排練.
排兩年都沒做到的MRI.
免費給他做.
當然有些.
做了都不知道是甚麼病.
他.
不是很輕鬆.
他一直做準備去面對.

$^{761}$他講的情況.
醫生沒有講.
你知道.
好像曾律師我們這樣.
在教會內大概知道.
這樣這樣就應該是.
那種病.
因為那種病是找不到的.
到找到的時候都太遲了.
然後.
我都不敢跟他講是甚麼.
我都不敢跟我的師兄講是甚麼.
我都為他祈禱.
但過程裡.
一路.
由復活節到現在.
覺得他有一件事.
他將最好的留給家人.
最好的精神.
和健康狀態留給家人.
覆診.
等公立醫院去安排.
因為私家實在太貴.
然後.
有些甚麼.
去深圳.
去其他地方驗一下.
他說.
先過暑假.
等女兒有個.
開心的暑假.
嫂子.
都是跟我們兩個講.
你們兩個勸勸她.
到.
應該是上個月.
一次家庭聚會.
其他幾個哥哥.
才知道她的情況.
其實她.

$^{801}$很輕鬆地去面對.
其中一些原因是.
她和神的.
關係很好.
她和神的關係一直都很好.
她很重視.
祈禱.
她這個病她都一直放在禱告裡面.
是驗不到就.
繼續先驗不到.
到時候要看醫生.
我們說陪她進去.
你坐車來陪我走進去.
看急症.
她真的十分鐘就走進去.
急症室 由她家.
你坐車來陪我走進去.
急症室.
沒有辦法.
她自己去.
我們都去急症室.
等她.
因為想聽聽醫生怎樣說.
她都想.
盡量這段時間.
都上班.
因為她都知道.
醫院驗不到.
那是什麼呢.
留一些東西給女兒.
女兒還在讀書.
生命將來去哪裡.
她很.
她很有把握.
她是教會的終身執事.
教會按立了她終身執事.
其實應該是二十年前.
教會要按立她做長老.
不過她說.
我年輕 還有很多試探.

$^{841}$不要現在按了我將來做錯事.
所以她那時候就沒有去.
還有她做長老.
要講道.
這些她OK.
她和神保持一個好的關係.
等於姐妹.
我們要預備好.
不知道什麼時候會有什麼事情發生.
可能是很開心.
可能都很辛苦.
就算經濟的情況.
都有時會很開心.
有時會很辛苦.
我們要預備好.
進入一個更美好的狀況.
進入靈生.
我想問問你.
如果我們這樣一起祈禱.
你願意嗎.
我們跟天父講.
求祂保守我們常在祂的愛中.
求祂提醒我們.
保守自己常在神的愛中.
仰望我們主耶穌基督的憐憫.
為了永生做好準備.
你願意這樣禱告嗎.
如果你願意.
我邀請你.
我們一起站起來.
一起祈禱.
我們可以站起來一起禱告.
我們合眼一起祈禱.
親愛的天父.
來到你的面前.
我們很祈求你幫我們.
祈求你幫我們.
我們很祈求你幫助我們.
求你時時保守我們.
在你的愛中.

$^{881}$在你的照顧中.
我們又祈求你提醒我們.
我們都很想保守自己.
常在你的愛中.
但有時我們會忘記或做得不好.
求主你積聚聖靈.
時時提醒我們.
能夠時時保守自己.
行在合你的心意中.
行在你的愛中.
盡自己的本份.
我們又求你提醒我們.
仰望主耶穌基督的憐憫.
為永生做好準備.
奉主耶穌基督的聖名禱告.
阿們.
大家請.
\newpage



\section{哥林多前書 13:1-13}
\label{sec:V1IfIE0yrWg}
\textbf{2025.07.13 《求主賜異象》 經文:哥林多前書13\_1-13 曾裔貴牧師}
\newline
\newline
連結: \href{https://youtube.com/watch?v=V1IfIE0yrWg}{\texttt{https://youtube.com/watch?v=V1IfIE0yrWg}} ~~~~ 語音日期: 2025-07-15
\newline
\newline
\hyperref[sec:m1thLVxgZE8]{\small{< < < PREV SERMON < < <}}
~
\hyperref[sec:index_chronic]{\small{[返順時目]}}
~
\hyperref[sec:index_scriptual]{\small{[返順卷目]}}
~
\hyperref[sec:tbFVD8f2sac]{\small{> > > NEXT SERMON > > >}}
\newline
\newline
哥林多前書 13:1-13
\newline
\begin{longtable}{cl}
\hline
\hline
章節 & 經文 (和合本修訂版)\\
\hline
13:1 & \begin{tabularx}{0.7\textwidth}{X} 我若能說人間的方言，甚至天使的語言，卻沒有愛，我就成為鳴的鑼、響的鈸一般。 \end{tabularx} \\ \\ \relax
13:2 & \begin{tabularx}{0.7\textwidth}{X} 我若有先知講道的能力，也明白各樣的奧祕，各樣的知識，而且有齊備的信心，使我能夠移山，卻沒有愛，我就算不了甚麼。 \end{tabularx} \\ \\ \relax
13:3 & \begin{tabularx}{0.7\textwidth}{X} 我若將所有的財產救濟窮人，又犧牲自己的身體讓人誇讚，卻沒有愛，仍然對我無益。 \end{tabularx} \\ \\ \relax
13:4 & \begin{tabularx}{0.7\textwidth}{X} 愛是恆久忍耐；又有恩慈；愛是不嫉妒；愛是不自誇，不張狂， \end{tabularx} \\ \\ \relax
13:5 & \begin{tabularx}{0.7\textwidth}{X} 不做害羞的事，不求自己的益處，不輕易發怒，不計算人的惡， \end{tabularx} \\ \\ \relax
13:6 & \begin{tabularx}{0.7\textwidth}{X} 不喜歡不義，只喜歡真理； \end{tabularx} \\ \\ \relax
13:7 & \begin{tabularx}{0.7\textwidth}{X} 凡事包容，凡事相信，凡事盼望，凡事忍耐。 \end{tabularx} \\ \\ \relax
13:8 & \begin{tabularx}{0.7\textwidth}{X} 愛是永不止息。先知講道之能終必歸於無有；說方言之能終必停止；知識也終必歸於無有。 \end{tabularx} \\ \\ \relax
13:9 & \begin{tabularx}{0.7\textwidth}{X} 我們現在所知道的有限，先知所講的也有限， \end{tabularx} \\ \\ \relax
13:10 & \begin{tabularx}{0.7\textwidth}{X} 等那完全的來到，這有限的必消逝。 \end{tabularx} \\ \\ \relax
13:11 & \begin{tabularx}{0.7\textwidth}{X} 我作孩子的時候，說話像孩子，心思像孩子，意念像孩子；既長大成人，就把孩子的事丟棄了。 \end{tabularx} \\ \\ \relax
13:12 & \begin{tabularx}{0.7\textwidth}{X} 我們現在是對著鏡子觀看，模糊不清；到那時，就要面對面了。我如今所認識的有限，到那時就全認識，如同主認識我一樣。 \end{tabularx} \\ \\ \relax
13:13 & \begin{tabularx}{0.7\textwidth}{X} 如今常存的有信，有望，有愛這三樣，其中最大的是愛。 \end{tabularx} \\ \\
[1ex]
\hline
\hline
\end{longtable}
$^{1}$感恩天父召集我們能夠在今天聖餐主日一同來敬拜你.
十字架是我們人生一個很重要的根基.
因為主將我們拯救.
我們是屬於主的兒女.
我們在地上成為神的兒女.
成為基督的身體.
我們一同為著要尊崇,傳揚,榮耀的主的十字架福音.
盼望我們能夠在不同的崗位,不同的人生都能夠看到.
唯有愛的事奉才能夠真正將人帶進永恆的救恩裡面.
求主賜下儀象在我們大家的心靈.
我們有初信的,或者在座的,有些可能是誤導.
也有信了一段時間,有相當的人生經歷.
我們都一同等候在主面前.
很盼望主的聖靈.
藉著今天我們看十三章.
這一段這麼重要的經文.
我們能夠得到心靈的啟發.
我們真的能夠在我們的生命裡面.
由量的改變成為質的質變.
我們能夠在一個有品質的事奉人生.
我們這樣仰望禱告.
奉耶穌基督的名.
阿們.
相信第十三章大家去分禮絕對不陌生.
對一對新人的訓練是非常恰當貼切的.
不過很可惜.
整段聖經從來不是講婚姻.
從頭到尾都不是講婚姻.
這段聖經是橫跨在十二,十四章.
保羅叫哥林多教會要追求恩賜.
但是在恩賜的事奉裡面.
就要留心有兩樣很重要的.
原來所有的恩賜都是來自聖靈.
隨他的意思來給每一個信徒.
每一個信徒都會有神聖靈所賜的恩賜.
每一個都有.
這是十二章所講.
接著到十四章.
每一個信徒要追求最大的恩賜.
就是先知的講道.

$^{41}$這是十四章重點的總結.
十四章恩賜方面收集了兩個恩賜.
就是先知講道和方言的比較.
你可以想像到.
第十三章的重要.
就是介乎兩者之間的最重要.
保羅稱為最妙的道.
求恩賜求儀仗從來都不陌生.
7月5日日本世紀大地震.
言之鑿鑿的這位龍樹亮.
日語就是那個字.
我也不太懂得讀.
他所寫的一本虛構小說.
我所看見的未來.
很早很早期.
預言2025年7月5日凌晨四點.
日本會有史無前例的大地震.
所以在當日凌晨四點.
在Youtube同時間有25萬人.
不知道是來自世界各地.
還是日本本土.
25萬人同時等候大地震.
如果在日本就不用等了.
你一定死了.
如果真的有大地震的話.
可想而知原來世人對於預言和意象.
是趨之若鶩的.
所以連帶旅遊.
也有飛機說沒有人乘坐.
因為不敢去.
不知道你敢不敢.
也過了.
我們看見沒有事發生.
作者甚至說.
其實不是我的意思.
是出版社誤會了我的意思.
這就是意象.
但原來當我們看聖經.
不是未信的人才相信意象.
原來當時哥林多教會.

$^{81}$他們是追求恩賜,異能和異象.
所以教會追求異象本來不是錯事.
因為保羅就是要告訴頂子梅.
所有的恩賜是來自神的.
如果是神賜的真正的異象恩賜.
是來自神的.
因為當時哥林多教會.
他們有些未信耶穌之前是拜鬼的.
是拜偶像的.
是去偶像的廟.
去參拜和吃那些祭偶像之物.
所以保羅很想要他們弄清楚.
一切的美善恩賜都是來自神的.
異象不單止是夢境和啟示.
聖經保羅告訴我們.
真正的異象是神對我們的生命和使命.
安排我們人生的使命的一種光照.
所以今天我們說的求異象.
我們就要在這個部分.
我們的內心被光照這個部分.
再一次思想我們所得到的.
所擁有的神賜的恩賜和異象.
我們一起進入經文的時候.
和大家看看第十二章.
這裡告訴我們原來每一個基督徒.
都領受聖靈的恩賜.
這是第十二章.
我忘記了是第八節說的.
你可以翻看.
不會有人沒有恩賜.
這是很重要的.
恩賜的目的.
保羅說不是為了炫耀自己.
而是為了造就整個教會.
我們要把握這些大原則.
否則我們就會走錯了.
你越有恩賜.
而你只是炫耀你自己.
這個恩賜就會收回.
甚至這個恩賜會影響.

$^{121}$是一種壞影響對人.
因為你炫耀你自己.
真正的恩賜是服侍教會.
叫教會得到益處.
這是十二章很重要的總結.
接著保羅也在同一時間.
十二章告訴我們.
教會就是基督的身體.
需要彼此的配搭.
不同的器官.
讓身體的整體.
發揮正常的功用.
手的功用就是拿起一件物件.
眼的功用就是看前面.
腳就是站在這裡.
現在我需要站在台上.
所以需要用腳支撐我這個人.
還有我的腦在運作.
我的口在說話.
由我的腦指揮我的口去說.
今天我預備好的訊息.
所以是整個人.
整體運作.
由醞釀訊息.
等候的題目.
到預備的材料.
全部在這個準備過程.
是為著去造就教會.
希望今天這個題目.
神給我感動去祝福和給錢免禮.
這就是基督的身體.
弟兄姊妹.
你知不知道你是其中一個身子.
其中身子裡的器官.
我們每一個基督徒.
你回來.
你成為神的兒女.
你就是基督身體裡的一員.
很重要.
沒有你不行的.

$^{161}$每個人有自己的崗位.
有自己的作用.
每個人.
所以最大的奧秘.
就是你知道自己是基督徒.
和你發現.
神安排你在這個身體裡.
你有什麼功用.
那個功用就是很好的一件事.
因為你知道你在發揮.
你試想一下.
如果我的右腳跌撞到.
我現在很痛.
我要靠著左腳.
我根本就站不住半個小時.
你可以想像到.
教會就是每一個被主.
來刺下恩賜,裝備,訓練.
令你的人生.
可以在整個教會裡.
成為合用的器皿.
真的要記得這件事.
每個人你有用.
每個人都有恩賜.
這樣繼續下去的時候.
你就會發現一件事很特別.
你看看保羅這麼誇張的修辭.
你就可以想像到.
當時哥林多教會追求.
渴慕其中一兩樣的恩賜.
而忽略,貶低其他的恩賜.
你就會變成教會很不正常.
真正的正常就是.
臣隨己而將各樣恩賜賜給每一個信徒.
你看看保羅的修辭.
由28至30節.
本來筆字在原文是沒有的.
但你可以肯定你加上去是對的.
因為不斷地說了7次.
個個都是使徒嗎?.

$^{201}$個個都是先知嗎?.
個個都是教師嗎?.
個個都是行異能嗎?.
個個都是得恩賜治病嗎?.
個個都是可以講方言嗎?.
個個都能夠翻方言嗎?.
每一句可以加一個就是.
不可能,不是,很清楚,缺不.
你可以想像到.
原來保羅這個恩賜的清單.
是有這麼多種的恩賜的.
但是原來每個人.
神會安排不同的恩賜.
使整個教會.
得到健全的運作.
每一個教會就是基督的身體.
每一個教會就有神安排給的不同的恩賜.
為了我們能夠在教會裡.
豐富我們對舊恩,對耶穌,對傳福音的認識和參與.
所以這裡的恩賜的多樣性.
保羅在這裡好像不厭其煩.
這樣來做修辭去講.
證明了他內心裡很激動.
為什麼你們會追求單一的某一樣很明顯的恩賜.
來排擠其他人呢?.
來覺得自己很厲害?.
保羅說不應該這樣的.
應該按照聖靈在每個人的心裡的感動.
彼此欣賞,彼此合作,彼此的同工.
使整個教會得到建立和造就.
繼續下去.
你會想到保羅的誇張還要延伸到剛才主席所讀的經文.
一到第六節.
講的很誇張.
中間那裡.
如果無礙我有這些言語,奧秘的能力.
都是與我這個講的人無益的.
能夠講萬人的方言.
所有當時你想到的不同國家語系的方言.
講到天使的話語.

$^{241}$相當奧秘了.
再高層次一點.
靈界的言語.
先知的講道.
保羅當然有了.
他明白各種各樣的奧秘.
舊約能夠解夢的人.
譬如說但爾尼.
解了兩個大的夢.
尼布甲尼撒王.
發了兩個很大的夢.
第一個夢.
就是他夢見一個金像.
然後就什麼都忘記了.
要術士解那個夢給他聽.
術士說沒問題.
你就講出來.
那個王說.
那個夢我已經忘記了.
你們幫我解吧.
很誇張很離譜.
所有解不了的.
全部準備拿去殺了.
他們就很急找了但爾尼.
當時叫做伯提沙薩.
巴比倫的名字.
伯提沙薩就祈禱.
得到那個異象的解釋.
就去跟王講.
另一個異象.
是直接.
又是一樣.
這次王解了那個夢的原因.
那個內容.
但所有術士都解不了.
巴比倫的術士.
最終又是要找但爾尼.
但爾尼就幫他解.
接著講了一句話.
王啊.

$^{281}$求你接納我的諫言.
使你能夠好好的.
神憐憫你.
你行公義.
你行義行.
巴比倫王聽到夢解得很周詳.
OK OK 沒問題沒問題.
一年之後.
有一天.
在一個大巴比倫的空中花園.
整個國家.
豈不是我尼布甲尼薩王所創辦的嗎.
還沒講完這句話.
天上有聲音跟他說.
你的刑罰就到了.
就是照著但爾尼幫他解夢的那個內容.
他這樣受到懲罰.
整個人被人突然間瘋癲了.
送到野外.
像一隻牛一樣吃草.
一段很長的時間.
才讓他有智慧.
因為他突然間患了瘋狂症.
那些人要把他趕出.
困住他.
用鎖鏈鎖住他.
但爾尼解了夢.
這就是各樣的奧秘.
普羅說如果我有各樣的奧秘可以解.
我有各樣的知識.
我有傳遞的信心.
就是說.
尼祿王我要你死.
神要你死.
因為你現在逼迫我們基督徒.
他可以這樣做.
神就可以透過普羅來懲罰當時的羅馬皇帝.
但普羅有沒有這樣做呢.
不過普羅在這裡很誇張地.
將一個重點襯托出來.

$^{321}$就是如果在教會.
在事奉.
在運用恩賜上.
這一切你已經擁有了.
你沒有愛去運用這些恩賜.
第二個部分就是.
對別人的一個施捨.
周濟和捨身.
好像用他自己的一個精神去感染其他人.
普羅說這一切如果沒有愛.
一切都是沒有用沒有益.
這就是第十二章和十三章.
普羅說的誇張的修辭.
為著一個原因.
不是婚姻的勸勉.
是教會的事奉.
怎麼能夠由量的變成為質變.
個人和教會.
怎麼能夠有質的變化.
每個人的事奉是有很高的質量.
這個質量就是在神裡面的愛.
如果不是的話.
哥林多教會.
正正就是一個互相競爭.
彼此來結黨的氛圍.
以事奉為競爭.
以自己的恩賜來體小.
體低.
體高自己.
排擠其他的弟兄姊妹.
普羅見到這樣的情況.
他就告訴我們.
事奉我們要追求的.
當然最好的恩賜就是先知的講道.
但是最妙的道.
就是神的愛在我們眾人的心裡.
開始十三章的七至八節.
普羅開始說.
怎麼能夠具備擁有質量的事奉.
凡是包容.

$^{361}$凡是相信.
凡是盼望.
凡是忍耐.
而且因為裡面有一個很大的泉源.
就是愛.
所以在這樣的裡面.
是不會止息的.
來包容.
相信.
盼望.
和忍耐.
一直預備的時候.
看到一個例子.
有一個在醫院裡面.
管理層的護士長.
他不是基督徒.
很專業.
也很繁忙.
所以他對他的下屬.
對他的病房.
那個管理.
極其的嚴格.
要求的.
只是不要做錯.
有一天.
有一個老婆婆.
年紀很大了.
是一個很有愛心的基督徒.
患了病.
年紀很大.
進來.
看到這個護士長.
這樣的嚴格.
嚴厲.
對下屬.
對其他病人.
只有鐵面.
沒有愛心.
沒有笑容.
做好自己的份內事.

$^{401}$千萬不要做錯.
這個婆婆.
每天為這個護士長祈禱.
希望她能夠有愛心.
去工作.
每次跟她說話.
都帶著一份對她的愛的接納.
久而久之.
這個婆婆.
這樣的說話.
這樣的笑容.
感染了這個護士長.
她自己再看自己的工作.
自己每天的那種嚴厲.
一換上制服之後.
整個人鐵面無私.
人人都怕她.
她就開始有一個反思.
為什麼這個婆婆這麼重病.
她仍然有笑容.
有平安.
慢慢感染了她.
她真的有一個改變.
她真的發現.
原來她過去只是做了一些.
表面是對的.
但是對人的工作.
她沒有做到.
她沒有作為一個護士.
應該對一個病人的.
接納愛.
她有一種反省.
所以她自己也有一個悔改.
後來有沒有信主.
我就不知道了.
那個故事就停在那裡.
那個例子.
一個老婆婆.
後來都真的離開世界了.
但是令到這個護士長.

$^{441}$有一個生命的反思.
親愛的弟兄姊妹.
十三章.
可以說是非常重要的.
聖經的一章聖經.
是告訴我們.
我們原來已經擁有.
這份十字架耶穌的大愛.
而且還是聖靈.
在我們大家的心裡面.
將恩賜給我們.
本來我們根本連信耶穌都不懂.
都不會.
但是當我們信了耶穌之後.
我們就得到聖靈.
按照著我們的才幹.
我們過去的成長.
甚至是度身訂做的.
恩賜給我們每一個人.
為著的是一個整體.
為著我們能夠有一個合一的服事.
使到神的教會被建立.
所以每個人.
都有你自己的功用.
神給你的崗位.
但是我們發現之後.
我們還不能夠再進一步的轉化.
就是不知不覺.
我們擁有了恩賜.
我們很開心地.
用我們的恩賜.
但是忘記了.
跟其他人一起.
在愛和合一裡面運用.
這就是哥林多教會.
當時出現嚴重的問題.
有恩賜.
是神賜給你的恩賜.
但是不是用在神的愛裡面.
去建立教會.

$^{481}$而是來到不知不覺.
跟別人有比較.
在這個比較裡面.
就會帶來傷害.
會帶來一些人受到排擠.
這就是在十三章.
我們需要避免的地方.
異象必須是有愛作為動力.
不能夠是冷漠.
像剛才我說的護士長.
帶著一種冷漠的鐵面.
來做好他的專業工作.
但是教會.
或者任何的地方.
我們怎麼能夠只是做好自己的工作呢.
做工作的背後.
如果不是一份很大的愛心.
我們怎麼能夠讓其他人.
得到那個建立.
得到那個益處呢.
神賜給我們的異象.
是附帶著一個.
能夠改變生命的能力.
真的.
神賜給我們的異象.
恩賜,才幹.
其實是伴隨著.
我們有改變生命的能力.
讓我們能夠慢慢得到很大的啟發.
我們的生命得到一個轉化.
使我們所做的不是做對,做好,不錯.
這樣就叫做解決了那件事.
而是要讓整個有機體.
就是和其他人一起在愛裡面成長.
相信我們真的必須要由量到質的改變.
剛才主席Wallace 讀了第十二,十三節.
是告訴我們在人間還沒見耶穌之前.
還沒回到永恆天家之前.
現在我們有信,忘,愛.
都是同時間在我們每一個基督徒的內心.

$^{521}$在整個運作裡面.
最重要的保羅話.
最大,最重要的.
就是在愛裡面行使我們的恩賜.
來服侍其他人.
最後我想用一個例子來結束.
我漏了這個.
沒加進去.
不要緊,這個就是了.
就是這個很重要.
對我的生命有很大的影響的真實故事.
我在2016年的時候.
無意中在美國一份雜誌看到.
有一個小學的老師.
他教五年班.
他儘管用他的名字叫.
因為我忘記了他的姓氏.
儘管叫他Kitty.
這位老師.
他在五年班開學的時候.
跟全班的同學說.
美國來的.
跟全班的同學說.
我作為老師,我的專業.
我是一視同仁.
我不會偏心,我不會對某個同學有針對.
講得很清楚.
不過很快他就發現其中一個同學.
我記得他的名字Teddy.
是真人真事.
這個同學上課睡覺.
毫無精力.
做功課很馬虎很馬虎.
老師說明了,我不會針對任何一個同學.
我會因材施教.
但是令他忍不住.
用紅筆寫一個大的D.
括著,然後就X,英文的.
令這個同學,算了,無所謂了.
你喜歡怎樣就怎樣.

$^{561}$上課完全是這樣睡覺.
但是這個老師發現.
原來學校有一個規矩.
就是他必須看回每個同學.
過去的班主任給的評價.
他就如是者逐個同學去看.
最終去到Teddy.
一年班的班主任就說.
Teddy是一個很活潑很乖.
很得到同學歡心的小朋友.
如是者到二年班的班主任說.
Teddy最近不開心.
悶悶不樂.
因為他媽媽患了重病.
他很擔心.
他不知道怎麼辦.
偶然有時爸爸要照顧太太.
不可以帶他上學.
所以Teddy現在很需要有人關心.
三年班的班主任的評價.
Teddy媽媽的離世.
令Teddy現在極需要有人關心.
如果再沒有人關心.
他就會更加失落.
因為父親已經不能照顧他.
父親沉迷於喝酒.
借酒逃避太太突然離世.
四年班就是上一年班主任的評價.
Teddy徹底放棄.
他現在的表現.
令他不能再有任何學習感興趣.
因為他父親也放棄他.
他自己也放棄自己.
所以當老師見到這四年班主任的評價後.
他自己突然發現.
原來Teddy有今天.
不是因為他的教書.
而是因為他遇到一個很大的.
他自己成長招架不到的困難.
他哭得很厲害.

$^{601}$接著那年快到聖誕節.
每個同學會送一份禮物給老師.
每個送很好的禮物.
當然不一定很貴重.
送給老師.
包裝得很漂亮.
Teddy用一個很簡單的雞皮紙袋.
就是普通雜貨店美國人買東西的雞皮紙袋.
包著一個手握人造鑽石.
當然是假的.
而且那些鑽石都脫落了.
送給老師.
全班笑得Teddy臉都黃了.
還有半瓶用過的香水.
送給老師.
老師見到.
完全不同了.
這四份報告讓他的心靈發現.
教學不應該看表面.
他用心去接納Teddy.
帶著他的手握.
還有塗一點香水.
完了.
每個人都走了.
Teddy留下跟老師說.
老師你好像我媽媽.
你的香氣和你的手握.
讓我想起母親.
這位老師在Teddy走了之後.
哭了差不多一個小時.
終於發現原來他可以得到Teddy的心.
用心去教他.
去安慰他.
支持他.
Teddy六年級畢業.
考得很好.
他畢業禮的時候.
多謝老師對他的悉心照顧.
六年之後他中學畢業.
老師收到一封信.

$^{641}$你是我遇過最好的老師.
多謝你.
這個小學老師就是這樣去與事者.
後來.
這個真人故事告訴我們.
Teddy入大學.
最終他完成了博士學位.
很久很久之後.
收到一封信.
一份邀請.
問這位老師.
你能不能出席我的婚禮.
因為在讀博士期間.
遇到一個女同學.
大家非常投契.
要結婚成為終身伴侶.
而且還要邀請她.
因為我的爸爸早幾年已經死了.
你能不能作為我的姦會人.
成為幫我簽字的姦會人呢.
親愛的弟兄姊妹.
人間都有很多這樣的故事.
是量度質的改變.
生命被真誠的愛感化.
我們再看的是人.
不同的需要.
我們不是做對貼面無私.
而是我們要在生命裡.
反映十字架的愛.
在我們大家的生命裡.
很盼望這些感動的例子.
感動我們每一位.
仕途保羅告訴我們.
凡事包容.
凡事相信.
凡事盼望.
凡事忍耐.
愛是永不止息.
願主祝福大家.
也使我們洪恩家.

$^{681}$一起在主的恩典裡服侍.
我們祈禱.
親愛的主.
多謝你.
我們今早來思想求主賜儀仗.
我們需要的是愛在我們生命裡.
讓我們看到.
原來我們可以靠著主.
十架的大能.
和永恆不改變的愛.
首先改變的是我們的生命.
是我們對人.
對事物的看法.
幫助我們一起.
在主的面前.
我們去反思.
未必是一個悔改.
反而是一個沉澱.
我們在神面前.
我們再一次看不同的人的價值.
求你施恩.
我們是有不同的.
因為我們來自不同的背景.
更加主給我們有不同的恩賜.
但是不同的恩賜.
就是豐富神的家.
那個合一.
那個被主用.
奉耶穌基督的名.
阿們.
\newpage



\section{希伯來書 5:11-6:1}
\label{sec:tbFVD8f2sac}
\textbf{2025 7 27主日崇拜講道:《竭力成為一個成熟的基督徒》 經文:希伯來書5\_11-6\_1 曾敏姑娘}
\newline
\newline
連結: \href{https://youtube.com/watch?v=tbFVD8f2sac}{\texttt{https://youtube.com/watch?v=tbFVD8f2sac}} ~~~~ 語音日期: 2025-07-30
\newline
\newline
\hyperref[sec:V1IfIE0yrWg]{\small{< < < PREV SERMON < < <}}
~
\hyperref[sec:index_chronic]{\small{[返順時目]}}
~
\hyperref[sec:index_scriptual]{\small{[返順卷目]}}
~
\hyperref[sec:YkzDcCtXDZ8]{\small{> > > NEXT SERMON > > >}}
\newline
\newline
希伯來書 5:11-6:1
\newline
\begin{longtable}{cl}
\hline
\hline
章節 & 經文 (和合本修訂版)\\
\hline
5:11 & \begin{tabularx}{0.7\textwidth}{X} 論到這事，我們有好些話要說，可是很難解釋，因為你們聽不進去。 \end{tabularx} \\ \\ \relax
5:12 & \begin{tabularx}{0.7\textwidth}{X} 按時間說，你們早該作教師了，誰知還需要有人再將神聖言基礎的要道教導你們；你們成了那需要吃奶、不能吃乾糧的人。 \end{tabularx} \\ \\ \relax
5:13 & \begin{tabularx}{0.7\textwidth}{X} 凡只能吃奶的，就不熟練仁義的道理，因為他是嬰孩。 \end{tabularx} \\ \\ \relax
5:14 & \begin{tabularx}{0.7\textwidth}{X} 惟獨長大成人的才能吃乾糧，他們的心竅因練習而靈活，能分辨善惡了。 \end{tabularx} \\ \\ \relax
6:1 & \begin{tabularx}{0.7\textwidth}{X} 所以，我們應當離開基督道理的基礎，竭力進到成熟的地步；不必再立根基，就如懊悔致死的行為、信靠神、 \end{tabularx} \\ \\
[1ex]
\hline
\hline
\end{longtable}
$^{1}$今天這個題目我想先問一問.
我今天一早就有人跟我說.
嘩!這個很難做到的.
這個做不到的.
我想問有多少人覺得這個做不到的.
你覺得做不到的你舉手.
有人很誠實的舉手.
你不需要介意別人的看法.
你覺得做不到的就誠實表達.
我欣賞頂尖的坦誠.
有沒有.
不要受旁邊的人影響.
你覺得我做不到的就舉手.
還不知道是吧.
我很想說.
這件事我們一定做到的.
聖經裡面說的話.
不會令我們做不到的.
做不到的話.
神就不會在當中讓我們看到這段話.
不會只是放一大堆很高深的文字在那裡.
然後我感覺.
嘩!他真的很高深很崇高.
在天上空的.
等我欣賞完之後.
我要看地面的生活.
我感覺相差很遠.
然後經常都很沮喪.
不會的.
我想跟你說.
今天這段話.
在座所有人都做到的.
是所有人.
我們首先在這裡認識一下.
這段經文來自希伯來書.
我相信在座不是每個人都很熟悉.
希伯來書.
有誰都說自己看過這本書呢.
請大家舉手.
我看過希伯來書的.

$^{41}$看過.
是吧.
都有些人看.
當中有些人從來都沒有看過.
不要緊.
經過今天之後.
我鼓勵你.
可以去看一看.
我們看看希伯來書的背景是甚麼.
如果大家仔細翻開這本書的時候.
你會留意到.
裡面會提到很多一些舊約的人物.
很明顯.
那個作者.
是相當熟悉.
猶太人過去他們所相信的.
信仰的背景.
很熟悉舊約.
非常熟悉.
而他寫信.
這個希伯來書.
可以這樣說.
好像一封信.
寫給當時的人.
更貼切地說.
其實這是一篇講章.
是一篇講章.
這篇講章寫給甚麼人呢.
作者熟悉猶太人的背景.
熟悉舊約的很多聖經背景.
同樣地讀者也很熟悉.
讀者是有猶太教背景的.
基督徒.
所以希伯來書.
不是寫給那些未信的人.
是寫給基督徒.
這些基督徒有猶太教的背景.
當時在這個地方.
仍然是羅馬政權在掌管.
猶太教被視為合法的宗教.

$^{81}$所以你說自己是猶太教的信徒.
政府會放過你.
你仍然可以享受很多福利.
沒有人會逼迫你.
你可以自由來往.
自由在會堂裡敬拜.
但是基督徒不同.
基督徒在當時.
基督教不是一個合法的宗教.
所以基督徒會面對逼迫.
因為逼迫而來有很多不好處.
有很多生活上的壓力.
很多的掣肘.
很多的困難.
所以當時來說.
一群由猶太教轉為信耶穌的基督徒.
其實是很大的壓力.
如果他們回到之前.
猶太教擁有自由.
不會被逼迫.
沒有那麼多痛苦.
沒有那麼多壓力.
但是信了耶穌之後就不同了.
很多壓力,逼迫,困苦.
生活上有很多艱難.
所以經常有一個趨向.
很容易陷入一個試探.
走回頭.
不如離棄耶穌.
離棄耶穌.
我就不用受這些苦.
我不用受這些壓力.
我可以生活得很自由.
如果我看重生活的享受.
就更加是這樣.
為什麼要那麼辛苦呢?.
一直逼迫.
好像不能呼吸.
這是當時基督徒的情況.
在這個情況下.

$^{121}$我想多說一點.
在這個情況下會出現什麼問題呢?.
在裡面有些人.
信仰很容易有背道危機.
我會放棄我的信仰.
而有些人慢慢地.
開始懶散.
很散漫.
信耶穌那一刻很火熱.
我為了信仰赴湯蹈火.
上高山下油鍋.
什麼都可以.
過了一段時間.
一年兩年三年.
我開始很散漫.
很懶惰.
我一直停留在當初的階段.
其實不只是停留.
一停留就有危機.
倒退.
壓力來到.
捏著我.
逼迫我.
我想放棄.
我想背棄耶穌.
在這個過程中.
有很多很多的問題.
在這裡.
我們這一段經文.
一開始的時候.
一開始的經文在十一節說.
論到這事.
我們有好些話要說.
我就要說說.
五章在十一節之前.
在說什麼呢?.
一到十一節.
其實是在說耶穌基督.
作為天上的大祭司.
祂是一個非常超越的大祭司.

$^{161}$經盟約祂和默基世德.
作為一個比較.
說祂是照著默基世德的體系.
永遠為祭司.
說到默基世德.
大家知不知道祂是誰呢?.
是不是很害怕我問這些問題?.
默基世德是誰呢?.
是大祭司.
在哪裡出現過呢?.
好啊,這個是有看過《創世記》的.
在《創世記》的時候.
大家記不記得.
在四王和五王的戰爭裡.
亞伯拉罕的侄羅德被擄去.
接著亞伯拉罕捲入戰爭當中.
去拯救他的侄羅德.
在這場戰爭之後.
他贏了.
他回來的時候.
就說這個祭司默基世德.
去迎接他.
其實默基世德出現的地方不多.
反而他的名字出現最多.
就是希伯來書.
說他是非常非常不同.
不同地面上的大祭司.
是不同的.
他是超越了地面上的.
一般的大祭司.
在這裡.
作者就將他和默基世德.
作為一個比較.
耶穌基督是一個非常超越的大祭司.
他不是地面上普通的.
接著去到十一節.
說到論道者事.
就是在說.
作者很想再說.
關於默基世德的事情再說多一點.

$^{201}$我剛才就這樣和大家說.
默基世德大家都呆了.
不知道是誰來的.
我都沒有聽過他的名字.
怎麼辦呢?.
其實那些讀者都有這樣的情況.
論道者事.
我們有好些話要說.
作者都還想再說多一點.
默基世德的事情是怎樣的.
耶穌基督作為大祭司是怎樣的.
為什麼用他作為一個比較呢?.
但是當時的作者已經覺得.
那件事很難解釋.
那個問題不是默基世德.
背後一連串的一些事情.
有多難解釋.
而是對著那些受眾.
那些受眾不明白.
所以就說很難解釋.
很難解釋.
因為為什麼我們聽下去.
生命不成熟.
我們今天可以透過神的話語.
做一個檢視.
我們照鏡子.
照鏡子看屬靈的鏡子.
去看看我們的生命成不成熟.
我們一起看.
是不是?.
這個生命不成熟有什麼表現呢?.
表徵呢?.
第一件事是聽不進去.
不能夠聽進去.
因為你們聽不進去.
所以我再多說一點都說不出來了.
聽不進去.
這個聽不進去是什麼意思呢?.
原來他說聽不進去.
是直接翻譯來說.

$^{241}$耳朵很懶惰.
很遲鈍.
直接翻譯是這樣.
是因為你們的耳朵很懶惰.
很遲鈍.
這個作者說話很有骨.
我想說這段話很有刺激性.
是一段責備的說話.
責備受眾.
聽不進去.
耳朵很遲鈍.
很懶惰.
耳朵是什麼呢?.
是屬靈真理的理解.
是接收的能力.
我是否願意去聽真理呢?.
願意去聽屬靈的事情呢?.
我們的屬靈的理解能力遲鈍.
不是因為我天生就是這樣.
你不要要求我那麼多.
我不是很明白的.
很多事情.
我就是理解那麼極端.
但是作者用這個詞語.
是說屬靈的理解能力不是天生的.
不是天生的.
是怎樣來的?.
是後面的時候沒有努力去上進.
作者是這樣說的.
沒有上進.
所以我才會變得遲鈍.
我接收不到.
我很感恩.
神給我這個星期有很特別的經歷.
神真的很幽默.
我預備講章的時候.
也沒有想過和我的探訪有什麼關係.
有一天我去探訪一個姐妹的娘家.
姐妹的媽媽正在回教會.
她生活裡有很多經歷.

$^{281}$我去探訪她.
神在那過程中給我大開眼睛.
那位姐妹完全一粒字都不懂.
我記得我曾經在那裡推讀經的時候.
也有聲音說.
我們當中挺姐妹都不懂字.
怎樣推讀經?很難推的.
我去到那位姐妹的媽媽家.
她一粒字都不懂.
但是她信耶穌.
每個星期都堅持回教會.
她身體也不舒服.
也有不少問題.
要定時定候看醫生吃藥.
也不方便每天出門.
在這樣的情況下.
她又不懂字.
那怎樣看聖經呢?.
她的女兒隔天打電話給她.
讀經文給她聽.
就和她背誦聖經.
這位姐妹的媽媽很用心背誦.
一直背誦.
背完又背,背完又背.
一節又另一節又另一節.
背完又背.
我去探望她的時候.
聊天的時候.
她說那些經文上半節是怎樣讀的呢?.
好像是這樣,我記錯了.
又再來,又背.
又再背.
我心裡真的很佩服.
其實我應該這樣說.
我有點慚愧.
我做傳道人以前很努力地背經文.
所以現在某程度上吃老本.
背下的經文.
我就跟別人說.
最近好像很少背經文.

$^{321}$我和小朋友和主要學的也有背.
讀啊讀啊讀啊.
好像很不容易.
我也很佩服小朋友很厲害.
一直聽著我說的時候.
一直朗朗上口.
但這位姐妹已經七十多歲了.
不懂得字的.
就這樣背經文.
一直背.
她的熟齡生命非常單純.
是活潑的,有力的.
是非常有力的.
她跟我說什麼呢?.
她說我一直讀經文.
一直稱讚她的女兒很好.
就跟我去背聖經.
一直神的話語.
就令我有信心.
她跟我說.
你知不知道我人生這樣.
她真的很苦.
苦了很多年.
現在還是很苦.
我甚至乎.
夠膽這樣說.
她比我們在座大部分.
百分之九十九的人都要苦.
苦到現在還是很苦.
你想像不到她這麼苦.
每天都是煎熬.
她說我曾經想放棄.
都不知道該怎樣.
但我讀經的時候.
神的話語就提升我了.
神的話語就鼓勵我了.
就有信心了.
七十多歲.
有一位長者,不識字的.
但對淑玲的事很渴望.

$^{361}$一點都不遲鈍.
她的耳朵很渴望.
一直聽,聽到入心裡去.
所以我知道.
作者為什麼要這樣責備讀者.
就是覺得很艱難.
但她看到的是很多散漫.
很多懶惰.
甚至是背道.
是什麼呢?.
你以為是天生的嗎?.
不是天生的.
是後天做什麼事.
務求上進.
弟兄姊妹.
今天是神對我們說話.
我們是不是.
不知道多少年前.
停在那裡了.
務求上進.
以致我們的耳朵遲鈍了.
我們懶惰了.
這是第一個淑玲不成熟的問題.
聽不進去.
很遲鈍.
第二個問題就是.
原來我需要被教導.
什麼是基督的福音.
竟然.
我記得我曾經在這裡問過.
在座的弟兄姊妹.
我們信主多少年.
當然我們從剛回教會沒多久.
一定有的.
我們有新朋友.
不是很久.
有些回教會還沒信耶穌.
還在探索.
回到教會就開始信耶穌.
但可以這麼說.

$^{401}$在座大部分人.
都是信耶穌很多年.
超過五年.
十年.
二十年.
甚至上三十年.
你計算一下自己的時間.
大部分都不是初信.
但這裡的讀者們.
其實他們信了很長時間.
為什麼我知道.
經文是這麼說的.
大家聽我讀.
五章十二節說.
按時間說.
你們早該作教師了.
誰知還需要有人.
再將神聖言基礎的要道.
教導你們.
按時間說.
這個時間是什麼時間.
是信耶穌的時間.
按信耶穌的時間.
其實是不短.
一點都不短.
你們早該作教師了.
不短的時候.
其實應該對於我們.
信耶穌裡面的.
那些屬靈的真理.
有一定的認識.
可以跟別人分享.
雖然不會每一個.
走出來做主教學老師.
也不會每一個.
做團契導師.
但你一定有很多.
屬靈的真理.
可以跟別人分享.
對於教會很多新朋友.

$^{441}$或者有信心軟弱的.
你可以鼓勵他們.
你用神的話語鼓勵人.
不一定是需要傳道人才做到.
是每一個都可以的.
你是可以用神的話.
你應該成熟到.
可以教導別人.
屬靈的真理.
但竟然不僅不行.
還要什麼.
按時間說.
我還要什麼.
誰知還需要.
我再次需要有人.
再將聖賢基礎的要道.
教導你們.
我們需要有人告訴我們.
聖賢基礎是福音.
我們進入教會的門內.
是因為福音的緣故.
因為信耶穌.
我們進入教會.
一起參與教會的生活.
一同敬拜.
一同團契.
分享等等.
因為福音.
這是我和你的生命的基礎.
那是起步.
那是基督的福音.
是一個起步.
哇!信耶穌信了這麼久.
原來又要回到基礎.
還要有人再教導我們.
你們成了需要吃奶.
不能吃乾糧的人.
有沒有這個感覺.
信耶穌這麼久.
我們不能用聖經.

$^{481}$去鼓勵別人.
去教導別人.
還停留在最基本的階段.
作者責備讀者們.
這真是屬靈生命.
不成熟的情況.
我記得很久之前.
我曾經在一個群體問過.
如果別人問你.
信耶穌.
信什麼.
你會怎樣回答.
當時群體很害怕.
我問這些問題.
不知道怎樣回答.
信耶穌真的信什麼.
我怎樣回答你.
所以難怪這裡說.
誰知道還需要有人.
再將神聖賢基礎的要道.
教導你們.
再說一次福音.
我們從頭再來建立根基.
我本來應該做老師.
做不到.
到了這麼久之後.
這是不成熟的表現.
第三不成熟的地方.
就是不能分辨善惡.
不能分辨善惡.
這裡說你們成了.
需要吃奶不能吃乾糧的人.
一個比喻.
做嬰兒.
信耶穌第一年說做嬰兒.
我們很開心.
我們喜歡認小.
認小謙卑一點.
別人照顧我多一點.
第二年還是嬰兒.

$^{521}$第三年又是嬰兒.
四年又是嬰兒.
五年又是嬰兒.
十年還是屬靈的嬰兒.
整個教會都是嬰兒.
怎樣可以呢.
我們養不大.
如果是這樣.
真的很可怕.
所以你成了.
需要吃奶不能吃乾糧的人.
這個指責是不是很嚴厲.
是呀.
學嬰兒.
嬰兒有什麼問題.
有些事情分不清.
唯獨長大成人.
才可以吃乾糧.
吃奶有什麼問題.
不熟練仁義的道理.
仁義的道理是說什麼.
仁義的道理是說.
義的標準準則.
我熟練義的標準準則.
我才可以判斷是非.
分辨善惡.
所以經文是有一個平衡的.
下面說分辨善惡.
是一個平衡的作用.
我是嬰兒.
我吃奶.
我就不懂得分辨.
義的準則是什麼.
什麼是對.
什麼是不對.
所以這個有危險.
如果今天有異端.
潛伏在我們當中.
潛伏很久都可以的.
忽然間他有一天.

$^{561}$跟你說聖經.
你應該分辨不到.
我們分辨不到.
弟兄 姐妹 你很愛臣.
你真的說得很好.
你比講完更加仔細.
我不懂得.
是呀 不懂.
你再教我多一點.
就跟著他去.
我記得我以前在.
應該是一段時間.
那時候可能還在.
我忘記是感受堂.
是服侍的時候.
還是再早一點.
我仍然和我自己的母會.
上宣有聯絡.
那時候我還在住.
上水圍的地方.
我沒有想到.
那裡是其中一個異端.
當時叫東方閃電.
後來叫全能神教會.
其中一個很重要的根據地.
我因為一些很特別的原因.
認識了一個婆婆.
原來她是宗教領袖.
這個婆婆很熱情.
到處和人聊天.
我自己又想著.
我要到處認識人.
希望有機會講福音.
是否熱情就可以了.
認識了這個婆婆.
講著講著.
有一天.
她帶了另一個女士.
她後證我在哪裡住.
她走到我家樓下.

$^{601}$就要和我說她的和音.
異端的和音.
知道將來會發生什麼事嗎.
她說審判.
說不同的事情.
後來我很清楚知道.
她原來是異端.
東方閃電的一個宗教領袖.
我和自己的母會聯絡的時候.
我的牧師給我看.
是否這個.
我的牧師說是.
在我們的教會.
也有人潛伏了很多年.
是很多年.
弟兄姊妹也完全不察覺.
她就會到差不多的時候才起來.
然後和你分享.
你的心是完全不能分辨.
你會跟著她.
有一段時間.
東方閃電影響了很多人.
包括傳道人在內.
我記得短訊中心.
也有發一些訊息出來.
在過程中.
我們心裡有很多思考.
我們對神話語的認識有多少.
我們的生命.
是否很成熟,強壯,有力呢?.
異端能不能分辨呢?.
不要說到極端去到異端.
譬如有不同的人.
是作為教導的角色.
有教導的,有講道的.
你有多少能夠分辨呢?.
那些是出於聖經.
哪些不是出於聖經呢?.
那個講員真的很有魅力.
我真的很喜歡他.

$^{641}$他說什麼都是很好.
我照單傳授.
全部都聽他說就行了.
我就要你這樣.
是否這樣呢?.
你跟人.
還是你自己.
追求成熟.
好好地看聖經.
你懂得分辨.
曾姑娘今天說得怎樣.
這裡說得對不對.
你可以反過來跟我說.
要有這個信心才行.
要完全不懂得分辨.
人家說什麼.
這個一定好.
是呀.
只要是誰說的就一定好.
拍爛手掌我們跟他去.
只要是誰說的就一定好.
我們不跟他去.
完全不能分辨.
這樣是很危險.
完全沒有分辨能力的時候.
就被人影響.
飄來飄去.
這也是一個不成熟的情況.
不能夠分辨善惡.
我想說.
世代非常邪惡.
弟兄姊妹.
非常邪惡.
每一天.
是很多很多的訊息.
你手機很多訊息.
你上網聽廣道.
都很多.
你怎樣懂得分辨呢.
哪些說完是真理呢.

$^{681}$照單傳修.
怎樣懂得分辨.
哪些真的帶你走在正路上.
弟兄姊妹.
我們的生命.
去到什麼階段.
所以去到六章一節的時候.
很重要說.
我們應當離開基督道理的基礎.
去到竭力進度.
成熟的地步.
不要十年都做嬰兒.
還在說.
我不知道信耶穌是信什麼.
你再說一次.
十年前是這樣.
十年後還在說這件事.
是應該要成熟.
是要竭力.
神想我們努力.
不是靠有信心就會成熟.
是不是.
好像有些人在說.
我信耶穌之後.
讀書一定拿100分.
有信心就行.
不用讀書.
我就坐在這裡.
去到市場就給我100分.
弟兄姊妹.
你沒有信心才讀書.
是不是.
不是的.
有信心.
都要盡力竭力.
上帝給我們位置.
要做些事.
完全不做.
不參與.
只要上帝做就行.

$^{721}$我坐在這裡就行.
都要竭力.
要成長.
耳朵已經長了枕.
聽不到東西.
現在開始要用心聽.
用心看.
用心聽.
識字就看多些經文.
不識字就聽多些.
都可以.
你給自己很努力.
這樣去成熟.
要去成熟.
有一個屬靈的吸收能力.
要去分辨是非.
這是非常非常寶貴.
不只是吃奶.
是要吃乾糧.
什麼叫竭力呢.
好像很抽象.
神的話語非常寶貴.
我講的秘訣就在這裡.
好好地看聖經.
聽聖經.
這個星期.
神就讓我看到那個姐妹.
下次不要再跟我說.
我看不到聖經.
可以聽的.
家人跟你背經都可以.
家人跟你讀經都可以.
要有一個學武的心.
一路讀.
這是非常寶貴.
希伯來書也有講到神的話語的寶貴.
讀給大家聽.
在希伯來書裡.
他講到.
在四章十二節.

$^{761}$他這樣說.
神的道是活潑的.
是有功效的.
比一切兩刃的劍更鋒利.
甚至魂與靈.
骨節與骨髓.
都能刺入.
否開.
連心中的思念和主義.
都能辨明.
被造的.
沒有一樣在他面前.
不是顯露的.
萬物在他眼前.
都是赤露闖開的.
我們必須向他交仗.
神的道是活潑的.
這個活潑的意思是有生命的.
帶來生命的.
有上帝的話語.
我們就有生命.
就會成長.
神的話語就照到我們裡面的問題.
讓我們成長.
神留下了整本聖經.
信道是從聽道而來.
你看聖經.
就會相信.
就會改變.
信道是從聽道而來.
信心是從什麼而來?.
從神的話語而來.
不是空泛的.
只是在求神給我信心.
我不夠信心.
你加添給我.
神啊.
神可能更想跟我們說.
你都完全不認識我.
怎樣認識他?.

$^{801}$在神的話語裡.
認識神的話語.
我們就會認識神.
所以一定是做到的.
不會做不到的.
有一個不識字的.
七十歲的一位長者.
是七十多歲.
他都做到.
他的生命是活潑的.
單純的.
我想不到有什麼理由.
我們做不到呢?.
神的話語非常寶貴.
最後送上詩篇第一章給大家.
詩篇第一篇.
是非常非常珍貴的.
我們成熟的時候.
生命應該像是.
這個棘蚌的果樹一樣.
被人看到是有果實的.
結果是雷雷的.
很有力的.
而且不會灰心.
不會灰心.
在詩篇第一篇.
是這樣說的.
一開始的時候說.
我們不從惡人的計謀.
不站罪人的道路.
不撞傲慢人的座位.
為喜愛耶和華的律法.
晝夜思想他的律法.
這人便為有福.
我們要與罪惡劃清界線.
同時我們要愛慕神的話語.
是晝夜思想.
不是只聽到的時候就聽.
回去自由了.
我就不需要了.

$^{841}$有福的人是晝夜思想他的律法.
好像那位姐妹一樣.
他要像一棵樹栽在溪水旁.
按時候結果指.
葉子也不枯乾.
凡他所做的盡都順利.
聖禮的人就是會這樣.
但是需要神的話語.
這位姐妹在此挑戰大家.
挑戰大家好好地讀聖經.
你一定可以的.
是做到的.
所以我們可以一起成熟.
我相信紅印加將來會有很多人聚會.
但不只是多人.
是要一起成熟.
我們的生命如果不成熟.
就不能夠承載別人的生命.
人們在前門來.
也可以在後門走.
因為我們不能夠承載他們.
\newpage



\section{使徒行傳 26:19-32}
\label{sec:YkzDcCtXDZ8}
\textbf{2025.08.03:《回應天上的呼召》 經文:使徒行傳26\_19-32 講員 : 杜文瀚牧師}
\newline
\newline
連結: \href{https://youtube.com/watch?v=YkzDcCtXDZ8}{\texttt{https://youtube.com/watch?v=YkzDcCtXDZ8}} ~~~~ 語音日期: 2025-08-07
\newline
\newline
\hyperref[sec:tbFVD8f2sac]{\small{< < < PREV SERMON < < <}}
~
\hyperref[sec:index_chronic]{\small{[返順時目]}}
~
\hyperref[sec:index_scriptual]{\small{[返順卷目]}}
~
\hyperref[sec:6RiKDzQP6RM]{\small{> > > NEXT SERMON > > >}}
\newline
\newline
使徒行傳 26:19-32
\newline
\begin{longtable}{cl}
\hline
\hline
章節 & 經文 (和合本修訂版)\\
\hline
26:19 & \begin{tabularx}{0.7\textwidth}{X} 「因此，亞基帕王啊！我沒有違背那從天上來的異象； \end{tabularx} \\ \\ \relax
26:20 & \begin{tabularx}{0.7\textwidth}{X} 我先在大馬士革，後在耶路撒冷和猶太全地，以及外邦，勸勉他們應當悔改歸向神，行事與悔改的心相稱。 \end{tabularx} \\ \\ \relax
26:21 & \begin{tabularx}{0.7\textwidth}{X} 為這緣故，猶太人在聖殿裡拿住我，想要殺我。 \end{tabularx} \\ \\ \relax
26:22 & \begin{tabularx}{0.7\textwidth}{X} 然而，我蒙神的幫助，直到今日還站立得穩，向尊貴的和卑微的作見證。我所講的，並不外乎眾先知和摩西所說將來必成的事， \end{tabularx} \\ \\ \relax
26:23 & \begin{tabularx}{0.7\textwidth}{X} 就是基督必須受害，並且首先從死人中復活，把亮光傳給猶太人和外邦人。」 \end{tabularx} \\ \\ \relax
26:24 & \begin{tabularx}{0.7\textwidth}{X} 保羅這樣申訴時，非斯都大聲說：「保羅，你瘋了！你的學問太大，反使你瘋了！」 \end{tabularx} \\ \\ \relax
26:25 & \begin{tabularx}{0.7\textwidth}{X} 保羅說：「非斯都大人，我不是瘋了，我說的乃是真實和清醒的話。 \end{tabularx} \\ \\ \relax
26:26 & \begin{tabularx}{0.7\textwidth}{X} 王也知道這些事，所以對王大膽直言，我深信這些事沒有一件能向王隱瞞的，因為都不是在背地裡做的。 \end{tabularx} \\ \\ \relax
26:27 & \begin{tabularx}{0.7\textwidth}{X} 亞基帕王啊，你信先知嗎？我知道你是信的。」 \end{tabularx} \\ \\ \relax
26:28 & \begin{tabularx}{0.7\textwidth}{X} 亞基帕對保羅說：「你想稍微勸一勸就能說服我作基督徒了嗎？」 \end{tabularx} \\ \\ \relax
26:29 & \begin{tabularx}{0.7\textwidth}{X} 保羅說：「無論少勸還是多勸，我向神所求的，不但你一個人，就是今天所有聽我說話的人都要像我一樣，只是不要有這些鎖鏈。」 \end{tabularx} \\ \\ \relax
26:30 & \begin{tabularx}{0.7\textwidth}{X} 於是，王和總督以及百妮基跟同坐的人都站起來， \end{tabularx} \\ \\ \relax
26:31 & \begin{tabularx}{0.7\textwidth}{X} 退到裡面，彼此談論說：「這個人並沒有犯甚麼該死該監禁的罪。」 \end{tabularx} \\ \\ \relax
26:32 & \begin{tabularx}{0.7\textwidth}{X} 亞基帕對非斯都說：「這人若沒有向凱撒上訴，早就被釋放了。」 \end{tabularx} \\ \\
[1ex]
\hline
\hline
\end{longtable}
$^{1}$張國柱議員.
究竟孫教士頭上是否有光環呢?.
這個我不太肯定.
不過在華人社會或華人教會群體中.
對孫教士這個名字似乎把身份抬高了一點.
所以我當時在警察局服事時.
曾經幫一個信徒孫教士準備出去工場.
在職前有些培訓.
他跟我說自從他跟教會領袖說.
我要申請警察會成為孫教士.
那一刻之後.
教會領袖突然在他的稱呼中加了一句.
「xxx孫教士」.
而且還公公敬敬地說「xxx孫教士」.
我不知道這是個別教會狀況.
還是我們普遍華人教會的狀況.
似乎對一些人願意奉獻.
做警察和孫教士工作的人.
總覺得他們是高一點的.
不過亦可能因為這樣.
把某些孫教士的身份或名銜.
抬高了一點的時候.
變成了一種不明文的規定.
我們覺得我不是孫教士.
我不用參與孫教工作.
或者變得有些卻步.
即是沒有我的份了.
我便遠離了.
好像有兩極的看法.
一種是這些人要受尊崇.
我們這些普通人.
其實百分之九十多都是這些人.
普通的信徒.
是做不到甚麼孫教工作的.
但究竟是否真的這樣呢?.
當然我們可以說.
孫教工作.
暫時來說.
都是以孫教工作或孫教士的身份去做.
換句話說.

$^{41}$孫教士某程度上.
都是一些有血有肉的人.
其實和我和你都一樣.
只不過孫教士為何要做孫教士呢?.
其實就是他帶著一個呼召.
去回應這個呼召.
所以成為孫教士.
其實只是多了一點點的東西.
所以今天我們從這個角度.
看看我們如何看待孫教士.
在聖經裡面.
我們知道保羅是一個孫教士.
剛才我們所讀的.
《士徒人傳》的二十六章這段經文.
其實是一個公開場合.
保羅在這個公開場合裡面.
面對很多人.
這裡提到一些名字.
有阿基帕王,菲斯都,伯尼基.
還有天夫長,城裡的顯耀等等.
總督等等.
當時在座的人.
是很多不同的人.
但都是一些官員或士兵.
那邊的人比較多.
其實保羅說見證的內容.
在《士徒人傳》出現了很多次.
這裡不是第一次.
在第九章,二十六章是第三次.
所以其實同樣的見證.
已經在《士徒人傳》出現了三次.
留意在逢生出現三次.
即是重覆,重覆,又重覆.
都是一些重要的東西.
當我們看保羅見證的內涵的時候.
我們看到保羅說什麼呢?.
「我先在大馬士革」.
「後在耶路撒冷和猶太傳遞」.
「以及外邦」.
「勸勉他們應當悔改歸向神」.

$^{81}$「行事與悔改的心相稱」.
保羅說他在不同的地方.
就是大馬士革.
即是他蒙召的地方.
再去到耶路撒冷,猶太傳遞.
去到外邦不同的地區裡面.
作見證.
向他們傳福音.
希望聽到的人.
願意悔改歸向神.
其實當他這樣做的時候.
由大馬士革出發.
去到外邦這樣的行程的時候.
其實他是在回應.
在大馬士革.
主耶穌所給他的呼召.
當我們對照.
《史壇傳九章》十五至十六節.
當時保羅在大馬士革.
路上蒙召.
他看不到東西.
去到大馬士革城裡面.
有一個門徒叫做亞蘭尼亞.
在未見保羅之前.
主對亞蘭尼亞說.
「你只管去」.
「他」是指保羅.
「是我所揀選的器皿」.
「要在外邦人,君王和以色列人面前」.
「宣揚我的名」.
「我又要指示他」.
「為我的名必須受許多的苦難」.
這個是最初最初.
保羅去領受主耶穌.
藉著亞蘭尼亞所帶給他的呼召.
這個就是他的一生.
他要準備在外邦人,君王和以色列人面前.
去作見證.
而往後.
當保羅被差派出去的時候.

$^{121}$其實他先後在這三個群體作見證.
第一個.
他在外邦人中作見證.
保羅在三次宣教行程中.
接觸了很多不同的外邦人.
從他在行程中從安提拉出發.
到現在的土耳其.
再去到希臘.
再去到馬其頓.
再去到希臘.
在這個範疇裡面.
其實他接觸很多不同的城市裡面.
都是有外邦人的.
所以保羅確實在這些城市裡面.
路斯德,哥林多,菲拉比,提特羅尼加等地方.
向外邦人傳了福音.
他做了.
同時.
這些城市裡面有猶太的會堂.
所以保羅說.
去到這些城市的時候.
他先進入會堂.
向猶太人傳講.
然後當猶太人拒絕的時候.
他才出去向外邦人傳福音.
所以其實在保羅的三次宣教行程裡面.
都有向猶太人傳福音.
所以第二方面又成就了.
來到君王.
其實.
史唐人傳26章.
這個時刻.
他就是站在君王的面前受審.
所以這個場景正正就是.
回應第三方面的要求.
三樣東西他都做了.
雖然他這個處境是受審的處境.
在凱撒尼亞受.
阿基帕王.
菲斯多大人去審問.

$^{161}$好像在囚犯的身份.
來作見證.
不過.
縱然是以囚犯身份作見證.
他確實是向君王.
亦都正正綜合了這三樣東西的時候.
保羅是正在回應.
主耶穌在大馬士革.
所給他的呼召.
我們再看保羅究竟講什麼內容呢.
保羅所講的福音.
其實正正就是.
福音的核心.
因為他自己說.
我所講的並不外乎眾先知.
和摩西所說.
將來必成之事.
就是基督必須受害.
並且首先從死人中復活.
把亮光傳給猶太人和外邦人.
保羅有講舊約.
但他又講基督.
保羅所講的就是現在.
他所傳的這個基督.
耶穌就是基督.
而耶穌正正就是.
回應了舊約.
眾先知和摩西所講.
預言的那位基督.
所以無論對猶太人來說.
對外邦人來說.
他都是在講福音的核心.
他沒有離開了福音的核心.
告訴猶太人.
你所等候的基督.
就是耶穌,他已經來了.
所以從整個保羅的宣教行程中.
我們看到.
保羅確實沒有違背宣教的呼召.
更加大膽地說.

$^{201}$保羅以這個為榮.
因為神呼召了他.
他完成了.
他很開心.
能夠忠言.
受監禁.
囚犯的身份.
但他完成了神的呼召.
所以我們看到.
什麼叫宣教士呢?.
宣教士其實就是回應了呼召.
帶著呼召.
來完成使命.
我們現在的確回應了宣教呼召.
說起我最早的宣教呼召.
是怎樣來的呢?.
就是在我讀大學的時候.
三年級.
在一個全福音訓練營的晚會裡.
那個聚會的主持人.
將一個很大的世界地圖.
放在眾人的眼前.
然後就說出.
如果要將福音傳遍世界的每一個角落.
需要很多人.
那個時候他就問.
有什麼人願意回應呢?.
那時候我就很清楚.
在神面前做了一個禱告.
舉手.
站出來.
禱告.
神啊.
如果你願意.
你若猜我去宣教的話.
我願意去回應.
其實當時只是一個回應.
我不知道去哪裡.
不過我相信從那一刻開始.
我知道有一天.

$^{241}$我會離開香港去做一個宣教工作.
所以這一刻.
這個呼召.
很清楚仍然在心裡.
以至往後在多年的宣教行程裡.
常常去到起點.
是這樣產生宣教呼召的.
是從神而來的.
往後尋索宣教路上的時候.
經過不同的短宣.
去過台灣,蒙古,西伯利亞,日本等地.
看到在眾多群體中.
有一個群體叫做海外華人.
其實海外華人散佈在不同的地方.
而當這群人離開了中國的境界.
其實這群人的心是打開的.
無論是年輕人,農夫或者是做生意的.
其實都是願意去聽福音的.
所以那時候就產生了一個想法.
神會不會呼召我去服侍這些海外華人.
所以從世界各地縮窄到海外華人.
然後呢.
之後就到了.
2002年.
神帶我們去到一個探索工場.
我們去了南非探索了兩個星期.
然後回來後就申請了警察會.
最後2004年就被警察控告.
去到這個工場.
在這個過程中.
其實神在我們家庭中的說話.
也跟太太分享這個意象的時候.
神也有認同.
於是我們一家五口就出發了.
所以回想這17年.
在南非的事奉.
其實我正正就是可以說.
跟寶蘭這一段經文很互相呼應.
因為當時我的差遣禮.
就是說我是回應了天神的呼召.

$^{281}$我相信其他宣教士也一樣.
他就是帶著一個呼召.
清楚神的呼召.
有時候群體未必馬上很清楚.
是要漸漸漸漸清楚.
但是他是在回應的過程中.
一直去尋索.
以至他更加堅定.
因為上帝的呼召不會左而右擺的.
只是我們肯定地抓緊.
去回應就行了.
不過在這裡.
我們發覺回應搭上的時候.
其實是需要一些東西放下的.
一些東西要捨棄.
捨棄什麼呢?.
不同的人捨棄的東西是不同的.
對我來說.
2004年當時.
出發前我們是住在寶達村的.
20年前.
秀博屏,寶達村,在山上.
當時是新樓.
我們有5個人.
加上外母,6個人.
所以就配了一個比較大的單位.
那個大單位也挺漂亮的.
三邊都是窗.
風涼水冷.
而且夠高.
我們能夠看到那個維多利亞港.
放煙花的時候可以看到.
當時前面的樓還沒有建成.
但我們知道要去宣教.
所以在那個時候.
當我們一直尋索的時候就發現.
上帝一直在安排.
在2004年就知道.
外母因為身體不好.
她有癌病的復發.

$^{321}$然後2004年的8月.
6月還是8月.
回到天家.
所以剩下我們5個人.
我們5個人離開這個單位的時候.
其實就沒有人了.
可能有個建議.
其實你不要浪費這個單位.
這麼漂亮.
你反正有時候回來述職.
也需要地方住.
你就留住它.
或者找人幫你交租.
留住它.
你回來也有用.
但我心裡面知道.
這一次的出發是真的出去.
我們不知道什麼時候回來.
所以有時候說不定.
要等多久呢.
而且租給別人住好像不太好.
不太合法.
所以裡面有掙扎.
但這個地方又很漂亮.
又捨不得.
在裡面其實中間有很多不捨.
不願意捨去.
又好像家人的意見很好.
不過真的最後在眾多掙扎中.
當我們最後決定.
把鎖匙交給房署的時候.
我們就決定離開這個單位.
然後在交鎖匙的那一刻.
簽名正式退屋的時候.
其實內心就好像一個很大的石頭.
終於放下.
能夠輕散的上路.
這個輕散的上路.
代表著真的憑著信心去踏上.
因為不知道回來的時候.

$^{361}$會住在哪裡.
但深信上帝會預備.
有東西要放下的.
我不知道你要踏上宣教路的時候.
會放下什麼.
但正如保羅所說.
去回應這個呼召.
是一件有意義的事的時候.
是值得放下的.
第二方面.
保羅除了回應天上的呼召之後.
聽的人.
其實不是每個人都明白的.
因為當時.
在聽到的人當中.
有一個叫菲斯都大人.
菲斯都大人就有回應了.
菲斯都大人說.
保羅你風涼.
你的學問太大.
反使你風涼.
當我們說一個人風涼的時候.
其實有兩個意思的.
一個是真的精神病.
精神錯亂.
不知道自己在說什麼.
總之就是精神病那一類.
但另一個風涼也可以說.
這個人做的事.
似乎不是一個正常人做的事.
或者我們俗語所說.
你神經病的.
照他而已.
不是真的說這個人.
與無人似的那種風涼.
我想菲斯都大人.
回應的時候.
看到在他眼中.
保羅這個人.
其實很有學問.

$^{401}$他是法利賽人.
由師父教的.
很有學問.
但這一刻他所見到的保羅.
他是以一個囚犯的身份.
站在這個台前.
要被人審問.
菲斯都大人就覺得.
你不值得.
你這麼厲害.
做了很多偉大的事.
為什麼現在輪到囚犯戴著鎖鏈.
去被人審問.
你是不是瘋了.
保羅很清楚回應.
他說.
菲斯都大人.
我不是風涼.
我說的是真實和清醒的話.
保羅回應菲斯都大人.
他說.
我能夠很清楚在前面的時候.
將我見證.
由大馬士革.
一直走到今天.
這段見證.
很仔細地.
告訴你.
一個人能夠這麼清楚.
詳細地.
將他的過程.
告訴別人的時候.
他是一個清醒的人.
他不是瘋了.
所以保羅不是一個.
瘋了的人.
他竟然在菲斯都大人眼前.
覺得他好像瘋了.
今天.
孫格士的工作.

$^{441}$有時可能在我們眼中.
會覺得.
不太正常.
這個意思.
不是說孫格士說得不清楚.
他在做什麼.
因為孫格士會述職.
每三年回來.
就會將他的過去的事情.
有個回報.
講得很清楚.
只不過可能聽的人.
會覺得.
為什麼你還要去呢.
述職完.
繼續去這些工場.
繼續去服侍呢.
有時這些工場帶來很多挑戰.
譬如選舊工場.
往往帶著的就是.
有簽證的問題.
有文化問題.
有時氣候的適應.
或者在孫格工場裡面.
遇到很多難阻.
有時有反對者.
異教徒的反對.
甚至有異端.
都是這些難阻.
甚至有時孫格士在孫格工場.
有時會遇到突發事件.
有時會有車禍.
內戰這些危險限制.
聽的人會覺得.
這麼多危險.
這麼多限制.
為什麼這個人還是繼續去呢.
好像傻傻的.
甚至有時基督徒.
信徒或家人.

$^{481}$有時都不明白.
為什麼孫格士要這樣去.
冒著危險繼續去.
我們就去南非.
南非就是非洲最南的位置.
我們就去紐卡索和伊爾沙伯港.
兩個地方服侍.
先後紐卡索就做了六年.
然後伊爾沙伯港就做了十一年.
總共十七年.
我記得最早2004年.
當我們被去警察派的時候.
我們家人就跟我們說.
你們夫婦去冒險.
已經被警察派了.
我們就去.
冒險.
已經不容易了.
為什麼還要帶著.
三個這麼小的小孩.
去南非冒險呢.
2004年.
南非的情況是怎樣呢.
香港人都知道.
南非治安不好.
打劫很多情況出現.
其實20年後的今天.
南非還是一樣這麼差.
沒什麼改善.
縱然政府很多努力.
但為什麼要這麼危險.
去這些地方傳福音呢.
其實他們不太知道.
只知道為什麼要去.
這麼危險的地方工作.
沒錯.
有時候真的.
不容易講解清楚.
為什麼要去.
不是說.

$^{521}$我們不知道.
在那裡做什麼.
而是聽的人有時候.
不太理解不太明白.
為什麼要去這麼危險的地方.
為什麼要繼續去做.
所以保羅所講的.
或者被人感覺的.
風潦其實就是指.
這班人宣教士.
好像有些瘋狂的感覺.
不是正常人做的事.
接著.
保羅去和.
非洲的人講完之後.
他就轉向.
和阿基帕王講.
然後.
和阿基帕王講的時候.
因為他知道.
阿基帕王這個人.
他熟悉猶太的事件.
所以他講了很多猶太的背景.
他相信阿基帕王聽得明白.
然後當他講.
阿基帕王的.
這些事件的時候.
他提到過去有提到先知.
有提到摩西.
然後他馬上轉移.
阿基帕王說你信先知嗎.
我知道你是信的.
阿基帕王也很聰明.
知道他轉移話題.
你講先知.
你又講信.
那是不是想叫我做基督徒呢.
於是阿基帕王馬上回應.
你想稍微勸一勸.
就能說服我做基督徒嗎.

$^{561}$阿基帕王很醒目的.
聽得明白,不過保羅的回應更加精彩.
保羅說.
無論是小勸還是多勸.
信神所求的不但你一個人.
就是今天所有.
聽我說話的人.
都要像我一樣.
只是不要有這些所念.
保羅沒有錯.
回答阿基帕王這個問題.
好像是回答一個人的問題.
不過從保羅所講的內容.
他知道他心裡所講的.
或者當時.
他辯護的場合.
其實他眼中看到的.
就是我要告訴所有人.
所以不止阿基帕王.
或者非洲大人.
一個人去聽.
其實保羅是想在座的.
每一個人,總督,官長.
士兵,所有的人.
都要聽到保羅這個見證.
然後去做基督徒.
所以他說不是一個人.
其實我想在座每一個人.
不過希望你們做基督徒的時候.
不要有這些所念.
自由的.
來做基督徒.
所以保羅看到在這裡.
看到保羅在這個場合.
縱然在受審的場合.
不過他沒有遺.
沒有失去.
宣教士的特質.
有機會他就傳福音了.
正如.

$^{601}$保羅在那麼多前書所講的.
就是對猶太人.
我就作猶太人,為要贏得猶太人.
對律法以下的人.
我需不在律法以下.
還是作律法以下的人.
為要贏得律法以下的人.
無論如何,我總要救一些人.
這個就是宣教士的心態.
就是一有機會.
無論是什麼場合都好.
有機會做見證的時候.
他都願意去做的.
亦都回應著.
提摩太后書四章二字.
保羅所講的,他說務要全度.
無論得時不得時.
都要專心.
並以百般的忍耐.
和各樣的教導責備人.
警戒人,勸勉人.
保羅的做法是一貫的.
就是有機會的時候.
無論是什麼人,得與不得.
都要向人分享福音.
這個是宣教士的心態.
是他的目標.
他都有人找機會跟他說.
所以你看到他真的去三次傳道的時候.
猶太人他都說,愛邦人他都說.
是一個這樣的心靈特質.
所以在這個場合裡面.
正正流露到保羅的.
事奉的心靈.
我們在南非要尋找.
這三十萬的華人.
真的不容易.
因為南非很大.
地圖很大.
這三十萬人分散在不同的地方.

$^{641}$所以要去找.
有不同的機會.
中國人喜歡做生意.
所以我們就嘗試去商場找.
就知道.
中國人開一間店舖很有趣.
因為店舖容易找.
他就把衣服掛出來.
單車就放在前面.
雖然名字是英文字.
但你猜到這樣的設計.
大概是中國人開的店舖.
所以就去商場找.
這種店舖.
有時是批發有時是零售.
南非有很多商場.
我們一個城市裡面.
其實有不同的區.
每個區總有一些商場.
十幾間店舖都好.
我當時的習慣就是.
有機會就開車去.
看一下一個商場.
其實我不是去買東西.
我是去找中國人.
當然有時找到有時找不到.
所以其實就養成一個習慣.
不要緊.
有或沒有都去逛逛.
看看.
不過回到香港之後.
這個習慣好像沒有失去.
總是有習慣想去逛逛商場.
不過這次逛不完.
香港太多商場了.
但這次也不需要找中國人.
因為香港已經有這麼多教會.
處境不同了.
但在香港這幾年裡.
其實都是一樣.

$^{681}$一有機會.
作為宗教士.
都是很想向人傳福音.
但亦都感恩神真的預備了不同的機會.
我講其中一個例子.
見證.
我們住的房屋.
是租的.
其實通常我們租屋.
交錢給房東就搞定.
不用理那麼多.
他又不需要理我那麼多.
我又不需要理他那麼多.
總之大家沒有拖沒有欠就搞定.
不過有一次.
這個房東.
當然我們用微信.
微信裡面有時他知道我的名字.
我就用牧師的名字.
他知道我是牧師.
不要緊.
牧師在香港都可以.
他就問我.
牧師想問你一件事.
你的教會在哪裡.
當時我不太知道.
他的信仰背景.
不知道他是否信徒.
所以一個人問牧師在哪裡.
事奉.
他對信仰有興趣.
我當然跟他說解.
我是宗教士沒有固定教會聚會.
如果你想參加教會.
我可以跟你介紹.
在附近的區或你那區.
看看我認不認識.
後來他講了一個區.
我就說.
有兩間教會.

$^{721}$A和B.
後來他選擇了A.
我陪你去A教會.
約定時間.
去A教會崇拜.
其實我不是純粹去這間教會.
終於見到他.
他肯去.
入去之後覺得整個崇拜很好.
然後他說.
這個教會很好,我會繼續.
在路途裡面.
他講起為何他會找教會.
原來他在當時之前.
早一個月.
他做了一個心臟手術.
這個心臟手術裡面.
其實他有些擔心.
所以有一個基督徒朋友.
發了一個訊息給他.
他說你按照這個.
圖文不斷讀.
你便可以有平安.
這個就是基督徒.
介紹他向耶穌祈禱.
他不停地祈禱.
祈禱後他很平安.
很順利地離開.
離開後他覺得很好.
他便回教會.
於是他便穩穩地.
通過我去了這間教會.
後來也隔了一段日子.
現在也有聯絡.
他跟我說偶然他回到香港.
也會繼續回這間教會.
講起教會也好.
有幫他做一些栽培.
問他一些信仰問題.
我經常是不自覺的.

$^{761}$不過既然機會在眼前.
我們就不要錯過.
哪怕是一個問題.
你不覺意.
你不為意.
但少少為意.
多想一想.
可能是一個信仰機會.
今天我們知道.
突發事件是隨時會發生的.
世界不是越來越平靜.
其實是越來越亂.
災禍.
戰爭.
俄羅斯8.8級地震.
到處的戰亂.
其實是在傳送訊息給我們.
是一個提醒.
一方面說主耶穌快來.
另一方面.
更加提醒我們.
要抓住機會.
我不知道這些在災難中的人.
有多少機會.
還可以聽到福音.
今天當我們說.
我們知道信耶穌在哪裡.
我們知道信耶穌永生.
是一個把握的時候.
其實世界上有很多人.
是未聽到的.
這些事情.
正正就是不斷提醒.
作為神的兒女.
我們需要警醒.
世界上很多人是在水深火熱中.
很多人是未聽過福音.
在南非這30萬華人裡.
有很多我去接觸的人.
都未聽過福音.

$^{801}$未聽過耶穌.
他們只知道一個價值.
所以其實是一個很多的需要.
究竟.
誰去向他們.
傳達和平的福音呢.
越亂.
越需要和平的福音.
越亂.
越需要和平的福音.
那時我們去黑人區.
傳福音的時候.
開車去到一個區.
找一個店舖.
店舖的華人.
就招呼.
我們走的時候.
當時上了車.
不過沒有鎖上車門.
誰知.
那時有幾個黑人圍住.
推開門拿著槍.
準備上去打劫.
那時我們就和他們合作.
任你拿什麼.
你拿到就拿.
你只可以這樣做.
你還可以做什麼呢.
和他反抗嗎.
你怎知道那支槍是真是假.
當然當他是真的.
但在掙扎的過程中.
很奇怪.
搶來搶去.
搶不到東西.
而且掙扎了幾分鐘.
他知道不對路.
就只好離開.
結果拿了一大袋東西.
那一大袋是福音單張.

$^{841}$很可惜.
他又不懂中文.
拿單張沒什麼用.
不過這件事.
回來有個回響.
原來去黑人區那麼危險.
沒試過是第一次.
所以心裡就想.
可不可以以後不要去.
這些黑人區傳福音那麼危險.
不過師母在這個事件之後.
心裡有一個聲音提醒他.
他說如果你不去黑人區.
沒有人會再去.
有些地方.
當地只有我們.
我們是傳道人.
願意去這些地方.
傳福音.
一般的華人是不會去的.
回應天上的呼召.
似乎是一個很偉大的題目.
正如我們覺得宣教士.
是很偉大的人.
但我們一般信徒.
做不到甚麼.
不過我想有兩個方面的回應.
第一個.
對於那些曾經回應.
宣教呼召的弟兄姊妹.
我不知道你是甚麼場合.
粉兄會,盤鈴會或營會.
你曾經.
有感動.
填表,剔過或祈禱.
其實.
這個立志的回應.
是認真的.
上帝是看到的.
所以我們鼓勵.

$^{881}$這些弟兄姊妹.
繼續尋索.
在尋索的過程中.
需要行動,需要短訓.
需要問牧者意見.
需要猜拳聚會.
認真追尋.
確定是否真是神呼召你.
如果是神呼召你.
你就憑著信心去踏上.
這是一條無不復的路.
第二方面.
對於每一個弟兄姊妹.
我們可以學習.
保羅宣教的生命.
他得時不得時.
都要傳道.
我們從保羅宣教.
生命的特質學習.
究竟我們日常生活中.
有甚麼機會.
可以為耶穌作見證.
為耶穌多得一個人.
哪怕一句說話.
一個問候,一個回應.
我不知道你來教會聚會時.
路上會否遇到街坊.
街坊問你去哪裡.
去喝茶嗎?.
我不是.
我去教會.
你嘗試這樣回應.
可能對方很奇怪.
你去教會?我不知道.
教會是甚麼?.
可能從此就有機會.
你去介紹他.
教會是做甚麼?.
一起見拜這位主.
所以今天我們有很多機會.

$^{921}$在我們社區.
即使退休.
在這個社群.
有位長者帶著孫子來教會.
你也可以帶人來教會.
所以不要放過任何機會.
今天機會在眼前.
在乎我們如何把握.
求主幫助我們.
縱然好像很偉大的題目.
但其實和我們每一個人有份.
讓我們願意學習.
保羅的榜樣.
回應主耶穌給我們的呼召.
我們一起禱告.
主耶穌說.
撒帖伯利恭敬.
謙卑在你面前.
我們願意去聆聽你的聲音.
主耶穌有時我們聽過.
關於你的呼召的聲音.
但有時我們很懼怕回應.
有時我們懼怕放下.
我們很懼怕跟隨.
但我們知道.
從保羅身上我們看到.
他真是一生沒有後悔.
知道他回應了天上的呼召.
從起初夢召那一刻.
我們知道.
他知道要怎樣回應.
上帝要他走的路.
以致他在這一刻.
在《誓言傳》26章記載的時候.
他真的完成了.
以致他很開心.
很滿意地告訴人.
他完成了呼召.
縱然處於受審的處境.
所以求主也幫助我們.

$^{961}$藉著保羅的生命.
他如何回應審問的過程.
他也有幫助我們.
我們也有幫助.
在審問的過程中.
他仍然把握機會去見證你.
希望在座的每一個人.
都能夠聽到福音.
求主幫助我們.
讓我們能夠找到機會.
我們身邊可能有很多人在我們身邊測試過.
讓我們能夠學習敏銳.
如何去祝福他們.
讓他們能夠認識到耶穌.
無論來教會也好.
講福音也好,講見證也好.
都能夠讓他們聽到耶穌是誰.
求主使用我們教會.
願主幫助奉耶穌基督的使命..
阿們..
\newpage



\section{加拉太書 5:13-16}
\label{sec:6RiKDzQP6RM}
\textbf{2025.08.10 《跟隨,跟隨》 經文:加拉太書5\_13-16 曾裔貴牧師}
\newline
\newline
連結: \href{https://youtube.com/watch?v=6RiKDzQP6RM}{\texttt{https://youtube.com/watch?v=6RiKDzQP6RM}} ~~~~ 語音日期: 2025-08-13
\newline
\newline
\hyperref[sec:YkzDcCtXDZ8]{\small{< < < PREV SERMON < < <}}
~
\hyperref[sec:index_chronic]{\small{[返順時目]}}
~
\hyperref[sec:index_scriptual]{\small{[返順卷目]}}
~
\hyperref[sec:OFK7CEvBZU8]{\small{> > > NEXT SERMON > > >}}
\newline
\newline
加拉太書 5:13-16
\newline
\begin{longtable}{cl}
\hline
\hline
章節 & 經文 (和合本修訂版)\\
\hline
5:13 & \begin{tabularx}{0.7\textwidth}{X} 弟兄們，你們蒙召是要得自由；只是不可把這自由當作放縱情慾的機會，總要用愛心互相服侍。 \end{tabularx} \\ \\ \relax
5:14 & \begin{tabularx}{0.7\textwidth}{X} 因為全部律法都包括在「愛鄰如己」這一句話之內了。 \end{tabularx} \\ \\ \relax
5:15 & \begin{tabularx}{0.7\textwidth}{X} 你們要謹慎，你們若相咬相吞，恐怕要彼此消滅了。 \end{tabularx} \\ \\ \relax
5:16 & \begin{tabularx}{0.7\textwidth}{X} 我說，你們要順著聖靈而行，絕不可滿足肉體的情慾。 \end{tabularx} \\ \\
[1ex]
\hline
\hline
\end{longtable}
$^{1}$各位弟兄姊妹早晨.
我們一起同心祈禱.
感恩天父讓我們在這段時間.
可以聆聽你的真道.
多謝你讓我們過去十天.
有港九培靈賢經會.
有幾位的牧者.
他們用心的教導.
還有賢經 港道 還有粉輕會.
我們真的很盼望.
很多聽完的弟兄姊妹.
心靈都受感動.
也明白真理.
也清楚這條屬靈的道路.
是要擺上服侍.
今早我們聚集在你的面前.
我們很盼望.
我們能夠明白.
使徒保羅給我們的教導.
在當時對著加拉太教會的場景.
面對異端的危險.
要跟隨的到底是哪一個呢.
幫助我們能夠從經文去明白.
多謝你聽感恩禱告.
奉耶穌基督的名.
阿們.
上一個月都說了.
哥林多前書的十二十三章.
其中相當重要.
就是愛比一切的恩賜.
都來得重要.
是使徒看到哥林多教會.
正正面對著很大的危機.
就是爭誰的恩賜大一點.
厲害一點.
這樣的心態.
在教會裡面.
就會變成分黨分派分裂.
所以十三章的哥林多前書.
是一個很重要的提點.

$^{41}$今天我們透過這一段經文.
是的.
是這個嗎?.
我開了.
我再試一試.
不知道是不是不行呢?.
OK,現在可以了.
謝謝.
用一個例子開始我們今天的講道.
盡量不要讓人看得太清楚.
大家在七月五日.
無線的真理的猜情尋.
尋找真相.
石安教會梁逸華.
這一位的領袖.
提倡雙氧水可以治病.
能醫百病.
其中更嚴重的就是.
他們的末日觀.
說主耶穌回來了.
他們要避難.
要去到很遠的地方.
要買他們的求生包.
另外呢.
原來在教義裡面有一樣東西.
是很貫徹的.
他說教會的領袖就是神的受高者.
教徒只能夠服從.
不可以拒絕.
教會的領袖說什麼.
所有信徒都要跟隨.
後來大家看到那個新聞.
你們會看到有些前教徒.
就現身說法.
說到當時他們怎樣去面對教主的.
好像洗腦式的教導.
在崇拜裡面.
一個大型的聚會.
只是將一些經文.
很片段性.

$^{81}$很扭曲地去教導信徒.
所以信徒所收到的.
得到的資訊.
就是來自教主他自己個人的.
所謂的威望.
所謂的解經.
其實已經是當時.
基督教已經定性了他們是異端.
不過呢.
藉著這個.
大家剛剛都有一個新印象.
可能當時呢.
這件事在香港引起轟動.
很多在座的都未信耶穌.
你可以想像得到.
當時帶來的都是非常大的影響.
但是當今天救市重提.
就讓我們能夠透過這一位的梁日華牧師.
他的這樣的言論.
這樣的主張.
其實就算很多人.
他帶領那幫人.
如果是一直跟著他.
只會走向真的滅亡.
因為給不到人永恆的救恩.
不是因為真理.
使得人得到自由.
所以這是一個很重要.
到底我們在人生.
我們跟隨哪一位的領袖呢.
這是一個警惕的例子.
有另一個例子呢.
都很想和大家來看看.
這個是最近在美國.
多麼嚴重的一個案件.
當事人呢.
最下面的這位黑人.
他是一個很出名的.
在美國的一個搖攝的歌手.
很出名的.

$^{121}$在美國的超級盃的超級盃.
可以來做獻唱.
他是非常有錢的.
這位稱為P. Diddy的這位人士.
他的成名.
他的重要是什麼呢.
他每年會有連續幾天的狂歡.
邀請明星.
出名的歌星.
政界的名人.
去他的大宅.
來去狂歡.
是連續可以是幾天的.
那些狂歡.
後來被揭發被告.
因為他偷運一些的.
未成年的男女.
來供這些名人去玩樂.
在他家裡搜查到.
超過一千瓶的BB的潤滑油.
和一些可以有毒品.
有其他的東西.
是讓這個人收買不同的人.
大家看到.
那個是出名的明星.
迪卡比奧.
記得那套電影的男主角.
很多的名人.
包括特朗普總統.
曾經都參加過.
這樣的狂歡聚會.
為什麼這麼出名呢.
因為有不同的人.
開始去控告這個P. Diddy.
現在他仍然在紐約坐牢.
那個案件就快要宣判.
但是兩件嚴重的指控.
被宣告為無罪.
只是很輕微的.
一些的控告是定罪.

$^{161}$但是現在據聞.
特朗普總統.
是打算pardon他.
用總統的權力.
去特赦這位的人.
你可以難以想像.
在美國.
這些公開的事情.
是由不同的州.
騙取和用很高的價錢.
叫這些未成年的男女.
和一些男妓.
來到他的大宅.
那個叫做White Party.
白色聚會.
今晚你上網搜尋一下.
P. Diddy這個案件.
香港原來也有大幅的報告.
可以想像到.
這樣的道德淪亡的罪案.
可以說是行使總統的權.
考慮特赦這個名人.
你可以想像到.
我們跟的是什麼的領袖.
帶來什麼的後果.
美國現在已經是.
道德淪亡的一個下坡.
可以是這樣的事情.
完全不覺得有什麼問題.
是亂童.
是這樣去濫交.
是這樣去放縱.
吸毒.
完全是肆無忌憚地這樣做.
可以是公然的.
完全不會介意.
因為他的律師團隊.
為他來辯護.
但是證據確鑿了.
都可以無罪地去準備特赦.

$^{201}$各位朋友.
各位弟兄姊妹.
跟誰的領袖很重要.
我們人生的方向.
是跟誰呢.
進入我們的分享和解經.
我們就要留意這裡.
剛才13節已經說了.
這是第五章的一個很重要的總結句.
因為五章的第一節已經說了.
基督釋放了我們.
叫我們得以自由.
所以我們不要再被律法的欺騙去挾制著我們.
現在13節再繼續這樣說.
告訴我們.
其實我們成為基督徒.
我們有什麼的自由呢.
不過保羅說了.
得到生命的釋放的自由之後.
成為神的兒女之後.
他就要告誡加拉太教會.
在教會的生活.
漫長的人生.
彼此一起同心.
在教會服事的時候.
就要有一個很大的留心.
就是這裡告訴我們.
只是不可以將這個自由.
就是這個身份.
我們這個現在在神耶穌的救恩裡面的狀態.
我們就以為我們可以做什麼都可以.
因為不再受猶太教.
摩西的律法的枷鎖.
那我們是不是做什麼都可以呢.
保羅說不是.
似乎到底是自由還是不自由.
到底是釋放還是再一次被捆綁.
所以裡面是有一些事情我們要注意的.
今天集中來看.
有幾段加拉太書的經文.

$^{241}$就可以將這個信息清楚地跟大家分享.
這裡說.
此時不可把這自由.
當作放縱情慾的機會.
總要用愛心互相的服事.
一個的反一個的要.
兩個都是命令.
就是告訴我們.
原來自由必須放在一個這樣的心靈的心態.
以愛走我們人生的路.
以愛.
這個愛當然就是耶穌基督的十價犧牲的愛.
來作為我們的價值觀去行使我們認為的自由.
這個是大前提.
我們進入下面.
我們就來看看.
這裡說.
我說你們要順著聖靈而行.
十六節.
絕對不可以滿足肉體的情慾.
就是說再一次再一次.
將情慾和聖靈做一個對立.
做一個我們要來留心的.
兩樣東西同時間在信徒的內心.
情慾.
信了耶穌.
仍然有這個救人救我.
有這個所謂的我們與生俱來.
犯罪的人的品性和價值觀.
但是信了耶穌的弟兄姊妹.
有聖靈在我們內.
所以保羅告訴我們.
聖靈和肉體是相爭的.
是彼此為敵的.
我們不可以以為自己是完全自由.
中立是沒有的.
這是值得我們留心思想的.
我們的人生原來在耶穌的救恩裡面.
我們現在是有聖靈的內著.
但是保羅一而再告訴信徒.

$^{281}$就算你們是外邦人.
你們不是猶太教的信徒.
但是我們都是有肉體的情慾在我們裡面.
未死一日.
我們的情慾都不會離開我們.
伴隨我們一生.
所以我們在經文裡面.
可以留心.
有什麼出路.
有什麼解救的辦法.
很感恩.
我們看到保羅在這個課題上.
一點都不含糊.
讓我們看到.
這裡說順著聖靈而行.
就是說原來.
用一些大家能夠理解明白的.
例如GPS.
例如高德地圖.
我有一個經驗.
早幾個月.
我很久沒有去青年會小學天水圍.
由錦繡去.
就走最外圍的公路.
我就記得那條路.
那些彎沒那麼多.
路口沒那麼多.
我就記得.
我這次從紅恩堂過去.
我就以為一樣的了.
因為我走的路是新的.
我就沒有開地圖.
那個電話的地圖.
沒有開.
繞了天水圍之後.
哎呀慘了.
繞了15分鐘.
為什麼路不一樣了.
我一直都沒有留意.
接著不行了.

$^{321}$停下車.
馬上開電話.
高德地圖.
打青年會小學.
哎呀原來我繞來繞去.
原來路口就在那裡.
很快.
其實已經去了很遠了.
因為天水圍進入了一個迷宮.
去了很遠.
去了濕地公園.
接著就停下.
馬上調校.
很快.
都不用一分鐘.
就走回那條路.
就進去了.
這是一個大家在現代.
如果你懂得開車.
你就會知道.
這麼方便的導航系統.
生命的導航.
跟我們的生命對話.
我們成為基督徒.
保羅很清楚.
我們得到了耶穌.
釋放我們的自由.
但是這個自由.
有可能會帶來.
我們自己會受傷害.
也會傷害其他人.
就在教會裡.
我們的情慾.
會帶來很多.
不必要的那種影響.
因為情慾的事.
是顯而易見的.
仇恨.
憎恨.
嫉妒.

$^{361}$各方面的事.
但是保羅讓我們知道.
如果我們在聖靈的帶領.
感動下.
聖靈就會結出.
在我們生命裡的果子.
相信大家都知道.
人愛喜樂和平.
忍耐.
恩慈良善.
信實.
溫柔.
節制.
原來我們的人生.
我們基督徒的人生.
在主耶穌的裡面.
我們可以在我們的生命裡.
不斷可以發展.
這一份這樣的聖靈果子.
所以問題就是.
我們跟隨誰呢?.
再一次我們來想一想.
原來保羅在迦拉太書.
講聖靈講了四次.
而且是人生的.
由信耶穌到終極.
整個群體要跟聖靈.
可能大家沒有機會詳細去讀.
大家留意一下.
這裡說聖靈是我們信心的印記.
聖靈就是我們的印證.
我們受聖靈.
是因為行律法.
還是因為聽信福音呢?.
當然答案是很清楚的.
迦拉太教會就是聽明白了耶穌的救恩.
他們就立即得到聖靈作為印記.
在他們的生命裡面.
所以現在保羅是質疑他們.
為什麼你們靠聖靈入門.

$^{401}$現在你被那些傳異端的人影響.
說你要守律法呢?.
說受了聖靈.
信了耶穌.
都不夠完全要行律法.
其中就是受割禮.
被這些猶太教的知識份子這樣影響.
所以原來聖靈是我們得到生命的印記.
每一個基督徒.
由你決志那一刻.
不是洗禮那一刻.
是決志那一刻.
清楚向耶穌認罪.
接受耶穌做救主.
你就立即得到聖靈.
剛剛我在圖書館訂了一本書.
政府圖書館.
那個題目叫做after.
整本書很厚的.
是很多研究的個案.
是一位精神科醫生.
看到很多臨死的病人.
他們的經歷.
他們的一些分享.
就是死過翻生.
死不去.
臨死之前.
他們遇見大光.
離開自己的身體.
用了50年的時間.
這個精神科醫生.
他說這本書做了一個總結.
他說他不是一個信徒.
但是他開始相信有靈魂.
開始相信離開肉身的生命.
是有另外一個世界.
因為不約而同.
那些有臨死經驗的人.
講的見證.
雖然不一而足描述.

$^{441}$但是都告訴我們.
聖經所講的.
所以今天我們有聖靈.
神聖的靈住在我們內心.
其實就好像生命的GPS.
帶領我們人生朝向永生的路.
原來朝向永生的路.
不單只是去天堂.
而是我們這個人.
生命的成長.
是有神的性情.
在我們裡面不斷的更新.
滋生我們.
我們不應該再是.
剛剛初信時候的舊有生命.
而是一個不斷更新的生命.
所以這是一個入門.
保羅說.
我們一信耶穌.
在救恩的宣告裡面.
我們就得到聖靈.
沒有條件.
只是因為我們信耶穌.
聖靈就會住在我們裡面.
不再離開.
永遠在我們裡面.
所以弟兄姊妹.
多些尋求裡面的教導.
多些安靜.
默想.
聖靈在我們內心裡面的感動.
有時叫你做這件事.
我又想講一個例子.
這次講那個地方.
不用隱瞞.
在感受一位很熟的弟兄.
因為他跟我說這個見證.
他都到處跟人說.
他說在共事的時候.
有一個弟兄.

$^{481}$因為一些爭執.
然後成為陌路人.
見到他在那裡.
他就在另一個門口走.
維持了一段長時間.
後來這個弟兄就輾轉移民.
然後很多年之後回來聚會.
因為這麼多年沒見.
見到這個弟兄.
很開心.
走過去跟他握手.
這個弟兄以往.
見到他就離遠走了.
沒人知道.
只有他自己知道.
但是那一刻.
聖靈感動他.
他覺得他要衝破這個記憶.
仇恨.
嫉妒.
走過去握手.
不是裝模作樣.
是真心的跟他有一個復和.
這樣來到大家寒暄.
聊天.
他說他的內心.
在那個感動裡面.
開始衝破了這個隔膜.
他以後不再有這個內心的困鎖.
兄弟姐妹.
有沒有這些事情.
有時候可能是少到不得了的事情.
但總是在我們的心靈裡面.
就會有這些芥蒂.
有沒有聖靈的感動.
我們可以釋放自己呢?.
不要囚禁自己.
很無謂的.
我們應該很坦蕩蕩.
在耶穌的救恩裡面.

$^{521}$大家是神兒女.
沒錯.
大家不同的性格.
不同的看法.
是很正常.
很應該的正常.
但是在愛裡面.
不應該有隔膜.
在主耶穌裡面.
所以聖靈的引導.
往往就是帶我們.
面對我們內心的陰暗.
往往是令我們.
可以勝過我們自己.
內心的軟弱.
這是其中一個例子.
感受的一個弟兄.
很熟的弟兄.
跟我說.
分享他的感恩.
可以衝破自己的內心.
另外.
這裡告訴我們.
不單止我們有聖靈住在我們裡面.
神的兒女的記號.
你們既是兒子.
這是第四章加勒太書.
神就差他兒子的靈.
進入你們心裡.
呼叫阿爸父神.
神兒女的身份.
是確定了.
永恆的確定了.
我們在聖靈裡面.
我們越來越清楚.
我們是神的兒女.
哪一位是神呢?.
造物主.
創造的主.
天地的主的兒女.

$^{561}$所以你的身份.
應該要表達這份尊貴.
這不關乎你的學識.
不關乎你的家庭背景.
是我們要理解.
就算我們是出生基層.
就算我們天生有很多不同的缺陷.
但是.
這是神兒女的記號.
我們是造物主的兒女.
是神在他的愛子.
犧牲裡面.
將我們買熟回來.
保不寶貴自己的生命.
這個寶貴的身份呢?.
如果有一個很有錢.
很有影響力的人.
他很愛你.
他很想收養你.
你覺得求之不得.
我經常都跟我們團契的弟兄姊妹說笑.
我八歲那年.
有一個很有錢的姓梁的夫婦.
梁太.
想收養我.
帶我去加拿大.
可惜當時我不知道.
如果知道我就立刻舉手.
因為你住在我那區.
家裡有十三兄弟姊妹.
打風呢.
爭不夠.
那些午餐肉.
轉過來就沒有了.
十三個跟你爭.
那你可以想像到.
你在那個時候.
當然想去一些很好的地方.
很好的家庭.
受很好的教育.

$^{601}$沒有.
但是神的安排.
一定是最好的.
這個是收養.
最好的人.
最有影響力的.
最愛你的.
收養你.
但是現在造物主.
在耶穌的愛裡面.
將我們拯救.
很清楚告訴我們.
我們真是神的兒女.
天使都不可以這樣稱呼自己.
天使不過是受造的靈.
他們是瀑液.
我們才是神的兒女.
所以洪恩的弟兄姊妹.
我們就是神的子女.
我們一離開這個世界.
我們的靈魂就去到神那裡.
很多瀕死的人.
臨死的人.
他說他們見到不約而同見到大光.
他們朝向那個大光.
他說他去到.
完全看到自己一生.
但不是一種痛苦.
是開心平安的.
我等一下下午就會拿這本書.
翻譯成中文.
五十年的個案.
不約而同說有靈魂.
有靈界.
有永恆的居所.
所以我們是神的子女.
這是保羅告訴我們.
所以我們今天在教會裡面.
我們就是要發展.
怎樣讓神的聖靈.

$^{641}$帶著我們的人生.
一起去成長呢?.
剛才已經讀過了.
聖靈是自由生活的引導者.
哪些是真正的自由呢?.
原來在保羅所得到的啟示裡面.
原來在聖靈帶領感動的生活.
是這樣的.
就是結出聖靈的果子.
在我們的生命裡面.
所以我們應該不斷地跟自己比較.
人愛喜樂和平.
忍耐恩慈良善.
信實溫柔節制.
在我們的生命裡面.
是怎樣來到不斷的滋長呢?.
是教會生活具體的表現.
每個弟兄姊妹.
一同地讓聖靈帶領.
我們有這樣的表現.
譬如我們團契.
知道哪一位弟兄姊妹病入院.
我們就二話不說去探訪.
去關心,去禱告.
這些就是很明顯的一個愛的表現.
是由內心來引導我們.
這樣來服侍.
是彼此用愛心的服侍.
剛才已經講了這個例子.
這一位隨著聖靈感動的弟兄.
他來改變他自己的內心.
那份沒有人知道的不開心.
那種可能是仇恨.
可能是妒忌.
也沒有人知道.
也不需要理會.
因為那個弟兄已經移民了.
但是當他回來的時候.
原來就是聖靈感動他.
你能不能夠在這一次機會.

$^{681}$你去將你自己的囚禁的自己.
可以釋放呢?.
無罪釋放了.
你就可以想像得到.
那個開心,那個快樂.
那份平安就馬上出現了.
其實根本就不是什麼大事.
只是看事物不同的角度.
就這麼一切.
所以聖靈往往在教會.
就是做這些配對工作.
解除冤仇.
釋除誤會.
大家願意在愛裡面同心服事.
聖靈往往就是做這些工作.
真的.
聖靈就是這樣在教會裡面運行.
不是說外面的人不做.
而是外面的人根本就不相信有聖靈.
根本就不會覺得自己的生命是屬於神.
所以他們還沒到那些階段.
讓聖靈去帶領.
但是教會人士就不同了.
我們真的有聖靈在我們的內心.
還有.
保羅告訴我們.
我們信著聖靈撒種.
必從聖靈修永生.
這是第六章的經文告訴我們.
我們要以聖靈為中心.
是屬靈群體的動力.
所以我們一起來.
在每天每個星期的聚集裡面.
就要不斷在聖靈的帶領下去撒種.
例如福音的聚會.
同心的為未信的朋友.
聽明白救恩.
按感動去信主.
我們需要在背後有爭戰.
有祈禱.

$^{721}$解除魔鬼撒旦.
在人心裡面的阻擋.
這些是我們一起撒種.
我們在聖餐.
師班的合一的領唱.
帶動弟兄姊妹的心.
進入那個境界裡面.
每一次的詩歌都是一個教育.
都是一個撒種.
讓我們體驗到.
原來這個地方.
在這個時刻.
是有一份屬天的平安.
寧靜和溫暖.
神聖同在的溫暖.
讓你感受到這一個多小時.
你是神兒女的身份.
其他地方是不會有的.
不要以為這個地方這麼靜.
其實不是.
是有神聖的寧靜.
是有一份神聖的平安.
是在其他地方你感受不到的.
在一群人聚集的時候.
我們去讓主居首位的時候.
你就會感受到一份屬天的平安.
教會.
真的.
你去一些很淑齡的陪靈會.
讓你感受到一份.
你自己很榮耀的身份.
我很多次.
在感受的聖樂會.
其實我又不會唱歌.
我坐在那裡聽師班的獻唱.
一個多小時.
你感受到.
那不只是一個很完美的聖樂的表演.
而是你感受到那裡有一份神聖的平安.
寧靜.

$^{761}$是其他地方沒有的.
所以我們要記得.
曾經在英國聖公會的一位牧師.
鍾馬田牧師.
他每個星期都有福音聚會.
有一天來做一個女士.
這個女士是有人帶她來.
原來這個女士是在外國.
教會的靈媒.
幫人占卜還是怎樣呢.
有人帶她來聽福音聚會.
聽完之後.
他跟鍾馬田牧師說.
他說我是經常接觸靈界的.
我知道有靈界.
因為我跟它有接觸.
但是我來到聖公會.
威斯敏斯特的教堂.
我感受到有靈界.
但是你們的靈是乾淨聖潔的靈.
它可以這樣分辨.
因為它是一個靈媒.
所以弟兄姊妹.
在神聖的聚會裡.
我們感受到一份聖潔.
平安.
溫暖.
這就是我們順著聖靈.
撒種.
我們每一個整個群體.
將來我們就收取.
永生的福樂.
永生的平安.
就是屬於我們每一個.
當然後面還要告訴我們.
這裡是信心作為我們救恩的起點.
身份的轉變.
在生活裡面的實踐.
和群體的一個建造.
我們一群人.

$^{801}$為著福音存在.
為著要敬拜這位造物主而存在.
為著我們一起同心.
在不同的事工上.
一起的努力.
在這個過程裡面.
保羅就提醒我們.
自由要小心.
因為我們每個人都有情慾.
我們如果太高舉自己的看法.
太主觀.
就會帶來一些紛爭.
但不是不應該有自己的見解.
不過要放在一個大家一起.
共傷的裡面.
誰的意見.
誰的角度.
是最好呢.
就應該大家都放下.
採納.
沒有一個意見.
是百分之一百perfect永恆的.
如果有的當然是最好.
但是.
正正神就是要我們不要這樣.
要我們大家都來彼此學習.
愛心互相服侍.
就令到群體不斷在這個大前提的價值觀下去成長.
所以今天我真的很感恩.
我們能夠可以透過聖經.
保羅提到四段關於聖靈.
要我們跟隨原來真正的最好的領袖.
不是哪個出色的演講者.
是聖靈在教會裡面帶領我們每一個聖徒.
所以盼望這節聖經.
是差不多全卷書裡面一個最重要的總結.
愛帶來的就是.
我們可以在全個律法裡面.
我們都能夠在這個大前提.
會做一件事.

$^{841}$我們想一下.
我們是不是出於愛心.
那就帶來很多的紛爭就會解除.
所以迦太教會他們當明白了自由之後.
保羅立刻給他們一個留意.
就是必須以愛心彼此的服侍.
因為全律法都包括在愛倫如己.
就是十誡的第二最重要的那一條.
這一句說話之內了.
盼望主來祝福我們每一位弟兄姊妹.
我們一起同心禱告.
感恩天父給我們能夠.
透過今天我們看的迦太書.
四段重要聖靈的工作的經文.
多謝主.
我們可以跟隨你.
我們可以看到.
我們在神的救恩裡面.
是多麼的有福和寶貴.
願你帶領感動弟兄姊妹.
明白聖靈在我們大家的心裡面的工作.
多謝你聽感恩禱告.
奉耶穌基督的名.
阿們.
\newpage



\section{馬太福音 22:15-22}
\label{sec:OFK7CEvBZU8}
\textbf{2025.08.17《這是誰的?》 經文:馬太福音22\_15-22 周克禮牧師}
\newline
\newline
連結: \href{https://youtube.com/watch?v=OFK7CEvBZU8}{\texttt{https://youtube.com/watch?v=OFK7CEvBZU8}} ~~~~ 語音日期: 2025-08-30
\newline
\newline
\hyperref[sec:6RiKDzQP6RM]{\small{< < < PREV SERMON < < <}}
~
\hyperref[sec:index_chronic]{\small{[返順時目]}}
~
\hyperref[sec:index_scriptual]{\small{[返順卷目]}}
~
\hyperref[sec:Lov4FCvlm6Q]{\small{> > > NEXT SERMON > > >}}
\newline
\newline
馬太福音 22:15-22
\newline
\begin{longtable}{cl}
\hline
\hline
章節 & 經文 (和合本修訂版)\\
\hline
22:15 & \begin{tabularx}{0.7\textwidth}{X} 於是，法利賽人出去商議，怎樣找話柄來陷害耶穌， \end{tabularx} \\ \\ \relax
22:16 & \begin{tabularx}{0.7\textwidth}{X} 就打發他們的門徒同希律黨人去見耶穌，說：「老師，我們知道你是誠實的，並且誠誠實實傳神的道，無論誰你都一視同仁，因為你不看人的面子。 \end{tabularx} \\ \\ \relax
22:17 & \begin{tabularx}{0.7\textwidth}{X} 請告訴我們，你的意見如何？納稅給凱撒合不合法？」 \end{tabularx} \\ \\ \relax
22:18 & \begin{tabularx}{0.7\textwidth}{X} 耶穌看出他們的惡意，就說：「假冒為善的人哪，為甚麼試探我？ \end{tabularx} \\ \\ \relax
22:19 & \begin{tabularx}{0.7\textwidth}{X} 拿一個納稅的錢給我看！」他們就拿一個銀幣來給他。 \end{tabularx} \\ \\ \relax
22:20 & \begin{tabularx}{0.7\textwidth}{X} 耶穌問他們：「這像和這名號是誰的？」 \end{tabularx} \\ \\ \relax
22:21 & \begin{tabularx}{0.7\textwidth}{X} 他們說：「是凱撒的。」於是耶穌說：「這樣，凱撒的歸凱撒；神的歸神。」 \end{tabularx} \\ \\ \relax
22:22 & \begin{tabularx}{0.7\textwidth}{X} 他們聽了十分驚訝，就離開他走了。 \end{tabularx} \\ \\
[1ex]
\hline
\hline
\end{longtable}
$^{1}$各位大家好,我是洪恩棠,很高興來到紅欣堂跟大家一起敬拜.
很投入,所以很興奮,我都聽到唱得很開心.
首先我講一個故事,是一宗新聞的改編.
我剛才開了,沒有反應,不是很行,你按,我按,我向這邊.
他不聽話,他不聽話,我按四五下,我做手勢好不好,介意嗎?.
我這樣,不好意思,我這樣做一下,好嗎?.
話說有一天,你經過深水Po ,去到深水Po 走,不知道做什麼,終於去到深水Po .
突然人頭湧湧,前面很轟動,大叫什麼?撿錢,你都在叫,好嗎?.
然後看到很多人湧向前面,衝前,是不是在拍戲?看周圍,什麼事?.
真還是假?抬高抬一看就看到這個景象,很多紅色的紙,不知道是什麼,在飄下來.
其實你應該知道是什麼,你猜到是什麼,在飄下來,有些人就爬上了地鐵站的上蓋.
為了可以高一點,先撿起來,希望撿多一點,然後當你呆住的時候,旁邊有人說,你還想快點過去.
身邊有幾個人跟你說,撿書,走路比嘟嘟還慘.
你都還沒踏出第一步,突然有張紅色的紙飄下來,飄落你的手中,逼你去接收.
你拿著的時候,心裡很大的疑問,這張銀紙,一百元當然是港幣,是誰的呢?.
旁邊有個大媽,她拿著一百元的大紙,哈哈哈哈,我撿到了,然後她的朋友更厲害,她說我撿到八張,拿來撥線,好像很興奮很開心,我更厲害.
然後突然,前面有警察,然後怎樣?那些人馬上向另一個方向走,再走前面的時候,看到他們收緊,要還錢,不是很行.
撿了八張的大媽,九秒鐘,已經不見了,有沒有人會問,那些錢是從哪裡來的?.
可能大家已經沒想過,已經不再問這個問題,很多人都會想地上撿到寶,問,是啦,謝謝,究竟我們又怎樣呢?我們先祈禱吧.
我們先看回經文,在這段經文裡面,似乎好像在討論交不交稅的問題,是吧?剛才說納稅還是不納稅呢?不過如果細心看的時候,其實我們知道重點不是在說交稅的.
不過不知道大家有沒有聽過提起交稅,曾經有些人說我不交錢給政府,看看會怎樣,不知道大家有沒有聽過,那出此真的有人試,知道後來是怎樣的嗎?無論寄什麼信都來,他不理,之後原來試的人發現銀行戶口自動被政府扣了錢,所以你不用想了,不交不行的,你放心了,只差你自動交,還是政府跟你拿錢.
我們看完經文的時候,開始的時候先留意一下耶穌身處在哪裡,可能要留意一下上文,追回上文的時候,21章23節的時候,耶穌當時在聖殿裡面,在一個敬拜神的地方裡面,甚至在教訓人,但是很快就有些祭司長和民間的長老去質詢他,劈頭底腳,這是之前的經文,所以你們聽得懂的了.
劈頭底腳是說你憑什麼權柄去教訓人,誰給權柄你的,不知道大家有沒有印象,那些民事法律上很喜歡質詢耶穌,一開始的時候已經充滿著火藥的味道,要大碰撞,耶穌反問了他們一句,令他們啞口無言,之後還說了一些比喻去教導人,為了要指證這些人,其實很淒涼的.
那些民事,那些宗教,聖殿裡面的工作者,事奉人員,明明他們天國裡面本來是有份的,但是很可惜他們不重視神,不重視這位爸爸,這位主人,這位天國的君王,所以耶穌說他們自己有一個結局,就是哀哭切齒.
來到這裡,22章的15節開始,一開始法利賽人商議,討論一些事情,開大會,他們討論什麼呢?不過在聖經記載裡面的時候,每次他們商議的時間都不是一些好東西,通常都是說我們要怎樣消滅耶穌,他們目的只有一個,就是很想製造機會,這次也不例外,就是怎樣去陷害耶穌.
這次不過他們的處理方法有點不同,一開口就不再是質問了,而是好聲好氣,裝得好像向耶穌請教一樣,第一句就說是夫子,很尊重,很禮貌,門徒也是這樣叫的,夫子,如果你不認識這些人,你認識他們,你知道他們的想法,知道他們心裡面的惡意的時候,你也不是說不害怕,是無管動,挺恐怖的.
他們活生生地展示什麼叫做少女藏刀,有把刀隨時拔出來,然後怎麼說呢?他說夫子,我知道你是一個誠實的人,並且誠誠實實地傳神的道,什麼情面你都不相信的了,因為你不看人的外貌.
明眼人也知道他們說的只是鋪排的廢話,只是一路鋪排,一路鋪排,讚耶穌誠誠實實的,因為如果真的這麼相信,他們心裡面的東西,他們早就相信了耶穌,怎麼會處處針對呢?但是他們說這些話出來,是想製造一個陷阱給耶穌踩,最後一句才是真正的話.
他們想請教耶穌,請問一下,你告訴我你的意見怎麼看,我們納粹比凱撒是可不可以的呢?他們這樣問出來,其實他們覺得自己很聰明,他們覺得耶穌怎麼解答都會死的,因為我講多一點.
當然很多人在場景裡面,很多人圍著耶穌,不只是法利賽,文士,長老那些,還有其他人在當中,聖殿很多人,當然自然是猶太人,他們很想認識真理,但很多時候普遍他們都會對當代羅馬帝國有一種負面的情緒在當中,尤其是主後第六年的時候,其實有些人已經確定了猶太人,他們設定了一些規矩.
就是納粹比帝國政府是一種褻瀆和不道德的行為,有這樣的觀念在當中,當然我不知道他們怎麼納粹,不納粹是不行的,如果耶穌要交稅的時候,無論如何都會令他們心情有一種失望,會減低耶穌對大家教導的那種公信力.
當然那時候也有他們自治區的君王,希律王就是當地的君王,這個自治區的統治權力是怎樣的?其實都是授權於整個羅馬帝國的,所以連投票都沒有的,他們已經準備好了,這班人,預備了,這句話,是怎樣呢?.
無論耶穌怎麼回答,要交還是不交稅,他們馬上會發難,如果耶穌說要交稅,他們會說:噢,你說要交給這個壞羅馬政府,如果說不用交稅,他們會說:噢,你說不交稅給我們的羅馬政府,所以他們好像很客氣去請教他們,但背後他們一點也不客氣動機.
他們會說:你死了,這個就是他們最真正的目的在背後,不過耶穌的回應是怎樣的?他第一句已經很不和他們客氣,不會像剛才那樣,像那些法利賽人,裝模作樣,他剛剛才用這些比喻去指明這些拒絕神的人是沒有什麼好下場的.
一開始一講完,這班人馬上又問耶穌,耶穌第一句就說,你這班假冒為善的人,為什麼要試探我呢?我們可以先想想,耶穌當時的心情是怎樣的呢?.
耶穌講比喻,作出一些警告,就是很想宗教的領袖人士去懂得珍惜,珍惜什麼?珍惜和神的那段關係,珍惜我們本身的身份,但這班人沒有當神是什麼一回事,他們一心只想剷除耶穌,無論用什麼計謀也好,即使用陷阱耕計.
試探這個詞有時候可以很正面的,例如測試,重點是試驗的目的是什麼?是想接受測試的人幫他,有時候考試也是幫我們,還是一個不良好,充滿敵意,充滿惡念的試探在當中呢?.
正如馬太科一開始的時候,在講耶穌的曠野是怎樣的呢?魔鬼四十天,就是試探耶穌,希望耶穌趁著軟弱的時候,前來的時候,就令到耶穌中招,今時今日這班宗教領袖的行徑,可能和魔鬼行徑是不遑多讓的.
耶穌說了一句話,他拿了一個上帥的錢來看看,他們就拿了一個小銀出來,按了一下,有個小銀掉了出來,值得留意的這兩個字,耶穌說拿了一個交稅的錢,他們就拿了一個銀錢出來,耶穌說的是一個叫做法定的貨幣.
通常他們當代的時候,運用的不止一種貨幣在當中,好像我們在香港,除了港幣之外,可能有些人都會收人民幣或者美金,如果我小時候去澳門的時候,或者今時今日都是,他們明明是有澳幣,澳門幣,但是你給港幣可以嗎?可以的,他照收,照找,很精細,還算了匯率,找回幾毛錢給你,很有趣.

$^{41}$所以就變成他們的貨幣算是法定的貨幣,有些可能日常用,有其他的貨幣都不奇怪,那時候他們很流通他們自己猶太人的貨幣,但是交稅的時候他們要交什麼?法定的貨幣,還不是用我的錢可以交,不可以的,澳門銀用美金去交稅,我們不可以的,要交港幣的稅,他們拿了一個小銀,這個小銀叫做什麼?這個小銀有名堂的,叫做大銀..
叫做大銀,有個名堂叫做大銀,這個大銀是什麼來的呢?今天我拿個兩元雞出來,你拿了兩元雞的時候腦海想起什麼?這個兩元雞有多重?是圓形還是什麼形狀?什麼顏色?兩元雞很特別的,我就特意拿個特別一點的,很有趣的,好像花花形的..
有個形狀,一講兩元雞的時候,你會有一種概念在腦海裡面,我試試在畫面上拿個兩元雞,因為有代表性一點,跟我們回歸之後的兩元雞有些不同的..
為什麼不同?因為最重要的是,你見到有一面有英女皇的頭像,是代表統治者的權力在這個地方,這個城市裡面..
其實當時那些人拿了一個Delarion,就是我剛才講的那隻小銀出來的時候,有些人會有反感的,因為他們的位置在聖殿裡面,暗了一個這樣的法定貨幣..
最主要是因為法定貨幣上面寫了一些東西,猶太人一般來說不是很接受,會覺得甚至有蝕俗的感覺,尤其是你在聖殿裡面拿出來..
當時竟然在一些法理上人的黨羽,他暗了一個這樣的小銀,可能有些人心裡面會想,有沒有搞錯?.
那麼Delarion是怎樣的?是怎樣的?是英女皇的頭,還是有什麼字在上面?在當代的時候,就是羅馬皇帝提比留的年代,Delarion就是這樣的,謝謝,很快,就是這樣的,不是很帥的,.
因為當時可能是鑄造技術的問題.我們先看背面,背面就是我們,如果我看就是右手邊那個,那個是什麼?就是一個人在椅子上,看不看得到,意思是他其實拿著一張橄欖枝,拿著權杖,兩隻手都拿著東西..
上面有拉丁文寫著大祭司這個名號,就是在說這個女士是大祭司利維亞,代表她是代表什麼?神明的身份,代表著神,當然這個神不是我們在說的真神野和話,代表他們在說他們所相信的神,去設立提比留為皇帝..
正面的內容又是怎樣呢?剛才的兩塊雞就是英女皇,現在這個Delarion就是凱撒提比留的圖像,字又快一點過下去,再下,這個就是他本身拉丁文的字,翻譯成英文,再按多一下,沒有聲音,不要緊,沒有聲音就算了,他只是讀英文出來,其實意思是怎樣呢?簡單地說,.
皇帝提比留凱撒神奧古士都之子,提比留是當代皇帝,他的父親是奧古士都,是羅馬的開國君王,在公元14年的時候,8月他逝世,羅馬的元老院決定封他為神的行列,.
你想想人多可怕,多大膽,這樣封一些人為神,為他建了神廟,有神廟供奉著,所以其實我們歷史到今時今日都是這樣,有影響的,我們8月叫August,其實是他的名字,不知道大家知不知道,August其實是這個皇帝的名字,即是奧古士都的名字,都是源於他..
當然他的兒子自不認清自己為神的兒子,是不是好像耶穌一樣,當然最不同的是奧古士都死了之後,是被一班小賊推他上神台,可能為了一些政治目的,為了人物崇拜,而被封為神的,其實連他都不知道,死了,他都不知道自己封我為神嗎?.
他都不知道這件事,而耶穌的父親是誰?是創天造地的天父,是我們宇宙萬物的創始者,為什麼猶太人會對這個小銀有反感呢?第一,當然他們對造像是一種禁忌在當中,因為敬畏耶和華,對很多有形象的物體,他們都有些忌諱,.
除了神叫他們造一些像之外,他們不敢再造任何其他的像.第二,更加剛才說了,這個小銀竟然將人字封為神,是一個很難以接受的事情..
那麼耶穌又介意什麼呢?耶穌拿出小銀,問了一個問題,祂說這個像,這個號,即是那些字,是誰的呢?其實本身這個不是問題,是人都知道的,我拿出兩塊錢,你都懂得回答我,是人都知道的..
或者交不交稅對耶穌來說都不是很重要,祂不是很重視的,大家當然很容易回答,這個銀子就是凱撒的,所以耶穌說凱撒的就給回他,交不交稅有什麼所謂,真的不是什麼問題..
在帝國管轄之下,其實是人都知道不能不交稅,迫於無奈,無論如何都要交稅,只是感情上他們不想,不喜悅交稅給羅馬政府,因為被統治..
反而耶穌提出一個更加重要的問題,就是在最後一句,祂說神的物當歸於神,和合本修定會有一個更好的翻譯,凱撒的凱撒,神的歸神.神的歸神這一句,首先什麼是神的?什麼是屬於神的物?.
如果我們以銀子為詢問,以銀子的圖像為詢問,就是說銀子的像是誰的?這個銀子上印著哪一個形象,就是屬於他的..
那有什麼是印了神的形象在上面的呢?大家覺得有什麼印了神的形象在上面呢?銀子,一個小銀子,一個天國金幣.我們眼見這個世界,有什麼是有神的形象展現出來的呢?.
聖經一早告訴我們,創世記的時候這麼說,神說,我們要照著我們的形象,按照我們的樣式做人,神就照著自己的形象做人,乃是照著他的形象做男做女..
那個就是我們,我們本身就已經有神的形象印在上面,我們本來就是用神的形象來做,我們本來就是屬於神..
如果我們照鏡子的時候,會不會覺得自己也很帥,很像神,當然不是,我不是說這些.神是個零,不是像我們這樣,有手有腳有眼有口,聖經當然是這樣形容的,讓我們形象化一點,明白多一點..
我和媽媽一起外出的時候,很多人看到我和媽媽的時候,就說,你們母子很相似,不過我老婆很多時候都笑,她說,你像爸爸多一點,我說,不像我爸爸胖一點,她說,不是,是你的動靜,你的行徑,你的說話,這樣的情況..
我就想,我不覺得,但是身邊的人就會這樣覺得,你會期望你的小朋友像你,還是你期望你的小朋友不要像我,千萬不要像我,我都惡習,你希望他們像你還是不像你呢?.
耶穌基督說,神的龜神,祂不是回答那些人的質詢,反過來,耶穌是邀請他們去詢問自己,究竟我們是屬於哪一個?.
我們是屬於自己,是屬於這個世界,還是屬於創造我們的神呢?法理賽人的形象像誰?他們外表像神的僕人,像眾人的老師,但是他們的內心像誰?.
他們心存惡毒要害死耶穌,他們究竟像神多一點,還是像蛇多一點呢?我不說他們,說說我們吧,我們詢問一下自己,我是誰的?我是誰的?我是屬於誰的?.
我們是屬於神,還是屬於這個世界呢?我們和主耶穌有什麼關係呢?我們的形象像不像呢?人類很有趣的,我這樣形容,我們人類既是貼地,同時也連於天上,貼地是因為神用泥土做我們,貼地也不得了,但我們同時有神的形象,神也出了一口氣在我們裡面而創造出來,.
我們同時連於天,今時今日,我們人生行事為人,我們能不能夠有神的形象,活出神的那份心意在當中,神的龜神,怎樣才算是龜神呢?我們是一個怎樣的基督徒呢?.
耶穌說完神的龜神這句話之後,我們回到那段經文,本來質疑他們人的感受是怎樣的呢?經文說他們覺得很稀奇,很驚訝,很訝異,然後還離開了,為什麼他們不說,哦,你說交稅就死了,他們沒有這樣說,他們完全閉嘴,他們想著乘機被一個陷阱耶穌踩,他們根本忘記了這件事,為什麼呢?.
因為他們聽得明白,他們知道耶穌在說什麼,如果他們不知道,他們傻乖乖地說,啊,耶穌,你死定了,還不可以騙你,還不可以騙你,但是聽得明白的時候,他們無言以對,銀子是屬於凱撒,你又是屬於哪一個呢?.
那時候他們收到這一句話,如果有盼望的想法,他們會慚愧,所以就回頭走了,或者不理想的想法,他們說,哎,真是失敗,又無法陷害耶穌,我們回去想個辦法再害他吧,作為基督徒,我們有沒有耶穌的身份和形象呢?.
歌羅西書這樣說,一起讀,預備起,穿上了新人,這新人在知識上漸漸更新,正如做他主的形象,當代很喜歡用衣著代表不同人的身份,所以保羅寫很多書信的時候,都很多時候說穿上脫下,代表身份,甚至是生命的轉化在當中,.
怎樣轉化呢?我想我們不會陌生的,首先要自死很多我們好不好的東西,淫亂污穢,邪情,惡慾和貪婪,同時要棄絕很多我們的行徑,老奴,憤怒,惡毒,畏羽,謊言等等,我們要脫去舊日那些不好惡意的行為當中,.
在之前疫情爆發的時候,我陪媽媽去看急症,因為之前她住的院舍房間有些確診者,我和我媽媽一起去排隊的時候,病房擠滿很多人,見到很多醫護人員出出入入,有些負壓房間,免得病菌流出,進行更換一些保護衣,要穿一些保護衣才能進去,.
如果醫護跟進完之後,確診人士穿保護衣四處走的時候,你會害怕嗎?他不懂得脫下那些東西怎麼辦?正如我們,如果我們帶著一些不好的意念在我們心中,我們不懂得脫下的時候,不僅影響自己,可能在感染或影響其他人,.
所以我們要學懂脫下細菌,有菌有毒不好的事情,在我們心中不好的意念,我們要脫下,當然還不夠,還要怎樣?第二步驟就是穿上,你們讀間線,我讀餘下,預備起.

$^{81}$要穿上衣物,因此兼許君有何忍耐?倘若這個人和那個人有嫌隙,總要彼此容忍,彼此饒恕,至於怎樣饒恕我們,我們都要怎樣學習去饒恕,除此之外就要穿上愛心,因為愛心是連貫全德的,你們要讓基督所賜的和平在你們心中,為此蒙召歸為一體,.
神的歸神,如果我們同意我們每一個都是屬於神的時候,就是以我們的生命改變成為一個最好的證明,不是我們口中說什麼,而是我們怎樣活現我們的生命,有主的形象,形象不是我們的外貌照鏡子的樣子,而形象是主的生命進入我們裡面,我們很渴望很想像主的那種情況,.
就像你剛才唱的Uber Uber,那些果子,很多顆子舉出來的形象情況,說回我一開始撿錢的故事為一個思考,無論那張一百元降落到我的手上,或是在地上撿到我拿著的銀紙有多少張也好,我要問的不是問這些錢是誰,我要問的反而是我屬於誰,我是誰的,我是屬於誰的,.
那我自然會怎樣,拿了錢之後,你會看到警察馬上走,九秒九,你不會的,你知道你屬於誰的時候你會怎樣,還給他,你根本不撿,你不走過去,不去臭熱鬧,被人踩到怎麼辦,那些人真的傻了,不要聽這些叫做,搶啊殺啊,或者搶啊,你會走都可能不走,.
因為看到他們這麼狂喜的狀態,你可能會說很危險,小心啊,小心受傷啊,你們,不是湧過去,袋多少袋到多少,先塞滿袋子,或者看到警察馬上閃,因為我們知道基督的心意,我們有基督的形象在我們裡面,面對這個世界,這個社會,很多很多很混亂,很紛擾的事情,.
不斷發生著,之前幾年都差不多排山倒海,或者搞到世界停頓的情況,生活很多問題,我們自己心裡面本身就算面對家庭,工作,壓力,家人等等,很多東西糾纏在一起的時候,我們更加要問,我們有沒有主的形象在我們裡面呢?.
我自己通常這樣想,我擔心不是這個世界,或者時局,或者局勢變得更好,還是更差,時局好的時候,有些人可以開心得太晚,夜夜生哥的情況,忘記了神,時局不好的時候,可能很多人很憂慮,很抑鬱,很難過,很無助,也都可能只會質問神,神啊,為什麼你不出手,為什麼你不幫我?.
所以很想大家再想多一次,我是誰的?我是屬於誰的?有沒有神的形象在裡面?我們學習去遷上憐憫恩慈,謙虛,溫柔,忍耐..
我這樣說,這份學習基本上是有些違反人性的,這個世界通常教我們,你敬我一尺,我就敬你一丈,你對我好,我當然對你好,你對我不好,我會對你好嗎?不好是很正常,我對你不好是很正常,因為是我們這個世界學回來的一套觀念..
耶穌基督的生命是怎樣的?當他甚至被釘在十字架上的時候,耶穌是不是第一句,他不是這樣說,他不是說,哇,你們這群惡人,負啊,擊殺他們,耶穌沒有這樣說過,是不是?.
耶穌掉進人面前,說一句違反人性的話,他一副赦免他們,他們做的,他們不知道的,可以說簡直是很沒有人性,不過應該這樣說,不是很沒有人性,而是超越人性,滿有憐憫恩慈的說話,來自主的說話..
當然我們不會因為有主的形象,我們日子變得好過一些,環境可能很多挑戰沒有變過,我們也不會因為有主的形象,我們自己變得百毒不侵,生活無憂,不一定,就算我們有病,有時候有菌感染的時候,或者生活開始越艱難的時候,我們都要面對,不過問題是我們靠著主面對,還是靠著這個世界的遊戲規則去面對呢?.
我想一個比喻,如果我們是一個代表耶穌的小銀子的時候,這個小銀子的樣子會是怎樣呢?我形容,我比喻一下,我形容有兩面的小銀子,正面就是耶穌基督的形象,就是我們記得他的憐憫恩慈,謙虛溫柔,當然你還可以說很多,我只說其中五樣,背後是主的十字架的形象,.
我們每一個人都要背起,而且看到有荊棘的冠冕,環繞著十字架,路途可能充滿崎嶇,艱難,甚至痛楚,但是原來是走向一條光榮的道路,回到羅馬政府當日,一個Dinarion就是一天的工資,給到盡,我當以現在工時這樣說,五十元,一百元,給你做足二十小時,.
都是一千多二千元,但是我們知道我們生命的價值是遠遠不止這個價值,因為耶穌基督用他的寶血,叫我們生命是價值連城,以至我們怎樣學習有神的形象,不單止他用他的形象打造我們出來,我們要活現出神的形象在當中,盼望我們珍惜,我們都願意努力..
願以我們的努力,保佑我們成為聖聖的聖徒,願以我們的努力,保佑我們成為聖聖的聖徒.阿們..
\newpage



\section{啟示錄 2:1-7}
\label{sec:Lov4FCvlm6Q}
\textbf{2025.08.24 《永恆的嘉獎-我知你的心》 經文:啟示錄2\_1-7 黃永康牧師}
\newline
\newline
連結: \href{https://youtube.com/watch?v=Lov4FCvlm6Q}{\texttt{https://youtube.com/watch?v=Lov4FCvlm6Q}} ~~~~ 語音日期: 2025-08-30
\newline
\newline
\hyperref[sec:OFK7CEvBZU8]{\small{< < < PREV SERMON < < <}}
~
\hyperref[sec:index_chronic]{\small{[返順時目]}}
~
\hyperref[sec:index_scriptual]{\small{[返順卷目]}}
~
\hyperref[sec:mR6Dgq2WzwQ]{\small{> > > NEXT SERMON > > >}}
\newline
\newline
啟示錄 2:1-7
\newline
\begin{longtable}{cl}
\hline
\hline
章節 & 經文 (和合本修訂版)\\
\hline
2:1 & \begin{tabularx}{0.7\textwidth}{X} 「你要寫信給以弗所教會的使者，說：『那右手拿著七顆星，在七個金燈臺中間行走的這樣說： \end{tabularx} \\ \\ \relax
2:2 & \begin{tabularx}{0.7\textwidth}{X} 我知道你的行為、勞碌、忍耐，也知道你不容忍惡人。你也曾察驗那自稱為使徒卻不是使徒的，看出他們是假的。 \end{tabularx} \\ \\ \relax
2:3 & \begin{tabularx}{0.7\textwidth}{X} 你能忍耐，曾為我的名勞苦而不困倦。 \end{tabularx} \\ \\ \relax
2:4 & \begin{tabularx}{0.7\textwidth}{X} 然而，有一件事我要責備你，就是你把起初的愛心拋棄了。 \end{tabularx} \\ \\ \relax
2:5 & \begin{tabularx}{0.7\textwidth}{X} 所以你要回想你是從哪裡墜落的，並且要悔改，做起初所做的工作。你若不悔改，我要到你那裡去，把你的燈臺從原處挪去。 \end{tabularx} \\ \\ \relax
2:6 & \begin{tabularx}{0.7\textwidth}{X} 然而你還有一件可取的事，就是你恨惡尼哥拉派的行為，這種行為也是我所恨惡的。 \end{tabularx} \\ \\ \relax
2:7 & \begin{tabularx}{0.7\textwidth}{X} 凡有耳朵的都應當聽聖靈向眾教會所說的話。得勝的，我必將神樂園中生命樹的果子賜給他吃。』」 \end{tabularx} \\ \\
[1ex]
\hline
\hline
\end{longtable}
$^{1}$剛才的詩歌,剛才的崇拜,不知道你們的感覺和心情如何?.
裡面的詩歌內容很有意思,是載著神的道,讓姐妹們投入去敬拜,去讚美,是美好的一個開始..
願以神都祝福我們今天的崇拜,我們一起禱告吧!.
慈愛天父,我們將今天的崇拜恭敬交託,懇求您臨到我們當中,願意透過您的說話,建立我們的生命,祝福我們在您面前的敬拜心,讓聖靈引導我們進入真理,讓我們打開屬靈耳朵,領受您的說話,禱告交託,是靠主耶穌基督得勝名求,阿們..
不知道你們有沒有定期做身體檢查,找醫生看看自己的身體近況如何?.
在過去一年有多少人去見過醫生做過身體檢查?都很多..
有沒有人在過去五年以外有做過,而在這五年之內沒有做過身體檢查?都有,因為你相當健康,所以在過去五年都沒有做過身體檢查..
當然,為什麼要做身體檢查,為什麼要找醫生做身體檢查,看看自己的身體?因為你只看不到自己身體有沒有不健康,隱然未現的隱疾或病患..
早些發現就早些醫治,病向前中醫.有些人可能長期都很健康,到後期才做身體檢查,可能發現很嚴重..
身體檢查是為了我們健康,如果醫生告訴你,你過去半年都做得不錯,可能你用藥,血壓方面很穩定,就繼續,鼓勵你繼續..
然後發現你吃的藥不適合,還是很高,他就會調教,建議你這樣這樣..
所以,定期有醫生看一看,看我們身體,都是安穩.如果突然間有一個醒覺,身體有醫生幫我們做檢查,我們屬靈的生命,我們的靈性,會不會有檢查呢?.
因為我們整個人除了身體之外,其實我們有靈魂的,聖經告訴我們,我們整個人有靈魂和身體,身體有醫生檢查,你會不會去找家庭醫生,幫我做一個靈性的檢查?.
我想做很多年,十幾二十年,三十幾四十年,我不知道我的靈性健不健康呢?有沒有人試過去找家庭醫生做這個邀請?.
他就說,我不是這科的,我轉介你吧.你猜他轉介你什麼科呢?可能是青山的吧?精神科可能適合你.你去全香港找哪個醫生幫你做靈性的身體檢查呢?.
耶穌,當然德國有研究四五十年自然教會增長,同教會或信徒把脈照X光,看教會健不健康,信徒健不健康,但更加暈眩的就是我們的大能醫生主耶穌,同教會同信徒做屬靈的檢查..
原來啟示錄第二章第三章,聖經裡面有對七人教會,七封書信做教會的一些靈明檢查,而每一封書信,七封書信,最後聖靈都會向眾教會說,有耳朵的就聽,除了給當時七間地方教會之外,聖經說聖靈對眾教會說,並且有耳朵的,.
即是任何信徒都要聽,即是說這七封書信啟示錄的第二章第三章,和這裡每一位弟兄姊妹都有關係,和錦書堂,紅宮恩堂,嘉靖堂都是有信息的..
當你預備了你的心去看一看耶穌和教會做身體檢查,包括在紅宮恩堂和靈明船做檢查,這個檢查是主耶穌做的,我們去評核的時候,主耶穌和我們去評估的時候,有兩方面今天要和你思想的,.
一方面就是我想你透過這一段經文,是以弗所這一封書信,你認識到上帝對教會有藍圖,有心意,有終極的願景,上帝有願景要和教會說,和你說,也有祂最終永恆的盼望,對教會,不知道你知不知道呢?.
突然間好像揭開主耶穌有一個永恆的藍圖,想給教會和你,然後好像醫生做了檢查,當時的以弗所教會,或者是地上很多的教會,有不足的地方,醫生就會給建議,你怎樣去修正,怎樣健康一點,.
這間教會主耶穌都有給勉勵,鼓勵他們怎樣修正,以至達到上帝的藍圖,我們就做檢查,洪恩堂或者我們個人的靈明,有沒有同樣的缺口,我們需要修正,所以第二段會和你思想,神對我們每一個人的方法建議,以至我們達到神的藍圖,.
意思是神不單止指出有毛病,並且給建議,給方法,給當時的教會,他們怎樣去修正,就好像醫生告訴你,你過肥了,你要做運動,留意飲食,這樣的向導,所以第二方面,上帝講了一些方法..
當然,剛才講的七封書信,不單止是講單方面,可能在教會的歷程發展,這七封書信裡面的靈明發展,都有和教會不同階段,可能洪恩堂十幾年,可能頭幾年,好像在書信裡面講第一方面,可能第二封書信,這個缺口,可能是這一方面,洪恩堂發展到一個階段,又或者信主二十幾三十年四十年,.
這七封書信都可以在不同階段成為你的肅靈提醒,好像我身體檢查,照心照肺,有七個報告出來,隨著你的身體,信主的年日,教會發展,都有不同階段,不同情況,所以這七封書信裡面都要留意一些內容,針對不同時段的發展..
在這二章第一節裡面提到,你要寫信給二佛所,這個教會,當時是使徒約翰在聖經一章,啟示錄一章二節,他在大概主后的五十年,寫成啟示錄,主后九十五年,當時羅馬迫北得很厲害,繼尼祿之後,鄧米田當時有些信徒,信道,.
而當時的處境是受迫北的時候,不是一個很安書,很安定的環境,而他在北魔島寫成啟示錄,是想讓信徒知道,君王雖然猛烈迫北基督徒,但是上帝仍然掌王權居首位,最終上帝會審判那些君王,君王都會過去..
然後在這裡面第一個描述,裡面都有書信的格式,書信格式第一方面,就是每一次都會介紹耶穌的尊稱,他的特質,他的身份,七封書信都有描述主耶穌的一些特質,然後講授書人,七間教會,他的名字,然後有稱讚,也有時有些評語,提醒他有什麼缺口,要修正,.
然後也有邀請呼召他們迴轉悔改,並且七間教會都有允許他們得勝的話,有神的獎賞.這幾方面的書信格式,一會兒我們都會留意到的..
而第一方面,一會兒我們仔細看,看看那個地方,當時就會有這些格式,剛才說授書者和主耶穌的身份尊稱,稱許,或者責備,警告,勸勉,如果不悔改,上帝會審判,然後有應許..
大概當時,這七間教會的地點,介紹少少給大家知道,在今天的土耳其,西邊,大概城與城之間,50到80公里,一個好像半圓形的馬蹄鐵的路.這間教會,如果你看到,由歐洲進入亞洲,一個入口就是爾弗所,爾弗所對上就是善美南,在這裡就好像一個倒轉的馬蹄形..
今天如果你去土耳其參觀七教會,你可能會由爾弗所,善美南,別加摩,一直到推亞推拉,撒地,飛拉飛鐵,老底加.這七間教會就是當時小亞細亞的一個地方,然後提醒當時的教會,有什麼地方做得好,有什麼地方需要努力改善..
這也是我們今天信徒的提醒寫照.這七間教會有些特別的地方,你要留意一下,七間教會,每一間教會都有它的特質,每一間教會都有特質,紅恩堂有紅恩堂的特質,錦宣堂有錦宣堂的特質..
而上帝欣賞第一間教會,爾弗所,一會兒我們詳細看,他有一些服侍很忍耐和很熱誠的事奉,他也有一些責備缺口,有一件事要責備他們,然後又勸勉,又加上..
你留意,七間教會裡面,是美娜是有讚武譚的,即是什麼呢?好像身體檢查,見到醫生,你確認指數很好,做得好,就繼續做..
師妹娜在患難,雖然很窮,但她做得很好,沒有什麼缺口,就鼓勵他們繼續忠心,至死忠心,繼續努力.但是呢,有兩間教會有讚武譚,是超級卓越,有三間有讚,也有提他們的缺口..
但是有兩間,請你留意,有兩間,薩迪和魯底加,我們聽慣魯底加,主耶穌想讚,想欣賞,想找他們的優點,都賺不到.咦,這樣很笨拙,咦,有這樣的缺口,原來他很多毛病,很嚴重的靈命不足,那要不要放棄呢?.
原來你看到每一封書信最後都有邀請和有應許,即是什麼意思?即是叫他們不要放棄,無論薩迪和魯底加多麼不足,多麼軟弱,都不要放棄..
所以上一次我在嘉欣堂分享聖經之前,有些人很謙卑,說我信了很多年,我就很沒用,我就很不足.有時信徒都會說,我靈命不好,我又沒有好處..
他一說到這時,我說你看看聖經,魯底加教會都沒有什麼讚他,但主耶穌有沒有放棄他?又沒有放棄他,薩迪在這裡都沒有留意到他有什麼優點,都沒有讚他,沒有嘉許他..

$^{41}$但是這兩間教會一樣有呼召,有邀請,有勉勵,仍然有應許,仍然上帝希望他得到上帝的永恆祝福..
即是說,起碼你有一個安慰,就是我們信主多久,你覺得自己靈命多軟弱,多不足,請你不要放棄.你回家再看第二章第三章,特別是薩迪和魯底加教會,神沒有放棄他們,神仍然很愛他們..
這個起步希望大家了解.剛才第一方面,很想你認識到上帝給我們的藍圖,就是他告訴你,他全然知道你,他所有都知道你,然後聖經描述他主耶穌的特質..
在啟示錄的二章一至三節,描述到要寫給已忽所教會的使者,七星代表當時的領袖,七個金燈台就是當時的七人教會,而描述耶穌是動態式的行走著,行走著這樣去看..
即是什麼意思?原來有時描述耶穌,可能坐著有權柄,亦有時描述,特別描述,他是行走著,然後說我知道你的行為,我知道你的行為,而原文是我一直知道已忽所教會的信徒的行為言語,他們的事奉質素,主耶穌一直知道..
而他知道是用動態,什麼意思呢?曾經很多年前,有一位屬靈領袖叫溫米耀博士,他分享,我們有時知道耶穌的臨在,需要有些幫助,他就說不如在家中放一張空的凳子,然後你跪在床邊祈禱,你知道耶穌坐在旁邊,你的祈禱是不一樣的..
有沒有人試過這種?沒有試過的話,今晚就可以實踐,晚上睡覺前不妨放一張空的凳子,然後跪在床邊祈禱,你的感覺是不一樣的..
上次分享完這件事之後,嘉欣堂的主席立即上來祈禱的時候,放一張空的凳子在床邊,立即祈禱,感覺都不一樣..
好了,你旁邊有張空的凳子,現在我告訴你,耶穌可以坐在你旁邊,可以坐在你旁邊嗎?可以坐在你旁邊..
如果耶穌坐在你旁邊有沒有不同?感覺是不同的..
有一次在團契裡面分享,有位神學老師也聽到,他說耶穌不止坐在旁邊,耶穌可以走路的..
不知道你有沒有印象,去到大型教會裡面,總是很多信徒坐在後面,前面那一排沒什麼人坐..
到處都是這樣的,這裡可能沒那麼長沒那麼寬,所以大家都很集中,但是越長的禮拜堂好像很天然的,相當多弟兄姊妹坐在後排,前排那些都是留著凳子的..
有沒有什麼原因?人的心理,坐在後面,有點累,說話有點悶,我閉著眼睛也不太覺,沒那麼呆眼..
在前排閉著眼睛,就不好意思了..
但如果你知道耶穌是走路的,坐到最後排,滾手機,他也看到的,如果不是滾聖經的話,手機在看著..
有時候信徒坐在後排,好像沒什麼,但是耶穌動態的,在禮拜堂前後,他都跟你一起..
並且你去到禮拜堂,在家裡,在職場,在家庭裡面,耶穌動態的知道你的行為,你的內心..
當然,人的心理很有趣,就是看不到,看不到..
其實學生也有時候是這樣的,特別是年輕人,心理上,你看到電視或東張成報告,夾公仔機,青少年總覺得好像沒人看到,.
但是那些CCTV同步,能不能拍清楚?.
但是他們想著沒人看到,他們就偷那些夾公仔機的公仔,是不是?.
年輕人,他不知道有人看到..
同樣地,學生有時候就那種挑戰,這種沒人看到的情況,特別是考試,公開試,或者是學校試,.
哇,那個廣場很長,老師在台前,或者在禮堂前,那麼遠,後排的學生看不到,看不到..
所以最近有些學生在五月,其中一間名校,他們學生就集體出貓,然後在七月四日,在黃大仙一間中三學生經濟科,他們又集體出貓..
現在他們成功出貓之後,最笨的就是什麼呢?.
他們發文給IG,或者是那些媒體,最笨的就是聰明過龍,我們集體成功出貓,他又以為其他人沒人看到他們的媒體..
這兩間學校出貓集體,就因為老師可能很遙遠看不到,但是公開試正常來說,可能不止一個老師在監考,老師還要巡下巡下,叫行動動態..
每個老師坐在前排,後面的學生就有試探..
我講到耶穌動態的意思,就是說,耶穌在我們任何時間,你一天廿四小時,任何一個處境,祂知道你的言語行為和內心..
當你知道耶穌講在第二節,我知道你一切,祂一直知道,過去,現在,甚至將來,祂全部知道..
你的優點,你的限制,你的不足,耶穌一直知道..
七間教會不斷重複,我知道你,我知道你..
當你知道耶穌知道你的言語行為,教會以外,在你家閉上房門,上網,又或者任何對人處事,祂知道你..
你坐下的時候的言語行為,你的道德操守,耶穌都知道..
若然你知道是這樣的,通常有兩個反應..
第一個反應,如果你一直做得好,做得為主義付出,為主義作見證,可能你很謙卑,在後面做影音,只是按part point,又或者你做接待,做一些隱言未言的事奉,很勞苦,人們不欣賞你,打錯字,或者按錯按鈕,或者你做一些教會掃奉獻,沒有人看到的,可能是你自己做的..
沒有人看到的,可能過幾個小時,有些人在大教會裡面,你幫忙掃奉獻,掃錯或漏,人們可能會提醒你,指責你,但是欣賞,可能欠奉..
你可能為主在教會,在家裡,付出很多,但是沒有人欣賞你的仇..

$^{81}$若然你知道今天有訊息,主知道你一直為他受苦,為主作見證,你為了傳福音愛家人,當如果你要持續很有心去愛人而付出的話,你聽到主耶穌知道你的辛苦,或者委屈,你怎樣?得到安慰..
特別是長期為主付出的犧牲,為了愛人的靈魂長期付出,沒有人欣賞,或者甚至有委屈的時候,耶穌一直知道你,你就會得到安慰,得到耶穌的明白..
相反,有第二類的情況,就是那些信徒長期是兩面人,兩面人的意思,這段經文裡面提到,他們的言語行為相當好的,以不鎖教會..
他們當時在大城市,一間教會裡面,他們的言語行為相當好,然後聖經在第二節說,他們勞碌,這間教會有很多人事奉,很願意事奉,並且有韌力,忍耐,不怕辛苦,哇,多好啊..
紅恩堂弟兄姊妹,又熱心事奉,遇到困難,又繼續事奉,又不累,接著他們不能夠容隱惡人,一些罪惡,他們不會跟他同流合污,他們的聖經神學相當好,能夠辨識當時自稱是使徒的,他們的信仰神學很穩實,辨識到那些假冒的使徒..
他們的忍耐,他們的勞苦,為主的名是不疲倦,不容易的,以前有些培訓,昨天聽到一位師弟說,我們培訓一些人沖奶茶,學的時候很多人學,到真的落手落腳做了兩次,不知為何好像很辛苦,受訓練的時候不少,接著就有些退去,因為慢慢就會累了,以及辛苦..
事奉,勞苦,又持續,相當好.而當時尼哥拉黨的人,他們的生活,行為,道德都是不妥當,他認同神惡的為惡,所以以不所計,教會的教會相當好,除了不同流合污,也恨神所惡的罪惡..
這些弟兄姊妹好不好?很好,這間教會好不好?很好,很卓越.但是,當然,地上教會稱讚他之後,上帝再盡心一點,他的願景不止地上稱讚以弗所教會,並且他還想講的就是在第七節,就是向宗教會說,好像以弗所教會一樣,.
他想得勝之後,會有神的樂園,生命的果子給他們.為甚麼呢?因為以弗所教會有很多優點,但是神提他們有一個缺口,他的缺口如果不修訂,他也沒有那種獎賞,或者叫做永恆的那份情況..
因為聖經說,間教會有時有他的優點,而以弗所教會有一個缺口,就是他的起初愛心凋棄了,如果他們不悔改,他就不得勝.地上我們信徒好像一個肅靈爭戰,我們要靠神靠聖靈,地上得勝,持續地得蒙神喜悅的話,上帝的永恆藍圖,是要告訴我們,神想我們永永遠遠,.
生命果樹的果子,就是啟示了和創世記的永恆生命,永恆和上帝的親密關係.上帝願景的藍牌,每一位信徒,永恆和主耶穌和神的親密關係,永恆生命..
這是他永恆的藍圖,但如果地上贊成了你,而你不生性,教會也不生性,那就失去了永恆的藍圖.特別要留意,剛才說的七間教會,有兩間不太好,但神也沒有放棄,一樣邀請他們悔改,一樣有應許,這個帶給我們任何光景,不要放棄..
然後,已不所有的無病,主耶穌給了一個方案,這個方案是什麼呢?在這個方案之前,提醒我們,有沒有一些地方我們不足,這些不足的地方,就是我們今天信徒沒有警覺性,以為我們有些行為,上帝主耶穌未必看到,有沒有哪方面,可能你覺得神也看不到的?.
有些人上巴士,可能他們看不到,我丟紙巾,我不小心把垃圾丟到地上,等別人清理,你會不會過去一直以為沒有人看到,司機是看不到的,但是今天告訴你,主耶穌,上帝會看到的,日常生活看不到的東西,耶穌也會看到..
而,祂一直告訴你,祂知道我們過去信仰,從何時信耶穌,我們的優點,我們的限制,祂全部知道,當你知道之後,對你的生活有沒有影響呢?你要想一想..
然後,如果你有一個願景,上帝終極想我們,好像父母對我們也有願景,你知道天父對你有願景的話,你願不願意去按著祂,我們的主耶穌,大能的醫生給你的指引,去圓滿上帝對我們的目標呢?吸不吸引到我們永恆的家長的虛應去呢?.
如果是的話,以下的重點告訴你,主耶穌給了方案,我們有不足的地方,我們可以努力..
耶穌給了以弗所教會的方案,他們的軟弱,有一件事是責備以弗所教會的,他們各樣東西都很好,很卓越,他們就丟失了起初的愛心,原來他們起初都是有愛心去發展以弗所教會的,他們起初的時候,行為與信仰是一致的,愛,內是一致的..
但是隨著日子,他們的愛心慢慢離棄了,這個離棄可以很理解的,初信主,或者是開頭事奉,很熱心,很愛神,但是隨著日子,會磨蝕的..
起初的愛心,就好像舊約聖經,新約聖經的初屬果子,是表示最純,最真,最好的那種表達,原來初屬的果子,或者叫初心,或者叫起初的愛心,那種愛神的心,如果能夠保持到見主面,這是天父的願景..
就正如進入婚姻,你拍拖戀愛的時候,熱戀,青年人很容易在拍拖之前熱戀,結婚初期或者後期,都開始仍然保持熱度,那種相愛,多好啊..
但是結婚十年,二十年,三十年,柴米油鹽之後,就開始有些不同了,之前他做的事情,多謝你,欣賞你,你煮得真是好,但是結婚之後,有機會十年,二十年,就開始有很多抱怨或者是埋怨,覺得應份,你做這個家務,或者你煮這餐飯菜,你洗衣服煮飯,你是應份,好像沒有了這個欣賞,多謝,感恩..
原來隨著日子,人有機會愛心會刁談,所以試過有些夫婦曾經說,不知道他配偶追求我之前,太太說先生很有創意,天天欣賞我,讚我,很多詩詞歌賦,先生很多創意,結婚幾年之後,不知為何先生的創意好像沒有了,沒有了以前那種欣賞,多謝,換來的有時是抱怨和不懂得感恩..
原來今天我們的行為,動機,態度,神是重視,我們過去剛剛信主,或者是水禮後,我們很感受到神的愛,我們回應神的愛,可能很有屬靈的胃口,聽到很有心,可能約上帝祈禱,不為甚麼,不只是求,喜歡和神聊天,祈禱,在熱戀和神,熱戀和耶穌,有愛的感情交流的時候,曾經可能很愛著神..
很熱心事奉,為主傳福音,可能有些信徒,這些日子是初信旨階段,信得久,結婚,有小朋友,慢慢對神的愛有機會冷淡..
所以好像以弗所教會有些信徒,他們起初愛心慢慢淡了,而主耶穌好像大能醫生給了一個方案,方案在第五節,在哪裡墜落,要在那裡悔改,而這個悔改就是透過回想,回想三個命令的導詞,要回想,要悔改,認同錯,知道錯,願意離開錯,是需要有悔改的..
是需要有行動,所以這三個命令動詞是給當時以弗所教會的信徒一個方案,醫生看完報告,讓他們知道,你冷淡了,你愛心,你是這樣重建的..
原來重溫回想是重建屬靈一個愛火的基礎,哪裡失足,哪裡跌倒,哪裡去回想一下.我太太有一次經歷,她跟我分享,她對這個回想失足跌倒有感受..
她以前住里園門道村,她有一次去宴會,穿著很漂亮的裙子,然後穿著新的絲襪,漂亮漂亮的從家裡準備去搭地鐵,在路途上到街上,前面有一灘污水,她因為急,快,沒有留意,踩到那些污水滑倒,她的衣服髒了,她的絲襪穿了洞,那你沒有辦法去宴會了..
立刻回到樓上,於是換回一件乾淨的衣服,新的絲襪.她說當有這個經驗之後,她說再到街上的時候,她已經想起之前在哪裡跌倒,她已經想起,當時跌倒的時候很尷尬,很狼狽,還沒到已經回想哪裡跌倒..
然後下去的時候,她的眼睛很留意著那灘污水,那灘污水她很小心注意,然後行動很小心慢慢走,經過,經過之後很安然,她很開心地跟我說,我明白了怎樣避開失足跌倒了,三部曲,回想,看著情景,然後小心翼翼行動避開..
這三部曲,若然你有知道你的靈命或不足跌倒,或者家人經常按著你的紅鍵,或者某些親友或同事不經意地有些話激起你的火,然後你的情緒很高漲,或者你的情緒不合基督徒的體統,你知道你經常在那裡跌倒,你都不想這樣發脾氣,還沒到這個失足跌倒,你覺得好像失見證的言語行為態度..
經常煩惱某些家人,可能你還沒發生的時候,你想改善,你就回想在哪裡被人激動,那些話怎樣激起你的火控制,你到的時候,即使他講,他無理罵你,或者他不經意地有些言語行為刺激你的情緒反應,你可能就留意,或者未見他的時候才祈禱,然後小心翼翼跟他相處等等..
其實日常生活,如果有些軟弱的基督徒,可能是一些色情網站,他跌倒的,還沒發生的時候,他已經知道在哪裡跌倒,哪個情景,什麼處境,早點去回想誰跌倒,到的時候,小心翼翼避開等等..
試過有一個家庭,他們的婚姻關係都很緊張,差不多破裂,但是現在的手機,有時候都會有些功能,這對夫婦已經去到差不多離婚的階段,先生就說手機原來是Facebook和Google Play,或者Google相簿,他會飛出五年前,十年前一些舊相給你..
這些舊相當我們一次一次,原來我們旅行是很開心的,我們曾經很甜蜜的,我們一家人那種開心甜蜜的時候,當他看到這些相片的時候,他回想那種感情,回想當日那種關係,他就不捨得,他就努力去還回..
是不是回想都是很重要?如果我們過去很愛主,可能曾經在哪個公園,哪個地方,你與主特別親密,你可以去那個場景的地方,與主聊天,可能你以前在公園坐下來,和神祈禱,或者在某個很安靜的地方,安靜讀聖經,或者你以前在港九年培靈經大會,或者某些聚的時候,你安靜聽道..
你可以回想一些與神經歷的愛的關係,重燃這份愛火..
有一位長者在疫情的時候,很欣賞,當時我們報養教會長者團契有七百多八百人,他們來自不同的年齡階段,有些五十多,有些六十多,有些七十多,有些八十多,九十多歲,不同的長者..
疫情的時候,教會崇拜停了,團契也停了,但是教會長者,他們經歷過神,有神的愛和恩典,不想就這樣與上帝失去了團體,團契和親近神的時間..

$^{121}$但是長者視像是有困難的,所以他們學視像比一般人都難,孫對他們來說是起步點,不斷開很多班,很多培訓,讓長者有一個屬靈的聯繫..
團契時間,或者是小組時間,他們用視像和孫來給大家一種屬靈的靈命,屬靈的一根稻草,這根稻草,這條線對他們來說是很重要的關係..
我記得有一個長者,八十多歲,在職場上是成功的,他會在一間相關公司的主席,簽很多支票,但是電腦不是他所長,他就來學,學電腦,學完就不記得,學完又不記得,學完又不記得..
他學到第四五次,突然夜晚八點多,兒子打電話來,他說媽媽想視像,他解釋給媽媽聽,視像不行,因為視像需要有一個主持人開了連結,他打開才行,下午練完,他在家夜晚八點多繼續想練習..
兒子聯絡,八點多打開給他,他懾而不捨,第五次都繼續要練習視像,是不是很欣賞他?很欣賞他..
雖然他真的會不記得,不是他擅長電腦的東西,不過幸好他每次都有姐姐在旁,姐姐最後成功,持續幫他,可以第六次,第七次,第八次,就連結弟兄姊妹團契,小組,他在肅寧的時候,經歷了疫情的時候的艱難..
看到弟兄姊妹那種懾而不捨,愛神追求神,弟兄姊妹跟神的那種愛的連結,也成為我們努力的開始..
愛神愛人的心如何保持?我們需要警覺,需要留心..
有一個家庭,他們一家都是基督教家庭,兒女都是清少,但是他們已經意覺到愛神愛家人的心是需要保持一份愛的關係,所以他們就邀請兩個兒女,當時他們都願意的,吃飯的時候一家人,四個人放下手機,都不容易的..
因為如果家長不刻意的話,有時兒女吃飯的時候都會滾手機,少了家庭的交通和親密的關係,所以愛需要培育,需要努力..
今天在這幾方面鼓勵大家去回想一下,檢視一下,第一方面,對神對人的那種愛是怎樣的呢?還有沒有以前的那份熱切呢?有沒有持續去感恩和讚美神呢?持續跟神聊天呢?.
而這個聊天是不為什麼的,我們很擅長求求,試試不為什麼,不求什麼,在家裡或在戶外約神,約神聊天..
如果你平時約人喝茶,或者喝高級茶,很投契,聊得很開心,很雀躍,試試約主耶穌,不是求什麼,跟主耶穌喝茶,喝東西的時候,站著聊天,讀聖經,一個小時,兩個小時,很安靜的環境,跟耶穌以前聊天的就聊天..
沒有啊,過去也沒有跟祂聊天的,未來你就跟祂約會一下,你拿手機約會主耶穌,我在哪個地方很安靜的,跟祂下午茶,不被人趕的,讀一下聖經,祈禱,去感恩,數算神恩典,可以不求什麼,不求什麼也可以,不需要為了飲食糊口,因為聖經說,先求祂的國,祂的義,祂會加給你所需要的東西,按著神的心意去感恩..
如果你過去日子有很多感恩和讚美的,你找一次約耶穌感恩讚美,錦宣堂,紅恩堂,奸堂都有個特質,師班敬拜詩歌,原來很強的,因為你留意到他們都是詩歌裡面載道,詩歌,剛才那三個詩歌,都是很有意思的,你拿著今天崇拜,約會主耶穌,自己唱詩歌,去禱告親人神,詩歌裡面有道的..
以色列人在曠野,有很多次麻怨,聖經描述有十次麻怨,然後他們二十歲以上的人又進不了迦南,只有約書亞和迦勒,還有二十歲以下的人進迦南,因為他們習慣了,像神很多恩典,反而不懂得感恩,很多麻怨,信得久的人有時就很多,像求,還有很多怨言..
但如果持續感恩,都是一個學習.所以有一次我分享完這個信息之後,我一向提,錦恩堂有個湖是很美的,平時我們習慣了湖很美,但外面的人進去就會很覺得這個湖很美,這個草地很美..
進到錦受堂那段路,其實你細心留意很多神創造的恩典.肅靈遺口有沒有愛慕神的話語呢?以前很愛主,很喜歡讀神話語,或者上一些主學,或者神學院有好的課程,去學慕神話語,你肅靈遺口有沒有呢?沒有的話,你比以前差了的話,求主給我們有肅靈遺口..
對弟兄姊妹的關愛呢?如果團契在小組,有弟兄姊妹提出需要,你信主久了,你缺乏了一份關愛,如果團契和小組是冷漠的話,你心裡會不舒服.所以主耶穌提醒我們,要有對神的愛,也要對人的愛,包括了《約翰福音》十五章,要彼此相愛,要服侍洗腳,這些想法是持續的,去愛神愛人..
對未信主的親友,救人靈魂的熱誠,對家人不要放棄,家人很長時間都拒絕你,吃檸檬邀請他聚會,都拒絕你的話,請你不要放棄..
最後講一個見證,在我們的人生裡面,向親友,向家人傳福音,是我們信耶穌的人責無旁貸,特別是你的家人,你第一個如果信主的話,就好像家庭的宣教士,為家人的靈魂守望,可能很難很難,我以前三十幾四十年前信耶穌的時候,都發覺到很難,特別爸爸,我有兩個哥哥和弟弟,於是就開始祈禱,.
十年後,都有點想放棄,爸爸都是很硬,有時都很抗拒,繼續祈禱,但爸爸當時有膽石,找了好的腳圖醫生和醫院治療,當時教會弟兄姊妹是一些閩南人探望他,他都聽不明白人家講什麼,即是那些福建人,但是他們的祈禱,他感受到愛心,感受到教會關心,好轉之後就去教會,然後信了耶穌,.
原來人生他病的時候是一個切入點,有機會神幫他,認識神,信了耶穌,後來我繼續禱告,為自己容易帶他信耶穌的,就是弟弟,所以早期就邀請弟弟去教會,去完一兩趟之後,教會很悶,然後就沒有了,沒有二十幾年之後,我弟弟突然人生有點閱歷,就問哥哥有沒有一些教會聚會,那時候聖誕節有報道會,又一次邀請他,.
他就去,並且人生有閱歷之後,聽福音已經不一樣,後來就信了耶穌,回教會到如今,爸爸等了他十年,弟弟等了二十年,然後有機會的時候,有些人問我弟弟為什麼信耶穌,回教會,除了聽信息以外,他說哥哥這二十幾年有教會聚會都沒有停過邀請他,雖然這二十幾年吃過很多次檸檬,拒絕過很多次,但是他說哥哥都不斷邀請,.
總有些好東西,總有些好東西,到有一天他悶悶地,人生有閱歷,他就想起可能教會有些特別,他就回教會,信了耶穌回到現在,然後三四年前,我大哥和媽媽又重病,在醫院的時候,大哥很清楚,想信耶穌,並且向家人表達信耶穌和水利,.
於是幫大哥信主水利,上年媽媽九十三歲,都信耶穌水利,雖然拜日正宗四十幾五十年,但是讓家人信耶穌可以是十年,可以二十年,可以三十幾四十年,然後一家七口,五個信耶穌,還有二哥和大家姐未信,在我人生裡繼續為家人信主受亡,.
而太太也是第一個信耶穌,大家姐,然後三個弟弟信耶穌,然後媽媽信耶穌,上年爸爸信耶穌水利,爸爸是很難,但是他一家六口都經不了上帝的祝福,持續地為家人的靈魂不要放棄,真的試過有時邀請一些等姐妹向家人傾呼音,.
有些等姐妹很容易放棄,她不會信的,再有些激烈一點的,她這麼差,她上天堂,為什麼她可以信耶穌上天堂?有些人已經下了閘,覺得不可能,但是邀請你,因著神的愛,去愛家人,為他受亡..
所以剛才那段經文告訴你,上帝對我們永恆的心意,就是藍圖,最終想我們的清許不只是地上,是想我們在世得勝持續裡外的愛主,信仰一致,我們的言語行為和我們愛主的心是一樣,那種永恆的家長得勝,這是上帝的藍圖..
方法方案就是你的愛,你的心思意念,常存有真信仰,裡面的動機愛神愛人,有盼望,不要放棄,不要放棄,無論你覺得多軟弱多困難,仍然有盼望,常存的信,盼望,愛,最大的是什麼?愛,請大家心裡持續愛上帝,好不好?我們一起禱告..
讓我們知道上帝的藍圖,也知道怎樣去滿足上帝祝福我們這份人生,就是永遠有一份和上帝親密的關係,願神祝福我們的人生常存有信,有望,有愛,我們的禱告交託是靠主耶穌基督得勝名求,阿們..
\newpage



\section{希伯來書 4:12-16}
\label{sec:mR6Dgq2WzwQ}
\textbf{2025 8 31主日講道2025.08.31 《真理的揭露與恩典的扶持》 經文:希伯來書4\_12-16 譚立晶傳道}
\newline
\newline
連結: \href{https://youtube.com/watch?v=mR6Dgq2WzwQ}{\texttt{https://youtube.com/watch?v=mR6Dgq2WzwQ}} ~~~~ 語音日期: 2025-09-06
\newline
\newline
\hyperref[sec:Lov4FCvlm6Q]{\small{< < < PREV SERMON < < <}}
~
\hyperref[sec:index_chronic]{\small{[返順時目]}}
~
\hyperref[sec:index_scriptual]{\small{[返順卷目]}}
~
\hyperref[sec:lUBiAGHTBWo]{\small{> > > NEXT SERMON > > >}}
\newline
\newline
希伯來書 4:12-16
\newline
\begin{longtable}{cl}
\hline
\hline
章節 & 經文 (和合本修訂版)\\
\hline
4:12 & \begin{tabularx}{0.7\textwidth}{X} 神的道是活潑的，是有功效的，比一切兩刃的劍更鋒利，甚至魂與靈、骨節與骨髓，都能刺入、剖開，連心中的思念和主意都能辨明。 \end{tabularx} \\ \\ \relax
4:13 & \begin{tabularx}{0.7\textwidth}{X} 被造的，沒有一樣在他面前不是顯露的；萬物在他眼前都是赤露敞開的，我們必須向他交賬。 \end{tabularx} \\ \\ \relax
4:14 & \begin{tabularx}{0.7\textwidth}{X} 既然我們有一位偉大、進入高天的大祭司，就是耶穌—神的兒子，我們應當持定所宣認的道。 \end{tabularx} \\ \\ \relax
4:15 & \begin{tabularx}{0.7\textwidth}{X} 因為我們的大祭司並非不能體恤我們的軟弱；他也在各方面受過試探，與我們一樣，只是他沒有犯罪。 \end{tabularx} \\ \\ \relax
4:16 & \begin{tabularx}{0.7\textwidth}{X} 所以，我們只管坦然無懼地來到施恩的寶座前，為要得憐憫，蒙恩惠，作及時的幫助。 \end{tabularx} \\ \\
[1ex]
\hline
\hline
\end{longtable}
$^{1}$各位紅人家的弟兄姊妹平安.
今次來講到特別有親切感.
因為之前帶完紅人家的弟兄姊妹一起去澳門.
所以來到個個都有點熟悉.
那些笑容特別不同.
因為我們經歷過那天在澳門周圍走.
身水身汗.
又很開心.
有些很甜美的時光.
小朋友又可以去玩的時光.
一起歡聚.
又可以認識澳門的文化和處境.
都是很感恩的.
今天特別選了這段希伯來書的經文.
想跟大家一起去領受上帝的話語.
這段經文對我來說是很深刻的.
特別是在我們信主.
我自己的經歷就是.
信主的頭五至十年都是不簡單.
不簡單是什麼呢?.
因為我是來自一個破碎的家庭.
所以發覺原來信主之後.
五至十年家裡不一定馬上有改變.
所以自己要靠著神的話語.
怎樣在家裡做好見證.
然後帶動家人去信主.
所以會有些很埋身肉搏的.
我們說的屬靈的爭戰.
或者自己要對付自己生命上很多的缺陷.
一些的留習.
我拿開這個讓我看不到這邊.
因為我高度有限制.
說一點自己的背景.
因為爸爸以前很喜歡去澳門賭錢.
所以就搞到傾家蕩產.
妻離子散.
在家裡的成長是怎樣的呢?.
就是沒有一口好吃.
因為原來他是夜班的.
所以我放學回家的時候.

$^{41}$就是他呼呼大睡的時間.
我肚子餓就不敢叫他起床.
因為有一次經歷就是我叫他起床.
想給我飯吃的時候.
他叫都不醒.
但當我嘗試打電話給我外婆.
說我很肚餓的時候.
我爸爸突然就醒過來了.
但就拿著一支雞毛掃.
那是怎樣呢?.
我以後肚子餓就不敢再出聲.
說我暴躁.
破壞了他在外婆心目中的形象.
所以為什麼生得這麼小呢?.
是有原因的.
快慾不良.
人們就很羨慕我.
說我吃多了都不肥.
這些其實是小時候的陰影.
所以從小在家裡.
爸爸教我很多旁門左道.
我不知道你家裡人教你善道還是旁門左道.
我覺得在家裡有兩個派別.
其實我外婆的幼稚園被外婆照顧.
小學被我父母照顧.
但因為媽媽要上班.
而爸爸又賭錢.
所以媽媽要很努力賺錢.
爸爸晚上做空間.
白天照顧家裡.
但其實他不是照顧家裡.
他只是在家裡睡覺.
通常看到媽媽新勞一日的工作後.
回來就很努力把所有事情都做好.
爸爸晚上上班後.
我們才覺得終於有個好日子.
爸爸在家裡的日子.
我們覺得是人間的煉獄.
所以親戚都知道.
我爸爸一出現.

$^{81}$我們原本很活潑.
但爸爸一出現.
他就形容我們.
為什麼好像老鼠見貓.
我們全部都不敢出聲.
因為知道這不是貓那麼簡單.
是一隻老虎.
所以小時候成長很不容易.
外婆教我們很好.
你看到有些人.
以前香港街道.
還是有很多要求黑的人.
外婆很熱心.
她說一定要幫他們.
不要管多少都要幫他們.
有人賣旗一定要幫.
外婆是一個很善心.
很願意幫所有人的人.
所以我就帶著這個心情.
但一到小學.
跟爸爸一起住.
好像完全不同.
爸爸就教我們.
你很多錢嗎?幫人.
不要幫.
他們顧他們.
我們顧我們.
然後他說留意地下.
地下有什麼?.
難道你跟我一樣教法?.
他就說留意地下.
一定有人掉錢.
你就撿起來.
不要出聲.
教導很不同.
星期六日.
沒有這一代那麼幸福.
就算他們不能上教會.
不能出街.
他們有東西玩.

$^{121}$無論電視.
ipad 智能手機.
我們那個年代.
星期六日不能出街.
因為爸爸喜歡賭錢.
所以我們看賽馬節目.
就是這樣.
所以你說.
我生不生爸爸的氣?.
我很生氣.
所以裡面有很多的積怨.
很多的仇恨.
很多千萬個為什麼.
也很不明白上帝.
所以初初.
去認識福音的時候.
我是很抗拒.
有上帝嗎?.
有耶穌嗎?.
為什麼我完全感受不到.
神的愛?.
為什麼我只覺得.
我來到這個世界很痛苦.
而我又覺得這個世界很不公平.
原本只有我一個女兒.
但是後來.
不知是不幸還是大幸.
我媽媽.
在我之後就生了一對雙胞胎.
開心嗎?.
原本很開心.
但是為什麼我這麼不開心.
就是因為爸爸對待很不公平.
爸爸是重男輕女.
所以爸爸說過一句話.
對我的童年帶來很大的傷害.
所以我之後要去服侍小朋友的時候.
我提醒自己千萬不要.
騙心.
一定要公平地對待.

$^{161}$不然有些事情.
在人的心裡.
要搞很久也搞不定.
我很記得那時候.
你知道小朋友打來打去很正常.
但怎知道我爸爸就說了一句.
我告訴你.
弟弟打你就可以.
你打弟弟就不行.
弟弟打你.
你不可以還手.
所以裡面有很多扭曲了的價值觀.
很多偏差.
會懷疑自己的價值.
會懷疑自己的生命.
但是感恩就是因為有神的真理.
幫助我成長.
不斷用神的話語.
糾正我裡面扭曲了的偏差.
我們一起看看經文.
PowerPoint可以出剛才讀經的經文.
我們一起看第十二到第十三節.
我們一起再讀一次.
第十二到第十三節.
預備起.
神的道是活潑的.
是有功效的.
比一切兩刃的點更鋒利.
甚至魂與靈.
骨節與骨髓都能切入開.
連心中的思念和主義都能辨明.
被造的沒有一樣在他面前不是顯露的.
萬物在他眼前都是赤露塘開的.
我們必須向他較仗.
這兩節的經文是什麼意思呢?.
他說神的道是活潑的.
怎樣形容他活潑呢?.
我用英文翻譯.
你們會更容易明白.
就是指Living.

$^{201}$如果有一灘水.
你會知道是活水還是死水.
活水就是不斷流動.
死水一潑就不動.
靜止的.
所以他說上帝的道是活潑的.
Living的.
意思是說.
原來神的道有持續的活性.
不是靜止的教條.
這對我很有作用.
雖然婆婆教我很多好的東西.
但我婆婆當時並不是信耶穌.
她可能受很多教條影響.
我記得有一次跟她去拜神.
不小心好像滑了一滑.
誰知我婆婆說了一句.
她說.
一定是你剛才沒有專心拜神.
所以出廟的時候就滑了一滑.
那一刻.
小朋友心裡就覺得.
這是什麼神.
我沒有這樣做.
就好像立即有些教條式.
有些.
(無語).
問號在心裡.
但我們的神不是這樣的.
然後他說上帝的道不單是活潑.
還是有功效的.
什麼是有功效呢.
其實是能夠產生一些實際的改變.
英文的翻譯是active.
能夠有能力幫助你的生命不斷運轉.
然後令你的生命有改變.
所以如果你是一個.
真的依靠神的道成長的人.
你是不會不變的.
你會變得更好.

$^{241}$變得更能夠結出聖靈的果子.
你會越來越像上帝.
會是這樣.
然後他再去形容.
上帝的道是活潑的.
有功效的.
還要比一切兩刃的劍更加鋒利.
什麼叫做兩刃的劍呢.
如果我們平時鋸牛排.
那些叫一刃的刀.
一刃的刀.
為什麼呢.
因為只有一邊可以切.
但是兩刃的劍.
就好像看武俠片.
這邊可以斬那邊也可以斬.
兩邊都可以很鋒利地去殺敵.
那些叫兩刃的劍.
原來神的道.
好像兩刃的劍更加鋒利.
怎樣去殺敵都可以.
但這裡不是說殺敵.
這裡反而是說上帝的話語.
很想進入我們的生命裡面.
幫助我們搞清楚我們的心思意念.
所以他說.
魂與靈.
骨節與骨髓.
都能夠刺入浮開.
連心裡面的思念和主義都能夠辨明.
什麼叫思念和主義呢.
思念就是我們想的東西.
我們想過什麼念頭.
主義是什麼呢.
是指意念.
我們的動機.
我們做這件事.
我們其實在想什麼.
我們想怎樣呢.
都能夠辨明.

$^{281}$也能夠判斷.
判斷對錯是非黑白.
所以原來有時候.
我們活在這個邪惡的世代.
可能如果我們不是一出生就信耶穌.
可能我們在這個世界都打滾過一段日子的時候.
其實我們很容易被很多世界的歪曲.
謬論.
一些風俗.
習慣.
會令我們以為是這樣.
但可能真的不是這樣.
我們就要重新給神的話語.
幫我們去分清那些是非黑白.
所以在我信主的頭五至十年.
我覺得是很不容易的.
因為可能我信了主.
我要在家做一個好的見證.
家裡沒有信主的人.
通常都會怎樣呢.
挑戰一下你.
又或者告訴你.
為什麼你信耶穌.
信耶穌是假的.
騙你的.
聖經不是真的.
拿很多道理來跟你辯論.
所以那段時間.
自己又要提醒自己.
不要說說變成什麼.
吵架.
又要做好見證.
然後可能家裡人.
未信主的.
有很多不好的習慣.
他們還沒有變.
去挑釁你.
怎樣去做好這些生命的堅持.
一定要繼續去看神的話語.
以及要祈禱去依靠神.

$^{321}$神的道這麼厲害.
但是我們怎樣去經歷它呢.
第十三節是說被造的.
我們每一個.
沒有一樣在他面前不是顯露的.
萬物都在他眼前.
都是赤露躺開的.
我們就要向他交賬.
交賬的意思.
你可以說.
其實我們要跟神交代.
所以我們有我們的賬要交.
我們未信主的家人.
有他們未信主的賬要交.
所以千萬不要覺得.
很詫異.
上帝這樣不公平.
為什麼我要堅持.
我這麼努力做.
但是家人可以繼續這樣.
但原來當我們不斷被神的話語改變的時候.
其實我們家人是知道的.
只是他還沒到那一刻.
就是時機去扭轉.
所以他可能心裡要很硬.
因為他不容易放下他們過往固有所信.
但可能就是到你那一刻.
對了.
他就會認不順了.
我很見證到我爸爸.
是很堅持自己的觀念.
我後來慢慢越來越長大.
甚至進入神學院裝備.
然後我慢慢回頭去看.
我爸爸整個人生的片段.
為什麼我爸爸會這樣.
讓我更加認識他的時候.
我覺得神正在幫我.
如何用上帝的愛或眼光去看待這個人.
不像從子女的下意上去看.

$^{361}$而是嘗試用神的眼光去看.
為什麼這個人會這樣.
我能不能愛他.
能不能像耶穌說的.
赦免他們吧.
因為他們所作的他們不曉得.
我會想.
如果我爸爸知道.
他是不是刻意這樣做.
還是他不想這樣做.
不過他被罪捆綁.
他沒有辦法.
他就泥足深陷.
我們就一起.
俗語說就是攬炒.
後來怎樣去經歷呢.
神就很奇妙.
我們剛才說的兩刃利劍.
就像一個屬靈的手術刀.
揭露自己的真相.
有時候也會用上帝的話語.
揭露別人生命上的問題.
以致知道.
原來種種人生命的不堪.
都是源於人的罪.
或者曾經經歷過的創傷.
以致導致的.
後來就認識我爸爸.
他來自中山石崎.
為什麼會來香港呢.
猜到嗎.
游泳下來的.
我爸爸早期游泳下來的.
但是很特別的機緣巧合.
我爺爺奶奶是當時的校長.
在國內發生這些事.
其實是很難處理的.
家人是否接納.
是否承認.
有些矛盾複雜的情況.

$^{401}$我爸爸是家裡的大哥.
所以那一刻.
就好像要和父母割裂.
但是我們回到家裡探親的時候.
用另一個形象去認識我爸爸.
因為我爸爸是大哥.
他下面還有五個弟弟妹妹.
我們上去後發覺.
那五個弟弟妹妹都很好.
我們小時候回鄉下.
都很不明白.
為什麼差這麼遠.
爸爸這樣.
但其他人是這樣的.
為什麼這樣.
真的很不明白.
但是從我叔叔姑姑的口中.
又認識了爸爸的另一面.
他跟我們說.
大哥.
他們叫我爸爸做大哥.
大哥很疼我們弟弟妹妹.
什麼都照顧好我們.
這麼大對比.
原來有很多人事的改變.
如果大家認識一些歷史.
那時候的知青.
要下放進農鄉.
要去做農奸.
他們要跟不同的領導.
原來跟一個好的領導.
就比較好.
但如果跟一個不好的領導.
有一句話你們應該有聽過.
做就36.
不做又36.
其實這個就是我爸爸的處境.
原來我爸爸經歷過很多.
其實我爸爸應該很能幹.
不過他跟了一些不好的領導.

$^{441}$所以令他原本是一個很積極的年青.
但在他的遭遇裡.
他覺得懷才不遇.
跟著不好的領導.
做又是這樣.
不做又是這樣.
整個人就好像很沮喪.
很挫敗.
又是懷疑人生.
然後到了香港.
他是新移民.
新移民當時找工作.
學歷是否承認.
完全不承認.
所以來到這裡.
又再一次被否定.
然後又是那個情況.
怎樣做都一樣.
都不會有人想認識你.
因為你是第一級.
所以慢慢到我成為他的女兒的時候.
跟他一起成長的時候.
就看到我認識的爸爸.
已經不是兄弟姐妹認識的那個大哥.
我認識的爸爸是一個.
既然做又36,不做又36.
所以就練正學懶.
用很多旁門左道去生活.
因為他多努力都覺得沒意思.
所以我應該見證到我爸爸.
進入人生最頹廢的時候.
連他自己都放棄他自己.
就是這個處境.
我覺得原來我一直看神的話語.
不單是幫我更加認識自己的罪性.
認識自己被扭曲的東西.
要怎樣跟著聖經跟著真理扭轉.
原來亦都更加幫我認識.
我爸爸被什麼罪.
被什麼扭曲了.

$^{481}$以至他不能成為一個好的男人.
好的丈夫.
甚至是一個好的爸爸.
所以我們一直看神話語的時候.
我們都要想一想.
如果當神用他的話語來管教我們的時候.
我們要跟他交仗的時候.
我們願不願意給神.
好像我們在一個屬靈的手術桌上.
讓神用他的屬靈手術刀幫我們解開.
然後看看我們心底有沒有陰暗扭曲.
未搞好的地方.
讓神幫我們治好自己.
所以神的道很想幫我們揭露真實的處境.
我們願不願意讓我們生命的深處.
被神治好.
我們就要認識幫我們下刀的是誰.
因為有句話是這樣說的.
畏疾是怎樣的.
忌醫.
有時候有些人很不想見醫生.
自己身體有點不妥.
不要看不要看.
沒事的沒事的.
有時候是因為你很怕.
做那個手術或下刀的人是怎樣的人.
但如果我們認識我們的主耶穌.
祂不是一個很嚴厲.
只是坐在一個審判台的醫生.
而是一個怎樣的神.
我們去認識一下.
就去到14至16節.
我們一起去讀.
預備起.
既然我們有一位偉大進入高天的大祭司.
就是耶穌神的兒子.
我們應當持定所宣認的道.
因為我們的大祭司並非不能體恤我們的軟弱.
他也在各方面受過試探.
與我們一樣.

$^{521}$只是他沒有犯罪.
所以我們只管坦然無懼地來到師恩的寶座前.
為要得憐憫蒙恩位作及時的幫助.
我們這位主耶穌偉大的大祭司.
他就是耶穌是上帝的兒子.
為什麼要特別去強調他.
然後再鼓勵我們.
應當要持定所宣認的道呢.
其實進入高天是在說.
耶穌基督已經完完全全為我們受死.
釘在十字架上.
埋葬第三天.
死而復活.
然後進入高天.
他已經在永恆裡成為我們的中寶.
透過耶穌這個中寶中間人.
成為一個很重要的橋樑.
令我們因著耶穌基督.
可以與父神建立永恆的關係.
這個耶穌是經歷受苦.
受過試探.
和我們一樣.
只是他沒有犯罪.
他那裡特別有個字去形容.
他說這個大祭司並非不能體恤.
其實體恤這個字.
在原本的文字.
其實是由兩個字組成的.
由什麼字組成呢?.
由同感和受苦組成.
就是說.
基督並非遙遠的審判者.
即是你有錯.
你有錯.
看清楚了.
指指點點.
不是的.
反而耶穌基督是親力試探.
而無罪的救主.
他不是在審判台前.

$^{561}$如果.
剛才用了什麼字去形容呢?.
他在哪裡呢?.
他在私恩寶座前.
這個就是一個很大的對比.
他不是一個高高在上的審判官.
怎樣怎樣.
而是他飽受這些痛苦.
試探的經歷.
他是以一個很慈悲.
很明白.
很同感.
很知道你怎樣受過這些艱難的.
這位大祭司.
是用他的慈情哀愁.
用他的慈悲的眼光去看著大家.
我明白你.
我知道你受這些苦.
被罪捆綁是很艱難.
或者你被家庭的創傷.
或者世界的一些不公平的對待.
令到你會這樣.
大祭司就是在私恩寶座前.
很關懷我們.
很想幫我們離開這些罪責捆綁的上帝.
所以他鼓勵我們要怎樣呢?.
他說.
在試探裡.
其實耶穌基督不是旁觀者.
而是我們的同行者.
所以他鼓勵我們要怎樣呢?.
坦然無懼地去到他的私恩寶座前.
坦然無懼.
英文的翻譯是.
BONUS.
即是大膽.
我們大膽地去到神面前.
不是說.
哎呀!施主叔叔.
神會罵我.

$^{601}$神會怎樣.
不是.
其實就好像.
爸爸媽媽.
很有愛心的爸爸媽媽.
做錯了事.
你做錯了什麼事呀?.
不是做錯了什麼事.
是不同的.
我發覺我做了媽媽之後.
我在這兩邊會游手.
因為我的DNA.
我的血給了我爸爸.
令我有點暴躁.
很容易看到小朋友.
一頑皮.
或者一不聽話.
就會著火.
但是我又信了耶穌.
我又要被神的話不斷調教.
所以我都有向我的女兒道歉.
我說有時候媽媽太兇.
我都要向你道歉.
因為媽媽也不想.
公公比媽媽更加澎湃.
我已經很淡了.
但是我都會有錯.
我會努力去承認.
因為我自己很深經歷.
我和爸爸的恩怨情仇.
爸爸從來犯錯都不會說對不起.
他怎樣錯.
我都會向他道歉.
我會努力去承認.
我覺得如果我有錯.
我寧願教會他.
要勇敢承認.
我都請求我女兒原諒我.
我印象很深刻.
記得有一次教她功課.

$^{641}$我真的越說話會越大聲.
我女兒很乖.
我覺得是上帝給我的小天使.
幫我成長.
她流著滴眼淚.
媽媽,你可不可以不要這麼大聲.
好不好?已經很好.
我說,女兒,對不起.
媽媽又來了.
我說,你讓媽媽休息一下.
一會兒再來.
我說,你做得好.
你夠膽跟我說.
如果你覺得媽媽越來越大聲.
越來越火.
你勇敢一點.
我鼓勵她.
你勇敢一點跟媽媽說.
我就收服自己.
因為我經歷到.
我的童年.
我已經不夠膽說.
爸爸的強勢.
暴龍式.
我寧願我們全都失宿一角.
不說.
但我想鼓勵我女兒.
你夠膽說.
其實媽媽想改變.
你夠膽說.
媽媽要變.
媽媽不要走回舊路.
所以這裡說.
我們只管坦然無懼.
大膽一點去到神面前.
其實上帝很想用祂的恩典.
用祂的恩惠.
用祂的憐憫.
用祂的幫手.
幫我們一起去改變.

$^{681}$所以我們要得憐憫無恩惠.
作及時的幫助.
雖然神的道是兩刃利劍.
很鋒利.
刺刺刺刺刺刺刺.
知道你的狀況.
但原來上帝不是為了審判我們.
祂為了要什麼.
私恩典要幫助我們.
所以得憐憫是什麼意思呢.
原來得憐憫是針對過去的失敗.
你做錯了事.
但上帝有憐憫.
憐憫的英文翻譯是mercy.
大家有沒有吃過那種朱古力.
見到那首廣告歌嗎.
很深入民心.
mercy輕鬆我柔軟的感激.
mercy永遠給你.
即是原來上帝很想給你憐憫.
你不要自己捱.
不要覺得自己一錯再錯.
回不了頭就不面對神.
不要怕上帝.
上帝說你來吧.
你錯了不要緊.
你這次不成功不要緊.
你這次失敗不要緊.
來吧來吧.
記不記得耶穌的門徒.
其實有很多都是失敗過的.
最經典的舉個例子.
誰三次不認主.
彼得.
耶穌早就預告了.
耶穌怎麼跟他說.
你會在雞叫以前三次不認我.
然後耶穌有沒有說.
你以後不要回來找我.
不是.

$^{721}$他後面那句說什麼.
你回頭之後.
堅固你的弟兄.
是誰給他機會.
耶穌.
所以原來耶穌去面對我們的失敗.
他不是要記仗.
反而是要幫我們.
去明白我們.
然後叫我們在跌倒的地方.
重新站起來.
然後再一次.
再去為神作見證.
所以我們這一班的基督徒.
甚至我是一個傳道人.
我是一個宣教士.
我是不是完美呢.
都不是.
但是我們有做得不好的地方.
大膽一點.
去到上帝面前.
上帝是想幫我們.
想扭轉我們.
想救我們.
我們還沒有再一次.
做好自己.
求神給恩.
給力量.
我們可以跨過那些艱難.
原來有些過錯.
有些失敗.
是需要很多很多很多次.
之後不放棄.
才有突破.
我們就給自己機會.
不要放棄.
蒙恩惠是什麼呢.
恩惠.
英文翻譯大家會理解.
Grace.

$^{761}$恩惠.
其實這個恩惠是說.
供應你當下的需要.
是一個及時足夠的供應.
而作及時的幫助是指什麼呢.
英文的及時是In time.
在你最需要的時候.
就幫助你.
所以是一個及時的援助.
在我人生裡.
都有經歷很多這些事.
我知道在在有在讀神學.
就要鼓勵一下.
我那時候.
在神學院之前.
有全時間裝備讀短宣中心.
然後在神學院裝備幾年.
其實全部都是在學習上帝的供應.
我記得讀短宣中心的時候.
那時候我已經出來社會工作了一年.
然後就有感動進去讀短宣中心.
我記得要很節儉.
因為我是大姐姐.
我已經沒有給家裡經濟來源收入.
媽媽就要靠她繼續供養弟弟.
爸爸又拿不到錢回家.
所以那時候學習的生活要很簡樸.
其實短宣中心也很好.
我們那個年代是有一個箱子.
那個箱子叫做耶和華爾納.
裡面有什麼呢.
大家猜猜.
有什麼錢呢.
瑞銀.
那時候我們一進去讀.
老師就說.
在隔壁房間裡.
有這個耶和華爾納箱.
裡面有很多瑞銀.
瑞銀箱旁邊.

$^{801}$也有一堆梳打餅.
如果同學有需要.
你就懷著感恩的心.
就在那裡拿.
你猜我有沒有拿過.
沒有.
那就是說.
上帝供應我.
我不需要透過這個箱子得供應.
在其他地方得供應.
那時候神感動了一對夫婦.
他們都是見證著我順主成長.
他們每個月.
很努力堅持那兩年沒有停止.
他們每個月.
支持我五百塊.
那時候讀短孫中心是很划算的.
我就算全時間去讀.
我一個月只付八百五十塊學費.
很深刻.
給五百塊.
教會也有些遵從支持.
然後就去經歷.
然後呢.
怎樣有及時的供應呢.
有一次我很深刻.
因為要節衣縮食.
有一天早上.
早了回去中心.
突然間.
在我腦海中浮現一個早晨全餐.
不是英文字那一間.
茶餐廳茶記那些全份的餐.
我就心想.
為什麼無端端想這些東西.
之後我就說.
真的很想吃.
因為很簡樸.
可能兩年內沒機會吃這些豐富的早餐.
然後呢.

$^{841}$怎料心裡有提醒.
吃早餐也有提醒.
為什麼你不向我求.
求得嗎.
為什麼不行.
我就覺得很奇妙.
那我求了.
他就說你求吧.
你想求什麼.
我就真的求過早餐.
有雞蛋.
有香腸.
有通粉.
有杯飲料.
阿們.
我求完了.
前面有沒有東西.
沒有.
沒有東西.
就走去餐廳.
走去飲食水的地方.
走廊有個小廚房.
走過去.
沒有早餐吃.
就泡杯水.
怎料.
我走過去的時候.
有另一個短孫學員.
就拿著兩大鍋.
很豐富的早餐.
就進去餐廳.
我就有很羨慕的眼光.
唉.
人家讀短孫.
我讀短孫.
為什麼這麼不同.
這樣.
這些都是生命的成長考驗.
但怎料.
最奇妙的地方是.

$^{881}$那個短孫學員跟我說.
立貞.
有份早餐是上帝今早叫我買給你吃的.
哇.
你說什麼.
是啊.
我今早祈禱的時候.
上帝是叫我要買一份早餐給你吃的.
哇.
真的嗎.
那個姐妹.
也是一個很祈禱的姐妹.
生命是經常祈禱.
經歷神的.
所以跟神很親密.
哇.
那我很大的經歷.
真的嗎.
真的.
你拿去吃吧.
這一份是你的.
我經歷了什麼.
上帝有聽我祈禱.
還要在我祈禱之前.
祂已經為我預備.
時間完全不是同一個時空.
是在另一個之前的時間已經發生.
我帶著很驚喜期待的心.
打開這份早餐.
就是有這些東西.
我那一刻是怎樣.
我很感動去吃這個早餐.
再一次看到.
原來上帝沒有忘記我.
上帝是認識我.
上帝是知道我.
就是這麼小的事.
一個早餐.
祂都會看顧我.
所以在我蒙召.

$^{921}$奉獻的路.
這麼多年.
都是經歷上帝親自的供應.
所以我們都是不斷磨練了.
對神的信心.
所以神是很願意幫我們.
到最後.
長話短說.
我爸爸到最後是信了主.
這個十惡不赦.
原本很想他落地獄的人.
到最後他悔改.
到他要認罪信耶穌的時候.
他就是人生.
願意認錯的那一次.
所以我們不要放棄我們自己.
我們就算信了主.
我們都有機會跌倒.
但是我們不要放棄.
上帝這個這麼寶貴的救恩.
我們要知道.
上帝的道是鋒利的.
他不是用審判官的角度.
去斬草除根.
而是用基督的柔軟的心靈慈愛.
他很想用他的恩典.
扶持我們成長.
神的道雖然多麼鋒利.
但他是用體恤的目光.
要展現的是神.
既是一個拆毀的神.
拆毀是為了什麼呢.
要重建你的生命.
重建我的生命.
令我們的生命更加亮麗.
更加被他使用.
所以我覺得.
在我的生命裡.
我要去經歷成為一個澳門傳福音的人.
就是神改寫了我們整家人的命脈.

$^{961}$我爸爸去澳門賭錢.
我去澳門不是賭錢.
我是澳門要將福音帶給那裡的人.
這已經是一個很大的扭轉.
你知道有很多賭的家族.
他們會牽連他們的子女也會賭.
我不用這樣.
然後我還可以扭轉.
靠著神的恩典.
可以將這個扭轉.
分享給其他澳門的信徒.
我們一起去做好這個工作.
所以盼望在洪恩家的每個特映節目.
我們大家一起興起.
一起被神建立.
一起去祝福和賜福這個社區.
那時候我是很感激.
有些很大膽的基督徒.
是夠膽去向我爸爸傳福音的.
我爸爸的狀況就是.
好像樓下的街坊.
擔著口煙.
穿著背心.
很有氣勢地坐在那裡.
誰夠膽靠近.
但是那些很有愛心的基督徒.
沒有嫌棄我爸爸.
很有愛的接觸我這位爸爸.
然後帶他回教會.
帶他信耶穌.
所以鼓勵我們一眾弟兄姊妹.
一起興起.
在香港好好做傳福音工作.
也在離開香港其他地方.
都可以好好去作見證.
我們一起同心同意.
慈悲仁愛天父.
多謝你留下你的話語.
多謝你願意來到世上.
多謝你願意讓我們可以大膽地.

$^{1001}$在你的私人寶座前.
去蒙恩慧.
得連率.
作隨時的幫助.
主啊求你讓我們每一個靠著你的恩典.
我們可以剛強起來.
我們可以堅持.
我們可以彼此的同行.
讓我們弟兄姊妹之間.
萬事互相效力.
叫愛神的人得益處.
多謝您 您在我們面前誠心禱告 奉主明求.
\newpage



\section{列王記下 5:1-7}
\label{sec:lUBiAGHTBWo}
\textbf{2025.09.07 《大時代中的福音見證》:列王記下5\_1-7 (和修版) 黃國樑傳道}
\newline
\newline
連結: \href{https://youtube.com/watch?v=lUBiAGHTBWo}{\texttt{https://youtube.com/watch?v=lUBiAGHTBWo}} ~~~~ 語音日期: 2025-09-11
\newline
\newline
\hyperref[sec:mR6Dgq2WzwQ]{\small{< < < PREV SERMON < < <}}
~
\hyperref[sec:index_chronic]{\small{[返順時目]}}
~
\hyperref[sec:index_scriptual]{\small{[返順卷目]}}
~
\hyperref[sec:pWTRHiVMfgo]{\small{> > > NEXT SERMON > > >}}
\newline
\newline
列王記下 5:1-7
\newline
\begin{longtable}{cl}
\hline
\hline
章節 & 經文 (和合本修訂版)\\
\hline
5:1 & \begin{tabularx}{0.7\textwidth}{X} 亞蘭王的元帥乃縵在他主人面前是一個偉大的人，得王的喜悅，因為耶和華曾藉他使亞蘭人得勝。他雖然是大能的勇士，卻染上了痲瘋。 \end{tabularx} \\ \\ \relax
5:2 & \begin{tabularx}{0.7\textwidth}{X} 亞蘭人成群出征的時候，從以色列地擄了一個小女孩，她就服事乃縵的妻子。 \end{tabularx} \\ \\ \relax
5:3 & \begin{tabularx}{0.7\textwidth}{X} 她對女主人說：「我希望主人去見撒瑪利亞的先知，他必能治好主人的痲瘋。」 \end{tabularx} \\ \\ \relax
5:4 & \begin{tabularx}{0.7\textwidth}{X} 乃縵去告訴他主人說，從以色列地來的女孩如此如此說。 \end{tabularx} \\ \\ \relax
5:5 & \begin{tabularx}{0.7\textwidth}{X} 亞蘭王說：「你可以去，我也會送信給以色列王。」於是乃縵手裡帶十他連得銀子、六千舍客勒金子和十套衣裳去了。 \end{tabularx} \\ \\ \relax
5:6 & \begin{tabularx}{0.7\textwidth}{X} 他帶著這信給以色列王，說：「現在你接到這信，看哪，我派臣僕乃縵到你這裡來，你要治好他的痲瘋。」 \end{tabularx} \\ \\ \relax
5:7 & \begin{tabularx}{0.7\textwidth}{X} 以色列王讀了信就撕裂衣服，說：「我豈是神，能使人死使人活呢？這人竟派人來，叫我治好一個人的痲瘋。你們要知道，看，這人是找機會來跟我吵架的。」 \end{tabularx} \\ \\
[1ex]
\hline
\hline
\end{longtable}
$^{1}$弟兄姊妹平安.
今天很感恩來到洪水橋洪恩堂.
和弟兄姊妹分享神話語.
今天是第一次在港台和弟兄姊妹分享神話語.
其實我是第二次來的.
因為第一次來的時候.
帶著我們社會福音堂的弟兄姊妹.
有些老友記我們去旅行的時候.
有機會去年你們的茶座開始的時候.
我們就成為第一批的弟兄姊妹.
有份見證你們茶座的開張.
代表我所服侍的教會.
是社會福音堂的弟兄姊妹.
和洪恩堂的弟兄姊妹來問安.
石澳在南區最端端的一個地方.
整個石澳範圍都叫石澳郊野公園.
裡面有一個村落.
我就在村落裡面.
一間浸信會的福音堂服侍.
來到洪水橋這個社區.
相信是一個新的社區.
因為北都會的發展.
有不少不同的部分的新的社區出現.
相信日後也會有商區.
在整個北都的計劃來發展.
我相信教會一定是充滿活力.
教會的發展一定是充滿了活力.
和弟兄姊妹一定是熱心為主.
因為是一個開發的發展中的社區.
不過在這個社區我們有機會做見證.
但是其實在這個大時代之中.
我們可能身處在不同的崗位上.
在社區裡面.
我們都用不同的方法.
我們怎樣來為主做見證呢?.
我相信是好看的.
我們生命是怎樣.
我們是用我們的生命來為主做見證.
今天和大家分享的這段經文.
大家不會陌生的.

$^{41}$因為連主日學的老師.
或者主日學的學生.
都有機會來聽過.
乃萬的《瑪豐德伊治的故事》.
我希望在這段經文裡面.
看到幾個主角.
其實在當時的大時代裡面.
怎樣透過他們的生命.
來為主做見證.
幫助我們今天.
其實我們怎樣.
神是怎樣透過我們每一位呢?.
我們有耶穌基督復活的生命.
的改變.
我相信大家都有耶穌基督復活的生命在裡面.
我們怎樣用不同的方式.
來見證福音的大能.
在人的身上.
讓人看到認識耶穌.
看到耶穌和經歷耶穌.
在繼續講道之前.
我們一起同心再次禱告.
賜寶貴的話語.
求聖靈你光照引導我們.
在你的話語之中.
我們有學習被教導.
我們的禱告.
奉耶穌基督的名祈求.
阿們.
今天我用兩個主題.
兩個跟大家分享.
第一就是我們身處在大時代之中.
其實對言福者.
是有很大的挑戰.
什麼叫言福者呢?.
現在我們很少講傳福音.
因為有很多敏感的字句.
特別是其他地區.
創體地區.
都開始慢慢改用一些沒有敏感的字句.

$^{81}$例如言福.
我們將神的福氣延伸出去.
所以我們多用言福這個詞語.
來形容傳福音的人.
其實我們都面對很多不同的挑戰.
特別是在現在的時代當中.
不同的環境.
世界的局勢.
香港的處境.
香港和周邊地區的那種的.
互動.
我們有時候傳福音很難.
有些地區我們不能夠直接講福音.
有些地區.
連你做見證.
你離開教會.
講都不行的.
要在教會範圍.
在經文的第一段裡.
記述了不同的主角.
故事的背景就是猶大國.
以色列國和亞蘭國.
這三國在中間有很大的互動的關係.
我們知道以色列國是由聯合王國之後.
北國,南國,猶大.
這個就是在大外王朝之後.
他們分裂之後.
變成了兩國.
亞蘭國其實在以色列北面的地方.
大約是現在北面的敘利亞和黎巴嫩.
那裡就叫亞蘭國.
當時亞蘭國的勢力是比較大的.
因為南北國分裂之後.
北國面對強敵.
就是亞蘭國.
再過一些就是亞述或巴比倫.
所以王國都是積弱的.
所以那些皇帝其實都是戰戰兢兢的.
第一他又拜偶像.
第二他又不跟隨神的吩咐去做.

$^{121}$所以中間有很多互動.
很多的角力.
今天講到的主角.
就是亞蘭王有一個元帥.
這個元帥叫乃曼.
這個元帥有什麼特別呢?.
在聖經形容他.
他是受王的重用的.
位高權重.
還有神是使他.
直去使亞蘭得勝.
神是使用他的.
亞蘭得勝就以色列他們.
神是使用他們來管教以色列.
他又是一個大能的勇士.
不過可惜.
他有什麼問題呢?.
長大麻風.
這個長大麻風未必像我們所想像.
後期的大麻風是腐爛的.
要隔離的.
他未必是隔離的那一種.
可能是類似我們現在所說的.
嚴重的皮膚病.
或者是牛皮癬之類.
他也很不舒服.
很痛苦的.
我相信我們認識過這方面的朋友.
你知道他很困擾的.
所以他就有這種大麻風的問題.
不過他沒有被隔離.
所以他仍然會跟人接觸.
這是第一個主角.
就是亞蘭王的元帥乃曼.
第二個主角.
大家都知道.
他是一個小女子.
他沒有名字.
現在找也找不到.
沒有記載.

$^{161}$因為通常聖經裡列出有名字的.
可能他也有意思.
但這個小女子只有小女子三個字.
他是被擄到亞蘭國.
他服侍亞蘭王的奶孖妻子.
因為以色列國.
以色列王沒有辦法保護好他們的百姓.
所以他被擄離開了.
去到敵國的地方.
離開家園和家人.
被迫服侍亞蘭人和奶孖將軍.
服侍他的妻子.
不過小女子的態度.
心態其實很單純.
因為在故事裡.
剛才讀到.
他在第五章第三節.
他說巴不得主人去見撒瑪利亞的先知.
必能治好他的大麻風.
這句小女子的說話.
他從心而發.
但產生了蝴蝶效應的結果.
蝴蝶效應.
蝴蝶不要以為是小小的翅膀.
一拍.
拍下拍下.
影響會越來越擴散出去.
擴散到很遠很遠很大的地方.
蝴蝶效應是形容這個.
由很小影響到很大.
他一句說話就是這樣.
不過這個原文的直譯.
他說我的主人.
如果可以去到撒瑪利亞先知的面前.
他就能夠除去他的大麻風.
如果廣東話.
我想大家就會投入一點.
廣東話會怎麼說呢.
如果我的主人.
可以去到撒瑪利亞的先知的話.

$^{201}$他的大麻風就一定可以好起來.
是不是全神一點.
廣東話又有些用詞.
有些感情投入去.
我想他說的時候.
可能也是這樣.
帶著那份感情.
在這裡.
小女子一句說話很簡單.
但我有個問題.
大家幫忙想一想.
她的信心從哪裡來的呢.
如果說她自己的國家是積弱.
經常被外敵欺負.
當然剛才說到.
神是籍其他國家的首就.
教訓以色列.
但她也被人欺負.
現在她在敵國.
敵人的國家裡面.
做了敵人將軍夫人的奴僕使女.
她最應該有的話.
是甚麼呢.
她可能會問.
耶和華在哪裡呀.
她的神在哪裡呀.
我的神在哪裡呀.
或者.
如果她是很有民族主義精神的.
應該不單說話.
又有些行動.
就是為國家效力.
取將軍的性命.
報被擄之仇.
這兩樣都可能有機會.
不過這個小女子通通都沒有做.
她心裡面很單純.
看事情.
不像以色列王的陳詞.
她看到的不是殺將軍的機會.

$^{241}$而是救將軍的機會.
我相信.
如果她知道自己的話.
帶來這麼大的影響.
迴響和效果的話.
她也會被自己的話震撼.
不過她從心裡面.
一段從心而發的話.
就帶著感情去說.
說就這麼輕.
但是意義就很重.
她的話是超越了民族和民族之間.
國家和國家之間.
友邦和敵國之間.
那一條界線.
指引了她的主人.
找到一條得救得醫治的路.
這個時候的一句.
撒瑪利亞的先知.
一定可以醫好她的.
這句話.
她的主母就相信了.
跟她老公將軍說.
將軍聽到之後.
就有行動了.
有行動就是.
不過你要想清楚.
將軍聽了之後.
聽了這句話.
她有沒有探聽一下情報.
有沒有查清楚事實.
她就報請皇上.
批准她保外就醫.
保證她會回來的.
她不是頭敵半邊.
這個保外就醫.
剛才這個問題.
她的信心從哪裡來.
大家有沒有答案.
如果撒瑪利亞的先知是假先知的話.

$^{281}$如果乃萬沒有得到醫治的話.
這個小女子的下場會是怎樣.
如果小女子的性格就好像最近.
前陣子有電影.
有很多角色.
有阿優,阿仇,阿嬌那些.
心中打架那些.
又很焦慮.
我想她一定會想.
多一事不如少一事.
免得找事.
真的不敢說.
醫不好會不會回來殺了我.
不過再想一想.
她的信心從哪裡來.
另外一個角色.
另一方面剛才大家也讀到.
以色列王的反應是怎樣.
相對於這位小女子.
信心的說話.
以色列王的反應.
我們就感受到.
他不單不知道.
仍然有可以醫治人的先知.
同時也不相信這位先知.
就在自己的國境之內.
因為他說.
我豈是神.
能使人死使人活.
這個人就打發來.
叫我治好他的大麻風.
他一定來尋釁攻擊我.
找方法找渣.
找我的藉口來攻擊.
當時其實亞蘭王.
是批准乃曼去的.
他又寫了這封信.
所以他才說.
有一封信帶到給以色列王.
不過.

$^{321}$他在這裡.
是用吩咐的語氣.
大家在經文裡面可以看.
我差我的神伯乃曼到你那裡.
你要除去他的大麻風.
這是第六節.
這是吩咐他.
請你快點做事.
你要除去他的大麻風.
所以以色列王的反應就很大.
撕裂了衣服.
當時人們為何會撕裂衣服.
就是在極其哀動的時候.
才會有這個反應.
或者在人離世.
或者面對極大的危難的時候.
才會這樣做.
當時以色列王就認為亞蘭王.
找渣.
借故攻擊.
不過至少在他的回應裡面.
肯定了有兩件事.
第一.
只有神能夠使人死.
使人活.
他都不行.
只有神可以.
第二.
他對於乃曼的大麻風.
是完全沒有辦法的.
做不到任何的事.
至少我們肯定了這件事.
兩件事.
身處在大時代之中.
我們現在都處於這個處境之中.
無論你在什麼崗位.
面對什麼情況.
如果我們在這幾個角色之中.
我們會看到.
這些大時代裡面的大人物也好.

$^{361}$小人物也好.
其實他面對的挑戰.
都是信心的挑戰.
是一個信心的挑戰.
備受重用的將軍.
他在疾病的折磨之下.
就要聽婢女之言.
向敵國尋找先知神醫.
都要有信心的.
誰知道他會不會作弄我.
被舉足輕重的亞蘭王.
願意要求一位他輕視的以色列王.
為自己的神僕找醫治之法.
雖然他的信是吩咐他做.
但他也要吩咐地求他.
你快點醫好他.
信心在哪裡.
昏庸的以色列王.
在強國的壓迫之下.
竟然忘記了.
那位能夠使人死,使人活的上帝.
就在他們的中間.
最後這位以色列國的小女子.
角色卑微,名字也沒有.
對話也不多,只有一句話.
她就被迫離開寧鄉別井.
在別井的異地過著三等的僕人的生活.
其實身處在大時代之中.
我們也同樣.
面對信心的考驗.
這是信心最佳的市場.
我們有多大的信心.
來跟隨神.
現在其實我們知道.
剛才說過在環境之中.
我們看到我們有很大的無力感.
無助.
無論在經濟,全球的政治.
中美也好,兩岸也好.
俄烏也好,中東也好.

$^{401}$其實香港以前.
我們也有一個比較好的環境.
我們可以左右逢源.
很多優惠政策都是因為這樣.
但現在似乎我們面對很多困難.
金飛色彼,種種的挑戰都是前所未有.
不過在這樣的大時代之中.
我們的信心怎樣幫助我們.
在生活上和信仰上.
面對各種的挑戰.
信心是很重要.
沒有了這份信心.
我們可能會驚惶失措.
不知道怎樣做,忘記了上帝.
但我希望我們的信心.
好像小女子一句簡單的話.
就顯示出她信仰的堅實.
第二個方面.
賢福者用生命來見證.
述說,福音.
接下來的八到十節.
以利沙聽到這個消息.
神人以利沙聽見以色列王撕裂衣服.
他打發人去見王.
你為什麼撕裂衣服呢?.
何時乃人到我這裡來?.
他就知道以色列中有善之料.
於是乃人就帶著車馬到以利沙的家站門前.
以利沙聽到.
他就派人去見以色列王.
其實以利沙這句話也有點諷刺.
現在不是阿難王和乃人不知道以色列中有善之料.
是誰不知道呢?.
是以色列王不知道.
王上你不知道而已.
現在我告訴你.
所以叫僕人告訴你.
叫人家知道以色列國有才知道.
是不是很諷刺呢?.
結果乃人就帶著車馬到以利沙門前.

$^{441}$在門前等.
剛才也沒有提到.
他帶了誠意金來醫藥費.
但醫藥費也很厲害.
剛才我們說到醫藥費.
保外就醫的醫療費很厲害.
譬如去瑞士那些.
但醫療費是天價的.
他帶了多少?.
他帶了十卡連德的銀紙.
金紙六千射克勒.
醫療費十套,很貴重的.
醫療費的誠意十足.
可以說是超級誠意.
因為銀紙十卡連德有多少重量呢?.
三百四十公斤.
公斤,不是公兩.
不是公克.
三百四十公斤.
金紙六千射克勒.
就是六十八公斤.
三十個卡連德的銀紙.
六千射克勒就等於三十個卡連德的銀紙.
還有一些貴重的衣服.
但要提醒大家.
亞蘭王的父親當時用了兩個卡連德銀紙.
就把整個薩馬里亞城買下了.
兩個卡連德,現在只帶十個卡連德.
這條數真的超級誠意.
如果以現在的價值來說.
金銀接近五千萬.
在當時和現在都是天文數字.
這麼有誠意.
當然是有期望的.
這麼厚禮.
他期望的是禮遇.
公開歡迎大家光臨.
多謝大家來到我們節目.
當然是希望有這樣的禮遇.
畢竟是亞蘭王的元帥.

$^{481}$就算沒有國家級的歡迎儀式.
至少先知本人都應該走出來.
幫我握握手,恭迎我來到.
好了,有沒有?.
沒有啊.
以利沙做什麼呢?.
就打發一個使者.
跟乃瑪說.
你去約旦河中.
沐浴七回.
你的浴就必復完.
而得潔淨.
什麼都沒有啊.
只是指示他去洗澡.
沐浴七回.
得潔淨.
乃瑪當然是.
什麼歡迎儀式都沒有.
帶著這麼厚禮.
連見都不見我.
無名火起三千丈.
他的反應是怎樣呢?.
乃瑪發怒走.
我想他必定出來見我.
暫著求告夜.
我話他的神的名.
在幻處以上搖手.
大馬士革的河阿巴那和.
以發起不比以色列的一切水更好嗎?.
我乃女沐浴不得潔淨嗎?.
於是氣憤憤地轉身去了.
他心上我這麼有誠意.
好歹你都跟我做場法事.
求告夜我話神.
然後用仙子的手在幻處搖一搖.
這樣就除去.
至少有一場法事.
有東西看.
什麼都沒得看.
如果只是洗澡的話.

$^{521}$就不用到水淵約旦河沖涼.
大馬士革有兩條河.
更好.
用現在的比喻.
廣東人有幾條江.
有東江.
大家在喝的水.
有西江.
有珠江.
如果要沖涼.
為什麼要到城門河沖涼呢?.
就是這個意思.
這麼多好水的河.
所以他要走.
氣憤憤地走了.
不過乃萬一埋伏人就知道.
勸他不要衝動.
他一埋伏人就說.
先知若吩咐你做一件大事.
你氣不做嗎?.
何況說你去沐浴而得潔淨呢?.
他又說了這句話.
反正都來了.
那就沖涼吧.
於是乃萬就下去.
照著神人的話.
約旦河沐浴七回.
他的肉就恢復.
好像小孩子的肉.
他就得潔淨.
他埋伏人很厲害.
他就說勸他.
大家都知道.
當他真的得到醫治的時候.
他就能夠看見.
見證原來這位真是先知.
乃萬得到醫治之後.
他就知道.
以利沙是神人.
是真的先知.

$^{561}$他怎麼說呢?.
乃萬帶著一切跟從他的人.
回到神人那裡.
站在他面前說.
如今我知道.
除了以色列之外.
普天下沒有神.
他宣認了.
除了耶和華說.
神是普天下的神之外.
沒有其他的神.
他找不到.
其他的神.
因為親身經歷了神跡醫治.
所以他就獻上禮物給神人.
以利沙怎麼回應呢?.
他就求他.
多少塔年都可以?.
十塔年都可以.
還有六千那些.
還有一些很漂亮的衣服.
你收了它吧.
我這麼有誠意.
以利沙說.
我指著所事奉永生的耶和華起誓.
我必收.
乃萬再求他.
你收吧.
好像以前的老人家一樣.
你等我給錢.
等我給錢.
他就不收.
乃萬再三請求.
他都不收.
他保持了先知應有的氣節.
因為他不像那些假先知.
在五章尾曾經出現的.
或者一會兒我們會說到.
他的僕人.
記哈西那樣.

$^{601}$他不是收錢做事的先知.
就好像文天祥所言.
是窮折乃見.
一日垂丹青.
有骨氣.
我不是收錢做事.
他就沒有收禮物.
不過乃萬有一個請求.
知不知道他請求什麼.
他沒有收錢.
請求先知去說這件事.
這段說話是我最感動的一段說話.
也是值得我們來到福音賢夫工作.
我們要留意.
我們一起讀這段乃萬說.
紅色字之後.
整段到以利沙那裡.
乃萬說.
你若不肯實.
請將兩類似陀的土賜給僕人.
從今以後.
僕人必不再將凡濟或平安濟獻與別神.
只獻給耶和華.
惟有一件事.
願耶和華要恕你僕人.
我主人進臨門廟求拜的時候.
我用手參他在臨門廟.
我也屈身.
我在臨門廟屈身的者.
願耶和華要恕我.
以利沙對他說.
你可以平平安安地回去.
乃萬他的生命改變了.
因為經歷了神跡醫治之後.
他完全不一樣.
才不收那些禮物.
但是他說.
你可不可以將兩類似陀的土賜給他.
他要將聖地的泥土帶回國.
預備建祭壇.

$^{641}$今後他只向耶和華獻祭.
不會向別神獻祭.
不過他有一件事求求耶和華.
他會陪他的王去廟裡拜其他偶像.
他是王的左右手.
在公務上他無法推遲.
不能不陪他.
所以王去臨門廟拜的時候.
他也要扶一扶.
要屈身.
好,屈身.
就像向偶像屈身敬拜.
他說我都無辦法.
因為王上都屈身.
我又要扶著他.
所以求耶和華要恕我做這件事.
他事先說了.
他知道這件事.
對於敬拜耶和華是真神的那一個.
是有抵觸的.
他自己知道.
所以他事先告訴以利沙.
他有這樣的事情.
乃曼他知道.
如果他這樣屈身.
就是違背了他之前的承諾.
只是敬拜耶和華.
不過他沒有辦法不做.
職責所在.
那以利沙的答案是怎樣呢.
他的答案是否像某些宗教塔利班那樣.
你不可以.
你要在眾人面前見證你只是敬拜耶和華.
又或者.
有些人可能會很嚴肅.
你已經是敬拜真神的人.
你不可以再為那些敬拜假神或無神論的人服務.
辭工.
沒有啊.
以利沙也不是這樣說.

$^{681}$他只是說.
可以平平安安地回去.
這個意思就是承認.
以利沙和乃曼和耶和華之間.
建立了一個約的關係.
這個約的關係.
就是在不得已的這件事上.
你平平安安地.
來到去做.
即是乃曼的請求.
獲得了答應.
大家有沒有看到這裡呢.
我們就差不多了.
其實在整個故事裡面.
有三個主角.
剛才說的這三個主角.
他們都是用他們的生命的表現.
去成為傳福音的人.
燕福的使者.
使人認識普天下的神耶和華.
第一個是哪位呢.
就是他說的話產生了蝴蝶效應的.
以色列的小女子.
我們稱之為信心的燕福者.
她說話雖然不多.
沒有人知道她的名字.
但是用信心默默地.
指引人尋找真實.
認識大能者.
得到生命並且得著豐盛的生命.
她的一句話很簡單.
可能你們也會說一句很簡單的話.
如果你有什麼問題.
你來找真牧師就好了.
他一定可以幫到你.
我曾經和醫院的弟兄姊妹.
醫護人員來分享.
有時候他們在病房裡.
真的不能夠傳福音.
但是如果有人有需要的時候.

$^{721}$他都可以像這位小女子那樣.
說:你不如試試找我們的遠牧吧.
他一定可以幫到你的.
是不是.
有時候我們不能夠酸之於口.
但是我們總能夠帶給人幫助.
可能你一句簡單的話.
成為了蝴蝶效應的結果.
相比之下.
那位勳勇的敬拜偶像的以色列王.
就成為了強烈的對比.
以色列國雖然有寶貝放在雅氣裡.
王不但不能帶領百姓認識神.
連他自己也不知道有神同在.
難怪受到亞蘭,亞述,巴比倫等烈強欺壓.
所以信心的賢福者.
不在於我說是不是全套的福音.
而是有沒有信心將人帶到神的面前.
指引他找到得救的路.
第二位就是.
不是收錢做事的神人以利沙.
我們稱他為使命的賢福者.
他沒有因為世俗的金錢名利.
失去了氣節,使命和神的能力.
表明耶和華是大能的神.
藉著神跡來見證神與人同在.
同時又看到生命改變之後.
那種確信和無奈.
以祝福代替了責備.
以同行代替了割席.
我相信今天我們有很多服事主的弟兄姊妹僕人.
我們都是追求.
我們不是收錢做事的.
你們很多服事的弟兄姊妹.
都不是收錢做事的.
我們是希望能夠在服事當中.
使人認識真神.
這種氣節不是在於我們得到多少世俗的回報.
反觀.
因為受到世俗的影響和貪腐的僕人基哈西.

$^{761}$因為今天沒有時間講,大家可以看經文.
他收了兩他年德.
真的可以買起整個三馬利亞城.
他的禮物就受到奏詛.
不單止不能購買橄欖圓,葡萄圓,牛羊,瀑皮.
還得了大麻風.
付出了貪財的極大代價.
今天我們每人都可以成為使命的燃福者.
不是因為物質的緣故.
我們將人帶到神的面前.
而是神差遣我們.
給我們這個使命.
來帶領人認識耶穌.
第三個.
就是經歷了神職醫治的亞蘭將軍乃曼.
我們稱之為德恩的燃福者.
乃曼由高傲的元帥.
改變成為謙卑的認信者.
由信心的觀眾改變成為信心的參與者.
由敬拜多神改變成為敬拜獨一真神.
他雖然在職場上受盡各種限制.
或許未能像眾人所願.
拋開一切顧慮.
勇敢地表達信仰.
但卻發自內心地承認他的不足求主饒恕.
其實在他的身上.
我相信今天我們的弟兄姊妹.
很多弟兄姊妹.
都在不同的崗位上.
受到我們的限制.
很想將福音告訴別人.
但是你卻真的有事不行.
很想做到福音的見證.
但是你卻偏偏人在江湖.
不過我們知道.
我們仍然要謙卑在神的面前.
神仍然有恩典.
我們的信心沒有動搖.
我們的信仰沒有動搖.
我們仍然做一個忠實的生命見證者.

$^{801}$服務的耶穌仍然在我們的身上.
讓人看到我們的改變.
看到我們的不一樣.
弟兄姊妹聽完這三位.
在大時代之中.
他們用自己的生命來做福音見證.
不知道你是屬於怎樣.
不過無論你是屬於哪一種.
我希望你們不是那些.
身在福中不知福的.
而是賢福者.
賢氣的賢.
不要給我那麼多福氣.
領受了很多福氣.
然後你去延伸出去.
讓人得到福氣.
最後我想做一個我的經歷.
在最後的講道.
我最近有一個很特別的經歷.
因為傳道的牙齒不是那麼好.
牙齒不好當然要看.
但是傳道的困擾.
在香港看牙醫.
香港牙醫的質素很高.
但傳道的負擔能力有限.
所以有人介紹我去深圳.
有一位牙醫.
都是輾轉介紹.
我去到的時候認識了.
這裡是牙醫的一個很大的醫療中心.
他叫曾醫生.
和曾牧師一樣.
曾牧師也認識的.
不過這位曾醫生.
他爺爺是牧師.
他爸爸也是牙醫.
我去到的時候去看我的情況.
他說黃大兄.
他叫黃先生.
你的牙齒真的有些問題.

$^{841}$有需要.
前前後後可能要植五隻牙.
五隻腳.
五隻腳而已.
還有些牙齒.
植五隻牙.
他說可以分階段.
看個人負擔能力.
他就和我說.
分階段.
最初我先做了兩隻.
他就和我說一個見證.
他的行醫理念.
我覺得很認同.
他說我要在你身上賺最多的錢.
賺最多的錢.
我怕了.
所費不菲.
但是.
我又不要你付出很多錢.
那是怎樣呢?.
又要賺最多的錢.
但又不要你付出錢.
原來是這樣.
他給我最合理和最高質素的牙醫服務.
我得到這個服務.
我會告訴我身邊的人.
你去找他吧.
他真的很好.
在我的介紹裡.
他就可以幫助不同的人.
在我身上.
他就賺最多的錢.
因為有介紹.
但他又不想我付出多錢.
結果是怎樣呢?.
我做第一次的時候有兩隻.
第二次的時候應該做三隻.
那天我去到的時候.
我跟醫生說.

$^{881}$醫生,我先做一隻.
因為負擔能力的問題.
我先做一隻.
他就準備好東西.
院長就親自幫我做.
王生.
今天就幫你做三隻.
痛就一次過痛.
你先給一隻錢.
如果你有能力有錢的時候.
你再給那兩隻.
我心想他會不會有事呢?.
這樣行不行呢?.
我覺得他真的不是為了賺錢.
來行醫.
還有他們那裡也有成立一個弱勢社群的基金.
如果真的有需要的弱勢社群.
我們也有牧者知道.
會帶他去告訴他.
幫助他.
可以得到很低廉的牙醫服務.
我們也見證著他真的用這樣的方式.
來行醫.
使那些人的牙齒獲得新生.
他可以吃東西.
牙齒有問題吃不到東西.
是最痛苦的.
我自己也經歷過.
院長雖然不會直接跟人講福音.
不過在他的牙醫中心.
你會見到有很多跟信仰有關的裝飾.
讓人看到就知道.
在深圳這個地區.
有這樣的信仰宣告.
不是容易的.
我相信他也會受壓力.
求助.
結果.
這段消息我也在我的教會有分享.
真的有很多人.

$^{921}$完了之後就問我.
傳道說你可不可以給我地址.
給我資料.
有需要去找他.
原來真的有很多人有需要.
真醫生的見證不是他獲得生意.
而是他真的願意付出.
去為主做見證.
很多北區的街坊.
可能你問起也有人去過.
是不是.
弟兄姊妹.
作為一位元福者.
我們用自己的生命來見證主.
其實聖靈也會引領我們.
用最合適的方式來述說福音.
不知你用什麼方法.
來講述福音.
但你願不願意.
無論什麼環境.
用什麼方式.
你願意去見證耶穌嗎.
願不願意.
你願意跟弟兄姊妹分享一下.
你是怎樣去見證到.
認識到.
體會到.
經歷到耶穌嗎.
人家可以透過你的身上.
認識這位主.
求主幫助我們.
使我們能夠在這個大時代裡.
仍然憑著信心.
我們去為福音.
作美好的見證.
我們現在同心同意.
天主上對我們多謝讚美你.
願你今天透過你的說話.
提醒我們.
無論在什麼艱難的環境之中.

$^{961}$我們都可以見證你的福音.
哪怕像小女子一句這麼簡單的話.
又或者像乃萬元帥一樣.
經歷你神跡的醫治.
又或者成為你使命的僕人.
被你所差遣.
真的去使人認識.
要畫獨一的主.
無論什麼方式.
主啊,你都在我們心裡.
帶領我們.
讓我們整個生命.
我們能夠彰顯耶穌基督.
求主幫助.
我們在你面前的祈禱.
奉耶穌基督的名祈求.
阿們.
祝福大家.
\newpage



\section{馬太福音 25:31-46}
\label{sec:pWTRHiVMfgo}
\textbf{2025.09.14《綿羊定山羊》:馬太福音25\_31-46 (和修版) 講員 :盧思源姑娘}
\newline
\newline
連結: \href{https://youtube.com/watch?v=pWTRHiVMfgo}{\texttt{https://youtube.com/watch?v=pWTRHiVMfgo}} ~~~~ 語音日期: 2025-09-20
\newline
\newline
\hyperref[sec:lUBiAGHTBWo]{\small{< < < PREV SERMON < < <}}
~
\hyperref[sec:index_chronic]{\small{[返順時目]}}
~
\hyperref[sec:index_scriptual]{\small{[返順卷目]}}
~
\hyperref[sec:pPe9_H3XXfI]{\small{> > > NEXT SERMON > > >}}
\newline
\newline
馬太福音 25:31-46
\newline
\begin{longtable}{cl}
\hline
\hline
章節 & 經文 (和合本修訂版)\\
\hline
25:31 & \begin{tabularx}{0.7\textwidth}{X} 「當人子在他榮耀裡，同著眾天使來臨的時候，要坐在他榮耀的寶座上。 \end{tabularx} \\ \\ \relax
25:32 & \begin{tabularx}{0.7\textwidth}{X} 萬民都要聚集在他面前。他要把他們分別出來，好像牧人分別綿羊、山羊一般， \end{tabularx} \\ \\ \relax
25:33 & \begin{tabularx}{0.7\textwidth}{X} 把綿羊安置在右邊，山羊在左邊。 \end{tabularx} \\ \\ \relax
25:34 & \begin{tabularx}{0.7\textwidth}{X} 於是王要向他右邊的說：『你們這蒙我父賜福的，可來承受那創世以來為你們所預備的國。 \end{tabularx} \\ \\ \relax
25:35 & \begin{tabularx}{0.7\textwidth}{X} 因為我餓了，你們給我吃；渴了，你們給我喝；我流浪在外，你們留我住； \end{tabularx} \\ \\ \relax
25:36 & \begin{tabularx}{0.7\textwidth}{X} 我赤身露體，你們給我穿；我病了，你們看顧我；我在監獄裡，你們來看我。』 \end{tabularx} \\ \\ \relax
25:37 & \begin{tabularx}{0.7\textwidth}{X} 義人就回答：『主啊，我們甚麼時候見你餓了，給你吃；渴了，給你喝？ \end{tabularx} \\ \\ \relax
25:38 & \begin{tabularx}{0.7\textwidth}{X} 甚麼時候見你流浪在外，留你住；或是赤身露體，給你穿？ \end{tabularx} \\ \\ \relax
25:39 & \begin{tabularx}{0.7\textwidth}{X} 又甚麼時候見你病了，或是在監獄裡，來看你呢？』 \end{tabularx} \\ \\ \relax
25:40 & \begin{tabularx}{0.7\textwidth}{X} 王回答他們說：『我實在告訴你們，這些事你們做在我弟兄中一個最小的身上，就是做在我身上了。』 \end{tabularx} \\ \\ \relax
25:41 & \begin{tabularx}{0.7\textwidth}{X} 「王又要向那左邊的說：『你們這被詛咒的人，離開我！進入那為魔鬼和他的使者所預備的永火裡去！ \end{tabularx} \\ \\ \relax
25:42 & \begin{tabularx}{0.7\textwidth}{X} 因為我餓了，你們沒有給我吃；渴了，你們沒有給我喝； \end{tabularx} \\ \\ \relax
25:43 & \begin{tabularx}{0.7\textwidth}{X} 我流浪在外，你們沒有留我住；我赤身露體，你們沒有給我穿；我病了，我在監獄裡，你們沒有來看顧我。』 \end{tabularx} \\ \\ \relax
25:44 & \begin{tabularx}{0.7\textwidth}{X} 他們也要回答：『主啊，我們甚麼時候見你餓了，或渴了，或流浪在外，或赤身露體，或病了，或在監獄裡，沒有伺候你呢？』 \end{tabularx} \\ \\ \relax
25:45 & \begin{tabularx}{0.7\textwidth}{X} 王要回答：『我實在告訴你們，這些事你們沒有做在任何一個最小的弟兄身上，就是沒有做在我身上了。』 \end{tabularx} \\ \\ \relax
25:46 & \begin{tabularx}{0.7\textwidth}{X} 這些人要往永刑裡去；那些義人要往永生裡去。」 \end{tabularx} \\ \\
[1ex]
\hline
\hline
\end{longtable}
$^{1}$弟兄姊妹平安.
感恩來到紅人堂中分享上帝的話.
願意在開始之前我們都有個祈禱.
可以回到你的殿裡一同敬拜你.
當我們一瞬間要打開上帝你自己的話的時候.
願你的靈與我們同在.
願真理的聖靈帶領我們進入上帝你自己的話裡.
以致我們被主你自己的話親自去得著.
讓我們真的聽到你要對我們眾人所說的話.
我們恭敬的將以下時間交上.
奉主耶穌基督名字祈求.
阿們.
弟兄姊妹.
不知道大家有沒有認識一些Twins.
很有趣吧.
有些Twins是雙胞胎.
有些叫做同卵雙生.
樣子就很像.
有些就叫做異卵雙生.
可能外表沒那麼像.
這些照片大部分都是identical twins.
多數他們很相似.
我以前讀書也認識一對姐妹.
她們都是雙胞胎.
初初認識她的時候.
其實也不太能分辨誰是大誰是小.
很多時候都認錯.
但當認識久了.
就知道原來她們有些細微的地方是不同的.
開始分辨到誰是姐姐.
誰是妹妹.
回到今天的經文.
今天我講到的題目是綿羊和山羊.
大家能分辨到哪一隻是綿羊哪一隻是山羊嗎.
右邊是哪一隻.
右邊是綿羊.
左邊是山羊.
大家都見過.
其實在外形上.
綿羊的體型就肥肥的.

$^{41}$肥一點.
毛又厚一點.
角有些彎位.
彎彎曲曲的.
尾巴又短一點.
而山羊相反.
普遍就比較大.
也有點像綿羊.
但看上去瘦一點.
毛又長一點.
頸部感覺短一點.
還有執羊揹鬚.
大家見到山羊就有執羊揹鬚.
如果我們就這樣以外形去看.
雖然剛才說綿羊和山羊都有點相似.
但其實綿羊和山羊都不像剛才的雙胞胎那麼相似.
為什麼耶穌要在這段經文裡.
專門提到要將綿羊和山羊分開呢.
原來在昔日猶太人裡.
他們通常在白天.
就是將綿羊和山羊一起放牧.
一起去吃草.
但到了晚上.
因為山羊的禦寒能力比較低.
所以他們晚上要收起山羊.
和綿羊分開.
特別是在冬天的晚上.
當時巴勒斯坦地區的溫差很大.
晝夜的溫差很大.
所以晚上會將他們分開.
而且在性格裡.
山羊比較喜歡撞.
綿羊相對地比較溫馴.
雖然他們是同類.
但嚴格來說.
又不是放在一起.
所以耶穌用分別綿羊和山羊.
去做了這個比喻.
如果我們打開今天的經文.
鼓勵大家如果有聖經也可以去翻譯一下.

$^{81}$其實今天講的整段經文.
是由二十四章第一節開始講的.
當時二十四章一節的時候.
經文記載耶穌出了聖殿.
他和門徒預言聖殿將會被毀.
而且我們一路讀下去整章的二十四章.
我們看到當時耶穌在橄欖山上.
和門徒講論有關末世的預兆.
當中提到末世的時間.
民要攻打民.
國要攻打國.
多處必有饑荒地震.
而且有很多假先知去迷惑人.
所以耶穌在當中提醒門徒要警醒.
免得入了迷惑.
我們也知道.
但日子的時辰是沒有人知道的.
接著到今天.
我們揭開二十五章的時候去看聖經.
我們剛才讀三十一至四十六節.
但在一至三十節裡.
我們看到耶穌繼續講.
他用了兩個比喻.
去比喻天國.
我想大家都熟悉的.
第一個是十個童女的比喻.
當中提到有五個就有預備好.
有五個就沒有預備.
被稱為愚拙的.
五個聰明的童女.
五個愚拙的童女.
聰明的就預備好遊.
等新郎回來.
五個愚拙的就沒有預備好.
所以當新郎回來的時候.
他們才去買油.
最後我們知道.
他進不了筵席.
最後被拒於門外.
而另一個比喻.

$^{121}$其實就是一個才幹的比喻.
亦提到耶穌說.
天國就好像比喻主人去愛國.
之後他叫僕人來.
因為他要去愛國.
就要將產業分給三個僕人.
有些是五千.
有個是二千.
有個是一千.
我們在經文裡看到.
五千二千的.
他就是將主人給他的錢.
再賺到五千二千回來.
相反一千的.
他就是將一千埋在地裡.
當主人回來的時候.
五千二千的就給主人稱讚.
而一千的主人就稱他為.
又惡又懶的僕人.
趕了他出去.
要在外面哀哭切齒.
如果我們這樣去讀聖經的時候.
我們看到這兩個比喻.
雖然所說的都有不同.
但兩個比喻都帶出一個重點.
就是無論當新郎或者主人回來的時候.
我們準備好了嗎?.
所以聖經告訴我們.
耶穌是會再回來的.
而耶穌第一次來到這個世界的時候.
他是以一個很謙卑.
很卑微的方式來到這個世界.
他以人的樣式降生在馬槽裡.
耶穌第一次回來的時候.
是為我們人類帶來一份救恩.
但當他第二次回來的時候.
今天的經文讓我們知道.
耶穌是以榮耀的方式再來的.
耶穌說當人子在他榮耀裡.
同著眾天使來臨的時候.

$^{161}$要坐在他榮耀的寶座上.
而且我們看回馬太福音24章.
前一點點裡30至31節.
也更加提到.
那時人子的預兆要顯在天上.
地上的萬族都要哀哭.
他們要看見人子帶著能力和大榮耀.
駕駛天上的雲降來臨.
他要參見天使.
用大聲的號筒.
從四方從天者邊直到天那邊.
召集他的選民.
如果我們這樣看的時候.
我們看到耶穌再回來的時候.
他不再是以一個卑微的方式來到.
而是以一個榮耀君王的方式來到.
而這次再來的時候.
耶穌再回來的時候.
第一次回來.
剛才說是帶來救恩.
但聖經告訴我們.
耶穌再回來的時候.
是帶來審判.
而我們今天讀的經文.
有些聖經的曰本.
如果大家有看過一些聖經曰本.
給了這個段落一個主題.
它叫做論審判的日子.
沒錯.
剛才我們說.
耶穌再回來的時候.
是帶來審判.
而今天我們原來剛才讀的經文.
是讓我們知道一些.
耶穌再回來的時候.
審判的一些原則.
經文32至33節.
萬民都要聚集在他面前.
他要把他們分別出來.
好在牧人分別綿羊山羊一般.

$^{201}$把綿羊安置在右邊.
山羊在左邊.
在經文裡我們看到.
耶穌以綿羊和山羊.
去比喻了兩類人.
之前有聽過一些人解釋.
這段經文的時候.
可能很多時候會說.
是不是就是綿羊和山羊.
其實是在說信主和不信主的人.
因為經文說綿羊永生.
是不是信主的人我們就可以永生.
不信主.
就好像山羊的人.
是不是就要下去永垷的裡面.
不過如果我們讀得深入一點.
我們會問.
如果只是說信和不信.
其實可以很簡單就說.
是不是.
你信就只有永生.
不信就永垷.
耶穌不需要說這麼多的東西.
其實耶穌不需要用綿羊和山羊去比喻.
而且如果剛才讀經文裡.
我們看到經文的內容.
是在說一些行為.
內容裡面提到綿羊有做一些事.
而山羊沒有做.
假如是在說信和不信的話.
耶穌是否意味著.
我們因行為去得救.
做了就得救了.
不做就不得救了.
是不是這樣.
我們知道不是.
因為聖經在《耳窟所書》裡說.
「你們得救是本乎因,也因著信,.
這並不是出於自己,而是神所賜的,.
也不是出於行為,免得有人自誇.」.

$^{241}$我們知道我們是因信被稱為義.
並不是因為我們做了什麼好事.
所以我們去得救.
全是因為神的恩典.
所以當我們願意去相信耶穌的時間.
救恩就是伯伯給我們的.
我們就問.
耶穌在這裡說什麼呢.
綿羊和山羊究竟在說什麼呢.
當時耶穌的對象是跟誰說呢.
剛才提到我們看聖經.
這段經文一開始的時間.
是在24章1節.
耶穌當時是跟門徒說的.
如果剛才提過25章的比喻.
其實對象都是一些門徒.
所以我們相信.
耶穌在說關於審判的日子.
是跟門徒說的.
也就是跟今天我們每一個相信主的人說.
但是可能我們有人會問.
審判不是應該所有人的嗎.
為什麼只跟信徒說呢.
沒錯.
聖經告訴我們.
耶穌再回來的時候.
是審判所有人.
正如我們剛才讀使徒信經.
審判活人死人.
不過如果我們記得.
在彼得前書4章17節看到.
因為時候到了.
耶穌再回來的時候到了.
審判是要從神的家開始.
如果我們這樣理解的話.
審判是從神的家開始.
我們會明白.
為什麼耶穌會以牧羊人.
去分別綿羊和山羊去做這個比喻.
原來綿羊和山羊.

$^{281}$是在比喻今天我們在教會裡的人.
是今天我們一群基督徒.
即是耶穌再回來的時候.
祂會先審判.
今天我們說是跟從祂的人.
即是你和我.
到這裡可能有人會問.
我們一群信主的人.
剛才不是已經說了.
我們是因信稱義的嗎.
我們不是信耶穌.
就能永生嗎.
為什麼仍然有綿羊和山羊的分別呢.
而且我們看到.
綿羊和山羊的結局是不同的.
綿羊在往永生.
相反山羊就在永型裡.
我們不是應該所有信耶穌的人.
都是有永生嗎.
所以這裡說的是.
沒錯 聖經說我們信耶穌.
得救是因為我們信耶穌.
但一個很重要的重點就是.
我們這份信究竟是怎樣的信呢.
當我們口裡承認.
心裡相信.
沒錯 聖經告訴我們.
我們被稱為義.
但這只是一個起點.
因為亞國書裡也提到.
信心沒有行為是死的.
原來真實的信心.
是必然帶有行為的.
行為是信心表達的一種延續.
換句話說.
真正的信心.
不只是停留在起點的信.
真正的信心.
是必然會帶來有行為去伴隨.
但如果沒有行為的信心.

$^{321}$聖經告訴我們.
這不是真正的信心.
所以為什麼耶穌會用綿羊和山羊去比喻.
祂對當時的門徒.
也是今天告訴我們.
當祂回來的時候.
祂在審判的原則當中.
祂想跟我們說什麼呢.
我們接著看.
耶穌想說什麼.
回到經文的時候.
耶穌說當祂永耀降臨.
剛才也讀過.
祂將牧羊人分開綿羊和山羊.
綿羊在右邊.
對著綿羊說.
「王耀對右邊說.
你這蒙我父賜福的.
可來承受那創世以來為你們所預備的禍.
因為我學 我餓了 你們給我吃.
學了 你們給我喝.
我流浪在外 你們給我住.
我赤身露體 你們給我穿.
我病了 你們看顧我.
我在監獄裡 你們來看我」.
相反 左邊的.
很像的.
「王耀向左邊說.
你們這被咒詛的人.
離開我 進入那為魔鬼和他的使者所預備的永火裡去.
因為我餓了 你們沒有給我吃.
學了 你們沒有給我喝.
我流浪在外 你們沒有留我住.
我赤身露體 你們沒有給我穿.
我病了 我在監獄裡.
你們沒有來看顧我」.
剛才這兩段經文有很多重複.
簡單來說.
簡單說 綿羊就是要做到六樣東西.
就是餓了 你們給我吃.

$^{361}$學了 你們給我喝.
流浪 你們留我住.
赤身露體 你們給我穿.
病了 看顧我.
監獄裡來看我.
相反 山羊就沒有做到那些東西.
如果就這樣看.
其實 剛才說如果假如這是一個審判的原則的話.
很容易分辨.
就是你有做和沒有做.
這些行為.
有做和沒有做的分別.
很簡單.
我們之後就組織一個隊伍.
我們去找找街上有沒有流浪漢.
我們就去給他吃.
給他穿衣服.
我們組織探訪.
我們去醫院探訪.
有時我們開始監獄就去探家.
我們是不是這樣呢?.
當我們做了這些的時候.
將來見到耶穌的時候.
耶穌問起的時候.
就回答「做了 我做了」.
聖經經文是不是這麼簡單.
只是想讓我們聽說這件事呢?.
所以當我們看完兩班人的回應的時候.
我們可以更加多思想.
究竟耶穌想說什麼.
「二人就回答:.
「綿羊,他們那邊就說:.
主啊!我什麼時候見你餓,.
給你吃,給你喝,.
學了,給你學,.
什麼時候見你流浪在外,留你住,.
或是赤身露體,給你穿,.
又什麼時候見你病了,.
或是在監獄內看你呢?」.
簡單來說,他們的回應是什麼呢?.

$^{401}$其實他們都不知道自己什麼時候做了.
再看山羊,山羊那邊又怎麼回答呢?.
他們都回答:.
「主啊!我什麼時候見你餓,.
餓了,學了,或者流浪在外,.
或赤身露體,或病了,或在監獄內,.
沒有時候你呢?」.
同樣,他們都在說同一件事.
我也不知道自己什麼時候沒有做到.
總括來說,兩班人都不知道自己什麼時候做了.
也不知道自己什麼時候沒有做.
最後,黃怎樣回答這兩班人呢?.
對著綿羊的說:.
「我實在告訴你,.
這些事你做不做在我弟兄中最小的身上,.
就是要做在我身上了.」.
對山羊的說:.
「我實在告訴你們,.
這些事你們沒有做在任何一個最小的弟兄身上,.
就是沒有做在我身上.」.
為什麼我要重覆讀這段經文呢?.
就這樣看這段經文,.
雖然好像說了為什麼綿羊有做,山羊沒有做.
但是,如果你想深一層,.
我們會問三個問題.
第一個問題,剛才說.
為什麼右邊的綿羊會不知道自己做了呢?.
正常,正常自己做過什麼,大家知不知道?.
正常人都知道,是吧?.
你今天吃了什麼,大家都知道.
為什麼做了什麼,會不知道呢?.
我就在想,.
有什麼自己做了,但是我不知道呢?.
我就想,可能會不會是一些很習以為常的事情.
你都做了,不覺得的.
舉個例子,.
我相信大家每天都會刷牙,是吧?.
你刷牙,每天早上都刷牙.
我問你,你會不會記得自己刷了多少次牙?.
你當然會笑,怎麼會記呢?.

$^{441}$沒有人會計算,相信大家都不會計算.
因為你刷牙是一件很習以為常的事情.
你不會刻意自己刷了40次.
你不會這樣計算.
因為基本上你每天都會做這些事情.
而剛才這段經文,第二個問題.
就是說,那山羊為什麼又會不知道自己沒有做呢?.
這個就容易理解了.
可能一開始,牠根本沒有留意要做.
所以就沒有做.
沒有留意,或者就算留意到,看到都好像看不到的時間.
最後沒有做,也很正常,是吧?.
第三個問題,看這段經文的時候.
就會想,耶穌在這裡說.
我實在告訴你,這些事你們做在我弟兄中最少的一個身上.
就是做在我身上.
另外,對答就是,你沒有做最少的身上,就是沒有做在我身上.
那你會問,什麼是最少?.
是不是年紀最少?.
如果計算年紀,誰最少應該是在樓上,是吧?.
最少的那些.
那我們就應該上去搶著找他們吃飯.
請他們吃飯才行,是不是這樣?.
我相信聖經不是這個意思.
這個最少不是說年紀上的最少.
聖經裡面的意思是說.
地位最低,最不起眼,或者最不重要的一個.
那我們就會問,那就是哪一個?.
是不是坐在角落那些,不是很起眼.
隱藏自己的那些,我們是不是要找那個出來呢?.
然後我們同樣的,找了那個出來.
我們就去請他吃飯,請他喝東西.
做了耶穌教我們的那些事.
我們是不是這樣呢?.
我相信,這裡耶穌的意思.
祂不是說叫我們找最少的一個出來.
然後去做就行.
如果我們去看清楚.
或者我們去看其他的譯本.
這句話是這麼說的.

$^{481}$在和諧本,祂說.
你們既做在我這些弟兄中最卑微的身上.
就是做在我身上.
耶穌說,只要你做在我弟兄中最少最卑微的身上.
就是做在我身上.
我們嘗試轉一轉去想.
就是如果做在一個最少的身上.
都已經是做在耶穌身上.
那其他不是最少的人呢?.
如果做在最少的身上.
耶穌都已經算在做在祂身上.
就好像算在祂身上.
何況是其他人呢?.
如果做在那些不是最少的身上.
豈不是更加就是做在耶穌身上呢?.
大家聽得明白嗎?.
換句話說,當中的意思是.
其實無論祂是不是最少最卑微的一個.
只要當你看到任何一個有需要的人.
你去幫他.
你去幫.
你見到他有需要,你去幫他的時間.
其實就是做在耶穌的身上.
問題是我們有沒有見到我們身邊這些人呢?.
所以回到經文的時候.
我們會更加明白.
為什麼綿羊和山羊.
其實不是很清楚自己什麼時候做了什麼時候沒有做.
因為耶穌上面提到.
剛才提到的六種情況.
餓,渴,流浪在外,赤身怒體,病了,監獄.
其實除了這六種情況.
是可以有第七,第八,第九,第十種的情況.
所以耶穌不是單單說有沒有做這幾種行為這麼簡單.
相反是有做和沒有做.
背後我們作為一群信徒.
我們一群跟從主的人.
我們的生命究竟是怎樣.
在我們生命裡面.
有沒有真的活出那份愛.

$^{521}$我們有沒有留心.
留心身邊的人呢?.
有沒有關心任何一個神.
放在你身邊有需要的人呢?.
在我們的群體.
我們是否真的在實踐愛在我們當中.
當見到你身邊有需要的人.
無論他是不是最少的一個.
我們有沒有主動去關心對方.
還是好像山羊我們看不到.
或者根本沒有留意過周圍的人呢?.
弟兄姊妹.
聖經告訴我們.
最大的誡命.
我相信很多弟兄姊妹都會識背.
你要盡心盡性盡義盡力愛主你的神.
其次就是要愛人如己.
再沒有這比兩條誡命更大的了.
耶穌說最大的誡命就是愛神愛人.
但是甚麼是愛神呢?.
是否我們每星期準時回來教會.
崇拜就是等於愛神呢?.
約翰一書提到.
人若說我愛神.
卻恨他的弟兄就是說謊了.
不愛他看得見的弟兄.
就不能愛看不見的神.
愛神的也要愛弟兄.
這是我們從神所受的命令.
這裡說得很清楚.
原來我們愛神的話.
我們是要愛身邊的弟兄.
如果我們說愛神.
但不愛眼前的人.
這是說謊.
因為一個人若不愛身旁看得見的弟兄.
你如何說愛看不見的神呢?.
所以一個愛神的人.
是會愛你的弟兄.
而這份愛不只是言語或洩頭.

$^{561}$而是具體行動.
就如我們今天讀的經文一樣.
有行動的.
如果我們總結今天的經文.
讓我們知道.
原來當耶穌再回來的時候.
祂是帶來審判.
而審判的標準.
不只是單單說我們是否信祂.
正如在馬太福音7:21也說.
「不是每個稱呼我主啊主啊的人.
都能進天國.
唯有遵行我天父旨意的人才能進去」.
原來耶穌所看重的是.
我們這份信.
有沒有帶來我們生命的改變呢?.
在我們的生活裡.
我們有沒有遵守神的吩咐.
神的誡命.
我們有沒有去愛身邊的人呢?.
誰是你身邊的人呢?.
大家可以看看你的左右旁邊.
這些就是你身邊的人.
當你看到他們有需要的時間.
我們是否以行動去關心對方呢?.
還是我們視而不見呢?.
最後.
我今天的講題是.
綿羊定山羊.
有個問號.
其實這一刻.
都不知道我們自己是綿羊定山羊.
不過聖經告訴我們.
如果我們遵行神的話.
神的吩咐.
神的誡命.
我們去愛神.
愛人.
而這份愛.
不只是單單用口去說.

$^{601}$而是用我們的生命.
去活出神的愛的時候.
相信.
相信當我們真的如此行.
我們真的這樣做的時候.
神是知道的.
盼望.
即是神的話.
都成為我們彼此的提醒.
讓我們真的活出那份愛的群體.
充滿愛的群體.
而且不單止是我們.
而是我們怎樣可以將這份愛.
帶給身邊有需要的人.
我們一同祈禱.
當我們從聖經裡知道.
主耶穌你會再回來.
當你再回來的時候.
你是以一個榮耀的君王.
這樣回到這裡的世界.
而且我們知道你回來的時候.
是帶來審判.
主耶穌當我們面對審判的時候.
我們不是那種恐懼.
相反我們在聖經裡知道.
上帝你吩咐我們所做的.
我們怎樣去行.
正如剛才提到綿羊和山羊.
縱然今天我們真的不知道.
我們自己是綿羊還是山羊.
但我們知道的是.
當我們願意聽從上帝的吩咐.
我們去愛神愛人.
我們以你的愛去關心.
身邊有需要的人.
無論那些是否最少最卑微.
而是當我們看到的時候.
我們就是這樣去行的時候.
我們就是這樣去活出.
上帝你這份的愛.

$^{641}$求主你幫助我們.
將你的話放在心裡.
以致當我們將來去見主面的時候.
我們每個人都可以被主.
你自己稱讚為有良善.
有忠心的僕人.
願意你帶領著我們每一個.
多謝你聽我們的祈禱.
奉主耶穌基督明智祈求.
阿們.
\newpage



\section{使徒行傳 6:1-7:60}
\label{sec:pPe9_H3XXfI}
\textbf{2025.09.21《作主忠心的僕人》使徒行傳6\_1-7\_60 曾敏姑娘}
\newline
\newline
連結: \href{https://youtube.com/watch?v=pPe9-H3XXfI}{\texttt{https://youtube.com/watch?v=pPe9-H3XXfI}} ~~~~ 語音日期: 2025-09-23
\newline
\newline
\hyperref[sec:pWTRHiVMfgo]{\small{< < < PREV SERMON < < <}}
~
\hyperref[sec:index_chronic]{\small{[返順時目]}}
~
\hyperref[sec:index_scriptual]{\small{[返順卷目]}}
~
\hyperref[sec:WCmVYVwec7g]{\small{> > > NEXT SERMON > > >}}
\newline
\newline
使徒行傳 6:1-7:60
\newline
\begin{longtable}{cl}
\hline
\hline
章節 & 經文 (和合本修訂版)\\
\hline
6:1 & \begin{tabularx}{0.7\textwidth}{X} 那些日子，門徒增多，有說希臘話的猶太人向希伯來人發怨言，因為在日常的供給上忽略了他們的寡婦。 \end{tabularx} \\ \\ \relax
6:2 & \begin{tabularx}{0.7\textwidth}{X} 十二使徒叫眾門徒來，說：「我們撇下神的道去管理飯食，是不合宜的。 \end{tabularx} \\ \\ \relax
6:3 & \begin{tabularx}{0.7\textwidth}{X} 所以弟兄們，當從你們中間選出七個有好名聲、滿有聖靈和智慧，我們派他們管理這事。 \end{tabularx} \\ \\ \relax
6:4 & \begin{tabularx}{0.7\textwidth}{X} 至於我們，我們要專注於祈禱和傳道的事奉。」 \end{tabularx} \\ \\ \relax
6:5 & \begin{tabularx}{0.7\textwidth}{X} 這話使全會眾都喜悅，就揀選了司提反—他是一個滿有信心和聖靈的人；他們又揀選了腓利、伯羅哥羅、尼迦挪、提門、巴米拿，並皈依猶太教的安提阿人尼哥拉， \end{tabularx} \\ \\ \relax
6:6 & \begin{tabularx}{0.7\textwidth}{X} 叫他們站在使徒面前，使徒禱告後，就為他們按手。 \end{tabularx} \\ \\ \relax
6:7 & \begin{tabularx}{0.7\textwidth}{X} 神的道興旺起來；在耶路撒冷門徒數目增加得很多，也有許多祭司聽從了這信仰。 \end{tabularx} \\ \\ \relax
6:8 & \begin{tabularx}{0.7\textwidth}{X} 司提反滿有恩惠和能力，在民間行了大奇事和神蹟。 \end{tabularx} \\ \\ \relax
6:9 & \begin{tabularx}{0.7\textwidth}{X} 當時有從稱為「自由人」會堂，並古利奈、亞歷山大會堂來的人，還有些從基利家、亞細亞來的人，起來和司提反辯論。 \end{tabularx} \\ \\ \relax
6:10 & \begin{tabularx}{0.7\textwidth}{X} 司提反是以智慧和聖靈說話，眾人抵擋不住， \end{tabularx} \\ \\ \relax
6:11 & \begin{tabularx}{0.7\textwidth}{X} 就收買人來說：「我們聽見他說褻瀆摩西和神的話。」 \end{tabularx} \\ \\ \relax
6:12 & \begin{tabularx}{0.7\textwidth}{X} 他們又煽動百姓、長老和文士，就突然來捉拿他，把他帶到議會去， \end{tabularx} \\ \\ \relax
6:13 & \begin{tabularx}{0.7\textwidth}{X} 設下假見證，說：「這個人不斷地說話，侮辱神聖的地方和律法。 \end{tabularx} \\ \\ \relax
6:14 & \begin{tabularx}{0.7\textwidth}{X} 我們曾聽見他說，這拿撒勒人耶穌要毀壞這地方，也要改變摩西所交給我們的規矩。」 \end{tabularx} \\ \\ \relax
6:15 & \begin{tabularx}{0.7\textwidth}{X} 在議會裡坐著的人都定睛看他，見他的面貌好像天使的面貌。 \end{tabularx} \\ \\ \relax
7:1 & \begin{tabularx}{0.7\textwidth}{X} 大祭司說：「果真有這些事嗎？」 \end{tabularx} \\ \\ \relax
7:2 & \begin{tabularx}{0.7\textwidth}{X} 司提反說：「諸位父老弟兄請聽！從前我們的祖宗亞伯拉罕在美索不達米亞，還沒有住在哈蘭的時候，榮耀的神向他顯現， \end{tabularx} \\ \\ \relax
7:3 & \begin{tabularx}{0.7\textwidth}{X} 對他說：『你要離開本地和親族，往我所要指示你的地去。』 \end{tabularx} \\ \\ \relax
7:4 & \begin{tabularx}{0.7\textwidth}{X} 他就離開迦勒底人的地方，住在哈蘭。他父親死了以後，神使他從那裡搬到你們現在所住的地方。 \end{tabularx} \\ \\ \relax
7:5 & \begin{tabularx}{0.7\textwidth}{X} 在這裡神並沒有給他產業，連立足的地方都沒有，但應許要將這地賜給他和他的後裔為業，雖然那時他還沒有兒子。 \end{tabularx} \\ \\ \relax
7:6 & \begin{tabularx}{0.7\textwidth}{X} 神這樣說：『他的後裔必寄居外邦，那裡的人要使他們作奴隸，苦待他們四百年。』 \end{tabularx} \\ \\ \relax
7:7 & \begin{tabularx}{0.7\textwidth}{X} 神又說：『但我要懲罰使他們作奴隸的那國。以後他們要出來，在這地方事奉我。』 \end{tabularx} \\ \\ \relax
7:8 & \begin{tabularx}{0.7\textwidth}{X} 神又賜他割禮的約。於是亞伯拉罕生了以撒，在第八日給他行了割禮；後來以撒生雅各，雅各生十二位先祖。 \end{tabularx} \\ \\ \relax
7:9 & \begin{tabularx}{0.7\textwidth}{X} 「先祖嫉妒約瑟，把他賣到埃及去，神卻與他同在， \end{tabularx} \\ \\ \relax
7:10 & \begin{tabularx}{0.7\textwidth}{X} 救他脫離一切苦難，又使他在埃及王法老面前蒙恩，又有智慧。法老派他作埃及國的宰相兼管法老的全家。 \end{tabularx} \\ \\ \relax
7:11 & \begin{tabularx}{0.7\textwidth}{X} 後來全埃及和迦南遭遇饑荒和大災難，我們的祖宗絕了糧。 \end{tabularx} \\ \\ \relax
7:12 & \begin{tabularx}{0.7\textwidth}{X} 雅各聽見在埃及有糧，就打發我們的祖宗初次往那裡去。 \end{tabularx} \\ \\ \relax
7:13 & \begin{tabularx}{0.7\textwidth}{X} 第二次約瑟與兄弟們相認，法老才認識他的家族。 \end{tabularx} \\ \\ \relax
7:14 & \begin{tabularx}{0.7\textwidth}{X} 約瑟就打發人，請父親雅各和全族七十五個人都來。 \end{tabularx} \\ \\ \relax
7:15 & \begin{tabularx}{0.7\textwidth}{X} 於是雅各下了埃及，後來他和我們的祖宗都死在那裡； \end{tabularx} \\ \\ \relax
7:16 & \begin{tabularx}{0.7\textwidth}{X} 他們又被遷到示劍，葬於亞伯拉罕在示劍用銀子從哈抹子孫買來的墳墓裡。 \end{tabularx} \\ \\ \relax
7:17 & \begin{tabularx}{0.7\textwidth}{X} 「當神應許亞伯拉罕的日期將到的時候，以色列人在埃及人丁興旺， \end{tabularx} \\ \\ \relax
7:18 & \begin{tabularx}{0.7\textwidth}{X} 直到另一位不認識約瑟的王興起統治埃及。 \end{tabularx} \\ \\ \relax
7:19 & \begin{tabularx}{0.7\textwidth}{X} 他用詭計待我們的宗族，苦待我們的祖宗，強迫他們丟棄嬰孩，使嬰孩不能存活。 \end{tabularx} \\ \\ \relax
7:20 & \begin{tabularx}{0.7\textwidth}{X} 就在那時，摩西生了下來，神看為俊美，在父親家裡被撫養了三個月。 \end{tabularx} \\ \\ \relax
7:21 & \begin{tabularx}{0.7\textwidth}{X} 他被丟棄的時候，法老的女兒拾了去，當自己的兒子撫養。 \end{tabularx} \\ \\ \relax
7:22 & \begin{tabularx}{0.7\textwidth}{X} 摩西學了埃及人一切的學問，說話辦事都有才能。 \end{tabularx} \\ \\ \relax
7:23 & \begin{tabularx}{0.7\textwidth}{X} 「他到了四十歲，心中起意去看望他的弟兄以色列人。 \end{tabularx} \\ \\ \relax
7:24 & \begin{tabularx}{0.7\textwidth}{X} 他見他們中的一個人受冤屈，就庇護他，為那被壓迫的人報仇，打死了那埃及人。 \end{tabularx} \\ \\ \relax
7:25 & \begin{tabularx}{0.7\textwidth}{X} 他以為他的弟兄們必明白神是藉他的手搭救他們，他們卻不明白。 \end{tabularx} \\ \\ \relax
7:26 & \begin{tabularx}{0.7\textwidth}{X} 第二天，他遇見有人在打架，就想勸他們和好，說：『二位，你們是弟兄，為甚麼彼此欺負呢？』 \end{tabularx} \\ \\ \relax
7:27 & \begin{tabularx}{0.7\textwidth}{X} 那欺負鄰舍的人把他推開，說：『誰立你作我們的領袖和審判官呢？ \end{tabularx} \\ \\ \relax
7:28 & \begin{tabularx}{0.7\textwidth}{X} 難道你要殺我像昨天殺那埃及人一樣嗎？』 \end{tabularx} \\ \\ \relax
7:29 & \begin{tabularx}{0.7\textwidth}{X} 摩西聽見這話就逃走了，寄居於米甸地，在那裡生了兩個兒子。 \end{tabularx} \\ \\ \relax
7:30 & \begin{tabularx}{0.7\textwidth}{X} 「過了四十年，在西奈山的曠野，有一位天使在荊棘的火焰中向摩西顯現。 \end{tabularx} \\ \\ \relax
7:31 & \begin{tabularx}{0.7\textwidth}{X} 摩西見了那異象，覺得很驚訝，正往前觀看的時候，有主的聲音說： \end{tabularx} \\ \\ \relax
7:32 & \begin{tabularx}{0.7\textwidth}{X} 『我是你列祖的神，就是亞伯拉罕、以撒、雅各的神。』摩西戰戰兢兢，不敢觀看。 \end{tabularx} \\ \\ \relax
7:33 & \begin{tabularx}{0.7\textwidth}{X} 主對他說：『把你腳上的鞋脫下來，因為你所站的地方是聖地。 \end{tabularx} \\ \\ \relax
7:34 & \begin{tabularx}{0.7\textwidth}{X} 我的百姓在埃及所受的困苦，我確實看見了；他們悲嘆的聲音，我也聽見了。我下來要救他們。現在，你來，我要差你往埃及去。』 \end{tabularx} \\ \\ \relax
7:35 & \begin{tabularx}{0.7\textwidth}{X} 「這摩西就是有人曾棄絕他說『誰立你作我們的領袖和審判官』的，神卻藉那在荊棘中顯現的天使的手差派他作領袖，作解救者。 \end{tabularx} \\ \\ \relax
7:36 & \begin{tabularx}{0.7\textwidth}{X} 這人領以色列人出來，在埃及地，在紅海，在曠野的四十年間行了奇事神蹟。 \end{tabularx} \\ \\ \relax
7:37 & \begin{tabularx}{0.7\textwidth}{X} 這人是摩西，就是那曾對以色列人說『神要從你們弟兄中給你們興起一位先知像我』的。 \end{tabularx} \\ \\ \relax
7:38 & \begin{tabularx}{0.7\textwidth}{X} 這人是那曾在曠野的會眾中和西奈山上，與那對他說話的天使同在，又與我們祖宗同在的，他領受了活潑的聖言傳給我們。 \end{tabularx} \\ \\ \relax
7:39 & \begin{tabularx}{0.7\textwidth}{X} 我們的祖宗不肯聽從，反棄絕他，他們的心轉向埃及， \end{tabularx} \\ \\ \relax
7:40 & \begin{tabularx}{0.7\textwidth}{X} 對亞倫說：『你為我們造神明，在我們前面引路，因為領我們出埃及地的這個摩西，我們不知道他遭遇了甚麼事。』 \end{tabularx} \\ \\ \relax
7:41 & \begin{tabularx}{0.7\textwidth}{X} 那時，他們造了一個牛犢，又拿祭物獻給那像，為自己手所做的工作歡躍。 \end{tabularx} \\ \\ \relax
7:42 & \begin{tabularx}{0.7\textwidth}{X} 但是神轉臉不顧，任憑他們祭拜天上的日月星辰，正如先知書上所寫的：『以色列家啊，你們四十年間在曠野，何曾將犧牲和祭物獻給我？ \end{tabularx} \\ \\ \relax
7:43 & \begin{tabularx}{0.7\textwidth}{X} 你們抬著摩洛的帳幕和理番---你們神明的星，就是你們所造為要敬拜的像。因此，我要把你們遷到巴比倫外去。』 \end{tabularx} \\ \\ \relax
7:44 & \begin{tabularx}{0.7\textwidth}{X} 「我們的祖宗在曠野，有作證的會幕，是神吩咐摩西照著他所看見的樣式做的。 \end{tabularx} \\ \\ \relax
7:45 & \begin{tabularx}{0.7\textwidth}{X} 這帳幕，我們的祖宗同約書亞相繼承受了，當神在他們面前趕走外邦人的時候，他們把這帳幕搬進承受為業之地，直存到大衛的日子。 \end{tabularx} \\ \\ \relax
7:46 & \begin{tabularx}{0.7\textwidth}{X} 大衛在神面前蒙恩，祈求為雅各的家預備居所。 \end{tabularx} \\ \\ \relax
7:47 & \begin{tabularx}{0.7\textwidth}{X} 但卻是所羅門為神造成殿宇。 \end{tabularx} \\ \\ \relax
7:48 & \begin{tabularx}{0.7\textwidth}{X} 其實，至高者並不住人手所造的，就如先知所言： \end{tabularx} \\ \\ \relax
7:49 & \begin{tabularx}{0.7\textwidth}{X} 『主說：天是我的寶座，地是我的腳凳。你們要為我造怎樣的殿宇？哪裡是我安歇的地方呢？ \end{tabularx} \\ \\ \relax
7:50 & \begin{tabularx}{0.7\textwidth}{X} 這一切不都是我手所造的嗎？』 \end{tabularx} \\ \\ \relax
7:51 & \begin{tabularx}{0.7\textwidth}{X} 「你們這硬著頸項，心與耳未受割禮的人哪，時常抗拒聖靈！你們的祖宗怎樣，你們也怎樣。 \end{tabularx} \\ \\ \relax
7:52 & \begin{tabularx}{0.7\textwidth}{X} 先知中有哪一個不是受你們祖宗的迫害呢？他們把預先宣告那義者要來的人殺了。如今你們成了那義者的出賣者和兇手了。 \end{tabularx} \\ \\ \relax
7:53 & \begin{tabularx}{0.7\textwidth}{X} 你們領受了天使所傳佈的律法，竟不遵守。」 \end{tabularx} \\ \\ \relax
7:54 & \begin{tabularx}{0.7\textwidth}{X} 眾人聽見這些話，心中極其惱怒，向司提反咬牙切齒。 \end{tabularx} \\ \\ \relax
7:55 & \begin{tabularx}{0.7\textwidth}{X} 但司提反滿有聖靈，定睛望天，看見神的榮耀，又看見耶穌站在神的右邊， \end{tabularx} \\ \\ \relax
7:56 & \begin{tabularx}{0.7\textwidth}{X} 就說：「我看見天開了，人子站在神的右邊。」 \end{tabularx} \\ \\ \relax
7:57 & \begin{tabularx}{0.7\textwidth}{X} 眾人大聲喊叫，摀著耳朵，齊心衝向他， \end{tabularx} \\ \\ \relax
7:58 & \begin{tabularx}{0.7\textwidth}{X} 把他推到城外，用石頭打他。作見證的人把他們的衣裳放在一個名叫掃羅的青年腳前。 \end{tabularx} \\ \\ \relax
7:59 & \begin{tabularx}{0.7\textwidth}{X} 他們正用石頭打司提反的時候，他呼求說：「主耶穌啊，求你接納我的靈魂！」 \end{tabularx} \\ \\ \relax
7:60 & \begin{tabularx}{0.7\textwidth}{X} 然後他跪下來，大聲喊著：「主啊，不要將這罪歸於他們！」說了這話，就長眠了。 \end{tabularx} \\ \\
[1ex]
\hline
\hline
\end{longtable}
$^{1}$弟兄姊妹平安.
每個星期日能夠回到神的家.
去敬拜祂.
都是一個很大的福氣.
我們一同禱告.
仍然讓我們活在世上.
能夠親近你 敬拜你.
今天當我們聽信息之前.
求你光照我們每一個的心.
讓我們看到我們自己生命的光景.
我們有什麼地方得罪了你 得罪人.
求你也去寬恕 潔淨我們.
願意讓耶穌基督的寶血.
大大的潔淨我們.
復辟和保護我們.
願意聖靈對我們眾人說話.
在當中讓我們明白神的心意.
讓我們作出悉切的回應.
奉耶穌基督的名字祈禱.
阿們.
今天我想說的唐道裡的兩位人物.
我都非常欣賞.
一直看的時候.
雖然以前其中一個.
其中一個聖經的人物.
以前都看過他人生的經歷.
以前都欣賞.
但這次更加欣賞.
很想和大家一同認識他多一點.
作主中心的僕人.
我和你都可以朝這個目標去成長.
不只是某個人才可以成為主中心的僕人.
我們都可以.
我們看看一些很好的榜樣.
這一個在《使徒行傳》裡.
一個很重要的時刻.
我們看看在這一段時間.
在《使徒行傳》的六章裡.
這段經文提及早期教會的時間.
提及當時的門徒越來越多.

$^{41}$教會裡有提及希臘的猶太人.
向希伯來人發怨言.
這些希臘的猶太人是指什麼人呢?.
原來在當地耶路撒冷的猶太人.
他們很熟習阿難民.
也很熟習希伯來民.
當然他們也會說希臘語.
但是不在耶路撒冷的猶太人.
例如來自其他地方的猶太人.
他們會說希臘語.
未必會說阿難民.
未必會說希伯來語.
現在提及這些來自其他地方的猶太人.
他們向當地的希伯來人發怨言.
因為在日常生活的飯食供應上.
照顧上忽略了他們當中的寡婦.
大家都知道寡婦無依無靠.
更加需要教會的照顧.
所以這些來自外地的猶太人.
就有這個抱怨.
在這裡就提及到當時教會的處理非常好.
大家記得有十二個使徒.
他們是教會很重要的領袖.
其實這些領袖到了這個時候.
在教會有很重要的使命.
任務他們的職份.
他們是以什麼為重要的呢?.
就是要祈禱和傳道.
其實十二使徒可以這樣說.
當時就像傳道童工.
他們重要的不是管飯食.
最重要的是祈禱和傳道.
但是這個時候教會人數不斷增加.
是需要很多很多的照顧.
那怎麼辦呢?.
每天的飯食也要處理.
如果這些傳道童工使徒.
全部放下傳道和祈禱的事情.
只去處理飯食.
這不是神的心意.

$^{81}$所以當時他們就叫了一些門徒來.
就來到這裡有這樣的畫面.
是做什麼呢?.
他們說我們要在他們當中.
選七個有好明星,滿有聖靈和智慧的人.
派他們管理飯食.
這七個當時最早期的什麼呢?.
執事.
原來早期的執事是做什麼的?.
管理飯食.
所以就有了這個層面.
十二使徒去選選和猜派.
選選的人從這裡開始.
我們來看看.
在教會裡面這麼重要的執事和領袖.
可以說領袖的第二梯隊的人.
應該有什麼生命的質素呢?.
是作主的僕人.
有什麼質素呢?.
我們從這些人身上都看到.
有什麼呢?.
我們馬上在經文裡面就看到.
原來這些人的七個都有幾個特色.
有好的明星.
什麼叫做有好的明星呢?.
不是自己說自己很好.
不是自己說自己.
而是被見證過的.
在群體裡面被周圍的人見證.
他是值得信賴的.
他是有很好的明星的.
你想想管理飯食不簡單的.
管理飯食不只是說食材的問題.
或者是煮食,派.
還牽涉到什麼?.
金錢.
這些人要是很值得信賴.
是有很好的明星.
大家一起見證.
一個忠心的僕人都應該是這樣.

$^{121}$他在教會的圈子裡面.
或者在外面也好.
別人見證他.
生命的榜樣是很好的.
我們開始去看一下.
自己有沒有一個很好的明星.
是別人一起見證我的.
值得信賴的.
第二樣東西是滿有聖靈.
這個並不是我們經常提到的.
我們很多時候會看那個人.
是什麼?.
忠心的僕人.
最重要的是他有能力.
辦事能力.
經常強調他有口才.
他有能力.
他就值得信賴了.
其實你看回《時圖行傳》.
它說的是什麼?.
是滿有聖靈.
是一個很重要的條件.
滿有聖靈.
你在哪裡看到他?.
你真的見證到他.
是真的看到他有聖靈的同在.
他生命裡面有聖靈的工作.
他是結滿聖靈的果子.
他的生命的果子讓你看到.
他整個出來的氣質不一樣.
原來我真的想說.
在教會裡面.
作為忠心的僕人.
這件事是很重要的.
要去看的.
我們不要像世界上的人那樣去看.
世界上的人就是看你有沒有口才.
你有什麼背景.
你的能力是怎麼樣.
我不是說在教會處事不需要能力.

$^{161}$但是你有沒有看過.
什麼叫滿有聖靈呢?.
滿有聖靈都會帶來能力.
一定會.
神同在.
怎麼會沒有能力呢?.
但這不只是能力這麼簡單.
而是神自己那種氣質.
那種質素.
神的同在是不一樣的.
忠心的僕人也應該有這一點.
那怎麼會有呢?.
總之很自然地.
不知道怎麼會有.
還是這個人是一個常常親近神的人.
是常常親近神.
常常禱告,讀經.
明白神的心意.
也不要只停在這裡.
我現在發覺我們弟兄姊妹都很熱心.
有不同的人都很喜歡查經.
不只是停在這裡.
你查經之後.
你聽了神的話語之後.
你去進行.
你走出來.
你很體貼神的心意.
你常常都很想知道.
上帝你的心意是什麼呢?.
我跟著你走.
你這樣帶領我.
就算我覺得不太合理.
或者我覺得有困難什麼都好.
但是我信服,我相信你.
我是你的僕人.
你是我生命的主.
上帝是我們的主.
不是我作主.
一個生命是否以神為中心.
有聖靈的人跟自我中心的信徒不一樣.

$^{201}$信徒也可以很自我中心.
我這麼久做人的經歷.
這麼多年.
這麼多歲數.
我比他們吃米吃鹽多.
我有很多經驗.
我何必多說.
我現在出來就做給你看.
這些不是那種屬靈的氣質.
不是那種滿有聖靈的.
這種生命.
僕人需要有.
很需要.
很需要有智慧.
需要有智慧.
是滿有智慧.
不是亂來橫衝直撞.
有熱誠就去.
隨便怎麼去都可以.
是不是這樣呢?.
不是,這些不是智慧.
闖禍多於吃飯.
那怎麼會是智慧呢?.
如果這樣去管理飯食.
一定會搞起很多風波.
是不是?這個不是.
中心的僕人不是這樣.
是滿有聖靈.
也滿有智慧.
非常非常寶貴.
在聖經裡提到一個人物.
今天我們的主角.
就是屍體犯.
他是第一個.
經文裡所提到的執事.
第一個.
第一個提他.
真的很重要.
在這個階段裡.
他具有一個很重要的影響性.

$^{241}$他生命的見證也非常好.
經文裡形容屍體犯很有信心.
滿有聖靈.
剛才說到.
選那七個執事.
真的要滿有聖靈.
有智慧等等.
你看屍體犯非常符合這個資格.
很有聖靈.
滿有信心.
不是他自己口說.
而是他的生命.
直接流出這種.
整個生命的特質.
他也滿有恩惠和能力.
在民間行大事和神職.
他的信心.
他很相信神.
很相信神的能力.
很相信神自己的聖情.
慈愛的公義.
他相信神有能力行神職祈事.
神事實上也透過他.
在民間行大事和神職.
做什麼呢?.
是見證福音.
通常神職祈事的出現.
神也要透過這樣.
透過他的僕人.
去見證福音.
讓人更加容易.
去相信耶穌基督.
這個很重要.
原來在屍體犯身上.
見到這樣.
其實當時的執事.
不只是管飯吃.
他們也是傳福音的人.
特別是在名單裡.
有一個肥利.

$^{281}$說明他是傳福音的肥利.
他是傳福音的人.
屍體犯也是傳福音的人.
所以不簡單.
今天我們當中也有堂委.
我也很想說.
作為堂委.
也可以這樣挑戰自己.
挑戰自己是什麼呢?.
滿有聖靈和智慧.
作傳福音的人.
靠著聖靈的大能.
是很奇妙的.
是有好的明星.
今天我們的堂委.
也需要這樣.
今天說到屍體犯.
後面也有說到.
他的面貌好像.
天使的面貌.
很厲害.
發光的.
有時我看到.
弟兄姊妹會說笑.
嘩!他頭上有光環.
這不只是頭上有光環.
大家看到他.
敵對他的人看到他.
也覺得他.
面貌好像天使的面貌.
是一個跟神很近的.
他的樣子.
在裡面的氣質滲出來.
整個樣子.
讓你看到他是一個.
完全屬於神的人.
是天使.
多厲害.
不簡單.
做教會的執事.

$^{321}$有這樣的氣質.
整個生命的表現.
這個屍體犯.
他經歷了什麼事呢?.
我們來看看.
經文說到神的道很興旺.
在耶路撒冷這個地方.
門徒的數目增加很多.
當時使徒.
去揀選了這些執事之後.
執事幫忙分擔了派飯的工作.
照顧不同的人.
使徒可以專注祈禱.
傳道的事.
門徒的數目增加.
教會又增長.
非常好.
甚至當中有很多的祭司.
都去信了耶穌.
當時這個屍體犯.
很忠心的.
為主工作的時候.
就有敵對的聲音.
什麼敵對的聲音呢?.
我們來看看.
原來有稱為自由人會堂的.
和古里內亞歷山大會堂來的人.
從不同的地方來.
我們來看看.
哪些地方呢?.
其實古里內亞是非洲的北面.
亞歷山大是埃及的.
一個很有名的城市.
基里加是古里內亞北面的位置.
這些地方的人.
什麼人呢?.
猶太人.
他們信非常傳統的猶太教.
非常傳統.
其實這些人.

$^{361}$是什麼人呢?.
不信耶穌的人.
覺得基督徒信耶穌的.
是正義異端.
背叛上帝.
什麼耶穌是神的兒子.
他是彌賽亞人嗎?.
他們一直迫使基督徒.
這些人起來和屍體犯辯論.
我們來看.
在這個時候.
屍體犯就去回應.
回應的時候.
面對這些反對的聲音.
他不是一個潑婦罵街式.
他不是以惡報惡.
他不是罵人.
他卻是以智慧和聖靈說話.
去到這些地方.
遠遠寸寸靠著聖靈.
是智慧.
哇!經文裡形容.
眾人抵擋不住.
勵志神的智慧.
勵志神的幫助.
人是抵擋不住.
屍體犯當時可以.
這樣去回應他們.
你看到那些人一直迫使他.
一直煽動百姓.
又作假見證.
控告他.
這個人不斷說話.
侮辱神聖的地方.
律法.
我們也聽到他說.
拿撒拉人耶穌.
要毀壞這個地方.
又要改變摩西.
所交給我們的規矩.

$^{401}$一邊說話一邊逼迫他.
但在過程中.
大家.
大家也無法否認這件事.
哇!他和神真的很近.
完全是一個屬於上帝的人.
這是非常非常寶貴.
在這裡我們看到很重要的.
第七章今天並沒有.
跟大家一起讀.
整個第七章的經文.
為什麼呢?.
第七章的頭部.
是一篇很長的講章.
非常寶貴的一篇講章.
在裡面你看到.
屍體犯的中心是什麼?.
他要講一些唐道.
去和那些來自不同地方的猶太人.
很傳統的猶太人.
就是那些不信耶穌的猶太人.
說什麼?.
其實神和整個民族的關係是什麼?.
神對整個民族的恩典是怎樣的?.
但我想說.
他在講這些唐道的時候.
其實他自己是一個什麼樣的處境?.
他自己是一個什麼樣的處境?.
他是被捉拿的.
他不是一個自由四處走動的人.
他是被捉拿了.
他現在是面對眾人.
有一個很大的危險.
眾人根本下定決心要害死他.
根本不是想著.
我聽你講完之後.
都講得很有道理.
我打算改變.
不是的.
是想害死他這群人.

$^{441}$他不怕危險.
在這裡忠心為主講信息.
我只是簡短在裡面說.
這篇信息.
他一直在講神對整個民族的恩典.
講到神怎樣呼召亞伯拉罕.
一直和他立約.
他的恩典怎樣從亞伯拉罕到以撒.
一直到雅各.
上帝在裡面有很多很奇妙的作為.
一直到約瑟.
約瑟的經歷是約瑟被他的哥哥們.
去質盜.
被賣到埃及.
約瑟在那裡受了很多苦.
但是神怎樣在約瑟的身上工作.
令他蒙恩.
作宰相.
然後在饑荒的時候.
雅各知道埃及有糧.
就派他的兒子去扔糧.
於是最後有機會.
整個雅各的家族去了埃及生活.
一直到以色列整個家族.
就在埃及的地方定居下來.
然後這個民族在埃及發展.
去到後來的時候.
很多年之後.
有一個不認識約瑟的王.
興起來統治埃及.
輔待當時以色列民族的祖宗.
說這麼多.
是想說.
其實猶太人他們的祖先.
經歷了很多困難.
但是有更多的神的恩典.
神的英雄要賜福亞伯拉罕和他的後裔.
他的英雄從來都沒有落空過.
一代又一代的保守.
拯救.

$^{481}$幫助他們.
到後來摩西出生.
神要透過摩西拯救整個以色列民族.
離開埃及為奴之地.
說到摩西的經歷.
很多很多.
一路說.
我相信猶太人心裡很有感受.
因為他們非常熟悉自己民族的歷史.
但是.
大家看歷史就很不同.
猶太人可能只是覺得自己.
很有民族的自豪感.
我是屬於亞伯拉罕的子孫.
我們多麼的尊貴.
是神的選民.
但是當中.
他們有沒有細細想想.
神這麼多年的保守.
這麼多年的拯救.
這麼多年數之不盡的恩典.
在這裡屍體版帶他們看.
一路帶他們看.
這是很寶貴.
所以整篇講章是很寶貴.
這裡說到摩西的夢召過紅海.
神救贖的恩典多麼大.
神不單救贖他們.
還賜下十誡有律法.
成為以色列民一個很大的指引.
應該怎樣去作為神的百姓生活.
這些都是神給以色列民的福氣.
但是在這裡.
以色列民的回應是什麼呢?.
在經文裡看到.
屍體版的中心位置在哪裡?.
一點都不客氣.
一點都不隱藏.
不需要婉轉.
直接說.

$^{521}$其實反而是以色列民很活躍上帝.
拜偶像.
拜這個 拜那個偶像.
撥翼的施恩典給他們的上帝.
所以在這裡.
屍體版是指責他們.
但這個指責其實.
不是純粹為罵而罵.
這個指責是很希望.
這些猶太人都能悔改.
你看看以你們的祖宗為尷尬.
他們這樣的撥翼.
他們的下場是什麼呢?.
今天你們是否要悔改呢?.
其實今天你們是如何對待耶穌基督呢?.
那二者其實是在說耶穌基督.
你們是如何對待耶穌基督呢?.
是不是?.
在經文裡是如何說的呢?.
是非常不客氣.
你們是如何對待祂呢?.
你們是否真的好好對待耶穌基督呢?.
不是的.
如今你們成了那.
二者的出賣者和兇手.
猶太人你們.
你們有份殺害耶穌基督.
你們出賣耶穌基督.
以前你們的祖宗是這樣背叛上帝.
現在你們是這樣背叛神的兒子.
都是背叛神.
你們說你們有從天使而來的律法.
但你們有沒有遵守呢?.
沒有.
司祭犯.
用了很長的篇幅.
去說這一切.
這裡令我們有兩點的思考.
第一點的思考就是.
為神發聲.

$^{561}$我經常會去想我們在教會當中.
有多少去述說.
神的恩典,神的好呢?.
這是整個民族.
神在整個民族為以色列民做了什麼.
我們有多少去數算.
神在我們個人一生裡對我們的恩典.
你說我很遲才信耶穌.
你知不知道還沒信耶穌的時候.
神都對我們家裡對我們自己很多恩典.
有沒有一起數算神的恩典.
一生的作為.
我們教會有多少去數算.
上帝對我們的作為呢?.
接下來是十四週年的堂慶.
在裡面有一個很好的機會.
數算神的恩典.
逐樣逐樣去看.
神在當中為我們做了什麼.
很重要.
有沒有為神說話.
在人群中.
我們在人群中說的是什麼呢?.
你有沒有留意?.
我們很多時候因為經歷一些很困難的事.
身體的疾病.
經濟的困難.
婚姻的問題.
很多的艱難.
我們很難去思考神的恩典.
很多時候我們會說很多負面的話.
我們甚至會埋怨神.
有些人離開了教會.
不再回教會.
人生這麼多困難.
神對不住我們.
在這個時候.
有沒有想起神的恩典.
在這個時候有沒有人出來為神說話.
有人一直在說神這樣這樣的時候.

$^{601}$我們有沒有出來說神對我們的恩典.
何其多呢?.
斯提法這個講座很重要.
中心的僕人就是有這個很重要.
第二樣很重要的就是.
不要搏逆上帝.
我們真的很容易搏逆.
很容易搏逆.
我們拜的偶像是有形的.
有些是無形的.
我們總有些東西自己很喜歡.
超越上帝.
有些什麼生命不如意.
我就離開了.
我就不相信了.
我要每樣東西都順利.
每樣東西都合我的意.
我就相信你了.
我就跟著你了.
這個不是這樣的.
還有我們有沒有.
好像斯提法一樣.
不怕死.
提醒別人.
不要搏逆.
忠於上帝.
有沒有這樣的勇氣.
有多少人在自己的家裡.
勇於為耶穌基督作見證呢?.
不要說這麼多.
一說我家人就不喜歡了.
破壞關係.
吃飯都吃得不開心.
很重要的.
吃飯開心關係和好.
不要說.
不夠膽在家裡見證耶穌基督.
不夠膽在鄰舍見證耶穌基督.
不夠膽在公司見證耶穌基督.
很怕.

$^{641}$我們看重別人多於看重上帝.
但斯提法是什麼呢?.
我不但不怕.
我現在被你抓了.
我面對這樣的審判.
我面對要迫害.
快要死了我也不怕.
弟兄姊妹我和你又怎樣?.
我和你又怎樣?.
昔日的以色列文就是這樣.
很多個拜偶像.
斯提法到最後是致死忠心.
剛才最後一段經文說到.
眾人一聽到那些說話.
指責自己拜偶像.
又潑逆.
甚至我們現在.
還要是出賣那位義者.
我們又殺了耶穌.
越聽越氣憤.
斯提法所說的訊息裡.
其實說到.
為什麼以色列會成為亡國呢?.
就是因為他們背叛上帝.
其實當時.
這群猶太人聽的時候很有感覺.
為什麼?.
當時他們聽的時候.
他們還沒有復國.
斯提法說.
我今天還沒有復國.
我今天是亡國奴.
都是因為我背叛了.
是不是很不中聽?.
他們真的很想.
你這樣說我.
我就要殺死你.
所以他們後面.
眾人聽見這些說話.
心中極其惱怒.

$^{681}$向他們咬牙切齒.
就要去殺死他們.
但是斯提法呢.
他很傷心.
在這個時候.
他所做的不是一連串的咒罵.
反而是將他的眼目.
放在望天.
神讓他看到.
神自己的榮耀.
又讓他看到耶穌站在神的右邊.
這個真是.
跟神很近的人.
神給他的福氣.
就在他離開之前.
那個時間.
給他一個很大的鼓勵.
我看到天開了.
孕子站在神的右邊.
這個時候.
大家非常憤怒.
一擁而上.
用石頭打死他.
但是我很想說.
到最後.
他不是這樣說.
如果我們通常被人這樣傷害.
這樣殺死.
我們會咒罵的.
你不得好死.
你一定死得很慘.
神一定會懲罰你.
但他不是這樣說.
他說什麼呢?.
他很像主耶穌.
他說:主啊!.
不要將這罪歸於他們.
好像主耶穌最後的日子.
不要將這罪歸於他們.
他不是以他的憤怒.

$^{721}$去回應那些傷害他的人.
忠心的僕人.
至死忠心.
是至死忠心.
他是寬恕傷害他的人.
跟從主耶穌的人.
就是這樣.
寬恕傷害他的人.
到底.
這是非常非常寶貴.
司提反的死.
在《使徒行傳》裡.
是一個轉捩點.
很重要.
他的死.
在後面.
教會要面對一個很大的逼迫.
滿途四射.
福音工作.
向不同的地方展開.
當時.
有一個年輕人.
就是素羅.
後來很有名的保羅.
看到司提反的殉道.
司提反對他.
有很深的影響.
很深很深的影響.
對保羅後來.
有很大的影響.
所以你想想.
司提反是一個很重要.
很重要的人物.
在這裡.
他為主的福音.
為主的道.
至死忠心.
最後仍然寬恕.
這個人.
近期熱話.

$^{761}$這個人.
是甚麼?.
查理柯克.
我想說一件事.
你知不知道司提反是多少歲?.
原來司提反很年輕.
我查過.
司提反很年輕大概27至31歲.
查理柯克離開的時候.
你知不知道他多少歲?.
31歲.
兩個都是年輕人.
兩個都是很熱血的年輕人.
我自己也很感動.
當我去了解查理柯克.
他短短的人生.
他對世界的影響.
心裡也很感動.
他說查理柯克.
是一個為基督.
站立的人.
很寶貴.
他在18歲的時候.
成立了.
Turning Point USA.
他覺得美國的文化.
和道德基礎.
正在被侵蝕.
這也是事實.
特別是基督教的價值觀.
是很大被否定.
所以他成立的機構.
他很想將年輕人.
引向聖經的原則.
和原先的價值.
但更重要的是.
柯克很重要的一件事.
他會公開見證自己的信仰.
很坦誠地.
分享自己的信仰.

$^{801}$「是的,我是相信耶穌的」.
「主耶穌拯救了我的生命」.
「我這個罪人是需要耶穌基督的」.
他本身成長在.
一個非常普通的美國家庭.
在大學的時候.
遇到一個改變一切的真理.
不過初期.
他仍然是野心勃勃.
但內心卻很空虛.
直到後來他看聖經.
他知道.
耶穌基督的復活.
並非只是神話.
而是真實的.
是真實的.
耶穌基督的復活.
於是上帝就改變了他一生.
從那時候開始.
他就特別注重.
在校園傳福音.
在校園裡.
他面對.
你想想,美國的校園和我們不一樣.
美國的校園.
其實都很反對福音.
很強調.
美國的精神.
美國的文化.
強調自己的人權.
等等.
但他就是要在校園傳福音.
你說他孤不孤單.
很孤單.
需要勇氣嗎.
非常需要.
面對的很多嘲笑.
威脅.
甚至死亡的陰影.
但他從來都沒有打算退縮.

$^{841}$去鼓勵年輕人.
告訴其他人.
耶穌的事情.
這個機構影響了幾萬名年輕人.
讓他們歸向上帝.
教會有成長.
他所說的訊息.
是非常道地的.
是說政治.
說婚姻,說祈禱的力量.
在這裡.
他不斷提醒基督徒.
不要做旁觀者.
弟兄姊妹,我們不要做旁觀者.
要做上帝的精兵.
是勇敢的.
要去作見證.
曾經有人問他.
將來有一天.
如果你離開.
你最希望別人記得你什麼.
離開世界.
你猜他怎麼說.
希望別人記得他什麼.
他的信仰,他的勇氣.
是他做到的.
在一個公開的場合.
他在說信仰.
和年輕人辯論的時候.
就有一個子彈.
打到他身上.
取走了他的生命.
在眾人面前.
取走了他的命.
31歲.
但是他離開.
激勵了很多信徒.
面對職場的壓力.
面對家庭的挑戰.
社會的風暴.

$^{881}$一定要堅守真理,不要退卻.
他用他的生命.
告訴你.
致死忠心.
是致死忠心.
為真道.
擺上自己的生命.
弟兄姊妹,我們要起來.
為主致死忠心.
不要退縮.
起來.
不要做旁觀者.
投身在這場仗當中.
見證耶穌基督.
當然,我們的名字.
不是每個人都叫屍體犯.
屍體犯的名字叫.
「冠冕」.
是一個非常漂亮的名字.
他真的得到生命的冠冕.
我們也不是叫「阿克」.
但我和你.
可以做忠心的僕人.
我們可以.
靠上帝.
親近神.
成為一個有好明星的人.
滿有聖靈.
我們是致死忠心.
有一天,我和你.
都由福氣在神的面前.
得到主的稱讚.
我們一同禱告.
天父上帝.
透過屍體犯的經歷.
透過「阿克」.
自己生命的見證.
讓我們了解.
神你真的.
很想我們.

$^{921}$成為一個忠心的僕人.
一生忠心為你.
一生要討你喜悅.
屍體犯為你發聲.
叫人離開.
搏役.
離開拜偶像.
「阿克」.
也挑戰人.
堅定地.
相信你.
跟從你.
做你的忠心的跟從者.
最後,今天求你.
給我們這份忠心.
在我們所生活的地方.
我們的家.
我們的職場.
在我們的論社.
在朋友圈當中.
讓我們忠心.
為主發聲.
讓我們鼓勵人.
離開罪惡.
鼓勵人離開偶像.
一生只要你.
為他們生命的主.
求你賜福我們.
復興我們.
奉耶穌基督的名字.
\newpage



\section{馬太福音 9:35}
\label{sec:WCmVYVwec7g}
\textbf{2025.09.28《耶穌的「三濟」事業:教導的「理濟」傳道的「道濟」,醫治的「利濟」》馬太福音9\_35 (和修版) 郭鴻標牧師}
\newline
\newline
連結: \href{https://youtube.com/watch?v=WCmVYVwec7g}{\texttt{https://youtube.com/watch?v=WCmVYVwec7g}} ~~~~ 語音日期: 2025-10-01
\newline
\newline
\hyperref[sec:pPe9_H3XXfI]{\small{< < < PREV SERMON < < <}}
~
\hyperref[sec:index_chronic]{\small{[返順時目]}}
~
\hyperref[sec:index_scriptual]{\small{[返順卷目]}}
~
\hyperref[sec:ZmW9HFGb5CU]{\small{> > > NEXT SERMON > > >}}
\newline
\newline
馬太福音 9:35
\newline
\begin{longtable}{cl}
\hline
\hline
章節 & 經文 (和合本修訂版)\\
\hline
9:35 & \begin{tabularx}{0.7\textwidth}{X} 耶穌走遍各城各鄉，在他們的會堂裡教導人，宣講天國的福音，又醫治各樣的病症。 \end{tabularx} \\ \\
[1ex]
\hline
\hline
\end{longtable}
$^{1}$各位弟兄姊妹早安.
我們一起同心祈禱.
親愛的天父上帝.
願你接納我們從心裡發出的敬拜.
願你的聖靈在我們心裡面.
向我們說話.
將你的話語植根在我們心靈裡.
亦都讓我們在今天聽到的時候.
是可以更加親近我們的神.
祈禱奉耶穌基督的名,阿們.
話說上個星期曾銳貴牧師聯絡我.
他說今個星期你提到的題目和經文是甚麼.
我就看看日曆上沒有這個紀錄.
然後就一直思想答應還是不答應呢.
曾牧師說你不答應都無所謂.
但我又覺得不好.
問題就來了,我是找甚麼講道理來紅恩堂呢.
不怕告訴大家我抽屜底有很多舊貨.
在電腦裡面儲起來.
但我又覺得這樣不是很尊重曾銳貴牧師.
神給我一個感動.
因為我們都有靈修,祈禱,有閱讀.
剛好我正在看一本書,看完了.
我心裡有感動.
不如就將那些靈感整理成為講章.
如果日後用,第一次都會在紅恩堂用.
這個講道理的起源是甚麼呢.
就是在疫情之後.
我們其實都在問一個問題.
香港的教會何去何從.
我們每個教會都問這個問題.
我估計和很多同道都有這個特點.
我是先從我自己的教會處境開始思考.
我教會有甚麼難處,有甚麼需要.
然後再看看我見到的那些教會情況.
很正常的,如果在洪水橋區.
你就想想洪水橋區的環境.
然後再想想教會的前路.
我看完那本書,我發覺完全改變了我.
那本書是甚麼呢.

$^{41}$那本書我沒有帶來,都幾後行.
就是香港基督教史1907到1997.
大家喜不喜歡歷史.
我一會兒就和你說故事.
就是2024年出版的.
其中有作者送了給我.
你知道有人送書給你,你都多多謝謝.
有時不經意就送了去不知哪本書.
但神又感動我,我真的拿來看這麼後的書.
我真的逐一看完.
這本書是講宣教士來中國的歷史.
如果你有少少印象,你知道應該何時開始說.
是主後七世紀的列斯多流.
來到中國,就是大秦景教碑開始說.
不是這樣說,你都要說16世紀的時候天主教的傳教士來華.
就是李馬竇,湯若望,南懷仁那些.
但他都沒有提.
他也沒有再講97年之後香港教會的情況.
但這本書都是有很特別的.
我自己覺得是,我得到很大的益處.
他就嘗試將香港的教會擺在.
由1807年馬勒信來華的脈絡底下去理解.
來展現一幅很宏觀的圖畫.
我今天沒有powerpoint,所以你要聽我說話了.
這本書有七章,我就用我自己的文字去表達.
我先講一下大概七部分是怎樣.
就是早期的我們叫更正教的宣教.
即是我們中文叫基督教,其實是Protestant.
就是從馬勒信來華開始,1807年.
然後就是鴉片戰爭之後的宣教.
就是1842年到1883年.
這些是很久的歷史.
然後這本書很有意思.
他就說,不只是講西教事.
中國教會領袖的出現.
就是由1884年到1914年.
即是那時候已經有華人的牧師.
然後去到進入自治的時期.
1915年到1940年.
即是華人的教會一直成長.

$^{81}$去到日本侵華時期.
就是1941到1945年.
不知道在座中有沒有人經歷過這段時間呢?.
你可能就會明白多一點.
然後就是香港教會本土化的時期.
1945到1973.
在座中可能有些人都出世了.
然後去到香港教會的增長期.
1974到1997年.
我們很多人都是在正正香港教會增長期的時候.
來去成長,來去信主.
而我將這七個時期變成兩個部分.
就是今天講到的第一第二部分.
就是說培養華人的領袖.
第二就是戰後對有需要的人的服侍.
這兩樣東西就構成今天我們香港萬人教會的基礎.
然後在祈禱和思考裡面.
我發覺一句話很特別.
就是作者,其中一個作者.
應該就是李志剛牧師,老前輩.
他就說,他用馬太方第九章三十五節的經文.
去講過去的中國和香港華人教會的宣教士和歷史.
可以說是根據著耶穌的三濟事業.
這些字眼就很深很古老.
我問一個問題.
有沒有人聽過香港有一間道濟會堂.
就是中華基督教會合一堂.
就是潘咸道那間合一堂.
他的前身就叫道濟會堂.
其實你不太知道是什麼.
你可能以為是道堂.
其實不是.
因為原來在早期的宣教士在那裡講.
耶穌三濟事業就包括教道的理濟.
教道教道道理.
用道理來濟世.
傳道.
那個是道濟.
即是我們傳福音.
你回帶回去一百四十幾年的時候.

$^{121}$原來那些宣教士是在講傳道.
就是我們去道濟.
讓其他人得益.
醫治就是利濟.
即是辦醫療所那些.
原來就是這些西教士一直下來就在做這些事.
所以我講到的三部分.
第一就是怎樣培養華人的領袖.
那時候沒有什麼華人基督徒.
第二就是戰後對於有需要的人的服侍.
然後第三就是.
今天我們要關心的更大的異象.
當我們問今天香港教會何去何從的時候.
你就要從根開始講起.
首先第一要講培養華人教會的領袖.
馬來西亞是一百零七年來.
他來的時候主要在澳門和廣州.
廣州其實是不讓人進去的.
所以那些西教士就停留在澳門.
當時是清朝的時代.
就禁止傳教.
馬來西亞就轉移到南洋宣教.
開展教育出版.
所以馬六甲那裡做這些事.
然後一百四十二年有南京條約.
割讓了香港.
就開放五口通商.
大家都記得這些.
還有一百四十四年望下條約和黃埔條約.
就容許宣教士在這五個口岸裡面.
建教堂.
辦醫院.
傳福音.
就是說五口通商不是只做生意.
而是後面望下條約.
黃埔條約就讓他們去買地建教堂.
辦醫院.
傳福音.
所以歷史是這麼走的.
然後一百四十二年的宣教士來香港.

$^{161}$他們的目的不是只在香港做生意.
而是他們向內地傳福音.
你看見這個歷史的時候.
原來澳門也是一樣的.
當南京條約之後.
駐澳門的所有的天主教宣教士.
基督教宣教士.
大部分都調走了.
所以澳門的基督教的發展是比較弱.
原因是警察會將他們調走了.
去上海那裡.
所以香港也是一樣.
香港的宣教士.
他們是踏足香港作為一個跳板.
而去國內傳福音.
你要明白這個歷史.
我看完之後也很震撼.
原來我們的角色就是這樣.
大家知不知道香港第一間教會是哪一間.
很有趣.
就是1843年7月17日.
在皇后.
皇后浸信會.
就是皇后大道中那裡.
就是今天的建浸.
香港浸信會.
長洲浸信會也是1843年開始主日崇拜.
故事就是美國浸信會的一個叔尾士穆斯夫婦.
在1836年去了澳門.
然後1842年.
不知道他們是否收到消息.
所以搬來香港.
準備建教堂.
為什麼1843年可以買到地建教堂.
早一年已經有這個計劃.
1843年建教堂.
今天有些泰國的姐妹在這裡.
有一個叫倫維仁穆斯.
叫William Dean.
1834年去到泰國曼谷宣教.

$^{201}$你們聽過吧.
真的很厲害.
在曼谷建立第一間華人教會.
這個真是歷史.
當時在泰國有很多潮州人.
在座中有沒有人是潮州人.
有潮州人.
很明顯你們知道這個歷史.
倫維仁穆斯就在泰國曼谷學潮州話.
這個只是故事的開始.
1842年.
倫維仁穆斯就來香港.
1843年在街市街的浸信會.
就開始第一間潮語浸信會.
不只是這樣.
他帶著潮州基督徒去長洲.
向潮州人傳福音.
所以為什麼1843年長洲浸信會開始有聚會.
因為泰國有些基督徒來香港宣教.
所以剛好有泰國的姐妹.
這個歷史不是假的.
可能很多人是不會知道這件事.
故事就是香港還沒有成為一個可以傳教的地方的時候.
宣教士就在東南亞做了很多事.
一時到那裡那些基督徒就來香港幫忙宣教.
你想一下西人怎麼會這麼容易做到這麼多事.
因為他們是華人的助手.
另外一個叫倫敦傳道會的馬勵信穆斯.
在馬六甲辦英華書院.
不知道在座有沒有人是英華的校友.
成為宣教的總部.
印中文聖經.
所以其實宣教士做不到事就是印聖經.
只是他做預備功夫.
1843年倫敦傳道會駐馬六甲的九位宣教士.
全部調過來香港和中國.
兩個去了福州.
兩個去廈門.
三個去上海和寧波.
有兩個去香港.

$^{241}$所以你看到部署就是.
九位宣教士只有兩個在香港.
香港只是一個其中的宣教點.
其實目標是全中國的.
1843年倫敦傳道會在上環成立道濟會堂.
就是今天的中華基督教會合一堂.
然後1863年又成立了後來叫中華基督教會灣仔堂.
就是在合會中心那一間.
可能你會說還有什麼呢.
原來有間叫公理堂.
是的,這個由於他們的差會就是美部會.
美部會在1883年就成立公理堂.
聖公會在香港第一間華人教會就是聖士提反堂.
就在1864年.
所以你看看華人教會成立就是這樣的局力程.
1852年巴釋會就在西應盤開機.
就是今天的崇禎會舊恩堂.
在座有沒有一些客家人在呢.
崇禎會就是客家人的.
1862年就在筲箕灣開始.
又開了一間教堂就是崇禎會筲箕灣堂.
那1842到1883年呢.
這40年倫敦傳道會成立了後來的中華基督教會合一堂.
中華基督教會灣仔堂.
還有聖公會聖士提反堂.
巴釋會成立的崇禎會舊恩堂.
筲箕灣堂.
估計華人信徒只有1000人.
40年只有1000個基督徒.
當時有三大公會.
倫敦傳道會.
聖公會和巴釋會.
都被香港政府批准在博庫林.
這個現在叫華人基督教永遠墳場.
就是你要多謝那些前輩.
就是我們的上一代.
信了主然後就葬在博庫林華人基督教永遠墳場.
就是因為當時他們三大公會和政府.
說不如拿個山頭給我們.
宣教士的宣教工作其實做很多事.

$^{281}$又要翻譯聖經.
又要開醫院.
又要辦學校.
又要搞社會服務.
其實真的很多事要搞.
而宣教士還有呢.
我要數給大家聽.
你知不知道中國過去有13間基督教大學.
很厲害的.
我說些名字出來可能你聽過.
金陵大學.
金陵女子文理學院.
開女子大學.
你知不知道是一個什麼意思.
是想男女平等.
所以基督教的宣教士做了很多事.
聖約翰大學.
上海的很出名.
滬江大學.
東吳大學.
之江大學.
燕京大學.
就是後來的北京大學.
華中大學.
齊魯大學.
華西協合大學.
福建協和大學.
華南女子文理學院.
還有嶺南大學.
嶺南就下來香港了.
13間.
1884到1914.
現在說到另一邊.
那30年又看看香港.
有以下的教會成立了.
我只是說華人教會.
只是說幾十年才有這麼多間.
不是變數字出來.
1884年有中華巡渡會.
就是巡渡會的人來.

$^{321}$1890年聖公會聖三一堂.
就是九龍城那間.
1890年就是九龍城浸信會.
經常去陪靈會那裡.
1890年就是聖公會朱聖堂.
在旺角那裡.
1892年就是中華基督教會深水寶堂.
然後1896年就是皇兒周宣崇真會.
1898年就是中華基督教會元朗堂.
這是全新界第一間華人教會.
就是搭西鐵是前幾個站.
然後1899年就是禮賢會香港堂.
就是很後才開的.
1901年尖沙咀浸信會.
1903年就是中華基督教會聖光堂.
1903年就是中國基督徒會.
1905年就是中華基督教會荃灣堂.
1905年就是香港仔浸信會.
你數下來,嘩,這麼厲害.
1905年就是粉嶺崇天堂.
1905年就是窩尾崇真會.
1907年就是中華基督教會大埔堂.
1907年就是五秦節會.
1907年就是呂廣朝人中華基督教會.
1911年就是聖公會聖瑪利亞堂.
好像廟那間在銅鑼灣那裡.
1911年就是聖公會聖保樂.
就是,是啦.
1912年就是海面傳道會.
專向漁民傳福音的.
在在中有沒有是,就是,水上人.
有,你們知道,海面傳道會.
我聽說那些牧師很有型.
那些船民,不是,那些漁民.
出海可能是一個月.
去很遠的地方打魚.
那些牧師就上去跟他們祈禱.
打完有魚獲回來.
就一打,其實牧師有條鹹魚給你.
很靚的鹹魚,要醃了的.

$^{361}$那些人是出很多風險的.
他們不拜其他的偶像.
而信耶穌,你明不明.
他是很不容易的.
好了,還有.
1914年是便義理會.
有沒有人聽過便義理會.
然後就,就是,這些教會.
說的也是過了30年.
40年,30年,70年.
到1915年,香港華人基督教會成立.
其實,原來是慢慢才有的.
培靈會什麼時候開的?1927年.
在廣州開,然後香港跟著做.
這樣的故事.
如果你說1843年是香港第一間華人教會.
到現在2025年.
說的是182年歷史.
大家有沒有想過,我們是在承傳182年的歷史.
包袱重了很多.
上帝提醒我,當我去問教會未來方向的問題.
就不只是好像我開始說.
我自己教會個別的堂會的難題.
或者這幾年的危機去想.
這樣想是斬斷了182年的歷史.
你明不明白我說什麼?.
所以我發覺,神提醒我.
我們是要有歷史意識.
香港教會不是今天才存在的.
而是已經有182年的歷史.
當你有這個思維的時候.
你就會知道.
我們的香港教會存在.
是一直跟中國的宣教是結連的.
宣教是這樣的.
教會從來不是為自己存在的.
教會也不只是提供宗教活動.
使我們覺得自己有存在感.
你明不明白?.
我們不是為自己存在的.

$^{401}$我們是有一個更大的使命.
究竟香港教會何去何從呢?.
就要看個別情況了.
不過你回顧過去的香港教會歷史.
是可以看到上帝是怎樣在這裡工作.
我們可以怎樣去跟隨.
過去宣教士做了很多事情.
同時培訓了非常多的華人牧者接班.
過去182年香港教會開展了很多事工.
栽培領袖.
直到今天這個地步.
你想想前人做的事情.
1883年香港有三大工會.
1914年加上浸信會.
就有四大工會.
即是什麼呢?.
當時宣道會,播道會都沒有名字.
這樣說是歷史事實.
在這樣的大光譜下.
去了解不同宗教的發展歷史.
不是要比較.
不是誰厲害,誰不厲害.
不是角力.
而是要承傳前人開方,火熱,學苦的精神.
如果我們沒有這個眼界.
我們其實是搞來搞去都是搞一下.
我們要做一些事情是承先啟後的.
第二,講戰後對有需要人的服侍.
講近一點.
1937年7月7日發生盧溝橋事變.
大家都應該知道這些中國歷史.
往後八年裡面中國很多的城市都淪陷.
1937年是天津,上海和南京.
1938年是廣州和廈門.
1842年是溫州和南昌.
1945年是長沙.
到1941年香港淪陷.
日治時期的香港社會出現很嚴重的經濟壓力.
原來1941年推出軍票.
在座中不知道有沒有人有軍票.

$^{441}$不知道怎麼搞的.
當時是說一元兌兩元港幣.
即是初時.
然後1842年又推出軍票.
今次是一元兌四元港幣.
究竟早換還是遲換.
然後到1843年.
日軍總督部禁止軍票以外的一切港元.
不可以用來做貿易.
你家裡有港幣都抓你.
所有香港市民一定要把港幣換成軍票.
很多人的資產一夜貶值.
所以現在很多人都說軍票放在那裡不知道做什麼.
要求賠償.
日治時期的香港教會出現大量信徒流失.
你都想到的.
肯定要走.
香港牧者走不走.
走.
所以那時候很艱難.
而1942年日軍總督部下令所有宗教團體要登記.
什麼教都要登記.
基督教成立了香港基督教的總會.
在那段日子很艱難.
社會不穩定.
人心惶惶.
教會領袖受到監控.
在各種壓力下.
繼續做關心牧羊的工作.
我看了另外一些傳記.
當時的牧者和信徒.
他出去探訪多數是危難的.
第二就是無端端被炸彈炸死.
要替他辦喪禮.
你想想.
那時候做這些事情很艱難.
然後呢.
1945年8月香港重光.
先不要說那麼多黑暗的事情.
1948年香港華人基督教會成立.

$^{481}$講一下國光譜又不同了.
1949年之前香港有七大工會.
倫敦會,聖公會,巴士會,浸信會,禮賢會.
突然間上了榜.
巡渡會又上榜.
公理會.
由於倫敦會和公理會加入了中華基督教會.
所以從此就有六大工會.
好像武林,武當,少林那些.
這個六大工會的格局就出來了.
成為中華基督教會,聖公會,巴士會,浸信會,禮賢會,巡渡會.
到1949年之後又怎樣呢?.
國內很多的.
就是警察會就沒有撤退.
但是就將他們的檔案或者資產那些東西.
在香港成立另一個辦事處.
他們是很想仍然留在中國的.
但是1950年.
警察會西人宣教士撤出大陸來香港是大勢所趨的.
你不走都不行的.
1951年香港的中派教會就要脫離和國內教會的關係.
所以大家劃清界線了.
然後封關那些大家有些印象了.
然後1949年大量的難民湧入香港.
就是我上一代的父母都是這樣的.
很多來到香港的宣教士覺得香港政局不穩.
有些警察會就說.
其實我準備調走所有香港的宣教士.
就是1949年之後應該是全部調走的.
但是有些宣教士就覺得不行.
我的心是很難過的.
我看到條頸嶺這麼多難民我要走去開荒的.
所以他們就在條頸嶺去做福音工作.
嚴格來說是這班不聽話的宣教士留在香港來做福音工作的.
歷史是真的這樣的.
就是警察會調你去印尼你不肯去.
你自己去做事.
我回鄉下籌款來做.
是這樣的.
這個動盪的歲月.

$^{521}$當然在韓戰之後.
中國是沒有用武力收回香港.
其實1949年的兵去到羅湖沒有過來.
警察會就覺得.
似乎都可以在香港宣教.
就開始派一些宣教士過來香港.
就是1953年之後才開始有這樣的情況.
所以在1956年.
香港教會名錄有50個警察會.
在香港有辦事處.
估計在香港的宣教士有300個.
就是未到60年代的時候有300個西教士.
好了.
就是1950年的時候.
新中國那個.
就是50年底.
他要求國內教會一定要肅清帝國主義的餘毒.
就斷絕和帝國主義的一切關係.
所以全部都變形了.
香港教會就變成什麼.
他可以繼續發展.
1962年有60個中派和獨立的堂會.
你看看,起飛了.
然後有344間教會.
大概有11萬的信徒.
剛才說了1000,現在11萬.
很特別,是不是.
到了1960年.
這班難民變成了新來港人士.
面對著大批的新來港人士.
教會就要問,我在這裡有什麼使命.
就是你做不做難民工作.
是有分別的.
58年大躍進出現大饑荒.
大家都記得了,有逃亡潮.
這個時候香港的社會是風雨飄搖.
經歷了50年,51年的木屋區火災.
36000人沒有家園.
石硤尾大火.
有6萬災民.

$^{561}$你可以記得這些東西.
大量學童是失學.
流落街頭.
1961年大概有8萬個6到11歲的失學兒童.
1956年,九龍的荃灣.
發生暴動,有激烈的衝突.
這些是故事.
在這樣的情況下.
社會福利方面就因為要推動救濟工作.
所以政府就容許教會去辦學校.
1966年,在香港教會名錄裡面.
更正教的小學和幼稚園有225間.
所以可能我們讀書就在這些學校讀.
1955年.
香港政府批准.
宗教和志願團體在.
七層高的寫字大廈辦天台學校.
有沒有人讀過天台學校.
有,我也是.
我讀幼稚園.
我那時候是8元學費.
我家人就說.
其實我們給不起學費.
基本上你沒得讀書.
不過教會辦學很好.
他給我十幾元.
他叫獎學金.
哪是獎學金.
叫助學金.
沒那十幾元,交不到學費我就不用讀書.
所以很多西教士就做了這些事.
當時有11間由教會辦的天台學校.
所以1955到1958年.
香港教會增加到294間.
去到7萬4470個弟兄姊妹.
而到58到62年.
去到344間.
有11萬2200個基督徒.
所以你看到那時候.
香港教會是這樣發展的.

$^{601}$而62年香港教會的光譜又變了.
出現了八大工會.
很有趣的.
你數一下有多少間堂會.
排頭的就是浸信會.
當時有54間.
會有1萬5千多.
為什麼突然間戰後浸信會這麼厲害呢.
中華基督教會只有24間教會.
只有1萬多個教友.
聖姑爺廟23間堂.
就1萬3千多.
信義會是33間堂.
就1萬多個教友.
路德會18間堂.
有9千多人.
秦道公會和惠利公會.
8間就有7千多人.
崇禎會11間就有5千多人.
禮賢會6間就有5千多人.
你發覺這個光譜又變化了.
你要追下去問.
為什麼浸信會這麼厲害呢.
這些就是另一個題目.
今天就講不到了.
八大工會已經佔全港堂會51\%.
會有人數是73\%.
當然他們很勤力去傳福音.
才有這樣的結果.
而浸信會的增長是最特別.
成為香港最大的中派.
信義會路德會的發展.
是因為在難民裡面有很多外省人.
他們說不懂廣東話.
所以信義會和路德會.
外省牧師帶了很多外省信徒.
信主.
60年代的中期香港經濟好轉.
你還記不記得.
什麼時候有家可以買到電視.

$^{641}$可以買到洗衣機.
我都是六幾年的時候.
發覺以為可以.
我小時候去隔壁看別人的電視.
你們不知道什麼.
原來是這樣的.
60年代中期香港轉型了.
不是再活在救濟品年代.
國際救援機構都撤離香港了.
烏溪沙都不再做孤兒院了.
到了67年香港暴動.
有什麼呢.
貧富距離很嚴重.
開始香港華人教會注重基層勞工的工作了.
你看到的.
70年代是一個脫邊的年代.
因為戰後出生的那批人.
本土意識比上一代高.
我們上一代是過客.
但是另一代他們對香港有承擔.
所以大專的基督徒.
對香港和中國有承擔.
所以是不同的.
有人將當時的教會分成兩派.
一派叫普世派.
一派叫福音派.
普世派就說天國在人間.
我們做多點服務.
關心社會公義.
福音派就注重個人得救.
1966年平安福音堂成立.
中華基督教會就70年代大量建立學校.
其實福音派的基督徒都有關心社會文化的.
譬如有關心弱勢的群體.
譬如突破.
有人聽過突破機構.
工業福音團契.
香港基督徒福音學生福音團契.
神學訓練方面.
也有重視高等神學訓練的中國神學研究院成立.

$^{681}$所以其實香港的變化就是這樣.
神給我們有機會去做這麼多事.
97年內臨.
有二百多個教會在1984年簽署信念書.
而1997年到2025年.
這28年香港華人教會的光譜又有變化.
我講得很快.
目前會有人數最多的是哪一間呢.
播道會恩福堂在長沙灣那間.
之前你都未聽過它的名字.
宣道會北角堂都有好幾千人.
還有荃灣的611靈良堂.
這些全部是新興的.
九龍城浸信會人數就跌了一點.
你會發覺光譜是不斷變的.
簡單來說你做事就有進步.
你不做事就停滯.
這個道理來的.
我就無意將人數多少.
來作為衡量一個教會是否屬靈上健康.
我不做這件事.
不過我們一定要承認.
時代是在變的.
你的屬靈需要都在變的.
你不迎合那個需要就不行.
當然耶穌基督的福音內容是不應該變的.
你不是去賣東西.
我們是傳福音.
但你傳息的方法.
你切入點.
都因應福音的對象.
你有調節的.
你在洪水橋做的事就不同了.
剛才周牧師跟我說.
他們有很多新的事工在做.
我想這些是我們要想的.
究竟香港華人教會怎樣走下去呢.
聽了很多東西.
但我有感動跟大家說這些委婉的事.
原因是你要跟我一起承繼182年的歷史.

$^{721}$你才可以承先啟後.
第三就要說更大的異象.
我未讀這本書之前.
我就是剛才這樣想.
如果我幫孫德堂.
我就想孫德堂的需要.
全世界或全香港教會.
就是用孫德堂做起點.
或者從元朗教會來想.
或者從宣道會來想.
但這幾年.
我們香港整體教會有很多變化.
但我沒有宏觀的視野.
這就是我看完這本書之後.
我最大的發現.
香港華人教會由1843到2025年.
182年之間.
都經歷了無數的危機.
現在仍然繼續前進.
現在我們都前進.
你不前進就應該是.
你在歷史裡面成為一個.
你不能夠扭轉乾坤的人.
《馬太福音》第九章三十五節說.
耶穌走遍各城各鄉.
在他們的會堂裡面教導人.
宣講天國的福音.
有醫治各樣的病症.
道濟就是傳道.
理濟就是教導.
醫治就是理濟.
這三樣東西都是要做的.
變成你發覺.
做教會的工作這麼多東西要做.
是的.
過去香港和中國的教會前輩.
就是在解釋耶穌三濟的事業.
他就是身體力行.
所以他們很沒空.
他們很多事情要做.

$^{761}$而且做得很開心.
我就認為香港華人教會.
要真的承先啟後.
你繼續做多些傳道工作.
教導工作.
還有支持醫治的工作.
香港的醫療水平很高.
但是疾病有身體有心靈.
這麼多醫生可以治身體的疾病.
但是心靈的病誰去治呢?.
輔導和屬靈指導是很需要的.
我們應該鼓勵那些有恩賜的弟兄姊妹.
有能力的同道.
你用耶穌基督的心腸去服侍有需要的人.
其實到處都有需要的.
我看回我們宣導會.
宣導會香港區年會是2024年的年報.
我們有121間的堂會.
都不算少的.
我們也就是49年之後才出現的.
會有人數名單上有五萬多人.
很威風的.
不過平均崇拜出席是二萬多.
二萬八就是打半了.
所以我們其實應該要想這些問題.
宣導會香港區年會在2021到2025年.
有一個五年計劃的主題.
承擔時代使命開拓福音疆界.
每一個宗派都有自己的獨特性.
大家各有特色.
可以彼此學習.
每間堂會的牧師和信徒領袖.
當他們覺得看不清楚前路的時候.
反過來想這些宗派領袖都不是看得很清.
大家一樣的模糊.
一起在香港互嚇互嚇.
但這個時候信心和意象是很重要的.
舉一個例子.
我們見到神學院在1997年之前.
當時的院長決定立根在香港.

$^{801}$所以就要重修禮堂.
現在你去長洲.
我們的禮堂是後來建造出來的.
這是他們的行動.
表明做這件事.
在疫情期間.
大家可能有上過神學院的無牆教室.
沒有上過都聽過.
這件事是什麼呢.
當時很多同道覺得要服侍弟兄姊妹.
究竟香港教會應該有什麼異象呢.
我從這本書得到提醒.
香港就是中國宣教的踏腳石.
大家覺得這個題目很敏感.
我就說定點.
實踐異象的方式可以很有彈性.
你不用太害怕.
不過我們是不可以沒有異象的.
你沒有異象你為什麼存在呢.
你明白我說什麼嗎.
這是最根本的問題.
你搞這麼多東西.
原來你沒有一個從天而來的異象.
你可不可以持續呢.
不可以的.
所以凡是從神而來的異象.
可能是很困難很具挑戰性.
但是如果是神給你的.
你就要去擔起它.
如果香港教會變成提供屬靈娛樂.
或者變成同鄉會.
那會怎樣呢.
是失去感染力的.
我2000年開始幫忙在錦首堂.
做義務的傳道人.
我記得在那裡.
在2001年已經開始說.
我們不應該推崇明星效應.
我們不應該走屬靈娛樂化的道路.
我太太就跟我說.

$^{841}$你這次又是說這些嗎.
我說是呀.
因為我覺得底因子媒好像還沒有學到.
她就說其實你怕不怕人家會悶.
我說不是呀.
他們學到我就不說了.
但是我自己都掙扎.
我太太每次都提醒我.
你又說這些.
你經常都說這些.
你明白我說什麼嗎.
所以後來我也很不懂怎樣說.
但是我自己是想說一個例子.
很想給大家去想一想.
當我.
我以前也有去做戒毒實習.
當我是一個吸毒者.
我信了主之後.
可能以前我的手下有幾十人幾百人.
但是我信了主之後.
我很謙卑.
我去學校做校工.
很多這樣的情況.
做校工就被人.
三叔你去廁所洗乾淨.
你明白嗎.
我叫我的弟子去做.
現在我自己下手去做.
那你是不是要很有生命.
我又不可以發火.
你叫我就好.
我馬上去做.
你明白嗎.
那在報道會.
如果教會辦報道會.
會不會請這個悔改了的郭鴻彪去做見證的弟兄呢.
那就有人打個問號.
不是不好.
不過寫他的名字可能沒什麼人來.
你明白嗎.

$^{881}$不是的.
我寫某個人的名字.
就有收視率的保證.
那為什麼呢.
我們教會辦這些活動.
是不是想多些人來.
是嘛.
所以我就找個出名的人.
大家明白我說什麼嗎.
這個做法很自然.
但是你背後就有一些價值取向.
價值取向的問題出在哪裡呢.
我就將問題攤出來.
究竟一個以前是很敗壞的人.
他信主之後生命改變了的人.
他的屬靈生命.
和這個很有名望的人的屬靈生命.
在神的面前.
神怎麼看呢.
你明白我說什麼嗎.
問題的嚴重性就在這裡.
在神的眼中是一樣的.
不過在人的眼中就可以不一樣.
一做起事來.
我當然要過做得野的.
你明白我的矛盾就在這裡.
其實如果這樣說的時候.
我們上帝就不應該.
導成人生在馬槽出世.
在牧羊人.
一個牧工的家庭裡面.
我們怎麼不找大祭司的家就好了.
這個是困擾我的問題.
所以如果上帝要做事.
要靠這樣的話.
就不是福音.
明明福音不是這樣的.
所以我就很煩躁.
不過我馬上告訴大家.
我們真的要明白.

$^{921}$上帝做事是不需要很厲害的人.
我沒有任何敵意對任何厲害的人.
而是說上帝做事.
就是在鷹械和吃奶的人.
去顯出他的能力.
所以我就想.
如果是今天的時候.
我怎麼看教會的未來呢.
我怎麼去做呢.
唯有學耶穌基督的人生.
就去走遍各城各鄉.
你去報道傳道.
你去教人.
你能夠治療人就治療.
你明不明白我講什麼.
就是做這些前線的事情.
所以我希望大家再次去想.
有些事情我們是解決不了的.
有些事情我們真的不太知道.
但是神感動我們去做的時候.
就快點去做.
因為你想走遍各城各鄉.
你以為你想走就行了嗎.
你明不明白.
我今天坐曾祿師的車.
我上車下車都要像個阿伯一樣.
不過我真的是阿伯.
我拿了旅遊卡.
為什麼呢.
因為我的筋腱.
就是髖關節那裡.
是要慢的.
大家明白嗎.
有很多大家的病友.
很明顯.
我太大的動作.
我是不可以快的.
我看到有人上車.
就一下子坐下.
哇 你這麼厲害.

$^{961}$我年輕都可以.
不過現在不行的意思.
你明白嗎.
你想走動.
Hallelujah.
不到你去勉強的.
所以我們可以走得走得的時候.
就做多一點.
是不是這樣說.
如果你健康許可.
就做多一點.
所以曾祿師.
他真的很好.
他知道我長洲出來.
我經常都說洪恩堂是很好的.
我都經常來的.
不過現在年紀大.
又長洲出.
又要搭去元朗.
又要怎麼辦.
他就義無反顧.
天水圍站我來接你.
我聽完就很開心.
你明白嗎.
我都會有軟弱的.
這麼長時間.
我老實說.
交通時間長過港島時間.
我就來回交通時間.
你明白嗎.
真的 不是開玩笑的.
所以我會有掙扎的.
其實我會有掙扎.
其實誰不想舒服一點.
是不是這樣說.
我就開始想問題.
看完那本書.
那些宣教士其實都很辛苦的.
當年沒有車坐.
或者怎樣說.

$^{1001}$你明白我說什麼嗎.
如果去到保羅的時候.
更加沒有.
是不是要行山越嶺.
你怎麼搞.
所以我就會想.
其實我們真的可以做的.
就盡做.
剛才曾祿師帶我車走.
看到新的屋落成了.
在附近.
其實當年是.
很感恩我們.
錦繡堂有些牧者.
信徒領袖.
覺得要在洪水橋直走.
其實是很艱難的.
很不容易.
洪水橋沒什麼人.
但曾祿師跟我說.
沿途找個地方租來做教會.
不是這麼簡單.
不是你沒本事租到.
沒錢也租不到.
何況沒錢.
所以.
你明白我說什麼嗎.
你能夠有一個這樣的地方.
這麼好的.
就在這裡去崇拜.
全福音.
有個社區中心.
我進來之後.
有東西喝.
羅漢果水.
你明白嗎.
還想要求什麼呢.
很難.
這裡不是旺角.
是一個很難的情況.

$^{1041}$但究竟神給大家有什麼託付呢.
這個就留給.
曾祿師和牧者.
和信徒領袖.
一起去想.
我真的覺得.
這次來我是很開心的.
一點不覺得怨.
一點不覺得累.
老實說.
還有我會覺得.
前人的努力.
是一定要去肯定的.
你不肯定前人的努力.
其實是很不應該的.
所以同一道理.
為什麼我要說182年的歷史.
這麼長氣呢.
就是因為沒有第一個來香港.
或者來中國的宣教士.
你哪有今天這麼多的基督徒呢.
如果我們忘記了他們的辛苦.
我覺得這樣對不上上帝.
我很想和大家說這些.
一起有.
希望神和你們說話.
就是讓大家.
你一起負起.
神託付給你們.
在這裡的使命.
我們一起同心祈禱.
親愛的天父上帝.
因為你是願意拯救世人的主.
你差遣你的獨生子.
耶穌基督來到世上.
並且他走遍各城各鄉.
傳道教道去醫治.
他甚至是甘心樂意走上十字架.
為人去寫命.
我們的主我們的神.

$^{1081}$過去世界有很多宣教士.
來將福音傳開.
我們得到這個福分.
到今天我們求你.
讓我們不要忘記他們的精神.
求你幫助我們去學習他們的精神.
堅持一生去跟隨他們.
有些甚至是死在中國的土地.
墳墓也在中國.
或者他們到老都沒有回家鄉.
亦希望我們仍然有這種心思意念.
我們學習跟隨.
哪怕我們做一些一般的工作.
但求主讓我們的信仰堅定.
讓我們不要白費時間.
不要白佔土地.
讓我們不要空手回天府.
好像太羽的姐妹唱的詩歌提醒我們.
不要空手回去見上帝.
我們的禱告.
奉耶穌基督的名.
阿們.
\newpage



\section{馬太福音 9:35}
\label{sec:ZmW9HFGb5CU}
\textbf{2025.09.28《耶穌的「三濟」事業:教導的「理濟」傳道的「道濟」,醫治的「利濟」》馬太福音9\_35 (和修版) 郭鴻標牧師}
\newline
\newline
連結: \href{https://youtube.com/watch?v=ZmW9HFGb5CU}{\texttt{https://youtube.com/watch?v=ZmW9HFGb5CU}} ~~~~ 語音日期: 2025-10-01
\newline
\newline
\hyperref[sec:WCmVYVwec7g]{\small{< < < PREV SERMON < < <}}
~
\hyperref[sec:index_chronic]{\small{[返順時目]}}
~
\hyperref[sec:index_scriptual]{\small{[返順卷目]}}
~
\hyperref[sec:j0UisY02_wA]{\small{> > > NEXT SERMON > > >}}
\newline
\newline
馬太福音 9:35
\newline
\begin{longtable}{cl}
\hline
\hline
章節 & 經文 (和合本修訂版)\\
\hline
9:35 & \begin{tabularx}{0.7\textwidth}{X} 耶穌走遍各城各鄉，在他們的會堂裡教導人，宣講天國的福音，又醫治各樣的病症。 \end{tabularx} \\ \\
[1ex]
\hline
\hline
\end{longtable}
$^{1}$各位弟兄姊妹早安.
我們一起同心祈禱.
親愛的天父上帝.
願你接納我們從心裡發出的敬拜.
願你的聖靈在我們心裡面.
向我們說話.
將你的話語植根在我們心靈裡.
亦都讓我們在今天聽到的時候.
是可以更加親近我們的神.
祈禱奉耶穌基督的名,阿們.
話說上個星期曾銳貴牧師聯絡我.
他說今個星期你提到的題目和經文是甚麼.
我就看看日曆上沒有這個紀錄.
然後就一直思想答應還是不答應呢.
曾牧師說你不答應都無所謂.
但我又覺得不好.
問題就來了,我是找甚麼講道理來紅恩堂呢.
不怕告訴大家我抽屜底有很多舊貨.
在電腦裡面儲起來.
但我又覺得這樣不是很尊重曾銳貴牧師.
神給我一個感動.
因為我們都有靈修,祈禱,有閱讀.
剛好我正在看一本書,看完了.
我心裡有感動.
不如就將那些靈感整理成為講章.
如果日後用,第一次都會在紅恩堂用.
這個講道理的起源是甚麼呢.
就是在疫情之後.
我們其實都在問一個問題.
香港的教會何去何從.
我們每個教會都問這個問題.
我估計和很多同道都有這個特點.
我是先從我自己的教會處境開始思考.
我教會有甚麼難處,有甚麼需要.
然後再看看我見到的那些教會情況.
很正常的,如果在洪水橋區.
你就想想洪水橋區的環境.
然後再想想教會的前路.
我看完那本書,我發覺完全改變了我.
那本書是甚麼呢.

$^{41}$那本書我沒有帶來,都幾後行.
就是香港基督教史1907到1997.
大家喜不喜歡歷史.
我一會兒就和你說故事.
就是2024年出版的.
其中有作者送了給我.
你知道有人送書給你,你都多多謝謝.
有時不經意就送了去不知哪本書.
但神又感動我,我真的拿來看這麼後的書.
我真的逐一看完.
這本書是講宣教士來中國的歷史.
如果你有少少印象,你知道應該何時開始說.
是主後七世紀的列斯多流.
來到中國,就是大秦景教碑開始說.
不是這樣說,你都要說16世紀的時候天主教的傳教士來華.
就是李馬竇,湯若望,南懷仁那些.
但他都沒有提.
他也沒有再講97年之後香港教會的情況.
但這本書都是有很特別的.
我自己覺得是,我得到很大的益處.
他就嘗試將香港的教會擺在.
由1807年馬勒信來華的脈絡底下去理解.
來展現一幅很宏觀的圖畫.
我今天沒有powerpoint,所以你要聽我說話了.
這本書有七章,我就用我自己的文字去表達.
我先講一下大概七部分是怎樣.
就是早期的我們叫更正教的宣教.
即是我們中文叫基督教,其實是Protestant.
就是從馬勒信來華開始,1807年.
然後就是鴉片戰爭之後的宣教.
就是1842年到1883年.
這些是很久的歷史.
然後這本書很有意思.
他就說,不只是講西教事.
中國教會領袖的出現.
就是由1884年到1914年.
即是那時候已經有華人的牧師.
然後去到進入自治的時期.
1915年到1940年.
即是華人的教會一直成長.

$^{81}$去到日本侵華時期.
就是1941到1945年.
不知道在座中有沒有人經歷過這段時間呢?.
你可能就會明白多一點.
然後就是香港教會本土化的時期.
1945到1973.
在座中可能有些人都出世了.
然後去到香港教會的增長期.
1974到1997年.
我們很多人都是在正正香港教會增長期的時候.
來去成長,來去信主.
而我將這七個時期變成兩個部分.
就是今天講到的第一第二部分.
就是說培養華人的領袖.
第二就是戰後對有需要的人的服侍.
這兩樣東西就構成今天我們香港萬人教會的基礎.
然後在祈禱和思考裡面.
我發覺一句話很特別.
就是作者,其中一個作者.
應該就是李志剛牧師,老前輩.
他就說,他用馬太方第九章三十五節的經文.
去講過去的中國和香港華人教會的宣教士和歷史.
可以說是根據著耶穌的三濟事業.
這些字眼就很深很古老.
我問一個問題.
有沒有人聽過香港有一間道濟會堂.
就是中華基督教會合一堂.
就是潘咸道那間合一堂.
他的前身就叫道濟會堂.
其實你不太知道是什麼.
你可能以為是道堂.
其實不是.
因為原來在早期的宣教士在那裡講.
耶穌三濟事業就包括教道的理濟.
教道教道道理.
用道理來濟世.
傳道.
那個是道濟.
即是我們傳福音.
你回帶回去一百四十幾年的時候.

$^{121}$原來那些宣教士是在講傳道.
就是我們去道濟.
讓其他人得益.
醫治就是利濟.
即是辦醫療所那些.
原來就是這些西教士一直下來就在做這些事.
所以我講到的三部分.
第一就是怎樣培養華人的領袖.
那時候沒有什麼華人基督徒.
第二就是戰後對於有需要的人的服侍.
然後第三就是.
今天我們要關心的更大的異象.
當我們問今天香港教會何去何從的時候.
你就要從根開始講起.
首先第一要講培養華人教會的領袖.
馬來西亞是一百零七年來.
他來的時候主要在澳門和廣州.
廣州其實是不讓人進去的.
所以那些西教士就停留在澳門.
當時是清朝的時代.
就禁止傳教.
馬來西亞就轉移到南洋宣教.
開展教育出版.
所以馬六甲那裡做這些事.
然後一百四十二年有南京條約.
割讓了香港.
就開放五口通商.
大家都記得這些.
還有一百四十四年望下條約和黃埔條約.
就容許宣教士在這五個口岸裡面.
建教堂.
辦醫院.
傳福音.
就是說五口通商不是只做生意.
而是後面望下條約.
黃埔條約就讓他們去買地建教堂.
辦醫院.
傳福音.
所以歷史是這麼走的.
然後一百四十二年的宣教士來香港.

$^{161}$他們的目的不是只在香港做生意.
而是他們向內地傳福音.
你看見這個歷史的時候.
原來澳門也是一樣的.
當南京條約之後.
駐澳門的所有的天主教宣教士.
基督教宣教士.
大部分都調走了.
所以澳門的基督教的發展是比較弱.
原因是警察會將他們調走了.
去上海那裡.
所以香港也是一樣.
香港的宣教士.
他們是踏足香港作為一個跳板.
而去國內傳福音.
你要明白這個歷史.
我看完之後也很震撼.
原來我們的角色就是這樣.
大家知不知道香港第一間教會是哪一間.
很有趣.
就是1843年7月17日.
在皇后.
皇后浸信會.
就是皇后大道中那裡.
就是今天的建浸.
香港浸信會.
長洲浸信會也是1843年開始主日崇拜.
故事就是美國浸信會的一個叔尾士穆斯夫婦.
在1836年去了澳門.
然後1842年.
不知道他們是否收到消息.
所以搬來香港.
準備建教堂.
為什麼1843年可以買到地建教堂.
早一年已經有這個計劃.
1843年建教堂.
今天有些泰國的姐妹在這裡.
有一個叫倫維仁穆斯.
叫William Dean.
1834年去到泰國曼谷宣教.

$^{201}$你們聽過吧.
真的很厲害.
在曼谷建立第一間華人教會.
這個真是歷史.
當時在泰國有很多潮州人.
在座中有沒有人是潮州人.
有潮州人.
很明顯你們知道這個歷史.
倫維仁穆斯就在泰國曼谷學潮州話.
這個只是故事的開始.
1842年.
倫維仁穆斯就來香港.
1843年在街市街的浸信會.
就開始第一間潮語浸信會.
不只是這樣.
他帶著潮州基督徒去長洲.
向潮州人傳福音.
所以為什麼1843年長洲浸信會開始有聚會.
因為泰國有些基督徒來香港宣教.
所以剛好有泰國的姐妹.
這個歷史不是假的.
可能很多人是不會知道這件事.
故事就是香港還沒有成為一個可以傳教的地方的時候.
宣教士就在東南亞做了很多事.
一時到那裡那些基督徒就來香港幫忙宣教.
你想一下西人怎麼會這麼容易做到這麼多事.
因為他們是華人的助手.
另外一個叫倫敦傳道會的馬勵信穆斯.
在馬六甲辦英華書院.
不知道在座有沒有人是英華的校友.
成為宣教的總部.
印中文聖經.
所以其實宣教士做不到事就是印聖經.
只是他做預備功夫.
1843年倫敦傳道會駐馬六甲的九位宣教士.
全部調過來香港和中國.
兩個去了福州.
兩個去廈門.
三個去上海和寧波.
有兩個去香港.

$^{241}$所以你看到部署就是.
九位宣教士只有兩個在香港.
香港只是一個其中的宣教點.
其實目標是全中國的.
1843年倫敦傳道會在上環成立道濟會堂.
就是今天的中華基督教會合一堂.
然後1863年又成立了後來叫中華基督教會灣仔堂.
就是在合會中心那一間.
可能你會說還有什麼呢.
原來有間叫公理堂.
是的,這個由於他們的差會就是美部會.
美部會在1883年就成立公理堂.
聖公會在香港第一間華人教會就是聖士提反堂.
就在1864年.
所以你看看華人教會成立就是這樣的局力程.
1852年巴釋會就在西應盤開機.
就是今天的崇禎會舊恩堂.
在座有沒有一些客家人在呢.
崇禎會就是客家人的.
1862年就在筲箕灣開始.
又開了一間教堂就是崇禎會筲箕灣堂.
那1842到1883年呢.
這40年倫敦傳道會成立了後來的中華基督教會合一堂.
中華基督教會灣仔堂.
還有聖公會聖士提反堂.
巴釋會成立的崇禎會舊恩堂.
筲箕灣堂.
估計華人信徒只有1000人.
40年只有1000個基督徒.
當時有三大公會.
倫敦傳道會.
聖公會和巴釋會.
都被香港政府批准在博庫林.
這個現在叫華人基督教永遠墳場.
就是你要多謝那些前輩.
就是我們的上一代.
信了主然後就葬在博庫林華人基督教永遠墳場.
就是因為當時他們三大公會和政府.
說不如拿個山頭給我們.
宣教士的宣教工作其實做很多事.

$^{281}$又要翻譯聖經.
又要開醫院.
又要辦學校.
又要搞社會服務.
其實真的很多事要搞.
而宣教士還有呢.
我要數給大家聽.
你知不知道中國過去有13間基督教大學.
很厲害的.
我說些名字出來可能你聽過.
金陵大學.
金陵女子文理學院.
開女子大學.
你知不知道是一個什麼意思.
是想男女平等.
所以基督教的宣教士做了很多事.
聖約翰大學.
上海的很出名.
滬江大學.
東吳大學.
之江大學.
燕京大學.
就是後來的北京大學.
華中大學.
齊魯大學.
華西協合大學.
福建協和大學.
華南女子文理學院.
還有嶺南大學.
嶺南就下來香港了.
13間.
1884到1914.
現在說到另一邊.
那30年又看看香港.
有以下的教會成立了.
我只是說華人教會.
只是說幾十年才有這麼多間.
不是變數字出來.
1884年有中華巡渡會.
就是巡渡會的人來.

$^{321}$1890年聖公會聖三一堂.
就是九龍城那間.
1890年就是九龍城浸信會.
經常去陪靈會那裡.
1890年就是聖公會朱聖堂.
在旺角那裡.
1892年就是中華基督教會深水寶堂.
然後1896年就是皇兒周宣崇真會.
1898年就是中華基督教會元朗堂.
這是全新界第一間華人教會.
就是搭西鐵是前幾個站.
然後1899年就是禮賢會香港堂.
就是很後才開的.
1901年尖沙咀浸信會.
1903年就是中華基督教會聖光堂.
1903年就是中國基督徒會.
1905年就是中華基督教會荃灣堂.
1905年就是香港仔浸信會.
你數下來,嘩,這麼厲害.
1905年就是粉嶺崇天堂.
1905年就是窩尾崇真會.
1907年就是中華基督教會大埔堂.
1907年就是五秦節會.
1907年就是呂廣朝人中華基督教會.
1911年就是聖公會聖瑪利亞堂.
好像廟那間在銅鑼灣那裡.
1911年就是聖公會聖保樂.
就是,是啦.
1912年就是海面傳道會.
專向漁民傳福音的.
在在中有沒有是,就是,水上人.
有,你們知道,海面傳道會.
我聽說那些牧師很有型.
那些船民,不是,那些漁民.
出海可能是一個月.
去很遠的地方打魚.
那些牧師就上去跟他們祈禱.
打完有魚獲回來.
就一打,其實牧師有條鹹魚給你.
很靚的鹹魚,要醃了的.

$^{361}$那些人是出很多風險的.
他們不拜其他的偶像.
而信耶穌,你明不明.
他是很不容易的.
好了,還有.
1914年是便義理會.
有沒有人聽過便義理會.
然後就,就是,這些教會.
說的也是過了30年.
40年,30年,70年.
到1915年,香港華人基督教會成立.
其實,原來是慢慢才有的.
培靈會什麼時候開的?1927年.
在廣州開,然後香港跟著做.
這樣的故事.
如果你說1843年是香港第一間華人教會.
到現在2025年.
說的是182年歷史.
大家有沒有想過,我們是在承傳182年的歷史.
包袱重了很多.
上帝提醒我,當我去問教會未來方向的問題.
就不只是好像我開始說.
我自己教會個別的堂會的難題.
或者這幾年的危機去想.
這樣想是斬斷了182年的歷史.
你明不明白我說什麼?.
所以我發覺,神提醒我.
我們是要有歷史意識.
香港教會不是今天才存在的.
而是已經有182年的歷史.
當你有這個思維的時候.
你就會知道.
我們的香港教會存在.
是一直跟中國的宣教是結連的.
宣教是這樣的.
教會從來不是為自己存在的.
教會也不只是提供宗教活動.
使我們覺得自己有存在感.
你明不明白?.
我們不是為自己存在的.

$^{401}$我們是有一個更大的使命.
究竟香港教會何去何從呢?.
就要看個別情況了.
不過你回顧過去的香港教會歷史.
是可以看到上帝是怎樣在這裡工作.
我們可以怎樣去跟隨.
過去宣教士做了很多事情.
同時培訓了非常多的華人牧者接班.
過去182年香港教會開展了很多事工.
栽培領袖.
直到今天這個地步.
你想想前人做的事情.
1883年香港有三大工會.
1914年加上浸信會.
就有四大工會.
即是什麼呢?.
當時宣道會,播道會都沒有名字.
這樣說是歷史事實.
在這樣的大光譜下.
去了解不同宗教的發展歷史.
不是要比較.
不是誰厲害,誰不厲害.
不是角力.
而是要承傳前人開方,火熱,學苦的精神.
如果我們沒有這個眼界.
我們其實是搞來搞去都是搞一下.
我們要做一些事情是承先啟後的.
第二,講戰後對有需要人的服侍.
講近一點.
1937年7月7日發生盧溝橋事變.
大家都應該知道這些中國歷史.
往後八年裡面中國很多的城市都淪陷.
1937年是天津,上海和南京.
1938年是廣州和廈門.
1842年是溫州和南昌.
1945年是長沙.
到1941年香港淪陷.
日治時期的香港社會出現很嚴重的經濟壓力.
原來1941年推出軍票.
在座中不知道有沒有人有軍票.

$^{441}$不知道怎麼搞的.
當時是說一元兌兩元港幣.
即是初時.
然後1842年又推出軍票.
今次是一元兌四元港幣.
究竟早換還是遲換.
然後到1843年.
日軍總督部禁止軍票以外的一切港元.
不可以用來做貿易.
你家裡有港幣都抓你.
所有香港市民一定要把港幣換成軍票.
很多人的資產一夜貶值.
所以現在很多人都說軍票放在那裡不知道做什麼.
要求賠償.
日治時期的香港教會出現大量信徒流失.
你都想到的.
肯定要走.
香港牧者走不走.
走.
所以那時候很艱難.
而1942年日軍總督部下令所有宗教團體要登記.
什麼教都要登記.
基督教成立了香港基督教的總會.
在那段日子很艱難.
社會不穩定.
人心惶惶.
教會領袖受到監控.
在各種壓力下.
繼續做關心牧羊的工作.
我看了另外一些傳記.
當時的牧者和信徒.
他出去探訪多數是危難的.
第二就是無端端被炸彈炸死.
要替他辦喪禮.
你想想.
那時候做這些事情很艱難.
然後呢.
1945年8月香港重光.
先不要說那麼多黑暗的事情.
1948年香港華人基督教會成立.

$^{481}$講一下國光譜又不同了.
1949年之前香港有七大工會.
倫敦會,聖公會,巴士會,浸信會,禮賢會.
突然間上了榜.
巡渡會又上榜.
公理會.
由於倫敦會和公理會加入了中華基督教會.
所以從此就有六大工會.
好像武林,武當,少林那些.
這個六大工會的格局就出來了.
成為中華基督教會,聖公會,巴士會,浸信會,禮賢會,巡渡會.
到1949年之後又怎樣呢?.
國內很多的.
就是警察會就沒有撤退.
但是就將他們的檔案或者資產那些東西.
在香港成立另一個辦事處.
他們是很想仍然留在中國的.
但是1950年.
警察會西人宣教士撤出大陸來香港是大勢所趨的.
你不走都不行的.
1951年香港的中派教會就要脫離和國內教會的關係.
所以大家劃清界線了.
然後封關那些大家有些印象了.
然後1949年大量的難民湧入香港.
就是我上一代的父母都是這樣的.
很多來到香港的宣教士覺得香港政局不穩.
有些警察會就說.
其實我準備調走所有香港的宣教士.
就是1949年之後應該是全部調走的.
但是有些宣教士就覺得不行.
我的心是很難過的.
我看到條頸嶺這麼多難民我要走去開荒的.
所以他們就在條頸嶺去做福音工作.
嚴格來說是這班不聽話的宣教士留在香港來做福音工作的.
歷史是真的這樣的.
就是警察會調你去印尼你不肯去.
你自己去做事.
我回鄉下籌款來做.
是這樣的.
這個動盪的歲月.

$^{521}$當然在韓戰之後.
中國是沒有用武力收回香港.
其實1949年的兵去到羅湖沒有過來.
警察會就覺得.
似乎都可以在香港宣教.
就開始派一些宣教士過來香港.
就是1953年之後才開始有這樣的情況.
所以在1956年.
香港教會名錄有50個警察會.
在香港有辦事處.
估計在香港的宣教士有300個.
就是未到60年代的時候有300個西教士.
好了.
就是1950年的時候.
新中國那個.
就是50年底.
他要求國內教會一定要肅清帝國主義的餘毒.
就斷絕和帝國主義的一切關係.
所以全部都變形了.
香港教會就變成什麼.
他可以繼續發展.
1962年有60個中派和獨立的堂會.
你看看,起飛了.
然後有344間教會.
大概有11萬的信徒.
剛才說了1000,現在11萬.
很特別,是不是.
到了1960年.
這班難民變成了新來港人士.
面對著大批的新來港人士.
教會就要問,我在這裡有什麼使命.
就是你做不做難民工作.
是有分別的.
58年大躍進出現大饑荒.
大家都記得了,有逃亡潮.
這個時候香港的社會是風雨飄搖.
經歷了50年,51年的木屋區火災.
36000人沒有家園.
石硤尾大火.
有6萬災民.

$^{561}$你可以記得這些東西.
大量學童是失學.
流落街頭.
1961年大概有8萬個6到11歲的失學兒童.
1956年,九龍的荃灣.
發生暴動,有激烈的衝突.
這些是故事.
在這樣的情況下.
社會福利方面就因為要推動救濟工作.
所以政府就容許教會去辦學校.
1966年,在香港教會名錄裡面.
更正教的小學和幼稚園有225間.
所以可能我們讀書就在這些學校讀.
1955年.
香港政府批准.
宗教和志願團體在.
七層高的寫字大廈辦天台學校.
有沒有人讀過天台學校.
有,我也是.
我讀幼稚園.
我那時候是8元學費.
我家人就說.
其實我們給不起學費.
基本上你沒得讀書.
不過教會辦學很好.
他給我十幾元.
他叫獎學金.
哪是獎學金.
叫助學金.
沒那十幾元,交不到學費我就不用讀書.
所以很多西教士就做了這些事.
當時有11間由教會辦的天台學校.
所以1955到1958年.
香港教會增加到294間.
去到7萬4470個弟兄姊妹.
而到58到62年.
去到344間.
有11萬2200個基督徒.
所以你看到那時候.
香港教會是這樣發展的.

$^{601}$而62年香港教會的光譜又變了.
出現了八大工會.
很有趣的.
你數一下有多少間堂會.
排頭的就是浸信會.
當時有54間.
會有1萬5千多.
為什麼突然間戰後浸信會這麼厲害呢.
中華基督教會只有24間教會.
只有1萬多個教友.
聖姑爺廟23間堂.
就1萬3千多.
信義會是33間堂.
就1萬多個教友.
路德會18間堂.
有9千多人.
秦道公會和惠利公會.
8間就有7千多人.
崇禎會11間就有5千多人.
禮賢會6間就有5千多人.
你發覺這個光譜又變化了.
你要追下去問.
為什麼浸信會這麼厲害呢.
這些就是另一個題目.
今天就講不到了.
八大工會已經佔全港堂會51\%.
會有人數是73\%.
當然他們很勤力去傳福音.
才有這樣的結果.
而浸信會的增長是最特別.
成為香港最大的中派.
信義會路德會的發展.
是因為在難民裡面有很多外省人.
他們說不懂廣東話.
所以信義會和路德會.
外省牧師帶了很多外省信徒.
信主.
60年代的中期香港經濟好轉.
你還記不記得.
什麼時候有家可以買到電視.

$^{641}$可以買到洗衣機.
我都是六幾年的時候.
發覺以為可以.
我小時候去隔壁看別人的電視.
你們不知道什麼.
原來是這樣的.
60年代中期香港轉型了.
不是再活在救濟品年代.
國際救援機構都撤離香港了.
烏溪沙都不再做孤兒院了.
到了67年香港暴動.
有什麼呢.
貧富距離很嚴重.
開始香港華人教會注重基層勞工的工作了.
你看到的.
70年代是一個脫邊的年代.
因為戰後出生的那批人.
本土意識比上一代高.
我們上一代是過客.
但是另一代他們對香港有承擔.
所以大專的基督徒.
對香港和中國有承擔.
所以是不同的.
有人將當時的教會分成兩派.
一派叫普世派.
一派叫福音派.
普世派就說天國在人間.
我們做多點服務.
關心社會公義.
福音派就注重個人得救.
1966年平安福音堂成立.
中華基督教會就70年代大量建立學校.
其實福音派的基督徒都有關心社會文化的.
譬如有關心弱勢的群體.
譬如突破.
有人聽過突破機構.
工業福音團契.
香港基督徒福音學生福音團契.
神學訓練方面.
也有重視高等神學訓練的中國神學研究院成立.

$^{681}$所以其實香港的變化就是這樣.
神給我們有機會去做這麼多事.
97年內臨.
有二百多個教會在1984年簽署信念書.
而1997年到2025年.
這28年香港華人教會的光譜又有變化.
我講得很快.
目前會有人數最多的是哪一間呢.
播道會恩福堂在長沙灣那間.
之前你都未聽過它的名字.
宣道會北角堂都有好幾千人.
還有荃灣的611靈良堂.
這些全部是新興的.
九龍城浸信會人數就跌了一點.
你會發覺光譜是不斷變的.
簡單來說你做事就有進步.
你不做事就停滯.
這個道理來的.
我就無意將人數多少.
來作為衡量一個教會是否屬靈上健康.
我不做這件事.
不過我們一定要承認.
時代是在變的.
你的屬靈需要都在變的.
你不迎合那個需要就不行.
當然耶穌基督的福音內容是不應該變的.
你不是去賣東西.
我們是傳福音.
但你傳息的方法.
你切入點.
都因應福音的對象.
你有調節的.
你在洪水橋做的事就不同了.
剛才周牧師跟我說.
他們有很多新的事工在做.
我想這些是我們要想的.
究竟香港華人教會怎樣走下去呢.
聽了很多東西.
但我有感動跟大家說這些委婉的事.
原因是你要跟我一起承繼182年的歷史.

$^{721}$你才可以承先啟後.
第三就要說更大的異象.
我未讀這本書之前.
我就是剛才這樣想.
如果我幫孫德堂.
我就想孫德堂的需要.
全世界或全香港教會.
就是用孫德堂做起點.
或者從元朗教會來想.
或者從宣道會來想.
但這幾年.
我們香港整體教會有很多變化.
但我沒有宏觀的視野.
這就是我看完這本書之後.
我最大的發現.
香港華人教會由1843到2025年.
182年之間.
都經歷了無數的危機.
現在仍然繼續前進.
現在我們都前進.
你不前進就應該是.
你在歷史裡面成為一個.
你不能夠扭轉乾坤的人.
《馬太福音》第九章三十五節說.
耶穌走遍各城各鄉.
在他們的會堂裡面教導人.
宣講天國的福音.
有醫治各樣的病症.
道濟就是傳道.
理濟就是教導.
醫治就是理濟.
這三樣東西都是要做的.
變成你發覺.
做教會的工作這麼多東西要做.
是的.
過去香港和中國的教會前輩.
就是在解釋耶穌三濟的事業.
他就是身體力行.
所以他們很沒空.
他們很多事情要做.

$^{761}$而且做得很開心.
我就認為香港華人教會.
要真的承先啟後.
你繼續做多些傳道工作.
教導工作.
還有支持醫治的工作.
香港的醫療水平很高.
但是疾病有身體有心靈.
這麼多醫生可以治身體的疾病.
但是心靈的病誰去治呢?.
輔導和屬靈指導是很需要的.
我們應該鼓勵那些有恩賜的弟兄姊妹.
有能力的同道.
你用耶穌基督的心腸去服侍有需要的人.
其實到處都有需要的.
我看回我們宣導會.
宣導會香港區年會是2024年的年報.
我們有121間的堂會.
都不算少的.
我們也就是49年之後才出現的.
會有人數名單上有五萬多人.
很威風的.
不過平均崇拜出席是二萬多.
二萬八就是打半了.
所以我們其實應該要想這些問題.
宣導會香港區年會在2021到2025年.
有一個五年計劃的主題.
承擔時代使命開拓福音疆界.
每一個宗派都有自己的獨特性.
大家各有特色.
可以彼此學習.
每間堂會的牧師和信徒領袖.
當他們覺得看不清楚前路的時候.
反過來想這些宗派領袖都不是看得很清.
大家一樣的模糊.
一起在香港互嚇互嚇.
但這個時候信心和意象是很重要的.
舉一個例子.
我們見到神學院在1997年之前.
當時的院長決定立根在香港.

$^{801}$所以就要重修禮堂.
現在你去長洲.
我們的禮堂是後來建造出來的.
這是他們的行動.
表明做這件事.
在疫情期間.
大家可能有上過神學院的無牆教室.
沒有上過都聽過.
這件事是什麼呢.
當時很多同道覺得要服侍弟兄姊妹.
究竟香港教會應該有什麼異象呢.
我從這本書得到提醒.
香港就是中國宣教的踏腳石.
大家覺得這個題目很敏感.
我就說定點.
實踐異象的方式可以很有彈性.
你不用太害怕.
不過我們是不可以沒有異象的.
你沒有異象你為什麼存在呢.
你明白我說什麼嗎.
這是最根本的問題.
你搞這麼多東西.
原來你沒有一個從天而來的異象.
你可不可以持續呢.
不可以的.
所以凡是從神而來的異象.
可能是很困難很具挑戰性.
但是如果是神給你的.
你就要去擔起它.
如果香港教會變成提供屬靈娛樂.
或者變成同鄉會.
那會怎樣呢.
是失去感染力的.
我2000年開始幫忙在錦首堂.
做義務的傳道人.
我記得在那裡.
在2001年已經開始說.
我們不應該推崇明星效應.
我們不應該走屬靈娛樂化的道路.
我太太就跟我說.

$^{841}$你這次又是說這些嗎.
我說是呀.
因為我覺得底因子媒好像還沒有學到.
她就說其實你怕不怕人家會悶.
我說不是呀.
他們學到我就不說了.
但是我自己都掙扎.
我太太每次都提醒我.
你又說這些.
你經常都說這些.
你明白我說什麼嗎.
所以後來我也很不懂怎樣說.
但是我自己是想說一個例子.
很想給大家去想一想.
當我.
我以前也有去做戒毒實習.
當我是一個吸毒者.
我信了主之後.
可能以前我的手下有幾十人幾百人.
但是我信了主之後.
我很謙卑.
我去學校做校工.
很多這樣的情況.
做校工就被人.
三叔你去廁所洗乾淨.
你明白嗎.
我叫我的弟子去做.
現在我自己下手去做.
那你是不是要很有生命.
我又不可以發火.
你叫我就好.
我馬上去做.
你明白嗎.
那在報道會.
如果教會辦報道會.
會不會請這個悔改了的郭鴻彪去做見證的弟兄呢.
那就有人打個問號.
不是不好.
不過寫他的名字可能沒什麼人來.
你明白嗎.

$^{881}$不是的.
我寫某個人的名字.
就有收視率的保證.
那為什麼呢.
我們教會辦這些活動.
是不是想多些人來.
是嘛.
所以我就找個出名的人.
大家明白我說什麼嗎.
這個做法很自然.
但是你背後就有一些價值取向.
價值取向的問題出在哪裡呢.
我就將問題攤出來.
究竟一個以前是很敗壞的人.
他信主之後生命改變了的人.
他的屬靈生命.
和這個很有名望的人的屬靈生命.
在神的面前.
神怎麼看呢.
你明白我說什麼嗎.
問題的嚴重性就在這裡.
在神的眼中是一樣的.
不過在人的眼中就可以不一樣.
一做起事來.
我當然要過做得野的.
你明白我的矛盾就在這裡.
其實如果這樣說的時候.
我們上帝就不應該.
導成人生在馬槽出世.
在牧羊人.
一個牧工的家庭裡面.
我們怎麼不找大祭司的家就好了.
這個是困擾我的問題.
所以如果上帝要做事.
要靠這樣的話.
就不是福音.
明明福音不是這樣的.
所以我就很煩躁.
不過我馬上告訴大家.
我們真的要明白.

$^{921}$上帝做事是不需要很厲害的人.
我沒有任何敵意對任何厲害的人.
而是說上帝做事.
就是在鷹械和吃奶的人.
去顯出他的能力.
所以我就想.
如果是今天的時候.
我怎麼看教會的未來呢.
我怎麼去做呢.
唯有學耶穌基督的人生.
就去走遍各城各鄉.
你去報道傳道.
你去教人.
你能夠治療人就治療.
你明不明白我講什麼.
就是做這些前線的事情.
所以我希望大家再次去想.
有些事情我們是解決不了的.
有些事情我們真的不太知道.
但是神感動我們去做的時候.
就快點去做.
因為你想走遍各城各鄉.
你以為你想走就行了嗎.
你明不明白.
我今天坐曾祿師的車.
我上車下車都要像個阿伯一樣.
不過我真的是阿伯.
我拿了旅遊卡.
為什麼呢.
因為我的筋腱.
就是髖關節那裡.
是要慢的.
大家明白嗎.
有很多大家的病友.
很明顯.
我太大的動作.
我是不可以快的.
我看到有人上車.
就一下子坐下.
哇 你這麼厲害.

$^{961}$我年輕都可以.
不過現在不行的意思.
你明白嗎.
你想走動.
Hallelujah.
不到你去勉強的.
所以我們可以走得走得的時候.
就做多一點.
是不是這樣說.
如果你健康許可.
就做多一點.
所以曾祿師.
他真的很好.
他知道我長洲出來.
我經常都說洪恩堂是很好的.
我都經常來的.
不過現在年紀大.
又長洲出.
又要搭去元朗.
又要怎麼辦.
他就義無反顧.
天水圍站我來接你.
我聽完就很開心.
你明白嗎.
我都會有軟弱的.
這麼長時間.
我老實說.
交通時間長過港島時間.
我就來回交通時間.
你明白嗎.
真的 不是開玩笑的.
所以我會有掙扎的.
其實我會有掙扎.
其實誰不想舒服一點.
是不是這樣說.
我就開始想問題.
看完那本書.
那些宣教士其實都很辛苦的.
當年沒有車坐.
或者怎樣說.

$^{1001}$你明白我說什麼嗎.
如果去到保羅的時候.
更加沒有.
是不是要行山越嶺.
你怎麼搞.
所以我就會想.
其實我們真的可以做的.
就盡做.
剛才曾祿師帶我車走.
看到新的屋落成了.
在附近.
其實當年是.
很感恩我們.
錦繡堂有些牧者.
信徒領袖.
覺得要在洪水橋直走.
其實是很艱難的.
很不容易.
洪水橋沒什麼人.
但曾祿師跟我說.
沿途找個地方租來做教會.
不是這麼簡單.
不是你沒本事租到.
沒錢也租不到.
何況沒錢.
所以.
你明白我說什麼嗎.
你能夠有一個這樣的地方.
這麼好的.
就在這裡去崇拜.
全福音.
有個社區中心.
我進來之後.
有東西喝.
羅漢果水.
你明白嗎.
還想要求什麼呢.
很難.
這裡不是旺角.
是一個很難的情況.

$^{1041}$但究竟神給大家有什麼託付呢.
這個就留給.
曾祿師和牧者.
和信徒領袖.
一起去想.
我真的覺得.
這次來我是很開心的.
一點不覺得怨.
一點不覺得累.
老實說.
還有我會覺得.
前人的努力.
是一定要去肯定的.
你不肯定前人的努力.
其實是很不應該的.
所以同一道理.
為什麼我要說182年的歷史.
這麼長氣呢.
就是因為沒有第一個來香港.
或者來中國的宣教士.
你哪有今天這麼多的基督徒呢.
如果我們忘記了他們的辛苦.
我覺得這樣對不上上帝.
我很想和大家說這些.
一起有.
希望神和你們說話.
就是讓大家.
你一起負起.
神託付給你們.
在這裡的使命.
我們一起同心祈禱.
親愛的天父上帝.
因為你是願意拯救世人的主.
你差遣你的獨生子.
耶穌基督來到世上.
並且他走遍各城各鄉.
傳道教道去醫治.
他甚至是甘心樂意走上十字架.
為人去寫命.
我們的主我們的神.

$^{1081}$過去世界有很多宣教士.
來將福音傳開.
我們得到這個福分.
到今天我們求你.
讓我們不要忘記他們的精神.
求你幫助我們去學習他們的精神.
堅持一生去跟隨他們.
有些甚至是死在中國的土地.
墳墓也在中國.
或者他們到老都沒有回家鄉.
亦希望我們仍然有這種心思意念.
我們學習跟隨.
哪怕我們做一些一般的工作.
但求主讓我們的信仰堅定.
讓我們不要白費時間.
不要白佔土地.
讓我們不要空手回天府.
好像太羽的姐妹唱的詩歌提醒我們.
不要空手回去見上帝.
我們的禱告.
奉耶穌基督的名.
阿們.
\newpage



\section{以賽亞書 54:1-10}
\label{sec:j0UisY02_wA}
\textbf{2025.10.05《神在你荒涼中說話》:以賽亞書54\_1-10(和修版) 曾裔貴牧師}
\newline
\newline
連結: \href{https://youtube.com/watch?v=j0UisY02_wA}{\texttt{https://youtube.com/watch?v=j0UisY02\_wA}} ~~~~ 語音日期: 2025-10-07
\newline
\newline
\hyperref[sec:ZmW9HFGb5CU]{\small{< < < PREV SERMON < < <}}
~
\hyperref[sec:index_chronic]{\small{[返順時目]}}
~
\hyperref[sec:index_scriptual]{\small{[返順卷目]}}
~
\hyperref[sec:Bm5kMjeNAdQ]{\small{> > > NEXT SERMON > > >}}
\newline
\newline
以賽亞書 54:1-10
\newline
\begin{longtable}{cl}
\hline
\hline
章節 & 經文 (和合本修訂版)\\
\hline
54:1 & \begin{tabularx}{0.7\textwidth}{X} 你這不懷孕、不生育的，要歡呼；你這未曾經過產難的，要歡呼，揚聲呼喊；因為被遺棄的婦人，比有丈夫的人兒女更多；這是耶和華說的。 \end{tabularx} \\ \\ \relax
54:2 & \begin{tabularx}{0.7\textwidth}{X} 要擴張你帳幕之地，伸展你居所的幔子，不要縮回；要放長你的繩子，堅固你的橛子。 \end{tabularx} \\ \\ \relax
54:3 & \begin{tabularx}{0.7\textwidth}{X} 因為你要向左向右開展，你的後裔必得列國為業，又使荒廢的城鎮有人居住。 \end{tabularx} \\ \\ \relax
54:4 & \begin{tabularx}{0.7\textwidth}{X} 不要懼怕，因你必不致蒙羞；不要抱愧，因你必不致受辱。你必忘記年輕時的羞愧，不再記得守寡的恥辱。 \end{tabularx} \\ \\ \relax
54:5 & \begin{tabularx}{0.7\textwidth}{X} 因為造你的是你的丈夫，萬軍之耶和華是他的名；救贖你的是以色列的聖者，他必稱為全地之神。 \end{tabularx} \\ \\ \relax
54:6 & \begin{tabularx}{0.7\textwidth}{X} 耶和華召你，如同召回心中憂傷遭遺棄的婦人，就是年輕時所娶被遺棄的妻子；這是你的神說的。 \end{tabularx} \\ \\ \relax
54:7 & \begin{tabularx}{0.7\textwidth}{X} 我離棄你不過片時，卻要大施憐憫將你尋回。 \end{tabularx} \\ \\ \relax
54:8 & \begin{tabularx}{0.7\textwidth}{X} 我因漲溢的怒氣，一時向你轉臉，但我要以永遠的慈愛憐憫你；這是耶和華－你的救贖主說的。 \end{tabularx} \\ \\ \relax
54:9 & \begin{tabularx}{0.7\textwidth}{X} 這事於我有如挪亞的洪水；我怎樣起誓不再使挪亞的洪水淹沒全地，也照樣起誓不再向你發怒，且不斥責你。 \end{tabularx} \\ \\ \relax
54:10 & \begin{tabularx}{0.7\textwidth}{X} 大山可以挪開，小山可以遷移，但我的慈愛必不離開你，我平安的約也不遷移；這是憐憫你的耶和華說的。 \end{tabularx} \\ \\
[1ex]
\hline
\hline
\end{longtable}
$^{1}$請使每一主裡面平安.
平安.
今天很想用這個題目.
跟大家彼此的勉勵.
這一章不只是十節的.
是有十七節.
但是後面不夠時間講了.
所以只能夠用三個重點.
來將這個題目.
來跟大家一起.
透過這十節聖經來認識.
大家可以看到.
在這一章裡面.
是有相當大的勉勵的.
從荒涼,羞愧,恥辱.
到被神接納.
第十節更加寶貴.
大山可以挪開.
有小朋友的聖詩.
有這樣的詩歌.
小山也可以遷移.
但是神永遠不改變的愛.
對這個民族的重新收納.
是會恢復的.
所以今天很想透過五十四章.
一到第十節.
跟大家彼此的勉勵.
教會是需要講安慰的訊息.
尤其是一個在現在社會.
人心都是荒亂的社會.
我們需要的是安慰.
神安慰的訊息.
當然.
那個安慰.
為什麼會有這樣呢?.
就要明白那些背景.
我們就有一個簡單的祈禱.
我們仰望主.
透過這十節的聖經.
尤其是經過了五十三章.

$^{41}$受苦僕人的被掛.
吸引眾人歸向你的前提.
來看到這首僕人之歌.
我們是如何看到.
當時以色列這個民族.
是已經在神過去的刑罰.
將他們被擄之後.
來恢復的這一份安慰.
而且是要透過先知.
來宣告安慰的訊息.
所以我們很盼望.
五十四章成為我們現代版的.
一個安慰.
讓我們在我們的生活層面.
我們能夠感受到.
原來神從來沒有改變對我們的.
愛和接納.
求主帶領我們.
盼望透過舊約的經文.
我們不是一些避氣的婦人.
我們更加不是一些守寡.
甚至不能夠生育的婦人.
但我們知道經文的擴展性.
延伸性.
是隨著時代.
可以有不同的.
所以我們很盼望.
我們能夠明白整個的處境.
以及對經文.
當我們去實際應用的時候.
滿有聖靈的教導.
在我們大家的心靈裡面.
我們這樣仰望禱告.
奉耶穌基督的名.
阿們.
我們一同看有三個.
在十節裡面.
有三類的人.
是神透過先知以賽亞.
來勸勉他們.

$^{81}$就是心靈在荒涼底下的人.
這是一到三節.
第四到第八節.
就是羞愧.
面對很大的羞愧的人.
更加重要就是.
有些已經好像被神離棄的人.
他們的光景.
心靈是怎樣的呢.
我們一同進入這裡.
去看一看.
一到第三節.
我們先看看經文.
在荒涼裡面.
強調.
一到三節強調.
不斷要擴張你的靈脈.
擴張你的地圖.
這裡告訴我們.
怎樣擴張呢.
我們看看這裡.
剛才讀了.
不懷孕不生育的要歡呼.
未曾經過產難的.
都要歡呼.
要揚聲的呼喊.
因為被遺棄的婦人.
比有丈夫的人的兒女更加多.
相信不是真的照字面來講述.
因為是在講整個民族.
所以.
神看自己是以色列人的丈夫.
被遺棄就是他們被神懲罰.
甚至要被擄到巴比倫.
但是.
祝福和安慰就是.
下面這裡.
要擴張你的將惡之地.
擴張.
本來是一種這樣的荒涼.

$^{121}$古代.
婦人.
結了婚.
生不了孩子是大事.
今天我們要鼓勵人們生.
但是舊約古代的時候.
生不了孩子.
是某程度上.
與神的咒詛拉上關係的.
記得是古代.
你可以想像得到.
聖經往往是用他們能夠明白的處境.
他們的時空.
來講述恩典.
在這樣的荒涼下.
神現在的工作.
那個應許還沒有兌現.
但是他們現在已經感受到.
他們自己活在荒涼的咒詛下的時候.
應該怎樣呢?.
但是神的話語.
就比兌現的處境.
要他們擴張.
準備擴張.
準備你們要向左向右.
開展你們的將惡.
意味著這個民族.
經過神的懲罰之後.
是會帶來一個很大的復興的.
如果回到我們個人.
又會怎樣呢?.
在荒涼裡面.
神呼召.
不生養.
對應我們現代.
我們在什麼處境呢?.
你說我都沒有結婚.
我是男士.
生什麼呢?.
暗瘡有時候會.

$^{161}$那怎樣呢?.
怎樣對應那些經民呢?.
相信那個擴張又是怎樣呢?.
相信在這樣的所謂的荒涼.
現代版.
我想我們在我們的生活.
我們的壓力.
我們面對工作的那種掙扎.
我家裡就有一個現省的.
因為他上班都要用一個半.
一個小時零九個字.
才從家裡出門口.
去到九龍塘上班.
其實已經很舒服了.
比起我們以前.
但是呢.
他都覺得很辛苦.
每天只是上班的車程.
要一個小時九個字.
接著很密集的工作.
就已經很辛苦了.
下班呢.
也是一樣.
回到元朗之後.
再要坐巴士.
那個K66呢.
也是爆棚的.
只是從元朗.
由朗平村.
回到大堂.
都要差不多九個字.
因為站站停.
教育路呢.
只是教育路那條馬路.
都停你半個小時的.
你可以想象.
我想我們每個人.
每天形形異異.
有時候可能你真的感覺自己.
好像一個荒島.

$^{201}$好像一個很荒涼的人.
我嘗試用一個例子.
今天是大雄大子的.
奧斯卡的影帝.
但是原來他有一段相當長的時間.
左至古來.
他的事業停滯得很厲害.
他不斷去試鏡.
去不同的導演試鏡.
他渾身解數.
展現他的演出才華.
但是每一次.
沒有回音.
是說很多次.
沒有回音.
後來呢.
他就有一個反省.
他說他不再用自己的角度.
去看那個角色.
而應該就是.
到底那個導演.
要求這個角色是怎樣呢.
他這次嘗試用導演的角度.
去看那個角色.
不是盡量表現自己怎樣有才華.
怎樣出眾.
去演出那個角色.
而是那個劇本.
到底那個導演想要求.
什麼樣的演員呢.
他從那個角度去看的時候.
整個局面.
他的事業.
就從那裡開始.
好像一個零的突破.
開始他越演越多角色.
自然就更多好的劇本.
送到他的面前.
就造就了我們所認識到的.
左至古來.

$^{241}$很多電影都是很拒絕人口的.
但是.
原來在他事業很低潮的時候.
他往往從自身去看.
以至他雖然很有演技.
但不是對導演認為的那個角色的那個度身訂做.
跟我們有沒有什麼特別的關係呢.
我們會埋葬在我們自己的方良.
在我們自己的缺乏.
在我們自己的那種空虛裡面.
我們不懂得去擴張呢.
神的應許.
是不是一直都好像不對焉.
往往是不是我們不斷在我們裡面.
去找我們自己的才華.
我們自己的那個角色.
而忘記了呢.
其實神期望.
他的應許就是.
刑罰已經過去了.
你現在要擴張你的地圖了.
以色列人.
我們個人的方良.
有什麼可以令我們擴張呢.
有什麼大家仍然停留在你自己那種兜兜轉轉的裡面.
而你忘記了呢.
神的同在恩典.
因為我們還沒講下去.
其實每一節都互有相關的.
因為講到第九節到第十節的時候.
神最後一次肯定就是.
他的應許.
好像羅亞的彩虹之約.
永遠不會改變.
正如羅亞的洪水.
那你可以想象得到.
如果弟兄姊妹.
你在你現在自己的事業.
在你自己現在的感情世界.
在你自己個人的一些自身的問題.

$^{281}$還是在一種這樣的方良呢.
每個人都不同的.
每個人都有自己的方良.
但是如果你繼續在這種方良感裡面.
看不到有神的應許.
我們是神兒女.
我們是有十字架.
成為我們很大的信主.
以及倚靠主的力量.
是五十三章之後的.
「理者不懷孕不生養」的.
要生養眾多.
是來自五十三章.
這一首獨人受苦的詩歌.
這樣開始的神的安慰.
以及神的應許.
第二方面.
剛才提到的就是收攜.
要除去.
恢復我們尊貴的身份.
就是神現在要他們除去他們的收攜.
經文是這樣說的.
「不要懼怕,因你必不至蒙羞,不要抱攜」.
這些相關的詞語.
蒙羞,抱攜,受辱.
因你必不至受辱.
你必忘記年輕時候的收攜.
不再記得守寡的恥辱.
這些話出現了五次了.
「因為做你的事你的丈夫,萬君之耶和華是他的命」.
這樣的比喻.
猶太人就是辛苦.
但是被離棄的.
因為在荷西亞書的概念.
就是他們對神不忠不貞.
就是拜偶像,離棄神.
所以神就要離棄這個婦人.
就是休了她.
但是現在在五十四章的處境.
神就說我現在要重新收納你.

$^{321}$所以你年輕時候的恥辱.
就是被神的懲罰.
現在可以恢復這個身份了.
「但我要以永遠的慈愛憐憫你」.
「這是耶和華你的救贖主說的」.
耶和華已經足夠那個身份了.
再加上耶和華救贖主.
就更加肯定.
神現在是用救恩.
將他們挽回這份祝福.
羞恥.
年輕時的羞恥.
跟大家看一個人物.
都挺好笑的.
這位年輕人.
經過無數的挫敗.
這個年輕人.
他畢業的時候.
原來他高考考了三次.
兩次都是很低分的.
第三次勉強進了杭州的師範大學.
出來就想著教書了.
輾轉出來.
還沒教書之前.
他就應徵一些工作.
他應徵KFC.
有24個人申請那個工作.
有23個人被取錄.
他就是唯一一個不被取錄.
接著他考警察.
五個警察.
請了四個.
但是呢.
講明了為什麼不請他.
因為你的樣子太醜陋了.
太突出了.
如果你是警察.
所有人都知道.
就不會找你了.
又是落選了.

$^{361}$輾轉呢.
他就是這樣經過了不斷的挫敗.
你可以想像得到.
今天的馬雲.
是一個真是.
大家都明白認識的一個.
很出色的人.
但是原來他經過無數的挫敗.
不過他提到他自己回想這些挫敗.
他說挫敗只不過是起點.
不是結果.
挫敗令他更加知道.
我下一次一定會成功.
世間其實這樣的人.
是少之又少的.
經過挫敗.
其實可以是內捲.
可以是很奮勢疾逐.
可以是很放棄自己的.
更多的人可能根本不會能夠重新起來.
不過.
好像這樣開玩笑.
這個人這樣能夠重新起來.
但其實.
在神的真正的應許裡面.
不斷叫我們.
你感受到的恥辱.
那個羞愧.
神說現在我要用新的身份.
來去裝飾你.
這個十一節後.
就不行了.
要用彩色的石頭.
來去裝飾.
這樣的一個安慰和應許.
這裡告訴我們.
神兒女.
如果現在是在一個刑罰裡面.
刑罰背後都是愛和管教.
都是一份要使我們靈命成長.

$^{401}$更加認識神恩典的一個人.
所以神就在這段經文.
第四到第八節.
就是安慰他們.
在羞愧裡面.
要除去這個羞愧.
恢復你自己在神面前的身份.
除去羞愧.
不是一種阿Q精神.
而是真的相信.
神有祂的安排.
在我們人生不同的際遇.
使我們能夠更加感受到.
我們是神兒女身份的尊貴.
最後.
動盪裡面的一個不動搖的藥.
剛才第九到十節大家已經讀了.
這事好像我們在羅亞的洪水.
神自己親自說.
所以真的有洪水滅世.
因為二次教書都肯定這件事.
我怎樣起誓.
不再使羅亞的洪水.
淹沒全地.
也照樣起誓.
不再向你發怒.
而且不會再斥責你.
大山.
這就是第十節.
可以挪開.
小山可以遷移.
但是我的慈愛.
必定不會離開你.
我平安的藥.
是平安的藥.
都不會遷移.
這是連文裡的耶和華說的.
各位.
這裡是一個很重要的應許.
神給我們立的這個永約.

$^{441}$是永遠不改變.
就是祂的慈愛.
對於我們每一個人.
羅亞的洪水.
只剩下八個人在世界.
全部其他都死了.
完全滅絕.
但是神的應許.
當你看到彩虹的時候.
你要紀念神的憐憫.
神不會改變的平安的永約.
跟大家玩一個遊戲.
很特別的.
我發現有兩個題目.
大家一定要回答.
現在藍的先回答.
神很愛她.
很美.
那個形容真的很美麗.
彎彎的鵝眉.
是怎樣的呢?.
臉色很白.
不是睡眠不足的那種.
是真的美白.
美麗動人.
不過.
突然有一天.
遇上了車禍.
當他康復之後.
那個臉色.
有大大的疤痕.
很醜的樣子.
大家弟兄男士.
你表態一下.
他會不會.
說的是他,不是說你.
他會不會一如既往地.
愛這個女孩子呢?.
一定會的.
都舉手好嗎?.

$^{481}$一定會.
一定不會.
沒有人舉手.
可能會的.
可能會好像多一點點.
等下再開估.
到女士了.
他很愛她.
是商界的精英.
儒雅沉穩.
敢打敢拼.
不過突然之間.
完全的破產.
你覺得呢?.
姐妹們.
他會不會像以前一樣.
愛她呢?.
已經破產了.
答A的請舉手.
因為和你無關.
一定不會的.
一定不會也有.
一定不會也沒有.
可能會呢.
個個都很肯定.
但我們還沒問這兩個心理測驗的問題.
是一個教授在課堂裡面.
問他的學生.
幫一些機構做一個問卷.
但沒有說到受眾是誰.
然後教授讓那些人填了之後.
那些年輕人就填.
女的通常不會來接納男朋友破產.
反而男的對有疤痕的女朋友會好一點.
但教授後面開估.
因為很多人都說可能會.
因為一定不會.
相對是少數的.
可能會大多數一點.
但教授再重新問一次.

$^{521}$他要開估答案.
其實大家不知道他代表誰.
如果男性他代表他爸爸.
大家不用想.
一定是所有人都舉手.
一定會.
如果下一個是他的媽媽.
也不需要想了.
所有人都舉手.
一定會.
不會因為她的臉發了不愛她.
不會因為她破產不愛她.
沒有人會.
我想古今中外都是.
一定會.
無論怎樣都會愛.
你可以想像得到.
這豈不是神的慈愛嗎?.
神的忠貞信實嗎?.
有沒有看過這套電影?.
李卓基演的《鐘犬白宮》.
其實叫八千.
Hachi.
日文的秋田犬.
秋田犬本身已經很忠誠.
但這隻狗好像有靈性.
每天因為主人帶牠去車站.
牠在車站送別主人.
等主人下班.
是一位教授.
日本東京大學教授.
真人真事.
所以現在你去到車站.
會有這隻狗去銅像.
後來這位教授突然教書.
心臟病死了.
狗不知道.
繼續每天在車站等主人.
等不到自己回家.
如是者每天連續十年.

$^{561}$後來那些人看到這隻狗這麼忠誠.
就幫牠立了一個像.
親愛的弟兄姊妹.
我們說父母的愛.
不需要問牠變成怎樣.
牠怎樣我們都會愛我們的兒女.
一隻這麼特別的狗.
尚且對主人有忠貞.
我們今天面對的是這位愛我們的造物主.
他說我不會再改變對以色列民族的愛.
你好像已經失去了那種信心.
但是神的慈愛是不會離開.
大山可以挪開.
小山可以遷移.
但是神的慈愛是永遠不會離開.
祂平安的約都不會改變.
弟兄姊妹.
你不是落在荒涼.
羞愧.
或者感受不到神的慈愛.
回頭看回53章.
這位耶和華受苦的僕人.
他被棄絕.
常經憂患的這位耶穌基督.
被人掩面不看一樣的.
這位常經憂患的主.
就是這樣為我們默默獨自走上十字架的孤苦的路.
就是成就神救贖的大恩.
所以在54章.
神的安慰就因為在主的救恩下.
就來向他本來已經責備到體無完膚的民族.
來重新收納.
重新安慰他們.
所以如果你已經落在這樣的荒涼.
教會就是一個醫治心靈的地方.
教會就是一個能夠彼此安慰的地方.
教會就是一個.
不只是說一些心靈雞湯.
是神要肯定你.
盼望今天.

$^{601}$神不只是說加油這麼簡單.
而是很想透過聖經.
透過聚會告訴你.
神仍然愛你.
祂不是要你撐下去.
而是要你擴張你的地圖.
恢復你能夠成為可以被神用的人.
在不同的領域去服侍主.
你不是避氣.
而是在神的救恩裡面.
你是祂所喜悅的兒女.
我們一起同心禱告.
再一次我們將剛才提到的三點.
荒涼者.
羞愧者.
感覺避氣.
得不到愛的群體.
天父啊.
我們雖然不是生在古代.
我們更加不是太明白那些語境.
那些好像避氣的婦人.
不能夠生育的那種在古代的羞恥.
還有因為在神的刑罰.
而被擄到他鄉的以色列人.
我們今天.
我們都是在上一代來到香港.
我們在今天.
現在面對相當大的疫情後的經濟壓力.
有些地方.
有些行業.
甚至失業率是十幾個百分比.
我們仰望主.
讓我們無論有工作.
或者我們都面對我們自己的心靈.
或者生活的荒涼的時候.
我們去倚靠你.
仰望你.
說的是很有限制.
盼望神自己透過聖靈.
在每一位弟兄姊妹的內心.

$^{641}$來工作.
有些人可能每天要面對無盡.
無休止的那種慢性的食藥.
好像越吃越多的情況.
我們都需要主給我們有信心.
有一個信靠的心靈.
或者有些老人家.
開始連行動都不是太方便.
甚至出門口都有困難.
但是我們都盼望主給我們大家.
都能夠透過主聖靈的安慰.
得著你話語的提醒.
我們再一次獻上感恩.
奉耶穌基督的名.
阿們.
祝福大家.
\newpage



\section{約翰福音 5:24}
\label{sec:Bm5kMjeNAdQ}
\textbf{2025 10 12《惶恐中的平安》:約翰福音5\_24 (和修版) 講員 : 樓恩德牧師}
\newline
\newline
連結: \href{https://youtube.com/watch?v=Bm5kMjeNAdQ}{\texttt{https://youtube.com/watch?v=Bm5kMjeNAdQ}} ~~~~ 語音日期: 2025-10-15
\newline
\newline
\hyperref[sec:j0UisY02_wA]{\small{< < < PREV SERMON < < <}}
~
\hyperref[sec:index_chronic]{\small{[返順時目]}}
~
\hyperref[sec:index_scriptual]{\small{[返順卷目]}}
~
\hyperref[sec:RkNxGeCkTHA]{\small{> > > NEXT SERMON > > >}}
\newline
\newline
約翰福音 5:24
\newline
\begin{longtable}{cl}
\hline
\hline
章節 & 經文 (和合本修訂版)\\
\hline
5:24 & \begin{tabularx}{0.7\textwidth}{X} 「我實實在在地告訴你們，那聽我話又信差我來那位的，就有永生，不至於被定罪，而是已經出死入生了。 \end{tabularx} \\ \\
[1ex]
\hline
\hline
\end{longtable}
$^{1}$各位朋友 各位弟兄姊妹.
剛才大家聽了惠妮姐妹過去的經歷.
我相信坐在這裡聽的時候.
你很想聽前半部分.
20歲不到就找到一桶金.
那個時候是一百萬.
不是現在.
現在一百萬 三十歲.
但是後期的你想不想過.
又不想過 那就很矛盾.
到底想前還是想後.
我今天想跟大家透過王孔這兩個字來看你的心態.
什麼叫王.
經常不上不下.
心怯怯.
什麼叫安.
或者孔.
驚 很驚青.
無論你現在是什麼處境 什麼環境.
你現在坐在這裡聽.
都算是很平安的了.
當我這幾天看到.
電視畫面上我偶爾看一下.
那些難民要回家沙地帶的時候.
請問你回不回去.
我想我都不會回去了.
回去不知道怎麼收拾產局.
爛了地方.
不像實質區那樣的環境.
有瓦遮不住頭.
恐懼.
如果說得貼切一點.
各位朋友.
其實你現在的人生.
是很不放心.
什麼叫放心.
睡得著就放心.
睡不著就不放心.
不要說打仗那些.
社會環境.

$^{41}$留心自己健康都不放心.
是不是.
那怎麼辦.
平安.
這兩個字更加矛盾.
希不希望和約簽成功.
不打仗了.
行不行.
和約簽了.
大家可以.
不打仗.
你信不信.
不信那簽了幹什麼.
簽了又不好.
簽了又不信.
用一個很直率的話.
和約就會簽好.
和好就永遠不會發生.
這個不單單戰場.
你的人生都是這樣.
請問在座各位.
不要說你和同事.
舊同事.
請問你和左右鄰舍.
甚至家裡的親人.
親戚.
有沒有一些.
未和解的問題.
請記住.
所以這個世界不單單戰爭.
令到惶恐不安.
你個人.
人際關係.
你個人的健康.
都令到你不放心.
一個不放心的人.
怎麼生活下去.
唯有.
每天過完就算了.
又要過又不想過.

$^{81}$矛盾.
約翰福音第五章二十四節.
聖經這樣記載.
他怎麼說呢.
我讀給大家聽.
如果你聽著都可以.
我實實在在告訴你們.
那聽我話.
猶信差我來的那位.
就是那位神.
有永生.
不至於定罪.
是已經出死入生了.
聖經就說.
你現在活著的.
是一個面對死亡的人生.
每個人都是.
沒有人能逃得了.
但這個面對死亡的人生.
是你最大的陰影.
也是你最大的掛心.
令到你不安.
不放心的地方.
我們俗語說.
死都要死.
還怕什麼.
你這樣說.
不要說死都要死怕什麼.
你今天有些不舒服的檢查.
醫生說看不出什麼病.
不如去化驗所驗一驗.
抽血驗東西.
驗完之後.
下個禮拜一回去.
今天禮拜幾.
今天禮拜一.
明天回去看報告.
今晚你怎麼睡覺.
很多人都睡不著.
為什麼睡不著.

$^{121}$為什麼不可以想.
醫生說沒事.
沒事就是下次又來.
又沒事.
下次又來又沒事.
幾十次沒事之後.
會不會有一次有事.
會不會.
會啊.
那你怕什麼.
這就是人生.
出生入死.
每個人都出世.
但每個人都要去那條路.
我和大家看.
記住聖經.
因為時間不夠.
我透過聖經和大家一起看下去.
我按就行了.
耶穌基督在二千幾年前.
耶穌曾經在地上.
活動過.
他怎麼說.
他說了很多話.
要你們在我裡面有平安.
世上你們有苦難.
但你們要有勇氣.
我已經勝過這個世界.
耶穌基督二千幾年前.
在地上說了這句話.
不是當時的人有苦難.
我告訴你和我聽.
今時今日這個世界仍然有苦難.
因為我很好.
你擔不擔心.
你放心不放心.
你放心不到.
聖經說.
我們一生因怕死.
而成為奴隸.

$^{161}$其實我們是怕死.
因為我不怕.
你以為自殺不怕死嗎.
他不是不怕死.
他不怕死.
他只是不知死.
剛才姊妹做見證.
我以前在學校工作.
有個學生中一中二中三.
有一天.
他在書桌上.
自己拿一張紙畫畫.
有個老師走過.
一看到之後.
他問我先生.
不理他.
先生說.
你畫什麼.
一看下去.
畫一棟高樓大廈.
然後有個窗.
有一層樓.
書桌上畫的東西.
有個人從空中飛人.
飛下來.
那位老師看到之後.
你做什麼.
不要說了.
那位老師知道出問題了.
馬上第一時間.
和他的輔導主任說.
你這幅畫.
輔導主任一看.
不行.
馬上跟進.
原來他想放學後.
走上天台.
在上面跳下去.
經過自己書房.
然後落地.

$^{201}$他不知道死.
你跌下來死了.
沒問題.
不死又怎樣.
當然後來.
這個小孩撿回來了.
各位.
朋友.
各位弟兄姊妹.
不是像我們這樣年紀想死.
很多年輕人.
都想死.
什麼問題.
不是像當時.
姐妹說.
要吃皇家飯.
才辛苦.
人原來.
在這個苦難的世界.
是不安的.
是不放心的.
是沒有平安的.
說的是一些東西.
和約可以簽好.
獎可以拿了.
如果可以拿的話.
但有一件事實.
和好不用想.
但是.
在我們的信仰裡.
有解答方案.
人將心中.
的仇恨.
將心中的.
結怨.
最大的問題.
罪要除去.
我和大家看聖經.
耶穌.
當時對這班門徒們說.

$^{241}$你們心裡不要憂愁.
你們信神吧.
信我吧.
因為很多人以為一條路是正路.
始終成為死亡之路.
不知道大家.
這個年紀.
有沒有想過這個問題.
到我這個年紀.
我已經離開了學校.
退休.
再離開教會.
工作崗位.
所以我兩次.
不是退休.
是下崗.
不用上班.
所以不是退休.
是下崗.
因為我們的工作沒有休好退.
我到今天.
還看我以前的學生.
他們都醉到下崗.
有很多問題出來了.
他發覺真的解答不了.
唯一.
是正路.
每個人都以為自己走正路.
沒有人想走歪路.
我想姊妹剛才見證.
他要在四年坐牢的過程裡.
他不會想做那些事.
不知道為什麼會這樣.
我們不會去理解.
總之但是有條路.
每個人都以為走正路.
結果成為死亡之路.
對了.
按不上去.
退一張.

$^{281}$聖經寫記載.
有一樣東西.
大家有沒有留意到.
施育懷胎就生出罪來.
罪記長成就生出死來.
人會死.
因為人有罪.
罪是怎樣來的.
罪不是人家加在你身上.
不是人引誘你去做壞事.
是你裡面生出來的.
施育懷胎就生出罪來.
罪記長成就生出死來.
這個聖經很棒.
你只要肯想.
肯想要.
就會有.
舉例來說.
我想加一件衣服.
口渴想喝一口水.
都是一種想法.
但問題是.
正常的想法.
滿足了嗎.
永遠不會滿足.
如果你問.
惠尼茲.
孫教士.
一桶金夠不夠.
你都會搖頭.
兩桶金夠不夠.
被你做到馬斯克.
夠不夠.
全世界首富.
沒人夠.
施育懷胎就生出罪來.
你會不會生氣.
會不會嫉妒人.
很多事.
但問題在這裡.

$^{321}$當罪生出來之後.
為甚麼.
有些事不能榮耀神.
神將你放在這個世界上.
是要你反映出.
祂的榮耀.
在這裡我講給弟兄姊妹聽.
基督徒.
發不到光.
不是說話.
基督徒只可以反光.
都不明白.
剛剛過了中秋.
明不明.
會不會發光.
不會的.
月亮只可以做甚麼.
只要沒有月蝕.
沒有被東西遮住.
你不斷有光.
一有東西遮住你便發不到光.
反不到光.
所以不要說像蠟燭一樣.
燃燒.
犧牲自己.
做不到的.
基督徒只可以做一件事.
將上帝的榮耀反射出去.
這不是福音信息.
這是人生的信息.
世人都犯了罪.
規節了神的榮耀.
好.
我跟大家講講甚麼是罪.
這本聖經很有意思.
很多時候.
我每次報道福音.
我都一定講.
很多年前我做行政工作.
在學校做校長的時候.

$^{361}$遇到一宗案件.
有一個學生.
因為在日間集會的時候.
講台講話.
學生在後面.
經常滋擾其他同學.
所以訓導主任不見了之後.
叫了名字.
叫他走出來.
不要坐在那裡.
站在旁邊.
站五分鐘.
他覺得.
落了我的臉.
趁沒有人知道.
在停車場.
把教師的車.
車輪中間.
刺了幾下.
刺了之後.
他就走了.
因為沒有人知道.
是表示自己的不滿.
發洩而已.
沒甚麼的.
那一刻老師開車回家的時候.
住在錦繡花園.
那個老師.
我告訴你 你們都認識的.
所以人死我不會說的.
不是到胡震.
因為後來他安息了一星期.
我跟他說的.
他星期一回到學校.
他就說那天.
晚會.
星期五晚上回家.
那麼吵 搞了我換的車.
換到十一點多才回家.
他旁邊的同事.

$^{401}$坐在旁邊.
偉哥的同事.
你的車.
發生甚麼事了.
他說漏氣了.
他說我星期六開車出街的時候.
車也漏氣了.
於是兩個多口說一句.
你的甚麼輪胎.
我的左前輪胎 你的右前輪胎.
兩個車是隔離輪車.
所以他們兩個想.
不對辦.
沒理由這麼巧.
舊經驗老師.
潘Sir很有經驗.
馬上拿名單一找.
找了幾個同學來問話.
其中一個就是這個小孩.
沒有痕跡 甚麼都沒有.
後來這些同事說不是了.
他說如果是的話.
一定要矯正他的想法.
結果.
我走過走廊.
看到他在房間裡面審問學生.
到最後還有一個學生在.
那時已經過了一至兩堂時間.
我說甚麼事.
他說我們快要完成了.
他說沒甚麼事.
我們問問話.
當後期的時候.
他們快要結束的時候.
我說如果你們都找不到甚麼特別的.
他沒有說以前的事給我聽.
我說沒甚麼特別的事.
讓他上課吧.
免得.
學生將來怪責我們.

$^{441}$冤枉他.
很奇妙的事.
我知道案件很麻煩.
沒人知道 沒人證明.
我只有禱告床等我一會.
如果這個小孩.
是真的做錯事.
給他機會悔改.
但如果沒有的話.
希望你給老師們放過他.
老師不放過是很麻煩的.
你的案件是老師不放過你.
我一聽就明白.
老師是做甚麼的.
半小時後.
訓導主任.
來到我的房間.
把張紙給我.
他說校長.
這就是同學寫的悔改書.
我一看就.
嘩.
這樣的事件.
都會這樣做.
我第一句問潘Sir.
我說潘Sir.
這件是刑事案.
你想怎樣處理.
我說我可以報警.
他說報警不需要了.
他說他是無知無聊.
我說你.
你要我怎樣處理.
我說你是訓導主任.
他說你處理好.
我處理好就不行了.
我怎可以代替你.
因為車輪要修好.
也要付錢的.
我就問他怎樣處理.

$^{481}$他說這樣吧.
我們記下兩個小學的案件.
當時我們記下兩個小學的案件.
學期尾沒有事.
那年就取消了.
我們很仁慈的.
不會像你這樣.
記下一個小學的案件.
不可能的.
我們很仁慈的.
真的沒有.
所以我聽二十多個一定趕.
不趕的.
結果呢.
他說我們要上課了.
他上課了.
他叫學生進來.
學生進來之後.
看著我.
很低.
臉白白灰灰的.
挺高大的.
我就問他.
我說你自己做過什麼.
他不出聲.
我問他.
這張東西是不是你寫的.
我說你知不知道自己做過什麼.
他沒有出聲.
沒有行動來回應我.
我再問他.
是不是你寫的.
我說你知不知道自己做過什麼.
一次兩次.
到第三次我問得特別一點.
是不是你寫的.
是不是逼供.
你知不知道自己做過什麼.
三次之後我想死了.
這次交給我要我處理.

$^{521}$我問他.
是他寫的.
知不知道自己做過什麼.
我沒辦法.
我心中只有默默禱告.
上帝幫助我.
怎樣可以處理這件事.
我就想起這句話.
當時我不敢肯定.
我立刻抽筋拿出聖經.
當我拿出聖經的時候.
那個小朋友.
他瞅瞅我.
奇怪.
問一問自己看聖經.
我為了肯定.
因為我怕自己看錯.
我看清楚.
是真的這樣說.
於是乎.
我看完之後.
我就跟他說.
同學.
你知不知道自己做過什麼.
我說我告訴你.
我知道你為什麼會告訴我.
你知道自己做過什麼.
他聽完我講之後.
他莫名其妙知道不知道不知道.
到底知不知道.
我就說.
你讀吧.
一起讀.
我所不願意的事.
我不去做.
我所不願意的惡.
我反而去做.
如我去做.
我不願意做的.
就實是我做的.

$^{561}$(以示處在對面的罪行).
是不是答案.
我不是被校台揭開的.
是他寫的.
寫完原因.
為什麼.
他揭開車輪.
揭開旁邊的車輪.
(笑聲).
所以兩輛車.
都同時出問題.
他揭開車輪.
車輪旁邊的車輪比較遠.
不是我願意做的.
是我裡面的.
罪行.
各位朋友.
各位頂尖妹.
你想不想罵人.
想不想激氣.
你不想的.
不是你做的.
是肢肉做的.
人家厲害.
你沒那麼厲害.
沒問題的.
為什麼到你有一天厲害一點.
那個厲害的.
就會插你背.
我講反話.
你厲害.
人家沒那麼厲害.
到有一天那個人厲害了.
你會怎樣.
裡面就有東西出來了.
那些是什麼.
不要說同事.
不要說同學.
這些我見得多.
連家裡兄弟姊妹.

$^{601}$都是這樣.
小的圈子叫家庭.
大的圈子叫社會.
再大一點叫國際.
國際是不是這樣.
都是這樣.
你不厲害就行.
我厲害就行.
看到嗎.
人裡面的罪.
做的.
所以聖經說.
世人到.
罪的功架乃死.
唯有神的恩賜.
在基督耶穌裡乃永生.
人因為有肢肉.
人會有罪.
有罪就會死.
是結局.
所以按著定命人.
人都有一死.
死後且有審判.
有一條路人以為正.
結果成為死亡之路.
所以聖經說.
世人都犯了罪.
虧缺了神的榮耀.
各位朋友.
各位弟妹.
我不是說你要坐牢.
達官司脫罪.
脫罪了.
無論妒忌驕傲自大.
狂妄.
什麼都好.
仇恨在這裡.
脫不了.
只有在耶穌裡面.
你因為有罪.

$^{641}$所以一直不安.
不但心靈不安.
連身體都不安.
所以你的血液有問題.
你的睡眠有問題.
甚至乎.
當醫生檢查不到你的病因的時候.
只能說這可能是遺傳有關.
真的.
只有神才能監察.
你現在.
以前看到外貌.
中醫很厲害.
有把脈.
加上x-ray.
越照越入.
但越來越入.
神可以看到.
人心靈中間有的就是罪.
罪要得赦免.
只有靠耶穌基督.
因為聖經這樣說.
唯有基督在我們還作罪.
一起讀.
唯有基督在我們還作罪人的時候.
為我們死.
上帝的愛.
是我們還作罪人的時候.
為我們死.
神的愛就是向我們顯明.
我不想你再心怪怪.
想著剛才審訊的結果.
怎樣告訴你.
結果我把聖經給小孩.
要讀的時候.
他眼泛淚光.
面紅紅.
低頭.
我說你不知道自己有罪.
叫你這樣做嗎.

$^{681}$老師的車.
他們不需要你賠償.
本來按道理.
我有責任要報警.
但現在給你機會.
你寫了你認了嗎.
我相信很多時候.
在教育過程中間.
很多寬恕.
很多愛.
是透過.
一些受害者.
他願意給出來.
這不是一單.
當然不是這麼嚴重的問題.
我有一次見到一個老師.
搭著肩膀.
和一個學生.
在走廊上走.
我知道那個學生.
被老師說的時候.
我沒有留意跟著他聽.
我認同不見你.
只是老師說了一句話.
以後不要說.
他說了這句話.
我小時候哭了.
今天你和我.
不安.
是因為你還沒哭.
你還沒悔改.
唯有基督在我們還作罪的時候.
為我們死.
為什麼? 耶穌為你.
為我所罪的懲罰.
為我們死在十字架上.
他不是為好人死.
不是為義人死.
是為罪人而死.
所以聖經說.

$^{721}$神的愛向我們顯明.
是透過十字架上顯明.
這句話是可信的.
是值得完全接受的.
因為耶穌基督.
來到世上是要拯救罪人.
所有宗教.
都是勸人為善.
這是另一個範圍.
你被拯救之後.
你得到饒恕之後.
你就會饒恕得罪你的人.
你被愛之後.
你會更加愛人.
我和曾志偉.
今天第一次見面.
但我看完他見面之後.
他的生命的改變.
因為耶穌基督的愛.
臨到他.
他能夠在他的.
監獄裡面.
去幫助人.
在那裡得到救恩之後.
他願意出來幫助人.
這不是他天生的.
不是他小時候.
記住那些東西.
給這個發給那個發.
不是那些東西.
是媽媽的行動.
對他的關懷.
愛心行動.
教會裡面的牧師.
對他的接納.
這種領受.
是愛人的.
再講一次.
你不是一個在被愛的環境中.
成長你不會愛人的.

$^{761}$但世間的愛有條件的.
會過去的.
但神的愛是永遠的.
沒有條件的.
這就是神忠誠之愛.
所以當耶穌在世.
為你為我死的時候.
各位朋友.
他不單止為我們的罪死了.
他因著聖靈的工作.
神讓耶穌基督從死裡復活.
使我們恩順清怡.
不單止罪得赦免.
我們任何的蒙恩德的人.
是沒有虧沒欠的.
不會欠耶穌的人情.
耶穌的愛是公開的.
是自然的.
是無捨的.
是無條件的.
在這個愛裡面你被愛之後.
你就會愛其他的人.
你會信其他的人.
否則.
請問這個世界上有什麼值得你信.
下一節聖經這樣記載.
你若考慮認耶穌為主.
心裡信神.
叫他從死人中復活.
就必得救.
很多人以為.
我知道了我信了.
信不是說我接受耶穌.
做救主那麼簡單.
是心裡相信.
這位耶穌.
現今仍然活著.
你看.
他現在的活著.
是超自然.

$^{801}$等於你和我將來會復活.
得永生的活.
我們不是將來.
領魂上天堂那麼簡單.
不是的.
我們會復活的.
誰知道耶穌已經在歷史上降生.
事實不能改.
現在是十月.
再過兩個月一定放假.
放什麼假.
關你什麼事.
耶穌.
降生關你什麼事.
基督耶穌降世.
為了什麼.
拯救罪人.
如果你不是罪人.
你不和他拯救降生關你什麼事.
多個假.
可以放假.
不單止.
上帝叫他降生.
還叫他從死裡復活.
這個復活.
我們才可以得到永生.
耶穌沒有復活.
即是說沒有人會復活.
唯有信耶穌的人.
在基督耶穌裡面才會復活.
所以耶穌降生是事實.
耶穌受死是事實.
有復活節有聖誕節.
耶穌升天也是事實.
耶穌.
再來.
耶穌復活.
耶穌升天.
還有基督耶穌再來.
所以當你口裡.

$^{841}$肯承認他做主.
心裡信信他從死復活就必得救.
聖經第最後一節.
他口袋求告.
他每個人.
所以每個人都可以得到救恩.
只要你肯求告主名.
那怎麼求告呢.
聖經告訴我們.
當我們要求告耶穌的時候.
我們要禱告.
你問我怎麼祈禱.
我也不懂.
我現在會和大家一起.
做一個禱告.
這個禱告就是.
你剛才所聽到的.
無論是姐妹的見證.
無論你聽到聖經的說話.
你承認.
承不承認自己是罪人.
需不需要得到救恩.
只要你肯.
口裡承認耶穌是救主.
心裡相信.
神已經叫耶穌從死復活.
這個復活的大能.
不單單在我身上.
在姐妹身上.
在每一個弟兄姊妹身上.
今天我沒有夾到.
如果大家會唱奇異恩典的話.
我想蕭詩歌是一個很好的見證.
奇異恩典.
何等今天.
我罪已得赦免.
前我失喪.
今被塵埃.
核眼今得看見.
這是基督徒的見證.

$^{881}$我們今天只能夠有平安.
不是因為環境上.
最重要的一件事.
是心裡有放心和平安.
平安.
喜樂是一種盼望.
這個盼望只有在耶穌基督裡面才有.
你要考慮認耶穌為主.
心裡相信神.
叫祂從死復活.
就必得.
當你求告的時候.
這個平安就會想到你.
所以我會和大家一起做一個禱告.
希望在這個時候.
你跟著我一起祈禱.
我講一句請你跟著我講一句.
一起低頭我們一起祈禱.
基督徒都可以和我一起禱告.
作為你願意在神面前.
接受祂的愛.
領受祂的信心.
得著這個盼望的禱告.
我講一句請你跟著我講一句.
親愛的天父.
今天早上.
讓我聽到福音.
讓我們看到.
我們姐妹的見證.
也看到你自己.
華語的應許.
凡求告主明的.
就必得救.
我現在口裡承認.
我是一個罪人.
求你赦免我罪.
也求你賜給我永生.
活得豐盛.
有喜樂平安.
過一個榮神益人的生活.

$^{921}$求你誰聽我禱告.
豐盛.
我求你.
祂是神給我的.
奉耶穌名字祈求.
阿們.
剛才大家是實體說話的 是不是?.
不是虛擬的嘛.
不是AI來的嘛.
是不是? 是真的嘛.
所以有請舉手給我看.
頂尖妹妹 你有沒有帶朋友來呀?.
有沒有舉手呀? 有舉手的話 幫她握握手好不好?.
幫她握握手 恭喜她.
我相信我們有陪談員的 是不是?.
我想請剛才第一次這樣祈禱的.
應該很多人都是第一次這樣祈禱.
因為你從來不是跟我一起祈禱過的嘛.
不過你是未信主 未受洗的.
請你跟座位起立.
頂尖姐妹 你陪她一起起立好不好?.
一起起立 還有沒有?.
你剛才聽她說的嘛.
請出來 陪談員呀.
陪談員出來好不好?.
帶她出來 我們一起祈禱好不好?.
出來.
出來 出來 出來.
頂尖妹妹 陪她一起走出來啦.
你這麼辛苦 帶她來這裡了.
走那一兩步路.
恭喜你們 恭喜你們.
願意今天禱告.
剛才我聽到你說的.
你的起身呢.
聖經說.
凡在人面前肯承認耶穌是救主的.
將來在永遠的審判裡面的時候.
耶穌基督都承認我們是他的兒女.
所以你得享的平安喜樂是有盼望的.

$^{961}$還有沒有其他人呢?.
有請你到台前來我們一起禱告.
我們一起禱告.
你用奇妙的恩典.
帶領我們這群朋友.
來到你的面前.
主當他們聽見見證.
當他們看到你自己奇妙的作為.
也看到你聽到你自己說話的時候.
讓他們口裡肯認耶穌為主.
心裡相信神使耶穌重生你復活.
主你應許賜給他們平安.
賜給他們永生.
求主賜福恩代.
讓他們往後日子裡面.
得到你所賜的喜樂.
得到你所賜的平安.
滿有信望愛.
活在地上過豐盛的生活.
能夠得享永生.
一人得享永生.
奉耶穌名字祈求.
\newpage



\section{以賽亞書 43:18-19}
\label{sec:RkNxGeCkTHA}
\textbf{2025.10.19《信仰中不斷變新的東西》:以賽亞書43\_18-19(和修版) 陳韋安牧師}
\newline
\newline
連結: \href{https://youtube.com/watch?v=RkNxGeCkTHA}{\texttt{https://youtube.com/watch?v=RkNxGeCkTHA}} ~~~~ 語音日期: 2025-11-03
\newline
\newline
\hyperref[sec:Bm5kMjeNAdQ]{\small{< < < PREV SERMON < < <}}
~
\hyperref[sec:index_chronic]{\small{[返順時目]}}
~
\hyperref[sec:index_scriptual]{\small{[返順卷目]}}
~
\hyperref[sec:Uvg7nmv7Igo]{\small{> > > NEXT SERMON > > >}}
\newline
\newline
以賽亞書 43:18-19
\newline
\begin{longtable}{cl}
\hline
\hline
章節 & 經文 (和合本修訂版)\\
\hline
43:18 & \begin{tabularx}{0.7\textwidth}{X} 「你們不要追念從前的事，也不要思想古時的事。 \end{tabularx} \\ \\ \relax
43:19 & \begin{tabularx}{0.7\textwidth}{X} 看哪，我要行一件新事，如今就要顯明，你們豈不知道嗎？我必在曠野開道路，在沙漠開江河。 \end{tabularx} \\ \\
[1ex]
\hline
\hline
\end{longtable}
$^{1}$各位觀眾早安.
今天跟大家一起看經文.
《易賽亞書》第43章的經文.
經文這樣說.
「耶和華如此說,你們不要紀念從前的事,也不要思想古時的事..
看!我要做一件新事,如今要發言,你們豈不知道嗎?.
我必在曠野開道路,在沙漠開江河.」.
這段經文是關於我們如何思想信仰中的一些新事.
上次我來到港道時也有提及類似的話題.
提及我們生命中的根深.
今天我們會繼續談論類似的話題.
《易賽亞書》第43章.
我們首先要認識這段經文.
這段經文被稱為第二易賽亞的部分.
什麼是第二易賽亞呢?.
發現《易賽亞書》第43章.
經文有一個很明顯的時代轉換.
原來有的.
世聲有時空轉換的東西.
在第一章第43章.
易賽亞說的是以色列人被擄的前景.
他們將要被擄.
耶路撒冷將要被毀.
國家滅亡前的一個狀況.
但如果打開第43章之後.
發現整個時空突然轉換150年之後.
說的是滅亡之後的情況.
耶路撒冷已經被毀.
以色列已經被擄到巴比倫.
國家已經滅亡.
所以我們將第43章稱為第二易賽亞.
第二段時空的景象.
跟第一易賽亞相差150年.
不再預告耶路撒冷滅亡.
而是對於身處滅亡後的巴比倫人.
一群將來的人來說話.
不再勸導人悔改.
而是一群被擄的人.
他們怎樣歸回.
回到自己的家鄉信息.

$^{41}$因此《易賽亞書》第43章.
其實是一段以色列人被擄的時候.
上帝呼籲他們離開巴比倫.
回到自己的家鄉.
呼籲以色列人不要再逗留在巴比倫.
快點離開.
離開現在居住的地方.
不要再留戀這個異邦城市.
這個城市不再屬於你.
你應該回到自己的家鄉.
基本上整個第41到48章.
都不斷重複類似的信息.
不過對於這麼充滿盼望的喜訊.
這群身處巴比倫的人.
他們不覺得有什麼慶賀.
大概是一聲噢.
就算數.
他們不斷推搪去砌詞.
回到耶路撒冷的路途很遙遠.
回到家鄉找不到食物.
其實我們很好的來信服.
這個異邦的皇帝就好了.
甚至乎去享受被擄的快樂.
尋求苟且偷生的幸福.
坦白說就是習慣了.
已經習慣了.
習慣了被擄的生活.
習慣了自己的家鄉已經滅亡.
習慣了一個沒有未來.
沒有什麼期盼的信仰.
相比起耶路撒冷那個未知之數.
他們寧願活在一個沒有理想的狀態裡面.
認定自己一輩子就是這樣.
去容養自己在一個被擄的裡面.
認命不服妥協.
何錫爾曾在第一次會議.
《以上下書》第43章這段經文.
正是針對這個問題來寫的.
又說新的事情將要發生.
不要再停留在過去.

$^{81}$快點離開.
所以又說.
你們不要紀念從前的事.
也不要思想古時的事.
看來我要做一件新事.
如今要發現你們不知道.
我必在曠野開道路.
在沙漠開江河.
又說不要再記住那些舊事.
不要停留在以前的事.
英文更加直接翻譯.
Forget the former things.
Do not dwell on the past.
不要停留在過去.
不要思念舊事.
忘記以前的事.
因為現在將會有一件新事發生.
是的.
如果你有機會去外國.
去看朋友.
加拿大,英國.
北美的地方.
你會發現很多時候那些華僑.
其實他們是在這裡.
如果你去荷蘭的華僑.
去德美的華僑.
基本上整個是放棄了.
不知道大家有沒有去過北美的.
例如加拿大,溫哥華,多倫多.
的唐人街.
基本上找回很多以前的歌.
聽以前的歌.
還在播張學友,孫之鋒,宇宙.
林子祥的歌.
整個社區用80年代.
那時不久出香港為主題.
電視上仍然有Doodle姐.
年輕的Doodle姐.
Doodle在這裡.
萬千星星在學台慶.

$^{121}$整個唐人街是一個.
80年代的香港的狀況.
吃我們以前.
古典80年代的早餐.
很難找到.
很漂亮的煎蛋.
多士.
牛油.
很正宗的80年代早餐.
現在找不到.
飲館是監館.
是這樣的情況.
完全保留了.
香港那個年代的特色.
為什麼?因為他們移民的就是那個年代.
現在都是一樣.
英國移民都是一樣.
停在2000年代的香港.
都會保留在這裡.
究竟這班在巴比倫.
他們所謂的舊事.
是說什麼呢?.
我們再看前一篇文章.
第17節說.
在滄海中開渡.
在大水中開路.
使車輛.
馬匹軍兵.
湧出來一同躺下不再起來.
他們淹沒.
好像熄滅的燈火.
再看一次.
在滄海中開渡.
在大水中開路.
使車輛馬匹軍兵.
湧出來一同躺下不再起來.
他們淹沒好像熄滅的燈火.
大家如果有些印象.
覺得這篇文章說什麼?.
這篇文章說什麼?.

$^{161}$就是以前的事.
就是出埃及的時候的情況.
雖然沒有記載很具體的時間.
但聽的情景都能記住.
最經典的那一幕.
就是摩西分洪海的廣播.
前無去路.
大堆的埃及軍隊追回來的時候.
一下子一按.
整個海就分開了兩邊.
突然間風浪吹過來.
海就突然間分開兩邊.
這樣就變成一個非常壯觀的情況.
當年N年前.
出埃及的情景.
出埃及是所有以色列人的集體回憶.
出埃及正正是以色列人的相思風雨中.
大家不要小看出埃及的三個字.
不要低估這個字的形式.
雖然出埃及已經是他們千年的古代的事情.
不過對於這班人來說.
這是一個世代相傳.
不停傳頌滾瓜爛熟的故事.
出埃及是小朋友的Bedtime story.
是一個每年大家新年茶餘飯後的以前的話題.
出埃及是以色列人的血液.
就像我們的金庸小說.
大家有沒有看過神鵰俠侶.
大家有吧.
應該問你看什麼版本的神鵰俠侶.
你們看的是劉德華版本.
還是古天樂版本.
還是近期的黃曉明那些.
你們看什麼版本的.
我是古天樂.
我也是看劉德華的.
重播的時候看的是劉德華.
基本上是一個每個人都認識的故事.
不過比起神鵰俠侶.
出了一個更加重要的意義.

$^{201}$他們不單單是熟悉這個故事的情況.
他們更加是擁有這個故事.
他們是屬於這個故事的.
他們不單單是聽這個故事.
更加是投入自己去這個故事裡面.
相信自己是置身於這個紅海裡面.
投入自己在這個曠野裡面的人.
相信是帶領他們.
即是自己離開法老的手.
耶和華是領我們出埃及的神.
我們是從埃及被召出來的人.
整個傳統告訴你.
他們是一個出埃及的群體.
始終埃及是這群被老式的民的血液裡面.
那個家傳戶曉的傳統.
我們也有很多這樣的傳統教會.
由我們錦繡堂到紅恩堂.
這麼多年來.
那些東西是沒有什麼變化的.
大家都會敬拜.
唱兩三首歌.
然後就會站起來祈禱.
然後就笑笑順經.
我公經也是這樣.
基本上紅恩宣德和錦宣都是差不多的情況.
總是有星期三祈禱會.
不是星期四.
總是有星期三祈禱會.
之後有牧師祝福.
很多東西大家都是很習慣的.
有師班.
又有很多不同的師班婆.
大家都很熟悉的情況.
不過因為這樣.
這群在巴比倫的幾十年.
出現了一個很尷尬.
很矛盾的現象.
基本上他們是一群堅守著出埃及傳統的幾十年.
大家都很尊重.
很熟悉這個故事.

$^{241}$但其實他們一點都不出埃及.
為什麼呢.
他們不願意離開巴比倫.
他們很虔誠地重複出埃及的傳統.
熟悉每一個細節.
認同裡面的精神.
不過他們竟然甘願容讓自己在另一個埃及生活.
明白嗎.
在巴比倫是最不出埃及的事情.
他們心裡承認.
他們相信我們是一群出埃及的群體.
但他們一點都不願意離開另一個埃及.
最後成為了一個很熟悉.
大家都很清楚.
但沒有靈魂的故事.
就像一個公園.
大家都很熟悉.
家媳戶曉.
很受歡迎.
很多主題公園的現象.
分紅海盜船.
雲柱跳樓機.
很多這些都會玩耍.
大家都會買些紀念品來紀念.
但沒有一個真正認真看待這些事情.
成為一個生活中真正要行走的細節.
這正正是耶和華上帝呼籲這些人背後的原因.
耶和華說不要再紀念從前的事.
也不要思想古時的事.
要做一件新事.
上帝提醒這些不認識的人.
不要再活在昔日的陰影.
不要停留在傳統的家喪.
我們要在這一代人做一件新的事情.
上帝在每一個年代都做一些新事.
新是一個上帝的形容詞.
所以變成新的狀態.
是一個上帝很專有的狀態.
每天的精神都是新的恩典.
全都要向耶和華唱新歌.

$^{281}$新酒,新聖殿,新耶和華撒冷,新天,新地,新造的人.
不過一提起新造人.
我們都說過.
我們這群人都是新造人.
但又好像有點舊舊.
洪恩堂是不是也叫新教會?.
都不知道了.
多少年了?14年了.
新但又不算很新.
不可以說自己是很新的教會.
都14年了.
不知道大家有沒有在這裡開始.
跟洪恩堂一起開始的弟兄姊妹.
很多年前在這裡租的地方.
兩層高.
有些新的屋苑在這裡.
有些新的學校,新的屋苑在這裡.
很多新的事情發生在這裡.
大家在這裡建立一間教會出來.
這麼多年來.
有些東西開始有點舊了.
港台,我們的牆.
開始不算很新.
很多新的東西都不再是新的東西.
所以我們問究竟我們新事是什麼呢?.
我們這群基督徒.
在這裡一個新的地方,新的時宮.
我們問究竟怎樣可以找回新的方向呢?.
上帝仍然是一個行新事的上帝.
我要向耶和華唱新歌.
這個說法仍然是真的.
剛才那個師班唱的也是一首新歌.
很多年前.
其實是很突破.
唱一些粵曲的方式.
在師班裡面唱出來.
這是一個很新的破格的方法.
不是一些西樂.
用一些粵曲來獻詩給上帝.
本身是一件很新的事情.

$^{321}$雖然粵曲好像很舊.
但我們仍然有一個很突破的精神.
來將上帝用一個形式去讚美他.
每一個年代的基督徒.
都屬於他有一個自己的新的新歌.
我剛剛才幫一個青年聖歌的聚會做講院.
大家知道青年聖歌是吧?.
青年聖歌是很諷刺的.
青年聖歌已經不是青年了.
是60年代的青年.
60年代青年一班青年人去唱的詩歌.
聽說以前是牧師不讓他們在崇拜唱.
逼著團契只能唱紅青藍青綠青.
其實是很反叛的.
因為崇拜不讓唱.
我就要在團契裡面唱藍青紅青的歌.
今天我就去那些聚會做講道.
基本上是一個青年聖歌的金曲之夜.
有詩班有班的會眾.
大家都去唱一年多.
唱15首青年聖歌.
我講15分鐘左右.
我講少一點.
最後是一個配角.
唱15首這樣的歌.
有詩班唱就唱.
大家懷緬以前的信仰.
那個美好.
我自己很支持青年聖歌.
我自己很喜歡.
因為是很好的歌詞.
雖然我教會是很多Band隊.
現代歌.
但我覺得都是好東西.
有好地方.
所以很有趣.
青年聖歌.
以前的新歌.
今天被稱之為.
好像不是很新的東西.

$^{361}$所以每個年代的人.
都有屬於自己的新事情.
上帝一次又一次.
讓每一代人發現那個年代的新事.
但實際上出埃及.
本來就是一件上帝的新事.
意思是在一千多年前的出埃及.
本來是一件新事.
誰想過住了幾百年.
埃及竟然可以離開.
誰想過前面有一個大海.
是很會分開的.
誰想過天會掉下來.
這些麻辣是可以吃的.
上帝做了一些很破格的事情.
做一些人沒想過的事情.
但上帝這些事情.
是給每一代人的那一代的新事.
我們不能夠活在別人的新事裡.
那些就變成不是你的新事.
你只能找回你這一代人的新事.
上帝不是叫我們推翻過去.
不是叫我們零起爐灶.
要推翻別人的東西.
但我們這一代人.
要找回我們這一代人.
屬於我們的新事情.
很有趣.
如果你看回經文.
經文是非常矛盾的.
上帝明明叫以色列人忘記過去.
叫他們不要再提起以前的舊事.
他要故意去舊事重提.
很有趣.
又要你忘記它.
不斷叫你提起它.
特意貼一張保紙.
叫你不要忘記.
這樣就記得.
所以其實不是去忘記.

$^{401}$其實是忘記不了的.
不是叫我們忘記.
實際上忘記是不需要的.
所以說一個情況.
我們需要暫時去忘記.
就是這個舊事的東西.
去窒礙我們去發現上帝新事的時候.
我們需要暫時放下去忘記這些舊的東西.
這些傳統.
教會是很重要的傳統.
由古時到現在.
很多傳統都很重要.
傳統是有意義的.
傳統是幫助我們迎面將來.
一旦傳統埋沒我們的將來.
成為我們面對將來的枷鎖的時候.
上帝就說你們不要再去思念古時的事.
這個時候就要.
因此新的事情不是要去推翻傳統.
也不是跟傳統對立.
而是我們要去找回我們這一代人.
新的事情.
這時候就明白第十九節的意思.
又說漢蘭我要做一件新事.
如今要發現你們是不知道嗎.
我必在曠野開道路.
在沙漠開江河.
上帝說那件新事是什麼.
他剛才說要在曠野開道路.
在沙漠開江河.
發現這句話的景象有點熟悉.
其實這是另一次出埃及的過程.
一件全新版本的出埃及.
很有趣的是.
這件事跟昏龍海是剛剛相反的.
以前是怎樣.
昏龍海分開兩邊見到陸地.
現在是剛好相反.
在陸地見到海洋見到江河.
上帝刻意把以前的反調.

$^{441}$以前的昏龍海見到海洋見到地.
現在是在地上見到江河.
是剛好相反的事情.
上帝要在你沒有經驗的過去經驗裡.
去跟以前的事情作出反調.
上帝要求斯蒂文不要停留在以前的做法裡.
再一次經歷新一次的出埃及.
一個憑著信心要經歷的旅程.
忘記以前的神鵰俠侶.
自己經歷一次新一次的出埃及經驗.
忘記以前列祖的故事.
經歷一次自己神給你的故事.
上帝要稱之為新出埃及.
New Ecclesiastical.
兩個唯一不同的地方.
很多不同的地方.
唯一相同的地方是什麼.
就是上帝估你不到.
兩個都是估不到的.
我們以前叫那些人估不到很多的東西.
今年也是一樣.
上帝要他們去到洞穴都估不到上帝的座位.
新出埃及是這班人為上帝經歷的新經驗.
因此對我們頂尖的朋友來說.
一個教會.
我們怎樣去看見上帝給我們的新事呢.
在這個地方裡.
雖然是十四年的時間.
但仍然有很多新的事情可以去做.
我經常說.
我教的受教會有一個很重要的精神.
就是開荒.
開荒吃苦火熱.
開荒是一件很特別的事情.
因為它正正是一個.
由零到一.
在一個沙漠裡開江河的過程.
香港裡好像已經沒有什麼人去做開荒的事情.
我們做一些新的事情.
但是開荒的精神.

$^{481}$永遠都是我們在裡面很重要的.
因為上帝最重要的就是那位開荒的人.
他叫我們能夠去冒著一定的風險.
懷著信心.
為上帝做一些新的事情.
縱然這個地方已經是.
一段時間.
我們的年紀.
我們的外貌.
雖然都是一段時間.
但是上帝給我們裡面的心.
仍然是可以改變.
這就是我上次說到的訊息.
我們裡面的人仍然不斷的更新.
這個地方也一樣.
教會仍然可以為上帝做一些新的事情.
這半年.
有一個奶茶茶水檔.
我們看到很多很多新的可以做出來.
可以不斷的在這裡想想.
有什麼新的事情可以做.
傳統是重要的.
但是我們仍然可以為上帝做一些新的事情.
所以我們經常說.
有一個詞語叫做Start-Upper.
就是起初的東西.
開一間餐廳.
開一個茶水檔.
開一個店鋪.
其實很多不信主的人都在做這些事情.
很多香港人都開很多創業.
新的事情.
基督徒反而沒有這樣的膽量.
明明有神在背後.
明明可以憑信心去做.
但是反而是太過不敢去冒那個風險.
我們作為基督徒.
是一個做新事的人.
更加值得懷著信心去為神做一件新事.
是有機會會倒閉的.

$^{521}$外面的店鋪.
我們知道教會是不會倒閉的.
起碼上帝的國度是不會倒閉的.
只要是願意.
是神願意帶領我們的時候.
我們仍然可以懷著一點勇氣.
為神做一些新的事情.
我們一起祈禱.
看看神怎樣來到我們身邊.
讓我們有新的事情可以做.
因為你是那位叫做以色列人.
出埃及的上帝.
你帶領不同年代的.
屬於你的子民.
去衝破我們的界限.
去打破我們舊有的一些枷鎖.
去為你發展一些新的事情.
主要求你幫我們能夠在裡面.
仍然懷著這份勇氣.
這份信心.
去為你做新的事情.
保守我們的教會.
我們的教會仍然在當中.
是一間很新的教會.
因為我們的心.
我們裡面的信心仍然是新的.
我們仍然願意在裡面尋求.
你的帶領.
讓我們教會裡面不同的事工.
我們為你所做的事情.
都仍然在裡面有信心.
憑著勇氣去做.
更新我們的教會.
讓我們能夠不斷去.
有創新的想法.
去為你做事情.
讓我們一起同心合意的.
去更新我們的教會.
在這裡看到新的景象.
奉主明求 阿們.

$^{561}$謝謝大家.
\newpage



\section{馬太福音 16:13-19}
\label{sec:Uvg7nmv7Igo}
\textbf{2025.11.02 《教會的基石》 經文:馬太福音16\_13-19 陳麗蟬牧師}
\newline
\newline
連結: \href{https://youtube.com/watch?v=Uvg7nmv7Igo}{\texttt{https://youtube.com/watch?v=Uvg7nmv7Igo}} ~~~~ 語音日期: 2025-11-04
\newline
\newline
\hyperref[sec:RkNxGeCkTHA]{\small{< < < PREV SERMON < < <}}
~
\hyperref[sec:index_chronic]{\small{[返順時目]}}
~
\hyperref[sec:index_scriptual]{\small{[返順卷目]}}
~
\hyperref[sec:7BCvZD9NU64]{\small{> > > NEXT SERMON > > >}}
\newline
\newline
馬太福音 16:13-19
\newline
\begin{longtable}{cl}
\hline
\hline
章節 & 經文 (和合本修訂版)\\
\hline
16:13 & \begin{tabularx}{0.7\textwidth}{X} 耶穌到了凱撒利亞．腓立比的境內，就問門徒：「人們說人子是誰？」 \end{tabularx} \\ \\ \relax
16:14 & \begin{tabularx}{0.7\textwidth}{X} 他們說：「有人說是施洗的約翰；有人說是以利亞；又有人說是耶利米或是先知中的一位。」 \end{tabularx} \\ \\ \relax
16:15 & \begin{tabularx}{0.7\textwidth}{X} 耶穌問他們：「你們說我是誰？」 \end{tabularx} \\ \\ \relax
16:16 & \begin{tabularx}{0.7\textwidth}{X} 西門‧彼得回答說：「你是基督，是永生神的兒子。」 \end{tabularx} \\ \\ \relax
16:17 & \begin{tabularx}{0.7\textwidth}{X} 耶穌回答他說：「約拿的兒子西門，你是有福的！因為這不是屬血肉的啟示你的，而是我在天上的父啟示的。 \end{tabularx} \\ \\ \relax
16:18 & \begin{tabularx}{0.7\textwidth}{X} 我還告訴你，你是彼得，我要把我的教會建造在這磐石上，陰間的權柄不能勝過它。 \end{tabularx} \\ \\ \relax
16:19 & \begin{tabularx}{0.7\textwidth}{X} 我要把天國的鑰匙給你，凡你在地上所捆綁的，在天上也要捆綁；凡你在地上所釋放的，在天上也要釋放。」 \end{tabularx} \\ \\
[1ex]
\hline
\hline
\end{longtable}
$^{1}$主席.
陳婉嫻議員.
主席.
馬太福音第1至16章.
主席.
當時耶穌問那群門徒.
你心目當中覺得我是誰呢?.
我一直想邀請弟兄姊妹.
你去想一想.
在你心目當中.
你覺得耶穌是你的誰呢?.
每一次當我要問準備受浸的弟兄姊妹.
我一定問一個問題.
你相信耶穌.
你覺得耶穌在你的生命當中.
成就了一件很大的事.
你可不可以和我分享呢?.
通常他們會說.
耶穌聽了我的禱告.
那時我怎樣怎樣.
耶穌聽了我的禱告.
或者本來有一個大病.
耶穌又醫治了我.
有很多很多的見證.
我說在你的生命當中.
耶穌為你做了一件很大的事.
是什麼事呢?.
你總是說很多見證.
所以和我們聊天.
很多見證聽.
但其實我最想聽的.
就是他們要告訴我.
耶穌基督為我釘十字架.
祂是基督.
所以有時在一些報道會.
或者一些方式聚會裡面.
我們同樣聽到很多很多人說見證.
但當一個報道會裡面.
不說耶穌基督被釘十字架.
不說耶穌基督為我們寫生.

$^{41}$我都不會覺得那是一個真正的報道會.
因為所有的神蹟奇事.
所有宗教都可以.
甚至他們說有求必應.
說這麼多的神蹟.
但其他宗教都可以.
唯有耶穌基督被釘.
其他的宗教都沒有.
當彼得這樣去承認.
耶穌就是基督.
就是永生神的兒子的時候.
我們就看看.
基督這個字是甚麼意思呢.
原文的意思是一個受高者.
這個受高者在舊約裡面.
受高的是甚麼呢.
就是先知,祭司和君王.
是神自己所選立的.
所以那班猶太人知道.
受高者基督.
就是神自己親自選立的君王,祭司和先知.
另外他說永生神的兒子.
就是這個神不只是這段時間有在.
而是昔在,今在,以後永在的上帝.
我們知道耶穌基督就是那位.
為我們去被釘的救世主.
這個彌賽亞.
我們多不多覺得自己是罪人呢.
少少吧.
有些人說多少少.
我自己覺得自己是罪人的時候.
就是當我決志那一刻.
我當然很記得.
很多很多年前.
有同學帶我去一個報道會.
那位牧師一直說.
我們每一個都罪人.
我立即在心裡反駁他.
我是好好小姐.
所有同學都說我是好好小姐.

$^{81}$我不是罪人.
當我這樣想的時候.
裡面就開始出現我自己平時的心思,疑念.
在外表來說.
很多同學都覺得我是很好的.
有求必應的.
但在裡面拍出來.
當有些同學有很多提示,考試.
好筆記.
別人分享出來給我.
但當我自己有的時候.
放在書包底.
一定不會分出來.
我看到我自己的自私.
雖然外表很好.
我記得有同學約我去街.
我在電車站等了一個小時.
心裡想著.
豈有此理.
我跟他去街第二次.
但當同學來到的時候.
「對不起,對不起,利是員」.
我有表現.
「不要緊,不要緊,你有重要事」.
我看到我裡面的自己的虛偽.
在當中,我就很清楚知道.
我是一個罪人.
我是耶穌基督的教授.
所以當我們要知道.
我們所信的耶穌.
是那位用寶血洗清我們罪孽的耶穌.
我們要很清楚知道.
另外,祂是永生神的兒子.
我們是否感覺到.
我們上帝是一個永生神.
是永遠永遠都不改變的.
何時你會感覺到.
強烈一點,祂是永生的呢?.
大家都拍過拖吧.
大家失過戀了嗎?.

$^{121}$我告訴你.
雖然我沒有結婚.
但我都失過戀不止一次.
其實那次是第一次.
我掙扎了大半年的感情.
才投入和一位弟兄拍拖.
拍拖時日夜聊電話.
那時我已經在教書.
上學放學又經常走在一起.
當他決定了.
我們一起去看看將來的日子.
拍拖一年後.
我發覺他放學沒有等我一起就走了.
平時晚上打電話的時候.
他又不打電話.
我打電話給他的時候.
他又不接聽.
我一路追查.
原來他在教會裡認識了另一個人.
當時我當然很生氣.
人都是這樣.
一定會很生氣.
但當中我覺得要接受.
我記得在某個星期日下午.
平時的時間.
我們會一起去吃飯,逛街.
那天在家裡很悶.
不斷地祈禱.
在祈禱當中.
上帝讓我領受到一件事.
人變得很快.
當我們決定了.
不如我們試試.
看看是否真的.
當時我們都願意投入這段感情.
向對方說一說.
我們一起去觀察.
但一年後就變了.
變了的時候.
竟然不敢跟我說找了另一個.

$^{161}$靜靜地就走了.
我覺得人的變.
今天可以跟你山盟海誓.
但過了不知多少天之後.
他要變起來就變了.
在那天一直祈禱的當中.
我真的很清楚地感受到.
只有我們的上帝.
祂是永遠都不會變.
我未出生的時候.
祂已經知道我.
我出生的時候.
祂很愛我.
祂不單今天很愛我.
祂也愛我到永恆.
我覺得在這個世界上.
只有上帝永恆不變.
所以在那次裡.
我是第一次跟上帝有一個承諾.
我要一生一世地去事奉祂.
人的不可靠.
讓我找到上帝.
才可以終生地去投靠祂.
第二方面.
我們要看看耶穌基督.
在這個裡面.
說祂是永生神的兒子.
耶穌基督怎樣去回答彼得呢?.
要讚賞祂.
我要把我的教會建在盤石上.
陰間的權柄不能勝過祂.
這幅圖畫.
就是我們剛才看的.
在菲勒比大山的裡面.
其中一個洞.
這個洞是拜潘神的洞.
他們覺得洞裡面.
就是魔鬼撒旦住的地方.
他說進入這個洞裡就是死亡.
所以他們拜祭的祭物.

$^{201}$會被扔進這個洞裡.
當裡面的洞有血流出來的時候.
裡面的潘神就接納了那個祭物.
其實你殺了一隻雞.
不是殺雞.
不知道是殺什麼.
當你扔進洞裡的時候.
自然就會有血流出來.
他們覺得這是陰間的入口.
其實也很恐怖.
陰間的入口.
陰間就是死了的靈魂的居所.
因為陰間的權柄.
代表了魔鬼的權勢.
我們一說到魔鬼.
一說到陰間有沒有覺得毛骨悚然?.
有一點吧.
因為冷氣很冷.
但有時也是這樣的.
前幾天是萬聖節.
當你晚上出去地鐵的時候.
看到有人扮鬼扮馬.
我們在錦繡堂走的時候.
九點多十點的時候.
很精彩.
但我們要想一想.
這些東西.
我們看的時候很害怕.
但我們一定要知道一件事.
這些魔鬼殺蛋的東西.
他們不是全能的.
他們只是不斷地嚇人.
我們要記住.
面對的時候.
我們是有驚無險.
千萬不要被它嚇倒.
它永遠都不會勝過耶穌基督.
耶穌基督才是那位全能的.
耶穌基督說.
我要將天國的弱士賜給你.

$^{241}$凡你在地上高所捆綁的.
在天上也要被捆綁.
在地上高所釋放的.
在天上高也要被釋放.
我們知道彼得列這個名字.
這個名字的意思是石頭.
石頭不是大石頭.
是小石頭.
剛才耶穌說.
我要將我的教會建造在盆石上.
我們很清楚知道.
耶穌基督為何帶這班門徒.
去到凱撒尼亞肥納比這個地方.
去介紹自己真正的身份.
凱撒尼亞肥納比這個地方.
剛才我們說是一座石山.
而彼得這個名字.
是一個石頭的意思.
他要這班門徒知道.
這個真正穩定的大盆石.
他跟彼得說.
既然你知道耶穌基督就是彌賽亞.
就是永生神的兒子.
他就宣告.
要將教會建造在盆石上.
天主教誤會了.
以為這個教會就是彼得.
將權柄給了彼得.
然後給予教皇.
其實不是.
他說從彼得的口中.
剛才所說到我們信仰的基礎.
就是耶穌基督是彌賽亞.
耶穌基督就是永生神的兒子.
所以耶穌基督是說.
彼得很清楚知道信仰的根基.
所以這個就是一個大盆石.
是我們教會的根基.
所以當我們很清楚知道.
我們所信的不是有求必應的上帝.

$^{281}$是一個為我們犧牲.
為我們捨己的上帝.
並且是永恆不變的上帝.
唯有這樣的時候.
這個才是實證.
這個才是穩定.
上星期是我們錦宣間的堂慶.
寫《牧者之言》的時候.
找回我們以前的歷史.
以前又舉辦很多音樂會.
籌款音樂會.
然後很多旅行.
又健康檢查.
又老人院探訪.
又讚美操.
又排排舞.
又拉丁舞.
然後又朝聖茶室.
又好像紅恩堂一樣.
有很多咖啡奶茶.
現在還有什麼爵士鼓.
又舉辦古班.
然後我們又去新田部屋.
或者簡約公屋.
每逢有什麼聚會.
我們都會舉辦咖啡奶茶.
又交流美食.
又烹飪心得.
又端節製作鹹肉粽.
中秋節剛製作完冰皮月餅.
一口氣聽完都是什麼.
喝喝吃吃.
全部都喝喝吃吃.
但我們感恩.
當我們開堂董會.
開執事會的時候.
我們都會提醒一件事.
每一個這些活動.
都必須加上福音元素在當中.
如果沒有了這個元素的時候.

$^{321}$我們只是一些俱樂部.
無論我們錦繡堂.
或者紅恩堂.
我們都知道.
有時候去嶺南大學.
又有這些聚會.
有時候聽到弟兄姊妹回來.
我感恩的.
不是告訴我.
我沖了多少杯咖啡.
沖了多少杯奶茶.
他們不是說沖到手軟.
而是告訴我.
他們跟人聊天.
聊到口水都乾.
這個就對了.
因為沖奶茶沖咖啡.
但原來他們在當中.
一直跟人聊天.
有機會就說說耶穌.
這個才是真的.
所以凡你在地上所捆綁的.
在天上耶穌基督會捆綁.
凡我們在地上所釋放的.
在天上耶穌基督都會釋放.
這個解釋了.
耶穌基督在地上的時候.
祂就自己來做.
但當祂現在已經升了天.
已經升了天.
祂就將釋放人的生命的權柄.
使命來交給我們.
我們需要來運用.
我經常按錯鍵.
《以賽亞先知書》第61章第一至二節.
預言耶穌基督要做的工作.
我們一起讀.
1,2,3.
主耶和華的靈在我身上.
因為耶和華用膏告我.

$^{361}$叫我報好信息給貧窮的人.
差遣我醫好傷心的人.
報告被擄得釋放.
被捆綁得自由.
宣告耶和華的恩賢.
望我們上帝報仇的日子.
安慰所有被愛的人.
所以耶穌基督已經做完了.
去升了天.
然後在地上我們就要去承擔使命.
所以在教會上.
我們的目標.
就是將人從罪惡捆綁中釋放出來.
無論我們做什麼.
任何活動都不能沒有福音的元素.
我們需要釋放人的時候.
我們要想一想.
我們信了耶穌.
我們自己的生命.
有沒有脫離了罪惡的捆綁呢?.
我們要想一想.
我們自己有沒有脫離了罪惡的捆綁呢?.
在情慾上.
在人與人之間的紛爭.
仇恨.
嫉妒的裡面.
有沒有釋放了出來呢?.
在我家庭裡面見證又如何呢?.
我與我的丈夫.
我與我的妻子.
我與我的兒女的關係又如何呢?.
我有沒有從貪婪中釋放出來呢?.
耶穌基督給我們一個榜樣.
就是要犧牲.
要捨己.
而不是別人來犧牲.
我自己就不斷去奪取.
所以當我們說.
我們要將人從罪惡中釋放出來的時候.
我們必須要時時檢討我們自己.

$^{401}$我們的生命都要從罪惡中釋放出來.
否則我們沒有能力去做到.
求主幫助我們.
所以我們要在這裡.
很想在今天大家十四歲生日的時候.
我們與上帝都有一個立志.
我們一定要活出一個應該有信徒的角色.
我們就用這兩句話.
我們一起來彼此的香神.
來立約.
我們就讀後面那四行就行了.
1,2,3.
我要為主作世上的光.
照亮在黑暗地的人.
我要為主作世上的業.
生命聖潔.
生活和平.
我們同心吐告.
親愛的天父上帝.
因為是你先拯救我們.
以致我們能夠從罪惡的裡面.
我們被釋放出來.
主大在世上實在有很多的隱憂.
很多的事情令我們有時會軟弱.
求主聖靈幫助我們.
以致我們能夠去懸崖勒馬.
求主親自到來的工作.
幫助幫助我們.
真是努力在方的工作裡面.
去榮耀主的名字.
祝福每一位洪恩家的弟兄姊妹.
英國的生命.
榮耀你的生命.
我們的生活就要不斷見證主你的榮美.
求你使用.
以致在一年一年一年.
我們就看到主你不斷將福音的果子.
加添到洪恩家.
求主賜福.
我們多謝你垂聽禱告.

$^{441}$奉靠主基督聖名.
阿姆.
\newpage



\section{傳道書 7:1-18}
\label{sec:7BCvZD9NU64}
\textbf{2025.11.9《義人的滅亡》:傳道書7\_1-18 (和修版) 講員 :張雲開老師}
\newline
\newline
連結: \href{https://youtube.com/watch?v=7BCvZD9NU64}{\texttt{https://youtube.com/watch?v=7BCvZD9NU64}} ~~~~ 語音日期: 2025-11-12
\newline
\newline
\hyperref[sec:Uvg7nmv7Igo]{\small{< < < PREV SERMON < < <}}
~
\hyperref[sec:index_chronic]{\small{[返順時目]}}
~
\hyperref[sec:index_scriptual]{\small{[返順卷目]}}
~
\hyperref[sec:7vPzFHXDGtY]{\small{> > > NEXT SERMON > > >}}
\newline
\newline
傳道書 7:1-18
\newline
\begin{longtable}{cl}
\hline
\hline
章節 & 經文 (和合本修訂版)\\
\hline
7:1 & \begin{tabularx}{0.7\textwidth}{X} 名譽強如美好的膏油，人死去的日子勝過他出生的日子。 \end{tabularx} \\ \\ \relax
7:2 & \begin{tabularx}{0.7\textwidth}{X} 往喪家去，強如往宴樂的家，因為死是眾人的結局，活人必將這事放在心上。 \end{tabularx} \\ \\ \relax
7:3 & \begin{tabularx}{0.7\textwidth}{X} 憂愁強如喜笑，因為面帶愁容，終必使心喜樂。 \end{tabularx} \\ \\ \relax
7:4 & \begin{tabularx}{0.7\textwidth}{X} 智慧人的心在遭喪之家；愚昧人的心在快樂之家。 \end{tabularx} \\ \\ \relax
7:5 & \begin{tabularx}{0.7\textwidth}{X} 聽智慧人的責備，強如聽愚昧人歌唱； \end{tabularx} \\ \\ \relax
7:6 & \begin{tabularx}{0.7\textwidth}{X} 因為愚昧人的笑聲，好像鍋子下面燒荊棘的爆聲，這也是虛空。 \end{tabularx} \\ \\ \relax
7:7 & \begin{tabularx}{0.7\textwidth}{X} 勒索使智慧人變為愚妄，賄賂能敗壞人的心。 \end{tabularx} \\ \\ \relax
7:8 & \begin{tabularx}{0.7\textwidth}{X} 事情的終局強如它的起頭；存心忍耐的，勝過居心驕傲的。 \end{tabularx} \\ \\ \relax
7:9 & \begin{tabularx}{0.7\textwidth}{X} 你的心不要急躁惱怒，因為惱怒存在愚昧人的懷中。 \end{tabularx} \\ \\ \relax
7:10 & \begin{tabularx}{0.7\textwidth}{X} 不要說：為甚麼先前的日子強過現今的日子呢？你這樣問不是出於智慧。 \end{tabularx} \\ \\ \relax
7:11 & \begin{tabularx}{0.7\textwidth}{X} 智慧加上產業是美好的，對見天日的人都有益處。 \end{tabularx} \\ \\ \relax
7:12 & \begin{tabularx}{0.7\textwidth}{X} 因為智慧庇護人，好像金錢庇護人一樣；智慧能保全智慧者的生命，這就是知識的益處。 \end{tabularx} \\ \\ \relax
7:13 & \begin{tabularx}{0.7\textwidth}{X} 你要觀看神的作為，誰能使他所彎曲的變直呢？ \end{tabularx} \\ \\ \relax
7:14 & \begin{tabularx}{0.7\textwidth}{X} 順利時要喜樂；患難時當思考。神使這兩樣都發生，因此，人不知將會發生甚麼事。 \end{tabularx} \\ \\ \relax
7:15 & \begin{tabularx}{0.7\textwidth}{X} 在虛度的日子裡，我見過各樣的事情，義人在他的義中滅亡，惡人在他的惡中倒享長壽。 \end{tabularx} \\ \\ \relax
7:16 & \begin{tabularx}{0.7\textwidth}{X} 不要行義過分，也不要過於自逞智慧，何必自取敗亡呢？ \end{tabularx} \\ \\ \relax
7:17 & \begin{tabularx}{0.7\textwidth}{X} 不要行惡過分，也不要為人愚昧，何必未到期而死呢？ \end{tabularx} \\ \\ \relax
7:18 & \begin{tabularx}{0.7\textwidth}{X} 你持守這個，那個也不要鬆手才好。敬畏神的人，這一切都能兼得。 \end{tabularx} \\ \\
[1ex]
\hline
\hline
\end{longtable}
$^{1}$各位觀眾早安.
很高興再一次在主日的時候跟大家一起敬拜上帝.
分享聖經裡面的話語.
今天跟大家讀剛才那段經文.
我不知道大家有沒有機會來過這個講座.
還沒有嗎?.
他是香港華人界的猶太人智慧文學的專家.
所謂智慧文學.
大家知道舊約哪幾本書是屬於猶太人智慧文學的例子?.
在詩篇之前的《約伯記》.
詩篇之後的《箴言傳道書》.
這裡算是起碼這三本書.
但是其實在猶太人裡面智慧的說話.
智慧的內容不只是這三本書.
如果你看的話.
其實在先知書,詩歌書,詩篇裡面也有很多這類的說話.
大小先知書裡面也有.
我們今天思考的這段經文.
也是一段時間裡面對我比較困擾.
我不知道大家之前有沒有讀過.
如果你讀過或者今天才第一次接觸這段經文.
有沒有困擾你?.
對我的困擾.
有些事情已經知道了.
一輩子都困擾的.
所以如果不困擾你.
我也覺得有一點奇怪.
因為就算你今天第一次接觸這段經文.
其實你過去應該也被我們人生當中某些事情困擾的.
可能也見過.
或者甚至你自己經歷過.
就在第十五節所說.
「二人行義,反智滅亡,惡人行惡,倒響長壽」.
我們可以將經文.
其實一直都會看著經文的.
這是一件跟我們直覺來說不應該發生的事情.
我們說「善有善報,惡有惡報」.
這裡有,我可以自己控制.
忘記了.
惡有惡報.

$^{41}$還有兩句,我們等一下再處理那兩句.
中國人的智慧言語.
大家要明白.
我們說到智慧的文學.
或者智慧的說話.
不是猶太人獨有的.
也不是聖經獨有的.
是每個人類的文化.
他自己群體的文化裡面.
都會在他歷史的過程當中.
都會產生的一些文化產物的說話.
所以我們說.
「真言」,「言語」,「警句」.
警告人的說話.
這些全部都是同一類型的說話.
為什麼會在人類文化裡面產生這些說話呢?.
一般來說.
如果人類不是群居的.
一個人會不會自己.
也會像泰山那樣.
在森林裡面自己一個人.
還有一些猩猩.
其他的森林動物.
起碼一起住.
他會不會自己也得出一些智慧的說話呢?.
其實應該也有的.
因為智慧的說話.
不單止是人群一起共居.
一起生活的時候.
你要遷就我.
我要遷就你.
結果我們的生活.
就不可以你想做什麼就做什麼的.
是不是?.
我們做人.
待人接物.
小時候爸爸媽媽教我們.
要聰明一點.
說話要想過才說.
不要得罪人多.

$^{81}$然後呢.
稱呼人少.
得罪人多.
稱呼人少.
也要有人.
這句話才會有效.
泰山自己一個人.
自己森林.
他說什麼都可以.
他會得罪誰呢?.
得罪猩猩嗎?.
猩猩也不知道他說什麼.
那是怎樣呢?.
但是他也需要智慧的.
智慧的文學.
猶太人智慧的文學.
其實所有的文化.
不單止是人群共處.
人與人之間住在一起.
或者是格里倫舍.
或者是有家庭.
家庭裡面有夫婦.
夫婦兩個人也是.
不單止是你跟其他人一起生活的時候.
你就不能夠做.
你說想怎樣就怎樣.
人又想怎樣就怎樣.
你又想怎樣就怎樣.
大家一起.
什麼都沒有了.
一下子就爆炸.
不單止是這樣.
就算一個人也好.
因為他生活的世界.
除了他之外.
是沒有其他人了.
泰山在森林裡面生活.
是沒有其他人了.
但是還有大自然.
是不是?.

$^{121}$有些花草樹木.
昆蟲,蛇蟲,樹蟻.
有些有毒的.
有些有害的.
他自己衣不閉體.
可能是跟猩猩一樣.
但是猩猩多毛.
調節溫度.
可能是夏天比較辛苦一點.
但是冬天.
冷一點的時候.
即使在森林裡面也會冷.
我們假設.
他也是有保暖.
泰山也要找一層被遮一下自己.
找些什麼東西.
讓自己晚上的時候.
或者是吹風的時候.
不要太涼.
你會發覺.
總之一個人要生活.
他只要有一個環境.
他也需要有智慧的.
需要知道.
這些是浮沙.
不要踩下去.
我不知道有沒有踩過一次.
但是踩過一次可能也再見了.
他可能見過一隻東西掉下去.
這些不能踩下去.
這些是浮沙.
這些昆蟲.
雖然是很漂亮.
但是不要玩.
去針一下就很糟糕.
這些植物.
走過就好了.
碰一下.
全身都一粒粒的.
就長出來.

$^{161}$敏感的.
Poison ivy.
這些植物.
我試過.
有一次是替教會.
一對老人家夫婦在美國的時候.
回美國教會.
美國人是獨立屋.
很多時候都住獨立屋.
他自己有一個花園.
長了很多東西.
年紀大了.
自己不整理花園.
又不想花錢請人來.
因為都很貴的.
我們幾個學生.
就很合適.
跟他又熟悉.
找個週末幫他清理花園.
但是花園那裡.
我們不懂.
誰知道.
總之就斬斬斬.
剪剪剪.
推車.
但是有poison ivy.
那些纏著.
我都不知道中文是不是.
又不是獨場味.
不是exactly獨場味.
碰到.
我和另外一個同學.
去幫忙的時候.
兩個碰到.
碰一碰.
然後一粒粒就長出來了.
poison ivy很容易認識.
你知道就很容易認識.
三塊葉.
三分的葉.

$^{201}$但是那時候不知道.
所以中招了.
這些都需要智慧.
一種我們說常識.
你跟人相處是有種常識.
常常都要知道.
或者你常常做得多.
你就自然會知道.
滾水.
火不要鬥.
滾水.
煮滾了的東西不要碰.
之類之類.
所有這些東西是用來做什麼.
就是讓我們過一個相對平穩.
安穩.
和越.
和就是.
加減乘除.
和口的和.
越就是越下的越.
就是跟人能夠和平相處.
應該和平的和.
能夠是比較閱絡的生活.
如果你和你隔離.
常常都是嘟口嘟鼻.
常常兩個住在隔離.
兩個互相是頂心杉的話.
你住在那裡都很麻煩.
就是.
一會兒第二天你一出門.
一踩他就弄一塊狗糞.
踩就踩到.
他耍你.
你又耍他.
是不是.
這些東西你不想發生.
為什麼這樣.
令你的生活不能夠生活下去.
所以人類就會發現.

$^{241}$發明.
想出.
經驗的累積.
失敗乃成功之母.
在這些位置.
失敗過一次.
就發覺.
下次學乖了.
就不要這麼做.
慢慢慢慢就成為一大堆的.
我們所謂的智慧語言.
真言.
諺語.
聖經的智慧文學都是屬於類似.
不過.
有一個很不同的地方.
就是我們中國人普通的智慧說話.
基本上沒有一個對上帝的信仰在裡面.
中國人是.
我們智慧說話當中又不是說.
不牽涉及神明.
都會有些東西說天打雷劈.
怎樣怎樣.
都可能會有一些類似牽涉及神明的說話.
但是整體.
我們中國人生存智慧的基礎.
其實跟神明沒有關係.
就真的平時生活裡面的人生經驗.
智慧的累積.
不過猶太人就不同了.
猶太人的智慧文學的基礎.
就是敬畏耶和華.
所以我們都很熟悉.
箴言一開始就說敬畏耶和華.
第一章就已經是.
敬畏耶和華是智慧的開端.
後來第八章又再第九章又重覆.
連傳道書我們今天所讀的都是這樣.
不過他就放在後面.
放到十二章最後那章才說.

$^{281}$哎呀說了這麼多話.
所羅梅王傳道者.
所羅梅王說了這麼多話.
總之就是要我們敬畏耶和華.
要敬畏上帝.
箴言事實上他三十一章裡面.
投了九章的基礎.
奠定智慧說話後面的基礎.
就是要敬畏上帝.
這個就是猶太人.
和所有其他民族.
很不同的地方的智慧語言.
也是我們讀的時候.
必須要意識到的一件事情.
但是所有的東西.
敬畏耶和華都是讓我們.
猶太人因為他們相信上帝.
認識上帝.
都是讓我們能夠過一個和悅的生活.
過一個順境的生活.
你得罪人.
你跟奸人論事是不妥的.
你也知道.
每天生活都很艱難的.
何況你跟上帝不妥.
你的生命都可以想像到.
是很艱難的.
或者你根本不知道明天會發生什麼事.
因為你的生命都是掌控在上帝的手裡.
所以這個前提就變成很重要.
所以這件事我們明白.
我們讀箴言.
學箴言.
在箴言裡面.
所羅門叫人默想.
要思想箴言.
思想這些說話.
智慧人就會默想.
愚昧人就不理會.
都是為了我們能夠跟上帝建立一個正常的關係.

$^{321}$跟人與人之間建立一個正常的關係.
跟大自然.
包括大自然.
所以猶太人心目中的工匠.
手作的工匠.
比如說窯匠,木匠,鐵匠.
這些都是智慧人.
不同類別的智慧人.
因為他們懂得運用大自然.
野金之術.
能夠武來武去.
武了大自然.
製造一些器物,器具出來幫助自己.
在這些過程裡面.
就發現很多問題.
今天我們所讀的一段經文.
就是所羅門自己所注意到的.
其實在傳道書裡面.
智慧人會默想,觀察,吸收經驗的教訓.
他們其中一個很重要的教導.
當然就是我剛才所說的因果關係.
善有善報,惡有惡報.
但是我們自己也知道.
善有善報,惡有惡報.
我不應該問有沒有.
是不是每一次都是現眼報的?.
當然不是.
不會每一次都是現眼報.
你做了一件好事.
突然間就有一個好處.
一聲飛就飛回來了.
會不會呢?.
很難的.
我剛才在地鐵讓位.
讓位給了肚子已經大幅變變.
我想已經八,九個月的孕婦.
我就一直站著.
沒有人給我位子.
沒有好處就繞回來.
我就繼續站著.

$^{361}$等到有人下車讓位.
不是讓位.
下車有空位.
有位子你坐到就坐.
所以行惡也是這樣.
你會發覺惡是有惡報.
但是大家見過惡人是.
這裡正正是剛才所說的.
惡人行惡倒向長壽.
見過沒有?.
見過.
你看電影看過了.
看電影就肯定看過.
很多劇都是這樣.
但是見過沒有?.
見過.
為什麼還沒見過?.
你也覺得起碼.
就算你不認識他.
看新聞.
哇!十惡不赦的.
但是也好.
也很修美.
有些是.
肯定不是十惡不赦.
但是普通人.
就是普通人.
而且也很友善.
但是又會是.
這裡就是未到期而死.
是不是?.
我剛剛出來的時候.
有位校友.
WhatsApp給我.
他說張牧師.
我記得你以前說過張牧師.
在你的教會那裡.
他說張牧師.
他截了WhatsApp的圖給我.
我找來找去也不知道是誰張牧師.

$^{401}$我說張牧師.
我又不是張牧師.
他在跟我說誰呢?.
原來在頂那裡.
WhatsApp截了頂.
張百笠牧師.
張百笠.
不知道大家有沒有印象.
他是民運人士.
是八九民運.
是逃脫了.
輾轉.
真的很輾轉.
叮叮咚咚輾轉.
就過了美國.
當時我在美國.
牧會讀神學.
牧會的時候.
他就在我的教會那裡.
剛剛收到訊息.
我說什麼事呢?.
看來看去.
原來最底層才看到發生什麼事.
就是他33歲的兒子.
剛剛在睡覺夢中的時候.
心肌梗塞.
心臟病褪逝.
白頭人送黑頭人真的很慘.
33歲.
而且他可能自己結婚的話.
還有他自己的家人.
他一家.
那時候還沒有信主.
我第一次認識他的時候.
他可能剛到美國沒多久.
搞著搞著.
就在美國.
那時候是90年代初.
90年代中.
還沒有到95年.

$^{441}$是93,94的時候.
他還沒有信主.
就在我們福音組裡面.
查經組裡面.
參加一些聚會.
我也知道他是民運人士.
是偷走出來的.
很多人幫過他之類.
但是他也很專心.
後來他信了主.
讀神學.
現在在加州做牧師.
他的兒子33歲.
心臟病睡覺的時候就去世了.
所以義人恆二.
我假設他的兒子也信主.
我看他兒子也是信主.
因為他說在天上再見.
我一開始說他會有去世.
但是他會有去世.
為什麼這個校友又發一個截圖給我.
原來他的兒子去世.
信主的人為什麼會經歷這樣的事情呢.
一對牧師夫婦.
為什麼需要經歷白頭人送黑頭人的事情呢.
這些事情在我們觀察世界的時候.
都是困擾著我們的問題.
如果有因果的話.
為什麼好像看不到那個因.
不是不是.
有個因.
或者是見到那個因.
但是看不到那個果.
所以我們中國人很聰明.
善有善報.
惡有惡報.
觀察了一段時間.
發覺不能報.
不報.
再加兩句.

$^{481}$哪兩句呢.
若然未到.
未報.
時辰未到.
很聰明.
中國人很聰明.
時辰未到.
如果到死都未到呢.
所以所羅門還是比較聰明.
義人恆二反哲滅亡.
假設滅亡是死了.
惡人恆惡倒響長壽.
他不會衰收尾.
惡人都要死的.
長壽的意思是倒響長壽.
假設他一直好到最後.
如果死了.
時辰未到.
到了.
每個人都死了.
死了就沒得追了.
還追什麼.
總不能追下一代.
這就變成犯罪.
所以在這些位置.
都是聖經比較聰明.
因為.
所羅門想多一層.
所羅門作為皇帝.
他傳道書一開始就說.
我們已經知道所羅門是智慧人.
他本身就很聰明.
因為大衛是聰明人.
他自己生了一個兒子.
他認為這個兒子應該繼承皇位.
所以所羅門做皇帝.
而且大衛吩咐.
他的兒子所羅門.
我做不到的事你要做.
上帝不讓我建聖殿.

$^{521}$你要建聖殿.
以色列人將上帝的居所建立起來.
不要再到處走.
以前還是會幕.
會到處搬走.
所羅門答應了.
大衛也預備了很多前期的資源.
但是所羅門.
上帝也喜歡他.
正如上帝說.
大衛是合他心意的皇.
所羅門上帝也愛他.
所以上帝在他年輕的時候.
向他兩次顯現.
第一次顯現就是他當皇帝的時候.
登基.
登基沒多久.
第二次就是他完成.
聖殿的建設的工作.
奉獻聖殿的時候.
上帝又第二次.
向他顯現.
第一次向他顯現就是我們很熟悉的故事.
上帝問他.
哎呀!所羅門.
你也知道我疼你.
你做皇帝.
你想要什麼你告訴我.
我什麼都給你.
在異象裡面.
上帝居然這樣跟他說.
你也想上帝這樣跟你說.
不過上帝這樣問你.
你就要很小心.
不要答錯.
一旦答錯.
以前那個不知道是波斯人的寓言還是伊索的寓言.
都是的.
可能是阿拉伯.
一千零一夜.

$^{561}$神是問一個人給三個.
你可以許三個願.
你問我三樣東西.
他問得不好.
結果到頭來什麼都沒有.
所羅門就求上帝.
上帝說你求我什麼我都給你.
求上帝給他有智慧.
能夠好好地管治這個民.
上帝託付他這個子民.
上帝一聽見就很開心.
如此可教.
你本來可以問我什麼我都會給你.
你要是財富,女人什麼的.
你只要出聲.
我已經答應給你一定給你.
你居然這些都不求.
只求能夠好好管治這個子民.
這個答案就很託上帝喜悅.
上帝說好這件事我給你.
給你聰明智慧.
但是你沒有求的那些我都給你.
所以說對了上帝什麼都給你.
說錯了就什麼都沒有.
所以所羅門是一個有智慧的王.
他的智慧是甚至外族周圍附近的外族君主都走來和他建立外交關係.
所羅門的武功都很厲害.
雖然打仗大衛.
總體上大衛是打仗打生打死.
滿手是鮮血.
雖然他也流過無辜人的血.
上帝不讓他見聖殿.
大衛基本上是將巴勒斯坦地.
即是猶大地 迦南地裡面的那些外族平定了.
所以交位給所羅門的時候.
就有一點像是.
不應該用那個做對比.
倒過來.
我本來說文景二帝交給漢武帝的時候.
倒過來就敗光了.

$^{601}$交給所羅門他本來就不需要打太多.
但他仍然要平定一些東西.
所以所羅門的武功 軍事都厲害.
但所羅門最出名的就是他的建築.
建築聖殿.
他自己的殿.
和不同南面和北面的城市.
到今天的考古都好.
有些所羅門的建築還遺留在那裡.
有些城牆的城門口都遺留在那裡.
但所羅門做皇帝就不單止是做皇帝管治子民.
他還將他自己有的智慧寫下來給我們.
真言不是全部是所羅門寫的.
大部分是.
傳道書應該全部都是傳道者.
是皇帝.
我們後世猶太人都認為是所羅門寫的.
全部都是他寫的.
他作為皇帝一開始就立意要觀察生命裡面的事情.
他看著看著.
他就發覺.
生命裡面的事情有三件事很困擾他.
其中一件就是我們剛才說過的.
善有善報.
我報讓未報.
時辰未到.
死.
人都死了.
還未報.
那怎麼辦.
就是報應這件事.
這件事困擾他.
死亡.
死亡好像將一切的東西都拉平了.
扯平了.
扯.
如果是這樣.
惡人就是這樣.
義人就是這樣.
惡人會不會是黑黑白白的比較有好處呢.

$^{641}$一會兒我會處理這件事.
所以其實所羅門有答案的.
這就是後面那句的答案.
但後面那句答案都困擾我們.
但無論如何.
這是一件事.
第二件事.
他一早已經觀察到.
就是做來做去都是沒用的.
半斤八兩.
出了半斤力.
下面那句是什麼.
怎麼會拿回八兩這麼理想.
後面那兩個字就不要說了.
唱歌是雙聲報喜的歌.
許氏兄弟.
這些年輕人都不知道.
可能不知道我說什麼.
是一些歌詞.
出了半斤力.
怎麼會拿回八兩這麼理想呢.
你會發覺.
都是某程度上和因果有關係.
但是疏離一點.
因為這是人際.
你本身做了這麼多事.
本身就應該有差不多的成果.
你預期的成果.
就是說有時候做了一些事.
是沒什麼用.
我們人說是白做.
我不知道你怕什麼.
我有一件事很怕的就是白做.
做了一大輪.
什麼用都沒有.
其實牧師都很怕白做.
牧師會怎麼白做呢.
一年講52堂道.
講完之後那些會眾一模一樣的.
白講.

$^{681}$沒有獎項的.
我不知道牧師.
有沒有這樣的感覺.
不是說大家.
每個牧師都有這樣的感覺.
父母親都會有.
公書教學.
壞人結果.
白樣子.
家裡都會有的.
他很困擾.
但是你做了一些事.
為什麼沒什麼效用呢.
又或者.
做了一些事.
結果不一致.
一事急.
做同樣的事.
結果得來的.
每次都不一樣.
甚至是令你遭遇損失.
所以他考察當中.
在這句.
我們剛才所讀的第七章.
十五至十八字.
主要.
其實三樣東西都包含在裡面.
惡人行惡反至滅亡.
惡人行惡.
倒享長壽.
最後死亡又拉扯在一起.
所以他要解決.
聰明人智慧人.
有一個人生的問題.
一定要解決這個問題.
怎麼解決呢.
就是第十六節.
第十六節.
我們看經文的第十六節.
第十六節是什麼呢.

$^{721}$這是他的.
十六節又是.
十六十七節.
我們先看第十六節.
他得出一個結論.
既然如此.
就不要行義過分.
也不要自稱智慧.
何必自取滅亡呢.
自取敗亡呢.
他說.
惡人行義反至滅亡.
那不要行更多義了.
是不是這麼說.
不過他不是說不要行更多義.
而是不要過分行義.
行義也要過分嗎.
行義可不可以過分.
或者過分行義又如何.
做好事會不會.
做得過分呢.
做得過分好事又如何.
有什麼問題呢.
想想中國人有沒有類似的觀察.
或者說法呢.
你會發覺.
也有的.
我們說.
有風就不要.
洗盡你.
這句說話是什麼意思呢.
有風不要洗盡你.
就是說道理在你那邊.
有風嘛.
你穿.
但是你穿又如何.
那去到盡嗎.
人家你說你不穿我穿.
那就去到盡.
將對方按到沒有彎.

$^{761}$因為我穿嘛.
你不穿嘛.
是不是.
做人留一線.
日後好相見.
做人留一線背後的假設.
都是我穿的.
我可以.
你欠我錢.
我追追追追到你天腳底.
我追死你.
這個也是這樣的意思.
又或者是你得罪我.
我現在壯大了.
我就將你滅族滅了.
但是日後.
做人留一線意思就是說.
不要做到那麼絕.
就算你有這樣的能力.
又或者道理上你是穿的.
你有意義.
這個是意人嘛.
你有意義也好.
都要留一點空間.
因為日後好相見.
還有後來的.
所以中文也有類似的說法.
即使我們做事是對的.
都不可以去到那麼盡.
香港人喜不喜歡排隊.
你喜不喜歡排隊.
你當然不喜歡啦.
但是你會不會排隊呢.
我希望你會排隊.
大家排隊的時候.
我希望你也會排隊.
香港是出了名的.
有很多遊客來香港.
尤其是內地同胞.
早期的.

$^{801}$現在他們都習慣了.
早期是九零年代.
尤其是落到來的時候.
他們會笑的.
香港人在排隊.
早期回過內地.
八零年代九零年代.
二千年之後.
我比較頻密一點.
你回內地想著.
飛機場也好.
人家宣佈了.
五分鐘內關閘口.
你還站在那裡.
一勇如上.
現在我們的文化.
在這個位置上也好了很多.
所以他們見怪不怪.
來到香港人排隊.
排隊是不是美德.
中文有句說話.
香港人執書行頭.
慘過敗家.
這些才是智慧.
B仔.
你記住.
不要一輩子被人叫B仔.
叫A仔.
執書行頭慘過敗家.
有沒有A在.
但我們又要排隊.
很奇怪.
有矛盾.
民間智慧.
有時真的會有矛盾.
所羅門正正就是.
見識到矛盾.
惡人行惡道.
是現實中的矛盾.
我們要排隊.

$^{841}$先到先得.
保障自己利益.
不想被人搶了.
被人逼走.
所以也是保障自己利益的排隊.
你喜歡就早點過去.
有機會成本.
就早點去.
但不想大家一窩蜂.
搞到出事.
但是.
又有執書行頭慘過敗家.
排隊.
排隊是不是很絕對.
是不是絕對.
一定要排隊.
又不是.
因為還有兩個字.
還有禮.
陽.
在排隊當中.
你發覺排隊這條規矩.
或者執書行頭慘過敗家.
這條規矩.
是有例外的.
有人有需要的時候.
你會讓他.
是不是.
我們車那裡.
紅色位.
已經很清楚.
參與拐杖.
剛才我這樣的大肚子.
或者老人家.
我現在開始被人讓位.
通常我執得這麼好.
就沒有人讓位給我.
但我穿得殘破一點.
就有人讓位給我.
你想有人讓位給你.

$^{881}$穿得殘破一點就行.
你會讓的.
是不是.
所以.
你會發覺有些事都是有例外的.
慢慢我們開始.
明白.
第十五節.
你行義過分.
不要自請自為.
何必自吹拜網.
你什麼都不讓.
排隊.
道善德按次序.
都很金科玉律.
但你可以去到盡.
我理得你.
坐輪椅.
拐杖.
排隊.
出什麼事.
未必一定是個人出事.
如果其他人肯讓.
可能會說你不要這麼.
將心比己.
但出事其實是整個社會.
一個不懂禮讓的社會.
一定出事.
我們是法家.
不懂禮讓.
一定出事.
整個社會.
你會發覺先到先得.
本來是好的.
亦是令社會能夠運作.
但你沒有例外.
社會都會沉淪.
都會出事.
都會滅亡.
自取滅亡.

$^{921}$但這是說社會.
你會說不是.
張文開說的是.
一個人.
我又舉個例子.
我嘗試去理解.
一直都困擾我.
直到我找到一個例子.
那個例子是跟我現在做的有關係.
突然間醒悟了.
我講到.
是否是異行.
是否行義.
講到宣講上帝話語是否一件好事.
好事是上帝.
稱之為有異的行為.
是宣講上帝的話語.
現在是11點45分.
你不要走.
我講到12點.
1點.
2點.
2點你才走.
剛剛區先生告訴我.
在八達利亞有周年慶典.
你今天做主席.
不是的.
2點你不要走.
我在講宣講上帝話語.
你對上帝話語不恭不敬.
真的.
我可不可以這樣說.
你會不會聽我這支笛.
去到2點你猜還有沒有人坐在那裡.
應該沒有.
曾牧師也要走.
為什麼出問題.
為什麼.
很明顯是過份.
為什麼過份.

$^{961}$在哪個位置過份.
講得太久過份.
太久就有一點主觀.
我不覺得過份.
可能有一兩個都覺得過份.
所以一定有客觀的原因是過份.
不是單純主觀.
長短的問題.
多一秒鐘算不算過份.
一定有客觀的因素.
我在想.
我也要小心一點.
因為我也不知道.
一次是叫人講到.
下次過來講到.
我們控制時間.
好一點.
不然不找你過來.
有一次婚禮.
被人圈我.
婚禮我講了9個字.
現在婚禮最短的歷史.
我去過那麼多婚禮.
最短的那個是7分鐘.
講完了.
9個字真的過份.
因為我關心少了我的學生.
其他同學做幫手.
姐妹兄弟.
已經圈我了.
我還沒講完.
問題出在哪裡.
問題就是.
大家都會承受.
甚至大家都會公認.
我做這件事是義行.
是好事.
但如果我一直講.
還沒講完大家都不能走.
我中間有個假設.

$^{1001}$假設是.
這段時間裡.
只有我做的事才是行義.
其他人如果要去買菜.
要去煮飯.
要去做其他事.
就無意可言.
但其實無意可言.
我的義行是.
trump(取消)所有的義行.
總之我trump(取消).
出黃牌港島.
全部都要坐下來聽我說完.
這個就是問題.
這個就是行義過份.
客觀的問題.
就是你以為.
你做的是好事.
是義行.
所有的義都在你那裡.
別人就不會再有義了.
兩夫妻吵架.
吵架通常都有人.
有對有錯.
有時沒對沒錯都會叫起來.
不知做什麼.
但是我們都吵過架.
你吵架的時候.
你都覺得自己是對的.
因為你覺得.
對方都對的時候.
就叫不出來.
你都對的就完了.
問題正正就是這裡.
你吵架的時候.
總是覺得你對.
所有的問題.
出在所有的.
行義過份.
你可能是對的.

$^{1041}$我不會否認你的錯.
但你能夠對所有.
假設對方一定錯.
對方是不著的.
一不著就沒一點著.
就起在你那裡.
你對所有就容易去到盡.
因為你覺得你對所有.
如果你知道你自己不是對所有.
一般來說.
除非你火遮眼.
一般來說你不會去到太盡.
你都會留一個空位.
我有一次.
我跟大家confess.
我已經認過.
找不到對方.
但我記得之前都和大家分享過.
因為我去到盡.
一個很典型的.
對我來說是警告自己的.
典型的做法.
我去長洲.
長洲山頂是我的地頭.
鶴嶼附近也是我的地頭.
有一次應該是周末.
也是周末.
因為很多遊人.
現在很多人去長洲帶著小狗.
現在連快船都可以帶寵物上快船.
以前快船不能帶寵物.
現在快船都可以.
很多人都帶小狗去.
但小狗去.
你也知道周圍會有.
留下黃金.
我剛從宿舍.
走上來.
剛好有對年輕人.
他們的小狗.

$^{1081}$剛好門口旁邊.
一戳金就走開.
我就立刻喝住.
他們.
有沒有搞錯.
人做狗主你做狗主.
有一點公德心.
我沒有說錯.
但說話都很難聽.
你很小就不要在自己門口.
不要在別人門口.
或者你小了沒問題.
因為小狗你控制不到.
男士篇.
沒得男.
你怎樣.
你出來永遠都要預備好.
要收拾.
住長洲和不住長洲.
很多時候都會分到.
長洲人.
一定帶著水起身.
一定有膠袋.
但那兩位年輕人.
一男一女打橫打點.
想走的時候我叫住他們.
叫住他們又沒有反應.
我開始把火.
火遮眼.
有其他途人.
都開始斜路上.
有人在看.
你給我收拾.
你怎樣給我收拾.
沒東西收拾.
那裡有樹葉.
我逼他.
結果逼.
男士可以不理我.
坐在那裡就走.

$^{1121}$我都下去跟他打架.
但我又逼得太盡.
那個男士又.
可能覺得自己有誤著.
結果就.
把樹葉撿了.
弄得手髒了.
其實在那裡我已經開始覺得.
自己有點不妥.
但完了整件事之後.
我就覺得自己很不妥.
他們沒有知.
沒有業.
他們沒有預備.
你罵他們一輪.
他們已經領了教訓.
有人在看著被罵.
已經領了教訓.
下一次他們應該.
不會走同一條路.
撞到這個惡野.
我懷疑他們.
應該會比較聰明.
我應該放手就算了.
我逼他.
去撿樹葉.
我其實那時候.
應該幫他.
隔壁的門就是家.
我幫他撿.
幫他解決了這個問題.
就搞定了.
我就過份逼他撿.
撿就拿上去.
斜路很短一段.
二三十步有個垃圾桶.
有個狗糞池.
所以氣就在這裡.
狗要拉你沒辦法.
所以我逼得太盡.

$^{1161}$我覺得到現在反省.
一直都是我去得太盡.
沒必要去到這麼盡.
真的沒這個必要.
我又不是沒撿過垃圾.
撿過了.
有次男生宿舍.
不知道為什麼男生宿舍廁所地上有垃圾.
我那時候做社監.
叔母就說張sir有奇觀.
什麼奇觀.
叫我來廁所幹什麼.
你們自己管理.
我去廁所看.
有兩塊垃圾在地上.
很奇怪.
男生宿舍.
怎麼會有兩塊垃圾在地上.
沒有狗.
誰拉的.
我就撿.
叔母不撿.
我撿沒問題.
可以做的.
為什麼不能做.
但那時候氣上心頭.
就逼他下牆角.
我覺得這樣做是過份.
現在回想就過份.
我沒滅亡就已經很好.
因為對方比我年輕.
被人打一拳.
你以為很好玩嗎.
打一拳他還是不撿.
所以大家知道.
這是智慧.
不要行義過份.
行義應該的.
但行義過份的意思.
不是說對的事不要做.

$^{1201}$而是不要以為你對的事.
每次以為你對的事.
你就開始過份.
沒什麼幾可會對的事.
所以大家在這裡.
你還說下一句.
不要行惡過份.
即是可以行惡.
不要行惡過份.
和可不可以行惡.
是兩回事.
但都這麼拼.
他也看到第二層面的事情.
因為惡人是反倒美到期.
惡人是反倒美到期而死.
不要行惡過份.
如何理解.
不要行惡過份.
又不要為人愚昧.
何必美到期而死呢.
兩家結果都是不好的.
所以.
所羅門承認.
有些惡人真的倒向長壽.
他發覺.
他就說.
如果惡人倒向長壽.
他就可以行惡.
在所羅門的價值取向裡.
在我們價值取向裡.
當然不可以行惡.
那你有沒有行惡呢.
有沒有.
我都不想叫你舉手了.
你舉手就是惡事.
你講謊.
我們都行過惡的.
大大小小.
我們盡量不做.
尤其是在主面前悔改.

$^{1241}$我們學習主的樣式.
但我們都會的.
再加上.
保羅講得很清楚.
世上有沒有義人.
沒有義人.
第三章.
我們信主之後.
稱義.
得稱為義.
不代表我們突然變成.
不懂行惡的義人.
不是.
我們信主之後.
都可能會行惡.
我那次做得太過份.
已經是惡了.
不單做人做得惡.
其實那件事都不是好事.
所以在這些位置.
我們就要承認.
我們都是有罪的人.
而且.
在一個不完全的世界裡.
往往.
所羅門說不要行惡過份.
我都在想.
但你都要想.
仍然想這些例子.
行惡過份.
其實在人生當中.
往往會面對一些兩難的事情.
什麼叫兩難.
即是你手臂又是肉.
手掌又是肉.
手背.
都是肉.
兩面你都很難做.
但你一定要做.
你不做不行.

$^{1281}$但兩邊都很難.
兩邊都不討好.
這些事.
你遭遇過嗎.
可能.
我都希望你不會遭遇.
但人生.
生活得越來越長.
你會發覺.
真的會越來越碰到兩難的事.
我舉個例子.
大家看過鐵達尼號嗎.
鐵達尼號.
晚上行船撞冰山.
很窄.
很多人扔救生艇.
你又黑漆漆.
電又沒了.
居然找到救生艇.
你上了救生艇.
但下去水浮的時候.
聽到遠處叫救命.
黑漆漆.
但燈又未熄.
鐵達尼號下面的路都斷了.
還未熄.
沉下去才熄.
見到老婆和媽媽.
捧著木頭浮.
浮下.
就離開了.
居然是浮下.
救生艇只剩下一個位置.
只剩下一個位置.
這是典型的道德倫理問題.
你給個位置.
救你老婆.
還是救你媽媽.
這叫兩難.
很難.

$^{1321}$做一個決定.
都是一件惡事.
因為另一個會死.
好.
這兩個都是我親愛的女人.
我就跳下去.
給多個位置.
救了兩個女人.
都可以.
但這都是一件惡事.
他們死了兒子或老公.
都是一件惡事.
都不是一件好事.
三難.
其實整件事是四難.
你看看.
我和我老婆.
一腳踩下水.
我們三個在水上.
旁邊的女人也踩下水.
是一件惡事.
都是惡事.
四難.
你有四個選擇.
你救一個不救另一個.
你可以自己跳下水.
等他們兩個上來.
你可以踢走旁邊的女人.
等他們兩個上來.
四件都是惡事.
是很大的事.
你從這個角度看.
人格上是偉大.
整件事都是惡事.
就是有人死.
你就是不想人死.
就是這麼簡單.
你沒有辦法.
你做醫護.
資源就是這麼多.

$^{1361}$所以我們現在.
在急症室.
以前急症室真的先到先得.
你排隊的就是排隊.
先到先得.
在急症室的層面可能是一件壞事.
因為人家重症.
要排到後面.
那怎麼辦.
大家去急症室就知道.
如果有救傷車來.
一定橫進.
但你坐在那裡的也有急症.
受傷.
你見過人全身血嗎.
但我小時候.
先到先得.
拿了籌就慢慢等.
總體上是這樣.
有些病徵不是太明顯.
但其實是很急症.
現在就沒有這件事.
現在第一個.
他要審視一下你.
那些叫做分流.
要分流.
你打球要分流.
你扔三個小時.
有很久都還不到你.
本來不應該來急症.
看他跌打就行.
快要濫用醫院資源.
因為他們每天都要衡量.
就是這麼多資源.
就是這麼多時間.
就是這麼多醫生.
就是這麼多藥.
你給誰先給誰後.
你給這個先那個就要等.
如果他做錯決定.

$^{1401}$就應該是倒轉的.
另一個先的.
那就有人出事.
很多事是沒有辦法避免的.
行惡在這個世界上.
經常是沒有辦法避免.
因為這個世界有很多兩難三難四難的事情.
只不過你不能行惡.
你盡量不行惡.
而沒有辦法避免的情況.
智慧人所羅門的勸誡.
就是不要行惡過分.
不要為人愚昧.
你聰明一點.
想一點.
要義惡取其.
輕.
你只能選擇你自己的智慧知識範圍以內.
你認為輕一點的惡的事.
因為你沒有選擇.
有選擇你就不要選擇惡.
但你沒有選擇的時候.
你一定只能選擇輕的那個.
你選擇重的那個.
就是行惡過分.
就未到期而死.
就會出事.
所以如果你整體這樣看.
你就發覺所羅門真的很有智慧.
真的很聰明.
第十八節.
大家的版本是和修版.
我仍然讀和修版.
很少和修版翻譯得那麼好.
我覺得這個和修版翻譯得不那麼好.
和修版就是你持守.
那個也沒有鬆手.
才能敬畏上帝的人.
這一切都能兼得.
上面其實也沒有說你得著什麼.

$^{1441}$他叫你不要這樣不要那樣.
免得你要這樣.
所以我覺得和修版的翻譯不好.
反而回到和合本.
這裡就容易一點.
和合本怎麼說呢.
前面三句一樣.
你持守這個唯美那個也不要鬆手.
因為敬畏神的人.
必從這兩樣出來.
意思就是說必然會經歷這兩樣東西.
行義過分.
有可能你行義過分.
行惡過分你可能會行惡過分.
你會遇到這些情況.
就不要行惡過分.
敬畏上帝的人.
一定要在這些事情裡面.
出來作對的決定.
這樣走出來.
所以我覺得反而和合本在這個案件.
是翻譯得更好一點.
所以我們要領受智慧.
我們在上帝面前.
認真讀.
是所羅門的.
雖然他雖修美.
但是我們不能夠因人廢言.
他雖修美.
他沒有服藥.
沒有領受教訓.
很可惜.
扮偶像老婆太多.
王祖皇帝.
外族女子太多.
但是他說的話.
是對的.
我們仍然要聽.
記載在聖經裡面.
仍然是教導我們.

$^{1481}$依照我們的說話.
讓我們在這個世界上生活.
面對上帝.
都能夠好好地過.
到最後死亡都不是一個問題.
我們面對上帝的時候.
我們都能夠在他面前.
憑著主耶穌基督.
憑著我們自己過往的經歷.
在他面前.
是站立得穩.
一起低頭祈禱.
因為感恩.
因為你在耶穌基督裡面.
拯救了我們.
你稱我們為義.
看在基督的份上.
看在我們悔改.
在承認基督為主的份上.
你稱我們為義.
但我們都知道.
我們一直生活在這個世界上的話.
一直仍然是有各式的試探.
有難.
有困難的決定.
是很難做的.
我們做人本來就是這樣.
求天父讓我有聰明智慧.
在各樣的事情上.
在面對艱難的決定的時候.
你為我們開一條出路.
你給我們智慧.
但又沒有自稱智慧.
以為自己全對.
卻也沒有為人愚昧.
以致做錯決定.
求天父幫助我們在每日的生活當中.
看見你自己人靈的手.
帶領我們每日的生活.
奉主耶穌基督的名求.

$^{1521}$謝謝大家收看 再見.
\newpage



\section{馬太福音 10:34-39}
\label{sec:7vPzFHXDGtY}
\textbf{2025.11.16《與神同行,超越困難》:馬太福音10\_34-39 (和修版) 講員 :許淑芬牧師}
\newline
\newline
連結: \href{https://youtube.com/watch?v=7vPzFHXDGtY}{\texttt{https://youtube.com/watch?v=7vPzFHXDGtY}} ~~~~ 語音日期: 2025-11-19
\newline
\newline
\hyperref[sec:7BCvZD9NU64]{\small{< < < PREV SERMON < < <}}
~
\hyperref[sec:index_chronic]{\small{[返順時目]}}
~
\hyperref[sec:index_scriptual]{\small{[返順卷目]}}
~
\hyperref[sec:11YyQ2efYE8]{\small{> > > NEXT SERMON > > >}}
\newline
\newline
馬太福音 10:34-39
\newline
\begin{longtable}{cl}
\hline
\hline
章節 & 經文 (和合本修訂版)\\
\hline
10:34 & \begin{tabularx}{0.7\textwidth}{X} 「你們不要以為我來是帶給地上和平，我來並不是帶來和平，而是刀劍。 \end{tabularx} \\ \\ \relax
10:35 & \begin{tabularx}{0.7\textwidth}{X} 因為我來是要叫『人與父親對立，女兒與母親對立，媳婦與婆婆對立。 \end{tabularx} \\ \\ \relax
10:36 & \begin{tabularx}{0.7\textwidth}{X} 人的仇敵就是自己家裡的人。』 \end{tabularx} \\ \\ \relax
10:37 & \begin{tabularx}{0.7\textwidth}{X} 愛父母勝過愛我的，不配作我的門徒；愛兒女勝過愛我的，不配作我的門徒。 \end{tabularx} \\ \\ \relax
10:38 & \begin{tabularx}{0.7\textwidth}{X} 不背自己的十字架跟從我的，不配作我的門徒。 \end{tabularx} \\ \\ \relax
10:39 & \begin{tabularx}{0.7\textwidth}{X} 得著性命的，要喪失性命；為我喪失性命的，要得著性命。」 \end{tabularx} \\ \\
[1ex]
\hline
\hline
\end{longtable}
$^{1}$各位弟兄姊妹平安.
我想問問大家.
在信主的年日.
可能每個人都不同.
有沒有誰覺得上了耶穌之後.
一帆風順的.
沒有任何困難的.
舉手.
不敢舉手還是真的沒有.
但可能不舉手的.
心裡有數.
上了耶穌之後好像沒有什麼困難.
也有很多.
我有時候會想.
如果信耶穌.
一信耶穌,每件事都順利.
多好.
讓我們見證.
見證也好一點.
或者告訴別人.
我們後面有些力量.
信耶穌之後.
一帆風順的.
你相信吧.
但偏偏在我們信仰裡不是這樣.
不要說信仰的路.
就算人生的路.
也沒有什麼.
順境的情況.
但我們就說.
為什麼我們信耶穌之後.
仍然有這麼多困難.
我們和一個不信的人.
有什麼分別呢?.
分別就大了.
但我們今天想透過.
剛才大家讀的經文.
主席讀了一段經文.
你們覺得刺不刺激?.
信耶穌這麼大的問題.

$^{41}$如果我用一個字更刺激.
就叫做革命.
信耶穌就像革命.
完全改變.
為什麼會這樣呢?.
如果我們去看.
期望人生的路當然好.
信耶穌的路當然好.
當然是最好的.
但原來在這段經文.
信耶穌不單不是一帆風順.
還要有很多的衝擊.
或者我們說.
是一個革命性的改變.
為什麼耶穌用這麼強烈.
為什麼用這麼強烈的字眼呢?.
我們好像有一個很大的改變.
如果我們有人.
你會說.
改革和革命.
有些不同的地方.
改革可能是外在的東西.
去改變.
改革就是將原有的東西.
或者系統,體系.
改一改.
但如果革命.
其實是說由內而外.
由內面,由裡面.
完全有一個.
根深格深的情況.
甚至推翻舊的東西.
舊的體系.
如果不好的話,就推翻.
如果要成立一個新的體系.
其實我們信仰.
帶給我們一個革命性的改變.
應該由我們由內而外.
一個完全的.
根深改變.

$^{81}$為什麼福音.
是一個革命的影響呢?.
我們最近去談.
我們有一個報道會.
展秋也在.
我們報道會選主題.
其實想主題很重要.
想了之後就是宣傳的口號.
想了很多的內容.
然後我們就定了一個.
方向,叫做「重生性原生」.
不知道大家.
原生家庭對大家是不是很大影響呢?.
很多人都說.
現在我們經常都賴原生家庭.
為什麼我今天這樣呢?.
跟老婆,老公說的時候.
可能是因為我原生家庭.
可能都成為一個藉口.
但是,如果我們在.
福音的生命的改變裡面.
我們的重生.
是不是勝過.
我們原生家庭帶給我們.
一些破壞的影響.
反而重生是帶來一個.
完全的醫治,將我們生命改變.
如果是這樣,恭喜你.
你的信仰才是真的.
我們的信仰,基督的復活大能.
就是將我們由內.
由裡而外.
去完全改變.
將我們的價值體系.
將我們的生命.
裡面的內容全部改變.
所以馬太福音.
可以說是整卷的.
馬太福音.
有五篇的講章.

$^{121}$如果大家翻著學.
或者查經的時候都會知道.
五篇的大講章就連繫成.
馬太福音的一卷.
很重要的書卷,方生性的書卷.
而耶穌生平的.
其中一卷的書.
我們去到這一段.
講這篇的內容.
其實是第二篇的講章.
由第十章第一到.
十一章的第一節.
上文講什麼呢?.
當然我沒有時間和大家講上文.
上文講到差遣十二個門徒.
然後耶穌給福音的權柄他們.
但是已經預告.
門徒面對傳福音的時候.
一定會有面對.
一些迫不及待.
遭遇害的情況.
但是鼓勵他們.
不要怕,要怕的就怕誰呢?.
就要怕永生的神.
那些垃圾身體的.
不用怕他.
意思是魔鬼都會害我們身體.
但是我們要怕.
那個敬畏的.
就會將我們身體和靈魂.
掉進陰間.
或者地獄裡面.
這個才要怕.
因為我們神是擁有.
生命的主權.
就是因為這樣.
我們需要更加的敬畏神.
所以最後在這一段的.
內容.
四十到四十二節就講到.

$^{161}$門徒有一個很崇高的地位.
不過我這裡都不講.
我特別要集中講.
第三十四到第三十九節.
就是告訴我們.
當我們信耶穌之後.
耶穌就吩咐門徒.
我們怎樣在種種的.
迫害艱難裡面.
我們好好預備自己.
這是一個心智的準備.
我不是說大家現在馬上走出去.
就會面對很多的苦難.
不是的,但是我們怎樣在心智上.
做好一個準備.
讓我們的信仰連繫於.
我們的生活.
當面對衝擊來臨的時候.
我們能夠配得上.
成為主的門徒.
而我們這個心智的準備.
無論是心智還是實際的.
面對困難.
我們都能夠面對.
今天種種的難處.
這就是信仰.
所以我們要.
我喜歡這樣.
跟旁邊的人說.
小聲提醒他.
要準備好.
福音是你生命的革命.
小聲一點.
不要太大聲.
告訴對方,這是一個革命.
這段經文的主題.
是要告訴我們.
我們的信仰.
如果與神同行.
我們就能夠超越困難.

$^{201}$這不是告訴我們.
我們沒有困難.
也不是告訴我們我們得勝困難.
困難我們得勝不了.
我們只能超越它.
或是越過它.
或是跟它打招呼.
但我們怎樣才能夠.
面對這一場.
因為福音的緣故.
帶來的革命.
我用三個方向.
跟大家思考.
究竟福音是怎樣的一回事.
怎樣去革命.
有三個很重要的.
其實這是一個本質的問題.
本質的革命.
或者是堅忍的革命.
或者是信服的革命.
所以在這三方面.
我們先要打開.
這個問題來看看.
為什麼今天我們要藉著福音.
帶來一個改變.
第一樣我們要知道.
福音本來.
就是要讓我們知道.
福音是一個本質的改變.
福音就是對我們.
生命本質的改變.
所以在這段經文第34節.
主耶穌一開門見山.
就已經說了.
你不要想錯了.
你不要以為我是帶來給地上太平.
而是告訴大家.
我叫地上動刀兵.
大家看到好像很害怕.
千萬不要拿著刀兵.

$^{241}$我們要改革.
耶穌這句話.
很重要的意思是什麼.
或者我們很多時候去想.
可能有一個困擾.
有一個衝突.
耶穌明明是和平之君.
耶穌明明是要我們與神和好.
與人和諧.
為什麼會令我們動刀兵呢.
所以我們不要看錯.
因為耶穌就是和平之君.
但這個和平之君.
來到這個世界上.
不是每個人都接納的.
既然不是每個人都接納.
我們就跟著和平之君.
可能都要動刀兵.
這個刀兵我們要知道.
不是靠著我們自己的力量.
我們也不是靠著自己.
可以面對.
我們需要靠著和平之君.
怎樣讓我們的生命.
在面對的時候有力量.
究竟我們正在面對.
一個怎樣的世界呢.
其實我們說一個很簡單的事.
就是我們大家都知道.
我們今天正在面對一個.
是謊言的世界.
不只是今天詐騙的事件.
其實可能今天.
大家一人一句謊言都不靈.
我們有些專家.
要研究.
每十句說話.
有七句是謊言.
但與外教徒不是這樣.
如果我們不過一個.

$^{281}$說謊的生活.
其實我和大家姐妹.
說笑也是說真的.
有時明明餓得不得了.
有人問你吃飯了嗎.
你為了方便.
答他方便.
但其實是肚餓.
或者我們聽到.
現在WhatsApp很方便.
有人提出問題.
你馬上舉起祈禱手.
我要勸告教徒.
不要說謊.
當你舉起祈禱手的時候.
真的要祈禱.
這個我也要操練自己.
當我這樣舉起祈禱手的時候.
即時為他祈禱.
但即刻為他祈禱.
至少我也要對得起自己的信仰.
我今天所說的.
就是我活響.
我就是為他祈禱.
所以如果大家看不到我這樣.
即是我還未祈禱.
但這樣.
我們知道有人為我們祈禱.
這也是我們一個很重要的.
基督徒的生命信仰.
和我們生活是一致的情況.
所以我們福音是要改變我們.
讓我們有一個新的生命.
要讓我們.
是本質的完全的改變.
就好像光進入這個世界的時候.
讓我們生命改變.
所以耶穌說了一句很簡單的說話.
祂說我們敬拜天父.
我們敬拜天父.

$^{321}$只有一個關口.
一個關鍵是甚麼呢?.
我們用心靈誠實.
然後大家都想人見人愛.
車見車載.
就是要我們有慈愛.
和誠實的生活.
如果大家用聖經一些很重要的.
關鍵字.
當我們有一個很重要的.
鎖匙告訴我們.
究竟我們今天怎樣做基督徒的時候.
其實聖經是說得非常清楚.
不過.
有時候我們自己都很迷糊.
究竟上帝是怎樣說的呢?.
其實為甚麼這麼迷糊呢?.
因為我們裡面都很迷糊.
因為我們裡面有私心.
有私心的時候.
我們就很迷糊.
正如我們.
曾經有個弟兄和太太說.
兩個都是基督徒.
我都會祈禱.
看看上帝是不是叫我們離婚.
這些應該獎他兩巴掌.
用不著祈禱.
聖經告訴我們是不行的.
但為甚麼這麼迷糊呢?.
因為裡面沒有.
究竟是怎樣呢?.
還要在那裡撕花瓣.
或者在那裡抽籤.
或是公定志.
聖經很清楚說得.
但如果我們沒有心靈誠實.
去面對神的話語.
我們根本沒有辦法去明白.
神的心意是怎樣.

$^{361}$所以耶穌的光.
神是愛.
當光進入我們生命.
其實為甚麼.
今天我們要面對這麼大的衝擊.
就是因為這個世界.
不跟你說真話.
我們基督徒就要.
立志說真話.
我常常都要問自己.
今天你做的事是不是真的.
不過我又相信.
你沒有吃這麼多年飯.
你都很難說真話.
如果我做基督徒這麼久.
要操練說真話.
是真的需要花功夫.
你需要花時間.
你怎樣過誠實的生活.
信仰跟你生活是言行合一的生活.
是一個操練的過程.
所以為甚麼有些調查.
你告訴我.
恭喜各位.
有些信主中年的.
越信主越順.
越信主越覺得信仰很容易去遵守.
但最難的是哪些呢?.
年輕人.
年輕人這麼容易.
說兩句話.
這麼容易衝擊就離開.
因為他們面對實在太大的困難.
因為大家已經有經驗.
你吃鹽多過他吃米.
但我又看見.
當我們經歷了這麼多.
事實上.
我們的位置高了.
所以當我們見到年輕人的時候.

$^{401}$我們不要常常戴著眼鏡.
一孩不如一孩.
我們需要多些用心力.
去幫助他們如何面對.
就像光要進入一個黑暗的房間.
這當然不是我們近日.
包裝的那些.
因為我們近日團契.
有一半的地方是裝修.
裝修其實要去收拾東西.
都收拾得很難.
我們有八十多張東西.
那些師傅說.
你這麼多.
在說我.
有些聲音微言.
你十六噸.
我只有十二噸車.
但我都要車滿.
車都要車滿.
我們本來辦公室.
其實沒有什麼.
禮堂只有椅子.
但原來一到雜物的時候.
是很多東西要拿出來.
所以我都有些悔改.
我們今天.
我去扔的東西.
我自己房間的東西.
在我門口都有一大堆.
那些是什麼呢.
特別多.
我都會想我講章需要印嗎.
有個同工提我.
有時電腦不聽話.
你突然間沒了怎麼辦.
都要印出來.
但同工的會議記錄.
導師的會議記錄.
我們那些.

$^{441}$董事會的會議記錄.
多到離譜.
所以那天.
我就去跟董事會說.
我們今天不派紙.
董事會都不給紙.
大家用開電腦去看.
省紙.
我都覺得需要.
我們砍那麼多樹.
有些我都要悔改的地方.
但這些都是小事.
小事的意思是你破壞我.
我破壞你.
但有些東西很重要.
我們的心靈裡面.
會不會有很多雜物.
心靈裡面要堆砌什麼.
會不會驕傲.
或者我們裡面的仇恨.
我們用假的愛心.
去面對仇恨.
剛才詩歌有一句.
我覺得很感動.
這個世界充滿冷漠.
這個世界愛的相反.
不是恨.
恨都說有感情.
如果有些夫婦吵架.
大家唇槍舌劍.
都恭喜她.
有些激情.
但你不出聲.
才難搞.
婚姻.
冷漠才難搞.
所以我們心靈裡.
必須要大裝修.
拆開門窗.
看看裡面有沒有毒素.

$^{481}$讓我們的身心.
活得健康.
讓我們的生命在神面前誠實.
在人面前.
都能夠坦誠相對.
而這樣的生命.
讓我們有更多空間.
騰出來.
讓神的話語裝載進去.
否則.
我們今天聽也聽不到.
看也看不到.
因為裡面太多毒素.
所以我們求主幫助大家.
又和旁邊的小聲說.
小聲說.
快點清理愛心的垃圾.
屬靈的垃圾一定要清理.
讓我們這個革命.
真正重生的生命.
讓我們的生命.
發生果效.
我知道大家可能也聽過新聞.
二十多年前.
911事件.
911事件很深刻.
我兒子剛出生不久.
我剛和女兒在旁邊.
突然看到發生了什麼事.
那個.
衝擊很大.
美國那兩座大樓.
塌了.
看回新聞.
人生最重要的一句說話.
你不會告訴他銀行存摺號碼.
一定不會.
你會告訴他.
我愛你.
可能你心裡最想說.

$^{521}$你很重要.
我很愛你.
911事件.
很多人說最後的遺言.
就是一句我愛你.
如果這句這麼重要.
大家一起去想.
你想和你最愛的人說什麼.
千萬不要說我很生氣.
你又扔了我的東西.
你為什麼不尊重我.
不是說這些.
而是今天如果知道.
我們想最後的說話.
想和對方說的時候.
不如今天珍惜.
多說一點.
不要說後悔沒有和你喝多點茶.
不如現在就喝茶.
或者現在.
可以更多相聚的時候.
所以福音帶來.
給我們一個生命的改變.
生命的改變就是要扔掉我們的垃圾.
要清走我們的垃圾.
讓我們能夠分清楚.
什麼叫短暫.
什麼叫永恆.
當我們能夠分清楚什麼叫短暫.
什麼叫永恆的時候.
生活很有力量.
我記得當我最早聽.
什麼叫永恆的時候.
我去種緩糞.
去學去帶領種緩糞的時候.
有一個永恆的觀念.
帶來給我很大的影響.
就是我們怎樣去告訴.
一群初信主的人.
你今天為什麼這麼重視福音.

$^{561}$就是因為有永恆價值的東西.
和沒有永恆價值的東西.
你需要去分辨.
有永恆價值的東西只有三樣東西.
神是永恆的.
所以我們要終生了認識神.
人是永恆的 人的靈魂是永恆的.
所以我們要投資在搶救靈魂的事業上.
神的話語是永恆的.
神的話是一點一滑都不會廢去.
所以我們遵行神的話語.
是永遠長存的.
所以今天當你吵架的時候.
這些是不是永恆.
不是永恆就不要罵那麼多.
今天你對老闆很不開心的時候.
老闆這樣作弄我的時候.
你就想一下這些.
你升職加薪是不是永恆.
不是永恆就滾聲吞了算了.
這些不是永恆的事.
我們能夠懂得分辨.
以致我們生命的所出來的毒素就少了.
毒素少.
我們健康的生命.
就會滋養我們整個人生.
這是第一方面.
我們要思想的.
第二方面就是.
福音改變我們本質.
所以我們整個生命是一個翻天覆地的轉變.
第二方面.
福音我們需要因為基督的緣故.
在面對困難的時候.
我們需要依靠神去面對.
所以福音.
帶來給我們.
是更加的堅忍.
患難身忍耐.
忍耐就是.

$^{601}$老練成熟.
成熟就會生盼望.
盼望就不至於羞愧.
所以讓我們看到聖經.
在這段經文描繪一幅.
很生動的圖畫.
耶穌來.
是要讓我們從外圍的社會.
到親密的家庭.
到我們的內心.
都帶來一個.
這是一個和話本.
又伸梳.
又會動刀兵.
剛才帶來不和的.
無論怎樣翻譯.
其實帶來給我們有爭論.
是一個戰場.
原來福音.
帶來給我們至少有三樣東西.
是無可避免的.
衝突.
第一是社會的衝突.
社會的衝突就是因為.
我們面對這個世界.
有很多不公不義的事.
不過我要說.
我們不是改變了世界.
我們只是改變了自己的生命.
耶穌基督差點要見證.
做一個見證.
讓我們這個見證.
讓人能夠去改變.
人能否改變是和神之間的關係.
但這個福音在我們生命裡.
是事實有一個很大的改變.
我記得我以前讀教會歷史的時候.
教會歷史的講師.
第一句上課的時候.
教會歷史是用殉道的血所寫的.

$^{641}$我當時聽到的時候也覺得很震撼.
所以就要洗耳恭聽.
如何用殉道士的血來寫成.
所以我們今天要知道.
今天有很多的前人.
付上他們生命的代價.
來付出來.
所以我們說福音是免費的.
大家信主決志.
但傳福音者.
是要付上代價的.
因為耶穌也付上了生命的代價.
所以今天我們不是.
特別是我們基督徒.
在座的大部分都是基督徒.
我覺得這張椅子很舒服.
大家不要硬硬腳.
坐在這裡就受人服侍.
我們每一個基督徒.
都應該上戰場.
我們為基督作見證.
所以我們在教會歷史裡.
告訴我們是這樣.
或者如果我們在美國.
最出名的總統.
或者仍然今天.
最受人尊重的總統.
就是1861年.
很久以前.
八十多年前.
林肯總統.
林肯最大的改變是什麼呢.
他一句名言.
我想對今天的時代.
都有一個很大的影響.
就是我要建立.
民有,民治,民享的社會.
所以他就要解放.
黑奴.
他說要推翻黑奴.

$^{681}$黑奴是不應該的.
但是當要推翻黑奴.
當要將這個信念.
去堅持的時候.
其實一點都不容易.
因為當你要堅持你的信念的時候.
很多記得利益者.
當時南北角.
或者特別南角那邊.
這麼大的問題.
覺得不應該這樣.
因為是記得利益者.
我有一個黑奴幫我洗碗煮飯.
不好嗎.
免費來為我做勞工不好嗎.
但是當他要推翻.
這個信念人生而平等的時候.
整個世界就因此.
要改變.
當然林肯這個革命.
是付了代價.
我們知道他成功推翻黑奴.
不過最後還是被.
一個傻傻的人.
殺了他.
但是我們知道這個生命代價.
是無可厚非.
我們沒辦法決定別人怎麼做.
但我今天仍然有生命氣息的時候.
我今天是怎麼為主.
作見證.
你不要說這麼遠.
我們為主作見證.
我們為主早點起床行不行.
你說為主犧牲.
連犧牲時間都不行.
你為主遲點喝茶行不行.
如果我們今天.
我們說在這個社會裡.
我們要做見證的時候.

$^{721}$其實我們的生命是需要下苦功.
每一個教徒的.
美德.
是需要我們操練的.
這是第一方面.
第二就是要我們看家庭的衝突.
這句話可能令到在座父母.
或者你以前的父母.
都很大的衝突.
耶穌是要帶來給我們.
父子.
或者是婆婆與息父.
生疏的.
兩代的溝通可能很大的.
衝突.
而事實是.
當我們曾經有一個弟兄.
為什麼他是大兒子.
獨子.
當父母離世的時候.
他要做什麼.
如果是傳統的.
他一定要做的.
你如何去選擇.
我記得我以前.
初上學.
我有一個同學.
有一個姐妹.
那時候分享.
有一天回家的時候.
為什麼我的筆記.
都掉在地上.
在家樓下.
後來知道原來我爸爸.
都飛到街上.
你信耶穌.
都飛到街上.
但沒有人抓他垃圾蟲.
當時她很生氣.
她信耶穌很生氣.

$^{761}$甚至逢星期六星期日.
就把她關在家裡.
不准她出門.
當然她分享的時候已經過了這個階段.
不知道關了幾個星期.
你怎會每天都看著她.
幾個星期後你仍然堅持.
算了,放手吧.
這就是她的堅忍.
她爸爸也沒辦法不放手.
我最近聽到一個姐妹.
分享.
她有一個姐妹.
信了耶穌之後.
在家裡不裝香.
不拜偶像.
不拜偶像也是很大的困難.
尤其是過年過節.
雞也拜過.
燒肉也拜過.
她還要做一個見證.
她要學操練做一個見證.
為了表示自己不拜.
吃也不吃.
雖然我們可以用哥林多前衰說.
是可以吃的.
但她也有一個堅持.
我不吃.
表示一個見證.
很奇妙.
兩次之後.
有一次她奶奶.
煮了一隻雞.
告訴她.
這隻雞是沒有拜過的.
這就是一個.
快樂的結局.
但實際上.
也不是一個快樂的結局.
很多時候.

$^{801}$要爭執很大的衝突.
但今天當我們面對.
我們願不願意.
付出這些代價.
可能社會沒有那麼軟.
家庭.
現在一家之主.
沒有家庭信主的煩惱.
但可能也會有個人的掙扎.
個人上的價值觀.
也有很大的衝突.
我們自己和自己.
也有很大的爭鬥.
我記得有個弟兄.
他有一次.
在茶餐廳吃東西的時候.
他知道.
多士加一元.
就有一個多士.
現在可能加兩元.
當時加一元就有多士.
當他付款的時候.
他心裡很清晰.
為什麼少了一元.
不用我付呢.
最後他想起.
一元多士應該要加.
但那個同事沒有加.
他就掙扎.
付不付一元.
於是他就付出一元.
其實這個也是個人信用.
有很大的衝突.
有時候.
電子會好一點.
有時候付多錢.
有一次我付多錢的時候.
我打算跟售貨員說.
付多錢給我.
我還沒開口.

$^{841}$他就想說不對嗎.
現在付這麼多錢.
我說不是.
整個面容都不同了.
於是當然要拿回錢.
但原來我們面對這些.
很大掙扎.
有些人說.
這個超市是大超市.
我都預料到.
我不是自己.
有些人覺得不好意思.
下次多買兩樣東西.
可能這些藉口有太多.
但我告訴你.
今天不操練.
明天也不會操練.
正如減肥是要今天開始.
明天才減肥.
一定是減不成.
所以我們今天要做的事.
就要今天去做.
無論社會也好.
要面對我們.
其實由女到愛.
雖然剛才我們說是由愛到女.
但我們整個生命.
面對每一個處境.
我們是否這麼忠於.
剛才大家唱詩歌很激昂.
我們是否這麼忠於上帝.
唱歌很激昂.
但生活上是否一樣激昂.
我們是否這樣選擇.
同神同行.
讓我們能夠配得成為.
我們主的門徒.
所以我們跟旁邊的人說.
無論遇到什麼.
人一定主會幫助我們.

$^{881}$主一定會幫助我們.
讓我們旁邊的人.
彼此去激勵.
我們今天每一樣都要堅忍.
堅忍的過程我們要實踐.
在實踐的過程中我們就會經歷.
神是真實的.
第三方面我們要思想.
不單止這個福音.
為什麼會有這麼大的改變.
就是將我們這個.
我完全改變.
雖然我們的樣子沒有馬上改變.
我們的家庭沒有馬上改變.
但原來你心裡的信念.
帶來給你整個人.
整個家庭甚至你處身.
在社會裡面.
某個角度上你都是在改革.
這個社會.
所以帶來第三方面.
信服是一個很重要的行動.
很多人誤解衝突的意思.
所以耶穌就放在.
第一位是說不通的.
就以為是沒有可能的.
放在第一位好像是一個很大的衝突.
所以我最近去一個小組的時候.
那個同事.
是非常聖經的.
他講的話.
他帶同事.
全部很合乎聖經.
他還說.
如果要信一個宗教.
信耶穌是最好的.
我問他為何不信呢.
現在就信了.
為何現在不信呢.
我相信是因為背後家庭.

$^{921}$很大的壓力.
很多人都想著.
父母先死了.
你為何要等他們先死.
你應該先信主.
等他們活生生的時候.
你才信耶穌.
但這個我們相信是很大的壓力.
父母是有很大的壓力.
信了耶穌之後真的會動刀兵.
不是開玩笑的.
但我今天是否堅持呢.
我講一個例子.
一個真正的見證給大家聽.
我有一個姐妹.
她信了耶穌之後.
有一個很大的改變.
她覺得耶穌才是真實的.
所以她就要認定這個信仰.
所以當她信了耶穌之後.
回教會生命改變.
但她老公很反對她.
她老公目露凶光地看著她.
告訴她.
你聽我說.
你一定不可以信耶穌.
但這個姐妹很堅定地看著她.
告訴她.
我一定不可以聽你的話.
她老公就說.
如果你信耶穌.
我家用都不給你.
她說你不給家用是你的事.
我都要信耶穌.
信仰是否很堅定.
她因為信服主的緣故.
她生命有這麼大的選擇.
結果是.
她老公真的不給她家用.
你不要以為是快樂的結局.

$^{961}$她老公真的不給她家用.
於是她有兩個小孩子.
她就拿東西回來.
這兩個小孩子又不可以出來工作.
她又真的這麼堅定.
她老公一路看著她改變的時候.
她自己心裡也有改變.
這個耶穌是誰.
搶走我老婆.
這個耶穌竟然是這樣.
她老婆一開始就說.
她要回教會.
一開始就在門口.
然後就坐在後面.
然後就坐在裡面.
這是一個.
生命信服主的時候.
帶來的改變.
但我們說很多的經歷.
很多的見證.
只是我們今天不相信.
所以我們不肯信服.
我們以為用自己的智慧暫緩.
我們堅忍.
暫時做臥底教徒.
遲些就會表露身份.
其實當我們今天不表露的時候.
我們今天啞忍的時候.
我們明天也會繼續啞忍.
所以我們求主幫助我們.
一時的衝突.
總比一世的內疚好.
很多時候.
有些人問我.
爸爸媽媽很麻煩.
我那時候裝完香.
我認罪.
耶穌會不會原諒我呢?.
不知道你怎樣回答他.
我說如果按著聖經.

$^{1001}$耶穌也會原諒你.
不過如果按著我的經驗.
告訴你.
你一世也會這樣問我.
你一世也會內疚.
裝完香之後就會問我.
不如一次過.
我今天就告訴爸爸媽媽.
不過我希望大家不要這麼兇.
告訴爸爸媽媽.
我們有一個溫和的心.
告訴他我真的相信耶穌.
我不再拜祭.
所以我們有一個弟兄.
跟家人說.
他媽媽有一個很大的反對.
信耶穌.
有一天他就很心平氣靜.
開著桌子就跟媽媽說.
媽媽我們現在相信耶穌.
弟弟和妹妹也相信了.
現在我們很清楚.
你死了之後.
我一定不會裝香.
當然他們很溫和.
很溫柔.
你現在怎樣決定.
他媽媽真的.
如果這樣的話.
我也相信耶穌.
有這樣的改變.
我們不相信.
我們以為我說的全是真事.
不是哄大家快點去見證.
而是這些見證.
是告訴我們是可以的.
只不過你今天肯不肯.
正如我們有些弟兄姊妹.
在職場上.
職場上告訴別人.

$^{1041}$是很難的.
但當他告訴別人.
同事也會幫他.
幫他什麼呢.
他是基督徒 星期天不要騷擾他.
是這樣的.
所以最好一開始.
我們鼓勵同事.
信主之後回到工作崗位.
第一件認知是基督徒.
你沒有機會說的時候.
你就第一件.
你擔帳.
或者你找些金具.
貼在桌子上.
擺在那裡.
有的不用擺.
擺在那裡.
這才是表明身份.
雖然沒有機會去表明.
不要緊 你就先擺著.
所以我經常鼓勵弟兄姊妹.
去找基督徒很容易.
去飯堂那裡.
看哪個基督徒是祈禱吃飯.
除非你是隱藏基督徒.
不然基本上.
你都是會祈禱吃飯.
借飯祈禱的就是基督徒.
你就先向他握手.
大家一起在職場上.
去見證神.
所以我都求主幫助我們.
如果我們看今天.
我們的生活.
阿傑果我都告訴大家.
那對夫婦是怎樣.
不單是她丈夫信了主.
她有兩個兒女.
其實是很大的風浪.

$^{1081}$我沒有時間和大家一一講.
這個見證.
有機會的話請我們來講.
但她們兩個兒女.
分別結婚.
分別信了耶穌.
而建立一個人信主.
三個基督教家庭.
其實是吧.
有三倍,三百倍.
不是八倍增長.
所以我求主幫助我們.
讓我們有這個信心.
來面對.
我告訴大家.
你不孤單的.
耶穌要和你同行.
和我們一起面對.
我們可能會說.
信仰是一條孤單的路.
有時我們聽到耶穌說.
要我們走窄路.
走窄路都是排隊走.
所以我們都知道.
走窄路.
每一個人的重擔.
我們都要自己擔.
但我們可以在前後.
互相彼此.
分擔壓力.
所以我求主幫助我們.
這條路.
信仰福音的路.
看起來好像很難走.
但福音原來帶來給我們力量.
福音就是我們的力量.
我們不是獨自行走.
因為上帝告訴我們.
我們靠著祂.
凡事都能.

$^{1121}$我們要信著神的話語.
凡事亨通.
是凡事.
聖經從來都不會騙我們.
因為我們的神是信實的.
是誠實的神.
我們要用誠實的生命.
和神一起同行.
我們就能體會.
當我們和神同行的時候.
我們就超越困難.
願神幫助我們.
我們不只是聽就算.
我們試試用神的話語.
幫助我們生活.
我們實踐的時候.
原來是真的.
某個角度上.
今天你就是我們年輕的大英姐妹的盼望.
你的見證就是激勵他們.
如何在屬靈生命路上成長.
願神祝福我們每一個人.
我們一起祈禱.
秋天天父上.
我們要多謝你.
多謝你因為你愛我們.
你選擇我們在這個世代裡.
成為你的見證者.
我們要體會福音.
原來在我們生命裡.
有一個改變.
讓我們知道.
我們不是表面看到這句話.
我們很害怕.
我們覺得很震撼.
我們覺得要革命是一件無能為力的事.
但原來當我們信了你.
接受你之後.
我們的生命就會自然有這個反應.
我們都會自然結出.

$^{1161}$像豐盛的果子.
像聖靈的果子.
人愛喜樂和平的果子.
所以我們祈求你幫助我們.
讓我們更加要浸在你裡面.
讓我們浸在你的話語裡面.
讓我們的生命能夠在你的改變裡面.
讓我們一點一滴.
讓我們的生活慢慢地改變.
讓我們的生命回顧過去的時候.
我們都會驚訝.
原來神是這樣改變我們.
也都會覺得原來你這樣.
你會藉著我們改變我們的家人.
改變我們的職場.
改變這個社會.
所以我們祈求你幫助我們.
讓我們今天.
我們不是立一個很宏大的志願.
我們怎樣改變這個世界.
而是我們求你幫助我們.
立一個很簡單的志願.
就是主啊無論是怎樣.
我們都願意.
我們能夠讓這個福音.
在我們的生命裡面.
來發揮它的果效.
主啊願你幫助我們.
多謝你聽我們祈禱.
也都願意祝福紅恩堂的弟兄姊妹.
來到這個社區裡面.
我們能夠發揮教會的影響力.
讓這個燈台.
不斷的照亮這個社區.
願你親自來祝福.
無論是祝福這張單張也好.
祝福一杯羅漢果茶.
或者奶茶也好.
讓這些點點滴滴.
弟兄姊妹的愛.

$^{1201}$能夠在這個社區裡面.
有一個好的明星.
願你來祝福.
帶領引導.
聽我們祈禱奉耶穌的名.
\newpage



\section{馬太福音 28:16-20}
\label{sec:11YyQ2efYE8}
\textbf{2025.11.23《尚待完成的使命》:馬太福音28\_16-20 (和修版) 講員 :張美薇牧師}
\newline
\newline
連結: \href{https://youtube.com/watch?v=11YyQ2efYE8}{\texttt{https://youtube.com/watch?v=11YyQ2efYE8}} ~~~~ 語音日期: 2025-11-29
\newline
\newline
\hyperref[sec:7vPzFHXDGtY]{\small{< < < PREV SERMON < < <}}
~
\hyperref[sec:index_chronic]{\small{[返順時目]}}
~
\hyperref[sec:index_scriptual]{\small{[返順卷目]}}
~
\hyperref[sec:L0es3lIs6g8]{\small{> > > NEXT SERMON > > >}}
\newline
\newline
馬太福音 28:16-20
\newline
\begin{longtable}{cl}
\hline
\hline
章節 & 經文 (和合本修訂版)\\
\hline
28:16 & \begin{tabularx}{0.7\textwidth}{X} 十一個門徒往加利利去，到了耶穌指定他們去的山上。 \end{tabularx} \\ \\ \relax
28:17 & \begin{tabularx}{0.7\textwidth}{X} 他們見了耶穌就拜他，然而還有人疑惑。 \end{tabularx} \\ \\ \relax
28:18 & \begin{tabularx}{0.7\textwidth}{X} 耶穌進前來，對他們說：「天上地下所有的權柄都賜給我了。 \end{tabularx} \\ \\ \relax
28:19 & \begin{tabularx}{0.7\textwidth}{X} 所以，你們要去，使萬民作我的門徒，奉父、子、聖靈的名給他們施洗， \end{tabularx} \\ \\ \relax
28:20 & \begin{tabularx}{0.7\textwidth}{X} 凡我所吩咐你們的，都教導他們遵守。看哪，我天天與你們同在，直到世代的終結。」 \end{tabularx} \\ \\
[1ex]
\hline
\hline
\end{longtable}
$^{1}$弟兄姊妹平安.
很高興今天有機會再一次回來.
跟大家分享神的話.
我想起碼一年半以上了.
我病之前應該來過.
那次好像是三月來過.
我四月就病了.
所以一年半沒有來.
不過也很高興.
當然與你們無關.
絕對沒有這個意思.
病在裡面.
遲早也會出來.
只差幾時發.
首先要恭喜你們唐慶.
你們這個月是唐慶月.
你們是2011年開始的.
我仍然記得.
十多年前的時候.
跟李牧師四處去看.
來到洪水橋.
那時候洪福村還沒有.
現在有洪福村.
有很多人搬進來.
對你們來說是一個很大的和唱.
所以你們那隊的隊聯.
也很有福音的意識.
懷抱報道初心.
福音遍及社群.
你們是一個報道的初心.
因為當直堂出來的時候.
通常人數也不多.
所以很著重傳福音去報道.
我們開始的時候.
首先祈禱.
天主上願你使用孩子今天所說的.
叫孩子所說的.
都是你要孩子說的.
也幫助會眾.
他們有一個能聽到的耳朵.

$^{41}$有一個能明白道理的頭腦.
更加有一個心.
能夠將你的話藏在他們的心裡.
有手有腳.
去行出你吩咐他們行的.
使他們能夠聽到.
行到土里喜悅.
做一個土里喜悅.
合理心意的信徒.
聽我們的祈禱.
奉耶穌基督的名求.
阿們.
耶穌基督升天已經兩千年有多了.
福音也由發源地耶路撒冷.
可以說是傳到地極.
現在世界上每一個地方.
七大洲五大洋.
都好像是傳遍了.
如果我們看回這張聖經的地圖.
比較小一點.
我們可以從耶路撒冷那裡圈住那裡.
我們會發覺.
由耶路撒冷然後猶大地.
去到撒瑪利亞.
然後去到安提阿.
然後去到亞細亞.
去到歐洲.
去到羅馬.
在保羅的時候.
他還說想去到西班牙.
其實那裡已經看到.
其實是不夠一百年.
福音已經去到當時.
聖經所提到的世界各地.
我不知道大家看這張圖的時候.
有沒有想過.
我們中國在哪裡呢.
我就找了一張很古老的地圖.
如果大家看看這裡說的話.
是主後19年的地圖.

$^{81}$所以樣子不是很正宗.
大家看不看到.
這裡是哪裡.
這裡是非洲.
歐洲就在上方.
這裡應該是中東亞洲.
這裡就是印度.
有沒有看到中國在哪裡呢.
原來中國在這裡.
塞內斯在那裡.
根據傳說.
多瑪去了哪裡呢.
這裡就是我們剛才看到的地方.
就是剛才現代化地圖的地方.
多瑪去了印度.
傳說說巴多羅馬去了中國.
所以主後一百年左右.
福音都去到地極.
連印度和中國都去到.
但又有多少人信耶穌呢.
大使命又執行到哪裡呢.
我們今天就好好地去看這五節經文.
這五節經文.
我會分五個範疇去看.
首先就是背景.
當十一個門徒去了加利利的時候.
耶穌約定他們.
耶穌在他們被賣之前的晚上設立聖餐.
然後在他們去赫西瑪尼之前.
就跟他們說.
今晚你們為我的緣故都要跌倒.
因為經上記著說.
我要擊打牧人.
羊就分散了.
然後耶穌就說.
但我復活以後.
要在你們以先往加利利去.
他們其實去加利利那裡.
就是耶穌在那一個晚上.
就跟他們約好了.

$^{121}$而在那個山上.
耶穌就向他們頒佈大使命.
當時有人拜祂.
拜祂的意思是尊崇祂.
當祂是神聖.
因為祂復活了.
但仍然有人疑惑.
不是很清楚.
其實當時十一個門徒都在那裡.
十二個已經死了.
但可能不止這個數.
因為保羅在哥林多前書十五章六節.
就說曾經有五百人見過復活的耶穌.
很可能說的就是這裡.
是也好不是也好.
這裡應該不止是那十一個門徒.
但當他們去到加利利的山上的時候.
耶穌就向他們頒佈了大使命.
他頒佈大使命之前.
他就對他自己有一些陳述.
我稱之為陳述.
就是說耶穌進前來對他們說.
天上地下所有的權柄都賜給我了.
耶穌首先去陳述一下.
他現在的境況.
他的境況就是天上地下所有的權柄都給了他.
大家先記住這裡.
我們一會兒到最後的時候.
我們都會再說的.
但耶穌的陳述就是說.
天上地下所有的權柄都給了他了.
他告訴這十一個門徒.
和很可能當時都在的人.
然後他就向他們.
充滿了權柄的耶穌.
向他們將使命頒佈給他們.
使命是什麼呢?.
就是第十九節裡的那一句.
使萬民作我的門徒.
所以耶穌是吩咐門徒們.

$^{161}$要去使萬民去做耶穌的門徒.
他對那十一個門徒說.
你們要使萬民去做我的門徒.
我們就會看萬民是什麼意思呢?.
可能從地方來說.
萬民就是世界各地.
剛才都說過了.
七大洲,五大洋.
或者各國.
好像聯合國的中英美法蘇.
還有其他的很多很多的國家.
現在差不多有二百多個國家.
也都包括沿海的,內陸的.
或者山區的各國.
地方上都可以這樣分.
又或者當我們用人的族群來分.
就是國民,國族,各族群.
包括中國人.
其實中國人都很多的.
大家有沒有想過.
中國人我們有香港人.
還有大陸人.
還有台灣人.
還有澳門人.
我們說這個叫做兩岸四地.
然後世界各地的華人.
中國人他們說的話.
剛才都聽到了.
後面有些印尼的華僑.
他們就說一些.
不知道什麼印尼話.
不同的地區的話.
有說普通話.
有說廣東話.
有說福建話.
有說潮州話.
有說客家話.
或者在外地的第二代.
說當地話的土生華人.
有說英文.

$^{201}$有說法文.
有說德文.
有說各方面.
有做餐館的.
有做的士司機的.
有做生意的.
或者做各行各業.
各民.
每一個人.
這個就叫做萬民.
使萬民做我的門徒.
我們又看一看呂鎮忠譯本.
有沒有聽過呂鎮忠譯本.
呂鎮忠譯本就將這裡說成.
他說天上地下的一切權柄.
都給了我了.
耶穌說的.
所以你們要去.
使一切外國人都做門徒.
奉父子聖靈的名.
給他們施洗.
教訓他們遵守.
我所吩咐你們的一切事.
看罷一切的日子.
我都和你們同在.
直到今天的完結.
呂鎮忠譯本用很清楚的.
四個一切.
將這個大使命連在一起.
不過他那個一切外國人.
大家覺得準不準確.
好像漏了一些東西.
包括中國人嗎.
其實是包括的.
萬國萬民.
不只是外國人.
我們都在裡面.
所以中國人都在.
不是一切的外國人.
他應該要說.

$^{241}$我不知道應該怎樣說.
將外國人.
將我們都包括在裡面.
保羅是外邦的使徒.
他關心猶太的同胞.
正如保羅一樣.
我們香港的信徒.
萬民都是包括外國人.
也包括香港人.
中國人.
大陸人.
台灣人.
澳門人.
或者灣區的人.
最近十五大運動會.
我們灣區人.
是很集中在一起.
是中國人.
外國人.
和香港人.
我們全部都包括在一起.
當我們又再看一看.
現在來說.
萬民又有些什麼呢.
我是教宣教的.
所以都想和大家說一下.
2025年.
全球人口.
我在.
什麼時候.
應該是上個.
應該是上個禮拜.
上個禮拜二.
的早上.
大概下午五點四十六分就看到.
當時來說.
全世界的人口.
應該有82.59億人.
不過這個數字是2025年初.
因為無論如何都要截的.

$^{281}$不可能永遠這樣數下去.
先計算.
全球人口是81.92億.
全世界萬民.
有多少人是信耶穌的呢.
那麼.
廣義的基督徒.
大概有32.3\%.
世界上高.
而客義來說.
即是福音派的基督徒.
大概佔10\%.
這10\%.
包括一切福音派的信徒.
而廣義的基督徒.
包括接受英鞋洗的信徒.
一些自認是基督徒的人.
包括天主教徒.
東正教徒.
科林教徒.
光榮教徒.
掛名信徒等等.
所以大家可以記住.
我們大概有三分之一人.
被稱為基督徒.
但如果是福音派.
就大概有10\%.
我們又可以看看.
我們說洗萬民.
去做耶穌的門徒.
我們在過去的日子裡.
去傳福音的結果是怎樣呢.
我們可以去看.
我們看第二行.
看基督徒.
我們比較1910年和2025年.
即是大概100年.
大概100年前.
全球人口只有17.58億.
去到2025年.

$^{321}$就已經有81.92億.
在1910年的時候.
基督徒廣義來說.
就有6.12\%.
是6.12億的基督徒.
佔全球人口的34.8\%.
去到2025年年初的時候.
基督徒廣義來說.
就是26.45億.
多不多呢.
其實都多的.
是佔全球人口的32.3\%.
即是說在100年之內.
人口多了20億的基督徒.
多不多呢.
多.
因為聖經說什麼.
有一個人在地上.
是願意信耶穌的話.
天上的天使都會拍手歡呼.
何況有20億是很多.
但我們再看小心一點.
好像不是很對路.
因為100年前是34.8\%是基督徒.
但現在32.3\%.
即是說什麼.
我們跑輸了出生率.
出生的人可能比我們傳福音更快.
所以我們要留心.
這裡耶穌給我們的使命.
是叫萬民做主的門徒.
我們好像有一點點.
我們又看看其他數字.
很少的.
我們隨便看看.
其實呢.
如果我們看中國民間宗教.
或者無宗教信仰者.
在這100年裡.
也轉變很大.

$^{361}$因為中國民間信仰.
由22.2\%去到5.5\%.
為什麼呢.
原因是國內在過去100年.
從一個有宗教自由的信仰.
變成什麼.
一個無神論的國家.
所以有很多國內的中國民間信仰.
他們就叫做轉了.
其實他們實際點.
未必一定是轉的.
而無宗教信仰.
由0.2\%變成12.8\%.
因為國內是無神論.
這是政治的原因.
令到有轉變.
但我們也會看到穆斯林.
百分比是雙倍.
由12.6\%去到25.2\%.
而按照估計.
去到2050年.
穆斯林的人口.
按照宣教學者估計.
可能超越廣義基督徒的數字.
大於廣義基督徒.
所以我們要留心.
要努力.
真的要熱心地.
向這些人傳福音.
我們又看看.
其實10\%是福音派的信徒.
32\%是廣義基督徒.
他們不是平均分配在世界各地.
是不平均的.
我們又看一幅圖.
這幅應該是我給大家看.
最後的一幅.
不是平均分配.
大家會看到.
橙色是福音派信徒.

$^{401}$橙色比較多的是拉丁美洲.
北美洲.
或者是非回教的非洲.
非洲有一部分是回教國家.
有一部分是非回教.
非回教的非洲國家.
也有多福音派的信徒.
其實中國也不少.
如果大家這樣看.
大家會看到橙色的是福音派信徒.
而淺黃色是其他的基督徒.
廣義的基督徒.
其實廣義的基督徒.
也需要傳福音.
但如果我們看.
中國就不是那麼多廣義的基督徒.
因為在中國.
你做基督徒是要付代價的.
所以不是那麼多掛名的基督徒.
但綠色又是什麼呢?.
綠色是近文化的非信徒.
綠色的不是信徒.
是未信主的.
但是他們在附近.
已經有跟他們同文化的教會.
例如洪水橋這裡.
我們今天來教會聚會.
我們是橙色的.
但在我們附近有很多人.
他們仍然未信主.
但我們比較容易去到他們的裡面.
跟他們傳福音.
我們又不用學語言.
又不用學各樣的東西.
又不用跨任何的界限.
我們都可以去跟他們傳福音.
所以綠色的意思是.
可以跟我們同文化的.
他是非信徒.
我們看中國就有很多.

$^{441}$這些同文化的非信徒.
是等待著我們去將福音傳給他們.
而藍色是什麼呢?.
藍色就是遠離文化的非信徒.
或者我們簡單來說.
叫做未得之物.
就是說在他們當中.
可能未必有教會.
跟他們文化的教會.
又或者在他們當中.
他們自己出不到去傳福音.
因為他們的信徒很少.
力量很小.
所以很需要有人去到他們那裡.
去傳福音給他們.
大家記住這張圖.
少少.
大家都不會記得很多.
大家記住三個色.
黃色就是基督徒.
綠色就是近文化的非基督徒.
而藍色就是未得之物.
其實藍色和綠色都未信主.
他們大概佔67\%左右.
這些就是藍色和綠色的萬民.
不是平均的分配在各國裡面.
我們又回到耶穌的使命.
耶穌的使命就是說.
我們要使萬民做耶穌的門徒.
英國的神學家維格斯韋.
在1923年他說.
世人未能夠看見耶穌的原因.
是基督徒沒有被耶穌所充滿.
他們以為每個星期三加上崇拜.
偶然讀聖經.
有時祈禱一下就夠了.
他說我最難過的是.
看見自稱是基督徒的人.
卻沒有生命,沒有能力.
生活方式跟非信徒差不多.

$^{481}$以至難以分辨他們是屬血氣的.
還是屬聖靈的.
這些是當時神學家這樣說.
他說他覺得很可惜.
因為基督徒沒有被耶穌的充滿.
但耶穌吩咐我們.
使萬民做我的門徒.
這是一個動詞.
是一個命令式的動詞.
我們使一個人成為基督的門徒.
是什麼意思呢?.
就是將他帶入一個師生的關係當中.
接受基督的話.
負起基督給我們的使命.
或者我們用一個形象.
負起基督給我們的額.
信服,接受基督的話是真理.
執行基督的話.
耶穌的門徒必須使別人.
像他們一樣成為耶穌基督的門徒.
耶穌對十一個門徒說的.
而這十一個門徒對他們的門徒.
一直說,直到現在來到我們這裡.
就像耶穌現在對大家說.
你們要使萬民做耶穌的門徒.
你們會說.
我們怎樣可以使萬民做耶穌的門徒呢?.
其實耶穌有說的.
耶穌在大使命裡有說的.
祂給了三個方法.
我相信不止三個方法.
但耶穌在這裡給了三個方法.
第一就是說.
所以你們要去使萬民做我的門徒.
去,要出去.
要出去,要放下一些東西.
才可以出去.
我想大家都知道亞伯拉罕.
當神父召亞伯拉罕從吾爾到迦南地時.
他要離開本地本族父家.

$^{521}$往神所要指示他的地方去.
所以亞伯拉罕要放下他的東西.
去,去有什麼要做呢?.
就是跨越.
去的意思是要越過一些障礙.
跨過一些不同的東西.
障礙是很多的.
我們常常會想.
如果要去宣教,要去外國.
我們最多,最流行的.
做宣教士一定要去外國.
是海外宣教.
跨越各界,離開我們的家人.
但現在不是.
現在在附近.
我們剛才聽到後面有些說印尼話的人.
我們附近也有很多印尼話的人.
我們也有些印尼姐姐.
現在我們當中也有印尼姐姐.
我們可以跨過一些言語的障礙.
一些文化的障礙.
一些社會階層的障礙.
一些生活習慣的障礙.
地域的障礙.
如果我們要去傳福音.
一定要跨過一些障礙.
所以大家去想.
我幫旁邊的人傳福音.
有什麼障礙是要跨過的呢?.
其實當我們去看.
當耶穌基督.
其實耶穌基督就是一個跨過障礙的宣教士.
祂道成肉身.
即是祂從天上來到地上.
然後祂從無限變成有限.
從全知變成一個像學生一樣.
從小時候開始.
再重新學習.
再重新來到我們當中.
祂全然聖潔.

$^{561}$來到一個骯髒的地方.
祂從安定變成漂流.
祂從喜樂的天堂.
來到痛苦流淚的世界.
祂離開天父的懷抱.
成為人子.
越過自己的界限.
來到我們中間.
用人能夠聽得懂的言語.
來與我們相交.
降到我們的層面.
來與我們說話.
其實很多人都說過.
你見到一隻螞蟻.
你無法與牠說話.
你要變成一隻螞蟻才能與牠說話.
耶穌就是這樣.
來到我們的層面.
與我們說話.
否則我們怎樣能夠明白祂呢?.
其實除了耶穌是這樣之外.
早期教會都是一群去的人.
他們要越過自己的界限去傳福音.
好像巴拿巴原本是一個富有的人.
他變賣了他的田產.
到有需要的人當中.
保羅越過了地域.
與法利賽人的幽穴.
做一個福音使者.
西方的宣教士.
去到中國,印度,非洲.
穿上中國服裝的.
有很多不同宗教的階層工作.
有連邊都紮了.
又有在非洲部落為主作供.
我還見到.
曾牧師有位老師.
他的照片是沒有穿衣服的.
因為他去到佩波尼爾堅尼.
在那群不穿衣服的土人工作.

$^{601}$他是一個美國的宣教士.
但他卻願意這樣去做.
有數之不盡的例子.
但重點是甚麼?.
我們要去.
去都不夠.
還有呢?.
就是私洗.
你們要去洗萬民作惡的門徒.
奉父子聖靈的名.
給他們私洗.
去私洗.
其實去到私洗已經是走了很大步.
很感恩.
你們上星期有洗禮.
一會兒也有見證.
很開心的事.
但我們知道.
葉松茂博士說.
我們要吸引到人回來.
留下來.
讓他繼續成長.
然後帶領他信主.
信了主之後.
教導他.
一直上.
我不知道你們要上甚麼班.
以前的感受就是要上初信班.
要上基督生平班.
然後上洗禮班.
然後才可以洗禮.
要上一年的主一學.
然後才可以私洗.
私洗之後都不停.
還要上事奉學習班.
繼續上下去.
所以私洗.
私洗是得到栽培.
得到幫助.
得到教導.

$^{641}$還有剛才也說過.
私洗之後不是停下去.
第三個.
原來我漏了一個.
使萬民作我的門徒.
去私洗.
教訓他們遵守.
所以你們要去.
使萬民作我的門徒.
奉父子聖靈的名.
給他們私洗.
凡我所吩咐你們的.
都教訓他們遵守.
所以私洗完之後還未停.
待會兒你們講完見證.
要繼續上主一學.
要繼續上.
而且不只是上.
我們要遵守.
我們知還要行.
這才叫學識.
不然只知道只聽.
頭就越來越大.
但沒有行動.
沒有實際的影響.
是遵守.
教訓他們遵守.
所以學的就要虛心學習.
教的就要教訓人遵守.
不只是把知識教.
要教他們去遵守.
教人祈禱.
讓他們懂得祈禱.
能和你一起祈禱.
教人傳福音.
就要和他們一起傳福音.
這樣才行.
不單是教導.
還要遵守基督的說話.
做一個基督的門徒.

$^{681}$是很切實的生命改變.
我們做門徒.
要接受賢教.
也要接受新教.
然後再去做別人的賢教,新教.
我們看回過去兩千年.
有人去大人信主.
有人去施洗.
叫人的生命改變.
為耶穌增加門徒.
但基督徒的人口也有很大的改變.
我們看回這張圖是1978年.
我們不看那麼多.
我們只看中國.
看不看到中國那時是怎樣.
1978年.
1978年即是四十五年前.
即是四十五,六年前.
那時橙色有多少.
看不看到? 看不到.
綠色有多少? 也看不到.
大部分是甚麼? 是藍色.
但現在的中國.
現在的中國我們的橙色是好.
我們的綠色是全世界最多綠色的地方.
就是中國.
所以中國人真的要努力傳福音給我們的同胞.
因為神明已經開了很多很多的和祥.
在我們的附近.
反而中國裡面的藍色就不是那麼多.
所以有人說到可能去到2030年.
中國是全球最多基督徒的國家.
雖然仍然受迫迫.
但很可能是全球最多基督徒的國家.
所以當我們看到,真是不簡單.
但當我們看看我們的資源是怎樣用在傳福音上呢?.
2025年全球猜全數據那裡拿出來的.
全世界大概有四十五萬個宣教士.
用的錢,放在宣教上的錢.
大概是四百五十億.

$^{721}$在世界,我們這次就把它們混在一起.
我們有A世界,B世界,C世界.
記不記得那三個顏色?.
C世界是甚麼?.
是廣義基督徒,佔百分之三十二.
B世界就是福音已達,但未信的基督徒.
大概有三十四個百分比.
A世界就是美德之民.
也大概有三十四個百分比.
那宣教的資源是怎樣分配的呢?.
首先,我們有百分之七十二的宣教士.
就在C世界裡面.
大家要明白C世界是甚麼?.
廣義基督徒的世界裡面.
為甚麼還要有那麼多宣教士?.
有原因,因為我們要教他們出去宣教.
我們要教他們去做各樣的事.
去跟不同的地方去聯繫.
所以有百分之七十二的宣教士就在C世界裡面.
有百分之八十七的宣教費用就用在C世界裡面.
我們看看B世界,B世界有多少呢?.
B世界有百分之二十五的宣教士.
因為B世界是福音已達,但他們未信主.
還有很多人未信主.
所以都有百分之二十五的宣教士在他們當中.
有百分之十二的宣教費用就在B世界裡面.
A世界是甚麼?.
未得知民.
意思是A世界是甚麼?.
其實是最需要宣教士.
A世界是最需要宣教的經費.
但我們看看,原來只有百分之三的宣教士在當中.
也只有百分之一的宣教費用在當中.
大家看到的時候會覺得怎樣?.
好像有點錯配.
我們未必需要所有的資源都在A世界.
但A世界可能需要的比較多.
大家都需要為這件事祈禱,思想.
當我們這樣再問的時候.
我們會問教會又應該怎樣去用我們的資源呢?.

$^{761}$很簡單地說明.
林安國牧師有一個叫五十五十教會的理論.
是他的理論.
他說五十五十的教會.
他說五十五十教會中間的五十就是教會的本身.
教會本身就會有向上敬拜.
因為我們重點是大家星期日回來都要來敬拜.
還有向內的強化.
譬如我們需要用錢在港台供應.
有門徒訓練.
系統的教導.
肢體的生活.
這個是教會本身.
林安國牧師說最好是百分之五十.
將教會的人力物力財力用百分之五十在教會裡面.
另外百分之五十用在哪裡呢?.
他說百分之二十五用在我們本身這個社區.
百分之二十五用在向外的跨越.
向外的社區.
本身的社區有什麼可以做呢?.
有社區的報道.
救靈魂的行動.
譬如有些叫吸力的報道.
動力的報道.
直堂的報道.
又或者有些是報道以前的愛心行動.
有些是剛嚴的報道.
橋樑.
即是跟人的關係.
傳人的報道等等.
當在社區裡面有的報道的工作.
有的事情做到有成果的時候.
那些人信得住就會回來.
就會回來教會.
所以就會擴大的增長.
數量上的增長.
都會回到教會增長.
當我們做社區的工作.
另外百分之二十五.
就會向外的跨越.

$^{801}$有近文化的報道直堂.
或者跨文化的報道直堂等等.
即是做宣教的工作.
其實林匡所牧師就是說.
這個五十教會其實不是一個死的東西.
也不是一個必然的東西.
其實給大家做一些參考.
當如果大家在教會裡面.
你們是有這個崗位.
是要去負責去看教會.
怎樣用你的人力物力財力的話.
最好就是這樣.
如果教會是九十十的話.
很可能教會就會是一間很內聚的教會.
沒有什麼太外向的東西.
我記得我以前一向都是用北宣做我的例子.
就經常都會很希望我們會好像北宣那樣.
大概他們是五十五十的.
即是向外和對內.
其實給大家去做一些參考.
這裡是說教會.
我們個人又怎樣呢?.
我們個人是怎樣去傳福音.
我們怎樣去帶門徒.
又或者我們個人的奉獻.
當然是要透過教會去使用.
所以教會就好像用了那裡.
但是我們個人在傳福音上是怎樣呢?.
我知道大家都有出去傳福音.
你們都有些弟兄姊妹傳福音很厲害.
假如你能夠一天帶一個人順著.
三年你帶了多少人順著呢?.
三年當你一千人吧.
做一些好的做法.
每三年帶一千人順著.
從今年開始.
2025年開始.
2028年就有一千零一人.
因為你自己還順著.
然後呢.

$^{841}$2031年就有二千零一人.
2058年就有一萬一千零一人.
多不多呢?.
多!.
要不要鼓掌呢?.
一定要鼓掌.
但是好像都不是很夠.
但假如你一年,兩年,三年去帶門徒.
耶穌說你們要去使萬民做我的門徒.
三年帶出十一個門徒.
誰是這樣呢?.
耶穌就是這樣.
祂三年帶出十一個門徒.
第一個三年.
很少人.
一加十一.
十二.
第二個三年.
他們再去門徒的意思是.
他們會重覆去做你教他的.
所以十二乘十二.
因為你和另外十一個門徒.
就十二乘十二.
就一百四十四.
都不多的.
但當你一路一路地乘下去的時候.
去到2058年的時間.
就會是七百四十三億人.
其實那時候呢.
去到2080年.
人口大概是超過了一百億.
那時候其實已經是超過了人口.
而2058年距離現在多少年呢?.
三十三年.
就是耶穌再世的日子.
原來是可以的.
如果就這樣計算來看.
是可以的.
是我計算出來的.
你們可能幫我看看.

$^{881}$我有沒有計錯.
我的數字都算不錯.
是我計算出來的.
當然我們說不是的.
實質上不是.
第一我們不是耶穌.
耶穌三年可以帶出十一個.
我們未必帶出十一個.
還有困難的.
有夭折.
其實耶穌是帶出十二個的.
為什麼會少了一個呢?.
因為猶大背叛了.
我們未必像耶穌一樣.
帶出十一個.
一代一代傳下去.
會變樣的.
會變歪了.
但不要忘記.
我們還有耶穌給我們的應許.
我們回到.
耶穌說.
「我常與你們同在」.
只要我們去洗萬民.
做耶穌的門徒.
奉父子聖靈的名.
給他們施洗.
凡我所吩咐你們的.
都教訓他們遵守.
我常與你們同在.
直到世界的末了.
耶穌是有應許的.
祂說只要我們去的話.
只要我們去教導.
我們去帶人信主.
祂就會與我們同在.
直到世界的末了.
祂是誰呢?.
祂是天上地下.
所有的權柄.

$^{921}$都在祂那裡.
我們還需要害怕嗎?.
還害不害怕?.
不需要害怕.
所以要不要去?.
去!.
所以我們要去.
我通常說到.
都是這樣結束的.
最近都是這樣.
與我何干?.
與你有什麼關係?.
去!.
去帶人信主.
如果你說:我好像不太行.
我怎樣可以叫人跟隨我呢?.
那你就做好一點.
將你學到的去用出來.
如果你說:我好像有些東西不太懂.
回來讀主一學.
叫老師去幫你.
去找教會的導師.
去成為你的幫助.
讓我們每一個都出去.
我相信當我們去的時候.
主就會祝福.
主就會使用.
我們低頭祈禱.
天父上帝.
你給我們這個使命.
我們初時看的時候.
好像覺得我們的使命.
好像很難遵守.
又或者好像很難達到.
但原來看著算著.
又好像不是真的難得這麼厲害.
我們就求主親自的幫助.
帶領.
叫洪映嘉的每一個弟兄姊妹.
都回應你的說話.

$^{961}$願意成為你的門徒.
遵行你給我們的命令.
我們今天禱告.
是奉我主耶穌基督的名為求.
阿們.
\newpage



\section{馬太福音 5:43-48}
\label{sec:L0es3lIs6g8}
\textbf{2025.11.30《我們要完全如同天父》:馬太福音5\_43-48 (和修版)}
\newline
\newline
連結: \href{https://youtube.com/watch?v=L0es3lIs6g8}{\texttt{https://youtube.com/watch?v=L0es3lIs6g8}} ~~~~ 語音日期: 2025-12-03
\newline
\newline
\hyperref[sec:11YyQ2efYE8]{\small{< < < PREV SERMON < < <}}
~
\hyperref[sec:index_chronic]{\small{[返順時目]}}
~
\hyperref[sec:index_scriptual]{\small{[返順卷目]}}
~
\hyperref[sec:WQHTSm8b4dk]{\small{> > > NEXT SERMON > > >}}
\newline
\newline
馬太福音 5:43-48
\newline
\begin{longtable}{cl}
\hline
\hline
章節 & 經文 (和合本修訂版)\\
\hline
5:43 & \begin{tabularx}{0.7\textwidth}{X} 「你們聽過有話說：『要愛你的鄰舍，恨你的仇敵。』 \end{tabularx} \\ \\ \relax
5:44 & \begin{tabularx}{0.7\textwidth}{X} 但是我告訴你們：要愛你們的仇敵，為那迫害你們的禱告。 \end{tabularx} \\ \\ \relax
5:45 & \begin{tabularx}{0.7\textwidth}{X} 這樣，你們就可以作天父的兒女了。因為他叫太陽照好人，也照壞人；降雨給義人，也給不義的人。 \end{tabularx} \\ \\ \relax
5:46 & \begin{tabularx}{0.7\textwidth}{X} 你們若只愛那愛你們的人，有甚麼賞賜呢？就是稅吏不也是這樣做嗎？ \end{tabularx} \\ \\ \relax
5:47 & \begin{tabularx}{0.7\textwidth}{X} 你們若只請你弟兄的安，有甚麼比別人強呢？就是外邦人不也是這樣做嗎？ \end{tabularx} \\ \\ \relax
5:48 & \begin{tabularx}{0.7\textwidth}{X} 所以，你們要完全，如同你們的天父是完全的。」 \end{tabularx} \\ \\
[1ex]
\hline
\hline
\end{longtable}
$^{1}$弟兄姊妹平安.
聽到嗎?.
好.
很感恩可以再來紅恩堂分享神的話語.
讓我們透過神的話語.
可以彼此激勵.
不知道大家還記不記得我上次說的那一套.
就是真理的揭露.
和恩典的扶持.
就是說原來神的話語.
是很想幫我們更加認識了解自己.
或者是勇敢去面對自己一些陰暗面.
一些罪性.
同而不是神是一個高高在上的審判官去對待我們.
而是祂是我們的同行者.
陪伴者.
一直沒有離開我們.
對我們不離不棄的陪伴.
以致我們能夠去衝過那些艱難.
我們可以在主裡面去被愛.
被蒙恩受德蒙的連戳.
可以繼續向前走.
這個星期發生的事其實都很震撼.
我不知道大家有時不想看到這些新聞.
但是又很想知道新聞發生了什麼.
不知道可以怎樣去走.
我都有想過.
那我需要改變今天所說的訊息嗎?.
其實都很不簡單.
我先說了今天這個訊息.
其實有些事情都可以有些共鳴或者回響.
我們就看看怎樣求聖靈.
幫助我們去消化.
去明白祂的話語.
我們去同感在災場的每一個.
或者受災的周圍的人.
其實都有感覺的.
我們怎樣去一起去面對去處理.
我們先一起祈禱好不好?.
親愛的天父求你的話語.

$^{41}$求你的聖靈.
主耶穌求你愛的陪伴.
去跟我們一起.
因為這個星期香港發生了很大很震撼的事.
我們大家都是.
這些景象歷歷在目.
是很多的不容易.
我們都真的很不想再有下一次發生這些情況.
主啊求你真的去憐憫香港這個城市.
你也去醫治香港這個城市.
特別是你的愛.
你的恩典.
你的醫治.
你的同在.
是在大埔區.
裡面你給大埔區的教會.
能夠同心協力的.
去與這些居民去同行.
很多的不簡單.
求主也給我們.
雖然不是在那個社區裡面.
但是我們是同一個城市.
讓我們不忘去為他們禱告.
甚至在經濟上我們可以作出支援.
甚至我們預備自己的狀態.
如果冬天有一天.
你讓我們遇見這些災情中的朋友.
可能是他們的家人.
可能是他們的至親.
就是因為這一場火災而離世.
很多的傷痛.
很多的損失.
主啊求你幫助我們.
能夠有聆聽的耳朵.
有關懷的態度.
願意去與他們一起同哀哭.
與他們一起通過這些艱難.
主啊求你賜力量給我們.
讓我們去能夠與他們同行.
讓我們不要被香港很忙碌.

$^{81}$那個很快速的節奏.
讓我們好像很快就可以將這件事淡忘.
主啊求你與我們同行.
我們誠心禱告交託.
奉主明求.
阿們.
今天選擇這段經文的時候.
其實是有一點都很想.
我們與天父一樣的心腸屬性去連結.
以致我們才能夠活出神的模樣.
如果大家深刻記得.
我曾經說過.
我是由破碎家庭而來.
而我的破碎又與澳門有關係.
我們說.
好像災情,火災.
我們一看就看到很靠近.
但大家有沒有想過.
其實有些家庭都可以是一個災區.
但不是一些很顯而易見的事.
有些是關上門才會發生的事.
所以我們都會想.
究竟我們怎樣去到.
可以好像神一樣.
如同天父完全一樣.
可能這個講題已經覺得.
我們要如何完全如同父.
什麼是完全呢.
有些人誤會了.
以為完全是完美.
大家知不知道什麼是完美.
如果用英文去說.
怎樣說完美.
Perfect.
我們就會想.
我們怎可能有人是十全十美呢.
這就是指標我們永遠達不到.
原來這是我們錯誤解讀了.
原來神想我們好像他一樣.
完全不是要我們完美.

$^{121}$完全是什麼意思呢.
其實有四方面.
第一方面是說道德和屬靈上的成熟.
其實完全不是說沒有罪.
或者是完美無缺.
而是說生命的成熟.
品格的完整.
就好像天父的屬性一般.
其實天父.
我們說三位一體神.
天父,聖父,聖子,聖靈.
耶穌,聖子來到人間.
道成肉身的時候.
其實是很想展露一個完全的人性.
耶穌的完全人性.
和昔日亞當夏娃.
因為罪而扭壞了的人性是不同的.
所以有時我們會罵人.
如果用人性來形容.
是怎樣呢.
你都沒有人性.
有沒有記得會這樣罵.
即是說他缺乏了神起初做人的人性.
如果真是上帝起初創造的人性.
其實沒有被罪影響的時候.
其實就是耶穌的道成肉身的展露.
耶穌有憐憫,有愛顧.
不會被一些死胡胡的律法框架.
就框死.
也不是會有那些階級觀念.
因為耶穌面前每個人都是平等的.
所以你會見到昔日有些敗壞了的人性.
或者一些有缺乏的人性.
耶穌經常罵.
罵那些法利賽人,文士.
因為他們懂一點,理解一點.
但其實不是做到真正神想人去發揮那種人性.
耶穌來到是要展現.
究竟真是要守住神的律法.
最重要的是兩樣.

$^{161}$哪兩樣?.
愛神,愛人.
就是這兩樣.
你說很多話.
但原來你不愛神,不愛人.
那就是浪費力氣.
耶穌是這樣去展述的.
所以其實在這段經文.
我們看到一些很有趣的東西.
我最初看的時候不是很明白在說什麼.
但慢慢在生命裡面.
去到讓神的話語.
去到在自己生命裡面.
怎樣去實踐的時候.
就拳拳到肉的時候.
就明白更加多了.
求聖靈幫我們去理解.
43節.
「你們聽說過」.
這是耶穌去說天國的理論.
如果大家有留意.
《馬太福音》第五章一直說下去.
其實就是想說明.
所有的律例典章其實應該是怎樣的.
因為有些聖經學者.
當時已經曲解了很多神的話語.
令到當時的以色列文.
都很被受壓制或者受捆綁.
耶穌反而怎樣去教呢.
祂說要愛你的鄰舍.
祂說你們聽說過.
這個聽說過就是當時的律法師這樣說的.
要愛你的鄰舍.
恨你的仇敵.
這個教導聽起來好像很合理.
所以你就明白為什麼有些舊約.
有些人拿來講的時候.
就說要怎樣.
以牙還牙.
然後以眼還眼.

$^{201}$但是說真的.
我們如果人有罪性.
我們是不是真的可以.
「塌塌點」.
以牙還牙.
以眼還眼呢.
還是你多贈一批簪給他.
多踢一腳給他.
你的仇恨反而是不斷地加增.
這些是舊約.
有時候可能有些規矩就這樣定了.
但是人性很難很有標準.
做得到這些事.
所以在第44字.
耶穌反而怎樣說呢.
祂說要愛你們的仇敵.
為那迫害你們的禱告.
我們信神的人就覺得很難消化.
怎樣做.
怎樣搞.
其實這就說到.
完全的第二點.
就是愛的超然性.
天父的完全.
其實是體現了.
祂對所有人的普遍恩典.
什麼叫普遍恩典呢.
就是人人有份.
永不落空.
什麼叫人人有份.
永不落空.
就是第45節.
有一句這樣說.
因為他叫太陽照好人.
也照壞人.
降雨給異人.
也給不異的人.
這就是普遍的恩典.
人人有份 永不落空.
耶穌其實要求我們門徒.

$^{241}$突破當時的猶太文化.
當時的猶太文化就是第43節.
愛你的鄰舍.
恨你的仇敵.
這反而會局限了.
我們跟人與人之間的關係.
愛的人就繼續愛.
但不喜歡的人就繼續不喜歡.
有什麼分別呢.
沒有分別.
所以耶穌要傳遞的.
我們要活出一個超越世俗標準的愛.
要反映神的聖情.
這段經文.
我想說是救了我一命.
為什麼呢.
我原本是想著.
有仇報仇.
我以為是大義滅親.
我以為不是你死就我亡.
但當我要這樣做的時候.
神就出現了.
告訴我其實祂很愛我.
其實祂沒有離開過我.
其實祂知道我家裡的災情.
其實祂很願意陪我.
醫治我 憐憫我.
使我得到幫助.
所以我信了主.
信了主也不是馬上扭曲了內心的思想.
或者一些創傷可以馬上處理.
我們說災後要重建.
不只是重建建築物.
是重建心靈.
重建.
有些創傷不是那一刻.
有些創傷是你做夢會出來.
或者在你之後的人生不同的場景會出來.
然後你就要面對那些痛.
那些創傷什麼時候才會成為疤痕不痛呢?.

$^{281}$真的每個人都不痛.
因為你要知道究竟他有沒有繼續受創傷.
還是他已經離開創傷.
也夠膽被醫治.
以致那個不再是一個傷口.
而只是一個疤痕.
這些全部都是要過程.
我們說我們要像上帝的愛.
就是我們願不願意去明白.
他這一份恨久忍耐的愛.
說比做簡單.
做或者持續去做.
這才是最艱難.
所以其實這個愛的超然性.
是在說這種愛不是情感上的偏好.
我選擇我愛你還是我恨你.
而是意志上的選擇.
即是面對敵對者.
都仍然選擇我要以善去回應惡.
我信了耶穌之後.
是不是就不用和我爸爸一起住呢?.
不是.
我信了耶穌之後.
都是要對著惡魔.
我形容他就是走火入魔.
當你和一個走火入魔的人同住的時候.
我稱之為20年的長征.
特別有些年份覺得很特別.
我數算一下我爸爸何時再極致一點.
那個走火入魔的狀態.
去到92年走火入魔.
20年的長征.
即是去到多少年?.
12年.
12年他終於安息主懷.
對我來說.
20年的痛苦歲月.
終於劃上了一個句號.
但他走了.
代不代表我的創傷就沒有了呢?.

$^{321}$不會的.
我記得我上次的結尾有說.
對我來說最不容易的就是親子的關係.
原來當我很想教好我的小朋友.
但我每次一生氣.
一著火.
就不小心用我爸爸對我的說話.
語氣和神態.
就走出來了.
原來我也有些什麼?.
不是罪性.
我也有魔性.
但我想不想?.
我不想的.
求主幫我們.
但我想先說.
原來我怎樣經歷這些.
被神的話語去塗灶.
我很想激勵免麗丁姐妹.
我們大家一起.
以神為我們的榜樣.
我們怎樣可以靠主去完全.
了解上帝的屬性和心態.
或者期待我們這群信耶穌的.
要成為怎樣的人.
其實第45節當他要說.
他的普遍恩典是.
他叫太陽照好人也照壞人.
他幹預彼二人也給不二人的時候.
他上一節是怎樣說的?.
我們看45節那段經文.
第一句我們一起讀.
這樣你們就可以作天父的兒女了.
(天父的兒女).
所以原來神很想我們成為他的真兒女.
要成為他的真兒女就要像他.
而我們的天父就是完全.
他也很呼召我們.
要成為好像天父完全一樣的兒女.
所以他給了我們一個挑戰.

$^{361}$要成為生命成熟,品格完整.
我們要超越世俗標準的外.
我們要反映神的性情.
我們要一起去效法神的榜樣.
他那裡說.
像天父完全一樣.
在48節的尾.
所以你們要完全.
如同你們的天父完全一樣.
如果我們教子女.
就是你像爸爸的好.
你像媽媽的好.
你把你爸爸媽媽的好都學了.
有很多人都說.
我的女兒很像我.
因為完全不害羞.
很能見大場面.
又做小天使去幫忙.
我都很記得.
我帶小孫隊也有小朋友一起來.
有很多小朋友一起來.
其實都是.
不少不少反映父母的特質.
沒有害怕,很淡定.
社交能力很高,很強.
或者情緒很穩定.
這些全部都是像爸爸媽媽.
我們天父很想我們每一個兒女像他.
但此去學習進程是不簡單的.
我們要持續追求這個目標.
就是我們每一個門徒.
其實蒙召在這個世上.
要活出神國度的價值觀.
透過實踐無條件的愛.
逐漸跟神的旨意可以對齊.
我那時說.
有一段經文救了我.
其實我那時初信主.
我已經很熱心傳福音.
未學傳福音.

$^{401}$我就已經傳福音.
但我裡面還有一件扭曲的事.
未完.
我一定要將福音傳給我所有家人.
除了一個.
除了誰?.
因為天父死有餘辜.
不要讓他上天堂.
要讓他落地獄.
那時我已經信了耶穌.
初信.
很熱心傳福音.
但原來我這個位置未完成.
因為原來我還在生氣.
我還在傷.
他還在惡意對待我們家人.
所以很難.
這個結很難解.
但後來上帝不斷用這段經文.
勸導我.
說服我.
幫助我.
捍衛我.
如果你跟你爸爸一樣.
仍然繼續受恨.
其實你像你爸爸.
而不是你像天父.
當頭大汗.
聖靈提醒我.
你要像天父還是像你爸爸.
我懂不懂選擇?.
我當然不想像走火入魔的爸爸.
我當然想像天父.
那你想像我嗎?.
那你想像我.
你不要恨他.
你饒恕他.
你愛他.
你幫他.
怎麼幫?.

$^{441}$可能每天你都原諒了很多次.
第二天又來了.
所以我很記得這段經文.
救了我很多次.
每次我都說.
上帝啊.
我不要再去愛這個人.
我愛不了.
我受不了他.
我不行了.
但耶穌說.
女兒,我知道.
女兒,你為我做的事,我明白.
你的眼淚,我陪伴你.
我知道.
但你可以的.
你繼續去愛他.
你繼續幫他.
但上帝你要繼續給我你的話語.
去扶持我.
以至我可以繼續去打這場仗.
很奇妙.
上帝給了我另一段經文.
我想大家都有讀過這段經文.
而剛才有首詩歌.
也有提及過這個場景.
就是耶穌釘十字架的時候.
我很想大家可以閉上眼睛.
我們一起閉上眼睛.
我用我的聲音.
跟大家進入一些場景.
看看耶穌如何成為我們的榜樣.
他昔日是如何和我們同行.
而裡面可能出現過很多不同的人.
我們也看看聖靈.
讓我們更加被哪些狀態.
哪些人與事去觸動我們.
我們就求聖靈幫助我們.
我們一起閉上眼睛.
耶穌被釘十字架.

$^{481}$在路加福音23章記述.
他們把耶穌帶去的時候.
有一個古尼來人西門從鄉下來.
他們就拿著他.
把十字架擱在他身上.
叫他背著跟在耶穌後面.
有許多百姓跟隨耶穌.
其中有些婦女為他號嘯痛哭.
耶穌轉身對他們說.
耶路撒冷的女子不要為我哭.
要為你們自己和你們的兒女哭.
因為日子將到.
人要說不生育的.
未曾懷孕的.
和未曾與補孩子的有福了.
那時人要向大山說.
倒在我們身上.
向小山說.
遮蓋我們.
他們若在樹木青綠的時候做這事.
那麼在枯乾的時候將會怎麼樣呢?.
另外有兩個犯人.
也被帶來和耶穌一同處死.
到了一個地方.
名叫竹樓地.
他們就在那裡把耶穌釘在十字架上.
又釘了兩個犯人.
一個在右邊.
一個在左邊.
這時耶穌說.
父啊赦免他們.
因為他們所做的.
他們不知道.
士兵就抽籤分他的衣服.
官長也嘲笑他.
說他救了別人.
他若是基督.
是神所揀選的.
救救他自己吧.
士兵也戲弄他.

$^{521}$上前那處送給他學.
說你若是猶太人的王.
救救你自己吧.
在耶穌上方有一個牌子寫著.
這是猶太人的王.
同釘的犯人中.
有一個譏笑說.
你不是基督嗎?.
救救你自己和我們吧.
另一個應聲責備他說.
你是一樣受刑的.
還不怕神嗎?.
我們是應當的.
因為我們是自作自受.
但這個人沒有做過一件不對的事.
他對耶穌說.
耶穌啊.
你進入你國的時候.
求你紀念我.
耶穌對他說.
我實在告訴你.
今天你要和我再下元旅了.
大家可以張開眼睛.
當你聽我讀經文的場景.
你有沒有想過.
當耶穌基督的普遍恩典.
是為全人類釘十字架.
為任何災情的人釘十字架.
甚至為那些作惡的人釘十字架.
其實這些全部都是很.
盜慾的.
很貼身的.
每個人都有份.
都是耶穌為我們每一個釘十字架.
但耶穌受苦的時候.
他說了一句絕句.
父啊.
赦免他們.
因為他們所做的.
他們不知道.

$^{561}$他們不知道什麼呢.
他們不知道耶穌是神.
竟然釘神的兒子.
他們不知道自己的罪.
他們不知道自己有些被罪.
弄得走火入魔.
有些宗教領袖已經成為.
以色列人很多百姓的壓制者.
被捆綁.
他們不知道.
所以每次神幫我去想.
其實你爸爸所做的.
是怎樣的.
他不知道.
所以每一次.
耶穌就跟我說.
你要完全像我一樣.
你可不可以原諒你爸爸.
我也願意赦免你爸爸.
他所做的.
他不知道.
你愛他.
你幫他.
你原諒他.
你一起去救他.
其實在這段經文.
我剛才說了.
其實是耶穌登山寶訓.
一路一路的講論.
耶穌在第五至七章.
是重新詮釋律法的正義.
48節的完全.
亦是視為對前民教導.
例如不要發怒.
不要報復.
要忠誠婚姻等等.
作為一個總結.
要求門徒整體性.
活出神國子民的樣式.
所以我說長征了二十年.

$^{601}$但其實.
頭五年我還未信主.
後面十五年我已經信了主.
然後就用不一樣的方式.
去打理場象.
以前的爭鬥可能是.
你一言我一語.
大家一起在戰場裡.
但後面的時候.
神幫我的戰略就不同了.
我要用寬恕.
用愛.
用明白.
去體會我爸爸.
為何他今時今日.
會變成這樣的惡魔.
如果大家有聽過.
創傷知情的理論.
其實有時是說.
究竟這個人發生了甚麼事.
以致他會變成這樣.
其實是關心他.
而不是只問他.
你做了甚麼事.
有時我們人都不想活得那麼老.
或者活得不能自制.
但原來有時我們因著.
很多處境擺在一起.
可能他都是一個很受創的人.
他不能用其他方法.
去處理創傷.
然後就延續創傷給下一代.
因為我覺得後來不斷成長.
我自己都有進步.
亦學聖經.
亦被教會的傳道人.
去醫治牧羊.
去療傷.
是很不容易的.
去慢慢認識我爸爸.

$^{641}$因為我爸爸在大陸.
所有他的兄弟姊妹.
都沒有偷渡來香港.
其他都繼續留在國內.
但原來他們的際遇.
和我爸爸完全不一樣.
所以可能在這些際遇不同.
原來都會消磨了一個人.
或者是腐化了一個人.
就要幫他怎樣可以歸根究底.
或者尋遠.
怎樣可以令他處理那些結.
要去解結.
原來長話短說.
我後來理解.
我爸爸都有很多創傷.
但一個男人.
很少會告訴別人.
我受創傷.
很多男人最痛.
都是放在自己心裡.
或者是放在自己屋裡.
我看到我爸爸那種表裡不一致.
他在屋裡面.
關上門是另一個他.
但他出門的他.
是另一個他.
我們小時候就覺得.
為什麼家裡的爸爸.
和出門的爸爸.
是兩個樣子呢.
原來我爸爸很努力.
將所有傷勢都隱藏.
但最容易發出的傷勢.
就是在家裡.
但在外面.
都是很努力去保護自己的傷.
不讓別人觸碰.
所以很展現我沒事.
我沒事.

$^{681}$我很好.
所以如果我們說.
我們要去幫這些受傷的人.
有沒有樣子看的.
有些是沒有樣子看的.
就算我那時候.
我是年青時期.
最受家裡的影響.
我上學的樣子照常.
我繼續對人歡笑.
不過後面那幾個字是什麼.
背人仇.
我對人就是笑.
很開心的.
沒事的.
沒有人在我面前.
我就自己傷心.
是這樣的.
所以有些人能夠坦露.
告訴你他的情況.
這些已經是很好了.
你真的幫到他.
但大部分的受傷者.
他不一定要展露給你看.
因為他太沉重.
他都很害怕.
他表露出來的時候.
你能不能處理得到.
你聽完之後.
你會不會還很震驚.
我真的很深刻記得.
那時候中學的學生.
開始看報紙.
有個同學在旁邊拿起報紙.
他就說.
為什麼這麼殘忍的.
家裡發生命案.
可以殺了家人.
為什麼會這樣.
我心裡怎麼說.

$^{721}$其實坐在你旁邊的.
都打算這樣做.
有什麼奇怪的.
不過你告訴他嗎.
所以我們說.
我們每個人都盡量好樣回來見大家.
但真的有多傷.
或者我們要去服侍我們的社區.
我們很感恩.
聽到曾牧師說.
有十幾位弟兄姊妹.
一起去讀社關課程.
其實我們就要去學這些東西.
原來有很多東西.
我們真的要關心到入獄.
你一定要聽一下.
原來有些家庭是可以很爛的.
有些破碎家庭.
真的可以很傷.
有些創傷的後遺症是可以持續的.
我們怎樣可以陪一個人.
或者孕育一個人.
可以離開.
是不能完全離開.
不過怎樣可以對自己好一點.
或者對下一代再好一點.
其實我都是不斷成長當中.
但是我很感激的是什麼呢.
我很感激我信了主之後.
我一直在教會這個氛圍成長.
我自己還沒學會的.
我學一下別人的父母怎樣對子女說話.
子女會開心一點.
或者父母的樣子不一定要那麼憎恨.
可以持長一點.
你說他都可以笑著說.
這些全部都要學的.
我覺得我最艱難的是.
什麼高峰時期.
最不容易呢.

$^{761}$這個親子關係.
就是測驗考試.
對我們來說測驗考試.
要跟小朋友同行是最艱難的.
其實我們很想他好.
不過裡面有很多東西還沒完全調教得到.
有時候女兒都會哭著.
你不要那麼大聲.
你可不可以不要那麼兇.
這些自己要調整.
所以求主幫助我們.
完全在這裡的核心意義.
是我們以天父的愛為標準.
去追求生命和品格的成熟.
特別在愛仇敵的實踐上.
更加顯明神兒女的身份.
所以我們問問自己.
我們是不是像神呢.
我們像神有多少呢.
有時候我們很流行.
有些小朋友出生.
就像爸爸一樣.
特別像媽媽.
但如果其他人.
未信主的人見到我們的時候.
有多覺得我們像上帝呢.
這個不是靠我們自我的努力.
可以達成的.
而是要藉著跟基督聯合.
信服聖靈.
而是一個成長的過程.
要求神幫我們.
你們紅恩堂坐立在洪水橋區.
我們怎樣去愛這裡的居民呢.
其實我覺得你們附近也有一個災場.
馬會是不是一個災場呢.
可不可以是一個災場呢.
我們要想想.
上帝說太陽照好人.
也照壞人.

$^{801}$降雨給異人也給不異的人.
我們看到.
舉例說一些災的狀態.
我不說那麼近.
因為有時候那麼近發生的事情.
有很多事情還沒消化好.
有些複雜的情緒.
我講遠一點.
汶川大地震那時候.
其實我不知道大家有沒有去過那個遺址.
他真的特意保留了所有倒塌的房子.
但他沒有危險.
因為每個人都支持著他.
但成為一幅當時的模樣.
很想這個位置給其他人去參觀.
是一個借鑒.
是一個警惕.
也說的是不去忘記這個創傷.
其實有時候我們說.
都過去了,你忘記吧.
其實有些事情不是用忘記這個方法.
反而是用接受這個方法.
但要接受就不是一下子.
就是要長時間重建和復甦.
我們說離開的人是很不幸.
但我很想說.
最不容易的是什麼.
在這裡的人.
最不容易的是還在這裡的人.
最不容易的是還在生存的人.
最不容易的是那些在醫院.
不知道能不能怎樣醫治.
或者他之後能否復康.
能否回復以前的.
很不容易.
但那時候有什麼暴露了出來呢.
就是很有名的豆腐渣工程.
但有時候就是這個情況彰顯了.
其他樓房都要去監管.
不要容讓有第二次地震再來的時候.

$^{841}$再發生同樣的情況.
所以你可以說.
這個天災,但也包括人禍.
但經過了四方八面.
當時如果你有印象.
香港也出動了很多人.
去救災.
我很想說.
我也很感恩.
這一兩年.
我真的有機會去過那附近.
原來那附近的教會很興旺.
三自教會很興旺.
為什麼會這麼興旺呢.
他們說那時候.
很多香港熱心的人.
全部湧到那裡幫忙鎮災.
也有很多信主的人.
所以在那裡特別幫助那裡的三自教會.
做很多災後重建的事.
所以那裡是很精彩.
你也看到那城市的規劃.
越來越有規模.
我們在那次重大災情之後.
我們發現好多了.
再沒有這麼嚴峻的事發生.
但可能在其他城市裡.
未處理的時候就發生這件事.
但有時候我們也會很弔詭地去想.
有時候我們說.
捉不到的人是不是那個人要去背負責任呢.
捉不到的人是不是代表他不用去背負責任呢.
我講得比較俗一點.
去承擔那個責任.
上帝令我想起一段經文.
大家有沒有想起.
聖經有單案.
叫做行淫時被拿的婦人.
行淫可以一個人做嗎.
不可以 但為什麼只捉到一個女人呢.

$^{881}$那個男人呢.
是不是.
但耶穌怎樣跟被捉到的人說.
你以後不要再做了.
我也不去拿石頭扔你.
我也不去定你的罪了.
耶穌也叫其他人去定別人的罪.
但耶穌說了.
究竟我們有沒有資格呢.
我們沒有資格去定別人的罪.
因為我們每一個人都有罪.
所以有時候我們在這些處境的時候.
求神幫助我們.
我們要將所有審判的權柄.
最終交給上帝.
當然神在地上設立律法.
然後每個國家不同的城市.
都有不同的法制去規管.
我們都很尊重這些.
因為大家互相保障.
但你說是不是真的最終那個.
有些是捉到 有些是捉不到.
譬如我舉個最實例.
其實我爸爸走火入魔.
其實我們也有考慮過報不報警.
我可以告他虐兒 可以告他虐妻.
我們有沒有報警.
沒有 為什麼我們不報.
都是想給他機會.
所以有時候有些關口是很難 很複雜的.
所以我們說我們要去做社會關懷.
要做社區關懷.
其實是每一個家庭都有每一個故事.
家家有本難圓的經.
但我們說聖經是我們最好的一個標準.
我們怎樣去用上帝的標準.
去超越我們所有不同的限制和框架.
所以我很感恩.
我說這段經文救了我.
我也很想這段經文可以救到大家.

$^{921}$或者大家也去幫忙用這段經文.
去救其他社區的人.
昔日如果不是這段經文救了我.
可能我就是其中一個囚犯.
但這段經文救了我.
我真的竭力地跪下祈禱.
上帝你再給我力量去愛我爸爸.
但我又要懂得保護自己.
因為不可以再被他傷害.
所以祈禱又要懂得保護自己.
然後到最後就是他一場大病的時候.
趁他病 救他命.
記住坊間有很多俗語術語.
但那些就是人性的標準.
抵他死 罪有餘辜.
應該打入十八層地獄.
但聖經完全不是這一套.
趁他病 救他命.
當他最無力的時候.
你還去探望他.
當他最照顧不了的時候.
你照顧他.
然後你告訴他.
我記得那時候我含淚.
經常跟爸爸說.
爸爸說真的.
如果不是我真的相信耶穌.
我不會來探望你的.
如果不是我真的相信耶穌.
我不會這樣愛你的.
你做過什麼 你自己知道吧.
求聖靈提醒他 光照他.
誰知道真的持續這樣做.
然後到最後他願意悔改 信主.
我不是預先信主才做.
我只是說我要像天父.
我不是看結果而去做.
我要像你 所以我要去做.
結果交給誰.
交給主 亦交給他自己去負責.

$^{961}$因為這麼多神的愛.
這麼多基督徒的愛都在你身邊.
但我們說神和基督徒.
都不是要逼人信主.
其實都是讓人選擇.
但你會堅持流露上帝那份普遍的恩典.
跟他們去同行.
所以我剛才舉例.
賽馬會海是一個災情.
你想想 澳門有很多災情.
我們說要過大海 坐金巴 坐船才能去.
澳門過隔壁馬路 過隔壁街就有了.
有多少不容易的家庭.
所以我們說我們能否看到.
有些災難是立即看到的.
但有一種災叫「溫水煮蛙」的災.
我們就要看看魔鬼在用什麼方法.
要潛移默化扭曲我們的人生.
扭曲我們的價值.
我們要將神的話語 將神的真理.
怎樣去幫助那些人.
求主與我們同行.
讓祂的話語調整我們 教導我們.
讓我們更像上帝 完全像祂.
我們同心同功.
慈悲仁愛的天父.
懇求你幫助我們.
我們要像你完全一樣.
這個功課不容易學.
但你懂得的是.
做好整個很好的版 很好的榜樣.
可能在這次災情的時候.
我們會想起我們有受害的居民.
受害的無辜喪生的人.
亦有殉職的人.
我們亦會想在這個世代裡.
其實耶穌為我們釘十字架.
你也是殉職嗎?.
你也是殉道嗎?.
你也是為了這麼大的救贖計劃.

$^{1001}$就是要救我們而犧牲自己嗎?.
所以懇求主幫助我們.
當我們仍然活著 在這個世代的時候.
讓我們像你 讓我們跟你的腳步去行.
讓我們可以以上帝你自己的普遍恩典.
可以與身邊的人同行.
甚至可能就是那些.
在這次事件裡.
可能要承擔一切法律責任的人.
其實可能他們都很需要耶穌基督的愛.
所以幫助我們.
為不只是災民禱告.
亦為那些要受法律責任的人.
我們都去祈禱.
或者有些是逃之夭夭的人.
主啊 我們都求你.
你親自審判 你親自管教.
你親自去做你奇妙的工作.
讓你的基督徒遍佈在這個世界裡.
都能夠與不同的人同行.
讓人知道有神你自己.
我們誠心禱告教徒.
奉主名求.
阿們.
\newpage



\section{腓立比書 2:21}
\label{sec:WQHTSm8b4dk}
\textbf{2025.12.07《求耶穌基督的事》:腓立比書2\_21 (和修版) 講員 : 謝靄薇牧師}
\newline
\newline
連結: \href{https://youtube.com/watch?v=WQHTSm8b4dk}{\texttt{https://youtube.com/watch?v=WQHTSm8b4dk}} ~~~~ 語音日期: 2025-12-10
\newline
\newline
\hyperref[sec:L0es3lIs6g8]{\small{< < < PREV SERMON < < <}}
~
\hyperref[sec:index_chronic]{\small{[返順時目]}}
~
\hyperref[sec:index_scriptual]{\small{[返順卷目]}}
~
\hyperref[sec:moCCv_ijPQc]{\small{> > > NEXT SERMON > > >}}
\newline
\newline
腓立比書 2:21
\newline
\begin{longtable}{cl}
\hline
\hline
章節 & 經文 (和合本修訂版)\\
\hline
2:21 & \begin{tabularx}{0.7\textwidth}{X} 其他的人都求自己的事，並不求耶穌基督的事。 \end{tabularx} \\ \\
[1ex]
\hline
\hline
\end{longtable}
$^{1}$各位頂智的平安.
很久沒見你們了.
上一次颱風也來不了.
今天很開心見到你們.
今天充滿陽光.
希望我們的心靈能夠打開開放.
願意神在當中保守我們.
讓我們在神的懷疑之前一起祈禱.
讓我們在這麼好的環境裡.
我們敬拜你.
給我們恩典.
求你的話語進入我們的心.
讓我們願意聽你的旨意.
以致我們所行的.
是行在你的心意當中.
求你保守.
求你帶領某樣孩子.
在當中所說的.
能夠蒙你月立.
願你讓我們眾人.
我們聽到的時候.
若然有不合理心意的.
求你剔除.
以致我們能夠尊心一意的.
只見耶穌.
奉主耶穌基督的名求.
我們看《主經文》的時候.
好像很簡單.
其他人都是求自己的事.
並不求耶穌基督的事.
我們今天就去看看.
如何求耶穌基督的事.
其實很簡單.
大家都在做的.
就是服侍主.
服侍主必須承受主的託付.
然後就會跟從主的腳蹤.
最後就發現.
你真正在做耶穌基督的事.
我們想起耶穌基督.

$^{41}$祂有三年半的時間.
跟門徒一起.
第一年祂可能教他.
如何做得人如膚.
然後進一步教他.
要寫記.
最後再進一步.
原來要為教會死.
這是一個不容易的進程.
當我們要去事奉主.
跟從主.
主首先讓門徒了解.
祂的事工.
跟祂的事工有份的.
是一個極榮耀的使命.
當門徒跟著主耶穌去傳講.
聽到主耶穌說.
「天國近了,你們應當悔改」.
之後主就呼召彼得跟從我.
要祂得人如得魚.
在馬科音或魯加福音.
都有同樣的記載.
主耶穌沿著加利利海邊行走.
看到兩兄弟.
彼得西門和弟弟.
本來是打魚的.
但耶穌說要跟從我.
他們很順利地跟從了耶穌.
原來他們跟從耶穌.
這麼快就願意捨棄亡.
跟從祂.
見一次面.
好像今天來教會.
見到弟兄姊妹.
這麼順利地放下.
自己很重要的求生工具.
回來教會事奉.
是不是這麼簡單呢?.
當我們仔細看四福音的時候.
會發現.

$^{81}$其實不是第一次見耶穌.
他們遇到耶穌.
差不多.
你看看耶穌開始工作.
差不多跟他們相近.
有一年的時間.
耶穌受了針.
工作了一年的時間.
他在加利利傳天國福音.
他呼召彼得做他的工作.
而在約翰福音.
也記載一些.
在馬太福音沒有記載的.
在約翰福音記載的.
是施洗約翰福音出現了.
他出來傳天國近了.
你們要悔改.
當時的人很渴望天國降臨.
猶太人被外邦人統治.
也有六七八年的時間.
北國在公元722年被亞述俘虜.
南國在公元前586年.
最後也沒有了.
被巴比倫滅亡了.
他們也是在外邦的地方.
被人俘虜了.
你看到這個時候.
施洗約翰出現的時候.
他說天國近了你們要悔改.
那些人不知道等了多久.
等彌賽亞來.
請教他們脫離苦難.
當施洗約翰一說的時候.
每個人都呼白應.
都願意悔改.
當然有些人等不及.
就變成了分裂黨.
他們就靠武力.
靠他們的生命去打拼.
要去推翻羅馬帝國.

$^{121}$當然他們的作風.
後來就影響著羅馬帝國.
他們的大軍毀滅了耶路撒冷.
另外有一批人.
他們走相反的路線.
他們不是用武力.
他們就是所謂的愛色尼人.
他們在抗疫的地方.
這些人就居住在死海附近.
所以你看死海古卷.
挖掘的地方就在那裡.
你又看到安德烈和彼得.
你不要小看他們.
雖然他們是打魚的.
他們都很盼望彌創亞快點來.
希望天國快點來.
所以他們雖然是打魚的漁民.
但是他們都很希望天國的降臨.
希望神的恩典臨到他們的百姓.
所以安德烈和彼得.
老早就聽了施拉約翰的信息.
跟著他.
可以說是他的徒弟.
好了.
到約翰福音.
你慢慢看的時候.
你會看到施拉約翰.
他出來做見證.
他見證什麼呢?.
他說這個在我以後來的.
是在我之前的.
為什麼呢?.
其實耶穌基督是神.
他還沒出生之前.
他已經在這裡了.
只不過是我們道成肉身.
所以他說.
我連和他解鞋帶都不配.
他看到這是真正的神的羔羊.
所以他去做見證.

$^{161}$當他做見證的時候.
很多人就認識到.
原來耶穌基督是真正的羔羊.
所以這個.
等一下.
退後一點.
幫我退後.
再退後.
施拉約翰去見證耶穌的時候.
安德烈和彼得馬上跟從主.
他聽到這是真正的羔羊.
所以他馬上就去跟從主.
他去跟從主的階段.
你會發現.
在聖經裡.
看到耶穌的婚宴.
就是迦納的婚宴.
他一路去看.
看到在撒瑪利亞的婦人.
耶穌和她傳福音.
看到很多很多的東西.
耶穌和她談道.
然後告訴她裝嫁耶穌.
要收裝嫁.
她一路跟著耶穌.
然後你又看到.
她最後走了一圈.
又回到加利利.
回到加利利做什麼呢?.
原來她回到加利利.
是從操故業.
她繼續打魚.
他們那時候.
很知道天國的事.
又很羨慕.
領這麼多人歸主.
不過他們跟了一段時間.
回到加利利.
又去打魚.
又去打魚的時候.

$^{201}$這個時候.
我們看到耶穌呼召他們.
因為有一天耶穌來到湖邊.
看到他們打魚.
耶穌叫他們將船駛開一點.
然後跟著他們跟著跟他們講信息.
講完信息之後.
耶穌就說.
你要船開到海深之處.
就怎樣?.
下網打魚.
彼得可能心想.
不用打了.
我們搞了一整晚.
整夜勞累.
打不到一條魚.
簡直是沒有辦法.
但耶穌叫他們.
儘管聽我的話.
你下網吧.
他們一下網.
看到神跡.
滿滿的魚.
差點網都裂了.
然後你又看到.
他們兩個.
跟從耶穌.
在神跡裡.
讓他們看到.
耶穌真的有能力.
他說那裡有就是有.
他立刻跪下.
說我是罪人.
我真的不配.
你離開我吧.
他發現神是這麼聖潔的.
但竟然鄙視了神的吩咐.
然後耶穌立刻.
跟從他.
呼召他.

$^{241}$我要叫你們.
德人與德魚一樣.
弟兄姊妹.
當你了解背景的時候.
你看到彼得,約翰,安德烈,雅各.
他們對於主來到地上的工作.
主說天國近了.
神的國臨到.
尼昌雅來到.
他們真的知道.
這是一個很榮耀的使命.
他們打魚打了一輩子.
他們知道打魚算不上什麼.
但最尊貴,最榮耀的就是.
能夠在耶穌的帶領下.
指揮之下.
能夠是.
德人與德魚一樣.
如果靠我們自己就不行.
靠主耶穌說可以就可以.
所以他們看到這是很崇高的使命.
我們有沒有看到呢.
他們看到了世上很多人流離失所.
裝假熟了.
但沒有人去修國.
但主這樣呼召彼得,雅各,約翰.
安德烈的時候.
他們認清了.
能夠為基督做事.
求基督的事.
能夠與他的事公有份.
帶領人歸向耶穌基督.
是極為榮耀的職份.
但不一定每個基督徒.
都像保羅一樣.
放下職業.
專心做全途的工作.
記住.
無論你做什麼.
哪個行業,哪個職業.

$^{281}$但我都要盡心去做.
將最好的獻給主.
其實任何職業.
都不是你的正職.
我們的正職反而是.
在任何環境裡.
無論在哪個崗位.
在哪個機構工作.
你要成為一個得人如膚.
主託付你.
一個重要的使命.
就是將天國這麼偉大的事業.
帶到地上.
所以天國近了.
你們要悔改.
這是主耶穌一路傳的訊息.
拓展天國的事業.
就是藉著你的工作.
將主的託付傳開.
我們看到.
每個基督徒.
無論從事什麼行業.
記住你的正職.
就是耶穌基督.
這個天國的事業.
是藉著你.
才可以推廣.
因為你是順主耶穌.
你是基督徒.
你就要藉著你的身份去傳開.
現在你所做的工作.
只不過是副業.
藉著這個工作.
將主的託付傳開的時候.
弟兄姊妹.
你一方面做正職.
一方面做副業.
有一天.
如果主給你的負擔越來越重.
由聖經的話語.

$^{321}$神的話語提醒你.
你的心跳得很快.
越來越感動你.
神的呼聲說.
「裝教很多,但工人很少」.
求主打發工人.
去求他的裝教.
神在天上說.
「我可以拆遣誰呢?」.
「誰肯為我們去呢?」.
當聖經有很多很多的話語.
在你的心中越來越成為衝擊的時候.
神去引導你.
神在教會的全道同工裡.
阿都越來越清晰地.
和你聊天的時候.
讓你看到這條路.
那時候你就不得不.
放下現在的副業.
然後跟從主.
去從事正業.
今天我們為世上.
我們花很多時間.
是應該的.
因為你在地上.
無論你做什麼職業.
我們都要忠心.
我們要盡心竭力地去做.
但有一天.
神感動你.
神的話語打動你的時候.
去激勵你.
如果在教會裡工作的時候.
再加上你跟全道同工聊天.
他們的印證之下.
你到了那個程度.
自然就會將你的副業放下.
你會專心.
以祈禱全道為事.
這是一個更榮耀的積分.

$^{361}$主將全部的責任.
交給基督徒.
我們看到很多例子.
這是什麼例子?.
這是埃及阿伯.
就是埃塞俄比亞那個地方.
那麼遠的.
有一個人.
他竟然去到耶路撒冷.
去敬拜神.
去尋求神.
一個那麼有權位的太監.
他是掌管錢的.
在他的國家裡.
掌管銀庫的.
在財政大城.
但是這個人.
他竟然來到聖殿.
看到那麼宏偉的建築物.
看到祭司獻祭.
看到拉比穿著服裝.
但是看不到主.
他就回去了.
一路上.
他在讀以賽瓦書.
但是他不明白.
主著急.
就派了天使去.
各位大姐姐妹.
你覺得天使是不是直接.
跟這個太監.
說明聖經呢?.
是不是?.
派了天使直接去找他.
是不是呢?.
你說是不是呢?.
他去找誰?.
他去找肥力.
找肥力跟他說.
有個人.

$^{401}$他在讀以賽瓦書.
他感動他.
你要跟廣黎.
下加薩的路.
怎樣去找到那個人.
但肥力.
可能像今天禮拜.
每個人都很多事情要做.
看著時風.
可能下午還有事情.
叫你去.
沒空.
去那麼遠.
我這裡很多事情.
可能下午還要.
傳呼音.
很多事情.
但肥力就放下.
他手上的東西.
憑著聖靈的感動.
他就去了.
向南走.
他去到那個地方.
聖靈就催迫他.
要靠近車.
我不知道車子走多慢.
他可以靠近.
還上了車.
去跟他講解聖經.
讓他明白.
基督徒.
在座的基督徒.
有沒有留意到.
有時神的事公.
就是那麼突然.
臨到.
但我們可能.
沒怎麼留意.
沒怎麼留意就錯過機會.
你看到.

$^{441}$天使去催促.
神透過使者去催促.
肥力向南走.
就得著了.
他跟他傳呼音.
原來以賽瓦書.
講到耶穌.
他明白了.
另一個人.
另一個人其實是.
一個白夫長.
是哥尼劉.
也是.
天使跟他說了一段時間.
但不是跟他傳呼音.
只不過告訴他.
你的州制.
神已經看到了.
但天使跟他說.
你去找彼得.
叫他去找彼得.
彼得是誰?.
彼得這個時候.
也是雲遊外將.
他在天板裡.
可能那時很餓.
但他看到一個異象.
看到一塊布掉下來.
但裡面的食物.
是他從小到大都不會吃的.
因為猶太人很清楚.
神的教導.
哪些是潔淨的.
哪些是不潔淨的.
很清楚的.
他看到那些食物.
平時是不會吃的.
但有幾次叫他.
你摘下來吃吧.
連續三次.

$^{481}$很奇怪.
但他終於明白.
是神的吩咐.
當他明白的時候.
這些人就拍門了.
這些人是甚麼?.
就是哥尼劉派來的人.
就是要找彼得.
這些是外邦人.
原來神看為潔淨的.
就是潔淨的.
所以這班人.
可以藉著彼得.
去到哥尼劉當中.
有機會傳福音.
看到天使.
只不過是.
他的工作.
是為了承受.
救恩的人的效力.
但真正傳福音.
就是你們每一個基督徒.
神要你去告訴別人.
得不得人就是主.
主給你得人.
你就得到了.
不過門徒跟來跟去.
耶穌也不滿意.
因為門徒.
好像明白一點.
又好像不明白一點.
他知道這是一個.
很榮耀的即是.
但他們沒有依靠主.
也沒有信服主.
而是靠自己的意思.
所以我們看到聖經裡.
門徒在爭誰最大.
在爭耶穌左右.
每個人都不是想.

$^{521}$為耶穌做事.
而是想權位.
到底為耶穌做事.
要完成什麼事工呢?.
其實彼得當時.
也很模糊.
不太清楚.
我坐在你左右.
就在做榮耀的事.
還有什麼呢?.
但是主馬上就說.
他要上耶路撒冷.
他要受祭司,文士.
很多的苦要被受害.
他馬上說了幾次.
他大三天會復活.
但是彼得說什麼呢?.
千萬不要.
我們要跟你在榮耀裡.
你怎樣去受苦呢?.
他還是不明白.
他這樣說的.
你怎樣去上耶路撒冷?.
他受人羞辱.
他阻止他.
但是耶穌就知道他不明白.
他知道這是來自撒旦的.
聲音去教他的.
所以就在攔阻.
主做一些很偉大的事.
這是走十字架的道路.
是不容易的.
但是撒旦的聲音.
就去教彼得.
這是受苦的道路.
我們不應該走受羞辱的道路.
但是耶穌就是要這麼艱辛地.
走這個十字架的道路.
才可以帶來救恩.
才可以拯救我們.

$^{561}$所以.
我們背起十字架來跟從他.
如果有人要跟從我.
就當寫記.
只是教他們要寫記.
要背十字架去跟從他.
每個人都有不同的十字架.
我們不要羨慕別人的十字架.
每個人都有神給你的十字架.
是適合你去背上的.
但是.
我們不要.
體貼人的意思.
我們要了解神的意思.
了解人的意思.
我們就以自我為中心.
所以最後.
你看到他去哥尼納.
帶了全家都信主.
這是很奇妙的事.
其實我們帶了一個已經很厲害.
還帶了全家.
當主耶穌.
要復活升天.
向他們顯現.
復活升天之前.
有四十天在地上.
他不斷地向這些門徒顯現.
他向他們顯現的時候.
主有三次問彼得.
「你愛我比這些更深嗎?」.
主要讓彼得知道.
如果他的工作.
跟他的牧羊有份.
是出自愛主的心態.
要愛主的羊群.
這樣擺上去.
就跟基督的事有份.
跟基督的事有關.
因為你要餵養主的群羊.

$^{601}$要牧養牠的群羊.
當他想到.
主是一個好牧人.
好牧人餵羊寫命.
但他自己.
他自己也是一隻羊.
但他自己卻是很敗壞的羊.
他三次不認主.
覺得自己很背逆.
這麼愛他為他寫命的主.
他竟然說不認識他.
還發愁起誓.
但這個時候.
當耶穌三次問他.
「你愛我比這些更深嗎?」.
再次提醒他.
他想他自己是這麼背逆的.
但主耶穌沒有放棄他.
主耶穌還交代瑪利亞.
去找他.
跟他說耶穌已經復活了.
去提醒他.
提醒他的名字.
但這個時候他很軟弱.
他去打魚.
他帶了六個人去打魚.
整晚都打不到什麼.
但主在岸邊看著他們.
主看著他們做什麼呢?.
他們還在打魚.
但那時候他們的心情.
不是很好.
因為又不成功.
但主耶穌又叫他們.
在岸邊失望.
結果聖經說他們釣上來.
有153條魚.
然後他們看到.
耶穌為他們預備了早餐.
有魚.

$^{641}$看到主耶穌這麼愛他們.
主耶穌看到他們的面.
沒有責備他們.
「你為什麼不理我?」.
「為什麼不認識我?」.
沒有說這些.
他只是說.
「你愛我比這些更深嗎?」.
然後當他去問他的時候.
他說「我要把工作給你」.
「我要託付你」.
可想而知.
再一次體會到.
主多麼的愛他們.
沒有責備他們.
主甘願為他們寫命.
主問他們.
「你愛我比這些更深嗎?」.
連續三次的問答.
然後主就把羊.
託付給彼得.
交給他了.
他真正的愛主.
他願意承擔這個工作.
看到在他的位身裡.
這個路程.
發現他長大了.
他開始知道.
走主的道路.
不是一個容易的路.
但耶穌這麼愛他.
他也學習.
把耶穌那份愛愛人.
去牧養他的群若.
這樣就和主的事供有關.
但主耶穌又帶來一個經文.
很要命的.
「我實實在在地告訴你」.
「你年輕的時候」.
「自己戳上帶子」.

$^{681}$「隨意往來」.
「但年老的時候」.
「你要伸出手來」.
「別人要把你戳上」.
「帶你到不願意去的地方」.
耶穌說這句話.
就是指彼得會怎樣死.
來榮耀上帝.
他說了這句話之後.
耶穌對他說.
「你跟從我吧」.
將來會被帶到這樣的地步.
如果你是彼得.
會有什麼反應呢?.
你還願意走嗎?.
你走的路是死的.
早晚都會碰到這條路.
怎麼辦呢?.
不如我們現在拒絕了託付.
我們避開萬年.
不用被人綁著手.
綁手綁腳.
沒有自由.
但他完全沒有.
他懂得主的意思.
他願意跟隨主.
這是一個真正完全的.
有愛的擺相.
我們去事奉.
都要有愛.
我們做.
不論是接待.
或試探訪.
都要看到這個人.
是神所愛的.
你要愛他.
你才可以做到.
如果我們沒有愛心.
這個星期又排到我去事奉了.
又是我了.

$^{721}$連續幾個星期.
你就會不太願意.
因為沒有愛心.
但當你看到.
耶穌那麼愛你的時候.
你願意走前.
爭著去做這個事奉.
這個事奉不是每個人都可以做.
神說你可以才可以.
但很多時候我們.
不得不去.
覺得有很多時間.
這個星期不做.
總有人做.
但我們不知道.
我們的時間是多少.
你只是一個財主.
今天就要取你的命.
你就不能事奉了.
你雖然很盼望.
但當神叫你的時候.
你就快點去做.
這是很重要的.
給他,他懂得.
什麼是事奉主的道路.
所以無論耶穌剛才說了什麼.
他沒有反駁.
他很願意.
但他晚年.
他怎樣?.
他真的為主被釘十字架.
要倒釘.
一個完全真正的.
願意將愛擺上生命.
主那麼愛他.
他用愛去回應主.
今天我們都沒有受什麼苦.
是不是?.
可能你的苦就是你的子女.
或者你的工作.

$^{761}$其實我們沒有像耶穌那樣.
流血的苦.
但受苦的原則.
是我們每一個基督徒.
我們走上這條路的時候.
我們都要預備受苦的心智.
這是耶穌交代給我們的.
所以我們和主的事工有份.
要去服侍主.
我們必須懂得.
耶穌基督的十字架路.
你肯不肯走?.
你今天聽了之後.
最後是要為教會去死的.
你肯不肯呢?.
主被捉的時候.
他的門徒都走了.
彼得還奏了.
佛奏起了.
我不認識這個人的.
但後來.
很大的轉變.
很放膽的.
如意跟從主.
你看了幾次.
他們都被捉了.
但又看到很神奇的.
天使又不知道怎樣.
又帶了他出來.
很神奇的.
你看《福音》裡面.
他就這樣走上這條道路.
相信大家生.
大家都認識.
他十七歲就被神緊抱.
他聽到中國報道那裡沒有福音.
雖然馬里遜都去過.
但內陸仍然很多人沒有福音.
當時的統計.
一年有1200萬人去世.

$^{801}$一個月有100萬人去世.
每一星期有25萬人去世.
他聽到中國沒有福音.
沒有神的盼望.
他就對主說.
主啊,我相信你.
我自己都有福音.
已經得到福音的好處了.
他就對主說.
我要做一個宣教士.
我要去中國為你傳福音.
他的心智一直是慢慢的培養.
到他十九歲的時候.
他立志要學習怎樣去中國傳福音.
所以他專門去住一些很簡陋的地方.
預備受苦.
因為中國的地方.
當時的經濟沒有現在那麼蓬勃.
環境都差現在很多.
他要預備受苦的心智.
就去專門去找一些比較簡樸的地方.
學習怎樣去.
比較困苦艱難的地方去住.
希望自己在生活程度較低的地方.
他都可以生存.
去培養自己.
才有這樣的基礎.
去操練自己.
他有去學醫的.
他十九歲那年.
他妹妹帶了一個同學來.
這個同學很漂亮.
彈了一首好的琴.
他們兩個一見鍾情.
他們經常約會.
但他忽然間就想.
萬一這個女孩知道我要去中國宣教.
他會不會不肯呢?.
那怎麼辦呢?.
他想到這個問題.

$^{841}$於是他就跟這個女孩說.
這個女孩跟他爸爸說.
你知道他爸爸的反應嗎?.
他爸爸是父親.
他愛惜他的兒女.
怎會讓他去中國.
當時很艱難的地方.
不像現在強國.
以前不是.
他爸爸一聽到就火冒身上.
然後就告訴他.
你要好好考慮.
你要中國就不要玩我女兒.
於是他就回家.
日記上寫.
他真的不知道如何祈禱.
到底祈禱是要改變神的心意.
還是改變我的心意呢?.
他每次去祈禱都不知道要說什麼.
他真的不懂祈禱.
他覺得很痛苦.
最後他就跟主說.
怎麼辦呢?.
我們怎樣去回答他父親呢?.
他寫他很痛苦.
他不知道要怎樣祈禱.
他從早上到下午.
都不知道該怎麼辦.
他有很大的掙扎.
在那個時候聖經的話語就進入他心裡.
就是保羅所說的.
我已經凋棄萬事.
看作糞土.
在《肥臘比書》三章八節.
說要得著基督.
他要認識耶穌為至寶.
他在日記上寫.
主,還有誰比你更加寶貴呢?.
他這樣寫的.
他說,我不用祈禱了.

$^{881}$我已經有答案了.
於是他就找了他女朋友的爸爸.
他跟她說.
他說,中國我要.
你的女兒我也要.
她爸爸說,不要這回事.
怎麼可能呢?.
於是.
從那天開始.
他就沒有再見這個女朋友了.
你看到就是.
神都沒有虧待他.
在21歲.
他準備去中國了.
他的父母.
還有教會的弟妹.
就會送行.
但是他的傳奇就延誤了.
他爸爸有工作.
於是就留下媽媽.
陪伴他.
他的媽媽.
幫他安裝床鋪.
放好行李.
不敢流一滴眼淚.
她很怕影響她的兒子.
那天跟他祈禱,唱詩,讀經.
然後.
他喜歡下船.
他下船到岸邊的時候.
他看到.
在船裡面.
只有一個客人.
就是他的兒子.
其他都是在水上,沒有別人.
他很傷心.
怎麼辦呢?.
將來去中國的地方.
是很艱難的.
在船上這段日子.

$^{921}$天氣的轉變.
有很多的問題會出現.
怎麼辦呢?.
他媽媽真的忍不住流眼淚.
看著船航行.
跟著船去追.
追著船.
追到不行為止.
他很傷心.
他一邊追一邊哭.
他兒子又將這件事情.
寫在日記上.
他說.
我現在才真正體會.
《約翰福音》三章十六字.
神愛世人.
神將他去獨身寺.
其一切信他的.
不自滅亡,反得永生.
他說他看到媽媽.
他要去中國.
媽媽那麼難過.
為我流淚.
我知道媽媽很痛苦.
雖然我到中國.
不一定會死.
有很多人照顧我.
但媽媽的心是那麼難過.
上天府.
她猜過她獨身寺.
目的就是為我們死.
所以.
在裡面.
耶穌降世.
是為死而來.
神的愛是多麼大.
祂多麼愛世人.
所以今天他開始明白經文.
這麼奇妙的愛.
去感動他.

$^{961}$他說今天.
我們踏上中國.
將主的愛傳出去.
將十教奇妙的救恩.
傳出去給中國.
去了中國之後.
他寫信給媽媽.
他問媽媽那天的感受.
他說我的感受.
是《約翰福音》三章十六節.
他媽媽的回信裡.
她說:兒啊.
我那時的心情.
都是想到這一節.
弟兄姊妹.
神真的賜他獨身兒子.
耶穌給我們.
神有賜福.
神的僕人.
所以他很愛中國.
他說:我若有一千英鎊.
都給中國.
若有一千條性命.
不留一條不給中國.
所以神恩待他.
那幾代.
都在中國.
這個就是曾孫.
戴孝曾.
曾經在中國.
但也被逮捕.
1929年.
他出生在開封.
珍珠港事變後.
被關在山東維縣.
集中營裡.
他曾經為主受苦.
他下一代.
就是戴繼宗.
戴繼宗是1959年.

$^{1001}$出生在台灣.
他現在是台灣人的女婿.
他的正名.
就是繼宗.
繼承.
祀奉祖宗的神.
繼續.
很開心地去事奉主.
他很關心.
中國.
所以他的子女的名字.
很特別.
兒子叫戴承若.
女兒叫成書.
另一個女兒.
叫成雅.
很容易記.
很特別.
我們來看看他們的生活.
雖然我們只說了一會.
其實是有血有淚的.
他們真的願意為主擺上.
弟兄姊妹我們都.
領受神的救恩這麼久.
盼望我們在生活裡.
我們都不忘記神這麼愛我們.
我們都願意盡力.
雖然我們未知是很苦.
但我們預備一個受苦的心智.
去跟奉主.
一起禱告.
你的兒女.
是何等的福分.
我們原是不配的罪人.
但你憐憫我們.
在我們是罪人的時候.
你已經救我們.
今天在座的每一個.
我們已經成為你的兒女.
求你在聖靈時提醒我們.

$^{1041}$身邊有很多人擦過.
他們都未認識主.
甚至我們的親人.
我們的家人.
求你幫助我們有勇氣.
離開忠於主的托付.
求主在當中帶領我們.
教我們信心.
奉主耶穌基督的名求.
\newpage



\section{馬太福音 6:25-34}
\label{sec:moCCv_ijPQc}
\textbf{2025.12.14《今日的恩典》:馬太福音6\_25-34(和修版) 曾裔貴牧師}
\newline
\newline
連結: \href{https://youtube.com/watch?v=moCCv_ijPQc}{\texttt{https://youtube.com/watch?v=moCCv\_ijPQc}} ~~~~ 語音日期: 2025-12-16
\newline
\newline
\hyperref[sec:WQHTSm8b4dk]{\small{< < < PREV SERMON < < <}}
~
\hyperref[sec:index_chronic]{\small{[返順時目]}}
~
\hyperref[sec:index_scriptual]{\small{[返順卷目]}}
~
\hyperref[sec:fk3jYby_MrQ]{\small{> > > NEXT SERMON > > >}}
\newline
\newline
馬太福音 6:25-34
\newline
\begin{longtable}{cl}
\hline
\hline
章節 & 經文 (和合本修訂版)\\
\hline
6:25 & \begin{tabularx}{0.7\textwidth}{X} 「所以，我告訴你們，不要為你們的生命憂慮吃甚麼喝甚麼，或為你們的身體憂慮穿甚麼。生命不勝於飲食嗎？身體不勝於衣裳嗎？ \end{tabularx} \\ \\ \relax
6:26 & \begin{tabularx}{0.7\textwidth}{X} 你們看一看那天上的飛鳥，也不種也不收，也不在倉裡存糧，你們的天父尚且養活牠們。你們不比飛鳥貴重得多嗎？ \end{tabularx} \\ \\ \relax
6:27 & \begin{tabularx}{0.7\textwidth}{X} 你們哪一個能藉著憂慮使壽數多加一刻呢？ \end{tabularx} \\ \\ \relax
6:28 & \begin{tabularx}{0.7\textwidth}{X} 何必為衣裳憂慮呢？你們想一想野地裡的百合花是怎麼長起來的：它也不勞動也不紡線。 \end{tabularx} \\ \\ \relax
6:29 & \begin{tabularx}{0.7\textwidth}{X} 然而我告訴你們，就是所羅門極榮華的時候，他所穿戴的還不如這些花的一朵呢！ \end{tabularx} \\ \\ \relax
6:30 & \begin{tabularx}{0.7\textwidth}{X} 你們這小信的人哪！野地裡的草今天還在，明天就丟在爐裡，神還給它這樣的妝飾，何況你們呢？ \end{tabularx} \\ \\ \relax
6:31 & \begin{tabularx}{0.7\textwidth}{X} 所以，不要憂慮，說：『我們吃甚麼？喝甚麼？穿甚麼？』 \end{tabularx} \\ \\ \relax
6:32 & \begin{tabularx}{0.7\textwidth}{X} 這都是外邦人所求的。你們需要這一切東西，你們的天父都知道。 \end{tabularx} \\ \\ \relax
6:33 & \begin{tabularx}{0.7\textwidth}{X} 你們要先求神的國和他的義，這些東西都要加給你們了。 \end{tabularx} \\ \\ \relax
6:34 & \begin{tabularx}{0.7\textwidth}{X} 所以，不要為明天憂慮，因為明天自有明天的憂慮；一天的難處一天當就夠了。」 \end{tabularx} \\ \\
[1ex]
\hline
\hline
\end{longtable}
$^{1}$我們一起同心祈禱.
天父啊 我們等候在你的面前.
我們知道.
這一段經文.
其實是相當不容易學習.
如果不是.
根本就不會有人.
需要 可能要見心理醫生.
因為憂慮會帶給我們.
種種的新心靈的疾病.
我們知道主耶穌勸勉當時的門徒.
在登山寶分上.
來教導.
證明這是困擾我們.
無論是信還是不信的人.
都是相當大的挑戰.
就是我們人類.
有一個好像與生俱來的習慣.
個性.
就是會憂慮.
有時計劃和憂慮.
好像有時我們都很難去界分.
天父啊盼望這段經文.
給我們再一次彼此有提醒.
特別我們在香港.
經歷了這麼特大的火災.
很多人的心都是受影響.
雖然已經過了差不多一個月了.
但是後續的問題.
可能接種的需要去關心.
需要去支援.
但是這個正正告訴我們.
這段經文叫我們真的不可以.
再沉溺在憂慮裡面.
反而我們要靠著主.
給我們實用的.
現在的.
每天的恩典.
這樣來生活.
尤其是我們是屬主的兒女.

$^{41}$更加盼望透過主.
今天給我們的教導.
我們去學習.
我們一起等候在你的面前.
我們這樣禱告.
奉耶穌基督的名.
阿們.
我就沒有一些具體的數據.
香港也有不同的.
類似是心理衛生會.
或者一些精神健康的組織.
每年都有一些的調查.
男性.
女性.
有抑鬱.
有嚴重的困擾.
譬如睡得不好.
或者有心靈的問題.
由這樣困擾而帶來的身體的毛病.
每年都有這些數據.
也不少的.
最近知道就是原來美國.
多人有困擾是女性.
香港呢.
好像剛剛反過來.
男性會比較多.
我不知道在座.
憂慮.
是不是我們每天生活.
起床睡醒之後的常態呢?.
我們在座也有一些精神科的醫生,護士.
都應該有一定的數據.
知道這些問題的嚴重性.
所以今天很想在12月最後一個月.
我們大家一起去想一想.
我們神的兒女.
我們如何學習.
在主裡面去知取這份恩典呢?.
剛才讀了經文.
我是刻意沒有將二十四節擺出來.

$^{81}$因為二十四節原來是相當重要的.
因為由這個段落.
稱為登山寶訓是由第五章到第七章.
在中間這一段六章.
主耶穌突然間叫門徒不要憂慮.
但是原來上文.
這個憂慮.
跟下面二十四節.
接著二十五到三十一節.
又有上下文的關聯.
這裡很簡單的分段.
如果我們首先不看下面.
十九到二十三節.
就是錢財的討論.
就是不要積贊財寶在地上.
因為地上有蟲子咬.
能夠受壞.
也有賊挖的洞來偷.
要在天上積蓄財寶.
當然用的描述.
是二千年前耶穌的時代.
不是今天我們擺在定期.
擺在銀行.
擺在保險箱.
或者其他各類的投資.
這裡告訴我們.
很清楚告訴我們.
要注目永恆.
不只是物質的世界.
因為第二.
二十一節之後.
二十二二十三節.
很難去解釋的.
就是說.
你心裡面的光如果黑暗了.
那個黑暗就何等大.
其實很簡單的.
主耶穌很想門徒.
留意不只是物質.
目前的情況.

$^{121}$當時的門徒剛剛跟隨耶穌.
可能也有很多不同的思想.
就是他們要放下.
不再打魚.
甚至約翰要放下他的漁船.
這樣來離別自己的父親,家人.
這樣來跟隨耶穌.
或者還有其他的門徒.
馬太要放下他不再做稅吏了.
各樣的情況.
可能在門徒當中.
開始會想.
如果我這樣跟隨耶穌.
會怎麼辦呢?.
可能就在這樣的場景.
這樣的背景下.
耶穌意識到門徒.
有這種小信.
所以就來跟他們說.
錢財的問題.
很重要的就是.
將眼目.
心靈所關注的.
放在永恆裡面.
所以那個比喻就是說.
要將錢財積蓄在天上.
就是說為將來.
為永恆.
為神的國去效力.
這正正就是三十四節.
三十三,三十四節.
教我們為什麼.
怎麼可以不憂慮的方法.
接著.
廈文的二十五至三十一節.
就開始說了兩個比喻.
天空的飛鳥和野地的百合花的比喻.
來講到我們不要憂慮.
憂慮也沒用.
因為不會幫助我們.

$^{161}$使得我們的生命.
我們的壽數.
可以多加一刻.
或者一爪.
所以你會看到.
上下文.
原來中間有一個很重要的重點.
中心.
就是二十四節.
我們先看這裡.
所以十九到二十三節.
大概可以歸納為.
門徒要注意的就是自己的心.
有沒有體重.
永恆.
將我們的心思意念.
心所是的最重要的.
擺在這個永恆的天上.
天上是不是我們會關注的重點呢?.
所以.
下面就是這麼簡單的說.
二十五,二十六節.
就告訴我們.
所以我告訴你們.
不要為你們的生命憂慮吃什麼喝什麼.
或為你們的身體憂慮穿什麼.
強調生命和生活.
不是說我們不注意飲食.
生命的健康.
透過飲食,生活來維持的.
但是如果那個焦點.
只擺在地上的生活.
整個生命的聚焦.
就是我要吃什麼穿什麼.
沒有想到永恆.
沒有想到天上的天家.
這個問題.
正正就是.
跟剛才上文所說的.
有一個很大的連貫.

$^{201}$積蓄永恆.
以及不要憂慮.
接著那個中心思想.
就停留在這一段很重要的.
就是二十四節.
一個人不能夠服侍兩個主.
他不能夠恨這個愛那個.
如果當時的服侍.
就應該是奴隸制度的比喻.
當然門徒就不是來做奴隸.
不過跟隨耶穌.
就有這樣的含義.
有這樣的原則.
這裡不只是對著這十二個門徒.
當時已經呼召他們是使徒了.
這樣來說.
不過這裡強調一個很重要的原則.
耶穌來教導他的門徒.
就是說我們不可以縱輕另外一個.
你們不能夠又服侍神.
又服侍馬門.
馬門中文的翻譯就是財利.
所以不僅是財利.
應該說.
整個物質世界.
維繫我們的生活的一切.
物質主義的大層面.
如果我們在這裡人生.
根本就沒有永恆.
沒有神.
沒有這位獨一的主.
是我們的生命的焦點.
這就是馬門主義.
敬拜馬門.
所以這裡耶穌也告訴了門徒.
只能夠二選其一.
只能夠敬拜這位獨一的主.
只能夠拜這位神.
所以耶穌對他們的教導和吩咐.
相當要緊的就是告訴我們.

$^{241}$在兩個段落的中間.
這個中心思想.
也是針對憂慮.
就是說有這樣的含義.
如果拜馬門.
一心只為地上的事物去掛慮.
就會出現種種的憂慮.
就會在神以外.
不依靠神的恩典,力量,教導.
和看自己在地上的人生.
就是一切這樣的心思.
正正就是這大段聖經裡面的教導和吩咐.
因為在二十四節.
本來跟上下文都沒有關聯的.
明明說要積財寶在天上.
二十五節立即說.
所以你們不要憂慮.
但是二十四節突然說.
只能夠敬拜一位神.
這樣的拆解.
就真的給我們一個很大的思想.
原來我們只是為自己的憂慮去找辦法.
那不是最終極解決的辦法.
我們一定會有不同的擔心.
昨天我們都很感恩.
透過駱喬留的社工.
他們呼籲有什麼物資可以給寶寶.
我們立即鼓勵弟兄姊妹.
有的就帶回教會.
有些媽媽就拿了很多衣服.
拿了一些寶寶的床.
我們立即送過去.
因為有二百五十多戶.
在火災.
政府編配了駱喬留.
就是說是我們的鄰居.
我們可能有一段比較長的時間.
我們都可以來關心,支援.
總有各樣的憂慮.
開始考試了.

$^{281}$媽媽又有憂慮.
小朋友又有憂慮.
聽一些在這方面的弟兄姊妹說.
考DSE的學生.
去到精神崩潰.
因為父母逼他讀書.
逼到一個地步.
他完全崩潰.
是真實的個案.
很多這些事情.
我想我們每個人.
都有我們自己需要處理的事情.
有我們要憂慮的地方.
很多事情每個人都不同.
每個家庭都不同.
我們感受有一位姐妹.
媽媽又有精神困擾.
那種痛苦是難以想像.
突然間會失控.
要來到吵吵鬧鬧.
令到姐妹自己.
日間的工作已經很繁重了.
都要這樣來照顧的時候.
那個照顧者本身自己都有困擾.
要我們為他祈禱.
有時工作太忙碌.
都很多容易病痛.
所以我想每個弟兄姊妹.
每個家庭都有自己要困擾.
要面對的困難.
耶穌在這裡不是禁止我們不要憂慮.
耶穌告訴我們.
為什麼我們可以不憂慮.
還有到底憂慮的定義.
我們怎樣去處理.
如果我們用錯了方法.
可能我們持續都有困難.
在這樣的困境.
我希望弟兄姊妹.
你的人生不是沒有神.

$^{321}$我們屬主的兒女.
我們一樣會有不同的憂慮.
如果不用憂慮.
即是說我們根本不需要計劃.
明天就有明天的憂慮了.
這是34節.
是不是這樣呢?.
肯定不是.
肯定我們每件事要有很周詳的計劃.
我們要認真面對不同的事情.
我們要有準備.
包括一個powerpoint(力量點).
都用了起碼三個小時去準備.
是需要每件事都有規劃的.
所以不是逃避那個憂慮.
什麼都不要想.
叫做不憂慮了.
不是這樣的.
但是你看看這個結構.
在24節.
會不會真是這個最重要的重點.
就是我們到底我們的生命.
歸於哪一位呢?.
我們是不是認定這一位愛我們的真神.
是我們生命的主.
我們就將我們整個的生命投向祂.
為祂而活呢?.
以致我們要面對的事情.
我們是讓祂來跟我們一起去參與.
來一起的計劃.
在祂的裡面.
我們來做這些種種的規劃.
規劃以外我們承擔不了的.
就是稱為災難.
都有祂的管理.
保守.
和祂背後的意思.
在這樣的時候.
會不會就是我們基督徒.
需要去想.

$^{361}$不只是面向憂慮本身.
例如剛剛生了孩子.
很忙.
睡眠不足.
很疲倦.
會很燥.
我們就想想用什麼辦法.
或者我的性格本身是怕事的.
我們要讀一些課程.
去處理我的性格.
積極一點.
開心一點.
這些本身是正面的.
不是錯事.
了解了很多事情.
但是.
耶穌既然叫我們不要憂慮.
憂慮也沒用.
因為明天自然有明天的憂慮.
那應該怎麼樣呢?.
會不會二十四節.
正正就是真真正正的解約呢?.
就是我們認定.
我們有一位這麼愛我們的主.
我們就交託我們的生命給主.
以致既然交託的時候.
我的家庭.
我的事業.
我的家人.
我都是要交託給主.
在這樣的情況下.
我們來到去學習的.
不憂慮.
以致我們經驗到.
經歷到的.
就是神給我們今天的恩典.
繼續下去.
我們看看下面的經文.
在這個結構裡面.
剛才提到.

$^{401}$生命要注目在永恆.
就是說我們人生地上的生活.
每天是應該有永恆的.
每天我們有沒有為永恆而努力.
來準備.
來將我們的生命.
放在關注.
在神的國度裡面呢?.
以致我們能夠在這個事奉獨一的主裡面.
我們就是這樣來到每天.
不斷地進心和學習.
接著.
主耶穌的教導二十五節.
到三十一節.
就是告訴我們.
不要憂慮.
憂慮也是沒有用的.
在這樣的時候.
就很不一樣了.
因為我們在做的事情.
每天都要面對不同的決定.
不同的煩惱.
但是.
我們整個的焦點.
就不是我們自己去煩惱.
我們自己去單打獨鬥.
而是有神跟我們一起.
因為我們是在事奉獨一的主.
剛才讀過了.
第四個.
似乎就是跟上文耶穌說的.
是完全一致的.
他說.
這都是外邦人所求的.
你們需要的這一切的東西.
跟這些門徒說.
你們的天父都是知道.
你們要先求神的國和他的義.
這些東西都要加給你們了.
這些東西.

$^{441}$如果上文耶穌用兩個的比喻.
食物.
衣著.
就是維繫著我們生活最基本的東西.
神都已經為我們有準備了.
我們不需要為那些來去憂慮.
不過呢.
耶穌在這裡給我們一個很清晰.
跟上文似乎相當一致的.
就是先求神的國和他的義.
以致我們不需要為明天來去憂慮.
只需要在今天我們活著.
就是為神的國來去效力.
整個的心思意念.
是為著神的國度去考慮.
這裡同時間也說.
不代表我們的生活.
我們不需要去計劃.
兩者是沒有衝突.
但是.
如果看整段聖經.
沒有了單單倚靠和敬拜獨一的主.
而將焦點不自覺放在.
有另外一位神.
就是馬門.
在我們的生命.
在我們的心靈裡面.
這就是一個.
令我們很多憂慮的原因所在.
就是經文給我們的解釋.
以及透視.
這裡很重要.
告訴我們.
主耶穌教我們不要憂慮.
不是面對憂慮本身.
是我們有沒有注目永恆.
為神的國.
我們去擺上.
為神的義.
在我們的生活裡面來去落實.

$^{481}$我曾經在.
不知是在《團契》說過.
我曾經在屯門醫院工作了兩年.
差不多到我要清除自己的蒙召.
去讀神學的時間.
那時候差不多接近一個終身制.
政府公務員的終身制.
兩年的 probation(懲罰)就過了.
有一次.
本來我有一個上司.
我對上的上司.
比我高級的.
他剛好放大假.
我的大老闆.
就是有房間在裡面的大老闆.
那天是下午三點多.
他走出來拿著他的手提袋.
就跟我說.
我叫 Calvin(凱爾文).
他就說 Calvin(凱爾文).
我出去做治療.
如果總部的大老闆打電話來.
你說我不在就行了.
那慘了.
我怎麼辦呢?.
因為那時候剛剛信耶穌.
信了一年多了.
當然是祈禱希望不要有電話響.
過了不久之後.
他的房間的電話響了.
原來是一個回電.
因為我的大老闆.
打完電話去總部的大老闆之後.
就是 headquarter(總部)的總部.
就交代了那些工作.
交代完以為沒事了.
就走了.
早點去做治療.
後來呢.
這個大老闆在 headquarter(總部).

$^{521}$在 AntiVaxx(反毒藥)的.
葵興那裡.
那個 headquarter(總部).
就打來找我的大老闆.
問一些事情.
一些細節.
那麼.
沒有遇到.
當時沒有手提電話那些的.
那麼.
裡面的電話響了很久.
那麼我的下屬就說.
你還不去聽電話?.
我說我聽電話.
我不知道是不是大老闆打來的.
他打來那我怎麼回答呢?.
那個下屬又不肯進去聽.
他說應該是你聽的.
那麼死死的氣進去聽.
果然是這個大老闆打來.
我擔心的事情真的發生了.
就問.
叫誰聽電話?.
阿黃先生.
就是我們的大老闆.
那一秒鐘.
我不知道怎麼回答.
因為老闆說.
你說我不在就行了.
對了.
那麼我直接說.
阿黃先生不在.
他剛剛才談完電話.
他去哪了?.
那麼我又不能夠說謊.
他說他去做物理治療.
當然好像很好笑.
弟兄姊妹.
這些是其中很少的一些事情.
我們有沒有去行神的義呢?.

$^{561}$我們有沒有避免了衝突呢?.
我們就說一個善意的謊言.
其實完全不關我的事.
那麼我們怎麼去處理呢?.
當然還有另外一些再大一點的.
是有一些人來看醫生.
因為延誤了.
他的小朋友死了.
去不了急症.
因為他沒有帶身份證.
我的下屬的同事.
我那時候已經出去吃飯了.
我的下屬就說.
沒有身份證.
不能賣.
那麼他很急.
他又不懂得去急症室.
只是一直在搖擺.
衝回家拿身份證.
接著他的小朋友就救不了.
那麼就投訴我們的部門.
那麼我的上司.
就很不喜歡信耶穌的人.
就當我是有份的.
其實我根本就不在現場.
我根本就去了吃飯.
我是剛剛午餐時間.
那個時間.
有一個HR的同事.
來跟我說.
喂.
你被上司投訴得很厲害.
說你跟那個小朋友死了有關.
那麼我說.
真的不關我的事.
後來輾轉的調查.
後來都平息了.
所以兄弟姐妹.
耶穌教我們.
不只是針對憂慮本身.

$^{601}$不是我們要做好一點自己.
就解決了.
而是真的告訴我們.
我們是屬主的兒女.
我們要追求的是.
在神的國度裡面.
在神的義裡面.
我們是有份的.
還有我們要看重的.
我們只能敬拜獨一的主.
不能夠有兩個主.
在我們自己的生命裡面.
更加不能夠.
我們整個生命的心思意念.
只是放在地上的短暫.
我能夠拿捏的.
我能夠抓住的.
這樣來令自己有保障.
永遠不能夠滿足的.
永遠不能夠改變到.
我們憂慮的本身.
所以耶穌基督似乎告訴我們.
向永恆而活的生命.
就是專注在神的國度裡面.
我們就能夠有滿足.
我們能夠面對不同的事情的時候.
絕對不是保證沒有困難.
我們要決策.
我們要煩惱的事情.
可能每天都有很多的.
我們作為牧者都要面對不同的人.
剛剛有一位社工轉介.
有一個癌症很嚴重的一個家庭.
他的兒子失業了三年了.
很想我們去關心.
你可以想像得到.
我們一定要面對不同的困難的事情.
我們要幫助有這樣的困難的人.
所以一定會有的.
但是這段經文告訴我們.

$^{641}$我們有的恩典就是.
基督徒我們只需要專一在神的裡面.
以及我們積蓄財寶.
那就是為神的國.
為神的義去效力.
我們就能夠有真正的平安.
真正的恩典.
去逃避我們每天.
心思所掛念的.
懼怕的.
憂慮的事情.
用一些例子跟大家去剖析一下.
這些的總結.
想跟大家看看這個人.
貴格物片可能大家都吃過了.
這位的赫魯威爾先生.
也是相當敬虔.
他第一幅貴格物片的圖像.
就是貴格會的清教徒的裝扮.
他17歲的時候.
其實他爸爸也是死於費潔克.
當時一百多年的時間.
本來醫生也說.
你時日無多了.
當時他只不過17歲而已.
但是他很虛心地向神禱告.
然後因為他爸爸死的時候.
也遺下了一大筆遺產.
給他很多的時間.
他有七年的時間去休養去旅行.
他整個人生閱歷更加豐富.
在這樣的情況下.
居然奇蹟地康復.
他就有一個很大的立志.
他說我由今天開始.
徹底地將自己降服.
敬拜事奉神.
因為當時他聽到.
慕迪這位報道家的訊息.
他的生命完全交託給主.

$^{681}$然後他一生就是做生意的.
他成立的貴格的物品.
希望美國人能夠吃得健康.
然後他就有一個立志.
一生為神的國度.
為傳揚福音來效力.
所以他在四十年裡面.
一直都是慕迪神學院的主席.
他奉獻很多的金錢.
在神學院裡面.
每年他的總收入的六成到七成.
是奉獻給神的福音工作.
在1927年的時候.
他就成立了一個基金.
到今天都有的.
這個基金的宗旨.
唯一的使用.
只能夠是為福音來使用.
1927年已經成立了的基金.
這一位一生為神來效力的生意人.
他是二十世紀最令人敬佩的.
一個做生意的基督徒.
這個赫魯威爾.
當然不是賣告白.
叫你回去吃貴格物品.
不過真的有人這樣以聖經的教導.
來作為他做生意的原則.
賺的錢為神而活.
我們都有退休的弟兄姊妹.
在神學院裡面有不同的參與.
有不同的服事.
有不同的機構.
或者在教會裡面有不同的服事.
我們有沒有將我們的人生.
我們平日的生活完了.
我們的工作完了.
我將我的有限的人生.
盡力為神的國度去效力呢?.
另外有一位姊妹.
這些字很小.

$^{721}$我特意不給大家看.
這位黃勳傳道.
其實是一位很出色的無線電視的編劇.
他的工作做得很好.
不過他很無奈地.
他在做工作的時候已經患了癌症.
是一個很罕有的癌症.
輾轉他的同學.
大學的同學.
帶他去聽報道會.
他真的信主信得很好.
後來經過了不同的事件.
他自己就切心地去奉獻自己做傳道.
我跟他同過兩屆的.
他是2008年來畢業的.
我是2009年畢業.
他是2008年畢業的.
但是在讀神學的期間.
癌症再度的復發.
令到他真的很痛苦.
延遲了一年.
輾轉到了2011年.
他就要安息住了.
但是在他這個抗癌的路上.
他得到的就是.
每天很多的痛苦.
無盡的治療.
頭髮差不多都掉了.
很多的作嘔.
很多的事情令到他痛苦.
但是他展現的一種生命力.
居然畢業之後.
有三年時間做半職的傳道.
這樣的生命反而激勵很多弟兄姊妹.
看到神的作為.
是多麼的奇妙.
在這位姊妹的身上.
稱為黃昏傳道.
這位姊妹.
這位牧者.

$^{761}$就是這樣走完她的人生.
最後一句話.
在她的Facebook的留言.
她說天父啊.
我已經盡力了.
就是完結了她的人生.
一位這麼出色的一個劇作家.
無線電視的編劇.
在傳道的生涯.
做不了多少事情.
進神學院已經再次.
他的病再復發.
成為學院裡面.
一個需要照顧的同學.
很多他的同班同學.
老師來照顧他.
幫他.
但是他展現到的生命力.
他展現到的就是無論多困難.
他的心.
因耶和華.
歡欣.
我的神.
喜樂.
詩篇的二十七篇.
就是他自己生命的一種見證.
用他自己的生命.
來反照出.
原來我們能夠事奉一位主.
是最寶貴的一件事.
弟兄姊妹.
每個人都有你的困難.
每個人都有你自己的憂慮.
或大或小.
但是.
有沒有將這一切.
交託在神的手中.
有沒有將自己的生命.
交託給神.
為主.

$^{801}$為永恆.
來效力呢?.
我們的兒女.
我們的職業.
等等.
我們有沒有將它交託給主.
來這樣一起服侍呢?.
盼望這一段經文.
成為我們在整個今年裡面.
我們大家都來到不斷地思想.
我是不是已經注目了.
在那些難處.
而忘記了神呢?.
我有沒有真的在神的裡面.
去看我在永恆裡面的投資呢?.
願主幫助和激勵我們每一位.
因為其實.
我們有很多事情.
是很需要大家一起去做.
周邊有很多人.
他們落在那個困難.
但是問題就是.
他們是沒有神.
他們還沒有認識神.
所以是很需要我們大家一起.
將這個寶貴的福音.
跟他們去分享.
當他們有了神之後.
他知道.
主是他的獨一的主的時候.
他們就有力量面對他們的困難.
否則的話.
當他們自己這樣投入那個困難裡面.
其實是沒有出路的.
願主幫助.
我們一起禱告.
我們感恩天父給我們能夠.
有主耶穌剛才的教導.
我們可以來到去單一的.
來到去事奉我們的主.

$^{841}$因為你是那一位永遠不改變.
而且是愛我們.
野地的花草.
野地的飛鳥.
都不夠我們這麼貴重.
這個是主耶穌說的.
多謝主讓我們看到.
原來我們在你的眼中.
是這麼貴重.
幫助我們能夠認識.
我們是這麼寶貴.
願神施恩給我們.
以致我們屬主的兒女.
我們不要落在一種小信.
懼怕裡面.
反而我們能夠不斷地去操練.
學習.
好像剛才我的學子.
黃昏傳道那樣.
無論多麼困難.
我們能夠因耶和華歡恩.
因我的神而快樂.
求主幫助.
也都讓我們定睛.
在永恆不改變的真實.
以致我們所擺上的.
將來一定會蒙主的恩典賞賜.
我們這樣感恩禱告.
奉耶穌基督的名.
阿們.
\newpage



\section{約翰福音 10:7-18}
\label{sec:fk3jYby_MrQ}
\textbf{2025.12.21《祂來了,是叫人得生命》約翰福音10\_7-18(和修版) 李炳光牧師}
\newline
\newline
連結: \href{https://youtube.com/watch?v=fk3jYby_MrQ}{\texttt{https://youtube.com/watch?v=fk3jYby\_MrQ}} ~~~~ 語音日期: 2025-12-23
\newline
\newline
\hyperref[sec:moCCv_ijPQc]{\small{< < < PREV SERMON < < <}}
~
\hyperref[sec:index_chronic]{\small{[返順時目]}}
~
\hyperref[sec:index_scriptual]{\small{[返順卷目]}}
~
\hyperref[sec:0LnkonZdJmk]{\small{> > > NEXT SERMON > > >}}
\newline
\newline
約翰福音 10:7-18
\newline
\begin{longtable}{cl}
\hline
\hline
章節 & 經文 (和合本修訂版)\\
\hline
10:7 & \begin{tabularx}{0.7\textwidth}{X} 所以，耶穌又對他們說：「我實實在在地告訴你們，我就是羊的門。 \end{tabularx} \\ \\ \relax
10:8 & \begin{tabularx}{0.7\textwidth}{X} 凡在我以前來的都是賊，是強盜；羊沒有聽從他們。 \end{tabularx} \\ \\ \relax
10:9 & \begin{tabularx}{0.7\textwidth}{X} 我就是門，凡從我進來的，必得安全，並且可進出，找到草吃。 \end{tabularx} \\ \\ \relax
10:10 & \begin{tabularx}{0.7\textwidth}{X} 盜賊來，無非要偷竊、殺害、毀壞；我來了，是要羊得生命，並且得的更豐盛。 \end{tabularx} \\ \\ \relax
10:11 & \begin{tabularx}{0.7\textwidth}{X} 「我是好牧人，好牧人為羊捨命。 \end{tabularx} \\ \\ \relax
10:12 & \begin{tabularx}{0.7\textwidth}{X} 雇工不是牧人，羊不是他自己的，他一看見狼來，就撇下羊群逃跑；狼抓住羊，把牠們趕散。 \end{tabularx} \\ \\ \relax
10:13 & \begin{tabularx}{0.7\textwidth}{X} 雇工逃走，因為他是雇工，對羊毫不關心。 \end{tabularx} \\ \\ \relax
10:14 & \begin{tabularx}{0.7\textwidth}{X} 我是好牧人；我認識我的羊，我的羊也認識我， \end{tabularx} \\ \\ \relax
10:15 & \begin{tabularx}{0.7\textwidth}{X} 正如父認識我，我也認識父一樣；並且我為羊捨命。 \end{tabularx} \\ \\ \relax
10:16 & \begin{tabularx}{0.7\textwidth}{X} 我另外有羊，不屬這圈裡的，我必須領牠們來，牠們也要聽我的聲音，並且要合成一群，歸一個牧人。 \end{tabularx} \\ \\ \relax
10:17 & \begin{tabularx}{0.7\textwidth}{X} 為此，我父愛我，因為我把命捨去，好再取回來。 \end{tabularx} \\ \\ \relax
10:18 & \begin{tabularx}{0.7\textwidth}{X} 沒有人奪去我的命，是我自己捨的；我有權捨棄，也有權再取回。這是我從我父所受的命令。」 \end{tabularx} \\ \\
[1ex]
\hline
\hline
\end{longtable}
$^{1}$各位大家剛才聽到葉華先生和他太太的分享.
大家有什麼感受?.
很坦白說.
我自己坐在這裡.
是非常的感動.
我不止聽葉華先生和他太太分享一次.
很多次.
其實過去有一段很長的時間.
我常常和葉華先生和他太太一起舉行報道會.
這次我沒有約他們很久了.
這次特別約他們來.
特別在貴堂舉行一個很好的分享會.
老實說.
我在下面聽.
我都流起眼淚.
因為這是真實的見證.
不是開玩笑的.
我今天是非常的開心.
來到這個地方.
這次離差很遠.
為什麼?.
今天坐滿了人.
這是主的吸引力.
還是葉華先生夫人的吸引力呢?.
或者說.
是真正上帝的呼召.
各位朋友.
各位弟兄姊妹.
這不是一個簡單的聚會.
因為.
這是真實的見證.
不是說故事.
我告訴你們.
多年前.
我坐的士.
坐的士的時候.
看到前面的的士車.
車上掛了很多.
教會的飾物.
有十字架.

$^{41}$很多東西放在那裡.
我.
你知道我做牧師.
任何場合.
我都希望可以做見證.
傳福音.
咦?這個司機.
可能是基督徒.
我說:朋友,你是不是基督徒?.
他說:是的.
原來這個就是葉華.
前面開車.
我和他聊天.
聊了一會.
突然他說:牧師.
我請你喝一杯,好不好?.
我問請我喝什麼?.
第一次和你坐車.
第二次請我喝.
我再去結婚.
我說:再結婚.
我就不來了.
為什麼?.
原來他說.
我和我太太再結婚.
各位.
你覺得這句話有沒有意思呢?.
從太太再結婚.
為什麼再結婚呢?.
原來是一個這麼感人的故事.
各位.
這是很真實的.
我告訴大家.
我做牧師這麼久.
我做了很多年牧師.
今年我做了六十多年牧師.
咦?.
在當中還沒有出生.
我多少歲呢?.
保守秘密.

$^{81}$大家想想.
各位.
這是一個這樣的情況.
我告訴大家.
聖誕節快到了.
幾十年前.
聖誕節.
很高興.
很多人洗禮.
那一年.
有八十多人洗禮.
我做禮拜的時候.
不能夠八十多人洗禮.
花了很多時間.
我特別舉行了一個.
洗禮的聚會.
在聖誕節的下午.
那天.
我的禮拜堂坐滿了人.
我禮拜堂從哪裡來呢?.
是灣仔一間紅磚屋.
三角禮拜堂.
坐滿了人.
那.
我說為什麼這麼重要呢?.
原來.
那一年.
我便要退休了.
二十多年前.
我退休了.
那些人知道.
李牧師退休了.
我便趕快給他洗禮.
二十多人.
其中有一個.
像葉華生這樣的病態賭徒.
故事.
真的差不多像葉華.
他.
憑著.

$^{121}$在我門口的電子板.
寫著.
一段聖經.
這段聖經說.
凡路苦擔重擔的人.
可以到我這裡來.
我就使得安息.
這個人進入了我的禮拜堂.
我建了一間禮拜堂.
重建的.
我的規矩就是.
打開禮拜堂的門.
禮拜堂的門常要打開.
而且禮拜堂是祈禱的電.
要隨時被人祈禱.
現在你去到灣仔的禮拜堂.
樓下有個房間.
是祈禱室.
直到現在每天都有幾十人.
不是我們教會的.
有些人在裡面祈禱.
我打開門.
被人進去.
這位賭徒.
因為看到電子板.
又被老婆趕出去.
又收到球場.
這是我的禮拜堂.
他那天比以為更慘.
想自殺.
來到門口.
看到這個字.
為什麼這麼奇怪.
進去坐在那裡.
為什麼祈禱室這麼奇怪.
這麼多聖經在那裡.
怎知道他一打開聖經.
你猜他站在哪裡.
約翰福音.
十五章.

$^{161}$講浪子回家的故事.
他第一次看.
他哭了起來.
哭了起來.
我有個童工走過.
看到他.
就跟他說福音.
他很感動.
跟他祈禱.
結果.
那一次.
我和我的童工.
跟他說福音的時候.
一定要放下.
這個壞習慣.
過一次洗禮.
老婆又離婚了.
那一次洗禮.
我今天沒有準備講章.
因為我今天不用講章.
特別聽了見證之後.
我就隨心而發.
希望神的靈動功.
感動在在各位.
為什麼呢.
那一次.
洗禮的時候.
因為很多人.
很多人.
所以我就要下午.
我就有八十多人洗禮.
我就找了幾個.
做見證.
我當然是找這個.
這位姓謝的.
不是姓葉的.
不是.
其他的姓.
也不是.
華哥.

$^{201}$竟然.
他出來做見證.
在整個洗禮當中.
他就說.
他自己.
一樣像自己的見證一樣.
那一刻他停了.
原來他看著頭.
後面知道醒水了.
當然是他太太.
來了.
那我就走出去.
按了按鈕.
等一下.
我問一下.
是不是謝太太來了.
看那邊.
如果是.
你太太來了.
你叫出來好不好.
這個做見證的按了按鈕.
我說.
是謝太太.
姓謝的謝太太.
我叫了幾次.
都沒有人認.
我以為他不肯.
或者沒有來到.
誰知道突然間.
我看到後面的女人舉了手.
我說.
謝太太你來了.
真的好.
可不可以出來.
當她說到離婚的時候.
離婚的時候.
由於是聖誕節.
那張離婚書還沒有簽.
說到這時.
他自己流眼淚.

$^{241}$原來他太太坐在後面.
我就請她出來.
我告訴大家.
這是我做牧師一個最難忘的經驗.
我就把他們兩個拉起來.
各位.
這個就是謝太太.
這位是今天洗禮的謝先生.
我在各位的面前.
想問問她.
謝太太.
你是不是真的要離婚.
你肯不肯說不離婚.
你肯不肯和她在一起.
在這個時候.
那個謝太太.
點頭.
全場鼓掌.
哇.
真是很重要.
在離婚的時候.
我親自把他們兩個拉起來.
兩個抱著.
哭了起來.
我主持禮拜.
主持報道會.
未曾見過一個人哭.
一個聚會哭.
這麼多人哭.
很多人在聚會中流眼淚.
看見.
罪人悔改.
一個罪人悔改.
每次都為他們歡喜快樂.
這件事是真實的事.
拉起來.
各位.
今天我再聽.
耀華先生和他的太太.
這麼好的見證.

$^{281}$各位.
他請我喝酒.
我說.
當然我那晚沒有去.
但我怎麼好意思.
第一次見面.
當然要喝一杯.
但是我和他認識了.
多次的見證.
不是隨口而出.
現在已經十年了.
我遇到他的時候.
遇到他的時候.
已經二十多年了.
他們已經悔改了.
這麼久的時間.
仍然是神的恩典.
淋到他們的家庭當中.
當我開始看他們的家庭的時候.
有這麼多的孫子.
兒女.
老實說我真的忍不住眼淚.
因為看見.
神在他們家庭當中的祝福.
所以我看到.
一個信耶穌的人.
神一定特別祝福他們的家庭.
今天我們很感動.
不是那麼容易在的士上遇到我的朋友.
幾十年了.
看到他仍然是信耶穌.
仍然是很堅持的.
不斷的作見證.
不斷的為他們感恩.
各位.
我今天帶了兩位來.
在各位面前.
各位.
這不是講故事.
是真實的見證.

$^{321}$真實的見證就是.
在我們的生命當中.
有特別的事情發生了.
今年聖誕節快到了.
曾經有人問.
大埔發生了這麼可怕的事情.
今年你們基督徒還慶祝聖誕.
是不是行裝落伍.
我說不是.
因為我們基督徒慶祝聖誕.
不是喝喝吃吃.
不是歡歌作樂.
而是我們基督徒慶祝聖誕.
紀念耶穌基督在我們生命當中.
所做的偉大的事.
耶穌來到世界上.
是為了什麼呢?.
使人得生命.
並且得到更豐盛的生命.
你看看.
我們這兩位做見證的嘉賓.
是不是生命改變了?.
是不是?.
各位.
我們每個人都說.
我們不是病態賭徒.
那需要改變嗎?.
不是的.
各位朋友.
我發現有很多人.
各人有各人的問題.
我們的家庭.
我們的事業.
我們的兒女.
我們的生活.
還有你知不知道.
活在世界上.
每個人都在困難中.
痛苦中.
掙扎中.

$^{361}$你知不知道.
現在香港有很多人.
是抑鬱的.
什麼叫抑鬱病?.
抑鬱病.
不是說那麼容易.
抑鬱病.
就是你看見.
世界上很多東西.
都沒有興趣.
有些花不香.
有些鳥不唱歌.
有些人不笑.
為什麼?.
因為有很多困難.
有很多人睡不著覺.
有很多人靠安眠藥.
才可以睡得著.
我不叫大家舉手.
你有沒有.
發惡夢?.
葉太太.
剛才說.
我發了惡夢.
但是.
有很多人睡不著覺.
你知不知道.
有很多人自殺.
就是因為失眠.
你不要以為失眠是小事.
我研究過.
有幾個電影明星.
他早上自殺.
幾個都是早上.
服藥.
原來早上是最危險的.
整晚睡不著.
早上天空光.
死了又睡不著.
很空虛.

$^{401}$很難過.
很失望.
他想自己死.
很可憐.
雙人同情.
有一個電影明星.
自殺.
服藥後.
雙鳥自殺.
開了煤氣.
服了安眠藥.
打電話給老公.
來不及.
老公來了.
他死了.
原來.
今天有很多學生自殺.
原因是家庭壓力.
有很多老人家自殺.
各位不要輕看這件事.
很多老人家.
不知道將來怎麼辦.
特別年老的時候.
有病.
沒有人理.
所以今天.
各位朋友.
各位兄弟姐妹.
祂的恩典.
今天在紅恩堂.
是為兩位.
他們的見證.
他們家庭的改變.
他們生命的改變.
而得到神的恩典.
真實的見證.
希望大家.
都能夠感動.
不是說說.
聽聽.

$^{441}$而是真實的.
耶穌基督說.
我是好牧人.
好牧人是為羊而命.
祂來到世界上.
不是吃喝喝.
不是喝喝玩樂.
祂來到世界上.
是為了你和我的緣故.
所以聖誕節.
今天很多人過.
沒有耶穌的聖誕.
就是說你和人家做生日.
每個人.
是受聖公上聖婆不在.
自己飲飲食食.
自己交換禮物.
但祂沒有份.
現在我們的主耶穌.
都一樣來到世界上.
我們將祂放在旁邊.
沒有人理.
為什麼?我們只有慶祝.
今年.
我們全香港.
都遺有大埔的事件.
我們難過.
我們不想.
用飲飲食食.
來慶祝聖誕.
我希望大家再一次思想.
生命的價值.
生命的意義.
並且看到我們為何而生.
為何而活.
以及我們應該怎樣.
面對人生的.
不同的處境.
我認識一位傳教士.
這位傳教士.

$^{481}$英文叫Stanley Jones.
他在印度做傳教.
他做得非常非常好.
但有一次.
他坐飛機.
飛機的機師.
突然宣布.
他說.
他要去.
印度.
他就說.
我去印度.
他就突然宣布.
他說各位.
你快點幫我救生帶.
還有.
把氧氣罩拔下來.
因為我的飛機.
有問題.
你準備好了.
就在這個時候.
報告之後.
你知道坐飛機.
有這樣的報告.
真的會失魂.
很多人大叫大喊.
但這位牧師坐在那裡.
他說我是牧師.
我不是像他這樣.
這樣.
我看上帝怎樣跟我說話.
就在公事包裡.
拿一本聖經出來.
但是.
他看到一張白紙在那裡.
或者可能上帝.
寫給後人.
可能這次我不能錄了.
於是那張紙上寫了.
這次我們遭到.

$^{521}$飛機的作用問題.
我呢.
麥克風不夠好.
還不夠好.
那.
那白紙上寫了什麼呢.
他說這次我在飛機上遇到.
是有這樣的問題.
現在呢.
我們在飛機上繞了很久.
都不能下降.
不如這樣吧.
我寫的字留給後人.
誰能收拾到.
希望有幫助.
他說我呢.
傳福音已經有幾十年了.
我只是想留一句.
說話給世上的人.
這句什麼說話呢.
我自己呢.
記憶很深入.
為什麼呢.
我看了之後很感動.
這句說話寫著.
生活在耶穌裡面.
是生活唯一的方法.
這個老人家Stanley Jones.
就將他一生的傳教.
的經驗.
將這句說話綜合.
生活在耶穌裡面.
是生活唯一的方法.
他因為經驗到.
原來在基督裡面.
我們有平安.
在基督裡面.
我們有希望.
在基督裡面我們有新的生命.
各位朋友.

$^{561}$我今天將這句說話.
分享給各位.
不是在乎多長.
而是在乎我們有多好.
這句說話.
不是我說的.
又是一個著名的.
牧者.
這個牧者是誰呢.
就是Peter Marshall.
又是一個很有名的.
美國牧師.
他做到國會的牧師.
他怎樣做國會的牧師呢.
他是一個記者.
他未曾做牧師之前.
怎樣能夠吸引他做牧師呢.
他未曾做牧師之前.
他做牧師.
有一次.
這個記者接到一個消息.
說有一個樓快要塌了.
他要去採訪.
去到現場的時候.
看到消防員個個都很緊張.
叫人家全部離開.
全層樓.
都出來了.
以為口淡大氣.
突然間.
發現到.
在不知幾多層樓的騎樓.
有個小朋友在那裡.
嘩!不得了.
這個時候還有個小朋友在騎樓.
不得了.
這樣.
小朋友.
現在我拉開救生網.
你跳下來.

$^{601}$誰知.
越叫.
真的,你和我都會.
你看上去.
救生網好像大餅.
你敢不敢跳下去.
不敢跳.
而且又不認識.
最危急的時候.
突然有個聲音.
說什麼呢.
他說叫個小朋友的名字.
什麼名字我們不知道.
小強.
爸爸在下面.
你跳下來.
一說完.
爸爸在下面.
這個小朋友就突然間跳下去.
為什麼.
因為爸爸在下面.
他可以放心.
朋友.
你知不知道.
聖誕節真正的意義.
不是在乎吃喝玩樂.
而是在乎讓我們知道.
再次檢視我們的生命.
究竟我們為何而活.
為何而生.
這個.
It is how long but how well.
這個不是在乎多長是好.
就是這個Peter Marshall.
他後來是做了牧師.
做了國會牧師.
去到西點軍校.
大戰就快開始的時候.
他對軍官學生說話.
一進到西點軍校的禮堂.

$^{641}$個個坐滿了不停的軍服.
紅斗斗氣昂昂.
準備出戰場.
個個都很緊張.
這位牧師上到台的時候.
就跟校長說.
西點軍校校長說.
我害怕.
為什麼.
有些怯場.
不是說笑的.
見到這樣的情況.
但竟然.
西點軍校校長就說.
牧師.
你隨便說.
他有個題目.
就叫做.
叫做.
什麼呢.
叫做.
死好一點.
嘩.
真是可怕.
還在雙專場.
就說死好一點.
嘩.
這篇講章.
我很詳細讀.
讀了好幾次.
好精彩.
裡面說了什麼.
他說.
今天我們在這個世界上.
我們的生命.
是在神的手中.
不是在乎我們活得多長.
而是活得多好.
不是活得多長.
而是活得多好.

$^{681}$嘩,這句說話.
幫助我.
能夠再一次想到.
今天我們很多人活得長.
你猜活得好不好.
有些人有病.
又不知道怎樣過.
最慘就是.
臨去的時候.
不知道怎樣去.
真是很辛苦.
但這句說話.
It is not how long.
but how well.
嘩,這句說話.
提醒我.
今天我們生活在世界上.
我和你都有生命.
但耶穌說.
我要來給你新的生命.
原來為什麼.
現在這個不是生命嗎.
不是啊.
原來這個只是肉體的生命.
原來神給我們的生命.
不單止是肉體的生命.
耶穌說.
我來使你得到永生.
這個永生的意思就是說.
豐富的生命.
什麼叫豐富的生命呢.
不是在乎多長.
而是在乎多好.
喂,我們活得好不好.
活得好不好.
聖誕節就提醒我們.
怎樣叫做好.
我們有希望.
我們有平安.
我們有喜樂.

$^{721}$我們有信心.
這是聖誕節的訊息.
各位.
我希望大家.
都能夠在聖誕節.
今年的聖誕節.
能夠得到耶穌基督.
給我們最好的生命.
好不好.
這個生命.
耶穌親自答應過的.
我來.
不是做什麼.
是給大家生命.
什麼我們沒有生命嗎.
不是,我要更豐盛的生命.
有喜樂.
有平安.
有盼望.
還有.
生活得有意義.
有一個很有名的商人.
這個商人.
很成功.
但怎料有一天.
他的太太病得很嚴重.
他祈禱上天.
為何我好多東西都有了.
為何我太太有病.
就在這個時候.
他覺得不行,要想想.
結果.
他在神的面前祈禱.
神啊,我一生當中.
我營營役役.
做這麼多事,但我忘記了.
我為了什麼而活.
結果.
他寫了一本書.
這本書如果大家有機會.

$^{761}$跟曾牧師說說.
可以買到的.
就是《人生下半場》,看過沒有.
很多人沒看過.
這本書應該看過.
他發現原來我人生.
好像打了一場球.
我打完上半場.
我贏得很厲害.
入了很多球.
但還有下半場.
我打不打呢.
打啊.
還要入多少球呢.
已經贏了,還要入多少球.
不如這樣.
我要打一場有意義的球.
《人生下半場》.
各位.
在座當中.
年輕人不是太多.
每個人都差不多.
已經打完上半場了嗎.
打完上半場了嗎.
你究竟.
是贏還是輸.
我不知道.
但是下半場.
我們一起打一場有意義的球.
好不好.
今天我們打一場有意義的球.
有意義的球是什麼呢.
生命有價值.
活得有盼望.
活得喜樂.
活得平安.
葉華先生和他的太太.
這麼好的見證.
引發了我和大家.
不要說槍.

$^{801}$就是和大家從心裡.
說出分享的說話.
各位.
我.
為主工作.
已經.
六多年了.
今年.
現在我宣佈.
我87歲了.
有人會說.
你還出來說什麼.
但是我覺得.
能夠.
使大家有機會.
活得有意義.
活得好,我就很開心.
上帝讓我說了多久.
我就說多久.
不用怕,我就快.
我就快.
我說就快不說字.
我就快說完了.
但是很坦白說.
無論說多久.
多長.
我覺得最重要.
我們有沒有浪費了我們的生命.
有沒有.
今天我們浪費了很多.
我們浪費了時間.
氣人家.
浪費了很多時間在做無聊的事.
浪費了很多時間.
我們自己掛機.
浪費了很多時間.
自己憂愁.
憂愁一些我們得不到的東西.
各位朋友.
聖誕節提醒我們.

$^{841}$耶穌零.
是他人的生命.
這個生命.
是豐富的生命.
寶貴的生命.
我盼望各位都能夠得著.
好不好.
我現在.
想和大家一起祈禱.
如果.
你覺得.
受感動.
你覺得回應神對你的呼召.
你覺得.
希望得著這個豐富的生命的話.
請你.
和我一起祈禱好不好.
好,如果你想和我一起祈禱.
並且.
接受這個賜給我們生命的主.
耶穌基督.
請你在座位當中.
舉高手來好不好.
我為你祈禱.
舉高手來好不好.
舉高手來為你祈禱.
哇.
好.
這樣.
這裡有一條柱子.
我看不到.
各位不如這樣.
你們走出來好不好.
走出來.
站在下面.
看著我.
是.
葉維先生和太太請上來好不好.
好.
為了他們兩個今天的見證.

$^{881}$一個很好的.
一個真正的.
生命的故事.
我希望大家.
都能夠像他們這樣.
將來有八個孫.
(笑聲).
而且.
每一個都孝順.
每一個讀書都很厲害.
很重要的是.
上帝給你.
有恩典,有慈愛.
而且給你有希望.
給你有力量.
給你過得更加豐富.
還有誰呢.
不要再等了.
我們一起祈禱好不好.
還有誰.
還有.
各位.
在座當中.
是基督徒的.
請在座位當中起身.
是基督徒的.
在座位當中起身.
差不多全部.
還有誰坐在那裡.
沒有起身的.
即是他不是基督徒.
(笑聲).
請你幫助我.
鼓勵他出來.
還有.
還有陪他出來.
好不好.
是.
是,謝謝.
大家鼓掌好不好.

$^{921}$是.
好.
現在希望大家和我一起祈禱.
有些是第一次祈禱的.
沒有試過祈禱.
祈禱就是我們向天上的父.
講說話.
我們沒有試過.
為何天父都聽到我說話.
我們祈禱就是向祂講說話.
祂是聽到的.
祂是知道的.
請你試一試.
我講一句,你跟我講一句.
你想想,全間堂.
每個人都開聲講.
你不是天父上帝的感動.
我講一句,你跟我講一句.
天父上帝.
(天父上帝).
多謝你.
因為你賜給我們.
(因為你賜給我們).
主耶穌基督.
耶穌答應過.
祂來到世界上.
(祂來到世界上).
賜給我們有新的生命.
(賜給我們有新的生命).
主耶穌,我們願意得著你的生命.
(我們願意得著你的生命).
我們願意得著無憂無慮的生命.
(我們願意得著無憂無慮的生命).
得著豐富的生命.
(得著豐富的生命).
不是金錢的豐富.
(不是金錢的豐富).
不是物質的豐富.
(不是物質的豐富).
而是心靈的豐富.

$^{961}$(而是心靈的豐富).
多謝天父.
願你做我們的救主.
(願你做我們的救主).
帶領我們前面的道路.
(帶領我們前面的道路).
使我們越走越好.
(使我們越走越好).
越走越光明.
現在我們有憂愁的.
(現在我們有憂愁的).
我們願意放下.
(我們願意放下).
我們有失望的.
(我們有失望的).
我們願意放下.
(我們願意放下).
現在我們很開心.
(現在我們很開心).
因為耶穌基督.
接納我們.
做你的兒女.
多謝天父.
求你幫助我們.
(求你幫助我們).
在今年的聖誕節.
特別有意思.
因為有耶穌在我們當中.
(因為有耶穌在我們當中).
我們有耶穌做我們的救主.
(我們有耶穌做我們的救主).
多謝你.
求你保守我們.
從今天開始.
(從今天開始).
活得更有意義.
得著豐盛的生命.
(得著豐盛的生命).
多謝天父.
願你聽我們的禱告.

$^{1001}$奉主耶穌基督的聖名祈禱.
(奉主耶穌基督的聖名祈禱).
Amen.
Hallelujah.
各位朋友.
我們多謝天父.
今天我們這間堂.
充滿著喜樂.
因為很多人得著盼望.
得著喜樂 得著平安.
我希望大家繼續學習.
今天在神的面前.
有個決志的禱告.
希望你們以後.
繼續學習.
讀聖經 祈禱.
而且繼續回教會.
並在教會中學習更多.
聖經的道理.
請你們相信我.
從今天開始.
你們.
將會很開心.
因為有喜樂.
你記不記得.
我很多時候都說.
一天笑幾笑.
醫生不用.
是不是.
因為聖經說.
喜樂的心是什麼.
是兩樣.
不要傻傻的.
笑一下.
你看多漂亮.
原來上帝創造我們.
有個很漂亮的笑容.
有誰說.
笑不漂亮.
如果笑都不漂亮.

$^{1041}$就不用笑了.
我們多謝葉環生夫婦.
好嗎.
好 也多謝.
周牧師安排.
給我們機會.
周牧師宣佈.
大家先坐下.
在前面剛才決志的弟兄姊妹.
新朋友.
我們在203室.
我們預備了一個房間.
和大家聊聊天.
和講解一下.
信耶穌的一些細節.
現在我們過去203室.
謝謝大家 再見.
\newpage



\section{哥林多後書 8:1-5}
\label{sec:0LnkonZdJmk}
\textbf{2025.12.28《我拿甚麼奉獻我的主》哥林多後書8\_1-5(和修版) 李富成牧師}
\newline
\newline
連結: \href{https://youtube.com/watch?v=0LnkonZdJmk}{\texttt{https://youtube.com/watch?v=0LnkonZdJmk}} ~~~~ 語音日期: 2025-12-30
\newline
\newline
\hyperref[sec:fk3jYby_MrQ]{\small{< < < PREV SERMON < < <}}
~
\hyperref[sec:index_chronic]{\small{[返順時目]}}
~
\hyperref[sec:index_scriptual]{\small{[返順卷目]}}
~
\hyperref[sec:VGwZW_K3or4]{\small{> > > NEXT SERMON > > >}}
\newline
\newline
哥林多後書 8:1-5
\newline
\begin{longtable}{cl}
\hline
\hline
章節 & 經文 (和合本修訂版)\\
\hline
8:1 & \begin{tabularx}{0.7\textwidth}{X} 弟兄們，我們要把神賜給馬其頓眾教會的恩惠告訴你們： \end{tabularx} \\ \\ \relax
8:2 & \begin{tabularx}{0.7\textwidth}{X} 他們在患難中受大考驗的時候，仍然滿有喜樂，在極度貧窮中還格外顯出他們樂捐的慷慨。 \end{tabularx} \\ \\ \relax
8:3 & \begin{tabularx}{0.7\textwidth}{X} 我可以證明，他們是按著能力，而且超過了能力來捐助，主動 \end{tabularx} \\ \\ \relax
8:4 & \begin{tabularx}{0.7\textwidth}{X} 再三懇求我們，准他們在這供給聖徒的善事上有份； \end{tabularx} \\ \\ \relax
8:5 & \begin{tabularx}{0.7\textwidth}{X} 並且他們所做的，不但照我們所期望的，更照神的旨意先把自己獻給主，又給了我們。 \end{tabularx} \\ \\
[1ex]
\hline
\hline
\end{longtable}
$^{1}$今早有無講過呢D 說話.
喂 唔該.
麻煩曬.
有勞.
Thank you.
今早我由牛頭角ge 九龍灣站遠征入黎.
受惠於香港ge 鐵路網ge 完善.
其實所謂好遠ge 地方都係50分鐘ge 交通車程.
就係天水圍站gwo 度專車接送.
我過黎洪水橋呢個地方.
如果唔係我入jor 呢個迷陣裡面.
我唔知點樣搵黎架啦.
所以我好多謝曾牧師ge 安排.
好多謝劉弟兄ge 接送.
一黎到ge 時候就撞到展秋弟兄.
牧師喝杯奶茶先.
弟兄姊妹.
當別人係你身上做jor 一D 恩惠.
幫jor 你.
我地通常都會講.
唔該 麻煩曬 有勞 Thank you.
我ge 大孫女就黎五歲啦.
屋企係有工人姐姐ge .
D 細蚊仔都好叻ge .
好識洗工人姐姐做ye .
但係我見到我ge 女同女婿都識教佢.
我去食飯ge 時候.
我見到我大孫女又去洗工人姐姐.
lor ye  唔知lor 奶定lor 乜ye .
呢對我ge 女同女婿就一定stop佢.
停佢.
你漏jor 講乜ye 呀.
漏jor 講乜ye 呀.
英文呀嘛要.
Please 咁啦.
跟住我大孫女就好勉強咁.
Please 咁啦.
我地會教我地ge 孩子.
對gwo D 對我地有幫助ge 人.
講唔該 講多謝.

$^{41}$有禮貌 哪怕gwo 個係.
黎聘請ge 家庭傭工.
我地都要對佢十分ge 尊重.
同埋表示多謝.
今天下晝貴教會有感恩會.
你有冇諗住黎同上帝講多謝.
如果你話我想返去訓眼覺.
咁你就晏D 先訓.
應該黎到弟兄姊妹ge 群體裡面.
數算主因.
羅馬書第十二章第一節.
可能大家都好熟悉.
亦都可能識背.
你識背ge 就背.
你唔識背就一齊讀 開始.
所以弟兄們.
我以神的慈悲勸你們.
將身體獻上當作活祭.
是聖潔的 是神所喜悅的.
你們如此事奉乃是理所當然的.
史陶保羅去寫羅馬書.
係一卷非常之有邏輯性.
嚴謹ge 神學ge 論文.
佢去要討論一個好重要ge 信仰教義.
就係恩順 清義.
而先要指出世人都犯jor 罪.
虧缺jor 上帝ge 榮耀.
咁點算呀.
唔使驚.
當我地還作罪人ge 時候.
耶穌基督就為我地死.
係咪D 經文好熟呀.
係羅馬書黎架喎.
跟住講完整個救贖ge 大恩.
咁偉大 咁寶貴.
去到第十二章.
佢不足一轉就進入去.
弟兄姊妹呀.
既然上帝咁有恩典.
係讓我地白白清義.

$^{81}$咁我地應該點呀.
保羅就教導修信ge 教會.
要將每一個弟兄姊妹都要將自己.
黎到獻祭.
唔再好似舊約咁.
同洋生畜黎代替.
係自己黎到獻上.
而且係唔需要討論架啦.
咩叫理所當然呀.
劉弟兄黎接我係咪理所當然.
我應該架.
唔係架喎.
佢可以唔黎接我.
佢話比我聽穆斯林.
我搭咩車咩車咩車咪黎到囉.
所以當佢黎接我ge 時候.
我真係好多謝佢.
因為唔係奉旨架嘛.
有乜ye 理所當然架.
我1984年畢業.
收第一份人工即係教書.
我就比D 家用我阿媽.
使咩討論架呢件事.
大家認為呢.
仲使討論.
梗係要比家用啦係咪呀.
咩叫理所當然.
理所當然就係.
自然係咁樣做.
寶郎係度講緊我地弟兄姊妹.
點解我地會將身體獻上當做活祭呀.
係同上帝ge 慈悲.
呢個慈悲意思就係講到上帝.
將我地挽回.
福音本係神ge 大能要救.
一切相信ge 包括你同埋我.
我地都已經得到啦.
以至我地今日一齊係度聚集.
黎到去敬拜主.
咁我就想問大家ge 問題.

$^{121}$將身體獻上當做活祭.
包唔包金錢奉獻架.
講多一次包唔包先.
講得咁爽.
諗清楚先.
頂轉回.
金錢奉獻.
事實上係對好多基督徒黎講.
都係一個.
即係點呢.
我知架啦.
應該ge .
不過.
今日我地就同大家去破除下呢個.
應該ge .
不過ge 掙扎好唔好.
弟兄姊妹你係教會裡面.
常唔常見呢D 熟口熟面ge 東西.
奉獻箱啦.
全袋奉獻ge 奉獻袋啦.
呢個白色gwo 個係我成長舞會ge 月捐封黎架.
頭先我arm arm 撞到阿展秋弟兄.
我見佢zar 住個月捐封係貴教會ge .
好多教會都會透過月捐封ge 方式.
黎到鼓勵信眾.
特別係已經受浸ge 弟兄姊妹.
履行呢一個奉獻ge 責任或者權利.
通常一D 未信主ge 父母呢.
佢地好擔心返教會D 兒女有兩樣ye ge .
我自細呢就爸爸媽媽係傳統拜偶像呀.
即係逢年過節即係人地燒咩佢就拜乜gwo D 啦.
咁到我小五返教會ge 時候呢.
我爸爸媽媽都同意我去教會呀點解呢.
佢覺得去教會學乖呀嘛.
我媽第一句ge 回應就係.
跟你老師去教會好過你爛街呀衰仔咁樣.
因為我一時爛波地大架嘛.
爛波地大大家知係幾危險架啦.
咁所以華人ge 父母都知呢.
送D 仔女去教校讀書啦.

$^{161}$返教會都冇所謂啦.
不過一黎到呢兩件事就實係假啦.
第一樣叫受浸.
因為對華人父母黎講呢.
佢覺得受浸係真係入教啦.
你去教會玩下lor 教會著數冇乜所謂.
阿仔返教會睇黎都係鼓勵讀書ge .
讀書乖D 咁樣.
你讀唔掂呢D 哥哥姐姐會幫你免費補習ge .
咁所以D 父母覺得返教會冇問題呀.
等於去教會學讀書好ye 黎架.
不過當佢返返下返到有一個位.
有D 青年人呀特別返去同爸爸媽媽講.
我想受浸啦.
D 父母呢就可能有好大ge 反應.
差唔多呢如果你要受浸呢.
就世係假.
你要受浸呢你就唔好做我仔女咁.
差唔多要割裂個關係.
因為有D 父母會擔心乜ye 呢.
佢覺得你正式入教之後呢.
到佢百年歸老ge 時候呢.
聽講就唔裝香架嘛.
你D 教徒唔裝香架嘛.
話我淨係食你D 鮮花呀食唔夠食架喎.
我要裝香先夠食架嘛.
又唔燒ye 比我咁我又冇bans落去洗架嘛.
又唔燒gwo D 樓比我又冇麻雀燒比我咁.
無得打.
佢地好多佢地ge 迷思覺得.
我ge 仔女入jor 教呢就禮真的.
所以你一同佢講受浸呢佢地會有抗拒.
所以當我作為一個執事去默會ge 時候呢.
即係參與配合同工去默會呢.
有時同工都會分享.
有D 年青人都想受浸.
不過父母有好大ge 抗拒.
通常我地都會用緩兵之計ge .
即係鼓勵個青年人係屋企做好D 見證啦.
感化個父母啦.

$^{201}$無謂係受浸ge 問題上搞到好似即係分裂咁啦.
第二樣你估佢地最擔心係咩啦.
大家諗到啦.
奉獻呀.
香港ge 人呢就對基督教呢就識D 唔識D 架喎.
點解識D 唔識D 呢.
聽講你返教會要抽佣架喎.
即係要添香油架喎.
佢地會用D 觀念去諗.
就係好似你賺十蚊要比一蚊架喎.
佢地聽過架喎.
佢地知架喎.
佢地聽過咩十捐一架喎.
我地講ge 十一奉獻喎.
佢就覺得阿仔你讀書ge 時候.
返教會食教會ge 玩教會冇所謂呀.
到你做ye ge 時候.
我記得我1984年arm arm 出黎教書.
第一個月ge 月薪係四千七百蚊高唔高呀.
高呀係囉.
所以教書gwo 陣時真係好正架.
因為呢就有穩定收入仲係半日安黎架喎.
即係1984年仲係上下晝班ge 日子.
而家你D 書全部讀全日制啦.
gwo 陣時我D 同事真係一點鐘就走曬架啦.
所以我三點鐘走.
我校長已經覺得阿富城真係一個好勤力ge 同事.
我返屋企訓眼教架咋其實.
gwo D 日子.
如果十分一奉獻咁咪抽jor 四百七十五蚊去囉.
佢地好擔心兒女將薪金lor 返教會.
唔比家用.
唔供養佢地.
教會做jor 佢地ge 父母咁.
所以佢地好擔心.
我唔知你信主年青ge 時候有冇遇過D 咁ge 情況.
又或者你係未信主之前.
你D 仔女信jor 耶穌啦.
跟住佢做ye 呢你真係好擔心佢呢.
拎jor D 錢返教會.

$^{241}$當然我地都有聽過一D 比較極端ge 例子.
真係有D 教派呢.
係即係即係好似洗腦咁呢叫D 信眾呢.
真係出糧呢唔比家用父母拎曬D 錢去教會咁.
當然gwo D 係極端ge 例子.
弟兄姊妹我地中國人有句說話叫做.
講錢就.
你唔好同我借錢呀.
講錢就失感情啦.
耶穌講唔講錢架.
耶穌講.
係聖經裡面耶穌講ge 大概三十八個比喻裡面.
有十六個都同錢有關ge .
我地可以睇到耶穌對重視gwo D 跟隨者.
即係你同我作為耶穌ge 門徒.
佢對錢ge 態度係非常之重要.
例如我地好熟悉ge 呢一節ge 經文.
耶穌係登山寶訓呢個講道天國子民gwo 個大獻章ge 時候.
佢細不兩立ge 講jor 呢段說話.
一齊讀藍色ge 紙.
開始.
一個人不能事奉兩個主.
不是惡者過愛那個.
就是眾者過興那個.
你們不能又事奉神.
又事奉馬門.
馬門ge 意思即係財力.
咁係咪基督教係抗拒人地擁有金錢ge 呢.
唔係.
基督信仰對人擁有財富係體為上帝ge 賜福.
不過gwo 個擁有ge 方法必須要合實ge 心意.
舉一個好簡單ge 例.
我92年做校長.
個風聲都未公佈.
但係呢.
有一個書商已經收到風聲.
我會做校長.
佢第一時間黎要請我食飯.
咁老友你地.
大家諗到咩原因佢第一時間要黎請我食飯.

$^{281}$佢話你主任 唔係你主任 你校長.
我話你事慣一句.
請我食飯 好呀 過隔離gwo 間茶餐廳食.
佢話唔得.
校長嘛 梗係要去酒店食.
gwo 陣時我原任gwo 間校係沙田.
沙田有間叫帝都.
gwo 度食 咁好啦去啦去啦.
去到打開餐牌Set Lunch 囉係咪.
即係gwo D 標準ge 午餐.
佢話點得呀校長梗係食牛扒啦.
咁我話好啦好啦牛扒啦.
唔得呀 淨係牛扒要加隻龍蝦.
即係牛扒龍蝦餐.
咁樣款待我.
其中得一個原因ge 咋嘛.
你估佢真係同你老友呀.
咩原因呀.
因為希望同我建立呢種關係.
因為我係新教校.
我係最高領導人.
仲係講緊92年.
gwo 個選書ge 程序都未咁嚴謹ge 時候.
校長絕對有影響力.
用乜ye 書.
所以如果佢同我建立jor 呢種.
即係走肉朋友ge 關係呢.
我又好撚咁樣去享受佢提供ge 酒水呀.
甚至更豪華ge 款待.
咁佢個出版社.
我用佢出版社ge 書咪即係印銀紙囉.
丙姐妹.
當我如果用一D 唔合神ge 心意.
就係我地不能護者愛那一個.
而擁有ge 財富.
gwo D 就係呢節經文所講ge 事情.
但係如果話你做jor 校長啦.
於是你薪勞jor 啦.
於是你ge 人工呢就高jor 啦.
於是我收入高jor .

$^{321}$我奉獻多jor 呢個正唔正常呀.
我鼓勵你地代代當中.
你地如果係做生意ge .
你有好ge 業績.
你係工作上有好ge 升遷.
呢個絕對都係上帝ge 祝福.
錢同信仰有密切ge 關係.
我地對錢ge 態度係點.
我地同神ge 關係係點.
咁抽象ge .
具體D 呢.
具體D .
我地就用路加福音第十二章三十四節.
一齊讀開始.
因為你門的財寶在哪裡.
你門的心也在哪裡.
我有一個好心儀ge 男歌星.
叫做林子祥.
識唔識呀.
識呀.
睇你年紀個個都係林子祥gwo D 歌迷掛.
即係分分鐘需要你呀gwo D 咁.
咁呢.
係十月初佢開兩開幾場演唱會.
佢七十八歲架喇喎.
咁呢.
我就因為心儀ge 歌星.
咁所以我就去買飛.
買jor 一場星期六ge .
諗一諗.
聽得幾次得一次啦七十八歲嘛.
點知仲入得開呀.
唔好啦.
唔夠後ge 聽一場.
買多場.
星期日.
都比我買到.
而且每一場呢.
都買多一張請自己D 老友ge 一齊去.
買jor 四張.

$^{361}$咁呢.
我平時去聽演唱會呢係買最平gwo D 飛ge .
因為我淨係享受gwo 種感覺ge je .
前後冇乜所謂ge .
咁今次唔係.
聽得一次得一次.
想買最貴ge 飛買曬.
你唔好以為你有錢好多人都有錢呀.
跟住買九百八十蚊ge 飛.
嘩對我黎講都已經好貴呀.
因為我由gwo D 388咁呀.
一路買架咋嘛.
九百八十.
嘩咁九百八十我買四張喎.
洗幾多錢呀.
接近四千蚊啦.
其實都係坐山頂.
咁呢入到去之後.
第二晚.
坐我後面有大概六個人.
你知啦林子祥D 歌大家好熟嘛.
一起前奏呢.
其實全場都係度唱ge je .
點解後面gwo 咋人唱得好興奮.
不過全部都走音ge .
但佢地唱得好投入.
我唔好意思拎轉頭話唔該.
算啦人地入黎都係想高興下je .
咁就頂住自己唱大聲D 囉唯有.
我點解咁樂意呀.
點解咁肯比呀.
點解我冇肉痛呀.
點解我撳完gwo 個買飛.
過我信用卡ge 時候.
我一D 都冇摃心口呀.
因為我個心係度嘛.
我個心好想入場聽到林子祥用78歲.
仲唱gwo D 好快ge 歌好lum ge 歌嘛.
我個心係gwo 度.
所以我洗gwo 四千蚊面不改容.

$^{401}$我大孫女年幾之前就開始讀幼稚園.
咁我就問我個女.
佢報咩幼稚園呀.
我女就話.
我最後決定唔去英基啦.
去國際法國學校讀幼兒班.
法國國際學校.
係咪讀完幼稚園就出返黎讀中學呀.
唔係一路讀完中學架啦.
一路讀完中學.
咁即係洗幾多錢呀.
幾百萬啦.
我見到我女面不改容.
我心諗你細個ge 時候.
我係送你讀宣道會沙角幼兒中心.
沙前ge 沙角村.
差唔多平到差唔走.
好似唔洗錢咁.
咁你咪讀得幾好.
跟住你咪讀樽小.
跟住你咪讀樽中.
都係唔洗比學費架嘛.
你依家洗落個女gwo 度.
未讀大學已經洗幾百萬.
D 說話我全部係個套裡面講架je .
因為你知依家做四大長老呢.
千祈唔好對D 仔女ge .
佢地認為合理ge 行為.
又唔係違法ge .
佢又唔係要你比.
佢又比得起ge .
你咁多意見架嘛係咪.
你做長老ge 即係要識做啦.
咁我話我比jor 咁ge 表情之後.
我就冇講到我頭先套裡面ge 說話.
你細個咪讀沙閣幼兒中心.
我就話哦咁樣呀.
點解我女肯比呀.
佢個心擺晒係個女度囉.
佢覺得佢將佢所有比晒佢架啦.

$^{441}$依家我女其實後生D ge 時候呢.
即係大概七年前佢結婚ge 時候.
真係好似香港小姐咁靚架.
我女婿真係好似D 好型ge .
即係高過我ge 我女婿.
依家一個就變jor 阿叔.
一個就變jor 阿嬸.
因為呢個佢地已經唔再再注重自己啦.
佢將所有ge 資源放在一個大孫.
依家仲有個歲幾ge 細孫.
兩個就係佢ge 心之所在.
所以佢洗錢佢ge 時間安排.
佢做乜都好.
都係圍繞住呢兩個孫黎轉.
佢ge 心係gwo 度.
佢就擺啦佢ge 費用.
佢ge 金錢係gwo 度.
弟兄姊妹.
耶穌呢一個ge 比喻.
正正話jor 比我地聽.
你點樣去衡量你對耶穌ge 救贖.
佢比你ge 恩典.
去顯示gwo 種ge 溫度呢.
從你ge 錢點樣運用.
就係一個好好ge 指揮棒.
好好ge 溫度計.
剛才主席帶領我讀ge 呢段經文.
呢個用jor 和修版.
佢話比我地聽.
當時保羅寫信比一個教會.
叫哥林多教會.
其實係借用jor 馬其頓宗教會ge 奉獻ge 表現.
黎潤哥林多教會.
潤ge 意思即係責備佢地.
責備佢地乜ye 呢.
當時耶路撒冷教會大饑荒.
所以使徒保羅就發起呢個.
愛心奉獻ge 賑災大行動.
雖然愛邦人ge 教會.
同耶路撒冷以猶太人為主ge 教會.

$^{481}$種族上係有分別.
但係係主裡面係一家.
所以剛才太極ge 姐妹唱歌.
我地同感一鳴.
所以保羅就發動gwo D .
愛邦人ge 教會黎到去賑災.
但係哥林多教會呢間.
稍為具備能力ge 教會.
竟然係懶懶行.
所以保羅寫信去潤佢地.
去質疑佢地.
就用jor 馬其頓宗教會.
佢地樂捐ge 榜樣.
黎到刺激佢地.
馬其頓宗教會點解咁窮呀.
有可能係因為佢地當時係.
係呢D 愛邦城市裡面ge 信小眾.
呢D 信小眾唔單止係信仰上.
比人地定位為.
即係D 似一等ge 宗教信仰.
因為耶穌係當時係.
佢地認為係猶太教gwo 度演變過黎ge .
呢D 低劣ge 宗教黎ge .
相對佢地希臘ge 情況呢D 低等ge .
而呢種ge 排斥可能係.
涉及到佢地日常ge 生計ge .
譬如唔同你做生意啦唔請你啦.
所以可能當時ge 信眾為jor 要堅持信仰.
佢地會面對經濟ge 壓力.
所以佢地窮.
但係佢地窮jor 喎.
佢地係奉獻ge 事情上.
保羅話佢地係樂捐.
而且呢個樂捐係超過佢地應該能夠捐ge .
即係話你以為我捐jor 十蚊咩.
我捐十二蚊.
不過呢個提醒裡面有一個好重要ge 觀念.
係尾二gwo 行佢話.
佢地呢一個ge 奉獻ge 行為.
係照上帝ge 旨意先將自己獻給主.

$^{521}$如果弟兄姊妹你唔係諗住將自己同主有個咁樣ge 關係.
教會各種奉獻ge 呼籲對你黎講都係例行公事.
甚至你覺得厭煩.
又叫我捐錢.
如果我地唔先同上帝建立一個奉獻ge 關係.
教會各種ge 呼籲.
我大膽講.
你可以唔洗立.
因為保羅話佢地肯咁樣做.
係先將自己獻給主.
然後歸附我地.
即係意思佢知道奉獻唔係因為聽保羅ge 募捐ge 大行動呼籲.
係因為佢先獻給主.
然後對保羅ge 呢一個奉獻.
賑災ge 奉獻有感動.
然後作出一個具體ge 行為.
呢個次序係非常之ge 重要.
弟兄姊妹我地返教會返得耐.
因為我地對於奉獻可能都已經進入jor 例行公事.
而忘記jor 你奉獻.
如果你唔係感應同上帝有一個密切ge 關係.
你放落去ge 時候.
有陣時覺得係咁易.
有陣時覺得擺多jor .
早知唔係擺咁多.
你有冇試過賣旗.
gwo 個細路拎張賣旗.
諗住十蚊點.
咁就唔知點解十蚊點.
咁就有返肉赤.
如果弟兄姊妹係呢個奉獻ge 教導裡面.
知道我地係先黎到去回應.
感應主對我地ge 厚愛.
而我地作出回應.
呢個先至係我地奉獻ge 起點.
保羅隨後係哥林多後書九章五至七節裡面.
對奉獻再作進一步ge 教導.
我地只讀紅色ge 字一齊黎.
各人要隨本心所著定的.
不要作難不要勉強.

$^{561}$因為捐得樂意的人.
是神所喜愛的.
我地讀呢節經文呢.
有D 弟兄姊妹就誤解jor .
哦咁呀即係我鍾意捐幾多咪捐幾多囉.
唔好作難唔好勉強.
捐得樂意我唔樂意嘛.
你唔樂意真係唔好捐.
隨本心著定.
即係我諗住幾多咪幾多囉.
即係我咁呀.
如果你ge 奉獻.
保羅係度所鼓勵ge .
唔係叫我地你話點就點.
係你感應上帝ge 愛.
然後你按照呢個感應.
你作出ge 一個甘心樂意ge 行動.
所以你奉獻ge 時候冇肉痛.
你奉獻ge 時候冇為難.
呢個先係上帝喜悅.
你有冇留意黎到呢度.
保羅係兩段ge 教導奉獻經文.
都冇提到金額呢樣ye 架.
咩叫金額呀.
五百定一千.
冇提.
保羅淨係提一個奉獻ge 原則.
或者一個叫奉獻ge 動機.
或者再直白D 講.
係一個奉獻ge 態度.
如果今日A君奉獻一千蚊.
我奉獻一百蚊.
除jor 教會應該進行呢個保密.
唔應該用黎比較之外.
係上帝ge 眼裡面.
奉獻一千蚊同奉獻一百蚊ge .
只要佢地都係甘心樂意ge .
上帝都喜愛.
如果奉獻一千ge gwo 個.
好似比人扔心口咁先奉獻ge .

$^{601}$gwo 個奉獻一百甘心樂意.
好明顯個金額佢係大過佢十倍.
但上帝只會喜悅個個甘心樂意.
大家即刻會聯想.
耶穌係《福音書記》.
耶穌賺過一個寡婦ge 奉獻.
耶穌去到聖殿ge 時候.
見到有一個寡婦ge 奉獻.
奉獻jor 幾多呀.
兩個小池.
叮叮咁樣.
咁呀耶穌賺喎.
咁D 人覺得好奇怪.
點計都係好少je .
耶穌賺ge 係賺乜ye .
耶穌冇賺過金額.
耶穌係賺過比例.
耶穌ge 回應係.
佢已經將佢養生ge 全部奉獻.
金額係好少.
比例係好大.
係教會ge 奉獻裡面.
同樣弟兄姊妹.
我地奉獻唔需要問隔離gwo 個.
喂你幾多呀.
當你甘心樂意奉獻ge 時候.
哪怕你奉獻ge 係一個幾號.
你奉獻ge 係一千幾百.
或者你有某一個感動.
有一個大額ge 奉獻.
只要你係甘心樂意ge .
上帝都願立.
記住.
聖經係度講.
冇講金額ge .
我工作.
我以前工作係仁濟醫院ge 學校.
我依家仲幫緊仁濟醫院ge 學校.
做緊佢D 校監校董.
係佢地ge 慈善團體ge 文化裡面.

$^{641}$全部都要講金額ge .
做主席捐幾多.
做副主席捐幾多.
做總理捐幾多.
有D 人做總理唔做因為冇錢呀.
係教會ge 奉獻裡面.
從來都唔講金額.
講你ge 動機.
你ge 態度.
有一次係團契聚會裡面.
我當時我記得.
我應該係.
我仲係做緊執事.
即係D 教會D 大哥哥.
做緊團契導師.
有一次講奉獻ge 聚會.
有個青年人就問我.
我當時四十幾啦.
佢大概廿歲出頭.
開始做ye 啦.
佢話富哥富哥我地奉獻.
聖經講十一奉獻呀嘛.
我就想問呀.
究竟係除歲前十一呀.
定係比完家用ge 十一呀.
定係我供完樓ge 十一呀.
定係我供完車ge 十一呀.
定係我去完旅行ge 十一呀.
我話得呀弟兄.
呢個問題你收返.
我唔會答你.
你唔使奉獻.
你同唔同意我ge 回應呀.
當佢要諗咁樣ge 諗法ge 時候.
佢根本都唔想奉獻.
佢淨係想滿足教會一種形象.
就係人地問有冇十一奉獻.
我有.
但佢冇講到除歲前除歲後呀.
又比完家用係點樣.

$^{681}$佢只係滿足一種地上教會ge 一種提問.
黎到呢度.
我地有三個反思.
第一 頂子梅會諗.
係呀 十一奉獻.
呢個係舊約ge 標準.
我地新約啦.
進入恩典時代.
仲要履行十一奉獻咩.
申命記第十四章講gwo D 係.
律法書.
舊約.
仲要咩.
我地黎到新約仲要十一奉獻.
請你留意.
舊約為jor 比上帝ge 百姓.
有一個奉獻ge 基準.
就提出jor 十一奉獻係咪.
頭先同保羅gwo 兩段經文.
或者加埋所有保羅教導有關奉獻ge ye .
包括埋羅馬書十二章第一節.
黎到新約ge 奉獻ge 比例係幾多.
好似冇喎.
好似冇提到十一喎.
冇話十二喎.
冇話十三喎.
你有冇搞錯jor 呀.
保羅講緊咩ye 呀.
全部呀.
百分之一百呀.
新約難D 定舊約難D 呀.
如果淨係新約信徒進行十分一.
你可能好過D 呀.
我地就唔會唱gwo D 咩.
全所有奉獻ge 歌啦.
淨係唱十一奉獻ge 歌得啦.
我地唱緊ge 係全所有奉獻ge 詩歌.
呢個黎自新約聖經ge 教導.
我地係將身體獻上當作活祭.
黎到奉獻.

$^{721}$但係當我地話全所有奉獻.
我地都要有個起點架.
我地都有一個奉獻ge 操練ge 指標架係咪.
所以舊約裡面比神ge 百姓.
十一ge 奉獻仍然可以成為.
我地今日信徒奉獻ge 一個.
符合聖經ge 起點同埋指標.
當你ge 奉獻超越十一ge 時候.
你唔需要質疑我幾勁.
呢個係你理所當然架je .
我arm arm 受浸ge 時候.
就通常有D 叫迎新聚會架嘛.
你新抱歡迎會.
我仲好記得當時舞會一個大哥哥.
就同我地講.
聽住呀細路.
你而家未做ye 我都我受洗時中四ge .
所以十一奉獻對我黎講一D 壓迫力都冇呀.
即係冇乜收入呀嘛.
好少咋嘛.
佢話你而家呢就做學生就唔覺架咋.
日後你做ye 呢你就好覺架啦.
你賺一萬要捐一千架.
咁同我地講.
我當時個心諗咁傻架咩你.
鄭英前輩.
你估黎到今日我奉獻幾多呀.
我梗係唔會話比你聽.
因為呢個應該係私隱黎架.
不過我肯定話比你聽.
我而家所奉獻ge 比例.
比我做校長ge 時候更多.
點解.
其實好簡單咋嘛.
因為我離開校長工作之後.
我全面投入係事奉裡面.
我見到ge ye 機構.
教會信徒有需要ge .
我見得更多.
我亦都已經將自己擺jor 去備主所使用.

$^{761}$所以我奉獻ge 比率呢.
係高過我以前做校長.
咁好明顯我做校長gwo 陣時.
人工一定多過依家啦係咪.
起碼多一倍以上啦.
但係我而家發覺我呢件西裝.
係我做校長gwo 陣時穿到依家.
比你睇埋.
我呢條褲又係做校長gwo 陣時穿到依家.
我呢件牧師衫係我按立ge 時候D 頂尖會送ge .
咁我今日企出黎.
斯唔斯文呀.
使唔使錢呀.
唔使錢ge .
即係話.
原來你ge 物質生活.
唔係唔需要我地都要起個飲食.
我地作為牧師唔好出著到出黎烏烏ge 咁.
人地以為你爛仔.
但係係咪需要好名貴.
好多錢呢.
唔係.
一件深藍色ge 西裝.
黎港島又得.
去婚禮又得.
去喪禮都得架.
咩聚會.
真牧師燒.
我地一件深藍色或者一件黑色ge 西裝.
可以中環四海架啦.
我唔需要買到成個屋櫃都係西裝架嘛.
再加多套灰色咁都基本上夠架.
原來當我去進入jor .
一個甘心樂意去服侍主ge 時候.
金錢對我黎講已經係一個.
更加明顯睇到係上帝託管我.
黎到去為佢使用ge 一個器皿.
我係上帝私隱ge 器皿.
上禮拜五去一個戶外團契港島.
咁佢有居馬費比我架嘛.

$^{801}$港元費好笑呀gwo 日.
佢地話元旦崇拜呢佢地有步行呀即係入去浪劇.
咁我元旦呢就要去我舞會ge 崇拜我去唔到.
咁我完jor 有個感動.
雷牧師雷牧師我支持你步行.
咁我係銀包銀jor 舊錢出黎比佢.
唔係一舊即係一D 錢比佢.
咁跟住我返去呢對下我收到個港元費呢.
原來我奉獻比佢ge 係多過佢比我港元費ge 一倍.
咁你傻jor 呀你.
我冇傻我好開心喎點解呀.
神仲使用我有能力奉獻.
奉完獻之後你冇肉痛之餘仲好感恩.
呢個已經係從十一奉獻呢個操練開始.
黎到演化第二.
奉獻同眾籌有咩分別呀.
咩叫眾籌呢.
昨日星期六.
你有冇出街有人同你賣旗呀香港青年協會.
我係九龍區我唔知新界區有冇.
依家社福利署就係每每星期三星期六都比人賣旗ge .
以前係得星期六ge je .
咁呢就賣旗好肯定係一個眾籌ge 行動.
因為賣旗ge 人就係諗住做善事.
有陣時大家都唔係望清楚係咩機構.
總之有個細路仔拎住個奉獻袋.
唔係奉獻袋即係gwo 個賣旗袋呢.
特別細蚊仔呢.
我地呢D 長輩呢就好樂意去.
哄下佢幫佢買.
我母會都有一個自負盈虧ge 社福中心.
好似大家係呢度做ge 一樣.
咁我地呢就用母會ge 名義呢就申請賣旗.
當我母會賣旗ge 時候.
佢固然真係眾籌啦因為係社會服務.
但當我要去幫我母會買金旗ge 時候.
我係睇為奉獻ge .
因為對我黎講呢個已經成立jor 接近20年ge .
我地係何文田村裡面ge 自負盈虧ge 社福中心.
係教會一個進入社區服事ge 一個延伸.

$^{841}$係需要教會ge 弟兄姊妹去奉獻去支持.
我睇呢個賣旗對我黎講係一個金旗ge 奉獻.
我對佢係一個奉獻.
奉獻同埋眾籌一個最大ge 分別就係眾籌係咩呢.
我鍾意咪比囉.
我心情唔好唔比囉.
李富城個樣衰呀咁唔奉獻買你支金旗囉.
呢個小朋友好可愛我幫你買囉.
隨心所欲係咩呀.
我呢輪同教會D 弟兄姊妹開會嗌交.
我咪唔奉獻.
佢得罪我.
我發財立品我咪捐D 囉.
我做善長仁翁咪捐D 囉.
記住奉獻我地絕對唔可以叫自己做善長仁翁ge .
眾籌可以.
仁濟醫院gwo D 主席總理gwo D 叫善長仁翁.
如果你係紅印堂ge 奉獻你千祈唔好當自己善長仁翁.
你只係一個無用ge 僕人對處一個應該ge 擺相.
所以唔算得上唔可以稱自己係善長仁翁.
大家要睇自己係金錢ge 管家.
最後.
牧師我冇收入架.
我低收入架咁都要奉獻架咩.
你地教會係咪想將我捏呀.
渣幹渣正係咪係咪你D 牧師全都冇糧出呀.
弟兄姊妹.
奉獻再去返經文ge 解釋.
係冇金額ge 限制.
如果你依家係低收入你甚至冇收入.
你只係由D 兒孫比一D 零用錢你生活.
你甚至係lor 緊縱活.
你都可以參與奉獻.
奉獻係講.
態度唔係講金額.
我都參與不少基層企構或者教會.
佢地都係服侍緊好多老弱傷殘.
有疾病.
勞縮者或者係智障ge 朋友.
我都同樣同我地講奉獻ge 訊息.

$^{881}$就算你今日係受人接濟ge .
你受社會福利所保障ge .
你都有一個...你記住奉獻係你ge 權利.
教會唔可以剝奪你奉獻ge 權利.
李富城你lor 縱活奉咩獻呀你.
對唔住.
奉獻係我作為基督徒應該有ge 權利.
牧師比我奉獻.
呢個係我成長ge 舞會靚唔靚呀.
靚啦係咪.
我舞會明年100週年.
對我黎講一個好大ge 震撼.
因為我返jor 超過50年啦.
我一返我舞會兩個禮拜就搬黎呢一個座堂.
係1972年.
明年100週年.
黎緊下個禮拜四係我地進入100週年ge 元旦崇拜.
非常之感恩.
我問返D 前輩.
當時舞會都係深水Po 一D 即係基層教會.
點解會有能力.
係70年代能夠起一個禮拜堂呢.
佢話比我聽.
我地舞會係戰後.
即係四五年戰後重光之後慢慢D 弟兄姊妹聚集返啦.
開始有成型D ge 教會.
係50年代初期有個弟兄傻傻地.
奉獻jor 十蚊入黎執事會.
佢話我有個感動要建堂.
執事會D 人暈jor .
你捐十蚊入黎要建堂.
當時ge 十蚊可能都好大啦.
我就當你一千蚊啦.
十蚊等於一千蚊.
你捐一千蚊要我建堂.
所以當時執事會都唔知點算.
但又唔可以即係消滅jor 呢個弟兄ge 感動.
所以個建堂ge 意象係咁樣開始.
經歷jor 20年ge 過程.
裡面有不同年代ge 執事同工.

$^{921}$去為建堂黎到奔走.
最後建成.
前輩話我聽.
阿富城.
我地全部係弟兄姊妹一分一毫ge 奉獻.
冇大額奉獻.
冇gwo D 叫冠名奉獻.
冇全部都冇.
一分一毫代表乜ye 呢.
當時五六十年代香港經濟都仲係好差.
我地教會都未有所謂中產階級ge 出現.
全部有D 阿毛係收係金山寄錢返黎.
D 老公gwo D 錢.
係咁樣.
所以到奉獻禮gwo 日我已經返jor 教會啦.
我見到好多人流曬眼淚.
因為佢地ge 經歷ge 正如大衛王.
係歷代至上廿九章十四節.
大家知道大衛就想幫上帝起聖殿.
但係上帝就唔容許佢起.
就留比佢仔起.
不過大衛就預備好多材料.
係呢個預備材料ge 過程裡面.
佢發出jor 一次聚集ge 圖文.
佢話我算乜.
我ge 百姓算乜.
竟然能夠如此樂意奉獻.
因為萬物都從你而來.
我把從你的手得來的獻比你.
大衛呢個圖文正好表達.
好似係大衛同佢gwo 班百姓.
為上帝做jor 好多ye 係咪.
究竟邊個比面邊個.
大衛話係上帝你比面我地.
比我地有份參與.
係呢個見面ge 過程.
我地算得係乜ye 呢.
有個牧師就搵jor 三個信徒黎.
就覺得佢地呢.
係奉獻ge 事情上要教導一下.

$^{961}$咁於是阿A信徒就聽完牧師ge 教導.
阿A信徒就話咁啦牧師.
聽完你教導.
我係地下劃個圈.
然後我就同對又講啦.
全能ge 上帝你決定我應該捐幾多啦.
於是佢就將D 錢向上掉.
佢話如果D 錢掉返落黎個圈裡面gwo D 呢.
就俾我用ge .
掉jor 個圈外面gwo D 呢就捐比上帝.
咁你諗下佢會點掉.
佢會將D 錢集中咁向上掉係咪.
咁於是佢就諗住咁樣.
敷衍jor 牧師ge 教導.
我有做啊上帝唔收咋嘛.
咁阿B就睇唔過眼啦.
阿B就話咁搞錯啊你.
你睇下我啦.
我又劃個圈.
我就同你調轉ge .
上帝如果我掉ge 時候.
D 錢係落返黎ge .
垂直落返黎gwo D 奉獻比你.
外面gwo D 即係外圍gwo D 呢.
就俾我用黎開飯啦咁.
佢咪好叻D 啊.
咁你又諗佢點掉.
佢會仙女散花咁掉.
冇乜ye 係度落返黎.
咁於是阿C就話你兩個傻ge .
咁樣呃啊牧師.
你唔好呃自己欺人啦.
你睇下我啦.
唔好劃咩圈啊.
我就同上帝講.
我大聲偉大的上帝.
我將D 錢拋上黎.
你係偉大的上帝.
你要lor 幾多你咪lor 囉.
跌返落黎gwo D 比我用ge .

$^{1001}$阿C係咪最乖啊.
佢唔劃圈啊.
佢話上帝係全能ge .
上帝要lor 幾多咪lor 囉.
當然佢冇讀過牛頓ge 定理.
就係有萬地心吸力.
係咪啊.
掉上去落返黎gwo D .
鄭小沐議員.
我係1972年返舞會.
73年我就返兒童團契.
當時我小六.
兒童團契呢每次都有收奉獻ge .
當時我唔明咩叫奉獻.
就係覺得教會要捐錢.
但係我地奉獻係咁ge .
當時有個錢罌.
係間屋黎ge .
有個煙通ge .
上屋有個窿ge .
將D 錢係個煙通度掉落去ge .
一到奉獻時間呢.
個導師就話.
邊個小朋友想去收奉獻.
我最藝架嘛.
咁就舉手.
我放手.
富仔出黎收奉獻.
好.
我最鍾意拎住個錢罌呢.
就係度繞過哂D 團友.
全部人都要係我面前掟錢.
好高興咁樣自己.
叮咁樣時係抖鈴一號咁奉獻架嘛.
叮叮叮叮叮叮.
完一輪之後呢.
我又交返比.
將個錢罌交返比導師.
我地當時就唔係奉獻祈禱.
導師就帶我黎唱奉獻歌.

$^{1041}$就係呢一首.
呢首詩歌跟jor 我50年.
(唱).
當我小六ge 時候.
我唔知詩歌ge 意思.
我就係咁唱.
當我越長大.
從聖經ge 教導.
越知人生ge 經歷越豐富.
影響到我全職事奉.
過程裡面.
呢首兒歌成為我好重要ge 提醒.
當我地講奉獻ge 訊息ge 時候.
必定係講奉獻你ge 心.
你ge 心係邊度.
你ge 金錢係邊度.
跟我唱一次.
預備起.
(唱).
第三行佢話我若是博士.
我細個ge 時候就以為真係gwo D 學士碩士博士.
我後來知道呢.
我相信呢個典故係應該引用緊.
有幾個博士跟住粒星黎尋找耶穌ge gwo 個博士.
因為gwo D 博士黎到見到耶穌.
佢地做過咩行動呀.
獻上禮物.
黃金 乳香 抹藥.
弟兄姊妹.
我地今日黎朝見主.
我地係需要大禮物.
禮物按你ge 心.
同主ge 關係.
樂而 甘心 主必定喜悅.
我地同心祈禱.
多謝主你帶領我地今天早上.
特別係一年ge 年終.
讓我地知道你對我地ge 愛.
唔單止係拯救.
仲係有引導.

$^{1081}$同埋護臂.
讓我地深深去察驗主ge 愛ge 時候.
我地曉得點樣黎到將身體獻上.
當作活祭.
特別係金錢ge 奉獻上.
成為主所喜悅ge 管家.
黎到去拓展上帝ge 國.
我地祈禱奉主ge 名及.
阿們.
\newpage



\section{創世記 12:1-9}
\label{sec:VGwZW_K3or4}
\textbf{2026.01.04《起來》:創世記12\_1-9 (和修版) 曾敏姑娘}
\newline
\newline
連結: \href{https://youtube.com/watch?v=VGwZW-K3or4}{\texttt{https://youtube.com/watch?v=VGwZW-K3or4}} ~~~~ 語音日期: 2026-01-07
\newline
\newline
\hyperref[sec:0LnkonZdJmk]{\small{< < < PREV SERMON < < <}}
~
\hyperref[sec:index_chronic]{\small{[返順時目]}}
~
\hyperref[sec:index_scriptual]{\small{[返順卷目]}}
~
\hyperref[sec:8tVCcvQzY68]{\small{> > > NEXT SERMON > > >}}
\newline
\newline
創世記 12:1-9
\newline
\begin{longtable}{cl}
\hline
\hline
章節 & 經文 (和合本修訂版)\\
\hline
12:1 & \begin{tabularx}{0.7\textwidth}{X} 耶和華對亞伯蘭說：「你要離開本地、本族、父家，往我所要指示你的地去。 \end{tabularx} \\ \\ \relax
12:2 & \begin{tabularx}{0.7\textwidth}{X} 我必使你成為大國，我必賜福給你，使你的名為大；你要使別人得福。 \end{tabularx} \\ \\ \relax
12:3 & \begin{tabularx}{0.7\textwidth}{X} 為你祝福的，我必賜福給他；詛咒你的，我必詛咒他。地上的萬族都必因你得福。」 \end{tabularx} \\ \\ \relax
12:4 & \begin{tabularx}{0.7\textwidth}{X} 亞伯蘭就遵照耶和華的吩咐去了；羅得也和他同去。亞伯蘭離開哈蘭的時候年七十五歲。 \end{tabularx} \\ \\ \relax
12:5 & \begin{tabularx}{0.7\textwidth}{X} 亞伯蘭帶著他妻子撒萊和姪兒羅得，以及他們在哈蘭積蓄的財物、獲得的人口，往迦南地去。他們就來到了迦南地。 \end{tabularx} \\ \\ \relax
12:6 & \begin{tabularx}{0.7\textwidth}{X} 亞伯蘭經過那地，直到示劍地方，摩利橡樹那裡；當時迦南人住在那地。 \end{tabularx} \\ \\ \relax
12:7 & \begin{tabularx}{0.7\textwidth}{X} 耶和華向亞伯蘭顯現，說：「我要把這地賜給你的後裔。」亞伯蘭就在那裡為向他顯現的耶和華築了一座壇。 \end{tabularx} \\ \\ \relax
12:8 & \begin{tabularx}{0.7\textwidth}{X} 從那裡他又遷到伯特利東邊的山，支搭帳棚；西邊是伯特利，東邊是艾。他在那裡又為耶和華築了一座壇，求告耶和華的名。 \end{tabularx} \\ \\ \relax
12:9 & \begin{tabularx}{0.7\textwidth}{X} 後來亞伯蘭漸漸遷往尼革夫去。 \end{tabularx} \\ \\
[1ex]
\hline
\hline
\end{longtable}
$^{1}$各位弟兄姊妹平安.
我們再一次一同禱告.
今早真的有一個很大的意外收穫,我參加這個我們的教會一個月開始的星期日的祈禱會,好像平常一樣,我們都會分小組,大家一起祈禱,弟兄姊妹有什麼需要,都會為他們祈禱,互相代禱,其實是沒有什麼特別的事情..
但是今早,神透過我的小組給我一個很大的鼓勵,肯定,我之前要去預備這篇信息的時候,神讓我知道要說什麼,但是今早在小組裡面,很奇妙地,上帝透過一個很熟悉的長者,一個長輩,祂的禱告,祂的分享,再一次肯定,.
神要我去分享今天的信息.我心裡面很感動,這個長者去禱告的時候,一邊禱告就說,啊,神啊,今年我們不要再是一樣的了,你記住是一個長者,一個熟齡長者的禱告,我們不要一樣的,雖然我們的腳很無力,但不知道為什麼,我們可以多走一些路,為上帝你的緣故,我們的手可以多做一些工作,哇,很激動啊..
你知道嗎?我早上禱告,六點半到幾點,他回到那間教會,有神禱,在外面,他去不到,只能夠透過Zoom,你知道多甜蜜啊,我早上在禱告裡面,經歷上帝,哇,我想一下,神真的要振奮我的心,再一次,再一次的肯定,上帝要我去說,起來,起來什麼呢?.
是不是金堂的崇拜,大家全部站起來,我們一起做什麼呢?起來是做什麼呢?起來要回應神的呼召,要活出神所託付我們的使命,很多時候一說這些,就覺得,哎呀,那些在空中,某些人會說的,這些跟我們沒有關係,不是啊,弟兄姊妹,是跟我們每一個都有關係,.
是上帝今天很想跟我們每一個說的,洪恩嘉在洪水橋這個地方,我最近去上課的時候,聽到一個很振奮人心的一個訊息,洪水橋位於北部都會區,我們是其中一個發展的城市,.
北部都會區將來的時間會有250萬人口,全港,這個數字是全港人口的三分之一,而且這幾年,很多去申請公屋的人都會被派到這一區,所以我們當中將會有很多年輕的夫婦,很多的兒童,很多的青少年,這一區將會不一樣,.
在這個過程裡面,我想一下洪水橋有幸成為其中一間去盛載這一區的福音工作的教會,不會只是對著幾個弟兄姊妹說,哎呀,大家起來了,那幾個弟兄姊妹,十多歲的我們就出去傳福音了,我們現在好好預備,不是啊,是每一個弟兄姊妹都要預備,.
就好像今早那個屬靈長者說的,我們要不一樣,今年是要不一樣,不可以只是坐在這裡,等一下吧,坐在這裡都是挺舒服的,這些椅子是很舒服的,坐在這裡就是每一個星期,又一個星期這樣過,不是的,這個不是神對我們的心意,要起來,我們要明白神對我們的呼召是什麼,.
我們從一個人物的身上去看他怎樣起來的,起來是什麼意思呢?是不是一些口號呢?起來了,我們等一下完了那個堂道之後,哼幾個口號,是不是呢?亞伯蘭怎樣起來的呢?我們來看看他蒙召的經歷,.
亞伯蘭首先他一直都是在一個地方生活,是加勒底的吾爾這個地方,這個地方本身是拜月神的,所以可以這樣說,亞伯蘭絕對不是生於一個屬於上帝的一個這樣的家族裡面,他小時候應該和家人一樣都是拜月神的,.
在這裡,原來我們看回《使徒行傳》的時候,神真的不止一次去呼召亞伯蘭,第一次的呼召在《使徒行傳》七章二十四節那裡這樣說,當時那個時候亞伯蘭還沒有去到哈蘭,.
還在吾爾的時候,神就呼召他,第一次的呼召他要離開本地親族,往往是上帝要指示他的地方去,然後他和家人,那時候和他爸爸都有去到哈蘭這個地方,在這個地方住了一段時間,然後他爸爸就死了在那裡,這個是第一次的呼召,.
亞伯蘭也回應了,出去了吾爾,你看回這裡地圖上說是一千公里,一千公里當時沒有那麼多很先進的交通工具,並沒有,也沒有真的很容易,哎呀,又搬屋公司,.
亞伯蘭首先把我們家裡所有東西打包了,裝好了,搬屋公司先行,可能這個搬屋公司是駱駝隊所組成的,他們走著走著走著走著,於是亞伯蘭舒舒服服的去到哈蘭,不是,沒有這一切,不是一般的搬屋,絕對不是,.
在當中說的是家裡所有家當,連同生畜的財產,連同家人浩浩蕩蕩的,他從吾爾這個地方拔根而起,整家人經歷了一千公里,途中經過一些很艱難的地方,這樣去到哈蘭已經很不容易,.
好了,上帝的呼召沒有停在那裡,再來一次,第二次,這裡去到《創世記》,十一章三十二節開始,第二次的時候,亞伯蘭的爸爸過世之後,上帝對他說,你要離開本地,本族,父家,往我所要指示你的地去,第二次要再去遠一點,.
好了,這個時候我們想問,上帝這個呼召說什麼呢?你要離開本地是說什麼呢?很徹底的,本地是說他的國家,你要離開本族,是亞伯蘭的家鄉,然後要離開父家,就是他和直系親屬一起住的地方,由大的到一個很小的範圍,.
一直亞伯蘭都要離開,然後往上帝所要指示他的地去,他的國家,他的家鄉,他自己所住的地方,親友的群體代表什麼呢?這裡很重要的是什麼呢?.
從熟悉到陌生,從安穩到不肯定,幾十年時間在一個地方生活,吃的,住的,是一切生活裡所需要的一切,生活的習慣,我平時和什麼人一起相處,朋友圈,全部在一個地方,.
我不是只是普通的朋友圈,普通的親友圈,在那裡是很有感情的,很深厚的感情的,這個地方閉上眼睛,我都懂得怎樣去生活的,一個這麼熟悉的地方,神就是要他離開,.
不只是地方,要離開網絡,離開人群,離開親友圈子,去到一個完全陌生的地方,在那裡無親無故,現在的人,科技各方面都很發達,交通很方便,.
我想香港人移民都是很多了解的,多方了解,先看看我在當地有沒有親友,有沒有什麼資訊,上網看一大輪,去到中間的經理公司又了解一大輪,最好在那邊的親友又幫我們了解一大輪,.
我鋪好路了,我一定是搭好橋才去的,但是這個時候的阿伯賴沒有這些,什麼都不知道,完全去一個很陌生的地方,他由一個很安穩的地方,去到他完全不肯定,那裡的人怎樣的呢?.
會不會對我好呢?不好呢?很奸詐的,歡迎我去嗎?我到時去到會面對一個什麼樣的處境,天氣怎樣,各方面怎樣,我適不適應那裡的生活,我可以做些什麼?.
這個是一個很大的,很重要的考驗,非常艱難,今天絕對不是我們香港人,不是現代人所可以明白,所可以了解,還有說搬遷的日子非常艱難,.
一千公里之後還有下一步,很長很長的路途,經過曠野,曠野有什麼?有野獸,你不要想,晚上有街燈,有很多酒店投宿,沒有的,.
每去一個地方怎樣生活?支搭帳篷,一個地方支搭完帳篷,去下一站的時候拆了帳篷,是整個家當,一點都不好玩,很艱難,很累的,.
你兼顧的人很多,不只是自己一個人,兩袖清風,很容易去,整個過程很辛苦,神就是這樣叫他去,你說神知不知道?知,知他將會面對這個局面,知他這麼艱難,知他的考驗很大,但上帝就叫他去,.
神,這麼艱難都有的,還有這個呼召不單止是要離開自己的本地本族父家,很重要的重點是阿巴蘭,要使別人得福,.
在這裡,神不是說我呼召你,你離開本地本族父家,我給你很多很多數不盡的福氣,你接著就享福,你人生會很美好,你每一天都和你的家人享福,今天享這些福,明天享那些福,明天就這樣享受,你自己很滿足,你很快樂,什麼都是你自己,.
這個不是神所說的,神所說的很重要,在這個呼召裡面,直接是一個呼召是什麼?要使別人得福,阿巴蘭的後裔都要使別人得福,這個是神的心意,.
神要透過他的子民百姓,使別人得福,是很重要的,在這裡我們看回阿巴蘭的回應是什麼?.
英文並沒有說,她有很多掙扎,和上帝討價還價,跟著說我不想去,上帝這樣很艱難,英文很直接,就遵照耶和華的吩咐去了,就去了,很撇脫,就去了,阿巴蘭透過行動表明她的信服,阿巴蘭也透過她的行動表明她的信靠,.
她信得過這位上帝,就是上帝,可以帶領她一生的上帝,帶領她一家的上帝,上帝你吩咐過,雖然前面有很多不肯定的因素,很多艱難,很多困難,上帝你說過的,但是你是一個可信的上帝,我就跟你,就開了,很直接,在這裡,還有,我們看到兩次阿巴蘭是捉壇的,.
有一次我們看回第七節這個地方,耶和華向阿巴蘭顯現,就說我要將這個地賜給你的後裔,當阿巴蘭去到一個地方的時候,上帝說要將這個地賜給他,就是迦南地這個地方,阿巴蘭就在那裡,向他顯現的耶和華,捉了一座壇,.
捉壇就是我們親近神,向神去獻祭,去求問神,去敬拜神,在這裡也是尊崇上帝,這是第一次的捉壇,第二次的捉壇就是當阿巴蘭從那裡,他又遷到伯特利東邊的山,支搭帳棚的時候,在那裡他又為耶和華捉了一座壇,求告耶和華的名,.
在這裡他再一次的捉壇,我想說迦南地這個地方,是不是迦南地的人信上帝的,上帝就這樣呼召他,由一個派偶像的地方去到一個信上帝的地方,對不起,不是的,迦南地本身的人有自己的偶像,他們不是信上帝的,.

$^{41}$在這裡阿巴蘭捉壇的意義非常重要,不單止是他代表他的家人向上帝獻上感恩,向上帝去禱告,去求告上帝,去倚靠他,敬拜他,也是一個公開的舉動,在一個拜偶像的地方,讓眾人看到,他信的神是不一樣的,.
他不會敬拜當地人的神,是一個公開的宣告,很重要,阿巴拉罕去回應神的呼召,不單止是去,上帝你帶我去哪裡,我都去,但這個過程,他直接去尊崇神,有些人去回應神,都很多埋怨,都很多怨毒,很多苦澀,.
哎呀,好淒涼呀,哎呀,這樣那樣,那些不信神的人看到我們,真是很慘,信耶穌信得像你一樣,真是很慘,但在這裡,要讓人看到,我尊崇我的上帝,我上帝是偉大的,這個尊崇不容易,為什麼?.
神說,叫他離開本地本族富家,往我所要指示你的地去,往我所要指示你的地,即是哪裡呢?雖然阿巴拉罕大概知道方向,向著迦南地的方向,但實際是哪裡是不知道的,這個真是很辛苦,當你不知道哪裡,你還要走上去,真是很困難,需要一個很大的信心的擺設,.
需要一個對神很大的信靠,阿巴拉在這個竹壇裡,就是一個宣告,我信我的上帝,就算我不知道去到哪裡都好,我現在跟著他去,我高舉他的名字,這個是很重要..
我在這裡想起,我們的弟兄姊妹,每一個人背景都不一樣,我相信上帝對大家的大齡呼召都不一樣,真是每一個不一樣,每一個讀書,讀什麼科,讀什麼學校,接受什麼教育都不一樣,.
我們的身體不一樣,我們的性情不一樣,我們的恩賜不一樣,但我相信上帝呼召我們每一個,千萬不要說我現在收奴,沒有收奴這回事,我們回去見天父之前都沒有收奴這件事,上帝呼召我們每一個,那個召命還沒完結,.
要找回我們的召命,我今早在祈禱會很感動,就是我那組,有一個弟兄很感動,就是他原來一直一直有做義工,九龍醫院為此都要頒一個金獎給他,就是做義工的金獎,他還告訴我,是因為做義工的過程,他跌倒了,手受傷,我真的很佩服,.
神是用我們,很多時候我們經常看的那些,一定是揭光顛,一定是那些讀很多書,說很多話的人,特別是口才好的人,他就很多恩賜,上帝用他,所以當一講到服侍,你不要找我,我沒有什麼口才,我字都不識多,你不要找我,我不識的,不是的,弟兄姊妹,.
神給我們每一個都有很多恩賜,只不過是恩賜不同的意思,上帝今早就震撼我的身,你看看我們的弟兄,我們的姊妹就在做,在祝福人家,不只是在教會裡面,在外面不同的地方在祝福人家,有的,神有呼召給我們,.
在過去這幾個月,我是跟不同的弟兄姊妹合作,真的多了很多機會,特別是多了跟一些新的弟兄姊妹去聊天,去認識,一同合作去服侍,我心裡面很感動,原來神給我們教會來了很多弟兄姊妹,真的各自都有不同的恩賜,而大家都很願意擺上,.
我看到有弟兄姊妹是很有服侍的恩賜,譬如我會去,我知道有弟兄姊妹很想過來我們紅印堂聚會,但是行動未必很方便,我願意去接送她,甚至我崇拜的時候,我都配合她一起崇拜,當然這個服侍都不是今時今日才有,在之前都有弟兄是很忠心的,在這裡服侍,.
但是我想說,大家未必知道,有很多事情在背後默默做了,真的做了,大家未必知道,因為大家只看到燈光下的,你看到的人,但是其實原來我們教會有很多弟兄姊妹默默做很多工作,有弟兄姊妹是經歷神很豐厚的恩典,在那種苦難裡面,自己的生命裡面很多眼淚,在這一刻,她會願意錄下她的見證,.
很想和弟兄姊妹分享,神怎樣改變她,改變她的一家,也有弟兄姊妹在最近,我在感恩分享會也有分享過,我未必真的,你講策劃,很多計劃,但是你一吹哨,我們回來包東西,都可以的,很多東西,一起做,很同心,.
真的很多弟兄姊妹,很忠心的服侍,大家都未必知道原來他們擺上了這麼多,我最近和不同的人包過東西,和一班婦女包過東西,和兒童包過東西,和青少年包過東西,大家就是這樣,一路擺檔,我們成為一個包裝的工場,.
我都很肯定,如果有一天我們要建什麼社企,我們有老少中青,大家都可以來幫忙,我看到那個情況,就是大家都很忠心,沒有人說,哎呀,這麼辛苦的,這麼麻煩的,個個都很樂意,因為知道為神的家,為福音,為不同的緣故,.
在當中,我看到,大家怎樣擺上自己的能力恩賜,去祝福別人,上帝叫我們來做,我們就去做,我們當中有弟兄姊妹是很喜歡探訪,真的做很多很多探訪的工作,在不同的老人院,在不同的群體裡面,做很多工作,.
但我真的覺得神很想我們更多的人要起來,是每一個都要起來,將你的恩賜擺上,祝福別人,我們就想一下,祝福別人有哪幾個不同的層面呢?.
第一個層面,可以在教會裡面,教會有不同的服事,可能你說我嘴巴不行,如果要去參與教聖經,分享聖經,不是我強項,但是我很擅長維修,我們的燈,我們的水喉,廁所,這個都是我們教會非常需要,是很需要的,.
我有力氣,是很需要,大家回來搬燈,我有美感,我懂得排燈,我們很需要,你知道嗎?每個星期都要排燈,是很需要,.
我又很感動,我知道太與姊妹,將會一同去參與,每個月都有一次機會,一同去幫忙排燈,大家知不知道坐在這裡,能夠好好地去敬拜我們的上帝,舒舒服服,都不是必然的,都有人為這件事去付上代價的,原來將會有弟兄姊妹這樣,.
很多層面,你對我們的兒童有沒有負擔?你對我們的青少年有沒有負擔?你有沒有感動去做一個代禱用士?真是為不同的人需要,為教會的發展,為洪水橋區的科學工作,為神的國去禱告,不是禱告一會兒,.
是上上上上的禱告,如果有一個大班的禱告團隊,整間教會就不一樣,弟兄姊妹,在新的一年,我鼓勵你找回神給你的呼召,你生命裡面有什麼恩賜,你千萬不要走來跟我說你沒有恩賜,在座每一個人都有很多恩賜,所以下一次不要再跟我說你不行,.
只不過我們參與的時侯是不同的,真是很多,真是過去的日子,我要為大家獻上感恩,但需要更多的人起來,我們一起也要尊崇我們的上帝,在別人面前也要讓人認識我們的上帝是奇妙,是可畏的,.
所以在新的一年,我誠意邀請大家,也鼓勵大家脫離我們的負面情緒,脫離我們的沮喪,灰暗,怨毒,是非,我們在別人面前的敬拜不是一個造樣,是內內外外都是真誠的讓人看到,.
我站在台上,我一敬拜,之後我們很負面,很灰心,人們一進來的群體感受,咦,這樣的,很不真實,生命需要更新改變,今年神給我一個很大的字,在我預備這篇講章時,神放一個很大的字在我生命裡面,就是新字,新救的新,要做新人,一個新的人,.
無論我們因為什麼,令我們很沮喪,我們很灰心,我們很失意,在新的一年,要做新人,放下罪的重擔,因為我們所服侍的是耶穌基督,我們不是服侍人,我們服侍耶穌基督,將我們的生命交給主,為主而活,.
是新的人,內內外外一致,在人面前的見證,才會有真正的感染力..
在這裡我們看到,你說,哎呀,好像要靠自己做很多事情,不是,我們看回我們的上帝多好,上帝怎樣好?上帝親自鼓勵亞伯蘭,在這裡,剛才我讀過這段經文,第七節,要向亞伯蘭顯現,是親自的鼓勵,.
為什麼亞伯蘭領受神對他的呼召,一會兒我們會看神的應許,他接著去到迦南地這個地方,去到這個地方,他難免會受到影響,為什麼迦南地的人有什麼特色?.
高大畏盲,一去到,這個地方,上帝不是空城,一坐下,我就擁有這個地方,神都沒有說這個地方立刻是屬於他的,沒有的,這個地方有人住的,不是你一去到,立刻人們就四散,然後你擁有的地方,不是的,.
那個地方還是有人住,有人在管理,迦南地的人,迦南地的人高大畏盲,不容易,你一去到,看到現實環境是什麼?我又陌生,來到這裡,那些人又這樣,神所應承的是不是真的呢?.
這個對自己信心都是一個考驗,上帝在這裡親自向他顯現,鼓勵他,很重要,還要說,我要將這個地賜給你的後裔,即是現在亞伯蘭還沒有得著,給你的後裔,哇,上帝要他去的地方,他不是即時得到的,還要等,不知道要等到什麼時候,這是一個信心的考驗,.
亞伯蘭擺上,真是一個信心的考驗,在新一年我們都需要對神有這個信心,很重要的一件事,神的鼓勵很重要,神賜下英雄,這個非常寶貴,神呼召,神也會給英雄,.
神說,我必使你成為大國,我必賜福給你,使你的名為大,這個很重要,使你成為大國代表什麼?一個國家有什麼?有國民,即是亞伯拉罕會有後裔,會有國民,會有國土,有他自己的語言,有他自己的組織,大國,即是說他會有很多國民,.
上帝答應他,很多國土,他的組織很強大,他的軍事很強大,神答應會給他這一切,以色列國有他自己的語言,我必賜福給你,這個福是什麼福?很大的福,神給的福是不一樣的,.
一方面這個福是按照我們所需要的,給我們,阿伯拉這一刻心裡也在想,我都沒有後裔,你說應許我什麼?沒有後裔,神在這個賜福裡一定包括後裔,上帝呼召他離開以前的地方,那他現在都需要什麼?需要地方,需要土地,這個一定包括了土地,.
但不只是這麼少,是更多的上帝要給他,還有他的後裔一切的福氣,從這裡開始,經文裡所說的,不只是阿伯拉自己,也包括將來阿伯拉從他而出的國度,那個國家,神要他大大的賜福,使你的名為大,這個很重要,名為大是什麼?.
他的名聲很有名望,別人很尊重他,其實直到今天,亞伯拉罕的名字都很響,在以色列當中,在世界不同的地方,在不同的宗教當中,他的名字都很響,上帝所答應的,都是有實現的,不只是使他的名,還有阿伯拉後裔的名都是偉大的..
這裡做一個對比,大家記得在十一章的時候,提到巴別塔的事件,當時那些人很想建巴別塔,是做什麼?他們說,來呀,讓我們建造一座城和一座塔,塔頂通天,哇,很高,我們要為自己立命,.

$^{81}$紀念我們自己,免得我們分散在全地面上,我們的名字很重要,我們要立一個塔,讓別人想起我們的名字,我們人都很喜歡自己,很喜歡自己的名聲,被標榜的,被高舉的,我們很自我的,在過程中,人是這樣想的時候,最終的結局是什麼?臭名遠播,整個巴別塔的事件是失敗的事件,.
是人類歷史中的一個污點,但反而在阿伯拉後裔的時期,他沒有求神讓自己很出名,讓自己的名字被萬國萬邦高舉,沒有,上帝卻要他的名顯偉大,所以有時候,我一直在想一件事,我覺得這成為我們很大的鼓勵,弟兄姊妹,我們不要為自己求名求利,.
不用,謙卑在神面前,上帝自會讓我們高聲,說真的,也不應該為名為利去服侍,這是沒有意思的,但神就是這樣,要去高舉亞伯拉罕的名字,這是很大的應許,.
當亞伯拉罕回應神的呼召的時候,神就給他這麼大的祝福,今天當我們願意回應神對我們每一個的呼召,神都會賜福給我們,就好像今早那位長者,那位姐妹所說的,當我們都很願意去賜福別人,很願意去為教會擺上的時候,.
上帝賜福,我們腳不太行動,變得很有力,我的手不舒服,神又讓我行動,我感覺上不清醒,忽然間很精靈,神賜福,上帝給我們的福,不單給我們個人,也給我們的家庭,我們的子女,父母很多,神的應許很寶貴..
在這裡,神也讓我們看到他多麼愛我們,神呼召我們,神在當中也有祝福和奏訟,大家不要輕看自己的身份,我們是屬於上帝,上帝很愛我們,我們受傷被迫害,被攻擊的時候,上帝也憤怒,上帝也憤怒..
上帝也憤怒,上帝會為我們出手,傷害我們的人,也就是傷害我們背後的上帝,這句話真的要記住,誰傷害我們,上帝也不會放過他..
在這裡,神提出祝福和奏訟,為亞伯蘭尼祝福,我必賜福給他.這個祝福是說善待,好好對待他,對他有很多好的地方,誰對亞伯蘭尼好,神會賜福給他..
這兩個福是不一樣的,人比人的福氣有限,但是神的賜福是比人的祝福大,所以誰對你好,上帝就會對他很好,福氣很大..
但是,誰詛咒你,我必詛咒他.這個詛咒是說輕視,凌辱,對他做些不好的事,人也是會這樣對人的,但是上帝是會詛咒那個人,詛咒是比之前的詛咒大,上帝的詛咒是致命的..
我不希望我們被上帝詛咒,我只希望我們被上帝祝福,所以讓我們去祝福別人,讓我們不要去詛咒,這是非常非常寶貴的..
在當中,值得弟兄姊妹去想一想,在新的一年,你怎樣去祝福別人的生命呢?新的開始,放下罪的重擔,走前面的路,靠主賜福給我們..
由禱告開始,不是你覺得那件事很好,我又跟你做那件事,一會兒那個人又覺得那件事很好,我又跟你做,上帝是對我們每一個的旨意的..
有些事我們一起做,但有些事是你個人做的,問神.今早有些福事,但是我先說,你先禱告,你先禱告,你回來告訴我,上帝有沒有感動你做?沒有,你跟我說,我禱告完之後就沒有了..
是神分派不同的恩賜給我們,要讓我們用在神的家裡,有幾樣東西在新的一年,起來是起什麼來?.
回應神的召命不只是在教會裡,在教會裡固然我們找回我們可以服侍神的地方,盡情發揮我們的恩賜.第二樣東西,在我們家裡,神有沒有召命給你?.
因為神給你在你的朋友圈當中,在你家裡當中,你是首先信主的,請你回到你的朋友圈,你家裡當中,帶他們信主,我又很盼望,在你年底的時候,在你隔壁看到你的親友坐在你隔壁,這裡應該不夠坐了,這是第二樣東西,在你家裡..
第三樣東西,對於我們這一區,每一個我們都有份,我們每一個都有份,前一陣子有機會和一些兒童,青少年,都出去在附近派單張,在報佳音那晚,也有很多弟兄姊妹,成年的弟兄姊妹,長者的弟兄姊妹出去派不同的禮物,單張,這是一個新的開始..
大家有機會,我鼓勵你們走一走,在駱橋樓,Leo那邊,一路走過來,你感受一下整個洪水橋區是不一樣的,將來不只是那裡不一樣,整個區都不一樣,大家一路走一路禱告..
大家有沒有想過,出一分什麼力量?經過黃福苑的火災之後,我看到神預備人心,這個事件令人很痛心,我記得那晚我看新聞的時候,覺得很痛苦,很難過,我經常都很想那場火快點停,很擔心裡面還有多少人..
我相信很多香港市民和我的心情一樣,那晚很難過,但那一晚之後,上帝很大的工作,叫人心柔軟了..
我之前在洪水橋區派單張的時候,真的感受到人群很硬心,但最近派單張不一樣,人心柔軟了,上帝在這裡工作,弟兄姊妹,讓我們配合神的工作,起來祝福這一區的居民,祝福我們的街坊,不只是洪水橋區,.
洪水橋所處的地方,我們都可以接觸團圓天的街坊,居民,弟兄姊妹起來,整個教會起來的感覺是不一樣的,讓我們一同禱告..
天父啊,求你讓我們每一個看到,是我們每一個都有份,你很愛我們每一個,我們每一個都有很多恩賜,我們不需要在這裡和別人比較,因為上帝給我們恩賜,每一樣都很好..
我們可以用你給我們的恩賜建立神的家,我們可以在我們家裡傳福音,我們可以叫洪水橋區的居民去聽福音,可以叫團圓天的街坊朋友們都可以去相信你..
天父啊,這個禍場何其大啊,求你讓我們起來,天父啊,你知道我們很軟弱,我們有很多罪的重擔,求你寬恕我們在罪惡當中,讓耶穌基督以你的寶血去潔淨我們的罪惡,洗淨我們的罪惡,求你讓我們很有力,放下這些罪的重擔,.
走前面的路,讓我們做新的人,去成就上帝在我們生命裡的旨意.天父啊,求你賜給我們每一個都有亞伯拉罕那種信心,我們不知道前面會怎樣,不知道自己的經歷會怎樣,很多事情都不肯定,.
阿滿.
\newpage



\section{歌羅西書 4:7-9}
\label{sec:8tVCcvQzY68}
\textbf{2026.01.11《當保羅被囚》:歌羅西書4\_7-9 (和修版) 曾裔貴牧師}
\newline
\newline
連結: \href{https://youtube.com/watch?v=8tVCcvQzY68}{\texttt{https://youtube.com/watch?v=8tVCcvQzY68}} ~~~~ 語音日期: 2026-01-14
\newline
\newline
\hyperref[sec:VGwZW_K3or4]{\small{< < < PREV SERMON < < <}}
~
\hyperref[sec:index_chronic]{\small{[返順時目]}}
~
\hyperref[sec:index_scriptual]{\small{[返順卷目]}}
~
\hyperref[sec:sojB_TREQFc]{\small{> > > NEXT SERMON > > >}}
\newline
\newline
歌羅西書 4:7-9
\newline
\begin{longtable}{cl}
\hline
\hline
章節 & 經文 (和合本修訂版)\\
\hline
4:7 & \begin{tabularx}{0.7\textwidth}{X} 推基古是我親愛的弟兄，忠心的僕役，和我一同作主的僕人；他要把我一切的事都告訴你們。 \end{tabularx} \\ \\ \relax
4:8 & \begin{tabularx}{0.7\textwidth}{X} 我特意打發他到你們那裡去，好讓你們知道我們的情況，又讓他安慰你們的心。 \end{tabularx} \\ \\ \relax
4:9 & \begin{tabularx}{0.7\textwidth}{X} 我又打發一位親愛忠心的弟兄阿尼西謀同去；他也是你們那裡的人。他們會把這裡一切的事都告訴你們。 \end{tabularx} \\ \\
[1ex]
\hline
\hline
\end{longtable}
$^{1}$我們一起同心祈禱.
再一次我們將整個聚會.
剛才的訊息.
我們知道有宣教士.
有帶這個訊息過來我們弟兄姊妹當中的.
這樣一個人.
我們想道主.
你就是很想我們有訊息.
這個世代需要真確的.
和從主而來的那個召命訊息.
我們很盼望主幫助我們.
感動我們每一位弟兄姊妹.
以致我們能夠.
真的將我們的心靈.
來擺在主的救恩裡面.
神國度的那個擴展.
人能夠悔改信主.
這個重大的事業上面.
幫助我們能夠在主面前.
透過剛才兩位的見證.
親身的見證.
和我們在這個時間.
一個短短的訊息的回應.
求主帶領我們.
我們這樣禱告.
奉耶穌基督的名.
阿們.
很想和大家做一個這樣的回應.
這段聖經剛才William已經讀了.
很想和大家再分享一下.
這兩位年輕的弟兄.
保羅在他傳道的時候.
神讓他能夠行很大的神蹟.
在《小藤傳》第一次.
他在路斯德的時候.
他和阿巴拿巴見到一個人.
兩腳從來都沒有走過路的.
生下來是有腿的.
但是保羅突然間.
神給他有那種感動.

$^{41}$見到他.
有感動神要他治好這個生下來是有腿的.
在路斯德的人.
保羅就說我吩咐你站起來.
那個人就立即站起來.
來行走.
然後.
能夠叫死人復活.
《小藤傳》的二十章.
第九節告訴我們.
保羅在特羅亞港島的時候.
有一個少年人叫尤推古.
在三樓上面.
一直聽保羅講道.
睡著了.
就跌下來地下.
跌死了.
真的沒了氣了.
保羅抱著他之後.
神使他復活.
路加醫生記載好幾次保羅的行神蹟.
告訴我們.
他在菲律賓的時候.
第二次的傳道.
他奉耶穌的名.
來到有一個侍女.
就整天跟著保羅.
來騷擾他.
不斷告訴人.
這個保羅就是神的僕人.
講的是救人的道理.
由一個被鬼附的侍女.
來介紹保羅.
保羅為什麼不幫助你.
介紹你.
不用你講.
別人都告訴別人.
保羅可能覺得.
這樣可能對他一個騷擾.
一個妨礙.

$^{81}$保羅就來叫這個侍女.
裡面的邪靈.
要立即離開他.
這個侍女.
當那個鬼離開之後.
他就不能夠再幫助主人發財.
所以其實好幾段聖經.
都是事同行轉的.
保羅能夠行神蹟.
我們要問一個問題.
為什麼現在歌羅西書.
保羅要坐牢呢?.
為什麼不可以道門自己打開.
可以走出去繼續傳道呢?.
為什麼在這個時候的保羅.
甘於成為一個囚犯.
福音就受難了.
為什麼他不行神蹟.
讓事情繼續超自然地發生呢?.
不過要留意一件事很寶貴.
原來路加醫生告訴我們.
保羅在傳道行神蹟.
是很偶然的.
保羅在他的傳道裡面.
這一段聖經.
我們真的很需要大家一起學習.
跟今天兩位的見證.
是相當重要的聯合.
就是差遣和傳承.
保羅在這裡告訴我們.
這個二十章的第四節.
其實當時是一個大騷亂.
因為信了耶穌的.
在當時的地方.
很多拜偶像的那些人.
都信了耶穌.
製造偶像.
做生意的人失業.
所以工會就集體控訴保羅.
要抓保羅.

$^{121}$就去地方官告保羅.
終於找到一些年輕的弟兄.
就帶到他去那個已婚所的地方.
已婚所是一個很大的廣場.
這裡告訴我們.
勉強那些官員平定了這一次的暴亂.
保羅不在其中.
但是有一個管會堂的人被抓了去.
這個所提尼.
這個時候.
聖經告訴我們.
保羅決意去這個地方.
這裡說有一些人跟他一起.
大家看看有多少人.
比利亞人德羅斯的兒子索巴特.
帖撒羅尼迦人阿里達古.
和西宮都已經三個了.
特比人蓋尤和提摩泰.
還有就是我們剛才讀的第九節.
又有亞細亞人推基古.
和這個特羅菲摩.
總共已經有七個的年輕人.
在《少林寺傳》20章.
原來我們不為意.
原來歐洲不同國籍.
土耳其不同的國籍的民族的年輕人.
成為保羅的團隊.
這個才是宣教的恆常.
大大使用保羅的方法.
就是福音不是只扣實在一個人.
他的一生.
福音是透過他去傳遞給下一代.
下一代的中心的人.
繼續來接續福音的棒.
接續福音的工作.
所以教會.
無論是本地的.
或者Chris或者Kurt他們去做的.
就是要傳承給下一代新的信主.
這次有機會說一點點在埃及的見聞.

$^{161}$因為在土耳其.
他們來自不同的工作室裡的.
製作的年輕弟兄姊妹.
分享有些是埃及的.
有些是土耳其的.
其中一位埃及的姐妹.
是聽到福音廣播.
信主.
然後加入這個事工.
成為埃及的一位製作的節目.
兒童節目的姐妹.
好像如果沒有記錯.
二十歲多一點.
很年輕.
她就是聽到福音廣播.
他們的節目.
不是聽.
是看到這些節目.
來信主.
然後來參加這個事工.
很令我感動.
就是今天我們很需要有傳承的弟兄姊妹.
我們需要有一個聽到信息.
有這個信息.
剛才Chris差不多想哭.
其實我也想哭了.
聽到那個儀仗的需要.
其實我們真的需要傳遞儀仗.
教會是一個傳遞儀仗的地方.
讓人可以看到.
原來福音不是僅僅保羅.
所以他被囚.
不要緊.
保羅不祈求.
我要行神蹟.
好像彼得那樣.
門要打開.
天使幫我開門.
他不需要了.
因為他外面有忠心的弟兄.

$^{201}$推基古是忠心的弟兄.
還有亞尼西茂.
原來保羅一直有做的是栽培門徒.
他一生忠心的.
看重的他們都看重.
他們都看重.
感恩我們在去年開始.
我們開始有很多年輕人.
昨晚我們一起聚會.
都有十幾個年輕人.
我們需要一直將福音的真理.
傳承到下一代.
讓更多的人看到.
原來服事主是一件很寶貴的事.
Chris 當然他回教會比我久一點.
在感受我的時候.
但是其實我還是看著他.
由在職場到現在.
來參與宣教的事工.
真的看到福音.
可以慢慢改變一個人.
真的看到神的大恩和愛.
弟兄姊妹.
當我們看到世界動盪.
宣教的環境越來越困難.
教會也面對世代的交替.
神仍然在興起新一代.
讓福音繼續奔跑.
Z7的事工.
就是在這幅圖畫裡面的現代版本.
上一次真的有不同的國家.
製作中心的代表.
輪番來分享他們今年的異象.
很感動.
因為真的不同國籍的弟兄姊妹.
來分享他們的服事.
大家不要忘記剛才已經說了.
是一些很危險的地方.
來傳福音.
不像我們香港那麼容易.

$^{241}$真的很危險.
但是他們都願意這麼做.
因為他們真的看到福音.
改變他們的生命.
所以很盼望弟兄姊妹.
我們在今年一起來服事.
又說一點.
剛才我在地上.
跟街坊聊天的時候.
又收到一個短信.
邀請教會.
我們現在可以去探望.
去到家裡面探望.
在駱橋樓.
這些大埔的災民.
因為有250夥的居民.
分配了來我們駱橋樓.
所以我們已經開始可以去.
不只是在街上跟他們聊天.
而是我們由下個星期開始.
我們可以去到他們家裡.
去探望,聊天,聆聽他們的痛苦.
所以弟兄姊妹.
大家一起祈禱.
我們可以安排人去家訪.
我們等一下會報告.
2月7日.
我們就跟駱橋樓的經理合作.
我們在他門口.
會擺攤子,寫揮春.
讓一些街坊.
可以認識我們教會.
所以其實我們一直要將這些儀仗.
將這些訊息.
將這些救人福音的美麗訊息.
傳遞給弟兄姊妹.
所以有孫教士在我們中間.
就是傳遞.
傳遞之後我們就要傳承了.
所以我們大家一起學習.

$^{281}$所以在困難的地方.
興起新一代的信徒.
讓福音不被封鎖.
不被囚禁.
最後很想講這個亞尼西謀.
他曾經失敗逃跑.
他的心靈破碎.
原來他是一個奴隸.
他的主人就是這個肥利門.
他就是在哥羅西.
肥利門就是哥羅西教會的一位領袖.
但是這個亞尼西謀.
偷了主人的錢.
跑到了羅馬.
在羅馬坐牢.
誰知道坐牢的時候.
就遇見保羅坐牢.
保羅就帶著亞尼西謀信主.
然後告訴我們.
保羅差遣他和推基古兩個.
回到二佛所.
所以一個本來人生很混亂的.
現在被神將他的人生重新反轉.
成為新一代的見證人.
福音不單改變他.
更加使用他.
彭王弟兄姊妹我們一起學習.
最後有一個見證很想跟大家分享.
因為最近我都來看了保羅班德醫生的書.
了解他的情況.
其實認識這位醫生很多年.
他的父母是在印度相當虔誠的宣教士.
一生就在印度服侍.
保羅班德醫生也是在印度出生.
後來回到英國讀書.
一直學習學醫.
最終因為他從小看著父母.
對待麻瘋病人的那種愛心付出.
最後感動他.
專攻對麻瘋病的專門研究.

$^{321}$最終他是第一個發現.
原來麻瘋病人.
為什麼手腳會扭曲.
會有這麼多問題.
原來是因為麻瘋病毒.
侵害了痛覺神經.
就是說所有麻瘋病人.
是不會感覺痛的.
你踩到一塊石頭.
很尖的.
都不覺得痛.
不斷流血都不覺得痛.
所以是很恐怖的事情.
而且因為他扭曲的時候.
一般人就會很害怕.
就會離開他.
所以原來班德醫生.
他有一次很深的經歷.
有一個麻瘋病人康復了.
不再感染不再傳染了.
他信了主.
很想去教會.
班德醫生就帶他去教會.
大家可以想像得到.
他坐在那裡.
教友就一直往後退.
個個都不敢見他.
而且很害怕.
原來班德醫生說.
最大的痛楚不是病毒.
是人的冷漠歧視.
令病人的心靈更深的痛.
所以班德醫生一生就是這樣.
用他的精力去研究.
怎樣幫這些麻瘋病人.
做一個外科的矯形復健.
一個醫生要學習怎樣製造鞋子.
怎樣製造鞋子.
最好不要令那些麻瘋病人.
穿著就不舒服.

$^{361}$因為他不會痛的.
所以可以想像得到.
原來可以這樣為自己的生命擺上.
當然醫生已經回天家了.
一位很敬虔的醫生.
我想說的是.
保羅看到的是一個不是壞透的生命.
看到的是可以改變的亞尼西牟.
帶著他信主.
成為福音的見證人.
這就是聖經福音的偉大之處.
我們一起禱告.
天父啊.
多謝你給我們洪恩家.
真是很有恩典.
我們需要有很多的學習.
當然在本地.
我們現在有很多新的屋苑.
在我們附近落成.
我們真的好像做都做不及.
剛剛才說我們下個禮拜.
可以去探望這些災民的家庭.
天父啊我們懇求你.
藉著今天的訊息.
讓我們看到.
原來我們就算多難的生命.
都可以被福音改變.
可以被主的愛來感動.
使我們能夠傳福音給有需要的人.
多謝你讓我們不單止本地有這個意象.
我們更加看到弟兄姊妹來自世界各地的.
當他們明白了福音之後.
就是這樣很想在他們自己的地方.
來傳福音.
我們很需要.
我們大家一起有這樣的等候.
禱告,裝備.
我們真的很想能夠一起參與.
願主聖靈感動我們每一位.
我們再一次獻上感恩.

$^{401}$給我們有很豐富的屬靈的筵席.
這樣禱告奉耶穌基督的名.
\newpage



\section{路加福音 10:25-37}
\label{sec:sojB_TREQFc}
\textbf{2026.01.18《走近你不想走近的人》:路加福音10\_25-37 (和修版) 何世傑牧師}
\newline
\newline
連結: \href{https://youtube.com/watch?v=sojB-TREQFc}{\texttt{https://youtube.com/watch?v=sojB-TREQFc}} ~~~~ 語音日期: 2026-01-20
\newline
\newline
\hyperref[sec:8tVCcvQzY68]{\small{< < < PREV SERMON < < <}}
~
\hyperref[sec:index_chronic]{\small{[返順時目]}}
~
\hyperref[sec:index_scriptual]{\small{[返順卷目]}}
~
\hyperref[sec:wcZQh6png8E]{\small{> > > NEXT SERMON > > >}}
\newline
\newline
路加福音 10:25-37
\newline
\begin{longtable}{cl}
\hline
\hline
章節 & 經文 (和合本修訂版)\\
\hline
10:25 & \begin{tabularx}{0.7\textwidth}{X} 有一個律法師起來試探耶穌，說：「老師！我該做甚麼才可以承受永生？」 \end{tabularx} \\ \\ \relax
10:26 & \begin{tabularx}{0.7\textwidth}{X} 耶穌對他說：「律法上寫的是甚麼？你是怎樣念的呢？」 \end{tabularx} \\ \\ \relax
10:27 & \begin{tabularx}{0.7\textwidth}{X} 他回答說：「你要盡心、盡性、盡力、盡意愛主—你的神，又要愛鄰如己。」 \end{tabularx} \\ \\ \relax
10:28 & \begin{tabularx}{0.7\textwidth}{X} 耶穌對他說：「你回答得正確，你這樣做就會得永生。」 \end{tabularx} \\ \\ \relax
10:29 & \begin{tabularx}{0.7\textwidth}{X} 那人要證明自己有理，就對耶穌說：「誰是我的鄰舍呢？」 \end{tabularx} \\ \\ \relax
10:30 & \begin{tabularx}{0.7\textwidth}{X} 耶穌回答：「有一個人從耶路撒冷下耶利哥去，落在強盜手中。他們剝去他的衣裳，把他打個半死，丟下他走了。 \end{tabularx} \\ \\ \relax
10:31 & \begin{tabularx}{0.7\textwidth}{X} 偶然有一個祭司從那條路下來，看見他就從另一邊過去了。 \end{tabularx} \\ \\ \relax
10:32 & \begin{tabularx}{0.7\textwidth}{X} 又有一個利未人來到那裡，看見他，也照樣從另一邊過去了。 \end{tabularx} \\ \\ \relax
10:33 & \begin{tabularx}{0.7\textwidth}{X} 可是，有一個撒瑪利亞人路過那裡，看見他就動了慈心， \end{tabularx} \\ \\ \relax
10:34 & \begin{tabularx}{0.7\textwidth}{X} 上前用油和酒倒在他的傷處，包裹好了，扶他騎上自己的牲口，帶他到旅店裡去，照應他。 \end{tabularx} \\ \\ \relax
10:35 & \begin{tabularx}{0.7\textwidth}{X} 第二天，他拿出兩個銀幣來，交給店主，說：『請你照應他，額外的費用，我回來時會還你。』 \end{tabularx} \\ \\ \relax
10:36 & \begin{tabularx}{0.7\textwidth}{X} 你想，這三個人哪一個是落在強盜手中那人的鄰舍呢？」 \end{tabularx} \\ \\ \relax
10:37 & \begin{tabularx}{0.7\textwidth}{X} 他說：「是憐憫他的。」耶穌對他說：「你去，照樣做吧！」 \end{tabularx} \\ \\
[1ex]
\hline
\hline
\end{longtable}
$^{1}$弟兄姊妹平安.
平安.
很久不見了.
我們先有個祈禱.
我們可以一同敬拜讚美你.
願主幫助我們.
因著你的話語.
叫我們生命得以改變.
叫我們對你的信心得以增加.
求你幫助我們.
你的話語實在太寶貴.
願主祝福.
祈禱.
奉耶穌名求.
阿們.
今天講題叫做.
走近你不想走近的人.
是的.
路加福音第十章的故事.
其實我們應該也很熟悉的.
就是那個律法師要試探耶穌.
問.
誰是我的鄰舍?.
耶穌沒有給他一個定義.
而是講了一個故事.
關於一個看見和走近的故事.
就是一個人從耶路撒冷到耶利哥.
途中遇見強盜.
被人打得半死就死了.
吊在路旁.
耶穌安排了三個人出現.
祭司.
利未人.
撒瑪利亞人.
他們也是看得見受傷的人.
但是他們的反應是不同的.
祭司看見.
咦?.
啊.
咦?.

$^{41}$哦.
走了.
OK.
另外那個是利未人.
也看見.
哦.
又在那裡.
又走了.
OK.
但是撒瑪利亞人看見.
哇.
看見了.
動了慈心.
就走到他身邊.
這三個人.
也看見了.
但是只有一個人走近他.
看見不是難事.
我們看見很多事情.
走近.
是需要勇氣.
在生活裡面.
我們每天也看見很多需要.
我們一打開電視.
其實你看見很多人很有需要.
那個地方有意外.
那個地方有災難.
看見很多需要.
回到家裡也可以.
又看見很多需要.
家人.
哇.
為什麼兒子回來是這樣的呢?.
我女兒也大學三年級.
她回來的時候.
她養完黑了的時候.
發生什麼事情.
看見了.
或者你上班的時候.
你的同事.

$^{81}$啊.
也是很不開心的時候.
你也看見他可能有些壓力.
可能他收到大信封之後.
有什麼壓力呢?.
你看見了.
教會弟兄姊妹.
有些人是默默坐在那裡而已.
我不知道是誰.
每一間教會也有幾個人坐在那裡.
他是很孤單的.
很有趣的.
不知道為什麼.
他會回來的.
但是他靜靜坐在那裡.
很孤單的.
又看得見他.
在長洲的裡面.
很多涼亭.
很多伯伯.
我也快要伯伯了.
我還沒有伯伯的.
看見他們坐在那裡看著海.
其實也挺寂寞的.
就是說你覺得他很享受.
其實是很寂寞的.
看見不難.
走過去才難.
為什麼?.
因為走過去.
需要有承擔.
你走過去.
可能會有麻煩出現.
你走過去.
可能要付出你的時間.
你走過去.
可能會被拒絕.
你走過去.
可能覺得不舒服.
所以我會說.

$^{121}$如果我們有一個人在那裡坐著看著海.
他會覺得很寂寞的.
但是他坐在那裡看著海.
他會覺得很寂寞的.
但是他坐在那裡看著海.
他會覺得很寂寞的.
但是他坐在那裡看著海.
他會覺得很寂寞的.
但是他坐在那裡看著海.
他會覺得很寂寞的.
我不是在聖殿.
你假設是在聖殿.
做完事情了.
很想回家休息.
是不是?.
又講完.
OK,有很多服侍.
又唱歌,自招.
加上主日學.
還有祈禱會.
還有團契.
一天到尾.
很累了.
我要回家休息.
很怕捲入那麼麻煩的那些.
也可能認為.
這些不是自己的責任.
看過就走了.
唯有那個撒瑪利亞人.
看見之後.
沒有理會風俗.
如果你知道他有什麼背景.
撒瑪利亞人和那些人.
不可以在一起的.
碰也碰不碰.
沒有計算風俗.
沒有計算代價.
沒有退縮.
而是靠近那個傷者.
觸摸他.

$^{161}$冒著危險.
為什麼會冒著危險呢?.
那個人被人打到趴在那裡.
那條街其實也挺危險的.
OK,是挺危險的.
可能那些強盜.
可能躲起來看看是誰去再搞的.
也還可以.
那裡是危險地區.
是危險街.
他就這樣.
肯停在那裡.
然後就幫他.
甚至花錢.
花時間來陪伴.
耶穌清楚提醒我們.
信仰不是背多少經文.
或者讀多少神學.
就像我讀很多書.
不是的.
而是在問我們願意走近多少.
你願意走近嗎?.
走近呢.
又不是什麼英雄行為.
就是他走過去幫他.
有沒有人拍照?.
沒有的.
沒有人理會他.
不是一個英雄行為.
只不過是一個憐憫的表現.
耶穌要我們做的就是走近他.
為什麼撒瑪利亞人願意走近他呢?.
走近他呢?.
因為什麼呢?.
動了慈心.
什麼來的呢?.
什麼叫做動了慈心呢?.
其實他的字大概是這樣的.
就是他的心腸.
什麼呢?.

$^{201}$震動.
OK.
震啊震啊震啊震啊.
不是說腦袋.
我看見他.
我先找AI分析一下.
不是的.
這裡是說他的心腸有震動.
譬如有時候我們也會震震.
是心有震動.
憐憫不是單單知道他有多慘.
他很慘.
而是心被觸動.
自己也有一點痛.
憐憫不是理性的.
而是生命和生命裡面的相連.
撒瑪利亞人看見傷者的時候.
不是想自己會不會有麻煩.
而是想.
如果躺在那裡的是我.
我也很希望有人來幫我.
這是憐憫.
就是愛淪捨如己.
耶穌要一種是要將別人的痛.
看成自己的痛的生命.
因此.
他願意走近.
包紮.
扶起.
陪伴.
照顧.
弟兄姊妹.
走近他.
不是一種什麼技巧.
不是什麼特別技巧.
是不是要受過一些訓練呢?.
不是.
走近他.
是我們的選擇.
你選擇而已.

$^{241}$就是那些很厲害的人看見.
他選擇離開.
他選擇而已.
我們也是選擇.
走近他是一個選擇.
選擇我們的心成為柔軟的心.
不要選擇一個爛漠的心.
走近他不是很華麗.
但是很寶貴.
瑪利亞人做的事情其實很普通的.
其實他做的事情很簡單而已.
他有沒有施行神蹟呢?.
一走近他.
我奉的命.
你會好起來.
說了很多事情.
沒有.
有沒有說神學呢?.
根據創世記的說法.
人被創造的說法.
然後說了詩篇.
又說了先知書.
然後大先知書,小先知書.
說了新約.
那一堆事情給他聽.
沒有.
十屆也沒有說.
什麼也沒有說.
沒有說神學.
其實他沒有什麼說的.
沒有說過什麼.
一走近他.
做了一個很簡單的事情.
就是走近他.
包閘.
送到他去旅店.
用支付寶.
是不是?.
就是幫他付錢.
然後呢.

$^{281}$還用了信用卡.
就是怎樣呢?.
我回來之後.
如果他不救你.
你已經幫我刷了卡了.
才能結帳.
厲害.
其實這些事情.
任何人都做得到.
不是很難的事情.
今天我們看到新聞或者世界.
很慘很慘.
我們可能無能為力.
我們也不能夠改變.
那些不知道什麼打來打去的政策.
那些也不知道怎麼搞.
也不能夠處理很多問題.
其實我們是有限的.
我們怎樣去處理問題呢?.
但是我們可以做一些微小的.
但是可以很深刻的事情.
陪伴.
聆聽.
祈禱.
探訪.
發一個WhatsApp.
幫助一下他.
有一些愛心的奉獻.
神體重.
不是你做得多少.
而是在問你願意走近有多少.
耶穌沒有說.
你用港片道來安慰他.
沒有.
耶穌只是說.
他怎麼做.
你就走過去就行了.
我們不是被呼召去解決所有痛苦的問題.
我們沒有的.
不是這樣呼召的.

$^{321}$而是被呼召去靠近那些受傷的人而已.
信耶穌的忠心.
不是你知道有多少.
而是你願意走近有多少.
那些律法師很懂得律法.
又懂得神學.
又懂得答案.
OK,你問我多一點我也會回答你的.
是的.
我也會回答你的問題的.
但是耶穌最後問的是.
不是問他.
喂.
你問那個.
他說李孜陽.
不是問他懂得什麼.
而是說.
沒有問他你知道多少.
只是問.
你應該怎麼做.
你是誰.
信仰不是一個資訊.
就是回教會學習多一點資訊.
而是你的生命有沒有改變.
你可能懂得很多道理.
但是頭很大.
OK.
但是沒有走到那裡.
那些道理就好像懸空在上面一樣.
撒瑪利亞人可能沒有讀過神學.
但是實踐律法裡面.
最核心的命令.
就是.
愛神.
愛人.
愛神.
愛人.
走近.
是信仰的一個真實的樣子.
走近不是你要去做一些偉大的事情.

$^{361}$走近是從你現在的地方開始.
走近你家人那裡.
是不是有些家人是不想走過去的.
我也有些.
有的.
走近被忽略的弟兄姊妹那裡.
你不是很想見到他們.
有沒有?.
我也有些.
有的.
走近一些孤單的鄰居那裡.
你知道他們一個人住.
走過去.
走近一些.
情緒裡可能有掙扎的人的裡面.
走近教會.
坐在那裡.
正正坐在那裡的人.
你走過去.
走過去.
其實沒什麼的.
可能只是坐在他旁邊而已.
有時候我們會覺得.
我們幫不了他們.
但是耶穌基督說.
你能夠給的.
不是解決方案.
而是你自己.
坐在他旁邊.
在別人的苦難裡面.
你能夠帶去的是一份陪伴.
陪伴就是走近.
走近就是一個有愛的樣子.
走近就是憐憫.
我們可能會問.
神呢.
要我們怎樣去愛倫社.
可能會問一個問題.
怎樣去愛倫社.
其實聖經裡面的舊約利未記第十九章.

$^{401}$早就已經說了一些事情了.
他講了一個故事.
講了一些怎樣去做.
教我們怎樣去做.
不是很抽象的理論而已.
就是要憐憫.
就是愛愛倫社.
不是很複雜的.
在利未記的第十九章那裡.
九至十八節那裡.
其實講了一個故事而已.
那個內容大概是講一個田.
田地.
有些果子.
有些十歲.
生活.
的經文.
就是在講命令以色列人.
去修國的時候.
就是他們有很多地.
又種了很多東西.
修國的時候.
要求他們.
田葛不可割盡.
一塊田.
OK.
那個葛不可以全部割掉.
葡萄園不可以摘乾淨.
掉到地上的那些不可以撿起來.
要留給貧窮的.
寄居的.
弱勢的.
為什麼呢?.
因為神說.
我是耶和華.
你們的神.
為什麼要這麼做.
答案就是.
因為我是你們的神.
挺有趣的.

$^{441}$就是沒得問為什麼的.
為什麼我要愛人.
不是的.
因為我是你的神.
你就要愛人.
因為你相信耶穌.
你就要愛人.
很簡單而已.
因為你的神.
就是你的神.
你的神是這樣.
所以你的需要.
是這麼做的.
田葛不可割盡.
就是不可以割盡.
是什麼呢?.
就是在問.
我們有沒有預留一些空間給弱勢的社群.
田葛是什麼呢?.
是神要我們預留空間.
為什麼我們要預留空間呢?.
因為神是這樣.
那時候在一個農業的社會.
耕田的社會.
田葛可不可以修國呢?.
當然可以了.
如果我是地主.
所有的東西都是我的.
是不是?.
教會那塊地.
就是沒有地了.
教會所有的地方都是我的.
這個國.
我也要用盡它.
是錢嘛.
這一格多少錢呢?.
就是一萬塊的東西.
五千塊一尺.
五千塊.
如果種東西更過癮.

$^{481}$怎麼樣?.
你丟一些種子下去就能蓋起來了.
每年收國.
收租也不行.
你可以租一個格子給一尺的人.
那些迷你倉是值錢的.
全部都是錢.
可以賺錢.
因為是他的財產.
但是神說.
那些不能搞.
不能動.
要留著.
在農業社會.
當然不可以留著.
因為那些東西不是多出來的.
就是那個國不是多出來的.
是屬於田主的.
是他的部分.
在這裡教導我們有兩件事情.
就是愛倫社.
是需要刻意預留的.
不是剩下給別人的.
如果我們的生活排得滿滿.
一起來你就趕上班.
一起來你就找位置喝茶.
OK.
排得滿滿要坐到十二點的.
喝到十一點的.
沒有空間.
沒有餘地.
沒有時間.
那麼我們就沒有能力.
去走近任何人.
預留.
什麼叫預留呢?.
在我年輕的時候.
我信耶穌.
有一個導師.
他很有趣.

$^{521}$是工程師.
工程師當然是有一點錢.
他教我們.
他沒有教我們.
他說他怎樣去實踐.
信耶穌.
他說經常教什麼十分之一.
什麼十分之一奉獻.
他說好.
我說十分之一.
做完我就不做了.
他說.
十分之一.
那包什麼呢?.
當然是你的金錢.
第二.
你的時間.
你的生命.
他就算了.
好啊.
十分之一.
他二十歲才出來工作.
他假設自己.
做到五十歲退休了.
那就是有三十年了.
要上班三十年.
那就是.
要奉獻三年出來.
三年.
整整三年.
做什麼呢?.
他是工程師.
經常去地盤.
我們知道他每逢六至七年.
就會轉工一次.
剛剛轉工的時候.
他就找一年的時間.
不知道去一些很落後的地方.
真的去搬磚.
真的跟他抿石屎蓋房子.

$^{561}$做足一年.
做完回來.
那他就是工程師了.
是不是?.
一定有很多工資找他.
他很厲害的.
他有買車.
買很多東西的.
做一些最簡單的服事.
用了一年的時間.
那又做到不知道六七年之後.
他又轉工了.
就是又不做了.
那又做一年了.
他總共用了三年時間.
然後他跟上帝說.
那.
我這二十七年我不管你的了.
OK.
這個是他.
那崇拜他會回來的.
那是他對上帝的十分一的奉獻.
不單止是他的人工.
而是他的生命.
哇.
很大.
那麼.
只是大聖經教我們預留田國而已.
那我們先學習一下預留田國.
以及預留田國.
是很刻意的留下來的.
譬如你買房子.
如果你要置業了.
年輕人置業了.
或者成年人就置業了.
或者你要換房子了.
都盼望你們在家裡面留一個角.
給一些有需要的人.
不知道什麼時候用的.
總之.

$^{601}$有百分之五的地方.
是留給有需要的人用的.
不知道什麼時候會用.
預留田國.
是信仰的責任.
不是一個慈善的活動.
神不是問你想不想幫人.
神是問.
既然你也是屬我.
你應該這麼做.
不是慈善活動.
也不是慈善興趣.
幫人不是一個慈善興趣.
知道什麼叫慈善興趣嗎?.
知道什麼叫拉頭馬嗎?.
我們沒有買馬仔.
看電視賽馬的時候.
那隻馬仔贏了.
然後怎麼樣呢?.
會拍照的.
就會拉出一個 banner(橫額).
然後我們就拍照.
贏了.
有很多人做慈善也是這樣的.
做完之後拉出一個 banner(橫額).
做了是不是?.
這些慈善興趣.
我也做很多.
我以前也做很多的.
是興趣.
很開心的.
好像幫了人一樣.
心裡面也歡暢.
是興趣.
但是預留田角.
是一個信仰的責任.
因為神是這樣.
我們就要預留田角給人.
神將我們的愛鄰舍行動.
放在一個生活秩序裡面.

$^{641}$甚至放在一個修葛的規則裡面.
要提醒我們.
愛鄰舍不是偶然.
而是日常的生活.
有人說.
我也不可以做什麼.
我不懂得做.
但是神說.
你能夠做的.
就是你田角那一點點的地方.
田角不是一個很巨大的行動.
而只不過是一個預留而已.
預留你的眼睛.
看一些有需要的東西.
我們只是滾手機而已.
是不是?.
有時候很模糊.
虛空來滾手機的嘛.
你留一點時間.
你也可以滾多一點也可以.
留一些看周圍有需要的東西.
預留心力承載別人的情緒.
就是不要那麼興奮.
OK.
不要去到那麼蠢.
留一些東西給別人.
預留時間陪伴別人的脆弱.
你不要排得那麼密.
留一點點而已.
預留金錢.
支持別人的困境.
留一點點.
放在那裡.
預備去幫人.
他要求我們成為願意看見.
以及預留的信徒.
還有在十九章裡面.
你還沒記十九章的裡面.
是說愛倫社之前.
有一句話經常出現的.

$^{681}$不知道為什麼.
他說.
你們要聖潔.
因為我耶和華你們的神是聖潔的.
在講這一段的.
剛才在講田國之前.
在講他們是聖潔.
為什麼會這麼說呢?.
我們經常以為聖潔跟宗教的儀式.
敬拜的感覺.
屬靈的氣氛.
是有關的.
但是在你還沒記十九章告訴你.
聖潔是你每天怎樣對待你身邊的人.
你想聖潔嗎?.
就要看你平時怎樣對待別人.
很實踐,很實際.
聖潔不是主日唱詩的感動這麼簡單.
而是星期一至星期日.
你有沒有願意放慢腳步.
陪伴一個受傷的人.
聖潔不是在崇拜的裡面.
舉高雙手.
主啊!主啊!主啊!.
而是在生活的裡面.
為人預留多少.
為人預留多少田國.
以色列人以為.
去敬拜.
就是走入回眸敬拜.
但是神說.
你怎樣去對待窮人.
寡婦.
寄居者.
這就是敬拜.
我們打開聖經以為.
我讀很多書.
認識很多經文.
但是神說.
你願意走近那些人嗎?.

$^{721}$你未記十九章十八節.
我們也很熟悉的.
那裡一樣是那四個字.
叫做愛人如己.
愛人如己.
但我們常常誤會.
如己.
不是情緒.
其實是一個標準.
意思是.
你願意別人怎樣對待你.
你就怎樣去對待別人呢?.
如果你家裡有事.
你也很希望有人安慰你.
如果你受傷.
你也很希望有人聽聽你的話.
如果你孤單.
你也很希望有人陪伴你.
神說.
這樣子.
你就去實踐你的愛.
上主給你的愛.
愛人如己.
是一條命令.
不是提議.
不是說提議你去愛一下別人.
不是的.
是一條命令.
不是有空才去做.
不是等你情緒好才去做.
是神要他的百姓時刻記得.
一個生命的字彈.
我們的田國可以在哪裡呢?.
每個人也可以有的.
你可不可以問牧師.
問牧師.
誰是牧師.
有的.
你問牧師.
牧師.

$^{761}$我時間很少.
哪裡有田國呢?.
牧師.
我自己也有很多壓力.
怎樣預留給別人呢?.
牧師.
走過去那些人那裡.
很累的.
我自己也怕掉下來.
弟兄姊妹.
神沒有要求你預留全部.
只是要你的田國而已.
田國可能就是你生活裡面.
百分之五的空間.
可能是十分鐘.
可能是一點點錢.
可能是一句安慰說話.
可能是一個擁抱.
可能是一個代禱.
田國不是奢侈的付出.
而是你的信仰.
生活.
你可能覺得很累.
但是仍然願意聽朋友說五分鐘.
那個就是田國.
你可能經濟緊絀.
但是仍然奉獻一點點.
這個就是田國.
你可能很忙.
但是仍然願意為人禱告.
這個是你的田國.
神不是要你成為英雄.
而是要求我們成為忠心的鄰舍.
所以聖經提醒我們.
不要將生命.
我們的生活.
放得太滿.
沒有留空間去愛其他人.
你的田國.
是神給你保留的祝福.

$^{801}$你為別人預留.
神也會為你預留一些東西.
你有沒有發現.
當以色列人願意預留田國的時候.
他們不只是在幫助那些窮人.
有需要的人.
而是他們在與神同行.
與神同行.
因為神曾經跟他們說.
你們在埃及是寄居的.
我救你們.
所以你們要記得那些寄居的.
就是說.
你怎樣對待那些弱勢.
就是你怎樣回應神曾經給你的恩典.
你們信耶穌為什麼呢?.
開始的時候.
那些呼召是怎樣的.
我自己也是.
是有些事情不好.
是真的處理不來.
主就在那裡拯救我.
所以我們也要好像主一樣.
來幫助其他人.
你要記得主怎樣去救你.
曾經怎樣救你.
怎樣幫助你.
今天我們可以平安.
在這裡敬拜上主.
敬拜上帝.
我們可以在家裡睡覺.
跟家人吃飯.
這些都不是理所當然.
因為我們也需要靠神的恩典.
才可以得著這些生活.
因此我們也需要將恩典活出來.
預留田園.
走近那些有需要的人.
求主幫助我們.
我們同心同意.

$^{841}$你賜下聖經給我們.
你賜下很好的教導給我們.
讓我們可以遵行上主的心意.
上主的旨意.
願主祝福.
祈禱奉耶穌名求.
阿們.
\newpage



\section{民數記 20:14-21}
\label{sec:wcZQh6png8E}
\textbf{2026.01.25《生命的關口》:民數記20\_14-21 (和修版) 陳韋安牧師}
\newline
\newline
連結: \href{https://youtube.com/watch?v=wcZQh6png8E}{\texttt{https://youtube.com/watch?v=wcZQh6png8E}} ~~~~ 語音日期: 2026-01-26
\newline
\newline
\hyperref[sec:sojB_TREQFc]{\small{< < < PREV SERMON < < <}}
~
\hyperref[sec:index_chronic]{\small{[返順時目]}}
~
\hyperref[sec:index_scriptual]{\small{[返順卷目]}}
~
\hyperref[sec:w_aZ5m39lOc]{\small{> > > NEXT SERMON > > >}}
\newline
\newline
民數記 20:14-21
\newline
\begin{longtable}{cl}
\hline
\hline
章節 & 經文 (和合本修訂版)\\
\hline
20:14 & \begin{tabularx}{0.7\textwidth}{X} 摩西從加低斯差遣使者到以東王那裡，說：「你的弟兄以色列這樣說：『你知道我們所遭遇的一切困難。 \end{tabularx} \\ \\ \relax
20:15 & \begin{tabularx}{0.7\textwidth}{X} 我們的祖先曾下到埃及，我們也在埃及住了很多年。然而，埃及人卻惡待我們和我們的祖先。 \end{tabularx} \\ \\ \relax
20:16 & \begin{tabularx}{0.7\textwidth}{X} 我們哀求耶和華，他垂聽了我們的聲音，差遣使者把我們從埃及領出來。看哪，我們到了你邊界的加低斯城。 \end{tabularx} \\ \\ \relax
20:17 & \begin{tabularx}{0.7\textwidth}{X} 求你讓我們穿越你的地。我們不走田間和葡萄園，也不喝井裡的水。我們只走王的大道，不偏左右，直到過了你的邊界。』」 \end{tabularx} \\ \\ \relax
20:18 & \begin{tabularx}{0.7\textwidth}{X} 但是，以東對他說：「你不可從我這裡穿越！否則，我要帶刀出去攻擊你。」 \end{tabularx} \\ \\ \relax
20:19 & \begin{tabularx}{0.7\textwidth}{X} 以色列人對他說：「我們只上大道。如果我和我的牲畜喝了你的水，我必付錢給你。我不求別的事，只求讓我步行過去。」 \end{tabularx} \\ \\ \relax
20:20 & \begin{tabularx}{0.7\textwidth}{X} 以東說：「你不可經過！」他就率領一大群軍隊，以強硬的手出來攻擊以色列。 \end{tabularx} \\ \\ \relax
20:21 & \begin{tabularx}{0.7\textwidth}{X} 這樣，以東不肯讓以色列穿越他的境內，以色列就轉去，離開他了。 \end{tabularx} \\ \\
[1ex]
\hline
\hline
\end{longtable}
$^{1}$關於過關.
過關通常想起兩款.
一款就是打遊戲機的過關.
另一款就是回鄉症的過關.
就是今天說的回鄉症的過關.
我想大家如果住洪水橋.
大都過關應該都會過吧.
應該是深圳灣過吧.
我媽媽也是住天水圍的.
她也是在深圳灣過去.
發覺聖經裡面有沒有一些講過關的經文呢?.
發覺原來有一段是說以色列人過關.
就是到了一個邊境要過關的經文.
和大家一起來看這段經文.
剛才主席庫奧德他讀出經文.
大家都大概知道整個經文裡面的內容.
我們開始說,我們一起祈禱.
天父多謝你讓我們今天來敬拜你.
我們在崇拜裡面.
我們來獻上我們的身心靈.
我們願意獻上我們的專注.
我們的耳朵.
我們來認識,明白你的說話.
求主你幫助孩子有清晰的說話.
能夠說出你的恩言.
奉主名求.
阿們.
今天的過關經文.
是發生在以色列人在抗疫十年裡面.
最後一年的時間.
大家都知道以色列人出埃及之後.
漂流在曠野有四十年那麼久.
經文正正是說他們整個曠野漂流四十年裡面.
最後一年發生的事情.
有幾件事大家要知道.
其實這篇經文沒有詳細記載整個四十年的事情.
整個摩西五經裡面大概都是記載了頭幾年.
開始漂流的那幾年和最後那幾年的事情.
中間那三十多年基本上是跳過了.
是沒有說過.

$^{41}$技術上就是在昭合記的下半部.
整個書卷的下半部分都是那個曠野開始的時間.
和民數記頭十四章的經文.
就是開頭那幾年.
去到二十章之後就發覺是在說最後那一兩年的時間.
所以你會發現.
所謂四十年其實就是一個這樣的記載.
中間是沒有怎麼提起中間那三十多年那麼長的時間.
這個大家要知道的就是.
四十年那麼長的時間其實他們是怎麼走的呢?.
走了那麼久為什麼走不完呢?.
其實就是由西邊開始走起.
埃及大概有個想像.
就是由埃及慢慢走到東邊.
如果打開民數記三十三章的話.
有四十二個地方名.
一個很詳細的四十二章.
就是他們在這十年去紮營的地點.
你想一下大概每年都會轉一個地方.
一個這樣的概念.
所以這十年在這塊曠野裡面.
每年遊牧民族不斷去漂流.
找地方有水源有草就紮營.
然後再遊牧.
所以有點像那些現在都流行的.
車中泊露營那些.
其實是這樣的.
大家就是一個郊外地方裡面來生活.
不過就不是說三五知己幾個家庭.
而是整個民族.
想一下有二三百萬人.
這樣的一個露營會是一個.
所以其實也是一個挺大陣仗的規模.
在曠野裡面一排連營.
一公里紮營.
一個非常非常厲害的景象.
一個移動的城市可以說是.
所以這個就是我們今天要明白.
整個曠野漂流的小小的背景.
二百萬人去露營.

$^{81}$露營了三十多年.
去到一個叫迦底斯的地方.
他們這樣說是回到迦底斯的地方.
為什麼說是回到迦底斯呢?.
因為迦底斯是什麼地方來的呢?.
就是三十多年前十二探子去窺探迦南地的地方.
大家記得是嗎?.
就是約書亞·迦勒有信心去窺探迦南地.
後來伊斯蘭人就沒有信心.
然後就被上帝懲罰.
就開始這個曠野漂流的故事.
迦底斯是一個曠野漂流的起點.
他們經過了四十年就回到這個起點.
你見到經文是這樣說的.
這個十六節.
摩西那一段的說話.
這個就是他們回到迦底斯裡面的請求.
我們一起讀一讀.
預備,一二三.
十一節.
我們不守天監和奉告員.
就不喝這類的水.
我們只走王的大道.
不變左右.
直到過了你的邊界.
謝謝.
摩西就要求以東王就容許他們一班的色人入境.
就是從迦底斯去到這個以東的國境.
再穿越到迦底斯的國境.
然後進入迦底斯的國境.
從迦底斯去到這個以東的國境,再穿越到以東的領土.
直接就去到迦南地.
這個就是最終的目標.
這個所以是一段很正宗的過關經文.
過關的是以色列文.
他們是拿著以色列護照,是一個真正的以色列人.
他們要穿過關口進入到以東的境界裡面.
不過可惜這次的過關是失敗的.
他們被拒絕入境.
以東王說你不可從我的地經過.

$^{121}$免得我大道出去攻擊你.
以東王拒絕他們入境.
甚至沒有交代任何的理由.
不過這個都正常.
通常國家要求你不讓你入境都沒什麼理由解.
不讓就不讓.
直到十九節,以色列人就再一次.
這次就不是摩西,這次是一班以色列人.
就再一次要求以東王批准入境.
我們要走大道上去.
我們和伸出約翰尼的水必給你加值.
不求別的,只求你容我們步行過去.
基本上和摩西的說法是差不多的.
不過加一條要求,就是要喝水.
這個就是之前他們因為喝水吵架的那些相關的事情.
不過我就不說了.
所以以色列人就過關入境失敗.
他們不能夠穿過以東去到迦南地那裡.
他們就要繞道而行.
花多了半年的時間.
離開迦底斯,走上山.
去一個很陡峭的荷爾山.
再去到摩納平原.
才去到最後的迦南.
對經文論有一個最終極的問題.
就是為什麼這班人一定要經過以東的地方.
去到別人國家,去到入境呢?.
最直接的答案是什麼呢?.
因為就是快.
他們可以直接穿過,直線穿過以東.
是節省很多時間的.
如果他們直接經過以東去迦南的話.
基本上是十天的路程.
就可以直接去到迦南.
所以差很遠的.
十天就快速地過.
而且花了半年的時間.
又走山路,又走很窄的路去到迦南地.
所以這個就是大概以色列人.
就是最後一年發生的事情.

$^{161}$一月他們就被驅逐入境.
同年米利欽就離開世界.
五月他們去到荷爾山.
亞倫就死了.
然後繞過以東出現了銅蛇事件.
巴蘭獅子事件.
去到摩納平原.
十月就開始分河,河東的地方.
十一月在摩納平原,摩西就講申命記.
十二月底,摩西都離開世界.
這個就是他們整個礦業裡面.
最後那一年,每個月份的事情.
所以說如果他們直接去到迦南地.
可能就不會發生這些事情.
沒有所謂銅蛇事件.
也沒有那麼多以東事件.
整本文書記就薄很多.
沒有那麼多星篤.
所以最後這個經驗告訴我們什麼呢.
就是原來世上有些關.
不是說難還是辛苦.
也不是說很辛苦很慘.
而是純粹你過不了.
有些關純粹不讓你過就不讓你過.
或者卡了你很久.
所謂卡關就是這個意思.
有些關是你過不了的.
不是你的能力問題.
也不是你能不能做到問題.
純粹別人不讓你過.
不知道大家有沒有試過申請簽證.
說的不是入台證那麼簡單.
而是那些比較長期的入境簽證.
以前去德國的書就是這樣.
是一個很大的惡夢.
對於去留學德國的人.
根本上德國領事館有一個職員.
是很德國人的.
那些方法就是那些壞事.
他是很嚴格但是很官僚.

$^{201}$根本上這是每一個香港人去德國的一個共同話題.
必須要過的一個關.
怎樣卡你呢?.
首先黑面就不說了.
黑面就只是痛苦.
這些不是卡關的痛苦.
假設你有十個文件需要你去申請簽證.
這十個文件是不會寫出來的.
是一些隱藏文件.
基本上文件之外的十個文件.
你是不知道要給的.
每一次去到就黑面.
說你欠這個.
回去就要再交.
要再排期再預約.
回去就還欠這個.
這個不夠還欠這個.
十件你走十次.
十個文件你要交.
所以去到去留學.
花了大半年時間.
交十次文件交十次不同文件.
再補交.
所以就十次黑面十次補交.
十次這樣的卡關.
所以就是這樣.
人家有時候是國家不讓你去那個地方.
你沒有任何的發言權.
只能夠滿足.
或者不讓你去.
當然香港人都是急的.
香港人都是比較心急的人.
我們很怕被人排隊卡關等很久.
不知道大家是不是這樣.
你有沒有發覺通常你過關.
人特別會心跳加速.
人特別會快很多.
你看到哪條隊.
在哪裡哪裡.
你就馬上豎過去.

$^{241}$前面那位搞了很久.
然後就掃不到他的回鄉證.
下飛機也是這樣.
香港人你下飛機.
還沒有下飛機都已經可以換SIM卡.
然後就會一豎就放下那個.
那個儲物櫃.
馬上就拿行李.
十點半下飛機.
十點九點去到的士位.
就說自己已經過關.
就回家.
很快的.
通常我們那條隊是.
很快的那些.
走電梯.
我們會急速的走電梯.
所以是快上加快就可以去到.
所以就是我們香港人的習慣.
這件事情就是.
由羅湖到落馬洲還是都一樣.
全部都是這樣的景象.
你問就是.
為什麼這些人要這麼快.
為什麼他們在那時間.
原來他們.
就算他們十天之內.
可以去到迦南地.
很快的去到迦南地.
他們是不是真的可以去到迦南地.
你看回十四章經文裡面.
當時他們被上帝去懲罰.
經文多讀出來.
十四章第三十三節.
你們的兒女必在曠野漂流四十年.
擔當你們仁恨的罪.
直到你們的屍首在曠野消滅.
我們們窺探那地的四十天.
一年頂一日.
你們要擔當罪孽四十年.

$^{281}$直到我與你們疏遠了.
我爺爺話說過.
我總要這樣做得以我的惡護眾.
他們必在這曠野消滅.
在這裡死亡.
經文說什麼.
你們窺探那地四十天.
一年頂一日.
你們要擔當罪孽四十年.
所以這四十年是一個真實的數字.
不是一個虛數.
他們是需要.
Literally住了四十年曠野.
才能夠離開這個曠野.
經文告訴我們.
這些人在曠野.
可以離開曠野有兩個客觀條件.
要去滿足的.
第一我們要停在裡面四十年.
這麼久.
第二是什麼.
就是上一代的人要壽終正寢.
他們要全部死去.
才能夠進入迦南地.
所以我們問.
其實回來做什麼.
過了以東境也好.
過不了也好.
他們這個快.
其實不會真正快到.
他們想成功快點過關.
快點離開曠野.
其實不會真正離開曠野.
因為真正離開曠野的條件.
不是一個地理問題.
而是一個上帝的時間問題.
也是一個壽命問題.
他們只能夠真的坐足四十年.
真是能夠所有的人都死了.
才能夠離開曠野.

$^{321}$因此兩姐妹.
這些人真真正正要面對的那一關.
最後那一關.
其實從來都不是以東王的關口.
而是上帝自己的關口.
我們人生裡面其實也是這樣.
我們面對著很多不同人生的關口.
有時過得到.
有時過不了.
我們想辦法希望可以順利過到.
無論是我們自己小時候的考試.
到面試.
到我們人生不同的關口.
我們都有很多不同人生的關口.
但我們都要面對一個最終極的關口.
就是上帝的關口.
幾個月前.
大埔的居民就面對著一個很大的關口.
有些人過得到.
有些人過不了.
幾年前的香港都有一個很大的關口.
有些人就選擇移民.
離開這個關口.
今天想告訴你.
就算我們過不了某一關.
不讓過某一關也好.
我們真正要面對的那一關.
其實不是我們所面對的目前的關口.
真正的那一關.
其實是上帝更加宏觀我們心靈裡面的那一關口.
我們怎樣面對上帝這一關口.
大家有沒有想過.
整個曠野故事裡面.
摩西,亞倫這班人.
他們怎樣面對曠野.
大家有沒有想過.
這班人是怎樣.
整世人是注定過不了關口.
他們知道自己是新人.
他們肯定會過不了曠野.

$^{361}$所以你見到最後一年不是利欽一月死.
五月份亞倫死.
十二月就是摩西死.
然後才真正離開這個曠野地方.
所以每個人都說走出曠野.
邁向迦南去到神所應去之地.
其實不是每個人都是這樣.
有些人一生人都去不了迦南地.
一世人都去不了應去之地.
所以這班人作為以色列人的屬靈領袖.
他們怎樣理解自己的生命.
怎樣理解自己和上帝的關係.
一世人都過不了關.
一世人都去不了應去之地.
難道他們一世人就沒有盼望嗎.
我們問這班一世人都注定離開不了曠野.
他們怎樣理解他們自己和神的關係.
怎樣去追求他們人生的目標.
這是我們今天要學習的一個功課.
不論不說我們很多時候都是一個.
在教會裡面叫做Baby Boomer的一代.
Baby Boomer就是二戰之後出世的一代.
他們都是一個百廢待興.
重新建立的一代.
他們的成長是怎樣.
只要你肯努力就能成功的年代.
只要你肯努力就必定有建設.
必定見到果後的年代.
所以他們總是有一種思維.
就是要不斷去發展.
不斷去向上.
我自己不是Baby Boomer的一代.
但我的前輩是.
基本上很多人都是這樣想.
從小就被教導要上進.
要努力.
要不斷去奮鬥.
教會也是一樣.
教會不斷去擴張.
不斷去人生過關闖關.

$^{401}$是對的.
但又不是完全唯一的人生的道理.
現在職場裡面見到很多二十多歲的Gen Z.
他們不是這樣想的.
他們基本上不一定要和你衝.
不是要你不斷去好盡去做一些事.
他們會講保留自己的生活質素.
怎樣好好的去想其他的事.
所以人生的意義不是純粹奮不顧身.
去出一個龐大客觀的所謂成就.
這就叫成功.
所以成功的定義.
不一定要那些所謂的成功.
而是那些人生裡面很大的成就才叫成功.
這點其實和摩西亞倫的想法是相似的.
對他們來說過關不是他們人生的終極目標.
過不了曠野這一關其實不重要.
進不了迦南也不是最重要的問題.
對他們來說上帝才是他們最終極的那一關.
就算過不了迦南.
他們仍然要去過的是上帝這一關.
對他們來說信仰的目標.
不是要去衝到一個很偉大的成就.
而相反是每一天他們所面對的生活.
就是中間那三十多年聖經沒有記載.
那些每一天很悶.
好像沒有什麼特別.
沒有什麼出色的生活.
他們就可以好好地去過好每一天.
要遵行耶和華上帝的律例典章.
每一天發現上帝給他們的瑪雅.
發現瑪雅.
享受瑪雅.
收取瑪雅.
正正就是他們每一天信仰的目標.
甚至乎就是他們信仰的高峰.
好像沒有什麼特別.
但是這些事情正正就是他們要去做.
一直要做.
做足四十年就離開世界的目標.

$^{441}$無論人生遇到什麼關口.
能否過得了以東這個關口.
能否去到迦南.
今天是菲律賓人也好.
這些不是最重要的事情.
最重要的反而是.
我們在每一天能否過得了這個關口.
我自己也是.
自己人生去到一半.
發覺要能夠成就的大概都是這些.
大概都不會去到一個完全沒想過的.
想不到的目標.
你開始想想剩下的二十幾年.
十幾年 三十幾年.
你再想想剩下的.
大概都算是自限.
有些事情其實都不是小時候去幻想的成就.
但不代表你人生就沒有成就.
因為那個最重要的是上帝的成就.
正如摩西.
我們說詩篇九十篇是一段摩西的詩篇.
大家都熟識經文.
我們生的年日是七十歲.
若是強壯可到八十歲.
求你指教我們怎樣算自己的子.
可叫我們得著智慧的心.
一生的年日有多少.
有數計.
七十歲大概是二萬六千幾日.
八十歲就二萬九千幾日.
所以一生的年日大概是二萬六千幾日.
若是強壯就二萬九千幾日.
其實不是很多.
二萬幾日.
今日你回完崇拜.
吃個飯.
回家做些事.
再煮個飯或看一場戲.
晚上回家早點睡.
就有一天了.

$^{481}$這可能就是很多人退休生活的情境.
我媽媽也是.
現在基本上每天都是自己.
我爸爸離開後就更加荒廢.
總之每天都在外面做早餐.
有時回家打麻將.
回家吃點東西.
就有一天了.
偶爾有一天會去深圳.
人生裡面有很多不同的日子.
有些日子是特別重要.
有些日子是更加多份量.
我想大家也是.
我自己也記得二十多年前結婚那天.
很記得自己在德國讀完論文考試.
由西德坐火車回家.
那一天的情境.
很記得幾年前開了自己的教會.
第一次聚會的情境.
在這二萬幾日裡面.
有些日子對我來說是不一樣的.
特別重要.
特別深刻.
不過你的人生卻不是只有那幾天.
你要相反.
你要好好去度過上帝給你的每一天.
不是去追求更加美好的那幾天.
而是去享受上帝給你的每一天.
哪怕那天好像沒什麼特別.
好像很平淡.
但我們更加懂得去數算自己的日子.
這就是我自己對於自己未來人生.
基本上很簡單.
每天做兩件事.
這也是聖經所說.
第一就是每天發現天父上帝的恩典.
發現嗎哪.
每天都同樣地要去.
特意去想想今天有什麼恩典.
第二就是主動去遵從天父上帝的命令.

$^{521}$這兩件事.
你做足,做夠.
就是每天要做的事情.
我們就是這樣去活下去.
其餘的事.
過不過關.
是否做了什麼大事.
如何去理解自己.
這些其實都不是最重要.
剛才才說過.
稍後新年就去旅行.
大家在每一年裡.
有一兩個星期去旅行.
或者偶爾去一些特別好玩的事情.
但我們更加要懂得.
我們人生不只是那些日子才值得精彩.
每一天好好地去發現上帝的恩典.
每一天好好地去做好.
我們天父要幫我們做的事情.
所以很喜歡一首我們以前經常唱的詩歌.
就叫做《每一天》.
就是每一天所度過的每一刻.
我都得著能敵勝過善言.
我依靠天父周詳的供應.
不用再恐慌與掛念.
他保護他的兒女如珍寶.
他熱心必要成全這事.
你日子如何,力量也如何.
視乎我是否有所應許.
我們人生去到下半截的時候.
更加懂得去活好我們的每一天.
知道每一天.
就算發生的事情不一樣.
有些很平淡.
但仍然是值得我們快樂地去過活.
多謝祈禱.
因為你幫我們生命.
幫我們能夠在生命裡數算自己的日子.
我們雖然地上的日子不斷地減少.
這是我們人生的定律.

$^{561}$但求主你讓我們更加有智慧.
去學習去面對我們生命的每一天.
去享受我們生命的每一天.
我們能夠來數算你的恩典.
去服從你的命令.
做好我們每一天應該有的本分.
求主你這樣幫助我們.
我們能夠知道人生裡有很多不同的關口.
我們盡力地去過.
即使過不到的時候.
我們仍然感到滿足.
因為我們最終要向你去過的.
就是你自己的關口.
你能夠在你裡面得著這個智慧的生.
奉主名求.
阿們.
\newpage



\section{箴言 3:5-8}
\label{sec:w_aZ5m39lOc}
\textbf{2026.02.01《謙卑事奉---生命的轉向》:箴言3\_5-8(和修版) 楊明麗姑娘}
\newline
\newline
連結: \href{https://youtube.com/watch?v=w-aZ5m39lOc}{\texttt{https://youtube.com/watch?v=w-aZ5m39lOc}} ~~~~ 語音日期: 2026-02-04
\newline
\newline
\hyperref[sec:wcZQh6png8E]{\small{< < < PREV SERMON < < <}}
~
\hyperref[sec:index_chronic]{\small{[返順時目]}}
~
\hyperref[sec:index_scriptual]{\small{[返順卷目]}}
~
\hyperref[sec:4qbR6KQJgy4]{\small{> > > NEXT SERMON > > >}}
\newline
\newline
箴言 3:5-8
\newline
\begin{longtable}{cl}
\hline
\hline
章節 & 經文 (和合本修訂版)\\
\hline
3:5 & \begin{tabularx}{0.7\textwidth}{X} 你要專心仰賴耶和華，不可倚靠自己的聰明， \end{tabularx} \\ \\ \relax
3:6 & \begin{tabularx}{0.7\textwidth}{X} 在你一切所行的路上都要認定他，他必使你的道路平直。 \end{tabularx} \\ \\ \relax
3:7 & \begin{tabularx}{0.7\textwidth}{X} 不要自以為有智慧；要敬畏耶和華，遠離惡事。 \end{tabularx} \\ \\ \relax
3:8 & \begin{tabularx}{0.7\textwidth}{X} 這便醫治你的肉體，滋潤你的百骨。 \end{tabularx} \\ \\
[1ex]
\hline
\hline
\end{longtable}
$^{1}$謝主有機會在這裡和大家分享.
我代表母堂問候每一位.
我們每一次都說到.
紅欣堂現在事工非常之有進展.
很多元化.
我們都為紅欣堂很感恩.
我們都一起祈禱.
還有我們一起聽道.
親愛的天父我們來到你的庚前.
我們祈求你慰養小羊.
無論聽的說的都蒙你的恩典.
得你的教導.
奉主基督 聖名 阿們.
今天和大家分享一段大家很熟悉的經文.
《箴言三章》五至八節.
為什麼在這一段呢.
最近在讀《箴言》的時候.
我發覺《箴言》有很多很好的教導.
但是都想自己有去專心學習和大家一起分享.
我發覺《箴言》裡面有很多地方講驕傲.
《新約》也有很多地方講驕傲.
我又相信其實神不是教我們驕傲的.
不過祂想教我們謙卑.
所以我就由驕傲我們開始去看驕傲.
從驕傲一直看下去的時候.
看看人的生命要從驕傲轉去學習謙卑的時候.
我們要怎樣去演進呢.
讓我們一直學習的時候.
我們的生命又能夠有一個轉向.
由我們本性上可能有的驕傲.
轉而有一個謙卑馴服的心來事奉神.
有一節的聖經.
我覺得要很留神的去提醒自己.
驕傲在敗壞以先.
耐心高傲的在跌倒之前.
原來驕傲是很嚴重的.
驕傲會帶來失敗的.
我立刻想起誰呢?.
我想起大胃王.
大胃王很成功的.

$^{41}$我們每個人都覺得他很成功的.
很成功的大胃王就做了一件事.
他立國之後.
有一次他就說.
我要先知道我有多少國民.
這個其實都是一個心裡面的驕傲.
看看我的國家有多大.
我的人民有多多.
最終這件事.
他也知道神不想他去做.
但是有那種隱憂.
我還是想知道.
於是引來一個大的禍患.
死了很多國民.
有一個瘟疫來到.
令到他們受到很大的災難.
於是大胃就謙卑,凝聚,悔改.
才止息了這一次的瘟疫.
原來敗壞就是在驕傲之後.
躺在謙卑的時候.
很多事情是可以預防了的.
《針苑》裡面有很多關於驕傲的經文.
這裡只是部分.
還未列得完.
如果再看《舊約聖經》.
和《新約聖經》的教導的時候.
更長.
驕傲很多的.
但是驕傲說出來的時候.
一直提醒我們一件事.
驕傲是沒有智慧的.
驕傲是哪些人有呢?.
是愚妄人才有的.
驕傲是什麼呢?.
是不受責備已經是驕傲了.
驕傲是帶來敗壞.
驕傲是為了一些傲慢人預留了的刑罰.
裡面說了很多驕傲不好的事情.
驕傲這麼多的不好.
但是社會裡面怎麼看呢?.

$^{81}$社會裡面不是叫我們不要驕傲.
社會裡面怎麼看人呢?.
社會裡面看人是看一件事.
你有沒有成就啊?.
說真的.
我們出去找一間茶餐廳吃飯的時候.
多不多客人啊?.
沒有客人的那一間.
都不要想了.
那是那一間茶餐廳的成就.
如果那個廚師出來告訴你.
你叫牛扒.
其實牛扒要怎麼煎呢?.
你還進不進去吃啊?.
你想一下.
真的.
將來有機會分享牛扒怎麼煮.
我去隔壁那間餐廳.
我們會有這樣的情況出現.
做事.
人家看效率的.
而且你跟別人應對.
人家都看你的效率的.
是不是啊?.
跟著就會看看.
這個人走出來有沒有自信啊.
如果頭低低的話.
人家就覺得都不知道是不是.
找不找他工作呢?.
好像社會裡面.
要求並不是說他謙不謙卑啊?.
不要看這些.
反而看他出來的樣子.
他都驕傲得起.
他能夠有優越感.
他還有些優越.
才能夠有優越感.
於是就引起我們.
會朝向這個方向走.
甚至我們生活裡.

$^{121}$我們自己都會.
你有沒有讀聖經啊?.
讀完新舊約全書沒有呢?.
你讀了一次嗎?.
我新約讀了兩次.
我比你好.
我們會比較.
比較帶來.
我比較聰明.
跟著有少少傲慢.
有少少對隔壁那個人.
這樣的態度.
這個都不是神想的.
神不希望我們這樣的.
神也不想我們.
開會都說是這樣.
我說就是了.
我說是就可以了.
全票通過,拜託.
我們應該堅持自己的意見.
不是.
我們應該堅持的是神的意見.
我們大家一起開會.
一起談的時候.
我們就會想多一些.
思想到闊一些.
看看神對我們.
整體有什麼共同的.
帶領.
在聖經裡面.
我就選擇這段經文.
這段經文.
大家可能背過.
大家去看一看.
裡面有多少次提到驕傲呢?.
你要專心仰賴耶和華.
不可以靠自己的聰明.
在你一切所行的路上.
都要認定他.
他必使你的道路.

$^{161}$平直.
不要自以為有智慧.
要敬畏耶和華.
遠離惡事.
這便醫治你的肉身.
肉體.
滋潤你的白骨.
現在一個字關於驕傲謙卑.
都沒有.
好像與主題不太相關.
我又覺得.
他好像提醒我們.
不要驕傲.
不要裝聰明.
以為自己聰明.
你走的路是怎樣呢?.
要認住看著哪條路.
是神帶領你的.
不要靠自己的智慧.
但要做一件事.
要敬畏神.
敬畏神.
祂就保守我們.
這是我選擇這段經文的目的.
祂好像一個字都沒有說驕傲.
就是提醒了我們.
要謙卑敬拜神.
我們謙卑敬拜神.
這件事與社會中所說的.
未必完全一樣.
但我們值得思想.
社會中的一些觀念.
其實深入地影響我們的價值觀.
有時我們會覺得.
需要堅持.
不堅持.
好像沒有主見.
不表達自信.
就好像.
別人會覺得.

$^{201}$我沒有把握.
我出來能夠表達.
我很有自信.
靠自己很努力.
一定做得到.
是不是呢?.
不一定.
我上星期.
沒有去崇拜.
我請了假.
為什麼呢?.
我去了我外甥女婿的喪禮.
是外甥女婿.
不是祂的長輩.
是外甥女婿的喪禮.
即是年輕人.
年輕一點的.
祂又.
很正常地成長.
跟我們每一個人一樣.
自小受教育.
成長.
出來工作.
結婚,生小孩.
顧著家庭.
很乖巧,很認真.
為什麼這麼年輕就走了?.
這些我們真的自己沒有把握.
不會知道這條路.
是否可以平順地走.
祂只是四,五年前.
有胃癌.
治好了.
才復發.
祂還很乖巧.
祂復發之後.
外父八十歲大壽.
他也來賀壽.
其實是很辛苦的.
但是原來.

$^{241}$我們不能夠靠自己.
靠著神無論怎樣.
這條路祂都帶領.
家人祂都是看顧的.
我們只是將一件事要做.
就是將我們前面的道路.
每一天的事情.
交託給神.
我們知與不知的都交託給神.
你要專心仰賴耶和華.
不可以依靠自己的聰明.
在我們每一條.
要做的事情.
要行的道路上.
我們認定了一個方向.
就是以耶穌基督的引領.
天父的引領.
祂帶我們走那條路.
無論我們看上去很光.
還是我們覺得不夠光.
祂都賜恩給我們.
一條平直的道路.
這是神的恩典.
經文當中.
曾賢.
也說了很多關於謙卑.
關於謙卑.
他似乎沒有說什麼叫謙卑.
只是說了謙卑是有智慧的.
我們會覺得.
有智慧就驕傲得起.
原來.
謙卑是有智慧才謙卑得起.
弟兄姊妹.
我們誤解了.
有智慧才謙卑得起.
當我們沒有智慧的時候.
我們就自卑.
自卑也是將眼睛.
看著自己.

$^{281}$並不是看著神.
看著自卑就會跟別人比較.
然後他覺得我很差嗎?.
你是不是比我還差?.
是不是?.
但不是.
神是欣賞.
懂得謙卑的人.
能夠正視自己的人.
這種人.
神將榮耀賜給他.
還有一個承諾.
能夠謙卑的人.
有財富,有尊榮.
有生命.
做賞賜.
神原來要這樣賜予我們.
有一節經文.
三節.
我覺得挺可怕的.
不知道你讀下去的時候.
是否要很小心.
我覺得這幾節經文.
一定要記住它.
作自己的提醒.
就是《箴言》第六章十六到十九節.
它說到.
神有些事情.
祂很憎恨.
祂非常不喜歡.
怎樣去形容呢?.
祂是恨惡這些事,憎惡這些事.
在神的角度.
這麼憎恨的事情是很嚴重的.
是什麼呢?.
就是那些口吐謊言的假證人.
做假見證.
陷害人的人.
我們都怕了他.
我們都不喜歡.

$^{321}$我們都不喜歡這樣的人.
用假見證去陷害人.
在弟兄姊妹當中.
他散佈一些謠言.
或者是一些排斥的說話.
讓大家彼此紛爭.
製造紛爭.
我們都不喜歡.
讓大家分開.
一黨一黨.
有什麼好處呢?.
大家都沒有好處.
說話的人都沒有好處.
還有一些什麼人.
神不喜歡,很憎恨的呢?.
殺荒的人.
危害別人.
殺害無辜的人.
無辜的人被殺害.
我們都會覺得很同情.
那個被殺害的人.
很不喜歡那個謀殺的.
很憎惡他.
但接著.
其實在這些之上.
第一件事說了一件事.
除了是.
說謊.
去做假見證的人.
殺害別人的人.
心裡面圖謀惡計的人.
要做壞事.
跑那麼快去做壞事的人.
有一件事.
更加提醒我們.
就是高傲的眼睛.
高傲.
是神所不喜歡的.
我用眼睛去高傲一下.
神都不喜歡.

$^{361}$神是不喜歡人高傲.
神不喜歡人高傲.
但為什麼?.
就是不喜歡.
為什麼這麼憎恨呢?.
不喜歡就不喜歡.
都不用去到憎恨.
厭惡.
但神就真的憎恨,厭惡,驕傲.
因為.
有些事.
神都經歷了.
神創造了天地.
神創造了天使.
神創造了天使之後.
又有一件事發生.
就是天使當中.
有一個.
連同他的下屬.
做了一件事.
他心裡面想.
因為天使想.
我要升上天上.
高舉我的寶座.
我要坐在.
會中聚集的山上.
在極北的地方.
我要升到.
高雲之上.
我要.
和至高者.
同等.
我們驕傲都未驕傲到.
神呀!我和你一樣.
未去到這樣的地步.
原來天使驕傲起來的時候.
神呀!我和你差不多.
不如.
你給我位子.
這就是天使怎樣.

$^{401}$變成魔鬼的經過.
記載在以賽亞書.
以賽亞書的十四章.
十三到十五節.
以西結書二十八章.
十三到十五節.
大家可以看看.
他覺得我就是天上.
明亮的神星.
我們說主耶穌基督.
是代表神.
不是我們.
天使為何這樣看自己呢?.
因為他自覺聰明.
他真的比我們聰明.
他自覺有能力.
他真的比我們大能力.
很多對翅膀.
我們一對都沒有.
我們要飛要坐飛機.
他不用的.
他又能夠.
有很多不同的面目.
我們要變面.
他有幾個面罩.
他有四個面前後左右.
這個是.
他很高傲的.
他又很有能力.
神差派他來.
保護我們.
原來.
聰明可以令到人.
甚至令到天使.
驕傲.
不過神不喜歡驕傲.
還有一件事發生.
當天使變成魔鬼之後.
他以蛇的形象.
去找夏娃.

$^{441}$跟夏娃說.
神真的跟你說.
不要吃這棵樹的果子嗎?.
說了一段時間之後.
蛇就跟他說.
其實不是的.
神怕你像他那麼聰明.
你吃了果子之後.
你的眼睛就開了.
沒有東西遮住你.
你什麼都看得很明白很清楚.
你就好像神一樣.
能知善惡.
不一定很吸引.
好像神一樣就很吸引.
夏娃再看看.
這棵天天走過的樹.
平時都不覺得它很漂亮.
今天特別漂亮.
那些果實.
今天才熟了.
好像特別香甜.
真的要先吃掉它.
為什麼?很好吃嗎?.
沒試過的.
很甜的.
不過.
吃了就像神一樣有智慧.
這成為一個吸引.
驕傲是一種吸引.
甚至吸引人到一個地步.
我要跟神.
平等.
其實他跟亞當.
已經在伊甸園生活了一段日子.
沒有遇過什麼問題.
沒有遇過什麼.
但他覺得自己沒有智慧.
只不過這一句說話.
你就好像神一樣.

$^{481}$嘩!多吸引.
當有一個人跟你說.
嘩!你要了這東西.
你就好像神一樣漂亮.
這樣也很吸引.
就是好像神一樣.
好像神一樣.
就吸引了他.
無論那個是蘋果還是蛇果.
他都覺得吃了它就對了.
於是.
發生什麼事呢?.
於是懲罰來到.
每個人都會犯罪.
每個人跟神的關係.
都分離了.
每個人都要面對一樣東西.
就是罪的刑罰.
生老病死.
生命的結束.
和永恆的懲罰.
這是對人的一個很嚴重的後果.
但這件事.
也是神很不愉快的.
他也因為這樣.
他要拯救人的時候.
他應許了救恩.
聖經跟我們說.
第十六章十八節.
驕傲是敗壞之先.
人就是想要.
更加高的位置.
於是想了一些東西出來.
想了一些東西出來.
令自己受害.
好像第一個受害的.
亞當夏娃.
他們開始要勞苦工作.
他們開始.
生育的痛苦.

$^{521}$他們開始要面對.
子女的不平和.
於是.
該隱殺了阿伯.
因為罪進入人間.
所以家庭破裂.
原來驕傲.
可以帶來.
很多的後果.
驕傲會不會危害到教會呢?.
如果我們將.
社會裡面.
對我們是很有自信的.
我要能夠.
控制到一切.
才叫帶領的話.
無論你在哪個部門事奉.
在哪個小組團契當中.
你都會覺得自己.
要做到這件事的時候.
就會引起.
引起紛爭.
引起紛爭什麼呢?.
就是比較.
我覺得你這樣說.
是很不尊重.
你應該用另一個方式.
才尊重.
尊重團友,尊重導師.
其實.
比較之下.
引來的是分裂.
亦有些時候.
我們亦會覺得.
《聖經》我都讀過.
這段我不用聽了.
引導我們.
學我用不著翻.
崇拜.
我用不著聽.

$^{561}$算了,我坐下來扮演一下.
然後我們就會出現問題.
出現了我們.
自身成長上的問題.
我們依靠神的心.
有時都會在這方面.
受到影響.
據說,我相信大家.
大部份人都認識有一個.
演藝界的.
名人.
已經離開了.
叫做石堅.
大家都認識他.
據說石堅.
他已經成年之後.
才加入演藝界.
所以他做得辛苦.
但他很努力.
不過他一起拍戲的人.
就說他好人.
他跟人相處沒有驕傲.
不過.
就因為他很努力.
我很努力,什麼事我都做到.
他太太.
很想他信主.
為他祈禱了很久.
據說三,四十年之後.
他覺得.
為什麼我腳又不太行.
為什麼我手又不太行.
眼睛又不太行.
最終都覺得.
不行了,我要信主.
我們都不要弄到自己.
我都不行了,我才信主.
我要早點信主.
是吧?.
神是聰明的,是有智慧的.

$^{601}$祂還有很多能力.
讓神可以使用你.
讓你可以在教會當中事奉.
不要等到我們.
我手又不太行.
腳又不太行.
這個又不太行,那個又不太行.
你跟我說了.
我一會兒就忘記了.
你寫下了,我不懂開電話來看.
那怎麼辦?.
我早點信主.
告訴神.
我真的需要你.
你賜給我救恩.
你接受我.
其實我們經常說.
我們接受救恩.
我們只是將救恩放得很低.
應該是基督說.
你拯救我,接受我.
求主幫助我們.
賜給我們智慧.
我們步向謙卑.
而不是步向驕傲.
雖然如此.
但驕傲時常埋伏在那裡.
引誘我們.
求主幫助我們.
耶穌基督跟我們說.
祂是葡萄樹,我們是枝子.
祂也跟我們說.
離開了祂,我們什麼都做不到.
我們不能夠離開祂.
離開祂,我們方向都失去了.
離開祂.
我們毅力也會失去了.
我們視野也會失去了.
看不清楚.
神想怎樣.

$^{641}$看不清楚神想我們怎樣.
主耶穌基督.
祂經常說要引領我們.
要我們跟著祂,黏著祂.
這樣祂就很驕傲.
是不是?.
腓立比書不是這樣說.
腓立比書二章提醒我們.
一件事.
提醒我們要以耶穌基督的心.
為心.
耶穌基督的心是怎樣的呢?.
耶穌基督的心是很謙卑.
保羅也很.
學校祂的謙卑.
他在腓立比書的第四章也說.
我們現在.
無論處於好的環境.
或者不好的環境.
我都可以應付.
為什麼呢?.
他不是很驕傲地說.
他只是說我靠著那家.
給我力量的.
凡事都能做.
無論什麼情況.
什麼情況.
我們都可以經過.
無論是喜樂的環境.
或者不喜樂的環境.
我們都可以經過.
我記得我媽媽離世之前.
她說過.
她放心走了.
為什麼呢?.
雖然我們還小.
她離世的時候我們還小.
但是她說.
她知道誰引導著我們.
有神看著我們.

$^{681}$神看著我們.
讓祂看著我們.
更加周到.
所以她放心.
原來有神在我們心裡.
我們就得到.
真正的力量.
這種不是體力.
而是心靈裡.
真正的生命的力量.
謙卑.
是怎樣呢?.
謙卑是不要自卑.
謙卑不是自卑.
謙卑是我們要.
看得正確.
看得適當.
在《羅馬書》.
十二章三節提醒我們.
保羅提醒我們.
我們要看自己.
看得合乎中度.
看得適中.
不看自己太高.
也不用看自己太低.
看自己適中的時候.
神就很能夠用我們.
我們不會被差遣的時候.
覺得我們做不到.
我們也不會被差遣的時候.
覺得當然是我.
而是我們能夠存著謙卑的心.
跟著神的帶領去做.
存著謙卑的心的時候.
大家更加能夠.
同心合意一起去做.
因為大家都覺得.
其實神把我們一起.
我們都是一樣.
一樣是神的.

$^{721}$所揀選,事奉祂的人.
我們是同一組人.
我們是一起去工作.
求主去幫助我們.
我們能夠有這種堅定的.
信靠祂的心.
謙卑在教會裡面的影響.
不只是在.
我們事奉上.
其實謙卑在教會裡面的影響.
直接影響到.
我們信不信靠主.
直接影響我們信不信靠主.
當我們.
覺得自己很行的時候.
其實也想一下.
我有什麼做錯.
為什麼說我是一個罪人.
為什麼說.
我要罪得赦免.
我們看自己.
好端端的.
為什麼說自己是罪人.
沒錯,我們看自己的確是.
好端端的.
好像不是一個罪人.
不過我們好端端的一個人.
我們說話又會得罪人.
我們看人.
可能是鄙視下人.
如果我是他做這件事.
我就不會這樣做.
我們雖然沒有說出口.
對方也聽不到.
可能我們的眼神.
他也看不到.
心裡面做了.
我都沒有做出來.
沒有做出來也算嗎.
應該做的好事.

$^{761}$沒有做出來就好像不算.
沒有做出來的壞事.
心裡面做了.
好像很真心地做了.
對.
我們真的要認真一點.
認真一點去想.
原來我們要由心裡面.
不驕傲出來.
由心裡面不驕傲出來的時候.
我們就真正會欠被獨.
怎樣能夠由心裡面.
不驕傲出來呢.
我們能夠自視得正確.
不要看自己太高.
不要看自己太低.
看自己是正確地.
我需要救恩.
看自己是正確地.
我不至於低級到一個地步.
主耶穌基督不拯救我.
不會的.
神就是拯救我們.
神就是要愛我們.
拯救我們.
祂就是這樣.
我們放下自己願意接受救恩.
神.
祂就是很寬恕我們.
很開心地.
願意迎接我們.
進入天堂的境界.
我們不只是.
要回來崇拜.
不用回來崇拜.
不是,你聽錯了.
不只是回來崇拜.
我們不只是回來崇拜.
我們回來是要信主.
我們信主.

$^{801}$我們就應該要.
接受洗禮.
認真地去接受主耶穌基督.
做我們的救主.
讓祂進入我們的生命.
引導我們.
讓祂建立我們.
我們信主,我們要親近主.
我們親近主.
好好禱告,祈求.
「祈禱我知道,我有」.
「請姐妹問問你」.
「不要舉手」.
「問問而已」.
「你自己想想吧」.
「你自己回答自己」.
「你一天祈禱多少次?」.
「有沒有借飯祈禱?」.
「我們小時候」.
「借飯祈禱的時候」.
「會發生兩件事」.
「一件事就是父母說」.
「我賺錢回來養你的」.
「你祈禱,多謝你的神」.
「第二件事就是」.
「隔壁同學夾了你的菜」.
「現在沒有人做這件事」.
「可能你祈禱借飯的時候」.
「有人幫你夾點心」.
「不過你自己就好像覺得」.
「我也是賺錢回來養我自己」.
「接著吃飯就忘記了」.
「祈禱感恩,對不對?」.
「是我們祈禱感恩」.
「神給我們能力」.
「可以工作」.
「神給我們機會工作」.
「我們祈禱感恩」.
「我們能夠吃得下」.
「能夠有這個」.

$^{841}$「恩典有牙」.
「吞得到」.
「還有一個恩典」.
「你想去吃」.
「說不想吃」.
「你擺多少東西在祂面前」.
「都覺得我預備得這麼辛苦」.
「你看都不看」.
「祂真的不想吃」.
「能夠想吃也是一個恩典」.
「我們要感恩」.
「借飯感恩」.
「早上起床」.
「為了一天要走的路」.
「去求問神引導」.
「每一天」.
「晚上」.
「去思想今天的工作」.
「做得不好的」.
「認罪悔改」.
「求主繼續引導」.
「祈禱」.
「很小事而已」.
「不過我們和神談談」.
「去旅行」.
「我們有機會再去旅行」.
「求主保守」.
「旅途平安」.
「到處都有不同的事情」.
「求主幫助」.
「不受感染」.
「求主幫助」.
「每一步路都有祂看顧」.
「飲食又平安」.
「我們需要親近神」.
「很多事我們都想的」.
「但不一定我們自己做到的」.
「誰會想去旅行之後」.
「回來感冒呢?」.
「沒有人」.

$^{881}$「但去完旅行回來感冒的有多少人呢?」.
「求主幫助我們」.
在耶穌基督.
在寫腓立比書的時候.
說耶穌基督是很謙卑的.
叫我們學習他的謙卑.
在腓立比書的第二章.
當中他提到.
「他本有神的形象」.
「不以自己與神同等」.
「為強奪的」.
「反倒虛己」.
「取了奴僕的形象」.
「成為人的樣式」.
「既有人的樣子」.
「就自己謙卑」.
「全心馴服」.
「以致於死」.
「且死在十字架上」.
主耶穌基督本來是神.
不需要來為我們死.
弟兄姊妹.
你有沒有為一隻你打爛的碗.
哭死.
真的想去死嗎?.
不會的.
買一隻碗.
就算那隻是你自己做的.
做那隻.
但主耶穌基督.
天父上帝.
聖靈.
他們就是用泥和水.
造了我們.
造了我們之後.
竟然要為我們死.
在十字架上.
為甚麼?.
我們有甚麼珍貴?.
不是的.

$^{921}$祂就是這麼愛我們.
祂沒有將自己視為.
比我們高多少倍.
是計算不到的.
我們也不會將自己和泥.
覺得是同等.
但我們就是泥.
天父祂是神.
主耶穌基督是神.
聖靈是神.
為甚麼要拯救我們?.
到一刻我們認罪悔改信主.
聖靈還要來到我們的心中.
帶領我們.
我的心很美.
有多美呢?.
我們的心有多久.
不想做壞事.
不偷懶.
有時我們也想找著數.
但聖靈竟然.
屈居當中.
弟兄姊妹.
我們有一個這樣的神.
祂是神.
祂還謙卑.
跟我們連繫.
我們為甚麼.
不學習謙卑?.
更加像祂.
求主幫助我們.
我們今天看的聖經裡有說.
祂說.
你要專心仰賴耶和華.
不要靠自己的聰明.
是的,我們如果看著自己的聰明.
有時我們會走一些.
不太適合的路.
有沒有聽人說過.
我以為.

$^{961}$走山路是可以的.
誰知走上去後.
要下來.
你不是搬到山上住嗎?.
下來的時候.
就走到腳受傷.
是的.
我們就是這樣.
我們不要想全局.
但神就是帶領著我們.
祂願意在我們.
一切的道路上.
引領著我們.
給我們平直的道路.
祂要我們做些甚麼呢?.
要我們.
不要想著自己很聰明.
祂不是說我們沒有智慧.
祂只是說我們不要想著.
自己很有智慧.
神是賜智慧給我們.
賜給我們不同的才幹.
賜給我們不同的恩典.
但祂不想我們持著這些.
接著我們不懂得.
去敬畏神.
不懂得離開.
一些不好的事.
反而自己會成為.
一些自己的陷阱.
被自己墮落在罪惡當中.
我們要學習一些事.
不要堅持自己的意見.
反而多些去禱告.
神一定會處理的.
引導我們的道路.
我們需要神的帶領.
就好像約書亞.
約書亞?.
約書亞記裡面.

$^{1001}$記載約書亞.
在摩西離世之後.
神揀選他帶領.
以色列人進入迦南地.
其實也不是摩西死後.
摩西一直傳遞給他.
也跟別人說.
他就是未來的領袖.
但是.
他要帶領以色列人.
攻打迦南地之前.
就在第五章的時候.
在一個晚上.
他看到一個.
穿著軍裝的人.
在他附近.
他去問他.
到底你是來幫我們.
還是幫迦南人?.
這個穿著軍裝的人.
怎麼回答他?.
他說.
我是耶和華軍隊的元帥.
我不是來幫你.
也不是來幫他.
我是耶和華軍隊的元帥.
約書亞應該告訴他.
耶和華指定了.
是我帶領以色列人.
去打仗的.
我才是耶和華軍隊的元帥.
不是.
約書亞聽了之後.
他就照著這位元帥所說.
脫下他的鞋.
敬拜這位元帥.
為什麼呢?.
他知道誰才是真正的元帥.
真正的元帥是神.
是神猜派使者.

$^{1041}$與他同在.
一起來對爭戰.
在我們每一條道路上.
神都猜派使者與我們同行.
很多事情.
好像是我們做的.
其實是神一路引導我們.
給我們機會去做.
還有一件事.
我們一路看約書亞記的時候.
去到第24章.
大概第十四節.
十四十五節的時候.
就說.
約書亞.
叫齊.
十二個支派的首領.
長老來.
跟他們說.
我年紀大了.
他當時應該近一百歲.
九十多一百.
他說我都年紀大了.
那麼.
你們要想一想.
想什麼呢?.
你們以前在埃及地.
有埃及地的神.
河東有河東的神.
迦南地.
有迦南地的神.
吾爾那邊有吾爾的神.
但是.
你要想清楚自己要信誰.
不要信著信著.
就轉去信另一個.
這樣是不行的.
你要想清楚自己要信誰.
也不是多信一個.
信神.

$^{1081}$不是我們信的.
以往每一個信的.
之上再加多一個.
不是.
是我們單單只信靠神.
接著約書亞就跟他說.
我看你們都做不到.
那麼你們不如不要信.
你不是要多信一個神.
你一是就只是.
信神.
接著他用了一句.
說話去結束.
他說至於我和我家.
我們都必.
事奉耶和華.
至於我和我家.
我們都必定事奉耶和華.
他不是叫以色列人.
不要信神.
而是叫他們想清想楚.
你是否信神.
是就要認認真真.
去相信神.
依靠神.
他說你就要想清想楚.
我已經想清楚了.
雖然我一把年紀.
我決定事奉神.
但他一生都事奉神.
他說也不是這樣.
我的全家都會事奉神.
他鼓勵他的家人都事奉神.
鼓勵不只是口說.
他培育他們.
敬拜神.
培育他們事奉神.
帶領他們事奉神.
這是約書亞.
他最後跟那群.

$^{1121}$以色列人所說的話.
由第五章.
他承認.
耶和華是軍隊的元帥.
去到最後的一章.
他就說.
他和他的家人.
我們都必定事奉耶和華.
這樣來帶領.
以色列人.
弟兄姊妹.
我們都要很認真地去想.
我們是要.
真真正正去事奉神.
不靠自己的聰明.
跟隨他的時候.
原來他還賜給我們恩典.
醫治我們的肉體.
滋潤我們的白骨.
這不是說.
信主之後你不會病.
而是說無論什麼環境.
他和我們同在.
我們的身心靈.
就好像很滋潤那麼舒服.
又得到他的看顧.
靠著主.
我們可以走過.
無論是多麼幽暗的幽谷.
謙卑.
馴服.
交託給主.
放開我們的心懷.
讓主來對我們帶領.
成敗得失.
是我們看的.
神看我們的成績表.
不是這一張.
乃是祂的教導我們有沒有遵行.
而祂又能夠令我們真實的生命.

$^{1161}$健健康康.
喜喜樂樂.
我很想送這一節經文.
給大家.
敬威耶和華.
心存謙卑.
就得財富.
尊榮.
生命為賞詞.
《僧言二十二章》第四節.
你願意得到嗎?.
我們一起禱告吧.
親愛的天父上帝.
我們來到你的面前.
天父下我們.
要向你認罪悔改.
又求你保守.
因為我們真的不想驕傲.
主下我們真的不想驕傲.
不過驕傲的確是.
埋伏在這裡.
不斷地引誘我們.
在我們不覺意的時候.
祂就會浮現出來.
求主幫助我們.
赦免我們.
帶領我們防備驕傲.
賜給我們恩典.
讓我們能夠立志謙卑馴服.
敬畏事奉你.
求你幫助我們.
建立我們.
培育我們有一個謙卑馴服的心.
求你帶領.
看顧我們.
讓我們能夠得著你的恩典.
靠著你立志.
一生好好事奉你.
真是能夠學習.
約書亞所說的.

$^{1201}$我和我家.
必定事奉耶和華.
祈禱仰望您奉主耶穌基督的聖名.
阿們.
\newpage



\section{哥林多前書 9:19-27}
\label{sec:4qbR6KQJgy4}
\textbf{2026.02.15《不再「人有我有」的事奉》:哥林多前書9\_19-27) 曾裔貴牧師}
\newline
\newline
連結: \href{https://youtube.com/watch?v=4qbR6KQJgy4}{\texttt{https://youtube.com/watch?v=4qbR6KQJgy4}} ~~~~ 語音日期: 2026-02-19
\newline
\newline
\hyperref[sec:w_aZ5m39lOc]{\small{< < < PREV SERMON < < <}}
~
\hyperref[sec:index_chronic]{\small{[返順時目]}}
~
\hyperref[sec:index_scriptual]{\small{[返順卷目]}}
~
\hyperref[sec:POs_KWCBzOE]{\small{> > > NEXT SERMON > > >}}
\newline
\newline
哥林多前書 9:19-27
\newline
\begin{longtable}{cl}
\hline
\hline
章節 & 經文 (和合本修訂版)\\
\hline
9:19 & \begin{tabularx}{0.7\textwidth}{X} 我雖然是自由的，不受人管轄，但我甘心作了眾人的僕人，為贏得更多的人。 \end{tabularx} \\ \\ \relax
9:20 & \begin{tabularx}{0.7\textwidth}{X} 對猶太人，我就作猶太人，為要贏得猶太人；對律法以下的人，我雖不在律法以下，還是作律法以下的人，為要贏得律法以下的人。 \end{tabularx} \\ \\ \relax
9:21 & \begin{tabularx}{0.7\textwidth}{X} 對沒有律法的人，我就作沒有律法的人，為要贏得沒有律法的人；其實我在神面前，不是沒有律法，而是在基督的律法之下。 \end{tabularx} \\ \\ \relax
9:22 & \begin{tabularx}{0.7\textwidth}{X} 對軟弱的人，我就作軟弱的人，為要贏得軟弱的人。對甚麼樣的人，我就作甚麼樣的人。無論如何我總要救一些人。 \end{tabularx} \\ \\ \relax
9:23 & \begin{tabularx}{0.7\textwidth}{X} 凡我所做的，都是為福音的緣故，為要與人共享這福音的好處。 \end{tabularx} \\ \\ \relax
9:24 & \begin{tabularx}{0.7\textwidth}{X} 你們不知道在運動場上賽跑的，大家都跑，但得獎賞的只有一人？你們也要這樣跑，好使你們得著獎賞。 \end{tabularx} \\ \\ \relax
9:25 & \begin{tabularx}{0.7\textwidth}{X} 凡參加競賽的，在各方面都要有節制，他們不過是要得會朽壞的冠冕；我們卻是要得不會朽壞的冠冕。 \end{tabularx} \\ \\ \relax
9:26 & \begin{tabularx}{0.7\textwidth}{X} 所以，我奔跑，不像無目標的；我鬥拳，不像打空氣的。 \end{tabularx} \\ \\ \relax
9:27 & \begin{tabularx}{0.7\textwidth}{X} 我克制己身，使它完全順服，免得我傳福音給別人，自己反而被淘汰了。 \end{tabularx} \\ \\
[1ex]
\hline
\hline
\end{longtable}
$^{1}$請我們一齊同心祈禱.
多謝主.
我們能夠在不同的崗位.
在教會這個崇拜公共的空間.
我們能夠事奉我們這位永生的主.
剛才的聖師能夠事奉這位的主.
是最尊貴的.
所以我們需要忠於我們的使命.
今天早上.
使徒保羅這一段這麼發人深省的經文.
就真的要讓我們思想.
他是如何不斷面對不同的群體的時候.
他可以放下他自己的身段.
來圍著要將人帶進福音裡面.
與人同得福音的好處.
所以我們盼望洪恩家的弟兄姊妹.
因為我們在這過去幾個月.
真的有不同的群體.
落橋樓的居民.
青年會小學的學生.
或者有我們洪福村的過去一路的居民.
以及在洪水橋社區.
都有不同的居民.
在我們中間來聚會.
所以我們很盼望.
我們每一位弟兄姊妹.
透過今早我們的淨道.
主的聖靈來感動.
來呼召我們.
叫我們渴慕成為一個服侍主的人.
今早我們等候在你的面前.
這樣感恩禱告.
奉耶穌基督的命.
阿們.
講這一堂道.
相信是一個幾刻意想出來的一個淨道.
就是很想鼓勵弟兄姊妹.
來參與在我們洪恩家.
不同的崗位的事奉.
但是當一想到事奉的時候.

$^{41}$怎樣的事奉才能夠在神的面前.
或者在你自己面對不同的場景.
不同的工作的時候.
有時總會人與人之間會衍生一些張力.
誤會.
在事奉裡面會有忙碌.
會有疲乏等等.
怎樣能夠讓大家都能夠有一個持續性的.
來去服侍呢?.
相信最簡單的定義大家都一定知道.
事奉不僅僅是一個工作.
不僅僅是一個上班.
尤其是大家根本就是義務.
無所謂上班.
但是事奉怎樣能夠體會是一個.
在神的面前.
你能夠不斷.
就算面對一些難處.
你都能夠可以去克服呢?.
我相信這段經文.
就已經是日於言表告訴我們.
保羅就是因為這樣看到.
那個服侍神給他的崗位.
那個寶貴.
而他能夠這樣持續的事奉.
而且.
今天這一段經文.
你想想他能夠對著不同的群體.
他有一個很重要的秘訣.
其中這個秘訣.
就是他能夠隨時的變身.
這個變身.
值得我們大家一起去思想一下.
為什麼要用這個題目呢?.
人有我有.
後天就是初一.
我不知道大家小時候.
媽媽是最忙的.
一去到年三十晚.
她就會說.

$^{81}$年晚煎堆.
當然我們不炸煎堆.
太複雜了.
就買回來.
就會跟我們說.
家家戶戶都要有煎堆.
還有人有我有.
一定要有.
我一直去思考事奉的時候.
突然想起媽媽小時候說過.
年晚煎堆.
總之要有.
不過弟兄姊妹.
這個題目其實是正正告訴我們.
保羅不是這樣.
保羅看到事奉很重要的心態.
這段經文頭四節至二十三節.
我們看到一個很重要的秘訣.
試圖保羅可以不斷的變身.
變身其實是不是代表虛偽呢?.
大家不知道你有沒有試過WhatsApp呢?.
很激氣.
在打字.
其實在說的話是想罵上司的.
當然不敢.
打完之後馬上刪除.
然後就說一些正常的話.
有時候激氣到一個地步.
可能都會這樣.
因為沒有人知道.
曾經在基督教的一個節目.
在YouTube裡面.
那個女傳道.
真的示範給年輕人看.
她自己明明很激氣.
上司唐主任叫她做事.
她很不喜歡.
她打了一句差不多近乎粗口.
去罵上司.
當然不敢.

$^{121}$馬上刪除.
她說很想反映.
有時候年輕人.
或者你很激氣的時候.
就會有這種心情.
是一種發洩.
當然拍這樣的片很爭議.
但是有時候反映到我們.
大家有時候面對一些挫折的時候.
可能都會有這種處境.
心態.
但是我們今天想的.
就是試圖她的變身.
記住.
她的變身.
是隨著不同的群體.
她要面對.
怎樣讓這些群體的人.
來聽得懂保羅的訊息.
所以原來很重要的.
保羅從來沒有改變她的內心.
對那些群體的愛.
待會裂開了好幾個群體.
我們一起來.
大家先思考一下.
剛才說了.
其實真正的切換模式.
我們其實不是虛偽.
如果我們的動機.
是真的這樣來到.
好像變色龍一樣.
見到人就說人話.
見到鬼就說鬼話.
這樣.
這樣你就是一個變色龍.
一個虛偽的社會的人.
但是保羅告訴我們.
其實不是.
所以我們每個人.
大家都有不同的群體.

$^{161}$我們在不同的群體裡面.
那個關係.
是怎樣來適應.
你在父母面前.
你當然不會說一些.
很粗俗的話.
但在自己的朋友圈.
可能有時候你又會很粗俗地.
來彼此的分享.
有時候好像在這個關係上面.
問題就是怎樣拿捏那個程度.
怎樣能夠不是真的這樣.
來做一個虛偽的人的變色呢?.
保羅就給了我們很好的教導.
我們一起進入去.
來看一看.
他說對什麼樣的人.
這個九章的二十二節.
其中他就說.
對軟弱的人.
我就作軟弱的人.
為要贏得軟弱的人.
對什麼樣的人.
我就作什麼樣的人.
無論如何.
我總要救一些人.
在上文.
他還說了不同的群體.
有他自己的同胞.
猶太人.
有一些不是在律法底下的外邦人.
有一些是有其他情況的人.
你可以想像一下.
保羅真真實實面對著.
在哥林多.
在他牧養的時候.
是有不同的群體的人.
那他為什麼.
怎樣來變身呢?.
這裡就告訴我們.

$^{201}$他有一個很大.
很重要的原則.
就是他很想.
分享基督的救恩.
來讓這些不同的群體的人.
認識他要傳遞的福音.
這個是很重要的.
怎樣能夠讓軟弱的人.
來到去.
保羅又成為了軟弱的人呢?.
就是說他不會是一套的標準.
來到去將好像一本的經書.
一本的標準.
扔到一個人的當中.
使得他.
你跟著做吧.
你做不到我的標準.
你就不能夠合格了.
不是這樣的.
如果用一些再簡單的例子.
來說的時候.
可能.
有些朋友.
他開始很想明白福音.
但他仍然吸煙.
你說對不起了.
你走進我的門口.
我聞到煙味.
你不准進來.
那你就很難叫他信主了.
因為他還沒放下.
以為要靠自己.
逼自己戒了煙才回教會.
你這樣這個標準.
就不對了.
因為他現在正正在那個軟弱裡面.
他現在是在那個.
這樣不好的嗜好裡面.
被捆綁住.
你怎樣能夠.

$^{241}$既沒有放下那個真理.
能夠對這個軟弱的人.
有一個接納呢?.
讓他明白.
你的心是想他.
慢慢認識耶穌.
得到耶穌.
聖靈就會幫他.
來徹底的改變他那個壞習慣.
我們仍然認為那個是壞習慣.
是不應該鼓勵的.
但不等於.
你要將他那個壞習慣.
你不改好就不要來教會了.
這樣不行的.
所以.
單單這樣已經讓我們看到.
相信我們每個弟兄姊妹.
在我們的生命裡面.
都有我們自己成長的時候.
需要不斷克服的那些問題.
哥林多教會.
你想一下.
當時.
如果用我們今天的術語.
當時的網紅.
是什麼人呢?.
是在哥林多那些廟宇裡面的那些廟祭.
有哪些網紅呢?.
是有很多在哥林多和當地的人.
可以去講不同的哲學.
這些巡遊的.
這些講學的人.
有以彼古龍.
有斯多瓦這些門派的人.
到處去講學.
這些是當時出色的網紅.
也有不同的人.
譬如猶太教.
也有很多人到處去招搖過市.

$^{281}$逼猶太人要根據摩西律法生活.
他們要去看看他們有沒有守安息日等等.
你想一下.
保羅就要面對著不同的群體.
對他這個單純的耶穌福音的攻擊.
挑戰的.
但是保羅能夠很了解福音的那個唯獨.
那個必須.
那個唯一.
但是對著不同的群體.
他就能夠因應那個群體的情況.
慢慢地改變他們的心靈.
其中有一樣東西很重要.
就是首先要讓群體對保羅所講的.
不會帶來一種抗拒.
就是有一個關係上的.
首先的建立.
這個我們真的要在這裡學習.
我們面對著不同的群體的時候.
怎樣能夠讓對方感受到我們不是一種壓迫.
一種強硬.
而是一種對他們的一個接納呢?.
相信真的愛能夠靠近這樣的人.
靠近到一個地步.
小朋友聽得明白你說的話.
老人家又能夠感覺到你有安全感.
年輕人覺得你不是跟他們.
好像不能夠對話的硬邦邦那樣.
而是對他們有一個這樣的接納.
開放.
不同的群體.
你就能夠可以跟他們有一個對話的空間.
相信這個保羅用了兩三節的經文.
來這樣濃縮的去講.
他是怎樣對著這些群體.
對猶太人.
我就作猶太人.
是多麼的鏗鏘有力.
你們要守節旗.
你們要來做猶太教要做的事.

$^{321}$你們的食物條例.
我比你們更加的嚴格的遵守.
但是不是這些的節旗不是這些的食物.
令到我可以在神面前蒙恩.
是耶穌的救恩.
他自己很清楚.
但是要得到這些猶太教的猶太人.
他要怎樣來靠近他們呢?.
跟他們有共同的話題之餘.
更重要告訴他們.
你做的一切不可以使你得救.
唯有來到耶穌的面前.
這一位原來就是我們舊約.
應許的彌賽亞.
這樣的時候.
猶太人就慢慢放下他們對保羅的誤解.
那個敵視.
慢慢就開放自己.
原來那個十架上面的耶穌.
就是我們舊約等待的彌賽亞.
這樣一個忍辱負重的心情.
要慢慢進到猶太社群當中.
所以保羅在第三次宣教旅程完畢之後.
帶著馬其頓教會的捐項金錢.
上去周濟耶路撒冷的窮人的時候.
保羅在《四途行傳》的二十二章.
他去到聖殿還願.
那西爾人的願.
路加醫生這樣記載.
一個已經是做外邦人使徒的保羅.
上到耶路撒冷.
他有願在身.
讓猶太人看到.
其實他是循規蹈矩的猶太人.
不過他不同的.
和其他猶太人不同的.
就是他有真正的救世的福音.
在他的內心裡面.
所以他在猶太人當中.
他就作猶太人.

$^{361}$為要贏得猶太人.
這個就是保羅.
這個就是事奉官.
原來不是自己不斷的自我膨脹.
個個都要聽我.
因為我是使徒.
不是的.
他是將自己放下了.
他的身段.
所以這種變身.
是一個很大的.
一個生命的成熟的事奉者.
必須要具備.
他說對律法以下的人.
我雖然已經不再律法以下了.
但還是作律法以下的人.
為什麼呢?.
為要贏得律法以下的人.
這個就是保羅.
他說我已經不再律法以下了.
基督耶穌釋放了他.
使得他得自由.
不再被奴僕的欺騙來協制.
格拉提書五章一節這樣說.
基督耶穌將每一個人.
無論你是猶太人.
外邦人.
無論你是一個奴隸的身份.
或者一個自主的人的身份.
或者一個男的.
一個女的.
無論是什麼身份.
與生俱來.
無法改變的.
但是我們就是泰國人.
我們有福建人.
我們有不同的.
我們的族群.
在主救恩裡面.
每一個都因為基督耶穌的救恩.

$^{401}$來得到釋放.
所以保羅說.
他很清楚自己不再律法以下了.
但是他要在律法以下.
來幫助那些在律法以下.
以為只要做好自己循規蹈矩.
就可以得救的律法以下的人.
要帶他們離開律法.
帶他們不需要再怕律法.
帶他們明白.
耶穌已經為我們完成了.
整個的救恩.
所以我們不是要守十誡的誡條.
就可以得救.
是在救恩裡面得救.
保羅說對沒有律法的人.
我就作沒有律法的人.
為要贏得沒有律法的人.
當然不是教大家過年.
現在三缺一.
你走過去蹬腳.
我要得到這些打牌的人.
就不是這個意思.
我們不認同.
但是我們要讓人知道.
我們未必是這樣參與賭博.
但是我們是關愛他們的內心.
很多不同的例子.
我們應該怎樣來激放我們自己的身段.
但是同一時間.
我們讓他們感受到.
我們對他們的愛.
這一份這樣的接納.
各位弟兄姊妹.
其實有時候要舉一反三.
我們有時候要有不同的改變.
兩件事是完全不同.
我們不是要扮演角色.
不是角色扮演.
保羅現在不是在做一個角色扮演.

$^{441}$一個變色龍.
到處去逢迎不同的人.
不是這樣將福音完全放下.
沒有任何的立場.
這樣不是在說一個傳福音的人.
事奉的人的心態.
這是一個能夠像Crossover一樣.
我和你一樣這樣來面對.
我們人生不同的掙扎.
有些人覺得你們基督徒.
說就好像很偉大.
很公義.
很敬虔.
但實際做出來就很不對.
如果你的同事這樣來敵視基督教的時候.
你應該怎樣來回答他呢?.
是不是這樣大家決裂就算了?.
你這麼偏激.
我不能夠跟你說話了.
是不是這樣呢?.
很多時候能不能夠用我們的溫柔.
慢慢的聆聽對方的憤怒.
可能人家真的有問題.
可能真的有過一些遭遇.
是真的有一些不好的基督徒.
曾經對他的一個傷害.
我們是不是因為這樣就等於.
我們要不斷為一些不對的東西.
都要來辯護呢?.
來跟他爭辯呢?.
抑或我們能夠靜靜的來聆聽他的憤怒.
來告訴他.
分享我們怎樣在神裡面.
來得到我們自己的救恩呢?.
能不能夠將那個話題.
不只是一種對立式的來抗爭.
而是一個我們能夠體會對方.
這個時間的一種憤怒.
來慢慢的改變他呢?.
所以不是角色扮演.

$^{481}$不是一個辯識龍.
我們是要來到.
真的可以從對方的角度去看事情.
更重要當然就是我們的生命的見證.
怎樣能夠改變對方.
剛才也說了.
保羅說他甘心作眾人的僕人.
這個就是第十九節.
這個相信是我們面對是逢最困難的地方.
就是我們很多時候一面對是逢的難處.
我們就開始多愁善感.
開始自我膨脹.
開始永遠是對方是錯的.
保羅告訴我們.
我雖然是自由的.
不受任何人的管轄.
但我甘心作了眾人的僕人.
為贏得更多的人.
相信這個價值觀.
就是相信非常明白自己.
這個不是一個工作.
可做可不做的義工.
這是神給他的一個很清楚的照明.
當有了這個照明的時候.
面對的挫折.
他會願意自己自我調整心態.
各位應該是我們的心態要調整.
那個環境自然就會因為我們的心態調整.
而會得到很大的改變.
所以有些長輩有時候教我們.
當你面對挫折的時候.
或者有些不開心的時候.
或者想不通的時候.
不要想它了.
去走一走.
睡一場覺.
睡醒覺第二天可能沒事了.
因為你的心態已經有點不同了.
你不再在那個事件上緊緊地綁著.
你慢慢放鬆自己.

$^{521}$其實更多的時候就是聖靈讓你去反思.
可能根本連你自己都有錯.
根本可能連你自己都面對的一些.
可能在那些張力裡面.
有時那些言語之間.
都有一些角力.
告訴我們.
祂能夠在事奉裡面.
看到自己能夠有這麼尊貴的事奉身份.
因為沒有人克制祂.
祂甘心樂意.
在主的呼召下.
做眾人的僕人.
可以想像得到.
這就是最重要的事奉秘訣.
讓我們持續不斷地更新.
不斷地令自己的事奉更新.
的一個很重要的秘訣.
因為我們一面對事奉.
一點點不滿意.
我們就立即很多很多的怨言.
不是說我們有正常的表達是錯的.
不是說我們要去調整一些事奉上的.
做法上的意見是錯的.
當然不是這樣.
因為不斷地去改良.
不斷地去優化每一件事.
這是我們應有的責任.
但是心態呢?.
這裡告訴我們.
保羅很清楚自己在主的救恩裡面.
其實外邦人你能不能救.
根本與我無關.
我是猶太人的拉比.
我是一個很尊貴的.
一個猶太教裡面的一個法利賽人.
我的身份.
我的成長.
都是在猶太教的氛圍.
你能不能救你們拜偶像.

$^{561}$關我什麼事?.
不過有一個很重要的托付.
在保羅的心靈裡面.
所以他不斷調整自己的心靈.
裝備自己成為外邦人的使徒.
你想想西方的宣教士.
最早期清末民初的時候.
來中國.
大得生要留鞭.
為著貼近我們中國人.
去到我們當中.
你可以想像得到.
這一種的背後.
是不是要逢迎呢?.
變色龍呢?.
其實不是.
是要讓中國人在這麼大拒絕.
這麼大的誤解基督教的時候.
完全看為是一個很抗拒的一個宗教的時候.
讓他們感受到這一份的愛.
這樣的溫柔.
讓他們來明白.
基督差遣我們這些西方的宣教士.
飄洋過海來到你們的中間.
來跟你們傳講福音.
有時候還不明白的時候.
就開始要辦一些醫療.
辦一些教學等等.
來慢慢讓中國人不感覺到那種噁心.
用了很多的用心良苦.
原因就是.
甘心作眾人的僕人.
各位弟兄姊妹.
參與著事奉的弟兄姊妹.
甘心.
就是說你很清楚.
主在你生命裡面的呼召.
你願意來到去甘心樂意地去做.
有怨言.
有不同的意見.

$^{601}$你仍然虛心的聆聽.
就是讓你自己的事奉.
能夠做在神的面前.
這是非常寶貴的教會的資產.
如果有弟兄姊妹參與事奉.
你能夠不斷地看到.
是神在你生命裡面的一個感動.
來去參與.
是義工.
沒有人逼你.
你不回去你不做.
好像沒有人跟你關心.
其實不是.
每一個參與事奉都是主愛的呼召.
每一個你在你的崗位忠於使命.
你就能夠真的看到.
神在我們每一個弟兄姊妹當中的工作.
所以今天這段經文.
我自己很珍惜.
很寶貴.
我們看到這件事情.
他面對的比我們更加嚴重.
面對的是整個社會的敗壞.
他要面對這些人.
還有哥林多教會在第八章的臨前.
有些人走去廟裡.
來吃廟宇裡面的食物.
之後拜完那些偶像.
就跟那些廟祭來發生不好的關係.
保羅面對這樣的事情.
你可以想像到還有很多錯誤的哲學.
都是這樣來推翻福音的信息.
但是保羅要面對的就是.
這些人為了要分享福音.
總要救一些人.
就是剛才我們說的.
為了愛可以自願地降級自己的.
無論身份.
自己的舒適.
自己的自由.

$^{641}$你願不願意這樣放下呢?.
你願不願意為主.
來將自己的一些舒適降低呢?.
當然不是說.
什麼假期都不准放.
長假期都要繼續當值.
不是這個意思.
但是每個人自己一定有自己的心靈調整.
剛才我們說到咖啡,奶茶.
因為嶺南大學又再邀請我們了.
我又擔心弟兄姊妹平日上班很忙.
我也不敢馬上答應這麼多.
但是剛才對話的時候.
弟兄說不是的.
我可以調整一下我的工作.
也不行的.
這樣的時候.
如果出自弟兄姊妹內心對這些大學生的感動.
就完全不同了.
整件事變成一個學務服事.
就很不同了.
原來可以調整.
工作的期間.
原來在工作上都可以有些調整.
那就很不同了.
弟兄姊妹.
完全是出於愛可以自我的降級.
愛可以帶來很大的感動.
愛可以引發很多的服事.
當一無了愛.
就這樣計較.
那樣就為自己去想.
本來不是錯.
不過那個計較.
有時候就變成擴大了自己的心靈.
有時候慢慢不知不覺.
原來自己開始有了一個玻璃心都不知道.
很容易就會破碎.
很少就會破碎.
你看保羅.

$^{681}$面對這樣的人.
他能夠將自己擺在不同的群體當中.
不是變身.
是他的心靈從來都沒有變.
他整個人就是這樣為神的福音.
嘗試用一個例子來思想一下.
剛才提過.
在這兩段經文裡面.
對猶太人.
他就作猶太人.
說要得到猶太人.
對律法以下的人.
他就在律法以下.
來得到這班人.
即是說未必是一些宗教領袖.
但是自小就在猶太教的氛圍成長.
保羅都願意來去輔奏這些人.
對軟弱的人.
有些人真的好像在羅馬書第十四章.
有些人覺得日日都一樣.
有些人就覺得某些日子他們要牽手.
保羅就是這樣來去.
在他們當中.
慢慢引導他們明白.
神的真理.
對他們的釋放.
對沒有律法的人.
保羅就是在神.
在基督的律法底下.
來帶領他們.
好像很抽象.
其實一點都不抽象.
就是他能夠面對不同的人.
他願意放下他自己的身段.
放下自己的身段.
來幫助這樣的人.
惠著的是福音.
所以九章的二十三節.
相信是這一大段聖經.
最重要的保羅內心的價值.

$^{721}$凡我所行的.
都是為福音的緣故.
為要與人同得.
這福音的好處.
他看到福音在人生命裡面的更新改變.
這個的好處.
比他個人的自由.
個人所應該做.
不應該做的.
來得更加重要.
願意放下自己.
貼近那個人的需要.
弟兄姊妹.
這一節聖經.
保羅很清楚自己能夠做什麼.
但是為了福音的緣故.
他就願意放下他自己的一些自由.
這是一個很成熟的表現.
最後都要用一個例子.
跟大家一起分享.
所以為了什麼而變呢?.
不是生命的改變.
只是表面的一個.
好像包裝上面.
實際上內裡是沒有變過.
我們不是需要一個完美的自己.
而是我們需要去靠近有需要的人.
讓他們能夠看到.
我們內心裡面的福音.
大家認不認識這位張平儀?.
在公元已經是二千年.
他是感動中國的一個人物.
這位台灣相當出色的記者.
他有一段時間.
他專做一些專題的訪問.
有一次有一年.
他去了四川梁山的大營盤村.
來做一個明查暗訪.
因為他有一段時間.
開始針對麻風的族群.

$^{761}$他在台灣或者不同的地方.
去訪問這些麻風村.
令他很震驚的.
就是他去到這些宜族的大營盤村.
四川.
發現裡面有一個麻風村.
裡面是完全與世隔絕.
正常人是不會進去.
那些麻風病人是不會有自己的身份.
然後更慘的就是.
原來他們的下一代.
其實是99\%沒有麻風病毒感染的.
但是被整個當時的社區所遺棄.
就是說原來那個麻風的觀念.
在很多人的內心.
還有很多錯誤的理解.
所以他見到這些小朋友.
完全沒有人理.
骯髒骯髒到都沒有人去關心.
只是好像只有一雙眼睛.
是可以見到的.
其他就是骯髒骯髒.
根本就沒有人去關心他們的情況.
更加不要說讀書.
是不可以讀書.
是沒有人會理他.
他是一個記者.
訪問了之後.
就將這個信息帶回台灣.
找到一些贊助人.
他就決心回去和當地的領導去說.
他要開一個小學.
去幫這些麻風的第二代.
將來融入那個社會.
完全沒有麻風病毒的.
但是居然人對這些族群的歧視.
那個厭棄.
是難以想像.
所以張萍儀這位姊妹.
也是信主相當好的姊妹.

$^{801}$就來到這個麻風的村裡面.
待了好幾年.
就這樣建立了一個小學.
經過這麼多年.
最近我再搜尋.
她已經不在這裡了.
在四川.
回到台灣.
她有另外的工作.
繼續參與.
心裡面有基督的愛.
一些不可能的事工.
可以變成為事工.
一個小女子.
當然了.
因為她先生也是一些跨國企業的CEO.
一定的資源.
也有一定的人際的網絡.
動員一些人的奉獻.
就可以做成一件事.
但是仍然源自什麼呢?.
內心裡不值得這麼需要的人.
沒有人關心.
當然我們不需要再追問.
為什麼社會.
為什麼福利的制度.
幫不了這些小朋友.
那些另算吧.
但是她能夠喚醒很多人的心靈.
成為感動中國的一個得獎人.
一個台灣的記者.
親愛的弟兄姊妹.
我們現在面對的是很多的福音機遇.
剛好星期四.
上水浸信會的唐主任牧師和執事.
來跟我們談合作.
因為過去他們在大埔上水的時候.
幫這個那個雲福縣.
後來有250戶已經搬了來.
我們下橋樓.

$^{841}$就是後面了.
他過去一直都有幫忙.
他們現在接觸的有差不多100戶.
他們很想跟教會一起合作.
所以過了一年之後.
我們就會談一起怎樣合作.
在平日的日子.
我們就讓這些居民過來.
參與一些活動.
一些聚會等等.
他很想將來這班朋友家庭.
信了之後.
就留在我們教會.
因為他們始終太遠在上水.
將來他們在未來幾年之後.
他們就會在古洞北來直塘.
因為北區的那個大都會的發展.
是相當大的那個幅員.
弟兄姊妹.
那個環境上面已經具備.
未來是一個很大的福音和祥.
但是我們今天需要的就是.
先裝備弟兄姊妹.
參與事奉是最尊貴的.
在我們人生裡面.
神給你什麼的感動呢?.
榮將肯定就是唱聖詩.
卡洛比也是.
有不同的弟兄姊妹.
有些人可能在管理上面.
有些人可能在其他的.
譬如美蓮關顧的.
是很有心的.
慧玲姊妹.
真是天天去探望老人家.
探望到他們能夠清楚信主.
接著我是幫他們洗禮.
其實也不是我帶他們信主.
每個月在尹健老人院.
很多這些事工.

$^{881}$兩位老師.
在福音樂曲已經不是今天了.
十幾二十年了.
弟兄姊妹.
你要在神面前開始祈禱.
有什麼神給你的負擔呢?.
你來到去裝備好.
教會一起迎接很多新來的朋友.
他們來認識主.
得到福音的好處.
所以很想在過年的時候.
跟大家用這段經文.
彼此來學習.
大家都知道.
我這兩年來了.
小朋友又很喜歡我.
老人家又很喜歡我.
其實我在感受的時候.
我根本連望客的機會都沒有.
老人家不當我做.
因為有楊姑娘,有謝牧師.
小朋友更加不用想了.
禮拜他們在幼稚園.
見都見不到.
一個都見不到.
見到的就是那些手抱的BB.
都不會認識我.
其他那些很活潑,很開心的.
全部在幼稚園.
平日是不會在感受出現的.
青年人更加.
有時候就在伯大利中學.
後來輾轉.
他們好像遊牧民族一樣.
最終回來感受.
是我爭取.
星期六一定要青少年的崇拜.
在自己的禮堂.
他們要用現代敬拜.
在台上一定要擺一套鼓.

$^{921}$是他們來帶領的.
現在的人數已經增加到四十多人.
青少年崇拜.
以往不好意思了.
在幼稚園.
每個禮拜朝行晚策.
都沒有自己的歸屬.
各位弟兄姊妹.
其實我們每個人.
都真的要參與事奉.
剛才我這樣說.
就不是說我在感受的時候不開心.
不事奉.
不是這個意思.
而是有不同的崗位.
不同的角色.
但來到一個很親密的小堂會.
每個弟兄姊妹每個禮拜都見到.
那種開心是另一種層次.
讓我能夠和弟兄姊妹有這樣的認識.
所以其實很開心.
我見到我們的茶座的成長.
我們的關故部的弟兄姊妹.
每個禮拜來探訪的成長.
另外我們現在星期四有的福音團契.
又見到一群老人家.
很開心地來參加崇拜.
參加聚會團契.
團契完之後.
一起有愛言.
一起有吃飯.
不同的群體.
所以很需要弟兄姊妹一起努力.
一起同心.
在神的家裡.
一起為著福音的緣故.
這就是第23節.
凡我所做的.
都是為福音的緣故.
為要與人共享.

$^{961}$這福音的好處.
與我們一起同心地來參與服事.
我們同心禱告.
因為其實我們根本就沒有資格去服事.
這位永生的神.
多謝你在救恩.
在耶穌的大愛.
將我們拯救.
使我們能夠在不同的崗位.
在不同的恩賜上.
一起來同心.
更重要的是我們能夠服事我們這位主.
將人帶進福音的美好.
那個生命改變.
願你得盡一切的榮耀.
也願弟兄姊妹都不同的裝備.
以致我們在神的面前.
等候主的差遣,揀選.
這樣來同心地服事.
我們獻上感恩奉耶穌基督的名.
阿們.
非常多謝周牧師的分享.
令我在事奉的生涯裡.
其實有兩段經文很影響我.
一段就是《羅馬書》十二章二節.
心意更生而變化.
叫你們測驗何謂神的善良,純善,可喜悅的旨意.
我見到周牧師剛才提及.
有很多零碎的東西.
不是做變色龍.
我們要心意更生.
可能我們會遇到很多困難挫折.
但我們要心意更生.
因為目的是什麼呢?.
福音的好處.
上帝的好處.
這個非常之好.
第二段經文在前一段很重要.
祂說:神的慈悲勸你們將身體獻上當作活祭.
是聖潔的是神所喜悅的.

$^{1001}$你們如此事奉乃是理所當然的.
我以前一直覺得奉獻是可以的.
瑪麗基書說十分之一.
三章八字十節.
現在可以了.
但新約告訴我不是.
要全人的獻上.
還要當作活祭.
將我們的身體 將我們事奉.
將我們的能力 將我們的時間獻上.
因為這是理所當然的.
這也影響我事奉的心態.
將這兩段經文和大家分享.
以下是一個奉獻的時間.
我希望大家都記得.
當你將你的金錢放進袋子的時候.
其實是你對上帝愛的一種回應.
這只是回應瑪麗基書十分之一的奉獻.
但有一個更重要背後的奉獻.
那是全人的奉獻.
那是將我們的才幹 能力奉獻給上帝.
剛才曾牧師說了很多.
一些我們事奉的機會和廣大的禍祥.
希望待會奉獻完畢之後.
請大家在你的座位思考一下.
剛才所聽到的和你生命有什麼關係.
我們現在開始奉獻.
當你不是基督徒的時候.
你可以輕輕將奉獻袋子全開.
將思想愛著靠著.
天天活在著面前.
全所有奉獻.
願如愛我尊貴救主.
全所有奉獻.
我向所有羨慕耶穌.
謙卑呼福著面前.
世俗一諾敢於撇棄.
求主圓立我心願.
全所有奉獻.
願如愛我尊貴救主.

$^{1041}$全所有奉獻.
好,請各位起來.
我們一起唱奉獻詩.
萬物都是從主而來.
我們把從主而來的應給主.
阿們.
請坐.
再一次在主裡給錢文安.
想問一問今早有沒有新朋友在我們中間聚會.
今天因為過年.
有些弟兄姊妹可能要回鄉下.
或者有其他外遊等等.
如果想問一問有沒有新朋友呢.
最近我們有一些下橋樓的弟兄姊妹.
在我們中間聚會.
很感恩.
我們簡單給錢文安.
我們有一些報告.
我們先背完今句.
第二月份的今句.
約翰福音的十章十節.
我們的主耶穌使得我們得到的是豐盛的生命.
我們一起讀一詞.
我來了是要揚得生命並且得的更豐盛.
約翰福音十章十節.
相當清楚主耶穌是要使得我們的生命.
不只是物質的這個軀殼的生命.
是整個的心靈.
來跟神有一個美麗的契合形式.
使得我們得到的是新的豐盛的生命.
我們有一些報告.
我們很感恩.
上個星期六日.
我們有美好的服侍.
弟兄姊妹在不同的地方.
剛才我都說了.
有羅牧師和王執事.
在我們中間星期四的開會.
我們已經談好了初步的合作.
我們過了年之後就會開展.

$^{1081}$如果到時需要弟兄姊妹平日來幫忙.
到時我們會再呼籲弟兄姊妹可以參加.
一起去服侍.
我們會去家房.
因為接近一百個家庭.
他們好像是幫他們安裝一些隔簾.
因為暫時還是清水樓.
是甚麼都沒有的.
有些家庭有父母和子女.
所以需要有一些簡單的間隔.
因為暫時他們簽約是兩年.
都未知道終極的方案是怎樣的.
所以暫時他們都會居住住.
所以就不會大裝修.
所以就會有這些隔簾.
和一些基本的設施.
所以過了年之後就會安裝.
所以我們雄恩的弟兄姊妹.
我們就會一起去探訪這些家庭.
就建立這個關係.
上個星期我們無論在駱橋樓.
和在教會的門口外.
我們聽志偉說.
超過五百張的揮春.
我們寫好了派發.
和很多的小禮物.
我們能夠一直都可以在社區裡面.
成為一個很好的見證.
努力的參加擺相和服侍.
我們整理好之後.
我們都會將這些相片.
我們在想怎樣發放會好.
因為有些未必有參加的.
就變成我們都要留意一下.
還有我們上個星期日.
有青年會小學的同學.
五年班的.
就在我們樓上的兒童崇拜.
所以都很感恩.
未來我們會繼續.

$^{1121}$一直跟到他們畢業為止.
希望透過這些關係.
繼續我們可以來到去讓.
更多這些小朋友.
住天水圍的家庭.
能夠來參加教會的聚會.
因為這個都是校長.
馬校長很想我們教會.
和他們一起有一個合作.
所以現在我們是有小學的合作.
也都有我們在駱橋樓的.
和上水的浸訓會.
和香港的朝人生命堂的合作.
所以這些跨宗派的合作.
似乎現在是在香港.
新的一個事奉的形態.
所以盼望弟兄姊妹都能夠.
有這樣的服侍.
譬如雷偉聰弟兄夫婦.
他是播道會的.
他又會去不同的教會教主學的.
所以開始湧現這些跨宗派的合作.
應該就是最近這幾年.
開始有一些課程.
一些神學的教導.
例如伯特利神學院的.
一個城市轉化的課程.
還有陳淑娟牧師.
她三月會再來講的.
我們現在加入的.
這個行在社區的課程.
正正就是一個轉化.
所以現在不同的資源這樣共享.
這樣一起服侍的一個國道觀的.
One Church的概念.
是很重要的.
所以盼望弟兄姊妹都來學習.
開放自己的心靈.
讓不同的資源一起去服侍.
讓基督的教會.

$^{1161}$一路來壯大.
更多人信著.
我想我們都有不同的服侍.
例如我們即將在三月底.
我們都會有一個.
現代敬拜的崇拜模式.
弟兄姊妹現在都在預備,籌備.
所以我們在洪恩家.
我們很多年都已經有福音的樂曲.
有我們的傳統的詩班.
所以我們都希望再加上.
有現代的讓年青人.
都可以參與的現代敬拜.
所以大家都為這樣的組合去祈禱.
當然泰語的弟兄姊妹.
他們用泰文獻唱.
也是非常好.
所以盼望我們更多元化.
這樣連在一起的服侍.
另外這個星期.
都是農曆年長假.
因為社區中心.
由今天開始到下個星期一.
都是放假的.
因為社區都會暫停.
所以接下來的聚會都會暫停.
教會就不開放.
要到接下來的星期六,日.
就會開放.
大家留意一下.
我們一起祈禱.
親愛的主.
我們獻上感恩.
因為你給我們有家庭的天倫之樂.
我們很盼望.
不單是地上的這份平安.
而是我們希望.
更多未得到救恩的人.
他們能夠得著.
你自己永遠的救恩.

$^{1201}$生命改變.
得到主在他們生命裡面的.
帶領.
人生的召命.
以至每一個基督徒.
都能夠知道自己.
在神的面前的崗位.
安排.
神給的恩賜.
服侍,願主施恩.
接下來這個星期日.
這個星期我們都會有.
農曆年的長假.
盼望神給弟兄姊妹.
有很好的.
家庭的溫暖.
天倫之樂.
也能夠和一些.
可能少見面的親戚.
遠房的親戚.
都有很好的相交.
真的帶出我們這一份愛的禮物.
就是我們生命裡面.
對耶穌的迎訊.
求主幫助.
也為了上個星期.
和我們一直以來.
在不同的事奉裡面的參與.
在每個月老人院的.
福音崇拜.
也讓我們能夠有很多機會.
接觸到街坊.
多謝主讓我們能夠.
參與在駱橋樓.
一些居民的探訪.
願你保守.
我們過完年之後.
和上水浸信會的合作.
真的能夠將一個跨宗派.
我們大家一起.

$^{1241}$同心服事的.
美好的一個圖畫.
在很多居民的內心.
都留下一個印象.
也讓他們對福音.
開放,慢慢地.
走出一些困難.
傷痛,或者.
對前面的安排.
還是忐忑的.
都能夠有在主裡面的.
盼望和這一份.
安穩的心靈.
求主幫助,特別在過年.
我們要記掛一些.
身體有軟弱的肢體.
高伯伯.
和玉清婆婆.
都還是進了醫院.
盼望他們,雖然都年紀大.
但是盼望他們都能夠.
可以早日的.
痊癒,出院.
能夠和家人有過年的.
那個歡聚,也都保守.
阿冰的.
乾媽.
她.
昨日都表示決志.
但是我們希望她.
很清楚地相信,然後.
真的快些的接受水禮.
求主紀念這些老人家.
親自求主.
的安慰臨到他們身上.
也都保守.
我們有很多的青少年和小朋友.
在長假裡面.
他們能夠有好的休息.
有太繁重的.

$^{1281}$功課的壓力.
之外,能夠有一個.
很好的享受.
長假給他們的.
新心靈的調整,特別中六的同學.
好像Morris.
和其他的同學.
要面對非常緊張的.
DSE的考試.
盼望這段時間他們都能夠.
有好的調整,求神幫助.
我們將這一切都仰望.
在主的手中.
這樣禱告奉耶穌.
基督的名,阿們.
\newpage



\section{哥林多後書 12:7-11}
\label{sec:POs_KWCBzOE}
\textbf{2026.02.22《因祂用厚恩待我》 哥林多後書12\_7-11 (和修版) 李富成牧師}
\newline
\newline
連結: \href{https://youtube.com/watch?v=POs-KWCBzOE}{\texttt{https://youtube.com/watch?v=POs-KWCBzOE}} ~~~~ 語音日期: 2026-02-25
\newline
\newline
\hyperref[sec:4qbR6KQJgy4]{\small{< < < PREV SERMON < < <}}
~
\hyperref[sec:index_chronic]{\small{[返順時目]}}
~
\hyperref[sec:index_scriptual]{\small{[返順卷目]}}
~
\hyperref[sec:S89_Yf51Xes]{\small{> > > NEXT SERMON > > >}}
\newline
\newline
哥林多後書 12:7-11
\newline
\begin{longtable}{cl}
\hline
\hline
章節 & 經文 (和合本修訂版)\\
\hline
12:7 & \begin{tabularx}{0.7\textwidth}{X} 又恐怕我因所得的啟示太高深，就過於高抬自己，所以有一根刺加在我身上，就是撒但的差役來折磨我，免得我過於高抬自己。 \end{tabularx} \\ \\ \relax
12:8 & \begin{tabularx}{0.7\textwidth}{X} 為了這事，我曾三次求主使這根刺離開我。 \end{tabularx} \\ \\ \relax
12:9 & \begin{tabularx}{0.7\textwidth}{X} 他對我說：「我的恩典是夠你用的，因為我的能力是在人的軟弱上顯得完全。」所以，我更喜歡誇耀自己的軟弱，好使基督的能力覆庇我。 \end{tabularx} \\ \\ \relax
12:10 & \begin{tabularx}{0.7\textwidth}{X} 為基督的緣故，我以軟弱、凌辱、艱難、迫害、困苦為可喜樂的事；因為我甚麼時候軟弱，甚麼時候就剛強了。 \end{tabularx} \\ \\ \relax
12:11 & \begin{tabularx}{0.7\textwidth}{X} 我成了愚蠢人，是被你們逼出來的，因為我本該被你們讚許才是。雖然我算不了甚麼，卻沒有一件事在那些超級使徒以下。 \end{tabularx} \\ \\
[1ex]
\hline
\hline
\end{longtable}
$^{1}$可能弟兄姊妹覺得很奇怪.
李牧師你不是12月28日來過一次嗎.
這麼快又來?.
兩個月了.
其實大家要留意.
上一次12月28日是屬於2025年.
今天2月22日是屬於2026年.
曾牧師每年會邀請我來一次.
下一次再見可能是2027年.
有很大具體的困難.
因為人只有一個.
主日有52個星期.
香港有1300間教會.
雖然我不會去1300間教會.
但我會去40間教會.
所以很感恩能夠在紅金塘一齊敬拜師父.
如果除了主日崇拜.
有其他聚會可以用得著小弟.
可以合作.
香港有很多稱號.
香港其中一個稱號叫美食天堂.
你看看這六幅圖.
都是食物.
有什麼共通點.
好吃.
好吃 邪惡.
少吃一點.
還有什麼共通點.
很有吸引力.
待會崇拜完後.
先做一道厚切牛扒.
今早有人吃菠蘿油嗎.
有青年人說「咪」.
大家還沒想到我想的謎底.
這不是智商題.
是真事.
這六樣食物.
都有一個共通點.
叫做厚切.
厚切牛扒.

$^{41}$厚切牛舌.
當我年輕的時候.
家裡開一罐午餐肉.
一家六口吃的時候.
媽媽會將那罐梅林牌圓咕嚕的午餐肉.
切得很薄.
然後讓大家可以夾.
如果切得很厚.
一人夾一次就沒有了.
所以小時候經常想.
長大後賺到錢.
自己開罐午餐肉.
切多厚就多厚.
現在去吃茶餐.
盡量不吃醃製性食品.
做兩件厚切的午餐肉.
有時候不好意思.
吃一件就不吃了.
所以現在盡量吃清淡一點的東西.
厚切.
我參與幫忙成立了一個.
齊唱四個詠團.
有七年了.
去年我們有週年的音樂會.
在荔枝角有一間宣道國際學校裡的禮堂.
舉行的.
場面都挺好看.
我們的主題叫一生為你歌唱.
我們定這個主題其實有一個想法.
有一個想法.
就是這一節經文.
詩篇十三篇六節.
我們一起讀好不好.
我要向耶和華歌唱.
因太容厚恩待我.
十三篇六節.
詩人說到他會很樂意的向神歌唱.
為什麼呢?.
是不是很喜歡唱歌這麼簡單呢?.
不是,他說有個原因.

$^{81}$這個原因是很清楚的.
是因他容厚恩待我.
在新約聖經裡.
保羅說你心被恩感歌頌神.
剛才這群哥哥姐姐.
厲害了.
我很少聽到福曲隊.
超過差不多二十人唱得這麼整齊.
指喉和喉的平均.
都是唱得很好.
在歌藝上都非常出色.
不過更加吸引我的地方.
不知道是不是你的同感.
他們是心被恩感歌頌神.
怎麼表達出來呢?.
除了歌藝是一種媒介表達.
他們差不多識背.
九成都看到你.
他用他的面容表達給你聽.
他很想將這個恩典這個福氣傳遞給你.
所以福曲隊是一個.
將上帝的恩典先領受.
透過音樂和練習.
呈獻給上帝.
亦在牧養弟兄姊妹.
因他容厚恩待我.
厚切厚恩.
當然我只是在玩食字.
不過你要想.
你何時覺得上帝對你的恩典是厚恩的時候.
厚到一個地步.
你吃得差不多滯.
你覺得滴到滿地都是滋油.
的時候.
你會洗澡又唱歌.
回教會又唱歌.
乘坐很長的地鐵.
你的心裡也在哼.
因為你很不知道怎樣表達的時候.
音樂和詩歌正在幫你.

$^{121}$向上帝發出讚歎.
這間是我安立的教會.
叫做同人浸信會.
他們有一個傳統.
每年農曆前或後.
有一個主日.
不是牧師懶惰不肯說.
他們在崇拜裡停止說.
就是一早鼓勵弟兄姊妹.
因為他們不是一間很大型的教會.
聚會人數只有三至四十人.
太大型的教會是做不到的.
例如錦宣家那些是做不到的.
因為會很亂.
但是比較小型的家庭式教會會做到.
就是整家人圍成一圈之後.
就好像團計分享見證.
叫做見證主日.
見證主日最重要的目的.
大家一起來數算主因.
好不好?數算主因.
好.
我們通常數算主因的時候.
就會說上帝怎樣幫助我們.
引導我們.
保守,看顧,帶領,扶持.
無論在病患得醫治.
在困境得到指引.
心靈閉閉得到一個復甦.
人際關係破碎.
但是上帝給你的安穩.
我們通常在弟兄姊妹的見證分享裡面.
都會聽到這些很美麗的見證.
大概一個多月前.
我和我女兒WhatsApp了一些過年的東西之後.
女兒突然又再WhatsApp我.
拍了一張照片給我.
是我的小孫.
我大女孫五歲多.
小男孫歲幾.

$^{161}$很可愛的一個.
會走來走去的.
拍了一張什麼照片給我呢?.
拍了我的孫子的胸口.
好像被人行了酷刑一樣.
有一塊紅色的.
然後她說.
阿康拍六親.
哇!六親.
哇!做阿公這些.
平時他乖就當粒糖.
但是他受傷害就插一支.
我說發生什麼事?.
原來工人不小心.
沖了一杯很熱的咖啡.
放在一個很矮的機器裡.
是工人喝的不是我孫子喝的.
我孫子已經在移動.
他移動了之後.
就推回自己那裡.
我說怎樣?.
他說我現在馬上回家.
我女婿兩個都分別在工作地點.
你知道現在父母多緊張.
我回家.
接著我說怎樣?.
去看醫生.
不幸中的大幸.
不是很嚴重.
塗藥膏就行了.
當然要洗傷口.
我說他有沒有哭?.
哭了45分鐘.
我說厲害啊.
將來如果被人折磨.
要供出一些不應該說的話.
他面對酷刑都不用怕.
哭了45分鐘.
過了幾天.
再問怎樣?.

$^{201}$再拍一張照片給我.
傷口已經好很多了.
有沒有結痂?有沒有咬?.
大家好,我是小梅.
當我們分享.
譬如兒孫遇到困難.
你的心都插一把刀.
沒什麼大礙了.
我們說這些見證好不好?.
全部都好.
全部都沒問題.
全部都值得說.
不過.
史陶保羅卻在.
《哥林多後書》第12章.
說了一個很感恩的見證.
卻聽起來.
咦?.
有些不是我們平時的套路.
牧師剛才你說你的孫子.
又沒什麼大礙了.
我們都會分享這種.
不幸中的大幸.
蒙主的保守.
很自然.
當女兒告訴我她沒事了.
都是好的.
為什麼史陶保羅在.
《哥林多後書》第12章.
分享他個人的一個.
感恩見證的前提.
比起來.
不是很棒的.
他說有根刺在他身上.
大家刺過沒有?.
刺是個比喻.
當然大家知道.
就是有東西卡住你.
你覺得不爽.
總是麻麻煩煩.

$^{241}$總是不舒服.
如果小時候我刺到.
我向我媽求救.
她就會幫我施手術.
我不知道你小時候.
你媽媽會幫你施什麼手術.
拿一支針出來.
沒有經過消毒.
就幫你挑一根針出來.
挑一根刺出來.
有沒有人玩過這種東西.
我們當時的民間藝術.
就等於我小時候.
拆選了香爐灰.
用來蓋著那些.
OK.
我跟哥林多教會的弟子們說.
弟子們我講個見證給你們聽.
我身上有根刺.
而且這根刺.
我真的覺得很不爽.
我求主將它拿走.
但是三次都沒有回應.
當然這三次.
是不是一二三.
不一定.
可能是一個虛數.
意思就是說.
多次求主將我這根刺拿走.
但是主沒有回應.
這個見證的起題.
好像有點.
咦 即是怎樣呢.
我們先問這根是什麼.
大家知道.
使徒保羅是一個翻天覆地的使徒.
因著他的勇毅.
因著他打不死的精神.
因著他經歷復活.
主向他顯現的大能.

$^{281}$他一生拼搏宣揚福音.
特別是將福音.
向外邦人傳揚.
以至這個本來最初.
只是在耶路撒冷.
或者猶太人的一個基督信仰.
能夠擴散水銀射地.
為了整個教會歷史.
這二千年來的宣教.
他可以說是樹立了一個先驅.
如果說影響二千年人類歷史的.
除了耶穌基督之外.
如果以人來說.
使徒保羅都應該排行很前.
最低限度.
我們今天特別在基督新教.
或者叫更正教.
堅持的一個基要教義.
叫做恩信甚麼?.
恩信清義是從.
保羅的書信.
即是保羅的神學裡面提煉而來.
你和我都受影響.
這個這麼厲害.
上帝使用的僕人.
他說他有根刺卡在這裡.
求主多次母下文.
你有沒有經歷過這種掙扎?.
你有些事情.
祈禱多次母下文.
保羅這根刺是什麼?.
在經文裡面.
自己沒有解釋.
保羅在經文裡面.
在書信裡面沒有解釋.
因此有個空間.
我們來思考.
按照我們認識保羅的生平.
和保羅的書信.
第一.

$^{321}$他是不是因為不斷受到迫迫.
他很想不再面對迫迫.
而就是這根刺呢?.
這個有可能.
因為你讀史唐行傳的時候.
你發現保羅去到哪裡傳福音.
都有班人追著他.
那班不是粉絲.
那班是猶太律法主義者.
他們說這個保羅.
本來以前叫素羅的.
是和我們同一個陣營.
去打壓耶穌.
去捉拿基督徒.
現在他調轉槍頭.
成為了我們的背叛者.
還到處為那個什麼耶穌大發熱心.
我們誓要將他.
剷除.
誓要將他殺害.
史唐行傳記載過一次.
猶太有四十多個殺手.
真是殺手.
因為他們立了誓.
不殺保羅.
他們不吃不喝.
不保羅死就是他們死.
他們立了一個死誓.
一定要殺保羅.
不過這個伏擊行動.
外洩.
他們的情報做得很差.
所以羅馬軍官.
將他移送去一個高設防的監獄.
變相保護了保羅.
所以你看到保羅.
面對這種的逼迫.
是靠近的.
你有沒有想過.
你信耶穌每天都被逼迫.

$^{361}$你會不會求主將這根刺拿走?.
會的.
很正常.
所以有可能是這個原因.
第二個原因有可能是.
保羅未信主之前.
其實他是不斷殘害信徒.
即是說他滿手基督徒鮮血.
他雖然信了主.
悔改.
但這種殺害基督徒的陰影.
會不會留在他的心裡.
他有時會觸動.
可能見到信徒的家人.
原來你就是我殺父仇人.
會不會是這個原因.
令到他這份傷痕.
磨滅不了呢?.
以至好像一根刺那樣.
刺住他呢?.
亦都非常有可能.
不過近代的研究.
更加認為比較可信的可能.
是第三個.
是什麼呢?.
是攻擊.
是教會內的.
假教師對他的攻擊.
我相信在紅恩堂.
除了曾牧師傳道人外.
沒有假教師.
因為在今時今日.
我們很容易分辨到.
因為第一.
拿一本聖經出來量度一下.
第二拿一個使徒信經出來量度一下.
跟著再觀察他的行為量度一下.
我們很快很容易將一些.
混亂信仰或者.
在教會所謂拿油水攪攪陣的人.

$^{401}$將他處置.
不過在初期教會.
都未有一個叫做.
成文的聖經的權威.
李保羅最厲害嗎?.
還有其他人的.
有些人不會服於保羅的權柄.
你說你經歷過主向你顯現嗎?.
我都經歷過了.
怎知道?.
所以在保羅的事奉生涯裡面.
其實他都有對頭人.
其實哥林多後書整卷.
哥林多後書你會看到.
保羅和哥林多教會.
有部分的人是有張力的.
有些人會幫他搞破壞.
有些人會製造麻煩.
有些人可能會挑釁他.
有些人擺明車馬就不信你.
當你作為一個傳道人.
一個牧者要去帶領.
要去服侍的時候.
你遇到這些所謂內亂的時候.
那根刺可能比起外在的攻擊.
更令你睡不著.
我們不肯定哪一個是他那根刺.
不過他卻在第九和第十節.
講了一個逆轉的見證.
「丁姊妹,我有根刺插在我身上」.
「多次求主都不拿開」.
「不過那又怎樣?」.
就是這樣.
一起讀第九和第十節.
「他對我說:我的恩典救你容的」.
「因為我的能力是在人的軟弱上」.
「顯得完全」.
「所以我更喜歡誇自己的軟弱」.
「好叫基督的能力覆蔽我」.
「我為基督的緣故」.

$^{441}$「就以軟弱,凌辱,急難,逼白」.
「困苦為可喜樂的」.
「因我在麼時候軟弱」.
「甚麼時候就剛強了」.
史蒂芙·保羅跟哥林多教會.
說了一個他感恩.
他說「神用厚恩待我的見證」.
這個方程式是.
「縱使他有根刺」.
「主沒有按照他的心願將它拿開」.
「他仍然見到他用厚恩待我」.
他說「神的」.
你看到我為什麼用顏色的字.
「我」綠色就是保羅.
有些「我」紫色就是神.
紅色就是恩典和逆境的對比.
最後藍色就是結論.
你會看到這段經文很有趣的地方.
逆轉了甚至顛覆了我們對感恩見證的套路.
我再說一次不是那些有問題.
不過我們應該分享一種見證.
「縱使那個刺沒有離開過」.
可能頂梓湖在事奉中覺得自己撲心撲命.
但不被人肯定和欣賞.
甚至有時被人誤會,遺漏,埋怨.
「最壞是你」.
那一刻最壞是你插到想死.
自己的兒孫面對很多困難.
自己好像無能為力.
自己的靈命有時跌入一個幽谷.
好像起不來.
有時你會覺得要經歷一段不知多長的時間.
那個黑夜究竟有多長.
你不知道.
你都會求主一次兩次.
有時你都會灰心.
不過史淘寶跟哥林多教會說.
「那又怎樣?」.
「縱使這根刺沒有離開過」.
「但原來在這個逆境裡」.

$^{481}$「我很肯定神會這樣回應我」.
「祂說我的恩典夠你洗」.
大家想起「夠你洗」會想起什麼?.
夠不夠洗?.
有些事原來是夠不夠洗?.
「恩典」神說「夠洗」.
「夠洗」的恩典不在於.
神將你開的支票全部幫你簽了.
「夠洗」的地方是.
祂與你同在.
陪你經歷困難.
所以在軟弱的人身上顯得完全.
所以保羅說經歷過這個之後.
他不是有自虐狂或被虐狂.
因為他說以什麼為可喜樂?.
軟弱,凌辱,急難,迫不及待,困苦.
這些是喜樂.
這完全逆轉我們正常.
祝你順風順水,家肥屋潤,盆滿缽滿.
我們祝福人有這些好東西都不是有問題.
不過我總不會對周牧師說.
祝你軟弱,凌辱,急難,困苦.
為周牧師今年的喜樂.
不過我們說笑我們不會這樣說.
但我們在主裡面生活一段時間.
事奉主一段時間.
我們知道根本沒有一條一帆風順的基督徒人生.
如果有人這樣跟你說.
他除了是祝福語之外.
他有機會說錯話.
他會誤導一些初信者.
信耶穌就大魚大肉.
信耶穌就家肥屋潤.
信耶穌說病就一定沒有死.
不一定醫治你.
如果我們祝福人沒有問題.
但我們不要將他變成教義.
所以有些人信耶穌就覺得被人騙.
信耶穌就不是這樣.
原來真相不是這樣的.

$^{521}$如果信耶穌就像我們祝福人那樣.
每個人都家肥屋潤.
每個人牧師不用去探病.
牧師不用主持安息禮拜.
信耶穌不用死嗎.
死得好而已.
因為我們安息主懷.
頂主妹.
什麼叫做因他用後因代我.
從保羅這個見證.
詮釋了一個這樣的角度.
在軟弱裡靠主剛強.
這就是和一般沒有耶穌生命在我們裡面的人.
最大的分別.
大家知道每逢年初一.
一踏入年初一.
有一個地方很多人爭著去插東西.
那是什麼地方.
哇 多誠心.
我們教會這麼誠心就好了.
每個人都來插位坐.
電視就播放黃夏蕙.
為什麼那些人這麼誠心.
因為他們要求什麼.
他們要求一個.
今年黃大仙照著我.
成功了.
我下年就來玩神.
不成功就不會告訴你.
如果黃大仙真的有求必應.
你不用讀書.
我不用努力工作.
只要找黃大仙就成功.
當然有一萬人去拜黃大仙.
今年一定有一千人順風順水.
不過有九千人不會跟你說.
基督徒不是這樣.
基督徒要告訴信主和未信的朋友.
我們一樣面對軟弱.
我們一樣面對逆境.

$^{561}$不過我們多了一個寶貝.
那樣叫抗逆力.
不過這種抗逆力.
現在當然在一般學校都會說抗逆力.
因為現在的小朋友.
就像士多啤梨.
什麼叫士多啤梨.
士多啤梨一壓就變成怎樣.
扁就不夠全神.
一壓就變成啪.
承擔不起.
我小時候.
我們父母和學校的老師.
任意抽擊我們.
我們那時候是怎樣的.
打打打.
最初當然是哭.
痛的打打.
之後會怎樣.
打打打.
我們的抗逆力很強.
我們抗逆力是怎樣強呢.
所以旁邊的黃師奶會告訴我.
李師奶不要打兒子.
不要打得那麼硬.
打得很沒用.
打得很硬就麻木了.
打打打.
我們又不會跳樓.
又不會死.
我們只會想.
你打吧打吧.
打到我讀完小學就去找工作.
找工作我就自己搞定自己.
那時候的思考是這樣.
所以那時候的抗逆力就是這樣.
現在的小朋友因為在太過風雨和呵護的環境長大.
等於我小時候.
我們不用去學攀石.
自己會去周圍的山那邊爬.

$^{601}$我們踢球不用找阿仙奴足球學校.
我們走去球地穿著白飯魚就能踢到一身球技.
不過現在不踢而已.
聽著小妹.
我們的抗逆力信徒一樣需要從神起的這種我們叫做屬靈的抗逆力.
這個抗逆力其實我們借用保羅在同樣哥林多後書四章第八第九節.
我們很熟悉的這一句話.
我們四面受敵.
卻不被困住.
心裡作難卻不自失望.
做壁伯卻不被凋棄.
打倒了卻不自死亡.
這兩節的經文真的翻譯得非常好.
是不是很公正.
應該這樣就玩完了.
誰知道又不是.
一個這樣的對比.
其實四樣都是在說同一樣東西.
當然很可能我們從保羅生平我們可能找到一些更具體的東西.
例如打倒了卻不自死亡.
在《時代行傳》記載過有一次保羅去到外邦傳道.
被那些猶太主義者攻擊.
他們用石頭扔他.
到一個地步那些攻擊的人以為他死了.
如果我要攻擊曾牧師.
搞到大家以為死了.
他應該的面貌是怎樣的呢.
爆發了.
奄奄一息.
那些人以為他死了就把他拖出城.
打算棄屍.
《路加》記載很幽默的.
那些保羅也有他的追隨者.
那些門徒走過去看看師父死了沒有.
誰知道走過去看的時候.
聖經描述保羅站起來.
然後第二天就繼續去傳道.
好像什麼事都沒有.
你讀到這裡覺得《路加》是亂作的嗎.
以為死了又站起來.

$^{641}$我讀這段聖經.
其實有些想像空間.
總之不要當它是教義就行了.
有兩個可能.
第一個可能.
神蹟.
上帝親自介入.
打到死.
馬上變他.
沒事 站起來.
有沒有可能.
有 神蹟.
神蹟就是超越我們能夠想像.
所以在小行傳真的記載.
使徒保羅有行神蹟的能力.
上帝的主權如果讓他.
從一個將死的狀態.
一下子站起來.
可以.
第二個從比較人性的角度去想.
可能是什麼.
他真的被人打到軟軟的.
門徒走過去.
被紙巾的棉花頂住.
然後他慢慢被人扶著.
站起來.
一下子一拐一拐.
回到客棧休息一晚.
第二天.
他繼續.
打甩門牙和血吞地.
去傳道.
聽不明白我剛才那句言語.
聽過沒有.
打甩門牙和血吞是什麼意思.
即是我和這位弟兄鬥拳.
我被他打甩了門牙.
我應該棄權了.
不是.
我繼續打甩門牙.

$^{681}$吞了又怎麼辦.
我繼續和他競賽.
這是說一個人面對迫不得力的時候.
他沒有放棄.
打甩門牙和血吞地.
站起來.
繼續傳道.
打倒了卻不至死亡.
弟兄姊妹.
我們讀這節經文之前.
我們必須要讀第七節.
因為我們不讀第七節.
我們會過於將士徒保羅神化.
他就行的.
有什麼可能.
應該這樣就這樣.
不是這樣就這樣.
我不行的.
打甩門牙和血吞.
他摸我一下.
我已經不相信耶穌了.
攻擊來的時候.
打甩門牙.
保羅這麼厲害.
只有一個就行.
不是.
保羅在說第八第九節之前.
他說了一個秘訣.
而且這個秘訣的主語.
請留意我一直藍色的字.
叫做「我們」.
這個「我們」最低限度.
不只是保羅一個.
是說他和宣教的團隊.
我們借用這個引申的意思.
這個「我們」.
包括你今天和我.
去跟隨耶穌的弟兄姊妹.
我們都可以經歷這種.
在軟弱裡靠著.

$^{721}$剛強的一種生命力.
秘訣在第七節.
我們一起讀.
我們有者寶貝放在瓦器裡.
要顯明者莫大的能力.
是出於神.
不是出於我們.
保羅也是一個有血有肉的人.
和我們一樣血肉之軀.
被人打回痛.
被人擊會睡不著.
弟兄姊妹.
什麼令他有這種力量去抗疫呢.
他說有寶貝在他裡面.
保羅這個藍色的字的一個比喻.
大家留意「我們」的比喻是什麼.
「瓦器」.
你和我都是瓦器.
什麼叫瓦器.
一碰就怎樣.
難啊.
我今天有多剛強.
可能我明天就軟弱了.
這個弟兄在教會很多年了.
為什麼這段時間.
他好像消聲匿跡.
我們只不過是瓦器.
保羅很明白.
他和他的事奉團隊.
面對這麼大的逆境去事奉主.
他們只不過是一個器皿.
而且要形容這種器皿是瓦器.
一碰其實都受不了.
我想你和我在事奉歷程裡.
都經歷過有些臨界點想放棄.
真的想放棄.
我為什麼還要這樣做.
一年我在我舞會做執事會主席.
大概我四十歲出頭.
那年我舞會八十週年.

$^{761}$很大的慶典.
組織了一個除了標準的崇拜.
我們橫社會都一定會做什麼.
吃一頓厲害的.
叫做感恩聚餐.
大概二十多年前.
二十年前.
當時在澳海城還有一間叫聯邦酒樓.
當時的酒樓還可以容納很多位.
八十週年.
大家認為擺多少位.
當然是八十位.
我舞會明年就一百週年了.
要找到一間酒樓擺一百位很難.
我們現在只是預算六十位.
八十位.
很大場面.
我是執事會主席.
當然是全力去做.
拼命去預備.
我就跟執事說.
各位執事.
記得散席的時候.
要在門口做什麼.
送蝦.
因為那天八十席.
除了鄧英姐妹.
還有很多鄧英姐妹邀請來的親朋戚友.
我們作為主人家.
我們都要多謝人.
我們一起送蝦.
我們當時的舞會的執事會.
有十五個執事.
我加起來十四個.
我心想沒有一半.
起碼有七八個動在這裡.
多謝啦.
但實在太開心了那天.
送蝦的時候.
我一支公棋在那裡.

$^{801}$每個人都想著拍照.
之前大家的默契都忘記了.
我就一個人在送蝦.
我不覺得有什麼問題.
我都不到處找他們.
他們都是開心.
我就代表教會送蝦.
一個星期之後.
我在教會.
當時我們有些執事信格.
有些東西放在這裡.
有張便條.
便條是匿名的.
你猜他寫什麼呢.
他說李富城.
我認為你那天.
騎劫了教會的八十週年聚餐.
好像你自己搞聚餐.
我馬上爆哭.
為什麼呢.
因為我覺得很委屈.
我做到盡.
那天很開心.
收不到一句半句的讚賞.
這些都預料到.
竟然收到一個.
你可以當是一個誤會.
認為我霸佔送蝦.
不像你擺酒一樣.
但第一他不明白.
有14個人甩我底.
製造這個誤會.
我覺得最大困難是什麼呢.
匿名.
我無從解釋.
即是說這個死貓.
我吃硬了.
吃到今天.
我不可以利用教會的程序表.
登一段廣告.

$^{841}$就說李富城那天.
有14個執事甩我底.
令我製造了一個誤會.
好像我搞一個宴會為自己.
我不能夠這樣做.
我和當時的主任牧師說.
主任牧師說你吃硬了.
我說我明白了.
我做執事都知道.
我們經常吃黃年.
什麼吃黃年.
啞仔吃黃年有苦.
那次去到林大點.
林大點就是牧師我不做了.
我當時在仁濟做校長.
我說在仁濟我有很多世界.
你這裡搞.
我做到這樣.
都收到這樣的東西.
如果我真的處心積慮.
要騎劫教會.
抬高自己.
我活該.
但我說這裡不是教會.
因為那個人匿名.
匿名就不是弟兄姊妹.
無法溝通.
無法調解.
我吃一輩子.
說到現在都差點眼淚泛光.
我當時牧師比較胖.
他叫我進來.
牧師只和我做什麼.
只抱著我哭.
他哭我哭.
弟兄姊妹.
當你去面對.
一些這樣的.
不明白.
委屈.

$^{881}$你這次是什麼.
我們想想.
原來我從小回到這教會.
我都只是一個啞氣.
原來去到這個臨界點.
就是我都想放棄.
經過調整.
我繼續做.
今天仍然很愛我的.
母會.
如果我們不將.
這個寶貝放在我們裡面.
你和我都沒有能力.
去面對這次.
我們只能夠說一些.
不幸中大幸的見證.
我們不能夠.
和弟兄姊妹分享.
弟兄姊妹.
我求主幫我拿開這個病很久.
不過都沒有.
我是很辛苦.
你可不可以為我祈禱.
這個年輕的牧師.
他的年紀和我女兒一樣大.
所以我當他是我的兒子.
我認識了他很久.
有一年.
美師來香港踢足球.
大家知道是誰嗎.
國際最好球的人.
我認為.
竟然來香港現場.
可以看.
那時候他是受一間.
保險公司贊助的.
我認識那間保險公司的.
一個從業員.
請幫我頒兩張票.
880元,不是很貴.

$^{921}$看到國際級球星.
我不用飛去看.
拿880元在香港大球場.
看到.
我找了這個.
他是一個運動員.
年輕的時候.
曾經是香港田徑隊青年軍的代表.
黃勝.
我找了兩張票一起看球賽.
他和我一起去看.
我拍下了這張照片.
那天我和他當然非常失望.
因為美師那天.
只是在球場亮相.
沒有踢球.
所以我們在等.
等了上半場45分鐘.
下半場等了44分鐘.
我們說糟了,球證快吹哨完場.
他說是.
於是全場都一個聲音.
你們做得非常好.
這是我很好的弟兄.
這是他.
按木慢現影.
他成長的過程.
其實他在做青年軍的時候.
因為一次意外.
他要離隊.
做了很嚴重的.
手術.
他經歷過一個人生很大的身體上的和緩.
但神卻這樣.
將他從運動員的生涯.
帶到教會做幹事.
他很熱心事奉.
我在一間教會.
做半職傳道的時候.
他在教會做福音幹事.

$^{961}$他其實是我的前輩.
我就留意到這個青年人.
很熱心愛主.
我就和他說.
如果將來你離開這間教會.
記住我不會滑角.
你離開的時候.
記得通知我.
因為我母會很需要青年同工.
你記得.
你離開這間教會.
不要告訴我.
總之你離開.
你告訴我.
我不可以搶別人的精兵.
過了幾年後.
他跟我說.
我離開了某些教會.
我當時已經全職事奉.
我當時已經沒有一官半職.
我找了執事會主席.
大家坐下來聊天.
他就來了我成長的教會.
做青年牧者.
他結婚.
年輕人都很懂事.
很喜歡搞一生一世.
找了一個日子.
誰知遇到疫情.
他訂了長沙灣.
以馬來尼浸信會.
坐400人.
因為網絡很大.
他預計會坐滿.
那時疫情.
結婚的場面.
400人禮堂.
坐了20人.
男家一行.
女家一行.

$^{1001}$當時的政策是可以行婚禮.
20人.
女家有個舅父不知病.
沒有報名.
被教會幹事趕走.
如果被檢查違法.
很大件事.
20人.
我因為負責領事才能進去.
這算什麼疫情.
大衛疫情.
400人禮堂.
只有20人.
那晚喝酒.
4人一桌.
很過癮.
4人一桌.
真是糟糕.
他走來走去.
我們只是開玩笑.
不過.
這位牧師.
年輕傳道人.
牧師遇到一個很大型的.
不是啊.
我都勉強回到.
舞會唱斯班.
看到我斯班的白頭.
出現在這個聚會.
我特意圈了.
綠色的小朋友.
是他的愛情.
結晶品.
豪言.
為什麼在他按牧禮.
豪言.
我這麼觸動.
當他知道準備要.
按立.
帶領舞會踏入100週年.

$^{1041}$大家都很開心.
有個年輕的牧師.
帶領自己的教會進入100週年.
他也太太懷孕了.
很開心.
因為大概小朋友.
去年4月出生.
按立.
住進教會宿舍.
一切美好.
在去年.
農曆年之前.
他這個.
小朋友早產.
27週出生.
我們當時收到這個消息.
28週出生.
我們收到這個消息.
心如刀割.
早產有什麼後遺症.
接下來的消息.
就是這麼小的小朋友.
已經要做手術.
大家想想.
這麼小的一團肉.
醫生的手法.
真是很好.
怎樣開刀呢?.
做腸的手術.
做完一次之後.
再過兩個星期.
要再做.
當時有一個晚上.
星期五黃昏.
我發訊息給父母.
我的小朋友又要再做手術.
那次我有一個很大的感動.
我說我來兒童醫院.
一起吧.
去到醫院.

$^{1081}$他的外父外母.
是基督徒.
太太還在休養.
邱宏勝.
Vincent.
牧師你來了.
你和我外父外母.
到手術室等.
小朋友將要推到手術室.
我說好.
那次我看到.
小朋友從手術室.
那張床推進來.
那次我看到.
那團肉有多大.
爸爸要跟進手術室.
因為要解釋風險.
接著爸爸出來.
外父外母是基督徒.
有牧師在.
自然叫我做什麼.
牧師帶我們祈禱.
你祈禱嗎.
根本四個人.
牧師反應.
祈禱.
我們低頭祈禱.
都說不出就哭.
怎麼說.
哭完一輪.
都要祈禱.
哭完一輪都要祈禱.
最後都要哭.
在整個教會裡.
大家和牧師一起.
去面對這件事.
救得回來.
養不養得大.
害不害怕.
長話短說.

$^{1121}$香港的兒童醫院.
非常好.
住到.
六月.
才能夠離開.
兒童醫院.
慢慢穩定.
在《愛牧禮》裡.
這個小朋友第一次出場.
全教會的弟兄姊妹.
除了為.
幽牧師感恩.
他安立之外.
更加為他這次.
上帝在未來的日子.
如何照顧他的孩子.
我們一起教堂.
當然還有些笑話.
他按牧禮那天打風.
按不到.
移到八月.
晚宴.
按牧禮.
但酒樓處理不來.
有些牧師.
按牧團也處理不來.
於是按牧禮.
結了婚才擺酒.
九月才按牧禮.
按牧晚宴.
其實連按牧這麼簡單的事.
都打風打中你.
其實他也有仇家.
有些人在Facebook.
輿論他.
有間教會有人按.
按不到.
原來可以這樣.
頂主妹.
我為他感恩的地方.

$^{1161}$他大概三十七歲.
未來.
我舞會一百週年的日子.
由一個有經歷的牧者.
雖然年輕.
但是.
他明白.
打倒了卻不知死亡.
去帶領.
就讓我想起.
這首詩歌.
錯節有理許可.
要讓.
祝福萬國.
至死痛楚.
要將生命加多.
你破碎我一切.
卻換上.
更多恩惠.
榮耀路徑.
必有雷蹄.
我開頭有說.
我幫忙.
建立了一個詠團.
其中有一對夫婦.
應該.
可能很大機會.
是這個詠團年紀最大.
那位男士.
你見他穿得像尖士邦.
他在.
那天的聚會.
他被訪問.
他說我.
七十五歲了.
我們才第一次知道他七十五歲.
神的恩典.
夠這位陳兆權弟兄用.
因為神的能力.
在這個七十五歲的.

$^{1201}$事奉者身上.
顯得完全.
他跟我一樣站了兩個多小時.
在台上.
他為基督的緣故.
就以站到腳軟.
唱到口乾.
為可喜樂地.
因為他什麼時候.
軟弱.
什麼時候.
就剛強了.
他怎樣得到.
神的厚恩.
他在詠團的服事裡.
得到神賜的厚恩.
基督徒不能躺平.
你躺平.
你永遠經歷不到這種厚恩.
你要經歷厚恩.
你必須要.
參與在整個.
事奉和上帝同行的.
經歷裡.
你才會.
經歷到這份厚恩.
他七十五歲.
可以站兩個多小時.
我們今年.
在香港大會堂.
1400人的演奏廳.
有今年的演唱會.
我相信他會用最大的動力.
再動多一年.
動多一年我不知道.
他會以此為.
他的喜樂.
我見到他和他太太.
每次來練習.
我就很感動.

$^{1241}$這首詩歌.
值得我們一起.
作一個回應.
來反思.
你有沒有一根刺.
今天幫你.
去經歷.
上帝的厚恩.
和我一起唱這首.
《有你許可》.
(歌詞).
挫折有你許可.
要讓.
祝福萬過.
至死同挫.
要將.
生命.
加多.
你破碎我一切.
卻換上更多恩惠.
榮耀路徑.
必有擂台.
誰肯放棄自救.
才將舊恩釋透.
葡萄被壓才得釋放.
讓初心走.
人生到了絕處.
才知道愛的深處.
才能勝在神恩豐滿傾注.
誰肯放棄自救.
才將舊恩釋透.
葡萄被壓才得釋放.
讓初心走.
人生到了絕處.
才知道愛的深處.
才能勝在神恩豐滿傾注.
人生到了絕處.
才知道愛的深處.
才能勝在神恩豐滿傾注.
讓我們同心禱告.

$^{1281}$多謝主你透過使徒保羅的見證.
讓我們看到我們都只不過是啞氣.
卻是有寶貝在我們裡面.
以致我們能夠在軟弱裡靠主剛強.
讓頌鼎前輩給你的靈繼續啟迪.
我們心裡有沒有一根刺.
但上帝你卻沒有拿開.
成為我們繼續經歷神恩典的一個途徑.
因為你的恩典救我們肉.
求主幫助我們.
祈禱奉耶穌基督的聖名.
阿們.
\newpage



\section{馬可福音路加福音 6:34}
\label{sec:S89_Yf51Xes}
\textbf{2026.03.08 馬可福音6\_34;路加福音15\_20 (和修版) 曾裔貴牧師}
\newline
\newline
連結: \href{https://youtube.com/watch?v=S89-Yf51Xes}{\texttt{https://youtube.com/watch?v=S89-Yf51Xes}} ~~~~ 語音日期: 2026-03-11
\newline
\newline
\hyperref[sec:POs_KWCBzOE]{\small{< < < PREV SERMON < < <}}
~
\hyperref[sec:index_chronic]{\small{[返順時目]}}
~
\hyperref[sec:index_scriptual]{\small{[返順卷目]}}
~
\hyperref[sec:Jj4ZioVnCDo]{\small{> > > NEXT SERMON > > >}}
\newline
\newline
$^{1}$各位兄弟姐妹,主佬們平安.
我要等一等powerpoint.
我們先一齊同心的禱告.
親愛的天父.
我們今早能夠在聖餐的桌前.
去默想你的生命.
賦予我們永恆的價值.
也讓我們能夠出死入生.
得著的是永恆的救贖.
而且更加寶貴的.
我們在神面前.
是有一個神兒女的身份.
是永遠不會被剝奪.
因為我們是屬於你.
也在你的救恩裡面有份.
我們真的很盼望.
我們能夠以此為我們很重要的根基.
我們去將救恩帶給有需要的人.
今早的正道求主.
透過聖靈.
感動我們每一個.
都讓我們能夠看到.
原來救恩可以帶給人.
生命很大的釋放.
拯救和夢想.
我們今早就等候在你的面前.
小小的訊息.
能夠成為我們彼此的激勵.
我們這樣的感恩禱告.
奉耶穌基督的名.
多謝弟兄姊妹.
我們有八位.
在上個禮拜.
星期四到星期一.
我們在菲律賓.
參加野地牧人的.
這個短孫的體驗.
那個考察.
去了解到底當地信徒領袖.
他們接受連續六天.

$^{41}$早午晚的培訓.
當然我們沒有參加那個培訓.
我們只是觀察了一個很短的一個上午.
然後我們主要兩天.
就會在一些叫做垃圾村.
來到去走訪這些民居.
去到裡面和這些小朋友.
來到去run一些program.
來到看到他們這些小朋友的情況.
給我們一班的短孫隊的弟兄姊妹.
看到什麼是一個很大的生命的落差.
所以我今天很想在今早的正道.
在我的內心裡面仍然激盪著的那個感動.
和大家來到去作出一個分享.
當然這個不是一個整體詳細的分享.
只能夠是透過兩段經文.
和大家來到去.
分享一下我內心裡面的那個激動.
那個內心的那個.
在這個幾天裡面.
如果一個禮拜前.
你問我對菲律賓的印象.
一個罪惡的世界.
一個不會去的地方.
從來不會踏足的地方.
從來不會羨慕想去的旅遊的國家.
我昨晚就上網去搜尋一下.
日本香港人每年在疫情之後.
大概有二百多萬人一年去日本.
當然有些人是一年去五次六次.
變成那個量更大.
但我再問去菲律賓的香港人有多少.
對不起沒有這個數.
因為根本沒有統計.
根本沒有什麼人會去菲律賓.
菲律賓來香港旅遊的.
要累積了很久才有一百萬.
你可以想像到這個國家對我們來說.
其實一點都不遙遠.
但是很陌生.

$^{81}$所以在我去這個短訊之前.
這個國家在我的印象.
是非常之惡劣和差的.
隨時會擔心在市中心.
就會被警察叫停你的車.
拿錢等等.
隨時會有聽聞.
你可以想像到對我的印象是極差的一個地方.
不過這裡給我們一個很大的思想.
一個我最不想去的國家.
帶著的真是有一個偏見.
帶著的真是有一個很大的誤解.
甚至乎應該是一個以為自己的優越.
好像這個地方沒有任何的可取.
沒有任何的值得我們留下來觀察.
但是當我們一群.
我們都是大概八個短訊隊的弟兄姊妹.
敢宣的和我們洪恩家的幾位.
我們去到見到的.
有四十幾個信徒領袖.
來自菲律賓不同的地方.
偏遠的很遙遠的小島的領袖.
不辭勞苦坐船.
很貧窮的地方.
來到馬尼拉市.
來接受那個裝備.
連續六天接受裝備.
和他們本身都有自己的工作.
來做小生意.
來謀生.
但是放下那些工作.
就來到裝備.
就是希望回到他們自己的地方.
來做福音的使者.
來去有系統去教導聖經.
這個是我們很短時間的觀察.
但是我們更多的就是連續兩天.
是早午來到這些貧窮的垃圾村.
來去看一下裡面的人.
聽裡面的人的故事.

$^{121}$和聽宣教士.
為什麼會留下來七年.
來義務去幫助這些小朋友.
來希望他們能夠生命有改變.
當然他們都安排了有好幾位.
可以接受幾年的.
已經幾年的資助的小朋友.
當然已經中學了.
來去接受那個學校的教育.
他們來講述他們的故事.
其中一位的姊妹.
她已經是去到高中了.
是三姊妹同一個媽媽生的.
其中兩位接受那個助學的.
三個姊妹.
親姊妹來的.
不同爸爸的.
就是三個爸爸一個媽媽.
還要是在菲律賓是普遍.
是很難以想像一個這樣的結構的社會.
令到這些下一代那個的痛苦.
和很多的無助.
但是聽這一位的小姊妹.
用很流利的英文.
來講述她的成長故事.
其中一件事令到我和我們的隊員.
是非常之擇心.
她說我們能夠進了這個學校之後.
我們接受幾年的幫助.
令我可以對前途有少許的夢想.
可以想一下.
我接受完這些技能訓練之後.
我可以找到一份很好的工作.
來去謀生.
孫教士再和我們說.
他說能夠讀到書有助學.
他們在學校有裝備.
有教育就不會學壞.
接著原來第一天帶我們去看的那些.
是很整齊的但是很骯髒的民居.

$^{161}$帶我們去看.
已經算是很好.
和很多的小朋友是可以.
進入那間基督教學校受教育.
有一些是助學得到教育的.
原來如果沒有得到教育.
就流連在街頭.
就由得你.
根本就沒有人會理你.
政府也不會強制你讀書.
這樣就會變成街童.
變成沒有人理的.
他們的前途就是長大.
女兒盡快找個男朋友.
有人照顧就算了.
男的很快就學壞.
開始去販毒入黑社會.
而且是一個很結構性的問題.
孫教士真的很好.
第二天就真的帶我們去貧窮的垃圾村.
我們看到的景象.
令我們感覺到很大的痛心.
因為原來貧窮.
是會令到他們完全被整個社會遺忘.
大家看到的這個地方.
是每家每戶都是這樣.
你難以想像.
我們今天的香港已經不會再有這樣的貧窮.
是我們50年前.
我們小時候長大的時候.
沒有屋區長大的時候.
那個環境.
是我們小時候重要惡劣的環境.
就是他們的日常.
就是他們每天生活的那個場景.
和小時候的小朋友.
根本是不會走出他們的垃圾村.
他們所撿到的只是很少很少的那些膠樽.
來希望換取十塊八塊的Peso.
這樣來買日用的飲食.

$^{201}$有時孫教士說.
每天可能只吃一餐.
一餐.
所以你很難想像.
我們想到的小朋友.
你會幻想一下他們是會怎樣的呢?.
但是最震動我們大家的心靈.
就是每個小朋友.
充滿了童真的歡笑.
喜樂和天真.
對著我們這些外國人.
個個湧過來和我們聊天.
擁抱我們.
等一下我給大家看一張照片.
有一個其中的小朋友.
我們要走了.
抱著我不放.
要我抱他.
這樣的一個場景.
我想起.
其實我當時已經想起這一段聖經.
告訴我們.
主耶穌如果見到這樣的環境.
祂的內心是怎樣呢?.
這裡告訴我們.
耶穌走遍各城各鄉之後.
祂見到很多猶太人.
前來找祂.
祂說祂憐憫他們.
祂見到的是他們的困苦.
是他們的流離.
他們好像沒有目人.
沒有領袖.
沒有明天.
更加最重要的是.
根本沒有盼望.
根本不存在對將來的夢想.
你有沒有想像一個沒有夢想的小朋友.
他現在在一個村裡面.
他不知道世界的凶險.

$^{241}$他還是在那個村裡面.
因為是一個垃圾村.
是不會有人去的.
但是原來他一路長大.
是沒有夢想的.
整個貧窮令他們不敢夢想.
但是耶穌見到的.
是猶太人的困苦.
是流離.
是不知道自己將來要去到神的天家.
他們是猶太教徒.
他們不是沒有吃東西的地方.
不是沒有錢.
但是當時耶穌見到的.
是困苦流離的人.
是沒有了目人去牧羊的群羊.
走散的迷失的羊.
但是我們見到的.
是一群自己坐下來.
開始聚會了.
小朋友就自己自發跑來坐下.
準備聚會.
來聽那些姐姐帶的唱詩歌.
做手工,做遊戲和分組.
就是這些小朋友.
這張照片是很有代表性.
因為真的好像坐在監獄裡一樣.
看著我們的眼神.
這樣告訴我們.
好像述說他們的無助一樣.
是令到我們的心靈.
在現場的弟兄姊妹.
是很觸動我們的內心.
耶穌見到的不是很多的問題.
我帶著的是一路坐車.
一路在控訴那些交通.
五十年都沒有變.
由下飛機在馬尼拉的大城市機場.
就已經塞車了.
塞了一個小時.

$^{281}$才去到很近的地方.
你難以想像一個大城市.
為什麼會這樣.
你可以想像到.
偏遠的地方.
哪裡會有人來呢?.
你見到這些無助的.
是他們的痛苦.
還不懂得什麼叫痛苦的小朋友.
但是他們的內心.
原來是充滿了童真和快樂的.
這個落差是我們難以想像的.
任何一個住一學的兒童班.
都不會有這樣的快樂.
是充滿了完全好像沒有任何的煩惱.
壓力的一個小朋友的聚會.
我們想像得到.
你可以見到.
這個已經可以給我們影像的一張床.
其實是用一些很簡陋的布.
搭成的一張床.
就是全家人在那裡睡.
你看多麼的骯髒.
是多麼的難以想像.
每天的生活就在那裡.
那些父母就這樣坐在那裡.
沒有事情做.
沒有工作.
那些小朋友就到處走.
見到我們的車來了.
很開心.
因為他們宣教士已經建了一個地方.
只是一個上蓋.
沒有牆.
沒有什麼都沒有的.
只是一個上蓋.
就是來給小朋友聚會的地方.
就在那裡聚會.
見到我們的車來了.
知道有聚會了.

$^{321}$幾十個小朋友就自己這樣蜂擁.
跑來坐下準備來聚會.
唱聖詩和做手工遊戲.
然後分組查經.
一群小朋友就這麼開心地.
過了他們一個整個的上午.
親愛的兄弟姐妹.
我們見到的是什麼呢?.
我們見到的是不是很多的髒亂.
是的.
真的是骯髒的.
有其中的組員.
可不可以快點走.
很臭.
受不了.
真的很臭.
戴了口罩.
都受不了.
但是我們從小在火區長大.
就習慣了.
沒問題.
我戴了三個口罩.
希望是借給別人.
我自己沒有用.
因為我們都有心理準備.
是真的會很臭的.
但果然真的臭.
果然真的很貧窮.
貧窮到家徒四壁.
就是這樣.
但是我們有沒有這一份.
主耶穌見到的.
不是很多的問題.
是見到這些人其實如果福音.
去到他們的生命.
他們就會得到改變.
我們會不會有時見到的.
直接見到的是很多的問題.
勾起我們很多很多的控訴.
所以坐在我旁邊的弟兄姊妹就很慘.

$^{361}$不斷罵有沒有搞錯.
為什麼還是容許.
為什麼沒有公共的交通.
集體的運輸.
香港都已經不斷有改變.
輕鐵,地鐵什麼都有.
但馬尼拉是完全沒有的.
地鐵是完全不會有的.
一直都不想著會有的.
你很難想像.
那些人就是這樣去.
就是個個開電單車.
左穿右擦.
所以難免是會不斷塞車.
由早到晚塞車.
所以你在車裡的人.
是不斷控訴的.
但原來耶穌見到的.
不是問題的本身.
而是人的內心的需要.
所以宣教士在找一個.
很重要的機遇.
就是小朋友.
他們真的是一張.
完全未被污染的白紙.
如果改變了他們的命運.
改變了他們現在.
只是為生存掙扎.
就真的會有盼望有明天.
宣教士原來告訴我們.
他只是很希望將那個盼望.
放在他們的心裡.
所以第三位的女士.
那個中學生.
高中的了.
講她的見證的時候.
用流利的英文來講她的見證.
她說我進了這個.
CLC這個中心之後.
得到那個助學之後.

$^{401}$她說我現在真的對前面有盼望.
她敢去盼望前面的人生.
這句話就令到我們一班義工隊.
是很感動.
因為一個沒有可能有盼望的國家.
可以給到一點點的光.
這些小朋友.
親愛的兄弟姐妹.
這個好像是一個對我自己的.
一個靈魂的敲問.
就是這位小朋友.
不讓我走.
抱著我.
那我就抱起你.
很開心.
終於到處和別人拍照.
看到的是什麼呢?.
你看到這些這樣的爛地.
由入村.
整條村都是這樣的.
是不會有水泥地.
你可以想像得到.
就快到雨季了.
四月開始雨季了.
你可以想像得到.
那種的泥板.
就是這樣.
片山都是這樣.
現在剛剛還沒有到雨季之前.
所以很乾爽.
但是他們的家更加骯髒.
他們就是這樣長大了.
但是我們看到的.
到底是耶穌眼裡面的珍珠.
還是貧窮.
還是一代一代的惡性循環的貧窮.
無助和絕望呢?.
耶穌看到的就是憐憫帶來的.
愛的生命改變.
另外一段聖經.

$^{441}$大家應該都相當熟悉.
小兒子突然間有一天他醒悟.
我得罪了天.
又得罪了你.
你就是爸爸.
我不配做你的兒子.
你把我當作一個故宮吧.
但是他想好的台詞.
都沒有機會說的時候.
爸爸很遠就看見這個小兒子回來.
他就衝出去.
動了這個池心.
跑去.
所以抱著他的頸項親嘴.
根本小兒子沒有機會說這個台詞.
說你讓我做一個故宮就行了.
因為你家裡.
爸爸家裡有這麼多產業.
這就是小兒子的經歷.
親愛的弟兄姊妹.
愛可以帶來很大的生命改變.
小兒子根本沒有錢再穿新衣服.
洗乾淨自己的身體.
爸爸完全沒有對他有任何責備.
因為他已經內心知道.
他已經有很大的悔改,悔悟.
這就是主耶穌所說的.
稅吏和罪人靠近耶穌的時候.
耶穌向那些法利賽人展述.
為什麼神會接納罪人到祂面前呢?.
所以這是一個無條件的擁抱.
原來垃圾山的小朋友.
垃圾村的小朋友.
看到我們外國人來.
這些小朋友就很開心地衝過來.
就是等待我們的.
不是求我們給錢.
沒有一個小朋友要問我們拿錢,拿糖,拿玩具.
什麼都沒有.
就是準備好坐下來等著聽故事.

$^{481}$我們很難以想像.
因為其實我們真的每個人都穿得很漂亮.
我們可以給錢他們.
但宣教士說不需要.
因為他們根本不稀罕這些東西.
但是沒有可能的.
窮的人當然希望有外國人來.
當然要問他們拿錢.
小朋友就是一個很大的,很好的機會.
由小朋友開口問拿錢.
但是完全沒有.
我們覺得很奇怪.
我們只是為他們預備了一桶朱古力粥.
聽完故事之後.
每個小朋友就快點跑回家拿個碗.
每人拿一個碗來.
就是裝著那碗粥.
然後就坐下來慢慢吃.
就是這麼多.
他的要求就是那碗粥.
但更多的是.
很開心見到我們和他聊天.
和他抱著拍照.
他們很快就會做那個手勢.
個個就已經準備好.
就是這麼多.
他們很喜歡,需要的.
其實只不過是愛的擁抱.
親愛的弟兄姊妹.
今天很想在這個講道裡.
來鼓勵大家.
我自己的那個改變.
我自己對菲律賓人的那個誤解.
那個偏見.
那個敵視.
很多的事件.
有旅遊社的那個年輕的導遊.
因為藏毒.
其實是被人陷害.
終身監禁.

$^{521}$試過又有槍擊事件.
人質的事件.
令到我們香港人對菲律賓.
是有很大的那種反感的.
相信在座.
我不用舉手問.
大家不會想.
今年我們希望去菲律賓旅遊.
應該沒有人這樣做的.
做的人都很傻.
除非孫教士.
那個孫教士就是.
好像他很傻那樣.
來到他在菲律賓已經七年.
他賣了四層樓在中國.
來去做這個兒童之家.
幫助這些窮苦的兒童.
所以有幾個中學生.
是看著他們長大.
你可以想像得到.
這種無私的愛.
其實就是一個愛的接納和擁抱.
弟兄姊妹.
盼望這段的經文.
盼望這個的講道.
是和大家一個彼此的勉勵.
我們是要怎樣去做這件事呢?.
我們怎樣能夠讓我們的進入.
成為他們可以追夢的一個開始呢?.
聽孫教士講.
他們已經計算了所有的行政費用.
一個小朋友.
無論小學中學.
只要你讓他入學.
入中學或者入小學.
基督教的.
我們參觀過那個學校.
那個校長是一個很好的基督徒.
一年的那個需要的生活費.
只不過是二千塊人民幣.

$^{561}$一個人一年.
因為原來政府的讀書是不用錢的.
但是書簿費,生活費.
就麻煩你自己搞定.
你搞不定.
那你就沒有書讀了.
政府是不會理你的.
是不會有強制的入學的.
就是說你給不到學生的生活費.
書簿費.
家裡人沒有錢.
何來拿書簿費去給小朋友讀書呢?.
所以大部分的小朋友.
在垃圾村是不能讀書的.
因為只是沒有書簿費.
一年二千塊人民幣.
你可以想像得到.
那個需要是多麼微不足道.
我們香港人能夠想像到的.
有姐妹已經說.
二千塊?.
你問清楚了沒有?.
曾牧師.
我說不是.
我確定了.
孫教授說.
一年二千塊人民幣.
一個人.
他馬上說.
那我支持十個.
可以改變十個兒童可以有書讀.
送他們進基督教學校.
在學校裡面.
他們就很開心地有機會讀書.
有很多自然團體.
奉獻很多的電腦.
給他們來學習.
原來韓國的孫教.
十幾二十年前已經很大力地做.
韓國的教會.

$^{601}$已經很多年前在菲律賓.
做很多孫教工作.
我們住的孫教士的宿舍.
就是韓國人來開的.
有他們的營地等等.
中國的教會.
香港就更加少.
來到當地做福音工作.
還是很後起.
最近都多了.
是內地過去的.
親愛的兄弟姐妹.
很想和大家一起.
彼此來思想.
或者作為一個反思.
我們要支持孫教的士工.
讓更加多的孩子.
有機會讀書.
他們的生命有改變.
他們認識主耶穌之後.
他們就能夠開始.
隨著他們的成長.
他們成為一個敢於.
思想和夢想的人生.
願主祝福大家.
我們一起有一個簡單的.
我們默想這一段經文.
耶穌走遍各城各鄉.
見到的人他憐憫他們.
然後我帶領大家一起有一個祈禱.
(祈禱).
親愛的主.
段孫隊和錦孫.
和其他教會的弟兄姊妹.
都有差不多十個人.
我們去到菲律賓馬尼拉.
一個大城市裡面的垃圾村.
裡面的小朋友那種天真.
那種流露出來的喜悅.
平安.

$^{641}$是我們難以想像的落差.
因為我們認為物質的貧窮.
一定會帶來心靈的.
面容的.
內心的那種痛苦.
但是我們見不到這些的痛苦.
在他們的內心.
在他們的面容.
告訴我們原來就算貧窮.
他們都是這麼喜樂.
但是我們更加擔心的就是.
隨著時間的長大.
他們一直升上去.
悉靈的學習的時候.
而失學.
而來到街頭流連.
那種無助.
那種學壞.
就成為真正的社會問題.
所以我們真的很盼望.
既然主感動教會.
我們的弟兄姊妹.
有這個野地牧人的事工.
來支持當地的宣教組織.
和當地的信徒領袖.
菲律賓的信徒.
來在不同的地方傳福音.
求主給我們.
都一起有這個聖靈的感動.
我們大家都來到去.
一起在主面前思想.
我們是不是都需要.
負起我們洪恩家.
錦宣家的行動.
我們一起為這些垃圾村的村民.
這些小朋友.
來去改變他們的命運.
讓他們能夠有.
接受基督教的教育和學習.
使他們的生命成長的時候.

$^{681}$隨著他們的年紀.
他們的學識.
他們的學問.
都一起長大的時候.
將來他們能夠改變他們的命運.
求主幫助.
我們很盼望今早的講道.
成為一個彼此分享.
來到去一起.
來到去思考.
應該怎樣來到去按照神.
在我們大家的心靈裡面的感動.
來到去行事.
求主幫助.
我們這樣禱告.
奉耶穌基督的命.
阿們.
\newpage



\section{約翰福音 13:1-17}
\label{sec:Jj4ZioVnCDo}
\textbf{2026.03.15 《做騷的耶穌》:約翰福音13\_1-17 (和修版) 張雲開老師}
\newline
\newline
連結: \href{https://youtube.com/watch?v=Jj4ZioVnCDo}{\texttt{https://youtube.com/watch?v=Jj4ZioVnCDo}} ~~~~ 語音日期: 2026-03-18
\newline
\newline
\hyperref[sec:S89_Yf51Xes]{\small{< < < PREV SERMON < < <}}
~
\hyperref[sec:index_chronic]{\small{[返順時目]}}
~
\hyperref[sec:index_scriptual]{\small{[返順卷目]}}
~
\hyperref[sec:zmeyo_PNxjc]{\small{> > > NEXT SERMON > > >}}
\newline
\newline
約翰福音 13:1-17
\newline
\begin{longtable}{cl}
\hline
\hline
章節 & 經文 (和合本修訂版)\\
\hline
13:1 & \begin{tabularx}{0.7\textwidth}{X} 逾越節以前，耶穌知道自己離世歸父的時候到了。他一向愛世間屬自己的人，就愛他們到底。 \end{tabularx} \\ \\ \relax
13:2 & \begin{tabularx}{0.7\textwidth}{X} 晚餐的時候，魔鬼已把出賣耶穌的意思放在加略人西門的兒子猶大心裡。 \end{tabularx} \\ \\ \relax
13:3 & \begin{tabularx}{0.7\textwidth}{X} 耶穌知道父已把萬有交在他手裡，且知道自己是從神出來的，又要回到神那裡去， \end{tabularx} \\ \\ \relax
13:4 & \begin{tabularx}{0.7\textwidth}{X} 就離席站起來，脫了衣服，拿一條手巾束腰， \end{tabularx} \\ \\ \relax
13:5 & \begin{tabularx}{0.7\textwidth}{X} 隨後把水倒在盆裡，開始洗門徒的腳，並用束腰的手巾擦乾。 \end{tabularx} \\ \\ \relax
13:6 & \begin{tabularx}{0.7\textwidth}{X} 到了西門‧彼得跟前，彼得對他說：「主啊，你洗我的腳嗎？」 \end{tabularx} \\ \\ \relax
13:7 & \begin{tabularx}{0.7\textwidth}{X} 耶穌回答他說：「我所做的，你現在不知道，但以後會明白。」 \end{tabularx} \\ \\ \relax
13:8 & \begin{tabularx}{0.7\textwidth}{X} 彼得對他說：「你絕對不可以洗我的腳！」耶穌回答他：「我若不洗你，你就與我無份了。」 \end{tabularx} \\ \\ \relax
13:9 & \begin{tabularx}{0.7\textwidth}{X} 西門‧彼得對他說：「主啊，不僅是我的腳，連手和頭也要洗！」 \end{tabularx} \\ \\ \relax
13:10 & \begin{tabularx}{0.7\textwidth}{X} 耶穌對他說：「凡洗過澡的人不需要再洗，只要把腳一洗，全身就乾淨了。你們是乾淨的，然而不都是乾淨的。」 \end{tabularx} \\ \\ \relax
13:11 & \begin{tabularx}{0.7\textwidth}{X} 耶穌已知道要出賣他的是誰，因此說「你們不都是乾淨的」。 \end{tabularx} \\ \\ \relax
13:12 & \begin{tabularx}{0.7\textwidth}{X} 耶穌洗完了他們的腳，就穿上衣服，又坐下，對他們說：「我為你們所做的，你們明白嗎？ \end{tabularx} \\ \\ \relax
13:13 & \begin{tabularx}{0.7\textwidth}{X} 你們稱呼我老師，稱呼我主，你們說的不錯，我本來就是。 \end{tabularx} \\ \\ \relax
13:14 & \begin{tabularx}{0.7\textwidth}{X} 我是你們的主，你們的老師，尚且洗你們的腳，你們也應當彼此洗腳。 \end{tabularx} \\ \\ \relax
13:15 & \begin{tabularx}{0.7\textwidth}{X} 我給你們作了榜樣，為要你們照著我為你們所做的去做。 \end{tabularx} \\ \\ \relax
13:16 & \begin{tabularx}{0.7\textwidth}{X} 我實實在在地告訴你們，僕人不大於主人；奉差的人也不大於差他的人。 \end{tabularx} \\ \\ \relax
13:17 & \begin{tabularx}{0.7\textwidth}{X} 你們既知道這些事，若是去實行就有福了。 \end{tabularx} \\ \\
[1ex]
\hline
\hline
\end{longtable}
$^{1}$Thank you 早晨.
多謝教會邀請我再一次來到洪水橋跟大家分享上帝的話語.
我們都已經進入主誕在十二月過後.
其實很快就進入所謂的御苦期.
如果我們看一年主耶穌基督.
他在福音書裡面記載他所做的所經歷的一些重要的事件.
一路排列我們的教會的年齡.
再過幾多個星期就是復活節.
壽難復活的主日.
再過幾多個星期知道嗎?.
三個星期.
所以兩個星期之後就是中樹節.
就是主騎驢入耶路撒冷.
那一個星期就是主最後一個星期.
在耶路撒冷最終在星期五被出賣.
被審訊被敲打.
被釘十字架.
所有的事情都會在兩個星期之後.
在兩個星期之後到三個星期之間發生.
今天跟大家所分享的主席讀的經文.
就是在最後一個星期發生的.
按道理應該說在前面的事情.
因為畢竟現在還是三個星期之前.
應該說可能是在蓋薩利亞,菲勒比.
在北面向南走的時候主問門徒.
你們說我是誰?.
別人說我是誰?.
你們又說我是誰?.
那段經文.
不過我還是選擇這段在最後一個星期.
主跟門徒吃御節晚餐.
事實上就是星期五晚上的時候了.
但是有幾件事情我們要先講清楚的.
如果我們看.
我們可能都需要PowerPoint的弟兄姊妹.
跟我們打回經文出來.
如果我們看第十五,第十六節.
第十五,第十六節我用舊和學本讀的.
剛才大家讀的是和修版.
新一點的.

$^{41}$我還是老一代的舊和學本.
但是意思都是一樣的.
第十五節他說我給你們做了榜樣.
是叫你們照著我向你們所做的去做.
耶穌多沉悶.
叫你們照著我所做的去做.
意思就是說主耶穌在這段故事.
這段記載裡面向門徒做了些什麼?.
洗腳.
如果主說我給你們做了榜樣.
叫你們照著我向你們所做的去做.
就是他在吩咐門徒做什麼?.
彼此洗腳.
第十六節我實實在在的告訴你們.
僕人不能夠大於主人.
差人的也不能夠大於差他的人.
你們既然知道這件事.
如果去做就有福了.
但是第十六節很奇怪的.
第十六節又說僕人不能夠高於主人.
所謂差人就是被差的人.
不能夠大於差他出去做事的那個人.
一方面主耶穌就說我洗你們的腳.
是一個服侍人的舉動.
另一個就是說我其實比你們大.
因為他吩咐門徒.
你要照著我向你們所做的去做.
所以到底耶穌是比門徒大.
還是比門徒卑微?.
很難說.
其實兩邊都說得通的.
起碼他是願意卑微服侍門徒.
做連門徒都不會彼此做起對方身上的事情.
由於如此我們看完之後.
我先給幾個聲稱.
第一就是耶穌和門徒洗腳的行動.
只是一個象徵而已.
什麼叫象徵?什麼叫符號?.
我們馬路上有個路標指向這裡.
洪水橋就在這裡.

$^{81}$洪水橋站有個牌子告訴你.
那個牌子不是洪水橋的本身.
這個地方才是洪水橋.
那個牌子告訴你站在這個地方就是洪水橋.
所以那個牌子是一個象徵.
是一個符號.
就等於紅綠燈是一個符號.
紅色就是紅色.
為什麼紅色不要動?.
紅色和不要動是沒有關係的.
你不會看到所有東西紅色就立刻站在這裡.
所以紅色和不要動是沒有必然的關係.
但是我們約定俗成.
就說紅色是一個記號.
是一個符號.
是一個象徵符號.
如果在路邊有三個圓圈是盞燈的話.
紅色你就不要動了.
所以符號有意義是指向某些現實.
所以在這裡第一件事就是說.
耶穌洗腳的行動只是指向現實.
不是現實的本身.
我們知道耶穌替人洗腳替門徒洗腳.
只不過是象徵的行動.
為什麼我們會這樣說呢.
就因為這個並不是主耶穌服侍門徒的常態.
福音書裡面記載.
是主耶穌為期幾年的公開服侍.
自從祂焦聚了門徒那十二個才算了.
還有七十個還有五百個.
有多少次是祂和門徒洗腳.
按照福音書裡面記載.
一千零一次.
其實是一次不是一千零一次.
是一次.
所以不是常態.
如果不是常態.
突然間耶穌做了一件不是常態的事.
我們就要想清楚.
接著耶穌就說.

$^{121}$我向你們所做的你們要照著去做.
你就要想清想錯了.
到底是要去做還是.
還是祂在說其他東西呢.
這個是要思考的.
我在這裡先說.
其實不是需要照著字面.
祂去和人洗腳.
我們就和人洗腳的.
第二.
就是第二個原因.
我們不需要按照耶穌的具體行動去做.
原因就是早期耶路撒冷教會.
包括新約裡面.
我們所認識的第一代的外邦人教會.
有沒有按照著主耶穌這個榜樣去做呢.
換句話說.
早期教會有沒有一個洗腳.
互相洗腳的文化.
早期教會我們怎麼知道呢.
看《士徒人傳》.
或者看《保羅書信》.
有沒有這種.
保羅有沒有教導.
愛是恆久忍耐.
彼此洗腳.
有沒有一個洗腳禮儀.
也沒有.
所以我們又多一個原因.
就是說其實主耶穌的吩咐.
我向你們所做的.
我們要按照著去做.
就不是叫我們具體的彼此洗腳.
還有第三個原因.
第三個原因可能有些信徒覺得不是一個很大的原因.
嚴格來說.
主耶穌當天替門徒洗腳.
這種行為.
當時候的意義.
今天已經不存在了.

$^{161}$當時候洗腳有某種意義.
今天你洗腳.
別人以為你是做足浴.
就沒有了以前那種意義.
一個行動的意義.
有它自己的歷史文化場景.
是一個約定俗成的東西.
當天主耶穌對門徒做這樣的事的意義.
和今天我們同樣的做法所產生的意義.
會很不一樣.
主耶穌洗門徒的腳.
那時候是在扮演一個什麼角色.
奴僕的角色.
奴僕做什麼的.
奴隸以前做什麼的.
服侍他的主人.
包括是按照著他的主人的意願.
來服侍他.
對猶太人來說.
連學生徒弟都不會洗.
都不需要洗老師的腳.
洗腳的工作被視為低下到一個地步.
連猶太人當中.
都是讓非猶太裔的奴僕去做的.
所以耶穌以一個老師的身份.
做出這種行動.
自然就有一個震驚的價值.
讓門徒震驚.
被突然間這樣做.
不知道他在做什麼.
突然間打大赤膊.
是不是要打人呢.
到底要做什麼.
紮起毛巾.
紮起腰.
然後就拿一盆水.
開始和門徒洗腳.
嚇死他們.
因為他們的師父.
他們的拉比.

$^{201}$成為了他們的奴僕.
我們現在說僕人.
都高了一格.
我們僕人就是想外傭.
今天是星期日.
全香港在室外.
都有很多外傭.
他們想離開自己工作的環境.
一方面.
我們就會想.
外傭.
傭工.
但是奴僕就低很多.
奴僕是賣身給主人的.
他在羅馬世界裡面.
在羅馬的法例裡面.
奴隸不是人來的.
沒有一般人享有的權利.
他是主人的財物.
物件.
差很遠.
所以當你想.
讀新約.
耶穌基督的僕人.
耶穌基督的奴僕.
保羅寫信給哥林多教會.
你一定要想.
他在說羅馬世界的時候.
大家一看.
都明白他和主耶穌的關係.
他就是主耶穌的奴隸.
他的命是屬於主耶穌的.
在主耶穌面前.
他連東西都不是.
人不是人.
是一件東西來的.
如果你是嚴格地.
按照羅馬的法律去看的話.
所以對他們來說.
耶穌所做的做法.

$^{241}$突然間由師父變成了.
不可能的.
變成了他們對師父的那種關係.
還低一格的奴僕.
但是他們當然知道.
耶穌作為師父不是他們的奴隸.
他們當然知道.
那十二個.
所以他們就懂得去詮釋這件事.
我們的師父放下身段.
卑微服侍我們.
這個是當時候的意義.
所以耶穌吩咐門徒彼此洗腳.
要吩咐他們放下身段.
彼此的服侍.
倒過來我們今天的社會文化裡面.
洗腳就已經沒有這個意義了.
洗腳剛才說過.
可能是某種的工種.
某種職業的工作範圍.
也因為這樣.
洗腳的服侍意義.
已經失去了.
你說不是.
我們那些足浴.
那些都是服務性行業.
那也是.
香港是主打服務性行業.
除了金融.
地產之外.
服務性行業.
但是每個人都做服務性行業.
他收你的錢.
你不給錢誰服務你.
當然沒有人服務你.
所以也不行的.
不收錢的.
而且收了錢就名聲言順.
他要滿足你的要求.
在這個位置上.

$^{281}$主耶穌有沒有問.
跟門徒洗腳的時候有沒有問.
你想不想把甲收起來.
你想用冷水洗還是熱水洗.
有沒有這樣來問門徒.
當然沒有.
因為他不是在服務他們.
所以.
三個原因是告訴我們.
今天我們需不需要.
做彼此洗腳的事情.
不需要.
誰說需要呢.
他是騙你的.
耶穌也是這樣說的.
第十五節說.
我給你們作的榜樣.
你們照著.
我向你們所作的去作.
他是不明白聖經.
他是斷章取義.
不是斷章取義.
字面的意思的確是這樣.
不是斷章.
的確是這樣的意思.
但是其實他誤解了意思.
不是字面的意思.
所以我們不必彼此洗腳.
基督教有沒有洗腳的一些做法呢.
其實有的.
一年一次.
有些地方還會這樣做.
如果是犯基督宗教的傳統.
很明顯是天主教這樣做的.
教皇一年一次.
在壽難節期間.
也是在禮拜五晚上.
替一些事先安排好的人洗腳.
作為一個服侍的象徵.
那些人的腳.

$^{321}$你猜乾不乾淨.
很乾淨.
肯定沒有味道.
教皇每年做一次的洗腳.
這個就是洗腳.
我們的主也在他被賣的前夕.
為他的門徒和我們後來讀這個故事的人.
也做了一場洗腳.
所以題目叫做洗腳的耶穌.
他洗腳而已.
不是叫你照足百分之一百.
表面這樣去做.
所以當我們問.
可能大家都聽過.
這是二三十年前很流行的一個說法.
就是我們要遵循上帝的旨意去做事.
我們就要問一個問題.
在英文是WWJD.
What would Jesus do?.
耶穌會怎麼做?.
當我們在排隊.
見到有個我們叫大肚婆.
孕婦.
大肚婆好像粗俗一點.
孕婦的時候排在後面.
你心裡突然間有個想法.
是不是需要繼續照常排隊.
什麼都不做呢?.
你就會問.
耶穌會怎麼做?.
What would Jesus do?.
你覺得耶穌會這樣做.
那你照著他這樣做吧.
這樣的態度其實也是很值得欣賞的.
但我們需要很小心.
耶穌可能在某個場合.
以替門徒洗腳來傳遞一個重要的訊息.
但我們就不是東施效顰.
就照著耶穌的表面洗腳的做法去做.
機械化來模仿.

$^{361}$無視了耶穌這樣做背後的意義.
反過來我們應該問另外一個問題.
不是問耶穌會怎麼做.
而是問耶穌會想我們怎麼做.
耶穌會怎麼做和耶穌會想我們怎麼做是兩回事.
剛才我們福曲的詩班裡.
有姐妹一下子掉了一個杯子.
那個不知道是什麼東西.
一彈彈過去.
我也在想我要不要過去撿.
我在想耶穌會怎麼做.
還是耶穌想要我怎麼做.
我在想會不會有人過去撿.
我先看看誰過去撿.
很多事情都在想.
其實我比較近.
結果主席過去撿.
當然看到.
突然間.
現在要很警醒.
砰一聲可能是有什麼炸彈.
什麼事.
不知道.
所以突然間有個聲音一定很警醒.
還有東西到處飛.
主席就去做.
他沒有想.
我猜而已.
我自己當場馬上閃過.
因為今天講到.
馬上閃過幾個選擇.
耶穌會怎麼做.
耶穌會想我怎麼做.
會不會有人其他去做.
我想等其他人做好.
我不要做了.
一大堆的問題.
所以我就覺得與其問耶穌會怎麼做.
你就應該問耶穌會想我怎麼做.
這個需要一些想像.

$^{401}$因為你不是耶穌.
你要想像一下主會想我怎麼做.
這個我覺得總是比較有效.
因為你是透過全面了解耶穌的行事為人.
你才能回答耶穌會想我怎麼做.
祂的做人宗旨是怎麼樣.
祂會想我也有類似的做人宗旨.
所以就主替門徒洗腳這件事來說.
耶穌給門徒的榜樣是象徵著一種.
願意放下身段去服侍人的卑微的態度.
你發覺我不講謙卑的態度.
一會兒會解釋為什麼我不用謙卑.
用卑微.
可能在另外一個場合.
耶穌在平息門徒之間的紛爭.
他們爭論誰為大的過程裡面.
在另外一個場合主耶穌就指出.
人子尼不是受人的服侍.
是要服侍人.
而祂也說你們中間誰想做大衛首的.
就必作你們自己的僕人.
這是馬太福音20章.
路加福音22章都有說.
我們明顯可以看到.
其實主教導門徒在不同的場合裡面.
一直都教門徒.
門徒總是想爭大.
12個要排位.
我讀小學的中學也會這樣.
起碼初中也會這樣.
因為我們上課之前要在操場排隊.
日曬雨淋.
雨淋就不會.
曬到燒焦.
夏天曬到燒焦.
曬到頭髮都掉了.
應該是那時候造成的.
排隊是怎樣排的.
就是矮的先排.
因為高的排前面.

$^{441}$好像我剛才站在前面遮住了.
你遮住我看螢幕.
所以排隊有方法排的.
由矮至高.
我們讀中小學的時候也有學號.
按照著英文姓氏的ABCD排.
我姓張的.
CH字頭.
C就永遠排在前面.
我沒試過是十號之後.
七號五號都試過.
十號都好像試過.
總是有個方法排.
這些方法大家不用爭.
你出生是怎樣就怎樣.
你現在的樣子是怎樣.
客觀標準.
而且有道理.
你站在前面不要遮住後面.
這些道理為了幫後面的人方便.
所以大家都信服.
所以大家不會爭的.
但如果沒有這些規矩排.
那不是先後嗎.
你大一點.
為什麼那十二個經常爭大小.
爭大小有它的政治原因.
因為他們是跟著主耶穌.
譬如大家在.
我這個不是正確的.
應該是馬可福音.
都記載過爭大小的時候.
是約翰和雅各.
雅各約翰這兩兄弟.
媽媽都出聲.
媽媽就跟主說.
主啊你將來得國的時候.
得到國道真是有權.
我都想著主耶穌要搞革命.
趕走羅馬人.

$^{481}$復興以色列國.
讓我兩個兒子.
將來一個坐你左邊.
一個坐你右邊.
做丞相.
左右帶丞相.
會爭權的.
真的會爭權的.
在盧家坊第九章.
盧家坊記載了兩次.
那些門徒爭權.
第一次第九章.
第二次是二十二章.
二十二章正正就是這個場合.
約翰沒有記載的.
盧家有記載.
所以要卑微服事.
不要好真偉大.
這個就是主耶穌.
洗腳行動的合理詮釋.
也是我們回答.
主耶穌想我做什麼的時候.
恆常需要考慮的一個基礎答案.
就是我們不要好真偉大.
是願意卑微服事.
為什麼我不說謙卑服事.
現在就應該提了.
因為謙卑謙卑.
中文不是很好.
聖經常常翻譯做謙卑.
謙卑不好的原因就是多了個謙字.
謙是翻譯錯的.
在聖經起立就翻譯錯了.
奴隸要不要謙.
奴隸不用謙的.
他本來就是卑.
他是地底泥.
在法律面前不是人.
所以你用謙.
謙是什麼意思.

$^{521}$謙就是我不行的.
你去做吧.
你彈琴比我好.
我五級你八級.
你做吧.
這叫謙.
謙就是其實我可以的.
不過我讓給你.
孔庸讓你.
這就叫謙.
謙謙君子.
就是不爭.
跟不爭是有關係的.
他可以爭的.
不過他放棄不跟你爭.
這叫謙.
但主耶穌還多走一步.
極端很多.
你謙有時候是做謙謙君子.
有時候是偽君子.
你心裡面也有不同的人爭.
這裡我不跟你爭.
對我沒什麼價值.
轉個彎.
起另外一些事情.
我跟你爭到癢.
往往人都是這樣.
這些吃虧的東西.
我當然不跟你爭.
其他東西有好處的.
我就去爭.
但爭又不明目的爭.
又暗地裡搞一大輪手腳.
背後做很多事就去爭.
明爭就暗地裡鬥.
所以謙就不好.
有些人太謙.
我們反而覺得他好像很虛偽.
我不行的.
其實是推三推四.

$^{561}$我彈琴.
他比較厲害.
他彈吧.
我不想彈了.
不要教我.
煩著我.
我們說小孩子最厲害.
他做吧.
因為他發覺沒什麼爵士.
我讀裕科的時候.
裕科我的同學.
裕科的時候.
我入了一間很好的中學讀裕科.
我的同學全部比我厲害.
你知道以前A級裕科只有兩年.
一入去就要馬上做試題.
預備兩年之後考試.
我有些不懂得做.
我就問.
大概讀了一年之後.
哪個同學厲害.
哪個同學沒那麼厲害.
都比我厲害.
全部都知道.
哪個同學比較友善.
哪個同學沒那麼友善.
都知道.
大家相處一班二十幾個人.
我問其中一個同學.
你可不可以教我做這個數.
這個數搞不定.
怎麼做呢.
他就說我不懂的.
其實他後來是A級.
他是A級四條A.
考五科四條A.
打橫入港大那一種.
他不懂.
對我來說.
我吃了閉門根.

$^{601}$明明你懂的.
只是幫我.
我不懂的.
好像很謙虛.
其實是你不想幫.
所以謙虛不是一個好的翻譯.
記住.
謙虛是很多東西都可以.
很謙虛.
其實是幾格下的.
可以幾格下的.
其實是卑.
你要拿個卑字不要謙.
謙卑.
如果我們講謙卑核心意義.
其實是卑微的部份.
是願意服侍.
社會的常態.
少的服侍大的.
做小的當然服侍大的.
服侍老細.
老細就不是很小.
老細其實是老大.
我從來都不知道為什麼叫老細.
你知不知道為什麼叫老細.
為什麼加個小字.
老細扮謙卑.
我是老細.
其實他是大.
所以卑微一詞就準確一點.
一般都是用起.
一貫古時都是用起身份低下的人身上.
是貶詞.
不是褒詞.
褒貶褒貶.
不是褒詞是貶詞.
耶穌用卑微.
用謙卑來形容自己.
我心裡面柔和謙卑.
我額是輕的.

$^{641}$擔子是輕的.
額是容易負的.
新約亦都將謙卑納入信徒美德當中.
羅馬書法所書很多地方都講我們要謙卑.
這亦都是主體門徒洗腳的核心意義.
但其實是卑微.
古時在羅馬世界裡面.
甚麼人是必須服侍人.
沒有甚麼會接受人的服侍.
就是奴隸.
所以是社會裡面最低層.
奴隸是不能謙的.
主人叫你做甚麼你就要做得到.
你不可以跟主人說.
因為有阿一阿二阿三阿四.
有四個奴隸.
主人說你跟我去甚麼甚麼.
你不可以跟主人說.
我不行的.
你找阿三去.
你找阿四去.
沒有的.
主人叫你做甚麼.
不能謙的.
他比我厲害.
他做得比我好.
主人說了算.
所以是聖經裡面.
尤其是肥納比書第二章.
講到謙卑.
用卑微來形容主耶穌選擇.
來到世界上的行動.
就是心得其要領.
他是降卑為人.
放下自己的身段.
不以自己和上帝同等為強奪的.
他本來是上帝的獨身子.
是有和上帝一樣的能力.
是本性.
但是他放下一切.

$^{681}$成為人的樣式.
所以他自己本身就做了這件事.
而在最後御節的晚餐裡面.
他就用另外一個方式表現出.
好像奴隸.
其實他不是奴隸.
他是師父.
但是他最終都是用師父的口氣.
吩咐門徒.
你要照著我向你們所做的去做.
不是彼此洗腳.
他們後來都沒有彼此洗腳.
而是彼此的服侍.
將自己視為卑微.
是各人看對方比自己強.
保羅在肥納比書第二章裡面.
他就看自己比自己強.
我剛才進來的時候.
是和主席和歐sir握手.
和主席和歐sir握手.
看到他正在賺爆西裝.
我做侵略性的行動.
我沒有徵求他的同意.
我捏他的手臂.
不是很熟悉.
你捏人做什麼.
你不應該這樣做.
真的有實力.
但是我和歐sir熟悉一點.
我又捏他的手臂.
他和歐sir說你要努力一點.
如果他們又捏我的手臂.
我就太瘦了.
肌肉流失.
各人看別人比自己強.
主席肯定強過我.
歐sir的肌肉都比我多.
都比我強.
什麼意思.
就是你比我強.

$^{721}$我很羨慕你.
我們現在會這樣說.
沒有什麼特別關係.
你比我強.
但是在古時.
看別人比自己強.
不是說你比我強.
那你去做吧.
我不用做了.
不是.
而是對方是強人.
所以你就服侍他.
你作為奴隸.
在羅馬社會最弱的就是奴隸.
奴隸的特徵就是要服侍人.
就是這樣.
所以當普羅在肥納比書那裡.
各人都看別人比自己強.
各人不要單顧自己的事.
也要顧別人的事.
不是叫你去八卦.
是你要顧及別人的事情.
掉了一顆東西出來.
你幫人撿起來.
他可能不知道.
可以撿起來.
剛才姐妹說還可以修理.
用一支大陸針弄一下就可以了.
不需要再買那個.
你要服侍人.
你幫人.
你視對方是你服侍的對象.
即是他比你大你才服侍他.
所以在這些位置.
連爸爸媽媽都會扮小.
爸爸媽媽對他生出來的嬰兒.
面乾碎濕.
換尿片.
他就好像變成一個小孩子.
他在服侍那個嬰兒.

$^{761}$他視那個嬰兒是大人.
如珠如寶.
你碰到嬰兒都有罪.
他保護他.
他的地位比他高.
如果嬰兒的地位不高.
他跟他換尿片.
你去換.
我就不換了.
是很難的一件事情.
有一次.
我自己以前的關心小組.
他們畢業在一間教會的幕會.
結了婚.
那間教會就借了我們的禮堂做營會.
他進來的時候生了兩個小孩.
大的就自己勉強行到.
小的還在抱.
進來又看看我.
看到嬰兒很開心.
突然間在聚會的中間.
那位太太 媽媽就說.
張Sir 嬰兒嗡嗡.
禮堂嗡嗡 搞不定.
我宿舍就在禮堂旁邊.
過來在我宿舍那裡搞定他.
大人宿舍.
有張擺桌的桌面.
吃飯桌面.
一解開.
我馬上 我先走了.
你搞定.
我說嬰兒太好吃了.
影片那裡.
旁邊已經嗶了出來.
你見過嗶了出來沒有.
真的搞不定.
你搞定他.
要洗什麼.
你去廁所你自己搞定.

$^{801}$沒問題 我先走了.
不是我嗶 我不理.
你要我服侍嗶.
媽媽 爸爸願意服侍嗶.
這種就是大的服侍小的.
大的願意服侍小.
最典型的例子.
爸爸媽媽是大過嗶.
一定大過嗶.
身份 地位 經驗 能力.
什麼都大過嗶.
偏偏大的就願意服侍小的.
旁人沒有關係就不服侍.
所以這個服侍也是一種愛的服侍.
所以我們剛才讀第十三章.
第一第二節那裡.
預節以前耶穌知道自己離世歸父的日子到了.
祂既然愛世間屬自己人.
就愛他們到底.
是愛的一種表現.
不要想著你是很屈就.
現在很勉為其難.
是你視人比自己強.
其實是愛對方.
是幫助對方.
願意服侍對方的表現.
所以.
這幾個聲稱完了.
我就繼續講另外幾點.
第一點已經講了.
耶穌這樣做是出於.
對門徒的愛.
愛是彼此服侍的基礎.
如果不是.
你憑什麼叫我這樣做.
我為什麼要這樣做.
你問很多問題.
一定是出於對對方的關懷.
對對方的愛.
正如主愛我們一樣.

$^{841}$第二主耶穌這樣做.
有另外一個原因.
也是第二節到第五節.
那裡主耶穌所說.
剛才我讀的.
因為時間少了.
沒什麼時間.
祂知道自己在地上的時間沒什麼.
其實就是那一晚.
那一晚完了.
祂就會和親自挑選的十二位使徒.
一個出賣他.
其中一個.
就會分開.
再晚一點.
再過幾個小時就永遠分開.
復活之後當然又再見面.
所以在這裡祂也出了絕招.
祂看到那十二位使徒冥頑不寧.
我剛才說約翰沒有記載.
路加有記載.
你看22章正是最後.
還在爭大小.
路加沒有記載這段洗腳的事件.
有記載吵架.
爭大小事件.
主耶穌有教訓他們.
約翰有記載洗腳事件.
但沒有記載吵架事件.
很奇怪.
這是同一晚發生.
而且性質也是一樣.
應該是主耶穌罵完他們之後.
然後就做事了.
所以在這些位置.
主耶穌是申明時間沒什麼.
所以祂要震撼門徒.
這些你永遠記得.
因為平時不做的.
突然間做一次這樣的事.

$^{881}$門徒會一輩子記得.
嚇死我們.
你看見彼得的反應.
我剛才讀彼得的反應.
嚇死他.
他好像明白又好像不明白.
他以為自己明白.
洗了我全身.
他以為洗那件事很重要.
其實不是洗這件事.
是洗背後的異議.
所以在這些位置.
門徒還沒掌握到.
但主耶穌這樣做.
起碼他們一定會一輩子記得.
開始能夠實踐.
這是第二點.
所以一方面出於祂的愛.
另一方面出於時間少.
第三.
這個也是很特別的.
服事就是幫人.
等人不用做.
等人有空.
等人漏了的東西.
你幫他補償.
始終是幫人.
但主耶穌的服事.
是很特別的.
在第21節.
服事還有潔淨的功能.
我們說洗.
主耶穌又揀選了洗這件事.
主耶穌不揀選.
我幫你擺桌子.
我勞碌擺桌子.
他之前可以擺桌子.
又或者完結了.
收桌子.
也是奴隸做的.

$^{921}$不做這些不做那些.
也是小的服事大.
他用洗.
做這件事還有多一層的象徵意義.
就是他的服事.
也有潔淨的作用.
我們服事人.
在主面前服事的確.
是有潔淨的功能.
如果我們從大的服事.
小的角度來看.
這一點可能更加清楚.
譬如說以老師服事.
學生為例.
老師的成熟經驗智慧.
往往能夠幫助學生.
在學習上甚至.
在生命上避免.
一些他們.
看不到的一些陷阱.
導學生於善.
一個服事.
學生的老師不單只是.
教育工作者.
而是一個生命工作者.
他的教育工作.
有潔淨的作用.
他在引領學生走一條路.
不只是幫他考A級.
以前教會考.
現在全部都DSE.
或者IB.
不只是幫他考好試.
而是幫他.
做一個人.
幫他預備.
將來進入社會.
能夠順利一點.
就不單純是讀書厲害.
讀書不厲害的問題.

$^{961}$所以我發覺有很多老師.
有些老師的確是.
你考試最厲害的那幾個.
他特別喜歡的.
你有沒有見過這些老師.
花多點時間讀書.
但也有些老師.
偏偏不是讀書最厲害那些.
可能中間是尾巴那些.
很奇怪.
但是無論如何.
老師.
其實是教化人.
我們現在慢慢就變成.
幫他做考試.
老師變成補習社.
補習社就不會做人的工作.
我理你死,看電視.
我賺錢.
我是金牌.
補習社的老師.
我幫助.
今年在我班上.
出了21個五星星.
賺錢.
老闆要賺錢要招來.
就不是教化.
學校的老師.
除了是教書.
真的教化.
起碼我自己的年代.
我受我的老師,好幾個老師影響很大.
我英文名叫Paul.
也是受我中二的老師影響.
因為他英文名是Paul.
那時候成績表.
還要是.
每個禮拜派成績表有老師簽名.
接著我們回去給家長看.
家長簽名回去.

$^{1001}$所以我知道老師簽名.
那年我們全班同學都受我們班主任老師影響很大.
我自己那時候沒有英文名.
升中三就叫自己Paul.
簽名都跟老師簽名.
我的Paul張.
簽名是因為他特別有個簽法.
我就照他簽法.
到今天都是這樣簽法.
我的銀行戶口都是這樣簽法.
影響很大.
教化,他教化了我有些東西.
所以服事就有潔淨的功能.
在我生命裡面有個正面的影響.
耶穌以奴僕身份.
服事門徒.
選擇了洗腳為榜樣.
有潔淨的功能.
耶穌以道成肉身來到世界.
以道教化他的門徒.
亦都教化我們正正就是.
以道潔淨門徒的生命.
我們的服事亦都不應該只是停留.
在丹丹台台這些層面之上.
別人受惠於我們的服事.
亦都可以從得到表面上的幫助.
進而成為得到生命上的幫助.
被建立,被潔淨.
使都讓主能夠使用.
第四點.
主耶穌這場的表演.
是那種顛覆性.
十二節至十四節.
耶穌作為門徒的拉比.
師父就有主的身份.
居然成為奴僕.
一樣替門徒洗腳.
自然嚇倒門徒.
有一定顛覆他們一貫的思想.
作用的思想.

$^{1041}$耶穌的顛覆並不是顛覆自己的身份.
顛覆的是社會上一貫的做法.
一個大的管小.
小的服事大的不變之律.
社會上一定是這樣的.
大的一定是管小的.
小的一定是服事大的.
因為主的門徒不能照樣辦事.
所以主才特別突破這個框架.
破格地做.
羅馬社會是一個絕對階級的社會.
耶穌的教導和堅持正正顛覆了當時社會最大的生存基石.
就是小的要服事大的要聽話.
而今天照樣.
在我們現在這個時候顛覆我們這個生活.
兩千年之後.
自己以為我們已經解放了.
我們已經沒有奴隸制度.
我們看事情已經很進步.
這些人最有開明頭腦的現代人.
因為中觀這個世界就算是最先進的國家.
還是奉行大魚吃小魚.
小魚吃蝦 蝦吃沙.
這樣的金火玉律.
生存定律 森林生存定律.
所以還是一樣.
大魚吃小魚 小魚吃蝦.
蝦吃沙.
我們要問的就是在教會裡面.
在聖徒當中.
我們和我們的社會有什麼分別.
一定要問這件事.
要不然我們就是一個NGO.
就是一個普通組織.
街坊組織.
地方勢力.
洪水橋我們也是一個地方勢力.
跟其他的地方人士一樣.
跟我們沒有分別.
我們一定要問.

$^{1081}$我們和社會上一貫的做法有什麼分別.
不過當然我們也不要太簡單.
就說什麼都要顛覆.
什麼都要顛覆.
年輕人最喜歡什麼都要顛覆.
你說去左去右.
後現代的我們.
尤其是年輕的一群.
往往就是對顛覆這個顛覆那個.
有一定的興趣.
但是我們不能夠為顛覆而顛覆.
當然.
有推翻而沒有建設.
要注意的是耶穌.
他作為主.
作為大那個 作為拉比的權力.
底下來給予門徒.
一個榜樣耶穌.
是使用傳統.
他作為拉比 作為師父.
賦予他的權力來顛覆.
傳統對權力的使用.
聽清楚這句.
他是使用.
傳統賦予他的權力.
作為一個師父來顛覆.
傳統一直下來.
認為權力是應該這樣使用的.
作為拉比.
他是並非.
從一個無有者.
做出發 如果耶穌本來就是奴隸.
我替你洗腳.
和你做個榜樣.
你本來就是奴隸.
本來就是洗腳.
沒有意思的 同樣的動作.
另外一個人做.
是完全沒有意義的.
正因為主耶穌.

$^{1121}$不是一個無有者.
他是他們的師父.
所以要推倒的.
反而不是主耶穌.
他自己.
作為師父的身份.
門徒拉比的身份.
主乃是說.
我能夠以奴僕的樣式.
來服侍我的門生.
我亦都願意以朋友的身份.
來對待我的僕人.
你們亦都要.
彼此這樣對待.
在我們當中的社會.
當今的社會.
在家庭裡面.
我們自己身為社會 身為家庭裡面.
在社會裡面.
習慣做大事的人.
為高權重.
其實.
主耶穌的榜樣就意味著.
一件很重要的事.
就是我們每個人的服侍潛力.
尤其是如果你有遺訓.
你做大的那個.
服侍的潛力.
小的就本來要服侍大的.
是不是.
主耶穌有沒有說.
你大的要服侍小的.
以後都不用服侍大的.
主耶穌有沒有這樣說.
從來沒有的.
你有沒有想過這件事很奇怪.
主耶穌經常說大的要服侍小的.
但是從來沒有說.
小的不用服侍大的.
主耶穌沒有顛覆.

$^{1161}$小的要服侍大的這條規矩.
顛覆不了.
所以小的要服侍大的.
但是社會就是大的不需要服侍小的.
所以他顛覆就是顛覆這件事.
所以在這個情況裡面.
既然你是做大的.
那你就要再想一下.
爸爸媽媽相對於嬰兒那種情況.
主從來都沒有嘗試.
推翻小的服侍大的這個次序.
問題是世界上做大的那些.
多數都不願意.
又或者不懂得.
去服侍其他人.
而從主的眼光看.
其實你做大的.
是更加有服侍的潛力.
這個才是重點.
因為我們能夠.
用自己做大的.
無論是怎樣大.
大的身份和能力.
來建立更加多.
更加深和更加大的服侍文化.
一般小的做小的做不到.
做大的是有機會做到.
問題是願不願意做.
和怎樣去做.
這個才是其中一個.
最重要的精髓.
就是往往大的能夠做的事.
是小的做不到.
媽媽能夠做的事.
是嬰兒自己做不到.
所以媽媽不敢.
卑躬屈膝.
面乾睡濕.
我們說神昏定醒.
睡都睡.

$^{1201}$照顧嬰兒 嬰兒自己照顧不到自己.
所以大的有些事做到.
小的自己做不到.
我又拿現今的一個例子.
政治的例子.
不要說香港.
不要說我們的中國祖國.
我們拿別人來說.
侵侵.
無論你同不同意侵侵的政見.
他做法.
無論你怎樣看.
問題就是侵侵作為美國總統.
最大粒那個.
他做到的事是否其他人能夠做到.
就是做不到.
你一定要明白這一點.
是超級重要的這一點.
你以為你是大的.
Good.
意思就是你做到的事.
其他人小一點的人是沒辦法做到.
你能夠服侍的空間更加大.
看你願不願意去做.
還有你怎樣去做.
又回來近一點不要說侵侵.
香港的高官.
中國的高官.
如果他們願意.
仍然是公僕.
不管現在的意義是怎樣.
公僕這個字很聰明.
因為他做大.
他有權.
富裕無論是基本法.
富裕的權力他就有權.
所以如果要做好.
做福利.
是他們才能做到的.
平民老百姓自己在自己身上做不到.

$^{1241}$我們只會吵架.
是他一直下來.
推行又有工具又有整個團隊.
公務員隊伍又可以有.
立法執行這樣他就做到.
所以你一定要明白.
大主耶穌.
做最大那個這樣服侍.
門徒的意義.
尤其是對大人說的.
因為他們做到.
很多其他人都.
做不到的事.
所以我們更加要反思.
你覺得做大的你是大.
你更加應該想想你怎樣服侍.
小的.
還是要小的服侍你.
小的是不是都會服侍你.
不用想的.
是不是都會服侍你.
問題是你會不會服侍小的.
問題在這裡.
所以在.
謙卑服侍的事情上.
我們這樣看的時候.
就不是盲目的.
抄襲我們去和人洗腳.
而是我們做二次創作.
不同的做法不是洗腳.
我們是要清出於藍.
主耶穌和人洗腳.
我再進一步.
耶穌從來都不擔心我們技術上的表現.
祂是要求.
我們有祂的態度.
做show有個惡名.
做show就是假的.
做著做著演著表演.
有時我們會說那些官.

$^{1281}$實事不放.
只是懂得落驅做show.
做show好像很假.
但問題就是.
主耶穌這場show是很重要.
做得很真.
是祂練整個.
服侍生命裡面的.
一個結晶來的這場show.
所以我們就不要.
看輕做show這件事.
以後我見到某些人做show.
你都不要用貶義說他在做show.
你要看清楚.
他真的在做什麼.
可能有些事情是很重要的.
正如我們主一樣.
做show是可以的.
又或者說如果沒有.
現實的情況.
如果沒有掩飾我們自己裡面的居心.
如果沒有只是利己.
做這場show就是為了益自己.
收門票賺錢.
如果不是只是為了這樣.
可以為了這樣.
做show又有何不可呢.
我們的主在來世前就為門徒.
做了一場show.
show給他們看.
祂對他們的榜樣.
祂對他們的心意.
祂對他們的命令.
這個他們.
就應該變成我們.
做給我們看.
讓我們去學習.
我們一起去祈禱.
頂福上帝我們.
紀念我們.

$^{1321}$救主耶穌基督.
二千年前.
祂就快要面臨的.
一個犧牲.
祂的.
服侍的.
工作做到.
頂峰.
將自己擺上.
而且服侍我們這班.
頂撞祂本來就很多.
本來就不願意相信祂的人.
本來也沒有辦法自救的人.
我們回想.
我們主所擺上的人.
我們感恩.
我們被感動.
我們知道主.
的教導裡面.
我們應該也是照著祂去做.
所以頂福上帝當我們思考.
在這段日子.
我們又再回來思考.
主的榜樣的時候.
你也加添給我們力量.
你給我們有眼光.
你給我們有能力.
聰明智慧.
懂得如何彼此服侍.
尤其是我們當中.
有身份地位.
有能力的.
更加能夠在我們自己.
特殊的崗位.
所被賦予的.
一些位置上.
能夠做到.
更加大的服侍.
更加有意義的.
服侍的事情.

$^{1361}$求天父幫助我們.
讓我們在教會裡面.
能夠成為彼此服侍的一群.
誰聽我們禱告.
奉救主耶穌基督命求.
\newpage



\section{哈該書 1:1-15}
\label{sec:zmeyo_PNxjc}
\textbf{2026.03.22 《你們要省察自己的行為》:哈該書1\_1-15 (和修版) 陳淑娟牧師}
\newline
\newline
連結: \href{https://youtube.com/watch?v=zmeyo_PNxjc}{\texttt{https://youtube.com/watch?v=zmeyo\_PNxjc}} ~~~~ 語音日期: 2026-03-25
\newline
\newline
\hyperref[sec:Jj4ZioVnCDo]{\small{< < < PREV SERMON < < <}}
~
\hyperref[sec:index_chronic]{\small{[返順時目]}}
~
\hyperref[sec:index_scriptual]{\small{[返順卷目]}}
~
\hyperref[sec:_kJpwvPJ95U]{\small{> > > NEXT SERMON > > >}}
\newline
\newline
哈該書 1:1-15
\newline
\begin{longtable}{cl}
\hline
\hline
章節 & 經文 (和合本修訂版)\\
\hline
1:1 & \begin{tabularx}{0.7\textwidth}{X} 大流士王第二年六月初一，耶和華的話藉哈該先知向撒拉鐵的兒子猶大省長所羅巴伯和約撒答的兒子約書亞大祭司傳講，說： \end{tabularx} \\ \\ \relax
1:2 & \begin{tabularx}{0.7\textwidth}{X} 「萬軍之耶和華如此說，這百姓說，建造耶和華殿的時候還沒有到。」 \end{tabularx} \\ \\ \relax
1:3 & \begin{tabularx}{0.7\textwidth}{X} 耶和華的話藉哈該先知傳講，說： \end{tabularx} \\ \\ \relax
1:4 & \begin{tabularx}{0.7\textwidth}{X} 「這殿荒涼，你們自己還住天花板的房屋嗎？ \end{tabularx} \\ \\ \relax
1:5 & \begin{tabularx}{0.7\textwidth}{X} 現在，萬軍之耶和華如此說，你們要省察自己的行為。 \end{tabularx} \\ \\ \relax
1:6 & \begin{tabularx}{0.7\textwidth}{X} 你們撒的種多，收的卻少；你們吃，卻不得飽；喝，卻不得足；穿衣服，卻不得暖；領工錢的，領了工錢卻裝入有破洞的袋中。 \end{tabularx} \\ \\ \relax
1:7 & \begin{tabularx}{0.7\textwidth}{X} 「萬軍之耶和華如此說，你們要省察自己的行為。 \end{tabularx} \\ \\ \relax
1:8 & \begin{tabularx}{0.7\textwidth}{X} 你們要上山取木料，建造這殿，我就因此喜樂，且得榮耀。這是耶和華說的。 \end{tabularx} \\ \\ \relax
1:9 & \begin{tabularx}{0.7\textwidth}{X} 你們盼望多得，看哪，所得的卻少；你們收到家中，我就吹去。這是為甚麼呢？因為我的殿荒涼，你們各人卻只為自己的房屋奔走。這是萬軍之耶和華說的。 \end{tabularx} \\ \\ \relax
1:10 & \begin{tabularx}{0.7\textwidth}{X} 所以，因你們的緣故，天不降甘露，地也不出土產。 \end{tabularx} \\ \\ \relax
1:11 & \begin{tabularx}{0.7\textwidth}{X} 我命令乾旱臨到土地、山岡、五穀、新酒、新油和地上的出產，也臨到人和牲畜，以及一切人手勞碌得來的。」 \end{tabularx} \\ \\ \relax
1:12 & \begin{tabularx}{0.7\textwidth}{X} 那時，撒拉鐵的兒子所羅巴伯、約撒答的兒子約書亞大祭司，和所有倖存的百姓都聽從耶和華－他們神的話，就是哈該先知奉耶和華他們神差遣所說的話；百姓在耶和華面前存敬畏的心。 \end{tabularx} \\ \\ \relax
1:13 & \begin{tabularx}{0.7\textwidth}{X} 耶和華的使者哈該奉耶和華差遣對百姓說：「我與你們同在。這是耶和華說的。」 \end{tabularx} \\ \\ \relax
1:14 & \begin{tabularx}{0.7\textwidth}{X} 耶和華激發撒拉鐵的兒子猶大省長所羅巴伯、約撒答的兒子約書亞大祭司，和所有倖存百姓的心，他們就來為萬軍之耶和華－他們神的殿做工。 \end{tabularx} \\ \\ \relax
1:15 & \begin{tabularx}{0.7\textwidth}{X} 這是在大流士王第二年六月二十四日。 \end{tabularx} \\ \\
[1ex]
\hline
\hline
\end{longtable}
$^{1}$洪恩堂弟兄姊妹主禮平安.
感恩又能夠來到你們的教會.
很溫馨的來到.
特別一來到就是門口有茶飲.
又有咖啡飲.
好一個家的感覺.
真是很好的一個地方.
今日我選取的經文是.
《哈該書》第一章.
《哈該書》主要的內容.
其實是講神要求.
當時猶大的百姓重建聖殿.
一開始很多人看《哈該書》就會想.
我整間教會還這麼好.
裝修又不錯.
地方又不錯.
都還未倒塌.
為什麼要重建聖殿呢?.
我聽《哈該書》對我來說有什麼意義呢?.
首先我想大家知道.
當主耶穌為我們的罪釘在十字架.
死了埋葬復活升天之後.
教會就是耶穌的身體.
也在說地上的教會就是代表主的同在.
所以我要說一件事要搞清楚.
原來神的殿.
這裡就是神的殿.
這個地方就是我們要預備.
復興和重建的地方.
不是說還在這裡.
不是這樣.
另外我們的身體也是.
也是上帝的殿.
所以今天我想從《哈該書》裡.
跟大家反思一個重要的問題.
就是教會和你個人.
我們每一個信徒.
我們應該做些什麼.
才是合神心意.
另外就是我們怎樣復興教會.

$^{41}$也怎樣復興自己.
教會復興.
不是說人數多寡.
不是說禮堂有多美.
我們生命的復興.
也不是說我們只讀多少次經.
祈禱多少次禱.
而是究竟這個地方和我們的生命.
有沒有很好的被主使用.
所以今天很想這樣去探討.
我會分了幾個分題.
不過在讀《哈該書》之前.
我很想給大家了解一些背景.
因為可能有些人未必常常去看舊約.
因為我們很少看舊約.
《哈該書》其實是寫於大約在主前的520年.
其實很短的.
只有兩章.
所以很容易讀.
不過裡面總共有四個神諭.
每一個神諭的開始.
都是由一個日子和一句句句式開始.
所以讀《哈該書》是非常容易分的.
一開始都是耶和華的「華臨到哈該」.
今天我很想先交代簡單的歷史背景.
否則稍後聽起來你會蒙嘈嘈.
讓大家容易一點.
首先大家知道.
以色列人曾經三次被擄到巴比倫.
接著他們又回歸.
我們先看看那三次被擄的情況.
有不同的時間,很簡單的說.
譬如主前在605年.
一擄就擄了皇宮的精兵.
擄你當然是擄你勁的那些.
接著在主前597年.
就擄了約翰勁和一些精兵.
約翰勁王.
接著之後就出現了.
聖承被焚.

$^{81}$猶大亡國.
接著就是白聖和王被擄到巴比倫.
所以為何聖殿被毀就在這個時段.
接著到主前的538年.
第一次的回歸.
為何會回歸呢?.
就是後來波斯崛起.
波斯滅了巴比倫.
成為紅霸近東的一個帝國.
當時波斯王一個.
你應該聽過的.
聖經裡面有一個叫古列.
在主前538年下詔.
允許以色列人回到耶路撒冷重建聖殿.
這樣就出現了第一批以色列人回歸.
哈該書提到回歸有兩個很重要的人物.
剛才你讀了.
一個叫做所羅巴伯.
一個叫做約書亞.
這兩個是在波斯統治之下.
在猶太人來說最高的領袖.
第一這個所羅巴伯是最高的行政領袖.
而約書亞就是最高的宗教領袖.
當時他帶著五萬人左右去回歸.
回到耶路撒冷這個地方去重建聖殿.
恢復敬拜的中心.
當然之後還有第二第三次的回歸.
在第一次回歸之後.
主前的535年2月.
他們其實已經開始動工建這個聖殿.
根據前面的以色列記講.
他們已經立了根基.
建東西當然要立了根基.
然後他們就面對很多的困難.
很多的問題.
加上敵人開始反對.
最後達薩西王下令.
你們不要搞那麼多事情了.
麻煩停工.
沒得搞.

$^{121}$直到又有另一個王.
波斯王叫做大利烏一世.
他重新找來以色列人建聖殿.
不過很可惜.
因為之前的經驗.
這些以色列人對聖殿的重建一點都不起勁.
可能心想你不要搞我了.
結果又停工.
直到主前的520年.
就是哈該書的一章一節.
他是這麼說的.
他和哈該和兩位領袖對以色列人展開責備和勸勉.
我想大家想一想.
主前538年回歸.
535年立了根基.
接著哈該在520年宣講.
究竟他停工停了多少年.
大家已經睡了吧.
幾年了.
誰這麼聰明.
主席厲害.
只有主席不能睡.
其他人可以睡.
停了工15年.
你想一想.
這麼重要的事情.
重建聖殿.
他們為什麼竟然停工了15年.
接著哈該書的第一章一到五節.
就是第一個神諭.
神要跟他們說什麼.
今天我從這個神諭裡有幾點跟大家分享.
第一點.
我叫他的主題叫做.
主次不分.
在第一到第四節.
剛才讀了不用特別讀了.
如果你有留意第二節.
萬軍之耶和華如此說.
這些百姓對於建造耶和華的殿.

$^{161}$他們認為是什麼.
時候未到.
如果用我們現在的言語說很簡單.
就是時間不對.
時間不對.
不是不做.
你要弄清楚他們沒有那麼大膽.
神叫他們建聖殿.
他們敢說不要搞我.
不會的.
今天神跟你說你要讀經祈禱.
你也不會說不要搞我.
我們也是這樣回答上帝.
時間不對.
這一句表明以色列民.
不是不知道重建聖殿的重要性.
是不知道的.
甚至他們覺得上帝要我做.
不過他們就說.
我沒有說不做.
不過就是時間未到.
百姓用什麼理據.
認為時間未到.
我們先看看百姓.
看看神是不是冤枉他們.
可能真的時間未到.
在回歸之前.
其實他們曾經展開過一輪工程.
不過就是因為出了很多問題.
他們先建了祭壇.
又立了根基.
他們確實面對著很多問題.
他們沒有說謊.
第一當時耶路撒冷.
土地非常荒涼.
經濟環境很差.
第二在建造工程.
附近各國族群.
住在那裡的人.
當然你走了那麼久.

$^{201}$還有波斯官員不斷批評和反對.
還有這個你聖經會知道.
撒瑪利亞的貴族.
對於回歸的猶太人非常有敵意.
這些貴族覺得.
我們已經在那裡那麼久.
你走來搗亂.
他們不願意將這個特權.
給回猶太人.
這些都是利益問題.
很多時候地上打仗都有這些問題.
如果是這樣.
你看一看.
經濟又不理想.
又有人在那裡搞來搞去.
這些百姓說.
時間未到.
合不合理呢.
理論上是合理的.
不過我想大家看第三第四節.
第三第四節這樣說.
耶和華的華籍哈該.
先知傳講說.
這殿荒涼.
你們自己還住在.
天花板的房屋嗎.
為何神這樣說他們呢.
其實神直接指出.
生命中很核心的問題.
這一群屬神的子民.
他認為建聖殿.
就未是時候.
面對這麼多事.
不過這樣建自己的房屋.
就沒有什麼問題.
什麼是天花板的房屋呢.
我想澄清一下.
不要想著.
為何上帝會這樣.
建房屋不讓建天花板.

$^{241}$不是要建茅屋嗎.
其實天花板的房屋.
是指上等的木材.
裝修得.
一間可以叫豪華的大屋.
為何這樣說呢.
因為這個原文的字.
經常用來形容皇的宮殿.
即是豪宅.
我想不止.
即是漂亮的東西.
神不是批評人住.
有天花板的房屋.
神是說你們既然.
豪裝自己的房屋.
但是你們可以任上帝的.
電荒涼.
請問你擺了神的.
位置在哪裡.
如果當時.
政局真的這麼不穩定.
經濟又是.
百廢待興的話.
不見上帝的電.
不是時候不出奇.
但是百姓卻自己.
去想著去建立.
自己的豪華的房屋.
百姓花盡心思.
時間的.
自己在這個地方.
謀一個好的環境.
將來的生活去籌算.
神的事.
他卻擺到最後.
所以.
神重重的責備.
神席在哈該.
才在糾正以色列民.
一個錯誤的觀念.

$^{281}$就是主持不分.
輕從不分.
做事永遠有先後.
他們用了.
自己的金錢.
時間來處理.
自己的事.
將神的事擺到最後.
甚至拖下線.
生活.
享受才是為重.
神的事.
為輕.
所以聖殿的工程.
足足停了十五年.
至九.
弟兄姊妹.
第一樣我很想大家反思.
我們今天無論是在神國的事.
無論是我們自己.
屬靈的生命.
在教會這個家裡面.
我們有沒有主持不分.
你有沒有急神所急.
還是你只是喜歡.
做我自己的事.
神的事.
我先拖下線.
我會不會是這樣呢?.
聖經怎樣告訴我們?.
聖經清晰告訴我們.
你們要謹慎行事.
不要像愚昧人.
要做智慧人.
愛惜光陰.
因為現今的世代邪惡.
我們做.
上帝做我們是很公平的.
每個人都是七成二十四.
沒有一個人多一點.

$^{321}$是一樣的.
但問題是我們怎樣有智慧.
去用我們這二十四小時.
如果我們不謹慎.
我們就會浪費.
一些沒有益處的事.
我們要愛惜光陰.
善用時間.
造神喜悅.
要我們做的事.
善後次序和主持.
是非常重要的.
你懂不懂得.
運用好你每一天呢?.
我以前在商界的時候.
很喜歡講時間管理.
時間管理.
這個圖沒有什麼人陌生.
這個比較簡化.
理論上可以做到七至八格.
我做了一個比較簡單的.
就是.
通常我經常都講.
我有時候對同工說.
你不要先做一些不要緊的.
不要緊的你不做.
因為我們經常都主持不分.
什麼是時間管理呢?.
就是你看到有四格.
通常我們會做什麼呢?.
就是做緊急和重要的.
最後是什麼呢?.
不緊急又不重要.
當然中間就會有.
緊急不重要的拉扯.
和重要不緊急的拉扯.
我們很多時候出了兩個問題.
第一個問題.
我們連放這些格都不懂.
什麼叫放這些格都不懂呢?.

$^{361}$就是究竟什麼是最重要的呢?.
有時候那些人會擺了.
我經常都跟年輕人說笑.
如果給你們填呢.
我懂得填.
最緊急最重要當然是爆機.
是不是?.
否則你不會用最多時間打機.
對不對?.
你不會飯都不吃去打機.
最緊急一定不是吃飯的.
你們餓不死的.
全部都是神仙.
我經常笑他們.
如果家庭主婦.
通常最緊急是什麼呢?.
想想今晚煮什麼菜.
因為真的很頭痛.
你可能坐著聽牧師講道.
都想一想.
一會兒買什麼好呢?.
那麻煩你扔下這個.
這個其實不是最緊急的.
我們就將那些緊急重要的.
其實是放錯了.
當然有些人就說.
不是,牧師,我真的有很多事要做.
我又要打機,又要煮飯.
我又要陪兒女溫書.
我又要OT博殺.
我又要應酬.
我又要追韓劇.
我也知道要讀經祈禱.
那很簡單.
全部擺了緊急重要就行了.
那就慘了.
又是主持不分.
你只有24小時.
所以第一.
我們出現的問題就是.

$^{401}$我們不懂得分什麼是主.
什麼是次.
什麼又輕.
什麼又重.
你寫回你一天24小時.
你發覺你就不會分了.
你就一定像這群以色列民那樣.
時機還沒到.
什麼時機還沒到.
還沒滾完手機.
所以還沒有時間靈修.
時機還沒到.
還沒時間祈禱.
因為還在追韓劇.
我們不懂得分.
第二種就是.
可能有些人都懂得分的.
你又沒有那麼愚昧的.
我知道很重要的.
傳福音重要.
重要和緊急.
末世了.
是不是.
讀經祈禱.
也沒有那麼緊急.
所以我們放在重要.
但不緊急.
也對的.
不一定很緊急.
你放在正確的位置.
但你忽略.
什麼忽略呢.
最多人忽略的.
就是剛才我說的.
重要.
但未必很緊急.
其實靈修祈禱.
就是其中一個.
重要嗎.
你們知道重要嗎.

$^{441}$但不是一定是.
飯都不吃.
我就做.
不是的.
未必很緊急.
但往往我們忽略.
就是這一格.
我們將很多.
屬靈的事情.
我們知道重要.
因為你覺得.
當下不緊急.
你放在這一格.
接著呢.
你就通常不做.
我們最喜歡做的是什麼.
知不知道.
不緊急又不重要那些.
不知道為什麼.
我們的生命.
亂七八糟.
所以我們真的要.
求主幫助.
今天.
教會的復興.
都是一樣.
我很多時候.
和教會.
我和你們教會.
都是很熟悉.
因為你們教會.
即將和我一起上課.
所以剛才我見到.
我是其中一個學生.
平時原來這麼漂亮.
平時可以上課.
戴口罩.
上來獻唱.
我好感動.
在教會復興.

$^{481}$說要重建教會.
讓教會進入社區裡.
去做科研工作.
這是沒有一間教會.
可以推諉的.
但很多時候.
我發覺當我去做.
教會的同行的時候.
教會很多時候.
就是主持不分.
很多教會和我說.
牧師牧師.
我們真的很想進入社區.
不過時間還未到.
為什麼時間還未到.
我要先搞定弟兄姊妹.
這麼多適合的牧養弟兄姊妹.
怎麼搞.
搞定.
我真的很想問.
怎樣才能搞定.
我沒有叫你放下弟兄姊妹.
去做傳福音.
你要和弟兄姊妹一起傳福音.
不過我們有時就是這樣.
我們都和上帝說.
時間還未到.
時候還未到.
因為我覺得這樣.
我們放錯了位置.
包括我們的生命也是.
我們很多時候覺得.
我經常都聽到.
退休後才回來服侍.
因為退休前的時間還未到.
是不是這樣.
如果從今天的經文看.
其實不是時候還未到.
你有時間去旅行.
你有時間去做很多事.

$^{521}$你有時間去看那些.
十五秒的短劇.
但你沒有時間回來.
去服侍多一刻.
這個不是你沒有時間.
是你主持不分.
所以弟兄姊妹.
第一點.
我很想大家認真去思考.
正確的目標.
決定了你時間的優先.
今天甚麼是你正確的目標.
是不是上帝的事.
還是你自己的事.
如果根據神執秘.
所羅巴伯和約書亞和百姓.
他們的主持不分.
就是將神的事放到最後.
今天為甚麼我們的格數.
會放到亂七八糟.
是因為我們用我們自己的事.
去決定怎樣放這些格數.
如果你有一個正確的目標.
你知道神的事.
才是你生命最首要的.
你就知道甚麼是最重要.
所以第一個生命的反思.
我想大家再想一想.
在你的生命的次序裡.
神是不是為首.
第二.
在第二段的經文.
在五到十一節.
我叫它做倒果為因.
我想大家先留意.
第五到第七節.
它不斷重覆了一個模式.
就是你們要審測自己的行為.
亦都是今天我講題的主旨.
你們要審測.

$^{561}$其實原文是安置你們的心.
即是你們要安置好你們的心.
當人的心放在哪裡.
你的心思意念就會放在哪裡.
跟著你的行動.
自然就會跟著你的心.
放在哪裡的方向去跑.
所以在經文裡面.
亦有一個是你會看他們奔向哪裡.
安置心是非常重要的.
究竟我們的位置放在哪裡.
剛才講到如果你願意.
去放好你的心在哪裡.
你的行為就會在哪裡.
行為原文是In your way.
在你的道路裡.
首先神要求他們在五到十一節.
神先要求他們注意當時他們惡劣的環境.
我想大家看一下.
神是沒有否定百姓當時遇到的經濟狀況.
反而指出它的真相.
剛才主席讀的時候.
他們當下的真相是怎樣.
種得多收得少.
吃又吃不飽.
喝又喝不足.
穿不暖其實都很慘.
他們的狀況是修吃喝穿.
這些是基本人的生活條件.
他們都不足夠.
神將問題提出來.
是讓以色列百姓警惕自己錯誤的行為.
第一包括了甚麼呢.
耶路撒冷經濟的確是困難.
不過他們卻願意在這個難處當中.
依然合理化地去建立自己的房屋.
卻不願意建聖殿.
另外百姓他們在第九節.
特別我想給大家看第九節這個字.
見不見到有個「各人顧自己的房屋」.

$^{601}$這個顧字原文是奔向奔.
即是百姓為了自己.
這麼惡劣的環境也好.
他為了自己的居所.
自己的安樂而奔波勞碌.
卻讓上帝的殿的荒涼.
下蓋教導百姓.
請你們對準焦點.
注意自己的處境.
就算他們為自己的房屋.
終日奔波勞碌.
是不代表他們獲得安全和保障的.
因為結果經文都說他們是徒勞無功.
以色列民因為現在的環境和經濟.
所以不積極做神所吩咐的功.
但這段經文卻剛剛相反地告訴他們.
其實是因為他們忽視了神的吩咐.
所以今天帶來這個經濟的困局.
相反的,倒果為因.
舊約其實其他先知同樣有這樣的教導.
如果你記得,何西亞先知指出.
以色列民即使播種也不能夠收成.
就算收成也會被外邦人所吞吃.
原因是因為他們拜偶像.
接著尼迦先知教導他們.
告訴他們吃不飽,殺種也沒有收割.
原因是因為百姓作惡,受到神的懲罰.
耶利米先知則宣告.
他們做公務好處,種的麥子會收到荊棘.
因為他們要接受神的憤怒.
今日新約同樣有這樣的教導.
在《馬太福音》六章三十三節.
大家應該很熟悉這段經文.
請大家讀,預備起.
「這些東西都要加給你們了」.
「原文是你需用的,他都給你」.
弟兄姊妹,以色列民用了結果.
現在經濟不好,收也收不到.
吃也吃不飽,所以做藉口.
環境不好,見聖殿未是時候.

$^{641}$經濟好些再想.
但神卻告訴他們.
「你以為收得多嗎?卻收得少」.
「你以為收到家中嗎?卻催去」.
為甚麼?是因為他們搞錯了一件事.
他們不是先求神的國和神的義.
然後相信他們需用的,上帝也給.
弟兄姊妹,我們會否犯這個錯誤?.
我們經常覺得我們有難處.
所以我們不去遵行神的旨意.
但我們沒有想過.
原來當我們願意遵行神的旨意.
這些難處都會迎刃而解.
你有否想過是這樣?.
如果沒有,我想問這段經文又如何解釋?.
「你相信先求他的國和他的義」.
「你需用的東西,他都給你」.
有時我們背金句很厲害.
但我們走出來就很弱.
我不是叫大家不關心弟兄姊妹的難處.
但當我們知道有難處時.
同時亦表現了一件事.
究竟我們覺得難處較大.
還是神的應許較大?.
舉兩個例子.
我最近舊的不說了.
這個都頗新的,可能學生聽過.
我最近在上水建了一個新的中心.
要請人.
其中有一個婦女.
其實不是信了很久.
一年多,洗了泥.
我見到她平時很勤力回來做義工.
她是一個單親媽媽.
有兩個小朋友.
兩個都是SEN.
回來幫忙,很勤力.
真是天天都見到她教會.
搬東西又做,倒垃圾又做.
幫忙小朋友,照顧小朋友又做.

$^{681}$接送癌症病人又做.
只是義工,甚麼都做.
我祈禱當初感動.
不如請她回來.
因為開了新中心,需要人手.
我需要一個行政同工或福音幹事.
我問她.
我有感動,不如一起祈禱.
祈禱了一個多月.
她說牧師,我願意.
我願意回來工作.
然後我就請她.
很便宜的工資請她.
可以說是一萬多元.
一萬二千元加上信用卡請她.
臨請她前,就要填寫.
要填寫保險和信用卡.
有一份很簡單,好像履歷表.
填寫完後.
我收了之後,打算給我財務.
一收,我看了之後.
哎呀,為何我哀怨呢?.
原來她之前拿了綜援.
為何我哀怨呢?.
我看了,她的綜援是一萬六千二千元.
哎呀,你給我嗎?.
哎呀,糟糕了.
我馬上不敢交給財務部.
我先找傅蕾瓊.
你拿綜援的嗎?.
她很輕言淡語.
是呀,你說請我.
一月一日開始,我上個月就取消了.
我說,我給你一萬二千元,還要連同NPF.
她說,我知道了.
我說,你兩個小孩要交租.
交租好像是五千八百多元.
五千八百多元.
兩個學生有很多訓練.
你又單親.

$^{721}$我說,一萬二千元夠不夠用?.
她說,夠了,牧師.
教會有飯吃,我和婆婆吃就行了.
她真的這樣說.
她說,沒問題的.
那一刻,我都不知道可以說什麼.
我跟她說,其實你有很大的難處.
為什麼你不跟我說?.
其實你真的不夠用.
因為我跟她計算過,每個月只剩下二三百元.
你怎麼夠用?.
我說,你明明每個月有三四千元.
有什麼意欲?.
我們的人,經常都覺得沒有錢在口袋裡.
很沒有安全感.
她回答,我也回答得很好笑.
我說,我習慣了.
她當然開玩笑.
她身上有幾千元.
我說,那怎麼辦?.
她說,怎麼辦?你不是不聘請我吧?.
我說,不是.
我說,為什麼你會答應呢?.
對我來說,我是很迷惑的.
第一,綜援需要做嗎?.
不用做.
回到我們那裡很慘的.
七成二十四的.
第一天放假,但他們都會回來.
這樣選擇.
她說,先求他的國和他的義.
我需用的,上帝都給我.
你教的,牧師.
我回到家裡.
我哭了很久.
我死了不少錢.
我跟我的教會商量.
當然也要公平.
因為我其他的行政同工都是這樣.
但她真的很甘心樂意.

$^{761}$真的很甘心樂意.
現在已經過了試用期.
做得非常好.
我在試用期前跟她聊天.
她說,牧師,是真的.
原來上帝說.
我需要用的東西都加給我.
我沒有加工資給你.
加給你.
她說,原來她做了這幾個月.
她的子女乖了很多.
她說,給我幾千元.
根本不是.
不是我想要的東西.
原來上帝才知道.
什麼是她需要的.
她的子女在教會收拾桌椅.
真是乖得.
有一個經常打人.
最近加入了我的小六門徒訓練.
自己來的.
牧師,你是否有門徒訓練.
我是,我行不行.
現在說吧.
她的子女乖得.
她看到媽媽的生命.
媽媽以前經常在家裡哭.
回到現在,經常說.
我今天又陪過癌症病人去哪裡.
我很開心,我們要珍惜生命,我們要祈禱.
兩個子女很大的轉變.
所以弟兄姊妹,我很想大家.
去從這個生命.
去看看,我們有沒有這樣的信心.
先求他的國,和他的業.
教會一樣.
很多時候,我們很擔心去做很多事情.
究竟有沒有果效.
要花多少錢.
教會剩下多少錢.

$^{801}$我們很多時候都是這樣.
但是.
我的教會經驗.
我經常說.
十一年前,蒙召回到我的教會.
教會原本是中產教會.
每個月.
收支平衡,還有一點收入.
我去了一年之後.
全程都是紅色的.
每個月都是負數.
我的執事就開始有點害怕.
有一次開會.
牧師,怎樣教牧師.
我全道.
我不出糧.
我都這樣跟他們說.
你做上帝的工作.
上帝不會虧負我們.
你相不相信.
如果我們的信心這麼弱,真是差到不得了.
我告訴你.
現在十一年過後.
我的教會.
當然我們都很堅持.
只是儲三個月的薪水.
十一年前.
我們負數只有養一個童工.
我們現在有十二個童工.
第二.
我們在COVID的時候,一次過五百萬買一層樓.
現在.
除了教會七千呎.
我有一層五百多呎的.
做Footband,我還租了四個單位.
接下來會租地舖.
先求他的國.
和他的義.
這些錢不是我們的.
是上帝給我們去做他的工作.

$^{841}$你相不相信呢?弟兄姊妹.
所以我很想大家搞清楚.
很多時候我們倒過為一.
覺得我們現在什麼都不行.
所以什麼都做不到.
其實是因為你什麼都做不到,所以不行.
你沒有做上帝的事,所以不行.
搞清楚.
所以今天,我很想問大家.
你的心在哪裡?.
你的心.
是不是在教會的復興裡面?.
你想不想有一天.
不單坐滿這個地方.
我經常和我的同工這樣祈禱.
我們由六十人開始祈禱.
連同小朋友.
我們祈禱要有.
坐滿堂,結果坐滿了.
我們要分兩堂.
我們的堂只可以坐.
二百人,結果我們現在四百多人崇拜.
要分兩堂.
我們又沒有青少年.
我們沒有到只有三個小朋友.
有一個是我兒子.
他死都不回來教會.
因為他說,媽媽.
你怎麼回來團契?.
Do Re Mi.
然後導師又是你,我不想見到你.
又是Do Re Mi.
我是很明白的.
他的團契是在他自己.
中學的團契長大的.
他說,我真的不想見到你.
一直都是你.
現在我們青少年崇拜有幾十人.
我們又沒有小學生.
只有幾個,我們現在.

$^{881}$只是移崇八十人.
我們又沒有錢.
去到今天,不是叫做.
豐祝,但我們可以支援其他教會.
弟兄姊妹.
今天你有沒有.
這個心去放在.
正確的位置,而相信.
神的供應.
最後.
行動帶來.
復興,在十二到十五節.
當.
下蓋宣告神的.
曉遇之後.
所羅伯伯.
約書亞和百姓.
從先知下居領受了.
神這麼重要.
的一個警醒.
責備,在.
第一章十二到十五節.
記載各人的回應.
由政治上的首領.
所羅伯伯,宗教上的大祭司.
約書亞帶領群眾.
重建聖殿.
所羅伯伯.
約書亞是最高的行政領袖.
對於重建.
聖殿其實責無旁貸.
這也提醒.
我們作為領袖.
不應該.
被百姓誤導.
也不可以.
放任百姓.
不聽上帝的說話.
這給我自己作為教睦.
或者這段.

$^{921}$宣講也在我自己的教會講過.
我也提醒我教會的執事.
這是非常重要的提醒.
當下蓋.
宣告神的話之後.
神就激動眾人的心.
眾人的心.
就放回正確的位置.
那個位置是什麼?.
敬畏的心.
剛才說.
他們的心放錯了.
他現在放回正確的位置.
放在神的裡面.
帶著敬畏的心.
不單止,他還悔改回轉.
並且行動.
重新重建聖殿.
在很多時候.
聽到神的教導.
我們被聖靈感動.
我們都會知道悔改.
但我很想大家記住.
今天你聽完之後.
悔改的重要不只是悔改.
不只是悔.
後面是什麼?.
改,一定配合你.
有一個立志的行動.
生命有改變.
所羅,巴伯和約書亞.
甚至百姓看到這裡.
我非常欣賞.
因為他們知道自己的行為.
他們立即悔改.
而且帶有行動.
聽到這裡.
你可能有以下的疑問.
神要以色列民.
重建聖殿的意義在哪裡?.

$^{961}$為何這麼緊張的上帝?.
在舊約當中.
從會幕開始.
神帶領以色列人出埃及.
在曠野漂流.
在他們每一天不斷前往.
帳幕是怎樣?會幕是怎樣?.
跟著移動.
當以色列民進到迦南地.
帳幕就不再移動.
就改為固定的建築物.
帳幕和聖殿.
有一個很重要的意義.
請大家記住.
就是神與人同在.
這是神與人的立約.
在新約.
耶穌基督就是上帝的同在.
正如我一開始.
給大家看那節經文.
在舊約中.
重建聖殿的意義是.
恢復以神為中心的事奉和敬拜.
恢復人以神為中心的生活.
神就與人同在.
使神得著榮耀.
所以當時.
以色列民只是專注在自己的事中.
把重建聖殿的事放下.
他們失去了.
上帝的同在.
以色列民聽到神的教導.
知道自己的問題.
主持不分倒果為因.
諸多的藉口推搪重建聖殿.
現在他們悔改.
帶著敬畏的心.
再次的來到.
所以在第十三節.
非常之重要.

$^{1001}$第十三節清楚講明.
耶和華說什麼?.
我與.
你們同在.
他們重新恢復.
上帝的同在和祝福.
因此土地乾旱.
田產再沒有.
生產的時候.
神今次再次.
應許我與你們同在.
聖殿和教會.
最重要的是.
神與我們同在.
我們和祂納了約.
我們的生命一樣.
神與我們同在.
我們有沒有.
深思這份關係.
而我們遵行上帝的旨意.
在哈該書最後.
結果在六月二十四日.
距離先知宣講.
有多少日.
大家有沒有睡著.
第一節是六月初一.
現在是.
即是.
二十三天.
他們重建聖殿.
大家想想他們聽完教導.
二十三天就開始重建聖殿.
快還是慢?.
絕對快.
還快得很離譜.
為什麼我這麼說呢?.
因為重建聖殿是需要木材和材料.
他們已經停工了.
十五年.
你想像一下又要找木材.

$^{1041}$又要清那個藍地.
但是他們二十三天已經可以.
開始重新啟動.
重建聖殿.
為什麼可以這麼快速?.
是因為他們帶著敬畏.
和迴轉的心.
如果應用在今天教會.
我們教會如何合神的心意呢?.
教會不是一座建築物.
而已.
這座建築物.
不需要有多宏偉.
裡面的活動不用多姿多彩.
這個並不重要.
重要的是.
神有沒有臨在.
在這裡.
是不是與我們同在.
如果我們這座建築物.
沒有真正的信仰實踐.
沒有真正的努力.
傳福音.
愛倫如己.
彼此相愛.
進行上帝的使命.
我們沒有實踐到信仰.
只剩下各樣的活動.
信徒沒有背上自己的十字架.
跟隨主.
我們的生活.
走的是世界價值觀.
我們的心還在世界裡面.
為了那份高薪厚職的工作.
寬敞的家居.
沉醉於各樣的享樂.
這個敬拜是虛假的.
信徒失去了以神為中心的生命.
沒有盡心盡義.
盡聖的愛神.

$^{1081}$並且愛倫如己.
這一切所謂教會的活動.
都是枉然的.
有時候一間裝修宏偉.
冷氣舒服.
設備齊全的教會.
往往成為了我們的.
安書區.
成為了我們一個信徒的.
俱樂部.
很感恩.
洪恩堂.
我看到你們的橫額.
其實讓我很感動.
我們這個社區.
將福音.
遍傳與人同行.
這是非常重要.
願主真的帶領你們.
繼續在這個洪水橋的地方.
成為真正的燈台.
而每一個信徒.
無論在這裡.
在你的家庭.
在你的學校.
在你的生活.
在你的工作場所.
你都要成為那個光和炎.
不僅是在這裡自嗨.
四面牆.
在這裡說「哈利路亞」.
讚美神是重要的.
但生命的讚美也更加重要.
最後我想做一個.
今天的總結.
今天你聽到什麼呢?.
你願不願意.
調整你整個生活的.
優先秩序.
被神差遣.

$^{1121}$以神為首呢?.
願意參與.
重建聖殿.
讓神的榮光.
繼續在這裡彰顯嗎?.
我們一起祈禱.
天上的巴父.
我要獻上感恩.
我看到洪恩堂.
在這個社區.
他們很願意.
跟那些失喪的靈魂.
跟那些還在水深火熱的.
街坊.
關愛他們.
去傳福音.
神啊!我們看到.
這一群的以色列民.
他們被索羅巴伯.
被約書亞所帶領.
但我更加看到.
這個悔改.
不僅是領袖.
還有一群百姓.
如果只是靠索羅巴伯.
和約書亞重建聖殿.
是不可能的.
但神啊!你的說話.
藉著蝦蓋.
激動每一個人的人心.
同樣今天洪恩堂.
要復興.
要回到上帝.
正確的位置的時候.
不僅是一群.
執事.
有心人或牧者所做.
而是需要.
每一個弟兄姊妹.
都願意重建他的優先次序.

$^{1161}$認罪.
悔改.
去你的裡面.
你說我們的心在哪裡.
我們的財寶就在哪裡.
我們要將我們的心.
安置在我們的使命裡面.
願.
你大大的去復興.
這裡.
不僅是復興教會.
也是復興我們每一個弟兄姊妹的心.
讓他們.
成為你.
貴重的器皿.
以上是禱告.
阿們.
\newpage



\section{約翰福音 12:23-28}
\label{sec:_kJpwvPJ95U}
\textbf{2026.03.29 《得榮耀的時候》:約翰福音12\_23-28 (和修版) 黃國維牧師}
\newline
\newline
連結: \href{https://youtube.com/watch?v=-kJpwvPJ95U}{\texttt{https://youtube.com/watch?v=-kJpwvPJ95U}} ~~~~ 語音日期: 2026-04-02
\newline
\newline
\hyperref[sec:zmeyo_PNxjc]{\small{< < < PREV SERMON < < <}}
~
\hyperref[sec:index_chronic]{\small{[返順時目]}}
~
\hyperref[sec:index_scriptual]{\small{[返順卷目]}}
~
\hyperref[sec:ke2vMZ9uUT4]{\small{> > > NEXT SERMON > > >}}
\newline
\newline
約翰福音 12:23-28
\newline
\begin{longtable}{cl}
\hline
\hline
章節 & 經文 (和合本修訂版)\\
\hline
12:23 & \begin{tabularx}{0.7\textwidth}{X} 耶穌回答他們說：「人子得榮耀的時候到了。 \end{tabularx} \\ \\ \relax
12:24 & \begin{tabularx}{0.7\textwidth}{X} 我實實在在地告訴你們，一粒麥子不落在地裡死了，仍舊是一粒；若是死了，就結出許多子粒來。 \end{tabularx} \\ \\ \relax
12:25 & \begin{tabularx}{0.7\textwidth}{X} 愛惜自己性命的，就喪失性命；那恨惡自己在這世上的性命的，要保全性命到永生。 \end{tabularx} \\ \\ \relax
12:26 & \begin{tabularx}{0.7\textwidth}{X} 若有人服事我，就當跟從我；我在哪裡，服事我的人也要在哪裡；若有人服事我，我父必尊重他。」 \end{tabularx} \\ \\ \relax
12:27 & \begin{tabularx}{0.7\textwidth}{X} 「我現在心裡憂愁，我說甚麼才好呢？說『父啊，救我脫離這時候』嗎？但我正是為這時候來的。 \end{tabularx} \\ \\ \relax
12:28 & \begin{tabularx}{0.7\textwidth}{X} 父啊，願你榮耀你的名！」於是有聲音從天上來，說：「我已經榮耀了我的名，還要再榮耀。」 \end{tabularx} \\ \\
[1ex]
\hline
\hline
\end{longtable}
$^{1}$各位節目 早安.
很感恩今天能夠跟大家一起去到敬拜.
以及聆聽上帝的說話.
剛才那段經文也是耶穌一個跟別人說的話.
說祂自己得榮耀的時候到了.
簡而言之這段經文也是跟這個星期很有關係.
大家知道下個星期就是復活節.
所以祈禱很多假期之餘.
其實都是紀念主耶穌基督的受死和復活.
耶穌基督那時候在耶路撒冷釘十架之前的那個星期.
就進入耶路撒冷.
那個教會傳統就叫做中枝主人或者中樹主人.
就是說耶穌進入耶路撒冷的時候.
很多人就用中樹枝去歡迎祂.
就是和撒拿和撒拿.
應當稱頌的救主來的.
所以這段經文也是出自那裡.
因為約翰福音 耶穌就對著那些人說.
人子得榮耀的時候到了.
得榮耀聽起來很好 每個人都喜歡榮耀的.
不過在約翰福音裡面講到榮耀.
未必是我們人想像的那種意思.
例如到了約翰福音最後的時候.
就是耶穌最後了 復活之後.
他就跟彼得有一個對話.
耶穌也有講到榮耀.
這裡耶穌跟彼得說.
我實實在在地告訴你 你年輕的時候.
自己速上帶自己隨意往來.
但年老的時候你要伸出手來.
別人要把你速上 帶你到不願意去的地方.
這裡耶穌跟他的門徒彼得講他將來會發生的事.
最後一句是說什麼.
耶穌說者說是指著彼得會怎樣死來榮耀神.
當我們讀到這裡 原來約翰福音將死跟榮耀放在一起.
將受苦跟榮耀放在一起.
有時候對我們來說.
都是一個很特別或者很難了解的狀況.
不過我想在教會群體裡面.
我們做基督徒很多時候都會了解.

$^{41}$其實榮耀是一個很被尊重的人.
還有一個人受苦.
其實很多時候都出現在同一個人身上.
我自己其實都是宣導會的人.
我自己的舞會是北角宣導會.
跟大家有一點遠.
港島東和新界西北有一點遠.
不過大家都是宣導會.
那時候我在北角宣導會聚會的時候.
入神學院讀福兆 奉獻.
我一做神學生 發覺原來在北角宣導會北宣有一個傳統.
弟兄姊妹知道我讀神學的時候.
每次吃飯都不想給錢.
挺好的 大家應該學一下.
為什麼呢.
原來他們就說.
做神學生 將來做傳達人 牧者.
就很窮的了.
不要吃飯不用給錢.
我都挺感受到榮耀和受苦放在一起.
他尊重你.
覺得你是一個被神用的僕人.
不過平時同時知道你受苦.
所以教會裡面都有這樣的說話.
都會這樣去理解這個事情.
不過回到經文裡面.
耶穌在約翰福音第十二章.
為人子得榮耀的時候到了.
究竟那時候發生了什麼事.
叫到耶穌講這句.
我們就要回到早一點的經文.
這裡記載.
耶穌入城.
約翰福音第十二章十二節就說.
第二天就有一大群上來過節的人.
聽見耶穌要來耶路撒冷.
就拿著鐘樹枝去迎接祂喊著.
我撒拿以色列的王.
奉主命來的是應當稱頌的.
那時候耶穌去耶路撒冷.

$^{81}$其實是過節.
禹節是猶太人最重要的節日.
他和很多其他人都住在.
以色列巴勒斯坦地.
很多不同的地方.
都要去耶路撒冷的聖殿過節.
所以很多人來.
不過特別的地方是很多人都認識耶穌.
因為知道耶穌那時候.
其實已經出來傳道.
做上帝的工作都三年時間.
所以他有些名氣.
那時候很多人都認為耶穌是敬人.
是救主.
所以那時候在過節.
很多人眼光聚集的時候.
就用鐘樹枝來迎接耶穌.
不過聖經不只是停留在那裡.
講多兩節就說.
當耶穌呼喚拉撒路.
使他從死人中復活出墳墓的時候.
和耶穌在那裡的眾人就作見證.
眾人應聽見耶穌這行蹟就去迎接他.
原來不只是耶穌那三年傳道.
因為《約翰福音》記載.
耶穌剛剛進入耶路撒冷之前.
就使一個死了的拉撒路復活.
我想是很大的新聞.
如果今天有這麼大的神蹟.
大家都會很留意那個人來了.
先看看他是誰.
即使之前沒有認識耶穌.
剛剛有新聞說那個人.
厲害到連死人都復活.
不如去看一下.
這是另一件事.
不過在這裡又不是每個人都歡迎耶穌.
法利賽人就說.
你們看你們一事無成.
世人都隨著他去了.

$^{121}$在那裡工作的時候.
那時候的宗教領袖.
法利賽人一直都不喜歡.
覺得耶穌太受歡迎了.
他們這班宗教領袖沒有人理會.
所以一直和耶穌有很多的爭拗.
到了那個時候.
你看到耶穌竟然有一班人.
用鐘樹枝歡迎他.
也使得一個死人復活.
很多人都很有興趣的時候.
這班法利賽人就說.
你看他現在很成功了.
你們即是說他們那班人.
每個人都跟了他.
但聖經都不停留在這裡.
繼續說些什麼.
按錯了 下面.
然後聖經再說什麼.
第二節就說.
那時上來過節禮拜的人.
有幾個希臘人.
他們來見加利利的白菜大人肥力.
請求他說.
先生我們想見耶穌.
甚至這裡說到.
不單止猶太人.
有些外邦的希臘人.
連耶穌的名聲外邦都聽到.
在這樣的狀況裡.
耶穌就說一句.
人子得榮耀的時候到了.
所以我們見到.
耶穌不是無端端說一句話.
跟門徒說的.
而是原來耶穌工作三年以來.
一路的鋪排.
越來越多人認識.
令到拉薩路復活.
到那些猶太人 法利賽人.

$^{161}$拒絕到一個地步.
快要將耶穌殺死.
外邦人都聽到的時候.
耶穌強調的是.
我得榮耀的時間到了.
他觀察到原來耶路撒冷.
一整城人這樣對他.
他知道時候到了.
時候到了這句話.
耶穌不是第一次說.
耶穌出來傳道的時候.
在其他福音書.
馬學福音那些會記載.
耶穌那時候出來.
做一些醫治的事情.
有些人說那個耶穌很厲害.
快點捧他出來.
耶穌那時候說什麼.
我時候還沒到.
還沒到.
但是耶穌到了這個時候.
終之主一的時候說.
時間到了.
這裡告訴你.
原來耶穌做事.
或者上帝做事.
其實是有一種時間.
耶穌知道.
當那麼多人的眼光.
注視在他身上的時候.
他最重要的任務.
受苦受死.
釘十字架復活.
時候到了.
在這個時間.
他一定要做這件事.
這個是上帝給他最大的任務.
他的時候到了.
所以我們發覺.
到了最後的時候.

$^{201}$這段經文之後也是.
耶穌其實上十字架.
每個人都不想的.
不過你看到後面第27節.
耶穌就說.
我知道時間到了.
我避不了.
我推不了.
但是上十字架要受苦.
耶穌也不願意.
我現在要做什麼.
我苦啊.
救我脫離這時候.
但我正是為這時候來的.
所以我們發覺.
原來耶穌是這樣.
我想我們也是.
原來我們生命裡面.
上帝是掌握我們的時間.
原來上帝說.
到了一個時間.
不是你選擇的.
當事情發展到這個時候.
給你的任務就很重要了.
這件事也提醒我們基督徒.
原來我們的生命掌握在神的手裡.
我們的問題就是.
究竟今天上帝給了我們什麼任務.
不僅是耶穌.
其實我們在教會的聖經歷史裡面.
我們也看到有這樣的人出現.
大家不知道有沒有聽到.
在舊約裡面.
有一個叫以斯帖的人的故事.
熟悉的弟兄姊妹就會知道.
以斯帖是波斯皇帝.
他是一個猶太人.
不過他是一個波斯王的皇后.
波斯是在哪裡呢.
現在大眼罩的伊拉克.

$^{241}$改名做伊拉克.
不是伊朗.
差點說錯了.
伊朗.
現在大眼罩的伊朗.
波斯是一個很強大的國家.
話說歷史裡面有一段時間.
他攻佔了很多民族.
包括以色列人.
所以以色列人在波斯王統治之下.
其實是很危險的.
為什麼呢.
因為他被人管治.
如果皇帝不開心趣事.
就殺光那些人都可以.
你知道是這樣.
很危險的時間.
而那個時候真的.
在波斯王的時候.
有一個皇帝聽到奸人之言.
說要殺光那些猶太人.
很危急的時候.
那個以斯帖是一個皇后.
他就是波斯王的皇后.
但是那個時候.
那個皇后都不是可以隨便跟皇帝說話.
不過有另一個猶太人.
是他的親戚.
就跟他說了一句話.
鼓勵他.
鼓勵他你在這個時間上世給你的任務.
就是要為到以色列人說話.
他這句話是什麼.
他說此時你若閉下不言.
猶太人必從別處得解決.
得夢拯救.
你和你父家必自滅亡.
焉知你現在這個位份.
不是為了現今的機會.
時間是一樣的.

$^{281}$上帝對於信他的人.
在時間上有一個任務.
你想想那個時候.
這個皇后.
做帝國皇帝的皇后.
其實真的很好.
很多好處.
榮華富貴.
生活無憂.
但沒想過在這個時間上.
上帝有話要跟他說.
其實突然之間我想.
回到我們那個環境裡面.
其實我想我們某程度上.
做香港的基督徒其實都很有福.
可能大家各個人的狀況不同.
不過起碼我們可以有一個.
自由去信耶穌的環境.
我們的城市某程度上.
都叫做富足.
相比很多的地方.
有沒有想過究竟.
我們的神給我們這個時刻.
有什麼任務呢.
有時候個人對基督徒來說.
我們領受上帝那麼多的恩典.
好像以色列皇后.
但有個時間他就要用我們.
可能對基督徒來說.
有時候我們都會面對這些掙扎.
例如如果你工作的環境裡面.
好好的.
突然間發覺.
你公司有些道德上的問題.
究竟我去不去做一個正直的人.
不到你選擇.
你可能好端端沒事.
在那些時刻.
究竟你知不知道上帝對你.
有些什麼心意呢.

$^{321}$不要只是說個人.
社會群體都是.
幾年前疫情的時候.
都差不多五六年前.
那時候我有個朋友.
他自己是一個傳道人.
不過他不是在西北.
他在油麻地那裡蒙會.
油麻地是另一個香港很窮的地方.
裡面很多基層的人可能住板間屋.
住tong 房.
守庭口庭那種.
做勞動工作那些.
他接觸街坊.
疫情的時候立刻有一個極大的困難.
平時上班.
譬如我都很有福.
平時上班起碼有積蓄.
那時候最多就是家裡工作.
辛苦一點.
但對於很多人來說.
他立刻明白守庭口庭.
沒得吃.
也對於疫苗.
對於物資很缺乏.
他那時刻很知道他的教會.
在座落在油麻地.
他不能避開.
他聯同一頂尖會.
聯同其他教會一起去.
我們找一些物資.
我們可不可以除了物資.
除了幫助一些防疫物品之外.
我們可不可以甚至食物都可以找到.
我們可不可以有罐頭.
可不可以怎樣去探望那些街坊.
讓他們很困難裡面得到一些幫助.
就因為他那麼激動幫助到人.
而慢慢慢慢這些受助的人.
也都明白原來他們這班弟兄姊妹.

$^{361}$這班基督徒.
慢慢慢慢認識信仰認識教會.
後來這位弟兄在油麻地區.
做了一個機構.
很感動的地方.
真的沒有想過.
還是疫情時突然想到怎樣自保.
原來上帝在那一刻對你.
對我們有一種任務.
我想究竟今天的香港怎樣呢.
或者你們在洪水橋又怎樣呢.
我經常聽到新聞說.
北都很多發展.
很感恩我跟牧師.
跟弟兄姊妹或者你們的教會領袖都是.
我們知道這裡會有很多的發展.
在這裡將來會有很多的轉變.
但當上帝擺我們在這裡的時間.
可能我們建立教會.
可能你信耶穌的時候.
你都沒有想像過會發生這樣的事.
好像耶穌那樣.
好像以基帖那樣.
不過對你來說.
對我們教會來說.
上帝給我們什麼任務.
怎樣在改變了的環境裡面.
去傳福音.
你的人不同了.
跟以往不同了.
狀態不同了.
究竟我今天可以做些什麼.
我想首先是求天父開我們的眼睛.
好像我在油麻地的朋友.
天父開了我們的眼睛.
看到當地的需要.
就起來服侍.
成為一個榮耀的見證.
這個是主耶穌.
都是看著那個時間.

$^{401}$把握機會做上帝的工作.
其實我們想像那時候.
耶穌在中之主日進入耶路撒冷.
好像一個明星.
今天好像一些追星.
很多人要來見祂.
又用中之去歡迎祂.
然後又說這個勁人來看看祂.
很多注意力在祂的身上.
大家知道耶穌剛才有說過.
得榮耀的時間.
其實原來是要死的.
對於那時候很多人還沒認識耶穌.
甚至門徒都不太清楚為什麼.
你明不明白.
記得耶穌跟門徒說.
我要去耶路撒冷受死.
彼得都說不要.
門徒都還沒明白.
所以耶穌知道很多人看著祂.
耶穌也知道一個星期之後.
祂要上十課.
同一班人要看著祂死.
祂要開始預備那些人.
去解釋給人聽.
為什麼榮耀要死.
所以後面那句話.
其實祂是跟眾人宣告.
解釋為什麼基督徒.
會把榮耀和死放在一起.
接著那句話其實很特別.
我們想清楚一點.
祂怎麼說的.
我們一起讀廿四節.
廿四節讀到粗體紅色.
我實實在在地告訴你們.
一粒麥子不落在地裡死了.
仍舊是一粒.
若是死了就結出許多子粒來.
平常讀《慣性經》真的不要覺得什麼.

$^{441}$不過這句話耶穌將自己比喻為一粒麥子.
原來是很重要的一個比喻.
我們想一想.
我們說榮耀.
榮耀通常說什麼.
什麼叫榮耀.
譬如你很辛苦讀書大學畢業.
畢業那一刻就有榮耀.
又或者我們工作了一輩子.
到退休的時候那些叫什麼.
榮休.
牧師榮休.
很多人就說.
為什麼榮耀.
因為做得很好.
當你做那件事做得很成功.
大家都認同你的成功.
這樣就叫做榮耀.
為什麼耶穌說榮耀是死的.
因為原來祂說自己是麥子.
說自己是麥子.
麥子其實對於我們來說.
麥或者米.
其實什麼叫做一粒好的麥子.
和一些好的米.
麥和米是用來做什麼的.
給人吃的.
不是像花一樣擺著給人看的.
所以原來為什麼要死.
是因為一粒麥子.
祂做得好.
祂滿足了祂存在的目的.
就是給人吃的.
所以這是耶穌說的話.
為什麼用麥子去解釋自己.
因為原來祂說.
做天父的工作.
信這位上帝的人.
原來在世上要成為一粒麥子.
給其他人吃得到好處.

$^{481}$讓其他人得到好處.
這樣去想.
其實對我們世界上世俗的人.
或者香港人來說.
其實沒有人會覺得成功的是一粒麥子.
是不是.
你想想對香港人來說.
什麼叫做榮耀.
什麼叫做做得好.
出生那一刻.
或者進入幼稚園那一刻.
都已經叫做什麼.
大家都說得很厲害.
贏在起跑線.
其實今天說多少都好.
這些都說不出來.
這個社會裡面說的.
第一天出生就跟別人爭了.
我兩個小孩十幾年前出生的時候.
帶著肚子碰到朋友.
說你沒有了幼稚園味.
那些氣氛大家都很明白.
我們都覺得很慘.
但是就是這樣.
我們會覺得怎樣.
我們會覺得原來什麼叫榮耀.
就是你考第一.
你入到名校.
第二天你找到一份很有錢.
很有前途的工作.
而背後是說我打低其他人.
我成為出眾的人.
原來我們今天的社會.
是告訴你怎樣做到墨子.
不做墨子應該做什麼.
做獅子.
我們其實全部人都覺得.
你做萬獸之王就厲害了.
我們都是這樣稱讚人.
你多厲害.

$^{521}$什麼虎父無犬子就是那種.
你明白嗎.
所以這麼難接受.
我想清楚耶穌這句話.
他說我得榮耀的時候到了.
接著就說墨子.
對人來說是多麼的震撼.
我們是不是都這樣.
當我們基督徒.
我們是不是看我們的生命.
真的有了墨子呢.
還是我們仍然追求到.
做一隻獅子.
其實想想聖經裡面.
對於門徒的比喻.
除了墨子之外還有什麼.
葡萄.
還有什麼.
果子.
還有什麼.
羊.
我們煮的羊.
特別的地方全部都是食物.
你有沒有發覺這個東西.
很有趣.
全部都給人吃的.
全部都是祝福這個世界.
沒有什麼叫我們做一個獅子.
做一隻老虎.
做一條鯊魚去殺人.
沒有的.
所以耶穌其實跟我們說.
當我們信耶穌的時候.
我們眼光要改變.
我們追求的東西要改變.
我不是說讀書好 入到名校是差.
我沒有這個意思.
是當我們可以做到這些東西.
當我們賺到錢.
當我們可以成功也好.

$^{561}$我們為了什麼呢.
我不是說不需要卓越.
我不是說不需要努力.
我們努力為了什麼.
我們努力是令自己成為一個.
很好的墨子.
成為一個很有營養.
很好吃的麵包給人吃.
原來這個就是耶穌說.
他死的意思是什麼.
他要犧牲自己.
捨棄自己.
去讓其他人得到益處.
這個就是做基督徒不同的地方.
教會耶穌.
過去二千年.
頭三百年左右.
教會很困難.
很多Big Bang.
大家可能都聽過.
頭幾百年在羅馬帝國裡面.
很多的迫害.
很多基督徒因為迫害而死.
不過最特別的地方是.
原來頭三百年.
同時是教會發展得最好的時間.
很好的時間.
福音傳得很快.
為什麼呢.
那時候人們見到.
原來信耶穌被迫要死.
為什麼要信耶穌呢.
為什麼有什麼吸引力.
很多歷史研究做出來.
不過有一樣東西.
他說很重要.
有一個因素為什麼吸引人信耶穌.
在那時候.
教會初期那幾百年.
都好像今天這樣.

$^{601}$很多的疫情.
很多瘟疫.
你知道過去都是沒有醫藥的.
一有那些疫情.
就會很快很多人死.
那時候人們都知道.
其實瘟疫就是傳染.
所以人們在一個地方.
發覺有瘟疫的自然反應是什麼.
就跑去那地方.
一定是這樣的.
但基督徒很特別.
基督徒是一群正正有瘟疫出現.
就願意留下.
去服侍那些跑不掉和生病的人.
其實不等於一定個個有神蹟保守基督徒.
都可以因為染病而死.
但是這件事讓羅馬帝國的人很好奇.
為什麼這群人可以這樣.
當個個走的時候.
他們寧願犧牲自己.
都服侍那些病了的人.
那些病的人也不是有錢有權的人.
為什麼這樣.
他們發現原來對基督徒來說.
他們知道永生的盼望已經得到了.
他們不害怕死.
他們不需要意外.
他們發覺原來.
跟從耶穌就是要服侍世界.
服侍有需要的人.
這件事讓很多人發覺.
這個信仰很特別.
對有興趣.
發覺原來基督徒愛的群體彼此相愛.
有永生的盼望.
吸引人.
很多人信耶穌.
其實我們都會問自己.
今天其實同樣的事.

$^{641}$同樣的道理都是一樣.
如果我們很渴望在社區裡面工作.
很渴望讓人見到耶穌基督.
原來最有力的見證.
就是一個服侍的見證.
不是看自己為重要.
不是做隻獅子去到哈哈爸爸.
而是成為墨子.
祝福這個世界.
祝福倫社.
早期的教會.
成為一個非常美好的見證.
讓很多人認識耶穌.
今天我們是不是都同樣去做呢.
不過困難的.
剛才都有講.
耶穌說這麼多人看著我.
我得榮耀的時候來了.
我們說耶穌在受難的時候.
祂自己很多的掙扎.
因為耶穌說父的旨意還是我的旨意.
其實當然.
我想耶穌的掙扎就是.
究竟祂是不是聽父對祂的旨意.
不過其實耶穌有另一個壓力.
祂發覺全城人看著祂.
你明不明.
祂怎樣去選擇.
就決定究竟祂是不是真的為父親作見證.
或者沒有作見證.
所以你看看那句.
耶穌說我現在心裡憂愁.
我說什麼才好.
說父啊.
救我脫離這時候嗎.
你看到耶穌是真正的掙扎.
對祂來說.
祂不是輕輕鬆鬆.
祂知道自己要受苦.
要痛.

$^{681}$要被眾叛親離.
要被羞辱的釘十字架.
是很難的事情.
不過祂也明白.
但我正是為這時候來的.
父說願你榮耀你的名.
祂說我願意生命去榮耀你.
於是有聲音從天上來.
父去答應祂.
去說一句話.
說我已經榮耀了我的榮.
還要再榮耀.
這句話是什麼.
天父跟耶穌說.
你是我兒子.
你的生命已經榮耀了我.
你已經跟從了我.
你已經為我作見證.
還要再榮耀就是說.
你願意去上十架跟從我.
是要更加榮耀天父的名.
所以原來今天我們在地上.
怎樣去選擇.
是不是做墨子呢.
跟神的榮耀是有關係的.
原來不只是自己的事.
是跟上帝的榮耀是有關係.
這也給了另一個動力.
我們再發覺.
世界的人未必看得起我們.
但是最重要的是.
我們天上的父看得起我們.
而天上的父真是最終掌權的位置.
我們基督徒也知道.
今天在地上受苦.
為了服侍人.
為了信仰的緣故.
受逼迫也好.
將來當上帝真正地上掌權的時候.
我們會得到榮耀.

$^{721}$得到福樂.
這些痛苦是短暫的.
剛才說初期教會.
很多的逼迫.
很多的困難.
很多人死亡.
甚至很多人因為.
不是個個殉道.
不過也有很多折磨.
令到可能因為信仰的緣故被迫害.
所以身體傷殘.
有這樣的狀況.
所以那時候的人就很渴望.
我為了信仰的緣故受那麼多苦.
可能身體傷殘.
可能沒得隻手沒得隻腳.
被逼迫.
將來當我復活的時候.
當耶穌再來的時候.
究竟我這些傷殘會不會復原呢.
那時候他們的答案是會的.
因為耶穌來工作的時候.
他醫病趕鬼.
他將到盲的醫好.
將到聾的醫好.
所以一定將我們那些傷殘.
不好的.
因為信仰的緣故.
因為逼迫的緣故.
那些傷殘全部醫好.
這個是盼望.
不過帶來一個問題就是.
大家知道在聖經裡面記載.
耶穌復活之後.
他手上仍然有釘痕.
他肋旁仍然有傷痕.
所以那時候的人就問.
為什麼如果逼迫的緣故.
將來那些傷痕.
是會醫好.

$^{761}$為什麼耶穌的釘痕.
和肋旁的傷痕都仍然在.
你會怎樣回答這個問題.
很難回答.
神學家真的很厲害.
有個叫奧古斯丁的很出名的神學家.
他回答.
其實耶穌復活之後.
天父要醫好他的釘痕和傷痕.
多容易.
復活也可以.
醫好他的傷痕很容易.
不過他說.
耶穌身上那種釘痕.
是天父故意留下.
故意留下是什麼意思.
是讓我們看到.
那個是一個榮耀的記號.
我們因為見到耶穌基督為人犧牲.
所以讓我們見到這個.
人看似是羞辱的記號.
是成為一個.
在神眼中榮耀的記號.
我在這個講法提醒我們基督徒.
我們其實原來要.
用一個不同的眼光去看我們自己.
其實服侍上帝的人.
很多時候在地上真的未必有好處.
有時候我認識一些傳統人.
認識一些.
就算我學院.
神學院裡面的老師也好.
學歷其實很高的.
不過也很辛苦的.
租一些很小的地方.
家裡也住得不是很好.
有時候可能弟兄姊妹.
也認識一些宣教士.
去到一些很困難的地方.
身體上可能快點老.

$^{801}$香港那些又有保養.
又吃得好住得好.
沒那麼快老.
可能真的很多這些東西.
不過我想這樣的看法.
耶穌的釘痕提醒我們.
其實我們要看那些.
所謂受傷.
那些人看為很不好的東西.
不過為著服侍神的緣故.
我們看那些傷痕.
是一個榮耀的記號.
所以弟兄姊妹.
我們怎樣去看自己呢.
我們願不願意因為信仰的緣故.
我們可能沒那麼光鮮.
我們可能為了回應上帝的呼召的緣故.
我們住得沒那麼好.
我們吃得沒那麼好.
我們願不願意這件事呢.
無論你可能服侍.
無論你可能參加一些短宣.
甚至上帝呼召你.
去獻你的生命給祂.
服侍祂.
在香港.
在不同的地方服侍祂.
是的.
其他人不認識可能覺得.
嘩 什麼這麼蠢的.
這樣都可以.
但是原來在天父的眼中.
是一個最榮耀的記號.
所以當我們今天去紀念主耶穌為我們受苦.
去思想什麼叫榮耀的時候.
讓我們都好像耶穌有同一個眼光.
看看什麼叫榮耀.
成為一粒墨紙.
好好地將自己擺上.
服侍身邊的人.

$^{841}$叫到天父能夠得到榮耀.
我們一起去祈禱.
因為我們知道你愛我們的緣故.
即使我們仍然是罪人.
仍然是敵黨上帝的時候.
你都為我們犧牲.
天父的愛因為這樣向我們獻命.
耶穌記得我們.
甚至透過你的死去教導我們.
原來天父所看什麼叫榮耀.
和世界真的很不同.
讓我們都好像你一樣去看我們.
是讓我們知道原來珍藏你的時間.
放下自己的好處.
讓世界的靈蛇得到好處.
是一個最榮耀的事情.
求你讓我們深刻地去到.
不單止明白.
亦都將自己的生命去到獻上.
像耶穌基督一樣.
去看看今天天父對我們有什麼心意.
在此時此刻.
給我們什麼任務.
讓我們祝福人.
廣傳福音.
關愛論社.
榮耀天父.
聽我們祈禱.
禱告奉耶穌基督的名.
阿們.
\newpage



\section{哥林多前書 15:14-19}
\label{sec:ke2vMZ9uUT4}
\textbf{2026.04.05 《復活,不僅是認同》:哥林多前書15\_14-19(和修版) 曾裔貴牧師}
\newline
\newline
連結: \href{https://youtube.com/watch?v=ke2vMZ9uUT4}{\texttt{https://youtube.com/watch?v=ke2vMZ9uUT4}} ~~~~ 語音日期: 2026-04-07
\newline
\newline
\hyperref[sec:_kJpwvPJ95U]{\small{< < < PREV SERMON < < <}}
~
\hyperref[sec:index_chronic]{\small{[返順時目]}}
~
\hyperref[sec:index_scriptual]{\small{[返順卷目]}}
~
\hyperref[sec:y558B5MlYc4]{\small{> > > NEXT SERMON > > >}}
\newline
\newline
哥林多前書 15:14-19
\newline
\begin{longtable}{cl}
\hline
\hline
章節 & 經文 (和合本修訂版)\\
\hline
15:14 & \begin{tabularx}{0.7\textwidth}{X} 基督若沒有復活，我們所傳的就是枉然，你們所信的也是枉然。 \end{tabularx} \\ \\ \relax
15:15 & \begin{tabularx}{0.7\textwidth}{X} 這樣，我們甚至被當作是為神妄作見證的，因為我們見證神是使基督復活了。如果死人真的沒有復活，神就沒有使基督復活了。 \end{tabularx} \\ \\ \relax
15:16 & \begin{tabularx}{0.7\textwidth}{X} 因為死人若不復活，基督也就沒有復活了。 \end{tabularx} \\ \\ \relax
15:17 & \begin{tabularx}{0.7\textwidth}{X} 基督若沒有復活，你們的信就是徒然，你們仍活在罪裡。 \end{tabularx} \\ \\ \relax
15:18 & \begin{tabularx}{0.7\textwidth}{X} 就是在基督裡睡了的人也滅亡了。 \end{tabularx} \\ \\ \relax
15:19 & \begin{tabularx}{0.7\textwidth}{X} 我們若靠基督只在今生有指望，就比所有的人更可憐了。 \end{tabularx} \\ \\
[1ex]
\hline
\hline
\end{longtable}
$^{1}$請姐妹再一次在主裡面行安.
打開手機看看.
神感動兩位同工牧者.
在感宣永康牧師.
這個時間傳遞的訊息就是.
使徒三重漢見.
他引用的經文就是約翰福音的二十章.
就是第三到第十節.
另外同樣都是約翰福音.
不過陳欣宏牧師在嘉欣堂.
這個時候傳遞的訊息就是.
破格的平安和喜樂.
在今早大家在現場和在紅欣的直播.
我就會用哥林多前書.
神感動幾位的牧者.
都有不同的向度.
去將復活節這個很重要的意義.
和大家分享.
我想用一個例子.
一個真人真事.
但絕對不是一件開心的例子.
很深刻.
因為整個過程陪伴著這位姐妹.
經歷那種傷痛.
在我很多年前.
這位姐妹.
她那個時候是讀中學.
大概中四還是中五.
姐姐和我們是同一個團契.
那個時候我剛剛傳道.
我的母堂,母會.
這位姐妹是信主的.
跟姐姐回教會,都自小回教會.
但上到中四開始無心向學.
就開始跟朋友玩.
有一次帶了一個朋友.
男孩來的,應該是男朋友.
來教會聚會.
我們一見到這個男孩.
就已經勸他.

$^{41}$因為都熟的.
不要跟這個男孩在一起.
當然不聽.
男孩比她大一點.
已經在社會工作.
但大家都知道.
這個男孩可能有些不好的嗜好.
吸毒的.
但這個姐妹完全不察覺.
可能墮入了愛河.
看不到這些現實.
不聽.
後來甚至不再讀書.
也跟家人關係很差.
因為姐姐對他都很嚴厲.
叫他盡快離開這個男孩.
但不聽.
輾轉一段時間都過去了.
後來他肯定沒有再回教會一段很長時間.
到再見他的時間.
他已經有了一個寶寶.
輾轉一段時間.
小朋友去到三年級小學.
為什麼再回教會呢?.
姐姐叫我們去探望她妹妹的兒子.
即是兒甥.
因為小朋友有血癌.
在瑪嘉烈.
我們趕去跟他祈禱.
當然經過輾轉一段時間.
因為我們知道的時候已經是後期.
那種痛苦.
他自己認為.
絕對不是我們跟他說.
他自己認為是神對他的懲罰.
他決心離開這個丈夫.
因為真的知道他沒有戒毒不得已.
再變本加厲.
遊手好閒.
不務正業.

$^{81}$和這個兒子去到這樣的階段.
這個男朋友或這個丈夫.
就離開他.
他雙重的打擊.
最後很深刻.
我們深夜趕去醫院.
因為真的救不了.
這個很乖的小朋友.
他的名字叫阿謙.
很乖的小朋友.
血癌就在瑪嘉烈醫院.
安詳的離世.
當然也有同工.
跟這個小朋友講福音,信主.
後來有一段時間沒有再見這位姐妹.
過了一段時間.
因為她沒有再回我們教會.
可能不好意思,也不想.
見到她有一個新的婚姻.
有一位很好的弟兄愛她.
和他告訴我們.
當然因為真的很熟.
不用擔心我.
我在主耶穌裡面.
我有盼望.
阿謙在醫院的時候.
離世之前.
這個小朋友信主.
這個故事是我自己.
初出來傳道的時候.
很深的經歷.
其實這麼多年了.
在傳道的過程.
很多這些偶然的例子.
這些真人的故事.
讓我感受到.
復活不僅僅是世人的一種狂歡.
大家每天看一看.
第一天有多少十萬人離開香港.
很開心,很興奮.

$^{121}$每個人都拿著行李箱去玩.
這些本來不是錯事.
要減壓.
但是基督徒的復活節.
不僅僅是一種只有快樂.
因為我們確實知道.
在我們的人生.
甚至在我們每一天.
我們經驗到的.
絕對不是想像中的快樂,圓滿.
而是往往有很多很多的困難.
但反而在這些困難裡面.
我們看到的.
就是復活的生命.
我選取這段聖經.
其實全部是否定句.
是保羅要去駁斥哥林多教會.
認為沒有人會死了之後復活.
他們認同耶穌復活.
但是他們已經沉醉於.
聖靈的充滿.
要爭在哪個教會裡面的人.
哪些恩賜會大.
所以他們爭的是聖靈的充滿.
爭的是講方言.
爭的是其他的東西.
甚至第一章.
他們已經是分門結黨.
他們沉醉在教會裡面.
這些分裂裡面.
哪個人是屬於哪一個.
所以他們根本不會太介意.
真的會有死人的復活.
所以保羅在這段聖經.
其實是很重的責備他們.
而保羅很直接了當.
來告訴我們.
是一個很重要的本體的生命.
復活了.
所以每一個信主.

$^{161}$在耶穌裡面相信的人.
其實就在這個復活的永恆真實的事件裡面.
我們就要找到我們人生的盼望.
所以.
我比較少寫這麼詳盡的講章.
有差不多七頁.
一般我就是做part point(部分點).
但是因為這幾天真的忙.
還有昨天有數目.
不夠時間去整理.
十四節.
如果這個復活沒有的時候.
保羅是在駁斥一些空洞的.
一些講道德的寓言.
一個令人很鼓舞的哲學故事.
是沒有改變整個宇宙的生命的力量.
所以也不能夠真正面對死亡和虛無.
所以如果沒有復活.
我們的講道只不過是一場情感充沛的演講.
跟一個沒有信仰的演講.
可能差不多.
有些演講者.
他們的感染力.
可能比我們講道的人.
來得更加有感染力.
第十五節.
也是駁斥那些虛假,虛無的見證.
甚至我們被當作是為神妄作見證.
所以使徒們可能比KK園區詐騙的地方.
更加的去欺騙人的騙子.
因為如果沒有復活的時候.
使徒居然用他們的生命去殉道.
圍著一個虛假的事實.
來這樣去死.
騙了歷史歷代很多的人投入在一個不實在的信仰裡面.
所以保羅說這是一個虛構的.
如果沒有復活.
但是整個教會.
我們建立在主耶穌的復活裡面.
第十七節.

$^{201}$因為十六節.
如果基督沒有復活.
就不會再有死人的復活.
十七節說.
駁斥的就是無效的信心.
因為你們的信就是謊言.
只是一種心理的作用.
自我的安慰.
不能夠觸及我們心靈最深處的困境.
所以保羅說.
我們仍然在一個罪惡的困鎖裡面.
不能夠有任何的釋放.
意味著罪的權勢.
死亡的毒鉤根本就未打破.
這個告訴我們.
保羅是在駁斥一些.
如果你是沒有信心的基督徒.
你只是認同基督還他.
第十八節.
駁斥一個絕望的無助的滅亡.
就是在基督裡面已經睡了的人.
都滅亡了.
這個就是一個真的很恐怖的絕望.
我近這幾年.
每次主禮一些安息禮拜.
內心裡面越來越有一種立體感.
我們幫他辦理的弟兄姊妹.
或者弟兄姊妹的家人.
他安躺在那個棺木裡面.
我們要來到去送別最後一程.
但是牧者我自己個人的內心.
是越來越感受到一種的立體.
他真的只是安睡.
那個棺木進入那個火化.
變成為一堆的灰.
但是那個我們所愛所認識的朋友.
弟兄姊妹.
他在主耶穌的救恩裡面.
不是一個絕望的永別.
是暫時的分開.

$^{241}$而來得著這個永恆的復活.
所以十八節保羅已經是很斬釘截鐵.
駁斥如果沒有復活.
任何一個在主裡面死的人.
無論是我們中國人說的好死.
不好死.
慘死.
都沒有意義.
都是滅亡.
是很重的駁斥.
任何一個人.
如果主耶穌沒有復活.
我們己有的信念.
在基督裡面睡的人.
都是滅亡.
這一切是全部駁斥那些.
不相信復活的事件的人.
最後第十九節.
保羅告訴我們.
我們在今生比任何一個人都可憐.
我們若靠基督.
旨在今生有指望.
就比所有的人更加可憐.
為什麼要這樣描述這個比較.
是不是因為保羅特別要受到羅馬的政權.
彼拉多和之後的官員.
對他的鞭打.
對他的追逼.
和猶太教他自己的以前的同伴.
他自己的一些法利賽人.
一些的祭司長.
一些在聖殿裡面.
大家一起服侍的拉比的同伴.
成為了敵人.
很好的朋友成為了攻擊保羅的敵人.
而令他感覺到.
他如果只是今生有一個指望.
因為都有一些年青的外邦的人.
跟隨保羅去傳道.
是不是就是這樣可以有一個指望呢?.

$^{281}$保羅說如果這樣的指望.
如果真的沒有基督復活.
我們這樣去幫助其他的年青人都沒有用.
跟隨他去傳道也都沒有用.
無論身邊有提摩太.
有推基古.
有路加.
有其他的人.
有提多等等.
都沒有用.
因為今生的指望.
不能夠帶來永恆的盼望.
比眾人更加的可憐.
五個駁斥.
對於不相信復活.
所以今日的主題是.
我們不僅僅是認同.
我們每一個星期講的.
我信身體復活.
這個信念.
這個真理的宣告.
不僅僅停留在一個認同就夠.
是要我們每一個弟兄姊妹.
由我們知道原來我們會有永恆的盼望.
原因是因為耶穌基督這一位的造物主.
來這個世界.
而祂勝過死亡.
勝過魔鬼.
勝過罪惡之後.
就使我們每一個在祂裡面的基督徒.
我們就有永恆的復活的盼望.
所以原來由我們信那一刻.
我們那個盼望.
那個復活的能力.
就已經在我們的生命裡面.
來到開始.
這個開始.
當然不是好像一無掛慮.
沒有任何的痛苦.
沒有任何的難處.

$^{321}$不是.
嘗試用一個例子.
有一位母親失去了她的兒子.
那個兒子是很年輕就離開了她.
對這個母親來說.
因為她是有信主.
她是有相信永恆的盼望.
她這樣的描述.
她的悲傷是真實的.
每一天好像都被那個空洞吞噬.
她的兒子失去了.
但是復活的生命告訴她.
不是立刻就拿走她的悲傷.
而是使那個悲傷好像埋藏在肚子裡面.
種下一個種子.
所以改變不是不再難過,不再悲傷.
她說在她兒子的忌日.
她仍然會哭泣.
但是當她開始整理她兒子的遺物的時候.
她說將一些她還能夠用的.
這些兒子的遺物.
送給有需要的家庭.
她就這樣來跟自己說這個安慰.
我兒子的一部分.
還能夠令到這個世界上.
可以去愛有需要的人.
沒錯,你可以說這是一個信念.
但是它就是基於復活的盼望和事實.
悲傷沒有減少.
仍然充斥著在她內心對兒子的掛念.
也不會因為復活的盼望.
兒子在她內心就慢慢地消失.
不是,繼續有這一份對兒子的懷念.
但是換來的.
就是這個盼望的種子埋藏在那個土壤裡面.
讓這個盼望成為更加具體的將來的重聚.
弟兄姊妹.
你有沒有你所愛的人離你而去?.
你要面對這種生離死別.
面對這種失去的撕裂.

$^{361}$內心的空洞.
好像吞噬一樣.
有,因為我們正常的人.
活在一個生老病死的世界.
有時候我們看到的.
比香港可能正常的疾病.
長期的慢性的疾病來得更加慘烈.
我們在幾個月之前.
親眼目睹著.
榮福縣的大火.
背後可能有很多現在在調查當中.
我們不知道的問題.
是不是浮現一些結構性的.
有結構有系統的罪惡.
都不知道.
更加荒謬的.
現在在伊朗,以色列,美國的戰爭.
俄烏的戰爭.
在這個世界好像還沒有任何的平息.
平安的時候.
我們每天看到.
是不是令到我們都已經好像麻木一樣.
但是親愛的弟兄姊妹.
如果我們是有盼望.
確信我們就是在這個永恆復活裡面的其中一份子.
我們就仍然面對著這些的地上.
悲傷,罪惡,痛苦.
而且有時候很靠近的.
撕裂的痛苦的時候.
我們仍然是有那個永恆盼望.
因為基督已經從死裡面復活.
基督已經是勝過了死亡.
我又用另一個例子是近期.
去年.
一位弟兄的爸爸是很原居民的一些拜偶像.
拜得很厲害.
他很難信主的.
因為他的身份,他的位份在村裡面.
令到他根本上是很難.
包括他自己的迷信觀念.

$^{401}$都很難信主.
兒子都不敢跟他講福音.
叫教會去探望他.
因為真的會罵得很厲害.
但是到了一個階段.
那個病已經很嚴重了.
太太勸丈夫.
你還不叫牧師去探望你的爸爸?.
已經這麼差了他的身體.
他真的好像勉為其難一樣.
叫我去探望.
你可不可以去?.
不過你記住我爸爸是原居民.
我去到的時候.
他那些原居民的村的叔伯兄弟.
都在病房外面等待.
第一次去到天水圍醫院.
世伯叫他信主.
很清楚.
灑手伶頭叫我快點走.
當時有一點點的感動.
不要緊.
不要急.
不如我為你的病祈禱.
他勉為其難.
點了點頭.
為他祈禱.
祈禱後.
都會聊天.
世伯祈禱你又肯了.
不如信主吧.
他的手已經綁了儀器.
拇指都包了.
馬上大力地搖.
還用他僅有的微弱聲音.
說這樣不行.
你說我祈禱OK.
原來你哄我.
然後叫我信主.
真的很堅定地拒絕.

$^{441}$如是者.
再過一個星期.
兒子就說不行.
一定要來.
因為已經很差.
要再去.
很奇妙地.
這次.
祈禱後.
再跟世伯說.
不如真的信主.
你現在這麼辛苦.
不知道什麼原因.
他當時已經沒有說話的能力.
他居然.
很清楚地.
捉緊我的手.
給我一個最清楚的表示.
你捉緊我的手.
如果你能夠.
你就搖搖頭.
他搖搖頭的能力都不夠.
但居然很大力地.
捉緊我的手.
很清楚.
很堅定.
而且當時他的叔伯兄弟.
已經在床邊.
圍村裡.
這樣來信主.
見證到復活的能力.
改變一個這麼重的迷信.
他要在.
他的村裡.
辦很多.
打造儀式.
他是主事的人.
這樣信主.
所以他的喪禮.
安息禮拜.

$^{481}$在圍村裡.
先准許我們做安息禮拜.
然後他們圍村.
做他們的儀式.
是一個很大的.
見證.
弟兄.
見證著爸爸.
經歷到這個死而復活.
生命的改變.
弟兄姊妹.
我們有很多.
這樣的經驗.
和經歷.
陪伴著不同的家庭.
面對不同的人生.
的遭遇.
很肯定就是.
信了主.
真的不會令到你.
完全沒有任何的苦痛.
沒有任何的困難.
但是.
我們今天的訊息.
不僅僅是認同.
基督的復活.
而是我們要講.
我們每一個弟兄姊妹.
無論你現在的.
情況.
是怎樣的光景.
有弟兄姊妹可能婚姻上.
有時很多的困難.
都很難和別人.
講他內心的.
痛苦.
婚姻變成一個.
好像囚禁著一樣.
每個人都有自己.
不同的故事.

$^{521}$不過信仰的.
寶貴和真實.
就是你真的可以將你的生命.
投向基督.
而他慢慢會給我們.
很大的力量改變.
勝過死亡.
勝過.
我們自己舊有的舊惡.
好像剛才.
靈暢的弟兄說.
我們在基督裡面.
這個新的生命.
吞噬罪惡.
我們自己的舊惡.
我們可以宣認.
我們與基督同釘十字架.
活著的.
不再是我們.
是基督在我們裡面.
復活.
有一個.
愚道故事不是真的.
有一個.
村的.
一個家長.
有三個兒子.
有一天.
這個村長.
就要去到遠方.
做一些買賣.
和他三個兒子說.
他的大兒子.
種一些橄欖樹.
第二兒子.
就是做牧羊.
養羊.
第三個兒子.
是一個藝術家.
表演歌舞.

$^{561}$所以每逢.
村裡面的節慶.
這個第三個兒子.
就會做一個領導.
引導村民.
籌備節慶.
跳舞.
或者辦一些禮儀.
這個家長.
這村長去了之後.
就和他三個兒子.
臨去之前就說.
你們要照顧我們這條村的村民.
有什麼需要.
盡量要配合那個需要.
來支持他們.
我很快就會回來.
但是爸爸一離開之後.
那年的冬天.
非常之嚴寒.
而且慢慢那些.
漫天的風雪.
令到那個村都受到很大影響.
到了一個階段.
他們那些木材的燃料.
都沒有了.
天氣實在太冷了.
所以那些村民.
就來和他們幾個說.
能不能夠.
將村裡面.
這一棵很重要的.
橄欖樹都斬了.
讓我們暫時可以渡過.
這個寒冬.
新一年我們在春天之後.
再種植.
大兒子起初不肯.
因為這一棵橄欖樹.
很有感情.

$^{601}$也有生長的橄欖.
給他們有一個食用.
但是無奈.
真的非常之大的雪.
所以呢.
大兒子都沒有辦法.
就斬了它.
將那些柴火.
成為大家在冬天的時候.
僅有的保暖.
過了一段時間.
就完全停止.
春天也沒有來.
爸爸也沒有回來.
已經沒有足夠的糧食.
去和其他的村.
也去不到其他的村.
來交換食物.
所以呢.
僅有的羊.
那些村民就說.
能不能夠讓我們在最後的時間.
渡過這個嚴寒的冬天呢.
又是輾轉.
沒有辦法.
將這些羊.
來成為大家在寒冬的時候.
最後的食物.
如是者.
那個爸爸.
回來了.
但還沒回來之前.
因為那個冬天實在太漫長.
很多的村民.
都忍不住.
大家的內心都很冰冷.
不知道怎麼辦.
就希望能夠盡快離開.
去到其他地方.
還可以得到.

$^{641}$那個拯救.
爸爸終於在春天.
回來.
看到.
只是自己的家庭.
有煙火.
很多的村民.
都已經離開了.
去了其他地方.
就問為什麼.
第三個兒子.
剛才說的那個.
藝術家.
在整個村裡.
有什麼節慶呢?.
就是籌備那些儀式.
令到村裡的人.
有那種盼望.
這個第三個兒子.
很開心看到爸爸回來了.
就很想.
歡迎爸爸跳舞.
爸爸就問.
大兒子和二兒子.
也說了他們的困難.
將他們僅有的.
無論是橄欖樹.
還是家中的羊.
都給村民用.
但是爸爸問.
第三個兒子.
為什麼在村民.
心靈最需要盼望.
最需要鼓勵的時候.
你不來.
用你的藝術.
用你的舞蹈.
用你的恩賜.
來鼓舞他們的心靈.
讓他們的心靈能夠有盼望.

$^{681}$去支持寒冬呢?.
第三個兒子.
無言以對.
因為他真的沒有在村民.
最需要盼望的時候.
提供給他們.
有這個鼓勵.
我們走開.
去到不同的地方.
我們會見到不同的人.
的需要.
好像我們上一個月.
三月二月底.
我們去到菲律賓.
見到有很多活潑的小朋友.
他們還未被世界的困難.
社會的罪惡.
來污染.
還是很開心很活潑.
見到我們就抱著我們.
不是想要錢.
不是想要物質.
很想是一個對他們的關心.
這樣而已.
很多的人.
需要不同的鼓勵.
生命的意義.
但是基督徒.
我們要給予人的是什麼呢?.
如果我們自己只是認同.
這個復活.
而我們自己沒有參與.
在耶穌的救恩裡面.
我們沒有在主耶穌的.
死而復活的生命裡面.
我們去投入.
我們自己的生命.
準備我們自己整個人.
為著福音.
復活.

$^{721}$這個世界見不到的.
是基督的復活.
帶給人的新生命的意義.
第三個兒子.
他本應在人的心靈最冰冷的時候.
來鼓動他們的心.
將盼望.
放在他們的心靈裡面.
所以.
一家人在這個時候.
悲哀哭泣.
因為所有的村民.
都離他們而去.
這是一個預道故事.
是告訴我們.
那個一千銀子.
沒有用的僕人.
是又懶.
又惡的僕人.
他就要面對.
主人回來的時候.
來問他.
為什麼.
你將你的恩賜.
你的才幹.
收起來不用呢?.
很盼望我們這一段.
很簡短的訊息.
成為我們來投入.
所以.
復活不是一個.
句號.
它只不過是一個冒號.
一切都更新了.
死亡.
你所謂的終極的武器.
已經被.
基督的復活所廢去了.
罪惡.
好像是堅固的.

$^{761}$鎖鏈.
但是已經被打破這個權勢.
所以.
我們不僅僅是認同一個.
的教義.
我們是活在復活這個事實.
的光芒裡面的.
每一個.
元主祝福大家.
我們一起祈禱.
親愛的主,多謝你.
因為復活的開始.
是由你.
來將這個永恆的生命.
帶進.
這個本來是罪惡的世界.
罪惡.
不能夠再狂傲.
因為耶穌基督.
你已經勝過了死亡.
就算我們看到.
有很多不公不義.
好像是結構性.
好像是.
永遠無法打破.
但絕對不是.
基督的復活.
已經告訴我們.
我們每一個小小的善行.
為主.
做在主裡面的.
每一個善行.
在永恆裡面.
都是一筆帳.
已經被記錄.
因為我們就是為主.
為復活的主而做.
也都是宣認這個世界.
再不是魔鬼掌權.
再不是罪惡掌管.

$^{801}$我們每一個基督徒.
我們見證和行使.
基督的復活能力.
盼望我們大家一起.
在整個復活節.
我們去思想.
真的復活的生命.
如何流露在我們.
軟弱的心靈.
我們這個平凡的人.
因為我們從來.
都是平凡的人.
但是基督已經.
住在我們的內心.
所以我們能夠有力量.
面對不單純是明天.
而是每分每秒.
我們都是來彰顯.
基督復活的能力.
在我們的生命裡面.
我們獻上感恩.
奉主耶穌基督的名.
阿們.
祝福各位.
\newpage



\section{哈巴谷書 3:16-19}
\label{sec:y558B5MlYc4}
\textbf{2026.04.12 《黑夜中的信心之歌》:哈巴谷書3\_16-19(和修版) 曾敏傳道}
\newline
\newline
連結: \href{https://youtube.com/watch?v=y558B5MlYc4}{\texttt{https://youtube.com/watch?v=y558B5MlYc4}} ~~~~ 語音日期: 2026-04-14
\newline
\newline
\hyperref[sec:ke2vMZ9uUT4]{\small{< < < PREV SERMON < < <}}
~
\hyperref[sec:index_chronic]{\small{[返順時目]}}
~
\hyperref[sec:index_scriptual]{\small{[返順卷目]}}
~
\hyperref[sec:R7nxbZBAviM]{\small{> > > NEXT SERMON > > >}}
\newline
\newline
哈巴谷書 3:16-19
\newline
\begin{longtable}{cl}
\hline
\hline
章節 & 經文 (和合本修訂版)\\
\hline
3:16 & \begin{tabularx}{0.7\textwidth}{X} 我聽見這聲音，身體戰兢，嘴唇發顫，骨中朽爛，在所立之處戰兢；但我安靜等候災難之日臨到那上來侵犯我們的民。 \end{tabularx} \\ \\ \relax
3:17 & \begin{tabularx}{0.7\textwidth}{X} 雖然無花果樹不發旺，葡萄樹不結果，橄欖樹也不收成，田地不出糧食，圈中絕了羊，棚內也沒有牛； \end{tabularx} \\ \\ \relax
3:18 & \begin{tabularx}{0.7\textwidth}{X} 然而，我要因耶和華歡欣，因救我的神喜樂。 \end{tabularx} \\ \\ \relax
3:19 & \begin{tabularx}{0.7\textwidth}{X} 主耶和華是我的力量，他使我的腳快如母鹿，又使我穩行在高處。這歌交給聖詠團長，用絲弦的樂器。 \end{tabularx} \\ \\
[1ex]
\hline
\hline
\end{longtable}
$^{1}$弟兄姊妹平安.
平安.
我們開始時再一同禱告.
上上要賜下你的話語.
成為我們腳前的燈,路上的光.
你的話語最為寶貴.
求你讓聖靈充滿我們每一個的生命.
願意聖靈幫助我們.
真的挪走所有破壞了我們聽到的一切元素.
讓我們的心全然向上帝你開放.
我們聽到神的教導.
我們又會去作出適切的回應.
願意神在我們當中得盡一切當得的榮耀和頌讚.
奉耶穌基督的名字祈禱.
阿們.
今天這段經文很難就這樣說幾節.
如果就這樣說幾節.
大家也不知道他在說什麼.
一開始就這樣說.
所以我有一個邀請.
平時我不太想大家看手機.
因為怕大家看了他的IG.
Facebook,是不是自己朋友的訊息.
今天卻是很想大家開一開你的手機.
你的APP是有經文的.
我們開去就是哈巴谷書.
因為也需要有少許上文下理.
看看在說什麼.
我們稍稍看一看.
哈巴谷先知所服侍的時間.
大概是什麼時間呢.
首先他服侍的對象.
他對著什麼人去說呢.
就是對著猶大的子民百姓.
而當時的時間是很後期的.
已經是在猶大的百姓被擄去巴比倫之前.
所以是比較後的時間.
我們都知道當時以色列國一分為二有南北.
北國很快就已經滅亡了.
後期就剩下南國.

$^{41}$你看到哈巴谷就在南國這一邊.
當時應該是約克斯打後的時間.
在那裡約瓦京的時間.
他在當中去服侍.
我們去這裡看看.
之後你看到巴比倫就上來毀滅了南國.
哈巴谷在這個形勢下去服侍.
面對一個非常惡劣的處境.
你說當日的處境和我們今天好像完全是兩回事.
也不是的.
我們試試在當中去思考一下.
當日的處境和我們今天有什麼地方相似呢.
我們一路去看看.
在這裡我想看看整個哈巴谷書的結構.
要拿這個結構作為一個很重要的骨幹.
首先在第一章的時間裡.
先知要問上帝.
他要問上帝為什麼你縱容.
在我們整個猶大國裡有這麼多的強暴.
有這麼多的罪孽.
我們看第一章的哈巴谷書時很快看到.
先知的心很痛.
看到自己國家裡有很多的罪孽.
很多的暴力.
所以他一路在哭喊.
雖然你看著手機.
我們很快去看一次.
第二節的時候.
先知向上帝哭喊.
暴力啊.
在國家裡有暴力.
你為什麼還不拯救呢.
為什麼你讓我看到這麼多的罪孽呢.
在這裡有毀滅.
有兇暴在我面前.
有爭執.
有紛爭不斷.
現在你知不知道什麼形勢.
律法是沒有效的.
以前的什麼什麼真理的教導.

$^{81}$全部沒有意思的.
公理從來都沒有彰顯.
有什麼意思啊.
還有啊.
惡人是圍困二人.
惡人當道.
現在我們國家是這樣的處境.
我們回想.
其實現在我們的世界.
是不是都很亂啊.
大家有沒有想過.
無論我是不是關心政治呢.
大家一定每天新聞.
看很多是什麼.
現在的戰爭.
你都很難說.
哎呀這些很遙遠的.
不關我的事.
一定關你事的.
然後價格上升了.
一定有東西是影響了.
我那天去看的時候.
影響一些很重要的救災的工作.
不只是這些都有很多東西影響.
原來是很多方面.
所以你不要想.
那裡遠方的戰爭.
還他的戰爭.
我還我的.
不是的.
一定是受到波及.
我記得去年年底的時候.
我們香港也發生了一個很震驚的.
是什麼啊.
望福縣的火災.
那一晚香港人的心.
真的凝結了.
真的感到很冰凍.
這樣的事都可以發生在我們當中.
還有很多.

$^{121}$其實每一天.
大家留意新聞.
是很多.
令到人心裡.
真的很不舒服的案件.
很多案件就在我們身邊發生.
不遙遠的.
可能就在你個人.
在你家裡.
自己都有經歷一些事.
不過你未必這樣.
向整間教會都說.
你知不知道.
我現在經歷著這些困境.
有這些不公義的人這樣對我.
有這些惡人就在我身邊.
就在我家裡.
你知不知道.
你未必這樣向所有人說.
我們當中.
也不只是退休的弟兄姊妹.
也有在工作當中的弟兄姊妹.
有做生意的弟兄姊妹.
可能在你的職場.
你也經歷過很多黑暗.
你不會經常跟別人說.
黑暗在我們身邊.
有時看電視.
我們看《東張西望》.
也有很多案件.
就在我們住的範圍.
住的範圍也可以有很多事情發生.
黑暗就在當中.
先知在哭喊.
暴力.
你為什麼不拯救.
他去展開和上帝的對話.
到了五到十一節.
神就回答他.
神回應.

$^{161}$祂說.
祂會興起巴比倫.
去擄掠聶邦.
祂會興起巴比倫.
來審判猶大國.
興起巴比倫.
巴比倫是什麼人呢.
我們看回一章.
在六節那裡.
在和修版這樣說.
祂說我必興起迦勒底人.
就是巴比倫.
他是那殘忍暴躁之民.
通行遍地.
這些人的特色.
他們的速度非常快.
他們的軍事很強壯.
很殘暴.
他們會霸佔不屬於自己的住處.
他威武可畏.
審判和威權都由他而出.
其實巴比倫的審判權.
他的威權.
都是上帝給他的.
上帝給這個民族.
要去審判聶邦.
更加要去審判猶大國.
經文形容.
他的馬比炮更快.
比晚上的野狼更猛.
野狼的特色.
撕裂生畜各方面.
很殘暴.
可以這樣說.
接著說他的戰馬.
說他的軍事強盛.
他們真的引以為自豪.
第九節說.
他們都為私行殘暴而來.
他們聚集俘虜多如塵沙.

$^{201}$甚至乎那種自豪去到怎樣.
他譏笑列王.
根本沒有國家.
在他們面前站立得住.
就是殘暴.
大家有沒有想過.
巴比倫是否屬於上帝.
信神的嗎.
不是的.
不是信神的.
你興起一個這樣外邦的民族.
來審判.
還有.
這個民族是殘暴的.
殘暴到怎樣.
之前說猶大是惡人圍困二人.
但這個殘暴者是怎樣.
十三節說.
惡人是吞滅比自己公義的人更壞.
不是圍困這麼簡單.
簡直是吞滅.
他是猛烈.
他很驕傲.
在過程中.
當然他們驕傲到第十一節.
他以自己的力量為神明.
實際這節的解釋.
他覺得自己的力量.
是來自巴比倫的上帝.
巴比倫的神明.
很厲害.
很勇猛.
哪個國家在我們面前站立得住.
上帝興起一個這樣的外邦人.
在這裡才知道.
真的很難受.
神.
我有沒有聽錯.
你在說什麼.
你興起.

$^{241}$不是相信你的.
你興起外邦的民族.
而且是殘暴的.
不是公義的.
上帝你不是公義的嗎.
你用一些不公義的人去審判我們.
比我們更加不公義.
我們差而已.
你用一些比我們更加不公義的來審判我們.
所以十二十七節才問上帝.
你為什麼用更惡的人去刑罰我們猶大人呢.
在二章一節才登上望樓.
還要等候神回答他.
二章二到二十節.
神回應他.
巴比倫到最後.
也會受到他的報應.
在裡面的詳細內容.
大家可以慢慢回去看.
二章的內容.
二章裡面才從五個向度.
說到巴比倫的罪惡.
他的差勁之處.
這裡.
神回應巴比倫一定會受到報應.
因為上帝是公義的神.
在這裡進入到第三章的時候.
詩人就不再一樣.
詩人開始轉了.
他說話的整個語氣.
他整個的內容不再一樣.
到了第三章的時候.
他開始進入到頌讚裡面.
他的語氣改變.
頌讚,賤敬,順心.
今天我們說的是黑夜中的順心之歌.
人在光明之處.
人在順境.
要說一些很有信心的話.
總是很容易的.

$^{281}$不用別人教.
我們每個人都懂得說.
而且說得特別大聲.
我們自己在順境當中.
對著那些不在順境當中的人.
特別大聲.
懂得去安慰別人.
你不要這樣想.
逆境會怎樣怎樣過去.
但是在過程中.
我想說.
如果自己也在逆境中.
我還能否這麼大聲,這麼有信心.
去說我信靠這位上帝呢.
就在這個結構裡面.
我也想說.
信心的秘訣在哪裡呢.
第一件很重要的是什麼.
我和你的生命都有什麼.
大聲一點說.
我和你的生命都有位上帝.
秘訣第一件事.
不是走去找誰屬靈為人.
我覺得.
阿A禱告很有能力.
祈禱很快就會實現.
所以我走去找阿A.
他有智慧.
所以我走去問阿A.
我遇到這個困難.
我覺得如何處理.
你為我祈禱.
讓阿A為你解決生命的問題.
我們每一個人都很容易陷入這樣的光景.
在教會裡面也是這樣.
阿A大別屬靈.
經常親近神.
他經常禱告.
他經常告訴我.
他禱告.

$^{321}$神怎樣回應他.
所以我有問題是怎樣.
第一時間不是找上帝.
我先找阿A.
阿A.
我告訴你.
我現在有這個困擾.
我有這個困難.
你和我問上帝.
不是自己問上帝.
你和我問上帝.
你問上帝.
我是什麼情況.
來到這個時間.
有耶穌基督在我們生命裡面.
我們每一個人都相信了耶穌.
在我們裡面有聖靈的人.
竟然.
我們自己不會第一時間找上帝.
就是找那些很屬靈的人.
先問問他.
阿A回答不了什麼.
還有一個阿B.
那個很出名的牧者.
排著隊.
整條行.
每個人都去找那個牧者.
聽聞很厲害.
他很厲害.
好像把脈.
中醫把脈.
西醫一看就知道什麼問題.
我又去找問問那個牧者.
我現在遇到這個困難.
你又為我問上帝.
究竟我做什麼.
我們是不是這樣.
有些更厲害.
魚翁撒網.
我問了ABCDE.

$^{361}$12345678910.
我重複這個故事很多次.
我說我多困擾.
說完之後又說.
說完又說.
我都是找人.
其實我今天說這堂道.
都是提醒我自己.
我第一個找誰.
才知道不是去找人傾訴.
才知道心裡為自己的百姓很難過.
不是第一時間去找人說.
你知不知道多黑暗.
多黑暗.
多黑暗我們這個地方.
而是去問上帝.
秘訣.
真是依靠上帝的第一個秘訣.
都要回到上帝那裡.
找上帝.
我們不一定明白.
為什麼會發生這些事.
我不相信.
只是一個半個人.
經歷黑暗.
可能我們每個人都經歷.
不同程度的黑暗.
甚至我們的人際關係.
甚至我們在教會裡.
我們和身邊的人的關係.
可能都在一個黑暗裡.
在這個黑暗裡.
第一時間可能我們很多事都不明白.
我們很不開心.
在裡面要找的.
就要找上帝.
神是很好.
神在過程中.
不要去責備哈巴菊.
有沒有搞錯.

$^{401}$你來質問我.
我是創天造地的上帝.
你竟然問我.
為什麼容許國宗有強暴罪孽.
我不去處理這些強暴罪孽.
神沒有責備他.
其實神很想我們去找他.
我們很多時候都會有幾個傾向.
一是我們自己很想做審判官.
我們見到有很多黑暗.
很多罪惡.
我們就會問神.
神你不出手的.
其實我們都可能像哈巴菊.
你不出手的.
你不是很恨惡罪惡的嗎.
為什麼你見到罪惡你不出手.
我們很容易做審判官.
為什麼你不出手對付罪惡.
有些人就說事不關己.
不要搞我.
不要緊.
會處理到的.
這些都不對.
要回到神身上.
要有真正的信心去問神.
神為什麼會有這樣的情況.
跟神對話.
跟神對話最終不是質問神.
而是你很想跟神有互動.
這是最寶貴的.
神正在等我們.
然後你看到神真的有回應.
神回應是多久多快.
神自己定.
很多時候我們的信心在哪裡失落.
其中這個位置都會失落.
我現在問神.
你叫我問就問.
不過神沒有回答我.

$^{441}$你怎樣計算.
神沒有回答你是什麼時候.
神回答你可以即時回答你.
一天回答你.
兩天回答你.
十天回答你.
很久都不出聲.
就是這樣.
你叫我問他又不回答我.
怎樣行呢.
神一定回答.
在神的時間他會回答你.
這個神真的會回答他.
但我知道在先知書裡.
大家有機會回去看.
不是即時回答.
我記得有一個等了十天.
有些可能更久.
他不是不回答.
在他的時間一定回答.
但他想你先去回應.
先去問.
在過程裡神回答了之後.
先知就停下了.
可能有些事我更加不明白.
你這樣安排的.
上帝在當中顯明什麼.
祂才是王 祂才是主.
我用什麼人去審判我的百姓.
我用什麼人去管教我的百姓.
我用什麼人去伸張正義.
神有主權.
不到我和你說.
神有神的時間.
不到我和你說.
所以先知在問神的時候.
其實就開啟了這條路去尋覓.
神自己怎樣去看.
然後你看到.
先知做了一個很好的榜樣.

$^{481}$他一直等神回應的時候.
他不會放棄.
他的心仍然很關注他的百姓.
先知登上守望樓.
守望樓是一個很重要的地方.
是他去為他的子民百姓去守望.
讓他們得到提醒.
得到幫助的地方.
先知登上的地方要等神回應.
我和你都應該有這個決心.
如果你問為什麼要有信心.
就是我們一定要等上帝.
等到為止.
他一天不回答 兩天不回答.
等吧.
一個月不回答 兩個月不回答.
等吧 等到為止.
神一定回答.
怎樣回答是他的選擇.
然後神都會回答.
這個過程巴比倫最後都會受到報應.
上帝顯明他自己公義的性情.
巴比倫是這麼兇.
他們自以為是.
神使用他們審判以色列民的時候.
他們自以為是.
以為自己天下無敵.
以為我們的神明那麼厲害.
他們驕傲.
審判亦都會臨到.
在三章的頭部.
先知很能看到上帝的威榮.
我們看看上帝的威榮有多厲害.
神回答.
當神降臨的時候.
我們很快看看.
神會讓我們看到什麼.
他的榮光是會遮蔽住天.
還稱讚要充滿大地.
在過程中.

$^{521}$神是很榮耀.
很光輝.
很全能.
在三章頭的位置.
一直讓大家看到.
神的那種厲害.
在過程中.
亦都讓人看到.
除了他的力量之外.
他的能力大到超越萬有.
大自然.
全宇宙都在上帝的掌管下.
在三章五節.
前面有瘟疫的流行.
腳下有熱症的發出.
疾病都在神的掌管下.
在COVID期間.
我們很難明白.
我們覺得很辛苦.
經常屈在屋裡.
有很多條條框框.
我們在這裡生活.
但你知道嗎.
其實神都在當中工作.
神在當中工作.
其中一件事.
就是有些地方的污染.
竟然在那幾年間.
大大減少.
原來大地都得以休息.
在那幾年間.
神都有他很奇妙的工作.
那幾年COVID期間.
不是失控.
神啊.
你是不是不知道有這個疫情.
神啊.
是不是這個病你搞不定.
不是.
我昨天和青少年人.

$^{561}$一起慶祝生日的時候.
我講詩篇139篇.
我們的日子.
尚未到一日.
都寫在哪裡.
寫在神的冊子上.
你記住.
我們的日子尚未到一日.
神都寫在冊子上.
我們今天就覺得很驚訝.
今天我會發生這件事.
對上帝來講從來都沒有意外.
祂寫下了.
祂根本就知道你會發生什麼事.
這個疫情對於神來講.
也不是意外.
祂早就知道.
今天我們生命裡面.
發生的幽暗的事.
神早就知道.
從來沒有一件事是失控的.
在神掌管裡面.
所以在過程裡面.
我們不得不去讚美這位上帝.
是這麼厲害的.
疫情.
疫症也是在祂掌管.
接著我們繼續看經文.
講到三張六箭的時候.
祂站立的時候.
震動大地.
整個地球都受祂影響.
還要永久的山崩裂.
長長的嶺是塌陷.
影響是何等的大.
到我們的神.
要發怒的時候.
全世界都要害怕.
我們有時在地面上見到惡人.
都覺得他們很恐怖.

$^{601}$但神的憤怒才會令人害怕.
為什麼?.
八節要說.
你豈是向江河發怒.
向江河生氣嗎?.
神生氣.
誰不害怕?.
真的害怕.
這真是很恐怖.
接著說.
我們上帝打仗的時候一定贏.
經文裡面說.
祂騎在馬上.
在戰車上.
一直得勝.
回到我們今天的經文的時候.
才說我聽見的聲音.
身體漸徑.
什麼聲音?.
是上帝臨在的聲音.
上帝懲罰的聲音.
懲罰誰?.
不只是懲罰巴比倫.
也懲罰自己百姓.
你記住.
當時猶大國很黑暗.
很多罪惡.
當上帝的聲音臨在的時候.
懲罰他的百姓.
懲罰巴比倫的聲音臨到的時候.
先知很害怕.
真的害怕.
整個人震動了.
很深很深的影響.
上帝要去懲罰的時候.
誰可以在神面前站立得住呢?.
但在過程中.
先知的心情不只是害怕.
反而是有盼望.
但我安靜等候.

$^{641}$安靜等候災難之日.
臨到上帝侵犯我們的民.
我等候神.
等候神的工作臨在.
我可以做什麼?.
我是否可以叫上帝不要管教百姓?.
我是否可以叫神不要去懲罰巴比倫?.
絕對不可以.
我可否叫上帝不要做什麼?.
不可以.
我根本沒有權.
上帝作為臨到的時候.
我可以做什麼?.
我就是在等候.
而且要有盼望.
那裡說了一堆.
很多基督徒很喜歡的金句.
如果無花果樹不發黃.
葡萄樹不結果.
橄欖樹也不收成.
整個國家在饑荒當中.
這肯定是出於神的管教.
是饑荒.
沒有糧食.
很重要的.
無花果樹,葡萄樹,橄欖樹,田地.
什麼收成都沒有.
沒有糧食.
卷中是羊卷.
沒有羊.
這些很重要.
然而我要因耶和華歡欣.
因救我的神喜樂.
這一個很大的逆轉的信心.
怎麼來的?.
一直和神都有對話.
沒有停過.
我不能阻止神做什麼.
但是.
他記起很重要的一件事.

$^{681}$我和神的關係在這裡.
一直我去問神.
我回應神.
一直神都和我對話.
我和神有關係.
不單止這樣.
還有什麼呢?.
很重要的.
在《哈巴谷書》二章四節說.
二人必因信得生.
這一句在新約也很喜歡引用.
二人必因他的信得到生命.
那個信是對上帝的信.
對上帝的忠誠.
我一直去信神.
因為什麼得生?.
不是因為我有好的行為.
不是因為我不犯罪.
沒有人不犯罪.
你說有好行為.
同時也會有罪惡.
但卻是因為我們的信而得生.
先知對自己這方面很有信心.
直到那一刻.
我堅定地去相信這位上帝.
所以祂要因耶和華.
歡欣欣救我的神而很喜樂.
有信心.
我因這份信.
我會得救.
無論外在環境如何.
我會得救.
我和上帝的關係不會改變.
再惡劣也好.
我還可以開心.
還要喜樂.
主耶和華是我的力量.
因為我知道祂會拯救我.
我不知道祂的拯救何時臨到.
我不知道最後會怎樣.

$^{721}$但我和神的關係不改變.
就什麼都不用怕.
上帝是我的力量.
祂令我的腳像舞鹿般快.
又令我穩恆在高處.
這是很重要的一件事.
到後面.
祂令我穩恆在高處.
上帝就是我的力量.
上帝就是我的詩歌.
我和神的關係不改變.
祂一定會拯救我.
最重要的.
不是你今天外在環境有多少依靠.
我有很多很多錢.
我有間漂亮的房子.
我有漂亮的車.
最重要不是這些.
最重要不是你什麼順利.
你擁有什麼.
有什麼朋友圍繞著你.
不是這些.
最重要是你和神的關係.
你因信得生.
因為信你和神的關係.
在這些惡劣的環境裡不會改變.
而且上帝給你力量.
上帝在災難裡拯救你.
我前陣子在一個很特別的場合.
認識一個黃福縣的住客.
一個很特別的情況.
這完全不是因為要關顧他們.
很特別的是.
這個家庭應該是屬於.
大埔一間教會的家人.
我記得那次認識他們的時候.
聽他們說整個經歷很奇妙.
大埔教會的姐妹.
一直都在上教會.
一直都很想自己的家人信耶穌.

$^{761}$當然我接觸這個家庭的時候.
他們應該是未信的.
但是都覺得有些事情很奇妙.
然後這個爸爸和媽媽.
都是住在黃福縣裡.
他說當日是很特別的.
火災剛剛起的時候.
這個爸爸趕得及回去.
就跟媽媽說我們一起走.
其實不是拿到太多東西.
不是拿到太多東西.
但是趕得及.
剛剛好一起走.
一離開就大火.
全部燒了那個大廈.
所有東西都沒有了.
但是他們的生命是保存了.
而他的兒子因為當日要上班.
所以就沒有牽涉在內.
我就在想.
就是那幾分鐘.
生命就是那幾分鐘.
就已經扭轉了局勢.
只要遲那幾分鐘.
他們兩個都會身陷在火海當中.
這個姐妹.
保守她爸爸媽媽.
在裡面你覺得人可以為自己.
家人做多少事呢.
你可以掌控什麼呢.
沒有什麼可以掌控.
你不能掌控交通.
你不能掌控自己的健康.
你不能掌控人際關係.
你的際遇.
但是有一樣東西不會改變.
因順得生.
上帝去拯救我們.
甚至拯救我們家人.
弟兄姊妹.

$^{801}$我很想鼓勵大家.
不要總是希望人生只有這些.
不要動我.
所有東西都是平平安安.
我們的生命就是這樣過日子.
反而要讓自己和神的關係.
要去盡心.
以致人生經歷各種風浪.
黑暗的時候.
我都不被打倒.
以後講這段經文不是空蕩蕩.
我送給你.
弟兄姊妹.
雖然無花果樹不發黃.
葡萄樹不結果.
橄欖樹日不收成.
雖然這樣.
我們要因我們的神喜樂.
我見很多弟兄姊妹.
很輕易地送這段經文出去.
我心裡面問一個問題.
你真的這樣經歷嗎?.
你在橄欖裡是否真的夠僵?.
你對神的信是否可以超越.
現有的環境裡一切的黑暗呢?.
記住.
第一時間去找上帝.
和上帝對話.
我很渴望你們未來的日子.
聽到更多弟兄姊妹和我說.
我和神對話.
你知道神跟我說什麼嗎?.
這個很棒.
我也很渴望自己這樣跟你說.
我和神對話.
神跟我說什麼.
你知道神怎樣帶領我經歷這件事嗎?.
這個才是最寶貴.
不是第一時間去找那些.
所謂的屬靈偉人.

$^{841}$找那些很厲害的.
禱告很重要的人.
然後他代你禱告.
他代你求問上帝.
不是這樣的.
我們一同禱告.
天父上帝.
我們在哈巴谷書當中.
真的看到.
先知那個心情很大的逆轉.
在一個很大的.
很痛心的處境裡.
在一個很黑暗的處境裡.
先知最後.
看到出路.
就在神的身上.
上帝你是我們唯一的拯救.
上帝你是創天造地的神.
我們還不到一天.
我們的日子.
都寫在你的冊子上.
不僅如此.
這個世界還沒有經歷一天前.
其實你已經知道將會發生什麼事.
在這個世界上所有的事情都沒有失控.
你在掌管當中.
你有你的心意.
你行奇妙事.
亦不是我們能夠明白的.
可能我們很多時候都會質問你.
為什麼會讓苦難發生.
為什麼會讓戰爭發生.
為什麼還要容隱惡人在身邊.
但是天父.
你有你的心意.
讓我們每天把握.
抓緊機會和你對話.
讓我們很重視和你的關係.
就在你的回應中.
就在你的帶領中.

$^{881}$甚至你的拯救中.
我們走過這一生.
向你致福.
紅印堂的眾弟兄姊妹.
讓大家有一個豐盛的人生.
豐盛不是因為我們今天生活的條件有多豐厚.
不是因為大家的日子有多順利.
而是正正在困難中.
大家經歷了多少.
向你致福我們.
奉耶穌基督的名字祈禱.
阿們.
\newpage



\section{歌羅西書 2:6-15}
\label{sec:R7nxbZBAviM}
\textbf{2026.04.19 《在基督裡生根建造滿豐盛》:歌羅西書2\_6-15(和修版) 講員 : 譚立晶傳道}
\newline
\newline
連結: \href{https://youtube.com/watch?v=R7nxbZBAviM}{\texttt{https://youtube.com/watch?v=R7nxbZBAviM}} ~~~~ 語音日期: 2026-04-21
\newline
\newline
\hyperref[sec:y558B5MlYc4]{\small{< < < PREV SERMON < < <}}
~
\hyperref[sec:index_chronic]{\small{[返順時目]}}
~
\hyperref[sec:index_scriptual]{\small{[返順卷目]}}
~
\hyperref[sec:WokQkPhl0sU]{\small{> > > NEXT SERMON > > >}}
\newline
\newline
歌羅西書 2:6-15
\newline
\begin{longtable}{cl}
\hline
\hline
章節 & 經文 (和合本修訂版)\\
\hline
2:6 & \begin{tabularx}{0.7\textwidth}{X} 既然你們接受了主基督耶穌，就要靠著他而生活， \end{tabularx} \\ \\ \relax
2:7 & \begin{tabularx}{0.7\textwidth}{X} 照著你們所領受的教導，在他裡面生根建造，信心堅固，充滿著感謝的心。 \end{tabularx} \\ \\ \relax
2:8 & \begin{tabularx}{0.7\textwidth}{X} 你們要謹慎，免得有人用他的哲學和虛空的廢話，不照著基督，而是照人間的傳統和世上粗淺的學說，把你們擄去。 \end{tabularx} \\ \\ \relax
2:9 & \begin{tabularx}{0.7\textwidth}{X} 因為神本性一切的豐盛都有形有體地居住在基督裡面； \end{tabularx} \\ \\ \relax
2:10 & \begin{tabularx}{0.7\textwidth}{X} 你們在他裡面也已經成為豐盛。他是所有執政掌權者的元首。 \end{tabularx} \\ \\ \relax
2:11 & \begin{tabularx}{0.7\textwidth}{X} 你們也在他裡面受了不是人手所行的割禮，而是使你們脫去肉體情慾的基督的割禮。 \end{tabularx} \\ \\ \relax
2:12 & \begin{tabularx}{0.7\textwidth}{X} 你們既受洗與他一同埋葬，也就在此禮上，因信那使他從死人中復活的神的作為跟他一同復活。 \end{tabularx} \\ \\ \relax
2:13 & \begin{tabularx}{0.7\textwidth}{X} 你們從前在過犯和未受割禮的肉體中死了，神卻赦免了你們一切的過犯，使你們與基督一同活過來， \end{tabularx} \\ \\ \relax
2:14 & \begin{tabularx}{0.7\textwidth}{X} 塗去了在律例上所寫、敵對我們、束縛我們的字據，把它撤去，釘在十字架上。 \end{tabularx} \\ \\ \relax
2:15 & \begin{tabularx}{0.7\textwidth}{X} 基督既將一切執政者、掌權者的權勢解除了，就在凱旋的行列中，將他們公開示眾，仗著十字架誇勝。 \end{tabularx} \\ \\
[1ex]
\hline
\hline
\end{longtable}
$^{1}$大家好,我由短頭髮變成長頭髮.
很感動,剛才聽完孫教士的分享.
其實我覺得我們的生命很豐盛.
是因著我們在基督裡面生根建造.
以致不只是我們信徒自己的家可以豐盛.
也可以將這個豐盛繼續帶給其他人.
我剛才一直聽孫教士的分享,我都覺得很感恩.
我都覺得有時我們在港道的服侍.
其實是萬事互相效力.
無論是詩歌,經文,大家互相輝映.
希望可以激勵,鼓勵信徒.
怎樣繼續讓神去豐富我們自己.
其實剛才孫教士分享的東西.
對自己都是一個很大的鼓舞.
亦都希望成為大家的激勵.
就是我以前都是不喜歡回家的.
因為家裡都很混亂.
如果大家有印象都知道.
因為爸爸賭錢.
媽媽要很顧家.
因為爸爸不能給家裡穩定的收入.
所以媽媽就很努力工作.
所以對我們的教養都沒有什麼.
所以我很明白他說在社區裡面.
很多人要流浪在街頭.
我自己就不想回家.
我那時的抉擇就是.
總之我中學的時候留在學校.
我是留到學校關門.
我才會想回家.
我又不想到處飄蕩.
因為有時都怕遇到不好的人.
所以就很不想回家.
因為回家就是另一個壓力的延伸.
但是很感恩.
就是上帝去更新我們的生命.
讓我們不再一樣.
以致當我回顧過去.
我都分享.
其實我頭十七,八年.

$^{41}$我的人生都不知道怎樣過就過了.
我不知道你們的人生是怎樣.
你們會不會家裡都很好.
所以小康之間長大.
沒有什麼優秀.
但是我自己是渾渾噩噩.
就十七,八歲了.
那時十七,八歲就已經要面對人生一些考驗.
會考那些.
那時考完試才知道.
原來考不到是不能繼續讀書的.
可能那一刻才很驚訝.
原來是這樣的.
早知道就用心讀書.
但是有早知道就.
沒有什麼大家都知道.
所以有時覺得原來有很多東西都不懂.
為什麼呢?.
因為爸爸只顧著賭錢.
他也不會教你.
媽媽努力上班.
也將不了很多東西告訴你.
所以我覺得有時我們做社區的服務.
很重要的是.
教會是另一個家.
我覺得這個是很重要的.
所以那時我信了主之後.
我是最喜歡回教會的.
我覺得教會這個家.
是給我最安全的.
給我最輕鬆的.
所以不用留在學校之後.
一有什麼就是回教會.
那時我教會的傳道目的都很好.
它都是打開教會的門.
你隨時回來.
你想回來休息也可以.
你回來看聖經也可以.
你回來做功課也可以.
總之教會的門就是為你開.

$^{81}$我覺得這種生命成長.
是帶給我人生很大的改變.
所以到我又再過十七,八年.
那就是三十多歲了.
再過十七,八年的時候.
就很感恩了.
原來上帝是很有耐性.
培育我們再次成長.
那時到了這個年紀.
我就成為了一個高齡產婦了.
神賜我有家庭.
賜我可以生小孩子.
就到現在了.
我女兒八歲了.
大家大概知道我幾多歲了.
所以我覺得神是很奇妙的.
有時候我們好像活不到現在.
因為我們經常被過去拉扯.
但是原來上帝是很豐盛的.
根本祂是很有耐性.
給我們時間去生根建造.
有時候我們會有一種自卑.
或者有一種比較的心態.
會令到我們.
別人很好.
自己一直都是不好.
會有一種拉力.
但是上帝其實是超乎我們所想像.
是很豐盛,很豐富.
也願意給我們時間去更生,改變.
在祂裡面得蒙,造就.
以致很豐盛的時候.
讓我們可以進入人生不同的階段.
繼續去不單祝福自己.
也去祝福身邊的人.
所以我覺得很好.
洪恩棠也是有很大的進步.
因為你們原本由自己去做社區工作.
你們去到做菲律賓的社區工作.
洪恩棠很厲害.

$^{121}$鼓勵一下自己.
原來你們跨出海外去服侍.
為你們感恩.
和大家一起看看今天的經文.
我們希望透過主的話語.
我們彼此提醒.
讓我們得蒙,造就.
在第六節到第七節.
我們待會兒再讀.
讓我們可以抓緊上來的話語.
預備起.
既然你們接受了主基督耶穌.
就要靠著祂而生活.
照著你們所領受的教導.
在祂裡面生根建造.
信心堅固.
這就是我們自從信耶穌之後.
會有的不一樣的生命.
所以祂說.
既然你們接受了主基督耶穌.
讓祂居守.
讓祂跟你有關係.
祂說就要靠著祂而生活.
所以我們說基督教的信仰.
不只是信條,思想上,知識上.
不是的.
原來是靠著祂而生活.
這是一個與主同活的經驗.
是一個很真實的互動.
可以互動的關係.
然後第七節說.
照著你們所領受的教導.
在祂裡面生根建造.
其實這個接受和領受.
我們看回當時的希臘文化.
其實是說傳承口述的傳統.
是什麼呢?.
就是你真的接受了的時候.
就不只是停在你那裡.
如果你真的接受了的時候.

$^{161}$你是可以說出來的.
然後再給下一個.
我現在教女兒就是.
我經常跟她說.
其實如果你想你測驗考試成績好.
你怎樣證明你是懂得呢?.
我做了媽媽之後.
真的有很多東西要學.
不只是小朋友要學東西.
媽媽也要不斷進步.
才可以跟這個時代的小朋友同行.
我看一些說讀書很厲害的人.
他們有一個秘技.
他們說原來他們.
可以隨便擺一個玩偶.
擺一個玩偶出來.
就當是學生.
他說如果那個很厲害的人.
他就是將他上完課後學的東西.
他說這個就當是你的學生.
然後你可以將你學的東西全部教會他.
就是說那樣東西是不是你的東西.
就是你的東西了.
你說學了學了.
但是你說不出來.
你傳遞不下去.
那你是不是真的學到了呢?.
就真的不知道剩下多少了.
我就很鼓勵我的女兒.
去幫她在校車上的其他同學.
因為她認識到.
她一起坐校巴.
其實也有半個多小時.
那沒有了.
當然可以聊天.
但是我說到測驗考試墨書的時候.
我說不如你也幫一下那些同學.
因為她有時候會跟我說.
媽媽.
他們什麼都不懂的.

$^{201}$哈哈哈.
但是我就跟她說.
我說其實你在恩典之下.
我說如果爸爸媽媽.
我說你爸爸媽媽沒有相信.
可能你也是什麼都不懂的.
我說是因為你爸爸媽媽相信了.
生命有改變.
我們能夠將比較好的東西給你.
以致你可以有一個比較好的環境去學東西.
我說如果你真的見到其他人不是很懂的.
你會不會可以幫一下他們呢?.
我說我以前讀書.
就是家裡的緣故.
其實真的不是很能夠用心向學的.
因為裡面有很多很複雜的情緒.
令到沒什麼心機讀書.
我說但是媽媽是很感恩.
回到教會之後.
是很多叔叔姐姐.
哥哥.
就是姨姨這樣去教媽媽.
很有耐性去教媽媽.
所以媽媽可以重新做人.
重新活過來.
我說你試一下這樣吧.
很好啊.
她真的教她的同學.
用不同的方法去教她的同學.
一起進步.
我說你做得很好.
媽媽很欣賞你.
就是這樣一個傳承的階段.
所以當我們這一群跟隨耶穌基督的人.
我們是基督徒的時候.
我們就要想.
我們怎樣是因著我們所接受,所領受的東西.
我們要活出來.
活出來就是將我們領受了的東西.
不要放在心裡.

$^{241}$用我們的言行舉止.
用我們的生活去流露出來.
將我們所信的東西去表達出來.
所以很感恩.
剛才我一來到樓下.
是先接受睿寧的羅漢果茶.
然後進來看到奶茶的團隊.
其實真的覺得這幅圖畫很美麗.
大家學而致用.
將你學到的東西.
去化成祝福.
去給街坊.
給主樂的弟兄姊妹.
一起去互相建立.
所以她就說.
如果我們真的這樣在這個傳承的口熟傳統裡面.
其實就是在說我們很鄭重地去領受了.
並且繼續地去傳遞開去.
當我們真的這樣做的時候.
她就說在她裡面生根建造.
會有什麼效果呢?.
就是信心會堅固.
我剛才看到這群團隊去服侍的時候.
每個人的樣子都是.
嘩!.
好像天使般美麗.
大家認不認同?.
認同.
對.
不是又要做了.
做又36.
不做又36.
不是這個樣子.
而是來喝杯茶.
來喝杯奶茶.
是很開心地去服侍的.
這種就是你對神有信心.
所以你願意抽你的時間去服侍.
因為神更新了你.
你很想其他人也繼續改變.

$^{281}$這種關係的互動就來得很重要了.
而說到生根建造.
這是一個很重要的詞語.
我很想跟大家解構一下這個詞語.
我們不可能用完整篇原文來跟大家說.
但是某些字眼很重要.
我又想跟大家探究一下.
生根其實有很深層的意義.
紮根的過程.
生根就是慢慢地滲入泥土裡面.
然後紮穩了.
那麼植物就可以蓋得很高.
或者它可以得到很多養分的時候.
就可以有很多收穫.
其實生根這個希臘文.
是用了一個「完成時被動分詞」.
大家會說「你說什麼?我會慢慢解釋的」.
「完成時態」是什麼呢?.
原來是在表示一個已經完成和持續存在的狀態.
所以不是只是嘗試紮根.
而是已經紮根並且持續如此.
所以當我們接受了基督的時候.
其實這個紮根的行動已經發生了.
現在是活在這個現實裡面.
所以我們全部信主的人.
其實我們已經紮根了在基督裡面的時候.
我們已經發生了和活在這個現實裡面.
而為什麼說「被動語態」呢?.
其實是在表示這個行動的主體是誰呢?.
神.
所以不是我們要把自己紮根在基督裡面.
而是我們信主的那一刻.
神已經把我們紮根了,栽種了在這裡.
所以這個就破除了我們經常有的迷思.
就是我們要靠自己的力量去紮根.
其實不是這個意思.
而這個「根」這個字.
其實描繪了植物深深向下紮根的景象.
這個不只是表面黏著.
而是深入土壤吸取氧化獲得一個穩固的過程.

$^{321}$大家有種過東西就知道了.
種植物不是一朝一夕就開花結果的.
如果有陪小朋友種植的學習.
通常那時候學校會叫大家回家種什麼呢?.
對,綠豆,蠶豆.
然後讓小朋友看到原來會有更新出來.
然後慢慢又會向上.
所以其實上帝是願意我們透過祂.
祂其實已經把我們紮根在主基督耶穌裡面.
其實我們肯有耐性願意被神繼續栽種.
其實我們的根就會越來越深.
並且我們開始的生命是有東西看的.
我不知道大家有沒有聽過一句話.
就是「忍耐到底,必有東西看」.
(笑聲).
原來有些東西不是立刻看得到的.
我也很記得自己初信主頭五年.
我覺得是最艱難的.
因為那時候的根其實還不夠深.
而且初信主有很多考驗.
我不知道你們有沒有遇過這些考驗.
我講一下我以前未信主的生活習慣給大家聽.
這個我就很明白.
剛才孫教士所說在社區的街童,流浪那些.
那種生活和你回到教會的生活是完全兩回事.
真的可以改變人.
一個好的習慣,一個持續的習慣.
是很改變一個人的人生.
我以前被灌輸.
我不知道你們有沒有經歷過跟我類同的事.
如有類同實屬巧合.
但我們都是在神的恩典下.
我記得以前家裡被灌輸.
灌輸我們一個觀念.
星期六,日最重要是用來做什麼?.
睡覺.
還要睡到什麼時候?.
睡到黃昏白夜.
真的可以睡到中午才起床.
然後睡醒之後.

$^{361}$一個很重要的任務.
知道要做什麼嗎?.
去喝茶.
現在回想起來就覺得很匪夷所思.
原來時間就是這樣浪費了.
喝茶就是那間茶樓超多人的.
大家很不介意用一個多兩個小時去排隊喝茶.
我那時候就是在這個時代去生活.
所以以前拜神.
很簡單幾句就拜完了.
因為通常說什麼?.
「保佑平安,身體健康」.
如果喜歡打牌的話就「求多陽」.
四個字的.
「橫財就手」.
就是這幾句.
但你想想信了主之後.
「早起身」是第一個挑戰.
然後就是去祈禱也是一個挑戰.
以前祈禱幾句而已.
就是求假神幾句.
我記得第一次被教會的弟兄姊妹很熱誠的邀請.
「不如你也試試回教會的祈禱會吧」.
很單純的我.
就應邀了.
「好啊!好啊!回來祈禱會吧」.
之後每個人都說完了.
然後就說「我們現在一起祈禱了」.
「不如我們就輪著祈禱了」.
「待會到完祂就到你了,好不好?」.
很單純的我很快就應了.
「好啊!」.
然後就開始了.
哇!怎料一期又自可.
哇!為甚麼他們的祈禱這麼豐富?.
哈哈哈.
然後到自己的時候.
糟糕了!不懂得說.
那怎麼辦?又不敢說話.
又不知道怎麼說.

$^{401}$但他們是聰明的.
他們就馬上幫我補位.
跟我一起祈禱.
然後就完結了.
也沒有怪我.
「你為甚麼剛才不說話?」.
也沒有怪我.
但你會看到原來就是要重新學習.
那個屬靈的生活.
原來跟神的關係.
未夠親密的時候.
真的不懂得怎麼說話.
但就是日積月累回教會.
學習如何跟神建立很深入的關係.
原來就發覺.
祈禱真的可以很豐富,說不遠的.
這就是生命的突破.
所以我很記得.
我覺得頭五年信主是很艱難,很挑戰.
但對我生命有益.
所以有時候我遲了賴床,未起床.
是我媽媽.
那時候她還未信主.
她也會主動拍醒我.
「你不是要回教會嗎?」.
「你起床吧!」.
「你再不起床就來不及了」.
這也是一個很特別的經驗.
就是媽媽未信主.
但她看到你回教會是好的.
也會叫你回去.
我就是這樣一步一步地成長.
頭五年是很艱難.
但頭五年.
我覺得我的生命是來一個很大的逆轉.
我每次都很珍貴.
你們手上拿著的那份東西.
那份甚麼?.
程序表.
原來我會帶這張程序表回家.

$^{441}$因為我當作是一個小小的功課.
小小的教科書.
小小的學習.
因為完全對聖經是零知識.
對詩歌是零知識.
但這份東西就很好.
有時候裡面會教我們如何祈禱.
我就是每個星期.
將這份東西帶回家.
那個星期我會不斷重溫這份東西.
唱過的詩歌我繼續唱.
看過的經文我繼續看.
慢慢地深化我自己.
我真的發覺原來.
根就是這樣慢慢地進心.
以致我的生命真的逐一改變.
所以很想鼓勵大家.
其實我們真的要很謹慎.
在第八節提醒我們.
第八節我想大家一起讀.
預備起.
你們要謹慎.
免得有用他的哲學.
和虛空的廢話.
不照著基督.
而是照人間的傳統.
和世上粗淺的學說.
把你們老去.
過往未信主前十七,八年.
就是被這些東西老去了.
以為這樣就叫做人生.
但原來信了耶穌的時候.
就發覺原來那些只是不同的學說.
不同的歪理.
不同的東西影響我們的生活為人.
當我信了主之後.
我發覺原來我的時間可以好好用.
和可以好好榮耀神.
所以那時候一邊成長一邊裝備.
教會有主教學上.

$^{481}$教會有訓練班上.
教會有機會讓我事奉.
我去事奉.
就是一步一腳印地建立了我的生命.
扭轉了我之後十七,八年的成長.
到我真的有機會再建立我的家庭.
我有兒女.
我就可以把不是我.
十七,八年的東西灌輸給她.
而是把我信了主之後的十七,八年.
寶貴的東西傳授給她.
以致不只是我成長.
我也希望她成長.
我之前有機會去國內.
做一些書法交流的工作.
我要離開她幾天.
當然女兒一定會說.
「我不捨得你」.
很不想媽媽離開.
媽媽是去做正經事.
我們很想幫人信耶穌.
你乖乖地等媽媽回來.
我覺得每一次信心的學習.
就是父母要學習信得過神.
也要讓小朋友信得過神.
因為父母不可以永遠陪在小朋友身邊.
但上帝可以永遠陪在小朋友身邊.
所以我很希望培育我的女兒.
知道最重要的是有上帝陪在身邊.
而不是最重要的是有媽媽陪在身邊.
當然不是這樣就不理會她.
「你就是和上帝一起」.
不是這個意思.
要學習.
讓她知道媽媽也很信得過她.
其實她是可以獨立的.
試試媽媽離開幾天.
其實也是幫你做一些你以為自己做不到.
但其實你是有能力做到的.
然後你又會對自己很有自信.

$^{521}$所以我出發前向上帝交代.
因為正正是放復活節假前的時間.
還有幾天要上學.
我就說「你努力」.
「你試試媽媽不在」.
「你全部事情自己處理好」.
「然後爸爸回來幫你處理好」.
「你又可以輕鬆了」.
「可以自己玩了」.
我們和傭人姐姐有一個群組.
我也很甜蜜.
我去公幹.
她會在群組說.
「爸爸,你放心」.
「我放學回來了」.
「我已經在放學前做好所有功課了」.
「我不安心,很開心」.
然後我在那次交流團.
主要是因為我是隨團的牧者.
所以我會負責所有祈禱的環節.
都是交給我祈禱.
而正正我的女兒又被揀選了.
我也覺得很奇妙.
讓我和女兒有機會好像是兵分兩路去服侍.
因為他們的學校是基督教的學校.
他們也會有復活節崇拜.
竟然這次老師選了她做領土.
他們會分開的.
一至三年班一場.
四至六年班一場.
我女兒就是三年班.
老師竟然選了她做領土.
我也相信老師看到她也有生命力.
給她機會.
之後她就和我分享.
其實我在那邊工作時不太記得那些時日.
是女兒在群組裡說.
「媽媽,剛才我領土」.
「另外兩位司儀看到我的手在顫抖」.
「很緊張」.

$^{561}$我說「那你有沒有吃螺絲?」.
「你能不能說出來?」.
她說「我也沒有吃螺絲」.
「也能說出來」.
我說「哇,那你很厲害」.
「媽媽在國內也是祈禱」.
「你也在香港學校裡祈禱」.
「哇,我們一起事奉神」.
「這幅圖畫很美麗」.
我說「如果下次老師再邀請你」.
「你還會祈禱嗎?」.
「會」.
哇,這就是互相輝映.
就是生命彼此的建造.
所以我很鼓勵我們一眾弟兄姊妹.
我們不要怕讓我們的根要時日去生.
要在神裡面去生.
讓我們的生命可以豐盛.
而且不要讓舊有的東西繼續影響我們.
因為我們真的要知道.
在我們未信主前.
真的會有一堆東西.
是潛移默化影響了我們很多.
我們要不斷看神的話語.
我們才可以成為一個很好的.
讓我們有一個真理的分辨能力.
以致我們不要將舊有的東西.
影響我們活出基督信仰的生命.
所以到了第九節.
我們一起讀吧.
第九節,預備起.
「因為神本性一切的豐盛」.
「都有形有體地居住在基督裡面」.
其實保羅要去寫《歌羅西書》的時候.
是因為當時他們在希羅文化的底下.
其實就是希臘和羅馬的文化裡面.
如果大家有去過希臘.
你會看到有很多的神像.
阿波羅,太陽神,雅典娜.
很多的神.

$^{601}$其實他們是拜很多泛神論的祭祀.
所以其實當時的基督徒.
他要不要受那些影響.
而且當時也有很多異端邪說出來.
他要去否定耶穌基督真實的人性和神性.
他很想將耶穌基督降格.
所以保羅就要寫給信徒知道.
我們真的千萬不要受其他影響.
以為信一個耶穌不夠.
不是的,其實我們信一個耶穌就很足夠.
所以到了第十節那裡很重要.
第十節那裡我們又一起讀.
預備起.
「你們在他裡面也已經成為豐盛」.
「是所有執政掌權者的元首」.
所以這裡說其實我們的山根建築.
就在耶穌基督裡面.
而在祂裡面就已經是一切的豐盛.
所以我們就要提醒自己.
不要那麼容易被周圍的東西.
轉移了我們的焦點和視線.
這個世代有很多東西會吸引我們的注意力.
但是就讓我們離開了真理.
求主幫助我們要認定.
其實在主裡面是已經得著最豐盛的.
所以我們要將元首這個守衛.
將上帝優先.
其實我女兒經常都會給我一些考驗.
但是我就要讓她知道媽媽信仰的根本.
她小時候經常會問.
「媽媽,你有多喜歡我?」.
她是女孩子,可能會比較甜.
會比較喜歡情感上的東西.
「你喜歡我多一點,還是喜歡爸爸多一點?」.
她經常都很想要做第一個.
但是我狠心一點都要跟她說.
我說其實媽媽是要學習最愛的是上帝.
我說第二個是爸爸.
我說第三個才是你.
但是全部都是我所愛的.

$^{641}$她有時候之前好像很接受不到.
為什麼我不是第一?.
我不知道你們做過媽媽就知道了.
有時候子女就很想獨霸著你.
但是我覺得我不想讓她有錯誤對信仰的投放.
其實我就要跟她慢慢去解釋.
我覺得這個耐性是需要花的.
也要讓她經歷到.
我說因為你知道嗎?.
媽媽如果沒有上帝.
我說可能你都離不開這個世界.
我說如果沒有上帝.
你就沒有媽媽,也沒有你.
所以我說最重要的是上帝.
我說第二,你想一下.
如果你沒有爸爸,也沒有你.
我說你想一下是你先出現還是你爸爸先出現.
我知道有學輔導的我們都知道.
家庭裡面真的要很小心不要失衡.
有時候就是夫妻關係不小心.
讓子女優先過配偶.
其實關係是會暖和的.
但是如果站穩夫妻關係.
怎樣都很好.
子女才是真正的好.
所以我就要慢慢很有耐性跟她講解.
雖然你不是排在第一,第二.
但是你都是排在最重要的.
我說沒有其他人會取代到你們家人這個位置.
這些都是聖經教的.
聖經除了很多神學觀念.
其實也說了很多人倫的價值觀.
我們應該怎樣排次序.
怎樣愛人,但是又要如己.
所以你是很懂得愛自己.
你會更加懂得愛護你身邊的人.
但是你都不是很懂得愛護自己的時候.
很多時候你會對待對方都不是很好.
因為有時候我們對自己很嚴苛.
然後你就怎樣?對別人都很嚴苛.

$^{681}$可能就不是很平衡的狀態.
這些我覺得都是一路學聖經.
是神的話語不斷將我以前舊有的觀念.
去到慢慢清除.
所以求主幫助我們記認.
在祂裡面是最豐盛.
還有祂是所有執政掌權者的元首.
有時候我們以前是被電視帶大的.
另一個狀態.
我們小時候是看很多劇集的.
你不要小看那些劇集.
其實那些劇集是怎樣?.
不斷在潛移默化不同的信念.
不同的價值觀給你.
我心想我就是看得那些多.
所以我就是有那種要報仇雪恨的心.
所以我爸爸不好.
我就要長大衛了報仇.
但是我回到教會學聖經.
聖經教我什麼?.
要去愛仇敵.
然後我又要學寬恕.
反而原來學回聖經真理那一套.
我才真的對我的生命是最好的.
有時候原來我們為了報仇.
都會令到自己變質.
但是如果我們是要恃上帝.
我們反而是學了耶穌基督的品格.
以基督的心為心.
願意謙卑,願意服侍,願意尊重神.
這一切的種種都是神對我們的生命的磨練.
所以我都很感恩.
就是在我和我的女兒相處的時候.
其實我就是不斷做那種接受,領受.
以及傳承的口述傳統.
將我得著的東西告訴她.
我都很記得有一次.
我們要帶她的同學回教會.
新朋友.
我就和我的女兒說.

$^{721}$「你帶這個同學回去」.
我說「你可不可以堅持一件事」.
我以前是在短孫中心受訓的.
我們受訓是怎樣可以幫助一個新朋友.
可以直跟在教會裡面呢?.
就是兩個月帶他,陪他回五次教會.
他就應該是熟悉教會的人.
他就可以繼續在裡面.
然後你又可以再帶其他新朋友回教會了.
我就說「你可不可以這兩個月,這五個星期」.
「你盡量去帶他回教會」.
「以及坐在他旁邊」.
之後我的女兒就說.
「o下!這樣是不可以的」.
「我教會的其他朋友」.
「經常都說要跟我坐的」.
「這樣是不可以的」.
她們好像爭著要和他坐一樣.
真是很厲害.
我就和她說.
「但是你這樣是最能幫到他」.
她就覺得有點不耐煩.
「那我試試吧」.
我也在想究竟是否能做到呢?.
是否可行呢?.
之後她那天完了聚會之後.
回到家裡.
她就告訴我.
「我今天陪了他」.
我說「你今天陪了他?」.
「誰沒有說你沒有和他坐嗎?」.
「我和他說了」.
「他知道了」.
所以有時候我們信仰生命都是這樣的.
有時候可能我們一開始要接受上帝的真理.
「o下!這樣嗎?」.
「嘩!很難」.
但是如果你肯嘗試堅持去做.
其實又不是不可以的.
然後你做到的時候.

$^{761}$其實你又在豐富別人的生命.
你又覺得你做得到這件事.
你又慢慢慢慢成長.
所以我很感恩.
有時候在學校的家長日見老師.
我覺得成績是另一個關注的地方.
但我另一個更加關注的地方就是.
究竟我的女兒怎樣實踐她的信仰生命呢?.
老師給她的評價都高.
都會說「其實你的女兒」.
「是很關注其他同學的需要」.
「她也是很能幫忙的」.
他們學校又很好.
給他們一人一職.
我記得一年級.
她是被選派做老師的秘書.
我心想「現在的學校這麼棒」.
「我們以前是流行甚麼?」.
「就是頒獎」.
「嘩!有秘書這麼誇張」.
但是很感恩.
就是賦予她可以做.
所以我們不要小看我們年輕這一代.
我們說這一代叫「Gen Z」.
其實他們年紀很小,很多東西都比你厲害.
所以是有不同的機會讓他們發揮潛能.
所以我很感激基督教學校.
讓小朋友發揮潛能.
我的女兒亦很有自信.
她說「老師會派我上四樓」.
「又要上五樓,又要上六樓」.
好像幫她做跑腿.
「有時真的很辛苦」.
「令到我都不能放小息」.
我便跟她說「但是你想想」.
「老師要你這樣做」.
「其實老師是怎樣看你?」.
她說「即是甚麼?即是甚麼?」.
我說「其實就是老師很信任你」.
我說「是否每個同學都得到老師的信任?」.

$^{801}$她說「不是,有些不知道去了哪裡」.
我說「就是了」.
我說「但是你每次老師交給你的東西」.
「你走了多少層樓都做得好」.
然後我說「你有幸成為老師的跑腿」.
「是你的榮譽來的」.
我說「所以你很厲害」.
她又飄飄然說.
「雖然沒有消息」.
「但是被人讚完,肯定完又很開心」.
所以其實我想說.
我們的生命是可以成為上帝的跑腿.
我們是可以與神同工.
神就是很想給我們一份尊榮.
讓我們可以與神一起受尊榮.
所以如果大家記得有些聖經是說.
我們是君尊的族裔.
即是君尊的祭司.
屬神的子民.
其實這一切都是說.
我們本身是很有地位的.
上帝給我們這個地位身份.
所以我們怎樣好好在基督裡面建立根基.
然後我們在裡面得著豐盛.
以致我們不再被舊有的東西.
或者現時的東西.
每一個時代都有不同的歪風.
但是求主幫助我們用真理去抗衡.
也不要被它打擾,影響.
我們要記得我們原本是怎樣的人.
但是因著耶穌我們不一樣.
在第十一,第十二節.
我們一起讀這兩節經文.
「預備起,你們也在他裡面受了不是人手所行的國禮.
而是使你們脫去肉體情慾的基督的國禮.
你們既受洗與他一同埋葬.
也就在此例上.
因信那使他從死人中復活的神的作為.
跟他一同復活」.
所以我們知道為何洗禮是要浸過下水.

$^{841}$然後又要上來.
其實是在宣認我們所信的東西.
我們能夠宣認我們與主同死同埋葬.
也是同復活.
所以我們要跟耶穌一起活過來.
我們基督徒的生命是活生生去見證神.
所以我們每一個人有生命氣息.
都要去讚美,都要去事奉.
都要去分享主裡面的豐盛.
所以求主幫助我們去記認我們生根的過程.
很多時候生根的過程是沒有人看得到.
有時候可能連自己也不太知道.
但是原來在過程裡面.
我們慢慢地被栽培,被囂張.
然後我們的生命不再一樣.
如果大家記得在《馬太福音》第十三章.
講到殺種的比喻.
其實是在說要落在好土裡面的種子.
人聽了做.
其實就是持續在誠實,善良的心裡面.
並且忍耐著結實.
所以紮根是隱藏的.
是向下的.
是需要時間的.
也是生命成長不可或缺的基礎.
所以我們不需要急於跟身邊的人去比較.
而是跟自己的過去比較就可以了.
因為每個人都不一樣.
到最後的三節.
我也想請弟兄姊妹我們一起去讀.
預備起.
你們從前在過犯和未受割禮的肉體中死了.
神卻赦免了你們一切的過犯.
使你們與基督一同活過來.
逃去了在律例上所寫的對我們.
束縛我們的字句.
把它撤去.
釘在十字架上.
基督既將一切執政者.
掌權者的權勢解除.

$^{881}$就在凱旋的行列中.
將他們公開示眾.
穿著十字架跨性.
所以最後這幾節其實是一個德性的宣告.
就是從死亡到生命.
回想我未信主之前.
我的信念就是.
唉,無論如何都是死.
人生有什麼意義.
是這麼灰的.
但信了主之後.
原來生命是祂賜給我的禮物.
我要活好我的每一天.
就算我的人生活得不好.
但原來我有永恆的生命.
有另一站的.
所以我不要太被我現存的生命完全影響.
就算是很艱苦也好.
但如果我堅持到底的時候.
上帝是會存留一個新天新地給我的.
所以令我對生命現存的.
或者將來的人生觀是完全不同的.
然後另一個德性的宣告是什麼呢?.
就是在說「同末聚齋」.
基督的十字架是處理我們聚齋的終極場所.
其實第十四節有寫「致據」.
大家知不知道是什麼來的?.
「致據」.
其實就是說要親手簽署的借據和罪狀.
以前譬如去衙門看劇集.
你真的認罪怎麼辦?.
那個法官就砰地拍.
然後那個罪犯就要「嗯」.
其實就是在說原來我們未信主之前.
其實是在說我們舊約律法的要求是達不到的.
還有我們已經累積了很多的罪責.
所以我都很深刻.
很記得我信主那一刻.
我是覺得其實我也有很多罪.
耶穌真的可以救我.

$^{921}$真的可以原諒我.
跟我傳福音的人就很有耐性.
他慢慢跟我講解.
他說耶穌是有能力赦免你一切的罪.
只要你願意認罪的時候.
耶穌是可以徹底既往不咎.
不去記住你一切的罪狀.
所以我最深刻的就是.
我決定立志跟隨基督.
一旦祈願出來.
我簡直覺得整個身體都變輕了.
我背負著很多的罪擔.
背負著很多人生的壓力.
但是那一刻原來有神跟我一起同行.
一起去背負的時候就不再一樣了.
所以其實這個字句就是.
神在我們面前怎樣?.
撕掉它.
當我們說.
「耶穌,我們願意相信你」.
耶穌就說「嗯,這個是你的字句」.
沒有了.
繼續努力去成長.
不要再犯罪.
那時候讀神學院的時候.
我有一個點也不太明白.
但是我相信了耶穌.
但我還繼續犯罪.
講到罪論的時候.
就是說其實耶穌是有能力.
將你繼續一生裡面犯的罪.
都可以去處理.
其實神是讓我們不斷去學習.
這個清二成聖的過程.
所以我們不是放縱地去犯罪.
而是我們正在改變.
其實上帝也願意饒恕我們.
信了主之後還有的罪惡.
其實我覺得在這個世上.
是很難完全聖潔不犯罪.

$^{961}$但是我們是怎樣?.
是追求聖潔.
那個是我們的目標.
是我們的焦點.
所以求主幫助我們.
我們要有一個遠觀.
有時候我們說.
看得遠.
你就會走得堅定.
你只是看眼前.
就會令你失腳.
所以他要提到.
第十五節的跨性.
那個海船.
其實就是說我們要慶祝勝利.
跨性是說一個大遊行.
基督是德性的元帥.
十字架不是失敗的記號.
而是神勝過一切靈界敵對的勢力.
祂要拿起十字架成為祂德性的升旗.
祂要公開海船.
所以我很期待.
將來我們等到神回來的時候.
我們所有一起事奉神的人.
一起海船而歸.
一起去榮耀神.
神也將我們提升.
去接我們回天.
所以在這段經文的短短幾節裡.
其實保羅很想為我們.
畫出一個完整的圖畫.
是什麼呢?.
有四方面.
大家記住今天所說的話.
就是在說根基.
我們已經被紮根在基督裡面.
其實就是要在裡面穩固.
所以我們要讀聖經.
我們要去祈禱.
我們要去回教會.

$^{1001}$我們要跟弟兄姊妹相愛.
我們要去學習.
也要去事奉.
我們的根基就會越紮越穩當.
然後我們也知道.
上帝是我們的供應.
我們已經在祂裡面得著豐盛滿足.
其實唯有耶穌基督才能夠滿足我們靈魂的需要.
不是其他的物質.
不是其他的知識.
但是沒有衝突的.
你得著上帝很多豐盛.
不代表其他東西不要讀.
其他東西不要學.
不是這個意思.
但是你要把位置放得對.
然後就是在說我們的地位.
我們已經受了國禮.
脫去肉體.
與它同死同活.
所以我們不再被綑綁.
我們的生命是可以自由做決定的.
我很記得我改變最大的就是.
如何看待自己的價值.
以及別人如何看待我.
以前我很自卑.
經常覺得別人看到我就不喜歡我了.
但是我改變了價值觀.
我不再看別人如何看待我.
我看誰如何看待我.
Yes.
看神如何看待我.
在聖經裡面就是說.
原來我們是祂的寶貝.
是祂如珠如寶.
神會提拔貧寒人.
所以不要覺得自己可能出於汗味.
身世不好.
家庭背景不好.
但是原來神是重視我們每一個人的性命.

$^{1041}$所以我們要在祂裡面得到自由.
神看我們為好.
我們就是好的了.
不用看其他人.
或者魔鬼如何對我們的說話和控訴.
然後第四就是德性.
我們的罪債已經被塗抹.
受敵已經被跨性.
他們其實沒有權力或任何狀況去說我們好不好.
我們只看上帝如何說話就可以了.
所以求主幫助我們.
透過上帝.
我們在祂裡面活好.
讓我們的根基建立得好.
以致我們很豐盛的時候.
也將這份豐盛傳遞給身邊的人.
我們做好我們的社區.
我們還去菲律賓的社區.
去到其他地方的社區.
有機會我們還去澳門的社區.
我們讓更多生命很豐盛的一起彼此去建立活動.
我們一起同心圖告.
慈悲仁愛天父.
因為你是很愛我們的.
其實你是很願意給我們時間在你裡面生根建造.
你也很願意我們去探索你一切的豐盛.
原來真的不斷去讀你的話語,不斷去祈禱.
去經歷你的同在,你的回應.
你一一的應允.
你給我們的盼望.
你讓我們真的在裡面要站立得穩.
不要那麼容易被外面動盪的時勢影響我們的心思意念.
主啊,求你幫助我們.
讓我們依然看得到.
你是那位充滿愛的上帝.
你是願意尋找失喪靈魂的上帝.
你是那個願人人都悔改的上帝.
主啊,你是那位不擔言的上帝.
你就是想給我們機會繼續去成長.
也將我們身邊所愛的親友都帶到你面前.

$^{1081}$主啊,為到無論我們在香港這個地方.
在菲律賓這個地方.
在澳門這個地方.
在世界各地戰亂中的地方.
我們都仰望教堂在你手中.
因為每個地方我都相信是你放了你自己的人在裡面.
主啊,你就讓每一個基督徒都在你面前.
能夠活出你的豐盛.
能夠繼續在這個全城裡面.
對你的話語有把握.
相信你的應許.
好好去活出不一樣的生命.
求你繼續去塗做我們每一個.
我們同心禱告.
願你賜福我們的社區.
讓我們去祝福我們的社區.
我們同心禱告,教託奉主名求.
Amen.
\newpage



\section{希伯來書 12:14-28}
\label{sec:WokQkPhl0sU}
\textbf{2026.04.26 《存感恩與神同行》:希伯來書12\_14-28(和修版) 許淑芬牧師}
\newline
\newline
連結: \href{https://youtube.com/watch?v=WokQkPhl0sU}{\texttt{https://youtube.com/watch?v=WokQkPhl0sU}} ~~~~ 語音日期: 2026-05-01
\newline
\newline
\hyperref[sec:R7nxbZBAviM]{\small{< < < PREV SERMON < < <}}
~
\hyperref[sec:index_chronic]{\small{[返順時目]}}
~
\hyperref[sec:index_scriptual]{\small{[返順卷目]}}
~
\hyperref[sec:DhmQL4xjGtQ]{\small{> > > NEXT SERMON > > >}}
\newline
\newline
希伯來書 12:14-28
\newline
\begin{longtable}{cl}
\hline
\hline
章節 & 經文 (和合本修訂版)\\
\hline
12:14 & \begin{tabularx}{0.7\textwidth}{X} 你們要追求與眾人和睦，並要追求聖潔；人非聖潔不能見主。 \end{tabularx} \\ \\ \relax
12:15 & \begin{tabularx}{0.7\textwidth}{X} 要謹慎，免得有人失去了神的恩典；免得有毒根生出來擾亂你們，因而使許多人沾染污穢， \end{tabularx} \\ \\ \relax
12:16 & \begin{tabularx}{0.7\textwidth}{X} 免得有人淫亂，或不敬虔如以掃，他因一點點食物把自己長子的名分賣了。 \end{tabularx} \\ \\ \relax
12:17 & \begin{tabularx}{0.7\textwidth}{X} 後來你們知道，他想要承受父親的祝福，竟被拒絕，雖然流著淚苦求，卻得不著門路使他父親回心轉意。 \end{tabularx} \\ \\ \relax
12:18 & \begin{tabularx}{0.7\textwidth}{X} 你們不是來到那可觸摸的山，那裡有火焰、密雲、黑暗、暴風、 \end{tabularx} \\ \\ \relax
12:19 & \begin{tabularx}{0.7\textwidth}{X} 角聲，和說話的聲音；當時那些聽見這聲音的，都求不要再向他們說話， \end{tabularx} \\ \\ \relax
12:20 & \begin{tabularx}{0.7\textwidth}{X} 因為他們擔當不起所命令他們的話，說：「靠近這山的，即使是走獸，也要用石頭打死。」 \end{tabularx} \\ \\ \relax
12:21 & \begin{tabularx}{0.7\textwidth}{X} 所見的景象極其可怕，以致摩西說：「我恐懼戰兢。」 \end{tabularx} \\ \\ \relax
12:22 & \begin{tabularx}{0.7\textwidth}{X} 但是你們是來到錫安山，永生神的城，就是天上的耶路撒冷，那裡有千千萬萬的天使， \end{tabularx} \\ \\ \relax
12:23 & \begin{tabularx}{0.7\textwidth}{X} 有名字記錄在天上眾長子的盛會，有審判眾人的神和成為完全的義人的靈魂， \end{tabularx} \\ \\ \relax
12:24 & \begin{tabularx}{0.7\textwidth}{X} 並新約的中保耶穌，以及所灑的血；這血所說的信息比亞伯的血所說的更美。 \end{tabularx} \\ \\ \relax
12:25 & \begin{tabularx}{0.7\textwidth}{X} 你們總要謹慎，不可拒絕那向你們說話的，因為那些拒絕了在地上警戒他們的，尚且不能逃罪，何況我們違背那從天上警戒我們的呢？ \end{tabularx} \\ \\ \relax
12:26 & \begin{tabularx}{0.7\textwidth}{X} 當時他的聲音震動了地，但如今他應許說：「再一次我不單要震動地，還要震動天。」 \end{tabularx} \\ \\ \relax
12:27 & \begin{tabularx}{0.7\textwidth}{X} 這「再一次」的話是指明被震動的要像受造之物一樣被挪去，使那不被震動的能常存。 \end{tabularx} \\ \\ \relax
12:28 & \begin{tabularx}{0.7\textwidth}{X} 所以，既然我們得了不能被震動的國度，就要感恩，照著神所喜悅的，用虔誠、敬畏的心事奉神， \end{tabularx} \\ \\
[1ex]
\hline
\hline
\end{longtable}
$^{1}$各位姐妹平安.
不知道大家今天我們要與神同行.
不知道大家會不會常存一個感恩的心來生活呢.
記得美國一個擔任最長的總統叫羅斯福.
因為戰爭或者其他原因.
他就不斷延長他總統的任務.
有一天他在家裡被偷東西.
失去了財物之後.
他的朋友就問他.
朋友慰問他.
但是他回應的朋友就不是有什麼抱怨.
他反而很感恩.
他講出三件事.
很感恩雖然被偷東西.
不過第一件事很感恩.
那個賊只是偷了我的東西.
不是偷了我的性命.
第二件事就是.
賊人只是偷了我部分的東西.
又不是全部.
第三件事最值得慶幸的就是.
那個賊就是他不是我.
這些都是我們可以數一下.
值得感恩的地方.
不知道大家會不會在逆境的時候.
都能夠懂得感恩.
如果我們基督徒懂得感恩的話.
可能我們的生命會更加燦爛.
會更加令人羨慕.
所以我們求主要給我們.
去過一個懂得感恩的生活.
很多人說.
人生不如意事是十常八九.
在困難中如何真正學懂感恩呢?.
我相信羅師傅能夠感恩.
其實不是因為他多偉大.
而是因為他的思想上.
他懂得或學懂三件事.
令他懂得感恩.
這亦是我們今天在心理學上.

$^{41}$讓我們學習的地方.
第一就是要正面思考.
每件事當然有正面的一面.
去埋怨.
第二就是要接受現實.
既然已經偷了.
我們當然需要從正面的角度.
積極的角度去相信.
既然有問題.
但需要找解決難處的方法去處理.
第三就是需要憑信心.
特別是基督徒.
憑信心相信神容許這些事發生.
一定有他的心意.
你可能會問.
不幸的事竟然能夠感恩.
是否真的呢?.
我不知道你有沒有試過.
我曾經探訪過一個弟兄.
這個弟兄不幸地三次中風.
原來中風是可以不斷處理的.
但當探訪他的時候.
他很感恩.
第三次中風.
很感恩.
他為什麼感恩呢?.
他第一次中風的時候.
他就很輕看問題.
仍然大飲大食.
當沒事發生.
因為第一次很輕微.
到第二次中風的時候.
又差不多.
覺得不要緊.
很快就好了.
但第三次中風的時候.
他開始清醒.
知道自己會傾向容易中風.
他就覺得這是神的管教.
他就要更加戒掉壞習慣.

$^{81}$倚靠神.
藉著神的醫治.
他能夠重新生活.
原來當我們面對這些逆境的時候.
我們可以用另一個角度去想.
去看.
這就是能夠幫助我們積極處理.
或者過一個很感恩的生活.
剛才我讀的一段希伯來書.
你應該用新譯本.
用新譯本去讀.
不好意思.
我仍然很舊.
我用和學本的內容.
剛才我們讀的那段經文.
我想大家特別注意.
第28節.
28節是核心.
有以上的內容.
我們就應該要傳一個感恩的心.
有一個敬畏神去感恩的心.
去度日.
上文就講了很多不同的內容.
特別講到基督的超越性.
基督超越天使.
超越摩西.
超越祭司.
超越舊約.
所有的律法和所有的憲制.
既然基督超越一切的話.
他就跟著在中間11章那裡.
就去講了很多的信心偉人.
他就指出由舊約到整段歷史裡面.
你要看見很多信心偉人.
他怎樣面對碧白苦難呢?.
他仍然有一個很堅定的心.
依靠上帝去面對.
第12,13章就特別提醒信徒.
我們因為得著基督這麼多東西.
我們應該要有正確的生活態度.

$^{121}$所以這裡特別要指出.
我們應該要有感恩的生活.
真正感恩的人.
是必定活在生活裡面.
能夠活出來的.
你沒理由說我心裡面感恩.
但我臉上有疤痕.
不可能的.
我們心裡面去感恩的時候.
我們在身上或是在行動上.
都能夠去發出感恩的表現.
所以當我們今天要去思想感恩的時候.
我希望用三個角度和大家去思想.
第一個就是我們向外應該會怎樣.
我們就過一個和睦的生活.
向內我們怎樣呢?.
我們持守上帝要我們的聖潔.
然後向上就是怎樣呢?.
我們繼續不斷持守上帝給我們的恩典.
這三樣東西你要帶來給我們.
就是因為這樣.
所以我們要過一個感恩的生活.
第一方面我們要思想.
我們怎樣過一個追求和睦的生活呢?.
追求和睦的生活.
向外的.
我們對身邊的人.
我們怎樣去過一個和睦.
我們說和睦不等於我們和其他人一致.
不可能我們和其他人一致.
但是聖經就常常去教導我們.
那個和睦的生活.
其實和諧的生活.
是必須要附上行動.
必須要學習和去努力的.
所以在聖經裡面.
不斷地要教導我們.
所以我問.
為什麼人與人之間.
你有眼耳口鼻.

$^{161}$我有眼耳口鼻.
為什麼我們會不和呢?.
所以我就要推翻到.
之前就是.
我們知道人類犯罪之後.
亞當夏娃犯罪之後.
所以我經常查經的時候.
跟弟兄姐說.
回到天家.
第一個要找誰呢?.
當然除了見基督之外.
第一個就要找亞當夏娃出來.
怎樣去找呢?.
有個爛笑話大家都知道.
亞當夏娃怎樣找的?.
沒錯.
因為他有肚臍.
所以一定找到他.
亞當夏娃犯罪之後.
當然我的前設就是.
亞當夏娃是得救的.
我在天上見到他.
但當我們去問亞當夏娃犯罪之後.
有一個我們知道是副產品.
就是推卸責任.
除了看到自己不好的東西之外.
他還會帶來一個很重要的.
就是推卸責任.
而推卸責任的背後.
其實是一個自我保護.
他們自我保護.
所以因為自我保護的過程.
人與人之間有很多的衝突.
不和諧.
因為你最關心的當然不是我.
不是隔壁那個.
你最關心的當然是自己.
當人有利益的時候.
那個和諧就會有衝突.
所以當我們今天面對這個.

$^{201}$嬌徒不要有太多本能的反應.
因為聖經反而教我們.
很多時候當我們面對這些問題.
面對人生這些問題的時候.
我們就需要反本能反應去面對.
你說和睦.
我們說對好的人.
人家對你好.
你就很自然和諧.
人家對你好沒理由罵他.
但人家對你不好.
很多人很誠實地說.
我當然對他不好.
我怎會對他好.
你很誠實地面對.
但基督教的信仰就是帶來.
給我們一個比較反本性.
本能反應的一個情況.
當我們面對人與人之間的衝突的時候.
首先我們要明白.
這個根本是正常的.
是人類犯罪之後一個正常的表現.
也是因為這樣.
所以我們不要幻想.
我想大家回教會也不會幻想.
教徒是好人.
大家坐在這裡.
大家都知道.
看看旁邊那個.
也是壞人一個.
只不過是我們蒙恩的罪人.
所以我們人類沒什麼好做.
我們通常都做壞事.
但既然是這樣的時候.
我們不要常常都期望有太完美的人.
或者我們不要期望.
對方為什麼會這樣.
我們覺得失望.
其實不需要.
這個正正是反映.

$^{241}$人面對人際關係的時候.
最大的衝突.
但是我們也不需要害怕.
反而聖經教導我們.
今天我們怎樣面對.
這個人際衝突的問題.
如果我回想.
結婚三十多年了.
眨一眼.
時間真的過得很快.
坦白說.
你問牧師會不會打算離婚.
有些人聽我分享的時候.
說牧師也打算離婚.
有多出奇.
想想而已.
不是真做.
但為什麼會有離婚的想法呢.
當然一定是.
你會覺得有衝突的情況.
誰結婚的.
都不要拒絕.
因為你旁邊那個看到.
但你心想.
你會不會有一刻覺得.
這個人真的一般般相處.
可能會好一點.
可能你會一閃閃過.
當然我們讀過書的人.
就會知道.
離婚這個字不能輕易出口.
所以我當然.
不要理他說.
忍一會就搞定.
我也相信結婚三十多年.
也相信忍一會就搞定.
忍一會就不會再想離婚.
這是一個經歷.
所以我要去面對.
家庭就算只有兩個人.

$^{281}$既然有這麼大的衝突.
我們面對一起.
一定有很多問題出現.
也都是教會.
教徒之間有這些問題出現.
我們都仍然要說.
是正常的.
但為什麼我們今天仍然可以一起合作.
事奉呢.
我就要明白到不是因為我們合得來.
不是因為我們性格一致.
而是因為基督.
是我們的生命.
因為有基督的緣故.
我們可以繼續合作.
繼續同行.
繼續忍你.
這是我們生活最好的相處之道.
希伯來書第十三章.
繼續回應一個.
很重要.
這章之後.
很重要的主題.
既然要相愛.
就要提醒弟兄們要怎樣相愛.
包括怎樣接待客人.
怎樣紀念受苦的人.
怎樣尊重婚姻.
怎樣有知足的心.
怎樣紀念傳道人.
後文就繼續來發展.
我們因為基督的緣故.
我們要過一個和諧的生活.
我們要怎樣展現.
在我們的生活裡.
後文來講.
但我們今天要看這裡.
如果我們說.
我們要傳一個感恩的心.
因為我們想與神同行.

$^{321}$我們要傳一個感恩的心.
第一樣就要學習.
我們怎樣過一個和諧的生活.
亦都我們剛才講.
人與人之間.
最不能和諧的原因.
有很多.
可能有很多人說.
基本都是價值觀不對.
亦都很多性格不對.
亦都做事方法有差異.
又絕對是.
但那個最大的核心.
其實都是差一個字.
就是愛字.
大家都可能會很明白.
愛的力量很大.
大到一個地步.
可以引起很多.
沒有愛.
甚麼事都成為一個衝突.
但有愛就能夠去包容.
就能夠和諧相處.
所以無論是親人也好.
夫婦之間也好.
弟兄姊妹也好.
愛其實是一個最大的核心.
特別就算在婚姻最.
關係最.
人與人之間最親密的關係裡面.
都離不開愛.
才能夠去面對.
我記得有.
一些生活例子.
大家可以思想一下.
有其中兩個例子.
我想和大家分享.
我見過一對夫妻.
結婚第一年.
很多衝突.

$^{361}$我都想離婚.
結婚第一年最大衝突.
這對夫婦結婚第一年.
很多衝突.
我星期一煮飯.
星期二煮飯.
星期三煮飯.
星期四也應該是她煮飯.
為甚麼她不煮呢.
她覺得很激氣.
很生氣.
覺得丈夫.
很不愛這個家.
不要不愛她.
但丈夫和我說的時候.
很冤枉,他都沒有說.
星期一煮飯,星期二煮飯.
經常煮飯,他以為她喜歡煮飯.
於是就煮.
沒有想過自己要煮.
衝突的原因.
簡單來說.
就是沒有溝通.
沒有告訴別人.
別人怎麼知道.
所以婚姻關係也是.
做婚前輔導要教導弟兄姊妹.
畫公仔要畫出腸.
你想怎樣,要說清楚給對方聽.
對方不會猜到.
猜也可能會猜錯.
所以不如.
正確地告訴對方.
不然就會有很多.
你覺得不合理,他又覺得.
你為甚麼會這麼生氣.
有很多不明的地方.
另一個例子就是.
不知道有沒有跟大家說過.
有個弟兄,你好分享.

$^{401}$有一天和他太太有爭執.
有爭執.
他百分之八都覺得自己對.
很厲害吧,百分之八都覺得對.
於是就生氣他老婆兩天.
有些冷戰的情況.
但那兩天之後.
他太太就說.
本來已經康復了.
但他太太說那兩天真的想死.
其實他突然間有醒覺.
為他醒覺而感恩.
他說雖然百分之八都覺得自己對.
但令他太太.
這麼不開心兩天.
他又覺得自己.
就算對,也是錯的.
不應該是這樣.
所以他就開始有.
更多的改變.
我講過都教你.
贏了長交就怎樣.
輸了就要投家.
所以很多婚姻的問題.
如果日積月累.
一些衝突不去講.
就會有很多問題.
所以如果有喊殺不要緊.
回到聖經講.
不要含怒到日落.
日落之前就解決.
日日清的問題.
就會處理關係更好.
我們怎樣去.
追求和睦.
我們再講.
我們不是追求一致的生活.
也不是追求百分之八.
覺得是正確的.
反而我們追求.

$^{441}$大家彼此之間怎樣去.
因為愛的緣故.
建立一個新的關係.
所以未必一定要常常拿一個.
最好的地方.
反而次好的.
能夠關係好的其實無妨.
我記得曾經去過一個婚禮.
這個婚禮.
新人邀請我講幾句.
講五分鐘.
講五分鐘之前.
那個爸爸.
他們結婚幾十年.
他叮囑一對新人.
我覺得父母講的更好.
父母和這對新人了講.
你們的夫婦關係.
要白頭到老送他們一個字.
你知道什麼字嗎?.
有人說愛.
有人說忍.
其實就是忍字.
尤其是結婚幾十年.
忍是很重要的.
他奉勸.
給這個新人.
你要去忍.
其實忍在聖經裡面就是愛.
如果講哥林多前書十三章.
大家都很熟悉的愛篇.
這個愛篇最出名.
頭那字就是.
恆久忍耐.
到最後那句是.
凡事忍耐.
這是整個關係的核心.
所以當我.
三十多年前結婚的時候.
我媽媽送了一句說話.

$^{481}$叫相敬如賓.
什麼意思呢?.
就是你要對對方像嘉賓一樣.
嘉賓就不要罵他.
要照顧他.
要服侍他.
這句說話其實有很大的影響.
當我和我先生有衝突的時候.
這句說話就幫我提醒我.
要用什麼態度去對待他.
願上帝幫助我們.
特別在這例子裡.
特別是兩個人的相處.
願上帝幫助我們.
用感恩的心.
要和睦的生活.
我們要感恩.
怎樣去靠著基督和諧的生活.
我又想和旁邊那個.
來講一下.
我們要有愛.
有愛的家才和諧.
和旁邊那個講一下.
他可能在很多時候.
我們忘記了.
我們以為是事件的問題.
其實是因為你不夠愛.
所以你又會有.
很多的抽青.
很多的衝突.
如果我們夠愛的話.
我們很多事都會忍到.
和誰相處都一樣.
愛扮演一個很重要的角色.
第二方面我們要思想.
除了我們向外之外.
其實我們裡面都要去建立.
我們裡面要過一個.
追求聖潔的生活.
這個篇幅.

$^{521}$我特別給大家.
十四召十六節這個篇幅.
特別用了一個例子.
耳掃的例子.
耳掃貪戀世俗的心.
用了這個.
來做對比.
或者來建立聖潔的生活.
所以聖經.
經文裡有講.
如果不是聖潔.
就沒有人能見主.
也有講.
清心人有福.
因為我們能夠見神.
可想而知.
聖潔和清心.
或者聖潔和不貪戀.
是掛上一個.
一層又一層的不同角度的.
關係.
我們怎樣去操練一個聖潔的生活.
聖潔的生活不單止.
是講到我們的行為上面.
我們不要過犯淫亂.
或者奸淫的生活.
不單止是講行為上面.
其實聖潔的生活.
更加要知道.
我們的生命的重點.
我們的生命方向.
你的生命佔據的位置.
是甚麼.
所以今天我們要去省思的時候.
究竟有甚麼.
是影響我們.
那個聖潔的生活.
或者有甚麼在影響你的生活.
如果錢是在影響你的生活.
你要知道.

$^{561}$可能這是在影響你和上帝的關係.
如果你的兒女是在影響你的生活.
可能這件事.
是阻擋了你和神和諧的關係.
所以這個.
聖潔的生活意思就是說.
我要追求和上帝.
親密的.
關係的生活.
我以上帝為中心.
的關係.
這才叫真正的聖潔生活.
當然他就舉出一個例子.
就是耳掃的例子.
耳掃貪戀.
耳掃貪戀為甚麼是一個.
不懂得感恩的表現.
在十六節你告訴我.
耳掃因為一點的食物.
將自己的長子名份賣掉.
耳掃的問題.
不只是貪食.
耳掃的問題.
就是他輕看.
他的長子名份.
這才是他容易出賣了.
他的長子名份.
他為了一個短暫的祝福.
就放棄了.
本來是屬於他的福分.
作者將這個耳掃.
作者如果在希伯來書.
新約是特別.
為了希伯來書沒有作者.
告訴我們.
但是我們知道作者.
將這個耳掃放在這段經文.
正正就告訴我們.
耳掃是不懂得感恩.
耳掃不懂得珍惜.

$^{601}$他所擁有的.
不知道他所擁有的.
是多麼重要.
原來當一個不懂得.
珍惜上帝所賜給他.
獨特的東西的時候.
這個都叫不勝結.
或者這個都被指.
指責為.
我們貪婪.
我們有貪心.
貪心和拜偶像一樣.
拜偶像和邪淫一樣.
邪淫就是不勝結.
我們這樣去推論.
今天有什麼.
影響我們沒有了感恩的心呢?.
我們會不會同樣.
像耳掃一樣.
我們沒有了感恩.
沒有了珍惜.
上帝所賜給你的東西.
所以今天為什麼很多人.
貪戀別人的東西呢?.
就是因為我們不夠.
看得到自己所擁有的東西.
我們不會懂得.
去享受上帝所賜給你的東西.
我們只是覺得隔離飯香.
隔離的那個比較好.
隔離的才是好的.
我不知道隔離的兒子會不會比你兒子好呢?.
隔離的配偶會不會比你配偶好呢?.
隔離的生活方式.
物質享受.
會不會比你自己好呢?.
如果是這樣的話.
我們要反思.
我們是不是缺少了.
一個對上帝感恩的心呢?.

$^{641}$因為對上帝缺少了感恩的心.
我們很難與神同行.
所以我們需要.
去珍惜.
上帝賜給我們所有的東西.
用另一個角度.
第二個例子.
我們不要任由這個世界的.
吸引.
去攔阻或拉走.
我們遠離上帝.
所以聖經很清楚地說.
我們不能夠事奉神.
又事奉馬門.
你說這樣賺錢是不是不對呢?.
很多弟兄姊妹會說.
賺錢真的不對.
賺錢有什麼不對呢?.
你升職有什麼不對呢?.
沒有什麼不對的.
但如果這個成為你生活的核心.
才是不對的.
所以問題不在於.
那些中性的東西.
問題在於你的心.
你的心有沒有被這些.
去勾引了.
吸引了.
以致我們開始.
那個不像基督徒的生活.
或者不像以基督為中心的生活.
這個才是值得我們思考的.
所以今天我們去問.
這個世界有沒有吸引力呢?.
肯定有.
你不要告訴我沒有吸引力.
這個世界有吸引力.
否則今天就沒有市場了.
這個世界就這麼多吸引力.
每天你看廣告.

$^{681}$就算你正經看YouTube.
也有很多廣告給你看.
很多東西吸引你.
廣告的目的就是.
五秒吸引你.
吸引了我們.
影響了我們的視線.
所以今天我們要.
不斷去耐性.
耐性是我們基督徒.
一個很重要的寶藏.
當我們越懂得.
去耐性的時候.
我們的生命就會越校正.
回到神那裡.
所以今天我就問自己.
有什麼東西在影響我們的視線.
名牌會不會呢?.
我們的面子,尊嚴.
會不會影響我們呢?.
或者無論身邊的人也好.
事也好.
今天我們要去學習.
如何將基督放在首位.
我們的生命見證.
讓不讓人看到.
基督信耶穌是好呢?.
這個就是.
你生命的本質.
讓人有沒有體會,吸引力.
我記得很多年前.
四川大地震的時候.
有災民來分享.
當時有兩類人.
來幫助他們,他覺得很感動.
第一類當然是公安.
不過他會知道.
公安要來幫助他們是應該的.
因為政府差他們來的.
但是另一類他覺得很特別.

$^{721}$就是基督徒.
不斷地來幫助他們.
而基督徒不斷幫助他們.
全部都是發自內心的.
要從全世界角度.
都來幫助他們.
那個無私.
在他們身上看到.
那個無私的愛是怎樣.
好像上年文福縣.
火災的時候.
我們見到也有很多自願的團體.
進去幫助.
我相信宣導會也做很多.
有些基金.
特別幫助災後的人民.
這個也是基督徒.
在這個世界上.
要告訴我們.
我們吸引人的地方.
不是因為你有錢.
而是因為你的愛心.
你那份的愛.
怎樣讓人感動.
就是因為你的真誠.
所以我們要很小心.
這個世界吸引我們.
我們回到神那裡.
怎樣為主要過一個生活.
但是我們也要知道.
魔鬼的伎倆從古至今.
都是一樣的.
魔鬼最要我們達到一個目的.
是甚麼呢?.
就是要我們離開神.
教導我們.
不可脫離凶惡.
原文上因為過渡權柄.
本來是加上去的.
但是最後那個不可脫離凶惡.

$^{761}$最大的凶惡.
是甚麼呢?.
就是離開神.
所以魔鬼用很多方法.
疾病的方法也好.
很多人與事的問題也好.
去要我們.
陷入一個最大的試探.
就是要我們離開上帝.
所以有些人.
自己做得不好.
上帝也救不了我.
我沒有面子來見上帝.
我說是誰說的.
如果我們是這樣的話.
你就中了魔鬼的計.
如果今天我們需要貼地也好.
我們相信在上帝那裡.
有很大的恩典.
但是我們也要知道.
上帝也要教導我們.
怎樣持守聖潔.
所以聖潔的生活.
我們要特別注意.
因為聖經教導我們.
所有的試探都在身體以外.
唯獨行淫.
得罪自己的身體.
所以我們要不斷地知道.
怎樣去保守自己.
我們看舊約也是.
巴蘭教巴勒.
也有一個很重要的誘惑.
如果你.
不斷地攻擊他.
或者用一些什麼方法.
是沒有用的.
巴蘭教了巴勒一個很毒的方法.
是什麼呢.
就是要推以色列人自己犯罪.

$^{801}$自己犯罪.
上帝就會審判他.
今天從古至今.
魔鬼的手段都是一樣.
魔鬼也要推我們犯罪.
讓我們被神罰.
因為神不容許.
怎麼會罰我們呢.
就是因為當我們要跌下的時候.
就會用這些方式.
讓我們跌倒.
所以我們要小心.
所以我們要保持眼睛清醒.
我們剛才唱的那首歌.
也是.
我們要求主開我們的眼睛.
免得我們得罪神.
得罪神.
神會審判我們.
或者.
如果我們因為魔鬼的緣故.
生命沮喪.
也會退後的時候.
所以我們要求主幫助我們.
上上都教證.
眼光,我們的焦點.
回到神那裡.
物質的東西,今生的東西.
我們要求主給我們一個新的眼光.
怎麼看永恆的東西.
永恆才值得我們去追求的東西.
願上帝幫助我們.
一個真正認識上帝的救恩.
使得我們得了這個.
不能震動的國.
這個作者再一次教導我們.
我們應該傳感恩的心.
過一個敬虔的生活.
所以過敬虔的生活.
必須要追求和睦.

$^{841}$也需要追求聖潔.
而這個聖潔的起點.
就是我們拒絕.
任何的貪戀世俗.
而這個拒絕貪戀任何的世俗.
不是我們自然會懂得的.
我們不能夠在那裡.
祈禱說信.
求你將試探離開.
聖經從來都不是這麼說的.
聖經教導我們.
第一個方法就是逃避.
面對魔鬼攻擊的時候.
我們就需要一個方法.
抵擋它.
如果在聖經裡面是這樣教導我們.
所以當我們面對試探的時候.
我們不要這麼相信自己.
我們需要逃避這些試探.
逃避私慾.
以致我們能夠保守自己.
這也是聖經教導我們.
所以跟旁邊那個小聲說.
我們追求聖潔.
保守自己.
小聲鼓勵對方.
我們各個追求聖潔.
保守自己的生活.
第三方面你要怎麼去看.
我們需要去.
持守上帝給我們的.
恩典.
特別在這裡講到.
恐怕有人失了.
上帝的恩典.
或者失了對上帝恩典的持守.
以致.
身出獨根.
所以這裡有兩件事要提醒我們.
第一,我們要持守上帝的恩典.

$^{881}$第二,我們需要除掉.
這些獨根.
我們先來看看.
我們怎麼去持守上帝給我們的恩典.
我們就需要反省.
在我們的生活裡面.
我們究竟是站在.
恩典裡面.
還是站在.
律法之中.
這兩個站立是很不同的.
當如果我們今天.
我們以上帝的恩典.
為我們的核心的話.
我們就需要過一個.
對人有恩典的生活.
但是如果我們想.
只是持守律法.
我們經常看到這個對那個不對.
如果是這樣的話.
上帝都以一個律法來審判我們.
所以在聖經裡面講得很嚴肅.
特別是《雅各書》講得很嚴肅.
這件事.
我們今天怎麼持守上帝的恩典呢.
在《雅各書》說.
憐憫原是向審判.
誇盛.
今天如果我們大家都自知.
自己都有很多問題出現.
所以我們都求主給我們.
我們今天真的要讚享.
用恩典與人相處.
用恩典與人.
去面對.
所以我們就用恩典.
我們持守恩典的時候.
其實就自然將這個毒根.
去除開.
但是如果我們不持守這個恩典的時候.

$^{921}$當人得罪你的時候.
你不斷記住它.
你不斷埋怨.
不斷伸出這個毒根的時候.
就會妖害自己的生活.
如果大家都能夠體會.
我們知道我們不想.
毒根要伸出.
我們就要學習.
聖經教導我們.
怎樣用恩典來待人.
恩典待人就是我們用憐憫待人.
我們用寬恕待人.
如果大家去賞心一層.
有個人得罪你.
你會覺得好一點.
舒服一點.
但是有沒有幫助.
其實沒有幫助.
如果我們以惡還惡.
沒有幫助.
但是我們怎樣用恩典來待人.
才能夠幫助我們去面對.
我們怎樣去持守這個恩典.
如果我用一些.
聖經例子.
我用一個生活例子.
來講.
我記得曾經有一個弟兄.
這個弟兄在職場上.
面臨升職的機會.
或者.
他覺得.
這個升職機會很渺茫.
為什麼呢.
特別當時.
如果警隊裡面都是比較救濟的時候.
埋堆.
或者一個小圈子.
是你升職的.

$^{961}$一些比較捷徑的地方.
但是他又堅持不埋堆.
因為他知道埋堆.
沒有什麼好處.
所以一直等.
等他自己做好自己的工作.
但也不需要和別人埋堆.
每一次.
等了三年都落空.
有些人就覺得.
你傻了,埋一堆沒所謂.
有什麼問題.
你和他一起做同一件事有什麼所謂.
但是他自己的堅持.
反而有一次.
發現了一個問題.
就是保護到自己.
他堅持的時候.
某年的廉政公署.
來展開對某些.
他埋堆的群組.
來調查的時候.
個個都出事.
個個都.
俗稱被折濕.
個個都有問題.
但是這個弟兄因為沒有埋堆.
他查也查不到任何的東西.
所以他在當年.
那一年就是成績最高分.
所以那一年就升了職.
如果這個弟兄.
等了三年一念之差.
他就覺得都沒有用.
上帝都不聽祈禱.
他自己就跌入另一個陷阱.
根本就沒有這個經歷.
沒有這個見證.
可以分享.
所以今天我們要去看.

$^{1001}$持續去相信上帝的恩典.
是必須要.
我們對人要忍耐,對上帝都要忍耐.
我們對上帝都要去.
忍耐,等候.
上帝的心意.
我們相信上帝一定有他的方法.
我自己也有一些弟兄姊妹.
他們結婚.
有些丈夫.
有些配偶不是基督徒.
其中有一個姊妹.
她結婚她的丈夫不是基督徒.
她一直耿耿於懷.
她覺得為什麼丈夫不相信.
經常都逼她回教會.
你知道越逼就越怎樣.
越逼越走.
通常越逼沒什麼用.
所以我有一天跟她說.
不如你代表.
耶穌自己.
用愛愛去救.
你不要逼她.
其實我們真的是耶穌.
耶穌在我們心靈裡面.
我們真的代表耶穌.
無條件去愛她.
愛她不是因為你回教會.
愛她不是因為你好.
愛她就是因為你愛她.
就是這麼簡單.
所以你不是逼迫她.
去做你喜歡做的事情.
但是你用愛來融化她.
這個反而才是我們基督徒.
可以去做的.
你不能改變一個人.
聽完之後她就突然間醒悟.
她覺得自己這樣逼迫丈夫.

$^{1041}$事實都沒有什麼好處.
所以她就放手.
放手就奇蹟又來了.
真的當她放手之後.
有一個奇蹟出現.
她突然間找一些東西的時候.
找到她先生一本日記.
她先生日記.
她先生一本寶裡面.
夾了一個有人祝賀她.
決志信主的書籤.
她很奇妙.
她說為什麼會這樣.
於是問她先生.
問她丈夫的時候.
她丈夫說不止爺爺就算了.
但是在她心裡面有一個很大的釋放.
她就知道原來她先生是信了主.
不過她就離開神.
所以逼迫她沒有用.
因為她跟神已經有關係.
她就因為這樣.
這個學習就反而放手.
大家去學一下放手.
我們反而覺得上帝有很多恩典.
讓我們看到上帝的工作是怎樣.
我們不要比上帝快.
我們拿著.
我們根本不能看到上帝的作為.
我們不如放開手.
我們用信心交託給上帝.
讓我們跟神之間的關係.
是更加清晰.
我們不是要逼上帝.
做一些你想做的事.
也不是逼人.
做一些你期望.
他跟神的關係是怎樣.
因為我知道我們活了這麼久.
大家都知道我們人生.

$^{1081}$最能夠改變的就是你自己.
你如果自己都改變不了.
你何必要改變人呢?何苦拿來新呢?.
我們不如學習去改變自己.
不去操控,不去強迫.
將恩典了待人.
我們反而就去經歷.
那份的釋放.
不單止是自己釋放.
你會越來越人見人愛.
因為你不會用迫不得已.
去操控別人.
所以我們要持守上帝給我們的恩典.
持守這個恩典之外.
我們再說.
我們怎樣去除掉這些毒根呢?.
這些毒根.
唯有一個方法.
就是遙恕這個方法.
去處理.
我們無論毒根是.
憎恨也好,怨恨也好.
不肯遙恕.
覺得侮辱了別人也好.
這些都是令我們.
放不下手.
聖經或耶穌說得很清楚.
當我們不肯原諒別人的時候.
我們天父又怎樣呢?.
很嚴重的.
都不原諒我們.
上帝這麼小器的.
不是的,原來上帝是在說.
當我們心裡有這條毒根的時候.
你和我和神的關係.
是不會有好結果的.
所以當我們除掉這個毒根.
就是用上帝的方法.
因為上帝的方法就是遙恕.
遙恕不是因為我們偉大.

$^{1121}$遙恕是因為我們.
跟從神的方法.
神既然都遙恕了自己.
我們都看到自己不是好人.
旁邊那個都是不好人.
但既然上帝都這樣遙恕我.
我怎樣可以遙恕旁邊那個.
不好人呢?其實大家都彼此彼此.
如果我們用這樣的方法.
我們的毒根.
就會除去.
所以我們求主要幫助我們.
我還記得羅斯福.
既然被人偷東西之後.
他能夠用正面的思想.
去接受現實,他能夠憑信心.
去面對生活的時候.
我們今日的教徒.
更加需要有這樣的態度.
去面對,所以我們再跟.
旁邊那個說.
記住.
記住什麼呢?.
毒根不用記了,最重要是.
持守神的恩典.
我們要說.
所以當我們回到希伯來書.
十二章二十八節,你要告訴我.
所以我們記得.
這個不能震動的國.
我記得很多年前.
你看這句話的時候,我覺得很震撼.
我們現在得到一個不能.
震動的國.
這個國度是永遠都不會改變.
所以我們應該怎樣生活呢?.
再提醒,我們應該要過一個.
感恩的生活.
照上帝所喜悅的.
方式,我們用虔誠.

$^{1161}$用敬畏神的心.
來去事奉神.
所以上帝既然賜我們一個這麼大的.
恩典,我們今天.
其實最大的恩典就是.
我們不用怕死,我們不怕死亡.
不怕死亡.
以前那些打仗的時候.
不怕死是最厲害的.
所以我們今天基督徒應該最厲害的.
我們連死都不怕.
所以我們今天面對這個生活的時候.
即是說身體有疾病也好.
人生不如意也好.
和人有衝突也好,面對種種.
苦難也好,事與願違也好.
我們都能夠去體會.
這個根本就是人生的.
常態,自從人犯罪之後.
就是這樣.
但是都不會影響.
我們有一個感恩的心.
這個感恩的心就是我們跳出來.
我們再說就是我們和上帝的.
關係,我們追求和諧.
不是因為我們有多厲害.
而是因為上帝.
怎樣改變我們的生命.
以致我們可以用愛.
來去忍耐那些我們覺得.
不可愛的人,我們能夠.
用和諧來去面對.
那些我們以為不和諧.
本來有衝突的一些人士.
所以我們能夠靠著主.
過一個和睦的生活.
我們都能夠靠著主,我們不被.
這個世界誘惑.
如果我們大家不信主,這個世界一定.
可以誘惑我們,但是既然我們.

$^{1201}$信主,我們有永生的眼光.
這個眼光能夠帶來給我們.
今天面對今生的時候.
我們不會為了.
短視的事情而放棄.
一些永恆的.
福樂,所以當我們.
連魔鬼都不怕.
我們今天看到永恆的時候.
魔鬼都不怕.
你就很強了,魔鬼都.
沒有你扶,我們經常都.
叫魔鬼退到後面去就行了.
有時候我們經常叫他捆綁.
但是我覺得很奇妙.
上帝在天庭裡面.
沒有捆綁魔鬼,可以讓他走.
為什麼?因為上帝.
很相信當我們願意.
去跟從他的時候,魔鬼不用.
被捆綁,因為其實.
魔鬼根本就不會在我們.
這些懂得感恩,懂得.
去倚靠神的人,那裡有什麼.
可以有好處的地方.
最多他可以在那些苦毒.
埋怨,痛苦.
那些人身上去.
做工,所以我們求主給我們.
我們可以因為追求.
聖祭儀師,我們追求同神.
一個親密無間的關係.
的時候,魔鬼就不可以.
在我們生命裡面做任何的事情.
我們再重新說一下.
我們需要持守上帝.
給我們恩典,以致.
毒根不會出來,一出來就.
剪掉,用感恩的心剪掉.
這些毒根不會.

$^{1241}$蠶食我們的生命.
所以我們今天要記住,即使我們.
跟別人意見不一致也好,其實不要緊.
或者環境不理想也好.
也不要緊,最重要是什麼呢?.
目標是.
神自己,我們今天的目標.
我們是不是討神喜悅.
這個才是我們的目標,所以我們求主.
幫助我們,不要被這個世界引誘.
反而我們成為.
這個世界的人的祝福.
這個世代的人的見證.
讓人見到我們.
覺得很棒,基督徒.
是好的,願上帝.
感動我們,幫助我們.
願我們都能夠成為一個感恩的人.
我們跟神同行的人.
我們一起去祈禱吧.
七夕天上的爸爸.
我們要多謝你.
因為你用你何等大的.
恩典來代我們.
你願意寫下你獨身的愛旨.
為我們寫在十字架上.
我知道你付的代價.
是何等的大.
今天我們能夠白白的領受你的恩典.
我們為了你給我們這個白白的恩典.
但今天我們領受你的恩典之後.
我們很願意來事奉你的時候.
耶穌我們求你都光照我們.
讓我們知道.
我們都要把這個恩典.
白白的告訴別人.
但當我們白白告訴別人的時候.
我們就需要付代價.
我們需要付上一個見證的代價.
我們需要付上一個忍耐的代價.

$^{1281}$我們需要付上我們倚靠你的時候.
我們可能面對人事不順的時候.
我們仍然相信你.
這個堅持的代價.
主要我們求你來幫助我們.
有時候我們知道我們自己很眼淺.
耶穌我們求你.
差遣你的聖靈.
你要住在我們心裡面.
你就是要不斷的保衛我們.
不斷的教導我們.
不斷的讓我們能夠進入你的真理.
明白你的真理.
所以我們祈求你來光照我們.
常常的來開我們眼睛.
來明白你話語的寶貴.
讓我們能夠明白我們和你的關係的重要.
主要我們願意.
我們每個人都能夠用心底來講.
我們愛你.
我們願意來隨從你.
跟從你.
讓你的旨意在我們身上來成就.
讓你的旨意在地上行走的時候.
就直著我們去行走.
主要願意來幫助我們.
來帶領我們.
多謝你給洪恩堂不斷的人數增長.
我們都願意.
我們的生命都同樣不斷的扎實.
我們求你來祝福.
無論牧者,無論童工,無論所有的領袖.
都能夠站出來彼此來服侍.
讓新的人都能夠在這個地方.
來看見基督在我們生命裡的作為.
願你親自來祝福.
讓洪恩的教會都成為洪水橋的地方.
一個很大的祝福.
願你帶領.
多謝你聽我們祈禱.

$^{1321}$奉耶穌基督的名.
阿們.
\newpage



\section{羅馬書 12:1-3}
\label{sec:DhmQL4xjGtQ}
\textbf{2026.05.03 《合乎中道的自我認識》:羅馬書12\_1-3節(和修版) 曾裔貴牧師}
\newline
\newline
連結: \href{https://youtube.com/watch?v=DhmQL4xjGtQ}{\texttt{https://youtube.com/watch?v=DhmQL4xjGtQ}} ~~~~ 語音日期: 2026-05-07
\newline
\newline
\hyperref[sec:WokQkPhl0sU]{\small{< < < PREV SERMON < < <}}
~
\hyperref[sec:index_chronic]{\small{[返順時目]}}
~
\hyperref[sec:index_scriptual]{\small{[返順卷目]}}
~
\hyperref[sec:WnFpUR4Lsgw]{\small{> > > NEXT SERMON > > >}}
\newline
\newline
羅馬書 12:1-3
\newline
\begin{longtable}{cl}
\hline
\hline
章節 & 經文 (和合本修訂版)\\
\hline
12:1 & \begin{tabularx}{0.7\textwidth}{X} 所以，弟兄們，我以神的慈悲勸你們，將身體獻上當作活祭，是聖潔的，是神所喜悅的，你們如此事奉乃是理所當然的。 \end{tabularx} \\ \\ \relax
12:2 & \begin{tabularx}{0.7\textwidth}{X} 不要效法這個世界，只要心意更新而變化，叫你們察驗何為神的善良、純全、可喜悅的旨意。 \end{tabularx} \\ \\ \relax
12:3 & \begin{tabularx}{0.7\textwidth}{X} 我憑著所賜我的恩對你們每一位說：不要把自己看得太高，要照著神所分給各人的信心來衡量，看得合乎中道。 \end{tabularx} \\ \\
[1ex]
\hline
\hline
\end{longtable}
$^{1}$請姐妹主女們平安.
和大家一起有祈禱.
天父我們在這個時候安靜我們的心靈.
聆聽你的真道的教誨.
我們真的很需要在神的豐厚恩典.
和你的慈悲憐憫裡面.
去認識我們自己.
求神幫助.
以致我們每一位都能夠.
真是合乎忠道來看自己.
在神面前所領受的信心.
以及我們能夠在神的揀選.
拯救的選照裡面.
來在教會.
在神安排的不同領域.
來成為一個合神使用的器皿.
願神帶領我們.
特別我們今天下午.
都會有我們在洪恩家有三位弟兄姊妹.
接受水禮.
正式來向世人宣認.
是歸屬主.
成為主的兒女.
亦有我們母堂.
嘉欣堂.
總共十幾位弟兄姊妹.
我們見到這是一個出死入生的過程.
藉著一個水禮.
好像浸下水.
然後有新生的復活的樣式.
這本身就是告訴我們.
我們是新造的人.
所以我們很盼望.
神自己來聖靈的工作.
在弟兄姊妹的內心.
亦在我們今早.
我們集中思想這個課題的時候.
讓我們能夠從聖經保羅的教導.
我們去認識.
我們這樣禱告.

$^{41}$奉耶穌基督的名.
阿們.
這個題目也是重要.
今日的訊息很重要的中心思想.
是怎樣來認識自己呢?.
必須要從神的眼中去認識.
我們在昨日.
我們有堂位.
有牧者.
來跟弟兄姊妹.
準備申請洗禮的弟兄姊妹.
來接見,聊天.
我們叫做「民心事」.
就是他內心的信德.
我們每一個堂都有這樣的安排.
如果過不了這個信德.
就不能洗禮.
當然我們沒有這樣跟他們說.
怕嚇到他們.
但如果我們不通過就不能洗禮.
當然我們一致內心都說阿們.
都能夠很清楚聽到弟兄姊妹的生命的改變.
但有一件事很重要.
就是我們每一次聽弟兄姊妹洗禮的時候.
連他們自己都發現.
他們過去舊有的生命.
未認識耶穌之前.
他們無論怎樣掙扎.
怎樣想活.
有時都沒有力量去改變.
剛剛相反.
就是信了主之後.
不知何原因.
就是他們裡面有一個新的生命力.
讓他們看到.
他們就是這樣很容易.
可以貼近神的心意.
來生活.
有很多具體的例子.
是由他們親口跟我們說.

$^{81}$我們聽到的時候都很感動.
看到他們的生命或多或少.
有神在他們內心的工作.
今天差不多是這一堂聚會的訊息.
很想給每一位新洗禮的弟兄姊妹.
再兼顧神的教導.
在神學上.
怎樣能夠看自己是一個新生命.
新的創造.
怎樣在神的眼中去看自己舊有的生命已經過去了.
新的生命怎樣接上.
跟耶穌這個永恆的生命.
在我們的生活裡.
彰顯出來新生命的動力.
所以這裡有幾節聖經.
其實是相當重要.
跟大家再一次重溫這段聖經.
其實羅馬書.
我已經講到在母堂.
我們每個星期的星期五.
就由verse by verse逐節去解釋.
已經講了三十八堂.
還有大概兩堂.
就會完結羅馬書的查經.
接著我們就會開始在《出埃及記》.
很詳細地逐段跟弟兄姊妹講解.
我記得.
我在講到第十二章.
第二十七講的時候.
我正正就用了這個題目.
合乎中道的自我認識.
這段經文的重要之處.
就是它連繫上文的十二章第一節.
和十二章的第四到第九節.
所以第三節是一個相當重要的中心的要點.
就是我們如果在教會裡面.
我們這個新的生命.
我們是一個已經奉獻給神的活祭.
如果我們在第三節出了差池.
出了問題.

$^{121}$教會就有一個很大的虧損.
因為第四節到第九節.
就立即講.
我們要按神給我們各人不同的恩賜.
來彼此的服侍.
我們還要記住.
我們眾多的弟兄姊妹.
但只是基督一個的身體.
我們是在這個身體裡面.
各自作肢體的.
所以你可以想像得到.
保羅他很清楚人性.
包括他自己就是過來人.
他自己擁有非常深厚的.
舊約拉比的訓練.
對舊約的聖經.
對舊約的詮釋.
對律法.
他是非常到位,到家的.
他這樣一個人.
就導致他對新信主的基督徒.
有一個很大的影響.
就是一個很大的錯誤.
就是以為這些基督徒是背叛猶太教.
因為那些是猶太人.
先是猶太人信主.
所以又信一個木匠的兒子.
又信一個釘在十字架上的死囚.
又來到他不再根據摩西的律法守十戒.
這樣的情況下.
保羅就自己都混亂了.
就激發他的內心.
原來是一種出於自然的驕傲.
來到他要將這些基督徒捉拿.
所以當他真真正正明白.
主耶穌為他犧牲.
神在他的愛子生命裡所成就的救恩.
保羅就真的有一個很大的生命扭轉.
他終於認識自己.
要在神的眼中看他自己.

$^{161}$這裡告訴我們.
我們可以將第三節一起讀一讀.
我就會再詳細講解.
我們一起讀這十二章第三節.
「我藉著神賜我的恩典.
對你們各人說.
不要對自己評價過高.
超過所應有的評價.
都要按著神量給各人的信心.
中肯地評價自己」.
他告訴我們有一把尺.
不是我們自己的尺.
不是我們和其他人量度的尺.
是神的尺.
這個尺是神給我們的信心.
原來是有大有小的.
所以神給你的信心的尺去量度自己.
你在神的面前.
如果你已經照著保羅的教導.
你已經來將自己的生命.
因為是屬於神的生命.
來奉獻.
作為一個活祭.
奉獻了給神.
你又開始來看.
當你的眼光開始集中.
要在教會裡面參與不同的服事.
開始你在教會也來得久了.
你開始會有不同的崗位的服事.
和無論你與生俱來的傾向.
你原先有的天分.
後天的訓練.
不同的恩賜.
你開始很想在教會大展拳腳的時候.
這樣本來是很好的.
是神閱立的.
但是保羅才從神學那裡.
給我們一個很重要的提醒和裝備.
如果我們沒有這個神學的醒覺.
沒有對自己為何會在教會坐下.

$^{201}$為何會每個星期來教會聚會.
我們沒有很認真的這套神學.
去讓我們的內心有這樣的認知.
很多時候往往在教會裡面一相碰.
我們就會很容易產生很多火花.
因為第四節之後.
就立即告訴我們.
我們在主的裡面.
各人得到不同的恩賜.
在不同的恩賜裡面.
保羅教導我們要怎樣專一.
在這些恩賜上面去服事.
所以我們求主讓我們能夠看到.
這個合乎中道.
就是說你記住.
不是看自己過去舊有的那把尺.
去量度自己,量度他人.
量度自己和別人的不同.
而是神給我們的信心的尺.
去量度自己.
這個就很抽象了.
這個就很違反我們本來.
信耶穌就是信耶穌.
我回來教會.
我熱心幫忙做事.
你也有這麼多要求嗎?.
其實不是.
原來真正的能夠長久事奉.
而且不斷在事奉裡面帶來.
對自己,對其他人都有助益的.
那個原先最重要的.
就是對這個真理的認識.
這個合乎中道.
這個字是很普遍.
在希臘文的原典.
以前的哲學家.
柏拉圖等等.
都用這個字.
合乎中道.
就是你要很有理智地.

$^{241}$很有sense地.
去謹慎地思想.
有自制地去判斷.
就是說強調.
這種是一個很清醒.
很謙卑.
很合宜的對自己的認識.
這個希臘文.
這個我不認識了.
已經放下很多年了.
希臘文.
合乎中道.
你要這樣去理解自己.
相信我們能夠達到.
這樣的一種認識.
必須像保羅所說的.
就是我們要在神的救恩裡面.
去看自己.
這個正確的認識.
所以首先.
必須要從神學的部分.
去理解自己.
什麼叫從神學的部分呢?.
又說得這麼深奧嗎?.
其實不深奧.
保羅一直在說.
十二章之前.
大家有沒有留意.
十二章的第一節.
中文已經告訴我們.
是有一個「所以弟兄們」.
是有一個很重要的連接詞.
連接十一章.
保羅告訴我們.
救恩的偉大.
是神自己的發明.
是神自己犧牲的愛子.
是神自己在萬世以前.
已經預定了我們.
得救,揀選我們.

$^{281}$使我們今天能夠.
接受神一個救恩的邀請.
來得救.
所以不是我們的努力.
是神的恩典.
是神一個很重要的慈悲.
是神的憐憫.
所以我們就能夠.
在我們實際生活的場景.
我們就能夠將自己.
願意在救恩底下.
來成為一個活祭.
獻給神.
所以要告訴我們.
真正正確的認識.
是要讓我們能夠認識到.
神到底為我們這個人.
做了什麼呢?.
今天三位弟兄姊妹.
暫時見到Amy.
Morris在不在樓上呢?.
另外就是周仔.
又未見.
崇拜都不回來.
就真的要罵了.
真的要認識.
必須.
待會你們說見證.
必須有一樣東西是很重要.
神到底為你們做了什麼呢?.
不是哪個教師.
不是哪個主學老師.
點明了你們一些道理.
而是我們要記得.
一件很重要的事.
造物主到底為我們.
這個準備洗禮的人.
做了什麼呢?.
是因為他做了.
這件影響整個世界.

$^{321}$人類的命運的事.
就是我們的救主.
為我們白白犧牲在十字架.
是這個原因.
使我們成為他的兒女.
所以保羅告訴我們.
是出於神的憐憫.
羅馬書九章到第十一章.
完全不是我們的努力.
是神的憐憫.
他要憐憫誰就憐憫誰.
他要恩代誰就恩代誰.
沒有人可以向神.
作出抗議反對.
他是用出埃及的摩西和法老.
在十災裡面的對話.
保羅是這樣說的.
他要令到法老的心乾硬.
這裡告訴我們.
神就是這樣為我們人.
預備了一個救恩.
所以救恩完全是出於神的憐憫.
不是人的行為.
所以如果我們明白.
這是神的憐憫的時候.
我們就不能夠自告.
更加不能夠在神的面前自告.
當這樣的心態.
裝備好神學的認識的時候.
每一個弟兄姊妹.
再用自己在神給你的信心的尺.
去量度自己.
那就不同了.
教會就非常精彩.
因為不同的弟兄姊妹.
能夠成為一個合一的團隊.
能夠在不同的崗位.
不同的恩賜.
各自發揮神給的服侍機會.
所以我們一想到服侍.

$^{361}$我們首先大前提.
就必須要看到.
神為我們每一個弟兄姊妹.
做了什麼.
這是很重要的.
包括對我們自己一生的服侍.
都是很重要的.
大家認識這個人.
有沒有人知道他叫什麼名字.
幫我打出來.
其實不是.
他應該叫做放天佑.
是這個劇集.
令到他名利相收.
大家聽過那個劇集沒有.
我和殭屍有個約會.
令到他大紅大紫.
令到他完全在亞視成為亞視的一哥.
亞視是一個很弱的台.
他能夠全香港人都認識他.
追看這套劇集.
但在這個時候.
有人跟他說福音.
他講見證.
當然他是順了主之後講見證.
他說我那時候紅到一個地步.
那些片藥那些錢.
加人工.
他說誰跟我說福音.
我甚至用粗口罵他.
嘲笑他.
這個世界為什麼需要這個神耶穌.
但好景不常.
在他最紅的時候.
太太染出血癌.
他整個人崩潰.
完全不知道怎麼辦.
不停流眼淚.
他很愛自己的太太.
不知道為什麼.

$^{401}$後來同樣都是那些人向他說傳福音.
他突然發現自己.
原來一點都不可以靠自己.
儘管聽一下那個教友.
去講去聽.
接著有一天在醫院的房間.
他無助地跪下向神祈禱.
一祈禱.
他說地上整個地板.
那個階磚就是他的眼淚.
他向這位真神祈禱.
他說我不認識你.
如果你是真的.
你醫治我的太太.
我甚至連命都願意奉獻.
接著第二天.
他打開窗簾.
在醫院裡面.
看到有很漂亮的一條大彩虹.
他覺得很奇怪.
就問那個教友.
為什麼會這樣.
其實是很自然的現象.
他自己可能想多了.
那個教友就跟他說.
是神昨晚聽了你祈禱.
應允你.
接著沒多久.
他太太的血癌.
經過很多標靶藥.
真的完全康復了.
令到他發現.
這個神我不能不信他.
他從此.
由他太太開始康復之後.
到今天.
他不斷跟人說.
我信這位主.
我不會再改變.
去到一個地步.

$^{441}$他作見證.
那些人在街上.
叫他的名字.
放天佑.
那個電影太出名了.
他說你不要叫我這個名字.
我不要這個名字.
我的名字叫尹天照.
或者叫我阿趙也好.
他說你不可以再叫我這個名字.
為什麼呢.
他說他覺得這些電影.
會影響別人的迷信.
他從此不再拍這些殭屍片.
不再拍這些鬼怪的電影.
剛好有一些片商.
邀請他拍電影.
給天價.
甚至乎說.
你開個價吧.
因為太紅了.
他說很簡單.
決絕.
他說這些錢我不賺.
我答應了神.
親愛弟兄姊妹.
生命的改變.
就是認識自己.
認識自己原來在神的眼中.
無論你多麼大紅大紫.
你多麼的.
好像你能夠駕馭自己的生命.
原來很少的東西.
一個家人.
你都不能夠照管.
你都不能夠保護.
但是感恩.
神改變他.
讓他有一個很堅定的信心.
來去依靠神.

$^{481}$他今天都不斷告訴別人.
這個耶穌.
他願意付一切的代價.
去跟隨.
去相信.
去見證.
就是這樣一個人.
你可以想像得到.
有時明星為何會有這種效力.
因為他們自己的生命.
他是一個公眾的人物.
他們自己的經歷.
神給他們有一個作見證的機會.
每個人的信心大小都不同.
神量給每一個人.
每個人有你自己信心的尺.
所以你的尺有多少.
你是跟自己去量.
在神給你的恩典.
恩賜裡面.
所以保羅說.
我照著神給我的恩典.
來向你們眾人說.
不要看自己過於所當看.
不要看高自己.
往往會看高自己.
就會有一種自負.
就會有一種.
不知不覺跟其他人做比較.
這樣是最不好的.
我初信的時候.
根本不懂這些真理.
信得很熱心.
很深刻.
剛剛初信主二十多歲.
很大的一個劇烈的轉變.
不吸煙.
不再賭錢等等.
跟家人關係又復和.
有一次去完太子道的種子書室.

$^{521}$十二樓上面.
舊樓.
下來地下準備過馬路.
過對面旺角警署.
明明有一個跟我一起乘電梯.
在種子書室裡面的弟兄.
不認識他的.
一起下電梯.
一起過馬路.
那時候我們很嚴格.
到現在應該都是.
紅燈要氣定不准動.
等紅燈過.
到綠燈你就過.
誰知我見到那個弟兄.
紅燈照過.
對我來說就是一個很晴天霹靂.
信耶穌可以衝紅燈嗎?.
在我的心裡有沒有搞錯?.
剛剛跟我一起乘電梯下來.
這樣的.
我初信的時候就是這樣.
當然不是說我現在可以放鬆.
紅燈我都過了.
不是這個意思.
而是原來我不知不覺就有一把尺.
就去量度別人.
原來慢慢我信主內.
我發現原來保羅告訴我們.
的而且確我們是需要量度.
是神賜給我們信心的大小.
去度量自己.
自己的人生.
自己在神面前所領受的恩典.
我們有沒有盡力為主而活呢?.
譬如教會.
我們實在有很多事工正在開展.
稍後報告的時候.
和剛才我們祈禱會的時候都有提到.
我們有很多事工正在開展.

$^{561}$是在祝福街坊.
駱橋樓,洪福村.
和現在我們青年會小學.
在三樓上面.
今天已經有九個同學來了.
我們有很多事工正在開展.
我們需要很多弟兄姊妹.
我們呼籲不同弟兄姊妹.
按你的才幹.
因此去服事.
但是大前提.
如果你根本從來沒有想過.
你要有在神面前領受的.
信心的大小.
來合乎忠道去看自己.
其實你根本不適合事奉.
因為你根本從來沒有想過.
原來我們能夠事奉.
是要放下我們自己的尺.
來在神面前.
領受祂給我們的恩典.
我們去服事.
再用另一個例子.
大家認不認識這個人?.
這是一個黑人.
其實他是海地裔的後裔.
他父親是海地裔的美國人.
但是很無奈地.
他父親娶了一個日本女子.
所以她的名字叫大阪直美.
你也沒什麼.
明明是日本的名字.
也就是說她是入了日本籍.
這樣就出事了.
原來日本人的文化.
這樣的膚色.
就導致她不斷受歧視.
去到一個地步.
她有一個很出色的網球的見象.
曾經也拿了大冠冠.

$^{601}$就導致日本人看她.
只是叫做half.
她就背著這樣的身份.
入籍了日本之後.
代表日本出戰奧運.
東京的奧運.
但是去到一個地步.
就是抑鬱.
出現了嚴重的抑鬱病.
在一些比賽.
規定那些落敗的球員.
都要出席記者會.
她居然抗議不出席.
她提出的理由就是.
我失敗已經很慘了.
還要我出席記者會.
是由她第一個發起.
她不出席.
那就更加帶來.
她這個日本籍的身份.
導致她面對很大的身份危機.
去到嚴重的抑鬱.
甚至比賽的水平.
下降得很厲害.
原因就是她爸爸跟隨媽媽.
入了日本籍.
本來可以選擇美國籍.
還是日本籍.
還有一段長時間.
她入了日本籍之後.
但是爸爸回了美國.
她跟隨爸爸.
從五歲開始.
一直在美國成長.
所以全部說的是英文.
日文完全是「啦啦卡卡」.
說不到日文.
所以更加令日本主流的文化.
對這個代表日本出戰奧運的網球選手.
有很大的要求.

$^{641}$過分的苛刻對她的要求.
導致她對自己的身份.
出現了一個很大的危機.
所以去到嚴重的抑鬱.
也在美國的時候.
她認識了高比拜人.
但是因為有一次高比拜人的飛機禍害.
當場死了.
對大阪直美是一個更嚴重的打擊.
所以她就退出了網球的比賽.
甚至打算放棄打網球.
什麼可以救她呢?.
因為她看打網球是比賽賺錢.
還要讓日本針對她的人無話可說.
但這樣帶來的只是更大的衝擊和抑鬱.
後來她在美國結婚生了兒子之後.
她終於發現原來我不需要活在別人的陰影.
這樣的無理的要求下.
我是黑色的皮膚.
不關我的事.
我是海地裔.
我爸爸是美國人.
美國籍.
我媽媽是日本人.
不是我的罪.
她終於想通了.
她終於能夠因為她的小朋友.
她能夠重拾打網球的熱誠.
原來走過了一段很長的抑鬱的路.
終於發現原來我不需要活在別人的影響下.
一個日本籍的黑人都可以.
親愛的弟兄姊妹.
她未必相信耶穌.
但如果我們有主耶穌.
有聖靈在我們的生命裡.
我們必定會經歷身份的轉移.
我們就是神的兒女.
我們的生活.
可能周圍的人告訴我們.
我們未必像十字架上的主這樣的生命.

$^{681}$但是我們不要忘記.
因為一個千真萬確的事實.
我們已經成為了神國度的兒女.
我們已經有這個永恆的護照.
是神的聖靈已經落印在我們的內心.
我們是屬神的.
所以我們應該怎樣來到去.
不要理會周邊的人怎樣看我們.
反而我們要活得像神給我們的恩典的生命.
最後的總結.
我想我們一定要明白.
我們要意識到.
我們的信心.
我們的地位.
我們的恩賜.
我們的才幹.
全部都是神憐憫的禮物.
是一份憐憫的禮物.
不是我們自己努力而得到的獎賞.
不是我們去換取回來的一個獎品.
所以我們更加沒有可能.
看自己過於所當看.
驕傲的本質.
本來就是忘記恩典.
忘記憐憫.
所以求神幫助我們一段很短的彼此的分享.
今天主要真的要向洗禮的弟兄姊妹.
有一個很重要的提醒.
就是真的不要太過於看自己過高.
我們必須要知道.
我們有一把尺.
是神已經量給我們.
用信心量給我們.
所以我們要慢慢在神的恩典裡.
去專注神給我們的事奉.
每個人有神給不同的事奉.
各自去作之體.
但是在一個合一的神的教會.
即是耶穌基督的身體.
但願神祝福我們.

$^{721}$今天是一個很短的訊息.
希望大家能夠一起彼此的互勉.
在教會裡面.
仍然我們說.
我們需要不同恩賜的弟兄姊妹.
去讓神的家得到各方面的豐富成長.
我們一起禱告.
天父讓我們能夠在你的豐厚恩典裡.
你的慈悲.
你的憐憫裡.
我們能夠成為你的兒女.
多謝主.
我們好像保羅的教導.
我們真的需要將身體獻上給神.
做一個活祭.
而且這個活祭.
是我們每天都需要這樣思想.
這樣擺上的.
願神幫助我們.
但願我們能夠在保羅的教導底下.
認識我們要看自己合乎中度.
是要很理智地.
很謙卑地.
這樣來看自己.
今天所得到的福氣.
所應該在教會裡面.
和其他不同的弟兄姊妹.
一起配合來服侍神.
願主帶領我們.
再一次我們慰問到今天.
我們有三位弟兄姊妹.
真的不同的年紀.
不同的背景.
不同的經歷.
但同樣得著這個永遠改變不了的救恩事實.
願你得到一切的榮耀.
也願我們能夠繼續在神的裡面.
這樣來服侍.
多謝你聽我們的感恩禱告.
奉耶穌基督的名.

$^{761}$阿們.
多謝各位.
\newpage



\section{彼得前書 3:13-22}
\label{sec:WnFpUR4Lsgw}
\textbf{2026.05.10 《聖道為軸的崇拜》:彼得前書3\_13-22(和修版) 譚靜芝博士}
\newline
\newline
連結: \href{https://youtube.com/watch?v=WnFpUR4Lsgw}{\texttt{https://youtube.com/watch?v=WnFpUR4Lsgw}} ~~~~ 語音日期: 2026-05-14
\newline
\newline
\hyperref[sec:DhmQL4xjGtQ]{\small{< < < PREV SERMON < < <}}
~
\hyperref[sec:index_chronic]{\small{[返順時目]}}
~
\hyperref[sec:index_scriptual]{\small{[返順卷目]}}
~
\hyperref[sec:W5xhymK0FGw]{\small{> > > NEXT SERMON > > >}}
\newline
\newline
彼得前書 3:13-22
\newline
\begin{longtable}{cl}
\hline
\hline
章節 & 經文 (和合本修訂版)\\
\hline
3:13 & \begin{tabularx}{0.7\textwidth}{X} 你們若熱心行善，有誰會害你們呢？ \end{tabularx} \\ \\ \relax
3:14 & \begin{tabularx}{0.7\textwidth}{X} 即使你們為義受苦，也是有福的。不要怕人的威嚇，也不要驚慌； \end{tabularx} \\ \\ \relax
3:15 & \begin{tabularx}{0.7\textwidth}{X} 只要心裡奉主基督為聖，尊他為主。有人問你們心中盼望的理由，要隨時準備答覆； \end{tabularx} \\ \\ \relax
3:16 & \begin{tabularx}{0.7\textwidth}{X} 不過，要以溫柔、敬畏的心回答。要存無虧的良心，使你們在何事上被毀謗，就在何事上使那些凌辱你們在基督裡有好品行的人自覺羞愧。 \end{tabularx} \\ \\ \relax
3:17 & \begin{tabularx}{0.7\textwidth}{X} 神的旨意若是要你們因行善受苦，這總比因行惡受苦好。 \end{tabularx} \\ \\ \relax
3:18 & \begin{tabularx}{0.7\textwidth}{X} 因為基督也曾一次為罪受苦 ，就是義的代替不義的，為要引領你們到神面前。在肉體裡，他被治死；但在靈裡，他復活了。 \end{tabularx} \\ \\ \relax
3:19 & \begin{tabularx}{0.7\textwidth}{X} 他藉這靈也曾去向那些在監獄裡的靈傳道， \end{tabularx} \\ \\ \relax
3:20 & \begin{tabularx}{0.7\textwidth}{X} 就是那些從前在挪亞預備方舟、神容忍等待的時候不信從的人。當時進入方舟，藉著水得救的不多，只有八個人。 \end{tabularx} \\ \\ \relax
3:21 & \begin{tabularx}{0.7\textwidth}{X} 這水所預表的洗禮，現在藉著耶穌基督的復活拯救你們，不是除掉肉體的污穢，而是向神懇求有無虧的良心。 \end{tabularx} \\ \\ \relax
3:22 & \begin{tabularx}{0.7\textwidth}{X} 耶穌已經到天上去，在神的右邊，眾天使、有權柄的、有權能的都服從了他。 \end{tabularx} \\ \\
[1ex]
\hline
\hline
\end{longtable}
$^{1}$和大家一起思想復活耶穌核心的訊息.
聖餐其實為我們擺開了整個救贖計劃.
並且保羅指出聖餐不只是一個紀念.
聖餐是一個宣告我們主的福音.
因為聖餐也指明了最後還有一餐.
所以每一餐我們更近一餐.
用一個歡喜快樂的心.
在在世盼望主的應許要再來.
並且我們要祂一同坐直.
所以我很怕基督教經常把復活當成是復活主日就完了.
其實祂不斷地在四十天裡.
向祂的門徒多次顯現.
日漢話也記載不完.
為的是什麼呢?.
要我們信.
信祂為我們死.
為我們復活.
並且升天坐在父的右邊.
所以復活這個訊息.
福音的訊息.
我們要天天聚在一起的時候.
不只是奉主的名.
而且是再一次放在我們面前.
教我們如何在世上活.
其實這三段經文包括詩篇.
是普世教會在復活期.
在今個主日裡讀的經課.
但是對於沒有經課.
什麼叫經課呢?不知道.
我們就會覺得.
只講一點點就行了.
如果講幾節經文.
你知不知道要講八十年才能講完整本聖經?.
就是說在基督教教會.
很多時候舊約都被忽略了.
其實舊約是佔了整個聖經的三分之二.
慢慢地我們就會遺忘了.
甚至我們和以色列的群體脫離了.
而詩篇就成為一個貫穿舊約到新約的主軸.
幫助我們再記得耶穌對他的門徒.

$^{41}$很迷糊的門徒.
正是離開耶路撒冷.
在二萬五師路上.
重新來跟他提起.
摩西的律法.
先知的書.
詩篇都有指著我的話.
在今天的基督教教會.
這些經文去了哪裡呢?.
我們只是著重新約.
甚至我們聖餐裡只讀某一節的經文.
我們慢慢地遺忘了.
甚至他都不公開宣讀.
我們會不會不見了主軸呢?.
所以我就和大家分享一下.
崇拜是什麼.
很多時候我們覺得我們的崇拜在於.
那個講員講得好不好.
我的崇拜就好不好.
或者那個詩班.
或者各樣東西在服侍弟兄姊妹.
他們的態度.
他們做得到位嗎?.
於是我們就覺得.
這就是一個好的崇拜.
崇拜在神的眼裡.
從來都是人.
不是事.
他向摩西.
向法老.
要這些人去事奉我.
這個事奉就是我們今天所說的敬拜.
敬拜主.
敬拜是一個動詞.
不是一個名詞.
所以我們不是參加一個崇拜.
我們去崇拜.
所以如果我沒有崇拜.
其實就沒有崇拜發生.
但是很多詩歌又說.

$^{81}$我要唱歌.
我要做什麼.
只有我又沒有用.
所以耶穌就和這個井旁婦人.
一個很普通的人.
甚至一個在社會階層裡.
就是隱蔽的人.
來探索崇拜之道.
也就是說.
告訴我們.
我們每個人.
都是祂宣講的對象.
祂要告訴我們.
什麼是崇拜.
崇拜第一要認識神是靈.
所以不是憑案見到的東西.
或者什麼物件.
或者什麼地方.
或者什麼祭物.
來進行崇拜.
所以崇拜的對象是神.
而貫穿整個聖經.
從舊約到新約.
都是說這個關係的建立.
是與神關係的建立.
那誰做主動的呢?.
崇拜是神自己.
好像在伊甸園裡面.
神是叫亞當.
亞當你在哪裡?.
祂是叫我們的名字.
摩世摩世.
薩姆爾薩姆爾.
我們是被召的.
我們來.
所以我們是回應上帝的召喚.
我們被聚在一起.
我們提我們是一個身體.
是很不同的肢體.
各有不同.

$^{121}$所以不同不是問題.
主是召我們那位.
就將我們放在一起.
所以我高興不是因為我認識你.
是因為我們在主裡面.
我們是一體的.
而在聖經的從頭到末.
天上的敬拜.
都是圍繞在主為中心.
如果主不說話.
我們是沒得回答的.
所以你的唱歌其實是來回應祂的.
所以我們要述說.
究竟祂要重新記憶.
祂做過什麼.
祂的作為.
祂的歧視.
要談論.
我們會由中來回應.
讚美祂.
高舉祂.
尊崇祂.
所以.
聖經其實是說一個題目.
就是萬有都屬祂.
祂是萬有之主.
我們的敬拜無論到什麼程度.
甚至付出多少.
都不配.
但是我們要盡心盡性.
盡意盡力地做.
如果我們記得崇拜是一個動詞.
於是自己就是崇拜的人.
我們要找到對象.
很清晰的.
祂叫我們成為有靈的活人.
在基督律裡面.
叫我們甦醒.
醒覺祂是神.
祂是我們的主.

$^{161}$我們就有對話的可能性.
而且層次就對了.
我們以為回來滿足.
牧師叫我們.
某個傳道叫我們.
或者某個教友.
我們熟悉的.
我們回來社交一下.
我就忘記了什麼叫做靈交.
靈交是找對象.
是神.
我們又知道我們熟悉祂.
又知道祂找我.
耶穌最終終生目的就是說.
神是靈.
所以.
我們在看這種對話是什麼呢?.
首先主角就是神.
然後我們才知道.
因為祂是這樣.
所以我們在靈裡面跟祂相交.
所以在這個過程裡面.
又看見神不單止是神.
不只是被動地坐著.
祂是耶穌說.
「負要者人的人」.
這個「要」其實背後是說.
原文是尋找.
去找真正拜祂的人.
所以我希望我們每個人一生都要認證.
神啊,你找到我了.
我願意成為一個真正拜你的人.
所以當你走在這個路上.
說要去教會了.
我們來崇拜了.
大衛提醒我們我們心裡是歡喜.
因為有人跟我們一起.
而且一切都是被照的.
來成為一體.
來向這個世界做見證.

$^{201}$其實這一班人個個都不同.
但是聖靈可以這樣來愛我們.
貫穿我們.
將我們連在一起.
成為一個合一的見證.
彼此相愛就使人認出我們是主的門徒.
所以我們一同讚美.
一同手撐.
這是一個很重要的標記.
不是因為愛足球.
不是因為愛乒乓球.
要一起歡呼.
不是的.
那些是過去的.
但是我們的「一」是無處不在的.
我們的連結.
我們的相愛.
是不會被限制和停止的.
因為聖道神的兒子在我們裡面.
聖靈在我們裡面運行工作.
因此崇拜本身就是一個見證.
你不要說它只是靠說.
那個本身就是一個相愛合一的行動.
以致我們能夠成為世界的光.
世界的見證.
它向卑微的人揭露一個這麼高的奧秘.
今天我們也一同來領受這個奧秘.
尤其是在聖餐裡面.
不過我今天很想強調.
不只是聖餐.
是聖道.
舊約有一個以色列認為的忠心.
就是「以色列,你要聽」.
不是口語.
是書面上「你要聽」.
這個「聽」不是聽聞.
是舊時的「聽」.
是沒有文本的.
是聲音.
你回到家裡.

$^{241}$不知道有沒有人.
你是不是會問.
有沒有人?.
那你想要聽什麼?.
聲音.
因為聲音就是臨在.
就是在.
有人回應你.
我在這裡.
我在房裡.
我在廚房裡.
我在哪裡.
有聲音就是有人.
神跟我們說是聲音.
你知不知道剛才我們一同讀經文.
其實也很麻煩.
因為我們沒有預備.
真正應該預備的是誰呢?.
我們要將文本變成聲音.
因為本來就是聲音.
聽到了嗎?.
所以第一個聲音不是講道.
第一個聲音是神的話.
被讀出來.
所以神的話被讀出來.
我們一同來聽.
就是我們最核心的崇拜.
以色列,你要聽.
來什麼?.
發現神是怎樣的.
神跟我們說什麼呢?.
所以不是一定要經過一個人.
來告訴你.
神跟你說話了.
所以讀經.
讀經源的預備是很重要的.
因為你要先讓神的話跟你說話.
然後你成為一個聖器.
將這個聲音,語氣,態度,量,質,圖畫.
全部將它用你的聲音演繹出來.

$^{281}$讓人聽見神憤怒.
聽見神的憐憫.
聽見神恩待我們.
聽見神斥責我們.
聽見神等候我們.
聽見神叫我們.
所以我覺得讀經源是一個.
哇,神很憐憫,恩待的即是.
我希望洪恩堂會建立一個讀經組.
真的願意聽神的話語.
願意將自己奉獻給神.
成為神話語的出口.
這種人聚在一起.
你讀給我聽,我讀給你聽.
學聽.
首先要學聽.
你聽到了.
你才可以讓人聽到神的聲音.
所以第一個崇拜是要醒覺神是誰.
然後我們來回應他.
然後學聽他的聲音.
崇拜的內容是神主動來尋找我們.
神在主導的位置.
不是我.
所以發現有很多古老的聖詩.
都不出現「我」這個字.
要麼是「我們」,要麼是「上帝是怎樣的」.
在詩篇裡才發現.
詩人在雙交的位置.
不單是為以色列宣告.
這位以色列的聖者是怎樣.
讓他看見一個神觀.
很清晰的神觀.
所以崇拜的開始.
最重要是要選擇甚麼詩.
祈禱是怎樣.
全部都是神觀.
神是怎樣.
神說甚麼.
神怎樣招聚他的兒女一起.

$^{321}$我們才醒覺.
我們不是找人.
我們不是要見哪個人或約哪個人.
和哪個人說話.
是和神.
以致我們整個氣氛不需要安靜.
因為我們等候.
聽見的是甚麼呢.
想聽見神說甚麼.
這個會改變我們的生命.
替我們的方向有指導性.
我們能夠懂得反省自己.
將自己真正地帶到神面前.
讓祂光照我們.
所以我們的目的一定很清楚.
然後呢.
自古以來.
神就藉著眾先知曉與列祖.
這些列祖要存在.
否則慢慢地就變成了我的教會.
我們投入到我來管的.
我捐錢多就大聲一點.
不是.
這是神的教會.
而且是普世超越時代的.
而且是一路匯聚到天上的.
在這個行列裡.
但當我們不讓眾先知說話的時候.
我們同時也不讓神的兒子說話.
因為希伯來書第十一章.
已經很清楚地告訴我們.
古時候神藉著眾先知多次多方向列祖說話.
末世藉著祂的兒子向我們說話.
所以耶穌在很多講話裡.
你會發覺祂是實實在在地告訴你.
是甚麼?.
是想你分辨出其實是誰在講話.
以及認真度需要聽的人一定要提高醒覺.
所以讀經後都會說.
這是上主的話.

$^{361}$不是我講的.
所以你不要看講的那個人或讀的那個人.
要知道分辨.
這個不是人講的.
是神在講話.
神對我們世人講話.
所以我們將人的話和神的話分開是很重要的.
神原有的從聖經出來的話.
高人一等的話.
應該備受尊重.
何況六十六卷裡面.
福音書包括《使徒行傳》.
那個是福音來的.
是給世界的福音.
天使說今天要報好的訊息.
是給世人的.
是給你們的.
這個訊息我們要很醒覺.
不是給我們某某堂的人.
不是給某個團體的人.
是給世界的.
所以教會有這樣的責任.
向世界宣告神講甚麼.
因此六十六卷裡面的福音書和《使徒行傳》.
是記載了道成肉身的主.
在我們中間行動,教訓,醫治.
因此普世教會有特別做尊重.
站起來聽.
摩西在《申命記》很清楚.
所有南丁以色列人.
聽得能明白的都來聽.
要聽甚麼.
要聽神的律法.
他當時不是坐在那裡聽.
是站在那裡聽.
當年希米恢復.
找到神的律法書的時候.
是站在那裡聽.
所以這是一個很重要的尊敬.
捍衛生命的中心.

$^{401}$那是生命之言.
耶穌自己這樣說.
我的話就是靈,就是生命.
所以我用一個身體語言來表達.
這是道中之道.
你能夠想像一個教會的團體.
認識聖道是很被尊重,很寶貴.
他們聆聽,要提醒自己.
這是神的話.
不是人說的話.
所以我就說.
讀經是比講道更高的.
英文說是primary word.
最基礎,最基本,最核心的.
聖言.
然後講道只不過是將它剪開來.
讓我們能夠明白.
所以我們首先要學聽.
就不是聽道.
是聽聖經.
所以我就覺得.
讀經完就有責任了.
就要好好地預備.
好好地聆聽.
首先不是讀,就是聽.
你看文字的時候要有聲音.
如果你這樣看不到.
裡面沒有聲音的話.
其實你不是進去的.
你回去試試.
無論你來靈修.
或者讀什麼文字.
其實裡面有聲音.
讀給你聽.
當你知道這是聖靈的工作.
聖經是聖靈的工作.
在眾人的心中來指示他們.
你也祈求這個真理的靈在你裡面.
剛才不是說派來保衛師嗎.
是不是.

$^{441}$祂就會指導你進入一切的真理.
多重要啊.
我們有個老師教我們的.
他會讀給我聽的.
有人問我你讀經是誰教你的.
我說是聖靈教我的.
我聽見祂的聲音.
我等待祂的聲音.
懇求祂教我.
讀給我聽.
所以最重要的就是.
我們有沒有聽的耳朵.
我們有沒有聽的圓耳的禱告.
然後我們就慢慢更明白.
我們的耶穌.
正正是不斷在腹的右邊.
為我們祈求.
你看見天上的禱告不斷地.
但是地上的禱告.
也是因為真理的聖靈在我們裡面.
是這樣平衡地去.
互相呼應.
我們更加明白.
耶穌是怎樣馴服神.
我們又怎樣馴服主.
使得我們更加.
不是從外表.
不是藉著詩歌這麼簡單.
是從裡面馴服神的聖道.
當我們知道神是靈.
祂又呼喚我們.
我們應該怎樣呢.
而塞瓦先知就這樣告訴我們.
當耶和華這樣去問祂的問題.
我來的時候為何沒有人等候.
我呼喚你們.
為何你們不回答.
原來神是在等待我們.
不單是尋求我們.
成為真正拜祂的人.

$^{481}$是在看我們有沒有等祂.
我們沒有來回答祂.
我帶人查經.
快四十年了.
我只有一個很清晰的目標.
我們小心地來聆聽.
幾節的經文.
對我們說什麼.
但最重要的.
其實是我查完經之後的禱告.
我教人怎樣聽了這個經文之後.
我們應該怎樣回應上帝呢.
怎樣回答祂呢.
怎樣像一面鏡照見自己呢.
怎樣告訴祂我聽明白呢.
然後告訴祂我應該怎樣做呢.
承諾祂.
表示我聽明白了.
你有沒有知道其實你讀經最重要.
其實神在等什麼.
在等你認為如何呢.
如果我們沒有回答祂.
我看完了.
今天做了五分鐘.
雞精了.
靈修五分鐘.
什麼叫雞精呢.
二十四小時給五分鐘.
就是只有五分鐘.
然後讀完之後.
轉面離去.
就等於我們崇拜完一樣.
如果我們聽見上帝跟我們說話.
我們不能一言不答.
就轉身離去.
你知不知道養一個.
年輕人在家裡.
吃飯吧,不理你.
吃飯吧.
第三次不見人就自己吃飯.

$^{521}$然後你要做什麼就轉頭離去.
一言不答就轉頭離去.
你認為他有沒有聽見.
多數我媽媽會覺得他沒有聽見.
因為他不回答.
所以我們也不能不回答.
所以崇拜.
聽見神的話之後應該怎樣.
是回應.
用什麼回應呢.
用禱告回應.
用代求回應.
我們明白上帝要在這個世界做什麼.
又已經在我身上延展這個工作.
我應該在我的家.
我的網絡.
我工作的地方.
在這個社會.
為這個國家.
為這個世界.
萬邦萬國.
我們唱歌.
祂是萬國的主.
那些國在哪裡.
是否在我的禱告裡.
因為祂自己管.
所以在我們代求就會展示出來.
我們究竟聽明什麼.
從看這個代求就知道.
究竟上帝的管治.
祂的國度.
在基督的復活裡面.
展明一個什麼國度呢.
然後教會對所有天災人禍.
戰爭.
究竟他應該用什麼立場.
什麼態度.
將他們抱過來.
好像耶穌將我們抱過來一樣.
擁抱過來.

$^{561}$來重回祂的平安.
重歸祂的安樂.
重歸祂的應許.
祂願意祝福我們.
但是我們的翼.
是否好像鷹一樣展開.
來撫慰所有需要的人呢.
所以我們不單在禮儀裡面.
可以在代求裡面表達.
在差遣裡面領命.
你派我去吧.
他不用再向我們問.
我可以派誰去呢.
我們就要認明.
所以到生活的場面.
崇拜和生活是沒有分割的.
所有大小先知都指責.
以色列最大的問題.
是有祭祀.
有節期.
什麼都有紀念.
有煮餐.
但是回去就做自己的事.
就是信仰和生活是割裂的.
但是崇拜就不是這樣的.
當我們的代求.
當我們領命很清晰的時候.
我們就知道我們有使命.
在這裡聽了上帝講話.
我就被猜出去.
我就是實踐在我的生活裡面.
我才在神面前是一個義人.
是一個完全人.
不單止從耶穌領受了.
這個恩順清義的恩.
而且我也要活像一個.
上帝國度的子民.
因此崇拜的差遣祝福.
為什麼要祝福呢.
祝福不是真的叫你.

$^{601}$祝福到自己的袋裡.
這麼簡單.
而是因為你有使命.
有差遣.
所以你需要三一的同在.
聖父的慈愛.
基督的恩惠.
聖靈的團契能力.
和你同在.
來完成祂給你的使命.
所以你是一個被增強的人.
被托住的人.
天祖的人.
去將生活與你的信仰.
落實在你的生活裡面.
所以你才需要祝福.
而且祝福基本上.
就是從早.
從前面.
舊約.
神要祝福亞伯拉罕.
叫祂祝福萬民.
如果我們是亞伯拉罕的子孫.
我們更要祝福.
所有我們見到的人.
能夠到之處.
能夠上到之處.
今天的新聞.
不只是說香港.
就告訴我們.
有很多的資訊.
很多的需要.
很多的災難.
很多的創惡.
為的是什麼?.
叫我們成為祝福萬處的人.
使得神的治權.
藉著我們的禱告代求.
臨到那些有需要的地方.
所以我們的使命.

$^{641}$不只是聽神的話語.
來祝福自己的靈魂.
而是叫我們.
不只是因為在這個平安裡面.
很高興生活.
平平凡凡.
無穿無爛.
不是.
叫我們成為一個.
這個世界需要的人.
好像希伯來書說.
這些信心的人.
是這個世界不配有的.
但是亦都是.
亞伯拉罕子孫的標記.
亞伯蘭是因為他的信.
被稱譽的.
因此我們必須在這個地上.
亦都因為信.
因著聽到的道.
產生信心.
以致我們唱.
所以我唱就是一個宣告.
所以.
中禮詩應該是怎樣的呢.
是動詞的詩.
我要去做什麼.
或者我要宣告上帝是我的主.
是我的王.
是我的神.
所以我不懼怕.
所以它可能有很多變數的.
有難處的.
但是耶穌吩咐他的門徒.
你不用怕.
因為我已經勝過世界.
他為我們勝過試探.
他為我們打低了.
死亡的毒鉤.
墳墓的門為我們闖開了.

$^{681}$所以我們不怕.
因此我們用口唱.
用口講.
用見證.
或者我們伸出關懷的手.
安慰的話語.
都是代表了我們的主.
我們的神.
所以千萬不要.
嘴唇尊敬我.
心卻遠離我.
耶穌指責這些.
早已指責的以色列人.
在他的時代.
一樣是這樣.
給我們看一些神蹟.
我就相信你了.
除了若拿的神蹟還要看什麼.
是不是.
今天我們已經知道.
主餐其實不只是關於耶穌的死.
而是關於祂再來.
應許,等著我們坐直.
我們走完這條路.
用信心,用盼望.
所以崇拜.
絕不與生活隔離.
差遣就是有力地去.
如果你很沮喪地進來.
你應該.
昂首挺胸地走出去.
一起來領受.
主榮耀的恩典.
生命的恩典.
所以聖賢.
應該有.
五經.
有詩篇幫助我們.
來回應律法.
當時在以色列的崇拜裡.

$^{721}$詩篇的角色.
就是回應律法書.
所以你才會有.
詩篇十九篇.
無言無語,有無聲音可聽.
是他量大.
通遍天下.
傳到地極.
教會更有這樣的責任.
傳到地極.
人人都聽到,會捲捲聖經.
都是重要的.
因為他們有任務指向.
榮耀的基督.
然後你會看見.
詩篇也使我們與.
以色列的.
聖文連在一起.
我們再不是說.
以色列.
不是用指責的口吻.
西人是會用指責的口吻.
但你知道根在那裡.
上帝一定會.
拯救他們的愚民.
所以.
我們要次次的禱告.
更加看見.
上帝在這個時代.
將人人的目光.
轉向以色列.
以以色列為中心.
不然誰會認識以色列.
那些國家不是讀聖經的.
你看見上帝在做什麼.
如果你知道結局是這樣.
他已經預言了是這樣.
你應許是這樣的時候.
你就已經看得明白.
現在的局是怎麼擺的.

$^{761}$重要的是.
你要去宣告.
好像詩篇29篇說.
上帝作王.
所以詩篇66篇.
今天就讓我們.
來有這樣的宣告.
我為什麼要尊崇他.
他考驗我們.
他熬練我們.
但我們的責任是什麼.
要稱頌我們的神.
使人得聽讚美他的聲音.
這是教會的責任.
不可以推給哪個機構.
推給政府.
一定是蒙恩.
相信的人才會有這樣的宣告.
這是責無旁貸.
不可以給誰.
只有你.
今天可能你會覺得.
像詩班.
喜歡唱歌就去.
不是,到你喜不喜歡.
是你的責任.
你們要讚美耶和華.
不是,你好不好讚美.
所以你要用信心.
來回應你的上帝.
除非你不認識你的上帝.
你越是認識.
越是知道是一個蒙恩的人.
你就用你的力.
用你的義.
來讚美他.
每個人都有責任.
這樣就是你的敬拜.
你嘴唇所發的.
你的嘴唇的見證.

$^{801}$要結果子.
不可以沒有這樣的果子.
然後對神的.
恩惠和敏銳.
是應該.
隨著我們的信心.
經歷上帝.
被熬煉,被考驗.
更加在末世的時候.
堅定.
堅定地祈求.
堅定地讚美.
這就是教會最重要的任務.
他在敬拜裡面的責任.
沒有別的敬拜團體.
他在讚美裡面.
他在禱告裡面.
如何將這個殿成為一個禱告的殿.
是耶穌最肉緊的.
也是要清理的.
那些不是的.
那些雜質的.
憑著其他動機的.
將它清除.
為的是什麼.
純正的.
祝福萬國的.
為萬國祈禱的.
容讓萬國來祈禱的.
萬國的殿.
萬國禱告的殿.
所以我們.
沒有一樣.
是可以沒有信心的.
讚美需要信心.
禱告需要信心.
但這個信心不是你的感覺.
有很多人以為.
我感覺我軟弱.
你為我祈禱啦.

$^{841}$你讚美啦.
這個信心是恩賜.
給了我們.
就等於我們有.
受洗.
歸入基督的人.
這個信已經踏出了一步.
上帝已經給了我們.
標記.
這個印記是使我們.
一直走到底的.
信就是所盼望的事.
有把握.
對未見的事.
有確據.
這個信.
這個歌.
這個禱告.
你們看見後面.
耶穌已經對我們顯明.
事實.
讓我們知道.
上帝會因為我們這個信.
報償我們.
這個世界.
滿了沒有信心的人.
是不是.
又或者信了別人的樣子.
很多偶像.
或者信自己.
但是這個世界.
很缺的是什麼呢.
單一信神的人.
所以.
你要為你的教會來祈禱.
興起我們對神的信心.
但是.
從哪裡來興起呢.
對神的話的態度.
愛我的就遵守我的道.

$^{881}$以其為寶貴.
所以.
你一定要好好學習.
詩篇一百一十九篇.
當它是珍寶.
如珠如寶.
當它是糧食.
當它是生命的指導.
當它是生命最重要的中心.
愛慕神的話語.
從你的讀經開始.
我希望.
有人會認證.
好,我來讀經.
語文老師可以.
自學最好.
因為你中文.
起碼知道不會說.
耶和華是我的墓者.
什麼是重點.
什麼是焦點.
什麼是抑揚頓挫.
將它應該的訊息出來.
而不是像打字機那樣讀出來.
我們不要因為有文字.
就很均勻地讀出來.
其實沒有讀進去.
叫過門不入.
所以我們一定要.
做好一點.
這是對世界的責任.
對我們牧羊.
整個群體的責任.
我們好好地.
將整個聖經的全貌.
藉著這個循環.
一路將這個.
深化.
不要片面.
所以我們要.

$^{921}$相信這一位.
回應我們信心的主.
所以起碼我們今天.
要立志我用信心.
宣告我的主.
在歌中我要用.
五成歌唱.
又用零歌唱.
然後我禱告呢.
我又用五成禱告.
用零禱告.
成為這個.
區域的光.
成為香港的光.
成為國家的光.
成為.
祝福萬民的渠道.
我們一起祈禱.
親愛的天父.
多謝你呼召我們每一個人.
成為敬拜你的人.
多謝你.
今天已經將不同的敬拜的.
事奉崗位派給我們.
當中每一個人.
但其實我們每一個人.
都是一個敬拜你的人.
你在等待我們成為真正.
求你今天.
對我們的靈說話.
開通我們的耳朵.
讓我們知道我們的投入.
我們的參與.
我們的信心都要.
付出力量代價.
教我們.
好好地來預備自己.
我們的生活.
要成為一個祭物.
我們在這裡的信心的.

$^{961}$長進.
從你道裡面的長進.
下雨的長進.
是在生命,生活裡面.
實踐出來的.
求主幫助我們.
不要只做表面的功夫.
教我們怎樣從裡面.
成為一個真正.
將自己呈現給你的人.
因為耶穌.
已經完全的犧牲.
為我們的罪而死.
也為我們的義來復活.
為叫我們清義.
求主叫我們.
挺胸昂首的.
在祝福同在的裡面.
歡喜的讚美你.
有動力的.
來禱告.
因為我們的禱告.
我們彭遜的禱告.
就在高舉你.
在大水之上作王的主.
就在洪濤中.
你仍然.
坐著為王.
求你.
在我們生命裡面掌權.
在我們這個家裡面掌權.
叫我們人人都在.
你的道中.
領受你的良.
領受你的生命之言.
以致我們的生命有力.
信心強健.
盼望充足.
以致我們彼此相愛.
更加是實踐.

$^{1001}$實在.
但願主.
聽到我們心中.
向你的呈獻.
獨到.
一路來引導我們.
在崇拜裡面的長進.
願主你見頌我們.
訴頌我們成為一個.
敬虔的人.
好像你昔日拯救挪亞.
和他一家人一樣.
大水中.
你要將我們保在.
基督耶穌裡面.
叫我們是穩道.
波圖.
但他看見你應許的實現.
就是那個天雄.
要祝福萬人都成為萬人的標記.
看見你的愛.
你的憐憫.
你的容忍.
最後是引我們進入你的美地.
願主親自來引導我們.
整個信心的前途.
走到最後.
向著標竿直跑.
從沒有放棄.
求主激勵我們的心.
聽我們禱告奉耶穌得勝的名求.
Amen.
\newpage



\section{創世記 13:8-9}
\label{sec:W5xhymK0FGw}
\textbf{2026.05.17 《有情義的人蒙福》:創世記13\_8-9(和修版) 郭鴻標牧師}
\newline
\newline
連結: \href{https://youtube.com/watch?v=W5xhymK0FGw}{\texttt{https://youtube.com/watch?v=W5xhymK0FGw}} ~~~~ 語音日期: 2026-05-20
\newline
\newline
\hyperref[sec:WnFpUR4Lsgw]{\small{< < < PREV SERMON < < <}}
~
\hyperref[sec:index_chronic]{\small{[返順時目]}}
~
\hyperref[sec:index_scriptual]{\small{[返順卷目]}}
~
\hyperref[sec:3yIpsNheSeo]{\small{> > > NEXT SERMON > > >}}
\newline
\newline
創世記 13:8-9
\newline
\begin{longtable}{cl}
\hline
\hline
章節 & 經文 (和合本修訂版)\\
\hline
13:8 & \begin{tabularx}{0.7\textwidth}{X} 亞伯蘭就對羅得說：「你我不可以相爭，你的牧人和我的牧人也不可以相爭，因為我們是一家人。 \end{tabularx} \\ \\ \relax
13:9 & \begin{tabularx}{0.7\textwidth}{X} 遍地不都在你眼前嗎？請你離開我吧！你向左，我就向右；你向右，我就向左。」 \end{tabularx} \\ \\
[1ex]
\hline
\hline
\end{longtable}
$^{1}$各位弟兄姊妹早晨.
我們一起感恩祈禱.
親愛的天父上帝.
我們今天來到你的面前去敬拜你.
願你的聖靈在我們的心裡面工作.
你去激勵我們.
也都打開我們屬靈的眼睛.
聽得明聖經的道理.
也都讓我們渴慕明白聖經的那個說話.
讓我們那個屬靈生命一路成長.
我們繼續在我們人生的路裡面.
抓緊我們的時間去做神要我們做的工作.
這樣祈禱奉耶穌基督的名.
阿們.
很開心又和大家見面了.
我自己靈修的時候.
我都看了《創世記》的經文.
我今天講的人物.
可能你們在主日學都聽過的.
就是亞伯拉罕.
他未改名之前就叫亞伯蘭.
大家可能都有一些印象.
不過我自己為什麼會特別要講這一段的經文.
就因為如果我們看《創世記》第十五章.
就開始亞伯蘭有另一些的經歷.
和十七章都是上帝和他立約的.
但是如果我們看回第十二章開始的時候.
就是上帝對亞伯蘭說.
你要離開戊爾去哈蘭.
你可以說對聖經的次序.
你想一下十二章十三章十四章.
其實中間發生了什麼事呢.
然後去到十五章上帝和他立約改名.
其實我自己一路再看的時候.
我發覺有很多東西.
我是得到一些啟發的.
所以我就很希望和大家分享一下.
我讀經的體會.
我就是跟經文的結構和大家講的.
有三個部分.

$^{41}$我今天就沒有做powerpoint.
讓大家可以專心一點.
第一個部分就是講信靠上帝.
其實都是從經文去講.
看亞伯蘭他怎樣去信靠上帝.
第二是知識紛爭.
第三就是他怎樣得到更大的福氣.
這些全部都在經文裡面出來的.
講回亞伯蘭或者叫亞伯拉罕.
我們經常都叫他做什麼.
信心之父.
你猜信心之父有沒有污點.
有的,姐妹們馬上說.
哪件事呢.
去埃及的時候做什麼.
叫他太太做什麼.
做妹妹.
是不是大家聽完之後已經是摩拳擦掌呢.
你太太叫他做妹妹.
然後就被法老王帶去皇宮.
你會怎麼想.
我們中國人有句話.
夫妻如衣服.
大能臨囚.
各自飛.
亞伯蘭你都很會做事.
有一次我講到.
有人說其實亞伯蘭沒有說謊.
為什麼沒有說謊.
因為他都是堂妹.
這樣都可以.
是妹不是妹都不要緊.
但是你那個時刻是什麼.
很明顯你是想保你自己.
是不是有污點.
有污點.
有污點為什麼還可以叫他做信心之父.
所以我今天是嘗試看回聖經的經文.
其實他還有另外一些地方是值得學習的.
不過那件事真的做得差.

$^{81}$這個真的要公道一點說話.
我們看回第十二章第一節的時候.
就是上帝呼喚亞伯蘭.
你可以有手機可以看回經文.
聖經你打開我接下來說的.
就叫他離鄉別井.
還要跟他說.
十二章的時候跟他說.
你將來會成為一個龐大的民族.
你的名聲會響遍天下.
你想像一下.
有朝一日上帝跟你這樣說.
將來你就會怎樣怎樣.
你可能想將來.
我現在好像都是差不多.
耶和華就應許.
亞伯蘭萬國是因你得福.
其實很開心聽到之後.
不過這個應許.
是好像一張支票.
不過還沒有兌現.
他說那些祝福你的.
你都得到祝福.
就坐你的那些都會給我就坐.
你聽起來真的很舒服.
神對他這麼好.
然後怎樣呢.
聖經就說亞伯蘭.
按著上帝的吩咐去出發.
不簡單的.
他要帶他的妻子撒拉.
或者撒萊.
還有他的兒子.
他的兒子羅特.
就是說一家人之外.
還有一族人這麼大.
還有一切的積蓄.
還有那些生畜.
生口.
去到這個迦南地.

$^{121}$而當時在士劍這個地方.
你聽到地名.
我們沒有地圖給你看.
但是他去到那裡.
就有迦南人居住.
經文有些很細微的地方.
在那個時候.
耶和華向他顯現.
很有趣.
我們沒有經歷過耶和華向我們顯現.
你明白嗎.
我們沒有這個體會.
然後就告訴他.
要將那塊地賜給他們.
這個這樣的感受.
是很特別的.
就好像剛才姐妹說.
走到洪恩家的門口.
就覺得這裡是他屬靈的家.
亞伯蘭有一點.
就去到士劍的地方.
耶和華向他顯現.
這個地方一定是給你的.
這些屬靈的經歷.
不是每個人都有.
但是亞伯蘭有這個情況.
但是你想一想.
亞伯蘭當時是怎麼樣.
他還未可以定居下來.
就好像那些人搬來搬去.
都未定居.
但是他做了一件事.
是很特別很特別.
就是在耶和華向他顯現的地方.
捉一座壇.
來做什麼.
去敬拜他.
你想一想.
你遇到上帝的時候.
你會不會當那個地方.

$^{161}$是聖地.
或者神聖的地方.
你有沒有這麼強烈.
我們未必會捉座壇.
但是你會不會為上帝.
在向你顯現.
向你說話的地方.
開一個查經班.
開一個祈禱會.
你明不明白我在說什麼.
其實我們未必會這樣做.
就算上帝給我有這個福份.
我都不一定想到我要做一些事.
亞伯蘭就有做這件事.
他然後再去往南邊的地方去.
就是說他對上帝和他說話.
向他顯現這件事.
他是真而眾知的.
他是做一些事的.
在那個時候.
你想想亞伯蘭.
他會怎麼看上帝的應許.
上帝你說的那些事很大.
不過和我現在的情況.
是很不合比例的.
我現在真的又沒有什麼明星.
我又沒有兒子.
怎麼搞呢.
他在迦南地遇到饑荒.
然後他就走了去埃及.
剛才我一開始就說了這件事.
他在埃及為了求生存.
假意說他太太是他妹妹.
在這個情況下.
很明顯是自保的.
怎麼萬國可以因這個人得福呢.
你明不明白我說什麼.
這個都很難想得明白.
他的明星怎麼會怨大呢.
這些這樣的應許.

$^{201}$真是很遙遠的.
結果法老王就將撒萊帶入皇宮.
結果上帝就降災給法老王.
然後他就查一下.
為什麼會搞成這樣呢.
他發覺原來亞伯蘭說謊.
那個不是他的妹妹.
是他的太太.
然後法老王說你們走吧.
不要搞到我們了.
亞伯蘭離開埃及.
就回去南地.
就再去伯特利和艾城中間的地方.
經文就說是他以前支搭帳棚的地方.
不知道有沒有人背背包去旅行那些.
在那裡起個帳幕.
然後他在那裡做什麼呢.
他再次求告耶和的名.
這些經文的記載.
這一小句是很短的.
還有在這些經文裡面穿插.
你不留意都會走眼的.
原來他是這樣的.
剛才要公道一點說的話.
他在埃及那裡不認他老婆.
但是他出來之後.
他去到不同的地方.
伯特利和艾城之間.
他又在以前支搭帳棚的地方.
求告耶和的名.
我們要這樣看一個人.
立體一點看.
亞伯蘭是這樣的.
有些事情是不好的.
不過耶和說向他顯現的地方.
他又起了壇去敬拜他.
第二.
他在那裡支搭帳棚的地方.
又去求告耶和的名.
看來他都算挺虔誠的.

$^{241}$你不要看他污點那些東西.
其實他都有虔誠的一面.
他在這些轉變裡面.
他仍然相信上帝的大靈.
他在飄泊著.
他沒有發出上帝對上帝的埋怨.
你看《創世記》是沒有這些的.
他也沒有懷疑上帝的應許.
說得那麼大.
都沒有經文沒有記載.
亞伯蘭說.
你在說的東西我都不信.
沒有這些.
他向神禱告.
他祈求他前面要走的路.
其實你想一下.
他這樣飄泊的時候.
他的體會是很深刻的.
他會更加體會到.
上帝是不是真的.
所以亞伯蘭後來被稱為信心之父.
當然看經文來說.
主要就是說他願意馴服上帝.
獻上他的兒子以撒.
這個是大家都知道的.
但是當你立體一點看經文的時候.
你會發覺原來他是從離開吳爾的時候.
他是正在經歷顛沛流離.
我們這一代人就沒有什麼經歷顛沛流離.
我們上一代的人.
不知道大家小時候會不會.
經常都說你們吃飯.
叫我們不要有飯剩在碗裡.
接著對我的兒子說.
你這樣就是草個豆皮婆.
不知道大家還記不記得.
那些童年回憶.
接著不是說這些的.
接著就說走難的時候是不能吃的.
吃樹皮吃老鼠.

$^{281}$這樣說完之後.
我們已經形成了一種習慣.
凡是碗裡面的飯都吃光.
不知道大家有沒有這樣的習慣.
我現在形容這些是強迫症.
到現在老了就發覺.
不知道大家會不會有一個感受.
你的家人跟你說.
你吃了肚子你的胃病更浪費.
不要要了.
會不會是這樣.
現在是我們進化了.
現在不是我們不想吃.
而是吃不到.
不是說要浪費東西.
而是根本沒有這個本事.
你想想顛沛流離的人.
他當然很珍惜他僅有的東西.
他明白到神是很真實.
和去磨練他的信心.
我們這樣看.
他沒有放棄前進.
你看經文他仍然前進.
縱使他在埃及真的軟弱.
這是沒有假的.
聖經不會隱瞞這些東西.
他說謊自保.
他對上帝的吩咐.
但另一方面他仍然深信不疑.
這是很特別的.
他馴服進行.
所以你看到一個屬靈的人.
你不要將他偶像化.
你可以看到他有弱點.
但有另一些東西你可以學他.
亞伯蘭對上帝的信靠.
給我們有很寶貴的提示.
其實今天我們又怎樣可以學習呢.
當然我讀到這個地步的時候.
我都開始想很多問題.

$^{321}$一個對信仰認真的基督徒.
應該在生活方面.
有很多自己的宗旨.
譬如說年輕的時候.
擇偶,擇業.
我們都會去問上帝的旨意.
這是一定的.
上帝的應許是什麼呢?.
上帝的帶領是什麼呢?.
一個找到上帝的應許的人.
他是會活得有方向的.
大家現在是不是會覺得.
方向是很重要呢?.
我自己就很體會.
就是說退休人士.
那個方向是很重要.
因為我真的退休了.
我今年就要搬出我的辦公室了.
其實退休人士是有很多潛在危機的.
有什麼危機呢?.
就是我們現在還有卡片可以派.
但是八月之後我就沒有片可以派了.
我認識一些人.
他就要令自己覺得我有東西可以派.
他就自己印自己的地址和電郵.
你派一張來,我派一張來.
我就不打算做這些事了.
其實是很不容易的.
為什麼呢?.
就是你習慣了上班.
突然間不上班其實是很艱難的.
所以究竟我的方向是會怎樣呢?.
尤其是現在這個時代.
就是醫學比較好.
以前的人說人生七十古來稀.
其實我和一些親戚聚在一起.
每個人都是.
你就是次次外來坐在一桌.
每個人一千歲.
你說我六十幾歲.

$^{361}$你最年輕了.
每個人七十七十幾.
八十都你.
你明白嗎?.
現在很多人真的是.
九十幾一百歲都很精壯.
這樣就問題了.
現在退休之後.
那些日子怎麼搞呢?.
難道每天都在淋花嗎?.
不行的.
所以方向你明白嗎?.
方向是很重要的.
所以在這裡我就想一下.
究竟我們有上帝應許的人.
怎樣活得有方向有意義.
我覺得我不是浪費時間.
很多時候我們都不太知道.
自己的長處或者恩賜.
但是我們有時候會感覺到.
退休後其實我可以做我以前想做.
但沒有時間做的事情.
這個是很自然的.
所以我都會發現有些弟兄姊妹.
他們是轉型的.
我都看到的.
希望我們自己都是這樣.
阿巴蘭就是這樣.
她的人生是.
讓我們看到很多事情是幫助我們.
越過了我們慣用的想法.
或者我們看自己只有這麼多.
我們不敢去突破.
我們要憑著上帝給我們的信心.
去走前.
做一些事情是有貢獻的.
或者是可以去.
做到上帝要我們做的事情.
我們可以看到有些人.
他曾經有決心做一些事情.

$^{401}$但是因為他不夠忍耐.
所以他試一下就放棄了.
有些人是這樣的.
然後又試另一個目標.
但是很可惜.
試來試去就是兜來兜去.
但是我們怎樣可以.
能夠在這些例子裡面學習.
我能夠立定了目標.
我可以不拖泥帶水.
我可以破釜沉舟.
我可以全力以赴.
我可以堅忍下去.
等到有出路.
這個其實是很不容易學的功課.
你看看舊約聖經.
很多人物都是在艱難中.
磨練了這種性格出來.
那你可以說什麼叫成功.
你想一下什麼叫成功.
最近我和我太太談了一個問題.
有些人就厲害了.
贏在起飽善.
你猜我太太回答我一句什麼.
最重要是贏在終點.
這給我很大安慰.
就是說贏在起飽善.
很多人都是贏在起飽善.
是的.
那個結果不是他一定贏的.
但是譬如我們這些都是.
不一定贏在起飽善的人.
但是如果我好像有聖經人物.
那種的堅毅.
對神的信靠.
對神的應許那種的依賴.
那你可以想像你會怎樣呢.
你是可以好像保羅那樣說.
我當跑的路.
我已經跑過了.

$^{441}$你可以是贏在終點.
或者我也可以贏在終點.
所以當我們看聖經人物的時候.
我們發覺原來有很多的提醒.
是可以讓我們知道.
就是說原來他們的人生的路.
是有很多啟迪的.
他們不一定是偉人.
他也不是十全十美.
但是他有一些屬靈的功課.
是值得我們今天去學習.
最低限度對我來說.
當我看亞伯蘭的時候.
給我有很多的安慰.
讓我知道我跟著走的路.
可以說我不一定全部知道是怎樣.
但是我都要憑信心走下去.
就算我走了幾十年.
但是還有很多的路.
我是不可以知道.
我又跟我太太談這個問題.
其實如果我知道了我的結果.
我會怎樣呢.
你可以想像.
如果上帝讓你知道.
你二十年後是怎樣的時候.
你是不是會現在開心一點呢.
那我就說不一定的.
那我就沒有了經歷.
看來上帝也是這樣的.
他不會在這一刻告訴你.
二十年後是怎樣的.
他會給你一些的感動.
要你走一步一步.
有時候你會覺得.
吊著一條癮.
究竟怎樣.
但是原來屬靈的經歷就是這樣的.
是你慢慢憑著信心去走.
我現在回看我二十歲的時候.

$^{481}$我就會發覺原來是這樣.
或者你到了八十歲的時候.
你看到六十歲的時候.
你可能又會更加清楚.
所以這個是第一部分.
信靠上帝.
亞伯蘭有這個心.
是繼續去走.
但是這個都不是經文裡面.
最多的觸動我的.
接下來就去到第十三章.
就是紫色紛爭的地方.
剛才讀了兩次經文.
正正就是說.
亞伯蘭和他的侄子羅德.
就是兩族人一起走的.
而他們有很多山口的.
而在伯特利和艾城中間的水源和草原.
就有限的.
你想都想到了.
地方那裡資源不夠.
兩家人都有他的牧羊人.
那些牧人.
那些牧人當然是.
當然是要找到水源草原先吃的.
對不對.
所以是不是有東西爭呢.
一定有東西爭.
大家這個例子很簡單.
如果你家裡兄弟姐妹多的.
譬如生十個以上一代的.
一隻雞只有兩隻雞腿.
通常有些人是吃不到的.
你明不明白.
十兄弟姐妹有八個吃不到.
你明不明白那個問題是什麼.
就是爭.
手快有手慢無.
當時都是一樣的.
亞伯蘭他和羅德的牧人去爭水源和草原.

$^{521}$爭水源草原.
我們可以出一些經文.
麻煩我們的弟兄姊妹.
你想一下.
中東的人或者是以色列人和華人都有點類似.
敬老尊賢.
對不對.
現在是誰和誰爭先.
是亞伯蘭的牧人和羅德的牧人去爭.
羅德是他的侄子.
遇到這樣的事.
你猜這件事應該怎麼擺平.
是不是應該羅德讓先.
對不對.
按我們的文化來想.
小的讓大.
那當然是羅德去讓.
但經文記載又不是這樣的.
我們可以一起去讀第十三章第八節.
亞伯蘭就對羅德說.
你我不可以相爭.
你的牧人和我的牧人也不可以相爭.
因為我們是一家人.
這是什麼道理.
亞伯蘭開口對羅德說.
不要爭了.
你的牧人和我的牧人都不要爭了.
我們是一家人.
這是什麼意思.
亞伯蘭退後了.
還有.
後面那個經文.
就是剛才第九節.
後面那句你再看看.
究竟亞伯蘭和羅德說完這番話之後.
他提議一個什麼處理方法.
我們一起讀.
遍地不數在你眼前嗎.
請你離開我吧.
你向左我就向右.

$^{561}$你向右我就向左.
這是什麼意思.
羅德.
你先選.
你去哪裡我就去旁邊.
那你會想想.
如果要處理這個問題.
會不會反過來.
大那個先選呢.
亞伯蘭說我們不要爭了.
我選東你就向西.
這樣也通不通呢.
通嘛.
但亞伯蘭採取一個什麼態度.
羅德.
你先選.
好的東西給他選.
不是說笑的.
給他選了就你沒得選了.
你想到這裡的時候.
你就會發覺.
他採取一個很特別的處理方法.
換成今天你如果有類似的情況.
你會怎樣.
亞伯蘭以他大家長的身份.
可以說羅德.
你好好管理一下你的牧人.
我的牧人帶點生畜.
吃完東西你才好來.
這樣說也不過分.
如果在中東人的文化來說.
是不是.
但亞伯蘭又不是.
不要爭.
你看這麼大地方你先選了.
然後就說到.
這個問題就是.
大家記不記得.
上帝對亞伯蘭說.
他將來成為一個大國.

$^{601}$你這樣讓法.
你怎樣做到大國.
你明不明.
真的很矛盾.
真的不行.
好了然後我們再看創世記.
我就讀給你聽.
創世記第十三章十到十一節.
這裡沒有打出來.
然後羅德就是聰明人.
經文說.
羅德舉目看見約旦河的全平原.
直到所以.
都是滋潤的.
在那裡.
就是在耶和華.
未滅所多瑪摩拉之前.
那裡地方好像耶和華的原址.
又好像埃及地.
於是羅德選了約旦河的全平原.
然後往東去遷移.
他就分開了.
很簡單的經文就是說這些.
如果你細心看的時候.
第一那裡地方好像耶和華的原址.
好像伊甸園.
又好像埃及肥沃.
又提到兩個城市.
所多瑪摩拉.
向東面.
這兩個在創世記裡面說是罪惡之城.
但是如果我們將這個罪惡之城.
和前面.
未到罪惡之城之前.
那裡是肥沃的土地.
你去連在一起的時候.
你會發覺羅德.
他是選對了地方.
那裡地方是很美的地方.
不過就包含了一個危機.

$^{641}$這個危機就是你往東去遷移.
這個好地方.
你不知道什麼時候會越走越遠.
走到所多瑪摩拉.
你明白嗎.
經文就沒有做任何判斷.
但是就說到他先選了好地方.
然後亞伯蘭就沒有反口.
你明白嗎.
他沒有反口.
好奇.
那你選了就算了.
亞伯蘭很真誠的對待羅德.
用我們今天的說話.
是因為他有愛.
不要爭.
他對神有信心.
他知道給羅德你選了最好的.
都不要緊.
上帝都一樣對我有好的帶領.
你看不看到亞伯蘭的這種做法.
然後最有趣的地方就是.
凡是亞伯蘭做了一些事情之後.
上帝又跟他說話.
在《創世記》第十三章十四到十八節.
我講一個關係次序.
上帝對他說.
從你所在的地方.
你抬頭望一下.
凡你所看見的一切地方.
我都要賜給你和你的後裔.
直到永遠.
我也都要使你的後裔.
好像地上的塵沙那麼多.
人如果能夠數到地上的塵沙.
都才可以數到你的後裔.
你起身.
你縱橫走遍這個地方.
因為我必定將這個地賜給你.
你發覺原來有個現象很有趣的.

$^{681}$就是上帝對亞伯蘭說的話.
是每次比每次具體的.
初初只是說你離開.
即吾爾去哈蘭.
比較大方向一點的.
接著他又做一些事情.
上帝又多說一些事情給他聽.
我發覺這個情況是一會兒都會出現的.
我一會兒再說給你聽.
好了 上帝在不同的階段對亞伯蘭.
說他應許的發展.
落實是越來越具體的.
這個給我很大的感動.
就是原來我們跟上帝的那個應許.
你是會越來越清晰的.
而倒過來你看到.
經文記載亞伯蘭.
去到哪裡見到上帝.
聽到上帝說話.
他又會捉個壇.
又去求告耶和華的名.
但經文是沒有記載過.
羅德做同樣的事情的.
你細心看一下.
原來是這樣的事情.
原來亞伯蘭他那個內在的屬靈生命.
他是有這些好的地方的.
他以神為他人生的中心.
你看到羅德其實他很懂得選擇.
從我們世界的角度來說.
他選擇了一些最好的東西.
羅德他看重的是水源 草源.
但是他不自覺的.
自己漸漸拉近了自己.
跟索多瑪 跟俄摩拉的距離.
你們有沒有看到.
好的地方.
不過那個地方是引領人不自覺.
走去墮落的.
相反亞伯蘭他看重神的帶領.

$^{721}$去到哪裡他都要為耶和華捉壇.
求問我們上帝的心意.
你想一下.
這兩種的生活方式.
是兩個不同的結果.
而兩個人的分別.
就在《創世記》14章1到16節表現出來.
有什麼事情發生了呢.
就是羅德被侵略索多瑪.
和俄摩拉的王捉了去.
即是那個好地方再過那兩個大城.
但是有其他的王就已經想侵佔他們.
結果就一併捉了羅德.
好了 故事到這裡.
你是亞伯蘭.
你聽到你的侄子選擇了好地方.
結果又被人捉了去.
你會怎麼想呢.
好多不同的可能性.
第一個可能性.
都預了你.
心驕異傲.
真是急於上位.
那你就吃個教訓.
可以理解.
好了.
《聖經》怎樣記載呢.
《聖經》記載亞伯蘭知道之後.
很特別的記載是什麼.
他帶同精煉的壯丁.
還要有具體數字.
318個人去追戰.
我就問這個問題.
為什麼《聖經》有些地方不記這麼多數字.
這些記載得這麼清楚.
亞伯蘭二話不說.
知道羅德有事.
二話不說.
不要理這麼多.
叫齊人出去救人.

$^{761}$《聖經》又沒有說他用什麼戰術.
最後打敗了那些侵略的王.
救回羅德和他的家人.
這個只是故事的開始.
你想想其實在這樣的情況下.
亞伯蘭和羅德分開發展.
他心底裡其實沒有什麼是假的.
大家都有一條刺.
但是亞伯蘭沒有和羅德反目成仇.
你有沒有看到這件事.
反過來他有骨肉的親情.
而他希望羅德.
你有更大的發展我就安樂了.
不過當羅德遇難的時候.
他真的顯出他的那種情義.
他就竭盡所能去救他.
其實你想到這個位置.
其實都可以有掙扎.
其實可以去到一個地步.
我救你不就是害你.
我這次救你.
下次你又來.
我什麼時候可以全力救你.
你明不明白掙扎就是.
我為什麼要這樣出手.
但是你看到亞伯蘭.
經文沒有記載這種矛盾.
他二話不說就去做.
如果亞伯蘭有這樣的實力.
當日我根本就不需要理會羅德.
我說什麼你都要聽.
但是亞伯蘭沒有用這一套.
而是他用一個非常之.
引用成人之美的方法.
但是倒過來當羅德有難的時候.
他就不堆相了.
你覺得亞伯蘭這個人是不是有情有義.
所以我就改了題目叫有情義的人.
聖經沒有說亞伯蘭為什麼會被神揀選.
就是主要說他有信心.

$^{801}$對神有信靠.
不過他對神有信心這個優點之外.
其實他人生都有些特點.
就是他樂於知識紛爭.
和對神的信靠是很大.
但這個都不是故事的結束.
我再看下去就發覺又學到很多東西.
第三就是他獲得更大的福氣.
當亞伯蘭打敗了入侵索多瑪.
和摩拉的那些基大老馬的王.
和其他王之後.
他回到索多瑪.
索多瑪王和撒冷王.
叫默基西德.
拿些餅和酒出來祝福.
這一段很有趣.
默基西德當時是另一個祭司的至高神.
在迦南地的至高神.
他們是不認識耶和華.
我們亞伯拉罕的神.
默基西德特別為亞伯蘭祝福.
就說了一些話很有趣.
願天地的主至高的神.
不是說耶和華.
他是外邦祭司.
他的神叫至高的神.
願至高的神是天地的主至高的神.
賜福給你亞伯蘭.
至高的神將敵人交在你手裡.
真是應當稱頌.
他很正常.
他用他的語言.
然後經文記載什麼.
亞伯蘭得到默基西德祝福之後.
他就拿十分一的財物出來.
十分一奉獻就是這個背景出來.
就去給祭司.
然後所多瑪王就出來截住他.
他說你只需要將人交給他.
他從基大魯瑪王和所多嗎哪回來的財物.

$^{841}$亞伯蘭你可以自己帶走.
你明不明白這段經文的含義是什麼.
即是亞伯蘭按規矩.
你祝福我給十分一.
所多瑪王不用這麼客氣.
凡是你在基大魯瑪王那裡拿到的財物.
所多瑪這裡你拿到的.
我不跟你計較全部是你的.
你是亞伯蘭你會怎樣.
你想想你是亞伯蘭會怎樣.
精彩這裡我看到.
是《創世記》十四章二十二到二十四節.
我讀給大家聽.
亞伯蘭就很豪氣地回答.
我已經向天地的主至高的神耶和華.
這裡很有趣.
天地的主就是對應迦南的.
麥基西德說的話.
而這個他馬上在下面補了個名字.
他叫耶和華.
我對他起了誓.
凡是你們的東西就是一根線.
一根鞋帶.
我都不拿.
免得你們說我洗到亞伯蘭腹中.
豪不豪氣.
很豪氣.
你是不是這麼豪氣.
然後他還要說.
只有僕人所吃的.
並和我同行的那幾個.
阿賴 以實國 萬利.
即是這些他的打手.
所應該得到的.
就可以任他們去拿.
創世記十四章二十二到二十四節.
亞伯蘭是一個怎樣的人.
即是你給他有財物.
他在這一刻.
他很清楚說.

$^{881}$我來救人不是要拿錢.
不過我那些夥計.
他拿命搏的.
他那些要拿的.
我給他拿.
你覺得這個人是不是真的很特別呢.
我覺得他很特別.
白基使得那個為亞伯蘭祝福.
是他的職份 祭司的職份.
而亞伯蘭是要將那個禮物給他是正常的.
而他很清楚他打得贏的原因不是自己.
是耶和華的力量.
三百一十八個人.
其實可能是告訴我們很少數都不一定.
當所多瑪王說那些財物你可以拿走了.
他一點都沒有動心.
他說我不會改變我原來的立場.
因為我早已經向耶和華起誓.
我來救羅德不是為了錢.
也不是為了拿你所多瑪地方的財物.
更加不是要拿你所多瑪地方的東西.
又不是要趕走你的所多瑪王.
我完全沒有這個想法.
沒有半點的野心.
去到這個位置.
我才覺得亞伯蘭有這種氣質很特別.
他向那些迦南不認識耶和華的王.
去見證他們所拜的至高神.
就是叫他離開吾爾的那個耶和華.
你明白嗎.
他在做一個見證.
就是說我是這樣見證給你們知道.
在有財物的時候我不起貪婪的心.
我只是盡上我應有的本份.
去到這個位置.
我覺得這個是不是真的英雄.
經過我這樣演繹.
他們都覺得是英雄.
我再補充.
他在埃及真的輸過一次.

$^{921}$在其他方面我真的很欣賞他.
但最有趣的地方就是說.
這件事只有去到創世記十五章.
耶和華跟他說話說的東西又不同了.
我就看見這個很特別.
耶和華神在十五章裡面對他說.
亞伯蘭你不要懼怕.
我是你的盾牌.
必定大大的賞賜你.
跟著亞伯蘭就回應.
很少他回應上帝的.
前面十三十四章說的時候.
他只是跟上帝說話.
這次亞伯蘭就回嘴了.
他說他沒有兒子.
怎麼可以有後裔承繼他的家業呢.
這樣無論是神應賜給他什麼後裔土地財富.
都只能夠透過收養別人的兒子來實現.
他開始跟上帝說話了.
你這樣說得這麼大.
我根本現在什麼都沒有.
那怎麼辦.
然後耶和華就清楚跟他說.
在十五章的第六節.
他就帶他去到外面對他說.
數算眾星你數不數得過你.
以前都說過這類事情.
然後又對他說.
你的後裔都好像這樣.
再說一次.
當亞伯蘭聽耶和華說完這番話之後.
他又不知道為什麼又相信了.
他是一個這樣的人.
第十五章第六節記載說.
亞伯蘭信耶和華.
耶和華就以此為他的義.
大家明白我的意思嗎.
在創世記十二章十三章十四章裡面.
記載亞伯蘭的人生裡面.
他是這樣信耶和華說的話.

$^{961}$那耶和華就覺得你是一個.
在他眼中正義的人.
我看到這件事就是.
亞伯蘭得到上帝的祝福.
就因為他信上帝的應許.
信服上帝的引領.
和他帶領族人前進的時候.
他是堅信上帝的力量.
他是很深信上帝的力量.
他寧願用一個退讓的方法.
來去解決一些和羅德的紛爭.
當羅德被攻擊的時候.
被捉了去的時候.
他就即刻帶族人去救他.
幫他解決他的困難.
亞伯蘭他這種品格是很正直的.
令到人們祝福他.
墨基使得祝福他.
又得到所多馬王的欣賞.
但重要的地方就是.
他得著上帝再一次的應許.
大家明白嗎.
我看完這件事.
我就覺得很值得我們學習.
創世記十五章第六節記載說.
亞伯蘭相信耶和華的應許.
就算他是懷疑.
其實早在亞伯蘭被呼召離開吾爾的時候.
他一直都相信上帝.
當他踏出新的一步.
他就面對新一輪的考驗.
由吾爾到迦南.
到迦南去到埃及.
由埃及回到迦南.
再由伯特利和艾城之間轉向希伯倫.
他一直在遷移.
不知道大家有沒有感覺到.
你的人生一直在遷移.
我們上一代的人都是遷移的.
我們可能自己都有一點遷移.

$^{1001}$在這些日子裡面.
其實亞伯蘭根本不可能看到自己.
會成為一個龐大的帝國.
他根本不可能見到這件事會發生.
所以在我們的人生裡面.
我們就會面對一樣東西.
就是我相信上帝.
我見到的東西可能不是馬上實現的.
那我相信還是不相信.
他要避開迦南人,比利時人,埃及人.
和擔心法路的敵意.
他沒有辦法和周邊的強大的部族角力.
他也沒有打算用權力去駕馭他的親屬.
所以你看到.
原來這個亞伯蘭也挺有趣的.
原來這個人.
那你想想他.
對上帝的信.
就是他應許.
信還是不信呢.
怎麼信呢.
我就覺得他很多時候.
在一個肯定和不肯定之間.
他在矛盾,浮游,掙扎.
但是他想過之後.
他還是走的了.
當亞伯蘭繼續以信心向前的時候.
上帝將他應許的內容.
就講得更加清楚.
將來他會成為更大的國度.
所以在這件事裡面.
讓我看見一樣東西.
就是去到第十五章第四節.
上帝應許他有一個親生的兒子.
後來以撒就出生了.
所以你會發覺.
原來我們的信仰.
上帝的應許.
和我們眼前見到的情況.
不是馬上兌現的.

$^{1041}$你會發覺不是馬上兌現的.
在這個情況下.
我們要學習.
我們怎樣去學習.
去做一個基督徒.
我們基本上.
信仰的旅程就是這樣.
你想更加看到上帝的真實.
你就會發覺.
有些現實和你所信的東西.
上帝應許的東西.
不是馬上兌現的.
應許是需要時間去慢慢去實現的.
所以為什麼我們基督徒.
是很需要有信心.
你沒有了信心.
你基本上看不到.
我們看到的是目前的東西.
信心是那些未發生的東西.
是不是.
所以來到總結.
我們就看回《創世記》十二章到十五章.
是說亞伯蘭照著上帝的旨意前進.
他沒有想過為自己做什麼.
他只是按著神的帶領去做事.
他信服神的帶領.
他又樂於知識紛爭.
但是他和他的兒子不去爭奪.
他願意讓兒子羅德你先選擇.
你選擇完之後我選擇其他地方.
他心裡面沒有敵人.
你覺得他是不是.
他也樂於助人.
他得到別人欣賞他.
但是他又不會貪.
這個我覺得是很好的.
亞伯蘭出手救羅德.
不是為自己.
他是為家族整體的好處.
所以當我們看到一個這樣的人的時候.

$^{1081}$究竟我在今天可以怎樣去學習呢.
亞伯蘭這個故事.
他的品格是可以怎樣呢.
《創世記》十二到十五章.
亞伯蘭是主角.
羅德只是配角.
你想一下.
經文是沒有交代羅德有沒有自我反省.
他經歷了這些之後.
他怎樣想我們不知道.
亞伯蘭他幫羅德一把.
但是很可能羅德他也不明白.
也不奇怪.
所以當我們看到.
這些情況下的時候.
我們就會想原來我們也要學習一下.
有些時候很複雜的情況.
我們也要退讓一下.
好像亞伯蘭那樣退一退.
但是當你看到有人有艱難有困難的時候.
你應該要出手幫助.
你就不要遲了.
這個是我看亞伯蘭這個故事裡面.
我自己學到的東西.
我們很需要祈禱,安靜.
聽聖經裡面的教訓.
來幫我們有人生的智慧.
我們也需要這樣東西.
特別我們要常懷善意.
常懷善念.
不以惡報惡.
寧願有些時候都要退一退.
然後看上帝來做事.
退讓不等於是非不分.
更加不是縱容惡人.
而是我們給人一條出路.
這個我想我們真的要學習.
為什麼這樣做呢.
你想一下上帝都給我們一條出路.
上帝什麼時候都給我們一條出路.

$^{1121}$所以我們都盡量的能夠這樣去做.
以致我們體會到上帝仍然是工作.
我們都希望周圍的人有機會回到上帝那裡.
我們希望他們同樣的看到神沒有離棄他們.
我們都要學這個故事.
我就將我讀經的體會跟大家講一下.
我自己就發覺是得益的.
有些很長的故事.
有時候都很不容易看.
但我也覺得勝經是神的話語.
祂是給我們很多提醒.
很感恩我臨退休就學到這個功課.
真的很有意思.
讓我明白了很多道理.
所以我也鼓勵大家熱愛讀一下聖經.
和上一下教會的主一學.
不只是聽道.
聽道詞也只有幾十分鐘.
那些是有限的.
但希望大家能夠繼續蒙神的恩典.
我今天看到這裡.
又和我上次來的情況不同了.
又多了很多人.
我自己感覺到洪水橋是發展地區.
所以我也希望大家繼續去堅守這個陣地.
不容易的.
我也看到洪恩堂是不容易的.
但我們也看到.
我這樣講你們不會反對的.
神是沒有離棄大家的.
這是我自己看到的.
其實我相信任何一個教會都是屬於神的.
只要我們肯堅守陣地的時候.
神是會工作的.
你同意嗎?這是很重要的事情.
神是不會離開祂所愛的教會.
我們一起同心祈禱.
我們的天父上帝.
因為在聖經裡面很多的人物.
他們都在經歷著你.

$^{1161}$聖經的記載就是將他們體會上帝的過程告訴我們.
我們讀聖經的時候.
我們真的要細心的去讀.
慢一點下來去思考.
去問自己.
究竟我在那個情況下.
我是會怎樣呢?.
我做不做得到呢?.
究竟我可以怎樣學習呢?.
因為你的聖經的話語是我們生命的糧食.
是牧養我們的.
是你建立我們.
以致我們得到智慧.
得到我們一些提醒.
免致我們犯錯.
就算我們信了很久.
我們很可能有時會軟弱.
或者我們有機會衝動而犯錯.
但我們真的需要你提醒我們.
建立我們.
特別祝福紅印堂的牧者和信徒領袖和眾弟兄姊妹.
在這個地區可以成為明亮的燈台.
可以服侍更多的人.
因為剛才我們唱詩的時候都聽到.
有很多人都活在罪惡之中.
他們不認識神.
他們的人生是缺乏盼望.
缺乏了一種勇氣.
缺乏力量.
我們懇求主.
你親自來臨到每個人的生命裡面.
也讓他們歸入你的陽蘭.
可以經歷你在你的教導之下.
黎明成長.
恭敬阿彌陀佛祈禱.
奉耶穌基督的名.
阿們.
\newpage



\section{約翰福音 7:37-39}
\label{sec:3yIpsNheSeo}
\textbf{2026.05.24 《湧出活水的生命》:約翰福音7\_37-39(和修版) 陳麗蟬牧師}
\newline
\newline
連結: \href{https://youtube.com/watch?v=3yIpsNheSeo}{\texttt{https://youtube.com/watch?v=3yIpsNheSeo}} ~~~~ 語音日期: 2026-05-27
\newline
\newline
\hyperref[sec:W5xhymK0FGw]{\small{< < < PREV SERMON < < <}}
~
\hyperref[sec:index_chronic]{\small{[返順時目]}}
~
\hyperref[sec:index_scriptual]{\small{[返順卷目]}}
~
\hyperref[sec:kpqG_RsMI0U]{\small{> > > NEXT SERMON > > >}}
\newline
\newline
約翰福音 7:37-39
\newline
\begin{longtable}{cl}
\hline
\hline
章節 & 經文 (和合本修訂版)\\
\hline
7:37 & \begin{tabularx}{0.7\textwidth}{X} 節期的最後一天，就是最隆重的一天，耶穌站著，喊著說：「人若渴了，到我這裡來喝！ \end{tabularx} \\ \\ \relax
7:38 & \begin{tabularx}{0.7\textwidth}{X} 信我的人，就如經上所說：『從他腹中將流出活水的江河來。』」 \end{tabularx} \\ \\ \relax
7:39 & \begin{tabularx}{0.7\textwidth}{X} 耶穌這話是指信他的人要受聖靈說的；那時還沒有賜下聖靈，因為耶穌還沒有得到榮耀。 \end{tabularx} \\ \\
[1ex]
\hline
\hline
\end{longtable}
$^{1}$各位姐妹早晨.
我想請問今天是甚麼日子呀?.
好像今天也是.
但另外呢.
今天是呀.
那另外呢.
我當然不是問你這個日子啦.
今天是聖靈降臨日.
就是耶穌基督復活之後的五十日.
就是五旬節.
就是聖靈降臨日.
在舊約裡面.
就是如意節之後.
五十日之後.
就是五旬節.
所以在舊約裡面.
去指向新約的時候.
就是聖靈降臨日.
聖靈降臨日很重要.
因為在當日的時候.
聖靈臨到地上的時候.
講到三千人信主.
而這個教會.
亦都在當時開始.
一路一路發展下去.
所以一些禮儀教會就定了.
這個聖靈降臨日.
就是這個教會的生日.
所以可能一陣子完崇拜的時候.
我們要為教會唱一首生日歌.
這個聖靈降臨日.
有人說到聖靈充滿.
有一位姐妹.
她被人去按手祈禱之後.
她就說我被聖靈充滿了.
很開心.
她笑了七天.
她就說笑得我很辛苦.
她感恩那段時間沒有安息禮拜.
否則她安息禮拜仍然在笑就慘了.

$^{41}$有人說當她被聖靈充滿之後.
她又會哭了七天.
因為聖靈充滿的時候.
人就會為義為罪.
去極度難過.
所以要不斷哭.
又有些人在聖靈充滿之後.
又會講方言.
究竟當聖靈充滿我們的時候.
是不是真的要有一些特殊現象.
然後才叫做被聖靈充滿呢?.
耶穌基督在地上.
在祂世上最後的鑄棚節的時候.
祂有一個宣告.
就是剛才輝哥為我們讀出的那一段.
引若鶴鳥到我這裡來.
就如經上所說.
從他腹中將流出活水的江河.
這個活水的江河.
就是指到聖靈.
經文一開始的時候.
它就講到這個.
是不是這麼白色紙那個.
謝謝.
一開始講說節期的最後一天.
就是指到鑄棚節.
在這裡記載的鑄棚節.
就算到現在以色列人都是這樣的.
他們就不會睡在家裡.
他們在屋外就會搭棚.
搭了棚之後.
他們就會按照舊約裡面上帝的吩咐.
用橄欖樹枝,棕樹枝,藤樹枝來搭一個棚.
在那幾天裡面.
他們就會在這個棚裡面住.
特別搭這個棚的時候.
在頂部一定不可以搭得很密.
因為要讓光從上去而淋到.
不是因為沒有電燈,沒有蠟燭.
所以要這樣.

$^{81}$而是他們要.
當他們睡覺的時候.
看上去的時候月光又有光進來.
在白天的時候又有太陽光照射進去.
給他們一個提醒.
我們所相信的上帝.
他們就叫耶和華.
因為他們不信耶穌.
就是這個光之源.
所以就讓光一直透進棚子裡面.
他們為什麼會有鑄棚節呢?.
他們就是要紀念.
在出埃及之後.
一段很長的時間.
在這個曠野裡面漂流.
在當中他們就知道.
在這個曠野裡面.
他們沒有房子住.
所以他們到現在.
今天的以色列人.
仍然都有一個這樣的慣例.
他們就記得.
在當天沒有地方住.
到處走.
其實藉著這個節日.
都給我們一個提醒.
就是我們在地上.
只是一個寄居.
只是走來走去.
但是到有一天.
是天父上帝的日子的時候.
去到天家才是我們永恆的居所.
在這個鑄棚節.
除了有搭中樹枝.
橄欖樹枝的情況之外.
還有一個特別的儀式.
整個節期有七天.
他們每一天的祭司.
就用一個金色的瓶子.
很名貴的瓶子.

$^{121}$從聖殿去到西羅亞祠取水.
入滿了金瓶子之後.
這個瓶子就一直拿回到聖殿.
然後就倒進一個銀色的瓶子裡.
就當是一個獻祭.
因為當他們在曠野的時候.
他們沒有水喝.
摩西就用一支杖去擊打這個盆石.
或者他只需要指著盆石.
去吩咐盆石出水的時候.
上帝就會很奇妙地.
賜水給他們.
就是這個鑄棚節.
在七天慶祝的時候.
他們都會有這個儀式.
但是去到第八天的時候.
就不再做這個儀式.
耶穌基督就在第八天的時候.
在這個慶典裡面.
有這個催國.
是一個很歡樂的日子的時候.
耶穌基督就起來宣佈.
因為過往七天.
剛才去西羅亞祠取水.
都是一個很重要的儀式.
是一個以水來做主題.
所以就是那天沒有了.
沒有再這個儀式的時候.
他就起來宣佈.
人若渴了.
可以到我這裡來渴.
信我的人就如經上所說.
從他腹中要流出活水的江河來.
耶穌的宣告.
他可以把這個活水賜給每一個相信他的人.
他其實預言到將來.
他死了三天之後復活.
復活了之後.
接著五十天之後.
就是聖靈降臨.

$^{161}$就是說當我們有聖靈在我們生命裡面的時候.
我們就可以衝破我們生命裡面.
種種的難處.
當我們感到口渴的時候.
我們又怎樣?.
想喝水.
當我們吃了一些味精很重的東西.
或者在現在這個夏天的時候.
當我們行山.
或者當我們出街的時候.
如果沒有水的時候.
我們會覺得很辛苦.
但耶穌不是問我們是不是想喝那些水.
耶穌想問我們.
在心靈的裡面.
我們是不是會覺得很乾渴呢?.
好像剛才弟兄一直在見證和分享的裡面.
他一直就算信了之後.
有時候在心靈的裡面.
都會有些艱難.
有些掙扎.
也有些反省.
不只是弟兄.
而是我們在我們生命裡面.
就算我們信了耶穌多久.
十幾年,二十幾年,三十幾年都好.
有時候我們都會感到心裡面很乾.
當時為什麼耶穌基督這樣站起來說呢?.
就是那幫猶太人.
他們很慘的.
是受著羅馬政府的統治.
受著他們的逼迫.
很多無道理的一些規條.
也都受著當時的一些宗教規條.
特別那些法利塞人.
自己就定了不知道多少條多少條規矩出來.
錯了又說不行又說犯了規.
還有在當中生活非常的艱苦.
所以他們面對著很多生活的壓力.
但回到我們今天的時候.

$^{201}$我相信我們今天受的那些壓力.
不會比當日他們那些人所經歷的輕.
我們現在很多人上班的時候.
很多都告訴我們.
有上班時間沒有下班時間.
而整個經濟低迷.
我們有時候都會令我們擔心的.
特別我們這些退休的人士.
你想一下之前一直存下來的那些夠不夠用.
又不知道天父給我們有多長命.
不知道夠不夠用.
那些物價就一直飆升.
那些人去看醫生.
前一陣子我們有個同工去看完醫生回來.
他說現在看政府醫院都很貴.
以前那幾十塊就包了藥.
現在另外包藥.
他說六百多塊走.
嘩,現在這麼厲害.
所以當我們想到現在這樣的情況的時候.
我們心都會害怕.
存了那麼多錢夠不夠用呢?.
又或者當我們一班年輕人的時候.
我們在想我們的工作.
到大一點的時候.
就可能在想我們一些感情的問題.
有時候心裡面有些很特別的.
會覺得有些的厭悶.
有些覺得很乏力.
很孤單.
又覺得好像沒什麼盼望.
也都會覺得有時候好像沒什麼人關心.
這些種種的心靈軟弱.
就是因為缺少了這個水.
這個水不是我們飲用的水.
是我們生命裡面的水.
奧古斯丁.
他就說我們生命裡面.
生命裡面有很多的洞洞洞.
在這個洞洞裡面.

$^{241}$需要有不同的東西來填滿它.
有個洞是需要我們有健康.
沒有健康的時候.
如果你是長期病患者的時候.
你就知道裡面那種慘.
有些人是需要金錢.
你有金錢的時候.
你就覺得安定了一點.
有些人需要有個愛情.
所以為什麼每一個人.
就時時都要走去.
是有去拿.
就是多大都想找個男朋友.
女朋友.
結了婚的.
不單有自己的太太.
還要有第二第三個.
是不是?.
有些人需要有個成就感.
為什麼我們要走去.
就是不斷地去長進呢?.
但是我們發覺到.
就是當我們擁有了.
無論有了健康.
有了金錢.
有了愛情.
有了這個成就.
但是我們心靈裡面.
有時我們都感到不滿足.
奧古斯丁說.
就是我們心靈裡面.
我們其中一格很重要.
很重要的.
就是我們需要找到上帝.
而這個上帝.
是一個真真正正.
可以滿足到我們心靈裡面.
那個情況的上帝.
而當我們填了這個上帝.
在我們心靈裡面的洞的時候.

$^{281}$我們就可以面對這個生命裡面.
無論是不夠錢用.
或者我們一個長期病患者.
或者我們找不到愛情.
或者我們不是有很高的成就.
但是我們都可以有一個豐盛的人生.
我們都會覺得.
整個人生很有意義.
所以我們就想一下.
在這個我們心靈裡面.
我們的生命裡面.
我們究竟有沒有那份力量.
有沒有那份平靜.
有沒有那份覺得.
我們真的很滿足那個的境況.
特別是在這個黑暗的世代裡面.
我們能不能夠去為上帝.
去做一個光明子女呢?.
剛才弟兄一直去做這個見證的時候.
我很欣賞他.
他經常去反省他自己的生命.
覺得我這樣做得不是很好.
昨天又發了脾氣.
無論自己發也好.
別人發也好.
都是不對的.
所以我們會覺得.
我們經常都會這樣.
做錯這樣又做錯那樣.
又做錯了那樣的.
但是不要緊.
你當我發了脾氣.
你自己說我經常做錯事.
不要緊.
最重要的是什麼?.
是這個的.
肯認罪,肯悔改.
《馬太方》第五章第三節.
這八幅裡面.
第一幅就是說.

$^{321}$不是第一幅.
饑渴無意的人有福了.
因為他們必得飽足.
還有五章第三節.
這是第一幅.
貧窮的人有福了.
貧窮就是指我們心靈裡面.
覺得我們自己不足夠.
覺得我們裡面.
可能很多時候.
我們有些事是做得不好的.
這樣很重要.
當我們唯有.
當我們自己知道不行的時候.
我們向天父上帝認錯.
我們向天父上帝.
去求知取力量的時候.
上帝就會來到我們當中.
肯這樣認罪.
上帝的靈就會充滿著我們.
所以我們要被聖靈充滿的時候.
我們必須先將.
我們生命裡面的雜質.
全部倒出來.
將我們裡面的苦澀.
我們裡面的驕傲.
我們裡面的不好的東西.
全部倒出來.
倒出來的時候.
聖靈才可以完完全全地.
充滿著我們的生命.
然後耶穌就說.
當我們知道自己口渴的時候.
去到祂的面前的時候.
祂就會幫助我們.
叫我們腹裡面能夠流出活水的江河.
這個活水.
耶穌基督就指到這個聖靈.
有聖靈的同在.
聖靈.

$^{361}$你看到.
當祂說活水不是污糟水.
污糟水又臭又污糟.
這個活水就是清新的.
是說到聖潔.
活水也是一直不斷流的.
是有能力.
有智慧.
有活力的.
葛培里穆斯.
他講這段經文的時候.
他說上帝的祝福不是一點一滴.
不是滴滴那麼多.
逐滴流的.
而是上帝的靈.
所給我們的恩典.
是不斷不斷流的.
他說就算在世界上很多很出名的.
河流,墨西西比河,阿瑪遜河.
我們的長江,我們的黃河.
你不會說我每天去舀幾壺水喝.
喝了一年之後就乾了.
不會的.
因為是很多很多很多.
所以無論我們怎麼舀都好.
它仍然可以有很多很多.
何況是我們的上帝.
祂的能力,祂的恩典是源源不絕的.
所以說當聖靈流到我們的生命裡面的時候.
祂不斷令到我們生命裡面的雜質枯乾.
流出來的時候.
我們的生命不再感到枯乾.
無論我們擁有的東西.
有沒有比其他人多還是少.
或者我們自己的性情.
我們自己的惡方面的東西.
無論如何.
上帝可以叫我們生命裡面可以穩妥下來.
所以很多人在聖經裡面都講到很多.
他自己本來在生命裡面很多乾涸的人.

$^{401}$但是去到耶穌基督的根前.
他就有改變.
我們記得當在約翰福音的時候.
有一個婦人.
她中午的時候出去拿水.
她有五個丈夫.
就是說她是一個不太妥當的婦人.
她感到羞恥,感到罪惡.
才會在中午的時候才出去拿水.
但是當她和耶穌聊完天.
她就知道自己不對了.
她竟然可以告訴別人.
耶穌基督知道了她那些東西.
耶穌基督就是真光.
我們又看到那個瑞利長.
很矮很矮的那個薩蓋.
以前又貪錢.
又去欺詐.
但是當他請了耶穌基督回家的時候.
結果他竟然可以將自己一半的身家分給窮人.
另外以前欺詐過的那些人.
欺詐過的那些人.
四倍的還給了他.
所以今天我相信.
我們在這裡每一個人.
我們相信了耶穌基督.
有聖靈充滿了我們.
但是我都說了.
我們在裡面.
我是有多少自我那些東西倒了出來.
有多少倒了出來.
越倒得多的時候.
聖靈就可以在我們當中的生命裡面.
做更多的工作.
我們的榮譽顧問.
我們錦繡堂榮譽顧問陳一環牧師.
他的太太有柏金遜症.
他也很辛苦.
我記得有一天.
有一位姐妹告訴我.

$^{441}$她說師母很淒涼.
有一天可能她的傭人姐姐出去買菜.
在家裡的時候.
不經意就跌倒.
跌倒在地上.
因為她的情況就是.
她自己跌倒在那裡.
她自己不能站起來.
那天是冬天.
就在地上坐了四十五分鐘.
直到她的姐姐回來.
知道之後就去探望她.
跟她說那次很淒涼.
又這麼冷.
坐在地上.
她應該是睡在那裡了.
動不了很慘.
但是她仍然是嘻嘻哈哈.
嘻嘻哈哈.
有些上帝完全信服上帝給她.
聖靈在她當中給她力量.
我的愛人就覺得很慘很淒涼.
但她覺得她不能死.
我更加驚訝的是.
她自己搞了一個團契.
是柏金遜病人的團契.
她就說她自己知道.
柏金遜那種慘況.
有時候她一直都說.
她自己坐地鐵的時候.
那時候可以走路的時候.
坐地鐵.
她可以突然間不能動.
站在閘口那裡.
竟然過不去.
後面的人又罵她.
又催她.
她很知道那個辛苦.
於是她就想.
其他人可能都是這麼辛苦.

$^{481}$所以她就想.
要成立一個團體來彼此鼓勵.
但是周圍的人就說.
怎麼可以?.
你們每個人都不能走路.
去哪裡?.
家裡或者哪間教會聚集好?.
還有每個人覆診的時間都不同.
還有很多次覆診.
定哪個日子好?.
不要想了.
師母你不要搞那麼多事.
一盆一盆冷水淋在她身上.
但她裡面那份熱誠.
她要覺得她要聯絡其他.
所以開始找她的女婿.
找她其他的親戚.
來幫她開工.
就是因為有活水.
這個聖靈在她生命裡面.
一直給她力量.
給她方法.
然後她用行動.
雖然行動很少.
用她的思想.
用她的邀請.
誠意邀請.
無論多艱難都好.
她終於可以成立.
就是在網上.
聯絡柏金遜的弟兄姊妹.
在那裡大家彼此分享.
裡面的難處.
所以有聖靈在我們生命當中的時候.
在裡面就有這個.
湧流著活水的力量.
源源不絕.
所以師母說她有時會灰心.
有時會難過.
有時會沒力.

$^{521}$但是這個聖靈裡面.
就不斷不斷地給她力量.
在這裡也想說一套電影.
在2013年的.
叫做《1942》.
有沒有聽過?.
《1942》.
你上網去看一看.
這套電影就說到.
在1942年的時候.
在當中河南省.
遇到大饑荒.
同時間在那裡.
軍閥管理著那個地方.
這群人民就很慘.
這套電影就由馮小剛導演.
他也挺有名的導演.
他就拍的.
結果就拿了很多的獎.
當他拍這套電影的時候.
他就故意走回河南.
看看有沒有這個.
在1842年的時候.
河南大饑荒的時候.
有沒有人.
到2013年的時候.
仍然在那裡.
他就想做一些訪問.
能夠更加明白.
當中的感受.
去到河南的時候.
被他千辛萬苦.
找到兩位老人家.
就問他當日的情況.
這兩位老人家.
他就說.
他最初想著.
這麼慘的事情.
他一定會非常憂愁.
或者是非常苦.

$^{561}$但他形容這兩位老人家.
是很開心的.
馮小剛導演就問他.
你在那段時間這麼苦.
中國大陸之後一直面對大戰.
國民黨和共產黨的戰爭.
還有很多的艱難.
為什麼你仍然這麼開心呢?.
這位老人家就說.
因為我們是信耶穌的.
馮導演就說.
信耶穌就可以這麼開心了嗎?.
老人家就說.
信耶穌可以上天堂.
馮導演就問他.
上天堂就不會餓了嗎?.
上天堂就不用逃難了嗎?.
他們就回答得很爽.
他們說上天堂有水喝.
就不會餓死.
哪裡有水喝呢?.
這個水就是聖靈的水.
唱一唱那首歌好不好?.
你會不會唱一首《天上的歌》?.
會吧?.
1,2,3,4.
(唱歌).
很多人都會,我們唱快一點.
我們知道有聖靈在我們心靈裡面的時候.
無論頭上有多少的烏雲.
心裡面有多少的憂傷.
全都灑下來,好不好?.
(唱歌).
我要唱歌一首歌.
唱耶穌天上的歌.
頭上有多少的烏雲.
心裡面有多少的憂傷.
全都灑下來.
我們一起祈禱吧.
因為你很愛我們.

$^{601}$拆敗了你的獨生兒子來到這個世界上.
為了我們的罪孽,你去捨身.
因為當主耶穌基督去請你的時候.
你就將聖靈給過我們.
臨到我們當中.
在我們的心靈裡面.
成為了我們很多的鼓勵,很多的支持.
主祝福紅印堂裡面每一位弟兄姊妹.
在我們心靈裡面.
有時候我們真的會感到乾渴.
有時候我們感到孤單.
有時候感到無力.
主就求聖靈親自在當中.
加力量給我們.
也親自給我們看到.
聖靈在我們心靈裡面.
日日夜夜與我們同在.
我們不孤單.
當我們陷在這個罪孽裡面的時候.
願聖靈都提醒我們.
以致我們快快的脫離過犯.
我們要聖潔.
因為主你就是聖潔.
求主幫助我們.
求主賜恩.
祝福我們每一個人.
生命裡面就時時刻刻流著活水江河.
讓我們的生命就見證.
上帝你是那位恩典滿滿的上帝.
我們多謝你.
誰聽禱告.
是奉靠主基督的聖名.
\newpage



\section{創世記 39:1-23}
\label{sec:kpqG_RsMI0U}
\textbf{2026.05.31 《有神與無神之別》:創世記39\_1-23(和修版) 杜文翰牧師}
\newline
\newline
連結: \href{https://youtube.com/watch?v=kpqG-RsMI0U}{\texttt{https://youtube.com/watch?v=kpqG-RsMI0U}} ~~~~ 語音日期: 2026-06-05
\newline
\newline
\hyperref[sec:3yIpsNheSeo]{\small{< < < PREV SERMON < < <}}
~
\hyperref[sec:index_chronic]{\small{[返順時目]}}
~
\hyperref[sec:index_scriptual]{\small{[返順卷目]}}
~
\hyperref[sec:iob9mODZZqY]{\small{> > > NEXT SERMON > > >}}
\newline
\newline
創世記 39:1-23
\newline
\begin{longtable}{cl}
\hline
\hline
章節 & 經文 (和合本修訂版)\\
\hline
39:1 & \begin{tabularx}{0.7\textwidth}{X} 約瑟被帶下埃及去。有一個埃及人波提乏，是法老的官員，是護衛長，他從那些帶約瑟下來的以實瑪利人手中把約瑟買了去。 \end{tabularx} \\ \\ \relax
39:2 & \begin{tabularx}{0.7\textwidth}{X} 約瑟在他埃及主人的家中，耶和華與他同在，他是一個通達的人。 \end{tabularx} \\ \\ \relax
39:3 & \begin{tabularx}{0.7\textwidth}{X} 他主人見耶和華與他同在，又見耶和華使他手裡所辦的事都順利， \end{tabularx} \\ \\ \relax
39:4 & \begin{tabularx}{0.7\textwidth}{X} 約瑟就在主人眼前蒙恩，伺候他主人，主人派他管理家務，把一切所有的都交在他手裡。 \end{tabularx} \\ \\ \relax
39:5 & \begin{tabularx}{0.7\textwidth}{X} 自從主人派約瑟管理家務和他一切所有的，耶和華就因約瑟的緣故賜福給那埃及人的家；凡家裡和田間一切所有的，都蒙耶和華賜福。 \end{tabularx} \\ \\ \relax
39:6 & \begin{tabularx}{0.7\textwidth}{X} 波提乏把他一切所有的都交在約瑟手中，除了自己所吃的食物，其他的事一概不知。約瑟英俊健美。 \end{tabularx} \\ \\ \relax
39:7 & \begin{tabularx}{0.7\textwidth}{X} 這些事以後，約瑟主人的妻子以目送情給約瑟，說：「你與我同寢吧！」 \end{tabularx} \\ \\ \relax
39:8 & \begin{tabularx}{0.7\textwidth}{X} 約瑟拒絕，對他主人的妻子說：「看哪，一切家務我主人一概不知，他把所有的都交在我手裡。 \end{tabularx} \\ \\ \relax
39:9 & \begin{tabularx}{0.7\textwidth}{X} 在這家裡沒有人比我更大，除你以外，他也沒有留下一樣不交給我，因為你是他的妻子。我怎能行這麼大的惡，得罪神呢？」 \end{tabularx} \\ \\ \relax
39:10 & \begin{tabularx}{0.7\textwidth}{X} 她天天這樣對約瑟說，約瑟卻不聽從她，不與她同寢，也不和她在一起。 \end{tabularx} \\ \\ \relax
39:11 & \begin{tabularx}{0.7\textwidth}{X} 有一天，約瑟進屋裡去辦事，家裡沒有一個人在那屋子裡， \end{tabularx} \\ \\ \relax
39:12 & \begin{tabularx}{0.7\textwidth}{X} 婦人就拉住他的衣服，說：「你與我同寢吧！」約瑟把衣服留在她手裡，逃出外面去了。 \end{tabularx} \\ \\ \relax
39:13 & \begin{tabularx}{0.7\textwidth}{X} 婦人看見約瑟把衣服留在她手裡逃到外面， \end{tabularx} \\ \\ \relax
39:14 & \begin{tabularx}{0.7\textwidth}{X} 就叫了家裡的人來，對他們說：「看，他帶了一個希伯來人到我們這裡戲弄我們。他到我這裡來，要與我同寢，我就大聲喊叫。 \end{tabularx} \\ \\ \relax
39:15 & \begin{tabularx}{0.7\textwidth}{X} 他聽見我放聲大喊，就把他的衣服留在我這裡，逃出外面去了。」 \end{tabularx} \\ \\ \relax
39:16 & \begin{tabularx}{0.7\textwidth}{X} 婦人把約瑟的衣服放在身邊，直到他主人回家， \end{tabularx} \\ \\ \relax
39:17 & \begin{tabularx}{0.7\textwidth}{X} 就用這樣的話對他說：「你帶到我們這裡來的那希伯來僕人進來要調戲我， \end{tabularx} \\ \\ \relax
39:18 & \begin{tabularx}{0.7\textwidth}{X} 我放聲大喊，他就把衣服留在我身邊，逃到外面。」 \end{tabularx} \\ \\ \relax
39:19 & \begin{tabularx}{0.7\textwidth}{X} 主人聽見他妻子對他說的話，說：「你的僕人就是這樣對待我」，就非常生氣。 \end{tabularx} \\ \\ \relax
39:20 & \begin{tabularx}{0.7\textwidth}{X} 約瑟的主人把他抓起來，關在監獄裡，就是王的囚犯被關的地方。於是約瑟在那裡坐牢。 \end{tabularx} \\ \\ \relax
39:21 & \begin{tabularx}{0.7\textwidth}{X} 但耶和華與約瑟同在，向他施恩，使他在監獄長的眼前蒙恩。 \end{tabularx} \\ \\ \relax
39:22 & \begin{tabularx}{0.7\textwidth}{X} 監獄長就把監獄裡所有的囚犯都交在約瑟手下；在那裡的一切事都由他處理。 \end{tabularx} \\ \\ \relax
39:23 & \begin{tabularx}{0.7\textwidth}{X} 任何交在約瑟手中的事，監獄長一概不察，因為耶和華與約瑟同在，耶和華使他所做的都順利。 \end{tabularx} \\ \\
[1ex]
\hline
\hline
\end{longtable}
$^{1}$現在的題目就是「有神與無神之辯」.
和大家分享.
香港人其實對一些不公平的事.
現在多了一些渠道去發聲.
所以你會見到電視也好.
在網絡也好.
甚至有一些突發的電視節目也好.
都是去探討人們對一些不公平的事的發聲.
甚至有時可能你見到的場面.
都是一些很小的事.
資訊不多時.
哪怕可能是地鐵裡面讓位和不讓位.
有時都可以在媒體裡甚麼都可以見到.
當然大事更加可以在媒體裡見到更加多.
而且香港人同時能夠利用到不同的機制.
其實每個系統都有一些投訴機制.
譬如哪怕醫院也好.
政府部門也好.
屋苑也好.
甚至學校.
其實學校原來有投訴機制.
如果你不滿意某個老師所做的.
你可以用一些方法去投訴.
所以其實香港人現在懂得用各種方法去表達.
這些聲音.
叫「發聲」.
不過其實是在乎於你發聲還是不發聲.
發聲之後.
是否真的能夠表達得到.
結果是怎樣有時都不是在你掌握之內.
如果再拉近一點.
去到工作場所裡面.
其實有時我們都會有.
你會見到一些你不如意的事情.
或者你會覺得一些看不過眼的事情.
究竟你發聲還是不發聲.
那個處境又不一樣.
因為那個是你的工作你要賺錢的地方.
但你見到明明在裡面的爭鬥.
或者是有些人弄拳.

$^{41}$上層中層下層.
其實多多少少都會對員工有影響.
你是見到的.
你甚至身歷其境在其中.
究竟在當中你遇到這些事情的時候.
是否就能夠你有方法去表達.
即是表達的後果又會如何.
你又要想清楚.
所以其實來到工作場面的時候.
發聲還是不發聲.
那個想法又不一樣.
我們說基督徒是相信有神的.
當我們將信仰元素加入到我們所發生的每一件事上.
其實你又會發覺有時神不是每件事都要插手干預.
所以由得他這件事繼續發展下去.
所以在那個情況下的時候.
我們應不應該繼續去發聲.
還是好像覺得.
聲音就說了.
不過好像那件事沒有改變.
究竟神是否存在呢.
在基督徒的信仰裡面是有神的.
不過我們去體驗我們每一天生活的時候.
又好像不是經常覺得神在我們身邊.
好像神不在場.
所以這個是很多時候我們信徒會面對的問題.
剛才主席跟我們讀的經文裡面.
我們看到這件事件是約瑟在三十九章裡面所經歷的.
我又不是很清楚當時的社會有沒有投訴機制.
如果有的話應該有很多事件可以投訴.
很多不公義的事件是發生了.
不過可能是沒有.
因為當時那麼古老的社會.
又未必有那麼多機會讓你去發聲.
在這個處境裡面.
特別是早期約瑟在埃及的處境裡面.
很多時候如果我問你這件事究竟是有神還是沒有神.
在這整個三十六章裡面.
你的結論是有神還是沒有神呢.
神是否存在呢.

$^{81}$有人說沒有神的.
因為神不在場.
他好像離開了約瑟.
因為在第一件事件裡面我們看到的是.
神母難阻約瑟的兄長將他賣給二十萬里人.
當時其實如果你看回三十七章的背景.
約瑟很年輕,十七歲.
不過他父親雅各特別寵愛約瑟.
所以為他預備了一件很漂亮的衣服.
其實不一定是彩衣的.
但是引致約瑟哥哥的妒忌.
一開始由妒忌變成恨,再由恨變成想要殺他.
一路一路加深.
所以當約瑟帶領他.
奉他父親的名.
將物資送到他哥哥那裡的時候.
其實只是一個關懷的行動.
不過哥哥就已經開始遠遠望見他.
就開始想辦法了.
究竟可以怎樣對付他.
所以一開始的時候他們只是.
想賣掉他.
但後來有人想殺了他.
拯救他,就少了煩惱.
但後來就有哥哥出面.
裡面想去救他.
不過最後他們哥哥談好之後.
就由猶大的哥哥獻計.
將約瑟賣給剛剛經過的二十萬里人.
然後二十萬里人就帶了他去埃及的地方.
這件事明明是一個.
今天所理解.
販賣人口是大事.
為什麼神不攔阻這件事發生呢.
以至約瑟去到埃及失去了自由.
在這個過程裡面.
看不到神的任何的干預.
所以有人覺得好像沒有神在這件事裡面.
第二件事是約瑟被主母誣告.
更加看到好像上帝離場.

$^{121}$從整件事的過程中.
剛才我們讀的經文讀了.
我們就不詳細逐節解釋了.
其實起因就是約瑟曾經說.
他英俊健美,非常有吸引力.
不過吸引錯了.
反應錯了在不同的人身上.
反而是主母的太太對他有興趣.
主人的太太對他有興趣.
於是吸引了他.
先用各種方法.
就是用言語去說服他.
你要我同寢白.
約瑟一口拒絕之後.
在後面你會看到.
於是這個主人就天天跟他說.
所謂疲勞轟炸.
天天不斷跟他說.
其實也不容易熬的.
然後終於有一個機會.
就是只有他跟約瑟在一起的時候.
主母就拉著他的手.
約瑟就逃跑.
然後留下了一件衣服.
事件往後更加發生的就是.
主母反而改變了態度.
整件事扭轉,編了一個故事出來.
說是約瑟進來調戲他.
他留著一件衣服.
當主人回來的時候.
就拿著一件衣服告訴他.
你這個欺負內部人.
這樣這樣這樣.
主人又相信他.
結果主人就將約瑟關到埃及法老的監獄裡.
於是約瑟就變成階下囚.
一層一層越來越低.
就跌到谷底.
所以你看到約瑟在兩件事的過程上.
被埋到被主母誣告的過程上.

$^{161}$看不到有神的介入.
神好像沒有幫約瑟脫險.
一直留在監獄裡.
究竟神是否存在呢?.
當我們去探討這個問題的時候.
其實聖經給我們答案.
究竟是有還是沒有.
當你剛才讀的時候.
不知道你有沒有特別留意.
其實有幾個字眼給了一些線索.
首先表面上好像耶和華沒有立刻阻止這兩件事發生.
但其實當中有些線索就是.
原來第二節他說.
約瑟在他埃及主人的家中.
耶和華與他同在.
第二節.
然後第三節就說.
他主人見耶和華與他同在.
又見耶和華使他手裡所辦的事都順利.
第五節.
耶和華就因約瑟的緣故賜福給那埃及人的家.
凡家裡和田間一切所有的.
都蒙耶和華賜福.
還有.
但耶和華與約瑟同在.
向他施恩使他在監獄長的眼前蒙恩.
二十一節.
還有二十三節.
因為耶和華與約瑟同在.
耶華使他所做的事都順利.
你看到嗎?.
耶和華的字出來了.
聖經說耶和華與他同在.
首先是主人的家同在.
而且主人見到耶和華與約瑟同在.
這個同在家裡的人.
他的主人是見到的.
清清楚楚見到的.
而且主人更加將其他的事都交在他手上的時候.
見到結果是好的.

$^{201}$是順利的.
田間又順利.
家裡又順利.
即使後來約瑟被關到監獄裡的時候.
監獄長都看到耶和華與這個人同在.
所以更加放手讓約瑟去管監獄裡的其他人.
所以你看到.
耶和華同在有幾個結果讓我們看得到.
就是辦事順利,蒙此福,所做的事都順利.
在這個過程裡.
在耶和華與約瑟同在中間.
他的處境很差.
圍爐,甚至被關在監獄裡.
但我們總看到一些線出來.
一些好的事.
一些好的結果.
是蒙福的.
周圍的人是蒙福的.
辦的事是順利的.
這些就是好的結果.
所以這些好的結果,而且不單止從約瑟本人開始.
更加延伸到周圍的人.
他的主人都看到.
再延伸出去.
他的家裡的事.
都可以樣樣順利.
在他手中.
主人將所有的事都加在他手裡.
一切都是好的.
甚至在監獄長手中的時候.
都是好的.
所以你看到.
耶和華與約瑟同在.
是一件好事.
但是是好的結果.
所以在整個過程裡.
我們可以說.
其實耶和華沒有離場.
耶和華是和約瑟一起在圍爐的時候.
也一起在監獄裡的時候.

$^{241}$沒錯,約瑟跟以前的三位以色列的列祖真的很不一樣.
大家可能熟悉的亞伯拉罕以撒雅各.
耶和華與這三個人同在的方式很特別.
真的是面對面聊天.
而且有意象.
還有聲音.
很真實的.
就在眼前看到耶和華是怎樣的真實.
在亞伯拉罕以撒雅各的眼中.
是清清楚楚看到的.
不過來到約瑟,約瑟沒有這些經歷.
沒有耶和華在他面前顯現.
沒有耶和華跟他說話.
連聲音都沒有.
好像是沒有這些.
不過耶和華與約瑟同在的方式只是轉換了.
不是用顯現的方式出現.
而是用同在的方式出現.
這是一個特別的地方.
所以仍然我們可以說.
這個過程是有神跟約瑟在一起的.
在波提法家的家裡.
在主母孕幼之前和下監之後.
都是一個有神的經歷.
因為看到有神的印記.
耶和華在後面.
祝福著約瑟身邊的兩個人.
主人和教堂長.
以至連間所有的事.
還有耶和華與人同在的.
有兩方面我們值得注意.
第一,這個同在原來是一個超越環境的同在.
哪怕你不太順利.
偷偷碰著黑.
甚至可能被人欺負等等.
各種狀況裡面.
即使被賣.
受冤屈,坐牢.
原來神仍然是跟約瑟在一起的.
這個同在是超越了環境的.

$^{281}$正如保羅所說.
難道是患難嗎?是刀劍嗎?是碧白嗎?等等.
都不能夠使我們與神的愛隔絕.
其實都是在呼應.
《創世記》這段經文.
神的同在超越了這一切的環境.
這是第一方面.
第二方面我們看到.
神的同在是真實可見的.
起碼約瑟的主人看到.
起碼監獄長看到.
於是他有的結果.
給他一些好東西.
而且這個同在帶來的.
就是他的家人.
家裡面田間都受惠.
等等等等.
有耶和華的同在是真實的.
不是純粹說耶和華祝福你.
不是這樣的.
是落到你的生活中.
你見得到的.
我不知道你自己看到這個的時候.
你自己看.
耶和華是不是真的常常跟你在一起呢?.
你是不是真的見得到.
是不是可見的呢?.
哪怕你是初信.
盡了自己一段時間.
盡了自己很久的領袖.
你能不能仍然說得出.
耶和華是跟你在一起的.
而且你能夠說出.
神怎樣跟你在一起.
同在的方式是怎樣的.
每個人可以不同的.
不過仍然可以很肯定地告訴你.
神在這裡的.
我自己都經歷很多神同在的印記.
只是在這裡爭取一點時間.

$^{321}$跟大家分享一兩件.
我自己很深刻的就是.
去年的時候.
我太太的家人患病.
其實在這個過程上.
我們一直都很希望這位太太的二姐.
能夠有機會信耶穌.
過去我們都用不同的渠道.
來跟她分享福音.
甚至邀請傳祖人到她家探訪.
只不過這位二姐仍然很客氣.
笑笑口.
但她仍然是當她被邀請的時候.
都沒有什麼反應.
都沒有信耶穌.
所以有時特別是患病.
她是癌症的末期.
所以更加著急.
怎麼辦呢.
有沒有機會能夠答應她信了耶穌呢.
在心中只能默默地為她禱告.
不過神的工作很奇妙.
有一次在醫院.
她在病床的時候.
有一位她比較相熟的姐妹.
在另外一間教會.
就來探訪她.
很奇怪.
她一開口.
這位姐妹就馬上問二姐.
你願不願意開口接受耶穌.
有時我們都不夠膽直接開口.
因為她很熟悉.
熟悉就馬上開口問她.
竟然二姐的回應是點頭.
我又在場.
太太又在場.
我說是呀,好呀.
因為她當時已經病得很重.
還可以有些回應.

$^{361}$我就跟她說.
你點頭.
你要將你的決定告訴耶穌.
你要開口.
還說你要去承認你的罪.
然後將主權交給耶穌.
她都說願意.
於是我一句一句.
帶她到床邊去祈禱.
她真的很慢.
一句一句.
可能說得不太清楚.
但我知道她在說.
然後整個祈禱就完結了.
然後走的時候.
她甚至先跟我們揮手.
你知道病人動一動都會痛.
都很辛苦.
然後過了一兩天之後.
還安排了教會的牧者幫她灑水禮.
然後灑了禮之後.
好像一兩天就離開世界了.
回到天父那裡.
從整件事的過程中.
我看到很奇妙.
前面好像沒有什麼可能.
沒有什麼希望.
就是等到什麼時候才到呢.
原來神的時間到了.
就能夠在裡面帶領她去信耶穌.
所以用這件事.
近的.
印證到神的同在的印記.
在我們和家人的家中.
同時如果過去.
我們在南非的十七年的經歷裡面.
其實時間又不夠了.
不可以說太多.
當很多很多神同在的印記.
我有本寶.

$^{401}$真的去為各樣的事祈禱.
當我們越寫的時候.
越發覺這本寶真的很珍貴.
原來神真的答應很多的事.
有一件事沒有寫下去.
就是很早之前.
我記得我兒子去學琴.
這個.
應該五歲左右,很小.
當時有個老師教他的.
那個老師教到他差不多.
差不多就可以報考試了.
但很奇怪,不知道為什麼.
考試前一個月.
那個老師說我沒空.
我教不了你.
他就不教了.
突然間退出.
很緊張.
下個月就考試了,那怎麼辦呢.
但你又要照樣考.
於是就照樣考.
我又在家幫他.
考完之後發覺成績很好.
很好.
所以老師說了一句話.
他說他很詫異,很驚訝.
你成績很好.
這位老師就說.
可能是因為你爸爸是牧師.
所以神賜福給你.
這位是從老師眼中.
口中出來的,不是我自己說的.
所以更加看到.
原來整件事是有神同在的印記.
哪怕是一件很小的事.
只是按琴考試而已.
都有神的同在印記.
所以你看到.
等於你可以去找.

$^{441}$大事小事.
都是可以找到神同在印記.
哪怕是順境不順利.
都可以找到神同在印記.
從弱失的生命這件事件裡面.
我們很肯定可以看到.
這是第一方面.
當我們留意中間那一段.
特別是主母引誘弱失的那一段.
第七節至二十節的時候.
你會發覺沒有耶和華語.
與弱失同在的字眼.
可能你會說.
那就是沒有了.
那就是神沒有與弱失同在.
特別是在受試探的時候.
離開了祂.
可能你會這樣想.
但是.
只要我們留心一下.
整段經文.
有兩節是很突出的.
八至九節.
是弱失自己的說話.
「看!一切家務我主人一概不知.
他把所有的都交在我手裡.
在這家裡沒有人比我更大.
除你以外.
他沒有留下一羊不交給我.
因為你似他的妻子.
我怎能行這麼大的惡得罪神呢?」.
這句話的重點是.
弱失清楚地表白.
主母是主人的太太.
弱失不能夠與她同寢.
因為這是神眼中一段大的惡.
弱失不能夠做這件得罪神的事.
換言之,從這句話裡我們可以反映到.
弱失的內心是有一把尺子的.
這把尺子是一個標準.

$^{481}$這個標準是一個行事為人的標準.
人是不能夠越過這個神的標準.
換句話來說.
當他這樣去表達的時候.
其實神的標準,神的話.
原來在哪裡呢?.
原來藏在弱失的心靈裡.
就在他心靈裡一直與他同在.
弱失沒有忘記.
即使他在困難中.
在被壓迫的處境中.
仍然是堅守著.
神的話要行第一.
我甚至可以說這是弱失心靈的高峰.
因為這不是在順利的時候.
風調雨順.
當然容易說是好的見證.
不是.
弱失這個時候所表達的是.
在壓迫,受冤屈.
被冤枉.
甚至他不知道是什麼後果的情況下.
他仍然在短時間.
即是當他主母拉著他的衣服的時候.
短時間來作出決定.
這是有神的決定.
是真正在他心靈裡的決定.
而且甚至可以說是活出了一個有神的生命.
這種有神的生命不是口說的.
是他行動表達了的,就是離場.
縱然他不知道後面是怎樣.
他的手,他的腳.
一致地去配合.
弱失的內在生命.
這是一個知行合一的生命.
所以我們看到.
弱失是.
神是與弱失同在.
甚至可以說.
神的話成為弱失生命的一部份.

$^{521}$重要的元素.
因為雖然沒有一個叫神看著他們.
當時只有兩個人在.
但是弱失仍然是拒絕堅守.
神的話是他生命的最重要的部份.
神的話植入了弱失生命的DNA基因裡面.
進去了.
不少人我們說口裡有神.
究竟我相信神.
我去教會.
不過想深一層.
神是不是真的在每一天的生活裡面.
在你的生活裡面.
是不是表現了出來呢?.
有時可能是你表現出無神的生命.
意思就是說.
我們沒有按神的心意去做每一件事.
沒有去做決定的時候.
去依照神的標準去做.
有時我們可能說.
回教會是一個樣子.
每個人都說有神.
那我就要活出一個有神的樣子.
我要告訴別人我是相信有神的.
不過離開這裡之後.
沒有人看到.
我可以變成另一個樣子.
變成一個沒有神的樣子.
神不在.
沒有把神放在我們面前.
有一節經文.
我們很喜歡的.
就是詩篇第十六篇第八節.
有兩個版本.
我讓耶和華常擺在我面前.
因他在我右邊.
我就不自動搖和收盤.
新表典和和盤.
我將耶和華常擺在我面前.
因他在我右邊.

$^{561}$我便不自搖動.
其實分別就是.
耶和華常擺在我面前.
為什麼我們喜歡這句說話呢.
因為有耶華看著我們.
特別在我們需要他幫助的時候.
當然有耶華在我們身邊.
是一個很大的應許和祝福.
我們都很喜歡.
所以有神在看著就很好了.
不過不要忘記這句說話.
是應用在所有的情況裡.
包括好和包括你.
正在做罪惡試探的時候.
不要忘記耶和華在身邊.
意思就是.
耶和華其實是看著你.
究竟你如何決定.
你做還是不做.
你跟神的意思.
還是不跟神的意思.
耶和華都在身邊.
今天很多時候.
我們覺得很多事情.
只有你與我知道.
沒有其他人知道.
沒事的.
但是其實在兩個人之間以外.
其實有一個無形的耶和華在身邊.
聽到你們兩個人說的話.
不要以為沒有人知道.
今天我們社會有一個.
叫做僥倖的心態.
就是說做的事情.
特別是一些不對的事情.
我們常有一個僥倖的心態.
希望這一次不會被人發現.
或者是不要被人揭發.
或者做一些桌底交易.
或者在這樣的情況下.

$^{601}$不要緊的.
只有你與我知道.
沒有其他人知道.
甚至對基督徒來說.
可能都會說.
或者這件是小事.
即使我做了.
上帝都會饒恕我吧.
我又不是殺人放火.
不是什麼大事.
小事而已.
所以有時候我們有一種心態.
就是好像我們可以做一些.
好像給我們一種規範.
許可我們做一些違背良心的事情.
違背神的事情.
我不知道這些是不是常常.
我們有時間中會受到的考驗.
我認識一位弟兄.
他在工程裡面工作.
你知道工程.
他做到工程師比較高級.
他要簽一些文件.
他看完整份文件之後.
他發覺其實有一部分.
其實不符合條件的.
不符合規矩.
應該不可以簽的.
不過他把文件.
拿到上司面前的時候.
上司怎麼說呢.
你就照簽吧.
但是這個弟兄心裡面有一個掙扎.
是不是我簽了就上司會負責呢.
但其實不是.
拿著簽名是你的簽名.
是你要負責的.
就是說你同意這份文件所有的東西.
所以其實內心有掙扎的.
作基督徒在這些處境裡面.

$^{641}$簽還是不簽.
簽名而已.
小事而已.
不過帶來的後果是.
不知道什麼時候會有的.
你要去承擔.
所以有時候我們是不是抱僥倖的心態.
算了.
只是打份工而已.
賺過錢而已.
簽了就算了.
但會知道的.
在這麼艱難的情況下.
還是你仍然去表達.
不好意思.
我是基督徒.
我不能簽.
其實這個是抉擇裡面來的.
常常都發生的.
今天我們要小心這些處境.
另一個處境就是僥倖心態.
就是說當我們明明做了一件錯事.
被抓到了.
可能心念說.
這次不幸運而已.
被人抓到.
那下次死目一點.
繼續做.
死目一點.
想想辦法怎樣.
不要被人知道.
其實都是另類的僥倖心態.
我們內心仍然很想去做那些不應該做的事.
不過好像找很多藉口.
逃避著我們應有的責任.
其實究竟如果是基督徒的時候.
我們活出的是究竟是有神還是無神.
神是否在你每一件大大小小的事裡面.
其實今天我們每一天.
我們要去思想的決定.

$^{681}$神究竟在哪裡.
其實是在乎你的.
你怎樣表達出來.
周圍的人就看到.
這個人有沒有神.
他說他去教會.
他是不是真的有神.
其實是看得到的.
周圍的人看到的.
虐殺周圍的人看到他的.
所以今天我們真的要去看處境.
真的要去想清楚.
在有神和無神之間.
我們常常都要去做決定.
神是否真的常常放在第一位.
還是在某些處境裡面.
他轉了位.
轉了在第二位.
錢是第一.
或者職位是第一.
或者什麼什麼是第一.
立刻轉了.
然後後面又轉回來.
星期天又轉回來.
我們是不是這樣經常轉來轉去.
還是能夠一致地.
常常神在我們身邊.
套在虐殺的處境.
當主母拉著他的衣服的時候.
其實他有很多藉口可以說.
哎呀,這次我只是被動而已.
而且那個犯錯不是我.
是主母引誘我在先.
我只是一個受害者.
今天有時候受害者好像聲音大一點.
所以抱著受害者的心態.
不怕,他抓著我而已.
我是受害者.
或者逃脫得了.
等等.

$^{721}$不過.
約瑟沒有這樣想.
他仍然堅持著.
我要做.
神要我做的.
堅持.
走了才算.
所以約瑟當時仍然是清醒的.
他清清楚楚地知道.
他仍然將神放在第一位.
縱然那一刻只有他和主母在那裡.
神好像不在.
但是約瑟內心.
他知道神在那裡.
在約瑟每一天的生命裡面.
神在他周圍.
所以我可以說.
歸納約瑟去經歷神的同在有兩個層面.
第一個.
我們.
沒錯,看到單向的.
神賜福給人的.
約瑟經歷到的.
他指人也都經歷到的.
所以這個是.
就是祝福來到.
但是同時間也都雙向的.
因為約瑟以尊神為大的方式.
去表達著神和他在一起.
心裡面也是有神的.
約瑟的回應.
縱然主母引誘他的時候.
他仍然拒絕.
今天我們說.
我們的靈命.
究竟是表面的信.
還是真的有實質的信.
有時候我們很容易說.
表面信神就OK了.
但是真正一個有信仰的人.

$^{761}$是要落到你內心的深處裡面.
才能夠真的表達到.
時時刻刻都是有神.
我記得我2024年去南非.
因為我做功課研究.
原本是去問一些牧者,信徒.
那我問的問題就是.
你怎樣形容南非華人的信徒.
其實華人的信徒是少的.
於是他的答案是.
質又差,量又少.
好像沒有什麼希望.
南非有三十萬華人.
其實數一數.
真正信徒是一千左右.
或者一千以下,我不知道.
你數一數也不夠1\%,真是太少了.
而且很多都是剛剛信耶穌.
質又少,質又差.
他甚至形容質又差的意思是.
很多時候信徒做錯了事.
他就會走來跟牧師說.
牧師,我有罪了.
因為我得罪了神.
我今天做了什麼什麼事.
請你為我禱告.
求神赦免我的罪.
所以有時信徒都是這樣去做.
我們有時在南非常常都見到.
華人,信徒.
有時都會面對當時的處境.
特別在貪污文化,賄賂.
甚至是有錢可以解決的處境下.
有時他們都說.
即使牧師見到,牧師知道.
他都說牧師你要明白.
我們真的沒有辦法了,還是要這樣做.
因為有一次我陪一個人去翻譯.
他好像是因為什麼簽證問題被警察扣留.
很遠的,我開車跟他一起去,幾個人.

$^{801}$去到差不多的時候.
就跟相關的職員說.
究竟為什麼要扣留他.
說到差不多的時候.
其中一個華人就轉換話題.
我們可不可以在外面說.
即是離開辦公室.
離開辦公室就不是公事,即是私事.
我立刻離開,我知道他想準備做什麼.
搞完之後,其中一個人跟我說.
牧師你要明白,我們這些處境.
沒有辦法,一定要這樣做.
其實我不同意的.
所以我沒有在場,我不知道他們最後談成怎樣.
我們要表達.
沒錯,質是差,量是少.
但我們仍然是需要去找回.
我們建立信徒的靈魂.
就是要針對他的弱點,他的深處.
希望香港信徒不是像南非信徒那樣.
應該好一點的,香港有這麼多教會.
不過仍然關鍵的是.
有時候我們回教會很久.
有時候都會覺得.
這些我都知道,不用說了.
變成了腦身常談,我們有時候忽略了.
所以我們教會裡有時候都要去想.
不要只想著牧師要求我們做的事.
我都做了,我又去教會.
我又有奉獻,我又有祈禱.
我又有參與事奉,我甚至有傳福音.
不錯,很好,還想要什麼.
我又不是做什麼壞事.
這些不是差的,這些不是錯的.
應該要做的,我鼓勵的,要做的.
不過是不是做完這些之後.
你的內心生命就是相對地一樣呢?.
我們不要以為自動地.
我信主多一年,我的靈魂就深了一年.
不是的,有時候.

$^{841}$可能就停留,甚至下跌.
所以我們真的要小心.
我們特別信了主久的時候.
我們都要想一下是不是因此.
信了主久,自動我們就知道.
懂得調教,凡事都將神放在第一位.
所以這是一個深層的問題.
不是一個表面的問題.
於是有時候我們遇到大衝擊的時候.
有時候我們站不穩.
為什麼?神好像沒有跟我一起.
以前的小說就是拿一朵花,撕一些葉子.
神同在,不同在,神同在,不同在.
看看最後,神是不是跟我一起呢?.
就好像博一下運氣那樣.
你是不是這麼疑惑呢?.
已經捉不到神是不是跟你一起.
還是其實祂在這裡,不過你不察覺.
所以有神和無神,不是未信主之前.
所謂無神論者和有神論者的爭論.
不是單單在那個層面去討論有神和無神.
其實有神和無神更加落實到信主之後.
在你的不同階段裡面.
你能不能夠同樣去用你的口.
和你的生命去表達.
我現在是一個有神的生命.
當然這個過程需要時間.
有時候不能一下子就變成.
凡事都交錯主.
但問題是你願不願意走到那個方向.
還是你退後.
哎呀,很難呀.
做基督徒很難呀,我不做了.
回到原來的那裡.
是你的選擇.
在《創世記》第39章裡面.
其實很少記錄他所說的話.
比較多記錄的是其他人所說的話.
不過他顯示的那句話就精彩了.
原來是顯示了他生命的高峰.

$^{881}$正正從他的八至九節這句話裡面.
看到原來神是已經深藏在心裡面.
他沒有表達出來.
不過從這句話裡面.
他告訴我們.
神是在虐室的.
生命中一直在和他同在.
甚至已經進入到他生命的深處.
生命的基因.
有些詩歌說.
有耶穌就有一切.
不容易唱.
字就簡單.
幾個字.
你要用生命唱出來的時候.
我就要看你了.
是不是真的有耶穌就有一切.
還是你去到這樣的場景.
你就「騰雞」了.
所以其實真的是.
上上給我們很多的考驗.
很多的提醒.
求神幫助我們.
哪怕我們尋求中.
信了主一段日子.
信了主很久.
我們都要上上去問自己.
今天此時此刻這個決定.
是有神還是無神呢.
求主幫助我們.
或者和祂同居.
聖靈求你繼續向我們說話.
無論我們在什麼處境.
我們在慕道中.
我們在認識耶穌初階段.
或者我們已經信了耶穌很久了.
原來這個問題都是很根本.
都是要在我們內心去自己回答.
別人不能夠代我們去回答.
究竟我們心中是有神還是無神.

$^{921}$神究竟在哪裡呢.
還是有時我們喜歡的時候.
就將祂拿出來.
有時覺得不好的時候就放在一邊.
所以求主去憐憫我們.
幫助我們.
我們能夠願意有個渴慕的心.
我們真的願意去成長.
透過約瑟的生命.
表面上祂的經歷是一個悽慘的經歷.
但我們確實見到.
在整個過程中.
見到耶和華與祂同在.
處處與祂同在.
無論在主人的家中.
在監獄中.
甚至在被引誘的過程中.
都見到神同在的足跡.
所以求主幫助我們.
堅固我們的信心.
讓我們能夠在香港這個地方立足.
在這個地方面對很多不同的挑戰.
引誘的時候.
我們仍然能夠站穩立場.
我們知道神會賜福我們.
哪怕我們做的結果.
眾人不認同.
但我們知道神是與我們同在.
求主幫助我們.
我們禱告奉耶穌基督的聖名.
阿們.
\newpage



\section{路得記 2:4-18}
\label{sec:iob9mODZZqY}
\textbf{2026.06.07 《跟從神!向有需要者施恩》:路得記2\_4-18(和修版) 翟凱欣傳道}
\newline
\newline
連結: \href{https://youtube.com/watch?v=iob9mODZZqY}{\texttt{https://youtube.com/watch?v=iob9mODZZqY}} ~~~~ 語音日期: 2026-06-10
\newline
\newline
\hyperref[sec:kpqG_RsMI0U]{\small{< < < PREV SERMON < < <}}
~
\hyperref[sec:index_chronic]{\small{[返順時目]}}
~
\hyperref[sec:index_scriptual]{\small{[返順卷目]}}
~
\hyperref[sec:csMea85HLY0]{\small{> > > NEXT SERMON > > >}}
\newline
\newline
路得記 2:4-18
\newline
\begin{longtable}{cl}
\hline
\hline
章節 & 經文 (和合本修訂版)\\
\hline
2:4 & \begin{tabularx}{0.7\textwidth}{X} 看哪，波阿斯正從伯利恆來，對收割的人說：「願耶和華與你們同在！」他們對他說：「願耶和華賜福給你！」 \end{tabularx} \\ \\ \relax
2:5 & \begin{tabularx}{0.7\textwidth}{X} 波阿斯對監督收割的僕人說：「那是誰家的女子？」 \end{tabularx} \\ \\ \relax
2:6 & \begin{tabularx}{0.7\textwidth}{X} 監督收割的僕人回答說：「她是摩押女子，跟隨拿娥米從摩押地回來的。 \end{tabularx} \\ \\ \relax
2:7 & \begin{tabularx}{0.7\textwidth}{X} 她說：『請你容許我拾取麥穗，在收割的人身後撿禾捆中掉落的麥穗。』她就來了，從早晨直到如今，除了在屋子裡坐一會兒，她都留在這裡。」 \end{tabularx} \\ \\ \relax
2:8 & \begin{tabularx}{0.7\textwidth}{X} 波阿斯對路得說：「女兒啊，聽我說，不要到別人田裡去拾取麥穗，也不要離開這裡，要緊跟著我的女僕們。 \end{tabularx} \\ \\ \relax
2:9 & \begin{tabularx}{0.7\textwidth}{X} 你要看好我的僕人正在哪塊田收割，就跟著女僕們去。我已經吩咐僕人不可侵犯你。你渴了，可以到水缸那裡喝僕人打來的水。」 \end{tabularx} \\ \\ \relax
2:10 & \begin{tabularx}{0.7\textwidth}{X} 路得就臉伏於地叩拜，對他說：「我既是外邦女子，怎麼會在你眼中蒙恩，使你這樣照顧我呢？」 \end{tabularx} \\ \\ \relax
2:11 & \begin{tabularx}{0.7\textwidth}{X} 波阿斯回答她說：「自從你丈夫死後，凡你向婆婆所行的，以及你離開父母和你的出生地，到素不相識的百姓中，這些事人都告訴我了。 \end{tabularx} \\ \\ \relax
2:12 & \begin{tabularx}{0.7\textwidth}{X} 願耶和華照你所行的報償你。你來投靠在耶和華－以色列神的翅膀下，願你滿得他的報償。」 \end{tabularx} \\ \\ \relax
2:13 & \begin{tabularx}{0.7\textwidth}{X} 路得說：「我主啊，願我在你眼前蒙恩。我雖然不及你的一個婢女，你還安慰我，對你的婢女說關心的話。」 \end{tabularx} \\ \\ \relax
2:14 & \begin{tabularx}{0.7\textwidth}{X} 吃飯的時候，波阿斯對路得說：「你到這裡來吃些餅，把你的一塊蘸在醋裡。」路得就在收割的人旁邊坐下。波阿斯把烘了的穗子遞給她。她吃飽了，還有剩餘的。 \end{tabularx} \\ \\ \relax
2:15 & \begin{tabularx}{0.7\textwidth}{X} 她又起來拾取麥穗，波阿斯吩咐僕人說：「她即使在禾捆中拾取麥穗，也不可羞辱她。 \end{tabularx} \\ \\ \relax
2:16 & \begin{tabularx}{0.7\textwidth}{X} 你們還要從捆裡抽一些出來，留給她拾取，不可責備她。」 \end{tabularx} \\ \\ \relax
2:17 & \begin{tabularx}{0.7\textwidth}{X} 這樣，路得在田間拾取麥穗，直到晚上。她把所拾取的麥穗打了約有一伊法的大麥。 \end{tabularx} \\ \\ \relax
2:18 & \begin{tabularx}{0.7\textwidth}{X} 路得把所拾取的帶進城去給婆婆看，又把她吃飽所剩的拿出來，給了婆婆。 \end{tabularx} \\ \\
[1ex]
\hline
\hline
\end{longtable}
$^{1}$各位訂閱姊妹,平安.
不知道大家有沒有試過幫助別人.
我猜應該有吧.
但你有沒有試過在你.
正在經歷很艱難的時侯.
你仍然有心力去關心或幫助別人.
又或者相反,你知道有人正在.
經歷很艱難的時侯.
但他仍然會過來關心你或問你的需要.
不知道大家有沒有試過.
今天和大家分享的訊息.
這段經文是講向有需要者施恩有關.
開始之前,我們有個祈禱.
因為你賜下你的話語.
讓我們真的能夠知道.
我們在地上如何去做.
才能得蒙你的喜悅.
主啊,可能我們在生活裡面.
都有很多不同的困難.
或者有很多的不足.
求你幫助我們,讓我們.
藉著你的話語.
我們能夠.
可以被你更新.
主啊,求你的聖靈幫助我們.
讓我們有一個柔軟的心.
去接受你話語的提醒.
吩咐和激勵.
我們這樣祈禱,奉主耶穌基德聖名,致祈求.
阿們.
今天和大家分享一個熟悉的故事.
就是《魯德記》.
《魯德記》背景是一個很混亂的時代.
當時正值饑荒.
正值饑荒.
有一個主人翁叫做伊麗米勒.
她是一個猶太人.
因為饑荒,她就帶了她的妻子蘭娥米.
和她的兩個兒子離開了伯利恆.
去了摩亞.

$^{41}$今天想和大家分享一下.
這個故事的開頭.
其實引申了我們這個故事.
繼續會有五個人物或者五組人物.
在裡面經歷著神的恩典.
又會向有需要的人施恩.
讓我們一起聽下去.
伊麗米勒帶著兩個兒子去了摩亞之後.
摩亞落地生根.
不過之後伊麗米勒死了.
不過她的兩個兒子各自娶了媳婦.
好景不常.
十年之後,她的兩個兒子又死了.
剩下婆媳三個人.
當時蘭娥米也聽到.
神眷顧了她的家鄉伯利恆那邊.
猶太地那邊.
有糧食.
所以她和她的兩個媳婦一起去起行.
離開了她的一組人的地方.
回到她的娘家.
蘭娥米和她的兩個媳婦說.
你們走吧,回到你們自己的地方.
大家知不知道那兩個媳婦的反應是怎樣的?.
她們哭了一口.
如果是現代的話.
不知道大家有沒有聽過婆媳糾紛.
婆媳糾紛就是.
如果關係不太好的時候.
大家也不想太常見到.
但這個不是的.
現在蘭娥米和她的兩個媳婦說.
你們走吧,回到你們自己的地方.
她們是放聲大哭.
蘭娥米看著她們說.
其實你們跟著我是沒有用的.
我又不能生孩子讓你們再結婚.
你們又沒有人,沒有物種.
跟著我做甚麼呢?.
我甚麼也給不了你.

$^{81}$其中一個媳婦.
就聽她的勸,真的離開了.
但她離開也不是說.
好吧,我走了,再見.
不是的,她仍然是放聲大哭.
所以我們可以從.
剛才說的這一段.
在《魯德記》的第一章看到.
其實她和她的兩個媳婦的關係.
真的不差,是挺好的.
也可以相信.
蘭娥米對媳婦是挺好的.
所以她們才會這麼不捨.
她不捨到甚麼程度呢?.
剛才說有一個人離開了.
另一個媳婦就叫做魯德.
魯德卻沒有離開.
他還要說甚麼呢?.
他還要說,我不會離開的.
你的百姓就是我的百姓.
你的神就是我的神.
你死在哪裡,我就死在哪裡.
這個世界沒有東西可以隔絕我和你的.
除了死之外.
否則就願神降佛於我.
你想想,那個感情是多麼深厚.
魯德是多麼愛她.
在這個難捨難離的狀態下.
我就帶魯德繼續回去.
今天就回到我們主要所說的第二章.
剛才第一章所說的過程.
其實你剛才聽到的.
起碼十年以上的一章.
但轉入第二章的時候.
其實只是一天的時間.
但他很詳細地記錄了發生的事.
我們想今天從一個恩惠的角度去看看.
剛才所說的五組人在裡面參與的角色.
是怎樣的.
我們首先看看第二章的第二至三節.

$^{121}$摩訶女子魯德對南迺說.
「不如你讓我到田裡去拾墨水」.
「我看看我可以在誰眼前蒙恩」.
「我就能夠跟在他後面」.
魯德就去了.
他去了田裡.
在修葛的人後面拾墨水.
我們可以看到.
他們回到白梨行.
就算神賜福給他們有糧食.
其實這兩位婆媳是無依無靠的.
他們不會坐在那裡自動有飯吃.
其實他們都需要付出一些東西.
如果魯德沒有跟南迺回來.
其實南迺就要出去拾東西.
現在魯德跟了回來.
魯德其實是一個寡婦.
他還是寄居的.
但他仍然沒有想太多.
雖然按照舊約的律法.
他可以被照顧.
需要被憐憫.
但他沒有想這件事.
他沒有覺得「我現在真的慘了」.
「我老公死了」.
「我又沒有兒子」.
「還要照顧奶奶」.
他沒有這樣.
他反而很全力以赴.
他用他的努力去想.
「我可以怎樣做呢?」.
他亦沒有因為.
「我現在人生路不熟」.
「奶奶,不如你跟我一起出去吧」.
他沒有.
他很努力地嘗試去做.
去照顧這位被困境.
被情緒困擾住的奶奶.
他主動地去勇於承擔.
到城外的一些田去撿墨水.

$^{161}$他還要跟當地人溝通.
看看可以跟誰去撿墨水.
魯德剛好去到哪裡呢?.
他去到以利物力.
即是他老爺的族人.
波瓦斯的那塊田.
波瓦斯經過那塊田時.
他去看那塊田.
他看到有個人不認識.
是誰來的呢?.
他看到他.
看到他在做甚麼.
波瓦斯問監督.
「這是誰來的?」.
「原來魯德從早到晚在那裡撿東西」.
「很勤力」.
「而且撿東西不要緊」.
「他還要吃飽又繼續撿」.
「最終他撿了22公升的大麥」.
「我們相信挺重的」.
「他要搬回城裡」.
「古時的路一定沒有現在那麼平坦」.
「他還要拿著這麼重的東西回去」.
「其實一點都不容易」.
「但魯德很甘心樂意送這些東西回去」.
「給那奧米」.
我們可以看到.
魯德對於那奧米.
真的有很大的恩惠在裡面.
他自願去做這件事.
又盡力去做.
我們可以看到.
魯德他做這些事的時候.
他的愛是很大的.
但我不知道大家有沒有想過.
其實他這麼愛的程度.
我們能否做到呢?.
又或者他容易做到呢?.
剛才說過.
魯德沒有權沒有勢.

$^{201}$他寄居的.
他在這裡人生路不熟.
但是他真的沒有想太多.
也沒有想自己的限制.
而是全心全意去對待這個奶奶.
今天其實和我們差不多.
可能我們不是每一個都可以很有錢.
可能我們都沒有特別大的權力.
沒有什麼權勢.
不過即使我們像魯德那樣.
只有體力.
其實我們只要有一個願意.
去為身邊有需要的人去施恩惠的話.
神都會好像祝福魯德一樣.
去藉著他去賜福給那奧米.
今天不知道我們會不會去想一想.
我們會不會都願意像魯德那樣.
去關心一下我們有需要的家人.
我們有需要的朋友呢?.
看完第二至三節之後.
我們又再看看.
神的恩典其實是沒有分高低的.
是平等的.
為什麼這樣說呢?.
我們又再去看看.
在三節過了之後.
魯德就突然間不再是主角.
誰去了主角呢?.
就是波亞斯.
他又登場了.
剛才有提及過一點.
波亞斯真的去了伯利恆的田地.
看見有一個人他不認識的.
他就問誰來的呢?.
但在之前他首先去問一問.
跟他的僕人是有互動的.
有什麼互動呢?.
我們看看經文.
波亞斯在伯利恆那裡回來.
對修國的人說.

$^{241}$「願耶和華與你同在,與你們同在」.
然後他們又回答他.
「願耶和華賜福給你」.
其實我們想一想.
如果魯德是主角的話.
其實是否需要第四節呢?.
其實是不用的.
可以省略掉.
但為什麼要有第四節呢?.
其實發生了什麼事呢?.
原來作者很想藉著這一節.
讓讀者知道.
波亞斯是一個很敬虔的人.
還有他的手下的僕人.
他們都很著重神.
所以他們會互相以神.
作為一個祝福對方的對象.
除了波亞斯和他的僕人的關係很好之外.
其實我們看得到.
他們所說的內容,大家都知道.
能夠有一個好的修成.
一個順利的修成.
其實都是要有神的賜恩蒙福才行.
即是這樣.
其實表達了.
無論波亞斯是主人.
或者那些僕人.
他們的身份高低也好.
其實都是需要神的恩惠.
我們又再看看他們的僕人.
他們的僕人和監督是做什麼事呢?.
他們的僕人和監督.
其實剛才說過.
他們都不是主人.
有什麼可以做呢?.
我們又看一看.
剛才說到主僕彼此問候之後.
波亞斯就留意到那個不認識的女人.
就問了.
其實做什麼事呢?.

$^{281}$波亞斯就問了.
這個是哪個女子來的?.
監督就告訴他.
她就是那個摩訶的女子.
就是跟著拿俄米從摩訶地回來的那個.
她就說.
請你容許我撿那些墨水.
在修葛的人身後掉了的墨水.
我就會撿.
然後問完之後.
她就來了.
由早上到現在.
除了在屋裡坐了一會之後.
其實她還是留在那裡.
即是沒有停過手.
這段經文可能大家都很熟悉.
除了反映了魯德很勤力之外.
甚至之後波亞斯會對魯德好.
大家都是很熟悉的.
但是大家有沒有想過.
其實原來在波亞斯來視察之前.
原來波亞斯是不知道魯德是來拾水的.
為什麼他可以拾水呢?.
原來是他的僕人監督允許了魯德拾水.
我們可能會想很正常.
很理所當然.
但是不是的.
如果我們留意下第八至九節的經文.
我們會看到原來不是那麼必然.
當中其實監督和僕人給了一些恩惠給魯德.
他顧及他這個有需要的人.
在波亞斯未吩咐他之前.
原來藍僕已經節制了自己.
沒有欺負魯德.
而且在這件事.
波亞斯我們待會會看到.
告訴我們原來不是那麼必然.
所以就算我們剛才說的僕人和監督.
他沒有什麼權勢.
他沒有什麼錢.

$^{321}$但是他仍然可以成為一個向人施恩的人.
就像我們這樣.
可能我們生活上.
可能我們在職場上.
我們可能位置都不是很高.
我們沒有什麼特別很大權力.
就算我們有權力也好.
又或者我們在教會.
我們教會可能有不同的崗位.
可能你會覺得有些崗位是高一些.
有些崗位是低一些.
但是其實原來經文提醒我們.
其實我們不要只在意那份權力高低.
或者我們有沒有這個能力.
其實原來只要我們有一些態度.
我們有一份微笑.
我們有一份接納.
我們有一份寬容.
其實我們都能夠令到.
在外面來的人感受到一份恩情.
令人感受到那份重視神的人.
是怎樣對人友善的.
試想一下.
如果僕人一開始就在想.
這個女人都不是猶太人.
她是外邦人來的.
走過來做什麼.
當她這樣想的時候.
其實她錯過了神去期望.
信她的人怎樣去愛倫社.
她錯失了她去經歷上帝心意的機會.
所以剛才說的波亞斯的僕人和監督.
都是很值得我們去學校.
我們又再看看剛才說到波亞斯.
波亞斯其實是一個.
他所做的一切.
我們看英文.
他超越了律法的叫他做的事.
我們看看第八至九節和十四至十五節.
波亞斯對魯德和他說.

$^{361}$女兒聽我說.
你不要去別人的田裡去拾墨水.
也不要離開這裡.
要緊緊跟著我的女僕.
你要看好我的僕人.
去了哪塊田去過.
你就跟著女僕去吧.
我已經吩咐了僕人.
不可以去侵犯你的.
你頸鶴可以去水缸.
喝我僕人打的水.
吃飯的時候.
波亞斯又對魯德說.
你來這裡吃點餅.
你喝點醋吃.
魯德就在修局的人旁邊坐下.
波亞斯又給了一些紅了的墨水紙.
遞給他吃.
他又吃飽了.
還要又靜.
他又收起了墨水.
波亞斯又吩咐僕人.
就算他在和菌那裡拿了墨水出來.
你都不可以羞辱他.
你們還要是從和菌裡面.
特意抽一些出來.
留給他.
讓他拿.
不可以責備他.
我們怎樣看到.
波亞斯是超越了.
律法的要求呢.
我們可以去看看.
舊約利未記.
十九章九至十節.
他說什麼呢.
他說你們要在自己的地收莊稼的時候.
不可以割盡田野的角落.
也不可以拾取莊稼所掉下來的.
也不可以摘盡葡萄園的葡萄.

$^{401}$也不可以摘了在葡萄園掉下來的葡萄.
要留給窮人和寄居的.
我是耶和華.
你們的神.
新穎記有二十四章十九節有提到.
你們在田裡面收莊稼的時候.
如果有一種菌在田裡忘記了.
不可以再回去拿.
你要留給寄居的.
孤兒和寡婦.
好讓耶和華.
你的神在你手裡所做的.
賜福給你.
路德又是寄居.
又是寡婦.
現在他真的在收割的人後面.
拾墨水.
所以波亞斯其實也跟隨了神所說.
跟隨了律法所說.
允許了路德繼續這樣做.
已經合乎了律法的義.
但是律法裡面.
還有沒有很詳細的說.
你之後又要怎樣對待他.
又要預備什麼.
其實沒有很詳細地寫.
但是我們剛才從第八節到第九節.
我們可以看到.
你不要去其他地方.
你要緊緊跟隨我們的女僕.
還有我已經吩咐了我的僕人.
不可以去侵犯你.
如果你拿了東西.
不可以去責備你.
你想想.
其實為什麼要這樣說呢.
就是說.
其實他可能知道.
他去其他田裡.
可能有機會被人這樣對待他.

$^{441}$所以他真的很恩待他.
告訴他.
你留在這裡.
在我的實力範圍底下.
我可以保護你的.
我就保護你.
事實上.
波亞斯如果純粹為了不犯律法的話.
其實剛才那些東西.
他根本不需要做.
但是他不是.
他不視律法是一個例行公事.
因為當如果我們真的做了.
我就做了律法所說的東西.
其實只是一個例行公事.
可能不需要有憐憫的心.
我做了.
神說我做了.
但是他不只這樣.
你看到他為了路德.
預備了很多東西.
無論吃東西.
還要讓他咬醋.
還要叫僕人給他什麼.
又要不要罵他.
預備得很細心.
你會看到.
他是超過了.
波亞斯是深深地經歷神在律法背後的精神.
那份愛淪捨的心.
那份照顧孤寡的心.
我們也可以從而看到.
波亞斯是真正地跟從神.
向有需要的人施憐憫.
而路德也因為這樣.
稱得上在波亞斯的眼中蒙恩.
為什麼路德可以在波亞斯的眼中蒙恩呢?.
波亞斯是怎麼說的呢?.
波亞斯說.
自從你丈夫死了之後.

$^{481}$凡你向婆婆所行的.
以及你離開你父母.
你的出生地.
到數不相識的百姓這裡.
這些事人都告訴我了.
現在願耶和華照你所行.
去報償你.
你來投靠在耶和華.
以色列神的翅膀下.
願你滿滿的得到他的報償.
原來波亞斯很欣賞路德對婆婆忠誠的愛.
以及他所做的一切.
他不是很簡單地看.
是不是有個媳婦跟著婆婆離鄉別井.
他不只是這樣看.
因為他覺得他放下父母.
放下家鄉.
甚至放下了他家鄉的神.
他來到這裡人生路不熟的伯利恆.
為了什麼?.
原來他是為了去投靠這位又真又活.
永生的神耶和華的翅膀下.
雖然波亞斯給了路德的恩惠.
已經遠超過律法的要求.
但是你有沒有發現.
波亞斯沒有強調.
沒有說我現在做了很多事.
其實律法沒有說.
他沒有.
他沒有很強調.
他也沒有很高調高舉自己所做.
他反而是祈求耶和華向路德施更多的恩賜賞賜.
事實上波亞斯的祈願.
其實是藉著他自己所做的.
在承受神對路德的祝福.
難怪我們可以看到.
波亞斯在吃飯的時候.
在一切事中.
給了很多恩典給路德和罵我們.
我們又看看我們自己.

$^{521}$教會今年在推動恆在社區.
我們今天又是怎樣對我們的鄰舍呢?.
我們是怎樣向他們施恩惠呢?.
哪些是我們的鄰舍呢?.
我們有沒有留意一下呢?.
我們會不會做的時候.
只是為了牧師這麼好.
我就聽他說一點.
教會現在一些普遍的指引和吩咐.
做一下.
還是我們是很積極的.
很投入的.
我們是由心去發出的那份憐憫.
去留意我們社區有需要的人.
我們會願意施恩給我們的鄰舍呢?.
我們會真切地去關心他們.
透過聆聽.
我們了解他們的需要.
我們嘗試盡上一份愛心.
如果你說我們正在做.
那就太好了.
那你們就繼續加油.
如果沒有的話.
你說我之前真的沒有.
不要緊.
因為神經物藉著它的經文.
它的道.
可能是激勵了你的話.
你說我現在被提醒了.
我已經開始了.
那也是好的.
還未遲.
我相信如果我們願意踏出一步的時候.
神仍然會一樣這麼開心.
這麼喜悅我們所做.
接著我們又說一下撈米.
撈米.
它的接受恩惠.
也就是施恩惠.
第十七節和十八節是這麼說的.

$^{561}$他說老爹在田間拾取了墨水.
直至到晚上.
他所拾取的墨水.
大約有22公升的大麥.
老爹就將所拾取的.
帶進城裡給婆婆看.
也不只給她看.
他還將今天剩下的.
帶給婆婆.
給婆婆看,給婆婆吃.
我們剛才一直數下來.
可能我們會想.
人物關係好像一層層的.
首先.
神是最大的.
給了波亞斯主人.
可以去祝福.
下面有監督.
有僕人.
寄居孤寡的老爹.
接著就是撈米.
咦,這樣看會不會有些想法.
撈米是最底層的.
他什麼都沒有做.
只是接受恩惠.
我們可能會想.
既然他很慘.
有一把年紀.
失去老公不久.
又失去兩個兒子.
又失去孫子.
你想想,什麼都沒有.
一無所有.
對於舊約的女性.
真的什麼都沒有.
她正在一個情緒低谷.
我們是否不能期望.
她成為一個施恩者.
是不是呢.
如果我們這樣想.

$^{601}$可能我們就要調教一下.
正如我剛才說的第二點.
在神的恩典下.
其實人人都是平等的.
他們去到糧食之城.
百利行裡面.
我們看到.
不會沒有領袖的成份.
其實在裡面.
每個人都有領袖的成份.
只不過在資源上.
可能各有不同.
波亞斯可能多些.
撈米可能少些.
但這個恩惠之間是流動的.
你會看到.
波亞斯對監督僕人很好.
監督僕人又會知道.
原來神是這樣吩咐.
我的上司也是這樣.
我的老闆也是這樣做.
我的主人也是這樣做.
他又這樣對盧德.
盧德當然不是因為他們這樣做.
他本身對撈米很好.
撈米有沒有對盧德好.
我們可以看看.
撈米本身的名字是有意思的.
大家知不知道.
它的名字本身的意思是什麼.
是的,我聽到有人說甜.
它本身是甜的.
但它現在正在經歷人生最苦的.
所以它回到巴黎的時候.
它就對人說.
你不要再叫我甜.
你叫我馬拉苦.
現在很苦.
但是.
如果我們再慢慢發現.

$^{641}$它經歷了第二章.
它聽了盧德.
講了它整天的經歷的時候.
它經歷了盧德對它的照顧.
它經歷了波亞斯.
給盧德和撈米的一些恩惠的時候.
它好像開始振作起來.
今天不會講到.
不過大家應該知道.
在第三章的時候.
它開始振作起來.
它要為盧德的前路.
它的歸宿去鋪路.
這種恩惠雖然還沒有發生.
在第二章的時候還沒有發生.
不過撈米對盧德.
是不是真的沒有做事.
原來不是的.
我們可以想一下.
其實它雖然正在受師傅的供養.
但其實剛才說過.
態度原來是很重要的.
如果撈米這麼苦.
我不知道大家有沒有看過電視劇.
那些很慘的人.
很苦的那些人.
他們會說什麼.
我不用你幫.
你不要可憐我.
就算是他身處很慘的家婆.
他都要挑剔家嫂.
很淒涼的.
如果你想一下.
盧德跟著撈米.
他又去幫她.
其實他可以挑剔她.
你跟著我幹嘛.
妨礙我.
或者是你很厲害嗎.
可以是這樣.

$^{681}$但他沒有.
他用一個願意領受恩典的心.
去接受盧德給他的私恩.
他是因然地接受.
所以我們可以看到.
撈米從盧德入門那刻開始.
未那麼苦的時候.
還是甜的時候.
他應該對盧德很好.
正如剛才所說.
他們哭著不捨得走.
其實可能在當中.
撈米雖然有些時候.
他會很沮喪的時候.
他會怨.
但無論如何.
在當中.
可能他的一言一行.
或者他的信心.
他的信念.
其實都讓盧德看到.
他信忍的神是實在的.
亦可能因為這樣.
不知不覺地體驗了.
令到盧德願意跟婆婆說.
你的神就是我的神.
他願意離鄉別井.
放下自己家鄉的神.
然後來到這裡.
去投靠撈米的神.
原來撈米其實都藉著.
他對神的認識.
他的見證.
去見證這個會私恩的神.
聽過今天的訊息.
我們可以知道.
無論我們任何身份,地位,狀況.
富有的,富足的.
還是貧乏的.
我們其實都可以成為一個.

$^{721}$跟從神吩咐.
成為一個可以私恩的人.
因為我們在乎的.
不是私恩.
不是我們只是給錢.
而是可能我們是一份關心.
我們一份付出.
我們一份同行.
我們一份在意.
其實都可以是一個恩惠.
而最重要的這些.
就是那個恩惠.
不是一個例行公事.
亦不是讓人看到.
我在幾艱難中.
我還是對你這麼好.
不是,不是看到我們這份慷慨.
要看的,就是看那份.
神藉著我們,帶著憐憫的心.
去對有需要人的私恩的時候.
人們可以看得見.
我們所信靠的主.
有多真實.
和祂給我們的恩典.
有多豐盛.
在生活裡,在社區中.
我們可能有跟人不少相處的機會.
我們不只是跟我們的親戚朋友.
我們可能接觸的人.
都有不同類型的關係.
不過其實只要有一個接觸點.
其實就可能是神要使用我們的時候.
使用我們的機會.
希望藉著我們去賜恩福.
給這些有需要的人.
學彼得所說.
金和銀我們可能未必有.
但是我們希望將我們所認識的.
那位永活的神.
我們將那麼好的永生的主.

$^{761}$我們可以去介紹給他們認識.
讓他們都經歷到耶穌的那份真實和愛.
我邀請大家一起去.
有一個沉澱的時間.
大家一起去合上眼睛.
今日藉著神的訊息或經文.
你覺得神在這裡跟你說了些甚麼呢?.
有沒有一些其實有需要的人.
早就在你們身邊出現了呢?.
今日主有沒有給你感動.
要你給更多的支持.
或者鼓勵那位有需要的人呢?.
在這個星期裡.
有沒有一些關愛的行動.
你可以去做呢?.
我們一起有個祈禱吧.
你讓我們見到.
在路得記裡面.
被你所祝福的人.
無論是波克斯.
無論是他的僕人.
無論是路德還是那奧米.
他們都看得見你的實在和恩典.
他們彼此之間的恩待.
他們那份不計較.
讓彼此都能見到你的實在.
亦都因為那份願意去愛.
和對你的信靠.
我們可以知道.
在路得記後面那一章.
第四章的時候.
我們見到你有更大的計劃.
要成就在波克斯和路德身上.
主也求你幫助我們.
讓我們真的可以學校.
波克斯和路德.
那份超越了一般的道德.
一般的規範.
對人的愛.
是由心去生出的時候.

$^{801}$我們可以看見.
主你所喜悅的時候.
我們亦能夠願意去學校他們.
主也可能我們信心少.
可能有我們的難處.
主也大求你.
藉著今天的經文來鼓勵我們.
讓我們願意一小步一小步地踏出來.
願意留意我們身邊有需要的人.
讓我們成為與你同工的施恩者.
讓別人看見.
我們信靠你這位又真又活的神.
你的恩典是實在的大.
求你幫助我們.
我們這樣祈禱.
是奉主耶穌基督得勝名字祈求.
阿們.
\newpage



\section{馬太福音 19:27-30}
\label{sec:csMea85HLY0}
\textbf{2026.06.14 《恩典,超越計算》:馬太福音19\_27-30 曾裔貴牧師}
\newline
\newline
連結: \href{https://youtube.com/watch?v=csMea85HLY0}{\texttt{https://youtube.com/watch?v=csMea85HLY0}} ~~~~ 語音日期: 2026-06-19
\newline
\newline
\hyperref[sec:iob9mODZZqY]{\small{< < < PREV SERMON < < <}}
~
\hyperref[sec:index_chronic]{\small{[返順時目]}}
~
\hyperref[sec:index_scriptual]{\small{[返順卷目]}}
~
\hyperref[sec:V7GlEPh1gqw]{\small{> > > NEXT SERMON > > >}}
\newline
\newline
馬太福音 19:27-30
\newline
\begin{longtable}{cl}
\hline
\hline
章節 & 經文 (和合本修訂版)\\
\hline
19:27 & \begin{tabularx}{0.7\textwidth}{X} 於是彼得回應，對他說：「看哪，我們已經撇下一切跟從你了，我們會得到甚麼呢？」 \end{tabularx} \\ \\ \relax
19:28 & \begin{tabularx}{0.7\textwidth}{X} 耶穌對他們說：「我實在告訴你們，你們這些跟從我的人，到了萬物更新、人子坐在他榮耀寶座上的時候，你們也要坐在十二個寶座上，審判以色列十二個支派。 \end{tabularx} \\ \\ \relax
19:29 & \begin{tabularx}{0.7\textwidth}{X} 凡為我的名撇下房屋，或是兄弟、姊妹、父親、母親、兒女、田地的，將得著百倍，並且承受永生。 \end{tabularx} \\ \\ \relax
19:30 & \begin{tabularx}{0.7\textwidth}{X} 然而，有許多在前的，將要在後；在後的，將要在前。」 \end{tabularx} \\ \\
[1ex]
\hline
\hline
\end{longtable}
$^{1}$弟兄姊妹 再一次講 主女便平安.
我們一起同心禱告.
天父啊 我們剛才.
謙守聖餐.
為了你的愛子.
為我們人類犧牲釘十字架.
成為了我們的救贖.
成為了我們的榮耀.
我們的冠冕.
使我們能夠在地上.
以榮耀你的聖名來生活.
多謝你讓我們能夠有聖經的教導.
這一段這麼長的經文.
這一個令我們有反思空間的比喻.
也讓我們能夠看到.
真是在整個的比喻裡面.
那個意思是相當清楚.
盼望我們能夠真是在地上的生活.
我們常常有一個警醒 謹慎的心.
因為主耶穌.
你來到去警告當時的門徒.
在前的真是會在後.
在後的又可以能夠在愛恩典裡面.
來站到前.
我們仰望主.
教導我們.
我們都是來自不同的背景.
性格.
我們更加是有不同的.
可能教會的經驗生活.
但是我們能夠在洪恩家.
一起這樣聚集.
我們可以來到去每個星期敬拜你.
這個是神冥冥中有你的帶領.
有你的感動.
也都有你的差遣.
所以我們很盼望.
我們從一個很宏觀的眼光看.
其實整個的世界.
就是神的葡萄園.

$^{41}$我們在主的角度裡面.
其實我們就能夠在不同的地方.
一同的敬拜.
幫助我們能夠透過這個.
今天我們要看的經文.
來彼此有勉勵提醒.
我們這樣仰望禱告.
奉耶穌基督的名.
阿們.
在今個月.
1號.
就收到一個邀請.
就是.
見道神學院.
對不起.
播道神學院88週年.
有一個弟兄居然邀請我.
其實我就不太認識他.
邀請我去參加他的畢業典禮.
很高興很榮幸.
很早去到5點半.
就要坐定了.
因為沒有位子的.
很早去到.
安靜等候.
接著.
整體.
很多的畢業生.
都有二十多人.
這個弟兄當然是其中一個.
正道的牧師.
非常親切.
是2004年的時間.
我申請播道神學院的時候.
就是這位院長.
楊院長接見我.
大家彼此交談.
接著.
那一年就進不了播道神學院.
第二年就進了見道神學院.

$^{81}$我變成見道的同學.
楊院長特地從蘇格蘭回來.
講這一堂道.
非常之感動.
原來他說每一個能夠事奉主.
長時間.
能夠仍然懷著這一份感動的初心.
只有一個原因.
就是認識主耶穌的大愛.
很感動在場的每一位畢業生.
畢業有很多事情要做.
又要上台.
又要拿獎.
可能都沒有什麼心情聽道.
但是我們坐在這裡.
特別我們服侍一段時間的弟兄姊妹.
牧者.
我們就憤外有一份的感動.
而且.
那個感動.
不單止是聽那一堂道.
而是能夠每一年出席畢業典禮.
見到有很多的牧者.
不同的年齡的.
來到去受栽培.
三年,四年.
能夠出工場.
記得當晚那個.
每一個畢業生.
都有一個代表.
來到去支持.
多謝學院.
多謝老師.
多謝很多的奉獻者.
支持者,同行者.
那位弟兄年紀相當大了.
可以看到.
原來讀神學.
不需要有一個法定年齡的.
有些很大年紀都可以讀神學的.

$^{121}$當然我的弟兄很年輕而已.
其實就是錦繡堂未來九月份的傳道牧者.
很感動.
讓我在一路的思想裡面.
我想起這一段我很喜歡.
亦是我自己經常提醒自己的經文.
很想今天跟大家分享.
剛剛過去的星期四.
我們每兩個月一次.
屯門,元朗,天水圍的所有宣道會.
我們會聚集一起禱告.
今次去到.
因為我星期四.
很多時候都是我們週間有福音團契.
不是太想來到去停.
停了太多.
但今次很有感動.
跟長者,弟兄姊妹說.
今次要停一次.
我要去參加祈禱會.
又帶給我有很大的感動.
因為又看到不同的宣道會.
又有新的神學生畢業.
加入宣道會的事奉的行列.
看到有很多年輕的牧者.
那個感動又是很大的.
就是說神不斷有呼召.
不斷有邀請不同的年輕的弟兄姊妹.
來奉獻自己的人生.
去讀神學.
所以在座的弟兄姊妹.
如果你實際已經很大年紀了.
也不要覺得自己沒用.
年輕的更加要思考.
神有沒有在今天的講道裡面.
仍然給大家有一個呼召呢?.
開始我們今天這個題目.
相信經文的重點.
已經呼之欲出了.
19章的30節.

$^{161}$和20章的16節.
很明顯是一個首尾的呼應.
耶穌在這裡這麼說.
祂說:這讓在後的將要在前.
在前的將要在後.
這一節聖經是一個總結句.
是將19章的第30節.
那一節其實是不需要說的.
說多就錯多.
說多之後.
就要用16節的聖經來解釋.
為什麼要說這一句.
所以這個最後的總結.
這樣.
這樣就是說.
如果你在這個比喻裡面.
你get不到神給你的訊息.
有哪些人get不到呢?.
就是九點鐘進去葡萄園做工的那群人.
其實很少數也說不定.
當然那裡沒有說那個數字有多少.
那個比例有多少.
但是.
整個葡萄園.
由早到晚準備下班.
來到去發工錢的時候.
其實每個人都很開心.
因為本來沒有想過可以進去葡萄園的.
誰知道被邀請進去了.
做了一會之後.
哇!立刻真的有錢了.
不是假的.
這樣就很開心.
帶著一個銀幣就回家了.
一路一路這樣去.
但是有為數.
不知道多少的數量.
九點鐘進去工作.
簽好了合同.
當然是口頭的.

$^{201}$來到去廣告了.
下班.
發工資.
一個銀幣.
哇!發脾氣了.
埋怨家主.
這個比喻就是這個訊息.
但是這個訊息背後.
耶穌就說得很清楚了.
就是這樣.
就是說如果你這樣發脾氣.
你沒有看到葡萄園園主的愛.
那個恩典.
這樣的人.
就會在後了.
他會慢慢地.
在工作裡面.
得不到真正的滿足喜樂.
我們一起進入經文裡面.
跟大家來看看.
在還沒說這個比喻之前.
耶穌就向門徒預告了一件事.
當然更上的上文.
就是有一個很厲害.
很有錢的財主.
來找耶穌.
主耶穌.
如何可以承受永生.
你告訴我.
你教我.
我就會去做的了.
耶穌就跟他說.
摩西的十誡.
從小已經遵守了.
還有什麼我做得不夠你告訴我.
耶穌說這樣嗎?.
那你把你的錢周濟窮人.
然後來跟隨我.
這個少年的財主.
聽完耶穌這樣說.

$^{241}$聖經告訴我們.
他悠悠愁愁地回家.
悠悠愁愁地走.
即是他不想變賣他的財富.
去周濟窮人.
這樣的一個財主.
那耶穌看著這個財主走的時候.
就說有錢財進神的國的人就難了.
記住不是說所有有錢的人.
都不可以進神的國.
是就著那個財主.
是就著他根本沒有想過要跟隨主.
所以大家不要搞錯經文.
耶穌這樣說完之後.
再加多一個這樣的例子.
他說那些駱駝穿過針眼.
比財主進神的國度還來得容易.
周圍的門徒彼得和約翰就問.
那誰可以得救啊?耶穌.
誰可以上天堂啊?.
主耶穌說了一句話.
在人是不能夠的.
但是在神凡事都能.
就是說神可以改變人心.
我們認為沒有可能的事.
沒有可能的人.
他都可以在神的恩典裡面.
得到那個生命改變.
當然那一刻那個財主還沒有改變.
還沒有肯將他自己的生命的主權交出來.
所以就導致這樣的情況.
接著本來就應該結束了.
這個只是偶遇.
但是彼得聽完耶穌這樣說之後.
突然間加了一句說話.
剛才二十七節已經說了.
耶穌說我們已經撇下一切跟隨你了.
將來在天堂得到什麼?.
可以問耶穌想換取那個獎賞.
那個兌換券.

$^{281}$將來我可以換取到什麼?.
耶穌的說話真是很快人心醒.
告訴我們耶穌預告這個葡萄園比喻之前.
就是使徒彼得的這個提問.
這個提問其實很有意思.
他真的問一下.
我今生由現在跟隨你之後.
我努力的來為你打拼.
為你來服侍.
你要我放下什麼?.
那個財主不肯放下這麼蠢的.
我們當然是跟你.
我們當然是撇下.
耶穌說.
我實在告訴你們.
即時的就是那幫門徒.
你們這些跟隨我的人.
到了萬物更生的時候.
人子,主耶穌.
坐在神的榮耀寶座上的時候.
你們十二個使徒.
就會坐在十二個寶座上.
審判以色列十二個之派.
你想想提升到有多高層次的榮耀.
還不夠.
29節.
凡為我的名撇下房屋.
或是兄弟姐妹父親母親兒女田地的.
將得著百倍.
並且承受永生.
不僅僅是十二使徒.
凡是在這樣的邏輯.
在這樣的心境裡面.
聽到明白耶穌說的天堂的永恆福樂.
又願意這樣來跟隨的人.
或多或少為了信仰.
為了永恆生命.
你放下.
所以這裡說得很盡.
房屋.

$^{321}$真的不是生物這麼簡單.
是你居住的.
兄弟姐妹.
如果不合適的話.
長大之後大家都不是這麼熟的.
父母又已經不在世了.
那也無所謂了.
大家都各有自己的家庭.
撇下就撇下.
兒女.
田地.
耶穌說將來會得到一百倍的回報獎賞.
如果你是彼得.
你聽著耶穌這麼說.
你看著那個財主.
悠悠愁愁地走.
你們是十二個的門徒.
跟隨耶穌.
就站在耶穌旁邊.
耶穌這樣公開地說.
來給大家一個這樣的暗示.
將來有這麼大的榮耀.
和承受永生.
你是彼得.
你是約翰.
或者你是猶大.
更加清清楚楚.
不過很可惜.
耶穌真的好像說多錯多.
突然間說第三十節.
那怎麼辦?.
這一節聖經是什麼意思?.
他說然而.
然而就是說.
原來剛才說的.
可能不能夠發生.
有些人可能得不到.
有許多在前的將要在後.
在後的將要在前.
我們不要用地上的排隊.

$^{361}$這樣來打尖的概念.
是說在永恆裡面.
一個本來很勤奮.
工作為神的.
慢慢因為各種的原因.
那個比喻裡面就說了原因.
說了他為什麼會被放在最後的原因.
所以這一節聖經.
其實真的說多錯多.
你說來幹什麼?.
剛才那裡已經很美麗.
我能夠清清楚楚.
現在我放下我的家人.
我的田地.
我將來就有一百倍.
還可以承受永生.
是不是不需要這一節呢?.
如果沒有了這一節.
就不需要那個比喻.
解釋這一節了.
你說是不是很好呢?.
大家清清楚楚.
計算一下.
A弟兄.
我信心也不是很夠.
我奉獻一點點.
將來也有收成的.
有些人說.
不是啊,我比你強.
我無論如何.
我全力以赴.
一定將來有百倍耶穌的應許.
但是這一節聖經就卡住了.
帶來一個很要停頓的反思.
耶穌也不等待門徒.
在那裡胡思亂想.
就立即去說圖圖員的比喻.
所以耶穌基督不想門徒誤會.
就說了這個.
這裡很清楚告訴我們.

$^{401}$弟兄姊妹.
如果我們只是停留在27,28和29節.
門徒之間就會用一種比較的心態.
來看誰撇得多.
誰奉獻得多.
誰來照著耶穌的教導去做.
而慢慢忘記了.
這不是恩典.
而走向了一種撇多少的競賽.
成為基督教各自好像爭嚴鬥厲.
好像爭相誰做得多誰做得少.
只是門徒中間已經是這樣.
很快沒多久.
當耶穌快要上十字架的時候.
門徒已經來爭大.
將來會坐哪個位置.
所以這一節聖經.
原來真的非常寶貴.
是要叫門徒.
每一個後來的跟隨者.
要認清到底我們跟隨的這位主.
這位神.
是一位怎樣的主.
怎樣的神.
所以耶穌基督立即說的這個比喻.
就真的告訴我們一件很重要的事.
這個葡萄園的比喻.
其實很簡單.
九點,十二點,三點,五點.
都有人被邀請.
被招聘進去葡萄園工作.
葡萄園無限這麼大.
外面的人.
有一些一直都行站.
所以這是一個比喻.
可能實際真的有這樣的場景.
當時都是每天發工錢.
耶穌都不是說一些天馬行空他們不懂的生活片段.
可能就是每天都會請一些臨時工進去.
但是真正的場景.

$^{441}$是不會有人五點就要放工.
還要請人.
沒有東西給他做.
只有一個小時.
所以我們千萬不要用這樣的實際.
來駁斥這個比喻.
這個比喻要帶出的.
就是到底這一位的原主所代表的神.
是一位怎樣的神.
這個很重要.
當然整個比喻最突出的.
就是已經很清楚了.
是有人不開心.
埋怨家主.
埋怨的人.
是有point的.
因為我九點鐘入園.
我看著那個五點鐘在這裡入園的.
興高采烈地坐一會就做完了.
又沒有東西給他做.
真的很不開心.
所以如果當時真是這樣的場景.
是不可能的.
這個是一個誇大的比喻.
突顯的就是這個原主.
真是有一個很大的慷慨.
那個心胸大到一個地步.
是可以海納百川.
每個人他都能夠容納.
人人都可以被罩.
不會說.
麻煩你給我看你的履歷表.
對不起.
原來你不會耕種.
對不起.
不請了.
從來沒有這樣的比喻.
每個人都請.
技術性的我們不要計算了.
因為比喻不是想講這麼技術性的.

$^{481}$不是那個工頭就教你了.
麻煩你在那裡幫我掘那些泥.
就不是這樣.
你幫我在那裡收割.
不是.
是說進去在葡萄園裡面的工人.
人人都是得到一個銀幣.
當然聖經沒有說.
九點鐘進去的要不要那個銀幣.
可能連那個銀幣都扔給主人.
自己走了.
聖經沒有說.
我們不需要去.
只知道他是很埋怨.
深深埋怨.
因為他們一路的考慮.
想法就是.
以為必要多得.
這個點就是扣住他們的心.
令他們不開心的地方.
就是以為必要多得.
因為他看著每個人都拿到一個銀幣.
他就說這次厲害了.
這個園主肯定會更加愛我們.
給多一點我們.
誰知道都是一個銀幣.
所以重點就是在這個圖裡面已經清楚.
整個比喻已經講清楚.
親愛的兄弟姐妹.
我們要思想.
這個比喻裡面要大進的.
他們領了工錢.
就埋怨那家的主人說.
我們整天勞苦受熱.
那些後來的只做了一小時.
當然是直接針對五點鐘那些.
他們未必敢說十二點鐘那些.
因為差不多相差三個小時.
又沒有提到三點的班.
一提就馬上提五點.

$^{521}$因為那個對比是太強烈.
九點和五點就是很強烈的對比.
就是說振振有詞都有point.
我是九點.
你是五點.
為什麼你們會得到那個錢幣呢?.
為什麼我們不行呢?.
你竟待他們和我們一樣呢?.
開始這樣有比較.
開始好像上文所說的.
彼得的心境.
如果我們撇得多一點.
將來就會更加多一點.
相信這一段聖經.
給我們看到的就是.
原來我們在神的國度裡面.
慢慢忘記了歡樂的原因.
慢慢忘記了原來葡萄園裡面.
我們可以看到不同的眾生相.
我們有很大的喜樂.
而且是日於言表的.
內心的高興.
有不同的階層.
有不同的能力.
有不同的性格.
都看到這樣的弟兄姊妹.
被召入葡萄園.
在葡萄園裡面.
大家是歡天喜地準備等待.
發工前的時間.
就是這個比喻最精粹的地方.
就是我們需要來到上上.
讓我們在勤奮工作的時候.
我們想一下身邊的人事物.
我們想一想不同的弟兄姊妹.
有些可能真是所謂五點鐘.
如果用那個生命的階段.
可能就是一些完全沒有那麼多.
跟我們一樣恩賜的人.
他都可以被邀請入葡萄園.

$^{561}$我們有恩賜有呼召.
你就盡你自己的責任.
在你的恩賜.
在你的呼召裡面.
去忠心忠於你的所託.
這樣去做.
但也有一些人.
根本沒有這樣的能力.
所以相信園主的慷慨.
是九點鐘入園的.
比喻裡面的那些故宮.
他一下想不通.
園主的慷慨.
其實正正就是耶穌告訴我們.
神的國度救恩.
不是用我們去換取.
用多少的撇下去換取.
而是在神的慷慨裡面.
我們看到的是神很大的恩典.
心存感恩.
這個比喻肯定不是告訴我們.
這樣我們要心存僥倖.
我們要在我們的生活.
得過且過就行了.
完全不是這樣的意思.
但是心存感恩的弟兄姊妹.
是會有很不同的眼光.
看待其他的人.
因為園主的慷慨.
正正就是葡萄園的法則.
園主的慷慨.
就是能夠看到.
有些人行站的那種無奈.
那種貧乏.
很想邀請他進入這個園裡面.
欣賞園主的慷慨和喜樂.
弟兄姊妹.
來到教會.
我們看到的是什麼呢?.
我們能不能夠慢慢建立一份心靈的感動.

$^{601}$看到這位主.
祂的大愛.
祂的恩典.
祂的厚恩.
我們能夠常常心存一份感恩的心靈呢?.
昨天一直預備的時候.
我想起一位弟兄.
學父牧師.
其實他比我高一班.
他是2008年畢業的.
也是見到的傳奇.
因為2008年那班.
就是四年的神學.
有些是二十多歲.
但是學父牧師.
進入神學院的時候.
已經差不多六十歲.
你想一下.
差不多五點鐘就入園的人.
他在見證也告訴我們.
他是全校最大年紀.
比那些老師還大.
很辛苦又要讀希臘文.
很辛苦很辛苦就這樣完了.
其實.
原來他的見證是很感動的.
他說他1975年已經移民在溫哥華.
他是一個政府裡面的工程師.
還有幾年就可以退休領長俸畢業.
退休拿長俸.
但是他說他終於感動.
很清楚神的印證.
毅然辭職.
當然不會有長俸.
毅然辭職之後.
就回來香港.
闊別了29年的香港.
回來讀神學.
上長洲見到神學院.
畢業的時候.

$^{641}$已經是接近退休.
接著回去溫哥華.
最近退休了.
他現在和太太去到Alberta.
那些渺無人煙的地方.
和一些華人做餐館的.
來傳福音.
做宣教士.
做自由的宣教士.
原來在神的國度裡面.
根本是每一個人都能夠有用.
每一個都不會覺得是遲的.
這位家主.
我們難以想像.
耶穌說這樣的比喻.
難道耶穌不知道.
違反了那些市場的公平原則嗎?.
有什麼理由.
這樣的不對的待遇.
我這麼勤力.
他又不用做.
又一樣和我一樣的工資.
耶穌基督說的神的國度.
不是地上的市場學.
不是地上公平的原則.
絕對不是耶穌不公平.
而是這個比喻.
有一個很平衡的地方.
就是神的國度.
不可以用世界的眼光去看.
有很多人.
其實根本沒有資格去信耶穌.
但是耶穌的國度裡面.
包含的有很多.
曾經犯很大罪的人.
曾經很抵擋神的人.
都可以得到主耶穌的救恩.
每一個弟兄姊妹.
耶穌基督都有祂的呼召.
有祂的邀請.

$^{681}$我們要在這個葡萄園裡面.
在我們的心靈.
在我們的眼界.
我們要常常有葡萄園的觀念.
你和我.
不過有些是祖先.
有些是稍遲.
在葡萄園裡面工作的工人.
我們每一個等待的.
在未到去派工前之前.
我們要等待.
要看的.
就是做物主那份慷慨.
我們就能夠深深感恩.
我們就能夠很賣力地服侍.
剛才那節聖經.
剛才提到的九點鐘入園勤奮工作的.
他們埋怨家主一句話.
你有沒有看到我們九點鐘入園的.
整天在做什麼?.
他們不是坐在那裡嘆世界.
喝冰水.
我們整天在這裡勞苦受熱.
是事實.
他們在說的事實.
但這個事實和耶穌.
這個世界的國度很不同.
祂的勞苦受熱.
是忠於祂自己所託的去工作.
但是家主難道.
不看到你九點鐘就勞苦受熱?.
難道這位做物主.
不知道我們每一個人.
盡心盡力地服侍.
來鞠躬盡瘁地服侍?.
難道家主只看到五點鐘.
就很開心地等收工的人.
看不到九點鐘入園.
最大的謬誤.
就是他以為主耶穌看不到他.

$^{721}$他以為神這個園主.
看不到他的勞苦受熱.
弟兄姊妹.
我們要記得.
我們在主裡面服侍.
我們要記得.
我們所做的.
在主的眼中.
祂真的會看在眼內.
這是將來主耶穌賞賜的時候.
來跟不同的牧者.
弟兄姊妹去說.
某某人.
你來盡心的勞苦受熱.
你想想如果由家主.
由這個園主.
由主耶穌.
由將來永恆的主.
祂的口.
這樣來讚賞.
答謝這樣的弟兄姊妹.
這樣的牧者.
你是勞苦受熱.
你想想那個讚賞.
那個稱讚.
有多大的力量呢?.
但是由自己的口說出來.
就變成埋怨.
變成對主安排的挑戰.
不甘心.
以為這位主.
為什麼你會這樣計算呢?.
所以弟兄姊妹.
今天的題目是.
恩典超越了計算.
記住.
天國的比喻有很多.
這是其中一個比喻.
但是這個比喻.
的確要制止.

$^{761}$以為要跟別人比較的撇下.
是唯一的天國的原則.
天國的原則.
這裡很清楚.
是以恩典.
是以接納.
是以欣賞.
是以一種的歡欣.
來看這個世界.
剛才說了.
受到的邀請.
在這裡.
八十八屆的畢業典禮.
楊院長給我們的教訓.
和弟兄的邀請.
和看到眾多弟兄姊妹.
能夠完成神學.
進入神的國度.
服事.
有些是宣教士.
有些是機構的同工.
也有些進入教會.
文輝弟兄.
我沒有把我和他拍的照片.
放給大家看.
文輝弟兄.
非常有感動.
來到我們錦宣服事.
九月他就會到任我們的牧者.
其實我不認識他.
因為我在雄恩.
也沒有回去.
只是有時候打個招呼.
點個點頭.
他居然給我一個邀請卡.
曾牧師.
你可不可以來參加我的畢業典禮.
我當然很開心.
因為沒有請柬.
去不到.

$^{801}$很感恩.
看到弟兄姊妹.
他們熱切服事的心.
進入葡萄園工作的.
這一份心靈.
大家如果有機會去看看.
不同的神學院畢業.
在座我們也有一些.
在神學教育方面有幫手.
你每一年見到畢業生.
他們進入工場.
那種雀躍.
那種眼的發光.
準備去幫助神的國度.
去建立不同的弟兄姊妹.
那種雀躍的心靈.
我們會看到.
是很興奮的.
很開心.
弟兄姊妹.
葡萄園最值得看的.
就是那些.
以為很僥倖.
但其實是園主的慷慨.
弟兄姊妹我們要學習.
看神的恩典.
在不同弟兄姊妹身上的作為.
以致我們真的帶著一份感恩.
能夠同主同感一靈.
在座有些是老師.
在師班.
在福曲.
來訓練弟兄姊妹.
你看到他們有進步.
其實你的心是很興奮.
很雀躍.
有神學院的老師.
展秋.
你看到神學生畢業.
其實是很興奮.

$^{841}$因為神國度的葡萄園.
是收割不盡的.
需要不同的人.
來做主的工作.
盼望我們大家一起去學習.
園主慷慨的心.
就是他看到有閒暫的人.
以致我們能夠感受一下.
園主的慷慨.
去領取獎賞的喜悅.
你看到別人做這麼少.
就拿到一個銀幣.
而你真的能夠完成地上的工作.
又是得到那個銀幣.
那個銀幣是多麼珍貴.
因為主在主的眼裡.
看到的你.
是勞苦受熱的你.
不是一個整天埋怨比較的你.
很難學習的.
如果不是耶穌.
都不會這樣告誡他的門徒.
最後我們要來做一個總結.
所以如果我們一直都沒有用.
神角度的眼光去看教會.
看神的角度的工場.
我們就很容易會落入一種.
埋怨比較容易產生驕傲.
恣意的心.
但是真正的福氣.
就是我們很開心.
很天真地在神的角度裡.
用神的愛去愛身邊的人.
我們一起禱告.
親愛的主.
我們要將榮耀歸給你.
因為你就是這樣來帶領世界的教會.
走進不同的歷史.
也告訴我們你是那位歷史的主.
前面的路是怎樣.

$^{881}$神啊,你是知道的.
我們現在的宗教會.
是等候主的再臨.
在萬物復興的時候.
當你坐在榮耀的寶座.
我們就會得到真正永恆的稱讚和獎賞.
多謝你讓我們知道.
在這個世界就是屬於你的.
我們在地上為你而活.
為你作見證.
以致我們在地上.
我們真的可能有不同的人.
在不同的階段.
信主和服侍.
求主發出你的邀請.
感動弟兄姊妹.
能夠在你的安排.
在你的聖靈的恩賜.
以致在不同的地方來服侍.
我們也有回流的宣教士.
在我們中間一起聚集.
親愛的主.
但願你讓我們在這個比喻裡面.
深深地來認識.
一個不是說得少的比喻.
因為耶穌是很認真地說.
在前的會在後.
所以我們盼望主給我們.
上上懷著感恩的心.
上上都是忠於我們自己的託付.
我們去殷勤.
不去有一個僥倖的心態.
而去等待將來獎賞的日子.
我們獻上感恩.
奉耶穌基督的名.
阿們.
\newpage



\section{路加福音 8:40-56}
\label{sec:V7GlEPh1gqw}
\textbf{2026.06.21 《父.愛》:路加福音8\_40-56(和修版) 曾敬宗牧師}
\newline
\newline
連結: \href{https://youtube.com/watch?v=V7GlEPh1gqw}{\texttt{https://youtube.com/watch?v=V7GlEPh1gqw}} ~~~~ 語音日期: 2026-06-25
\newline
\newline
\hyperref[sec:csMea85HLY0]{\small{< < < PREV SERMON < < <}}
~
\hyperref[sec:index_chronic]{\small{[返順時目]}}
~
\hyperref[sec:index_scriptual]{\small{[返順卷目]}}
~
\hyperref[sec:code]{\small{> > > NEXT SERMON > > >}}
\newline
\newline
路加福音 8:40-56
\newline
\begin{longtable}{cl}
\hline
\hline
章節 & 經文 (和合本修訂版)\\
\hline
8:40 & \begin{tabularx}{0.7\textwidth}{X} 耶穌回來的時候，眾人迎接他，因為他們都等候著他。 \end{tabularx} \\ \\ \relax
8:41 & \begin{tabularx}{0.7\textwidth}{X} 有一個會堂主管，名叫葉魯，來俯伏在耶穌腳前，求耶穌到他家裡去， \end{tabularx} \\ \\ \relax
8:42 & \begin{tabularx}{0.7\textwidth}{X} 因為他有一個獨生女，約十二歲，快要死了。耶穌去的時候，眾人簇擁著他。 \end{tabularx} \\ \\ \relax
8:43 & \begin{tabularx}{0.7\textwidth}{X} 有一個女人，患了經血不止的病有十二年，在醫生手裡花盡了一生所有的，但沒有人能治好她。 \end{tabularx} \\ \\ \relax
8:44 & \begin{tabularx}{0.7\textwidth}{X} 她來到耶穌背後，摸他的衣裳繸子，經血立刻止住了。 \end{tabularx} \\ \\ \relax
8:45 & \begin{tabularx}{0.7\textwidth}{X} 耶穌說：「摸我的是誰？」眾人都不承認。彼得說：「老師，眾人擁擁擠擠緊靠著你。」 \end{tabularx} \\ \\ \relax
8:46 & \begin{tabularx}{0.7\textwidth}{X} 耶穌說：「有人摸了我，因為我覺得有能力從我身上出去。」 \end{tabularx} \\ \\ \relax
8:47 & \begin{tabularx}{0.7\textwidth}{X} 那女人知道瞞不住了，就戰戰兢兢地俯伏在耶穌跟前，把摸他的緣故和怎樣立刻痊癒的事，當著眾人都說出來。 \end{tabularx} \\ \\ \relax
8:48 & \begin{tabularx}{0.7\textwidth}{X} 耶穌對她說：「女兒，你的信救了你。平安地回去吧！」 \end{tabularx} \\ \\ \relax
8:49 & \begin{tabularx}{0.7\textwidth}{X} 耶穌還在說話的時候，有人從會堂主管的家裡來，說：「你的女兒死了，不要勞駕老師了。」 \end{tabularx} \\ \\ \relax
8:50 & \begin{tabularx}{0.7\textwidth}{X} 耶穌聽見就對他說：「不要怕，只要信！她必得痊癒。」 \end{tabularx} \\ \\ \relax
8:51 & \begin{tabularx}{0.7\textwidth}{X} 耶穌到了他的家，除了彼得、約翰、雅各，和女兒的父母，不許別人同他進去。 \end{tabularx} \\ \\ \relax
8:52 & \begin{tabularx}{0.7\textwidth}{X} 眾人都在為這女孩哀哭捶胸。耶穌說：「不要哭，她不是死了，是睡著了。」 \end{tabularx} \\ \\ \relax
8:53 & \begin{tabularx}{0.7\textwidth}{X} 他們知道她已經死了，就嘲笑耶穌。 \end{tabularx} \\ \\ \relax
8:54 & \begin{tabularx}{0.7\textwidth}{X} 耶穌拉著她的手，呼叫著：「孩子，起來吧！」 \end{tabularx} \\ \\ \relax
8:55 & \begin{tabularx}{0.7\textwidth}{X} 她的靈魂就回來了，她立刻起來。耶穌吩咐給她東西吃。 \end{tabularx} \\ \\ \relax
8:56 & \begin{tabularx}{0.7\textwidth}{X} 她的父母非常驚奇；耶穌吩咐他們不要把所發生的事告訴任何人。 \end{tabularx} \\ \\
[1ex]
\hline
\hline
\end{longtable}
$^{1}$各位朋友各位弟兄姊妹平安.
平安.
不如開始的時候我都想先問一問大家.
究竟在座當中我聽說今天是福音主日.
在座當中有多少弟兄姊妹.
你邀請了你的家人或者朋友或者是同事.
什麼都好.
你邀請了朋友來參加今天的崇拜.
你可不可以舉手讓我知道.
有邀請到 非常之好.
如果是這樣我們更加要問的就是.
不知道在座當中有哪一個.
你今天是被邀請來到紅欣堂參加崇拜的呢.
歡迎你 歡迎你.
很開心可以見到有朋友來到我們當中.
我不知道有沒有一些朋友其實你都是被邀請的.
不過你不好意思舉手.
但是無論如何我們都是很歡迎你的.
我都想問一下大家.
你覺得今天是不是一個特別的日子.
是 有沒有人說不是的.
如果你說是的話.
可能因為今天是父親節.
但是如果你說不是的.
可能你心裡面在想我每天都是父親節.
我都是這樣對我的爸爸媽媽.
我不知道你是不是這樣在想.
對我來說今天都是一個特別的日子.
不單單因為我是一個爸爸.
我有兩個兒子.
一個升中三一個升中二.
更加重要的是我自己都是一個兒子.
我想每一個人都有自己的父母.
特別是今天來到洪水橋.
洪水橋都是我的一些回憶.
因為四十多年前.
應該未夠五十年前.
當時我爸爸都帶我來洪水橋住了一晚.
當時還記得我們是去到一間三角瓦頂的村屋住.
我不知道在座當中有多少位是在洪水橋長大的.

$^{41}$我還記得當時來到的時候.
那間屋是很特別的.
可以有一個閣樓爬上去.
上去之後你很近看到三角形的屋頂.
一條條一條條的橫木橫樑.
感覺很特別.
外面很難聞到大自然的味道.
很香.
有很多果樹在外面.
然後進入屋內.
晚上的時候又發覺很安靜.
但安靜得來其實又不安靜.
因為有很多叫聲.
當然今天來到洪水橋.
甚麼都不同了.
真的完全不同了.
我都沒有想過四十多年後的今天.
這個變化是這麼大.
我也想像不到未來的日子.
洪水橋這裡會變成怎樣.
元朗這裡會變成怎樣.
我都發覺人生真的有很多很多的變化.
這些變化或者甚至乎是.
無論這些變化是好是壞都好.
我都發覺我們自己是掌握不到.
我們很想去掌握.
甚至乎我想預計一下.
我都很難去預計.
說起小時候我就想起.
我告訴大家.
我頭上後面這裡有一條兩吋長的疤痕.
這條疤痕是怎樣來的呢.
我記得大概兩歲左右的時候.
我住在油麻地.
在舊樓那裡.
跌了一個球進去洗手間.
然後我就走進去洗手間撿那個球.
當我撿了那個球之後.
因為洗手間很滑.
那時候洗手間地上是鋪了甚麼.

$^{81}$瓷磚.
瓷磚真是好.
我只鋪了紙皮石.
很滑的紙皮石.
然後一滑向後.
就跌了下去.
同樣都是鋪了紙皮石的門杉那裡.
叫門欖.
那時候不知道發生甚麼事.
只知道.
我媽媽就很緊張.
抱著我.
不管三七二十一就衝下樓.
我們住七樓.
跑下樓梯.
然後走出去街口就馬上截的士.
跟的士司機說要去醫院.
我只記得.
那時候我媽媽還穿著一件很長的神勾.
那時候我媽媽告訴我.
因為太緊急.
所以去到醫院.
醫生就馬上幫我連針.
他告訴我連麻醉藥都沒有放.
然後.
因為她穿著那件.
很長的神勾.
棉質的神勾.
當她去到醫院的時候.
其實整件神勾已經染了血.
但是.
對我來說.
說真的.
今天.
我只記得幾幕的情景.
但是我連當日痛我都不記得.
沒有放麻醉藥連針.
感覺是怎樣我都不記得.
所以很感恩.
今天可以在這裡跟大家繼續分享.

$^{121}$是的.
在這件事發生的時候.
我想我爸爸媽媽都沒有想過.
在他們這麼小的兒子身上.
會發生一件這樣的事.
這麼突然的事情.
當這件事發生的時候.
我相信我的爸爸媽媽.
他們都完全沒有能力.
不知道可以怎樣處理.
總之做什麼就做什麼.
想起做什麼就要去做什麼.
我們今天一起看路加福音的經文.
路加福音都是神藉著一位醫生.
路加所寫的經文.
正正記載了耶穌生命的一件事蹟.
在這件事蹟裡面同樣都有一個爸爸.
他為他所愛的獨生女.
非常非常地擔心.
在今天我們讀的經文的開頭的時候.
很簡單地交代了.
有很多人迎接耶穌.
不過他們不是真的認識耶穌.
他們純粹可能是聽聞.
耶穌之前在不同的地方.
他似乎教導神的話語非常好的一個老師.
又或者他們想起.
耶穌聽聞他是一個很厲害.
可以治病趕鬼的一個醫生.
他們不是真的認識耶穌.
我想我們香港人可能都會很熟悉.
就是聽聞聽到一些人說.
他追星的時候是什麼一個狀況.
很多人很哄動.
當我們看到經文第四十一節的時候.
路加記載他在發生什麼事.
他在記載的時候突然間加了兩個字.
可能在我們看的中文譯本裡面未必看到.
但是如果按著聖經原本的翻譯.
直接去翻譯的話.

$^{161}$他是有留心看清楚這個字.
他說留心看清楚.
有一個男人來到求耶穌.
來到他家.
在經文這裡.
好像鏡頭突然間一轉.
路加突然間要我們留意這個男人.
他不只是為了一睹耶穌的風采而來.
一方面聖經告訴我們這個男人是怎樣的.
他是一個有身份有名望的人.
他是一個會堂主管.
會堂是當時猶太人聚集.
講解和學習神話語的一個地方.
他有個名字叫做「葉老」.
「葉老」這個字有一個意思.
他說他是一個被神光照的人.
這個就是「葉老」這個名字的意思.
聖經好像告訴我們這個男人.
他應該是一個相信神.
很認識神話語的人.
他還是一個領袖.
是被人尊重的一個領袖.
另一方面.
這個男人和體熱鬧的群眾很不同.
因為他來到耶穌的面前.
他是有目的的.
他去到耶穌的腳前.
他就趴在地上去求耶穌.
我們知道他的誠意真是十足.
他很懇切去求耶穌的一個表達.
然後在第四十二節就告訴我們.
這個男人原來是一個爸爸.
他有一個十二歲的獨生女.
不過這個獨生女.
聖經說她正步向死亡.
她快要死了.
我們相信這個爸爸去到耶穌的面前.
是很想求耶穌去到他的家.
來救他的獨生女.
救他這個所愛的女兒.

$^{201}$但是正當耶穌要回應這個爸爸的請求.
去到他家的時候.
在經文的第四十三到四十八節.
一個最危急的關頭.
耶穌應該要趕去救人的時候.
他竟然停下來.
去問誰摸我.
發生什麼事了.
這個時間救人要緊.
這麼多人逼著你.
壞壞地這麼說.
你管得了誰摸你呢.
聖經就是這麼說.
在第四十九節.
有一個人在耶穌還在說話的時候.
就在這個會堂主管的家.
同樣都說.
聖經說他來到他們那裡.
然後他說什麼.
他說你的女兒剛剛就死了.
你不需要再麻煩老師了.
我想意思大概就是.
人都已經死了.
這個老師無論醫生也好.
老師也好.
他多有名也好.
你叫他去你家也沒有用.
我們試想一下這個會堂主管.
他也是一個相信神.
認識神的人.
如果家裡突然發生一些不幸的事情.
你覺得他在心裡會跟神說些什麼呢.
你覺得他會說些什麼呢.
他可能也會問.
信神的人為什麼會遇到這麼不幸的事情呢.
甚至我也會想.
他會不會問.
神啊 你是不是不愛我呢.
但是今天各位朋友各位弟兄姊妹.
我也要毫不避忌地告訴大家.

$^{241}$信神的人在人生裡面.
其實和所有人都是一樣.
我們同樣都是要經歷世界的變遷.
我們都是要經歷人生的生老病死.
我們都是要面對人生裡面很多.
突如其來 意想不到的事情.
其實沒有分別的.
雖然我是一個傳道人.
但是我同樣和每一個人都是一樣.
同樣要面對不能夠預測不能夠掌握的人生.
我太太也是一個傳道人.
我太太曾經因為身體上有一個腫瘤.
而要做一個婦科的切除手術.
她很怕看醫生.
所以每一次去看醫生.
我都會陪她去看醫生.
每一次我記得.
她每次去看醫生的醫生都是很熟悉的.
他每次都是說同一句話.
我聽了很多次.
他說什麼呢.
他跟我們說.
你知道自己身體的狀況.
你知道自己什麼年紀.
我結婚的時候已經是四字頭.
然後他再加多一句.
再清楚一點 怕我們不清楚.
你都以為自己未必有小朋友.
就算有 你們都未必托得穩.
他每次都是這樣重複這兩句話.
但是很感恩.
太太之後懷孕.
雖然曾經作過小產.
但在臥床一段時間之後又穩定下來.
但我最記得的時候.
是我太太去懷孕五個月的時候.
我記得有一晚.
那一晚為什麼我那麼記得呢.
因為我生日.
和太太在外面吃完飯回家.

$^{281}$然後我太太就說她的胸口開始痛.
越來越痛 很痛.
接著我們唯有馬上趕去急症室.
去到急症室的時候.
急症室醫生就和她做檢查.
要照X-ray.
照完之後.
醫生就告訴我.
我們要再等一下.
因為我看不到.
不知道究竟你太太發生什麼事.
所以我叫了我們急症室的主管回來.
再照多一次.
又等.
再等了一段時間之後.
醫生就再告訴我.
主管還是照完之後.
不知道發生了什麼事.
所以我們已經叫了婦產科的醫生回來.
再照第三次.
然後他就告訴我.
不過你不可以留在這裡等.
你要先回家.
我記得我那時候是一個人在晚上.
黑漆漆的一個晚上.
兩三點的時候.
我一個人離開醫院.
走回家.
那時候我應該有什麼心情.
我記得.
我一路一路走出來的時候.
周圍一個人都沒有.
簡單來說.
我無言.
我沒有話可以說.
心裡面有一個很大的無力感.
因為無論是肚子裡的嬰兒.
還是我太太的生命.
我什麼都做不到.
我一路走出來的時候.

$^{321}$我心裡面好像默默地跟天父說.
或者問天父.
天父,我可以怎麼樣.
天父,你可不可以保守我一家呢.
當我們回看今天的經文.
這個會堂主管的爸爸.
他不但要面對自己內心的掙扎.
他還要面對身邊一群幫不了忙的群眾.
因為當耶穌對這位爸爸說話.
來鼓勵他.
來發出一個醫治的應許的時候.
當時那些群眾就在發生什麼事呢.
他們哀哭隨哄.
他們甚至嘲笑耶穌.
因為耶穌說.
這個女孩只是睡了覺.
這群群眾認定了.
這個女孩子已經死了.
他們其實在宣告.
其實這個爸爸再沒有任何盼望.
沒有任何的希望了.
群眾不單止阻礙了耶穌去救這個女兒.
他們更加是這個爸爸.
信心上的一個阻力和障礙.
耶穌在第五十節的經文.
耶穌說:他聽見就對這個爸爸說.
他說:不要怕.
只要信.
他必得拯救.
傳語這個字本身的意思是拯救的意思.
原來耶穌不單單是看重那個快要死的女兒的生命.
或者是那個死了的女兒的生命.
他更加好體恤到那個爸爸那一刻的心情.
他所受到的打擊.
所以耶穌即時很有力地鼓勵他.
不要怕.
耶穌特別給了這個爸爸一個的應許.
他說:只要你相信神.
你的女兒必得拯救.
不是在疾病裡面拯救.

$^{361}$是在死亡裡面拯救出來.
耶穌再有力地跟這個爸爸說.
不要哭.
你所愛的女兒不是死了.
只是睡著了.
我們知道的.
如果死了.
按照常理來說.
她不會再起來.
但是如果耶穌真是一個掌管生命的主.
她死了之後.
其實只不過人死了之後.
只是好像睡了覺一樣.
有一天她都會復活.
起來.
所以耶穌要幫這個爸爸.
隔絕去被群眾的騷擾.
不讓其他人跟她.
跟這個父母一起進入屋裡.
耶穌很親切地拉著.
這個爸爸的女兒的手.
來叫她.
小女孩.
起來.
很奇妙地這樣.
這個女兒的靈魂.
聖經說.
她就回去.
她就馬上起來.
在經文裡面.
我不知道大家是否留意到.
還有一個彩蛋.
彩蛋的意思就是.
記載了一些.
讓我們可以認識到耶穌.
認識到神的一些很特別的地方.
就是.
耶穌不單止體貼這個爸爸的心靈.
祂也很體貼女兒身體的需要.
所以祂特別吩咐人.

$^{401}$要給這個剛剛復活的女兒吃東西.
耶穌.
祂竟然考慮到這個女兒.
可能從重病到死亡.
到復活.
祂什麼都沒有吃過.
祂的身體可能很虛弱.
所以耶穌特別提醒.
大家要給女兒吃東西.
經文最後說.
這個女兒的父母見證事情的發生.
就覺得非常驚奇.
在耶穌醫治女兒的記載裡面.
我們可以看到.
耶穌是一個掌管生命.
可以叫人復活的一個真神.
但是耶穌祂不是一個.
在天上高高在上.
遙不可及的主.
祂是一個去憐恤人.
信心軟弱的一個神.
祂是一個體恤人身體需要.
很愛人.
愛到底的一個天父爸爸.
耶穌不單止沒有責怪葉老.
為什麼你這麼軟弱.
你的信心這麼容易被動搖.
祂也沒有任由他的信心.
被那些群眾影響.
雖然葉老.
仍然要經歷他所愛的女兒的死亡.
但是耶穌卻可以叫他.
叫葉老和他的太太非常驚奇.
各位朋友各位弟兄姊妹.
我不知道在你的人生裡面.
有沒有經歷過一些事情.
可能你覺得天父不知道.
你以為耶穌不知道.
以為祂不理你.
又或者以為祂不體恤你.

$^{441}$說真的.
當我們面對人生的起跌.
人生的風浪.
我們在人生的艱難裡面.
我們有這種感覺.
其實說真的.
真是人之常情.
我們有這樣的問題.
其實都是人之常情.
但是經文裡面.
它正正提醒我們.
耶穌基督其實很明白.
我們每一個.
耶穌明白你.
祂明白你的軟弱.
你的艱難.
就算事情的發生.
不是按著我們所祈求.
但是耶穌基督仍然可以叫我們.
非常驚奇.
我們不能夠去估計得到.
耶穌基督去為我們之後的安排.
如此的奇妙.
經文來到這裡.
其實還未完滿結束.
因為我們中間跳過了一個小段.
一個小插曲.
但是這個小插曲.
更加是這段經文的重點所在.
經文裡面.
小插曲講到一個女人.
去到耶穌的後面去摸祂.
其實這個記載.
是很刻意的一個編排.
來跟前面說的.
你要留心有一個男人.
來做一個很大的.
極端的對比.
我們看到這裡說.
有一個女人來到耶穌的背後.

$^{481}$摸祂的衣裳睡紙.
其實正正對應著的就是.
第四十一節.
有一個男人來到耶穌腳前.
不同的是.
這是一個女人.
當時是一個男尊女卑的時代.
所以女人是一個.
被人看低的角色.
這個女人不是一個男人.
她患了經血不止的病.
這個病在當時宗教上.
被視為不潔症.
是大家很厭棄要跟她保持距離.
她不是一個會堂主管.
一個在宗教上得高望重的人.
她患病十二年.
她向來都不會被人尊重.
患了這個病.
相信她也很難懷有小朋友.
她不是集萬千寵愛於一身.
十二歲的獨生女.
她有的只是患了十二年的絕症.
這個女人是沒有名字的.
她是一個被捨棄的低身體.
沒有人當她存在.
她不是叫葉佬.
一個別人覺得她是被神光照的人.
她只是敢去到耶穌的背後.
她不敢去到耶穌的面前.
她只能夠很羞恥隱藏地.
在耶穌的背後.
她連摸耶穌的衣服都不敢.
她摸耶穌的哪裡?.
衣裳的睡紙.
當時的猶太人.
他們的衣服綁著一條睡紙.
你想想要摸那裡要怎樣?.
其實她同樣都是要俯伏下來.
趴在地上才能摸到.

$^{521}$但是她不敢說.
好像葉佬一樣.
去到耶穌的面前.
俯伏在祂的面前來求祂.
她不敢.
這段經文.
路加夫音第八章的經文.
如果我們完全不記載這個女人的事情.
就這樣說耶穌遇到一個管會堂的男人.
然後去醫治他的女兒.
叫他的女兒復活.
很通順.
其實不記載中間這段也可以.
就算要記載這個女人的事跡.
簡單說她碰上去就康復了.
這樣就可以了.
耶穌為何偏偏要在這個時候.
故意停下來?.
聖經想我們認識到.
我們所信的神是怎樣的一個神.
耶穌是刻意圍著.
這個被視為最卑微的女人而停下來.
所以耶穌特別去問.
誰摸我呢?.
就算耶穌最信得過的三個門徒之一彼得.
他的回應也是.
三番四次這樣說.
你這麼勇勇擠擠.
這麼多人迫著你.
所以耶穌說你這位老師.
有什麼這麼出奇呢?.
連彼得都不能夠明白.
耶穌的心意.
可能彼得就好像.
當時在會堂主管的家裡.
走出來通風報信的那個人一樣.
他以為耶穌只不過是一個普通人.
是一個老師.
可能是一個醫生.
他們未真真正正認知到.

$^{561}$耶穌是神的兒子.
他未真的可以完全相信.
他不單止可以.
醫治教導神的話語.
他更加可以拯救.
去掌管生命.
他可以叫人復活.
但是耶穌沒有放棄.
這個被人排擠.
被人拒絕的no body.
耶穌再次表明.
確實是有人摸他.
如果我們直譯耶穌的說話.
耶穌是這樣說的.
他知道有人摸他.
因為耶穌說.
我知道權柄能力.
是由我發出來的.
按照經文直譯.
耶穌是說.
我知道因為權柄能力.
是由我發出去的.
所以我知道.
一定是有人特別地摸我.
最後這個女人避無可避.
她唯有戰戰兢兢地.
俯伏在耶穌的腳前.
沒有想過.
這個女人竟然.
由耶穌的背後.
她可以去俯伏在耶穌的腳前.
在耶穌和門徒.
和所有群眾的面前.
將她摸耶穌的原因.
過程.
她如何馬上得到醫治.
去說出來.
變成一個生命公開的見證.
我們要留意的是.
耶穌如何跟這個女人說話.

$^{601}$耶穌主動稱呼這個女人為女兒.
女兒.
當這個女人懷著.
她很小很小的信心.
即管一試.
戰戰兢兢地.
去到耶穌那裡.
她竟然由一個.
no body.
沒有人理會的一個人.
成為天父所愛的女兒.
她竟然可以得到一個.
極之愛她.
願意醫治她.
拯救她的一個天父爸爸.
耶穌問.
誰摸我.
絕對不是要為難這個女人.
耶穌不單止希望這個.
患了十二年絕症的女人的病.
得到醫治.
耶穌更加想這個.
一直被人拒絕孤立了十二年的女人.
祂希望她的身心靈.
都可以同樣被釋放.
讓她有一個女兒的名分.
讓她成為一個聖潔的人.
耶穌因為這個女人.
很微小的信心.
而拯救她.
讓她以後可以得著永遠的平安.
我記得.
有一首詩歌說.
神未曾應許天色常藍.
神未曾應許.
在人生裡經常聞到花香.
但是耶穌應許過.
真正應許過的是什麼呢?.
耶穌真正應許過的是.
我留下平安給你們.

$^{641}$我把我的平安賜給你們.
我所賜給你們的.
不像世人所賜的.
你們心裡不要憂愁.
也不要膽怯.
經文告訴我們.
耶穌所賜的平安.
是祂自己的平安.
祂自己的平安是怎樣的平安呢?.
祂所賜給我們的平安.
是一個祂要為我們的罪.
被出賣.
被凌辱.
被釘在十字架上.
甚至流血.
致死都不怕的一個平安.
耶穌自己的平安.
是一個勝過死亡.
復活的平安.
所以耶穌提醒跟從祂的人.
既然他們有耶穌自己這麼特別的平安.
我們的心就不需要再不斷被煩擾.
我們不需要憂愁.
我們不需要再害怕.
無論你是一個.
人家覺得.
或者你覺得人家看你為一個最低賤的人也好.
你覺得人家很尊重你也好.
又或者你覺得自己是一個被人嫌棄的人.
又或者你覺得自己是一個摘萬千寵愛於一身的人也好.
無論你是什麼人.
我們都要面對人生.
面對生老病死.
患難是真的確實存在在這個世界.
問題不斷出現.
但是不要忘記.
神所賜的平安.
卻是可以叫你意想不到地平靜你的心.
可以釋去你一切的恐懼.
你一切的無力感.

$^{681}$平安是神賜給我們每一個跟從祂的人.
一個永遠無可替代的一份禮物.
我記得有一次.
我自己都經歷了.
進了一個營會開會.
開了兩天一夜.
很累很累.
那個營會是在上水.
我住在港島區.
很累.
我乘車那天回到家.
在黃昏的時候回到家.
回到家的時候.
其實一直已經有一種天旋地轉的感覺.
但是我進到家門口.
我發覺我支持不了.
我開始很想吐.
覺得自己呼吸不了.
我一進屋我就坐在家門口.
即是換鞋的地方.
那張椅子那裡.
然後我開始覺得胸口壓著.
慢慢地我好像感覺不到自己再有心跳.
然後我手腳開始麻痺.
然後越來越那個人很繃緊.
然後我的手就越來越繃緊.
繃緊了一個拳頭.
當我支持不了的時候.
我就跟我的家人說.
我就勉強地叫了他們一聲.
叫他們召救護車.
然後他們就真的召救護車.
當救護車來到的時候.
他叫我放鬆一點.
你鬆開你的拳頭.
我很辛苦地忍著不吐.
我告訴他.
我鬆不了.
然後救護員是強行用他的手.
拉開我的手指.

$^{721}$拉開我兩隻手的拳頭.
來幫我鬆開.
然後他要移動我.
送我去醫院的時候.
他一移動我就開始吐.
當我去到醫院的時候.
我終於感受到一份舒暢.
那份舒暢是.
你知道你怎樣吐都吐不出來的舒暢.
儘管有反應.
但是沒有東西吐出來.
就是這樣.
今天大家還看到我在這裡.
太太在我出院之後.
她就問我.
其實當你呼吸不了的時候.
救護員還沒有到的時候.
你在想什麼.
直接地說.
其實她想問我.
你有沒有想過你會死.
我在等救護員來的時候.
我記得我小兒子很擔心.
他在哭.
他第一次看到我有這樣的反應.
說真的.
我跟太太說.
我呼吸不了的時候.
我不能想太多.
我只能盡力呼吸.
我當時沒有特別想過會死.
我沒有慢慢去想死是什麼一回事.
不過其實我自己知道.
這個字是存在在我的心裡面.
不過.
當時不知為什麼我就想起詩篇23篇的經文.
詩篇23篇提醒我.
耶和華是一個會悉心照顧我.
保護我的一個牧者.
接著.

$^{761}$我心裡面想起的就是.
神在我的生命裡面.
很多很多的恩典.
其實有好幾次意外.
我兩個兒子.
分別都可以不再存在.
我大兒子曾經在小時候.
差點被巴士撞到.
小兒子試過在學校跌穿頭.
我沒有那麼害怕.
因為他去校服.
染血的程度沒有我當年那麼嚴重.
但是.
我想起神很多的恩典.
我開始心裡面去祈禱.
所以我告訴我太太.
很奇怪.
很特別.
我當時在感恩.
我在做感恩的祈禱.
一路感恩.
一切就好像變得很平靜.
時間好像停頓了下來.
我就可以繼續去等救護員來.
我再想想.
什麼是神的平安.
神的平安包括了.
神過去讓我曾經經歷過.
祂的保守.
祂的大能.
原來神的平安也包括了.
祂應許過將來每一個人.
在生命盡頭的時候.
讓我們有永生的盼望.
主耶穌的平安.
很奇妙地平靜了我內心的恐懼.
讓我可以繼續安然地.
去面對著身體的狀況.
各位朋友.
是的.

$^{801}$好像我剛才所說.
每一個人都要面對人生.
人生不如意事.
十常八九.
但是只要信靠耶穌.
我們不單單是得著.
主耶穌基督的愛.
我們還有從祂而來.
連死都不怕的平安.
我們還有那個.
復活得永生的盼望.
最後我分享多一個見證.
我認識一個七十多歲.
患了末期胰臟癌的姊妹.
她是我牧養的團契.
一個團友的媽媽.
我認識她的時候.
我發覺這個長者.
是一個很有智慧.
很優雅又堅毅.
做事盡善最美的一個女士.
她人生經歷過很多的事.
她雖然七十多歲.
但是她以前在國內.
是可以讀完中學畢業.
不過她接著去了上山下鄉.
她又經歷過.
放下當時只是一歲大的女兒.
她隻身去到自己完全陌生的地方.
其實來到香港.
一個人來打工.
要放下她的女兒.
其實是很艱難的事情.
之後雖然她們一家可以在香港團聚.
但是她又為了照顧女兒.
她當時又要做另一個決定.
就是放下了.
她自己一直在香港.
做的一個電子貿易的生意.
她和她的女兒分享過.

$^{841}$在過去的日子.
其實她的內心.
自己一個人的時候.
她一直都很害怕.
不過她靠著什麼.
她就是靠著自己.
做每件事都一定有要求.
做事一定要一絲不苟.
她要一直走下去.
所以她一個單身.
即是一個在香港的女士.
她卻可以做起自己的生意.
但到她患了胰臟癌的時候.
她仍然是一個很堅毅的人.
她忍痛撐了很長很長的時間.
胰臟癌是一個很辛苦的癌症.
她為什麼這樣撐著呢?.
她說她很想陪伴照顧她的家人.
她有三個很可愛的孫子.
她很想照顧她的女兒和三個孫子.
多一點的日子.
有一次.
我太太即是師母去探望她的時候.
她很認真地跟師母說.
因為這個女士.
她一直都不肯相信耶穌.
她跟我太太說.
我不知道哪個信仰才是真實可信的.
她說我需要時間再去認識.
我要思考多一點才可以做到決定.
過了一段日子之後.
我們再去探望她.
她提起在兩天前.
她晚上痛得很厲害.
痛到睡不著.
然後她就坐了起來.
她坐了起來.
她的腦袋就不斷地思考.
哪個信仰才是真實可信的呢?.
她很清楚自己的人生在步向一個什麼地方.

$^{881}$她想來想去.
她說我都很難抉擇.
她想到一個地步.
她自己哭了出來.
說到這裡的時候.
她的眼睛就開始流淚.
師母邀請她一起祈禱.
在祈禱裡面.
我們很大膽地向神祈禱.
神啊如果你是真神.
求你就讓我知道.
你是真實可信的神.
我們都不知道神會怎麼做.
不過我們就和她一起來祈禱.
最後.
這位女士在她家人的邀請之下.
她那年在聖誕節的平安夜參加了報道會.
在這位女士在她人生的病患和生死之間.
她當晚正正聽到很對應到她.
她怎樣形容呢?.
是一個人生的感悟和認知的智慧.
她一直很想聽這件事.
她那晚真的在報道會裡面.
她說她經過.
或者我們這樣去看她.
就是她經過一個很認真.
很謙卑的一個思考.
她那晚終於選擇相信耶穌基督.
之後她還接受洗禮.
她信主之後.
我們都有關心她.
姐妹你有什麼心願呢?.
她說她很希望有神蹟.
她希望神可以給她多一點生命連日.
希望她可以繼續陪伴她的家人.
接著我記得師母當時很大膽地問她.
如果這個神蹟不出現.
你還信不信主?.
但是我們沒有想過.
她毫不猶豫斬釘截鐵地說.

$^{921}$我當然相信.
我怎會因為沒有這個神蹟而不相信呢?.
她真的是一個很有智慧.
經過很多的思考.
深思熟慮去跟從耶穌的一個人.
在姐妹離樓的時候.
當師母去探望她的時候.
她雖然很軟弱的身體.
但是她仍然用盡她的氣力.
叫師母靠近她的身邊.
耳朵伸到她的口邊聽她說話.
她說什麼呢?.
我們完全沒有想過.
她跟師母說.
我的妹妹還沒有信耶穌.
我很想去教會可以信耶穌.
我們試想想.
如果這個姐妹不是很肯定.
耶穌是值得信靠的.
她不會在臨終的時候.
第一時間盡最後一口氣去提及.
去掛心的是她妹妹還沒有信耶穌.
我記得在姐妹離世之前.
當我們在醫院為她一家去祈禱的時候.
她的家人,女兒,女婿,幾個孫一起圍著她.
握著她的手跟她祈禱的時候.
在祈禱之後.
我們看到的是這位姐妹.
她的臉色.
我們說是很嫣紅,很漂亮.
然後她好像帶著一份很安穩的笑容.
她就離開了.
這位姐妹.
她找到她人生生老病死的感悟和認知的智慧.
她知道自己要去到那位愛她的主那裡.
她知道神也會親自帶領她去到天家.
有一天她會復活.
她會跟她的家人再見.
所以今天我也很想.
我也很想在座當中.

$^{961}$每一位都好像這位姐妹一樣.
在幻變無常的人生裡面.
得到那份永遠永遠的平安.
讓我們一起跟神祈禱.
這個時候.
我也很想先邀請在座當中.
無論你是被邀請來的朋友.
又或者是我們的弟兄姊妹.
我都同樣想邀請你.
你想不想在生命中得著神的更新,醫治?.
你想不想在生命裡面.
可以得著天父,爸爸那份愛到底的愛?.
你願不願意成為一個被主所愛的女兒或者兒子?.
如果你願意讓耶穌做你的救主.
你願意.
你很想得著天父的醫治.
得著祂的拯救.
在變幻的人生裡面.
得著從神而來的平安.
如果你願意相信耶穌基督.
讓祂將平安賜給你的話.
我想請你輕輕舉一舉你的手.
我很想請你輕輕舉一舉你的手.
就好像經文裡面那個女人.
去伸手輕輕一摸耶穌衣裳的垂直一樣.
請你舉一舉你的手.
如果你願意.
願意相信耶穌基督的話.
耶穌說:不要怕,只要信,必得拯救.
耶穌說:你的信能夠救你,讓你得著平安.
弟兄姊妹,如果你曾經拒絕過主.
但是今天天父仍然以他的愛再一次觸動你的話.
在這一刻仍然是一個機會.
如果你願意讓耶穌做你的救主.
我也請你將你的手輕輕舉起.
我們看到.
是的,如果你曾經相信主.
但是你也因為種種的原因離開了主.
離開了教會.
但是你今天也很想重新得著神所賜的平安.

$^{1001}$再回到主的裡面.
我也同樣地想邀請你.
輕輕舉一舉你的手.
是的,這是你向耶穌的表達.
耶穌知道.
就算你只是在你的心裡面.
向耶穌說:耶穌,我也想舉手.
可能這一刻你還不敢舉起你的手.
但是耶穌知道,知道你的回應.
耶穌在等你,祂一直在呼喚你.
所以我鼓勵你把今天的機會.
在你的心裡面去摸一摸耶穌.
去跟耶穌說:我願意.
讓我們一起祈禱.
我們天上的父上帝.
你明白我們的艱難.
你知道我們為什麼.
在痛苦之中會向你質問.
我們為什麼會懷疑你,拒絕你.
主你知道,主啊,多謝你.
主啊,今天我們願意在你的面前.
接受你做我們的救主.
主啊,是的.
我們知道我們過去一直拒絕你.
不承認你是我們生命的主.
我們的神的時候.
主啊,我們知道我們犯罪.
求你赦免我們的罪.
主啊,求你今天成為我們生命的主.
主啊,求你幫助我們.
將你的平安,將你的喜樂.
滿足來賜給我們.
讓我們真的能夠像今天.
我們一班弟兄姊妹獻唱的詩歌一樣.
我們可以經歷到在你裡面的平安喜樂.
我們可以在我們的艱難在世的日子.
我們可以感受到.
是有天父你和我們一起.
我們無所懼怕.
主啊,我們多謝你.

$^{1041}$求你兼顧我們的信心.
讓我們得著從你而來.
豐豐富富的恩典.
我們多謝你.
我們的祈禱交托.
奉主名求.
阿們.
\newpage

\allsectionsfont{\centering}

\setlength\parindent{0pt}
\setlength{\columnsep}{1.25em}
\setlength{\parfillskip}{0pt}
\setlength{\tabcolsep}{1em}
\raggedbottom

\pagenumbering{gobble}


\newfontfamily\leftfont[Path=../fonts/fell_french_canon/, Ligatures=TeX, ItalicFont=IMFeFCit29C.otf, BoldFont=AveriaLibre-Bold.ttf]{IMFeFCrm29C.otf}
\newfontfamily\leftcitationfont[Path=../fonts/frankruehl/]{FrankRuehlCLM-Medium.ttf}
\newfontfamily\centerfont[Path=../fonts/garamond/, Ligatures=TeX, ItalicFont=EBGaramond-SemiBoldItalic.ttf]{EBGaramond-SemiBold.ttf}
\newfontfamily\rightfont[Path=../fonts/averia/, Ligatures=TeX, ItalicFont=AveriaLibre-RegularItalic.ttf, BoldFont=AveriaLibre-Bold.ttf, BoldItalicFont=AveriaLibre-BoldItalic.ttf]{AveriaLibre-Light.ttf}
\newfontfamily\rightcitationfont[Path=../fonts/rashi/]{Mekorot-Rashi.ttf}
\definecolor{hcolor}{HTML}{D3230C}
\definecolor{rcolor}{HTML}{D36F0C}
\newcommand{\chfont}[1]{\centerfont{\huge\textcolor{hcolor}{#1}}}
\newcommand{\leftcitation}[1]{\leftcitationfont{\Large\textcolor{hcolor}{#1}}}
\newcommand{\rightcitation}[1]{\rightcitationfont{\normalsize\textcolor{rcolor}{#1}}}
\newfontfamily\flowerfont[Path=../fonts/fell_flowers/]{IMFeFlow2.otf}

\begin{sloppypar}

\chapter*{\chfont{編按結語}}

\columnratio{0.5,0.5}\begin{paracol}{2}

\fontsize{11}{13}\leftfont \Large \leftcitation{א} \leftfont 余少好文．宏志博覽群書而不忘．善存藏經籍文獻備後時之用。\leftcitation{ב} \leftfont 歸主年時．受友所薦．聞道網海．\switchcolumn\fontsize{11}{13}\rightfont \Large \leftcitation{ח} \rightfont 有見粵道之危．國之封講道千言亦將就至．急之何則為？\leftcitation{ט} \rightfont 嘗聞猶太者之傳承．在其力守口述之

\end{paracol}


\columnratio{0.32,0.32,0.32}\begin{paracol}{3}

\fontsize{11}{13}\leftfont \Large 尤以吳約翰遜者 \switchcolumn[2]\fontsize{11}{13}\rightfont \Large 統．以煉千載不

\end{paracol}

\columnratio{0.32,0.32,0.32}
\begin{paracol}{3}\fontsize{11}{13}\leftfont \Large 為重．其載上之粵語講道緩緩入耳．收之藏其音頻．善妥整存．反復而嚼．受益無窮。\leftcitation{ג} \leftfont 我城我國既限．歷一四一九之不測．肺疫延年．信徒靈長屢受圍創．神州燈臺數盡指日可待．粵道之求與日俱增。\leftcitation{ד} \leftfont 觀乎社、經、法、媒、言、信、網之地．愈趨受鋤．自翔不果．授受壓力．粵道聖言亦愈漸艱難。\leftcitation{ה} \leftfont 況崇基例乎．學苑講道屢逆權勢者．其言末強受壓．舊章盡刪以存其身。講道釋數失傳．徒嘆奈何。

\switchcolumn

\fontsize{11}{13}\centerfont 
\begin{tikzpicture}
    \node (0,0) [xshift=-0.10cm, yshift=-1.0cm, opacity=0.10]{\includegraphics[width=0.30\textwidth]{../ot_frontcover.png}} ;
    \node (0,0) [xshift=+0.20cm, yshift=+2.0cm, opacity=0.10]{\includegraphics[width=0.20\textwidth]{../christ_on_cross.png}} ;
\end{tikzpicture}
\Large 

\leftcitation{ס} \centerfont 詩百又廿七載：
\leftcitation{ע} \centerfont 非耶和華建屋宇.則匠人之經營徒.
\leftcitation{פ} \centerfont 非耶和華衛城邑.則守者之儆醒徒.
\leftcitation{צ} \centerfont 余獻是卷予華人社區．願為福音流通之器．願獻斯微材為祭榮耀上帝．
\leftcitation{ק} \centerfont 阿門

\switchcolumn

\fontsize{11}{13}\rightfont \Large 滅．時越次聖殿期及當今。\leftcitation{י} \rightfont 猶太者力廣納之．筆錄以卷軸．便以傳、閱、頌、攜、守、鎖、抄、譯、釋、編，得書塔木德、密示拿等經傳．家喻戶曉．傳流若芳。\leftcitation{כ} \rightfont 猶太者文以載道．傳其口述．今我輩粵道之傳應當作如是．遂力行粵音識辨之法．載言載道．以盡忠傳粵道以待興。\leftcitation{ל} \rightfont 蒙下賜恩惠．無畏海量字音文書．既馭上帝之道．今廣及粵語講道．重駛編程之技．匯導粵音遂字稿．重塑講道現場．以傚猶太卷軸之舉便以傳流。\leftcitation{מ} \rightfont 是卷乃粵音口述傳之屬．莫通華文白話之語．

\end{paracol}

\columnratio{0.5,0.5}
\begin{paracol}{2}\fontsize{11}{13}\leftfont \Large \leftcitation{ו} \leftfont 斯殺一違儆百逆．既禁壓之．我輩聞風無奈．在所難免。\leftcitation{ז} \leftfont 另有異人例乎．以版權之名．脅網絡頻道之舉．同授礙予粵道之存流。

\switchcolumn

\fontsize{11}{13}\rightfont \Large 惟待後繼來者之傚．以譯釋傳之於神州華文地。\leftcitation{נ} \rightfont 今能排程驅馭圖靈以編彙文檔，其碼長共數千千亦無逢大礙．全蒙上帝保守。

\end{paracol}



\columnratio{1}\begin{paracol}{1}

\fontsize{11}{13}\rightfont \Large
~~~~~~~~~~~~~~~~~~~~~~~~~~~~~~~~~~~~~~~~~~~~~~~~~~~~~~~~~~~~~~~~~~~~~~~~~~~~~~~\leftcitation{ר} \rightfont 二零二三年二月一日

~~~~~~~~~~~~~~~~~~~~~~~~~~~~~~~~~~~~~~~~~~~~~~~~~~~~~~~~~~~~~~~~~~~~~~~~~~~~~~~\leftcitation{ש} \rightfont 米迦勒

~~~~~~~~~~~~~~~~~~~~~~~~~~~~~~~~~~~~~~~~~~~~~~~~~~~~~~~~~~~~~~~~~~~~~~~~~~~~~~~\leftcitation{ת} \rightfont 書於香港

\end{paracol}

\end{sloppypar}
\end{document}
